% 自動生成 (prover/build.js): 第III部 証明木
\chapter{集合論と順序体 の証明木}
本章の定理数 87, ノード数の合計 11654.

\section{第3章・第4章 (ZFC): 空集合, 数学的帰納法}
% 自動生成 (prover/build.js): EmptyEx  ノード数 32
\begin{fullproof}{定理(空集合の存在)}{$\exists y\,\forall u\,(u \notin y)$}{32}\phantomsection\label{tree:EmptyEx}
\noindent$\Pi_{1}$:
\begin{ptree}
\AXC{$x = x$\lbl{A5}}
\UIC{$\forall x\,(x = x)$}
\AXC{$\forall x\,(x = x) \to u = u$\lbl{A2}}
\BIC{$u = u$}
\AXC{$\forall u\,(u \in y \leftrightarrow u \in a \land u \neq u) \to (u \in y \leftrightarrow u \in a \land u \neq u)$\lbl{A2}}
\AXC{$\varphi_{9}$\lbl{A1}}
\BIC{$u = u \to \forall u\,(u \in y \leftrightarrow u \in a \land u \neq u) \to u \notin y$}
\BIC{$\forall u\,(u \in y \leftrightarrow u \in a \land u \neq u) \to u \notin y$}
\UIC{$\forall u\,(\forall u\,(u \in y \leftrightarrow u \in a \land u \neq u) \to u \notin y)$}
\end{ptree}
\noindent$\Pi_{2}$:
\begin{ptree}
\AXC{$\varphi_{4}$\lbl{A4}}
\AXC{$\varphi_{7}$\lbl{A3}}
\AXC{$\stackrel{\Pi_{1}}{\vdots}$}\noLine
\UIC{$\forall u\,(\forall u\,(u \in y \leftrightarrow u \in a \land u \neq u) \to u \notin y)$}
\AXC{$\varphi_{10}$\lbl{A1}}
\BIC{$\varphi_{7} \to \varphi_{4} \to \forall u\,(u \in y \leftrightarrow u \in a \land u \neq u) \to \forall u\,(u \notin y)$}
\BIC{$\varphi_{4} \to \forall u\,(u \in y \leftrightarrow u \in a \land u \neq u) \to \forall u\,(u \notin y)$}
\BIC{$\forall u\,(u \in y \leftrightarrow u \in a \land u \neq u) \to \forall u\,(u \notin y)$}
\end{ptree}
\noindent$\Pi_{3}$:
\begin{ptree}
\AXC{$\stackrel{\Pi_{2}}{\vdots}$}\noLine
\UIC{$\forall u\,(u \in y \leftrightarrow u \in a \land u \neq u) \to \forall u\,(u \notin y)$}
\AXC{$(\forall u\,(u \in y \leftrightarrow u \in a \land u \neq u) \to \forall u\,(u \notin y)) \to \varphi_{6}$\lbl{A1}}
\BIC{$\varphi_{6}$}
\end{ptree}
\noindent$\Pi_{4}$:
\begin{ptree}
\AXC{$\forall a\,\exists y\,\forall u\,(u \in y \leftrightarrow u \in a \land u \neq u)$\lbl{Sep}}
\AXC{$\varphi_{5}$\lbl{A2}}
\BIC{$\exists y\,\forall u\,(u \in y \leftrightarrow u \in a \land u \neq u)$}
\AXC{$\exists y\,\forall u\,(u \in y \leftrightarrow u \in a \land u \neq u) \to \varphi_{1} \to (\varphi_{1} \to \varphi_{2}) \to \varphi_{3}$\lbl{A1}}
\BIC{$\varphi_{1} \to (\varphi_{1} \to \varphi_{2}) \to \varphi_{3}$}
\end{ptree}
\noindent$\Pi_{5}$:
\begin{ptree}
\AXC{$\varphi_{1} \to \varphi_{2}$\lbl{A3}}
\AXC{$\forall y\,(\neg \forall u\,(u \notin y)) \to \neg \forall u\,(u \notin y)$\lbl{A2}}
\AXC{$\stackrel{\Pi_{3}}{\vdots}$}\noLine
\UIC{$\varphi_{6}$}
\BIC{$\forall u\,(u \in y \leftrightarrow u \in a \land u \neq u) \to \exists y\,\forall u\,(u \notin y)$}
\AXC{$\varphi_{8}$\lbl{A1}}
\BIC{$\neg \exists y\,\forall u\,(u \notin y) \to \neg \forall u\,(u \in y \leftrightarrow u \in a \land u \neq u)$}
\UIC{$\varphi_{1}$}
\AXC{$\stackrel{\Pi_{4}}{\vdots}$}\noLine
\UIC{$\varphi_{1} \to (\varphi_{1} \to \varphi_{2}) \to \varphi_{3}$}
\BIC{$(\varphi_{1} \to \varphi_{2}) \to \varphi_{3}$}
\BIC{$\varphi_{3}$}
\end{ptree}
\noindent 主木:
\begin{ptree}
\AXC{$\neg \exists y\,\forall u\,(u \notin y) \to \forall y\,(\neg \exists y\,\forall u\,(u \notin y))$\lbl{A4}}
\AXC{$\stackrel{\Pi_{5}}{\vdots}$}\noLine
\UIC{$\varphi_{3}$}
\BIC{$\exists y\,\forall u\,(u \notin y)$}
\end{ptree}
\begin{abbrevlist}
\item $\varphi_{1} :\equiv \forall y\,(\neg \exists y\,\forall u\,(u \notin y) \to \neg \forall u\,(u \in y \leftrightarrow u \in a \land u \neq u))$
\item $\varphi_{2} :\equiv \forall y\,(\neg \exists y\,\forall u\,(u \notin y)) \to \forall y\,(\neg \forall u\,(u \in y \leftrightarrow u \in a \land u \neq u))$
\item $\varphi_{3} :\equiv (\neg \exists y\,\forall u\,(u \notin y) \to \forall y\,(\neg \exists y\,\forall u\,(u \notin y))) \to \exists y\,\forall u\,(u \notin y)$
\item $\varphi_{4} :\equiv \forall u\,(u \in y \leftrightarrow u \in a \land u \neq u) \to \forall u\,\forall u\,(u \in y \leftrightarrow u \in a \land u \neq u)$
\item $\varphi_{5} :\equiv \forall a\,\exists y\,\forall u\,(u \in y \leftrightarrow u \in a \land u \neq u) \to \exists y\,\forall u\,(u \in y \leftrightarrow u \in a \land u \neq u)$
\item $\varphi_{6} :\equiv (\forall y\,(\neg \forall u\,(u \notin y)) \to \neg \forall u\,(u \notin y)) \to \forall u\,(u \in y \leftrightarrow u \in a \land u \neq u) \to \exists y\,\forall u\,(u \notin y)$
\item $\varphi_{7} :\equiv \forall u\,(\forall u\,(u \in y \leftrightarrow u \in a \land u \neq u) \to u \notin y) \to \forall u\,\forall u\,(u \in y \leftrightarrow u \in a \land u \neq u) \to \forall u\,(u \notin y)$
\item $\varphi_{8} :\equiv (\forall u\,(u \in y \leftrightarrow u \in a \land u \neq u) \to \exists y\,\forall u\,(u \notin y)) \to \neg \exists y\,\forall u\,(u \notin y) \to \neg \forall u\,(u \in y \leftrightarrow u \in a \land u \neq u)$
\item $\varphi_{9} :\equiv (\forall u\,(u \in y \leftrightarrow u \in a \land u \neq u) \to (u \in y \leftrightarrow u \in a \land u \neq u)) \to u = u \to \forall u\,(u \in y \leftrightarrow u \in a \land u \neq u) \to u \notin y$
\item $\varphi_{10} :\equiv \forall u\,(\forall u\,(u \in y \leftrightarrow u \in a \land u \neq u) \to u \notin y) \to \varphi_{7} \to \varphi_{4} \to \forall u\,(u \in y \leftrightarrow u \in a \land u \neq u) \to \forall u\,(u \notin y)$
\end{abbrevlist}
\end{fullproof}

% 自動生成 (prover/build.js): EmptyUnique  ノード数 52
\begin{fullproof}{定理(空集合の一意性)}{$\forall y\,\forall y_{2}\,(\forall u\,(u \notin y) \to \forall u\,(u \notin y_{2}) \to y = y_{2})$}{52}\phantomsection\label{tree:EmptyUnique}
\noindent$\Pi_{1}$:
\begin{ptree}
\AXC{$\varphi_{16}$\lbl{A4}}
\AXC{$\forall z\,(\varphi_{11}) \to \varphi_{18}$\lbl{A3}}
\AXC{$\varphi_{11}$\lbl{A2}}
\UIC{$\forall z\,(\varphi_{11})$}
\AXC{$\forall z\,(\varphi_{11}) \to (\forall z\,(\varphi_{11}) \to \varphi_{18}) \to \varphi_{16} \to \varphi_{14}$\lbl{A1}}
\BIC{$(\forall z\,(\varphi_{11}) \to \varphi_{18}) \to \varphi_{16} \to \varphi_{14}$}
\BIC{$\varphi_{16} \to \varphi_{14}$}
\BIC{$\varphi_{14}$}
\end{ptree}
\noindent$\Pi_{2}$:
\begin{ptree}
\AXC{$\varphi_{17}$\lbl{A4}}
\AXC{$\forall u\,(\varphi_{12}) \to \varphi_{19}$\lbl{A3}}
\AXC{$\varphi_{12}$\lbl{A2}}
\UIC{$\forall u\,(\varphi_{12})$}
\AXC{$\forall u\,(\varphi_{12}) \to (\forall u\,(\varphi_{12}) \to \varphi_{19}) \to \varphi_{17} \to \varphi_{15}$\lbl{A1}}
\BIC{$(\forall u\,(\varphi_{12}) \to \varphi_{19}) \to \varphi_{17} \to \varphi_{15}$}
\BIC{$\varphi_{17} \to \varphi_{15}$}
\BIC{$\varphi_{15}$}
\end{ptree}
\noindent$\Pi_{3}$:
\begin{ptree}
\AXC{$\forall u\,(u \notin y) \to \forall u\,(u \notin y_{2}) \to \forall u\,(u \notin y_{2})$\lbl{A1}}
\AXC{$(\forall u\,(u \notin y) \to \forall u\,(u \notin y_{2}) \to \forall u\,(u \notin y_{2})) \to \varphi_{8}$\lbl{A1}}
\BIC{$\varphi_{8}$}
\end{ptree}
\noindent$\Pi_{4}$:
\begin{ptree}
\AXC{$\forall u\,(u \notin y_{2}) \to u \notin y_{2}$\lbl{A2}}
\AXC{$\stackrel{\Pi_{3}}{\vdots}$}\noLine
\UIC{$\varphi_{8}$}
\BIC{$\forall u\,(u \notin y) \to \forall u\,(u \notin y_{2}) \to u \notin y_{2}$}
\AXC{$\forall u\,(u \notin y) \to u \notin y$\lbl{A2}}
\AXC{$(\forall u\,(u \notin y) \to u \notin y) \to \varphi_{13}$\lbl{A1}}
\BIC{$\varphi_{13}$}
\BIC{$\forall u\,(u \notin y) \to \forall u\,(u \notin y_{2}) \to (u \in y \leftrightarrow u \in y_{2})$}
\AXC{$\varphi_{20}$\lbl{A1}}
\BIC{$\forall u\,(u \notin y) \land \forall u\,(u \notin y_{2}) \to (u \in y \leftrightarrow u \in y_{2})$}
\UIC{$\varphi_{3}$}
\end{ptree}
\noindent$\Pi_{5}$:
\begin{ptree}
\AXC{$\forall x\,\forall y\,(\forall z\,(z \in x \leftrightarrow z \in y) \to x = y)$\lbl{Ext}}
\AXC{$\varphi_{9}$\lbl{A2}}
\BIC{$\forall y\,(\forall z\,(z \in v_{1} \leftrightarrow z \in y) \to v_{1} = y)$}
\AXC{$\varphi_{7}$\lbl{A2}}
\BIC{$\forall z\,(z \in v_{1} \leftrightarrow z \in v_{2}) \to v_{1} = v_{2}$}
\UIC{$\forall v_{2}\,(\forall z\,(z \in v_{1} \leftrightarrow z \in v_{2}) \to v_{1} = v_{2})$}
\UIC{$\forall v_{1}\,\forall v_{2}\,(\forall z\,(z \in v_{1} \leftrightarrow z \in v_{2}) \to v_{1} = v_{2})$}
\AXC{$\varphi_{10}$\lbl{A2}}
\BIC{$\forall v_{2}\,(\forall z\,(z \in y \leftrightarrow z \in v_{2}) \to y = v_{2})$}
\AXC{$\varphi_{6}$\lbl{A2}}
\BIC{$\forall z\,(z \in y \leftrightarrow z \in y_{2}) \to y = y_{2}$}
\end{ptree}
\noindent$\Pi_{6}$:
\begin{ptree}
\AXC{$\stackrel{\Pi_{5}}{\vdots}$}\noLine
\UIC{$\forall z\,(z \in y \leftrightarrow z \in y_{2}) \to y = y_{2}$}
\AXC{$(\forall z\,(z \in y \leftrightarrow z \in y_{2}) \to y = y_{2}) \to \varphi_{22}$\lbl{A1}}
\BIC{$\varphi_{22}$}
\end{ptree}
\noindent$\Pi_{7}$:
\begin{ptree}
\AXC{$\stackrel{\Pi_{2}}{\vdots}$}\noLine
\UIC{$\varphi_{15}$}
\AXC{$\varphi_{5}$\lbl{A4}}
\AXC{$\varphi_{3} \to \varphi_{4}$\lbl{A3}}
\AXC{$\stackrel{\Pi_{4}}{\vdots}$}\noLine
\UIC{$\varphi_{3}$}
\AXC{$\varphi_{3} \to (\varphi_{3} \to \varphi_{4}) \to \varphi_{21}$\lbl{A1}}
\BIC{$(\varphi_{3} \to \varphi_{4}) \to \varphi_{21}$}
\BIC{$\varphi_{21}$}
\BIC{$\forall u\,(u \notin y) \to \forall u\,(u \notin y_{2}) \to \forall u\,(u \in y \leftrightarrow u \in y_{2})$}
\AXC{$\stackrel{\Pi_{6}}{\vdots}$}\noLine
\UIC{$\varphi_{22}$}
\BIC{$\varphi_{15} \to \varphi_{14} \to \forall u\,(u \notin y) \to \forall u\,(u \notin y_{2}) \to y = y_{2}$}
\BIC{$\varphi_{14} \to \forall u\,(u \notin y) \to \forall u\,(u \notin y_{2}) \to y = y_{2}$}
\end{ptree}
\noindent 主木:
\begin{ptree}
\AXC{$\stackrel{\Pi_{1}}{\vdots}$}\noLine
\UIC{$\varphi_{14}$}
\AXC{$\stackrel{\Pi_{7}}{\vdots}$}\noLine
\UIC{$\varphi_{14} \to \forall u\,(u \notin y) \to \forall u\,(u \notin y_{2}) \to y = y_{2}$}
\BIC{$\forall u\,(u \notin y) \to \forall u\,(u \notin y_{2}) \to y = y_{2}$}
\UIC{$\forall y_{2}\,(\forall u\,(u \notin y) \to \forall u\,(u \notin y_{2}) \to y = y_{2})$}
\UIC{$\forall y\,\forall y_{2}\,(\forall u\,(u \notin y) \to \forall u\,(u \notin y_{2}) \to y = y_{2})$}
\end{ptree}
\begin{abbrevlist}
\item $\varphi_{1} :\equiv \forall z\,\forall u\,(\neg ((u \in y \to u \in y_{2}) \to \neg (u \in y_{2} \to u \in y)))$
\item $\varphi_{2} :\equiv \forall u\,\forall z\,(\neg ((z \in y \to z \in y_{2}) \to \neg (z \in y_{2} \to z \in y)))$
\item $\varphi_{3} :\equiv \forall u\,(\forall u\,(u \notin y) \land \forall u\,(u \notin y_{2}) \to (u \in y \leftrightarrow u \in y_{2}))$
\item $\varphi_{4} :\equiv \forall u\,(\forall u\,(u \notin y) \land \forall u\,(u \notin y_{2})) \to \forall u\,(u \in y \leftrightarrow u \in y_{2})$
\item $\varphi_{5} :\equiv \forall u\,(u \notin y) \land \forall u\,(u \notin y_{2}) \to \forall u\,(\forall u\,(u \notin y) \land \forall u\,(u \notin y_{2}))$
\item $\varphi_{6} :\equiv \forall v_{2}\,(\forall z\,(z \in y \leftrightarrow z \in v_{2}) \to y = v_{2}) \to \forall z\,(z \in y \leftrightarrow z \in y_{2}) \to y = y_{2}$
\item $\varphi_{7} :\equiv \forall y\,(\forall z\,(z \in v_{1} \leftrightarrow z \in y) \to v_{1} = y) \to \forall z\,(z \in v_{1} \leftrightarrow z \in v_{2}) \to v_{1} = v_{2}$
\item $\varphi_{8} :\equiv (\forall u\,(u \notin y_{2}) \to u \notin y_{2}) \to \forall u\,(u \notin y) \to \forall u\,(u \notin y_{2}) \to u \notin y_{2}$
\item $\varphi_{9} :\equiv \forall x\,\forall y\,(\forall z\,(z \in x \leftrightarrow z \in y) \to x = y) \to \forall y\,(\forall z\,(z \in v_{1} \leftrightarrow z \in y) \to v_{1} = y)$
\item $\varphi_{10} :\equiv \forall v_{1}\,\forall v_{2}\,(\forall z\,(z \in v_{1} \leftrightarrow z \in v_{2}) \to v_{1} = v_{2}) \to \forall v_{2}\,(\forall z\,(z \in y \leftrightarrow z \in v_{2}) \to y = v_{2})$
\item $\varphi_{11} :\equiv \forall u\,(\neg ((u \in y \to u \in y_{2}) \to \neg (u \in y_{2} \to u \in y))) \to \neg ((z \in y \to z \in y_{2}) \to \neg (z \in y_{2} \to z \in y))$
\item $\varphi_{12} :\equiv \forall z\,(\neg ((z \in y \to z \in y_{2}) \to \neg (z \in y_{2} \to z \in y))) \to \neg ((u \in y \to u \in y_{2}) \to \neg (u \in y_{2} \to u \in y))$
\item $\varphi_{13} :\equiv (\forall u\,(u \notin y) \to \forall u\,(u \notin y_{2}) \to u \notin y_{2}) \to \forall u\,(u \notin y) \to \forall u\,(u \notin y_{2}) \to (u \in y \leftrightarrow u \in y_{2})$
\item $\varphi_{14} :\equiv \forall u\,(\neg ((u \in y \to u \in y_{2}) \to \neg (u \in y_{2} \to u \in y))) \to \forall z\,(\neg ((z \in y \to z \in y_{2}) \to \neg (z \in y_{2} \to z \in y)))$
\item $\varphi_{15} :\equiv \forall z\,(\neg ((z \in y \to z \in y_{2}) \to \neg (z \in y_{2} \to z \in y))) \to \forall u\,(\neg ((u \in y \to u \in y_{2}) \to \neg (u \in y_{2} \to u \in y)))$
\item $\varphi_{16} :\equiv \forall u\,(\neg ((u \in y \to u \in y_{2}) \to \neg (u \in y_{2} \to u \in y))) \to \varphi_{1}$
\item $\varphi_{17} :\equiv \forall z\,(\neg ((z \in y \to z \in y_{2}) \to \neg (z \in y_{2} \to z \in y))) \to \varphi_{2}$
\item $\varphi_{18} :\equiv \varphi_{1} \to \forall z\,(\neg ((z \in y \to z \in y_{2}) \to \neg (z \in y_{2} \to z \in y)))$
\item $\varphi_{19} :\equiv \varphi_{2} \to \forall u\,(\neg ((u \in y \to u \in y_{2}) \to \neg (u \in y_{2} \to u \in y)))$
\item $\varphi_{20} :\equiv (\forall u\,(u \notin y) \to \forall u\,(u \notin y_{2}) \to (u \in y \leftrightarrow u \in y_{2})) \to \forall u\,(u \notin y) \land \forall u\,(u \notin y_{2}) \to (u \in y \leftrightarrow u \in y_{2})$
\item $\varphi_{21} :\equiv \varphi_{5} \to \forall u\,(u \notin y) \to \forall u\,(u \notin y_{2}) \to \forall u\,(u \in y \leftrightarrow u \in y_{2})$
\item $\varphi_{22} :\equiv (\forall u\,(u \notin y) \to \forall u\,(u \notin y_{2}) \to \forall u\,(u \in y \leftrightarrow u \in y_{2})) \to \varphi_{15} \to \varphi_{14} \to \forall u\,(u \notin y) \to \forall u\,(u \notin y_{2}) \to y = y_{2}$
\end{abbrevlist}
\end{fullproof}

% 自動生成 (prover/build.js): OmegaInd  ノード数 78
\begin{fullproof}{補題(ωは帰納的)}{$\mathrm{Ind}(\omega)$}{78}\phantomsection\label{tree:OmegaInd}
\noindent$\Pi_{1}$:
\begin{ptree}
\AXC{$\forall n \in \omega\,(S(n) \in \omega) \to \forall y\,\forall n \in \omega\,(S(n) \in \omega)$\lbl{A4}}
\AXC{$\varphi_{17}$\lbl{A3}}
\AXC{$\forall n \in \omega\,(S(n) \in \omega) \to y \in \omega \to S(y) \in \omega$\lbl{A2}}
\UIC{$\forall y\,(\forall n \in \omega\,(S(n) \in \omega) \to y \in \omega \to S(y) \in \omega)$}
\AXC{$\varphi_{23}$\lbl{A1}}
\BIC{$\varphi_{17} \to \varphi_{18}$}
\BIC{$\varphi_{18}$}
\BIC{$\forall n \in \omega\,(S(n) \in \omega) \to \forall y \in \omega\,(S(y) \in \omega)$}
\end{ptree}
\noindent$\Pi_{2}$:
\begin{ptree}
\AXC{$\forall y \in \omega\,(S(y) \in \omega) \to \forall n\,\forall y \in \omega\,(S(y) \in \omega)$\lbl{A4}}
\AXC{$\varphi_{19}$\lbl{A3}}
\AXC{$\forall y \in \omega\,(S(y) \in \omega) \to n \in \omega \to S(n) \in \omega$\lbl{A2}}
\UIC{$\forall n\,(\forall y \in \omega\,(S(y) \in \omega) \to n \in \omega \to S(n) \in \omega)$}
\AXC{$\varphi_{24}$\lbl{A1}}
\BIC{$\varphi_{19} \to \varphi_{20}$}
\BIC{$\varphi_{20}$}
\BIC{$\forall y \in \omega\,(S(y) \in \omega) \to \forall n \in \omega\,(S(n) \in \omega)$}
\end{ptree}
\noindent$\Pi_{3}$:
\begin{ptree}
\AXC{$\forall n\,(n \in \omega \leftrightarrow \forall x\,(\mathrm{Ind}(x) \to n \in x))$\lbl{Def$\omega$}}
\AXC{$\varphi_{6}$\lbl{A2}}
\BIC{$n \in \omega \leftrightarrow \forall x\,(\mathrm{Ind}(x) \to n \in x)$}
\AXC{$n \in \omega \to n \in \omega$\lbl{A1}}
\AXC{$(n \in \omega \to n \in \omega) \to \varphi_{3}$\lbl{A1}}
\BIC{$\varphi_{3}$}
\BIC{$n \in \omega \to \forall x\,(\mathrm{Ind}(x) \to n \in x)$}
\end{ptree}
\noindent$\Pi_{4}$:
\begin{ptree}
\AXC{$\stackrel{\Pi_{3}}{\vdots}$}\noLine
\UIC{$n \in \omega \to \forall x\,(\mathrm{Ind}(x) \to n \in x)$}
\AXC{$(n \in \omega \to \forall x\,(\mathrm{Ind}(x) \to n \in x)) \to \varphi_{5}$\lbl{A1}}
\BIC{$\varphi_{5}$}
\end{ptree}
\noindent$\Pi_{5}$:
\begin{ptree}
\AXC{$n \in \omega \to \mathrm{Ind}(x) \to \mathrm{Ind}(x)$\lbl{A1}}
\AXC{$(n \in \omega \to \mathrm{Ind}(x) \to \mathrm{Ind}(x)) \to \varphi_{14}$\lbl{A1}}
\BIC{$\varphi_{14}$}
\end{ptree}
\noindent$\Pi_{6}$:
\begin{ptree}
\AXC{$\forall x\,(\mathrm{Ind}(x) \leftrightarrow \emptyset \in x \land \forall y \in x\,(S(y) \in x))$\lbl{DefInd}}
\AXC{$\varphi_{15}$\lbl{A2}}
\BIC{$\mathrm{Ind}(x) \leftrightarrow \emptyset \in x \land \forall y \in x\,(S(y) \in x)$}
\AXC{$\stackrel{\Pi_{5}}{\vdots}$}\noLine
\UIC{$\varphi_{14}$}
\BIC{$n \in \omega \to \mathrm{Ind}(x) \to \forall y \in x\,(S(y) \in x)$}
\end{ptree}
\noindent$\Pi_{7}$:
\begin{ptree}
\AXC{$\stackrel{\Pi_{6}}{\vdots}$}\noLine
\UIC{$n \in \omega \to \mathrm{Ind}(x) \to \forall y \in x\,(S(y) \in x)$}
\AXC{$(n \in \omega \to \mathrm{Ind}(x) \to \forall y \in x\,(S(y) \in x)) \to \varphi_{12}$\lbl{A1}}
\BIC{$\varphi_{12}$}
\end{ptree}
\noindent$\Pi_{8}$:
\begin{ptree}
\AXC{$\forall y \in x\,(S(y) \in x) \to n \in x \to S(n) \in x$\lbl{A2}}
\AXC{$\stackrel{\Pi_{7}}{\vdots}$}\noLine
\UIC{$\varphi_{12}$}
\BIC{$n \in \omega \to \mathrm{Ind}(x) \to n \in x \to S(n) \in x$}
\AXC{$(n \in \omega \to \mathrm{Ind}(x) \to n \in x \to S(n) \in x) \to \varphi_{1}$\lbl{A1}}
\BIC{$\varphi_{1}$}
\end{ptree}
\noindent$\Pi_{9}$:
\begin{ptree}
\AXC{$\forall x\,(\mathrm{Ind}(x) \to n \in x) \to \mathrm{Ind}(x) \to n \in x$\lbl{A2}}
\AXC{$\stackrel{\Pi_{4}}{\vdots}$}\noLine
\UIC{$\varphi_{5}$}
\BIC{$n \in \omega \to \mathrm{Ind}(x) \to n \in x$}
\AXC{$\stackrel{\Pi_{8}}{\vdots}$}\noLine
\UIC{$\varphi_{1}$}
\BIC{$n \in \omega \to \mathrm{Ind}(x) \to S(n) \in x$}
\UIC{$\forall x\,(n \in \omega \to \mathrm{Ind}(x) \to S(n) \in x)$}
\AXC{$\forall x\,(n \in \omega \to \mathrm{Ind}(x) \to S(n) \in x) \to \varphi_{9} \to \varphi_{2}$\lbl{A1}}
\BIC{$\varphi_{9} \to \varphi_{2}$}
\end{ptree}
\noindent$\Pi_{10}$:
\begin{ptree}
\AXC{$n \in \omega \to \forall x\,(n \in \omega)$\lbl{A4}}
\AXC{$\varphi_{9}$\lbl{A3}}
\AXC{$\stackrel{\Pi_{9}}{\vdots}$}\noLine
\UIC{$\varphi_{9} \to \varphi_{2}$}
\BIC{$\varphi_{2}$}
\BIC{$n \in \omega \to \forall x\,(\mathrm{Ind}(x) \to S(n) \in x)$}
\AXC{$(n \in \omega \to \forall x\,(\mathrm{Ind}(x) \to S(n) \in x)) \to \varphi_{4}$\lbl{A1}}
\BIC{$\varphi_{4}$}
\end{ptree}
\noindent$\Pi_{11}$:
\begin{ptree}
\AXC{$\mathrm{Ind}(x) \to \mathrm{Ind}(x)$\lbl{A1}}
\AXC{$(\mathrm{Ind}(x) \to \mathrm{Ind}(x)) \to \varphi_{7}$\lbl{A1}}
\BIC{$\varphi_{7}$}
\end{ptree}
\noindent$\Pi_{12}$:
\begin{ptree}
\AXC{$\forall x\,(\mathrm{Ind}(x) \leftrightarrow \emptyset \in x \land \forall y \in x\,(S(y) \in x))$\lbl{DefInd}}
\AXC{$\varphi_{15}$\lbl{A2}}
\BIC{$\mathrm{Ind}(x) \leftrightarrow \emptyset \in x \land \forall y \in x\,(S(y) \in x)$}
\AXC{$\stackrel{\Pi_{11}}{\vdots}$}\noLine
\UIC{$\varphi_{7}$}
\BIC{$\mathrm{Ind}(x) \to \emptyset \in x$}
\UIC{$\forall x\,(\mathrm{Ind}(x) \to \emptyset \in x)$}
\end{ptree}
\noindent$\Pi_{13}$:
\begin{ptree}
\AXC{$\stackrel{\Pi_{12}}{\vdots}$}\noLine
\UIC{$\forall x\,(\mathrm{Ind}(x) \to \emptyset \in x)$}
\AXC{$\forall n\,(n \in \omega \leftrightarrow \forall x\,(\mathrm{Ind}(x) \to n \in x))$\lbl{Def$\omega$}}
\AXC{$\varphi_{11}$\lbl{A2}}
\BIC{$\emptyset \in \omega \leftrightarrow \forall x\,(\mathrm{Ind}(x) \to \emptyset \in x)$}
\AXC{$\varphi_{13}$\lbl{A1}}
\BIC{$\forall x\,(\mathrm{Ind}(x) \to \emptyset \in x) \to \emptyset \in \omega$}
\BIC{$\emptyset \in \omega$}
\end{ptree}
\noindent$\Pi_{14}$:
\begin{ptree}
\AXC{$\stackrel{\Pi_{13}}{\vdots}$}\noLine
\UIC{$\emptyset \in \omega$}
\AXC{$\forall x\,(\mathrm{Ind}(x) \leftrightarrow \emptyset \in x \land \forall y \in x\,(S(y) \in x))$\lbl{DefInd}}
\AXC{$\varphi_{16}$\lbl{A2}}
\BIC{$\mathrm{Ind}(\omega) \leftrightarrow \emptyset \in \omega \land \forall y \in \omega\,(S(y) \in \omega)$}
\AXC{$\varphi_{22}$\lbl{A1}}
\BIC{$\emptyset \in \omega \to \varphi_{21}$}
\BIC{$\varphi_{21}$}
\end{ptree}
\noindent$\Pi_{15}$:
\begin{ptree}
\AXC{$\forall n\,(n \in \omega \leftrightarrow \forall x\,(\mathrm{Ind}(x) \to n \in x))$\lbl{Def$\omega$}}
\AXC{$\varphi_{8}$\lbl{A2}}
\BIC{$S(n) \in \omega \leftrightarrow \forall x\,(\mathrm{Ind}(x) \to S(n) \in x)$}
\AXC{$\stackrel{\Pi_{10}}{\vdots}$}\noLine
\UIC{$\varphi_{4}$}
\BIC{$n \in \omega \to S(n) \in \omega$}
\UIC{$\forall n \in \omega\,(S(n) \in \omega)$}
\AXC{$\stackrel{\Pi_{14}}{\vdots}$}\noLine
\UIC{$\varphi_{21}$}
\BIC{$(\forall y \in \omega\,(S(y) \in \omega) \to \forall n \in \omega\,(S(n) \in \omega)) \to \varphi_{10}$}
\end{ptree}
\noindent$\Pi_{16}$:
\begin{ptree}
\AXC{$\stackrel{\Pi_{2}}{\vdots}$}\noLine
\UIC{$\forall y \in \omega\,(S(y) \in \omega) \to \forall n \in \omega\,(S(n) \in \omega)$}
\AXC{$\stackrel{\Pi_{15}}{\vdots}$}\noLine
\UIC{$(\forall y \in \omega\,(S(y) \in \omega) \to \forall n \in \omega\,(S(n) \in \omega)) \to \varphi_{10}$}
\BIC{$\varphi_{10}$}
\end{ptree}
\noindent 主木:
\begin{ptree}
\AXC{$\stackrel{\Pi_{1}}{\vdots}$}\noLine
\UIC{$\forall n \in \omega\,(S(n) \in \omega) \to \forall y \in \omega\,(S(y) \in \omega)$}
\AXC{$\stackrel{\Pi_{16}}{\vdots}$}\noLine
\UIC{$\varphi_{10}$}
\BIC{$\mathrm{Ind}(\omega)$}
\end{ptree}
\begin{abbrevlist}
\item $\varphi_{1} :\equiv (n \in \omega \to \mathrm{Ind}(x) \to n \in x) \to n \in \omega \to \mathrm{Ind}(x) \to S(n) \in x$
\item $\varphi_{2} :\equiv (n \in \omega \to \forall x\,(n \in \omega)) \to n \in \omega \to \forall x\,(\mathrm{Ind}(x) \to S(n) \in x)$
\item $\varphi_{3} :\equiv (n \in \omega \leftrightarrow \forall x\,(\mathrm{Ind}(x) \to n \in x)) \to n \in \omega \to \forall x\,(\mathrm{Ind}(x) \to n \in x)$
\item $\varphi_{4} :\equiv (S(n) \in \omega \leftrightarrow \forall x\,(\mathrm{Ind}(x) \to S(n) \in x)) \to n \in \omega \to S(n) \in \omega$
\item $\varphi_{5} :\equiv (\forall x\,(\mathrm{Ind}(x) \to n \in x) \to \mathrm{Ind}(x) \to n \in x) \to n \in \omega \to \mathrm{Ind}(x) \to n \in x$
\item $\varphi_{6} :\equiv \forall n\,(n \in \omega \leftrightarrow \forall x\,(\mathrm{Ind}(x) \to n \in x)) \to (n \in \omega \leftrightarrow \forall x\,(\mathrm{Ind}(x) \to n \in x))$
\item $\varphi_{7} :\equiv (\mathrm{Ind}(x) \leftrightarrow \emptyset \in x \land \forall y \in x\,(S(y) \in x)) \to \mathrm{Ind}(x) \to \emptyset \in x$
\item $\varphi_{8} :\equiv \forall n\,(n \in \omega \leftrightarrow \forall x\,(\mathrm{Ind}(x) \to n \in x)) \to (S(n) \in \omega \leftrightarrow \forall x\,(\mathrm{Ind}(x) \to S(n) \in x))$
\item $\varphi_{9} :\equiv \forall x\,(n \in \omega \to \mathrm{Ind}(x) \to S(n) \in x) \to \forall x\,(n \in \omega) \to \forall x\,(\mathrm{Ind}(x) \to S(n) \in x)$
\item $\varphi_{10} :\equiv (\forall n \in \omega\,(S(n) \in \omega) \to \forall y \in \omega\,(S(y) \in \omega)) \to \mathrm{Ind}(\omega)$
\item $\varphi_{11} :\equiv \forall n\,(n \in \omega \leftrightarrow \forall x\,(\mathrm{Ind}(x) \to n \in x)) \to (\emptyset \in \omega \leftrightarrow \forall x\,(\mathrm{Ind}(x) \to \emptyset \in x))$
\item $\varphi_{12} :\equiv (\forall y \in x\,(S(y) \in x) \to n \in x \to S(n) \in x) \to n \in \omega \to \mathrm{Ind}(x) \to n \in x \to S(n) \in x$
\item $\varphi_{13} :\equiv (\emptyset \in \omega \leftrightarrow \forall x\,(\mathrm{Ind}(x) \to \emptyset \in x)) \to \forall x\,(\mathrm{Ind}(x) \to \emptyset \in x) \to \emptyset \in \omega$
\item $\varphi_{14} :\equiv (\mathrm{Ind}(x) \leftrightarrow \emptyset \in x \land \forall y \in x\,(S(y) \in x)) \to n \in \omega \to \mathrm{Ind}(x) \to \forall y \in x\,(S(y) \in x)$
\item $\varphi_{15} :\equiv \forall x\,(\mathrm{Ind}(x) \leftrightarrow \emptyset \in x \land \forall y \in x\,(S(y) \in x)) \to (\mathrm{Ind}(x) \leftrightarrow \emptyset \in x \land \forall y \in x\,(S(y) \in x))$
\item $\varphi_{16} :\equiv \forall x\,(\mathrm{Ind}(x) \leftrightarrow \emptyset \in x \land \forall y \in x\,(S(y) \in x)) \to (\mathrm{Ind}(\omega) \leftrightarrow \emptyset \in \omega \land \forall y \in \omega\,(S(y) \in \omega))$
\item $\varphi_{17} :\equiv \forall y\,(\forall n \in \omega\,(S(n) \in \omega) \to y \in \omega \to S(y) \in \omega) \to \forall y\,\forall n \in \omega\,(S(n) \in \omega) \to \forall y \in \omega\,(S(y) \in \omega)$
\item $\varphi_{18} :\equiv (\forall n \in \omega\,(S(n) \in \omega) \to \forall y\,\forall n \in \omega\,(S(n) \in \omega)) \to \forall n \in \omega\,(S(n) \in \omega) \to \forall y \in \omega\,(S(y) \in \omega)$
\item $\varphi_{19} :\equiv \forall n\,(\forall y \in \omega\,(S(y) \in \omega) \to n \in \omega \to S(n) \in \omega) \to \forall n\,\forall y \in \omega\,(S(y) \in \omega) \to \forall n \in \omega\,(S(n) \in \omega)$
\item $\varphi_{20} :\equiv (\forall y \in \omega\,(S(y) \in \omega) \to \forall n\,\forall y \in \omega\,(S(y) \in \omega)) \to \forall y \in \omega\,(S(y) \in \omega) \to \forall n \in \omega\,(S(n) \in \omega)$
\item $\varphi_{21} :\equiv \forall n \in \omega\,(S(n) \in \omega) \to (\forall y \in \omega\,(S(y) \in \omega) \to \forall n \in \omega\,(S(n) \in \omega)) \to \varphi_{10}$
\item $\varphi_{22} :\equiv (\mathrm{Ind}(\omega) \leftrightarrow \emptyset \in \omega \land \forall y \in \omega\,(S(y) \in \omega)) \to \emptyset \in \omega \to \varphi_{21}$
\item $\varphi_{23} :\equiv \forall y\,(\forall n \in \omega\,(S(n) \in \omega) \to y \in \omega \to S(y) \in \omega) \to \varphi_{17} \to \varphi_{18}$
\item $\varphi_{24} :\equiv \forall n\,(\forall y \in \omega\,(S(y) \in \omega) \to n \in \omega \to S(n) \in \omega) \to \varphi_{19} \to \varphi_{20}$
\end{abbrevlist}
\end{fullproof}

% 自動生成 (prover/build.js): Induction  ノード数 111
\begin{fullproof}{定理(数学的帰納法)}{$\mathrm{Phi}(\emptyset) \to \forall n \in \omega\,(\mathrm{Phi}(n) \to \mathrm{Phi}(S(n))) \to \forall n \in \omega\,(\mathrm{Phi}(n))$}{111}\phantomsection\label{tree:Induction}
\noindent$\Pi_{1}$:
\begin{ptree}
\AXC{$\forall x\,(\mathrm{Ind}(x) \leftrightarrow \emptyset \in x \land \forall y \in x\,(S(y) \in x))$\lbl{DefInd}}
\AXC{$\varphi_{35}$\lbl{A2}}
\BIC{$\mathrm{Ind}(\omega) \leftrightarrow \emptyset \in \omega \land \forall y \in \omega\,(S(y) \in \omega)$}
\AXC{$\mathrm{Ind}(\omega)$\lbl{補題(ωは帰納的)}}
\AXC{$\varphi_{46}$\lbl{A1}}
\BIC{$(\mathrm{Ind}(\omega) \leftrightarrow \emptyset \in \omega \land \forall y \in \omega\,(S(y) \in \omega)) \to \varphi_{45}$}
\BIC{$\varphi_{45}$}
\end{ptree}
\noindent$\Pi_{2}$:
\begin{ptree}
\AXC{$\stackrel{\Pi_{1}}{\vdots}$}\noLine
\UIC{$\varphi_{45}$}
\AXC{$\varphi_{45} \to \varphi_{43}$\lbl{A1}}
\BIC{$\varphi_{43}$}
\AXC{$\varphi_{43} \to (\forall y \in \omega\,(S(y) \in \omega) \to y \in \omega \to S(y) \in \omega) \to \varphi_{41}$\lbl{A1}}
\BIC{$(\forall y \in \omega\,(S(y) \in \omega) \to y \in \omega \to S(y) \in \omega) \to \varphi_{41}$}
\end{ptree}
\noindent$\Pi_{3}$:
\begin{ptree}
\AXC{$\forall y \in \omega\,(S(y) \in \omega) \to y \in \omega \to S(y) \in \omega$\lbl{A2}}
\AXC{$\stackrel{\Pi_{2}}{\vdots}$}\noLine
\UIC{$(\forall y \in \omega\,(S(y) \in \omega) \to y \in \omega \to S(y) \in \omega) \to \varphi_{41}$}
\BIC{$\varphi_{41}$}
\end{ptree}
\noindent$\Pi_{4}$:
\begin{ptree}
\AXC{$\varphi_{10}$\lbl{A2}}
\AXC{$\mathrm{Phi}(\emptyset) \to \varphi_{12}$\lbl{A1}}
\AXC{$(\mathrm{Phi}(\emptyset) \to \varphi_{12}) \to \varphi_{10} \to \mathrm{Phi}(\emptyset) \to \varphi_{10}$\lbl{A1}}
\BIC{$\varphi_{10} \to \mathrm{Phi}(\emptyset) \to \varphi_{10}$}
\BIC{$\mathrm{Phi}(\emptyset) \to \varphi_{10}$}
\end{ptree}
\noindent$\Pi_{5}$:
\begin{ptree}
\AXC{$\mathrm{Phi}(\emptyset) \to \varphi_{27}$\lbl{A1}}
\AXC{$\varphi_{5}$\lbl{A2}}
\AXC{$\varphi_{39}$\lbl{A1}}
\AXC{$\varphi_{39} \to \varphi_{5} \to \varphi_{36}$\lbl{A1}}
\BIC{$\varphi_{5} \to \varphi_{36}$}
\BIC{$\varphi_{36}$}
\AXC{$\varphi_{36} \to (\mathrm{Phi}(\emptyset) \to \varphi_{27}) \to \varphi_{37}$\lbl{A1}}
\BIC{$(\mathrm{Phi}(\emptyset) \to \varphi_{27}) \to \varphi_{37}$}
\BIC{$\varphi_{37}$}
\end{ptree}
\noindent$\Pi_{6}$:
\begin{ptree}
\AXC{$\stackrel{\Pi_{3}}{\vdots}$}\noLine
\UIC{$\varphi_{41}$}
\AXC{$\stackrel{\Pi_{4}}{\vdots}$}\noLine
\UIC{$\mathrm{Phi}(\emptyset) \to \varphi_{10}$}
\AXC{$\stackrel{\Pi_{5}}{\vdots}$}\noLine
\UIC{$\varphi_{37}$}
\AXC{$\varphi_{37} \to (\mathrm{Phi}(\emptyset) \to \varphi_{10}) \to \varphi_{41} \to \varphi_{42} \to \mathrm{Phi}(\emptyset) \to \varphi_{28}$\lbl{A1}}
\BIC{$(\mathrm{Phi}(\emptyset) \to \varphi_{10}) \to \varphi_{41} \to \varphi_{42} \to \mathrm{Phi}(\emptyset) \to \varphi_{28}$}
\BIC{$\varphi_{41} \to \varphi_{42} \to \mathrm{Phi}(\emptyset) \to \varphi_{28}$}
\BIC{$\varphi_{42} \to \mathrm{Phi}(\emptyset) \to \varphi_{28}$}
\end{ptree}
\noindent$\Pi_{7}$:
\begin{ptree}
\AXC{$\varphi_{17}$\lbl{A2}}
\AXC{$\varphi_{39}$\lbl{A1}}
\AXC{$\varphi_{39} \to \varphi_{17} \to \varphi_{42}$\lbl{A1}}
\BIC{$\varphi_{17} \to \varphi_{42}$}
\BIC{$\varphi_{42}$}
\AXC{$\stackrel{\Pi_{6}}{\vdots}$}\noLine
\UIC{$\varphi_{42} \to \mathrm{Phi}(\emptyset) \to \varphi_{28}$}
\BIC{$\mathrm{Phi}(\emptyset) \to \varphi_{28}$}
\AXC{$(\mathrm{Phi}(\emptyset) \to \varphi_{28}) \to \varphi_{24} \to y \in B \to S(y) \in B$\lbl{A1}}
\BIC{$\varphi_{24} \to y \in B \to S(y) \in B$}
\UIC{$\forall y\,(\varphi_{24} \to y \in B \to S(y) \in B)$}
\end{ptree}
\noindent$\Pi_{8}$:
\begin{ptree}
\AXC{$\varphi_{48}$\lbl{A3}}
\AXC{$\stackrel{\Pi_{7}}{\vdots}$}\noLine
\UIC{$\forall y\,(\varphi_{24} \to y \in B \to S(y) \in B)$}
\AXC{$\forall y\,(\varphi_{24} \to y \in B \to S(y) \in B) \to \varphi_{48} \to (\varphi_{24} \to \forall y\,(\varphi_{24})) \to \varphi_{32}$\lbl{A1}}
\BIC{$\varphi_{48} \to (\varphi_{24} \to \forall y\,(\varphi_{24})) \to \varphi_{32}$}
\BIC{$(\varphi_{24} \to \forall y\,(\varphi_{24})) \to \varphi_{32}$}
\end{ptree}
\noindent$\Pi_{9}$:
\begin{ptree}
\AXC{$\stackrel{\Pi_{1}}{\vdots}$}\noLine
\UIC{$\varphi_{45}$}
\AXC{$\varphi_{23}$\lbl{A2}}
\AXC{$\varphi_{39}$\lbl{A1}}
\AXC{$\varphi_{39} \to \varphi_{23} \to \varphi_{44}$\lbl{A1}}
\BIC{$\varphi_{23} \to \varphi_{44}$}
\BIC{$\varphi_{44}$}
\AXC{$\varphi_{44} \to \varphi_{45} \to (\mathrm{Phi}(\emptyset) \to \mathrm{Phi}(\emptyset)) \to \mathrm{Phi}(\emptyset) \to \varphi_{26}$\lbl{A1}}
\BIC{$\varphi_{45} \to (\mathrm{Phi}(\emptyset) \to \mathrm{Phi}(\emptyset)) \to \mathrm{Phi}(\emptyset) \to \varphi_{26}$}
\BIC{$(\mathrm{Phi}(\emptyset) \to \mathrm{Phi}(\emptyset)) \to \mathrm{Phi}(\emptyset) \to \varphi_{26}$}
\end{ptree}
\noindent$\Pi_{10}$:
\begin{ptree}
\AXC{$\forall x\,(\mathrm{Ind}(x) \leftrightarrow \emptyset \in x \land \forall y \in x\,(S(y) \in x))$\lbl{DefInd}}
\AXC{$\varphi_{31}$\lbl{A2}}
\BIC{$\mathrm{Ind}(B) \leftrightarrow \emptyset \in B \land \forall y \in B\,(S(y) \in B)$}
\AXC{$\varphi_{52}$\lbl{A1}}
\BIC{$(\mathrm{Phi}(\emptyset) \to \varphi_{26}) \to \varphi_{32} \to \mathrm{Phi}(\emptyset) \to \varphi_{22}$}
\end{ptree}
\noindent$\Pi_{11}$:
\begin{ptree}
\AXC{$\mathrm{Phi}(\emptyset) \to \mathrm{Phi}(\emptyset)$\lbl{A1}}
\AXC{$\stackrel{\Pi_{9}}{\vdots}$}\noLine
\UIC{$(\mathrm{Phi}(\emptyset) \to \mathrm{Phi}(\emptyset)) \to \mathrm{Phi}(\emptyset) \to \varphi_{26}$}
\BIC{$\mathrm{Phi}(\emptyset) \to \varphi_{26}$}
\AXC{$\stackrel{\Pi_{10}}{\vdots}$}\noLine
\UIC{$(\mathrm{Phi}(\emptyset) \to \varphi_{26}) \to \varphi_{32} \to \mathrm{Phi}(\emptyset) \to \varphi_{22}$}
\BIC{$\varphi_{32} \to \mathrm{Phi}(\emptyset) \to \varphi_{22}$}
\end{ptree}
\noindent$\Pi_{12}$:
\begin{ptree}
\AXC{$\varphi_{24} \to \forall y\,(\varphi_{24})$\lbl{A4}}
\AXC{$\stackrel{\Pi_{8}}{\vdots}$}\noLine
\UIC{$(\varphi_{24} \to \forall y\,(\varphi_{24})) \to \varphi_{32}$}
\BIC{$\varphi_{32}$}
\AXC{$\stackrel{\Pi_{11}}{\vdots}$}\noLine
\UIC{$\varphi_{32} \to \mathrm{Phi}(\emptyset) \to \varphi_{22}$}
\BIC{$\mathrm{Phi}(\emptyset) \to \varphi_{22}$}
\end{ptree}
\noindent$\Pi_{13}$:
\begin{ptree}
\AXC{$\forall n\,(n \in \omega \leftrightarrow \forall x\,(\mathrm{Ind}(x) \to n \in x))$\lbl{Def$\omega$}}
\AXC{$\varphi_{11}$\lbl{A2}}
\BIC{$n \in \omega \leftrightarrow \forall x\,(\mathrm{Ind}(x) \to n \in x)$}
\AXC{$\mathrm{Phi}(\emptyset) \to \varphi_{30}$\lbl{A1}}
\AXC{$\varphi_{50}$\lbl{A1}}
\BIC{$(n \in \omega \leftrightarrow \forall x\,(\mathrm{Ind}(x) \to n \in x)) \to \varphi_{40}$}
\BIC{$\varphi_{40}$}
\AXC{$\varphi_{51}$\lbl{A1}}
\BIC{$(\forall x\,(\mathrm{Ind}(x) \to n \in x) \to \mathrm{Ind}(B) \to n \in B) \to \varphi_{38}$}
\end{ptree}
\noindent$\Pi_{14}$:
\begin{ptree}
\AXC{$\forall x\,(\mathrm{Ind}(x) \to n \in x) \to \mathrm{Ind}(B) \to n \in B$\lbl{A2}}
\AXC{$\stackrel{\Pi_{13}}{\vdots}$}\noLine
\UIC{$(\forall x\,(\mathrm{Ind}(x) \to n \in x) \to \mathrm{Ind}(B) \to n \in B) \to \varphi_{38}$}
\BIC{$\varphi_{38}$}
\end{ptree}
\noindent$\Pi_{15}$:
\begin{ptree}
\AXC{$\stackrel{\Pi_{12}}{\vdots}$}\noLine
\UIC{$\mathrm{Phi}(\emptyset) \to \varphi_{22}$}
\AXC{$\stackrel{\Pi_{14}}{\vdots}$}\noLine
\UIC{$\varphi_{38}$}
\AXC{$\varphi_{38} \to (\mathrm{Phi}(\emptyset) \to \varphi_{22}) \to \varphi_{34} \to \mathrm{Phi}(\emptyset) \to \varphi_{29}$\lbl{A1}}
\BIC{$(\mathrm{Phi}(\emptyset) \to \varphi_{22}) \to \varphi_{34} \to \mathrm{Phi}(\emptyset) \to \varphi_{29}$}
\BIC{$\varphi_{34} \to \mathrm{Phi}(\emptyset) \to \varphi_{29}$}
\end{ptree}
\noindent$\Pi_{16}$:
\begin{ptree}
\AXC{$\varphi_{4}$\lbl{A2}}
\AXC{$\varphi_{39}$\lbl{A1}}
\AXC{$\varphi_{39} \to \varphi_{4} \to \varphi_{34}$\lbl{A1}}
\BIC{$\varphi_{4} \to \varphi_{34}$}
\BIC{$\varphi_{34}$}
\AXC{$\stackrel{\Pi_{15}}{\vdots}$}\noLine
\UIC{$\varphi_{34} \to \mathrm{Phi}(\emptyset) \to \varphi_{29}$}
\BIC{$\mathrm{Phi}(\emptyset) \to \varphi_{29}$}
\AXC{$(\mathrm{Phi}(\emptyset) \to \varphi_{29}) \to \varphi_{24} \to n \in \omega \to \mathrm{Phi}(n)$\lbl{A1}}
\BIC{$\varphi_{24} \to n \in \omega \to \mathrm{Phi}(n)$}
\UIC{$\forall n\,(\varphi_{24} \to n \in \omega \to \mathrm{Phi}(n))$}
\end{ptree}
\noindent$\Pi_{17}$:
\begin{ptree}
\AXC{$\varphi_{47}$\lbl{A3}}
\AXC{$\stackrel{\Pi_{16}}{\vdots}$}\noLine
\UIC{$\forall n\,(\varphi_{24} \to n \in \omega \to \mathrm{Phi}(n))$}
\AXC{$\forall n\,(\varphi_{24} \to n \in \omega \to \mathrm{Phi}(n)) \to \varphi_{47} \to (\varphi_{24} \to \forall n\,(\varphi_{24})) \to \varphi_{33}$\lbl{A1}}
\BIC{$\varphi_{47} \to (\varphi_{24} \to \forall n\,(\varphi_{24})) \to \varphi_{33}$}
\BIC{$(\varphi_{24} \to \forall n\,(\varphi_{24})) \to \varphi_{33}$}
\end{ptree}
\noindent$\Pi_{18}$:
\begin{ptree}
\AXC{$\varphi_{24} \to \forall n\,(\varphi_{24})$\lbl{A4}}
\AXC{$\stackrel{\Pi_{17}}{\vdots}$}\noLine
\UIC{$(\varphi_{24} \to \forall n\,(\varphi_{24})) \to \varphi_{33}$}
\BIC{$\varphi_{33}$}
\AXC{$\varphi_{33} \to \varphi_{19} \to \neg \forall x\,(x \in B \leftrightarrow x \in \omega \land \mathrm{Phi}(x))$\lbl{A1}}
\BIC{$\varphi_{19} \to \neg \forall x\,(x \in B \leftrightarrow x \in \omega \land \mathrm{Phi}(x))$}
\UIC{$\forall B\,(\varphi_{19} \to \neg \forall x\,(x \in B \leftrightarrow x \in \omega \land \mathrm{Phi}(x)))$}
\end{ptree}
\noindent$\Pi_{19}$:
\begin{ptree}
\AXC{$\forall a\,\exists y\,\forall x\,(x \in y \leftrightarrow x \in a \land \mathrm{Phi}(x))$\lbl{Sep}}
\AXC{$\varphi_{14}$\lbl{A2}}
\BIC{$\exists y\,\forall x\,(x \in y \leftrightarrow x \in \omega \land \mathrm{Phi}(x))$}
\AXC{$\exists y\,\forall x\,(x \in y \leftrightarrow x \in \omega \land \mathrm{Phi}(x)) \to \mathrm{Phi}(\emptyset) \to \varphi_{16}$\lbl{A1}}
\BIC{$\mathrm{Phi}(\emptyset) \to \varphi_{16}$}
\end{ptree}
\noindent$\Pi_{20}$:
\begin{ptree}
\AXC{$\varphi_{49}$\lbl{A3}}
\AXC{$\stackrel{\Pi_{18}}{\vdots}$}\noLine
\UIC{$\forall B\,(\varphi_{19} \to \neg \forall x\,(x \in B \leftrightarrow x \in \omega \land \mathrm{Phi}(x)))$}
\AXC{$\stackrel{\Pi_{19}}{\vdots}$}\noLine
\UIC{$\mathrm{Phi}(\emptyset) \to \varphi_{16}$}
\AXC{$\varphi_{54}$\lbl{A1}}
\BIC{$\forall B\,(\varphi_{19} \to \neg \forall x\,(x \in B \leftrightarrow x \in \omega \land \mathrm{Phi}(x))) \to \varphi_{49} \to \varphi_{53}$}
\BIC{$\varphi_{49} \to \varphi_{53}$}
\BIC{$\varphi_{53}$}
\end{ptree}
\noindent 主木:
\begin{ptree}
\AXC{$\varphi_{20}$\lbl{A4}}
\AXC{$\forall y\,(\varphi_{13}) \to \varphi_{21}$\lbl{A3}}
\AXC{$\varphi_{13}$\lbl{A2}}
\UIC{$\forall y\,(\varphi_{13})$}
\AXC{$\varphi_{19} \to \forall B\,(\varphi_{19})$\lbl{A4}}
\AXC{$\stackrel{\Pi_{20}}{\vdots}$}\noLine
\UIC{$\varphi_{53}$}
\BIC{$\forall y\,(\varphi_{13}) \to (\forall y\,(\varphi_{13}) \to \varphi_{21}) \to \varphi_{20} \to \mathrm{Phi}(\emptyset) \to \varphi_{3}$}
\BIC{$(\forall y\,(\varphi_{13}) \to \varphi_{21}) \to \varphi_{20} \to \mathrm{Phi}(\emptyset) \to \varphi_{3}$}
\BIC{$\varphi_{20} \to \mathrm{Phi}(\emptyset) \to \varphi_{3}$}
\BIC{$\mathrm{Phi}(\emptyset) \to \varphi_{3}$}
\end{ptree}
\begin{abbrevlist}
\item $\varphi_{1} :\equiv \forall x\,(x \in B \leftrightarrow x \in \omega \land \mathrm{Phi}(x)) \to \forall y \in B\,(S(y) \in B)$
\item $\varphi_{2} :\equiv \forall x\,(x \in B \leftrightarrow x \in \omega \land \mathrm{Phi}(x)) \to \forall n \in \omega\,(\mathrm{Phi}(n))$
\item $\varphi_{3} :\equiv \forall n \in \omega\,(\mathrm{Phi}(n) \to \mathrm{Phi}(S(n))) \to \forall n \in \omega\,(\mathrm{Phi}(n))$
\item $\varphi_{4} :\equiv \forall x\,(x \in B \leftrightarrow x \in \omega \land \mathrm{Phi}(x)) \to (n \in B \leftrightarrow n \in \omega \land \mathrm{Phi}(n))$
\item $\varphi_{5} :\equiv \forall x\,(x \in B \leftrightarrow x \in \omega \land \mathrm{Phi}(x)) \to (y \in B \leftrightarrow y \in \omega \land \mathrm{Phi}(y))$
\item $\varphi_{6} :\equiv \forall x\,(x \in B \leftrightarrow x \in \omega \land \mathrm{Phi}(x)) \to y \in B \to y \in \omega \land \mathrm{Phi}(y)$
\item $\varphi_{7} :\equiv \forall x\,(x \in B \leftrightarrow x \in \omega \land \mathrm{Phi}(x)) \to n \in \omega \to \mathrm{Ind}(B) \to n \in B$
\item $\varphi_{8} :\equiv \forall x\,(x \in B \leftrightarrow x \in \omega \land \mathrm{Phi}(x)) \to \forall x\,(x \in B \leftrightarrow x \in \omega \land \mathrm{Phi}(x))$
\item $\varphi_{9} :\equiv \forall x\,(x \in B \leftrightarrow x \in \omega \land \mathrm{Phi}(x)) \to n \in \omega \to \forall x\,(\mathrm{Ind}(x) \to n \in x)$
\item $\varphi_{10} :\equiv \forall n \in \omega\,(\mathrm{Phi}(n) \to \mathrm{Phi}(S(n))) \to y \in \omega \to \mathrm{Phi}(y) \to \mathrm{Phi}(S(y))$
\item $\varphi_{11} :\equiv \forall n\,(n \in \omega \leftrightarrow \forall x\,(\mathrm{Ind}(x) \to n \in x)) \to (n \in \omega \leftrightarrow \forall x\,(\mathrm{Ind}(x) \to n \in x))$
\item $\varphi_{12} :\equiv \forall n \in \omega\,(\mathrm{Phi}(n) \to \mathrm{Phi}(S(n))) \to \forall n \in \omega\,(\mathrm{Phi}(n) \to \mathrm{Phi}(S(n)))$
\item $\varphi_{13} :\equiv \forall B\,(\neg \forall x\,(x \in B \leftrightarrow x \in \omega \land \mathrm{Phi}(x))) \to \neg \forall x\,(x \in y \leftrightarrow x \in \omega \land \mathrm{Phi}(x))$
\item $\varphi_{14} :\equiv \forall a\,\exists y\,\forall x\,(x \in y \leftrightarrow x \in a \land \mathrm{Phi}(x)) \to \exists y\,\forall x\,(x \in y \leftrightarrow x \in \omega \land \mathrm{Phi}(x))$
\item $\varphi_{15} :\equiv \forall x\,(x \in B \leftrightarrow x \in \omega \land \mathrm{Phi}(x)) \to y \in B \to y \in \omega \to S(y) \in \omega$
\item $\varphi_{16} :\equiv \forall n \in \omega\,(\mathrm{Phi}(n) \to \mathrm{Phi}(S(n))) \to \exists y\,\forall x\,(x \in y \leftrightarrow x \in \omega \land \mathrm{Phi}(x))$
\item $\varphi_{17} :\equiv \forall x\,(x \in B \leftrightarrow x \in \omega \land \mathrm{Phi}(x)) \to (S(y) \in B \leftrightarrow S(y) \in \omega \land \mathrm{Phi}(S(y)))$
\item $\varphi_{18} :\equiv \forall x\,(x \in B \leftrightarrow x \in \omega \land \mathrm{Phi}(x)) \to y \in B \to \forall y \in \omega\,(S(y) \in \omega)$
\item $\varphi_{19} :\equiv \mathrm{Phi}(\emptyset) \land \forall n \in \omega\,(\mathrm{Phi}(n) \to \mathrm{Phi}(S(n))) \land \neg \forall n \in \omega\,(\mathrm{Phi}(n))$
\item $\varphi_{20} :\equiv \forall B\,(\neg \forall x\,(x \in B \leftrightarrow x \in \omega \land \mathrm{Phi}(x))) \to \forall y\,\forall B\,(\neg \forall x\,(x \in B \leftrightarrow x \in \omega \land \mathrm{Phi}(x)))$
\item $\varphi_{21} :\equiv \forall y\,\forall B\,(\neg \forall x\,(x \in B \leftrightarrow x \in \omega \land \mathrm{Phi}(x))) \to \forall y\,(\neg \forall x\,(x \in y \leftrightarrow x \in \omega \land \mathrm{Phi}(x)))$
\item $\varphi_{22} :\equiv \forall n \in \omega\,(\mathrm{Phi}(n) \to \mathrm{Phi}(S(n))) \to \forall x\,(x \in B \leftrightarrow x \in \omega \land \mathrm{Phi}(x)) \to \mathrm{Ind}(B)$
\item $\varphi_{23} :\equiv \forall x\,(x \in B \leftrightarrow x \in \omega \land \mathrm{Phi}(x)) \to (\emptyset \in B \leftrightarrow \emptyset \in \omega \land \mathrm{Phi}(\emptyset))$
\item $\varphi_{24} :\equiv \mathrm{Phi}(\emptyset) \land \forall n \in \omega\,(\mathrm{Phi}(n) \to \mathrm{Phi}(S(n))) \land \forall x\,(x \in B \leftrightarrow x \in \omega \land \mathrm{Phi}(x))$
\item $\varphi_{25} :\equiv \forall x\,(x \in B \leftrightarrow x \in \omega \land \mathrm{Phi}(x)) \to \emptyset \in \omega \land \forall y \in \omega\,(S(y) \in \omega)$
\item $\varphi_{26} :\equiv \forall n \in \omega\,(\mathrm{Phi}(n) \to \mathrm{Phi}(S(n))) \to \forall x\,(x \in B \leftrightarrow x \in \omega \land \mathrm{Phi}(x)) \to \emptyset \in B$
\item $\varphi_{27} :\equiv \forall n \in \omega\,(\mathrm{Phi}(n) \to \mathrm{Phi}(S(n))) \to \forall x\,(x \in B \leftrightarrow x \in \omega \land \mathrm{Phi}(x)) \to y \in B \to y \in B$
\item $\varphi_{28} :\equiv \forall n \in \omega\,(\mathrm{Phi}(n) \to \mathrm{Phi}(S(n))) \to \forall x\,(x \in B \leftrightarrow x \in \omega \land \mathrm{Phi}(x)) \to y \in B \to S(y) \in B$
\item $\varphi_{29} :\equiv \forall n \in \omega\,(\mathrm{Phi}(n) \to \mathrm{Phi}(S(n))) \to \forall x\,(x \in B \leftrightarrow x \in \omega \land \mathrm{Phi}(x)) \to n \in \omega \to \mathrm{Phi}(n)$
\item $\varphi_{30} :\equiv \forall n \in \omega\,(\mathrm{Phi}(n) \to \mathrm{Phi}(S(n))) \to \forall x\,(x \in B \leftrightarrow x \in \omega \land \mathrm{Phi}(x)) \to n \in \omega \to n \in \omega$
\item $\varphi_{31} :\equiv \forall x\,(\mathrm{Ind}(x) \leftrightarrow \emptyset \in x \land \forall y \in x\,(S(y) \in x)) \to (\mathrm{Ind}(B) \leftrightarrow \emptyset \in B \land \forall y \in B\,(S(y) \in B))$
\item $\varphi_{32} :\equiv \mathrm{Phi}(\emptyset) \to \forall n \in \omega\,(\mathrm{Phi}(n) \to \mathrm{Phi}(S(n))) \to \varphi_{1}$
\item $\varphi_{33} :\equiv \mathrm{Phi}(\emptyset) \to \forall n \in \omega\,(\mathrm{Phi}(n) \to \mathrm{Phi}(S(n))) \to \varphi_{2}$
\item $\varphi_{34} :\equiv \mathrm{Phi}(\emptyset) \to \forall n \in \omega\,(\mathrm{Phi}(n) \to \mathrm{Phi}(S(n))) \to \varphi_{4}$
\item $\varphi_{35} :\equiv \forall x\,(\mathrm{Ind}(x) \leftrightarrow \emptyset \in x \land \forall y \in x\,(S(y) \in x)) \to (\mathrm{Ind}(\omega) \leftrightarrow \emptyset \in \omega \land \forall y \in \omega\,(S(y) \in \omega))$
\item $\varphi_{36} :\equiv \mathrm{Phi}(\emptyset) \to \forall n \in \omega\,(\mathrm{Phi}(n) \to \mathrm{Phi}(S(n))) \to \varphi_{5}$
\item $\varphi_{37} :\equiv \mathrm{Phi}(\emptyset) \to \forall n \in \omega\,(\mathrm{Phi}(n) \to \mathrm{Phi}(S(n))) \to \varphi_{6}$
\item $\varphi_{38} :\equiv \mathrm{Phi}(\emptyset) \to \forall n \in \omega\,(\mathrm{Phi}(n) \to \mathrm{Phi}(S(n))) \to \varphi_{7}$
\item $\varphi_{39} :\equiv \mathrm{Phi}(\emptyset) \to \forall n \in \omega\,(\mathrm{Phi}(n) \to \mathrm{Phi}(S(n))) \to \varphi_{8}$
\item $\varphi_{40} :\equiv \mathrm{Phi}(\emptyset) \to \forall n \in \omega\,(\mathrm{Phi}(n) \to \mathrm{Phi}(S(n))) \to \varphi_{9}$
\item $\varphi_{41} :\equiv \mathrm{Phi}(\emptyset) \to \forall n \in \omega\,(\mathrm{Phi}(n) \to \mathrm{Phi}(S(n))) \to \varphi_{15}$
\item $\varphi_{42} :\equiv \mathrm{Phi}(\emptyset) \to \forall n \in \omega\,(\mathrm{Phi}(n) \to \mathrm{Phi}(S(n))) \to \varphi_{17}$
\item $\varphi_{43} :\equiv \mathrm{Phi}(\emptyset) \to \forall n \in \omega\,(\mathrm{Phi}(n) \to \mathrm{Phi}(S(n))) \to \varphi_{18}$
\item $\varphi_{44} :\equiv \mathrm{Phi}(\emptyset) \to \forall n \in \omega\,(\mathrm{Phi}(n) \to \mathrm{Phi}(S(n))) \to \varphi_{23}$
\item $\varphi_{45} :\equiv \mathrm{Phi}(\emptyset) \to \forall n \in \omega\,(\mathrm{Phi}(n) \to \mathrm{Phi}(S(n))) \to \varphi_{25}$
\item $\varphi_{46} :\equiv \mathrm{Ind}(\omega) \to (\mathrm{Ind}(\omega) \leftrightarrow \emptyset \in \omega \land \forall y \in \omega\,(S(y) \in \omega)) \to \varphi_{45}$
\item $\varphi_{47} :\equiv \forall n\,(\varphi_{24} \to n \in \omega \to \mathrm{Phi}(n)) \to \forall n\,(\varphi_{24}) \to \forall n \in \omega\,(\mathrm{Phi}(n))$
\item $\varphi_{48} :\equiv \forall y\,(\varphi_{24} \to y \in B \to S(y) \in B) \to \forall y\,(\varphi_{24}) \to \forall y \in B\,(S(y) \in B)$
\item $\varphi_{49} :\equiv \forall B\,(\varphi_{19} \to \neg \forall x\,(x \in B \leftrightarrow x \in \omega \land \mathrm{Phi}(x))) \to \forall B\,(\varphi_{19}) \to \forall B\,(\neg \forall x\,(x \in B \leftrightarrow x \in \omega \land \mathrm{Phi}(x)))$
\item $\varphi_{50} :\equiv (\mathrm{Phi}(\emptyset) \to \varphi_{30}) \to (n \in \omega \leftrightarrow \forall x\,(\mathrm{Ind}(x) \to n \in x)) \to \varphi_{40}$
\item $\varphi_{51} :\equiv \varphi_{40} \to (\forall x\,(\mathrm{Ind}(x) \to n \in x) \to \mathrm{Ind}(B) \to n \in B) \to \varphi_{38}$
\item $\varphi_{52} :\equiv (\mathrm{Ind}(B) \leftrightarrow \emptyset \in B \land \forall y \in B\,(S(y) \in B)) \to (\mathrm{Phi}(\emptyset) \to \varphi_{26}) \to \varphi_{32} \to \mathrm{Phi}(\emptyset) \to \varphi_{22}$
\item $\varphi_{53} :\equiv (\varphi_{19} \to \forall B\,(\varphi_{19})) \to \forall y\,(\varphi_{13}) \to (\forall y\,(\varphi_{13}) \to \varphi_{21}) \to \varphi_{20} \to \mathrm{Phi}(\emptyset) \to \varphi_{3}$
\item $\varphi_{54} :\equiv (\mathrm{Phi}(\emptyset) \to \varphi_{16}) \to \forall B\,(\varphi_{19} \to \neg \forall x\,(x \in B \leftrightarrow x \in \omega \land \mathrm{Phi}(x))) \to \varphi_{49} \to \varphi_{53}$
\end{abbrevlist}
\end{fullproof}

\section{順序の一般的性質 (< は ⊊ として定義されているので, ℝ に限らずすべての集合で成り立つ)}
% 自動生成 (prover/build.js): LtIrrefl  ノード数 13
\begin{fullproof}{[O1]}{$\forall x\,(\neg x < x)$}{13}\phantomsection\label{tree:LtIrrefl}
\noindent$\Pi_{1}$:
\begin{ptree}
\AXC{$\forall x\,\forall y\,(x < y \leftrightarrow x \subseteq y \land x \neq y)$\lbl{Def$<$}}
\AXC{$\varphi_{1}$\lbl{A2}}
\BIC{$\forall y\,(x < y \leftrightarrow x \subseteq y \land x \neq y)$}
\AXC{$\forall y\,(x < y \leftrightarrow x \subseteq y \land x \neq y) \to (x < x \leftrightarrow x \subseteq x \land x \neq x)$\lbl{A2}}
\BIC{$x < x \leftrightarrow x \subseteq x \land x \neq x$}
\end{ptree}
\noindent 主木:
\begin{ptree}
\AXC{$x = x$\lbl{A5}}
\UIC{$\forall x\,(x = x)$}
\AXC{$\forall x\,(x = x) \to x = x$\lbl{A2}}
\BIC{$x = x$}
\AXC{$\stackrel{\Pi_{1}}{\vdots}$}\noLine
\UIC{$x < x \leftrightarrow x \subseteq x \land x \neq x$}
\AXC{$(x < x \leftrightarrow x \subseteq x \land x \neq x) \to x = x \to \neg x < x$\lbl{A1}}
\BIC{$x = x \to \neg x < x$}
\BIC{$\neg x < x$}
\UIC{$\forall x\,(\neg x < x)$}
\end{ptree}
\begin{abbrevlist}
\item $\varphi_{1} :\equiv \forall x\,\forall y\,(x < y \leftrightarrow x \subseteq y \land x \neq y) \to \forall y\,(x < y \leftrightarrow x \subseteq y \land x \neq y)$
\end{abbrevlist}
\end{fullproof}

% 自動生成 (prover/build.js): LtTrans  ノード数 173
\begin{fullproof}{[O2]}{$\forall x\,\forall y\,\forall z\,(x < y \to y < z \to x < z)$}{173}\phantomsection\label{tree:LtTrans}
\noindent$\Pi_{1}$:
\begin{ptree}
\AXC{$\forall x\,\forall y\,(x < y \leftrightarrow x \subseteq y \land x \neq y)$\lbl{Def$<$}}
\AXC{$\varphi_{3}$\lbl{A2}}
\BIC{$\forall y\,(x < y \leftrightarrow x \subseteq y \land x \neq y)$}
\AXC{$\forall y\,(x < y \leftrightarrow x \subseteq y \land x \neq y) \to (x < z \leftrightarrow x \subseteq z \land x \neq z)$\lbl{A2}}
\BIC{$x < z \leftrightarrow x \subseteq z \land x \neq z$}
\end{ptree}
\noindent$\Pi_{2}$:
\begin{ptree}
\AXC{$\forall x\,\forall y\,(x \subseteq y \leftrightarrow \forall u \in x\,(u \in y))$\lbl{Def$\subseteq$}}
\AXC{$\varphi_{18}$\lbl{A2}}
\BIC{$\forall y\,(x \subseteq y \leftrightarrow \forall u \in x\,(u \in y))$}
\AXC{$\varphi_{8}$\lbl{A2}}
\BIC{$x \subseteq z \leftrightarrow \forall u \in x\,(u \in z)$}
\end{ptree}
\noindent$\Pi_{3}$:
\begin{ptree}
\AXC{$\forall x\,\forall y\,(x < y \leftrightarrow x \subseteq y \land x \neq y)$\lbl{Def$<$}}
\AXC{$\varphi_{4}$\lbl{A2}}
\BIC{$\forall y\,(v_{3} < y \leftrightarrow v_{3} \subseteq y \land v_{3} \neq y)$}
\AXC{$\varphi_{5}$\lbl{A2}}
\BIC{$v_{3} < v_{4} \leftrightarrow v_{3} \subseteq v_{4} \land v_{3} \neq v_{4}$}
\UIC{$\forall v_{4}\,(v_{3} < v_{4} \leftrightarrow v_{3} \subseteq v_{4} \land v_{3} \neq v_{4})$}
\UIC{$\forall v_{3}\,\forall v_{4}\,(v_{3} < v_{4} \leftrightarrow v_{3} \subseteq v_{4} \land v_{3} \neq v_{4})$}
\AXC{$\varphi_{13}$\lbl{A2}}
\BIC{$\forall v_{4}\,(y < v_{4} \leftrightarrow y \subseteq v_{4} \land y \neq v_{4})$}
\AXC{$\varphi_{1}$\lbl{A2}}
\BIC{$y < z \leftrightarrow y \subseteq z \land y \neq z$}
\end{ptree}
\noindent$\Pi_{4}$:
\begin{ptree}
\AXC{$\stackrel{\Pi_{3}}{\vdots}$}\noLine
\UIC{$y < z \leftrightarrow y \subseteq z \land y \neq z$}
\AXC{$x < y \to y < z \to y < z$\lbl{A1}}
\AXC{$(x < y \to y < z \to y < z) \to (y < z \leftrightarrow y \subseteq z \land y \neq z) \to \varphi_{14}$\lbl{A1}}
\BIC{$(y < z \leftrightarrow y \subseteq z \land y \neq z) \to \varphi_{14}$}
\BIC{$\varphi_{14}$}
\end{ptree}
\noindent$\Pi_{5}$:
\begin{ptree}
\AXC{$\forall x\,\forall y\,(x \subseteq y \leftrightarrow \forall u \in x\,(u \in y))$\lbl{Def$\subseteq$}}
\AXC{$\varphi_{21}$\lbl{A2}}
\BIC{$\forall y\,(v_{5} \subseteq y \leftrightarrow \forall u \in v_{5}\,(u \in y))$}
\AXC{$\varphi_{19}$\lbl{A2}}
\BIC{$v_{5} \subseteq v_{6} \leftrightarrow \forall u \in v_{5}\,(u \in v_{6})$}
\UIC{$\forall v_{6}\,(v_{5} \subseteq v_{6} \leftrightarrow \forall u \in v_{5}\,(u \in v_{6}))$}
\UIC{$\forall v_{5}\,\forall v_{6}\,(v_{5} \subseteq v_{6} \leftrightarrow \forall u \in v_{5}\,(u \in v_{6}))$}
\AXC{$\varphi_{29}$\lbl{A2}}
\BIC{$\forall v_{6}\,(y \subseteq v_{6} \leftrightarrow \forall u \in y\,(u \in v_{6}))$}
\AXC{$\varphi_{11}$\lbl{A2}}
\BIC{$y \subseteq z \leftrightarrow \forall u \in y\,(u \in z)$}
\AXC{$\stackrel{\Pi_{4}}{\vdots}$}\noLine
\UIC{$\varphi_{14}$}
\BIC{$x < y \to y < z \to \forall u \in y\,(u \in z)$}
\end{ptree}
\noindent$\Pi_{6}$:
\begin{ptree}
\AXC{$\forall u \in y\,(u \in z) \to u \in y \to u \in z$\lbl{A2}}
\AXC{$\stackrel{\Pi_{5}}{\vdots}$}\noLine
\UIC{$x < y \to y < z \to \forall u \in y\,(u \in z)$}
\AXC{$(x < y \to y < z \to \forall u \in y\,(u \in z)) \to \varphi_{24}$\lbl{A1}}
\BIC{$\varphi_{24}$}
\BIC{$x < y \to y < z \to u \in y \to u \in z$}
\end{ptree}
\noindent$\Pi_{7}$:
\begin{ptree}
\AXC{$\forall x\,\forall y\,(x < y \leftrightarrow x \subseteq y \land x \neq y)$\lbl{Def$<$}}
\AXC{$\varphi_{3}$\lbl{A2}}
\BIC{$\forall y\,(x < y \leftrightarrow x \subseteq y \land x \neq y)$}
\AXC{$\forall y\,(x < y \leftrightarrow x \subseteq y \land x \neq y) \to (x < y \leftrightarrow x \subseteq y \land x \neq y)$\lbl{A2}}
\BIC{$x < y \leftrightarrow x \subseteq y \land x \neq y$}
\end{ptree}
\noindent$\Pi_{8}$:
\begin{ptree}
\AXC{$\stackrel{\Pi_{7}}{\vdots}$}\noLine
\UIC{$x < y \leftrightarrow x \subseteq y \land x \neq y$}
\AXC{$x < y \to x < y$\lbl{A1}}
\AXC{$(x < y \to x < y) \to (x < y \leftrightarrow x \subseteq y \land x \neq y) \to \varphi_{15}$\lbl{A1}}
\BIC{$(x < y \leftrightarrow x \subseteq y \land x \neq y) \to \varphi_{15}$}
\BIC{$\varphi_{15}$}
\end{ptree}
\noindent$\Pi_{9}$:
\begin{ptree}
\AXC{$\forall x\,\forall y\,(x \subseteq y \leftrightarrow \forall u \in x\,(u \in y))$\lbl{Def$\subseteq$}}
\AXC{$\varphi_{18}$\lbl{A2}}
\BIC{$\forall y\,(x \subseteq y \leftrightarrow \forall u \in x\,(u \in y))$}
\AXC{$\varphi_{9}$\lbl{A2}}
\BIC{$x \subseteq y \leftrightarrow \forall u \in x\,(u \in y)$}
\AXC{$\stackrel{\Pi_{8}}{\vdots}$}\noLine
\UIC{$\varphi_{15}$}
\BIC{$x < y \to y < z \to \forall u \in x\,(u \in y)$}
\end{ptree}
\noindent$\Pi_{10}$:
\begin{ptree}
\AXC{$\forall u \in x\,(u \in y) \to u \in x \to u \in y$\lbl{A2}}
\AXC{$\stackrel{\Pi_{9}}{\vdots}$}\noLine
\UIC{$x < y \to y < z \to \forall u \in x\,(u \in y)$}
\AXC{$(x < y \to y < z \to \forall u \in x\,(u \in y)) \to \varphi_{25}$\lbl{A1}}
\BIC{$\varphi_{25}$}
\BIC{$x < y \to y < z \to u \in x \to u \in y$}
\end{ptree}
\noindent$\Pi_{11}$:
\begin{ptree}
\AXC{$\stackrel{\Pi_{6}}{\vdots}$}\noLine
\UIC{$x < y \to y < z \to u \in y \to u \in z$}
\AXC{$\stackrel{\Pi_{10}}{\vdots}$}\noLine
\UIC{$x < y \to y < z \to u \in x \to u \in y$}
\AXC{$(x < y \to y < z \to u \in x \to u \in y) \to \varphi_{16}$\lbl{A1}}
\BIC{$\varphi_{16}$}
\BIC{$x < y \to y < z \to u \in x \to u \in z$}
\AXC{$\varphi_{17}$\lbl{A1}}
\BIC{$x < y \land y < z \to u \in x \to u \in z$}
\UIC{$\forall u\,(x < y \land y < z \to u \in x \to u \in z)$}
\end{ptree}
\noindent$\Pi_{12}$:
\begin{ptree}
\AXC{$\varphi_{28}$\lbl{A3}}
\AXC{$\stackrel{\Pi_{11}}{\vdots}$}\noLine
\UIC{$\forall u\,(x < y \land y < z \to u \in x \to u \in z)$}
\AXC{$\forall u\,(x < y \land y < z \to u \in x \to u \in z) \to \varphi_{28} \to \varphi_{23}$\lbl{A1}}
\BIC{$\varphi_{28} \to \varphi_{23}$}
\BIC{$\varphi_{23}$}
\end{ptree}
\noindent$\Pi_{13}$:
\begin{ptree}
\AXC{$\stackrel{\Pi_{2}}{\vdots}$}\noLine
\UIC{$x \subseteq z \leftrightarrow \forall u \in x\,(u \in z)$}
\AXC{$x < y \land y < z \to \forall u\,(x < y \land y < z)$\lbl{A4}}
\AXC{$\stackrel{\Pi_{12}}{\vdots}$}\noLine
\UIC{$\varphi_{23}$}
\BIC{$x < y \to y < z \to \forall u \in x\,(u \in z)$}
\AXC{$\varphi_{46}$\lbl{A1}}
\BIC{$(x \subseteq z \leftrightarrow \forall u \in x\,(u \in z)) \to x < y \to y < z \to x \subseteq z$}
\BIC{$x < y \to y < z \to x \subseteq z$}
\end{ptree}
\noindent$\Pi_{14}$:
\begin{ptree}
\AXC{$\varphi_{40}$\lbl{A4}}
\AXC{$\forall z\,(\varphi_{32}) \to \varphi_{42}$\lbl{A3}}
\AXC{$\varphi_{32}$\lbl{A2}}
\UIC{$\forall z\,(\varphi_{32})$}
\AXC{$\forall z\,(\varphi_{32}) \to (\forall z\,(\varphi_{32}) \to \varphi_{42}) \to \varphi_{40} \to \varphi_{38}$\lbl{A1}}
\BIC{$(\forall z\,(\varphi_{32}) \to \varphi_{42}) \to \varphi_{40} \to \varphi_{38}$}
\BIC{$\varphi_{40} \to \varphi_{38}$}
\BIC{$\varphi_{38}$}
\end{ptree}
\noindent$\Pi_{15}$:
\begin{ptree}
\AXC{$\forall x\,\forall y\,(x \subseteq y \leftrightarrow \forall u \in x\,(u \in y))$\lbl{Def$\subseteq$}}
\AXC{$\varphi_{22}$\lbl{A2}}
\BIC{$\forall y\,(v_{7} \subseteq y \leftrightarrow \forall u \in v_{7}\,(u \in y))$}
\AXC{$\varphi_{20}$\lbl{A2}}
\BIC{$v_{7} \subseteq v_{8} \leftrightarrow \forall u \in v_{7}\,(u \in v_{8})$}
\UIC{$\forall v_{8}\,(v_{7} \subseteq v_{8} \leftrightarrow \forall u \in v_{7}\,(u \in v_{8}))$}
\UIC{$\forall v_{7}\,\forall v_{8}\,(v_{7} \subseteq v_{8} \leftrightarrow \forall u \in v_{7}\,(u \in v_{8}))$}
\AXC{$\varphi_{30}$\lbl{A2}}
\BIC{$\forall v_{8}\,(y \subseteq v_{8} \leftrightarrow \forall u \in y\,(u \in v_{8}))$}
\AXC{$\varphi_{12}$\lbl{A2}}
\BIC{$y \subseteq x \leftrightarrow \forall u \in y\,(u \in x)$}
\end{ptree}
\noindent$\Pi_{16}$:
\begin{ptree}
\AXC{$z_{1} = w \to y \subseteq z_{1} \to y \subseteq w$\lbl{A6}}
\UIC{$\forall w\,(z_{1} = w \to y \subseteq z_{1} \to y \subseteq w)$}
\UIC{$\forall z_{1}\,\forall w\,(z_{1} = w \to y \subseteq z_{1} \to y \subseteq w)$}
\AXC{$\varphi_{7}$\lbl{A2}}
\BIC{$\forall w\,(z = w \to y \subseteq z \to y \subseteq w)$}
\AXC{$\forall w\,(z = w \to y \subseteq z \to y \subseteq w) \to z = x \to y \subseteq z \to y \subseteq x$\lbl{A2}}
\BIC{$z = x \to y \subseteq z \to y \subseteq x$}
\end{ptree}
\noindent$\Pi_{17}$:
\begin{ptree}
\AXC{$x = x$\lbl{A5}}
\AXC{$x = y \to x = x \to y = x$\lbl{A6}}
\AXC{$(x = y \to x = x \to y = x) \to x = x \to x = y \to y = x$\lbl{A1}}
\BIC{$x = x \to x = y \to y = x$}
\BIC{$x = y \to y = x$}
\UIC{$\forall y\,(x = y \to y = x)$}
\UIC{$\forall x\,\forall y\,(x = y \to y = x)$}
\end{ptree}
\noindent$\Pi_{18}$:
\begin{ptree}
\AXC{$\stackrel{\Pi_{17}}{\vdots}$}\noLine
\UIC{$\forall x\,\forall y\,(x = y \to y = x)$}
\AXC{$\forall x\,\forall y\,(x = y \to y = x) \to \forall y\,(x = y \to y = x)$\lbl{A2}}
\BIC{$\forall y\,(x = y \to y = x)$}
\AXC{$\forall y\,(x = y \to y = x) \to x = z \to z = x$\lbl{A2}}
\BIC{$x = z \to z = x$}
\end{ptree}
\noindent$\Pi_{19}$:
\begin{ptree}
\AXC{$\stackrel{\Pi_{18}}{\vdots}$}\noLine
\UIC{$x = z \to z = x$}
\AXC{$x < y \to y < z \to x = z \to x = z$\lbl{A1}}
\AXC{$\varphi_{10}$\lbl{A1}}
\BIC{$(x = z \to z = x) \to x < y \to y < z \to x = z \to z = x$}
\BIC{$x < y \to y < z \to x = z \to z = x$}
\AXC{$(x < y \to y < z \to x = z \to z = x) \to \varphi_{44}$\lbl{A1}}
\BIC{$\varphi_{44}$}
\end{ptree}
\noindent$\Pi_{20}$:
\begin{ptree}
\AXC{$\stackrel{\Pi_{16}}{\vdots}$}\noLine
\UIC{$z = x \to y \subseteq z \to y \subseteq x$}
\AXC{$\stackrel{\Pi_{3}}{\vdots}$}\noLine
\UIC{$y < z \leftrightarrow y \subseteq z \land y \neq z$}
\AXC{$x < y \to y < z \to y < z$\lbl{A1}}
\AXC{$\varphi_{36}$\lbl{A1}}
\BIC{$(y < z \leftrightarrow y \subseteq z \land y \neq z) \to x < y \to y < z \to x = z \to y \subseteq z$}
\BIC{$x < y \to y < z \to x = z \to y \subseteq z$}
\AXC{$\stackrel{\Pi_{19}}{\vdots}$}\noLine
\UIC{$\varphi_{44}$}
\BIC{$(z = x \to y \subseteq z \to y \subseteq x) \to x < y \to y < z \to x = z \to y \subseteq x$}
\BIC{$x < y \to y < z \to x = z \to y \subseteq x$}
\end{ptree}
\noindent$\Pi_{21}$:
\begin{ptree}
\AXC{$\stackrel{\Pi_{15}}{\vdots}$}\noLine
\UIC{$y \subseteq x \leftrightarrow \forall u \in y\,(u \in x)$}
\AXC{$\stackrel{\Pi_{20}}{\vdots}$}\noLine
\UIC{$x < y \to y < z \to x = z \to y \subseteq x$}
\AXC{$(x < y \to y < z \to x = z \to y \subseteq x) \to \varphi_{26}$\lbl{A1}}
\BIC{$\varphi_{26}$}
\BIC{$x < y \to y < z \to x = z \to \forall u \in y\,(u \in x)$}
\end{ptree}
\noindent$\Pi_{22}$:
\begin{ptree}
\AXC{$\stackrel{\Pi_{21}}{\vdots}$}\noLine
\UIC{$x < y \to y < z \to x = z \to \forall u \in y\,(u \in x)$}
\AXC{$(x < y \to y < z \to x = z \to \forall u \in y\,(u \in x)) \to \varphi_{33}$\lbl{A1}}
\BIC{$\varphi_{33}$}
\end{ptree}
\noindent$\Pi_{23}$:
\begin{ptree}
\AXC{$\forall u \in x\,(u \in y) \to u \in x \to u \in y$\lbl{A2}}
\AXC{$\stackrel{\Pi_{9}}{\vdots}$}\noLine
\UIC{$x < y \to y < z \to \forall u \in x\,(u \in y)$}
\AXC{$(x < y \to y < z \to \forall u \in x\,(u \in y)) \to \varphi_{25}$\lbl{A1}}
\BIC{$\varphi_{25}$}
\BIC{$x < y \to y < z \to u \in x \to u \in y$}
\end{ptree}
\noindent$\Pi_{24}$:
\begin{ptree}
\AXC{$\forall u \in y\,(u \in x) \to u \in y \to u \in x$\lbl{A2}}
\AXC{$\stackrel{\Pi_{22}}{\vdots}$}\noLine
\UIC{$\varphi_{33}$}
\BIC{$x < y \to y < z \to x = z \to u \in y \to u \in x$}
\AXC{$\stackrel{\Pi_{23}}{\vdots}$}\noLine
\UIC{$x < y \to y < z \to u \in x \to u \in y$}
\AXC{$(x < y \to y < z \to u \in x \to u \in y) \to \varphi_{34}$\lbl{A1}}
\BIC{$\varphi_{34}$}
\BIC{$x < y \to y < z \to x = z \to (u \in x \leftrightarrow u \in y)$}
\end{ptree}
\noindent$\Pi_{25}$:
\begin{ptree}
\AXC{$\varphi_{43}$\lbl{A3}}
\AXC{$\stackrel{\Pi_{24}}{\vdots}$}\noLine
\UIC{$x < y \to y < z \to x = z \to (u \in x \leftrightarrow u \in y)$}
\AXC{$\varphi_{35}$\lbl{A1}}
\BIC{$x < y \land y < z \land x = z \to (u \in x \leftrightarrow u \in y)$}
\UIC{$\forall u\,(x < y \land y < z \land x = z \to (u \in x \leftrightarrow u \in y))$}
\AXC{$\forall u\,(x < y \land y < z \land x = z \to (u \in x \leftrightarrow u \in y)) \to \varphi_{43} \to \varphi_{45}$\lbl{A1}}
\BIC{$\varphi_{43} \to \varphi_{45}$}
\BIC{$\varphi_{45}$}
\end{ptree}
\noindent$\Pi_{26}$:
\begin{ptree}
\AXC{$x < y \land y < z \land x = z \to \forall u\,(x < y \land y < z \land x = z)$\lbl{A4}}
\AXC{$\stackrel{\Pi_{25}}{\vdots}$}\noLine
\UIC{$\varphi_{45}$}
\BIC{$x < y \to y < z \to x = z \to \forall u\,(u \in x \leftrightarrow u \in y)$}
\AXC{$(x < y \to y < z \to x = z \to \forall u\,(u \in x \leftrightarrow u \in y)) \to \varphi_{48}$\lbl{A1}}
\BIC{$\varphi_{48}$}
\end{ptree}
\noindent$\Pi_{27}$:
\begin{ptree}
\AXC{$\stackrel{\Pi_{14}}{\vdots}$}\noLine
\UIC{$\varphi_{38}$}
\AXC{$\forall x\,\forall y\,(\forall z\,(z \in x \leftrightarrow z \in y) \to x = y)$\lbl{Ext}}
\AXC{$\varphi_{6}$\lbl{A2}}
\BIC{$\forall y\,(\forall z\,(z \in x \leftrightarrow z \in y) \to x = y)$}
\AXC{$\varphi_{2}$\lbl{A2}}
\BIC{$\forall z\,(z \in x \leftrightarrow z \in y) \to x = y$}
\AXC{$\stackrel{\Pi_{26}}{\vdots}$}\noLine
\UIC{$\varphi_{48}$}
\BIC{$\varphi_{38} \to \varphi_{37} \to x < y \to y < z \to x = z \to x = y$}
\BIC{$\varphi_{37} \to x < y \to y < z \to x = z \to x = y$}
\end{ptree}
\noindent$\Pi_{28}$:
\begin{ptree}
\AXC{$\varphi_{39}$\lbl{A4}}
\AXC{$\forall u\,(\varphi_{31}) \to \varphi_{41}$\lbl{A3}}
\AXC{$\varphi_{31}$\lbl{A2}}
\UIC{$\forall u\,(\varphi_{31})$}
\AXC{$\forall u\,(\varphi_{31}) \to (\forall u\,(\varphi_{31}) \to \varphi_{41}) \to \varphi_{39} \to \varphi_{37}$\lbl{A1}}
\BIC{$(\forall u\,(\varphi_{31}) \to \varphi_{41}) \to \varphi_{39} \to \varphi_{37}$}
\BIC{$\varphi_{39} \to \varphi_{37}$}
\BIC{$\varphi_{37}$}
\AXC{$\stackrel{\Pi_{27}}{\vdots}$}\noLine
\UIC{$\varphi_{37} \to x < y \to y < z \to x = z \to x = y$}
\BIC{$x < y \to y < z \to x = z \to x = y$}
\end{ptree}
\noindent$\Pi_{29}$:
\begin{ptree}
\AXC{$\stackrel{\Pi_{13}}{\vdots}$}\noLine
\UIC{$x < y \to y < z \to x \subseteq z$}
\AXC{$\stackrel{\Pi_{7}}{\vdots}$}\noLine
\UIC{$x < y \leftrightarrow x \subseteq y \land x \neq y$}
\AXC{$x < y \to x < y$\lbl{A1}}
\AXC{$\stackrel{\Pi_{28}}{\vdots}$}\noLine
\UIC{$x < y \to y < z \to x = z \to x = y$}
\AXC{$\varphi_{47}$\lbl{A1}}
\BIC{$(x < y \to x < y) \to (x < y \leftrightarrow x \subseteq y \land x \neq y) \to \varphi_{27}$}
\BIC{$(x < y \leftrightarrow x \subseteq y \land x \neq y) \to \varphi_{27}$}
\BIC{$\varphi_{27}$}
\BIC{$(x < z \leftrightarrow x \subseteq z \land x \neq z) \to x < y \to y < z \to x < z$}
\end{ptree}
\noindent 主木:
\begin{ptree}
\AXC{$\stackrel{\Pi_{1}}{\vdots}$}\noLine
\UIC{$x < z \leftrightarrow x \subseteq z \land x \neq z$}
\AXC{$\stackrel{\Pi_{29}}{\vdots}$}\noLine
\UIC{$(x < z \leftrightarrow x \subseteq z \land x \neq z) \to x < y \to y < z \to x < z$}
\BIC{$x < y \to y < z \to x < z$}
\UIC{$\forall z\,(x < y \to y < z \to x < z)$}
\UIC{$\forall y\,\forall z\,(x < y \to y < z \to x < z)$}
\UIC{$\forall x\,\forall y\,\forall z\,(x < y \to y < z \to x < z)$}
\end{ptree}
\begin{abbrevlist}
\item $\varphi_{1} :\equiv \forall v_{4}\,(y < v_{4} \leftrightarrow y \subseteq v_{4} \land y \neq v_{4}) \to (y < z \leftrightarrow y \subseteq z \land y \neq z)$
\item $\varphi_{2} :\equiv \forall y\,(\forall z\,(z \in x \leftrightarrow z \in y) \to x = y) \to \forall z\,(z \in x \leftrightarrow z \in y) \to x = y$
\item $\varphi_{3} :\equiv \forall x\,\forall y\,(x < y \leftrightarrow x \subseteq y \land x \neq y) \to \forall y\,(x < y \leftrightarrow x \subseteq y \land x \neq y)$
\item $\varphi_{4} :\equiv \forall x\,\forall y\,(x < y \leftrightarrow x \subseteq y \land x \neq y) \to \forall y\,(v_{3} < y \leftrightarrow v_{3} \subseteq y \land v_{3} \neq y)$
\item $\varphi_{5} :\equiv \forall y\,(v_{3} < y \leftrightarrow v_{3} \subseteq y \land v_{3} \neq y) \to (v_{3} < v_{4} \leftrightarrow v_{3} \subseteq v_{4} \land v_{3} \neq v_{4})$
\item $\varphi_{6} :\equiv \forall x\,\forall y\,(\forall z\,(z \in x \leftrightarrow z \in y) \to x = y) \to \forall y\,(\forall z\,(z \in x \leftrightarrow z \in y) \to x = y)$
\item $\varphi_{7} :\equiv \forall z_{1}\,\forall w\,(z_{1} = w \to y \subseteq z_{1} \to y \subseteq w) \to \forall w\,(z = w \to y \subseteq z \to y \subseteq w)$
\item $\varphi_{8} :\equiv \forall y\,(x \subseteq y \leftrightarrow \forall u \in x\,(u \in y)) \to (x \subseteq z \leftrightarrow \forall u \in x\,(u \in z))$
\item $\varphi_{9} :\equiv \forall y\,(x \subseteq y \leftrightarrow \forall u \in x\,(u \in y)) \to (x \subseteq y \leftrightarrow \forall u \in x\,(u \in y))$
\item $\varphi_{10} :\equiv (x < y \to y < z \to x = z \to x = z) \to (x = z \to z = x) \to x < y \to y < z \to x = z \to z = x$
\item $\varphi_{11} :\equiv \forall v_{6}\,(y \subseteq v_{6} \leftrightarrow \forall u \in y\,(u \in v_{6})) \to (y \subseteq z \leftrightarrow \forall u \in y\,(u \in z))$
\item $\varphi_{12} :\equiv \forall v_{8}\,(y \subseteq v_{8} \leftrightarrow \forall u \in y\,(u \in v_{8})) \to (y \subseteq x \leftrightarrow \forall u \in y\,(u \in x))$
\item $\varphi_{13} :\equiv \forall v_{3}\,\forall v_{4}\,(v_{3} < v_{4} \leftrightarrow v_{3} \subseteq v_{4} \land v_{3} \neq v_{4}) \to \forall v_{4}\,(y < v_{4} \leftrightarrow y \subseteq v_{4} \land y \neq v_{4})$
\item $\varphi_{14} :\equiv (y \subseteq z \leftrightarrow \forall u \in y\,(u \in z)) \to x < y \to y < z \to \forall u \in y\,(u \in z)$
\item $\varphi_{15} :\equiv (x \subseteq y \leftrightarrow \forall u \in x\,(u \in y)) \to x < y \to y < z \to \forall u \in x\,(u \in y)$
\item $\varphi_{16} :\equiv (x < y \to y < z \to u \in y \to u \in z) \to x < y \to y < z \to u \in x \to u \in z$
\item $\varphi_{17} :\equiv (x < y \to y < z \to u \in x \to u \in z) \to x < y \land y < z \to u \in x \to u \in z$
\item $\varphi_{18} :\equiv \forall x\,\forall y\,(x \subseteq y \leftrightarrow \forall u \in x\,(u \in y)) \to \forall y\,(x \subseteq y \leftrightarrow \forall u \in x\,(u \in y))$
\item $\varphi_{19} :\equiv \forall y\,(v_{5} \subseteq y \leftrightarrow \forall u \in v_{5}\,(u \in y)) \to (v_{5} \subseteq v_{6} \leftrightarrow \forall u \in v_{5}\,(u \in v_{6}))$
\item $\varphi_{20} :\equiv \forall y\,(v_{7} \subseteq y \leftrightarrow \forall u \in v_{7}\,(u \in y)) \to (v_{7} \subseteq v_{8} \leftrightarrow \forall u \in v_{7}\,(u \in v_{8}))$
\item $\varphi_{21} :\equiv \forall x\,\forall y\,(x \subseteq y \leftrightarrow \forall u \in x\,(u \in y)) \to \forall y\,(v_{5} \subseteq y \leftrightarrow \forall u \in v_{5}\,(u \in y))$
\item $\varphi_{22} :\equiv \forall x\,\forall y\,(x \subseteq y \leftrightarrow \forall u \in x\,(u \in y)) \to \forall y\,(v_{7} \subseteq y \leftrightarrow \forall u \in v_{7}\,(u \in y))$
\item $\varphi_{23} :\equiv (x < y \land y < z \to \forall u\,(x < y \land y < z)) \to x < y \to y < z \to \forall u \in x\,(u \in z)$
\item $\varphi_{24} :\equiv (\forall u \in y\,(u \in z) \to u \in y \to u \in z) \to x < y \to y < z \to u \in y \to u \in z$
\item $\varphi_{25} :\equiv (\forall u \in x\,(u \in y) \to u \in x \to u \in y) \to x < y \to y < z \to u \in x \to u \in y$
\item $\varphi_{26} :\equiv (y \subseteq x \leftrightarrow \forall u \in y\,(u \in x)) \to x < y \to y < z \to x = z \to \forall u \in y\,(u \in x)$
\item $\varphi_{27} :\equiv (x < y \to y < z \to x \subseteq z) \to (x < z \leftrightarrow x \subseteq z \land x \neq z) \to x < y \to y < z \to x < z$
\item $\varphi_{28} :\equiv \forall u\,(x < y \land y < z \to u \in x \to u \in z) \to \forall u\,(x < y \land y < z) \to \forall u \in x\,(u \in z)$
\item $\varphi_{29} :\equiv \forall v_{5}\,\forall v_{6}\,(v_{5} \subseteq v_{6} \leftrightarrow \forall u \in v_{5}\,(u \in v_{6})) \to \forall v_{6}\,(y \subseteq v_{6} \leftrightarrow \forall u \in y\,(u \in v_{6}))$
\item $\varphi_{30} :\equiv \forall v_{7}\,\forall v_{8}\,(v_{7} \subseteq v_{8} \leftrightarrow \forall u \in v_{7}\,(u \in v_{8})) \to \forall v_{8}\,(y \subseteq v_{8} \leftrightarrow \forall u \in y\,(u \in v_{8}))$
\item $\varphi_{31} :\equiv \forall z\,(\neg ((z \in x \to z \in y) \to \neg (z \in y \to z \in x))) \to \neg ((u \in x \to u \in y) \to \neg (u \in y \to u \in x))$
\item $\varphi_{32} :\equiv \forall u\,(\neg ((u \in x \to u \in y) \to \neg (u \in y \to u \in x))) \to \neg ((z \in x \to z \in y) \to \neg (z \in y \to z \in x))$
\item $\varphi_{33} :\equiv (\forall u \in y\,(u \in x) \to u \in y \to u \in x) \to x < y \to y < z \to x = z \to u \in y \to u \in x$
\item $\varphi_{34} :\equiv (x < y \to y < z \to x = z \to u \in y \to u \in x) \to x < y \to y < z \to x = z \to (u \in x \leftrightarrow u \in y)$
\item $\varphi_{35} :\equiv (x < y \to y < z \to x = z \to (u \in x \leftrightarrow u \in y)) \to x < y \land y < z \land x = z \to (u \in x \leftrightarrow u \in y)$
\item $\varphi_{36} :\equiv (x < y \to y < z \to y < z) \to (y < z \leftrightarrow y \subseteq z \land y \neq z) \to x < y \to y < z \to x = z \to y \subseteq z$
\item $\varphi_{37} :\equiv \forall z\,(\neg ((z \in x \to z \in y) \to \neg (z \in y \to z \in x))) \to \forall u\,(\neg ((u \in x \to u \in y) \to \neg (u \in y \to u \in x)))$
\item $\varphi_{38} :\equiv \forall u\,(\neg ((u \in x \to u \in y) \to \neg (u \in y \to u \in x))) \to \forall z\,(\neg ((z \in x \to z \in y) \to \neg (z \in y \to z \in x)))$
\item $\varphi_{39} :\equiv \forall z\,(\neg ((z \in x \to z \in y) \to \neg (z \in y \to z \in x))) \to \forall u\,\forall z\,(\neg ((z \in x \to z \in y) \to \neg (z \in y \to z \in x)))$
\item $\varphi_{40} :\equiv \forall u\,(\neg ((u \in x \to u \in y) \to \neg (u \in y \to u \in x))) \to \forall z\,\forall u\,(\neg ((u \in x \to u \in y) \to \neg (u \in y \to u \in x)))$
\item $\varphi_{41} :\equiv \forall u\,\forall z\,(\neg ((z \in x \to z \in y) \to \neg (z \in y \to z \in x))) \to \forall u\,(\neg ((u \in x \to u \in y) \to \neg (u \in y \to u \in x)))$
\item $\varphi_{42} :\equiv \forall z\,\forall u\,(\neg ((u \in x \to u \in y) \to \neg (u \in y \to u \in x))) \to \forall z\,(\neg ((z \in x \to z \in y) \to \neg (z \in y \to z \in x)))$
\item $\varphi_{43} :\equiv \forall u\,(x < y \land y < z \land x = z \to (u \in x \leftrightarrow u \in y)) \to \forall u\,(x < y \land y < z \land x = z) \to \forall u\,(u \in x \leftrightarrow u \in y)$
\item $\varphi_{44} :\equiv (x < y \to y < z \to x = z \to y \subseteq z) \to (z = x \to y \subseteq z \to y \subseteq x) \to x < y \to y < z \to x = z \to y \subseteq x$
\item $\varphi_{45} :\equiv (x < y \land y < z \land x = z \to \forall u\,(x < y \land y < z \land x = z)) \to x < y \to y < z \to x = z \to \forall u\,(u \in x \leftrightarrow u \in y)$
\item $\varphi_{46} :\equiv (x < y \to y < z \to \forall u \in x\,(u \in z)) \to (x \subseteq z \leftrightarrow \forall u \in x\,(u \in z)) \to x < y \to y < z \to x \subseteq z$
\item $\varphi_{47} :\equiv (x < y \to y < z \to x = z \to x = y) \to (x < y \to x < y) \to (x < y \leftrightarrow x \subseteq y \land x \neq y) \to \varphi_{27}$
\item $\varphi_{48} :\equiv (\forall z\,(z \in x \leftrightarrow z \in y) \to x = y) \to \varphi_{38} \to \varphi_{37} \to x < y \to y < z \to x = z \to x = y$
\end{abbrevlist}
\end{fullproof}

% 自動生成 (prover/build.js): LtLeTrans  ノード数 30
\begin{fullproof}{[O3]}{$\forall x\,\forall y\,\forall z\,(x < y \to y \leq z \to x < z)$}{30}\phantomsection\label{tree:LtLeTrans}
\noindent$\Pi_{1}$:
\begin{ptree}
\AXC{$z_{1} = w \to x < z_{1} \to x < w$\lbl{A6}}
\UIC{$\forall w\,(z_{1} = w \to x < z_{1} \to x < w)$}
\UIC{$\forall z_{1}\,\forall w\,(z_{1} = w \to x < z_{1} \to x < w)$}
\AXC{$\varphi_{1}$\lbl{A2}}
\BIC{$\forall w\,(y = w \to x < y \to x < w)$}
\AXC{$\forall w\,(y = w \to x < y \to x < w) \to y = z \to x < y \to x < z$\lbl{A2}}
\BIC{$y = z \to x < y \to x < z$}
\end{ptree}
\noindent$\Pi_{2}$:
\begin{ptree}
\AXC{$x < y \to y \leq z \to y = z \to x < y$\lbl{A1}}
\AXC{$x < y \to y \leq z \to y = z \to y = z$\lbl{A1}}
\AXC{$(x < y \to y \leq z \to y = z \to y = z) \to \varphi_{5}$\lbl{A1}}
\BIC{$\varphi_{5}$}
\BIC{$(y = z \to x < y \to x < z) \to x < y \to y \leq z \to y = z \to x < z$}
\end{ptree}
\noindent$\Pi_{3}$:
\begin{ptree}
\AXC{$\stackrel{\Pi_{1}}{\vdots}$}\noLine
\UIC{$y = z \to x < y \to x < z$}
\AXC{$\stackrel{\Pi_{2}}{\vdots}$}\noLine
\UIC{$(y = z \to x < y \to x < z) \to x < y \to y \leq z \to y = z \to x < z$}
\BIC{$x < y \to y \leq z \to y = z \to x < z$}
\end{ptree}
\noindent$\Pi_{4}$:
\begin{ptree}
\AXC{$\forall x\,\forall y\,\forall z\,(x < y \to y < z \to x < z)$\lbl{[O2]}}
\AXC{$\varphi_{3}$\lbl{A2}}
\BIC{$\forall y\,\forall z\,(x < y \to y < z \to x < z)$}
\AXC{$\varphi_{2}$\lbl{A2}}
\BIC{$\forall z\,(x < y \to y < z \to x < z)$}
\AXC{$\forall z\,(x < y \to y < z \to x < z) \to x < y \to y < z \to x < z$\lbl{A2}}
\BIC{$x < y \to y < z \to x < z$}
\end{ptree}
\noindent$\Pi_{5}$:
\begin{ptree}
\AXC{$x < y \to y \leq z \to y \leq z$\lbl{A1}}
\AXC{$x < y \to x < y$\lbl{A1}}
\AXC{$(x < y \to x < y) \to (x < y \to y \leq z \to y \leq z) \to \varphi_{4}$\lbl{A1}}
\BIC{$(x < y \to y \leq z \to y \leq z) \to \varphi_{4}$}
\BIC{$\varphi_{4}$}
\end{ptree}
\noindent 主木:
\begin{ptree}
\AXC{$\stackrel{\Pi_{3}}{\vdots}$}\noLine
\UIC{$x < y \to y \leq z \to y = z \to x < z$}
\AXC{$\stackrel{\Pi_{4}}{\vdots}$}\noLine
\UIC{$x < y \to y < z \to x < z$}
\AXC{$\stackrel{\Pi_{5}}{\vdots}$}\noLine
\UIC{$\varphi_{4}$}
\BIC{$(x < y \to y \leq z \to y = z \to x < z) \to x < y \to y \leq z \to x < z$}
\BIC{$x < y \to y \leq z \to x < z$}
\UIC{$\forall z\,(x < y \to y \leq z \to x < z)$}
\UIC{$\forall y\,\forall z\,(x < y \to y \leq z \to x < z)$}
\UIC{$\forall x\,\forall y\,\forall z\,(x < y \to y \leq z \to x < z)$}
\end{ptree}
\begin{abbrevlist}
\item $\varphi_{1} :\equiv \forall z_{1}\,\forall w\,(z_{1} = w \to x < z_{1} \to x < w) \to \forall w\,(y = w \to x < y \to x < w)$
\item $\varphi_{2} :\equiv \forall y\,\forall z\,(x < y \to y < z \to x < z) \to \forall z\,(x < y \to y < z \to x < z)$
\item $\varphi_{3} :\equiv \forall x\,\forall y\,\forall z\,(x < y \to y < z \to x < z) \to \forall y\,\forall z\,(x < y \to y < z \to x < z)$
\item $\varphi_{4} :\equiv (x < y \to y < z \to x < z) \to (x < y \to y \leq z \to y = z \to x < z) \to x < y \to y \leq z \to x < z$
\item $\varphi_{5} :\equiv (x < y \to y \leq z \to y = z \to x < y) \to (y = z \to x < y \to x < z) \to x < y \to y \leq z \to y = z \to x < z$
\end{abbrevlist}
\end{fullproof}

% 自動生成 (prover/build.js): LeLtTrans  ノード数 44
\begin{fullproof}{[O4]}{$\forall x\,\forall y\,\forall z\,(x \leq y \to y < z \to x < z)$}{44}\phantomsection\label{tree:LeLtTrans}
\noindent$\Pi_{1}$:
\begin{ptree}
\AXC{$z_{1} = w \to z_{1} < z \to w < z$\lbl{A6}}
\UIC{$\forall w\,(z_{1} = w \to z_{1} < z \to w < z)$}
\UIC{$\forall z_{1}\,\forall w\,(z_{1} = w \to z_{1} < z \to w < z)$}
\AXC{$\varphi_{1}$\lbl{A2}}
\BIC{$\forall w\,(y = w \to y < z \to w < z)$}
\AXC{$\forall w\,(y = w \to y < z \to w < z) \to y = x \to y < z \to x < z$\lbl{A2}}
\BIC{$y = x \to y < z \to x < z$}
\end{ptree}
\noindent$\Pi_{2}$:
\begin{ptree}
\AXC{$x = x$\lbl{A5}}
\AXC{$x = y \to x = x \to y = x$\lbl{A6}}
\AXC{$(x = y \to x = x \to y = x) \to x = x \to x = y \to y = x$\lbl{A1}}
\BIC{$x = x \to x = y \to y = x$}
\BIC{$x = y \to y = x$}
\UIC{$\forall y\,(x = y \to y = x)$}
\UIC{$\forall x\,\forall y\,(x = y \to y = x)$}
\end{ptree}
\noindent$\Pi_{3}$:
\begin{ptree}
\AXC{$\stackrel{\Pi_{2}}{\vdots}$}\noLine
\UIC{$\forall x\,\forall y\,(x = y \to y = x)$}
\AXC{$\forall x\,\forall y\,(x = y \to y = x) \to \forall y\,(x = y \to y = x)$\lbl{A2}}
\BIC{$\forall y\,(x = y \to y = x)$}
\AXC{$\forall y\,(x = y \to y = x) \to x = y \to y = x$\lbl{A2}}
\BIC{$x = y \to y = x$}
\end{ptree}
\noindent$\Pi_{4}$:
\begin{ptree}
\AXC{$\stackrel{\Pi_{3}}{\vdots}$}\noLine
\UIC{$x = y \to y = x$}
\AXC{$x \leq y \to y < z \to x = y \to x = y$\lbl{A1}}
\AXC{$\varphi_{3}$\lbl{A1}}
\BIC{$(x = y \to y = x) \to x \leq y \to y < z \to x = y \to y = x$}
\BIC{$x \leq y \to y < z \to x = y \to y = x$}
\AXC{$(x \leq y \to y < z \to x = y \to y = x) \to \varphi_{6}$\lbl{A1}}
\BIC{$\varphi_{6}$}
\end{ptree}
\noindent$\Pi_{5}$:
\begin{ptree}
\AXC{$\forall x\,\forall y\,\forall z\,(x < y \to y < z \to x < z)$\lbl{[O2]}}
\AXC{$\varphi_{4}$\lbl{A2}}
\BIC{$\forall y\,\forall z\,(x < y \to y < z \to x < z)$}
\AXC{$\varphi_{2}$\lbl{A2}}
\BIC{$\forall z\,(x < y \to y < z \to x < z)$}
\AXC{$\forall z\,(x < y \to y < z \to x < z) \to x < y \to y < z \to x < z$\lbl{A2}}
\BIC{$x < y \to y < z \to x < z$}
\end{ptree}
\noindent$\Pi_{6}$:
\begin{ptree}
\AXC{$x \leq y \to y < z \to y < z$\lbl{A1}}
\AXC{$x \leq y \to x \leq y$\lbl{A1}}
\AXC{$(x \leq y \to x \leq y) \to (x \leq y \to y < z \to y < z) \to \varphi_{5}$\lbl{A1}}
\BIC{$(x \leq y \to y < z \to y < z) \to \varphi_{5}$}
\BIC{$\varphi_{5}$}
\end{ptree}
\noindent 主木:
\begin{ptree}
\AXC{$\stackrel{\Pi_{1}}{\vdots}$}\noLine
\UIC{$y = x \to y < z \to x < z$}
\AXC{$x \leq y \to y < z \to x = y \to y < z$\lbl{A1}}
\AXC{$\stackrel{\Pi_{4}}{\vdots}$}\noLine
\UIC{$\varphi_{6}$}
\BIC{$(y = x \to y < z \to x < z) \to x \leq y \to y < z \to x = y \to x < z$}
\BIC{$x \leq y \to y < z \to x = y \to x < z$}
\AXC{$\stackrel{\Pi_{5}}{\vdots}$}\noLine
\UIC{$x < y \to y < z \to x < z$}
\AXC{$\stackrel{\Pi_{6}}{\vdots}$}\noLine
\UIC{$\varphi_{5}$}
\BIC{$(x \leq y \to y < z \to x = y \to x < z) \to x \leq y \to y < z \to x < z$}
\BIC{$x \leq y \to y < z \to x < z$}
\UIC{$\forall z\,(x \leq y \to y < z \to x < z)$}
\UIC{$\forall y\,\forall z\,(x \leq y \to y < z \to x < z)$}
\UIC{$\forall x\,\forall y\,\forall z\,(x \leq y \to y < z \to x < z)$}
\end{ptree}
\begin{abbrevlist}
\item $\varphi_{1} :\equiv \forall z_{1}\,\forall w\,(z_{1} = w \to z_{1} < z \to w < z) \to \forall w\,(y = w \to y < z \to w < z)$
\item $\varphi_{2} :\equiv \forall y\,\forall z\,(x < y \to y < z \to x < z) \to \forall z\,(x < y \to y < z \to x < z)$
\item $\varphi_{3} :\equiv (x \leq y \to y < z \to x = y \to x = y) \to (x = y \to y = x) \to x \leq y \to y < z \to x = y \to y = x$
\item $\varphi_{4} :\equiv \forall x\,\forall y\,\forall z\,(x < y \to y < z \to x < z) \to \forall y\,\forall z\,(x < y \to y < z \to x < z)$
\item $\varphi_{5} :\equiv (x < y \to y < z \to x < z) \to (x \leq y \to y < z \to x = y \to x < z) \to x \leq y \to y < z \to x < z$
\item $\varphi_{6} :\equiv (x \leq y \to y < z \to x = y \to y < z) \to (y = x \to y < z \to x < z) \to x \leq y \to y < z \to x = y \to x < z$
\end{abbrevlist}
\end{fullproof}

% 自動生成 (prover/build.js): LeTrans  ノード数 52
\begin{fullproof}{[O5]}{$\forall x\,\forall y\,\forall z\,(x \leq y \to y \leq z \to x \leq z)$}{52}\phantomsection\label{tree:LeTrans}
\noindent$\Pi_{1}$:
\begin{ptree}
\AXC{$z_{1} = w \to z_{1} = z \to w = z$\lbl{A6}}
\UIC{$\forall w\,(z_{1} = w \to z_{1} = z \to w = z)$}
\UIC{$\forall z_{1}\,\forall w\,(z_{1} = w \to z_{1} = z \to w = z)$}
\AXC{$\varphi_{1}$\lbl{A2}}
\BIC{$\forall w\,(y = w \to y = z \to w = z)$}
\AXC{$\forall w\,(y = w \to y = z \to w = z) \to y = x \to y = z \to x = z$\lbl{A2}}
\BIC{$y = x \to y = z \to x = z$}
\end{ptree}
\noindent$\Pi_{2}$:
\begin{ptree}
\AXC{$x = x$\lbl{A5}}
\AXC{$x = y \to x = x \to y = x$\lbl{A6}}
\AXC{$(x = y \to x = x \to y = x) \to x = x \to x = y \to y = x$\lbl{A1}}
\BIC{$x = x \to x = y \to y = x$}
\BIC{$x = y \to y = x$}
\UIC{$\forall y\,(x = y \to y = x)$}
\UIC{$\forall x\,\forall y\,(x = y \to y = x)$}
\end{ptree}
\noindent$\Pi_{3}$:
\begin{ptree}
\AXC{$\stackrel{\Pi_{2}}{\vdots}$}\noLine
\UIC{$\forall x\,\forall y\,(x = y \to y = x)$}
\AXC{$\forall x\,\forall y\,(x = y \to y = x) \to \forall y\,(x = y \to y = x)$\lbl{A2}}
\BIC{$\forall y\,(x = y \to y = x)$}
\AXC{$\forall y\,(x = y \to y = x) \to x = y \to y = x$\lbl{A2}}
\BIC{$x = y \to y = x$}
\end{ptree}
\noindent$\Pi_{4}$:
\begin{ptree}
\AXC{$x \leq y \to y \leq z \to x = y \to y \leq z$\lbl{A1}}
\AXC{$\stackrel{\Pi_{3}}{\vdots}$}\noLine
\UIC{$x = y \to y = x$}
\AXC{$x \leq y \to y \leq z \to x = y \to x = y$\lbl{A1}}
\AXC{$\varphi_{4}$\lbl{A1}}
\BIC{$(x = y \to y = x) \to x \leq y \to y \leq z \to x = y \to y = x$}
\BIC{$x \leq y \to y \leq z \to x = y \to y = x$}
\AXC{$\varphi_{8}$\lbl{A1}}
\BIC{$(x \leq y \to y \leq z \to x = y \to y \leq z) \to \varphi_{6}$}
\BIC{$\varphi_{6}$}
\end{ptree}
\noindent$\Pi_{5}$:
\begin{ptree}
\AXC{$z_{1} = w \to z_{1} < z \to w < z$\lbl{A6}}
\UIC{$\forall w\,(z_{1} = w \to z_{1} < z \to w < z)$}
\UIC{$\forall z_{1}\,\forall w\,(z_{1} = w \to z_{1} < z \to w < z)$}
\AXC{$\varphi_{3}$\lbl{A2}}
\BIC{$\forall w\,(y = w \to y < z \to w < z)$}
\AXC{$\forall w\,(y = w \to y < z \to w < z) \to y = x \to y < z \to x < z$\lbl{A2}}
\BIC{$y = x \to y < z \to x < z$}
\AXC{$\stackrel{\Pi_{4}}{\vdots}$}\noLine
\UIC{$\varphi_{6}$}
\BIC{$(y = x \to y = z \to x = z) \to x \leq y \to y \leq z \to x = y \to x \leq z$}
\end{ptree}
\noindent$\Pi_{6}$:
\begin{ptree}
\AXC{$\stackrel{\Pi_{1}}{\vdots}$}\noLine
\UIC{$y = x \to y = z \to x = z$}
\AXC{$\stackrel{\Pi_{5}}{\vdots}$}\noLine
\UIC{$(y = x \to y = z \to x = z) \to x \leq y \to y \leq z \to x = y \to x \leq z$}
\BIC{$x \leq y \to y \leq z \to x = y \to x \leq z$}
\end{ptree}
\noindent$\Pi_{7}$:
\begin{ptree}
\AXC{$\forall x\,\forall y\,\forall z\,(x < y \to y \leq z \to x < z)$\lbl{[O3]}}
\AXC{$\varphi_{5}$\lbl{A2}}
\BIC{$\forall y\,\forall z\,(x < y \to y \leq z \to x < z)$}
\AXC{$\varphi_{2}$\lbl{A2}}
\BIC{$\forall z\,(x < y \to y \leq z \to x < z)$}
\AXC{$\forall z\,(x < y \to y \leq z \to x < z) \to x < y \to y \leq z \to x < z$\lbl{A2}}
\BIC{$x < y \to y \leq z \to x < z$}
\end{ptree}
\noindent$\Pi_{8}$:
\begin{ptree}
\AXC{$x \leq y \to y \leq z \to y \leq z$\lbl{A1}}
\AXC{$x \leq y \to x \leq y$\lbl{A1}}
\AXC{$(x \leq y \to x \leq y) \to (x \leq y \to y \leq z \to y \leq z) \to \varphi_{7}$\lbl{A1}}
\BIC{$(x \leq y \to y \leq z \to y \leq z) \to \varphi_{7}$}
\BIC{$\varphi_{7}$}
\end{ptree}
\noindent 主木:
\begin{ptree}
\AXC{$\stackrel{\Pi_{6}}{\vdots}$}\noLine
\UIC{$x \leq y \to y \leq z \to x = y \to x \leq z$}
\AXC{$\stackrel{\Pi_{7}}{\vdots}$}\noLine
\UIC{$x < y \to y \leq z \to x < z$}
\AXC{$\stackrel{\Pi_{8}}{\vdots}$}\noLine
\UIC{$\varphi_{7}$}
\BIC{$(x \leq y \to y \leq z \to x = y \to x \leq z) \to x \leq y \to y \leq z \to x \leq z$}
\BIC{$x \leq y \to y \leq z \to x \leq z$}
\UIC{$\forall z\,(x \leq y \to y \leq z \to x \leq z)$}
\UIC{$\forall y\,\forall z\,(x \leq y \to y \leq z \to x \leq z)$}
\UIC{$\forall x\,\forall y\,\forall z\,(x \leq y \to y \leq z \to x \leq z)$}
\end{ptree}
\begin{abbrevlist}
\item $\varphi_{1} :\equiv \forall z_{1}\,\forall w\,(z_{1} = w \to z_{1} = z \to w = z) \to \forall w\,(y = w \to y = z \to w = z)$
\item $\varphi_{2} :\equiv \forall y\,\forall z\,(x < y \to y \leq z \to x < z) \to \forall z\,(x < y \to y \leq z \to x < z)$
\item $\varphi_{3} :\equiv \forall z_{1}\,\forall w\,(z_{1} = w \to z_{1} < z \to w < z) \to \forall w\,(y = w \to y < z \to w < z)$
\item $\varphi_{4} :\equiv (x \leq y \to y \leq z \to x = y \to x = y) \to (x = y \to y = x) \to x \leq y \to y \leq z \to x = y \to y = x$
\item $\varphi_{5} :\equiv \forall x\,\forall y\,\forall z\,(x < y \to y \leq z \to x < z) \to \forall y\,\forall z\,(x < y \to y \leq z \to x < z)$
\item $\varphi_{6} :\equiv (y = x \to y < z \to x < z) \to (y = x \to y = z \to x = z) \to x \leq y \to y \leq z \to x = y \to x \leq z$
\item $\varphi_{7} :\equiv (x < y \to y \leq z \to x < z) \to (x \leq y \to y \leq z \to x = y \to x \leq z) \to x \leq y \to y \leq z \to x \leq z$
\item $\varphi_{8} :\equiv (x \leq y \to y \leq z \to x = y \to y = x) \to (x \leq y \to y \leq z \to x = y \to y \leq z) \to \varphi_{6}$
\end{abbrevlist}
\end{fullproof}

% 自動生成 (prover/build.js): LeAntisym  ノード数 41
\begin{fullproof}{[O6]}{$\forall x\,\forall y\,(x \leq y \to y \leq x \to x = y)$}{41}\phantomsection\label{tree:LeAntisym}
\noindent$\Pi_{1}$:
\begin{ptree}
\AXC{$x = x$\lbl{A5}}
\AXC{$x = y \to x = x \to y = x$\lbl{A6}}
\AXC{$(x = y \to x = x \to y = x) \to x = x \to x = y \to y = x$\lbl{A1}}
\BIC{$x = x \to x = y \to y = x$}
\BIC{$x = y \to y = x$}
\UIC{$\forall y\,(x = y \to y = x)$}
\UIC{$\forall x\,\forall y\,(x = y \to y = x)$}
\end{ptree}
\noindent$\Pi_{2}$:
\begin{ptree}
\AXC{$\stackrel{\Pi_{1}}{\vdots}$}\noLine
\UIC{$\forall x\,\forall y\,(x = y \to y = x)$}
\AXC{$\forall x\,\forall y\,(x = y \to y = x) \to \forall y\,(v_{9} = y \to y = v_{9})$\lbl{A2}}
\BIC{$\forall y\,(v_{9} = y \to y = v_{9})$}
\end{ptree}
\noindent$\Pi_{3}$:
\begin{ptree}
\AXC{$\stackrel{\Pi_{2}}{\vdots}$}\noLine
\UIC{$\forall y\,(v_{9} = y \to y = v_{9})$}
\AXC{$\forall y\,(v_{9} = y \to y = v_{9}) \to v_{9} = v_{10} \to v_{10} = v_{9}$\lbl{A2}}
\BIC{$v_{9} = v_{10} \to v_{10} = v_{9}$}
\UIC{$\forall v_{10}\,(v_{9} = v_{10} \to v_{10} = v_{9})$}
\UIC{$\forall v_{9}\,\forall v_{10}\,(v_{9} = v_{10} \to v_{10} = v_{9})$}
\AXC{$\varphi_{1}$\lbl{A2}}
\BIC{$\forall v_{10}\,(y = v_{10} \to v_{10} = y)$}
\end{ptree}
\noindent$\Pi_{4}$:
\begin{ptree}
\AXC{$\stackrel{\Pi_{3}}{\vdots}$}\noLine
\UIC{$\forall v_{10}\,(y = v_{10} \to v_{10} = y)$}
\AXC{$\forall v_{10}\,(y = v_{10} \to v_{10} = y) \to y = x \to x = y$\lbl{A2}}
\BIC{$y = x \to x = y$}
\end{ptree}
\noindent$\Pi_{5}$:
\begin{ptree}
\AXC{$\stackrel{\Pi_{4}}{\vdots}$}\noLine
\UIC{$y = x \to x = y$}
\AXC{$x \leq y \to y \leq x \to y = x \to y = x$\lbl{A1}}
\AXC{$\varphi_{3}$\lbl{A1}}
\BIC{$(y = x \to x = y) \to x \leq y \to y \leq x \to y = x \to x = y$}
\BIC{$x \leq y \to y \leq x \to y = x \to x = y$}
\end{ptree}
\noindent$\Pi_{6}$:
\begin{ptree}
\AXC{$\forall x\,\forall y\,\forall z\,(x < y \to y < z \to x < z)$\lbl{[O2]}}
\AXC{$\varphi_{5}$\lbl{A2}}
\BIC{$\forall y\,\forall z\,(x < y \to y < z \to x < z)$}
\AXC{$\varphi_{4}$\lbl{A2}}
\BIC{$\forall z\,(x < y \to y < z \to x < z)$}
\AXC{$\forall z\,(x < y \to y < z \to x < z) \to x < y \to y < x \to x < x$\lbl{A2}}
\BIC{$x < y \to y < x \to x < x$}
\end{ptree}
\noindent$\Pi_{7}$:
\begin{ptree}
\AXC{$\stackrel{\Pi_{6}}{\vdots}$}\noLine
\UIC{$x < y \to y < x \to x < x$}
\AXC{$x \leq y \to y \leq x \to y \leq x$\lbl{A1}}
\AXC{$x \leq y \to x \leq y$\lbl{A1}}
\AXC{$\varphi_{6}$\lbl{A1}}
\BIC{$(x \leq y \to y \leq x \to y \leq x) \to (x < y \to y < x \to x < x) \to \varphi_{2}$}
\BIC{$(x < y \to y < x \to x < x) \to \varphi_{2}$}
\BIC{$\varphi_{2}$}
\end{ptree}
\noindent 主木:
\begin{ptree}
\AXC{$\stackrel{\Pi_{5}}{\vdots}$}\noLine
\UIC{$x \leq y \to y \leq x \to y = x \to x = y$}
\AXC{$\forall x\,(\neg x < x)$\lbl{[O1]}}
\AXC{$\forall x\,(\neg x < x) \to \neg x < x$\lbl{A2}}
\BIC{$\neg x < x$}
\AXC{$\stackrel{\Pi_{7}}{\vdots}$}\noLine
\UIC{$\varphi_{2}$}
\BIC{$(x \leq y \to y \leq x \to y = x \to x = y) \to x \leq y \to y \leq x \to x = y$}
\BIC{$x \leq y \to y \leq x \to x = y$}
\UIC{$\forall y\,(x \leq y \to y \leq x \to x = y)$}
\UIC{$\forall x\,\forall y\,(x \leq y \to y \leq x \to x = y)$}
\end{ptree}
\begin{abbrevlist}
\item $\varphi_{1} :\equiv \forall v_{9}\,\forall v_{10}\,(v_{9} = v_{10} \to v_{10} = v_{9}) \to \forall v_{10}\,(y = v_{10} \to v_{10} = y)$
\item $\varphi_{2} :\equiv \neg x < x \to (x \leq y \to y \leq x \to y = x \to x = y) \to x \leq y \to y \leq x \to x = y$
\item $\varphi_{3} :\equiv (x \leq y \to y \leq x \to y = x \to y = x) \to (y = x \to x = y) \to x \leq y \to y \leq x \to y = x \to x = y$
\item $\varphi_{4} :\equiv \forall y\,\forall z\,(x < y \to y < z \to x < z) \to \forall z\,(x < y \to y < z \to x < z)$
\item $\varphi_{5} :\equiv \forall x\,\forall y\,\forall z\,(x < y \to y < z \to x < z) \to \forall y\,\forall z\,(x < y \to y < z \to x < z)$
\item $\varphi_{6} :\equiv (x \leq y \to x \leq y) \to (x \leq y \to y \leq x \to y \leq x) \to (x < y \to y < x \to x < x) \to \varphi_{2}$
\end{abbrevlist}
\end{fullproof}

% 自動生成 (prover/build.js): LeNotLt  ノード数 48
\begin{fullproof}{[O7]}{$\forall x\,\forall y\,(x \leq y \to \neg y < x)$}{48}\phantomsection\label{tree:LeNotLt}
\noindent$\Pi_{1}$:
\begin{ptree}
\AXC{$z = w \to z < x \to w < x$\lbl{A6}}
\UIC{$\forall w\,(z = w \to z < x \to w < x)$}
\UIC{$\forall z\,\forall w\,(z = w \to z < x \to w < x)$}
\AXC{$\varphi_{2}$\lbl{A2}}
\BIC{$\forall w\,(y = w \to y < x \to w < x)$}
\AXC{$\forall w\,(y = w \to y < x \to w < x) \to y = x \to y < x \to x < x$\lbl{A2}}
\BIC{$y = x \to y < x \to x < x$}
\end{ptree}
\noindent$\Pi_{2}$:
\begin{ptree}
\AXC{$x = x$\lbl{A5}}
\AXC{$x = y \to x = x \to y = x$\lbl{A6}}
\AXC{$(x = y \to x = x \to y = x) \to x = x \to x = y \to y = x$\lbl{A1}}
\BIC{$x = x \to x = y \to y = x$}
\BIC{$x = y \to y = x$}
\UIC{$\forall y\,(x = y \to y = x)$}
\UIC{$\forall x\,\forall y\,(x = y \to y = x)$}
\end{ptree}
\noindent$\Pi_{3}$:
\begin{ptree}
\AXC{$\stackrel{\Pi_{2}}{\vdots}$}\noLine
\UIC{$\forall x\,\forall y\,(x = y \to y = x)$}
\AXC{$\forall x\,\forall y\,(x = y \to y = x) \to \forall y\,(x = y \to y = x)$\lbl{A2}}
\BIC{$\forall y\,(x = y \to y = x)$}
\AXC{$\forall y\,(x = y \to y = x) \to x = y \to y = x$\lbl{A2}}
\BIC{$x = y \to y = x$}
\end{ptree}
\noindent$\Pi_{4}$:
\begin{ptree}
\AXC{$\stackrel{\Pi_{3}}{\vdots}$}\noLine
\UIC{$x = y \to y = x$}
\AXC{$x \leq y \to x = y \to x = y$\lbl{A1}}
\AXC{$\varphi_{1}$\lbl{A1}}
\BIC{$(x = y \to y = x) \to x \leq y \to x = y \to y = x$}
\BIC{$x \leq y \to x = y \to y = x$}
\AXC{$(x \leq y \to x = y \to y = x) \to \varphi_{7}$\lbl{A1}}
\BIC{$\varphi_{7}$}
\end{ptree}
\noindent$\Pi_{5}$:
\begin{ptree}
\AXC{$\stackrel{\Pi_{1}}{\vdots}$}\noLine
\UIC{$y = x \to y < x \to x < x$}
\AXC{$x \leq y \to x = y \to y < x \to y < x$\lbl{A1}}
\AXC{$\stackrel{\Pi_{4}}{\vdots}$}\noLine
\UIC{$\varphi_{7}$}
\BIC{$(y = x \to y < x \to x < x) \to x \leq y \to x = y \to y < x \to x < x$}
\BIC{$x \leq y \to x = y \to y < x \to x < x$}
\AXC{$\varphi_{6}$\lbl{A1}}
\BIC{$\neg x < x \to x \leq y \to x = y \to y < x \to \neg y < x$}
\end{ptree}
\noindent$\Pi_{6}$:
\begin{ptree}
\AXC{$\forall x\,(\neg x < x)$\lbl{[O1]}}
\AXC{$\forall x\,(\neg x < x) \to \neg x < x$\lbl{A2}}
\BIC{$\neg x < x$}
\AXC{$\stackrel{\Pi_{5}}{\vdots}$}\noLine
\UIC{$\neg x < x \to x \leq y \to x = y \to y < x \to \neg y < x$}
\BIC{$x \leq y \to x = y \to y < x \to \neg y < x$}
\end{ptree}
\noindent$\Pi_{7}$:
\begin{ptree}
\AXC{$\forall x\,\forall y\,\forall z\,(x < y \to y < z \to x < z)$\lbl{[O2]}}
\AXC{$\varphi_{5}$\lbl{A2}}
\BIC{$\forall y\,\forall z\,(x < y \to y < z \to x < z)$}
\AXC{$\varphi_{4}$\lbl{A2}}
\BIC{$\forall z\,(x < y \to y < z \to x < z)$}
\AXC{$\forall z\,(x < y \to y < z \to x < z) \to x < y \to y < x \to x < x$\lbl{A2}}
\BIC{$x < y \to y < x \to x < x$}
\end{ptree}
\noindent$\Pi_{8}$:
\begin{ptree}
\AXC{$\stackrel{\Pi_{7}}{\vdots}$}\noLine
\UIC{$x < y \to y < x \to x < x$}
\AXC{$x \leq y \to x \leq y$\lbl{A1}}
\AXC{$(x \leq y \to x \leq y) \to (x < y \to y < x \to x < x) \to \varphi_{3}$\lbl{A1}}
\BIC{$(x < y \to y < x \to x < x) \to \varphi_{3}$}
\BIC{$\varphi_{3}$}
\end{ptree}
\noindent 主木:
\begin{ptree}
\AXC{$\stackrel{\Pi_{6}}{\vdots}$}\noLine
\UIC{$x \leq y \to x = y \to y < x \to \neg y < x$}
\AXC{$\forall x\,(\neg x < x)$\lbl{[O1]}}
\AXC{$\forall x\,(\neg x < x) \to \neg x < x$\lbl{A2}}
\BIC{$\neg x < x$}
\AXC{$\stackrel{\Pi_{8}}{\vdots}$}\noLine
\UIC{$\varphi_{3}$}
\BIC{$(x \leq y \to x = y \to y < x \to \neg y < x) \to x \leq y \to \neg y < x$}
\BIC{$x \leq y \to \neg y < x$}
\UIC{$\forall y\,(x \leq y \to \neg y < x)$}
\UIC{$\forall x\,\forall y\,(x \leq y \to \neg y < x)$}
\end{ptree}
\begin{abbrevlist}
\item $\varphi_{1} :\equiv (x \leq y \to x = y \to x = y) \to (x = y \to y = x) \to x \leq y \to x = y \to y = x$
\item $\varphi_{2} :\equiv \forall z\,\forall w\,(z = w \to z < x \to w < x) \to \forall w\,(y = w \to y < x \to w < x)$
\item $\varphi_{3} :\equiv \neg x < x \to (x \leq y \to x = y \to y < x \to \neg y < x) \to x \leq y \to \neg y < x$
\item $\varphi_{4} :\equiv \forall y\,\forall z\,(x < y \to y < z \to x < z) \to \forall z\,(x < y \to y < z \to x < z)$
\item $\varphi_{5} :\equiv \forall x\,\forall y\,\forall z\,(x < y \to y < z \to x < z) \to \forall y\,\forall z\,(x < y \to y < z \to x < z)$
\item $\varphi_{6} :\equiv (x \leq y \to x = y \to y < x \to x < x) \to \neg x < x \to x \leq y \to x = y \to y < x \to \neg y < x$
\item $\varphi_{7} :\equiv (x \leq y \to x = y \to y < x \to y < x) \to (y = x \to y < x \to x < x) \to x \leq y \to x = y \to y < x \to x < x$
\end{abbrevlist}
\end{fullproof}

% 自動生成 (prover/build.js): NotLtLe  ノード数 33
\begin{fullproof}{[O8]}{$\forall x\,\forall y\,(x \in \mathbb{R} \to y \in \mathbb{R} \to \neg x < y \to y \leq x)$}{33}\phantomsection\label{tree:NotLtLe}
\noindent$\Pi_{1}$:
\begin{ptree}
\AXC{$x = x$\lbl{A5}}
\AXC{$x = y \to x = x \to y = x$\lbl{A6}}
\AXC{$(x = y \to x = x \to y = x) \to x = x \to x = y \to y = x$\lbl{A1}}
\BIC{$x = x \to x = y \to y = x$}
\BIC{$x = y \to y = x$}
\UIC{$\forall y\,(x = y \to y = x)$}
\UIC{$\forall x\,\forall y\,(x = y \to y = x)$}
\end{ptree}
\noindent$\Pi_{2}$:
\begin{ptree}
\AXC{$\stackrel{\Pi_{1}}{\vdots}$}\noLine
\UIC{$\forall x\,\forall y\,(x = y \to y = x)$}
\AXC{$\forall x\,\forall y\,(x = y \to y = x) \to \forall y\,(x = y \to y = x)$\lbl{A2}}
\BIC{$\forall y\,(x = y \to y = x)$}
\AXC{$\forall y\,(x = y \to y = x) \to x = y \to y = x$\lbl{A2}}
\BIC{$x = y \to y = x$}
\end{ptree}
\noindent$\Pi_{3}$:
\begin{ptree}
\AXC{$\stackrel{\Pi_{2}}{\vdots}$}\noLine
\UIC{$x = y \to y = x$}
\AXC{$x \in \mathbb{R} \to y \in \mathbb{R} \to \neg x < y \to x = y \to x = y$\lbl{A1}}
\AXC{$(x \in \mathbb{R} \to y \in \mathbb{R} \to \neg x < y \to x = y \to x = y) \to \varphi_{1}$\lbl{A1}}
\BIC{$\varphi_{1}$}
\BIC{$x \in \mathbb{R} \to y \in \mathbb{R} \to \neg x < y \to x = y \to y = x$}
\end{ptree}
\noindent$\Pi_{4}$:
\begin{ptree}
\AXC{$\forall x\,\forall y\,(x \in \mathbb{R} \to y \in \mathbb{R} \to x < y \lor x = y \lor y < x)$\lbl{(R13)}}
\AXC{$\varphi_{4}$\lbl{A2}}
\BIC{$\forall y\,(x \in \mathbb{R} \to y \in \mathbb{R} \to x < y \lor x = y \lor y < x)$}
\AXC{$\varphi_{3}$\lbl{A2}}
\BIC{$x \in \mathbb{R} \to y \in \mathbb{R} \to x < y \lor x = y \lor y < x$}
\AXC{$\varphi_{7}$\lbl{A1}}
\BIC{$(x \in \mathbb{R} \to y \in \mathbb{R} \to \neg x < y \to x \in \mathbb{R}) \to \varphi_{5}$}
\end{ptree}
\noindent$\Pi_{5}$:
\begin{ptree}
\AXC{$x \in \mathbb{R} \to y \in \mathbb{R} \to \neg x < y \to y \in \mathbb{R}$\lbl{A1}}
\AXC{$x \in \mathbb{R} \to y \in \mathbb{R} \to \neg x < y \to x \in \mathbb{R}$\lbl{A1}}
\AXC{$\stackrel{\Pi_{4}}{\vdots}$}\noLine
\UIC{$(x \in \mathbb{R} \to y \in \mathbb{R} \to \neg x < y \to x \in \mathbb{R}) \to \varphi_{5}$}
\BIC{$\varphi_{5}$}
\BIC{$x \in \mathbb{R} \to y \in \mathbb{R} \to \neg x < y \to x < y \lor x = y \lor y < x$}
\end{ptree}
\noindent$\Pi_{6}$:
\begin{ptree}
\AXC{$\stackrel{\Pi_{5}}{\vdots}$}\noLine
\UIC{$x \in \mathbb{R} \to y \in \mathbb{R} \to \neg x < y \to x < y \lor x = y \lor y < x$}
\AXC{$x \in \mathbb{R} \to y \in \mathbb{R} \to \neg x < y \to \neg x < y$\lbl{A1}}
\AXC{$\varphi_{6}$\lbl{A1}}
\BIC{$(x \in \mathbb{R} \to y \in \mathbb{R} \to \neg x < y \to x < y \lor x = y \lor y < x) \to \varphi_{2}$}
\BIC{$\varphi_{2}$}
\end{ptree}
\noindent 主木:
\begin{ptree}
\AXC{$\stackrel{\Pi_{3}}{\vdots}$}\noLine
\UIC{$x \in \mathbb{R} \to y \in \mathbb{R} \to \neg x < y \to x = y \to y = x$}
\AXC{$\stackrel{\Pi_{6}}{\vdots}$}\noLine
\UIC{$\varphi_{2}$}
\BIC{$x \in \mathbb{R} \to y \in \mathbb{R} \to \neg x < y \to y \leq x$}
\UIC{$\forall y\,(x \in \mathbb{R} \to y \in \mathbb{R} \to \neg x < y \to y \leq x)$}
\UIC{$\forall x\,\forall y\,(x \in \mathbb{R} \to y \in \mathbb{R} \to \neg x < y \to y \leq x)$}
\end{ptree}
\begin{abbrevlist}
\item $\varphi_{1} :\equiv (x = y \to y = x) \to x \in \mathbb{R} \to y \in \mathbb{R} \to \neg x < y \to x = y \to y = x$
\item $\varphi_{2} :\equiv (x \in \mathbb{R} \to y \in \mathbb{R} \to \neg x < y \to x = y \to y = x) \to x \in \mathbb{R} \to y \in \mathbb{R} \to \neg x < y \to y \leq x$
\item $\varphi_{3} :\equiv \forall y\,(x \in \mathbb{R} \to y \in \mathbb{R} \to x < y \lor x = y \lor y < x) \to x \in \mathbb{R} \to y \in \mathbb{R} \to x < y \lor x = y \lor y < x$
\item $\varphi_{4} :\equiv \forall x\,\forall y\,(x \in \mathbb{R} \to y \in \mathbb{R} \to x < y \lor x = y \lor y < x) \to \forall y\,(x \in \mathbb{R} \to y \in \mathbb{R} \to x < y \lor x = y \lor y < x)$
\item $\varphi_{5} :\equiv (x \in \mathbb{R} \to y \in \mathbb{R} \to \neg x < y \to y \in \mathbb{R}) \to x \in \mathbb{R} \to y \in \mathbb{R} \to \neg x < y \to x < y \lor x = y \lor y < x$
\item $\varphi_{6} :\equiv (x \in \mathbb{R} \to y \in \mathbb{R} \to \neg x < y \to \neg x < y) \to (x \in \mathbb{R} \to y \in \mathbb{R} \to \neg x < y \to x < y \lor x = y \lor y < x) \to \varphi_{2}$
\item $\varphi_{7} :\equiv (x \in \mathbb{R} \to y \in \mathbb{R} \to x < y \lor x = y \lor y < x) \to (x \in \mathbb{R} \to y \in \mathbb{R} \to \neg x < y \to x \in \mathbb{R}) \to \varphi_{5}$
\end{abbrevlist}
\end{fullproof}

\section{順序体の基本性質 (公理 (R1)–(R16) からの完全展開した証明)}
% 自動生成 (prover/build.js): SubR  ノード数 55
\begin{fullproof}{[F1]}{$\forall x\,\forall y\,(x \in \mathbb{R} \to y \in \mathbb{R} \to x - y \in \mathbb{R})$}{55}\phantomsection\label{tree:SubR}
\noindent$\Pi_{1}$:
\begin{ptree}
\AXC{$\forall x \in \mathbb{R}\,(-x \in \mathbb{R})$\lbl{(R1)}}
\AXC{$\forall x \in \mathbb{R}\,(-x \in \mathbb{R}) \to y \in \mathbb{R} \to -y \in \mathbb{R}$\lbl{A2}}
\BIC{$y \in \mathbb{R} \to -y \in \mathbb{R}$}
\AXC{$(y \in \mathbb{R} \to -y \in \mathbb{R}) \to \varphi_{4}$\lbl{A1}}
\BIC{$\varphi_{4}$}
\end{ptree}
\noindent$\Pi_{2}$:
\begin{ptree}
\AXC{$x \in \mathbb{R} \to y \in \mathbb{R} \to x \in \mathbb{R}$\lbl{A1}}
\AXC{$\forall x\,\forall y\,(x \in \mathbb{R} \to y \in \mathbb{R} \to x + y \in \mathbb{R})$\lbl{(R1)}}
\AXC{$\varphi_{8}$\lbl{A2}}
\BIC{$\forall y\,(x \in \mathbb{R} \to y \in \mathbb{R} \to x + y \in \mathbb{R})$}
\AXC{$\varphi_{9}$\lbl{A2}}
\BIC{$x \in \mathbb{R} \to -y \in \mathbb{R} \to x + -y \in \mathbb{R}$}
\AXC{$\varphi_{12}$\lbl{A1}}
\BIC{$(x \in \mathbb{R} \to y \in \mathbb{R} \to x \in \mathbb{R}) \to \varphi_{7}$}
\BIC{$\varphi_{7}$}
\end{ptree}
\noindent$\Pi_{3}$:
\begin{ptree}
\AXC{$x = x$\lbl{A5}}
\AXC{$x = y \to x = x \to y = x$\lbl{A6}}
\AXC{$(x = y \to x = x \to y = x) \to x = x \to x = y \to y = x$\lbl{A1}}
\BIC{$x = x \to x = y \to y = x$}
\BIC{$x = y \to y = x$}
\UIC{$\forall y\,(x = y \to y = x)$}
\UIC{$\forall x\,\forall y\,(x = y \to y = x)$}
\end{ptree}
\noindent$\Pi_{4}$:
\begin{ptree}
\AXC{$\stackrel{\Pi_{3}}{\vdots}$}\noLine
\UIC{$\forall x\,\forall y\,(x = y \to y = x)$}
\AXC{$\forall x\,\forall y\,(x = y \to y = x) \to \forall y\,(v_{11} = y \to y = v_{11})$\lbl{A2}}
\BIC{$\forall y\,(v_{11} = y \to y = v_{11})$}
\AXC{$\varphi_{1}$\lbl{A2}}
\BIC{$v_{11} = v_{12} \to v_{12} = v_{11}$}
\UIC{$\forall v_{12}\,(v_{11} = v_{12} \to v_{12} = v_{11})$}
\UIC{$\forall v_{11}\,\forall v_{12}\,(v_{11} = v_{12} \to v_{12} = v_{11})$}
\AXC{$\varphi_{2}$\lbl{A2}}
\BIC{$\forall v_{12}\,(x - y = v_{12} \to v_{12} = x - y)$}
\AXC{$\varphi_{3}$\lbl{A2}}
\BIC{$x - y = x + -y \to x + -y = x - y$}
\end{ptree}
\noindent$\Pi_{5}$:
\begin{ptree}
\AXC{$\forall x\,\forall y\,(x - y = x + -y)$\lbl{Def$-$}}
\AXC{$\forall x\,\forall y\,(x - y = x + -y) \to \forall y\,(x - y = x + -y)$\lbl{A2}}
\BIC{$\forall y\,(x - y = x + -y)$}
\end{ptree}
\noindent$\Pi_{6}$:
\begin{ptree}
\AXC{$\stackrel{\Pi_{4}}{\vdots}$}\noLine
\UIC{$x - y = x + -y \to x + -y = x - y$}
\AXC{$\stackrel{\Pi_{5}}{\vdots}$}\noLine
\UIC{$\forall y\,(x - y = x + -y)$}
\AXC{$\forall y\,(x - y = x + -y) \to x - y = x + -y$\lbl{A2}}
\BIC{$x - y = x + -y$}
\AXC{$\varphi_{5}$\lbl{A1}}
\BIC{$(x - y = x + -y \to x + -y = x - y) \to x + -y = x - y$}
\BIC{$x + -y = x - y$}
\end{ptree}
\noindent$\Pi_{7}$:
\begin{ptree}
\AXC{$\stackrel{\Pi_{6}}{\vdots}$}\noLine
\UIC{$x + -y = x - y$}
\AXC{$x + -y = x - y \to (x \in \mathbb{R} \to y \in \mathbb{R} \to x + -y \in \mathbb{R}) \to \varphi_{10}$\lbl{A1}}
\BIC{$(x \in \mathbb{R} \to y \in \mathbb{R} \to x + -y \in \mathbb{R}) \to \varphi_{10}$}
\end{ptree}
\noindent$\Pi_{8}$:
\begin{ptree}
\AXC{$x \in \mathbb{R} \to y \in \mathbb{R} \to y \in \mathbb{R}$\lbl{A1}}
\AXC{$\stackrel{\Pi_{1}}{\vdots}$}\noLine
\UIC{$\varphi_{4}$}
\BIC{$x \in \mathbb{R} \to y \in \mathbb{R} \to -y \in \mathbb{R}$}
\AXC{$\stackrel{\Pi_{2}}{\vdots}$}\noLine
\UIC{$\varphi_{7}$}
\BIC{$x \in \mathbb{R} \to y \in \mathbb{R} \to x + -y \in \mathbb{R}$}
\AXC{$\stackrel{\Pi_{7}}{\vdots}$}\noLine
\UIC{$(x \in \mathbb{R} \to y \in \mathbb{R} \to x + -y \in \mathbb{R}) \to \varphi_{10}$}
\BIC{$\varphi_{10}$}
\end{ptree}
\noindent 主木:
\begin{ptree}
\AXC{$z = w \to z \in \mathbb{R} \to w \in \mathbb{R}$\lbl{A6}}
\UIC{$\forall w\,(z = w \to z \in \mathbb{R} \to w \in \mathbb{R})$}
\UIC{$\forall z\,\forall w\,(z = w \to z \in \mathbb{R} \to w \in \mathbb{R})$}
\AXC{$\varphi_{6}$\lbl{A2}}
\BIC{$\forall w\,(x + -y = w \to x + -y \in \mathbb{R} \to w \in \mathbb{R})$}
\AXC{$\varphi_{11}$\lbl{A2}}
\BIC{$x + -y = x - y \to x + -y \in \mathbb{R} \to x - y \in \mathbb{R}$}
\AXC{$\stackrel{\Pi_{8}}{\vdots}$}\noLine
\UIC{$\varphi_{10}$}
\BIC{$x \in \mathbb{R} \to y \in \mathbb{R} \to x - y \in \mathbb{R}$}
\UIC{$\forall y\,(x \in \mathbb{R} \to y \in \mathbb{R} \to x - y \in \mathbb{R})$}
\UIC{$\forall x\,\forall y\,(x \in \mathbb{R} \to y \in \mathbb{R} \to x - y \in \mathbb{R})$}
\end{ptree}
\begin{abbrevlist}
\item $\varphi_{1} :\equiv \forall y\,(v_{11} = y \to y = v_{11}) \to v_{11} = v_{12} \to v_{12} = v_{11}$
\item $\varphi_{2} :\equiv \forall v_{11}\,\forall v_{12}\,(v_{11} = v_{12} \to v_{12} = v_{11}) \to \forall v_{12}\,(x - y = v_{12} \to v_{12} = x - y)$
\item $\varphi_{3} :\equiv \forall v_{12}\,(x - y = v_{12} \to v_{12} = x - y) \to x - y = x + -y \to x + -y = x - y$
\item $\varphi_{4} :\equiv (x \in \mathbb{R} \to y \in \mathbb{R} \to y \in \mathbb{R}) \to x \in \mathbb{R} \to y \in \mathbb{R} \to -y \in \mathbb{R}$
\item $\varphi_{5} :\equiv x - y = x + -y \to (x - y = x + -y \to x + -y = x - y) \to x + -y = x - y$
\item $\varphi_{6} :\equiv \forall z\,\forall w\,(z = w \to z \in \mathbb{R} \to w \in \mathbb{R}) \to \forall w\,(x + -y = w \to x + -y \in \mathbb{R} \to w \in \mathbb{R})$
\item $\varphi_{7} :\equiv (x \in \mathbb{R} \to y \in \mathbb{R} \to -y \in \mathbb{R}) \to x \in \mathbb{R} \to y \in \mathbb{R} \to x + -y \in \mathbb{R}$
\item $\varphi_{8} :\equiv \forall x\,\forall y\,(x \in \mathbb{R} \to y \in \mathbb{R} \to x + y \in \mathbb{R}) \to \forall y\,(x \in \mathbb{R} \to y \in \mathbb{R} \to x + y \in \mathbb{R})$
\item $\varphi_{9} :\equiv \forall y\,(x \in \mathbb{R} \to y \in \mathbb{R} \to x + y \in \mathbb{R}) \to x \in \mathbb{R} \to -y \in \mathbb{R} \to x + -y \in \mathbb{R}$
\item $\varphi_{10} :\equiv (x + -y = x - y \to x + -y \in \mathbb{R} \to x - y \in \mathbb{R}) \to x \in \mathbb{R} \to y \in \mathbb{R} \to x - y \in \mathbb{R}$
\item $\varphi_{11} :\equiv \forall w\,(x + -y = w \to x + -y \in \mathbb{R} \to w \in \mathbb{R}) \to x + -y = x - y \to x + -y \in \mathbb{R} \to x - y \in \mathbb{R}$
\item $\varphi_{12} :\equiv (x \in \mathbb{R} \to -y \in \mathbb{R} \to x + -y \in \mathbb{R}) \to (x \in \mathbb{R} \to y \in \mathbb{R} \to x \in \mathbb{R}) \to \varphi_{7}$
\end{abbrevlist}
\end{fullproof}

% 自動生成 (prover/build.js): TwoR  ノード数 36
\begin{fullproof}{[F2]}{$2 \in \mathbb{R}$}{36}\phantomsection\label{tree:TwoR}
\noindent$\Pi_{1}$:
\begin{ptree}
\AXC{$1 \in \mathbb{R}$\lbl{(R1)}}
\AXC{$1 \in \mathbb{R}$\lbl{(R1)}}
\AXC{$\forall x\,\forall y\,(x \in \mathbb{R} \to y \in \mathbb{R} \to x + y \in \mathbb{R})$\lbl{(R1)}}
\AXC{$\varphi_{6}$\lbl{A2}}
\BIC{$\forall y\,(1 \in \mathbb{R} \to y \in \mathbb{R} \to 1 + y \in \mathbb{R})$}
\AXC{$\varphi_{5}$\lbl{A2}}
\BIC{$1 \in \mathbb{R} \to 1 \in \mathbb{R} \to 1 + 1 \in \mathbb{R}$}
\AXC{$\varphi_{4}$\lbl{A1}}
\BIC{$1 \in \mathbb{R} \to 1 \in \mathbb{R} \to 1 + 1 \in \mathbb{R}$}
\BIC{$1 \in \mathbb{R} \to 1 + 1 \in \mathbb{R}$}
\BIC{$1 + 1 \in \mathbb{R}$}
\end{ptree}
\noindent$\Pi_{2}$:
\begin{ptree}
\AXC{$x = x$\lbl{A5}}
\AXC{$x = y \to x = x \to y = x$\lbl{A6}}
\AXC{$(x = y \to x = x \to y = x) \to x = x \to x = y \to y = x$\lbl{A1}}
\BIC{$x = x \to x = y \to y = x$}
\BIC{$x = y \to y = x$}
\UIC{$\forall y\,(x = y \to y = x)$}
\UIC{$\forall x\,\forall y\,(x = y \to y = x)$}
\end{ptree}
\noindent$\Pi_{3}$:
\begin{ptree}
\AXC{$\stackrel{\Pi_{2}}{\vdots}$}\noLine
\UIC{$\forall x\,\forall y\,(x = y \to y = x)$}
\AXC{$\forall x\,\forall y\,(x = y \to y = x) \to \forall y\,(2 = y \to y = 2)$\lbl{A2}}
\BIC{$\forall y\,(2 = y \to y = 2)$}
\AXC{$\forall y\,(2 = y \to y = 2) \to 2 = 1 + 1 \to 1 + 1 = 2$\lbl{A2}}
\BIC{$2 = 1 + 1 \to 1 + 1 = 2$}
\end{ptree}
\noindent$\Pi_{4}$:
\begin{ptree}
\AXC{$\stackrel{\Pi_{3}}{\vdots}$}\noLine
\UIC{$2 = 1 + 1 \to 1 + 1 = 2$}
\AXC{$2 = 1 + 1$\lbl{Def$2$}}
\AXC{$2 = 1 + 1 \to (2 = 1 + 1 \to 1 + 1 = 2) \to 1 + 1 = 2$\lbl{A1}}
\BIC{$(2 = 1 + 1 \to 1 + 1 = 2) \to 1 + 1 = 2$}
\BIC{$1 + 1 = 2$}
\AXC{$1 + 1 = 2 \to \varphi_{1}$\lbl{A1}}
\BIC{$\varphi_{1}$}
\end{ptree}
\noindent 主木:
\begin{ptree}
\AXC{$z = w \to z \in \mathbb{R} \to w \in \mathbb{R}$\lbl{A6}}
\UIC{$\forall w\,(z = w \to z \in \mathbb{R} \to w \in \mathbb{R})$}
\UIC{$\forall z\,\forall w\,(z = w \to z \in \mathbb{R} \to w \in \mathbb{R})$}
\AXC{$\varphi_{2}$\lbl{A2}}
\BIC{$\forall w\,(1 + 1 = w \to 1 + 1 \in \mathbb{R} \to w \in \mathbb{R})$}
\AXC{$\varphi_{3}$\lbl{A2}}
\BIC{$1 + 1 = 2 \to 1 + 1 \in \mathbb{R} \to 2 \in \mathbb{R}$}
\AXC{$\stackrel{\Pi_{1}}{\vdots}$}\noLine
\UIC{$1 + 1 \in \mathbb{R}$}
\AXC{$\stackrel{\Pi_{4}}{\vdots}$}\noLine
\UIC{$\varphi_{1}$}
\BIC{$(1 + 1 = 2 \to 1 + 1 \in \mathbb{R} \to 2 \in \mathbb{R}) \to 2 \in \mathbb{R}$}
\BIC{$2 \in \mathbb{R}$}
\end{ptree}
\begin{abbrevlist}
\item $\varphi_{1} :\equiv 1 + 1 \in \mathbb{R} \to (1 + 1 = 2 \to 1 + 1 \in \mathbb{R} \to 2 \in \mathbb{R}) \to 2 \in \mathbb{R}$
\item $\varphi_{2} :\equiv \forall z\,\forall w\,(z = w \to z \in \mathbb{R} \to w \in \mathbb{R}) \to \forall w\,(1 + 1 = w \to 1 + 1 \in \mathbb{R} \to w \in \mathbb{R})$
\item $\varphi_{3} :\equiv \forall w\,(1 + 1 = w \to 1 + 1 \in \mathbb{R} \to w \in \mathbb{R}) \to 1 + 1 = 2 \to 1 + 1 \in \mathbb{R} \to 2 \in \mathbb{R}$
\item $\varphi_{4} :\equiv (1 \in \mathbb{R} \to 1 \in \mathbb{R} \to 1 + 1 \in \mathbb{R}) \to 1 \in \mathbb{R} \to 1 \in \mathbb{R} \to 1 + 1 \in \mathbb{R}$
\item $\varphi_{5} :\equiv \forall y\,(1 \in \mathbb{R} \to y \in \mathbb{R} \to 1 + y \in \mathbb{R}) \to 1 \in \mathbb{R} \to 1 \in \mathbb{R} \to 1 + 1 \in \mathbb{R}$
\item $\varphi_{6} :\equiv \forall x\,\forall y\,(x \in \mathbb{R} \to y \in \mathbb{R} \to x + y \in \mathbb{R}) \to \forall y\,(1 \in \mathbb{R} \to y \in \mathbb{R} \to 1 + y \in \mathbb{R})$
\end{abbrevlist}
\end{fullproof}

% 自動生成 (prover/build.js): AddCancel  ノード数 256
\begin{fullproof}{[F3]}{$\forall x\,\forall y\,\forall z\,(x \in \mathbb{R} \to y \in \mathbb{R} \to z \in \mathbb{R} \to x + z = y + z \to x = y)$}{256}\phantomsection\label{tree:AddCancel}
\noindent$\Pi_{1}$:
\begin{ptree}
\AXC{$y = z \to x = y \to x = z$\lbl{A6}}
\AXC{$(y = z \to x = y \to x = z) \to x = y \to y = z \to x = z$\lbl{A1}}
\BIC{$x = y \to y = z \to x = z$}
\UIC{$\forall z\,(x = y \to y = z \to x = z)$}
\UIC{$\forall y\,\forall z\,(x = y \to y = z \to x = z)$}
\UIC{$\forall x\,\forall y\,\forall z\,(x = y \to y = z \to x = z)$}
\AXC{$\varphi_{5}$\lbl{A2}}
\BIC{$\forall y\,\forall z\,(x = y \to y = z \to x = z)$}
\AXC{$\varphi_{9}$\lbl{A2}}
\BIC{$\forall z\,(x = y + 0 \to y + 0 = z \to x = z)$}
\AXC{$\varphi_{23}$\lbl{A2}}
\BIC{$x = y + 0 \to y + 0 = y \to x = y$}
\end{ptree}
\noindent$\Pi_{2}$:
\begin{ptree}
\AXC{$\forall x \in \mathbb{R}\,(x + 0 = x)$\lbl{(R4)}}
\AXC{$\forall x \in \mathbb{R}\,(x + 0 = x) \to y \in \mathbb{R} \to y + 0 = y$\lbl{A2}}
\BIC{$y \in \mathbb{R} \to y + 0 = y$}
\AXC{$(y \in \mathbb{R} \to y + 0 = y) \to \varphi_{78}$\lbl{A1}}
\BIC{$\varphi_{78}$}
\end{ptree}
\noindent$\Pi_{3}$:
\begin{ptree}
\AXC{$y = z \to x = y \to x = z$\lbl{A6}}
\AXC{$(y = z \to x = y \to x = z) \to x = y \to y = z \to x = z$\lbl{A1}}
\BIC{$x = y \to y = z \to x = z$}
\UIC{$\forall z\,(x = y \to y = z \to x = z)$}
\UIC{$\forall y\,\forall z\,(x = y \to y = z \to x = z)$}
\UIC{$\forall x\,\forall y\,\forall z\,(x = y \to y = z \to x = z)$}
\AXC{$\varphi_{10}$\lbl{A2}}
\BIC{$\forall y\,\forall z\,(v_{28} = y \to y = z \to v_{28} = z)$}
\AXC{$\varphi_{15}$\lbl{A2}}
\BIC{$\forall z\,(v_{28} = v_{29} \to v_{29} = z \to v_{28} = z)$}
\AXC{$\varphi_{24}$\lbl{A2}}
\BIC{$v_{28} = v_{29} \to v_{29} = v_{30} \to v_{28} = v_{30}$}
\UIC{$\forall v_{30}\,(v_{28} = v_{29} \to v_{29} = v_{30} \to v_{28} = v_{30})$}
\UIC{$\varphi_{1}$}
\UIC{$\forall v_{28}\,(\varphi_{1})$}
\end{ptree}
\noindent$\Pi_{4}$:
\begin{ptree}
\AXC{$\forall x \in \mathbb{R}\,(x + -x = 0)$\lbl{(R5)}}
\AXC{$\forall x \in \mathbb{R}\,(x + -x = 0) \to z \in \mathbb{R} \to z + -z = 0$\lbl{A2}}
\BIC{$z \in \mathbb{R} \to z + -z = 0$}
\AXC{$(z \in \mathbb{R} \to z + -z = 0) \to \varphi_{81}$\lbl{A1}}
\BIC{$\varphi_{81}$}
\end{ptree}
\noindent$\Pi_{5}$:
\begin{ptree}
\AXC{$x \in \mathbb{R} \to y \in \mathbb{R} \to z \in \mathbb{R} \to x + z = y + z \to z \in \mathbb{R}$\lbl{A1}}
\AXC{$\stackrel{\Pi_{4}}{\vdots}$}\noLine
\UIC{$\varphi_{81}$}
\BIC{$x \in \mathbb{R} \to y \in \mathbb{R} \to z \in \mathbb{R} \to x + z = y + z \to z + -z = 0$}
\AXC{$\varphi_{95}$\lbl{A1}}
\BIC{$y + (z + -z) = y + (z + -z) \to \varphi_{43} \to \varphi_{57}$}
\end{ptree}
\noindent$\Pi_{6}$:
\begin{ptree}
\AXC{$x = x$\lbl{A5}}
\UIC{$\forall x\,(x = x)$}
\AXC{$\forall x\,(x = x) \to y + (z + -z) = y + (z + -z)$\lbl{A2}}
\BIC{$y + (z + -z) = y + (z + -z)$}
\AXC{$\stackrel{\Pi_{5}}{\vdots}$}\noLine
\UIC{$y + (z + -z) = y + (z + -z) \to \varphi_{43} \to \varphi_{57}$}
\BIC{$\varphi_{43} \to \varphi_{57}$}
\end{ptree}
\noindent$\Pi_{7}$:
\begin{ptree}
\AXC{$y = z \to x = y \to x = z$\lbl{A6}}
\AXC{$(y = z \to x = y \to x = z) \to x = y \to y = z \to x = z$\lbl{A1}}
\BIC{$x = y \to y = z \to x = z$}
\UIC{$\forall z\,(x = y \to y = z \to x = z)$}
\UIC{$\forall y\,\forall z\,(x = y \to y = z \to x = z)$}
\UIC{$\forall x\,\forall y\,\forall z\,(x = y \to y = z \to x = z)$}
\AXC{$\varphi_{11}$\lbl{A2}}
\BIC{$\forall y\,\forall z\,(v_{25} = y \to y = z \to v_{25} = z)$}
\AXC{$\varphi_{17}$\lbl{A2}}
\BIC{$\forall z\,(v_{25} = v_{26} \to v_{26} = z \to v_{25} = z)$}
\AXC{$\varphi_{25}$\lbl{A2}}
\BIC{$v_{25} = v_{26} \to v_{26} = v_{27} \to v_{25} = v_{27}$}
\UIC{$\forall v_{27}\,(v_{25} = v_{26} \to v_{26} = v_{27} \to v_{25} = v_{27})$}
\UIC{$\varphi_{2}$}
\UIC{$\forall v_{25}\,(\varphi_{2})$}
\end{ptree}
\noindent$\Pi_{8}$:
\begin{ptree}
\AXC{$\stackrel{\Pi_{7}}{\vdots}$}\noLine
\UIC{$\forall v_{25}\,(\varphi_{2})$}
\AXC{$\varphi_{58}$\lbl{A2}}
\BIC{$\forall v_{26}\,\forall v_{27}\,(x = v_{26} \to v_{26} = v_{27} \to x = v_{27})$}
\AXC{$\varphi_{71}$\lbl{A2}}
\BIC{$\varphi_{18}$}
\AXC{$\varphi_{18} \to \varphi_{44}$\lbl{A2}}
\BIC{$\varphi_{44}$}
\end{ptree}
\noindent$\Pi_{9}$:
\begin{ptree}
\AXC{$\forall x \in \mathbb{R}\,(-x \in \mathbb{R})$\lbl{(R1)}}
\AXC{$\forall x \in \mathbb{R}\,(-x \in \mathbb{R}) \to z \in \mathbb{R} \to -z \in \mathbb{R}$\lbl{A2}}
\BIC{$z \in \mathbb{R} \to -z \in \mathbb{R}$}
\AXC{$(z \in \mathbb{R} \to -z \in \mathbb{R}) \to \varphi_{83}$\lbl{A1}}
\BIC{$\varphi_{83}$}
\end{ptree}
\noindent$\Pi_{10}$:
\begin{ptree}
\AXC{$x \in \mathbb{R} \to y \in \mathbb{R} \to z \in \mathbb{R} \to x + z = y + z \to z \in \mathbb{R}$\lbl{A1}}
\AXC{$\stackrel{\Pi_{9}}{\vdots}$}\noLine
\UIC{$\varphi_{83}$}
\BIC{$x \in \mathbb{R} \to y \in \mathbb{R} \to z \in \mathbb{R} \to x + z = y + z \to -z \in \mathbb{R}$}
\end{ptree}
\noindent$\Pi_{11}$:
\begin{ptree}
\AXC{$\varphi_{34}$\lbl{(R2)}}
\AXC{$\varphi_{34} \to \varphi_{41}$\lbl{A2}}
\BIC{$\varphi_{41}$}
\AXC{$\varphi_{41} \to \varphi_{51}$\lbl{A2}}
\BIC{$\varphi_{51}$}
\AXC{$\varphi_{51} \to v_{16} \in \mathbb{R} \to v_{17} \in \mathbb{R} \to \varphi_{7}$\lbl{A2}}
\BIC{$v_{16} \in \mathbb{R} \to v_{17} \in \mathbb{R} \to \varphi_{7}$}
\UIC{$\forall v_{18}\,(v_{16} \in \mathbb{R} \to v_{17} \in \mathbb{R} \to \varphi_{7})$}
\UIC{$\forall v_{17}\,\forall v_{18}\,(v_{16} \in \mathbb{R} \to v_{17} \in \mathbb{R} \to \varphi_{7})$}
\UIC{$\forall v_{16}\,\forall v_{17}\,\forall v_{18}\,(v_{16} \in \mathbb{R} \to v_{17} \in \mathbb{R} \to \varphi_{7})$}
\AXC{$\varphi_{89}$\lbl{A2}}
\BIC{$\varphi_{63}$}
\AXC{$\varphi_{63} \to \varphi_{35}$\lbl{A2}}
\BIC{$\varphi_{35}$}
\end{ptree}
\noindent$\Pi_{12}$:
\begin{ptree}
\AXC{$x \in \mathbb{R} \to y \in \mathbb{R} \to z \in \mathbb{R} \to x + z = y + z \to y \in \mathbb{R}$\lbl{A1}}
\AXC{$\stackrel{\Pi_{11}}{\vdots}$}\noLine
\UIC{$\varphi_{35}$}
\AXC{$\varphi_{35} \to \varphi_{52}$\lbl{A2}}
\BIC{$\varphi_{52}$}
\AXC{$\varphi_{52} \to \varphi_{98}$\lbl{A1}}
\BIC{$\varphi_{98}$}
\BIC{$\varphi_{93}$}
\end{ptree}
\noindent$\Pi_{13}$:
\begin{ptree}
\AXC{$\stackrel{\Pi_{10}}{\vdots}$}\noLine
\UIC{$x \in \mathbb{R} \to y \in \mathbb{R} \to z \in \mathbb{R} \to x + z = y + z \to -z \in \mathbb{R}$}
\AXC{$x \in \mathbb{R} \to y \in \mathbb{R} \to z \in \mathbb{R} \to x + z = y + z \to z \in \mathbb{R}$\lbl{A1}}
\AXC{$\stackrel{\Pi_{12}}{\vdots}$}\noLine
\UIC{$\varphi_{93}$}
\BIC{$\varphi_{86}$}
\BIC{$x \in \mathbb{R} \to \varphi_{47}$}
\end{ptree}
\noindent$\Pi_{14}$:
\begin{ptree}
\AXC{$y = z \to x = y \to x = z$\lbl{A6}}
\AXC{$(y = z \to x = y \to x = z) \to x = y \to y = z \to x = z$\lbl{A1}}
\BIC{$x = y \to y = z \to x = z$}
\UIC{$\forall z\,(x = y \to y = z \to x = z)$}
\UIC{$\forall y\,\forall z\,(x = y \to y = z \to x = z)$}
\UIC{$\forall x\,\forall y\,\forall z\,(x = y \to y = z \to x = z)$}
\AXC{$\varphi_{12}$\lbl{A2}}
\BIC{$\forall y\,\forall z\,(v_{22} = y \to y = z \to v_{22} = z)$}
\AXC{$\varphi_{19}$\lbl{A2}}
\BIC{$\forall z\,(v_{22} = v_{23} \to v_{23} = z \to v_{22} = z)$}
\AXC{$\varphi_{26}$\lbl{A2}}
\BIC{$v_{22} = v_{23} \to v_{23} = v_{24} \to v_{22} = v_{24}$}
\UIC{$\forall v_{24}\,(v_{22} = v_{23} \to v_{23} = v_{24} \to v_{22} = v_{24})$}
\UIC{$\varphi_{3}$}
\UIC{$\forall v_{22}\,(\varphi_{3})$}
\end{ptree}
\noindent$\Pi_{15}$:
\begin{ptree}
\AXC{$\stackrel{\Pi_{14}}{\vdots}$}\noLine
\UIC{$\forall v_{22}\,(\varphi_{3})$}
\AXC{$\varphi_{59}$\lbl{A2}}
\BIC{$\forall v_{23}\,\forall v_{24}\,(x = v_{23} \to v_{23} = v_{24} \to x = v_{24})$}
\AXC{$\varphi_{72}$\lbl{A2}}
\BIC{$\varphi_{20}$}
\AXC{$\varphi_{20} \to \varphi_{45}$\lbl{A2}}
\BIC{$\varphi_{45}$}
\end{ptree}
\noindent$\Pi_{16}$:
\begin{ptree}
\AXC{$x \in \mathbb{R} \to y \in \mathbb{R} \to z \in \mathbb{R} \to x + z = y + z \to x + z = y + z$\lbl{A1}}
\AXC{$\varphi_{97}$\lbl{A1}}
\BIC{$x + z + -z = x + z + -z \to \varphi_{65} \to x \in \mathbb{R} \to \varphi_{48}$}
\end{ptree}
\noindent$\Pi_{17}$:
\begin{ptree}
\AXC{$x = x$\lbl{A5}}
\UIC{$\forall x\,(x = x)$}
\AXC{$\forall x\,(x = x) \to x + z + -z = x + z + -z$\lbl{A2}}
\BIC{$x + z + -z = x + z + -z$}
\AXC{$\stackrel{\Pi_{16}}{\vdots}$}\noLine
\UIC{$x + z + -z = x + z + -z \to \varphi_{65} \to x \in \mathbb{R} \to \varphi_{48}$}
\BIC{$\varphi_{65} \to x \in \mathbb{R} \to \varphi_{48}$}
\end{ptree}
\noindent$\Pi_{18}$:
\begin{ptree}
\AXC{$z_{1} = w \to x + z + -z = z_{1} + -z \to x + z + -z = w + -z$\lbl{A6}}
\UIC{$\varphi_{33}$}
\UIC{$\forall z_{1}\,(\varphi_{33})$}
\AXC{$\forall z_{1}\,(\varphi_{33}) \to \varphi_{60}$\lbl{A2}}
\BIC{$\varphi_{60}$}
\AXC{$\varphi_{60} \to \varphi_{65}$\lbl{A2}}
\BIC{$\varphi_{65}$}
\AXC{$\stackrel{\Pi_{17}}{\vdots}$}\noLine
\UIC{$\varphi_{65} \to x \in \mathbb{R} \to \varphi_{48}$}
\BIC{$x \in \mathbb{R} \to \varphi_{48}$}
\end{ptree}
\noindent$\Pi_{19}$:
\begin{ptree}
\AXC{$y = z \to x = y \to x = z$\lbl{A6}}
\AXC{$(y = z \to x = y \to x = z) \to x = y \to y = z \to x = z$\lbl{A1}}
\BIC{$x = y \to y = z \to x = z$}
\UIC{$\forall z\,(x = y \to y = z \to x = z)$}
\UIC{$\forall y\,\forall z\,(x = y \to y = z \to x = z)$}
\UIC{$\forall x\,\forall y\,\forall z\,(x = y \to y = z \to x = z)$}
\AXC{$\varphi_{13}$\lbl{A2}}
\BIC{$\forall y\,\forall z\,(v_{19} = y \to y = z \to v_{19} = z)$}
\AXC{$\varphi_{21}$\lbl{A2}}
\BIC{$\forall z\,(v_{19} = v_{20} \to v_{20} = z \to v_{19} = z)$}
\AXC{$\varphi_{27}$\lbl{A2}}
\BIC{$v_{19} = v_{20} \to v_{20} = v_{21} \to v_{19} = v_{21}$}
\UIC{$\forall v_{21}\,(v_{19} = v_{20} \to v_{20} = v_{21} \to v_{19} = v_{21})$}
\UIC{$\varphi_{4}$}
\UIC{$\forall v_{19}\,(\varphi_{4})$}
\end{ptree}
\noindent$\Pi_{20}$:
\begin{ptree}
\AXC{$x = x$\lbl{A5}}
\AXC{$x = y \to x = x \to y = x$\lbl{A6}}
\AXC{$(x = y \to x = x \to y = x) \to x = x \to x = y \to y = x$\lbl{A1}}
\BIC{$x = x \to x = y \to y = x$}
\BIC{$x = y \to y = x$}
\UIC{$\forall y\,(x = y \to y = x)$}
\UIC{$\forall x\,\forall y\,(x = y \to y = x)$}
\end{ptree}
\noindent$\Pi_{21}$:
\begin{ptree}
\AXC{$\stackrel{\Pi_{20}}{\vdots}$}\noLine
\UIC{$\forall x\,\forall y\,(x = y \to y = x)$}
\AXC{$\varphi_{29}$\lbl{A2}}
\BIC{$\forall y\,(x + z + -z = y \to y = x + z + -z)$}
\AXC{$\forall y\,(x + z + -z = y \to y = x + z + -z) \to \varphi_{32}$\lbl{A2}}
\BIC{$\varphi_{32}$}
\end{ptree}
\noindent$\Pi_{22}$:
\begin{ptree}
\AXC{$\varphi_{34}$\lbl{(R2)}}
\AXC{$\varphi_{34} \to \varphi_{42}$\lbl{A2}}
\BIC{$\varphi_{42}$}
\AXC{$\varphi_{42} \to \varphi_{53}$\lbl{A2}}
\BIC{$\varphi_{53}$}
\AXC{$\varphi_{53} \to v_{13} \in \mathbb{R} \to v_{14} \in \mathbb{R} \to \varphi_{8}$\lbl{A2}}
\BIC{$v_{13} \in \mathbb{R} \to v_{14} \in \mathbb{R} \to \varphi_{8}$}
\UIC{$\forall v_{15}\,(v_{13} \in \mathbb{R} \to v_{14} \in \mathbb{R} \to \varphi_{8})$}
\UIC{$\forall v_{14}\,\forall v_{15}\,(v_{13} \in \mathbb{R} \to v_{14} \in \mathbb{R} \to \varphi_{8})$}
\UIC{$\forall v_{13}\,\forall v_{14}\,\forall v_{15}\,(v_{13} \in \mathbb{R} \to v_{14} \in \mathbb{R} \to \varphi_{8})$}
\AXC{$\varphi_{90}$\lbl{A2}}
\BIC{$\varphi_{64}$}
\AXC{$\varphi_{64} \to \varphi_{36}$\lbl{A2}}
\BIC{$\varphi_{36}$}
\end{ptree}
\noindent$\Pi_{23}$:
\begin{ptree}
\AXC{$x \in \mathbb{R} \to y \in \mathbb{R} \to z \in \mathbb{R} \to x + z = y + z \to x \in \mathbb{R}$\lbl{A1}}
\AXC{$\stackrel{\Pi_{22}}{\vdots}$}\noLine
\UIC{$\varphi_{36}$}
\AXC{$\varphi_{36} \to \varphi_{54}$\lbl{A2}}
\BIC{$\varphi_{54}$}
\AXC{$\varphi_{54} \to \varphi_{99}$\lbl{A1}}
\BIC{$\varphi_{99}$}
\BIC{$\varphi_{94}$}
\end{ptree}
\noindent$\Pi_{24}$:
\begin{ptree}
\AXC{$\stackrel{\Pi_{10}}{\vdots}$}\noLine
\UIC{$x \in \mathbb{R} \to y \in \mathbb{R} \to z \in \mathbb{R} \to x + z = y + z \to -z \in \mathbb{R}$}
\AXC{$x \in \mathbb{R} \to y \in \mathbb{R} \to z \in \mathbb{R} \to x + z = y + z \to z \in \mathbb{R}$\lbl{A1}}
\AXC{$\stackrel{\Pi_{23}}{\vdots}$}\noLine
\UIC{$\varphi_{94}$}
\BIC{$\varphi_{87}$}
\BIC{$x \in \mathbb{R} \to \varphi_{49}$}
\end{ptree}
\noindent$\Pi_{25}$:
\begin{ptree}
\AXC{$y = z \to x = y \to x = z$\lbl{A6}}
\AXC{$(y = z \to x = y \to x = z) \to x = y \to y = z \to x = z$\lbl{A1}}
\BIC{$x = y \to y = z \to x = z$}
\UIC{$\forall z\,(x = y \to y = z \to x = z)$}
\UIC{$\forall y\,\forall z\,(x = y \to y = z \to x = z)$}
\UIC{$\forall x\,\forall y\,\forall z\,(x = y \to y = z \to x = z)$}
\AXC{$\varphi_{5}$\lbl{A2}}
\BIC{$\forall y\,\forall z\,(x = y \to y = z \to x = z)$}
\AXC{$\varphi_{14}$\lbl{A2}}
\BIC{$\forall z\,(x = x + 0 \to x + 0 = z \to x = z)$}
\AXC{$\varphi_{68}$\lbl{A2}}
\BIC{$x = x + 0 \to x + 0 = x + (z + -z) \to x = x + (z + -z)$}
\end{ptree}
\noindent$\Pi_{26}$:
\begin{ptree}
\AXC{$z_{1} = w \to x + 0 = x + z_{1} \to x + 0 = x + w$\lbl{A6}}
\UIC{$\forall w\,(z_{1} = w \to x + 0 = x + z_{1} \to x + 0 = x + w)$}
\UIC{$\forall z_{1}\,\forall w\,(z_{1} = w \to x + 0 = x + z_{1} \to x + 0 = x + w)$}
\AXC{$\varphi_{66}$\lbl{A2}}
\BIC{$\forall w\,(0 = w \to x + 0 = x + 0 \to x + 0 = x + w)$}
\AXC{$\varphi_{74}$\lbl{A2}}
\BIC{$0 = z + -z \to x + 0 = x + 0 \to x + 0 = x + (z + -z)$}
\end{ptree}
\noindent$\Pi_{27}$:
\begin{ptree}
\AXC{$\forall x \in \mathbb{R}\,(x + -x = 0)$\lbl{(R5)}}
\AXC{$\forall x \in \mathbb{R}\,(x + -x = 0) \to z \in \mathbb{R} \to z + -z = 0$\lbl{A2}}
\BIC{$z \in \mathbb{R} \to z + -z = 0$}
\AXC{$(z \in \mathbb{R} \to z + -z = 0) \to \varphi_{81}$\lbl{A1}}
\BIC{$\varphi_{81}$}
\end{ptree}
\noindent$\Pi_{28}$:
\begin{ptree}
\AXC{$x \in \mathbb{R} \to y \in \mathbb{R} \to z \in \mathbb{R} \to x + z = y + z \to z \in \mathbb{R}$\lbl{A1}}
\AXC{$\stackrel{\Pi_{27}}{\vdots}$}\noLine
\UIC{$\varphi_{81}$}
\BIC{$x \in \mathbb{R} \to y \in \mathbb{R} \to z \in \mathbb{R} \to x + z = y + z \to z + -z = 0$}
\AXC{$\varphi_{88}$\lbl{A1}}
\BIC{$\varphi_{75}$}
\end{ptree}
\noindent$\Pi_{29}$:
\begin{ptree}
\AXC{$\stackrel{\Pi_{20}}{\vdots}$}\noLine
\UIC{$\forall x\,\forall y\,(x = y \to y = x)$}
\AXC{$\varphi_{6}$\lbl{A2}}
\BIC{$\forall y\,(z + -z = y \to y = z + -z)$}
\AXC{$\varphi_{30}$\lbl{A2}}
\BIC{$z + -z = 0 \to 0 = z + -z$}
\AXC{$\stackrel{\Pi_{28}}{\vdots}$}\noLine
\UIC{$\varphi_{75}$}
\BIC{$x \in \mathbb{R} \to y \in \mathbb{R} \to z \in \mathbb{R} \to x + z = y + z \to 0 = z + -z$}
\AXC{$\varphi_{92}$\lbl{A1}}
\BIC{$x + 0 = x + 0 \to \varphi_{82}$}
\end{ptree}
\noindent$\Pi_{30}$:
\begin{ptree}
\AXC{$\stackrel{\Pi_{26}}{\vdots}$}\noLine
\UIC{$0 = z + -z \to x + 0 = x + 0 \to x + 0 = x + (z + -z)$}
\AXC{$x = x$\lbl{A5}}
\UIC{$\forall x\,(x = x)$}
\AXC{$\forall x\,(x = x) \to x + 0 = x + 0$\lbl{A2}}
\BIC{$x + 0 = x + 0$}
\AXC{$\stackrel{\Pi_{29}}{\vdots}$}\noLine
\UIC{$x + 0 = x + 0 \to \varphi_{82}$}
\BIC{$\varphi_{82}$}
\BIC{$\varphi_{62}$}
\end{ptree}
\noindent$\Pi_{31}$:
\begin{ptree}
\AXC{$\stackrel{\Pi_{20}}{\vdots}$}\noLine
\UIC{$\forall x\,\forall y\,(x = y \to y = x)$}
\AXC{$\forall x\,\forall y\,(x = y \to y = x) \to \forall y\,(x + 0 = y \to y = x + 0)$\lbl{A2}}
\BIC{$\forall y\,(x + 0 = y \to y = x + 0)$}
\end{ptree}
\noindent$\Pi_{32}$:
\begin{ptree}
\AXC{$\stackrel{\Pi_{31}}{\vdots}$}\noLine
\UIC{$\forall y\,(x + 0 = y \to y = x + 0)$}
\AXC{$\forall y\,(x + 0 = y \to y = x + 0) \to x + 0 = x \to x = x + 0$\lbl{A2}}
\BIC{$x + 0 = x \to x = x + 0$}
\end{ptree}
\noindent$\Pi_{33}$:
\begin{ptree}
\AXC{$\forall x \in \mathbb{R}\,(x + 0 = x)$\lbl{(R4)}}
\AXC{$\forall x \in \mathbb{R}\,(x + 0 = x) \to x \in \mathbb{R} \to x + 0 = x$\lbl{A2}}
\BIC{$x \in \mathbb{R} \to x + 0 = x$}
\AXC{$(x \in \mathbb{R} \to x + 0 = x) \to \varphi_{79}$\lbl{A1}}
\BIC{$\varphi_{79}$}
\end{ptree}
\noindent$\Pi_{34}$:
\begin{ptree}
\AXC{$\stackrel{\Pi_{32}}{\vdots}$}\noLine
\UIC{$x + 0 = x \to x = x + 0$}
\AXC{$x \in \mathbb{R} \to y \in \mathbb{R} \to z \in \mathbb{R} \to x + z = y + z \to x \in \mathbb{R}$\lbl{A1}}
\AXC{$\stackrel{\Pi_{33}}{\vdots}$}\noLine
\UIC{$\varphi_{79}$}
\BIC{$x \in \mathbb{R} \to y \in \mathbb{R} \to z \in \mathbb{R} \to x + z = y + z \to x + 0 = x$}
\AXC{$\varphi_{84}$\lbl{A1}}
\BIC{$\varphi_{67}$}
\BIC{$x \in \mathbb{R} \to y \in \mathbb{R} \to z \in \mathbb{R} \to x + z = y + z \to x = x + 0$}
\AXC{$\varphi_{96}$\lbl{A1}}
\BIC{$\varphi_{62} \to \varphi_{80}$}
\end{ptree}
\noindent$\Pi_{35}$:
\begin{ptree}
\AXC{$\stackrel{\Pi_{25}}{\vdots}$}\noLine
\UIC{$x = x + 0 \to x + 0 = x + (z + -z) \to x = x + (z + -z)$}
\AXC{$\stackrel{\Pi_{30}}{\vdots}$}\noLine
\UIC{$\varphi_{62}$}
\AXC{$\stackrel{\Pi_{34}}{\vdots}$}\noLine
\UIC{$\varphi_{62} \to \varphi_{80}$}
\BIC{$\varphi_{80}$}
\BIC{$\varphi_{37}$}
\AXC{$\varphi_{37} \to (x \in \mathbb{R} \to \varphi_{50}) \to \varphi_{46} \to \varphi_{38}$\lbl{A1}}
\BIC{$(x \in \mathbb{R} \to \varphi_{50}) \to \varphi_{46} \to \varphi_{38}$}
\end{ptree}
\noindent$\Pi_{36}$:
\begin{ptree}
\AXC{$\stackrel{\Pi_{21}}{\vdots}$}\noLine
\UIC{$\varphi_{32}$}
\AXC{$\stackrel{\Pi_{24}}{\vdots}$}\noLine
\UIC{$x \in \mathbb{R} \to \varphi_{49}$}
\AXC{$(x \in \mathbb{R} \to \varphi_{49}) \to \varphi_{32} \to x \in \mathbb{R} \to \varphi_{50}$\lbl{A1}}
\BIC{$\varphi_{32} \to x \in \mathbb{R} \to \varphi_{50}$}
\BIC{$x \in \mathbb{R} \to \varphi_{50}$}
\AXC{$\stackrel{\Pi_{35}}{\vdots}$}\noLine
\UIC{$(x \in \mathbb{R} \to \varphi_{50}) \to \varphi_{46} \to \varphi_{38}$}
\BIC{$\varphi_{46} \to \varphi_{38}$}
\end{ptree}
\noindent$\Pi_{37}$:
\begin{ptree}
\AXC{$\stackrel{\Pi_{19}}{\vdots}$}\noLine
\UIC{$\forall v_{19}\,(\varphi_{4})$}
\AXC{$\varphi_{61}$\lbl{A2}}
\BIC{$\forall v_{20}\,\forall v_{21}\,(x = v_{20} \to v_{20} = v_{21} \to x = v_{21})$}
\AXC{$\varphi_{73}$\lbl{A2}}
\BIC{$\varphi_{22}$}
\AXC{$\varphi_{22} \to \varphi_{46}$\lbl{A2}}
\BIC{$\varphi_{46}$}
\AXC{$\stackrel{\Pi_{36}}{\vdots}$}\noLine
\UIC{$\varphi_{46} \to \varphi_{38}$}
\BIC{$\varphi_{38}$}
\end{ptree}
\noindent$\Pi_{38}$:
\begin{ptree}
\AXC{$\stackrel{\Pi_{15}}{\vdots}$}\noLine
\UIC{$\varphi_{45}$}
\AXC{$\stackrel{\Pi_{18}}{\vdots}$}\noLine
\UIC{$x \in \mathbb{R} \to \varphi_{48}$}
\AXC{$\stackrel{\Pi_{37}}{\vdots}$}\noLine
\UIC{$\varphi_{38}$}
\AXC{$\varphi_{38} \to (x \in \mathbb{R} \to \varphi_{48}) \to \varphi_{45} \to \varphi_{39}$\lbl{A1}}
\BIC{$(x \in \mathbb{R} \to \varphi_{48}) \to \varphi_{45} \to \varphi_{39}$}
\BIC{$\varphi_{45} \to \varphi_{39}$}
\BIC{$\varphi_{39}$}
\AXC{$\varphi_{39} \to (x \in \mathbb{R} \to \varphi_{47}) \to \varphi_{44} \to \varphi_{40}$\lbl{A1}}
\BIC{$(x \in \mathbb{R} \to \varphi_{47}) \to \varphi_{44} \to \varphi_{40}$}
\end{ptree}
\noindent$\Pi_{39}$:
\begin{ptree}
\AXC{$\stackrel{\Pi_{8}}{\vdots}$}\noLine
\UIC{$\varphi_{44}$}
\AXC{$\stackrel{\Pi_{13}}{\vdots}$}\noLine
\UIC{$x \in \mathbb{R} \to \varphi_{47}$}
\AXC{$\stackrel{\Pi_{38}}{\vdots}$}\noLine
\UIC{$(x \in \mathbb{R} \to \varphi_{47}) \to \varphi_{44} \to \varphi_{40}$}
\BIC{$\varphi_{44} \to \varphi_{40}$}
\BIC{$\varphi_{40}$}
\AXC{$\varphi_{40} \to \varphi_{57} \to \varphi_{77}$\lbl{A1}}
\BIC{$\varphi_{57} \to \varphi_{77}$}
\end{ptree}
\noindent$\Pi_{40}$:
\begin{ptree}
\AXC{$\varphi_{28}$\lbl{A6}}
\UIC{$\forall w\,(\varphi_{28})$}
\UIC{$\forall z_{1}\,\forall w\,(\varphi_{28})$}
\AXC{$\forall z_{1}\,\forall w\,(\varphi_{28}) \to \varphi_{56}$\lbl{A2}}
\BIC{$\varphi_{56}$}
\AXC{$\varphi_{56} \to \varphi_{43}$\lbl{A2}}
\BIC{$\varphi_{43}$}
\AXC{$\stackrel{\Pi_{6}}{\vdots}$}\noLine
\UIC{$\varphi_{43} \to \varphi_{57}$}
\BIC{$\varphi_{57}$}
\AXC{$\stackrel{\Pi_{39}}{\vdots}$}\noLine
\UIC{$\varphi_{57} \to \varphi_{77}$}
\BIC{$\varphi_{77}$}
\end{ptree}
\noindent$\Pi_{41}$:
\begin{ptree}
\AXC{$\stackrel{\Pi_{3}}{\vdots}$}\noLine
\UIC{$\forall v_{28}\,(\varphi_{1})$}
\AXC{$\varphi_{55}$\lbl{A2}}
\BIC{$\forall v_{29}\,\forall v_{30}\,(x = v_{29} \to v_{29} = v_{30} \to x = v_{30})$}
\AXC{$\varphi_{70}$\lbl{A2}}
\BIC{$\varphi_{16}$}
\AXC{$\varphi_{76}$\lbl{A2}}
\BIC{$x = y + (z + -z) \to y + (z + -z) = y + 0 \to x = y + 0$}
\AXC{$\stackrel{\Pi_{40}}{\vdots}$}\noLine
\UIC{$\varphi_{77}$}
\BIC{$x \in \mathbb{R} \to y \in \mathbb{R} \to z \in \mathbb{R} \to x + z = y + z \to x = y + 0$}
\AXC{$\varphi_{91}$\lbl{A1}}
\BIC{$\varphi_{85}$}
\end{ptree}
\noindent 主木:
\begin{ptree}
\AXC{$\stackrel{\Pi_{1}}{\vdots}$}\noLine
\UIC{$x = y + 0 \to y + 0 = y \to x = y$}
\AXC{$x \in \mathbb{R} \to y \in \mathbb{R} \to z \in \mathbb{R} \to x + z = y + z \to y \in \mathbb{R}$\lbl{A1}}
\AXC{$\stackrel{\Pi_{2}}{\vdots}$}\noLine
\UIC{$\varphi_{78}$}
\BIC{$x \in \mathbb{R} \to y \in \mathbb{R} \to z \in \mathbb{R} \to x + z = y + z \to y + 0 = y$}
\AXC{$\stackrel{\Pi_{41}}{\vdots}$}\noLine
\UIC{$\varphi_{85}$}
\BIC{$\varphi_{69}$}
\BIC{$x \in \mathbb{R} \to y \in \mathbb{R} \to z \in \mathbb{R} \to x + z = y + z \to x = y$}
\UIC{$\forall z\,(x \in \mathbb{R} \to y \in \mathbb{R} \to z \in \mathbb{R} \to x + z = y + z \to x = y)$}
\UIC{$\varphi_{31}$}
\UIC{$\forall x\,(\varphi_{31})$}
\end{ptree}
\begin{abbrevlist}
\item $\varphi_{1} :\equiv \forall v_{29}\,\forall v_{30}\,(v_{28} = v_{29} \to v_{29} = v_{30} \to v_{28} = v_{30})$
\item $\varphi_{2} :\equiv \forall v_{26}\,\forall v_{27}\,(v_{25} = v_{26} \to v_{26} = v_{27} \to v_{25} = v_{27})$
\item $\varphi_{3} :\equiv \forall v_{23}\,\forall v_{24}\,(v_{22} = v_{23} \to v_{23} = v_{24} \to v_{22} = v_{24})$
\item $\varphi_{4} :\equiv \forall v_{20}\,\forall v_{21}\,(v_{19} = v_{20} \to v_{20} = v_{21} \to v_{19} = v_{21})$
\item $\varphi_{5} :\equiv \forall x\,\forall y\,\forall z\,(x = y \to y = z \to x = z) \to \forall y\,\forall z\,(x = y \to y = z \to x = z)$
\item $\varphi_{6} :\equiv \forall x\,\forall y\,(x = y \to y = x) \to \forall y\,(z + -z = y \to y = z + -z)$
\item $\varphi_{7} :\equiv v_{18} \in \mathbb{R} \to v_{16} + v_{17} + v_{18} = v_{16} + (v_{17} + v_{18})$
\item $\varphi_{8} :\equiv v_{15} \in \mathbb{R} \to v_{13} + v_{14} + v_{15} = v_{13} + (v_{14} + v_{15})$
\item $\varphi_{9} :\equiv \forall y\,\forall z\,(x = y \to y = z \to x = z) \to \forall z\,(x = y + 0 \to y + 0 = z \to x = z)$
\item $\varphi_{10} :\equiv \forall x\,\forall y\,\forall z\,(x = y \to y = z \to x = z) \to \forall y\,\forall z\,(v_{28} = y \to y = z \to v_{28} = z)$
\item $\varphi_{11} :\equiv \forall x\,\forall y\,\forall z\,(x = y \to y = z \to x = z) \to \forall y\,\forall z\,(v_{25} = y \to y = z \to v_{25} = z)$
\item $\varphi_{12} :\equiv \forall x\,\forall y\,\forall z\,(x = y \to y = z \to x = z) \to \forall y\,\forall z\,(v_{22} = y \to y = z \to v_{22} = z)$
\item $\varphi_{13} :\equiv \forall x\,\forall y\,\forall z\,(x = y \to y = z \to x = z) \to \forall y\,\forall z\,(v_{19} = y \to y = z \to v_{19} = z)$
\item $\varphi_{14} :\equiv \forall y\,\forall z\,(x = y \to y = z \to x = z) \to \forall z\,(x = x + 0 \to x + 0 = z \to x = z)$
\item $\varphi_{15} :\equiv \forall y\,\forall z\,(v_{28} = y \to y = z \to v_{28} = z) \to \forall z\,(v_{28} = v_{29} \to v_{29} = z \to v_{28} = z)$
\item $\varphi_{16} :\equiv \forall v_{30}\,(x = y + (z + -z) \to y + (z + -z) = v_{30} \to x = v_{30})$
\item $\varphi_{17} :\equiv \forall y\,\forall z\,(v_{25} = y \to y = z \to v_{25} = z) \to \forall z\,(v_{25} = v_{26} \to v_{26} = z \to v_{25} = z)$
\item $\varphi_{18} :\equiv \forall v_{27}\,(x = y + z + -z \to y + z + -z = v_{27} \to x = v_{27})$
\item $\varphi_{19} :\equiv \forall y\,\forall z\,(v_{22} = y \to y = z \to v_{22} = z) \to \forall z\,(v_{22} = v_{23} \to v_{23} = z \to v_{22} = z)$
\item $\varphi_{20} :\equiv \forall v_{24}\,(x = x + z + -z \to x + z + -z = v_{24} \to x = v_{24})$
\item $\varphi_{21} :\equiv \forall y\,\forall z\,(v_{19} = y \to y = z \to v_{19} = z) \to \forall z\,(v_{19} = v_{20} \to v_{20} = z \to v_{19} = z)$
\item $\varphi_{22} :\equiv \forall v_{21}\,(x = x + (z + -z) \to x + (z + -z) = v_{21} \to x = v_{21})$
\item $\varphi_{23} :\equiv \forall z\,(x = y + 0 \to y + 0 = z \to x = z) \to x = y + 0 \to y + 0 = y \to x = y$
\item $\varphi_{24} :\equiv \forall z\,(v_{28} = v_{29} \to v_{29} = z \to v_{28} = z) \to v_{28} = v_{29} \to v_{29} = v_{30} \to v_{28} = v_{30}$
\item $\varphi_{25} :\equiv \forall z\,(v_{25} = v_{26} \to v_{26} = z \to v_{25} = z) \to v_{25} = v_{26} \to v_{26} = v_{27} \to v_{25} = v_{27}$
\item $\varphi_{26} :\equiv \forall z\,(v_{22} = v_{23} \to v_{23} = z \to v_{22} = z) \to v_{22} = v_{23} \to v_{23} = v_{24} \to v_{22} = v_{24}$
\item $\varphi_{27} :\equiv \forall z\,(v_{19} = v_{20} \to v_{20} = z \to v_{19} = z) \to v_{19} = v_{20} \to v_{20} = v_{21} \to v_{19} = v_{21}$
\item $\varphi_{28} :\equiv z_{1} = w \to y + (z + -z) = y + z_{1} \to y + (z + -z) = y + w$
\item $\varphi_{29} :\equiv \forall x\,\forall y\,(x = y \to y = x) \to \forall y\,(x + z + -z = y \to y = x + z + -z)$
\item $\varphi_{30} :\equiv \forall y\,(z + -z = y \to y = z + -z) \to z + -z = 0 \to 0 = z + -z$
\item $\varphi_{31} :\equiv \forall y\,\forall z\,(x \in \mathbb{R} \to y \in \mathbb{R} \to z \in \mathbb{R} \to x + z = y + z \to x = y)$
\item $\varphi_{32} :\equiv x + z + -z = x + (z + -z) \to x + (z + -z) = x + z + -z$
\item $\varphi_{33} :\equiv \forall w\,(z_{1} = w \to x + z + -z = z_{1} + -z \to x + z + -z = w + -z)$
\item $\varphi_{34} :\equiv \forall x\,\forall y\,\forall z\,(x \in \mathbb{R} \to y \in \mathbb{R} \to z \in \mathbb{R} \to x + y + z = x + (y + z))$
\item $\varphi_{35} :\equiv \forall v_{18}\,(y \in \mathbb{R} \to z \in \mathbb{R} \to v_{18} \in \mathbb{R} \to y + z + v_{18} = y + (z + v_{18}))$
\item $\varphi_{36} :\equiv \forall v_{15}\,(x \in \mathbb{R} \to z \in \mathbb{R} \to v_{15} \in \mathbb{R} \to x + z + v_{15} = x + (z + v_{15}))$
\item $\varphi_{37} :\equiv x \in \mathbb{R} \to y \in \mathbb{R} \to z \in \mathbb{R} \to x + z = y + z \to x = x + (z + -z)$
\item $\varphi_{38} :\equiv x \in \mathbb{R} \to y \in \mathbb{R} \to z \in \mathbb{R} \to x + z = y + z \to x = x + z + -z$
\item $\varphi_{39} :\equiv x \in \mathbb{R} \to y \in \mathbb{R} \to z \in \mathbb{R} \to x + z = y + z \to x = y + z + -z$
\item $\varphi_{40} :\equiv x \in \mathbb{R} \to y \in \mathbb{R} \to z \in \mathbb{R} \to x + z = y + z \to x = y + (z + -z)$
\item $\varphi_{41} :\equiv \forall y\,\forall z\,(v_{16} \in \mathbb{R} \to y \in \mathbb{R} \to z \in \mathbb{R} \to v_{16} + y + z = v_{16} + (y + z))$
\item $\varphi_{42} :\equiv \forall y\,\forall z\,(v_{13} \in \mathbb{R} \to y \in \mathbb{R} \to z \in \mathbb{R} \to v_{13} + y + z = v_{13} + (y + z))$
\item $\varphi_{43} :\equiv z + -z = 0 \to y + (z + -z) = y + (z + -z) \to y + (z + -z) = y + 0$
\item $\varphi_{44} :\equiv x = y + z + -z \to y + z + -z = y + (z + -z) \to x = y + (z + -z)$
\item $\varphi_{45} :\equiv x = x + z + -z \to x + z + -z = y + z + -z \to x = y + z + -z$
\item $\varphi_{46} :\equiv x = x + (z + -z) \to x + (z + -z) = x + z + -z \to x = x + z + -z$
\item $\varphi_{47} :\equiv y \in \mathbb{R} \to z \in \mathbb{R} \to x + z = y + z \to y + z + -z = y + (z + -z)$
\item $\varphi_{48} :\equiv y \in \mathbb{R} \to z \in \mathbb{R} \to x + z = y + z \to x + z + -z = y + z + -z$
\item $\varphi_{49} :\equiv y \in \mathbb{R} \to z \in \mathbb{R} \to x + z = y + z \to x + z + -z = x + (z + -z)$
\item $\varphi_{50} :\equiv y \in \mathbb{R} \to z \in \mathbb{R} \to x + z = y + z \to x + (z + -z) = x + z + -z$
\item $\varphi_{51} :\equiv \forall z\,(v_{16} \in \mathbb{R} \to v_{17} \in \mathbb{R} \to z \in \mathbb{R} \to v_{16} + v_{17} + z = v_{16} + (v_{17} + z))$
\item $\varphi_{52} :\equiv y \in \mathbb{R} \to z \in \mathbb{R} \to -z \in \mathbb{R} \to y + z + -z = y + (z + -z)$
\item $\varphi_{53} :\equiv \forall z\,(v_{13} \in \mathbb{R} \to v_{14} \in \mathbb{R} \to z \in \mathbb{R} \to v_{13} + v_{14} + z = v_{13} + (v_{14} + z))$
\item $\varphi_{54} :\equiv x \in \mathbb{R} \to z \in \mathbb{R} \to -z \in \mathbb{R} \to x + z + -z = x + (z + -z)$
\item $\varphi_{55} :\equiv \forall v_{28}\,(\varphi_{1}) \to \forall v_{29}\,\forall v_{30}\,(x = v_{29} \to v_{29} = v_{30} \to x = v_{30})$
\item $\varphi_{56} :\equiv \forall w\,(z + -z = w \to y + (z + -z) = y + (z + -z) \to y + (z + -z) = y + w)$
\item $\varphi_{57} :\equiv x \in \mathbb{R} \to y \in \mathbb{R} \to z \in \mathbb{R} \to x + z = y + z \to y + (z + -z) = y + 0$
\item $\varphi_{58} :\equiv \forall v_{25}\,(\varphi_{2}) \to \forall v_{26}\,\forall v_{27}\,(x = v_{26} \to v_{26} = v_{27} \to x = v_{27})$
\item $\varphi_{59} :\equiv \forall v_{22}\,(\varphi_{3}) \to \forall v_{23}\,\forall v_{24}\,(x = v_{23} \to v_{23} = v_{24} \to x = v_{24})$
\item $\varphi_{60} :\equiv \forall w\,(x + z = w \to x + z + -z = x + z + -z \to x + z + -z = w + -z)$
\item $\varphi_{61} :\equiv \forall v_{19}\,(\varphi_{4}) \to \forall v_{20}\,\forall v_{21}\,(x = v_{20} \to v_{20} = v_{21} \to x = v_{21})$
\item $\varphi_{62} :\equiv x \in \mathbb{R} \to y \in \mathbb{R} \to z \in \mathbb{R} \to x + z = y + z \to x + 0 = x + (z + -z)$
\item $\varphi_{63} :\equiv \forall v_{17}\,\forall v_{18}\,(y \in \mathbb{R} \to v_{17} \in \mathbb{R} \to v_{18} \in \mathbb{R} \to y + v_{17} + v_{18} = y + (v_{17} + v_{18}))$
\item $\varphi_{64} :\equiv \forall v_{14}\,\forall v_{15}\,(x \in \mathbb{R} \to v_{14} \in \mathbb{R} \to v_{15} \in \mathbb{R} \to x + v_{14} + v_{15} = x + (v_{14} + v_{15}))$
\item $\varphi_{65} :\equiv x + z = y + z \to x + z + -z = x + z + -z \to x + z + -z = y + z + -z$
\item $\varphi_{66} :\equiv \forall z_{1}\,\forall w\,(z_{1} = w \to x + 0 = x + z_{1} \to x + 0 = x + w) \to \forall w\,(0 = w \to x + 0 = x + 0 \to x + 0 = x + w)$
\item $\varphi_{67} :\equiv (x + 0 = x \to x = x + 0) \to x \in \mathbb{R} \to y \in \mathbb{R} \to z \in \mathbb{R} \to x + z = y + z \to x = x + 0$
\item $\varphi_{68} :\equiv \forall z\,(x = x + 0 \to x + 0 = z \to x = z) \to x = x + 0 \to x + 0 = x + (z + -z) \to x = x + (z + -z)$
\item $\varphi_{69} :\equiv (x = y + 0 \to y + 0 = y \to x = y) \to x \in \mathbb{R} \to y \in \mathbb{R} \to z \in \mathbb{R} \to x + z = y + z \to x = y$
\item $\varphi_{70} :\equiv \forall v_{29}\,\forall v_{30}\,(x = v_{29} \to v_{29} = v_{30} \to x = v_{30}) \to \varphi_{16}$
\item $\varphi_{71} :\equiv \forall v_{26}\,\forall v_{27}\,(x = v_{26} \to v_{26} = v_{27} \to x = v_{27}) \to \varphi_{18}$
\item $\varphi_{72} :\equiv \forall v_{23}\,\forall v_{24}\,(x = v_{23} \to v_{23} = v_{24} \to x = v_{24}) \to \varphi_{20}$
\item $\varphi_{73} :\equiv \forall v_{20}\,\forall v_{21}\,(x = v_{20} \to v_{20} = v_{21} \to x = v_{21}) \to \varphi_{22}$
\item $\varphi_{74} :\equiv \forall w\,(0 = w \to x + 0 = x + 0 \to x + 0 = x + w) \to 0 = z + -z \to x + 0 = x + 0 \to x + 0 = x + (z + -z)$
\item $\varphi_{75} :\equiv (z + -z = 0 \to 0 = z + -z) \to x \in \mathbb{R} \to y \in \mathbb{R} \to z \in \mathbb{R} \to x + z = y + z \to 0 = z + -z$
\item $\varphi_{76} :\equiv \varphi_{16} \to x = y + (z + -z) \to y + (z + -z) = y + 0 \to x = y + 0$
\item $\varphi_{77} :\equiv (x = y + (z + -z) \to y + (z + -z) = y + 0 \to x = y + 0) \to x \in \mathbb{R} \to y \in \mathbb{R} \to z \in \mathbb{R} \to x + z = y + z \to x = y + 0$
\item $\varphi_{78} :\equiv (x \in \mathbb{R} \to y \in \mathbb{R} \to z \in \mathbb{R} \to x + z = y + z \to y \in \mathbb{R}) \to x \in \mathbb{R} \to y \in \mathbb{R} \to z \in \mathbb{R} \to x + z = y + z \to y + 0 = y$
\item $\varphi_{79} :\equiv (x \in \mathbb{R} \to y \in \mathbb{R} \to z \in \mathbb{R} \to x + z = y + z \to x \in \mathbb{R}) \to x \in \mathbb{R} \to y \in \mathbb{R} \to z \in \mathbb{R} \to x + z = y + z \to x + 0 = x$
\item $\varphi_{80} :\equiv (x = x + 0 \to x + 0 = x + (z + -z) \to x = x + (z + -z)) \to \varphi_{37}$
\item $\varphi_{81} :\equiv (x \in \mathbb{R} \to y \in \mathbb{R} \to z \in \mathbb{R} \to x + z = y + z \to z \in \mathbb{R}) \to x \in \mathbb{R} \to y \in \mathbb{R} \to z \in \mathbb{R} \to x + z = y + z \to z + -z = 0$
\item $\varphi_{82} :\equiv (0 = z + -z \to x + 0 = x + 0 \to x + 0 = x + (z + -z)) \to \varphi_{62}$
\item $\varphi_{83} :\equiv (x \in \mathbb{R} \to y \in \mathbb{R} \to z \in \mathbb{R} \to x + z = y + z \to z \in \mathbb{R}) \to x \in \mathbb{R} \to y \in \mathbb{R} \to z \in \mathbb{R} \to x + z = y + z \to -z \in \mathbb{R}$
\item $\varphi_{84} :\equiv (x \in \mathbb{R} \to y \in \mathbb{R} \to z \in \mathbb{R} \to x + z = y + z \to x + 0 = x) \to \varphi_{67}$
\item $\varphi_{85} :\equiv (x \in \mathbb{R} \to y \in \mathbb{R} \to z \in \mathbb{R} \to x + z = y + z \to y + 0 = y) \to \varphi_{69}$
\item $\varphi_{86} :\equiv (x \in \mathbb{R} \to y \in \mathbb{R} \to z \in \mathbb{R} \to x + z = y + z \to -z \in \mathbb{R}) \to x \in \mathbb{R} \to \varphi_{47}$
\item $\varphi_{87} :\equiv (x \in \mathbb{R} \to y \in \mathbb{R} \to z \in \mathbb{R} \to x + z = y + z \to -z \in \mathbb{R}) \to x \in \mathbb{R} \to \varphi_{49}$
\item $\varphi_{88} :\equiv (x \in \mathbb{R} \to y \in \mathbb{R} \to z \in \mathbb{R} \to x + z = y + z \to z + -z = 0) \to \varphi_{75}$
\item $\varphi_{89} :\equiv \forall v_{16}\,\forall v_{17}\,\forall v_{18}\,(v_{16} \in \mathbb{R} \to v_{17} \in \mathbb{R} \to \varphi_{7}) \to \varphi_{63}$
\item $\varphi_{90} :\equiv \forall v_{13}\,\forall v_{14}\,\forall v_{15}\,(v_{13} \in \mathbb{R} \to v_{14} \in \mathbb{R} \to \varphi_{8}) \to \varphi_{64}$
\item $\varphi_{91} :\equiv (x \in \mathbb{R} \to y \in \mathbb{R} \to z \in \mathbb{R} \to x + z = y + z \to x = y + 0) \to \varphi_{85}$
\item $\varphi_{92} :\equiv (x \in \mathbb{R} \to y \in \mathbb{R} \to z \in \mathbb{R} \to x + z = y + z \to 0 = z + -z) \to x + 0 = x + 0 \to \varphi_{82}$
\item $\varphi_{93} :\equiv (x \in \mathbb{R} \to y \in \mathbb{R} \to z \in \mathbb{R} \to x + z = y + z \to z \in \mathbb{R}) \to \varphi_{86}$
\item $\varphi_{94} :\equiv (x \in \mathbb{R} \to y \in \mathbb{R} \to z \in \mathbb{R} \to x + z = y + z \to z \in \mathbb{R}) \to \varphi_{87}$
\item $\varphi_{95} :\equiv (x \in \mathbb{R} \to y \in \mathbb{R} \to z \in \mathbb{R} \to x + z = y + z \to z + -z = 0) \to y + (z + -z) = y + (z + -z) \to \varphi_{43} \to \varphi_{57}$
\item $\varphi_{96} :\equiv (x \in \mathbb{R} \to y \in \mathbb{R} \to z \in \mathbb{R} \to x + z = y + z \to x = x + 0) \to \varphi_{62} \to \varphi_{80}$
\item $\varphi_{97} :\equiv (x \in \mathbb{R} \to y \in \mathbb{R} \to z \in \mathbb{R} \to x + z = y + z \to x + z = y + z) \to x + z + -z = x + z + -z \to \varphi_{65} \to x \in \mathbb{R} \to \varphi_{48}$
\item $\varphi_{98} :\equiv (x \in \mathbb{R} \to y \in \mathbb{R} \to z \in \mathbb{R} \to x + z = y + z \to y \in \mathbb{R}) \to \varphi_{93}$
\item $\varphi_{99} :\equiv (x \in \mathbb{R} \to y \in \mathbb{R} \to z \in \mathbb{R} \to x + z = y + z \to x \in \mathbb{R}) \to \varphi_{94}$
\end{abbrevlist}
\end{fullproof}

% 自動生成 (prover/build.js): NegUnique  ノード数 281
\begin{fullproof}{[F4]}{$\forall x\,\forall y\,(x \in \mathbb{R} \to y \in \mathbb{R} \to x + y = 0 \to y = -x)$}{281}\phantomsection\label{tree:NegUnique}
\noindent$\Pi_{1}$:
\begin{ptree}
\AXC{$y = z \to x = y \to x = z$\lbl{A6}}
\AXC{$(y = z \to x = y \to x = z) \to x = y \to y = z \to x = z$\lbl{A1}}
\BIC{$x = y \to y = z \to x = z$}
\UIC{$\forall z\,(x = y \to y = z \to x = z)$}
\UIC{$\forall y\,\forall z\,(x = y \to y = z \to x = z)$}
\UIC{$\forall x\,\forall y\,\forall z\,(x = y \to y = z \to x = z)$}
\AXC{$\varphi_{12}$\lbl{A2}}
\BIC{$\forall y\,\forall z\,(v_{55} = y \to y = z \to v_{55} = z)$}
\AXC{$\varphi_{18}$\lbl{A2}}
\BIC{$\forall z\,(v_{55} = v_{56} \to v_{56} = z \to v_{55} = z)$}
\AXC{$\varphi_{27}$\lbl{A2}}
\BIC{$v_{55} = v_{56} \to v_{56} = v_{57} \to v_{55} = v_{57}$}
\UIC{$\forall v_{57}\,(v_{55} = v_{56} \to v_{56} = v_{57} \to v_{55} = v_{57})$}
\UIC{$\varphi_{3}$}
\UIC{$\forall v_{55}\,(\varphi_{3})$}
\end{ptree}
\noindent$\Pi_{2}$:
\begin{ptree}
\AXC{$\forall x \in \mathbb{R}\,(-x \in \mathbb{R})$\lbl{(R1)}}
\AXC{$\forall x \in \mathbb{R}\,(-x \in \mathbb{R}) \to x \in \mathbb{R} \to -x \in \mathbb{R}$\lbl{A2}}
\BIC{$x \in \mathbb{R} \to -x \in \mathbb{R}$}
\AXC{$(x \in \mathbb{R} \to -x \in \mathbb{R}) \to \varphi_{80}$\lbl{A1}}
\BIC{$\varphi_{80}$}
\end{ptree}
\noindent$\Pi_{3}$:
\begin{ptree}
\AXC{$x \in \mathbb{R} \to y \in \mathbb{R} \to x + y = 0 \to x \in \mathbb{R}$\lbl{A1}}
\AXC{$\stackrel{\Pi_{2}}{\vdots}$}\noLine
\UIC{$\varphi_{80}$}
\BIC{$x \in \mathbb{R} \to y \in \mathbb{R} \to x + y = 0 \to -x \in \mathbb{R}$}
\end{ptree}
\noindent$\Pi_{4}$:
\begin{ptree}
\AXC{$\forall x \in \mathbb{R}\,(x + 0 = x)$\lbl{(R4)}}
\AXC{$\forall x \in \mathbb{R}\,(x + 0 = x) \to -x \in \mathbb{R} \to -x + 0 = -x$\lbl{A2}}
\BIC{$-x \in \mathbb{R} \to -x + 0 = -x$}
\AXC{$(-x \in \mathbb{R} \to -x + 0 = -x) \to \varphi_{84}$\lbl{A1}}
\BIC{$\varphi_{84}$}
\end{ptree}
\noindent$\Pi_{5}$:
\begin{ptree}
\AXC{$y = z \to x = y \to x = z$\lbl{A6}}
\AXC{$(y = z \to x = y \to x = z) \to x = y \to y = z \to x = z$\lbl{A1}}
\BIC{$x = y \to y = z \to x = z$}
\UIC{$\forall z\,(x = y \to y = z \to x = z)$}
\UIC{$\forall y\,\forall z\,(x = y \to y = z \to x = z)$}
\UIC{$\forall x\,\forall y\,\forall z\,(x = y \to y = z \to x = z)$}
\AXC{$\varphi_{13}$\lbl{A2}}
\BIC{$\forall y\,\forall z\,(v_{52} = y \to y = z \to v_{52} = z)$}
\AXC{$\varphi_{19}$\lbl{A2}}
\BIC{$\forall z\,(v_{52} = v_{53} \to v_{53} = z \to v_{52} = z)$}
\AXC{$\varphi_{28}$\lbl{A2}}
\BIC{$v_{52} = v_{53} \to v_{53} = v_{54} \to v_{52} = v_{54}$}
\UIC{$\forall v_{54}\,(v_{52} = v_{53} \to v_{53} = v_{54} \to v_{52} = v_{54})$}
\UIC{$\varphi_{4}$}
\UIC{$\forall v_{52}\,(\varphi_{4})$}
\end{ptree}
\noindent$\Pi_{6}$:
\begin{ptree}
\AXC{$\forall x\,\forall y\,(x \in \mathbb{R} \to y \in \mathbb{R} \to x + y = y + x)$\lbl{(R3)}}
\AXC{$\varphi_{66}$\lbl{A2}}
\BIC{$\forall y\,(0 \in \mathbb{R} \to y \in \mathbb{R} \to 0 + y = y + 0)$}
\AXC{$\varphi_{77}$\lbl{A2}}
\BIC{$0 \in \mathbb{R} \to -x \in \mathbb{R} \to 0 + -x = -x + 0$}
\end{ptree}
\noindent$\Pi_{7}$:
\begin{ptree}
\AXC{$0 \in \mathbb{R}$\lbl{(R1)}}
\AXC{$\stackrel{\Pi_{6}}{\vdots}$}\noLine
\UIC{$0 \in \mathbb{R} \to -x \in \mathbb{R} \to 0 + -x = -x + 0$}
\AXC{$(0 \in \mathbb{R} \to -x \in \mathbb{R} \to 0 + -x = -x + 0) \to 0 \in \mathbb{R} \to \varphi_{87}$\lbl{A1}}
\BIC{$0 \in \mathbb{R} \to \varphi_{87}$}
\BIC{$\varphi_{87}$}
\end{ptree}
\noindent$\Pi_{8}$:
\begin{ptree}
\AXC{$y = z \to x = y \to x = z$\lbl{A6}}
\AXC{$(y = z \to x = y \to x = z) \to x = y \to y = z \to x = z$\lbl{A1}}
\BIC{$x = y \to y = z \to x = z$}
\UIC{$\forall z\,(x = y \to y = z \to x = z)$}
\UIC{$\forall y\,\forall z\,(x = y \to y = z \to x = z)$}
\UIC{$\forall x\,\forall y\,\forall z\,(x = y \to y = z \to x = z)$}
\AXC{$\varphi_{14}$\lbl{A2}}
\BIC{$\forall y\,\forall z\,(v_{49} = y \to y = z \to v_{49} = z)$}
\AXC{$\varphi_{20}$\lbl{A2}}
\BIC{$\forall z\,(v_{49} = v_{50} \to v_{50} = z \to v_{49} = z)$}
\AXC{$\varphi_{29}$\lbl{A2}}
\BIC{$v_{49} = v_{50} \to v_{50} = v_{51} \to v_{49} = v_{51}$}
\UIC{$\forall v_{51}\,(v_{49} = v_{50} \to v_{50} = v_{51} \to v_{49} = v_{51})$}
\UIC{$\varphi_{5}$}
\UIC{$\forall v_{49}\,(\varphi_{5})$}
\end{ptree}
\noindent$\Pi_{9}$:
\begin{ptree}
\AXC{$\stackrel{\Pi_{8}}{\vdots}$}\noLine
\UIC{$\forall v_{49}\,(\varphi_{5})$}
\AXC{$\varphi_{53}$\lbl{A2}}
\BIC{$\forall v_{50}\,\forall v_{51}\,(y = v_{50} \to v_{50} = v_{51} \to y = v_{51})$}
\AXC{$\varphi_{67}$\lbl{A2}}
\BIC{$\varphi_{21}$}
\AXC{$\varphi_{21} \to y = x + y + -x \to x + y + -x = 0 + -x \to y = 0 + -x$\lbl{A2}}
\BIC{$y = x + y + -x \to x + y + -x = 0 + -x \to y = 0 + -x$}
\end{ptree}
\noindent$\Pi_{10}$:
\begin{ptree}
\AXC{$z = w \to x + y + -x = z + -x \to x + y + -x = w + -x$\lbl{A6}}
\UIC{$\varphi_{37}$}
\UIC{$\forall z\,(\varphi_{37})$}
\AXC{$\forall z\,(\varphi_{37}) \to \varphi_{54}$\lbl{A2}}
\BIC{$\varphi_{54}$}
\AXC{$\varphi_{54} \to \varphi_{46}$\lbl{A2}}
\BIC{$\varphi_{46}$}
\end{ptree}
\noindent$\Pi_{11}$:
\begin{ptree}
\AXC{$x \in \mathbb{R} \to y \in \mathbb{R} \to x + y = 0 \to x + y = 0$\lbl{A1}}
\AXC{$(x \in \mathbb{R} \to y \in \mathbb{R} \to x + y = 0 \to x + y = 0) \to \varphi_{93}$\lbl{A1}}
\BIC{$\varphi_{93}$}
\end{ptree}
\noindent$\Pi_{12}$:
\begin{ptree}
\AXC{$y = z \to x = y \to x = z$\lbl{A6}}
\AXC{$(y = z \to x = y \to x = z) \to x = y \to y = z \to x = z$\lbl{A1}}
\BIC{$x = y \to y = z \to x = z$}
\UIC{$\forall z\,(x = y \to y = z \to x = z)$}
\UIC{$\forall y\,\forall z\,(x = y \to y = z \to x = z)$}
\UIC{$\forall x\,\forall y\,\forall z\,(x = y \to y = z \to x = z)$}
\AXC{$\varphi_{15}$\lbl{A2}}
\BIC{$\forall y\,\forall z\,(v_{46} = y \to y = z \to v_{46} = z)$}
\AXC{$\varphi_{22}$\lbl{A2}}
\BIC{$\forall z\,(v_{46} = v_{47} \to v_{47} = z \to v_{46} = z)$}
\AXC{$\varphi_{30}$\lbl{A2}}
\BIC{$v_{46} = v_{47} \to v_{47} = v_{48} \to v_{46} = v_{48}$}
\UIC{$\forall v_{48}\,(v_{46} = v_{47} \to v_{47} = v_{48} \to v_{46} = v_{48})$}
\UIC{$\varphi_{6}$}
\UIC{$\forall v_{46}\,(\varphi_{6})$}
\end{ptree}
\noindent$\Pi_{13}$:
\begin{ptree}
\AXC{$\stackrel{\Pi_{12}}{\vdots}$}\noLine
\UIC{$\forall v_{46}\,(\varphi_{6})$}
\AXC{$\varphi_{55}$\lbl{A2}}
\BIC{$\forall v_{47}\,\forall v_{48}\,(y = v_{47} \to v_{47} = v_{48} \to y = v_{48})$}
\AXC{$\varphi_{68}$\lbl{A2}}
\BIC{$\varphi_{23}$}
\AXC{$\varphi_{23} \to \varphi_{47}$\lbl{A2}}
\BIC{$\varphi_{47}$}
\end{ptree}
\noindent$\Pi_{14}$:
\begin{ptree}
\AXC{$\forall x\,\forall y\,(x \in \mathbb{R} \to y \in \mathbb{R} \to x + y = y + x)$\lbl{(R3)}}
\AXC{$\varphi_{75}$\lbl{A2}}
\BIC{$\forall y\,(v_{34} \in \mathbb{R} \to y \in \mathbb{R} \to v_{34} + y = y + v_{34})$}
\AXC{$\varphi_{82}$\lbl{A2}}
\BIC{$v_{34} \in \mathbb{R} \to v_{35} \in \mathbb{R} \to v_{34} + v_{35} = v_{35} + v_{34}$}
\UIC{$\varphi_{31}$}
\UIC{$\forall v_{34}\,(\varphi_{31})$}
\AXC{$\varphi_{86}$\lbl{A2}}
\BIC{$\forall v_{35}\,(y \in \mathbb{R} \to v_{35} \in \mathbb{R} \to y + v_{35} = v_{35} + y)$}
\AXC{$\varphi_{71}$\lbl{A2}}
\BIC{$y \in \mathbb{R} \to x \in \mathbb{R} \to y + x = x + y$}
\AXC{$\varphi_{97}$\lbl{A1}}
\BIC{$(x \in \mathbb{R} \to y \in \mathbb{R} \to x + y = 0 \to y \in \mathbb{R}) \to \varphi_{78}$}
\end{ptree}
\noindent$\Pi_{15}$:
\begin{ptree}
\AXC{$x \in \mathbb{R} \to y \in \mathbb{R} \to x + y = 0 \to y \in \mathbb{R}$\lbl{A1}}
\AXC{$\stackrel{\Pi_{14}}{\vdots}$}\noLine
\UIC{$(x \in \mathbb{R} \to y \in \mathbb{R} \to x + y = 0 \to y \in \mathbb{R}) \to \varphi_{78}$}
\BIC{$\varphi_{78}$}
\end{ptree}
\noindent$\Pi_{16}$:
\begin{ptree}
\AXC{$x \in \mathbb{R} \to y \in \mathbb{R} \to x + y = 0 \to x \in \mathbb{R}$\lbl{A1}}
\AXC{$\stackrel{\Pi_{15}}{\vdots}$}\noLine
\UIC{$\varphi_{78}$}
\BIC{$x \in \mathbb{R} \to y \in \mathbb{R} \to x + y = 0 \to y + x = x + y$}
\AXC{$\varphi_{101}$\lbl{A1}}
\BIC{$y + x + -x = y + x + -x \to \varphi_{90}$}
\end{ptree}
\noindent$\Pi_{17}$:
\begin{ptree}
\AXC{$x = x$\lbl{A5}}
\UIC{$\forall x\,(x = x)$}
\AXC{$\forall x\,(x = x) \to y + x + -x = y + x + -x$\lbl{A2}}
\BIC{$y + x + -x = y + x + -x$}
\AXC{$\stackrel{\Pi_{16}}{\vdots}$}\noLine
\UIC{$y + x + -x = y + x + -x \to \varphi_{90}$}
\BIC{$\varphi_{90}$}
\end{ptree}
\noindent$\Pi_{18}$:
\begin{ptree}
\AXC{$y = z \to x = y \to x = z$\lbl{A6}}
\AXC{$(y = z \to x = y \to x = z) \to x = y \to y = z \to x = z$\lbl{A1}}
\BIC{$x = y \to y = z \to x = z$}
\UIC{$\forall z\,(x = y \to y = z \to x = z)$}
\UIC{$\forall y\,\forall z\,(x = y \to y = z \to x = z)$}
\UIC{$\forall x\,\forall y\,\forall z\,(x = y \to y = z \to x = z)$}
\AXC{$\varphi_{16}$\lbl{A2}}
\BIC{$\forall y\,\forall z\,(v_{43} = y \to y = z \to v_{43} = z)$}
\AXC{$\varphi_{24}$\lbl{A2}}
\BIC{$\forall z\,(v_{43} = v_{44} \to v_{44} = z \to v_{43} = z)$}
\AXC{$\varphi_{32}$\lbl{A2}}
\BIC{$v_{43} = v_{44} \to v_{44} = v_{45} \to v_{43} = v_{45}$}
\UIC{$\forall v_{45}\,(v_{43} = v_{44} \to v_{44} = v_{45} \to v_{43} = v_{45})$}
\UIC{$\varphi_{7}$}
\UIC{$\forall v_{43}\,(\varphi_{7})$}
\end{ptree}
\noindent$\Pi_{19}$:
\begin{ptree}
\AXC{$\stackrel{\Pi_{18}}{\vdots}$}\noLine
\UIC{$\forall v_{43}\,(\varphi_{7})$}
\AXC{$\varphi_{57}$\lbl{A2}}
\BIC{$\forall v_{44}\,\forall v_{45}\,(y = v_{44} \to v_{44} = v_{45} \to y = v_{45})$}
\AXC{$\varphi_{69}$\lbl{A2}}
\BIC{$\varphi_{25}$}
\AXC{$\varphi_{25} \to \varphi_{48}$\lbl{A2}}
\BIC{$\varphi_{48}$}
\end{ptree}
\noindent$\Pi_{20}$:
\begin{ptree}
\AXC{$x = x$\lbl{A5}}
\AXC{$x = y \to x = x \to y = x$\lbl{A6}}
\AXC{$(x = y \to x = x \to y = x) \to x = x \to x = y \to y = x$\lbl{A1}}
\BIC{$x = x \to x = y \to y = x$}
\BIC{$x = y \to y = x$}
\UIC{$\forall y\,(x = y \to y = x)$}
\UIC{$\forall x\,\forall y\,(x = y \to y = x)$}
\end{ptree}
\noindent$\Pi_{21}$:
\begin{ptree}
\AXC{$\stackrel{\Pi_{20}}{\vdots}$}\noLine
\UIC{$\forall x\,\forall y\,(x = y \to y = x)$}
\AXC{$\forall x\,\forall y\,(x = y \to y = x) \to \forall y\,(v_{38} = y \to y = v_{38})$\lbl{A2}}
\BIC{$\forall y\,(v_{38} = y \to y = v_{38})$}
\AXC{$\varphi_{1}$\lbl{A2}}
\BIC{$v_{38} = v_{39} \to v_{39} = v_{38}$}
\UIC{$\forall v_{39}\,(v_{38} = v_{39} \to v_{39} = v_{38})$}
\UIC{$\forall v_{38}\,\forall v_{39}\,(v_{38} = v_{39} \to v_{39} = v_{38})$}
\AXC{$\varphi_{58}$\lbl{A2}}
\BIC{$\forall v_{39}\,(y + x + -x = v_{39} \to v_{39} = y + x + -x)$}
\end{ptree}
\noindent$\Pi_{22}$:
\begin{ptree}
\AXC{$\varphi_{43}$\lbl{(R2)}}
\AXC{$\varphi_{43} \to \varphi_{45}$\lbl{A2}}
\BIC{$\varphi_{45}$}
\AXC{$\varphi_{45} \to \varphi_{49}$\lbl{A2}}
\BIC{$\varphi_{49}$}
\AXC{$\varphi_{49} \to v_{31} \in \mathbb{R} \to v_{32} \in \mathbb{R} \to \varphi_{11}$\lbl{A2}}
\BIC{$v_{31} \in \mathbb{R} \to v_{32} \in \mathbb{R} \to \varphi_{11}$}
\UIC{$\forall v_{33}\,(v_{31} \in \mathbb{R} \to v_{32} \in \mathbb{R} \to \varphi_{11})$}
\UIC{$\forall v_{32}\,\forall v_{33}\,(v_{31} \in \mathbb{R} \to v_{32} \in \mathbb{R} \to \varphi_{11})$}
\UIC{$\forall v_{31}\,\forall v_{32}\,\forall v_{33}\,(v_{31} \in \mathbb{R} \to v_{32} \in \mathbb{R} \to \varphi_{11})$}
\AXC{$\varphi_{94}$\lbl{A2}}
\BIC{$\varphi_{64}$}
\AXC{$\varphi_{64} \to \varphi_{44}$\lbl{A2}}
\BIC{$\varphi_{44}$}
\end{ptree}
\noindent$\Pi_{23}$:
\begin{ptree}
\AXC{$x \in \mathbb{R} \to y \in \mathbb{R} \to x + y = 0 \to y \in \mathbb{R}$\lbl{A1}}
\AXC{$\stackrel{\Pi_{22}}{\vdots}$}\noLine
\UIC{$\varphi_{44}$}
\AXC{$\varphi_{44} \to \varphi_{50}$\lbl{A2}}
\BIC{$\varphi_{50}$}
\AXC{$\varphi_{50} \to (x \in \mathbb{R} \to y \in \mathbb{R} \to x + y = 0 \to y \in \mathbb{R}) \to \varphi_{91}$\lbl{A1}}
\BIC{$(x \in \mathbb{R} \to y \in \mathbb{R} \to x + y = 0 \to y \in \mathbb{R}) \to \varphi_{91}$}
\BIC{$\varphi_{91}$}
\end{ptree}
\noindent$\Pi_{24}$:
\begin{ptree}
\AXC{$\stackrel{\Pi_{3}}{\vdots}$}\noLine
\UIC{$x \in \mathbb{R} \to y \in \mathbb{R} \to x + y = 0 \to -x \in \mathbb{R}$}
\AXC{$x \in \mathbb{R} \to y \in \mathbb{R} \to x + y = 0 \to x \in \mathbb{R}$\lbl{A1}}
\AXC{$\stackrel{\Pi_{23}}{\vdots}$}\noLine
\UIC{$\varphi_{91}$}
\BIC{$(x \in \mathbb{R} \to y \in \mathbb{R} \to x + y = 0 \to -x \in \mathbb{R}) \to \varphi_{40}$}
\BIC{$\varphi_{40}$}
\AXC{$\varphi_{40} \to \varphi_{39} \to \varphi_{41}$\lbl{A1}}
\BIC{$\varphi_{39} \to \varphi_{41}$}
\end{ptree}
\noindent$\Pi_{25}$:
\begin{ptree}
\AXC{$\stackrel{\Pi_{21}}{\vdots}$}\noLine
\UIC{$\forall v_{39}\,(y + x + -x = v_{39} \to v_{39} = y + x + -x)$}
\AXC{$\forall v_{39}\,(y + x + -x = v_{39} \to v_{39} = y + x + -x) \to \varphi_{39}$\lbl{A2}}
\BIC{$\varphi_{39}$}
\AXC{$\stackrel{\Pi_{24}}{\vdots}$}\noLine
\UIC{$\varphi_{39} \to \varphi_{41}$}
\BIC{$\varphi_{41}$}
\end{ptree}
\noindent$\Pi_{26}$:
\begin{ptree}
\AXC{$y = z \to x = y \to x = z$\lbl{A6}}
\AXC{$(y = z \to x = y \to x = z) \to x = y \to y = z \to x = z$\lbl{A1}}
\BIC{$x = y \to y = z \to x = z$}
\UIC{$\forall z\,(x = y \to y = z \to x = z)$}
\UIC{$\forall y\,\forall z\,(x = y \to y = z \to x = z)$}
\UIC{$\forall x\,\forall y\,\forall z\,(x = y \to y = z \to x = z)$}
\AXC{$\varphi_{17}$\lbl{A2}}
\BIC{$\forall y\,\forall z\,(v_{40} = y \to y = z \to v_{40} = z)$}
\AXC{$\varphi_{26}$\lbl{A2}}
\BIC{$\forall z\,(v_{40} = v_{41} \to v_{41} = z \to v_{40} = z)$}
\AXC{$\varphi_{33}$\lbl{A2}}
\BIC{$v_{40} = v_{41} \to v_{41} = v_{42} \to v_{40} = v_{42}$}
\UIC{$\forall v_{42}\,(v_{40} = v_{41} \to v_{41} = v_{42} \to v_{40} = v_{42})$}
\UIC{$\varphi_{8}$}
\UIC{$\forall v_{40}\,(\varphi_{8})$}
\end{ptree}
\noindent$\Pi_{27}$:
\begin{ptree}
\AXC{$\forall x \in \mathbb{R}\,(x + -x = 0)$\lbl{(R5)}}
\AXC{$\forall x \in \mathbb{R}\,(x + -x = 0) \to x \in \mathbb{R} \to x + -x = 0$\lbl{A2}}
\BIC{$x \in \mathbb{R} \to x + -x = 0$}
\AXC{$(x \in \mathbb{R} \to x + -x = 0) \to \varphi_{79}$\lbl{A1}}
\BIC{$\varphi_{79}$}
\end{ptree}
\noindent$\Pi_{28}$:
\begin{ptree}
\AXC{$x \in \mathbb{R} \to y \in \mathbb{R} \to x + y = 0 \to x \in \mathbb{R}$\lbl{A1}}
\AXC{$\stackrel{\Pi_{27}}{\vdots}$}\noLine
\UIC{$\varphi_{79}$}
\BIC{$x \in \mathbb{R} \to y \in \mathbb{R} \to x + y = 0 \to x + -x = 0$}
\AXC{$(x \in \mathbb{R} \to y \in \mathbb{R} \to x + y = 0 \to x + -x = 0) \to \varphi_{63}$\lbl{A1}}
\BIC{$\varphi_{63}$}
\end{ptree}
\noindent$\Pi_{29}$:
\begin{ptree}
\AXC{$\stackrel{\Pi_{20}}{\vdots}$}\noLine
\UIC{$\forall x\,\forall y\,(x = y \to y = x)$}
\AXC{$\varphi_{9}$\lbl{A2}}
\BIC{$\forall y\,(x + -x = y \to y = x + -x)$}
\AXC{$\varphi_{35}$\lbl{A2}}
\BIC{$x + -x = 0 \to 0 = x + -x$}
\AXC{$\stackrel{\Pi_{28}}{\vdots}$}\noLine
\UIC{$\varphi_{63}$}
\BIC{$x \in \mathbb{R} \to y \in \mathbb{R} \to x + y = 0 \to 0 = x + -x$}
\AXC{$\varphi_{96}$\lbl{A1}}
\BIC{$y + 0 = y + 0 \to \varphi_{88}$}
\end{ptree}
\noindent$\Pi_{30}$:
\begin{ptree}
\AXC{$x = x$\lbl{A5}}
\UIC{$\forall x\,(x = x)$}
\AXC{$\forall x\,(x = x) \to y + 0 = y + 0$\lbl{A2}}
\BIC{$y + 0 = y + 0$}
\AXC{$\stackrel{\Pi_{29}}{\vdots}$}\noLine
\UIC{$y + 0 = y + 0 \to \varphi_{88}$}
\BIC{$\varphi_{88}$}
\end{ptree}
\noindent$\Pi_{31}$:
\begin{ptree}
\AXC{$\stackrel{\Pi_{20}}{\vdots}$}\noLine
\UIC{$\forall x\,\forall y\,(x = y \to y = x)$}
\AXC{$\forall x\,\forall y\,(x = y \to y = x) \to \forall y\,(v_{36} = y \to y = v_{36})$\lbl{A2}}
\BIC{$\forall y\,(v_{36} = y \to y = v_{36})$}
\AXC{$\varphi_{2}$\lbl{A2}}
\BIC{$v_{36} = v_{37} \to v_{37} = v_{36}$}
\UIC{$\forall v_{37}\,(v_{36} = v_{37} \to v_{37} = v_{36})$}
\UIC{$\forall v_{36}\,\forall v_{37}\,(v_{36} = v_{37} \to v_{37} = v_{36})$}
\AXC{$\varphi_{34}$\lbl{A2}}
\BIC{$\forall v_{37}\,(y + 0 = v_{37} \to v_{37} = y + 0)$}
\AXC{$\varphi_{10}$\lbl{A2}}
\BIC{$y + 0 = y \to y = y + 0$}
\end{ptree}
\noindent$\Pi_{32}$:
\begin{ptree}
\AXC{$\forall x \in \mathbb{R}\,(x + 0 = x)$\lbl{(R4)}}
\AXC{$\forall x \in \mathbb{R}\,(x + 0 = x) \to y \in \mathbb{R} \to y + 0 = y$\lbl{A2}}
\BIC{$y \in \mathbb{R} \to y + 0 = y$}
\AXC{$(y \in \mathbb{R} \to y + 0 = y) \to \varphi_{74}$\lbl{A1}}
\BIC{$\varphi_{74}$}
\end{ptree}
\noindent$\Pi_{33}$:
\begin{ptree}
\AXC{$x \in \mathbb{R} \to y \in \mathbb{R} \to x + y = 0 \to y \in \mathbb{R}$\lbl{A1}}
\AXC{$\stackrel{\Pi_{32}}{\vdots}$}\noLine
\UIC{$\varphi_{74}$}
\BIC{$x \in \mathbb{R} \to y \in \mathbb{R} \to x + y = 0 \to y + 0 = y$}
\AXC{$(x \in \mathbb{R} \to y \in \mathbb{R} \to x + y = 0 \to y + 0 = y) \to \varphi_{36}$\lbl{A1}}
\BIC{$\varphi_{36}$}
\end{ptree}
\noindent$\Pi_{34}$:
\begin{ptree}
\AXC{$\stackrel{\Pi_{31}}{\vdots}$}\noLine
\UIC{$y + 0 = y \to y = y + 0$}
\AXC{$\stackrel{\Pi_{33}}{\vdots}$}\noLine
\UIC{$\varphi_{36}$}
\BIC{$x \in \mathbb{R} \to y \in \mathbb{R} \to x + y = 0 \to y = y + 0$}
\AXC{$(x \in \mathbb{R} \to y \in \mathbb{R} \to x + y = 0 \to y = y + 0) \to \varphi_{92}$\lbl{A1}}
\BIC{$\varphi_{92}$}
\end{ptree}
\noindent$\Pi_{35}$:
\begin{ptree}
\AXC{$z = w \to y + 0 = y + z \to y + 0 = y + w$\lbl{A6}}
\UIC{$\forall w\,(z = w \to y + 0 = y + z \to y + 0 = y + w)$}
\UIC{$\forall z\,\forall w\,(z = w \to y + 0 = y + z \to y + 0 = y + w)$}
\AXC{$\varphi_{62}$\lbl{A2}}
\BIC{$\forall w\,(0 = w \to y + 0 = y + 0 \to y + 0 = y + w)$}
\AXC{$\varphi_{76}$\lbl{A2}}
\BIC{$0 = x + -x \to y + 0 = y + 0 \to y + 0 = y + (x + -x)$}
\AXC{$\stackrel{\Pi_{30}}{\vdots}$}\noLine
\UIC{$\varphi_{88}$}
\BIC{$x \in \mathbb{R} \to y \in \mathbb{R} \to x + y = 0 \to y + 0 = y + (x + -x)$}
\AXC{$\stackrel{\Pi_{34}}{\vdots}$}\noLine
\UIC{$\varphi_{92}$}
\BIC{$\varphi_{85}$}
\end{ptree}
\noindent$\Pi_{36}$:
\begin{ptree}
\AXC{$\stackrel{\Pi_{25}}{\vdots}$}\noLine
\UIC{$\varphi_{41}$}
\AXC{$\stackrel{\Pi_{26}}{\vdots}$}\noLine
\UIC{$\forall v_{40}\,(\varphi_{8})$}
\AXC{$\varphi_{59}$\lbl{A2}}
\BIC{$\forall v_{41}\,\forall v_{42}\,(y = v_{41} \to v_{41} = v_{42} \to y = v_{42})$}
\AXC{$\varphi_{42}$\lbl{A2}}
\BIC{$\forall v_{42}\,(y = y + 0 \to y + 0 = v_{42} \to y = v_{42})$}
\AXC{$\varphi_{72}$\lbl{A2}}
\BIC{$y = y + 0 \to y + 0 = y + (x + -x) \to y = y + (x + -x)$}
\AXC{$\stackrel{\Pi_{35}}{\vdots}$}\noLine
\UIC{$\varphi_{85}$}
\BIC{$x \in \mathbb{R} \to y \in \mathbb{R} \to x + y = 0 \to y = y + (x + -x)$}
\AXC{$\varphi_{102}$\lbl{A1}}
\BIC{$\varphi_{41} \to \varphi_{48} \to x \in \mathbb{R} \to y \in \mathbb{R} \to x + y = 0 \to y = y + x + -x$}
\BIC{$\varphi_{48} \to x \in \mathbb{R} \to y \in \mathbb{R} \to x + y = 0 \to y = y + x + -x$}
\end{ptree}
\noindent$\Pi_{37}$:
\begin{ptree}
\AXC{$\stackrel{\Pi_{19}}{\vdots}$}\noLine
\UIC{$\varphi_{48}$}
\AXC{$\stackrel{\Pi_{36}}{\vdots}$}\noLine
\UIC{$\varphi_{48} \to x \in \mathbb{R} \to y \in \mathbb{R} \to x + y = 0 \to y = y + x + -x$}
\BIC{$x \in \mathbb{R} \to y \in \mathbb{R} \to x + y = 0 \to y = y + x + -x$}
\AXC{$(x \in \mathbb{R} \to y \in \mathbb{R} \to x + y = 0 \to y = y + x + -x) \to \varphi_{99}$\lbl{A1}}
\BIC{$\varphi_{99}$}
\end{ptree}
\noindent$\Pi_{38}$:
\begin{ptree}
\AXC{$z = w \to y + x + -x = z + -x \to y + x + -x = w + -x$\lbl{A6}}
\UIC{$\varphi_{38}$}
\UIC{$\forall z\,(\varphi_{38})$}
\AXC{$\forall z\,(\varphi_{38}) \to \varphi_{56}$\lbl{A2}}
\BIC{$\varphi_{56}$}
\AXC{$\varphi_{56} \to \varphi_{65}$\lbl{A2}}
\BIC{$\varphi_{65}$}
\AXC{$\stackrel{\Pi_{17}}{\vdots}$}\noLine
\UIC{$\varphi_{90}$}
\BIC{$x \in \mathbb{R} \to y \in \mathbb{R} \to x + y = 0 \to y + x + -x = x + y + -x$}
\AXC{$\stackrel{\Pi_{37}}{\vdots}$}\noLine
\UIC{$\varphi_{99}$}
\BIC{$\varphi_{47} \to x \in \mathbb{R} \to y \in \mathbb{R} \to x + y = 0 \to y = x + y + -x$}
\end{ptree}
\noindent$\Pi_{39}$:
\begin{ptree}
\AXC{$\stackrel{\Pi_{13}}{\vdots}$}\noLine
\UIC{$\varphi_{47}$}
\AXC{$\stackrel{\Pi_{38}}{\vdots}$}\noLine
\UIC{$\varphi_{47} \to x \in \mathbb{R} \to y \in \mathbb{R} \to x + y = 0 \to y = x + y + -x$}
\BIC{$x \in \mathbb{R} \to y \in \mathbb{R} \to x + y = 0 \to y = x + y + -x$}
\AXC{$(x \in \mathbb{R} \to y \in \mathbb{R} \to x + y = 0 \to y = x + y + -x) \to \varphi_{95}$\lbl{A1}}
\BIC{$\varphi_{95}$}
\end{ptree}
\noindent$\Pi_{40}$:
\begin{ptree}
\AXC{$\stackrel{\Pi_{10}}{\vdots}$}\noLine
\UIC{$\varphi_{46}$}
\AXC{$x = x$\lbl{A5}}
\UIC{$\forall x\,(x = x)$}
\AXC{$\forall x\,(x = x) \to x + y + -x = x + y + -x$\lbl{A2}}
\BIC{$x + y + -x = x + y + -x$}
\AXC{$\stackrel{\Pi_{11}}{\vdots}$}\noLine
\UIC{$\varphi_{93}$}
\BIC{$\varphi_{46} \to x \in \mathbb{R} \to y \in \mathbb{R} \to x + y = 0 \to x + y + -x = 0 + -x$}
\BIC{$x \in \mathbb{R} \to y \in \mathbb{R} \to x + y = 0 \to x + y + -x = 0 + -x$}
\AXC{$\stackrel{\Pi_{39}}{\vdots}$}\noLine
\UIC{$\varphi_{95}$}
\BIC{$\varphi_{89}$}
\end{ptree}
\noindent$\Pi_{41}$:
\begin{ptree}
\AXC{$\stackrel{\Pi_{9}}{\vdots}$}\noLine
\UIC{$y = x + y + -x \to x + y + -x = 0 + -x \to y = 0 + -x$}
\AXC{$\stackrel{\Pi_{40}}{\vdots}$}\noLine
\UIC{$\varphi_{89}$}
\BIC{$x \in \mathbb{R} \to y \in \mathbb{R} \to x + y = 0 \to y = 0 + -x$}
\AXC{$\varphi_{100}$\lbl{A1}}
\BIC{$(x \in \mathbb{R} \to y \in \mathbb{R} \to x + y = 0 \to 0 + -x = -x + 0) \to \varphi_{83}$}
\end{ptree}
\noindent$\Pi_{42}$:
\begin{ptree}
\AXC{$\stackrel{\Pi_{3}}{\vdots}$}\noLine
\UIC{$x \in \mathbb{R} \to y \in \mathbb{R} \to x + y = 0 \to -x \in \mathbb{R}$}
\AXC{$\stackrel{\Pi_{7}}{\vdots}$}\noLine
\UIC{$\varphi_{87}$}
\BIC{$x \in \mathbb{R} \to y \in \mathbb{R} \to x + y = 0 \to 0 + -x = -x + 0$}
\AXC{$\stackrel{\Pi_{41}}{\vdots}$}\noLine
\UIC{$(x \in \mathbb{R} \to y \in \mathbb{R} \to x + y = 0 \to 0 + -x = -x + 0) \to \varphi_{83}$}
\BIC{$\varphi_{83}$}
\end{ptree}
\noindent$\Pi_{43}$:
\begin{ptree}
\AXC{$\stackrel{\Pi_{5}}{\vdots}$}\noLine
\UIC{$\forall v_{52}\,(\varphi_{4})$}
\AXC{$\varphi_{52}$\lbl{A2}}
\BIC{$\forall v_{53}\,\forall v_{54}\,(y = v_{53} \to v_{53} = v_{54} \to y = v_{54})$}
\AXC{$\varphi_{61}$\lbl{A2}}
\BIC{$\forall v_{54}\,(y = 0 + -x \to 0 + -x = v_{54} \to y = v_{54})$}
\AXC{$\varphi_{81}$\lbl{A2}}
\BIC{$y = 0 + -x \to 0 + -x = -x + 0 \to y = -x + 0$}
\AXC{$\stackrel{\Pi_{42}}{\vdots}$}\noLine
\UIC{$\varphi_{83}$}
\BIC{$x \in \mathbb{R} \to y \in \mathbb{R} \to x + y = 0 \to y = -x + 0$}
\AXC{$\varphi_{98}$\lbl{A1}}
\BIC{$(x \in \mathbb{R} \to y \in \mathbb{R} \to x + y = 0 \to -x + 0 = -x) \to \varphi_{73}$}
\end{ptree}
\noindent$\Pi_{44}$:
\begin{ptree}
\AXC{$\stackrel{\Pi_{3}}{\vdots}$}\noLine
\UIC{$x \in \mathbb{R} \to y \in \mathbb{R} \to x + y = 0 \to -x \in \mathbb{R}$}
\AXC{$\stackrel{\Pi_{4}}{\vdots}$}\noLine
\UIC{$\varphi_{84}$}
\BIC{$x \in \mathbb{R} \to y \in \mathbb{R} \to x + y = 0 \to -x + 0 = -x$}
\AXC{$\stackrel{\Pi_{43}}{\vdots}$}\noLine
\UIC{$(x \in \mathbb{R} \to y \in \mathbb{R} \to x + y = 0 \to -x + 0 = -x) \to \varphi_{73}$}
\BIC{$\varphi_{73}$}
\end{ptree}
\noindent 主木:
\begin{ptree}
\AXC{$\stackrel{\Pi_{1}}{\vdots}$}\noLine
\UIC{$\forall v_{55}\,(\varphi_{3})$}
\AXC{$\varphi_{51}$\lbl{A2}}
\BIC{$\forall v_{56}\,\forall v_{57}\,(y = v_{56} \to v_{56} = v_{57} \to y = v_{57})$}
\AXC{$\varphi_{60}$\lbl{A2}}
\BIC{$\forall v_{57}\,(y = -x + 0 \to -x + 0 = v_{57} \to y = v_{57})$}
\AXC{$\varphi_{70}$\lbl{A2}}
\BIC{$y = -x + 0 \to -x + 0 = -x \to y = -x$}
\AXC{$\stackrel{\Pi_{44}}{\vdots}$}\noLine
\UIC{$\varphi_{73}$}
\BIC{$x \in \mathbb{R} \to y \in \mathbb{R} \to x + y = 0 \to y = -x$}
\UIC{$\forall y\,(x \in \mathbb{R} \to y \in \mathbb{R} \to x + y = 0 \to y = -x)$}
\UIC{$\forall x\,\forall y\,(x \in \mathbb{R} \to y \in \mathbb{R} \to x + y = 0 \to y = -x)$}
\end{ptree}
\begin{abbrevlist}
\item $\varphi_{1} :\equiv \forall y\,(v_{38} = y \to y = v_{38}) \to v_{38} = v_{39} \to v_{39} = v_{38}$
\item $\varphi_{2} :\equiv \forall y\,(v_{36} = y \to y = v_{36}) \to v_{36} = v_{37} \to v_{37} = v_{36}$
\item $\varphi_{3} :\equiv \forall v_{56}\,\forall v_{57}\,(v_{55} = v_{56} \to v_{56} = v_{57} \to v_{55} = v_{57})$
\item $\varphi_{4} :\equiv \forall v_{53}\,\forall v_{54}\,(v_{52} = v_{53} \to v_{53} = v_{54} \to v_{52} = v_{54})$
\item $\varphi_{5} :\equiv \forall v_{50}\,\forall v_{51}\,(v_{49} = v_{50} \to v_{50} = v_{51} \to v_{49} = v_{51})$
\item $\varphi_{6} :\equiv \forall v_{47}\,\forall v_{48}\,(v_{46} = v_{47} \to v_{47} = v_{48} \to v_{46} = v_{48})$
\item $\varphi_{7} :\equiv \forall v_{44}\,\forall v_{45}\,(v_{43} = v_{44} \to v_{44} = v_{45} \to v_{43} = v_{45})$
\item $\varphi_{8} :\equiv \forall v_{41}\,\forall v_{42}\,(v_{40} = v_{41} \to v_{41} = v_{42} \to v_{40} = v_{42})$
\item $\varphi_{9} :\equiv \forall x\,\forall y\,(x = y \to y = x) \to \forall y\,(x + -x = y \to y = x + -x)$
\item $\varphi_{10} :\equiv \forall v_{37}\,(y + 0 = v_{37} \to v_{37} = y + 0) \to y + 0 = y \to y = y + 0$
\item $\varphi_{11} :\equiv v_{33} \in \mathbb{R} \to v_{31} + v_{32} + v_{33} = v_{31} + (v_{32} + v_{33})$
\item $\varphi_{12} :\equiv \forall x\,\forall y\,\forall z\,(x = y \to y = z \to x = z) \to \forall y\,\forall z\,(v_{55} = y \to y = z \to v_{55} = z)$
\item $\varphi_{13} :\equiv \forall x\,\forall y\,\forall z\,(x = y \to y = z \to x = z) \to \forall y\,\forall z\,(v_{52} = y \to y = z \to v_{52} = z)$
\item $\varphi_{14} :\equiv \forall x\,\forall y\,\forall z\,(x = y \to y = z \to x = z) \to \forall y\,\forall z\,(v_{49} = y \to y = z \to v_{49} = z)$
\item $\varphi_{15} :\equiv \forall x\,\forall y\,\forall z\,(x = y \to y = z \to x = z) \to \forall y\,\forall z\,(v_{46} = y \to y = z \to v_{46} = z)$
\item $\varphi_{16} :\equiv \forall x\,\forall y\,\forall z\,(x = y \to y = z \to x = z) \to \forall y\,\forall z\,(v_{43} = y \to y = z \to v_{43} = z)$
\item $\varphi_{17} :\equiv \forall x\,\forall y\,\forall z\,(x = y \to y = z \to x = z) \to \forall y\,\forall z\,(v_{40} = y \to y = z \to v_{40} = z)$
\item $\varphi_{18} :\equiv \forall y\,\forall z\,(v_{55} = y \to y = z \to v_{55} = z) \to \forall z\,(v_{55} = v_{56} \to v_{56} = z \to v_{55} = z)$
\item $\varphi_{19} :\equiv \forall y\,\forall z\,(v_{52} = y \to y = z \to v_{52} = z) \to \forall z\,(v_{52} = v_{53} \to v_{53} = z \to v_{52} = z)$
\item $\varphi_{20} :\equiv \forall y\,\forall z\,(v_{49} = y \to y = z \to v_{49} = z) \to \forall z\,(v_{49} = v_{50} \to v_{50} = z \to v_{49} = z)$
\item $\varphi_{21} :\equiv \forall v_{51}\,(y = x + y + -x \to x + y + -x = v_{51} \to y = v_{51})$
\item $\varphi_{22} :\equiv \forall y\,\forall z\,(v_{46} = y \to y = z \to v_{46} = z) \to \forall z\,(v_{46} = v_{47} \to v_{47} = z \to v_{46} = z)$
\item $\varphi_{23} :\equiv \forall v_{48}\,(y = y + x + -x \to y + x + -x = v_{48} \to y = v_{48})$
\item $\varphi_{24} :\equiv \forall y\,\forall z\,(v_{43} = y \to y = z \to v_{43} = z) \to \forall z\,(v_{43} = v_{44} \to v_{44} = z \to v_{43} = z)$
\item $\varphi_{25} :\equiv \forall v_{45}\,(y = y + (x + -x) \to y + (x + -x) = v_{45} \to y = v_{45})$
\item $\varphi_{26} :\equiv \forall y\,\forall z\,(v_{40} = y \to y = z \to v_{40} = z) \to \forall z\,(v_{40} = v_{41} \to v_{41} = z \to v_{40} = z)$
\item $\varphi_{27} :\equiv \forall z\,(v_{55} = v_{56} \to v_{56} = z \to v_{55} = z) \to v_{55} = v_{56} \to v_{56} = v_{57} \to v_{55} = v_{57}$
\item $\varphi_{28} :\equiv \forall z\,(v_{52} = v_{53} \to v_{53} = z \to v_{52} = z) \to v_{52} = v_{53} \to v_{53} = v_{54} \to v_{52} = v_{54}$
\item $\varphi_{29} :\equiv \forall z\,(v_{49} = v_{50} \to v_{50} = z \to v_{49} = z) \to v_{49} = v_{50} \to v_{50} = v_{51} \to v_{49} = v_{51}$
\item $\varphi_{30} :\equiv \forall z\,(v_{46} = v_{47} \to v_{47} = z \to v_{46} = z) \to v_{46} = v_{47} \to v_{47} = v_{48} \to v_{46} = v_{48}$
\item $\varphi_{31} :\equiv \forall v_{35}\,(v_{34} \in \mathbb{R} \to v_{35} \in \mathbb{R} \to v_{34} + v_{35} = v_{35} + v_{34})$
\item $\varphi_{32} :\equiv \forall z\,(v_{43} = v_{44} \to v_{44} = z \to v_{43} = z) \to v_{43} = v_{44} \to v_{44} = v_{45} \to v_{43} = v_{45}$
\item $\varphi_{33} :\equiv \forall z\,(v_{40} = v_{41} \to v_{41} = z \to v_{40} = z) \to v_{40} = v_{41} \to v_{41} = v_{42} \to v_{40} = v_{42}$
\item $\varphi_{34} :\equiv \forall v_{36}\,\forall v_{37}\,(v_{36} = v_{37} \to v_{37} = v_{36}) \to \forall v_{37}\,(y + 0 = v_{37} \to v_{37} = y + 0)$
\item $\varphi_{35} :\equiv \forall y\,(x + -x = y \to y = x + -x) \to x + -x = 0 \to 0 = x + -x$
\item $\varphi_{36} :\equiv (y + 0 = y \to y = y + 0) \to x \in \mathbb{R} \to y \in \mathbb{R} \to x + y = 0 \to y = y + 0$
\item $\varphi_{37} :\equiv \forall w\,(z = w \to x + y + -x = z + -x \to x + y + -x = w + -x)$
\item $\varphi_{38} :\equiv \forall w\,(z = w \to y + x + -x = z + -x \to y + x + -x = w + -x)$
\item $\varphi_{39} :\equiv y + x + -x = y + (x + -x) \to y + (x + -x) = y + x + -x$
\item $\varphi_{40} :\equiv x \in \mathbb{R} \to y \in \mathbb{R} \to x + y = 0 \to y + x + -x = y + (x + -x)$
\item $\varphi_{41} :\equiv x \in \mathbb{R} \to y \in \mathbb{R} \to x + y = 0 \to y + (x + -x) = y + x + -x$
\item $\varphi_{42} :\equiv \forall v_{41}\,\forall v_{42}\,(y = v_{41} \to v_{41} = v_{42} \to y = v_{42}) \to \forall v_{42}\,(y = y + 0 \to y + 0 = v_{42} \to y = v_{42})$
\item $\varphi_{43} :\equiv \forall x\,\forall y\,\forall z\,(x \in \mathbb{R} \to y \in \mathbb{R} \to z \in \mathbb{R} \to x + y + z = x + (y + z))$
\item $\varphi_{44} :\equiv \forall v_{33}\,(y \in \mathbb{R} \to x \in \mathbb{R} \to v_{33} \in \mathbb{R} \to y + x + v_{33} = y + (x + v_{33}))$
\item $\varphi_{45} :\equiv \forall y\,\forall z\,(v_{31} \in \mathbb{R} \to y \in \mathbb{R} \to z \in \mathbb{R} \to v_{31} + y + z = v_{31} + (y + z))$
\item $\varphi_{46} :\equiv x + y = 0 \to x + y + -x = x + y + -x \to x + y + -x = 0 + -x$
\item $\varphi_{47} :\equiv y = y + x + -x \to y + x + -x = x + y + -x \to y = x + y + -x$
\item $\varphi_{48} :\equiv y = y + (x + -x) \to y + (x + -x) = y + x + -x \to y = y + x + -x$
\item $\varphi_{49} :\equiv \forall z\,(v_{31} \in \mathbb{R} \to v_{32} \in \mathbb{R} \to z \in \mathbb{R} \to v_{31} + v_{32} + z = v_{31} + (v_{32} + z))$
\item $\varphi_{50} :\equiv y \in \mathbb{R} \to x \in \mathbb{R} \to -x \in \mathbb{R} \to y + x + -x = y + (x + -x)$
\item $\varphi_{51} :\equiv \forall v_{55}\,(\varphi_{3}) \to \forall v_{56}\,\forall v_{57}\,(y = v_{56} \to v_{56} = v_{57} \to y = v_{57})$
\item $\varphi_{52} :\equiv \forall v_{52}\,(\varphi_{4}) \to \forall v_{53}\,\forall v_{54}\,(y = v_{53} \to v_{53} = v_{54} \to y = v_{54})$
\item $\varphi_{53} :\equiv \forall v_{49}\,(\varphi_{5}) \to \forall v_{50}\,\forall v_{51}\,(y = v_{50} \to v_{50} = v_{51} \to y = v_{51})$
\item $\varphi_{54} :\equiv \forall w\,(x + y = w \to x + y + -x = x + y + -x \to x + y + -x = w + -x)$
\item $\varphi_{55} :\equiv \forall v_{46}\,(\varphi_{6}) \to \forall v_{47}\,\forall v_{48}\,(y = v_{47} \to v_{47} = v_{48} \to y = v_{48})$
\item $\varphi_{56} :\equiv \forall w\,(y + x = w \to y + x + -x = y + x + -x \to y + x + -x = w + -x)$
\item $\varphi_{57} :\equiv \forall v_{43}\,(\varphi_{7}) \to \forall v_{44}\,\forall v_{45}\,(y = v_{44} \to v_{44} = v_{45} \to y = v_{45})$
\item $\varphi_{58} :\equiv \forall v_{38}\,\forall v_{39}\,(v_{38} = v_{39} \to v_{39} = v_{38}) \to \forall v_{39}\,(y + x + -x = v_{39} \to v_{39} = y + x + -x)$
\item $\varphi_{59} :\equiv \forall v_{40}\,(\varphi_{8}) \to \forall v_{41}\,\forall v_{42}\,(y = v_{41} \to v_{41} = v_{42} \to y = v_{42})$
\item $\varphi_{60} :\equiv \forall v_{56}\,\forall v_{57}\,(y = v_{56} \to v_{56} = v_{57} \to y = v_{57}) \to \forall v_{57}\,(y = -x + 0 \to -x + 0 = v_{57} \to y = v_{57})$
\item $\varphi_{61} :\equiv \forall v_{53}\,\forall v_{54}\,(y = v_{53} \to v_{53} = v_{54} \to y = v_{54}) \to \forall v_{54}\,(y = 0 + -x \to 0 + -x = v_{54} \to y = v_{54})$
\item $\varphi_{62} :\equiv \forall z\,\forall w\,(z = w \to y + 0 = y + z \to y + 0 = y + w) \to \forall w\,(0 = w \to y + 0 = y + 0 \to y + 0 = y + w)$
\item $\varphi_{63} :\equiv (x + -x = 0 \to 0 = x + -x) \to x \in \mathbb{R} \to y \in \mathbb{R} \to x + y = 0 \to 0 = x + -x$
\item $\varphi_{64} :\equiv \forall v_{32}\,\forall v_{33}\,(y \in \mathbb{R} \to v_{32} \in \mathbb{R} \to v_{33} \in \mathbb{R} \to y + v_{32} + v_{33} = y + (v_{32} + v_{33}))$
\item $\varphi_{65} :\equiv y + x = x + y \to y + x + -x = y + x + -x \to y + x + -x = x + y + -x$
\item $\varphi_{66} :\equiv \forall x\,\forall y\,(x \in \mathbb{R} \to y \in \mathbb{R} \to x + y = y + x) \to \forall y\,(0 \in \mathbb{R} \to y \in \mathbb{R} \to 0 + y = y + 0)$
\item $\varphi_{67} :\equiv \forall v_{50}\,\forall v_{51}\,(y = v_{50} \to v_{50} = v_{51} \to y = v_{51}) \to \varphi_{21}$
\item $\varphi_{68} :\equiv \forall v_{47}\,\forall v_{48}\,(y = v_{47} \to v_{47} = v_{48} \to y = v_{48}) \to \varphi_{23}$
\item $\varphi_{69} :\equiv \forall v_{44}\,\forall v_{45}\,(y = v_{44} \to v_{44} = v_{45} \to y = v_{45}) \to \varphi_{25}$
\item $\varphi_{70} :\equiv \forall v_{57}\,(y = -x + 0 \to -x + 0 = v_{57} \to y = v_{57}) \to y = -x + 0 \to -x + 0 = -x \to y = -x$
\item $\varphi_{71} :\equiv \forall v_{35}\,(y \in \mathbb{R} \to v_{35} \in \mathbb{R} \to y + v_{35} = v_{35} + y) \to y \in \mathbb{R} \to x \in \mathbb{R} \to y + x = x + y$
\item $\varphi_{72} :\equiv \forall v_{42}\,(y = y + 0 \to y + 0 = v_{42} \to y = v_{42}) \to y = y + 0 \to y + 0 = y + (x + -x) \to y = y + (x + -x)$
\item $\varphi_{73} :\equiv (y = -x + 0 \to -x + 0 = -x \to y = -x) \to x \in \mathbb{R} \to y \in \mathbb{R} \to x + y = 0 \to y = -x$
\item $\varphi_{74} :\equiv (x \in \mathbb{R} \to y \in \mathbb{R} \to x + y = 0 \to y \in \mathbb{R}) \to x \in \mathbb{R} \to y \in \mathbb{R} \to x + y = 0 \to y + 0 = y$
\item $\varphi_{75} :\equiv \forall x\,\forall y\,(x \in \mathbb{R} \to y \in \mathbb{R} \to x + y = y + x) \to \forall y\,(v_{34} \in \mathbb{R} \to y \in \mathbb{R} \to v_{34} + y = y + v_{34})$
\item $\varphi_{76} :\equiv \forall w\,(0 = w \to y + 0 = y + 0 \to y + 0 = y + w) \to 0 = x + -x \to y + 0 = y + 0 \to y + 0 = y + (x + -x)$
\item $\varphi_{77} :\equiv \forall y\,(0 \in \mathbb{R} \to y \in \mathbb{R} \to 0 + y = y + 0) \to 0 \in \mathbb{R} \to -x \in \mathbb{R} \to 0 + -x = -x + 0$
\item $\varphi_{78} :\equiv (x \in \mathbb{R} \to y \in \mathbb{R} \to x + y = 0 \to x \in \mathbb{R}) \to x \in \mathbb{R} \to y \in \mathbb{R} \to x + y = 0 \to y + x = x + y$
\item $\varphi_{79} :\equiv (x \in \mathbb{R} \to y \in \mathbb{R} \to x + y = 0 \to x \in \mathbb{R}) \to x \in \mathbb{R} \to y \in \mathbb{R} \to x + y = 0 \to x + -x = 0$
\item $\varphi_{80} :\equiv (x \in \mathbb{R} \to y \in \mathbb{R} \to x + y = 0 \to x \in \mathbb{R}) \to x \in \mathbb{R} \to y \in \mathbb{R} \to x + y = 0 \to -x \in \mathbb{R}$
\item $\varphi_{81} :\equiv \forall v_{54}\,(y = 0 + -x \to 0 + -x = v_{54} \to y = v_{54}) \to y = 0 + -x \to 0 + -x = -x + 0 \to y = -x + 0$
\item $\varphi_{82} :\equiv \forall y\,(v_{34} \in \mathbb{R} \to y \in \mathbb{R} \to v_{34} + y = y + v_{34}) \to v_{34} \in \mathbb{R} \to v_{35} \in \mathbb{R} \to v_{34} + v_{35} = v_{35} + v_{34}$
\item $\varphi_{83} :\equiv (y = 0 + -x \to 0 + -x = -x + 0 \to y = -x + 0) \to x \in \mathbb{R} \to y \in \mathbb{R} \to x + y = 0 \to y = -x + 0$
\item $\varphi_{84} :\equiv (x \in \mathbb{R} \to y \in \mathbb{R} \to x + y = 0 \to -x \in \mathbb{R}) \to x \in \mathbb{R} \to y \in \mathbb{R} \to x + y = 0 \to -x + 0 = -x$
\item $\varphi_{85} :\equiv (y = y + 0 \to y + 0 = y + (x + -x) \to y = y + (x + -x)) \to x \in \mathbb{R} \to y \in \mathbb{R} \to x + y = 0 \to y = y + (x + -x)$
\item $\varphi_{86} :\equiv \forall v_{34}\,(\varphi_{31}) \to \forall v_{35}\,(y \in \mathbb{R} \to v_{35} \in \mathbb{R} \to y + v_{35} = v_{35} + y)$
\item $\varphi_{87} :\equiv (x \in \mathbb{R} \to y \in \mathbb{R} \to x + y = 0 \to -x \in \mathbb{R}) \to x \in \mathbb{R} \to y \in \mathbb{R} \to x + y = 0 \to 0 + -x = -x + 0$
\item $\varphi_{88} :\equiv (0 = x + -x \to y + 0 = y + 0 \to y + 0 = y + (x + -x)) \to x \in \mathbb{R} \to y \in \mathbb{R} \to x + y = 0 \to y + 0 = y + (x + -x)$
\item $\varphi_{89} :\equiv (y = x + y + -x \to x + y + -x = 0 + -x \to y = 0 + -x) \to x \in \mathbb{R} \to y \in \mathbb{R} \to x + y = 0 \to y = 0 + -x$
\item $\varphi_{90} :\equiv \varphi_{65} \to x \in \mathbb{R} \to y \in \mathbb{R} \to x + y = 0 \to y + x + -x = x + y + -x$
\item $\varphi_{91} :\equiv (x \in \mathbb{R} \to y \in \mathbb{R} \to x + y = 0 \to x \in \mathbb{R}) \to (x \in \mathbb{R} \to y \in \mathbb{R} \to x + y = 0 \to -x \in \mathbb{R}) \to \varphi_{40}$
\item $\varphi_{92} :\equiv (x \in \mathbb{R} \to y \in \mathbb{R} \to x + y = 0 \to y + 0 = y + (x + -x)) \to \varphi_{85}$
\item $\varphi_{93} :\equiv x + y + -x = x + y + -x \to \varphi_{46} \to x \in \mathbb{R} \to y \in \mathbb{R} \to x + y = 0 \to x + y + -x = 0 + -x$
\item $\varphi_{94} :\equiv \forall v_{31}\,\forall v_{32}\,\forall v_{33}\,(v_{31} \in \mathbb{R} \to v_{32} \in \mathbb{R} \to \varphi_{11}) \to \varphi_{64}$
\item $\varphi_{95} :\equiv (x \in \mathbb{R} \to y \in \mathbb{R} \to x + y = 0 \to x + y + -x = 0 + -x) \to \varphi_{89}$
\item $\varphi_{96} :\equiv (x \in \mathbb{R} \to y \in \mathbb{R} \to x + y = 0 \to 0 = x + -x) \to y + 0 = y + 0 \to \varphi_{88}$
\item $\varphi_{97} :\equiv (y \in \mathbb{R} \to x \in \mathbb{R} \to y + x = x + y) \to (x \in \mathbb{R} \to y \in \mathbb{R} \to x + y = 0 \to y \in \mathbb{R}) \to \varphi_{78}$
\item $\varphi_{98} :\equiv (x \in \mathbb{R} \to y \in \mathbb{R} \to x + y = 0 \to y = -x + 0) \to (x \in \mathbb{R} \to y \in \mathbb{R} \to x + y = 0 \to -x + 0 = -x) \to \varphi_{73}$
\item $\varphi_{99} :\equiv (x \in \mathbb{R} \to y \in \mathbb{R} \to x + y = 0 \to y + x + -x = x + y + -x) \to \varphi_{47} \to x \in \mathbb{R} \to y \in \mathbb{R} \to x + y = 0 \to y = x + y + -x$
\item $\varphi_{100} :\equiv (x \in \mathbb{R} \to y \in \mathbb{R} \to x + y = 0 \to y = 0 + -x) \to (x \in \mathbb{R} \to y \in \mathbb{R} \to x + y = 0 \to 0 + -x = -x + 0) \to \varphi_{83}$
\item $\varphi_{101} :\equiv (x \in \mathbb{R} \to y \in \mathbb{R} \to x + y = 0 \to y + x = x + y) \to y + x + -x = y + x + -x \to \varphi_{90}$
\item $\varphi_{102} :\equiv (x \in \mathbb{R} \to y \in \mathbb{R} \to x + y = 0 \to y = y + (x + -x)) \to \varphi_{41} \to \varphi_{48} \to x \in \mathbb{R} \to y \in \mathbb{R} \to x + y = 0 \to y = y + x + -x$
\end{abbrevlist}
\end{fullproof}

% 自動生成 (prover/build.js): NegNeg  ノード数 58
\begin{fullproof}{[F5]}{$\forall x \in \mathbb{R}\,(-(-x) = x)$}{58}\phantomsection\label{tree:NegNeg}
\noindent$\Pi_{1}$:
\begin{ptree}
\AXC{$x = x$\lbl{A5}}
\AXC{$x = y \to x = x \to y = x$\lbl{A6}}
\AXC{$(x = y \to x = x \to y = x) \to x = x \to x = y \to y = x$\lbl{A1}}
\BIC{$x = x \to x = y \to y = x$}
\BIC{$x = y \to y = x$}
\UIC{$\forall y\,(x = y \to y = x)$}
\UIC{$\forall x\,\forall y\,(x = y \to y = x)$}
\end{ptree}
\noindent$\Pi_{2}$:
\begin{ptree}
\AXC{$\stackrel{\Pi_{1}}{\vdots}$}\noLine
\UIC{$\forall x\,\forall y\,(x = y \to y = x)$}
\AXC{$\forall x\,\forall y\,(x = y \to y = x) \to \forall y\,(x = y \to y = x)$\lbl{A2}}
\BIC{$\forall y\,(x = y \to y = x)$}
\end{ptree}
\noindent$\Pi_{3}$:
\begin{ptree}
\AXC{$\stackrel{\Pi_{2}}{\vdots}$}\noLine
\UIC{$\forall y\,(x = y \to y = x)$}
\AXC{$\forall y\,(x = y \to y = x) \to x = -(-x) \to -(-x) = x$\lbl{A2}}
\BIC{$x = -(-x) \to -(-x) = x$}
\end{ptree}
\noindent$\Pi_{4}$:
\begin{ptree}
\AXC{$y = z \to x = y \to x = z$\lbl{A6}}
\AXC{$(y = z \to x = y \to x = z) \to x = y \to y = z \to x = z$\lbl{A1}}
\BIC{$x = y \to y = z \to x = z$}
\UIC{$\forall z\,(x = y \to y = z \to x = z)$}
\UIC{$\forall y\,\forall z\,(x = y \to y = z \to x = z)$}
\UIC{$\forall x\,\forall y\,\forall z\,(x = y \to y = z \to x = z)$}
\AXC{$\varphi_{1}$\lbl{A2}}
\BIC{$\forall y\,\forall z\,(-x + x = y \to y = z \to -x + x = z)$}
\AXC{$\varphi_{6}$\lbl{A2}}
\BIC{$\forall z\,(-x + x = x + -x \to x + -x = z \to -x + x = z)$}
\AXC{$\varphi_{9}$\lbl{A2}}
\BIC{$-x + x = x + -x \to x + -x = 0 \to -x + x = 0$}
\end{ptree}
\noindent$\Pi_{5}$:
\begin{ptree}
\AXC{$\forall x\,\forall y\,(x \in \mathbb{R} \to y \in \mathbb{R} \to x + y = y + x)$\lbl{(R3)}}
\AXC{$\varphi_{7}$\lbl{A2}}
\BIC{$\forall y\,(-x \in \mathbb{R} \to y \in \mathbb{R} \to -x + y = y + -x)$}
\AXC{$\varphi_{8}$\lbl{A2}}
\BIC{$-x \in \mathbb{R} \to x \in \mathbb{R} \to -x + x = x + -x$}
\AXC{$(-x \in \mathbb{R} \to x \in \mathbb{R} \to -x + x = x + -x) \to \varphi_{5}$\lbl{A1}}
\BIC{$\varphi_{5}$}
\end{ptree}
\noindent$\Pi_{6}$:
\begin{ptree}
\AXC{$x \in \mathbb{R} \to x \in \mathbb{R}$\lbl{A1}}
\AXC{$\forall x \in \mathbb{R}\,(-x \in \mathbb{R})$\lbl{(R1)}}
\AXC{$\forall x \in \mathbb{R}\,(-x \in \mathbb{R}) \to x \in \mathbb{R} \to -x \in \mathbb{R}$\lbl{A2}}
\BIC{$x \in \mathbb{R} \to -x \in \mathbb{R}$}
\AXC{$\stackrel{\Pi_{5}}{\vdots}$}\noLine
\UIC{$\varphi_{5}$}
\BIC{$(x \in \mathbb{R} \to x \in \mathbb{R}) \to x \in \mathbb{R} \to -x + x = x + -x$}
\BIC{$x \in \mathbb{R} \to -x + x = x + -x$}
\end{ptree}
\noindent$\Pi_{7}$:
\begin{ptree}
\AXC{$\forall x \in \mathbb{R}\,(x + -x = 0)$\lbl{(R5)}}
\AXC{$\forall x \in \mathbb{R}\,(x + -x = 0) \to x \in \mathbb{R} \to x + -x = 0$\lbl{A2}}
\BIC{$x \in \mathbb{R} \to x + -x = 0$}
\AXC{$\stackrel{\Pi_{6}}{\vdots}$}\noLine
\UIC{$x \in \mathbb{R} \to -x + x = x + -x$}
\AXC{$\varphi_{12}$\lbl{A1}}
\BIC{$(x \in \mathbb{R} \to x + -x = 0) \to \varphi_{3}$}
\BIC{$\varphi_{3}$}
\end{ptree}
\noindent$\Pi_{8}$:
\begin{ptree}
\AXC{$\forall x\,\forall y\,(x \in \mathbb{R} \to y \in \mathbb{R} \to x + y = 0 \to y = -x)$\lbl{[F4]}}
\AXC{$\varphi_{10}$\lbl{A2}}
\BIC{$\forall y\,(-x \in \mathbb{R} \to y \in \mathbb{R} \to -x + y = 0 \to y = -(-x))$}
\AXC{$\varphi_{11}$\lbl{A2}}
\BIC{$-x \in \mathbb{R} \to x \in \mathbb{R} \to -x + x = 0 \to x = -(-x)$}
\AXC{$\varphi_{13}$\lbl{A1}}
\BIC{$(x \in \mathbb{R} \to -x \in \mathbb{R}) \to \varphi_{2}$}
\end{ptree}
\noindent$\Pi_{9}$:
\begin{ptree}
\AXC{$\forall x \in \mathbb{R}\,(-x \in \mathbb{R})$\lbl{(R1)}}
\AXC{$\forall x \in \mathbb{R}\,(-x \in \mathbb{R}) \to x \in \mathbb{R} \to -x \in \mathbb{R}$\lbl{A2}}
\BIC{$x \in \mathbb{R} \to -x \in \mathbb{R}$}
\AXC{$\stackrel{\Pi_{8}}{\vdots}$}\noLine
\UIC{$(x \in \mathbb{R} \to -x \in \mathbb{R}) \to \varphi_{2}$}
\BIC{$\varphi_{2}$}
\end{ptree}
\noindent$\Pi_{10}$:
\begin{ptree}
\AXC{$\stackrel{\Pi_{4}}{\vdots}$}\noLine
\UIC{$-x + x = x + -x \to x + -x = 0 \to -x + x = 0$}
\AXC{$\stackrel{\Pi_{7}}{\vdots}$}\noLine
\UIC{$\varphi_{3}$}
\BIC{$x \in \mathbb{R} \to -x + x = 0$}
\AXC{$x \in \mathbb{R} \to x \in \mathbb{R}$\lbl{A1}}
\AXC{$\stackrel{\Pi_{9}}{\vdots}$}\noLine
\UIC{$\varphi_{2}$}
\BIC{$(x \in \mathbb{R} \to -x + x = 0) \to x \in \mathbb{R} \to x = -(-x)$}
\BIC{$x \in \mathbb{R} \to x = -(-x)$}
\AXC{$\varphi_{4}$\lbl{A1}}
\BIC{$(x = -(-x) \to -(-x) = x) \to x \in \mathbb{R} \to -(-x) = x$}
\end{ptree}
\noindent 主木:
\begin{ptree}
\AXC{$\stackrel{\Pi_{3}}{\vdots}$}\noLine
\UIC{$x = -(-x) \to -(-x) = x$}
\AXC{$\stackrel{\Pi_{10}}{\vdots}$}\noLine
\UIC{$(x = -(-x) \to -(-x) = x) \to x \in \mathbb{R} \to -(-x) = x$}
\BIC{$x \in \mathbb{R} \to -(-x) = x$}
\UIC{$\forall x \in \mathbb{R}\,(-(-x) = x)$}
\end{ptree}
\begin{abbrevlist}
\item $\varphi_{1} :\equiv \forall x\,\forall y\,\forall z\,(x = y \to y = z \to x = z) \to \forall y\,\forall z\,(-x + x = y \to y = z \to -x + x = z)$
\item $\varphi_{2} :\equiv (x \in \mathbb{R} \to x \in \mathbb{R}) \to (x \in \mathbb{R} \to -x + x = 0) \to x \in \mathbb{R} \to x = -(-x)$
\item $\varphi_{3} :\equiv (-x + x = x + -x \to x + -x = 0 \to -x + x = 0) \to x \in \mathbb{R} \to -x + x = 0$
\item $\varphi_{4} :\equiv (x \in \mathbb{R} \to x = -(-x)) \to (x = -(-x) \to -(-x) = x) \to x \in \mathbb{R} \to -(-x) = x$
\item $\varphi_{5} :\equiv (x \in \mathbb{R} \to -x \in \mathbb{R}) \to (x \in \mathbb{R} \to x \in \mathbb{R}) \to x \in \mathbb{R} \to -x + x = x + -x$
\item $\varphi_{6} :\equiv \forall y\,\forall z\,(-x + x = y \to y = z \to -x + x = z) \to \forall z\,(-x + x = x + -x \to x + -x = z \to -x + x = z)$
\item $\varphi_{7} :\equiv \forall x\,\forall y\,(x \in \mathbb{R} \to y \in \mathbb{R} \to x + y = y + x) \to \forall y\,(-x \in \mathbb{R} \to y \in \mathbb{R} \to -x + y = y + -x)$
\item $\varphi_{8} :\equiv \forall y\,(-x \in \mathbb{R} \to y \in \mathbb{R} \to -x + y = y + -x) \to -x \in \mathbb{R} \to x \in \mathbb{R} \to -x + x = x + -x$
\item $\varphi_{9} :\equiv \forall z\,(-x + x = x + -x \to x + -x = z \to -x + x = z) \to -x + x = x + -x \to x + -x = 0 \to -x + x = 0$
\item $\varphi_{10} :\equiv \forall x\,\forall y\,(x \in \mathbb{R} \to y \in \mathbb{R} \to x + y = 0 \to y = -x) \to \forall y\,(-x \in \mathbb{R} \to y \in \mathbb{R} \to -x + y = 0 \to y = -(-x))$
\item $\varphi_{11} :\equiv \forall y\,(-x \in \mathbb{R} \to y \in \mathbb{R} \to -x + y = 0 \to y = -(-x)) \to -x \in \mathbb{R} \to x \in \mathbb{R} \to -x + x = 0 \to x = -(-x)$
\item $\varphi_{12} :\equiv (x \in \mathbb{R} \to -x + x = x + -x) \to (x \in \mathbb{R} \to x + -x = 0) \to \varphi_{3}$
\item $\varphi_{13} :\equiv (-x \in \mathbb{R} \to x \in \mathbb{R} \to -x + x = 0 \to x = -(-x)) \to (x \in \mathbb{R} \to -x \in \mathbb{R}) \to \varphi_{2}$
\end{abbrevlist}
\end{fullproof}

% 自動生成 (prover/build.js): NegZero  ノード数 31
\begin{fullproof}{[F6]}{$-0 = 0$}{31}\phantomsection\label{tree:NegZero}
\noindent$\Pi_{1}$:
\begin{ptree}
\AXC{$x = x$\lbl{A5}}
\AXC{$x = y \to x = x \to y = x$\lbl{A6}}
\AXC{$(x = y \to x = x \to y = x) \to x = x \to x = y \to y = x$\lbl{A1}}
\BIC{$x = x \to x = y \to y = x$}
\BIC{$x = y \to y = x$}
\UIC{$\forall y\,(x = y \to y = x)$}
\UIC{$\forall x\,\forall y\,(x = y \to y = x)$}
\end{ptree}
\noindent$\Pi_{2}$:
\begin{ptree}
\AXC{$\stackrel{\Pi_{1}}{\vdots}$}\noLine
\UIC{$\forall x\,\forall y\,(x = y \to y = x)$}
\AXC{$\forall x\,\forall y\,(x = y \to y = x) \to \forall y\,(0 = y \to y = 0)$\lbl{A2}}
\BIC{$\forall y\,(0 = y \to y = 0)$}
\AXC{$\forall y\,(0 = y \to y = 0) \to 0 = -0 \to -0 = 0$\lbl{A2}}
\BIC{$0 = -0 \to -0 = 0$}
\end{ptree}
\noindent$\Pi_{3}$:
\begin{ptree}
\AXC{$\forall x \in \mathbb{R}\,(x + 0 = x)$\lbl{(R4)}}
\AXC{$\forall x \in \mathbb{R}\,(x + 0 = x) \to 0 \in \mathbb{R} \to 0 + 0 = 0$\lbl{A2}}
\BIC{$0 \in \mathbb{R} \to 0 + 0 = 0$}
\AXC{$(0 \in \mathbb{R} \to 0 + 0 = 0) \to 0 \in \mathbb{R} \to 0 + 0 = 0$\lbl{A1}}
\BIC{$0 \in \mathbb{R} \to 0 + 0 = 0$}
\end{ptree}
\noindent$\Pi_{4}$:
\begin{ptree}
\AXC{$0 \in \mathbb{R}$\lbl{(R1)}}
\AXC{$\forall x\,\forall y\,(x \in \mathbb{R} \to y \in \mathbb{R} \to x + y = 0 \to y = -x)$\lbl{[F4]}}
\AXC{$\varphi_{3}$\lbl{A2}}
\BIC{$\forall y\,(0 \in \mathbb{R} \to y \in \mathbb{R} \to 0 + y = 0 \to y = -0)$}
\AXC{$\varphi_{2}$\lbl{A2}}
\BIC{$0 \in \mathbb{R} \to 0 \in \mathbb{R} \to 0 + 0 = 0 \to 0 = -0$}
\AXC{$\varphi_{1}$\lbl{A1}}
\BIC{$0 \in \mathbb{R} \to 0 \in \mathbb{R} \to 0 + 0 = 0 \to 0 = -0$}
\BIC{$0 \in \mathbb{R} \to 0 + 0 = 0 \to 0 = -0$}
\end{ptree}
\noindent$\Pi_{5}$:
\begin{ptree}
\AXC{$0 \in \mathbb{R}$\lbl{(R1)}}
\AXC{$\stackrel{\Pi_{3}}{\vdots}$}\noLine
\UIC{$0 \in \mathbb{R} \to 0 + 0 = 0$}
\BIC{$0 + 0 = 0$}
\AXC{$0 \in \mathbb{R}$\lbl{(R1)}}
\AXC{$\stackrel{\Pi_{4}}{\vdots}$}\noLine
\UIC{$0 \in \mathbb{R} \to 0 + 0 = 0 \to 0 = -0$}
\BIC{$0 + 0 = 0 \to 0 = -0$}
\BIC{$0 = -0$}
\end{ptree}
\noindent 主木:
\begin{ptree}
\AXC{$\stackrel{\Pi_{2}}{\vdots}$}\noLine
\UIC{$0 = -0 \to -0 = 0$}
\AXC{$\stackrel{\Pi_{5}}{\vdots}$}\noLine
\UIC{$0 = -0$}
\AXC{$0 = -0 \to (0 = -0 \to -0 = 0) \to -0 = 0$\lbl{A1}}
\BIC{$(0 = -0 \to -0 = 0) \to -0 = 0$}
\BIC{$-0 = 0$}
\end{ptree}
\begin{abbrevlist}
\item $\varphi_{1} :\equiv (0 \in \mathbb{R} \to 0 \in \mathbb{R} \to 0 + 0 = 0 \to 0 = -0) \to 0 \in \mathbb{R} \to 0 \in \mathbb{R} \to 0 + 0 = 0 \to 0 = -0$
\item $\varphi_{2} :\equiv \forall y\,(0 \in \mathbb{R} \to y \in \mathbb{R} \to 0 + y = 0 \to y = -0) \to 0 \in \mathbb{R} \to 0 \in \mathbb{R} \to 0 + 0 = 0 \to 0 = -0$
\item $\varphi_{3} :\equiv \forall x\,\forall y\,(x \in \mathbb{R} \to y \in \mathbb{R} \to x + y = 0 \to y = -x) \to \forall y\,(0 \in \mathbb{R} \to y \in \mathbb{R} \to 0 + y = 0 \to y = -0)$
\end{abbrevlist}
\end{fullproof}

% 自動生成 (prover/build.js): NegAdd  ノード数 341
\begin{fullproof}{[F7]}{$\forall x\,\forall y\,(x \in \mathbb{R} \to y \in \mathbb{R} \to -(x + y) = -x + -y)$}{341}\phantomsection\label{tree:NegAdd}
\noindent$\Pi_{1}$:
\begin{ptree}
\AXC{$x = x$\lbl{A5}}
\AXC{$x = y \to x = x \to y = x$\lbl{A6}}
\AXC{$(x = y \to x = x \to y = x) \to x = x \to x = y \to y = x$\lbl{A1}}
\BIC{$x = x \to x = y \to y = x$}
\BIC{$x = y \to y = x$}
\UIC{$\forall y\,(x = y \to y = x)$}
\UIC{$\forall x\,\forall y\,(x = y \to y = x)$}
\end{ptree}
\noindent$\Pi_{2}$:
\begin{ptree}
\AXC{$\stackrel{\Pi_{1}}{\vdots}$}\noLine
\UIC{$\forall x\,\forall y\,(x = y \to y = x)$}
\AXC{$\forall x\,\forall y\,(x = y \to y = x) \to \forall y\,(v_{83} = y \to y = v_{83})$\lbl{A2}}
\BIC{$\forall y\,(v_{83} = y \to y = v_{83})$}
\AXC{$\varphi_{1}$\lbl{A2}}
\BIC{$v_{83} = v_{84} \to v_{84} = v_{83}$}
\UIC{$\forall v_{84}\,(v_{83} = v_{84} \to v_{84} = v_{83})$}
\UIC{$\forall v_{83}\,\forall v_{84}\,(v_{83} = v_{84} \to v_{84} = v_{83})$}
\AXC{$\varphi_{42}$\lbl{A2}}
\BIC{$\forall v_{84}\,(-x + -y = v_{84} \to v_{84} = -x + -y)$}
\AXC{$\varphi_{90}$\lbl{A2}}
\BIC{$-x + -y = -(x + y) \to -(x + y) = -x + -y$}
\end{ptree}
\noindent$\Pi_{3}$:
\begin{ptree}
\AXC{$y = z \to x = y \to x = z$\lbl{A6}}
\AXC{$(y = z \to x = y \to x = z) \to x = y \to y = z \to x = z$\lbl{A1}}
\BIC{$x = y \to y = z \to x = z$}
\UIC{$\forall z\,(x = y \to y = z \to x = z)$}
\UIC{$\forall y\,\forall z\,(x = y \to y = z \to x = z)$}
\UIC{$\forall x\,\forall y\,\forall z\,(x = y \to y = z \to x = z)$}
\AXC{$\varphi_{10}$\lbl{A2}}
\BIC{$\forall y\,\forall z\,(v_{78} = y \to y = z \to v_{78} = z)$}
\AXC{$\varphi_{16}$\lbl{A2}}
\BIC{$\forall z\,(v_{78} = v_{79} \to v_{79} = z \to v_{78} = z)$}
\AXC{$\varphi_{22}$\lbl{A2}}
\BIC{$v_{78} = v_{79} \to v_{79} = v_{80} \to v_{78} = v_{80}$}
\UIC{$\forall v_{80}\,(v_{78} = v_{79} \to v_{79} = v_{80} \to v_{78} = v_{80})$}
\UIC{$\varphi_{3}$}
\UIC{$\forall v_{78}\,(\varphi_{3})$}
\end{ptree}
\noindent$\Pi_{4}$:
\begin{ptree}
\AXC{$\stackrel{\Pi_{3}}{\vdots}$}\noLine
\UIC{$\forall v_{78}\,(\varphi_{3})$}
\AXC{$\forall v_{78}\,(\varphi_{3}) \to \varphi_{56}$\lbl{A2}}
\BIC{$\varphi_{56}$}
\AXC{$\varphi_{56} \to \varphi_{75}$\lbl{A2}}
\BIC{$\varphi_{75}$}
\AXC{$\varphi_{75} \to \varphi_{63}$\lbl{A2}}
\BIC{$\varphi_{63}$}
\end{ptree}
\noindent$\Pi_{5}$:
\begin{ptree}
\AXC{$\forall x \in \mathbb{R}\,(x + -x = 0)$\lbl{(R5)}}
\AXC{$\forall x \in \mathbb{R}\,(x + -x = 0) \to x \in \mathbb{R} \to x + -x = 0$\lbl{A2}}
\BIC{$x \in \mathbb{R} \to x + -x = 0$}
\AXC{$(x \in \mathbb{R} \to x + -x = 0) \to \varphi_{43}$\lbl{A1}}
\BIC{$\varphi_{43}$}
\end{ptree}
\noindent$\Pi_{6}$:
\begin{ptree}
\AXC{$y = z \to x = y \to x = z$\lbl{A6}}
\AXC{$(y = z \to x = y \to x = z) \to x = y \to y = z \to x = z$\lbl{A1}}
\BIC{$x = y \to y = z \to x = z$}
\UIC{$\forall z\,(x = y \to y = z \to x = z)$}
\UIC{$\forall y\,\forall z\,(x = y \to y = z \to x = z)$}
\UIC{$\forall x\,\forall y\,\forall z\,(x = y \to y = z \to x = z)$}
\AXC{$\varphi_{11}$\lbl{A2}}
\BIC{$\forall y\,\forall z\,(v_{75} = y \to y = z \to v_{75} = z)$}
\AXC{$\varphi_{17}$\lbl{A2}}
\BIC{$\forall z\,(v_{75} = v_{76} \to v_{76} = z \to v_{75} = z)$}
\AXC{$\varphi_{23}$\lbl{A2}}
\BIC{$v_{75} = v_{76} \to v_{76} = v_{77} \to v_{75} = v_{77}$}
\UIC{$\forall v_{77}\,(v_{75} = v_{76} \to v_{76} = v_{77} \to v_{75} = v_{77})$}
\UIC{$\varphi_{4}$}
\UIC{$\forall v_{75}\,(\varphi_{4})$}
\end{ptree}
\noindent$\Pi_{7}$:
\begin{ptree}
\AXC{$\stackrel{\Pi_{6}}{\vdots}$}\noLine
\UIC{$\forall v_{75}\,(\varphi_{4})$}
\AXC{$\forall v_{75}\,(\varphi_{4}) \to \varphi_{57}$\lbl{A2}}
\BIC{$\varphi_{57}$}
\AXC{$\varphi_{57} \to \varphi_{81}$\lbl{A2}}
\BIC{$\varphi_{81}$}
\AXC{$\varphi_{81} \to \varphi_{94}$\lbl{A2}}
\BIC{$\varphi_{94}$}
\end{ptree}
\noindent$\Pi_{8}$:
\begin{ptree}
\AXC{$\forall x \in \mathbb{R}\,(-x \in \mathbb{R})$\lbl{(R1)}}
\AXC{$\forall x \in \mathbb{R}\,(-x \in \mathbb{R}) \to x \in \mathbb{R} \to -x \in \mathbb{R}$\lbl{A2}}
\BIC{$x \in \mathbb{R} \to -x \in \mathbb{R}$}
\AXC{$(x \in \mathbb{R} \to -x \in \mathbb{R}) \to \varphi_{53}$\lbl{A1}}
\BIC{$\varphi_{53}$}
\end{ptree}
\noindent$\Pi_{9}$:
\begin{ptree}
\AXC{$x \in \mathbb{R} \to y \in \mathbb{R} \to x \in \mathbb{R}$\lbl{A1}}
\AXC{$\stackrel{\Pi_{8}}{\vdots}$}\noLine
\UIC{$\varphi_{53}$}
\BIC{$x \in \mathbb{R} \to y \in \mathbb{R} \to -x \in \mathbb{R}$}
\end{ptree}
\noindent$\Pi_{10}$:
\begin{ptree}
\AXC{$\forall x \in \mathbb{R}\,(x + 0 = x)$\lbl{(R4)}}
\AXC{$\forall x \in \mathbb{R}\,(x + 0 = x) \to -x \in \mathbb{R} \to -x + 0 = -x$\lbl{A2}}
\BIC{$-x \in \mathbb{R} \to -x + 0 = -x$}
\AXC{$(-x \in \mathbb{R} \to -x + 0 = -x) \to \varphi_{65}$\lbl{A1}}
\BIC{$\varphi_{65}$}
\end{ptree}
\noindent$\Pi_{11}$:
\begin{ptree}
\AXC{$x = x$\lbl{A5}}
\UIC{$\forall x\,(x = x)$}
\AXC{$\forall x\,(x = x) \to x + (-x + 0) = x + (-x + 0)$\lbl{A2}}
\BIC{$x + (-x + 0) = x + (-x + 0)$}
\AXC{$\stackrel{\Pi_{9}}{\vdots}$}\noLine
\UIC{$x \in \mathbb{R} \to y \in \mathbb{R} \to -x \in \mathbb{R}$}
\AXC{$\stackrel{\Pi_{10}}{\vdots}$}\noLine
\UIC{$\varphi_{65}$}
\BIC{$x \in \mathbb{R} \to y \in \mathbb{R} \to -x + 0 = -x$}
\AXC{$\varphi_{117}$\lbl{A1}}
\BIC{$x + (-x + 0) = x + (-x + 0) \to \varphi_{104}$}
\BIC{$\varphi_{104}$}
\end{ptree}
\noindent$\Pi_{12}$:
\begin{ptree}
\AXC{$z = w \to x + (-x + 0) = x + z \to x + (-x + 0) = x + w$\lbl{A6}}
\UIC{$\varphi_{29}$}
\UIC{$\forall z\,(\varphi_{29})$}
\AXC{$\forall z\,(\varphi_{29}) \to \varphi_{44}$\lbl{A2}}
\BIC{$\varphi_{44}$}
\AXC{$\varphi_{44} \to -x + 0 = -x \to \varphi_{34}$\lbl{A2}}
\BIC{$-x + 0 = -x \to \varphi_{34}$}
\AXC{$\stackrel{\Pi_{11}}{\vdots}$}\noLine
\UIC{$\varphi_{104}$}
\BIC{$x \in \mathbb{R} \to y \in \mathbb{R} \to x + (-x + 0) = x + -x$}
\end{ptree}
\noindent$\Pi_{13}$:
\begin{ptree}
\AXC{$y = z \to x = y \to x = z$\lbl{A6}}
\AXC{$(y = z \to x = y \to x = z) \to x = y \to y = z \to x = z$\lbl{A1}}
\BIC{$x = y \to y = z \to x = z$}
\UIC{$\forall z\,(x = y \to y = z \to x = z)$}
\UIC{$\forall y\,\forall z\,(x = y \to y = z \to x = z)$}
\UIC{$\forall x\,\forall y\,\forall z\,(x = y \to y = z \to x = z)$}
\AXC{$\varphi_{12}$\lbl{A2}}
\BIC{$\forall y\,\forall z\,(v_{72} = y \to y = z \to v_{72} = z)$}
\AXC{$\varphi_{18}$\lbl{A2}}
\BIC{$\forall z\,(v_{72} = v_{73} \to v_{73} = z \to v_{72} = z)$}
\AXC{$\varphi_{24}$\lbl{A2}}
\BIC{$v_{72} = v_{73} \to v_{73} = v_{74} \to v_{72} = v_{74}$}
\UIC{$\forall v_{74}\,(v_{72} = v_{73} \to v_{73} = v_{74} \to v_{72} = v_{74})$}
\UIC{$\varphi_{5}$}
\UIC{$\forall v_{72}\,(\varphi_{5})$}
\end{ptree}
\noindent$\Pi_{14}$:
\begin{ptree}
\AXC{$\forall x\,\forall y\,(x \in \mathbb{R} \to y \in \mathbb{R} \to x + y = y + x)$\lbl{(R3)}}
\AXC{$\varphi_{72}$\lbl{A2}}
\BIC{$\forall y\,(0 \in \mathbb{R} \to y \in \mathbb{R} \to 0 + y = y + 0)$}
\AXC{$\varphi_{78}$\lbl{A2}}
\BIC{$0 \in \mathbb{R} \to -x \in \mathbb{R} \to 0 + -x = -x + 0$}
\end{ptree}
\noindent$\Pi_{15}$:
\begin{ptree}
\AXC{$0 \in \mathbb{R}$\lbl{(R1)}}
\AXC{$\stackrel{\Pi_{14}}{\vdots}$}\noLine
\UIC{$0 \in \mathbb{R} \to -x \in \mathbb{R} \to 0 + -x = -x + 0$}
\AXC{$(0 \in \mathbb{R} \to -x \in \mathbb{R} \to 0 + -x = -x + 0) \to 0 \in \mathbb{R} \to \varphi_{74}$\lbl{A1}}
\BIC{$0 \in \mathbb{R} \to \varphi_{74}$}
\BIC{$\varphi_{74}$}
\end{ptree}
\noindent$\Pi_{16}$:
\begin{ptree}
\AXC{$x = x$\lbl{A5}}
\UIC{$\forall x\,(x = x)$}
\AXC{$\forall x\,(x = x) \to x + (0 + -x) = x + (0 + -x)$\lbl{A2}}
\BIC{$x + (0 + -x) = x + (0 + -x)$}
\AXC{$\stackrel{\Pi_{9}}{\vdots}$}\noLine
\UIC{$x \in \mathbb{R} \to y \in \mathbb{R} \to -x \in \mathbb{R}$}
\AXC{$\stackrel{\Pi_{15}}{\vdots}$}\noLine
\UIC{$\varphi_{74}$}
\BIC{$x \in \mathbb{R} \to y \in \mathbb{R} \to 0 + -x = -x + 0$}
\AXC{$\varphi_{120}$\lbl{A1}}
\BIC{$x + (0 + -x) = x + (0 + -x) \to \varphi_{106}$}
\BIC{$\varphi_{106}$}
\end{ptree}
\noindent$\Pi_{17}$:
\begin{ptree}
\AXC{$z = w \to x + (0 + -x) = x + z \to x + (0 + -x) = x + w$\lbl{A6}}
\UIC{$\varphi_{30}$}
\UIC{$\forall z\,(\varphi_{30})$}
\AXC{$\forall z\,(\varphi_{30}) \to \varphi_{45}$\lbl{A2}}
\BIC{$\varphi_{45}$}
\AXC{$\varphi_{45} \to 0 + -x = -x + 0 \to \varphi_{37}$\lbl{A2}}
\BIC{$0 + -x = -x + 0 \to \varphi_{37}$}
\AXC{$\stackrel{\Pi_{16}}{\vdots}$}\noLine
\UIC{$\varphi_{106}$}
\BIC{$x \in \mathbb{R} \to y \in \mathbb{R} \to x + (0 + -x) = x + (-x + 0)$}
\end{ptree}
\noindent$\Pi_{18}$:
\begin{ptree}
\AXC{$y = z \to x = y \to x = z$\lbl{A6}}
\AXC{$(y = z \to x = y \to x = z) \to x = y \to y = z \to x = z$\lbl{A1}}
\BIC{$x = y \to y = z \to x = z$}
\UIC{$\forall z\,(x = y \to y = z \to x = z)$}
\UIC{$\forall y\,\forall z\,(x = y \to y = z \to x = z)$}
\UIC{$\forall x\,\forall y\,\forall z\,(x = y \to y = z \to x = z)$}
\AXC{$\varphi_{13}$\lbl{A2}}
\BIC{$\forall y\,\forall z\,(v_{69} = y \to y = z \to v_{69} = z)$}
\AXC{$\varphi_{19}$\lbl{A2}}
\BIC{$\forall z\,(v_{69} = v_{70} \to v_{70} = z \to v_{69} = z)$}
\AXC{$\varphi_{25}$\lbl{A2}}
\BIC{$v_{69} = v_{70} \to v_{70} = v_{71} \to v_{69} = v_{71}$}
\UIC{$\forall v_{71}\,(v_{69} = v_{70} \to v_{70} = v_{71} \to v_{69} = v_{71})$}
\UIC{$\varphi_{6}$}
\UIC{$\forall v_{69}\,(\varphi_{6})$}
\end{ptree}
\noindent$\Pi_{19}$:
\begin{ptree}
\AXC{$\stackrel{\Pi_{18}}{\vdots}$}\noLine
\UIC{$\forall v_{69}\,(\varphi_{6})$}
\AXC{$\forall v_{69}\,(\varphi_{6}) \to \varphi_{59}$\lbl{A2}}
\BIC{$\varphi_{59}$}
\AXC{$\varphi_{59} \to \varphi_{100}$\lbl{A2}}
\BIC{$\varphi_{100}$}
\AXC{$\varphi_{100} \to x + y + (-x + -y) = x + (y + -y + -x) \to \varphi_{70}$\lbl{A2}}
\BIC{$x + y + (-x + -y) = x + (y + -y + -x) \to \varphi_{70}$}
\end{ptree}
\noindent$\Pi_{20}$:
\begin{ptree}
\AXC{$\forall x \in \mathbb{R}\,(x + -x = 0)$\lbl{(R5)}}
\AXC{$\forall x \in \mathbb{R}\,(x + -x = 0) \to y \in \mathbb{R} \to y + -y = 0$\lbl{A2}}
\BIC{$y \in \mathbb{R} \to y + -y = 0$}
\AXC{$(y \in \mathbb{R} \to y + -y = 0) \to \varphi_{46}$\lbl{A1}}
\BIC{$\varphi_{46}$}
\end{ptree}
\noindent$\Pi_{21}$:
\begin{ptree}
\AXC{$x = x$\lbl{A5}}
\UIC{$\forall x\,(x = x)$}
\AXC{$\varphi_{31}$\lbl{A2}}
\BIC{$x + (y + -y + -x) = x + (y + -y + -x)$}
\AXC{$x \in \mathbb{R} \to y \in \mathbb{R} \to y \in \mathbb{R}$\lbl{A1}}
\AXC{$\stackrel{\Pi_{20}}{\vdots}$}\noLine
\UIC{$\varphi_{46}$}
\BIC{$x \in \mathbb{R} \to y \in \mathbb{R} \to y + -y = 0$}
\AXC{$\varphi_{125}$\lbl{A1}}
\BIC{$x + (y + -y + -x) = x + (y + -y + -x) \to \varphi_{110}$}
\BIC{$\varphi_{110}$}
\end{ptree}
\noindent$\Pi_{22}$:
\begin{ptree}
\AXC{$z = w \to \varphi_{71}$\lbl{A6}}
\UIC{$\forall w\,(z = w \to \varphi_{71})$}
\UIC{$\forall z\,\forall w\,(z = w \to \varphi_{71})$}
\AXC{$\forall z\,\forall w\,(z = w \to \varphi_{71}) \to \varphi_{96}$\lbl{A2}}
\BIC{$\varphi_{96}$}
\AXC{$\varphi_{96} \to y + -y = 0 \to \varphi_{79}$\lbl{A2}}
\BIC{$y + -y = 0 \to \varphi_{79}$}
\AXC{$\stackrel{\Pi_{21}}{\vdots}$}\noLine
\UIC{$\varphi_{110}$}
\BIC{$x \in \mathbb{R} \to y \in \mathbb{R} \to x + (y + -y + -x) = x + (0 + -x)$}
\end{ptree}
\noindent$\Pi_{23}$:
\begin{ptree}
\AXC{$y = z \to x = y \to x = z$\lbl{A6}}
\AXC{$(y = z \to x = y \to x = z) \to x = y \to y = z \to x = z$\lbl{A1}}
\BIC{$x = y \to y = z \to x = z$}
\UIC{$\forall z\,(x = y \to y = z \to x = z)$}
\UIC{$\forall y\,\forall z\,(x = y \to y = z \to x = z)$}
\UIC{$\forall x\,\forall y\,\forall z\,(x = y \to y = z \to x = z)$}
\AXC{$\varphi_{14}$\lbl{A2}}
\BIC{$\forall y\,\forall z\,(v_{66} = y \to y = z \to v_{66} = z)$}
\AXC{$\varphi_{20}$\lbl{A2}}
\BIC{$\forall z\,(v_{66} = v_{67} \to v_{67} = z \to v_{66} = z)$}
\AXC{$\varphi_{26}$\lbl{A2}}
\BIC{$v_{66} = v_{67} \to v_{67} = v_{68} \to v_{66} = v_{68}$}
\UIC{$\forall v_{68}\,(v_{66} = v_{67} \to v_{67} = v_{68} \to v_{66} = v_{68})$}
\UIC{$\varphi_{7}$}
\UIC{$\forall v_{66}\,(\varphi_{7})$}
\end{ptree}
\noindent$\Pi_{24}$:
\begin{ptree}
\AXC{$\stackrel{\Pi_{23}}{\vdots}$}\noLine
\UIC{$\forall v_{66}\,(\varphi_{7})$}
\AXC{$\forall v_{66}\,(\varphi_{7}) \to \varphi_{60}$\lbl{A2}}
\BIC{$\varphi_{60}$}
\AXC{$\varphi_{60} \to \varphi_{101}$\lbl{A2}}
\BIC{$\varphi_{101}$}
\AXC{$\varphi_{101} \to x + y + (-x + -y) = x + (y + (-y + -x)) \to \varphi_{86}$\lbl{A2}}
\BIC{$x + y + (-x + -y) = x + (y + (-y + -x)) \to \varphi_{86}$}
\end{ptree}
\noindent$\Pi_{25}$:
\begin{ptree}
\AXC{$z = w \to \varphi_{38}$\lbl{A6}}
\UIC{$\forall w\,(z = w \to \varphi_{38})$}
\UIC{$\forall z\,\forall w\,(z = w \to \varphi_{38})$}
\AXC{$\forall z\,\forall w\,(z = w \to \varphi_{38}) \to \varphi_{97}$\lbl{A2}}
\BIC{$\varphi_{97}$}
\AXC{$\varphi_{97} \to y + (-y + -x) = y + -y + -x \to \varphi_{87}$\lbl{A2}}
\BIC{$y + (-y + -x) = y + -y + -x \to \varphi_{87}$}
\end{ptree}
\noindent$\Pi_{26}$:
\begin{ptree}
\AXC{$\stackrel{\Pi_{1}}{\vdots}$}\noLine
\UIC{$\forall x\,\forall y\,(x = y \to y = x)$}
\AXC{$\forall x\,\forall y\,(x = y \to y = x) \to \forall y\,(v_{61} = y \to y = v_{61})$\lbl{A2}}
\BIC{$\forall y\,(v_{61} = y \to y = v_{61})$}
\AXC{$\varphi_{2}$\lbl{A2}}
\BIC{$v_{61} = v_{62} \to v_{62} = v_{61}$}
\UIC{$\forall v_{62}\,(v_{61} = v_{62} \to v_{62} = v_{61})$}
\UIC{$\forall v_{61}\,\forall v_{62}\,(v_{61} = v_{62} \to v_{62} = v_{61})$}
\AXC{$\varphi_{66}$\lbl{A2}}
\BIC{$\forall v_{62}\,(y + -y + -x = v_{62} \to v_{62} = y + -y + -x)$}
\AXC{$\varphi_{105}$\lbl{A2}}
\BIC{$\varphi_{69}$}
\end{ptree}
\noindent$\Pi_{27}$:
\begin{ptree}
\AXC{$\forall x \in \mathbb{R}\,(-x \in \mathbb{R})$\lbl{(R1)}}
\AXC{$\forall x \in \mathbb{R}\,(-x \in \mathbb{R}) \to y \in \mathbb{R} \to -y \in \mathbb{R}$\lbl{A2}}
\BIC{$y \in \mathbb{R} \to -y \in \mathbb{R}$}
\AXC{$(y \in \mathbb{R} \to -y \in \mathbb{R}) \to \varphi_{54}$\lbl{A1}}
\BIC{$\varphi_{54}$}
\end{ptree}
\noindent$\Pi_{28}$:
\begin{ptree}
\AXC{$x \in \mathbb{R} \to y \in \mathbb{R} \to y \in \mathbb{R}$\lbl{A1}}
\AXC{$\stackrel{\Pi_{27}}{\vdots}$}\noLine
\UIC{$\varphi_{54}$}
\BIC{$x \in \mathbb{R} \to y \in \mathbb{R} \to -y \in \mathbb{R}$}
\end{ptree}
\noindent$\Pi_{29}$:
\begin{ptree}
\AXC{$\forall x\,\forall y\,(\varphi_{35})$\lbl{(R2)}}
\AXC{$\forall x\,\forall y\,(\varphi_{35}) \to \varphi_{40}$\lbl{A2}}
\BIC{$\varphi_{40}$}
\AXC{$\varphi_{40} \to \varphi_{41}$\lbl{A2}}
\BIC{$\varphi_{41}$}
\AXC{$\varphi_{41} \to v_{58} \in \mathbb{R} \to v_{59} \in \mathbb{R} \to \varphi_{9}$\lbl{A2}}
\BIC{$v_{58} \in \mathbb{R} \to v_{59} \in \mathbb{R} \to \varphi_{9}$}
\UIC{$\forall v_{60}\,(v_{58} \in \mathbb{R} \to v_{59} \in \mathbb{R} \to \varphi_{9})$}
\UIC{$\forall v_{59}\,\forall v_{60}\,(v_{58} \in \mathbb{R} \to v_{59} \in \mathbb{R} \to \varphi_{9})$}
\UIC{$\forall v_{58}\,\forall v_{59}\,\forall v_{60}\,(v_{58} \in \mathbb{R} \to v_{59} \in \mathbb{R} \to \varphi_{9})$}
\AXC{$\varphi_{108}$\lbl{A2}}
\BIC{$\varphi_{62}$}
\end{ptree}
\noindent$\Pi_{30}$:
\begin{ptree}
\AXC{$x \in \mathbb{R} \to y \in \mathbb{R} \to y \in \mathbb{R}$\lbl{A1}}
\AXC{$\stackrel{\Pi_{29}}{\vdots}$}\noLine
\UIC{$\varphi_{62}$}
\AXC{$\varphi_{62} \to \varphi_{67}$\lbl{A2}}
\BIC{$\varphi_{67}$}
\AXC{$\varphi_{67} \to \varphi_{73}$\lbl{A2}}
\BIC{$\varphi_{73}$}
\AXC{$\varphi_{73} \to \varphi_{111}$\lbl{A1}}
\BIC{$\varphi_{111}$}
\BIC{$(x \in \mathbb{R} \to y \in \mathbb{R} \to -y \in \mathbb{R}) \to \varphi_{92}$}
\end{ptree}
\noindent$\Pi_{31}$:
\begin{ptree}
\AXC{$\stackrel{\Pi_{9}}{\vdots}$}\noLine
\UIC{$x \in \mathbb{R} \to y \in \mathbb{R} \to -x \in \mathbb{R}$}
\AXC{$\stackrel{\Pi_{28}}{\vdots}$}\noLine
\UIC{$x \in \mathbb{R} \to y \in \mathbb{R} \to -y \in \mathbb{R}$}
\AXC{$\stackrel{\Pi_{30}}{\vdots}$}\noLine
\UIC{$(x \in \mathbb{R} \to y \in \mathbb{R} \to -y \in \mathbb{R}) \to \varphi_{92}$}
\BIC{$\varphi_{92}$}
\BIC{$x \in \mathbb{R} \to y \in \mathbb{R} \to y + -y + -x = y + (-y + -x)$}
\AXC{$\varphi_{118}$\lbl{A1}}
\BIC{$\varphi_{69} \to x \in \mathbb{R} \to y \in \mathbb{R} \to y + (-y + -x) = y + -y + -x$}
\end{ptree}
\noindent$\Pi_{32}$:
\begin{ptree}
\AXC{$x = x$\lbl{A5}}
\UIC{$\forall x\,(x = x)$}
\AXC{$\varphi_{32}$\lbl{A2}}
\BIC{$x + (y + (-y + -x)) = x + (y + (-y + -x))$}
\AXC{$\stackrel{\Pi_{26}}{\vdots}$}\noLine
\UIC{$\varphi_{69}$}
\AXC{$\stackrel{\Pi_{31}}{\vdots}$}\noLine
\UIC{$\varphi_{69} \to x \in \mathbb{R} \to y \in \mathbb{R} \to y + (-y + -x) = y + -y + -x$}
\BIC{$x \in \mathbb{R} \to y \in \mathbb{R} \to y + (-y + -x) = y + -y + -x$}
\AXC{$\varphi_{131}$\lbl{A1}}
\BIC{$\varphi_{127}$}
\BIC{$(y + (-y + -x) = y + -y + -x \to \varphi_{87}) \to \varphi_{47}$}
\end{ptree}
\noindent$\Pi_{33}$:
\begin{ptree}
\AXC{$\stackrel{\Pi_{25}}{\vdots}$}\noLine
\UIC{$y + (-y + -x) = y + -y + -x \to \varphi_{87}$}
\AXC{$\stackrel{\Pi_{32}}{\vdots}$}\noLine
\UIC{$(y + (-y + -x) = y + -y + -x \to \varphi_{87}) \to \varphi_{47}$}
\BIC{$\varphi_{47}$}
\end{ptree}
\noindent$\Pi_{34}$:
\begin{ptree}
\AXC{$y = z \to x = y \to x = z$\lbl{A6}}
\AXC{$(y = z \to x = y \to x = z) \to x = y \to y = z \to x = z$\lbl{A1}}
\BIC{$x = y \to y = z \to x = z$}
\UIC{$\forall z\,(x = y \to y = z \to x = z)$}
\UIC{$\forall y\,\forall z\,(x = y \to y = z \to x = z)$}
\UIC{$\forall x\,\forall y\,\forall z\,(x = y \to y = z \to x = z)$}
\AXC{$\varphi_{15}$\lbl{A2}}
\BIC{$\forall y\,\forall z\,(v_{63} = y \to y = z \to v_{63} = z)$}
\AXC{$\varphi_{21}$\lbl{A2}}
\BIC{$\forall z\,(v_{63} = v_{64} \to v_{64} = z \to v_{63} = z)$}
\AXC{$\varphi_{27}$\lbl{A2}}
\BIC{$v_{63} = v_{64} \to v_{64} = v_{65} \to v_{63} = v_{65}$}
\UIC{$\forall v_{65}\,(v_{63} = v_{64} \to v_{64} = v_{65} \to v_{63} = v_{65})$}
\UIC{$\varphi_{8}$}
\UIC{$\forall v_{63}\,(\varphi_{8})$}
\end{ptree}
\noindent$\Pi_{35}$:
\begin{ptree}
\AXC{$\stackrel{\Pi_{34}}{\vdots}$}\noLine
\UIC{$\forall v_{63}\,(\varphi_{8})$}
\AXC{$\forall v_{63}\,(\varphi_{8}) \to \varphi_{61}$\lbl{A2}}
\BIC{$\varphi_{61}$}
\AXC{$\varphi_{61} \to \varphi_{102}$\lbl{A2}}
\BIC{$\varphi_{102}$}
\AXC{$\varphi_{102} \to x + y + (-x + -y) = x + (y + (-x + -y)) \to \varphi_{88}$\lbl{A2}}
\BIC{$x + y + (-x + -y) = x + (y + (-x + -y)) \to \varphi_{88}$}
\end{ptree}
\noindent$\Pi_{36}$:
\begin{ptree}
\AXC{$z = w \to \varphi_{52}$\lbl{A6}}
\UIC{$\forall w\,(z = w \to \varphi_{52})$}
\UIC{$\forall z\,\forall w\,(z = w \to \varphi_{52})$}
\AXC{$\forall z\,\forall w\,(z = w \to \varphi_{52}) \to \varphi_{98}$\lbl{A2}}
\BIC{$\varphi_{98}$}
\AXC{$\varphi_{98} \to -x + -y = -y + -x \to \varphi_{89}$\lbl{A2}}
\BIC{$-x + -y = -y + -x \to \varphi_{89}$}
\end{ptree}
\noindent$\Pi_{37}$:
\begin{ptree}
\AXC{$\stackrel{\Pi_{9}}{\vdots}$}\noLine
\UIC{$x \in \mathbb{R} \to y \in \mathbb{R} \to -x \in \mathbb{R}$}
\AXC{$\forall x\,\forall y\,(x \in \mathbb{R} \to y \in \mathbb{R} \to x + y = y + x)$\lbl{(R3)}}
\AXC{$\varphi_{84}$\lbl{A2}}
\BIC{$\forall y\,(-x \in \mathbb{R} \to y \in \mathbb{R} \to -x + y = y + -x)$}
\AXC{$\varphi_{103}$\lbl{A2}}
\BIC{$-x \in \mathbb{R} \to -y \in \mathbb{R} \to -x + -y = -y + -x$}
\AXC{$\varphi_{116}$\lbl{A1}}
\BIC{$(x \in \mathbb{R} \to y \in \mathbb{R} \to -x \in \mathbb{R}) \to \varphi_{80}$}
\BIC{$\varphi_{80}$}
\end{ptree}
\noindent$\Pi_{38}$:
\begin{ptree}
\AXC{$\stackrel{\Pi_{28}}{\vdots}$}\noLine
\UIC{$x \in \mathbb{R} \to y \in \mathbb{R} \to -y \in \mathbb{R}$}
\AXC{$\stackrel{\Pi_{37}}{\vdots}$}\noLine
\UIC{$\varphi_{80}$}
\BIC{$x \in \mathbb{R} \to y \in \mathbb{R} \to -x + -y = -y + -x$}
\AXC{$(x \in \mathbb{R} \to y \in \mathbb{R} \to -x + -y = -y + -x) \to \varphi_{124}$\lbl{A1}}
\BIC{$\varphi_{124}$}
\end{ptree}
\noindent$\Pi_{39}$:
\begin{ptree}
\AXC{$\stackrel{\Pi_{9}}{\vdots}$}\noLine
\UIC{$x \in \mathbb{R} \to y \in \mathbb{R} \to -x \in \mathbb{R}$}
\AXC{$\forall x\,\forall y\,(x \in \mathbb{R} \to y \in \mathbb{R} \to x + y \in \mathbb{R})$\lbl{(R1)}}
\AXC{$\varphi_{83}$\lbl{A2}}
\BIC{$\forall y\,(-x \in \mathbb{R} \to y \in \mathbb{R} \to -x + y \in \mathbb{R})$}
\AXC{$\varphi_{99}$\lbl{A2}}
\BIC{$-x \in \mathbb{R} \to -y \in \mathbb{R} \to -x + -y \in \mathbb{R}$}
\AXC{$\varphi_{113}$\lbl{A1}}
\BIC{$(x \in \mathbb{R} \to y \in \mathbb{R} \to -x \in \mathbb{R}) \to \varphi_{77}$}
\BIC{$\varphi_{77}$}
\end{ptree}
\noindent$\Pi_{40}$:
\begin{ptree}
\AXC{$\stackrel{\Pi_{28}}{\vdots}$}\noLine
\UIC{$x \in \mathbb{R} \to y \in \mathbb{R} \to -y \in \mathbb{R}$}
\AXC{$\stackrel{\Pi_{39}}{\vdots}$}\noLine
\UIC{$\varphi_{77}$}
\BIC{$x \in \mathbb{R} \to y \in \mathbb{R} \to -x + -y \in \mathbb{R}$}
\end{ptree}
\noindent$\Pi_{41}$:
\begin{ptree}
\AXC{$\forall x\,\forall y\,(\varphi_{35})$\lbl{(R2)}}
\AXC{$\forall x\,\forall y\,(\varphi_{35}) \to \forall y\,(\varphi_{35})$\lbl{A2}}
\BIC{$\forall y\,(\varphi_{35})$}
\AXC{$\forall y\,(\varphi_{35}) \to \varphi_{35}$\lbl{A2}}
\BIC{$\varphi_{35}$}
\AXC{$\varphi_{35} \to x \in \mathbb{R} \to y \in \mathbb{R} \to \varphi_{55}$\lbl{A2}}
\BIC{$x \in \mathbb{R} \to y \in \mathbb{R} \to \varphi_{55}$}
\AXC{$\varphi_{129}$\lbl{A1}}
\BIC{$(x \in \mathbb{R} \to y \in \mathbb{R} \to x \in \mathbb{R}) \to \varphi_{107}$}
\end{ptree}
\noindent$\Pi_{42}$:
\begin{ptree}
\AXC{$x \in \mathbb{R} \to y \in \mathbb{R} \to y \in \mathbb{R}$\lbl{A1}}
\AXC{$x \in \mathbb{R} \to y \in \mathbb{R} \to x \in \mathbb{R}$\lbl{A1}}
\AXC{$\stackrel{\Pi_{41}}{\vdots}$}\noLine
\UIC{$(x \in \mathbb{R} \to y \in \mathbb{R} \to x \in \mathbb{R}) \to \varphi_{107}$}
\BIC{$\varphi_{107}$}
\BIC{$(x \in \mathbb{R} \to y \in \mathbb{R} \to -x + -y \in \mathbb{R}) \to \varphi_{48}$}
\end{ptree}
\noindent$\Pi_{43}$:
\begin{ptree}
\AXC{$\stackrel{\Pi_{40}}{\vdots}$}\noLine
\UIC{$x \in \mathbb{R} \to y \in \mathbb{R} \to -x + -y \in \mathbb{R}$}
\AXC{$\stackrel{\Pi_{42}}{\vdots}$}\noLine
\UIC{$(x \in \mathbb{R} \to y \in \mathbb{R} \to -x + -y \in \mathbb{R}) \to \varphi_{48}$}
\BIC{$\varphi_{48}$}
\AXC{$\varphi_{48} \to (x \in \mathbb{R} \to \varphi_{36}) \to \varphi_{121}$\lbl{A1}}
\BIC{$(x \in \mathbb{R} \to \varphi_{36}) \to \varphi_{121}$}
\end{ptree}
\noindent$\Pi_{44}$:
\begin{ptree}
\AXC{$\stackrel{\Pi_{36}}{\vdots}$}\noLine
\UIC{$-x + -y = -y + -x \to \varphi_{89}$}
\AXC{$x = x$\lbl{A5}}
\UIC{$\forall x\,(x = x)$}
\AXC{$\varphi_{33}$\lbl{A2}}
\BIC{$x + (y + (-x + -y)) = x + (y + (-x + -y))$}
\AXC{$\stackrel{\Pi_{38}}{\vdots}$}\noLine
\UIC{$\varphi_{124}$}
\BIC{$(-x + -y = -y + -x \to \varphi_{89}) \to x \in \mathbb{R} \to \varphi_{36}$}
\BIC{$x \in \mathbb{R} \to \varphi_{36}$}
\AXC{$\stackrel{\Pi_{43}}{\vdots}$}\noLine
\UIC{$(x \in \mathbb{R} \to \varphi_{36}) \to \varphi_{121}$}
\BIC{$\varphi_{121}$}
\end{ptree}
\noindent$\Pi_{45}$:
\begin{ptree}
\AXC{$\stackrel{\Pi_{33}}{\vdots}$}\noLine
\UIC{$\varphi_{47}$}
\AXC{$\stackrel{\Pi_{35}}{\vdots}$}\noLine
\UIC{$x + y + (-x + -y) = x + (y + (-x + -y)) \to \varphi_{88}$}
\AXC{$\stackrel{\Pi_{44}}{\vdots}$}\noLine
\UIC{$\varphi_{121}$}
\BIC{$\varphi_{49}$}
\AXC{$\varphi_{49} \to \varphi_{47} \to \varphi_{122}$\lbl{A1}}
\BIC{$\varphi_{47} \to \varphi_{122}$}
\BIC{$\varphi_{122}$}
\end{ptree}
\noindent$\Pi_{46}$:
\begin{ptree}
\AXC{$\stackrel{\Pi_{24}}{\vdots}$}\noLine
\UIC{$x + y + (-x + -y) = x + (y + (-y + -x)) \to \varphi_{86}$}
\AXC{$\stackrel{\Pi_{45}}{\vdots}$}\noLine
\UIC{$\varphi_{122}$}
\BIC{$\varphi_{50}$}
\AXC{$\varphi_{50} \to \varphi_{126}$\lbl{A1}}
\BIC{$\varphi_{126}$}
\end{ptree}
\noindent$\Pi_{47}$:
\begin{ptree}
\AXC{$\stackrel{\Pi_{19}}{\vdots}$}\noLine
\UIC{$x + y + (-x + -y) = x + (y + -y + -x) \to \varphi_{70}$}
\AXC{$\stackrel{\Pi_{22}}{\vdots}$}\noLine
\UIC{$x \in \mathbb{R} \to y \in \mathbb{R} \to x + (y + -y + -x) = x + (0 + -x)$}
\AXC{$\stackrel{\Pi_{46}}{\vdots}$}\noLine
\UIC{$\varphi_{126}$}
\BIC{$\varphi_{115}$}
\BIC{$x \in \mathbb{R} \to y \in \mathbb{R} \to x + y + (-x + -y) = x + (0 + -x)$}
\end{ptree}
\noindent$\Pi_{48}$:
\begin{ptree}
\AXC{$\stackrel{\Pi_{17}}{\vdots}$}\noLine
\UIC{$x \in \mathbb{R} \to y \in \mathbb{R} \to x + (0 + -x) = x + (-x + 0)$}
\AXC{$\stackrel{\Pi_{47}}{\vdots}$}\noLine
\UIC{$x \in \mathbb{R} \to y \in \mathbb{R} \to x + y + (-x + -y) = x + (0 + -x)$}
\AXC{$\varphi_{130}$\lbl{A1}}
\BIC{$\varphi_{123}$}
\BIC{$\varphi_{112}$}
\end{ptree}
\noindent$\Pi_{49}$:
\begin{ptree}
\AXC{$\stackrel{\Pi_{13}}{\vdots}$}\noLine
\UIC{$\forall v_{72}\,(\varphi_{5})$}
\AXC{$\forall v_{72}\,(\varphi_{5}) \to \varphi_{58}$\lbl{A2}}
\BIC{$\varphi_{58}$}
\AXC{$\varphi_{58} \to \varphi_{82}$\lbl{A2}}
\BIC{$\varphi_{82}$}
\AXC{$\varphi_{82} \to x + y + (-x + -y) = x + (0 + -x) \to \varphi_{51}$\lbl{A2}}
\BIC{$x + y + (-x + -y) = x + (0 + -x) \to \varphi_{51}$}
\AXC{$\stackrel{\Pi_{48}}{\vdots}$}\noLine
\UIC{$\varphi_{112}$}
\BIC{$x \in \mathbb{R} \to y \in \mathbb{R} \to x + y + (-x + -y) = x + (-x + 0)$}
\end{ptree}
\noindent$\Pi_{50}$:
\begin{ptree}
\AXC{$\stackrel{\Pi_{12}}{\vdots}$}\noLine
\UIC{$x \in \mathbb{R} \to y \in \mathbb{R} \to x + (-x + 0) = x + -x$}
\AXC{$\stackrel{\Pi_{49}}{\vdots}$}\noLine
\UIC{$x \in \mathbb{R} \to y \in \mathbb{R} \to x + y + (-x + -y) = x + (-x + 0)$}
\AXC{$\varphi_{128}$\lbl{A1}}
\BIC{$\varphi_{119}$}
\BIC{$\varphi_{94} \to x \in \mathbb{R} \to y \in \mathbb{R} \to x + y + (-x + -y) = x + -x$}
\end{ptree}
\noindent$\Pi_{51}$:
\begin{ptree}
\AXC{$\stackrel{\Pi_{7}}{\vdots}$}\noLine
\UIC{$\varphi_{94}$}
\AXC{$\stackrel{\Pi_{50}}{\vdots}$}\noLine
\UIC{$\varphi_{94} \to x \in \mathbb{R} \to y \in \mathbb{R} \to x + y + (-x + -y) = x + -x$}
\BIC{$x \in \mathbb{R} \to y \in \mathbb{R} \to x + y + (-x + -y) = x + -x$}
\AXC{$(x \in \mathbb{R} \to y \in \mathbb{R} \to x + y + (-x + -y) = x + -x) \to \varphi_{109}$\lbl{A1}}
\BIC{$\varphi_{109}$}
\end{ptree}
\noindent$\Pi_{52}$:
\begin{ptree}
\AXC{$\forall x\,\forall y\,(x \in \mathbb{R} \to y \in \mathbb{R} \to x + y \in \mathbb{R})$\lbl{(R1)}}
\AXC{$\varphi_{76}$\lbl{A2}}
\BIC{$\forall y\,(x \in \mathbb{R} \to y \in \mathbb{R} \to x + y \in \mathbb{R})$}
\AXC{$\varphi_{68}$\lbl{A2}}
\BIC{$x \in \mathbb{R} \to y \in \mathbb{R} \to x + y \in \mathbb{R}$}
\end{ptree}
\noindent$\Pi_{53}$:
\begin{ptree}
\AXC{$\forall x\,\forall y\,(x \in \mathbb{R} \to y \in \mathbb{R} \to x + y = 0 \to y = -x)$\lbl{[F4]}}
\AXC{$\varphi_{85}$\lbl{A2}}
\BIC{$\forall y\,(v_{81} \in \mathbb{R} \to y \in \mathbb{R} \to v_{81} + y = 0 \to y = -v_{81})$}
\AXC{$\varphi_{93}$\lbl{A2}}
\BIC{$\varphi_{28}$}
\UIC{$\forall v_{82}\,(\varphi_{28})$}
\UIC{$\forall v_{81}\,\forall v_{82}\,(\varphi_{28})$}
\AXC{$\forall v_{81}\,\forall v_{82}\,(\varphi_{28}) \to \varphi_{39}$\lbl{A2}}
\BIC{$\varphi_{39}$}
\end{ptree}
\noindent$\Pi_{54}$:
\begin{ptree}
\AXC{$\stackrel{\Pi_{52}}{\vdots}$}\noLine
\UIC{$x \in \mathbb{R} \to y \in \mathbb{R} \to x + y \in \mathbb{R}$}
\AXC{$\stackrel{\Pi_{53}}{\vdots}$}\noLine
\UIC{$\varphi_{39}$}
\AXC{$\varphi_{39} \to x + y \in \mathbb{R} \to \varphi_{64}$\lbl{A2}}
\BIC{$x + y \in \mathbb{R} \to \varphi_{64}$}
\AXC{$(x + y \in \mathbb{R} \to \varphi_{64}) \to \varphi_{114}$\lbl{A1}}
\BIC{$\varphi_{114}$}
\BIC{$(x \in \mathbb{R} \to y \in \mathbb{R} \to -x + -y \in \mathbb{R}) \to \varphi_{91}$}
\end{ptree}
\noindent$\Pi_{55}$:
\begin{ptree}
\AXC{$\stackrel{\Pi_{40}}{\vdots}$}\noLine
\UIC{$x \in \mathbb{R} \to y \in \mathbb{R} \to -x + -y \in \mathbb{R}$}
\AXC{$\stackrel{\Pi_{54}}{\vdots}$}\noLine
\UIC{$(x \in \mathbb{R} \to y \in \mathbb{R} \to -x + -y \in \mathbb{R}) \to \varphi_{91}$}
\BIC{$\varphi_{91}$}
\end{ptree}
\noindent$\Pi_{56}$:
\begin{ptree}
\AXC{$\stackrel{\Pi_{4}}{\vdots}$}\noLine
\UIC{$\varphi_{63}$}
\AXC{$x \in \mathbb{R} \to y \in \mathbb{R} \to x \in \mathbb{R}$\lbl{A1}}
\AXC{$\stackrel{\Pi_{5}}{\vdots}$}\noLine
\UIC{$\varphi_{43}$}
\BIC{$x \in \mathbb{R} \to y \in \mathbb{R} \to x + -x = 0$}
\AXC{$\stackrel{\Pi_{51}}{\vdots}$}\noLine
\UIC{$\varphi_{109}$}
\BIC{$\varphi_{63} \to x \in \mathbb{R} \to y \in \mathbb{R} \to x + y + (-x + -y) = 0$}
\BIC{$x \in \mathbb{R} \to y \in \mathbb{R} \to x + y + (-x + -y) = 0$}
\AXC{$\stackrel{\Pi_{55}}{\vdots}$}\noLine
\UIC{$\varphi_{91}$}
\BIC{$x \in \mathbb{R} \to y \in \mathbb{R} \to -x + -y = -(x + y)$}
\end{ptree}
\noindent$\Pi_{57}$:
\begin{ptree}
\AXC{$\stackrel{\Pi_{56}}{\vdots}$}\noLine
\UIC{$x \in \mathbb{R} \to y \in \mathbb{R} \to -x + -y = -(x + y)$}
\AXC{$(x \in \mathbb{R} \to y \in \mathbb{R} \to -x + -y = -(x + y)) \to \varphi_{95}$\lbl{A1}}
\BIC{$\varphi_{95}$}
\end{ptree}
\noindent 主木:
\begin{ptree}
\AXC{$\stackrel{\Pi_{2}}{\vdots}$}\noLine
\UIC{$-x + -y = -(x + y) \to -(x + y) = -x + -y$}
\AXC{$\stackrel{\Pi_{57}}{\vdots}$}\noLine
\UIC{$\varphi_{95}$}
\BIC{$x \in \mathbb{R} \to y \in \mathbb{R} \to -(x + y) = -x + -y$}
\UIC{$\forall y\,(x \in \mathbb{R} \to y \in \mathbb{R} \to -(x + y) = -x + -y)$}
\UIC{$\forall x\,\forall y\,(x \in \mathbb{R} \to y \in \mathbb{R} \to -(x + y) = -x + -y)$}
\end{ptree}
\begin{abbrevlist}
\item $\varphi_{1} :\equiv \forall y\,(v_{83} = y \to y = v_{83}) \to v_{83} = v_{84} \to v_{84} = v_{83}$
\item $\varphi_{2} :\equiv \forall y\,(v_{61} = y \to y = v_{61}) \to v_{61} = v_{62} \to v_{62} = v_{61}$
\item $\varphi_{3} :\equiv \forall v_{79}\,\forall v_{80}\,(v_{78} = v_{79} \to v_{79} = v_{80} \to v_{78} = v_{80})$
\item $\varphi_{4} :\equiv \forall v_{76}\,\forall v_{77}\,(v_{75} = v_{76} \to v_{76} = v_{77} \to v_{75} = v_{77})$
\item $\varphi_{5} :\equiv \forall v_{73}\,\forall v_{74}\,(v_{72} = v_{73} \to v_{73} = v_{74} \to v_{72} = v_{74})$
\item $\varphi_{6} :\equiv \forall v_{70}\,\forall v_{71}\,(v_{69} = v_{70} \to v_{70} = v_{71} \to v_{69} = v_{71})$
\item $\varphi_{7} :\equiv \forall v_{67}\,\forall v_{68}\,(v_{66} = v_{67} \to v_{67} = v_{68} \to v_{66} = v_{68})$
\item $\varphi_{8} :\equiv \forall v_{64}\,\forall v_{65}\,(v_{63} = v_{64} \to v_{64} = v_{65} \to v_{63} = v_{65})$
\item $\varphi_{9} :\equiv v_{60} \in \mathbb{R} \to v_{58} + v_{59} + v_{60} = v_{58} + (v_{59} + v_{60})$
\item $\varphi_{10} :\equiv \forall x\,\forall y\,\forall z\,(x = y \to y = z \to x = z) \to \forall y\,\forall z\,(v_{78} = y \to y = z \to v_{78} = z)$
\item $\varphi_{11} :\equiv \forall x\,\forall y\,\forall z\,(x = y \to y = z \to x = z) \to \forall y\,\forall z\,(v_{75} = y \to y = z \to v_{75} = z)$
\item $\varphi_{12} :\equiv \forall x\,\forall y\,\forall z\,(x = y \to y = z \to x = z) \to \forall y\,\forall z\,(v_{72} = y \to y = z \to v_{72} = z)$
\item $\varphi_{13} :\equiv \forall x\,\forall y\,\forall z\,(x = y \to y = z \to x = z) \to \forall y\,\forall z\,(v_{69} = y \to y = z \to v_{69} = z)$
\item $\varphi_{14} :\equiv \forall x\,\forall y\,\forall z\,(x = y \to y = z \to x = z) \to \forall y\,\forall z\,(v_{66} = y \to y = z \to v_{66} = z)$
\item $\varphi_{15} :\equiv \forall x\,\forall y\,\forall z\,(x = y \to y = z \to x = z) \to \forall y\,\forall z\,(v_{63} = y \to y = z \to v_{63} = z)$
\item $\varphi_{16} :\equiv \forall y\,\forall z\,(v_{78} = y \to y = z \to v_{78} = z) \to \forall z\,(v_{78} = v_{79} \to v_{79} = z \to v_{78} = z)$
\item $\varphi_{17} :\equiv \forall y\,\forall z\,(v_{75} = y \to y = z \to v_{75} = z) \to \forall z\,(v_{75} = v_{76} \to v_{76} = z \to v_{75} = z)$
\item $\varphi_{18} :\equiv \forall y\,\forall z\,(v_{72} = y \to y = z \to v_{72} = z) \to \forall z\,(v_{72} = v_{73} \to v_{73} = z \to v_{72} = z)$
\item $\varphi_{19} :\equiv \forall y\,\forall z\,(v_{69} = y \to y = z \to v_{69} = z) \to \forall z\,(v_{69} = v_{70} \to v_{70} = z \to v_{69} = z)$
\item $\varphi_{20} :\equiv \forall y\,\forall z\,(v_{66} = y \to y = z \to v_{66} = z) \to \forall z\,(v_{66} = v_{67} \to v_{67} = z \to v_{66} = z)$
\item $\varphi_{21} :\equiv \forall y\,\forall z\,(v_{63} = y \to y = z \to v_{63} = z) \to \forall z\,(v_{63} = v_{64} \to v_{64} = z \to v_{63} = z)$
\item $\varphi_{22} :\equiv \forall z\,(v_{78} = v_{79} \to v_{79} = z \to v_{78} = z) \to v_{78} = v_{79} \to v_{79} = v_{80} \to v_{78} = v_{80}$
\item $\varphi_{23} :\equiv \forall z\,(v_{75} = v_{76} \to v_{76} = z \to v_{75} = z) \to v_{75} = v_{76} \to v_{76} = v_{77} \to v_{75} = v_{77}$
\item $\varphi_{24} :\equiv \forall z\,(v_{72} = v_{73} \to v_{73} = z \to v_{72} = z) \to v_{72} = v_{73} \to v_{73} = v_{74} \to v_{72} = v_{74}$
\item $\varphi_{25} :\equiv \forall z\,(v_{69} = v_{70} \to v_{70} = z \to v_{69} = z) \to v_{69} = v_{70} \to v_{70} = v_{71} \to v_{69} = v_{71}$
\item $\varphi_{26} :\equiv \forall z\,(v_{66} = v_{67} \to v_{67} = z \to v_{66} = z) \to v_{66} = v_{67} \to v_{67} = v_{68} \to v_{66} = v_{68}$
\item $\varphi_{27} :\equiv \forall z\,(v_{63} = v_{64} \to v_{64} = z \to v_{63} = z) \to v_{63} = v_{64} \to v_{64} = v_{65} \to v_{63} = v_{65}$
\item $\varphi_{28} :\equiv v_{81} \in \mathbb{R} \to v_{82} \in \mathbb{R} \to v_{81} + v_{82} = 0 \to v_{82} = -v_{81}$
\item $\varphi_{29} :\equiv \forall w\,(z = w \to x + (-x + 0) = x + z \to x + (-x + 0) = x + w)$
\item $\varphi_{30} :\equiv \forall w\,(z = w \to x + (0 + -x) = x + z \to x + (0 + -x) = x + w)$
\item $\varphi_{31} :\equiv \forall x\,(x = x) \to x + (y + -y + -x) = x + (y + -y + -x)$
\item $\varphi_{32} :\equiv \forall x\,(x = x) \to x + (y + (-y + -x)) = x + (y + (-y + -x))$
\item $\varphi_{33} :\equiv \forall x\,(x = x) \to x + (y + (-x + -y)) = x + (y + (-x + -y))$
\item $\varphi_{34} :\equiv x + (-x + 0) = x + (-x + 0) \to x + (-x + 0) = x + -x$
\item $\varphi_{35} :\equiv \forall z\,(x \in \mathbb{R} \to y \in \mathbb{R} \to z \in \mathbb{R} \to x + y + z = x + (y + z))$
\item $\varphi_{36} :\equiv y \in \mathbb{R} \to x + (y + (-x + -y)) = x + (y + (-y + -x))$
\item $\varphi_{37} :\equiv x + (0 + -x) = x + (0 + -x) \to x + (0 + -x) = x + (-x + 0)$
\item $\varphi_{38} :\equiv x + (y + (-y + -x)) = x + z \to x + (y + (-y + -x)) = x + w$
\item $\varphi_{39} :\equiv \forall v_{82}\,(x + y \in \mathbb{R} \to v_{82} \in \mathbb{R} \to x + y + v_{82} = 0 \to v_{82} = -(x + y))$
\item $\varphi_{40} :\equiv \forall y\,\forall z\,(v_{58} \in \mathbb{R} \to y \in \mathbb{R} \to z \in \mathbb{R} \to v_{58} + y + z = v_{58} + (y + z))$
\item $\varphi_{41} :\equiv \forall z\,(v_{58} \in \mathbb{R} \to v_{59} \in \mathbb{R} \to z \in \mathbb{R} \to v_{58} + v_{59} + z = v_{58} + (v_{59} + z))$
\item $\varphi_{42} :\equiv \forall v_{83}\,\forall v_{84}\,(v_{83} = v_{84} \to v_{84} = v_{83}) \to \forall v_{84}\,(-x + -y = v_{84} \to v_{84} = -x + -y)$
\item $\varphi_{43} :\equiv (x \in \mathbb{R} \to y \in \mathbb{R} \to x \in \mathbb{R}) \to x \in \mathbb{R} \to y \in \mathbb{R} \to x + -x = 0$
\item $\varphi_{44} :\equiv \forall w\,(-x + 0 = w \to x + (-x + 0) = x + (-x + 0) \to x + (-x + 0) = x + w)$
\item $\varphi_{45} :\equiv \forall w\,(0 + -x = w \to x + (0 + -x) = x + (0 + -x) \to x + (0 + -x) = x + w)$
\item $\varphi_{46} :\equiv (x \in \mathbb{R} \to y \in \mathbb{R} \to y \in \mathbb{R}) \to x \in \mathbb{R} \to y \in \mathbb{R} \to y + -y = 0$
\item $\varphi_{47} :\equiv x \in \mathbb{R} \to y \in \mathbb{R} \to x + (y + (-y + -x)) = x + (y + -y + -x)$
\item $\varphi_{48} :\equiv x \in \mathbb{R} \to y \in \mathbb{R} \to x + y + (-x + -y) = x + (y + (-x + -y))$
\item $\varphi_{49} :\equiv x \in \mathbb{R} \to y \in \mathbb{R} \to x + y + (-x + -y) = x + (y + (-y + -x))$
\item $\varphi_{50} :\equiv x \in \mathbb{R} \to y \in \mathbb{R} \to x + y + (-x + -y) = x + (y + -y + -x)$
\item $\varphi_{51} :\equiv x + (0 + -x) = x + (-x + 0) \to x + y + (-x + -y) = x + (-x + 0)$
\item $\varphi_{52} :\equiv x + (y + (-x + -y)) = x + (y + z) \to x + (y + (-x + -y)) = x + (y + w)$
\item $\varphi_{53} :\equiv (x \in \mathbb{R} \to y \in \mathbb{R} \to x \in \mathbb{R}) \to x \in \mathbb{R} \to y \in \mathbb{R} \to -x \in \mathbb{R}$
\item $\varphi_{54} :\equiv (x \in \mathbb{R} \to y \in \mathbb{R} \to y \in \mathbb{R}) \to x \in \mathbb{R} \to y \in \mathbb{R} \to -y \in \mathbb{R}$
\item $\varphi_{55} :\equiv -x + -y \in \mathbb{R} \to x + y + (-x + -y) = x + (y + (-x + -y))$
\item $\varphi_{56} :\equiv \forall v_{79}\,\forall v_{80}\,(x + y + (-x + -y) = v_{79} \to v_{79} = v_{80} \to x + y + (-x + -y) = v_{80})$
\item $\varphi_{57} :\equiv \forall v_{76}\,\forall v_{77}\,(x + y + (-x + -y) = v_{76} \to v_{76} = v_{77} \to x + y + (-x + -y) = v_{77})$
\item $\varphi_{58} :\equiv \forall v_{73}\,\forall v_{74}\,(x + y + (-x + -y) = v_{73} \to v_{73} = v_{74} \to x + y + (-x + -y) = v_{74})$
\item $\varphi_{59} :\equiv \forall v_{70}\,\forall v_{71}\,(x + y + (-x + -y) = v_{70} \to v_{70} = v_{71} \to x + y + (-x + -y) = v_{71})$
\item $\varphi_{60} :\equiv \forall v_{67}\,\forall v_{68}\,(x + y + (-x + -y) = v_{67} \to v_{67} = v_{68} \to x + y + (-x + -y) = v_{68})$
\item $\varphi_{61} :\equiv \forall v_{64}\,\forall v_{65}\,(x + y + (-x + -y) = v_{64} \to v_{64} = v_{65} \to x + y + (-x + -y) = v_{65})$
\item $\varphi_{62} :\equiv \forall v_{59}\,\forall v_{60}\,(y \in \mathbb{R} \to v_{59} \in \mathbb{R} \to v_{60} \in \mathbb{R} \to y + v_{59} + v_{60} = y + (v_{59} + v_{60}))$
\item $\varphi_{63} :\equiv x + y + (-x + -y) = x + -x \to x + -x = 0 \to x + y + (-x + -y) = 0$
\item $\varphi_{64} :\equiv -x + -y \in \mathbb{R} \to x + y + (-x + -y) = 0 \to -x + -y = -(x + y)$
\item $\varphi_{65} :\equiv (x \in \mathbb{R} \to y \in \mathbb{R} \to -x \in \mathbb{R}) \to x \in \mathbb{R} \to y \in \mathbb{R} \to -x + 0 = -x$
\item $\varphi_{66} :\equiv \forall v_{61}\,\forall v_{62}\,(v_{61} = v_{62} \to v_{62} = v_{61}) \to \forall v_{62}\,(y + -y + -x = v_{62} \to v_{62} = y + -y + -x)$
\item $\varphi_{67} :\equiv \forall v_{60}\,(y \in \mathbb{R} \to -y \in \mathbb{R} \to v_{60} \in \mathbb{R} \to y + -y + v_{60} = y + (-y + v_{60}))$
\item $\varphi_{68} :\equiv \forall y\,(x \in \mathbb{R} \to y \in \mathbb{R} \to x + y \in \mathbb{R}) \to x \in \mathbb{R} \to y \in \mathbb{R} \to x + y \in \mathbb{R}$
\item $\varphi_{69} :\equiv y + -y + -x = y + (-y + -x) \to y + (-y + -x) = y + -y + -x$
\item $\varphi_{70} :\equiv x + (y + -y + -x) = x + (0 + -x) \to x + y + (-x + -y) = x + (0 + -x)$
\item $\varphi_{71} :\equiv x + (y + -y + -x) = x + (z + -x) \to x + (y + -y + -x) = x + (w + -x)$
\item $\varphi_{72} :\equiv \forall x\,\forall y\,(x \in \mathbb{R} \to y \in \mathbb{R} \to x + y = y + x) \to \forall y\,(0 \in \mathbb{R} \to y \in \mathbb{R} \to 0 + y = y + 0)$
\item $\varphi_{73} :\equiv y \in \mathbb{R} \to -y \in \mathbb{R} \to -x \in \mathbb{R} \to y + -y + -x = y + (-y + -x)$
\item $\varphi_{74} :\equiv (x \in \mathbb{R} \to y \in \mathbb{R} \to -x \in \mathbb{R}) \to x \in \mathbb{R} \to y \in \mathbb{R} \to 0 + -x = -x + 0$
\item $\varphi_{75} :\equiv \forall v_{80}\,(x + y + (-x + -y) = x + -x \to x + -x = v_{80} \to x + y + (-x + -y) = v_{80})$
\item $\varphi_{76} :\equiv \forall x\,\forall y\,(x \in \mathbb{R} \to y \in \mathbb{R} \to x + y \in \mathbb{R}) \to \forall y\,(x \in \mathbb{R} \to y \in \mathbb{R} \to x + y \in \mathbb{R})$
\item $\varphi_{77} :\equiv (x \in \mathbb{R} \to y \in \mathbb{R} \to -y \in \mathbb{R}) \to x \in \mathbb{R} \to y \in \mathbb{R} \to -x + -y \in \mathbb{R}$
\item $\varphi_{78} :\equiv \forall y\,(0 \in \mathbb{R} \to y \in \mathbb{R} \to 0 + y = y + 0) \to 0 \in \mathbb{R} \to -x \in \mathbb{R} \to 0 + -x = -x + 0$
\item $\varphi_{79} :\equiv x + (y + -y + -x) = x + (y + -y + -x) \to x + (y + -y + -x) = x + (0 + -x)$
\item $\varphi_{80} :\equiv (x \in \mathbb{R} \to y \in \mathbb{R} \to -y \in \mathbb{R}) \to x \in \mathbb{R} \to y \in \mathbb{R} \to -x + -y = -y + -x$
\item $\varphi_{81} :\equiv \forall v_{77}\,(x + y + (-x + -y) = x + (-x + 0) \to x + (-x + 0) = v_{77} \to x + y + (-x + -y) = v_{77})$
\item $\varphi_{82} :\equiv \forall v_{74}\,(x + y + (-x + -y) = x + (0 + -x) \to x + (0 + -x) = v_{74} \to x + y + (-x + -y) = v_{74})$
\item $\varphi_{83} :\equiv \forall x\,\forall y\,(x \in \mathbb{R} \to y \in \mathbb{R} \to x + y \in \mathbb{R}) \to \forall y\,(-x \in \mathbb{R} \to y \in \mathbb{R} \to -x + y \in \mathbb{R})$
\item $\varphi_{84} :\equiv \forall x\,\forall y\,(x \in \mathbb{R} \to y \in \mathbb{R} \to x + y = y + x) \to \forall y\,(-x \in \mathbb{R} \to y \in \mathbb{R} \to -x + y = y + -x)$
\item $\varphi_{85} :\equiv \forall x\,\forall y\,(x \in \mathbb{R} \to y \in \mathbb{R} \to x + y = 0 \to y = -x) \to \forall y\,(v_{81} \in \mathbb{R} \to y \in \mathbb{R} \to v_{81} + y = 0 \to y = -v_{81})$
\item $\varphi_{86} :\equiv x + (y + (-y + -x)) = x + (y + -y + -x) \to x + y + (-x + -y) = x + (y + -y + -x)$
\item $\varphi_{87} :\equiv x + (y + (-y + -x)) = x + (y + (-y + -x)) \to x + (y + (-y + -x)) = x + (y + -y + -x)$
\item $\varphi_{88} :\equiv x + (y + (-x + -y)) = x + (y + (-y + -x)) \to x + y + (-x + -y) = x + (y + (-y + -x))$
\item $\varphi_{89} :\equiv x + (y + (-x + -y)) = x + (y + (-x + -y)) \to x + (y + (-x + -y)) = x + (y + (-y + -x))$
\item $\varphi_{90} :\equiv \forall v_{84}\,(-x + -y = v_{84} \to v_{84} = -x + -y) \to -x + -y = -(x + y) \to -(x + y) = -x + -y$
\item $\varphi_{91} :\equiv (x \in \mathbb{R} \to y \in \mathbb{R} \to x + y + (-x + -y) = 0) \to x \in \mathbb{R} \to y \in \mathbb{R} \to -x + -y = -(x + y)$
\item $\varphi_{92} :\equiv (x \in \mathbb{R} \to y \in \mathbb{R} \to -x \in \mathbb{R}) \to x \in \mathbb{R} \to y \in \mathbb{R} \to y + -y + -x = y + (-y + -x)$
\item $\varphi_{93} :\equiv \forall y\,(v_{81} \in \mathbb{R} \to y \in \mathbb{R} \to v_{81} + y = 0 \to y = -v_{81}) \to \varphi_{28}$
\item $\varphi_{94} :\equiv x + y + (-x + -y) = x + (-x + 0) \to x + (-x + 0) = x + -x \to x + y + (-x + -y) = x + -x$
\item $\varphi_{95} :\equiv (-x + -y = -(x + y) \to -(x + y) = -x + -y) \to x \in \mathbb{R} \to y \in \mathbb{R} \to -(x + y) = -x + -y$
\item $\varphi_{96} :\equiv \forall w\,(y + -y = w \to x + (y + -y + -x) = x + (y + -y + -x) \to x + (y + -y + -x) = x + (w + -x))$
\item $\varphi_{97} :\equiv \forall w\,(y + (-y + -x) = w \to x + (y + (-y + -x)) = x + (y + (-y + -x)) \to x + (y + (-y + -x)) = x + w)$
\item $\varphi_{98} :\equiv \forall w\,(-x + -y = w \to x + (y + (-x + -y)) = x + (y + (-x + -y)) \to x + (y + (-x + -y)) = x + (y + w))$
\item $\varphi_{99} :\equiv \forall y\,(-x \in \mathbb{R} \to y \in \mathbb{R} \to -x + y \in \mathbb{R}) \to -x \in \mathbb{R} \to -y \in \mathbb{R} \to -x + -y \in \mathbb{R}$
\item $\varphi_{100} :\equiv \forall v_{71}\,(x + y + (-x + -y) = x + (y + -y + -x) \to x + (y + -y + -x) = v_{71} \to x + y + (-x + -y) = v_{71})$
\item $\varphi_{101} :\equiv \forall v_{68}\,(x + y + (-x + -y) = x + (y + (-y + -x)) \to x + (y + (-y + -x)) = v_{68} \to x + y + (-x + -y) = v_{68})$
\item $\varphi_{102} :\equiv \forall v_{65}\,(x + y + (-x + -y) = x + (y + (-x + -y)) \to x + (y + (-x + -y)) = v_{65} \to x + y + (-x + -y) = v_{65})$
\item $\varphi_{103} :\equiv \forall y\,(-x \in \mathbb{R} \to y \in \mathbb{R} \to -x + y = y + -x) \to -x \in \mathbb{R} \to -y \in \mathbb{R} \to -x + -y = -y + -x$
\item $\varphi_{104} :\equiv (-x + 0 = -x \to \varphi_{34}) \to x \in \mathbb{R} \to y \in \mathbb{R} \to x + (-x + 0) = x + -x$
\item $\varphi_{105} :\equiv \forall v_{62}\,(y + -y + -x = v_{62} \to v_{62} = y + -y + -x) \to \varphi_{69}$
\item $\varphi_{106} :\equiv (0 + -x = -x + 0 \to \varphi_{37}) \to x \in \mathbb{R} \to y \in \mathbb{R} \to x + (0 + -x) = x + (-x + 0)$
\item $\varphi_{107} :\equiv (x \in \mathbb{R} \to y \in \mathbb{R} \to y \in \mathbb{R}) \to (x \in \mathbb{R} \to y \in \mathbb{R} \to -x + -y \in \mathbb{R}) \to \varphi_{48}$
\item $\varphi_{108} :\equiv \forall v_{58}\,\forall v_{59}\,\forall v_{60}\,(v_{58} \in \mathbb{R} \to v_{59} \in \mathbb{R} \to \varphi_{9}) \to \varphi_{62}$
\item $\varphi_{109} :\equiv (x \in \mathbb{R} \to y \in \mathbb{R} \to x + -x = 0) \to \varphi_{63} \to x \in \mathbb{R} \to y \in \mathbb{R} \to x + y + (-x + -y) = 0$
\item $\varphi_{110} :\equiv (y + -y = 0 \to \varphi_{79}) \to x \in \mathbb{R} \to y \in \mathbb{R} \to x + (y + -y + -x) = x + (0 + -x)$
\item $\varphi_{111} :\equiv (x \in \mathbb{R} \to y \in \mathbb{R} \to y \in \mathbb{R}) \to (x \in \mathbb{R} \to y \in \mathbb{R} \to -y \in \mathbb{R}) \to \varphi_{92}$
\item $\varphi_{112} :\equiv (x + y + (-x + -y) = x + (0 + -x) \to \varphi_{51}) \to x \in \mathbb{R} \to y \in \mathbb{R} \to x + y + (-x + -y) = x + (-x + 0)$
\item $\varphi_{113} :\equiv (-x \in \mathbb{R} \to -y \in \mathbb{R} \to -x + -y \in \mathbb{R}) \to (x \in \mathbb{R} \to y \in \mathbb{R} \to -x \in \mathbb{R}) \to \varphi_{77}$
\item $\varphi_{114} :\equiv (x \in \mathbb{R} \to y \in \mathbb{R} \to x + y \in \mathbb{R}) \to (x \in \mathbb{R} \to y \in \mathbb{R} \to -x + -y \in \mathbb{R}) \to \varphi_{91}$
\item $\varphi_{115} :\equiv (x + y + (-x + -y) = x + (y + -y + -x) \to \varphi_{70}) \to x \in \mathbb{R} \to y \in \mathbb{R} \to x + y + (-x + -y) = x + (0 + -x)$
\item $\varphi_{116} :\equiv (-x \in \mathbb{R} \to -y \in \mathbb{R} \to -x + -y = -y + -x) \to (x \in \mathbb{R} \to y \in \mathbb{R} \to -x \in \mathbb{R}) \to \varphi_{80}$
\item $\varphi_{117} :\equiv (x \in \mathbb{R} \to y \in \mathbb{R} \to -x + 0 = -x) \to x + (-x + 0) = x + (-x + 0) \to \varphi_{104}$
\item $\varphi_{118} :\equiv (x \in \mathbb{R} \to y \in \mathbb{R} \to y + -y + -x = y + (-y + -x)) \to \varphi_{69} \to x \in \mathbb{R} \to y \in \mathbb{R} \to y + (-y + -x) = y + -y + -x$
\item $\varphi_{119} :\equiv (x \in \mathbb{R} \to y \in \mathbb{R} \to x + (-x + 0) = x + -x) \to \varphi_{94} \to x \in \mathbb{R} \to y \in \mathbb{R} \to x + y + (-x + -y) = x + -x$
\item $\varphi_{120} :\equiv (x \in \mathbb{R} \to y \in \mathbb{R} \to 0 + -x = -x + 0) \to x + (0 + -x) = x + (0 + -x) \to \varphi_{106}$
\item $\varphi_{121} :\equiv (x + y + (-x + -y) = x + (y + (-x + -y)) \to \varphi_{88}) \to \varphi_{49}$
\item $\varphi_{122} :\equiv (x + y + (-x + -y) = x + (y + (-y + -x)) \to \varphi_{86}) \to \varphi_{50}$
\item $\varphi_{123} :\equiv (x \in \mathbb{R} \to y \in \mathbb{R} \to x + (0 + -x) = x + (-x + 0)) \to \varphi_{112}$
\item $\varphi_{124} :\equiv x + (y + (-x + -y)) = x + (y + (-x + -y)) \to (-x + -y = -y + -x \to \varphi_{89}) \to x \in \mathbb{R} \to \varphi_{36}$
\item $\varphi_{125} :\equiv (x \in \mathbb{R} \to y \in \mathbb{R} \to y + -y = 0) \to x + (y + -y + -x) = x + (y + -y + -x) \to \varphi_{110}$
\item $\varphi_{126} :\equiv (x \in \mathbb{R} \to y \in \mathbb{R} \to x + (y + -y + -x) = x + (0 + -x)) \to \varphi_{115}$
\item $\varphi_{127} :\equiv x + (y + (-y + -x)) = x + (y + (-y + -x)) \to (y + (-y + -x) = y + -y + -x \to \varphi_{87}) \to \varphi_{47}$
\item $\varphi_{128} :\equiv (x \in \mathbb{R} \to y \in \mathbb{R} \to x + y + (-x + -y) = x + (-x + 0)) \to \varphi_{119}$
\item $\varphi_{129} :\equiv (x \in \mathbb{R} \to y \in \mathbb{R} \to \varphi_{55}) \to (x \in \mathbb{R} \to y \in \mathbb{R} \to x \in \mathbb{R}) \to \varphi_{107}$
\item $\varphi_{130} :\equiv (x \in \mathbb{R} \to y \in \mathbb{R} \to x + y + (-x + -y) = x + (0 + -x)) \to \varphi_{123}$
\item $\varphi_{131} :\equiv (x \in \mathbb{R} \to y \in \mathbb{R} \to y + (-y + -x) = y + -y + -x) \to \varphi_{127}$
\end{abbrevlist}
\end{fullproof}

% 自動生成 (prover/build.js): Mul0  ノード数 131
\begin{fullproof}{[F8]}{$\forall x \in \mathbb{R}\,(x \cdot 0 = 0)$}{131}\phantomsection\label{tree:Mul0}
\noindent$\Pi_{1}$:
\begin{ptree}
\AXC{$y = z \to x = y \to x = z$\lbl{A6}}
\AXC{$(y = z \to x = y \to x = z) \to x = y \to y = z \to x = z$\lbl{A1}}
\BIC{$x = y \to y = z \to x = z$}
\UIC{$\forall z\,(x = y \to y = z \to x = z)$}
\UIC{$\forall y\,\forall z\,(x = y \to y = z \to x = z)$}
\UIC{$\forall x\,\forall y\,\forall z\,(x = y \to y = z \to x = z)$}
\AXC{$\varphi_{20}$\lbl{A2}}
\BIC{$\forall y\,\forall z\,(x \cdot 0 + x \cdot 0 = y \to y = z \to x \cdot 0 + x \cdot 0 = z)$}
\AXC{$\varphi_{33}$\lbl{A2}}
\BIC{$\varphi_{10}$}
\AXC{$\varphi_{10} \to \varphi_{23}$\lbl{A2}}
\BIC{$\varphi_{23}$}
\end{ptree}
\noindent$\Pi_{2}$:
\begin{ptree}
\AXC{$x \in \mathbb{R} \to x \in \mathbb{R}$\lbl{A1}}
\AXC{$\forall x\,\forall y\,(x \in \mathbb{R} \to y \in \mathbb{R} \to x \cdot y \in \mathbb{R})$\lbl{(R1)}}
\AXC{$\varphi_{26}$\lbl{A2}}
\BIC{$\forall y\,(x \in \mathbb{R} \to y \in \mathbb{R} \to x \cdot y \in \mathbb{R})$}
\AXC{$\varphi_{24}$\lbl{A2}}
\BIC{$x \in \mathbb{R} \to 0 \in \mathbb{R} \to x \cdot 0 \in \mathbb{R}$}
\AXC{$\varphi_{29}$\lbl{A1}}
\BIC{$(x \in \mathbb{R} \to x \in \mathbb{R}) \to 0 \in \mathbb{R} \to x \in \mathbb{R} \to x \cdot 0 \in \mathbb{R}$}
\BIC{$0 \in \mathbb{R} \to x \in \mathbb{R} \to x \cdot 0 \in \mathbb{R}$}
\end{ptree}
\noindent$\Pi_{3}$:
\begin{ptree}
\AXC{$0 \in \mathbb{R}$\lbl{(R1)}}
\AXC{$\stackrel{\Pi_{2}}{\vdots}$}\noLine
\UIC{$0 \in \mathbb{R} \to x \in \mathbb{R} \to x \cdot 0 \in \mathbb{R}$}
\BIC{$x \in \mathbb{R} \to x \cdot 0 \in \mathbb{R}$}
\end{ptree}
\noindent$\Pi_{4}$:
\begin{ptree}
\AXC{$\forall x\,\forall y\,(x \in \mathbb{R} \to y \in \mathbb{R} \to x + y = y + x)$\lbl{(R3)}}
\AXC{$\varphi_{28}$\lbl{A2}}
\BIC{$\forall y\,(x \cdot 0 \in \mathbb{R} \to y \in \mathbb{R} \to x \cdot 0 + y = y + x \cdot 0)$}
\AXC{$\varphi_{31}$\lbl{A2}}
\BIC{$x \cdot 0 \in \mathbb{R} \to 0 \in \mathbb{R} \to x \cdot 0 + 0 = 0 + x \cdot 0$}
\AXC{$(x \cdot 0 \in \mathbb{R} \to 0 \in \mathbb{R} \to x \cdot 0 + 0 = 0 + x \cdot 0) \to \varphi_{22}$\lbl{A1}}
\BIC{$\varphi_{22}$}
\end{ptree}
\noindent$\Pi_{5}$:
\begin{ptree}
\AXC{$y = z \to x = y \to x = z$\lbl{A6}}
\AXC{$(y = z \to x = y \to x = z) \to x = y \to y = z \to x = z$\lbl{A1}}
\BIC{$x = y \to y = z \to x = z$}
\UIC{$\forall z\,(x = y \to y = z \to x = z)$}
\UIC{$\forall y\,\forall z\,(x = y \to y = z \to x = z)$}
\UIC{$\forall x\,\forall y\,\forall z\,(x = y \to y = z \to x = z)$}
\AXC{$\varphi_{20}$\lbl{A2}}
\BIC{$\forall y\,\forall z\,(x \cdot 0 + x \cdot 0 = y \to y = z \to x \cdot 0 + x \cdot 0 = z)$}
\AXC{$\varphi_{30}$\lbl{A2}}
\BIC{$\varphi_{5}$}
\AXC{$\varphi_{5} \to \varphi_{17}$\lbl{A2}}
\BIC{$\varphi_{17}$}
\end{ptree}
\noindent$\Pi_{6}$:
\begin{ptree}
\AXC{$x = x$\lbl{A5}}
\AXC{$x = y \to x = x \to y = x$\lbl{A6}}
\AXC{$(x = y \to x = x \to y = x) \to x = x \to x = y \to y = x$\lbl{A1}}
\BIC{$x = x \to x = y \to y = x$}
\BIC{$x = y \to y = x$}
\UIC{$\forall y\,(x = y \to y = x)$}
\UIC{$\forall x\,\forall y\,(x = y \to y = x)$}
\end{ptree}
\noindent$\Pi_{7}$:
\begin{ptree}
\AXC{$\stackrel{\Pi_{3}}{\vdots}$}\noLine
\UIC{$x \in \mathbb{R} \to x \cdot 0 \in \mathbb{R}$}
\AXC{$\forall x \in \mathbb{R}\,(x + 0 = x)$\lbl{(R4)}}
\AXC{$\forall x \in \mathbb{R}\,(x + 0 = x) \to x \cdot 0 \in \mathbb{R} \to x \cdot 0 + 0 = x \cdot 0$\lbl{A2}}
\BIC{$x \cdot 0 \in \mathbb{R} \to x \cdot 0 + 0 = x \cdot 0$}
\AXC{$\varphi_{27}$\lbl{A1}}
\BIC{$(x \in \mathbb{R} \to x \cdot 0 \in \mathbb{R}) \to x \in \mathbb{R} \to x \cdot 0 + 0 = x \cdot 0$}
\BIC{$x \in \mathbb{R} \to x \cdot 0 + 0 = x \cdot 0$}
\end{ptree}
\noindent$\Pi_{8}$:
\begin{ptree}
\AXC{$\stackrel{\Pi_{7}}{\vdots}$}\noLine
\UIC{$x \in \mathbb{R} \to x \cdot 0 + 0 = x \cdot 0$}
\AXC{$(x \in \mathbb{R} \to x \cdot 0 + 0 = x \cdot 0) \to \varphi_{14}$\lbl{A1}}
\BIC{$\varphi_{14}$}
\end{ptree}
\noindent$\Pi_{9}$:
\begin{ptree}
\AXC{$y = z \to x = y \to x = z$\lbl{A6}}
\AXC{$(y = z \to x = y \to x = z) \to x = y \to y = z \to x = z$\lbl{A1}}
\BIC{$x = y \to y = z \to x = z$}
\UIC{$\forall z\,(x = y \to y = z \to x = z)$}
\UIC{$\forall y\,\forall z\,(x = y \to y = z \to x = z)$}
\UIC{$\forall x\,\forall y\,\forall z\,(x = y \to y = z \to x = z)$}
\AXC{$\varphi_{20}$\lbl{A2}}
\BIC{$\forall y\,\forall z\,(x \cdot 0 + x \cdot 0 = y \to y = z \to x \cdot 0 + x \cdot 0 = z)$}
\AXC{$\varphi_{34}$\lbl{A2}}
\BIC{$\varphi_{11}$}
\AXC{$\varphi_{11} \to \varphi_{18}$\lbl{A2}}
\BIC{$\varphi_{18}$}
\end{ptree}
\noindent$\Pi_{10}$:
\begin{ptree}
\AXC{$z = w \to x \cdot (0 + 0) = x \cdot z \to x \cdot (0 + 0) = x \cdot w$\lbl{A6}}
\UIC{$\varphi_{1}$}
\UIC{$\forall z\,(\varphi_{1})$}
\AXC{$\forall z\,(\varphi_{1}) \to \varphi_{6}$\lbl{A2}}
\BIC{$\varphi_{6}$}
\AXC{$\varphi_{6} \to \varphi_{4}$\lbl{A2}}
\BIC{$\varphi_{4}$}
\end{ptree}
\noindent$\Pi_{11}$:
\begin{ptree}
\AXC{$\forall x \in \mathbb{R}\,(x + 0 = x)$\lbl{(R4)}}
\AXC{$\forall x \in \mathbb{R}\,(x + 0 = x) \to 0 \in \mathbb{R} \to 0 + 0 = 0$\lbl{A2}}
\BIC{$0 \in \mathbb{R} \to 0 + 0 = 0$}
\end{ptree}
\noindent$\Pi_{12}$:
\begin{ptree}
\AXC{$0 \in \mathbb{R}$\lbl{(R1)}}
\AXC{$\stackrel{\Pi_{11}}{\vdots}$}\noLine
\UIC{$0 \in \mathbb{R} \to 0 + 0 = 0$}
\AXC{$(0 \in \mathbb{R} \to 0 + 0 = 0) \to 0 \in \mathbb{R} \to x \in \mathbb{R} \to 0 + 0 = 0$\lbl{A1}}
\BIC{$0 \in \mathbb{R} \to x \in \mathbb{R} \to 0 + 0 = 0$}
\BIC{$x \in \mathbb{R} \to 0 + 0 = 0$}
\end{ptree}
\noindent$\Pi_{13}$:
\begin{ptree}
\AXC{$x = x$\lbl{A5}}
\UIC{$\forall x\,(x = x)$}
\AXC{$\forall x\,(x = x) \to x \cdot (0 + 0) = x \cdot (0 + 0)$\lbl{A2}}
\BIC{$x \cdot (0 + 0) = x \cdot (0 + 0)$}
\AXC{$\stackrel{\Pi_{12}}{\vdots}$}\noLine
\UIC{$x \in \mathbb{R} \to 0 + 0 = 0$}
\AXC{$(x \in \mathbb{R} \to 0 + 0 = 0) \to \varphi_{32}$\lbl{A1}}
\BIC{$\varphi_{32}$}
\BIC{$\varphi_{4} \to x \in \mathbb{R} \to x \cdot (0 + 0) = x \cdot 0$}
\end{ptree}
\noindent$\Pi_{14}$:
\begin{ptree}
\AXC{$\stackrel{\Pi_{6}}{\vdots}$}\noLine
\UIC{$\forall x\,\forall y\,(x = y \to y = x)$}
\AXC{$\varphi_{3}$\lbl{A2}}
\BIC{$\forall y\,(x \cdot (0 + 0) = y \to y = x \cdot (0 + 0))$}
\AXC{$\forall y\,(x \cdot (0 + 0) = y \to y = x \cdot (0 + 0)) \to \varphi_{7}$\lbl{A2}}
\BIC{$\varphi_{7}$}
\end{ptree}
\noindent$\Pi_{15}$:
\begin{ptree}
\AXC{$x \in \mathbb{R} \to x \in \mathbb{R}$\lbl{A1}}
\AXC{$\forall x\,(\varphi_{15})$\lbl{(R10)}}
\AXC{$\forall x\,(\varphi_{15}) \to \varphi_{15}$\lbl{A2}}
\BIC{$\varphi_{15}$}
\AXC{$\varphi_{15} \to \varphi_{12}$\lbl{A2}}
\BIC{$\varphi_{12}$}
\AXC{$\varphi_{35}$\lbl{A2}}
\BIC{$x \in \mathbb{R} \to 0 \in \mathbb{R} \to 0 \in \mathbb{R} \to x \cdot (0 + 0) = x \cdot 0 + x \cdot 0$}
\AXC{$\varphi_{40}$\lbl{A1}}
\BIC{$\varphi_{25}$}
\BIC{$0 \in \mathbb{R} \to 0 \in \mathbb{R} \to x \in \mathbb{R} \to x \cdot (0 + 0) = x \cdot 0 + x \cdot 0$}
\end{ptree}
\noindent$\Pi_{16}$:
\begin{ptree}
\AXC{$0 \in \mathbb{R}$\lbl{(R1)}}
\AXC{$0 \in \mathbb{R}$\lbl{(R1)}}
\AXC{$\stackrel{\Pi_{15}}{\vdots}$}\noLine
\UIC{$0 \in \mathbb{R} \to 0 \in \mathbb{R} \to x \in \mathbb{R} \to x \cdot (0 + 0) = x \cdot 0 + x \cdot 0$}
\BIC{$0 \in \mathbb{R} \to x \in \mathbb{R} \to x \cdot (0 + 0) = x \cdot 0 + x \cdot 0$}
\BIC{$x \in \mathbb{R} \to x \cdot (0 + 0) = x \cdot 0 + x \cdot 0$}
\AXC{$\varphi_{37}$\lbl{A1}}
\BIC{$\varphi_{7} \to x \in \mathbb{R} \to x \cdot 0 + x \cdot 0 = x \cdot (0 + 0)$}
\end{ptree}
\noindent$\Pi_{17}$:
\begin{ptree}
\AXC{$\stackrel{\Pi_{14}}{\vdots}$}\noLine
\UIC{$\varphi_{7}$}
\AXC{$\stackrel{\Pi_{16}}{\vdots}$}\noLine
\UIC{$\varphi_{7} \to x \in \mathbb{R} \to x \cdot 0 + x \cdot 0 = x \cdot (0 + 0)$}
\BIC{$x \in \mathbb{R} \to x \cdot 0 + x \cdot 0 = x \cdot (0 + 0)$}
\AXC{$(x \in \mathbb{R} \to x \cdot 0 + x \cdot 0 = x \cdot (0 + 0)) \to \varphi_{38}$\lbl{A1}}
\BIC{$\varphi_{38}$}
\end{ptree}
\noindent$\Pi_{18}$:
\begin{ptree}
\AXC{$\stackrel{\Pi_{9}}{\vdots}$}\noLine
\UIC{$\varphi_{18}$}
\AXC{$\stackrel{\Pi_{10}}{\vdots}$}\noLine
\UIC{$\varphi_{4}$}
\AXC{$\stackrel{\Pi_{13}}{\vdots}$}\noLine
\UIC{$\varphi_{4} \to x \in \mathbb{R} \to x \cdot (0 + 0) = x \cdot 0$}
\BIC{$x \in \mathbb{R} \to x \cdot (0 + 0) = x \cdot 0$}
\AXC{$\stackrel{\Pi_{17}}{\vdots}$}\noLine
\UIC{$\varphi_{38}$}
\BIC{$\varphi_{18} \to x \in \mathbb{R} \to x \cdot 0 + x \cdot 0 = x \cdot 0$}
\BIC{$x \in \mathbb{R} \to x \cdot 0 + x \cdot 0 = x \cdot 0$}
\AXC{$(x \in \mathbb{R} \to x \cdot 0 + x \cdot 0 = x \cdot 0) \to \varphi_{39}$\lbl{A1}}
\BIC{$\varphi_{39}$}
\end{ptree}
\noindent$\Pi_{19}$:
\begin{ptree}
\AXC{$\stackrel{\Pi_{5}}{\vdots}$}\noLine
\UIC{$\varphi_{17}$}
\AXC{$\stackrel{\Pi_{6}}{\vdots}$}\noLine
\UIC{$\forall x\,\forall y\,(x = y \to y = x)$}
\AXC{$\varphi_{2}$\lbl{A2}}
\BIC{$\forall y\,(x \cdot 0 + 0 = y \to y = x \cdot 0 + 0)$}
\AXC{$\varphi_{16}$\lbl{A2}}
\BIC{$x \cdot 0 + 0 = x \cdot 0 \to x \cdot 0 = x \cdot 0 + 0$}
\AXC{$\stackrel{\Pi_{8}}{\vdots}$}\noLine
\UIC{$\varphi_{14}$}
\BIC{$x \in \mathbb{R} \to x \cdot 0 = x \cdot 0 + 0$}
\AXC{$\stackrel{\Pi_{18}}{\vdots}$}\noLine
\UIC{$\varphi_{39}$}
\BIC{$\varphi_{17} \to x \in \mathbb{R} \to x \cdot 0 + x \cdot 0 = x \cdot 0 + 0$}
\BIC{$x \in \mathbb{R} \to x \cdot 0 + x \cdot 0 = x \cdot 0 + 0$}
\end{ptree}
\noindent$\Pi_{20}$:
\begin{ptree}
\AXC{$\stackrel{\Pi_{19}}{\vdots}$}\noLine
\UIC{$x \in \mathbb{R} \to x \cdot 0 + x \cdot 0 = x \cdot 0 + 0$}
\AXC{$(x \in \mathbb{R} \to x \cdot 0 + x \cdot 0 = x \cdot 0 + 0) \to \varphi_{41}$\lbl{A1}}
\BIC{$\varphi_{41}$}
\end{ptree}
\noindent$\Pi_{21}$:
\begin{ptree}
\AXC{$\varphi_{9}$\lbl{[F3]}}
\AXC{$\varphi_{9} \to \varphi_{21}$\lbl{A2}}
\BIC{$\varphi_{21}$}
\AXC{$\varphi_{21} \to \varphi_{19}$\lbl{A2}}
\BIC{$\varphi_{19}$}
\AXC{$\varphi_{19} \to x \cdot 0 \in \mathbb{R} \to \varphi_{13}$\lbl{A2}}
\BIC{$x \cdot 0 \in \mathbb{R} \to \varphi_{13}$}
\AXC{$(x \cdot 0 \in \mathbb{R} \to \varphi_{13}) \to \varphi_{36}$\lbl{A1}}
\BIC{$\varphi_{36}$}
\end{ptree}
\noindent$\Pi_{22}$:
\begin{ptree}
\AXC{$\stackrel{\Pi_{3}}{\vdots}$}\noLine
\UIC{$x \in \mathbb{R} \to x \cdot 0 \in \mathbb{R}$}
\AXC{$0 \in \mathbb{R}$\lbl{(R1)}}
\AXC{$\stackrel{\Pi_{3}}{\vdots}$}\noLine
\UIC{$x \in \mathbb{R} \to x \cdot 0 \in \mathbb{R}$}
\AXC{$\stackrel{\Pi_{21}}{\vdots}$}\noLine
\UIC{$\varphi_{36}$}
\BIC{$0 \in \mathbb{R} \to (x \in \mathbb{R} \to x \cdot 0 \in \mathbb{R}) \to \varphi_{8}$}
\BIC{$(x \in \mathbb{R} \to x \cdot 0 \in \mathbb{R}) \to \varphi_{8}$}
\BIC{$\varphi_{8}$}
\end{ptree}
\noindent 主木:
\begin{ptree}
\AXC{$\stackrel{\Pi_{1}}{\vdots}$}\noLine
\UIC{$\varphi_{23}$}
\AXC{$0 \in \mathbb{R}$\lbl{(R1)}}
\AXC{$\stackrel{\Pi_{3}}{\vdots}$}\noLine
\UIC{$x \in \mathbb{R} \to x \cdot 0 \in \mathbb{R}$}
\AXC{$\stackrel{\Pi_{4}}{\vdots}$}\noLine
\UIC{$\varphi_{22}$}
\BIC{$0 \in \mathbb{R} \to x \in \mathbb{R} \to x \cdot 0 + 0 = 0 + x \cdot 0$}
\BIC{$x \in \mathbb{R} \to x \cdot 0 + 0 = 0 + x \cdot 0$}
\AXC{$\stackrel{\Pi_{20}}{\vdots}$}\noLine
\UIC{$\varphi_{41}$}
\BIC{$\varphi_{23} \to x \in \mathbb{R} \to x \cdot 0 + x \cdot 0 = 0 + x \cdot 0$}
\BIC{$x \in \mathbb{R} \to x \cdot 0 + x \cdot 0 = 0 + x \cdot 0$}
\AXC{$\stackrel{\Pi_{22}}{\vdots}$}\noLine
\UIC{$\varphi_{8}$}
\BIC{$x \in \mathbb{R} \to x \cdot 0 = 0$}
\UIC{$\forall x \in \mathbb{R}\,(x \cdot 0 = 0)$}
\end{ptree}
\begin{abbrevlist}
\item $\varphi_{1} :\equiv \forall w\,(z = w \to x \cdot (0 + 0) = x \cdot z \to x \cdot (0 + 0) = x \cdot w)$
\item $\varphi_{2} :\equiv \forall x\,\forall y\,(x = y \to y = x) \to \forall y\,(x \cdot 0 + 0 = y \to y = x \cdot 0 + 0)$
\item $\varphi_{3} :\equiv \forall x\,\forall y\,(x = y \to y = x) \to \forall y\,(x \cdot (0 + 0) = y \to y = x \cdot (0 + 0))$
\item $\varphi_{4} :\equiv 0 + 0 = 0 \to x \cdot (0 + 0) = x \cdot (0 + 0) \to x \cdot (0 + 0) = x \cdot 0$
\item $\varphi_{5} :\equiv \forall z\,(x \cdot 0 + x \cdot 0 = x \cdot 0 \to x \cdot 0 = z \to x \cdot 0 + x \cdot 0 = z)$
\item $\varphi_{6} :\equiv \forall w\,(0 + 0 = w \to x \cdot (0 + 0) = x \cdot (0 + 0) \to x \cdot (0 + 0) = x \cdot w)$
\item $\varphi_{7} :\equiv x \cdot (0 + 0) = x \cdot 0 + x \cdot 0 \to x \cdot 0 + x \cdot 0 = x \cdot (0 + 0)$
\item $\varphi_{8} :\equiv (x \in \mathbb{R} \to x \cdot 0 + x \cdot 0 = 0 + x \cdot 0) \to x \in \mathbb{R} \to x \cdot 0 = 0$
\item $\varphi_{9} :\equiv \forall x\,\forall y\,\forall z\,(x \in \mathbb{R} \to y \in \mathbb{R} \to z \in \mathbb{R} \to x + z = y + z \to x = y)$
\item $\varphi_{10} :\equiv \forall z\,(x \cdot 0 + x \cdot 0 = x \cdot 0 + 0 \to x \cdot 0 + 0 = z \to x \cdot 0 + x \cdot 0 = z)$
\item $\varphi_{11} :\equiv \forall z\,(x \cdot 0 + x \cdot 0 = x \cdot (0 + 0) \to x \cdot (0 + 0) = z \to x \cdot 0 + x \cdot 0 = z)$
\item $\varphi_{12} :\equiv \forall z\,(x \in \mathbb{R} \to 0 \in \mathbb{R} \to z \in \mathbb{R} \to x \cdot (0 + z) = x \cdot 0 + x \cdot z)$
\item $\varphi_{13} :\equiv 0 \in \mathbb{R} \to x \cdot 0 \in \mathbb{R} \to x \cdot 0 + x \cdot 0 = 0 + x \cdot 0 \to x \cdot 0 = 0$
\item $\varphi_{14} :\equiv (x \cdot 0 + 0 = x \cdot 0 \to x \cdot 0 = x \cdot 0 + 0) \to x \in \mathbb{R} \to x \cdot 0 = x \cdot 0 + 0$
\item $\varphi_{15} :\equiv \forall y\,\forall z\,(x \in \mathbb{R} \to y \in \mathbb{R} \to z \in \mathbb{R} \to x \cdot (y + z) = x \cdot y + x \cdot z)$
\item $\varphi_{16} :\equiv \forall y\,(x \cdot 0 + 0 = y \to y = x \cdot 0 + 0) \to x \cdot 0 + 0 = x \cdot 0 \to x \cdot 0 = x \cdot 0 + 0$
\item $\varphi_{17} :\equiv x \cdot 0 + x \cdot 0 = x \cdot 0 \to x \cdot 0 = x \cdot 0 + 0 \to x \cdot 0 + x \cdot 0 = x \cdot 0 + 0$
\item $\varphi_{18} :\equiv x \cdot 0 + x \cdot 0 = x \cdot (0 + 0) \to x \cdot (0 + 0) = x \cdot 0 \to x \cdot 0 + x \cdot 0 = x \cdot 0$
\item $\varphi_{19} :\equiv \forall z\,(x \cdot 0 \in \mathbb{R} \to 0 \in \mathbb{R} \to z \in \mathbb{R} \to x \cdot 0 + z = 0 + z \to x \cdot 0 = 0)$
\item $\varphi_{20} :\equiv \forall x\,\forall y\,\forall z\,(x = y \to y = z \to x = z) \to \forall y\,\forall z\,(x \cdot 0 + x \cdot 0 = y \to y = z \to x \cdot 0 + x \cdot 0 = z)$
\item $\varphi_{21} :\equiv \forall y\,\forall z\,(x \cdot 0 \in \mathbb{R} \to y \in \mathbb{R} \to z \in \mathbb{R} \to x \cdot 0 + z = y + z \to x \cdot 0 = y)$
\item $\varphi_{22} :\equiv (x \in \mathbb{R} \to x \cdot 0 \in \mathbb{R}) \to 0 \in \mathbb{R} \to x \in \mathbb{R} \to x \cdot 0 + 0 = 0 + x \cdot 0$
\item $\varphi_{23} :\equiv x \cdot 0 + x \cdot 0 = x \cdot 0 + 0 \to x \cdot 0 + 0 = 0 + x \cdot 0 \to x \cdot 0 + x \cdot 0 = 0 + x \cdot 0$
\item $\varphi_{24} :\equiv \forall y\,(x \in \mathbb{R} \to y \in \mathbb{R} \to x \cdot y \in \mathbb{R}) \to x \in \mathbb{R} \to 0 \in \mathbb{R} \to x \cdot 0 \in \mathbb{R}$
\item $\varphi_{25} :\equiv (x \in \mathbb{R} \to x \in \mathbb{R}) \to 0 \in \mathbb{R} \to 0 \in \mathbb{R} \to x \in \mathbb{R} \to x \cdot (0 + 0) = x \cdot 0 + x \cdot 0$
\item $\varphi_{26} :\equiv \forall x\,\forall y\,(x \in \mathbb{R} \to y \in \mathbb{R} \to x \cdot y \in \mathbb{R}) \to \forall y\,(x \in \mathbb{R} \to y \in \mathbb{R} \to x \cdot y \in \mathbb{R})$
\item $\varphi_{27} :\equiv (x \cdot 0 \in \mathbb{R} \to x \cdot 0 + 0 = x \cdot 0) \to (x \in \mathbb{R} \to x \cdot 0 \in \mathbb{R}) \to x \in \mathbb{R} \to x \cdot 0 + 0 = x \cdot 0$
\item $\varphi_{28} :\equiv \forall x\,\forall y\,(x \in \mathbb{R} \to y \in \mathbb{R} \to x + y = y + x) \to \forall y\,(x \cdot 0 \in \mathbb{R} \to y \in \mathbb{R} \to x \cdot 0 + y = y + x \cdot 0)$
\item $\varphi_{29} :\equiv (x \in \mathbb{R} \to 0 \in \mathbb{R} \to x \cdot 0 \in \mathbb{R}) \to (x \in \mathbb{R} \to x \in \mathbb{R}) \to 0 \in \mathbb{R} \to x \in \mathbb{R} \to x \cdot 0 \in \mathbb{R}$
\item $\varphi_{30} :\equiv \forall y\,\forall z\,(x \cdot 0 + x \cdot 0 = y \to y = z \to x \cdot 0 + x \cdot 0 = z) \to \varphi_{5}$
\item $\varphi_{31} :\equiv \forall y\,(x \cdot 0 \in \mathbb{R} \to y \in \mathbb{R} \to x \cdot 0 + y = y + x \cdot 0) \to x \cdot 0 \in \mathbb{R} \to 0 \in \mathbb{R} \to x \cdot 0 + 0 = 0 + x \cdot 0$
\item $\varphi_{32} :\equiv x \cdot (0 + 0) = x \cdot (0 + 0) \to \varphi_{4} \to x \in \mathbb{R} \to x \cdot (0 + 0) = x \cdot 0$
\item $\varphi_{33} :\equiv \forall y\,\forall z\,(x \cdot 0 + x \cdot 0 = y \to y = z \to x \cdot 0 + x \cdot 0 = z) \to \varphi_{10}$
\item $\varphi_{34} :\equiv \forall y\,\forall z\,(x \cdot 0 + x \cdot 0 = y \to y = z \to x \cdot 0 + x \cdot 0 = z) \to \varphi_{11}$
\item $\varphi_{35} :\equiv \varphi_{12} \to x \in \mathbb{R} \to 0 \in \mathbb{R} \to 0 \in \mathbb{R} \to x \cdot (0 + 0) = x \cdot 0 + x \cdot 0$
\item $\varphi_{36} :\equiv (x \in \mathbb{R} \to x \cdot 0 \in \mathbb{R}) \to 0 \in \mathbb{R} \to (x \in \mathbb{R} \to x \cdot 0 \in \mathbb{R}) \to \varphi_{8}$
\item $\varphi_{37} :\equiv (x \in \mathbb{R} \to x \cdot (0 + 0) = x \cdot 0 + x \cdot 0) \to \varphi_{7} \to x \in \mathbb{R} \to x \cdot 0 + x \cdot 0 = x \cdot (0 + 0)$
\item $\varphi_{38} :\equiv (x \in \mathbb{R} \to x \cdot (0 + 0) = x \cdot 0) \to \varphi_{18} \to x \in \mathbb{R} \to x \cdot 0 + x \cdot 0 = x \cdot 0$
\item $\varphi_{39} :\equiv (x \in \mathbb{R} \to x \cdot 0 = x \cdot 0 + 0) \to \varphi_{17} \to x \in \mathbb{R} \to x \cdot 0 + x \cdot 0 = x \cdot 0 + 0$
\item $\varphi_{40} :\equiv (x \in \mathbb{R} \to 0 \in \mathbb{R} \to 0 \in \mathbb{R} \to x \cdot (0 + 0) = x \cdot 0 + x \cdot 0) \to \varphi_{25}$
\item $\varphi_{41} :\equiv (x \in \mathbb{R} \to x \cdot 0 + 0 = 0 + x \cdot 0) \to \varphi_{23} \to x \in \mathbb{R} \to x \cdot 0 + x \cdot 0 = 0 + x \cdot 0$
\end{abbrevlist}
\end{fullproof}

% 自動生成 (prover/build.js): NegMul  ノード数 235
\begin{fullproof}{[F9]}{$\forall x\,\forall y\,(x \in \mathbb{R} \to y \in \mathbb{R} \to -x \cdot y = -(x \cdot y))$}{235}\phantomsection\label{tree:NegMul}
\noindent$\Pi_{1}$:
\begin{ptree}
\AXC{$y = z \to x = y \to x = z$\lbl{A6}}
\AXC{$(y = z \to x = y \to x = z) \to x = y \to y = z \to x = z$\lbl{A1}}
\BIC{$x = y \to y = z \to x = z$}
\UIC{$\forall z\,(x = y \to y = z \to x = z)$}
\UIC{$\forall y\,\forall z\,(x = y \to y = z \to x = z)$}
\UIC{$\forall x\,\forall y\,\forall z\,(x = y \to y = z \to x = z)$}
\AXC{$\varphi_{7}$\lbl{A2}}
\BIC{$\forall y\,\forall z\,(v_{99} = y \to y = z \to v_{99} = z)$}
\AXC{$\varphi_{14}$\lbl{A2}}
\BIC{$\forall z\,(v_{99} = v_{100} \to v_{100} = z \to v_{99} = z)$}
\AXC{$\varphi_{20}$\lbl{A2}}
\BIC{$v_{99} = v_{100} \to v_{100} = v_{101} \to v_{99} = v_{101}$}
\UIC{$\forall v_{101}\,(v_{99} = v_{100} \to v_{100} = v_{101} \to v_{99} = v_{101})$}
\UIC{$\varphi_{5}$}
\UIC{$\forall v_{99}\,(\varphi_{5})$}
\end{ptree}
\noindent$\Pi_{2}$:
\begin{ptree}
\AXC{$\forall x \in \mathbb{R}\,(x \cdot 0 = 0)$\lbl{[F8]}}
\AXC{$\forall x \in \mathbb{R}\,(x \cdot 0 = 0) \to y \in \mathbb{R} \to y \cdot 0 = 0$\lbl{A2}}
\BIC{$y \in \mathbb{R} \to y \cdot 0 = 0$}
\AXC{$(y \in \mathbb{R} \to y \cdot 0 = 0) \to \varphi_{22}$\lbl{A1}}
\BIC{$\varphi_{22}$}
\end{ptree}
\noindent$\Pi_{3}$:
\begin{ptree}
\AXC{$y = z \to x = y \to x = z$\lbl{A6}}
\AXC{$(y = z \to x = y \to x = z) \to x = y \to y = z \to x = z$\lbl{A1}}
\BIC{$x = y \to y = z \to x = z$}
\UIC{$\forall z\,(x = y \to y = z \to x = z)$}
\UIC{$\forall y\,\forall z\,(x = y \to y = z \to x = z)$}
\UIC{$\forall x\,\forall y\,\forall z\,(x = y \to y = z \to x = z)$}
\AXC{$\varphi_{8}$\lbl{A2}}
\BIC{$\forall y\,\forall z\,(v_{96} = y \to y = z \to v_{96} = z)$}
\AXC{$\varphi_{11}$\lbl{A2}}
\BIC{$\forall z\,(v_{96} = v_{97} \to v_{97} = z \to v_{96} = z)$}
\AXC{$\varphi_{15}$\lbl{A2}}
\BIC{$v_{96} = v_{97} \to v_{97} = v_{98} \to v_{96} = v_{98}$}
\UIC{$\forall v_{98}\,(v_{96} = v_{97} \to v_{97} = v_{98} \to v_{96} = v_{98})$}
\UIC{$\varphi_{2}$}
\UIC{$\forall v_{96}\,(\varphi_{2})$}
\end{ptree}
\noindent$\Pi_{4}$:
\begin{ptree}
\AXC{$z = w \to y \cdot (x + -x) = y \cdot z \to y \cdot (x + -x) = y \cdot w$\lbl{A6}}
\UIC{$\varphi_{19}$}
\UIC{$\forall z\,(\varphi_{19})$}
\AXC{$\forall z\,(\varphi_{19}) \to \varphi_{31}$\lbl{A2}}
\BIC{$\varphi_{31}$}
\AXC{$\varphi_{31} \to \varphi_{26}$\lbl{A2}}
\BIC{$\varphi_{26}$}
\end{ptree}
\noindent$\Pi_{5}$:
\begin{ptree}
\AXC{$\forall x \in \mathbb{R}\,(x + -x = 0)$\lbl{(R5)}}
\AXC{$\forall x \in \mathbb{R}\,(x + -x = 0) \to x \in \mathbb{R} \to x + -x = 0$\lbl{A2}}
\BIC{$x \in \mathbb{R} \to x + -x = 0$}
\AXC{$(x \in \mathbb{R} \to x + -x = 0) \to \varphi_{32}$\lbl{A1}}
\BIC{$\varphi_{32}$}
\end{ptree}
\noindent$\Pi_{6}$:
\begin{ptree}
\AXC{$x \in \mathbb{R} \to y \in \mathbb{R} \to x \in \mathbb{R}$\lbl{A1}}
\AXC{$\stackrel{\Pi_{5}}{\vdots}$}\noLine
\UIC{$\varphi_{32}$}
\BIC{$x \in \mathbb{R} \to y \in \mathbb{R} \to x + -x = 0$}
\AXC{$(x \in \mathbb{R} \to y \in \mathbb{R} \to x + -x = 0) \to \varphi_{69}$\lbl{A1}}
\BIC{$\varphi_{69}$}
\end{ptree}
\noindent$\Pi_{7}$:
\begin{ptree}
\AXC{$\stackrel{\Pi_{4}}{\vdots}$}\noLine
\UIC{$\varphi_{26}$}
\AXC{$x = x$\lbl{A5}}
\UIC{$\forall x\,(x = x)$}
\AXC{$\forall x\,(x = x) \to y \cdot (x + -x) = y \cdot (x + -x)$\lbl{A2}}
\BIC{$y \cdot (x + -x) = y \cdot (x + -x)$}
\AXC{$\stackrel{\Pi_{6}}{\vdots}$}\noLine
\UIC{$\varphi_{69}$}
\BIC{$\varphi_{26} \to x \in \mathbb{R} \to y \in \mathbb{R} \to y \cdot (x + -x) = y \cdot 0$}
\BIC{$x \in \mathbb{R} \to y \in \mathbb{R} \to y \cdot (x + -x) = y \cdot 0$}
\end{ptree}
\noindent$\Pi_{8}$:
\begin{ptree}
\AXC{$y = z \to x = y \to x = z$\lbl{A6}}
\AXC{$(y = z \to x = y \to x = z) \to x = y \to y = z \to x = z$\lbl{A1}}
\BIC{$x = y \to y = z \to x = z$}
\UIC{$\forall z\,(x = y \to y = z \to x = z)$}
\UIC{$\forall y\,\forall z\,(x = y \to y = z \to x = z)$}
\UIC{$\forall x\,\forall y\,\forall z\,(x = y \to y = z \to x = z)$}
\AXC{$\varphi_{9}$\lbl{A2}}
\BIC{$\forall y\,\forall z\,(v_{93} = y \to y = z \to v_{93} = z)$}
\AXC{$\varphi_{12}$\lbl{A2}}
\BIC{$\forall z\,(v_{93} = v_{94} \to v_{94} = z \to v_{93} = z)$}
\AXC{$\varphi_{16}$\lbl{A2}}
\BIC{$v_{93} = v_{94} \to v_{94} = v_{95} \to v_{93} = v_{95}$}
\UIC{$\forall v_{95}\,(v_{93} = v_{94} \to v_{94} = v_{95} \to v_{93} = v_{95})$}
\UIC{$\varphi_{3}$}
\UIC{$\forall v_{93}\,(\varphi_{3})$}
\end{ptree}
\noindent$\Pi_{9}$:
\begin{ptree}
\AXC{$x = x$\lbl{A5}}
\AXC{$x = y \to x = x \to y = x$\lbl{A6}}
\AXC{$(x = y \to x = x \to y = x) \to x = x \to x = y \to y = x$\lbl{A1}}
\BIC{$x = x \to x = y \to y = x$}
\BIC{$x = y \to y = x$}
\UIC{$\forall y\,(x = y \to y = x)$}
\UIC{$\forall x\,\forall y\,(x = y \to y = x)$}
\end{ptree}
\noindent$\Pi_{10}$:
\begin{ptree}
\AXC{$\stackrel{\Pi_{9}}{\vdots}$}\noLine
\UIC{$\forall x\,\forall y\,(x = y \to y = x)$}
\AXC{$\forall x\,\forall y\,(x = y \to y = x) \to \forall y\,(v_{88} = y \to y = v_{88})$\lbl{A2}}
\BIC{$\forall y\,(v_{88} = y \to y = v_{88})$}
\AXC{$\varphi_{1}$\lbl{A2}}
\BIC{$v_{88} = v_{89} \to v_{89} = v_{88}$}
\UIC{$\forall v_{89}\,(v_{88} = v_{89} \to v_{89} = v_{88})$}
\UIC{$\forall v_{88}\,\forall v_{89}\,(v_{88} = v_{89} \to v_{89} = v_{88})$}
\end{ptree}
\noindent$\Pi_{11}$:
\begin{ptree}
\AXC{$\stackrel{\Pi_{10}}{\vdots}$}\noLine
\UIC{$\forall v_{88}\,\forall v_{89}\,(v_{88} = v_{89} \to v_{89} = v_{88})$}
\AXC{$\forall v_{88}\,\forall v_{89}\,(v_{88} = v_{89} \to v_{89} = v_{88}) \to \varphi_{6}$\lbl{A2}}
\BIC{$\varphi_{6}$}
\AXC{$\varphi_{6} \to \varphi_{35}$\lbl{A2}}
\BIC{$\varphi_{35}$}
\end{ptree}
\noindent$\Pi_{12}$:
\begin{ptree}
\AXC{$\forall x \in \mathbb{R}\,(-x \in \mathbb{R})$\lbl{(R1)}}
\AXC{$\forall x \in \mathbb{R}\,(-x \in \mathbb{R}) \to x \in \mathbb{R} \to -x \in \mathbb{R}$\lbl{A2}}
\BIC{$x \in \mathbb{R} \to -x \in \mathbb{R}$}
\AXC{$(x \in \mathbb{R} \to -x \in \mathbb{R}) \to \varphi_{38}$\lbl{A1}}
\BIC{$\varphi_{38}$}
\end{ptree}
\noindent$\Pi_{13}$:
\begin{ptree}
\AXC{$x \in \mathbb{R} \to y \in \mathbb{R} \to x \in \mathbb{R}$\lbl{A1}}
\AXC{$\stackrel{\Pi_{12}}{\vdots}$}\noLine
\UIC{$\varphi_{38}$}
\BIC{$x \in \mathbb{R} \to y \in \mathbb{R} \to -x \in \mathbb{R}$}
\end{ptree}
\noindent$\Pi_{14}$:
\begin{ptree}
\AXC{$\varphi_{29}$\lbl{(R10)}}
\AXC{$\varphi_{29} \to \varphi_{39}$\lbl{A2}}
\BIC{$\varphi_{39}$}
\AXC{$\varphi_{39} \to \varphi_{41}$\lbl{A2}}
\BIC{$\varphi_{41}$}
\AXC{$\varphi_{41} \to v_{85} \in \mathbb{R} \to v_{86} \in \mathbb{R} \to \varphi_{18}$\lbl{A2}}
\BIC{$v_{85} \in \mathbb{R} \to v_{86} \in \mathbb{R} \to \varphi_{18}$}
\UIC{$\forall v_{87}\,(v_{85} \in \mathbb{R} \to v_{86} \in \mathbb{R} \to \varphi_{18})$}
\UIC{$\forall v_{86}\,\forall v_{87}\,(v_{85} \in \mathbb{R} \to v_{86} \in \mathbb{R} \to \varphi_{18})$}
\UIC{$\forall v_{85}\,\forall v_{86}\,\forall v_{87}\,(v_{85} \in \mathbb{R} \to v_{86} \in \mathbb{R} \to \varphi_{18})$}
\AXC{$\varphi_{76}$\lbl{A2}}
\BIC{$\varphi_{46}$}
\AXC{$\varphi_{46} \to \varphi_{33}$\lbl{A2}}
\BIC{$\varphi_{33}$}
\end{ptree}
\noindent$\Pi_{15}$:
\begin{ptree}
\AXC{$x \in \mathbb{R} \to y \in \mathbb{R} \to x \in \mathbb{R}$\lbl{A1}}
\AXC{$x \in \mathbb{R} \to y \in \mathbb{R} \to y \in \mathbb{R}$\lbl{A1}}
\AXC{$\stackrel{\Pi_{14}}{\vdots}$}\noLine
\UIC{$\varphi_{33}$}
\AXC{$\varphi_{33} \to \varphi_{40}$\lbl{A2}}
\BIC{$\varphi_{40}$}
\AXC{$\varphi_{40} \to \varphi_{72}$\lbl{A1}}
\BIC{$\varphi_{72}$}
\BIC{$(x \in \mathbb{R} \to y \in \mathbb{R} \to x \in \mathbb{R}) \to \varphi_{62}$}
\BIC{$\varphi_{62}$}
\end{ptree}
\noindent$\Pi_{16}$:
\begin{ptree}
\AXC{$\stackrel{\Pi_{11}}{\vdots}$}\noLine
\UIC{$\varphi_{35}$}
\AXC{$\stackrel{\Pi_{13}}{\vdots}$}\noLine
\UIC{$x \in \mathbb{R} \to y \in \mathbb{R} \to -x \in \mathbb{R}$}
\AXC{$\stackrel{\Pi_{15}}{\vdots}$}\noLine
\UIC{$\varphi_{62}$}
\BIC{$x \in \mathbb{R} \to y \in \mathbb{R} \to y \cdot (x + -x) = y \cdot x + y \cdot -x$}
\AXC{$\varphi_{81}$\lbl{A1}}
\BIC{$\varphi_{35} \to x \in \mathbb{R} \to y \in \mathbb{R} \to y \cdot x + y \cdot -x = y \cdot (x + -x)$}
\BIC{$x \in \mathbb{R} \to y \in \mathbb{R} \to y \cdot x + y \cdot -x = y \cdot (x + -x)$}
\end{ptree}
\noindent$\Pi_{17}$:
\begin{ptree}
\AXC{$y = z \to x = y \to x = z$\lbl{A6}}
\AXC{$(y = z \to x = y \to x = z) \to x = y \to y = z \to x = z$\lbl{A1}}
\BIC{$x = y \to y = z \to x = z$}
\UIC{$\forall z\,(x = y \to y = z \to x = z)$}
\UIC{$\forall y\,\forall z\,(x = y \to y = z \to x = z)$}
\UIC{$\forall x\,\forall y\,\forall z\,(x = y \to y = z \to x = z)$}
\AXC{$\varphi_{10}$\lbl{A2}}
\BIC{$\forall y\,\forall z\,(v_{90} = y \to y = z \to v_{90} = z)$}
\AXC{$\varphi_{13}$\lbl{A2}}
\BIC{$\forall z\,(v_{90} = v_{91} \to v_{91} = z \to v_{90} = z)$}
\AXC{$\varphi_{17}$\lbl{A2}}
\BIC{$v_{90} = v_{91} \to v_{91} = v_{92} \to v_{90} = v_{92}$}
\UIC{$\forall v_{92}\,(v_{90} = v_{91} \to v_{91} = v_{92} \to v_{90} = v_{92})$}
\UIC{$\varphi_{4}$}
\UIC{$\forall v_{90}\,(\varphi_{4})$}
\end{ptree}
\noindent$\Pi_{18}$:
\begin{ptree}
\AXC{$\stackrel{\Pi_{17}}{\vdots}$}\noLine
\UIC{$\forall v_{90}\,(\varphi_{4})$}
\AXC{$\forall v_{90}\,(\varphi_{4}) \to \varphi_{25}$\lbl{A2}}
\BIC{$\varphi_{25}$}
\AXC{$\varphi_{25} \to \varphi_{59}$\lbl{A2}}
\BIC{$\varphi_{59}$}
\AXC{$\varphi_{59} \to x \cdot y + -x \cdot y = y \cdot x + -x \cdot y \to \varphi_{47}$\lbl{A2}}
\BIC{$x \cdot y + -x \cdot y = y \cdot x + -x \cdot y \to \varphi_{47}$}
\end{ptree}
\noindent$\Pi_{19}$:
\begin{ptree}
\AXC{$\stackrel{\Pi_{13}}{\vdots}$}\noLine
\UIC{$x \in \mathbb{R} \to y \in \mathbb{R} \to -x \in \mathbb{R}$}
\AXC{$\forall x\,\forall y\,(x \in \mathbb{R} \to y \in \mathbb{R} \to x \cdot y = y \cdot x)$\lbl{(R7)}}
\AXC{$\varphi_{61}$\lbl{A2}}
\BIC{$\forall y\,(-x \in \mathbb{R} \to y \in \mathbb{R} \to -x \cdot y = y \cdot -x)$}
\AXC{$\varphi_{65}$\lbl{A2}}
\BIC{$-x \in \mathbb{R} \to y \in \mathbb{R} \to -x \cdot y = y \cdot -x$}
\AXC{$\varphi_{74}$\lbl{A1}}
\BIC{$(x \in \mathbb{R} \to y \in \mathbb{R} \to -x \in \mathbb{R}) \to \varphi_{44}$}
\BIC{$\varphi_{44}$}
\end{ptree}
\noindent$\Pi_{20}$:
\begin{ptree}
\AXC{$x = x$\lbl{A5}}
\UIC{$\forall x\,(x = x)$}
\AXC{$\forall x\,(x = x) \to y \cdot x + -x \cdot y = y \cdot x + -x \cdot y$\lbl{A2}}
\BIC{$y \cdot x + -x \cdot y = y \cdot x + -x \cdot y$}
\AXC{$x \in \mathbb{R} \to y \in \mathbb{R} \to y \in \mathbb{R}$\lbl{A1}}
\AXC{$\stackrel{\Pi_{19}}{\vdots}$}\noLine
\UIC{$\varphi_{44}$}
\BIC{$x \in \mathbb{R} \to y \in \mathbb{R} \to -x \cdot y = y \cdot -x$}
\AXC{$\varphi_{86}$\lbl{A1}}
\BIC{$y \cdot x + -x \cdot y = y \cdot x + -x \cdot y \to \varphi_{73}$}
\BIC{$\varphi_{73}$}
\end{ptree}
\noindent$\Pi_{21}$:
\begin{ptree}
\AXC{$\forall x\,\forall y\,(x \in \mathbb{R} \to y \in \mathbb{R} \to x \cdot y = y \cdot x)$\lbl{(R7)}}
\AXC{$\varphi_{50}$\lbl{A2}}
\BIC{$\forall y\,(x \in \mathbb{R} \to y \in \mathbb{R} \to x \cdot y = y \cdot x)$}
\AXC{$\varphi_{42}$\lbl{A2}}
\BIC{$x \in \mathbb{R} \to y \in \mathbb{R} \to x \cdot y = y \cdot x$}
\AXC{$\varphi_{83}$\lbl{A1}}
\BIC{$x \cdot y + -x \cdot y = x \cdot y + -x \cdot y \to \varphi_{70}$}
\end{ptree}
\noindent$\Pi_{22}$:
\begin{ptree}
\AXC{$x = x$\lbl{A5}}
\UIC{$\forall x\,(x = x)$}
\AXC{$\forall x\,(x = x) \to x \cdot y + -x \cdot y = x \cdot y + -x \cdot y$\lbl{A2}}
\BIC{$x \cdot y + -x \cdot y = x \cdot y + -x \cdot y$}
\AXC{$\stackrel{\Pi_{21}}{\vdots}$}\noLine
\UIC{$x \cdot y + -x \cdot y = x \cdot y + -x \cdot y \to \varphi_{70}$}
\BIC{$\varphi_{70}$}
\end{ptree}
\noindent$\Pi_{23}$:
\begin{ptree}
\AXC{$z = w \to \varphi_{37}$\lbl{A6}}
\UIC{$\forall w\,(z = w \to \varphi_{37})$}
\UIC{$\forall z\,\forall w\,(z = w \to \varphi_{37})$}
\AXC{$\forall z\,\forall w\,(z = w \to \varphi_{37}) \to \varphi_{56}$\lbl{A2}}
\BIC{$\varphi_{56}$}
\AXC{$\varphi_{56} \to x \cdot y = y \cdot x \to \varphi_{49}$\lbl{A2}}
\BIC{$x \cdot y = y \cdot x \to \varphi_{49}$}
\AXC{$\stackrel{\Pi_{22}}{\vdots}$}\noLine
\UIC{$\varphi_{70}$}
\BIC{$x \in \mathbb{R} \to y \in \mathbb{R} \to x \cdot y + -x \cdot y = y \cdot x + -x \cdot y$}
\AXC{$\varphi_{88}$\lbl{A1}}
\BIC{$\varphi_{85}$}
\end{ptree}
\noindent$\Pi_{24}$:
\begin{ptree}
\AXC{$\varphi_{27}$\lbl{A6}}
\UIC{$\forall w\,(\varphi_{27})$}
\UIC{$\forall z\,\forall w\,(\varphi_{27})$}
\AXC{$\forall z\,\forall w\,(\varphi_{27}) \to \varphi_{55}$\lbl{A2}}
\BIC{$\varphi_{55}$}
\AXC{$\varphi_{55} \to -x \cdot y = y \cdot -x \to \varphi_{48}$\lbl{A2}}
\BIC{$-x \cdot y = y \cdot -x \to \varphi_{48}$}
\AXC{$\stackrel{\Pi_{20}}{\vdots}$}\noLine
\UIC{$\varphi_{73}$}
\BIC{$x \in \mathbb{R} \to y \in \mathbb{R} \to y \cdot x + -x \cdot y = y \cdot x + y \cdot -x$}
\AXC{$\stackrel{\Pi_{23}}{\vdots}$}\noLine
\UIC{$\varphi_{85}$}
\BIC{$\varphi_{79}$}
\end{ptree}
\noindent$\Pi_{25}$:
\begin{ptree}
\AXC{$\stackrel{\Pi_{16}}{\vdots}$}\noLine
\UIC{$x \in \mathbb{R} \to y \in \mathbb{R} \to y \cdot x + y \cdot -x = y \cdot (x + -x)$}
\AXC{$\stackrel{\Pi_{18}}{\vdots}$}\noLine
\UIC{$x \cdot y + -x \cdot y = y \cdot x + -x \cdot y \to \varphi_{47}$}
\AXC{$\stackrel{\Pi_{24}}{\vdots}$}\noLine
\UIC{$\varphi_{79}$}
\BIC{$x \in \mathbb{R} \to y \in \mathbb{R} \to x \cdot y + -x \cdot y = y \cdot x + y \cdot -x$}
\AXC{$\varphi_{87}$\lbl{A1}}
\BIC{$\varphi_{82}$}
\BIC{$\varphi_{75}$}
\end{ptree}
\noindent$\Pi_{26}$:
\begin{ptree}
\AXC{$\stackrel{\Pi_{8}}{\vdots}$}\noLine
\UIC{$\forall v_{93}\,(\varphi_{3})$}
\AXC{$\forall v_{93}\,(\varphi_{3}) \to \varphi_{24}$\lbl{A2}}
\BIC{$\varphi_{24}$}
\AXC{$\varphi_{24} \to \varphi_{58}$\lbl{A2}}
\BIC{$\varphi_{58}$}
\AXC{$\varphi_{58} \to x \cdot y + -x \cdot y = y \cdot x + y \cdot -x \to \varphi_{36}$\lbl{A2}}
\BIC{$x \cdot y + -x \cdot y = y \cdot x + y \cdot -x \to \varphi_{36}$}
\AXC{$\stackrel{\Pi_{25}}{\vdots}$}\noLine
\UIC{$\varphi_{75}$}
\BIC{$x \in \mathbb{R} \to y \in \mathbb{R} \to x \cdot y + -x \cdot y = y \cdot (x + -x)$}
\end{ptree}
\noindent$\Pi_{27}$:
\begin{ptree}
\AXC{$\stackrel{\Pi_{7}}{\vdots}$}\noLine
\UIC{$x \in \mathbb{R} \to y \in \mathbb{R} \to y \cdot (x + -x) = y \cdot 0$}
\AXC{$\stackrel{\Pi_{26}}{\vdots}$}\noLine
\UIC{$x \in \mathbb{R} \to y \in \mathbb{R} \to x \cdot y + -x \cdot y = y \cdot (x + -x)$}
\AXC{$\varphi_{84}$\lbl{A1}}
\BIC{$\varphi_{78}$}
\BIC{$\varphi_{54} \to x \in \mathbb{R} \to y \in \mathbb{R} \to x \cdot y + -x \cdot y = y \cdot 0$}
\end{ptree}
\noindent$\Pi_{28}$:
\begin{ptree}
\AXC{$\stackrel{\Pi_{3}}{\vdots}$}\noLine
\UIC{$\forall v_{96}\,(\varphi_{2})$}
\AXC{$\forall v_{96}\,(\varphi_{2}) \to \varphi_{23}$\lbl{A2}}
\BIC{$\varphi_{23}$}
\AXC{$\varphi_{23} \to \varphi_{52}$\lbl{A2}}
\BIC{$\varphi_{52}$}
\AXC{$\varphi_{52} \to \varphi_{54}$\lbl{A2}}
\BIC{$\varphi_{54}$}
\AXC{$\stackrel{\Pi_{27}}{\vdots}$}\noLine
\UIC{$\varphi_{54} \to x \in \mathbb{R} \to y \in \mathbb{R} \to x \cdot y + -x \cdot y = y \cdot 0$}
\BIC{$x \in \mathbb{R} \to y \in \mathbb{R} \to x \cdot y + -x \cdot y = y \cdot 0$}
\end{ptree}
\noindent$\Pi_{29}$:
\begin{ptree}
\AXC{$x \in \mathbb{R} \to y \in \mathbb{R} \to y \in \mathbb{R}$\lbl{A1}}
\AXC{$\stackrel{\Pi_{2}}{\vdots}$}\noLine
\UIC{$\varphi_{22}$}
\BIC{$x \in \mathbb{R} \to y \in \mathbb{R} \to y \cdot 0 = 0$}
\AXC{$\stackrel{\Pi_{28}}{\vdots}$}\noLine
\UIC{$x \in \mathbb{R} \to y \in \mathbb{R} \to x \cdot y + -x \cdot y = y \cdot 0$}
\AXC{$\varphi_{80}$\lbl{A1}}
\BIC{$(x \in \mathbb{R} \to y \in \mathbb{R} \to y \cdot 0 = 0) \to \varphi_{66}$}
\BIC{$\varphi_{66}$}
\end{ptree}
\noindent$\Pi_{30}$:
\begin{ptree}
\AXC{$\stackrel{\Pi_{1}}{\vdots}$}\noLine
\UIC{$\forall v_{99}\,(\varphi_{5})$}
\AXC{$\forall v_{99}\,(\varphi_{5}) \to \varphi_{34}$\lbl{A2}}
\BIC{$\varphi_{34}$}
\AXC{$\varphi_{34} \to \varphi_{30}$\lbl{A2}}
\BIC{$\varphi_{30}$}
\AXC{$\varphi_{68}$\lbl{A2}}
\BIC{$x \cdot y + -x \cdot y = y \cdot 0 \to y \cdot 0 = 0 \to x \cdot y + -x \cdot y = 0$}
\AXC{$\stackrel{\Pi_{29}}{\vdots}$}\noLine
\UIC{$\varphi_{66}$}
\BIC{$x \in \mathbb{R} \to y \in \mathbb{R} \to x \cdot y + -x \cdot y = 0$}
\end{ptree}
\noindent$\Pi_{31}$:
\begin{ptree}
\AXC{$\stackrel{\Pi_{13}}{\vdots}$}\noLine
\UIC{$x \in \mathbb{R} \to y \in \mathbb{R} \to -x \in \mathbb{R}$}
\AXC{$\forall x\,\forall y\,(x \in \mathbb{R} \to y \in \mathbb{R} \to x \cdot y \in \mathbb{R})$\lbl{(R1)}}
\AXC{$\varphi_{60}$\lbl{A2}}
\BIC{$\forall y\,(-x \in \mathbb{R} \to y \in \mathbb{R} \to -x \cdot y \in \mathbb{R})$}
\AXC{$\varphi_{63}$\lbl{A2}}
\BIC{$-x \in \mathbb{R} \to y \in \mathbb{R} \to -x \cdot y \in \mathbb{R}$}
\AXC{$\varphi_{71}$\lbl{A1}}
\BIC{$(x \in \mathbb{R} \to y \in \mathbb{R} \to -x \in \mathbb{R}) \to \varphi_{43}$}
\BIC{$\varphi_{43}$}
\end{ptree}
\noindent$\Pi_{32}$:
\begin{ptree}
\AXC{$\forall x\,\forall y\,(x \in \mathbb{R} \to y \in \mathbb{R} \to x + y = 0 \to y = -x)$\lbl{[F4]}}
\AXC{$\varphi_{64}$\lbl{A2}}
\BIC{$\forall y\,(v_{102} \in \mathbb{R} \to y \in \mathbb{R} \to v_{102} + y = 0 \to y = -v_{102})$}
\AXC{$\varphi_{67}$\lbl{A2}}
\BIC{$\varphi_{21}$}
\UIC{$\forall v_{103}\,(\varphi_{21})$}
\UIC{$\forall v_{102}\,\forall v_{103}\,(\varphi_{21})$}
\AXC{$\forall v_{102}\,\forall v_{103}\,(\varphi_{21}) \to \varphi_{28}$\lbl{A2}}
\BIC{$\varphi_{28}$}
\AXC{$\varphi_{28} \to \varphi_{51}$\lbl{A2}}
\BIC{$\varphi_{51}$}
\end{ptree}
\noindent$\Pi_{33}$:
\begin{ptree}
\AXC{$\forall x\,\forall y\,(x \in \mathbb{R} \to y \in \mathbb{R} \to x \cdot y \in \mathbb{R})$\lbl{(R1)}}
\AXC{$\varphi_{53}$\lbl{A2}}
\BIC{$\forall y\,(x \in \mathbb{R} \to y \in \mathbb{R} \to x \cdot y \in \mathbb{R})$}
\AXC{$\varphi_{45}$\lbl{A2}}
\BIC{$x \in \mathbb{R} \to y \in \mathbb{R} \to x \cdot y \in \mathbb{R}$}
\AXC{$\stackrel{\Pi_{32}}{\vdots}$}\noLine
\UIC{$\varphi_{51}$}
\AXC{$\varphi_{51} \to \varphi_{77}$\lbl{A1}}
\BIC{$\varphi_{77}$}
\BIC{$(x \in \mathbb{R} \to y \in \mathbb{R} \to -x \cdot y \in \mathbb{R}) \to \varphi_{57}$}
\end{ptree}
\noindent$\Pi_{34}$:
\begin{ptree}
\AXC{$x \in \mathbb{R} \to y \in \mathbb{R} \to y \in \mathbb{R}$\lbl{A1}}
\AXC{$\stackrel{\Pi_{31}}{\vdots}$}\noLine
\UIC{$\varphi_{43}$}
\BIC{$x \in \mathbb{R} \to y \in \mathbb{R} \to -x \cdot y \in \mathbb{R}$}
\AXC{$\stackrel{\Pi_{33}}{\vdots}$}\noLine
\UIC{$(x \in \mathbb{R} \to y \in \mathbb{R} \to -x \cdot y \in \mathbb{R}) \to \varphi_{57}$}
\BIC{$\varphi_{57}$}
\end{ptree}
\noindent 主木:
\begin{ptree}
\AXC{$\stackrel{\Pi_{30}}{\vdots}$}\noLine
\UIC{$x \in \mathbb{R} \to y \in \mathbb{R} \to x \cdot y + -x \cdot y = 0$}
\AXC{$\stackrel{\Pi_{34}}{\vdots}$}\noLine
\UIC{$\varphi_{57}$}
\BIC{$x \in \mathbb{R} \to y \in \mathbb{R} \to -x \cdot y = -(x \cdot y)$}
\UIC{$\forall y\,(x \in \mathbb{R} \to y \in \mathbb{R} \to -x \cdot y = -(x \cdot y))$}
\UIC{$\forall x\,\forall y\,(x \in \mathbb{R} \to y \in \mathbb{R} \to -x \cdot y = -(x \cdot y))$}
\end{ptree}
\begin{abbrevlist}
\item $\varphi_{1} :\equiv \forall y\,(v_{88} = y \to y = v_{88}) \to v_{88} = v_{89} \to v_{89} = v_{88}$
\item $\varphi_{2} :\equiv \forall v_{97}\,\forall v_{98}\,(v_{96} = v_{97} \to v_{97} = v_{98} \to v_{96} = v_{98})$
\item $\varphi_{3} :\equiv \forall v_{94}\,\forall v_{95}\,(v_{93} = v_{94} \to v_{94} = v_{95} \to v_{93} = v_{95})$
\item $\varphi_{4} :\equiv \forall v_{91}\,\forall v_{92}\,(v_{90} = v_{91} \to v_{91} = v_{92} \to v_{90} = v_{92})$
\item $\varphi_{5} :\equiv \forall v_{100}\,\forall v_{101}\,(v_{99} = v_{100} \to v_{100} = v_{101} \to v_{99} = v_{101})$
\item $\varphi_{6} :\equiv \forall v_{89}\,(y \cdot (x + -x) = v_{89} \to v_{89} = y \cdot (x + -x))$
\item $\varphi_{7} :\equiv \forall x\,\forall y\,\forall z\,(x = y \to y = z \to x = z) \to \forall y\,\forall z\,(v_{99} = y \to y = z \to v_{99} = z)$
\item $\varphi_{8} :\equiv \forall x\,\forall y\,\forall z\,(x = y \to y = z \to x = z) \to \forall y\,\forall z\,(v_{96} = y \to y = z \to v_{96} = z)$
\item $\varphi_{9} :\equiv \forall x\,\forall y\,\forall z\,(x = y \to y = z \to x = z) \to \forall y\,\forall z\,(v_{93} = y \to y = z \to v_{93} = z)$
\item $\varphi_{10} :\equiv \forall x\,\forall y\,\forall z\,(x = y \to y = z \to x = z) \to \forall y\,\forall z\,(v_{90} = y \to y = z \to v_{90} = z)$
\item $\varphi_{11} :\equiv \forall y\,\forall z\,(v_{96} = y \to y = z \to v_{96} = z) \to \forall z\,(v_{96} = v_{97} \to v_{97} = z \to v_{96} = z)$
\item $\varphi_{12} :\equiv \forall y\,\forall z\,(v_{93} = y \to y = z \to v_{93} = z) \to \forall z\,(v_{93} = v_{94} \to v_{94} = z \to v_{93} = z)$
\item $\varphi_{13} :\equiv \forall y\,\forall z\,(v_{90} = y \to y = z \to v_{90} = z) \to \forall z\,(v_{90} = v_{91} \to v_{91} = z \to v_{90} = z)$
\item $\varphi_{14} :\equiv \forall y\,\forall z\,(v_{99} = y \to y = z \to v_{99} = z) \to \forall z\,(v_{99} = v_{100} \to v_{100} = z \to v_{99} = z)$
\item $\varphi_{15} :\equiv \forall z\,(v_{96} = v_{97} \to v_{97} = z \to v_{96} = z) \to v_{96} = v_{97} \to v_{97} = v_{98} \to v_{96} = v_{98}$
\item $\varphi_{16} :\equiv \forall z\,(v_{93} = v_{94} \to v_{94} = z \to v_{93} = z) \to v_{93} = v_{94} \to v_{94} = v_{95} \to v_{93} = v_{95}$
\item $\varphi_{17} :\equiv \forall z\,(v_{90} = v_{91} \to v_{91} = z \to v_{90} = z) \to v_{90} = v_{91} \to v_{91} = v_{92} \to v_{90} = v_{92}$
\item $\varphi_{18} :\equiv v_{87} \in \mathbb{R} \to v_{85} \cdot (v_{86} + v_{87}) = v_{85} \cdot v_{86} + v_{85} \cdot v_{87}$
\item $\varphi_{19} :\equiv \forall w\,(z = w \to y \cdot (x + -x) = y \cdot z \to y \cdot (x + -x) = y \cdot w)$
\item $\varphi_{20} :\equiv \forall z\,(v_{99} = v_{100} \to v_{100} = z \to v_{99} = z) \to v_{99} = v_{100} \to v_{100} = v_{101} \to v_{99} = v_{101}$
\item $\varphi_{21} :\equiv v_{102} \in \mathbb{R} \to v_{103} \in \mathbb{R} \to v_{102} + v_{103} = 0 \to v_{103} = -v_{102}$
\item $\varphi_{22} :\equiv (x \in \mathbb{R} \to y \in \mathbb{R} \to y \in \mathbb{R}) \to x \in \mathbb{R} \to y \in \mathbb{R} \to y \cdot 0 = 0$
\item $\varphi_{23} :\equiv \forall v_{97}\,\forall v_{98}\,(x \cdot y + -x \cdot y = v_{97} \to v_{97} = v_{98} \to x \cdot y + -x \cdot y = v_{98})$
\item $\varphi_{24} :\equiv \forall v_{94}\,\forall v_{95}\,(x \cdot y + -x \cdot y = v_{94} \to v_{94} = v_{95} \to x \cdot y + -x \cdot y = v_{95})$
\item $\varphi_{25} :\equiv \forall v_{91}\,\forall v_{92}\,(x \cdot y + -x \cdot y = v_{91} \to v_{91} = v_{92} \to x \cdot y + -x \cdot y = v_{92})$
\item $\varphi_{26} :\equiv x + -x = 0 \to y \cdot (x + -x) = y \cdot (x + -x) \to y \cdot (x + -x) = y \cdot 0$
\item $\varphi_{27} :\equiv z = w \to y \cdot x + -x \cdot y = y \cdot x + z \to y \cdot x + -x \cdot y = y \cdot x + w$
\item $\varphi_{28} :\equiv \forall v_{103}\,(x \cdot y \in \mathbb{R} \to v_{103} \in \mathbb{R} \to x \cdot y + v_{103} = 0 \to v_{103} = -(x \cdot y))$
\item $\varphi_{29} :\equiv \forall x\,\forall y\,\forall z\,(x \in \mathbb{R} \to y \in \mathbb{R} \to z \in \mathbb{R} \to x \cdot (y + z) = x \cdot y + x \cdot z)$
\item $\varphi_{30} :\equiv \forall v_{101}\,(x \cdot y + -x \cdot y = y \cdot 0 \to y \cdot 0 = v_{101} \to x \cdot y + -x \cdot y = v_{101})$
\item $\varphi_{31} :\equiv \forall w\,(x + -x = w \to y \cdot (x + -x) = y \cdot (x + -x) \to y \cdot (x + -x) = y \cdot w)$
\item $\varphi_{32} :\equiv (x \in \mathbb{R} \to y \in \mathbb{R} \to x \in \mathbb{R}) \to x \in \mathbb{R} \to y \in \mathbb{R} \to x + -x = 0$
\item $\varphi_{33} :\equiv \forall v_{87}\,(y \in \mathbb{R} \to x \in \mathbb{R} \to v_{87} \in \mathbb{R} \to y \cdot (x + v_{87}) = y \cdot x + y \cdot v_{87})$
\item $\varphi_{34} :\equiv \forall v_{100}\,\forall v_{101}\,(x \cdot y + -x \cdot y = v_{100} \to v_{100} = v_{101} \to x \cdot y + -x \cdot y = v_{101})$
\item $\varphi_{35} :\equiv y \cdot (x + -x) = y \cdot x + y \cdot -x \to y \cdot x + y \cdot -x = y \cdot (x + -x)$
\item $\varphi_{36} :\equiv y \cdot x + y \cdot -x = y \cdot (x + -x) \to x \cdot y + -x \cdot y = y \cdot (x + -x)$
\item $\varphi_{37} :\equiv x \cdot y + -x \cdot y = z + -x \cdot y \to x \cdot y + -x \cdot y = w + -x \cdot y$
\item $\varphi_{38} :\equiv (x \in \mathbb{R} \to y \in \mathbb{R} \to x \in \mathbb{R}) \to x \in \mathbb{R} \to y \in \mathbb{R} \to -x \in \mathbb{R}$
\item $\varphi_{39} :\equiv \forall y\,\forall z\,(v_{85} \in \mathbb{R} \to y \in \mathbb{R} \to z \in \mathbb{R} \to v_{85} \cdot (y + z) = v_{85} \cdot y + v_{85} \cdot z)$
\item $\varphi_{40} :\equiv y \in \mathbb{R} \to x \in \mathbb{R} \to -x \in \mathbb{R} \to y \cdot (x + -x) = y \cdot x + y \cdot -x$
\item $\varphi_{41} :\equiv \forall z\,(v_{85} \in \mathbb{R} \to v_{86} \in \mathbb{R} \to z \in \mathbb{R} \to v_{85} \cdot (v_{86} + z) = v_{85} \cdot v_{86} + v_{85} \cdot z)$
\item $\varphi_{42} :\equiv \forall y\,(x \in \mathbb{R} \to y \in \mathbb{R} \to x \cdot y = y \cdot x) \to x \in \mathbb{R} \to y \in \mathbb{R} \to x \cdot y = y \cdot x$
\item $\varphi_{43} :\equiv (x \in \mathbb{R} \to y \in \mathbb{R} \to y \in \mathbb{R}) \to x \in \mathbb{R} \to y \in \mathbb{R} \to -x \cdot y \in \mathbb{R}$
\item $\varphi_{44} :\equiv (x \in \mathbb{R} \to y \in \mathbb{R} \to y \in \mathbb{R}) \to x \in \mathbb{R} \to y \in \mathbb{R} \to -x \cdot y = y \cdot -x$
\item $\varphi_{45} :\equiv \forall y\,(x \in \mathbb{R} \to y \in \mathbb{R} \to x \cdot y \in \mathbb{R}) \to x \in \mathbb{R} \to y \in \mathbb{R} \to x \cdot y \in \mathbb{R}$
\item $\varphi_{46} :\equiv \forall v_{86}\,\forall v_{87}\,(y \in \mathbb{R} \to v_{86} \in \mathbb{R} \to v_{87} \in \mathbb{R} \to y \cdot (v_{86} + v_{87}) = y \cdot v_{86} + y \cdot v_{87})$
\item $\varphi_{47} :\equiv y \cdot x + -x \cdot y = y \cdot x + y \cdot -x \to x \cdot y + -x \cdot y = y \cdot x + y \cdot -x$
\item $\varphi_{48} :\equiv y \cdot x + -x \cdot y = y \cdot x + -x \cdot y \to y \cdot x + -x \cdot y = y \cdot x + y \cdot -x$
\item $\varphi_{49} :\equiv x \cdot y + -x \cdot y = x \cdot y + -x \cdot y \to x \cdot y + -x \cdot y = y \cdot x + -x \cdot y$
\item $\varphi_{50} :\equiv \forall x\,\forall y\,(x \in \mathbb{R} \to y \in \mathbb{R} \to x \cdot y = y \cdot x) \to \forall y\,(x \in \mathbb{R} \to y \in \mathbb{R} \to x \cdot y = y \cdot x)$
\item $\varphi_{51} :\equiv x \cdot y \in \mathbb{R} \to -x \cdot y \in \mathbb{R} \to x \cdot y + -x \cdot y = 0 \to -x \cdot y = -(x \cdot y)$
\item $\varphi_{52} :\equiv \forall v_{98}\,(x \cdot y + -x \cdot y = y \cdot (x + -x) \to y \cdot (x + -x) = v_{98} \to x \cdot y + -x \cdot y = v_{98})$
\item $\varphi_{53} :\equiv \forall x\,\forall y\,(x \in \mathbb{R} \to y \in \mathbb{R} \to x \cdot y \in \mathbb{R}) \to \forall y\,(x \in \mathbb{R} \to y \in \mathbb{R} \to x \cdot y \in \mathbb{R})$
\item $\varphi_{54} :\equiv x \cdot y + -x \cdot y = y \cdot (x + -x) \to y \cdot (x + -x) = y \cdot 0 \to x \cdot y + -x \cdot y = y \cdot 0$
\item $\varphi_{55} :\equiv \forall w\,(-x \cdot y = w \to y \cdot x + -x \cdot y = y \cdot x + -x \cdot y \to y \cdot x + -x \cdot y = y \cdot x + w)$
\item $\varphi_{56} :\equiv \forall w\,(x \cdot y = w \to x \cdot y + -x \cdot y = x \cdot y + -x \cdot y \to x \cdot y + -x \cdot y = w + -x \cdot y)$
\item $\varphi_{57} :\equiv (x \in \mathbb{R} \to y \in \mathbb{R} \to x \cdot y + -x \cdot y = 0) \to x \in \mathbb{R} \to y \in \mathbb{R} \to -x \cdot y = -(x \cdot y)$
\item $\varphi_{58} :\equiv \forall v_{95}\,(x \cdot y + -x \cdot y = y \cdot x + y \cdot -x \to y \cdot x + y \cdot -x = v_{95} \to x \cdot y + -x \cdot y = v_{95})$
\item $\varphi_{59} :\equiv \forall v_{92}\,(x \cdot y + -x \cdot y = y \cdot x + -x \cdot y \to y \cdot x + -x \cdot y = v_{92} \to x \cdot y + -x \cdot y = v_{92})$
\item $\varphi_{60} :\equiv \forall x\,\forall y\,(x \in \mathbb{R} \to y \in \mathbb{R} \to x \cdot y \in \mathbb{R}) \to \forall y\,(-x \in \mathbb{R} \to y \in \mathbb{R} \to -x \cdot y \in \mathbb{R})$
\item $\varphi_{61} :\equiv \forall x\,\forall y\,(x \in \mathbb{R} \to y \in \mathbb{R} \to x \cdot y = y \cdot x) \to \forall y\,(-x \in \mathbb{R} \to y \in \mathbb{R} \to -x \cdot y = y \cdot -x)$
\item $\varphi_{62} :\equiv (x \in \mathbb{R} \to y \in \mathbb{R} \to -x \in \mathbb{R}) \to x \in \mathbb{R} \to y \in \mathbb{R} \to y \cdot (x + -x) = y \cdot x + y \cdot -x$
\item $\varphi_{63} :\equiv \forall y\,(-x \in \mathbb{R} \to y \in \mathbb{R} \to -x \cdot y \in \mathbb{R}) \to -x \in \mathbb{R} \to y \in \mathbb{R} \to -x \cdot y \in \mathbb{R}$
\item $\varphi_{64} :\equiv \forall x\,\forall y\,(x \in \mathbb{R} \to y \in \mathbb{R} \to x + y = 0 \to y = -x) \to \forall y\,(v_{102} \in \mathbb{R} \to y \in \mathbb{R} \to v_{102} + y = 0 \to y = -v_{102})$
\item $\varphi_{65} :\equiv \forall y\,(-x \in \mathbb{R} \to y \in \mathbb{R} \to -x \cdot y = y \cdot -x) \to -x \in \mathbb{R} \to y \in \mathbb{R} \to -x \cdot y = y \cdot -x$
\item $\varphi_{66} :\equiv (x \cdot y + -x \cdot y = y \cdot 0 \to y \cdot 0 = 0 \to x \cdot y + -x \cdot y = 0) \to x \in \mathbb{R} \to y \in \mathbb{R} \to x \cdot y + -x \cdot y = 0$
\item $\varphi_{67} :\equiv \forall y\,(v_{102} \in \mathbb{R} \to y \in \mathbb{R} \to v_{102} + y = 0 \to y = -v_{102}) \to \varphi_{21}$
\item $\varphi_{68} :\equiv \varphi_{30} \to x \cdot y + -x \cdot y = y \cdot 0 \to y \cdot 0 = 0 \to x \cdot y + -x \cdot y = 0$
\item $\varphi_{69} :\equiv y \cdot (x + -x) = y \cdot (x + -x) \to \varphi_{26} \to x \in \mathbb{R} \to y \in \mathbb{R} \to y \cdot (x + -x) = y \cdot 0$
\item $\varphi_{70} :\equiv (x \cdot y = y \cdot x \to \varphi_{49}) \to x \in \mathbb{R} \to y \in \mathbb{R} \to x \cdot y + -x \cdot y = y \cdot x + -x \cdot y$
\item $\varphi_{71} :\equiv (-x \in \mathbb{R} \to y \in \mathbb{R} \to -x \cdot y \in \mathbb{R}) \to (x \in \mathbb{R} \to y \in \mathbb{R} \to -x \in \mathbb{R}) \to \varphi_{43}$
\item $\varphi_{72} :\equiv (x \in \mathbb{R} \to y \in \mathbb{R} \to y \in \mathbb{R}) \to (x \in \mathbb{R} \to y \in \mathbb{R} \to x \in \mathbb{R}) \to \varphi_{62}$
\item $\varphi_{73} :\equiv (-x \cdot y = y \cdot -x \to \varphi_{48}) \to x \in \mathbb{R} \to y \in \mathbb{R} \to y \cdot x + -x \cdot y = y \cdot x + y \cdot -x$
\item $\varphi_{74} :\equiv (-x \in \mathbb{R} \to y \in \mathbb{R} \to -x \cdot y = y \cdot -x) \to (x \in \mathbb{R} \to y \in \mathbb{R} \to -x \in \mathbb{R}) \to \varphi_{44}$
\item $\varphi_{75} :\equiv (x \cdot y + -x \cdot y = y \cdot x + y \cdot -x \to \varphi_{36}) \to x \in \mathbb{R} \to y \in \mathbb{R} \to x \cdot y + -x \cdot y = y \cdot (x + -x)$
\item $\varphi_{76} :\equiv \forall v_{85}\,\forall v_{86}\,\forall v_{87}\,(v_{85} \in \mathbb{R} \to v_{86} \in \mathbb{R} \to \varphi_{18}) \to \varphi_{46}$
\item $\varphi_{77} :\equiv (x \in \mathbb{R} \to y \in \mathbb{R} \to x \cdot y \in \mathbb{R}) \to (x \in \mathbb{R} \to y \in \mathbb{R} \to -x \cdot y \in \mathbb{R}) \to \varphi_{57}$
\item $\varphi_{78} :\equiv (x \in \mathbb{R} \to y \in \mathbb{R} \to y \cdot (x + -x) = y \cdot 0) \to \varphi_{54} \to x \in \mathbb{R} \to y \in \mathbb{R} \to x \cdot y + -x \cdot y = y \cdot 0$
\item $\varphi_{79} :\equiv (x \cdot y + -x \cdot y = y \cdot x + -x \cdot y \to \varphi_{47}) \to x \in \mathbb{R} \to y \in \mathbb{R} \to x \cdot y + -x \cdot y = y \cdot x + y \cdot -x$
\item $\varphi_{80} :\equiv (x \in \mathbb{R} \to y \in \mathbb{R} \to x \cdot y + -x \cdot y = y \cdot 0) \to (x \in \mathbb{R} \to y \in \mathbb{R} \to y \cdot 0 = 0) \to \varphi_{66}$
\item $\varphi_{81} :\equiv (x \in \mathbb{R} \to y \in \mathbb{R} \to y \cdot (x + -x) = y \cdot x + y \cdot -x) \to \varphi_{35} \to x \in \mathbb{R} \to y \in \mathbb{R} \to y \cdot x + y \cdot -x = y \cdot (x + -x)$
\item $\varphi_{82} :\equiv (x \in \mathbb{R} \to y \in \mathbb{R} \to y \cdot x + y \cdot -x = y \cdot (x + -x)) \to \varphi_{75}$
\item $\varphi_{83} :\equiv (x \in \mathbb{R} \to y \in \mathbb{R} \to x \cdot y = y \cdot x) \to x \cdot y + -x \cdot y = x \cdot y + -x \cdot y \to \varphi_{70}$
\item $\varphi_{84} :\equiv (x \in \mathbb{R} \to y \in \mathbb{R} \to x \cdot y + -x \cdot y = y \cdot (x + -x)) \to \varphi_{78}$
\item $\varphi_{85} :\equiv (x \in \mathbb{R} \to y \in \mathbb{R} \to y \cdot x + -x \cdot y = y \cdot x + y \cdot -x) \to \varphi_{79}$
\item $\varphi_{86} :\equiv (x \in \mathbb{R} \to y \in \mathbb{R} \to -x \cdot y = y \cdot -x) \to y \cdot x + -x \cdot y = y \cdot x + -x \cdot y \to \varphi_{73}$
\item $\varphi_{87} :\equiv (x \in \mathbb{R} \to y \in \mathbb{R} \to x \cdot y + -x \cdot y = y \cdot x + y \cdot -x) \to \varphi_{82}$
\item $\varphi_{88} :\equiv (x \in \mathbb{R} \to y \in \mathbb{R} \to x \cdot y + -x \cdot y = y \cdot x + -x \cdot y) \to \varphi_{85}$
\end{abbrevlist}
\end{fullproof}

% 自動生成 (prover/build.js): MulNeg  ノード数 112
\begin{fullproof}{[F10]}{$\forall x\,\forall y\,(x \in \mathbb{R} \to y \in \mathbb{R} \to x \cdot -y = -(x \cdot y))$}{112}\phantomsection\label{tree:MulNeg}
\noindent$\Pi_{1}$:
\begin{ptree}
\AXC{$y = z \to x = y \to x = z$\lbl{A6}}
\AXC{$(y = z \to x = y \to x = z) \to x = y \to y = z \to x = z$\lbl{A1}}
\BIC{$x = y \to y = z \to x = z$}
\UIC{$\forall z\,(x = y \to y = z \to x = z)$}
\UIC{$\forall y\,\forall z\,(x = y \to y = z \to x = z)$}
\UIC{$\forall x\,\forall y\,\forall z\,(x = y \to y = z \to x = z)$}
\AXC{$\varphi_{3}$\lbl{A2}}
\BIC{$\forall y\,\forall z\,(v_{111} = y \to y = z \to v_{111} = z)$}
\AXC{$\varphi_{8}$\lbl{A2}}
\BIC{$\forall z\,(v_{111} = v_{112} \to v_{112} = z \to v_{111} = z)$}
\AXC{$\varphi_{14}$\lbl{A2}}
\BIC{$v_{111} = v_{112} \to v_{112} = v_{113} \to v_{111} = v_{113}$}
\UIC{$\forall v_{113}\,(v_{111} = v_{112} \to v_{112} = v_{113} \to v_{111} = v_{113})$}
\UIC{$\varphi_{1}$}
\UIC{$\forall v_{111}\,(\varphi_{1})$}
\end{ptree}
\noindent$\Pi_{2}$:
\begin{ptree}
\AXC{$z = w \to -(y \cdot x) = -z \to -(y \cdot x) = -w$\lbl{A6}}
\UIC{$\forall w\,(z = w \to -(y \cdot x) = -z \to -(y \cdot x) = -w)$}
\UIC{$\forall z\,\forall w\,(z = w \to -(y \cdot x) = -z \to -(y \cdot x) = -w)$}
\AXC{$\varphi_{30}$\lbl{A2}}
\BIC{$\varphi_{5}$}
\AXC{$\varphi_{5} \to \varphi_{13}$\lbl{A2}}
\BIC{$\varphi_{13}$}
\end{ptree}
\noindent$\Pi_{3}$:
\begin{ptree}
\AXC{$\forall x\,\forall y\,(x \in \mathbb{R} \to y \in \mathbb{R} \to x \cdot y = y \cdot x)$\lbl{(R7)}}
\AXC{$\varphi_{27}$\lbl{A2}}
\BIC{$\forall y\,(v_{106} \in \mathbb{R} \to y \in \mathbb{R} \to v_{106} \cdot y = y \cdot v_{106})$}
\AXC{$\varphi_{29}$\lbl{A2}}
\BIC{$v_{106} \in \mathbb{R} \to v_{107} \in \mathbb{R} \to v_{106} \cdot v_{107} = v_{107} \cdot v_{106}$}
\UIC{$\varphi_{11}$}
\UIC{$\forall v_{106}\,(\varphi_{11})$}
\AXC{$\varphi_{31}$\lbl{A2}}
\BIC{$\forall v_{107}\,(y \in \mathbb{R} \to v_{107} \in \mathbb{R} \to y \cdot v_{107} = v_{107} \cdot y)$}
\AXC{$\varphi_{26}$\lbl{A2}}
\BIC{$y \in \mathbb{R} \to x \in \mathbb{R} \to y \cdot x = x \cdot y$}
\AXC{$\varphi_{33}$\lbl{A1}}
\BIC{$(x \in \mathbb{R} \to y \in \mathbb{R} \to y \in \mathbb{R}) \to \varphi_{21}$}
\end{ptree}
\noindent$\Pi_{4}$:
\begin{ptree}
\AXC{$x \in \mathbb{R} \to y \in \mathbb{R} \to x \in \mathbb{R}$\lbl{A1}}
\AXC{$x \in \mathbb{R} \to y \in \mathbb{R} \to y \in \mathbb{R}$\lbl{A1}}
\AXC{$\stackrel{\Pi_{3}}{\vdots}$}\noLine
\UIC{$(x \in \mathbb{R} \to y \in \mathbb{R} \to y \in \mathbb{R}) \to \varphi_{21}$}
\BIC{$\varphi_{21}$}
\BIC{$x \in \mathbb{R} \to y \in \mathbb{R} \to y \cdot x = x \cdot y$}
\end{ptree}
\noindent$\Pi_{5}$:
\begin{ptree}
\AXC{$\stackrel{\Pi_{4}}{\vdots}$}\noLine
\UIC{$x \in \mathbb{R} \to y \in \mathbb{R} \to y \cdot x = x \cdot y$}
\AXC{$(x \in \mathbb{R} \to y \in \mathbb{R} \to y \cdot x = x \cdot y) \to \varphi_{32}$\lbl{A1}}
\BIC{$\varphi_{32}$}
\end{ptree}
\noindent$\Pi_{6}$:
\begin{ptree}
\AXC{$y = z \to x = y \to x = z$\lbl{A6}}
\AXC{$(y = z \to x = y \to x = z) \to x = y \to y = z \to x = z$\lbl{A1}}
\BIC{$x = y \to y = z \to x = z$}
\UIC{$\forall z\,(x = y \to y = z \to x = z)$}
\UIC{$\forall y\,\forall z\,(x = y \to y = z \to x = z)$}
\UIC{$\forall x\,\forall y\,\forall z\,(x = y \to y = z \to x = z)$}
\AXC{$\varphi_{4}$\lbl{A2}}
\BIC{$\forall y\,\forall z\,(v_{108} = y \to y = z \to v_{108} = z)$}
\AXC{$\varphi_{9}$\lbl{A2}}
\BIC{$\forall z\,(v_{108} = v_{109} \to v_{109} = z \to v_{108} = z)$}
\AXC{$\varphi_{15}$\lbl{A2}}
\BIC{$v_{108} = v_{109} \to v_{109} = v_{110} \to v_{108} = v_{110}$}
\UIC{$\forall v_{110}\,(v_{108} = v_{109} \to v_{109} = v_{110} \to v_{108} = v_{110})$}
\UIC{$\varphi_{2}$}
\UIC{$\forall v_{108}\,(\varphi_{2})$}
\end{ptree}
\noindent$\Pi_{7}$:
\begin{ptree}
\AXC{$\forall x\,\forall y\,(x \in \mathbb{R} \to y \in \mathbb{R} \to -x \cdot y = -(x \cdot y))$\lbl{[F9]}}
\AXC{$\forall x\,\forall y\,(x \in \mathbb{R} \to y \in \mathbb{R} \to -x \cdot y = -(x \cdot y)) \to \varphi_{10}$\lbl{A2}}
\BIC{$\varphi_{10}$}
\AXC{$\varphi_{10} \to \varphi_{18}$\lbl{A2}}
\BIC{$\varphi_{18}$}
\UIC{$\forall v_{105}\,(\varphi_{18})$}
\UIC{$\forall v_{104}\,\forall v_{105}\,(\varphi_{18})$}
\end{ptree}
\noindent$\Pi_{8}$:
\begin{ptree}
\AXC{$\stackrel{\Pi_{7}}{\vdots}$}\noLine
\UIC{$\forall v_{104}\,\forall v_{105}\,(\varphi_{18})$}
\AXC{$\forall v_{104}\,\forall v_{105}\,(\varphi_{18}) \to \varphi_{12}$\lbl{A2}}
\BIC{$\varphi_{12}$}
\AXC{$\varphi_{12} \to y \in \mathbb{R} \to x \in \mathbb{R} \to -y \cdot x = -(y \cdot x)$\lbl{A2}}
\BIC{$y \in \mathbb{R} \to x \in \mathbb{R} \to -y \cdot x = -(y \cdot x)$}
\AXC{$\varphi_{34}$\lbl{A1}}
\BIC{$(x \in \mathbb{R} \to y \in \mathbb{R} \to y \in \mathbb{R}) \to \varphi_{23}$}
\end{ptree}
\noindent$\Pi_{9}$:
\begin{ptree}
\AXC{$\forall x \in \mathbb{R}\,(-x \in \mathbb{R})$\lbl{(R1)}}
\AXC{$\forall x \in \mathbb{R}\,(-x \in \mathbb{R}) \to y \in \mathbb{R} \to -y \in \mathbb{R}$\lbl{A2}}
\BIC{$y \in \mathbb{R} \to -y \in \mathbb{R}$}
\AXC{$(y \in \mathbb{R} \to -y \in \mathbb{R}) \to \varphi_{22}$\lbl{A1}}
\BIC{$\varphi_{22}$}
\end{ptree}
\noindent$\Pi_{10}$:
\begin{ptree}
\AXC{$x \in \mathbb{R} \to y \in \mathbb{R} \to x \in \mathbb{R}$\lbl{A1}}
\AXC{$\forall x\,\forall y\,(x \in \mathbb{R} \to y \in \mathbb{R} \to x \cdot y = y \cdot x)$\lbl{(R7)}}
\AXC{$\varphi_{24}$\lbl{A2}}
\BIC{$\forall y\,(x \in \mathbb{R} \to y \in \mathbb{R} \to x \cdot y = y \cdot x)$}
\AXC{$\varphi_{28}$\lbl{A2}}
\BIC{$x \in \mathbb{R} \to -y \in \mathbb{R} \to x \cdot -y = -y \cdot x$}
\AXC{$\varphi_{37}$\lbl{A1}}
\BIC{$(x \in \mathbb{R} \to y \in \mathbb{R} \to x \in \mathbb{R}) \to \varphi_{25}$}
\BIC{$\varphi_{25}$}
\end{ptree}
\noindent$\Pi_{11}$:
\begin{ptree}
\AXC{$x \in \mathbb{R} \to y \in \mathbb{R} \to y \in \mathbb{R}$\lbl{A1}}
\AXC{$\stackrel{\Pi_{9}}{\vdots}$}\noLine
\UIC{$\varphi_{22}$}
\BIC{$x \in \mathbb{R} \to y \in \mathbb{R} \to -y \in \mathbb{R}$}
\AXC{$\stackrel{\Pi_{10}}{\vdots}$}\noLine
\UIC{$\varphi_{25}$}
\BIC{$x \in \mathbb{R} \to y \in \mathbb{R} \to x \cdot -y = -y \cdot x$}
\AXC{$(x \in \mathbb{R} \to y \in \mathbb{R} \to x \cdot -y = -y \cdot x) \to \varphi_{35}$\lbl{A1}}
\BIC{$\varphi_{35}$}
\end{ptree}
\noindent$\Pi_{12}$:
\begin{ptree}
\AXC{$x \in \mathbb{R} \to y \in \mathbb{R} \to x \in \mathbb{R}$\lbl{A1}}
\AXC{$x \in \mathbb{R} \to y \in \mathbb{R} \to y \in \mathbb{R}$\lbl{A1}}
\AXC{$\stackrel{\Pi_{8}}{\vdots}$}\noLine
\UIC{$(x \in \mathbb{R} \to y \in \mathbb{R} \to y \in \mathbb{R}) \to \varphi_{23}$}
\BIC{$\varphi_{23}$}
\BIC{$x \in \mathbb{R} \to y \in \mathbb{R} \to -y \cdot x = -(y \cdot x)$}
\AXC{$\stackrel{\Pi_{11}}{\vdots}$}\noLine
\UIC{$\varphi_{35}$}
\BIC{$\varphi_{20} \to x \in \mathbb{R} \to y \in \mathbb{R} \to x \cdot -y = -(y \cdot x)$}
\end{ptree}
\noindent$\Pi_{13}$:
\begin{ptree}
\AXC{$\stackrel{\Pi_{6}}{\vdots}$}\noLine
\UIC{$\forall v_{108}\,(\varphi_{2})$}
\AXC{$\forall v_{108}\,(\varphi_{2}) \to \varphi_{7}$\lbl{A2}}
\BIC{$\varphi_{7}$}
\AXC{$\varphi_{7} \to \varphi_{17}$\lbl{A2}}
\BIC{$\varphi_{17}$}
\AXC{$\varphi_{17} \to \varphi_{20}$\lbl{A2}}
\BIC{$\varphi_{20}$}
\AXC{$\stackrel{\Pi_{12}}{\vdots}$}\noLine
\UIC{$\varphi_{20} \to x \in \mathbb{R} \to y \in \mathbb{R} \to x \cdot -y = -(y \cdot x)$}
\BIC{$x \in \mathbb{R} \to y \in \mathbb{R} \to x \cdot -y = -(y \cdot x)$}
\end{ptree}
\noindent$\Pi_{14}$:
\begin{ptree}
\AXC{$\stackrel{\Pi_{13}}{\vdots}$}\noLine
\UIC{$x \in \mathbb{R} \to y \in \mathbb{R} \to x \cdot -y = -(y \cdot x)$}
\AXC{$(x \in \mathbb{R} \to y \in \mathbb{R} \to x \cdot -y = -(y \cdot x)) \to \varphi_{36}$\lbl{A1}}
\BIC{$\varphi_{36}$}
\end{ptree}
\noindent$\Pi_{15}$:
\begin{ptree}
\AXC{$\stackrel{\Pi_{2}}{\vdots}$}\noLine
\UIC{$\varphi_{13}$}
\AXC{$x = x$\lbl{A5}}
\UIC{$\forall x\,(x = x)$}
\AXC{$\forall x\,(x = x) \to -(y \cdot x) = -(y \cdot x)$\lbl{A2}}
\BIC{$-(y \cdot x) = -(y \cdot x)$}
\AXC{$\stackrel{\Pi_{5}}{\vdots}$}\noLine
\UIC{$\varphi_{32}$}
\BIC{$\varphi_{13} \to x \in \mathbb{R} \to y \in \mathbb{R} \to -(y \cdot x) = -(x \cdot y)$}
\BIC{$x \in \mathbb{R} \to y \in \mathbb{R} \to -(y \cdot x) = -(x \cdot y)$}
\AXC{$\stackrel{\Pi_{14}}{\vdots}$}\noLine
\UIC{$\varphi_{36}$}
\BIC{$\varphi_{19} \to x \in \mathbb{R} \to y \in \mathbb{R} \to x \cdot -y = -(x \cdot y)$}
\end{ptree}
\noindent 主木:
\begin{ptree}
\AXC{$\stackrel{\Pi_{1}}{\vdots}$}\noLine
\UIC{$\forall v_{111}\,(\varphi_{1})$}
\AXC{$\forall v_{111}\,(\varphi_{1}) \to \varphi_{6}$\lbl{A2}}
\BIC{$\varphi_{6}$}
\AXC{$\varphi_{6} \to \varphi_{16}$\lbl{A2}}
\BIC{$\varphi_{16}$}
\AXC{$\varphi_{16} \to \varphi_{19}$\lbl{A2}}
\BIC{$\varphi_{19}$}
\AXC{$\stackrel{\Pi_{15}}{\vdots}$}\noLine
\UIC{$\varphi_{19} \to x \in \mathbb{R} \to y \in \mathbb{R} \to x \cdot -y = -(x \cdot y)$}
\BIC{$x \in \mathbb{R} \to y \in \mathbb{R} \to x \cdot -y = -(x \cdot y)$}
\UIC{$\forall y\,(x \in \mathbb{R} \to y \in \mathbb{R} \to x \cdot -y = -(x \cdot y))$}
\UIC{$\forall x\,\forall y\,(x \in \mathbb{R} \to y \in \mathbb{R} \to x \cdot -y = -(x \cdot y))$}
\end{ptree}
\begin{abbrevlist}
\item $\varphi_{1} :\equiv \forall v_{112}\,\forall v_{113}\,(v_{111} = v_{112} \to v_{112} = v_{113} \to v_{111} = v_{113})$
\item $\varphi_{2} :\equiv \forall v_{109}\,\forall v_{110}\,(v_{108} = v_{109} \to v_{109} = v_{110} \to v_{108} = v_{110})$
\item $\varphi_{3} :\equiv \forall x\,\forall y\,\forall z\,(x = y \to y = z \to x = z) \to \forall y\,\forall z\,(v_{111} = y \to y = z \to v_{111} = z)$
\item $\varphi_{4} :\equiv \forall x\,\forall y\,\forall z\,(x = y \to y = z \to x = z) \to \forall y\,\forall z\,(v_{108} = y \to y = z \to v_{108} = z)$
\item $\varphi_{5} :\equiv \forall w\,(y \cdot x = w \to -(y \cdot x) = -(y \cdot x) \to -(y \cdot x) = -w)$
\item $\varphi_{6} :\equiv \forall v_{112}\,\forall v_{113}\,(x \cdot -y = v_{112} \to v_{112} = v_{113} \to x \cdot -y = v_{113})$
\item $\varphi_{7} :\equiv \forall v_{109}\,\forall v_{110}\,(x \cdot -y = v_{109} \to v_{109} = v_{110} \to x \cdot -y = v_{110})$
\item $\varphi_{8} :\equiv \forall y\,\forall z\,(v_{111} = y \to y = z \to v_{111} = z) \to \forall z\,(v_{111} = v_{112} \to v_{112} = z \to v_{111} = z)$
\item $\varphi_{9} :\equiv \forall y\,\forall z\,(v_{108} = y \to y = z \to v_{108} = z) \to \forall z\,(v_{108} = v_{109} \to v_{109} = z \to v_{108} = z)$
\item $\varphi_{10} :\equiv \forall y\,(v_{104} \in \mathbb{R} \to y \in \mathbb{R} \to -v_{104} \cdot y = -(v_{104} \cdot y))$
\item $\varphi_{11} :\equiv \forall v_{107}\,(v_{106} \in \mathbb{R} \to v_{107} \in \mathbb{R} \to v_{106} \cdot v_{107} = v_{107} \cdot v_{106})$
\item $\varphi_{12} :\equiv \forall v_{105}\,(y \in \mathbb{R} \to v_{105} \in \mathbb{R} \to -y \cdot v_{105} = -(y \cdot v_{105}))$
\item $\varphi_{13} :\equiv y \cdot x = x \cdot y \to -(y \cdot x) = -(y \cdot x) \to -(y \cdot x) = -(x \cdot y)$
\item $\varphi_{14} :\equiv \forall z\,(v_{111} = v_{112} \to v_{112} = z \to v_{111} = z) \to v_{111} = v_{112} \to v_{112} = v_{113} \to v_{111} = v_{113}$
\item $\varphi_{15} :\equiv \forall z\,(v_{108} = v_{109} \to v_{109} = z \to v_{108} = z) \to v_{108} = v_{109} \to v_{109} = v_{110} \to v_{108} = v_{110}$
\item $\varphi_{16} :\equiv \forall v_{113}\,(x \cdot -y = -(y \cdot x) \to -(y \cdot x) = v_{113} \to x \cdot -y = v_{113})$
\item $\varphi_{17} :\equiv \forall v_{110}\,(x \cdot -y = -y \cdot x \to -y \cdot x = v_{110} \to x \cdot -y = v_{110})$
\item $\varphi_{18} :\equiv v_{104} \in \mathbb{R} \to v_{105} \in \mathbb{R} \to -v_{104} \cdot v_{105} = -(v_{104} \cdot v_{105})$
\item $\varphi_{19} :\equiv x \cdot -y = -(y \cdot x) \to -(y \cdot x) = -(x \cdot y) \to x \cdot -y = -(x \cdot y)$
\item $\varphi_{20} :\equiv x \cdot -y = -y \cdot x \to -y \cdot x = -(y \cdot x) \to x \cdot -y = -(y \cdot x)$
\item $\varphi_{21} :\equiv (x \in \mathbb{R} \to y \in \mathbb{R} \to x \in \mathbb{R}) \to x \in \mathbb{R} \to y \in \mathbb{R} \to y \cdot x = x \cdot y$
\item $\varphi_{22} :\equiv (x \in \mathbb{R} \to y \in \mathbb{R} \to y \in \mathbb{R}) \to x \in \mathbb{R} \to y \in \mathbb{R} \to -y \in \mathbb{R}$
\item $\varphi_{23} :\equiv (x \in \mathbb{R} \to y \in \mathbb{R} \to x \in \mathbb{R}) \to x \in \mathbb{R} \to y \in \mathbb{R} \to -y \cdot x = -(y \cdot x)$
\item $\varphi_{24} :\equiv \forall x\,\forall y\,(x \in \mathbb{R} \to y \in \mathbb{R} \to x \cdot y = y \cdot x) \to \forall y\,(x \in \mathbb{R} \to y \in \mathbb{R} \to x \cdot y = y \cdot x)$
\item $\varphi_{25} :\equiv (x \in \mathbb{R} \to y \in \mathbb{R} \to -y \in \mathbb{R}) \to x \in \mathbb{R} \to y \in \mathbb{R} \to x \cdot -y = -y \cdot x$
\item $\varphi_{26} :\equiv \forall v_{107}\,(y \in \mathbb{R} \to v_{107} \in \mathbb{R} \to y \cdot v_{107} = v_{107} \cdot y) \to y \in \mathbb{R} \to x \in \mathbb{R} \to y \cdot x = x \cdot y$
\item $\varphi_{27} :\equiv \forall x\,\forall y\,(x \in \mathbb{R} \to y \in \mathbb{R} \to x \cdot y = y \cdot x) \to \forall y\,(v_{106} \in \mathbb{R} \to y \in \mathbb{R} \to v_{106} \cdot y = y \cdot v_{106})$
\item $\varphi_{28} :\equiv \forall y\,(x \in \mathbb{R} \to y \in \mathbb{R} \to x \cdot y = y \cdot x) \to x \in \mathbb{R} \to -y \in \mathbb{R} \to x \cdot -y = -y \cdot x$
\item $\varphi_{29} :\equiv \forall y\,(v_{106} \in \mathbb{R} \to y \in \mathbb{R} \to v_{106} \cdot y = y \cdot v_{106}) \to v_{106} \in \mathbb{R} \to v_{107} \in \mathbb{R} \to v_{106} \cdot v_{107} = v_{107} \cdot v_{106}$
\item $\varphi_{30} :\equiv \forall z\,\forall w\,(z = w \to -(y \cdot x) = -z \to -(y \cdot x) = -w) \to \varphi_{5}$
\item $\varphi_{31} :\equiv \forall v_{106}\,(\varphi_{11}) \to \forall v_{107}\,(y \in \mathbb{R} \to v_{107} \in \mathbb{R} \to y \cdot v_{107} = v_{107} \cdot y)$
\item $\varphi_{32} :\equiv -(y \cdot x) = -(y \cdot x) \to \varphi_{13} \to x \in \mathbb{R} \to y \in \mathbb{R} \to -(y \cdot x) = -(x \cdot y)$
\item $\varphi_{33} :\equiv (y \in \mathbb{R} \to x \in \mathbb{R} \to y \cdot x = x \cdot y) \to (x \in \mathbb{R} \to y \in \mathbb{R} \to y \in \mathbb{R}) \to \varphi_{21}$
\item $\varphi_{34} :\equiv (y \in \mathbb{R} \to x \in \mathbb{R} \to -y \cdot x = -(y \cdot x)) \to (x \in \mathbb{R} \to y \in \mathbb{R} \to y \in \mathbb{R}) \to \varphi_{23}$
\item $\varphi_{35} :\equiv (x \in \mathbb{R} \to y \in \mathbb{R} \to -y \cdot x = -(y \cdot x)) \to \varphi_{20} \to x \in \mathbb{R} \to y \in \mathbb{R} \to x \cdot -y = -(y \cdot x)$
\item $\varphi_{36} :\equiv (x \in \mathbb{R} \to y \in \mathbb{R} \to -(y \cdot x) = -(x \cdot y)) \to \varphi_{19} \to x \in \mathbb{R} \to y \in \mathbb{R} \to x \cdot -y = -(x \cdot y)$
\item $\varphi_{37} :\equiv (x \in \mathbb{R} \to -y \in \mathbb{R} \to x \cdot -y = -y \cdot x) \to (x \in \mathbb{R} \to y \in \mathbb{R} \to x \in \mathbb{R}) \to \varphi_{25}$
\end{abbrevlist}
\end{fullproof}

% 自動生成 (prover/build.js): NegNegMul  ノード数 94
\begin{fullproof}{[F11]}{$\forall x\,\forall y\,(x \in \mathbb{R} \to y \in \mathbb{R} \to -x \cdot -y = x \cdot y)$}{94}\phantomsection\label{tree:NegNegMul}
\noindent$\Pi_{1}$:
\begin{ptree}
\AXC{$y = z \to x = y \to x = z$\lbl{A6}}
\AXC{$(y = z \to x = y \to x = z) \to x = y \to y = z \to x = z$\lbl{A1}}
\BIC{$x = y \to y = z \to x = z$}
\UIC{$\forall z\,(x = y \to y = z \to x = z)$}
\UIC{$\forall y\,\forall z\,(x = y \to y = z \to x = z)$}
\UIC{$\forall x\,\forall y\,\forall z\,(x = y \to y = z \to x = z)$}
\AXC{$\varphi_{3}$\lbl{A2}}
\BIC{$\forall y\,\forall z\,(v_{117} = y \to y = z \to v_{117} = z)$}
\AXC{$\varphi_{5}$\lbl{A2}}
\BIC{$\forall z\,(v_{117} = v_{118} \to v_{118} = z \to v_{117} = z)$}
\AXC{$\varphi_{8}$\lbl{A2}}
\BIC{$v_{117} = v_{118} \to v_{118} = v_{119} \to v_{117} = v_{119}$}
\UIC{$\forall v_{119}\,(v_{117} = v_{118} \to v_{118} = v_{119} \to v_{117} = v_{119})$}
\UIC{$\varphi_{1}$}
\UIC{$\forall v_{117}\,(\varphi_{1})$}
\end{ptree}
\noindent$\Pi_{2}$:
\begin{ptree}
\AXC{$\forall x \in \mathbb{R}\,(-(-x) = x)$\lbl{[F5]}}
\AXC{$\varphi_{14}$\lbl{A2}}
\BIC{$x \cdot y \in \mathbb{R} \to -(-(x \cdot y)) = x \cdot y$}
\AXC{$(x \cdot y \in \mathbb{R} \to -(-(x \cdot y)) = x \cdot y) \to \varphi_{21}$\lbl{A1}}
\BIC{$\varphi_{21}$}
\end{ptree}
\noindent$\Pi_{3}$:
\begin{ptree}
\AXC{$y = z \to x = y \to x = z$\lbl{A6}}
\AXC{$(y = z \to x = y \to x = z) \to x = y \to y = z \to x = z$\lbl{A1}}
\BIC{$x = y \to y = z \to x = z$}
\UIC{$\forall z\,(x = y \to y = z \to x = z)$}
\UIC{$\forall y\,\forall z\,(x = y \to y = z \to x = z)$}
\UIC{$\forall x\,\forall y\,\forall z\,(x = y \to y = z \to x = z)$}
\AXC{$\varphi_{4}$\lbl{A2}}
\BIC{$\forall y\,\forall z\,(v_{114} = y \to y = z \to v_{114} = z)$}
\AXC{$\varphi_{6}$\lbl{A2}}
\BIC{$\forall z\,(v_{114} = v_{115} \to v_{115} = z \to v_{114} = z)$}
\AXC{$\varphi_{9}$\lbl{A2}}
\BIC{$v_{114} = v_{115} \to v_{115} = v_{116} \to v_{114} = v_{116}$}
\UIC{$\forall v_{116}\,(v_{114} = v_{115} \to v_{115} = v_{116} \to v_{114} = v_{116})$}
\UIC{$\varphi_{2}$}
\UIC{$\forall v_{114}\,(\varphi_{2})$}
\end{ptree}
\noindent$\Pi_{4}$:
\begin{ptree}
\AXC{$\stackrel{\Pi_{3}}{\vdots}$}\noLine
\UIC{$\forall v_{114}\,(\varphi_{2})$}
\AXC{$\forall v_{114}\,(\varphi_{2}) \to \varphi_{11}$\lbl{A2}}
\BIC{$\varphi_{11}$}
\AXC{$\varphi_{11} \to \varphi_{19}$\lbl{A2}}
\BIC{$\varphi_{19}$}
\AXC{$\varphi_{19} \to -x \cdot -y = -(x \cdot -y) \to \varphi_{12}$\lbl{A2}}
\BIC{$-x \cdot -y = -(x \cdot -y) \to \varphi_{12}$}
\end{ptree}
\noindent$\Pi_{5}$:
\begin{ptree}
\AXC{$z = w \to -(x \cdot -y) = -z \to -(x \cdot -y) = -w$\lbl{A6}}
\UIC{$\forall w\,(z = w \to -(x \cdot -y) = -z \to -(x \cdot -y) = -w)$}
\UIC{$\varphi_{7}$}
\AXC{$\varphi_{7} \to \varphi_{15}$\lbl{A2}}
\BIC{$\varphi_{15}$}
\AXC{$\varphi_{15} \to x \cdot -y = -(x \cdot y) \to \varphi_{13}$\lbl{A2}}
\BIC{$x \cdot -y = -(x \cdot y) \to \varphi_{13}$}
\end{ptree}
\noindent$\Pi_{6}$:
\begin{ptree}
\AXC{$\forall x\,\forall y\,(x \in \mathbb{R} \to y \in \mathbb{R} \to x \cdot -y = -(x \cdot y))$\lbl{[F10]}}
\AXC{$\varphi_{25}$\lbl{A2}}
\BIC{$\forall y\,(x \in \mathbb{R} \to y \in \mathbb{R} \to x \cdot -y = -(x \cdot y))$}
\AXC{$\varphi_{24}$\lbl{A2}}
\BIC{$x \in \mathbb{R} \to y \in \mathbb{R} \to x \cdot -y = -(x \cdot y)$}
\AXC{$\varphi_{33}$\lbl{A1}}
\BIC{$-(x \cdot -y) = -(x \cdot -y) \to \varphi_{28}$}
\end{ptree}
\noindent$\Pi_{7}$:
\begin{ptree}
\AXC{$x = x$\lbl{A5}}
\UIC{$\forall x\,(x = x)$}
\AXC{$\forall x\,(x = x) \to -(x \cdot -y) = -(x \cdot -y)$\lbl{A2}}
\BIC{$-(x \cdot -y) = -(x \cdot -y)$}
\AXC{$\stackrel{\Pi_{6}}{\vdots}$}\noLine
\UIC{$-(x \cdot -y) = -(x \cdot -y) \to \varphi_{28}$}
\BIC{$\varphi_{28}$}
\end{ptree}
\noindent$\Pi_{8}$:
\begin{ptree}
\AXC{$\forall x \in \mathbb{R}\,(-x \in \mathbb{R})$\lbl{(R1)}}
\AXC{$\forall x \in \mathbb{R}\,(-x \in \mathbb{R}) \to y \in \mathbb{R} \to -y \in \mathbb{R}$\lbl{A2}}
\BIC{$y \in \mathbb{R} \to -y \in \mathbb{R}$}
\AXC{$(y \in \mathbb{R} \to -y \in \mathbb{R}) \to \varphi_{16}$\lbl{A1}}
\BIC{$\varphi_{16}$}
\end{ptree}
\noindent$\Pi_{9}$:
\begin{ptree}
\AXC{$x \in \mathbb{R} \to y \in \mathbb{R} \to x \in \mathbb{R}$\lbl{A1}}
\AXC{$\forall x\,\forall y\,(x \in \mathbb{R} \to y \in \mathbb{R} \to -x \cdot y = -(x \cdot y))$\lbl{[F9]}}
\AXC{$\varphi_{26}$\lbl{A2}}
\BIC{$\forall y\,(x \in \mathbb{R} \to y \in \mathbb{R} \to -x \cdot y = -(x \cdot y))$}
\AXC{$\varphi_{27}$\lbl{A2}}
\BIC{$x \in \mathbb{R} \to -y \in \mathbb{R} \to -x \cdot -y = -(x \cdot -y)$}
\AXC{$\varphi_{31}$\lbl{A1}}
\BIC{$(x \in \mathbb{R} \to y \in \mathbb{R} \to x \in \mathbb{R}) \to \varphi_{23}$}
\BIC{$\varphi_{23}$}
\end{ptree}
\noindent$\Pi_{10}$:
\begin{ptree}
\AXC{$x \in \mathbb{R} \to y \in \mathbb{R} \to y \in \mathbb{R}$\lbl{A1}}
\AXC{$\stackrel{\Pi_{8}}{\vdots}$}\noLine
\UIC{$\varphi_{16}$}
\BIC{$x \in \mathbb{R} \to y \in \mathbb{R} \to -y \in \mathbb{R}$}
\AXC{$\stackrel{\Pi_{9}}{\vdots}$}\noLine
\UIC{$\varphi_{23}$}
\BIC{$x \in \mathbb{R} \to y \in \mathbb{R} \to -x \cdot -y = -(x \cdot -y)$}
\AXC{$(x \in \mathbb{R} \to y \in \mathbb{R} \to -x \cdot -y = -(x \cdot -y)) \to \varphi_{32}$\lbl{A1}}
\BIC{$\varphi_{32}$}
\end{ptree}
\noindent$\Pi_{11}$:
\begin{ptree}
\AXC{$\stackrel{\Pi_{4}}{\vdots}$}\noLine
\UIC{$-x \cdot -y = -(x \cdot -y) \to \varphi_{12}$}
\AXC{$\stackrel{\Pi_{5}}{\vdots}$}\noLine
\UIC{$x \cdot -y = -(x \cdot y) \to \varphi_{13}$}
\AXC{$\stackrel{\Pi_{7}}{\vdots}$}\noLine
\UIC{$\varphi_{28}$}
\BIC{$x \in \mathbb{R} \to y \in \mathbb{R} \to -(x \cdot -y) = -(-(x \cdot y))$}
\AXC{$\stackrel{\Pi_{10}}{\vdots}$}\noLine
\UIC{$\varphi_{32}$}
\BIC{$\varphi_{29}$}
\BIC{$x \in \mathbb{R} \to y \in \mathbb{R} \to -x \cdot -y = -(-(x \cdot y))$}
\end{ptree}
\noindent$\Pi_{12}$:
\begin{ptree}
\AXC{$\stackrel{\Pi_{11}}{\vdots}$}\noLine
\UIC{$x \in \mathbb{R} \to y \in \mathbb{R} \to -x \cdot -y = -(-(x \cdot y))$}
\AXC{$(x \in \mathbb{R} \to y \in \mathbb{R} \to -x \cdot -y = -(-(x \cdot y))) \to \varphi_{30}$\lbl{A1}}
\BIC{$\varphi_{30}$}
\end{ptree}
\noindent$\Pi_{13}$:
\begin{ptree}
\AXC{$\forall x\,\forall y\,(x \in \mathbb{R} \to y \in \mathbb{R} \to x \cdot y \in \mathbb{R})$\lbl{(R1)}}
\AXC{$\varphi_{22}$\lbl{A2}}
\BIC{$\forall y\,(x \in \mathbb{R} \to y \in \mathbb{R} \to x \cdot y \in \mathbb{R})$}
\AXC{$\varphi_{20}$\lbl{A2}}
\BIC{$x \in \mathbb{R} \to y \in \mathbb{R} \to x \cdot y \in \mathbb{R}$}
\AXC{$\stackrel{\Pi_{2}}{\vdots}$}\noLine
\UIC{$\varphi_{21}$}
\BIC{$x \in \mathbb{R} \to y \in \mathbb{R} \to -(-(x \cdot y)) = x \cdot y$}
\AXC{$\stackrel{\Pi_{12}}{\vdots}$}\noLine
\UIC{$\varphi_{30}$}
\BIC{$\varphi_{17} \to x \in \mathbb{R} \to y \in \mathbb{R} \to -x \cdot -y = x \cdot y$}
\end{ptree}
\noindent 主木:
\begin{ptree}
\AXC{$\stackrel{\Pi_{1}}{\vdots}$}\noLine
\UIC{$\forall v_{117}\,(\varphi_{1})$}
\AXC{$\forall v_{117}\,(\varphi_{1}) \to \varphi_{10}$\lbl{A2}}
\BIC{$\varphi_{10}$}
\AXC{$\varphi_{10} \to \varphi_{18}$\lbl{A2}}
\BIC{$\varphi_{18}$}
\AXC{$\varphi_{18} \to \varphi_{17}$\lbl{A2}}
\BIC{$\varphi_{17}$}
\AXC{$\stackrel{\Pi_{13}}{\vdots}$}\noLine
\UIC{$\varphi_{17} \to x \in \mathbb{R} \to y \in \mathbb{R} \to -x \cdot -y = x \cdot y$}
\BIC{$x \in \mathbb{R} \to y \in \mathbb{R} \to -x \cdot -y = x \cdot y$}
\UIC{$\forall y\,(x \in \mathbb{R} \to y \in \mathbb{R} \to -x \cdot -y = x \cdot y)$}
\UIC{$\forall x\,\forall y\,(x \in \mathbb{R} \to y \in \mathbb{R} \to -x \cdot -y = x \cdot y)$}
\end{ptree}
\begin{abbrevlist}
\item $\varphi_{1} :\equiv \forall v_{118}\,\forall v_{119}\,(v_{117} = v_{118} \to v_{118} = v_{119} \to v_{117} = v_{119})$
\item $\varphi_{2} :\equiv \forall v_{115}\,\forall v_{116}\,(v_{114} = v_{115} \to v_{115} = v_{116} \to v_{114} = v_{116})$
\item $\varphi_{3} :\equiv \forall x\,\forall y\,\forall z\,(x = y \to y = z \to x = z) \to \forall y\,\forall z\,(v_{117} = y \to y = z \to v_{117} = z)$
\item $\varphi_{4} :\equiv \forall x\,\forall y\,\forall z\,(x = y \to y = z \to x = z) \to \forall y\,\forall z\,(v_{114} = y \to y = z \to v_{114} = z)$
\item $\varphi_{5} :\equiv \forall y\,\forall z\,(v_{117} = y \to y = z \to v_{117} = z) \to \forall z\,(v_{117} = v_{118} \to v_{118} = z \to v_{117} = z)$
\item $\varphi_{6} :\equiv \forall y\,\forall z\,(v_{114} = y \to y = z \to v_{114} = z) \to \forall z\,(v_{114} = v_{115} \to v_{115} = z \to v_{114} = z)$
\item $\varphi_{7} :\equiv \forall z\,\forall w\,(z = w \to -(x \cdot -y) = -z \to -(x \cdot -y) = -w)$
\item $\varphi_{8} :\equiv \forall z\,(v_{117} = v_{118} \to v_{118} = z \to v_{117} = z) \to v_{117} = v_{118} \to v_{118} = v_{119} \to v_{117} = v_{119}$
\item $\varphi_{9} :\equiv \forall z\,(v_{114} = v_{115} \to v_{115} = z \to v_{114} = z) \to v_{114} = v_{115} \to v_{115} = v_{116} \to v_{114} = v_{116}$
\item $\varphi_{10} :\equiv \forall v_{118}\,\forall v_{119}\,(-x \cdot -y = v_{118} \to v_{118} = v_{119} \to -x \cdot -y = v_{119})$
\item $\varphi_{11} :\equiv \forall v_{115}\,\forall v_{116}\,(-x \cdot -y = v_{115} \to v_{115} = v_{116} \to -x \cdot -y = v_{116})$
\item $\varphi_{12} :\equiv -(x \cdot -y) = -(-(x \cdot y)) \to -x \cdot -y = -(-(x \cdot y))$
\item $\varphi_{13} :\equiv -(x \cdot -y) = -(x \cdot -y) \to -(x \cdot -y) = -(-(x \cdot y))$
\item $\varphi_{14} :\equiv \forall x \in \mathbb{R}\,(-(-x) = x) \to x \cdot y \in \mathbb{R} \to -(-(x \cdot y)) = x \cdot y$
\item $\varphi_{15} :\equiv \forall w\,(x \cdot -y = w \to -(x \cdot -y) = -(x \cdot -y) \to -(x \cdot -y) = -w)$
\item $\varphi_{16} :\equiv (x \in \mathbb{R} \to y \in \mathbb{R} \to y \in \mathbb{R}) \to x \in \mathbb{R} \to y \in \mathbb{R} \to -y \in \mathbb{R}$
\item $\varphi_{17} :\equiv -x \cdot -y = -(-(x \cdot y)) \to -(-(x \cdot y)) = x \cdot y \to -x \cdot -y = x \cdot y$
\item $\varphi_{18} :\equiv \forall v_{119}\,(-x \cdot -y = -(-(x \cdot y)) \to -(-(x \cdot y)) = v_{119} \to -x \cdot -y = v_{119})$
\item $\varphi_{19} :\equiv \forall v_{116}\,(-x \cdot -y = -(x \cdot -y) \to -(x \cdot -y) = v_{116} \to -x \cdot -y = v_{116})$
\item $\varphi_{20} :\equiv \forall y\,(x \in \mathbb{R} \to y \in \mathbb{R} \to x \cdot y \in \mathbb{R}) \to x \in \mathbb{R} \to y \in \mathbb{R} \to x \cdot y \in \mathbb{R}$
\item $\varphi_{21} :\equiv (x \in \mathbb{R} \to y \in \mathbb{R} \to x \cdot y \in \mathbb{R}) \to x \in \mathbb{R} \to y \in \mathbb{R} \to -(-(x \cdot y)) = x \cdot y$
\item $\varphi_{22} :\equiv \forall x\,\forall y\,(x \in \mathbb{R} \to y \in \mathbb{R} \to x \cdot y \in \mathbb{R}) \to \forall y\,(x \in \mathbb{R} \to y \in \mathbb{R} \to x \cdot y \in \mathbb{R})$
\item $\varphi_{23} :\equiv (x \in \mathbb{R} \to y \in \mathbb{R} \to -y \in \mathbb{R}) \to x \in \mathbb{R} \to y \in \mathbb{R} \to -x \cdot -y = -(x \cdot -y)$
\item $\varphi_{24} :\equiv \forall y\,(x \in \mathbb{R} \to y \in \mathbb{R} \to x \cdot -y = -(x \cdot y)) \to x \in \mathbb{R} \to y \in \mathbb{R} \to x \cdot -y = -(x \cdot y)$
\item $\varphi_{25} :\equiv \forall x\,\forall y\,(x \in \mathbb{R} \to y \in \mathbb{R} \to x \cdot -y = -(x \cdot y)) \to \forall y\,(x \in \mathbb{R} \to y \in \mathbb{R} \to x \cdot -y = -(x \cdot y))$
\item $\varphi_{26} :\equiv \forall x\,\forall y\,(x \in \mathbb{R} \to y \in \mathbb{R} \to -x \cdot y = -(x \cdot y)) \to \forall y\,(x \in \mathbb{R} \to y \in \mathbb{R} \to -x \cdot y = -(x \cdot y))$
\item $\varphi_{27} :\equiv \forall y\,(x \in \mathbb{R} \to y \in \mathbb{R} \to -x \cdot y = -(x \cdot y)) \to x \in \mathbb{R} \to -y \in \mathbb{R} \to -x \cdot -y = -(x \cdot -y)$
\item $\varphi_{28} :\equiv (x \cdot -y = -(x \cdot y) \to \varphi_{13}) \to x \in \mathbb{R} \to y \in \mathbb{R} \to -(x \cdot -y) = -(-(x \cdot y))$
\item $\varphi_{29} :\equiv (-x \cdot -y = -(x \cdot -y) \to \varphi_{12}) \to x \in \mathbb{R} \to y \in \mathbb{R} \to -x \cdot -y = -(-(x \cdot y))$
\item $\varphi_{30} :\equiv (x \in \mathbb{R} \to y \in \mathbb{R} \to -(-(x \cdot y)) = x \cdot y) \to \varphi_{17} \to x \in \mathbb{R} \to y \in \mathbb{R} \to -x \cdot -y = x \cdot y$
\item $\varphi_{31} :\equiv (x \in \mathbb{R} \to -y \in \mathbb{R} \to -x \cdot -y = -(x \cdot -y)) \to (x \in \mathbb{R} \to y \in \mathbb{R} \to x \in \mathbb{R}) \to \varphi_{23}$
\item $\varphi_{32} :\equiv (x \in \mathbb{R} \to y \in \mathbb{R} \to -(x \cdot -y) = -(-(x \cdot y))) \to \varphi_{29}$
\item $\varphi_{33} :\equiv (x \in \mathbb{R} \to y \in \mathbb{R} \to x \cdot -y = -(x \cdot y)) \to -(x \cdot -y) = -(x \cdot -y) \to \varphi_{28}$
\end{abbrevlist}
\end{fullproof}

% 自動生成 (prover/build.js): SubAdd  ノード数 181
\begin{fullproof}{[F12]}{$\forall x\,\forall y\,(x \in \mathbb{R} \to y \in \mathbb{R} \to x - y + y = x)$}{181}\phantomsection\label{tree:SubAdd}
\noindent$\Pi_{1}$:
\begin{ptree}
\AXC{$y = z \to x = y \to x = z$\lbl{A6}}
\AXC{$(y = z \to x = y \to x = z) \to x = y \to y = z \to x = z$\lbl{A1}}
\BIC{$x = y \to y = z \to x = z$}
\UIC{$\forall z\,(x = y \to y = z \to x = z)$}
\UIC{$\forall y\,\forall z\,(x = y \to y = z \to x = z)$}
\UIC{$\forall x\,\forall y\,\forall z\,(x = y \to y = z \to x = z)$}
\AXC{$\varphi_{5}$\lbl{A2}}
\BIC{$\forall y\,\forall z\,(v_{131} = y \to y = z \to v_{131} = z)$}
\AXC{$\varphi_{17}$\lbl{A2}}
\BIC{$\forall z\,(v_{131} = v_{132} \to v_{132} = z \to v_{131} = z)$}
\AXC{$\varphi_{23}$\lbl{A2}}
\BIC{$v_{131} = v_{132} \to v_{132} = v_{133} \to v_{131} = v_{133}$}
\UIC{$\forall v_{133}\,(v_{131} = v_{132} \to v_{132} = v_{133} \to v_{131} = v_{133})$}
\UIC{$\varphi_{1}$}
\UIC{$\forall v_{131}\,(\varphi_{1})$}
\end{ptree}
\noindent$\Pi_{2}$:
\begin{ptree}
\AXC{$\stackrel{\Pi_{1}}{\vdots}$}\noLine
\UIC{$\forall v_{131}\,(\varphi_{1})$}
\AXC{$\forall v_{131}\,(\varphi_{1}) \to \varphi_{13}$\lbl{A2}}
\BIC{$\varphi_{13}$}
\AXC{$\varphi_{13} \to \varphi_{9}$\lbl{A2}}
\BIC{$\varphi_{9}$}
\AXC{$\varphi_{9} \to x - y + y = x + 0 \to x + 0 = x \to x - y + y = x$\lbl{A2}}
\BIC{$x - y + y = x + 0 \to x + 0 = x \to x - y + y = x$}
\end{ptree}
\noindent$\Pi_{3}$:
\begin{ptree}
\AXC{$\forall x \in \mathbb{R}\,(x + 0 = x)$\lbl{(R4)}}
\AXC{$\forall x \in \mathbb{R}\,(x + 0 = x) \to x \in \mathbb{R} \to x + 0 = x$\lbl{A2}}
\BIC{$x \in \mathbb{R} \to x + 0 = x$}
\AXC{$(x \in \mathbb{R} \to x + 0 = x) \to \varphi_{30}$\lbl{A1}}
\BIC{$\varphi_{30}$}
\end{ptree}
\noindent$\Pi_{4}$:
\begin{ptree}
\AXC{$y = z \to x = y \to x = z$\lbl{A6}}
\AXC{$(y = z \to x = y \to x = z) \to x = y \to y = z \to x = z$\lbl{A1}}
\BIC{$x = y \to y = z \to x = z$}
\UIC{$\forall z\,(x = y \to y = z \to x = z)$}
\UIC{$\forall y\,\forall z\,(x = y \to y = z \to x = z)$}
\UIC{$\forall x\,\forall y\,\forall z\,(x = y \to y = z \to x = z)$}
\AXC{$\varphi_{6}$\lbl{A2}}
\BIC{$\forall y\,\forall z\,(v_{128} = y \to y = z \to v_{128} = z)$}
\AXC{$\varphi_{18}$\lbl{A2}}
\BIC{$\forall z\,(v_{128} = v_{129} \to v_{129} = z \to v_{128} = z)$}
\AXC{$\varphi_{24}$\lbl{A2}}
\BIC{$v_{128} = v_{129} \to v_{129} = v_{130} \to v_{128} = v_{130}$}
\UIC{$\forall v_{130}\,(v_{128} = v_{129} \to v_{129} = v_{130} \to v_{128} = v_{130})$}
\UIC{$\varphi_{2}$}
\UIC{$\forall v_{128}\,(\varphi_{2})$}
\end{ptree}
\noindent$\Pi_{5}$:
\begin{ptree}
\AXC{$z = w \to x + (y + -y) = x + z \to x + (y + -y) = x + w$\lbl{A6}}
\UIC{$\varphi_{10}$}
\UIC{$\forall z\,(\varphi_{10})$}
\AXC{$\forall z\,(\varphi_{10}) \to \varphi_{38}$\lbl{A2}}
\BIC{$\varphi_{38}$}
\AXC{$\varphi_{38} \to \varphi_{32}$\lbl{A2}}
\BIC{$\varphi_{32}$}
\end{ptree}
\noindent$\Pi_{6}$:
\begin{ptree}
\AXC{$\forall x \in \mathbb{R}\,(x + -x = 0)$\lbl{(R5)}}
\AXC{$\forall x \in \mathbb{R}\,(x + -x = 0) \to y \in \mathbb{R} \to y + -y = 0$\lbl{A2}}
\BIC{$y \in \mathbb{R} \to y + -y = 0$}
\AXC{$(y \in \mathbb{R} \to y + -y = 0) \to \varphi_{39}$\lbl{A1}}
\BIC{$\varphi_{39}$}
\end{ptree}
\noindent$\Pi_{7}$:
\begin{ptree}
\AXC{$x \in \mathbb{R} \to y \in \mathbb{R} \to y \in \mathbb{R}$\lbl{A1}}
\AXC{$\stackrel{\Pi_{6}}{\vdots}$}\noLine
\UIC{$\varphi_{39}$}
\BIC{$x \in \mathbb{R} \to y \in \mathbb{R} \to y + -y = 0$}
\AXC{$(x \in \mathbb{R} \to y \in \mathbb{R} \to y + -y = 0) \to \varphi_{53}$\lbl{A1}}
\BIC{$\varphi_{53}$}
\end{ptree}
\noindent$\Pi_{8}$:
\begin{ptree}
\AXC{$y = z \to x = y \to x = z$\lbl{A6}}
\AXC{$(y = z \to x = y \to x = z) \to x = y \to y = z \to x = z$\lbl{A1}}
\BIC{$x = y \to y = z \to x = z$}
\UIC{$\forall z\,(x = y \to y = z \to x = z)$}
\UIC{$\forall y\,\forall z\,(x = y \to y = z \to x = z)$}
\UIC{$\forall x\,\forall y\,\forall z\,(x = y \to y = z \to x = z)$}
\AXC{$\varphi_{7}$\lbl{A2}}
\BIC{$\forall y\,\forall z\,(v_{125} = y \to y = z \to v_{125} = z)$}
\AXC{$\varphi_{19}$\lbl{A2}}
\BIC{$\forall z\,(v_{125} = v_{126} \to v_{126} = z \to v_{125} = z)$}
\AXC{$\varphi_{25}$\lbl{A2}}
\BIC{$v_{125} = v_{126} \to v_{126} = v_{127} \to v_{125} = v_{127}$}
\UIC{$\forall v_{127}\,(v_{125} = v_{126} \to v_{126} = v_{127} \to v_{125} = v_{127})$}
\UIC{$\varphi_{3}$}
\UIC{$\forall v_{125}\,(\varphi_{3})$}
\end{ptree}
\noindent$\Pi_{9}$:
\begin{ptree}
\AXC{$\stackrel{\Pi_{8}}{\vdots}$}\noLine
\UIC{$\forall v_{125}\,(\varphi_{3})$}
\AXC{$\forall v_{125}\,(\varphi_{3}) \to \varphi_{15}$\lbl{A2}}
\BIC{$\varphi_{15}$}
\AXC{$\varphi_{15} \to \varphi_{35}$\lbl{A2}}
\BIC{$\varphi_{35}$}
\AXC{$\varphi_{35} \to x - y + y = x + (-y + y) \to \varphi_{22}$\lbl{A2}}
\BIC{$x - y + y = x + (-y + y) \to \varphi_{22}$}
\end{ptree}
\noindent$\Pi_{10}$:
\begin{ptree}
\AXC{$\forall x \in \mathbb{R}\,(-x \in \mathbb{R})$\lbl{(R1)}}
\AXC{$\forall x \in \mathbb{R}\,(-x \in \mathbb{R}) \to y \in \mathbb{R} \to -y \in \mathbb{R}$\lbl{A2}}
\BIC{$y \in \mathbb{R} \to -y \in \mathbb{R}$}
\AXC{$(y \in \mathbb{R} \to -y \in \mathbb{R}) \to \varphi_{41}$\lbl{A1}}
\BIC{$\varphi_{41}$}
\end{ptree}
\noindent$\Pi_{11}$:
\begin{ptree}
\AXC{$x \in \mathbb{R} \to y \in \mathbb{R} \to y \in \mathbb{R}$\lbl{A1}}
\AXC{$\stackrel{\Pi_{10}}{\vdots}$}\noLine
\UIC{$\varphi_{41}$}
\BIC{$x \in \mathbb{R} \to y \in \mathbb{R} \to -y \in \mathbb{R}$}
\end{ptree}
\noindent$\Pi_{12}$:
\begin{ptree}
\AXC{$\forall x\,\forall y\,(x \in \mathbb{R} \to y \in \mathbb{R} \to x + y = y + x)$\lbl{(R3)}}
\AXC{$\varphi_{47}$\lbl{A2}}
\BIC{$\forall y\,(v_{120} \in \mathbb{R} \to y \in \mathbb{R} \to v_{120} + y = y + v_{120})$}
\AXC{$\varphi_{48}$\lbl{A2}}
\BIC{$v_{120} \in \mathbb{R} \to v_{121} \in \mathbb{R} \to v_{120} + v_{121} = v_{121} + v_{120}$}
\UIC{$\varphi_{21}$}
\UIC{$\forall v_{120}\,(\varphi_{21})$}
\AXC{$\forall v_{120}\,(\varphi_{21}) \to \varphi_{28}$\lbl{A2}}
\BIC{$\varphi_{28}$}
\end{ptree}
\noindent$\Pi_{13}$:
\begin{ptree}
\AXC{$\stackrel{\Pi_{11}}{\vdots}$}\noLine
\UIC{$x \in \mathbb{R} \to y \in \mathbb{R} \to -y \in \mathbb{R}$}
\AXC{$\stackrel{\Pi_{12}}{\vdots}$}\noLine
\UIC{$\varphi_{28}$}
\AXC{$\varphi_{28} \to -y \in \mathbb{R} \to y \in \mathbb{R} \to -y + y = y + -y$\lbl{A2}}
\BIC{$-y \in \mathbb{R} \to y \in \mathbb{R} \to -y + y = y + -y$}
\AXC{$\varphi_{57}$\lbl{A1}}
\BIC{$(x \in \mathbb{R} \to y \in \mathbb{R} \to -y \in \mathbb{R}) \to \varphi_{44}$}
\BIC{$\varphi_{44}$}
\end{ptree}
\noindent$\Pi_{14}$:
\begin{ptree}
\AXC{$x = x$\lbl{A5}}
\UIC{$\forall x\,(x = x)$}
\AXC{$\forall x\,(x = x) \to x + (-y + y) = x + (-y + y)$\lbl{A2}}
\BIC{$x + (-y + y) = x + (-y + y)$}
\AXC{$x \in \mathbb{R} \to y \in \mathbb{R} \to y \in \mathbb{R}$\lbl{A1}}
\AXC{$\stackrel{\Pi_{13}}{\vdots}$}\noLine
\UIC{$\varphi_{44}$}
\BIC{$x \in \mathbb{R} \to y \in \mathbb{R} \to -y + y = y + -y$}
\AXC{$\varphi_{60}$\lbl{A1}}
\BIC{$x + (-y + y) = x + (-y + y) \to \varphi_{51}$}
\BIC{$\varphi_{51}$}
\end{ptree}
\noindent$\Pi_{15}$:
\begin{ptree}
\AXC{$z = w \to x + (-y + y) = x + z \to x + (-y + y) = x + w$\lbl{A6}}
\UIC{$\varphi_{11}$}
\UIC{$\forall z\,(\varphi_{11})$}
\AXC{$\forall z\,(\varphi_{11}) \to \varphi_{40}$\lbl{A2}}
\BIC{$\varphi_{40}$}
\AXC{$\varphi_{40} \to -y + y = y + -y \to \varphi_{29}$\lbl{A2}}
\BIC{$-y + y = y + -y \to \varphi_{29}$}
\AXC{$\stackrel{\Pi_{14}}{\vdots}$}\noLine
\UIC{$\varphi_{51}$}
\BIC{$x \in \mathbb{R} \to y \in \mathbb{R} \to x + (-y + y) = x + (y + -y)$}
\end{ptree}
\noindent$\Pi_{16}$:
\begin{ptree}
\AXC{$y = z \to x = y \to x = z$\lbl{A6}}
\AXC{$(y = z \to x = y \to x = z) \to x = y \to y = z \to x = z$\lbl{A1}}
\BIC{$x = y \to y = z \to x = z$}
\UIC{$\forall z\,(x = y \to y = z \to x = z)$}
\UIC{$\forall y\,\forall z\,(x = y \to y = z \to x = z)$}
\UIC{$\forall x\,\forall y\,\forall z\,(x = y \to y = z \to x = z)$}
\AXC{$\varphi_{8}$\lbl{A2}}
\BIC{$\forall y\,\forall z\,(v_{122} = y \to y = z \to v_{122} = z)$}
\AXC{$\varphi_{20}$\lbl{A2}}
\BIC{$\forall z\,(v_{122} = v_{123} \to v_{123} = z \to v_{122} = z)$}
\AXC{$\varphi_{26}$\lbl{A2}}
\BIC{$v_{122} = v_{123} \to v_{123} = v_{124} \to v_{122} = v_{124}$}
\UIC{$\forall v_{124}\,(v_{122} = v_{123} \to v_{123} = v_{124} \to v_{122} = v_{124})$}
\UIC{$\varphi_{4}$}
\UIC{$\forall v_{122}\,(\varphi_{4})$}
\end{ptree}
\noindent$\Pi_{17}$:
\begin{ptree}
\AXC{$\stackrel{\Pi_{16}}{\vdots}$}\noLine
\UIC{$\forall v_{122}\,(\varphi_{4})$}
\AXC{$\forall v_{122}\,(\varphi_{4}) \to \varphi_{16}$\lbl{A2}}
\BIC{$\varphi_{16}$}
\AXC{$\varphi_{16} \to \varphi_{36}$\lbl{A2}}
\BIC{$\varphi_{36}$}
\AXC{$\varphi_{36} \to \varphi_{45}$\lbl{A2}}
\BIC{$\varphi_{45}$}
\end{ptree}
\noindent$\Pi_{18}$:
\begin{ptree}
\AXC{$x \in \mathbb{R} \to y \in \mathbb{R} \to x \in \mathbb{R}$\lbl{A1}}
\AXC{$\forall x\,(\varphi_{27})$\lbl{(R2)}}
\AXC{$\forall x\,(\varphi_{27}) \to \varphi_{27}$\lbl{A2}}
\BIC{$\varphi_{27}$}
\AXC{$\varphi_{27} \to \varphi_{42}$\lbl{A2}}
\BIC{$\varphi_{42}$}
\AXC{$\varphi_{42} \to \varphi_{37}$\lbl{A2}}
\BIC{$\varphi_{37}$}
\AXC{$\varphi_{37} \to \varphi_{56}$\lbl{A1}}
\BIC{$\varphi_{56}$}
\BIC{$(x \in \mathbb{R} \to y \in \mathbb{R} \to -y \in \mathbb{R}) \to \varphi_{46}$}
\end{ptree}
\noindent$\Pi_{19}$:
\begin{ptree}
\AXC{$x \in \mathbb{R} \to y \in \mathbb{R} \to y \in \mathbb{R}$\lbl{A1}}
\AXC{$\stackrel{\Pi_{11}}{\vdots}$}\noLine
\UIC{$x \in \mathbb{R} \to y \in \mathbb{R} \to -y \in \mathbb{R}$}
\AXC{$\stackrel{\Pi_{18}}{\vdots}$}\noLine
\UIC{$(x \in \mathbb{R} \to y \in \mathbb{R} \to -y \in \mathbb{R}) \to \varphi_{46}$}
\BIC{$\varphi_{46}$}
\BIC{$x \in \mathbb{R} \to y \in \mathbb{R} \to x + -y + y = x + (-y + y)$}
\end{ptree}
\noindent$\Pi_{20}$:
\begin{ptree}
\AXC{$z = w \to x - y + y = z + y \to x - y + y = w + y$\lbl{A6}}
\UIC{$\forall w\,(z = w \to x - y + y = z + y \to x - y + y = w + y)$}
\UIC{$\forall z\,\forall w\,(z = w \to x - y + y = z + y \to x - y + y = w + y)$}
\AXC{$\varphi_{49}$\lbl{A2}}
\BIC{$\varphi_{12}$}
\AXC{$\varphi_{12} \to \varphi_{33}$\lbl{A2}}
\BIC{$\varphi_{33}$}
\end{ptree}
\noindent$\Pi_{21}$:
\begin{ptree}
\AXC{$\forall x\,\forall y\,(x - y = x + -y)$\lbl{Def$-$}}
\AXC{$\forall x\,\forall y\,(x - y = x + -y) \to \forall y\,(x - y = x + -y)$\lbl{A2}}
\BIC{$\forall y\,(x - y = x + -y)$}
\end{ptree}
\noindent$\Pi_{22}$:
\begin{ptree}
\AXC{$\stackrel{\Pi_{21}}{\vdots}$}\noLine
\UIC{$\forall y\,(x - y = x + -y)$}
\AXC{$\forall y\,(x - y = x + -y) \to x - y = x + -y$\lbl{A2}}
\BIC{$x - y = x + -y$}
\AXC{$x - y = x + -y \to \varphi_{52}$\lbl{A1}}
\BIC{$\varphi_{52}$}
\end{ptree}
\noindent$\Pi_{23}$:
\begin{ptree}
\AXC{$\stackrel{\Pi_{20}}{\vdots}$}\noLine
\UIC{$\varphi_{33}$}
\AXC{$x = x$\lbl{A5}}
\UIC{$\forall x\,(x = x)$}
\AXC{$\forall x\,(x = x) \to x - y + y = x - y + y$\lbl{A2}}
\BIC{$x - y + y = x - y + y$}
\AXC{$\stackrel{\Pi_{22}}{\vdots}$}\noLine
\UIC{$\varphi_{52}$}
\BIC{$\varphi_{33} \to x \in \mathbb{R} \to y \in \mathbb{R} \to x - y + y = x + -y + y$}
\BIC{$x \in \mathbb{R} \to y \in \mathbb{R} \to x - y + y = x + -y + y$}
\end{ptree}
\noindent$\Pi_{24}$:
\begin{ptree}
\AXC{$\stackrel{\Pi_{23}}{\vdots}$}\noLine
\UIC{$x \in \mathbb{R} \to y \in \mathbb{R} \to x - y + y = x + -y + y$}
\AXC{$(x \in \mathbb{R} \to y \in \mathbb{R} \to x - y + y = x + -y + y) \to \varphi_{58}$\lbl{A1}}
\BIC{$\varphi_{58}$}
\end{ptree}
\noindent$\Pi_{25}$:
\begin{ptree}
\AXC{$\stackrel{\Pi_{17}}{\vdots}$}\noLine
\UIC{$\varphi_{45}$}
\AXC{$\stackrel{\Pi_{19}}{\vdots}$}\noLine
\UIC{$x \in \mathbb{R} \to y \in \mathbb{R} \to x + -y + y = x + (-y + y)$}
\AXC{$\stackrel{\Pi_{24}}{\vdots}$}\noLine
\UIC{$\varphi_{58}$}
\BIC{$\varphi_{45} \to x \in \mathbb{R} \to y \in \mathbb{R} \to x - y + y = x + (-y + y)$}
\BIC{$x \in \mathbb{R} \to y \in \mathbb{R} \to x - y + y = x + (-y + y)$}
\AXC{$(x \in \mathbb{R} \to y \in \mathbb{R} \to x - y + y = x + (-y + y)) \to \varphi_{59}$\lbl{A1}}
\BIC{$\varphi_{59}$}
\end{ptree}
\noindent$\Pi_{26}$:
\begin{ptree}
\AXC{$\stackrel{\Pi_{9}}{\vdots}$}\noLine
\UIC{$x - y + y = x + (-y + y) \to \varphi_{22}$}
\AXC{$\stackrel{\Pi_{15}}{\vdots}$}\noLine
\UIC{$x \in \mathbb{R} \to y \in \mathbb{R} \to x + (-y + y) = x + (y + -y)$}
\AXC{$\stackrel{\Pi_{25}}{\vdots}$}\noLine
\UIC{$\varphi_{59}$}
\BIC{$\varphi_{50}$}
\BIC{$x \in \mathbb{R} \to y \in \mathbb{R} \to x - y + y = x + (y + -y)$}
\end{ptree}
\noindent$\Pi_{27}$:
\begin{ptree}
\AXC{$\stackrel{\Pi_{26}}{\vdots}$}\noLine
\UIC{$x \in \mathbb{R} \to y \in \mathbb{R} \to x - y + y = x + (y + -y)$}
\AXC{$(x \in \mathbb{R} \to y \in \mathbb{R} \to x - y + y = x + (y + -y)) \to \varphi_{55}$\lbl{A1}}
\BIC{$\varphi_{55}$}
\end{ptree}
\noindent$\Pi_{28}$:
\begin{ptree}
\AXC{$\stackrel{\Pi_{5}}{\vdots}$}\noLine
\UIC{$\varphi_{32}$}
\AXC{$x = x$\lbl{A5}}
\UIC{$\forall x\,(x = x)$}
\AXC{$\forall x\,(x = x) \to x + (y + -y) = x + (y + -y)$\lbl{A2}}
\BIC{$x + (y + -y) = x + (y + -y)$}
\AXC{$\stackrel{\Pi_{7}}{\vdots}$}\noLine
\UIC{$\varphi_{53}$}
\BIC{$\varphi_{32} \to x \in \mathbb{R} \to y \in \mathbb{R} \to x + (y + -y) = x + 0$}
\BIC{$x \in \mathbb{R} \to y \in \mathbb{R} \to x + (y + -y) = x + 0$}
\AXC{$\stackrel{\Pi_{27}}{\vdots}$}\noLine
\UIC{$\varphi_{55}$}
\BIC{$\varphi_{31} \to x \in \mathbb{R} \to y \in \mathbb{R} \to x - y + y = x + 0$}
\end{ptree}
\noindent$\Pi_{29}$:
\begin{ptree}
\AXC{$\stackrel{\Pi_{4}}{\vdots}$}\noLine
\UIC{$\forall v_{128}\,(\varphi_{2})$}
\AXC{$\forall v_{128}\,(\varphi_{2}) \to \varphi_{14}$\lbl{A2}}
\BIC{$\varphi_{14}$}
\AXC{$\varphi_{14} \to \varphi_{34}$\lbl{A2}}
\BIC{$\varphi_{34}$}
\AXC{$\varphi_{34} \to \varphi_{31}$\lbl{A2}}
\BIC{$\varphi_{31}$}
\AXC{$\stackrel{\Pi_{28}}{\vdots}$}\noLine
\UIC{$\varphi_{31} \to x \in \mathbb{R} \to y \in \mathbb{R} \to x - y + y = x + 0$}
\BIC{$x \in \mathbb{R} \to y \in \mathbb{R} \to x - y + y = x + 0$}
\end{ptree}
\noindent$\Pi_{30}$:
\begin{ptree}
\AXC{$x \in \mathbb{R} \to y \in \mathbb{R} \to x \in \mathbb{R}$\lbl{A1}}
\AXC{$\stackrel{\Pi_{3}}{\vdots}$}\noLine
\UIC{$\varphi_{30}$}
\BIC{$x \in \mathbb{R} \to y \in \mathbb{R} \to x + 0 = x$}
\AXC{$\stackrel{\Pi_{29}}{\vdots}$}\noLine
\UIC{$x \in \mathbb{R} \to y \in \mathbb{R} \to x - y + y = x + 0$}
\AXC{$\varphi_{54}$\lbl{A1}}
\BIC{$(x \in \mathbb{R} \to y \in \mathbb{R} \to x + 0 = x) \to \varphi_{43}$}
\BIC{$\varphi_{43}$}
\end{ptree}
\noindent 主木:
\begin{ptree}
\AXC{$\stackrel{\Pi_{2}}{\vdots}$}\noLine
\UIC{$x - y + y = x + 0 \to x + 0 = x \to x - y + y = x$}
\AXC{$\stackrel{\Pi_{30}}{\vdots}$}\noLine
\UIC{$\varphi_{43}$}
\BIC{$x \in \mathbb{R} \to y \in \mathbb{R} \to x - y + y = x$}
\UIC{$\forall y\,(x \in \mathbb{R} \to y \in \mathbb{R} \to x - y + y = x)$}
\UIC{$\forall x\,\forall y\,(x \in \mathbb{R} \to y \in \mathbb{R} \to x - y + y = x)$}
\end{ptree}
\begin{abbrevlist}
\item $\varphi_{1} :\equiv \forall v_{132}\,\forall v_{133}\,(v_{131} = v_{132} \to v_{132} = v_{133} \to v_{131} = v_{133})$
\item $\varphi_{2} :\equiv \forall v_{129}\,\forall v_{130}\,(v_{128} = v_{129} \to v_{129} = v_{130} \to v_{128} = v_{130})$
\item $\varphi_{3} :\equiv \forall v_{126}\,\forall v_{127}\,(v_{125} = v_{126} \to v_{126} = v_{127} \to v_{125} = v_{127})$
\item $\varphi_{4} :\equiv \forall v_{123}\,\forall v_{124}\,(v_{122} = v_{123} \to v_{123} = v_{124} \to v_{122} = v_{124})$
\item $\varphi_{5} :\equiv \forall x\,\forall y\,\forall z\,(x = y \to y = z \to x = z) \to \forall y\,\forall z\,(v_{131} = y \to y = z \to v_{131} = z)$
\item $\varphi_{6} :\equiv \forall x\,\forall y\,\forall z\,(x = y \to y = z \to x = z) \to \forall y\,\forall z\,(v_{128} = y \to y = z \to v_{128} = z)$
\item $\varphi_{7} :\equiv \forall x\,\forall y\,\forall z\,(x = y \to y = z \to x = z) \to \forall y\,\forall z\,(v_{125} = y \to y = z \to v_{125} = z)$
\item $\varphi_{8} :\equiv \forall x\,\forall y\,\forall z\,(x = y \to y = z \to x = z) \to \forall y\,\forall z\,(v_{122} = y \to y = z \to v_{122} = z)$
\item $\varphi_{9} :\equiv \forall v_{133}\,(x - y + y = x + 0 \to x + 0 = v_{133} \to x - y + y = v_{133})$
\item $\varphi_{10} :\equiv \forall w\,(z = w \to x + (y + -y) = x + z \to x + (y + -y) = x + w)$
\item $\varphi_{11} :\equiv \forall w\,(z = w \to x + (-y + y) = x + z \to x + (-y + y) = x + w)$
\item $\varphi_{12} :\equiv \forall w\,(x - y = w \to x - y + y = x - y + y \to x - y + y = w + y)$
\item $\varphi_{13} :\equiv \forall v_{132}\,\forall v_{133}\,(x - y + y = v_{132} \to v_{132} = v_{133} \to x - y + y = v_{133})$
\item $\varphi_{14} :\equiv \forall v_{129}\,\forall v_{130}\,(x - y + y = v_{129} \to v_{129} = v_{130} \to x - y + y = v_{130})$
\item $\varphi_{15} :\equiv \forall v_{126}\,\forall v_{127}\,(x - y + y = v_{126} \to v_{126} = v_{127} \to x - y + y = v_{127})$
\item $\varphi_{16} :\equiv \forall v_{123}\,\forall v_{124}\,(x - y + y = v_{123} \to v_{123} = v_{124} \to x - y + y = v_{124})$
\item $\varphi_{17} :\equiv \forall y\,\forall z\,(v_{131} = y \to y = z \to v_{131} = z) \to \forall z\,(v_{131} = v_{132} \to v_{132} = z \to v_{131} = z)$
\item $\varphi_{18} :\equiv \forall y\,\forall z\,(v_{128} = y \to y = z \to v_{128} = z) \to \forall z\,(v_{128} = v_{129} \to v_{129} = z \to v_{128} = z)$
\item $\varphi_{19} :\equiv \forall y\,\forall z\,(v_{125} = y \to y = z \to v_{125} = z) \to \forall z\,(v_{125} = v_{126} \to v_{126} = z \to v_{125} = z)$
\item $\varphi_{20} :\equiv \forall y\,\forall z\,(v_{122} = y \to y = z \to v_{122} = z) \to \forall z\,(v_{122} = v_{123} \to v_{123} = z \to v_{122} = z)$
\item $\varphi_{21} :\equiv \forall v_{121}\,(v_{120} \in \mathbb{R} \to v_{121} \in \mathbb{R} \to v_{120} + v_{121} = v_{121} + v_{120})$
\item $\varphi_{22} :\equiv x + (-y + y) = x + (y + -y) \to x - y + y = x + (y + -y)$
\item $\varphi_{23} :\equiv \forall z\,(v_{131} = v_{132} \to v_{132} = z \to v_{131} = z) \to v_{131} = v_{132} \to v_{132} = v_{133} \to v_{131} = v_{133}$
\item $\varphi_{24} :\equiv \forall z\,(v_{128} = v_{129} \to v_{129} = z \to v_{128} = z) \to v_{128} = v_{129} \to v_{129} = v_{130} \to v_{128} = v_{130}$
\item $\varphi_{25} :\equiv \forall z\,(v_{125} = v_{126} \to v_{126} = z \to v_{125} = z) \to v_{125} = v_{126} \to v_{126} = v_{127} \to v_{125} = v_{127}$
\item $\varphi_{26} :\equiv \forall z\,(v_{122} = v_{123} \to v_{123} = z \to v_{122} = z) \to v_{122} = v_{123} \to v_{123} = v_{124} \to v_{122} = v_{124}$
\item $\varphi_{27} :\equiv \forall y\,\forall z\,(x \in \mathbb{R} \to y \in \mathbb{R} \to z \in \mathbb{R} \to x + y + z = x + (y + z))$
\item $\varphi_{28} :\equiv \forall v_{121}\,(-y \in \mathbb{R} \to v_{121} \in \mathbb{R} \to -y + v_{121} = v_{121} + -y)$
\item $\varphi_{29} :\equiv x + (-y + y) = x + (-y + y) \to x + (-y + y) = x + (y + -y)$
\item $\varphi_{30} :\equiv (x \in \mathbb{R} \to y \in \mathbb{R} \to x \in \mathbb{R}) \to x \in \mathbb{R} \to y \in \mathbb{R} \to x + 0 = x$
\item $\varphi_{31} :\equiv x - y + y = x + (y + -y) \to x + (y + -y) = x + 0 \to x - y + y = x + 0$
\item $\varphi_{32} :\equiv y + -y = 0 \to x + (y + -y) = x + (y + -y) \to x + (y + -y) = x + 0$
\item $\varphi_{33} :\equiv x - y = x + -y \to x - y + y = x - y + y \to x - y + y = x + -y + y$
\item $\varphi_{34} :\equiv \forall v_{130}\,(x - y + y = x + (y + -y) \to x + (y + -y) = v_{130} \to x - y + y = v_{130})$
\item $\varphi_{35} :\equiv \forall v_{127}\,(x - y + y = x + (-y + y) \to x + (-y + y) = v_{127} \to x - y + y = v_{127})$
\item $\varphi_{36} :\equiv \forall v_{124}\,(x - y + y = x + -y + y \to x + -y + y = v_{124} \to x - y + y = v_{124})$
\item $\varphi_{37} :\equiv x \in \mathbb{R} \to -y \in \mathbb{R} \to y \in \mathbb{R} \to x + -y + y = x + (-y + y)$
\item $\varphi_{38} :\equiv \forall w\,(y + -y = w \to x + (y + -y) = x + (y + -y) \to x + (y + -y) = x + w)$
\item $\varphi_{39} :\equiv (x \in \mathbb{R} \to y \in \mathbb{R} \to y \in \mathbb{R}) \to x \in \mathbb{R} \to y \in \mathbb{R} \to y + -y = 0$
\item $\varphi_{40} :\equiv \forall w\,(-y + y = w \to x + (-y + y) = x + (-y + y) \to x + (-y + y) = x + w)$
\item $\varphi_{41} :\equiv (x \in \mathbb{R} \to y \in \mathbb{R} \to y \in \mathbb{R}) \to x \in \mathbb{R} \to y \in \mathbb{R} \to -y \in \mathbb{R}$
\item $\varphi_{42} :\equiv \forall z\,(x \in \mathbb{R} \to -y \in \mathbb{R} \to z \in \mathbb{R} \to x + -y + z = x + (-y + z))$
\item $\varphi_{43} :\equiv (x - y + y = x + 0 \to x + 0 = x \to x - y + y = x) \to x \in \mathbb{R} \to y \in \mathbb{R} \to x - y + y = x$
\item $\varphi_{44} :\equiv (x \in \mathbb{R} \to y \in \mathbb{R} \to y \in \mathbb{R}) \to x \in \mathbb{R} \to y \in \mathbb{R} \to -y + y = y + -y$
\item $\varphi_{45} :\equiv x - y + y = x + -y + y \to x + -y + y = x + (-y + y) \to x - y + y = x + (-y + y)$
\item $\varphi_{46} :\equiv (x \in \mathbb{R} \to y \in \mathbb{R} \to y \in \mathbb{R}) \to x \in \mathbb{R} \to y \in \mathbb{R} \to x + -y + y = x + (-y + y)$
\item $\varphi_{47} :\equiv \forall x\,\forall y\,(x \in \mathbb{R} \to y \in \mathbb{R} \to x + y = y + x) \to \forall y\,(v_{120} \in \mathbb{R} \to y \in \mathbb{R} \to v_{120} + y = y + v_{120})$
\item $\varphi_{48} :\equiv \forall y\,(v_{120} \in \mathbb{R} \to y \in \mathbb{R} \to v_{120} + y = y + v_{120}) \to v_{120} \in \mathbb{R} \to v_{121} \in \mathbb{R} \to v_{120} + v_{121} = v_{121} + v_{120}$
\item $\varphi_{49} :\equiv \forall z\,\forall w\,(z = w \to x - y + y = z + y \to x - y + y = w + y) \to \varphi_{12}$
\item $\varphi_{50} :\equiv (x - y + y = x + (-y + y) \to \varphi_{22}) \to x \in \mathbb{R} \to y \in \mathbb{R} \to x - y + y = x + (y + -y)$
\item $\varphi_{51} :\equiv (-y + y = y + -y \to \varphi_{29}) \to x \in \mathbb{R} \to y \in \mathbb{R} \to x + (-y + y) = x + (y + -y)$
\item $\varphi_{52} :\equiv x - y + y = x - y + y \to \varphi_{33} \to x \in \mathbb{R} \to y \in \mathbb{R} \to x - y + y = x + -y + y$
\item $\varphi_{53} :\equiv x + (y + -y) = x + (y + -y) \to \varphi_{32} \to x \in \mathbb{R} \to y \in \mathbb{R} \to x + (y + -y) = x + 0$
\item $\varphi_{54} :\equiv (x \in \mathbb{R} \to y \in \mathbb{R} \to x - y + y = x + 0) \to (x \in \mathbb{R} \to y \in \mathbb{R} \to x + 0 = x) \to \varphi_{43}$
\item $\varphi_{55} :\equiv (x \in \mathbb{R} \to y \in \mathbb{R} \to x + (y + -y) = x + 0) \to \varphi_{31} \to x \in \mathbb{R} \to y \in \mathbb{R} \to x - y + y = x + 0$
\item $\varphi_{56} :\equiv (x \in \mathbb{R} \to y \in \mathbb{R} \to x \in \mathbb{R}) \to (x \in \mathbb{R} \to y \in \mathbb{R} \to -y \in \mathbb{R}) \to \varphi_{46}$
\item $\varphi_{57} :\equiv (-y \in \mathbb{R} \to y \in \mathbb{R} \to -y + y = y + -y) \to (x \in \mathbb{R} \to y \in \mathbb{R} \to -y \in \mathbb{R}) \to \varphi_{44}$
\item $\varphi_{58} :\equiv (x \in \mathbb{R} \to y \in \mathbb{R} \to x + -y + y = x + (-y + y)) \to \varphi_{45} \to x \in \mathbb{R} \to y \in \mathbb{R} \to x - y + y = x + (-y + y)$
\item $\varphi_{59} :\equiv (x \in \mathbb{R} \to y \in \mathbb{R} \to x + (-y + y) = x + (y + -y)) \to \varphi_{50}$
\item $\varphi_{60} :\equiv (x \in \mathbb{R} \to y \in \mathbb{R} \to -y + y = y + -y) \to x + (-y + y) = x + (-y + y) \to \varphi_{51}$
\end{abbrevlist}
\end{fullproof}

% 自動生成 (prover/build.js): SubSelf  ノード数 25
\begin{fullproof}{[F13]}{$\forall x \in \mathbb{R}\,(x - x = 0)$}{25}\phantomsection\label{tree:SubSelf}
\noindent$\Pi_{1}$:
\begin{ptree}
\AXC{$y = z \to x = y \to x = z$\lbl{A6}}
\AXC{$(y = z \to x = y \to x = z) \to x = y \to y = z \to x = z$\lbl{A1}}
\BIC{$x = y \to y = z \to x = z$}
\UIC{$\forall z\,(x = y \to y = z \to x = z)$}
\UIC{$\forall y\,\forall z\,(x = y \to y = z \to x = z)$}
\UIC{$\forall x\,\forall y\,\forall z\,(x = y \to y = z \to x = z)$}
\AXC{$\varphi_{1}$\lbl{A2}}
\BIC{$\forall y\,\forall z\,(x - x = y \to y = z \to x - x = z)$}
\AXC{$\varphi_{3}$\lbl{A2}}
\BIC{$\forall z\,(x - x = x + -x \to x + -x = z \to x - x = z)$}
\AXC{$\varphi_{4}$\lbl{A2}}
\BIC{$x - x = x + -x \to x + -x = 0 \to x - x = 0$}
\end{ptree}
\noindent$\Pi_{2}$:
\begin{ptree}
\AXC{$\forall x\,\forall y\,(x - y = x + -y)$\lbl{Def$-$}}
\AXC{$\forall x\,\forall y\,(x - y = x + -y) \to \forall y\,(x - y = x + -y)$\lbl{A2}}
\BIC{$\forall y\,(x - y = x + -y)$}
\end{ptree}
\noindent$\Pi_{3}$:
\begin{ptree}
\AXC{$\stackrel{\Pi_{2}}{\vdots}$}\noLine
\UIC{$\forall y\,(x - y = x + -y)$}
\AXC{$\forall y\,(x - y = x + -y) \to x - x = x + -x$\lbl{A2}}
\BIC{$x - x = x + -x$}
\AXC{$x - x = x + -x \to (x \in \mathbb{R} \to x + -x = 0) \to \varphi_{2}$\lbl{A1}}
\BIC{$(x \in \mathbb{R} \to x + -x = 0) \to \varphi_{2}$}
\end{ptree}
\noindent$\Pi_{4}$:
\begin{ptree}
\AXC{$\forall x \in \mathbb{R}\,(x + -x = 0)$\lbl{(R5)}}
\AXC{$\forall x \in \mathbb{R}\,(x + -x = 0) \to x \in \mathbb{R} \to x + -x = 0$\lbl{A2}}
\BIC{$x \in \mathbb{R} \to x + -x = 0$}
\AXC{$\stackrel{\Pi_{3}}{\vdots}$}\noLine
\UIC{$(x \in \mathbb{R} \to x + -x = 0) \to \varphi_{2}$}
\BIC{$\varphi_{2}$}
\end{ptree}
\noindent 主木:
\begin{ptree}
\AXC{$\stackrel{\Pi_{1}}{\vdots}$}\noLine
\UIC{$x - x = x + -x \to x + -x = 0 \to x - x = 0$}
\AXC{$\stackrel{\Pi_{4}}{\vdots}$}\noLine
\UIC{$\varphi_{2}$}
\BIC{$x \in \mathbb{R} \to x - x = 0$}
\UIC{$\forall x \in \mathbb{R}\,(x - x = 0)$}
\end{ptree}
\begin{abbrevlist}
\item $\varphi_{1} :\equiv \forall x\,\forall y\,\forall z\,(x = y \to y = z \to x = z) \to \forall y\,\forall z\,(x - x = y \to y = z \to x - x = z)$
\item $\varphi_{2} :\equiv (x - x = x + -x \to x + -x = 0 \to x - x = 0) \to x \in \mathbb{R} \to x - x = 0$
\item $\varphi_{3} :\equiv \forall y\,\forall z\,(x - x = y \to y = z \to x - x = z) \to \forall z\,(x - x = x + -x \to x + -x = z \to x - x = z)$
\item $\varphi_{4} :\equiv \forall z\,(x - x = x + -x \to x + -x = z \to x - x = z) \to x - x = x + -x \to x + -x = 0 \to x - x = 0$
\end{abbrevlist}
\end{fullproof}

% 自動生成 (prover/build.js): SubNeg  ノード数 201
\begin{fullproof}{[F14]}{$\forall x\,\forall y\,(x \in \mathbb{R} \to y \in \mathbb{R} \to -(y - x) = x - y)$}{201}\phantomsection\label{tree:SubNeg}
\noindent$\Pi_{1}$:
\begin{ptree}
\AXC{$y = z \to x = y \to x = z$\lbl{A6}}
\AXC{$(y = z \to x = y \to x = z) \to x = y \to y = z \to x = z$\lbl{A1}}
\BIC{$x = y \to y = z \to x = z$}
\UIC{$\forall z\,(x = y \to y = z \to x = z)$}
\UIC{$\forall y\,\forall z\,(x = y \to y = z \to x = z)$}
\UIC{$\forall x\,\forall y\,\forall z\,(x = y \to y = z \to x = z)$}
\AXC{$\varphi_{6}$\lbl{A2}}
\BIC{$\forall y\,\forall z\,(v_{151} = y \to y = z \to v_{151} = z)$}
\AXC{$\varphi_{16}$\lbl{A2}}
\BIC{$\forall z\,(v_{151} = v_{152} \to v_{152} = z \to v_{151} = z)$}
\AXC{$\varphi_{23}$\lbl{A2}}
\BIC{$v_{151} = v_{152} \to v_{152} = v_{153} \to v_{151} = v_{153}$}
\UIC{$\forall v_{153}\,(v_{151} = v_{152} \to v_{152} = v_{153} \to v_{151} = v_{153})$}
\UIC{$\varphi_{2}$}
\UIC{$\forall v_{151}\,(\varphi_{2})$}
\end{ptree}
\noindent$\Pi_{2}$:
\begin{ptree}
\AXC{$x = x$\lbl{A5}}
\AXC{$x = y \to x = x \to y = x$\lbl{A6}}
\AXC{$(x = y \to x = x \to y = x) \to x = x \to x = y \to y = x$\lbl{A1}}
\BIC{$x = x \to x = y \to y = x$}
\BIC{$x = y \to y = x$}
\UIC{$\forall y\,(x = y \to y = x)$}
\UIC{$\forall x\,\forall y\,(x = y \to y = x)$}
\end{ptree}
\noindent$\Pi_{3}$:
\begin{ptree}
\AXC{$\stackrel{\Pi_{2}}{\vdots}$}\noLine
\UIC{$\forall x\,\forall y\,(x = y \to y = x)$}
\AXC{$\forall x\,\forall y\,(x = y \to y = x) \to \forall y\,(v_{140} = y \to y = v_{140})$\lbl{A2}}
\BIC{$\forall y\,(v_{140} = y \to y = v_{140})$}
\AXC{$\varphi_{1}$\lbl{A2}}
\BIC{$v_{140} = v_{141} \to v_{141} = v_{140}$}
\UIC{$\forall v_{141}\,(v_{140} = v_{141} \to v_{141} = v_{140})$}
\UIC{$\forall v_{140}\,\forall v_{141}\,(v_{140} = v_{141} \to v_{141} = v_{140})$}
\AXC{$\varphi_{29}$\lbl{A2}}
\BIC{$\forall v_{141}\,(x - y = v_{141} \to v_{141} = x - y)$}
\AXC{$\varphi_{34}$\lbl{A2}}
\BIC{$x - y = x + -y \to x + -y = x - y$}
\end{ptree}
\noindent$\Pi_{4}$:
\begin{ptree}
\AXC{$\forall x\,\forall y\,(x - y = x + -y)$\lbl{Def$-$}}
\AXC{$\forall x\,\forall y\,(x - y = x + -y) \to \forall y\,(x - y = x + -y)$\lbl{A2}}
\BIC{$\forall y\,(x - y = x + -y)$}
\end{ptree}
\noindent$\Pi_{5}$:
\begin{ptree}
\AXC{$\stackrel{\Pi_{3}}{\vdots}$}\noLine
\UIC{$x - y = x + -y \to x + -y = x - y$}
\AXC{$\stackrel{\Pi_{4}}{\vdots}$}\noLine
\UIC{$\forall y\,(x - y = x + -y)$}
\AXC{$\forall y\,(x - y = x + -y) \to x - y = x + -y$\lbl{A2}}
\BIC{$x - y = x + -y$}
\AXC{$\varphi_{45}$\lbl{A1}}
\BIC{$(x - y = x + -y \to x + -y = x - y) \to x + -y = x - y$}
\BIC{$x + -y = x - y$}
\end{ptree}
\noindent$\Pi_{6}$:
\begin{ptree}
\AXC{$y = z \to x = y \to x = z$\lbl{A6}}
\AXC{$(y = z \to x = y \to x = z) \to x = y \to y = z \to x = z$\lbl{A1}}
\BIC{$x = y \to y = z \to x = z$}
\UIC{$\forall z\,(x = y \to y = z \to x = z)$}
\UIC{$\forall y\,\forall z\,(x = y \to y = z \to x = z)$}
\UIC{$\forall x\,\forall y\,\forall z\,(x = y \to y = z \to x = z)$}
\AXC{$\varphi_{7}$\lbl{A2}}
\BIC{$\forall y\,\forall z\,(v_{148} = y \to y = z \to v_{148} = z)$}
\AXC{$\varphi_{17}$\lbl{A2}}
\BIC{$\forall z\,(v_{148} = v_{149} \to v_{149} = z \to v_{148} = z)$}
\AXC{$\varphi_{24}$\lbl{A2}}
\BIC{$v_{148} = v_{149} \to v_{149} = v_{150} \to v_{148} = v_{150}$}
\UIC{$\forall v_{150}\,(v_{148} = v_{149} \to v_{149} = v_{150} \to v_{148} = v_{150})$}
\UIC{$\varphi_{3}$}
\UIC{$\forall v_{148}\,(\varphi_{3})$}
\end{ptree}
\noindent$\Pi_{7}$:
\begin{ptree}
\AXC{$\stackrel{\Pi_{6}}{\vdots}$}\noLine
\UIC{$\forall v_{148}\,(\varphi_{3})$}
\AXC{$\forall v_{148}\,(\varphi_{3}) \to \varphi_{13}$\lbl{A2}}
\BIC{$\varphi_{13}$}
\AXC{$\varphi_{13} \to \varphi_{30}$\lbl{A2}}
\BIC{$\varphi_{30}$}
\AXC{$\varphi_{30} \to \varphi_{36}$\lbl{A2}}
\BIC{$\varphi_{36}$}
\end{ptree}
\noindent$\Pi_{8}$:
\begin{ptree}
\AXC{$\forall x \in \mathbb{R}\,(-x \in \mathbb{R})$\lbl{(R1)}}
\AXC{$\forall x \in \mathbb{R}\,(-x \in \mathbb{R}) \to y \in \mathbb{R} \to -y \in \mathbb{R}$\lbl{A2}}
\BIC{$y \in \mathbb{R} \to -y \in \mathbb{R}$}
\AXC{$(y \in \mathbb{R} \to -y \in \mathbb{R}) \to \varphi_{43}$\lbl{A1}}
\BIC{$\varphi_{43}$}
\end{ptree}
\noindent$\Pi_{9}$:
\begin{ptree}
\AXC{$\forall x\,\forall y\,(x \in \mathbb{R} \to y \in \mathbb{R} \to x + y = y + x)$\lbl{(R3)}}
\AXC{$\varphi_{50}$\lbl{A2}}
\BIC{$\forall y\,(v_{138} \in \mathbb{R} \to y \in \mathbb{R} \to v_{138} + y = y + v_{138})$}
\AXC{$\varphi_{52}$\lbl{A2}}
\BIC{$v_{138} \in \mathbb{R} \to v_{139} \in \mathbb{R} \to v_{138} + v_{139} = v_{139} + v_{138}$}
\UIC{$\varphi_{21}$}
\UIC{$\forall v_{138}\,(\varphi_{21})$}
\AXC{$\forall v_{138}\,(\varphi_{21}) \to \varphi_{31}$\lbl{A2}}
\BIC{$\varphi_{31}$}
\end{ptree}
\noindent$\Pi_{10}$:
\begin{ptree}
\AXC{$x \in \mathbb{R} \to y \in \mathbb{R} \to y \in \mathbb{R}$\lbl{A1}}
\AXC{$\stackrel{\Pi_{8}}{\vdots}$}\noLine
\UIC{$\varphi_{43}$}
\BIC{$x \in \mathbb{R} \to y \in \mathbb{R} \to -y \in \mathbb{R}$}
\AXC{$\stackrel{\Pi_{9}}{\vdots}$}\noLine
\UIC{$\varphi_{31}$}
\AXC{$\varphi_{31} \to -y \in \mathbb{R} \to x \in \mathbb{R} \to -y + x = x + -y$\lbl{A2}}
\BIC{$-y \in \mathbb{R} \to x \in \mathbb{R} \to -y + x = x + -y$}
\AXC{$\varphi_{59}$\lbl{A1}}
\BIC{$(x \in \mathbb{R} \to y \in \mathbb{R} \to -y \in \mathbb{R}) \to \varphi_{47}$}
\BIC{$\varphi_{47}$}
\end{ptree}
\noindent$\Pi_{11}$:
\begin{ptree}
\AXC{$y = z \to x = y \to x = z$\lbl{A6}}
\AXC{$(y = z \to x = y \to x = z) \to x = y \to y = z \to x = z$\lbl{A1}}
\BIC{$x = y \to y = z \to x = z$}
\UIC{$\forall z\,(x = y \to y = z \to x = z)$}
\UIC{$\forall y\,\forall z\,(x = y \to y = z \to x = z)$}
\UIC{$\forall x\,\forall y\,\forall z\,(x = y \to y = z \to x = z)$}
\AXC{$\varphi_{8}$\lbl{A2}}
\BIC{$\forall y\,\forall z\,(v_{145} = y \to y = z \to v_{145} = z)$}
\AXC{$\varphi_{18}$\lbl{A2}}
\BIC{$\forall z\,(v_{145} = v_{146} \to v_{146} = z \to v_{145} = z)$}
\AXC{$\varphi_{25}$\lbl{A2}}
\BIC{$v_{145} = v_{146} \to v_{146} = v_{147} \to v_{145} = v_{147}$}
\UIC{$\forall v_{147}\,(v_{145} = v_{146} \to v_{146} = v_{147} \to v_{145} = v_{147})$}
\UIC{$\varphi_{4}$}
\UIC{$\forall v_{145}\,(\varphi_{4})$}
\end{ptree}
\noindent$\Pi_{12}$:
\begin{ptree}
\AXC{$\forall x \in \mathbb{R}\,(-(-x) = x)$\lbl{[F5]}}
\AXC{$\forall x \in \mathbb{R}\,(-(-x) = x) \to x \in \mathbb{R} \to -(-x) = x$\lbl{A2}}
\BIC{$x \in \mathbb{R} \to -(-x) = x$}
\AXC{$(x \in \mathbb{R} \to -(-x) = x) \to \varphi_{41}$\lbl{A1}}
\BIC{$\varphi_{41}$}
\end{ptree}
\noindent$\Pi_{13}$:
\begin{ptree}
\AXC{$x = x$\lbl{A5}}
\UIC{$\forall x\,(x = x)$}
\AXC{$\forall x\,(x = x) \to -y + -(-x) = -y + -(-x)$\lbl{A2}}
\BIC{$-y + -(-x) = -y + -(-x)$}
\AXC{$x \in \mathbb{R} \to y \in \mathbb{R} \to x \in \mathbb{R}$\lbl{A1}}
\AXC{$\stackrel{\Pi_{12}}{\vdots}$}\noLine
\UIC{$\varphi_{41}$}
\BIC{$x \in \mathbb{R} \to y \in \mathbb{R} \to -(-x) = x$}
\AXC{$\varphi_{62}$\lbl{A1}}
\BIC{$-y + -(-x) = -y + -(-x) \to \varphi_{55}$}
\BIC{$\varphi_{55}$}
\end{ptree}
\noindent$\Pi_{14}$:
\begin{ptree}
\AXC{$y = z \to x = y \to x = z$\lbl{A6}}
\AXC{$(y = z \to x = y \to x = z) \to x = y \to y = z \to x = z$\lbl{A1}}
\BIC{$x = y \to y = z \to x = z$}
\UIC{$\forall z\,(x = y \to y = z \to x = z)$}
\UIC{$\forall y\,\forall z\,(x = y \to y = z \to x = z)$}
\UIC{$\forall x\,\forall y\,\forall z\,(x = y \to y = z \to x = z)$}
\AXC{$\varphi_{9}$\lbl{A2}}
\BIC{$\forall y\,\forall z\,(v_{142} = y \to y = z \to v_{142} = z)$}
\AXC{$\varphi_{19}$\lbl{A2}}
\BIC{$\forall z\,(v_{142} = v_{143} \to v_{143} = z \to v_{142} = z)$}
\AXC{$\varphi_{26}$\lbl{A2}}
\BIC{$v_{142} = v_{143} \to v_{143} = v_{144} \to v_{142} = v_{144}$}
\UIC{$\forall v_{144}\,(v_{142} = v_{143} \to v_{143} = v_{144} \to v_{142} = v_{144})$}
\UIC{$\varphi_{5}$}
\UIC{$\forall v_{142}\,(\varphi_{5})$}
\end{ptree}
\noindent$\Pi_{15}$:
\begin{ptree}
\AXC{$\forall x \in \mathbb{R}\,(-x \in \mathbb{R})$\lbl{(R1)}}
\AXC{$\forall x \in \mathbb{R}\,(-x \in \mathbb{R}) \to x \in \mathbb{R} \to -x \in \mathbb{R}$\lbl{A2}}
\BIC{$x \in \mathbb{R} \to -x \in \mathbb{R}$}
\AXC{$(x \in \mathbb{R} \to -x \in \mathbb{R}) \to \varphi_{44}$\lbl{A1}}
\BIC{$\varphi_{44}$}
\end{ptree}
\noindent$\Pi_{16}$:
\begin{ptree}
\AXC{$\forall x\,\forall y\,(x \in \mathbb{R} \to y \in \mathbb{R} \to -(x + y) = -x + -y)$\lbl{[F7]}}
\AXC{$\forall x\,\forall y\,(x \in \mathbb{R} \to y \in \mathbb{R} \to -(x + y) = -x + -y) \to \varphi_{27}$\lbl{A2}}
\BIC{$\varphi_{27}$}
\end{ptree}
\noindent$\Pi_{17}$:
\begin{ptree}
\AXC{$\stackrel{\Pi_{16}}{\vdots}$}\noLine
\UIC{$\varphi_{27}$}
\AXC{$\varphi_{27} \to \varphi_{33}$\lbl{A2}}
\BIC{$\varphi_{33}$}
\UIC{$\forall v_{137}\,(\varphi_{33})$}
\UIC{$\forall v_{136}\,\forall v_{137}\,(\varphi_{33})$}
\AXC{$\forall v_{136}\,\forall v_{137}\,(\varphi_{33}) \to \varphi_{32}$\lbl{A2}}
\BIC{$\varphi_{32}$}
\AXC{$\varphi_{32} \to y \in \mathbb{R} \to -x \in \mathbb{R} \to -(y + -x) = -y + -(-x)$\lbl{A2}}
\BIC{$y \in \mathbb{R} \to -x \in \mathbb{R} \to -(y + -x) = -y + -(-x)$}
\end{ptree}
\noindent$\Pi_{18}$:
\begin{ptree}
\AXC{$x \in \mathbb{R} \to y \in \mathbb{R} \to y \in \mathbb{R}$\lbl{A1}}
\AXC{$\stackrel{\Pi_{17}}{\vdots}$}\noLine
\UIC{$y \in \mathbb{R} \to -x \in \mathbb{R} \to -(y + -x) = -y + -(-x)$}
\AXC{$\varphi_{61}$\lbl{A1}}
\BIC{$(x \in \mathbb{R} \to y \in \mathbb{R} \to y \in \mathbb{R}) \to \varphi_{51}$}
\BIC{$\varphi_{51}$}
\end{ptree}
\noindent$\Pi_{19}$:
\begin{ptree}
\AXC{$z = w \to -(y - x) = -z \to -(y - x) = -w$\lbl{A6}}
\UIC{$\forall w\,(z = w \to -(y - x) = -z \to -(y - x) = -w)$}
\UIC{$\forall z\,\forall w\,(z = w \to -(y - x) = -z \to -(y - x) = -w)$}
\AXC{$\varphi_{53}$\lbl{A2}}
\BIC{$\varphi_{11}$}
\AXC{$\varphi_{11} \to \varphi_{38}$\lbl{A2}}
\BIC{$\varphi_{38}$}
\end{ptree}
\noindent$\Pi_{20}$:
\begin{ptree}
\AXC{$\forall x\,\forall y\,(x - y = x + -y)$\lbl{Def$-$}}
\AXC{$\varphi_{10}$\lbl{A2}}
\BIC{$\forall y\,(v_{134} - y = v_{134} + -y)$}
\AXC{$\varphi_{20}$\lbl{A2}}
\BIC{$v_{134} - v_{135} = v_{134} + -v_{135}$}
\UIC{$\forall v_{135}\,(v_{134} - v_{135} = v_{134} + -v_{135})$}
\UIC{$\forall v_{134}\,\forall v_{135}\,(v_{134} - v_{135} = v_{134} + -v_{135})$}
\AXC{$\varphi_{39}$\lbl{A2}}
\BIC{$\forall v_{135}\,(y - v_{135} = y + -v_{135})$}
\AXC{$\forall v_{135}\,(y - v_{135} = y + -v_{135}) \to y - x = y + -x$\lbl{A2}}
\BIC{$y - x = y + -x$}
\end{ptree}
\noindent$\Pi_{21}$:
\begin{ptree}
\AXC{$x = x$\lbl{A5}}
\UIC{$\forall x\,(x = x)$}
\AXC{$\forall x\,(x = x) \to -(y - x) = -(y - x)$\lbl{A2}}
\BIC{$-(y - x) = -(y - x)$}
\AXC{$\stackrel{\Pi_{20}}{\vdots}$}\noLine
\UIC{$y - x = y + -x$}
\AXC{$y - x = y + -x \to \varphi_{56}$\lbl{A1}}
\BIC{$\varphi_{56}$}
\BIC{$\varphi_{38} \to x \in \mathbb{R} \to y \in \mathbb{R} \to -(y - x) = -(y + -x)$}
\end{ptree}
\noindent$\Pi_{22}$:
\begin{ptree}
\AXC{$\stackrel{\Pi_{19}}{\vdots}$}\noLine
\UIC{$\varphi_{38}$}
\AXC{$\stackrel{\Pi_{21}}{\vdots}$}\noLine
\UIC{$\varphi_{38} \to x \in \mathbb{R} \to y \in \mathbb{R} \to -(y - x) = -(y + -x)$}
\BIC{$x \in \mathbb{R} \to y \in \mathbb{R} \to -(y - x) = -(y + -x)$}
\AXC{$\varphi_{63}$\lbl{A1}}
\BIC{$(x \in \mathbb{R} \to y \in \mathbb{R} \to -(y + -x) = -y + -(-x)) \to \varphi_{57}$}
\end{ptree}
\noindent$\Pi_{23}$:
\begin{ptree}
\AXC{$x \in \mathbb{R} \to y \in \mathbb{R} \to x \in \mathbb{R}$\lbl{A1}}
\AXC{$\stackrel{\Pi_{15}}{\vdots}$}\noLine
\UIC{$\varphi_{44}$}
\BIC{$x \in \mathbb{R} \to y \in \mathbb{R} \to -x \in \mathbb{R}$}
\AXC{$\stackrel{\Pi_{18}}{\vdots}$}\noLine
\UIC{$\varphi_{51}$}
\BIC{$x \in \mathbb{R} \to y \in \mathbb{R} \to -(y + -x) = -y + -(-x)$}
\AXC{$\stackrel{\Pi_{22}}{\vdots}$}\noLine
\UIC{$(x \in \mathbb{R} \to y \in \mathbb{R} \to -(y + -x) = -y + -(-x)) \to \varphi_{57}$}
\BIC{$\varphi_{57}$}
\end{ptree}
\noindent$\Pi_{24}$:
\begin{ptree}
\AXC{$\stackrel{\Pi_{14}}{\vdots}$}\noLine
\UIC{$\forall v_{142}\,(\varphi_{5})$}
\AXC{$\forall v_{142}\,(\varphi_{5}) \to \varphi_{15}$\lbl{A2}}
\BIC{$\varphi_{15}$}
\AXC{$\varphi_{15} \to \varphi_{40}$\lbl{A2}}
\BIC{$\varphi_{40}$}
\AXC{$\varphi_{40} \to -(y - x) = -(y + -x) \to \varphi_{35}$\lbl{A2}}
\BIC{$-(y - x) = -(y + -x) \to \varphi_{35}$}
\AXC{$\stackrel{\Pi_{23}}{\vdots}$}\noLine
\UIC{$\varphi_{57}$}
\BIC{$x \in \mathbb{R} \to y \in \mathbb{R} \to -(y - x) = -y + -(-x)$}
\end{ptree}
\noindent$\Pi_{25}$:
\begin{ptree}
\AXC{$\stackrel{\Pi_{24}}{\vdots}$}\noLine
\UIC{$x \in \mathbb{R} \to y \in \mathbb{R} \to -(y - x) = -y + -(-x)$}
\AXC{$(x \in \mathbb{R} \to y \in \mathbb{R} \to -(y - x) = -y + -(-x)) \to \varphi_{60}$\lbl{A1}}
\BIC{$\varphi_{60}$}
\end{ptree}
\noindent$\Pi_{26}$:
\begin{ptree}
\AXC{$\varphi_{37}$\lbl{A6}}
\UIC{$\forall w\,(\varphi_{37})$}
\UIC{$\forall z\,\forall w\,(\varphi_{37})$}
\AXC{$\forall z\,\forall w\,(\varphi_{37}) \to \varphi_{49}$\lbl{A2}}
\BIC{$\varphi_{49}$}
\AXC{$\varphi_{49} \to -(-x) = x \to \varphi_{42}$\lbl{A2}}
\BIC{$-(-x) = x \to \varphi_{42}$}
\AXC{$\stackrel{\Pi_{13}}{\vdots}$}\noLine
\UIC{$\varphi_{55}$}
\BIC{$x \in \mathbb{R} \to y \in \mathbb{R} \to -y + -(-x) = -y + x$}
\AXC{$\stackrel{\Pi_{25}}{\vdots}$}\noLine
\UIC{$\varphi_{60}$}
\BIC{$\varphi_{48} \to x \in \mathbb{R} \to y \in \mathbb{R} \to -(y - x) = -y + x$}
\end{ptree}
\noindent$\Pi_{27}$:
\begin{ptree}
\AXC{$\stackrel{\Pi_{11}}{\vdots}$}\noLine
\UIC{$\forall v_{145}\,(\varphi_{4})$}
\AXC{$\forall v_{145}\,(\varphi_{4}) \to \varphi_{14}$\lbl{A2}}
\BIC{$\varphi_{14}$}
\AXC{$\varphi_{14} \to \varphi_{46}$\lbl{A2}}
\BIC{$\varphi_{46}$}
\AXC{$\varphi_{46} \to \varphi_{48}$\lbl{A2}}
\BIC{$\varphi_{48}$}
\AXC{$\stackrel{\Pi_{26}}{\vdots}$}\noLine
\UIC{$\varphi_{48} \to x \in \mathbb{R} \to y \in \mathbb{R} \to -(y - x) = -y + x$}
\BIC{$x \in \mathbb{R} \to y \in \mathbb{R} \to -(y - x) = -y + x$}
\end{ptree}
\noindent$\Pi_{28}$:
\begin{ptree}
\AXC{$\stackrel{\Pi_{27}}{\vdots}$}\noLine
\UIC{$x \in \mathbb{R} \to y \in \mathbb{R} \to -(y - x) = -y + x$}
\AXC{$(x \in \mathbb{R} \to y \in \mathbb{R} \to -(y - x) = -y + x) \to \varphi_{58}$\lbl{A1}}
\BIC{$\varphi_{58}$}
\end{ptree}
\noindent$\Pi_{29}$:
\begin{ptree}
\AXC{$\stackrel{\Pi_{7}}{\vdots}$}\noLine
\UIC{$\varphi_{36}$}
\AXC{$x \in \mathbb{R} \to y \in \mathbb{R} \to x \in \mathbb{R}$\lbl{A1}}
\AXC{$\stackrel{\Pi_{10}}{\vdots}$}\noLine
\UIC{$\varphi_{47}$}
\BIC{$x \in \mathbb{R} \to y \in \mathbb{R} \to -y + x = x + -y$}
\AXC{$\stackrel{\Pi_{28}}{\vdots}$}\noLine
\UIC{$\varphi_{58}$}
\BIC{$\varphi_{36} \to x \in \mathbb{R} \to y \in \mathbb{R} \to -(y - x) = x + -y$}
\BIC{$x \in \mathbb{R} \to y \in \mathbb{R} \to -(y - x) = x + -y$}
\AXC{$(x \in \mathbb{R} \to y \in \mathbb{R} \to -(y - x) = x + -y) \to \varphi_{54}$\lbl{A1}}
\BIC{$\varphi_{54}$}
\end{ptree}
\noindent 主木:
\begin{ptree}
\AXC{$\stackrel{\Pi_{1}}{\vdots}$}\noLine
\UIC{$\forall v_{151}\,(\varphi_{2})$}
\AXC{$\forall v_{151}\,(\varphi_{2}) \to \varphi_{12}$\lbl{A2}}
\BIC{$\varphi_{12}$}
\AXC{$\varphi_{12} \to \varphi_{28}$\lbl{A2}}
\BIC{$\varphi_{28}$}
\AXC{$\varphi_{28} \to \varphi_{22}$\lbl{A2}}
\BIC{$\varphi_{22}$}
\AXC{$\stackrel{\Pi_{5}}{\vdots}$}\noLine
\UIC{$x + -y = x - y$}
\AXC{$\stackrel{\Pi_{29}}{\vdots}$}\noLine
\UIC{$\varphi_{54}$}
\BIC{$\varphi_{22} \to x \in \mathbb{R} \to y \in \mathbb{R} \to -(y - x) = x - y$}
\BIC{$x \in \mathbb{R} \to y \in \mathbb{R} \to -(y - x) = x - y$}
\UIC{$\forall y\,(x \in \mathbb{R} \to y \in \mathbb{R} \to -(y - x) = x - y)$}
\UIC{$\forall x\,\forall y\,(x \in \mathbb{R} \to y \in \mathbb{R} \to -(y - x) = x - y)$}
\end{ptree}
\begin{abbrevlist}
\item $\varphi_{1} :\equiv \forall y\,(v_{140} = y \to y = v_{140}) \to v_{140} = v_{141} \to v_{141} = v_{140}$
\item $\varphi_{2} :\equiv \forall v_{152}\,\forall v_{153}\,(v_{151} = v_{152} \to v_{152} = v_{153} \to v_{151} = v_{153})$
\item $\varphi_{3} :\equiv \forall v_{149}\,\forall v_{150}\,(v_{148} = v_{149} \to v_{149} = v_{150} \to v_{148} = v_{150})$
\item $\varphi_{4} :\equiv \forall v_{146}\,\forall v_{147}\,(v_{145} = v_{146} \to v_{146} = v_{147} \to v_{145} = v_{147})$
\item $\varphi_{5} :\equiv \forall v_{143}\,\forall v_{144}\,(v_{142} = v_{143} \to v_{143} = v_{144} \to v_{142} = v_{144})$
\item $\varphi_{6} :\equiv \forall x\,\forall y\,\forall z\,(x = y \to y = z \to x = z) \to \forall y\,\forall z\,(v_{151} = y \to y = z \to v_{151} = z)$
\item $\varphi_{7} :\equiv \forall x\,\forall y\,\forall z\,(x = y \to y = z \to x = z) \to \forall y\,\forall z\,(v_{148} = y \to y = z \to v_{148} = z)$
\item $\varphi_{8} :\equiv \forall x\,\forall y\,\forall z\,(x = y \to y = z \to x = z) \to \forall y\,\forall z\,(v_{145} = y \to y = z \to v_{145} = z)$
\item $\varphi_{9} :\equiv \forall x\,\forall y\,\forall z\,(x = y \to y = z \to x = z) \to \forall y\,\forall z\,(v_{142} = y \to y = z \to v_{142} = z)$
\item $\varphi_{10} :\equiv \forall x\,\forall y\,(x - y = x + -y) \to \forall y\,(v_{134} - y = v_{134} + -y)$
\item $\varphi_{11} :\equiv \forall w\,(y - x = w \to -(y - x) = -(y - x) \to -(y - x) = -w)$
\item $\varphi_{12} :\equiv \forall v_{152}\,\forall v_{153}\,(-(y - x) = v_{152} \to v_{152} = v_{153} \to -(y - x) = v_{153})$
\item $\varphi_{13} :\equiv \forall v_{149}\,\forall v_{150}\,(-(y - x) = v_{149} \to v_{149} = v_{150} \to -(y - x) = v_{150})$
\item $\varphi_{14} :\equiv \forall v_{146}\,\forall v_{147}\,(-(y - x) = v_{146} \to v_{146} = v_{147} \to -(y - x) = v_{147})$
\item $\varphi_{15} :\equiv \forall v_{143}\,\forall v_{144}\,(-(y - x) = v_{143} \to v_{143} = v_{144} \to -(y - x) = v_{144})$
\item $\varphi_{16} :\equiv \forall y\,\forall z\,(v_{151} = y \to y = z \to v_{151} = z) \to \forall z\,(v_{151} = v_{152} \to v_{152} = z \to v_{151} = z)$
\item $\varphi_{17} :\equiv \forall y\,\forall z\,(v_{148} = y \to y = z \to v_{148} = z) \to \forall z\,(v_{148} = v_{149} \to v_{149} = z \to v_{148} = z)$
\item $\varphi_{18} :\equiv \forall y\,\forall z\,(v_{145} = y \to y = z \to v_{145} = z) \to \forall z\,(v_{145} = v_{146} \to v_{146} = z \to v_{145} = z)$
\item $\varphi_{19} :\equiv \forall y\,\forall z\,(v_{142} = y \to y = z \to v_{142} = z) \to \forall z\,(v_{142} = v_{143} \to v_{143} = z \to v_{142} = z)$
\item $\varphi_{20} :\equiv \forall y\,(v_{134} - y = v_{134} + -y) \to v_{134} - v_{135} = v_{134} + -v_{135}$
\item $\varphi_{21} :\equiv \forall v_{139}\,(v_{138} \in \mathbb{R} \to v_{139} \in \mathbb{R} \to v_{138} + v_{139} = v_{139} + v_{138})$
\item $\varphi_{22} :\equiv -(y - x) = x + -y \to x + -y = x - y \to -(y - x) = x - y$
\item $\varphi_{23} :\equiv \forall z\,(v_{151} = v_{152} \to v_{152} = z \to v_{151} = z) \to v_{151} = v_{152} \to v_{152} = v_{153} \to v_{151} = v_{153}$
\item $\varphi_{24} :\equiv \forall z\,(v_{148} = v_{149} \to v_{149} = z \to v_{148} = z) \to v_{148} = v_{149} \to v_{149} = v_{150} \to v_{148} = v_{150}$
\item $\varphi_{25} :\equiv \forall z\,(v_{145} = v_{146} \to v_{146} = z \to v_{145} = z) \to v_{145} = v_{146} \to v_{146} = v_{147} \to v_{145} = v_{147}$
\item $\varphi_{26} :\equiv \forall z\,(v_{142} = v_{143} \to v_{143} = z \to v_{142} = z) \to v_{142} = v_{143} \to v_{143} = v_{144} \to v_{142} = v_{144}$
\item $\varphi_{27} :\equiv \forall y\,(v_{136} \in \mathbb{R} \to y \in \mathbb{R} \to -(v_{136} + y) = -v_{136} + -y)$
\item $\varphi_{28} :\equiv \forall v_{153}\,(-(y - x) = x + -y \to x + -y = v_{153} \to -(y - x) = v_{153})$
\item $\varphi_{29} :\equiv \forall v_{140}\,\forall v_{141}\,(v_{140} = v_{141} \to v_{141} = v_{140}) \to \forall v_{141}\,(x - y = v_{141} \to v_{141} = x - y)$
\item $\varphi_{30} :\equiv \forall v_{150}\,(-(y - x) = -y + x \to -y + x = v_{150} \to -(y - x) = v_{150})$
\item $\varphi_{31} :\equiv \forall v_{139}\,(-y \in \mathbb{R} \to v_{139} \in \mathbb{R} \to -y + v_{139} = v_{139} + -y)$
\item $\varphi_{32} :\equiv \forall v_{137}\,(y \in \mathbb{R} \to v_{137} \in \mathbb{R} \to -(y + v_{137}) = -y + -v_{137})$
\item $\varphi_{33} :\equiv v_{136} \in \mathbb{R} \to v_{137} \in \mathbb{R} \to -(v_{136} + v_{137}) = -v_{136} + -v_{137}$
\item $\varphi_{34} :\equiv \forall v_{141}\,(x - y = v_{141} \to v_{141} = x - y) \to x - y = x + -y \to x + -y = x - y$
\item $\varphi_{35} :\equiv -(y + -x) = -y + -(-x) \to -(y - x) = -y + -(-x)$
\item $\varphi_{36} :\equiv -(y - x) = -y + x \to -y + x = x + -y \to -(y - x) = x + -y$
\item $\varphi_{37} :\equiv z = w \to -y + -(-x) = -y + z \to -y + -(-x) = -y + w$
\item $\varphi_{38} :\equiv y - x = y + -x \to -(y - x) = -(y - x) \to -(y - x) = -(y + -x)$
\item $\varphi_{39} :\equiv \forall v_{134}\,\forall v_{135}\,(v_{134} - v_{135} = v_{134} + -v_{135}) \to \forall v_{135}\,(y - v_{135} = y + -v_{135})$
\item $\varphi_{40} :\equiv \forall v_{144}\,(-(y - x) = -(y + -x) \to -(y + -x) = v_{144} \to -(y - x) = v_{144})$
\item $\varphi_{41} :\equiv (x \in \mathbb{R} \to y \in \mathbb{R} \to x \in \mathbb{R}) \to x \in \mathbb{R} \to y \in \mathbb{R} \to -(-x) = x$
\item $\varphi_{42} :\equiv -y + -(-x) = -y + -(-x) \to -y + -(-x) = -y + x$
\item $\varphi_{43} :\equiv (x \in \mathbb{R} \to y \in \mathbb{R} \to y \in \mathbb{R}) \to x \in \mathbb{R} \to y \in \mathbb{R} \to -y \in \mathbb{R}$
\item $\varphi_{44} :\equiv (x \in \mathbb{R} \to y \in \mathbb{R} \to x \in \mathbb{R}) \to x \in \mathbb{R} \to y \in \mathbb{R} \to -x \in \mathbb{R}$
\item $\varphi_{45} :\equiv x - y = x + -y \to (x - y = x + -y \to x + -y = x - y) \to x + -y = x - y$
\item $\varphi_{46} :\equiv \forall v_{147}\,(-(y - x) = -y + -(-x) \to -y + -(-x) = v_{147} \to -(y - x) = v_{147})$
\item $\varphi_{47} :\equiv (x \in \mathbb{R} \to y \in \mathbb{R} \to x \in \mathbb{R}) \to x \in \mathbb{R} \to y \in \mathbb{R} \to -y + x = x + -y$
\item $\varphi_{48} :\equiv -(y - x) = -y + -(-x) \to -y + -(-x) = -y + x \to -(y - x) = -y + x$
\item $\varphi_{49} :\equiv \forall w\,(-(-x) = w \to -y + -(-x) = -y + -(-x) \to -y + -(-x) = -y + w)$
\item $\varphi_{50} :\equiv \forall x\,\forall y\,(x \in \mathbb{R} \to y \in \mathbb{R} \to x + y = y + x) \to \forall y\,(v_{138} \in \mathbb{R} \to y \in \mathbb{R} \to v_{138} + y = y + v_{138})$
\item $\varphi_{51} :\equiv (x \in \mathbb{R} \to y \in \mathbb{R} \to -x \in \mathbb{R}) \to x \in \mathbb{R} \to y \in \mathbb{R} \to -(y + -x) = -y + -(-x)$
\item $\varphi_{52} :\equiv \forall y\,(v_{138} \in \mathbb{R} \to y \in \mathbb{R} \to v_{138} + y = y + v_{138}) \to v_{138} \in \mathbb{R} \to v_{139} \in \mathbb{R} \to v_{138} + v_{139} = v_{139} + v_{138}$
\item $\varphi_{53} :\equiv \forall z\,\forall w\,(z = w \to -(y - x) = -z \to -(y - x) = -w) \to \varphi_{11}$
\item $\varphi_{54} :\equiv x + -y = x - y \to \varphi_{22} \to x \in \mathbb{R} \to y \in \mathbb{R} \to -(y - x) = x - y$
\item $\varphi_{55} :\equiv (-(-x) = x \to \varphi_{42}) \to x \in \mathbb{R} \to y \in \mathbb{R} \to -y + -(-x) = -y + x$
\item $\varphi_{56} :\equiv -(y - x) = -(y - x) \to \varphi_{38} \to x \in \mathbb{R} \to y \in \mathbb{R} \to -(y - x) = -(y + -x)$
\item $\varphi_{57} :\equiv (-(y - x) = -(y + -x) \to \varphi_{35}) \to x \in \mathbb{R} \to y \in \mathbb{R} \to -(y - x) = -y + -(-x)$
\item $\varphi_{58} :\equiv (x \in \mathbb{R} \to y \in \mathbb{R} \to -y + x = x + -y) \to \varphi_{36} \to x \in \mathbb{R} \to y \in \mathbb{R} \to -(y - x) = x + -y$
\item $\varphi_{59} :\equiv (-y \in \mathbb{R} \to x \in \mathbb{R} \to -y + x = x + -y) \to (x \in \mathbb{R} \to y \in \mathbb{R} \to -y \in \mathbb{R}) \to \varphi_{47}$
\item $\varphi_{60} :\equiv (x \in \mathbb{R} \to y \in \mathbb{R} \to -y + -(-x) = -y + x) \to \varphi_{48} \to x \in \mathbb{R} \to y \in \mathbb{R} \to -(y - x) = -y + x$
\item $\varphi_{61} :\equiv (y \in \mathbb{R} \to -x \in \mathbb{R} \to -(y + -x) = -y + -(-x)) \to (x \in \mathbb{R} \to y \in \mathbb{R} \to y \in \mathbb{R}) \to \varphi_{51}$
\item $\varphi_{62} :\equiv (x \in \mathbb{R} \to y \in \mathbb{R} \to -(-x) = x) \to -y + -(-x) = -y + -(-x) \to \varphi_{55}$
\item $\varphi_{63} :\equiv (x \in \mathbb{R} \to y \in \mathbb{R} \to -(y - x) = -(y + -x)) \to (x \in \mathbb{R} \to y \in \mathbb{R} \to -(y + -x) = -y + -(-x)) \to \varphi_{57}$
\end{abbrevlist}
\end{fullproof}

% 自動生成 (prover/build.js): SubAddAdd  ノード数 450
\begin{fullproof}{[F15]}{$\forall x\,\forall y\,\forall a\,\forall b\,(x \in \mathbb{R} \to y \in \mathbb{R} \to a \in \mathbb{R} \to b \in \mathbb{R} \to x + y - (a + b) = x - a + (y - b))$}{450}\phantomsection\label{tree:SubAddAdd}
\noindent$\Pi_{1}$:
\begin{ptree}
\AXC{$y = z \to x = y \to x = z$\lbl{A6}}
\AXC{$(y = z \to x = y \to x = z) \to x = y \to y = z \to x = z$\lbl{A1}}
\BIC{$x = y \to y = z \to x = z$}
\UIC{$\forall z\,(x = y \to y = z \to x = z)$}
\UIC{$\forall y\,\forall z\,(x = y \to y = z \to x = z)$}
\UIC{$\forall x\,\forall y\,\forall z\,(x = y \to y = z \to x = z)$}
\AXC{$\varphi_{12}$\lbl{A2}}
\BIC{$\forall y\,\forall z\,(v_{192} = y \to y = z \to v_{192} = z)$}
\AXC{$\varphi_{23}$\lbl{A2}}
\BIC{$\forall z\,(v_{192} = v_{193} \to v_{193} = z \to v_{192} = z)$}
\AXC{$\varphi_{41}$\lbl{A2}}
\BIC{$v_{192} = v_{193} \to v_{193} = v_{194} \to v_{192} = v_{194}$}
\UIC{$\forall v_{194}\,(v_{192} = v_{193} \to v_{193} = v_{194} \to v_{192} = v_{194})$}
\UIC{$\varphi_{4}$}
\UIC{$\forall v_{192}\,(\varphi_{4})$}
\end{ptree}
\noindent$\Pi_{2}$:
\begin{ptree}
\AXC{$z = w \to \varphi_{58}$\lbl{A6}}
\UIC{$\forall w\,(z = w \to \varphi_{58})$}
\UIC{$\forall z\,\forall w\,(z = w \to \varphi_{58})$}
\AXC{$\forall z\,\forall w\,(z = w \to \varphi_{58}) \to \varphi_{111}$\lbl{A2}}
\BIC{$\varphi_{111}$}
\AXC{$\varphi_{111} \to y + -b = y - b \to \varphi_{96}$\lbl{A2}}
\BIC{$y + -b = y - b \to \varphi_{96}$}
\end{ptree}
\noindent$\Pi_{3}$:
\begin{ptree}
\AXC{$x = x$\lbl{A5}}
\AXC{$x = y \to x = x \to y = x$\lbl{A6}}
\AXC{$(x = y \to x = x \to y = x) \to x = x \to x = y \to y = x$\lbl{A1}}
\BIC{$x = x \to x = y \to y = x$}
\BIC{$x = y \to y = x$}
\UIC{$\forall y\,(x = y \to y = x)$}
\UIC{$\forall x\,\forall y\,(x = y \to y = x)$}
\end{ptree}
\noindent$\Pi_{4}$:
\begin{ptree}
\AXC{$\stackrel{\Pi_{3}}{\vdots}$}\noLine
\UIC{$\forall x\,\forall y\,(x = y \to y = x)$}
\AXC{$\forall x\,\forall y\,(x = y \to y = x) \to \forall y\,(v_{169} = y \to y = v_{169})$\lbl{A2}}
\BIC{$\forall y\,(v_{169} = y \to y = v_{169})$}
\AXC{$\varphi_{1}$\lbl{A2}}
\BIC{$v_{169} = v_{170} \to v_{170} = v_{169}$}
\UIC{$\forall v_{170}\,(v_{169} = v_{170} \to v_{170} = v_{169})$}
\UIC{$\forall v_{169}\,\forall v_{170}\,(v_{169} = v_{170} \to v_{170} = v_{169})$}
\AXC{$\varphi_{49}$\lbl{A2}}
\BIC{$\forall v_{170}\,(y - b = v_{170} \to v_{170} = y - b)$}
\AXC{$\varphi_{70}$\lbl{A2}}
\BIC{$y - b = y + -b \to y + -b = y - b$}
\end{ptree}
\noindent$\Pi_{5}$:
\begin{ptree}
\AXC{$\forall x\,\forall y\,(x - y = x + -y)$\lbl{Def$-$}}
\AXC{$\varphi_{13}$\lbl{A2}}
\BIC{$\forall y\,(v_{163} - y = v_{163} + -y)$}
\AXC{$\varphi_{24}$\lbl{A2}}
\BIC{$v_{163} - v_{164} = v_{163} + -v_{164}$}
\UIC{$\forall v_{164}\,(v_{163} - v_{164} = v_{163} + -v_{164})$}
\UIC{$\forall v_{163}\,\forall v_{164}\,(v_{163} - v_{164} = v_{163} + -v_{164})$}
\AXC{$\varphi_{74}$\lbl{A2}}
\BIC{$\forall v_{164}\,(y - v_{164} = y + -v_{164})$}
\AXC{$\forall v_{164}\,(y - v_{164} = y + -v_{164}) \to y - b = y + -b$\lbl{A2}}
\BIC{$y - b = y + -b$}
\end{ptree}
\noindent$\Pi_{6}$:
\begin{ptree}
\AXC{$\stackrel{\Pi_{4}}{\vdots}$}\noLine
\UIC{$y - b = y + -b \to y + -b = y - b$}
\AXC{$\stackrel{\Pi_{5}}{\vdots}$}\noLine
\UIC{$y - b = y + -b$}
\AXC{$\varphi_{88}$\lbl{A1}}
\BIC{$(y - b = y + -b \to y + -b = y - b) \to y + -b = y - b$}
\BIC{$y + -b = y - b$}
\AXC{$y + -b = y - b \to \varphi_{164}$\lbl{A1}}
\BIC{$\varphi_{164}$}
\end{ptree}
\noindent$\Pi_{7}$:
\begin{ptree}
\AXC{$\stackrel{\Pi_{2}}{\vdots}$}\noLine
\UIC{$y + -b = y - b \to \varphi_{96}$}
\AXC{$x = x$\lbl{A5}}
\UIC{$\forall x\,(x = x)$}
\AXC{$\forall x\,(x = x) \to x - a + (y + -b) = x - a + (y + -b)$\lbl{A2}}
\BIC{$x - a + (y + -b) = x - a + (y + -b)$}
\AXC{$\stackrel{\Pi_{6}}{\vdots}$}\noLine
\UIC{$\varphi_{164}$}
\BIC{$(y + -b = y - b \to \varphi_{96}) \to x \in \mathbb{R} \to \varphi_{77}$}
\BIC{$x \in \mathbb{R} \to \varphi_{77}$}
\end{ptree}
\noindent$\Pi_{8}$:
\begin{ptree}
\AXC{$y = z \to x = y \to x = z$\lbl{A6}}
\AXC{$(y = z \to x = y \to x = z) \to x = y \to y = z \to x = z$\lbl{A1}}
\BIC{$x = y \to y = z \to x = z$}
\UIC{$\forall z\,(x = y \to y = z \to x = z)$}
\UIC{$\forall y\,\forall z\,(x = y \to y = z \to x = z)$}
\UIC{$\forall x\,\forall y\,\forall z\,(x = y \to y = z \to x = z)$}
\AXC{$\varphi_{14}$\lbl{A2}}
\BIC{$\forall y\,\forall z\,(v_{189} = y \to y = z \to v_{189} = z)$}
\AXC{$\varphi_{25}$\lbl{A2}}
\BIC{$\forall z\,(v_{189} = v_{190} \to v_{190} = z \to v_{189} = z)$}
\AXC{$\varphi_{42}$\lbl{A2}}
\BIC{$v_{189} = v_{190} \to v_{190} = v_{191} \to v_{189} = v_{191}$}
\UIC{$\forall v_{191}\,(v_{189} = v_{190} \to v_{190} = v_{191} \to v_{189} = v_{191})$}
\UIC{$\varphi_{5}$}
\UIC{$\forall v_{189}\,(\varphi_{5})$}
\end{ptree}
\noindent$\Pi_{9}$:
\begin{ptree}
\AXC{$\stackrel{\Pi_{8}}{\vdots}$}\noLine
\UIC{$\forall v_{189}\,(\varphi_{5})$}
\AXC{$\forall v_{189}\,(\varphi_{5}) \to \varphi_{51}$\lbl{A2}}
\BIC{$\varphi_{51}$}
\AXC{$\varphi_{51} \to \varphi_{123}$\lbl{A2}}
\BIC{$\varphi_{123}$}
\AXC{$\varphi_{123} \to x + y - (a + b) = x + -a + (y + -b) \to \varphi_{99}$\lbl{A2}}
\BIC{$x + y - (a + b) = x + -a + (y + -b) \to \varphi_{99}$}
\end{ptree}
\noindent$\Pi_{10}$:
\begin{ptree}
\AXC{$z = w \to \varphi_{100}$\lbl{A6}}
\UIC{$\forall w\,(z = w \to \varphi_{100})$}
\UIC{$\forall z\,\forall w\,(z = w \to \varphi_{100})$}
\AXC{$\forall z\,\forall w\,(z = w \to \varphi_{100}) \to \varphi_{137}$\lbl{A2}}
\BIC{$\varphi_{137}$}
\AXC{$\varphi_{137} \to x + -a = x - a \to \varphi_{121}$\lbl{A2}}
\BIC{$x + -a = x - a \to \varphi_{121}$}
\end{ptree}
\noindent$\Pi_{11}$:
\begin{ptree}
\AXC{$\stackrel{\Pi_{3}}{\vdots}$}\noLine
\UIC{$\forall x\,\forall y\,(x = y \to y = x)$}
\AXC{$\forall x\,\forall y\,(x = y \to y = x) \to \forall y\,(x - a = y \to y = x - a)$\lbl{A2}}
\BIC{$\forall y\,(x - a = y \to y = x - a)$}
\AXC{$\varphi_{38}$\lbl{A2}}
\BIC{$x - a = x + -a \to x + -a = x - a$}
\end{ptree}
\noindent$\Pi_{12}$:
\begin{ptree}
\AXC{$\forall x\,\forall y\,(x - y = x + -y)$\lbl{Def$-$}}
\AXC{$\forall x\,\forall y\,(x - y = x + -y) \to \forall y\,(x - y = x + -y)$\lbl{A2}}
\BIC{$\forall y\,(x - y = x + -y)$}
\end{ptree}
\noindent$\Pi_{13}$:
\begin{ptree}
\AXC{$\stackrel{\Pi_{11}}{\vdots}$}\noLine
\UIC{$x - a = x + -a \to x + -a = x - a$}
\AXC{$\stackrel{\Pi_{12}}{\vdots}$}\noLine
\UIC{$\forall y\,(x - y = x + -y)$}
\AXC{$\forall y\,(x - y = x + -y) \to x - a = x + -a$\lbl{A2}}
\BIC{$x - a = x + -a$}
\AXC{$\varphi_{89}$\lbl{A1}}
\BIC{$(x - a = x + -a \to x + -a = x - a) \to x + -a = x - a$}
\BIC{$x + -a = x - a$}
\end{ptree}
\noindent$\Pi_{14}$:
\begin{ptree}
\AXC{$x = x$\lbl{A5}}
\UIC{$\forall x\,(x = x)$}
\AXC{$\varphi_{33}$\lbl{A2}}
\BIC{$x + -a + (y + -b) = x + -a + (y + -b)$}
\AXC{$\stackrel{\Pi_{13}}{\vdots}$}\noLine
\UIC{$x + -a = x - a$}
\AXC{$x + -a = x - a \to \varphi_{171}$\lbl{A1}}
\BIC{$\varphi_{171}$}
\BIC{$(x + -a = x - a \to \varphi_{121}) \to x \in \mathbb{R} \to y \in \mathbb{R} \to \varphi_{72}$}
\end{ptree}
\noindent$\Pi_{15}$:
\begin{ptree}
\AXC{$y = z \to x = y \to x = z$\lbl{A6}}
\AXC{$(y = z \to x = y \to x = z) \to x = y \to y = z \to x = z$\lbl{A1}}
\BIC{$x = y \to y = z \to x = z$}
\UIC{$\forall z\,(x = y \to y = z \to x = z)$}
\UIC{$\forall y\,\forall z\,(x = y \to y = z \to x = z)$}
\UIC{$\forall x\,\forall y\,\forall z\,(x = y \to y = z \to x = z)$}
\AXC{$\varphi_{15}$\lbl{A2}}
\BIC{$\forall y\,\forall z\,(v_{186} = y \to y = z \to v_{186} = z)$}
\AXC{$\varphi_{26}$\lbl{A2}}
\BIC{$\forall z\,(v_{186} = v_{187} \to v_{187} = z \to v_{186} = z)$}
\AXC{$\varphi_{43}$\lbl{A2}}
\BIC{$v_{186} = v_{187} \to v_{187} = v_{188} \to v_{186} = v_{188}$}
\UIC{$\forall v_{188}\,(v_{186} = v_{187} \to v_{187} = v_{188} \to v_{186} = v_{188})$}
\UIC{$\varphi_{6}$}
\UIC{$\forall v_{186}\,(\varphi_{6})$}
\end{ptree}
\noindent$\Pi_{16}$:
\begin{ptree}
\AXC{$\stackrel{\Pi_{15}}{\vdots}$}\noLine
\UIC{$\forall v_{186}\,(\varphi_{6})$}
\AXC{$\forall v_{186}\,(\varphi_{6}) \to \varphi_{52}$\lbl{A2}}
\BIC{$\varphi_{52}$}
\AXC{$\varphi_{52} \to \varphi_{124}$\lbl{A2}}
\BIC{$\varphi_{124}$}
\AXC{$\varphi_{124} \to x + y - (a + b) = x + (-a + (y + -b)) \to \varphi_{114}$\lbl{A2}}
\BIC{$x + y - (a + b) = x + (-a + (y + -b)) \to \varphi_{114}$}
\end{ptree}
\noindent$\Pi_{17}$:
\begin{ptree}
\AXC{$\stackrel{\Pi_{3}}{\vdots}$}\noLine
\UIC{$\forall x\,\forall y\,(x = y \to y = x)$}
\AXC{$\forall x\,\forall y\,(x = y \to y = x) \to \forall y\,(v_{167} = y \to y = v_{167})$\lbl{A2}}
\BIC{$\forall y\,(v_{167} = y \to y = v_{167})$}
\AXC{$\varphi_{2}$\lbl{A2}}
\BIC{$v_{167} = v_{168} \to v_{168} = v_{167}$}
\UIC{$\forall v_{168}\,(v_{167} = v_{168} \to v_{168} = v_{167})$}
\UIC{$\forall v_{167}\,\forall v_{168}\,(v_{167} = v_{168} \to v_{168} = v_{167})$}
\end{ptree}
\noindent$\Pi_{18}$:
\begin{ptree}
\AXC{$\stackrel{\Pi_{17}}{\vdots}$}\noLine
\UIC{$\forall v_{167}\,\forall v_{168}\,(v_{167} = v_{168} \to v_{168} = v_{167})$}
\AXC{$\forall v_{167}\,\forall v_{168}\,(v_{167} = v_{168} \to v_{168} = v_{167}) \to \varphi_{68}$\lbl{A2}}
\BIC{$\varphi_{68}$}
\AXC{$\varphi_{68} \to \varphi_{130}$\lbl{A2}}
\BIC{$\varphi_{130}$}
\end{ptree}
\noindent$\Pi_{19}$:
\begin{ptree}
\AXC{$\forall x \in \mathbb{R}\,(-x \in \mathbb{R})$\lbl{(R1)}}
\AXC{$\forall x \in \mathbb{R}\,(-x \in \mathbb{R}) \to b \in \mathbb{R} \to -b \in \mathbb{R}$\lbl{A2}}
\BIC{$b \in \mathbb{R} \to -b \in \mathbb{R}$}
\AXC{$(b \in \mathbb{R} \to -b \in \mathbb{R}) \to \varphi_{144}$\lbl{A1}}
\BIC{$\varphi_{144}$}
\end{ptree}
\noindent$\Pi_{20}$:
\begin{ptree}
\AXC{$x \in \mathbb{R} \to y \in \mathbb{R} \to a \in \mathbb{R} \to b \in \mathbb{R} \to b \in \mathbb{R}$\lbl{A1}}
\AXC{$\stackrel{\Pi_{19}}{\vdots}$}\noLine
\UIC{$\varphi_{144}$}
\BIC{$x \in \mathbb{R} \to y \in \mathbb{R} \to a \in \mathbb{R} \to b \in \mathbb{R} \to -b \in \mathbb{R}$}
\end{ptree}
\noindent$\Pi_{21}$:
\begin{ptree}
\AXC{$\forall x\,\forall y\,(x \in \mathbb{R} \to y \in \mathbb{R} \to x + y \in \mathbb{R})$\lbl{(R1)}}
\AXC{$\varphi_{119}$\lbl{A2}}
\BIC{$\forall y\,(v_{161} \in \mathbb{R} \to y \in \mathbb{R} \to v_{161} + y \in \mathbb{R})$}
\AXC{$\varphi_{122}$\lbl{A2}}
\BIC{$v_{161} \in \mathbb{R} \to v_{162} \in \mathbb{R} \to v_{161} + v_{162} \in \mathbb{R}$}
\UIC{$\forall v_{162}\,(v_{161} \in \mathbb{R} \to v_{162} \in \mathbb{R} \to v_{161} + v_{162} \in \mathbb{R})$}
\UIC{$\varphi_{39}$}
\AXC{$\varphi_{39} \to \forall v_{162}\,(y \in \mathbb{R} \to v_{162} \in \mathbb{R} \to y + v_{162} \in \mathbb{R})$\lbl{A2}}
\BIC{$\forall v_{162}\,(y \in \mathbb{R} \to v_{162} \in \mathbb{R} \to y + v_{162} \in \mathbb{R})$}
\end{ptree}
\noindent$\Pi_{22}$:
\begin{ptree}
\AXC{$x \in \mathbb{R} \to y \in \mathbb{R} \to a \in \mathbb{R} \to b \in \mathbb{R} \to y \in \mathbb{R}$\lbl{A1}}
\AXC{$\stackrel{\Pi_{21}}{\vdots}$}\noLine
\UIC{$\forall v_{162}\,(y \in \mathbb{R} \to v_{162} \in \mathbb{R} \to y + v_{162} \in \mathbb{R})$}
\AXC{$\varphi_{125}$\lbl{A2}}
\BIC{$y \in \mathbb{R} \to -b \in \mathbb{R} \to y + -b \in \mathbb{R}$}
\AXC{$\varphi_{165}$\lbl{A1}}
\BIC{$(x \in \mathbb{R} \to y \in \mathbb{R} \to a \in \mathbb{R} \to b \in \mathbb{R} \to y \in \mathbb{R}) \to \varphi_{146}$}
\BIC{$\varphi_{146}$}
\end{ptree}
\noindent$\Pi_{23}$:
\begin{ptree}
\AXC{$\forall x \in \mathbb{R}\,(-x \in \mathbb{R})$\lbl{(R1)}}
\AXC{$\forall x \in \mathbb{R}\,(-x \in \mathbb{R}) \to a \in \mathbb{R} \to -a \in \mathbb{R}$\lbl{A2}}
\BIC{$a \in \mathbb{R} \to -a \in \mathbb{R}$}
\AXC{$(a \in \mathbb{R} \to -a \in \mathbb{R}) \to \varphi_{145}$\lbl{A1}}
\BIC{$\varphi_{145}$}
\end{ptree}
\noindent$\Pi_{24}$:
\begin{ptree}
\AXC{$x \in \mathbb{R} \to y \in \mathbb{R} \to a \in \mathbb{R} \to b \in \mathbb{R} \to a \in \mathbb{R}$\lbl{A1}}
\AXC{$\stackrel{\Pi_{23}}{\vdots}$}\noLine
\UIC{$\varphi_{145}$}
\BIC{$x \in \mathbb{R} \to y \in \mathbb{R} \to a \in \mathbb{R} \to b \in \mathbb{R} \to -a \in \mathbb{R}$}
\end{ptree}
\noindent$\Pi_{25}$:
\begin{ptree}
\AXC{$\forall x\,\forall y\,(\varphi_{40})$\lbl{(R2)}}
\AXC{$\forall x\,\forall y\,(\varphi_{40}) \to \forall y\,(\varphi_{40})$\lbl{A2}}
\BIC{$\forall y\,(\varphi_{40})$}
\AXC{$\forall y\,(\varphi_{40}) \to \varphi_{85}$\lbl{A2}}
\BIC{$\varphi_{85}$}
\AXC{$\varphi_{85} \to x \in \mathbb{R} \to -a \in \mathbb{R} \to \varphi_{75}$\lbl{A2}}
\BIC{$x \in \mathbb{R} \to -a \in \mathbb{R} \to \varphi_{75}$}
\end{ptree}
\noindent$\Pi_{26}$:
\begin{ptree}
\AXC{$x \in \mathbb{R} \to y \in \mathbb{R} \to a \in \mathbb{R} \to b \in \mathbb{R} \to x \in \mathbb{R}$\lbl{A1}}
\AXC{$\stackrel{\Pi_{25}}{\vdots}$}\noLine
\UIC{$x \in \mathbb{R} \to -a \in \mathbb{R} \to \varphi_{75}$}
\AXC{$(x \in \mathbb{R} \to -a \in \mathbb{R} \to \varphi_{75}) \to \varphi_{175}$\lbl{A1}}
\BIC{$\varphi_{175}$}
\BIC{$(x \in \mathbb{R} \to y \in \mathbb{R} \to a \in \mathbb{R} \to b \in \mathbb{R} \to -a \in \mathbb{R}) \to \varphi_{152}$}
\end{ptree}
\noindent$\Pi_{27}$:
\begin{ptree}
\AXC{$\stackrel{\Pi_{24}}{\vdots}$}\noLine
\UIC{$x \in \mathbb{R} \to y \in \mathbb{R} \to a \in \mathbb{R} \to b \in \mathbb{R} \to -a \in \mathbb{R}$}
\AXC{$\stackrel{\Pi_{26}}{\vdots}$}\noLine
\UIC{$(x \in \mathbb{R} \to y \in \mathbb{R} \to a \in \mathbb{R} \to b \in \mathbb{R} \to -a \in \mathbb{R}) \to \varphi_{152}$}
\BIC{$\varphi_{152}$}
\end{ptree}
\noindent$\Pi_{28}$:
\begin{ptree}
\AXC{$\stackrel{\Pi_{20}}{\vdots}$}\noLine
\UIC{$x \in \mathbb{R} \to y \in \mathbb{R} \to a \in \mathbb{R} \to b \in \mathbb{R} \to -b \in \mathbb{R}$}
\AXC{$\stackrel{\Pi_{22}}{\vdots}$}\noLine
\UIC{$\varphi_{146}$}
\BIC{$x \in \mathbb{R} \to y \in \mathbb{R} \to a \in \mathbb{R} \to b \in \mathbb{R} \to y + -b \in \mathbb{R}$}
\AXC{$\stackrel{\Pi_{27}}{\vdots}$}\noLine
\UIC{$\varphi_{152}$}
\BIC{$x \in \mathbb{R} \to y \in \mathbb{R} \to \varphi_{79}$}
\AXC{$(x \in \mathbb{R} \to y \in \mathbb{R} \to \varphi_{79}) \to \varphi_{130} \to x \in \mathbb{R} \to y \in \mathbb{R} \to \varphi_{80}$\lbl{A1}}
\BIC{$\varphi_{130} \to x \in \mathbb{R} \to y \in \mathbb{R} \to \varphi_{80}$}
\end{ptree}
\noindent$\Pi_{29}$:
\begin{ptree}
\AXC{$y = z \to x = y \to x = z$\lbl{A6}}
\AXC{$(y = z \to x = y \to x = z) \to x = y \to y = z \to x = z$\lbl{A1}}
\BIC{$x = y \to y = z \to x = z$}
\UIC{$\forall z\,(x = y \to y = z \to x = z)$}
\UIC{$\forall y\,\forall z\,(x = y \to y = z \to x = z)$}
\UIC{$\forall x\,\forall y\,\forall z\,(x = y \to y = z \to x = z)$}
\AXC{$\varphi_{16}$\lbl{A2}}
\BIC{$\forall y\,\forall z\,(v_{183} = y \to y = z \to v_{183} = z)$}
\AXC{$\varphi_{27}$\lbl{A2}}
\BIC{$\forall z\,(v_{183} = v_{184} \to v_{184} = z \to v_{183} = z)$}
\AXC{$\varphi_{44}$\lbl{A2}}
\BIC{$v_{183} = v_{184} \to v_{184} = v_{185} \to v_{183} = v_{185}$}
\UIC{$\forall v_{185}\,(v_{183} = v_{184} \to v_{184} = v_{185} \to v_{183} = v_{185})$}
\UIC{$\varphi_{7}$}
\UIC{$\forall v_{183}\,(\varphi_{7})$}
\end{ptree}
\noindent$\Pi_{30}$:
\begin{ptree}
\AXC{$\stackrel{\Pi_{29}}{\vdots}$}\noLine
\UIC{$\forall v_{183}\,(\varphi_{7})$}
\AXC{$\forall v_{183}\,(\varphi_{7}) \to \varphi_{53}$\lbl{A2}}
\BIC{$\varphi_{53}$}
\AXC{$\varphi_{53} \to \varphi_{126}$\lbl{A2}}
\BIC{$\varphi_{126}$}
\AXC{$\varphi_{126} \to x + y - (a + b) = x + (-a + y + -b) \to \varphi_{115}$\lbl{A2}}
\BIC{$x + y - (a + b) = x + (-a + y + -b) \to \varphi_{115}$}
\end{ptree}
\noindent$\Pi_{31}$:
\begin{ptree}
\AXC{$\forall x\,\forall y\,(\varphi_{40})$\lbl{(R2)}}
\AXC{$\forall x\,\forall y\,(\varphi_{40}) \to \forall y\,(\varphi_{86})$\lbl{A2}}
\BIC{$\forall y\,(\varphi_{86})$}
\AXC{$\forall y\,(\varphi_{86}) \to \varphi_{86}$\lbl{A2}}
\BIC{$\varphi_{86}$}
\AXC{$\varphi_{86} \to \varphi_{102}$\lbl{A2}}
\BIC{$\varphi_{102}$}
\AXC{$\varphi_{102} \to \varphi_{173}$\lbl{A1}}
\BIC{$\varphi_{173}$}
\end{ptree}
\noindent$\Pi_{32}$:
\begin{ptree}
\AXC{$x \in \mathbb{R} \to y \in \mathbb{R} \to a \in \mathbb{R} \to b \in \mathbb{R} \to y \in \mathbb{R}$\lbl{A1}}
\AXC{$\stackrel{\Pi_{24}}{\vdots}$}\noLine
\UIC{$x \in \mathbb{R} \to y \in \mathbb{R} \to a \in \mathbb{R} \to b \in \mathbb{R} \to -a \in \mathbb{R}$}
\AXC{$\stackrel{\Pi_{31}}{\vdots}$}\noLine
\UIC{$\varphi_{173}$}
\BIC{$(x \in \mathbb{R} \to y \in \mathbb{R} \to a \in \mathbb{R} \to b \in \mathbb{R} \to y \in \mathbb{R}) \to \varphi_{150}$}
\BIC{$\varphi_{150}$}
\end{ptree}
\noindent$\Pi_{33}$:
\begin{ptree}
\AXC{$x = x$\lbl{A5}}
\UIC{$\forall x\,(x = x)$}
\AXC{$\varphi_{34}$\lbl{A2}}
\BIC{$x + (-a + y + -b) = x + (-a + y + -b)$}
\AXC{$\stackrel{\Pi_{20}}{\vdots}$}\noLine
\UIC{$x \in \mathbb{R} \to y \in \mathbb{R} \to a \in \mathbb{R} \to b \in \mathbb{R} \to -b \in \mathbb{R}$}
\AXC{$\stackrel{\Pi_{32}}{\vdots}$}\noLine
\UIC{$\varphi_{150}$}
\BIC{$x \in \mathbb{R} \to \varphi_{92}$}
\AXC{$\varphi_{179}$\lbl{A1}}
\BIC{$x + (-a + y + -b) = x + (-a + y + -b) \to \varphi_{166}$}
\BIC{$\varphi_{166}$}
\end{ptree}
\noindent$\Pi_{34}$:
\begin{ptree}
\AXC{$z = w \to \varphi_{59}$\lbl{A6}}
\UIC{$\forall w\,(z = w \to \varphi_{59})$}
\UIC{$\forall z\,\forall w\,(z = w \to \varphi_{59})$}
\AXC{$\forall z\,\forall w\,(z = w \to \varphi_{59}) \to \varphi_{138}$\lbl{A2}}
\BIC{$\varphi_{138}$}
\AXC{$\varphi_{138} \to -a + y + -b = -a + (y + -b) \to \varphi_{131}$\lbl{A2}}
\BIC{$-a + y + -b = -a + (y + -b) \to \varphi_{131}$}
\AXC{$\stackrel{\Pi_{33}}{\vdots}$}\noLine
\UIC{$\varphi_{166}$}
\BIC{$x \in \mathbb{R} \to y \in \mathbb{R} \to \varphi_{81}$}
\end{ptree}
\noindent$\Pi_{35}$:
\begin{ptree}
\AXC{$y = z \to x = y \to x = z$\lbl{A6}}
\AXC{$(y = z \to x = y \to x = z) \to x = y \to y = z \to x = z$\lbl{A1}}
\BIC{$x = y \to y = z \to x = z$}
\UIC{$\forall z\,(x = y \to y = z \to x = z)$}
\UIC{$\forall y\,\forall z\,(x = y \to y = z \to x = z)$}
\UIC{$\forall x\,\forall y\,\forall z\,(x = y \to y = z \to x = z)$}
\AXC{$\varphi_{17}$\lbl{A2}}
\BIC{$\forall y\,\forall z\,(v_{180} = y \to y = z \to v_{180} = z)$}
\AXC{$\varphi_{28}$\lbl{A2}}
\BIC{$\forall z\,(v_{180} = v_{181} \to v_{181} = z \to v_{180} = z)$}
\AXC{$\varphi_{45}$\lbl{A2}}
\BIC{$v_{180} = v_{181} \to v_{181} = v_{182} \to v_{180} = v_{182}$}
\UIC{$\forall v_{182}\,(v_{180} = v_{181} \to v_{181} = v_{182} \to v_{180} = v_{182})$}
\UIC{$\varphi_{8}$}
\UIC{$\forall v_{180}\,(\varphi_{8})$}
\end{ptree}
\noindent$\Pi_{36}$:
\begin{ptree}
\AXC{$\stackrel{\Pi_{35}}{\vdots}$}\noLine
\UIC{$\forall v_{180}\,(\varphi_{8})$}
\AXC{$\forall v_{180}\,(\varphi_{8}) \to \varphi_{54}$\lbl{A2}}
\BIC{$\varphi_{54}$}
\AXC{$\varphi_{54} \to \varphi_{127}$\lbl{A2}}
\BIC{$\varphi_{127}$}
\AXC{$\varphi_{127} \to x + y - (a + b) = x + (y + -a + -b) \to \varphi_{116}$\lbl{A2}}
\BIC{$x + y - (a + b) = x + (y + -a + -b) \to \varphi_{116}$}
\end{ptree}
\noindent$\Pi_{37}$:
\begin{ptree}
\AXC{$z = w \to \varphi_{101}$\lbl{A6}}
\UIC{$\forall w\,(z = w \to \varphi_{101})$}
\UIC{$\forall z\,\forall w\,(z = w \to \varphi_{101})$}
\AXC{$\forall z\,\forall w\,(z = w \to \varphi_{101}) \to \varphi_{139}$\lbl{A2}}
\BIC{$\varphi_{139}$}
\AXC{$\varphi_{139} \to y + -a = -a + y \to \varphi_{132}$\lbl{A2}}
\BIC{$y + -a = -a + y \to \varphi_{132}$}
\end{ptree}
\noindent$\Pi_{38}$:
\begin{ptree}
\AXC{$\forall x\,\forall y\,(x \in \mathbb{R} \to y \in \mathbb{R} \to x + y = y + x)$\lbl{(R3)}}
\AXC{$\varphi_{113}$\lbl{A2}}
\BIC{$\forall y\,(v_{159} \in \mathbb{R} \to y \in \mathbb{R} \to v_{159} + y = y + v_{159})$}
\AXC{$\varphi_{134}$\lbl{A2}}
\BIC{$v_{159} \in \mathbb{R} \to v_{160} \in \mathbb{R} \to v_{159} + v_{160} = v_{160} + v_{159}$}
\UIC{$\varphi_{37}$}
\UIC{$\forall v_{159}\,(\varphi_{37})$}
\AXC{$\varphi_{143}$\lbl{A2}}
\BIC{$\forall v_{160}\,(y \in \mathbb{R} \to v_{160} \in \mathbb{R} \to y + v_{160} = v_{160} + y)$}
\AXC{$\varphi_{135}$\lbl{A2}}
\BIC{$y \in \mathbb{R} \to -a \in \mathbb{R} \to y + -a = -a + y$}
\AXC{$\varphi_{168}$\lbl{A1}}
\BIC{$(x \in \mathbb{R} \to y \in \mathbb{R} \to a \in \mathbb{R} \to b \in \mathbb{R} \to y \in \mathbb{R}) \to \varphi_{147}$}
\end{ptree}
\noindent$\Pi_{39}$:
\begin{ptree}
\AXC{$x \in \mathbb{R} \to y \in \mathbb{R} \to a \in \mathbb{R} \to b \in \mathbb{R} \to y \in \mathbb{R}$\lbl{A1}}
\AXC{$\stackrel{\Pi_{38}}{\vdots}$}\noLine
\UIC{$(x \in \mathbb{R} \to y \in \mathbb{R} \to a \in \mathbb{R} \to b \in \mathbb{R} \to y \in \mathbb{R}) \to \varphi_{147}$}
\BIC{$\varphi_{147}$}
\end{ptree}
\noindent$\Pi_{40}$:
\begin{ptree}
\AXC{$x = x$\lbl{A5}}
\UIC{$\forall x\,(x = x)$}
\AXC{$\varphi_{35}$\lbl{A2}}
\BIC{$x + (y + -a + -b) = x + (y + -a + -b)$}
\AXC{$\stackrel{\Pi_{24}}{\vdots}$}\noLine
\UIC{$x \in \mathbb{R} \to y \in \mathbb{R} \to a \in \mathbb{R} \to b \in \mathbb{R} \to -a \in \mathbb{R}$}
\AXC{$\stackrel{\Pi_{39}}{\vdots}$}\noLine
\UIC{$\varphi_{147}$}
\BIC{$x \in \mathbb{R} \to y \in \mathbb{R} \to a \in \mathbb{R} \to b \in \mathbb{R} \to y + -a = -a + y$}
\AXC{$\varphi_{177}$\lbl{A1}}
\BIC{$\varphi_{172}$}
\BIC{$(y + -a = -a + y \to \varphi_{132}) \to x \in \mathbb{R} \to y \in \mathbb{R} \to \varphi_{82}$}
\end{ptree}
\noindent$\Pi_{41}$:
\begin{ptree}
\AXC{$y = z \to x = y \to x = z$\lbl{A6}}
\AXC{$(y = z \to x = y \to x = z) \to x = y \to y = z \to x = z$\lbl{A1}}
\BIC{$x = y \to y = z \to x = z$}
\UIC{$\forall z\,(x = y \to y = z \to x = z)$}
\UIC{$\forall y\,\forall z\,(x = y \to y = z \to x = z)$}
\UIC{$\forall x\,\forall y\,\forall z\,(x = y \to y = z \to x = z)$}
\AXC{$\varphi_{18}$\lbl{A2}}
\BIC{$\forall y\,\forall z\,(v_{177} = y \to y = z \to v_{177} = z)$}
\AXC{$\varphi_{29}$\lbl{A2}}
\BIC{$\forall z\,(v_{177} = v_{178} \to v_{178} = z \to v_{177} = z)$}
\AXC{$\varphi_{46}$\lbl{A2}}
\BIC{$v_{177} = v_{178} \to v_{178} = v_{179} \to v_{177} = v_{179}$}
\UIC{$\forall v_{179}\,(v_{177} = v_{178} \to v_{178} = v_{179} \to v_{177} = v_{179})$}
\UIC{$\varphi_{9}$}
\UIC{$\forall v_{177}\,(\varphi_{9})$}
\end{ptree}
\noindent$\Pi_{42}$:
\begin{ptree}
\AXC{$\stackrel{\Pi_{41}}{\vdots}$}\noLine
\UIC{$\forall v_{177}\,(\varphi_{9})$}
\AXC{$\forall v_{177}\,(\varphi_{9}) \to \varphi_{55}$\lbl{A2}}
\BIC{$\varphi_{55}$}
\AXC{$\varphi_{55} \to \varphi_{128}$\lbl{A2}}
\BIC{$\varphi_{128}$}
\AXC{$\varphi_{128} \to x + y - (a + b) = x + (y + (-a + -b)) \to \varphi_{117}$\lbl{A2}}
\BIC{$x + y - (a + b) = x + (y + (-a + -b)) \to \varphi_{117}$}
\end{ptree}
\noindent$\Pi_{43}$:
\begin{ptree}
\AXC{$\stackrel{\Pi_{3}}{\vdots}$}\noLine
\UIC{$\forall x\,\forall y\,(x = y \to y = x)$}
\AXC{$\forall x\,\forall y\,(x = y \to y = x) \to \forall y\,(v_{165} = y \to y = v_{165})$\lbl{A2}}
\BIC{$\forall y\,(v_{165} = y \to y = v_{165})$}
\AXC{$\varphi_{3}$\lbl{A2}}
\BIC{$v_{165} = v_{166} \to v_{166} = v_{165}$}
\UIC{$\forall v_{166}\,(v_{165} = v_{166} \to v_{166} = v_{165})$}
\UIC{$\forall v_{165}\,\forall v_{166}\,(v_{165} = v_{166} \to v_{166} = v_{165})$}
\AXC{$\varphi_{109}$\lbl{A2}}
\BIC{$\forall v_{166}\,(y + -a + -b = v_{166} \to v_{166} = y + -a + -b)$}
\AXC{$\varphi_{149}$\lbl{A2}}
\BIC{$\varphi_{98}$}
\end{ptree}
\noindent$\Pi_{44}$:
\begin{ptree}
\AXC{$\forall x\,\forall y\,(\varphi_{40})$\lbl{(R2)}}
\AXC{$\forall x\,\forall y\,(\varphi_{40}) \to \varphi_{76}$\lbl{A2}}
\BIC{$\varphi_{76}$}
\AXC{$\varphi_{76} \to \varphi_{95}$\lbl{A2}}
\BIC{$\varphi_{95}$}
\AXC{$\varphi_{95} \to v_{156} \in \mathbb{R} \to v_{157} \in \mathbb{R} \to \varphi_{22}$\lbl{A2}}
\BIC{$v_{156} \in \mathbb{R} \to v_{157} \in \mathbb{R} \to \varphi_{22}$}
\UIC{$\forall v_{158}\,(v_{156} \in \mathbb{R} \to v_{157} \in \mathbb{R} \to \varphi_{22})$}
\UIC{$\forall v_{157}\,\forall v_{158}\,(v_{156} \in \mathbb{R} \to v_{157} \in \mathbb{R} \to \varphi_{22})$}
\UIC{$\forall v_{156}\,\forall v_{157}\,\forall v_{158}\,(v_{156} \in \mathbb{R} \to v_{157} \in \mathbb{R} \to \varphi_{22})$}
\AXC{$\varphi_{154}$\lbl{A2}}
\BIC{$\varphi_{103}$}
\end{ptree}
\noindent$\Pi_{45}$:
\begin{ptree}
\AXC{$x \in \mathbb{R} \to y \in \mathbb{R} \to a \in \mathbb{R} \to b \in \mathbb{R} \to y \in \mathbb{R}$\lbl{A1}}
\AXC{$\stackrel{\Pi_{44}}{\vdots}$}\noLine
\UIC{$\varphi_{103}$}
\AXC{$\varphi_{103} \to \varphi_{104}$\lbl{A2}}
\BIC{$\varphi_{104}$}
\AXC{$\varphi_{104} \to \varphi_{105}$\lbl{A2}}
\BIC{$\varphi_{105}$}
\AXC{$\varphi_{105} \to \varphi_{174}$\lbl{A1}}
\BIC{$\varphi_{174}$}
\BIC{$(x \in \mathbb{R} \to y \in \mathbb{R} \to a \in \mathbb{R} \to b \in \mathbb{R} \to -a \in \mathbb{R}) \to \varphi_{151}$}
\end{ptree}
\noindent$\Pi_{46}$:
\begin{ptree}
\AXC{$\stackrel{\Pi_{24}}{\vdots}$}\noLine
\UIC{$x \in \mathbb{R} \to y \in \mathbb{R} \to a \in \mathbb{R} \to b \in \mathbb{R} \to -a \in \mathbb{R}$}
\AXC{$\stackrel{\Pi_{45}}{\vdots}$}\noLine
\UIC{$(x \in \mathbb{R} \to y \in \mathbb{R} \to a \in \mathbb{R} \to b \in \mathbb{R} \to -a \in \mathbb{R}) \to \varphi_{151}$}
\BIC{$\varphi_{151}$}
\end{ptree}
\noindent$\Pi_{47}$:
\begin{ptree}
\AXC{$\stackrel{\Pi_{43}}{\vdots}$}\noLine
\UIC{$\varphi_{98}$}
\AXC{$\stackrel{\Pi_{20}}{\vdots}$}\noLine
\UIC{$x \in \mathbb{R} \to y \in \mathbb{R} \to a \in \mathbb{R} \to b \in \mathbb{R} \to -b \in \mathbb{R}$}
\AXC{$\stackrel{\Pi_{46}}{\vdots}$}\noLine
\UIC{$\varphi_{151}$}
\BIC{$x \in \mathbb{R} \to \varphi_{93}$}
\AXC{$(x \in \mathbb{R} \to \varphi_{93}) \to \varphi_{98} \to x \in \mathbb{R} \to \varphi_{94}$\lbl{A1}}
\BIC{$\varphi_{98} \to x \in \mathbb{R} \to \varphi_{94}$}
\BIC{$x \in \mathbb{R} \to \varphi_{94}$}
\AXC{$\varphi_{180}$\lbl{A1}}
\BIC{$x + (y + (-a + -b)) = x + (y + (-a + -b)) \to \varphi_{167}$}
\end{ptree}
\noindent$\Pi_{48}$:
\begin{ptree}
\AXC{$x = x$\lbl{A5}}
\UIC{$\forall x\,(x = x)$}
\AXC{$\varphi_{36}$\lbl{A2}}
\BIC{$x + (y + (-a + -b)) = x + (y + (-a + -b))$}
\AXC{$\stackrel{\Pi_{47}}{\vdots}$}\noLine
\UIC{$x + (y + (-a + -b)) = x + (y + (-a + -b)) \to \varphi_{167}$}
\BIC{$\varphi_{167}$}
\end{ptree}
\noindent$\Pi_{49}$:
\begin{ptree}
\AXC{$z = w \to \varphi_{60}$\lbl{A6}}
\UIC{$\forall w\,(z = w \to \varphi_{60})$}
\UIC{$\forall z\,\forall w\,(z = w \to \varphi_{60})$}
\AXC{$\forall z\,\forall w\,(z = w \to \varphi_{60}) \to \varphi_{140}$\lbl{A2}}
\BIC{$\varphi_{140}$}
\AXC{$\varphi_{140} \to y + (-a + -b) = y + -a + -b \to \varphi_{133}$\lbl{A2}}
\BIC{$y + (-a + -b) = y + -a + -b \to \varphi_{133}$}
\AXC{$\stackrel{\Pi_{48}}{\vdots}$}\noLine
\UIC{$\varphi_{167}$}
\BIC{$x \in \mathbb{R} \to y \in \mathbb{R} \to \varphi_{83}$}
\end{ptree}
\noindent$\Pi_{50}$:
\begin{ptree}
\AXC{$y = z \to x = y \to x = z$\lbl{A6}}
\AXC{$(y = z \to x = y \to x = z) \to x = y \to y = z \to x = z$\lbl{A1}}
\BIC{$x = y \to y = z \to x = z$}
\UIC{$\forall z\,(x = y \to y = z \to x = z)$}
\UIC{$\forall y\,\forall z\,(x = y \to y = z \to x = z)$}
\UIC{$\forall x\,\forall y\,\forall z\,(x = y \to y = z \to x = z)$}
\AXC{$\varphi_{19}$\lbl{A2}}
\BIC{$\forall y\,\forall z\,(v_{174} = y \to y = z \to v_{174} = z)$}
\AXC{$\varphi_{30}$\lbl{A2}}
\BIC{$\forall z\,(v_{174} = v_{175} \to v_{175} = z \to v_{174} = z)$}
\AXC{$\varphi_{47}$\lbl{A2}}
\BIC{$v_{174} = v_{175} \to v_{175} = v_{176} \to v_{174} = v_{176}$}
\UIC{$\forall v_{176}\,(v_{174} = v_{175} \to v_{175} = v_{176} \to v_{174} = v_{176})$}
\UIC{$\varphi_{10}$}
\UIC{$\forall v_{174}\,(\varphi_{10})$}
\end{ptree}
\noindent$\Pi_{51}$:
\begin{ptree}
\AXC{$\stackrel{\Pi_{50}}{\vdots}$}\noLine
\UIC{$\forall v_{174}\,(\varphi_{10})$}
\AXC{$\forall v_{174}\,(\varphi_{10}) \to \varphi_{56}$\lbl{A2}}
\BIC{$\varphi_{56}$}
\AXC{$\varphi_{56} \to \varphi_{129}$\lbl{A2}}
\BIC{$\varphi_{129}$}
\AXC{$\varphi_{129} \to x + y - (a + b) = x + y + (-a + -b) \to \varphi_{118}$\lbl{A2}}
\BIC{$x + y - (a + b) = x + y + (-a + -b) \to \varphi_{118}$}
\end{ptree}
\noindent$\Pi_{52}$:
\begin{ptree}
\AXC{$\forall x\,\forall y\,(x \in \mathbb{R} \to y \in \mathbb{R} \to x + y \in \mathbb{R})$\lbl{(R1)}}
\AXC{$\varphi_{120}$\lbl{A2}}
\BIC{$\forall y\,(-a \in \mathbb{R} \to y \in \mathbb{R} \to -a + y \in \mathbb{R})$}
\AXC{$\varphi_{141}$\lbl{A2}}
\BIC{$-a \in \mathbb{R} \to -b \in \mathbb{R} \to -a + -b \in \mathbb{R}$}
\AXC{$\varphi_{170}$\lbl{A1}}
\BIC{$(x \in \mathbb{R} \to y \in \mathbb{R} \to a \in \mathbb{R} \to b \in \mathbb{R} \to -a \in \mathbb{R}) \to \varphi_{148}$}
\end{ptree}
\noindent$\Pi_{53}$:
\begin{ptree}
\AXC{$\stackrel{\Pi_{24}}{\vdots}$}\noLine
\UIC{$x \in \mathbb{R} \to y \in \mathbb{R} \to a \in \mathbb{R} \to b \in \mathbb{R} \to -a \in \mathbb{R}$}
\AXC{$\stackrel{\Pi_{52}}{\vdots}$}\noLine
\UIC{$(x \in \mathbb{R} \to y \in \mathbb{R} \to a \in \mathbb{R} \to b \in \mathbb{R} \to -a \in \mathbb{R}) \to \varphi_{148}$}
\BIC{$\varphi_{148}$}
\end{ptree}
\noindent$\Pi_{54}$:
\begin{ptree}
\AXC{$\forall x\,\forall y\,(\varphi_{40})$\lbl{(R2)}}
\AXC{$\forall x\,\forall y\,(\varphi_{40}) \to \forall y\,(\varphi_{40})$\lbl{A2}}
\BIC{$\forall y\,(\varphi_{40})$}
\AXC{$\forall y\,(\varphi_{40}) \to \varphi_{40}$\lbl{A2}}
\BIC{$\varphi_{40}$}
\AXC{$\varphi_{40} \to x \in \mathbb{R} \to y \in \mathbb{R} \to \varphi_{91}$\lbl{A2}}
\BIC{$x \in \mathbb{R} \to y \in \mathbb{R} \to \varphi_{91}$}
\end{ptree}
\noindent$\Pi_{55}$:
\begin{ptree}
\AXC{$x \in \mathbb{R} \to y \in \mathbb{R} \to a \in \mathbb{R} \to b \in \mathbb{R} \to x \in \mathbb{R}$\lbl{A1}}
\AXC{$\stackrel{\Pi_{54}}{\vdots}$}\noLine
\UIC{$x \in \mathbb{R} \to y \in \mathbb{R} \to \varphi_{91}$}
\AXC{$(x \in \mathbb{R} \to y \in \mathbb{R} \to \varphi_{91}) \to \varphi_{176}$\lbl{A1}}
\BIC{$\varphi_{176}$}
\BIC{$(x \in \mathbb{R} \to y \in \mathbb{R} \to a \in \mathbb{R} \to b \in \mathbb{R} \to y \in \mathbb{R}) \to \varphi_{153}$}
\end{ptree}
\noindent$\Pi_{56}$:
\begin{ptree}
\AXC{$x \in \mathbb{R} \to y \in \mathbb{R} \to a \in \mathbb{R} \to b \in \mathbb{R} \to y \in \mathbb{R}$\lbl{A1}}
\AXC{$\stackrel{\Pi_{55}}{\vdots}$}\noLine
\UIC{$(x \in \mathbb{R} \to y \in \mathbb{R} \to a \in \mathbb{R} \to b \in \mathbb{R} \to y \in \mathbb{R}) \to \varphi_{153}$}
\BIC{$\varphi_{153}$}
\end{ptree}
\noindent$\Pi_{57}$:
\begin{ptree}
\AXC{$y = z \to x = y \to x = z$\lbl{A6}}
\AXC{$(y = z \to x = y \to x = z) \to x = y \to y = z \to x = z$\lbl{A1}}
\BIC{$x = y \to y = z \to x = z$}
\UIC{$\forall z\,(x = y \to y = z \to x = z)$}
\UIC{$\forall y\,\forall z\,(x = y \to y = z \to x = z)$}
\UIC{$\forall x\,\forall y\,\forall z\,(x = y \to y = z \to x = z)$}
\AXC{$\varphi_{20}$\lbl{A2}}
\BIC{$\forall y\,\forall z\,(v_{171} = y \to y = z \to v_{171} = z)$}
\AXC{$\varphi_{31}$\lbl{A2}}
\BIC{$\forall z\,(v_{171} = v_{172} \to v_{172} = z \to v_{171} = z)$}
\AXC{$\varphi_{48}$\lbl{A2}}
\BIC{$v_{171} = v_{172} \to v_{172} = v_{173} \to v_{171} = v_{173}$}
\UIC{$\forall v_{173}\,(v_{171} = v_{172} \to v_{172} = v_{173} \to v_{171} = v_{173})$}
\UIC{$\varphi_{11}$}
\UIC{$\forall v_{171}\,(\varphi_{11})$}
\end{ptree}
\noindent$\Pi_{58}$:
\begin{ptree}
\AXC{$\stackrel{\Pi_{57}}{\vdots}$}\noLine
\UIC{$\forall v_{171}\,(\varphi_{11})$}
\AXC{$\forall v_{171}\,(\varphi_{11}) \to \varphi_{57}$\lbl{A2}}
\BIC{$\varphi_{57}$}
\AXC{$\varphi_{57} \to \varphi_{110}$\lbl{A2}}
\BIC{$\varphi_{110}$}
\AXC{$\varphi_{110} \to x + y - (a + b) = x + y + -(a + b) \to \varphi_{106}$\lbl{A2}}
\BIC{$x + y - (a + b) = x + y + -(a + b) \to \varphi_{106}$}
\end{ptree}
\noindent$\Pi_{59}$:
\begin{ptree}
\AXC{$z = w \to \varphi_{61}$\lbl{A6}}
\UIC{$\forall w\,(z = w \to \varphi_{61})$}
\UIC{$\forall z\,\forall w\,(z = w \to \varphi_{61})$}
\AXC{$\forall z\,\forall w\,(z = w \to \varphi_{61}) \to \varphi_{112}$\lbl{A2}}
\BIC{$\varphi_{112}$}
\AXC{$\varphi_{112} \to -(a + b) = -a + -b \to \varphi_{107}$\lbl{A2}}
\BIC{$-(a + b) = -a + -b \to \varphi_{107}$}
\end{ptree}
\noindent$\Pi_{60}$:
\begin{ptree}
\AXC{$\forall x\,\forall y\,(x \in \mathbb{R} \to y \in \mathbb{R} \to -(x + y) = -x + -y)$\lbl{[F7]}}
\AXC{$\varphi_{142}$\lbl{A2}}
\BIC{$\forall y\,(a \in \mathbb{R} \to y \in \mathbb{R} \to -(a + y) = -a + -y)$}
\AXC{$\varphi_{136}$\lbl{A2}}
\BIC{$a \in \mathbb{R} \to b \in \mathbb{R} \to -(a + b) = -a + -b$}
\end{ptree}
\noindent$\Pi_{61}$:
\begin{ptree}
\AXC{$\stackrel{\Pi_{60}}{\vdots}$}\noLine
\UIC{$a \in \mathbb{R} \to b \in \mathbb{R} \to -(a + b) = -a + -b$}
\AXC{$(a \in \mathbb{R} \to b \in \mathbb{R} \to -(a + b) = -a + -b) \to \varphi_{156}$\lbl{A1}}
\BIC{$\varphi_{156}$}
\end{ptree}
\noindent$\Pi_{62}$:
\begin{ptree}
\AXC{$x \in \mathbb{R} \to y \in \mathbb{R} \to a \in \mathbb{R} \to b \in \mathbb{R} \to b \in \mathbb{R}$\lbl{A1}}
\AXC{$x \in \mathbb{R} \to y \in \mathbb{R} \to a \in \mathbb{R} \to b \in \mathbb{R} \to a \in \mathbb{R}$\lbl{A1}}
\AXC{$\stackrel{\Pi_{61}}{\vdots}$}\noLine
\UIC{$\varphi_{156}$}
\BIC{$(x \in \mathbb{R} \to y \in \mathbb{R} \to a \in \mathbb{R} \to b \in \mathbb{R} \to b \in \mathbb{R}) \to \varphi_{87}$}
\BIC{$\varphi_{87}$}
\AXC{$\varphi_{87} \to \varphi_{169}$\lbl{A1}}
\BIC{$\varphi_{169}$}
\end{ptree}
\noindent$\Pi_{63}$:
\begin{ptree}
\AXC{$\stackrel{\Pi_{59}}{\vdots}$}\noLine
\UIC{$-(a + b) = -a + -b \to \varphi_{107}$}
\AXC{$x = x$\lbl{A5}}
\UIC{$\forall x\,(x = x)$}
\AXC{$\forall x\,(x = x) \to x + y + -(a + b) = x + y + -(a + b)$\lbl{A2}}
\BIC{$x + y + -(a + b) = x + y + -(a + b)$}
\AXC{$\stackrel{\Pi_{62}}{\vdots}$}\noLine
\UIC{$\varphi_{169}$}
\BIC{$(-(a + b) = -a + -b \to \varphi_{107}) \to x \in \mathbb{R} \to y \in \mathbb{R} \to \varphi_{73}$}
\BIC{$x \in \mathbb{R} \to y \in \mathbb{R} \to \varphi_{73}$}
\end{ptree}
\noindent$\Pi_{64}$:
\begin{ptree}
\AXC{$\forall x\,\forall y\,(x - y = x + -y)$\lbl{Def$-$}}
\AXC{$\varphi_{21}$\lbl{A2}}
\BIC{$\forall y\,(v_{154} - y = v_{154} + -y)$}
\AXC{$\varphi_{32}$\lbl{A2}}
\BIC{$v_{154} - v_{155} = v_{154} + -v_{155}$}
\UIC{$\forall v_{155}\,(v_{154} - v_{155} = v_{154} + -v_{155})$}
\UIC{$\forall v_{154}\,\forall v_{155}\,(v_{154} - v_{155} = v_{154} + -v_{155})$}
\AXC{$\varphi_{97}$\lbl{A2}}
\BIC{$\forall v_{155}\,(x + y - v_{155} = x + y + -v_{155})$}
\AXC{$\varphi_{90}$\lbl{A2}}
\BIC{$x + y - (a + b) = x + y + -(a + b)$}
\AXC{$\varphi_{178}$\lbl{A1}}
\BIC{$(x \in \mathbb{R} \to y \in \mathbb{R} \to \varphi_{73}) \to \varphi_{158}$}
\end{ptree}
\noindent$\Pi_{65}$:
\begin{ptree}
\AXC{$\stackrel{\Pi_{58}}{\vdots}$}\noLine
\UIC{$x + y - (a + b) = x + y + -(a + b) \to \varphi_{106}$}
\AXC{$\stackrel{\Pi_{63}}{\vdots}$}\noLine
\UIC{$x \in \mathbb{R} \to y \in \mathbb{R} \to \varphi_{73}$}
\AXC{$\stackrel{\Pi_{64}}{\vdots}$}\noLine
\UIC{$(x \in \mathbb{R} \to y \in \mathbb{R} \to \varphi_{73}) \to \varphi_{158}$}
\BIC{$\varphi_{158}$}
\BIC{$x \in \mathbb{R} \to y \in \mathbb{R} \to \varphi_{62}$}
\end{ptree}
\noindent$\Pi_{66}$:
\begin{ptree}
\AXC{$\stackrel{\Pi_{65}}{\vdots}$}\noLine
\UIC{$x \in \mathbb{R} \to y \in \mathbb{R} \to \varphi_{62}$}
\AXC{$(x \in \mathbb{R} \to y \in \mathbb{R} \to \varphi_{62}) \to (x \in \mathbb{R} \to y \in \mathbb{R} \to \varphi_{84}) \to \varphi_{159}$\lbl{A1}}
\BIC{$(x \in \mathbb{R} \to y \in \mathbb{R} \to \varphi_{84}) \to \varphi_{159}$}
\end{ptree}
\noindent$\Pi_{67}$:
\begin{ptree}
\AXC{$\stackrel{\Pi_{20}}{\vdots}$}\noLine
\UIC{$x \in \mathbb{R} \to y \in \mathbb{R} \to a \in \mathbb{R} \to b \in \mathbb{R} \to -b \in \mathbb{R}$}
\AXC{$\stackrel{\Pi_{53}}{\vdots}$}\noLine
\UIC{$\varphi_{148}$}
\BIC{$x \in \mathbb{R} \to y \in \mathbb{R} \to a \in \mathbb{R} \to b \in \mathbb{R} \to -a + -b \in \mathbb{R}$}
\AXC{$\stackrel{\Pi_{56}}{\vdots}$}\noLine
\UIC{$\varphi_{153}$}
\BIC{$x \in \mathbb{R} \to y \in \mathbb{R} \to \varphi_{84}$}
\AXC{$\stackrel{\Pi_{66}}{\vdots}$}\noLine
\UIC{$(x \in \mathbb{R} \to y \in \mathbb{R} \to \varphi_{84}) \to \varphi_{159}$}
\BIC{$\varphi_{159}$}
\end{ptree}
\noindent$\Pi_{68}$:
\begin{ptree}
\AXC{$\stackrel{\Pi_{51}}{\vdots}$}\noLine
\UIC{$x + y - (a + b) = x + y + (-a + -b) \to \varphi_{118}$}
\AXC{$\stackrel{\Pi_{67}}{\vdots}$}\noLine
\UIC{$\varphi_{159}$}
\BIC{$x \in \mathbb{R} \to y \in \mathbb{R} \to \varphi_{63}$}
\AXC{$(x \in \mathbb{R} \to y \in \mathbb{R} \to \varphi_{63}) \to (x \in \mathbb{R} \to y \in \mathbb{R} \to \varphi_{83}) \to \varphi_{160}$\lbl{A1}}
\BIC{$(x \in \mathbb{R} \to y \in \mathbb{R} \to \varphi_{83}) \to \varphi_{160}$}
\end{ptree}
\noindent$\Pi_{69}$:
\begin{ptree}
\AXC{$\stackrel{\Pi_{42}}{\vdots}$}\noLine
\UIC{$x + y - (a + b) = x + (y + (-a + -b)) \to \varphi_{117}$}
\AXC{$\stackrel{\Pi_{49}}{\vdots}$}\noLine
\UIC{$x \in \mathbb{R} \to y \in \mathbb{R} \to \varphi_{83}$}
\AXC{$\stackrel{\Pi_{68}}{\vdots}$}\noLine
\UIC{$(x \in \mathbb{R} \to y \in \mathbb{R} \to \varphi_{83}) \to \varphi_{160}$}
\BIC{$\varphi_{160}$}
\BIC{$x \in \mathbb{R} \to y \in \mathbb{R} \to \varphi_{64}$}
\end{ptree}
\noindent$\Pi_{70}$:
\begin{ptree}
\AXC{$\stackrel{\Pi_{69}}{\vdots}$}\noLine
\UIC{$x \in \mathbb{R} \to y \in \mathbb{R} \to \varphi_{64}$}
\AXC{$(x \in \mathbb{R} \to y \in \mathbb{R} \to \varphi_{64}) \to (x \in \mathbb{R} \to y \in \mathbb{R} \to \varphi_{82}) \to \varphi_{161}$\lbl{A1}}
\BIC{$(x \in \mathbb{R} \to y \in \mathbb{R} \to \varphi_{82}) \to \varphi_{161}$}
\end{ptree}
\noindent$\Pi_{71}$:
\begin{ptree}
\AXC{$\stackrel{\Pi_{37}}{\vdots}$}\noLine
\UIC{$y + -a = -a + y \to \varphi_{132}$}
\AXC{$\stackrel{\Pi_{40}}{\vdots}$}\noLine
\UIC{$(y + -a = -a + y \to \varphi_{132}) \to x \in \mathbb{R} \to y \in \mathbb{R} \to \varphi_{82}$}
\BIC{$x \in \mathbb{R} \to y \in \mathbb{R} \to \varphi_{82}$}
\AXC{$\stackrel{\Pi_{70}}{\vdots}$}\noLine
\UIC{$(x \in \mathbb{R} \to y \in \mathbb{R} \to \varphi_{82}) \to \varphi_{161}$}
\BIC{$\varphi_{161}$}
\end{ptree}
\noindent$\Pi_{72}$:
\begin{ptree}
\AXC{$\stackrel{\Pi_{36}}{\vdots}$}\noLine
\UIC{$x + y - (a + b) = x + (y + -a + -b) \to \varphi_{116}$}
\AXC{$\stackrel{\Pi_{71}}{\vdots}$}\noLine
\UIC{$\varphi_{161}$}
\BIC{$x \in \mathbb{R} \to y \in \mathbb{R} \to \varphi_{65}$}
\AXC{$(x \in \mathbb{R} \to y \in \mathbb{R} \to \varphi_{65}) \to (x \in \mathbb{R} \to y \in \mathbb{R} \to \varphi_{81}) \to \varphi_{162}$\lbl{A1}}
\BIC{$(x \in \mathbb{R} \to y \in \mathbb{R} \to \varphi_{81}) \to \varphi_{162}$}
\end{ptree}
\noindent$\Pi_{73}$:
\begin{ptree}
\AXC{$\stackrel{\Pi_{30}}{\vdots}$}\noLine
\UIC{$x + y - (a + b) = x + (-a + y + -b) \to \varphi_{115}$}
\AXC{$\stackrel{\Pi_{34}}{\vdots}$}\noLine
\UIC{$x \in \mathbb{R} \to y \in \mathbb{R} \to \varphi_{81}$}
\AXC{$\stackrel{\Pi_{72}}{\vdots}$}\noLine
\UIC{$(x \in \mathbb{R} \to y \in \mathbb{R} \to \varphi_{81}) \to \varphi_{162}$}
\BIC{$\varphi_{162}$}
\BIC{$x \in \mathbb{R} \to y \in \mathbb{R} \to \varphi_{66}$}
\end{ptree}
\noindent$\Pi_{74}$:
\begin{ptree}
\AXC{$\stackrel{\Pi_{18}}{\vdots}$}\noLine
\UIC{$\varphi_{130}$}
\AXC{$\stackrel{\Pi_{28}}{\vdots}$}\noLine
\UIC{$\varphi_{130} \to x \in \mathbb{R} \to y \in \mathbb{R} \to \varphi_{80}$}
\BIC{$x \in \mathbb{R} \to y \in \mathbb{R} \to \varphi_{80}$}
\AXC{$\stackrel{\Pi_{73}}{\vdots}$}\noLine
\UIC{$x \in \mathbb{R} \to y \in \mathbb{R} \to \varphi_{66}$}
\AXC{$(x \in \mathbb{R} \to y \in \mathbb{R} \to \varphi_{66}) \to (x \in \mathbb{R} \to y \in \mathbb{R} \to \varphi_{80}) \to \varphi_{163}$\lbl{A1}}
\BIC{$(x \in \mathbb{R} \to y \in \mathbb{R} \to \varphi_{80}) \to \varphi_{163}$}
\BIC{$\varphi_{163}$}
\end{ptree}
\noindent$\Pi_{75}$:
\begin{ptree}
\AXC{$\stackrel{\Pi_{16}}{\vdots}$}\noLine
\UIC{$x + y - (a + b) = x + (-a + (y + -b)) \to \varphi_{114}$}
\AXC{$\stackrel{\Pi_{74}}{\vdots}$}\noLine
\UIC{$\varphi_{163}$}
\BIC{$x \in \mathbb{R} \to y \in \mathbb{R} \to \varphi_{67}$}
\AXC{$(x \in \mathbb{R} \to y \in \mathbb{R} \to \varphi_{67}) \to (x \in \mathbb{R} \to y \in \mathbb{R} \to \varphi_{72}) \to \varphi_{157}$\lbl{A1}}
\BIC{$(x \in \mathbb{R} \to y \in \mathbb{R} \to \varphi_{72}) \to \varphi_{157}$}
\end{ptree}
\noindent$\Pi_{76}$:
\begin{ptree}
\AXC{$\stackrel{\Pi_{10}}{\vdots}$}\noLine
\UIC{$x + -a = x - a \to \varphi_{121}$}
\AXC{$\stackrel{\Pi_{14}}{\vdots}$}\noLine
\UIC{$(x + -a = x - a \to \varphi_{121}) \to x \in \mathbb{R} \to y \in \mathbb{R} \to \varphi_{72}$}
\BIC{$x \in \mathbb{R} \to y \in \mathbb{R} \to \varphi_{72}$}
\AXC{$\stackrel{\Pi_{75}}{\vdots}$}\noLine
\UIC{$(x \in \mathbb{R} \to y \in \mathbb{R} \to \varphi_{72}) \to \varphi_{157}$}
\BIC{$\varphi_{157}$}
\end{ptree}
\noindent$\Pi_{77}$:
\begin{ptree}
\AXC{$\stackrel{\Pi_{7}}{\vdots}$}\noLine
\UIC{$x \in \mathbb{R} \to \varphi_{77}$}
\AXC{$\stackrel{\Pi_{9}}{\vdots}$}\noLine
\UIC{$x + y - (a + b) = x + -a + (y + -b) \to \varphi_{99}$}
\AXC{$\stackrel{\Pi_{76}}{\vdots}$}\noLine
\UIC{$\varphi_{157}$}
\BIC{$x \in \mathbb{R} \to \varphi_{78}$}
\AXC{$(x \in \mathbb{R} \to \varphi_{78}) \to (x \in \mathbb{R} \to \varphi_{77}) \to \varphi_{155}$\lbl{A1}}
\BIC{$(x \in \mathbb{R} \to \varphi_{77}) \to \varphi_{155}$}
\BIC{$\varphi_{155}$}
\end{ptree}
\noindent 主木:
\begin{ptree}
\AXC{$\stackrel{\Pi_{1}}{\vdots}$}\noLine
\UIC{$\forall v_{192}\,(\varphi_{4})$}
\AXC{$\forall v_{192}\,(\varphi_{4}) \to \varphi_{50}$\lbl{A2}}
\BIC{$\varphi_{50}$}
\AXC{$\varphi_{50} \to \varphi_{108}$\lbl{A2}}
\BIC{$\varphi_{108}$}
\AXC{$\varphi_{108} \to x + y - (a + b) = x - a + (y + -b) \to \varphi_{71}$\lbl{A2}}
\BIC{$x + y - (a + b) = x - a + (y + -b) \to \varphi_{71}$}
\AXC{$\stackrel{\Pi_{77}}{\vdots}$}\noLine
\UIC{$\varphi_{155}$}
\BIC{$x \in \mathbb{R} \to \varphi_{69}$}
\UIC{$\forall b\,(x \in \mathbb{R} \to \varphi_{69})$}
\UIC{$\forall a\,\forall b\,(x \in \mathbb{R} \to \varphi_{69})$}
\UIC{$\forall y\,\forall a\,\forall b\,(x \in \mathbb{R} \to \varphi_{69})$}
\UIC{$\forall x\,\forall y\,\forall a\,\forall b\,(x \in \mathbb{R} \to \varphi_{69})$}
\end{ptree}
\begin{abbrevlist}
\item $\varphi_{1} :\equiv \forall y\,(v_{169} = y \to y = v_{169}) \to v_{169} = v_{170} \to v_{170} = v_{169}$
\item $\varphi_{2} :\equiv \forall y\,(v_{167} = y \to y = v_{167}) \to v_{167} = v_{168} \to v_{168} = v_{167}$
\item $\varphi_{3} :\equiv \forall y\,(v_{165} = y \to y = v_{165}) \to v_{165} = v_{166} \to v_{166} = v_{165}$
\item $\varphi_{4} :\equiv \forall v_{193}\,\forall v_{194}\,(v_{192} = v_{193} \to v_{193} = v_{194} \to v_{192} = v_{194})$
\item $\varphi_{5} :\equiv \forall v_{190}\,\forall v_{191}\,(v_{189} = v_{190} \to v_{190} = v_{191} \to v_{189} = v_{191})$
\item $\varphi_{6} :\equiv \forall v_{187}\,\forall v_{188}\,(v_{186} = v_{187} \to v_{187} = v_{188} \to v_{186} = v_{188})$
\item $\varphi_{7} :\equiv \forall v_{184}\,\forall v_{185}\,(v_{183} = v_{184} \to v_{184} = v_{185} \to v_{183} = v_{185})$
\item $\varphi_{8} :\equiv \forall v_{181}\,\forall v_{182}\,(v_{180} = v_{181} \to v_{181} = v_{182} \to v_{180} = v_{182})$
\item $\varphi_{9} :\equiv \forall v_{178}\,\forall v_{179}\,(v_{177} = v_{178} \to v_{178} = v_{179} \to v_{177} = v_{179})$
\item $\varphi_{10} :\equiv \forall v_{175}\,\forall v_{176}\,(v_{174} = v_{175} \to v_{175} = v_{176} \to v_{174} = v_{176})$
\item $\varphi_{11} :\equiv \forall v_{172}\,\forall v_{173}\,(v_{171} = v_{172} \to v_{172} = v_{173} \to v_{171} = v_{173})$
\item $\varphi_{12} :\equiv \forall x\,\forall y\,\forall z\,(x = y \to y = z \to x = z) \to \forall y\,\forall z\,(v_{192} = y \to y = z \to v_{192} = z)$
\item $\varphi_{13} :\equiv \forall x\,\forall y\,(x - y = x + -y) \to \forall y\,(v_{163} - y = v_{163} + -y)$
\item $\varphi_{14} :\equiv \forall x\,\forall y\,\forall z\,(x = y \to y = z \to x = z) \to \forall y\,\forall z\,(v_{189} = y \to y = z \to v_{189} = z)$
\item $\varphi_{15} :\equiv \forall x\,\forall y\,\forall z\,(x = y \to y = z \to x = z) \to \forall y\,\forall z\,(v_{186} = y \to y = z \to v_{186} = z)$
\item $\varphi_{16} :\equiv \forall x\,\forall y\,\forall z\,(x = y \to y = z \to x = z) \to \forall y\,\forall z\,(v_{183} = y \to y = z \to v_{183} = z)$
\item $\varphi_{17} :\equiv \forall x\,\forall y\,\forall z\,(x = y \to y = z \to x = z) \to \forall y\,\forall z\,(v_{180} = y \to y = z \to v_{180} = z)$
\item $\varphi_{18} :\equiv \forall x\,\forall y\,\forall z\,(x = y \to y = z \to x = z) \to \forall y\,\forall z\,(v_{177} = y \to y = z \to v_{177} = z)$
\item $\varphi_{19} :\equiv \forall x\,\forall y\,\forall z\,(x = y \to y = z \to x = z) \to \forall y\,\forall z\,(v_{174} = y \to y = z \to v_{174} = z)$
\item $\varphi_{20} :\equiv \forall x\,\forall y\,\forall z\,(x = y \to y = z \to x = z) \to \forall y\,\forall z\,(v_{171} = y \to y = z \to v_{171} = z)$
\item $\varphi_{21} :\equiv \forall x\,\forall y\,(x - y = x + -y) \to \forall y\,(v_{154} - y = v_{154} + -y)$
\item $\varphi_{22} :\equiv v_{158} \in \mathbb{R} \to v_{156} + v_{157} + v_{158} = v_{156} + (v_{157} + v_{158})$
\item $\varphi_{23} :\equiv \forall y\,\forall z\,(v_{192} = y \to y = z \to v_{192} = z) \to \forall z\,(v_{192} = v_{193} \to v_{193} = z \to v_{192} = z)$
\item $\varphi_{24} :\equiv \forall y\,(v_{163} - y = v_{163} + -y) \to v_{163} - v_{164} = v_{163} + -v_{164}$
\item $\varphi_{25} :\equiv \forall y\,\forall z\,(v_{189} = y \to y = z \to v_{189} = z) \to \forall z\,(v_{189} = v_{190} \to v_{190} = z \to v_{189} = z)$
\item $\varphi_{26} :\equiv \forall y\,\forall z\,(v_{186} = y \to y = z \to v_{186} = z) \to \forall z\,(v_{186} = v_{187} \to v_{187} = z \to v_{186} = z)$
\item $\varphi_{27} :\equiv \forall y\,\forall z\,(v_{183} = y \to y = z \to v_{183} = z) \to \forall z\,(v_{183} = v_{184} \to v_{184} = z \to v_{183} = z)$
\item $\varphi_{28} :\equiv \forall y\,\forall z\,(v_{180} = y \to y = z \to v_{180} = z) \to \forall z\,(v_{180} = v_{181} \to v_{181} = z \to v_{180} = z)$
\item $\varphi_{29} :\equiv \forall y\,\forall z\,(v_{177} = y \to y = z \to v_{177} = z) \to \forall z\,(v_{177} = v_{178} \to v_{178} = z \to v_{177} = z)$
\item $\varphi_{30} :\equiv \forall y\,\forall z\,(v_{174} = y \to y = z \to v_{174} = z) \to \forall z\,(v_{174} = v_{175} \to v_{175} = z \to v_{174} = z)$
\item $\varphi_{31} :\equiv \forall y\,\forall z\,(v_{171} = y \to y = z \to v_{171} = z) \to \forall z\,(v_{171} = v_{172} \to v_{172} = z \to v_{171} = z)$
\item $\varphi_{32} :\equiv \forall y\,(v_{154} - y = v_{154} + -y) \to v_{154} - v_{155} = v_{154} + -v_{155}$
\item $\varphi_{33} :\equiv \forall x\,(x = x) \to x + -a + (y + -b) = x + -a + (y + -b)$
\item $\varphi_{34} :\equiv \forall x\,(x = x) \to x + (-a + y + -b) = x + (-a + y + -b)$
\item $\varphi_{35} :\equiv \forall x\,(x = x) \to x + (y + -a + -b) = x + (y + -a + -b)$
\item $\varphi_{36} :\equiv \forall x\,(x = x) \to x + (y + (-a + -b)) = x + (y + (-a + -b))$
\item $\varphi_{37} :\equiv \forall v_{160}\,(v_{159} \in \mathbb{R} \to v_{160} \in \mathbb{R} \to v_{159} + v_{160} = v_{160} + v_{159})$
\item $\varphi_{38} :\equiv \forall y\,(x - a = y \to y = x - a) \to x - a = x + -a \to x + -a = x - a$
\item $\varphi_{39} :\equiv \forall v_{161}\,\forall v_{162}\,(v_{161} \in \mathbb{R} \to v_{162} \in \mathbb{R} \to v_{161} + v_{162} \in \mathbb{R})$
\item $\varphi_{40} :\equiv \forall z\,(x \in \mathbb{R} \to y \in \mathbb{R} \to z \in \mathbb{R} \to x + y + z = x + (y + z))$
\item $\varphi_{41} :\equiv \forall z\,(v_{192} = v_{193} \to v_{193} = z \to v_{192} = z) \to v_{192} = v_{193} \to v_{193} = v_{194} \to v_{192} = v_{194}$
\item $\varphi_{42} :\equiv \forall z\,(v_{189} = v_{190} \to v_{190} = z \to v_{189} = z) \to v_{189} = v_{190} \to v_{190} = v_{191} \to v_{189} = v_{191}$
\item $\varphi_{43} :\equiv \forall z\,(v_{186} = v_{187} \to v_{187} = z \to v_{186} = z) \to v_{186} = v_{187} \to v_{187} = v_{188} \to v_{186} = v_{188}$
\item $\varphi_{44} :\equiv \forall z\,(v_{183} = v_{184} \to v_{184} = z \to v_{183} = z) \to v_{183} = v_{184} \to v_{184} = v_{185} \to v_{183} = v_{185}$
\item $\varphi_{45} :\equiv \forall z\,(v_{180} = v_{181} \to v_{181} = z \to v_{180} = z) \to v_{180} = v_{181} \to v_{181} = v_{182} \to v_{180} = v_{182}$
\item $\varphi_{46} :\equiv \forall z\,(v_{177} = v_{178} \to v_{178} = z \to v_{177} = z) \to v_{177} = v_{178} \to v_{178} = v_{179} \to v_{177} = v_{179}$
\item $\varphi_{47} :\equiv \forall z\,(v_{174} = v_{175} \to v_{175} = z \to v_{174} = z) \to v_{174} = v_{175} \to v_{175} = v_{176} \to v_{174} = v_{176}$
\item $\varphi_{48} :\equiv \forall z\,(v_{171} = v_{172} \to v_{172} = z \to v_{171} = z) \to v_{171} = v_{172} \to v_{172} = v_{173} \to v_{171} = v_{173}$
\item $\varphi_{49} :\equiv \forall v_{169}\,\forall v_{170}\,(v_{169} = v_{170} \to v_{170} = v_{169}) \to \forall v_{170}\,(y - b = v_{170} \to v_{170} = y - b)$
\item $\varphi_{50} :\equiv \forall v_{193}\,\forall v_{194}\,(x + y - (a + b) = v_{193} \to v_{193} = v_{194} \to x + y - (a + b) = v_{194})$
\item $\varphi_{51} :\equiv \forall v_{190}\,\forall v_{191}\,(x + y - (a + b) = v_{190} \to v_{190} = v_{191} \to x + y - (a + b) = v_{191})$
\item $\varphi_{52} :\equiv \forall v_{187}\,\forall v_{188}\,(x + y - (a + b) = v_{187} \to v_{187} = v_{188} \to x + y - (a + b) = v_{188})$
\item $\varphi_{53} :\equiv \forall v_{184}\,\forall v_{185}\,(x + y - (a + b) = v_{184} \to v_{184} = v_{185} \to x + y - (a + b) = v_{185})$
\item $\varphi_{54} :\equiv \forall v_{181}\,\forall v_{182}\,(x + y - (a + b) = v_{181} \to v_{181} = v_{182} \to x + y - (a + b) = v_{182})$
\item $\varphi_{55} :\equiv \forall v_{178}\,\forall v_{179}\,(x + y - (a + b) = v_{178} \to v_{178} = v_{179} \to x + y - (a + b) = v_{179})$
\item $\varphi_{56} :\equiv \forall v_{175}\,\forall v_{176}\,(x + y - (a + b) = v_{175} \to v_{175} = v_{176} \to x + y - (a + b) = v_{176})$
\item $\varphi_{57} :\equiv \forall v_{172}\,\forall v_{173}\,(x + y - (a + b) = v_{172} \to v_{172} = v_{173} \to x + y - (a + b) = v_{173})$
\item $\varphi_{58} :\equiv x - a + (y + -b) = x - a + z \to x - a + (y + -b) = x - a + w$
\item $\varphi_{59} :\equiv x + (-a + y + -b) = x + z \to x + (-a + y + -b) = x + w$
\item $\varphi_{60} :\equiv x + (y + (-a + -b)) = x + z \to x + (y + (-a + -b)) = x + w$
\item $\varphi_{61} :\equiv x + y + -(a + b) = x + y + z \to x + y + -(a + b) = x + y + w$
\item $\varphi_{62} :\equiv a \in \mathbb{R} \to b \in \mathbb{R} \to x + y - (a + b) = x + y + (-a + -b)$
\item $\varphi_{63} :\equiv a \in \mathbb{R} \to b \in \mathbb{R} \to x + y - (a + b) = x + (y + (-a + -b))$
\item $\varphi_{64} :\equiv a \in \mathbb{R} \to b \in \mathbb{R} \to x + y - (a + b) = x + (y + -a + -b)$
\item $\varphi_{65} :\equiv a \in \mathbb{R} \to b \in \mathbb{R} \to x + y - (a + b) = x + (-a + y + -b)$
\item $\varphi_{66} :\equiv a \in \mathbb{R} \to b \in \mathbb{R} \to x + y - (a + b) = x + (-a + (y + -b))$
\item $\varphi_{67} :\equiv a \in \mathbb{R} \to b \in \mathbb{R} \to x + y - (a + b) = x + -a + (y + -b)$
\item $\varphi_{68} :\equiv \forall v_{168}\,(x + -a + (y + -b) = v_{168} \to v_{168} = x + -a + (y + -b))$
\item $\varphi_{69} :\equiv y \in \mathbb{R} \to a \in \mathbb{R} \to b \in \mathbb{R} \to x + y - (a + b) = x - a + (y - b)$
\item $\varphi_{70} :\equiv \forall v_{170}\,(y - b = v_{170} \to v_{170} = y - b) \to y - b = y + -b \to y + -b = y - b$
\item $\varphi_{71} :\equiv x - a + (y + -b) = x - a + (y - b) \to x + y - (a + b) = x - a + (y - b)$
\item $\varphi_{72} :\equiv a \in \mathbb{R} \to b \in \mathbb{R} \to x + -a + (y + -b) = x - a + (y + -b)$
\item $\varphi_{73} :\equiv a \in \mathbb{R} \to b \in \mathbb{R} \to x + y + -(a + b) = x + y + (-a + -b)$
\item $\varphi_{74} :\equiv \forall v_{163}\,\forall v_{164}\,(v_{163} - v_{164} = v_{163} + -v_{164}) \to \forall v_{164}\,(y - v_{164} = y + -v_{164})$
\item $\varphi_{75} :\equiv y + -b \in \mathbb{R} \to x + -a + (y + -b) = x + (-a + (y + -b))$
\item $\varphi_{76} :\equiv \forall y\,\forall z\,(v_{156} \in \mathbb{R} \to y \in \mathbb{R} \to z \in \mathbb{R} \to v_{156} + y + z = v_{156} + (y + z))$
\item $\varphi_{77} :\equiv y \in \mathbb{R} \to a \in \mathbb{R} \to b \in \mathbb{R} \to x - a + (y + -b) = x - a + (y - b)$
\item $\varphi_{78} :\equiv y \in \mathbb{R} \to a \in \mathbb{R} \to b \in \mathbb{R} \to x + y - (a + b) = x - a + (y + -b)$
\item $\varphi_{79} :\equiv a \in \mathbb{R} \to b \in \mathbb{R} \to x + -a + (y + -b) = x + (-a + (y + -b))$
\item $\varphi_{80} :\equiv a \in \mathbb{R} \to b \in \mathbb{R} \to x + (-a + (y + -b)) = x + -a + (y + -b)$
\item $\varphi_{81} :\equiv a \in \mathbb{R} \to b \in \mathbb{R} \to x + (-a + y + -b) = x + (-a + (y + -b))$
\item $\varphi_{82} :\equiv a \in \mathbb{R} \to b \in \mathbb{R} \to x + (y + -a + -b) = x + (-a + y + -b)$
\item $\varphi_{83} :\equiv a \in \mathbb{R} \to b \in \mathbb{R} \to x + (y + (-a + -b)) = x + (y + -a + -b)$
\item $\varphi_{84} :\equiv a \in \mathbb{R} \to b \in \mathbb{R} \to x + y + (-a + -b) = x + (y + (-a + -b))$
\item $\varphi_{85} :\equiv \forall z\,(x \in \mathbb{R} \to -a \in \mathbb{R} \to z \in \mathbb{R} \to x + -a + z = x + (-a + z))$
\item $\varphi_{86} :\equiv \forall z\,(-a \in \mathbb{R} \to y \in \mathbb{R} \to z \in \mathbb{R} \to -a + y + z = -a + (y + z))$
\item $\varphi_{87} :\equiv x \in \mathbb{R} \to y \in \mathbb{R} \to a \in \mathbb{R} \to b \in \mathbb{R} \to -(a + b) = -a + -b$
\item $\varphi_{88} :\equiv y - b = y + -b \to (y - b = y + -b \to y + -b = y - b) \to y + -b = y - b$
\item $\varphi_{89} :\equiv x - a = x + -a \to (x - a = x + -a \to x + -a = x - a) \to x + -a = x - a$
\item $\varphi_{90} :\equiv \forall v_{155}\,(x + y - v_{155} = x + y + -v_{155}) \to x + y - (a + b) = x + y + -(a + b)$
\item $\varphi_{91} :\equiv -a + -b \in \mathbb{R} \to x + y + (-a + -b) = x + (y + (-a + -b))$
\item $\varphi_{92} :\equiv y \in \mathbb{R} \to a \in \mathbb{R} \to b \in \mathbb{R} \to -a + y + -b = -a + (y + -b)$
\item $\varphi_{93} :\equiv y \in \mathbb{R} \to a \in \mathbb{R} \to b \in \mathbb{R} \to y + -a + -b = y + (-a + -b)$
\item $\varphi_{94} :\equiv y \in \mathbb{R} \to a \in \mathbb{R} \to b \in \mathbb{R} \to y + (-a + -b) = y + -a + -b$
\item $\varphi_{95} :\equiv \forall z\,(v_{156} \in \mathbb{R} \to v_{157} \in \mathbb{R} \to z \in \mathbb{R} \to v_{156} + v_{157} + z = v_{156} + (v_{157} + z))$
\item $\varphi_{96} :\equiv x - a + (y + -b) = x - a + (y + -b) \to x - a + (y + -b) = x - a + (y - b)$
\item $\varphi_{97} :\equiv \forall v_{154}\,\forall v_{155}\,(v_{154} - v_{155} = v_{154} + -v_{155}) \to \forall v_{155}\,(x + y - v_{155} = x + y + -v_{155})$
\item $\varphi_{98} :\equiv y + -a + -b = y + (-a + -b) \to y + (-a + -b) = y + -a + -b$
\item $\varphi_{99} :\equiv x + -a + (y + -b) = x - a + (y + -b) \to x + y - (a + b) = x - a + (y + -b)$
\item $\varphi_{100} :\equiv x + -a + (y + -b) = z + (y + -b) \to x + -a + (y + -b) = w + (y + -b)$
\item $\varphi_{101} :\equiv x + (y + -a + -b) = x + (z + -b) \to x + (y + -a + -b) = x + (w + -b)$
\item $\varphi_{102} :\equiv -a \in \mathbb{R} \to y \in \mathbb{R} \to -b \in \mathbb{R} \to -a + y + -b = -a + (y + -b)$
\item $\varphi_{103} :\equiv \forall v_{157}\,\forall v_{158}\,(y \in \mathbb{R} \to v_{157} \in \mathbb{R} \to v_{158} \in \mathbb{R} \to y + v_{157} + v_{158} = y + (v_{157} + v_{158}))$
\item $\varphi_{104} :\equiv \forall v_{158}\,(y \in \mathbb{R} \to -a \in \mathbb{R} \to v_{158} \in \mathbb{R} \to y + -a + v_{158} = y + (-a + v_{158}))$
\item $\varphi_{105} :\equiv y \in \mathbb{R} \to -a \in \mathbb{R} \to -b \in \mathbb{R} \to y + -a + -b = y + (-a + -b)$
\item $\varphi_{106} :\equiv x + y + -(a + b) = x + y + (-a + -b) \to x + y - (a + b) = x + y + (-a + -b)$
\item $\varphi_{107} :\equiv x + y + -(a + b) = x + y + -(a + b) \to x + y + -(a + b) = x + y + (-a + -b)$
\item $\varphi_{108} :\equiv \forall v_{194}\,(x + y - (a + b) = x - a + (y + -b) \to x - a + (y + -b) = v_{194} \to x + y - (a + b) = v_{194})$
\item $\varphi_{109} :\equiv \forall v_{165}\,\forall v_{166}\,(v_{165} = v_{166} \to v_{166} = v_{165}) \to \forall v_{166}\,(y + -a + -b = v_{166} \to v_{166} = y + -a + -b)$
\item $\varphi_{110} :\equiv \forall v_{173}\,(x + y - (a + b) = x + y + -(a + b) \to x + y + -(a + b) = v_{173} \to x + y - (a + b) = v_{173})$
\item $\varphi_{111} :\equiv \forall w\,(y + -b = w \to x - a + (y + -b) = x - a + (y + -b) \to x - a + (y + -b) = x - a + w)$
\item $\varphi_{112} :\equiv \forall w\,(-(a + b) = w \to x + y + -(a + b) = x + y + -(a + b) \to x + y + -(a + b) = x + y + w)$
\item $\varphi_{113} :\equiv \forall x\,\forall y\,(x \in \mathbb{R} \to y \in \mathbb{R} \to x + y = y + x) \to \forall y\,(v_{159} \in \mathbb{R} \to y \in \mathbb{R} \to v_{159} + y = y + v_{159})$
\item $\varphi_{114} :\equiv x + (-a + (y + -b)) = x + -a + (y + -b) \to x + y - (a + b) = x + -a + (y + -b)$
\item $\varphi_{115} :\equiv x + (-a + y + -b) = x + (-a + (y + -b)) \to x + y - (a + b) = x + (-a + (y + -b))$
\item $\varphi_{116} :\equiv x + (y + -a + -b) = x + (-a + y + -b) \to x + y - (a + b) = x + (-a + y + -b)$
\item $\varphi_{117} :\equiv x + (y + (-a + -b)) = x + (y + -a + -b) \to x + y - (a + b) = x + (y + -a + -b)$
\item $\varphi_{118} :\equiv x + y + (-a + -b) = x + (y + (-a + -b)) \to x + y - (a + b) = x + (y + (-a + -b))$
\item $\varphi_{119} :\equiv \forall x\,\forall y\,(x \in \mathbb{R} \to y \in \mathbb{R} \to x + y \in \mathbb{R}) \to \forall y\,(v_{161} \in \mathbb{R} \to y \in \mathbb{R} \to v_{161} + y \in \mathbb{R})$
\item $\varphi_{120} :\equiv \forall x\,\forall y\,(x \in \mathbb{R} \to y \in \mathbb{R} \to x + y \in \mathbb{R}) \to \forall y\,(-a \in \mathbb{R} \to y \in \mathbb{R} \to -a + y \in \mathbb{R})$
\item $\varphi_{121} :\equiv x + -a + (y + -b) = x + -a + (y + -b) \to x + -a + (y + -b) = x - a + (y + -b)$
\item $\varphi_{122} :\equiv \forall y\,(v_{161} \in \mathbb{R} \to y \in \mathbb{R} \to v_{161} + y \in \mathbb{R}) \to v_{161} \in \mathbb{R} \to v_{162} \in \mathbb{R} \to v_{161} + v_{162} \in \mathbb{R}$
\item $\varphi_{123} :\equiv \forall v_{191}\,(x + y - (a + b) = x + -a + (y + -b) \to x + -a + (y + -b) = v_{191} \to x + y - (a + b) = v_{191})$
\item $\varphi_{124} :\equiv \forall v_{188}\,(x + y - (a + b) = x + (-a + (y + -b)) \to x + (-a + (y + -b)) = v_{188} \to x + y - (a + b) = v_{188})$
\item $\varphi_{125} :\equiv \forall v_{162}\,(y \in \mathbb{R} \to v_{162} \in \mathbb{R} \to y + v_{162} \in \mathbb{R}) \to y \in \mathbb{R} \to -b \in \mathbb{R} \to y + -b \in \mathbb{R}$
\item $\varphi_{126} :\equiv \forall v_{185}\,(x + y - (a + b) = x + (-a + y + -b) \to x + (-a + y + -b) = v_{185} \to x + y - (a + b) = v_{185})$
\item $\varphi_{127} :\equiv \forall v_{182}\,(x + y - (a + b) = x + (y + -a + -b) \to x + (y + -a + -b) = v_{182} \to x + y - (a + b) = v_{182})$
\item $\varphi_{128} :\equiv \forall v_{179}\,(x + y - (a + b) = x + (y + (-a + -b)) \to x + (y + (-a + -b)) = v_{179} \to x + y - (a + b) = v_{179})$
\item $\varphi_{129} :\equiv \forall v_{176}\,(x + y - (a + b) = x + y + (-a + -b) \to x + y + (-a + -b) = v_{176} \to x + y - (a + b) = v_{176})$
\item $\varphi_{130} :\equiv x + -a + (y + -b) = x + (-a + (y + -b)) \to x + (-a + (y + -b)) = x + -a + (y + -b)$
\item $\varphi_{131} :\equiv x + (-a + y + -b) = x + (-a + y + -b) \to x + (-a + y + -b) = x + (-a + (y + -b))$
\item $\varphi_{132} :\equiv x + (y + -a + -b) = x + (y + -a + -b) \to x + (y + -a + -b) = x + (-a + y + -b)$
\item $\varphi_{133} :\equiv x + (y + (-a + -b)) = x + (y + (-a + -b)) \to x + (y + (-a + -b)) = x + (y + -a + -b)$
\item $\varphi_{134} :\equiv \forall y\,(v_{159} \in \mathbb{R} \to y \in \mathbb{R} \to v_{159} + y = y + v_{159}) \to v_{159} \in \mathbb{R} \to v_{160} \in \mathbb{R} \to v_{159} + v_{160} = v_{160} + v_{159}$
\item $\varphi_{135} :\equiv \forall v_{160}\,(y \in \mathbb{R} \to v_{160} \in \mathbb{R} \to y + v_{160} = v_{160} + y) \to y \in \mathbb{R} \to -a \in \mathbb{R} \to y + -a = -a + y$
\item $\varphi_{136} :\equiv \forall y\,(a \in \mathbb{R} \to y \in \mathbb{R} \to -(a + y) = -a + -y) \to a \in \mathbb{R} \to b \in \mathbb{R} \to -(a + b) = -a + -b$
\item $\varphi_{137} :\equiv \forall w\,(x + -a = w \to x + -a + (y + -b) = x + -a + (y + -b) \to x + -a + (y + -b) = w + (y + -b))$
\item $\varphi_{138} :\equiv \forall w\,(-a + y + -b = w \to x + (-a + y + -b) = x + (-a + y + -b) \to x + (-a + y + -b) = x + w)$
\item $\varphi_{139} :\equiv \forall w\,(y + -a = w \to x + (y + -a + -b) = x + (y + -a + -b) \to x + (y + -a + -b) = x + (w + -b))$
\item $\varphi_{140} :\equiv \forall w\,(y + (-a + -b) = w \to x + (y + (-a + -b)) = x + (y + (-a + -b)) \to x + (y + (-a + -b)) = x + w)$
\item $\varphi_{141} :\equiv \forall y\,(-a \in \mathbb{R} \to y \in \mathbb{R} \to -a + y \in \mathbb{R}) \to -a \in \mathbb{R} \to -b \in \mathbb{R} \to -a + -b \in \mathbb{R}$
\item $\varphi_{142} :\equiv \forall x\,\forall y\,(x \in \mathbb{R} \to y \in \mathbb{R} \to -(x + y) = -x + -y) \to \forall y\,(a \in \mathbb{R} \to y \in \mathbb{R} \to -(a + y) = -a + -y)$
\item $\varphi_{143} :\equiv \forall v_{159}\,(\varphi_{37}) \to \forall v_{160}\,(y \in \mathbb{R} \to v_{160} \in \mathbb{R} \to y + v_{160} = v_{160} + y)$
\item $\varphi_{144} :\equiv (x \in \mathbb{R} \to y \in \mathbb{R} \to a \in \mathbb{R} \to b \in \mathbb{R} \to b \in \mathbb{R}) \to x \in \mathbb{R} \to y \in \mathbb{R} \to a \in \mathbb{R} \to b \in \mathbb{R} \to -b \in \mathbb{R}$
\item $\varphi_{145} :\equiv (x \in \mathbb{R} \to y \in \mathbb{R} \to a \in \mathbb{R} \to b \in \mathbb{R} \to a \in \mathbb{R}) \to x \in \mathbb{R} \to y \in \mathbb{R} \to a \in \mathbb{R} \to b \in \mathbb{R} \to -a \in \mathbb{R}$
\item $\varphi_{146} :\equiv (x \in \mathbb{R} \to y \in \mathbb{R} \to a \in \mathbb{R} \to b \in \mathbb{R} \to -b \in \mathbb{R}) \to x \in \mathbb{R} \to y \in \mathbb{R} \to a \in \mathbb{R} \to b \in \mathbb{R} \to y + -b \in \mathbb{R}$
\item $\varphi_{147} :\equiv (x \in \mathbb{R} \to y \in \mathbb{R} \to a \in \mathbb{R} \to b \in \mathbb{R} \to -a \in \mathbb{R}) \to x \in \mathbb{R} \to y \in \mathbb{R} \to a \in \mathbb{R} \to b \in \mathbb{R} \to y + -a = -a + y$
\item $\varphi_{148} :\equiv (x \in \mathbb{R} \to y \in \mathbb{R} \to a \in \mathbb{R} \to b \in \mathbb{R} \to -b \in \mathbb{R}) \to x \in \mathbb{R} \to y \in \mathbb{R} \to a \in \mathbb{R} \to b \in \mathbb{R} \to -a + -b \in \mathbb{R}$
\item $\varphi_{149} :\equiv \forall v_{166}\,(y + -a + -b = v_{166} \to v_{166} = y + -a + -b) \to \varphi_{98}$
\item $\varphi_{150} :\equiv (x \in \mathbb{R} \to y \in \mathbb{R} \to a \in \mathbb{R} \to b \in \mathbb{R} \to -b \in \mathbb{R}) \to x \in \mathbb{R} \to \varphi_{92}$
\item $\varphi_{151} :\equiv (x \in \mathbb{R} \to y \in \mathbb{R} \to a \in \mathbb{R} \to b \in \mathbb{R} \to -b \in \mathbb{R}) \to x \in \mathbb{R} \to \varphi_{93}$
\item $\varphi_{152} :\equiv (x \in \mathbb{R} \to y \in \mathbb{R} \to a \in \mathbb{R} \to b \in \mathbb{R} \to y + -b \in \mathbb{R}) \to x \in \mathbb{R} \to y \in \mathbb{R} \to \varphi_{79}$
\item $\varphi_{153} :\equiv (x \in \mathbb{R} \to y \in \mathbb{R} \to a \in \mathbb{R} \to b \in \mathbb{R} \to -a + -b \in \mathbb{R}) \to x \in \mathbb{R} \to y \in \mathbb{R} \to \varphi_{84}$
\item $\varphi_{154} :\equiv \forall v_{156}\,\forall v_{157}\,\forall v_{158}\,(v_{156} \in \mathbb{R} \to v_{157} \in \mathbb{R} \to \varphi_{22}) \to \varphi_{103}$
\item $\varphi_{155} :\equiv (x + y - (a + b) = x - a + (y + -b) \to \varphi_{71}) \to x \in \mathbb{R} \to \varphi_{69}$
\item $\varphi_{156} :\equiv (x \in \mathbb{R} \to y \in \mathbb{R} \to a \in \mathbb{R} \to b \in \mathbb{R} \to a \in \mathbb{R}) \to (x \in \mathbb{R} \to y \in \mathbb{R} \to a \in \mathbb{R} \to b \in \mathbb{R} \to b \in \mathbb{R}) \to \varphi_{87}$
\item $\varphi_{157} :\equiv (x + y - (a + b) = x + -a + (y + -b) \to \varphi_{99}) \to x \in \mathbb{R} \to \varphi_{78}$
\item $\varphi_{158} :\equiv (x + y - (a + b) = x + y + -(a + b) \to \varphi_{106}) \to x \in \mathbb{R} \to y \in \mathbb{R} \to \varphi_{62}$
\item $\varphi_{159} :\equiv (x + y - (a + b) = x + y + (-a + -b) \to \varphi_{118}) \to x \in \mathbb{R} \to y \in \mathbb{R} \to \varphi_{63}$
\item $\varphi_{160} :\equiv (x + y - (a + b) = x + (y + (-a + -b)) \to \varphi_{117}) \to x \in \mathbb{R} \to y \in \mathbb{R} \to \varphi_{64}$
\item $\varphi_{161} :\equiv (x + y - (a + b) = x + (y + -a + -b) \to \varphi_{116}) \to x \in \mathbb{R} \to y \in \mathbb{R} \to \varphi_{65}$
\item $\varphi_{162} :\equiv (x + y - (a + b) = x + (-a + y + -b) \to \varphi_{115}) \to x \in \mathbb{R} \to y \in \mathbb{R} \to \varphi_{66}$
\item $\varphi_{163} :\equiv (x + y - (a + b) = x + (-a + (y + -b)) \to \varphi_{114}) \to x \in \mathbb{R} \to y \in \mathbb{R} \to \varphi_{67}$
\item $\varphi_{164} :\equiv x - a + (y + -b) = x - a + (y + -b) \to (y + -b = y - b \to \varphi_{96}) \to x \in \mathbb{R} \to \varphi_{77}$
\item $\varphi_{165} :\equiv (y \in \mathbb{R} \to -b \in \mathbb{R} \to y + -b \in \mathbb{R}) \to (x \in \mathbb{R} \to y \in \mathbb{R} \to a \in \mathbb{R} \to b \in \mathbb{R} \to y \in \mathbb{R}) \to \varphi_{146}$
\item $\varphi_{166} :\equiv (-a + y + -b = -a + (y + -b) \to \varphi_{131}) \to x \in \mathbb{R} \to y \in \mathbb{R} \to \varphi_{81}$
\item $\varphi_{167} :\equiv (y + (-a + -b) = y + -a + -b \to \varphi_{133}) \to x \in \mathbb{R} \to y \in \mathbb{R} \to \varphi_{83}$
\item $\varphi_{168} :\equiv (y \in \mathbb{R} \to -a \in \mathbb{R} \to y + -a = -a + y) \to (x \in \mathbb{R} \to y \in \mathbb{R} \to a \in \mathbb{R} \to b \in \mathbb{R} \to y \in \mathbb{R}) \to \varphi_{147}$
\item $\varphi_{169} :\equiv x + y + -(a + b) = x + y + -(a + b) \to (-(a + b) = -a + -b \to \varphi_{107}) \to x \in \mathbb{R} \to y \in \mathbb{R} \to \varphi_{73}$
\item $\varphi_{170} :\equiv (-a \in \mathbb{R} \to -b \in \mathbb{R} \to -a + -b \in \mathbb{R}) \to (x \in \mathbb{R} \to y \in \mathbb{R} \to a \in \mathbb{R} \to b \in \mathbb{R} \to -a \in \mathbb{R}) \to \varphi_{148}$
\item $\varphi_{171} :\equiv x + -a + (y + -b) = x + -a + (y + -b) \to (x + -a = x - a \to \varphi_{121}) \to x \in \mathbb{R} \to y \in \mathbb{R} \to \varphi_{72}$
\item $\varphi_{172} :\equiv x + (y + -a + -b) = x + (y + -a + -b) \to (y + -a = -a + y \to \varphi_{132}) \to x \in \mathbb{R} \to y \in \mathbb{R} \to \varphi_{82}$
\item $\varphi_{173} :\equiv (x \in \mathbb{R} \to y \in \mathbb{R} \to a \in \mathbb{R} \to b \in \mathbb{R} \to -a \in \mathbb{R}) \to (x \in \mathbb{R} \to y \in \mathbb{R} \to a \in \mathbb{R} \to b \in \mathbb{R} \to y \in \mathbb{R}) \to \varphi_{150}$
\item $\varphi_{174} :\equiv (x \in \mathbb{R} \to y \in \mathbb{R} \to a \in \mathbb{R} \to b \in \mathbb{R} \to y \in \mathbb{R}) \to (x \in \mathbb{R} \to y \in \mathbb{R} \to a \in \mathbb{R} \to b \in \mathbb{R} \to -a \in \mathbb{R}) \to \varphi_{151}$
\item $\varphi_{175} :\equiv (x \in \mathbb{R} \to y \in \mathbb{R} \to a \in \mathbb{R} \to b \in \mathbb{R} \to x \in \mathbb{R}) \to (x \in \mathbb{R} \to y \in \mathbb{R} \to a \in \mathbb{R} \to b \in \mathbb{R} \to -a \in \mathbb{R}) \to \varphi_{152}$
\item $\varphi_{176} :\equiv (x \in \mathbb{R} \to y \in \mathbb{R} \to a \in \mathbb{R} \to b \in \mathbb{R} \to x \in \mathbb{R}) \to (x \in \mathbb{R} \to y \in \mathbb{R} \to a \in \mathbb{R} \to b \in \mathbb{R} \to y \in \mathbb{R}) \to \varphi_{153}$
\item $\varphi_{177} :\equiv (x \in \mathbb{R} \to y \in \mathbb{R} \to a \in \mathbb{R} \to b \in \mathbb{R} \to y + -a = -a + y) \to \varphi_{172}$
\item $\varphi_{178} :\equiv x + y - (a + b) = x + y + -(a + b) \to (x \in \mathbb{R} \to y \in \mathbb{R} \to \varphi_{73}) \to \varphi_{158}$
\item $\varphi_{179} :\equiv (x \in \mathbb{R} \to \varphi_{92}) \to x + (-a + y + -b) = x + (-a + y + -b) \to \varphi_{166}$
\item $\varphi_{180} :\equiv (x \in \mathbb{R} \to \varphi_{94}) \to x + (y + (-a + -b)) = x + (y + (-a + -b)) \to \varphi_{167}$
\end{abbrevlist}
\end{fullproof}

% 自動生成 (prover/build.js): SubSplit  ノード数 376
\begin{fullproof}{[F16]}{$\forall a\,\forall b\,\forall x\,(a \in \mathbb{R} \to b \in \mathbb{R} \to x \in \mathbb{R} \to a - b = a - x + (x - b))$}{376}\phantomsection\label{tree:SubSplit}
\noindent$\Pi_{1}$:
\begin{ptree}
\AXC{$x = x$\lbl{A5}}
\AXC{$x = y \to x = x \to y = x$\lbl{A6}}
\AXC{$(x = y \to x = x \to y = x) \to x = x \to x = y \to y = x$\lbl{A1}}
\BIC{$x = x \to x = y \to y = x$}
\BIC{$x = y \to y = x$}
\UIC{$\forall y\,(x = y \to y = x)$}
\UIC{$\forall x\,\forall y\,(x = y \to y = x)$}
\end{ptree}
\noindent$\Pi_{2}$:
\begin{ptree}
\AXC{$y = z \to x = y \to x = z$\lbl{A6}}
\AXC{$(y = z \to x = y \to x = z) \to x = y \to y = z \to x = z$\lbl{A1}}
\BIC{$x = y \to y = z \to x = z$}
\UIC{$\forall z\,(x = y \to y = z \to x = z)$}
\UIC{$\forall y\,\forall z\,(x = y \to y = z \to x = z)$}
\UIC{$\forall x\,\forall y\,\forall z\,(x = y \to y = z \to x = z)$}
\AXC{$\forall x\,\forall y\,\forall z\,(x = y \to y = z \to x = z) \to \varphi_{1}$\lbl{A2}}
\BIC{$\varphi_{1}$}
\end{ptree}
\noindent$\Pi_{3}$:
\begin{ptree}
\AXC{$\stackrel{\Pi_{1}}{\vdots}$}\noLine
\UIC{$\forall x\,\forall y\,(x = y \to y = x)$}
\AXC{$\forall x\,\forall y\,(x = y \to y = x) \to \forall y\,(a - b = y \to y = a - b)$\lbl{A2}}
\BIC{$\forall y\,(a - b = y \to y = a - b)$}
\AXC{$\varphi_{12}$\lbl{A2}}
\BIC{$a - b = a + -b \to a + -b = a - b$}
\end{ptree}
\noindent$\Pi_{4}$:
\begin{ptree}
\AXC{$\forall x\,\forall y\,(x - y = x + -y)$\lbl{Def$-$}}
\AXC{$\forall x\,\forall y\,(x - y = x + -y) \to \forall y\,(a - y = a + -y)$\lbl{A2}}
\BIC{$\forall y\,(a - y = a + -y)$}
\end{ptree}
\noindent$\Pi_{5}$:
\begin{ptree}
\AXC{$\stackrel{\Pi_{3}}{\vdots}$}\noLine
\UIC{$a - b = a + -b \to a + -b = a - b$}
\AXC{$\stackrel{\Pi_{4}}{\vdots}$}\noLine
\UIC{$\forall y\,(a - y = a + -y)$}
\AXC{$\forall y\,(a - y = a + -y) \to a - b = a + -b$\lbl{A2}}
\BIC{$a - b = a + -b$}
\AXC{$\varphi_{53}$\lbl{A1}}
\BIC{$(a - b = a + -b \to a + -b = a - b) \to a + -b = a - b$}
\BIC{$a + -b = a - b$}
\end{ptree}
\noindent$\Pi_{6}$:
\begin{ptree}
\AXC{$y = z \to x = y \to x = z$\lbl{A6}}
\AXC{$(y = z \to x = y \to x = z) \to x = y \to y = z \to x = z$\lbl{A1}}
\BIC{$x = y \to y = z \to x = z$}
\UIC{$\forall z\,(x = y \to y = z \to x = z)$}
\UIC{$\forall y\,\forall z\,(x = y \to y = z \to x = z)$}
\UIC{$\forall x\,\forall y\,\forall z\,(x = y \to y = z \to x = z)$}
\AXC{$\forall x\,\forall y\,\forall z\,(x = y \to y = z \to x = z) \to \varphi_{1}$\lbl{A2}}
\BIC{$\varphi_{1}$}
\end{ptree}
\noindent$\Pi_{7}$:
\begin{ptree}
\AXC{$z = w \to a + (-b + 0) = a + z \to a + (-b + 0) = a + w$\lbl{A6}}
\UIC{$\varphi_{3}$}
\UIC{$\forall z\,(\varphi_{3})$}
\AXC{$\forall z\,(\varphi_{3}) \to \varphi_{40}$\lbl{A2}}
\BIC{$\varphi_{40}$}
\AXC{$\varphi_{40} \to -b + 0 = -b \to \varphi_{13}$\lbl{A2}}
\BIC{$-b + 0 = -b \to \varphi_{13}$}
\end{ptree}
\noindent$\Pi_{8}$:
\begin{ptree}
\AXC{$\forall x \in \mathbb{R}\,(-x \in \mathbb{R})$\lbl{(R1)}}
\AXC{$\forall x \in \mathbb{R}\,(-x \in \mathbb{R}) \to b \in \mathbb{R} \to -b \in \mathbb{R}$\lbl{A2}}
\BIC{$b \in \mathbb{R} \to -b \in \mathbb{R}$}
\AXC{$(b \in \mathbb{R} \to -b \in \mathbb{R}) \to \varphi_{89}$\lbl{A1}}
\BIC{$\varphi_{89}$}
\end{ptree}
\noindent$\Pi_{9}$:
\begin{ptree}
\AXC{$a \in \mathbb{R} \to b \in \mathbb{R} \to x \in \mathbb{R} \to b \in \mathbb{R}$\lbl{A1}}
\AXC{$\stackrel{\Pi_{8}}{\vdots}$}\noLine
\UIC{$\varphi_{89}$}
\BIC{$a \in \mathbb{R} \to b \in \mathbb{R} \to x \in \mathbb{R} \to -b \in \mathbb{R}$}
\end{ptree}
\noindent$\Pi_{10}$:
\begin{ptree}
\AXC{$\forall x \in \mathbb{R}\,(x + 0 = x)$\lbl{(R4)}}
\AXC{$\forall x \in \mathbb{R}\,(x + 0 = x) \to -b \in \mathbb{R} \to -b + 0 = -b$\lbl{A2}}
\BIC{$-b \in \mathbb{R} \to -b + 0 = -b$}
\AXC{$(-b \in \mathbb{R} \to -b + 0 = -b) \to \varphi_{95}$\lbl{A1}}
\BIC{$\varphi_{95}$}
\end{ptree}
\noindent$\Pi_{11}$:
\begin{ptree}
\AXC{$\stackrel{\Pi_{9}}{\vdots}$}\noLine
\UIC{$a \in \mathbb{R} \to b \in \mathbb{R} \to x \in \mathbb{R} \to -b \in \mathbb{R}$}
\AXC{$\stackrel{\Pi_{10}}{\vdots}$}\noLine
\UIC{$\varphi_{95}$}
\BIC{$a \in \mathbb{R} \to b \in \mathbb{R} \to x \in \mathbb{R} \to -b + 0 = -b$}
\AXC{$\varphi_{123}$\lbl{A1}}
\BIC{$a + (-b + 0) = a + (-b + 0) \to \varphi_{106}$}
\end{ptree}
\noindent$\Pi_{12}$:
\begin{ptree}
\AXC{$x = x$\lbl{A5}}
\UIC{$\forall x\,(x = x)$}
\AXC{$\forall x\,(x = x) \to a + (-b + 0) = a + (-b + 0)$\lbl{A2}}
\BIC{$a + (-b + 0) = a + (-b + 0)$}
\AXC{$\stackrel{\Pi_{11}}{\vdots}$}\noLine
\UIC{$a + (-b + 0) = a + (-b + 0) \to \varphi_{106}$}
\BIC{$\varphi_{106}$}
\end{ptree}
\noindent$\Pi_{13}$:
\begin{ptree}
\AXC{$y = z \to x = y \to x = z$\lbl{A6}}
\AXC{$(y = z \to x = y \to x = z) \to x = y \to y = z \to x = z$\lbl{A1}}
\BIC{$x = y \to y = z \to x = z$}
\UIC{$\forall z\,(x = y \to y = z \to x = z)$}
\UIC{$\forall y\,\forall z\,(x = y \to y = z \to x = z)$}
\UIC{$\forall x\,\forall y\,\forall z\,(x = y \to y = z \to x = z)$}
\AXC{$\forall x\,\forall y\,\forall z\,(x = y \to y = z \to x = z) \to \varphi_{1}$\lbl{A2}}
\BIC{$\varphi_{1}$}
\end{ptree}
\noindent$\Pi_{14}$:
\begin{ptree}
\AXC{$z = w \to a + (0 + -b) = a + z \to a + (0 + -b) = a + w$\lbl{A6}}
\UIC{$\varphi_{4}$}
\UIC{$\forall z\,(\varphi_{4})$}
\AXC{$\forall z\,(\varphi_{4}) \to \varphi_{42}$\lbl{A2}}
\BIC{$\varphi_{42}$}
\AXC{$\varphi_{42} \to 0 + -b = -b + 0 \to \varphi_{21}$\lbl{A2}}
\BIC{$0 + -b = -b + 0 \to \varphi_{21}$}
\end{ptree}
\noindent$\Pi_{15}$:
\begin{ptree}
\AXC{$\forall x\,\forall y\,(x \in \mathbb{R} \to y \in \mathbb{R} \to x + y = y + x)$\lbl{(R3)}}
\AXC{$\varphi_{67}$\lbl{A2}}
\BIC{$\forall y\,(0 \in \mathbb{R} \to y \in \mathbb{R} \to 0 + y = y + 0)$}
\AXC{$\varphi_{81}$\lbl{A2}}
\BIC{$0 \in \mathbb{R} \to -b \in \mathbb{R} \to 0 + -b = -b + 0$}
\end{ptree}
\noindent$\Pi_{16}$:
\begin{ptree}
\AXC{$0 \in \mathbb{R}$\lbl{(R1)}}
\AXC{$\stackrel{\Pi_{15}}{\vdots}$}\noLine
\UIC{$0 \in \mathbb{R} \to -b \in \mathbb{R} \to 0 + -b = -b + 0$}
\AXC{$(0 \in \mathbb{R} \to -b \in \mathbb{R} \to 0 + -b = -b + 0) \to 0 \in \mathbb{R} \to \varphi_{101}$\lbl{A1}}
\BIC{$0 \in \mathbb{R} \to \varphi_{101}$}
\BIC{$\varphi_{101}$}
\end{ptree}
\noindent$\Pi_{17}$:
\begin{ptree}
\AXC{$\stackrel{\Pi_{9}}{\vdots}$}\noLine
\UIC{$a \in \mathbb{R} \to b \in \mathbb{R} \to x \in \mathbb{R} \to -b \in \mathbb{R}$}
\AXC{$\stackrel{\Pi_{16}}{\vdots}$}\noLine
\UIC{$\varphi_{101}$}
\BIC{$a \in \mathbb{R} \to b \in \mathbb{R} \to x \in \mathbb{R} \to 0 + -b = -b + 0$}
\AXC{$(a \in \mathbb{R} \to b \in \mathbb{R} \to x \in \mathbb{R} \to 0 + -b = -b + 0) \to \varphi_{115}$\lbl{A1}}
\BIC{$\varphi_{115}$}
\end{ptree}
\noindent$\Pi_{18}$:
\begin{ptree}
\AXC{$\stackrel{\Pi_{14}}{\vdots}$}\noLine
\UIC{$0 + -b = -b + 0 \to \varphi_{21}$}
\AXC{$x = x$\lbl{A5}}
\UIC{$\forall x\,(x = x)$}
\AXC{$\forall x\,(x = x) \to a + (0 + -b) = a + (0 + -b)$\lbl{A2}}
\BIC{$a + (0 + -b) = a + (0 + -b)$}
\AXC{$\stackrel{\Pi_{17}}{\vdots}$}\noLine
\UIC{$\varphi_{115}$}
\BIC{$(0 + -b = -b + 0 \to \varphi_{21}) \to \varphi_{31}$}
\BIC{$\varphi_{31}$}
\end{ptree}
\noindent$\Pi_{19}$:
\begin{ptree}
\AXC{$y = z \to x = y \to x = z$\lbl{A6}}
\AXC{$(y = z \to x = y \to x = z) \to x = y \to y = z \to x = z$\lbl{A1}}
\BIC{$x = y \to y = z \to x = z$}
\UIC{$\forall z\,(x = y \to y = z \to x = z)$}
\UIC{$\forall y\,\forall z\,(x = y \to y = z \to x = z)$}
\UIC{$\forall x\,\forall y\,\forall z\,(x = y \to y = z \to x = z)$}
\AXC{$\forall x\,\forall y\,\forall z\,(x = y \to y = z \to x = z) \to \varphi_{1}$\lbl{A2}}
\BIC{$\varphi_{1}$}
\end{ptree}
\noindent$\Pi_{20}$:
\begin{ptree}
\AXC{$\stackrel{\Pi_{19}}{\vdots}$}\noLine
\UIC{$\varphi_{1}$}
\AXC{$\varphi_{1} \to \varphi_{75}$\lbl{A2}}
\BIC{$\varphi_{75}$}
\AXC{$\varphi_{75} \to a - x + (x - b) = a + (x + -x + -b) \to \varphi_{44}$\lbl{A2}}
\BIC{$a - x + (x - b) = a + (x + -x + -b) \to \varphi_{44}$}
\end{ptree}
\noindent$\Pi_{21}$:
\begin{ptree}
\AXC{$z = w \to \varphi_{65}$\lbl{A6}}
\UIC{$\forall w\,(z = w \to \varphi_{65})$}
\UIC{$\forall z\,\forall w\,(z = w \to \varphi_{65})$}
\AXC{$\forall z\,\forall w\,(z = w \to \varphi_{65}) \to \varphi_{97}$\lbl{A2}}
\BIC{$\varphi_{97}$}
\AXC{$\varphi_{97} \to x + -x = 0 \to \varphi_{82}$\lbl{A2}}
\BIC{$x + -x = 0 \to \varphi_{82}$}
\end{ptree}
\noindent$\Pi_{22}$:
\begin{ptree}
\AXC{$\forall x \in \mathbb{R}\,(x + -x = 0)$\lbl{(R5)}}
\AXC{$\forall x \in \mathbb{R}\,(x + -x = 0) \to x \in \mathbb{R} \to x + -x = 0$\lbl{A2}}
\BIC{$x \in \mathbb{R} \to x + -x = 0$}
\AXC{$(x \in \mathbb{R} \to x + -x = 0) \to \varphi_{86}$\lbl{A1}}
\BIC{$\varphi_{86}$}
\end{ptree}
\noindent$\Pi_{23}$:
\begin{ptree}
\AXC{$a \in \mathbb{R} \to b \in \mathbb{R} \to x \in \mathbb{R} \to x \in \mathbb{R}$\lbl{A1}}
\AXC{$\stackrel{\Pi_{22}}{\vdots}$}\noLine
\UIC{$\varphi_{86}$}
\BIC{$a \in \mathbb{R} \to b \in \mathbb{R} \to x \in \mathbb{R} \to x + -x = 0$}
\AXC{$(a \in \mathbb{R} \to b \in \mathbb{R} \to x \in \mathbb{R} \to x + -x = 0) \to \varphi_{125}$\lbl{A1}}
\BIC{$\varphi_{125}$}
\end{ptree}
\noindent$\Pi_{24}$:
\begin{ptree}
\AXC{$\stackrel{\Pi_{21}}{\vdots}$}\noLine
\UIC{$x + -x = 0 \to \varphi_{82}$}
\AXC{$x = x$\lbl{A5}}
\UIC{$\forall x\,(x = x)$}
\AXC{$\varphi_{8}$\lbl{A2}}
\BIC{$a + (x + -x + -b) = a + (x + -x + -b)$}
\AXC{$\stackrel{\Pi_{23}}{\vdots}$}\noLine
\UIC{$\varphi_{125}$}
\BIC{$(x + -x = 0 \to \varphi_{82}) \to \varphi_{54}$}
\BIC{$\varphi_{54}$}
\end{ptree}
\noindent$\Pi_{25}$:
\begin{ptree}
\AXC{$y = z \to x = y \to x = z$\lbl{A6}}
\AXC{$(y = z \to x = y \to x = z) \to x = y \to y = z \to x = z$\lbl{A1}}
\BIC{$x = y \to y = z \to x = z$}
\UIC{$\forall z\,(x = y \to y = z \to x = z)$}
\UIC{$\forall y\,\forall z\,(x = y \to y = z \to x = z)$}
\UIC{$\forall x\,\forall y\,\forall z\,(x = y \to y = z \to x = z)$}
\AXC{$\forall x\,\forall y\,\forall z\,(x = y \to y = z \to x = z) \to \varphi_{1}$\lbl{A2}}
\BIC{$\varphi_{1}$}
\end{ptree}
\noindent$\Pi_{26}$:
\begin{ptree}
\AXC{$\stackrel{\Pi_{25}}{\vdots}$}\noLine
\UIC{$\varphi_{1}$}
\AXC{$\varphi_{1} \to \varphi_{76}$\lbl{A2}}
\BIC{$\varphi_{76}$}
\AXC{$\varphi_{76} \to a - x + (x - b) = a + (x + (-x + -b)) \to \varphi_{83}$\lbl{A2}}
\BIC{$a - x + (x - b) = a + (x + (-x + -b)) \to \varphi_{83}$}
\end{ptree}
\noindent$\Pi_{27}$:
\begin{ptree}
\AXC{$z = w \to \varphi_{22}$\lbl{A6}}
\UIC{$\forall w\,(z = w \to \varphi_{22})$}
\UIC{$\forall z\,\forall w\,(z = w \to \varphi_{22})$}
\AXC{$\forall z\,\forall w\,(z = w \to \varphi_{22}) \to \varphi_{98}$\lbl{A2}}
\BIC{$\varphi_{98}$}
\AXC{$\varphi_{98} \to x + (-x + -b) = x + -x + -b \to \varphi_{92}$\lbl{A2}}
\BIC{$x + (-x + -b) = x + -x + -b \to \varphi_{92}$}
\end{ptree}
\noindent$\Pi_{28}$:
\begin{ptree}
\AXC{$\forall x \in \mathbb{R}\,(-x \in \mathbb{R})$\lbl{(R1)}}
\AXC{$\forall x \in \mathbb{R}\,(-x \in \mathbb{R}) \to x \in \mathbb{R} \to -x \in \mathbb{R}$\lbl{A2}}
\BIC{$x \in \mathbb{R} \to -x \in \mathbb{R}$}
\AXC{$(x \in \mathbb{R} \to -x \in \mathbb{R}) \to \varphi_{90}$\lbl{A1}}
\BIC{$\varphi_{90}$}
\end{ptree}
\noindent$\Pi_{29}$:
\begin{ptree}
\AXC{$a \in \mathbb{R} \to b \in \mathbb{R} \to x \in \mathbb{R} \to x \in \mathbb{R}$\lbl{A1}}
\AXC{$\stackrel{\Pi_{28}}{\vdots}$}\noLine
\UIC{$\varphi_{90}$}
\BIC{$a \in \mathbb{R} \to b \in \mathbb{R} \to x \in \mathbb{R} \to -x \in \mathbb{R}$}
\end{ptree}
\noindent$\Pi_{30}$:
\begin{ptree}
\AXC{$\forall x\,(\varphi_{18})$\lbl{(R2)}}
\AXC{$\forall x\,(\varphi_{18}) \to \varphi_{18}$\lbl{A2}}
\BIC{$\varphi_{18}$}
\AXC{$\varphi_{18} \to \varphi_{50}$\lbl{A2}}
\BIC{$\varphi_{50}$}
\AXC{$\varphi_{50} \to \varphi_{68}$\lbl{A2}}
\BIC{$\varphi_{68}$}
\AXC{$\varphi_{68} \to (a \in \mathbb{R} \to b \in \mathbb{R} \to x \in \mathbb{R} \to x \in \mathbb{R}) \to \varphi_{110}$\lbl{A1}}
\BIC{$(a \in \mathbb{R} \to b \in \mathbb{R} \to x \in \mathbb{R} \to x \in \mathbb{R}) \to \varphi_{110}$}
\end{ptree}
\noindent$\Pi_{31}$:
\begin{ptree}
\AXC{$\stackrel{\Pi_{29}}{\vdots}$}\noLine
\UIC{$a \in \mathbb{R} \to b \in \mathbb{R} \to x \in \mathbb{R} \to -x \in \mathbb{R}$}
\AXC{$a \in \mathbb{R} \to b \in \mathbb{R} \to x \in \mathbb{R} \to x \in \mathbb{R}$\lbl{A1}}
\AXC{$\stackrel{\Pi_{30}}{\vdots}$}\noLine
\UIC{$(a \in \mathbb{R} \to b \in \mathbb{R} \to x \in \mathbb{R} \to x \in \mathbb{R}) \to \varphi_{110}$}
\BIC{$\varphi_{110}$}
\BIC{$(a \in \mathbb{R} \to b \in \mathbb{R} \to x \in \mathbb{R} \to -b \in \mathbb{R}) \to \varphi_{55}$}
\end{ptree}
\noindent$\Pi_{32}$:
\begin{ptree}
\AXC{$\stackrel{\Pi_{9}}{\vdots}$}\noLine
\UIC{$a \in \mathbb{R} \to b \in \mathbb{R} \to x \in \mathbb{R} \to -b \in \mathbb{R}$}
\AXC{$\stackrel{\Pi_{31}}{\vdots}$}\noLine
\UIC{$(a \in \mathbb{R} \to b \in \mathbb{R} \to x \in \mathbb{R} \to -b \in \mathbb{R}) \to \varphi_{55}$}
\BIC{$\varphi_{55}$}
\AXC{$\varphi_{55} \to \varphi_{64} \to \varphi_{56}$\lbl{A1}}
\BIC{$\varphi_{64} \to \varphi_{56}$}
\end{ptree}
\noindent$\Pi_{33}$:
\begin{ptree}
\AXC{$\stackrel{\Pi_{1}}{\vdots}$}\noLine
\UIC{$\forall x\,\forall y\,(x = y \to y = x)$}
\AXC{$\varphi_{30}$\lbl{A2}}
\BIC{$\forall y\,(x + -x + -b = y \to y = x + -x + -b)$}
\AXC{$\forall y\,(x + -x + -b = y \to y = x + -x + -b) \to \varphi_{64}$\lbl{A2}}
\BIC{$\varphi_{64}$}
\AXC{$\stackrel{\Pi_{32}}{\vdots}$}\noLine
\UIC{$\varphi_{64} \to \varphi_{56}$}
\BIC{$\varphi_{56}$}
\AXC{$\varphi_{56} \to \varphi_{131}$\lbl{A1}}
\BIC{$\varphi_{131}$}
\end{ptree}
\noindent$\Pi_{34}$:
\begin{ptree}
\AXC{$\stackrel{\Pi_{27}}{\vdots}$}\noLine
\UIC{$x + (-x + -b) = x + -x + -b \to \varphi_{92}$}
\AXC{$x = x$\lbl{A5}}
\UIC{$\forall x\,(x = x)$}
\AXC{$\varphi_{9}$\lbl{A2}}
\BIC{$a + (x + (-x + -b)) = a + (x + (-x + -b))$}
\AXC{$\stackrel{\Pi_{33}}{\vdots}$}\noLine
\UIC{$\varphi_{131}$}
\BIC{$(x + (-x + -b) = x + -x + -b \to \varphi_{92}) \to a \in \mathbb{R} \to \varphi_{45}$}
\BIC{$a \in \mathbb{R} \to \varphi_{45}$}
\end{ptree}
\noindent$\Pi_{35}$:
\begin{ptree}
\AXC{$y = z \to x = y \to x = z$\lbl{A6}}
\AXC{$(y = z \to x = y \to x = z) \to x = y \to y = z \to x = z$\lbl{A1}}
\BIC{$x = y \to y = z \to x = z$}
\UIC{$\forall z\,(x = y \to y = z \to x = z)$}
\UIC{$\forall y\,\forall z\,(x = y \to y = z \to x = z)$}
\UIC{$\forall x\,\forall y\,\forall z\,(x = y \to y = z \to x = z)$}
\AXC{$\forall x\,\forall y\,\forall z\,(x = y \to y = z \to x = z) \to \varphi_{1}$\lbl{A2}}
\BIC{$\varphi_{1}$}
\end{ptree}
\noindent$\Pi_{36}$:
\begin{ptree}
\AXC{$z = w \to \varphi_{46}$\lbl{A6}}
\UIC{$\forall w\,(z = w \to \varphi_{46})$}
\UIC{$\forall z\,\forall w\,(z = w \to \varphi_{46})$}
\AXC{$\forall z\,\forall w\,(z = w \to \varphi_{46}) \to \varphi_{99}$\lbl{A2}}
\BIC{$\varphi_{99}$}
\AXC{$\varphi_{99} \to -b + -x = -x + -b \to \varphi_{93}$\lbl{A2}}
\BIC{$-b + -x = -x + -b \to \varphi_{93}$}
\end{ptree}
\noindent$\Pi_{37}$:
\begin{ptree}
\AXC{$\forall x\,\forall y\,(x \in \mathbb{R} \to y \in \mathbb{R} \to x + y = y + x)$\lbl{(R3)}}
\AXC{$\varphi_{87}$\lbl{A2}}
\BIC{$\forall y\,(-b \in \mathbb{R} \to y \in \mathbb{R} \to -b + y = y + -b)$}
\AXC{$\varphi_{102}$\lbl{A2}}
\BIC{$-b \in \mathbb{R} \to -x \in \mathbb{R} \to -b + -x = -x + -b$}
\AXC{$\varphi_{127}$\lbl{A1}}
\BIC{$(a \in \mathbb{R} \to b \in \mathbb{R} \to x \in \mathbb{R} \to -b \in \mathbb{R}) \to \varphi_{103}$}
\end{ptree}
\noindent$\Pi_{38}$:
\begin{ptree}
\AXC{$\stackrel{\Pi_{9}}{\vdots}$}\noLine
\UIC{$a \in \mathbb{R} \to b \in \mathbb{R} \to x \in \mathbb{R} \to -b \in \mathbb{R}$}
\AXC{$\stackrel{\Pi_{37}}{\vdots}$}\noLine
\UIC{$(a \in \mathbb{R} \to b \in \mathbb{R} \to x \in \mathbb{R} \to -b \in \mathbb{R}) \to \varphi_{103}$}
\BIC{$\varphi_{103}$}
\end{ptree}
\noindent$\Pi_{39}$:
\begin{ptree}
\AXC{$\stackrel{\Pi_{29}}{\vdots}$}\noLine
\UIC{$a \in \mathbb{R} \to b \in \mathbb{R} \to x \in \mathbb{R} \to -x \in \mathbb{R}$}
\AXC{$\stackrel{\Pi_{38}}{\vdots}$}\noLine
\UIC{$\varphi_{103}$}
\BIC{$a \in \mathbb{R} \to b \in \mathbb{R} \to x \in \mathbb{R} \to -b + -x = -x + -b$}
\AXC{$(a \in \mathbb{R} \to b \in \mathbb{R} \to x \in \mathbb{R} \to -b + -x = -x + -b) \to \varphi_{130}$\lbl{A1}}
\BIC{$\varphi_{130}$}
\end{ptree}
\noindent$\Pi_{40}$:
\begin{ptree}
\AXC{$\stackrel{\Pi_{36}}{\vdots}$}\noLine
\UIC{$-b + -x = -x + -b \to \varphi_{93}$}
\AXC{$x = x$\lbl{A5}}
\UIC{$\forall x\,(x = x)$}
\AXC{$\varphi_{10}$\lbl{A2}}
\BIC{$a + (x + (-b + -x)) = a + (x + (-b + -x))$}
\AXC{$\stackrel{\Pi_{39}}{\vdots}$}\noLine
\UIC{$\varphi_{130}$}
\BIC{$(-b + -x = -x + -b \to \varphi_{93}) \to a \in \mathbb{R} \to b \in \mathbb{R} \to \varphi_{16}$}
\BIC{$a \in \mathbb{R} \to b \in \mathbb{R} \to \varphi_{16}$}
\end{ptree}
\noindent$\Pi_{41}$:
\begin{ptree}
\AXC{$y = z \to x = y \to x = z$\lbl{A6}}
\AXC{$(y = z \to x = y \to x = z) \to x = y \to y = z \to x = z$\lbl{A1}}
\BIC{$x = y \to y = z \to x = z$}
\UIC{$\forall z\,(x = y \to y = z \to x = z)$}
\UIC{$\forall y\,\forall z\,(x = y \to y = z \to x = z)$}
\UIC{$\forall x\,\forall y\,\forall z\,(x = y \to y = z \to x = z)$}
\AXC{$\forall x\,\forall y\,\forall z\,(x = y \to y = z \to x = z) \to \varphi_{1}$\lbl{A2}}
\BIC{$\varphi_{1}$}
\end{ptree}
\noindent$\Pi_{42}$:
\begin{ptree}
\AXC{$z = w \to \varphi_{23}$\lbl{A6}}
\UIC{$\forall w\,(z = w \to \varphi_{23})$}
\UIC{$\forall z\,\forall w\,(z = w \to \varphi_{23})$}
\AXC{$\forall z\,\forall w\,(z = w \to \varphi_{23}) \to \varphi_{100}$\lbl{A2}}
\BIC{$\varphi_{100}$}
\AXC{$\varphi_{100} \to x + -b + -x = x + (-b + -x) \to \varphi_{94}$\lbl{A2}}
\BIC{$x + -b + -x = x + (-b + -x) \to \varphi_{94}$}
\end{ptree}
\noindent$\Pi_{43}$:
\begin{ptree}
\AXC{$\forall x\,(\varphi_{18})$\lbl{(R2)}}
\AXC{$\forall x\,(\varphi_{18}) \to \varphi_{18}$\lbl{A2}}
\BIC{$\varphi_{18}$}
\AXC{$\varphi_{18} \to \varphi_{51}$\lbl{A2}}
\BIC{$\varphi_{51}$}
\AXC{$\varphi_{51} \to \varphi_{69}$\lbl{A2}}
\BIC{$\varphi_{69}$}
\AXC{$\varphi_{69} \to (a \in \mathbb{R} \to b \in \mathbb{R} \to x \in \mathbb{R} \to x \in \mathbb{R}) \to \varphi_{111}$\lbl{A1}}
\BIC{$(a \in \mathbb{R} \to b \in \mathbb{R} \to x \in \mathbb{R} \to x \in \mathbb{R}) \to \varphi_{111}$}
\end{ptree}
\noindent$\Pi_{44}$:
\begin{ptree}
\AXC{$\stackrel{\Pi_{9}}{\vdots}$}\noLine
\UIC{$a \in \mathbb{R} \to b \in \mathbb{R} \to x \in \mathbb{R} \to -b \in \mathbb{R}$}
\AXC{$a \in \mathbb{R} \to b \in \mathbb{R} \to x \in \mathbb{R} \to x \in \mathbb{R}$\lbl{A1}}
\AXC{$\stackrel{\Pi_{43}}{\vdots}$}\noLine
\UIC{$(a \in \mathbb{R} \to b \in \mathbb{R} \to x \in \mathbb{R} \to x \in \mathbb{R}) \to \varphi_{111}$}
\BIC{$\varphi_{111}$}
\BIC{$(a \in \mathbb{R} \to b \in \mathbb{R} \to x \in \mathbb{R} \to -x \in \mathbb{R}) \to \varphi_{57}$}
\end{ptree}
\noindent$\Pi_{45}$:
\begin{ptree}
\AXC{$\stackrel{\Pi_{29}}{\vdots}$}\noLine
\UIC{$a \in \mathbb{R} \to b \in \mathbb{R} \to x \in \mathbb{R} \to -x \in \mathbb{R}$}
\AXC{$\stackrel{\Pi_{44}}{\vdots}$}\noLine
\UIC{$(a \in \mathbb{R} \to b \in \mathbb{R} \to x \in \mathbb{R} \to -x \in \mathbb{R}) \to \varphi_{57}$}
\BIC{$\varphi_{57}$}
\AXC{$\varphi_{57} \to \varphi_{132}$\lbl{A1}}
\BIC{$\varphi_{132}$}
\end{ptree}
\noindent$\Pi_{46}$:
\begin{ptree}
\AXC{$\stackrel{\Pi_{42}}{\vdots}$}\noLine
\UIC{$x + -b + -x = x + (-b + -x) \to \varphi_{94}$}
\AXC{$x = x$\lbl{A5}}
\UIC{$\forall x\,(x = x)$}
\AXC{$\varphi_{11}$\lbl{A2}}
\BIC{$a + (x + -b + -x) = a + (x + -b + -x)$}
\AXC{$\stackrel{\Pi_{45}}{\vdots}$}\noLine
\UIC{$\varphi_{132}$}
\BIC{$(x + -b + -x = x + (-b + -x) \to \varphi_{94}) \to a \in \mathbb{R} \to \varphi_{47}$}
\BIC{$a \in \mathbb{R} \to \varphi_{47}$}
\end{ptree}
\noindent$\Pi_{47}$:
\begin{ptree}
\AXC{$y = z \to x = y \to x = z$\lbl{A6}}
\AXC{$(y = z \to x = y \to x = z) \to x = y \to y = z \to x = z$\lbl{A1}}
\BIC{$x = y \to y = z \to x = z$}
\UIC{$\forall z\,(x = y \to y = z \to x = z)$}
\UIC{$\forall y\,\forall z\,(x = y \to y = z \to x = z)$}
\UIC{$\forall x\,\forall y\,\forall z\,(x = y \to y = z \to x = z)$}
\AXC{$\forall x\,\forall y\,\forall z\,(x = y \to y = z \to x = z) \to \varphi_{1}$\lbl{A2}}
\BIC{$\varphi_{1}$}
\end{ptree}
\noindent$\Pi_{48}$:
\begin{ptree}
\AXC{$\stackrel{\Pi_{47}}{\vdots}$}\noLine
\UIC{$\varphi_{1}$}
\AXC{$\varphi_{1} \to \varphi_{60}$\lbl{A2}}
\BIC{$\varphi_{60}$}
\AXC{$\varphi_{60} \to a - x + (x - b) = a + (x - b + -x) \to \varphi_{72}$\lbl{A2}}
\BIC{$a - x + (x - b) = a + (x - b + -x) \to \varphi_{72}$}
\end{ptree}
\noindent$\Pi_{49}$:
\begin{ptree}
\AXC{$z = w \to \varphi_{48}$\lbl{A6}}
\UIC{$\forall w\,(z = w \to \varphi_{48})$}
\UIC{$\forall z\,\forall w\,(z = w \to \varphi_{48})$}
\AXC{$\forall z\,\forall w\,(z = w \to \varphi_{48}) \to \varphi_{79}$\lbl{A2}}
\BIC{$\varphi_{79}$}
\AXC{$\varphi_{79} \to x - b = x + -b \to \varphi_{73}$\lbl{A2}}
\BIC{$x - b = x + -b \to \varphi_{73}$}
\end{ptree}
\noindent$\Pi_{50}$:
\begin{ptree}
\AXC{$\forall x\,\forall y\,(x - y = x + -y)$\lbl{Def$-$}}
\AXC{$\forall x\,\forall y\,(x - y = x + -y) \to \forall y\,(x - y = x + -y)$\lbl{A2}}
\BIC{$\forall y\,(x - y = x + -y)$}
\end{ptree}
\noindent$\Pi_{51}$:
\begin{ptree}
\AXC{$\stackrel{\Pi_{50}}{\vdots}$}\noLine
\UIC{$\forall y\,(x - y = x + -y)$}
\AXC{$\forall y\,(x - y = x + -y) \to x - b = x + -b$\lbl{A2}}
\BIC{$x - b = x + -b$}
\AXC{$x - b = x + -b \to \varphi_{124}$\lbl{A1}}
\BIC{$\varphi_{124}$}
\end{ptree}
\noindent$\Pi_{52}$:
\begin{ptree}
\AXC{$\stackrel{\Pi_{49}}{\vdots}$}\noLine
\UIC{$x - b = x + -b \to \varphi_{73}$}
\AXC{$x = x$\lbl{A5}}
\UIC{$\forall x\,(x = x)$}
\AXC{$\forall x\,(x = x) \to a + (x - b + -x) = a + (x - b + -x)$\lbl{A2}}
\BIC{$a + (x - b + -x) = a + (x - b + -x)$}
\AXC{$\stackrel{\Pi_{51}}{\vdots}$}\noLine
\UIC{$\varphi_{124}$}
\BIC{$(x - b = x + -b \to \varphi_{73}) \to a \in \mathbb{R} \to \varphi_{34}$}
\BIC{$a \in \mathbb{R} \to \varphi_{34}$}
\end{ptree}
\noindent$\Pi_{53}$:
\begin{ptree}
\AXC{$y = z \to x = y \to x = z$\lbl{A6}}
\AXC{$(y = z \to x = y \to x = z) \to x = y \to y = z \to x = z$\lbl{A1}}
\BIC{$x = y \to y = z \to x = z$}
\UIC{$\forall z\,(x = y \to y = z \to x = z)$}
\UIC{$\forall y\,\forall z\,(x = y \to y = z \to x = z)$}
\UIC{$\forall x\,\forall y\,\forall z\,(x = y \to y = z \to x = z)$}
\AXC{$\forall x\,\forall y\,\forall z\,(x = y \to y = z \to x = z) \to \varphi_{1}$\lbl{A2}}
\BIC{$\varphi_{1}$}
\end{ptree}
\noindent$\Pi_{54}$:
\begin{ptree}
\AXC{$\stackrel{\Pi_{53}}{\vdots}$}\noLine
\UIC{$\varphi_{1}$}
\AXC{$\varphi_{1} \to \varphi_{61}$\lbl{A2}}
\BIC{$\varphi_{61}$}
\AXC{$\varphi_{61} \to a - x + (x - b) = a + (-x + (x - b)) \to \varphi_{58}$\lbl{A2}}
\BIC{$a - x + (x - b) = a + (-x + (x - b)) \to \varphi_{58}$}
\end{ptree}
\noindent$\Pi_{55}$:
\begin{ptree}
\AXC{$z = w \to \varphi_{5}$\lbl{A6}}
\UIC{$\forall w\,(z = w \to \varphi_{5})$}
\UIC{$\forall z\,\forall w\,(z = w \to \varphi_{5})$}
\AXC{$\forall z\,\forall w\,(z = w \to \varphi_{5}) \to \varphi_{80}$\lbl{A2}}
\BIC{$\varphi_{80}$}
\AXC{$\varphi_{80} \to -x + (x - b) = x - b + -x \to \varphi_{66}$\lbl{A2}}
\BIC{$-x + (x - b) = x - b + -x \to \varphi_{66}$}
\end{ptree}
\noindent$\Pi_{56}$:
\begin{ptree}
\AXC{$a \in \mathbb{R} \to b \in \mathbb{R} \to x \in \mathbb{R} \to x \in \mathbb{R}$\lbl{A1}}
\AXC{$\forall x\,\forall y\,(x \in \mathbb{R} \to y \in \mathbb{R} \to x - y \in \mathbb{R})$\lbl{[F1]}}
\AXC{$\varphi_{71}$\lbl{A2}}
\BIC{$\forall y\,(x \in \mathbb{R} \to y \in \mathbb{R} \to x - y \in \mathbb{R})$}
\AXC{$\varphi_{62}$\lbl{A2}}
\BIC{$x \in \mathbb{R} \to b \in \mathbb{R} \to x - b \in \mathbb{R}$}
\AXC{$\varphi_{113}$\lbl{A1}}
\BIC{$(a \in \mathbb{R} \to b \in \mathbb{R} \to x \in \mathbb{R} \to x \in \mathbb{R}) \to \varphi_{91}$}
\BIC{$\varphi_{91}$}
\end{ptree}
\noindent$\Pi_{57}$:
\begin{ptree}
\AXC{$a \in \mathbb{R} \to b \in \mathbb{R} \to x \in \mathbb{R} \to b \in \mathbb{R}$\lbl{A1}}
\AXC{$\stackrel{\Pi_{56}}{\vdots}$}\noLine
\UIC{$\varphi_{91}$}
\BIC{$a \in \mathbb{R} \to b \in \mathbb{R} \to x \in \mathbb{R} \to x - b \in \mathbb{R}$}
\end{ptree}
\noindent$\Pi_{58}$:
\begin{ptree}
\AXC{$\forall x\,\forall y\,(x \in \mathbb{R} \to y \in \mathbb{R} \to x + y = y + x)$\lbl{(R3)}}
\AXC{$\varphi_{88}$\lbl{A2}}
\BIC{$\forall y\,(-x \in \mathbb{R} \to y \in \mathbb{R} \to -x + y = y + -x)$}
\AXC{$\varphi_{104}$\lbl{A2}}
\BIC{$-x \in \mathbb{R} \to x - b \in \mathbb{R} \to -x + (x - b) = x - b + -x$}
\AXC{$\varphi_{128}$\lbl{A1}}
\BIC{$(a \in \mathbb{R} \to b \in \mathbb{R} \to x \in \mathbb{R} \to -x \in \mathbb{R}) \to \varphi_{105}$}
\end{ptree}
\noindent$\Pi_{59}$:
\begin{ptree}
\AXC{$\stackrel{\Pi_{29}}{\vdots}$}\noLine
\UIC{$a \in \mathbb{R} \to b \in \mathbb{R} \to x \in \mathbb{R} \to -x \in \mathbb{R}$}
\AXC{$\stackrel{\Pi_{58}}{\vdots}$}\noLine
\UIC{$(a \in \mathbb{R} \to b \in \mathbb{R} \to x \in \mathbb{R} \to -x \in \mathbb{R}) \to \varphi_{105}$}
\BIC{$\varphi_{105}$}
\end{ptree}
\noindent$\Pi_{60}$:
\begin{ptree}
\AXC{$x = x$\lbl{A5}}
\UIC{$\forall x\,(x = x)$}
\AXC{$\varphi_{2}$\lbl{A2}}
\BIC{$a + (-x + (x - b)) = a + (-x + (x - b))$}
\AXC{$\stackrel{\Pi_{57}}{\vdots}$}\noLine
\UIC{$a \in \mathbb{R} \to b \in \mathbb{R} \to x \in \mathbb{R} \to x - b \in \mathbb{R}$}
\AXC{$\stackrel{\Pi_{59}}{\vdots}$}\noLine
\UIC{$\varphi_{105}$}
\BIC{$a \in \mathbb{R} \to b \in \mathbb{R} \to x \in \mathbb{R} \to -x + (x - b) = x - b + -x$}
\AXC{$\varphi_{133}$\lbl{A1}}
\BIC{$\varphi_{126}$}
\BIC{$(-x + (x - b) = x - b + -x \to \varphi_{66}) \to a \in \mathbb{R} \to \varphi_{24}$}
\end{ptree}
\noindent$\Pi_{61}$:
\begin{ptree}
\AXC{$\stackrel{\Pi_{55}}{\vdots}$}\noLine
\UIC{$-x + (x - b) = x - b + -x \to \varphi_{66}$}
\AXC{$\stackrel{\Pi_{60}}{\vdots}$}\noLine
\UIC{$(-x + (x - b) = x - b + -x \to \varphi_{66}) \to a \in \mathbb{R} \to \varphi_{24}$}
\BIC{$a \in \mathbb{R} \to \varphi_{24}$}
\end{ptree}
\noindent$\Pi_{62}$:
\begin{ptree}
\AXC{$y = z \to x = y \to x = z$\lbl{A6}}
\AXC{$(y = z \to x = y \to x = z) \to x = y \to y = z \to x = z$\lbl{A1}}
\BIC{$x = y \to y = z \to x = z$}
\UIC{$\forall z\,(x = y \to y = z \to x = z)$}
\UIC{$\forall y\,\forall z\,(x = y \to y = z \to x = z)$}
\UIC{$\forall x\,\forall y\,\forall z\,(x = y \to y = z \to x = z)$}
\AXC{$\forall x\,\forall y\,\forall z\,(x = y \to y = z \to x = z) \to \varphi_{1}$\lbl{A2}}
\BIC{$\varphi_{1}$}
\end{ptree}
\noindent$\Pi_{63}$:
\begin{ptree}
\AXC{$\stackrel{\Pi_{62}}{\vdots}$}\noLine
\UIC{$\varphi_{1}$}
\AXC{$\varphi_{1} \to \varphi_{63}$\lbl{A2}}
\BIC{$\varphi_{63}$}
\AXC{$\varphi_{63} \to a - x + (x - b) = a + -x + (x - b) \to \varphi_{59}$\lbl{A2}}
\BIC{$a - x + (x - b) = a + -x + (x - b) \to \varphi_{59}$}
\end{ptree}
\noindent$\Pi_{64}$:
\begin{ptree}
\AXC{$\forall x\,(\varphi_{18})$\lbl{(R2)}}
\AXC{$\forall x\,(\varphi_{18}) \to \varphi_{19}$\lbl{A2}}
\BIC{$\varphi_{19}$}
\AXC{$\varphi_{19} \to \varphi_{52}$\lbl{A2}}
\BIC{$\varphi_{52}$}
\AXC{$\varphi_{52} \to a \in \mathbb{R} \to \varphi_{49}$\lbl{A2}}
\BIC{$a \in \mathbb{R} \to \varphi_{49}$}
\AXC{$\varphi_{134}$\lbl{A1}}
\BIC{$(a \in \mathbb{R} \to b \in \mathbb{R} \to x \in \mathbb{R} \to a \in \mathbb{R}) \to \varphi_{112}$}
\end{ptree}
\noindent$\Pi_{65}$:
\begin{ptree}
\AXC{$\stackrel{\Pi_{29}}{\vdots}$}\noLine
\UIC{$a \in \mathbb{R} \to b \in \mathbb{R} \to x \in \mathbb{R} \to -x \in \mathbb{R}$}
\AXC{$a \in \mathbb{R} \to b \in \mathbb{R} \to x \in \mathbb{R} \to a \in \mathbb{R}$\lbl{A1}}
\AXC{$\stackrel{\Pi_{64}}{\vdots}$}\noLine
\UIC{$(a \in \mathbb{R} \to b \in \mathbb{R} \to x \in \mathbb{R} \to a \in \mathbb{R}) \to \varphi_{112}$}
\BIC{$\varphi_{112}$}
\BIC{$(a \in \mathbb{R} \to b \in \mathbb{R} \to x \in \mathbb{R} \to x - b \in \mathbb{R}) \to a \in \mathbb{R} \to \varphi_{25}$}
\end{ptree}
\noindent$\Pi_{66}$:
\begin{ptree}
\AXC{$\stackrel{\Pi_{57}}{\vdots}$}\noLine
\UIC{$a \in \mathbb{R} \to b \in \mathbb{R} \to x \in \mathbb{R} \to x - b \in \mathbb{R}$}
\AXC{$\stackrel{\Pi_{65}}{\vdots}$}\noLine
\UIC{$(a \in \mathbb{R} \to b \in \mathbb{R} \to x \in \mathbb{R} \to x - b \in \mathbb{R}) \to a \in \mathbb{R} \to \varphi_{25}$}
\BIC{$a \in \mathbb{R} \to \varphi_{25}$}
\end{ptree}
\noindent$\Pi_{67}$:
\begin{ptree}
\AXC{$z = w \to \varphi_{6}$\lbl{A6}}
\UIC{$\forall w\,(z = w \to \varphi_{6})$}
\UIC{$\forall z\,\forall w\,(z = w \to \varphi_{6})$}
\AXC{$\forall z\,\forall w\,(z = w \to \varphi_{6}) \to \varphi_{43}$\lbl{A2}}
\BIC{$\varphi_{43}$}
\AXC{$\varphi_{43} \to a - x = a + -x \to \varphi_{35}$\lbl{A2}}
\BIC{$a - x = a + -x \to \varphi_{35}$}
\end{ptree}
\noindent$\Pi_{68}$:
\begin{ptree}
\AXC{$\forall x\,\forall y\,(x - y = x + -y)$\lbl{Def$-$}}
\AXC{$\forall x\,\forall y\,(x - y = x + -y) \to \forall y\,(a - y = a + -y)$\lbl{A2}}
\BIC{$\forall y\,(a - y = a + -y)$}
\end{ptree}
\noindent$\Pi_{69}$:
\begin{ptree}
\AXC{$\stackrel{\Pi_{68}}{\vdots}$}\noLine
\UIC{$\forall y\,(a - y = a + -y)$}
\AXC{$\forall y\,(a - y = a + -y) \to a - x = a + -x$\lbl{A2}}
\BIC{$a - x = a + -x$}
\AXC{$a - x = a + -x \to \varphi_{116}$\lbl{A1}}
\BIC{$\varphi_{116}$}
\end{ptree}
\noindent$\Pi_{70}$:
\begin{ptree}
\AXC{$\stackrel{\Pi_{67}}{\vdots}$}\noLine
\UIC{$a - x = a + -x \to \varphi_{35}$}
\AXC{$x = x$\lbl{A5}}
\UIC{$\forall x\,(x = x)$}
\AXC{$\forall x\,(x = x) \to a - x + (x - b) = a - x + (x - b)$\lbl{A2}}
\BIC{$a - x + (x - b) = a - x + (x - b)$}
\AXC{$\stackrel{\Pi_{69}}{\vdots}$}\noLine
\UIC{$\varphi_{116}$}
\BIC{$(a - x = a + -x \to \varphi_{35}) \to \varphi_{37}$}
\BIC{$\varphi_{37}$}
\end{ptree}
\noindent$\Pi_{71}$:
\begin{ptree}
\AXC{$\stackrel{\Pi_{63}}{\vdots}$}\noLine
\UIC{$a - x + (x - b) = a + -x + (x - b) \to \varphi_{59}$}
\AXC{$\stackrel{\Pi_{66}}{\vdots}$}\noLine
\UIC{$a \in \mathbb{R} \to \varphi_{25}$}
\AXC{$\stackrel{\Pi_{70}}{\vdots}$}\noLine
\UIC{$\varphi_{37}$}
\AXC{$\varphi_{37} \to (a \in \mathbb{R} \to \varphi_{25}) \to \varphi_{114}$\lbl{A1}}
\BIC{$(a \in \mathbb{R} \to \varphi_{25}) \to \varphi_{114}$}
\BIC{$\varphi_{114}$}
\BIC{$a \in \mathbb{R} \to \varphi_{14}$}
\end{ptree}
\noindent$\Pi_{72}$:
\begin{ptree}
\AXC{$\stackrel{\Pi_{54}}{\vdots}$}\noLine
\UIC{$a - x + (x - b) = a + (-x + (x - b)) \to \varphi_{58}$}
\AXC{$\stackrel{\Pi_{61}}{\vdots}$}\noLine
\UIC{$a \in \mathbb{R} \to \varphi_{24}$}
\AXC{$\stackrel{\Pi_{71}}{\vdots}$}\noLine
\UIC{$a \in \mathbb{R} \to \varphi_{14}$}
\AXC{$(a \in \mathbb{R} \to \varphi_{14}) \to \varphi_{129}$\lbl{A1}}
\BIC{$\varphi_{129}$}
\BIC{$(a - x + (x - b) = a + (-x + (x - b)) \to \varphi_{58}) \to \varphi_{38}$}
\BIC{$\varphi_{38}$}
\end{ptree}
\noindent$\Pi_{73}$:
\begin{ptree}
\AXC{$\stackrel{\Pi_{48}}{\vdots}$}\noLine
\UIC{$a - x + (x - b) = a + (x - b + -x) \to \varphi_{72}$}
\AXC{$\stackrel{\Pi_{52}}{\vdots}$}\noLine
\UIC{$a \in \mathbb{R} \to \varphi_{34}$}
\AXC{$\stackrel{\Pi_{72}}{\vdots}$}\noLine
\UIC{$\varphi_{38}$}
\AXC{$\varphi_{38} \to (a \in \mathbb{R} \to \varphi_{34}) \to \varphi_{117}$\lbl{A1}}
\BIC{$(a \in \mathbb{R} \to \varphi_{34}) \to \varphi_{117}$}
\BIC{$\varphi_{117}$}
\BIC{$a \in \mathbb{R} \to \varphi_{26}$}
\end{ptree}
\noindent$\Pi_{74}$:
\begin{ptree}
\AXC{$\stackrel{\Pi_{46}}{\vdots}$}\noLine
\UIC{$a \in \mathbb{R} \to \varphi_{47}$}
\AXC{$\stackrel{\Pi_{73}}{\vdots}$}\noLine
\UIC{$a \in \mathbb{R} \to \varphi_{26}$}
\AXC{$(a \in \mathbb{R} \to \varphi_{26}) \to (a \in \mathbb{R} \to \varphi_{47}) \to \varphi_{118}$\lbl{A1}}
\BIC{$(a \in \mathbb{R} \to \varphi_{47}) \to \varphi_{118}$}
\BIC{$\varphi_{118}$}
\end{ptree}
\noindent$\Pi_{75}$:
\begin{ptree}
\AXC{$\stackrel{\Pi_{41}}{\vdots}$}\noLine
\UIC{$\varphi_{1}$}
\AXC{$\varphi_{1} \to \varphi_{78}$\lbl{A2}}
\BIC{$\varphi_{78}$}
\AXC{$\varphi_{78} \to a - x + (x - b) = a + (x + -b + -x) \to \varphi_{85}$\lbl{A2}}
\BIC{$a - x + (x - b) = a + (x + -b + -x) \to \varphi_{85}$}
\AXC{$\stackrel{\Pi_{74}}{\vdots}$}\noLine
\UIC{$\varphi_{118}$}
\BIC{$a \in \mathbb{R} \to \varphi_{27}$}
\end{ptree}
\noindent$\Pi_{76}$:
\begin{ptree}
\AXC{$\stackrel{\Pi_{40}}{\vdots}$}\noLine
\UIC{$a \in \mathbb{R} \to b \in \mathbb{R} \to \varphi_{16}$}
\AXC{$\stackrel{\Pi_{75}}{\vdots}$}\noLine
\UIC{$a \in \mathbb{R} \to \varphi_{27}$}
\AXC{$(a \in \mathbb{R} \to \varphi_{27}) \to (a \in \mathbb{R} \to b \in \mathbb{R} \to \varphi_{16}) \to \varphi_{119}$\lbl{A1}}
\BIC{$(a \in \mathbb{R} \to b \in \mathbb{R} \to \varphi_{16}) \to \varphi_{119}$}
\BIC{$\varphi_{119}$}
\end{ptree}
\noindent$\Pi_{77}$:
\begin{ptree}
\AXC{$\stackrel{\Pi_{35}}{\vdots}$}\noLine
\UIC{$\varphi_{1}$}
\AXC{$\varphi_{1} \to \varphi_{77}$\lbl{A2}}
\BIC{$\varphi_{77}$}
\AXC{$\varphi_{77} \to a - x + (x - b) = a + (x + (-b + -x)) \to \varphi_{84}$\lbl{A2}}
\BIC{$a - x + (x - b) = a + (x + (-b + -x)) \to \varphi_{84}$}
\AXC{$\stackrel{\Pi_{76}}{\vdots}$}\noLine
\UIC{$\varphi_{119}$}
\BIC{$a \in \mathbb{R} \to \varphi_{28}$}
\end{ptree}
\noindent$\Pi_{78}$:
\begin{ptree}
\AXC{$\stackrel{\Pi_{34}}{\vdots}$}\noLine
\UIC{$a \in \mathbb{R} \to \varphi_{45}$}
\AXC{$\stackrel{\Pi_{77}}{\vdots}$}\noLine
\UIC{$a \in \mathbb{R} \to \varphi_{28}$}
\AXC{$(a \in \mathbb{R} \to \varphi_{28}) \to (a \in \mathbb{R} \to \varphi_{45}) \to \varphi_{120}$\lbl{A1}}
\BIC{$(a \in \mathbb{R} \to \varphi_{45}) \to \varphi_{120}$}
\BIC{$\varphi_{120}$}
\end{ptree}
\noindent$\Pi_{79}$:
\begin{ptree}
\AXC{$\stackrel{\Pi_{24}}{\vdots}$}\noLine
\UIC{$\varphi_{54}$}
\AXC{$\stackrel{\Pi_{26}}{\vdots}$}\noLine
\UIC{$a - x + (x - b) = a + (x + (-x + -b)) \to \varphi_{83}$}
\AXC{$\stackrel{\Pi_{78}}{\vdots}$}\noLine
\UIC{$\varphi_{120}$}
\BIC{$a \in \mathbb{R} \to \varphi_{29}$}
\AXC{$\varphi_{135}$\lbl{A1}}
\BIC{$\varphi_{54} \to (a - x + (x - b) = a + (x + -x + -b) \to \varphi_{44}) \to \varphi_{32}$}
\BIC{$(a - x + (x - b) = a + (x + -x + -b) \to \varphi_{44}) \to \varphi_{32}$}
\end{ptree}
\noindent$\Pi_{80}$:
\begin{ptree}
\AXC{$\stackrel{\Pi_{20}}{\vdots}$}\noLine
\UIC{$a - x + (x - b) = a + (x + -x + -b) \to \varphi_{44}$}
\AXC{$\stackrel{\Pi_{79}}{\vdots}$}\noLine
\UIC{$(a - x + (x - b) = a + (x + -x + -b) \to \varphi_{44}) \to \varphi_{32}$}
\BIC{$\varphi_{32}$}
\end{ptree}
\noindent$\Pi_{81}$:
\begin{ptree}
\AXC{$\stackrel{\Pi_{18}}{\vdots}$}\noLine
\UIC{$\varphi_{31}$}
\AXC{$\stackrel{\Pi_{80}}{\vdots}$}\noLine
\UIC{$\varphi_{32}$}
\AXC{$\varphi_{32} \to \varphi_{31} \to (a - x + (x - b) = a + (0 + -b) \to \varphi_{20}) \to \varphi_{33}$\lbl{A1}}
\BIC{$\varphi_{31} \to (a - x + (x - b) = a + (0 + -b) \to \varphi_{20}) \to \varphi_{33}$}
\BIC{$(a - x + (x - b) = a + (0 + -b) \to \varphi_{20}) \to \varphi_{33}$}
\end{ptree}
\noindent$\Pi_{82}$:
\begin{ptree}
\AXC{$\stackrel{\Pi_{13}}{\vdots}$}\noLine
\UIC{$\varphi_{1}$}
\AXC{$\varphi_{1} \to \varphi_{41}$\lbl{A2}}
\BIC{$\varphi_{41}$}
\AXC{$\varphi_{41} \to a - x + (x - b) = a + (0 + -b) \to \varphi_{20}$\lbl{A2}}
\BIC{$a - x + (x - b) = a + (0 + -b) \to \varphi_{20}$}
\AXC{$\stackrel{\Pi_{81}}{\vdots}$}\noLine
\UIC{$(a - x + (x - b) = a + (0 + -b) \to \varphi_{20}) \to \varphi_{33}$}
\BIC{$\varphi_{33}$}
\end{ptree}
\noindent$\Pi_{83}$:
\begin{ptree}
\AXC{$\stackrel{\Pi_{6}}{\vdots}$}\noLine
\UIC{$\varphi_{1}$}
\AXC{$\varphi_{1} \to \varphi_{39}$\lbl{A2}}
\BIC{$\varphi_{39}$}
\AXC{$\varphi_{39} \to \varphi_{74}$\lbl{A2}}
\BIC{$\varphi_{74}$}
\AXC{$\stackrel{\Pi_{7}}{\vdots}$}\noLine
\UIC{$-b + 0 = -b \to \varphi_{13}$}
\AXC{$\stackrel{\Pi_{12}}{\vdots}$}\noLine
\UIC{$\varphi_{106}$}
\BIC{$a \in \mathbb{R} \to b \in \mathbb{R} \to x \in \mathbb{R} \to a + (-b + 0) = a + -b$}
\AXC{$\stackrel{\Pi_{82}}{\vdots}$}\noLine
\UIC{$\varphi_{33}$}
\AXC{$\varphi_{33} \to \varphi_{122}$\lbl{A1}}
\BIC{$\varphi_{122}$}
\BIC{$\varphi_{108}$}
\BIC{$a \in \mathbb{R} \to b \in \mathbb{R} \to x \in \mathbb{R} \to a - x + (x - b) = a + -b$}
\end{ptree}
\noindent$\Pi_{84}$:
\begin{ptree}
\AXC{$\stackrel{\Pi_{5}}{\vdots}$}\noLine
\UIC{$a + -b = a - b$}
\AXC{$\stackrel{\Pi_{83}}{\vdots}$}\noLine
\UIC{$a \in \mathbb{R} \to b \in \mathbb{R} \to x \in \mathbb{R} \to a - x + (x - b) = a + -b$}
\AXC{$\varphi_{121}$\lbl{A1}}
\BIC{$\varphi_{107}$}
\BIC{$\varphi_{36} \to a \in \mathbb{R} \to b \in \mathbb{R} \to x \in \mathbb{R} \to a - x + (x - b) = a - b$}
\end{ptree}
\noindent$\Pi_{85}$:
\begin{ptree}
\AXC{$\stackrel{\Pi_{2}}{\vdots}$}\noLine
\UIC{$\varphi_{1}$}
\AXC{$\varphi_{1} \to \varphi_{17}$\lbl{A2}}
\BIC{$\varphi_{17}$}
\AXC{$\varphi_{17} \to \varphi_{36}$\lbl{A2}}
\BIC{$\varphi_{36}$}
\AXC{$\stackrel{\Pi_{84}}{\vdots}$}\noLine
\UIC{$\varphi_{36} \to a \in \mathbb{R} \to b \in \mathbb{R} \to x \in \mathbb{R} \to a - x + (x - b) = a - b$}
\BIC{$a \in \mathbb{R} \to b \in \mathbb{R} \to x \in \mathbb{R} \to a - x + (x - b) = a - b$}
\AXC{$\varphi_{109}$\lbl{A1}}
\BIC{$\varphi_{96}$}
\end{ptree}
\noindent 主木:
\begin{ptree}
\AXC{$\stackrel{\Pi_{1}}{\vdots}$}\noLine
\UIC{$\forall x\,\forall y\,(x = y \to y = x)$}
\AXC{$\varphi_{7}$\lbl{A2}}
\BIC{$\forall y\,(a - x + (x - b) = y \to y = a - x + (x - b))$}
\AXC{$\varphi_{70}$\lbl{A2}}
\BIC{$a - x + (x - b) = a - b \to a - b = a - x + (x - b)$}
\AXC{$\stackrel{\Pi_{85}}{\vdots}$}\noLine
\UIC{$\varphi_{96}$}
\BIC{$a \in \mathbb{R} \to b \in \mathbb{R} \to x \in \mathbb{R} \to a - b = a - x + (x - b)$}
\UIC{$\varphi_{15}$}
\UIC{$\forall b\,(\varphi_{15})$}
\UIC{$\forall a\,\forall b\,(\varphi_{15})$}
\end{ptree}
\begin{abbrevlist}
\item $\varphi_{1} :\equiv \forall y\,\forall z\,(a - x + (x - b) = y \to y = z \to a - x + (x - b) = z)$
\item $\varphi_{2} :\equiv \forall x\,(x = x) \to a + (-x + (x - b)) = a + (-x + (x - b))$
\item $\varphi_{3} :\equiv \forall w\,(z = w \to a + (-b + 0) = a + z \to a + (-b + 0) = a + w)$
\item $\varphi_{4} :\equiv \forall w\,(z = w \to a + (0 + -b) = a + z \to a + (0 + -b) = a + w)$
\item $\varphi_{5} :\equiv a + (-x + (x - b)) = a + z \to a + (-x + (x - b)) = a + w$
\item $\varphi_{6} :\equiv a - x + (x - b) = z + (x - b) \to a - x + (x - b) = w + (x - b)$
\item $\varphi_{7} :\equiv \forall x\,\forall y\,(x = y \to y = x) \to \forall y\,(a - x + (x - b) = y \to y = a - x + (x - b))$
\item $\varphi_{8} :\equiv \forall x\,(x = x) \to a + (x + -x + -b) = a + (x + -x + -b)$
\item $\varphi_{9} :\equiv \forall x\,(x = x) \to a + (x + (-x + -b)) = a + (x + (-x + -b))$
\item $\varphi_{10} :\equiv \forall x\,(x = x) \to a + (x + (-b + -x)) = a + (x + (-b + -x))$
\item $\varphi_{11} :\equiv \forall x\,(x = x) \to a + (x + -b + -x) = a + (x + -b + -x)$
\item $\varphi_{12} :\equiv \forall y\,(a - b = y \to y = a - b) \to a - b = a + -b \to a + -b = a - b$
\item $\varphi_{13} :\equiv a + (-b + 0) = a + (-b + 0) \to a + (-b + 0) = a + -b$
\item $\varphi_{14} :\equiv b \in \mathbb{R} \to x \in \mathbb{R} \to a - x + (x - b) = a + (-x + (x - b))$
\item $\varphi_{15} :\equiv \forall x\,(a \in \mathbb{R} \to b \in \mathbb{R} \to x \in \mathbb{R} \to a - b = a - x + (x - b))$
\item $\varphi_{16} :\equiv x \in \mathbb{R} \to a + (x + (-b + -x)) = a + (x + (-x + -b))$
\item $\varphi_{17} :\equiv \forall z\,(a - x + (x - b) = a + -b \to a + -b = z \to a - x + (x - b) = z)$
\item $\varphi_{18} :\equiv \forall y\,\forall z\,(x \in \mathbb{R} \to y \in \mathbb{R} \to z \in \mathbb{R} \to x + y + z = x + (y + z))$
\item $\varphi_{19} :\equiv \forall y\,\forall z\,(a \in \mathbb{R} \to y \in \mathbb{R} \to z \in \mathbb{R} \to a + y + z = a + (y + z))$
\item $\varphi_{20} :\equiv a + (0 + -b) = a + (-b + 0) \to a - x + (x - b) = a + (-b + 0)$
\item $\varphi_{21} :\equiv a + (0 + -b) = a + (0 + -b) \to a + (0 + -b) = a + (-b + 0)$
\item $\varphi_{22} :\equiv a + (x + (-x + -b)) = a + z \to a + (x + (-x + -b)) = a + w$
\item $\varphi_{23} :\equiv a + (x + -b + -x) = a + z \to a + (x + -b + -x) = a + w$
\item $\varphi_{24} :\equiv b \in \mathbb{R} \to x \in \mathbb{R} \to a + (-x + (x - b)) = a + (x - b + -x)$
\item $\varphi_{25} :\equiv b \in \mathbb{R} \to x \in \mathbb{R} \to a + -x + (x - b) = a + (-x + (x - b))$
\item $\varphi_{26} :\equiv b \in \mathbb{R} \to x \in \mathbb{R} \to a - x + (x - b) = a + (x + -b + -x)$
\item $\varphi_{27} :\equiv b \in \mathbb{R} \to x \in \mathbb{R} \to a - x + (x - b) = a + (x + (-b + -x))$
\item $\varphi_{28} :\equiv b \in \mathbb{R} \to x \in \mathbb{R} \to a - x + (x - b) = a + (x + (-x + -b))$
\item $\varphi_{29} :\equiv b \in \mathbb{R} \to x \in \mathbb{R} \to a - x + (x - b) = a + (x + -x + -b)$
\item $\varphi_{30} :\equiv \forall x\,\forall y\,(x = y \to y = x) \to \forall y\,(x + -x + -b = y \to y = x + -x + -b)$
\item $\varphi_{31} :\equiv a \in \mathbb{R} \to b \in \mathbb{R} \to x \in \mathbb{R} \to a + (0 + -b) = a + (-b + 0)$
\item $\varphi_{32} :\equiv a \in \mathbb{R} \to b \in \mathbb{R} \to x \in \mathbb{R} \to a - x + (x - b) = a + (0 + -b)$
\item $\varphi_{33} :\equiv a \in \mathbb{R} \to b \in \mathbb{R} \to x \in \mathbb{R} \to a - x + (x - b) = a + (-b + 0)$
\item $\varphi_{34} :\equiv b \in \mathbb{R} \to x \in \mathbb{R} \to a + (x - b + -x) = a + (x + -b + -x)$
\item $\varphi_{35} :\equiv a - x + (x - b) = a - x + (x - b) \to a - x + (x - b) = a + -x + (x - b)$
\item $\varphi_{36} :\equiv a - x + (x - b) = a + -b \to a + -b = a - b \to a - x + (x - b) = a - b$
\item $\varphi_{37} :\equiv a \in \mathbb{R} \to b \in \mathbb{R} \to x \in \mathbb{R} \to a - x + (x - b) = a + -x + (x - b)$
\item $\varphi_{38} :\equiv a \in \mathbb{R} \to b \in \mathbb{R} \to x \in \mathbb{R} \to a - x + (x - b) = a + (x - b + -x)$
\item $\varphi_{39} :\equiv \forall z\,(a - x + (x - b) = a + (-b + 0) \to a + (-b + 0) = z \to a - x + (x - b) = z)$
\item $\varphi_{40} :\equiv \forall w\,(-b + 0 = w \to a + (-b + 0) = a + (-b + 0) \to a + (-b + 0) = a + w)$
\item $\varphi_{41} :\equiv \forall z\,(a - x + (x - b) = a + (0 + -b) \to a + (0 + -b) = z \to a - x + (x - b) = z)$
\item $\varphi_{42} :\equiv \forall w\,(0 + -b = w \to a + (0 + -b) = a + (0 + -b) \to a + (0 + -b) = a + w)$
\item $\varphi_{43} :\equiv \forall w\,(a - x = w \to a - x + (x - b) = a - x + (x - b) \to a - x + (x - b) = w + (x - b))$
\item $\varphi_{44} :\equiv a + (x + -x + -b) = a + (0 + -b) \to a - x + (x - b) = a + (0 + -b)$
\item $\varphi_{45} :\equiv b \in \mathbb{R} \to x \in \mathbb{R} \to a + (x + (-x + -b)) = a + (x + -x + -b)$
\item $\varphi_{46} :\equiv a + (x + (-b + -x)) = a + (x + z) \to a + (x + (-b + -x)) = a + (x + w)$
\item $\varphi_{47} :\equiv b \in \mathbb{R} \to x \in \mathbb{R} \to a + (x + -b + -x) = a + (x + (-b + -x))$
\item $\varphi_{48} :\equiv a + (x - b + -x) = a + (z + -x) \to a + (x - b + -x) = a + (w + -x)$
\item $\varphi_{49} :\equiv -x \in \mathbb{R} \to x - b \in \mathbb{R} \to a + -x + (x - b) = a + (-x + (x - b))$
\item $\varphi_{50} :\equiv \forall z\,(x \in \mathbb{R} \to -x \in \mathbb{R} \to z \in \mathbb{R} \to x + -x + z = x + (-x + z))$
\item $\varphi_{51} :\equiv \forall z\,(x \in \mathbb{R} \to -b \in \mathbb{R} \to z \in \mathbb{R} \to x + -b + z = x + (-b + z))$
\item $\varphi_{52} :\equiv \forall z\,(a \in \mathbb{R} \to -x \in \mathbb{R} \to z \in \mathbb{R} \to a + -x + z = a + (-x + z))$
\item $\varphi_{53} :\equiv a - b = a + -b \to (a - b = a + -b \to a + -b = a - b) \to a + -b = a - b$
\item $\varphi_{54} :\equiv a \in \mathbb{R} \to b \in \mathbb{R} \to x \in \mathbb{R} \to a + (x + -x + -b) = a + (0 + -b)$
\item $\varphi_{55} :\equiv a \in \mathbb{R} \to b \in \mathbb{R} \to x \in \mathbb{R} \to x + -x + -b = x + (-x + -b)$
\item $\varphi_{56} :\equiv a \in \mathbb{R} \to b \in \mathbb{R} \to x \in \mathbb{R} \to x + (-x + -b) = x + -x + -b$
\item $\varphi_{57} :\equiv a \in \mathbb{R} \to b \in \mathbb{R} \to x \in \mathbb{R} \to x + -b + -x = x + (-b + -x)$
\item $\varphi_{58} :\equiv a + (-x + (x - b)) = a + (x - b + -x) \to a - x + (x - b) = a + (x - b + -x)$
\item $\varphi_{59} :\equiv a + -x + (x - b) = a + (-x + (x - b)) \to a - x + (x - b) = a + (-x + (x - b))$
\item $\varphi_{60} :\equiv \forall z\,(a - x + (x - b) = a + (x - b + -x) \to a + (x - b + -x) = z \to a - x + (x - b) = z)$
\item $\varphi_{61} :\equiv \forall z\,(a - x + (x - b) = a + (-x + (x - b)) \to a + (-x + (x - b)) = z \to a - x + (x - b) = z)$
\item $\varphi_{62} :\equiv \forall y\,(x \in \mathbb{R} \to y \in \mathbb{R} \to x - y \in \mathbb{R}) \to x \in \mathbb{R} \to b \in \mathbb{R} \to x - b \in \mathbb{R}$
\item $\varphi_{63} :\equiv \forall z\,(a - x + (x - b) = a + -x + (x - b) \to a + -x + (x - b) = z \to a - x + (x - b) = z)$
\item $\varphi_{64} :\equiv x + -x + -b = x + (-x + -b) \to x + (-x + -b) = x + -x + -b$
\item $\varphi_{65} :\equiv a + (x + -x + -b) = a + (z + -b) \to a + (x + -x + -b) = a + (w + -b)$
\item $\varphi_{66} :\equiv a + (-x + (x - b)) = a + (-x + (x - b)) \to a + (-x + (x - b)) = a + (x - b + -x)$
\item $\varphi_{67} :\equiv \forall x\,\forall y\,(x \in \mathbb{R} \to y \in \mathbb{R} \to x + y = y + x) \to \forall y\,(0 \in \mathbb{R} \to y \in \mathbb{R} \to 0 + y = y + 0)$
\item $\varphi_{68} :\equiv x \in \mathbb{R} \to -x \in \mathbb{R} \to -b \in \mathbb{R} \to x + -x + -b = x + (-x + -b)$
\item $\varphi_{69} :\equiv x \in \mathbb{R} \to -b \in \mathbb{R} \to -x \in \mathbb{R} \to x + -b + -x = x + (-b + -x)$
\item $\varphi_{70} :\equiv \forall y\,(a - x + (x - b) = y \to y = a - x + (x - b)) \to a - x + (x - b) = a - b \to a - b = a - x + (x - b)$
\item $\varphi_{71} :\equiv \forall x\,\forall y\,(x \in \mathbb{R} \to y \in \mathbb{R} \to x - y \in \mathbb{R}) \to \forall y\,(x \in \mathbb{R} \to y \in \mathbb{R} \to x - y \in \mathbb{R})$
\item $\varphi_{72} :\equiv a + (x - b + -x) = a + (x + -b + -x) \to a - x + (x - b) = a + (x + -b + -x)$
\item $\varphi_{73} :\equiv a + (x - b + -x) = a + (x - b + -x) \to a + (x - b + -x) = a + (x + -b + -x)$
\item $\varphi_{74} :\equiv a - x + (x - b) = a + (-b + 0) \to a + (-b + 0) = a + -b \to a - x + (x - b) = a + -b$
\item $\varphi_{75} :\equiv \forall z\,(a - x + (x - b) = a + (x + -x + -b) \to a + (x + -x + -b) = z \to a - x + (x - b) = z)$
\item $\varphi_{76} :\equiv \forall z\,(a - x + (x - b) = a + (x + (-x + -b)) \to a + (x + (-x + -b)) = z \to a - x + (x - b) = z)$
\item $\varphi_{77} :\equiv \forall z\,(a - x + (x - b) = a + (x + (-b + -x)) \to a + (x + (-b + -x)) = z \to a - x + (x - b) = z)$
\item $\varphi_{78} :\equiv \forall z\,(a - x + (x - b) = a + (x + -b + -x) \to a + (x + -b + -x) = z \to a - x + (x - b) = z)$
\item $\varphi_{79} :\equiv \forall w\,(x - b = w \to a + (x - b + -x) = a + (x - b + -x) \to a + (x - b + -x) = a + (w + -x))$
\item $\varphi_{80} :\equiv \forall w\,(-x + (x - b) = w \to a + (-x + (x - b)) = a + (-x + (x - b)) \to a + (-x + (x - b)) = a + w)$
\item $\varphi_{81} :\equiv \forall y\,(0 \in \mathbb{R} \to y \in \mathbb{R} \to 0 + y = y + 0) \to 0 \in \mathbb{R} \to -b \in \mathbb{R} \to 0 + -b = -b + 0$
\item $\varphi_{82} :\equiv a + (x + -x + -b) = a + (x + -x + -b) \to a + (x + -x + -b) = a + (0 + -b)$
\item $\varphi_{83} :\equiv a + (x + (-x + -b)) = a + (x + -x + -b) \to a - x + (x - b) = a + (x + -x + -b)$
\item $\varphi_{84} :\equiv a + (x + (-b + -x)) = a + (x + (-x + -b)) \to a - x + (x - b) = a + (x + (-x + -b))$
\item $\varphi_{85} :\equiv a + (x + -b + -x) = a + (x + (-b + -x)) \to a - x + (x - b) = a + (x + (-b + -x))$
\item $\varphi_{86} :\equiv (a \in \mathbb{R} \to b \in \mathbb{R} \to x \in \mathbb{R} \to x \in \mathbb{R}) \to a \in \mathbb{R} \to b \in \mathbb{R} \to x \in \mathbb{R} \to x + -x = 0$
\item $\varphi_{87} :\equiv \forall x\,\forall y\,(x \in \mathbb{R} \to y \in \mathbb{R} \to x + y = y + x) \to \forall y\,(-b \in \mathbb{R} \to y \in \mathbb{R} \to -b + y = y + -b)$
\item $\varphi_{88} :\equiv \forall x\,\forall y\,(x \in \mathbb{R} \to y \in \mathbb{R} \to x + y = y + x) \to \forall y\,(-x \in \mathbb{R} \to y \in \mathbb{R} \to -x + y = y + -x)$
\item $\varphi_{89} :\equiv (a \in \mathbb{R} \to b \in \mathbb{R} \to x \in \mathbb{R} \to b \in \mathbb{R}) \to a \in \mathbb{R} \to b \in \mathbb{R} \to x \in \mathbb{R} \to -b \in \mathbb{R}$
\item $\varphi_{90} :\equiv (a \in \mathbb{R} \to b \in \mathbb{R} \to x \in \mathbb{R} \to x \in \mathbb{R}) \to a \in \mathbb{R} \to b \in \mathbb{R} \to x \in \mathbb{R} \to -x \in \mathbb{R}$
\item $\varphi_{91} :\equiv (a \in \mathbb{R} \to b \in \mathbb{R} \to x \in \mathbb{R} \to b \in \mathbb{R}) \to a \in \mathbb{R} \to b \in \mathbb{R} \to x \in \mathbb{R} \to x - b \in \mathbb{R}$
\item $\varphi_{92} :\equiv a + (x + (-x + -b)) = a + (x + (-x + -b)) \to a + (x + (-x + -b)) = a + (x + -x + -b)$
\item $\varphi_{93} :\equiv a + (x + (-b + -x)) = a + (x + (-b + -x)) \to a + (x + (-b + -x)) = a + (x + (-x + -b))$
\item $\varphi_{94} :\equiv a + (x + -b + -x) = a + (x + -b + -x) \to a + (x + -b + -x) = a + (x + (-b + -x))$
\item $\varphi_{95} :\equiv (a \in \mathbb{R} \to b \in \mathbb{R} \to x \in \mathbb{R} \to -b \in \mathbb{R}) \to a \in \mathbb{R} \to b \in \mathbb{R} \to x \in \mathbb{R} \to -b + 0 = -b$
\item $\varphi_{96} :\equiv (a - x + (x - b) = a - b \to a - b = a - x + (x - b)) \to a \in \mathbb{R} \to b \in \mathbb{R} \to x \in \mathbb{R} \to a - b = a - x + (x - b)$
\item $\varphi_{97} :\equiv \forall w\,(x + -x = w \to a + (x + -x + -b) = a + (x + -x + -b) \to a + (x + -x + -b) = a + (w + -b))$
\item $\varphi_{98} :\equiv \forall w\,(x + (-x + -b) = w \to a + (x + (-x + -b)) = a + (x + (-x + -b)) \to a + (x + (-x + -b)) = a + w)$
\item $\varphi_{99} :\equiv \forall w\,(-b + -x = w \to a + (x + (-b + -x)) = a + (x + (-b + -x)) \to a + (x + (-b + -x)) = a + (x + w))$
\item $\varphi_{100} :\equiv \forall w\,(x + -b + -x = w \to a + (x + -b + -x) = a + (x + -b + -x) \to a + (x + -b + -x) = a + w)$
\item $\varphi_{101} :\equiv (a \in \mathbb{R} \to b \in \mathbb{R} \to x \in \mathbb{R} \to -b \in \mathbb{R}) \to a \in \mathbb{R} \to b \in \mathbb{R} \to x \in \mathbb{R} \to 0 + -b = -b + 0$
\item $\varphi_{102} :\equiv \forall y\,(-b \in \mathbb{R} \to y \in \mathbb{R} \to -b + y = y + -b) \to -b \in \mathbb{R} \to -x \in \mathbb{R} \to -b + -x = -x + -b$
\item $\varphi_{103} :\equiv (a \in \mathbb{R} \to b \in \mathbb{R} \to x \in \mathbb{R} \to -x \in \mathbb{R}) \to a \in \mathbb{R} \to b \in \mathbb{R} \to x \in \mathbb{R} \to -b + -x = -x + -b$
\item $\varphi_{104} :\equiv \forall y\,(-x \in \mathbb{R} \to y \in \mathbb{R} \to -x + y = y + -x) \to -x \in \mathbb{R} \to x - b \in \mathbb{R} \to -x + (x - b) = x - b + -x$
\item $\varphi_{105} :\equiv (a \in \mathbb{R} \to b \in \mathbb{R} \to x \in \mathbb{R} \to x - b \in \mathbb{R}) \to a \in \mathbb{R} \to b \in \mathbb{R} \to x \in \mathbb{R} \to -x + (x - b) = x - b + -x$
\item $\varphi_{106} :\equiv (-b + 0 = -b \to \varphi_{13}) \to a \in \mathbb{R} \to b \in \mathbb{R} \to x \in \mathbb{R} \to a + (-b + 0) = a + -b$
\item $\varphi_{107} :\equiv a + -b = a - b \to \varphi_{36} \to a \in \mathbb{R} \to b \in \mathbb{R} \to x \in \mathbb{R} \to a - x + (x - b) = a - b$
\item $\varphi_{108} :\equiv \varphi_{74} \to a \in \mathbb{R} \to b \in \mathbb{R} \to x \in \mathbb{R} \to a - x + (x - b) = a + -b$
\item $\varphi_{109} :\equiv (a \in \mathbb{R} \to b \in \mathbb{R} \to x \in \mathbb{R} \to a - x + (x - b) = a - b) \to \varphi_{96}$
\item $\varphi_{110} :\equiv (a \in \mathbb{R} \to b \in \mathbb{R} \to x \in \mathbb{R} \to -x \in \mathbb{R}) \to (a \in \mathbb{R} \to b \in \mathbb{R} \to x \in \mathbb{R} \to -b \in \mathbb{R}) \to \varphi_{55}$
\item $\varphi_{111} :\equiv (a \in \mathbb{R} \to b \in \mathbb{R} \to x \in \mathbb{R} \to -b \in \mathbb{R}) \to (a \in \mathbb{R} \to b \in \mathbb{R} \to x \in \mathbb{R} \to -x \in \mathbb{R}) \to \varphi_{57}$
\item $\varphi_{112} :\equiv (a \in \mathbb{R} \to b \in \mathbb{R} \to x \in \mathbb{R} \to -x \in \mathbb{R}) \to (a \in \mathbb{R} \to b \in \mathbb{R} \to x \in \mathbb{R} \to x - b \in \mathbb{R}) \to a \in \mathbb{R} \to \varphi_{25}$
\item $\varphi_{113} :\equiv (x \in \mathbb{R} \to b \in \mathbb{R} \to x - b \in \mathbb{R}) \to (a \in \mathbb{R} \to b \in \mathbb{R} \to x \in \mathbb{R} \to x \in \mathbb{R}) \to \varphi_{91}$
\item $\varphi_{114} :\equiv (a - x + (x - b) = a + -x + (x - b) \to \varphi_{59}) \to a \in \mathbb{R} \to \varphi_{14}$
\item $\varphi_{115} :\equiv a + (0 + -b) = a + (0 + -b) \to (0 + -b = -b + 0 \to \varphi_{21}) \to \varphi_{31}$
\item $\varphi_{116} :\equiv a - x + (x - b) = a - x + (x - b) \to (a - x = a + -x \to \varphi_{35}) \to \varphi_{37}$
\item $\varphi_{117} :\equiv (a - x + (x - b) = a + (x - b + -x) \to \varphi_{72}) \to a \in \mathbb{R} \to \varphi_{26}$
\item $\varphi_{118} :\equiv (a - x + (x - b) = a + (x + -b + -x) \to \varphi_{85}) \to a \in \mathbb{R} \to \varphi_{27}$
\item $\varphi_{119} :\equiv (a - x + (x - b) = a + (x + (-b + -x)) \to \varphi_{84}) \to a \in \mathbb{R} \to \varphi_{28}$
\item $\varphi_{120} :\equiv (a - x + (x - b) = a + (x + (-x + -b)) \to \varphi_{83}) \to a \in \mathbb{R} \to \varphi_{29}$
\item $\varphi_{121} :\equiv (a \in \mathbb{R} \to b \in \mathbb{R} \to x \in \mathbb{R} \to a - x + (x - b) = a + -b) \to \varphi_{107}$
\item $\varphi_{122} :\equiv (a \in \mathbb{R} \to b \in \mathbb{R} \to x \in \mathbb{R} \to a + (-b + 0) = a + -b) \to \varphi_{108}$
\item $\varphi_{123} :\equiv (a \in \mathbb{R} \to b \in \mathbb{R} \to x \in \mathbb{R} \to -b + 0 = -b) \to a + (-b + 0) = a + (-b + 0) \to \varphi_{106}$
\item $\varphi_{124} :\equiv a + (x - b + -x) = a + (x - b + -x) \to (x - b = x + -b \to \varphi_{73}) \to a \in \mathbb{R} \to \varphi_{34}$
\item $\varphi_{125} :\equiv a + (x + -x + -b) = a + (x + -x + -b) \to (x + -x = 0 \to \varphi_{82}) \to \varphi_{54}$
\item $\varphi_{126} :\equiv a + (-x + (x - b)) = a + (-x + (x - b)) \to (-x + (x - b) = x - b + -x \to \varphi_{66}) \to a \in \mathbb{R} \to \varphi_{24}$
\item $\varphi_{127} :\equiv (-b \in \mathbb{R} \to -x \in \mathbb{R} \to -b + -x = -x + -b) \to (a \in \mathbb{R} \to b \in \mathbb{R} \to x \in \mathbb{R} \to -b \in \mathbb{R}) \to \varphi_{103}$
\item $\varphi_{128} :\equiv (-x \in \mathbb{R} \to x - b \in \mathbb{R} \to -x + (x - b) = x - b + -x) \to (a \in \mathbb{R} \to b \in \mathbb{R} \to x \in \mathbb{R} \to -x \in \mathbb{R}) \to \varphi_{105}$
\item $\varphi_{129} :\equiv (a \in \mathbb{R} \to \varphi_{24}) \to (a - x + (x - b) = a + (-x + (x - b)) \to \varphi_{58}) \to \varphi_{38}$
\item $\varphi_{130} :\equiv a + (x + (-b + -x)) = a + (x + (-b + -x)) \to (-b + -x = -x + -b \to \varphi_{93}) \to a \in \mathbb{R} \to b \in \mathbb{R} \to \varphi_{16}$
\item $\varphi_{131} :\equiv a + (x + (-x + -b)) = a + (x + (-x + -b)) \to (x + (-x + -b) = x + -x + -b \to \varphi_{92}) \to a \in \mathbb{R} \to \varphi_{45}$
\item $\varphi_{132} :\equiv a + (x + -b + -x) = a + (x + -b + -x) \to (x + -b + -x = x + (-b + -x) \to \varphi_{94}) \to a \in \mathbb{R} \to \varphi_{47}$
\item $\varphi_{133} :\equiv (a \in \mathbb{R} \to b \in \mathbb{R} \to x \in \mathbb{R} \to -x + (x - b) = x - b + -x) \to \varphi_{126}$
\item $\varphi_{134} :\equiv (a \in \mathbb{R} \to \varphi_{49}) \to (a \in \mathbb{R} \to b \in \mathbb{R} \to x \in \mathbb{R} \to a \in \mathbb{R}) \to \varphi_{112}$
\item $\varphi_{135} :\equiv (a \in \mathbb{R} \to \varphi_{29}) \to \varphi_{54} \to (a - x + (x - b) = a + (x + -x + -b) \to \varphi_{44}) \to \varphi_{32}$
\end{abbrevlist}
\end{fullproof}

% 自動生成 (prover/build.js): MulSub  ノード数 130
\begin{fullproof}{[F17]}{$\forall c\,\forall x\,\forall a\,(c \in \mathbb{R} \to x \in \mathbb{R} \to a \in \mathbb{R} \to c \cdot x - c \cdot a = c \cdot (x - a))$}{130}\phantomsection\label{tree:MulSub}
\noindent$\Pi_{1}$:
\begin{ptree}
\AXC{$x = x$\lbl{A5}}
\AXC{$x = y \to x = x \to y = x$\lbl{A6}}
\AXC{$(x = y \to x = x \to y = x) \to x = x \to x = y \to y = x$\lbl{A1}}
\BIC{$x = x \to x = y \to y = x$}
\BIC{$x = y \to y = x$}
\UIC{$\forall y\,(x = y \to y = x)$}
\UIC{$\forall x\,\forall y\,(x = y \to y = x)$}
\end{ptree}
\noindent$\Pi_{2}$:
\begin{ptree}
\AXC{$\stackrel{\Pi_{1}}{\vdots}$}\noLine
\UIC{$\forall x\,\forall y\,(x = y \to y = x)$}
\AXC{$\varphi_{2}$\lbl{A2}}
\BIC{$\forall y\,(c \cdot (x - a) = y \to y = c \cdot (x - a))$}
\AXC{$\forall y\,(c \cdot (x - a) = y \to y = c \cdot (x - a)) \to \varphi_{6}$\lbl{A2}}
\BIC{$\varphi_{6}$}
\end{ptree}
\noindent$\Pi_{3}$:
\begin{ptree}
\AXC{$y = z \to x = y \to x = z$\lbl{A6}}
\AXC{$(y = z \to x = y \to x = z) \to x = y \to y = z \to x = z$\lbl{A1}}
\BIC{$x = y \to y = z \to x = z$}
\UIC{$\forall z\,(x = y \to y = z \to x = z)$}
\UIC{$\forall y\,\forall z\,(x = y \to y = z \to x = z)$}
\UIC{$\forall x\,\forall y\,\forall z\,(x = y \to y = z \to x = z)$}
\AXC{$\varphi_{9}$\lbl{A2}}
\BIC{$\forall y\,\forall z\,(c \cdot (x - a) = y \to y = z \to c \cdot (x - a) = z)$}
\AXC{$\varphi_{35}$\lbl{A2}}
\BIC{$\varphi_{22}$}
\end{ptree}
\noindent$\Pi_{4}$:
\begin{ptree}
\AXC{$\stackrel{\Pi_{1}}{\vdots}$}\noLine
\UIC{$\forall x\,\forall y\,(x = y \to y = x)$}
\AXC{$\varphi_{7}$\lbl{A2}}
\BIC{$\forall y\,(c \cdot x - c \cdot a = y \to y = c \cdot x - c \cdot a)$}
\AXC{$\forall y\,(c \cdot x - c \cdot a = y \to y = c \cdot x - c \cdot a) \to \varphi_{24}$\lbl{A2}}
\BIC{$\varphi_{24}$}
\end{ptree}
\noindent$\Pi_{5}$:
\begin{ptree}
\AXC{$\stackrel{\Pi_{4}}{\vdots}$}\noLine
\UIC{$\varphi_{24}$}
\AXC{$\forall x\,\forall y\,(x - y = x + -y)$\lbl{Def$-$}}
\AXC{$\varphi_{4}$\lbl{A2}}
\BIC{$\forall y\,(c \cdot x - y = c \cdot x + -y)$}
\AXC{$\varphi_{14}$\lbl{A2}}
\BIC{$c \cdot x - c \cdot a = c \cdot x + -(c \cdot a)$}
\AXC{$\varphi_{38}$\lbl{A1}}
\BIC{$\varphi_{24} \to c \cdot x + -(c \cdot a) = c \cdot x - c \cdot a$}
\BIC{$c \cdot x + -(c \cdot a) = c \cdot x - c \cdot a$}
\end{ptree}
\noindent$\Pi_{6}$:
\begin{ptree}
\AXC{$y = z \to x = y \to x = z$\lbl{A6}}
\AXC{$(y = z \to x = y \to x = z) \to x = y \to y = z \to x = z$\lbl{A1}}
\BIC{$x = y \to y = z \to x = z$}
\UIC{$\forall z\,(x = y \to y = z \to x = z)$}
\UIC{$\forall y\,\forall z\,(x = y \to y = z \to x = z)$}
\UIC{$\forall x\,\forall y\,\forall z\,(x = y \to y = z \to x = z)$}
\AXC{$\varphi_{9}$\lbl{A2}}
\BIC{$\forall y\,\forall z\,(c \cdot (x - a) = y \to y = z \to c \cdot (x - a) = z)$}
\AXC{$\varphi_{36}$\lbl{A2}}
\BIC{$\varphi_{23}$}
\end{ptree}
\noindent$\Pi_{7}$:
\begin{ptree}
\AXC{$\forall x\,\forall y\,(x \in \mathbb{R} \to y \in \mathbb{R} \to x \cdot -y = -(x \cdot y))$\lbl{[F10]}}
\AXC{$\varphi_{32}$\lbl{A2}}
\BIC{$\forall y\,(c \in \mathbb{R} \to y \in \mathbb{R} \to c \cdot -y = -(c \cdot y))$}
\AXC{$\varphi_{30}$\lbl{A2}}
\BIC{$c \in \mathbb{R} \to a \in \mathbb{R} \to c \cdot -a = -(c \cdot a)$}
\AXC{$\varphi_{43}$\lbl{A1}}
\BIC{$(c \in \mathbb{R} \to x \in \mathbb{R} \to a \in \mathbb{R} \to c \in \mathbb{R}) \to \varphi_{34}$}
\end{ptree}
\noindent$\Pi_{8}$:
\begin{ptree}
\AXC{$c \in \mathbb{R} \to x \in \mathbb{R} \to a \in \mathbb{R} \to a \in \mathbb{R}$\lbl{A1}}
\AXC{$c \in \mathbb{R} \to x \in \mathbb{R} \to a \in \mathbb{R} \to c \in \mathbb{R}$\lbl{A1}}
\AXC{$\stackrel{\Pi_{7}}{\vdots}$}\noLine
\UIC{$(c \in \mathbb{R} \to x \in \mathbb{R} \to a \in \mathbb{R} \to c \in \mathbb{R}) \to \varphi_{34}$}
\BIC{$\varphi_{34}$}
\BIC{$c \in \mathbb{R} \to x \in \mathbb{R} \to a \in \mathbb{R} \to c \cdot -a = -(c \cdot a)$}
\end{ptree}
\noindent$\Pi_{9}$:
\begin{ptree}
\AXC{$\stackrel{\Pi_{8}}{\vdots}$}\noLine
\UIC{$c \in \mathbb{R} \to x \in \mathbb{R} \to a \in \mathbb{R} \to c \cdot -a = -(c \cdot a)$}
\AXC{$(c \in \mathbb{R} \to x \in \mathbb{R} \to a \in \mathbb{R} \to c \cdot -a = -(c \cdot a)) \to \varphi_{44}$\lbl{A1}}
\BIC{$\varphi_{44}$}
\end{ptree}
\noindent$\Pi_{10}$:
\begin{ptree}
\AXC{$x = x$\lbl{A5}}
\UIC{$\forall x\,(x = x)$}
\AXC{$\forall x\,(x = x) \to c \cdot x + c \cdot -a = c \cdot x + c \cdot -a$\lbl{A2}}
\BIC{$c \cdot x + c \cdot -a = c \cdot x + c \cdot -a$}
\AXC{$\stackrel{\Pi_{9}}{\vdots}$}\noLine
\UIC{$\varphi_{44}$}
\BIC{$(c \cdot -a = -(c \cdot a) \to \varphi_{28}) \to \varphi_{26}$}
\end{ptree}
\noindent$\Pi_{11}$:
\begin{ptree}
\AXC{$\varphi_{19}$\lbl{A6}}
\UIC{$\forall w\,(\varphi_{19})$}
\UIC{$\forall z\,\forall w\,(\varphi_{19})$}
\AXC{$\forall z\,\forall w\,(\varphi_{19}) \to \varphi_{29}$\lbl{A2}}
\BIC{$\varphi_{29}$}
\AXC{$\varphi_{29} \to c \cdot -a = -(c \cdot a) \to \varphi_{28}$\lbl{A2}}
\BIC{$c \cdot -a = -(c \cdot a) \to \varphi_{28}$}
\AXC{$\stackrel{\Pi_{10}}{\vdots}$}\noLine
\UIC{$(c \cdot -a = -(c \cdot a) \to \varphi_{28}) \to \varphi_{26}$}
\BIC{$\varphi_{26}$}
\end{ptree}
\noindent$\Pi_{12}$:
\begin{ptree}
\AXC{$y = z \to x = y \to x = z$\lbl{A6}}
\AXC{$(y = z \to x = y \to x = z) \to x = y \to y = z \to x = z$\lbl{A1}}
\BIC{$x = y \to y = z \to x = z$}
\UIC{$\forall z\,(x = y \to y = z \to x = z)$}
\UIC{$\forall y\,\forall z\,(x = y \to y = z \to x = z)$}
\UIC{$\forall x\,\forall y\,\forall z\,(x = y \to y = z \to x = z)$}
\AXC{$\varphi_{9}$\lbl{A2}}
\BIC{$\forall y\,\forall z\,(c \cdot (x - a) = y \to y = z \to c \cdot (x - a) = z)$}
\AXC{$\varphi_{33}$\lbl{A2}}
\BIC{$\varphi_{10}$}
\end{ptree}
\noindent$\Pi_{13}$:
\begin{ptree}
\AXC{$\forall x \in \mathbb{R}\,(-x \in \mathbb{R})$\lbl{(R1)}}
\AXC{$\forall x \in \mathbb{R}\,(-x \in \mathbb{R}) \to a \in \mathbb{R} \to -a \in \mathbb{R}$\lbl{A2}}
\BIC{$a \in \mathbb{R} \to -a \in \mathbb{R}$}
\AXC{$(a \in \mathbb{R} \to -a \in \mathbb{R}) \to \varphi_{31}$\lbl{A1}}
\BIC{$\varphi_{31}$}
\end{ptree}
\noindent$\Pi_{14}$:
\begin{ptree}
\AXC{$\varphi_{20}$\lbl{(R10)}}
\AXC{$\varphi_{20} \to \varphi_{17}$\lbl{A2}}
\BIC{$\varphi_{17}$}
\AXC{$\varphi_{17} \to \varphi_{12}$\lbl{A2}}
\BIC{$\varphi_{12}$}
\AXC{$\varphi_{12} \to \varphi_{27}$\lbl{A2}}
\BIC{$\varphi_{27}$}
\AXC{$\varphi_{27} \to (c \in \mathbb{R} \to x \in \mathbb{R} \to a \in \mathbb{R} \to c \in \mathbb{R}) \to \varphi_{40}$\lbl{A1}}
\BIC{$(c \in \mathbb{R} \to x \in \mathbb{R} \to a \in \mathbb{R} \to c \in \mathbb{R}) \to \varphi_{40}$}
\end{ptree}
\noindent$\Pi_{15}$:
\begin{ptree}
\AXC{$c \in \mathbb{R} \to x \in \mathbb{R} \to a \in \mathbb{R} \to x \in \mathbb{R}$\lbl{A1}}
\AXC{$c \in \mathbb{R} \to x \in \mathbb{R} \to a \in \mathbb{R} \to c \in \mathbb{R}$\lbl{A1}}
\AXC{$\stackrel{\Pi_{14}}{\vdots}$}\noLine
\UIC{$(c \in \mathbb{R} \to x \in \mathbb{R} \to a \in \mathbb{R} \to c \in \mathbb{R}) \to \varphi_{40}$}
\BIC{$\varphi_{40}$}
\BIC{$(c \in \mathbb{R} \to x \in \mathbb{R} \to a \in \mathbb{R} \to -a \in \mathbb{R}) \to \varphi_{21}$}
\end{ptree}
\noindent$\Pi_{16}$:
\begin{ptree}
\AXC{$c \in \mathbb{R} \to x \in \mathbb{R} \to a \in \mathbb{R} \to a \in \mathbb{R}$\lbl{A1}}
\AXC{$\stackrel{\Pi_{13}}{\vdots}$}\noLine
\UIC{$\varphi_{31}$}
\BIC{$c \in \mathbb{R} \to x \in \mathbb{R} \to a \in \mathbb{R} \to -a \in \mathbb{R}$}
\AXC{$\stackrel{\Pi_{15}}{\vdots}$}\noLine
\UIC{$(c \in \mathbb{R} \to x \in \mathbb{R} \to a \in \mathbb{R} \to -a \in \mathbb{R}) \to \varphi_{21}$}
\BIC{$\varphi_{21}$}
\end{ptree}
\noindent$\Pi_{17}$:
\begin{ptree}
\AXC{$z = w \to c \cdot (x - a) = c \cdot z \to c \cdot (x - a) = c \cdot w$\lbl{A6}}
\UIC{$\varphi_{1}$}
\UIC{$\forall z\,(\varphi_{1})$}
\AXC{$\forall z\,(\varphi_{1}) \to \varphi_{5}$\lbl{A2}}
\BIC{$\varphi_{5}$}
\AXC{$\varphi_{5} \to x - a = x + -a \to \varphi_{3}$\lbl{A2}}
\BIC{$x - a = x + -a \to \varphi_{3}$}
\end{ptree}
\noindent$\Pi_{18}$:
\begin{ptree}
\AXC{$\forall x\,\forall y\,(x - y = x + -y)$\lbl{Def$-$}}
\AXC{$\forall x\,\forall y\,(x - y = x + -y) \to \forall y\,(x - y = x + -y)$\lbl{A2}}
\BIC{$\forall y\,(x - y = x + -y)$}
\end{ptree}
\noindent$\Pi_{19}$:
\begin{ptree}
\AXC{$\stackrel{\Pi_{18}}{\vdots}$}\noLine
\UIC{$\forall y\,(x - y = x + -y)$}
\AXC{$\forall y\,(x - y = x + -y) \to x - a = x + -a$\lbl{A2}}
\BIC{$x - a = x + -a$}
\AXC{$\varphi_{41}$\lbl{A1}}
\BIC{$c \cdot (x - a) = c \cdot (x - a) \to (x - a = x + -a \to \varphi_{3}) \to \varphi_{8}$}
\end{ptree}
\noindent$\Pi_{20}$:
\begin{ptree}
\AXC{$x = x$\lbl{A5}}
\UIC{$\forall x\,(x = x)$}
\AXC{$\forall x\,(x = x) \to c \cdot (x - a) = c \cdot (x - a)$\lbl{A2}}
\BIC{$c \cdot (x - a) = c \cdot (x - a)$}
\AXC{$\stackrel{\Pi_{19}}{\vdots}$}\noLine
\UIC{$c \cdot (x - a) = c \cdot (x - a) \to (x - a = x + -a \to \varphi_{3}) \to \varphi_{8}$}
\BIC{$(x - a = x + -a \to \varphi_{3}) \to \varphi_{8}$}
\end{ptree}
\noindent$\Pi_{21}$:
\begin{ptree}
\AXC{$\stackrel{\Pi_{17}}{\vdots}$}\noLine
\UIC{$x - a = x + -a \to \varphi_{3}$}
\AXC{$\stackrel{\Pi_{20}}{\vdots}$}\noLine
\UIC{$(x - a = x + -a \to \varphi_{3}) \to \varphi_{8}$}
\BIC{$\varphi_{8}$}
\AXC{$\varphi_{8} \to \varphi_{21} \to (c \cdot (x - a) = c \cdot (x + -a) \to \varphi_{18}) \to \varphi_{15}$\lbl{A1}}
\BIC{$\varphi_{21} \to (c \cdot (x - a) = c \cdot (x + -a) \to \varphi_{18}) \to \varphi_{15}$}
\end{ptree}
\noindent$\Pi_{22}$:
\begin{ptree}
\AXC{$\stackrel{\Pi_{12}}{\vdots}$}\noLine
\UIC{$\varphi_{10}$}
\AXC{$\varphi_{10} \to c \cdot (x - a) = c \cdot (x + -a) \to \varphi_{18}$\lbl{A2}}
\BIC{$c \cdot (x - a) = c \cdot (x + -a) \to \varphi_{18}$}
\AXC{$\stackrel{\Pi_{16}}{\vdots}$}\noLine
\UIC{$\varphi_{21}$}
\AXC{$\stackrel{\Pi_{21}}{\vdots}$}\noLine
\UIC{$\varphi_{21} \to (c \cdot (x - a) = c \cdot (x + -a) \to \varphi_{18}) \to \varphi_{15}$}
\BIC{$(c \cdot (x - a) = c \cdot (x + -a) \to \varphi_{18}) \to \varphi_{15}$}
\BIC{$\varphi_{15}$}
\end{ptree}
\noindent$\Pi_{23}$:
\begin{ptree}
\AXC{$\stackrel{\Pi_{11}}{\vdots}$}\noLine
\UIC{$\varphi_{26}$}
\AXC{$\stackrel{\Pi_{22}}{\vdots}$}\noLine
\UIC{$\varphi_{15}$}
\AXC{$\varphi_{15} \to \varphi_{26} \to (c \cdot (x - a) = c \cdot x + c \cdot -a \to \varphi_{25}) \to \varphi_{16}$\lbl{A1}}
\BIC{$\varphi_{26} \to (c \cdot (x - a) = c \cdot x + c \cdot -a \to \varphi_{25}) \to \varphi_{16}$}
\BIC{$(c \cdot (x - a) = c \cdot x + c \cdot -a \to \varphi_{25}) \to \varphi_{16}$}
\end{ptree}
\noindent$\Pi_{24}$:
\begin{ptree}
\AXC{$\stackrel{\Pi_{6}}{\vdots}$}\noLine
\UIC{$\varphi_{23}$}
\AXC{$\varphi_{23} \to c \cdot (x - a) = c \cdot x + c \cdot -a \to \varphi_{25}$\lbl{A2}}
\BIC{$c \cdot (x - a) = c \cdot x + c \cdot -a \to \varphi_{25}$}
\AXC{$\stackrel{\Pi_{23}}{\vdots}$}\noLine
\UIC{$(c \cdot (x - a) = c \cdot x + c \cdot -a \to \varphi_{25}) \to \varphi_{16}$}
\BIC{$\varphi_{16}$}
\end{ptree}
\noindent$\Pi_{25}$:
\begin{ptree}
\AXC{$\stackrel{\Pi_{5}}{\vdots}$}\noLine
\UIC{$c \cdot x + -(c \cdot a) = c \cdot x - c \cdot a$}
\AXC{$\stackrel{\Pi_{24}}{\vdots}$}\noLine
\UIC{$\varphi_{16}$}
\AXC{$\varphi_{16} \to c \cdot x + -(c \cdot a) = c \cdot x - c \cdot a \to \varphi_{39}$\lbl{A1}}
\BIC{$c \cdot x + -(c \cdot a) = c \cdot x - c \cdot a \to \varphi_{39}$}
\BIC{$\varphi_{39}$}
\end{ptree}
\noindent 主木:
\begin{ptree}
\AXC{$\stackrel{\Pi_{2}}{\vdots}$}\noLine
\UIC{$\varphi_{6}$}
\AXC{$\stackrel{\Pi_{3}}{\vdots}$}\noLine
\UIC{$\varphi_{22}$}
\AXC{$\varphi_{22} \to c \cdot (x - a) = c \cdot x + -(c \cdot a) \to \varphi_{11}$\lbl{A2}}
\BIC{$c \cdot (x - a) = c \cdot x + -(c \cdot a) \to \varphi_{11}$}
\AXC{$\stackrel{\Pi_{25}}{\vdots}$}\noLine
\UIC{$\varphi_{39}$}
\BIC{$c \in \mathbb{R} \to x \in \mathbb{R} \to a \in \mathbb{R} \to c \cdot (x - a) = c \cdot x - c \cdot a$}
\AXC{$\varphi_{42}$\lbl{A1}}
\BIC{$\varphi_{37}$}
\BIC{$c \in \mathbb{R} \to x \in \mathbb{R} \to a \in \mathbb{R} \to c \cdot x - c \cdot a = c \cdot (x - a)$}
\UIC{$\varphi_{13}$}
\UIC{$\forall x\,(\varphi_{13})$}
\UIC{$\forall c\,\forall x\,(\varphi_{13})$}
\end{ptree}
\begin{abbrevlist}
\item $\varphi_{1} :\equiv \forall w\,(z = w \to c \cdot (x - a) = c \cdot z \to c \cdot (x - a) = c \cdot w)$
\item $\varphi_{2} :\equiv \forall x\,\forall y\,(x = y \to y = x) \to \forall y\,(c \cdot (x - a) = y \to y = c \cdot (x - a))$
\item $\varphi_{3} :\equiv c \cdot (x - a) = c \cdot (x - a) \to c \cdot (x - a) = c \cdot (x + -a)$
\item $\varphi_{4} :\equiv \forall x\,\forall y\,(x - y = x + -y) \to \forall y\,(c \cdot x - y = c \cdot x + -y)$
\item $\varphi_{5} :\equiv \forall w\,(x - a = w \to c \cdot (x - a) = c \cdot (x - a) \to c \cdot (x - a) = c \cdot w)$
\item $\varphi_{6} :\equiv c \cdot (x - a) = c \cdot x - c \cdot a \to c \cdot x - c \cdot a = c \cdot (x - a)$
\item $\varphi_{7} :\equiv \forall x\,\forall y\,(x = y \to y = x) \to \forall y\,(c \cdot x - c \cdot a = y \to y = c \cdot x - c \cdot a)$
\item $\varphi_{8} :\equiv c \in \mathbb{R} \to x \in \mathbb{R} \to a \in \mathbb{R} \to c \cdot (x - a) = c \cdot (x + -a)$
\item $\varphi_{9} :\equiv \forall x\,\forall y\,\forall z\,(x = y \to y = z \to x = z) \to \forall y\,\forall z\,(c \cdot (x - a) = y \to y = z \to c \cdot (x - a) = z)$
\item $\varphi_{10} :\equiv \forall z\,(c \cdot (x - a) = c \cdot (x + -a) \to c \cdot (x + -a) = z \to c \cdot (x - a) = z)$
\item $\varphi_{11} :\equiv c \cdot x + -(c \cdot a) = c \cdot x - c \cdot a \to c \cdot (x - a) = c \cdot x - c \cdot a$
\item $\varphi_{12} :\equiv \forall z\,(c \in \mathbb{R} \to x \in \mathbb{R} \to z \in \mathbb{R} \to c \cdot (x + z) = c \cdot x + c \cdot z)$
\item $\varphi_{13} :\equiv \forall a\,(c \in \mathbb{R} \to x \in \mathbb{R} \to a \in \mathbb{R} \to c \cdot x - c \cdot a = c \cdot (x - a))$
\item $\varphi_{14} :\equiv \forall y\,(c \cdot x - y = c \cdot x + -y) \to c \cdot x - c \cdot a = c \cdot x + -(c \cdot a)$
\item $\varphi_{15} :\equiv c \in \mathbb{R} \to x \in \mathbb{R} \to a \in \mathbb{R} \to c \cdot (x - a) = c \cdot x + c \cdot -a$
\item $\varphi_{16} :\equiv c \in \mathbb{R} \to x \in \mathbb{R} \to a \in \mathbb{R} \to c \cdot (x - a) = c \cdot x + -(c \cdot a)$
\item $\varphi_{17} :\equiv \forall y\,\forall z\,(c \in \mathbb{R} \to y \in \mathbb{R} \to z \in \mathbb{R} \to c \cdot (y + z) = c \cdot y + c \cdot z)$
\item $\varphi_{18} :\equiv c \cdot (x + -a) = c \cdot x + c \cdot -a \to c \cdot (x - a) = c \cdot x + c \cdot -a$
\item $\varphi_{19} :\equiv z = w \to c \cdot x + c \cdot -a = c \cdot x + z \to c \cdot x + c \cdot -a = c \cdot x + w$
\item $\varphi_{20} :\equiv \forall x\,\forall y\,\forall z\,(x \in \mathbb{R} \to y \in \mathbb{R} \to z \in \mathbb{R} \to x \cdot (y + z) = x \cdot y + x \cdot z)$
\item $\varphi_{21} :\equiv c \in \mathbb{R} \to x \in \mathbb{R} \to a \in \mathbb{R} \to c \cdot (x + -a) = c \cdot x + c \cdot -a$
\item $\varphi_{22} :\equiv \forall z\,(c \cdot (x - a) = c \cdot x + -(c \cdot a) \to c \cdot x + -(c \cdot a) = z \to c \cdot (x - a) = z)$
\item $\varphi_{23} :\equiv \forall z\,(c \cdot (x - a) = c \cdot x + c \cdot -a \to c \cdot x + c \cdot -a = z \to c \cdot (x - a) = z)$
\item $\varphi_{24} :\equiv c \cdot x - c \cdot a = c \cdot x + -(c \cdot a) \to c \cdot x + -(c \cdot a) = c \cdot x - c \cdot a$
\item $\varphi_{25} :\equiv c \cdot x + c \cdot -a = c \cdot x + -(c \cdot a) \to c \cdot (x - a) = c \cdot x + -(c \cdot a)$
\item $\varphi_{26} :\equiv c \in \mathbb{R} \to x \in \mathbb{R} \to a \in \mathbb{R} \to c \cdot x + c \cdot -a = c \cdot x + -(c \cdot a)$
\item $\varphi_{27} :\equiv c \in \mathbb{R} \to x \in \mathbb{R} \to -a \in \mathbb{R} \to c \cdot (x + -a) = c \cdot x + c \cdot -a$
\item $\varphi_{28} :\equiv c \cdot x + c \cdot -a = c \cdot x + c \cdot -a \to c \cdot x + c \cdot -a = c \cdot x + -(c \cdot a)$
\item $\varphi_{29} :\equiv \forall w\,(c \cdot -a = w \to c \cdot x + c \cdot -a = c \cdot x + c \cdot -a \to c \cdot x + c \cdot -a = c \cdot x + w)$
\item $\varphi_{30} :\equiv \forall y\,(c \in \mathbb{R} \to y \in \mathbb{R} \to c \cdot -y = -(c \cdot y)) \to c \in \mathbb{R} \to a \in \mathbb{R} \to c \cdot -a = -(c \cdot a)$
\item $\varphi_{31} :\equiv (c \in \mathbb{R} \to x \in \mathbb{R} \to a \in \mathbb{R} \to a \in \mathbb{R}) \to c \in \mathbb{R} \to x \in \mathbb{R} \to a \in \mathbb{R} \to -a \in \mathbb{R}$
\item $\varphi_{32} :\equiv \forall x\,\forall y\,(x \in \mathbb{R} \to y \in \mathbb{R} \to x \cdot -y = -(x \cdot y)) \to \forall y\,(c \in \mathbb{R} \to y \in \mathbb{R} \to c \cdot -y = -(c \cdot y))$
\item $\varphi_{33} :\equiv \forall y\,\forall z\,(c \cdot (x - a) = y \to y = z \to c \cdot (x - a) = z) \to \varphi_{10}$
\item $\varphi_{34} :\equiv (c \in \mathbb{R} \to x \in \mathbb{R} \to a \in \mathbb{R} \to a \in \mathbb{R}) \to c \in \mathbb{R} \to x \in \mathbb{R} \to a \in \mathbb{R} \to c \cdot -a = -(c \cdot a)$
\item $\varphi_{35} :\equiv \forall y\,\forall z\,(c \cdot (x - a) = y \to y = z \to c \cdot (x - a) = z) \to \varphi_{22}$
\item $\varphi_{36} :\equiv \forall y\,\forall z\,(c \cdot (x - a) = y \to y = z \to c \cdot (x - a) = z) \to \varphi_{23}$
\item $\varphi_{37} :\equiv \varphi_{6} \to c \in \mathbb{R} \to x \in \mathbb{R} \to a \in \mathbb{R} \to c \cdot x - c \cdot a = c \cdot (x - a)$
\item $\varphi_{38} :\equiv c \cdot x - c \cdot a = c \cdot x + -(c \cdot a) \to \varphi_{24} \to c \cdot x + -(c \cdot a) = c \cdot x - c \cdot a$
\item $\varphi_{39} :\equiv (c \cdot (x - a) = c \cdot x + -(c \cdot a) \to \varphi_{11}) \to c \in \mathbb{R} \to x \in \mathbb{R} \to a \in \mathbb{R} \to c \cdot (x - a) = c \cdot x - c \cdot a$
\item $\varphi_{40} :\equiv (c \in \mathbb{R} \to x \in \mathbb{R} \to a \in \mathbb{R} \to x \in \mathbb{R}) \to (c \in \mathbb{R} \to x \in \mathbb{R} \to a \in \mathbb{R} \to -a \in \mathbb{R}) \to \varphi_{21}$
\item $\varphi_{41} :\equiv x - a = x + -a \to c \cdot (x - a) = c \cdot (x - a) \to (x - a = x + -a \to \varphi_{3}) \to \varphi_{8}$
\item $\varphi_{42} :\equiv (c \in \mathbb{R} \to x \in \mathbb{R} \to a \in \mathbb{R} \to c \cdot (x - a) = c \cdot x - c \cdot a) \to \varphi_{37}$
\item $\varphi_{43} :\equiv (c \in \mathbb{R} \to a \in \mathbb{R} \to c \cdot -a = -(c \cdot a)) \to (c \in \mathbb{R} \to x \in \mathbb{R} \to a \in \mathbb{R} \to c \in \mathbb{R}) \to \varphi_{34}$
\item $\varphi_{44} :\equiv c \cdot x + c \cdot -a = c \cdot x + c \cdot -a \to (c \cdot -a = -(c \cdot a) \to \varphi_{28}) \to \varphi_{26}$
\end{abbrevlist}
\end{fullproof}

% 自動生成 (prover/build.js): MulSubSplit  ノード数 308
\begin{fullproof}{[F18]}{$\forall x\,\forall y\,\forall a\,\forall b\,(x \in \mathbb{R} \to y \in \mathbb{R} \to a \in \mathbb{R} \to b \in \mathbb{R} \to x \cdot y - a \cdot b = x \cdot (y - b) + b \cdot (x - a))$}{308}\phantomsection\label{tree:MulSubSplit}
\noindent$\Pi_{1}$:
\begin{ptree}
\AXC{$x = x$\lbl{A5}}
\AXC{$x = y \to x = x \to y = x$\lbl{A6}}
\AXC{$(x = y \to x = x \to y = x) \to x = x \to x = y \to y = x$\lbl{A1}}
\BIC{$x = x \to x = y \to y = x$}
\BIC{$x = y \to y = x$}
\UIC{$\forall y\,(x = y \to y = x)$}
\UIC{$\forall x\,\forall y\,(x = y \to y = x)$}
\end{ptree}
\noindent$\Pi_{2}$:
\begin{ptree}
\AXC{$\stackrel{\Pi_{1}}{\vdots}$}\noLine
\UIC{$\forall x\,\forall y\,(x = y \to y = x)$}
\AXC{$\forall x\,\forall y\,(x = y \to y = x) \to \forall y\,(v_{217} = y \to y = v_{217})$\lbl{A2}}
\BIC{$\forall y\,(v_{217} = y \to y = v_{217})$}
\AXC{$\varphi_{1}$\lbl{A2}}
\BIC{$v_{217} = v_{218} \to v_{218} = v_{217}$}
\UIC{$\forall v_{218}\,(v_{217} = v_{218} \to v_{218} = v_{217})$}
\UIC{$\forall v_{217}\,\forall v_{218}\,(v_{217} = v_{218} \to v_{218} = v_{217})$}
\end{ptree}
\noindent$\Pi_{3}$:
\begin{ptree}
\AXC{$\stackrel{\Pi_{2}}{\vdots}$}\noLine
\UIC{$\forall v_{217}\,\forall v_{218}\,(v_{217} = v_{218} \to v_{218} = v_{217})$}
\AXC{$\forall v_{217}\,\forall v_{218}\,(v_{217} = v_{218} \to v_{218} = v_{217}) \to \varphi_{39}$\lbl{A2}}
\BIC{$\varphi_{39}$}
\AXC{$\varphi_{39} \to \varphi_{64}$\lbl{A2}}
\BIC{$\varphi_{64}$}
\end{ptree}
\noindent$\Pi_{4}$:
\begin{ptree}
\AXC{$y = z \to x = y \to x = z$\lbl{A6}}
\AXC{$(y = z \to x = y \to x = z) \to x = y \to y = z \to x = z$\lbl{A1}}
\BIC{$x = y \to y = z \to x = z$}
\UIC{$\forall z\,(x = y \to y = z \to x = z)$}
\UIC{$\forall y\,\forall z\,(x = y \to y = z \to x = z)$}
\UIC{$\forall x\,\forall y\,\forall z\,(x = y \to y = z \to x = z)$}
\AXC{$\varphi_{9}$\lbl{A2}}
\BIC{$\forall y\,\forall z\,(v_{214} = y \to y = z \to v_{214} = z)$}
\AXC{$\varphi_{17}$\lbl{A2}}
\BIC{$\forall z\,(v_{214} = v_{215} \to v_{215} = z \to v_{214} = z)$}
\AXC{$\varphi_{27}$\lbl{A2}}
\BIC{$v_{214} = v_{215} \to v_{215} = v_{216} \to v_{214} = v_{216}$}
\UIC{$\forall v_{216}\,(v_{214} = v_{215} \to v_{215} = v_{216} \to v_{214} = v_{216})$}
\UIC{$\varphi_{4}$}
\UIC{$\forall v_{214}\,(\varphi_{4})$}
\end{ptree}
\noindent$\Pi_{5}$:
\begin{ptree}
\AXC{$\stackrel{\Pi_{4}}{\vdots}$}\noLine
\UIC{$\forall v_{214}\,(\varphi_{4})$}
\AXC{$\forall v_{214}\,(\varphi_{4}) \to \varphi_{65}$\lbl{A2}}
\BIC{$\varphi_{65}$}
\AXC{$\varphi_{65} \to \varphi_{90}$\lbl{A2}}
\BIC{$\varphi_{90}$}
\AXC{$\varphi_{90} \to \varphi_{23} \to \varphi_{75}$\lbl{A2}}
\BIC{$\varphi_{23} \to \varphi_{75}$}
\end{ptree}
\noindent$\Pi_{6}$:
\begin{ptree}
\AXC{$\stackrel{\Pi_{1}}{\vdots}$}\noLine
\UIC{$\forall x\,\forall y\,(x = y \to y = x)$}
\AXC{$\forall x\,\forall y\,(x = y \to y = x) \to \forall y\,(v_{201} = y \to y = v_{201})$\lbl{A2}}
\BIC{$\forall y\,(v_{201} = y \to y = v_{201})$}
\AXC{$\varphi_{2}$\lbl{A2}}
\BIC{$v_{201} = v_{202} \to v_{202} = v_{201}$}
\UIC{$\forall v_{202}\,(v_{201} = v_{202} \to v_{202} = v_{201})$}
\UIC{$\forall v_{201}\,\forall v_{202}\,(v_{201} = v_{202} \to v_{202} = v_{201})$}
\AXC{$\varphi_{59}$\lbl{A2}}
\BIC{$\forall v_{202}\,(x \cdot y - a \cdot b = v_{202} \to v_{202} = x \cdot y - a \cdot b)$}
\AXC{$\varphi_{94}$\lbl{A2}}
\BIC{$\varphi_{78}$}
\end{ptree}
\noindent$\Pi_{7}$:
\begin{ptree}
\AXC{$\forall x\,\forall y\,(x \in \mathbb{R} \to y \in \mathbb{R} \to x \cdot y \in \mathbb{R})$\lbl{(R1)}}
\AXC{$\varphi_{72}$\lbl{A2}}
\BIC{$\forall y\,(x \in \mathbb{R} \to y \in \mathbb{R} \to x \cdot y \in \mathbb{R})$}
\AXC{$\varphi_{61}$\lbl{A2}}
\BIC{$x \in \mathbb{R} \to b \in \mathbb{R} \to x \cdot b \in \mathbb{R}$}
\AXC{$\varphi_{102}$\lbl{A1}}
\BIC{$(x \in \mathbb{R} \to y \in \mathbb{R} \to a \in \mathbb{R} \to b \in \mathbb{R} \to x \in \mathbb{R}) \to \varphi_{84}$}
\end{ptree}
\noindent$\Pi_{8}$:
\begin{ptree}
\AXC{$x \in \mathbb{R} \to y \in \mathbb{R} \to a \in \mathbb{R} \to b \in \mathbb{R} \to x \in \mathbb{R}$\lbl{A1}}
\AXC{$\stackrel{\Pi_{7}}{\vdots}$}\noLine
\UIC{$(x \in \mathbb{R} \to y \in \mathbb{R} \to a \in \mathbb{R} \to b \in \mathbb{R} \to x \in \mathbb{R}) \to \varphi_{84}$}
\BIC{$\varphi_{84}$}
\end{ptree}
\noindent$\Pi_{9}$:
\begin{ptree}
\AXC{$\forall x\,\forall y\,(x \in \mathbb{R} \to y \in \mathbb{R} \to x \cdot y \in \mathbb{R})$\lbl{(R1)}}
\AXC{$\varphi_{73}$\lbl{A2}}
\BIC{$\forall y\,(a \in \mathbb{R} \to y \in \mathbb{R} \to a \cdot y \in \mathbb{R})$}
\AXC{$\varphi_{62}$\lbl{A2}}
\BIC{$a \in \mathbb{R} \to b \in \mathbb{R} \to a \cdot b \in \mathbb{R}$}
\AXC{$\varphi_{103}$\lbl{A1}}
\BIC{$(x \in \mathbb{R} \to y \in \mathbb{R} \to a \in \mathbb{R} \to b \in \mathbb{R} \to a \in \mathbb{R}) \to \varphi_{85}$}
\end{ptree}
\noindent$\Pi_{10}$:
\begin{ptree}
\AXC{$x \in \mathbb{R} \to y \in \mathbb{R} \to a \in \mathbb{R} \to b \in \mathbb{R} \to a \in \mathbb{R}$\lbl{A1}}
\AXC{$\stackrel{\Pi_{9}}{\vdots}$}\noLine
\UIC{$(x \in \mathbb{R} \to y \in \mathbb{R} \to a \in \mathbb{R} \to b \in \mathbb{R} \to a \in \mathbb{R}) \to \varphi_{85}$}
\BIC{$\varphi_{85}$}
\end{ptree}
\noindent$\Pi_{11}$:
\begin{ptree}
\AXC{$\forall x\,\forall y\,(x \in \mathbb{R} \to y \in \mathbb{R} \to x \cdot y \in \mathbb{R})$\lbl{(R1)}}
\AXC{$\varphi_{72}$\lbl{A2}}
\BIC{$\forall y\,(x \in \mathbb{R} \to y \in \mathbb{R} \to x \cdot y \in \mathbb{R})$}
\AXC{$\varphi_{63}$\lbl{A2}}
\BIC{$x \in \mathbb{R} \to y \in \mathbb{R} \to x \cdot y \in \mathbb{R}$}
\AXC{$\varphi_{104}$\lbl{A1}}
\BIC{$(x \in \mathbb{R} \to y \in \mathbb{R} \to a \in \mathbb{R} \to b \in \mathbb{R} \to x \in \mathbb{R}) \to \varphi_{86}$}
\end{ptree}
\noindent$\Pi_{12}$:
\begin{ptree}
\AXC{$x \in \mathbb{R} \to y \in \mathbb{R} \to a \in \mathbb{R} \to b \in \mathbb{R} \to x \in \mathbb{R}$\lbl{A1}}
\AXC{$\stackrel{\Pi_{11}}{\vdots}$}\noLine
\UIC{$(x \in \mathbb{R} \to y \in \mathbb{R} \to a \in \mathbb{R} \to b \in \mathbb{R} \to x \in \mathbb{R}) \to \varphi_{86}$}
\BIC{$\varphi_{86}$}
\end{ptree}
\noindent$\Pi_{13}$:
\begin{ptree}
\AXC{$\varphi_{40}$\lbl{[F16]}}
\AXC{$\varphi_{40} \to \varphi_{47}$\lbl{A2}}
\BIC{$\varphi_{47}$}
\AXC{$\varphi_{47} \to \varphi_{55}$\lbl{A2}}
\BIC{$\varphi_{55}$}
\AXC{$\varphi_{55} \to v_{198} \in \mathbb{R} \to v_{199} \in \mathbb{R} \to \varphi_{13}$\lbl{A2}}
\BIC{$v_{198} \in \mathbb{R} \to v_{199} \in \mathbb{R} \to \varphi_{13}$}
\UIC{$\forall v_{200}\,(v_{198} \in \mathbb{R} \to v_{199} \in \mathbb{R} \to \varphi_{13})$}
\UIC{$\forall v_{199}\,\forall v_{200}\,(v_{198} \in \mathbb{R} \to v_{199} \in \mathbb{R} \to \varphi_{13})$}
\UIC{$\forall v_{198}\,\forall v_{199}\,\forall v_{200}\,(v_{198} \in \mathbb{R} \to v_{199} \in \mathbb{R} \to \varphi_{13})$}
\AXC{$\varphi_{99}$\lbl{A2}}
\BIC{$\varphi_{76}$}
\AXC{$\varphi_{76} \to \varphi_{77}$\lbl{A2}}
\BIC{$\varphi_{77}$}
\end{ptree}
\noindent$\Pi_{14}$:
\begin{ptree}
\AXC{$\stackrel{\Pi_{13}}{\vdots}$}\noLine
\UIC{$\varphi_{77}$}
\AXC{$\varphi_{77} \to x \cdot y \in \mathbb{R} \to a \cdot b \in \mathbb{R} \to \varphi_{41}$\lbl{A2}}
\BIC{$x \cdot y \in \mathbb{R} \to a \cdot b \in \mathbb{R} \to \varphi_{41}$}
\AXC{$(x \cdot y \in \mathbb{R} \to a \cdot b \in \mathbb{R} \to \varphi_{41}) \to \varphi_{109}$\lbl{A1}}
\BIC{$\varphi_{109}$}
\end{ptree}
\noindent$\Pi_{15}$:
\begin{ptree}
\AXC{$x \in \mathbb{R} \to y \in \mathbb{R} \to a \in \mathbb{R} \to b \in \mathbb{R} \to y \in \mathbb{R}$\lbl{A1}}
\AXC{$\stackrel{\Pi_{12}}{\vdots}$}\noLine
\UIC{$\varphi_{86}$}
\BIC{$x \in \mathbb{R} \to y \in \mathbb{R} \to a \in \mathbb{R} \to b \in \mathbb{R} \to x \cdot y \in \mathbb{R}$}
\AXC{$\stackrel{\Pi_{14}}{\vdots}$}\noLine
\UIC{$\varphi_{109}$}
\BIC{$(x \in \mathbb{R} \to y \in \mathbb{R} \to a \in \mathbb{R} \to b \in \mathbb{R} \to a \cdot b \in \mathbb{R}) \to \varphi_{95}$}
\end{ptree}
\noindent$\Pi_{16}$:
\begin{ptree}
\AXC{$x \in \mathbb{R} \to y \in \mathbb{R} \to a \in \mathbb{R} \to b \in \mathbb{R} \to b \in \mathbb{R}$\lbl{A1}}
\AXC{$\stackrel{\Pi_{10}}{\vdots}$}\noLine
\UIC{$\varphi_{85}$}
\BIC{$x \in \mathbb{R} \to y \in \mathbb{R} \to a \in \mathbb{R} \to b \in \mathbb{R} \to a \cdot b \in \mathbb{R}$}
\AXC{$\stackrel{\Pi_{15}}{\vdots}$}\noLine
\UIC{$(x \in \mathbb{R} \to y \in \mathbb{R} \to a \in \mathbb{R} \to b \in \mathbb{R} \to a \cdot b \in \mathbb{R}) \to \varphi_{95}$}
\BIC{$\varphi_{95}$}
\end{ptree}
\noindent$\Pi_{17}$:
\begin{ptree}
\AXC{$x \in \mathbb{R} \to y \in \mathbb{R} \to a \in \mathbb{R} \to b \in \mathbb{R} \to b \in \mathbb{R}$\lbl{A1}}
\AXC{$\stackrel{\Pi_{8}}{\vdots}$}\noLine
\UIC{$\varphi_{84}$}
\BIC{$x \in \mathbb{R} \to y \in \mathbb{R} \to a \in \mathbb{R} \to b \in \mathbb{R} \to x \cdot b \in \mathbb{R}$}
\AXC{$\stackrel{\Pi_{16}}{\vdots}$}\noLine
\UIC{$\varphi_{95}$}
\BIC{$x \in \mathbb{R} \to y \in \mathbb{R} \to \varphi_{48}$}
\AXC{$(x \in \mathbb{R} \to y \in \mathbb{R} \to \varphi_{48}) \to \varphi_{78} \to x \in \mathbb{R} \to y \in \mathbb{R} \to \varphi_{49}$\lbl{A1}}
\BIC{$\varphi_{78} \to x \in \mathbb{R} \to y \in \mathbb{R} \to \varphi_{49}$}
\end{ptree}
\noindent$\Pi_{18}$:
\begin{ptree}
\AXC{$y = z \to x = y \to x = z$\lbl{A6}}
\AXC{$(y = z \to x = y \to x = z) \to x = y \to y = z \to x = z$\lbl{A1}}
\BIC{$x = y \to y = z \to x = z$}
\UIC{$\forall z\,(x = y \to y = z \to x = z)$}
\UIC{$\forall y\,\forall z\,(x = y \to y = z \to x = z)$}
\UIC{$\forall x\,\forall y\,\forall z\,(x = y \to y = z \to x = z)$}
\AXC{$\varphi_{10}$\lbl{A2}}
\BIC{$\forall y\,\forall z\,(v_{211} = y \to y = z \to v_{211} = z)$}
\AXC{$\varphi_{18}$\lbl{A2}}
\BIC{$\forall z\,(v_{211} = v_{212} \to v_{212} = z \to v_{211} = z)$}
\AXC{$\varphi_{28}$\lbl{A2}}
\BIC{$v_{211} = v_{212} \to v_{212} = v_{213} \to v_{211} = v_{213}$}
\UIC{$\forall v_{213}\,(v_{211} = v_{212} \to v_{212} = v_{213} \to v_{211} = v_{213})$}
\UIC{$\varphi_{5}$}
\UIC{$\forall v_{211}\,(\varphi_{5})$}
\end{ptree}
\noindent$\Pi_{19}$:
\begin{ptree}
\AXC{$\stackrel{\Pi_{18}}{\vdots}$}\noLine
\UIC{$\forall v_{211}\,(\varphi_{5})$}
\AXC{$\forall v_{211}\,(\varphi_{5}) \to \varphi_{66}$\lbl{A2}}
\BIC{$\varphi_{66}$}
\AXC{$\varphi_{66} \to \varphi_{91}$\lbl{A2}}
\BIC{$\varphi_{91}$}
\AXC{$\varphi_{91} \to \varphi_{24} \to \varphi_{44} \to \varphi_{23}$\lbl{A2}}
\BIC{$\varphi_{24} \to \varphi_{44} \to \varphi_{23}$}
\end{ptree}
\noindent$\Pi_{20}$:
\begin{ptree}
\AXC{$z = w \to \varphi_{32} \to \varphi_{31}$\lbl{A6}}
\UIC{$\forall w\,(z = w \to \varphi_{32} \to \varphi_{31})$}
\UIC{$\forall z\,\forall w\,(z = w \to \varphi_{32} \to \varphi_{31})$}
\AXC{$\forall z\,\forall w\,(z = w \to \varphi_{32} \to \varphi_{31}) \to \forall w\,(b \cdot a = w \to \varphi_{42} \to \varphi_{31})$\lbl{A2}}
\BIC{$\forall w\,(b \cdot a = w \to \varphi_{42} \to \varphi_{31})$}
\end{ptree}
\noindent$\Pi_{21}$:
\begin{ptree}
\AXC{$\forall x\,\forall y\,(x \in \mathbb{R} \to y \in \mathbb{R} \to x \cdot y = y \cdot x)$\lbl{(R7)}}
\AXC{$\varphi_{70}$\lbl{A2}}
\BIC{$\forall y\,(b \in \mathbb{R} \to y \in \mathbb{R} \to b \cdot y = y \cdot b)$}
\AXC{$\varphi_{56}$\lbl{A2}}
\BIC{$b \in \mathbb{R} \to a \in \mathbb{R} \to b \cdot a = a \cdot b$}
\AXC{$\varphi_{100}$\lbl{A1}}
\BIC{$(x \in \mathbb{R} \to y \in \mathbb{R} \to a \in \mathbb{R} \to b \in \mathbb{R} \to b \in \mathbb{R}) \to \varphi_{82}$}
\end{ptree}
\noindent$\Pi_{22}$:
\begin{ptree}
\AXC{$x \in \mathbb{R} \to y \in \mathbb{R} \to a \in \mathbb{R} \to b \in \mathbb{R} \to b \in \mathbb{R}$\lbl{A1}}
\AXC{$\stackrel{\Pi_{21}}{\vdots}$}\noLine
\UIC{$(x \in \mathbb{R} \to y \in \mathbb{R} \to a \in \mathbb{R} \to b \in \mathbb{R} \to b \in \mathbb{R}) \to \varphi_{82}$}
\BIC{$\varphi_{82}$}
\end{ptree}
\noindent$\Pi_{23}$:
\begin{ptree}
\AXC{$x = x$\lbl{A5}}
\UIC{$\forall x\,(x = x)$}
\AXC{$\forall x\,(x = x) \to \varphi_{42}$\lbl{A2}}
\BIC{$\varphi_{42}$}
\AXC{$x \in \mathbb{R} \to y \in \mathbb{R} \to a \in \mathbb{R} \to b \in \mathbb{R} \to a \in \mathbb{R}$\lbl{A1}}
\AXC{$\stackrel{\Pi_{22}}{\vdots}$}\noLine
\UIC{$\varphi_{82}$}
\BIC{$x \in \mathbb{R} \to y \in \mathbb{R} \to a \in \mathbb{R} \to b \in \mathbb{R} \to b \cdot a = a \cdot b$}
\AXC{$\varphi_{114}$\lbl{A1}}
\BIC{$\varphi_{42} \to \varphi_{107}$}
\BIC{$\varphi_{107}$}
\end{ptree}
\noindent$\Pi_{24}$:
\begin{ptree}
\AXC{$y = z \to x = y \to x = z$\lbl{A6}}
\AXC{$(y = z \to x = y \to x = z) \to x = y \to y = z \to x = z$\lbl{A1}}
\BIC{$x = y \to y = z \to x = z$}
\UIC{$\forall z\,(x = y \to y = z \to x = z)$}
\UIC{$\forall y\,\forall z\,(x = y \to y = z \to x = z)$}
\UIC{$\forall x\,\forall y\,\forall z\,(x = y \to y = z \to x = z)$}
\AXC{$\varphi_{11}$\lbl{A2}}
\BIC{$\forall y\,\forall z\,(v_{208} = y \to y = z \to v_{208} = z)$}
\AXC{$\varphi_{19}$\lbl{A2}}
\BIC{$\forall z\,(v_{208} = v_{209} \to v_{209} = z \to v_{208} = z)$}
\AXC{$\varphi_{29}$\lbl{A2}}
\BIC{$v_{208} = v_{209} \to v_{209} = v_{210} \to v_{208} = v_{210}$}
\UIC{$\forall v_{210}\,(v_{208} = v_{209} \to v_{209} = v_{210} \to v_{208} = v_{210})$}
\UIC{$\varphi_{6}$}
\UIC{$\forall v_{208}\,(\varphi_{6})$}
\end{ptree}
\noindent$\Pi_{25}$:
\begin{ptree}
\AXC{$\stackrel{\Pi_{24}}{\vdots}$}\noLine
\UIC{$\forall v_{208}\,(\varphi_{6})$}
\AXC{$\forall v_{208}\,(\varphi_{6}) \to \varphi_{67}$\lbl{A2}}
\BIC{$\varphi_{67}$}
\AXC{$\varphi_{67} \to \varphi_{92}$\lbl{A2}}
\BIC{$\varphi_{92}$}
\AXC{$\varphi_{92} \to \varphi_{25} \to \varphi_{45} \to \varphi_{24}$\lbl{A2}}
\BIC{$\varphi_{25} \to \varphi_{45} \to \varphi_{24}$}
\end{ptree}
\noindent$\Pi_{26}$:
\begin{ptree}
\AXC{$z = w \to \varphi_{34} \to \varphi_{33}$\lbl{A6}}
\UIC{$\forall w\,(z = w \to \varphi_{34} \to \varphi_{33})$}
\UIC{$\forall z\,\forall w\,(z = w \to \varphi_{34} \to \varphi_{33})$}
\AXC{$\forall z\,\forall w\,(z = w \to \varphi_{34} \to \varphi_{33}) \to \forall w\,(b \cdot x = w \to \varphi_{43} \to \varphi_{33})$\lbl{A2}}
\BIC{$\forall w\,(b \cdot x = w \to \varphi_{43} \to \varphi_{33})$}
\end{ptree}
\noindent$\Pi_{27}$:
\begin{ptree}
\AXC{$\forall x\,\forall y\,(x \in \mathbb{R} \to y \in \mathbb{R} \to x \cdot y = y \cdot x)$\lbl{(R7)}}
\AXC{$\varphi_{70}$\lbl{A2}}
\BIC{$\forall y\,(b \in \mathbb{R} \to y \in \mathbb{R} \to b \cdot y = y \cdot b)$}
\AXC{$\varphi_{57}$\lbl{A2}}
\BIC{$b \in \mathbb{R} \to x \in \mathbb{R} \to b \cdot x = x \cdot b$}
\AXC{$\varphi_{101}$\lbl{A1}}
\BIC{$(x \in \mathbb{R} \to y \in \mathbb{R} \to a \in \mathbb{R} \to b \in \mathbb{R} \to b \in \mathbb{R}) \to \varphi_{83}$}
\end{ptree}
\noindent$\Pi_{28}$:
\begin{ptree}
\AXC{$x \in \mathbb{R} \to y \in \mathbb{R} \to a \in \mathbb{R} \to b \in \mathbb{R} \to b \in \mathbb{R}$\lbl{A1}}
\AXC{$\stackrel{\Pi_{27}}{\vdots}$}\noLine
\UIC{$(x \in \mathbb{R} \to y \in \mathbb{R} \to a \in \mathbb{R} \to b \in \mathbb{R} \to b \in \mathbb{R}) \to \varphi_{83}$}
\BIC{$\varphi_{83}$}
\end{ptree}
\noindent$\Pi_{29}$:
\begin{ptree}
\AXC{$x = x$\lbl{A5}}
\UIC{$\forall x\,(x = x)$}
\AXC{$\forall x\,(x = x) \to \varphi_{43}$\lbl{A2}}
\BIC{$\varphi_{43}$}
\AXC{$x \in \mathbb{R} \to y \in \mathbb{R} \to a \in \mathbb{R} \to b \in \mathbb{R} \to x \in \mathbb{R}$\lbl{A1}}
\AXC{$\stackrel{\Pi_{28}}{\vdots}$}\noLine
\UIC{$\varphi_{83}$}
\BIC{$x \in \mathbb{R} \to y \in \mathbb{R} \to a \in \mathbb{R} \to b \in \mathbb{R} \to b \cdot x = x \cdot b$}
\AXC{$\varphi_{115}$\lbl{A1}}
\BIC{$\varphi_{43} \to \varphi_{108}$}
\BIC{$\varphi_{108}$}
\end{ptree}
\noindent$\Pi_{30}$:
\begin{ptree}
\AXC{$y = z \to x = y \to x = z$\lbl{A6}}
\AXC{$(y = z \to x = y \to x = z) \to x = y \to y = z \to x = z$\lbl{A1}}
\BIC{$x = y \to y = z \to x = z$}
\UIC{$\forall z\,(x = y \to y = z \to x = z)$}
\UIC{$\forall y\,\forall z\,(x = y \to y = z \to x = z)$}
\UIC{$\forall x\,\forall y\,\forall z\,(x = y \to y = z \to x = z)$}
\AXC{$\varphi_{12}$\lbl{A2}}
\BIC{$\forall y\,\forall z\,(v_{205} = y \to y = z \to v_{205} = z)$}
\AXC{$\varphi_{20}$\lbl{A2}}
\BIC{$\forall z\,(v_{205} = v_{206} \to v_{206} = z \to v_{205} = z)$}
\AXC{$\varphi_{30}$\lbl{A2}}
\BIC{$v_{205} = v_{206} \to v_{206} = v_{207} \to v_{205} = v_{207}$}
\UIC{$\forall v_{207}\,(v_{205} = v_{206} \to v_{206} = v_{207} \to v_{205} = v_{207})$}
\UIC{$\varphi_{7}$}
\UIC{$\forall v_{205}\,(\varphi_{7})$}
\end{ptree}
\noindent$\Pi_{31}$:
\begin{ptree}
\AXC{$\stackrel{\Pi_{30}}{\vdots}$}\noLine
\UIC{$\forall v_{205}\,(\varphi_{7})$}
\AXC{$\forall v_{205}\,(\varphi_{7}) \to \varphi_{68}$\lbl{A2}}
\BIC{$\varphi_{68}$}
\AXC{$\varphi_{68} \to \varphi_{87}$\lbl{A2}}
\BIC{$\varphi_{87}$}
\AXC{$\varphi_{87} \to \varphi_{14} \to \varphi_{35} \to \varphi_{25}$\lbl{A2}}
\BIC{$\varphi_{14} \to \varphi_{35} \to \varphi_{25}$}
\end{ptree}
\noindent$\Pi_{32}$:
\begin{ptree}
\AXC{$z = w \to \varphi_{79}$\lbl{A6}}
\UIC{$\forall w\,(z = w \to \varphi_{79})$}
\UIC{$\forall z\,\forall w\,(z = w \to \varphi_{79})$}
\AXC{$\forall z\,\forall w\,(z = w \to \varphi_{79}) \to \varphi_{88}$\lbl{A2}}
\BIC{$\varphi_{88}$}
\AXC{$\varphi_{88} \to b \cdot (x - a) = b \cdot x - b \cdot a \to \varphi_{22} \to \varphi_{35}$\lbl{A2}}
\BIC{$b \cdot (x - a) = b \cdot x - b \cdot a \to \varphi_{22} \to \varphi_{35}$}
\end{ptree}
\noindent$\Pi_{33}$:
\begin{ptree}
\AXC{$\varphi_{46}$\lbl{[F17]}}
\AXC{$\varphi_{46} \to \forall x\,(\varphi_{38})$\lbl{A2}}
\BIC{$\forall x\,(\varphi_{38})$}
\AXC{$\forall x\,(\varphi_{38}) \to \varphi_{38}$\lbl{A2}}
\BIC{$\varphi_{38}$}
\AXC{$\varphi_{89}$\lbl{A2}}
\BIC{$b \in \mathbb{R} \to x \in \mathbb{R} \to a \in \mathbb{R} \to b \cdot x - b \cdot a = b \cdot (x - a)$}
\AXC{$\varphi_{110}$\lbl{A1}}
\BIC{$(x \in \mathbb{R} \to y \in \mathbb{R} \to a \in \mathbb{R} \to b \in \mathbb{R} \to b \in \mathbb{R}) \to \varphi_{96}$}
\end{ptree}
\noindent$\Pi_{34}$:
\begin{ptree}
\AXC{$x \in \mathbb{R} \to y \in \mathbb{R} \to a \in \mathbb{R} \to b \in \mathbb{R} \to b \in \mathbb{R}$\lbl{A1}}
\AXC{$\stackrel{\Pi_{33}}{\vdots}$}\noLine
\UIC{$(x \in \mathbb{R} \to y \in \mathbb{R} \to a \in \mathbb{R} \to b \in \mathbb{R} \to b \in \mathbb{R}) \to \varphi_{96}$}
\BIC{$\varphi_{96}$}
\end{ptree}
\noindent$\Pi_{35}$:
\begin{ptree}
\AXC{$x \in \mathbb{R} \to y \in \mathbb{R} \to a \in \mathbb{R} \to b \in \mathbb{R} \to a \in \mathbb{R}$\lbl{A1}}
\AXC{$x \in \mathbb{R} \to y \in \mathbb{R} \to a \in \mathbb{R} \to b \in \mathbb{R} \to x \in \mathbb{R}$\lbl{A1}}
\AXC{$\stackrel{\Pi_{34}}{\vdots}$}\noLine
\UIC{$\varphi_{96}$}
\BIC{$(x \in \mathbb{R} \to y \in \mathbb{R} \to a \in \mathbb{R} \to b \in \mathbb{R} \to a \in \mathbb{R}) \to \varphi_{50}$}
\BIC{$\varphi_{50}$}
\end{ptree}
\noindent$\Pi_{36}$:
\begin{ptree}
\AXC{$\stackrel{\Pi_{1}}{\vdots}$}\noLine
\UIC{$\forall x\,\forall y\,(x = y \to y = x)$}
\AXC{$\varphi_{21}$\lbl{A2}}
\BIC{$\forall y\,(b \cdot x - b \cdot a = y \to y = b \cdot x - b \cdot a)$}
\AXC{$\forall y\,(b \cdot x - b \cdot a = y \to y = b \cdot x - b \cdot a) \to \varphi_{15}$\lbl{A2}}
\BIC{$\varphi_{15}$}
\AXC{$\stackrel{\Pi_{35}}{\vdots}$}\noLine
\UIC{$\varphi_{50}$}
\AXC{$\varphi_{50} \to \varphi_{15} \to \varphi_{51}$\lbl{A1}}
\BIC{$\varphi_{15} \to \varphi_{51}$}
\BIC{$\varphi_{51}$}
\end{ptree}
\noindent$\Pi_{37}$:
\begin{ptree}
\AXC{$\stackrel{\Pi_{32}}{\vdots}$}\noLine
\UIC{$b \cdot (x - a) = b \cdot x - b \cdot a \to \varphi_{22} \to \varphi_{35}$}
\AXC{$x = x$\lbl{A5}}
\UIC{$\forall x\,(x = x)$}
\AXC{$\forall x\,(x = x) \to \varphi_{22}$\lbl{A2}}
\BIC{$\varphi_{22}$}
\AXC{$\stackrel{\Pi_{36}}{\vdots}$}\noLine
\UIC{$\varphi_{51}$}
\AXC{$\varphi_{51} \to \varphi_{22} \to \varphi_{106}$\lbl{A1}}
\BIC{$\varphi_{22} \to \varphi_{106}$}
\BIC{$\varphi_{106}$}
\BIC{$x \in \mathbb{R} \to y \in \mathbb{R} \to a \in \mathbb{R} \to b \in \mathbb{R} \to \varphi_{35}$}
\end{ptree}
\noindent$\Pi_{38}$:
\begin{ptree}
\AXC{$\stackrel{\Pi_{1}}{\vdots}$}\noLine
\UIC{$\forall x\,\forall y\,(x = y \to y = x)$}
\AXC{$\forall x\,\forall y\,(x = y \to y = x) \to \forall y\,(v_{203} = y \to y = v_{203})$\lbl{A2}}
\BIC{$\forall y\,(v_{203} = y \to y = v_{203})$}
\AXC{$\varphi_{3}$\lbl{A2}}
\BIC{$v_{203} = v_{204} \to v_{204} = v_{203}$}
\UIC{$\forall v_{204}\,(v_{203} = v_{204} \to v_{204} = v_{203})$}
\UIC{$\forall v_{203}\,\forall v_{204}\,(v_{203} = v_{204} \to v_{204} = v_{203})$}
\AXC{$\varphi_{60}$\lbl{A2}}
\BIC{$\forall v_{204}\,(x \cdot y - x \cdot b = v_{204} \to v_{204} = x \cdot y - x \cdot b)$}
\AXC{$\varphi_{80}$\lbl{A2}}
\BIC{$\varphi_{16}$}
\end{ptree}
\noindent$\Pi_{39}$:
\begin{ptree}
\AXC{$\varphi_{46}$\lbl{[F17]}}
\AXC{$\varphi_{46} \to \varphi_{58}$\lbl{A2}}
\BIC{$\varphi_{58}$}
\AXC{$\varphi_{58} \to \varphi_{71}$\lbl{A2}}
\BIC{$\varphi_{71}$}
\AXC{$\varphi_{71} \to v_{195} \in \mathbb{R} \to v_{196} \in \mathbb{R} \to \varphi_{26}$\lbl{A2}}
\BIC{$v_{195} \in \mathbb{R} \to v_{196} \in \mathbb{R} \to \varphi_{26}$}
\UIC{$\forall v_{197}\,(v_{195} \in \mathbb{R} \to v_{196} \in \mathbb{R} \to \varphi_{26})$}
\UIC{$\forall v_{196}\,\forall v_{197}\,(v_{195} \in \mathbb{R} \to v_{196} \in \mathbb{R} \to \varphi_{26})$}
\UIC{$\forall v_{195}\,\forall v_{196}\,\forall v_{197}\,(v_{195} \in \mathbb{R} \to v_{196} \in \mathbb{R} \to \varphi_{26})$}
\AXC{$\varphi_{98}$\lbl{A2}}
\BIC{$\varphi_{74}$}
\AXC{$\varphi_{74} \to \varphi_{54}$\lbl{A2}}
\BIC{$\varphi_{54}$}
\end{ptree}
\noindent$\Pi_{40}$:
\begin{ptree}
\AXC{$x \in \mathbb{R} \to y \in \mathbb{R} \to a \in \mathbb{R} \to b \in \mathbb{R} \to x \in \mathbb{R}$\lbl{A1}}
\AXC{$\stackrel{\Pi_{39}}{\vdots}$}\noLine
\UIC{$\varphi_{54}$}
\AXC{$\varphi_{93}$\lbl{A2}}
\BIC{$x \in \mathbb{R} \to y \in \mathbb{R} \to b \in \mathbb{R} \to x \cdot y - x \cdot b = x \cdot (y - b)$}
\AXC{$\varphi_{111}$\lbl{A1}}
\BIC{$(x \in \mathbb{R} \to y \in \mathbb{R} \to a \in \mathbb{R} \to b \in \mathbb{R} \to x \in \mathbb{R}) \to \varphi_{97}$}
\BIC{$\varphi_{97}$}
\end{ptree}
\noindent$\Pi_{41}$:
\begin{ptree}
\AXC{$x \in \mathbb{R} \to y \in \mathbb{R} \to a \in \mathbb{R} \to b \in \mathbb{R} \to b \in \mathbb{R}$\lbl{A1}}
\AXC{$x \in \mathbb{R} \to y \in \mathbb{R} \to a \in \mathbb{R} \to b \in \mathbb{R} \to y \in \mathbb{R}$\lbl{A1}}
\AXC{$\stackrel{\Pi_{40}}{\vdots}$}\noLine
\UIC{$\varphi_{97}$}
\BIC{$(x \in \mathbb{R} \to y \in \mathbb{R} \to a \in \mathbb{R} \to b \in \mathbb{R} \to b \in \mathbb{R}) \to \varphi_{52}$}
\BIC{$\varphi_{52}$}
\end{ptree}
\noindent$\Pi_{42}$:
\begin{ptree}
\AXC{$x = x$\lbl{A5}}
\UIC{$\forall x\,(x = x)$}
\AXC{$\forall x\,(x = x) \to \varphi_{8}$\lbl{A2}}
\BIC{$\varphi_{8}$}
\AXC{$\stackrel{\Pi_{38}}{\vdots}$}\noLine
\UIC{$\varphi_{16}$}
\AXC{$\stackrel{\Pi_{41}}{\vdots}$}\noLine
\UIC{$\varphi_{52}$}
\AXC{$\varphi_{52} \to \varphi_{16} \to \varphi_{53}$\lbl{A1}}
\BIC{$\varphi_{16} \to \varphi_{53}$}
\BIC{$\varphi_{53}$}
\AXC{$\varphi_{53} \to \varphi_{8} \to \varphi_{105}$\lbl{A1}}
\BIC{$\varphi_{8} \to \varphi_{105}$}
\BIC{$\varphi_{105}$}
\end{ptree}
\noindent$\Pi_{43}$:
\begin{ptree}
\AXC{$z = w \to \varphi_{69}$\lbl{A6}}
\UIC{$\forall w\,(z = w \to \varphi_{69})$}
\UIC{$\forall z\,\forall w\,(z = w \to \varphi_{69})$}
\AXC{$\forall z\,\forall w\,(z = w \to \varphi_{69}) \to \varphi_{81}$\lbl{A2}}
\BIC{$\varphi_{81}$}
\AXC{$\varphi_{81} \to x \cdot (y - b) = x \cdot y - x \cdot b \to \varphi_{8} \to \varphi_{14}$\lbl{A2}}
\BIC{$x \cdot (y - b) = x \cdot y - x \cdot b \to \varphi_{8} \to \varphi_{14}$}
\AXC{$\stackrel{\Pi_{42}}{\vdots}$}\noLine
\UIC{$\varphi_{105}$}
\BIC{$x \in \mathbb{R} \to y \in \mathbb{R} \to a \in \mathbb{R} \to b \in \mathbb{R} \to \varphi_{14}$}
\end{ptree}
\noindent$\Pi_{44}$:
\begin{ptree}
\AXC{$\stackrel{\Pi_{43}}{\vdots}$}\noLine
\UIC{$x \in \mathbb{R} \to y \in \mathbb{R} \to a \in \mathbb{R} \to b \in \mathbb{R} \to \varphi_{14}$}
\AXC{$(x \in \mathbb{R} \to y \in \mathbb{R} \to a \in \mathbb{R} \to b \in \mathbb{R} \to \varphi_{14}) \to \varphi_{113}$\lbl{A1}}
\BIC{$\varphi_{113}$}
\end{ptree}
\noindent$\Pi_{45}$:
\begin{ptree}
\AXC{$\stackrel{\Pi_{31}}{\vdots}$}\noLine
\UIC{$\varphi_{14} \to \varphi_{35} \to \varphi_{25}$}
\AXC{$\stackrel{\Pi_{37}}{\vdots}$}\noLine
\UIC{$x \in \mathbb{R} \to y \in \mathbb{R} \to a \in \mathbb{R} \to b \in \mathbb{R} \to \varphi_{35}$}
\AXC{$\stackrel{\Pi_{44}}{\vdots}$}\noLine
\UIC{$\varphi_{113}$}
\BIC{$(\varphi_{14} \to \varphi_{35} \to \varphi_{25}) \to x \in \mathbb{R} \to y \in \mathbb{R} \to a \in \mathbb{R} \to b \in \mathbb{R} \to \varphi_{25}$}
\BIC{$x \in \mathbb{R} \to y \in \mathbb{R} \to a \in \mathbb{R} \to b \in \mathbb{R} \to \varphi_{25}$}
\AXC{$(x \in \mathbb{R} \to y \in \mathbb{R} \to a \in \mathbb{R} \to b \in \mathbb{R} \to \varphi_{25}) \to \varphi_{116}$\lbl{A1}}
\BIC{$\varphi_{116}$}
\end{ptree}
\noindent$\Pi_{46}$:
\begin{ptree}
\AXC{$\stackrel{\Pi_{26}}{\vdots}$}\noLine
\UIC{$\forall w\,(b \cdot x = w \to \varphi_{43} \to \varphi_{33})$}
\AXC{$\forall w\,(b \cdot x = w \to \varphi_{43} \to \varphi_{33}) \to b \cdot x = x \cdot b \to \varphi_{43} \to \varphi_{45}$\lbl{A2}}
\BIC{$b \cdot x = x \cdot b \to \varphi_{43} \to \varphi_{45}$}
\AXC{$\stackrel{\Pi_{29}}{\vdots}$}\noLine
\UIC{$\varphi_{108}$}
\BIC{$x \in \mathbb{R} \to y \in \mathbb{R} \to a \in \mathbb{R} \to b \in \mathbb{R} \to \varphi_{45}$}
\AXC{$\stackrel{\Pi_{45}}{\vdots}$}\noLine
\UIC{$\varphi_{116}$}
\BIC{$(\varphi_{25} \to \varphi_{45} \to \varphi_{24}) \to x \in \mathbb{R} \to y \in \mathbb{R} \to a \in \mathbb{R} \to b \in \mathbb{R} \to \varphi_{24}$}
\end{ptree}
\noindent$\Pi_{47}$:
\begin{ptree}
\AXC{$\stackrel{\Pi_{25}}{\vdots}$}\noLine
\UIC{$\varphi_{25} \to \varphi_{45} \to \varphi_{24}$}
\AXC{$\stackrel{\Pi_{46}}{\vdots}$}\noLine
\UIC{$(\varphi_{25} \to \varphi_{45} \to \varphi_{24}) \to x \in \mathbb{R} \to y \in \mathbb{R} \to a \in \mathbb{R} \to b \in \mathbb{R} \to \varphi_{24}$}
\BIC{$x \in \mathbb{R} \to y \in \mathbb{R} \to a \in \mathbb{R} \to b \in \mathbb{R} \to \varphi_{24}$}
\AXC{$(x \in \mathbb{R} \to y \in \mathbb{R} \to a \in \mathbb{R} \to b \in \mathbb{R} \to \varphi_{24}) \to \varphi_{117}$\lbl{A1}}
\BIC{$\varphi_{117}$}
\end{ptree}
\noindent$\Pi_{48}$:
\begin{ptree}
\AXC{$\stackrel{\Pi_{20}}{\vdots}$}\noLine
\UIC{$\forall w\,(b \cdot a = w \to \varphi_{42} \to \varphi_{31})$}
\AXC{$\forall w\,(b \cdot a = w \to \varphi_{42} \to \varphi_{31}) \to b \cdot a = a \cdot b \to \varphi_{42} \to \varphi_{44}$\lbl{A2}}
\BIC{$b \cdot a = a \cdot b \to \varphi_{42} \to \varphi_{44}$}
\AXC{$\stackrel{\Pi_{23}}{\vdots}$}\noLine
\UIC{$\varphi_{107}$}
\BIC{$x \in \mathbb{R} \to y \in \mathbb{R} \to a \in \mathbb{R} \to b \in \mathbb{R} \to \varphi_{44}$}
\AXC{$\stackrel{\Pi_{47}}{\vdots}$}\noLine
\UIC{$\varphi_{117}$}
\BIC{$(\varphi_{24} \to \varphi_{44} \to \varphi_{23}) \to x \in \mathbb{R} \to y \in \mathbb{R} \to a \in \mathbb{R} \to b \in \mathbb{R} \to \varphi_{23}$}
\end{ptree}
\noindent$\Pi_{49}$:
\begin{ptree}
\AXC{$\stackrel{\Pi_{19}}{\vdots}$}\noLine
\UIC{$\varphi_{24} \to \varphi_{44} \to \varphi_{23}$}
\AXC{$\stackrel{\Pi_{48}}{\vdots}$}\noLine
\UIC{$(\varphi_{24} \to \varphi_{44} \to \varphi_{23}) \to x \in \mathbb{R} \to y \in \mathbb{R} \to a \in \mathbb{R} \to b \in \mathbb{R} \to \varphi_{23}$}
\BIC{$x \in \mathbb{R} \to y \in \mathbb{R} \to a \in \mathbb{R} \to b \in \mathbb{R} \to \varphi_{23}$}
\AXC{$(x \in \mathbb{R} \to y \in \mathbb{R} \to a \in \mathbb{R} \to b \in \mathbb{R} \to \varphi_{23}) \to \varphi_{112}$\lbl{A1}}
\BIC{$\varphi_{112}$}
\end{ptree}
\noindent$\Pi_{50}$:
\begin{ptree}
\AXC{$\stackrel{\Pi_{5}}{\vdots}$}\noLine
\UIC{$\varphi_{23} \to \varphi_{75}$}
\AXC{$\stackrel{\Pi_{6}}{\vdots}$}\noLine
\UIC{$\varphi_{78}$}
\AXC{$\stackrel{\Pi_{17}}{\vdots}$}\noLine
\UIC{$\varphi_{78} \to x \in \mathbb{R} \to y \in \mathbb{R} \to \varphi_{49}$}
\BIC{$x \in \mathbb{R} \to y \in \mathbb{R} \to \varphi_{49}$}
\AXC{$\stackrel{\Pi_{49}}{\vdots}$}\noLine
\UIC{$\varphi_{112}$}
\BIC{$(\varphi_{23} \to \varphi_{75}) \to x \in \mathbb{R} \to y \in \mathbb{R} \to \varphi_{36}$}
\BIC{$x \in \mathbb{R} \to y \in \mathbb{R} \to \varphi_{36}$}
\AXC{$(x \in \mathbb{R} \to y \in \mathbb{R} \to \varphi_{36}) \to \varphi_{64} \to x \in \mathbb{R} \to y \in \mathbb{R} \to \varphi_{37}$\lbl{A1}}
\BIC{$\varphi_{64} \to x \in \mathbb{R} \to y \in \mathbb{R} \to \varphi_{37}$}
\end{ptree}
\noindent 主木:
\begin{ptree}
\AXC{$\stackrel{\Pi_{3}}{\vdots}$}\noLine
\UIC{$\varphi_{64}$}
\AXC{$\stackrel{\Pi_{50}}{\vdots}$}\noLine
\UIC{$\varphi_{64} \to x \in \mathbb{R} \to y \in \mathbb{R} \to \varphi_{37}$}
\BIC{$x \in \mathbb{R} \to y \in \mathbb{R} \to \varphi_{37}$}
\UIC{$\forall b\,(x \in \mathbb{R} \to y \in \mathbb{R} \to \varphi_{37})$}
\UIC{$\forall a\,\forall b\,(x \in \mathbb{R} \to y \in \mathbb{R} \to \varphi_{37})$}
\UIC{$\forall y\,\forall a\,\forall b\,(x \in \mathbb{R} \to y \in \mathbb{R} \to \varphi_{37})$}
\UIC{$\forall x\,\forall y\,\forall a\,\forall b\,(x \in \mathbb{R} \to y \in \mathbb{R} \to \varphi_{37})$}
\end{ptree}
\begin{abbrevlist}
\item $\varphi_{1} :\equiv \forall y\,(v_{217} = y \to y = v_{217}) \to v_{217} = v_{218} \to v_{218} = v_{217}$
\item $\varphi_{2} :\equiv \forall y\,(v_{201} = y \to y = v_{201}) \to v_{201} = v_{202} \to v_{202} = v_{201}$
\item $\varphi_{3} :\equiv \forall y\,(v_{203} = y \to y = v_{203}) \to v_{203} = v_{204} \to v_{204} = v_{203}$
\item $\varphi_{4} :\equiv \forall v_{215}\,\forall v_{216}\,(v_{214} = v_{215} \to v_{215} = v_{216} \to v_{214} = v_{216})$
\item $\varphi_{5} :\equiv \forall v_{212}\,\forall v_{213}\,(v_{211} = v_{212} \to v_{212} = v_{213} \to v_{211} = v_{213})$
\item $\varphi_{6} :\equiv \forall v_{209}\,\forall v_{210}\,(v_{208} = v_{209} \to v_{209} = v_{210} \to v_{208} = v_{210})$
\item $\varphi_{7} :\equiv \forall v_{206}\,\forall v_{207}\,(v_{205} = v_{206} \to v_{206} = v_{207} \to v_{205} = v_{207})$
\item $\varphi_{8} :\equiv x \cdot (y - b) + b \cdot (x - a) = x \cdot (y - b) + b \cdot (x - a)$
\item $\varphi_{9} :\equiv \forall x\,\forall y\,\forall z\,(x = y \to y = z \to x = z) \to \forall y\,\forall z\,(v_{214} = y \to y = z \to v_{214} = z)$
\item $\varphi_{10} :\equiv \forall x\,\forall y\,\forall z\,(x = y \to y = z \to x = z) \to \forall y\,\forall z\,(v_{211} = y \to y = z \to v_{211} = z)$
\item $\varphi_{11} :\equiv \forall x\,\forall y\,\forall z\,(x = y \to y = z \to x = z) \to \forall y\,\forall z\,(v_{208} = y \to y = z \to v_{208} = z)$
\item $\varphi_{12} :\equiv \forall x\,\forall y\,\forall z\,(x = y \to y = z \to x = z) \to \forall y\,\forall z\,(v_{205} = y \to y = z \to v_{205} = z)$
\item $\varphi_{13} :\equiv v_{200} \in \mathbb{R} \to v_{198} - v_{199} = v_{198} - v_{200} + (v_{200} - v_{199})$
\item $\varphi_{14} :\equiv x \cdot (y - b) + b \cdot (x - a) = x \cdot y - x \cdot b + b \cdot (x - a)$
\item $\varphi_{15} :\equiv b \cdot x - b \cdot a = b \cdot (x - a) \to b \cdot (x - a) = b \cdot x - b \cdot a$
\item $\varphi_{16} :\equiv x \cdot y - x \cdot b = x \cdot (y - b) \to x \cdot (y - b) = x \cdot y - x \cdot b$
\item $\varphi_{17} :\equiv \forall y\,\forall z\,(v_{214} = y \to y = z \to v_{214} = z) \to \forall z\,(v_{214} = v_{215} \to v_{215} = z \to v_{214} = z)$
\item $\varphi_{18} :\equiv \forall y\,\forall z\,(v_{211} = y \to y = z \to v_{211} = z) \to \forall z\,(v_{211} = v_{212} \to v_{212} = z \to v_{211} = z)$
\item $\varphi_{19} :\equiv \forall y\,\forall z\,(v_{208} = y \to y = z \to v_{208} = z) \to \forall z\,(v_{208} = v_{209} \to v_{209} = z \to v_{208} = z)$
\item $\varphi_{20} :\equiv \forall y\,\forall z\,(v_{205} = y \to y = z \to v_{205} = z) \to \forall z\,(v_{205} = v_{206} \to v_{206} = z \to v_{205} = z)$
\item $\varphi_{21} :\equiv \forall x\,\forall y\,(x = y \to y = x) \to \forall y\,(b \cdot x - b \cdot a = y \to y = b \cdot x - b \cdot a)$
\item $\varphi_{22} :\equiv x \cdot y - x \cdot b + b \cdot (x - a) = x \cdot y - x \cdot b + b \cdot (x - a)$
\item $\varphi_{23} :\equiv x \cdot (y - b) + b \cdot (x - a) = x \cdot y - x \cdot b + (x \cdot b - a \cdot b)$
\item $\varphi_{24} :\equiv x \cdot (y - b) + b \cdot (x - a) = x \cdot y - x \cdot b + (x \cdot b - b \cdot a)$
\item $\varphi_{25} :\equiv x \cdot (y - b) + b \cdot (x - a) = x \cdot y - x \cdot b + (b \cdot x - b \cdot a)$
\item $\varphi_{26} :\equiv v_{197} \in \mathbb{R} \to v_{195} \cdot v_{196} - v_{195} \cdot v_{197} = v_{195} \cdot (v_{196} - v_{197})$
\item $\varphi_{27} :\equiv \forall z\,(v_{214} = v_{215} \to v_{215} = z \to v_{214} = z) \to v_{214} = v_{215} \to v_{215} = v_{216} \to v_{214} = v_{216}$
\item $\varphi_{28} :\equiv \forall z\,(v_{211} = v_{212} \to v_{212} = z \to v_{211} = z) \to v_{211} = v_{212} \to v_{212} = v_{213} \to v_{211} = v_{213}$
\item $\varphi_{29} :\equiv \forall z\,(v_{208} = v_{209} \to v_{209} = z \to v_{208} = z) \to v_{208} = v_{209} \to v_{209} = v_{210} \to v_{208} = v_{210}$
\item $\varphi_{30} :\equiv \forall z\,(v_{205} = v_{206} \to v_{206} = z \to v_{205} = z) \to v_{205} = v_{206} \to v_{206} = v_{207} \to v_{205} = v_{207}$
\item $\varphi_{31} :\equiv x \cdot y - x \cdot b + (x \cdot b - b \cdot a) = x \cdot y - x \cdot b + (x \cdot b - w)$
\item $\varphi_{32} :\equiv x \cdot y - x \cdot b + (x \cdot b - b \cdot a) = x \cdot y - x \cdot b + (x \cdot b - z)$
\item $\varphi_{33} :\equiv x \cdot y - x \cdot b + (b \cdot x - b \cdot a) = x \cdot y - x \cdot b + (w - b \cdot a)$
\item $\varphi_{34} :\equiv x \cdot y - x \cdot b + (b \cdot x - b \cdot a) = x \cdot y - x \cdot b + (z - b \cdot a)$
\item $\varphi_{35} :\equiv x \cdot y - x \cdot b + b \cdot (x - a) = x \cdot y - x \cdot b + (b \cdot x - b \cdot a)$
\item $\varphi_{36} :\equiv a \in \mathbb{R} \to b \in \mathbb{R} \to x \cdot (y - b) + b \cdot (x - a) = x \cdot y - a \cdot b$
\item $\varphi_{37} :\equiv a \in \mathbb{R} \to b \in \mathbb{R} \to x \cdot y - a \cdot b = x \cdot (y - b) + b \cdot (x - a)$
\item $\varphi_{38} :\equiv \forall a\,(b \in \mathbb{R} \to x \in \mathbb{R} \to a \in \mathbb{R} \to b \cdot x - b \cdot a = b \cdot (x - a))$
\item $\varphi_{39} :\equiv \forall v_{218}\,(x \cdot (y - b) + b \cdot (x - a) = v_{218} \to v_{218} = x \cdot (y - b) + b \cdot (x - a))$
\item $\varphi_{40} :\equiv \forall a\,\forall b\,\forall x\,(a \in \mathbb{R} \to b \in \mathbb{R} \to x \in \mathbb{R} \to a - b = a - x + (x - b))$
\item $\varphi_{41} :\equiv x \cdot b \in \mathbb{R} \to x \cdot y - a \cdot b = x \cdot y - x \cdot b + (x \cdot b - a \cdot b)$
\item $\varphi_{42} :\equiv x \cdot y - x \cdot b + (x \cdot b - b \cdot a) = x \cdot y - x \cdot b + (x \cdot b - b \cdot a)$
\item $\varphi_{43} :\equiv x \cdot y - x \cdot b + (b \cdot x - b \cdot a) = x \cdot y - x \cdot b + (b \cdot x - b \cdot a)$
\item $\varphi_{44} :\equiv x \cdot y - x \cdot b + (x \cdot b - b \cdot a) = x \cdot y - x \cdot b + (x \cdot b - a \cdot b)$
\item $\varphi_{45} :\equiv x \cdot y - x \cdot b + (b \cdot x - b \cdot a) = x \cdot y - x \cdot b + (x \cdot b - b \cdot a)$
\item $\varphi_{46} :\equiv \forall c\,\forall x\,\forall a\,(c \in \mathbb{R} \to x \in \mathbb{R} \to a \in \mathbb{R} \to c \cdot x - c \cdot a = c \cdot (x - a))$
\item $\varphi_{47} :\equiv \forall b\,\forall x\,(v_{198} \in \mathbb{R} \to b \in \mathbb{R} \to x \in \mathbb{R} \to v_{198} - b = v_{198} - x + (x - b))$
\item $\varphi_{48} :\equiv a \in \mathbb{R} \to b \in \mathbb{R} \to x \cdot y - a \cdot b = x \cdot y - x \cdot b + (x \cdot b - a \cdot b)$
\item $\varphi_{49} :\equiv a \in \mathbb{R} \to b \in \mathbb{R} \to x \cdot y - x \cdot b + (x \cdot b - a \cdot b) = x \cdot y - a \cdot b$
\item $\varphi_{50} :\equiv x \in \mathbb{R} \to y \in \mathbb{R} \to a \in \mathbb{R} \to b \in \mathbb{R} \to b \cdot x - b \cdot a = b \cdot (x - a)$
\item $\varphi_{51} :\equiv x \in \mathbb{R} \to y \in \mathbb{R} \to a \in \mathbb{R} \to b \in \mathbb{R} \to b \cdot (x - a) = b \cdot x - b \cdot a$
\item $\varphi_{52} :\equiv x \in \mathbb{R} \to y \in \mathbb{R} \to a \in \mathbb{R} \to b \in \mathbb{R} \to x \cdot y - x \cdot b = x \cdot (y - b)$
\item $\varphi_{53} :\equiv x \in \mathbb{R} \to y \in \mathbb{R} \to a \in \mathbb{R} \to b \in \mathbb{R} \to x \cdot (y - b) = x \cdot y - x \cdot b$
\item $\varphi_{54} :\equiv \forall v_{197}\,(x \in \mathbb{R} \to y \in \mathbb{R} \to v_{197} \in \mathbb{R} \to x \cdot y - x \cdot v_{197} = x \cdot (y - v_{197}))$
\item $\varphi_{55} :\equiv \forall x\,(v_{198} \in \mathbb{R} \to v_{199} \in \mathbb{R} \to x \in \mathbb{R} \to v_{198} - v_{199} = v_{198} - x + (x - v_{199}))$
\item $\varphi_{56} :\equiv \forall y\,(b \in \mathbb{R} \to y \in \mathbb{R} \to b \cdot y = y \cdot b) \to b \in \mathbb{R} \to a \in \mathbb{R} \to b \cdot a = a \cdot b$
\item $\varphi_{57} :\equiv \forall y\,(b \in \mathbb{R} \to y \in \mathbb{R} \to b \cdot y = y \cdot b) \to b \in \mathbb{R} \to x \in \mathbb{R} \to b \cdot x = x \cdot b$
\item $\varphi_{58} :\equiv \forall x\,\forall a\,(v_{195} \in \mathbb{R} \to x \in \mathbb{R} \to a \in \mathbb{R} \to v_{195} \cdot x - v_{195} \cdot a = v_{195} \cdot (x - a))$
\item $\varphi_{59} :\equiv \forall v_{201}\,\forall v_{202}\,(v_{201} = v_{202} \to v_{202} = v_{201}) \to \forall v_{202}\,(x \cdot y - a \cdot b = v_{202} \to v_{202} = x \cdot y - a \cdot b)$
\item $\varphi_{60} :\equiv \forall v_{203}\,\forall v_{204}\,(v_{203} = v_{204} \to v_{204} = v_{203}) \to \forall v_{204}\,(x \cdot y - x \cdot b = v_{204} \to v_{204} = x \cdot y - x \cdot b)$
\item $\varphi_{61} :\equiv \forall y\,(x \in \mathbb{R} \to y \in \mathbb{R} \to x \cdot y \in \mathbb{R}) \to x \in \mathbb{R} \to b \in \mathbb{R} \to x \cdot b \in \mathbb{R}$
\item $\varphi_{62} :\equiv \forall y\,(a \in \mathbb{R} \to y \in \mathbb{R} \to a \cdot y \in \mathbb{R}) \to a \in \mathbb{R} \to b \in \mathbb{R} \to a \cdot b \in \mathbb{R}$
\item $\varphi_{63} :\equiv \forall y\,(x \in \mathbb{R} \to y \in \mathbb{R} \to x \cdot y \in \mathbb{R}) \to x \in \mathbb{R} \to y \in \mathbb{R} \to x \cdot y \in \mathbb{R}$
\item $\varphi_{64} :\equiv x \cdot (y - b) + b \cdot (x - a) = x \cdot y - a \cdot b \to x \cdot y - a \cdot b = x \cdot (y - b) + b \cdot (x - a)$
\item $\varphi_{65} :\equiv \forall v_{215}\,\forall v_{216}\,(x \cdot (y - b) + b \cdot (x - a) = v_{215} \to v_{215} = v_{216} \to x \cdot (y - b) + b \cdot (x - a) = v_{216})$
\item $\varphi_{66} :\equiv \forall v_{212}\,\forall v_{213}\,(x \cdot (y - b) + b \cdot (x - a) = v_{212} \to v_{212} = v_{213} \to x \cdot (y - b) + b \cdot (x - a) = v_{213})$
\item $\varphi_{67} :\equiv \forall v_{209}\,\forall v_{210}\,(x \cdot (y - b) + b \cdot (x - a) = v_{209} \to v_{209} = v_{210} \to x \cdot (y - b) + b \cdot (x - a) = v_{210})$
\item $\varphi_{68} :\equiv \forall v_{206}\,\forall v_{207}\,(x \cdot (y - b) + b \cdot (x - a) = v_{206} \to v_{206} = v_{207} \to x \cdot (y - b) + b \cdot (x - a) = v_{207})$
\item $\varphi_{69} :\equiv x \cdot (y - b) + b \cdot (x - a) = z + b \cdot (x - a) \to x \cdot (y - b) + b \cdot (x - a) = w + b \cdot (x - a)$
\item $\varphi_{70} :\equiv \forall x\,\forall y\,(x \in \mathbb{R} \to y \in \mathbb{R} \to x \cdot y = y \cdot x) \to \forall y\,(b \in \mathbb{R} \to y \in \mathbb{R} \to b \cdot y = y \cdot b)$
\item $\varphi_{71} :\equiv \forall a\,(v_{195} \in \mathbb{R} \to v_{196} \in \mathbb{R} \to a \in \mathbb{R} \to v_{195} \cdot v_{196} - v_{195} \cdot a = v_{195} \cdot (v_{196} - a))$
\item $\varphi_{72} :\equiv \forall x\,\forall y\,(x \in \mathbb{R} \to y \in \mathbb{R} \to x \cdot y \in \mathbb{R}) \to \forall y\,(x \in \mathbb{R} \to y \in \mathbb{R} \to x \cdot y \in \mathbb{R})$
\item $\varphi_{73} :\equiv \forall x\,\forall y\,(x \in \mathbb{R} \to y \in \mathbb{R} \to x \cdot y \in \mathbb{R}) \to \forall y\,(a \in \mathbb{R} \to y \in \mathbb{R} \to a \cdot y \in \mathbb{R})$
\item $\varphi_{74} :\equiv \forall v_{196}\,\forall v_{197}\,(x \in \mathbb{R} \to v_{196} \in \mathbb{R} \to v_{197} \in \mathbb{R} \to x \cdot v_{196} - x \cdot v_{197} = x \cdot (v_{196} - v_{197}))$
\item $\varphi_{75} :\equiv x \cdot y - x \cdot b + (x \cdot b - a \cdot b) = x \cdot y - a \cdot b \to x \cdot (y - b) + b \cdot (x - a) = x \cdot y - a \cdot b$
\item $\varphi_{76} :\equiv \forall v_{199}\,\forall v_{200}\,(x \cdot y \in \mathbb{R} \to v_{199} \in \mathbb{R} \to v_{200} \in \mathbb{R} \to x \cdot y - v_{199} = x \cdot y - v_{200} + (v_{200} - v_{199}))$
\item $\varphi_{77} :\equiv \forall v_{200}\,(x \cdot y \in \mathbb{R} \to a \cdot b \in \mathbb{R} \to v_{200} \in \mathbb{R} \to x \cdot y - a \cdot b = x \cdot y - v_{200} + (v_{200} - a \cdot b))$
\item $\varphi_{78} :\equiv x \cdot y - a \cdot b = x \cdot y - x \cdot b + (x \cdot b - a \cdot b) \to x \cdot y - x \cdot b + (x \cdot b - a \cdot b) = x \cdot y - a \cdot b$
\item $\varphi_{79} :\equiv x \cdot y - x \cdot b + b \cdot (x - a) = x \cdot y - x \cdot b + z \to x \cdot y - x \cdot b + b \cdot (x - a) = x \cdot y - x \cdot b + w$
\item $\varphi_{80} :\equiv \forall v_{204}\,(x \cdot y - x \cdot b = v_{204} \to v_{204} = x \cdot y - x \cdot b) \to \varphi_{16}$
\item $\varphi_{81} :\equiv \forall w\,(x \cdot (y - b) = w \to \varphi_{8} \to x \cdot (y - b) + b \cdot (x - a) = w + b \cdot (x - a))$
\item $\varphi_{82} :\equiv (x \in \mathbb{R} \to y \in \mathbb{R} \to a \in \mathbb{R} \to b \in \mathbb{R} \to a \in \mathbb{R}) \to x \in \mathbb{R} \to y \in \mathbb{R} \to a \in \mathbb{R} \to b \in \mathbb{R} \to b \cdot a = a \cdot b$
\item $\varphi_{83} :\equiv (x \in \mathbb{R} \to y \in \mathbb{R} \to a \in \mathbb{R} \to b \in \mathbb{R} \to x \in \mathbb{R}) \to x \in \mathbb{R} \to y \in \mathbb{R} \to a \in \mathbb{R} \to b \in \mathbb{R} \to b \cdot x = x \cdot b$
\item $\varphi_{84} :\equiv (x \in \mathbb{R} \to y \in \mathbb{R} \to a \in \mathbb{R} \to b \in \mathbb{R} \to b \in \mathbb{R}) \to x \in \mathbb{R} \to y \in \mathbb{R} \to a \in \mathbb{R} \to b \in \mathbb{R} \to x \cdot b \in \mathbb{R}$
\item $\varphi_{85} :\equiv (x \in \mathbb{R} \to y \in \mathbb{R} \to a \in \mathbb{R} \to b \in \mathbb{R} \to b \in \mathbb{R}) \to x \in \mathbb{R} \to y \in \mathbb{R} \to a \in \mathbb{R} \to b \in \mathbb{R} \to a \cdot b \in \mathbb{R}$
\item $\varphi_{86} :\equiv (x \in \mathbb{R} \to y \in \mathbb{R} \to a \in \mathbb{R} \to b \in \mathbb{R} \to y \in \mathbb{R}) \to x \in \mathbb{R} \to y \in \mathbb{R} \to a \in \mathbb{R} \to b \in \mathbb{R} \to x \cdot y \in \mathbb{R}$
\item $\varphi_{87} :\equiv \forall v_{207}\,(\varphi_{14} \to x \cdot y - x \cdot b + b \cdot (x - a) = v_{207} \to x \cdot (y - b) + b \cdot (x - a) = v_{207})$
\item $\varphi_{88} :\equiv \forall w\,(b \cdot (x - a) = w \to \varphi_{22} \to x \cdot y - x \cdot b + b \cdot (x - a) = x \cdot y - x \cdot b + w)$
\item $\varphi_{89} :\equiv \varphi_{38} \to b \in \mathbb{R} \to x \in \mathbb{R} \to a \in \mathbb{R} \to b \cdot x - b \cdot a = b \cdot (x - a)$
\item $\varphi_{90} :\equiv \forall v_{216}\,(\varphi_{23} \to x \cdot y - x \cdot b + (x \cdot b - a \cdot b) = v_{216} \to x \cdot (y - b) + b \cdot (x - a) = v_{216})$
\item $\varphi_{91} :\equiv \forall v_{213}\,(\varphi_{24} \to x \cdot y - x \cdot b + (x \cdot b - b \cdot a) = v_{213} \to x \cdot (y - b) + b \cdot (x - a) = v_{213})$
\item $\varphi_{92} :\equiv \forall v_{210}\,(\varphi_{25} \to x \cdot y - x \cdot b + (b \cdot x - b \cdot a) = v_{210} \to x \cdot (y - b) + b \cdot (x - a) = v_{210})$
\item $\varphi_{93} :\equiv \varphi_{54} \to x \in \mathbb{R} \to y \in \mathbb{R} \to b \in \mathbb{R} \to x \cdot y - x \cdot b = x \cdot (y - b)$
\item $\varphi_{94} :\equiv \forall v_{202}\,(x \cdot y - a \cdot b = v_{202} \to v_{202} = x \cdot y - a \cdot b) \to \varphi_{78}$
\item $\varphi_{95} :\equiv (x \in \mathbb{R} \to y \in \mathbb{R} \to a \in \mathbb{R} \to b \in \mathbb{R} \to x \cdot b \in \mathbb{R}) \to x \in \mathbb{R} \to y \in \mathbb{R} \to \varphi_{48}$
\item $\varphi_{96} :\equiv (x \in \mathbb{R} \to y \in \mathbb{R} \to a \in \mathbb{R} \to b \in \mathbb{R} \to x \in \mathbb{R}) \to (x \in \mathbb{R} \to y \in \mathbb{R} \to a \in \mathbb{R} \to b \in \mathbb{R} \to a \in \mathbb{R}) \to \varphi_{50}$
\item $\varphi_{97} :\equiv (x \in \mathbb{R} \to y \in \mathbb{R} \to a \in \mathbb{R} \to b \in \mathbb{R} \to y \in \mathbb{R}) \to (x \in \mathbb{R} \to y \in \mathbb{R} \to a \in \mathbb{R} \to b \in \mathbb{R} \to b \in \mathbb{R}) \to \varphi_{52}$
\item $\varphi_{98} :\equiv \forall v_{195}\,\forall v_{196}\,\forall v_{197}\,(v_{195} \in \mathbb{R} \to v_{196} \in \mathbb{R} \to \varphi_{26}) \to \varphi_{74}$
\item $\varphi_{99} :\equiv \forall v_{198}\,\forall v_{199}\,\forall v_{200}\,(v_{198} \in \mathbb{R} \to v_{199} \in \mathbb{R} \to \varphi_{13}) \to \varphi_{76}$
\item $\varphi_{100} :\equiv (b \in \mathbb{R} \to a \in \mathbb{R} \to b \cdot a = a \cdot b) \to (x \in \mathbb{R} \to y \in \mathbb{R} \to a \in \mathbb{R} \to b \in \mathbb{R} \to b \in \mathbb{R}) \to \varphi_{82}$
\item $\varphi_{101} :\equiv (b \in \mathbb{R} \to x \in \mathbb{R} \to b \cdot x = x \cdot b) \to (x \in \mathbb{R} \to y \in \mathbb{R} \to a \in \mathbb{R} \to b \in \mathbb{R} \to b \in \mathbb{R}) \to \varphi_{83}$
\item $\varphi_{102} :\equiv (x \in \mathbb{R} \to b \in \mathbb{R} \to x \cdot b \in \mathbb{R}) \to (x \in \mathbb{R} \to y \in \mathbb{R} \to a \in \mathbb{R} \to b \in \mathbb{R} \to x \in \mathbb{R}) \to \varphi_{84}$
\item $\varphi_{103} :\equiv (a \in \mathbb{R} \to b \in \mathbb{R} \to a \cdot b \in \mathbb{R}) \to (x \in \mathbb{R} \to y \in \mathbb{R} \to a \in \mathbb{R} \to b \in \mathbb{R} \to a \in \mathbb{R}) \to \varphi_{85}$
\item $\varphi_{104} :\equiv (x \in \mathbb{R} \to y \in \mathbb{R} \to x \cdot y \in \mathbb{R}) \to (x \in \mathbb{R} \to y \in \mathbb{R} \to a \in \mathbb{R} \to b \in \mathbb{R} \to x \in \mathbb{R}) \to \varphi_{86}$
\item $\varphi_{105} :\equiv (x \cdot (y - b) = x \cdot y - x \cdot b \to \varphi_{8} \to \varphi_{14}) \to x \in \mathbb{R} \to y \in \mathbb{R} \to a \in \mathbb{R} \to b \in \mathbb{R} \to \varphi_{14}$
\item $\varphi_{106} :\equiv (b \cdot (x - a) = b \cdot x - b \cdot a \to \varphi_{22} \to \varphi_{35}) \to x \in \mathbb{R} \to y \in \mathbb{R} \to a \in \mathbb{R} \to b \in \mathbb{R} \to \varphi_{35}$
\item $\varphi_{107} :\equiv (b \cdot a = a \cdot b \to \varphi_{42} \to \varphi_{44}) \to x \in \mathbb{R} \to y \in \mathbb{R} \to a \in \mathbb{R} \to b \in \mathbb{R} \to \varphi_{44}$
\item $\varphi_{108} :\equiv (b \cdot x = x \cdot b \to \varphi_{43} \to \varphi_{45}) \to x \in \mathbb{R} \to y \in \mathbb{R} \to a \in \mathbb{R} \to b \in \mathbb{R} \to \varphi_{45}$
\item $\varphi_{109} :\equiv (x \in \mathbb{R} \to y \in \mathbb{R} \to a \in \mathbb{R} \to b \in \mathbb{R} \to x \cdot y \in \mathbb{R}) \to (x \in \mathbb{R} \to y \in \mathbb{R} \to a \in \mathbb{R} \to b \in \mathbb{R} \to a \cdot b \in \mathbb{R}) \to \varphi_{95}$
\item $\varphi_{110} :\equiv (b \in \mathbb{R} \to x \in \mathbb{R} \to a \in \mathbb{R} \to b \cdot x - b \cdot a = b \cdot (x - a)) \to (x \in \mathbb{R} \to y \in \mathbb{R} \to a \in \mathbb{R} \to b \in \mathbb{R} \to b \in \mathbb{R}) \to \varphi_{96}$
\item $\varphi_{111} :\equiv (x \in \mathbb{R} \to y \in \mathbb{R} \to b \in \mathbb{R} \to x \cdot y - x \cdot b = x \cdot (y - b)) \to (x \in \mathbb{R} \to y \in \mathbb{R} \to a \in \mathbb{R} \to b \in \mathbb{R} \to x \in \mathbb{R}) \to \varphi_{97}$
\item $\varphi_{112} :\equiv (x \in \mathbb{R} \to y \in \mathbb{R} \to \varphi_{49}) \to (\varphi_{23} \to \varphi_{75}) \to x \in \mathbb{R} \to y \in \mathbb{R} \to \varphi_{36}$
\item $\varphi_{113} :\equiv (x \in \mathbb{R} \to y \in \mathbb{R} \to a \in \mathbb{R} \to b \in \mathbb{R} \to \varphi_{35}) \to (\varphi_{14} \to \varphi_{35} \to \varphi_{25}) \to x \in \mathbb{R} \to y \in \mathbb{R} \to a \in \mathbb{R} \to b \in \mathbb{R} \to \varphi_{25}$
\item $\varphi_{114} :\equiv (x \in \mathbb{R} \to y \in \mathbb{R} \to a \in \mathbb{R} \to b \in \mathbb{R} \to b \cdot a = a \cdot b) \to \varphi_{42} \to \varphi_{107}$
\item $\varphi_{115} :\equiv (x \in \mathbb{R} \to y \in \mathbb{R} \to a \in \mathbb{R} \to b \in \mathbb{R} \to b \cdot x = x \cdot b) \to \varphi_{43} \to \varphi_{108}$
\item $\varphi_{116} :\equiv (x \in \mathbb{R} \to y \in \mathbb{R} \to a \in \mathbb{R} \to b \in \mathbb{R} \to \varphi_{45}) \to (\varphi_{25} \to \varphi_{45} \to \varphi_{24}) \to x \in \mathbb{R} \to y \in \mathbb{R} \to a \in \mathbb{R} \to b \in \mathbb{R} \to \varphi_{24}$
\item $\varphi_{117} :\equiv (x \in \mathbb{R} \to y \in \mathbb{R} \to a \in \mathbb{R} \to b \in \mathbb{R} \to \varphi_{44}) \to (\varphi_{24} \to \varphi_{44} \to \varphi_{23}) \to x \in \mathbb{R} \to y \in \mathbb{R} \to a \in \mathbb{R} \to b \in \mathbb{R} \to \varphi_{23}$
\end{abbrevlist}
\end{fullproof}

% 自動生成 (prover/build.js): LeAddRight  ノード数 39
\begin{fullproof}{[F19]}{$\forall x\,\forall y\,\forall z\,(x \in \mathbb{R} \to y \in \mathbb{R} \to z \in \mathbb{R} \to x \leq y \to x + z \leq y + z)$}{39}\phantomsection\label{tree:LeAddRight}
\noindent$\Pi_{1}$:
\begin{ptree}
\AXC{$x = x$\lbl{A5}}
\UIC{$\forall x\,(x = x)$}
\AXC{$\forall x\,(x = x) \to x + z = x + z$\lbl{A2}}
\BIC{$x + z = x + z$}
\AXC{$x \in \mathbb{R} \to y \in \mathbb{R} \to z \in \mathbb{R} \to x \leq y \to x = y \to x = y$\lbl{A1}}
\AXC{$\varphi_{10}$\lbl{A1}}
\BIC{$x + z = x + z \to \varphi_{5}$}
\BIC{$\varphi_{5}$}
\end{ptree}
\noindent$\Pi_{2}$:
\begin{ptree}
\AXC{$z_{1} = w \to x + z = z_{1} + z \to x + z = w + z$\lbl{A6}}
\UIC{$\forall w\,(z_{1} = w \to x + z = z_{1} + z \to x + z = w + z)$}
\UIC{$\forall z_{1}\,\forall w\,(z_{1} = w \to x + z = z_{1} + z \to x + z = w + z)$}
\AXC{$\varphi_{4}$\lbl{A2}}
\BIC{$\forall w\,(x = w \to x + z = x + z \to x + z = w + z)$}
\AXC{$\varphi_{3}$\lbl{A2}}
\BIC{$x = y \to x + z = x + z \to x + z = y + z$}
\AXC{$\stackrel{\Pi_{1}}{\vdots}$}\noLine
\UIC{$\varphi_{5}$}
\BIC{$x \in \mathbb{R} \to y \in \mathbb{R} \to z \in \mathbb{R} \to x \leq y \to x = y \to x + z = y + z$}
\end{ptree}
\noindent$\Pi_{3}$:
\begin{ptree}
\AXC{$\forall x\,(\varphi_{2})$\lbl{(R14)}}
\AXC{$\forall x\,(\varphi_{2}) \to \varphi_{2}$\lbl{A2}}
\BIC{$\varphi_{2}$}
\AXC{$\varphi_{9}$\lbl{A2}}
\BIC{$\forall z\,(x \in \mathbb{R} \to y \in \mathbb{R} \to z \in \mathbb{R} \to x < y \to x + z < y + z)$}
\AXC{$\varphi_{8}$\lbl{A2}}
\BIC{$x \in \mathbb{R} \to y \in \mathbb{R} \to z \in \mathbb{R} \to x < y \to x + z < y + z$}
\AXC{$(x \in \mathbb{R} \to y \in \mathbb{R} \to z \in \mathbb{R} \to x < y \to x + z < y + z) \to \varphi_{12}$\lbl{A1}}
\BIC{$\varphi_{12}$}
\end{ptree}
\noindent$\Pi_{4}$:
\begin{ptree}
\AXC{$x \in \mathbb{R} \to y \in \mathbb{R} \to z \in \mathbb{R} \to x \leq y \to y \in \mathbb{R}$\lbl{A1}}
\AXC{$x \in \mathbb{R} \to y \in \mathbb{R} \to z \in \mathbb{R} \to x \leq y \to x \in \mathbb{R}$\lbl{A1}}
\AXC{$\stackrel{\Pi_{3}}{\vdots}$}\noLine
\UIC{$\varphi_{12}$}
\BIC{$(x \in \mathbb{R} \to y \in \mathbb{R} \to z \in \mathbb{R} \to x \leq y \to y \in \mathbb{R}) \to \varphi_{7}$}
\BIC{$\varphi_{7}$}
\end{ptree}
\noindent$\Pi_{5}$:
\begin{ptree}
\AXC{$x \in \mathbb{R} \to y \in \mathbb{R} \to z \in \mathbb{R} \to x \leq y \to x \leq y$\lbl{A1}}
\AXC{$(x \in \mathbb{R} \to y \in \mathbb{R} \to z \in \mathbb{R} \to x \leq y \to x \leq y) \to \varphi_{11}$\lbl{A1}}
\BIC{$\varphi_{11}$}
\end{ptree}
\noindent 主木:
\begin{ptree}
\AXC{$\stackrel{\Pi_{2}}{\vdots}$}\noLine
\UIC{$x \in \mathbb{R} \to y \in \mathbb{R} \to z \in \mathbb{R} \to x \leq y \to x = y \to x + z = y + z$}
\AXC{$x \in \mathbb{R} \to y \in \mathbb{R} \to z \in \mathbb{R} \to x \leq y \to z \in \mathbb{R}$\lbl{A1}}
\AXC{$\stackrel{\Pi_{4}}{\vdots}$}\noLine
\UIC{$\varphi_{7}$}
\BIC{$x \in \mathbb{R} \to y \in \mathbb{R} \to z \in \mathbb{R} \to x \leq y \to x < y \to x + z < y + z$}
\AXC{$\stackrel{\Pi_{5}}{\vdots}$}\noLine
\UIC{$\varphi_{11}$}
\BIC{$\varphi_{6}$}
\BIC{$x \in \mathbb{R} \to y \in \mathbb{R} \to z \in \mathbb{R} \to x \leq y \to x + z \leq y + z$}
\UIC{$\forall z\,(x \in \mathbb{R} \to y \in \mathbb{R} \to z \in \mathbb{R} \to x \leq y \to x + z \leq y + z)$}
\UIC{$\varphi_{1}$}
\UIC{$\forall x\,(\varphi_{1})$}
\end{ptree}
\begin{abbrevlist}
\item $\varphi_{1} :\equiv \forall y\,\forall z\,(x \in \mathbb{R} \to y \in \mathbb{R} \to z \in \mathbb{R} \to x \leq y \to x + z \leq y + z)$
\item $\varphi_{2} :\equiv \forall y\,\forall z\,(x \in \mathbb{R} \to y \in \mathbb{R} \to z \in \mathbb{R} \to x < y \to x + z < y + z)$
\item $\varphi_{3} :\equiv \forall w\,(x = w \to x + z = x + z \to x + z = w + z) \to x = y \to x + z = x + z \to x + z = y + z$
\item $\varphi_{4} :\equiv \forall z_{1}\,\forall w\,(z_{1} = w \to x + z = z_{1} + z \to x + z = w + z) \to \forall w\,(x = w \to x + z = x + z \to x + z = w + z)$
\item $\varphi_{5} :\equiv (x = y \to x + z = x + z \to x + z = y + z) \to x \in \mathbb{R} \to y \in \mathbb{R} \to z \in \mathbb{R} \to x \leq y \to x = y \to x + z = y + z$
\item $\varphi_{6} :\equiv (x \in \mathbb{R} \to y \in \mathbb{R} \to z \in \mathbb{R} \to x \leq y \to x = y \to x + z = y + z) \to x \in \mathbb{R} \to y \in \mathbb{R} \to z \in \mathbb{R} \to x \leq y \to x + z \leq y + z$
\item $\varphi_{7} :\equiv (x \in \mathbb{R} \to y \in \mathbb{R} \to z \in \mathbb{R} \to x \leq y \to z \in \mathbb{R}) \to x \in \mathbb{R} \to y \in \mathbb{R} \to z \in \mathbb{R} \to x \leq y \to x < y \to x + z < y + z$
\item $\varphi_{8} :\equiv \forall z\,(x \in \mathbb{R} \to y \in \mathbb{R} \to z \in \mathbb{R} \to x < y \to x + z < y + z) \to x \in \mathbb{R} \to y \in \mathbb{R} \to z \in \mathbb{R} \to x < y \to x + z < y + z$
\item $\varphi_{9} :\equiv \varphi_{2} \to \forall z\,(x \in \mathbb{R} \to y \in \mathbb{R} \to z \in \mathbb{R} \to x < y \to x + z < y + z)$
\item $\varphi_{10} :\equiv (x \in \mathbb{R} \to y \in \mathbb{R} \to z \in \mathbb{R} \to x \leq y \to x = y \to x = y) \to x + z = x + z \to \varphi_{5}$
\item $\varphi_{11} :\equiv (x \in \mathbb{R} \to y \in \mathbb{R} \to z \in \mathbb{R} \to x \leq y \to x < y \to x + z < y + z) \to \varphi_{6}$
\item $\varphi_{12} :\equiv (x \in \mathbb{R} \to y \in \mathbb{R} \to z \in \mathbb{R} \to x \leq y \to x \in \mathbb{R}) \to (x \in \mathbb{R} \to y \in \mathbb{R} \to z \in \mathbb{R} \to x \leq y \to y \in \mathbb{R}) \to \varphi_{7}$
\end{abbrevlist}
\end{fullproof}

% 自動生成 (prover/build.js): AddLtLe  ノード数 109
\begin{fullproof}{[F20]}{$\forall x\,\forall y\,\forall u\,\forall v\,(x \in \mathbb{R} \to y \in \mathbb{R} \to u \in \mathbb{R} \to v \in \mathbb{R} \to x < y \to u \leq v \to x + u < y + v)$}{109}\phantomsection\label{tree:AddLtLe}
\noindent$\Pi_{1}$:
\begin{ptree}
\AXC{$\varphi_{9}$\lbl{A1}}
\AXC{$\varphi_{10}$\lbl{A1}}
\AXC{$\varphi_{11}$\lbl{A1}}
\AXC{$\varphi_{2}$\lbl{[F19]}}
\AXC{$\varphi_{2} \to \varphi_{1}$\lbl{A2}}
\BIC{$\varphi_{1}$}
\AXC{$\varphi_{33}$\lbl{A2}}
\BIC{$\forall z\,(u \in \mathbb{R} \to v \in \mathbb{R} \to z \in \mathbb{R} \to u \leq v \to u + z \leq v + z)$}
\AXC{$\varphi_{32}$\lbl{A2}}
\BIC{$u \in \mathbb{R} \to v \in \mathbb{R} \to y \in \mathbb{R} \to u \leq v \to u + y \leq v + y$}
\AXC{$\varphi_{39}$\lbl{A1}}
\BIC{$\varphi_{11} \to \varphi_{10} \to \varphi_{9} \to \varphi_{3} \to \varphi_{13}$}
\BIC{$\varphi_{10} \to \varphi_{9} \to \varphi_{3} \to \varphi_{13}$}
\BIC{$\varphi_{9} \to \varphi_{3} \to \varphi_{13}$}
\BIC{$\varphi_{3} \to \varphi_{13}$}
\end{ptree}
\noindent$\Pi_{2}$:
\begin{ptree}
\AXC{$\forall x\,\forall y\,(x \in \mathbb{R} \to y \in \mathbb{R} \to x + y = y + x)$\lbl{(R3)}}
\AXC{$\varphi_{27}$\lbl{A2}}
\BIC{$\forall y\,(u \in \mathbb{R} \to y \in \mathbb{R} \to u + y = y + u)$}
\AXC{$\varphi_{22}$\lbl{A2}}
\BIC{$u \in \mathbb{R} \to y \in \mathbb{R} \to u + y = y + u$}
\AXC{$(u \in \mathbb{R} \to y \in \mathbb{R} \to u + y = y + u) \to \varphi_{11} \to \varphi_{9} \to \varphi_{14}$\lbl{A1}}
\BIC{$\varphi_{11} \to \varphi_{9} \to \varphi_{14}$}
\end{ptree}
\noindent$\Pi_{3}$:
\begin{ptree}
\AXC{$\varphi_{3}$\lbl{A1}}
\AXC{$\stackrel{\Pi_{1}}{\vdots}$}\noLine
\UIC{$\varphi_{3} \to \varphi_{13}$}
\BIC{$\varphi_{13}$}
\AXC{$\varphi_{9}$\lbl{A1}}
\AXC{$\varphi_{11}$\lbl{A1}}
\AXC{$\stackrel{\Pi_{2}}{\vdots}$}\noLine
\UIC{$\varphi_{11} \to \varphi_{9} \to \varphi_{14}$}
\BIC{$\varphi_{9} \to \varphi_{14}$}
\BIC{$\varphi_{14}$}
\AXC{$\varphi_{14} \to \varphi_{13} \to \varphi_{36}$\lbl{A1}}
\BIC{$\varphi_{13} \to \varphi_{36}$}
\BIC{$\varphi_{36}$}
\end{ptree}
\noindent$\Pi_{4}$:
\begin{ptree}
\AXC{$z = w \to z < v + y \to w < v + y$\lbl{A6}}
\UIC{$\forall w\,(z = w \to z < v + y \to w < v + y)$}
\UIC{$\forall z\,\forall w\,(z = w \to z < v + y \to w < v + y)$}
\AXC{$\varphi_{21}$\lbl{A2}}
\BIC{$\forall w\,(u + y = w \to u + y < v + y \to w < v + y)$}
\AXC{$\varphi_{30}$\lbl{A2}}
\BIC{$u + y = y + u \to u + y < v + y \to y + u < v + y$}
\AXC{$\stackrel{\Pi_{3}}{\vdots}$}\noLine
\UIC{$\varphi_{36}$}
\BIC{$(u + y = y + u \to u + y = v + y \to y + u = v + y) \to \varphi_{15}$}
\end{ptree}
\noindent$\Pi_{5}$:
\begin{ptree}
\AXC{$z = w \to z = v + y \to w = v + y$\lbl{A6}}
\UIC{$\forall w\,(z = w \to z = v + y \to w = v + y)$}
\UIC{$\forall z\,\forall w\,(z = w \to z = v + y \to w = v + y)$}
\AXC{$\varphi_{7}$\lbl{A2}}
\BIC{$\forall w\,(u + y = w \to u + y = v + y \to w = v + y)$}
\AXC{$\varphi_{26}$\lbl{A2}}
\BIC{$u + y = y + u \to u + y = v + y \to y + u = v + y$}
\AXC{$\stackrel{\Pi_{4}}{\vdots}$}\noLine
\UIC{$(u + y = y + u \to u + y = v + y \to y + u = v + y) \to \varphi_{15}$}
\BIC{$\varphi_{15}$}
\end{ptree}
\noindent$\Pi_{6}$:
\begin{ptree}
\AXC{$\forall x\,\forall y\,(x \in \mathbb{R} \to y \in \mathbb{R} \to x + y = y + x)$\lbl{(R3)}}
\AXC{$\varphi_{28}$\lbl{A2}}
\BIC{$\forall y\,(v \in \mathbb{R} \to y \in \mathbb{R} \to v + y = y + v)$}
\AXC{$\varphi_{23}$\lbl{A2}}
\BIC{$v \in \mathbb{R} \to y \in \mathbb{R} \to v + y = y + v$}
\AXC{$(v \in \mathbb{R} \to y \in \mathbb{R} \to v + y = y + v) \to \varphi_{10} \to \varphi_{9} \to \varphi_{16}$\lbl{A1}}
\BIC{$\varphi_{10} \to \varphi_{9} \to \varphi_{16}$}
\end{ptree}
\noindent$\Pi_{7}$:
\begin{ptree}
\AXC{$\stackrel{\Pi_{5}}{\vdots}$}\noLine
\UIC{$\varphi_{15}$}
\AXC{$\varphi_{9}$\lbl{A1}}
\AXC{$\varphi_{10}$\lbl{A1}}
\AXC{$\stackrel{\Pi_{6}}{\vdots}$}\noLine
\UIC{$\varphi_{10} \to \varphi_{9} \to \varphi_{16}$}
\BIC{$\varphi_{9} \to \varphi_{16}$}
\BIC{$\varphi_{16}$}
\AXC{$\varphi_{16} \to \varphi_{15} \to \varphi_{37}$\lbl{A1}}
\BIC{$\varphi_{15} \to \varphi_{37}$}
\BIC{$\varphi_{37}$}
\end{ptree}
\noindent$\Pi_{8}$:
\begin{ptree}
\AXC{$z = w \to y + u < z \to y + u < w$\lbl{A6}}
\UIC{$\forall w\,(z = w \to y + u < z \to y + u < w)$}
\UIC{$\forall z\,\forall w\,(z = w \to y + u < z \to y + u < w)$}
\AXC{$\varphi_{20}$\lbl{A2}}
\BIC{$\forall w\,(v + y = w \to y + u < v + y \to y + u < w)$}
\AXC{$\varphi_{29}$\lbl{A2}}
\BIC{$v + y = y + v \to y + u < v + y \to y + u < y + v$}
\AXC{$\stackrel{\Pi_{7}}{\vdots}$}\noLine
\UIC{$\varphi_{37}$}
\BIC{$(v + y = y + v \to y + u = v + y \to y + u = y + v) \to \varphi_{17}$}
\end{ptree}
\noindent$\Pi_{9}$:
\begin{ptree}
\AXC{$z = w \to y + u = z \to y + u = w$\lbl{A6}}
\UIC{$\forall w\,(z = w \to y + u = z \to y + u = w)$}
\UIC{$\forall z\,\forall w\,(z = w \to y + u = z \to y + u = w)$}
\AXC{$\varphi_{6}$\lbl{A2}}
\BIC{$\forall w\,(v + y = w \to y + u = v + y \to y + u = w)$}
\AXC{$\varphi_{25}$\lbl{A2}}
\BIC{$v + y = y + v \to y + u = v + y \to y + u = y + v$}
\AXC{$\stackrel{\Pi_{8}}{\vdots}$}\noLine
\UIC{$(v + y = y + v \to y + u = v + y \to y + u = y + v) \to \varphi_{17}$}
\BIC{$\varphi_{17}$}
\end{ptree}
\noindent$\Pi_{10}$:
\begin{ptree}
\AXC{$\varphi_{11}$\lbl{A1}}
\AXC{$\varphi_{9}$\lbl{A1}}
\AXC{$\varphi_{12}$\lbl{A1}}
\AXC{$\forall x\,(\varphi_{4})$\lbl{(R14)}}
\AXC{$\forall x\,(\varphi_{4}) \to \varphi_{4}$\lbl{A2}}
\BIC{$\varphi_{4}$}
\AXC{$\varphi_{35}$\lbl{A2}}
\BIC{$\forall z\,(x \in \mathbb{R} \to y \in \mathbb{R} \to z \in \mathbb{R} \to x < y \to x + z < y + z)$}
\AXC{$\varphi_{34}$\lbl{A2}}
\BIC{$x \in \mathbb{R} \to y \in \mathbb{R} \to u \in \mathbb{R} \to x < y \to x + u < y + u$}
\AXC{$\varphi_{40}$\lbl{A1}}
\BIC{$\varphi_{12} \to \varphi_{9} \to \varphi_{11} \to \varphi_{8} \to \varphi_{18}$}
\BIC{$\varphi_{9} \to \varphi_{11} \to \varphi_{8} \to \varphi_{18}$}
\BIC{$\varphi_{11} \to \varphi_{8} \to \varphi_{18}$}
\BIC{$\varphi_{8} \to \varphi_{18}$}
\end{ptree}
\noindent$\Pi_{11}$:
\begin{ptree}
\AXC{$\varphi_{8}$\lbl{A1}}
\AXC{$\stackrel{\Pi_{10}}{\vdots}$}\noLine
\UIC{$\varphi_{8} \to \varphi_{18}$}
\BIC{$\varphi_{18}$}
\AXC{$\forall x\,\forall y\,\forall z\,(x < y \to y \leq z \to x < z)$\lbl{[O3]}}
\AXC{$\varphi_{5}$\lbl{A2}}
\BIC{$\forall y\,\forall z\,(x + u < y \to y \leq z \to x + u < z)$}
\AXC{$\varphi_{24}$\lbl{A2}}
\BIC{$\forall z\,(x + u < y + u \to y + u \leq z \to x + u < z)$}
\AXC{$\varphi_{31}$\lbl{A2}}
\BIC{$x + u < y + u \to y + u \leq y + v \to x + u < y + v$}
\AXC{$\varphi_{38}$\lbl{A1}}
\BIC{$\varphi_{18} \to \varphi_{17} \to \varphi_{19}$}
\BIC{$\varphi_{17} \to \varphi_{19}$}
\end{ptree}
\noindent 主木:
\begin{ptree}
\AXC{$\stackrel{\Pi_{9}}{\vdots}$}\noLine
\UIC{$\varphi_{17}$}
\AXC{$\stackrel{\Pi_{11}}{\vdots}$}\noLine
\UIC{$\varphi_{17} \to \varphi_{19}$}
\BIC{$\varphi_{19}$}
\UIC{$\forall v\,(\varphi_{19})$}
\UIC{$\forall u\,\forall v\,(\varphi_{19})$}
\UIC{$\forall y\,\forall u\,\forall v\,(\varphi_{19})$}
\UIC{$\forall x\,\forall y\,\forall u\,\forall v\,(\varphi_{19})$}
\end{ptree}
\begin{abbrevlist}
\item $\varphi_{1} :\equiv \forall y\,\forall z\,(u \in \mathbb{R} \to y \in \mathbb{R} \to z \in \mathbb{R} \to u \leq y \to u + z \leq y + z)$
\item $\varphi_{2} :\equiv \forall x\,\forall y\,\forall z\,(x \in \mathbb{R} \to y \in \mathbb{R} \to z \in \mathbb{R} \to x \leq y \to x + z \leq y + z)$
\item $\varphi_{3} :\equiv x \in \mathbb{R} \to y \in \mathbb{R} \to u \in \mathbb{R} \to v \in \mathbb{R} \to x < y \to u \leq v \to u \leq v$
\item $\varphi_{4} :\equiv \forall y\,\forall z\,(x \in \mathbb{R} \to y \in \mathbb{R} \to z \in \mathbb{R} \to x < y \to x + z < y + z)$
\item $\varphi_{5} :\equiv \forall x\,\forall y\,\forall z\,(x < y \to y \leq z \to x < z) \to \forall y\,\forall z\,(x + u < y \to y \leq z \to x + u < z)$
\item $\varphi_{6} :\equiv \forall z\,\forall w\,(z = w \to y + u = z \to y + u = w) \to \forall w\,(v + y = w \to y + u = v + y \to y + u = w)$
\item $\varphi_{7} :\equiv \forall z\,\forall w\,(z = w \to z = v + y \to w = v + y) \to \forall w\,(u + y = w \to u + y = v + y \to w = v + y)$
\item $\varphi_{8} :\equiv x \in \mathbb{R} \to y \in \mathbb{R} \to u \in \mathbb{R} \to v \in \mathbb{R} \to x < y \to u \leq v \to x < y$
\item $\varphi_{9} :\equiv x \in \mathbb{R} \to y \in \mathbb{R} \to u \in \mathbb{R} \to v \in \mathbb{R} \to x < y \to u \leq v \to y \in \mathbb{R}$
\item $\varphi_{10} :\equiv x \in \mathbb{R} \to y \in \mathbb{R} \to u \in \mathbb{R} \to v \in \mathbb{R} \to x < y \to u \leq v \to v \in \mathbb{R}$
\item $\varphi_{11} :\equiv x \in \mathbb{R} \to y \in \mathbb{R} \to u \in \mathbb{R} \to v \in \mathbb{R} \to x < y \to u \leq v \to u \in \mathbb{R}$
\item $\varphi_{12} :\equiv x \in \mathbb{R} \to y \in \mathbb{R} \to u \in \mathbb{R} \to v \in \mathbb{R} \to x < y \to u \leq v \to x \in \mathbb{R}$
\item $\varphi_{13} :\equiv x \in \mathbb{R} \to y \in \mathbb{R} \to u \in \mathbb{R} \to v \in \mathbb{R} \to x < y \to u \leq v \to u + y \leq v + y$
\item $\varphi_{14} :\equiv x \in \mathbb{R} \to y \in \mathbb{R} \to u \in \mathbb{R} \to v \in \mathbb{R} \to x < y \to u \leq v \to u + y = y + u$
\item $\varphi_{15} :\equiv x \in \mathbb{R} \to y \in \mathbb{R} \to u \in \mathbb{R} \to v \in \mathbb{R} \to x < y \to u \leq v \to y + u \leq v + y$
\item $\varphi_{16} :\equiv x \in \mathbb{R} \to y \in \mathbb{R} \to u \in \mathbb{R} \to v \in \mathbb{R} \to x < y \to u \leq v \to v + y = y + v$
\item $\varphi_{17} :\equiv x \in \mathbb{R} \to y \in \mathbb{R} \to u \in \mathbb{R} \to v \in \mathbb{R} \to x < y \to u \leq v \to y + u \leq y + v$
\item $\varphi_{18} :\equiv x \in \mathbb{R} \to y \in \mathbb{R} \to u \in \mathbb{R} \to v \in \mathbb{R} \to x < y \to u \leq v \to x + u < y + u$
\item $\varphi_{19} :\equiv x \in \mathbb{R} \to y \in \mathbb{R} \to u \in \mathbb{R} \to v \in \mathbb{R} \to x < y \to u \leq v \to x + u < y + v$
\item $\varphi_{20} :\equiv \forall z\,\forall w\,(z = w \to y + u < z \to y + u < w) \to \forall w\,(v + y = w \to y + u < v + y \to y + u < w)$
\item $\varphi_{21} :\equiv \forall z\,\forall w\,(z = w \to z < v + y \to w < v + y) \to \forall w\,(u + y = w \to u + y < v + y \to w < v + y)$
\item $\varphi_{22} :\equiv \forall y\,(u \in \mathbb{R} \to y \in \mathbb{R} \to u + y = y + u) \to u \in \mathbb{R} \to y \in \mathbb{R} \to u + y = y + u$
\item $\varphi_{23} :\equiv \forall y\,(v \in \mathbb{R} \to y \in \mathbb{R} \to v + y = y + v) \to v \in \mathbb{R} \to y \in \mathbb{R} \to v + y = y + v$
\item $\varphi_{24} :\equiv \forall y\,\forall z\,(x + u < y \to y \leq z \to x + u < z) \to \forall z\,(x + u < y + u \to y + u \leq z \to x + u < z)$
\item $\varphi_{25} :\equiv \forall w\,(v + y = w \to y + u = v + y \to y + u = w) \to v + y = y + v \to y + u = v + y \to y + u = y + v$
\item $\varphi_{26} :\equiv \forall w\,(u + y = w \to u + y = v + y \to w = v + y) \to u + y = y + u \to u + y = v + y \to y + u = v + y$
\item $\varphi_{27} :\equiv \forall x\,\forall y\,(x \in \mathbb{R} \to y \in \mathbb{R} \to x + y = y + x) \to \forall y\,(u \in \mathbb{R} \to y \in \mathbb{R} \to u + y = y + u)$
\item $\varphi_{28} :\equiv \forall x\,\forall y\,(x \in \mathbb{R} \to y \in \mathbb{R} \to x + y = y + x) \to \forall y\,(v \in \mathbb{R} \to y \in \mathbb{R} \to v + y = y + v)$
\item $\varphi_{29} :\equiv \forall w\,(v + y = w \to y + u < v + y \to y + u < w) \to v + y = y + v \to y + u < v + y \to y + u < y + v$
\item $\varphi_{30} :\equiv \forall w\,(u + y = w \to u + y < v + y \to w < v + y) \to u + y = y + u \to u + y < v + y \to y + u < v + y$
\item $\varphi_{31} :\equiv \forall z\,(x + u < y + u \to y + u \leq z \to x + u < z) \to x + u < y + u \to y + u \leq y + v \to x + u < y + v$
\item $\varphi_{32} :\equiv \forall z\,(u \in \mathbb{R} \to v \in \mathbb{R} \to z \in \mathbb{R} \to u \leq v \to u + z \leq v + z) \to u \in \mathbb{R} \to v \in \mathbb{R} \to y \in \mathbb{R} \to u \leq v \to u + y \leq v + y$
\item $\varphi_{33} :\equiv \varphi_{1} \to \forall z\,(u \in \mathbb{R} \to v \in \mathbb{R} \to z \in \mathbb{R} \to u \leq v \to u + z \leq v + z)$
\item $\varphi_{34} :\equiv \forall z\,(x \in \mathbb{R} \to y \in \mathbb{R} \to z \in \mathbb{R} \to x < y \to x + z < y + z) \to x \in \mathbb{R} \to y \in \mathbb{R} \to u \in \mathbb{R} \to x < y \to x + u < y + u$
\item $\varphi_{35} :\equiv \varphi_{4} \to \forall z\,(x \in \mathbb{R} \to y \in \mathbb{R} \to z \in \mathbb{R} \to x < y \to x + z < y + z)$
\item $\varphi_{36} :\equiv (u + y = y + u \to u + y < v + y \to y + u < v + y) \to (u + y = y + u \to u + y = v + y \to y + u = v + y) \to \varphi_{15}$
\item $\varphi_{37} :\equiv (v + y = y + v \to y + u < v + y \to y + u < y + v) \to (v + y = y + v \to y + u = v + y \to y + u = y + v) \to \varphi_{17}$
\item $\varphi_{38} :\equiv (x + u < y + u \to y + u \leq y + v \to x + u < y + v) \to \varphi_{18} \to \varphi_{17} \to \varphi_{19}$
\item $\varphi_{39} :\equiv (u \in \mathbb{R} \to v \in \mathbb{R} \to y \in \mathbb{R} \to u \leq v \to u + y \leq v + y) \to \varphi_{11} \to \varphi_{10} \to \varphi_{9} \to \varphi_{3} \to \varphi_{13}$
\item $\varphi_{40} :\equiv (x \in \mathbb{R} \to y \in \mathbb{R} \to u \in \mathbb{R} \to x < y \to x + u < y + u) \to \varphi_{12} \to \varphi_{9} \to \varphi_{11} \to \varphi_{8} \to \varphi_{18}$
\end{abbrevlist}
\end{fullproof}

% 自動生成 (prover/build.js): AddLeLe  ノード数 109
\begin{fullproof}{[F21]}{$\forall x\,\forall y\,\forall u\,\forall v\,(x \in \mathbb{R} \to y \in \mathbb{R} \to u \in \mathbb{R} \to v \in \mathbb{R} \to x \leq y \to u \leq v \to x + u \leq y + v)$}{109}\phantomsection\label{tree:AddLeLe}
\noindent$\Pi_{1}$:
\begin{ptree}
\AXC{$\varphi_{9}$\lbl{A1}}
\AXC{$\varphi_{10}$\lbl{A1}}
\AXC{$\varphi_{11}$\lbl{A1}}
\AXC{$\forall x\,(\varphi_{3})$\lbl{[F19]}}
\AXC{$\forall x\,(\varphi_{3}) \to \varphi_{2}$\lbl{A2}}
\BIC{$\varphi_{2}$}
\AXC{$\varphi_{33}$\lbl{A2}}
\BIC{$\forall z\,(u \in \mathbb{R} \to v \in \mathbb{R} \to z \in \mathbb{R} \to u \leq v \to u + z \leq v + z)$}
\AXC{$\varphi_{31}$\lbl{A2}}
\BIC{$u \in \mathbb{R} \to v \in \mathbb{R} \to y \in \mathbb{R} \to u \leq v \to u + y \leq v + y$}
\AXC{$\varphi_{38}$\lbl{A1}}
\BIC{$\varphi_{11} \to \varphi_{10} \to \varphi_{9} \to \varphi_{4} \to \varphi_{13}$}
\BIC{$\varphi_{10} \to \varphi_{9} \to \varphi_{4} \to \varphi_{13}$}
\BIC{$\varphi_{9} \to \varphi_{4} \to \varphi_{13}$}
\BIC{$\varphi_{4} \to \varphi_{13}$}
\end{ptree}
\noindent$\Pi_{2}$:
\begin{ptree}
\AXC{$\forall x\,\forall y\,(x \in \mathbb{R} \to y \in \mathbb{R} \to x + y = y + x)$\lbl{(R3)}}
\AXC{$\varphi_{27}$\lbl{A2}}
\BIC{$\forall y\,(u \in \mathbb{R} \to y \in \mathbb{R} \to u + y = y + u)$}
\AXC{$\varphi_{22}$\lbl{A2}}
\BIC{$u \in \mathbb{R} \to y \in \mathbb{R} \to u + y = y + u$}
\AXC{$(u \in \mathbb{R} \to y \in \mathbb{R} \to u + y = y + u) \to \varphi_{11} \to \varphi_{9} \to \varphi_{14}$\lbl{A1}}
\BIC{$\varphi_{11} \to \varphi_{9} \to \varphi_{14}$}
\end{ptree}
\noindent$\Pi_{3}$:
\begin{ptree}
\AXC{$\varphi_{4}$\lbl{A1}}
\AXC{$\stackrel{\Pi_{1}}{\vdots}$}\noLine
\UIC{$\varphi_{4} \to \varphi_{13}$}
\BIC{$\varphi_{13}$}
\AXC{$\varphi_{9}$\lbl{A1}}
\AXC{$\varphi_{11}$\lbl{A1}}
\AXC{$\stackrel{\Pi_{2}}{\vdots}$}\noLine
\UIC{$\varphi_{11} \to \varphi_{9} \to \varphi_{14}$}
\BIC{$\varphi_{9} \to \varphi_{14}$}
\BIC{$\varphi_{14}$}
\AXC{$\varphi_{14} \to \varphi_{13} \to \varphi_{35}$\lbl{A1}}
\BIC{$\varphi_{13} \to \varphi_{35}$}
\BIC{$\varphi_{35}$}
\end{ptree}
\noindent$\Pi_{4}$:
\begin{ptree}
\AXC{$z = w \to z < v + y \to w < v + y$\lbl{A6}}
\UIC{$\forall w\,(z = w \to z < v + y \to w < v + y)$}
\UIC{$\forall z\,\forall w\,(z = w \to z < v + y \to w < v + y)$}
\AXC{$\varphi_{21}$\lbl{A2}}
\BIC{$\forall w\,(u + y = w \to u + y < v + y \to w < v + y)$}
\AXC{$\varphi_{30}$\lbl{A2}}
\BIC{$u + y = y + u \to u + y < v + y \to y + u < v + y$}
\AXC{$\stackrel{\Pi_{3}}{\vdots}$}\noLine
\UIC{$\varphi_{35}$}
\BIC{$(u + y = y + u \to u + y = v + y \to y + u = v + y) \to \varphi_{15}$}
\end{ptree}
\noindent$\Pi_{5}$:
\begin{ptree}
\AXC{$z = w \to z = v + y \to w = v + y$\lbl{A6}}
\UIC{$\forall w\,(z = w \to z = v + y \to w = v + y)$}
\UIC{$\forall z\,\forall w\,(z = w \to z = v + y \to w = v + y)$}
\AXC{$\varphi_{7}$\lbl{A2}}
\BIC{$\forall w\,(u + y = w \to u + y = v + y \to w = v + y)$}
\AXC{$\varphi_{25}$\lbl{A2}}
\BIC{$u + y = y + u \to u + y = v + y \to y + u = v + y$}
\AXC{$\stackrel{\Pi_{4}}{\vdots}$}\noLine
\UIC{$(u + y = y + u \to u + y = v + y \to y + u = v + y) \to \varphi_{15}$}
\BIC{$\varphi_{15}$}
\end{ptree}
\noindent$\Pi_{6}$:
\begin{ptree}
\AXC{$\forall x\,\forall y\,(x \in \mathbb{R} \to y \in \mathbb{R} \to x + y = y + x)$\lbl{(R3)}}
\AXC{$\varphi_{28}$\lbl{A2}}
\BIC{$\forall y\,(v \in \mathbb{R} \to y \in \mathbb{R} \to v + y = y + v)$}
\AXC{$\varphi_{23}$\lbl{A2}}
\BIC{$v \in \mathbb{R} \to y \in \mathbb{R} \to v + y = y + v$}
\AXC{$(v \in \mathbb{R} \to y \in \mathbb{R} \to v + y = y + v) \to \varphi_{10} \to \varphi_{9} \to \varphi_{16}$\lbl{A1}}
\BIC{$\varphi_{10} \to \varphi_{9} \to \varphi_{16}$}
\end{ptree}
\noindent$\Pi_{7}$:
\begin{ptree}
\AXC{$\stackrel{\Pi_{5}}{\vdots}$}\noLine
\UIC{$\varphi_{15}$}
\AXC{$\varphi_{9}$\lbl{A1}}
\AXC{$\varphi_{10}$\lbl{A1}}
\AXC{$\stackrel{\Pi_{6}}{\vdots}$}\noLine
\UIC{$\varphi_{10} \to \varphi_{9} \to \varphi_{16}$}
\BIC{$\varphi_{9} \to \varphi_{16}$}
\BIC{$\varphi_{16}$}
\AXC{$\varphi_{16} \to \varphi_{15} \to \varphi_{36}$\lbl{A1}}
\BIC{$\varphi_{15} \to \varphi_{36}$}
\BIC{$\varphi_{36}$}
\end{ptree}
\noindent$\Pi_{8}$:
\begin{ptree}
\AXC{$z = w \to y + u < z \to y + u < w$\lbl{A6}}
\UIC{$\forall w\,(z = w \to y + u < z \to y + u < w)$}
\UIC{$\forall z\,\forall w\,(z = w \to y + u < z \to y + u < w)$}
\AXC{$\varphi_{20}$\lbl{A2}}
\BIC{$\forall w\,(v + y = w \to y + u < v + y \to y + u < w)$}
\AXC{$\varphi_{29}$\lbl{A2}}
\BIC{$v + y = y + v \to y + u < v + y \to y + u < y + v$}
\AXC{$\stackrel{\Pi_{7}}{\vdots}$}\noLine
\UIC{$\varphi_{36}$}
\BIC{$(v + y = y + v \to y + u = v + y \to y + u = y + v) \to \varphi_{17}$}
\end{ptree}
\noindent$\Pi_{9}$:
\begin{ptree}
\AXC{$z = w \to y + u = z \to y + u = w$\lbl{A6}}
\UIC{$\forall w\,(z = w \to y + u = z \to y + u = w)$}
\UIC{$\forall z\,\forall w\,(z = w \to y + u = z \to y + u = w)$}
\AXC{$\varphi_{6}$\lbl{A2}}
\BIC{$\forall w\,(v + y = w \to y + u = v + y \to y + u = w)$}
\AXC{$\varphi_{24}$\lbl{A2}}
\BIC{$v + y = y + v \to y + u = v + y \to y + u = y + v$}
\AXC{$\stackrel{\Pi_{8}}{\vdots}$}\noLine
\UIC{$(v + y = y + v \to y + u = v + y \to y + u = y + v) \to \varphi_{17}$}
\BIC{$\varphi_{17}$}
\end{ptree}
\noindent$\Pi_{10}$:
\begin{ptree}
\AXC{$\varphi_{11}$\lbl{A1}}
\AXC{$\varphi_{9}$\lbl{A1}}
\AXC{$\varphi_{12}$\lbl{A1}}
\AXC{$\forall x\,(\varphi_{3})$\lbl{[F19]}}
\AXC{$\forall x\,(\varphi_{3}) \to \varphi_{3}$\lbl{A2}}
\BIC{$\varphi_{3}$}
\AXC{$\varphi_{34}$\lbl{A2}}
\BIC{$\forall z\,(x \in \mathbb{R} \to y \in \mathbb{R} \to z \in \mathbb{R} \to x \leq y \to x + z \leq y + z)$}
\AXC{$\varphi_{32}$\lbl{A2}}
\BIC{$x \in \mathbb{R} \to y \in \mathbb{R} \to u \in \mathbb{R} \to x \leq y \to x + u \leq y + u$}
\AXC{$\varphi_{39}$\lbl{A1}}
\BIC{$\varphi_{12} \to \varphi_{9} \to \varphi_{11} \to \varphi_{5} \to \varphi_{18}$}
\BIC{$\varphi_{9} \to \varphi_{11} \to \varphi_{5} \to \varphi_{18}$}
\BIC{$\varphi_{11} \to \varphi_{5} \to \varphi_{18}$}
\BIC{$\varphi_{5} \to \varphi_{18}$}
\end{ptree}
\noindent$\Pi_{11}$:
\begin{ptree}
\AXC{$\varphi_{5}$\lbl{A1}}
\AXC{$\stackrel{\Pi_{10}}{\vdots}$}\noLine
\UIC{$\varphi_{5} \to \varphi_{18}$}
\BIC{$\varphi_{18}$}
\AXC{$\forall x\,\forall y\,\forall z\,(x \leq y \to y \leq z \to x \leq z)$\lbl{[O5]}}
\AXC{$\varphi_{1}$\lbl{A2}}
\BIC{$\forall y\,\forall z\,(x + u \leq y \to y \leq z \to x + u \leq z)$}
\AXC{$\varphi_{8}$\lbl{A2}}
\BIC{$\forall z\,(x + u \leq y + u \to y + u \leq z \to x + u \leq z)$}
\AXC{$\varphi_{26}$\lbl{A2}}
\BIC{$x + u \leq y + u \to y + u \leq y + v \to x + u \leq y + v$}
\AXC{$\varphi_{37}$\lbl{A1}}
\BIC{$\varphi_{18} \to \varphi_{17} \to \varphi_{19}$}
\BIC{$\varphi_{17} \to \varphi_{19}$}
\end{ptree}
\noindent 主木:
\begin{ptree}
\AXC{$\stackrel{\Pi_{9}}{\vdots}$}\noLine
\UIC{$\varphi_{17}$}
\AXC{$\stackrel{\Pi_{11}}{\vdots}$}\noLine
\UIC{$\varphi_{17} \to \varphi_{19}$}
\BIC{$\varphi_{19}$}
\UIC{$\forall v\,(\varphi_{19})$}
\UIC{$\forall u\,\forall v\,(\varphi_{19})$}
\UIC{$\forall y\,\forall u\,\forall v\,(\varphi_{19})$}
\UIC{$\forall x\,\forall y\,\forall u\,\forall v\,(\varphi_{19})$}
\end{ptree}
\begin{abbrevlist}
\item $\varphi_{1} :\equiv \forall x\,\forall y\,\forall z\,(x \leq y \to y \leq z \to x \leq z) \to \forall y\,\forall z\,(x + u \leq y \to y \leq z \to x + u \leq z)$
\item $\varphi_{2} :\equiv \forall y\,\forall z\,(u \in \mathbb{R} \to y \in \mathbb{R} \to z \in \mathbb{R} \to u \leq y \to u + z \leq y + z)$
\item $\varphi_{3} :\equiv \forall y\,\forall z\,(x \in \mathbb{R} \to y \in \mathbb{R} \to z \in \mathbb{R} \to x \leq y \to x + z \leq y + z)$
\item $\varphi_{4} :\equiv x \in \mathbb{R} \to y \in \mathbb{R} \to u \in \mathbb{R} \to v \in \mathbb{R} \to x \leq y \to u \leq v \to u \leq v$
\item $\varphi_{5} :\equiv x \in \mathbb{R} \to y \in \mathbb{R} \to u \in \mathbb{R} \to v \in \mathbb{R} \to x \leq y \to u \leq v \to x \leq y$
\item $\varphi_{6} :\equiv \forall z\,\forall w\,(z = w \to y + u = z \to y + u = w) \to \forall w\,(v + y = w \to y + u = v + y \to y + u = w)$
\item $\varphi_{7} :\equiv \forall z\,\forall w\,(z = w \to z = v + y \to w = v + y) \to \forall w\,(u + y = w \to u + y = v + y \to w = v + y)$
\item $\varphi_{8} :\equiv \forall y\,\forall z\,(x + u \leq y \to y \leq z \to x + u \leq z) \to \forall z\,(x + u \leq y + u \to y + u \leq z \to x + u \leq z)$
\item $\varphi_{9} :\equiv x \in \mathbb{R} \to y \in \mathbb{R} \to u \in \mathbb{R} \to v \in \mathbb{R} \to x \leq y \to u \leq v \to y \in \mathbb{R}$
\item $\varphi_{10} :\equiv x \in \mathbb{R} \to y \in \mathbb{R} \to u \in \mathbb{R} \to v \in \mathbb{R} \to x \leq y \to u \leq v \to v \in \mathbb{R}$
\item $\varphi_{11} :\equiv x \in \mathbb{R} \to y \in \mathbb{R} \to u \in \mathbb{R} \to v \in \mathbb{R} \to x \leq y \to u \leq v \to u \in \mathbb{R}$
\item $\varphi_{12} :\equiv x \in \mathbb{R} \to y \in \mathbb{R} \to u \in \mathbb{R} \to v \in \mathbb{R} \to x \leq y \to u \leq v \to x \in \mathbb{R}$
\item $\varphi_{13} :\equiv x \in \mathbb{R} \to y \in \mathbb{R} \to u \in \mathbb{R} \to v \in \mathbb{R} \to x \leq y \to u \leq v \to u + y \leq v + y$
\item $\varphi_{14} :\equiv x \in \mathbb{R} \to y \in \mathbb{R} \to u \in \mathbb{R} \to v \in \mathbb{R} \to x \leq y \to u \leq v \to u + y = y + u$
\item $\varphi_{15} :\equiv x \in \mathbb{R} \to y \in \mathbb{R} \to u \in \mathbb{R} \to v \in \mathbb{R} \to x \leq y \to u \leq v \to y + u \leq v + y$
\item $\varphi_{16} :\equiv x \in \mathbb{R} \to y \in \mathbb{R} \to u \in \mathbb{R} \to v \in \mathbb{R} \to x \leq y \to u \leq v \to v + y = y + v$
\item $\varphi_{17} :\equiv x \in \mathbb{R} \to y \in \mathbb{R} \to u \in \mathbb{R} \to v \in \mathbb{R} \to x \leq y \to u \leq v \to y + u \leq y + v$
\item $\varphi_{18} :\equiv x \in \mathbb{R} \to y \in \mathbb{R} \to u \in \mathbb{R} \to v \in \mathbb{R} \to x \leq y \to u \leq v \to x + u \leq y + u$
\item $\varphi_{19} :\equiv x \in \mathbb{R} \to y \in \mathbb{R} \to u \in \mathbb{R} \to v \in \mathbb{R} \to x \leq y \to u \leq v \to x + u \leq y + v$
\item $\varphi_{20} :\equiv \forall z\,\forall w\,(z = w \to y + u < z \to y + u < w) \to \forall w\,(v + y = w \to y + u < v + y \to y + u < w)$
\item $\varphi_{21} :\equiv \forall z\,\forall w\,(z = w \to z < v + y \to w < v + y) \to \forall w\,(u + y = w \to u + y < v + y \to w < v + y)$
\item $\varphi_{22} :\equiv \forall y\,(u \in \mathbb{R} \to y \in \mathbb{R} \to u + y = y + u) \to u \in \mathbb{R} \to y \in \mathbb{R} \to u + y = y + u$
\item $\varphi_{23} :\equiv \forall y\,(v \in \mathbb{R} \to y \in \mathbb{R} \to v + y = y + v) \to v \in \mathbb{R} \to y \in \mathbb{R} \to v + y = y + v$
\item $\varphi_{24} :\equiv \forall w\,(v + y = w \to y + u = v + y \to y + u = w) \to v + y = y + v \to y + u = v + y \to y + u = y + v$
\item $\varphi_{25} :\equiv \forall w\,(u + y = w \to u + y = v + y \to w = v + y) \to u + y = y + u \to u + y = v + y \to y + u = v + y$
\item $\varphi_{26} :\equiv \forall z\,(x + u \leq y + u \to y + u \leq z \to x + u \leq z) \to x + u \leq y + u \to y + u \leq y + v \to x + u \leq y + v$
\item $\varphi_{27} :\equiv \forall x\,\forall y\,(x \in \mathbb{R} \to y \in \mathbb{R} \to x + y = y + x) \to \forall y\,(u \in \mathbb{R} \to y \in \mathbb{R} \to u + y = y + u)$
\item $\varphi_{28} :\equiv \forall x\,\forall y\,(x \in \mathbb{R} \to y \in \mathbb{R} \to x + y = y + x) \to \forall y\,(v \in \mathbb{R} \to y \in \mathbb{R} \to v + y = y + v)$
\item $\varphi_{29} :\equiv \forall w\,(v + y = w \to y + u < v + y \to y + u < w) \to v + y = y + v \to y + u < v + y \to y + u < y + v$
\item $\varphi_{30} :\equiv \forall w\,(u + y = w \to u + y < v + y \to w < v + y) \to u + y = y + u \to u + y < v + y \to y + u < v + y$
\item $\varphi_{31} :\equiv \forall z\,(u \in \mathbb{R} \to v \in \mathbb{R} \to z \in \mathbb{R} \to u \leq v \to u + z \leq v + z) \to u \in \mathbb{R} \to v \in \mathbb{R} \to y \in \mathbb{R} \to u \leq v \to u + y \leq v + y$
\item $\varphi_{32} :\equiv \forall z\,(x \in \mathbb{R} \to y \in \mathbb{R} \to z \in \mathbb{R} \to x \leq y \to x + z \leq y + z) \to x \in \mathbb{R} \to y \in \mathbb{R} \to u \in \mathbb{R} \to x \leq y \to x + u \leq y + u$
\item $\varphi_{33} :\equiv \varphi_{2} \to \forall z\,(u \in \mathbb{R} \to v \in \mathbb{R} \to z \in \mathbb{R} \to u \leq v \to u + z \leq v + z)$
\item $\varphi_{34} :\equiv \varphi_{3} \to \forall z\,(x \in \mathbb{R} \to y \in \mathbb{R} \to z \in \mathbb{R} \to x \leq y \to x + z \leq y + z)$
\item $\varphi_{35} :\equiv (u + y = y + u \to u + y < v + y \to y + u < v + y) \to (u + y = y + u \to u + y = v + y \to y + u = v + y) \to \varphi_{15}$
\item $\varphi_{36} :\equiv (v + y = y + v \to y + u < v + y \to y + u < y + v) \to (v + y = y + v \to y + u = v + y \to y + u = y + v) \to \varphi_{17}$
\item $\varphi_{37} :\equiv (x + u \leq y + u \to y + u \leq y + v \to x + u \leq y + v) \to \varphi_{18} \to \varphi_{17} \to \varphi_{19}$
\item $\varphi_{38} :\equiv (u \in \mathbb{R} \to v \in \mathbb{R} \to y \in \mathbb{R} \to u \leq v \to u + y \leq v + y) \to \varphi_{11} \to \varphi_{10} \to \varphi_{9} \to \varphi_{4} \to \varphi_{13}$
\item $\varphi_{39} :\equiv (x \in \mathbb{R} \to y \in \mathbb{R} \to u \in \mathbb{R} \to x \leq y \to x + u \leq y + u) \to \varphi_{12} \to \varphi_{9} \to \varphi_{11} \to \varphi_{5} \to \varphi_{18}$
\end{abbrevlist}
\end{fullproof}

% 自動生成 (prover/build.js): NegLt  ノード数 365
\begin{fullproof}{[F22]}{$\forall x\,\forall y\,(x \in \mathbb{R} \to y \in \mathbb{R} \to x < y \to -y < -x)$}{365}\phantomsection\label{tree:NegLt}
\noindent$\Pi_{1}$:
\begin{ptree}
\AXC{$z = w \to -y < z \to -y < w$\lbl{A6}}
\UIC{$\forall w\,(z = w \to -y < z \to -y < w)$}
\UIC{$\forall z\,\forall w\,(z = w \to -y < z \to -y < w)$}
\AXC{$\forall z\,\forall w\,(z = w \to -y < z \to -y < w) \to \varphi_{36}$\lbl{A2}}
\BIC{$\varphi_{36}$}
\AXC{$\varphi_{36} \to \varphi_{55}$\lbl{A2}}
\BIC{$\varphi_{55}$}
\end{ptree}
\noindent$\Pi_{2}$:
\begin{ptree}
\AXC{$z = w \to z < y + (-x + -y) \to w < y + (-x + -y)$\lbl{A6}}
\UIC{$\varphi_{25}$}
\UIC{$\forall z\,(\varphi_{25})$}
\AXC{$\forall z\,(\varphi_{25}) \to \varphi_{89}$\lbl{A2}}
\BIC{$\varphi_{89}$}
\AXC{$\varphi_{89} \to x + (-x + -y) = -y \to \varphi_{54}$\lbl{A2}}
\BIC{$x + (-x + -y) = -y \to \varphi_{54}$}
\end{ptree}
\noindent$\Pi_{3}$:
\begin{ptree}
\AXC{$\forall x \in \mathbb{R}\,(-x \in \mathbb{R})$\lbl{(R1)}}
\AXC{$\forall x \in \mathbb{R}\,(-x \in \mathbb{R}) \to y \in \mathbb{R} \to -y \in \mathbb{R}$\lbl{A2}}
\BIC{$y \in \mathbb{R} \to -y \in \mathbb{R}$}
\AXC{$(y \in \mathbb{R} \to -y \in \mathbb{R}) \to \varphi_{82}$\lbl{A1}}
\BIC{$\varphi_{82}$}
\end{ptree}
\noindent$\Pi_{4}$:
\begin{ptree}
\AXC{$x \in \mathbb{R} \to y \in \mathbb{R} \to x < y \to y \in \mathbb{R}$\lbl{A1}}
\AXC{$\stackrel{\Pi_{3}}{\vdots}$}\noLine
\UIC{$\varphi_{82}$}
\BIC{$x \in \mathbb{R} \to y \in \mathbb{R} \to x < y \to -y \in \mathbb{R}$}
\end{ptree}
\noindent$\Pi_{5}$:
\begin{ptree}
\AXC{$\forall x \in \mathbb{R}\,(-x \in \mathbb{R})$\lbl{(R1)}}
\AXC{$\forall x \in \mathbb{R}\,(-x \in \mathbb{R}) \to x \in \mathbb{R} \to -x \in \mathbb{R}$\lbl{A2}}
\BIC{$x \in \mathbb{R} \to -x \in \mathbb{R}$}
\AXC{$(x \in \mathbb{R} \to -x \in \mathbb{R}) \to \varphi_{83}$\lbl{A1}}
\BIC{$\varphi_{83}$}
\end{ptree}
\noindent$\Pi_{6}$:
\begin{ptree}
\AXC{$x \in \mathbb{R} \to y \in \mathbb{R} \to x < y \to x \in \mathbb{R}$\lbl{A1}}
\AXC{$\stackrel{\Pi_{5}}{\vdots}$}\noLine
\UIC{$\varphi_{83}$}
\BIC{$x \in \mathbb{R} \to y \in \mathbb{R} \to x < y \to -x \in \mathbb{R}$}
\end{ptree}
\noindent$\Pi_{7}$:
\begin{ptree}
\AXC{$\stackrel{\Pi_{6}}{\vdots}$}\noLine
\UIC{$x \in \mathbb{R} \to y \in \mathbb{R} \to x < y \to -x \in \mathbb{R}$}
\AXC{$\forall x\,\forall y\,(x \in \mathbb{R} \to y \in \mathbb{R} \to x + y \in \mathbb{R})$\lbl{(R1)}}
\AXC{$\varphi_{90}$\lbl{A2}}
\BIC{$\forall y\,(-x \in \mathbb{R} \to y \in \mathbb{R} \to -x + y \in \mathbb{R})$}
\AXC{$\varphi_{102}$\lbl{A2}}
\BIC{$-x \in \mathbb{R} \to -y \in \mathbb{R} \to -x + -y \in \mathbb{R}$}
\AXC{$\varphi_{121}$\lbl{A1}}
\BIC{$(x \in \mathbb{R} \to y \in \mathbb{R} \to x < y \to -x \in \mathbb{R}) \to \varphi_{99}$}
\BIC{$\varphi_{99}$}
\end{ptree}
\noindent$\Pi_{8}$:
\begin{ptree}
\AXC{$\forall x\,(\varphi_{35})$\lbl{(R14)}}
\AXC{$\forall x\,(\varphi_{35}) \to \varphi_{35}$\lbl{A2}}
\BIC{$\varphi_{35}$}
\AXC{$\varphi_{105}$\lbl{A2}}
\BIC{$\forall z\,(x \in \mathbb{R} \to y \in \mathbb{R} \to z \in \mathbb{R} \to x < y \to x + z < y + z)$}
\AXC{$\varphi_{108}$\lbl{A2}}
\BIC{$x \in \mathbb{R} \to y \in \mathbb{R} \to \varphi_{58}$}
\AXC{$(x \in \mathbb{R} \to y \in \mathbb{R} \to \varphi_{58}) \to \varphi_{128}$\lbl{A1}}
\BIC{$\varphi_{128}$}
\end{ptree}
\noindent$\Pi_{9}$:
\begin{ptree}
\AXC{$x \in \mathbb{R} \to y \in \mathbb{R} \to x < y \to y \in \mathbb{R}$\lbl{A1}}
\AXC{$x \in \mathbb{R} \to y \in \mathbb{R} \to x < y \to x \in \mathbb{R}$\lbl{A1}}
\AXC{$\stackrel{\Pi_{8}}{\vdots}$}\noLine
\UIC{$\varphi_{128}$}
\BIC{$(x \in \mathbb{R} \to y \in \mathbb{R} \to x < y \to y \in \mathbb{R}) \to \varphi_{112}$}
\BIC{$\varphi_{112}$}
\end{ptree}
\noindent$\Pi_{10}$:
\begin{ptree}
\AXC{$x \in \mathbb{R} \to y \in \mathbb{R} \to x < y \to x < y$\lbl{A1}}
\AXC{$\stackrel{\Pi_{4}}{\vdots}$}\noLine
\UIC{$x \in \mathbb{R} \to y \in \mathbb{R} \to x < y \to -y \in \mathbb{R}$}
\AXC{$\stackrel{\Pi_{7}}{\vdots}$}\noLine
\UIC{$\varphi_{99}$}
\BIC{$x \in \mathbb{R} \to y \in \mathbb{R} \to x < y \to -x + -y \in \mathbb{R}$}
\AXC{$\stackrel{\Pi_{9}}{\vdots}$}\noLine
\UIC{$\varphi_{112}$}
\BIC{$(x \in \mathbb{R} \to y \in \mathbb{R} \to x < y \to x < y) \to \varphi_{56}$}
\BIC{$\varphi_{56}$}
\end{ptree}
\noindent$\Pi_{11}$:
\begin{ptree}
\AXC{$y = z \to x = y \to x = z$\lbl{A6}}
\AXC{$(y = z \to x = y \to x = z) \to x = y \to y = z \to x = z$\lbl{A1}}
\BIC{$x = y \to y = z \to x = z$}
\UIC{$\forall z\,(x = y \to y = z \to x = z)$}
\UIC{$\forall y\,\forall z\,(x = y \to y = z \to x = z)$}
\UIC{$\forall x\,\forall y\,\forall z\,(x = y \to y = z \to x = z)$}
\AXC{$\varphi_{10}$\lbl{A2}}
\BIC{$\forall y\,\forall z\,(v_{227} = y \to y = z \to v_{227} = z)$}
\AXC{$\varphi_{18}$\lbl{A2}}
\BIC{$\forall z\,(v_{227} = v_{228} \to v_{228} = z \to v_{227} = z)$}
\AXC{$\varphi_{27}$\lbl{A2}}
\BIC{$v_{227} = v_{228} \to v_{228} = v_{229} \to v_{227} = v_{229}$}
\UIC{$\forall v_{229}\,(v_{227} = v_{228} \to v_{228} = v_{229} \to v_{227} = v_{229})$}
\UIC{$\varphi_{3}$}
\UIC{$\forall v_{227}\,(\varphi_{3})$}
\end{ptree}
\noindent$\Pi_{12}$:
\begin{ptree}
\AXC{$\stackrel{\Pi_{11}}{\vdots}$}\noLine
\UIC{$\forall v_{227}\,(\varphi_{3})$}
\AXC{$\forall v_{227}\,(\varphi_{3}) \to \varphi_{46}$\lbl{A2}}
\BIC{$\varphi_{46}$}
\AXC{$\varphi_{46} \to \varphi_{61}$\lbl{A2}}
\BIC{$\varphi_{61}$}
\AXC{$\varphi_{61} \to \varphi_{59}$\lbl{A2}}
\BIC{$\varphi_{59}$}
\end{ptree}
\noindent$\Pi_{13}$:
\begin{ptree}
\AXC{$\forall x \in \mathbb{R}\,(x + 0 = x)$\lbl{(R4)}}
\AXC{$\forall x \in \mathbb{R}\,(x + 0 = x) \to -y \in \mathbb{R} \to -y + 0 = -y$\lbl{A2}}
\BIC{$-y \in \mathbb{R} \to -y + 0 = -y$}
\AXC{$(-y \in \mathbb{R} \to -y + 0 = -y) \to \varphi_{92}$\lbl{A1}}
\BIC{$\varphi_{92}$}
\end{ptree}
\noindent$\Pi_{14}$:
\begin{ptree}
\AXC{$y = z \to x = y \to x = z$\lbl{A6}}
\AXC{$(y = z \to x = y \to x = z) \to x = y \to y = z \to x = z$\lbl{A1}}
\BIC{$x = y \to y = z \to x = z$}
\UIC{$\forall z\,(x = y \to y = z \to x = z)$}
\UIC{$\forall y\,\forall z\,(x = y \to y = z \to x = z)$}
\UIC{$\forall x\,\forall y\,\forall z\,(x = y \to y = z \to x = z)$}
\AXC{$\varphi_{11}$\lbl{A2}}
\BIC{$\forall y\,\forall z\,(v_{224} = y \to y = z \to v_{224} = z)$}
\AXC{$\varphi_{19}$\lbl{A2}}
\BIC{$\forall z\,(v_{224} = v_{225} \to v_{225} = z \to v_{224} = z)$}
\AXC{$\varphi_{28}$\lbl{A2}}
\BIC{$v_{224} = v_{225} \to v_{225} = v_{226} \to v_{224} = v_{226}$}
\UIC{$\forall v_{226}\,(v_{224} = v_{225} \to v_{225} = v_{226} \to v_{224} = v_{226})$}
\UIC{$\varphi_{4}$}
\UIC{$\forall v_{224}\,(\varphi_{4})$}
\end{ptree}
\noindent$\Pi_{15}$:
\begin{ptree}
\AXC{$\stackrel{\Pi_{14}}{\vdots}$}\noLine
\UIC{$\forall v_{224}\,(\varphi_{4})$}
\AXC{$\forall v_{224}\,(\varphi_{4}) \to \varphi_{47}$\lbl{A2}}
\BIC{$\varphi_{47}$}
\AXC{$\varphi_{47} \to \varphi_{62}$\lbl{A2}}
\BIC{$\varphi_{62}$}
\AXC{$\varphi_{62} \to \varphi_{74}$\lbl{A2}}
\BIC{$\varphi_{74}$}
\end{ptree}
\noindent$\Pi_{16}$:
\begin{ptree}
\AXC{$\forall x\,\forall y\,(x \in \mathbb{R} \to y \in \mathbb{R} \to x + y = y + x)$\lbl{(R3)}}
\AXC{$\varphi_{69}$\lbl{A2}}
\BIC{$\forall y\,(0 \in \mathbb{R} \to y \in \mathbb{R} \to 0 + y = y + 0)$}
\AXC{$\varphi_{87}$\lbl{A2}}
\BIC{$0 \in \mathbb{R} \to -y \in \mathbb{R} \to 0 + -y = -y + 0$}
\end{ptree}
\noindent$\Pi_{17}$:
\begin{ptree}
\AXC{$0 \in \mathbb{R}$\lbl{(R1)}}
\AXC{$\stackrel{\Pi_{16}}{\vdots}$}\noLine
\UIC{$0 \in \mathbb{R} \to -y \in \mathbb{R} \to 0 + -y = -y + 0$}
\AXC{$(0 \in \mathbb{R} \to -y \in \mathbb{R} \to 0 + -y = -y + 0) \to 0 \in \mathbb{R} \to \varphi_{97}$\lbl{A1}}
\BIC{$0 \in \mathbb{R} \to \varphi_{97}$}
\BIC{$\varphi_{97}$}
\end{ptree}
\noindent$\Pi_{18}$:
\begin{ptree}
\AXC{$y = z \to x = y \to x = z$\lbl{A6}}
\AXC{$(y = z \to x = y \to x = z) \to x = y \to y = z \to x = z$\lbl{A1}}
\BIC{$x = y \to y = z \to x = z$}
\UIC{$\forall z\,(x = y \to y = z \to x = z)$}
\UIC{$\forall y\,\forall z\,(x = y \to y = z \to x = z)$}
\UIC{$\forall x\,\forall y\,\forall z\,(x = y \to y = z \to x = z)$}
\AXC{$\varphi_{12}$\lbl{A2}}
\BIC{$\forall y\,\forall z\,(v_{221} = y \to y = z \to v_{221} = z)$}
\AXC{$\varphi_{20}$\lbl{A2}}
\BIC{$\forall z\,(v_{221} = v_{222} \to v_{222} = z \to v_{221} = z)$}
\AXC{$\varphi_{29}$\lbl{A2}}
\BIC{$v_{221} = v_{222} \to v_{222} = v_{223} \to v_{221} = v_{223}$}
\UIC{$\forall v_{223}\,(v_{221} = v_{222} \to v_{222} = v_{223} \to v_{221} = v_{223})$}
\UIC{$\varphi_{5}$}
\UIC{$\forall v_{221}\,(\varphi_{5})$}
\end{ptree}
\noindent$\Pi_{19}$:
\begin{ptree}
\AXC{$\stackrel{\Pi_{18}}{\vdots}$}\noLine
\UIC{$\forall v_{221}\,(\varphi_{5})$}
\AXC{$\forall v_{221}\,(\varphi_{5}) \to \varphi_{48}$\lbl{A2}}
\BIC{$\varphi_{48}$}
\AXC{$\varphi_{48} \to \varphi_{94}$\lbl{A2}}
\BIC{$\varphi_{94}$}
\AXC{$\varphi_{94} \to \varphi_{100}$\lbl{A2}}
\BIC{$\varphi_{100}$}
\end{ptree}
\noindent$\Pi_{20}$:
\begin{ptree}
\AXC{$z = w \to x + -x + -y = z + -y \to x + -x + -y = w + -y$\lbl{A6}}
\UIC{$\varphi_{38}$}
\UIC{$\forall z\,(\varphi_{38})$}
\AXC{$\forall z\,(\varphi_{38}) \to \varphi_{84}$\lbl{A2}}
\BIC{$\varphi_{84}$}
\AXC{$\varphi_{84} \to \varphi_{75}$\lbl{A2}}
\BIC{$\varphi_{75}$}
\end{ptree}
\noindent$\Pi_{21}$:
\begin{ptree}
\AXC{$\forall x \in \mathbb{R}\,(x + -x = 0)$\lbl{(R5)}}
\AXC{$\forall x \in \mathbb{R}\,(x + -x = 0) \to x \in \mathbb{R} \to x + -x = 0$\lbl{A2}}
\BIC{$x \in \mathbb{R} \to x + -x = 0$}
\AXC{$(x \in \mathbb{R} \to x + -x = 0) \to \varphi_{78}$\lbl{A1}}
\BIC{$\varphi_{78}$}
\end{ptree}
\noindent$\Pi_{22}$:
\begin{ptree}
\AXC{$x \in \mathbb{R} \to y \in \mathbb{R} \to x < y \to x \in \mathbb{R}$\lbl{A1}}
\AXC{$\stackrel{\Pi_{21}}{\vdots}$}\noLine
\UIC{$\varphi_{78}$}
\BIC{$x \in \mathbb{R} \to y \in \mathbb{R} \to x < y \to x + -x = 0$}
\AXC{$(x \in \mathbb{R} \to y \in \mathbb{R} \to x < y \to x + -x = 0) \to \varphi_{117}$\lbl{A1}}
\BIC{$\varphi_{117}$}
\end{ptree}
\noindent$\Pi_{23}$:
\begin{ptree}
\AXC{$x = x$\lbl{A5}}
\AXC{$x = y \to x = x \to y = x$\lbl{A6}}
\AXC{$(x = y \to x = x \to y = x) \to x = x \to x = y \to y = x$\lbl{A1}}
\BIC{$x = x \to x = y \to y = x$}
\BIC{$x = y \to y = x$}
\UIC{$\forall y\,(x = y \to y = x)$}
\UIC{$\forall x\,\forall y\,(x = y \to y = x)$}
\end{ptree}
\noindent$\Pi_{24}$:
\begin{ptree}
\AXC{$\stackrel{\Pi_{23}}{\vdots}$}\noLine
\UIC{$\forall x\,\forall y\,(x = y \to y = x)$}
\AXC{$\forall x\,\forall y\,(x = y \to y = x) \to \forall y\,(v_{219} = y \to y = v_{219})$\lbl{A2}}
\BIC{$\forall y\,(v_{219} = y \to y = v_{219})$}
\AXC{$\varphi_{1}$\lbl{A2}}
\BIC{$v_{219} = v_{220} \to v_{220} = v_{219}$}
\UIC{$\forall v_{220}\,(v_{219} = v_{220} \to v_{220} = v_{219})$}
\UIC{$\forall v_{219}\,\forall v_{220}\,(v_{219} = v_{220} \to v_{220} = v_{219})$}
\AXC{$\varphi_{80}$\lbl{A2}}
\BIC{$\forall v_{220}\,(x + -x + -y = v_{220} \to v_{220} = x + -x + -y)$}
\AXC{$\varphi_{106}$\lbl{A2}}
\BIC{$\varphi_{65}$}
\end{ptree}
\noindent$\Pi_{25}$:
\begin{ptree}
\AXC{$\forall x\,(\varphi_{34})$\lbl{(R2)}}
\AXC{$\forall x\,(\varphi_{34}) \to \varphi_{34}$\lbl{A2}}
\BIC{$\varphi_{34}$}
\AXC{$\varphi_{34} \to \varphi_{53}$\lbl{A2}}
\BIC{$\varphi_{53}$}
\AXC{$\varphi_{53} \to \varphi_{70}$\lbl{A2}}
\BIC{$\varphi_{70}$}
\AXC{$\varphi_{70} \to (x \in \mathbb{R} \to y \in \mathbb{R} \to x < y \to x \in \mathbb{R}) \to \varphi_{110}$\lbl{A1}}
\BIC{$(x \in \mathbb{R} \to y \in \mathbb{R} \to x < y \to x \in \mathbb{R}) \to \varphi_{110}$}
\end{ptree}
\noindent$\Pi_{26}$:
\begin{ptree}
\AXC{$\stackrel{\Pi_{6}}{\vdots}$}\noLine
\UIC{$x \in \mathbb{R} \to y \in \mathbb{R} \to x < y \to -x \in \mathbb{R}$}
\AXC{$x \in \mathbb{R} \to y \in \mathbb{R} \to x < y \to x \in \mathbb{R}$\lbl{A1}}
\AXC{$\stackrel{\Pi_{25}}{\vdots}$}\noLine
\UIC{$(x \in \mathbb{R} \to y \in \mathbb{R} \to x < y \to x \in \mathbb{R}) \to \varphi_{110}$}
\BIC{$\varphi_{110}$}
\BIC{$(x \in \mathbb{R} \to y \in \mathbb{R} \to x < y \to -y \in \mathbb{R}) \to \varphi_{39}$}
\end{ptree}
\noindent$\Pi_{27}$:
\begin{ptree}
\AXC{$\stackrel{\Pi_{24}}{\vdots}$}\noLine
\UIC{$\varphi_{65}$}
\AXC{$\stackrel{\Pi_{4}}{\vdots}$}\noLine
\UIC{$x \in \mathbb{R} \to y \in \mathbb{R} \to x < y \to -y \in \mathbb{R}$}
\AXC{$\stackrel{\Pi_{26}}{\vdots}$}\noLine
\UIC{$(x \in \mathbb{R} \to y \in \mathbb{R} \to x < y \to -y \in \mathbb{R}) \to \varphi_{39}$}
\BIC{$\varphi_{39}$}
\AXC{$\varphi_{39} \to \varphi_{65} \to \varphi_{40}$\lbl{A1}}
\BIC{$\varphi_{65} \to \varphi_{40}$}
\BIC{$\varphi_{40}$}
\end{ptree}
\noindent$\Pi_{28}$:
\begin{ptree}
\AXC{$\stackrel{\Pi_{20}}{\vdots}$}\noLine
\UIC{$\varphi_{75}$}
\AXC{$x = x$\lbl{A5}}
\UIC{$\forall x\,(x = x)$}
\AXC{$\forall x\,(x = x) \to x + -x + -y = x + -x + -y$\lbl{A2}}
\BIC{$x + -x + -y = x + -x + -y$}
\AXC{$\stackrel{\Pi_{22}}{\vdots}$}\noLine
\UIC{$\varphi_{117}$}
\BIC{$\varphi_{75} \to x \in \mathbb{R} \to y \in \mathbb{R} \to x < y \to x + -x + -y = 0 + -y$}
\BIC{$x \in \mathbb{R} \to y \in \mathbb{R} \to x < y \to x + -x + -y = 0 + -y$}
\AXC{$\stackrel{\Pi_{27}}{\vdots}$}\noLine
\UIC{$\varphi_{40}$}
\AXC{$\varphi_{40} \to \varphi_{124}$\lbl{A1}}
\BIC{$\varphi_{124}$}
\BIC{$\varphi_{100} \to x \in \mathbb{R} \to y \in \mathbb{R} \to x < y \to x + (-x + -y) = 0 + -y$}
\end{ptree}
\noindent$\Pi_{29}$:
\begin{ptree}
\AXC{$\stackrel{\Pi_{19}}{\vdots}$}\noLine
\UIC{$\varphi_{100}$}
\AXC{$\stackrel{\Pi_{28}}{\vdots}$}\noLine
\UIC{$\varphi_{100} \to x \in \mathbb{R} \to y \in \mathbb{R} \to x < y \to x + (-x + -y) = 0 + -y$}
\BIC{$x \in \mathbb{R} \to y \in \mathbb{R} \to x < y \to x + (-x + -y) = 0 + -y$}
\AXC{$\varphi_{129}$\lbl{A1}}
\BIC{$\varphi_{119}$}
\end{ptree}
\noindent$\Pi_{30}$:
\begin{ptree}
\AXC{$\stackrel{\Pi_{15}}{\vdots}$}\noLine
\UIC{$\varphi_{74}$}
\AXC{$\stackrel{\Pi_{4}}{\vdots}$}\noLine
\UIC{$x \in \mathbb{R} \to y \in \mathbb{R} \to x < y \to -y \in \mathbb{R}$}
\AXC{$\stackrel{\Pi_{17}}{\vdots}$}\noLine
\UIC{$\varphi_{97}$}
\BIC{$x \in \mathbb{R} \to y \in \mathbb{R} \to x < y \to 0 + -y = -y + 0$}
\AXC{$\stackrel{\Pi_{29}}{\vdots}$}\noLine
\UIC{$\varphi_{119}$}
\BIC{$\varphi_{74} \to x \in \mathbb{R} \to y \in \mathbb{R} \to x < y \to x + (-x + -y) = -y + 0$}
\BIC{$x \in \mathbb{R} \to y \in \mathbb{R} \to x < y \to x + (-x + -y) = -y + 0$}
\AXC{$\varphi_{126}$\lbl{A1}}
\BIC{$\varphi_{115}$}
\end{ptree}
\noindent$\Pi_{31}$:
\begin{ptree}
\AXC{$\stackrel{\Pi_{10}}{\vdots}$}\noLine
\UIC{$\varphi_{56}$}
\AXC{$\stackrel{\Pi_{12}}{\vdots}$}\noLine
\UIC{$\varphi_{59}$}
\AXC{$\stackrel{\Pi_{4}}{\vdots}$}\noLine
\UIC{$x \in \mathbb{R} \to y \in \mathbb{R} \to x < y \to -y \in \mathbb{R}$}
\AXC{$\stackrel{\Pi_{13}}{\vdots}$}\noLine
\UIC{$\varphi_{92}$}
\BIC{$x \in \mathbb{R} \to y \in \mathbb{R} \to x < y \to -y + 0 = -y$}
\AXC{$\stackrel{\Pi_{30}}{\vdots}$}\noLine
\UIC{$\varphi_{115}$}
\BIC{$\varphi_{59} \to x \in \mathbb{R} \to y \in \mathbb{R} \to x < y \to x + (-x + -y) = -y$}
\BIC{$x \in \mathbb{R} \to y \in \mathbb{R} \to x < y \to x + (-x + -y) = -y$}
\AXC{$\varphi_{131}$\lbl{A1}}
\BIC{$\varphi_{56} \to \varphi_{109}$}
\BIC{$\varphi_{109}$}
\end{ptree}
\noindent$\Pi_{32}$:
\begin{ptree}
\AXC{$y = z \to x = y \to x = z$\lbl{A6}}
\AXC{$(y = z \to x = y \to x = z) \to x = y \to y = z \to x = z$\lbl{A1}}
\BIC{$x = y \to y = z \to x = z$}
\UIC{$\forall z\,(x = y \to y = z \to x = z)$}
\UIC{$\forall y\,\forall z\,(x = y \to y = z \to x = z)$}
\UIC{$\forall x\,\forall y\,\forall z\,(x = y \to y = z \to x = z)$}
\AXC{$\varphi_{13}$\lbl{A2}}
\BIC{$\forall y\,\forall z\,(v_{244} = y \to y = z \to v_{244} = z)$}
\AXC{$\varphi_{21}$\lbl{A2}}
\BIC{$\forall z\,(v_{244} = v_{245} \to v_{245} = z \to v_{244} = z)$}
\AXC{$\varphi_{30}$\lbl{A2}}
\BIC{$v_{244} = v_{245} \to v_{245} = v_{246} \to v_{244} = v_{246}$}
\UIC{$\forall v_{246}\,(v_{244} = v_{245} \to v_{245} = v_{246} \to v_{244} = v_{246})$}
\UIC{$\varphi_{6}$}
\UIC{$\forall v_{244}\,(\varphi_{6})$}
\end{ptree}
\noindent$\Pi_{33}$:
\begin{ptree}
\AXC{$\stackrel{\Pi_{32}}{\vdots}$}\noLine
\UIC{$\forall v_{244}\,(\varphi_{6})$}
\AXC{$\forall v_{244}\,(\varphi_{6}) \to \varphi_{49}$\lbl{A2}}
\BIC{$\varphi_{49}$}
\AXC{$\varphi_{49} \to \varphi_{63}$\lbl{A2}}
\BIC{$\varphi_{63}$}
\AXC{$\varphi_{63} \to \varphi_{60}$\lbl{A2}}
\BIC{$\varphi_{60}$}
\end{ptree}
\noindent$\Pi_{34}$:
\begin{ptree}
\AXC{$\forall x \in \mathbb{R}\,(x + 0 = x)$\lbl{(R4)}}
\AXC{$\forall x \in \mathbb{R}\,(x + 0 = x) \to -x \in \mathbb{R} \to -x + 0 = -x$\lbl{A2}}
\BIC{$-x \in \mathbb{R} \to -x + 0 = -x$}
\AXC{$(-x \in \mathbb{R} \to -x + 0 = -x) \to \varphi_{93}$\lbl{A1}}
\BIC{$\varphi_{93}$}
\end{ptree}
\noindent$\Pi_{35}$:
\begin{ptree}
\AXC{$y = z \to x = y \to x = z$\lbl{A6}}
\AXC{$(y = z \to x = y \to x = z) \to x = y \to y = z \to x = z$\lbl{A1}}
\BIC{$x = y \to y = z \to x = z$}
\UIC{$\forall z\,(x = y \to y = z \to x = z)$}
\UIC{$\forall y\,\forall z\,(x = y \to y = z \to x = z)$}
\UIC{$\forall x\,\forall y\,\forall z\,(x = y \to y = z \to x = z)$}
\AXC{$\varphi_{14}$\lbl{A2}}
\BIC{$\forall y\,\forall z\,(v_{241} = y \to y = z \to v_{241} = z)$}
\AXC{$\varphi_{22}$\lbl{A2}}
\BIC{$\forall z\,(v_{241} = v_{242} \to v_{242} = z \to v_{241} = z)$}
\AXC{$\varphi_{31}$\lbl{A2}}
\BIC{$v_{241} = v_{242} \to v_{242} = v_{243} \to v_{241} = v_{243}$}
\UIC{$\forall v_{243}\,(v_{241} = v_{242} \to v_{242} = v_{243} \to v_{241} = v_{243})$}
\UIC{$\varphi_{7}$}
\UIC{$\forall v_{241}\,(\varphi_{7})$}
\end{ptree}
\noindent$\Pi_{36}$:
\begin{ptree}
\AXC{$\stackrel{\Pi_{35}}{\vdots}$}\noLine
\UIC{$\forall v_{241}\,(\varphi_{7})$}
\AXC{$\forall v_{241}\,(\varphi_{7}) \to \varphi_{50}$\lbl{A2}}
\BIC{$\varphi_{50}$}
\AXC{$\varphi_{50} \to \varphi_{64}$\lbl{A2}}
\BIC{$\varphi_{64}$}
\AXC{$\varphi_{64} \to \varphi_{76}$\lbl{A2}}
\BIC{$\varphi_{76}$}
\end{ptree}
\noindent$\Pi_{37}$:
\begin{ptree}
\AXC{$\forall x\,\forall y\,(x \in \mathbb{R} \to y \in \mathbb{R} \to x + y = y + x)$\lbl{(R3)}}
\AXC{$\varphi_{69}$\lbl{A2}}
\BIC{$\forall y\,(0 \in \mathbb{R} \to y \in \mathbb{R} \to 0 + y = y + 0)$}
\AXC{$\varphi_{88}$\lbl{A2}}
\BIC{$0 \in \mathbb{R} \to -x \in \mathbb{R} \to 0 + -x = -x + 0$}
\end{ptree}
\noindent$\Pi_{38}$:
\begin{ptree}
\AXC{$0 \in \mathbb{R}$\lbl{(R1)}}
\AXC{$\stackrel{\Pi_{37}}{\vdots}$}\noLine
\UIC{$0 \in \mathbb{R} \to -x \in \mathbb{R} \to 0 + -x = -x + 0$}
\AXC{$(0 \in \mathbb{R} \to -x \in \mathbb{R} \to 0 + -x = -x + 0) \to 0 \in \mathbb{R} \to \varphi_{98}$\lbl{A1}}
\BIC{$0 \in \mathbb{R} \to \varphi_{98}$}
\BIC{$\varphi_{98}$}
\end{ptree}
\noindent$\Pi_{39}$:
\begin{ptree}
\AXC{$y = z \to x = y \to x = z$\lbl{A6}}
\AXC{$(y = z \to x = y \to x = z) \to x = y \to y = z \to x = z$\lbl{A1}}
\BIC{$x = y \to y = z \to x = z$}
\UIC{$\forall z\,(x = y \to y = z \to x = z)$}
\UIC{$\forall y\,\forall z\,(x = y \to y = z \to x = z)$}
\UIC{$\forall x\,\forall y\,\forall z\,(x = y \to y = z \to x = z)$}
\AXC{$\varphi_{15}$\lbl{A2}}
\BIC{$\forall y\,\forall z\,(v_{238} = y \to y = z \to v_{238} = z)$}
\AXC{$\varphi_{23}$\lbl{A2}}
\BIC{$\forall z\,(v_{238} = v_{239} \to v_{239} = z \to v_{238} = z)$}
\AXC{$\varphi_{32}$\lbl{A2}}
\BIC{$v_{238} = v_{239} \to v_{239} = v_{240} \to v_{238} = v_{240}$}
\UIC{$\forall v_{240}\,(v_{238} = v_{239} \to v_{239} = v_{240} \to v_{238} = v_{240})$}
\UIC{$\varphi_{8}$}
\UIC{$\forall v_{238}\,(\varphi_{8})$}
\end{ptree}
\noindent$\Pi_{40}$:
\begin{ptree}
\AXC{$\stackrel{\Pi_{39}}{\vdots}$}\noLine
\UIC{$\forall v_{238}\,(\varphi_{8})$}
\AXC{$\forall v_{238}\,(\varphi_{8}) \to \varphi_{51}$\lbl{A2}}
\BIC{$\varphi_{51}$}
\AXC{$\varphi_{51} \to \varphi_{95}$\lbl{A2}}
\BIC{$\varphi_{95}$}
\AXC{$\varphi_{95} \to \varphi_{101}$\lbl{A2}}
\BIC{$\varphi_{101}$}
\end{ptree}
\noindent$\Pi_{41}$:
\begin{ptree}
\AXC{$z = w \to y + -y + -x = z + -x \to y + -y + -x = w + -x$\lbl{A6}}
\UIC{$\varphi_{41}$}
\UIC{$\forall z\,(\varphi_{41})$}
\AXC{$\forall z\,(\varphi_{41}) \to \varphi_{85}$\lbl{A2}}
\BIC{$\varphi_{85}$}
\AXC{$\varphi_{85} \to \varphi_{77}$\lbl{A2}}
\BIC{$\varphi_{77}$}
\end{ptree}
\noindent$\Pi_{42}$:
\begin{ptree}
\AXC{$\forall x \in \mathbb{R}\,(x + -x = 0)$\lbl{(R5)}}
\AXC{$\forall x \in \mathbb{R}\,(x + -x = 0) \to y \in \mathbb{R} \to y + -y = 0$\lbl{A2}}
\BIC{$y \in \mathbb{R} \to y + -y = 0$}
\AXC{$(y \in \mathbb{R} \to y + -y = 0) \to \varphi_{79}$\lbl{A1}}
\BIC{$\varphi_{79}$}
\end{ptree}
\noindent$\Pi_{43}$:
\begin{ptree}
\AXC{$x \in \mathbb{R} \to y \in \mathbb{R} \to x < y \to y \in \mathbb{R}$\lbl{A1}}
\AXC{$\stackrel{\Pi_{42}}{\vdots}$}\noLine
\UIC{$\varphi_{79}$}
\BIC{$x \in \mathbb{R} \to y \in \mathbb{R} \to x < y \to y + -y = 0$}
\AXC{$(x \in \mathbb{R} \to y \in \mathbb{R} \to x < y \to y + -y = 0) \to \varphi_{118}$\lbl{A1}}
\BIC{$\varphi_{118}$}
\end{ptree}
\noindent$\Pi_{44}$:
\begin{ptree}
\AXC{$y = z \to x = y \to x = z$\lbl{A6}}
\AXC{$(y = z \to x = y \to x = z) \to x = y \to y = z \to x = z$\lbl{A1}}
\BIC{$x = y \to y = z \to x = z$}
\UIC{$\forall z\,(x = y \to y = z \to x = z)$}
\UIC{$\forall y\,\forall z\,(x = y \to y = z \to x = z)$}
\UIC{$\forall x\,\forall y\,\forall z\,(x = y \to y = z \to x = z)$}
\AXC{$\varphi_{16}$\lbl{A2}}
\BIC{$\forall y\,\forall z\,(v_{235} = y \to y = z \to v_{235} = z)$}
\AXC{$\varphi_{24}$\lbl{A2}}
\BIC{$\forall z\,(v_{235} = v_{236} \to v_{236} = z \to v_{235} = z)$}
\AXC{$\varphi_{33}$\lbl{A2}}
\BIC{$v_{235} = v_{236} \to v_{236} = v_{237} \to v_{235} = v_{237}$}
\UIC{$\forall v_{237}\,(v_{235} = v_{236} \to v_{236} = v_{237} \to v_{235} = v_{237})$}
\UIC{$\varphi_{9}$}
\UIC{$\forall v_{235}\,(\varphi_{9})$}
\end{ptree}
\noindent$\Pi_{45}$:
\begin{ptree}
\AXC{$\stackrel{\Pi_{44}}{\vdots}$}\noLine
\UIC{$\forall v_{235}\,(\varphi_{9})$}
\AXC{$\forall v_{235}\,(\varphi_{9}) \to \varphi_{52}$\lbl{A2}}
\BIC{$\varphi_{52}$}
\AXC{$\varphi_{52} \to \varphi_{96}$\lbl{A2}}
\BIC{$\varphi_{96}$}
\AXC{$\varphi_{96} \to y + (-x + -y) = y + (-y + -x) \to \varphi_{67}$\lbl{A2}}
\BIC{$y + (-x + -y) = y + (-y + -x) \to \varphi_{67}$}
\end{ptree}
\noindent$\Pi_{46}$:
\begin{ptree}
\AXC{$\stackrel{\Pi_{23}}{\vdots}$}\noLine
\UIC{$\forall x\,\forall y\,(x = y \to y = x)$}
\AXC{$\forall x\,\forall y\,(x = y \to y = x) \to \forall y\,(v_{233} = y \to y = v_{233})$\lbl{A2}}
\BIC{$\forall y\,(v_{233} = y \to y = v_{233})$}
\AXC{$\varphi_{2}$\lbl{A2}}
\BIC{$v_{233} = v_{234} \to v_{234} = v_{233}$}
\UIC{$\forall v_{234}\,(v_{233} = v_{234} \to v_{234} = v_{233})$}
\UIC{$\forall v_{233}\,\forall v_{234}\,(v_{233} = v_{234} \to v_{234} = v_{233})$}
\AXC{$\varphi_{81}$\lbl{A2}}
\BIC{$\forall v_{234}\,(y + -y + -x = v_{234} \to v_{234} = y + -y + -x)$}
\AXC{$\varphi_{107}$\lbl{A2}}
\BIC{$\varphi_{66}$}
\end{ptree}
\noindent$\Pi_{47}$:
\begin{ptree}
\AXC{$\forall x\,(\varphi_{34})$\lbl{(R2)}}
\AXC{$\forall x\,(\varphi_{34}) \to \varphi_{37}$\lbl{A2}}
\BIC{$\varphi_{37}$}
\AXC{$\varphi_{37} \to \varphi_{57}$\lbl{A2}}
\BIC{$\varphi_{57}$}
\AXC{$\varphi_{57} \to v_{230} \in \mathbb{R} \to v_{231} \in \mathbb{R} \to \varphi_{17}$\lbl{A2}}
\BIC{$v_{230} \in \mathbb{R} \to v_{231} \in \mathbb{R} \to \varphi_{17}$}
\UIC{$\forall v_{232}\,(v_{230} \in \mathbb{R} \to v_{231} \in \mathbb{R} \to \varphi_{17})$}
\UIC{$\forall v_{231}\,\forall v_{232}\,(v_{230} \in \mathbb{R} \to v_{231} \in \mathbb{R} \to \varphi_{17})$}
\UIC{$\forall v_{230}\,\forall v_{231}\,\forall v_{232}\,(v_{230} \in \mathbb{R} \to v_{231} \in \mathbb{R} \to \varphi_{17})$}
\AXC{$\varphi_{113}$\lbl{A2}}
\BIC{$\varphi_{71}$}
\AXC{$\varphi_{71} \to \varphi_{72}$\lbl{A2}}
\BIC{$\varphi_{72}$}
\end{ptree}
\noindent$\Pi_{48}$:
\begin{ptree}
\AXC{$x \in \mathbb{R} \to y \in \mathbb{R} \to x < y \to y \in \mathbb{R}$\lbl{A1}}
\AXC{$\stackrel{\Pi_{47}}{\vdots}$}\noLine
\UIC{$\varphi_{72}$}
\AXC{$\varphi_{72} \to \varphi_{73}$\lbl{A2}}
\BIC{$\varphi_{73}$}
\AXC{$\varphi_{73} \to (x \in \mathbb{R} \to y \in \mathbb{R} \to x < y \to y \in \mathbb{R}) \to \varphi_{111}$\lbl{A1}}
\BIC{$(x \in \mathbb{R} \to y \in \mathbb{R} \to x < y \to y \in \mathbb{R}) \to \varphi_{111}$}
\BIC{$\varphi_{111}$}
\end{ptree}
\noindent$\Pi_{49}$:
\begin{ptree}
\AXC{$\stackrel{\Pi_{46}}{\vdots}$}\noLine
\UIC{$\varphi_{66}$}
\AXC{$\stackrel{\Pi_{6}}{\vdots}$}\noLine
\UIC{$x \in \mathbb{R} \to y \in \mathbb{R} \to x < y \to -x \in \mathbb{R}$}
\AXC{$\stackrel{\Pi_{4}}{\vdots}$}\noLine
\UIC{$x \in \mathbb{R} \to y \in \mathbb{R} \to x < y \to -y \in \mathbb{R}$}
\AXC{$\stackrel{\Pi_{48}}{\vdots}$}\noLine
\UIC{$\varphi_{111}$}
\BIC{$(x \in \mathbb{R} \to y \in \mathbb{R} \to x < y \to -x \in \mathbb{R}) \to \varphi_{42}$}
\BIC{$\varphi_{42}$}
\AXC{$\varphi_{42} \to \varphi_{66} \to \varphi_{43}$\lbl{A1}}
\BIC{$\varphi_{66} \to \varphi_{43}$}
\BIC{$\varphi_{43}$}
\end{ptree}
\noindent$\Pi_{50}$:
\begin{ptree}
\AXC{$\stackrel{\Pi_{6}}{\vdots}$}\noLine
\UIC{$x \in \mathbb{R} \to y \in \mathbb{R} \to x < y \to -x \in \mathbb{R}$}
\AXC{$\forall x\,\forall y\,(x \in \mathbb{R} \to y \in \mathbb{R} \to x + y = y + x)$\lbl{(R3)}}
\AXC{$\varphi_{91}$\lbl{A2}}
\BIC{$\forall y\,(-x \in \mathbb{R} \to y \in \mathbb{R} \to -x + y = y + -x)$}
\AXC{$\varphi_{104}$\lbl{A2}}
\BIC{$-x \in \mathbb{R} \to -y \in \mathbb{R} \to -x + -y = -y + -x$}
\AXC{$\varphi_{122}$\lbl{A1}}
\BIC{$(x \in \mathbb{R} \to y \in \mathbb{R} \to x < y \to -x \in \mathbb{R}) \to \varphi_{103}$}
\BIC{$\varphi_{103}$}
\end{ptree}
\noindent$\Pi_{51}$:
\begin{ptree}
\AXC{$\stackrel{\Pi_{4}}{\vdots}$}\noLine
\UIC{$x \in \mathbb{R} \to y \in \mathbb{R} \to x < y \to -y \in \mathbb{R}$}
\AXC{$\stackrel{\Pi_{50}}{\vdots}$}\noLine
\UIC{$\varphi_{103}$}
\BIC{$x \in \mathbb{R} \to y \in \mathbb{R} \to x < y \to -x + -y = -y + -x$}
\AXC{$(x \in \mathbb{R} \to y \in \mathbb{R} \to x < y \to -x + -y = -y + -x) \to \varphi_{123}$\lbl{A1}}
\BIC{$\varphi_{123}$}
\end{ptree}
\noindent$\Pi_{52}$:
\begin{ptree}
\AXC{$x = x$\lbl{A5}}
\UIC{$\forall x\,(x = x)$}
\AXC{$\forall x\,(x = x) \to y + (-x + -y) = y + (-x + -y)$\lbl{A2}}
\BIC{$y + (-x + -y) = y + (-x + -y)$}
\AXC{$\stackrel{\Pi_{51}}{\vdots}$}\noLine
\UIC{$\varphi_{123}$}
\BIC{$(-x + -y = -y + -x \to \varphi_{68}) \to \varphi_{44}$}
\end{ptree}
\noindent$\Pi_{53}$:
\begin{ptree}
\AXC{$\varphi_{26}$\lbl{A6}}
\UIC{$\forall w\,(\varphi_{26})$}
\UIC{$\forall z\,\forall w\,(\varphi_{26})$}
\AXC{$\forall z\,\forall w\,(\varphi_{26}) \to \varphi_{86}$\lbl{A2}}
\BIC{$\varphi_{86}$}
\AXC{$\varphi_{86} \to -x + -y = -y + -x \to \varphi_{68}$\lbl{A2}}
\BIC{$-x + -y = -y + -x \to \varphi_{68}$}
\AXC{$\stackrel{\Pi_{52}}{\vdots}$}\noLine
\UIC{$(-x + -y = -y + -x \to \varphi_{68}) \to \varphi_{44}$}
\BIC{$\varphi_{44}$}
\end{ptree}
\noindent$\Pi_{54}$:
\begin{ptree}
\AXC{$\stackrel{\Pi_{45}}{\vdots}$}\noLine
\UIC{$y + (-x + -y) = y + (-y + -x) \to \varphi_{67}$}
\AXC{$\stackrel{\Pi_{49}}{\vdots}$}\noLine
\UIC{$\varphi_{43}$}
\AXC{$\stackrel{\Pi_{53}}{\vdots}$}\noLine
\UIC{$\varphi_{44}$}
\AXC{$\varphi_{44} \to \varphi_{43} \to (y + (-x + -y) = y + (-y + -x) \to \varphi_{67}) \to \varphi_{45}$\lbl{A1}}
\BIC{$\varphi_{43} \to (y + (-x + -y) = y + (-y + -x) \to \varphi_{67}) \to \varphi_{45}$}
\BIC{$(y + (-x + -y) = y + (-y + -x) \to \varphi_{67}) \to \varphi_{45}$}
\BIC{$\varphi_{45}$}
\end{ptree}
\noindent$\Pi_{55}$:
\begin{ptree}
\AXC{$\stackrel{\Pi_{41}}{\vdots}$}\noLine
\UIC{$\varphi_{77}$}
\AXC{$x = x$\lbl{A5}}
\UIC{$\forall x\,(x = x)$}
\AXC{$\forall x\,(x = x) \to y + -y + -x = y + -y + -x$\lbl{A2}}
\BIC{$y + -y + -x = y + -y + -x$}
\AXC{$\stackrel{\Pi_{43}}{\vdots}$}\noLine
\UIC{$\varphi_{118}$}
\BIC{$\varphi_{77} \to x \in \mathbb{R} \to y \in \mathbb{R} \to x < y \to y + -y + -x = 0 + -x$}
\BIC{$x \in \mathbb{R} \to y \in \mathbb{R} \to x < y \to y + -y + -x = 0 + -x$}
\AXC{$\stackrel{\Pi_{54}}{\vdots}$}\noLine
\UIC{$\varphi_{45}$}
\AXC{$\varphi_{45} \to \varphi_{125}$\lbl{A1}}
\BIC{$\varphi_{125}$}
\BIC{$\varphi_{101} \to x \in \mathbb{R} \to y \in \mathbb{R} \to x < y \to y + (-x + -y) = 0 + -x$}
\end{ptree}
\noindent$\Pi_{56}$:
\begin{ptree}
\AXC{$\stackrel{\Pi_{40}}{\vdots}$}\noLine
\UIC{$\varphi_{101}$}
\AXC{$\stackrel{\Pi_{55}}{\vdots}$}\noLine
\UIC{$\varphi_{101} \to x \in \mathbb{R} \to y \in \mathbb{R} \to x < y \to y + (-x + -y) = 0 + -x$}
\BIC{$x \in \mathbb{R} \to y \in \mathbb{R} \to x < y \to y + (-x + -y) = 0 + -x$}
\AXC{$\varphi_{130}$\lbl{A1}}
\BIC{$\varphi_{120}$}
\end{ptree}
\noindent$\Pi_{57}$:
\begin{ptree}
\AXC{$\stackrel{\Pi_{36}}{\vdots}$}\noLine
\UIC{$\varphi_{76}$}
\AXC{$\stackrel{\Pi_{6}}{\vdots}$}\noLine
\UIC{$x \in \mathbb{R} \to y \in \mathbb{R} \to x < y \to -x \in \mathbb{R}$}
\AXC{$\stackrel{\Pi_{38}}{\vdots}$}\noLine
\UIC{$\varphi_{98}$}
\BIC{$x \in \mathbb{R} \to y \in \mathbb{R} \to x < y \to 0 + -x = -x + 0$}
\AXC{$\stackrel{\Pi_{56}}{\vdots}$}\noLine
\UIC{$\varphi_{120}$}
\BIC{$\varphi_{76} \to x \in \mathbb{R} \to y \in \mathbb{R} \to x < y \to y + (-x + -y) = -x + 0$}
\BIC{$x \in \mathbb{R} \to y \in \mathbb{R} \to x < y \to y + (-x + -y) = -x + 0$}
\AXC{$\varphi_{127}$\lbl{A1}}
\BIC{$\varphi_{116}$}
\end{ptree}
\noindent$\Pi_{58}$:
\begin{ptree}
\AXC{$\stackrel{\Pi_{33}}{\vdots}$}\noLine
\UIC{$\varphi_{60}$}
\AXC{$\stackrel{\Pi_{6}}{\vdots}$}\noLine
\UIC{$x \in \mathbb{R} \to y \in \mathbb{R} \to x < y \to -x \in \mathbb{R}$}
\AXC{$\stackrel{\Pi_{34}}{\vdots}$}\noLine
\UIC{$\varphi_{93}$}
\BIC{$x \in \mathbb{R} \to y \in \mathbb{R} \to x < y \to -x + 0 = -x$}
\AXC{$\stackrel{\Pi_{57}}{\vdots}$}\noLine
\UIC{$\varphi_{116}$}
\BIC{$\varphi_{60} \to x \in \mathbb{R} \to y \in \mathbb{R} \to x < y \to y + (-x + -y) = -x$}
\BIC{$x \in \mathbb{R} \to y \in \mathbb{R} \to x < y \to y + (-x + -y) = -x$}
\end{ptree}
\noindent$\Pi_{59}$:
\begin{ptree}
\AXC{$\stackrel{\Pi_{58}}{\vdots}$}\noLine
\UIC{$x \in \mathbb{R} \to y \in \mathbb{R} \to x < y \to y + (-x + -y) = -x$}
\AXC{$(x \in \mathbb{R} \to y \in \mathbb{R} \to x < y \to y + (-x + -y) = -x) \to \varphi_{114}$\lbl{A1}}
\BIC{$\varphi_{114}$}
\end{ptree}
\noindent 主木:
\begin{ptree}
\AXC{$\stackrel{\Pi_{1}}{\vdots}$}\noLine
\UIC{$\varphi_{55}$}
\AXC{$\stackrel{\Pi_{2}}{\vdots}$}\noLine
\UIC{$x + (-x + -y) = -y \to \varphi_{54}$}
\AXC{$\stackrel{\Pi_{31}}{\vdots}$}\noLine
\UIC{$\varphi_{109}$}
\BIC{$x \in \mathbb{R} \to y \in \mathbb{R} \to x < y \to -y < y + (-x + -y)$}
\AXC{$\stackrel{\Pi_{59}}{\vdots}$}\noLine
\UIC{$\varphi_{114}$}
\BIC{$\varphi_{55} \to x \in \mathbb{R} \to y \in \mathbb{R} \to x < y \to -y < -x$}
\BIC{$x \in \mathbb{R} \to y \in \mathbb{R} \to x < y \to -y < -x$}
\UIC{$\forall y\,(x \in \mathbb{R} \to y \in \mathbb{R} \to x < y \to -y < -x)$}
\UIC{$\forall x\,\forall y\,(x \in \mathbb{R} \to y \in \mathbb{R} \to x < y \to -y < -x)$}
\end{ptree}
\begin{abbrevlist}
\item $\varphi_{1} :\equiv \forall y\,(v_{219} = y \to y = v_{219}) \to v_{219} = v_{220} \to v_{220} = v_{219}$
\item $\varphi_{2} :\equiv \forall y\,(v_{233} = y \to y = v_{233}) \to v_{233} = v_{234} \to v_{234} = v_{233}$
\item $\varphi_{3} :\equiv \forall v_{228}\,\forall v_{229}\,(v_{227} = v_{228} \to v_{228} = v_{229} \to v_{227} = v_{229})$
\item $\varphi_{4} :\equiv \forall v_{225}\,\forall v_{226}\,(v_{224} = v_{225} \to v_{225} = v_{226} \to v_{224} = v_{226})$
\item $\varphi_{5} :\equiv \forall v_{222}\,\forall v_{223}\,(v_{221} = v_{222} \to v_{222} = v_{223} \to v_{221} = v_{223})$
\item $\varphi_{6} :\equiv \forall v_{245}\,\forall v_{246}\,(v_{244} = v_{245} \to v_{245} = v_{246} \to v_{244} = v_{246})$
\item $\varphi_{7} :\equiv \forall v_{242}\,\forall v_{243}\,(v_{241} = v_{242} \to v_{242} = v_{243} \to v_{241} = v_{243})$
\item $\varphi_{8} :\equiv \forall v_{239}\,\forall v_{240}\,(v_{238} = v_{239} \to v_{239} = v_{240} \to v_{238} = v_{240})$
\item $\varphi_{9} :\equiv \forall v_{236}\,\forall v_{237}\,(v_{235} = v_{236} \to v_{236} = v_{237} \to v_{235} = v_{237})$
\item $\varphi_{10} :\equiv \forall x\,\forall y\,\forall z\,(x = y \to y = z \to x = z) \to \forall y\,\forall z\,(v_{227} = y \to y = z \to v_{227} = z)$
\item $\varphi_{11} :\equiv \forall x\,\forall y\,\forall z\,(x = y \to y = z \to x = z) \to \forall y\,\forall z\,(v_{224} = y \to y = z \to v_{224} = z)$
\item $\varphi_{12} :\equiv \forall x\,\forall y\,\forall z\,(x = y \to y = z \to x = z) \to \forall y\,\forall z\,(v_{221} = y \to y = z \to v_{221} = z)$
\item $\varphi_{13} :\equiv \forall x\,\forall y\,\forall z\,(x = y \to y = z \to x = z) \to \forall y\,\forall z\,(v_{244} = y \to y = z \to v_{244} = z)$
\item $\varphi_{14} :\equiv \forall x\,\forall y\,\forall z\,(x = y \to y = z \to x = z) \to \forall y\,\forall z\,(v_{241} = y \to y = z \to v_{241} = z)$
\item $\varphi_{15} :\equiv \forall x\,\forall y\,\forall z\,(x = y \to y = z \to x = z) \to \forall y\,\forall z\,(v_{238} = y \to y = z \to v_{238} = z)$
\item $\varphi_{16} :\equiv \forall x\,\forall y\,\forall z\,(x = y \to y = z \to x = z) \to \forall y\,\forall z\,(v_{235} = y \to y = z \to v_{235} = z)$
\item $\varphi_{17} :\equiv v_{232} \in \mathbb{R} \to v_{230} + v_{231} + v_{232} = v_{230} + (v_{231} + v_{232})$
\item $\varphi_{18} :\equiv \forall y\,\forall z\,(v_{227} = y \to y = z \to v_{227} = z) \to \forall z\,(v_{227} = v_{228} \to v_{228} = z \to v_{227} = z)$
\item $\varphi_{19} :\equiv \forall y\,\forall z\,(v_{224} = y \to y = z \to v_{224} = z) \to \forall z\,(v_{224} = v_{225} \to v_{225} = z \to v_{224} = z)$
\item $\varphi_{20} :\equiv \forall y\,\forall z\,(v_{221} = y \to y = z \to v_{221} = z) \to \forall z\,(v_{221} = v_{222} \to v_{222} = z \to v_{221} = z)$
\item $\varphi_{21} :\equiv \forall y\,\forall z\,(v_{244} = y \to y = z \to v_{244} = z) \to \forall z\,(v_{244} = v_{245} \to v_{245} = z \to v_{244} = z)$
\item $\varphi_{22} :\equiv \forall y\,\forall z\,(v_{241} = y \to y = z \to v_{241} = z) \to \forall z\,(v_{241} = v_{242} \to v_{242} = z \to v_{241} = z)$
\item $\varphi_{23} :\equiv \forall y\,\forall z\,(v_{238} = y \to y = z \to v_{238} = z) \to \forall z\,(v_{238} = v_{239} \to v_{239} = z \to v_{238} = z)$
\item $\varphi_{24} :\equiv \forall y\,\forall z\,(v_{235} = y \to y = z \to v_{235} = z) \to \forall z\,(v_{235} = v_{236} \to v_{236} = z \to v_{235} = z)$
\item $\varphi_{25} :\equiv \forall w\,(z = w \to z < y + (-x + -y) \to w < y + (-x + -y))$
\item $\varphi_{26} :\equiv z = w \to y + (-x + -y) = y + z \to y + (-x + -y) = y + w$
\item $\varphi_{27} :\equiv \forall z\,(v_{227} = v_{228} \to v_{228} = z \to v_{227} = z) \to v_{227} = v_{228} \to v_{228} = v_{229} \to v_{227} = v_{229}$
\item $\varphi_{28} :\equiv \forall z\,(v_{224} = v_{225} \to v_{225} = z \to v_{224} = z) \to v_{224} = v_{225} \to v_{225} = v_{226} \to v_{224} = v_{226}$
\item $\varphi_{29} :\equiv \forall z\,(v_{221} = v_{222} \to v_{222} = z \to v_{221} = z) \to v_{221} = v_{222} \to v_{222} = v_{223} \to v_{221} = v_{223}$
\item $\varphi_{30} :\equiv \forall z\,(v_{244} = v_{245} \to v_{245} = z \to v_{244} = z) \to v_{244} = v_{245} \to v_{245} = v_{246} \to v_{244} = v_{246}$
\item $\varphi_{31} :\equiv \forall z\,(v_{241} = v_{242} \to v_{242} = z \to v_{241} = z) \to v_{241} = v_{242} \to v_{242} = v_{243} \to v_{241} = v_{243}$
\item $\varphi_{32} :\equiv \forall z\,(v_{238} = v_{239} \to v_{239} = z \to v_{238} = z) \to v_{238} = v_{239} \to v_{239} = v_{240} \to v_{238} = v_{240}$
\item $\varphi_{33} :\equiv \forall z\,(v_{235} = v_{236} \to v_{236} = z \to v_{235} = z) \to v_{235} = v_{236} \to v_{236} = v_{237} \to v_{235} = v_{237}$
\item $\varphi_{34} :\equiv \forall y\,\forall z\,(x \in \mathbb{R} \to y \in \mathbb{R} \to z \in \mathbb{R} \to x + y + z = x + (y + z))$
\item $\varphi_{35} :\equiv \forall y\,\forall z\,(x \in \mathbb{R} \to y \in \mathbb{R} \to z \in \mathbb{R} \to x < y \to x + z < y + z)$
\item $\varphi_{36} :\equiv \forall w\,(y + (-x + -y) = w \to -y < y + (-x + -y) \to -y < w)$
\item $\varphi_{37} :\equiv \forall y\,\forall z\,(v_{230} \in \mathbb{R} \to y \in \mathbb{R} \to z \in \mathbb{R} \to v_{230} + y + z = v_{230} + (y + z))$
\item $\varphi_{38} :\equiv \forall w\,(z = w \to x + -x + -y = z + -y \to x + -x + -y = w + -y)$
\item $\varphi_{39} :\equiv x \in \mathbb{R} \to y \in \mathbb{R} \to x < y \to x + -x + -y = x + (-x + -y)$
\item $\varphi_{40} :\equiv x \in \mathbb{R} \to y \in \mathbb{R} \to x < y \to x + (-x + -y) = x + -x + -y$
\item $\varphi_{41} :\equiv \forall w\,(z = w \to y + -y + -x = z + -x \to y + -y + -x = w + -x)$
\item $\varphi_{42} :\equiv x \in \mathbb{R} \to y \in \mathbb{R} \to x < y \to y + -y + -x = y + (-y + -x)$
\item $\varphi_{43} :\equiv x \in \mathbb{R} \to y \in \mathbb{R} \to x < y \to y + (-y + -x) = y + -y + -x$
\item $\varphi_{44} :\equiv x \in \mathbb{R} \to y \in \mathbb{R} \to x < y \to y + (-x + -y) = y + (-y + -x)$
\item $\varphi_{45} :\equiv x \in \mathbb{R} \to y \in \mathbb{R} \to x < y \to y + (-x + -y) = y + -y + -x$
\item $\varphi_{46} :\equiv \forall v_{228}\,\forall v_{229}\,(x + (-x + -y) = v_{228} \to v_{228} = v_{229} \to x + (-x + -y) = v_{229})$
\item $\varphi_{47} :\equiv \forall v_{225}\,\forall v_{226}\,(x + (-x + -y) = v_{225} \to v_{225} = v_{226} \to x + (-x + -y) = v_{226})$
\item $\varphi_{48} :\equiv \forall v_{222}\,\forall v_{223}\,(x + (-x + -y) = v_{222} \to v_{222} = v_{223} \to x + (-x + -y) = v_{223})$
\item $\varphi_{49} :\equiv \forall v_{245}\,\forall v_{246}\,(y + (-x + -y) = v_{245} \to v_{245} = v_{246} \to y + (-x + -y) = v_{246})$
\item $\varphi_{50} :\equiv \forall v_{242}\,\forall v_{243}\,(y + (-x + -y) = v_{242} \to v_{242} = v_{243} \to y + (-x + -y) = v_{243})$
\item $\varphi_{51} :\equiv \forall v_{239}\,\forall v_{240}\,(y + (-x + -y) = v_{239} \to v_{239} = v_{240} \to y + (-x + -y) = v_{240})$
\item $\varphi_{52} :\equiv \forall v_{236}\,\forall v_{237}\,(y + (-x + -y) = v_{236} \to v_{236} = v_{237} \to y + (-x + -y) = v_{237})$
\item $\varphi_{53} :\equiv \forall z\,(x \in \mathbb{R} \to -x \in \mathbb{R} \to z \in \mathbb{R} \to x + -x + z = x + (-x + z))$
\item $\varphi_{54} :\equiv x + (-x + -y) < y + (-x + -y) \to -y < y + (-x + -y)$
\item $\varphi_{55} :\equiv y + (-x + -y) = -x \to -y < y + (-x + -y) \to -y < -x$
\item $\varphi_{56} :\equiv x \in \mathbb{R} \to y \in \mathbb{R} \to x < y \to x + (-x + -y) < y + (-x + -y)$
\item $\varphi_{57} :\equiv \forall z\,(v_{230} \in \mathbb{R} \to v_{231} \in \mathbb{R} \to z \in \mathbb{R} \to v_{230} + v_{231} + z = v_{230} + (v_{231} + z))$
\item $\varphi_{58} :\equiv -x + -y \in \mathbb{R} \to x < y \to x + (-x + -y) < y + (-x + -y)$
\item $\varphi_{59} :\equiv x + (-x + -y) = -y + 0 \to -y + 0 = -y \to x + (-x + -y) = -y$
\item $\varphi_{60} :\equiv y + (-x + -y) = -x + 0 \to -x + 0 = -x \to y + (-x + -y) = -x$
\item $\varphi_{61} :\equiv \forall v_{229}\,(x + (-x + -y) = -y + 0 \to -y + 0 = v_{229} \to x + (-x + -y) = v_{229})$
\item $\varphi_{62} :\equiv \forall v_{226}\,(x + (-x + -y) = 0 + -y \to 0 + -y = v_{226} \to x + (-x + -y) = v_{226})$
\item $\varphi_{63} :\equiv \forall v_{246}\,(y + (-x + -y) = -x + 0 \to -x + 0 = v_{246} \to y + (-x + -y) = v_{246})$
\item $\varphi_{64} :\equiv \forall v_{243}\,(y + (-x + -y) = 0 + -x \to 0 + -x = v_{243} \to y + (-x + -y) = v_{243})$
\item $\varphi_{65} :\equiv x + -x + -y = x + (-x + -y) \to x + (-x + -y) = x + -x + -y$
\item $\varphi_{66} :\equiv y + -y + -x = y + (-y + -x) \to y + (-y + -x) = y + -y + -x$
\item $\varphi_{67} :\equiv y + (-y + -x) = y + -y + -x \to y + (-x + -y) = y + -y + -x$
\item $\varphi_{68} :\equiv y + (-x + -y) = y + (-x + -y) \to y + (-x + -y) = y + (-y + -x)$
\item $\varphi_{69} :\equiv \forall x\,\forall y\,(x \in \mathbb{R} \to y \in \mathbb{R} \to x + y = y + x) \to \forall y\,(0 \in \mathbb{R} \to y \in \mathbb{R} \to 0 + y = y + 0)$
\item $\varphi_{70} :\equiv x \in \mathbb{R} \to -x \in \mathbb{R} \to -y \in \mathbb{R} \to x + -x + -y = x + (-x + -y)$
\item $\varphi_{71} :\equiv \forall v_{231}\,\forall v_{232}\,(y \in \mathbb{R} \to v_{231} \in \mathbb{R} \to v_{232} \in \mathbb{R} \to y + v_{231} + v_{232} = y + (v_{231} + v_{232}))$
\item $\varphi_{72} :\equiv \forall v_{232}\,(y \in \mathbb{R} \to -y \in \mathbb{R} \to v_{232} \in \mathbb{R} \to y + -y + v_{232} = y + (-y + v_{232}))$
\item $\varphi_{73} :\equiv y \in \mathbb{R} \to -y \in \mathbb{R} \to -x \in \mathbb{R} \to y + -y + -x = y + (-y + -x)$
\item $\varphi_{74} :\equiv x + (-x + -y) = 0 + -y \to 0 + -y = -y + 0 \to x + (-x + -y) = -y + 0$
\item $\varphi_{75} :\equiv x + -x = 0 \to x + -x + -y = x + -x + -y \to x + -x + -y = 0 + -y$
\item $\varphi_{76} :\equiv y + (-x + -y) = 0 + -x \to 0 + -x = -x + 0 \to y + (-x + -y) = -x + 0$
\item $\varphi_{77} :\equiv y + -y = 0 \to y + -y + -x = y + -y + -x \to y + -y + -x = 0 + -x$
\item $\varphi_{78} :\equiv (x \in \mathbb{R} \to y \in \mathbb{R} \to x < y \to x \in \mathbb{R}) \to x \in \mathbb{R} \to y \in \mathbb{R} \to x < y \to x + -x = 0$
\item $\varphi_{79} :\equiv (x \in \mathbb{R} \to y \in \mathbb{R} \to x < y \to y \in \mathbb{R}) \to x \in \mathbb{R} \to y \in \mathbb{R} \to x < y \to y + -y = 0$
\item $\varphi_{80} :\equiv \forall v_{219}\,\forall v_{220}\,(v_{219} = v_{220} \to v_{220} = v_{219}) \to \forall v_{220}\,(x + -x + -y = v_{220} \to v_{220} = x + -x + -y)$
\item $\varphi_{81} :\equiv \forall v_{233}\,\forall v_{234}\,(v_{233} = v_{234} \to v_{234} = v_{233}) \to \forall v_{234}\,(y + -y + -x = v_{234} \to v_{234} = y + -y + -x)$
\item $\varphi_{82} :\equiv (x \in \mathbb{R} \to y \in \mathbb{R} \to x < y \to y \in \mathbb{R}) \to x \in \mathbb{R} \to y \in \mathbb{R} \to x < y \to -y \in \mathbb{R}$
\item $\varphi_{83} :\equiv (x \in \mathbb{R} \to y \in \mathbb{R} \to x < y \to x \in \mathbb{R}) \to x \in \mathbb{R} \to y \in \mathbb{R} \to x < y \to -x \in \mathbb{R}$
\item $\varphi_{84} :\equiv \forall w\,(x + -x = w \to x + -x + -y = x + -x + -y \to x + -x + -y = w + -y)$
\item $\varphi_{85} :\equiv \forall w\,(y + -y = w \to y + -y + -x = y + -y + -x \to y + -y + -x = w + -x)$
\item $\varphi_{86} :\equiv \forall w\,(-x + -y = w \to y + (-x + -y) = y + (-x + -y) \to y + (-x + -y) = y + w)$
\item $\varphi_{87} :\equiv \forall y\,(0 \in \mathbb{R} \to y \in \mathbb{R} \to 0 + y = y + 0) \to 0 \in \mathbb{R} \to -y \in \mathbb{R} \to 0 + -y = -y + 0$
\item $\varphi_{88} :\equiv \forall y\,(0 \in \mathbb{R} \to y \in \mathbb{R} \to 0 + y = y + 0) \to 0 \in \mathbb{R} \to -x \in \mathbb{R} \to 0 + -x = -x + 0$
\item $\varphi_{89} :\equiv \forall w\,(x + (-x + -y) = w \to x + (-x + -y) < y + (-x + -y) \to w < y + (-x + -y))$
\item $\varphi_{90} :\equiv \forall x\,\forall y\,(x \in \mathbb{R} \to y \in \mathbb{R} \to x + y \in \mathbb{R}) \to \forall y\,(-x \in \mathbb{R} \to y \in \mathbb{R} \to -x + y \in \mathbb{R})$
\item $\varphi_{91} :\equiv \forall x\,\forall y\,(x \in \mathbb{R} \to y \in \mathbb{R} \to x + y = y + x) \to \forall y\,(-x \in \mathbb{R} \to y \in \mathbb{R} \to -x + y = y + -x)$
\item $\varphi_{92} :\equiv (x \in \mathbb{R} \to y \in \mathbb{R} \to x < y \to -y \in \mathbb{R}) \to x \in \mathbb{R} \to y \in \mathbb{R} \to x < y \to -y + 0 = -y$
\item $\varphi_{93} :\equiv (x \in \mathbb{R} \to y \in \mathbb{R} \to x < y \to -x \in \mathbb{R}) \to x \in \mathbb{R} \to y \in \mathbb{R} \to x < y \to -x + 0 = -x$
\item $\varphi_{94} :\equiv \forall v_{223}\,(x + (-x + -y) = x + -x + -y \to x + -x + -y = v_{223} \to x + (-x + -y) = v_{223})$
\item $\varphi_{95} :\equiv \forall v_{240}\,(y + (-x + -y) = y + -y + -x \to y + -y + -x = v_{240} \to y + (-x + -y) = v_{240})$
\item $\varphi_{96} :\equiv \forall v_{237}\,(y + (-x + -y) = y + (-y + -x) \to y + (-y + -x) = v_{237} \to y + (-x + -y) = v_{237})$
\item $\varphi_{97} :\equiv (x \in \mathbb{R} \to y \in \mathbb{R} \to x < y \to -y \in \mathbb{R}) \to x \in \mathbb{R} \to y \in \mathbb{R} \to x < y \to 0 + -y = -y + 0$
\item $\varphi_{98} :\equiv (x \in \mathbb{R} \to y \in \mathbb{R} \to x < y \to -x \in \mathbb{R}) \to x \in \mathbb{R} \to y \in \mathbb{R} \to x < y \to 0 + -x = -x + 0$
\item $\varphi_{99} :\equiv (x \in \mathbb{R} \to y \in \mathbb{R} \to x < y \to -y \in \mathbb{R}) \to x \in \mathbb{R} \to y \in \mathbb{R} \to x < y \to -x + -y \in \mathbb{R}$
\item $\varphi_{100} :\equiv x + (-x + -y) = x + -x + -y \to x + -x + -y = 0 + -y \to x + (-x + -y) = 0 + -y$
\item $\varphi_{101} :\equiv y + (-x + -y) = y + -y + -x \to y + -y + -x = 0 + -x \to y + (-x + -y) = 0 + -x$
\item $\varphi_{102} :\equiv \forall y\,(-x \in \mathbb{R} \to y \in \mathbb{R} \to -x + y \in \mathbb{R}) \to -x \in \mathbb{R} \to -y \in \mathbb{R} \to -x + -y \in \mathbb{R}$
\item $\varphi_{103} :\equiv (x \in \mathbb{R} \to y \in \mathbb{R} \to x < y \to -y \in \mathbb{R}) \to x \in \mathbb{R} \to y \in \mathbb{R} \to x < y \to -x + -y = -y + -x$
\item $\varphi_{104} :\equiv \forall y\,(-x \in \mathbb{R} \to y \in \mathbb{R} \to -x + y = y + -x) \to -x \in \mathbb{R} \to -y \in \mathbb{R} \to -x + -y = -y + -x$
\item $\varphi_{105} :\equiv \varphi_{35} \to \forall z\,(x \in \mathbb{R} \to y \in \mathbb{R} \to z \in \mathbb{R} \to x < y \to x + z < y + z)$
\item $\varphi_{106} :\equiv \forall v_{220}\,(x + -x + -y = v_{220} \to v_{220} = x + -x + -y) \to \varphi_{65}$
\item $\varphi_{107} :\equiv \forall v_{234}\,(y + -y + -x = v_{234} \to v_{234} = y + -y + -x) \to \varphi_{66}$
\item $\varphi_{108} :\equiv \forall z\,(x \in \mathbb{R} \to y \in \mathbb{R} \to z \in \mathbb{R} \to x < y \to x + z < y + z) \to x \in \mathbb{R} \to y \in \mathbb{R} \to \varphi_{58}$
\item $\varphi_{109} :\equiv (x + (-x + -y) = -y \to \varphi_{54}) \to x \in \mathbb{R} \to y \in \mathbb{R} \to x < y \to -y < y + (-x + -y)$
\item $\varphi_{110} :\equiv (x \in \mathbb{R} \to y \in \mathbb{R} \to x < y \to -x \in \mathbb{R}) \to (x \in \mathbb{R} \to y \in \mathbb{R} \to x < y \to -y \in \mathbb{R}) \to \varphi_{39}$
\item $\varphi_{111} :\equiv (x \in \mathbb{R} \to y \in \mathbb{R} \to x < y \to -y \in \mathbb{R}) \to (x \in \mathbb{R} \to y \in \mathbb{R} \to x < y \to -x \in \mathbb{R}) \to \varphi_{42}$
\item $\varphi_{112} :\equiv (x \in \mathbb{R} \to y \in \mathbb{R} \to x < y \to -x + -y \in \mathbb{R}) \to (x \in \mathbb{R} \to y \in \mathbb{R} \to x < y \to x < y) \to \varphi_{56}$
\item $\varphi_{113} :\equiv \forall v_{230}\,\forall v_{231}\,\forall v_{232}\,(v_{230} \in \mathbb{R} \to v_{231} \in \mathbb{R} \to \varphi_{17}) \to \varphi_{71}$
\item $\varphi_{114} :\equiv (x \in \mathbb{R} \to y \in \mathbb{R} \to x < y \to -y < y + (-x + -y)) \to \varphi_{55} \to x \in \mathbb{R} \to y \in \mathbb{R} \to x < y \to -y < -x$
\item $\varphi_{115} :\equiv (x \in \mathbb{R} \to y \in \mathbb{R} \to x < y \to -y + 0 = -y) \to \varphi_{59} \to x \in \mathbb{R} \to y \in \mathbb{R} \to x < y \to x + (-x + -y) = -y$
\item $\varphi_{116} :\equiv (x \in \mathbb{R} \to y \in \mathbb{R} \to x < y \to -x + 0 = -x) \to \varphi_{60} \to x \in \mathbb{R} \to y \in \mathbb{R} \to x < y \to y + (-x + -y) = -x$
\item $\varphi_{117} :\equiv x + -x + -y = x + -x + -y \to \varphi_{75} \to x \in \mathbb{R} \to y \in \mathbb{R} \to x < y \to x + -x + -y = 0 + -y$
\item $\varphi_{118} :\equiv y + -y + -x = y + -y + -x \to \varphi_{77} \to x \in \mathbb{R} \to y \in \mathbb{R} \to x < y \to y + -y + -x = 0 + -x$
\item $\varphi_{119} :\equiv (x \in \mathbb{R} \to y \in \mathbb{R} \to x < y \to 0 + -y = -y + 0) \to \varphi_{74} \to x \in \mathbb{R} \to y \in \mathbb{R} \to x < y \to x + (-x + -y) = -y + 0$
\item $\varphi_{120} :\equiv (x \in \mathbb{R} \to y \in \mathbb{R} \to x < y \to 0 + -x = -x + 0) \to \varphi_{76} \to x \in \mathbb{R} \to y \in \mathbb{R} \to x < y \to y + (-x + -y) = -x + 0$
\item $\varphi_{121} :\equiv (-x \in \mathbb{R} \to -y \in \mathbb{R} \to -x + -y \in \mathbb{R}) \to (x \in \mathbb{R} \to y \in \mathbb{R} \to x < y \to -x \in \mathbb{R}) \to \varphi_{99}$
\item $\varphi_{122} :\equiv (-x \in \mathbb{R} \to -y \in \mathbb{R} \to -x + -y = -y + -x) \to (x \in \mathbb{R} \to y \in \mathbb{R} \to x < y \to -x \in \mathbb{R}) \to \varphi_{103}$
\item $\varphi_{123} :\equiv y + (-x + -y) = y + (-x + -y) \to (-x + -y = -y + -x \to \varphi_{68}) \to \varphi_{44}$
\item $\varphi_{124} :\equiv (x \in \mathbb{R} \to y \in \mathbb{R} \to x < y \to x + -x + -y = 0 + -y) \to \varphi_{100} \to x \in \mathbb{R} \to y \in \mathbb{R} \to x < y \to x + (-x + -y) = 0 + -y$
\item $\varphi_{125} :\equiv (x \in \mathbb{R} \to y \in \mathbb{R} \to x < y \to y + -y + -x = 0 + -x) \to \varphi_{101} \to x \in \mathbb{R} \to y \in \mathbb{R} \to x < y \to y + (-x + -y) = 0 + -x$
\item $\varphi_{126} :\equiv (x \in \mathbb{R} \to y \in \mathbb{R} \to x < y \to x + (-x + -y) = -y + 0) \to \varphi_{115}$
\item $\varphi_{127} :\equiv (x \in \mathbb{R} \to y \in \mathbb{R} \to x < y \to y + (-x + -y) = -x + 0) \to \varphi_{116}$
\item $\varphi_{128} :\equiv (x \in \mathbb{R} \to y \in \mathbb{R} \to x < y \to x \in \mathbb{R}) \to (x \in \mathbb{R} \to y \in \mathbb{R} \to x < y \to y \in \mathbb{R}) \to \varphi_{112}$
\item $\varphi_{129} :\equiv (x \in \mathbb{R} \to y \in \mathbb{R} \to x < y \to x + (-x + -y) = 0 + -y) \to \varphi_{119}$
\item $\varphi_{130} :\equiv (x \in \mathbb{R} \to y \in \mathbb{R} \to x < y \to y + (-x + -y) = 0 + -x) \to \varphi_{120}$
\item $\varphi_{131} :\equiv (x \in \mathbb{R} \to y \in \mathbb{R} \to x < y \to x + (-x + -y) = -y) \to \varphi_{56} \to \varphi_{109}$
\end{abbrevlist}
\end{fullproof}

% 自動生成 (prover/build.js): NegLe  ノード数 48
\begin{fullproof}{[F23]}{$\forall x\,\forall y\,(x \in \mathbb{R} \to y \in \mathbb{R} \to x \leq y \to -y \leq -x)$}{48}\phantomsection\label{tree:NegLe}
\noindent$\Pi_{1}$:
\begin{ptree}
\AXC{$x = x$\lbl{A5}}
\AXC{$x = y \to x = x \to y = x$\lbl{A6}}
\AXC{$(x = y \to x = x \to y = x) \to x = x \to x = y \to y = x$\lbl{A1}}
\BIC{$x = x \to x = y \to y = x$}
\BIC{$x = y \to y = x$}
\UIC{$\forall y\,(x = y \to y = x)$}
\UIC{$\forall x\,\forall y\,(x = y \to y = x)$}
\end{ptree}
\noindent$\Pi_{2}$:
\begin{ptree}
\AXC{$\stackrel{\Pi_{1}}{\vdots}$}\noLine
\UIC{$\forall x\,\forall y\,(x = y \to y = x)$}
\AXC{$\forall x\,\forall y\,(x = y \to y = x) \to \forall y\,(-x = y \to y = -x)$\lbl{A2}}
\BIC{$\forall y\,(-x = y \to y = -x)$}
\end{ptree}
\noindent$\Pi_{3}$:
\begin{ptree}
\AXC{$x = x$\lbl{A5}}
\UIC{$\forall x\,(x = x)$}
\AXC{$\forall x\,(x = x) \to -x = -x$\lbl{A2}}
\BIC{$-x = -x$}
\AXC{$x \in \mathbb{R} \to y \in \mathbb{R} \to x \leq y \to x = y \to x = y$\lbl{A1}}
\AXC{$\varphi_{9}$\lbl{A1}}
\BIC{$-x = -x \to \varphi_{4}$}
\BIC{$\varphi_{4}$}
\end{ptree}
\noindent$\Pi_{4}$:
\begin{ptree}
\AXC{$z = w \to -x = -z \to -x = -w$\lbl{A6}}
\UIC{$\forall w\,(z = w \to -x = -z \to -x = -w)$}
\UIC{$\forall z\,\forall w\,(z = w \to -x = -z \to -x = -w)$}
\AXC{$\varphi_{2}$\lbl{A2}}
\BIC{$\forall w\,(x = w \to -x = -x \to -x = -w)$}
\AXC{$\varphi_{1}$\lbl{A2}}
\BIC{$x = y \to -x = -x \to -x = -y$}
\AXC{$\stackrel{\Pi_{3}}{\vdots}$}\noLine
\UIC{$\varphi_{4}$}
\BIC{$x \in \mathbb{R} \to y \in \mathbb{R} \to x \leq y \to x = y \to -x = -y$}
\end{ptree}
\noindent$\Pi_{5}$:
\begin{ptree}
\AXC{$\stackrel{\Pi_{4}}{\vdots}$}\noLine
\UIC{$x \in \mathbb{R} \to y \in \mathbb{R} \to x \leq y \to x = y \to -x = -y$}
\AXC{$(x \in \mathbb{R} \to y \in \mathbb{R} \to x \leq y \to x = y \to -x = -y) \to \varphi_{3}$\lbl{A1}}
\BIC{$\varphi_{3}$}
\end{ptree}
\noindent$\Pi_{6}$:
\begin{ptree}
\AXC{$\forall x\,\forall y\,(x \in \mathbb{R} \to y \in \mathbb{R} \to x < y \to -y < -x)$\lbl{[F22]}}
\AXC{$\varphi_{8}$\lbl{A2}}
\BIC{$\forall y\,(x \in \mathbb{R} \to y \in \mathbb{R} \to x < y \to -y < -x)$}
\AXC{$\varphi_{7}$\lbl{A2}}
\BIC{$x \in \mathbb{R} \to y \in \mathbb{R} \to x < y \to -y < -x$}
\AXC{$\varphi_{11}$\lbl{A1}}
\BIC{$(x \in \mathbb{R} \to y \in \mathbb{R} \to x \leq y \to x \in \mathbb{R}) \to \varphi_{6}$}
\end{ptree}
\noindent$\Pi_{7}$:
\begin{ptree}
\AXC{$x \in \mathbb{R} \to y \in \mathbb{R} \to x \leq y \to y \in \mathbb{R}$\lbl{A1}}
\AXC{$x \in \mathbb{R} \to y \in \mathbb{R} \to x \leq y \to x \in \mathbb{R}$\lbl{A1}}
\AXC{$\stackrel{\Pi_{6}}{\vdots}$}\noLine
\UIC{$(x \in \mathbb{R} \to y \in \mathbb{R} \to x \leq y \to x \in \mathbb{R}) \to \varphi_{6}$}
\BIC{$\varphi_{6}$}
\BIC{$x \in \mathbb{R} \to y \in \mathbb{R} \to x \leq y \to x < y \to -y < -x$}
\end{ptree}
\noindent$\Pi_{8}$:
\begin{ptree}
\AXC{$\stackrel{\Pi_{7}}{\vdots}$}\noLine
\UIC{$x \in \mathbb{R} \to y \in \mathbb{R} \to x \leq y \to x < y \to -y < -x$}
\AXC{$x \in \mathbb{R} \to y \in \mathbb{R} \to x \leq y \to x \leq y$\lbl{A1}}
\AXC{$\varphi_{10}$\lbl{A1}}
\BIC{$(x \in \mathbb{R} \to y \in \mathbb{R} \to x \leq y \to x < y \to -y < -x) \to \varphi_{5}$}
\BIC{$\varphi_{5}$}
\end{ptree}
\noindent 主木:
\begin{ptree}
\AXC{$\stackrel{\Pi_{2}}{\vdots}$}\noLine
\UIC{$\forall y\,(-x = y \to y = -x)$}
\AXC{$\forall y\,(-x = y \to y = -x) \to -x = -y \to -y = -x$\lbl{A2}}
\BIC{$-x = -y \to -y = -x$}
\AXC{$\stackrel{\Pi_{5}}{\vdots}$}\noLine
\UIC{$\varphi_{3}$}
\BIC{$x \in \mathbb{R} \to y \in \mathbb{R} \to x \leq y \to x = y \to -y = -x$}
\AXC{$\stackrel{\Pi_{8}}{\vdots}$}\noLine
\UIC{$\varphi_{5}$}
\BIC{$x \in \mathbb{R} \to y \in \mathbb{R} \to x \leq y \to -y \leq -x$}
\UIC{$\forall y\,(x \in \mathbb{R} \to y \in \mathbb{R} \to x \leq y \to -y \leq -x)$}
\UIC{$\forall x\,\forall y\,(x \in \mathbb{R} \to y \in \mathbb{R} \to x \leq y \to -y \leq -x)$}
\end{ptree}
\begin{abbrevlist}
\item $\varphi_{1} :\equiv \forall w\,(x = w \to -x = -x \to -x = -w) \to x = y \to -x = -x \to -x = -y$
\item $\varphi_{2} :\equiv \forall z\,\forall w\,(z = w \to -x = -z \to -x = -w) \to \forall w\,(x = w \to -x = -x \to -x = -w)$
\item $\varphi_{3} :\equiv (-x = -y \to -y = -x) \to x \in \mathbb{R} \to y \in \mathbb{R} \to x \leq y \to x = y \to -y = -x$
\item $\varphi_{4} :\equiv (x = y \to -x = -x \to -x = -y) \to x \in \mathbb{R} \to y \in \mathbb{R} \to x \leq y \to x = y \to -x = -y$
\item $\varphi_{5} :\equiv (x \in \mathbb{R} \to y \in \mathbb{R} \to x \leq y \to x = y \to -y = -x) \to x \in \mathbb{R} \to y \in \mathbb{R} \to x \leq y \to -y \leq -x$
\item $\varphi_{6} :\equiv (x \in \mathbb{R} \to y \in \mathbb{R} \to x \leq y \to y \in \mathbb{R}) \to x \in \mathbb{R} \to y \in \mathbb{R} \to x \leq y \to x < y \to -y < -x$
\item $\varphi_{7} :\equiv \forall y\,(x \in \mathbb{R} \to y \in \mathbb{R} \to x < y \to -y < -x) \to x \in \mathbb{R} \to y \in \mathbb{R} \to x < y \to -y < -x$
\item $\varphi_{8} :\equiv \forall x\,\forall y\,(x \in \mathbb{R} \to y \in \mathbb{R} \to x < y \to -y < -x) \to \forall y\,(x \in \mathbb{R} \to y \in \mathbb{R} \to x < y \to -y < -x)$
\item $\varphi_{9} :\equiv (x \in \mathbb{R} \to y \in \mathbb{R} \to x \leq y \to x = y \to x = y) \to -x = -x \to \varphi_{4}$
\item $\varphi_{10} :\equiv (x \in \mathbb{R} \to y \in \mathbb{R} \to x \leq y \to x \leq y) \to (x \in \mathbb{R} \to y \in \mathbb{R} \to x \leq y \to x < y \to -y < -x) \to \varphi_{5}$
\item $\varphi_{11} :\equiv (x \in \mathbb{R} \to y \in \mathbb{R} \to x < y \to -y < -x) \to (x \in \mathbb{R} \to y \in \mathbb{R} \to x \leq y \to x \in \mathbb{R}) \to \varphi_{6}$
\end{abbrevlist}
\end{fullproof}

% 自動生成 (prover/build.js): SubPos  ノード数 85
\begin{fullproof}{[F24]}{$\forall x\,\forall y\,(x \in \mathbb{R} \to y \in \mathbb{R} \to x < y \to 0 < y - x)$}{85}\phantomsection\label{tree:SubPos}
\noindent$\Pi_{1}$:
\begin{ptree}
\AXC{$\forall x \in \mathbb{R}\,(-x \in \mathbb{R})$\lbl{(R1)}}
\AXC{$\forall x \in \mathbb{R}\,(-x \in \mathbb{R}) \to x \in \mathbb{R} \to -x \in \mathbb{R}$\lbl{A2}}
\BIC{$x \in \mathbb{R} \to -x \in \mathbb{R}$}
\AXC{$(x \in \mathbb{R} \to -x \in \mathbb{R}) \to \varphi_{12}$\lbl{A1}}
\BIC{$\varphi_{12}$}
\end{ptree}
\noindent$\Pi_{2}$:
\begin{ptree}
\AXC{$x \in \mathbb{R} \to y \in \mathbb{R} \to x < y \to x \in \mathbb{R}$\lbl{A1}}
\AXC{$\forall x\,(\varphi_{5})$\lbl{(R14)}}
\AXC{$\forall x\,(\varphi_{5}) \to \varphi_{5}$\lbl{A2}}
\BIC{$\varphi_{5}$}
\AXC{$\varphi_{18}$\lbl{A2}}
\BIC{$\forall z\,(x \in \mathbb{R} \to y \in \mathbb{R} \to z \in \mathbb{R} \to x < y \to x + z < y + z)$}
\AXC{$\varphi_{19}$\lbl{A2}}
\BIC{$x \in \mathbb{R} \to y \in \mathbb{R} \to -x \in \mathbb{R} \to x < y \to x + -x < y + -x$}
\AXC{$\varphi_{23}$\lbl{A1}}
\BIC{$(x \in \mathbb{R} \to y \in \mathbb{R} \to x < y \to x \in \mathbb{R}) \to \varphi_{21}$}
\BIC{$\varphi_{21}$}
\end{ptree}
\noindent$\Pi_{3}$:
\begin{ptree}
\AXC{$x \in \mathbb{R} \to y \in \mathbb{R} \to x < y \to x \in \mathbb{R}$\lbl{A1}}
\AXC{$\stackrel{\Pi_{1}}{\vdots}$}\noLine
\UIC{$\varphi_{12}$}
\BIC{$x \in \mathbb{R} \to y \in \mathbb{R} \to x < y \to -x \in \mathbb{R}$}
\AXC{$x \in \mathbb{R} \to y \in \mathbb{R} \to x < y \to y \in \mathbb{R}$\lbl{A1}}
\AXC{$\stackrel{\Pi_{2}}{\vdots}$}\noLine
\UIC{$\varphi_{21}$}
\BIC{$(x \in \mathbb{R} \to y \in \mathbb{R} \to x < y \to -x \in \mathbb{R}) \to \varphi_{14}$}
\BIC{$\varphi_{14}$}
\end{ptree}
\noindent$\Pi_{4}$:
\begin{ptree}
\AXC{$\forall x \in \mathbb{R}\,(x + -x = 0)$\lbl{(R5)}}
\AXC{$\forall x \in \mathbb{R}\,(x + -x = 0) \to x \in \mathbb{R} \to x + -x = 0$\lbl{A2}}
\BIC{$x \in \mathbb{R} \to x + -x = 0$}
\AXC{$(x \in \mathbb{R} \to x + -x = 0) \to \varphi_{10}$\lbl{A1}}
\BIC{$\varphi_{10}$}
\end{ptree}
\noindent$\Pi_{5}$:
\begin{ptree}
\AXC{$x \in \mathbb{R} \to y \in \mathbb{R} \to x < y \to x < y$\lbl{A1}}
\AXC{$\stackrel{\Pi_{3}}{\vdots}$}\noLine
\UIC{$\varphi_{14}$}
\BIC{$x \in \mathbb{R} \to y \in \mathbb{R} \to x < y \to x + -x < y + -x$}
\AXC{$x \in \mathbb{R} \to y \in \mathbb{R} \to x < y \to x \in \mathbb{R}$\lbl{A1}}
\AXC{$\stackrel{\Pi_{4}}{\vdots}$}\noLine
\UIC{$\varphi_{10}$}
\BIC{$x \in \mathbb{R} \to y \in \mathbb{R} \to x < y \to x + -x = 0$}
\AXC{$\varphi_{22}$\lbl{A1}}
\BIC{$(x \in \mathbb{R} \to y \in \mathbb{R} \to x < y \to x + -x < y + -x) \to \varphi_{16}$}
\BIC{$\varphi_{16}$}
\end{ptree}
\noindent$\Pi_{6}$:
\begin{ptree}
\AXC{$x = x$\lbl{A5}}
\AXC{$x = y \to x = x \to y = x$\lbl{A6}}
\AXC{$(x = y \to x = x \to y = x) \to x = x \to x = y \to y = x$\lbl{A1}}
\BIC{$x = x \to x = y \to y = x$}
\BIC{$x = y \to y = x$}
\UIC{$\forall y\,(x = y \to y = x)$}
\UIC{$\forall x\,\forall y\,(x = y \to y = x)$}
\end{ptree}
\noindent$\Pi_{7}$:
\begin{ptree}
\AXC{$\stackrel{\Pi_{6}}{\vdots}$}\noLine
\UIC{$\forall x\,\forall y\,(x = y \to y = x)$}
\AXC{$\forall x\,\forall y\,(x = y \to y = x) \to \forall y\,(v_{249} = y \to y = v_{249})$\lbl{A2}}
\BIC{$\forall y\,(v_{249} = y \to y = v_{249})$}
\AXC{$\varphi_{1}$\lbl{A2}}
\BIC{$v_{249} = v_{250} \to v_{250} = v_{249}$}
\UIC{$\forall v_{250}\,(v_{249} = v_{250} \to v_{250} = v_{249})$}
\UIC{$\forall v_{249}\,\forall v_{250}\,(v_{249} = v_{250} \to v_{250} = v_{249})$}
\AXC{$\varphi_{4}$\lbl{A2}}
\BIC{$\forall v_{250}\,(y - x = v_{250} \to v_{250} = y - x)$}
\AXC{$\varphi_{6}$\lbl{A2}}
\BIC{$y - x = y + -x \to y + -x = y - x$}
\end{ptree}
\noindent$\Pi_{8}$:
\begin{ptree}
\AXC{$\forall x\,\forall y\,(x - y = x + -y)$\lbl{Def$-$}}
\AXC{$\varphi_{2}$\lbl{A2}}
\BIC{$\forall y\,(v_{247} - y = v_{247} + -y)$}
\AXC{$\varphi_{3}$\lbl{A2}}
\BIC{$v_{247} - v_{248} = v_{247} + -v_{248}$}
\UIC{$\forall v_{248}\,(v_{247} - v_{248} = v_{247} + -v_{248})$}
\UIC{$\forall v_{247}\,\forall v_{248}\,(v_{247} - v_{248} = v_{247} + -v_{248})$}
\AXC{$\varphi_{8}$\lbl{A2}}
\BIC{$\forall v_{248}\,(y - v_{248} = y + -v_{248})$}
\AXC{$\forall v_{248}\,(y - v_{248} = y + -v_{248}) \to y - x = y + -x$\lbl{A2}}
\BIC{$y - x = y + -x$}
\end{ptree}
\noindent$\Pi_{9}$:
\begin{ptree}
\AXC{$\stackrel{\Pi_{7}}{\vdots}$}\noLine
\UIC{$y - x = y + -x \to y + -x = y - x$}
\AXC{$\stackrel{\Pi_{8}}{\vdots}$}\noLine
\UIC{$y - x = y + -x$}
\AXC{$\varphi_{9}$\lbl{A1}}
\BIC{$(y - x = y + -x \to y + -x = y - x) \to y + -x = y - x$}
\BIC{$y + -x = y - x$}
\AXC{$\varphi_{20}$\lbl{A1}}
\BIC{$(x \in \mathbb{R} \to y \in \mathbb{R} \to x < y \to 0 < y + -x) \to \varphi_{13}$}
\end{ptree}
\noindent$\Pi_{10}$:
\begin{ptree}
\AXC{$z = w \to z < y + -x \to w < y + -x$\lbl{A6}}
\UIC{$\forall w\,(z = w \to z < y + -x \to w < y + -x)$}
\UIC{$\forall z\,\forall w\,(z = w \to z < y + -x \to w < y + -x)$}
\AXC{$\varphi_{15}$\lbl{A2}}
\BIC{$\forall w\,(x + -x = w \to x + -x < y + -x \to w < y + -x)$}
\AXC{$\varphi_{17}$\lbl{A2}}
\BIC{$x + -x = 0 \to x + -x < y + -x \to 0 < y + -x$}
\AXC{$\stackrel{\Pi_{5}}{\vdots}$}\noLine
\UIC{$\varphi_{16}$}
\BIC{$x \in \mathbb{R} \to y \in \mathbb{R} \to x < y \to 0 < y + -x$}
\AXC{$\stackrel{\Pi_{9}}{\vdots}$}\noLine
\UIC{$(x \in \mathbb{R} \to y \in \mathbb{R} \to x < y \to 0 < y + -x) \to \varphi_{13}$}
\BIC{$\varphi_{13}$}
\end{ptree}
\noindent 主木:
\begin{ptree}
\AXC{$z = w \to 0 < z \to 0 < w$\lbl{A6}}
\UIC{$\forall w\,(z = w \to 0 < z \to 0 < w)$}
\UIC{$\forall z\,\forall w\,(z = w \to 0 < z \to 0 < w)$}
\AXC{$\varphi_{7}$\lbl{A2}}
\BIC{$\forall w\,(y + -x = w \to 0 < y + -x \to 0 < w)$}
\AXC{$\varphi_{11}$\lbl{A2}}
\BIC{$y + -x = y - x \to 0 < y + -x \to 0 < y - x$}
\AXC{$\stackrel{\Pi_{10}}{\vdots}$}\noLine
\UIC{$\varphi_{13}$}
\BIC{$x \in \mathbb{R} \to y \in \mathbb{R} \to x < y \to 0 < y - x$}
\UIC{$\forall y\,(x \in \mathbb{R} \to y \in \mathbb{R} \to x < y \to 0 < y - x)$}
\UIC{$\forall x\,\forall y\,(x \in \mathbb{R} \to y \in \mathbb{R} \to x < y \to 0 < y - x)$}
\end{ptree}
\begin{abbrevlist}
\item $\varphi_{1} :\equiv \forall y\,(v_{249} = y \to y = v_{249}) \to v_{249} = v_{250} \to v_{250} = v_{249}$
\item $\varphi_{2} :\equiv \forall x\,\forall y\,(x - y = x + -y) \to \forall y\,(v_{247} - y = v_{247} + -y)$
\item $\varphi_{3} :\equiv \forall y\,(v_{247} - y = v_{247} + -y) \to v_{247} - v_{248} = v_{247} + -v_{248}$
\item $\varphi_{4} :\equiv \forall v_{249}\,\forall v_{250}\,(v_{249} = v_{250} \to v_{250} = v_{249}) \to \forall v_{250}\,(y - x = v_{250} \to v_{250} = y - x)$
\item $\varphi_{5} :\equiv \forall y\,\forall z\,(x \in \mathbb{R} \to y \in \mathbb{R} \to z \in \mathbb{R} \to x < y \to x + z < y + z)$
\item $\varphi_{6} :\equiv \forall v_{250}\,(y - x = v_{250} \to v_{250} = y - x) \to y - x = y + -x \to y + -x = y - x$
\item $\varphi_{7} :\equiv \forall z\,\forall w\,(z = w \to 0 < z \to 0 < w) \to \forall w\,(y + -x = w \to 0 < y + -x \to 0 < w)$
\item $\varphi_{8} :\equiv \forall v_{247}\,\forall v_{248}\,(v_{247} - v_{248} = v_{247} + -v_{248}) \to \forall v_{248}\,(y - v_{248} = y + -v_{248})$
\item $\varphi_{9} :\equiv y - x = y + -x \to (y - x = y + -x \to y + -x = y - x) \to y + -x = y - x$
\item $\varphi_{10} :\equiv (x \in \mathbb{R} \to y \in \mathbb{R} \to x < y \to x \in \mathbb{R}) \to x \in \mathbb{R} \to y \in \mathbb{R} \to x < y \to x + -x = 0$
\item $\varphi_{11} :\equiv \forall w\,(y + -x = w \to 0 < y + -x \to 0 < w) \to y + -x = y - x \to 0 < y + -x \to 0 < y - x$
\item $\varphi_{12} :\equiv (x \in \mathbb{R} \to y \in \mathbb{R} \to x < y \to x \in \mathbb{R}) \to x \in \mathbb{R} \to y \in \mathbb{R} \to x < y \to -x \in \mathbb{R}$
\item $\varphi_{13} :\equiv (y + -x = y - x \to 0 < y + -x \to 0 < y - x) \to x \in \mathbb{R} \to y \in \mathbb{R} \to x < y \to 0 < y - x$
\item $\varphi_{14} :\equiv (x \in \mathbb{R} \to y \in \mathbb{R} \to x < y \to x < y) \to x \in \mathbb{R} \to y \in \mathbb{R} \to x < y \to x + -x < y + -x$
\item $\varphi_{15} :\equiv \forall z\,\forall w\,(z = w \to z < y + -x \to w < y + -x) \to \forall w\,(x + -x = w \to x + -x < y + -x \to w < y + -x)$
\item $\varphi_{16} :\equiv (x + -x = 0 \to x + -x < y + -x \to 0 < y + -x) \to x \in \mathbb{R} \to y \in \mathbb{R} \to x < y \to 0 < y + -x$
\item $\varphi_{17} :\equiv \forall w\,(x + -x = w \to x + -x < y + -x \to w < y + -x) \to x + -x = 0 \to x + -x < y + -x \to 0 < y + -x$
\item $\varphi_{18} :\equiv \varphi_{5} \to \forall z\,(x \in \mathbb{R} \to y \in \mathbb{R} \to z \in \mathbb{R} \to x < y \to x + z < y + z)$
\item $\varphi_{19} :\equiv \forall z\,(x \in \mathbb{R} \to y \in \mathbb{R} \to z \in \mathbb{R} \to x < y \to x + z < y + z) \to x \in \mathbb{R} \to y \in \mathbb{R} \to -x \in \mathbb{R} \to x < y \to x + -x < y + -x$
\item $\varphi_{20} :\equiv y + -x = y - x \to (x \in \mathbb{R} \to y \in \mathbb{R} \to x < y \to 0 < y + -x) \to \varphi_{13}$
\item $\varphi_{21} :\equiv (x \in \mathbb{R} \to y \in \mathbb{R} \to x < y \to y \in \mathbb{R}) \to (x \in \mathbb{R} \to y \in \mathbb{R} \to x < y \to -x \in \mathbb{R}) \to \varphi_{14}$
\item $\varphi_{22} :\equiv (x \in \mathbb{R} \to y \in \mathbb{R} \to x < y \to x + -x = 0) \to (x \in \mathbb{R} \to y \in \mathbb{R} \to x < y \to x + -x < y + -x) \to \varphi_{16}$
\item $\varphi_{23} :\equiv (x \in \mathbb{R} \to y \in \mathbb{R} \to -x \in \mathbb{R} \to x < y \to x + -x < y + -x) \to (x \in \mathbb{R} \to y \in \mathbb{R} \to x < y \to x \in \mathbb{R}) \to \varphi_{21}$
\end{abbrevlist}
\end{fullproof}

% 自動生成 (prover/build.js): PosSub  ノード数 107
\begin{fullproof}{[F25]}{$\forall x\,\forall y\,(x \in \mathbb{R} \to y \in \mathbb{R} \to 0 < y - x \to x < y)$}{107}\phantomsection\label{tree:PosSub}
\noindent$\Pi_{1}$:
\begin{ptree}
\AXC{$\forall x\,\forall y\,(x \in \mathbb{R} \to y \in \mathbb{R} \to x - y \in \mathbb{R})$\lbl{[F1]}}
\AXC{$\varphi_{19}$\lbl{A2}}
\BIC{$\forall y\,(v_{251} \in \mathbb{R} \to y \in \mathbb{R} \to v_{251} - y \in \mathbb{R})$}
\AXC{$\varphi_{22}$\lbl{A2}}
\BIC{$v_{251} \in \mathbb{R} \to v_{252} \in \mathbb{R} \to v_{251} - v_{252} \in \mathbb{R}$}
\UIC{$\forall v_{252}\,(v_{251} \in \mathbb{R} \to v_{252} \in \mathbb{R} \to v_{251} - v_{252} \in \mathbb{R})$}
\UIC{$\varphi_{3}$}
\AXC{$\varphi_{3} \to \forall v_{252}\,(y \in \mathbb{R} \to v_{252} \in \mathbb{R} \to y - v_{252} \in \mathbb{R})$\lbl{A2}}
\BIC{$\forall v_{252}\,(y \in \mathbb{R} \to v_{252} \in \mathbb{R} \to y - v_{252} \in \mathbb{R})$}
\end{ptree}
\noindent$\Pi_{2}$:
\begin{ptree}
\AXC{$x \in \mathbb{R} \to y \in \mathbb{R} \to 0 < y - x \to y \in \mathbb{R}$\lbl{A1}}
\AXC{$\stackrel{\Pi_{1}}{\vdots}$}\noLine
\UIC{$\forall v_{252}\,(y \in \mathbb{R} \to v_{252} \in \mathbb{R} \to y - v_{252} \in \mathbb{R})$}
\AXC{$\varphi_{17}$\lbl{A2}}
\BIC{$y \in \mathbb{R} \to x \in \mathbb{R} \to y - x \in \mathbb{R}$}
\AXC{$\varphi_{33}$\lbl{A1}}
\BIC{$(x \in \mathbb{R} \to y \in \mathbb{R} \to 0 < y - x \to y \in \mathbb{R}) \to \varphi_{26}$}
\BIC{$\varphi_{26}$}
\end{ptree}
\noindent$\Pi_{3}$:
\begin{ptree}
\AXC{$0 \in \mathbb{R}$\lbl{(R1)}}
\AXC{$\varphi_{6}$\lbl{(R14)}}
\AXC{$\varphi_{6} \to \varphi_{4}$\lbl{A2}}
\BIC{$\varphi_{4}$}
\AXC{$\varphi_{4} \to \varphi_{10}$\lbl{A2}}
\BIC{$\varphi_{10}$}
\AXC{$\varphi_{10} \to \varphi_{9}$\lbl{A2}}
\BIC{$\varphi_{9}$}
\AXC{$\varphi_{9} \to 0 \in \mathbb{R} \to \varphi_{36}$\lbl{A1}}
\BIC{$0 \in \mathbb{R} \to \varphi_{36}$}
\BIC{$\varphi_{36}$}
\end{ptree}
\noindent$\Pi_{4}$:
\begin{ptree}
\AXC{$x \in \mathbb{R} \to y \in \mathbb{R} \to 0 < y - x \to x \in \mathbb{R}$\lbl{A1}}
\AXC{$x \in \mathbb{R} \to y \in \mathbb{R} \to 0 < y - x \to x \in \mathbb{R}$\lbl{A1}}
\AXC{$\stackrel{\Pi_{2}}{\vdots}$}\noLine
\UIC{$\varphi_{26}$}
\BIC{$x \in \mathbb{R} \to y \in \mathbb{R} \to 0 < y - x \to y - x \in \mathbb{R}$}
\AXC{$\stackrel{\Pi_{3}}{\vdots}$}\noLine
\UIC{$\varphi_{36}$}
\BIC{$(x \in \mathbb{R} \to y \in \mathbb{R} \to 0 < y - x \to x \in \mathbb{R}) \to \varphi_{29}$}
\BIC{$\varphi_{29}$}
\end{ptree}
\noindent$\Pi_{5}$:
\begin{ptree}
\AXC{$y = z \to x = y \to x = z$\lbl{A6}}
\AXC{$(y = z \to x = y \to x = z) \to x = y \to y = z \to x = z$\lbl{A1}}
\BIC{$x = y \to y = z \to x = z$}
\UIC{$\forall z\,(x = y \to y = z \to x = z)$}
\UIC{$\forall y\,\forall z\,(x = y \to y = z \to x = z)$}
\UIC{$\forall x\,\forall y\,\forall z\,(x = y \to y = z \to x = z)$}
\AXC{$\varphi_{1}$\lbl{A2}}
\BIC{$\forall y\,\forall z\,(0 + x = y \to y = z \to 0 + x = z)$}
\AXC{$\varphi_{5}$\lbl{A2}}
\BIC{$\forall z\,(0 + x = x + 0 \to x + 0 = z \to 0 + x = z)$}
\AXC{$\varphi_{8}$\lbl{A2}}
\BIC{$0 + x = x + 0 \to x + 0 = x \to 0 + x = x$}
\end{ptree}
\noindent$\Pi_{6}$:
\begin{ptree}
\AXC{$\forall x \in \mathbb{R}\,(x + 0 = x)$\lbl{(R4)}}
\AXC{$\forall x \in \mathbb{R}\,(x + 0 = x) \to x \in \mathbb{R} \to x + 0 = x$\lbl{A2}}
\BIC{$x \in \mathbb{R} \to x + 0 = x$}
\AXC{$(x \in \mathbb{R} \to x + 0 = x) \to \varphi_{20}$\lbl{A1}}
\BIC{$\varphi_{20}$}
\end{ptree}
\noindent$\Pi_{7}$:
\begin{ptree}
\AXC{$\forall x\,\forall y\,(x \in \mathbb{R} \to y \in \mathbb{R} \to x + y = y + x)$\lbl{(R3)}}
\AXC{$\varphi_{14}$\lbl{A2}}
\BIC{$\forall y\,(0 \in \mathbb{R} \to y \in \mathbb{R} \to 0 + y = y + 0)$}
\AXC{$\varphi_{11}$\lbl{A2}}
\BIC{$0 \in \mathbb{R} \to x \in \mathbb{R} \to 0 + x = x + 0$}
\AXC{$(0 \in \mathbb{R} \to x \in \mathbb{R} \to 0 + x = x + 0) \to 0 \in \mathbb{R} \to \varphi_{23}$\lbl{A1}}
\BIC{$0 \in \mathbb{R} \to \varphi_{23}$}
\end{ptree}
\noindent$\Pi_{8}$:
\begin{ptree}
\AXC{$x \in \mathbb{R} \to y \in \mathbb{R} \to 0 < y - x \to x \in \mathbb{R}$\lbl{A1}}
\AXC{$0 \in \mathbb{R}$\lbl{(R1)}}
\AXC{$\stackrel{\Pi_{7}}{\vdots}$}\noLine
\UIC{$0 \in \mathbb{R} \to \varphi_{23}$}
\BIC{$\varphi_{23}$}
\BIC{$x \in \mathbb{R} \to y \in \mathbb{R} \to 0 < y - x \to 0 + x = x + 0$}
\AXC{$\varphi_{31}$\lbl{A1}}
\BIC{$(x \in \mathbb{R} \to y \in \mathbb{R} \to 0 < y - x \to x + 0 = x) \to \varphi_{12}$}
\end{ptree}
\noindent$\Pi_{9}$:
\begin{ptree}
\AXC{$x \in \mathbb{R} \to y \in \mathbb{R} \to 0 < y - x \to x \in \mathbb{R}$\lbl{A1}}
\AXC{$\stackrel{\Pi_{6}}{\vdots}$}\noLine
\UIC{$\varphi_{20}$}
\BIC{$x \in \mathbb{R} \to y \in \mathbb{R} \to 0 < y - x \to x + 0 = x$}
\AXC{$\stackrel{\Pi_{8}}{\vdots}$}\noLine
\UIC{$(x \in \mathbb{R} \to y \in \mathbb{R} \to 0 < y - x \to x + 0 = x) \to \varphi_{12}$}
\BIC{$\varphi_{12}$}
\end{ptree}
\noindent$\Pi_{10}$:
\begin{ptree}
\AXC{$x \in \mathbb{R} \to y \in \mathbb{R} \to 0 < y - x \to 0 < y - x$\lbl{A1}}
\AXC{$\stackrel{\Pi_{4}}{\vdots}$}\noLine
\UIC{$\varphi_{29}$}
\BIC{$x \in \mathbb{R} \to y \in \mathbb{R} \to 0 < y - x \to 0 + x < y - x + x$}
\AXC{$\stackrel{\Pi_{5}}{\vdots}$}\noLine
\UIC{$0 + x = x + 0 \to x + 0 = x \to 0 + x = x$}
\AXC{$\stackrel{\Pi_{9}}{\vdots}$}\noLine
\UIC{$\varphi_{12}$}
\BIC{$x \in \mathbb{R} \to y \in \mathbb{R} \to 0 < y - x \to 0 + x = x$}
\AXC{$\varphi_{35}$\lbl{A1}}
\BIC{$(x \in \mathbb{R} \to y \in \mathbb{R} \to 0 < y - x \to 0 + x < y - x + x) \to \varphi_{27}$}
\BIC{$\varphi_{27}$}
\end{ptree}
\noindent$\Pi_{11}$:
\begin{ptree}
\AXC{$z = w \to z < y - x + x \to w < y - x + x$\lbl{A6}}
\UIC{$\forall w\,(z = w \to z < y - x + x \to w < y - x + x)$}
\UIC{$\forall z\,\forall w\,(z = w \to z < y - x + x \to w < y - x + x)$}
\AXC{$\varphi_{21}$\lbl{A2}}
\BIC{$\forall w\,(0 + x = w \to 0 + x < y - x + x \to w < y - x + x)$}
\AXC{$\varphi_{25}$\lbl{A2}}
\BIC{$0 + x = x \to 0 + x < y - x + x \to x < y - x + x$}
\AXC{$\stackrel{\Pi_{10}}{\vdots}$}\noLine
\UIC{$\varphi_{27}$}
\BIC{$x \in \mathbb{R} \to y \in \mathbb{R} \to 0 < y - x \to x < y - x + x$}
\end{ptree}
\noindent$\Pi_{12}$:
\begin{ptree}
\AXC{$\forall x\,\forall y\,(x \in \mathbb{R} \to y \in \mathbb{R} \to x - y + y = x)$\lbl{[F12]}}
\AXC{$\varphi_{18}$\lbl{A2}}
\BIC{$\forall y\,(v_{253} \in \mathbb{R} \to y \in \mathbb{R} \to v_{253} - y + y = v_{253})$}
\AXC{$\varphi_{28}$\lbl{A2}}
\BIC{$v_{253} \in \mathbb{R} \to v_{254} \in \mathbb{R} \to v_{253} - v_{254} + v_{254} = v_{253}$}
\UIC{$\varphi_{2}$}
\UIC{$\forall v_{253}\,(\varphi_{2})$}
\AXC{$\varphi_{30}$\lbl{A2}}
\BIC{$\forall v_{254}\,(y \in \mathbb{R} \to v_{254} \in \mathbb{R} \to y - v_{254} + v_{254} = y)$}
\AXC{$\varphi_{16}$\lbl{A2}}
\BIC{$y \in \mathbb{R} \to x \in \mathbb{R} \to y - x + x = y$}
\AXC{$\varphi_{32}$\lbl{A1}}
\BIC{$(x \in \mathbb{R} \to y \in \mathbb{R} \to 0 < y - x \to y \in \mathbb{R}) \to \varphi_{24}$}
\end{ptree}
\noindent$\Pi_{13}$:
\begin{ptree}
\AXC{$x \in \mathbb{R} \to y \in \mathbb{R} \to 0 < y - x \to y \in \mathbb{R}$\lbl{A1}}
\AXC{$\stackrel{\Pi_{12}}{\vdots}$}\noLine
\UIC{$(x \in \mathbb{R} \to y \in \mathbb{R} \to 0 < y - x \to y \in \mathbb{R}) \to \varphi_{24}$}
\BIC{$\varphi_{24}$}
\end{ptree}
\noindent$\Pi_{14}$:
\begin{ptree}
\AXC{$\stackrel{\Pi_{11}}{\vdots}$}\noLine
\UIC{$x \in \mathbb{R} \to y \in \mathbb{R} \to 0 < y - x \to x < y - x + x$}
\AXC{$x \in \mathbb{R} \to y \in \mathbb{R} \to 0 < y - x \to x \in \mathbb{R}$\lbl{A1}}
\AXC{$\stackrel{\Pi_{13}}{\vdots}$}\noLine
\UIC{$\varphi_{24}$}
\BIC{$x \in \mathbb{R} \to y \in \mathbb{R} \to 0 < y - x \to y - x + x = y$}
\AXC{$\varphi_{34}$\lbl{A1}}
\BIC{$(x \in \mathbb{R} \to y \in \mathbb{R} \to 0 < y - x \to x < y - x + x) \to \varphi_{15}$}
\BIC{$\varphi_{15}$}
\end{ptree}
\noindent 主木:
\begin{ptree}
\AXC{$z = w \to x < z \to x < w$\lbl{A6}}
\UIC{$\forall w\,(z = w \to x < z \to x < w)$}
\UIC{$\forall z\,\forall w\,(z = w \to x < z \to x < w)$}
\AXC{$\varphi_{7}$\lbl{A2}}
\BIC{$\forall w\,(y - x + x = w \to x < y - x + x \to x < w)$}
\AXC{$\varphi_{13}$\lbl{A2}}
\BIC{$y - x + x = y \to x < y - x + x \to x < y$}
\AXC{$\stackrel{\Pi_{14}}{\vdots}$}\noLine
\UIC{$\varphi_{15}$}
\BIC{$x \in \mathbb{R} \to y \in \mathbb{R} \to 0 < y - x \to x < y$}
\UIC{$\forall y\,(x \in \mathbb{R} \to y \in \mathbb{R} \to 0 < y - x \to x < y)$}
\UIC{$\forall x\,\forall y\,(x \in \mathbb{R} \to y \in \mathbb{R} \to 0 < y - x \to x < y)$}
\end{ptree}
\begin{abbrevlist}
\item $\varphi_{1} :\equiv \forall x\,\forall y\,\forall z\,(x = y \to y = z \to x = z) \to \forall y\,\forall z\,(0 + x = y \to y = z \to 0 + x = z)$
\item $\varphi_{2} :\equiv \forall v_{254}\,(v_{253} \in \mathbb{R} \to v_{254} \in \mathbb{R} \to v_{253} - v_{254} + v_{254} = v_{253})$
\item $\varphi_{3} :\equiv \forall v_{251}\,\forall v_{252}\,(v_{251} \in \mathbb{R} \to v_{252} \in \mathbb{R} \to v_{251} - v_{252} \in \mathbb{R})$
\item $\varphi_{4} :\equiv \forall y\,\forall z\,(0 \in \mathbb{R} \to y \in \mathbb{R} \to z \in \mathbb{R} \to 0 < y \to 0 + z < y + z)$
\item $\varphi_{5} :\equiv \forall y\,\forall z\,(0 + x = y \to y = z \to 0 + x = z) \to \forall z\,(0 + x = x + 0 \to x + 0 = z \to 0 + x = z)$
\item $\varphi_{6} :\equiv \forall x\,\forall y\,\forall z\,(x \in \mathbb{R} \to y \in \mathbb{R} \to z \in \mathbb{R} \to x < y \to x + z < y + z)$
\item $\varphi_{7} :\equiv \forall z\,\forall w\,(z = w \to x < z \to x < w) \to \forall w\,(y - x + x = w \to x < y - x + x \to x < w)$
\item $\varphi_{8} :\equiv \forall z\,(0 + x = x + 0 \to x + 0 = z \to 0 + x = z) \to 0 + x = x + 0 \to x + 0 = x \to 0 + x = x$
\item $\varphi_{9} :\equiv 0 \in \mathbb{R} \to y - x \in \mathbb{R} \to x \in \mathbb{R} \to 0 < y - x \to 0 + x < y - x + x$
\item $\varphi_{10} :\equiv \forall z\,(0 \in \mathbb{R} \to y - x \in \mathbb{R} \to z \in \mathbb{R} \to 0 < y - x \to 0 + z < y - x + z)$
\item $\varphi_{11} :\equiv \forall y\,(0 \in \mathbb{R} \to y \in \mathbb{R} \to 0 + y = y + 0) \to 0 \in \mathbb{R} \to x \in \mathbb{R} \to 0 + x = x + 0$
\item $\varphi_{12} :\equiv (0 + x = x + 0 \to x + 0 = x \to 0 + x = x) \to x \in \mathbb{R} \to y \in \mathbb{R} \to 0 < y - x \to 0 + x = x$
\item $\varphi_{13} :\equiv \forall w\,(y - x + x = w \to x < y - x + x \to x < w) \to y - x + x = y \to x < y - x + x \to x < y$
\item $\varphi_{14} :\equiv \forall x\,\forall y\,(x \in \mathbb{R} \to y \in \mathbb{R} \to x + y = y + x) \to \forall y\,(0 \in \mathbb{R} \to y \in \mathbb{R} \to 0 + y = y + 0)$
\item $\varphi_{15} :\equiv (y - x + x = y \to x < y - x + x \to x < y) \to x \in \mathbb{R} \to y \in \mathbb{R} \to 0 < y - x \to x < y$
\item $\varphi_{16} :\equiv \forall v_{254}\,(y \in \mathbb{R} \to v_{254} \in \mathbb{R} \to y - v_{254} + v_{254} = y) \to y \in \mathbb{R} \to x \in \mathbb{R} \to y - x + x = y$
\item $\varphi_{17} :\equiv \forall v_{252}\,(y \in \mathbb{R} \to v_{252} \in \mathbb{R} \to y - v_{252} \in \mathbb{R}) \to y \in \mathbb{R} \to x \in \mathbb{R} \to y - x \in \mathbb{R}$
\item $\varphi_{18} :\equiv \forall x\,\forall y\,(x \in \mathbb{R} \to y \in \mathbb{R} \to x - y + y = x) \to \forall y\,(v_{253} \in \mathbb{R} \to y \in \mathbb{R} \to v_{253} - y + y = v_{253})$
\item $\varphi_{19} :\equiv \forall x\,\forall y\,(x \in \mathbb{R} \to y \in \mathbb{R} \to x - y \in \mathbb{R}) \to \forall y\,(v_{251} \in \mathbb{R} \to y \in \mathbb{R} \to v_{251} - y \in \mathbb{R})$
\item $\varphi_{20} :\equiv (x \in \mathbb{R} \to y \in \mathbb{R} \to 0 < y - x \to x \in \mathbb{R}) \to x \in \mathbb{R} \to y \in \mathbb{R} \to 0 < y - x \to x + 0 = x$
\item $\varphi_{21} :\equiv \forall z\,\forall w\,(z = w \to z < y - x + x \to w < y - x + x) \to \forall w\,(0 + x = w \to 0 + x < y - x + x \to w < y - x + x)$
\item $\varphi_{22} :\equiv \forall y\,(v_{251} \in \mathbb{R} \to y \in \mathbb{R} \to v_{251} - y \in \mathbb{R}) \to v_{251} \in \mathbb{R} \to v_{252} \in \mathbb{R} \to v_{251} - v_{252} \in \mathbb{R}$
\item $\varphi_{23} :\equiv (x \in \mathbb{R} \to y \in \mathbb{R} \to 0 < y - x \to x \in \mathbb{R}) \to x \in \mathbb{R} \to y \in \mathbb{R} \to 0 < y - x \to 0 + x = x + 0$
\item $\varphi_{24} :\equiv (x \in \mathbb{R} \to y \in \mathbb{R} \to 0 < y - x \to x \in \mathbb{R}) \to x \in \mathbb{R} \to y \in \mathbb{R} \to 0 < y - x \to y - x + x = y$
\item $\varphi_{25} :\equiv \forall w\,(0 + x = w \to 0 + x < y - x + x \to w < y - x + x) \to 0 + x = x \to 0 + x < y - x + x \to x < y - x + x$
\item $\varphi_{26} :\equiv (x \in \mathbb{R} \to y \in \mathbb{R} \to 0 < y - x \to x \in \mathbb{R}) \to x \in \mathbb{R} \to y \in \mathbb{R} \to 0 < y - x \to y - x \in \mathbb{R}$
\item $\varphi_{27} :\equiv (0 + x = x \to 0 + x < y - x + x \to x < y - x + x) \to x \in \mathbb{R} \to y \in \mathbb{R} \to 0 < y - x \to x < y - x + x$
\item $\varphi_{28} :\equiv \forall y\,(v_{253} \in \mathbb{R} \to y \in \mathbb{R} \to v_{253} - y + y = v_{253}) \to v_{253} \in \mathbb{R} \to v_{254} \in \mathbb{R} \to v_{253} - v_{254} + v_{254} = v_{253}$
\item $\varphi_{29} :\equiv (x \in \mathbb{R} \to y \in \mathbb{R} \to 0 < y - x \to 0 < y - x) \to x \in \mathbb{R} \to y \in \mathbb{R} \to 0 < y - x \to 0 + x < y - x + x$
\item $\varphi_{30} :\equiv \forall v_{253}\,(\varphi_{2}) \to \forall v_{254}\,(y \in \mathbb{R} \to v_{254} \in \mathbb{R} \to y - v_{254} + v_{254} = y)$
\item $\varphi_{31} :\equiv (x \in \mathbb{R} \to y \in \mathbb{R} \to 0 < y - x \to 0 + x = x + 0) \to (x \in \mathbb{R} \to y \in \mathbb{R} \to 0 < y - x \to x + 0 = x) \to \varphi_{12}$
\item $\varphi_{32} :\equiv (y \in \mathbb{R} \to x \in \mathbb{R} \to y - x + x = y) \to (x \in \mathbb{R} \to y \in \mathbb{R} \to 0 < y - x \to y \in \mathbb{R}) \to \varphi_{24}$
\item $\varphi_{33} :\equiv (y \in \mathbb{R} \to x \in \mathbb{R} \to y - x \in \mathbb{R}) \to (x \in \mathbb{R} \to y \in \mathbb{R} \to 0 < y - x \to y \in \mathbb{R}) \to \varphi_{26}$
\item $\varphi_{34} :\equiv (x \in \mathbb{R} \to y \in \mathbb{R} \to 0 < y - x \to y - x + x = y) \to (x \in \mathbb{R} \to y \in \mathbb{R} \to 0 < y - x \to x < y - x + x) \to \varphi_{15}$
\item $\varphi_{35} :\equiv (x \in \mathbb{R} \to y \in \mathbb{R} \to 0 < y - x \to 0 + x = x) \to (x \in \mathbb{R} \to y \in \mathbb{R} \to 0 < y - x \to 0 + x < y - x + x) \to \varphi_{27}$
\item $\varphi_{36} :\equiv (x \in \mathbb{R} \to y \in \mathbb{R} \to 0 < y - x \to y - x \in \mathbb{R}) \to (x \in \mathbb{R} \to y \in \mathbb{R} \to 0 < y - x \to x \in \mathbb{R}) \to \varphi_{29}$
\end{abbrevlist}
\end{fullproof}

% 自動生成 (prover/build.js): LtNe  ノード数 35
\begin{fullproof}{[F26]}{$\forall x\,\forall y\,(x < y \to x \neq y)$}{35}\phantomsection\label{tree:LtNe}
\noindent$\Pi_{1}$:
\begin{ptree}
\AXC{$z = w \to x < z \to x < w$\lbl{A6}}
\UIC{$\forall w\,(z = w \to x < z \to x < w)$}
\UIC{$\forall z\,\forall w\,(z = w \to x < z \to x < w)$}
\AXC{$\varphi_{2}$\lbl{A2}}
\BIC{$\forall w\,(y = w \to x < y \to x < w)$}
\AXC{$\forall w\,(y = w \to x < y \to x < w) \to y = x \to x < y \to x < x$\lbl{A2}}
\BIC{$y = x \to x < y \to x < x$}
\end{ptree}
\noindent$\Pi_{2}$:
\begin{ptree}
\AXC{$x = x$\lbl{A5}}
\AXC{$x = y \to x = x \to y = x$\lbl{A6}}
\AXC{$(x = y \to x = x \to y = x) \to x = x \to x = y \to y = x$\lbl{A1}}
\BIC{$x = x \to x = y \to y = x$}
\BIC{$x = y \to y = x$}
\UIC{$\forall y\,(x = y \to y = x)$}
\UIC{$\forall x\,\forall y\,(x = y \to y = x)$}
\end{ptree}
\noindent$\Pi_{3}$:
\begin{ptree}
\AXC{$\stackrel{\Pi_{2}}{\vdots}$}\noLine
\UIC{$\forall x\,\forall y\,(x = y \to y = x)$}
\AXC{$\forall x\,\forall y\,(x = y \to y = x) \to \forall y\,(x = y \to y = x)$\lbl{A2}}
\BIC{$\forall y\,(x = y \to y = x)$}
\AXC{$\forall y\,(x = y \to y = x) \to x = y \to y = x$\lbl{A2}}
\BIC{$x = y \to y = x$}
\end{ptree}
\noindent$\Pi_{4}$:
\begin{ptree}
\AXC{$x < y \to x = y \to x < y$\lbl{A1}}
\AXC{$\stackrel{\Pi_{3}}{\vdots}$}\noLine
\UIC{$x = y \to y = x$}
\AXC{$x < y \to x = y \to x = y$\lbl{A1}}
\AXC{$\varphi_{1}$\lbl{A1}}
\BIC{$(x = y \to y = x) \to x < y \to x = y \to y = x$}
\BIC{$x < y \to x = y \to y = x$}
\AXC{$(x < y \to x = y \to y = x) \to \varphi_{3}$\lbl{A1}}
\BIC{$\varphi_{3}$}
\BIC{$(y = x \to x < y \to x < x) \to x < y \to x = y \to x < x$}
\end{ptree}
\noindent$\Pi_{5}$:
\begin{ptree}
\AXC{$\forall x\,(\neg x < x)$\lbl{[O1]}}
\AXC{$\forall x\,(\neg x < x) \to \neg x < x$\lbl{A2}}
\BIC{$\neg x < x$}
\AXC{$\neg x < x \to (x < y \to x = y \to x < x) \to x < y \to x \neq y$\lbl{A1}}
\BIC{$(x < y \to x = y \to x < x) \to x < y \to x \neq y$}
\end{ptree}
\noindent 主木:
\begin{ptree}
\AXC{$\stackrel{\Pi_{1}}{\vdots}$}\noLine
\UIC{$y = x \to x < y \to x < x$}
\AXC{$\stackrel{\Pi_{4}}{\vdots}$}\noLine
\UIC{$(y = x \to x < y \to x < x) \to x < y \to x = y \to x < x$}
\BIC{$x < y \to x = y \to x < x$}
\AXC{$\stackrel{\Pi_{5}}{\vdots}$}\noLine
\UIC{$(x < y \to x = y \to x < x) \to x < y \to x \neq y$}
\BIC{$x < y \to x \neq y$}
\UIC{$\forall y\,(x < y \to x \neq y)$}
\UIC{$\forall x\,\forall y\,(x < y \to x \neq y)$}
\end{ptree}
\begin{abbrevlist}
\item $\varphi_{1} :\equiv (x < y \to x = y \to x = y) \to (x = y \to y = x) \to x < y \to x = y \to y = x$
\item $\varphi_{2} :\equiv \forall z\,\forall w\,(z = w \to x < z \to x < w) \to \forall w\,(y = w \to x < y \to x < w)$
\item $\varphi_{3} :\equiv (x < y \to x = y \to x < y) \to (y = x \to x < y \to x < x) \to x < y \to x = y \to x < x$
\end{abbrevlist}
\end{fullproof}

% 自動生成 (prover/build.js): NegOneSq  ノード数 32
\begin{fullproof}{[F27]}{$-1 \cdot -1 = 1$}{32}\phantomsection\label{tree:NegOneSq}
\noindent$\Pi_{1}$:
\begin{ptree}
\AXC{$y = z \to x = y \to x = z$\lbl{A6}}
\AXC{$(y = z \to x = y \to x = z) \to x = y \to y = z \to x = z$\lbl{A1}}
\BIC{$x = y \to y = z \to x = z$}
\UIC{$\forall z\,(x = y \to y = z \to x = z)$}
\UIC{$\forall y\,\forall z\,(x = y \to y = z \to x = z)$}
\UIC{$\forall x\,\forall y\,\forall z\,(x = y \to y = z \to x = z)$}
\AXC{$\varphi_{2}$\lbl{A2}}
\BIC{$\forall y\,\forall z\,(-1 \cdot -1 = y \to y = z \to -1 \cdot -1 = z)$}
\AXC{$\varphi_{6}$\lbl{A2}}
\BIC{$\forall z\,(-1 \cdot -1 = 1 \cdot 1 \to 1 \cdot 1 = z \to -1 \cdot -1 = z)$}
\AXC{$\varphi_{7}$\lbl{A2}}
\BIC{$-1 \cdot -1 = 1 \cdot 1 \to 1 \cdot 1 = 1 \to -1 \cdot -1 = 1$}
\end{ptree}
\noindent$\Pi_{2}$:
\begin{ptree}
\AXC{$\forall x \in \mathbb{R}\,(x \cdot 1 = x)$\lbl{(R8)}}
\AXC{$\forall x \in \mathbb{R}\,(x \cdot 1 = x) \to 1 \in \mathbb{R} \to 1 \cdot 1 = 1$\lbl{A2}}
\BIC{$1 \in \mathbb{R} \to 1 \cdot 1 = 1$}
\AXC{$(1 \in \mathbb{R} \to 1 \cdot 1 = 1) \to 1 \in \mathbb{R} \to 1 \cdot 1 = 1$\lbl{A1}}
\BIC{$1 \in \mathbb{R} \to 1 \cdot 1 = 1$}
\end{ptree}
\noindent$\Pi_{3}$:
\begin{ptree}
\AXC{$1 \in \mathbb{R}$\lbl{(R1)}}
\AXC{$1 \in \mathbb{R}$\lbl{(R1)}}
\AXC{$\forall x\,\forall y\,(x \in \mathbb{R} \to y \in \mathbb{R} \to -x \cdot -y = x \cdot y)$\lbl{[F11]}}
\AXC{$\varphi_{5}$\lbl{A2}}
\BIC{$\forall y\,(1 \in \mathbb{R} \to y \in \mathbb{R} \to -1 \cdot -y = 1 \cdot y)$}
\AXC{$\varphi_{4}$\lbl{A2}}
\BIC{$1 \in \mathbb{R} \to 1 \in \mathbb{R} \to -1 \cdot -1 = 1 \cdot 1$}
\AXC{$\varphi_{3}$\lbl{A1}}
\BIC{$1 \in \mathbb{R} \to 1 \in \mathbb{R} \to -1 \cdot -1 = 1 \cdot 1$}
\BIC{$1 \in \mathbb{R} \to -1 \cdot -1 = 1 \cdot 1$}
\BIC{$-1 \cdot -1 = 1 \cdot 1$}
\end{ptree}
\noindent$\Pi_{4}$:
\begin{ptree}
\AXC{$1 \in \mathbb{R}$\lbl{(R1)}}
\AXC{$\stackrel{\Pi_{2}}{\vdots}$}\noLine
\UIC{$1 \in \mathbb{R} \to 1 \cdot 1 = 1$}
\BIC{$1 \cdot 1 = 1$}
\AXC{$\stackrel{\Pi_{3}}{\vdots}$}\noLine
\UIC{$-1 \cdot -1 = 1 \cdot 1$}
\AXC{$-1 \cdot -1 = 1 \cdot 1 \to 1 \cdot 1 = 1 \to \varphi_{1}$\lbl{A1}}
\BIC{$1 \cdot 1 = 1 \to \varphi_{1}$}
\BIC{$\varphi_{1}$}
\end{ptree}
\noindent 主木:
\begin{ptree}
\AXC{$\stackrel{\Pi_{1}}{\vdots}$}\noLine
\UIC{$-1 \cdot -1 = 1 \cdot 1 \to 1 \cdot 1 = 1 \to -1 \cdot -1 = 1$}
\AXC{$\stackrel{\Pi_{4}}{\vdots}$}\noLine
\UIC{$\varphi_{1}$}
\BIC{$-1 \cdot -1 = 1$}
\end{ptree}
\begin{abbrevlist}
\item $\varphi_{1} :\equiv (-1 \cdot -1 = 1 \cdot 1 \to 1 \cdot 1 = 1 \to -1 \cdot -1 = 1) \to -1 \cdot -1 = 1$
\item $\varphi_{2} :\equiv \forall x\,\forall y\,\forall z\,(x = y \to y = z \to x = z) \to \forall y\,\forall z\,(-1 \cdot -1 = y \to y = z \to -1 \cdot -1 = z)$
\item $\varphi_{3} :\equiv (1 \in \mathbb{R} \to 1 \in \mathbb{R} \to -1 \cdot -1 = 1 \cdot 1) \to 1 \in \mathbb{R} \to 1 \in \mathbb{R} \to -1 \cdot -1 = 1 \cdot 1$
\item $\varphi_{4} :\equiv \forall y\,(1 \in \mathbb{R} \to y \in \mathbb{R} \to -1 \cdot -y = 1 \cdot y) \to 1 \in \mathbb{R} \to 1 \in \mathbb{R} \to -1 \cdot -1 = 1 \cdot 1$
\item $\varphi_{5} :\equiv \forall x\,\forall y\,(x \in \mathbb{R} \to y \in \mathbb{R} \to -x \cdot -y = x \cdot y) \to \forall y\,(1 \in \mathbb{R} \to y \in \mathbb{R} \to -1 \cdot -y = 1 \cdot y)$
\item $\varphi_{6} :\equiv \forall y\,\forall z\,(-1 \cdot -1 = y \to y = z \to -1 \cdot -1 = z) \to \forall z\,(-1 \cdot -1 = 1 \cdot 1 \to 1 \cdot 1 = z \to -1 \cdot -1 = z)$
\item $\varphi_{7} :\equiv \forall z\,(-1 \cdot -1 = 1 \cdot 1 \to 1 \cdot 1 = z \to -1 \cdot -1 = z) \to -1 \cdot -1 = 1 \cdot 1 \to 1 \cdot 1 = 1 \to -1 \cdot -1 = 1$
\end{abbrevlist}
\end{fullproof}

% 自動生成 (prover/build.js): ZeroLtOne  ノード数 70
\begin{fullproof}{[F28]}{$0 < 1$}{70}\phantomsection\label{tree:ZeroLtOne}
\noindent$\Pi_{1}$:
\begin{ptree}
\AXC{$0 \in \mathbb{R}$\lbl{(R1)}}
\AXC{$1 \in \mathbb{R}$\lbl{(R1)}}
\AXC{$\forall x\,\forall y\,(x \in \mathbb{R} \to y \in \mathbb{R} \to x < y \to -y < -x)$\lbl{[F22]}}
\AXC{$\varphi_{12}$\lbl{A2}}
\BIC{$\forall y\,(1 \in \mathbb{R} \to y \in \mathbb{R} \to 1 < y \to -y < -1)$}
\AXC{$\varphi_{11}$\lbl{A2}}
\BIC{$1 \in \mathbb{R} \to 0 \in \mathbb{R} \to 1 < 0 \to -0 < -1$}
\AXC{$\varphi_{13}$\lbl{A1}}
\BIC{$1 \in \mathbb{R} \to 0 \in \mathbb{R} \to (1 < 0 \to 1 < 0) \to 1 < 0 \to -0 < -1$}
\BIC{$0 \in \mathbb{R} \to (1 < 0 \to 1 < 0) \to 1 < 0 \to -0 < -1$}
\BIC{$(1 < 0 \to 1 < 0) \to 1 < 0 \to -0 < -1$}
\end{ptree}
\noindent$\Pi_{2}$:
\begin{ptree}
\AXC{$1 < 0 \to 1 < 0$\lbl{A1}}
\AXC{$\stackrel{\Pi_{1}}{\vdots}$}\noLine
\UIC{$(1 < 0 \to 1 < 0) \to 1 < 0 \to -0 < -1$}
\BIC{$1 < 0 \to -0 < -1$}
\AXC{$-0 = 0$\lbl{[F6]}}
\AXC{$-0 = 0 \to \varphi_{5}$\lbl{A1}}
\BIC{$\varphi_{5}$}
\BIC{$(-0 = 0 \to -0 < -1 \to 0 < -1) \to 1 < 0 \to 0 < -1$}
\end{ptree}
\noindent$\Pi_{3}$:
\begin{ptree}
\AXC{$z = w \to z < -1 \to w < -1$\lbl{A6}}
\UIC{$\forall w\,(z = w \to z < -1 \to w < -1)$}
\UIC{$\forall z\,\forall w\,(z = w \to z < -1 \to w < -1)$}
\AXC{$\varphi_{3}$\lbl{A2}}
\BIC{$\forall w\,(-0 = w \to -0 < -1 \to w < -1)$}
\AXC{$\varphi_{4}$\lbl{A2}}
\BIC{$-0 = 0 \to -0 < -1 \to 0 < -1$}
\AXC{$\stackrel{\Pi_{2}}{\vdots}$}\noLine
\UIC{$(-0 = 0 \to -0 < -1 \to 0 < -1) \to 1 < 0 \to 0 < -1$}
\BIC{$1 < 0 \to 0 < -1$}
\end{ptree}
\noindent$\Pi_{4}$:
\begin{ptree}
\AXC{$\forall x \in \mathbb{R}\,(-x \in \mathbb{R})$\lbl{(R1)}}
\AXC{$\forall x \in \mathbb{R}\,(-x \in \mathbb{R}) \to 1 \in \mathbb{R} \to -1 \in \mathbb{R}$\lbl{A2}}
\BIC{$1 \in \mathbb{R} \to -1 \in \mathbb{R}$}
\end{ptree}
\noindent$\Pi_{5}$:
\begin{ptree}
\AXC{$1 \in \mathbb{R}$\lbl{(R1)}}
\AXC{$\stackrel{\Pi_{4}}{\vdots}$}\noLine
\UIC{$1 \in \mathbb{R} \to -1 \in \mathbb{R}$}
\AXC{$(1 \in \mathbb{R} \to -1 \in \mathbb{R}) \to 1 \in \mathbb{R} \to 1 < 0 \to -1 \in \mathbb{R}$\lbl{A1}}
\BIC{$1 \in \mathbb{R} \to 1 < 0 \to -1 \in \mathbb{R}$}
\BIC{$1 < 0 \to -1 \in \mathbb{R}$}
\end{ptree}
\noindent$\Pi_{6}$:
\begin{ptree}
\AXC{$\stackrel{\Pi_{5}}{\vdots}$}\noLine
\UIC{$1 < 0 \to -1 \in \mathbb{R}$}
\AXC{$\forall x\,\forall y\,(x \in \mathbb{R} \to y \in \mathbb{R} \to 0 < x \to 0 < y \to 0 < x \cdot y)$\lbl{(R15)}}
\AXC{$\varphi_{14}$\lbl{A2}}
\BIC{$\forall y\,(-1 \in \mathbb{R} \to y \in \mathbb{R} \to 0 < -1 \to 0 < y \to 0 < -1 \cdot y)$}
\AXC{$\varphi_{15}$\lbl{A2}}
\BIC{$-1 \in \mathbb{R} \to -1 \in \mathbb{R} \to 0 < -1 \to 0 < -1 \to 0 < -1 \cdot -1$}
\AXC{$\varphi_{16}$\lbl{A1}}
\BIC{$(1 < 0 \to -1 \in \mathbb{R}) \to (1 < 0 \to -1 \in \mathbb{R}) \to \varphi_{1}$}
\BIC{$(1 < 0 \to -1 \in \mathbb{R}) \to \varphi_{1}$}
\end{ptree}
\noindent$\Pi_{7}$:
\begin{ptree}
\AXC{$\stackrel{\Pi_{3}}{\vdots}$}\noLine
\UIC{$1 < 0 \to 0 < -1$}
\AXC{$\stackrel{\Pi_{3}}{\vdots}$}\noLine
\UIC{$1 < 0 \to 0 < -1$}
\AXC{$\stackrel{\Pi_{5}}{\vdots}$}\noLine
\UIC{$1 < 0 \to -1 \in \mathbb{R}$}
\AXC{$\stackrel{\Pi_{6}}{\vdots}$}\noLine
\UIC{$(1 < 0 \to -1 \in \mathbb{R}) \to \varphi_{1}$}
\BIC{$\varphi_{1}$}
\BIC{$(1 < 0 \to 0 < -1) \to 1 < 0 \to 0 < -1 \cdot -1$}
\BIC{$1 < 0 \to 0 < -1 \cdot -1$}
\end{ptree}
\noindent$\Pi_{8}$:
\begin{ptree}
\AXC{$\stackrel{\Pi_{7}}{\vdots}$}\noLine
\UIC{$1 < 0 \to 0 < -1 \cdot -1$}
\AXC{$-1 \cdot -1 = 1$\lbl{[F27]}}
\AXC{$-1 \cdot -1 = 1 \to \varphi_{8}$\lbl{A1}}
\BIC{$\varphi_{8}$}
\BIC{$(-1 \cdot -1 = 1 \to 0 < -1 \cdot -1 \to 0 < 1) \to 1 < 0 \to 0 < 1$}
\end{ptree}
\noindent$\Pi_{9}$:
\begin{ptree}
\AXC{$z = w \to 0 < z \to 0 < w$\lbl{A6}}
\UIC{$\forall w\,(z = w \to 0 < z \to 0 < w)$}
\UIC{$\forall z\,\forall w\,(z = w \to 0 < z \to 0 < w)$}
\AXC{$\varphi_{2}$\lbl{A2}}
\BIC{$\forall w\,(-1 \cdot -1 = w \to 0 < -1 \cdot -1 \to 0 < w)$}
\AXC{$\varphi_{10}$\lbl{A2}}
\BIC{$-1 \cdot -1 = 1 \to 0 < -1 \cdot -1 \to 0 < 1$}
\AXC{$\stackrel{\Pi_{8}}{\vdots}$}\noLine
\UIC{$(-1 \cdot -1 = 1 \to 0 < -1 \cdot -1 \to 0 < 1) \to 1 < 0 \to 0 < 1$}
\BIC{$1 < 0 \to 0 < 1$}
\end{ptree}
\noindent$\Pi_{10}$:
\begin{ptree}
\AXC{$0 \in \mathbb{R}$\lbl{(R1)}}
\AXC{$\forall x\,\forall y\,(x \in \mathbb{R} \to y \in \mathbb{R} \to x < y \lor x = y \lor y < x)$\lbl{(R13)}}
\AXC{$\varphi_{9}$\lbl{A2}}
\BIC{$\forall y\,(0 \in \mathbb{R} \to y \in \mathbb{R} \to 0 < y \lor 0 = y \lor y < 0)$}
\AXC{$\varphi_{7}$\lbl{A2}}
\BIC{$0 \in \mathbb{R} \to 1 \in \mathbb{R} \to 0 < 1 \lor 0 = 1 \lor 1 < 0$}
\AXC{$\varphi_{6}$\lbl{A1}}
\BIC{$0 \in \mathbb{R} \to 1 \in \mathbb{R} \to 0 < 1 \lor 0 = 1 \lor 1 < 0$}
\BIC{$1 \in \mathbb{R} \to 0 < 1 \lor 0 = 1 \lor 1 < 0$}
\end{ptree}
\noindent$\Pi_{11}$:
\begin{ptree}
\AXC{$0 \neq 1$\lbl{(R8)}}
\AXC{$1 \in \mathbb{R}$\lbl{(R1)}}
\AXC{$\stackrel{\Pi_{10}}{\vdots}$}\noLine
\UIC{$1 \in \mathbb{R} \to 0 < 1 \lor 0 = 1 \lor 1 < 0$}
\BIC{$0 < 1 \lor 0 = 1 \lor 1 < 0$}
\AXC{$0 < 1 \lor 0 = 1 \lor 1 < 0 \to 0 \neq 1 \to (1 < 0 \to 0 < 1) \to 0 < 1$\lbl{A1}}
\BIC{$0 \neq 1 \to (1 < 0 \to 0 < 1) \to 0 < 1$}
\BIC{$(1 < 0 \to 0 < 1) \to 0 < 1$}
\end{ptree}
\noindent 主木:
\begin{ptree}
\AXC{$\stackrel{\Pi_{9}}{\vdots}$}\noLine
\UIC{$1 < 0 \to 0 < 1$}
\AXC{$\stackrel{\Pi_{11}}{\vdots}$}\noLine
\UIC{$(1 < 0 \to 0 < 1) \to 0 < 1$}
\BIC{$0 < 1$}
\end{ptree}
\begin{abbrevlist}
\item $\varphi_{1} :\equiv (1 < 0 \to 0 < -1) \to (1 < 0 \to 0 < -1) \to 1 < 0 \to 0 < -1 \cdot -1$
\item $\varphi_{2} :\equiv \forall z\,\forall w\,(z = w \to 0 < z \to 0 < w) \to \forall w\,(-1 \cdot -1 = w \to 0 < -1 \cdot -1 \to 0 < w)$
\item $\varphi_{3} :\equiv \forall z\,\forall w\,(z = w \to z < -1 \to w < -1) \to \forall w\,(-0 = w \to -0 < -1 \to w < -1)$
\item $\varphi_{4} :\equiv \forall w\,(-0 = w \to -0 < -1 \to w < -1) \to -0 = 0 \to -0 < -1 \to 0 < -1$
\item $\varphi_{5} :\equiv (1 < 0 \to -0 < -1) \to (-0 = 0 \to -0 < -1 \to 0 < -1) \to 1 < 0 \to 0 < -1$
\item $\varphi_{6} :\equiv (0 \in \mathbb{R} \to 1 \in \mathbb{R} \to 0 < 1 \lor 0 = 1 \lor 1 < 0) \to 0 \in \mathbb{R} \to 1 \in \mathbb{R} \to 0 < 1 \lor 0 = 1 \lor 1 < 0$
\item $\varphi_{7} :\equiv \forall y\,(0 \in \mathbb{R} \to y \in \mathbb{R} \to 0 < y \lor 0 = y \lor y < 0) \to 0 \in \mathbb{R} \to 1 \in \mathbb{R} \to 0 < 1 \lor 0 = 1 \lor 1 < 0$
\item $\varphi_{8} :\equiv (1 < 0 \to 0 < -1 \cdot -1) \to (-1 \cdot -1 = 1 \to 0 < -1 \cdot -1 \to 0 < 1) \to 1 < 0 \to 0 < 1$
\item $\varphi_{9} :\equiv \forall x\,\forall y\,(x \in \mathbb{R} \to y \in \mathbb{R} \to x < y \lor x = y \lor y < x) \to \forall y\,(0 \in \mathbb{R} \to y \in \mathbb{R} \to 0 < y \lor 0 = y \lor y < 0)$
\item $\varphi_{10} :\equiv \forall w\,(-1 \cdot -1 = w \to 0 < -1 \cdot -1 \to 0 < w) \to -1 \cdot -1 = 1 \to 0 < -1 \cdot -1 \to 0 < 1$
\item $\varphi_{11} :\equiv \forall y\,(1 \in \mathbb{R} \to y \in \mathbb{R} \to 1 < y \to -y < -1) \to 1 \in \mathbb{R} \to 0 \in \mathbb{R} \to 1 < 0 \to -0 < -1$
\item $\varphi_{12} :\equiv \forall x\,\forall y\,(x \in \mathbb{R} \to y \in \mathbb{R} \to x < y \to -y < -x) \to \forall y\,(1 \in \mathbb{R} \to y \in \mathbb{R} \to 1 < y \to -y < -1)$
\item $\varphi_{13} :\equiv (1 \in \mathbb{R} \to 0 \in \mathbb{R} \to 1 < 0 \to -0 < -1) \to 1 \in \mathbb{R} \to 0 \in \mathbb{R} \to (1 < 0 \to 1 < 0) \to 1 < 0 \to -0 < -1$
\item $\varphi_{14} :\equiv \forall x\,\forall y\,(x \in \mathbb{R} \to y \in \mathbb{R} \to 0 < x \to 0 < y \to 0 < x \cdot y) \to \forall y\,(-1 \in \mathbb{R} \to y \in \mathbb{R} \to 0 < -1 \to 0 < y \to 0 < -1 \cdot y)$
\item $\varphi_{15} :\equiv \forall y\,(-1 \in \mathbb{R} \to y \in \mathbb{R} \to 0 < -1 \to 0 < y \to 0 < -1 \cdot y) \to -1 \in \mathbb{R} \to -1 \in \mathbb{R} \to 0 < -1 \to 0 < -1 \to 0 < -1 \cdot -1$
\item $\varphi_{16} :\equiv (-1 \in \mathbb{R} \to -1 \in \mathbb{R} \to 0 < -1 \to 0 < -1 \to 0 < -1 \cdot -1) \to (1 < 0 \to -1 \in \mathbb{R}) \to (1 < 0 \to -1 \in \mathbb{R}) \to \varphi_{1}$
\end{abbrevlist}
\end{fullproof}

% 自動生成 (prover/build.js): TwoPos  ノード数 96
\begin{fullproof}{[F29]}{$0 < 2$}{96}\phantomsection\label{tree:TwoPos}
\noindent$\Pi_{1}$:
\begin{ptree}
\AXC{$1 \in \mathbb{R}$\lbl{(R1)}}
\AXC{$0 \in \mathbb{R}$\lbl{(R1)}}
\AXC{$\varphi_{10}$\lbl{(R14)}}
\AXC{$\varphi_{10} \to \varphi_{8}$\lbl{A2}}
\BIC{$\varphi_{8}$}
\AXC{$\varphi_{20}$\lbl{A2}}
\BIC{$\forall z\,(0 \in \mathbb{R} \to 1 \in \mathbb{R} \to z \in \mathbb{R} \to 0 < 1 \to 0 + z < 1 + z)$}
\AXC{$\varphi_{19}$\lbl{A2}}
\BIC{$0 \in \mathbb{R} \to 1 \in \mathbb{R} \to 1 \in \mathbb{R} \to 0 < 1 \to 0 + 1 < 1 + 1$}
\AXC{$\varphi_{18}$\lbl{A1}}
\BIC{$0 \in \mathbb{R} \to 1 \in \mathbb{R} \to 1 \in \mathbb{R} \to 0 < 1 \to 0 + 1 < 1 + 1$}
\BIC{$1 \in \mathbb{R} \to 1 \in \mathbb{R} \to 0 < 1 \to 0 + 1 < 1 + 1$}
\BIC{$1 \in \mathbb{R} \to 0 < 1 \to 0 + 1 < 1 + 1$}
\end{ptree}
\noindent$\Pi_{2}$:
\begin{ptree}
\AXC{$y = z \to x = y \to x = z$\lbl{A6}}
\AXC{$(y = z \to x = y \to x = z) \to x = y \to y = z \to x = z$\lbl{A1}}
\BIC{$x = y \to y = z \to x = z$}
\UIC{$\forall z\,(x = y \to y = z \to x = z)$}
\UIC{$\forall y\,\forall z\,(x = y \to y = z \to x = z)$}
\UIC{$\forall x\,\forall y\,\forall z\,(x = y \to y = z \to x = z)$}
\AXC{$\varphi_{2}$\lbl{A2}}
\BIC{$\forall y\,\forall z\,(0 + 1 = y \to y = z \to 0 + 1 = z)$}
\AXC{$\varphi_{9}$\lbl{A2}}
\BIC{$\forall z\,(0 + 1 = 1 + 0 \to 1 + 0 = z \to 0 + 1 = z)$}
\AXC{$\varphi_{12}$\lbl{A2}}
\BIC{$0 + 1 = 1 + 0 \to 1 + 0 = 1 \to 0 + 1 = 1$}
\end{ptree}
\noindent$\Pi_{3}$:
\begin{ptree}
\AXC{$\forall x \in \mathbb{R}\,(x + 0 = x)$\lbl{(R4)}}
\AXC{$\forall x \in \mathbb{R}\,(x + 0 = x) \to 1 \in \mathbb{R} \to 1 + 0 = 1$\lbl{A2}}
\BIC{$1 \in \mathbb{R} \to 1 + 0 = 1$}
\AXC{$(1 \in \mathbb{R} \to 1 + 0 = 1) \to 1 \in \mathbb{R} \to 1 + 0 = 1$\lbl{A1}}
\BIC{$1 \in \mathbb{R} \to 1 + 0 = 1$}
\end{ptree}
\noindent$\Pi_{4}$:
\begin{ptree}
\AXC{$1 \in \mathbb{R}$\lbl{(R1)}}
\AXC{$0 \in \mathbb{R}$\lbl{(R1)}}
\AXC{$\forall x\,\forall y\,(x \in \mathbb{R} \to y \in \mathbb{R} \to x + y = y + x)$\lbl{(R3)}}
\AXC{$\varphi_{17}$\lbl{A2}}
\BIC{$\forall y\,(0 \in \mathbb{R} \to y \in \mathbb{R} \to 0 + y = y + 0)$}
\AXC{$\varphi_{15}$\lbl{A2}}
\BIC{$0 \in \mathbb{R} \to 1 \in \mathbb{R} \to 0 + 1 = 1 + 0$}
\AXC{$\varphi_{13}$\lbl{A1}}
\BIC{$0 \in \mathbb{R} \to 1 \in \mathbb{R} \to 0 + 1 = 1 + 0$}
\BIC{$1 \in \mathbb{R} \to 0 + 1 = 1 + 0$}
\BIC{$0 + 1 = 1 + 0$}
\end{ptree}
\noindent$\Pi_{5}$:
\begin{ptree}
\AXC{$1 \in \mathbb{R}$\lbl{(R1)}}
\AXC{$\stackrel{\Pi_{3}}{\vdots}$}\noLine
\UIC{$1 \in \mathbb{R} \to 1 + 0 = 1$}
\BIC{$1 + 0 = 1$}
\AXC{$\stackrel{\Pi_{4}}{\vdots}$}\noLine
\UIC{$0 + 1 = 1 + 0$}
\AXC{$0 + 1 = 1 + 0 \to \varphi_{1}$\lbl{A1}}
\BIC{$\varphi_{1}$}
\BIC{$(0 + 1 = 1 + 0 \to 1 + 0 = 1 \to 0 + 1 = 1) \to 0 + 1 = 1$}
\end{ptree}
\noindent$\Pi_{6}$:
\begin{ptree}
\AXC{$\stackrel{\Pi_{2}}{\vdots}$}\noLine
\UIC{$0 + 1 = 1 + 0 \to 1 + 0 = 1 \to 0 + 1 = 1$}
\AXC{$\stackrel{\Pi_{5}}{\vdots}$}\noLine
\UIC{$(0 + 1 = 1 + 0 \to 1 + 0 = 1 \to 0 + 1 = 1) \to 0 + 1 = 1$}
\BIC{$0 + 1 = 1$}
\AXC{$0 + 1 = 1 \to \varphi_{11}$\lbl{A1}}
\BIC{$\varphi_{11}$}
\end{ptree}
\noindent$\Pi_{7}$:
\begin{ptree}
\AXC{$0 < 1$\lbl{[F28]}}
\AXC{$1 \in \mathbb{R}$\lbl{(R1)}}
\AXC{$\stackrel{\Pi_{1}}{\vdots}$}\noLine
\UIC{$1 \in \mathbb{R} \to 0 < 1 \to 0 + 1 < 1 + 1$}
\BIC{$0 < 1 \to 0 + 1 < 1 + 1$}
\BIC{$0 + 1 < 1 + 1$}
\AXC{$\stackrel{\Pi_{6}}{\vdots}$}\noLine
\UIC{$\varphi_{11}$}
\BIC{$(0 + 1 = 1 \to 0 + 1 < 1 + 1 \to 1 < 1 + 1) \to 1 < 1 + 1$}
\end{ptree}
\noindent$\Pi_{8}$:
\begin{ptree}
\AXC{$z = w \to z < 1 + 1 \to w < 1 + 1$\lbl{A6}}
\UIC{$\forall w\,(z = w \to z < 1 + 1 \to w < 1 + 1)$}
\UIC{$\forall z\,\forall w\,(z = w \to z < 1 + 1 \to w < 1 + 1)$}
\AXC{$\varphi_{14}$\lbl{A2}}
\BIC{$\forall w\,(0 + 1 = w \to 0 + 1 < 1 + 1 \to w < 1 + 1)$}
\AXC{$\varphi_{16}$\lbl{A2}}
\BIC{$0 + 1 = 1 \to 0 + 1 < 1 + 1 \to 1 < 1 + 1$}
\AXC{$\stackrel{\Pi_{7}}{\vdots}$}\noLine
\UIC{$(0 + 1 = 1 \to 0 + 1 < 1 + 1 \to 1 < 1 + 1) \to 1 < 1 + 1$}
\BIC{$1 < 1 + 1$}
\end{ptree}
\noindent$\Pi_{9}$:
\begin{ptree}
\AXC{$x = x$\lbl{A5}}
\AXC{$x = y \to x = x \to y = x$\lbl{A6}}
\AXC{$(x = y \to x = x \to y = x) \to x = x \to x = y \to y = x$\lbl{A1}}
\BIC{$x = x \to x = y \to y = x$}
\BIC{$x = y \to y = x$}
\UIC{$\forall y\,(x = y \to y = x)$}
\UIC{$\forall x\,\forall y\,(x = y \to y = x)$}
\end{ptree}
\noindent$\Pi_{10}$:
\begin{ptree}
\AXC{$\stackrel{\Pi_{9}}{\vdots}$}\noLine
\UIC{$\forall x\,\forall y\,(x = y \to y = x)$}
\AXC{$\forall x\,\forall y\,(x = y \to y = x) \to \forall y\,(2 = y \to y = 2)$\lbl{A2}}
\BIC{$\forall y\,(2 = y \to y = 2)$}
\AXC{$\forall y\,(2 = y \to y = 2) \to 2 = 1 + 1 \to 1 + 1 = 2$\lbl{A2}}
\BIC{$2 = 1 + 1 \to 1 + 1 = 2$}
\end{ptree}
\noindent$\Pi_{11}$:
\begin{ptree}
\AXC{$\stackrel{\Pi_{10}}{\vdots}$}\noLine
\UIC{$2 = 1 + 1 \to 1 + 1 = 2$}
\AXC{$2 = 1 + 1$\lbl{Def$2$}}
\AXC{$2 = 1 + 1 \to (2 = 1 + 1 \to 1 + 1 = 2) \to 1 + 1 = 2$\lbl{A1}}
\BIC{$(2 = 1 + 1 \to 1 + 1 = 2) \to 1 + 1 = 2$}
\BIC{$1 + 1 = 2$}
\AXC{$\varphi_{5}$\lbl{A1}}
\BIC{$1 < 1 + 1 \to (1 + 1 = 2 \to 1 < 1 + 1 \to 1 < 2) \to 1 < 2$}
\end{ptree}
\noindent$\Pi_{12}$:
\begin{ptree}
\AXC{$\stackrel{\Pi_{8}}{\vdots}$}\noLine
\UIC{$1 < 1 + 1$}
\AXC{$\stackrel{\Pi_{11}}{\vdots}$}\noLine
\UIC{$1 < 1 + 1 \to (1 + 1 = 2 \to 1 < 1 + 1 \to 1 < 2) \to 1 < 2$}
\BIC{$(1 + 1 = 2 \to 1 < 1 + 1 \to 1 < 2) \to 1 < 2$}
\end{ptree}
\noindent$\Pi_{13}$:
\begin{ptree}
\AXC{$z = w \to 1 < z \to 1 < w$\lbl{A6}}
\UIC{$\forall w\,(z = w \to 1 < z \to 1 < w)$}
\UIC{$\forall z\,\forall w\,(z = w \to 1 < z \to 1 < w)$}
\AXC{$\varphi_{4}$\lbl{A2}}
\BIC{$\forall w\,(1 + 1 = w \to 1 < 1 + 1 \to 1 < w)$}
\AXC{$\varphi_{7}$\lbl{A2}}
\BIC{$1 + 1 = 2 \to 1 < 1 + 1 \to 1 < 2$}
\AXC{$\stackrel{\Pi_{12}}{\vdots}$}\noLine
\UIC{$(1 + 1 = 2 \to 1 < 1 + 1 \to 1 < 2) \to 1 < 2$}
\BIC{$1 < 2$}
\end{ptree}
\noindent$\Pi_{14}$:
\begin{ptree}
\AXC{$\forall x\,\forall y\,\forall z\,(x < y \to y < z \to x < z)$\lbl{[O2]}}
\AXC{$\varphi_{6}$\lbl{A2}}
\BIC{$\forall y\,\forall z\,(0 < y \to y < z \to 0 < z)$}
\AXC{$\varphi_{3}$\lbl{A2}}
\BIC{$\forall z\,(0 < 1 \to 1 < z \to 0 < z)$}
\AXC{$\forall z\,(0 < 1 \to 1 < z \to 0 < z) \to 0 < 1 \to 1 < 2 \to 0 < 2$\lbl{A2}}
\BIC{$0 < 1 \to 1 < 2 \to 0 < 2$}
\end{ptree}
\noindent 主木:
\begin{ptree}
\AXC{$\stackrel{\Pi_{13}}{\vdots}$}\noLine
\UIC{$1 < 2$}
\AXC{$0 < 1$\lbl{[F28]}}
\AXC{$\stackrel{\Pi_{14}}{\vdots}$}\noLine
\UIC{$0 < 1 \to 1 < 2 \to 0 < 2$}
\AXC{$(0 < 1 \to 1 < 2 \to 0 < 2) \to 0 < 1 \to 1 < 2 \to 0 < 2$\lbl{A1}}
\BIC{$0 < 1 \to 1 < 2 \to 0 < 2$}
\BIC{$1 < 2 \to 0 < 2$}
\BIC{$0 < 2$}
\end{ptree}
\begin{abbrevlist}
\item $\varphi_{1} :\equiv 1 + 0 = 1 \to (0 + 1 = 1 + 0 \to 1 + 0 = 1 \to 0 + 1 = 1) \to 0 + 1 = 1$
\item $\varphi_{2} :\equiv \forall x\,\forall y\,\forall z\,(x = y \to y = z \to x = z) \to \forall y\,\forall z\,(0 + 1 = y \to y = z \to 0 + 1 = z)$
\item $\varphi_{3} :\equiv \forall y\,\forall z\,(0 < y \to y < z \to 0 < z) \to \forall z\,(0 < 1 \to 1 < z \to 0 < z)$
\item $\varphi_{4} :\equiv \forall z\,\forall w\,(z = w \to 1 < z \to 1 < w) \to \forall w\,(1 + 1 = w \to 1 < 1 + 1 \to 1 < w)$
\item $\varphi_{5} :\equiv 1 + 1 = 2 \to 1 < 1 + 1 \to (1 + 1 = 2 \to 1 < 1 + 1 \to 1 < 2) \to 1 < 2$
\item $\varphi_{6} :\equiv \forall x\,\forall y\,\forall z\,(x < y \to y < z \to x < z) \to \forall y\,\forall z\,(0 < y \to y < z \to 0 < z)$
\item $\varphi_{7} :\equiv \forall w\,(1 + 1 = w \to 1 < 1 + 1 \to 1 < w) \to 1 + 1 = 2 \to 1 < 1 + 1 \to 1 < 2$
\item $\varphi_{8} :\equiv \forall y\,\forall z\,(0 \in \mathbb{R} \to y \in \mathbb{R} \to z \in \mathbb{R} \to 0 < y \to 0 + z < y + z)$
\item $\varphi_{9} :\equiv \forall y\,\forall z\,(0 + 1 = y \to y = z \to 0 + 1 = z) \to \forall z\,(0 + 1 = 1 + 0 \to 1 + 0 = z \to 0 + 1 = z)$
\item $\varphi_{10} :\equiv \forall x\,\forall y\,\forall z\,(x \in \mathbb{R} \to y \in \mathbb{R} \to z \in \mathbb{R} \to x < y \to x + z < y + z)$
\item $\varphi_{11} :\equiv 0 + 1 < 1 + 1 \to (0 + 1 = 1 \to 0 + 1 < 1 + 1 \to 1 < 1 + 1) \to 1 < 1 + 1$
\item $\varphi_{12} :\equiv \forall z\,(0 + 1 = 1 + 0 \to 1 + 0 = z \to 0 + 1 = z) \to 0 + 1 = 1 + 0 \to 1 + 0 = 1 \to 0 + 1 = 1$
\item $\varphi_{13} :\equiv (0 \in \mathbb{R} \to 1 \in \mathbb{R} \to 0 + 1 = 1 + 0) \to 0 \in \mathbb{R} \to 1 \in \mathbb{R} \to 0 + 1 = 1 + 0$
\item $\varphi_{14} :\equiv \forall z\,\forall w\,(z = w \to z < 1 + 1 \to w < 1 + 1) \to \forall w\,(0 + 1 = w \to 0 + 1 < 1 + 1 \to w < 1 + 1)$
\item $\varphi_{15} :\equiv \forall y\,(0 \in \mathbb{R} \to y \in \mathbb{R} \to 0 + y = y + 0) \to 0 \in \mathbb{R} \to 1 \in \mathbb{R} \to 0 + 1 = 1 + 0$
\item $\varphi_{16} :\equiv \forall w\,(0 + 1 = w \to 0 + 1 < 1 + 1 \to w < 1 + 1) \to 0 + 1 = 1 \to 0 + 1 < 1 + 1 \to 1 < 1 + 1$
\item $\varphi_{17} :\equiv \forall x\,\forall y\,(x \in \mathbb{R} \to y \in \mathbb{R} \to x + y = y + x) \to \forall y\,(0 \in \mathbb{R} \to y \in \mathbb{R} \to 0 + y = y + 0)$
\item $\varphi_{18} :\equiv (0 \in \mathbb{R} \to 1 \in \mathbb{R} \to 1 \in \mathbb{R} \to 0 < 1 \to 0 + 1 < 1 + 1) \to 0 \in \mathbb{R} \to 1 \in \mathbb{R} \to 1 \in \mathbb{R} \to 0 < 1 \to 0 + 1 < 1 + 1$
\item $\varphi_{19} :\equiv \forall z\,(0 \in \mathbb{R} \to 1 \in \mathbb{R} \to z \in \mathbb{R} \to 0 < 1 \to 0 + z < 1 + z) \to 0 \in \mathbb{R} \to 1 \in \mathbb{R} \to 1 \in \mathbb{R} \to 0 < 1 \to 0 + 1 < 1 + 1$
\item $\varphi_{20} :\equiv \varphi_{8} \to \forall z\,(0 \in \mathbb{R} \to 1 \in \mathbb{R} \to z \in \mathbb{R} \to 0 < 1 \to 0 + z < 1 + z)$
\end{abbrevlist}
\end{fullproof}

% 自動生成 (prover/build.js): TwoNe  ノード数 23
\begin{fullproof}{[F30]}{$2 \neq 0$}{23}\phantomsection\label{tree:TwoNe}
\noindent$\Pi_{1}$:
\begin{ptree}
\AXC{$x = x$\lbl{A5}}
\AXC{$x = y \to x = x \to y = x$\lbl{A6}}
\AXC{$(x = y \to x = x \to y = x) \to x = x \to x = y \to y = x$\lbl{A1}}
\BIC{$x = x \to x = y \to y = x$}
\BIC{$x = y \to y = x$}
\UIC{$\forall y\,(x = y \to y = x)$}
\UIC{$\forall x\,\forall y\,(x = y \to y = x)$}
\end{ptree}
\noindent$\Pi_{2}$:
\begin{ptree}
\AXC{$\stackrel{\Pi_{1}}{\vdots}$}\noLine
\UIC{$\forall x\,\forall y\,(x = y \to y = x)$}
\AXC{$\forall x\,\forall y\,(x = y \to y = x) \to \forall y\,(2 = y \to y = 2)$\lbl{A2}}
\BIC{$\forall y\,(2 = y \to y = 2)$}
\AXC{$\forall y\,(2 = y \to y = 2) \to 2 = 0 \to 0 = 2$\lbl{A2}}
\BIC{$2 = 0 \to 0 = 2$}
\end{ptree}
\noindent$\Pi_{3}$:
\begin{ptree}
\AXC{$\forall x\,\forall y\,(x < y \to x \neq y)$\lbl{[F26]}}
\AXC{$\forall x\,\forall y\,(x < y \to x \neq y) \to \forall y\,(0 < y \to 0 \neq y)$\lbl{A2}}
\BIC{$\forall y\,(0 < y \to 0 \neq y)$}
\AXC{$\forall y\,(0 < y \to 0 \neq y) \to 0 < 2 \to 0 \neq 2$\lbl{A2}}
\BIC{$0 < 2 \to 0 \neq 2$}
\end{ptree}
\noindent$\Pi_{4}$:
\begin{ptree}
\AXC{$0 < 2$\lbl{[F29]}}
\AXC{$\stackrel{\Pi_{3}}{\vdots}$}\noLine
\UIC{$0 < 2 \to 0 \neq 2$}
\AXC{$(0 < 2 \to 0 \neq 2) \to 0 < 2 \to 0 \neq 2$\lbl{A1}}
\BIC{$0 < 2 \to 0 \neq 2$}
\BIC{$0 \neq 2$}
\AXC{$0 \neq 2 \to (2 = 0 \to 0 = 2) \to 2 \neq 0$\lbl{A1}}
\BIC{$(2 = 0 \to 0 = 2) \to 2 \neq 0$}
\end{ptree}
\noindent 主木:
\begin{ptree}
\AXC{$\stackrel{\Pi_{2}}{\vdots}$}\noLine
\UIC{$2 = 0 \to 0 = 2$}
\AXC{$\stackrel{\Pi_{4}}{\vdots}$}\noLine
\UIC{$(2 = 0 \to 0 = 2) \to 2 \neq 0$}
\BIC{$2 \neq 0$}
\end{ptree}
\end{fullproof}

% 自動生成 (prover/build.js): InvTwoR  ノード数 9
\begin{fullproof}{[F31]}{$2^{-1} \in \mathbb{R}$}{9}\phantomsection\label{tree:InvTwoR}
\noindent 主木:
\begin{ptree}
\AXC{$2 \neq 0$\lbl{[F30]}}
\AXC{$2 \in \mathbb{R}$\lbl{[F2]}}
\AXC{$\forall x \in \mathbb{R}\,(x \neq 0 \to x^{-1} \in \mathbb{R})$\lbl{(R1)}}
\AXC{$\varphi_{2}$\lbl{A2}}
\BIC{$2 \in \mathbb{R} \to 2 \neq 0 \to 2^{-1} \in \mathbb{R}$}
\AXC{$\varphi_{1}$\lbl{A1}}
\BIC{$2 \in \mathbb{R} \to 2 \neq 0 \to 2^{-1} \in \mathbb{R}$}
\BIC{$2 \neq 0 \to 2^{-1} \in \mathbb{R}$}
\BIC{$2^{-1} \in \mathbb{R}$}
\end{ptree}
\begin{abbrevlist}
\item $\varphi_{1} :\equiv (2 \in \mathbb{R} \to 2 \neq 0 \to 2^{-1} \in \mathbb{R}) \to 2 \in \mathbb{R} \to 2 \neq 0 \to 2^{-1} \in \mathbb{R}$
\item $\varphi_{2} :\equiv \forall x \in \mathbb{R}\,(x \neq 0 \to x^{-1} \in \mathbb{R}) \to 2 \in \mathbb{R} \to 2 \neq 0 \to 2^{-1} \in \mathbb{R}$
\end{abbrevlist}
\end{fullproof}

% 自動生成 (prover/build.js): MulLtPos  ノード数 265
\begin{fullproof}{[F32]}{$\forall x\,\forall y\,\forall z\,(x \in \mathbb{R} \to y \in \mathbb{R} \to z \in \mathbb{R} \to 0 < z \to x < y \to x \cdot z < y \cdot z)$}{265}\phantomsection\label{tree:MulLtPos}
\noindent$\Pi_{1}$:
\begin{ptree}
\AXC{$\forall x\,\forall y\,(x \in \mathbb{R} \to y \in \mathbb{R} \to x < y \to 0 < y - x)$\lbl{[F24]}}
\AXC{$\varphi_{67}$\lbl{A2}}
\BIC{$\forall y\,(x \in \mathbb{R} \to y \in \mathbb{R} \to x < y \to 0 < y - x)$}
\AXC{$\varphi_{63}$\lbl{A2}}
\BIC{$x \in \mathbb{R} \to y \in \mathbb{R} \to x < y \to 0 < y - x$}
\AXC{$(x \in \mathbb{R} \to y \in \mathbb{R} \to x < y \to 0 < y - x) \to \varphi_{98}$\lbl{A1}}
\BIC{$\varphi_{98}$}
\end{ptree}
\noindent$\Pi_{2}$:
\begin{ptree}
\AXC{$x \in \mathbb{R} \to y \in \mathbb{R} \to z \in \mathbb{R} \to 0 < z \to x < y \to y \in \mathbb{R}$\lbl{A1}}
\AXC{$x \in \mathbb{R} \to y \in \mathbb{R} \to z \in \mathbb{R} \to 0 < z \to x < y \to x \in \mathbb{R}$\lbl{A1}}
\AXC{$\stackrel{\Pi_{1}}{\vdots}$}\noLine
\UIC{$\varphi_{98}$}
\BIC{$\varphi_{85}$}
\BIC{$\varphi_{74}$}
\end{ptree}
\noindent$\Pi_{3}$:
\begin{ptree}
\AXC{$\forall x\,\forall y\,(x \in \mathbb{R} \to y \in \mathbb{R} \to x - y \in \mathbb{R})$\lbl{[F1]}}
\AXC{$\varphi_{65}$\lbl{A2}}
\BIC{$\forall y\,(v_{257} \in \mathbb{R} \to y \in \mathbb{R} \to v_{257} - y \in \mathbb{R})$}
\AXC{$\varphi_{68}$\lbl{A2}}
\BIC{$v_{257} \in \mathbb{R} \to v_{258} \in \mathbb{R} \to v_{257} - v_{258} \in \mathbb{R}$}
\UIC{$\forall v_{258}\,(v_{257} \in \mathbb{R} \to v_{258} \in \mathbb{R} \to v_{257} - v_{258} \in \mathbb{R})$}
\UIC{$\varphi_{19}$}
\AXC{$\varphi_{19} \to \forall v_{258}\,(y \in \mathbb{R} \to v_{258} \in \mathbb{R} \to y - v_{258} \in \mathbb{R})$\lbl{A2}}
\BIC{$\forall v_{258}\,(y \in \mathbb{R} \to v_{258} \in \mathbb{R} \to y - v_{258} \in \mathbb{R})$}
\end{ptree}
\noindent$\Pi_{4}$:
\begin{ptree}
\AXC{$\stackrel{\Pi_{3}}{\vdots}$}\noLine
\UIC{$\forall v_{258}\,(y \in \mathbb{R} \to v_{258} \in \mathbb{R} \to y - v_{258} \in \mathbb{R})$}
\AXC{$\varphi_{61}$\lbl{A2}}
\BIC{$y \in \mathbb{R} \to x \in \mathbb{R} \to y - x \in \mathbb{R}$}
\AXC{$(y \in \mathbb{R} \to x \in \mathbb{R} \to y - x \in \mathbb{R}) \to \varphi_{89}$\lbl{A1}}
\BIC{$\varphi_{89}$}
\end{ptree}
\noindent$\Pi_{5}$:
\begin{ptree}
\AXC{$x \in \mathbb{R} \to y \in \mathbb{R} \to z \in \mathbb{R} \to 0 < z \to x < y \to x \in \mathbb{R}$\lbl{A1}}
\AXC{$x \in \mathbb{R} \to y \in \mathbb{R} \to z \in \mathbb{R} \to 0 < z \to x < y \to y \in \mathbb{R}$\lbl{A1}}
\AXC{$\stackrel{\Pi_{4}}{\vdots}$}\noLine
\UIC{$\varphi_{89}$}
\BIC{$\varphi_{78}$}
\BIC{$x \in \mathbb{R} \to y \in \mathbb{R} \to z \in \mathbb{R} \to 0 < z \to x < y \to y - x \in \mathbb{R}$}
\end{ptree}
\noindent$\Pi_{6}$:
\begin{ptree}
\AXC{$\forall x\,\forall y\,(x \in \mathbb{R} \to y \in \mathbb{R} \to 0 < x \to 0 < y \to 0 < x \cdot y)$\lbl{(R15)}}
\AXC{$\varphi_{73}$\lbl{A2}}
\BIC{$\varphi_{24}$}
\AXC{$\varphi_{24} \to \varphi_{40}$\lbl{A2}}
\BIC{$\varphi_{40}$}
\UIC{$\forall v_{256}\,(\varphi_{40})$}
\UIC{$\forall v_{255}\,\forall v_{256}\,(\varphi_{40})$}
\AXC{$\forall v_{255}\,\forall v_{256}\,(\varphi_{40}) \to \varphi_{56}$\lbl{A2}}
\BIC{$\varphi_{56}$}
\end{ptree}
\noindent$\Pi_{7}$:
\begin{ptree}
\AXC{$\stackrel{\Pi_{5}}{\vdots}$}\noLine
\UIC{$x \in \mathbb{R} \to y \in \mathbb{R} \to z \in \mathbb{R} \to 0 < z \to x < y \to y - x \in \mathbb{R}$}
\AXC{$\stackrel{\Pi_{6}}{\vdots}$}\noLine
\UIC{$\varphi_{56}$}
\AXC{$\varphi_{56} \to \varphi_{32}$\lbl{A2}}
\BIC{$\varphi_{32}$}
\AXC{$\varphi_{32} \to \varphi_{107}$\lbl{A1}}
\BIC{$\varphi_{107}$}
\BIC{$\varphi_{101}$}
\end{ptree}
\noindent$\Pi_{8}$:
\begin{ptree}
\AXC{$x \in \mathbb{R} \to y \in \mathbb{R} \to z \in \mathbb{R} \to 0 < z \to x < y \to z \in \mathbb{R}$\lbl{A1}}
\AXC{$\stackrel{\Pi_{7}}{\vdots}$}\noLine
\UIC{$\varphi_{101}$}
\BIC{$\varphi_{86}$}
\end{ptree}
\noindent$\Pi_{9}$:
\begin{ptree}
\AXC{$x \in \mathbb{R} \to y \in \mathbb{R} \to z \in \mathbb{R} \to 0 < z \to x < y \to 0 < z$\lbl{A1}}
\AXC{$x \in \mathbb{R} \to y \in \mathbb{R} \to z \in \mathbb{R} \to 0 < z \to x < y \to x < y$\lbl{A1}}
\AXC{$\stackrel{\Pi_{2}}{\vdots}$}\noLine
\UIC{$\varphi_{74}$}
\BIC{$x \in \mathbb{R} \to y \in \mathbb{R} \to z \in \mathbb{R} \to 0 < z \to x < y \to 0 < y - x$}
\AXC{$\stackrel{\Pi_{8}}{\vdots}$}\noLine
\UIC{$\varphi_{86}$}
\BIC{$(x \in \mathbb{R} \to y \in \mathbb{R} \to z \in \mathbb{R} \to 0 < z \to x < y \to 0 < z) \to \varphi_{41}$}
\BIC{$\varphi_{41}$}
\end{ptree}
\noindent$\Pi_{10}$:
\begin{ptree}
\AXC{$y = z \to x = y \to x = z$\lbl{A6}}
\AXC{$(y = z \to x = y \to x = z) \to x = y \to y = z \to x = z$\lbl{A1}}
\BIC{$x = y \to y = z \to x = z$}
\UIC{$\forall z\,(x = y \to y = z \to x = z)$}
\UIC{$\forall y\,\forall z\,(x = y \to y = z \to x = z)$}
\UIC{$\forall x\,\forall y\,\forall z\,(x = y \to y = z \to x = z)$}
\AXC{$\varphi_{5}$\lbl{A2}}
\BIC{$\forall y\,\forall z\,(v_{269} = y \to y = z \to v_{269} = z)$}
\AXC{$\varphi_{13}$\lbl{A2}}
\BIC{$\forall z\,(v_{269} = v_{270} \to v_{270} = z \to v_{269} = z)$}
\AXC{$\varphi_{21}$\lbl{A2}}
\BIC{$v_{269} = v_{270} \to v_{270} = v_{271} \to v_{269} = v_{271}$}
\UIC{$\forall v_{271}\,(v_{269} = v_{270} \to v_{270} = v_{271} \to v_{269} = v_{271})$}
\UIC{$\varphi_{2}$}
\UIC{$\forall v_{269}\,(\varphi_{2})$}
\end{ptree}
\noindent$\Pi_{11}$:
\begin{ptree}
\AXC{$\stackrel{\Pi_{10}}{\vdots}$}\noLine
\UIC{$\forall v_{269}\,(\varphi_{2})$}
\AXC{$\forall v_{269}\,(\varphi_{2}) \to \varphi_{8}$\lbl{A2}}
\BIC{$\varphi_{8}$}
\AXC{$\varphi_{8} \to \varphi_{47}$\lbl{A2}}
\BIC{$\varphi_{47}$}
\AXC{$\varphi_{47} \to (y - x) \cdot z = y \cdot z - z \cdot x \to \varphi_{17}$\lbl{A2}}
\BIC{$(y - x) \cdot z = y \cdot z - z \cdot x \to \varphi_{17}$}
\end{ptree}
\noindent$\Pi_{12}$:
\begin{ptree}
\AXC{$\forall x\,\forall y\,(x \in \mathbb{R} \to y \in \mathbb{R} \to x \cdot y = y \cdot x)$\lbl{(R7)}}
\AXC{$\varphi_{59}$\lbl{A2}}
\BIC{$\forall y\,(z \in \mathbb{R} \to y \in \mathbb{R} \to z \cdot y = y \cdot z)$}
\AXC{$\varphi_{54}$\lbl{A2}}
\BIC{$z \in \mathbb{R} \to x \in \mathbb{R} \to z \cdot x = x \cdot z$}
\AXC{$(z \in \mathbb{R} \to x \in \mathbb{R} \to z \cdot x = x \cdot z) \to \varphi_{87}$\lbl{A1}}
\BIC{$\varphi_{87}$}
\end{ptree}
\noindent$\Pi_{13}$:
\begin{ptree}
\AXC{$x \in \mathbb{R} \to y \in \mathbb{R} \to z \in \mathbb{R} \to 0 < z \to x < y \to x \in \mathbb{R}$\lbl{A1}}
\AXC{$x \in \mathbb{R} \to y \in \mathbb{R} \to z \in \mathbb{R} \to 0 < z \to x < y \to z \in \mathbb{R}$\lbl{A1}}
\AXC{$\stackrel{\Pi_{12}}{\vdots}$}\noLine
\UIC{$\varphi_{87}$}
\BIC{$\varphi_{75}$}
\BIC{$x \in \mathbb{R} \to y \in \mathbb{R} \to z \in \mathbb{R} \to 0 < z \to x < y \to z \cdot x = x \cdot z$}
\end{ptree}
\noindent$\Pi_{14}$:
\begin{ptree}
\AXC{$\stackrel{\Pi_{13}}{\vdots}$}\noLine
\UIC{$x \in \mathbb{R} \to y \in \mathbb{R} \to z \in \mathbb{R} \to 0 < z \to x < y \to z \cdot x = x \cdot z$}
\AXC{$\varphi_{103}$\lbl{A1}}
\BIC{$\varphi_{92}$}
\end{ptree}
\noindent$\Pi_{15}$:
\begin{ptree}
\AXC{$x = x$\lbl{A5}}
\UIC{$\forall x\,(x = x)$}
\AXC{$\forall x\,(x = x) \to y \cdot z - z \cdot x = y \cdot z - z \cdot x$\lbl{A2}}
\BIC{$y \cdot z - z \cdot x = y \cdot z - z \cdot x$}
\AXC{$\stackrel{\Pi_{14}}{\vdots}$}\noLine
\UIC{$\varphi_{92}$}
\BIC{$(z \cdot x = x \cdot z \to \varphi_{30}) \to x \in \mathbb{R} \to \varphi_{49}$}
\end{ptree}
\noindent$\Pi_{16}$:
\begin{ptree}
\AXC{$\varphi_{25}$\lbl{A6}}
\UIC{$\forall w\,(\varphi_{25})$}
\UIC{$\forall z_{1}\,\forall w\,(\varphi_{25})$}
\AXC{$\forall z_{1}\,\forall w\,(\varphi_{25}) \to \varphi_{51}$\lbl{A2}}
\BIC{$\varphi_{51}$}
\AXC{$\varphi_{51} \to z \cdot x = x \cdot z \to \varphi_{30}$\lbl{A2}}
\BIC{$z \cdot x = x \cdot z \to \varphi_{30}$}
\AXC{$\stackrel{\Pi_{15}}{\vdots}$}\noLine
\UIC{$(z \cdot x = x \cdot z \to \varphi_{30}) \to x \in \mathbb{R} \to \varphi_{49}$}
\BIC{$x \in \mathbb{R} \to \varphi_{49}$}
\end{ptree}
\noindent$\Pi_{17}$:
\begin{ptree}
\AXC{$y = z \to x = y \to x = z$\lbl{A6}}
\AXC{$(y = z \to x = y \to x = z) \to x = y \to y = z \to x = z$\lbl{A1}}
\BIC{$x = y \to y = z \to x = z$}
\UIC{$\forall z\,(x = y \to y = z \to x = z)$}
\UIC{$\forall y\,\forall z\,(x = y \to y = z \to x = z)$}
\UIC{$\forall x\,\forall y\,\forall z\,(x = y \to y = z \to x = z)$}
\AXC{$\varphi_{6}$\lbl{A2}}
\BIC{$\forall y\,\forall z\,(v_{266} = y \to y = z \to v_{266} = z)$}
\AXC{$\varphi_{14}$\lbl{A2}}
\BIC{$\forall z\,(v_{266} = v_{267} \to v_{267} = z \to v_{266} = z)$}
\AXC{$\varphi_{22}$\lbl{A2}}
\BIC{$v_{266} = v_{267} \to v_{267} = v_{268} \to v_{266} = v_{268}$}
\UIC{$\forall v_{268}\,(v_{266} = v_{267} \to v_{267} = v_{268} \to v_{266} = v_{268})$}
\UIC{$\varphi_{3}$}
\UIC{$\forall v_{266}\,(\varphi_{3})$}
\end{ptree}
\noindent$\Pi_{18}$:
\begin{ptree}
\AXC{$\stackrel{\Pi_{17}}{\vdots}$}\noLine
\UIC{$\forall v_{266}\,(\varphi_{3})$}
\AXC{$\forall v_{266}\,(\varphi_{3}) \to \varphi_{9}$\lbl{A2}}
\BIC{$\varphi_{9}$}
\AXC{$\varphi_{9} \to \varphi_{48}$\lbl{A2}}
\BIC{$\varphi_{48}$}
\AXC{$\varphi_{48} \to (y - x) \cdot z = z \cdot y - z \cdot x \to \varphi_{18}$\lbl{A2}}
\BIC{$(y - x) \cdot z = z \cdot y - z \cdot x \to \varphi_{18}$}
\end{ptree}
\noindent$\Pi_{19}$:
\begin{ptree}
\AXC{$\forall x\,\forall y\,(x \in \mathbb{R} \to y \in \mathbb{R} \to x \cdot y = y \cdot x)$\lbl{(R7)}}
\AXC{$\varphi_{59}$\lbl{A2}}
\BIC{$\forall y\,(z \in \mathbb{R} \to y \in \mathbb{R} \to z \cdot y = y \cdot z)$}
\AXC{$\varphi_{55}$\lbl{A2}}
\BIC{$z \in \mathbb{R} \to y \in \mathbb{R} \to z \cdot y = y \cdot z$}
\AXC{$(z \in \mathbb{R} \to y \in \mathbb{R} \to z \cdot y = y \cdot z) \to \varphi_{88}$\lbl{A1}}
\BIC{$\varphi_{88}$}
\end{ptree}
\noindent$\Pi_{20}$:
\begin{ptree}
\AXC{$x \in \mathbb{R} \to y \in \mathbb{R} \to z \in \mathbb{R} \to 0 < z \to x < y \to y \in \mathbb{R}$\lbl{A1}}
\AXC{$x \in \mathbb{R} \to y \in \mathbb{R} \to z \in \mathbb{R} \to 0 < z \to x < y \to z \in \mathbb{R}$\lbl{A1}}
\AXC{$\stackrel{\Pi_{19}}{\vdots}$}\noLine
\UIC{$\varphi_{88}$}
\BIC{$\varphi_{76}$}
\BIC{$x \in \mathbb{R} \to y \in \mathbb{R} \to z \in \mathbb{R} \to 0 < z \to x < y \to z \cdot y = y \cdot z$}
\end{ptree}
\noindent$\Pi_{21}$:
\begin{ptree}
\AXC{$\stackrel{\Pi_{20}}{\vdots}$}\noLine
\UIC{$x \in \mathbb{R} \to y \in \mathbb{R} \to z \in \mathbb{R} \to 0 < z \to x < y \to z \cdot y = y \cdot z$}
\AXC{$\varphi_{104}$\lbl{A1}}
\BIC{$\varphi_{93}$}
\end{ptree}
\noindent$\Pi_{22}$:
\begin{ptree}
\AXC{$x = x$\lbl{A5}}
\UIC{$\forall x\,(x = x)$}
\AXC{$\forall x\,(x = x) \to z \cdot y - z \cdot x = z \cdot y - z \cdot x$\lbl{A2}}
\BIC{$z \cdot y - z \cdot x = z \cdot y - z \cdot x$}
\AXC{$\stackrel{\Pi_{21}}{\vdots}$}\noLine
\UIC{$\varphi_{93}$}
\BIC{$(z \cdot y = y \cdot z \to \varphi_{31}) \to x \in \mathbb{R} \to \varphi_{50}$}
\end{ptree}
\noindent$\Pi_{23}$:
\begin{ptree}
\AXC{$\varphi_{26}$\lbl{A6}}
\UIC{$\forall w\,(\varphi_{26})$}
\UIC{$\forall z_{1}\,\forall w\,(\varphi_{26})$}
\AXC{$\forall z_{1}\,\forall w\,(\varphi_{26}) \to \varphi_{52}$\lbl{A2}}
\BIC{$\varphi_{52}$}
\AXC{$\varphi_{52} \to z \cdot y = y \cdot z \to \varphi_{31}$\lbl{A2}}
\BIC{$z \cdot y = y \cdot z \to \varphi_{31}$}
\AXC{$\stackrel{\Pi_{22}}{\vdots}$}\noLine
\UIC{$(z \cdot y = y \cdot z \to \varphi_{31}) \to x \in \mathbb{R} \to \varphi_{50}$}
\BIC{$x \in \mathbb{R} \to \varphi_{50}$}
\end{ptree}
\noindent$\Pi_{24}$:
\begin{ptree}
\AXC{$y = z \to x = y \to x = z$\lbl{A6}}
\AXC{$(y = z \to x = y \to x = z) \to x = y \to y = z \to x = z$\lbl{A1}}
\BIC{$x = y \to y = z \to x = z$}
\UIC{$\forall z\,(x = y \to y = z \to x = z)$}
\UIC{$\forall y\,\forall z\,(x = y \to y = z \to x = z)$}
\UIC{$\forall x\,\forall y\,\forall z\,(x = y \to y = z \to x = z)$}
\AXC{$\varphi_{7}$\lbl{A2}}
\BIC{$\forall y\,\forall z\,(v_{263} = y \to y = z \to v_{263} = z)$}
\AXC{$\varphi_{15}$\lbl{A2}}
\BIC{$\forall z\,(v_{263} = v_{264} \to v_{264} = z \to v_{263} = z)$}
\AXC{$\varphi_{23}$\lbl{A2}}
\BIC{$v_{263} = v_{264} \to v_{264} = v_{265} \to v_{263} = v_{265}$}
\UIC{$\forall v_{265}\,(v_{263} = v_{264} \to v_{264} = v_{265} \to v_{263} = v_{265})$}
\UIC{$\varphi_{4}$}
\UIC{$\forall v_{263}\,(\varphi_{4})$}
\end{ptree}
\noindent$\Pi_{25}$:
\begin{ptree}
\AXC{$\stackrel{\Pi_{24}}{\vdots}$}\noLine
\UIC{$\forall v_{263}\,(\varphi_{4})$}
\AXC{$\forall v_{263}\,(\varphi_{4}) \to \varphi_{10}$\lbl{A2}}
\BIC{$\varphi_{10}$}
\AXC{$\varphi_{10} \to \varphi_{27}$\lbl{A2}}
\BIC{$\varphi_{27}$}
\AXC{$\varphi_{27} \to (y - x) \cdot z = z \cdot (y - x) \to \varphi_{12}$\lbl{A2}}
\BIC{$(y - x) \cdot z = z \cdot (y - x) \to \varphi_{12}$}
\end{ptree}
\noindent$\Pi_{26}$:
\begin{ptree}
\AXC{$x = x$\lbl{A5}}
\AXC{$x = y \to x = x \to y = x$\lbl{A6}}
\AXC{$(x = y \to x = x \to y = x) \to x = x \to x = y \to y = x$\lbl{A1}}
\BIC{$x = x \to x = y \to y = x$}
\BIC{$x = y \to y = x$}
\UIC{$\forall y\,(x = y \to y = x)$}
\UIC{$\forall x\,\forall y\,(x = y \to y = x)$}
\end{ptree}
\noindent$\Pi_{27}$:
\begin{ptree}
\AXC{$\stackrel{\Pi_{26}}{\vdots}$}\noLine
\UIC{$\forall x\,\forall y\,(x = y \to y = x)$}
\AXC{$\forall x\,\forall y\,(x = y \to y = x) \to \forall y\,(v_{261} = y \to y = v_{261})$\lbl{A2}}
\BIC{$\forall y\,(v_{261} = y \to y = v_{261})$}
\AXC{$\varphi_{1}$\lbl{A2}}
\BIC{$v_{261} = v_{262} \to v_{262} = v_{261}$}
\UIC{$\forall v_{262}\,(v_{261} = v_{262} \to v_{262} = v_{261})$}
\UIC{$\forall v_{261}\,\forall v_{262}\,(v_{261} = v_{262} \to v_{262} = v_{261})$}
\AXC{$\varphi_{57}$\lbl{A2}}
\BIC{$\forall v_{262}\,(z \cdot y - z \cdot x = v_{262} \to v_{262} = z \cdot y - z \cdot x)$}
\AXC{$\varphi_{71}$\lbl{A2}}
\BIC{$\varphi_{11}$}
\end{ptree}
\noindent$\Pi_{28}$:
\begin{ptree}
\AXC{$x \in \mathbb{R} \to y \in \mathbb{R} \to z \in \mathbb{R} \to 0 < z \to x < y \to z \in \mathbb{R}$\lbl{A1}}
\AXC{$\varphi_{43}$\lbl{[F17]}}
\AXC{$\varphi_{43} \to \varphi_{42}$\lbl{A2}}
\BIC{$\varphi_{42}$}
\AXC{$\varphi_{42} \to \varphi_{33}$\lbl{A2}}
\BIC{$\varphi_{33}$}
\AXC{$\varphi_{77}$\lbl{A2}}
\BIC{$z \in \mathbb{R} \to y \in \mathbb{R} \to x \in \mathbb{R} \to z \cdot y - z \cdot x = z \cdot (y - x)$}
\AXC{$\varphi_{108}$\lbl{A1}}
\BIC{$\varphi_{105}$}
\BIC{$\varphi_{94}$}
\end{ptree}
\noindent$\Pi_{29}$:
\begin{ptree}
\AXC{$x \in \mathbb{R} \to y \in \mathbb{R} \to z \in \mathbb{R} \to 0 < z \to x < y \to x \in \mathbb{R}$\lbl{A1}}
\AXC{$x \in \mathbb{R} \to y \in \mathbb{R} \to z \in \mathbb{R} \to 0 < z \to x < y \to y \in \mathbb{R}$\lbl{A1}}
\AXC{$\stackrel{\Pi_{28}}{\vdots}$}\noLine
\UIC{$\varphi_{94}$}
\BIC{$\varphi_{84}$}
\BIC{$x \in \mathbb{R} \to \varphi_{35}$}
\end{ptree}
\noindent$\Pi_{30}$:
\begin{ptree}
\AXC{$\forall x\,\forall y\,(x \in \mathbb{R} \to y \in \mathbb{R} \to x \cdot y = y \cdot x)$\lbl{(R7)}}
\AXC{$\varphi_{64}$\lbl{A2}}
\BIC{$\forall y\,(v_{259} \in \mathbb{R} \to y \in \mathbb{R} \to v_{259} \cdot y = y \cdot v_{259})$}
\AXC{$\varphi_{70}$\lbl{A2}}
\BIC{$v_{259} \in \mathbb{R} \to v_{260} \in \mathbb{R} \to v_{259} \cdot v_{260} = v_{260} \cdot v_{259}$}
\UIC{$\varphi_{16}$}
\UIC{$\forall v_{259}\,(\varphi_{16})$}
\AXC{$\forall v_{259}\,(\varphi_{16}) \to \varphi_{29}$\lbl{A2}}
\BIC{$\varphi_{29}$}
\end{ptree}
\noindent$\Pi_{31}$:
\begin{ptree}
\AXC{$\stackrel{\Pi_{30}}{\vdots}$}\noLine
\UIC{$\varphi_{29}$}
\AXC{$\varphi_{29} \to y - x \in \mathbb{R} \to z \in \mathbb{R} \to (y - x) \cdot z = z \cdot (y - x)$\lbl{A2}}
\BIC{$y - x \in \mathbb{R} \to z \in \mathbb{R} \to (y - x) \cdot z = z \cdot (y - x)$}
\AXC{$\varphi_{102}$\lbl{A1}}
\BIC{$\varphi_{95}$}
\end{ptree}
\noindent$\Pi_{32}$:
\begin{ptree}
\AXC{$x \in \mathbb{R} \to y \in \mathbb{R} \to z \in \mathbb{R} \to 0 < z \to x < y \to z \in \mathbb{R}$\lbl{A1}}
\AXC{$\stackrel{\Pi_{5}}{\vdots}$}\noLine
\UIC{$x \in \mathbb{R} \to y \in \mathbb{R} \to z \in \mathbb{R} \to 0 < z \to x < y \to y - x \in \mathbb{R}$}
\AXC{$\stackrel{\Pi_{31}}{\vdots}$}\noLine
\UIC{$\varphi_{95}$}
\BIC{$\varphi_{82}$}
\BIC{$x \in \mathbb{R} \to \varphi_{28}$}
\end{ptree}
\noindent$\Pi_{33}$:
\begin{ptree}
\AXC{$\stackrel{\Pi_{27}}{\vdots}$}\noLine
\UIC{$\varphi_{11}$}
\AXC{$\stackrel{\Pi_{29}}{\vdots}$}\noLine
\UIC{$x \in \mathbb{R} \to \varphi_{35}$}
\AXC{$(x \in \mathbb{R} \to \varphi_{35}) \to \varphi_{11} \to x \in \mathbb{R} \to \varphi_{36}$\lbl{A1}}
\BIC{$\varphi_{11} \to x \in \mathbb{R} \to \varphi_{36}$}
\BIC{$x \in \mathbb{R} \to \varphi_{36}$}
\AXC{$\stackrel{\Pi_{32}}{\vdots}$}\noLine
\UIC{$x \in \mathbb{R} \to \varphi_{28}$}
\AXC{$(x \in \mathbb{R} \to \varphi_{28}) \to \varphi_{97}$\lbl{A1}}
\BIC{$\varphi_{97}$}
\BIC{$((y - x) \cdot z = z \cdot (y - x) \to \varphi_{12}) \to x \in \mathbb{R} \to \varphi_{37}$}
\end{ptree}
\noindent$\Pi_{34}$:
\begin{ptree}
\AXC{$\stackrel{\Pi_{25}}{\vdots}$}\noLine
\UIC{$(y - x) \cdot z = z \cdot (y - x) \to \varphi_{12}$}
\AXC{$\stackrel{\Pi_{33}}{\vdots}$}\noLine
\UIC{$((y - x) \cdot z = z \cdot (y - x) \to \varphi_{12}) \to x \in \mathbb{R} \to \varphi_{37}$}
\BIC{$x \in \mathbb{R} \to \varphi_{37}$}
\AXC{$(x \in \mathbb{R} \to \varphi_{37}) \to \varphi_{99}$\lbl{A1}}
\BIC{$\varphi_{99}$}
\end{ptree}
\noindent$\Pi_{35}$:
\begin{ptree}
\AXC{$\stackrel{\Pi_{18}}{\vdots}$}\noLine
\UIC{$(y - x) \cdot z = z \cdot y - z \cdot x \to \varphi_{18}$}
\AXC{$\stackrel{\Pi_{23}}{\vdots}$}\noLine
\UIC{$x \in \mathbb{R} \to \varphi_{50}$}
\AXC{$\stackrel{\Pi_{34}}{\vdots}$}\noLine
\UIC{$\varphi_{99}$}
\BIC{$((y - x) \cdot z = z \cdot y - z \cdot x \to \varphi_{18}) \to x \in \mathbb{R} \to \varphi_{38}$}
\BIC{$x \in \mathbb{R} \to \varphi_{38}$}
\AXC{$(x \in \mathbb{R} \to \varphi_{38}) \to \varphi_{100}$\lbl{A1}}
\BIC{$\varphi_{100}$}
\end{ptree}
\noindent$\Pi_{36}$:
\begin{ptree}
\AXC{$\stackrel{\Pi_{9}}{\vdots}$}\noLine
\UIC{$\varphi_{41}$}
\AXC{$\stackrel{\Pi_{11}}{\vdots}$}\noLine
\UIC{$(y - x) \cdot z = y \cdot z - z \cdot x \to \varphi_{17}$}
\AXC{$\stackrel{\Pi_{16}}{\vdots}$}\noLine
\UIC{$x \in \mathbb{R} \to \varphi_{49}$}
\AXC{$\stackrel{\Pi_{35}}{\vdots}$}\noLine
\UIC{$\varphi_{100}$}
\BIC{$((y - x) \cdot z = y \cdot z - z \cdot x \to \varphi_{17}) \to x \in \mathbb{R} \to \varphi_{39}$}
\BIC{$x \in \mathbb{R} \to \varphi_{39}$}
\AXC{$(x \in \mathbb{R} \to \varphi_{39}) \to \varphi_{41} \to \varphi_{34} \to \varphi_{53}$\lbl{A1}}
\BIC{$\varphi_{41} \to \varphi_{34} \to \varphi_{53}$}
\BIC{$\varphi_{34} \to \varphi_{53}$}
\end{ptree}
\noindent$\Pi_{37}$:
\begin{ptree}
\AXC{$z_{1} = w \to 0 < z_{1} \to 0 < w$\lbl{A6}}
\UIC{$\forall w\,(z_{1} = w \to 0 < z_{1} \to 0 < w)$}
\UIC{$\forall z_{1}\,\forall w\,(z_{1} = w \to 0 < z_{1} \to 0 < w)$}
\AXC{$\varphi_{46}$\lbl{A2}}
\BIC{$\forall w\,((y - x) \cdot z = w \to 0 < (y - x) \cdot z \to 0 < w)$}
\AXC{$\forall w\,((y - x) \cdot z = w \to 0 < (y - x) \cdot z \to 0 < w) \to \varphi_{34}$\lbl{A2}}
\BIC{$\varphi_{34}$}
\AXC{$\stackrel{\Pi_{36}}{\vdots}$}\noLine
\UIC{$\varphi_{34} \to \varphi_{53}$}
\BIC{$\varphi_{53}$}
\end{ptree}
\noindent$\Pi_{38}$:
\begin{ptree}
\AXC{$\forall x\,\forall y\,(x \in \mathbb{R} \to y \in \mathbb{R} \to x \cdot y \in \mathbb{R})$\lbl{(R1)}}
\AXC{$\varphi_{66}$\lbl{A2}}
\BIC{$\forall y\,(v_{272} \in \mathbb{R} \to y \in \mathbb{R} \to v_{272} \cdot y \in \mathbb{R})$}
\AXC{$\varphi_{69}$\lbl{A2}}
\BIC{$v_{272} \in \mathbb{R} \to v_{273} \in \mathbb{R} \to v_{272} \cdot v_{273} \in \mathbb{R}$}
\UIC{$\forall v_{273}\,(v_{272} \in \mathbb{R} \to v_{273} \in \mathbb{R} \to v_{272} \cdot v_{273} \in \mathbb{R})$}
\UIC{$\varphi_{20}$}
\AXC{$\varphi_{20} \to \forall v_{273}\,(y \in \mathbb{R} \to v_{273} \in \mathbb{R} \to y \cdot v_{273} \in \mathbb{R})$\lbl{A2}}
\BIC{$\forall v_{273}\,(y \in \mathbb{R} \to v_{273} \in \mathbb{R} \to y \cdot v_{273} \in \mathbb{R})$}
\end{ptree}
\noindent$\Pi_{39}$:
\begin{ptree}
\AXC{$\stackrel{\Pi_{38}}{\vdots}$}\noLine
\UIC{$\forall v_{273}\,(y \in \mathbb{R} \to v_{273} \in \mathbb{R} \to y \cdot v_{273} \in \mathbb{R})$}
\AXC{$\varphi_{62}$\lbl{A2}}
\BIC{$y \in \mathbb{R} \to z \in \mathbb{R} \to y \cdot z \in \mathbb{R}$}
\AXC{$(y \in \mathbb{R} \to z \in \mathbb{R} \to y \cdot z \in \mathbb{R}) \to \varphi_{90}$\lbl{A1}}
\BIC{$\varphi_{90}$}
\end{ptree}
\noindent$\Pi_{40}$:
\begin{ptree}
\AXC{$x \in \mathbb{R} \to y \in \mathbb{R} \to z \in \mathbb{R} \to 0 < z \to x < y \to z \in \mathbb{R}$\lbl{A1}}
\AXC{$x \in \mathbb{R} \to y \in \mathbb{R} \to z \in \mathbb{R} \to 0 < z \to x < y \to y \in \mathbb{R}$\lbl{A1}}
\AXC{$\stackrel{\Pi_{39}}{\vdots}$}\noLine
\UIC{$\varphi_{90}$}
\BIC{$\varphi_{79}$}
\BIC{$x \in \mathbb{R} \to y \in \mathbb{R} \to z \in \mathbb{R} \to 0 < z \to x < y \to y \cdot z \in \mathbb{R}$}
\end{ptree}
\noindent$\Pi_{41}$:
\begin{ptree}
\AXC{$\forall x\,\forall y\,(x \in \mathbb{R} \to y \in \mathbb{R} \to x \cdot y \in \mathbb{R})$\lbl{(R1)}}
\AXC{$\varphi_{60}$\lbl{A2}}
\BIC{$\forall y\,(x \in \mathbb{R} \to y \in \mathbb{R} \to x \cdot y \in \mathbb{R})$}
\AXC{$\varphi_{58}$\lbl{A2}}
\BIC{$x \in \mathbb{R} \to z \in \mathbb{R} \to x \cdot z \in \mathbb{R}$}
\AXC{$(x \in \mathbb{R} \to z \in \mathbb{R} \to x \cdot z \in \mathbb{R}) \to \varphi_{91}$\lbl{A1}}
\BIC{$\varphi_{91}$}
\end{ptree}
\noindent$\Pi_{42}$:
\begin{ptree}
\AXC{$x \in \mathbb{R} \to y \in \mathbb{R} \to z \in \mathbb{R} \to 0 < z \to x < y \to z \in \mathbb{R}$\lbl{A1}}
\AXC{$x \in \mathbb{R} \to y \in \mathbb{R} \to z \in \mathbb{R} \to 0 < z \to x < y \to x \in \mathbb{R}$\lbl{A1}}
\AXC{$\stackrel{\Pi_{41}}{\vdots}$}\noLine
\UIC{$\varphi_{91}$}
\BIC{$\varphi_{80}$}
\BIC{$x \in \mathbb{R} \to y \in \mathbb{R} \to z \in \mathbb{R} \to 0 < z \to x < y \to x \cdot z \in \mathbb{R}$}
\end{ptree}
\noindent$\Pi_{43}$:
\begin{ptree}
\AXC{$\forall x\,\forall y\,(x \in \mathbb{R} \to y \in \mathbb{R} \to 0 < y - x \to x < y)$\lbl{[F25]}}
\AXC{$\varphi_{72}$\lbl{A2}}
\BIC{$\forall y\,(x \cdot z \in \mathbb{R} \to y \in \mathbb{R} \to 0 < y - x \cdot z \to x \cdot z < y)$}
\AXC{$\varphi_{81}$\lbl{A2}}
\BIC{$\varphi_{44}$}
\AXC{$\varphi_{44} \to \varphi_{106}$\lbl{A1}}
\BIC{$\varphi_{106}$}
\end{ptree}
\noindent$\Pi_{44}$:
\begin{ptree}
\AXC{$\stackrel{\Pi_{40}}{\vdots}$}\noLine
\UIC{$x \in \mathbb{R} \to y \in \mathbb{R} \to z \in \mathbb{R} \to 0 < z \to x < y \to y \cdot z \in \mathbb{R}$}
\AXC{$\stackrel{\Pi_{42}}{\vdots}$}\noLine
\UIC{$x \in \mathbb{R} \to y \in \mathbb{R} \to z \in \mathbb{R} \to 0 < z \to x < y \to x \cdot z \in \mathbb{R}$}
\AXC{$\stackrel{\Pi_{43}}{\vdots}$}\noLine
\UIC{$\varphi_{106}$}
\BIC{$\varphi_{96}$}
\BIC{$\varphi_{83}$}
\end{ptree}
\noindent 主木:
\begin{ptree}
\AXC{$\stackrel{\Pi_{37}}{\vdots}$}\noLine
\UIC{$\varphi_{53}$}
\AXC{$\stackrel{\Pi_{44}}{\vdots}$}\noLine
\UIC{$\varphi_{83}$}
\BIC{$x \in \mathbb{R} \to y \in \mathbb{R} \to z \in \mathbb{R} \to 0 < z \to x < y \to x \cdot z < y \cdot z$}
\UIC{$\varphi_{45}$}
\UIC{$\forall y\,(\varphi_{45})$}
\UIC{$\forall x\,\forall y\,(\varphi_{45})$}
\end{ptree}
\begin{abbrevlist}
\item $\varphi_{1} :\equiv \forall y\,(v_{261} = y \to y = v_{261}) \to v_{261} = v_{262} \to v_{262} = v_{261}$
\item $\varphi_{2} :\equiv \forall v_{270}\,\forall v_{271}\,(v_{269} = v_{270} \to v_{270} = v_{271} \to v_{269} = v_{271})$
\item $\varphi_{3} :\equiv \forall v_{267}\,\forall v_{268}\,(v_{266} = v_{267} \to v_{267} = v_{268} \to v_{266} = v_{268})$
\item $\varphi_{4} :\equiv \forall v_{264}\,\forall v_{265}\,(v_{263} = v_{264} \to v_{264} = v_{265} \to v_{263} = v_{265})$
\item $\varphi_{5} :\equiv \forall x\,\forall y\,\forall z\,(x = y \to y = z \to x = z) \to \forall y\,\forall z\,(v_{269} = y \to y = z \to v_{269} = z)$
\item $\varphi_{6} :\equiv \forall x\,\forall y\,\forall z\,(x = y \to y = z \to x = z) \to \forall y\,\forall z\,(v_{266} = y \to y = z \to v_{266} = z)$
\item $\varphi_{7} :\equiv \forall x\,\forall y\,\forall z\,(x = y \to y = z \to x = z) \to \forall y\,\forall z\,(v_{263} = y \to y = z \to v_{263} = z)$
\item $\varphi_{8} :\equiv \forall v_{270}\,\forall v_{271}\,((y - x) \cdot z = v_{270} \to v_{270} = v_{271} \to (y - x) \cdot z = v_{271})$
\item $\varphi_{9} :\equiv \forall v_{267}\,\forall v_{268}\,((y - x) \cdot z = v_{267} \to v_{267} = v_{268} \to (y - x) \cdot z = v_{268})$
\item $\varphi_{10} :\equiv \forall v_{264}\,\forall v_{265}\,((y - x) \cdot z = v_{264} \to v_{264} = v_{265} \to (y - x) \cdot z = v_{265})$
\item $\varphi_{11} :\equiv z \cdot y - z \cdot x = z \cdot (y - x) \to z \cdot (y - x) = z \cdot y - z \cdot x$
\item $\varphi_{12} :\equiv z \cdot (y - x) = z \cdot y - z \cdot x \to (y - x) \cdot z = z \cdot y - z \cdot x$
\item $\varphi_{13} :\equiv \forall y\,\forall z\,(v_{269} = y \to y = z \to v_{269} = z) \to \forall z\,(v_{269} = v_{270} \to v_{270} = z \to v_{269} = z)$
\item $\varphi_{14} :\equiv \forall y\,\forall z\,(v_{266} = y \to y = z \to v_{266} = z) \to \forall z\,(v_{266} = v_{267} \to v_{267} = z \to v_{266} = z)$
\item $\varphi_{15} :\equiv \forall y\,\forall z\,(v_{263} = y \to y = z \to v_{263} = z) \to \forall z\,(v_{263} = v_{264} \to v_{264} = z \to v_{263} = z)$
\item $\varphi_{16} :\equiv \forall v_{260}\,(v_{259} \in \mathbb{R} \to v_{260} \in \mathbb{R} \to v_{259} \cdot v_{260} = v_{260} \cdot v_{259})$
\item $\varphi_{17} :\equiv y \cdot z - z \cdot x = y \cdot z - x \cdot z \to (y - x) \cdot z = y \cdot z - x \cdot z$
\item $\varphi_{18} :\equiv z \cdot y - z \cdot x = y \cdot z - z \cdot x \to (y - x) \cdot z = y \cdot z - z \cdot x$
\item $\varphi_{19} :\equiv \forall v_{257}\,\forall v_{258}\,(v_{257} \in \mathbb{R} \to v_{258} \in \mathbb{R} \to v_{257} - v_{258} \in \mathbb{R})$
\item $\varphi_{20} :\equiv \forall v_{272}\,\forall v_{273}\,(v_{272} \in \mathbb{R} \to v_{273} \in \mathbb{R} \to v_{272} \cdot v_{273} \in \mathbb{R})$
\item $\varphi_{21} :\equiv \forall z\,(v_{269} = v_{270} \to v_{270} = z \to v_{269} = z) \to v_{269} = v_{270} \to v_{270} = v_{271} \to v_{269} = v_{271}$
\item $\varphi_{22} :\equiv \forall z\,(v_{266} = v_{267} \to v_{267} = z \to v_{266} = z) \to v_{266} = v_{267} \to v_{267} = v_{268} \to v_{266} = v_{268}$
\item $\varphi_{23} :\equiv \forall z\,(v_{263} = v_{264} \to v_{264} = z \to v_{263} = z) \to v_{263} = v_{264} \to v_{264} = v_{265} \to v_{263} = v_{265}$
\item $\varphi_{24} :\equiv \forall y\,(v_{255} \in \mathbb{R} \to y \in \mathbb{R} \to 0 < v_{255} \to 0 < y \to 0 < v_{255} \cdot y)$
\item $\varphi_{25} :\equiv z_{1} = w \to y \cdot z - z \cdot x = y \cdot z - z_{1} \to y \cdot z - z \cdot x = y \cdot z - w$
\item $\varphi_{26} :\equiv z_{1} = w \to z \cdot y - z \cdot x = z_{1} - z \cdot x \to z \cdot y - z \cdot x = w - z \cdot x$
\item $\varphi_{27} :\equiv \forall v_{265}\,((y - x) \cdot z = z \cdot (y - x) \to z \cdot (y - x) = v_{265} \to (y - x) \cdot z = v_{265})$
\item $\varphi_{28} :\equiv y \in \mathbb{R} \to z \in \mathbb{R} \to 0 < z \to x < y \to (y - x) \cdot z = z \cdot (y - x)$
\item $\varphi_{29} :\equiv \forall v_{260}\,(y - x \in \mathbb{R} \to v_{260} \in \mathbb{R} \to (y - x) \cdot v_{260} = v_{260} \cdot (y - x))$
\item $\varphi_{30} :\equiv y \cdot z - z \cdot x = y \cdot z - z \cdot x \to y \cdot z - z \cdot x = y \cdot z - x \cdot z$
\item $\varphi_{31} :\equiv z \cdot y - z \cdot x = z \cdot y - z \cdot x \to z \cdot y - z \cdot x = y \cdot z - z \cdot x$
\item $\varphi_{32} :\equiv y - x \in \mathbb{R} \to z \in \mathbb{R} \to 0 < y - x \to 0 < z \to 0 < (y - x) \cdot z$
\item $\varphi_{33} :\equiv \forall a\,(z \in \mathbb{R} \to y \in \mathbb{R} \to a \in \mathbb{R} \to z \cdot y - z \cdot a = z \cdot (y - a))$
\item $\varphi_{34} :\equiv (y - x) \cdot z = y \cdot z - x \cdot z \to 0 < (y - x) \cdot z \to 0 < y \cdot z - x \cdot z$
\item $\varphi_{35} :\equiv y \in \mathbb{R} \to z \in \mathbb{R} \to 0 < z \to x < y \to z \cdot y - z \cdot x = z \cdot (y - x)$
\item $\varphi_{36} :\equiv y \in \mathbb{R} \to z \in \mathbb{R} \to 0 < z \to x < y \to z \cdot (y - x) = z \cdot y - z \cdot x$
\item $\varphi_{37} :\equiv y \in \mathbb{R} \to z \in \mathbb{R} \to 0 < z \to x < y \to (y - x) \cdot z = z \cdot y - z \cdot x$
\item $\varphi_{38} :\equiv y \in \mathbb{R} \to z \in \mathbb{R} \to 0 < z \to x < y \to (y - x) \cdot z = y \cdot z - z \cdot x$
\item $\varphi_{39} :\equiv y \in \mathbb{R} \to z \in \mathbb{R} \to 0 < z \to x < y \to (y - x) \cdot z = y \cdot z - x \cdot z$
\item $\varphi_{40} :\equiv v_{255} \in \mathbb{R} \to v_{256} \in \mathbb{R} \to 0 < v_{255} \to 0 < v_{256} \to 0 < v_{255} \cdot v_{256}$
\item $\varphi_{41} :\equiv x \in \mathbb{R} \to y \in \mathbb{R} \to z \in \mathbb{R} \to 0 < z \to x < y \to 0 < (y - x) \cdot z$
\item $\varphi_{42} :\equiv \forall x\,\forall a\,(z \in \mathbb{R} \to x \in \mathbb{R} \to a \in \mathbb{R} \to z \cdot x - z \cdot a = z \cdot (x - a))$
\item $\varphi_{43} :\equiv \forall c\,\forall x\,\forall a\,(c \in \mathbb{R} \to x \in \mathbb{R} \to a \in \mathbb{R} \to c \cdot x - c \cdot a = c \cdot (x - a))$
\item $\varphi_{44} :\equiv x \cdot z \in \mathbb{R} \to y \cdot z \in \mathbb{R} \to 0 < y \cdot z - x \cdot z \to x \cdot z < y \cdot z$
\item $\varphi_{45} :\equiv \forall z\,(x \in \mathbb{R} \to y \in \mathbb{R} \to z \in \mathbb{R} \to 0 < z \to x < y \to x \cdot z < y \cdot z)$
\item $\varphi_{46} :\equiv \forall z_{1}\,\forall w\,(z_{1} = w \to 0 < z_{1} \to 0 < w) \to \forall w\,((y - x) \cdot z = w \to 0 < (y - x) \cdot z \to 0 < w)$
\item $\varphi_{47} :\equiv \forall v_{271}\,((y - x) \cdot z = y \cdot z - z \cdot x \to y \cdot z - z \cdot x = v_{271} \to (y - x) \cdot z = v_{271})$
\item $\varphi_{48} :\equiv \forall v_{268}\,((y - x) \cdot z = z \cdot y - z \cdot x \to z \cdot y - z \cdot x = v_{268} \to (y - x) \cdot z = v_{268})$
\item $\varphi_{49} :\equiv y \in \mathbb{R} \to z \in \mathbb{R} \to 0 < z \to x < y \to y \cdot z - z \cdot x = y \cdot z - x \cdot z$
\item $\varphi_{50} :\equiv y \in \mathbb{R} \to z \in \mathbb{R} \to 0 < z \to x < y \to z \cdot y - z \cdot x = y \cdot z - z \cdot x$
\item $\varphi_{51} :\equiv \forall w\,(z \cdot x = w \to y \cdot z - z \cdot x = y \cdot z - z \cdot x \to y \cdot z - z \cdot x = y \cdot z - w)$
\item $\varphi_{52} :\equiv \forall w\,(z \cdot y = w \to z \cdot y - z \cdot x = z \cdot y - z \cdot x \to z \cdot y - z \cdot x = w - z \cdot x)$
\item $\varphi_{53} :\equiv x \in \mathbb{R} \to y \in \mathbb{R} \to z \in \mathbb{R} \to 0 < z \to x < y \to 0 < y \cdot z - x \cdot z$
\item $\varphi_{54} :\equiv \forall y\,(z \in \mathbb{R} \to y \in \mathbb{R} \to z \cdot y = y \cdot z) \to z \in \mathbb{R} \to x \in \mathbb{R} \to z \cdot x = x \cdot z$
\item $\varphi_{55} :\equiv \forall y\,(z \in \mathbb{R} \to y \in \mathbb{R} \to z \cdot y = y \cdot z) \to z \in \mathbb{R} \to y \in \mathbb{R} \to z \cdot y = y \cdot z$
\item $\varphi_{56} :\equiv \forall v_{256}\,(y - x \in \mathbb{R} \to v_{256} \in \mathbb{R} \to 0 < y - x \to 0 < v_{256} \to 0 < (y - x) \cdot v_{256})$
\item $\varphi_{57} :\equiv \forall v_{261}\,\forall v_{262}\,(v_{261} = v_{262} \to v_{262} = v_{261}) \to \forall v_{262}\,(z \cdot y - z \cdot x = v_{262} \to v_{262} = z \cdot y - z \cdot x)$
\item $\varphi_{58} :\equiv \forall y\,(x \in \mathbb{R} \to y \in \mathbb{R} \to x \cdot y \in \mathbb{R}) \to x \in \mathbb{R} \to z \in \mathbb{R} \to x \cdot z \in \mathbb{R}$
\item $\varphi_{59} :\equiv \forall x\,\forall y\,(x \in \mathbb{R} \to y \in \mathbb{R} \to x \cdot y = y \cdot x) \to \forall y\,(z \in \mathbb{R} \to y \in \mathbb{R} \to z \cdot y = y \cdot z)$
\item $\varphi_{60} :\equiv \forall x\,\forall y\,(x \in \mathbb{R} \to y \in \mathbb{R} \to x \cdot y \in \mathbb{R}) \to \forall y\,(x \in \mathbb{R} \to y \in \mathbb{R} \to x \cdot y \in \mathbb{R})$
\item $\varphi_{61} :\equiv \forall v_{258}\,(y \in \mathbb{R} \to v_{258} \in \mathbb{R} \to y - v_{258} \in \mathbb{R}) \to y \in \mathbb{R} \to x \in \mathbb{R} \to y - x \in \mathbb{R}$
\item $\varphi_{62} :\equiv \forall v_{273}\,(y \in \mathbb{R} \to v_{273} \in \mathbb{R} \to y \cdot v_{273} \in \mathbb{R}) \to y \in \mathbb{R} \to z \in \mathbb{R} \to y \cdot z \in \mathbb{R}$
\item $\varphi_{63} :\equiv \forall y\,(x \in \mathbb{R} \to y \in \mathbb{R} \to x < y \to 0 < y - x) \to x \in \mathbb{R} \to y \in \mathbb{R} \to x < y \to 0 < y - x$
\item $\varphi_{64} :\equiv \forall x\,\forall y\,(x \in \mathbb{R} \to y \in \mathbb{R} \to x \cdot y = y \cdot x) \to \forall y\,(v_{259} \in \mathbb{R} \to y \in \mathbb{R} \to v_{259} \cdot y = y \cdot v_{259})$
\item $\varphi_{65} :\equiv \forall x\,\forall y\,(x \in \mathbb{R} \to y \in \mathbb{R} \to x - y \in \mathbb{R}) \to \forall y\,(v_{257} \in \mathbb{R} \to y \in \mathbb{R} \to v_{257} - y \in \mathbb{R})$
\item $\varphi_{66} :\equiv \forall x\,\forall y\,(x \in \mathbb{R} \to y \in \mathbb{R} \to x \cdot y \in \mathbb{R}) \to \forall y\,(v_{272} \in \mathbb{R} \to y \in \mathbb{R} \to v_{272} \cdot y \in \mathbb{R})$
\item $\varphi_{67} :\equiv \forall x\,\forall y\,(x \in \mathbb{R} \to y \in \mathbb{R} \to x < y \to 0 < y - x) \to \forall y\,(x \in \mathbb{R} \to y \in \mathbb{R} \to x < y \to 0 < y - x)$
\item $\varphi_{68} :\equiv \forall y\,(v_{257} \in \mathbb{R} \to y \in \mathbb{R} \to v_{257} - y \in \mathbb{R}) \to v_{257} \in \mathbb{R} \to v_{258} \in \mathbb{R} \to v_{257} - v_{258} \in \mathbb{R}$
\item $\varphi_{69} :\equiv \forall y\,(v_{272} \in \mathbb{R} \to y \in \mathbb{R} \to v_{272} \cdot y \in \mathbb{R}) \to v_{272} \in \mathbb{R} \to v_{273} \in \mathbb{R} \to v_{272} \cdot v_{273} \in \mathbb{R}$
\item $\varphi_{70} :\equiv \forall y\,(v_{259} \in \mathbb{R} \to y \in \mathbb{R} \to v_{259} \cdot y = y \cdot v_{259}) \to v_{259} \in \mathbb{R} \to v_{260} \in \mathbb{R} \to v_{259} \cdot v_{260} = v_{260} \cdot v_{259}$
\item $\varphi_{71} :\equiv \forall v_{262}\,(z \cdot y - z \cdot x = v_{262} \to v_{262} = z \cdot y - z \cdot x) \to \varphi_{11}$
\item $\varphi_{72} :\equiv \forall x\,\forall y\,(x \in \mathbb{R} \to y \in \mathbb{R} \to 0 < y - x \to x < y) \to \forall y\,(x \cdot z \in \mathbb{R} \to y \in \mathbb{R} \to 0 < y - x \cdot z \to x \cdot z < y)$
\item $\varphi_{73} :\equiv \forall x\,\forall y\,(x \in \mathbb{R} \to y \in \mathbb{R} \to 0 < x \to 0 < y \to 0 < x \cdot y) \to \varphi_{24}$
\item $\varphi_{74} :\equiv (x \in \mathbb{R} \to y \in \mathbb{R} \to z \in \mathbb{R} \to 0 < z \to x < y \to x < y) \to x \in \mathbb{R} \to y \in \mathbb{R} \to z \in \mathbb{R} \to 0 < z \to x < y \to 0 < y - x$
\item $\varphi_{75} :\equiv (x \in \mathbb{R} \to y \in \mathbb{R} \to z \in \mathbb{R} \to 0 < z \to x < y \to x \in \mathbb{R}) \to x \in \mathbb{R} \to y \in \mathbb{R} \to z \in \mathbb{R} \to 0 < z \to x < y \to z \cdot x = x \cdot z$
\item $\varphi_{76} :\equiv (x \in \mathbb{R} \to y \in \mathbb{R} \to z \in \mathbb{R} \to 0 < z \to x < y \to y \in \mathbb{R}) \to x \in \mathbb{R} \to y \in \mathbb{R} \to z \in \mathbb{R} \to 0 < z \to x < y \to z \cdot y = y \cdot z$
\item $\varphi_{77} :\equiv \varphi_{33} \to z \in \mathbb{R} \to y \in \mathbb{R} \to x \in \mathbb{R} \to z \cdot y - z \cdot x = z \cdot (y - x)$
\item $\varphi_{78} :\equiv (x \in \mathbb{R} \to y \in \mathbb{R} \to z \in \mathbb{R} \to 0 < z \to x < y \to x \in \mathbb{R}) \to x \in \mathbb{R} \to y \in \mathbb{R} \to z \in \mathbb{R} \to 0 < z \to x < y \to y - x \in \mathbb{R}$
\item $\varphi_{79} :\equiv (x \in \mathbb{R} \to y \in \mathbb{R} \to z \in \mathbb{R} \to 0 < z \to x < y \to z \in \mathbb{R}) \to x \in \mathbb{R} \to y \in \mathbb{R} \to z \in \mathbb{R} \to 0 < z \to x < y \to y \cdot z \in \mathbb{R}$
\item $\varphi_{80} :\equiv (x \in \mathbb{R} \to y \in \mathbb{R} \to z \in \mathbb{R} \to 0 < z \to x < y \to z \in \mathbb{R}) \to x \in \mathbb{R} \to y \in \mathbb{R} \to z \in \mathbb{R} \to 0 < z \to x < y \to x \cdot z \in \mathbb{R}$
\item $\varphi_{81} :\equiv \forall y\,(x \cdot z \in \mathbb{R} \to y \in \mathbb{R} \to 0 < y - x \cdot z \to x \cdot z < y) \to \varphi_{44}$
\item $\varphi_{82} :\equiv (x \in \mathbb{R} \to y \in \mathbb{R} \to z \in \mathbb{R} \to 0 < z \to x < y \to z \in \mathbb{R}) \to x \in \mathbb{R} \to \varphi_{28}$
\item $\varphi_{83} :\equiv \varphi_{53} \to x \in \mathbb{R} \to y \in \mathbb{R} \to z \in \mathbb{R} \to 0 < z \to x < y \to x \cdot z < y \cdot z$
\item $\varphi_{84} :\equiv (x \in \mathbb{R} \to y \in \mathbb{R} \to z \in \mathbb{R} \to 0 < z \to x < y \to x \in \mathbb{R}) \to x \in \mathbb{R} \to \varphi_{35}$
\item $\varphi_{85} :\equiv (x \in \mathbb{R} \to y \in \mathbb{R} \to z \in \mathbb{R} \to 0 < z \to x < y \to y \in \mathbb{R}) \to \varphi_{74}$
\item $\varphi_{86} :\equiv (x \in \mathbb{R} \to y \in \mathbb{R} \to z \in \mathbb{R} \to 0 < z \to x < y \to 0 < y - x) \to (x \in \mathbb{R} \to y \in \mathbb{R} \to z \in \mathbb{R} \to 0 < z \to x < y \to 0 < z) \to \varphi_{41}$
\item $\varphi_{87} :\equiv (x \in \mathbb{R} \to y \in \mathbb{R} \to z \in \mathbb{R} \to 0 < z \to x < y \to z \in \mathbb{R}) \to \varphi_{75}$
\item $\varphi_{88} :\equiv (x \in \mathbb{R} \to y \in \mathbb{R} \to z \in \mathbb{R} \to 0 < z \to x < y \to z \in \mathbb{R}) \to \varphi_{76}$
\item $\varphi_{89} :\equiv (x \in \mathbb{R} \to y \in \mathbb{R} \to z \in \mathbb{R} \to 0 < z \to x < y \to y \in \mathbb{R}) \to \varphi_{78}$
\item $\varphi_{90} :\equiv (x \in \mathbb{R} \to y \in \mathbb{R} \to z \in \mathbb{R} \to 0 < z \to x < y \to y \in \mathbb{R}) \to \varphi_{79}$
\item $\varphi_{91} :\equiv (x \in \mathbb{R} \to y \in \mathbb{R} \to z \in \mathbb{R} \to 0 < z \to x < y \to x \in \mathbb{R}) \to \varphi_{80}$
\item $\varphi_{92} :\equiv y \cdot z - z \cdot x = y \cdot z - z \cdot x \to (z \cdot x = x \cdot z \to \varphi_{30}) \to x \in \mathbb{R} \to \varphi_{49}$
\item $\varphi_{93} :\equiv z \cdot y - z \cdot x = z \cdot y - z \cdot x \to (z \cdot y = y \cdot z \to \varphi_{31}) \to x \in \mathbb{R} \to \varphi_{50}$
\item $\varphi_{94} :\equiv (x \in \mathbb{R} \to y \in \mathbb{R} \to z \in \mathbb{R} \to 0 < z \to x < y \to y \in \mathbb{R}) \to \varphi_{84}$
\item $\varphi_{95} :\equiv (x \in \mathbb{R} \to y \in \mathbb{R} \to z \in \mathbb{R} \to 0 < z \to x < y \to y - x \in \mathbb{R}) \to \varphi_{82}$
\item $\varphi_{96} :\equiv (x \in \mathbb{R} \to y \in \mathbb{R} \to z \in \mathbb{R} \to 0 < z \to x < y \to y \cdot z \in \mathbb{R}) \to \varphi_{83}$
\item $\varphi_{97} :\equiv (x \in \mathbb{R} \to \varphi_{36}) \to ((y - x) \cdot z = z \cdot (y - x) \to \varphi_{12}) \to x \in \mathbb{R} \to \varphi_{37}$
\item $\varphi_{98} :\equiv (x \in \mathbb{R} \to y \in \mathbb{R} \to z \in \mathbb{R} \to 0 < z \to x < y \to x \in \mathbb{R}) \to \varphi_{85}$
\item $\varphi_{99} :\equiv (x \in \mathbb{R} \to \varphi_{50}) \to ((y - x) \cdot z = z \cdot y - z \cdot x \to \varphi_{18}) \to x \in \mathbb{R} \to \varphi_{38}$
\item $\varphi_{100} :\equiv (x \in \mathbb{R} \to \varphi_{49}) \to ((y - x) \cdot z = y \cdot z - z \cdot x \to \varphi_{17}) \to x \in \mathbb{R} \to \varphi_{39}$
\item $\varphi_{101} :\equiv (x \in \mathbb{R} \to y \in \mathbb{R} \to z \in \mathbb{R} \to 0 < z \to x < y \to z \in \mathbb{R}) \to \varphi_{86}$
\item $\varphi_{102} :\equiv (y - x \in \mathbb{R} \to z \in \mathbb{R} \to (y - x) \cdot z = z \cdot (y - x)) \to \varphi_{95}$
\item $\varphi_{103} :\equiv (x \in \mathbb{R} \to y \in \mathbb{R} \to z \in \mathbb{R} \to 0 < z \to x < y \to z \cdot x = x \cdot z) \to \varphi_{92}$
\item $\varphi_{104} :\equiv (x \in \mathbb{R} \to y \in \mathbb{R} \to z \in \mathbb{R} \to 0 < z \to x < y \to z \cdot y = y \cdot z) \to \varphi_{93}$
\item $\varphi_{105} :\equiv (x \in \mathbb{R} \to y \in \mathbb{R} \to z \in \mathbb{R} \to 0 < z \to x < y \to z \in \mathbb{R}) \to \varphi_{94}$
\item $\varphi_{106} :\equiv (x \in \mathbb{R} \to y \in \mathbb{R} \to z \in \mathbb{R} \to 0 < z \to x < y \to x \cdot z \in \mathbb{R}) \to \varphi_{96}$
\item $\varphi_{107} :\equiv (x \in \mathbb{R} \to y \in \mathbb{R} \to z \in \mathbb{R} \to 0 < z \to x < y \to y - x \in \mathbb{R}) \to \varphi_{101}$
\item $\varphi_{108} :\equiv (z \in \mathbb{R} \to y \in \mathbb{R} \to x \in \mathbb{R} \to z \cdot y - z \cdot x = z \cdot (y - x)) \to \varphi_{105}$
\end{abbrevlist}
\end{fullproof}

% 自動生成 (prover/build.js): MulLeNonneg  ノード数 177
\begin{fullproof}{[F33]}{$\forall x\,\forall y\,\forall z\,(x \in \mathbb{R} \to y \in \mathbb{R} \to z \in \mathbb{R} \to 0 \leq z \to x \leq y \to x \cdot z \leq y \cdot z)$}{177}\phantomsection\label{tree:MulLeNonneg}
\noindent$\Pi_{1}$:
\begin{ptree}
\AXC{$x = x$\lbl{A5}}
\UIC{$\forall x\,(x = x)$}
\AXC{$\forall x\,(x = x) \to x \cdot z = x \cdot z$\lbl{A2}}
\BIC{$x \cdot z = x \cdot z$}
\AXC{$x \in \mathbb{R} \to y \in \mathbb{R} \to z \in \mathbb{R} \to 0 \leq z \to x \leq y \to x = y \to x = y$\lbl{A1}}
\AXC{$\varphi_{54}$\lbl{A1}}
\BIC{$\varphi_{47}$}
\BIC{$(x = y \to x \cdot z = x \cdot z \to x \cdot z = y \cdot z) \to \varphi_{26}$}
\end{ptree}
\noindent$\Pi_{2}$:
\begin{ptree}
\AXC{$z_{1} = w \to x \cdot z = z_{1} \cdot z \to x \cdot z = w \cdot z$\lbl{A6}}
\UIC{$\forall w\,(z_{1} = w \to x \cdot z = z_{1} \cdot z \to x \cdot z = w \cdot z)$}
\UIC{$\forall z_{1}\,\forall w\,(z_{1} = w \to x \cdot z = z_{1} \cdot z \to x \cdot z = w \cdot z)$}
\AXC{$\varphi_{40}$\lbl{A2}}
\BIC{$\forall w\,(x = w \to x \cdot z = x \cdot z \to x \cdot z = w \cdot z)$}
\AXC{$\varphi_{36}$\lbl{A2}}
\BIC{$x = y \to x \cdot z = x \cdot z \to x \cdot z = y \cdot z$}
\AXC{$\stackrel{\Pi_{1}}{\vdots}$}\noLine
\UIC{$(x = y \to x \cdot z = x \cdot z \to x \cdot z = y \cdot z) \to \varphi_{26}$}
\BIC{$\varphi_{26}$}
\end{ptree}
\noindent$\Pi_{3}$:
\begin{ptree}
\AXC{$y = z \to x = y \to x = z$\lbl{A6}}
\AXC{$(y = z \to x = y \to x = z) \to x = y \to y = z \to x = z$\lbl{A1}}
\BIC{$x = y \to y = z \to x = z$}
\UIC{$\forall z\,(x = y \to y = z \to x = z)$}
\UIC{$\forall y\,\forall z\,(x = y \to y = z \to x = z)$}
\UIC{$\forall x\,\forall y\,\forall z\,(x = y \to y = z \to x = z)$}
\AXC{$\varphi_{9}$\lbl{A2}}
\BIC{$\forall y\,\forall z\,(v_{282} = y \to y = z \to v_{282} = z)$}
\AXC{$\varphi_{12}$\lbl{A2}}
\BIC{$\forall z\,(v_{282} = v_{283} \to v_{283} = z \to v_{282} = z)$}
\AXC{$\varphi_{19}$\lbl{A2}}
\BIC{$v_{282} = v_{283} \to v_{283} = v_{284} \to v_{282} = v_{284}$}
\UIC{$\forall v_{284}\,(v_{282} = v_{283} \to v_{283} = v_{284} \to v_{282} = v_{284})$}
\UIC{$\varphi_{2}$}
\UIC{$\forall v_{282}\,(\varphi_{2})$}
\end{ptree}
\noindent$\Pi_{4}$:
\begin{ptree}
\AXC{$\stackrel{\Pi_{3}}{\vdots}$}\noLine
\UIC{$\forall v_{282}\,(\varphi_{2})$}
\AXC{$\forall v_{282}\,(\varphi_{2}) \to \varphi_{5}$\lbl{A2}}
\BIC{$\varphi_{5}$}
\AXC{$\varphi_{45}$\lbl{A2}}
\BIC{$\forall v_{284}\,(x \cdot z = y \cdot 0 \to y \cdot 0 = v_{284} \to x \cdot z = v_{284})$}
\AXC{$\varphi_{44}$\lbl{A2}}
\BIC{$x \cdot z = y \cdot 0 \to y \cdot 0 = y \cdot z \to x \cdot z = y \cdot z$}
\end{ptree}
\noindent$\Pi_{5}$:
\begin{ptree}
\AXC{$x = x$\lbl{A5}}
\UIC{$\forall x\,(x = x)$}
\AXC{$\forall x\,(x = x) \to y \cdot 0 = y \cdot 0$\lbl{A2}}
\BIC{$y \cdot 0 = y \cdot 0$}
\AXC{$x \in \mathbb{R} \to y \in \mathbb{R} \to z \in \mathbb{R} \to 0 \leq z \to x \leq y \to 0 = z \to 0 = z$\lbl{A1}}
\AXC{$\varphi_{55}$\lbl{A1}}
\BIC{$\varphi_{48}$}
\BIC{$(0 = z \to y \cdot 0 = y \cdot 0 \to y \cdot 0 = y \cdot z) \to \varphi_{28}$}
\end{ptree}
\noindent$\Pi_{6}$:
\begin{ptree}
\AXC{$z_{1} = w \to y \cdot 0 = y \cdot z_{1} \to y \cdot 0 = y \cdot w$\lbl{A6}}
\UIC{$\forall w\,(z_{1} = w \to y \cdot 0 = y \cdot z_{1} \to y \cdot 0 = y \cdot w)$}
\UIC{$\forall z_{1}\,\forall w\,(z_{1} = w \to y \cdot 0 = y \cdot z_{1} \to y \cdot 0 = y \cdot w)$}
\AXC{$\varphi_{41}$\lbl{A2}}
\BIC{$\forall w\,(0 = w \to y \cdot 0 = y \cdot 0 \to y \cdot 0 = y \cdot w)$}
\AXC{$\varphi_{37}$\lbl{A2}}
\BIC{$0 = z \to y \cdot 0 = y \cdot 0 \to y \cdot 0 = y \cdot z$}
\AXC{$\stackrel{\Pi_{5}}{\vdots}$}\noLine
\UIC{$(0 = z \to y \cdot 0 = y \cdot 0 \to y \cdot 0 = y \cdot z) \to \varphi_{28}$}
\BIC{$\varphi_{28}$}
\end{ptree}
\noindent$\Pi_{7}$:
\begin{ptree}
\AXC{$y = z \to x = y \to x = z$\lbl{A6}}
\AXC{$(y = z \to x = y \to x = z) \to x = y \to y = z \to x = z$\lbl{A1}}
\BIC{$x = y \to y = z \to x = z$}
\UIC{$\forall z\,(x = y \to y = z \to x = z)$}
\UIC{$\forall y\,\forall z\,(x = y \to y = z \to x = z)$}
\UIC{$\forall x\,\forall y\,\forall z\,(x = y \to y = z \to x = z)$}
\AXC{$\varphi_{10}$\lbl{A2}}
\BIC{$\forall y\,\forall z\,(v_{279} = y \to y = z \to v_{279} = z)$}
\AXC{$\varphi_{13}$\lbl{A2}}
\BIC{$\forall z\,(v_{279} = v_{280} \to v_{280} = z \to v_{279} = z)$}
\AXC{$\varphi_{20}$\lbl{A2}}
\BIC{$v_{279} = v_{280} \to v_{280} = v_{281} \to v_{279} = v_{281}$}
\UIC{$\forall v_{281}\,(v_{279} = v_{280} \to v_{280} = v_{281} \to v_{279} = v_{281})$}
\UIC{$\varphi_{3}$}
\UIC{$\forall v_{279}\,(\varphi_{3})$}
\end{ptree}
\noindent$\Pi_{8}$:
\begin{ptree}
\AXC{$\stackrel{\Pi_{7}}{\vdots}$}\noLine
\UIC{$\forall v_{279}\,(\varphi_{3})$}
\AXC{$\forall v_{279}\,(\varphi_{3}) \to \varphi_{6}$\lbl{A2}}
\BIC{$\varphi_{6}$}
\AXC{$\varphi_{6} \to \forall v_{281}\,(x \cdot z = 0 \to 0 = v_{281} \to x \cdot z = v_{281})$\lbl{A2}}
\BIC{$\forall v_{281}\,(x \cdot z = 0 \to 0 = v_{281} \to x \cdot z = v_{281})$}
\AXC{$\varphi_{35}$\lbl{A2}}
\BIC{$x \cdot z = 0 \to 0 = y \cdot 0 \to x \cdot z = y \cdot 0$}
\end{ptree}
\noindent$\Pi_{9}$:
\begin{ptree}
\AXC{$x = x$\lbl{A5}}
\AXC{$x = y \to x = x \to y = x$\lbl{A6}}
\AXC{$(x = y \to x = x \to y = x) \to x = x \to x = y \to y = x$\lbl{A1}}
\BIC{$x = x \to x = y \to y = x$}
\BIC{$x = y \to y = x$}
\UIC{$\forall y\,(x = y \to y = x)$}
\UIC{$\forall x\,\forall y\,(x = y \to y = x)$}
\end{ptree}
\noindent$\Pi_{10}$:
\begin{ptree}
\AXC{$\stackrel{\Pi_{9}}{\vdots}$}\noLine
\UIC{$\forall x\,\forall y\,(x = y \to y = x)$}
\AXC{$\forall x\,\forall y\,(x = y \to y = x) \to \forall y\,(v_{274} = y \to y = v_{274})$\lbl{A2}}
\BIC{$\forall y\,(v_{274} = y \to y = v_{274})$}
\AXC{$\varphi_{1}$\lbl{A2}}
\BIC{$v_{274} = v_{275} \to v_{275} = v_{274}$}
\UIC{$\forall v_{275}\,(v_{274} = v_{275} \to v_{275} = v_{274})$}
\UIC{$\forall v_{274}\,\forall v_{275}\,(v_{274} = v_{275} \to v_{275} = v_{274})$}
\AXC{$\varphi_{22}$\lbl{A2}}
\BIC{$\forall v_{275}\,(y \cdot 0 = v_{275} \to v_{275} = y \cdot 0)$}
\AXC{$\varphi_{8}$\lbl{A2}}
\BIC{$y \cdot 0 = 0 \to 0 = y \cdot 0$}
\end{ptree}
\noindent$\Pi_{11}$:
\begin{ptree}
\AXC{$\varphi_{23}$\lbl{A1}}
\AXC{$\forall x \in \mathbb{R}\,(x \cdot 0 = 0)$\lbl{[F8]}}
\AXC{$\forall x \in \mathbb{R}\,(x \cdot 0 = 0) \to y \in \mathbb{R} \to y \cdot 0 = 0$\lbl{A2}}
\BIC{$y \in \mathbb{R} \to y \cdot 0 = 0$}
\AXC{$(y \in \mathbb{R} \to y \cdot 0 = 0) \to \varphi_{23} \to \varphi_{15}$\lbl{A1}}
\BIC{$\varphi_{23} \to \varphi_{15}$}
\BIC{$\varphi_{15}$}
\end{ptree}
\noindent$\Pi_{12}$:
\begin{ptree}
\AXC{$\stackrel{\Pi_{10}}{\vdots}$}\noLine
\UIC{$y \cdot 0 = 0 \to 0 = y \cdot 0$}
\AXC{$\stackrel{\Pi_{11}}{\vdots}$}\noLine
\UIC{$\varphi_{15}$}
\AXC{$\varphi_{15} \to (y \cdot 0 = 0 \to 0 = y \cdot 0) \to \varphi_{16}$\lbl{A1}}
\BIC{$(y \cdot 0 = 0 \to 0 = y \cdot 0) \to \varphi_{16}$}
\BIC{$\varphi_{16}$}
\end{ptree}
\noindent$\Pi_{13}$:
\begin{ptree}
\AXC{$y = z \to x = y \to x = z$\lbl{A6}}
\AXC{$(y = z \to x = y \to x = z) \to x = y \to y = z \to x = z$\lbl{A1}}
\BIC{$x = y \to y = z \to x = z$}
\UIC{$\forall z\,(x = y \to y = z \to x = z)$}
\UIC{$\forall y\,\forall z\,(x = y \to y = z \to x = z)$}
\UIC{$\forall x\,\forall y\,\forall z\,(x = y \to y = z \to x = z)$}
\AXC{$\varphi_{11}$\lbl{A2}}
\BIC{$\forall y\,\forall z\,(v_{276} = y \to y = z \to v_{276} = z)$}
\AXC{$\varphi_{14}$\lbl{A2}}
\BIC{$\forall z\,(v_{276} = v_{277} \to v_{277} = z \to v_{276} = z)$}
\AXC{$\varphi_{21}$\lbl{A2}}
\BIC{$v_{276} = v_{277} \to v_{277} = v_{278} \to v_{276} = v_{278}$}
\UIC{$\forall v_{278}\,(v_{276} = v_{277} \to v_{277} = v_{278} \to v_{276} = v_{278})$}
\UIC{$\varphi_{4}$}
\UIC{$\forall v_{276}\,(\varphi_{4})$}
\end{ptree}
\noindent$\Pi_{14}$:
\begin{ptree}
\AXC{$\varphi_{24}$\lbl{A1}}
\AXC{$\forall x \in \mathbb{R}\,(x \cdot 0 = 0)$\lbl{[F8]}}
\AXC{$\forall x \in \mathbb{R}\,(x \cdot 0 = 0) \to x \in \mathbb{R} \to x \cdot 0 = 0$\lbl{A2}}
\BIC{$x \in \mathbb{R} \to x \cdot 0 = 0$}
\AXC{$(x \in \mathbb{R} \to x \cdot 0 = 0) \to \varphi_{24} \to \varphi_{17}$\lbl{A1}}
\BIC{$\varphi_{24} \to \varphi_{17}$}
\BIC{$\varphi_{17}$}
\end{ptree}
\noindent$\Pi_{15}$:
\begin{ptree}
\AXC{$z_{1} = w \to x \cdot z = x \cdot z_{1} \to x \cdot z = x \cdot w$\lbl{A6}}
\UIC{$\forall w\,(z_{1} = w \to x \cdot z = x \cdot z_{1} \to x \cdot z = x \cdot w)$}
\UIC{$\forall z_{1}\,\forall w\,(z_{1} = w \to x \cdot z = x \cdot z_{1} \to x \cdot z = x \cdot w)$}
\AXC{$\varphi_{43}$\lbl{A2}}
\BIC{$\forall w\,(z = w \to x \cdot z = x \cdot z \to x \cdot z = x \cdot w)$}
\AXC{$\varphi_{38}$\lbl{A2}}
\BIC{$z = 0 \to x \cdot z = x \cdot z \to x \cdot z = x \cdot 0$}
\end{ptree}
\noindent$\Pi_{16}$:
\begin{ptree}
\AXC{$\stackrel{\Pi_{9}}{\vdots}$}\noLine
\UIC{$\forall x\,\forall y\,(x = y \to y = x)$}
\AXC{$\forall x\,\forall y\,(x = y \to y = x) \to \forall y\,(0 = y \to y = 0)$\lbl{A2}}
\BIC{$\forall y\,(0 = y \to y = 0)$}
\AXC{$\forall y\,(0 = y \to y = 0) \to 0 = z \to z = 0$\lbl{A2}}
\BIC{$0 = z \to z = 0$}
\end{ptree}
\noindent$\Pi_{17}$:
\begin{ptree}
\AXC{$\stackrel{\Pi_{16}}{\vdots}$}\noLine
\UIC{$0 = z \to z = 0$}
\AXC{$x \in \mathbb{R} \to y \in \mathbb{R} \to z \in \mathbb{R} \to 0 \leq z \to x \leq y \to 0 = z \to 0 = z$\lbl{A1}}
\AXC{$\varphi_{50}$\lbl{A1}}
\BIC{$\varphi_{33}$}
\BIC{$x \in \mathbb{R} \to y \in \mathbb{R} \to z \in \mathbb{R} \to 0 \leq z \to x \leq y \to 0 = z \to z = 0$}
\AXC{$\varphi_{56}$\lbl{A1}}
\BIC{$\varphi_{49}$}
\end{ptree}
\noindent$\Pi_{18}$:
\begin{ptree}
\AXC{$\stackrel{\Pi_{15}}{\vdots}$}\noLine
\UIC{$z = 0 \to x \cdot z = x \cdot z \to x \cdot z = x \cdot 0$}
\AXC{$x = x$\lbl{A5}}
\UIC{$\forall x\,(x = x)$}
\AXC{$\forall x\,(x = x) \to x \cdot z = x \cdot z$\lbl{A2}}
\BIC{$x \cdot z = x \cdot z$}
\AXC{$\stackrel{\Pi_{17}}{\vdots}$}\noLine
\UIC{$\varphi_{49}$}
\BIC{$(z = 0 \to x \cdot z = x \cdot z \to x \cdot z = x \cdot 0) \to \varphi_{29}$}
\BIC{$\varphi_{29}$}
\end{ptree}
\noindent$\Pi_{19}$:
\begin{ptree}
\AXC{$\stackrel{\Pi_{14}}{\vdots}$}\noLine
\UIC{$\varphi_{17}$}
\AXC{$\stackrel{\Pi_{18}}{\vdots}$}\noLine
\UIC{$\varphi_{29}$}
\AXC{$\varphi_{29} \to \varphi_{17} \to (x \cdot z = x \cdot 0 \to x \cdot 0 = 0 \to x \cdot z = 0) \to \varphi_{18}$\lbl{A1}}
\BIC{$\varphi_{17} \to (x \cdot z = x \cdot 0 \to x \cdot 0 = 0 \to x \cdot z = 0) \to \varphi_{18}$}
\BIC{$(x \cdot z = x \cdot 0 \to x \cdot 0 = 0 \to x \cdot z = 0) \to \varphi_{18}$}
\end{ptree}
\noindent$\Pi_{20}$:
\begin{ptree}
\AXC{$\stackrel{\Pi_{13}}{\vdots}$}\noLine
\UIC{$\forall v_{276}\,(\varphi_{4})$}
\AXC{$\forall v_{276}\,(\varphi_{4}) \to \varphi_{7}$\lbl{A2}}
\BIC{$\varphi_{7}$}
\AXC{$\varphi_{46}$\lbl{A2}}
\BIC{$\forall v_{278}\,(x \cdot z = x \cdot 0 \to x \cdot 0 = v_{278} \to x \cdot z = v_{278})$}
\AXC{$\varphi_{42}$\lbl{A2}}
\BIC{$x \cdot z = x \cdot 0 \to x \cdot 0 = 0 \to x \cdot z = 0$}
\AXC{$\stackrel{\Pi_{19}}{\vdots}$}\noLine
\UIC{$(x \cdot z = x \cdot 0 \to x \cdot 0 = 0 \to x \cdot z = 0) \to \varphi_{18}$}
\BIC{$\varphi_{18}$}
\end{ptree}
\noindent$\Pi_{21}$:
\begin{ptree}
\AXC{$\stackrel{\Pi_{8}}{\vdots}$}\noLine
\UIC{$x \cdot z = 0 \to 0 = y \cdot 0 \to x \cdot z = y \cdot 0$}
\AXC{$\stackrel{\Pi_{12}}{\vdots}$}\noLine
\UIC{$\varphi_{16}$}
\AXC{$\stackrel{\Pi_{20}}{\vdots}$}\noLine
\UIC{$\varphi_{18}$}
\AXC{$\varphi_{18} \to \varphi_{16} \to (x \cdot z = 0 \to 0 = y \cdot 0 \to x \cdot z = y \cdot 0) \to \varphi_{30}$\lbl{A1}}
\BIC{$\varphi_{16} \to (x \cdot z = 0 \to 0 = y \cdot 0 \to x \cdot z = y \cdot 0) \to \varphi_{30}$}
\BIC{$(x \cdot z = 0 \to 0 = y \cdot 0 \to x \cdot z = y \cdot 0) \to \varphi_{30}$}
\BIC{$\varphi_{30}$}
\end{ptree}
\noindent$\Pi_{22}$:
\begin{ptree}
\AXC{$\stackrel{\Pi_{4}}{\vdots}$}\noLine
\UIC{$x \cdot z = y \cdot 0 \to y \cdot 0 = y \cdot z \to x \cdot z = y \cdot z$}
\AXC{$\stackrel{\Pi_{6}}{\vdots}$}\noLine
\UIC{$\varphi_{28}$}
\AXC{$\stackrel{\Pi_{21}}{\vdots}$}\noLine
\UIC{$\varphi_{30}$}
\AXC{$\varphi_{58}$\lbl{A1}}
\BIC{$\varphi_{28} \to (x \cdot z = y \cdot 0 \to y \cdot 0 = y \cdot z \to x \cdot z = y \cdot z) \to \varphi_{31}$}
\BIC{$(x \cdot z = y \cdot 0 \to y \cdot 0 = y \cdot z \to x \cdot z = y \cdot z) \to \varphi_{31}$}
\BIC{$\varphi_{31}$}
\AXC{$\varphi_{31} \to \varphi_{32}$\lbl{A1}}
\BIC{$\varphi_{32}$}
\end{ptree}
\noindent$\Pi_{23}$:
\begin{ptree}
\AXC{$\forall x\,\forall y\,(\varphi_{39})$\lbl{[F32]}}
\AXC{$\forall x\,\forall y\,(\varphi_{39}) \to \forall y\,(\varphi_{39})$\lbl{A2}}
\BIC{$\forall y\,(\varphi_{39})$}
\AXC{$\forall y\,(\varphi_{39}) \to \varphi_{39}$\lbl{A2}}
\BIC{$\varphi_{39}$}
\AXC{$\varphi_{53}$\lbl{A2}}
\BIC{$x \in \mathbb{R} \to y \in \mathbb{R} \to z \in \mathbb{R} \to 0 < z \to x < y \to x \cdot z < y \cdot z$}
\AXC{$\varphi_{60}$\lbl{A1}}
\BIC{$\varphi_{59}$}
\end{ptree}
\noindent$\Pi_{24}$:
\begin{ptree}
\AXC{$x \in \mathbb{R} \to y \in \mathbb{R} \to z \in \mathbb{R} \to 0 \leq z \to x \leq y \to y \in \mathbb{R}$\lbl{A1}}
\AXC{$x \in \mathbb{R} \to y \in \mathbb{R} \to z \in \mathbb{R} \to 0 \leq z \to x \leq y \to x \in \mathbb{R}$\lbl{A1}}
\AXC{$\stackrel{\Pi_{23}}{\vdots}$}\noLine
\UIC{$\varphi_{59}$}
\BIC{$\varphi_{57}$}
\BIC{$\varphi_{52}$}
\end{ptree}
\noindent$\Pi_{25}$:
\begin{ptree}
\AXC{$x \in \mathbb{R} \to y \in \mathbb{R} \to z \in \mathbb{R} \to 0 \leq z \to x \leq y \to z \in \mathbb{R}$\lbl{A1}}
\AXC{$\stackrel{\Pi_{24}}{\vdots}$}\noLine
\UIC{$\varphi_{52}$}
\BIC{$x \in \mathbb{R} \to \varphi_{34}$}
\end{ptree}
\noindent$\Pi_{26}$:
\begin{ptree}
\AXC{$x \in \mathbb{R} \to y \in \mathbb{R} \to z \in \mathbb{R} \to 0 \leq z \to x \leq y \to 0 \leq z$\lbl{A1}}
\AXC{$(x \in \mathbb{R} \to y \in \mathbb{R} \to z \in \mathbb{R} \to 0 \leq z \to x \leq y \to 0 \leq z) \to \varphi_{61}$\lbl{A1}}
\BIC{$\varphi_{61}$}
\end{ptree}
\noindent 主木:
\begin{ptree}
\AXC{$\stackrel{\Pi_{2}}{\vdots}$}\noLine
\UIC{$\varphi_{26}$}
\AXC{$\varphi_{26} \to \varphi_{27}$\lbl{A1}}
\BIC{$\varphi_{27}$}
\AXC{$\stackrel{\Pi_{22}}{\vdots}$}\noLine
\UIC{$\varphi_{32}$}
\AXC{$\stackrel{\Pi_{25}}{\vdots}$}\noLine
\UIC{$x \in \mathbb{R} \to \varphi_{34}$}
\AXC{$x \in \mathbb{R} \to y \in \mathbb{R} \to z \in \mathbb{R} \to 0 \leq z \to x \leq y \to x \leq y$\lbl{A1}}
\AXC{$\stackrel{\Pi_{26}}{\vdots}$}\noLine
\UIC{$\varphi_{61}$}
\BIC{$(x \in \mathbb{R} \to \varphi_{34}) \to \varphi_{32} \to \varphi_{51}$}
\BIC{$\varphi_{32} \to \varphi_{51}$}
\BIC{$\varphi_{51}$}
\BIC{$x \in \mathbb{R} \to y \in \mathbb{R} \to z \in \mathbb{R} \to 0 \leq z \to x \leq y \to x \cdot z \leq y \cdot z$}
\UIC{$\varphi_{25}$}
\UIC{$\forall y\,(\varphi_{25})$}
\UIC{$\forall x\,\forall y\,(\varphi_{25})$}
\end{ptree}
\begin{abbrevlist}
\item $\varphi_{1} :\equiv \forall y\,(v_{274} = y \to y = v_{274}) \to v_{274} = v_{275} \to v_{275} = v_{274}$
\item $\varphi_{2} :\equiv \forall v_{283}\,\forall v_{284}\,(v_{282} = v_{283} \to v_{283} = v_{284} \to v_{282} = v_{284})$
\item $\varphi_{3} :\equiv \forall v_{280}\,\forall v_{281}\,(v_{279} = v_{280} \to v_{280} = v_{281} \to v_{279} = v_{281})$
\item $\varphi_{4} :\equiv \forall v_{277}\,\forall v_{278}\,(v_{276} = v_{277} \to v_{277} = v_{278} \to v_{276} = v_{278})$
\item $\varphi_{5} :\equiv \forall v_{283}\,\forall v_{284}\,(x \cdot z = v_{283} \to v_{283} = v_{284} \to x \cdot z = v_{284})$
\item $\varphi_{6} :\equiv \forall v_{280}\,\forall v_{281}\,(x \cdot z = v_{280} \to v_{280} = v_{281} \to x \cdot z = v_{281})$
\item $\varphi_{7} :\equiv \forall v_{277}\,\forall v_{278}\,(x \cdot z = v_{277} \to v_{277} = v_{278} \to x \cdot z = v_{278})$
\item $\varphi_{8} :\equiv \forall v_{275}\,(y \cdot 0 = v_{275} \to v_{275} = y \cdot 0) \to y \cdot 0 = 0 \to 0 = y \cdot 0$
\item $\varphi_{9} :\equiv \forall x\,\forall y\,\forall z\,(x = y \to y = z \to x = z) \to \forall y\,\forall z\,(v_{282} = y \to y = z \to v_{282} = z)$
\item $\varphi_{10} :\equiv \forall x\,\forall y\,\forall z\,(x = y \to y = z \to x = z) \to \forall y\,\forall z\,(v_{279} = y \to y = z \to v_{279} = z)$
\item $\varphi_{11} :\equiv \forall x\,\forall y\,\forall z\,(x = y \to y = z \to x = z) \to \forall y\,\forall z\,(v_{276} = y \to y = z \to v_{276} = z)$
\item $\varphi_{12} :\equiv \forall y\,\forall z\,(v_{282} = y \to y = z \to v_{282} = z) \to \forall z\,(v_{282} = v_{283} \to v_{283} = z \to v_{282} = z)$
\item $\varphi_{13} :\equiv \forall y\,\forall z\,(v_{279} = y \to y = z \to v_{279} = z) \to \forall z\,(v_{279} = v_{280} \to v_{280} = z \to v_{279} = z)$
\item $\varphi_{14} :\equiv \forall y\,\forall z\,(v_{276} = y \to y = z \to v_{276} = z) \to \forall z\,(v_{276} = v_{277} \to v_{277} = z \to v_{276} = z)$
\item $\varphi_{15} :\equiv x \in \mathbb{R} \to y \in \mathbb{R} \to z \in \mathbb{R} \to 0 \leq z \to x \leq y \to 0 = z \to y \cdot 0 = 0$
\item $\varphi_{16} :\equiv x \in \mathbb{R} \to y \in \mathbb{R} \to z \in \mathbb{R} \to 0 \leq z \to x \leq y \to 0 = z \to 0 = y \cdot 0$
\item $\varphi_{17} :\equiv x \in \mathbb{R} \to y \in \mathbb{R} \to z \in \mathbb{R} \to 0 \leq z \to x \leq y \to 0 = z \to x \cdot 0 = 0$
\item $\varphi_{18} :\equiv x \in \mathbb{R} \to y \in \mathbb{R} \to z \in \mathbb{R} \to 0 \leq z \to x \leq y \to 0 = z \to x \cdot z = 0$
\item $\varphi_{19} :\equiv \forall z\,(v_{282} = v_{283} \to v_{283} = z \to v_{282} = z) \to v_{282} = v_{283} \to v_{283} = v_{284} \to v_{282} = v_{284}$
\item $\varphi_{20} :\equiv \forall z\,(v_{279} = v_{280} \to v_{280} = z \to v_{279} = z) \to v_{279} = v_{280} \to v_{280} = v_{281} \to v_{279} = v_{281}$
\item $\varphi_{21} :\equiv \forall z\,(v_{276} = v_{277} \to v_{277} = z \to v_{276} = z) \to v_{276} = v_{277} \to v_{277} = v_{278} \to v_{276} = v_{278}$
\item $\varphi_{22} :\equiv \forall v_{274}\,\forall v_{275}\,(v_{274} = v_{275} \to v_{275} = v_{274}) \to \forall v_{275}\,(y \cdot 0 = v_{275} \to v_{275} = y \cdot 0)$
\item $\varphi_{23} :\equiv x \in \mathbb{R} \to y \in \mathbb{R} \to z \in \mathbb{R} \to 0 \leq z \to x \leq y \to 0 = z \to y \in \mathbb{R}$
\item $\varphi_{24} :\equiv x \in \mathbb{R} \to y \in \mathbb{R} \to z \in \mathbb{R} \to 0 \leq z \to x \leq y \to 0 = z \to x \in \mathbb{R}$
\item $\varphi_{25} :\equiv \forall z\,(x \in \mathbb{R} \to y \in \mathbb{R} \to z \in \mathbb{R} \to 0 \leq z \to x \leq y \to x \cdot z \leq y \cdot z)$
\item $\varphi_{26} :\equiv x \in \mathbb{R} \to y \in \mathbb{R} \to z \in \mathbb{R} \to 0 \leq z \to x \leq y \to x = y \to x \cdot z = y \cdot z$
\item $\varphi_{27} :\equiv x \in \mathbb{R} \to y \in \mathbb{R} \to z \in \mathbb{R} \to 0 \leq z \to x \leq y \to x = y \to x \cdot z \leq y \cdot z$
\item $\varphi_{28} :\equiv x \in \mathbb{R} \to y \in \mathbb{R} \to z \in \mathbb{R} \to 0 \leq z \to x \leq y \to 0 = z \to y \cdot 0 = y \cdot z$
\item $\varphi_{29} :\equiv x \in \mathbb{R} \to y \in \mathbb{R} \to z \in \mathbb{R} \to 0 \leq z \to x \leq y \to 0 = z \to x \cdot z = x \cdot 0$
\item $\varphi_{30} :\equiv x \in \mathbb{R} \to y \in \mathbb{R} \to z \in \mathbb{R} \to 0 \leq z \to x \leq y \to 0 = z \to x \cdot z = y \cdot 0$
\item $\varphi_{31} :\equiv x \in \mathbb{R} \to y \in \mathbb{R} \to z \in \mathbb{R} \to 0 \leq z \to x \leq y \to 0 = z \to x \cdot z = y \cdot z$
\item $\varphi_{32} :\equiv x \in \mathbb{R} \to y \in \mathbb{R} \to z \in \mathbb{R} \to 0 \leq z \to x \leq y \to 0 = z \to x \cdot z \leq y \cdot z$
\item $\varphi_{33} :\equiv (0 = z \to z = 0) \to x \in \mathbb{R} \to y \in \mathbb{R} \to z \in \mathbb{R} \to 0 \leq z \to x \leq y \to 0 = z \to z = 0$
\item $\varphi_{34} :\equiv y \in \mathbb{R} \to z \in \mathbb{R} \to 0 \leq z \to x \leq y \to 0 < z \to x < y \to x \cdot z < y \cdot z$
\item $\varphi_{35} :\equiv \forall v_{281}\,(x \cdot z = 0 \to 0 = v_{281} \to x \cdot z = v_{281}) \to x \cdot z = 0 \to 0 = y \cdot 0 \to x \cdot z = y \cdot 0$
\item $\varphi_{36} :\equiv \forall w\,(x = w \to x \cdot z = x \cdot z \to x \cdot z = w \cdot z) \to x = y \to x \cdot z = x \cdot z \to x \cdot z = y \cdot z$
\item $\varphi_{37} :\equiv \forall w\,(0 = w \to y \cdot 0 = y \cdot 0 \to y \cdot 0 = y \cdot w) \to 0 = z \to y \cdot 0 = y \cdot 0 \to y \cdot 0 = y \cdot z$
\item $\varphi_{38} :\equiv \forall w\,(z = w \to x \cdot z = x \cdot z \to x \cdot z = x \cdot w) \to z = 0 \to x \cdot z = x \cdot z \to x \cdot z = x \cdot 0$
\item $\varphi_{39} :\equiv \forall z\,(x \in \mathbb{R} \to y \in \mathbb{R} \to z \in \mathbb{R} \to 0 < z \to x < y \to x \cdot z < y \cdot z)$
\item $\varphi_{40} :\equiv \forall z_{1}\,\forall w\,(z_{1} = w \to x \cdot z = z_{1} \cdot z \to x \cdot z = w \cdot z) \to \forall w\,(x = w \to x \cdot z = x \cdot z \to x \cdot z = w \cdot z)$
\item $\varphi_{41} :\equiv \forall z_{1}\,\forall w\,(z_{1} = w \to y \cdot 0 = y \cdot z_{1} \to y \cdot 0 = y \cdot w) \to \forall w\,(0 = w \to y \cdot 0 = y \cdot 0 \to y \cdot 0 = y \cdot w)$
\item $\varphi_{42} :\equiv \forall v_{278}\,(x \cdot z = x \cdot 0 \to x \cdot 0 = v_{278} \to x \cdot z = v_{278}) \to x \cdot z = x \cdot 0 \to x \cdot 0 = 0 \to x \cdot z = 0$
\item $\varphi_{43} :\equiv \forall z_{1}\,\forall w\,(z_{1} = w \to x \cdot z = x \cdot z_{1} \to x \cdot z = x \cdot w) \to \forall w\,(z = w \to x \cdot z = x \cdot z \to x \cdot z = x \cdot w)$
\item $\varphi_{44} :\equiv \forall v_{284}\,(x \cdot z = y \cdot 0 \to y \cdot 0 = v_{284} \to x \cdot z = v_{284}) \to x \cdot z = y \cdot 0 \to y \cdot 0 = y \cdot z \to x \cdot z = y \cdot z$
\item $\varphi_{45} :\equiv \varphi_{5} \to \forall v_{284}\,(x \cdot z = y \cdot 0 \to y \cdot 0 = v_{284} \to x \cdot z = v_{284})$
\item $\varphi_{46} :\equiv \varphi_{7} \to \forall v_{278}\,(x \cdot z = x \cdot 0 \to x \cdot 0 = v_{278} \to x \cdot z = v_{278})$
\item $\varphi_{47} :\equiv x \cdot z = x \cdot z \to (x = y \to x \cdot z = x \cdot z \to x \cdot z = y \cdot z) \to \varphi_{26}$
\item $\varphi_{48} :\equiv y \cdot 0 = y \cdot 0 \to (0 = z \to y \cdot 0 = y \cdot 0 \to y \cdot 0 = y \cdot z) \to \varphi_{28}$
\item $\varphi_{49} :\equiv x \cdot z = x \cdot z \to (z = 0 \to x \cdot z = x \cdot z \to x \cdot z = x \cdot 0) \to \varphi_{29}$
\item $\varphi_{50} :\equiv (x \in \mathbb{R} \to y \in \mathbb{R} \to z \in \mathbb{R} \to 0 \leq z \to x \leq y \to 0 = z \to 0 = z) \to \varphi_{33}$
\item $\varphi_{51} :\equiv \varphi_{27} \to x \in \mathbb{R} \to y \in \mathbb{R} \to z \in \mathbb{R} \to 0 \leq z \to x \leq y \to x \cdot z \leq y \cdot z$
\item $\varphi_{52} :\equiv (x \in \mathbb{R} \to y \in \mathbb{R} \to z \in \mathbb{R} \to 0 \leq z \to x \leq y \to z \in \mathbb{R}) \to x \in \mathbb{R} \to \varphi_{34}$
\item $\varphi_{53} :\equiv \varphi_{39} \to x \in \mathbb{R} \to y \in \mathbb{R} \to z \in \mathbb{R} \to 0 < z \to x < y \to x \cdot z < y \cdot z$
\item $\varphi_{54} :\equiv (x \in \mathbb{R} \to y \in \mathbb{R} \to z \in \mathbb{R} \to 0 \leq z \to x \leq y \to x = y \to x = y) \to \varphi_{47}$
\item $\varphi_{55} :\equiv (x \in \mathbb{R} \to y \in \mathbb{R} \to z \in \mathbb{R} \to 0 \leq z \to x \leq y \to 0 = z \to 0 = z) \to \varphi_{48}$
\item $\varphi_{56} :\equiv (x \in \mathbb{R} \to y \in \mathbb{R} \to z \in \mathbb{R} \to 0 \leq z \to x \leq y \to 0 = z \to z = 0) \to \varphi_{49}$
\item $\varphi_{57} :\equiv (x \in \mathbb{R} \to y \in \mathbb{R} \to z \in \mathbb{R} \to 0 \leq z \to x \leq y \to y \in \mathbb{R}) \to \varphi_{52}$
\item $\varphi_{58} :\equiv \varphi_{30} \to \varphi_{28} \to (x \cdot z = y \cdot 0 \to y \cdot 0 = y \cdot z \to x \cdot z = y \cdot z) \to \varphi_{31}$
\item $\varphi_{59} :\equiv (x \in \mathbb{R} \to y \in \mathbb{R} \to z \in \mathbb{R} \to 0 \leq z \to x \leq y \to x \in \mathbb{R}) \to \varphi_{57}$
\item $\varphi_{60} :\equiv (x \in \mathbb{R} \to y \in \mathbb{R} \to z \in \mathbb{R} \to 0 < z \to x < y \to x \cdot z < y \cdot z) \to \varphi_{59}$
\item $\varphi_{61} :\equiv (x \in \mathbb{R} \to y \in \mathbb{R} \to z \in \mathbb{R} \to 0 \leq z \to x \leq y \to x \leq y) \to (x \in \mathbb{R} \to \varphi_{34}) \to \varphi_{32} \to \varphi_{51}$
\end{abbrevlist}
\end{fullproof}

% 自動生成 (prover/build.js): InvPos  ノード数 263
\begin{fullproof}{[F34]}{$\forall x \in \mathbb{R}\,(0 < x \to 0 < x^{-1})$}{263}\phantomsection\label{tree:InvPos}
\noindent$\Pi_{1}$:
\begin{ptree}
\AXC{$\forall x\,\forall y\,(x \in \mathbb{R} \to y \in \mathbb{R} \to x < y \to -y < -x)$\lbl{[F22]}}
\AXC{$\varphi_{54}$\lbl{A2}}
\BIC{$\forall y\,(0 \in \mathbb{R} \to y \in \mathbb{R} \to 0 < y \to -y < -0)$}
\AXC{$\varphi_{46}$\lbl{A2}}
\BIC{$0 \in \mathbb{R} \to 1 \in \mathbb{R} \to 0 < 1 \to -1 < -0$}
\AXC{$(0 \in \mathbb{R} \to 1 \in \mathbb{R} \to 0 < 1 \to -1 < -0) \to \varphi_{31}$\lbl{A1}}
\BIC{$\varphi_{31}$}
\end{ptree}
\noindent$\Pi_{2}$:
\begin{ptree}
\AXC{$0 < 1$\lbl{[F28]}}
\AXC{$1 \in \mathbb{R}$\lbl{(R1)}}
\AXC{$0 \in \mathbb{R}$\lbl{(R1)}}
\AXC{$\stackrel{\Pi_{1}}{\vdots}$}\noLine
\UIC{$\varphi_{31}$}
\BIC{$1 \in \mathbb{R} \to 0 < 1 \to x \in \mathbb{R} \to 0 < x \to x^{-1} < 0 \to -1 < -0$}
\BIC{$0 < 1 \to x \in \mathbb{R} \to 0 < x \to x^{-1} < 0 \to -1 < -0$}
\BIC{$x \in \mathbb{R} \to 0 < x \to x^{-1} < 0 \to -1 < -0$}
\end{ptree}
\noindent$\Pi_{3}$:
\begin{ptree}
\AXC{$\stackrel{\Pi_{2}}{\vdots}$}\noLine
\UIC{$x \in \mathbb{R} \to 0 < x \to x^{-1} < 0 \to -1 < -0$}
\AXC{$-0 = 0$\lbl{[F6]}}
\AXC{$-0 = 0 \to (x \in \mathbb{R} \to 0 < x \to x^{-1} < 0 \to -1 < -0) \to \varphi_{37}$\lbl{A1}}
\BIC{$(x \in \mathbb{R} \to 0 < x \to x^{-1} < 0 \to -1 < -0) \to \varphi_{37}$}
\BIC{$\varphi_{37}$}
\end{ptree}
\noindent$\Pi_{4}$:
\begin{ptree}
\AXC{$x = x$\lbl{A5}}
\AXC{$x = y \to x = x \to y = x$\lbl{A6}}
\AXC{$(x = y \to x = x \to y = x) \to x = x \to x = y \to y = x$\lbl{A1}}
\BIC{$x = x \to x = y \to y = x$}
\BIC{$x = y \to y = x$}
\UIC{$\forall y\,(x = y \to y = x)$}
\UIC{$\forall x\,\forall y\,(x = y \to y = x)$}
\end{ptree}
\noindent$\Pi_{5}$:
\begin{ptree}
\AXC{$\stackrel{\Pi_{4}}{\vdots}$}\noLine
\UIC{$\forall x\,\forall y\,(x = y \to y = x)$}
\AXC{$\forall x\,\forall y\,(x = y \to y = x) \to \forall y\,(x = y \to y = x)$\lbl{A2}}
\BIC{$\forall y\,(x = y \to y = x)$}
\AXC{$\forall y\,(x = y \to y = x) \to x = 0 \to 0 = x$\lbl{A2}}
\BIC{$x = 0 \to 0 = x$}
\end{ptree}
\noindent$\Pi_{6}$:
\begin{ptree}
\AXC{$\forall x\,\forall y\,(x < y \to x \neq y)$\lbl{[F26]}}
\AXC{$\forall x\,\forall y\,(x < y \to x \neq y) \to \forall y\,(0 < y \to 0 \neq y)$\lbl{A2}}
\BIC{$\forall y\,(0 < y \to 0 \neq y)$}
\AXC{$\forall y\,(0 < y \to 0 \neq y) \to 0 < x \to 0 \neq x$\lbl{A2}}
\BIC{$0 < x \to 0 \neq x$}
\end{ptree}
\noindent$\Pi_{7}$:
\begin{ptree}
\AXC{$\stackrel{\Pi_{5}}{\vdots}$}\noLine
\UIC{$x = 0 \to 0 = x$}
\AXC{$x \in \mathbb{R} \to 0 < x \to 0 < x$\lbl{A1}}
\AXC{$\stackrel{\Pi_{6}}{\vdots}$}\noLine
\UIC{$0 < x \to 0 \neq x$}
\AXC{$\varphi_{17}$\lbl{A1}}
\BIC{$(x \in \mathbb{R} \to 0 < x \to 0 < x) \to x \in \mathbb{R} \to 0 < x \to 0 \neq x$}
\BIC{$x \in \mathbb{R} \to 0 < x \to 0 \neq x$}
\AXC{$\varphi_{13}$\lbl{A1}}
\BIC{$(x = 0 \to 0 = x) \to x \in \mathbb{R} \to 0 < x \to x \neq 0$}
\BIC{$x \in \mathbb{R} \to 0 < x \to x \neq 0$}
\end{ptree}
\noindent$\Pi_{8}$:
\begin{ptree}
\AXC{$x \in \mathbb{R} \to 0 < x \to x^{-1} < 0 \to x \in \mathbb{R}$\lbl{A1}}
\AXC{$\forall x \in \mathbb{R}\,(x \neq 0 \to x^{-1} \in \mathbb{R})$\lbl{(R1)}}
\AXC{$\varphi_{20}$\lbl{A2}}
\BIC{$x \in \mathbb{R} \to x \neq 0 \to x^{-1} \in \mathbb{R}$}
\AXC{$\varphi_{64}$\lbl{A1}}
\BIC{$(x \in \mathbb{R} \to 0 < x \to x^{-1} < 0 \to x \in \mathbb{R}) \to \varphi_{24}$}
\BIC{$\varphi_{24}$}
\end{ptree}
\noindent$\Pi_{9}$:
\begin{ptree}
\AXC{$\stackrel{\Pi_{7}}{\vdots}$}\noLine
\UIC{$x \in \mathbb{R} \to 0 < x \to x \neq 0$}
\AXC{$\stackrel{\Pi_{8}}{\vdots}$}\noLine
\UIC{$\varphi_{24}$}
\BIC{$x \in \mathbb{R} \to 0 < x \to x^{-1} < 0 \to x^{-1} \in \mathbb{R}$}
\end{ptree}
\noindent$\Pi_{10}$:
\begin{ptree}
\AXC{$\forall x\,\forall y\,(x \in \mathbb{R} \to y \in \mathbb{R} \to x < y \to -y < -x)$\lbl{[F22]}}
\AXC{$\varphi_{60}$\lbl{A2}}
\BIC{$\forall y\,(x^{-1} \in \mathbb{R} \to y \in \mathbb{R} \to x^{-1} < y \to -y < -x^{-1})$}
\AXC{$\varphi_{62}$\lbl{A2}}
\BIC{$x^{-1} \in \mathbb{R} \to 0 \in \mathbb{R} \to x^{-1} < 0 \to -0 < -x^{-1}$}
\AXC{$\varphi_{76}$\lbl{A1}}
\BIC{$(x \in \mathbb{R} \to 0 < x \to x^{-1} < 0 \to x^{-1} \in \mathbb{R}) \to 0 \in \mathbb{R} \to \varphi_{45}$}
\end{ptree}
\noindent$\Pi_{11}$:
\begin{ptree}
\AXC{$0 \in \mathbb{R}$\lbl{(R1)}}
\AXC{$\stackrel{\Pi_{9}}{\vdots}$}\noLine
\UIC{$x \in \mathbb{R} \to 0 < x \to x^{-1} < 0 \to x^{-1} \in \mathbb{R}$}
\AXC{$\stackrel{\Pi_{10}}{\vdots}$}\noLine
\UIC{$(x \in \mathbb{R} \to 0 < x \to x^{-1} < 0 \to x^{-1} \in \mathbb{R}) \to 0 \in \mathbb{R} \to \varphi_{45}$}
\BIC{$0 \in \mathbb{R} \to \varphi_{45}$}
\BIC{$\varphi_{45}$}
\end{ptree}
\noindent$\Pi_{12}$:
\begin{ptree}
\AXC{$x \in \mathbb{R} \to 0 < x \to x^{-1} < 0 \to x^{-1} < 0$\lbl{A1}}
\AXC{$\stackrel{\Pi_{11}}{\vdots}$}\noLine
\UIC{$\varphi_{45}$}
\BIC{$x \in \mathbb{R} \to 0 < x \to x^{-1} < 0 \to -0 < -x^{-1}$}
\AXC{$-0 = 0$\lbl{[F6]}}
\AXC{$\varphi_{66}$\lbl{A1}}
\BIC{$(x \in \mathbb{R} \to 0 < x \to x^{-1} < 0 \to -0 < -x^{-1}) \to \varphi_{43}$}
\BIC{$\varphi_{43}$}
\end{ptree}
\noindent$\Pi_{13}$:
\begin{ptree}
\AXC{$\forall x \in \mathbb{R}\,(-x \in \mathbb{R})$\lbl{(R1)}}
\AXC{$\forall x \in \mathbb{R}\,(-x \in \mathbb{R}) \to x^{-1} \in \mathbb{R} \to -x^{-1} \in \mathbb{R}$\lbl{A2}}
\BIC{$x^{-1} \in \mathbb{R} \to -x^{-1} \in \mathbb{R}$}
\AXC{$(x^{-1} \in \mathbb{R} \to -x^{-1} \in \mathbb{R}) \to \varphi_{49}$\lbl{A1}}
\BIC{$\varphi_{49}$}
\end{ptree}
\noindent$\Pi_{14}$:
\begin{ptree}
\AXC{$\forall x\,\forall y\,(x \in \mathbb{R} \to y \in \mathbb{R} \to 0 < x \to 0 < y \to 0 < x \cdot y)$\lbl{(R15)}}
\AXC{$\varphi_{58}$\lbl{A2}}
\BIC{$\forall y\,(x \in \mathbb{R} \to y \in \mathbb{R} \to 0 < x \to 0 < y \to 0 < x \cdot y)$}
\AXC{$\forall y\,(x \in \mathbb{R} \to y \in \mathbb{R} \to 0 < x \to 0 < y \to 0 < x \cdot y) \to \varphi_{28}$\lbl{A2}}
\BIC{$\varphi_{28}$}
\end{ptree}
\noindent$\Pi_{15}$:
\begin{ptree}
\AXC{$x \in \mathbb{R} \to 0 < x \to x^{-1} < 0 \to x \in \mathbb{R}$\lbl{A1}}
\AXC{$\stackrel{\Pi_{14}}{\vdots}$}\noLine
\UIC{$\varphi_{28}$}
\AXC{$\varphi_{28} \to (x \in \mathbb{R} \to 0 < x \to x^{-1} < 0 \to x \in \mathbb{R}) \to \varphi_{72}$\lbl{A1}}
\BIC{$(x \in \mathbb{R} \to 0 < x \to x^{-1} < 0 \to x \in \mathbb{R}) \to \varphi_{72}$}
\BIC{$\varphi_{72}$}
\end{ptree}
\noindent$\Pi_{16}$:
\begin{ptree}
\AXC{$x \in \mathbb{R} \to 0 < x \to 0 < x$\lbl{A1}}
\AXC{$\stackrel{\Pi_{9}}{\vdots}$}\noLine
\UIC{$x \in \mathbb{R} \to 0 < x \to x^{-1} < 0 \to x^{-1} \in \mathbb{R}$}
\AXC{$\stackrel{\Pi_{13}}{\vdots}$}\noLine
\UIC{$\varphi_{49}$}
\BIC{$x \in \mathbb{R} \to 0 < x \to x^{-1} < 0 \to -x^{-1} \in \mathbb{R}$}
\AXC{$\stackrel{\Pi_{15}}{\vdots}$}\noLine
\UIC{$\varphi_{72}$}
\BIC{$(x \in \mathbb{R} \to 0 < x \to 0 < x) \to \varphi_{52}$}
\BIC{$\varphi_{52}$}
\end{ptree}
\noindent$\Pi_{17}$:
\begin{ptree}
\AXC{$z = w \to z < -x^{-1} \to w < -x^{-1}$\lbl{A6}}
\UIC{$\forall w\,(z = w \to z < -x^{-1} \to w < -x^{-1})$}
\UIC{$\forall z\,\forall w\,(z = w \to z < -x^{-1} \to w < -x^{-1})$}
\AXC{$\varphi_{42}$\lbl{A2}}
\BIC{$\forall w\,(-0 = w \to -0 < -x^{-1} \to w < -x^{-1})$}
\AXC{$\varphi_{47}$\lbl{A2}}
\BIC{$-0 = 0 \to -0 < -x^{-1} \to 0 < -x^{-1}$}
\AXC{$\stackrel{\Pi_{12}}{\vdots}$}\noLine
\UIC{$\varphi_{43}$}
\BIC{$x \in \mathbb{R} \to 0 < x \to x^{-1} < 0 \to 0 < -x^{-1}$}
\AXC{$\stackrel{\Pi_{16}}{\vdots}$}\noLine
\UIC{$\varphi_{52}$}
\BIC{$x \in \mathbb{R} \to 0 < x \to x^{-1} < 0 \to 0 < x \cdot -x^{-1}$}
\end{ptree}
\noindent$\Pi_{18}$:
\begin{ptree}
\AXC{$y = z \to x = y \to x = z$\lbl{A6}}
\AXC{$(y = z \to x = y \to x = z) \to x = y \to y = z \to x = z$\lbl{A1}}
\BIC{$x = y \to y = z \to x = z$}
\UIC{$\forall z\,(x = y \to y = z \to x = z)$}
\UIC{$\forall y\,\forall z\,(x = y \to y = z \to x = z)$}
\UIC{$\forall x\,\forall y\,\forall z\,(x = y \to y = z \to x = z)$}
\AXC{$\varphi_{21}$\lbl{A2}}
\BIC{$\forall y\,\forall z\,(x \cdot -x^{-1} = y \to y = z \to x \cdot -x^{-1} = z)$}
\end{ptree}
\noindent$\Pi_{19}$:
\begin{ptree}
\AXC{$\stackrel{\Pi_{18}}{\vdots}$}\noLine
\UIC{$\forall y\,\forall z\,(x \cdot -x^{-1} = y \to y = z \to x \cdot -x^{-1} = z)$}
\AXC{$\forall y\,\forall z\,(x \cdot -x^{-1} = y \to y = z \to x \cdot -x^{-1} = z) \to \varphi_{22}$\lbl{A2}}
\BIC{$\varphi_{22}$}
\AXC{$\varphi_{22} \to \varphi_{29}$\lbl{A2}}
\BIC{$\varphi_{29}$}
\end{ptree}
\noindent$\Pi_{20}$:
\begin{ptree}
\AXC{$z = w \to -(x \cdot x^{-1}) = -z \to -(x \cdot x^{-1}) = -w$\lbl{A6}}
\UIC{$\varphi_{5}$}
\UIC{$\forall z\,(\varphi_{5})$}
\AXC{$\forall z\,(\varphi_{5}) \to \varphi_{23}$\lbl{A2}}
\BIC{$\varphi_{23}$}
\AXC{$\varphi_{23} \to \varphi_{19}$\lbl{A2}}
\BIC{$\varphi_{19}$}
\end{ptree}
\noindent$\Pi_{21}$:
\begin{ptree}
\AXC{$\stackrel{\Pi_{7}}{\vdots}$}\noLine
\UIC{$x \in \mathbb{R} \to 0 < x \to x \neq 0$}
\AXC{$x \in \mathbb{R} \to 0 < x \to x \in \mathbb{R}$\lbl{A1}}
\AXC{$\forall x \in \mathbb{R}\,(x \neq 0 \to x \cdot x^{-1} = 1)$\lbl{(R9)}}
\AXC{$\varphi_{16}$\lbl{A2}}
\BIC{$x \in \mathbb{R} \to x \neq 0 \to x \cdot x^{-1} = 1$}
\AXC{$\varphi_{61}$\lbl{A1}}
\BIC{$(x \in \mathbb{R} \to 0 < x \to x \in \mathbb{R}) \to \varphi_{9}$}
\BIC{$\varphi_{9}$}
\BIC{$x \in \mathbb{R} \to 0 < x \to x \cdot x^{-1} = 1$}
\end{ptree}
\noindent$\Pi_{22}$:
\begin{ptree}
\AXC{$\stackrel{\Pi_{21}}{\vdots}$}\noLine
\UIC{$x \in \mathbb{R} \to 0 < x \to x \cdot x^{-1} = 1$}
\AXC{$(x \in \mathbb{R} \to 0 < x \to x \cdot x^{-1} = 1) \to \varphi_{65}$\lbl{A1}}
\BIC{$\varphi_{65}$}
\end{ptree}
\noindent$\Pi_{23}$:
\begin{ptree}
\AXC{$\stackrel{\Pi_{20}}{\vdots}$}\noLine
\UIC{$\varphi_{19}$}
\AXC{$x = x$\lbl{A5}}
\UIC{$\forall x\,(x = x)$}
\AXC{$\forall x\,(x = x) \to -(x \cdot x^{-1}) = -(x \cdot x^{-1})$\lbl{A2}}
\BIC{$-(x \cdot x^{-1}) = -(x \cdot x^{-1})$}
\AXC{$\stackrel{\Pi_{22}}{\vdots}$}\noLine
\UIC{$\varphi_{65}$}
\BIC{$\varphi_{19} \to x \in \mathbb{R} \to 0 < x \to x^{-1} < 0 \to -(x \cdot x^{-1}) = -1$}
\BIC{$x \in \mathbb{R} \to 0 < x \to x^{-1} < 0 \to -(x \cdot x^{-1}) = -1$}
\end{ptree}
\noindent$\Pi_{24}$:
\begin{ptree}
\AXC{$\forall x\,\forall y\,(x \in \mathbb{R} \to y \in \mathbb{R} \to x \cdot -y = -(x \cdot y))$\lbl{[F10]}}
\AXC{$\varphi_{50}$\lbl{A2}}
\BIC{$\forall y\,(x \in \mathbb{R} \to y \in \mathbb{R} \to x \cdot -y = -(x \cdot y))$}
\AXC{$\varphi_{56}$\lbl{A2}}
\BIC{$x \in \mathbb{R} \to x^{-1} \in \mathbb{R} \to x \cdot -x^{-1} = -(x \cdot x^{-1})$}
\AXC{$\varphi_{75}$\lbl{A1}}
\BIC{$(x \in \mathbb{R} \to 0 < x \to x^{-1} < 0 \to x \in \mathbb{R}) \to \varphi_{59}$}
\end{ptree}
\noindent$\Pi_{25}$:
\begin{ptree}
\AXC{$x \in \mathbb{R} \to 0 < x \to x^{-1} < 0 \to x \in \mathbb{R}$\lbl{A1}}
\AXC{$\stackrel{\Pi_{24}}{\vdots}$}\noLine
\UIC{$(x \in \mathbb{R} \to 0 < x \to x^{-1} < 0 \to x \in \mathbb{R}) \to \varphi_{59}$}
\BIC{$\varphi_{59}$}
\end{ptree}
\noindent$\Pi_{26}$:
\begin{ptree}
\AXC{$\stackrel{\Pi_{23}}{\vdots}$}\noLine
\UIC{$x \in \mathbb{R} \to 0 < x \to x^{-1} < 0 \to -(x \cdot x^{-1}) = -1$}
\AXC{$\stackrel{\Pi_{9}}{\vdots}$}\noLine
\UIC{$x \in \mathbb{R} \to 0 < x \to x^{-1} < 0 \to x^{-1} \in \mathbb{R}$}
\AXC{$\stackrel{\Pi_{25}}{\vdots}$}\noLine
\UIC{$\varphi_{59}$}
\BIC{$x \in \mathbb{R} \to 0 < x \to x^{-1} < 0 \to x \cdot -x^{-1} = -(x \cdot x^{-1})$}
\AXC{$\varphi_{77}$\lbl{A1}}
\BIC{$\varphi_{73}$}
\BIC{$\varphi_{29} \to x \in \mathbb{R} \to 0 < x \to x^{-1} < 0 \to x \cdot -x^{-1} = -1$}
\end{ptree}
\noindent$\Pi_{27}$:
\begin{ptree}
\AXC{$\stackrel{\Pi_{17}}{\vdots}$}\noLine
\UIC{$x \in \mathbb{R} \to 0 < x \to x^{-1} < 0 \to 0 < x \cdot -x^{-1}$}
\AXC{$\stackrel{\Pi_{19}}{\vdots}$}\noLine
\UIC{$\varphi_{29}$}
\AXC{$\stackrel{\Pi_{26}}{\vdots}$}\noLine
\UIC{$\varphi_{29} \to x \in \mathbb{R} \to 0 < x \to x^{-1} < 0 \to x \cdot -x^{-1} = -1$}
\BIC{$x \in \mathbb{R} \to 0 < x \to x^{-1} < 0 \to x \cdot -x^{-1} = -1$}
\AXC{$\varphi_{74}$\lbl{A1}}
\BIC{$(x \in \mathbb{R} \to 0 < x \to x^{-1} < 0 \to 0 < x \cdot -x^{-1}) \to \varphi_{51}$}
\BIC{$\varphi_{51}$}
\end{ptree}
\noindent$\Pi_{28}$:
\begin{ptree}
\AXC{$\forall x\,\forall y\,\forall z\,(x < y \to y < z \to x < z)$\lbl{[O2]}}
\AXC{$\varphi_{12}$\lbl{A2}}
\BIC{$\forall y\,\forall z\,(0 < y \to y < z \to 0 < z)$}
\AXC{$\varphi_{15}$\lbl{A2}}
\BIC{$\forall z\,(0 < -1 \to -1 < z \to 0 < z)$}
\AXC{$\varphi_{18}$\lbl{A2}}
\BIC{$0 < -1 \to -1 < 0 \to 0 < 0$}
\AXC{$\varphi_{67}$\lbl{A1}}
\BIC{$(x \in \mathbb{R} \to 0 < x \to x^{-1} < 0 \to 0 < -1) \to \varphi_{38}$}
\end{ptree}
\noindent$\Pi_{29}$:
\begin{ptree}
\AXC{$z = w \to 0 < z \to 0 < w$\lbl{A6}}
\UIC{$\forall w\,(z = w \to 0 < z \to 0 < w)$}
\UIC{$\forall z\,\forall w\,(z = w \to 0 < z \to 0 < w)$}
\AXC{$\varphi_{27}$\lbl{A2}}
\BIC{$\forall w\,(x \cdot -x^{-1} = w \to 0 < x \cdot -x^{-1} \to 0 < w)$}
\AXC{$\varphi_{55}$\lbl{A2}}
\BIC{$x \cdot -x^{-1} = -1 \to 0 < x \cdot -x^{-1} \to 0 < -1$}
\AXC{$\stackrel{\Pi_{27}}{\vdots}$}\noLine
\UIC{$\varphi_{51}$}
\BIC{$x \in \mathbb{R} \to 0 < x \to x^{-1} < 0 \to 0 < -1$}
\AXC{$\stackrel{\Pi_{28}}{\vdots}$}\noLine
\UIC{$(x \in \mathbb{R} \to 0 < x \to x^{-1} < 0 \to 0 < -1) \to \varphi_{38}$}
\BIC{$\varphi_{38}$}
\end{ptree}
\noindent$\Pi_{30}$:
\begin{ptree}
\AXC{$z = w \to -1 < z \to -1 < w$\lbl{A6}}
\UIC{$\forall w\,(z = w \to -1 < z \to -1 < w)$}
\UIC{$\forall z\,\forall w\,(z = w \to -1 < z \to -1 < w)$}
\AXC{$\varphi_{26}$\lbl{A2}}
\BIC{$\forall w\,(-0 = w \to -1 < -0 \to -1 < w)$}
\AXC{$\varphi_{34}$\lbl{A2}}
\BIC{$-0 = 0 \to -1 < -0 \to -1 < 0$}
\AXC{$\stackrel{\Pi_{3}}{\vdots}$}\noLine
\UIC{$\varphi_{37}$}
\BIC{$x \in \mathbb{R} \to 0 < x \to x^{-1} < 0 \to -1 < 0$}
\AXC{$\stackrel{\Pi_{29}}{\vdots}$}\noLine
\UIC{$\varphi_{38}$}
\BIC{$x \in \mathbb{R} \to 0 < x \to x^{-1} < 0 \to 0 < 0$}
\AXC{$\varphi_{41}$\lbl{A1}}
\BIC{$\neg 0 < 0 \to x \in \mathbb{R} \to 0 < x \to x^{-1} < 0 \to 0 < x^{-1}$}
\end{ptree}
\noindent$\Pi_{31}$:
\begin{ptree}
\AXC{$\forall x\,(\neg x < x)$\lbl{[O1]}}
\AXC{$\forall x\,(\neg x < x) \to \neg 0 < 0$\lbl{A2}}
\BIC{$\neg 0 < 0$}
\AXC{$\stackrel{\Pi_{30}}{\vdots}$}\noLine
\UIC{$\neg 0 < 0 \to x \in \mathbb{R} \to 0 < x \to x^{-1} < 0 \to 0 < x^{-1}$}
\BIC{$x \in \mathbb{R} \to 0 < x \to x^{-1} < 0 \to 0 < x^{-1}$}
\end{ptree}
\noindent$\Pi_{32}$:
\begin{ptree}
\AXC{$\stackrel{\Pi_{4}}{\vdots}$}\noLine
\UIC{$\forall x\,\forall y\,(x = y \to y = x)$}
\AXC{$\forall x\,\forall y\,(x = y \to y = x) \to \forall y\,(1 = y \to y = 1)$\lbl{A2}}
\BIC{$\forall y\,(1 = y \to y = 1)$}
\AXC{$\forall y\,(1 = y \to y = 1) \to 1 = 0 \to 0 = 1$\lbl{A2}}
\BIC{$1 = 0 \to 0 = 1$}
\end{ptree}
\noindent$\Pi_{33}$:
\begin{ptree}
\AXC{$y = z \to x = y \to x = z$\lbl{A6}}
\AXC{$(y = z \to x = y \to x = z) \to x = y \to y = z \to x = z$\lbl{A1}}
\BIC{$x = y \to y = z \to x = z$}
\UIC{$\forall z\,(x = y \to y = z \to x = z)$}
\UIC{$\forall y\,\forall z\,(x = y \to y = z \to x = z)$}
\UIC{$\forall x\,\forall y\,\forall z\,(x = y \to y = z \to x = z)$}
\AXC{$\varphi_{1}$\lbl{A2}}
\BIC{$\forall y\,\forall z\,(1 = y \to y = z \to 1 = z)$}
\AXC{$\varphi_{3}$\lbl{A2}}
\BIC{$\forall z\,(1 = x \cdot 0 \to x \cdot 0 = z \to 1 = z)$}
\AXC{$\varphi_{4}$\lbl{A2}}
\BIC{$1 = x \cdot 0 \to x \cdot 0 = 0 \to 1 = 0$}
\end{ptree}
\noindent$\Pi_{34}$:
\begin{ptree}
\AXC{$\forall x \in \mathbb{R}\,(x \cdot 0 = 0)$\lbl{[F8]}}
\AXC{$\forall x \in \mathbb{R}\,(x \cdot 0 = 0) \to x \in \mathbb{R} \to x \cdot 0 = 0$\lbl{A2}}
\BIC{$x \in \mathbb{R} \to x \cdot 0 = 0$}
\AXC{$(x \in \mathbb{R} \to x \cdot 0 = 0) \to \varphi_{32}$\lbl{A1}}
\BIC{$\varphi_{32}$}
\end{ptree}
\noindent$\Pi_{35}$:
\begin{ptree}
\AXC{$y = z \to x = y \to x = z$\lbl{A6}}
\AXC{$(y = z \to x = y \to x = z) \to x = y \to y = z \to x = z$\lbl{A1}}
\BIC{$x = y \to y = z \to x = z$}
\UIC{$\forall z\,(x = y \to y = z \to x = z)$}
\UIC{$\forall y\,\forall z\,(x = y \to y = z \to x = z)$}
\UIC{$\forall x\,\forall y\,\forall z\,(x = y \to y = z \to x = z)$}
\AXC{$\varphi_{1}$\lbl{A2}}
\BIC{$\forall y\,\forall z\,(1 = y \to y = z \to 1 = z)$}
\AXC{$\varphi_{7}$\lbl{A2}}
\BIC{$\forall z\,(1 = x \cdot x^{-1} \to x \cdot x^{-1} = z \to 1 = z)$}
\AXC{$\varphi_{30}$\lbl{A2}}
\BIC{$1 = x \cdot x^{-1} \to x \cdot x^{-1} = x \cdot 0 \to 1 = x \cdot 0$}
\end{ptree}
\noindent$\Pi_{36}$:
\begin{ptree}
\AXC{$z = w \to x \cdot x^{-1} = x \cdot z \to x \cdot x^{-1} = x \cdot w$\lbl{A6}}
\UIC{$\forall w\,(z = w \to x \cdot x^{-1} = x \cdot z \to x \cdot x^{-1} = x \cdot w)$}
\UIC{$\forall z\,\forall w\,(z = w \to x \cdot x^{-1} = x \cdot z \to x \cdot x^{-1} = x \cdot w)$}
\AXC{$\varphi_{53}$\lbl{A2}}
\BIC{$\varphi_{6}$}
\AXC{$\varphi_{6} \to x^{-1} = 0 \to x \cdot x^{-1} = x \cdot x^{-1} \to x \cdot x^{-1} = x \cdot 0$\lbl{A2}}
\BIC{$x^{-1} = 0 \to x \cdot x^{-1} = x \cdot x^{-1} \to x \cdot x^{-1} = x \cdot 0$}
\end{ptree}
\noindent$\Pi_{37}$:
\begin{ptree}
\AXC{$\stackrel{\Pi_{4}}{\vdots}$}\noLine
\UIC{$\forall x\,\forall y\,(x = y \to y = x)$}
\AXC{$\forall x\,\forall y\,(x = y \to y = x) \to \forall y\,(0 = y \to y = 0)$\lbl{A2}}
\BIC{$\forall y\,(0 = y \to y = 0)$}
\AXC{$\forall y\,(0 = y \to y = 0) \to 0 = x^{-1} \to x^{-1} = 0$\lbl{A2}}
\BIC{$0 = x^{-1} \to x^{-1} = 0$}
\end{ptree}
\noindent$\Pi_{38}$:
\begin{ptree}
\AXC{$\stackrel{\Pi_{37}}{\vdots}$}\noLine
\UIC{$0 = x^{-1} \to x^{-1} = 0$}
\AXC{$x \in \mathbb{R} \to 0 < x \to 0 = x^{-1} \to 0 = x^{-1}$\lbl{A1}}
\AXC{$(x \in \mathbb{R} \to 0 < x \to 0 = x^{-1} \to 0 = x^{-1}) \to \varphi_{8}$\lbl{A1}}
\BIC{$\varphi_{8}$}
\BIC{$x \in \mathbb{R} \to 0 < x \to 0 = x^{-1} \to x^{-1} = 0$}
\end{ptree}
\noindent$\Pi_{39}$:
\begin{ptree}
\AXC{$x = x$\lbl{A5}}
\UIC{$\forall x\,(x = x)$}
\AXC{$\forall x\,(x = x) \to x \cdot x^{-1} = x \cdot x^{-1}$\lbl{A2}}
\BIC{$x \cdot x^{-1} = x \cdot x^{-1}$}
\AXC{$\stackrel{\Pi_{38}}{\vdots}$}\noLine
\UIC{$x \in \mathbb{R} \to 0 < x \to 0 = x^{-1} \to x^{-1} = 0$}
\AXC{$\varphi_{68}$\lbl{A1}}
\BIC{$x \cdot x^{-1} = x \cdot x^{-1} \to \varphi_{48}$}
\BIC{$\varphi_{48}$}
\end{ptree}
\noindent$\Pi_{40}$:
\begin{ptree}
\AXC{$\stackrel{\Pi_{21}}{\vdots}$}\noLine
\UIC{$x \in \mathbb{R} \to 0 < x \to x \cdot x^{-1} = 1$}
\AXC{$(x \in \mathbb{R} \to 0 < x \to x \cdot x^{-1} = 1) \to \varphi_{14}$\lbl{A1}}
\BIC{$\varphi_{14}$}
\end{ptree}
\noindent$\Pi_{41}$:
\begin{ptree}
\AXC{$\stackrel{\Pi_{4}}{\vdots}$}\noLine
\UIC{$\forall x\,\forall y\,(x = y \to y = x)$}
\AXC{$\varphi_{2}$\lbl{A2}}
\BIC{$\forall y\,(x \cdot x^{-1} = y \to y = x \cdot x^{-1})$}
\AXC{$\varphi_{10}$\lbl{A2}}
\BIC{$x \cdot x^{-1} = 1 \to 1 = x \cdot x^{-1}$}
\AXC{$\stackrel{\Pi_{40}}{\vdots}$}\noLine
\UIC{$\varphi_{14}$}
\BIC{$x \in \mathbb{R} \to 0 < x \to 1 = x \cdot x^{-1}$}
\AXC{$\varphi_{69}$\lbl{A1}}
\BIC{$(x \in \mathbb{R} \to 0 < x \to 0 = x^{-1} \to x \cdot x^{-1} = x \cdot 0) \to \varphi_{35}$}
\end{ptree}
\noindent$\Pi_{42}$:
\begin{ptree}
\AXC{$\stackrel{\Pi_{36}}{\vdots}$}\noLine
\UIC{$x^{-1} = 0 \to x \cdot x^{-1} = x \cdot x^{-1} \to x \cdot x^{-1} = x \cdot 0$}
\AXC{$\stackrel{\Pi_{39}}{\vdots}$}\noLine
\UIC{$\varphi_{48}$}
\BIC{$x \in \mathbb{R} \to 0 < x \to 0 = x^{-1} \to x \cdot x^{-1} = x \cdot 0$}
\AXC{$\stackrel{\Pi_{41}}{\vdots}$}\noLine
\UIC{$(x \in \mathbb{R} \to 0 < x \to 0 = x^{-1} \to x \cdot x^{-1} = x \cdot 0) \to \varphi_{35}$}
\BIC{$\varphi_{35}$}
\end{ptree}
\noindent$\Pi_{43}$:
\begin{ptree}
\AXC{$x \in \mathbb{R} \to 0 < x \to 0 = x^{-1} \to x \in \mathbb{R}$\lbl{A1}}
\AXC{$\stackrel{\Pi_{34}}{\vdots}$}\noLine
\UIC{$\varphi_{32}$}
\BIC{$x \in \mathbb{R} \to 0 < x \to 0 = x^{-1} \to x \cdot 0 = 0$}
\AXC{$\stackrel{\Pi_{35}}{\vdots}$}\noLine
\UIC{$1 = x \cdot x^{-1} \to x \cdot x^{-1} = x \cdot 0 \to 1 = x \cdot 0$}
\AXC{$\stackrel{\Pi_{42}}{\vdots}$}\noLine
\UIC{$\varphi_{35}$}
\BIC{$x \in \mathbb{R} \to 0 < x \to 0 = x^{-1} \to 1 = x \cdot 0$}
\AXC{$\varphi_{63}$\lbl{A1}}
\BIC{$(x \in \mathbb{R} \to 0 < x \to 0 = x^{-1} \to x \cdot 0 = 0) \to \varphi_{11}$}
\BIC{$\varphi_{11}$}
\end{ptree}
\noindent$\Pi_{44}$:
\begin{ptree}
\AXC{$0 \neq 1$\lbl{(R8)}}
\AXC{$\stackrel{\Pi_{32}}{\vdots}$}\noLine
\UIC{$1 = 0 \to 0 = 1$}
\AXC{$\stackrel{\Pi_{33}}{\vdots}$}\noLine
\UIC{$1 = x \cdot 0 \to x \cdot 0 = 0 \to 1 = 0$}
\AXC{$\stackrel{\Pi_{43}}{\vdots}$}\noLine
\UIC{$\varphi_{11}$}
\BIC{$x \in \mathbb{R} \to 0 < x \to 0 = x^{-1} \to 1 = 0$}
\AXC{$\varphi_{33}$\lbl{A1}}
\BIC{$(1 = 0 \to 0 = 1) \to x \in \mathbb{R} \to 0 < x \to 0 = x^{-1} \to 0 = 1$}
\BIC{$x \in \mathbb{R} \to 0 < x \to 0 = x^{-1} \to 0 = 1$}
\AXC{$\varphi_{36}$\lbl{A1}}
\BIC{$0 \neq 1 \to x \in \mathbb{R} \to 0 < x \to 0 = x^{-1} \to 0 < x^{-1}$}
\BIC{$x \in \mathbb{R} \to 0 < x \to 0 = x^{-1} \to 0 < x^{-1}$}
\end{ptree}
\noindent$\Pi_{45}$:
\begin{ptree}
\AXC{$x \in \mathbb{R} \to 0 < x \to x \in \mathbb{R}$\lbl{A1}}
\AXC{$\forall x \in \mathbb{R}\,(x \neq 0 \to x^{-1} \in \mathbb{R})$\lbl{(R1)}}
\AXC{$\varphi_{20}$\lbl{A2}}
\BIC{$x \in \mathbb{R} \to x \neq 0 \to x^{-1} \in \mathbb{R}$}
\AXC{$(x \in \mathbb{R} \to x \neq 0 \to x^{-1} \in \mathbb{R}) \to \varphi_{40}$\lbl{A1}}
\BIC{$\varphi_{40}$}
\BIC{$(x \in \mathbb{R} \to 0 < x \to x \neq 0) \to x \in \mathbb{R} \to 0 < x \to x^{-1} \in \mathbb{R}$}
\end{ptree}
\noindent$\Pi_{46}$:
\begin{ptree}
\AXC{$0 \in \mathbb{R}$\lbl{(R1)}}
\AXC{$\forall x\,\forall y\,(x \in \mathbb{R} \to y \in \mathbb{R} \to x < y \lor x = y \lor y < x)$\lbl{(R13)}}
\AXC{$\varphi_{44}$\lbl{A2}}
\BIC{$\forall y\,(0 \in \mathbb{R} \to y \in \mathbb{R} \to 0 < y \lor 0 = y \lor y < 0)$}
\AXC{$\varphi_{57}$\lbl{A2}}
\BIC{$0 \in \mathbb{R} \to x^{-1} \in \mathbb{R} \to 0 < x^{-1} \lor 0 = x^{-1} \lor x^{-1} < 0$}
\AXC{$\varphi_{70}$\lbl{A1}}
\BIC{$0 \in \mathbb{R} \to \varphi_{39}$}
\BIC{$\varphi_{39}$}
\end{ptree}
\noindent$\Pi_{47}$:
\begin{ptree}
\AXC{$\stackrel{\Pi_{7}}{\vdots}$}\noLine
\UIC{$x \in \mathbb{R} \to 0 < x \to x \neq 0$}
\AXC{$\stackrel{\Pi_{45}}{\vdots}$}\noLine
\UIC{$(x \in \mathbb{R} \to 0 < x \to x \neq 0) \to x \in \mathbb{R} \to 0 < x \to x^{-1} \in \mathbb{R}$}
\BIC{$x \in \mathbb{R} \to 0 < x \to x^{-1} \in \mathbb{R}$}
\AXC{$\stackrel{\Pi_{46}}{\vdots}$}\noLine
\UIC{$\varphi_{39}$}
\BIC{$x \in \mathbb{R} \to 0 < x \to 0 < x^{-1} \lor 0 = x^{-1} \lor x^{-1} < 0$}
\AXC{$\varphi_{71}$\lbl{A1}}
\BIC{$(x \in \mathbb{R} \to 0 < x \to 0 = x^{-1} \to 0 < x^{-1}) \to \varphi_{25}$}
\end{ptree}
\noindent$\Pi_{48}$:
\begin{ptree}
\AXC{$\stackrel{\Pi_{44}}{\vdots}$}\noLine
\UIC{$x \in \mathbb{R} \to 0 < x \to 0 = x^{-1} \to 0 < x^{-1}$}
\AXC{$\stackrel{\Pi_{47}}{\vdots}$}\noLine
\UIC{$(x \in \mathbb{R} \to 0 < x \to 0 = x^{-1} \to 0 < x^{-1}) \to \varphi_{25}$}
\BIC{$\varphi_{25}$}
\end{ptree}
\noindent 主木:
\begin{ptree}
\AXC{$\stackrel{\Pi_{31}}{\vdots}$}\noLine
\UIC{$x \in \mathbb{R} \to 0 < x \to x^{-1} < 0 \to 0 < x^{-1}$}
\AXC{$\stackrel{\Pi_{48}}{\vdots}$}\noLine
\UIC{$\varphi_{25}$}
\BIC{$x \in \mathbb{R} \to 0 < x \to 0 < x^{-1}$}
\UIC{$\forall x \in \mathbb{R}\,(0 < x \to 0 < x^{-1})$}
\end{ptree}
\begin{abbrevlist}
\item $\varphi_{1} :\equiv \forall x\,\forall y\,\forall z\,(x = y \to y = z \to x = z) \to \forall y\,\forall z\,(1 = y \to y = z \to 1 = z)$
\item $\varphi_{2} :\equiv \forall x\,\forall y\,(x = y \to y = x) \to \forall y\,(x \cdot x^{-1} = y \to y = x \cdot x^{-1})$
\item $\varphi_{3} :\equiv \forall y\,\forall z\,(1 = y \to y = z \to 1 = z) \to \forall z\,(1 = x \cdot 0 \to x \cdot 0 = z \to 1 = z)$
\item $\varphi_{4} :\equiv \forall z\,(1 = x \cdot 0 \to x \cdot 0 = z \to 1 = z) \to 1 = x \cdot 0 \to x \cdot 0 = 0 \to 1 = 0$
\item $\varphi_{5} :\equiv \forall w\,(z = w \to -(x \cdot x^{-1}) = -z \to -(x \cdot x^{-1}) = -w)$
\item $\varphi_{6} :\equiv \forall w\,(x^{-1} = w \to x \cdot x^{-1} = x \cdot x^{-1} \to x \cdot x^{-1} = x \cdot w)$
\item $\varphi_{7} :\equiv \forall y\,\forall z\,(1 = y \to y = z \to 1 = z) \to \forall z\,(1 = x \cdot x^{-1} \to x \cdot x^{-1} = z \to 1 = z)$
\item $\varphi_{8} :\equiv (0 = x^{-1} \to x^{-1} = 0) \to x \in \mathbb{R} \to 0 < x \to 0 = x^{-1} \to x^{-1} = 0$
\item $\varphi_{9} :\equiv (x \in \mathbb{R} \to 0 < x \to x \neq 0) \to x \in \mathbb{R} \to 0 < x \to x \cdot x^{-1} = 1$
\item $\varphi_{10} :\equiv \forall y\,(x \cdot x^{-1} = y \to y = x \cdot x^{-1}) \to x \cdot x^{-1} = 1 \to 1 = x \cdot x^{-1}$
\item $\varphi_{11} :\equiv (1 = x \cdot 0 \to x \cdot 0 = 0 \to 1 = 0) \to x \in \mathbb{R} \to 0 < x \to 0 = x^{-1} \to 1 = 0$
\item $\varphi_{12} :\equiv \forall x\,\forall y\,\forall z\,(x < y \to y < z \to x < z) \to \forall y\,\forall z\,(0 < y \to y < z \to 0 < z)$
\item $\varphi_{13} :\equiv (x \in \mathbb{R} \to 0 < x \to 0 \neq x) \to (x = 0 \to 0 = x) \to x \in \mathbb{R} \to 0 < x \to x \neq 0$
\item $\varphi_{14} :\equiv (x \cdot x^{-1} = 1 \to 1 = x \cdot x^{-1}) \to x \in \mathbb{R} \to 0 < x \to 1 = x \cdot x^{-1}$
\item $\varphi_{15} :\equiv \forall y\,\forall z\,(0 < y \to y < z \to 0 < z) \to \forall z\,(0 < -1 \to -1 < z \to 0 < z)$
\item $\varphi_{16} :\equiv \forall x \in \mathbb{R}\,(x \neq 0 \to x \cdot x^{-1} = 1) \to x \in \mathbb{R} \to x \neq 0 \to x \cdot x^{-1} = 1$
\item $\varphi_{17} :\equiv (0 < x \to 0 \neq x) \to (x \in \mathbb{R} \to 0 < x \to 0 < x) \to x \in \mathbb{R} \to 0 < x \to 0 \neq x$
\item $\varphi_{18} :\equiv \forall z\,(0 < -1 \to -1 < z \to 0 < z) \to 0 < -1 \to -1 < 0 \to 0 < 0$
\item $\varphi_{19} :\equiv x \cdot x^{-1} = 1 \to -(x \cdot x^{-1}) = -(x \cdot x^{-1}) \to -(x \cdot x^{-1}) = -1$
\item $\varphi_{20} :\equiv \forall x \in \mathbb{R}\,(x \neq 0 \to x^{-1} \in \mathbb{R}) \to x \in \mathbb{R} \to x \neq 0 \to x^{-1} \in \mathbb{R}$
\item $\varphi_{21} :\equiv \forall x\,\forall y\,\forall z\,(x = y \to y = z \to x = z) \to \forall y\,\forall z\,(x \cdot -x^{-1} = y \to y = z \to x \cdot -x^{-1} = z)$
\item $\varphi_{22} :\equiv \forall z\,(x \cdot -x^{-1} = -(x \cdot x^{-1}) \to -(x \cdot x^{-1}) = z \to x \cdot -x^{-1} = z)$
\item $\varphi_{23} :\equiv \forall w\,(x \cdot x^{-1} = w \to -(x \cdot x^{-1}) = -(x \cdot x^{-1}) \to -(x \cdot x^{-1}) = -w)$
\item $\varphi_{24} :\equiv (x \in \mathbb{R} \to 0 < x \to x \neq 0) \to x \in \mathbb{R} \to 0 < x \to x^{-1} < 0 \to x^{-1} \in \mathbb{R}$
\item $\varphi_{25} :\equiv (x \in \mathbb{R} \to 0 < x \to x^{-1} < 0 \to 0 < x^{-1}) \to x \in \mathbb{R} \to 0 < x \to 0 < x^{-1}$
\item $\varphi_{26} :\equiv \forall z\,\forall w\,(z = w \to -1 < z \to -1 < w) \to \forall w\,(-0 = w \to -1 < -0 \to -1 < w)$
\item $\varphi_{27} :\equiv \forall z\,\forall w\,(z = w \to 0 < z \to 0 < w) \to \forall w\,(x \cdot -x^{-1} = w \to 0 < x \cdot -x^{-1} \to 0 < w)$
\item $\varphi_{28} :\equiv x \in \mathbb{R} \to -x^{-1} \in \mathbb{R} \to 0 < x \to 0 < -x^{-1} \to 0 < x \cdot -x^{-1}$
\item $\varphi_{29} :\equiv x \cdot -x^{-1} = -(x \cdot x^{-1}) \to -(x \cdot x^{-1}) = -1 \to x \cdot -x^{-1} = -1$
\item $\varphi_{30} :\equiv \forall z\,(1 = x \cdot x^{-1} \to x \cdot x^{-1} = z \to 1 = z) \to 1 = x \cdot x^{-1} \to x \cdot x^{-1} = x \cdot 0 \to 1 = x \cdot 0$
\item $\varphi_{31} :\equiv 0 \in \mathbb{R} \to 1 \in \mathbb{R} \to 0 < 1 \to x \in \mathbb{R} \to 0 < x \to x^{-1} < 0 \to -1 < -0$
\item $\varphi_{32} :\equiv (x \in \mathbb{R} \to 0 < x \to 0 = x^{-1} \to x \in \mathbb{R}) \to x \in \mathbb{R} \to 0 < x \to 0 = x^{-1} \to x \cdot 0 = 0$
\item $\varphi_{33} :\equiv (x \in \mathbb{R} \to 0 < x \to 0 = x^{-1} \to 1 = 0) \to (1 = 0 \to 0 = 1) \to x \in \mathbb{R} \to 0 < x \to 0 = x^{-1} \to 0 = 1$
\item $\varphi_{34} :\equiv \forall w\,(-0 = w \to -1 < -0 \to -1 < w) \to -0 = 0 \to -1 < -0 \to -1 < 0$
\item $\varphi_{35} :\equiv (1 = x \cdot x^{-1} \to x \cdot x^{-1} = x \cdot 0 \to 1 = x \cdot 0) \to x \in \mathbb{R} \to 0 < x \to 0 = x^{-1} \to 1 = x \cdot 0$
\item $\varphi_{36} :\equiv (x \in \mathbb{R} \to 0 < x \to 0 = x^{-1} \to 0 = 1) \to 0 \neq 1 \to x \in \mathbb{R} \to 0 < x \to 0 = x^{-1} \to 0 < x^{-1}$
\item $\varphi_{37} :\equiv (-0 = 0 \to -1 < -0 \to -1 < 0) \to x \in \mathbb{R} \to 0 < x \to x^{-1} < 0 \to -1 < 0$
\item $\varphi_{38} :\equiv (x \in \mathbb{R} \to 0 < x \to x^{-1} < 0 \to -1 < 0) \to x \in \mathbb{R} \to 0 < x \to x^{-1} < 0 \to 0 < 0$
\item $\varphi_{39} :\equiv (x \in \mathbb{R} \to 0 < x \to x^{-1} \in \mathbb{R}) \to x \in \mathbb{R} \to 0 < x \to 0 < x^{-1} \lor 0 = x^{-1} \lor x^{-1} < 0$
\item $\varphi_{40} :\equiv (x \in \mathbb{R} \to 0 < x \to x \in \mathbb{R}) \to (x \in \mathbb{R} \to 0 < x \to x \neq 0) \to x \in \mathbb{R} \to 0 < x \to x^{-1} \in \mathbb{R}$
\item $\varphi_{41} :\equiv (x \in \mathbb{R} \to 0 < x \to x^{-1} < 0 \to 0 < 0) \to \neg 0 < 0 \to x \in \mathbb{R} \to 0 < x \to x^{-1} < 0 \to 0 < x^{-1}$
\item $\varphi_{42} :\equiv \forall z\,\forall w\,(z = w \to z < -x^{-1} \to w < -x^{-1}) \to \forall w\,(-0 = w \to -0 < -x^{-1} \to w < -x^{-1})$
\item $\varphi_{43} :\equiv (-0 = 0 \to -0 < -x^{-1} \to 0 < -x^{-1}) \to x \in \mathbb{R} \to 0 < x \to x^{-1} < 0 \to 0 < -x^{-1}$
\item $\varphi_{44} :\equiv \forall x\,\forall y\,(x \in \mathbb{R} \to y \in \mathbb{R} \to x < y \lor x = y \lor y < x) \to \forall y\,(0 \in \mathbb{R} \to y \in \mathbb{R} \to 0 < y \lor 0 = y \lor y < 0)$
\item $\varphi_{45} :\equiv (x \in \mathbb{R} \to 0 < x \to x^{-1} < 0 \to x^{-1} < 0) \to x \in \mathbb{R} \to 0 < x \to x^{-1} < 0 \to -0 < -x^{-1}$
\item $\varphi_{46} :\equiv \forall y\,(0 \in \mathbb{R} \to y \in \mathbb{R} \to 0 < y \to -y < -0) \to 0 \in \mathbb{R} \to 1 \in \mathbb{R} \to 0 < 1 \to -1 < -0$
\item $\varphi_{47} :\equiv \forall w\,(-0 = w \to -0 < -x^{-1} \to w < -x^{-1}) \to -0 = 0 \to -0 < -x^{-1} \to 0 < -x^{-1}$
\item $\varphi_{48} :\equiv (x^{-1} = 0 \to x \cdot x^{-1} = x \cdot x^{-1} \to x \cdot x^{-1} = x \cdot 0) \to x \in \mathbb{R} \to 0 < x \to 0 = x^{-1} \to x \cdot x^{-1} = x \cdot 0$
\item $\varphi_{49} :\equiv (x \in \mathbb{R} \to 0 < x \to x^{-1} < 0 \to x^{-1} \in \mathbb{R}) \to x \in \mathbb{R} \to 0 < x \to x^{-1} < 0 \to -x^{-1} \in \mathbb{R}$
\item $\varphi_{50} :\equiv \forall x\,\forall y\,(x \in \mathbb{R} \to y \in \mathbb{R} \to x \cdot -y = -(x \cdot y)) \to \forall y\,(x \in \mathbb{R} \to y \in \mathbb{R} \to x \cdot -y = -(x \cdot y))$
\item $\varphi_{51} :\equiv (x \cdot -x^{-1} = -1 \to 0 < x \cdot -x^{-1} \to 0 < -1) \to x \in \mathbb{R} \to 0 < x \to x^{-1} < 0 \to 0 < -1$
\item $\varphi_{52} :\equiv (x \in \mathbb{R} \to 0 < x \to x^{-1} < 0 \to 0 < -x^{-1}) \to x \in \mathbb{R} \to 0 < x \to x^{-1} < 0 \to 0 < x \cdot -x^{-1}$
\item $\varphi_{53} :\equiv \forall z\,\forall w\,(z = w \to x \cdot x^{-1} = x \cdot z \to x \cdot x^{-1} = x \cdot w) \to \varphi_{6}$
\item $\varphi_{54} :\equiv \forall x\,\forall y\,(x \in \mathbb{R} \to y \in \mathbb{R} \to x < y \to -y < -x) \to \forall y\,(0 \in \mathbb{R} \to y \in \mathbb{R} \to 0 < y \to -y < -0)$
\item $\varphi_{55} :\equiv \forall w\,(x \cdot -x^{-1} = w \to 0 < x \cdot -x^{-1} \to 0 < w) \to x \cdot -x^{-1} = -1 \to 0 < x \cdot -x^{-1} \to 0 < -1$
\item $\varphi_{56} :\equiv \forall y\,(x \in \mathbb{R} \to y \in \mathbb{R} \to x \cdot -y = -(x \cdot y)) \to x \in \mathbb{R} \to x^{-1} \in \mathbb{R} \to x \cdot -x^{-1} = -(x \cdot x^{-1})$
\item $\varphi_{57} :\equiv \forall y\,(0 \in \mathbb{R} \to y \in \mathbb{R} \to 0 < y \lor 0 = y \lor y < 0) \to 0 \in \mathbb{R} \to x^{-1} \in \mathbb{R} \to 0 < x^{-1} \lor 0 = x^{-1} \lor x^{-1} < 0$
\item $\varphi_{58} :\equiv \forall x\,\forall y\,(x \in \mathbb{R} \to y \in \mathbb{R} \to 0 < x \to 0 < y \to 0 < x \cdot y) \to \forall y\,(x \in \mathbb{R} \to y \in \mathbb{R} \to 0 < x \to 0 < y \to 0 < x \cdot y)$
\item $\varphi_{59} :\equiv (x \in \mathbb{R} \to 0 < x \to x^{-1} < 0 \to x^{-1} \in \mathbb{R}) \to x \in \mathbb{R} \to 0 < x \to x^{-1} < 0 \to x \cdot -x^{-1} = -(x \cdot x^{-1})$
\item $\varphi_{60} :\equiv \forall x\,\forall y\,(x \in \mathbb{R} \to y \in \mathbb{R} \to x < y \to -y < -x) \to \forall y\,(x^{-1} \in \mathbb{R} \to y \in \mathbb{R} \to x^{-1} < y \to -y < -x^{-1})$
\item $\varphi_{61} :\equiv (x \in \mathbb{R} \to x \neq 0 \to x \cdot x^{-1} = 1) \to (x \in \mathbb{R} \to 0 < x \to x \in \mathbb{R}) \to \varphi_{9}$
\item $\varphi_{62} :\equiv \forall y\,(x^{-1} \in \mathbb{R} \to y \in \mathbb{R} \to x^{-1} < y \to -y < -x^{-1}) \to x^{-1} \in \mathbb{R} \to 0 \in \mathbb{R} \to x^{-1} < 0 \to -0 < -x^{-1}$
\item $\varphi_{63} :\equiv (x \in \mathbb{R} \to 0 < x \to 0 = x^{-1} \to 1 = x \cdot 0) \to (x \in \mathbb{R} \to 0 < x \to 0 = x^{-1} \to x \cdot 0 = 0) \to \varphi_{11}$
\item $\varphi_{64} :\equiv (x \in \mathbb{R} \to x \neq 0 \to x^{-1} \in \mathbb{R}) \to (x \in \mathbb{R} \to 0 < x \to x^{-1} < 0 \to x \in \mathbb{R}) \to \varphi_{24}$
\item $\varphi_{65} :\equiv -(x \cdot x^{-1}) = -(x \cdot x^{-1}) \to \varphi_{19} \to x \in \mathbb{R} \to 0 < x \to x^{-1} < 0 \to -(x \cdot x^{-1}) = -1$
\item $\varphi_{66} :\equiv -0 = 0 \to (x \in \mathbb{R} \to 0 < x \to x^{-1} < 0 \to -0 < -x^{-1}) \to \varphi_{43}$
\item $\varphi_{67} :\equiv (0 < -1 \to -1 < 0 \to 0 < 0) \to (x \in \mathbb{R} \to 0 < x \to x^{-1} < 0 \to 0 < -1) \to \varphi_{38}$
\item $\varphi_{68} :\equiv (x \in \mathbb{R} \to 0 < x \to 0 = x^{-1} \to x^{-1} = 0) \to x \cdot x^{-1} = x \cdot x^{-1} \to \varphi_{48}$
\item $\varphi_{69} :\equiv (x \in \mathbb{R} \to 0 < x \to 1 = x \cdot x^{-1}) \to (x \in \mathbb{R} \to 0 < x \to 0 = x^{-1} \to x \cdot x^{-1} = x \cdot 0) \to \varphi_{35}$
\item $\varphi_{70} :\equiv (0 \in \mathbb{R} \to x^{-1} \in \mathbb{R} \to 0 < x^{-1} \lor 0 = x^{-1} \lor x^{-1} < 0) \to 0 \in \mathbb{R} \to \varphi_{39}$
\item $\varphi_{71} :\equiv (x \in \mathbb{R} \to 0 < x \to 0 < x^{-1} \lor 0 = x^{-1} \lor x^{-1} < 0) \to (x \in \mathbb{R} \to 0 < x \to 0 = x^{-1} \to 0 < x^{-1}) \to \varphi_{25}$
\item $\varphi_{72} :\equiv (x \in \mathbb{R} \to 0 < x \to x^{-1} < 0 \to -x^{-1} \in \mathbb{R}) \to (x \in \mathbb{R} \to 0 < x \to 0 < x) \to \varphi_{52}$
\item $\varphi_{73} :\equiv (x \in \mathbb{R} \to 0 < x \to x^{-1} < 0 \to -(x \cdot x^{-1}) = -1) \to \varphi_{29} \to x \in \mathbb{R} \to 0 < x \to x^{-1} < 0 \to x \cdot -x^{-1} = -1$
\item $\varphi_{74} :\equiv (x \in \mathbb{R} \to 0 < x \to x^{-1} < 0 \to x \cdot -x^{-1} = -1) \to (x \in \mathbb{R} \to 0 < x \to x^{-1} < 0 \to 0 < x \cdot -x^{-1}) \to \varphi_{51}$
\item $\varphi_{75} :\equiv (x \in \mathbb{R} \to x^{-1} \in \mathbb{R} \to x \cdot -x^{-1} = -(x \cdot x^{-1})) \to (x \in \mathbb{R} \to 0 < x \to x^{-1} < 0 \to x \in \mathbb{R}) \to \varphi_{59}$
\item $\varphi_{76} :\equiv (x^{-1} \in \mathbb{R} \to 0 \in \mathbb{R} \to x^{-1} < 0 \to -0 < -x^{-1}) \to (x \in \mathbb{R} \to 0 < x \to x^{-1} < 0 \to x^{-1} \in \mathbb{R}) \to 0 \in \mathbb{R} \to \varphi_{45}$
\item $\varphi_{77} :\equiv (x \in \mathbb{R} \to 0 < x \to x^{-1} < 0 \to x \cdot -x^{-1} = -(x \cdot x^{-1})) \to \varphi_{73}$
\end{abbrevlist}
\end{fullproof}

% 自動生成 (prover/build.js): HalfPos  ノード数 24
\begin{fullproof}{[F35]}{$\forall e \in \mathbb{R}\,(0 < e \to 0 < e \cdot 2^{-1})$}{24}\phantomsection\label{tree:HalfPos}
\noindent$\Pi_{1}$:
\begin{ptree}
\AXC{$0 < 2$\lbl{[F29]}}
\AXC{$2 \in \mathbb{R}$\lbl{[F2]}}
\AXC{$\forall x \in \mathbb{R}\,(0 < x \to 0 < x^{-1})$\lbl{[F34]}}
\AXC{$\varphi_{1}$\lbl{A2}}
\BIC{$2 \in \mathbb{R} \to 0 < 2 \to 0 < 2^{-1}$}
\AXC{$\varphi_{3}$\lbl{A1}}
\BIC{$2 \in \mathbb{R} \to 0 < 2 \to e \in \mathbb{R} \to 0 < e \to 0 < 2^{-1}$}
\BIC{$0 < 2 \to e \in \mathbb{R} \to 0 < e \to 0 < 2^{-1}$}
\BIC{$e \in \mathbb{R} \to 0 < e \to 0 < 2^{-1}$}
\end{ptree}
\noindent$\Pi_{2}$:
\begin{ptree}
\AXC{$\forall x\,\forall y\,(x \in \mathbb{R} \to y \in \mathbb{R} \to 0 < x \to 0 < y \to 0 < x \cdot y)$\lbl{(R15)}}
\AXC{$\varphi_{4}$\lbl{A2}}
\BIC{$\forall y\,(e \in \mathbb{R} \to y \in \mathbb{R} \to 0 < e \to 0 < y \to 0 < e \cdot y)$}
\AXC{$\varphi_{5}$\lbl{A2}}
\BIC{$e \in \mathbb{R} \to 2^{-1} \in \mathbb{R} \to 0 < e \to 0 < 2^{-1} \to 0 < e \cdot 2^{-1}$}
\AXC{$\varphi_{7}$\lbl{A1}}
\BIC{$\varphi_{6}$}
\end{ptree}
\noindent 主木:
\begin{ptree}
\AXC{$\stackrel{\Pi_{1}}{\vdots}$}\noLine
\UIC{$e \in \mathbb{R} \to 0 < e \to 0 < 2^{-1}$}
\AXC{$e \in \mathbb{R} \to 0 < e \to 0 < e$\lbl{A1}}
\AXC{$2^{-1} \in \mathbb{R}$\lbl{[F31]}}
\AXC{$e \in \mathbb{R} \to 0 < e \to e \in \mathbb{R}$\lbl{A1}}
\AXC{$\stackrel{\Pi_{2}}{\vdots}$}\noLine
\UIC{$\varphi_{6}$}
\BIC{$2^{-1} \in \mathbb{R} \to (e \in \mathbb{R} \to 0 < e \to 0 < e) \to \varphi_{2}$}
\BIC{$(e \in \mathbb{R} \to 0 < e \to 0 < e) \to \varphi_{2}$}
\BIC{$\varphi_{2}$}
\BIC{$e \in \mathbb{R} \to 0 < e \to 0 < e \cdot 2^{-1}$}
\UIC{$\forall e \in \mathbb{R}\,(0 < e \to 0 < e \cdot 2^{-1})$}
\end{ptree}
\begin{abbrevlist}
\item $\varphi_{1} :\equiv \forall x \in \mathbb{R}\,(0 < x \to 0 < x^{-1}) \to 2 \in \mathbb{R} \to 0 < 2 \to 0 < 2^{-1}$
\item $\varphi_{2} :\equiv (e \in \mathbb{R} \to 0 < e \to 0 < 2^{-1}) \to e \in \mathbb{R} \to 0 < e \to 0 < e \cdot 2^{-1}$
\item $\varphi_{3} :\equiv (2 \in \mathbb{R} \to 0 < 2 \to 0 < 2^{-1}) \to 2 \in \mathbb{R} \to 0 < 2 \to e \in \mathbb{R} \to 0 < e \to 0 < 2^{-1}$
\item $\varphi_{4} :\equiv \forall x\,\forall y\,(x \in \mathbb{R} \to y \in \mathbb{R} \to 0 < x \to 0 < y \to 0 < x \cdot y) \to \forall y\,(e \in \mathbb{R} \to y \in \mathbb{R} \to 0 < e \to 0 < y \to 0 < e \cdot y)$
\item $\varphi_{5} :\equiv \forall y\,(e \in \mathbb{R} \to y \in \mathbb{R} \to 0 < e \to 0 < y \to 0 < e \cdot y) \to e \in \mathbb{R} \to 2^{-1} \in \mathbb{R} \to 0 < e \to 0 < 2^{-1} \to 0 < e \cdot 2^{-1}$
\item $\varphi_{6} :\equiv (e \in \mathbb{R} \to 0 < e \to e \in \mathbb{R}) \to 2^{-1} \in \mathbb{R} \to (e \in \mathbb{R} \to 0 < e \to 0 < e) \to \varphi_{2}$
\item $\varphi_{7} :\equiv (e \in \mathbb{R} \to 2^{-1} \in \mathbb{R} \to 0 < e \to 0 < 2^{-1} \to 0 < e \cdot 2^{-1}) \to \varphi_{6}$
\end{abbrevlist}
\end{fullproof}

% 自動生成 (prover/build.js): HalfSum  ノード数 250
\begin{fullproof}{[F36]}{$\forall e \in \mathbb{R}\,(e \cdot 2^{-1} + e \cdot 2^{-1} = e)$}{250}\phantomsection\label{tree:HalfSum}
\noindent$\Pi_{1}$:
\begin{ptree}
\AXC{$y = z \to x = y \to x = z$\lbl{A6}}
\AXC{$(y = z \to x = y \to x = z) \to x = y \to y = z \to x = z$\lbl{A1}}
\BIC{$x = y \to y = z \to x = z$}
\UIC{$\forall z\,(x = y \to y = z \to x = z)$}
\UIC{$\forall y\,\forall z\,(x = y \to y = z \to x = z)$}
\UIC{$\forall x\,\forall y\,\forall z\,(x = y \to y = z \to x = z)$}
\AXC{$\forall x\,\forall y\,\forall z\,(x = y \to y = z \to x = z) \to \varphi_{10}$\lbl{A2}}
\BIC{$\varphi_{10}$}
\end{ptree}
\noindent$\Pi_{2}$:
\begin{ptree}
\AXC{$y = z \to x = y \to x = z$\lbl{A6}}
\AXC{$(y = z \to x = y \to x = z) \to x = y \to y = z \to x = z$\lbl{A1}}
\BIC{$x = y \to y = z \to x = z$}
\UIC{$\forall z\,(x = y \to y = z \to x = z)$}
\UIC{$\forall y\,\forall z\,(x = y \to y = z \to x = z)$}
\UIC{$\forall x\,\forall y\,\forall z\,(x = y \to y = z \to x = z)$}
\AXC{$\forall x\,\forall y\,\forall z\,(x = y \to y = z \to x = z) \to \varphi_{10}$\lbl{A2}}
\BIC{$\varphi_{10}$}
\end{ptree}
\noindent$\Pi_{3}$:
\begin{ptree}
\AXC{$x = x$\lbl{A5}}
\UIC{$\forall x\,(x = x)$}
\AXC{$\forall x\,(x = x) \to e \cdot (2 \cdot 2^{-1}) = e \cdot (2 \cdot 2^{-1})$\lbl{A2}}
\BIC{$e \cdot (2 \cdot 2^{-1}) = e \cdot (2 \cdot 2^{-1})$}
\end{ptree}
\noindent$\Pi_{4}$:
\begin{ptree}
\AXC{$2 \neq 0$\lbl{[F30]}}
\AXC{$2 \in \mathbb{R}$\lbl{[F2]}}
\AXC{$\forall x \in \mathbb{R}\,(x \neq 0 \to x \cdot x^{-1} = 1)$\lbl{(R9)}}
\AXC{$\varphi_{15}$\lbl{A2}}
\BIC{$2 \in \mathbb{R} \to 2 \neq 0 \to 2 \cdot 2^{-1} = 1$}
\AXC{$\varphi_{33}$\lbl{A1}}
\BIC{$2 \in \mathbb{R} \to 2 \neq 0 \to e \in \mathbb{R} \to 2 \cdot 2^{-1} = 1$}
\BIC{$2 \neq 0 \to e \in \mathbb{R} \to 2 \cdot 2^{-1} = 1$}
\BIC{$e \in \mathbb{R} \to 2 \cdot 2^{-1} = 1$}
\end{ptree}
\noindent$\Pi_{5}$:
\begin{ptree}
\AXC{$\stackrel{\Pi_{3}}{\vdots}$}\noLine
\UIC{$e \cdot (2 \cdot 2^{-1}) = e \cdot (2 \cdot 2^{-1})$}
\AXC{$\stackrel{\Pi_{4}}{\vdots}$}\noLine
\UIC{$e \in \mathbb{R} \to 2 \cdot 2^{-1} = 1$}
\AXC{$(e \in \mathbb{R} \to 2 \cdot 2^{-1} = 1) \to \varphi_{65}$\lbl{A1}}
\BIC{$\varphi_{65}$}
\BIC{$(2 \cdot 2^{-1} = 1 \to \varphi_{5}) \to e \in \mathbb{R} \to e \cdot (2 \cdot 2^{-1}) = e \cdot 1$}
\end{ptree}
\noindent$\Pi_{6}$:
\begin{ptree}
\AXC{$\varphi_{3}$\lbl{A6}}
\UIC{$\forall w\,(\varphi_{3})$}
\UIC{$\forall z\,\forall w\,(\varphi_{3})$}
\AXC{$\forall z\,\forall w\,(\varphi_{3}) \to \varphi_{25}$\lbl{A2}}
\BIC{$\varphi_{25}$}
\AXC{$\varphi_{25} \to 2 \cdot 2^{-1} = 1 \to \varphi_{5}$\lbl{A2}}
\BIC{$2 \cdot 2^{-1} = 1 \to \varphi_{5}$}
\AXC{$\stackrel{\Pi_{5}}{\vdots}$}\noLine
\UIC{$(2 \cdot 2^{-1} = 1 \to \varphi_{5}) \to e \in \mathbb{R} \to e \cdot (2 \cdot 2^{-1}) = e \cdot 1$}
\BIC{$e \in \mathbb{R} \to e \cdot (2 \cdot 2^{-1}) = e \cdot 1$}
\end{ptree}
\noindent$\Pi_{7}$:
\begin{ptree}
\AXC{$y = z \to x = y \to x = z$\lbl{A6}}
\AXC{$(y = z \to x = y \to x = z) \to x = y \to y = z \to x = z$\lbl{A1}}
\BIC{$x = y \to y = z \to x = z$}
\UIC{$\forall z\,(x = y \to y = z \to x = z)$}
\UIC{$\forall y\,\forall z\,(x = y \to y = z \to x = z)$}
\UIC{$\forall x\,\forall y\,\forall z\,(x = y \to y = z \to x = z)$}
\AXC{$\forall x\,\forall y\,\forall z\,(x = y \to y = z \to x = z) \to \varphi_{10}$\lbl{A2}}
\BIC{$\varphi_{10}$}
\end{ptree}
\noindent$\Pi_{8}$:
\begin{ptree}
\AXC{$2 \in \mathbb{R}$\lbl{[F2]}}
\AXC{$2^{-1} \in \mathbb{R}$\lbl{[F31]}}
\AXC{$\forall x\,\forall y\,(x \in \mathbb{R} \to y \in \mathbb{R} \to x \cdot y = y \cdot x)$\lbl{(R7)}}
\AXC{$\varphi_{46}$\lbl{A2}}
\BIC{$\forall y\,(2^{-1} \in \mathbb{R} \to y \in \mathbb{R} \to 2^{-1} \cdot y = y \cdot 2^{-1})$}
\AXC{$\varphi_{52}$\lbl{A2}}
\BIC{$2^{-1} \in \mathbb{R} \to 2 \in \mathbb{R} \to 2^{-1} \cdot 2 = 2 \cdot 2^{-1}$}
\AXC{$\varphi_{56}$\lbl{A1}}
\BIC{$2^{-1} \in \mathbb{R} \to 2 \in \mathbb{R} \to e \in \mathbb{R} \to 2^{-1} \cdot 2 = 2 \cdot 2^{-1}$}
\BIC{$2 \in \mathbb{R} \to e \in \mathbb{R} \to 2^{-1} \cdot 2 = 2 \cdot 2^{-1}$}
\BIC{$e \in \mathbb{R} \to 2^{-1} \cdot 2 = 2 \cdot 2^{-1}$}
\end{ptree}
\noindent$\Pi_{9}$:
\begin{ptree}
\AXC{$x = x$\lbl{A5}}
\UIC{$\forall x\,(x = x)$}
\AXC{$\forall x\,(x = x) \to e \cdot (2^{-1} \cdot 2) = e \cdot (2^{-1} \cdot 2)$\lbl{A2}}
\BIC{$e \cdot (2^{-1} \cdot 2) = e \cdot (2^{-1} \cdot 2)$}
\AXC{$\stackrel{\Pi_{8}}{\vdots}$}\noLine
\UIC{$e \in \mathbb{R} \to 2^{-1} \cdot 2 = 2 \cdot 2^{-1}$}
\AXC{$\varphi_{71}$\lbl{A1}}
\BIC{$e \cdot (2^{-1} \cdot 2) = e \cdot (2^{-1} \cdot 2) \to \varphi_{63}$}
\BIC{$\varphi_{63}$}
\end{ptree}
\noindent$\Pi_{10}$:
\begin{ptree}
\AXC{$\varphi_{4}$\lbl{A6}}
\UIC{$\forall w\,(\varphi_{4})$}
\UIC{$\forall z\,\forall w\,(\varphi_{4})$}
\AXC{$\forall z\,\forall w\,(\varphi_{4}) \to \varphi_{26}$\lbl{A2}}
\BIC{$\varphi_{26}$}
\AXC{$\varphi_{26} \to 2^{-1} \cdot 2 = 2 \cdot 2^{-1} \to \varphi_{9}$\lbl{A2}}
\BIC{$2^{-1} \cdot 2 = 2 \cdot 2^{-1} \to \varphi_{9}$}
\AXC{$\stackrel{\Pi_{9}}{\vdots}$}\noLine
\UIC{$\varphi_{63}$}
\BIC{$e \in \mathbb{R} \to e \cdot (2^{-1} \cdot 2) = e \cdot (2 \cdot 2^{-1})$}
\end{ptree}
\noindent$\Pi_{11}$:
\begin{ptree}
\AXC{$y = z \to x = y \to x = z$\lbl{A6}}
\AXC{$(y = z \to x = y \to x = z) \to x = y \to y = z \to x = z$\lbl{A1}}
\BIC{$x = y \to y = z \to x = z$}
\UIC{$\forall z\,(x = y \to y = z \to x = z)$}
\UIC{$\forall y\,\forall z\,(x = y \to y = z \to x = z)$}
\UIC{$\forall x\,\forall y\,\forall z\,(x = y \to y = z \to x = z)$}
\AXC{$\forall x\,\forall y\,\forall z\,(x = y \to y = z \to x = z) \to \varphi_{10}$\lbl{A2}}
\BIC{$\varphi_{10}$}
\end{ptree}
\noindent$\Pi_{12}$:
\begin{ptree}
\AXC{$e \in \mathbb{R} \to e \in \mathbb{R}$\lbl{A1}}
\AXC{$\varphi_{11}$\lbl{(R6)}}
\AXC{$\varphi_{11} \to \varphi_{8}$\lbl{A2}}
\BIC{$\varphi_{8}$}
\AXC{$\varphi_{8} \to \varphi_{31}$\lbl{A2}}
\BIC{$\varphi_{31}$}
\AXC{$\varphi_{31} \to \varphi_{20}$\lbl{A2}}
\BIC{$\varphi_{20}$}
\AXC{$\varphi_{20} \to (e \in \mathbb{R} \to e \in \mathbb{R}) \to \varphi_{21}$\lbl{A1}}
\BIC{$(e \in \mathbb{R} \to e \in \mathbb{R}) \to \varphi_{21}$}
\BIC{$\varphi_{21}$}
\end{ptree}
\noindent$\Pi_{13}$:
\begin{ptree}
\AXC{$y = z \to x = y \to x = z$\lbl{A6}}
\AXC{$(y = z \to x = y \to x = z) \to x = y \to y = z \to x = z$\lbl{A1}}
\BIC{$x = y \to y = z \to x = z$}
\UIC{$\forall z\,(x = y \to y = z \to x = z)$}
\UIC{$\forall y\,\forall z\,(x = y \to y = z \to x = z)$}
\UIC{$\forall x\,\forall y\,\forall z\,(x = y \to y = z \to x = z)$}
\AXC{$\forall x\,\forall y\,\forall z\,(x = y \to y = z \to x = z) \to \varphi_{10}$\lbl{A2}}
\BIC{$\varphi_{10}$}
\end{ptree}
\noindent$\Pi_{14}$:
\begin{ptree}
\AXC{$x = x$\lbl{A5}}
\AXC{$x = y \to x = x \to y = x$\lbl{A6}}
\AXC{$(x = y \to x = x \to y = x) \to x = x \to x = y \to y = x$\lbl{A1}}
\BIC{$x = x \to x = y \to y = x$}
\BIC{$x = y \to y = x$}
\UIC{$\forall y\,(x = y \to y = x)$}
\UIC{$\forall x\,\forall y\,(x = y \to y = x)$}
\end{ptree}
\noindent$\Pi_{15}$:
\begin{ptree}
\AXC{$\stackrel{\Pi_{14}}{\vdots}$}\noLine
\UIC{$\forall x\,\forall y\,(x = y \to y = x)$}
\AXC{$\forall x\,\forall y\,(x = y \to y = x) \to \forall y\,(2 = y \to y = 2)$\lbl{A2}}
\BIC{$\forall y\,(2 = y \to y = 2)$}
\AXC{$\forall y\,(2 = y \to y = 2) \to 2 = 1 + 1 \to 1 + 1 = 2$\lbl{A2}}
\BIC{$2 = 1 + 1 \to 1 + 1 = 2$}
\end{ptree}
\noindent$\Pi_{16}$:
\begin{ptree}
\AXC{$\stackrel{\Pi_{15}}{\vdots}$}\noLine
\UIC{$2 = 1 + 1 \to 1 + 1 = 2$}
\AXC{$2 = 1 + 1$\lbl{Def$2$}}
\AXC{$2 = 1 + 1 \to (2 = 1 + 1 \to 1 + 1 = 2) \to 1 + 1 = 2$\lbl{A1}}
\BIC{$(2 = 1 + 1 \to 1 + 1 = 2) \to 1 + 1 = 2$}
\BIC{$1 + 1 = 2$}
\AXC{$\varphi_{70}$\lbl{A1}}
\BIC{$e \cdot 2^{-1} \cdot (1 + 1) = e \cdot 2^{-1} \cdot (1 + 1) \to \varphi_{64}$}
\end{ptree}
\noindent$\Pi_{17}$:
\begin{ptree}
\AXC{$x = x$\lbl{A5}}
\UIC{$\forall x\,(x = x)$}
\AXC{$\varphi_{1}$\lbl{A2}}
\BIC{$e \cdot 2^{-1} \cdot (1 + 1) = e \cdot 2^{-1} \cdot (1 + 1)$}
\AXC{$\stackrel{\Pi_{16}}{\vdots}$}\noLine
\UIC{$e \cdot 2^{-1} \cdot (1 + 1) = e \cdot 2^{-1} \cdot (1 + 1) \to \varphi_{64}$}
\BIC{$\varphi_{64}$}
\end{ptree}
\noindent$\Pi_{18}$:
\begin{ptree}
\AXC{$z = w \to \varphi_{30}$\lbl{A6}}
\UIC{$\forall w\,(z = w \to \varphi_{30})$}
\UIC{$\forall z\,\forall w\,(z = w \to \varphi_{30})$}
\AXC{$\forall z\,\forall w\,(z = w \to \varphi_{30}) \to \varphi_{41}$\lbl{A2}}
\BIC{$\varphi_{41}$}
\AXC{$\varphi_{41} \to 1 + 1 = 2 \to \varphi_{35}$\lbl{A2}}
\BIC{$1 + 1 = 2 \to \varphi_{35}$}
\AXC{$\stackrel{\Pi_{17}}{\vdots}$}\noLine
\UIC{$\varphi_{64}$}
\BIC{$e \in \mathbb{R} \to e \cdot 2^{-1} \cdot (1 + 1) = e \cdot 2^{-1} \cdot 2$}
\end{ptree}
\noindent$\Pi_{19}$:
\begin{ptree}
\AXC{$y = z \to x = y \to x = z$\lbl{A6}}
\AXC{$(y = z \to x = y \to x = z) \to x = y \to y = z \to x = z$\lbl{A1}}
\BIC{$x = y \to y = z \to x = z$}
\UIC{$\forall z\,(x = y \to y = z \to x = z)$}
\UIC{$\forall y\,\forall z\,(x = y \to y = z \to x = z)$}
\UIC{$\forall x\,\forall y\,\forall z\,(x = y \to y = z \to x = z)$}
\AXC{$\forall x\,\forall y\,\forall z\,(x = y \to y = z \to x = z) \to \varphi_{10}$\lbl{A2}}
\BIC{$\varphi_{10}$}
\end{ptree}
\noindent$\Pi_{20}$:
\begin{ptree}
\AXC{$\forall x\,\forall y\,(x \in \mathbb{R} \to y \in \mathbb{R} \to x \cdot y \in \mathbb{R})$\lbl{(R1)}}
\AXC{$\varphi_{37}$\lbl{A2}}
\BIC{$\forall y\,(e \in \mathbb{R} \to y \in \mathbb{R} \to e \cdot y \in \mathbb{R})$}
\AXC{$\varphi_{42}$\lbl{A2}}
\BIC{$e \in \mathbb{R} \to 2^{-1} \in \mathbb{R} \to e \cdot 2^{-1} \in \mathbb{R}$}
\AXC{$(e \in \mathbb{R} \to 2^{-1} \in \mathbb{R} \to e \cdot 2^{-1} \in \mathbb{R}) \to \varphi_{22}$\lbl{A1}}
\BIC{$\varphi_{22}$}
\end{ptree}
\noindent$\Pi_{21}$:
\begin{ptree}
\AXC{$2^{-1} \in \mathbb{R}$\lbl{[F31]}}
\AXC{$e \in \mathbb{R} \to e \in \mathbb{R}$\lbl{A1}}
\AXC{$\stackrel{\Pi_{20}}{\vdots}$}\noLine
\UIC{$\varphi_{22}$}
\BIC{$2^{-1} \in \mathbb{R} \to e \in \mathbb{R} \to e \cdot 2^{-1} \in \mathbb{R}$}
\BIC{$e \in \mathbb{R} \to e \cdot 2^{-1} \in \mathbb{R}$}
\end{ptree}
\noindent$\Pi_{22}$:
\begin{ptree}
\AXC{$\varphi_{18}$\lbl{(R10)}}
\AXC{$\varphi_{18} \to \varphi_{51}$\lbl{A2}}
\BIC{$\varphi_{51}$}
\AXC{$\varphi_{51} \to \varphi_{50}$\lbl{A2}}
\BIC{$\varphi_{50}$}
\AXC{$\varphi_{50} \to e \cdot 2^{-1} \in \mathbb{R} \to 1 \in \mathbb{R} \to \varphi_{14}$\lbl{A2}}
\BIC{$e \cdot 2^{-1} \in \mathbb{R} \to 1 \in \mathbb{R} \to \varphi_{14}$}
\AXC{$\varphi_{74}$\lbl{A1}}
\BIC{$(e \in \mathbb{R} \to e \cdot 2^{-1} \in \mathbb{R}) \to 1 \in \mathbb{R} \to 1 \in \mathbb{R} \to \varphi_{12}$}
\end{ptree}
\noindent$\Pi_{23}$:
\begin{ptree}
\AXC{$1 \in \mathbb{R}$\lbl{(R1)}}
\AXC{$1 \in \mathbb{R}$\lbl{(R1)}}
\AXC{$\stackrel{\Pi_{21}}{\vdots}$}\noLine
\UIC{$e \in \mathbb{R} \to e \cdot 2^{-1} \in \mathbb{R}$}
\AXC{$\stackrel{\Pi_{22}}{\vdots}$}\noLine
\UIC{$(e \in \mathbb{R} \to e \cdot 2^{-1} \in \mathbb{R}) \to 1 \in \mathbb{R} \to 1 \in \mathbb{R} \to \varphi_{12}$}
\BIC{$1 \in \mathbb{R} \to 1 \in \mathbb{R} \to \varphi_{12}$}
\BIC{$1 \in \mathbb{R} \to \varphi_{12}$}
\BIC{$\varphi_{12}$}
\end{ptree}
\noindent$\Pi_{24}$:
\begin{ptree}
\AXC{$\stackrel{\Pi_{14}}{\vdots}$}\noLine
\UIC{$\forall x\,\forall y\,(x = y \to y = x)$}
\AXC{$\forall x\,\forall y\,(x = y \to y = x) \to \varphi_{2}$\lbl{A2}}
\BIC{$\varphi_{2}$}
\AXC{$\varphi_{2} \to \varphi_{55}$\lbl{A2}}
\BIC{$\varphi_{55}$}
\AXC{$\stackrel{\Pi_{23}}{\vdots}$}\noLine
\UIC{$\varphi_{12}$}
\AXC{$\varphi_{12} \to \varphi_{55} \to \varphi_{13}$\lbl{A1}}
\BIC{$\varphi_{55} \to \varphi_{13}$}
\BIC{$\varphi_{13}$}
\end{ptree}
\noindent$\Pi_{25}$:
\begin{ptree}
\AXC{$y = z \to x = y \to x = z$\lbl{A6}}
\AXC{$(y = z \to x = y \to x = z) \to x = y \to y = z \to x = z$\lbl{A1}}
\BIC{$x = y \to y = z \to x = z$}
\UIC{$\forall z\,(x = y \to y = z \to x = z)$}
\UIC{$\forall y\,\forall z\,(x = y \to y = z \to x = z)$}
\UIC{$\forall x\,\forall y\,\forall z\,(x = y \to y = z \to x = z)$}
\AXC{$\forall x\,\forall y\,\forall z\,(x = y \to y = z \to x = z) \to \varphi_{10}$\lbl{A2}}
\BIC{$\varphi_{10}$}
\end{ptree}
\noindent$\Pi_{26}$:
\begin{ptree}
\AXC{$z = w \to \varphi_{49}$\lbl{A6}}
\UIC{$\forall w\,(z = w \to \varphi_{49})$}
\UIC{$\forall z\,\forall w\,(z = w \to \varphi_{49})$}
\AXC{$\forall z\,\forall w\,(z = w \to \varphi_{49}) \to \varphi_{61}$\lbl{A2}}
\BIC{$\varphi_{61}$}
\AXC{$\varphi_{61} \to e \cdot 2^{-1} = e \cdot 2^{-1} \cdot 1 \to \varphi_{59}$\lbl{A2}}
\BIC{$e \cdot 2^{-1} = e \cdot 2^{-1} \cdot 1 \to \varphi_{59}$}
\end{ptree}
\noindent$\Pi_{27}$:
\begin{ptree}
\AXC{$\stackrel{\Pi_{21}}{\vdots}$}\noLine
\UIC{$e \in \mathbb{R} \to e \cdot 2^{-1} \in \mathbb{R}$}
\AXC{$\forall x \in \mathbb{R}\,(x \cdot 1 = x)$\lbl{(R8)}}
\AXC{$\varphi_{27}$\lbl{A2}}
\BIC{$e \cdot 2^{-1} \in \mathbb{R} \to e \cdot 2^{-1} \cdot 1 = e \cdot 2^{-1}$}
\AXC{$(e \cdot 2^{-1} \in \mathbb{R} \to e \cdot 2^{-1} \cdot 1 = e \cdot 2^{-1}) \to \varphi_{23}$\lbl{A1}}
\BIC{$\varphi_{23}$}
\BIC{$e \in \mathbb{R} \to e \cdot 2^{-1} \cdot 1 = e \cdot 2^{-1}$}
\end{ptree}
\noindent$\Pi_{28}$:
\begin{ptree}
\AXC{$\stackrel{\Pi_{27}}{\vdots}$}\noLine
\UIC{$e \in \mathbb{R} \to e \cdot 2^{-1} \cdot 1 = e \cdot 2^{-1}$}
\AXC{$(e \in \mathbb{R} \to e \cdot 2^{-1} \cdot 1 = e \cdot 2^{-1}) \to \varphi_{43}$\lbl{A1}}
\BIC{$\varphi_{43}$}
\end{ptree}
\noindent$\Pi_{29}$:
\begin{ptree}
\AXC{$\stackrel{\Pi_{14}}{\vdots}$}\noLine
\UIC{$\forall x\,\forall y\,(x = y \to y = x)$}
\AXC{$\varphi_{6}$\lbl{A2}}
\BIC{$\forall y\,(e \cdot 2^{-1} \cdot 1 = y \to y = e \cdot 2^{-1} \cdot 1)$}
\AXC{$\varphi_{44}$\lbl{A2}}
\BIC{$e \cdot 2^{-1} \cdot 1 = e \cdot 2^{-1} \to e \cdot 2^{-1} = e \cdot 2^{-1} \cdot 1$}
\AXC{$\stackrel{\Pi_{28}}{\vdots}$}\noLine
\UIC{$\varphi_{43}$}
\BIC{$e \in \mathbb{R} \to e \cdot 2^{-1} = e \cdot 2^{-1} \cdot 1$}
\AXC{$(e \in \mathbb{R} \to e \cdot 2^{-1} = e \cdot 2^{-1} \cdot 1) \to \varphi_{79}$\lbl{A1}}
\BIC{$\varphi_{79}$}
\end{ptree}
\noindent$\Pi_{30}$:
\begin{ptree}
\AXC{$\stackrel{\Pi_{21}}{\vdots}$}\noLine
\UIC{$e \in \mathbb{R} \to e \cdot 2^{-1} \in \mathbb{R}$}
\AXC{$\forall x \in \mathbb{R}\,(x \cdot 1 = x)$\lbl{(R8)}}
\AXC{$\varphi_{27}$\lbl{A2}}
\BIC{$e \cdot 2^{-1} \in \mathbb{R} \to e \cdot 2^{-1} \cdot 1 = e \cdot 2^{-1}$}
\AXC{$(e \cdot 2^{-1} \in \mathbb{R} \to e \cdot 2^{-1} \cdot 1 = e \cdot 2^{-1}) \to \varphi_{23}$\lbl{A1}}
\BIC{$\varphi_{23}$}
\BIC{$e \in \mathbb{R} \to e \cdot 2^{-1} \cdot 1 = e \cdot 2^{-1}$}
\end{ptree}
\noindent$\Pi_{31}$:
\begin{ptree}
\AXC{$\stackrel{\Pi_{30}}{\vdots}$}\noLine
\UIC{$e \in \mathbb{R} \to e \cdot 2^{-1} \cdot 1 = e \cdot 2^{-1}$}
\AXC{$(e \in \mathbb{R} \to e \cdot 2^{-1} \cdot 1 = e \cdot 2^{-1}) \to \varphi_{43}$\lbl{A1}}
\BIC{$\varphi_{43}$}
\end{ptree}
\noindent$\Pi_{32}$:
\begin{ptree}
\AXC{$x = x$\lbl{A5}}
\UIC{$\forall x\,(x = x)$}
\AXC{$\varphi_{7}$\lbl{A2}}
\BIC{$e \cdot 2^{-1} + e \cdot 2^{-1} = e \cdot 2^{-1} + e \cdot 2^{-1}$}
\AXC{$\stackrel{\Pi_{14}}{\vdots}$}\noLine
\UIC{$\forall x\,\forall y\,(x = y \to y = x)$}
\AXC{$\varphi_{6}$\lbl{A2}}
\BIC{$\forall y\,(e \cdot 2^{-1} \cdot 1 = y \to y = e \cdot 2^{-1} \cdot 1)$}
\AXC{$\varphi_{44}$\lbl{A2}}
\BIC{$e \cdot 2^{-1} \cdot 1 = e \cdot 2^{-1} \to e \cdot 2^{-1} = e \cdot 2^{-1} \cdot 1$}
\AXC{$\stackrel{\Pi_{31}}{\vdots}$}\noLine
\UIC{$\varphi_{43}$}
\BIC{$e \in \mathbb{R} \to e \cdot 2^{-1} = e \cdot 2^{-1} \cdot 1$}
\AXC{$\varphi_{81}$\lbl{A1}}
\BIC{$e \cdot 2^{-1} + e \cdot 2^{-1} = e \cdot 2^{-1} + e \cdot 2^{-1} \to \varphi_{72}$}
\BIC{$\varphi_{72}$}
\end{ptree}
\noindent$\Pi_{33}$:
\begin{ptree}
\AXC{$z = w \to \varphi_{36}$\lbl{A6}}
\UIC{$\forall w\,(z = w \to \varphi_{36})$}
\UIC{$\forall z\,\forall w\,(z = w \to \varphi_{36})$}
\AXC{$\forall z\,\forall w\,(z = w \to \varphi_{36}) \to \varphi_{54}$\lbl{A2}}
\BIC{$\varphi_{54}$}
\AXC{$\varphi_{54} \to e \cdot 2^{-1} = e \cdot 2^{-1} \cdot 1 \to \varphi_{53}$\lbl{A2}}
\BIC{$e \cdot 2^{-1} = e \cdot 2^{-1} \cdot 1 \to \varphi_{53}$}
\AXC{$\stackrel{\Pi_{32}}{\vdots}$}\noLine
\UIC{$\varphi_{72}$}
\BIC{$e \in \mathbb{R} \to e \cdot 2^{-1} + e \cdot 2^{-1} = e \cdot 2^{-1} \cdot 1 + e \cdot 2^{-1}$}
\AXC{$\varphi_{82}$\lbl{A1}}
\BIC{$\varphi_{32} \to \varphi_{75}$}
\end{ptree}
\noindent$\Pi_{34}$:
\begin{ptree}
\AXC{$\stackrel{\Pi_{26}}{\vdots}$}\noLine
\UIC{$e \cdot 2^{-1} = e \cdot 2^{-1} \cdot 1 \to \varphi_{59}$}
\AXC{$x = x$\lbl{A5}}
\UIC{$\forall x\,(x = x)$}
\AXC{$\varphi_{16}$\lbl{A2}}
\BIC{$e \cdot 2^{-1} \cdot 1 + e \cdot 2^{-1} = e \cdot 2^{-1} \cdot 1 + e \cdot 2^{-1}$}
\AXC{$\stackrel{\Pi_{29}}{\vdots}$}\noLine
\UIC{$\varphi_{79}$}
\BIC{$(e \cdot 2^{-1} = e \cdot 2^{-1} \cdot 1 \to \varphi_{59}) \to \varphi_{32}$}
\BIC{$\varphi_{32}$}
\AXC{$\stackrel{\Pi_{33}}{\vdots}$}\noLine
\UIC{$\varphi_{32} \to \varphi_{75}$}
\BIC{$\varphi_{75}$}
\end{ptree}
\noindent$\Pi_{35}$:
\begin{ptree}
\AXC{$\stackrel{\Pi_{25}}{\vdots}$}\noLine
\UIC{$\varphi_{10}$}
\AXC{$\varphi_{10} \to \varphi_{57}$\lbl{A2}}
\BIC{$\varphi_{57}$}
\AXC{$\varphi_{57} \to e \cdot 2^{-1} + e \cdot 2^{-1} = e \cdot 2^{-1} \cdot 1 + e \cdot 2^{-1} \to \varphi_{58}$\lbl{A2}}
\BIC{$e \cdot 2^{-1} + e \cdot 2^{-1} = e \cdot 2^{-1} \cdot 1 + e \cdot 2^{-1} \to \varphi_{58}$}
\AXC{$\stackrel{\Pi_{34}}{\vdots}$}\noLine
\UIC{$\varphi_{75}$}
\BIC{$\varphi_{19}$}
\AXC{$\varphi_{19} \to \varphi_{13} \to \varphi_{73}$\lbl{A1}}
\BIC{$\varphi_{13} \to \varphi_{73}$}
\end{ptree}
\noindent$\Pi_{36}$:
\begin{ptree}
\AXC{$\stackrel{\Pi_{19}}{\vdots}$}\noLine
\UIC{$\varphi_{10}$}
\AXC{$\varphi_{10} \to \varphi_{60}$\lbl{A2}}
\BIC{$\varphi_{60}$}
\AXC{$\varphi_{80}$\lbl{A2}}
\BIC{$e \cdot 2^{-1} + e \cdot 2^{-1} = e \cdot 2^{-1} \cdot 1 + e \cdot 2^{-1} \cdot 1 \to \varphi_{48}$}
\AXC{$\stackrel{\Pi_{24}}{\vdots}$}\noLine
\UIC{$\varphi_{13}$}
\AXC{$\stackrel{\Pi_{35}}{\vdots}$}\noLine
\UIC{$\varphi_{13} \to \varphi_{73}$}
\BIC{$\varphi_{73}$}
\BIC{$e \in \mathbb{R} \to e \cdot 2^{-1} + e \cdot 2^{-1} = e \cdot 2^{-1} \cdot (1 + 1)$}
\AXC{$\varphi_{78}$\lbl{A1}}
\BIC{$(e \in \mathbb{R} \to e \cdot 2^{-1} \cdot (1 + 1) = e \cdot 2^{-1} \cdot 2) \to \varphi_{69}$}
\end{ptree}
\noindent$\Pi_{37}$:
\begin{ptree}
\AXC{$\stackrel{\Pi_{18}}{\vdots}$}\noLine
\UIC{$e \in \mathbb{R} \to e \cdot 2^{-1} \cdot (1 + 1) = e \cdot 2^{-1} \cdot 2$}
\AXC{$\stackrel{\Pi_{36}}{\vdots}$}\noLine
\UIC{$(e \in \mathbb{R} \to e \cdot 2^{-1} \cdot (1 + 1) = e \cdot 2^{-1} \cdot 2) \to \varphi_{69}$}
\BIC{$\varphi_{69}$}
\end{ptree}
\noindent$\Pi_{38}$:
\begin{ptree}
\AXC{$\stackrel{\Pi_{13}}{\vdots}$}\noLine
\UIC{$\varphi_{10}$}
\AXC{$\varphi_{10} \to \varphi_{47}$\lbl{A2}}
\BIC{$\varphi_{47}$}
\AXC{$\varphi_{47} \to e \cdot 2^{-1} + e \cdot 2^{-1} = e \cdot 2^{-1} \cdot (1 + 1) \to \varphi_{34}$\lbl{A2}}
\BIC{$e \cdot 2^{-1} + e \cdot 2^{-1} = e \cdot 2^{-1} \cdot (1 + 1) \to \varphi_{34}$}
\AXC{$\stackrel{\Pi_{37}}{\vdots}$}\noLine
\UIC{$\varphi_{69}$}
\BIC{$e \in \mathbb{R} \to e \cdot 2^{-1} + e \cdot 2^{-1} = e \cdot 2^{-1} \cdot 2$}
\AXC{$\varphi_{76}$\lbl{A1}}
\BIC{$(e \in \mathbb{R} \to e \cdot 2^{-1} \cdot 2 = e \cdot (2^{-1} \cdot 2)) \to \varphi_{66}$}
\end{ptree}
\noindent$\Pi_{39}$:
\begin{ptree}
\AXC{$2 \in \mathbb{R}$\lbl{[F2]}}
\AXC{$2^{-1} \in \mathbb{R}$\lbl{[F31]}}
\AXC{$\stackrel{\Pi_{12}}{\vdots}$}\noLine
\UIC{$\varphi_{21}$}
\BIC{$2 \in \mathbb{R} \to e \in \mathbb{R} \to e \cdot 2^{-1} \cdot 2 = e \cdot (2^{-1} \cdot 2)$}
\BIC{$e \in \mathbb{R} \to e \cdot 2^{-1} \cdot 2 = e \cdot (2^{-1} \cdot 2)$}
\AXC{$\stackrel{\Pi_{38}}{\vdots}$}\noLine
\UIC{$(e \in \mathbb{R} \to e \cdot 2^{-1} \cdot 2 = e \cdot (2^{-1} \cdot 2)) \to \varphi_{66}$}
\BIC{$\varphi_{66}$}
\end{ptree}
\noindent$\Pi_{40}$:
\begin{ptree}
\AXC{$\stackrel{\Pi_{11}}{\vdots}$}\noLine
\UIC{$\varphi_{10}$}
\AXC{$\varphi_{10} \to \varphi_{40}$\lbl{A2}}
\BIC{$\varphi_{40}$}
\AXC{$\varphi_{40} \to e \cdot 2^{-1} + e \cdot 2^{-1} = e \cdot 2^{-1} \cdot 2 \to \varphi_{29}$\lbl{A2}}
\BIC{$e \cdot 2^{-1} + e \cdot 2^{-1} = e \cdot 2^{-1} \cdot 2 \to \varphi_{29}$}
\AXC{$\stackrel{\Pi_{39}}{\vdots}$}\noLine
\UIC{$\varphi_{66}$}
\BIC{$e \in \mathbb{R} \to e \cdot 2^{-1} + e \cdot 2^{-1} = e \cdot (2^{-1} \cdot 2)$}
\AXC{$\varphi_{77}$\lbl{A1}}
\BIC{$(e \in \mathbb{R} \to e \cdot (2^{-1} \cdot 2) = e \cdot (2 \cdot 2^{-1})) \to \varphi_{67}$}
\end{ptree}
\noindent$\Pi_{41}$:
\begin{ptree}
\AXC{$\stackrel{\Pi_{10}}{\vdots}$}\noLine
\UIC{$e \in \mathbb{R} \to e \cdot (2^{-1} \cdot 2) = e \cdot (2 \cdot 2^{-1})$}
\AXC{$\stackrel{\Pi_{40}}{\vdots}$}\noLine
\UIC{$(e \in \mathbb{R} \to e \cdot (2^{-1} \cdot 2) = e \cdot (2 \cdot 2^{-1})) \to \varphi_{67}$}
\BIC{$\varphi_{67}$}
\end{ptree}
\noindent$\Pi_{42}$:
\begin{ptree}
\AXC{$\stackrel{\Pi_{7}}{\vdots}$}\noLine
\UIC{$\varphi_{10}$}
\AXC{$\varphi_{10} \to \varphi_{39}$\lbl{A2}}
\BIC{$\varphi_{39}$}
\AXC{$\varphi_{39} \to e \cdot 2^{-1} + e \cdot 2^{-1} = e \cdot (2^{-1} \cdot 2) \to \varphi_{28}$\lbl{A2}}
\BIC{$e \cdot 2^{-1} + e \cdot 2^{-1} = e \cdot (2^{-1} \cdot 2) \to \varphi_{28}$}
\AXC{$\stackrel{\Pi_{41}}{\vdots}$}\noLine
\UIC{$\varphi_{67}$}
\BIC{$e \in \mathbb{R} \to e \cdot 2^{-1} + e \cdot 2^{-1} = e \cdot (2 \cdot 2^{-1})$}
\end{ptree}
\noindent$\Pi_{43}$:
\begin{ptree}
\AXC{$\stackrel{\Pi_{42}}{\vdots}$}\noLine
\UIC{$e \in \mathbb{R} \to e \cdot 2^{-1} + e \cdot 2^{-1} = e \cdot (2 \cdot 2^{-1})$}
\AXC{$(e \in \mathbb{R} \to e \cdot 2^{-1} + e \cdot 2^{-1} = e \cdot (2 \cdot 2^{-1})) \to \varphi_{68}$\lbl{A1}}
\BIC{$\varphi_{68}$}
\end{ptree}
\noindent$\Pi_{44}$:
\begin{ptree}
\AXC{$\stackrel{\Pi_{2}}{\vdots}$}\noLine
\UIC{$\varphi_{10}$}
\AXC{$\varphi_{10} \to \varphi_{38}$\lbl{A2}}
\BIC{$\varphi_{38}$}
\AXC{$\varphi_{38} \to \varphi_{45}$\lbl{A2}}
\BIC{$\varphi_{45}$}
\AXC{$\stackrel{\Pi_{6}}{\vdots}$}\noLine
\UIC{$e \in \mathbb{R} \to e \cdot (2 \cdot 2^{-1}) = e \cdot 1$}
\AXC{$\stackrel{\Pi_{43}}{\vdots}$}\noLine
\UIC{$\varphi_{68}$}
\BIC{$\varphi_{45} \to e \in \mathbb{R} \to e \cdot 2^{-1} + e \cdot 2^{-1} = e \cdot 1$}
\BIC{$e \in \mathbb{R} \to e \cdot 2^{-1} + e \cdot 2^{-1} = e \cdot 1$}
\end{ptree}
\noindent$\Pi_{45}$:
\begin{ptree}
\AXC{$\stackrel{\Pi_{44}}{\vdots}$}\noLine
\UIC{$e \in \mathbb{R} \to e \cdot 2^{-1} + e \cdot 2^{-1} = e \cdot 1$}
\AXC{$(e \in \mathbb{R} \to e \cdot 2^{-1} + e \cdot 2^{-1} = e \cdot 1) \to \varphi_{62}$\lbl{A1}}
\BIC{$\varphi_{62}$}
\end{ptree}
\noindent$\Pi_{46}$:
\begin{ptree}
\AXC{$\forall x \in \mathbb{R}\,(x \cdot 1 = x)$\lbl{(R8)}}
\AXC{$\forall x \in \mathbb{R}\,(x \cdot 1 = x) \to e \in \mathbb{R} \to e \cdot 1 = e$\lbl{A2}}
\BIC{$e \in \mathbb{R} \to e \cdot 1 = e$}
\AXC{$\stackrel{\Pi_{45}}{\vdots}$}\noLine
\UIC{$\varphi_{62}$}
\BIC{$\varphi_{17} \to e \in \mathbb{R} \to e \cdot 2^{-1} + e \cdot 2^{-1} = e$}
\end{ptree}
\noindent 主木:
\begin{ptree}
\AXC{$\stackrel{\Pi_{1}}{\vdots}$}\noLine
\UIC{$\varphi_{10}$}
\AXC{$\varphi_{10} \to \varphi_{24}$\lbl{A2}}
\BIC{$\varphi_{24}$}
\AXC{$\varphi_{24} \to \varphi_{17}$\lbl{A2}}
\BIC{$\varphi_{17}$}
\AXC{$\stackrel{\Pi_{46}}{\vdots}$}\noLine
\UIC{$\varphi_{17} \to e \in \mathbb{R} \to e \cdot 2^{-1} + e \cdot 2^{-1} = e$}
\BIC{$e \in \mathbb{R} \to e \cdot 2^{-1} + e \cdot 2^{-1} = e$}
\UIC{$\forall e \in \mathbb{R}\,(e \cdot 2^{-1} + e \cdot 2^{-1} = e)$}
\end{ptree}
\begin{abbrevlist}
\item $\varphi_{1} :\equiv \forall x\,(x = x) \to e \cdot 2^{-1} \cdot (1 + 1) = e \cdot 2^{-1} \cdot (1 + 1)$
\item $\varphi_{2} :\equiv \forall y\,(e \cdot 2^{-1} \cdot (1 + 1) = y \to y = e \cdot 2^{-1} \cdot (1 + 1))$
\item $\varphi_{3} :\equiv z = w \to e \cdot (2 \cdot 2^{-1}) = e \cdot z \to e \cdot (2 \cdot 2^{-1}) = e \cdot w$
\item $\varphi_{4} :\equiv z = w \to e \cdot (2^{-1} \cdot 2) = e \cdot z \to e \cdot (2^{-1} \cdot 2) = e \cdot w$
\item $\varphi_{5} :\equiv e \cdot (2 \cdot 2^{-1}) = e \cdot (2 \cdot 2^{-1}) \to e \cdot (2 \cdot 2^{-1}) = e \cdot 1$
\item $\varphi_{6} :\equiv \forall x\,\forall y\,(x = y \to y = x) \to \forall y\,(e \cdot 2^{-1} \cdot 1 = y \to y = e \cdot 2^{-1} \cdot 1)$
\item $\varphi_{7} :\equiv \forall x\,(x = x) \to e \cdot 2^{-1} + e \cdot 2^{-1} = e \cdot 2^{-1} + e \cdot 2^{-1}$
\item $\varphi_{8} :\equiv \forall y\,\forall z\,(e \in \mathbb{R} \to y \in \mathbb{R} \to z \in \mathbb{R} \to e \cdot y \cdot z = e \cdot (y \cdot z))$
\item $\varphi_{9} :\equiv e \cdot (2^{-1} \cdot 2) = e \cdot (2^{-1} \cdot 2) \to e \cdot (2^{-1} \cdot 2) = e \cdot (2 \cdot 2^{-1})$
\item $\varphi_{10} :\equiv \forall y\,\forall z\,(e \cdot 2^{-1} + e \cdot 2^{-1} = y \to y = z \to e \cdot 2^{-1} + e \cdot 2^{-1} = z)$
\item $\varphi_{11} :\equiv \forall x\,\forall y\,\forall z\,(x \in \mathbb{R} \to y \in \mathbb{R} \to z \in \mathbb{R} \to x \cdot y \cdot z = x \cdot (y \cdot z))$
\item $\varphi_{12} :\equiv e \in \mathbb{R} \to e \cdot 2^{-1} \cdot (1 + 1) = e \cdot 2^{-1} \cdot 1 + e \cdot 2^{-1} \cdot 1$
\item $\varphi_{13} :\equiv e \in \mathbb{R} \to e \cdot 2^{-1} \cdot 1 + e \cdot 2^{-1} \cdot 1 = e \cdot 2^{-1} \cdot (1 + 1)$
\item $\varphi_{14} :\equiv 1 \in \mathbb{R} \to e \cdot 2^{-1} \cdot (1 + 1) = e \cdot 2^{-1} \cdot 1 + e \cdot 2^{-1} \cdot 1$
\item $\varphi_{15} :\equiv \forall x \in \mathbb{R}\,(x \neq 0 \to x \cdot x^{-1} = 1) \to 2 \in \mathbb{R} \to 2 \neq 0 \to 2 \cdot 2^{-1} = 1$
\item $\varphi_{16} :\equiv \forall x\,(x = x) \to e \cdot 2^{-1} \cdot 1 + e \cdot 2^{-1} = e \cdot 2^{-1} \cdot 1 + e \cdot 2^{-1}$
\item $\varphi_{17} :\equiv e \cdot 2^{-1} + e \cdot 2^{-1} = e \cdot 1 \to e \cdot 1 = e \to e \cdot 2^{-1} + e \cdot 2^{-1} = e$
\item $\varphi_{18} :\equiv \forall x\,\forall y\,\forall z\,(x \in \mathbb{R} \to y \in \mathbb{R} \to z \in \mathbb{R} \to x \cdot (y + z) = x \cdot y + x \cdot z)$
\item $\varphi_{19} :\equiv e \in \mathbb{R} \to e \cdot 2^{-1} + e \cdot 2^{-1} = e \cdot 2^{-1} \cdot 1 + e \cdot 2^{-1} \cdot 1$
\item $\varphi_{20} :\equiv e \in \mathbb{R} \to 2^{-1} \in \mathbb{R} \to 2 \in \mathbb{R} \to e \cdot 2^{-1} \cdot 2 = e \cdot (2^{-1} \cdot 2)$
\item $\varphi_{21} :\equiv 2^{-1} \in \mathbb{R} \to 2 \in \mathbb{R} \to e \in \mathbb{R} \to e \cdot 2^{-1} \cdot 2 = e \cdot (2^{-1} \cdot 2)$
\item $\varphi_{22} :\equiv (e \in \mathbb{R} \to e \in \mathbb{R}) \to 2^{-1} \in \mathbb{R} \to e \in \mathbb{R} \to e \cdot 2^{-1} \in \mathbb{R}$
\item $\varphi_{23} :\equiv (e \in \mathbb{R} \to e \cdot 2^{-1} \in \mathbb{R}) \to e \in \mathbb{R} \to e \cdot 2^{-1} \cdot 1 = e \cdot 2^{-1}$
\item $\varphi_{24} :\equiv \forall z\,(e \cdot 2^{-1} + e \cdot 2^{-1} = e \cdot 1 \to e \cdot 1 = z \to e \cdot 2^{-1} + e \cdot 2^{-1} = z)$
\item $\varphi_{25} :\equiv \forall w\,(2 \cdot 2^{-1} = w \to e \cdot (2 \cdot 2^{-1}) = e \cdot (2 \cdot 2^{-1}) \to e \cdot (2 \cdot 2^{-1}) = e \cdot w)$
\item $\varphi_{26} :\equiv \forall w\,(2^{-1} \cdot 2 = w \to e \cdot (2^{-1} \cdot 2) = e \cdot (2^{-1} \cdot 2) \to e \cdot (2^{-1} \cdot 2) = e \cdot w)$
\item $\varphi_{27} :\equiv \forall x \in \mathbb{R}\,(x \cdot 1 = x) \to e \cdot 2^{-1} \in \mathbb{R} \to e \cdot 2^{-1} \cdot 1 = e \cdot 2^{-1}$
\item $\varphi_{28} :\equiv e \cdot (2^{-1} \cdot 2) = e \cdot (2 \cdot 2^{-1}) \to e \cdot 2^{-1} + e \cdot 2^{-1} = e \cdot (2 \cdot 2^{-1})$
\item $\varphi_{29} :\equiv e \cdot 2^{-1} \cdot 2 = e \cdot (2^{-1} \cdot 2) \to e \cdot 2^{-1} + e \cdot 2^{-1} = e \cdot (2^{-1} \cdot 2)$
\item $\varphi_{30} :\equiv e \cdot 2^{-1} \cdot (1 + 1) = e \cdot 2^{-1} \cdot z \to e \cdot 2^{-1} \cdot (1 + 1) = e \cdot 2^{-1} \cdot w$
\item $\varphi_{31} :\equiv \forall z\,(e \in \mathbb{R} \to 2^{-1} \in \mathbb{R} \to z \in \mathbb{R} \to e \cdot 2^{-1} \cdot z = e \cdot (2^{-1} \cdot z))$
\item $\varphi_{32} :\equiv e \in \mathbb{R} \to e \cdot 2^{-1} \cdot 1 + e \cdot 2^{-1} = e \cdot 2^{-1} \cdot 1 + e \cdot 2^{-1} \cdot 1$
\item $\varphi_{33} :\equiv (2 \in \mathbb{R} \to 2 \neq 0 \to 2 \cdot 2^{-1} = 1) \to 2 \in \mathbb{R} \to 2 \neq 0 \to e \in \mathbb{R} \to 2 \cdot 2^{-1} = 1$
\item $\varphi_{34} :\equiv e \cdot 2^{-1} \cdot (1 + 1) = e \cdot 2^{-1} \cdot 2 \to e \cdot 2^{-1} + e \cdot 2^{-1} = e \cdot 2^{-1} \cdot 2$
\item $\varphi_{35} :\equiv e \cdot 2^{-1} \cdot (1 + 1) = e \cdot 2^{-1} \cdot (1 + 1) \to e \cdot 2^{-1} \cdot (1 + 1) = e \cdot 2^{-1} \cdot 2$
\item $\varphi_{36} :\equiv e \cdot 2^{-1} + e \cdot 2^{-1} = z + e \cdot 2^{-1} \to e \cdot 2^{-1} + e \cdot 2^{-1} = w + e \cdot 2^{-1}$
\item $\varphi_{37} :\equiv \forall x\,\forall y\,(x \in \mathbb{R} \to y \in \mathbb{R} \to x \cdot y \in \mathbb{R}) \to \forall y\,(e \in \mathbb{R} \to y \in \mathbb{R} \to e \cdot y \in \mathbb{R})$
\item $\varphi_{38} :\equiv \forall z\,(e \cdot 2^{-1} + e \cdot 2^{-1} = e \cdot (2 \cdot 2^{-1}) \to e \cdot (2 \cdot 2^{-1}) = z \to e \cdot 2^{-1} + e \cdot 2^{-1} = z)$
\item $\varphi_{39} :\equiv \forall z\,(e \cdot 2^{-1} + e \cdot 2^{-1} = e \cdot (2^{-1} \cdot 2) \to e \cdot (2^{-1} \cdot 2) = z \to e \cdot 2^{-1} + e \cdot 2^{-1} = z)$
\item $\varphi_{40} :\equiv \forall z\,(e \cdot 2^{-1} + e \cdot 2^{-1} = e \cdot 2^{-1} \cdot 2 \to e \cdot 2^{-1} \cdot 2 = z \to e \cdot 2^{-1} + e \cdot 2^{-1} = z)$
\item $\varphi_{41} :\equiv \forall w\,(1 + 1 = w \to e \cdot 2^{-1} \cdot (1 + 1) = e \cdot 2^{-1} \cdot (1 + 1) \to e \cdot 2^{-1} \cdot (1 + 1) = e \cdot 2^{-1} \cdot w)$
\item $\varphi_{42} :\equiv \forall y\,(e \in \mathbb{R} \to y \in \mathbb{R} \to e \cdot y \in \mathbb{R}) \to e \in \mathbb{R} \to 2^{-1} \in \mathbb{R} \to e \cdot 2^{-1} \in \mathbb{R}$
\item $\varphi_{43} :\equiv (e \cdot 2^{-1} \cdot 1 = e \cdot 2^{-1} \to e \cdot 2^{-1} = e \cdot 2^{-1} \cdot 1) \to e \in \mathbb{R} \to e \cdot 2^{-1} = e \cdot 2^{-1} \cdot 1$
\item $\varphi_{44} :\equiv \forall y\,(e \cdot 2^{-1} \cdot 1 = y \to y = e \cdot 2^{-1} \cdot 1) \to e \cdot 2^{-1} \cdot 1 = e \cdot 2^{-1} \to e \cdot 2^{-1} = e \cdot 2^{-1} \cdot 1$
\item $\varphi_{45} :\equiv e \cdot 2^{-1} + e \cdot 2^{-1} = e \cdot (2 \cdot 2^{-1}) \to e \cdot (2 \cdot 2^{-1}) = e \cdot 1 \to e \cdot 2^{-1} + e \cdot 2^{-1} = e \cdot 1$
\item $\varphi_{46} :\equiv \forall x\,\forall y\,(x \in \mathbb{R} \to y \in \mathbb{R} \to x \cdot y = y \cdot x) \to \forall y\,(2^{-1} \in \mathbb{R} \to y \in \mathbb{R} \to 2^{-1} \cdot y = y \cdot 2^{-1})$
\item $\varphi_{47} :\equiv \forall z\,(e \cdot 2^{-1} + e \cdot 2^{-1} = e \cdot 2^{-1} \cdot (1 + 1) \to e \cdot 2^{-1} \cdot (1 + 1) = z \to e \cdot 2^{-1} + e \cdot 2^{-1} = z)$
\item $\varphi_{48} :\equiv e \cdot 2^{-1} \cdot 1 + e \cdot 2^{-1} \cdot 1 = e \cdot 2^{-1} \cdot (1 + 1) \to e \cdot 2^{-1} + e \cdot 2^{-1} = e \cdot 2^{-1} \cdot (1 + 1)$
\item $\varphi_{49} :\equiv e \cdot 2^{-1} \cdot 1 + e \cdot 2^{-1} = e \cdot 2^{-1} \cdot 1 + z \to e \cdot 2^{-1} \cdot 1 + e \cdot 2^{-1} = e \cdot 2^{-1} \cdot 1 + w$
\item $\varphi_{50} :\equiv \forall z\,(e \cdot 2^{-1} \in \mathbb{R} \to 1 \in \mathbb{R} \to z \in \mathbb{R} \to e \cdot 2^{-1} \cdot (1 + z) = e \cdot 2^{-1} \cdot 1 + e \cdot 2^{-1} \cdot z)$
\item $\varphi_{51} :\equiv \forall y\,\forall z\,(e \cdot 2^{-1} \in \mathbb{R} \to y \in \mathbb{R} \to z \in \mathbb{R} \to e \cdot 2^{-1} \cdot (y + z) = e \cdot 2^{-1} \cdot y + e \cdot 2^{-1} \cdot z)$
\item $\varphi_{52} :\equiv \forall y\,(2^{-1} \in \mathbb{R} \to y \in \mathbb{R} \to 2^{-1} \cdot y = y \cdot 2^{-1}) \to 2^{-1} \in \mathbb{R} \to 2 \in \mathbb{R} \to 2^{-1} \cdot 2 = 2 \cdot 2^{-1}$
\item $\varphi_{53} :\equiv e \cdot 2^{-1} + e \cdot 2^{-1} = e \cdot 2^{-1} + e \cdot 2^{-1} \to e \cdot 2^{-1} + e \cdot 2^{-1} = e \cdot 2^{-1} \cdot 1 + e \cdot 2^{-1}$
\item $\varphi_{54} :\equiv \forall w\,(e \cdot 2^{-1} = w \to e \cdot 2^{-1} + e \cdot 2^{-1} = e \cdot 2^{-1} + e \cdot 2^{-1} \to e \cdot 2^{-1} + e \cdot 2^{-1} = w + e \cdot 2^{-1})$
\item $\varphi_{55} :\equiv e \cdot 2^{-1} \cdot (1 + 1) = e \cdot 2^{-1} \cdot 1 + e \cdot 2^{-1} \cdot 1 \to e \cdot 2^{-1} \cdot 1 + e \cdot 2^{-1} \cdot 1 = e \cdot 2^{-1} \cdot (1 + 1)$
\item $\varphi_{56} :\equiv (2^{-1} \in \mathbb{R} \to 2 \in \mathbb{R} \to 2^{-1} \cdot 2 = 2 \cdot 2^{-1}) \to 2^{-1} \in \mathbb{R} \to 2 \in \mathbb{R} \to e \in \mathbb{R} \to 2^{-1} \cdot 2 = 2 \cdot 2^{-1}$
\item $\varphi_{57} :\equiv \forall z\,(e \cdot 2^{-1} + e \cdot 2^{-1} = e \cdot 2^{-1} \cdot 1 + e \cdot 2^{-1} \to e \cdot 2^{-1} \cdot 1 + e \cdot 2^{-1} = z \to e \cdot 2^{-1} + e \cdot 2^{-1} = z)$
\item $\varphi_{58} :\equiv e \cdot 2^{-1} \cdot 1 + e \cdot 2^{-1} = e \cdot 2^{-1} \cdot 1 + e \cdot 2^{-1} \cdot 1 \to e \cdot 2^{-1} + e \cdot 2^{-1} = e \cdot 2^{-1} \cdot 1 + e \cdot 2^{-1} \cdot 1$
\item $\varphi_{59} :\equiv e \cdot 2^{-1} \cdot 1 + e \cdot 2^{-1} = e \cdot 2^{-1} \cdot 1 + e \cdot 2^{-1} \to e \cdot 2^{-1} \cdot 1 + e \cdot 2^{-1} = e \cdot 2^{-1} \cdot 1 + e \cdot 2^{-1} \cdot 1$
\item $\varphi_{60} :\equiv \forall z\,(e \cdot 2^{-1} + e \cdot 2^{-1} = e \cdot 2^{-1} \cdot 1 + e \cdot 2^{-1} \cdot 1 \to e \cdot 2^{-1} \cdot 1 + e \cdot 2^{-1} \cdot 1 = z \to e \cdot 2^{-1} + e \cdot 2^{-1} = z)$
\item $\varphi_{61} :\equiv \forall w\,(e \cdot 2^{-1} = w \to e \cdot 2^{-1} \cdot 1 + e \cdot 2^{-1} = e \cdot 2^{-1} \cdot 1 + e \cdot 2^{-1} \to e \cdot 2^{-1} \cdot 1 + e \cdot 2^{-1} = e \cdot 2^{-1} \cdot 1 + w)$
\item $\varphi_{62} :\equiv (e \in \mathbb{R} \to e \cdot 1 = e) \to \varphi_{17} \to e \in \mathbb{R} \to e \cdot 2^{-1} + e \cdot 2^{-1} = e$
\item $\varphi_{63} :\equiv (2^{-1} \cdot 2 = 2 \cdot 2^{-1} \to \varphi_{9}) \to e \in \mathbb{R} \to e \cdot (2^{-1} \cdot 2) = e \cdot (2 \cdot 2^{-1})$
\item $\varphi_{64} :\equiv (1 + 1 = 2 \to \varphi_{35}) \to e \in \mathbb{R} \to e \cdot 2^{-1} \cdot (1 + 1) = e \cdot 2^{-1} \cdot 2$
\item $\varphi_{65} :\equiv e \cdot (2 \cdot 2^{-1}) = e \cdot (2 \cdot 2^{-1}) \to (2 \cdot 2^{-1} = 1 \to \varphi_{5}) \to e \in \mathbb{R} \to e \cdot (2 \cdot 2^{-1}) = e \cdot 1$
\item $\varphi_{66} :\equiv (e \cdot 2^{-1} + e \cdot 2^{-1} = e \cdot 2^{-1} \cdot 2 \to \varphi_{29}) \to e \in \mathbb{R} \to e \cdot 2^{-1} + e \cdot 2^{-1} = e \cdot (2^{-1} \cdot 2)$
\item $\varphi_{67} :\equiv (e \cdot 2^{-1} + e \cdot 2^{-1} = e \cdot (2^{-1} \cdot 2) \to \varphi_{28}) \to e \in \mathbb{R} \to e \cdot 2^{-1} + e \cdot 2^{-1} = e \cdot (2 \cdot 2^{-1})$
\item $\varphi_{68} :\equiv (e \in \mathbb{R} \to e \cdot (2 \cdot 2^{-1}) = e \cdot 1) \to \varphi_{45} \to e \in \mathbb{R} \to e \cdot 2^{-1} + e \cdot 2^{-1} = e \cdot 1$
\item $\varphi_{69} :\equiv (e \cdot 2^{-1} + e \cdot 2^{-1} = e \cdot 2^{-1} \cdot (1 + 1) \to \varphi_{34}) \to e \in \mathbb{R} \to e \cdot 2^{-1} + e \cdot 2^{-1} = e \cdot 2^{-1} \cdot 2$
\item $\varphi_{70} :\equiv 1 + 1 = 2 \to e \cdot 2^{-1} \cdot (1 + 1) = e \cdot 2^{-1} \cdot (1 + 1) \to \varphi_{64}$
\item $\varphi_{71} :\equiv (e \in \mathbb{R} \to 2^{-1} \cdot 2 = 2 \cdot 2^{-1}) \to e \cdot (2^{-1} \cdot 2) = e \cdot (2^{-1} \cdot 2) \to \varphi_{63}$
\item $\varphi_{72} :\equiv (e \cdot 2^{-1} = e \cdot 2^{-1} \cdot 1 \to \varphi_{53}) \to e \in \mathbb{R} \to e \cdot 2^{-1} + e \cdot 2^{-1} = e \cdot 2^{-1} \cdot 1 + e \cdot 2^{-1}$
\item $\varphi_{73} :\equiv (e \cdot 2^{-1} + e \cdot 2^{-1} = e \cdot 2^{-1} \cdot 1 + e \cdot 2^{-1} \cdot 1 \to \varphi_{48}) \to e \in \mathbb{R} \to e \cdot 2^{-1} + e \cdot 2^{-1} = e \cdot 2^{-1} \cdot (1 + 1)$
\item $\varphi_{74} :\equiv (e \cdot 2^{-1} \in \mathbb{R} \to 1 \in \mathbb{R} \to \varphi_{14}) \to (e \in \mathbb{R} \to e \cdot 2^{-1} \in \mathbb{R}) \to 1 \in \mathbb{R} \to 1 \in \mathbb{R} \to \varphi_{12}$
\item $\varphi_{75} :\equiv (e \cdot 2^{-1} + e \cdot 2^{-1} = e \cdot 2^{-1} \cdot 1 + e \cdot 2^{-1} \to \varphi_{58}) \to \varphi_{19}$
\item $\varphi_{76} :\equiv (e \in \mathbb{R} \to e \cdot 2^{-1} + e \cdot 2^{-1} = e \cdot 2^{-1} \cdot 2) \to (e \in \mathbb{R} \to e \cdot 2^{-1} \cdot 2 = e \cdot (2^{-1} \cdot 2)) \to \varphi_{66}$
\item $\varphi_{77} :\equiv (e \in \mathbb{R} \to e \cdot 2^{-1} + e \cdot 2^{-1} = e \cdot (2^{-1} \cdot 2)) \to (e \in \mathbb{R} \to e \cdot (2^{-1} \cdot 2) = e \cdot (2 \cdot 2^{-1})) \to \varphi_{67}$
\item $\varphi_{78} :\equiv (e \in \mathbb{R} \to e \cdot 2^{-1} + e \cdot 2^{-1} = e \cdot 2^{-1} \cdot (1 + 1)) \to (e \in \mathbb{R} \to e \cdot 2^{-1} \cdot (1 + 1) = e \cdot 2^{-1} \cdot 2) \to \varphi_{69}$
\item $\varphi_{79} :\equiv e \cdot 2^{-1} \cdot 1 + e \cdot 2^{-1} = e \cdot 2^{-1} \cdot 1 + e \cdot 2^{-1} \to (e \cdot 2^{-1} = e \cdot 2^{-1} \cdot 1 \to \varphi_{59}) \to \varphi_{32}$
\item $\varphi_{80} :\equiv \varphi_{60} \to e \cdot 2^{-1} + e \cdot 2^{-1} = e \cdot 2^{-1} \cdot 1 + e \cdot 2^{-1} \cdot 1 \to \varphi_{48}$
\item $\varphi_{81} :\equiv (e \in \mathbb{R} \to e \cdot 2^{-1} = e \cdot 2^{-1} \cdot 1) \to e \cdot 2^{-1} + e \cdot 2^{-1} = e \cdot 2^{-1} + e \cdot 2^{-1} \to \varphi_{72}$
\item $\varphi_{82} :\equiv (e \in \mathbb{R} \to e \cdot 2^{-1} + e \cdot 2^{-1} = e \cdot 2^{-1} \cdot 1 + e \cdot 2^{-1}) \to \varphi_{32} \to \varphi_{75}$
\end{abbrevlist}
\end{fullproof}

\section{順序体の補題(追加): 逆数, 不等式の操作}
% 自動生成 (prover/build.js): LeAddNonneg  ノード数 107
\begin{fullproof}{[E1]}{$\forall x\,\forall y\,(x \in \mathbb{R} \to y \in \mathbb{R} \to 0 \leq y \to x \leq x + y)$}{107}\phantomsection\label{tree:LeAddNonneg}
\noindent$\Pi_{1}$:
\begin{ptree}
\AXC{$0 \in \mathbb{R}$\lbl{(R1)}}
\AXC{$\varphi_{7}$\lbl{[F19]}}
\AXC{$\varphi_{7} \to \varphi_{5}$\lbl{A2}}
\BIC{$\varphi_{5}$}
\AXC{$\varphi_{29}$\lbl{A2}}
\BIC{$\forall z\,(0 \in \mathbb{R} \to y \in \mathbb{R} \to z \in \mathbb{R} \to 0 \leq y \to 0 + z \leq y + z)$}
\AXC{$\varphi_{27}$\lbl{A2}}
\BIC{$0 \in \mathbb{R} \to y \in \mathbb{R} \to x \in \mathbb{R} \to 0 \leq y \to 0 + x \leq y + x$}
\AXC{$\varphi_{35}$\lbl{A1}}
\BIC{$0 \in \mathbb{R} \to \varphi_{33}$}
\BIC{$\varphi_{33}$}
\end{ptree}
\noindent$\Pi_{2}$:
\begin{ptree}
\AXC{$x \in \mathbb{R} \to y \in \mathbb{R} \to 0 \leq y \to x \in \mathbb{R}$\lbl{A1}}
\AXC{$x \in \mathbb{R} \to y \in \mathbb{R} \to 0 \leq y \to y \in \mathbb{R}$\lbl{A1}}
\AXC{$\stackrel{\Pi_{1}}{\vdots}$}\noLine
\UIC{$\varphi_{33}$}
\BIC{$(x \in \mathbb{R} \to y \in \mathbb{R} \to 0 \leq y \to x \in \mathbb{R}) \to \varphi_{16}$}
\BIC{$\varphi_{16}$}
\end{ptree}
\noindent$\Pi_{3}$:
\begin{ptree}
\AXC{$y = z \to x = y \to x = z$\lbl{A6}}
\AXC{$(y = z \to x = y \to x = z) \to x = y \to y = z \to x = z$\lbl{A1}}
\BIC{$x = y \to y = z \to x = z$}
\UIC{$\forall z\,(x = y \to y = z \to x = z)$}
\UIC{$\forall y\,\forall z\,(x = y \to y = z \to x = z)$}
\UIC{$\forall x\,\forall y\,\forall z\,(x = y \to y = z \to x = z)$}
\AXC{$\varphi_{2}$\lbl{A2}}
\BIC{$\forall y\,\forall z\,(0 + x = y \to y = z \to 0 + x = z)$}
\AXC{$\varphi_{9}$\lbl{A2}}
\BIC{$\forall z\,(0 + x = x + 0 \to x + 0 = z \to 0 + x = z)$}
\AXC{$\varphi_{11}$\lbl{A2}}
\BIC{$0 + x = x + 0 \to x + 0 = x \to 0 + x = x$}
\end{ptree}
\noindent$\Pi_{4}$:
\begin{ptree}
\AXC{$\forall x \in \mathbb{R}\,(x + 0 = x)$\lbl{(R4)}}
\AXC{$\forall x \in \mathbb{R}\,(x + 0 = x) \to x \in \mathbb{R} \to x + 0 = x$\lbl{A2}}
\BIC{$x \in \mathbb{R} \to x + 0 = x$}
\AXC{$(x \in \mathbb{R} \to x + 0 = x) \to \varphi_{19}$\lbl{A1}}
\BIC{$\varphi_{19}$}
\end{ptree}
\noindent$\Pi_{5}$:
\begin{ptree}
\AXC{$\forall x\,\forall y\,(x \in \mathbb{R} \to y \in \mathbb{R} \to x + y = y + x)$\lbl{(R3)}}
\AXC{$\varphi_{21}$\lbl{A2}}
\BIC{$\forall y\,(0 \in \mathbb{R} \to y \in \mathbb{R} \to 0 + y = y + 0)$}
\AXC{$\varphi_{18}$\lbl{A2}}
\BIC{$0 \in \mathbb{R} \to x \in \mathbb{R} \to 0 + x = x + 0$}
\AXC{$(0 \in \mathbb{R} \to x \in \mathbb{R} \to 0 + x = x + 0) \to 0 \in \mathbb{R} \to \varphi_{22}$\lbl{A1}}
\BIC{$0 \in \mathbb{R} \to \varphi_{22}$}
\end{ptree}
\noindent$\Pi_{6}$:
\begin{ptree}
\AXC{$x \in \mathbb{R} \to y \in \mathbb{R} \to 0 \leq y \to x \in \mathbb{R}$\lbl{A1}}
\AXC{$0 \in \mathbb{R}$\lbl{(R1)}}
\AXC{$\stackrel{\Pi_{5}}{\vdots}$}\noLine
\UIC{$0 \in \mathbb{R} \to \varphi_{22}$}
\BIC{$\varphi_{22}$}
\BIC{$x \in \mathbb{R} \to y \in \mathbb{R} \to 0 \leq y \to 0 + x = x + 0$}
\AXC{$\varphi_{31}$\lbl{A1}}
\BIC{$(x \in \mathbb{R} \to y \in \mathbb{R} \to 0 \leq y \to x + 0 = x) \to \varphi_{13}$}
\end{ptree}
\noindent$\Pi_{7}$:
\begin{ptree}
\AXC{$x \in \mathbb{R} \to y \in \mathbb{R} \to 0 \leq y \to x \in \mathbb{R}$\lbl{A1}}
\AXC{$\stackrel{\Pi_{4}}{\vdots}$}\noLine
\UIC{$\varphi_{19}$}
\BIC{$x \in \mathbb{R} \to y \in \mathbb{R} \to 0 \leq y \to x + 0 = x$}
\AXC{$\stackrel{\Pi_{6}}{\vdots}$}\noLine
\UIC{$(x \in \mathbb{R} \to y \in \mathbb{R} \to 0 \leq y \to x + 0 = x) \to \varphi_{13}$}
\BIC{$\varphi_{13}$}
\end{ptree}
\noindent$\Pi_{8}$:
\begin{ptree}
\AXC{$\stackrel{\Pi_{3}}{\vdots}$}\noLine
\UIC{$0 + x = x + 0 \to x + 0 = x \to 0 + x = x$}
\AXC{$\stackrel{\Pi_{7}}{\vdots}$}\noLine
\UIC{$\varphi_{13}$}
\BIC{$x \in \mathbb{R} \to y \in \mathbb{R} \to 0 \leq y \to 0 + x = x$}
\AXC{$(x \in \mathbb{R} \to y \in \mathbb{R} \to 0 \leq y \to 0 + x = x) \to \varphi_{32}$\lbl{A1}}
\BIC{$\varphi_{32}$}
\end{ptree}
\noindent$\Pi_{9}$:
\begin{ptree}
\AXC{$z = w \to z < y + x \to w < y + x$\lbl{A6}}
\UIC{$\forall w\,(z = w \to z < y + x \to w < y + x)$}
\UIC{$\forall z\,\forall w\,(z = w \to z < y + x \to w < y + x)$}
\AXC{$\varphi_{17}$\lbl{A2}}
\BIC{$\forall w\,(0 + x = w \to 0 + x < y + x \to w < y + x)$}
\AXC{$\varphi_{20}$\lbl{A2}}
\BIC{$0 + x = x \to 0 + x < y + x \to x < y + x$}
\AXC{$x \in \mathbb{R} \to y \in \mathbb{R} \to 0 \leq y \to 0 \leq y$\lbl{A1}}
\AXC{$\stackrel{\Pi_{2}}{\vdots}$}\noLine
\UIC{$\varphi_{16}$}
\BIC{$x \in \mathbb{R} \to y \in \mathbb{R} \to 0 \leq y \to 0 + x \leq y + x$}
\AXC{$\stackrel{\Pi_{8}}{\vdots}$}\noLine
\UIC{$\varphi_{32}$}
\BIC{$(0 + x = x \to 0 + x < y + x \to x < y + x) \to \varphi_{14}$}
\BIC{$\varphi_{14}$}
\end{ptree}
\noindent$\Pi_{10}$:
\begin{ptree}
\AXC{$\forall x\,\forall y\,(x \in \mathbb{R} \to y \in \mathbb{R} \to x + y = y + x)$\lbl{(R3)}}
\AXC{$\varphi_{25}$\lbl{A2}}
\BIC{$\forall y\,(v_{422} \in \mathbb{R} \to y \in \mathbb{R} \to v_{422} + y = y + v_{422})$}
\AXC{$\varphi_{26}$\lbl{A2}}
\BIC{$v_{422} \in \mathbb{R} \to v_{423} \in \mathbb{R} \to v_{422} + v_{423} = v_{423} + v_{422}$}
\UIC{$\varphi_{3}$}
\UIC{$\forall v_{422}\,(\varphi_{3})$}
\AXC{$\varphi_{28}$\lbl{A2}}
\BIC{$\forall v_{423}\,(y \in \mathbb{R} \to v_{423} \in \mathbb{R} \to y + v_{423} = v_{423} + y)$}
\AXC{$\varphi_{24}$\lbl{A2}}
\BIC{$y \in \mathbb{R} \to x \in \mathbb{R} \to y + x = x + y$}
\AXC{$\varphi_{34}$\lbl{A1}}
\BIC{$(x \in \mathbb{R} \to y \in \mathbb{R} \to 0 \leq y \to y \in \mathbb{R}) \to \varphi_{23}$}
\end{ptree}
\noindent$\Pi_{11}$:
\begin{ptree}
\AXC{$x \in \mathbb{R} \to y \in \mathbb{R} \to 0 \leq y \to x \in \mathbb{R}$\lbl{A1}}
\AXC{$x \in \mathbb{R} \to y \in \mathbb{R} \to 0 \leq y \to y \in \mathbb{R}$\lbl{A1}}
\AXC{$\stackrel{\Pi_{10}}{\vdots}$}\noLine
\UIC{$(x \in \mathbb{R} \to y \in \mathbb{R} \to 0 \leq y \to y \in \mathbb{R}) \to \varphi_{23}$}
\BIC{$\varphi_{23}$}
\BIC{$x \in \mathbb{R} \to y \in \mathbb{R} \to 0 \leq y \to y + x = x + y$}
\end{ptree}
\noindent$\Pi_{12}$:
\begin{ptree}
\AXC{$\stackrel{\Pi_{11}}{\vdots}$}\noLine
\UIC{$x \in \mathbb{R} \to y \in \mathbb{R} \to 0 \leq y \to y + x = x + y$}
\AXC{$(x \in \mathbb{R} \to y \in \mathbb{R} \to 0 \leq y \to y + x = x + y) \to \varphi_{30}$\lbl{A1}}
\BIC{$\varphi_{30}$}
\end{ptree}
\noindent$\Pi_{13}$:
\begin{ptree}
\AXC{$z = w \to z = y + x \to w = y + x$\lbl{A6}}
\UIC{$\forall w\,(z = w \to z = y + x \to w = y + x)$}
\UIC{$\forall z\,\forall w\,(z = w \to z = y + x \to w = y + x)$}
\AXC{$\varphi_{8}$\lbl{A2}}
\BIC{$\forall w\,(0 + x = w \to 0 + x = y + x \to w = y + x)$}
\AXC{$\varphi_{10}$\lbl{A2}}
\BIC{$0 + x = x \to 0 + x = y + x \to x = y + x$}
\AXC{$\stackrel{\Pi_{9}}{\vdots}$}\noLine
\UIC{$\varphi_{14}$}
\BIC{$x \in \mathbb{R} \to y \in \mathbb{R} \to 0 \leq y \to x \leq y + x$}
\AXC{$\stackrel{\Pi_{12}}{\vdots}$}\noLine
\UIC{$\varphi_{30}$}
\BIC{$(y + x = x + y \to x < y + x \to x < x + y) \to \varphi_{15}$}
\end{ptree}
\noindent$\Pi_{14}$:
\begin{ptree}
\AXC{$z = w \to x < z \to x < w$\lbl{A6}}
\UIC{$\forall w\,(z = w \to x < z \to x < w)$}
\UIC{$\forall z\,\forall w\,(z = w \to x < z \to x < w)$}
\AXC{$\varphi_{4}$\lbl{A2}}
\BIC{$\forall w\,(y + x = w \to x < y + x \to x < w)$}
\AXC{$\varphi_{12}$\lbl{A2}}
\BIC{$y + x = x + y \to x < y + x \to x < x + y$}
\AXC{$\stackrel{\Pi_{13}}{\vdots}$}\noLine
\UIC{$(y + x = x + y \to x < y + x \to x < x + y) \to \varphi_{15}$}
\BIC{$\varphi_{15}$}
\end{ptree}
\noindent 主木:
\begin{ptree}
\AXC{$z = w \to x = z \to x = w$\lbl{A6}}
\UIC{$\forall w\,(z = w \to x = z \to x = w)$}
\UIC{$\forall z\,\forall w\,(z = w \to x = z \to x = w)$}
\AXC{$\varphi_{1}$\lbl{A2}}
\BIC{$\forall w\,(y + x = w \to x = y + x \to x = w)$}
\AXC{$\varphi_{6}$\lbl{A2}}
\BIC{$y + x = x + y \to x = y + x \to x = x + y$}
\AXC{$\stackrel{\Pi_{14}}{\vdots}$}\noLine
\UIC{$\varphi_{15}$}
\BIC{$x \in \mathbb{R} \to y \in \mathbb{R} \to 0 \leq y \to x \leq x + y$}
\UIC{$\forall y\,(x \in \mathbb{R} \to y \in \mathbb{R} \to 0 \leq y \to x \leq x + y)$}
\UIC{$\forall x\,\forall y\,(x \in \mathbb{R} \to y \in \mathbb{R} \to 0 \leq y \to x \leq x + y)$}
\end{ptree}
\begin{abbrevlist}
\item $\varphi_{1} :\equiv \forall z\,\forall w\,(z = w \to x = z \to x = w) \to \forall w\,(y + x = w \to x = y + x \to x = w)$
\item $\varphi_{2} :\equiv \forall x\,\forall y\,\forall z\,(x = y \to y = z \to x = z) \to \forall y\,\forall z\,(0 + x = y \to y = z \to 0 + x = z)$
\item $\varphi_{3} :\equiv \forall v_{423}\,(v_{422} \in \mathbb{R} \to v_{423} \in \mathbb{R} \to v_{422} + v_{423} = v_{423} + v_{422})$
\item $\varphi_{4} :\equiv \forall z\,\forall w\,(z = w \to x < z \to x < w) \to \forall w\,(y + x = w \to x < y + x \to x < w)$
\item $\varphi_{5} :\equiv \forall y\,\forall z\,(0 \in \mathbb{R} \to y \in \mathbb{R} \to z \in \mathbb{R} \to 0 \leq y \to 0 + z \leq y + z)$
\item $\varphi_{6} :\equiv \forall w\,(y + x = w \to x = y + x \to x = w) \to y + x = x + y \to x = y + x \to x = x + y$
\item $\varphi_{7} :\equiv \forall x\,\forall y\,\forall z\,(x \in \mathbb{R} \to y \in \mathbb{R} \to z \in \mathbb{R} \to x \leq y \to x + z \leq y + z)$
\item $\varphi_{8} :\equiv \forall z\,\forall w\,(z = w \to z = y + x \to w = y + x) \to \forall w\,(0 + x = w \to 0 + x = y + x \to w = y + x)$
\item $\varphi_{9} :\equiv \forall y\,\forall z\,(0 + x = y \to y = z \to 0 + x = z) \to \forall z\,(0 + x = x + 0 \to x + 0 = z \to 0 + x = z)$
\item $\varphi_{10} :\equiv \forall w\,(0 + x = w \to 0 + x = y + x \to w = y + x) \to 0 + x = x \to 0 + x = y + x \to x = y + x$
\item $\varphi_{11} :\equiv \forall z\,(0 + x = x + 0 \to x + 0 = z \to 0 + x = z) \to 0 + x = x + 0 \to x + 0 = x \to 0 + x = x$
\item $\varphi_{12} :\equiv \forall w\,(y + x = w \to x < y + x \to x < w) \to y + x = x + y \to x < y + x \to x < x + y$
\item $\varphi_{13} :\equiv (0 + x = x + 0 \to x + 0 = x \to 0 + x = x) \to x \in \mathbb{R} \to y \in \mathbb{R} \to 0 \leq y \to 0 + x = x$
\item $\varphi_{14} :\equiv (0 + x = x \to 0 + x = y + x \to x = y + x) \to x \in \mathbb{R} \to y \in \mathbb{R} \to 0 \leq y \to x \leq y + x$
\item $\varphi_{15} :\equiv (y + x = x + y \to x = y + x \to x = x + y) \to x \in \mathbb{R} \to y \in \mathbb{R} \to 0 \leq y \to x \leq x + y$
\item $\varphi_{16} :\equiv (x \in \mathbb{R} \to y \in \mathbb{R} \to 0 \leq y \to 0 \leq y) \to x \in \mathbb{R} \to y \in \mathbb{R} \to 0 \leq y \to 0 + x \leq y + x$
\item $\varphi_{17} :\equiv \forall z\,\forall w\,(z = w \to z < y + x \to w < y + x) \to \forall w\,(0 + x = w \to 0 + x < y + x \to w < y + x)$
\item $\varphi_{18} :\equiv \forall y\,(0 \in \mathbb{R} \to y \in \mathbb{R} \to 0 + y = y + 0) \to 0 \in \mathbb{R} \to x \in \mathbb{R} \to 0 + x = x + 0$
\item $\varphi_{19} :\equiv (x \in \mathbb{R} \to y \in \mathbb{R} \to 0 \leq y \to x \in \mathbb{R}) \to x \in \mathbb{R} \to y \in \mathbb{R} \to 0 \leq y \to x + 0 = x$
\item $\varphi_{20} :\equiv \forall w\,(0 + x = w \to 0 + x < y + x \to w < y + x) \to 0 + x = x \to 0 + x < y + x \to x < y + x$
\item $\varphi_{21} :\equiv \forall x\,\forall y\,(x \in \mathbb{R} \to y \in \mathbb{R} \to x + y = y + x) \to \forall y\,(0 \in \mathbb{R} \to y \in \mathbb{R} \to 0 + y = y + 0)$
\item $\varphi_{22} :\equiv (x \in \mathbb{R} \to y \in \mathbb{R} \to 0 \leq y \to x \in \mathbb{R}) \to x \in \mathbb{R} \to y \in \mathbb{R} \to 0 \leq y \to 0 + x = x + 0$
\item $\varphi_{23} :\equiv (x \in \mathbb{R} \to y \in \mathbb{R} \to 0 \leq y \to x \in \mathbb{R}) \to x \in \mathbb{R} \to y \in \mathbb{R} \to 0 \leq y \to y + x = x + y$
\item $\varphi_{24} :\equiv \forall v_{423}\,(y \in \mathbb{R} \to v_{423} \in \mathbb{R} \to y + v_{423} = v_{423} + y) \to y \in \mathbb{R} \to x \in \mathbb{R} \to y + x = x + y$
\item $\varphi_{25} :\equiv \forall x\,\forall y\,(x \in \mathbb{R} \to y \in \mathbb{R} \to x + y = y + x) \to \forall y\,(v_{422} \in \mathbb{R} \to y \in \mathbb{R} \to v_{422} + y = y + v_{422})$
\item $\varphi_{26} :\equiv \forall y\,(v_{422} \in \mathbb{R} \to y \in \mathbb{R} \to v_{422} + y = y + v_{422}) \to v_{422} \in \mathbb{R} \to v_{423} \in \mathbb{R} \to v_{422} + v_{423} = v_{423} + v_{422}$
\item $\varphi_{27} :\equiv \forall z\,(0 \in \mathbb{R} \to y \in \mathbb{R} \to z \in \mathbb{R} \to 0 \leq y \to 0 + z \leq y + z) \to 0 \in \mathbb{R} \to y \in \mathbb{R} \to x \in \mathbb{R} \to 0 \leq y \to 0 + x \leq y + x$
\item $\varphi_{28} :\equiv \forall v_{422}\,(\varphi_{3}) \to \forall v_{423}\,(y \in \mathbb{R} \to v_{423} \in \mathbb{R} \to y + v_{423} = v_{423} + y)$
\item $\varphi_{29} :\equiv \varphi_{5} \to \forall z\,(0 \in \mathbb{R} \to y \in \mathbb{R} \to z \in \mathbb{R} \to 0 \leq y \to 0 + z \leq y + z)$
\item $\varphi_{30} :\equiv (x \in \mathbb{R} \to y \in \mathbb{R} \to 0 \leq y \to x \leq y + x) \to (y + x = x + y \to x < y + x \to x < x + y) \to \varphi_{15}$
\item $\varphi_{31} :\equiv (x \in \mathbb{R} \to y \in \mathbb{R} \to 0 \leq y \to 0 + x = x + 0) \to (x \in \mathbb{R} \to y \in \mathbb{R} \to 0 \leq y \to x + 0 = x) \to \varphi_{13}$
\item $\varphi_{32} :\equiv (x \in \mathbb{R} \to y \in \mathbb{R} \to 0 \leq y \to 0 + x \leq y + x) \to (0 + x = x \to 0 + x < y + x \to x < y + x) \to \varphi_{14}$
\item $\varphi_{33} :\equiv (x \in \mathbb{R} \to y \in \mathbb{R} \to 0 \leq y \to y \in \mathbb{R}) \to (x \in \mathbb{R} \to y \in \mathbb{R} \to 0 \leq y \to x \in \mathbb{R}) \to \varphi_{16}$
\item $\varphi_{34} :\equiv (y \in \mathbb{R} \to x \in \mathbb{R} \to y + x = x + y) \to (x \in \mathbb{R} \to y \in \mathbb{R} \to 0 \leq y \to y \in \mathbb{R}) \to \varphi_{23}$
\item $\varphi_{35} :\equiv (0 \in \mathbb{R} \to y \in \mathbb{R} \to x \in \mathbb{R} \to 0 \leq y \to 0 + x \leq y + x) \to 0 \in \mathbb{R} \to \varphi_{33}$
\end{abbrevlist}
\end{fullproof}

% 自動生成 (prover/build.js): LeAddNonnegL  ノード数 37
\begin{fullproof}{[E2]}{$\forall x\,\forall y\,(x \in \mathbb{R} \to y \in \mathbb{R} \to 0 \leq y \to x \leq y + x)$}{37}\phantomsection\label{tree:LeAddNonnegL}
\noindent$\Pi_{1}$:
\begin{ptree}
\AXC{$x \in \mathbb{R} \to y \in \mathbb{R} \to 0 \leq y \to x \in \mathbb{R}$\lbl{A1}}
\AXC{$\forall x\,\forall y\,(x \in \mathbb{R} \to y \in \mathbb{R} \to x + y = y + x)$\lbl{(R3)}}
\AXC{$\varphi_{8}$\lbl{A2}}
\BIC{$\forall y\,(x \in \mathbb{R} \to y \in \mathbb{R} \to x + y = y + x)$}
\AXC{$\varphi_{6}$\lbl{A2}}
\BIC{$x \in \mathbb{R} \to y \in \mathbb{R} \to x + y = y + x$}
\AXC{$\varphi_{12}$\lbl{A1}}
\BIC{$(x \in \mathbb{R} \to y \in \mathbb{R} \to 0 \leq y \to x \in \mathbb{R}) \to \varphi_{9}$}
\BIC{$\varphi_{9}$}
\end{ptree}
\noindent$\Pi_{2}$:
\begin{ptree}
\AXC{$x \in \mathbb{R} \to y \in \mathbb{R} \to 0 \leq y \to y \in \mathbb{R}$\lbl{A1}}
\AXC{$\stackrel{\Pi_{1}}{\vdots}$}\noLine
\UIC{$\varphi_{9}$}
\BIC{$x \in \mathbb{R} \to y \in \mathbb{R} \to 0 \leq y \to x + y = y + x$}
\AXC{$(x \in \mathbb{R} \to y \in \mathbb{R} \to 0 \leq y \to x + y = y + x) \to \varphi_{11}$\lbl{A1}}
\BIC{$\varphi_{11}$}
\end{ptree}
\noindent$\Pi_{3}$:
\begin{ptree}
\AXC{$\forall x\,\forall y\,(x \in \mathbb{R} \to y \in \mathbb{R} \to 0 \leq y \to x \leq x + y)$\lbl{[E1]}}
\AXC{$\varphi_{10}$\lbl{A2}}
\BIC{$\forall y\,(x \in \mathbb{R} \to y \in \mathbb{R} \to 0 \leq y \to x \leq x + y)$}
\AXC{$\varphi_{7}$\lbl{A2}}
\BIC{$x \in \mathbb{R} \to y \in \mathbb{R} \to 0 \leq y \to x \leq x + y$}
\AXC{$\stackrel{\Pi_{2}}{\vdots}$}\noLine
\UIC{$\varphi_{11}$}
\BIC{$(x + y = y + x \to x < x + y \to x < y + x) \to \varphi_{5}$}
\end{ptree}
\noindent$\Pi_{4}$:
\begin{ptree}
\AXC{$z = w \to x < z \to x < w$\lbl{A6}}
\UIC{$\forall w\,(z = w \to x < z \to x < w)$}
\UIC{$\forall z\,\forall w\,(z = w \to x < z \to x < w)$}
\AXC{$\varphi_{2}$\lbl{A2}}
\BIC{$\forall w\,(x + y = w \to x < x + y \to x < w)$}
\AXC{$\varphi_{4}$\lbl{A2}}
\BIC{$x + y = y + x \to x < x + y \to x < y + x$}
\AXC{$\stackrel{\Pi_{3}}{\vdots}$}\noLine
\UIC{$(x + y = y + x \to x < x + y \to x < y + x) \to \varphi_{5}$}
\BIC{$\varphi_{5}$}
\end{ptree}
\noindent 主木:
\begin{ptree}
\AXC{$z = w \to x = z \to x = w$\lbl{A6}}
\UIC{$\forall w\,(z = w \to x = z \to x = w)$}
\UIC{$\forall z\,\forall w\,(z = w \to x = z \to x = w)$}
\AXC{$\varphi_{1}$\lbl{A2}}
\BIC{$\forall w\,(x + y = w \to x = x + y \to x = w)$}
\AXC{$\varphi_{3}$\lbl{A2}}
\BIC{$x + y = y + x \to x = x + y \to x = y + x$}
\AXC{$\stackrel{\Pi_{4}}{\vdots}$}\noLine
\UIC{$\varphi_{5}$}
\BIC{$x \in \mathbb{R} \to y \in \mathbb{R} \to 0 \leq y \to x \leq y + x$}
\UIC{$\forall y\,(x \in \mathbb{R} \to y \in \mathbb{R} \to 0 \leq y \to x \leq y + x)$}
\UIC{$\forall x\,\forall y\,(x \in \mathbb{R} \to y \in \mathbb{R} \to 0 \leq y \to x \leq y + x)$}
\end{ptree}
\begin{abbrevlist}
\item $\varphi_{1} :\equiv \forall z\,\forall w\,(z = w \to x = z \to x = w) \to \forall w\,(x + y = w \to x = x + y \to x = w)$
\item $\varphi_{2} :\equiv \forall z\,\forall w\,(z = w \to x < z \to x < w) \to \forall w\,(x + y = w \to x < x + y \to x < w)$
\item $\varphi_{3} :\equiv \forall w\,(x + y = w \to x = x + y \to x = w) \to x + y = y + x \to x = x + y \to x = y + x$
\item $\varphi_{4} :\equiv \forall w\,(x + y = w \to x < x + y \to x < w) \to x + y = y + x \to x < x + y \to x < y + x$
\item $\varphi_{5} :\equiv (x + y = y + x \to x = x + y \to x = y + x) \to x \in \mathbb{R} \to y \in \mathbb{R} \to 0 \leq y \to x \leq y + x$
\item $\varphi_{6} :\equiv \forall y\,(x \in \mathbb{R} \to y \in \mathbb{R} \to x + y = y + x) \to x \in \mathbb{R} \to y \in \mathbb{R} \to x + y = y + x$
\item $\varphi_{7} :\equiv \forall y\,(x \in \mathbb{R} \to y \in \mathbb{R} \to 0 \leq y \to x \leq x + y) \to x \in \mathbb{R} \to y \in \mathbb{R} \to 0 \leq y \to x \leq x + y$
\item $\varphi_{8} :\equiv \forall x\,\forall y\,(x \in \mathbb{R} \to y \in \mathbb{R} \to x + y = y + x) \to \forall y\,(x \in \mathbb{R} \to y \in \mathbb{R} \to x + y = y + x)$
\item $\varphi_{9} :\equiv (x \in \mathbb{R} \to y \in \mathbb{R} \to 0 \leq y \to y \in \mathbb{R}) \to x \in \mathbb{R} \to y \in \mathbb{R} \to 0 \leq y \to x + y = y + x$
\item $\varphi_{10} :\equiv \forall x\,\forall y\,(x \in \mathbb{R} \to y \in \mathbb{R} \to 0 \leq y \to x \leq x + y) \to \forall y\,(x \in \mathbb{R} \to y \in \mathbb{R} \to 0 \leq y \to x \leq x + y)$
\item $\varphi_{11} :\equiv (x \in \mathbb{R} \to y \in \mathbb{R} \to 0 \leq y \to x \leq x + y) \to (x + y = y + x \to x < x + y \to x < y + x) \to \varphi_{5}$
\item $\varphi_{12} :\equiv (x \in \mathbb{R} \to y \in \mathbb{R} \to x + y = y + x) \to (x \in \mathbb{R} \to y \in \mathbb{R} \to 0 \leq y \to x \in \mathbb{R}) \to \varphi_{9}$
\end{abbrevlist}
\end{fullproof}

% 自動生成 (prover/build.js): LtAddPos  ノード数 91
\begin{fullproof}{[E3]}{$\forall x\,\forall y\,(x \in \mathbb{R} \to y \in \mathbb{R} \to 0 < y \to x < x + y)$}{91}\phantomsection\label{tree:LtAddPos}
\noindent$\Pi_{1}$:
\begin{ptree}
\AXC{$0 \in \mathbb{R}$\lbl{(R1)}}
\AXC{$\varphi_{6}$\lbl{(R14)}}
\AXC{$\varphi_{6} \to \varphi_{4}$\lbl{A2}}
\BIC{$\varphi_{4}$}
\AXC{$\varphi_{25}$\lbl{A2}}
\BIC{$\forall z\,(0 \in \mathbb{R} \to y \in \mathbb{R} \to z \in \mathbb{R} \to 0 < y \to 0 + z < y + z)$}
\AXC{$\varphi_{24}$\lbl{A2}}
\BIC{$0 \in \mathbb{R} \to y \in \mathbb{R} \to x \in \mathbb{R} \to 0 < y \to 0 + x < y + x$}
\AXC{$\varphi_{31}$\lbl{A1}}
\BIC{$0 \in \mathbb{R} \to \varphi_{30}$}
\BIC{$\varphi_{30}$}
\end{ptree}
\noindent$\Pi_{2}$:
\begin{ptree}
\AXC{$x \in \mathbb{R} \to y \in \mathbb{R} \to 0 < y \to x \in \mathbb{R}$\lbl{A1}}
\AXC{$x \in \mathbb{R} \to y \in \mathbb{R} \to 0 < y \to y \in \mathbb{R}$\lbl{A1}}
\AXC{$\stackrel{\Pi_{1}}{\vdots}$}\noLine
\UIC{$\varphi_{30}$}
\BIC{$(x \in \mathbb{R} \to y \in \mathbb{R} \to 0 < y \to x \in \mathbb{R}) \to \varphi_{17}$}
\BIC{$\varphi_{17}$}
\end{ptree}
\noindent$\Pi_{3}$:
\begin{ptree}
\AXC{$y = z \to x = y \to x = z$\lbl{A6}}
\AXC{$(y = z \to x = y \to x = z) \to x = y \to y = z \to x = z$\lbl{A1}}
\BIC{$x = y \to y = z \to x = z$}
\UIC{$\forall z\,(x = y \to y = z \to x = z)$}
\UIC{$\forall y\,\forall z\,(x = y \to y = z \to x = z)$}
\UIC{$\forall x\,\forall y\,\forall z\,(x = y \to y = z \to x = z)$}
\AXC{$\varphi_{1}$\lbl{A2}}
\BIC{$\forall y\,\forall z\,(0 + x = y \to y = z \to 0 + x = z)$}
\AXC{$\varphi_{5}$\lbl{A2}}
\BIC{$\forall z\,(0 + x = x + 0 \to x + 0 = z \to 0 + x = z)$}
\AXC{$\varphi_{7}$\lbl{A2}}
\BIC{$0 + x = x + 0 \to x + 0 = x \to 0 + x = x$}
\end{ptree}
\noindent$\Pi_{4}$:
\begin{ptree}
\AXC{$\forall x \in \mathbb{R}\,(x + 0 = x)$\lbl{(R4)}}
\AXC{$\forall x \in \mathbb{R}\,(x + 0 = x) \to x \in \mathbb{R} \to x + 0 = x$\lbl{A2}}
\BIC{$x \in \mathbb{R} \to x + 0 = x$}
\AXC{$(x \in \mathbb{R} \to x + 0 = x) \to \varphi_{16}$\lbl{A1}}
\BIC{$\varphi_{16}$}
\end{ptree}
\noindent$\Pi_{5}$:
\begin{ptree}
\AXC{$\forall x\,\forall y\,(x \in \mathbb{R} \to y \in \mathbb{R} \to x + y = y + x)$\lbl{(R3)}}
\AXC{$\varphi_{13}$\lbl{A2}}
\BIC{$\forall y\,(0 \in \mathbb{R} \to y \in \mathbb{R} \to 0 + y = y + 0)$}
\AXC{$\varphi_{11}$\lbl{A2}}
\BIC{$0 \in \mathbb{R} \to x \in \mathbb{R} \to 0 + x = x + 0$}
\AXC{$(0 \in \mathbb{R} \to x \in \mathbb{R} \to 0 + x = x + 0) \to 0 \in \mathbb{R} \to \varphi_{18}$\lbl{A1}}
\BIC{$0 \in \mathbb{R} \to \varphi_{18}$}
\end{ptree}
\noindent$\Pi_{6}$:
\begin{ptree}
\AXC{$x \in \mathbb{R} \to y \in \mathbb{R} \to 0 < y \to x \in \mathbb{R}$\lbl{A1}}
\AXC{$0 \in \mathbb{R}$\lbl{(R1)}}
\AXC{$\stackrel{\Pi_{5}}{\vdots}$}\noLine
\UIC{$0 \in \mathbb{R} \to \varphi_{18}$}
\BIC{$\varphi_{18}$}
\BIC{$x \in \mathbb{R} \to y \in \mathbb{R} \to 0 < y \to 0 + x = x + 0$}
\AXC{$\varphi_{26}$\lbl{A1}}
\BIC{$(x \in \mathbb{R} \to y \in \mathbb{R} \to 0 < y \to x + 0 = x) \to \varphi_{9}$}
\end{ptree}
\noindent$\Pi_{7}$:
\begin{ptree}
\AXC{$x \in \mathbb{R} \to y \in \mathbb{R} \to 0 < y \to x \in \mathbb{R}$\lbl{A1}}
\AXC{$\stackrel{\Pi_{4}}{\vdots}$}\noLine
\UIC{$\varphi_{16}$}
\BIC{$x \in \mathbb{R} \to y \in \mathbb{R} \to 0 < y \to x + 0 = x$}
\AXC{$\stackrel{\Pi_{6}}{\vdots}$}\noLine
\UIC{$(x \in \mathbb{R} \to y \in \mathbb{R} \to 0 < y \to x + 0 = x) \to \varphi_{9}$}
\BIC{$\varphi_{9}$}
\end{ptree}
\noindent$\Pi_{8}$:
\begin{ptree}
\AXC{$x \in \mathbb{R} \to y \in \mathbb{R} \to 0 < y \to 0 < y$\lbl{A1}}
\AXC{$\stackrel{\Pi_{2}}{\vdots}$}\noLine
\UIC{$\varphi_{17}$}
\BIC{$x \in \mathbb{R} \to y \in \mathbb{R} \to 0 < y \to 0 + x < y + x$}
\AXC{$\stackrel{\Pi_{3}}{\vdots}$}\noLine
\UIC{$0 + x = x + 0 \to x + 0 = x \to 0 + x = x$}
\AXC{$\stackrel{\Pi_{7}}{\vdots}$}\noLine
\UIC{$\varphi_{9}$}
\BIC{$x \in \mathbb{R} \to y \in \mathbb{R} \to 0 < y \to 0 + x = x$}
\AXC{$\varphi_{28}$\lbl{A1}}
\BIC{$(x \in \mathbb{R} \to y \in \mathbb{R} \to 0 < y \to 0 + x < y + x) \to \varphi_{14}$}
\BIC{$\varphi_{14}$}
\end{ptree}
\noindent$\Pi_{9}$:
\begin{ptree}
\AXC{$z = w \to z < y + x \to w < y + x$\lbl{A6}}
\UIC{$\forall w\,(z = w \to z < y + x \to w < y + x)$}
\UIC{$\forall z\,\forall w\,(z = w \to z < y + x \to w < y + x)$}
\AXC{$\varphi_{10}$\lbl{A2}}
\BIC{$\forall w\,(0 + x = w \to 0 + x < y + x \to w < y + x)$}
\AXC{$\varphi_{12}$\lbl{A2}}
\BIC{$0 + x = x \to 0 + x < y + x \to x < y + x$}
\AXC{$\stackrel{\Pi_{8}}{\vdots}$}\noLine
\UIC{$\varphi_{14}$}
\BIC{$x \in \mathbb{R} \to y \in \mathbb{R} \to 0 < y \to x < y + x$}
\end{ptree}
\noindent$\Pi_{10}$:
\begin{ptree}
\AXC{$\forall x\,\forall y\,(x \in \mathbb{R} \to y \in \mathbb{R} \to x + y = y + x)$\lbl{(R3)}}
\AXC{$\varphi_{21}$\lbl{A2}}
\BIC{$\forall y\,(v_{424} \in \mathbb{R} \to y \in \mathbb{R} \to v_{424} + y = y + v_{424})$}
\AXC{$\varphi_{22}$\lbl{A2}}
\BIC{$v_{424} \in \mathbb{R} \to v_{425} \in \mathbb{R} \to v_{424} + v_{425} = v_{425} + v_{424}$}
\UIC{$\varphi_{2}$}
\UIC{$\forall v_{424}\,(\varphi_{2})$}
\AXC{$\varphi_{23}$\lbl{A2}}
\BIC{$\forall v_{425}\,(y \in \mathbb{R} \to v_{425} \in \mathbb{R} \to y + v_{425} = v_{425} + y)$}
\AXC{$\varphi_{19}$\lbl{A2}}
\BIC{$y \in \mathbb{R} \to x \in \mathbb{R} \to y + x = x + y$}
\AXC{$\varphi_{27}$\lbl{A1}}
\BIC{$(x \in \mathbb{R} \to y \in \mathbb{R} \to 0 < y \to y \in \mathbb{R}) \to \varphi_{20}$}
\end{ptree}
\noindent$\Pi_{11}$:
\begin{ptree}
\AXC{$x \in \mathbb{R} \to y \in \mathbb{R} \to 0 < y \to x \in \mathbb{R}$\lbl{A1}}
\AXC{$x \in \mathbb{R} \to y \in \mathbb{R} \to 0 < y \to y \in \mathbb{R}$\lbl{A1}}
\AXC{$\stackrel{\Pi_{10}}{\vdots}$}\noLine
\UIC{$(x \in \mathbb{R} \to y \in \mathbb{R} \to 0 < y \to y \in \mathbb{R}) \to \varphi_{20}$}
\BIC{$\varphi_{20}$}
\BIC{$x \in \mathbb{R} \to y \in \mathbb{R} \to 0 < y \to y + x = x + y$}
\end{ptree}
\noindent$\Pi_{12}$:
\begin{ptree}
\AXC{$\stackrel{\Pi_{9}}{\vdots}$}\noLine
\UIC{$x \in \mathbb{R} \to y \in \mathbb{R} \to 0 < y \to x < y + x$}
\AXC{$\stackrel{\Pi_{11}}{\vdots}$}\noLine
\UIC{$x \in \mathbb{R} \to y \in \mathbb{R} \to 0 < y \to y + x = x + y$}
\AXC{$\varphi_{29}$\lbl{A1}}
\BIC{$(x \in \mathbb{R} \to y \in \mathbb{R} \to 0 < y \to x < y + x) \to \varphi_{15}$}
\BIC{$\varphi_{15}$}
\end{ptree}
\noindent 主木:
\begin{ptree}
\AXC{$z = w \to x < z \to x < w$\lbl{A6}}
\UIC{$\forall w\,(z = w \to x < z \to x < w)$}
\UIC{$\forall z\,\forall w\,(z = w \to x < z \to x < w)$}
\AXC{$\varphi_{3}$\lbl{A2}}
\BIC{$\forall w\,(y + x = w \to x < y + x \to x < w)$}
\AXC{$\varphi_{8}$\lbl{A2}}
\BIC{$y + x = x + y \to x < y + x \to x < x + y$}
\AXC{$\stackrel{\Pi_{12}}{\vdots}$}\noLine
\UIC{$\varphi_{15}$}
\BIC{$x \in \mathbb{R} \to y \in \mathbb{R} \to 0 < y \to x < x + y$}
\UIC{$\forall y\,(x \in \mathbb{R} \to y \in \mathbb{R} \to 0 < y \to x < x + y)$}
\UIC{$\forall x\,\forall y\,(x \in \mathbb{R} \to y \in \mathbb{R} \to 0 < y \to x < x + y)$}
\end{ptree}
\begin{abbrevlist}
\item $\varphi_{1} :\equiv \forall x\,\forall y\,\forall z\,(x = y \to y = z \to x = z) \to \forall y\,\forall z\,(0 + x = y \to y = z \to 0 + x = z)$
\item $\varphi_{2} :\equiv \forall v_{425}\,(v_{424} \in \mathbb{R} \to v_{425} \in \mathbb{R} \to v_{424} + v_{425} = v_{425} + v_{424})$
\item $\varphi_{3} :\equiv \forall z\,\forall w\,(z = w \to x < z \to x < w) \to \forall w\,(y + x = w \to x < y + x \to x < w)$
\item $\varphi_{4} :\equiv \forall y\,\forall z\,(0 \in \mathbb{R} \to y \in \mathbb{R} \to z \in \mathbb{R} \to 0 < y \to 0 + z < y + z)$
\item $\varphi_{5} :\equiv \forall y\,\forall z\,(0 + x = y \to y = z \to 0 + x = z) \to \forall z\,(0 + x = x + 0 \to x + 0 = z \to 0 + x = z)$
\item $\varphi_{6} :\equiv \forall x\,\forall y\,\forall z\,(x \in \mathbb{R} \to y \in \mathbb{R} \to z \in \mathbb{R} \to x < y \to x + z < y + z)$
\item $\varphi_{7} :\equiv \forall z\,(0 + x = x + 0 \to x + 0 = z \to 0 + x = z) \to 0 + x = x + 0 \to x + 0 = x \to 0 + x = x$
\item $\varphi_{8} :\equiv \forall w\,(y + x = w \to x < y + x \to x < w) \to y + x = x + y \to x < y + x \to x < x + y$
\item $\varphi_{9} :\equiv (0 + x = x + 0 \to x + 0 = x \to 0 + x = x) \to x \in \mathbb{R} \to y \in \mathbb{R} \to 0 < y \to 0 + x = x$
\item $\varphi_{10} :\equiv \forall z\,\forall w\,(z = w \to z < y + x \to w < y + x) \to \forall w\,(0 + x = w \to 0 + x < y + x \to w < y + x)$
\item $\varphi_{11} :\equiv \forall y\,(0 \in \mathbb{R} \to y \in \mathbb{R} \to 0 + y = y + 0) \to 0 \in \mathbb{R} \to x \in \mathbb{R} \to 0 + x = x + 0$
\item $\varphi_{12} :\equiv \forall w\,(0 + x = w \to 0 + x < y + x \to w < y + x) \to 0 + x = x \to 0 + x < y + x \to x < y + x$
\item $\varphi_{13} :\equiv \forall x\,\forall y\,(x \in \mathbb{R} \to y \in \mathbb{R} \to x + y = y + x) \to \forall y\,(0 \in \mathbb{R} \to y \in \mathbb{R} \to 0 + y = y + 0)$
\item $\varphi_{14} :\equiv (0 + x = x \to 0 + x < y + x \to x < y + x) \to x \in \mathbb{R} \to y \in \mathbb{R} \to 0 < y \to x < y + x$
\item $\varphi_{15} :\equiv (y + x = x + y \to x < y + x \to x < x + y) \to x \in \mathbb{R} \to y \in \mathbb{R} \to 0 < y \to x < x + y$
\item $\varphi_{16} :\equiv (x \in \mathbb{R} \to y \in \mathbb{R} \to 0 < y \to x \in \mathbb{R}) \to x \in \mathbb{R} \to y \in \mathbb{R} \to 0 < y \to x + 0 = x$
\item $\varphi_{17} :\equiv (x \in \mathbb{R} \to y \in \mathbb{R} \to 0 < y \to 0 < y) \to x \in \mathbb{R} \to y \in \mathbb{R} \to 0 < y \to 0 + x < y + x$
\item $\varphi_{18} :\equiv (x \in \mathbb{R} \to y \in \mathbb{R} \to 0 < y \to x \in \mathbb{R}) \to x \in \mathbb{R} \to y \in \mathbb{R} \to 0 < y \to 0 + x = x + 0$
\item $\varphi_{19} :\equiv \forall v_{425}\,(y \in \mathbb{R} \to v_{425} \in \mathbb{R} \to y + v_{425} = v_{425} + y) \to y \in \mathbb{R} \to x \in \mathbb{R} \to y + x = x + y$
\item $\varphi_{20} :\equiv (x \in \mathbb{R} \to y \in \mathbb{R} \to 0 < y \to x \in \mathbb{R}) \to x \in \mathbb{R} \to y \in \mathbb{R} \to 0 < y \to y + x = x + y$
\item $\varphi_{21} :\equiv \forall x\,\forall y\,(x \in \mathbb{R} \to y \in \mathbb{R} \to x + y = y + x) \to \forall y\,(v_{424} \in \mathbb{R} \to y \in \mathbb{R} \to v_{424} + y = y + v_{424})$
\item $\varphi_{22} :\equiv \forall y\,(v_{424} \in \mathbb{R} \to y \in \mathbb{R} \to v_{424} + y = y + v_{424}) \to v_{424} \in \mathbb{R} \to v_{425} \in \mathbb{R} \to v_{424} + v_{425} = v_{425} + v_{424}$
\item $\varphi_{23} :\equiv \forall v_{424}\,(\varphi_{2}) \to \forall v_{425}\,(y \in \mathbb{R} \to v_{425} \in \mathbb{R} \to y + v_{425} = v_{425} + y)$
\item $\varphi_{24} :\equiv \forall z\,(0 \in \mathbb{R} \to y \in \mathbb{R} \to z \in \mathbb{R} \to 0 < y \to 0 + z < y + z) \to 0 \in \mathbb{R} \to y \in \mathbb{R} \to x \in \mathbb{R} \to 0 < y \to 0 + x < y + x$
\item $\varphi_{25} :\equiv \varphi_{4} \to \forall z\,(0 \in \mathbb{R} \to y \in \mathbb{R} \to z \in \mathbb{R} \to 0 < y \to 0 + z < y + z)$
\item $\varphi_{26} :\equiv (x \in \mathbb{R} \to y \in \mathbb{R} \to 0 < y \to 0 + x = x + 0) \to (x \in \mathbb{R} \to y \in \mathbb{R} \to 0 < y \to x + 0 = x) \to \varphi_{9}$
\item $\varphi_{27} :\equiv (y \in \mathbb{R} \to x \in \mathbb{R} \to y + x = x + y) \to (x \in \mathbb{R} \to y \in \mathbb{R} \to 0 < y \to y \in \mathbb{R}) \to \varphi_{20}$
\item $\varphi_{28} :\equiv (x \in \mathbb{R} \to y \in \mathbb{R} \to 0 < y \to 0 + x = x) \to (x \in \mathbb{R} \to y \in \mathbb{R} \to 0 < y \to 0 + x < y + x) \to \varphi_{14}$
\item $\varphi_{29} :\equiv (x \in \mathbb{R} \to y \in \mathbb{R} \to 0 < y \to y + x = x + y) \to (x \in \mathbb{R} \to y \in \mathbb{R} \to 0 < y \to x < y + x) \to \varphi_{15}$
\item $\varphi_{30} :\equiv (x \in \mathbb{R} \to y \in \mathbb{R} \to 0 < y \to y \in \mathbb{R}) \to (x \in \mathbb{R} \to y \in \mathbb{R} \to 0 < y \to x \in \mathbb{R}) \to \varphi_{17}$
\item $\varphi_{31} :\equiv (0 \in \mathbb{R} \to y \in \mathbb{R} \to x \in \mathbb{R} \to 0 < y \to 0 + x < y + x) \to 0 \in \mathbb{R} \to \varphi_{30}$
\end{abbrevlist}
\end{fullproof}

% 自動生成 (prover/build.js): LtAddCancel  ノード数 267
\begin{fullproof}{[E4]}{$\forall x\,\forall y\,\forall z\,(x \in \mathbb{R} \to y \in \mathbb{R} \to z \in \mathbb{R} \to x + z < y + z \to x < y)$}{267}\phantomsection\label{tree:LtAddCancel}
\noindent$\Pi_{1}$:
\begin{ptree}
\AXC{$z_{1} = w \to z_{1} < y + z + -z \to w < y + z + -z$\lbl{A6}}
\UIC{$\forall w\,(z_{1} = w \to z_{1} < y + z + -z \to w < y + z + -z)$}
\UIC{$\forall z_{1}\,\forall w\,(z_{1} = w \to z_{1} < y + z + -z \to w < y + z + -z)$}
\AXC{$\varphi_{81}$\lbl{A2}}
\BIC{$\varphi_{56}$}
\AXC{$\varphi_{56} \to \varphi_{48}$\lbl{A2}}
\BIC{$\varphi_{48}$}
\end{ptree}
\noindent$\Pi_{2}$:
\begin{ptree}
\AXC{$\forall x \in \mathbb{R}\,(-x \in \mathbb{R})$\lbl{(R1)}}
\AXC{$\forall x \in \mathbb{R}\,(-x \in \mathbb{R}) \to z \in \mathbb{R} \to -z \in \mathbb{R}$\lbl{A2}}
\BIC{$z \in \mathbb{R} \to -z \in \mathbb{R}$}
\AXC{$(z \in \mathbb{R} \to -z \in \mathbb{R}) \to \varphi_{78}$\lbl{A1}}
\BIC{$\varphi_{78}$}
\end{ptree}
\noindent$\Pi_{3}$:
\begin{ptree}
\AXC{$x \in \mathbb{R} \to y \in \mathbb{R} \to z \in \mathbb{R} \to x + z < y + z \to z \in \mathbb{R}$\lbl{A1}}
\AXC{$\stackrel{\Pi_{2}}{\vdots}$}\noLine
\UIC{$\varphi_{78}$}
\BIC{$x \in \mathbb{R} \to y \in \mathbb{R} \to z \in \mathbb{R} \to x + z < y + z \to -z \in \mathbb{R}$}
\end{ptree}
\noindent$\Pi_{4}$:
\begin{ptree}
\AXC{$\forall x\,\forall y\,(x \in \mathbb{R} \to y \in \mathbb{R} \to x + y \in \mathbb{R})$\lbl{(R1)}}
\AXC{$\varphi_{69}$\lbl{A2}}
\BIC{$\forall y\,(v_{429} \in \mathbb{R} \to y \in \mathbb{R} \to v_{429} + y \in \mathbb{R})$}
\AXC{$\varphi_{70}$\lbl{A2}}
\BIC{$v_{429} \in \mathbb{R} \to v_{430} \in \mathbb{R} \to v_{429} + v_{430} \in \mathbb{R}$}
\UIC{$\forall v_{430}\,(v_{429} \in \mathbb{R} \to v_{430} \in \mathbb{R} \to v_{429} + v_{430} \in \mathbb{R})$}
\UIC{$\varphi_{17}$}
\AXC{$\varphi_{17} \to \forall v_{430}\,(y \in \mathbb{R} \to v_{430} \in \mathbb{R} \to y + v_{430} \in \mathbb{R})$\lbl{A2}}
\BIC{$\forall v_{430}\,(y \in \mathbb{R} \to v_{430} \in \mathbb{R} \to y + v_{430} \in \mathbb{R})$}
\end{ptree}
\noindent$\Pi_{5}$:
\begin{ptree}
\AXC{$\stackrel{\Pi_{4}}{\vdots}$}\noLine
\UIC{$\forall v_{430}\,(y \in \mathbb{R} \to v_{430} \in \mathbb{R} \to y + v_{430} \in \mathbb{R})$}
\AXC{$\varphi_{68}$\lbl{A2}}
\BIC{$y \in \mathbb{R} \to z \in \mathbb{R} \to y + z \in \mathbb{R}$}
\AXC{$(y \in \mathbb{R} \to z \in \mathbb{R} \to y + z \in \mathbb{R}) \to \varphi_{89}$\lbl{A1}}
\BIC{$\varphi_{89}$}
\end{ptree}
\noindent$\Pi_{6}$:
\begin{ptree}
\AXC{$x \in \mathbb{R} \to y \in \mathbb{R} \to z \in \mathbb{R} \to x + z < y + z \to z \in \mathbb{R}$\lbl{A1}}
\AXC{$x \in \mathbb{R} \to y \in \mathbb{R} \to z \in \mathbb{R} \to x + z < y + z \to y \in \mathbb{R}$\lbl{A1}}
\AXC{$\stackrel{\Pi_{5}}{\vdots}$}\noLine
\UIC{$\varphi_{89}$}
\BIC{$\varphi_{79}$}
\BIC{$x \in \mathbb{R} \to y \in \mathbb{R} \to z \in \mathbb{R} \to x + z < y + z \to y + z \in \mathbb{R}$}
\end{ptree}
\noindent$\Pi_{7}$:
\begin{ptree}
\AXC{$\forall x\,\forall y\,(x \in \mathbb{R} \to y \in \mathbb{R} \to x + y \in \mathbb{R})$\lbl{(R1)}}
\AXC{$\varphi_{67}$\lbl{A2}}
\BIC{$\forall y\,(x \in \mathbb{R} \to y \in \mathbb{R} \to x + y \in \mathbb{R})$}
\AXC{$\varphi_{63}$\lbl{A2}}
\BIC{$x \in \mathbb{R} \to z \in \mathbb{R} \to x + z \in \mathbb{R}$}
\AXC{$(x \in \mathbb{R} \to z \in \mathbb{R} \to x + z \in \mathbb{R}) \to \varphi_{90}$\lbl{A1}}
\BIC{$\varphi_{90}$}
\end{ptree}
\noindent$\Pi_{8}$:
\begin{ptree}
\AXC{$x \in \mathbb{R} \to y \in \mathbb{R} \to z \in \mathbb{R} \to x + z < y + z \to z \in \mathbb{R}$\lbl{A1}}
\AXC{$x \in \mathbb{R} \to y \in \mathbb{R} \to z \in \mathbb{R} \to x + z < y + z \to x \in \mathbb{R}$\lbl{A1}}
\AXC{$\stackrel{\Pi_{7}}{\vdots}$}\noLine
\UIC{$\varphi_{90}$}
\BIC{$\varphi_{80}$}
\BIC{$x \in \mathbb{R} \to y \in \mathbb{R} \to z \in \mathbb{R} \to x + z < y + z \to x + z \in \mathbb{R}$}
\end{ptree}
\noindent$\Pi_{9}$:
\begin{ptree}
\AXC{$\varphi_{30}$\lbl{(R14)}}
\AXC{$\varphi_{30} \to \varphi_{46}$\lbl{A2}}
\BIC{$\varphi_{46}$}
\AXC{$\varphi_{46} \to \varphi_{57}$\lbl{A2}}
\BIC{$\varphi_{57}$}
\AXC{$\varphi_{57} \to v_{426} \in \mathbb{R} \to \varphi_{33}$\lbl{A2}}
\BIC{$v_{426} \in \mathbb{R} \to \varphi_{33}$}
\UIC{$\forall v_{428}\,(v_{426} \in \mathbb{R} \to \varphi_{33})$}
\UIC{$\forall v_{427}\,\forall v_{428}\,(v_{426} \in \mathbb{R} \to \varphi_{33})$}
\UIC{$\forall v_{426}\,\forall v_{427}\,\forall v_{428}\,(v_{426} \in \mathbb{R} \to \varphi_{33})$}
\AXC{$\forall v_{426}\,\forall v_{427}\,\forall v_{428}\,(v_{426} \in \mathbb{R} \to \varphi_{33}) \to \varphi_{72}$\lbl{A2}}
\BIC{$\varphi_{72}$}
\end{ptree}
\noindent$\Pi_{10}$:
\begin{ptree}
\AXC{$\stackrel{\Pi_{9}}{\vdots}$}\noLine
\UIC{$\varphi_{72}$}
\AXC{$\varphi_{72} \to \varphi_{73}$\lbl{A2}}
\BIC{$\varphi_{73}$}
\AXC{$\varphi_{73} \to x + z \in \mathbb{R} \to \varphi_{64}$\lbl{A2}}
\BIC{$x + z \in \mathbb{R} \to \varphi_{64}$}
\AXC{$(x + z \in \mathbb{R} \to \varphi_{64}) \to \varphi_{99}$\lbl{A1}}
\BIC{$\varphi_{99}$}
\end{ptree}
\noindent$\Pi_{11}$:
\begin{ptree}
\AXC{$\stackrel{\Pi_{6}}{\vdots}$}\noLine
\UIC{$x \in \mathbb{R} \to y \in \mathbb{R} \to z \in \mathbb{R} \to x + z < y + z \to y + z \in \mathbb{R}$}
\AXC{$\stackrel{\Pi_{8}}{\vdots}$}\noLine
\UIC{$x \in \mathbb{R} \to y \in \mathbb{R} \to z \in \mathbb{R} \to x + z < y + z \to x + z \in \mathbb{R}$}
\AXC{$\stackrel{\Pi_{10}}{\vdots}$}\noLine
\UIC{$\varphi_{99}$}
\BIC{$\varphi_{98}$}
\BIC{$\varphi_{93}$}
\end{ptree}
\noindent$\Pi_{12}$:
\begin{ptree}
\AXC{$x \in \mathbb{R} \to y \in \mathbb{R} \to z \in \mathbb{R} \to x + z < y + z \to x + z < y + z$\lbl{A1}}
\AXC{$\stackrel{\Pi_{3}}{\vdots}$}\noLine
\UIC{$x \in \mathbb{R} \to y \in \mathbb{R} \to z \in \mathbb{R} \to x + z < y + z \to -z \in \mathbb{R}$}
\AXC{$\stackrel{\Pi_{11}}{\vdots}$}\noLine
\UIC{$\varphi_{93}$}
\BIC{$\varphi_{84}$}
\BIC{$x \in \mathbb{R} \to \varphi_{53}$}
\end{ptree}
\noindent$\Pi_{13}$:
\begin{ptree}
\AXC{$y = z \to x = y \to x = z$\lbl{A6}}
\AXC{$(y = z \to x = y \to x = z) \to x = y \to y = z \to x = z$\lbl{A1}}
\BIC{$x = y \to y = z \to x = z$}
\UIC{$\forall z\,(x = y \to y = z \to x = z)$}
\UIC{$\forall y\,\forall z\,(x = y \to y = z \to x = z)$}
\UIC{$\forall x\,\forall y\,\forall z\,(x = y \to y = z \to x = z)$}
\AXC{$\varphi_{5}$\lbl{A2}}
\BIC{$\forall y\,\forall z\,(v_{446} = y \to y = z \to v_{446} = z)$}
\AXC{$\varphi_{13}$\lbl{A2}}
\BIC{$\forall z\,(v_{446} = v_{447} \to v_{447} = z \to v_{446} = z)$}
\AXC{$\varphi_{18}$\lbl{A2}}
\BIC{$v_{446} = v_{447} \to v_{447} = v_{448} \to v_{446} = v_{448}$}
\UIC{$\forall v_{448}\,(v_{446} = v_{447} \to v_{447} = v_{448} \to v_{446} = v_{448})$}
\UIC{$\varphi_{1}$}
\UIC{$\forall v_{446}\,(\varphi_{1})$}
\end{ptree}
\noindent$\Pi_{14}$:
\begin{ptree}
\AXC{$\stackrel{\Pi_{13}}{\vdots}$}\noLine
\UIC{$\forall v_{446}\,(\varphi_{1})$}
\AXC{$\forall v_{446}\,(\varphi_{1}) \to \varphi_{24}$\lbl{A2}}
\BIC{$\varphi_{24}$}
\AXC{$\varphi_{24} \to \varphi_{22}$\lbl{A2}}
\BIC{$\varphi_{22}$}
\AXC{$\varphi_{22} \to x + z + -z = x + 0 \to x + 0 = x \to x + z + -z = x$\lbl{A2}}
\BIC{$x + z + -z = x + 0 \to x + 0 = x \to x + z + -z = x$}
\end{ptree}
\noindent$\Pi_{15}$:
\begin{ptree}
\AXC{$\forall x \in \mathbb{R}\,(x + 0 = x)$\lbl{(R4)}}
\AXC{$\forall x \in \mathbb{R}\,(x + 0 = x) \to x \in \mathbb{R} \to x + 0 = x$\lbl{A2}}
\BIC{$x \in \mathbb{R} \to x + 0 = x$}
\AXC{$(x \in \mathbb{R} \to x + 0 = x) \to \varphi_{75}$\lbl{A1}}
\BIC{$\varphi_{75}$}
\end{ptree}
\noindent$\Pi_{16}$:
\begin{ptree}
\AXC{$y = z \to x = y \to x = z$\lbl{A6}}
\AXC{$(y = z \to x = y \to x = z) \to x = y \to y = z \to x = z$\lbl{A1}}
\BIC{$x = y \to y = z \to x = z$}
\UIC{$\forall z\,(x = y \to y = z \to x = z)$}
\UIC{$\forall y\,\forall z\,(x = y \to y = z \to x = z)$}
\UIC{$\forall x\,\forall y\,\forall z\,(x = y \to y = z \to x = z)$}
\AXC{$\varphi_{6}$\lbl{A2}}
\BIC{$\forall y\,\forall z\,(v_{443} = y \to y = z \to v_{443} = z)$}
\AXC{$\varphi_{14}$\lbl{A2}}
\BIC{$\forall z\,(v_{443} = v_{444} \to v_{444} = z \to v_{443} = z)$}
\AXC{$\varphi_{19}$\lbl{A2}}
\BIC{$v_{443} = v_{444} \to v_{444} = v_{445} \to v_{443} = v_{445}$}
\UIC{$\forall v_{445}\,(v_{443} = v_{444} \to v_{444} = v_{445} \to v_{443} = v_{445})$}
\UIC{$\varphi_{2}$}
\UIC{$\forall v_{443}\,(\varphi_{2})$}
\end{ptree}
\noindent$\Pi_{17}$:
\begin{ptree}
\AXC{$\stackrel{\Pi_{16}}{\vdots}$}\noLine
\UIC{$\forall v_{443}\,(\varphi_{2})$}
\AXC{$\forall v_{443}\,(\varphi_{2}) \to \varphi_{25}$\lbl{A2}}
\BIC{$\varphi_{25}$}
\AXC{$\varphi_{25} \to \varphi_{61}$\lbl{A2}}
\BIC{$\varphi_{61}$}
\AXC{$\varphi_{61} \to \varphi_{58}$\lbl{A2}}
\BIC{$\varphi_{58}$}
\end{ptree}
\noindent$\Pi_{18}$:
\begin{ptree}
\AXC{$\forall x \in \mathbb{R}\,(x + -x = 0)$\lbl{(R5)}}
\AXC{$\forall x \in \mathbb{R}\,(x + -x = 0) \to z \in \mathbb{R} \to z + -z = 0$\lbl{A2}}
\BIC{$z \in \mathbb{R} \to z + -z = 0$}
\AXC{$(z \in \mathbb{R} \to z + -z = 0) \to \varphi_{77}$\lbl{A1}}
\BIC{$\varphi_{77}$}
\end{ptree}
\noindent$\Pi_{19}$:
\begin{ptree}
\AXC{$x \in \mathbb{R} \to y \in \mathbb{R} \to z \in \mathbb{R} \to x + z < y + z \to z \in \mathbb{R}$\lbl{A1}}
\AXC{$\stackrel{\Pi_{18}}{\vdots}$}\noLine
\UIC{$\varphi_{77}$}
\BIC{$x \in \mathbb{R} \to y \in \mathbb{R} \to z \in \mathbb{R} \to x + z < y + z \to z + -z = 0$}
\AXC{$\varphi_{94}$\lbl{A1}}
\BIC{$x + (z + -z) = x + (z + -z) \to \varphi_{31} \to \varphi_{49}$}
\end{ptree}
\noindent$\Pi_{20}$:
\begin{ptree}
\AXC{$x = x$\lbl{A5}}
\UIC{$\forall x\,(x = x)$}
\AXC{$\forall x\,(x = x) \to x + (z + -z) = x + (z + -z)$\lbl{A2}}
\BIC{$x + (z + -z) = x + (z + -z)$}
\AXC{$\stackrel{\Pi_{19}}{\vdots}$}\noLine
\UIC{$x + (z + -z) = x + (z + -z) \to \varphi_{31} \to \varphi_{49}$}
\BIC{$\varphi_{31} \to \varphi_{49}$}
\end{ptree}
\noindent$\Pi_{21}$:
\begin{ptree}
\AXC{$\varphi_{9}$\lbl{A6}}
\UIC{$\forall w\,(\varphi_{9})$}
\UIC{$\forall z_{1}\,\forall w\,(\varphi_{9})$}
\AXC{$\forall z_{1}\,\forall w\,(\varphi_{9}) \to \varphi_{42}$\lbl{A2}}
\BIC{$\varphi_{42}$}
\AXC{$\varphi_{42} \to \varphi_{31}$\lbl{A2}}
\BIC{$\varphi_{31}$}
\AXC{$\stackrel{\Pi_{20}}{\vdots}$}\noLine
\UIC{$\varphi_{31} \to \varphi_{49}$}
\BIC{$\varphi_{49}$}
\end{ptree}
\noindent$\Pi_{22}$:
\begin{ptree}
\AXC{$\varphi_{29}$\lbl{(R2)}}
\AXC{$\varphi_{29} \to \varphi_{36}$\lbl{A2}}
\BIC{$\varphi_{36}$}
\AXC{$\varphi_{36} \to \varphi_{54}$\lbl{A2}}
\BIC{$\varphi_{54}$}
\AXC{$\varphi_{54} \to v_{440} \in \mathbb{R} \to v_{441} \in \mathbb{R} \to \varphi_{11}$\lbl{A2}}
\BIC{$v_{440} \in \mathbb{R} \to v_{441} \in \mathbb{R} \to \varphi_{11}$}
\UIC{$\forall v_{442}\,(v_{440} \in \mathbb{R} \to v_{441} \in \mathbb{R} \to \varphi_{11})$}
\UIC{$\forall v_{441}\,\forall v_{442}\,(v_{440} \in \mathbb{R} \to v_{441} \in \mathbb{R} \to \varphi_{11})$}
\UIC{$\forall v_{440}\,\forall v_{441}\,\forall v_{442}\,(v_{440} \in \mathbb{R} \to v_{441} \in \mathbb{R} \to \varphi_{11})$}
\AXC{$\varphi_{85}$\lbl{A2}}
\BIC{$\varphi_{65}$}
\AXC{$\varphi_{65} \to \varphi_{37}$\lbl{A2}}
\BIC{$\varphi_{37}$}
\end{ptree}
\noindent$\Pi_{23}$:
\begin{ptree}
\AXC{$x \in \mathbb{R} \to y \in \mathbb{R} \to z \in \mathbb{R} \to x + z < y + z \to x \in \mathbb{R}$\lbl{A1}}
\AXC{$\stackrel{\Pi_{22}}{\vdots}$}\noLine
\UIC{$\varphi_{37}$}
\AXC{$\varphi_{37} \to \varphi_{38}$\lbl{A2}}
\BIC{$\varphi_{38}$}
\AXC{$\varphi_{38} \to \varphi_{96}$\lbl{A1}}
\BIC{$\varphi_{96}$}
\BIC{$\varphi_{91}$}
\end{ptree}
\noindent$\Pi_{24}$:
\begin{ptree}
\AXC{$\stackrel{\Pi_{3}}{\vdots}$}\noLine
\UIC{$x \in \mathbb{R} \to y \in \mathbb{R} \to z \in \mathbb{R} \to x + z < y + z \to -z \in \mathbb{R}$}
\AXC{$x \in \mathbb{R} \to y \in \mathbb{R} \to z \in \mathbb{R} \to x + z < y + z \to z \in \mathbb{R}$\lbl{A1}}
\AXC{$\stackrel{\Pi_{23}}{\vdots}$}\noLine
\UIC{$\varphi_{91}$}
\BIC{$\varphi_{82}$}
\BIC{$x \in \mathbb{R} \to \varphi_{44}$}
\end{ptree}
\noindent$\Pi_{25}$:
\begin{ptree}
\AXC{$\stackrel{\Pi_{17}}{\vdots}$}\noLine
\UIC{$\varphi_{58}$}
\AXC{$\stackrel{\Pi_{21}}{\vdots}$}\noLine
\UIC{$\varphi_{49}$}
\AXC{$\stackrel{\Pi_{24}}{\vdots}$}\noLine
\UIC{$x \in \mathbb{R} \to \varphi_{44}$}
\AXC{$(x \in \mathbb{R} \to \varphi_{44}) \to \varphi_{49} \to \varphi_{58} \to \varphi_{50}$\lbl{A1}}
\BIC{$\varphi_{49} \to \varphi_{58} \to \varphi_{50}$}
\BIC{$\varphi_{58} \to \varphi_{50}$}
\BIC{$\varphi_{50}$}
\AXC{$\varphi_{50} \to \varphi_{87}$\lbl{A1}}
\BIC{$\varphi_{87}$}
\end{ptree}
\noindent$\Pi_{26}$:
\begin{ptree}
\AXC{$\stackrel{\Pi_{14}}{\vdots}$}\noLine
\UIC{$x + z + -z = x + 0 \to x + 0 = x \to x + z + -z = x$}
\AXC{$x \in \mathbb{R} \to y \in \mathbb{R} \to z \in \mathbb{R} \to x + z < y + z \to x \in \mathbb{R}$\lbl{A1}}
\AXC{$\stackrel{\Pi_{15}}{\vdots}$}\noLine
\UIC{$\varphi_{75}$}
\BIC{$x \in \mathbb{R} \to y \in \mathbb{R} \to z \in \mathbb{R} \to x + z < y + z \to x + 0 = x$}
\AXC{$\stackrel{\Pi_{25}}{\vdots}$}\noLine
\UIC{$\varphi_{87}$}
\BIC{$(x + z + -z = x + 0 \to x + 0 = x \to x + z + -z = x) \to \varphi_{34}$}
\BIC{$\varphi_{34}$}
\end{ptree}
\noindent$\Pi_{27}$:
\begin{ptree}
\AXC{$y = z \to x = y \to x = z$\lbl{A6}}
\AXC{$(y = z \to x = y \to x = z) \to x = y \to y = z \to x = z$\lbl{A1}}
\BIC{$x = y \to y = z \to x = z$}
\UIC{$\forall z\,(x = y \to y = z \to x = z)$}
\UIC{$\forall y\,\forall z\,(x = y \to y = z \to x = z)$}
\UIC{$\forall x\,\forall y\,\forall z\,(x = y \to y = z \to x = z)$}
\AXC{$\varphi_{7}$\lbl{A2}}
\BIC{$\forall y\,\forall z\,(v_{437} = y \to y = z \to v_{437} = z)$}
\AXC{$\varphi_{15}$\lbl{A2}}
\BIC{$\forall z\,(v_{437} = v_{438} \to v_{438} = z \to v_{437} = z)$}
\AXC{$\varphi_{20}$\lbl{A2}}
\BIC{$v_{437} = v_{438} \to v_{438} = v_{439} \to v_{437} = v_{439}$}
\UIC{$\forall v_{439}\,(v_{437} = v_{438} \to v_{438} = v_{439} \to v_{437} = v_{439})$}
\UIC{$\varphi_{3}$}
\UIC{$\forall v_{437}\,(\varphi_{3})$}
\end{ptree}
\noindent$\Pi_{28}$:
\begin{ptree}
\AXC{$\stackrel{\Pi_{27}}{\vdots}$}\noLine
\UIC{$\forall v_{437}\,(\varphi_{3})$}
\AXC{$\forall v_{437}\,(\varphi_{3}) \to \varphi_{26}$\lbl{A2}}
\BIC{$\varphi_{26}$}
\AXC{$\varphi_{26} \to \varphi_{23}$\lbl{A2}}
\BIC{$\varphi_{23}$}
\AXC{$\varphi_{23} \to y + z + -z = y + 0 \to y + 0 = y \to y + z + -z = y$\lbl{A2}}
\BIC{$y + z + -z = y + 0 \to y + 0 = y \to y + z + -z = y$}
\end{ptree}
\noindent$\Pi_{29}$:
\begin{ptree}
\AXC{$\forall x \in \mathbb{R}\,(x + 0 = x)$\lbl{(R4)}}
\AXC{$\forall x \in \mathbb{R}\,(x + 0 = x) \to y \in \mathbb{R} \to y + 0 = y$\lbl{A2}}
\BIC{$y \in \mathbb{R} \to y + 0 = y$}
\AXC{$(y \in \mathbb{R} \to y + 0 = y) \to \varphi_{76}$\lbl{A1}}
\BIC{$\varphi_{76}$}
\end{ptree}
\noindent$\Pi_{30}$:
\begin{ptree}
\AXC{$y = z \to x = y \to x = z$\lbl{A6}}
\AXC{$(y = z \to x = y \to x = z) \to x = y \to y = z \to x = z$\lbl{A1}}
\BIC{$x = y \to y = z \to x = z$}
\UIC{$\forall z\,(x = y \to y = z \to x = z)$}
\UIC{$\forall y\,\forall z\,(x = y \to y = z \to x = z)$}
\UIC{$\forall x\,\forall y\,\forall z\,(x = y \to y = z \to x = z)$}
\AXC{$\varphi_{8}$\lbl{A2}}
\BIC{$\forall y\,\forall z\,(v_{434} = y \to y = z \to v_{434} = z)$}
\AXC{$\varphi_{16}$\lbl{A2}}
\BIC{$\forall z\,(v_{434} = v_{435} \to v_{435} = z \to v_{434} = z)$}
\AXC{$\varphi_{21}$\lbl{A2}}
\BIC{$v_{434} = v_{435} \to v_{435} = v_{436} \to v_{434} = v_{436}$}
\UIC{$\forall v_{436}\,(v_{434} = v_{435} \to v_{435} = v_{436} \to v_{434} = v_{436})$}
\UIC{$\varphi_{4}$}
\UIC{$\forall v_{434}\,(\varphi_{4})$}
\end{ptree}
\noindent$\Pi_{31}$:
\begin{ptree}
\AXC{$\stackrel{\Pi_{30}}{\vdots}$}\noLine
\UIC{$\forall v_{434}\,(\varphi_{4})$}
\AXC{$\forall v_{434}\,(\varphi_{4}) \to \varphi_{27}$\lbl{A2}}
\BIC{$\varphi_{27}$}
\AXC{$\varphi_{27} \to \varphi_{62}$\lbl{A2}}
\BIC{$\varphi_{62}$}
\AXC{$\varphi_{62} \to \varphi_{59}$\lbl{A2}}
\BIC{$\varphi_{59}$}
\end{ptree}
\noindent$\Pi_{32}$:
\begin{ptree}
\AXC{$\forall x \in \mathbb{R}\,(x + -x = 0)$\lbl{(R5)}}
\AXC{$\forall x \in \mathbb{R}\,(x + -x = 0) \to z \in \mathbb{R} \to z + -z = 0$\lbl{A2}}
\BIC{$z \in \mathbb{R} \to z + -z = 0$}
\AXC{$(z \in \mathbb{R} \to z + -z = 0) \to \varphi_{77}$\lbl{A1}}
\BIC{$\varphi_{77}$}
\end{ptree}
\noindent$\Pi_{33}$:
\begin{ptree}
\AXC{$x \in \mathbb{R} \to y \in \mathbb{R} \to z \in \mathbb{R} \to x + z < y + z \to z \in \mathbb{R}$\lbl{A1}}
\AXC{$\stackrel{\Pi_{32}}{\vdots}$}\noLine
\UIC{$\varphi_{77}$}
\BIC{$x \in \mathbb{R} \to y \in \mathbb{R} \to z \in \mathbb{R} \to x + z < y + z \to z + -z = 0$}
\AXC{$\varphi_{95}$\lbl{A1}}
\BIC{$y + (z + -z) = y + (z + -z) \to \varphi_{32} \to \varphi_{51}$}
\end{ptree}
\noindent$\Pi_{34}$:
\begin{ptree}
\AXC{$x = x$\lbl{A5}}
\UIC{$\forall x\,(x = x)$}
\AXC{$\forall x\,(x = x) \to y + (z + -z) = y + (z + -z)$\lbl{A2}}
\BIC{$y + (z + -z) = y + (z + -z)$}
\AXC{$\stackrel{\Pi_{33}}{\vdots}$}\noLine
\UIC{$y + (z + -z) = y + (z + -z) \to \varphi_{32} \to \varphi_{51}$}
\BIC{$\varphi_{32} \to \varphi_{51}$}
\end{ptree}
\noindent$\Pi_{35}$:
\begin{ptree}
\AXC{$\varphi_{10}$\lbl{A6}}
\UIC{$\forall w\,(\varphi_{10})$}
\UIC{$\forall z_{1}\,\forall w\,(\varphi_{10})$}
\AXC{$\forall z_{1}\,\forall w\,(\varphi_{10}) \to \varphi_{43}$\lbl{A2}}
\BIC{$\varphi_{43}$}
\AXC{$\varphi_{43} \to \varphi_{32}$\lbl{A2}}
\BIC{$\varphi_{32}$}
\AXC{$\stackrel{\Pi_{34}}{\vdots}$}\noLine
\UIC{$\varphi_{32} \to \varphi_{51}$}
\BIC{$\varphi_{51}$}
\end{ptree}
\noindent$\Pi_{36}$:
\begin{ptree}
\AXC{$\varphi_{29}$\lbl{(R2)}}
\AXC{$\varphi_{29} \to \varphi_{39}$\lbl{A2}}
\BIC{$\varphi_{39}$}
\AXC{$\varphi_{39} \to \varphi_{55}$\lbl{A2}}
\BIC{$\varphi_{55}$}
\AXC{$\varphi_{55} \to v_{431} \in \mathbb{R} \to v_{432} \in \mathbb{R} \to \varphi_{12}$\lbl{A2}}
\BIC{$v_{431} \in \mathbb{R} \to v_{432} \in \mathbb{R} \to \varphi_{12}$}
\UIC{$\forall v_{433}\,(v_{431} \in \mathbb{R} \to v_{432} \in \mathbb{R} \to \varphi_{12})$}
\UIC{$\forall v_{432}\,\forall v_{433}\,(v_{431} \in \mathbb{R} \to v_{432} \in \mathbb{R} \to \varphi_{12})$}
\UIC{$\forall v_{431}\,\forall v_{432}\,\forall v_{433}\,(v_{431} \in \mathbb{R} \to v_{432} \in \mathbb{R} \to \varphi_{12})$}
\AXC{$\varphi_{86}$\lbl{A2}}
\BIC{$\varphi_{66}$}
\AXC{$\varphi_{66} \to \varphi_{40}$\lbl{A2}}
\BIC{$\varphi_{40}$}
\end{ptree}
\noindent$\Pi_{37}$:
\begin{ptree}
\AXC{$x \in \mathbb{R} \to y \in \mathbb{R} \to z \in \mathbb{R} \to x + z < y + z \to y \in \mathbb{R}$\lbl{A1}}
\AXC{$\stackrel{\Pi_{36}}{\vdots}$}\noLine
\UIC{$\varphi_{40}$}
\AXC{$\varphi_{40} \to \varphi_{41}$\lbl{A2}}
\BIC{$\varphi_{41}$}
\AXC{$\varphi_{41} \to \varphi_{97}$\lbl{A1}}
\BIC{$\varphi_{97}$}
\BIC{$\varphi_{92}$}
\end{ptree}
\noindent$\Pi_{38}$:
\begin{ptree}
\AXC{$\stackrel{\Pi_{3}}{\vdots}$}\noLine
\UIC{$x \in \mathbb{R} \to y \in \mathbb{R} \to z \in \mathbb{R} \to x + z < y + z \to -z \in \mathbb{R}$}
\AXC{$x \in \mathbb{R} \to y \in \mathbb{R} \to z \in \mathbb{R} \to x + z < y + z \to z \in \mathbb{R}$\lbl{A1}}
\AXC{$\stackrel{\Pi_{37}}{\vdots}$}\noLine
\UIC{$\varphi_{92}$}
\BIC{$\varphi_{83}$}
\BIC{$x \in \mathbb{R} \to \varphi_{45}$}
\end{ptree}
\noindent$\Pi_{39}$:
\begin{ptree}
\AXC{$\stackrel{\Pi_{31}}{\vdots}$}\noLine
\UIC{$\varphi_{59}$}
\AXC{$\stackrel{\Pi_{35}}{\vdots}$}\noLine
\UIC{$\varphi_{51}$}
\AXC{$\stackrel{\Pi_{38}}{\vdots}$}\noLine
\UIC{$x \in \mathbb{R} \to \varphi_{45}$}
\AXC{$(x \in \mathbb{R} \to \varphi_{45}) \to \varphi_{51} \to \varphi_{59} \to \varphi_{52}$\lbl{A1}}
\BIC{$\varphi_{51} \to \varphi_{59} \to \varphi_{52}$}
\BIC{$\varphi_{59} \to \varphi_{52}$}
\BIC{$\varphi_{52}$}
\AXC{$\varphi_{52} \to \varphi_{88}$\lbl{A1}}
\BIC{$\varphi_{88}$}
\end{ptree}
\noindent$\Pi_{40}$:
\begin{ptree}
\AXC{$\stackrel{\Pi_{28}}{\vdots}$}\noLine
\UIC{$y + z + -z = y + 0 \to y + 0 = y \to y + z + -z = y$}
\AXC{$x \in \mathbb{R} \to y \in \mathbb{R} \to z \in \mathbb{R} \to x + z < y + z \to y \in \mathbb{R}$\lbl{A1}}
\AXC{$\stackrel{\Pi_{29}}{\vdots}$}\noLine
\UIC{$\varphi_{76}$}
\BIC{$x \in \mathbb{R} \to y \in \mathbb{R} \to z \in \mathbb{R} \to x + z < y + z \to y + 0 = y$}
\AXC{$\stackrel{\Pi_{39}}{\vdots}$}\noLine
\UIC{$\varphi_{88}$}
\BIC{$(y + z + -z = y + 0 \to y + 0 = y \to y + z + -z = y) \to \varphi_{35}$}
\BIC{$\varphi_{35}$}
\end{ptree}
\noindent$\Pi_{41}$:
\begin{ptree}
\AXC{$\stackrel{\Pi_{1}}{\vdots}$}\noLine
\UIC{$\varphi_{48}$}
\AXC{$\stackrel{\Pi_{12}}{\vdots}$}\noLine
\UIC{$x \in \mathbb{R} \to \varphi_{53}$}
\AXC{$\stackrel{\Pi_{26}}{\vdots}$}\noLine
\UIC{$\varphi_{34}$}
\AXC{$\varphi_{34} \to (x \in \mathbb{R} \to \varphi_{53}) \to \varphi_{48} \to \varphi_{47}$\lbl{A1}}
\BIC{$(x \in \mathbb{R} \to \varphi_{53}) \to \varphi_{48} \to \varphi_{47}$}
\BIC{$\varphi_{48} \to \varphi_{47}$}
\BIC{$\varphi_{47}$}
\AXC{$\stackrel{\Pi_{40}}{\vdots}$}\noLine
\UIC{$\varphi_{35}$}
\AXC{$\varphi_{35} \to \varphi_{47} \to \varphi_{74}$\lbl{A1}}
\BIC{$\varphi_{47} \to \varphi_{74}$}
\BIC{$\varphi_{74}$}
\end{ptree}
\noindent 主木:
\begin{ptree}
\AXC{$z_{1} = w \to x < z_{1} \to x < w$\lbl{A6}}
\UIC{$\forall w\,(z_{1} = w \to x < z_{1} \to x < w)$}
\UIC{$\forall z_{1}\,\forall w\,(z_{1} = w \to x < z_{1} \to x < w)$}
\AXC{$\varphi_{60}$\lbl{A2}}
\BIC{$\forall w\,(y + z + -z = w \to x < y + z + -z \to x < w)$}
\AXC{$\varphi_{71}$\lbl{A2}}
\BIC{$y + z + -z = y \to x < y + z + -z \to x < y$}
\AXC{$\stackrel{\Pi_{41}}{\vdots}$}\noLine
\UIC{$\varphi_{74}$}
\BIC{$x \in \mathbb{R} \to y \in \mathbb{R} \to z \in \mathbb{R} \to x + z < y + z \to x < y$}
\UIC{$\forall z\,(x \in \mathbb{R} \to y \in \mathbb{R} \to z \in \mathbb{R} \to x + z < y + z \to x < y)$}
\UIC{$\varphi_{28}$}
\UIC{$\forall x\,(\varphi_{28})$}
\end{ptree}
\begin{abbrevlist}
\item $\varphi_{1} :\equiv \forall v_{447}\,\forall v_{448}\,(v_{446} = v_{447} \to v_{447} = v_{448} \to v_{446} = v_{448})$
\item $\varphi_{2} :\equiv \forall v_{444}\,\forall v_{445}\,(v_{443} = v_{444} \to v_{444} = v_{445} \to v_{443} = v_{445})$
\item $\varphi_{3} :\equiv \forall v_{438}\,\forall v_{439}\,(v_{437} = v_{438} \to v_{438} = v_{439} \to v_{437} = v_{439})$
\item $\varphi_{4} :\equiv \forall v_{435}\,\forall v_{436}\,(v_{434} = v_{435} \to v_{435} = v_{436} \to v_{434} = v_{436})$
\item $\varphi_{5} :\equiv \forall x\,\forall y\,\forall z\,(x = y \to y = z \to x = z) \to \forall y\,\forall z\,(v_{446} = y \to y = z \to v_{446} = z)$
\item $\varphi_{6} :\equiv \forall x\,\forall y\,\forall z\,(x = y \to y = z \to x = z) \to \forall y\,\forall z\,(v_{443} = y \to y = z \to v_{443} = z)$
\item $\varphi_{7} :\equiv \forall x\,\forall y\,\forall z\,(x = y \to y = z \to x = z) \to \forall y\,\forall z\,(v_{437} = y \to y = z \to v_{437} = z)$
\item $\varphi_{8} :\equiv \forall x\,\forall y\,\forall z\,(x = y \to y = z \to x = z) \to \forall y\,\forall z\,(v_{434} = y \to y = z \to v_{434} = z)$
\item $\varphi_{9} :\equiv z_{1} = w \to x + (z + -z) = x + z_{1} \to x + (z + -z) = x + w$
\item $\varphi_{10} :\equiv z_{1} = w \to y + (z + -z) = y + z_{1} \to y + (z + -z) = y + w$
\item $\varphi_{11} :\equiv v_{442} \in \mathbb{R} \to v_{440} + v_{441} + v_{442} = v_{440} + (v_{441} + v_{442})$
\item $\varphi_{12} :\equiv v_{433} \in \mathbb{R} \to v_{431} + v_{432} + v_{433} = v_{431} + (v_{432} + v_{433})$
\item $\varphi_{13} :\equiv \forall y\,\forall z\,(v_{446} = y \to y = z \to v_{446} = z) \to \forall z\,(v_{446} = v_{447} \to v_{447} = z \to v_{446} = z)$
\item $\varphi_{14} :\equiv \forall y\,\forall z\,(v_{443} = y \to y = z \to v_{443} = z) \to \forall z\,(v_{443} = v_{444} \to v_{444} = z \to v_{443} = z)$
\item $\varphi_{15} :\equiv \forall y\,\forall z\,(v_{437} = y \to y = z \to v_{437} = z) \to \forall z\,(v_{437} = v_{438} \to v_{438} = z \to v_{437} = z)$
\item $\varphi_{16} :\equiv \forall y\,\forall z\,(v_{434} = y \to y = z \to v_{434} = z) \to \forall z\,(v_{434} = v_{435} \to v_{435} = z \to v_{434} = z)$
\item $\varphi_{17} :\equiv \forall v_{429}\,\forall v_{430}\,(v_{429} \in \mathbb{R} \to v_{430} \in \mathbb{R} \to v_{429} + v_{430} \in \mathbb{R})$
\item $\varphi_{18} :\equiv \forall z\,(v_{446} = v_{447} \to v_{447} = z \to v_{446} = z) \to v_{446} = v_{447} \to v_{447} = v_{448} \to v_{446} = v_{448}$
\item $\varphi_{19} :\equiv \forall z\,(v_{443} = v_{444} \to v_{444} = z \to v_{443} = z) \to v_{443} = v_{444} \to v_{444} = v_{445} \to v_{443} = v_{445}$
\item $\varphi_{20} :\equiv \forall z\,(v_{437} = v_{438} \to v_{438} = z \to v_{437} = z) \to v_{437} = v_{438} \to v_{438} = v_{439} \to v_{437} = v_{439}$
\item $\varphi_{21} :\equiv \forall z\,(v_{434} = v_{435} \to v_{435} = z \to v_{434} = z) \to v_{434} = v_{435} \to v_{435} = v_{436} \to v_{434} = v_{436}$
\item $\varphi_{22} :\equiv \forall v_{448}\,(x + z + -z = x + 0 \to x + 0 = v_{448} \to x + z + -z = v_{448})$
\item $\varphi_{23} :\equiv \forall v_{439}\,(y + z + -z = y + 0 \to y + 0 = v_{439} \to y + z + -z = v_{439})$
\item $\varphi_{24} :\equiv \forall v_{447}\,\forall v_{448}\,(x + z + -z = v_{447} \to v_{447} = v_{448} \to x + z + -z = v_{448})$
\item $\varphi_{25} :\equiv \forall v_{444}\,\forall v_{445}\,(x + z + -z = v_{444} \to v_{444} = v_{445} \to x + z + -z = v_{445})$
\item $\varphi_{26} :\equiv \forall v_{438}\,\forall v_{439}\,(y + z + -z = v_{438} \to v_{438} = v_{439} \to y + z + -z = v_{439})$
\item $\varphi_{27} :\equiv \forall v_{435}\,\forall v_{436}\,(y + z + -z = v_{435} \to v_{435} = v_{436} \to y + z + -z = v_{436})$
\item $\varphi_{28} :\equiv \forall y\,\forall z\,(x \in \mathbb{R} \to y \in \mathbb{R} \to z \in \mathbb{R} \to x + z < y + z \to x < y)$
\item $\varphi_{29} :\equiv \forall x\,\forall y\,\forall z\,(x \in \mathbb{R} \to y \in \mathbb{R} \to z \in \mathbb{R} \to x + y + z = x + (y + z))$
\item $\varphi_{30} :\equiv \forall x\,\forall y\,\forall z\,(x \in \mathbb{R} \to y \in \mathbb{R} \to z \in \mathbb{R} \to x < y \to x + z < y + z)$
\item $\varphi_{31} :\equiv z + -z = 0 \to x + (z + -z) = x + (z + -z) \to x + (z + -z) = x + 0$
\item $\varphi_{32} :\equiv z + -z = 0 \to y + (z + -z) = y + (z + -z) \to y + (z + -z) = y + 0$
\item $\varphi_{33} :\equiv v_{427} \in \mathbb{R} \to v_{428} \in \mathbb{R} \to v_{426} < v_{427} \to v_{426} + v_{428} < v_{427} + v_{428}$
\item $\varphi_{34} :\equiv x \in \mathbb{R} \to y \in \mathbb{R} \to z \in \mathbb{R} \to x + z < y + z \to x + z + -z = x$
\item $\varphi_{35} :\equiv x \in \mathbb{R} \to y \in \mathbb{R} \to z \in \mathbb{R} \to x + z < y + z \to y + z + -z = y$
\item $\varphi_{36} :\equiv \forall y\,\forall z\,(v_{440} \in \mathbb{R} \to y \in \mathbb{R} \to z \in \mathbb{R} \to v_{440} + y + z = v_{440} + (y + z))$
\item $\varphi_{37} :\equiv \forall v_{442}\,(x \in \mathbb{R} \to z \in \mathbb{R} \to v_{442} \in \mathbb{R} \to x + z + v_{442} = x + (z + v_{442}))$
\item $\varphi_{38} :\equiv x \in \mathbb{R} \to z \in \mathbb{R} \to -z \in \mathbb{R} \to x + z + -z = x + (z + -z)$
\item $\varphi_{39} :\equiv \forall y\,\forall z\,(v_{431} \in \mathbb{R} \to y \in \mathbb{R} \to z \in \mathbb{R} \to v_{431} + y + z = v_{431} + (y + z))$
\item $\varphi_{40} :\equiv \forall v_{433}\,(y \in \mathbb{R} \to z \in \mathbb{R} \to v_{433} \in \mathbb{R} \to y + z + v_{433} = y + (z + v_{433}))$
\item $\varphi_{41} :\equiv y \in \mathbb{R} \to z \in \mathbb{R} \to -z \in \mathbb{R} \to y + z + -z = y + (z + -z)$
\item $\varphi_{42} :\equiv \forall w\,(z + -z = w \to x + (z + -z) = x + (z + -z) \to x + (z + -z) = x + w)$
\item $\varphi_{43} :\equiv \forall w\,(z + -z = w \to y + (z + -z) = y + (z + -z) \to y + (z + -z) = y + w)$
\item $\varphi_{44} :\equiv y \in \mathbb{R} \to z \in \mathbb{R} \to x + z < y + z \to x + z + -z = x + (z + -z)$
\item $\varphi_{45} :\equiv y \in \mathbb{R} \to z \in \mathbb{R} \to x + z < y + z \to y + z + -z = y + (z + -z)$
\item $\varphi_{46} :\equiv \forall y\,\forall z\,(v_{426} \in \mathbb{R} \to y \in \mathbb{R} \to z \in \mathbb{R} \to v_{426} < y \to v_{426} + z < y + z)$
\item $\varphi_{47} :\equiv x \in \mathbb{R} \to y \in \mathbb{R} \to z \in \mathbb{R} \to x + z < y + z \to x < y + z + -z$
\item $\varphi_{48} :\equiv x + z + -z = x \to x + z + -z < y + z + -z \to x < y + z + -z$
\item $\varphi_{49} :\equiv x \in \mathbb{R} \to y \in \mathbb{R} \to z \in \mathbb{R} \to x + z < y + z \to x + (z + -z) = x + 0$
\item $\varphi_{50} :\equiv x \in \mathbb{R} \to y \in \mathbb{R} \to z \in \mathbb{R} \to x + z < y + z \to x + z + -z = x + 0$
\item $\varphi_{51} :\equiv x \in \mathbb{R} \to y \in \mathbb{R} \to z \in \mathbb{R} \to x + z < y + z \to y + (z + -z) = y + 0$
\item $\varphi_{52} :\equiv x \in \mathbb{R} \to y \in \mathbb{R} \to z \in \mathbb{R} \to x + z < y + z \to y + z + -z = y + 0$
\item $\varphi_{53} :\equiv y \in \mathbb{R} \to z \in \mathbb{R} \to x + z < y + z \to x + z + -z < y + z + -z$
\item $\varphi_{54} :\equiv \forall z\,(v_{440} \in \mathbb{R} \to v_{441} \in \mathbb{R} \to z \in \mathbb{R} \to v_{440} + v_{441} + z = v_{440} + (v_{441} + z))$
\item $\varphi_{55} :\equiv \forall z\,(v_{431} \in \mathbb{R} \to v_{432} \in \mathbb{R} \to z \in \mathbb{R} \to v_{431} + v_{432} + z = v_{431} + (v_{432} + z))$
\item $\varphi_{56} :\equiv \forall w\,(x + z + -z = w \to x + z + -z < y + z + -z \to w < y + z + -z)$
\item $\varphi_{57} :\equiv \forall z\,(v_{426} \in \mathbb{R} \to v_{427} \in \mathbb{R} \to z \in \mathbb{R} \to v_{426} < v_{427} \to v_{426} + z < v_{427} + z)$
\item $\varphi_{58} :\equiv x + z + -z = x + (z + -z) \to x + (z + -z) = x + 0 \to x + z + -z = x + 0$
\item $\varphi_{59} :\equiv y + z + -z = y + (z + -z) \to y + (z + -z) = y + 0 \to y + z + -z = y + 0$
\item $\varphi_{60} :\equiv \forall z_{1}\,\forall w\,(z_{1} = w \to x < z_{1} \to x < w) \to \forall w\,(y + z + -z = w \to x < y + z + -z \to x < w)$
\item $\varphi_{61} :\equiv \forall v_{445}\,(x + z + -z = x + (z + -z) \to x + (z + -z) = v_{445} \to x + z + -z = v_{445})$
\item $\varphi_{62} :\equiv \forall v_{436}\,(y + z + -z = y + (z + -z) \to y + (z + -z) = v_{436} \to y + z + -z = v_{436})$
\item $\varphi_{63} :\equiv \forall y\,(x \in \mathbb{R} \to y \in \mathbb{R} \to x + y \in \mathbb{R}) \to x \in \mathbb{R} \to z \in \mathbb{R} \to x + z \in \mathbb{R}$
\item $\varphi_{64} :\equiv y + z \in \mathbb{R} \to -z \in \mathbb{R} \to x + z < y + z \to x + z + -z < y + z + -z$
\item $\varphi_{65} :\equiv \forall v_{441}\,\forall v_{442}\,(x \in \mathbb{R} \to v_{441} \in \mathbb{R} \to v_{442} \in \mathbb{R} \to x + v_{441} + v_{442} = x + (v_{441} + v_{442}))$
\item $\varphi_{66} :\equiv \forall v_{432}\,\forall v_{433}\,(y \in \mathbb{R} \to v_{432} \in \mathbb{R} \to v_{433} \in \mathbb{R} \to y + v_{432} + v_{433} = y + (v_{432} + v_{433}))$
\item $\varphi_{67} :\equiv \forall x\,\forall y\,(x \in \mathbb{R} \to y \in \mathbb{R} \to x + y \in \mathbb{R}) \to \forall y\,(x \in \mathbb{R} \to y \in \mathbb{R} \to x + y \in \mathbb{R})$
\item $\varphi_{68} :\equiv \forall v_{430}\,(y \in \mathbb{R} \to v_{430} \in \mathbb{R} \to y + v_{430} \in \mathbb{R}) \to y \in \mathbb{R} \to z \in \mathbb{R} \to y + z \in \mathbb{R}$
\item $\varphi_{69} :\equiv \forall x\,\forall y\,(x \in \mathbb{R} \to y \in \mathbb{R} \to x + y \in \mathbb{R}) \to \forall y\,(v_{429} \in \mathbb{R} \to y \in \mathbb{R} \to v_{429} + y \in \mathbb{R})$
\item $\varphi_{70} :\equiv \forall y\,(v_{429} \in \mathbb{R} \to y \in \mathbb{R} \to v_{429} + y \in \mathbb{R}) \to v_{429} \in \mathbb{R} \to v_{430} \in \mathbb{R} \to v_{429} + v_{430} \in \mathbb{R}$
\item $\varphi_{71} :\equiv \forall w\,(y + z + -z = w \to x < y + z + -z \to x < w) \to y + z + -z = y \to x < y + z + -z \to x < y$
\item $\varphi_{72} :\equiv \forall v_{427}\,\forall v_{428}\,(x + z \in \mathbb{R} \to v_{427} \in \mathbb{R} \to v_{428} \in \mathbb{R} \to x + z < v_{427} \to x + z + v_{428} < v_{427} + v_{428})$
\item $\varphi_{73} :\equiv \forall v_{428}\,(x + z \in \mathbb{R} \to y + z \in \mathbb{R} \to v_{428} \in \mathbb{R} \to x + z < y + z \to x + z + v_{428} < y + z + v_{428})$
\item $\varphi_{74} :\equiv (y + z + -z = y \to x < y + z + -z \to x < y) \to x \in \mathbb{R} \to y \in \mathbb{R} \to z \in \mathbb{R} \to x + z < y + z \to x < y$
\item $\varphi_{75} :\equiv (x \in \mathbb{R} \to y \in \mathbb{R} \to z \in \mathbb{R} \to x + z < y + z \to x \in \mathbb{R}) \to x \in \mathbb{R} \to y \in \mathbb{R} \to z \in \mathbb{R} \to x + z < y + z \to x + 0 = x$
\item $\varphi_{76} :\equiv (x \in \mathbb{R} \to y \in \mathbb{R} \to z \in \mathbb{R} \to x + z < y + z \to y \in \mathbb{R}) \to x \in \mathbb{R} \to y \in \mathbb{R} \to z \in \mathbb{R} \to x + z < y + z \to y + 0 = y$
\item $\varphi_{77} :\equiv (x \in \mathbb{R} \to y \in \mathbb{R} \to z \in \mathbb{R} \to x + z < y + z \to z \in \mathbb{R}) \to x \in \mathbb{R} \to y \in \mathbb{R} \to z \in \mathbb{R} \to x + z < y + z \to z + -z = 0$
\item $\varphi_{78} :\equiv (x \in \mathbb{R} \to y \in \mathbb{R} \to z \in \mathbb{R} \to x + z < y + z \to z \in \mathbb{R}) \to x \in \mathbb{R} \to y \in \mathbb{R} \to z \in \mathbb{R} \to x + z < y + z \to -z \in \mathbb{R}$
\item $\varphi_{79} :\equiv (x \in \mathbb{R} \to y \in \mathbb{R} \to z \in \mathbb{R} \to x + z < y + z \to z \in \mathbb{R}) \to x \in \mathbb{R} \to y \in \mathbb{R} \to z \in \mathbb{R} \to x + z < y + z \to y + z \in \mathbb{R}$
\item $\varphi_{80} :\equiv (x \in \mathbb{R} \to y \in \mathbb{R} \to z \in \mathbb{R} \to x + z < y + z \to z \in \mathbb{R}) \to x \in \mathbb{R} \to y \in \mathbb{R} \to z \in \mathbb{R} \to x + z < y + z \to x + z \in \mathbb{R}$
\item $\varphi_{81} :\equiv \forall z_{1}\,\forall w\,(z_{1} = w \to z_{1} < y + z + -z \to w < y + z + -z) \to \varphi_{56}$
\item $\varphi_{82} :\equiv (x \in \mathbb{R} \to y \in \mathbb{R} \to z \in \mathbb{R} \to x + z < y + z \to -z \in \mathbb{R}) \to x \in \mathbb{R} \to \varphi_{44}$
\item $\varphi_{83} :\equiv (x \in \mathbb{R} \to y \in \mathbb{R} \to z \in \mathbb{R} \to x + z < y + z \to -z \in \mathbb{R}) \to x \in \mathbb{R} \to \varphi_{45}$
\item $\varphi_{84} :\equiv (x \in \mathbb{R} \to y \in \mathbb{R} \to z \in \mathbb{R} \to x + z < y + z \to x + z < y + z) \to x \in \mathbb{R} \to \varphi_{53}$
\item $\varphi_{85} :\equiv \forall v_{440}\,\forall v_{441}\,\forall v_{442}\,(v_{440} \in \mathbb{R} \to v_{441} \in \mathbb{R} \to \varphi_{11}) \to \varphi_{65}$
\item $\varphi_{86} :\equiv \forall v_{431}\,\forall v_{432}\,\forall v_{433}\,(v_{431} \in \mathbb{R} \to v_{432} \in \mathbb{R} \to \varphi_{12}) \to \varphi_{66}$
\item $\varphi_{87} :\equiv (x \in \mathbb{R} \to y \in \mathbb{R} \to z \in \mathbb{R} \to x + z < y + z \to x + 0 = x) \to (x + z + -z = x + 0 \to x + 0 = x \to x + z + -z = x) \to \varphi_{34}$
\item $\varphi_{88} :\equiv (x \in \mathbb{R} \to y \in \mathbb{R} \to z \in \mathbb{R} \to x + z < y + z \to y + 0 = y) \to (y + z + -z = y + 0 \to y + 0 = y \to y + z + -z = y) \to \varphi_{35}$
\item $\varphi_{89} :\equiv (x \in \mathbb{R} \to y \in \mathbb{R} \to z \in \mathbb{R} \to x + z < y + z \to y \in \mathbb{R}) \to \varphi_{79}$
\item $\varphi_{90} :\equiv (x \in \mathbb{R} \to y \in \mathbb{R} \to z \in \mathbb{R} \to x + z < y + z \to x \in \mathbb{R}) \to \varphi_{80}$
\item $\varphi_{91} :\equiv (x \in \mathbb{R} \to y \in \mathbb{R} \to z \in \mathbb{R} \to x + z < y + z \to z \in \mathbb{R}) \to \varphi_{82}$
\item $\varphi_{92} :\equiv (x \in \mathbb{R} \to y \in \mathbb{R} \to z \in \mathbb{R} \to x + z < y + z \to z \in \mathbb{R}) \to \varphi_{83}$
\item $\varphi_{93} :\equiv (x \in \mathbb{R} \to y \in \mathbb{R} \to z \in \mathbb{R} \to x + z < y + z \to -z \in \mathbb{R}) \to \varphi_{84}$
\item $\varphi_{94} :\equiv (x \in \mathbb{R} \to y \in \mathbb{R} \to z \in \mathbb{R} \to x + z < y + z \to z + -z = 0) \to x + (z + -z) = x + (z + -z) \to \varphi_{31} \to \varphi_{49}$
\item $\varphi_{95} :\equiv (x \in \mathbb{R} \to y \in \mathbb{R} \to z \in \mathbb{R} \to x + z < y + z \to z + -z = 0) \to y + (z + -z) = y + (z + -z) \to \varphi_{32} \to \varphi_{51}$
\item $\varphi_{96} :\equiv (x \in \mathbb{R} \to y \in \mathbb{R} \to z \in \mathbb{R} \to x + z < y + z \to x \in \mathbb{R}) \to \varphi_{91}$
\item $\varphi_{97} :\equiv (x \in \mathbb{R} \to y \in \mathbb{R} \to z \in \mathbb{R} \to x + z < y + z \to y \in \mathbb{R}) \to \varphi_{92}$
\item $\varphi_{98} :\equiv (x \in \mathbb{R} \to y \in \mathbb{R} \to z \in \mathbb{R} \to x + z < y + z \to y + z \in \mathbb{R}) \to \varphi_{93}$
\item $\varphi_{99} :\equiv (x \in \mathbb{R} \to y \in \mathbb{R} \to z \in \mathbb{R} \to x + z < y + z \to x + z \in \mathbb{R}) \to \varphi_{98}$
\end{abbrevlist}
\end{fullproof}

% 自動生成 (prover/build.js): AddPosPos  ノード数 27
\begin{fullproof}{[E5]}{$\forall x\,\forall y\,(x \in \mathbb{R} \to y \in \mathbb{R} \to 0 < x \to 0 < y \to 0 < x + y)$}{27}\phantomsection\label{tree:AddPosPos}
\noindent$\Pi_{1}$:
\begin{ptree}
\AXC{$\forall x\,\forall y\,(x \in \mathbb{R} \to y \in \mathbb{R} \to 0 < y \to x < x + y)$\lbl{[E3]}}
\AXC{$\varphi_{5}$\lbl{A2}}
\BIC{$\forall y\,(x \in \mathbb{R} \to y \in \mathbb{R} \to 0 < y \to x < x + y)$}
\AXC{$\varphi_{4}$\lbl{A2}}
\BIC{$x \in \mathbb{R} \to y \in \mathbb{R} \to 0 < y \to x < x + y$}
\AXC{$(x \in \mathbb{R} \to y \in \mathbb{R} \to 0 < y \to x < x + y) \to \varphi_{9}$\lbl{A1}}
\BIC{$\varphi_{9}$}
\end{ptree}
\noindent$\Pi_{2}$:
\begin{ptree}
\AXC{$x \in \mathbb{R} \to y \in \mathbb{R} \to 0 < x \to 0 < y \to y \in \mathbb{R}$\lbl{A1}}
\AXC{$x \in \mathbb{R} \to y \in \mathbb{R} \to 0 < x \to 0 < y \to x \in \mathbb{R}$\lbl{A1}}
\AXC{$\stackrel{\Pi_{1}}{\vdots}$}\noLine
\UIC{$\varphi_{9}$}
\BIC{$(x \in \mathbb{R} \to y \in \mathbb{R} \to 0 < x \to 0 < y \to y \in \mathbb{R}) \to \varphi_{6}$}
\BIC{$\varphi_{6}$}
\end{ptree}
\noindent$\Pi_{3}$:
\begin{ptree}
\AXC{$\forall x\,\forall y\,\forall z\,(x < y \to y < z \to x < z)$\lbl{[O2]}}
\AXC{$\varphi_{3}$\lbl{A2}}
\BIC{$\forall y\,\forall z\,(0 < y \to y < z \to 0 < z)$}
\AXC{$\varphi_{1}$\lbl{A2}}
\BIC{$\forall z\,(0 < x \to x < z \to 0 < z)$}
\AXC{$\varphi_{2}$\lbl{A2}}
\BIC{$0 < x \to x < x + y \to 0 < x + y$}
\AXC{$\varphi_{8}$\lbl{A1}}
\BIC{$(x \in \mathbb{R} \to y \in \mathbb{R} \to 0 < x \to 0 < y \to 0 < x) \to \varphi_{7}$}
\end{ptree}
\noindent$\Pi_{4}$:
\begin{ptree}
\AXC{$x \in \mathbb{R} \to y \in \mathbb{R} \to 0 < x \to 0 < y \to 0 < x$\lbl{A1}}
\AXC{$\stackrel{\Pi_{3}}{\vdots}$}\noLine
\UIC{$(x \in \mathbb{R} \to y \in \mathbb{R} \to 0 < x \to 0 < y \to 0 < x) \to \varphi_{7}$}
\BIC{$\varphi_{7}$}
\end{ptree}
\noindent 主木:
\begin{ptree}
\AXC{$x \in \mathbb{R} \to y \in \mathbb{R} \to 0 < x \to 0 < y \to 0 < y$\lbl{A1}}
\AXC{$\stackrel{\Pi_{2}}{\vdots}$}\noLine
\UIC{$\varphi_{6}$}
\BIC{$x \in \mathbb{R} \to y \in \mathbb{R} \to 0 < x \to 0 < y \to x < x + y$}
\AXC{$\stackrel{\Pi_{4}}{\vdots}$}\noLine
\UIC{$\varphi_{7}$}
\BIC{$x \in \mathbb{R} \to y \in \mathbb{R} \to 0 < x \to 0 < y \to 0 < x + y$}
\UIC{$\forall y\,(x \in \mathbb{R} \to y \in \mathbb{R} \to 0 < x \to 0 < y \to 0 < x + y)$}
\UIC{$\forall x\,\forall y\,(x \in \mathbb{R} \to y \in \mathbb{R} \to 0 < x \to 0 < y \to 0 < x + y)$}
\end{ptree}
\begin{abbrevlist}
\item $\varphi_{1} :\equiv \forall y\,\forall z\,(0 < y \to y < z \to 0 < z) \to \forall z\,(0 < x \to x < z \to 0 < z)$
\item $\varphi_{2} :\equiv \forall z\,(0 < x \to x < z \to 0 < z) \to 0 < x \to x < x + y \to 0 < x + y$
\item $\varphi_{3} :\equiv \forall x\,\forall y\,\forall z\,(x < y \to y < z \to x < z) \to \forall y\,\forall z\,(0 < y \to y < z \to 0 < z)$
\item $\varphi_{4} :\equiv \forall y\,(x \in \mathbb{R} \to y \in \mathbb{R} \to 0 < y \to x < x + y) \to x \in \mathbb{R} \to y \in \mathbb{R} \to 0 < y \to x < x + y$
\item $\varphi_{5} :\equiv \forall x\,\forall y\,(x \in \mathbb{R} \to y \in \mathbb{R} \to 0 < y \to x < x + y) \to \forall y\,(x \in \mathbb{R} \to y \in \mathbb{R} \to 0 < y \to x < x + y)$
\item $\varphi_{6} :\equiv (x \in \mathbb{R} \to y \in \mathbb{R} \to 0 < x \to 0 < y \to 0 < y) \to x \in \mathbb{R} \to y \in \mathbb{R} \to 0 < x \to 0 < y \to x < x + y$
\item $\varphi_{7} :\equiv (x \in \mathbb{R} \to y \in \mathbb{R} \to 0 < x \to 0 < y \to x < x + y) \to x \in \mathbb{R} \to y \in \mathbb{R} \to 0 < x \to 0 < y \to 0 < x + y$
\item $\varphi_{8} :\equiv (0 < x \to x < x + y \to 0 < x + y) \to (x \in \mathbb{R} \to y \in \mathbb{R} \to 0 < x \to 0 < y \to 0 < x) \to \varphi_{7}$
\item $\varphi_{9} :\equiv (x \in \mathbb{R} \to y \in \mathbb{R} \to 0 < x \to 0 < y \to x \in \mathbb{R}) \to (x \in \mathbb{R} \to y \in \mathbb{R} \to 0 < x \to 0 < y \to y \in \mathbb{R}) \to \varphi_{6}$
\end{abbrevlist}
\end{fullproof}

% 自動生成 (prover/build.js): AddNonnegNonneg  ノード数 27
\begin{fullproof}{[E6]}{$\forall x\,\forall y\,(x \in \mathbb{R} \to y \in \mathbb{R} \to 0 \leq x \to 0 \leq y \to 0 \leq x + y)$}{27}\phantomsection\label{tree:AddNonnegNonneg}
\noindent$\Pi_{1}$:
\begin{ptree}
\AXC{$\forall x\,\forall y\,(x \in \mathbb{R} \to y \in \mathbb{R} \to 0 \leq y \to x \leq x + y)$\lbl{[E1]}}
\AXC{$\varphi_{6}$\lbl{A2}}
\BIC{$\forall y\,(x \in \mathbb{R} \to y \in \mathbb{R} \to 0 \leq y \to x \leq x + y)$}
\AXC{$\varphi_{4}$\lbl{A2}}
\BIC{$x \in \mathbb{R} \to y \in \mathbb{R} \to 0 \leq y \to x \leq x + y$}
\AXC{$(x \in \mathbb{R} \to y \in \mathbb{R} \to 0 \leq y \to x \leq x + y) \to \varphi_{9}$\lbl{A1}}
\BIC{$\varphi_{9}$}
\end{ptree}
\noindent$\Pi_{2}$:
\begin{ptree}
\AXC{$x \in \mathbb{R} \to y \in \mathbb{R} \to 0 \leq x \to 0 \leq y \to y \in \mathbb{R}$\lbl{A1}}
\AXC{$x \in \mathbb{R} \to y \in \mathbb{R} \to 0 \leq x \to 0 \leq y \to x \in \mathbb{R}$\lbl{A1}}
\AXC{$\stackrel{\Pi_{1}}{\vdots}$}\noLine
\UIC{$\varphi_{9}$}
\BIC{$(x \in \mathbb{R} \to y \in \mathbb{R} \to 0 \leq x \to 0 \leq y \to y \in \mathbb{R}) \to \varphi_{5}$}
\BIC{$\varphi_{5}$}
\end{ptree}
\noindent$\Pi_{3}$:
\begin{ptree}
\AXC{$\forall x\,\forall y\,\forall z\,(x \leq y \to y \leq z \to x \leq z)$\lbl{[O5]}}
\AXC{$\varphi_{3}$\lbl{A2}}
\BIC{$\forall y\,\forall z\,(0 \leq y \to y \leq z \to 0 \leq z)$}
\AXC{$\varphi_{1}$\lbl{A2}}
\BIC{$\forall z\,(0 \leq x \to x \leq z \to 0 \leq z)$}
\AXC{$\varphi_{2}$\lbl{A2}}
\BIC{$0 \leq x \to x \leq x + y \to 0 \leq x + y$}
\AXC{$\varphi_{8}$\lbl{A1}}
\BIC{$(x \in \mathbb{R} \to y \in \mathbb{R} \to 0 \leq x \to 0 \leq y \to 0 \leq x) \to \varphi_{7}$}
\end{ptree}
\noindent$\Pi_{4}$:
\begin{ptree}
\AXC{$x \in \mathbb{R} \to y \in \mathbb{R} \to 0 \leq x \to 0 \leq y \to 0 \leq x$\lbl{A1}}
\AXC{$\stackrel{\Pi_{3}}{\vdots}$}\noLine
\UIC{$(x \in \mathbb{R} \to y \in \mathbb{R} \to 0 \leq x \to 0 \leq y \to 0 \leq x) \to \varphi_{7}$}
\BIC{$\varphi_{7}$}
\end{ptree}
\noindent 主木:
\begin{ptree}
\AXC{$x \in \mathbb{R} \to y \in \mathbb{R} \to 0 \leq x \to 0 \leq y \to 0 \leq y$\lbl{A1}}
\AXC{$\stackrel{\Pi_{2}}{\vdots}$}\noLine
\UIC{$\varphi_{5}$}
\BIC{$x \in \mathbb{R} \to y \in \mathbb{R} \to 0 \leq x \to 0 \leq y \to x \leq x + y$}
\AXC{$\stackrel{\Pi_{4}}{\vdots}$}\noLine
\UIC{$\varphi_{7}$}
\BIC{$x \in \mathbb{R} \to y \in \mathbb{R} \to 0 \leq x \to 0 \leq y \to 0 \leq x + y$}
\UIC{$\forall y\,(x \in \mathbb{R} \to y \in \mathbb{R} \to 0 \leq x \to 0 \leq y \to 0 \leq x + y)$}
\UIC{$\forall x\,\forall y\,(x \in \mathbb{R} \to y \in \mathbb{R} \to 0 \leq x \to 0 \leq y \to 0 \leq x + y)$}
\end{ptree}
\begin{abbrevlist}
\item $\varphi_{1} :\equiv \forall y\,\forall z\,(0 \leq y \to y \leq z \to 0 \leq z) \to \forall z\,(0 \leq x \to x \leq z \to 0 \leq z)$
\item $\varphi_{2} :\equiv \forall z\,(0 \leq x \to x \leq z \to 0 \leq z) \to 0 \leq x \to x \leq x + y \to 0 \leq x + y$
\item $\varphi_{3} :\equiv \forall x\,\forall y\,\forall z\,(x \leq y \to y \leq z \to x \leq z) \to \forall y\,\forall z\,(0 \leq y \to y \leq z \to 0 \leq z)$
\item $\varphi_{4} :\equiv \forall y\,(x \in \mathbb{R} \to y \in \mathbb{R} \to 0 \leq y \to x \leq x + y) \to x \in \mathbb{R} \to y \in \mathbb{R} \to 0 \leq y \to x \leq x + y$
\item $\varphi_{5} :\equiv (x \in \mathbb{R} \to y \in \mathbb{R} \to 0 \leq x \to 0 \leq y \to 0 \leq y) \to x \in \mathbb{R} \to y \in \mathbb{R} \to 0 \leq x \to 0 \leq y \to x \leq x + y$
\item $\varphi_{6} :\equiv \forall x\,\forall y\,(x \in \mathbb{R} \to y \in \mathbb{R} \to 0 \leq y \to x \leq x + y) \to \forall y\,(x \in \mathbb{R} \to y \in \mathbb{R} \to 0 \leq y \to x \leq x + y)$
\item $\varphi_{7} :\equiv (x \in \mathbb{R} \to y \in \mathbb{R} \to 0 \leq x \to 0 \leq y \to x \leq x + y) \to x \in \mathbb{R} \to y \in \mathbb{R} \to 0 \leq x \to 0 \leq y \to 0 \leq x + y$
\item $\varphi_{8} :\equiv (0 \leq x \to x \leq x + y \to 0 \leq x + y) \to (x \in \mathbb{R} \to y \in \mathbb{R} \to 0 \leq x \to 0 \leq y \to 0 \leq x) \to \varphi_{7}$
\item $\varphi_{9} :\equiv (x \in \mathbb{R} \to y \in \mathbb{R} \to 0 \leq x \to 0 \leq y \to x \in \mathbb{R}) \to (x \in \mathbb{R} \to y \in \mathbb{R} \to 0 \leq x \to 0 \leq y \to y \in \mathbb{R}) \to \varphi_{5}$
\end{abbrevlist}
\end{fullproof}

% 自動生成 (prover/build.js): AddLtLt  ノード数 29
\begin{fullproof}{[E7]}{$\forall x\,\forall y\,\forall u\,\forall v\,(x \in \mathbb{R} \to y \in \mathbb{R} \to u \in \mathbb{R} \to v \in \mathbb{R} \to x < y \to u < v \to x + u < y + v)$}{29}\phantomsection\label{tree:AddLtLt}
\noindent$\Pi_{1}$:
\begin{ptree}
\AXC{$\forall x\,\forall y\,\forall u\,\forall v\,(\varphi_{8})$\lbl{[F20]}}
\AXC{$\forall x\,\forall y\,\forall u\,\forall v\,(\varphi_{8}) \to \forall y\,\forall u\,\forall v\,(\varphi_{8})$\lbl{A2}}
\BIC{$\forall y\,\forall u\,\forall v\,(\varphi_{8})$}
\AXC{$\forall y\,\forall u\,\forall v\,(\varphi_{8}) \to \forall u\,\forall v\,(\varphi_{8})$\lbl{A2}}
\BIC{$\forall u\,\forall v\,(\varphi_{8})$}
\end{ptree}
\noindent$\Pi_{2}$:
\begin{ptree}
\AXC{$\varphi_{7}$\lbl{A1}}
\AXC{$\stackrel{\Pi_{1}}{\vdots}$}\noLine
\UIC{$\forall u\,\forall v\,(\varphi_{8})$}
\AXC{$\forall u\,\forall v\,(\varphi_{8}) \to \forall v\,(\varphi_{8})$\lbl{A2}}
\BIC{$\forall v\,(\varphi_{8})$}
\AXC{$\forall v\,(\varphi_{8}) \to \varphi_{8}$\lbl{A2}}
\BIC{$\varphi_{8}$}
\AXC{$\varphi_{8} \to \varphi_{7} \to \varphi_{6} \to \varphi_{5} \to \varphi_{4} \to \varphi_{3} \to \varphi_{1} \to \varphi_{9}$\lbl{A1}}
\BIC{$\varphi_{7} \to \varphi_{6} \to \varphi_{5} \to \varphi_{4} \to \varphi_{3} \to \varphi_{1} \to \varphi_{9}$}
\BIC{$\varphi_{6} \to \varphi_{5} \to \varphi_{4} \to \varphi_{3} \to \varphi_{1} \to \varphi_{9}$}
\end{ptree}
\noindent 主木:
\begin{ptree}
\AXC{$\varphi_{2}$\lbl{A1}}
\AXC{$\varphi_{2} \to \varphi_{1}$\lbl{A1}}
\BIC{$\varphi_{1}$}
\AXC{$\varphi_{3}$\lbl{A1}}
\AXC{$\varphi_{4}$\lbl{A1}}
\AXC{$\varphi_{5}$\lbl{A1}}
\AXC{$\varphi_{6}$\lbl{A1}}
\AXC{$\stackrel{\Pi_{2}}{\vdots}$}\noLine
\UIC{$\varphi_{6} \to \varphi_{5} \to \varphi_{4} \to \varphi_{3} \to \varphi_{1} \to \varphi_{9}$}
\BIC{$\varphi_{5} \to \varphi_{4} \to \varphi_{3} \to \varphi_{1} \to \varphi_{9}$}
\BIC{$\varphi_{4} \to \varphi_{3} \to \varphi_{1} \to \varphi_{9}$}
\BIC{$\varphi_{3} \to \varphi_{1} \to \varphi_{9}$}
\BIC{$\varphi_{1} \to \varphi_{9}$}
\BIC{$\varphi_{9}$}
\UIC{$\forall v\,(\varphi_{9})$}
\UIC{$\forall u\,\forall v\,(\varphi_{9})$}
\UIC{$\forall y\,\forall u\,\forall v\,(\varphi_{9})$}
\UIC{$\forall x\,\forall y\,\forall u\,\forall v\,(\varphi_{9})$}
\end{ptree}
\begin{abbrevlist}
\item $\varphi_{1} :\equiv x \in \mathbb{R} \to y \in \mathbb{R} \to u \in \mathbb{R} \to v \in \mathbb{R} \to x < y \to u < v \to u \leq v$
\item $\varphi_{2} :\equiv x \in \mathbb{R} \to y \in \mathbb{R} \to u \in \mathbb{R} \to v \in \mathbb{R} \to x < y \to u < v \to u < v$
\item $\varphi_{3} :\equiv x \in \mathbb{R} \to y \in \mathbb{R} \to u \in \mathbb{R} \to v \in \mathbb{R} \to x < y \to u < v \to x < y$
\item $\varphi_{4} :\equiv x \in \mathbb{R} \to y \in \mathbb{R} \to u \in \mathbb{R} \to v \in \mathbb{R} \to x < y \to u < v \to v \in \mathbb{R}$
\item $\varphi_{5} :\equiv x \in \mathbb{R} \to y \in \mathbb{R} \to u \in \mathbb{R} \to v \in \mathbb{R} \to x < y \to u < v \to u \in \mathbb{R}$
\item $\varphi_{6} :\equiv x \in \mathbb{R} \to y \in \mathbb{R} \to u \in \mathbb{R} \to v \in \mathbb{R} \to x < y \to u < v \to y \in \mathbb{R}$
\item $\varphi_{7} :\equiv x \in \mathbb{R} \to y \in \mathbb{R} \to u \in \mathbb{R} \to v \in \mathbb{R} \to x < y \to u < v \to x \in \mathbb{R}$
\item $\varphi_{8} :\equiv x \in \mathbb{R} \to y \in \mathbb{R} \to u \in \mathbb{R} \to v \in \mathbb{R} \to x < y \to u \leq v \to x + u < y + v$
\item $\varphi_{9} :\equiv x \in \mathbb{R} \to y \in \mathbb{R} \to u \in \mathbb{R} \to v \in \mathbb{R} \to x < y \to u < v \to x + u < y + v$
\end{abbrevlist}
\end{fullproof}

% 自動生成 (prover/build.js): MulLtPosL  ノード数 60
\begin{fullproof}{[E8]}{$\forall x\,\forall y\,\forall z\,(x \in \mathbb{R} \to y \in \mathbb{R} \to z \in \mathbb{R} \to 0 < z \to x < y \to z \cdot x < z \cdot y)$}{60}\phantomsection\label{tree:MulLtPosL}
\noindent$\Pi_{1}$:
\begin{ptree}
\AXC{$\forall x\,\forall y\,(\varphi_{2})$\lbl{[F32]}}
\AXC{$\forall x\,\forall y\,(\varphi_{2}) \to \forall y\,(\varphi_{2})$\lbl{A2}}
\BIC{$\forall y\,(\varphi_{2})$}
\AXC{$\forall y\,(\varphi_{2}) \to \varphi_{2}$\lbl{A2}}
\BIC{$\varphi_{2}$}
\AXC{$\varphi_{18}$\lbl{A2}}
\BIC{$x \in \mathbb{R} \to y \in \mathbb{R} \to z \in \mathbb{R} \to 0 < z \to x < y \to x \cdot z < y \cdot z$}
\end{ptree}
\noindent$\Pi_{2}$:
\begin{ptree}
\AXC{$\forall x\,\forall y\,(x \in \mathbb{R} \to y \in \mathbb{R} \to x \cdot y = y \cdot x)$\lbl{(R7)}}
\AXC{$\varphi_{7}$\lbl{A2}}
\BIC{$\forall y\,(x \in \mathbb{R} \to y \in \mathbb{R} \to x \cdot y = y \cdot x)$}
\AXC{$\varphi_{4}$\lbl{A2}}
\BIC{$x \in \mathbb{R} \to z \in \mathbb{R} \to x \cdot z = z \cdot x$}
\AXC{$(x \in \mathbb{R} \to z \in \mathbb{R} \to x \cdot z = z \cdot x) \to \varphi_{21}$\lbl{A1}}
\BIC{$\varphi_{21}$}
\end{ptree}
\noindent$\Pi_{3}$:
\begin{ptree}
\AXC{$x \in \mathbb{R} \to y \in \mathbb{R} \to z \in \mathbb{R} \to 0 < z \to x < y \to z \in \mathbb{R}$\lbl{A1}}
\AXC{$x \in \mathbb{R} \to y \in \mathbb{R} \to z \in \mathbb{R} \to 0 < z \to x < y \to x \in \mathbb{R}$\lbl{A1}}
\AXC{$\stackrel{\Pi_{2}}{\vdots}$}\noLine
\UIC{$\varphi_{21}$}
\BIC{$\varphi_{16}$}
\BIC{$x \in \mathbb{R} \to y \in \mathbb{R} \to z \in \mathbb{R} \to 0 < z \to x < y \to x \cdot z = z \cdot x$}
\end{ptree}
\noindent$\Pi_{4}$:
\begin{ptree}
\AXC{$\stackrel{\Pi_{1}}{\vdots}$}\noLine
\UIC{$x \in \mathbb{R} \to y \in \mathbb{R} \to z \in \mathbb{R} \to 0 < z \to x < y \to x \cdot z < y \cdot z$}
\AXC{$\stackrel{\Pi_{3}}{\vdots}$}\noLine
\UIC{$x \in \mathbb{R} \to y \in \mathbb{R} \to z \in \mathbb{R} \to 0 < z \to x < y \to x \cdot z = z \cdot x$}
\AXC{$\varphi_{23}$\lbl{A1}}
\BIC{$\varphi_{19}$}
\BIC{$\varphi_{14}$}
\end{ptree}
\noindent$\Pi_{5}$:
\begin{ptree}
\AXC{$z_{1} = w \to z_{1} < y \cdot z \to w < y \cdot z$\lbl{A6}}
\UIC{$\forall w\,(z_{1} = w \to z_{1} < y \cdot z \to w < y \cdot z)$}
\UIC{$\forall z_{1}\,\forall w\,(z_{1} = w \to z_{1} < y \cdot z \to w < y \cdot z)$}
\AXC{$\varphi_{6}$\lbl{A2}}
\BIC{$\forall w\,(x \cdot z = w \to x \cdot z < y \cdot z \to w < y \cdot z)$}
\AXC{$\varphi_{10}$\lbl{A2}}
\BIC{$x \cdot z = z \cdot x \to x \cdot z < y \cdot z \to z \cdot x < y \cdot z$}
\AXC{$\stackrel{\Pi_{4}}{\vdots}$}\noLine
\UIC{$\varphi_{14}$}
\BIC{$x \in \mathbb{R} \to y \in \mathbb{R} \to z \in \mathbb{R} \to 0 < z \to x < y \to z \cdot x < y \cdot z$}
\end{ptree}
\noindent$\Pi_{6}$:
\begin{ptree}
\AXC{$\forall x\,\forall y\,(x \in \mathbb{R} \to y \in \mathbb{R} \to x \cdot y = y \cdot x)$\lbl{(R7)}}
\AXC{$\varphi_{11}$\lbl{A2}}
\BIC{$\forall y\,(v_{449} \in \mathbb{R} \to y \in \mathbb{R} \to v_{449} \cdot y = y \cdot v_{449})$}
\AXC{$\varphi_{12}$\lbl{A2}}
\BIC{$v_{449} \in \mathbb{R} \to v_{450} \in \mathbb{R} \to v_{449} \cdot v_{450} = v_{450} \cdot v_{449}$}
\UIC{$\varphi_{1}$}
\UIC{$\forall v_{449}\,(\varphi_{1})$}
\AXC{$\varphi_{13}$\lbl{A2}}
\BIC{$\forall v_{450}\,(y \in \mathbb{R} \to v_{450} \in \mathbb{R} \to y \cdot v_{450} = v_{450} \cdot y)$}
\AXC{$\varphi_{8}$\lbl{A2}}
\BIC{$y \in \mathbb{R} \to z \in \mathbb{R} \to y \cdot z = z \cdot y$}
\end{ptree}
\noindent$\Pi_{7}$:
\begin{ptree}
\AXC{$\stackrel{\Pi_{6}}{\vdots}$}\noLine
\UIC{$y \in \mathbb{R} \to z \in \mathbb{R} \to y \cdot z = z \cdot y$}
\AXC{$(y \in \mathbb{R} \to z \in \mathbb{R} \to y \cdot z = z \cdot y) \to \varphi_{22}$\lbl{A1}}
\BIC{$\varphi_{22}$}
\end{ptree}
\noindent$\Pi_{8}$:
\begin{ptree}
\AXC{$x \in \mathbb{R} \to y \in \mathbb{R} \to z \in \mathbb{R} \to 0 < z \to x < y \to z \in \mathbb{R}$\lbl{A1}}
\AXC{$x \in \mathbb{R} \to y \in \mathbb{R} \to z \in \mathbb{R} \to 0 < z \to x < y \to y \in \mathbb{R}$\lbl{A1}}
\AXC{$\stackrel{\Pi_{7}}{\vdots}$}\noLine
\UIC{$\varphi_{22}$}
\BIC{$\varphi_{17}$}
\BIC{$x \in \mathbb{R} \to y \in \mathbb{R} \to z \in \mathbb{R} \to 0 < z \to x < y \to y \cdot z = z \cdot y$}
\end{ptree}
\noindent$\Pi_{9}$:
\begin{ptree}
\AXC{$\stackrel{\Pi_{5}}{\vdots}$}\noLine
\UIC{$x \in \mathbb{R} \to y \in \mathbb{R} \to z \in \mathbb{R} \to 0 < z \to x < y \to z \cdot x < y \cdot z$}
\AXC{$\stackrel{\Pi_{8}}{\vdots}$}\noLine
\UIC{$x \in \mathbb{R} \to y \in \mathbb{R} \to z \in \mathbb{R} \to 0 < z \to x < y \to y \cdot z = z \cdot y$}
\AXC{$\varphi_{24}$\lbl{A1}}
\BIC{$\varphi_{20}$}
\BIC{$\varphi_{15}$}
\end{ptree}
\noindent 主木:
\begin{ptree}
\AXC{$z_{1} = w \to z \cdot x < z_{1} \to z \cdot x < w$\lbl{A6}}
\UIC{$\forall w\,(z_{1} = w \to z \cdot x < z_{1} \to z \cdot x < w)$}
\UIC{$\forall z_{1}\,\forall w\,(z_{1} = w \to z \cdot x < z_{1} \to z \cdot x < w)$}
\AXC{$\varphi_{5}$\lbl{A2}}
\BIC{$\forall w\,(y \cdot z = w \to z \cdot x < y \cdot z \to z \cdot x < w)$}
\AXC{$\varphi_{9}$\lbl{A2}}
\BIC{$y \cdot z = z \cdot y \to z \cdot x < y \cdot z \to z \cdot x < z \cdot y$}
\AXC{$\stackrel{\Pi_{9}}{\vdots}$}\noLine
\UIC{$\varphi_{15}$}
\BIC{$x \in \mathbb{R} \to y \in \mathbb{R} \to z \in \mathbb{R} \to 0 < z \to x < y \to z \cdot x < z \cdot y$}
\UIC{$\varphi_{3}$}
\UIC{$\forall y\,(\varphi_{3})$}
\UIC{$\forall x\,\forall y\,(\varphi_{3})$}
\end{ptree}
\begin{abbrevlist}
\item $\varphi_{1} :\equiv \forall v_{450}\,(v_{449} \in \mathbb{R} \to v_{450} \in \mathbb{R} \to v_{449} \cdot v_{450} = v_{450} \cdot v_{449})$
\item $\varphi_{2} :\equiv \forall z\,(x \in \mathbb{R} \to y \in \mathbb{R} \to z \in \mathbb{R} \to 0 < z \to x < y \to x \cdot z < y \cdot z)$
\item $\varphi_{3} :\equiv \forall z\,(x \in \mathbb{R} \to y \in \mathbb{R} \to z \in \mathbb{R} \to 0 < z \to x < y \to z \cdot x < z \cdot y)$
\item $\varphi_{4} :\equiv \forall y\,(x \in \mathbb{R} \to y \in \mathbb{R} \to x \cdot y = y \cdot x) \to x \in \mathbb{R} \to z \in \mathbb{R} \to x \cdot z = z \cdot x$
\item $\varphi_{5} :\equiv \forall z_{1}\,\forall w\,(z_{1} = w \to z \cdot x < z_{1} \to z \cdot x < w) \to \forall w\,(y \cdot z = w \to z \cdot x < y \cdot z \to z \cdot x < w)$
\item $\varphi_{6} :\equiv \forall z_{1}\,\forall w\,(z_{1} = w \to z_{1} < y \cdot z \to w < y \cdot z) \to \forall w\,(x \cdot z = w \to x \cdot z < y \cdot z \to w < y \cdot z)$
\item $\varphi_{7} :\equiv \forall x\,\forall y\,(x \in \mathbb{R} \to y \in \mathbb{R} \to x \cdot y = y \cdot x) \to \forall y\,(x \in \mathbb{R} \to y \in \mathbb{R} \to x \cdot y = y \cdot x)$
\item $\varphi_{8} :\equiv \forall v_{450}\,(y \in \mathbb{R} \to v_{450} \in \mathbb{R} \to y \cdot v_{450} = v_{450} \cdot y) \to y \in \mathbb{R} \to z \in \mathbb{R} \to y \cdot z = z \cdot y$
\item $\varphi_{9} :\equiv \forall w\,(y \cdot z = w \to z \cdot x < y \cdot z \to z \cdot x < w) \to y \cdot z = z \cdot y \to z \cdot x < y \cdot z \to z \cdot x < z \cdot y$
\item $\varphi_{10} :\equiv \forall w\,(x \cdot z = w \to x \cdot z < y \cdot z \to w < y \cdot z) \to x \cdot z = z \cdot x \to x \cdot z < y \cdot z \to z \cdot x < y \cdot z$
\item $\varphi_{11} :\equiv \forall x\,\forall y\,(x \in \mathbb{R} \to y \in \mathbb{R} \to x \cdot y = y \cdot x) \to \forall y\,(v_{449} \in \mathbb{R} \to y \in \mathbb{R} \to v_{449} \cdot y = y \cdot v_{449})$
\item $\varphi_{12} :\equiv \forall y\,(v_{449} \in \mathbb{R} \to y \in \mathbb{R} \to v_{449} \cdot y = y \cdot v_{449}) \to v_{449} \in \mathbb{R} \to v_{450} \in \mathbb{R} \to v_{449} \cdot v_{450} = v_{450} \cdot v_{449}$
\item $\varphi_{13} :\equiv \forall v_{449}\,(\varphi_{1}) \to \forall v_{450}\,(y \in \mathbb{R} \to v_{450} \in \mathbb{R} \to y \cdot v_{450} = v_{450} \cdot y)$
\item $\varphi_{14} :\equiv (x \cdot z = z \cdot x \to x \cdot z < y \cdot z \to z \cdot x < y \cdot z) \to x \in \mathbb{R} \to y \in \mathbb{R} \to z \in \mathbb{R} \to 0 < z \to x < y \to z \cdot x < y \cdot z$
\item $\varphi_{15} :\equiv (y \cdot z = z \cdot y \to z \cdot x < y \cdot z \to z \cdot x < z \cdot y) \to x \in \mathbb{R} \to y \in \mathbb{R} \to z \in \mathbb{R} \to 0 < z \to x < y \to z \cdot x < z \cdot y$
\item $\varphi_{16} :\equiv (x \in \mathbb{R} \to y \in \mathbb{R} \to z \in \mathbb{R} \to 0 < z \to x < y \to z \in \mathbb{R}) \to x \in \mathbb{R} \to y \in \mathbb{R} \to z \in \mathbb{R} \to 0 < z \to x < y \to x \cdot z = z \cdot x$
\item $\varphi_{17} :\equiv (x \in \mathbb{R} \to y \in \mathbb{R} \to z \in \mathbb{R} \to 0 < z \to x < y \to z \in \mathbb{R}) \to x \in \mathbb{R} \to y \in \mathbb{R} \to z \in \mathbb{R} \to 0 < z \to x < y \to y \cdot z = z \cdot y$
\item $\varphi_{18} :\equiv \varphi_{2} \to x \in \mathbb{R} \to y \in \mathbb{R} \to z \in \mathbb{R} \to 0 < z \to x < y \to x \cdot z < y \cdot z$
\item $\varphi_{19} :\equiv (x \in \mathbb{R} \to y \in \mathbb{R} \to z \in \mathbb{R} \to 0 < z \to x < y \to x \cdot z < y \cdot z) \to \varphi_{14}$
\item $\varphi_{20} :\equiv (x \in \mathbb{R} \to y \in \mathbb{R} \to z \in \mathbb{R} \to 0 < z \to x < y \to z \cdot x < y \cdot z) \to \varphi_{15}$
\item $\varphi_{21} :\equiv (x \in \mathbb{R} \to y \in \mathbb{R} \to z \in \mathbb{R} \to 0 < z \to x < y \to x \in \mathbb{R}) \to \varphi_{16}$
\item $\varphi_{22} :\equiv (x \in \mathbb{R} \to y \in \mathbb{R} \to z \in \mathbb{R} \to 0 < z \to x < y \to y \in \mathbb{R}) \to \varphi_{17}$
\item $\varphi_{23} :\equiv (x \in \mathbb{R} \to y \in \mathbb{R} \to z \in \mathbb{R} \to 0 < z \to x < y \to x \cdot z = z \cdot x) \to \varphi_{19}$
\item $\varphi_{24} :\equiv (x \in \mathbb{R} \to y \in \mathbb{R} \to z \in \mathbb{R} \to 0 < z \to x < y \to y \cdot z = z \cdot y) \to \varphi_{20}$
\end{abbrevlist}
\end{fullproof}

% 自動生成 (prover/build.js): MulLeNonnegL  ノード数 76
\begin{fullproof}{[E9]}{$\forall x\,\forall y\,\forall z\,(x \in \mathbb{R} \to y \in \mathbb{R} \to z \in \mathbb{R} \to 0 \leq z \to x \leq y \to z \cdot x \leq z \cdot y)$}{76}\phantomsection\label{tree:MulLeNonnegL}
\noindent$\Pi_{1}$:
\begin{ptree}
\AXC{$\forall x\,\forall y\,(\varphi_{2})$\lbl{[F33]}}
\AXC{$\forall x\,\forall y\,(\varphi_{2}) \to \forall y\,(\varphi_{2})$\lbl{A2}}
\BIC{$\forall y\,(\varphi_{2})$}
\AXC{$\forall y\,(\varphi_{2}) \to \varphi_{2}$\lbl{A2}}
\BIC{$\varphi_{2}$}
\AXC{$\varphi_{22}$\lbl{A2}}
\BIC{$x \in \mathbb{R} \to y \in \mathbb{R} \to z \in \mathbb{R} \to 0 \leq z \to x \leq y \to x \cdot z \leq y \cdot z$}
\end{ptree}
\noindent$\Pi_{2}$:
\begin{ptree}
\AXC{$\forall x\,\forall y\,(x \in \mathbb{R} \to y \in \mathbb{R} \to x \cdot y = y \cdot x)$\lbl{(R7)}}
\AXC{$\varphi_{11}$\lbl{A2}}
\BIC{$\forall y\,(x \in \mathbb{R} \to y \in \mathbb{R} \to x \cdot y = y \cdot x)$}
\AXC{$\varphi_{6}$\lbl{A2}}
\BIC{$x \in \mathbb{R} \to z \in \mathbb{R} \to x \cdot z = z \cdot x$}
\AXC{$(x \in \mathbb{R} \to z \in \mathbb{R} \to x \cdot z = z \cdot x) \to \varphi_{23}$\lbl{A1}}
\BIC{$\varphi_{23}$}
\end{ptree}
\noindent$\Pi_{3}$:
\begin{ptree}
\AXC{$x \in \mathbb{R} \to y \in \mathbb{R} \to z \in \mathbb{R} \to 0 \leq z \to x \leq y \to z \in \mathbb{R}$\lbl{A1}}
\AXC{$x \in \mathbb{R} \to y \in \mathbb{R} \to z \in \mathbb{R} \to 0 \leq z \to x \leq y \to x \in \mathbb{R}$\lbl{A1}}
\AXC{$\stackrel{\Pi_{2}}{\vdots}$}\noLine
\UIC{$\varphi_{23}$}
\BIC{$\varphi_{20}$}
\BIC{$x \in \mathbb{R} \to y \in \mathbb{R} \to z \in \mathbb{R} \to 0 \leq z \to x \leq y \to x \cdot z = z \cdot x$}
\end{ptree}
\noindent$\Pi_{4}$:
\begin{ptree}
\AXC{$\stackrel{\Pi_{1}}{\vdots}$}\noLine
\UIC{$x \in \mathbb{R} \to y \in \mathbb{R} \to z \in \mathbb{R} \to 0 \leq z \to x \leq y \to x \cdot z \leq y \cdot z$}
\AXC{$\stackrel{\Pi_{3}}{\vdots}$}\noLine
\UIC{$x \in \mathbb{R} \to y \in \mathbb{R} \to z \in \mathbb{R} \to 0 \leq z \to x \leq y \to x \cdot z = z \cdot x$}
\AXC{$\varphi_{27}$\lbl{A1}}
\BIC{$\varphi_{25}$}
\BIC{$(x \cdot z = z \cdot x \to x \cdot z < y \cdot z \to z \cdot x < y \cdot z) \to \varphi_{17}$}
\end{ptree}
\noindent$\Pi_{5}$:
\begin{ptree}
\AXC{$z_{1} = w \to z_{1} < y \cdot z \to w < y \cdot z$\lbl{A6}}
\UIC{$\forall w\,(z_{1} = w \to z_{1} < y \cdot z \to w < y \cdot z)$}
\UIC{$\forall z_{1}\,\forall w\,(z_{1} = w \to z_{1} < y \cdot z \to w < y \cdot z)$}
\AXC{$\varphi_{10}$\lbl{A2}}
\BIC{$\forall w\,(x \cdot z = w \to x \cdot z < y \cdot z \to w < y \cdot z)$}
\AXC{$\varphi_{14}$\lbl{A2}}
\BIC{$x \cdot z = z \cdot x \to x \cdot z < y \cdot z \to z \cdot x < y \cdot z$}
\AXC{$\stackrel{\Pi_{4}}{\vdots}$}\noLine
\UIC{$(x \cdot z = z \cdot x \to x \cdot z < y \cdot z \to z \cdot x < y \cdot z) \to \varphi_{17}$}
\BIC{$\varphi_{17}$}
\end{ptree}
\noindent$\Pi_{6}$:
\begin{ptree}
\AXC{$z_{1} = w \to z_{1} = y \cdot z \to w = y \cdot z$\lbl{A6}}
\UIC{$\forall w\,(z_{1} = w \to z_{1} = y \cdot z \to w = y \cdot z)$}
\UIC{$\forall z_{1}\,\forall w\,(z_{1} = w \to z_{1} = y \cdot z \to w = y \cdot z)$}
\AXC{$\varphi_{5}$\lbl{A2}}
\BIC{$\forall w\,(x \cdot z = w \to x \cdot z = y \cdot z \to w = y \cdot z)$}
\AXC{$\varphi_{8}$\lbl{A2}}
\BIC{$x \cdot z = z \cdot x \to x \cdot z = y \cdot z \to z \cdot x = y \cdot z$}
\AXC{$\stackrel{\Pi_{5}}{\vdots}$}\noLine
\UIC{$\varphi_{17}$}
\BIC{$x \in \mathbb{R} \to y \in \mathbb{R} \to z \in \mathbb{R} \to 0 \leq z \to x \leq y \to z \cdot x \leq y \cdot z$}
\end{ptree}
\noindent$\Pi_{7}$:
\begin{ptree}
\AXC{$\forall x\,\forall y\,(x \in \mathbb{R} \to y \in \mathbb{R} \to x \cdot y = y \cdot x)$\lbl{(R7)}}
\AXC{$\varphi_{15}$\lbl{A2}}
\BIC{$\forall y\,(v_{451} \in \mathbb{R} \to y \in \mathbb{R} \to v_{451} \cdot y = y \cdot v_{451})$}
\AXC{$\varphi_{16}$\lbl{A2}}
\BIC{$v_{451} \in \mathbb{R} \to v_{452} \in \mathbb{R} \to v_{451} \cdot v_{452} = v_{452} \cdot v_{451}$}
\UIC{$\varphi_{1}$}
\UIC{$\forall v_{451}\,(\varphi_{1})$}
\AXC{$\varphi_{19}$\lbl{A2}}
\BIC{$\forall v_{452}\,(y \in \mathbb{R} \to v_{452} \in \mathbb{R} \to y \cdot v_{452} = v_{452} \cdot y)$}
\AXC{$\varphi_{12}$\lbl{A2}}
\BIC{$y \in \mathbb{R} \to z \in \mathbb{R} \to y \cdot z = z \cdot y$}
\end{ptree}
\noindent$\Pi_{8}$:
\begin{ptree}
\AXC{$\stackrel{\Pi_{7}}{\vdots}$}\noLine
\UIC{$y \in \mathbb{R} \to z \in \mathbb{R} \to y \cdot z = z \cdot y$}
\AXC{$(y \in \mathbb{R} \to z \in \mathbb{R} \to y \cdot z = z \cdot y) \to \varphi_{24}$\lbl{A1}}
\BIC{$\varphi_{24}$}
\end{ptree}
\noindent$\Pi_{9}$:
\begin{ptree}
\AXC{$x \in \mathbb{R} \to y \in \mathbb{R} \to z \in \mathbb{R} \to 0 \leq z \to x \leq y \to z \in \mathbb{R}$\lbl{A1}}
\AXC{$x \in \mathbb{R} \to y \in \mathbb{R} \to z \in \mathbb{R} \to 0 \leq z \to x \leq y \to y \in \mathbb{R}$\lbl{A1}}
\AXC{$\stackrel{\Pi_{8}}{\vdots}$}\noLine
\UIC{$\varphi_{24}$}
\BIC{$\varphi_{21}$}
\BIC{$x \in \mathbb{R} \to y \in \mathbb{R} \to z \in \mathbb{R} \to 0 \leq z \to x \leq y \to y \cdot z = z \cdot y$}
\end{ptree}
\noindent$\Pi_{10}$:
\begin{ptree}
\AXC{$\stackrel{\Pi_{6}}{\vdots}$}\noLine
\UIC{$x \in \mathbb{R} \to y \in \mathbb{R} \to z \in \mathbb{R} \to 0 \leq z \to x \leq y \to z \cdot x \leq y \cdot z$}
\AXC{$\stackrel{\Pi_{9}}{\vdots}$}\noLine
\UIC{$x \in \mathbb{R} \to y \in \mathbb{R} \to z \in \mathbb{R} \to 0 \leq z \to x \leq y \to y \cdot z = z \cdot y$}
\AXC{$\varphi_{28}$\lbl{A1}}
\BIC{$\varphi_{26}$}
\BIC{$(y \cdot z = z \cdot y \to z \cdot x < y \cdot z \to z \cdot x < z \cdot y) \to \varphi_{18}$}
\end{ptree}
\noindent$\Pi_{11}$:
\begin{ptree}
\AXC{$z_{1} = w \to z \cdot x < z_{1} \to z \cdot x < w$\lbl{A6}}
\UIC{$\forall w\,(z_{1} = w \to z \cdot x < z_{1} \to z \cdot x < w)$}
\UIC{$\forall z_{1}\,\forall w\,(z_{1} = w \to z \cdot x < z_{1} \to z \cdot x < w)$}
\AXC{$\varphi_{9}$\lbl{A2}}
\BIC{$\forall w\,(y \cdot z = w \to z \cdot x < y \cdot z \to z \cdot x < w)$}
\AXC{$\varphi_{13}$\lbl{A2}}
\BIC{$y \cdot z = z \cdot y \to z \cdot x < y \cdot z \to z \cdot x < z \cdot y$}
\AXC{$\stackrel{\Pi_{10}}{\vdots}$}\noLine
\UIC{$(y \cdot z = z \cdot y \to z \cdot x < y \cdot z \to z \cdot x < z \cdot y) \to \varphi_{18}$}
\BIC{$\varphi_{18}$}
\end{ptree}
\noindent 主木:
\begin{ptree}
\AXC{$z_{1} = w \to z \cdot x = z_{1} \to z \cdot x = w$\lbl{A6}}
\UIC{$\forall w\,(z_{1} = w \to z \cdot x = z_{1} \to z \cdot x = w)$}
\UIC{$\forall z_{1}\,\forall w\,(z_{1} = w \to z \cdot x = z_{1} \to z \cdot x = w)$}
\AXC{$\varphi_{4}$\lbl{A2}}
\BIC{$\forall w\,(y \cdot z = w \to z \cdot x = y \cdot z \to z \cdot x = w)$}
\AXC{$\varphi_{7}$\lbl{A2}}
\BIC{$y \cdot z = z \cdot y \to z \cdot x = y \cdot z \to z \cdot x = z \cdot y$}
\AXC{$\stackrel{\Pi_{11}}{\vdots}$}\noLine
\UIC{$\varphi_{18}$}
\BIC{$x \in \mathbb{R} \to y \in \mathbb{R} \to z \in \mathbb{R} \to 0 \leq z \to x \leq y \to z \cdot x \leq z \cdot y$}
\UIC{$\varphi_{3}$}
\UIC{$\forall y\,(\varphi_{3})$}
\UIC{$\forall x\,\forall y\,(\varphi_{3})$}
\end{ptree}
\begin{abbrevlist}
\item $\varphi_{1} :\equiv \forall v_{452}\,(v_{451} \in \mathbb{R} \to v_{452} \in \mathbb{R} \to v_{451} \cdot v_{452} = v_{452} \cdot v_{451})$
\item $\varphi_{2} :\equiv \forall z\,(x \in \mathbb{R} \to y \in \mathbb{R} \to z \in \mathbb{R} \to 0 \leq z \to x \leq y \to x \cdot z \leq y \cdot z)$
\item $\varphi_{3} :\equiv \forall z\,(x \in \mathbb{R} \to y \in \mathbb{R} \to z \in \mathbb{R} \to 0 \leq z \to x \leq y \to z \cdot x \leq z \cdot y)$
\item $\varphi_{4} :\equiv \forall z_{1}\,\forall w\,(z_{1} = w \to z \cdot x = z_{1} \to z \cdot x = w) \to \forall w\,(y \cdot z = w \to z \cdot x = y \cdot z \to z \cdot x = w)$
\item $\varphi_{5} :\equiv \forall z_{1}\,\forall w\,(z_{1} = w \to z_{1} = y \cdot z \to w = y \cdot z) \to \forall w\,(x \cdot z = w \to x \cdot z = y \cdot z \to w = y \cdot z)$
\item $\varphi_{6} :\equiv \forall y\,(x \in \mathbb{R} \to y \in \mathbb{R} \to x \cdot y = y \cdot x) \to x \in \mathbb{R} \to z \in \mathbb{R} \to x \cdot z = z \cdot x$
\item $\varphi_{7} :\equiv \forall w\,(y \cdot z = w \to z \cdot x = y \cdot z \to z \cdot x = w) \to y \cdot z = z \cdot y \to z \cdot x = y \cdot z \to z \cdot x = z \cdot y$
\item $\varphi_{8} :\equiv \forall w\,(x \cdot z = w \to x \cdot z = y \cdot z \to w = y \cdot z) \to x \cdot z = z \cdot x \to x \cdot z = y \cdot z \to z \cdot x = y \cdot z$
\item $\varphi_{9} :\equiv \forall z_{1}\,\forall w\,(z_{1} = w \to z \cdot x < z_{1} \to z \cdot x < w) \to \forall w\,(y \cdot z = w \to z \cdot x < y \cdot z \to z \cdot x < w)$
\item $\varphi_{10} :\equiv \forall z_{1}\,\forall w\,(z_{1} = w \to z_{1} < y \cdot z \to w < y \cdot z) \to \forall w\,(x \cdot z = w \to x \cdot z < y \cdot z \to w < y \cdot z)$
\item $\varphi_{11} :\equiv \forall x\,\forall y\,(x \in \mathbb{R} \to y \in \mathbb{R} \to x \cdot y = y \cdot x) \to \forall y\,(x \in \mathbb{R} \to y \in \mathbb{R} \to x \cdot y = y \cdot x)$
\item $\varphi_{12} :\equiv \forall v_{452}\,(y \in \mathbb{R} \to v_{452} \in \mathbb{R} \to y \cdot v_{452} = v_{452} \cdot y) \to y \in \mathbb{R} \to z \in \mathbb{R} \to y \cdot z = z \cdot y$
\item $\varphi_{13} :\equiv \forall w\,(y \cdot z = w \to z \cdot x < y \cdot z \to z \cdot x < w) \to y \cdot z = z \cdot y \to z \cdot x < y \cdot z \to z \cdot x < z \cdot y$
\item $\varphi_{14} :\equiv \forall w\,(x \cdot z = w \to x \cdot z < y \cdot z \to w < y \cdot z) \to x \cdot z = z \cdot x \to x \cdot z < y \cdot z \to z \cdot x < y \cdot z$
\item $\varphi_{15} :\equiv \forall x\,\forall y\,(x \in \mathbb{R} \to y \in \mathbb{R} \to x \cdot y = y \cdot x) \to \forall y\,(v_{451} \in \mathbb{R} \to y \in \mathbb{R} \to v_{451} \cdot y = y \cdot v_{451})$
\item $\varphi_{16} :\equiv \forall y\,(v_{451} \in \mathbb{R} \to y \in \mathbb{R} \to v_{451} \cdot y = y \cdot v_{451}) \to v_{451} \in \mathbb{R} \to v_{452} \in \mathbb{R} \to v_{451} \cdot v_{452} = v_{452} \cdot v_{451}$
\item $\varphi_{17} :\equiv (x \cdot z = z \cdot x \to x \cdot z = y \cdot z \to z \cdot x = y \cdot z) \to x \in \mathbb{R} \to y \in \mathbb{R} \to z \in \mathbb{R} \to 0 \leq z \to x \leq y \to z \cdot x \leq y \cdot z$
\item $\varphi_{18} :\equiv (y \cdot z = z \cdot y \to z \cdot x = y \cdot z \to z \cdot x = z \cdot y) \to x \in \mathbb{R} \to y \in \mathbb{R} \to z \in \mathbb{R} \to 0 \leq z \to x \leq y \to z \cdot x \leq z \cdot y$
\item $\varphi_{19} :\equiv \forall v_{451}\,(\varphi_{1}) \to \forall v_{452}\,(y \in \mathbb{R} \to v_{452} \in \mathbb{R} \to y \cdot v_{452} = v_{452} \cdot y)$
\item $\varphi_{20} :\equiv (x \in \mathbb{R} \to y \in \mathbb{R} \to z \in \mathbb{R} \to 0 \leq z \to x \leq y \to z \in \mathbb{R}) \to x \in \mathbb{R} \to y \in \mathbb{R} \to z \in \mathbb{R} \to 0 \leq z \to x \leq y \to x \cdot z = z \cdot x$
\item $\varphi_{21} :\equiv (x \in \mathbb{R} \to y \in \mathbb{R} \to z \in \mathbb{R} \to 0 \leq z \to x \leq y \to z \in \mathbb{R}) \to x \in \mathbb{R} \to y \in \mathbb{R} \to z \in \mathbb{R} \to 0 \leq z \to x \leq y \to y \cdot z = z \cdot y$
\item $\varphi_{22} :\equiv \varphi_{2} \to x \in \mathbb{R} \to y \in \mathbb{R} \to z \in \mathbb{R} \to 0 \leq z \to x \leq y \to x \cdot z \leq y \cdot z$
\item $\varphi_{23} :\equiv (x \in \mathbb{R} \to y \in \mathbb{R} \to z \in \mathbb{R} \to 0 \leq z \to x \leq y \to x \in \mathbb{R}) \to \varphi_{20}$
\item $\varphi_{24} :\equiv (x \in \mathbb{R} \to y \in \mathbb{R} \to z \in \mathbb{R} \to 0 \leq z \to x \leq y \to y \in \mathbb{R}) \to \varphi_{21}$
\item $\varphi_{25} :\equiv (x \in \mathbb{R} \to y \in \mathbb{R} \to z \in \mathbb{R} \to 0 \leq z \to x \leq y \to x \cdot z \leq y \cdot z) \to (x \cdot z = z \cdot x \to x \cdot z < y \cdot z \to z \cdot x < y \cdot z) \to \varphi_{17}$
\item $\varphi_{26} :\equiv (x \in \mathbb{R} \to y \in \mathbb{R} \to z \in \mathbb{R} \to 0 \leq z \to x \leq y \to z \cdot x \leq y \cdot z) \to (y \cdot z = z \cdot y \to z \cdot x < y \cdot z \to z \cdot x < z \cdot y) \to \varphi_{18}$
\item $\varphi_{27} :\equiv (x \in \mathbb{R} \to y \in \mathbb{R} \to z \in \mathbb{R} \to 0 \leq z \to x \leq y \to x \cdot z = z \cdot x) \to \varphi_{25}$
\item $\varphi_{28} :\equiv (x \in \mathbb{R} \to y \in \mathbb{R} \to z \in \mathbb{R} \to 0 \leq z \to x \leq y \to y \cdot z = z \cdot y) \to \varphi_{26}$
\end{abbrevlist}
\end{fullproof}

% 自動生成 (prover/build.js): PosNonneg  ノード数 4
\begin{fullproof}{[E10]}{$\forall x\,(0 < x \to 0 \leq x)$}{4}\phantomsection\label{tree:PosNonneg}
\noindent 主木:
\begin{ptree}
\AXC{$0 < x \to 0 < x$\lbl{A1}}
\AXC{$(0 < x \to 0 < x) \to 0 < x \to 0 \leq x$\lbl{A1}}
\BIC{$0 < x \to 0 \leq x$}
\UIC{$\forall x\,(0 < x \to 0 \leq x)$}
\end{ptree}
\end{fullproof}

% 自動生成 (prover/build.js): MulPosPos  ノード数 7
\begin{fullproof}{[E11]}{$\forall x\,\forall y\,(x \in \mathbb{R} \to y \in \mathbb{R} \to 0 < x \to 0 < y \to 0 < x \cdot y)$}{7}\phantomsection\label{tree:MulPosPos}
\noindent 主木:
\begin{ptree}
\AXC{$\forall x\,\forall y\,(x \in \mathbb{R} \to y \in \mathbb{R} \to 0 < x \to 0 < y \to 0 < x \cdot y)$\lbl{(R15)}}
\AXC{$\varphi_{2}$\lbl{A2}}
\BIC{$\forall y\,(x \in \mathbb{R} \to y \in \mathbb{R} \to 0 < x \to 0 < y \to 0 < x \cdot y)$}
\AXC{$\varphi_{1}$\lbl{A2}}
\BIC{$x \in \mathbb{R} \to y \in \mathbb{R} \to 0 < x \to 0 < y \to 0 < x \cdot y$}
\UIC{$\forall y\,(x \in \mathbb{R} \to y \in \mathbb{R} \to 0 < x \to 0 < y \to 0 < x \cdot y)$}
\UIC{$\forall x\,\forall y\,(x \in \mathbb{R} \to y \in \mathbb{R} \to 0 < x \to 0 < y \to 0 < x \cdot y)$}
\end{ptree}
\begin{abbrevlist}
\item $\varphi_{1} :\equiv \forall y\,(x \in \mathbb{R} \to y \in \mathbb{R} \to 0 < x \to 0 < y \to 0 < x \cdot y) \to x \in \mathbb{R} \to y \in \mathbb{R} \to 0 < x \to 0 < y \to 0 < x \cdot y$
\item $\varphi_{2} :\equiv \forall x\,\forall y\,(x \in \mathbb{R} \to y \in \mathbb{R} \to 0 < x \to 0 < y \to 0 < x \cdot y) \to \forall y\,(x \in \mathbb{R} \to y \in \mathbb{R} \to 0 < x \to 0 < y \to 0 < x \cdot y)$
\end{abbrevlist}
\end{fullproof}

% 自動生成 (prover/build.js): MulLeLe  ノード数 64
\begin{fullproof}{[E12]}{$\forall x\,\forall y\,\forall u\,\forall v\,(x \in \mathbb{R} \to y \in \mathbb{R} \to u \in \mathbb{R} \to v \in \mathbb{R} \to 0 \leq x \to 0 \leq u \to x \leq y \to u \leq v \to x \cdot u \leq y \cdot v)$}{64}\phantomsection\label{tree:MulLeLe}
\noindent$\Pi_{1}$:
\begin{ptree}
\AXC{$\forall x\,\forall y\,\forall z\,(x \leq y \to y \leq z \to x \leq z)$\lbl{[O5]}}
\AXC{$\varphi_{2}$\lbl{A2}}
\BIC{$\forall y\,\forall z\,(0 \leq y \to y \leq z \to 0 \leq z)$}
\AXC{$\varphi_{1}$\lbl{A2}}
\BIC{$\forall z\,(0 \leq x \to x \leq z \to 0 \leq z)$}
\AXC{$\forall z\,(0 \leq x \to x \leq z \to 0 \leq z) \to 0 \leq x \to x \leq y \to 0 \leq y$\lbl{A2}}
\BIC{$0 \leq x \to x \leq y \to 0 \leq y$}
\end{ptree}
\noindent$\Pi_{2}$:
\begin{ptree}
\AXC{$x \in \mathbb{R} \to \varphi_{8}$\lbl{A1}}
\AXC{$x \in \mathbb{R} \to \varphi_{9}$\lbl{A1}}
\AXC{$\stackrel{\Pi_{1}}{\vdots}$}\noLine
\UIC{$0 \leq x \to x \leq y \to 0 \leq y$}
\AXC{$\varphi_{24}$\lbl{A1}}
\BIC{$(x \in \mathbb{R} \to \varphi_{9}) \to (x \in \mathbb{R} \to \varphi_{8}) \to x \in \mathbb{R} \to \varphi_{10}$}
\BIC{$(x \in \mathbb{R} \to \varphi_{8}) \to x \in \mathbb{R} \to \varphi_{10}$}
\BIC{$x \in \mathbb{R} \to \varphi_{10}$}
\end{ptree}
\noindent$\Pi_{3}$:
\begin{ptree}
\AXC{$x \in \mathbb{R} \to \varphi_{18}$\lbl{A1}}
\AXC{$x \in \mathbb{R} \to \varphi_{19}$\lbl{A1}}
\AXC{$\varphi_{16}$\lbl{[E9]}}
\AXC{$\varphi_{16} \to \varphi_{14}$\lbl{A2}}
\BIC{$\varphi_{14}$}
\AXC{$\varphi_{14} \to \varphi_{12}$\lbl{A2}}
\BIC{$\varphi_{12}$}
\AXC{$\varphi_{22}$\lbl{A2}}
\BIC{$u \in \mathbb{R} \to v \in \mathbb{R} \to y \in \mathbb{R} \to 0 \leq y \to u \leq v \to y \cdot u \leq y \cdot v$}
\AXC{$\varphi_{28}$\lbl{A1}}
\BIC{$(x \in \mathbb{R} \to \varphi_{19}) \to (x \in \mathbb{R} \to \varphi_{18}) \to \varphi_{26}$}
\BIC{$(x \in \mathbb{R} \to \varphi_{18}) \to \varphi_{26}$}
\BIC{$\varphi_{26}$}
\end{ptree}
\noindent$\Pi_{4}$:
\begin{ptree}
\AXC{$x \in \mathbb{R} \to \varphi_{7}$\lbl{A1}}
\AXC{$\stackrel{\Pi_{2}}{\vdots}$}\noLine
\UIC{$x \in \mathbb{R} \to \varphi_{10}$}
\AXC{$x \in \mathbb{R} \to \varphi_{17}$\lbl{A1}}
\AXC{$\stackrel{\Pi_{3}}{\vdots}$}\noLine
\UIC{$\varphi_{26}$}
\BIC{$(x \in \mathbb{R} \to \varphi_{10}) \to (x \in \mathbb{R} \to \varphi_{7}) \to x \in \mathbb{R} \to y \in \mathbb{R} \to \varphi_{4}$}
\BIC{$(x \in \mathbb{R} \to \varphi_{7}) \to x \in \mathbb{R} \to y \in \mathbb{R} \to \varphi_{4}$}
\BIC{$x \in \mathbb{R} \to y \in \mathbb{R} \to \varphi_{4}$}
\end{ptree}
\noindent$\Pi_{5}$:
\begin{ptree}
\AXC{$x \in \mathbb{R} \to \varphi_{20}$\lbl{A1}}
\AXC{$\forall x\,\forall y\,(\varphi_{13})$\lbl{[F33]}}
\AXC{$\forall x\,\forall y\,(\varphi_{13}) \to \forall y\,(\varphi_{13})$\lbl{A2}}
\BIC{$\forall y\,(\varphi_{13})$}
\AXC{$\forall y\,(\varphi_{13}) \to \varphi_{13}$\lbl{A2}}
\BIC{$\varphi_{13}$}
\AXC{$\varphi_{23}$\lbl{A2}}
\BIC{$x \in \mathbb{R} \to y \in \mathbb{R} \to u \in \mathbb{R} \to 0 \leq u \to x \leq y \to x \cdot u \leq y \cdot u$}
\AXC{$\varphi_{29}$\lbl{A1}}
\BIC{$(x \in \mathbb{R} \to \varphi_{20}) \to (x \in \mathbb{R} \to \varphi_{17}) \to \varphi_{27}$}
\BIC{$(x \in \mathbb{R} \to \varphi_{17}) \to \varphi_{27}$}
\end{ptree}
\noindent$\Pi_{6}$:
\begin{ptree}
\AXC{$x \in \mathbb{R} \to \varphi_{8}$\lbl{A1}}
\AXC{$x \in \mathbb{R} \to \varphi_{11}$\lbl{A1}}
\AXC{$x \in \mathbb{R} \to \varphi_{19}$\lbl{A1}}
\AXC{$x \in \mathbb{R} \to \varphi_{17}$\lbl{A1}}
\AXC{$\stackrel{\Pi_{5}}{\vdots}$}\noLine
\UIC{$(x \in \mathbb{R} \to \varphi_{17}) \to \varphi_{27}$}
\BIC{$\varphi_{27}$}
\BIC{$(x \in \mathbb{R} \to \varphi_{11}) \to (x \in \mathbb{R} \to \varphi_{8}) \to x \in \mathbb{R} \to y \in \mathbb{R} \to \varphi_{5}$}
\BIC{$(x \in \mathbb{R} \to \varphi_{8}) \to x \in \mathbb{R} \to y \in \mathbb{R} \to \varphi_{5}$}
\BIC{$x \in \mathbb{R} \to y \in \mathbb{R} \to \varphi_{5}$}
\end{ptree}
\noindent$\Pi_{7}$:
\begin{ptree}
\AXC{$\forall x\,\forall y\,\forall z\,(x \leq y \to y \leq z \to x \leq z)$\lbl{[O5]}}
\AXC{$\varphi_{3}$\lbl{A2}}
\BIC{$\forall y\,\forall z\,(x \cdot u \leq y \to y \leq z \to x \cdot u \leq z)$}
\AXC{$\varphi_{15}$\lbl{A2}}
\BIC{$\forall z\,(x \cdot u \leq y \cdot u \to y \cdot u \leq z \to x \cdot u \leq z)$}
\AXC{$\varphi_{21}$\lbl{A2}}
\BIC{$x \cdot u \leq y \cdot u \to y \cdot u \leq y \cdot v \to x \cdot u \leq y \cdot v$}
\end{ptree}
\noindent$\Pi_{8}$:
\begin{ptree}
\AXC{$\stackrel{\Pi_{6}}{\vdots}$}\noLine
\UIC{$x \in \mathbb{R} \to y \in \mathbb{R} \to \varphi_{5}$}
\AXC{$\stackrel{\Pi_{7}}{\vdots}$}\noLine
\UIC{$x \cdot u \leq y \cdot u \to y \cdot u \leq y \cdot v \to x \cdot u \leq y \cdot v$}
\AXC{$(x \cdot u \leq y \cdot u \to y \cdot u \leq y \cdot v \to x \cdot u \leq y \cdot v) \to \varphi_{25}$\lbl{A1}}
\BIC{$\varphi_{25}$}
\BIC{$(x \in \mathbb{R} \to y \in \mathbb{R} \to \varphi_{4}) \to x \in \mathbb{R} \to y \in \mathbb{R} \to \varphi_{6}$}
\end{ptree}
\noindent 主木:
\begin{ptree}
\AXC{$\stackrel{\Pi_{4}}{\vdots}$}\noLine
\UIC{$x \in \mathbb{R} \to y \in \mathbb{R} \to \varphi_{4}$}
\AXC{$\stackrel{\Pi_{8}}{\vdots}$}\noLine
\UIC{$(x \in \mathbb{R} \to y \in \mathbb{R} \to \varphi_{4}) \to x \in \mathbb{R} \to y \in \mathbb{R} \to \varphi_{6}$}
\BIC{$x \in \mathbb{R} \to y \in \mathbb{R} \to \varphi_{6}$}
\UIC{$\forall v\,(x \in \mathbb{R} \to y \in \mathbb{R} \to \varphi_{6})$}
\UIC{$\forall u\,\forall v\,(x \in \mathbb{R} \to y \in \mathbb{R} \to \varphi_{6})$}
\UIC{$\forall y\,\forall u\,\forall v\,(x \in \mathbb{R} \to y \in \mathbb{R} \to \varphi_{6})$}
\UIC{$\forall x\,\forall y\,\forall u\,\forall v\,(x \in \mathbb{R} \to y \in \mathbb{R} \to \varphi_{6})$}
\end{ptree}
\begin{abbrevlist}
\item $\varphi_{1} :\equiv \forall y\,\forall z\,(0 \leq y \to y \leq z \to 0 \leq z) \to \forall z\,(0 \leq x \to x \leq z \to 0 \leq z)$
\item $\varphi_{2} :\equiv \forall x\,\forall y\,\forall z\,(x \leq y \to y \leq z \to x \leq z) \to \forall y\,\forall z\,(0 \leq y \to y \leq z \to 0 \leq z)$
\item $\varphi_{3} :\equiv \forall x\,\forall y\,\forall z\,(x \leq y \to y \leq z \to x \leq z) \to \forall y\,\forall z\,(x \cdot u \leq y \to y \leq z \to x \cdot u \leq z)$
\item $\varphi_{4} :\equiv u \in \mathbb{R} \to v \in \mathbb{R} \to 0 \leq x \to 0 \leq u \to x \leq y \to u \leq v \to y \cdot u \leq y \cdot v$
\item $\varphi_{5} :\equiv u \in \mathbb{R} \to v \in \mathbb{R} \to 0 \leq x \to 0 \leq u \to x \leq y \to u \leq v \to x \cdot u \leq y \cdot u$
\item $\varphi_{6} :\equiv u \in \mathbb{R} \to v \in \mathbb{R} \to 0 \leq x \to 0 \leq u \to x \leq y \to u \leq v \to x \cdot u \leq y \cdot v$
\item $\varphi_{7} :\equiv y \in \mathbb{R} \to u \in \mathbb{R} \to v \in \mathbb{R} \to 0 \leq x \to 0 \leq u \to x \leq y \to u \leq v \to u \leq v$
\item $\varphi_{8} :\equiv y \in \mathbb{R} \to u \in \mathbb{R} \to v \in \mathbb{R} \to 0 \leq x \to 0 \leq u \to x \leq y \to u \leq v \to x \leq y$
\item $\varphi_{9} :\equiv y \in \mathbb{R} \to u \in \mathbb{R} \to v \in \mathbb{R} \to 0 \leq x \to 0 \leq u \to x \leq y \to u \leq v \to 0 \leq x$
\item $\varphi_{10} :\equiv y \in \mathbb{R} \to u \in \mathbb{R} \to v \in \mathbb{R} \to 0 \leq x \to 0 \leq u \to x \leq y \to u \leq v \to 0 \leq y$
\item $\varphi_{11} :\equiv y \in \mathbb{R} \to u \in \mathbb{R} \to v \in \mathbb{R} \to 0 \leq x \to 0 \leq u \to x \leq y \to u \leq v \to 0 \leq u$
\item $\varphi_{12} :\equiv \forall z\,(u \in \mathbb{R} \to v \in \mathbb{R} \to z \in \mathbb{R} \to 0 \leq z \to u \leq v \to z \cdot u \leq z \cdot v)$
\item $\varphi_{13} :\equiv \forall z\,(x \in \mathbb{R} \to y \in \mathbb{R} \to z \in \mathbb{R} \to 0 \leq z \to x \leq y \to x \cdot z \leq y \cdot z)$
\item $\varphi_{14} :\equiv \forall y\,\forall z\,(u \in \mathbb{R} \to y \in \mathbb{R} \to z \in \mathbb{R} \to 0 \leq z \to u \leq y \to z \cdot u \leq z \cdot y)$
\item $\varphi_{15} :\equiv \forall y\,\forall z\,(x \cdot u \leq y \to y \leq z \to x \cdot u \leq z) \to \forall z\,(x \cdot u \leq y \cdot u \to y \cdot u \leq z \to x \cdot u \leq z)$
\item $\varphi_{16} :\equiv \forall x\,\forall y\,\forall z\,(x \in \mathbb{R} \to y \in \mathbb{R} \to z \in \mathbb{R} \to 0 \leq z \to x \leq y \to z \cdot x \leq z \cdot y)$
\item $\varphi_{17} :\equiv y \in \mathbb{R} \to u \in \mathbb{R} \to v \in \mathbb{R} \to 0 \leq x \to 0 \leq u \to x \leq y \to u \leq v \to y \in \mathbb{R}$
\item $\varphi_{18} :\equiv y \in \mathbb{R} \to u \in \mathbb{R} \to v \in \mathbb{R} \to 0 \leq x \to 0 \leq u \to x \leq y \to u \leq v \to v \in \mathbb{R}$
\item $\varphi_{19} :\equiv y \in \mathbb{R} \to u \in \mathbb{R} \to v \in \mathbb{R} \to 0 \leq x \to 0 \leq u \to x \leq y \to u \leq v \to u \in \mathbb{R}$
\item $\varphi_{20} :\equiv y \in \mathbb{R} \to u \in \mathbb{R} \to v \in \mathbb{R} \to 0 \leq x \to 0 \leq u \to x \leq y \to u \leq v \to x \in \mathbb{R}$
\item $\varphi_{21} :\equiv \forall z\,(x \cdot u \leq y \cdot u \to y \cdot u \leq z \to x \cdot u \leq z) \to x \cdot u \leq y \cdot u \to y \cdot u \leq y \cdot v \to x \cdot u \leq y \cdot v$
\item $\varphi_{22} :\equiv \varphi_{12} \to u \in \mathbb{R} \to v \in \mathbb{R} \to y \in \mathbb{R} \to 0 \leq y \to u \leq v \to y \cdot u \leq y \cdot v$
\item $\varphi_{23} :\equiv \varphi_{13} \to x \in \mathbb{R} \to y \in \mathbb{R} \to u \in \mathbb{R} \to 0 \leq u \to x \leq y \to x \cdot u \leq y \cdot u$
\item $\varphi_{24} :\equiv (0 \leq x \to x \leq y \to 0 \leq y) \to (x \in \mathbb{R} \to \varphi_{9}) \to (x \in \mathbb{R} \to \varphi_{8}) \to x \in \mathbb{R} \to \varphi_{10}$
\item $\varphi_{25} :\equiv (x \in \mathbb{R} \to y \in \mathbb{R} \to \varphi_{5}) \to (x \in \mathbb{R} \to y \in \mathbb{R} \to \varphi_{4}) \to x \in \mathbb{R} \to y \in \mathbb{R} \to \varphi_{6}$
\item $\varphi_{26} :\equiv (x \in \mathbb{R} \to \varphi_{17}) \to (x \in \mathbb{R} \to \varphi_{10}) \to (x \in \mathbb{R} \to \varphi_{7}) \to x \in \mathbb{R} \to y \in \mathbb{R} \to \varphi_{4}$
\item $\varphi_{27} :\equiv (x \in \mathbb{R} \to \varphi_{19}) \to (x \in \mathbb{R} \to \varphi_{11}) \to (x \in \mathbb{R} \to \varphi_{8}) \to x \in \mathbb{R} \to y \in \mathbb{R} \to \varphi_{5}$
\item $\varphi_{28} :\equiv (u \in \mathbb{R} \to v \in \mathbb{R} \to y \in \mathbb{R} \to 0 \leq y \to u \leq v \to y \cdot u \leq y \cdot v) \to (x \in \mathbb{R} \to \varphi_{19}) \to (x \in \mathbb{R} \to \varphi_{18}) \to \varphi_{26}$
\item $\varphi_{29} :\equiv (x \in \mathbb{R} \to y \in \mathbb{R} \to u \in \mathbb{R} \to 0 \leq u \to x \leq y \to x \cdot u \leq y \cdot u) \to (x \in \mathbb{R} \to \varphi_{20}) \to (x \in \mathbb{R} \to \varphi_{17}) \to \varphi_{27}$
\end{abbrevlist}
\end{fullproof}

% 自動生成 (prover/build.js): InvUnique  ノード数 284
\begin{fullproof}{[E13]}{$\forall x\,\forall y\,(x \in \mathbb{R} \to y \in \mathbb{R} \to x \neq 0 \to x \cdot y = 1 \to y = x^{-1})$}{284}\phantomsection\label{tree:InvUnique}
\noindent$\Pi_{1}$:
\begin{ptree}
\AXC{$y = z \to x = y \to x = z$\lbl{A6}}
\AXC{$(y = z \to x = y \to x = z) \to x = y \to y = z \to x = z$\lbl{A1}}
\BIC{$x = y \to y = z \to x = z$}
\UIC{$\forall z\,(x = y \to y = z \to x = z)$}
\UIC{$\forall y\,\forall z\,(x = y \to y = z \to x = z)$}
\UIC{$\forall x\,\forall y\,\forall z\,(x = y \to y = z \to x = z)$}
\AXC{$\varphi_{11}$\lbl{A2}}
\BIC{$\forall y\,\forall z\,(v_{477} = y \to y = z \to v_{477} = z)$}
\AXC{$\varphi_{23}$\lbl{A2}}
\BIC{$\forall z\,(v_{477} = v_{478} \to v_{478} = z \to v_{477} = z)$}
\AXC{$\varphi_{37}$\lbl{A2}}
\BIC{$v_{477} = v_{478} \to v_{478} = v_{479} \to v_{477} = v_{479}$}
\UIC{$\forall v_{479}\,(v_{477} = v_{478} \to v_{478} = v_{479} \to v_{477} = v_{479})$}
\UIC{$\varphi_{3}$}
\UIC{$\forall v_{477}\,(\varphi_{3})$}
\end{ptree}
\noindent$\Pi_{2}$:
\begin{ptree}
\AXC{$\stackrel{\Pi_{1}}{\vdots}$}\noLine
\UIC{$\forall v_{477}\,(\varphi_{3})$}
\AXC{$\varphi_{71}$\lbl{A2}}
\BIC{$\forall v_{478}\,\forall v_{479}\,(y = v_{478} \to v_{478} = v_{479} \to y = v_{479})$}
\AXC{$\varphi_{66}$\lbl{A2}}
\BIC{$\forall v_{479}\,(y = x^{-1} \cdot 1 \to x^{-1} \cdot 1 = v_{479} \to y = v_{479})$}
\AXC{$\varphi_{77}$\lbl{A2}}
\BIC{$y = x^{-1} \cdot 1 \to x^{-1} \cdot 1 = x^{-1} \to y = x^{-1}$}
\end{ptree}
\noindent$\Pi_{3}$:
\begin{ptree}
\AXC{$x \in \mathbb{R} \to y \in \mathbb{R} \to x \neq 0 \to x \cdot y = 1 \to x \in \mathbb{R}$\lbl{A1}}
\AXC{$\forall x \in \mathbb{R}\,(x \neq 0 \to x^{-1} \in \mathbb{R})$\lbl{(R1)}}
\AXC{$\varphi_{55}$\lbl{A2}}
\BIC{$x \in \mathbb{R} \to x \neq 0 \to x^{-1} \in \mathbb{R}$}
\AXC{$\varphi_{107}$\lbl{A1}}
\BIC{$(x \in \mathbb{R} \to y \in \mathbb{R} \to x \neq 0 \to x \cdot y = 1 \to x \in \mathbb{R}) \to \varphi_{89}$}
\BIC{$\varphi_{89}$}
\end{ptree}
\noindent$\Pi_{4}$:
\begin{ptree}
\AXC{$x \in \mathbb{R} \to y \in \mathbb{R} \to x \neq 0 \to x \cdot y = 1 \to x \neq 0$\lbl{A1}}
\AXC{$\stackrel{\Pi_{3}}{\vdots}$}\noLine
\UIC{$\varphi_{89}$}
\BIC{$x \in \mathbb{R} \to y \in \mathbb{R} \to x \neq 0 \to x \cdot y = 1 \to x^{-1} \in \mathbb{R}$}
\end{ptree}
\noindent$\Pi_{5}$:
\begin{ptree}
\AXC{$\forall x \in \mathbb{R}\,(x \cdot 1 = x)$\lbl{(R8)}}
\AXC{$\forall x \in \mathbb{R}\,(x \cdot 1 = x) \to x^{-1} \in \mathbb{R} \to x^{-1} \cdot 1 = x^{-1}$\lbl{A2}}
\BIC{$x^{-1} \in \mathbb{R} \to x^{-1} \cdot 1 = x^{-1}$}
\end{ptree}
\noindent$\Pi_{6}$:
\begin{ptree}
\AXC{$\stackrel{\Pi_{4}}{\vdots}$}\noLine
\UIC{$x \in \mathbb{R} \to y \in \mathbb{R} \to x \neq 0 \to x \cdot y = 1 \to x^{-1} \in \mathbb{R}$}
\AXC{$\stackrel{\Pi_{5}}{\vdots}$}\noLine
\UIC{$x^{-1} \in \mathbb{R} \to x^{-1} \cdot 1 = x^{-1}$}
\AXC{$(x^{-1} \in \mathbb{R} \to x^{-1} \cdot 1 = x^{-1}) \to \varphi_{94}$\lbl{A1}}
\BIC{$\varphi_{94}$}
\BIC{$x \in \mathbb{R} \to y \in \mathbb{R} \to x \neq 0 \to x \cdot y = 1 \to x^{-1} \cdot 1 = x^{-1}$}
\end{ptree}
\noindent$\Pi_{7}$:
\begin{ptree}
\AXC{$y = z \to x = y \to x = z$\lbl{A6}}
\AXC{$(y = z \to x = y \to x = z) \to x = y \to y = z \to x = z$\lbl{A1}}
\BIC{$x = y \to y = z \to x = z$}
\UIC{$\forall z\,(x = y \to y = z \to x = z)$}
\UIC{$\forall y\,\forall z\,(x = y \to y = z \to x = z)$}
\UIC{$\forall x\,\forall y\,\forall z\,(x = y \to y = z \to x = z)$}
\AXC{$\varphi_{12}$\lbl{A2}}
\BIC{$\forall y\,\forall z\,(v_{474} = y \to y = z \to v_{474} = z)$}
\AXC{$\varphi_{24}$\lbl{A2}}
\BIC{$\forall z\,(v_{474} = v_{475} \to v_{475} = z \to v_{474} = z)$}
\AXC{$\varphi_{38}$\lbl{A2}}
\BIC{$v_{474} = v_{475} \to v_{475} = v_{476} \to v_{474} = v_{476}$}
\UIC{$\forall v_{476}\,(v_{474} = v_{475} \to v_{475} = v_{476} \to v_{474} = v_{476})$}
\UIC{$\varphi_{4}$}
\UIC{$\forall v_{474}\,(\varphi_{4})$}
\end{ptree}
\noindent$\Pi_{8}$:
\begin{ptree}
\AXC{$\forall x\,\forall y\,(x \in \mathbb{R} \to y \in \mathbb{R} \to x \cdot y = y \cdot x)$\lbl{(R7)}}
\AXC{$\varphi_{69}$\lbl{A2}}
\BIC{$\forall y\,(1 \in \mathbb{R} \to y \in \mathbb{R} \to 1 \cdot y = y \cdot 1)$}
\AXC{$\varphi_{80}$\lbl{A2}}
\BIC{$1 \in \mathbb{R} \to x^{-1} \in \mathbb{R} \to 1 \cdot x^{-1} = x^{-1} \cdot 1$}
\AXC{$(1 \in \mathbb{R} \to x^{-1} \in \mathbb{R} \to 1 \cdot x^{-1} = x^{-1} \cdot 1) \to \varphi_{100}$\lbl{A1}}
\BIC{$\varphi_{100}$}
\end{ptree}
\noindent$\Pi_{9}$:
\begin{ptree}
\AXC{$\stackrel{\Pi_{4}}{\vdots}$}\noLine
\UIC{$x \in \mathbb{R} \to y \in \mathbb{R} \to x \neq 0 \to x \cdot y = 1 \to x^{-1} \in \mathbb{R}$}
\AXC{$1 \in \mathbb{R}$\lbl{(R1)}}
\AXC{$\stackrel{\Pi_{8}}{\vdots}$}\noLine
\UIC{$\varphi_{100}$}
\BIC{$(x \in \mathbb{R} \to y \in \mathbb{R} \to x \neq 0 \to x \cdot y = 1 \to x^{-1} \in \mathbb{R}) \to \varphi_{43}$}
\BIC{$\varphi_{43}$}
\end{ptree}
\noindent$\Pi_{10}$:
\begin{ptree}
\AXC{$y = z \to x = y \to x = z$\lbl{A6}}
\AXC{$(y = z \to x = y \to x = z) \to x = y \to y = z \to x = z$\lbl{A1}}
\BIC{$x = y \to y = z \to x = z$}
\UIC{$\forall z\,(x = y \to y = z \to x = z)$}
\UIC{$\forall y\,\forall z\,(x = y \to y = z \to x = z)$}
\UIC{$\forall x\,\forall y\,\forall z\,(x = y \to y = z \to x = z)$}
\AXC{$\varphi_{13}$\lbl{A2}}
\BIC{$\forall y\,\forall z\,(v_{471} = y \to y = z \to v_{471} = z)$}
\AXC{$\varphi_{25}$\lbl{A2}}
\BIC{$\forall z\,(v_{471} = v_{472} \to v_{472} = z \to v_{471} = z)$}
\AXC{$\varphi_{39}$\lbl{A2}}
\BIC{$v_{471} = v_{472} \to v_{472} = v_{473} \to v_{471} = v_{473}$}
\UIC{$\forall v_{473}\,(v_{471} = v_{472} \to v_{472} = v_{473} \to v_{471} = v_{473})$}
\UIC{$\varphi_{5}$}
\UIC{$\forall v_{471}\,(\varphi_{5})$}
\end{ptree}
\noindent$\Pi_{11}$:
\begin{ptree}
\AXC{$\stackrel{\Pi_{10}}{\vdots}$}\noLine
\UIC{$\forall v_{471}\,(\varphi_{5})$}
\AXC{$\varphi_{73}$\lbl{A2}}
\BIC{$\forall v_{472}\,\forall v_{473}\,(y = v_{472} \to v_{472} = v_{473} \to y = v_{473})$}
\AXC{$\varphi_{81}$\lbl{A2}}
\BIC{$\varphi_{18}$}
\AXC{$\varphi_{18} \to \varphi_{34}$\lbl{A2}}
\BIC{$\varphi_{34}$}
\end{ptree}
\noindent$\Pi_{12}$:
\begin{ptree}
\AXC{$x = x$\lbl{A5}}
\UIC{$\forall x\,(x = x)$}
\AXC{$\forall x\,(x = x) \to x \cdot y \cdot x^{-1} = x \cdot y \cdot x^{-1}$\lbl{A2}}
\BIC{$x \cdot y \cdot x^{-1} = x \cdot y \cdot x^{-1}$}
\AXC{$x \in \mathbb{R} \to y \in \mathbb{R} \to x \neq 0 \to x \cdot y = 1 \to x \cdot y = 1$\lbl{A1}}
\AXC{$\varphi_{110}$\lbl{A1}}
\BIC{$x \cdot y \cdot x^{-1} = x \cdot y \cdot x^{-1} \to \varphi_{52} \to \varphi_{49}$}
\BIC{$\varphi_{52} \to \varphi_{49}$}
\end{ptree}
\noindent$\Pi_{13}$:
\begin{ptree}
\AXC{$y = z \to x = y \to x = z$\lbl{A6}}
\AXC{$(y = z \to x = y \to x = z) \to x = y \to y = z \to x = z$\lbl{A1}}
\BIC{$x = y \to y = z \to x = z$}
\UIC{$\forall z\,(x = y \to y = z \to x = z)$}
\UIC{$\forall y\,\forall z\,(x = y \to y = z \to x = z)$}
\UIC{$\forall x\,\forall y\,\forall z\,(x = y \to y = z \to x = z)$}
\AXC{$\varphi_{14}$\lbl{A2}}
\BIC{$\forall y\,\forall z\,(v_{468} = y \to y = z \to v_{468} = z)$}
\AXC{$\varphi_{26}$\lbl{A2}}
\BIC{$\forall z\,(v_{468} = v_{469} \to v_{469} = z \to v_{468} = z)$}
\AXC{$\varphi_{40}$\lbl{A2}}
\BIC{$v_{468} = v_{469} \to v_{469} = v_{470} \to v_{468} = v_{470}$}
\UIC{$\forall v_{470}\,(v_{468} = v_{469} \to v_{469} = v_{470} \to v_{468} = v_{470})$}
\UIC{$\varphi_{6}$}
\UIC{$\forall v_{468}\,(\varphi_{6})$}
\end{ptree}
\noindent$\Pi_{14}$:
\begin{ptree}
\AXC{$\stackrel{\Pi_{13}}{\vdots}$}\noLine
\UIC{$\forall v_{468}\,(\varphi_{6})$}
\AXC{$\varphi_{74}$\lbl{A2}}
\BIC{$\forall v_{469}\,\forall v_{470}\,(y = v_{469} \to v_{469} = v_{470} \to y = v_{470})$}
\AXC{$\varphi_{82}$\lbl{A2}}
\BIC{$\varphi_{19}$}
\AXC{$\varphi_{19} \to \varphi_{53}$\lbl{A2}}
\BIC{$\varphi_{53}$}
\end{ptree}
\noindent$\Pi_{15}$:
\begin{ptree}
\AXC{$\forall x\,\forall y\,(x \in \mathbb{R} \to y \in \mathbb{R} \to x \cdot y = y \cdot x)$\lbl{(R7)}}
\AXC{$\varphi_{83}$\lbl{A2}}
\BIC{$\forall y\,(v_{460} \in \mathbb{R} \to y \in \mathbb{R} \to v_{460} \cdot y = y \cdot v_{460})$}
\AXC{$\varphi_{90}$\lbl{A2}}
\BIC{$v_{460} \in \mathbb{R} \to v_{461} \in \mathbb{R} \to v_{460} \cdot v_{461} = v_{461} \cdot v_{460}$}
\UIC{$\varphi_{30}$}
\UIC{$\forall v_{460}\,(\varphi_{30})$}
\AXC{$\varphi_{96}$\lbl{A2}}
\BIC{$\forall v_{461}\,(y \in \mathbb{R} \to v_{461} \in \mathbb{R} \to y \cdot v_{461} = v_{461} \cdot y)$}
\AXC{$\varphi_{78}$\lbl{A2}}
\BIC{$y \in \mathbb{R} \to x \in \mathbb{R} \to y \cdot x = x \cdot y$}
\AXC{$\varphi_{108}$\lbl{A1}}
\BIC{$(x \in \mathbb{R} \to y \in \mathbb{R} \to x \neq 0 \to x \cdot y = 1 \to y \in \mathbb{R}) \to \varphi_{91}$}
\end{ptree}
\noindent$\Pi_{16}$:
\begin{ptree}
\AXC{$x \in \mathbb{R} \to y \in \mathbb{R} \to x \neq 0 \to x \cdot y = 1 \to y \in \mathbb{R}$\lbl{A1}}
\AXC{$\stackrel{\Pi_{15}}{\vdots}$}\noLine
\UIC{$(x \in \mathbb{R} \to y \in \mathbb{R} \to x \neq 0 \to x \cdot y = 1 \to y \in \mathbb{R}) \to \varphi_{91}$}
\BIC{$\varphi_{91}$}
\end{ptree}
\noindent$\Pi_{17}$:
\begin{ptree}
\AXC{$x \in \mathbb{R} \to y \in \mathbb{R} \to x \neq 0 \to x \cdot y = 1 \to x \in \mathbb{R}$\lbl{A1}}
\AXC{$\stackrel{\Pi_{16}}{\vdots}$}\noLine
\UIC{$\varphi_{91}$}
\BIC{$x \in \mathbb{R} \to y \in \mathbb{R} \to x \neq 0 \to x \cdot y = 1 \to y \cdot x = x \cdot y$}
\AXC{$\varphi_{114}$\lbl{A1}}
\BIC{$y \cdot x \cdot x^{-1} = y \cdot x \cdot x^{-1} \to (y \cdot x = x \cdot y \to \varphi_{47}) \to \varphi_{56}$}
\end{ptree}
\noindent$\Pi_{18}$:
\begin{ptree}
\AXC{$x = x$\lbl{A5}}
\UIC{$\forall x\,(x = x)$}
\AXC{$\forall x\,(x = x) \to y \cdot x \cdot x^{-1} = y \cdot x \cdot x^{-1}$\lbl{A2}}
\BIC{$y \cdot x \cdot x^{-1} = y \cdot x \cdot x^{-1}$}
\AXC{$\stackrel{\Pi_{17}}{\vdots}$}\noLine
\UIC{$y \cdot x \cdot x^{-1} = y \cdot x \cdot x^{-1} \to (y \cdot x = x \cdot y \to \varphi_{47}) \to \varphi_{56}$}
\BIC{$(y \cdot x = x \cdot y \to \varphi_{47}) \to \varphi_{56}$}
\end{ptree}
\noindent$\Pi_{19}$:
\begin{ptree}
\AXC{$\varphi_{36}$\lbl{A6}}
\UIC{$\forall w\,(\varphi_{36})$}
\UIC{$\forall z\,\forall w\,(\varphi_{36})$}
\AXC{$\forall z\,\forall w\,(\varphi_{36}) \to \varphi_{61}$\lbl{A2}}
\BIC{$\varphi_{61}$}
\AXC{$\varphi_{61} \to y \cdot x = x \cdot y \to \varphi_{47}$\lbl{A2}}
\BIC{$y \cdot x = x \cdot y \to \varphi_{47}$}
\AXC{$\stackrel{\Pi_{18}}{\vdots}$}\noLine
\UIC{$(y \cdot x = x \cdot y \to \varphi_{47}) \to \varphi_{56}$}
\BIC{$\varphi_{56}$}
\end{ptree}
\noindent$\Pi_{20}$:
\begin{ptree}
\AXC{$y = z \to x = y \to x = z$\lbl{A6}}
\AXC{$(y = z \to x = y \to x = z) \to x = y \to y = z \to x = z$\lbl{A1}}
\BIC{$x = y \to y = z \to x = z$}
\UIC{$\forall z\,(x = y \to y = z \to x = z)$}
\UIC{$\forall y\,\forall z\,(x = y \to y = z \to x = z)$}
\UIC{$\forall x\,\forall y\,\forall z\,(x = y \to y = z \to x = z)$}
\AXC{$\varphi_{15}$\lbl{A2}}
\BIC{$\forall y\,\forall z\,(v_{465} = y \to y = z \to v_{465} = z)$}
\AXC{$\varphi_{27}$\lbl{A2}}
\BIC{$\forall z\,(v_{465} = v_{466} \to v_{466} = z \to v_{465} = z)$}
\AXC{$\varphi_{41}$\lbl{A2}}
\BIC{$v_{465} = v_{466} \to v_{466} = v_{467} \to v_{465} = v_{467}$}
\UIC{$\forall v_{467}\,(v_{465} = v_{466} \to v_{466} = v_{467} \to v_{465} = v_{467})$}
\UIC{$\varphi_{7}$}
\UIC{$\forall v_{465}\,(\varphi_{7})$}
\end{ptree}
\noindent$\Pi_{21}$:
\begin{ptree}
\AXC{$x = x$\lbl{A5}}
\AXC{$x = y \to x = x \to y = x$\lbl{A6}}
\AXC{$(x = y \to x = x \to y = x) \to x = x \to x = y \to y = x$\lbl{A1}}
\BIC{$x = x \to x = y \to y = x$}
\BIC{$x = y \to y = x$}
\UIC{$\forall y\,(x = y \to y = x)$}
\UIC{$\forall x\,\forall y\,(x = y \to y = x)$}
\end{ptree}
\noindent$\Pi_{22}$:
\begin{ptree}
\AXC{$\stackrel{\Pi_{21}}{\vdots}$}\noLine
\UIC{$\forall x\,\forall y\,(x = y \to y = x)$}
\AXC{$\forall x\,\forall y\,(x = y \to y = x) \to \forall y\,(v_{458} = y \to y = v_{458})$\lbl{A2}}
\BIC{$\forall y\,(v_{458} = y \to y = v_{458})$}
\AXC{$\varphi_{1}$\lbl{A2}}
\BIC{$v_{458} = v_{459} \to v_{459} = v_{458}$}
\UIC{$\forall v_{459}\,(v_{458} = v_{459} \to v_{459} = v_{458})$}
\UIC{$\forall v_{458}\,\forall v_{459}\,(v_{458} = v_{459} \to v_{459} = v_{458})$}
\AXC{$\varphi_{65}$\lbl{A2}}
\BIC{$\forall v_{459}\,(y \cdot x \cdot x^{-1} = v_{459} \to v_{459} = y \cdot x \cdot x^{-1})$}
\AXC{$\varphi_{93}$\lbl{A2}}
\BIC{$\varphi_{46}$}
\end{ptree}
\noindent$\Pi_{23}$:
\begin{ptree}
\AXC{$\varphi_{48}$\lbl{(R6)}}
\AXC{$\varphi_{48} \to \varphi_{57}$\lbl{A2}}
\BIC{$\varphi_{57}$}
\AXC{$\varphi_{57} \to \varphi_{64}$\lbl{A2}}
\BIC{$\varphi_{64}$}
\AXC{$\varphi_{64} \to v_{455} \in \mathbb{R} \to v_{456} \in \mathbb{R} \to \varphi_{21}$\lbl{A2}}
\BIC{$v_{455} \in \mathbb{R} \to v_{456} \in \mathbb{R} \to \varphi_{21}$}
\UIC{$\forall v_{457}\,(v_{455} \in \mathbb{R} \to v_{456} \in \mathbb{R} \to \varphi_{21})$}
\UIC{$\forall v_{456}\,\forall v_{457}\,(v_{455} \in \mathbb{R} \to v_{456} \in \mathbb{R} \to \varphi_{21})$}
\UIC{$\forall v_{455}\,\forall v_{456}\,\forall v_{457}\,(v_{455} \in \mathbb{R} \to v_{456} \in \mathbb{R} \to \varphi_{21})$}
\AXC{$\varphi_{106}$\lbl{A2}}
\BIC{$\varphi_{70}$}
\AXC{$\varphi_{70} \to \varphi_{58}$\lbl{A2}}
\BIC{$\varphi_{58}$}
\end{ptree}
\noindent$\Pi_{24}$:
\begin{ptree}
\AXC{$x \in \mathbb{R} \to y \in \mathbb{R} \to x \neq 0 \to x \cdot y = 1 \to y \in \mathbb{R}$\lbl{A1}}
\AXC{$\stackrel{\Pi_{23}}{\vdots}$}\noLine
\UIC{$\varphi_{58}$}
\AXC{$\varphi_{58} \to \varphi_{59}$\lbl{A2}}
\BIC{$\varphi_{59}$}
\AXC{$\varphi_{59} \to \varphi_{113}$\lbl{A1}}
\BIC{$\varphi_{113}$}
\BIC{$(x \in \mathbb{R} \to y \in \mathbb{R} \to x \neq 0 \to x \cdot y = 1 \to x \in \mathbb{R}) \to \varphi_{97}$}
\end{ptree}
\noindent$\Pi_{25}$:
\begin{ptree}
\AXC{$x \in \mathbb{R} \to y \in \mathbb{R} \to x \neq 0 \to x \cdot y = 1 \to x \in \mathbb{R}$\lbl{A1}}
\AXC{$\stackrel{\Pi_{24}}{\vdots}$}\noLine
\UIC{$(x \in \mathbb{R} \to y \in \mathbb{R} \to x \neq 0 \to x \cdot y = 1 \to x \in \mathbb{R}) \to \varphi_{97}$}
\BIC{$\varphi_{97}$}
\end{ptree}
\noindent$\Pi_{26}$:
\begin{ptree}
\AXC{$\stackrel{\Pi_{22}}{\vdots}$}\noLine
\UIC{$\varphi_{46}$}
\AXC{$\stackrel{\Pi_{4}}{\vdots}$}\noLine
\UIC{$x \in \mathbb{R} \to y \in \mathbb{R} \to x \neq 0 \to x \cdot y = 1 \to x^{-1} \in \mathbb{R}$}
\AXC{$\stackrel{\Pi_{25}}{\vdots}$}\noLine
\UIC{$\varphi_{97}$}
\BIC{$x \in \mathbb{R} \to \varphi_{32}$}
\AXC{$(x \in \mathbb{R} \to \varphi_{32}) \to \varphi_{46} \to x \in \mathbb{R} \to \varphi_{33}$\lbl{A1}}
\BIC{$\varphi_{46} \to x \in \mathbb{R} \to \varphi_{33}$}
\BIC{$x \in \mathbb{R} \to \varphi_{33}$}
\end{ptree}
\noindent$\Pi_{27}$:
\begin{ptree}
\AXC{$y = z \to x = y \to x = z$\lbl{A6}}
\AXC{$(y = z \to x = y \to x = z) \to x = y \to y = z \to x = z$\lbl{A1}}
\BIC{$x = y \to y = z \to x = z$}
\UIC{$\forall z\,(x = y \to y = z \to x = z)$}
\UIC{$\forall y\,\forall z\,(x = y \to y = z \to x = z)$}
\UIC{$\forall x\,\forall y\,\forall z\,(x = y \to y = z \to x = z)$}
\AXC{$\varphi_{16}$\lbl{A2}}
\BIC{$\forall y\,\forall z\,(v_{462} = y \to y = z \to v_{462} = z)$}
\AXC{$\varphi_{28}$\lbl{A2}}
\BIC{$\forall z\,(v_{462} = v_{463} \to v_{463} = z \to v_{462} = z)$}
\AXC{$\varphi_{42}$\lbl{A2}}
\BIC{$v_{462} = v_{463} \to v_{463} = v_{464} \to v_{462} = v_{464}$}
\UIC{$\forall v_{464}\,(v_{462} = v_{463} \to v_{463} = v_{464} \to v_{462} = v_{464})$}
\UIC{$\varphi_{8}$}
\UIC{$\forall v_{462}\,(\varphi_{8})$}
\end{ptree}
\noindent$\Pi_{28}$:
\begin{ptree}
\AXC{$\stackrel{\Pi_{27}}{\vdots}$}\noLine
\UIC{$\forall v_{462}\,(\varphi_{8})$}
\AXC{$\varphi_{76}$\lbl{A2}}
\BIC{$\forall v_{463}\,\forall v_{464}\,(y = v_{463} \to v_{463} = v_{464} \to y = v_{464})$}
\AXC{$\varphi_{62}$\lbl{A2}}
\BIC{$\forall v_{464}\,(y = y \cdot 1 \to y \cdot 1 = v_{464} \to y = v_{464})$}
\AXC{$\forall v_{464}\,(y = y \cdot 1 \to y \cdot 1 = v_{464} \to y = v_{464}) \to \varphi_{17}$\lbl{A2}}
\BIC{$\varphi_{17}$}
\end{ptree}
\noindent$\Pi_{29}$:
\begin{ptree}
\AXC{$x \in \mathbb{R} \to y \in \mathbb{R} \to x \neq 0 \to x \cdot y = 1 \to x \in \mathbb{R}$\lbl{A1}}
\AXC{$\forall x \in \mathbb{R}\,(x \neq 0 \to x \cdot x^{-1} = 1)$\lbl{(R9)}}
\AXC{$\varphi_{50}$\lbl{A2}}
\BIC{$x \in \mathbb{R} \to x \neq 0 \to x \cdot x^{-1} = 1$}
\AXC{$\varphi_{104}$\lbl{A1}}
\BIC{$(x \in \mathbb{R} \to y \in \mathbb{R} \to x \neq 0 \to x \cdot y = 1 \to x \in \mathbb{R}) \to \varphi_{87}$}
\BIC{$\varphi_{87}$}
\end{ptree}
\noindent$\Pi_{30}$:
\begin{ptree}
\AXC{$x \in \mathbb{R} \to y \in \mathbb{R} \to x \neq 0 \to x \cdot y = 1 \to x \neq 0$\lbl{A1}}
\AXC{$\stackrel{\Pi_{29}}{\vdots}$}\noLine
\UIC{$\varphi_{87}$}
\BIC{$x \in \mathbb{R} \to y \in \mathbb{R} \to x \neq 0 \to x \cdot y = 1 \to x \cdot x^{-1} = 1$}
\AXC{$\varphi_{101}$\lbl{A1}}
\BIC{$\varphi_{68}$}
\end{ptree}
\noindent$\Pi_{31}$:
\begin{ptree}
\AXC{$\stackrel{\Pi_{21}}{\vdots}$}\noLine
\UIC{$\forall x\,\forall y\,(x = y \to y = x)$}
\AXC{$\varphi_{9}$\lbl{A2}}
\BIC{$\forall y\,(x \cdot x^{-1} = y \to y = x \cdot x^{-1})$}
\AXC{$\varphi_{31}$\lbl{A2}}
\BIC{$x \cdot x^{-1} = 1 \to 1 = x \cdot x^{-1}$}
\AXC{$\stackrel{\Pi_{30}}{\vdots}$}\noLine
\UIC{$\varphi_{68}$}
\BIC{$x \in \mathbb{R} \to y \in \mathbb{R} \to x \neq 0 \to x \cdot y = 1 \to 1 = x \cdot x^{-1}$}
\AXC{$\varphi_{105}$\lbl{A1}}
\BIC{$y \cdot 1 = y \cdot 1 \to \varphi_{95}$}
\end{ptree}
\noindent$\Pi_{32}$:
\begin{ptree}
\AXC{$x = x$\lbl{A5}}
\UIC{$\forall x\,(x = x)$}
\AXC{$\forall x\,(x = x) \to y \cdot 1 = y \cdot 1$\lbl{A2}}
\BIC{$y \cdot 1 = y \cdot 1$}
\AXC{$\stackrel{\Pi_{31}}{\vdots}$}\noLine
\UIC{$y \cdot 1 = y \cdot 1 \to \varphi_{95}$}
\BIC{$\varphi_{95}$}
\end{ptree}
\noindent$\Pi_{33}$:
\begin{ptree}
\AXC{$z = w \to y \cdot 1 = y \cdot z \to y \cdot 1 = y \cdot w$\lbl{A6}}
\UIC{$\forall w\,(z = w \to y \cdot 1 = y \cdot z \to y \cdot 1 = y \cdot w)$}
\UIC{$\forall z\,\forall w\,(z = w \to y \cdot 1 = y \cdot z \to y \cdot 1 = y \cdot w)$}
\AXC{$\varphi_{63}$\lbl{A2}}
\BIC{$\forall w\,(1 = w \to y \cdot 1 = y \cdot 1 \to y \cdot 1 = y \cdot w)$}
\AXC{$\varphi_{79}$\lbl{A2}}
\BIC{$1 = x \cdot x^{-1} \to y \cdot 1 = y \cdot 1 \to y \cdot 1 = y \cdot (x \cdot x^{-1})$}
\AXC{$\stackrel{\Pi_{32}}{\vdots}$}\noLine
\UIC{$\varphi_{95}$}
\BIC{$\varphi_{44}$}
\end{ptree}
\noindent$\Pi_{34}$:
\begin{ptree}
\AXC{$\stackrel{\Pi_{21}}{\vdots}$}\noLine
\UIC{$\forall x\,\forall y\,(x = y \to y = x)$}
\AXC{$\forall x\,\forall y\,(x = y \to y = x) \to \forall y\,(v_{453} = y \to y = v_{453})$\lbl{A2}}
\BIC{$\forall y\,(v_{453} = y \to y = v_{453})$}
\AXC{$\varphi_{2}$\lbl{A2}}
\BIC{$v_{453} = v_{454} \to v_{454} = v_{453}$}
\UIC{$\forall v_{454}\,(v_{453} = v_{454} \to v_{454} = v_{453})$}
\UIC{$\forall v_{453}\,\forall v_{454}\,(v_{453} = v_{454} \to v_{454} = v_{453})$}
\AXC{$\varphi_{45}$\lbl{A2}}
\BIC{$\forall v_{454}\,(y \cdot 1 = v_{454} \to v_{454} = y \cdot 1)$}
\AXC{$\varphi_{10}$\lbl{A2}}
\BIC{$y \cdot 1 = y \to y = y \cdot 1$}
\end{ptree}
\noindent$\Pi_{35}$:
\begin{ptree}
\AXC{$\forall x \in \mathbb{R}\,(x \cdot 1 = x)$\lbl{(R8)}}
\AXC{$\forall x \in \mathbb{R}\,(x \cdot 1 = x) \to y \in \mathbb{R} \to y \cdot 1 = y$\lbl{A2}}
\BIC{$y \in \mathbb{R} \to y \cdot 1 = y$}
\AXC{$(y \in \mathbb{R} \to y \cdot 1 = y) \to \varphi_{88}$\lbl{A1}}
\BIC{$\varphi_{88}$}
\end{ptree}
\noindent$\Pi_{36}$:
\begin{ptree}
\AXC{$x \in \mathbb{R} \to y \in \mathbb{R} \to x \neq 0 \to x \cdot y = 1 \to y \in \mathbb{R}$\lbl{A1}}
\AXC{$\stackrel{\Pi_{35}}{\vdots}$}\noLine
\UIC{$\varphi_{88}$}
\BIC{$x \in \mathbb{R} \to y \in \mathbb{R} \to x \neq 0 \to x \cdot y = 1 \to y \cdot 1 = y$}
\AXC{$(x \in \mathbb{R} \to y \in \mathbb{R} \to x \neq 0 \to x \cdot y = 1 \to y \cdot 1 = y) \to \varphi_{51}$\lbl{A1}}
\BIC{$\varphi_{51}$}
\end{ptree}
\noindent$\Pi_{37}$:
\begin{ptree}
\AXC{$\stackrel{\Pi_{28}}{\vdots}$}\noLine
\UIC{$\varphi_{17}$}
\AXC{$\stackrel{\Pi_{33}}{\vdots}$}\noLine
\UIC{$\varphi_{44}$}
\AXC{$\stackrel{\Pi_{34}}{\vdots}$}\noLine
\UIC{$y \cdot 1 = y \to y = y \cdot 1$}
\AXC{$\stackrel{\Pi_{36}}{\vdots}$}\noLine
\UIC{$\varphi_{51}$}
\BIC{$x \in \mathbb{R} \to y \in \mathbb{R} \to x \neq 0 \to x \cdot y = 1 \to y = y \cdot 1$}
\AXC{$\varphi_{111}$\lbl{A1}}
\BIC{$\varphi_{44} \to \varphi_{17} \to \varphi_{29}$}
\BIC{$\varphi_{17} \to \varphi_{29}$}
\BIC{$\varphi_{29}$}
\AXC{$\varphi_{29} \to (x \in \mathbb{R} \to \varphi_{33}) \to \varphi_{98}$\lbl{A1}}
\BIC{$(x \in \mathbb{R} \to \varphi_{33}) \to \varphi_{98}$}
\end{ptree}
\noindent$\Pi_{38}$:
\begin{ptree}
\AXC{$\stackrel{\Pi_{20}}{\vdots}$}\noLine
\UIC{$\forall v_{465}\,(\varphi_{7})$}
\AXC{$\varphi_{75}$\lbl{A2}}
\BIC{$\forall v_{466}\,\forall v_{467}\,(y = v_{466} \to v_{466} = v_{467} \to y = v_{467})$}
\AXC{$\varphi_{84}$\lbl{A2}}
\BIC{$\varphi_{20}$}
\AXC{$\varphi_{20} \to \varphi_{54}$\lbl{A2}}
\BIC{$\varphi_{54}$}
\AXC{$\stackrel{\Pi_{26}}{\vdots}$}\noLine
\UIC{$x \in \mathbb{R} \to \varphi_{33}$}
\AXC{$\stackrel{\Pi_{37}}{\vdots}$}\noLine
\UIC{$(x \in \mathbb{R} \to \varphi_{33}) \to \varphi_{98}$}
\BIC{$\varphi_{98}$}
\BIC{$x \in \mathbb{R} \to y \in \mathbb{R} \to x \neq 0 \to x \cdot y = 1 \to y = y \cdot x \cdot x^{-1}$}
\end{ptree}
\noindent$\Pi_{39}$:
\begin{ptree}
\AXC{$\stackrel{\Pi_{14}}{\vdots}$}\noLine
\UIC{$\varphi_{53}$}
\AXC{$\stackrel{\Pi_{19}}{\vdots}$}\noLine
\UIC{$\varphi_{56}$}
\AXC{$\stackrel{\Pi_{38}}{\vdots}$}\noLine
\UIC{$x \in \mathbb{R} \to y \in \mathbb{R} \to x \neq 0 \to x \cdot y = 1 \to y = y \cdot x \cdot x^{-1}$}
\AXC{$\varphi_{116}$\lbl{A1}}
\BIC{$\varphi_{56} \to \varphi_{99}$}
\BIC{$\varphi_{99}$}
\BIC{$x \in \mathbb{R} \to y \in \mathbb{R} \to x \neq 0 \to x \cdot y = 1 \to y = x \cdot y \cdot x^{-1}$}
\AXC{$\varphi_{115}$\lbl{A1}}
\BIC{$\varphi_{103}$}
\end{ptree}
\noindent$\Pi_{40}$:
\begin{ptree}
\AXC{$\varphi_{35}$\lbl{A6}}
\UIC{$\forall w\,(\varphi_{35})$}
\UIC{$\forall z\,\forall w\,(\varphi_{35})$}
\AXC{$\forall z\,\forall w\,(\varphi_{35}) \to \varphi_{60}$\lbl{A2}}
\BIC{$\varphi_{60}$}
\AXC{$\varphi_{60} \to \varphi_{52}$\lbl{A2}}
\BIC{$\varphi_{52}$}
\AXC{$\stackrel{\Pi_{12}}{\vdots}$}\noLine
\UIC{$\varphi_{52} \to \varphi_{49}$}
\BIC{$\varphi_{49}$}
\AXC{$\stackrel{\Pi_{39}}{\vdots}$}\noLine
\UIC{$\varphi_{103}$}
\BIC{$\varphi_{34} \to x \in \mathbb{R} \to y \in \mathbb{R} \to x \neq 0 \to x \cdot y = 1 \to y = 1 \cdot x^{-1}$}
\end{ptree}
\noindent$\Pi_{41}$:
\begin{ptree}
\AXC{$\stackrel{\Pi_{9}}{\vdots}$}\noLine
\UIC{$\varphi_{43}$}
\AXC{$\stackrel{\Pi_{11}}{\vdots}$}\noLine
\UIC{$\varphi_{34}$}
\AXC{$\stackrel{\Pi_{40}}{\vdots}$}\noLine
\UIC{$\varphi_{34} \to x \in \mathbb{R} \to y \in \mathbb{R} \to x \neq 0 \to x \cdot y = 1 \to y = 1 \cdot x^{-1}$}
\BIC{$x \in \mathbb{R} \to y \in \mathbb{R} \to x \neq 0 \to x \cdot y = 1 \to y = 1 \cdot x^{-1}$}
\AXC{$\varphi_{112}$\lbl{A1}}
\BIC{$\varphi_{43} \to \varphi_{92}$}
\BIC{$\varphi_{92}$}
\end{ptree}
\noindent$\Pi_{42}$:
\begin{ptree}
\AXC{$\stackrel{\Pi_{7}}{\vdots}$}\noLine
\UIC{$\forall v_{474}\,(\varphi_{4})$}
\AXC{$\varphi_{72}$\lbl{A2}}
\BIC{$\forall v_{475}\,\forall v_{476}\,(y = v_{475} \to v_{475} = v_{476} \to y = v_{476})$}
\AXC{$\varphi_{67}$\lbl{A2}}
\BIC{$\forall v_{476}\,(y = 1 \cdot x^{-1} \to 1 \cdot x^{-1} = v_{476} \to y = v_{476})$}
\AXC{$\varphi_{86}$\lbl{A2}}
\BIC{$y = 1 \cdot x^{-1} \to 1 \cdot x^{-1} = x^{-1} \cdot 1 \to y = x^{-1} \cdot 1$}
\AXC{$\stackrel{\Pi_{41}}{\vdots}$}\noLine
\UIC{$\varphi_{92}$}
\BIC{$x \in \mathbb{R} \to y \in \mathbb{R} \to x \neq 0 \to x \cdot y = 1 \to y = x^{-1} \cdot 1$}
\AXC{$\varphi_{109}$\lbl{A1}}
\BIC{$\varphi_{102}$}
\end{ptree}
\noindent 主木:
\begin{ptree}
\AXC{$\stackrel{\Pi_{2}}{\vdots}$}\noLine
\UIC{$y = x^{-1} \cdot 1 \to x^{-1} \cdot 1 = x^{-1} \to y = x^{-1}$}
\AXC{$\stackrel{\Pi_{6}}{\vdots}$}\noLine
\UIC{$x \in \mathbb{R} \to y \in \mathbb{R} \to x \neq 0 \to x \cdot y = 1 \to x^{-1} \cdot 1 = x^{-1}$}
\AXC{$\stackrel{\Pi_{42}}{\vdots}$}\noLine
\UIC{$\varphi_{102}$}
\BIC{$\varphi_{85}$}
\BIC{$x \in \mathbb{R} \to y \in \mathbb{R} \to x \neq 0 \to x \cdot y = 1 \to y = x^{-1}$}
\UIC{$\forall y\,(x \in \mathbb{R} \to y \in \mathbb{R} \to x \neq 0 \to x \cdot y = 1 \to y = x^{-1})$}
\UIC{$\varphi_{22}$}
\end{ptree}
\begin{abbrevlist}
\item $\varphi_{1} :\equiv \forall y\,(v_{458} = y \to y = v_{458}) \to v_{458} = v_{459} \to v_{459} = v_{458}$
\item $\varphi_{2} :\equiv \forall y\,(v_{453} = y \to y = v_{453}) \to v_{453} = v_{454} \to v_{454} = v_{453}$
\item $\varphi_{3} :\equiv \forall v_{478}\,\forall v_{479}\,(v_{477} = v_{478} \to v_{478} = v_{479} \to v_{477} = v_{479})$
\item $\varphi_{4} :\equiv \forall v_{475}\,\forall v_{476}\,(v_{474} = v_{475} \to v_{475} = v_{476} \to v_{474} = v_{476})$
\item $\varphi_{5} :\equiv \forall v_{472}\,\forall v_{473}\,(v_{471} = v_{472} \to v_{472} = v_{473} \to v_{471} = v_{473})$
\item $\varphi_{6} :\equiv \forall v_{469}\,\forall v_{470}\,(v_{468} = v_{469} \to v_{469} = v_{470} \to v_{468} = v_{470})$
\item $\varphi_{7} :\equiv \forall v_{466}\,\forall v_{467}\,(v_{465} = v_{466} \to v_{466} = v_{467} \to v_{465} = v_{467})$
\item $\varphi_{8} :\equiv \forall v_{463}\,\forall v_{464}\,(v_{462} = v_{463} \to v_{463} = v_{464} \to v_{462} = v_{464})$
\item $\varphi_{9} :\equiv \forall x\,\forall y\,(x = y \to y = x) \to \forall y\,(x \cdot x^{-1} = y \to y = x \cdot x^{-1})$
\item $\varphi_{10} :\equiv \forall v_{454}\,(y \cdot 1 = v_{454} \to v_{454} = y \cdot 1) \to y \cdot 1 = y \to y = y \cdot 1$
\item $\varphi_{11} :\equiv \forall x\,\forall y\,\forall z\,(x = y \to y = z \to x = z) \to \forall y\,\forall z\,(v_{477} = y \to y = z \to v_{477} = z)$
\item $\varphi_{12} :\equiv \forall x\,\forall y\,\forall z\,(x = y \to y = z \to x = z) \to \forall y\,\forall z\,(v_{474} = y \to y = z \to v_{474} = z)$
\item $\varphi_{13} :\equiv \forall x\,\forall y\,\forall z\,(x = y \to y = z \to x = z) \to \forall y\,\forall z\,(v_{471} = y \to y = z \to v_{471} = z)$
\item $\varphi_{14} :\equiv \forall x\,\forall y\,\forall z\,(x = y \to y = z \to x = z) \to \forall y\,\forall z\,(v_{468} = y \to y = z \to v_{468} = z)$
\item $\varphi_{15} :\equiv \forall x\,\forall y\,\forall z\,(x = y \to y = z \to x = z) \to \forall y\,\forall z\,(v_{465} = y \to y = z \to v_{465} = z)$
\item $\varphi_{16} :\equiv \forall x\,\forall y\,\forall z\,(x = y \to y = z \to x = z) \to \forall y\,\forall z\,(v_{462} = y \to y = z \to v_{462} = z)$
\item $\varphi_{17} :\equiv y = y \cdot 1 \to y \cdot 1 = y \cdot (x \cdot x^{-1}) \to y = y \cdot (x \cdot x^{-1})$
\item $\varphi_{18} :\equiv \forall v_{473}\,(y = x \cdot y \cdot x^{-1} \to x \cdot y \cdot x^{-1} = v_{473} \to y = v_{473})$
\item $\varphi_{19} :\equiv \forall v_{470}\,(y = y \cdot x \cdot x^{-1} \to y \cdot x \cdot x^{-1} = v_{470} \to y = v_{470})$
\item $\varphi_{20} :\equiv \forall v_{467}\,(y = y \cdot (x \cdot x^{-1}) \to y \cdot (x \cdot x^{-1}) = v_{467} \to y = v_{467})$
\item $\varphi_{21} :\equiv v_{457} \in \mathbb{R} \to v_{455} \cdot v_{456} \cdot v_{457} = v_{455} \cdot (v_{456} \cdot v_{457})$
\item $\varphi_{22} :\equiv \forall x\,\forall y\,(x \in \mathbb{R} \to y \in \mathbb{R} \to x \neq 0 \to x \cdot y = 1 \to y = x^{-1})$
\item $\varphi_{23} :\equiv \forall y\,\forall z\,(v_{477} = y \to y = z \to v_{477} = z) \to \forall z\,(v_{477} = v_{478} \to v_{478} = z \to v_{477} = z)$
\item $\varphi_{24} :\equiv \forall y\,\forall z\,(v_{474} = y \to y = z \to v_{474} = z) \to \forall z\,(v_{474} = v_{475} \to v_{475} = z \to v_{474} = z)$
\item $\varphi_{25} :\equiv \forall y\,\forall z\,(v_{471} = y \to y = z \to v_{471} = z) \to \forall z\,(v_{471} = v_{472} \to v_{472} = z \to v_{471} = z)$
\item $\varphi_{26} :\equiv \forall y\,\forall z\,(v_{468} = y \to y = z \to v_{468} = z) \to \forall z\,(v_{468} = v_{469} \to v_{469} = z \to v_{468} = z)$
\item $\varphi_{27} :\equiv \forall y\,\forall z\,(v_{465} = y \to y = z \to v_{465} = z) \to \forall z\,(v_{465} = v_{466} \to v_{466} = z \to v_{465} = z)$
\item $\varphi_{28} :\equiv \forall y\,\forall z\,(v_{462} = y \to y = z \to v_{462} = z) \to \forall z\,(v_{462} = v_{463} \to v_{463} = z \to v_{462} = z)$
\item $\varphi_{29} :\equiv x \in \mathbb{R} \to y \in \mathbb{R} \to x \neq 0 \to x \cdot y = 1 \to y = y \cdot (x \cdot x^{-1})$
\item $\varphi_{30} :\equiv \forall v_{461}\,(v_{460} \in \mathbb{R} \to v_{461} \in \mathbb{R} \to v_{460} \cdot v_{461} = v_{461} \cdot v_{460})$
\item $\varphi_{31} :\equiv \forall y\,(x \cdot x^{-1} = y \to y = x \cdot x^{-1}) \to x \cdot x^{-1} = 1 \to 1 = x \cdot x^{-1}$
\item $\varphi_{32} :\equiv y \in \mathbb{R} \to x \neq 0 \to x \cdot y = 1 \to y \cdot x \cdot x^{-1} = y \cdot (x \cdot x^{-1})$
\item $\varphi_{33} :\equiv y \in \mathbb{R} \to x \neq 0 \to x \cdot y = 1 \to y \cdot (x \cdot x^{-1}) = y \cdot x \cdot x^{-1}$
\item $\varphi_{34} :\equiv y = x \cdot y \cdot x^{-1} \to x \cdot y \cdot x^{-1} = 1 \cdot x^{-1} \to y = 1 \cdot x^{-1}$
\item $\varphi_{35} :\equiv z = w \to x \cdot y \cdot x^{-1} = z \cdot x^{-1} \to x \cdot y \cdot x^{-1} = w \cdot x^{-1}$
\item $\varphi_{36} :\equiv z = w \to y \cdot x \cdot x^{-1} = z \cdot x^{-1} \to y \cdot x \cdot x^{-1} = w \cdot x^{-1}$
\item $\varphi_{37} :\equiv \forall z\,(v_{477} = v_{478} \to v_{478} = z \to v_{477} = z) \to v_{477} = v_{478} \to v_{478} = v_{479} \to v_{477} = v_{479}$
\item $\varphi_{38} :\equiv \forall z\,(v_{474} = v_{475} \to v_{475} = z \to v_{474} = z) \to v_{474} = v_{475} \to v_{475} = v_{476} \to v_{474} = v_{476}$
\item $\varphi_{39} :\equiv \forall z\,(v_{471} = v_{472} \to v_{472} = z \to v_{471} = z) \to v_{471} = v_{472} \to v_{472} = v_{473} \to v_{471} = v_{473}$
\item $\varphi_{40} :\equiv \forall z\,(v_{468} = v_{469} \to v_{469} = z \to v_{468} = z) \to v_{468} = v_{469} \to v_{469} = v_{470} \to v_{468} = v_{470}$
\item $\varphi_{41} :\equiv \forall z\,(v_{465} = v_{466} \to v_{466} = z \to v_{465} = z) \to v_{465} = v_{466} \to v_{466} = v_{467} \to v_{465} = v_{467}$
\item $\varphi_{42} :\equiv \forall z\,(v_{462} = v_{463} \to v_{463} = z \to v_{462} = z) \to v_{462} = v_{463} \to v_{463} = v_{464} \to v_{462} = v_{464}$
\item $\varphi_{43} :\equiv x \in \mathbb{R} \to y \in \mathbb{R} \to x \neq 0 \to x \cdot y = 1 \to 1 \cdot x^{-1} = x^{-1} \cdot 1$
\item $\varphi_{44} :\equiv x \in \mathbb{R} \to y \in \mathbb{R} \to x \neq 0 \to x \cdot y = 1 \to y \cdot 1 = y \cdot (x \cdot x^{-1})$
\item $\varphi_{45} :\equiv \forall v_{453}\,\forall v_{454}\,(v_{453} = v_{454} \to v_{454} = v_{453}) \to \forall v_{454}\,(y \cdot 1 = v_{454} \to v_{454} = y \cdot 1)$
\item $\varphi_{46} :\equiv y \cdot x \cdot x^{-1} = y \cdot (x \cdot x^{-1}) \to y \cdot (x \cdot x^{-1}) = y \cdot x \cdot x^{-1}$
\item $\varphi_{47} :\equiv y \cdot x \cdot x^{-1} = y \cdot x \cdot x^{-1} \to y \cdot x \cdot x^{-1} = x \cdot y \cdot x^{-1}$
\item $\varphi_{48} :\equiv \forall x\,\forall y\,\forall z\,(x \in \mathbb{R} \to y \in \mathbb{R} \to z \in \mathbb{R} \to x \cdot y \cdot z = x \cdot (y \cdot z))$
\item $\varphi_{49} :\equiv x \in \mathbb{R} \to y \in \mathbb{R} \to x \neq 0 \to x \cdot y = 1 \to x \cdot y \cdot x^{-1} = 1 \cdot x^{-1}$
\item $\varphi_{50} :\equiv \forall x \in \mathbb{R}\,(x \neq 0 \to x \cdot x^{-1} = 1) \to x \in \mathbb{R} \to x \neq 0 \to x \cdot x^{-1} = 1$
\item $\varphi_{51} :\equiv (y \cdot 1 = y \to y = y \cdot 1) \to x \in \mathbb{R} \to y \in \mathbb{R} \to x \neq 0 \to x \cdot y = 1 \to y = y \cdot 1$
\item $\varphi_{52} :\equiv x \cdot y = 1 \to x \cdot y \cdot x^{-1} = x \cdot y \cdot x^{-1} \to x \cdot y \cdot x^{-1} = 1 \cdot x^{-1}$
\item $\varphi_{53} :\equiv y = y \cdot x \cdot x^{-1} \to y \cdot x \cdot x^{-1} = x \cdot y \cdot x^{-1} \to y = x \cdot y \cdot x^{-1}$
\item $\varphi_{54} :\equiv y = y \cdot (x \cdot x^{-1}) \to y \cdot (x \cdot x^{-1}) = y \cdot x \cdot x^{-1} \to y = y \cdot x \cdot x^{-1}$
\item $\varphi_{55} :\equiv \forall x \in \mathbb{R}\,(x \neq 0 \to x^{-1} \in \mathbb{R}) \to x \in \mathbb{R} \to x \neq 0 \to x^{-1} \in \mathbb{R}$
\item $\varphi_{56} :\equiv x \in \mathbb{R} \to y \in \mathbb{R} \to x \neq 0 \to x \cdot y = 1 \to y \cdot x \cdot x^{-1} = x \cdot y \cdot x^{-1}$
\item $\varphi_{57} :\equiv \forall y\,\forall z\,(v_{455} \in \mathbb{R} \to y \in \mathbb{R} \to z \in \mathbb{R} \to v_{455} \cdot y \cdot z = v_{455} \cdot (y \cdot z))$
\item $\varphi_{58} :\equiv \forall v_{457}\,(y \in \mathbb{R} \to x \in \mathbb{R} \to v_{457} \in \mathbb{R} \to y \cdot x \cdot v_{457} = y \cdot (x \cdot v_{457}))$
\item $\varphi_{59} :\equiv y \in \mathbb{R} \to x \in \mathbb{R} \to x^{-1} \in \mathbb{R} \to y \cdot x \cdot x^{-1} = y \cdot (x \cdot x^{-1})$
\item $\varphi_{60} :\equiv \forall w\,(x \cdot y = w \to x \cdot y \cdot x^{-1} = x \cdot y \cdot x^{-1} \to x \cdot y \cdot x^{-1} = w \cdot x^{-1})$
\item $\varphi_{61} :\equiv \forall w\,(y \cdot x = w \to y \cdot x \cdot x^{-1} = y \cdot x \cdot x^{-1} \to y \cdot x \cdot x^{-1} = w \cdot x^{-1})$
\item $\varphi_{62} :\equiv \forall v_{463}\,\forall v_{464}\,(y = v_{463} \to v_{463} = v_{464} \to y = v_{464}) \to \forall v_{464}\,(y = y \cdot 1 \to y \cdot 1 = v_{464} \to y = v_{464})$
\item $\varphi_{63} :\equiv \forall z\,\forall w\,(z = w \to y \cdot 1 = y \cdot z \to y \cdot 1 = y \cdot w) \to \forall w\,(1 = w \to y \cdot 1 = y \cdot 1 \to y \cdot 1 = y \cdot w)$
\item $\varphi_{64} :\equiv \forall z\,(v_{455} \in \mathbb{R} \to v_{456} \in \mathbb{R} \to z \in \mathbb{R} \to v_{455} \cdot v_{456} \cdot z = v_{455} \cdot (v_{456} \cdot z))$
\item $\varphi_{65} :\equiv \forall v_{458}\,\forall v_{459}\,(v_{458} = v_{459} \to v_{459} = v_{458}) \to \forall v_{459}\,(y \cdot x \cdot x^{-1} = v_{459} \to v_{459} = y \cdot x \cdot x^{-1})$
\item $\varphi_{66} :\equiv \forall v_{478}\,\forall v_{479}\,(y = v_{478} \to v_{478} = v_{479} \to y = v_{479}) \to \forall v_{479}\,(y = x^{-1} \cdot 1 \to x^{-1} \cdot 1 = v_{479} \to y = v_{479})$
\item $\varphi_{67} :\equiv \forall v_{475}\,\forall v_{476}\,(y = v_{475} \to v_{475} = v_{476} \to y = v_{476}) \to \forall v_{476}\,(y = 1 \cdot x^{-1} \to 1 \cdot x^{-1} = v_{476} \to y = v_{476})$
\item $\varphi_{68} :\equiv (x \cdot x^{-1} = 1 \to 1 = x \cdot x^{-1}) \to x \in \mathbb{R} \to y \in \mathbb{R} \to x \neq 0 \to x \cdot y = 1 \to 1 = x \cdot x^{-1}$
\item $\varphi_{69} :\equiv \forall x\,\forall y\,(x \in \mathbb{R} \to y \in \mathbb{R} \to x \cdot y = y \cdot x) \to \forall y\,(1 \in \mathbb{R} \to y \in \mathbb{R} \to 1 \cdot y = y \cdot 1)$
\item $\varphi_{70} :\equiv \forall v_{456}\,\forall v_{457}\,(y \in \mathbb{R} \to v_{456} \in \mathbb{R} \to v_{457} \in \mathbb{R} \to y \cdot v_{456} \cdot v_{457} = y \cdot (v_{456} \cdot v_{457}))$
\item $\varphi_{71} :\equiv \forall v_{477}\,(\varphi_{3}) \to \forall v_{478}\,\forall v_{479}\,(y = v_{478} \to v_{478} = v_{479} \to y = v_{479})$
\item $\varphi_{72} :\equiv \forall v_{474}\,(\varphi_{4}) \to \forall v_{475}\,\forall v_{476}\,(y = v_{475} \to v_{475} = v_{476} \to y = v_{476})$
\item $\varphi_{73} :\equiv \forall v_{471}\,(\varphi_{5}) \to \forall v_{472}\,\forall v_{473}\,(y = v_{472} \to v_{472} = v_{473} \to y = v_{473})$
\item $\varphi_{74} :\equiv \forall v_{468}\,(\varphi_{6}) \to \forall v_{469}\,\forall v_{470}\,(y = v_{469} \to v_{469} = v_{470} \to y = v_{470})$
\item $\varphi_{75} :\equiv \forall v_{465}\,(\varphi_{7}) \to \forall v_{466}\,\forall v_{467}\,(y = v_{466} \to v_{466} = v_{467} \to y = v_{467})$
\item $\varphi_{76} :\equiv \forall v_{462}\,(\varphi_{8}) \to \forall v_{463}\,\forall v_{464}\,(y = v_{463} \to v_{463} = v_{464} \to y = v_{464})$
\item $\varphi_{77} :\equiv \forall v_{479}\,(y = x^{-1} \cdot 1 \to x^{-1} \cdot 1 = v_{479} \to y = v_{479}) \to y = x^{-1} \cdot 1 \to x^{-1} \cdot 1 = x^{-1} \to y = x^{-1}$
\item $\varphi_{78} :\equiv \forall v_{461}\,(y \in \mathbb{R} \to v_{461} \in \mathbb{R} \to y \cdot v_{461} = v_{461} \cdot y) \to y \in \mathbb{R} \to x \in \mathbb{R} \to y \cdot x = x \cdot y$
\item $\varphi_{79} :\equiv \forall w\,(1 = w \to y \cdot 1 = y \cdot 1 \to y \cdot 1 = y \cdot w) \to 1 = x \cdot x^{-1} \to y \cdot 1 = y \cdot 1 \to y \cdot 1 = y \cdot (x \cdot x^{-1})$
\item $\varphi_{80} :\equiv \forall y\,(1 \in \mathbb{R} \to y \in \mathbb{R} \to 1 \cdot y = y \cdot 1) \to 1 \in \mathbb{R} \to x^{-1} \in \mathbb{R} \to 1 \cdot x^{-1} = x^{-1} \cdot 1$
\item $\varphi_{81} :\equiv \forall v_{472}\,\forall v_{473}\,(y = v_{472} \to v_{472} = v_{473} \to y = v_{473}) \to \varphi_{18}$
\item $\varphi_{82} :\equiv \forall v_{469}\,\forall v_{470}\,(y = v_{469} \to v_{469} = v_{470} \to y = v_{470}) \to \varphi_{19}$
\item $\varphi_{83} :\equiv \forall x\,\forall y\,(x \in \mathbb{R} \to y \in \mathbb{R} \to x \cdot y = y \cdot x) \to \forall y\,(v_{460} \in \mathbb{R} \to y \in \mathbb{R} \to v_{460} \cdot y = y \cdot v_{460})$
\item $\varphi_{84} :\equiv \forall v_{466}\,\forall v_{467}\,(y = v_{466} \to v_{466} = v_{467} \to y = v_{467}) \to \varphi_{20}$
\item $\varphi_{85} :\equiv (y = x^{-1} \cdot 1 \to x^{-1} \cdot 1 = x^{-1} \to y = x^{-1}) \to x \in \mathbb{R} \to y \in \mathbb{R} \to x \neq 0 \to x \cdot y = 1 \to y = x^{-1}$
\item $\varphi_{86} :\equiv \forall v_{476}\,(y = 1 \cdot x^{-1} \to 1 \cdot x^{-1} = v_{476} \to y = v_{476}) \to y = 1 \cdot x^{-1} \to 1 \cdot x^{-1} = x^{-1} \cdot 1 \to y = x^{-1} \cdot 1$
\item $\varphi_{87} :\equiv (x \in \mathbb{R} \to y \in \mathbb{R} \to x \neq 0 \to x \cdot y = 1 \to x \neq 0) \to x \in \mathbb{R} \to y \in \mathbb{R} \to x \neq 0 \to x \cdot y = 1 \to x \cdot x^{-1} = 1$
\item $\varphi_{88} :\equiv (x \in \mathbb{R} \to y \in \mathbb{R} \to x \neq 0 \to x \cdot y = 1 \to y \in \mathbb{R}) \to x \in \mathbb{R} \to y \in \mathbb{R} \to x \neq 0 \to x \cdot y = 1 \to y \cdot 1 = y$
\item $\varphi_{89} :\equiv (x \in \mathbb{R} \to y \in \mathbb{R} \to x \neq 0 \to x \cdot y = 1 \to x \neq 0) \to x \in \mathbb{R} \to y \in \mathbb{R} \to x \neq 0 \to x \cdot y = 1 \to x^{-1} \in \mathbb{R}$
\item $\varphi_{90} :\equiv \forall y\,(v_{460} \in \mathbb{R} \to y \in \mathbb{R} \to v_{460} \cdot y = y \cdot v_{460}) \to v_{460} \in \mathbb{R} \to v_{461} \in \mathbb{R} \to v_{460} \cdot v_{461} = v_{461} \cdot v_{460}$
\item $\varphi_{91} :\equiv (x \in \mathbb{R} \to y \in \mathbb{R} \to x \neq 0 \to x \cdot y = 1 \to x \in \mathbb{R}) \to x \in \mathbb{R} \to y \in \mathbb{R} \to x \neq 0 \to x \cdot y = 1 \to y \cdot x = x \cdot y$
\item $\varphi_{92} :\equiv (y = 1 \cdot x^{-1} \to 1 \cdot x^{-1} = x^{-1} \cdot 1 \to y = x^{-1} \cdot 1) \to x \in \mathbb{R} \to y \in \mathbb{R} \to x \neq 0 \to x \cdot y = 1 \to y = x^{-1} \cdot 1$
\item $\varphi_{93} :\equiv \forall v_{459}\,(y \cdot x \cdot x^{-1} = v_{459} \to v_{459} = y \cdot x \cdot x^{-1}) \to \varphi_{46}$
\item $\varphi_{94} :\equiv (x \in \mathbb{R} \to y \in \mathbb{R} \to x \neq 0 \to x \cdot y = 1 \to x^{-1} \in \mathbb{R}) \to x \in \mathbb{R} \to y \in \mathbb{R} \to x \neq 0 \to x \cdot y = 1 \to x^{-1} \cdot 1 = x^{-1}$
\item $\varphi_{95} :\equiv (1 = x \cdot x^{-1} \to y \cdot 1 = y \cdot 1 \to y \cdot 1 = y \cdot (x \cdot x^{-1})) \to \varphi_{44}$
\item $\varphi_{96} :\equiv \forall v_{460}\,(\varphi_{30}) \to \forall v_{461}\,(y \in \mathbb{R} \to v_{461} \in \mathbb{R} \to y \cdot v_{461} = v_{461} \cdot y)$
\item $\varphi_{97} :\equiv (x \in \mathbb{R} \to y \in \mathbb{R} \to x \neq 0 \to x \cdot y = 1 \to x^{-1} \in \mathbb{R}) \to x \in \mathbb{R} \to \varphi_{32}$
\item $\varphi_{98} :\equiv \varphi_{54} \to x \in \mathbb{R} \to y \in \mathbb{R} \to x \neq 0 \to x \cdot y = 1 \to y = y \cdot x \cdot x^{-1}$
\item $\varphi_{99} :\equiv \varphi_{53} \to x \in \mathbb{R} \to y \in \mathbb{R} \to x \neq 0 \to x \cdot y = 1 \to y = x \cdot y \cdot x^{-1}$
\item $\varphi_{100} :\equiv 1 \in \mathbb{R} \to (x \in \mathbb{R} \to y \in \mathbb{R} \to x \neq 0 \to x \cdot y = 1 \to x^{-1} \in \mathbb{R}) \to \varphi_{43}$
\item $\varphi_{101} :\equiv (x \in \mathbb{R} \to y \in \mathbb{R} \to x \neq 0 \to x \cdot y = 1 \to x \cdot x^{-1} = 1) \to \varphi_{68}$
\item $\varphi_{102} :\equiv (x \in \mathbb{R} \to y \in \mathbb{R} \to x \neq 0 \to x \cdot y = 1 \to x^{-1} \cdot 1 = x^{-1}) \to \varphi_{85}$
\item $\varphi_{103} :\equiv \varphi_{49} \to \varphi_{34} \to x \in \mathbb{R} \to y \in \mathbb{R} \to x \neq 0 \to x \cdot y = 1 \to y = 1 \cdot x^{-1}$
\item $\varphi_{104} :\equiv (x \in \mathbb{R} \to x \neq 0 \to x \cdot x^{-1} = 1) \to (x \in \mathbb{R} \to y \in \mathbb{R} \to x \neq 0 \to x \cdot y = 1 \to x \in \mathbb{R}) \to \varphi_{87}$
\item $\varphi_{105} :\equiv (x \in \mathbb{R} \to y \in \mathbb{R} \to x \neq 0 \to x \cdot y = 1 \to 1 = x \cdot x^{-1}) \to y \cdot 1 = y \cdot 1 \to \varphi_{95}$
\item $\varphi_{106} :\equiv \forall v_{455}\,\forall v_{456}\,\forall v_{457}\,(v_{455} \in \mathbb{R} \to v_{456} \in \mathbb{R} \to \varphi_{21}) \to \varphi_{70}$
\item $\varphi_{107} :\equiv (x \in \mathbb{R} \to x \neq 0 \to x^{-1} \in \mathbb{R}) \to (x \in \mathbb{R} \to y \in \mathbb{R} \to x \neq 0 \to x \cdot y = 1 \to x \in \mathbb{R}) \to \varphi_{89}$
\item $\varphi_{108} :\equiv (y \in \mathbb{R} \to x \in \mathbb{R} \to y \cdot x = x \cdot y) \to (x \in \mathbb{R} \to y \in \mathbb{R} \to x \neq 0 \to x \cdot y = 1 \to y \in \mathbb{R}) \to \varphi_{91}$
\item $\varphi_{109} :\equiv (x \in \mathbb{R} \to y \in \mathbb{R} \to x \neq 0 \to x \cdot y = 1 \to y = x^{-1} \cdot 1) \to \varphi_{102}$
\item $\varphi_{110} :\equiv (x \in \mathbb{R} \to y \in \mathbb{R} \to x \neq 0 \to x \cdot y = 1 \to x \cdot y = 1) \to x \cdot y \cdot x^{-1} = x \cdot y \cdot x^{-1} \to \varphi_{52} \to \varphi_{49}$
\item $\varphi_{111} :\equiv (x \in \mathbb{R} \to y \in \mathbb{R} \to x \neq 0 \to x \cdot y = 1 \to y = y \cdot 1) \to \varphi_{44} \to \varphi_{17} \to \varphi_{29}$
\item $\varphi_{112} :\equiv (x \in \mathbb{R} \to y \in \mathbb{R} \to x \neq 0 \to x \cdot y = 1 \to y = 1 \cdot x^{-1}) \to \varphi_{43} \to \varphi_{92}$
\item $\varphi_{113} :\equiv (x \in \mathbb{R} \to y \in \mathbb{R} \to x \neq 0 \to x \cdot y = 1 \to y \in \mathbb{R}) \to (x \in \mathbb{R} \to y \in \mathbb{R} \to x \neq 0 \to x \cdot y = 1 \to x \in \mathbb{R}) \to \varphi_{97}$
\item $\varphi_{114} :\equiv (x \in \mathbb{R} \to y \in \mathbb{R} \to x \neq 0 \to x \cdot y = 1 \to y \cdot x = x \cdot y) \to y \cdot x \cdot x^{-1} = y \cdot x \cdot x^{-1} \to (y \cdot x = x \cdot y \to \varphi_{47}) \to \varphi_{56}$
\item $\varphi_{115} :\equiv (x \in \mathbb{R} \to y \in \mathbb{R} \to x \neq 0 \to x \cdot y = 1 \to y = x \cdot y \cdot x^{-1}) \to \varphi_{103}$
\item $\varphi_{116} :\equiv (x \in \mathbb{R} \to y \in \mathbb{R} \to x \neq 0 \to x \cdot y = 1 \to y = y \cdot x \cdot x^{-1}) \to \varphi_{56} \to \varphi_{99}$
\end{abbrevlist}
\end{fullproof}

% 自動生成 (prover/build.js): InvNe0  ノード数 88
\begin{fullproof}{[E14]}{$\forall x \in \mathbb{R}\,(x \neq 0 \to x^{-1} \neq 0)$}{88}\phantomsection\label{tree:InvNe0}
\noindent$\Pi_{1}$:
\begin{ptree}
\AXC{$x = x$\lbl{A5}}
\AXC{$x = y \to x = x \to y = x$\lbl{A6}}
\AXC{$(x = y \to x = x \to y = x) \to x = x \to x = y \to y = x$\lbl{A1}}
\BIC{$x = x \to x = y \to y = x$}
\BIC{$x = y \to y = x$}
\UIC{$\forall y\,(x = y \to y = x)$}
\UIC{$\forall x\,\forall y\,(x = y \to y = x)$}
\end{ptree}
\noindent$\Pi_{2}$:
\begin{ptree}
\AXC{$\stackrel{\Pi_{1}}{\vdots}$}\noLine
\UIC{$\forall x\,\forall y\,(x = y \to y = x)$}
\AXC{$\forall x\,\forall y\,(x = y \to y = x) \to \forall y\,(1 = y \to y = 1)$\lbl{A2}}
\BIC{$\forall y\,(1 = y \to y = 1)$}
\AXC{$\forall y\,(1 = y \to y = 1) \to 1 = 0 \to 0 = 1$\lbl{A2}}
\BIC{$1 = 0 \to 0 = 1$}
\end{ptree}
\noindent$\Pi_{3}$:
\begin{ptree}
\AXC{$\stackrel{\Pi_{2}}{\vdots}$}\noLine
\UIC{$1 = 0 \to 0 = 1$}
\AXC{$0 \neq 1$\lbl{(R8)}}
\AXC{$0 \neq 1 \to (1 = 0 \to 0 = 1) \to 1 \neq 0$\lbl{A1}}
\BIC{$(1 = 0 \to 0 = 1) \to 1 \neq 0$}
\BIC{$1 \neq 0$}
\AXC{$1 \neq 0 \to x \in \mathbb{R} \to x \neq 0 \to \neg x^{-1} \neq 0 \to 1 \neq 0$\lbl{A1}}
\BIC{$x \in \mathbb{R} \to x \neq 0 \to \neg x^{-1} \neq 0 \to 1 \neq 0$}
\end{ptree}
\noindent$\Pi_{4}$:
\begin{ptree}
\AXC{$y = z \to x = y \to x = z$\lbl{A6}}
\AXC{$(y = z \to x = y \to x = z) \to x = y \to y = z \to x = z$\lbl{A1}}
\BIC{$x = y \to y = z \to x = z$}
\UIC{$\forall z\,(x = y \to y = z \to x = z)$}
\UIC{$\forall y\,\forall z\,(x = y \to y = z \to x = z)$}
\UIC{$\forall x\,\forall y\,\forall z\,(x = y \to y = z \to x = z)$}
\AXC{$\varphi_{1}$\lbl{A2}}
\BIC{$\forall y\,\forall z\,(1 = y \to y = z \to 1 = z)$}
\AXC{$\varphi_{3}$\lbl{A2}}
\BIC{$\forall z\,(1 = x \cdot 0 \to x \cdot 0 = z \to 1 = z)$}
\AXC{$\varphi_{4}$\lbl{A2}}
\BIC{$1 = x \cdot 0 \to x \cdot 0 = 0 \to 1 = 0$}
\end{ptree}
\noindent$\Pi_{5}$:
\begin{ptree}
\AXC{$\forall x \in \mathbb{R}\,(x \cdot 0 = 0)$\lbl{[F8]}}
\AXC{$\forall x \in \mathbb{R}\,(x \cdot 0 = 0) \to x \in \mathbb{R} \to x \cdot 0 = 0$\lbl{A2}}
\BIC{$x \in \mathbb{R} \to x \cdot 0 = 0$}
\AXC{$(x \in \mathbb{R} \to x \cdot 0 = 0) \to \varphi_{14}$\lbl{A1}}
\BIC{$\varphi_{14}$}
\end{ptree}
\noindent$\Pi_{6}$:
\begin{ptree}
\AXC{$y = z \to x = y \to x = z$\lbl{A6}}
\AXC{$(y = z \to x = y \to x = z) \to x = y \to y = z \to x = z$\lbl{A1}}
\BIC{$x = y \to y = z \to x = z$}
\UIC{$\forall z\,(x = y \to y = z \to x = z)$}
\UIC{$\forall y\,\forall z\,(x = y \to y = z \to x = z)$}
\UIC{$\forall x\,\forall y\,\forall z\,(x = y \to y = z \to x = z)$}
\AXC{$\varphi_{1}$\lbl{A2}}
\BIC{$\forall y\,\forall z\,(1 = y \to y = z \to 1 = z)$}
\AXC{$\varphi_{6}$\lbl{A2}}
\BIC{$\forall z\,(1 = x \cdot x^{-1} \to x \cdot x^{-1} = z \to 1 = z)$}
\AXC{$\varphi_{13}$\lbl{A2}}
\BIC{$1 = x \cdot x^{-1} \to x \cdot x^{-1} = x \cdot 0 \to 1 = x \cdot 0$}
\end{ptree}
\noindent$\Pi_{7}$:
\begin{ptree}
\AXC{$z = w \to x \cdot x^{-1} = x \cdot z \to x \cdot x^{-1} = x \cdot w$\lbl{A6}}
\UIC{$\forall w\,(z = w \to x \cdot x^{-1} = x \cdot z \to x \cdot x^{-1} = x \cdot w)$}
\UIC{$\forall z\,\forall w\,(z = w \to x \cdot x^{-1} = x \cdot z \to x \cdot x^{-1} = x \cdot w)$}
\AXC{$\varphi_{18}$\lbl{A2}}
\BIC{$\varphi_{5}$}
\AXC{$\varphi_{5} \to x^{-1} = 0 \to x \cdot x^{-1} = x \cdot x^{-1} \to x \cdot x^{-1} = x \cdot 0$\lbl{A2}}
\BIC{$x^{-1} = 0 \to x \cdot x^{-1} = x \cdot x^{-1} \to x \cdot x^{-1} = x \cdot 0$}
\end{ptree}
\noindent$\Pi_{8}$:
\begin{ptree}
\AXC{$x \in \mathbb{R} \to x \neq 0 \to \neg x^{-1} \neq 0 \to \neg x^{-1} \neq 0$\lbl{A1}}
\AXC{$\varphi_{15}$\lbl{A1}}
\BIC{$x \in \mathbb{R} \to x \neq 0 \to \neg x^{-1} \neq 0 \to x^{-1} = 0$}
\AXC{$\varphi_{21}$\lbl{A1}}
\BIC{$x \cdot x^{-1} = x \cdot x^{-1} \to \varphi_{17}$}
\end{ptree}
\noindent$\Pi_{9}$:
\begin{ptree}
\AXC{$x = x$\lbl{A5}}
\UIC{$\forall x\,(x = x)$}
\AXC{$\forall x\,(x = x) \to x \cdot x^{-1} = x \cdot x^{-1}$\lbl{A2}}
\BIC{$x \cdot x^{-1} = x \cdot x^{-1}$}
\AXC{$\stackrel{\Pi_{8}}{\vdots}$}\noLine
\UIC{$x \cdot x^{-1} = x \cdot x^{-1} \to \varphi_{17}$}
\BIC{$\varphi_{17}$}
\end{ptree}
\noindent$\Pi_{10}$:
\begin{ptree}
\AXC{$x \in \mathbb{R} \to x \neq 0 \to \neg x^{-1} \neq 0 \to x \in \mathbb{R}$\lbl{A1}}
\AXC{$\forall x \in \mathbb{R}\,(x \neq 0 \to x \cdot x^{-1} = 1)$\lbl{(R9)}}
\AXC{$\varphi_{10}$\lbl{A2}}
\BIC{$x \in \mathbb{R} \to x \neq 0 \to x \cdot x^{-1} = 1$}
\AXC{$\varphi_{20}$\lbl{A1}}
\BIC{$(x \in \mathbb{R} \to x \neq 0 \to \neg x^{-1} \neq 0 \to x \in \mathbb{R}) \to \varphi_{11}$}
\BIC{$\varphi_{11}$}
\end{ptree}
\noindent$\Pi_{11}$:
\begin{ptree}
\AXC{$x \in \mathbb{R} \to x \neq 0 \to x \neq 0$\lbl{A1}}
\AXC{$\stackrel{\Pi_{10}}{\vdots}$}\noLine
\UIC{$\varphi_{11}$}
\BIC{$x \in \mathbb{R} \to x \neq 0 \to \neg x^{-1} \neq 0 \to x \cdot x^{-1} = 1$}
\AXC{$(x \in \mathbb{R} \to x \neq 0 \to \neg x^{-1} \neq 0 \to x \cdot x^{-1} = 1) \to \varphi_{12}$\lbl{A1}}
\BIC{$\varphi_{12}$}
\end{ptree}
\noindent$\Pi_{12}$:
\begin{ptree}
\AXC{$\stackrel{\Pi_{1}}{\vdots}$}\noLine
\UIC{$\forall x\,\forall y\,(x = y \to y = x)$}
\AXC{$\varphi_{2}$\lbl{A2}}
\BIC{$\forall y\,(x \cdot x^{-1} = y \to y = x \cdot x^{-1})$}
\AXC{$\varphi_{7}$\lbl{A2}}
\BIC{$x \cdot x^{-1} = 1 \to 1 = x \cdot x^{-1}$}
\AXC{$\stackrel{\Pi_{11}}{\vdots}$}\noLine
\UIC{$\varphi_{12}$}
\BIC{$x \in \mathbb{R} \to x \neq 0 \to \neg x^{-1} \neq 0 \to 1 = x \cdot x^{-1}$}
\AXC{$\varphi_{22}$\lbl{A1}}
\BIC{$(x \in \mathbb{R} \to x \neq 0 \to \neg x^{-1} \neq 0 \to x \cdot x^{-1} = x \cdot 0) \to \varphi_{16}$}
\end{ptree}
\noindent$\Pi_{13}$:
\begin{ptree}
\AXC{$\stackrel{\Pi_{7}}{\vdots}$}\noLine
\UIC{$x^{-1} = 0 \to x \cdot x^{-1} = x \cdot x^{-1} \to x \cdot x^{-1} = x \cdot 0$}
\AXC{$\stackrel{\Pi_{9}}{\vdots}$}\noLine
\UIC{$\varphi_{17}$}
\BIC{$x \in \mathbb{R} \to x \neq 0 \to \neg x^{-1} \neq 0 \to x \cdot x^{-1} = x \cdot 0$}
\AXC{$\stackrel{\Pi_{12}}{\vdots}$}\noLine
\UIC{$(x \in \mathbb{R} \to x \neq 0 \to \neg x^{-1} \neq 0 \to x \cdot x^{-1} = x \cdot 0) \to \varphi_{16}$}
\BIC{$\varphi_{16}$}
\end{ptree}
\noindent$\Pi_{14}$:
\begin{ptree}
\AXC{$x \in \mathbb{R} \to x \neq 0 \to \neg x^{-1} \neq 0 \to x \in \mathbb{R}$\lbl{A1}}
\AXC{$\stackrel{\Pi_{5}}{\vdots}$}\noLine
\UIC{$\varphi_{14}$}
\BIC{$x \in \mathbb{R} \to x \neq 0 \to \neg x^{-1} \neq 0 \to x \cdot 0 = 0$}
\AXC{$\stackrel{\Pi_{6}}{\vdots}$}\noLine
\UIC{$1 = x \cdot x^{-1} \to x \cdot x^{-1} = x \cdot 0 \to 1 = x \cdot 0$}
\AXC{$\stackrel{\Pi_{13}}{\vdots}$}\noLine
\UIC{$\varphi_{16}$}
\BIC{$x \in \mathbb{R} \to x \neq 0 \to \neg x^{-1} \neq 0 \to 1 = x \cdot 0$}
\AXC{$\varphi_{19}$\lbl{A1}}
\BIC{$(x \in \mathbb{R} \to x \neq 0 \to \neg x^{-1} \neq 0 \to x \cdot 0 = 0) \to \varphi_{8}$}
\BIC{$\varphi_{8}$}
\end{ptree}
\noindent$\Pi_{15}$:
\begin{ptree}
\AXC{$\stackrel{\Pi_{4}}{\vdots}$}\noLine
\UIC{$1 = x \cdot 0 \to x \cdot 0 = 0 \to 1 = 0$}
\AXC{$\stackrel{\Pi_{14}}{\vdots}$}\noLine
\UIC{$\varphi_{8}$}
\BIC{$x \in \mathbb{R} \to x \neq 0 \to \neg x^{-1} \neq 0 \to 1 = 0$}
\AXC{$(x \in \mathbb{R} \to x \neq 0 \to \neg x^{-1} \neq 0 \to 1 = 0) \to \varphi_{9}$\lbl{A1}}
\BIC{$\varphi_{9}$}
\end{ptree}
\noindent 主木:
\begin{ptree}
\AXC{$\stackrel{\Pi_{3}}{\vdots}$}\noLine
\UIC{$x \in \mathbb{R} \to x \neq 0 \to \neg x^{-1} \neq 0 \to 1 \neq 0$}
\AXC{$\stackrel{\Pi_{15}}{\vdots}$}\noLine
\UIC{$\varphi_{9}$}
\BIC{$x \in \mathbb{R} \to x \neq 0 \to x^{-1} \neq 0$}
\UIC{$\forall x \in \mathbb{R}\,(x \neq 0 \to x^{-1} \neq 0)$}
\end{ptree}
\begin{abbrevlist}
\item $\varphi_{1} :\equiv \forall x\,\forall y\,\forall z\,(x = y \to y = z \to x = z) \to \forall y\,\forall z\,(1 = y \to y = z \to 1 = z)$
\item $\varphi_{2} :\equiv \forall x\,\forall y\,(x = y \to y = x) \to \forall y\,(x \cdot x^{-1} = y \to y = x \cdot x^{-1})$
\item $\varphi_{3} :\equiv \forall y\,\forall z\,(1 = y \to y = z \to 1 = z) \to \forall z\,(1 = x \cdot 0 \to x \cdot 0 = z \to 1 = z)$
\item $\varphi_{4} :\equiv \forall z\,(1 = x \cdot 0 \to x \cdot 0 = z \to 1 = z) \to 1 = x \cdot 0 \to x \cdot 0 = 0 \to 1 = 0$
\item $\varphi_{5} :\equiv \forall w\,(x^{-1} = w \to x \cdot x^{-1} = x \cdot x^{-1} \to x \cdot x^{-1} = x \cdot w)$
\item $\varphi_{6} :\equiv \forall y\,\forall z\,(1 = y \to y = z \to 1 = z) \to \forall z\,(1 = x \cdot x^{-1} \to x \cdot x^{-1} = z \to 1 = z)$
\item $\varphi_{7} :\equiv \forall y\,(x \cdot x^{-1} = y \to y = x \cdot x^{-1}) \to x \cdot x^{-1} = 1 \to 1 = x \cdot x^{-1}$
\item $\varphi_{8} :\equiv (1 = x \cdot 0 \to x \cdot 0 = 0 \to 1 = 0) \to x \in \mathbb{R} \to x \neq 0 \to \neg x^{-1} \neq 0 \to 1 = 0$
\item $\varphi_{9} :\equiv (x \in \mathbb{R} \to x \neq 0 \to \neg x^{-1} \neq 0 \to 1 \neq 0) \to x \in \mathbb{R} \to x \neq 0 \to x^{-1} \neq 0$
\item $\varphi_{10} :\equiv \forall x \in \mathbb{R}\,(x \neq 0 \to x \cdot x^{-1} = 1) \to x \in \mathbb{R} \to x \neq 0 \to x \cdot x^{-1} = 1$
\item $\varphi_{11} :\equiv (x \in \mathbb{R} \to x \neq 0 \to x \neq 0) \to x \in \mathbb{R} \to x \neq 0 \to \neg x^{-1} \neq 0 \to x \cdot x^{-1} = 1$
\item $\varphi_{12} :\equiv (x \cdot x^{-1} = 1 \to 1 = x \cdot x^{-1}) \to x \in \mathbb{R} \to x \neq 0 \to \neg x^{-1} \neq 0 \to 1 = x \cdot x^{-1}$
\item $\varphi_{13} :\equiv \forall z\,(1 = x \cdot x^{-1} \to x \cdot x^{-1} = z \to 1 = z) \to 1 = x \cdot x^{-1} \to x \cdot x^{-1} = x \cdot 0 \to 1 = x \cdot 0$
\item $\varphi_{14} :\equiv (x \in \mathbb{R} \to x \neq 0 \to \neg x^{-1} \neq 0 \to x \in \mathbb{R}) \to x \in \mathbb{R} \to x \neq 0 \to \neg x^{-1} \neq 0 \to x \cdot 0 = 0$
\item $\varphi_{15} :\equiv (x \in \mathbb{R} \to x \neq 0 \to \neg x^{-1} \neq 0 \to \neg x^{-1} \neq 0) \to x \in \mathbb{R} \to x \neq 0 \to \neg x^{-1} \neq 0 \to x^{-1} = 0$
\item $\varphi_{16} :\equiv (1 = x \cdot x^{-1} \to x \cdot x^{-1} = x \cdot 0 \to 1 = x \cdot 0) \to x \in \mathbb{R} \to x \neq 0 \to \neg x^{-1} \neq 0 \to 1 = x \cdot 0$
\item $\varphi_{17} :\equiv (x^{-1} = 0 \to x \cdot x^{-1} = x \cdot x^{-1} \to x \cdot x^{-1} = x \cdot 0) \to x \in \mathbb{R} \to x \neq 0 \to \neg x^{-1} \neq 0 \to x \cdot x^{-1} = x \cdot 0$
\item $\varphi_{18} :\equiv \forall z\,\forall w\,(z = w \to x \cdot x^{-1} = x \cdot z \to x \cdot x^{-1} = x \cdot w) \to \varphi_{5}$
\item $\varphi_{19} :\equiv (x \in \mathbb{R} \to x \neq 0 \to \neg x^{-1} \neq 0 \to 1 = x \cdot 0) \to (x \in \mathbb{R} \to x \neq 0 \to \neg x^{-1} \neq 0 \to x \cdot 0 = 0) \to \varphi_{8}$
\item $\varphi_{20} :\equiv (x \in \mathbb{R} \to x \neq 0 \to x \cdot x^{-1} = 1) \to (x \in \mathbb{R} \to x \neq 0 \to \neg x^{-1} \neq 0 \to x \in \mathbb{R}) \to \varphi_{11}$
\item $\varphi_{21} :\equiv (x \in \mathbb{R} \to x \neq 0 \to \neg x^{-1} \neq 0 \to x^{-1} = 0) \to x \cdot x^{-1} = x \cdot x^{-1} \to \varphi_{17}$
\item $\varphi_{22} :\equiv (x \in \mathbb{R} \to x \neq 0 \to \neg x^{-1} \neq 0 \to 1 = x \cdot x^{-1}) \to (x \in \mathbb{R} \to x \neq 0 \to \neg x^{-1} \neq 0 \to x \cdot x^{-1} = x \cdot 0) \to \varphi_{16}$
\end{abbrevlist}
\end{fullproof}

% 自動生成 (prover/build.js): MulNe0  ノード数 242
\begin{fullproof}{[E15]}{$\forall x\,\forall y\,(x \in \mathbb{R} \to y \in \mathbb{R} \to x \neq 0 \to y \neq 0 \to x \cdot y \neq 0)$}{242}\phantomsection\label{tree:MulNe0}
\noindent$\Pi_{1}$:
\begin{ptree}
\AXC{$y = z \to x = y \to x = z$\lbl{A6}}
\AXC{$(y = z \to x = y \to x = z) \to x = y \to y = z \to x = z$\lbl{A1}}
\BIC{$x = y \to y = z \to x = z$}
\UIC{$\forall z\,(x = y \to y = z \to x = z)$}
\UIC{$\forall y\,\forall z\,(x = y \to y = z \to x = z)$}
\UIC{$\forall x\,\forall y\,\forall z\,(x = y \to y = z \to x = z)$}
\AXC{$\varphi_{10}$\lbl{A2}}
\BIC{$\forall y\,\forall z\,(v_{494} = y \to y = z \to v_{494} = z)$}
\AXC{$\varphi_{20}$\lbl{A2}}
\BIC{$\forall z\,(v_{494} = v_{495} \to v_{495} = z \to v_{494} = z)$}
\AXC{$\varphi_{32}$\lbl{A2}}
\BIC{$v_{494} = v_{495} \to v_{495} = v_{496} \to v_{494} = v_{496}$}
\UIC{$\forall v_{496}\,(v_{494} = v_{495} \to v_{495} = v_{496} \to v_{494} = v_{496})$}
\UIC{$\varphi_{2}$}
\UIC{$\forall v_{494}\,(\varphi_{2})$}
\end{ptree}
\noindent$\Pi_{2}$:
\begin{ptree}
\AXC{$\forall x \in \mathbb{R}\,(x \neq 0 \to x^{-1} \in \mathbb{R})$\lbl{(R1)}}
\AXC{$\varphi_{51}$\lbl{A2}}
\BIC{$x \in \mathbb{R} \to x \neq 0 \to x^{-1} \in \mathbb{R}$}
\AXC{$(x \in \mathbb{R} \to x \neq 0 \to x^{-1} \in \mathbb{R}) \to \varphi_{91}$\lbl{A1}}
\BIC{$\varphi_{91}$}
\end{ptree}
\noindent$\Pi_{3}$:
\begin{ptree}
\AXC{$x \in \mathbb{R} \to y \in \mathbb{R} \to x \neq 0 \to y \neq 0 \to x \neq 0$\lbl{A1}}
\AXC{$x \in \mathbb{R} \to y \in \mathbb{R} \to x \neq 0 \to y \neq 0 \to \neg x \cdot y \neq 0 \to x \in \mathbb{R}$\lbl{A1}}
\AXC{$\stackrel{\Pi_{2}}{\vdots}$}\noLine
\UIC{$\varphi_{91}$}
\BIC{$(x \in \mathbb{R} \to y \in \mathbb{R} \to x \neq 0 \to y \neq 0 \to x \neq 0) \to \varphi_{44}$}
\BIC{$\varphi_{44}$}
\end{ptree}
\noindent$\Pi_{4}$:
\begin{ptree}
\AXC{$\forall x \in \mathbb{R}\,(x \cdot 0 = 0)$\lbl{[F8]}}
\AXC{$\forall x \in \mathbb{R}\,(x \cdot 0 = 0) \to x^{-1} \in \mathbb{R} \to x^{-1} \cdot 0 = 0$\lbl{A2}}
\BIC{$x^{-1} \in \mathbb{R} \to x^{-1} \cdot 0 = 0$}
\AXC{$(x^{-1} \in \mathbb{R} \to x^{-1} \cdot 0 = 0) \to \varphi_{44} \to \varphi_{33}$\lbl{A1}}
\BIC{$\varphi_{44} \to \varphi_{33}$}
\end{ptree}
\noindent$\Pi_{5}$:
\begin{ptree}
\AXC{$y = z \to x = y \to x = z$\lbl{A6}}
\AXC{$(y = z \to x = y \to x = z) \to x = y \to y = z \to x = z$\lbl{A1}}
\BIC{$x = y \to y = z \to x = z$}
\UIC{$\forall z\,(x = y \to y = z \to x = z)$}
\UIC{$\forall y\,\forall z\,(x = y \to y = z \to x = z)$}
\UIC{$\forall x\,\forall y\,\forall z\,(x = y \to y = z \to x = z)$}
\AXC{$\varphi_{11}$\lbl{A2}}
\BIC{$\forall y\,\forall z\,(v_{491} = y \to y = z \to v_{491} = z)$}
\AXC{$\varphi_{21}$\lbl{A2}}
\BIC{$\forall z\,(v_{491} = v_{492} \to v_{492} = z \to v_{491} = z)$}
\AXC{$\varphi_{34}$\lbl{A2}}
\BIC{$v_{491} = v_{492} \to v_{492} = v_{493} \to v_{491} = v_{493}$}
\UIC{$\forall v_{493}\,(v_{491} = v_{492} \to v_{492} = v_{493} \to v_{491} = v_{493})$}
\UIC{$\varphi_{3}$}
\UIC{$\forall v_{491}\,(\varphi_{3})$}
\end{ptree}
\noindent$\Pi_{6}$:
\begin{ptree}
\AXC{$x = x$\lbl{A5}}
\UIC{$\forall x\,(x = x)$}
\AXC{$\forall x\,(x = x) \to x^{-1} \cdot (x \cdot y) = x^{-1} \cdot (x \cdot y)$\lbl{A2}}
\BIC{$x^{-1} \cdot (x \cdot y) = x^{-1} \cdot (x \cdot y)$}
\end{ptree}
\noindent$\Pi_{7}$:
\begin{ptree}
\AXC{$\stackrel{\Pi_{6}}{\vdots}$}\noLine
\UIC{$x^{-1} \cdot (x \cdot y) = x^{-1} \cdot (x \cdot y)$}
\AXC{$\varphi_{26}$\lbl{A1}}
\AXC{$\varphi_{80}$\lbl{A1}}
\BIC{$x \in \mathbb{R} \to y \in \mathbb{R} \to x \neq 0 \to y \neq 0 \to \neg x \cdot y \neq 0 \to x \cdot y = 0$}
\AXC{$\varphi_{94}$\lbl{A1}}
\BIC{$\varphi_{92}$}
\BIC{$(x \cdot y = 0 \to \varphi_{30}) \to x \in \mathbb{R} \to \varphi_{46}$}
\end{ptree}
\noindent$\Pi_{8}$:
\begin{ptree}
\AXC{$z = w \to \varphi_{19}$\lbl{A6}}
\UIC{$\forall w\,(z = w \to \varphi_{19})$}
\UIC{$\forall z\,\forall w\,(z = w \to \varphi_{19})$}
\AXC{$\forall z\,\forall w\,(z = w \to \varphi_{19}) \to \varphi_{53}$\lbl{A2}}
\BIC{$\varphi_{53}$}
\AXC{$\varphi_{53} \to x \cdot y = 0 \to \varphi_{30}$\lbl{A2}}
\BIC{$x \cdot y = 0 \to \varphi_{30}$}
\AXC{$\stackrel{\Pi_{7}}{\vdots}$}\noLine
\UIC{$(x \cdot y = 0 \to \varphi_{30}) \to x \in \mathbb{R} \to \varphi_{46}$}
\BIC{$x \in \mathbb{R} \to \varphi_{46}$}
\end{ptree}
\noindent$\Pi_{9}$:
\begin{ptree}
\AXC{$y = z \to x = y \to x = z$\lbl{A6}}
\AXC{$(y = z \to x = y \to x = z) \to x = y \to y = z \to x = z$\lbl{A1}}
\BIC{$x = y \to y = z \to x = z$}
\UIC{$\forall z\,(x = y \to y = z \to x = z)$}
\UIC{$\forall y\,\forall z\,(x = y \to y = z \to x = z)$}
\UIC{$\forall x\,\forall y\,\forall z\,(x = y \to y = z \to x = z)$}
\AXC{$\varphi_{12}$\lbl{A2}}
\BIC{$\forall y\,\forall z\,(v_{488} = y \to y = z \to v_{488} = z)$}
\AXC{$\varphi_{22}$\lbl{A2}}
\BIC{$\forall z\,(v_{488} = v_{489} \to v_{489} = z \to v_{488} = z)$}
\AXC{$\varphi_{35}$\lbl{A2}}
\BIC{$v_{488} = v_{489} \to v_{489} = v_{490} \to v_{488} = v_{490}$}
\UIC{$\forall v_{490}\,(v_{488} = v_{489} \to v_{489} = v_{490} \to v_{488} = v_{490})$}
\UIC{$\varphi_{4}$}
\UIC{$\forall v_{488}\,(\varphi_{4})$}
\end{ptree}
\noindent$\Pi_{10}$:
\begin{ptree}
\AXC{$\stackrel{\Pi_{3}}{\vdots}$}\noLine
\UIC{$\varphi_{44}$}
\AXC{$\varphi_{47}$\lbl{(R6)}}
\AXC{$\varphi_{47} \to \varphi_{57}$\lbl{A2}}
\BIC{$\varphi_{57}$}
\AXC{$\varphi_{57} \to \varphi_{55}$\lbl{A2}}
\BIC{$\varphi_{55}$}
\AXC{$\varphi_{55} \to \varphi_{52}$\lbl{A2}}
\BIC{$\varphi_{52}$}
\AXC{$\varphi_{52} \to \varphi_{44} \to \varphi_{93}$\lbl{A1}}
\BIC{$\varphi_{44} \to \varphi_{93}$}
\BIC{$\varphi_{93}$}
\end{ptree}
\noindent$\Pi_{11}$:
\begin{ptree}
\AXC{$x \in \mathbb{R} \to y \in \mathbb{R} \to x \neq 0 \to y \neq 0 \to \neg x \cdot y \neq 0 \to y \in \mathbb{R}$\lbl{A1}}
\AXC{$x \in \mathbb{R} \to y \in \mathbb{R} \to x \neq 0 \to y \neq 0 \to \neg x \cdot y \neq 0 \to x \in \mathbb{R}$\lbl{A1}}
\AXC{$\stackrel{\Pi_{10}}{\vdots}$}\noLine
\UIC{$\varphi_{93}$}
\BIC{$\varphi_{87}$}
\BIC{$x \in \mathbb{R} \to y \in \mathbb{R} \to \varphi_{25}$}
\end{ptree}
\noindent$\Pi_{12}$:
\begin{ptree}
\AXC{$y = z \to x = y \to x = z$\lbl{A6}}
\AXC{$(y = z \to x = y \to x = z) \to x = y \to y = z \to x = z$\lbl{A1}}
\BIC{$x = y \to y = z \to x = z$}
\UIC{$\forall z\,(x = y \to y = z \to x = z)$}
\UIC{$\forall y\,\forall z\,(x = y \to y = z \to x = z)$}
\UIC{$\forall x\,\forall y\,\forall z\,(x = y \to y = z \to x = z)$}
\AXC{$\varphi_{13}$\lbl{A2}}
\BIC{$\forall y\,\forall z\,(v_{485} = y \to y = z \to v_{485} = z)$}
\AXC{$\varphi_{23}$\lbl{A2}}
\BIC{$\forall z\,(v_{485} = v_{486} \to v_{486} = z \to v_{485} = z)$}
\AXC{$\varphi_{36}$\lbl{A2}}
\BIC{$v_{485} = v_{486} \to v_{486} = v_{487} \to v_{485} = v_{487}$}
\UIC{$\forall v_{487}\,(v_{485} = v_{486} \to v_{486} = v_{487} \to v_{485} = v_{487})$}
\UIC{$\varphi_{5}$}
\UIC{$\forall v_{485}\,(\varphi_{5})$}
\end{ptree}
\noindent$\Pi_{13}$:
\begin{ptree}
\AXC{$x = x$\lbl{A5}}
\AXC{$x = y \to x = x \to y = x$\lbl{A6}}
\AXC{$(x = y \to x = x \to y = x) \to x = x \to x = y \to y = x$\lbl{A1}}
\BIC{$x = x \to x = y \to y = x$}
\BIC{$x = y \to y = x$}
\UIC{$\forall y\,(x = y \to y = x)$}
\UIC{$\forall x\,\forall y\,(x = y \to y = x)$}
\end{ptree}
\noindent$\Pi_{14}$:
\begin{ptree}
\AXC{$y = z \to x = y \to x = z$\lbl{A6}}
\AXC{$(y = z \to x = y \to x = z) \to x = y \to y = z \to x = z$\lbl{A1}}
\BIC{$x = y \to y = z \to x = z$}
\UIC{$\forall z\,(x = y \to y = z \to x = z)$}
\UIC{$\forall y\,\forall z\,(x = y \to y = z \to x = z)$}
\UIC{$\forall x\,\forall y\,\forall z\,(x = y \to y = z \to x = z)$}
\AXC{$\varphi_{45}$\lbl{A2}}
\BIC{$\forall y\,\forall z\,(x^{-1} \cdot x = y \to y = z \to x^{-1} \cdot x = z)$}
\end{ptree}
\noindent$\Pi_{15}$:
\begin{ptree}
\AXC{$\stackrel{\Pi_{14}}{\vdots}$}\noLine
\UIC{$\forall y\,\forall z\,(x^{-1} \cdot x = y \to y = z \to x^{-1} \cdot x = z)$}
\AXC{$\forall y\,\forall z\,(x^{-1} \cdot x = y \to y = z \to x^{-1} \cdot x = z) \to \varphi_{17}$\lbl{A2}}
\BIC{$\varphi_{17}$}
\end{ptree}
\noindent$\Pi_{16}$:
\begin{ptree}
\AXC{$\forall x \in \mathbb{R}\,(x \neq 0 \to x \cdot x^{-1} = 1)$\lbl{(R9)}}
\AXC{$\varphi_{49}$\lbl{A2}}
\BIC{$x \in \mathbb{R} \to x \neq 0 \to x \cdot x^{-1} = 1$}
\AXC{$(x \in \mathbb{R} \to x \neq 0 \to x \cdot x^{-1} = 1) \to \varphi_{90}$\lbl{A1}}
\BIC{$\varphi_{90}$}
\end{ptree}
\noindent$\Pi_{17}$:
\begin{ptree}
\AXC{$x \in \mathbb{R} \to y \in \mathbb{R} \to x \neq 0 \to y \neq 0 \to x \neq 0$\lbl{A1}}
\AXC{$x \in \mathbb{R} \to y \in \mathbb{R} \to x \neq 0 \to y \neq 0 \to \neg x \cdot y \neq 0 \to x \in \mathbb{R}$\lbl{A1}}
\AXC{$\stackrel{\Pi_{16}}{\vdots}$}\noLine
\UIC{$\varphi_{90}$}
\BIC{$(x \in \mathbb{R} \to y \in \mathbb{R} \to x \neq 0 \to y \neq 0 \to x \neq 0) \to \varphi_{37}$}
\BIC{$\varphi_{37}$}
\end{ptree}
\noindent$\Pi_{18}$:
\begin{ptree}
\AXC{$\forall x\,\forall y\,(x \in \mathbb{R} \to y \in \mathbb{R} \to x \cdot y = y \cdot x)$\lbl{(R7)}}
\AXC{$\varphi_{76}$\lbl{A2}}
\BIC{$\forall y\,(x^{-1} \in \mathbb{R} \to y \in \mathbb{R} \to x^{-1} \cdot y = y \cdot x^{-1})$}
\AXC{$\varphi_{78}$\lbl{A2}}
\BIC{$x^{-1} \in \mathbb{R} \to x \in \mathbb{R} \to x^{-1} \cdot x = x \cdot x^{-1}$}
\end{ptree}
\noindent$\Pi_{19}$:
\begin{ptree}
\AXC{$\stackrel{\Pi_{3}}{\vdots}$}\noLine
\UIC{$\varphi_{44}$}
\AXC{$\stackrel{\Pi_{18}}{\vdots}$}\noLine
\UIC{$x^{-1} \in \mathbb{R} \to x \in \mathbb{R} \to x^{-1} \cdot x = x \cdot x^{-1}$}
\AXC{$(x^{-1} \in \mathbb{R} \to x \in \mathbb{R} \to x^{-1} \cdot x = x \cdot x^{-1}) \to \varphi_{44} \to \varphi_{85}$\lbl{A1}}
\BIC{$\varphi_{44} \to \varphi_{85}$}
\BIC{$\varphi_{85}$}
\end{ptree}
\noindent$\Pi_{20}$:
\begin{ptree}
\AXC{$\stackrel{\Pi_{17}}{\vdots}$}\noLine
\UIC{$\varphi_{37}$}
\AXC{$x \in \mathbb{R} \to y \in \mathbb{R} \to x \neq 0 \to y \neq 0 \to \neg x \cdot y \neq 0 \to x \in \mathbb{R}$\lbl{A1}}
\AXC{$\stackrel{\Pi_{19}}{\vdots}$}\noLine
\UIC{$\varphi_{85}$}
\BIC{$x \in \mathbb{R} \to \varphi_{27}$}
\AXC{$(x \in \mathbb{R} \to \varphi_{27}) \to \varphi_{37} \to \varphi_{79}$\lbl{A1}}
\BIC{$\varphi_{37} \to \varphi_{79}$}
\BIC{$\varphi_{79}$}
\end{ptree}
\noindent$\Pi_{21}$:
\begin{ptree}
\AXC{$\stackrel{\Pi_{15}}{\vdots}$}\noLine
\UIC{$\varphi_{17}$}
\AXC{$\varphi_{17} \to x^{-1} \cdot x = x \cdot x^{-1} \to x \cdot x^{-1} = 1 \to x^{-1} \cdot x = 1$\lbl{A2}}
\BIC{$x^{-1} \cdot x = x \cdot x^{-1} \to x \cdot x^{-1} = 1 \to x^{-1} \cdot x = 1$}
\AXC{$\stackrel{\Pi_{20}}{\vdots}$}\noLine
\UIC{$\varphi_{79}$}
\BIC{$\varphi_{38}$}
\AXC{$\varphi_{38} \to (x^{-1} \cdot x = 1 \to 1 = x^{-1} \cdot x) \to \varphi_{39}$\lbl{A1}}
\BIC{$(x^{-1} \cdot x = 1 \to 1 = x^{-1} \cdot x) \to \varphi_{39}$}
\end{ptree}
\noindent$\Pi_{22}$:
\begin{ptree}
\AXC{$\stackrel{\Pi_{13}}{\vdots}$}\noLine
\UIC{$\forall x\,\forall y\,(x = y \to y = x)$}
\AXC{$\varphi_{8}$\lbl{A2}}
\BIC{$\forall y\,(x^{-1} \cdot x = y \to y = x^{-1} \cdot x)$}
\AXC{$\varphi_{29}$\lbl{A2}}
\BIC{$x^{-1} \cdot x = 1 \to 1 = x^{-1} \cdot x$}
\AXC{$\stackrel{\Pi_{21}}{\vdots}$}\noLine
\UIC{$(x^{-1} \cdot x = 1 \to 1 = x^{-1} \cdot x) \to \varphi_{39}$}
\BIC{$\varphi_{39}$}
\AXC{$\varphi_{39} \to 1 \cdot y = 1 \cdot y \to \varphi_{84}$\lbl{A1}}
\BIC{$1 \cdot y = 1 \cdot y \to \varphi_{84}$}
\end{ptree}
\noindent$\Pi_{23}$:
\begin{ptree}
\AXC{$x = x$\lbl{A5}}
\UIC{$\forall x\,(x = x)$}
\AXC{$\forall x\,(x = x) \to 1 \cdot y = 1 \cdot y$\lbl{A2}}
\BIC{$1 \cdot y = 1 \cdot y$}
\AXC{$\stackrel{\Pi_{22}}{\vdots}$}\noLine
\UIC{$1 \cdot y = 1 \cdot y \to \varphi_{84}$}
\BIC{$\varphi_{84}$}
\end{ptree}
\noindent$\Pi_{24}$:
\begin{ptree}
\AXC{$z = w \to 1 \cdot y = z \cdot y \to 1 \cdot y = w \cdot y$\lbl{A6}}
\UIC{$\forall w\,(z = w \to 1 \cdot y = z \cdot y \to 1 \cdot y = w \cdot y)$}
\UIC{$\forall z\,\forall w\,(z = w \to 1 \cdot y = z \cdot y \to 1 \cdot y = w \cdot y)$}
\AXC{$\varphi_{56}$\lbl{A2}}
\BIC{$\forall w\,(1 = w \to 1 \cdot y = 1 \cdot y \to 1 \cdot y = w \cdot y)$}
\AXC{$\varphi_{70}$\lbl{A2}}
\BIC{$1 = x^{-1} \cdot x \to 1 \cdot y = 1 \cdot y \to 1 \cdot y = x^{-1} \cdot x \cdot y$}
\AXC{$\stackrel{\Pi_{23}}{\vdots}$}\noLine
\UIC{$\varphi_{84}$}
\BIC{$x \in \mathbb{R} \to \varphi_{28}$}
\end{ptree}
\noindent$\Pi_{25}$:
\begin{ptree}
\AXC{$\stackrel{\Pi_{13}}{\vdots}$}\noLine
\UIC{$\forall x\,\forall y\,(x = y \to y = x)$}
\AXC{$\forall x\,\forall y\,(x = y \to y = x) \to \forall y\,(v_{483} = y \to y = v_{483})$\lbl{A2}}
\BIC{$\forall y\,(v_{483} = y \to y = v_{483})$}
\AXC{$\varphi_{1}$\lbl{A2}}
\BIC{$v_{483} = v_{484} \to v_{484} = v_{483}$}
\UIC{$\forall v_{484}\,(v_{483} = v_{484} \to v_{484} = v_{483})$}
\UIC{$\forall v_{483}\,\forall v_{484}\,(v_{483} = v_{484} \to v_{484} = v_{483})$}
\AXC{$\varphi_{43}$\lbl{A2}}
\BIC{$\forall v_{484}\,(1 \cdot y = v_{484} \to v_{484} = 1 \cdot y)$}
\AXC{$\varphi_{9}$\lbl{A2}}
\BIC{$1 \cdot y = y \to y = 1 \cdot y$}
\end{ptree}
\noindent$\Pi_{26}$:
\begin{ptree}
\AXC{$y = z \to x = y \to x = z$\lbl{A6}}
\AXC{$(y = z \to x = y \to x = z) \to x = y \to y = z \to x = z$\lbl{A1}}
\BIC{$x = y \to y = z \to x = z$}
\UIC{$\forall z\,(x = y \to y = z \to x = z)$}
\UIC{$\forall y\,\forall z\,(x = y \to y = z \to x = z)$}
\UIC{$\forall x\,\forall y\,\forall z\,(x = y \to y = z \to x = z)$}
\AXC{$\varphi_{14}$\lbl{A2}}
\BIC{$\forall y\,\forall z\,(v_{480} = y \to y = z \to v_{480} = z)$}
\AXC{$\varphi_{24}$\lbl{A2}}
\BIC{$\forall z\,(v_{480} = v_{481} \to v_{481} = z \to v_{480} = z)$}
\AXC{$\varphi_{40}$\lbl{A2}}
\BIC{$v_{480} = v_{481} \to v_{481} = v_{482} \to v_{480} = v_{482}$}
\UIC{$\forall v_{482}\,(v_{480} = v_{481} \to v_{481} = v_{482} \to v_{480} = v_{482})$}
\UIC{$\varphi_{6}$}
\UIC{$\forall v_{480}\,(\varphi_{6})$}
\end{ptree}
\noindent$\Pi_{27}$:
\begin{ptree}
\AXC{$\forall x \in \mathbb{R}\,(x \cdot 1 = x)$\lbl{(R8)}}
\AXC{$\forall x \in \mathbb{R}\,(x \cdot 1 = x) \to y \in \mathbb{R} \to y \cdot 1 = y$\lbl{A2}}
\BIC{$y \in \mathbb{R} \to y \cdot 1 = y$}
\AXC{$(y \in \mathbb{R} \to y \cdot 1 = y) \to \varphi_{81}$\lbl{A1}}
\BIC{$\varphi_{81}$}
\end{ptree}
\noindent$\Pi_{28}$:
\begin{ptree}
\AXC{$\forall x\,\forall y\,(x \in \mathbb{R} \to y \in \mathbb{R} \to x \cdot y = y \cdot x)$\lbl{(R7)}}
\AXC{$\varphi_{63}$\lbl{A2}}
\BIC{$\forall y\,(1 \in \mathbb{R} \to y \in \mathbb{R} \to 1 \cdot y = y \cdot 1)$}
\AXC{$\varphi_{59}$\lbl{A2}}
\BIC{$1 \in \mathbb{R} \to y \in \mathbb{R} \to 1 \cdot y = y \cdot 1$}
\AXC{$(1 \in \mathbb{R} \to y \in \mathbb{R} \to 1 \cdot y = y \cdot 1) \to 1 \in \mathbb{R} \to \varphi_{83}$\lbl{A1}}
\BIC{$1 \in \mathbb{R} \to \varphi_{83}$}
\end{ptree}
\noindent$\Pi_{29}$:
\begin{ptree}
\AXC{$x \in \mathbb{R} \to y \in \mathbb{R} \to x \neq 0 \to y \neq 0 \to \neg x \cdot y \neq 0 \to y \in \mathbb{R}$\lbl{A1}}
\AXC{$1 \in \mathbb{R}$\lbl{(R1)}}
\AXC{$\stackrel{\Pi_{28}}{\vdots}$}\noLine
\UIC{$1 \in \mathbb{R} \to \varphi_{83}$}
\BIC{$\varphi_{83}$}
\BIC{$\varphi_{41}$}
\AXC{$\varphi_{41} \to \varphi_{88}$\lbl{A1}}
\BIC{$\varphi_{88}$}
\end{ptree}
\noindent$\Pi_{30}$:
\begin{ptree}
\AXC{$x \in \mathbb{R} \to y \in \mathbb{R} \to x \neq 0 \to y \neq 0 \to \neg x \cdot y \neq 0 \to y \in \mathbb{R}$\lbl{A1}}
\AXC{$\stackrel{\Pi_{27}}{\vdots}$}\noLine
\UIC{$\varphi_{81}$}
\BIC{$x \in \mathbb{R} \to y \in \mathbb{R} \to x \neq 0 \to y \neq 0 \to \neg x \cdot y \neq 0 \to y \cdot 1 = y$}
\AXC{$\stackrel{\Pi_{29}}{\vdots}$}\noLine
\UIC{$\varphi_{88}$}
\BIC{$\varphi_{74}$}
\end{ptree}
\noindent$\Pi_{31}$:
\begin{ptree}
\AXC{$\stackrel{\Pi_{25}}{\vdots}$}\noLine
\UIC{$1 \cdot y = y \to y = 1 \cdot y$}
\AXC{$\stackrel{\Pi_{26}}{\vdots}$}\noLine
\UIC{$\forall v_{480}\,(\varphi_{6})$}
\AXC{$\forall v_{480}\,(\varphi_{6}) \to \varphi_{7}$\lbl{A2}}
\BIC{$\varphi_{7}$}
\AXC{$\varphi_{73}$\lbl{A2}}
\BIC{$\forall v_{482}\,(1 \cdot y = y \cdot 1 \to y \cdot 1 = v_{482} \to 1 \cdot y = v_{482})$}
\AXC{$\varphi_{61}$\lbl{A2}}
\BIC{$1 \cdot y = y \cdot 1 \to y \cdot 1 = y \to 1 \cdot y = y$}
\AXC{$\stackrel{\Pi_{30}}{\vdots}$}\noLine
\UIC{$\varphi_{74}$}
\BIC{$x \in \mathbb{R} \to y \in \mathbb{R} \to x \neq 0 \to y \neq 0 \to \neg x \cdot y \neq 0 \to 1 \cdot y = y$}
\AXC{$\varphi_{86}$\lbl{A1}}
\BIC{$\varphi_{58}$}
\BIC{$x \in \mathbb{R} \to y \in \mathbb{R} \to x \neq 0 \to y \neq 0 \to \neg x \cdot y \neq 0 \to y = 1 \cdot y$}
\end{ptree}
\noindent$\Pi_{32}$:
\begin{ptree}
\AXC{$\stackrel{\Pi_{24}}{\vdots}$}\noLine
\UIC{$x \in \mathbb{R} \to \varphi_{28}$}
\AXC{$\stackrel{\Pi_{31}}{\vdots}$}\noLine
\UIC{$x \in \mathbb{R} \to y \in \mathbb{R} \to x \neq 0 \to y \neq 0 \to \neg x \cdot y \neq 0 \to y = 1 \cdot y$}
\AXC{$\varphi_{95}$\lbl{A1}}
\BIC{$(x \in \mathbb{R} \to \varphi_{28}) \to \varphi_{82}$}
\BIC{$\varphi_{82}$}
\end{ptree}
\noindent$\Pi_{33}$:
\begin{ptree}
\AXC{$\stackrel{\Pi_{12}}{\vdots}$}\noLine
\UIC{$\forall v_{485}\,(\varphi_{5})$}
\AXC{$\varphi_{67}$\lbl{A2}}
\BIC{$\forall v_{486}\,\forall v_{487}\,(y = v_{486} \to v_{486} = v_{487} \to y = v_{487})$}
\AXC{$\varphi_{54}$\lbl{A2}}
\BIC{$\forall v_{487}\,(y = 1 \cdot y \to 1 \cdot y = v_{487} \to y = v_{487})$}
\AXC{$\varphi_{69}$\lbl{A2}}
\BIC{$y = 1 \cdot y \to 1 \cdot y = x^{-1} \cdot x \cdot y \to y = x^{-1} \cdot x \cdot y$}
\AXC{$\stackrel{\Pi_{32}}{\vdots}$}\noLine
\UIC{$\varphi_{82}$}
\BIC{$\varphi_{48}$}
\end{ptree}
\noindent$\Pi_{34}$:
\begin{ptree}
\AXC{$\stackrel{\Pi_{11}}{\vdots}$}\noLine
\UIC{$x \in \mathbb{R} \to y \in \mathbb{R} \to \varphi_{25}$}
\AXC{$\stackrel{\Pi_{33}}{\vdots}$}\noLine
\UIC{$\varphi_{48}$}
\AXC{$\varphi_{48} \to (x \in \mathbb{R} \to y \in \mathbb{R} \to \varphi_{25}) \to \varphi_{50} \to x \in \mathbb{R} \to \varphi_{18}$\lbl{A1}}
\BIC{$(x \in \mathbb{R} \to y \in \mathbb{R} \to \varphi_{25}) \to \varphi_{50} \to x \in \mathbb{R} \to \varphi_{18}$}
\BIC{$\varphi_{50} \to x \in \mathbb{R} \to \varphi_{18}$}
\end{ptree}
\noindent$\Pi_{35}$:
\begin{ptree}
\AXC{$\stackrel{\Pi_{9}}{\vdots}$}\noLine
\UIC{$\forall v_{488}\,(\varphi_{4})$}
\AXC{$\varphi_{66}$\lbl{A2}}
\BIC{$\forall v_{489}\,\forall v_{490}\,(y = v_{489} \to v_{489} = v_{490} \to y = v_{490})$}
\AXC{$\varphi_{72}$\lbl{A2}}
\BIC{$\varphi_{16}$}
\AXC{$\varphi_{16} \to \varphi_{50}$\lbl{A2}}
\BIC{$\varphi_{50}$}
\AXC{$\stackrel{\Pi_{34}}{\vdots}$}\noLine
\UIC{$\varphi_{50} \to x \in \mathbb{R} \to \varphi_{18}$}
\BIC{$x \in \mathbb{R} \to \varphi_{18}$}
\end{ptree}
\noindent$\Pi_{36}$:
\begin{ptree}
\AXC{$\stackrel{\Pi_{8}}{\vdots}$}\noLine
\UIC{$x \in \mathbb{R} \to \varphi_{46}$}
\AXC{$\stackrel{\Pi_{35}}{\vdots}$}\noLine
\UIC{$x \in \mathbb{R} \to \varphi_{18}$}
\AXC{$(x \in \mathbb{R} \to \varphi_{18}) \to (x \in \mathbb{R} \to \varphi_{46}) \to \varphi_{31} \to \varphi_{42}$\lbl{A1}}
\BIC{$(x \in \mathbb{R} \to \varphi_{46}) \to \varphi_{31} \to \varphi_{42}$}
\BIC{$\varphi_{31} \to \varphi_{42}$}
\end{ptree}
\noindent$\Pi_{37}$:
\begin{ptree}
\AXC{$\stackrel{\Pi_{5}}{\vdots}$}\noLine
\UIC{$\forall v_{491}\,(\varphi_{3})$}
\AXC{$\varphi_{65}$\lbl{A2}}
\BIC{$\forall v_{492}\,\forall v_{493}\,(y = v_{492} \to v_{492} = v_{493} \to y = v_{493})$}
\AXC{$\varphi_{71}$\lbl{A2}}
\BIC{$\varphi_{15}$}
\AXC{$\varphi_{15} \to \varphi_{31}$\lbl{A2}}
\BIC{$\varphi_{31}$}
\AXC{$\stackrel{\Pi_{36}}{\vdots}$}\noLine
\UIC{$\varphi_{31} \to \varphi_{42}$}
\BIC{$\varphi_{42}$}
\AXC{$\varphi_{42} \to \varphi_{33} \to \varphi_{68}$\lbl{A1}}
\BIC{$\varphi_{33} \to \varphi_{68}$}
\end{ptree}
\noindent$\Pi_{38}$:
\begin{ptree}
\AXC{$\stackrel{\Pi_{1}}{\vdots}$}\noLine
\UIC{$\forall v_{494}\,(\varphi_{2})$}
\AXC{$\varphi_{64}$\lbl{A2}}
\BIC{$\forall v_{495}\,\forall v_{496}\,(y = v_{495} \to v_{495} = v_{496} \to y = v_{496})$}
\AXC{$\varphi_{62}$\lbl{A2}}
\BIC{$\forall v_{496}\,(y = x^{-1} \cdot 0 \to x^{-1} \cdot 0 = v_{496} \to y = v_{496})$}
\AXC{$\varphi_{60}$\lbl{A2}}
\BIC{$y = x^{-1} \cdot 0 \to x^{-1} \cdot 0 = 0 \to y = 0$}
\AXC{$\stackrel{\Pi_{3}}{\vdots}$}\noLine
\UIC{$\varphi_{44}$}
\AXC{$\stackrel{\Pi_{4}}{\vdots}$}\noLine
\UIC{$\varphi_{44} \to \varphi_{33}$}
\BIC{$\varphi_{33}$}
\AXC{$\stackrel{\Pi_{37}}{\vdots}$}\noLine
\UIC{$\varphi_{33} \to \varphi_{68}$}
\BIC{$\varphi_{68}$}
\BIC{$x \in \mathbb{R} \to y \in \mathbb{R} \to x \neq 0 \to y \neq 0 \to \neg x \cdot y \neq 0 \to y = 0$}
\end{ptree}
\noindent 主木:
\begin{ptree}
\AXC{$x \in \mathbb{R} \to y \in \mathbb{R} \to x \neq 0 \to y \neq 0 \to y \neq 0$\lbl{A1}}
\AXC{$\varphi_{75}$\lbl{A1}}
\BIC{$x \in \mathbb{R} \to y \in \mathbb{R} \to x \neq 0 \to y \neq 0 \to \neg x \cdot y \neq 0 \to y \neq 0$}
\AXC{$\stackrel{\Pi_{38}}{\vdots}$}\noLine
\UIC{$x \in \mathbb{R} \to y \in \mathbb{R} \to x \neq 0 \to y \neq 0 \to \neg x \cdot y \neq 0 \to y = 0$}
\AXC{$\varphi_{89}$\lbl{A1}}
\BIC{$\varphi_{77}$}
\BIC{$x \in \mathbb{R} \to y \in \mathbb{R} \to x \neq 0 \to y \neq 0 \to x \cdot y \neq 0$}
\UIC{$\forall y\,(x \in \mathbb{R} \to y \in \mathbb{R} \to x \neq 0 \to y \neq 0 \to x \cdot y \neq 0)$}
\UIC{$\forall x\,\forall y\,(x \in \mathbb{R} \to y \in \mathbb{R} \to x \neq 0 \to y \neq 0 \to x \cdot y \neq 0)$}
\end{ptree}
\begin{abbrevlist}
\item $\varphi_{1} :\equiv \forall y\,(v_{483} = y \to y = v_{483}) \to v_{483} = v_{484} \to v_{484} = v_{483}$
\item $\varphi_{2} :\equiv \forall v_{495}\,\forall v_{496}\,(v_{494} = v_{495} \to v_{495} = v_{496} \to v_{494} = v_{496})$
\item $\varphi_{3} :\equiv \forall v_{492}\,\forall v_{493}\,(v_{491} = v_{492} \to v_{492} = v_{493} \to v_{491} = v_{493})$
\item $\varphi_{4} :\equiv \forall v_{489}\,\forall v_{490}\,(v_{488} = v_{489} \to v_{489} = v_{490} \to v_{488} = v_{490})$
\item $\varphi_{5} :\equiv \forall v_{486}\,\forall v_{487}\,(v_{485} = v_{486} \to v_{486} = v_{487} \to v_{485} = v_{487})$
\item $\varphi_{6} :\equiv \forall v_{481}\,\forall v_{482}\,(v_{480} = v_{481} \to v_{481} = v_{482} \to v_{480} = v_{482})$
\item $\varphi_{7} :\equiv \forall v_{481}\,\forall v_{482}\,(1 \cdot y = v_{481} \to v_{481} = v_{482} \to 1 \cdot y = v_{482})$
\item $\varphi_{8} :\equiv \forall x\,\forall y\,(x = y \to y = x) \to \forall y\,(x^{-1} \cdot x = y \to y = x^{-1} \cdot x)$
\item $\varphi_{9} :\equiv \forall v_{484}\,(1 \cdot y = v_{484} \to v_{484} = 1 \cdot y) \to 1 \cdot y = y \to y = 1 \cdot y$
\item $\varphi_{10} :\equiv \forall x\,\forall y\,\forall z\,(x = y \to y = z \to x = z) \to \forall y\,\forall z\,(v_{494} = y \to y = z \to v_{494} = z)$
\item $\varphi_{11} :\equiv \forall x\,\forall y\,\forall z\,(x = y \to y = z \to x = z) \to \forall y\,\forall z\,(v_{491} = y \to y = z \to v_{491} = z)$
\item $\varphi_{12} :\equiv \forall x\,\forall y\,\forall z\,(x = y \to y = z \to x = z) \to \forall y\,\forall z\,(v_{488} = y \to y = z \to v_{488} = z)$
\item $\varphi_{13} :\equiv \forall x\,\forall y\,\forall z\,(x = y \to y = z \to x = z) \to \forall y\,\forall z\,(v_{485} = y \to y = z \to v_{485} = z)$
\item $\varphi_{14} :\equiv \forall x\,\forall y\,\forall z\,(x = y \to y = z \to x = z) \to \forall y\,\forall z\,(v_{480} = y \to y = z \to v_{480} = z)$
\item $\varphi_{15} :\equiv \forall v_{493}\,(y = x^{-1} \cdot (x \cdot y) \to x^{-1} \cdot (x \cdot y) = v_{493} \to y = v_{493})$
\item $\varphi_{16} :\equiv \forall v_{490}\,(y = x^{-1} \cdot x \cdot y \to x^{-1} \cdot x \cdot y = v_{490} \to y = v_{490})$
\item $\varphi_{17} :\equiv \forall z\,(x^{-1} \cdot x = x \cdot x^{-1} \to x \cdot x^{-1} = z \to x^{-1} \cdot x = z)$
\item $\varphi_{18} :\equiv y \in \mathbb{R} \to x \neq 0 \to y \neq 0 \to \neg x \cdot y \neq 0 \to y = x^{-1} \cdot (x \cdot y)$
\item $\varphi_{19} :\equiv x^{-1} \cdot (x \cdot y) = x^{-1} \cdot z \to x^{-1} \cdot (x \cdot y) = x^{-1} \cdot w$
\item $\varphi_{20} :\equiv \forall y\,\forall z\,(v_{494} = y \to y = z \to v_{494} = z) \to \forall z\,(v_{494} = v_{495} \to v_{495} = z \to v_{494} = z)$
\item $\varphi_{21} :\equiv \forall y\,\forall z\,(v_{491} = y \to y = z \to v_{491} = z) \to \forall z\,(v_{491} = v_{492} \to v_{492} = z \to v_{491} = z)$
\item $\varphi_{22} :\equiv \forall y\,\forall z\,(v_{488} = y \to y = z \to v_{488} = z) \to \forall z\,(v_{488} = v_{489} \to v_{489} = z \to v_{488} = z)$
\item $\varphi_{23} :\equiv \forall y\,\forall z\,(v_{485} = y \to y = z \to v_{485} = z) \to \forall z\,(v_{485} = v_{486} \to v_{486} = z \to v_{485} = z)$
\item $\varphi_{24} :\equiv \forall y\,\forall z\,(v_{480} = y \to y = z \to v_{480} = z) \to \forall z\,(v_{480} = v_{481} \to v_{481} = z \to v_{480} = z)$
\item $\varphi_{25} :\equiv x \neq 0 \to y \neq 0 \to \neg x \cdot y \neq 0 \to x^{-1} \cdot x \cdot y = x^{-1} \cdot (x \cdot y)$
\item $\varphi_{26} :\equiv x \in \mathbb{R} \to y \in \mathbb{R} \to x \neq 0 \to y \neq 0 \to \neg x \cdot y \neq 0 \to \neg x \cdot y \neq 0$
\item $\varphi_{27} :\equiv y \in \mathbb{R} \to x \neq 0 \to y \neq 0 \to \neg x \cdot y \neq 0 \to x^{-1} \cdot x = x \cdot x^{-1}$
\item $\varphi_{28} :\equiv y \in \mathbb{R} \to x \neq 0 \to y \neq 0 \to \neg x \cdot y \neq 0 \to 1 \cdot y = x^{-1} \cdot x \cdot y$
\item $\varphi_{29} :\equiv \forall y\,(x^{-1} \cdot x = y \to y = x^{-1} \cdot x) \to x^{-1} \cdot x = 1 \to 1 = x^{-1} \cdot x$
\item $\varphi_{30} :\equiv x^{-1} \cdot (x \cdot y) = x^{-1} \cdot (x \cdot y) \to x^{-1} \cdot (x \cdot y) = x^{-1} \cdot 0$
\item $\varphi_{31} :\equiv y = x^{-1} \cdot (x \cdot y) \to x^{-1} \cdot (x \cdot y) = x^{-1} \cdot 0 \to y = x^{-1} \cdot 0$
\item $\varphi_{32} :\equiv \forall z\,(v_{494} = v_{495} \to v_{495} = z \to v_{494} = z) \to v_{494} = v_{495} \to v_{495} = v_{496} \to v_{494} = v_{496}$
\item $\varphi_{33} :\equiv x \in \mathbb{R} \to y \in \mathbb{R} \to x \neq 0 \to y \neq 0 \to \neg x \cdot y \neq 0 \to x^{-1} \cdot 0 = 0$
\item $\varphi_{34} :\equiv \forall z\,(v_{491} = v_{492} \to v_{492} = z \to v_{491} = z) \to v_{491} = v_{492} \to v_{492} = v_{493} \to v_{491} = v_{493}$
\item $\varphi_{35} :\equiv \forall z\,(v_{488} = v_{489} \to v_{489} = z \to v_{488} = z) \to v_{488} = v_{489} \to v_{489} = v_{490} \to v_{488} = v_{490}$
\item $\varphi_{36} :\equiv \forall z\,(v_{485} = v_{486} \to v_{486} = z \to v_{485} = z) \to v_{485} = v_{486} \to v_{486} = v_{487} \to v_{485} = v_{487}$
\item $\varphi_{37} :\equiv x \in \mathbb{R} \to y \in \mathbb{R} \to x \neq 0 \to y \neq 0 \to \neg x \cdot y \neq 0 \to x \cdot x^{-1} = 1$
\item $\varphi_{38} :\equiv x \in \mathbb{R} \to y \in \mathbb{R} \to x \neq 0 \to y \neq 0 \to \neg x \cdot y \neq 0 \to x^{-1} \cdot x = 1$
\item $\varphi_{39} :\equiv x \in \mathbb{R} \to y \in \mathbb{R} \to x \neq 0 \to y \neq 0 \to \neg x \cdot y \neq 0 \to 1 = x^{-1} \cdot x$
\item $\varphi_{40} :\equiv \forall z\,(v_{480} = v_{481} \to v_{481} = z \to v_{480} = z) \to v_{480} = v_{481} \to v_{481} = v_{482} \to v_{480} = v_{482}$
\item $\varphi_{41} :\equiv x \in \mathbb{R} \to y \in \mathbb{R} \to x \neq 0 \to y \neq 0 \to \neg x \cdot y \neq 0 \to 1 \cdot y = y \cdot 1$
\item $\varphi_{42} :\equiv x \in \mathbb{R} \to y \in \mathbb{R} \to x \neq 0 \to y \neq 0 \to \neg x \cdot y \neq 0 \to y = x^{-1} \cdot 0$
\item $\varphi_{43} :\equiv \forall v_{483}\,\forall v_{484}\,(v_{483} = v_{484} \to v_{484} = v_{483}) \to \forall v_{484}\,(1 \cdot y = v_{484} \to v_{484} = 1 \cdot y)$
\item $\varphi_{44} :\equiv x \in \mathbb{R} \to y \in \mathbb{R} \to x \neq 0 \to y \neq 0 \to \neg x \cdot y \neq 0 \to x^{-1} \in \mathbb{R}$
\item $\varphi_{45} :\equiv \forall x\,\forall y\,\forall z\,(x = y \to y = z \to x = z) \to \forall y\,\forall z\,(x^{-1} \cdot x = y \to y = z \to x^{-1} \cdot x = z)$
\item $\varphi_{46} :\equiv y \in \mathbb{R} \to x \neq 0 \to y \neq 0 \to \neg x \cdot y \neq 0 \to x^{-1} \cdot (x \cdot y) = x^{-1} \cdot 0$
\item $\varphi_{47} :\equiv \forall x\,\forall y\,\forall z\,(x \in \mathbb{R} \to y \in \mathbb{R} \to z \in \mathbb{R} \to x \cdot y \cdot z = x \cdot (y \cdot z))$
\item $\varphi_{48} :\equiv x \in \mathbb{R} \to y \in \mathbb{R} \to x \neq 0 \to y \neq 0 \to \neg x \cdot y \neq 0 \to y = x^{-1} \cdot x \cdot y$
\item $\varphi_{49} :\equiv \forall x \in \mathbb{R}\,(x \neq 0 \to x \cdot x^{-1} = 1) \to x \in \mathbb{R} \to x \neq 0 \to x \cdot x^{-1} = 1$
\item $\varphi_{50} :\equiv y = x^{-1} \cdot x \cdot y \to x^{-1} \cdot x \cdot y = x^{-1} \cdot (x \cdot y) \to y = x^{-1} \cdot (x \cdot y)$
\item $\varphi_{51} :\equiv \forall x \in \mathbb{R}\,(x \neq 0 \to x^{-1} \in \mathbb{R}) \to x \in \mathbb{R} \to x \neq 0 \to x^{-1} \in \mathbb{R}$
\item $\varphi_{52} :\equiv x^{-1} \in \mathbb{R} \to x \in \mathbb{R} \to y \in \mathbb{R} \to x^{-1} \cdot x \cdot y = x^{-1} \cdot (x \cdot y)$
\item $\varphi_{53} :\equiv \forall w\,(x \cdot y = w \to x^{-1} \cdot (x \cdot y) = x^{-1} \cdot (x \cdot y) \to x^{-1} \cdot (x \cdot y) = x^{-1} \cdot w)$
\item $\varphi_{54} :\equiv \forall v_{486}\,\forall v_{487}\,(y = v_{486} \to v_{486} = v_{487} \to y = v_{487}) \to \forall v_{487}\,(y = 1 \cdot y \to 1 \cdot y = v_{487} \to y = v_{487})$
\item $\varphi_{55} :\equiv \forall z\,(x^{-1} \in \mathbb{R} \to x \in \mathbb{R} \to z \in \mathbb{R} \to x^{-1} \cdot x \cdot z = x^{-1} \cdot (x \cdot z))$
\item $\varphi_{56} :\equiv \forall z\,\forall w\,(z = w \to 1 \cdot y = z \cdot y \to 1 \cdot y = w \cdot y) \to \forall w\,(1 = w \to 1 \cdot y = 1 \cdot y \to 1 \cdot y = w \cdot y)$
\item $\varphi_{57} :\equiv \forall y\,\forall z\,(x^{-1} \in \mathbb{R} \to y \in \mathbb{R} \to z \in \mathbb{R} \to x^{-1} \cdot y \cdot z = x^{-1} \cdot (y \cdot z))$
\item $\varphi_{58} :\equiv (1 \cdot y = y \to y = 1 \cdot y) \to x \in \mathbb{R} \to y \in \mathbb{R} \to x \neq 0 \to y \neq 0 \to \neg x \cdot y \neq 0 \to y = 1 \cdot y$
\item $\varphi_{59} :\equiv \forall y\,(1 \in \mathbb{R} \to y \in \mathbb{R} \to 1 \cdot y = y \cdot 1) \to 1 \in \mathbb{R} \to y \in \mathbb{R} \to 1 \cdot y = y \cdot 1$
\item $\varphi_{60} :\equiv \forall v_{496}\,(y = x^{-1} \cdot 0 \to x^{-1} \cdot 0 = v_{496} \to y = v_{496}) \to y = x^{-1} \cdot 0 \to x^{-1} \cdot 0 = 0 \to y = 0$
\item $\varphi_{61} :\equiv \forall v_{482}\,(1 \cdot y = y \cdot 1 \to y \cdot 1 = v_{482} \to 1 \cdot y = v_{482}) \to 1 \cdot y = y \cdot 1 \to y \cdot 1 = y \to 1 \cdot y = y$
\item $\varphi_{62} :\equiv \forall v_{495}\,\forall v_{496}\,(y = v_{495} \to v_{495} = v_{496} \to y = v_{496}) \to \forall v_{496}\,(y = x^{-1} \cdot 0 \to x^{-1} \cdot 0 = v_{496} \to y = v_{496})$
\item $\varphi_{63} :\equiv \forall x\,\forall y\,(x \in \mathbb{R} \to y \in \mathbb{R} \to x \cdot y = y \cdot x) \to \forall y\,(1 \in \mathbb{R} \to y \in \mathbb{R} \to 1 \cdot y = y \cdot 1)$
\item $\varphi_{64} :\equiv \forall v_{494}\,(\varphi_{2}) \to \forall v_{495}\,\forall v_{496}\,(y = v_{495} \to v_{495} = v_{496} \to y = v_{496})$
\item $\varphi_{65} :\equiv \forall v_{491}\,(\varphi_{3}) \to \forall v_{492}\,\forall v_{493}\,(y = v_{492} \to v_{492} = v_{493} \to y = v_{493})$
\item $\varphi_{66} :\equiv \forall v_{488}\,(\varphi_{4}) \to \forall v_{489}\,\forall v_{490}\,(y = v_{489} \to v_{489} = v_{490} \to y = v_{490})$
\item $\varphi_{67} :\equiv \forall v_{485}\,(\varphi_{5}) \to \forall v_{486}\,\forall v_{487}\,(y = v_{486} \to v_{486} = v_{487} \to y = v_{487})$
\item $\varphi_{68} :\equiv (y = x^{-1} \cdot 0 \to x^{-1} \cdot 0 = 0 \to y = 0) \to x \in \mathbb{R} \to y \in \mathbb{R} \to x \neq 0 \to y \neq 0 \to \neg x \cdot y \neq 0 \to y = 0$
\item $\varphi_{69} :\equiv \forall v_{487}\,(y = 1 \cdot y \to 1 \cdot y = v_{487} \to y = v_{487}) \to y = 1 \cdot y \to 1 \cdot y = x^{-1} \cdot x \cdot y \to y = x^{-1} \cdot x \cdot y$
\item $\varphi_{70} :\equiv \forall w\,(1 = w \to 1 \cdot y = 1 \cdot y \to 1 \cdot y = w \cdot y) \to 1 = x^{-1} \cdot x \to 1 \cdot y = 1 \cdot y \to 1 \cdot y = x^{-1} \cdot x \cdot y$
\item $\varphi_{71} :\equiv \forall v_{492}\,\forall v_{493}\,(y = v_{492} \to v_{492} = v_{493} \to y = v_{493}) \to \varphi_{15}$
\item $\varphi_{72} :\equiv \forall v_{489}\,\forall v_{490}\,(y = v_{489} \to v_{489} = v_{490} \to y = v_{490}) \to \varphi_{16}$
\item $\varphi_{73} :\equiv \varphi_{7} \to \forall v_{482}\,(1 \cdot y = y \cdot 1 \to y \cdot 1 = v_{482} \to 1 \cdot y = v_{482})$
\item $\varphi_{74} :\equiv (1 \cdot y = y \cdot 1 \to y \cdot 1 = y \to 1 \cdot y = y) \to x \in \mathbb{R} \to y \in \mathbb{R} \to x \neq 0 \to y \neq 0 \to \neg x \cdot y \neq 0 \to 1 \cdot y = y$
\item $\varphi_{75} :\equiv (x \in \mathbb{R} \to y \in \mathbb{R} \to x \neq 0 \to y \neq 0 \to y \neq 0) \to x \in \mathbb{R} \to y \in \mathbb{R} \to x \neq 0 \to y \neq 0 \to \neg x \cdot y \neq 0 \to y \neq 0$
\item $\varphi_{76} :\equiv \forall x\,\forall y\,(x \in \mathbb{R} \to y \in \mathbb{R} \to x \cdot y = y \cdot x) \to \forall y\,(x^{-1} \in \mathbb{R} \to y \in \mathbb{R} \to x^{-1} \cdot y = y \cdot x^{-1})$
\item $\varphi_{77} :\equiv (x \in \mathbb{R} \to y \in \mathbb{R} \to x \neq 0 \to y \neq 0 \to \neg x \cdot y \neq 0 \to y \neq 0) \to x \in \mathbb{R} \to y \in \mathbb{R} \to x \neq 0 \to y \neq 0 \to x \cdot y \neq 0$
\item $\varphi_{78} :\equiv \forall y\,(x^{-1} \in \mathbb{R} \to y \in \mathbb{R} \to x^{-1} \cdot y = y \cdot x^{-1}) \to x^{-1} \in \mathbb{R} \to x \in \mathbb{R} \to x^{-1} \cdot x = x \cdot x^{-1}$
\item $\varphi_{79} :\equiv (x^{-1} \cdot x = x \cdot x^{-1} \to x \cdot x^{-1} = 1 \to x^{-1} \cdot x = 1) \to \varphi_{38}$
\item $\varphi_{80} :\equiv \varphi_{26} \to x \in \mathbb{R} \to y \in \mathbb{R} \to x \neq 0 \to y \neq 0 \to \neg x \cdot y \neq 0 \to x \cdot y = 0$
\item $\varphi_{81} :\equiv (x \in \mathbb{R} \to y \in \mathbb{R} \to x \neq 0 \to y \neq 0 \to \neg x \cdot y \neq 0 \to y \in \mathbb{R}) \to x \in \mathbb{R} \to y \in \mathbb{R} \to x \neq 0 \to y \neq 0 \to \neg x \cdot y \neq 0 \to y \cdot 1 = y$
\item $\varphi_{82} :\equiv (y = 1 \cdot y \to 1 \cdot y = x^{-1} \cdot x \cdot y \to y = x^{-1} \cdot x \cdot y) \to \varphi_{48}$
\item $\varphi_{83} :\equiv (x \in \mathbb{R} \to y \in \mathbb{R} \to x \neq 0 \to y \neq 0 \to \neg x \cdot y \neq 0 \to y \in \mathbb{R}) \to \varphi_{41}$
\item $\varphi_{84} :\equiv (1 = x^{-1} \cdot x \to 1 \cdot y = 1 \cdot y \to 1 \cdot y = x^{-1} \cdot x \cdot y) \to x \in \mathbb{R} \to \varphi_{28}$
\item $\varphi_{85} :\equiv (x \in \mathbb{R} \to y \in \mathbb{R} \to x \neq 0 \to y \neq 0 \to \neg x \cdot y \neq 0 \to x \in \mathbb{R}) \to x \in \mathbb{R} \to \varphi_{27}$
\item $\varphi_{86} :\equiv (x \in \mathbb{R} \to y \in \mathbb{R} \to x \neq 0 \to y \neq 0 \to \neg x \cdot y \neq 0 \to 1 \cdot y = y) \to \varphi_{58}$
\item $\varphi_{87} :\equiv (x \in \mathbb{R} \to y \in \mathbb{R} \to x \neq 0 \to y \neq 0 \to \neg x \cdot y \neq 0 \to y \in \mathbb{R}) \to x \in \mathbb{R} \to y \in \mathbb{R} \to \varphi_{25}$
\item $\varphi_{88} :\equiv (x \in \mathbb{R} \to y \in \mathbb{R} \to x \neq 0 \to y \neq 0 \to \neg x \cdot y \neq 0 \to y \cdot 1 = y) \to \varphi_{74}$
\item $\varphi_{89} :\equiv (x \in \mathbb{R} \to y \in \mathbb{R} \to x \neq 0 \to y \neq 0 \to \neg x \cdot y \neq 0 \to y = 0) \to \varphi_{77}$
\item $\varphi_{90} :\equiv (x \in \mathbb{R} \to y \in \mathbb{R} \to x \neq 0 \to y \neq 0 \to \neg x \cdot y \neq 0 \to x \in \mathbb{R}) \to (x \in \mathbb{R} \to y \in \mathbb{R} \to x \neq 0 \to y \neq 0 \to x \neq 0) \to \varphi_{37}$
\item $\varphi_{91} :\equiv (x \in \mathbb{R} \to y \in \mathbb{R} \to x \neq 0 \to y \neq 0 \to \neg x \cdot y \neq 0 \to x \in \mathbb{R}) \to (x \in \mathbb{R} \to y \in \mathbb{R} \to x \neq 0 \to y \neq 0 \to x \neq 0) \to \varphi_{44}$
\item $\varphi_{92} :\equiv x^{-1} \cdot (x \cdot y) = x^{-1} \cdot (x \cdot y) \to (x \cdot y = 0 \to \varphi_{30}) \to x \in \mathbb{R} \to \varphi_{46}$
\item $\varphi_{93} :\equiv (x \in \mathbb{R} \to y \in \mathbb{R} \to x \neq 0 \to y \neq 0 \to \neg x \cdot y \neq 0 \to x \in \mathbb{R}) \to \varphi_{87}$
\item $\varphi_{94} :\equiv (x \in \mathbb{R} \to y \in \mathbb{R} \to x \neq 0 \to y \neq 0 \to \neg x \cdot y \neq 0 \to x \cdot y = 0) \to \varphi_{92}$
\item $\varphi_{95} :\equiv (x \in \mathbb{R} \to y \in \mathbb{R} \to x \neq 0 \to y \neq 0 \to \neg x \cdot y \neq 0 \to y = 1 \cdot y) \to (x \in \mathbb{R} \to \varphi_{28}) \to \varphi_{82}$
\end{abbrevlist}
\end{fullproof}

% 自動生成 (prover/build.js): InvMul  ノード数 359
\begin{fullproof}{[E16]}{$\forall x\,\forall y\,(x \in \mathbb{R} \to y \in \mathbb{R} \to x \neq 0 \to y \neq 0 \to (x \cdot y)^{-1} = x^{-1} \cdot y^{-1})$}{359}\phantomsection\label{tree:InvMul}
\noindent$\Pi_{1}$:
\begin{ptree}
\AXC{$x = x$\lbl{A5}}
\AXC{$x = y \to x = x \to y = x$\lbl{A6}}
\AXC{$(x = y \to x = x \to y = x) \to x = x \to x = y \to y = x$\lbl{A1}}
\BIC{$x = x \to x = y \to y = x$}
\BIC{$x = y \to y = x$}
\UIC{$\forall y\,(x = y \to y = x)$}
\UIC{$\forall x\,\forall y\,(x = y \to y = x)$}
\end{ptree}
\noindent$\Pi_{2}$:
\begin{ptree}
\AXC{$\stackrel{\Pi_{1}}{\vdots}$}\noLine
\UIC{$\forall x\,\forall y\,(x = y \to y = x)$}
\AXC{$\forall x\,\forall y\,(x = y \to y = x) \to \forall y\,(v_{522} = y \to y = v_{522})$\lbl{A2}}
\BIC{$\forall y\,(v_{522} = y \to y = v_{522})$}
\AXC{$\varphi_{1}$\lbl{A2}}
\BIC{$v_{522} = v_{523} \to v_{523} = v_{522}$}
\UIC{$\forall v_{523}\,(v_{522} = v_{523} \to v_{523} = v_{522})$}
\UIC{$\forall v_{522}\,\forall v_{523}\,(v_{522} = v_{523} \to v_{523} = v_{522})$}
\AXC{$\varphi_{74}$\lbl{A2}}
\BIC{$\forall v_{523}\,(x^{-1} \cdot y^{-1} = v_{523} \to v_{523} = x^{-1} \cdot y^{-1})$}
\AXC{$\varphi_{107}$\lbl{A2}}
\BIC{$\varphi_{26}$}
\end{ptree}
\noindent$\Pi_{3}$:
\begin{ptree}
\AXC{$y = z \to x = y \to x = z$\lbl{A6}}
\AXC{$(y = z \to x = y \to x = z) \to x = y \to y = z \to x = z$\lbl{A1}}
\BIC{$x = y \to y = z \to x = z$}
\UIC{$\forall z\,(x = y \to y = z \to x = z)$}
\UIC{$\forall y\,\forall z\,(x = y \to y = z \to x = z)$}
\UIC{$\forall x\,\forall y\,\forall z\,(x = y \to y = z \to x = z)$}
\AXC{$\varphi_{12}$\lbl{A2}}
\BIC{$\forall y\,\forall z\,(v_{517} = y \to y = z \to v_{517} = z)$}
\AXC{$\varphi_{28}$\lbl{A2}}
\BIC{$\forall z\,(v_{517} = v_{518} \to v_{518} = z \to v_{517} = z)$}
\AXC{$\varphi_{44}$\lbl{A2}}
\BIC{$v_{517} = v_{518} \to v_{518} = v_{519} \to v_{517} = v_{519}$}
\UIC{$\forall v_{519}\,(v_{517} = v_{518} \to v_{518} = v_{519} \to v_{517} = v_{519})$}
\UIC{$\varphi_{3}$}
\UIC{$\forall v_{517}\,(\varphi_{3})$}
\end{ptree}
\noindent$\Pi_{4}$:
\begin{ptree}
\AXC{$\stackrel{\Pi_{3}}{\vdots}$}\noLine
\UIC{$\forall v_{517}\,(\varphi_{3})$}
\AXC{$\forall v_{517}\,(\varphi_{3}) \to \varphi_{77}$\lbl{A2}}
\BIC{$\varphi_{77}$}
\AXC{$\varphi_{77} \to \varphi_{90}$\lbl{A2}}
\BIC{$\varphi_{90}$}
\AXC{$\varphi_{90} \to \varphi_{73}$\lbl{A2}}
\BIC{$\varphi_{73}$}
\end{ptree}
\noindent$\Pi_{5}$:
\begin{ptree}
\AXC{$x \in \mathbb{R} \to y \in \mathbb{R} \to x \neq 0 \to y \neq 0 \to x \in \mathbb{R}$\lbl{A1}}
\AXC{$\forall x \in \mathbb{R}\,(x \neq 0 \to x \cdot x^{-1} = 1)$\lbl{(R9)}}
\AXC{$\varphi_{60}$\lbl{A2}}
\BIC{$x \in \mathbb{R} \to x \neq 0 \to x \cdot x^{-1} = 1$}
\AXC{$\varphi_{123}$\lbl{A1}}
\BIC{$(x \in \mathbb{R} \to y \in \mathbb{R} \to x \neq 0 \to y \neq 0 \to x \in \mathbb{R}) \to \varphi_{93}$}
\BIC{$\varphi_{93}$}
\end{ptree}
\noindent$\Pi_{6}$:
\begin{ptree}
\AXC{$y = z \to x = y \to x = z$\lbl{A6}}
\AXC{$(y = z \to x = y \to x = z) \to x = y \to y = z \to x = z$\lbl{A1}}
\BIC{$x = y \to y = z \to x = z$}
\UIC{$\forall z\,(x = y \to y = z \to x = z)$}
\UIC{$\forall y\,\forall z\,(x = y \to y = z \to x = z)$}
\UIC{$\forall x\,\forall y\,\forall z\,(x = y \to y = z \to x = z)$}
\AXC{$\varphi_{13}$\lbl{A2}}
\BIC{$\forall y\,\forall z\,(v_{514} = y \to y = z \to v_{514} = z)$}
\AXC{$\varphi_{29}$\lbl{A2}}
\BIC{$\forall z\,(v_{514} = v_{515} \to v_{515} = z \to v_{514} = z)$}
\AXC{$\varphi_{45}$\lbl{A2}}
\BIC{$v_{514} = v_{515} \to v_{515} = v_{516} \to v_{514} = v_{516}$}
\UIC{$\forall v_{516}\,(v_{514} = v_{515} \to v_{515} = v_{516} \to v_{514} = v_{516})$}
\UIC{$\varphi_{4}$}
\UIC{$\forall v_{514}\,(\varphi_{4})$}
\end{ptree}
\noindent$\Pi_{7}$:
\begin{ptree}
\AXC{$\stackrel{\Pi_{6}}{\vdots}$}\noLine
\UIC{$\forall v_{514}\,(\varphi_{4})$}
\AXC{$\forall v_{514}\,(\varphi_{4}) \to \varphi_{78}$\lbl{A2}}
\BIC{$\varphi_{78}$}
\AXC{$\varphi_{78} \to \varphi_{100}$\lbl{A2}}
\BIC{$\varphi_{100}$}
\AXC{$\varphi_{100} \to x \cdot y \cdot (x^{-1} \cdot y^{-1}) = x \cdot (x^{-1} \cdot 1) \to \varphi_{55}$\lbl{A2}}
\BIC{$x \cdot y \cdot (x^{-1} \cdot y^{-1}) = x \cdot (x^{-1} \cdot 1) \to \varphi_{55}$}
\end{ptree}
\noindent$\Pi_{8}$:
\begin{ptree}
\AXC{$x = x$\lbl{A5}}
\UIC{$\forall x\,(x = x)$}
\AXC{$\forall x\,(x = x) \to x \cdot (x^{-1} \cdot 1) = x \cdot (x^{-1} \cdot 1)$\lbl{A2}}
\BIC{$x \cdot (x^{-1} \cdot 1) = x \cdot (x^{-1} \cdot 1)$}
\end{ptree}
\noindent$\Pi_{9}$:
\begin{ptree}
\AXC{$x \in \mathbb{R} \to y \in \mathbb{R} \to x \neq 0 \to y \neq 0 \to x \in \mathbb{R}$\lbl{A1}}
\AXC{$\forall x \in \mathbb{R}\,(x \neq 0 \to x^{-1} \in \mathbb{R})$\lbl{(R1)}}
\AXC{$\varphi_{63}$\lbl{A2}}
\BIC{$x \in \mathbb{R} \to x \neq 0 \to x^{-1} \in \mathbb{R}$}
\AXC{$\varphi_{128}$\lbl{A1}}
\BIC{$(x \in \mathbb{R} \to y \in \mathbb{R} \to x \neq 0 \to y \neq 0 \to x \in \mathbb{R}) \to \varphi_{95}$}
\BIC{$\varphi_{95}$}
\end{ptree}
\noindent$\Pi_{10}$:
\begin{ptree}
\AXC{$x \in \mathbb{R} \to y \in \mathbb{R} \to x \neq 0 \to y \neq 0 \to x \neq 0$\lbl{A1}}
\AXC{$\stackrel{\Pi_{9}}{\vdots}$}\noLine
\UIC{$\varphi_{95}$}
\BIC{$x \in \mathbb{R} \to y \in \mathbb{R} \to x \neq 0 \to y \neq 0 \to x^{-1} \in \mathbb{R}$}
\end{ptree}
\noindent$\Pi_{11}$:
\begin{ptree}
\AXC{$\forall x \in \mathbb{R}\,(x \cdot 1 = x)$\lbl{(R8)}}
\AXC{$\forall x \in \mathbb{R}\,(x \cdot 1 = x) \to x^{-1} \in \mathbb{R} \to x^{-1} \cdot 1 = x^{-1}$\lbl{A2}}
\BIC{$x^{-1} \in \mathbb{R} \to x^{-1} \cdot 1 = x^{-1}$}
\end{ptree}
\noindent$\Pi_{12}$:
\begin{ptree}
\AXC{$\stackrel{\Pi_{10}}{\vdots}$}\noLine
\UIC{$x \in \mathbb{R} \to y \in \mathbb{R} \to x \neq 0 \to y \neq 0 \to x^{-1} \in \mathbb{R}$}
\AXC{$\stackrel{\Pi_{11}}{\vdots}$}\noLine
\UIC{$x^{-1} \in \mathbb{R} \to x^{-1} \cdot 1 = x^{-1}$}
\AXC{$(x^{-1} \in \mathbb{R} \to x^{-1} \cdot 1 = x^{-1}) \to \varphi_{108}$\lbl{A1}}
\BIC{$\varphi_{108}$}
\BIC{$x \in \mathbb{R} \to y \in \mathbb{R} \to x \neq 0 \to y \neq 0 \to x^{-1} \cdot 1 = x^{-1}$}
\end{ptree}
\noindent$\Pi_{13}$:
\begin{ptree}
\AXC{$\stackrel{\Pi_{8}}{\vdots}$}\noLine
\UIC{$x \cdot (x^{-1} \cdot 1) = x \cdot (x^{-1} \cdot 1)$}
\AXC{$\stackrel{\Pi_{12}}{\vdots}$}\noLine
\UIC{$x \in \mathbb{R} \to y \in \mathbb{R} \to x \neq 0 \to y \neq 0 \to x^{-1} \cdot 1 = x^{-1}$}
\AXC{$\varphi_{141}$\lbl{A1}}
\BIC{$\varphi_{127}$}
\BIC{$(x^{-1} \cdot 1 = x^{-1} \to \varphi_{41}) \to \varphi_{51}$}
\end{ptree}
\noindent$\Pi_{14}$:
\begin{ptree}
\AXC{$\varphi_{18}$\lbl{A6}}
\UIC{$\forall w\,(\varphi_{18})$}
\UIC{$\forall z\,\forall w\,(\varphi_{18})$}
\AXC{$\forall z\,\forall w\,(\varphi_{18}) \to \varphi_{66}$\lbl{A2}}
\BIC{$\varphi_{66}$}
\AXC{$\varphi_{66} \to x^{-1} \cdot 1 = x^{-1} \to \varphi_{41}$\lbl{A2}}
\BIC{$x^{-1} \cdot 1 = x^{-1} \to \varphi_{41}$}
\AXC{$\stackrel{\Pi_{13}}{\vdots}$}\noLine
\UIC{$(x^{-1} \cdot 1 = x^{-1} \to \varphi_{41}) \to \varphi_{51}$}
\BIC{$\varphi_{51}$}
\end{ptree}
\noindent$\Pi_{15}$:
\begin{ptree}
\AXC{$y = z \to x = y \to x = z$\lbl{A6}}
\AXC{$(y = z \to x = y \to x = z) \to x = y \to y = z \to x = z$\lbl{A1}}
\BIC{$x = y \to y = z \to x = z$}
\UIC{$\forall z\,(x = y \to y = z \to x = z)$}
\UIC{$\forall y\,\forall z\,(x = y \to y = z \to x = z)$}
\UIC{$\forall x\,\forall y\,\forall z\,(x = y \to y = z \to x = z)$}
\AXC{$\varphi_{14}$\lbl{A2}}
\BIC{$\forall y\,\forall z\,(v_{511} = y \to y = z \to v_{511} = z)$}
\AXC{$\varphi_{30}$\lbl{A2}}
\BIC{$\forall z\,(v_{511} = v_{512} \to v_{512} = z \to v_{511} = z)$}
\AXC{$\varphi_{46}$\lbl{A2}}
\BIC{$v_{511} = v_{512} \to v_{512} = v_{513} \to v_{511} = v_{513}$}
\UIC{$\forall v_{513}\,(v_{511} = v_{512} \to v_{512} = v_{513} \to v_{511} = v_{513})$}
\UIC{$\varphi_{5}$}
\UIC{$\forall v_{511}\,(\varphi_{5})$}
\end{ptree}
\noindent$\Pi_{16}$:
\begin{ptree}
\AXC{$\stackrel{\Pi_{15}}{\vdots}$}\noLine
\UIC{$\forall v_{511}\,(\varphi_{5})$}
\AXC{$\forall v_{511}\,(\varphi_{5}) \to \varphi_{79}$\lbl{A2}}
\BIC{$\varphi_{79}$}
\AXC{$\varphi_{79} \to \varphi_{101}$\lbl{A2}}
\BIC{$\varphi_{101}$}
\AXC{$\varphi_{101} \to x \cdot y \cdot (x^{-1} \cdot y^{-1}) = x \cdot (1 \cdot x^{-1}) \to \varphi_{68}$\lbl{A2}}
\BIC{$x \cdot y \cdot (x^{-1} \cdot y^{-1}) = x \cdot (1 \cdot x^{-1}) \to \varphi_{68}$}
\end{ptree}
\noindent$\Pi_{17}$:
\begin{ptree}
\AXC{$1 \in \mathbb{R}$\lbl{(R1)}}
\AXC{$\forall x\,\forall y\,(x \in \mathbb{R} \to y \in \mathbb{R} \to x \cdot y = y \cdot x)$\lbl{(R7)}}
\AXC{$\varphi_{86}$\lbl{A2}}
\BIC{$\forall y\,(1 \in \mathbb{R} \to y \in \mathbb{R} \to 1 \cdot y = y \cdot 1)$}
\AXC{$\varphi_{91}$\lbl{A2}}
\BIC{$1 \in \mathbb{R} \to x^{-1} \in \mathbb{R} \to 1 \cdot x^{-1} = x^{-1} \cdot 1$}
\AXC{$\varphi_{126}$\lbl{A1}}
\BIC{$1 \in \mathbb{R} \to \varphi_{113}$}
\BIC{$\varphi_{113}$}
\end{ptree}
\noindent$\Pi_{18}$:
\begin{ptree}
\AXC{$\stackrel{\Pi_{10}}{\vdots}$}\noLine
\UIC{$x \in \mathbb{R} \to y \in \mathbb{R} \to x \neq 0 \to y \neq 0 \to x^{-1} \in \mathbb{R}$}
\AXC{$\stackrel{\Pi_{17}}{\vdots}$}\noLine
\UIC{$\varphi_{113}$}
\BIC{$x \in \mathbb{R} \to y \in \mathbb{R} \to x \neq 0 \to y \neq 0 \to 1 \cdot x^{-1} = x^{-1} \cdot 1$}
\AXC{$\varphi_{146}$\lbl{A1}}
\BIC{$\varphi_{133}$}
\end{ptree}
\noindent$\Pi_{19}$:
\begin{ptree}
\AXC{$x = x$\lbl{A5}}
\UIC{$\forall x\,(x = x)$}
\AXC{$\forall x\,(x = x) \to x \cdot (1 \cdot x^{-1}) = x \cdot (1 \cdot x^{-1})$\lbl{A2}}
\BIC{$x \cdot (1 \cdot x^{-1}) = x \cdot (1 \cdot x^{-1})$}
\AXC{$\stackrel{\Pi_{18}}{\vdots}$}\noLine
\UIC{$\varphi_{133}$}
\BIC{$(1 \cdot x^{-1} = x^{-1} \cdot 1 \to \varphi_{56}) \to x \in \mathbb{R} \to \varphi_{34}$}
\end{ptree}
\noindent$\Pi_{20}$:
\begin{ptree}
\AXC{$\varphi_{19}$\lbl{A6}}
\UIC{$\forall w\,(\varphi_{19})$}
\UIC{$\forall z\,\forall w\,(\varphi_{19})$}
\AXC{$\forall z\,\forall w\,(\varphi_{19}) \to \varphi_{67}$\lbl{A2}}
\BIC{$\varphi_{67}$}
\AXC{$\varphi_{67} \to 1 \cdot x^{-1} = x^{-1} \cdot 1 \to \varphi_{56}$\lbl{A2}}
\BIC{$1 \cdot x^{-1} = x^{-1} \cdot 1 \to \varphi_{56}$}
\AXC{$\stackrel{\Pi_{19}}{\vdots}$}\noLine
\UIC{$(1 \cdot x^{-1} = x^{-1} \cdot 1 \to \varphi_{56}) \to x \in \mathbb{R} \to \varphi_{34}$}
\BIC{$x \in \mathbb{R} \to \varphi_{34}$}
\end{ptree}
\noindent$\Pi_{21}$:
\begin{ptree}
\AXC{$y = z \to x = y \to x = z$\lbl{A6}}
\AXC{$(y = z \to x = y \to x = z) \to x = y \to y = z \to x = z$\lbl{A1}}
\BIC{$x = y \to y = z \to x = z$}
\UIC{$\forall z\,(x = y \to y = z \to x = z)$}
\UIC{$\forall y\,\forall z\,(x = y \to y = z \to x = z)$}
\UIC{$\forall x\,\forall y\,\forall z\,(x = y \to y = z \to x = z)$}
\AXC{$\varphi_{15}$\lbl{A2}}
\BIC{$\forall y\,\forall z\,(v_{508} = y \to y = z \to v_{508} = z)$}
\AXC{$\varphi_{31}$\lbl{A2}}
\BIC{$\forall z\,(v_{508} = v_{509} \to v_{509} = z \to v_{508} = z)$}
\AXC{$\varphi_{47}$\lbl{A2}}
\BIC{$v_{508} = v_{509} \to v_{509} = v_{510} \to v_{508} = v_{510}$}
\UIC{$\forall v_{510}\,(v_{508} = v_{509} \to v_{509} = v_{510} \to v_{508} = v_{510})$}
\UIC{$\varphi_{6}$}
\UIC{$\forall v_{508}\,(\varphi_{6})$}
\end{ptree}
\noindent$\Pi_{22}$:
\begin{ptree}
\AXC{$z = w \to \varphi_{85}$\lbl{A6}}
\UIC{$\forall w\,(z = w \to \varphi_{85})$}
\UIC{$\forall z\,\forall w\,(z = w \to \varphi_{85})$}
\AXC{$\forall z\,\forall w\,(z = w \to \varphi_{85}) \to \varphi_{109}$\lbl{A2}}
\BIC{$\varphi_{109}$}
\AXC{$\varphi_{109} \to y \cdot y^{-1} = 1 \to \varphi_{92}$\lbl{A2}}
\BIC{$y \cdot y^{-1} = 1 \to \varphi_{92}$}
\end{ptree}
\noindent$\Pi_{23}$:
\begin{ptree}
\AXC{$x \in \mathbb{R} \to y \in \mathbb{R} \to x \neq 0 \to y \neq 0 \to y \in \mathbb{R}$\lbl{A1}}
\AXC{$\forall x \in \mathbb{R}\,(x \neq 0 \to x \cdot x^{-1} = 1)$\lbl{(R9)}}
\AXC{$\varphi_{61}$\lbl{A2}}
\BIC{$y \in \mathbb{R} \to y \neq 0 \to y \cdot y^{-1} = 1$}
\AXC{$\varphi_{124}$\lbl{A1}}
\BIC{$(x \in \mathbb{R} \to y \in \mathbb{R} \to x \neq 0 \to y \neq 0 \to y \in \mathbb{R}) \to \varphi_{94}$}
\BIC{$\varphi_{94}$}
\end{ptree}
\noindent$\Pi_{24}$:
\begin{ptree}
\AXC{$x \in \mathbb{R} \to y \in \mathbb{R} \to x \neq 0 \to y \neq 0 \to y \neq 0$\lbl{A1}}
\AXC{$\stackrel{\Pi_{23}}{\vdots}$}\noLine
\UIC{$\varphi_{94}$}
\BIC{$x \in \mathbb{R} \to y \in \mathbb{R} \to x \neq 0 \to y \neq 0 \to y \cdot y^{-1} = 1$}
\AXC{$(x \in \mathbb{R} \to y \in \mathbb{R} \to x \neq 0 \to y \neq 0 \to y \cdot y^{-1} = 1) \to \varphi_{142}$\lbl{A1}}
\BIC{$\varphi_{142}$}
\end{ptree}
\noindent$\Pi_{25}$:
\begin{ptree}
\AXC{$y = z \to x = y \to x = z$\lbl{A6}}
\AXC{$(y = z \to x = y \to x = z) \to x = y \to y = z \to x = z$\lbl{A1}}
\BIC{$x = y \to y = z \to x = z$}
\UIC{$\forall z\,(x = y \to y = z \to x = z)$}
\UIC{$\forall y\,\forall z\,(x = y \to y = z \to x = z)$}
\UIC{$\forall x\,\forall y\,\forall z\,(x = y \to y = z \to x = z)$}
\AXC{$\varphi_{16}$\lbl{A2}}
\BIC{$\forall y\,\forall z\,(v_{505} = y \to y = z \to v_{505} = z)$}
\AXC{$\varphi_{32}$\lbl{A2}}
\BIC{$\forall z\,(v_{505} = v_{506} \to v_{506} = z \to v_{505} = z)$}
\AXC{$\varphi_{48}$\lbl{A2}}
\BIC{$v_{505} = v_{506} \to v_{506} = v_{507} \to v_{505} = v_{507}$}
\UIC{$\forall v_{507}\,(v_{505} = v_{506} \to v_{506} = v_{507} \to v_{505} = v_{507})$}
\UIC{$\varphi_{7}$}
\UIC{$\forall v_{505}\,(\varphi_{7})$}
\end{ptree}
\noindent$\Pi_{26}$:
\begin{ptree}
\AXC{$\stackrel{\Pi_{25}}{\vdots}$}\noLine
\UIC{$\forall v_{505}\,(\varphi_{7})$}
\AXC{$\forall v_{505}\,(\varphi_{7}) \to \varphi_{81}$\lbl{A2}}
\BIC{$\varphi_{81}$}
\AXC{$\varphi_{81} \to \varphi_{116}$\lbl{A2}}
\BIC{$\varphi_{116}$}
\AXC{$\varphi_{116} \to \varphi_{121}$\lbl{A2}}
\BIC{$\varphi_{121}$}
\end{ptree}
\noindent$\Pi_{27}$:
\begin{ptree}
\AXC{$z = w \to \varphi_{57}$\lbl{A6}}
\UIC{$\forall w\,(z = w \to \varphi_{57})$}
\UIC{$\forall z\,\forall w\,(z = w \to \varphi_{57})$}
\AXC{$\forall z\,\forall w\,(z = w \to \varphi_{57}) \to \varphi_{110}$\lbl{A2}}
\BIC{$\varphi_{110}$}
\AXC{$\varphi_{110} \to y \cdot (y^{-1} \cdot x^{-1}) = y \cdot y^{-1} \cdot x^{-1} \to \varphi_{104}$\lbl{A2}}
\BIC{$y \cdot (y^{-1} \cdot x^{-1}) = y \cdot y^{-1} \cdot x^{-1} \to \varphi_{104}$}
\end{ptree}
\noindent$\Pi_{28}$:
\begin{ptree}
\AXC{$\stackrel{\Pi_{1}}{\vdots}$}\noLine
\UIC{$\forall x\,\forall y\,(x = y \to y = x)$}
\AXC{$\forall x\,\forall y\,(x = y \to y = x) \to \forall y\,(v_{500} = y \to y = v_{500})$\lbl{A2}}
\BIC{$\forall y\,(v_{500} = y \to y = v_{500})$}
\AXC{$\varphi_{2}$\lbl{A2}}
\BIC{$v_{500} = v_{501} \to v_{501} = v_{500}$}
\UIC{$\forall v_{501}\,(v_{500} = v_{501} \to v_{501} = v_{500})$}
\UIC{$\forall v_{500}\,\forall v_{501}\,(v_{500} = v_{501} \to v_{501} = v_{500})$}
\end{ptree}
\noindent$\Pi_{29}$:
\begin{ptree}
\AXC{$\stackrel{\Pi_{28}}{\vdots}$}\noLine
\UIC{$\forall v_{500}\,\forall v_{501}\,(v_{500} = v_{501} \to v_{501} = v_{500})$}
\AXC{$\forall v_{500}\,\forall v_{501}\,(v_{500} = v_{501} \to v_{501} = v_{500}) \to \varphi_{39}$\lbl{A2}}
\BIC{$\varphi_{39}$}
\AXC{$\varphi_{39} \to \varphi_{82}$\lbl{A2}}
\BIC{$\varphi_{82}$}
\end{ptree}
\noindent$\Pi_{30}$:
\begin{ptree}
\AXC{$x \in \mathbb{R} \to y \in \mathbb{R} \to x \neq 0 \to y \neq 0 \to y \in \mathbb{R}$\lbl{A1}}
\AXC{$\forall x \in \mathbb{R}\,(x \neq 0 \to x^{-1} \in \mathbb{R})$\lbl{(R1)}}
\AXC{$\varphi_{64}$\lbl{A2}}
\BIC{$y \in \mathbb{R} \to y \neq 0 \to y^{-1} \in \mathbb{R}$}
\AXC{$\varphi_{129}$\lbl{A1}}
\BIC{$(x \in \mathbb{R} \to y \in \mathbb{R} \to x \neq 0 \to y \neq 0 \to y \in \mathbb{R}) \to \varphi_{96}$}
\BIC{$\varphi_{96}$}
\end{ptree}
\noindent$\Pi_{31}$:
\begin{ptree}
\AXC{$x \in \mathbb{R} \to y \in \mathbb{R} \to x \neq 0 \to y \neq 0 \to y \neq 0$\lbl{A1}}
\AXC{$\stackrel{\Pi_{30}}{\vdots}$}\noLine
\UIC{$\varphi_{96}$}
\BIC{$x \in \mathbb{R} \to y \in \mathbb{R} \to x \neq 0 \to y \neq 0 \to y^{-1} \in \mathbb{R}$}
\end{ptree}
\noindent$\Pi_{32}$:
\begin{ptree}
\AXC{$\forall x\,\forall y\,(\varphi_{42})$\lbl{(R6)}}
\AXC{$\forall x\,\forall y\,(\varphi_{42}) \to \varphi_{65}$\lbl{A2}}
\BIC{$\varphi_{65}$}
\AXC{$\varphi_{65} \to \varphi_{72}$\lbl{A2}}
\BIC{$\varphi_{72}$}
\AXC{$\varphi_{72} \to v_{497} \in \mathbb{R} \to v_{498} \in \mathbb{R} \to \varphi_{21}$\lbl{A2}}
\BIC{$v_{497} \in \mathbb{R} \to v_{498} \in \mathbb{R} \to \varphi_{21}$}
\UIC{$\forall v_{499}\,(v_{497} \in \mathbb{R} \to v_{498} \in \mathbb{R} \to \varphi_{21})$}
\UIC{$\forall v_{498}\,\forall v_{499}\,(v_{497} \in \mathbb{R} \to v_{498} \in \mathbb{R} \to \varphi_{21})$}
\UIC{$\forall v_{497}\,\forall v_{498}\,\forall v_{499}\,(v_{497} \in \mathbb{R} \to v_{498} \in \mathbb{R} \to \varphi_{21})$}
\AXC{$\varphi_{131}$\lbl{A2}}
\BIC{$\varphi_{87}$}
\end{ptree}
\noindent$\Pi_{33}$:
\begin{ptree}
\AXC{$x \in \mathbb{R} \to y \in \mathbb{R} \to x \neq 0 \to y \neq 0 \to y \in \mathbb{R}$\lbl{A1}}
\AXC{$\stackrel{\Pi_{32}}{\vdots}$}\noLine
\UIC{$\varphi_{87}$}
\AXC{$\varphi_{87} \to \varphi_{88}$\lbl{A2}}
\BIC{$\varphi_{88}$}
\AXC{$\varphi_{88} \to y \in \mathbb{R} \to \varphi_{69}$\lbl{A2}}
\BIC{$y \in \mathbb{R} \to \varphi_{69}$}
\AXC{$(y \in \mathbb{R} \to \varphi_{69}) \to \varphi_{140}$\lbl{A1}}
\BIC{$\varphi_{140}$}
\BIC{$(x \in \mathbb{R} \to y \in \mathbb{R} \to x \neq 0 \to y \neq 0 \to y^{-1} \in \mathbb{R}) \to \varphi_{119}$}
\end{ptree}
\noindent$\Pi_{34}$:
\begin{ptree}
\AXC{$\stackrel{\Pi_{31}}{\vdots}$}\noLine
\UIC{$x \in \mathbb{R} \to y \in \mathbb{R} \to x \neq 0 \to y \neq 0 \to y^{-1} \in \mathbb{R}$}
\AXC{$\stackrel{\Pi_{33}}{\vdots}$}\noLine
\UIC{$(x \in \mathbb{R} \to y \in \mathbb{R} \to x \neq 0 \to y \neq 0 \to y^{-1} \in \mathbb{R}) \to \varphi_{119}$}
\BIC{$\varphi_{119}$}
\end{ptree}
\noindent$\Pi_{35}$:
\begin{ptree}
\AXC{$\stackrel{\Pi_{29}}{\vdots}$}\noLine
\UIC{$\varphi_{82}$}
\AXC{$\stackrel{\Pi_{10}}{\vdots}$}\noLine
\UIC{$x \in \mathbb{R} \to y \in \mathbb{R} \to x \neq 0 \to y \neq 0 \to x^{-1} \in \mathbb{R}$}
\AXC{$\stackrel{\Pi_{34}}{\vdots}$}\noLine
\UIC{$\varphi_{119}$}
\BIC{$x \in \mathbb{R} \to \varphi_{58}$}
\AXC{$(x \in \mathbb{R} \to \varphi_{58}) \to \varphi_{82} \to x \in \mathbb{R} \to \varphi_{59}$\lbl{A1}}
\BIC{$\varphi_{82} \to x \in \mathbb{R} \to \varphi_{59}$}
\BIC{$x \in \mathbb{R} \to \varphi_{59}$}
\end{ptree}
\noindent$\Pi_{36}$:
\begin{ptree}
\AXC{$x = x$\lbl{A5}}
\UIC{$\forall x\,(x = x)$}
\AXC{$\forall x\,(x = x) \to \varphi_{9}$\lbl{A2}}
\BIC{$\varphi_{9}$}
\AXC{$\stackrel{\Pi_{35}}{\vdots}$}\noLine
\UIC{$x \in \mathbb{R} \to \varphi_{59}$}
\AXC{$(x \in \mathbb{R} \to \varphi_{59}) \to \varphi_{9} \to \varphi_{139}$\lbl{A1}}
\BIC{$\varphi_{9} \to \varphi_{139}$}
\BIC{$\varphi_{139}$}
\end{ptree}
\noindent$\Pi_{37}$:
\begin{ptree}
\AXC{$y = z \to x = y \to x = z$\lbl{A6}}
\AXC{$(y = z \to x = y \to x = z) \to x = y \to y = z \to x = z$\lbl{A1}}
\BIC{$x = y \to y = z \to x = z$}
\UIC{$\forall z\,(x = y \to y = z \to x = z)$}
\UIC{$\forall y\,\forall z\,(x = y \to y = z \to x = z)$}
\UIC{$\forall x\,\forall y\,\forall z\,(x = y \to y = z \to x = z)$}
\AXC{$\varphi_{17}$\lbl{A2}}
\BIC{$\forall y\,\forall z\,(v_{502} = y \to y = z \to v_{502} = z)$}
\AXC{$\varphi_{33}$\lbl{A2}}
\BIC{$\forall z\,(v_{502} = v_{503} \to v_{503} = z \to v_{502} = z)$}
\AXC{$\varphi_{49}$\lbl{A2}}
\BIC{$v_{502} = v_{503} \to v_{503} = v_{504} \to v_{502} = v_{504}$}
\UIC{$\forall v_{504}\,(v_{502} = v_{503} \to v_{503} = v_{504} \to v_{502} = v_{504})$}
\UIC{$\varphi_{8}$}
\UIC{$\forall v_{502}\,(\varphi_{8})$}
\end{ptree}
\noindent$\Pi_{38}$:
\begin{ptree}
\AXC{$z = w \to \varphi_{70}$\lbl{A6}}
\UIC{$\forall w\,(z = w \to \varphi_{70})$}
\UIC{$\forall z\,\forall w\,(z = w \to \varphi_{70})$}
\AXC{$\forall z\,\forall w\,(z = w \to \varphi_{70}) \to \varphi_{111}$\lbl{A2}}
\BIC{$\varphi_{111}$}
\AXC{$\varphi_{111} \to x^{-1} \cdot y^{-1} = y^{-1} \cdot x^{-1} \to \varphi_{10} \to \varphi_{11}$\lbl{A2}}
\BIC{$x^{-1} \cdot y^{-1} = y^{-1} \cdot x^{-1} \to \varphi_{10} \to \varphi_{11}$}
\end{ptree}
\noindent$\Pi_{39}$:
\begin{ptree}
\AXC{$\forall x\,\forall y\,(x \in \mathbb{R} \to y \in \mathbb{R} \to x \cdot y = y \cdot x)$\lbl{(R7)}}
\AXC{$\varphi_{99}$\lbl{A2}}
\BIC{$\forall y\,(x^{-1} \in \mathbb{R} \to y \in \mathbb{R} \to x^{-1} \cdot y = y \cdot x^{-1})$}
\AXC{$\varphi_{118}$\lbl{A2}}
\BIC{$x^{-1} \in \mathbb{R} \to y^{-1} \in \mathbb{R} \to x^{-1} \cdot y^{-1} = y^{-1} \cdot x^{-1}$}
\AXC{$\varphi_{144}$\lbl{A1}}
\BIC{$\varphi_{125}$}
\end{ptree}
\noindent$\Pi_{40}$:
\begin{ptree}
\AXC{$\stackrel{\Pi_{31}}{\vdots}$}\noLine
\UIC{$x \in \mathbb{R} \to y \in \mathbb{R} \to x \neq 0 \to y \neq 0 \to y^{-1} \in \mathbb{R}$}
\AXC{$\stackrel{\Pi_{10}}{\vdots}$}\noLine
\UIC{$x \in \mathbb{R} \to y \in \mathbb{R} \to x \neq 0 \to y \neq 0 \to x^{-1} \in \mathbb{R}$}
\AXC{$\stackrel{\Pi_{39}}{\vdots}$}\noLine
\UIC{$\varphi_{125}$}
\BIC{$(x \in \mathbb{R} \to y \in \mathbb{R} \to x \neq 0 \to y \neq 0 \to y^{-1} \in \mathbb{R}) \to \varphi_{62}$}
\BIC{$\varphi_{62}$}
\AXC{$\varphi_{62} \to \varphi_{10} \to \varphi_{137}$\lbl{A1}}
\BIC{$\varphi_{10} \to \varphi_{137}$}
\end{ptree}
\noindent$\Pi_{41}$:
\begin{ptree}
\AXC{$\forall x\,\forall y\,(x \in \mathbb{R} \to y \in \mathbb{R} \to x \cdot y \in \mathbb{R})$\lbl{(R1)}}
\AXC{$\varphi_{97}$\lbl{A2}}
\BIC{$\forall y\,(x^{-1} \in \mathbb{R} \to y \in \mathbb{R} \to x^{-1} \cdot y \in \mathbb{R})$}
\AXC{$\varphi_{112}$\lbl{A2}}
\BIC{$x^{-1} \in \mathbb{R} \to y^{-1} \in \mathbb{R} \to x^{-1} \cdot y^{-1} \in \mathbb{R}$}
\AXC{$\varphi_{138}$\lbl{A1}}
\BIC{$(x \in \mathbb{R} \to y \in \mathbb{R} \to x \neq 0 \to y \neq 0 \to x^{-1} \in \mathbb{R}) \to \varphi_{114}$}
\end{ptree}
\noindent$\Pi_{42}$:
\begin{ptree}
\AXC{$\stackrel{\Pi_{10}}{\vdots}$}\noLine
\UIC{$x \in \mathbb{R} \to y \in \mathbb{R} \to x \neq 0 \to y \neq 0 \to x^{-1} \in \mathbb{R}$}
\AXC{$\stackrel{\Pi_{41}}{\vdots}$}\noLine
\UIC{$(x \in \mathbb{R} \to y \in \mathbb{R} \to x \neq 0 \to y \neq 0 \to x^{-1} \in \mathbb{R}) \to \varphi_{114}$}
\BIC{$\varphi_{114}$}
\end{ptree}
\noindent$\Pi_{43}$:
\begin{ptree}
\AXC{$\stackrel{\Pi_{31}}{\vdots}$}\noLine
\UIC{$x \in \mathbb{R} \to y \in \mathbb{R} \to x \neq 0 \to y \neq 0 \to y^{-1} \in \mathbb{R}$}
\AXC{$\stackrel{\Pi_{42}}{\vdots}$}\noLine
\UIC{$\varphi_{114}$}
\BIC{$x \in \mathbb{R} \to y \in \mathbb{R} \to x \neq 0 \to y \neq 0 \to x^{-1} \cdot y^{-1} \in \mathbb{R}$}
\end{ptree}
\noindent$\Pi_{44}$:
\begin{ptree}
\AXC{$\forall x\,\forall y\,(\varphi_{42})$\lbl{(R6)}}
\AXC{$\forall x\,\forall y\,(\varphi_{42}) \to \forall y\,(\varphi_{42})$\lbl{A2}}
\BIC{$\forall y\,(\varphi_{42})$}
\AXC{$\forall y\,(\varphi_{42}) \to \varphi_{42}$\lbl{A2}}
\BIC{$\varphi_{42}$}
\AXC{$\varphi_{42} \to x \in \mathbb{R} \to y \in \mathbb{R} \to \varphi_{71}$\lbl{A2}}
\BIC{$x \in \mathbb{R} \to y \in \mathbb{R} \to \varphi_{71}$}
\end{ptree}
\noindent$\Pi_{45}$:
\begin{ptree}
\AXC{$x \in \mathbb{R} \to y \in \mathbb{R} \to x \neq 0 \to y \neq 0 \to x \in \mathbb{R}$\lbl{A1}}
\AXC{$\stackrel{\Pi_{44}}{\vdots}$}\noLine
\UIC{$x \in \mathbb{R} \to y \in \mathbb{R} \to \varphi_{71}$}
\AXC{$(x \in \mathbb{R} \to y \in \mathbb{R} \to \varphi_{71}) \to \varphi_{145}$\lbl{A1}}
\BIC{$\varphi_{145}$}
\BIC{$(x \in \mathbb{R} \to y \in \mathbb{R} \to x \neq 0 \to y \neq 0 \to y \in \mathbb{R}) \to \varphi_{120}$}
\end{ptree}
\noindent$\Pi_{46}$:
\begin{ptree}
\AXC{$x \in \mathbb{R} \to y \in \mathbb{R} \to x \neq 0 \to y \neq 0 \to y \in \mathbb{R}$\lbl{A1}}
\AXC{$\stackrel{\Pi_{45}}{\vdots}$}\noLine
\UIC{$(x \in \mathbb{R} \to y \in \mathbb{R} \to x \neq 0 \to y \neq 0 \to y \in \mathbb{R}) \to \varphi_{120}$}
\BIC{$\varphi_{120}$}
\end{ptree}
\noindent$\Pi_{47}$:
\begin{ptree}
\AXC{$\stackrel{\Pi_{43}}{\vdots}$}\noLine
\UIC{$x \in \mathbb{R} \to y \in \mathbb{R} \to x \neq 0 \to y \neq 0 \to x^{-1} \cdot y^{-1} \in \mathbb{R}$}
\AXC{$\stackrel{\Pi_{46}}{\vdots}$}\noLine
\UIC{$\varphi_{120}$}
\BIC{$x \in \mathbb{R} \to y \in \mathbb{R} \to x \neq 0 \to \varphi_{36}$}
\AXC{$(x \in \mathbb{R} \to y \in \mathbb{R} \to x \neq 0 \to \varphi_{36}) \to \varphi_{148}$\lbl{A1}}
\BIC{$\varphi_{148}$}
\end{ptree}
\noindent$\Pi_{48}$:
\begin{ptree}
\AXC{$\stackrel{\Pi_{38}}{\vdots}$}\noLine
\UIC{$x^{-1} \cdot y^{-1} = y^{-1} \cdot x^{-1} \to \varphi_{10} \to \varphi_{11}$}
\AXC{$x = x$\lbl{A5}}
\UIC{$\forall x\,(x = x)$}
\AXC{$\forall x\,(x = x) \to \varphi_{10}$\lbl{A2}}
\BIC{$\varphi_{10}$}
\AXC{$\stackrel{\Pi_{40}}{\vdots}$}\noLine
\UIC{$\varphi_{10} \to \varphi_{137}$}
\BIC{$\varphi_{137}$}
\BIC{$x \in \mathbb{R} \to y \in \mathbb{R} \to x \neq 0 \to y \neq 0 \to \varphi_{11}$}
\AXC{$\stackrel{\Pi_{47}}{\vdots}$}\noLine
\UIC{$\varphi_{148}$}
\BIC{$\varphi_{122} \to x \in \mathbb{R} \to y \in \mathbb{R} \to x \neq 0 \to \varphi_{37}$}
\end{ptree}
\noindent$\Pi_{49}$:
\begin{ptree}
\AXC{$\stackrel{\Pi_{37}}{\vdots}$}\noLine
\UIC{$\forall v_{502}\,(\varphi_{8})$}
\AXC{$\forall v_{502}\,(\varphi_{8}) \to \varphi_{83}$\lbl{A2}}
\BIC{$\varphi_{83}$}
\AXC{$\varphi_{83} \to \varphi_{117}$\lbl{A2}}
\BIC{$\varphi_{117}$}
\AXC{$\varphi_{117} \to \varphi_{122}$\lbl{A2}}
\BIC{$\varphi_{122}$}
\AXC{$\stackrel{\Pi_{48}}{\vdots}$}\noLine
\UIC{$\varphi_{122} \to x \in \mathbb{R} \to y \in \mathbb{R} \to x \neq 0 \to \varphi_{37}$}
\BIC{$x \in \mathbb{R} \to y \in \mathbb{R} \to x \neq 0 \to \varphi_{37}$}
\end{ptree}
\noindent$\Pi_{50}$:
\begin{ptree}
\AXC{$\stackrel{\Pi_{49}}{\vdots}$}\noLine
\UIC{$x \in \mathbb{R} \to y \in \mathbb{R} \to x \neq 0 \to \varphi_{37}$}
\AXC{$(x \in \mathbb{R} \to y \in \mathbb{R} \to x \neq 0 \to \varphi_{37}) \to \varphi_{149}$\lbl{A1}}
\BIC{$\varphi_{149}$}
\end{ptree}
\noindent$\Pi_{51}$:
\begin{ptree}
\AXC{$\stackrel{\Pi_{26}}{\vdots}$}\noLine
\UIC{$\varphi_{121}$}
\AXC{$\stackrel{\Pi_{27}}{\vdots}$}\noLine
\UIC{$y \cdot (y^{-1} \cdot x^{-1}) = y \cdot y^{-1} \cdot x^{-1} \to \varphi_{104}$}
\AXC{$\stackrel{\Pi_{36}}{\vdots}$}\noLine
\UIC{$\varphi_{139}$}
\BIC{$x \in \mathbb{R} \to y \in \mathbb{R} \to x \neq 0 \to \varphi_{35}$}
\AXC{$\stackrel{\Pi_{50}}{\vdots}$}\noLine
\UIC{$\varphi_{149}$}
\BIC{$\varphi_{121} \to x \in \mathbb{R} \to y \in \mathbb{R} \to x \neq 0 \to \varphi_{38}$}
\BIC{$x \in \mathbb{R} \to y \in \mathbb{R} \to x \neq 0 \to \varphi_{38}$}
\AXC{$\varphi_{150}$\lbl{A1}}
\BIC{$(x \in \mathbb{R} \to y \in \mathbb{R} \to \varphi_{20}) \to \varphi_{136}$}
\end{ptree}
\noindent$\Pi_{52}$:
\begin{ptree}
\AXC{$\stackrel{\Pi_{22}}{\vdots}$}\noLine
\UIC{$y \cdot y^{-1} = 1 \to \varphi_{92}$}
\AXC{$x = x$\lbl{A5}}
\UIC{$\forall x\,(x = x)$}
\AXC{$\varphi_{40}$\lbl{A2}}
\BIC{$x \cdot (y \cdot y^{-1} \cdot x^{-1}) = x \cdot (y \cdot y^{-1} \cdot x^{-1})$}
\AXC{$\stackrel{\Pi_{24}}{\vdots}$}\noLine
\UIC{$\varphi_{142}$}
\BIC{$(y \cdot y^{-1} = 1 \to \varphi_{92}) \to x \in \mathbb{R} \to y \in \mathbb{R} \to \varphi_{20}$}
\BIC{$x \in \mathbb{R} \to y \in \mathbb{R} \to \varphi_{20}$}
\AXC{$\stackrel{\Pi_{51}}{\vdots}$}\noLine
\UIC{$(x \in \mathbb{R} \to y \in \mathbb{R} \to \varphi_{20}) \to \varphi_{136}$}
\BIC{$\varphi_{136}$}
\end{ptree}
\noindent$\Pi_{53}$:
\begin{ptree}
\AXC{$\stackrel{\Pi_{21}}{\vdots}$}\noLine
\UIC{$\forall v_{508}\,(\varphi_{6})$}
\AXC{$\forall v_{508}\,(\varphi_{6}) \to \varphi_{80}$\lbl{A2}}
\BIC{$\varphi_{80}$}
\AXC{$\varphi_{80} \to \varphi_{115}$\lbl{A2}}
\BIC{$\varphi_{115}$}
\AXC{$\varphi_{147}$\lbl{A2}}
\BIC{$x \cdot y \cdot (x^{-1} \cdot y^{-1}) = x \cdot (y \cdot y^{-1} \cdot x^{-1}) \to \varphi_{84}$}
\AXC{$\stackrel{\Pi_{52}}{\vdots}$}\noLine
\UIC{$\varphi_{136}$}
\BIC{$x \in \mathbb{R} \to y \in \mathbb{R} \to \varphi_{22}$}
\end{ptree}
\noindent$\Pi_{54}$:
\begin{ptree}
\AXC{$\stackrel{\Pi_{20}}{\vdots}$}\noLine
\UIC{$x \in \mathbb{R} \to \varphi_{34}$}
\AXC{$\stackrel{\Pi_{53}}{\vdots}$}\noLine
\UIC{$x \in \mathbb{R} \to y \in \mathbb{R} \to \varphi_{22}$}
\AXC{$(x \in \mathbb{R} \to y \in \mathbb{R} \to \varphi_{22}) \to (x \in \mathbb{R} \to \varphi_{34}) \to \varphi_{134}$\lbl{A1}}
\BIC{$(x \in \mathbb{R} \to \varphi_{34}) \to \varphi_{134}$}
\BIC{$\varphi_{134}$}
\end{ptree}
\noindent$\Pi_{55}$:
\begin{ptree}
\AXC{$\stackrel{\Pi_{14}}{\vdots}$}\noLine
\UIC{$\varphi_{51}$}
\AXC{$\stackrel{\Pi_{16}}{\vdots}$}\noLine
\UIC{$x \cdot y \cdot (x^{-1} \cdot y^{-1}) = x \cdot (1 \cdot x^{-1}) \to \varphi_{68}$}
\AXC{$\stackrel{\Pi_{54}}{\vdots}$}\noLine
\UIC{$\varphi_{134}$}
\BIC{$x \in \mathbb{R} \to y \in \mathbb{R} \to \varphi_{23}$}
\AXC{$(x \in \mathbb{R} \to y \in \mathbb{R} \to \varphi_{23}) \to \varphi_{51} \to \varphi_{130}$\lbl{A1}}
\BIC{$\varphi_{51} \to \varphi_{130}$}
\BIC{$\varphi_{130}$}
\end{ptree}
\noindent$\Pi_{56}$:
\begin{ptree}
\AXC{$\stackrel{\Pi_{7}}{\vdots}$}\noLine
\UIC{$x \cdot y \cdot (x^{-1} \cdot y^{-1}) = x \cdot (x^{-1} \cdot 1) \to \varphi_{55}$}
\AXC{$\stackrel{\Pi_{55}}{\vdots}$}\noLine
\UIC{$\varphi_{130}$}
\BIC{$x \in \mathbb{R} \to \varphi_{43}$}
\AXC{$(x \in \mathbb{R} \to \varphi_{43}) \to \varphi_{135}$\lbl{A1}}
\BIC{$\varphi_{135}$}
\end{ptree}
\noindent$\Pi_{57}$:
\begin{ptree}
\AXC{$\forall x\,\forall y\,(x \in \mathbb{R} \to y \in \mathbb{R} \to x \neq 0 \to y \neq 0 \to x \cdot y \neq 0)$\lbl{[E15]}}
\AXC{$\varphi_{106}$\lbl{A2}}
\BIC{$\forall y\,(x \in \mathbb{R} \to y \in \mathbb{R} \to x \neq 0 \to y \neq 0 \to x \cdot y \neq 0)$}
\AXC{$\varphi_{98}$\lbl{A2}}
\BIC{$x \in \mathbb{R} \to y \in \mathbb{R} \to x \neq 0 \to y \neq 0 \to x \cdot y \neq 0$}
\end{ptree}
\noindent$\Pi_{58}$:
\begin{ptree}
\AXC{$\forall x\,\forall y\,(x \in \mathbb{R} \to y \in \mathbb{R} \to x \cdot y \in \mathbb{R})$\lbl{(R1)}}
\AXC{$\varphi_{89}$\lbl{A2}}
\BIC{$\forall y\,(x \in \mathbb{R} \to y \in \mathbb{R} \to x \cdot y \in \mathbb{R})$}
\AXC{$\varphi_{76}$\lbl{A2}}
\BIC{$x \in \mathbb{R} \to y \in \mathbb{R} \to x \cdot y \in \mathbb{R}$}
\AXC{$\varphi_{132}$\lbl{A1}}
\BIC{$(x \in \mathbb{R} \to y \in \mathbb{R} \to x \neq 0 \to y \neq 0 \to x \in \mathbb{R}) \to \varphi_{102}$}
\end{ptree}
\noindent$\Pi_{59}$:
\begin{ptree}
\AXC{$x \in \mathbb{R} \to y \in \mathbb{R} \to x \neq 0 \to y \neq 0 \to x \in \mathbb{R}$\lbl{A1}}
\AXC{$\stackrel{\Pi_{58}}{\vdots}$}\noLine
\UIC{$(x \in \mathbb{R} \to y \in \mathbb{R} \to x \neq 0 \to y \neq 0 \to x \in \mathbb{R}) \to \varphi_{102}$}
\BIC{$\varphi_{102}$}
\end{ptree}
\noindent$\Pi_{60}$:
\begin{ptree}
\AXC{$\varphi_{25}$\lbl{[E13]}}
\AXC{$\varphi_{25} \to \varphi_{50}$\lbl{A2}}
\BIC{$\varphi_{50}$}
\AXC{$\varphi_{50} \to v_{520} \in \mathbb{R} \to \varphi_{24}$\lbl{A2}}
\BIC{$v_{520} \in \mathbb{R} \to \varphi_{24}$}
\UIC{$\forall v_{521}\,(v_{520} \in \mathbb{R} \to \varphi_{24})$}
\UIC{$\forall v_{520}\,\forall v_{521}\,(v_{520} \in \mathbb{R} \to \varphi_{24})$}
\AXC{$\forall v_{520}\,\forall v_{521}\,(v_{520} \in \mathbb{R} \to \varphi_{24}) \to \varphi_{75}$\lbl{A2}}
\BIC{$\varphi_{75}$}
\end{ptree}
\noindent$\Pi_{61}$:
\begin{ptree}
\AXC{$\stackrel{\Pi_{60}}{\vdots}$}\noLine
\UIC{$\varphi_{75}$}
\AXC{$\varphi_{75} \to x \cdot y \in \mathbb{R} \to x^{-1} \cdot y^{-1} \in \mathbb{R} \to x \cdot y \neq 0 \to \varphi_{27}$\lbl{A2}}
\BIC{$x \cdot y \in \mathbb{R} \to x^{-1} \cdot y^{-1} \in \mathbb{R} \to x \cdot y \neq 0 \to \varphi_{27}$}
\AXC{$\varphi_{151}$\lbl{A1}}
\BIC{$(x \in \mathbb{R} \to y \in \mathbb{R} \to x \neq 0 \to y \neq 0 \to x \cdot y \in \mathbb{R}) \to \varphi_{143}$}
\end{ptree}
\noindent$\Pi_{62}$:
\begin{ptree}
\AXC{$x \in \mathbb{R} \to y \in \mathbb{R} \to x \neq 0 \to y \neq 0 \to y \in \mathbb{R}$\lbl{A1}}
\AXC{$\stackrel{\Pi_{59}}{\vdots}$}\noLine
\UIC{$\varphi_{102}$}
\BIC{$x \in \mathbb{R} \to y \in \mathbb{R} \to x \neq 0 \to y \neq 0 \to x \cdot y \in \mathbb{R}$}
\AXC{$\stackrel{\Pi_{61}}{\vdots}$}\noLine
\UIC{$(x \in \mathbb{R} \to y \in \mathbb{R} \to x \neq 0 \to y \neq 0 \to x \cdot y \in \mathbb{R}) \to \varphi_{143}$}
\BIC{$\varphi_{143}$}
\end{ptree}
\noindent$\Pi_{63}$:
\begin{ptree}
\AXC{$\stackrel{\Pi_{57}}{\vdots}$}\noLine
\UIC{$x \in \mathbb{R} \to y \in \mathbb{R} \to x \neq 0 \to y \neq 0 \to x \cdot y \neq 0$}
\AXC{$\stackrel{\Pi_{43}}{\vdots}$}\noLine
\UIC{$x \in \mathbb{R} \to y \in \mathbb{R} \to x \neq 0 \to y \neq 0 \to x^{-1} \cdot y^{-1} \in \mathbb{R}$}
\AXC{$\stackrel{\Pi_{62}}{\vdots}$}\noLine
\UIC{$\varphi_{143}$}
\BIC{$(x \in \mathbb{R} \to y \in \mathbb{R} \to x \neq 0 \to y \neq 0 \to x \cdot y \neq 0) \to \varphi_{52} \to \varphi_{53}$}
\BIC{$\varphi_{52} \to \varphi_{53}$}
\end{ptree}
\noindent 主木:
\begin{ptree}
\AXC{$\stackrel{\Pi_{2}}{\vdots}$}\noLine
\UIC{$\varphi_{26}$}
\AXC{$\stackrel{\Pi_{4}}{\vdots}$}\noLine
\UIC{$\varphi_{73}$}
\AXC{$x \in \mathbb{R} \to y \in \mathbb{R} \to x \neq 0 \to y \neq 0 \to x \neq 0$\lbl{A1}}
\AXC{$\stackrel{\Pi_{5}}{\vdots}$}\noLine
\UIC{$\varphi_{93}$}
\BIC{$x \in \mathbb{R} \to y \in \mathbb{R} \to x \neq 0 \to y \neq 0 \to x \cdot x^{-1} = 1$}
\AXC{$\stackrel{\Pi_{56}}{\vdots}$}\noLine
\UIC{$\varphi_{135}$}
\BIC{$\varphi_{73} \to \varphi_{52}$}
\BIC{$\varphi_{52}$}
\AXC{$\stackrel{\Pi_{63}}{\vdots}$}\noLine
\UIC{$\varphi_{52} \to \varphi_{53}$}
\BIC{$\varphi_{53}$}
\AXC{$\varphi_{53} \to \varphi_{26} \to \varphi_{54}$\lbl{A1}}
\BIC{$\varphi_{26} \to \varphi_{54}$}
\BIC{$\varphi_{54}$}
\UIC{$\forall y\,(\varphi_{54})$}
\UIC{$\forall x\,\forall y\,(\varphi_{54})$}
\end{ptree}
\begin{abbrevlist}
\item $\varphi_{1} :\equiv \forall y\,(v_{522} = y \to y = v_{522}) \to v_{522} = v_{523} \to v_{523} = v_{522}$
\item $\varphi_{2} :\equiv \forall y\,(v_{500} = y \to y = v_{500}) \to v_{500} = v_{501} \to v_{501} = v_{500}$
\item $\varphi_{3} :\equiv \forall v_{518}\,\forall v_{519}\,(v_{517} = v_{518} \to v_{518} = v_{519} \to v_{517} = v_{519})$
\item $\varphi_{4} :\equiv \forall v_{515}\,\forall v_{516}\,(v_{514} = v_{515} \to v_{515} = v_{516} \to v_{514} = v_{516})$
\item $\varphi_{5} :\equiv \forall v_{512}\,\forall v_{513}\,(v_{511} = v_{512} \to v_{512} = v_{513} \to v_{511} = v_{513})$
\item $\varphi_{6} :\equiv \forall v_{509}\,\forall v_{510}\,(v_{508} = v_{509} \to v_{509} = v_{510} \to v_{508} = v_{510})$
\item $\varphi_{7} :\equiv \forall v_{506}\,\forall v_{507}\,(v_{505} = v_{506} \to v_{506} = v_{507} \to v_{505} = v_{507})$
\item $\varphi_{8} :\equiv \forall v_{503}\,\forall v_{504}\,(v_{502} = v_{503} \to v_{503} = v_{504} \to v_{502} = v_{504})$
\item $\varphi_{9} :\equiv x \cdot (y \cdot (y^{-1} \cdot x^{-1})) = x \cdot (y \cdot (y^{-1} \cdot x^{-1}))$
\item $\varphi_{10} :\equiv x \cdot (y \cdot (x^{-1} \cdot y^{-1})) = x \cdot (y \cdot (x^{-1} \cdot y^{-1}))$
\item $\varphi_{11} :\equiv x \cdot (y \cdot (x^{-1} \cdot y^{-1})) = x \cdot (y \cdot (y^{-1} \cdot x^{-1}))$
\item $\varphi_{12} :\equiv \forall x\,\forall y\,\forall z\,(x = y \to y = z \to x = z) \to \forall y\,\forall z\,(v_{517} = y \to y = z \to v_{517} = z)$
\item $\varphi_{13} :\equiv \forall x\,\forall y\,\forall z\,(x = y \to y = z \to x = z) \to \forall y\,\forall z\,(v_{514} = y \to y = z \to v_{514} = z)$
\item $\varphi_{14} :\equiv \forall x\,\forall y\,\forall z\,(x = y \to y = z \to x = z) \to \forall y\,\forall z\,(v_{511} = y \to y = z \to v_{511} = z)$
\item $\varphi_{15} :\equiv \forall x\,\forall y\,\forall z\,(x = y \to y = z \to x = z) \to \forall y\,\forall z\,(v_{508} = y \to y = z \to v_{508} = z)$
\item $\varphi_{16} :\equiv \forall x\,\forall y\,\forall z\,(x = y \to y = z \to x = z) \to \forall y\,\forall z\,(v_{505} = y \to y = z \to v_{505} = z)$
\item $\varphi_{17} :\equiv \forall x\,\forall y\,\forall z\,(x = y \to y = z \to x = z) \to \forall y\,\forall z\,(v_{502} = y \to y = z \to v_{502} = z)$
\item $\varphi_{18} :\equiv z = w \to x \cdot (x^{-1} \cdot 1) = x \cdot z \to x \cdot (x^{-1} \cdot 1) = x \cdot w$
\item $\varphi_{19} :\equiv z = w \to x \cdot (1 \cdot x^{-1}) = x \cdot z \to x \cdot (1 \cdot x^{-1}) = x \cdot w$
\item $\varphi_{20} :\equiv x \neq 0 \to y \neq 0 \to x \cdot (y \cdot y^{-1} \cdot x^{-1}) = x \cdot (1 \cdot x^{-1})$
\item $\varphi_{21} :\equiv v_{499} \in \mathbb{R} \to v_{497} \cdot v_{498} \cdot v_{499} = v_{497} \cdot (v_{498} \cdot v_{499})$
\item $\varphi_{22} :\equiv x \neq 0 \to y \neq 0 \to x \cdot y \cdot (x^{-1} \cdot y^{-1}) = x \cdot (1 \cdot x^{-1})$
\item $\varphi_{23} :\equiv x \neq 0 \to y \neq 0 \to x \cdot y \cdot (x^{-1} \cdot y^{-1}) = x \cdot (x^{-1} \cdot 1)$
\item $\varphi_{24} :\equiv v_{521} \in \mathbb{R} \to v_{520} \neq 0 \to v_{520} \cdot v_{521} = 1 \to v_{521} = v_{520}^{-1}$
\item $\varphi_{25} :\equiv \forall x\,\forall y\,(x \in \mathbb{R} \to y \in \mathbb{R} \to x \neq 0 \to x \cdot y = 1 \to y = x^{-1})$
\item $\varphi_{26} :\equiv x^{-1} \cdot y^{-1} = (x \cdot y)^{-1} \to (x \cdot y)^{-1} = x^{-1} \cdot y^{-1}$
\item $\varphi_{27} :\equiv x \cdot y \cdot (x^{-1} \cdot y^{-1}) = 1 \to x^{-1} \cdot y^{-1} = (x \cdot y)^{-1}$
\item $\varphi_{28} :\equiv \forall y\,\forall z\,(v_{517} = y \to y = z \to v_{517} = z) \to \forall z\,(v_{517} = v_{518} \to v_{518} = z \to v_{517} = z)$
\item $\varphi_{29} :\equiv \forall y\,\forall z\,(v_{514} = y \to y = z \to v_{514} = z) \to \forall z\,(v_{514} = v_{515} \to v_{515} = z \to v_{514} = z)$
\item $\varphi_{30} :\equiv \forall y\,\forall z\,(v_{511} = y \to y = z \to v_{511} = z) \to \forall z\,(v_{511} = v_{512} \to v_{512} = z \to v_{511} = z)$
\item $\varphi_{31} :\equiv \forall y\,\forall z\,(v_{508} = y \to y = z \to v_{508} = z) \to \forall z\,(v_{508} = v_{509} \to v_{509} = z \to v_{508} = z)$
\item $\varphi_{32} :\equiv \forall y\,\forall z\,(v_{505} = y \to y = z \to v_{505} = z) \to \forall z\,(v_{505} = v_{506} \to v_{506} = z \to v_{505} = z)$
\item $\varphi_{33} :\equiv \forall y\,\forall z\,(v_{502} = y \to y = z \to v_{502} = z) \to \forall z\,(v_{502} = v_{503} \to v_{503} = z \to v_{502} = z)$
\item $\varphi_{34} :\equiv y \in \mathbb{R} \to x \neq 0 \to y \neq 0 \to x \cdot (1 \cdot x^{-1}) = x \cdot (x^{-1} \cdot 1)$
\item $\varphi_{35} :\equiv y \neq 0 \to x \cdot (y \cdot (y^{-1} \cdot x^{-1})) = x \cdot (y \cdot y^{-1} \cdot x^{-1})$
\item $\varphi_{36} :\equiv y \neq 0 \to x \cdot y \cdot (x^{-1} \cdot y^{-1}) = x \cdot (y \cdot (x^{-1} \cdot y^{-1}))$
\item $\varphi_{37} :\equiv y \neq 0 \to x \cdot y \cdot (x^{-1} \cdot y^{-1}) = x \cdot (y \cdot (y^{-1} \cdot x^{-1}))$
\item $\varphi_{38} :\equiv y \neq 0 \to x \cdot y \cdot (x^{-1} \cdot y^{-1}) = x \cdot (y \cdot y^{-1} \cdot x^{-1})$
\item $\varphi_{39} :\equiv \forall v_{501}\,(y \cdot y^{-1} \cdot x^{-1} = v_{501} \to v_{501} = y \cdot y^{-1} \cdot x^{-1})$
\item $\varphi_{40} :\equiv \forall x\,(x = x) \to x \cdot (y \cdot y^{-1} \cdot x^{-1}) = x \cdot (y \cdot y^{-1} \cdot x^{-1})$
\item $\varphi_{41} :\equiv x \cdot (x^{-1} \cdot 1) = x \cdot (x^{-1} \cdot 1) \to x \cdot (x^{-1} \cdot 1) = x \cdot x^{-1}$
\item $\varphi_{42} :\equiv \forall z\,(x \in \mathbb{R} \to y \in \mathbb{R} \to z \in \mathbb{R} \to x \cdot y \cdot z = x \cdot (y \cdot z))$
\item $\varphi_{43} :\equiv y \in \mathbb{R} \to x \neq 0 \to y \neq 0 \to x \cdot y \cdot (x^{-1} \cdot y^{-1}) = x \cdot x^{-1}$
\item $\varphi_{44} :\equiv \forall z\,(v_{517} = v_{518} \to v_{518} = z \to v_{517} = z) \to v_{517} = v_{518} \to v_{518} = v_{519} \to v_{517} = v_{519}$
\item $\varphi_{45} :\equiv \forall z\,(v_{514} = v_{515} \to v_{515} = z \to v_{514} = z) \to v_{514} = v_{515} \to v_{515} = v_{516} \to v_{514} = v_{516}$
\item $\varphi_{46} :\equiv \forall z\,(v_{511} = v_{512} \to v_{512} = z \to v_{511} = z) \to v_{511} = v_{512} \to v_{512} = v_{513} \to v_{511} = v_{513}$
\item $\varphi_{47} :\equiv \forall z\,(v_{508} = v_{509} \to v_{509} = z \to v_{508} = z) \to v_{508} = v_{509} \to v_{509} = v_{510} \to v_{508} = v_{510}$
\item $\varphi_{48} :\equiv \forall z\,(v_{505} = v_{506} \to v_{506} = z \to v_{505} = z) \to v_{505} = v_{506} \to v_{506} = v_{507} \to v_{505} = v_{507}$
\item $\varphi_{49} :\equiv \forall z\,(v_{502} = v_{503} \to v_{503} = z \to v_{502} = z) \to v_{502} = v_{503} \to v_{503} = v_{504} \to v_{502} = v_{504}$
\item $\varphi_{50} :\equiv \forall y\,(v_{520} \in \mathbb{R} \to y \in \mathbb{R} \to v_{520} \neq 0 \to v_{520} \cdot y = 1 \to y = v_{520}^{-1})$
\item $\varphi_{51} :\equiv x \in \mathbb{R} \to y \in \mathbb{R} \to x \neq 0 \to y \neq 0 \to x \cdot (x^{-1} \cdot 1) = x \cdot x^{-1}$
\item $\varphi_{52} :\equiv x \in \mathbb{R} \to y \in \mathbb{R} \to x \neq 0 \to y \neq 0 \to x \cdot y \cdot (x^{-1} \cdot y^{-1}) = 1$
\item $\varphi_{53} :\equiv x \in \mathbb{R} \to y \in \mathbb{R} \to x \neq 0 \to y \neq 0 \to x^{-1} \cdot y^{-1} = (x \cdot y)^{-1}$
\item $\varphi_{54} :\equiv x \in \mathbb{R} \to y \in \mathbb{R} \to x \neq 0 \to y \neq 0 \to (x \cdot y)^{-1} = x^{-1} \cdot y^{-1}$
\item $\varphi_{55} :\equiv x \cdot (x^{-1} \cdot 1) = x \cdot x^{-1} \to x \cdot y \cdot (x^{-1} \cdot y^{-1}) = x \cdot x^{-1}$
\item $\varphi_{56} :\equiv x \cdot (1 \cdot x^{-1}) = x \cdot (1 \cdot x^{-1}) \to x \cdot (1 \cdot x^{-1}) = x \cdot (x^{-1} \cdot 1)$
\item $\varphi_{57} :\equiv x \cdot (y \cdot (y^{-1} \cdot x^{-1})) = x \cdot z \to x \cdot (y \cdot (y^{-1} \cdot x^{-1})) = x \cdot w$
\item $\varphi_{58} :\equiv y \in \mathbb{R} \to x \neq 0 \to y \neq 0 \to y \cdot y^{-1} \cdot x^{-1} = y \cdot (y^{-1} \cdot x^{-1})$
\item $\varphi_{59} :\equiv y \in \mathbb{R} \to x \neq 0 \to y \neq 0 \to y \cdot (y^{-1} \cdot x^{-1}) = y \cdot y^{-1} \cdot x^{-1}$
\item $\varphi_{60} :\equiv \forall x \in \mathbb{R}\,(x \neq 0 \to x \cdot x^{-1} = 1) \to x \in \mathbb{R} \to x \neq 0 \to x \cdot x^{-1} = 1$
\item $\varphi_{61} :\equiv \forall x \in \mathbb{R}\,(x \neq 0 \to x \cdot x^{-1} = 1) \to y \in \mathbb{R} \to y \neq 0 \to y \cdot y^{-1} = 1$
\item $\varphi_{62} :\equiv x \in \mathbb{R} \to y \in \mathbb{R} \to x \neq 0 \to y \neq 0 \to x^{-1} \cdot y^{-1} = y^{-1} \cdot x^{-1}$
\item $\varphi_{63} :\equiv \forall x \in \mathbb{R}\,(x \neq 0 \to x^{-1} \in \mathbb{R}) \to x \in \mathbb{R} \to x \neq 0 \to x^{-1} \in \mathbb{R}$
\item $\varphi_{64} :\equiv \forall x \in \mathbb{R}\,(x \neq 0 \to x^{-1} \in \mathbb{R}) \to y \in \mathbb{R} \to y \neq 0 \to y^{-1} \in \mathbb{R}$
\item $\varphi_{65} :\equiv \forall y\,\forall z\,(v_{497} \in \mathbb{R} \to y \in \mathbb{R} \to z \in \mathbb{R} \to v_{497} \cdot y \cdot z = v_{497} \cdot (y \cdot z))$
\item $\varphi_{66} :\equiv \forall w\,(x^{-1} \cdot 1 = w \to x \cdot (x^{-1} \cdot 1) = x \cdot (x^{-1} \cdot 1) \to x \cdot (x^{-1} \cdot 1) = x \cdot w)$
\item $\varphi_{67} :\equiv \forall w\,(1 \cdot x^{-1} = w \to x \cdot (1 \cdot x^{-1}) = x \cdot (1 \cdot x^{-1}) \to x \cdot (1 \cdot x^{-1}) = x \cdot w)$
\item $\varphi_{68} :\equiv x \cdot (1 \cdot x^{-1}) = x \cdot (x^{-1} \cdot 1) \to x \cdot y \cdot (x^{-1} \cdot y^{-1}) = x \cdot (x^{-1} \cdot 1)$
\item $\varphi_{69} :\equiv y^{-1} \in \mathbb{R} \to x^{-1} \in \mathbb{R} \to y \cdot y^{-1} \cdot x^{-1} = y \cdot (y^{-1} \cdot x^{-1})$
\item $\varphi_{70} :\equiv x \cdot (y \cdot (x^{-1} \cdot y^{-1})) = x \cdot (y \cdot z) \to x \cdot (y \cdot (x^{-1} \cdot y^{-1})) = x \cdot (y \cdot w)$
\item $\varphi_{71} :\equiv x^{-1} \cdot y^{-1} \in \mathbb{R} \to x \cdot y \cdot (x^{-1} \cdot y^{-1}) = x \cdot (y \cdot (x^{-1} \cdot y^{-1}))$
\item $\varphi_{72} :\equiv \forall z\,(v_{497} \in \mathbb{R} \to v_{498} \in \mathbb{R} \to z \in \mathbb{R} \to v_{497} \cdot v_{498} \cdot z = v_{497} \cdot (v_{498} \cdot z))$
\item $\varphi_{73} :\equiv x \cdot y \cdot (x^{-1} \cdot y^{-1}) = x \cdot x^{-1} \to x \cdot x^{-1} = 1 \to x \cdot y \cdot (x^{-1} \cdot y^{-1}) = 1$
\item $\varphi_{74} :\equiv \forall v_{522}\,\forall v_{523}\,(v_{522} = v_{523} \to v_{523} = v_{522}) \to \forall v_{523}\,(x^{-1} \cdot y^{-1} = v_{523} \to v_{523} = x^{-1} \cdot y^{-1})$
\item $\varphi_{75} :\equiv \forall v_{521}\,(x \cdot y \in \mathbb{R} \to v_{521} \in \mathbb{R} \to x \cdot y \neq 0 \to x \cdot y \cdot v_{521} = 1 \to v_{521} = (x \cdot y)^{-1})$
\item $\varphi_{76} :\equiv \forall y\,(x \in \mathbb{R} \to y \in \mathbb{R} \to x \cdot y \in \mathbb{R}) \to x \in \mathbb{R} \to y \in \mathbb{R} \to x \cdot y \in \mathbb{R}$
\item $\varphi_{77} :\equiv \forall v_{518}\,\forall v_{519}\,(x \cdot y \cdot (x^{-1} \cdot y^{-1}) = v_{518} \to v_{518} = v_{519} \to x \cdot y \cdot (x^{-1} \cdot y^{-1}) = v_{519})$
\item $\varphi_{78} :\equiv \forall v_{515}\,\forall v_{516}\,(x \cdot y \cdot (x^{-1} \cdot y^{-1}) = v_{515} \to v_{515} = v_{516} \to x \cdot y \cdot (x^{-1} \cdot y^{-1}) = v_{516})$
\item $\varphi_{79} :\equiv \forall v_{512}\,\forall v_{513}\,(x \cdot y \cdot (x^{-1} \cdot y^{-1}) = v_{512} \to v_{512} = v_{513} \to x \cdot y \cdot (x^{-1} \cdot y^{-1}) = v_{513})$
\item $\varphi_{80} :\equiv \forall v_{509}\,\forall v_{510}\,(x \cdot y \cdot (x^{-1} \cdot y^{-1}) = v_{509} \to v_{509} = v_{510} \to x \cdot y \cdot (x^{-1} \cdot y^{-1}) = v_{510})$
\item $\varphi_{81} :\equiv \forall v_{506}\,\forall v_{507}\,(x \cdot y \cdot (x^{-1} \cdot y^{-1}) = v_{506} \to v_{506} = v_{507} \to x \cdot y \cdot (x^{-1} \cdot y^{-1}) = v_{507})$
\item $\varphi_{82} :\equiv y \cdot y^{-1} \cdot x^{-1} = y \cdot (y^{-1} \cdot x^{-1}) \to y \cdot (y^{-1} \cdot x^{-1}) = y \cdot y^{-1} \cdot x^{-1}$
\item $\varphi_{83} :\equiv \forall v_{503}\,\forall v_{504}\,(x \cdot y \cdot (x^{-1} \cdot y^{-1}) = v_{503} \to v_{503} = v_{504} \to x \cdot y \cdot (x^{-1} \cdot y^{-1}) = v_{504})$
\item $\varphi_{84} :\equiv x \cdot (y \cdot y^{-1} \cdot x^{-1}) = x \cdot (1 \cdot x^{-1}) \to x \cdot y \cdot (x^{-1} \cdot y^{-1}) = x \cdot (1 \cdot x^{-1})$
\item $\varphi_{85} :\equiv x \cdot (y \cdot y^{-1} \cdot x^{-1}) = x \cdot (z \cdot x^{-1}) \to x \cdot (y \cdot y^{-1} \cdot x^{-1}) = x \cdot (w \cdot x^{-1})$
\item $\varphi_{86} :\equiv \forall x\,\forall y\,(x \in \mathbb{R} \to y \in \mathbb{R} \to x \cdot y = y \cdot x) \to \forall y\,(1 \in \mathbb{R} \to y \in \mathbb{R} \to 1 \cdot y = y \cdot 1)$
\item $\varphi_{87} :\equiv \forall v_{498}\,\forall v_{499}\,(y \in \mathbb{R} \to v_{498} \in \mathbb{R} \to v_{499} \in \mathbb{R} \to y \cdot v_{498} \cdot v_{499} = y \cdot (v_{498} \cdot v_{499}))$
\item $\varphi_{88} :\equiv \forall v_{499}\,(y \in \mathbb{R} \to y^{-1} \in \mathbb{R} \to v_{499} \in \mathbb{R} \to y \cdot y^{-1} \cdot v_{499} = y \cdot (y^{-1} \cdot v_{499}))$
\item $\varphi_{89} :\equiv \forall x\,\forall y\,(x \in \mathbb{R} \to y \in \mathbb{R} \to x \cdot y \in \mathbb{R}) \to \forall y\,(x \in \mathbb{R} \to y \in \mathbb{R} \to x \cdot y \in \mathbb{R})$
\item $\varphi_{90} :\equiv \forall v_{519}\,(x \cdot y \cdot (x^{-1} \cdot y^{-1}) = x \cdot x^{-1} \to x \cdot x^{-1} = v_{519} \to x \cdot y \cdot (x^{-1} \cdot y^{-1}) = v_{519})$
\item $\varphi_{91} :\equiv \forall y\,(1 \in \mathbb{R} \to y \in \mathbb{R} \to 1 \cdot y = y \cdot 1) \to 1 \in \mathbb{R} \to x^{-1} \in \mathbb{R} \to 1 \cdot x^{-1} = x^{-1} \cdot 1$
\item $\varphi_{92} :\equiv x \cdot (y \cdot y^{-1} \cdot x^{-1}) = x \cdot (y \cdot y^{-1} \cdot x^{-1}) \to x \cdot (y \cdot y^{-1} \cdot x^{-1}) = x \cdot (1 \cdot x^{-1})$
\item $\varphi_{93} :\equiv (x \in \mathbb{R} \to y \in \mathbb{R} \to x \neq 0 \to y \neq 0 \to x \neq 0) \to x \in \mathbb{R} \to y \in \mathbb{R} \to x \neq 0 \to y \neq 0 \to x \cdot x^{-1} = 1$
\item $\varphi_{94} :\equiv (x \in \mathbb{R} \to y \in \mathbb{R} \to x \neq 0 \to y \neq 0 \to y \neq 0) \to x \in \mathbb{R} \to y \in \mathbb{R} \to x \neq 0 \to y \neq 0 \to y \cdot y^{-1} = 1$
\item $\varphi_{95} :\equiv (x \in \mathbb{R} \to y \in \mathbb{R} \to x \neq 0 \to y \neq 0 \to x \neq 0) \to x \in \mathbb{R} \to y \in \mathbb{R} \to x \neq 0 \to y \neq 0 \to x^{-1} \in \mathbb{R}$
\item $\varphi_{96} :\equiv (x \in \mathbb{R} \to y \in \mathbb{R} \to x \neq 0 \to y \neq 0 \to y \neq 0) \to x \in \mathbb{R} \to y \in \mathbb{R} \to x \neq 0 \to y \neq 0 \to y^{-1} \in \mathbb{R}$
\item $\varphi_{97} :\equiv \forall x\,\forall y\,(x \in \mathbb{R} \to y \in \mathbb{R} \to x \cdot y \in \mathbb{R}) \to \forall y\,(x^{-1} \in \mathbb{R} \to y \in \mathbb{R} \to x^{-1} \cdot y \in \mathbb{R})$
\item $\varphi_{98} :\equiv \forall y\,(x \in \mathbb{R} \to y \in \mathbb{R} \to x \neq 0 \to y \neq 0 \to x \cdot y \neq 0) \to x \in \mathbb{R} \to y \in \mathbb{R} \to x \neq 0 \to y \neq 0 \to x \cdot y \neq 0$
\item $\varphi_{99} :\equiv \forall x\,\forall y\,(x \in \mathbb{R} \to y \in \mathbb{R} \to x \cdot y = y \cdot x) \to \forall y\,(x^{-1} \in \mathbb{R} \to y \in \mathbb{R} \to x^{-1} \cdot y = y \cdot x^{-1})$
\item $\varphi_{100} :\equiv \forall v_{516}\,(x \cdot y \cdot (x^{-1} \cdot y^{-1}) = x \cdot (x^{-1} \cdot 1) \to x \cdot (x^{-1} \cdot 1) = v_{516} \to x \cdot y \cdot (x^{-1} \cdot y^{-1}) = v_{516})$
\item $\varphi_{101} :\equiv \forall v_{513}\,(x \cdot y \cdot (x^{-1} \cdot y^{-1}) = x \cdot (1 \cdot x^{-1}) \to x \cdot (1 \cdot x^{-1}) = v_{513} \to x \cdot y \cdot (x^{-1} \cdot y^{-1}) = v_{513})$
\item $\varphi_{102} :\equiv (x \in \mathbb{R} \to y \in \mathbb{R} \to x \neq 0 \to y \neq 0 \to y \in \mathbb{R}) \to x \in \mathbb{R} \to y \in \mathbb{R} \to x \neq 0 \to y \neq 0 \to x \cdot y \in \mathbb{R}$
\item $\varphi_{103} :\equiv x \cdot (y \cdot (y^{-1} \cdot x^{-1})) = x \cdot (y \cdot y^{-1} \cdot x^{-1}) \to x \cdot y \cdot (x^{-1} \cdot y^{-1}) = x \cdot (y \cdot y^{-1} \cdot x^{-1})$
\item $\varphi_{104} :\equiv \varphi_{9} \to x \cdot (y \cdot (y^{-1} \cdot x^{-1})) = x \cdot (y \cdot y^{-1} \cdot x^{-1})$
\item $\varphi_{105} :\equiv \varphi_{11} \to x \cdot y \cdot (x^{-1} \cdot y^{-1}) = x \cdot (y \cdot (y^{-1} \cdot x^{-1}))$
\item $\varphi_{106} :\equiv \forall x\,\forall y\,(x \in \mathbb{R} \to y \in \mathbb{R} \to x \neq 0 \to y \neq 0 \to x \cdot y \neq 0) \to \forall y\,(x \in \mathbb{R} \to y \in \mathbb{R} \to x \neq 0 \to y \neq 0 \to x \cdot y \neq 0)$
\item $\varphi_{107} :\equiv \forall v_{523}\,(x^{-1} \cdot y^{-1} = v_{523} \to v_{523} = x^{-1} \cdot y^{-1}) \to \varphi_{26}$
\item $\varphi_{108} :\equiv (x \in \mathbb{R} \to y \in \mathbb{R} \to x \neq 0 \to y \neq 0 \to x^{-1} \in \mathbb{R}) \to x \in \mathbb{R} \to y \in \mathbb{R} \to x \neq 0 \to y \neq 0 \to x^{-1} \cdot 1 = x^{-1}$
\item $\varphi_{109} :\equiv \forall w\,(y \cdot y^{-1} = w \to x \cdot (y \cdot y^{-1} \cdot x^{-1}) = x \cdot (y \cdot y^{-1} \cdot x^{-1}) \to x \cdot (y \cdot y^{-1} \cdot x^{-1}) = x \cdot (w \cdot x^{-1}))$
\item $\varphi_{110} :\equiv \forall w\,(y \cdot (y^{-1} \cdot x^{-1}) = w \to \varphi_{9} \to x \cdot (y \cdot (y^{-1} \cdot x^{-1})) = x \cdot w)$
\item $\varphi_{111} :\equiv \forall w\,(x^{-1} \cdot y^{-1} = w \to \varphi_{10} \to x \cdot (y \cdot (x^{-1} \cdot y^{-1})) = x \cdot (y \cdot w))$
\item $\varphi_{112} :\equiv \forall y\,(x^{-1} \in \mathbb{R} \to y \in \mathbb{R} \to x^{-1} \cdot y \in \mathbb{R}) \to x^{-1} \in \mathbb{R} \to y^{-1} \in \mathbb{R} \to x^{-1} \cdot y^{-1} \in \mathbb{R}$
\item $\varphi_{113} :\equiv (x \in \mathbb{R} \to y \in \mathbb{R} \to x \neq 0 \to y \neq 0 \to x^{-1} \in \mathbb{R}) \to x \in \mathbb{R} \to y \in \mathbb{R} \to x \neq 0 \to y \neq 0 \to 1 \cdot x^{-1} = x^{-1} \cdot 1$
\item $\varphi_{114} :\equiv (x \in \mathbb{R} \to y \in \mathbb{R} \to x \neq 0 \to y \neq 0 \to y^{-1} \in \mathbb{R}) \to x \in \mathbb{R} \to y \in \mathbb{R} \to x \neq 0 \to y \neq 0 \to x^{-1} \cdot y^{-1} \in \mathbb{R}$
\item $\varphi_{115} :\equiv \forall v_{510}\,(x \cdot y \cdot (x^{-1} \cdot y^{-1}) = x \cdot (y \cdot y^{-1} \cdot x^{-1}) \to x \cdot (y \cdot y^{-1} \cdot x^{-1}) = v_{510} \to x \cdot y \cdot (x^{-1} \cdot y^{-1}) = v_{510})$
\item $\varphi_{116} :\equiv \forall v_{507}\,(x \cdot y \cdot (x^{-1} \cdot y^{-1}) = x \cdot (y \cdot (y^{-1} \cdot x^{-1})) \to x \cdot (y \cdot (y^{-1} \cdot x^{-1})) = v_{507} \to x \cdot y \cdot (x^{-1} \cdot y^{-1}) = v_{507})$
\item $\varphi_{117} :\equiv \forall v_{504}\,(x \cdot y \cdot (x^{-1} \cdot y^{-1}) = x \cdot (y \cdot (x^{-1} \cdot y^{-1})) \to x \cdot (y \cdot (x^{-1} \cdot y^{-1})) = v_{504} \to x \cdot y \cdot (x^{-1} \cdot y^{-1}) = v_{504})$
\item $\varphi_{118} :\equiv \forall y\,(x^{-1} \in \mathbb{R} \to y \in \mathbb{R} \to x^{-1} \cdot y = y \cdot x^{-1}) \to x^{-1} \in \mathbb{R} \to y^{-1} \in \mathbb{R} \to x^{-1} \cdot y^{-1} = y^{-1} \cdot x^{-1}$
\item $\varphi_{119} :\equiv (x \in \mathbb{R} \to y \in \mathbb{R} \to x \neq 0 \to y \neq 0 \to x^{-1} \in \mathbb{R}) \to x \in \mathbb{R} \to \varphi_{58}$
\item $\varphi_{120} :\equiv (x \in \mathbb{R} \to y \in \mathbb{R} \to x \neq 0 \to y \neq 0 \to x^{-1} \cdot y^{-1} \in \mathbb{R}) \to x \in \mathbb{R} \to y \in \mathbb{R} \to x \neq 0 \to \varphi_{36}$
\item $\varphi_{121} :\equiv x \cdot y \cdot (x^{-1} \cdot y^{-1}) = x \cdot (y \cdot (y^{-1} \cdot x^{-1})) \to \varphi_{103}$
\item $\varphi_{122} :\equiv x \cdot y \cdot (x^{-1} \cdot y^{-1}) = x \cdot (y \cdot (x^{-1} \cdot y^{-1})) \to \varphi_{105}$
\item $\varphi_{123} :\equiv (x \in \mathbb{R} \to x \neq 0 \to x \cdot x^{-1} = 1) \to (x \in \mathbb{R} \to y \in \mathbb{R} \to x \neq 0 \to y \neq 0 \to x \in \mathbb{R}) \to \varphi_{93}$
\item $\varphi_{124} :\equiv (y \in \mathbb{R} \to y \neq 0 \to y \cdot y^{-1} = 1) \to (x \in \mathbb{R} \to y \in \mathbb{R} \to x \neq 0 \to y \neq 0 \to y \in \mathbb{R}) \to \varphi_{94}$
\item $\varphi_{125} :\equiv (x \in \mathbb{R} \to y \in \mathbb{R} \to x \neq 0 \to y \neq 0 \to x^{-1} \in \mathbb{R}) \to (x \in \mathbb{R} \to y \in \mathbb{R} \to x \neq 0 \to y \neq 0 \to y^{-1} \in \mathbb{R}) \to \varphi_{62}$
\item $\varphi_{126} :\equiv (1 \in \mathbb{R} \to x^{-1} \in \mathbb{R} \to 1 \cdot x^{-1} = x^{-1} \cdot 1) \to 1 \in \mathbb{R} \to \varphi_{113}$
\item $\varphi_{127} :\equiv x \cdot (x^{-1} \cdot 1) = x \cdot (x^{-1} \cdot 1) \to (x^{-1} \cdot 1 = x^{-1} \to \varphi_{41}) \to \varphi_{51}$
\item $\varphi_{128} :\equiv (x \in \mathbb{R} \to x \neq 0 \to x^{-1} \in \mathbb{R}) \to (x \in \mathbb{R} \to y \in \mathbb{R} \to x \neq 0 \to y \neq 0 \to x \in \mathbb{R}) \to \varphi_{95}$
\item $\varphi_{129} :\equiv (y \in \mathbb{R} \to y \neq 0 \to y^{-1} \in \mathbb{R}) \to (x \in \mathbb{R} \to y \in \mathbb{R} \to x \neq 0 \to y \neq 0 \to y \in \mathbb{R}) \to \varphi_{96}$
\item $\varphi_{130} :\equiv (x \cdot y \cdot (x^{-1} \cdot y^{-1}) = x \cdot (x^{-1} \cdot 1) \to \varphi_{55}) \to x \in \mathbb{R} \to \varphi_{43}$
\item $\varphi_{131} :\equiv \forall v_{497}\,\forall v_{498}\,\forall v_{499}\,(v_{497} \in \mathbb{R} \to v_{498} \in \mathbb{R} \to \varphi_{21}) \to \varphi_{87}$
\item $\varphi_{132} :\equiv (x \in \mathbb{R} \to y \in \mathbb{R} \to x \cdot y \in \mathbb{R}) \to (x \in \mathbb{R} \to y \in \mathbb{R} \to x \neq 0 \to y \neq 0 \to x \in \mathbb{R}) \to \varphi_{102}$
\item $\varphi_{133} :\equiv x \cdot (1 \cdot x^{-1}) = x \cdot (1 \cdot x^{-1}) \to (1 \cdot x^{-1} = x^{-1} \cdot 1 \to \varphi_{56}) \to x \in \mathbb{R} \to \varphi_{34}$
\item $\varphi_{134} :\equiv (x \cdot y \cdot (x^{-1} \cdot y^{-1}) = x \cdot (1 \cdot x^{-1}) \to \varphi_{68}) \to x \in \mathbb{R} \to y \in \mathbb{R} \to \varphi_{23}$
\item $\varphi_{135} :\equiv (x \in \mathbb{R} \to y \in \mathbb{R} \to x \neq 0 \to y \neq 0 \to x \cdot x^{-1} = 1) \to \varphi_{73} \to \varphi_{52}$
\item $\varphi_{136} :\equiv (x \cdot y \cdot (x^{-1} \cdot y^{-1}) = x \cdot (y \cdot y^{-1} \cdot x^{-1}) \to \varphi_{84}) \to x \in \mathbb{R} \to y \in \mathbb{R} \to \varphi_{22}$
\item $\varphi_{137} :\equiv (x^{-1} \cdot y^{-1} = y^{-1} \cdot x^{-1} \to \varphi_{10} \to \varphi_{11}) \to x \in \mathbb{R} \to y \in \mathbb{R} \to x \neq 0 \to y \neq 0 \to \varphi_{11}$
\item $\varphi_{138} :\equiv (x^{-1} \in \mathbb{R} \to y^{-1} \in \mathbb{R} \to x^{-1} \cdot y^{-1} \in \mathbb{R}) \to (x \in \mathbb{R} \to y \in \mathbb{R} \to x \neq 0 \to y \neq 0 \to x^{-1} \in \mathbb{R}) \to \varphi_{114}$
\item $\varphi_{139} :\equiv (y \cdot (y^{-1} \cdot x^{-1}) = y \cdot y^{-1} \cdot x^{-1} \to \varphi_{104}) \to x \in \mathbb{R} \to y \in \mathbb{R} \to x \neq 0 \to \varphi_{35}$
\item $\varphi_{140} :\equiv (x \in \mathbb{R} \to y \in \mathbb{R} \to x \neq 0 \to y \neq 0 \to y \in \mathbb{R}) \to (x \in \mathbb{R} \to y \in \mathbb{R} \to x \neq 0 \to y \neq 0 \to y^{-1} \in \mathbb{R}) \to \varphi_{119}$
\item $\varphi_{141} :\equiv (x \in \mathbb{R} \to y \in \mathbb{R} \to x \neq 0 \to y \neq 0 \to x^{-1} \cdot 1 = x^{-1}) \to \varphi_{127}$
\item $\varphi_{142} :\equiv x \cdot (y \cdot y^{-1} \cdot x^{-1}) = x \cdot (y \cdot y^{-1} \cdot x^{-1}) \to (y \cdot y^{-1} = 1 \to \varphi_{92}) \to x \in \mathbb{R} \to y \in \mathbb{R} \to \varphi_{20}$
\item $\varphi_{143} :\equiv (x \in \mathbb{R} \to y \in \mathbb{R} \to x \neq 0 \to y \neq 0 \to x^{-1} \cdot y^{-1} \in \mathbb{R}) \to (x \in \mathbb{R} \to y \in \mathbb{R} \to x \neq 0 \to y \neq 0 \to x \cdot y \neq 0) \to \varphi_{52} \to \varphi_{53}$
\item $\varphi_{144} :\equiv (x^{-1} \in \mathbb{R} \to y^{-1} \in \mathbb{R} \to x^{-1} \cdot y^{-1} = y^{-1} \cdot x^{-1}) \to \varphi_{125}$
\item $\varphi_{145} :\equiv (x \in \mathbb{R} \to y \in \mathbb{R} \to x \neq 0 \to y \neq 0 \to x \in \mathbb{R}) \to (x \in \mathbb{R} \to y \in \mathbb{R} \to x \neq 0 \to y \neq 0 \to y \in \mathbb{R}) \to \varphi_{120}$
\item $\varphi_{146} :\equiv (x \in \mathbb{R} \to y \in \mathbb{R} \to x \neq 0 \to y \neq 0 \to 1 \cdot x^{-1} = x^{-1} \cdot 1) \to \varphi_{133}$
\item $\varphi_{147} :\equiv \varphi_{115} \to x \cdot y \cdot (x^{-1} \cdot y^{-1}) = x \cdot (y \cdot y^{-1} \cdot x^{-1}) \to \varphi_{84}$
\item $\varphi_{148} :\equiv (x \in \mathbb{R} \to y \in \mathbb{R} \to x \neq 0 \to y \neq 0 \to \varphi_{11}) \to \varphi_{122} \to x \in \mathbb{R} \to y \in \mathbb{R} \to x \neq 0 \to \varphi_{37}$
\item $\varphi_{149} :\equiv (x \in \mathbb{R} \to y \in \mathbb{R} \to x \neq 0 \to \varphi_{35}) \to \varphi_{121} \to x \in \mathbb{R} \to y \in \mathbb{R} \to x \neq 0 \to \varphi_{38}$
\item $\varphi_{150} :\equiv (x \in \mathbb{R} \to y \in \mathbb{R} \to x \neq 0 \to \varphi_{38}) \to (x \in \mathbb{R} \to y \in \mathbb{R} \to \varphi_{20}) \to \varphi_{136}$
\item $\varphi_{151} :\equiv (x \cdot y \in \mathbb{R} \to x^{-1} \cdot y^{-1} \in \mathbb{R} \to x \cdot y \neq 0 \to \varphi_{27}) \to (x \in \mathbb{R} \to y \in \mathbb{R} \to x \neq 0 \to y \neq 0 \to x \cdot y \in \mathbb{R}) \to \varphi_{143}$
\end{abbrevlist}
\end{fullproof}

% 自動生成 (prover/build.js): Abs1  ノード数 9
\begin{fullproof}{[E17]}{$\lvert 1 \rvert = 1$}{9}\phantomsection\label{tree:Abs1}
\noindent$\Pi_{1}$:
\begin{ptree}
\AXC{$\forall x\,(0 \leq x \to \lvert x \rvert = x)$\lbl{[A1]}}
\AXC{$\forall x\,(0 \leq x \to \lvert x \rvert = x) \to 0 \leq 1 \to \lvert 1 \rvert = 1$\lbl{A2}}
\BIC{$0 \leq 1 \to \lvert 1 \rvert = 1$}
\AXC{$(0 \leq 1 \to \lvert 1 \rvert = 1) \to 0 \leq 1 \to \lvert 1 \rvert = 1$\lbl{A1}}
\BIC{$0 \leq 1 \to \lvert 1 \rvert = 1$}
\end{ptree}
\noindent 主木:
\begin{ptree}
\AXC{$0 < 1$\lbl{[F28]}}
\AXC{$0 < 1 \to 0 \leq 1$\lbl{A1}}
\BIC{$0 \leq 1$}
\AXC{$\stackrel{\Pi_{1}}{\vdots}$}\noLine
\UIC{$0 \leq 1 \to \lvert 1 \rvert = 1$}
\BIC{$\lvert 1 \rvert = 1$}
\end{ptree}
\end{fullproof}

% 自動生成 (prover/build.js): AbsInv  ノード数 115
\begin{fullproof}{[E18]}{$\forall x \in \mathbb{R}\,(x \neq 0 \to \lvert x^{-1} \rvert = \lvert x \rvert^{-1})$}{115}\phantomsection\label{tree:AbsInv}
\noindent$\Pi_{1}$:
\begin{ptree}
\AXC{$y = z \to x = y \to x = z$\lbl{A6}}
\AXC{$(y = z \to x = y \to x = z) \to x = y \to y = z \to x = z$\lbl{A1}}
\BIC{$x = y \to y = z \to x = z$}
\UIC{$\forall z\,(x = y \to y = z \to x = z)$}
\UIC{$\forall y\,\forall z\,(x = y \to y = z \to x = z)$}
\UIC{$\forall x\,\forall y\,\forall z\,(x = y \to y = z \to x = z)$}
\AXC{$\forall x\,\forall y\,\forall z\,(x = y \to y = z \to x = z) \to \varphi_{4}$\lbl{A2}}
\BIC{$\varphi_{4}$}
\end{ptree}
\noindent$\Pi_{2}$:
\begin{ptree}
\AXC{$y = z \to x = y \to x = z$\lbl{A6}}
\AXC{$(y = z \to x = y \to x = z) \to x = y \to y = z \to x = z$\lbl{A1}}
\BIC{$x = y \to y = z \to x = z$}
\UIC{$\forall z\,(x = y \to y = z \to x = z)$}
\UIC{$\forall y\,\forall z\,(x = y \to y = z \to x = z)$}
\UIC{$\forall x\,\forall y\,\forall z\,(x = y \to y = z \to x = z)$}
\AXC{$\forall x\,\forall y\,\forall z\,(x = y \to y = z \to x = z) \to \varphi_{4}$\lbl{A2}}
\BIC{$\varphi_{4}$}
\end{ptree}
\noindent$\Pi_{3}$:
\begin{ptree}
\AXC{$z = w \to \lvert x \cdot x^{-1} \rvert = \lvert z \rvert \to \lvert x \cdot x^{-1} \rvert = \lvert w \rvert$\lbl{A6}}
\UIC{$\varphi_{1}$}
\UIC{$\forall z\,(\varphi_{1})$}
\AXC{$\forall z\,(\varphi_{1}) \to \varphi_{14}$\lbl{A2}}
\BIC{$\varphi_{14}$}
\AXC{$\varphi_{14} \to \varphi_{11}$\lbl{A2}}
\BIC{$\varphi_{11}$}
\end{ptree}
\noindent$\Pi_{4}$:
\begin{ptree}
\AXC{$\forall x \in \mathbb{R}\,(x \neq 0 \to x \cdot x^{-1} = 1)$\lbl{(R9)}}
\AXC{$\varphi_{10}$\lbl{A2}}
\BIC{$x \in \mathbb{R} \to x \neq 0 \to x \cdot x^{-1} = 1$}
\AXC{$(x \in \mathbb{R} \to x \neq 0 \to x \cdot x^{-1} = 1) \to \varphi_{25}$\lbl{A1}}
\BIC{$\varphi_{25}$}
\end{ptree}
\noindent$\Pi_{5}$:
\begin{ptree}
\AXC{$x = x$\lbl{A5}}
\AXC{$x = y \to x = x \to y = x$\lbl{A6}}
\AXC{$(x = y \to x = x \to y = x) \to x = x \to x = y \to y = x$\lbl{A1}}
\BIC{$x = x \to x = y \to y = x$}
\BIC{$x = y \to y = x$}
\UIC{$\forall y\,(x = y \to y = x)$}
\UIC{$\forall x\,\forall y\,(x = y \to y = x)$}
\end{ptree}
\noindent$\Pi_{6}$:
\begin{ptree}
\AXC{$\stackrel{\Pi_{5}}{\vdots}$}\noLine
\UIC{$\forall x\,\forall y\,(x = y \to y = x)$}
\AXC{$\varphi_{3}$\lbl{A2}}
\BIC{$\forall y\,(\lvert x \cdot x^{-1} \rvert = y \to y = \lvert x \cdot x^{-1} \rvert)$}
\AXC{$\forall y\,(\lvert x \cdot x^{-1} \rvert = y \to y = \lvert x \cdot x^{-1} \rvert) \to \varphi_{16}$\lbl{A2}}
\BIC{$\varphi_{16}$}
\end{ptree}
\noindent$\Pi_{7}$:
\begin{ptree}
\AXC{$\forall x\,\forall y\,(x \in \mathbb{R} \to y \in \mathbb{R} \to \lvert x \cdot y \rvert = \lvert x \rvert \cdot \lvert y \rvert)$\lbl{[A17]}}
\AXC{$\varphi_{23}$\lbl{A2}}
\BIC{$\forall y\,(x \in \mathbb{R} \to y \in \mathbb{R} \to \lvert x \cdot y \rvert = \lvert x \rvert \cdot \lvert y \rvert)$}
\AXC{$\varphi_{24}$\lbl{A2}}
\BIC{$x \in \mathbb{R} \to x^{-1} \in \mathbb{R} \to \lvert x \cdot x^{-1} \rvert = \lvert x \rvert \cdot \lvert x^{-1} \rvert$}
\AXC{$\varphi_{29}$\lbl{A1}}
\BIC{$(x \in \mathbb{R} \to x \neq 0 \to x \in \mathbb{R}) \to \varphi_{22}$}
\end{ptree}
\noindent$\Pi_{8}$:
\begin{ptree}
\AXC{$\forall x \in \mathbb{R}\,(x \neq 0 \to x^{-1} \in \mathbb{R})$\lbl{(R1)}}
\AXC{$\varphi_{13}$\lbl{A2}}
\BIC{$x \in \mathbb{R} \to x \neq 0 \to x^{-1} \in \mathbb{R}$}
\AXC{$x \in \mathbb{R} \to x \neq 0 \to x \in \mathbb{R}$\lbl{A1}}
\AXC{$\stackrel{\Pi_{7}}{\vdots}$}\noLine
\UIC{$(x \in \mathbb{R} \to x \neq 0 \to x \in \mathbb{R}) \to \varphi_{22}$}
\BIC{$\varphi_{22}$}
\BIC{$x \in \mathbb{R} \to x \neq 0 \to \lvert x \cdot x^{-1} \rvert = \lvert x \rvert \cdot \lvert x^{-1} \rvert$}
\end{ptree}
\noindent$\Pi_{9}$:
\begin{ptree}
\AXC{$\stackrel{\Pi_{6}}{\vdots}$}\noLine
\UIC{$\varphi_{16}$}
\AXC{$\stackrel{\Pi_{8}}{\vdots}$}\noLine
\UIC{$x \in \mathbb{R} \to x \neq 0 \to \lvert x \cdot x^{-1} \rvert = \lvert x \rvert \cdot \lvert x^{-1} \rvert$}
\AXC{$\varphi_{30}$\lbl{A1}}
\BIC{$\varphi_{16} \to x \in \mathbb{R} \to x \neq 0 \to \lvert x \rvert \cdot \lvert x^{-1} \rvert = \lvert x \cdot x^{-1} \rvert$}
\BIC{$x \in \mathbb{R} \to x \neq 0 \to \lvert x \rvert \cdot \lvert x^{-1} \rvert = \lvert x \cdot x^{-1} \rvert$}
\end{ptree}
\noindent$\Pi_{10}$:
\begin{ptree}
\AXC{$\stackrel{\Pi_{9}}{\vdots}$}\noLine
\UIC{$x \in \mathbb{R} \to x \neq 0 \to \lvert x \rvert \cdot \lvert x^{-1} \rvert = \lvert x \cdot x^{-1} \rvert$}
\AXC{$(x \in \mathbb{R} \to x \neq 0 \to \lvert x \rvert \cdot \lvert x^{-1} \rvert = \lvert x \cdot x^{-1} \rvert) \to \varphi_{28}$\lbl{A1}}
\BIC{$\varphi_{28}$}
\end{ptree}
\noindent$\Pi_{11}$:
\begin{ptree}
\AXC{$\stackrel{\Pi_{3}}{\vdots}$}\noLine
\UIC{$\varphi_{11}$}
\AXC{$x = x$\lbl{A5}}
\UIC{$\forall x\,(x = x)$}
\AXC{$\forall x\,(x = x) \to \lvert x \cdot x^{-1} \rvert = \lvert x \cdot x^{-1} \rvert$\lbl{A2}}
\BIC{$\lvert x \cdot x^{-1} \rvert = \lvert x \cdot x^{-1} \rvert$}
\AXC{$\stackrel{\Pi_{4}}{\vdots}$}\noLine
\UIC{$\varphi_{25}$}
\BIC{$\varphi_{11} \to x \in \mathbb{R} \to x \neq 0 \to \lvert x \cdot x^{-1} \rvert = \lvert 1 \rvert$}
\BIC{$x \in \mathbb{R} \to x \neq 0 \to \lvert x \cdot x^{-1} \rvert = \lvert 1 \rvert$}
\AXC{$\stackrel{\Pi_{10}}{\vdots}$}\noLine
\UIC{$\varphi_{28}$}
\BIC{$\varphi_{21} \to x \in \mathbb{R} \to x \neq 0 \to \lvert x \rvert \cdot \lvert x^{-1} \rvert = \lvert 1 \rvert$}
\end{ptree}
\noindent$\Pi_{12}$:
\begin{ptree}
\AXC{$\lvert 1 \rvert = 1$\lbl{[E17]}}
\AXC{$\stackrel{\Pi_{2}}{\vdots}$}\noLine
\UIC{$\varphi_{4}$}
\AXC{$\varphi_{4} \to \varphi_{19}$\lbl{A2}}
\BIC{$\varphi_{19}$}
\AXC{$\varphi_{19} \to \varphi_{21}$\lbl{A2}}
\BIC{$\varphi_{21}$}
\AXC{$\stackrel{\Pi_{11}}{\vdots}$}\noLine
\UIC{$\varphi_{21} \to x \in \mathbb{R} \to x \neq 0 \to \lvert x \rvert \cdot \lvert x^{-1} \rvert = \lvert 1 \rvert$}
\BIC{$x \in \mathbb{R} \to x \neq 0 \to \lvert x \rvert \cdot \lvert x^{-1} \rvert = \lvert 1 \rvert$}
\AXC{$\varphi_{27}$\lbl{A1}}
\BIC{$\lvert 1 \rvert = 1 \to \varphi_{6} \to x \in \mathbb{R} \to x \neq 0 \to \lvert x \rvert \cdot \lvert x^{-1} \rvert = 1$}
\BIC{$\varphi_{6} \to x \in \mathbb{R} \to x \neq 0 \to \lvert x \rvert \cdot \lvert x^{-1} \rvert = 1$}
\end{ptree}
\noindent$\Pi_{13}$:
\begin{ptree}
\AXC{$\stackrel{\Pi_{1}}{\vdots}$}\noLine
\UIC{$\varphi_{4}$}
\AXC{$\varphi_{4} \to \varphi_{7}$\lbl{A2}}
\BIC{$\varphi_{7}$}
\AXC{$\varphi_{7} \to \varphi_{6}$\lbl{A2}}
\BIC{$\varphi_{6}$}
\AXC{$\stackrel{\Pi_{12}}{\vdots}$}\noLine
\UIC{$\varphi_{6} \to x \in \mathbb{R} \to x \neq 0 \to \lvert x \rvert \cdot \lvert x^{-1} \rvert = 1$}
\BIC{$x \in \mathbb{R} \to x \neq 0 \to \lvert x \rvert \cdot \lvert x^{-1} \rvert = 1$}
\end{ptree}
\noindent$\Pi_{14}$:
\begin{ptree}
\AXC{$\stackrel{\Pi_{5}}{\vdots}$}\noLine
\UIC{$\forall x\,\forall y\,(x = y \to y = x)$}
\AXC{$\forall x\,\forall y\,(x = y \to y = x) \to \forall y\,(\lvert x \rvert = y \to y = \lvert x \rvert)$\lbl{A2}}
\BIC{$\forall y\,(\lvert x \rvert = y \to y = \lvert x \rvert)$}
\end{ptree}
\noindent$\Pi_{15}$:
\begin{ptree}
\AXC{$\stackrel{\Pi_{14}}{\vdots}$}\noLine
\UIC{$\forall y\,(\lvert x \rvert = y \to y = \lvert x \rvert)$}
\AXC{$\forall y\,(\lvert x \rvert = y \to y = \lvert x \rvert) \to \lvert x \rvert = 0 \to 0 = \lvert x \rvert$\lbl{A2}}
\BIC{$\lvert x \rvert = 0 \to 0 = \lvert x \rvert$}
\end{ptree}
\noindent$\Pi_{16}$:
\begin{ptree}
\AXC{$\forall x\,\forall y\,(x < y \to x \neq y)$\lbl{[F26]}}
\AXC{$\forall x\,\forall y\,(x < y \to x \neq y) \to \forall y\,(0 < y \to 0 \neq y)$\lbl{A2}}
\BIC{$\forall y\,(0 < y \to 0 \neq y)$}
\end{ptree}
\noindent$\Pi_{17}$:
\begin{ptree}
\AXC{$\stackrel{\Pi_{16}}{\vdots}$}\noLine
\UIC{$\forall y\,(0 < y \to 0 \neq y)$}
\AXC{$\forall y\,(0 < y \to 0 \neq y) \to 0 < \lvert x \rvert \to 0 \neq \lvert x \rvert$\lbl{A2}}
\BIC{$0 < \lvert x \rvert \to 0 \neq \lvert x \rvert$}
\AXC{$(0 < \lvert x \rvert \to 0 \neq \lvert x \rvert) \to \varphi_{5}$\lbl{A1}}
\BIC{$\varphi_{5}$}
\end{ptree}
\noindent$\Pi_{18}$:
\begin{ptree}
\AXC{$\stackrel{\Pi_{15}}{\vdots}$}\noLine
\UIC{$\lvert x \rvert = 0 \to 0 = \lvert x \rvert$}
\AXC{$\forall x \in \mathbb{R}\,(x \neq 0 \to 0 < \lvert x \rvert)$\lbl{[A14]}}
\AXC{$\varphi_{8}$\lbl{A2}}
\BIC{$x \in \mathbb{R} \to x \neq 0 \to 0 < \lvert x \rvert$}
\AXC{$\stackrel{\Pi_{17}}{\vdots}$}\noLine
\UIC{$\varphi_{5}$}
\BIC{$x \in \mathbb{R} \to x \neq 0 \to 0 \neq \lvert x \rvert$}
\AXC{$\varphi_{17}$\lbl{A1}}
\BIC{$(\lvert x \rvert = 0 \to 0 = \lvert x \rvert) \to x \in \mathbb{R} \to x \neq 0 \to \lvert x \rvert \neq 0$}
\BIC{$x \in \mathbb{R} \to x \neq 0 \to \lvert x \rvert \neq 0$}
\end{ptree}
\noindent$\Pi_{19}$:
\begin{ptree}
\AXC{$\forall x \in \mathbb{R}\,(\lvert x \rvert \in \mathbb{R})$\lbl{[A5]}}
\AXC{$\forall x \in \mathbb{R}\,(\lvert x \rvert \in \mathbb{R}) \to x^{-1} \in \mathbb{R} \to \lvert x^{-1} \rvert \in \mathbb{R}$\lbl{A2}}
\BIC{$x^{-1} \in \mathbb{R} \to \lvert x^{-1} \rvert \in \mathbb{R}$}
\end{ptree}
\noindent$\Pi_{20}$:
\begin{ptree}
\AXC{$\forall x \in \mathbb{R}\,(x \neq 0 \to x^{-1} \in \mathbb{R})$\lbl{(R1)}}
\AXC{$\varphi_{13}$\lbl{A2}}
\BIC{$x \in \mathbb{R} \to x \neq 0 \to x^{-1} \in \mathbb{R}$}
\AXC{$\stackrel{\Pi_{19}}{\vdots}$}\noLine
\UIC{$x^{-1} \in \mathbb{R} \to \lvert x^{-1} \rvert \in \mathbb{R}$}
\AXC{$(x^{-1} \in \mathbb{R} \to \lvert x^{-1} \rvert \in \mathbb{R}) \to \varphi_{15}$\lbl{A1}}
\BIC{$\varphi_{15}$}
\BIC{$x \in \mathbb{R} \to x \neq 0 \to \lvert x^{-1} \rvert \in \mathbb{R}$}
\end{ptree}
\noindent$\Pi_{21}$:
\begin{ptree}
\AXC{$\forall x \in \mathbb{R}\,(\lvert x \rvert \in \mathbb{R})$\lbl{[A5]}}
\AXC{$\forall x \in \mathbb{R}\,(\lvert x \rvert \in \mathbb{R}) \to x \in \mathbb{R} \to \lvert x \rvert \in \mathbb{R}$\lbl{A2}}
\BIC{$x \in \mathbb{R} \to \lvert x \rvert \in \mathbb{R}$}
\AXC{$(x \in \mathbb{R} \to \lvert x \rvert \in \mathbb{R}) \to \varphi_{9}$\lbl{A1}}
\BIC{$\varphi_{9}$}
\end{ptree}
\noindent$\Pi_{22}$:
\begin{ptree}
\AXC{$\varphi_{2}$\lbl{[E13]}}
\AXC{$\varphi_{2} \to \varphi_{12}$\lbl{A2}}
\BIC{$\varphi_{12}$}
\AXC{$\varphi_{12} \to \lvert x \rvert \in \mathbb{R} \to \varphi_{18}$\lbl{A2}}
\BIC{$\lvert x \rvert \in \mathbb{R} \to \varphi_{18}$}
\AXC{$(\lvert x \rvert \in \mathbb{R} \to \varphi_{18}) \to (x \in \mathbb{R} \to x \neq 0 \to \lvert x \rvert \in \mathbb{R}) \to \varphi_{26}$\lbl{A1}}
\BIC{$(x \in \mathbb{R} \to x \neq 0 \to \lvert x \rvert \in \mathbb{R}) \to \varphi_{26}$}
\end{ptree}
\noindent$\Pi_{23}$:
\begin{ptree}
\AXC{$\stackrel{\Pi_{20}}{\vdots}$}\noLine
\UIC{$x \in \mathbb{R} \to x \neq 0 \to \lvert x^{-1} \rvert \in \mathbb{R}$}
\AXC{$x \in \mathbb{R} \to x \neq 0 \to x \in \mathbb{R}$\lbl{A1}}
\AXC{$\stackrel{\Pi_{21}}{\vdots}$}\noLine
\UIC{$\varphi_{9}$}
\BIC{$x \in \mathbb{R} \to x \neq 0 \to \lvert x \rvert \in \mathbb{R}$}
\AXC{$\stackrel{\Pi_{22}}{\vdots}$}\noLine
\UIC{$(x \in \mathbb{R} \to x \neq 0 \to \lvert x \rvert \in \mathbb{R}) \to \varphi_{26}$}
\BIC{$\varphi_{26}$}
\BIC{$(x \in \mathbb{R} \to x \neq 0 \to \lvert x \rvert \neq 0) \to \varphi_{20}$}
\end{ptree}
\noindent 主木:
\begin{ptree}
\AXC{$\stackrel{\Pi_{13}}{\vdots}$}\noLine
\UIC{$x \in \mathbb{R} \to x \neq 0 \to \lvert x \rvert \cdot \lvert x^{-1} \rvert = 1$}
\AXC{$\stackrel{\Pi_{18}}{\vdots}$}\noLine
\UIC{$x \in \mathbb{R} \to x \neq 0 \to \lvert x \rvert \neq 0$}
\AXC{$\stackrel{\Pi_{23}}{\vdots}$}\noLine
\UIC{$(x \in \mathbb{R} \to x \neq 0 \to \lvert x \rvert \neq 0) \to \varphi_{20}$}
\BIC{$\varphi_{20}$}
\BIC{$x \in \mathbb{R} \to x \neq 0 \to \lvert x^{-1} \rvert = \lvert x \rvert^{-1}$}
\UIC{$\forall x \in \mathbb{R}\,(x \neq 0 \to \lvert x^{-1} \rvert = \lvert x \rvert^{-1})$}
\end{ptree}
\begin{abbrevlist}
\item $\varphi_{1} :\equiv \forall w\,(z = w \to \lvert x \cdot x^{-1} \rvert = \lvert z \rvert \to \lvert x \cdot x^{-1} \rvert = \lvert w \rvert)$
\item $\varphi_{2} :\equiv \forall x\,\forall y\,(x \in \mathbb{R} \to y \in \mathbb{R} \to x \neq 0 \to x \cdot y = 1 \to y = x^{-1})$
\item $\varphi_{3} :\equiv \forall x\,\forall y\,(x = y \to y = x) \to \forall y\,(\lvert x \cdot x^{-1} \rvert = y \to y = \lvert x \cdot x^{-1} \rvert)$
\item $\varphi_{4} :\equiv \forall y\,\forall z\,(\lvert x \rvert \cdot \lvert x^{-1} \rvert = y \to y = z \to \lvert x \rvert \cdot \lvert x^{-1} \rvert = z)$
\item $\varphi_{5} :\equiv (x \in \mathbb{R} \to x \neq 0 \to 0 < \lvert x \rvert) \to x \in \mathbb{R} \to x \neq 0 \to 0 \neq \lvert x \rvert$
\item $\varphi_{6} :\equiv \lvert x \rvert \cdot \lvert x^{-1} \rvert = \lvert 1 \rvert \to \lvert 1 \rvert = 1 \to \lvert x \rvert \cdot \lvert x^{-1} \rvert = 1$
\item $\varphi_{7} :\equiv \forall z\,(\lvert x \rvert \cdot \lvert x^{-1} \rvert = \lvert 1 \rvert \to \lvert 1 \rvert = z \to \lvert x \rvert \cdot \lvert x^{-1} \rvert = z)$
\item $\varphi_{8} :\equiv \forall x \in \mathbb{R}\,(x \neq 0 \to 0 < \lvert x \rvert) \to x \in \mathbb{R} \to x \neq 0 \to 0 < \lvert x \rvert$
\item $\varphi_{9} :\equiv (x \in \mathbb{R} \to x \neq 0 \to x \in \mathbb{R}) \to x \in \mathbb{R} \to x \neq 0 \to \lvert x \rvert \in \mathbb{R}$
\item $\varphi_{10} :\equiv \forall x \in \mathbb{R}\,(x \neq 0 \to x \cdot x^{-1} = 1) \to x \in \mathbb{R} \to x \neq 0 \to x \cdot x^{-1} = 1$
\item $\varphi_{11} :\equiv x \cdot x^{-1} = 1 \to \lvert x \cdot x^{-1} \rvert = \lvert x \cdot x^{-1} \rvert \to \lvert x \cdot x^{-1} \rvert = \lvert 1 \rvert$
\item $\varphi_{12} :\equiv \forall y\,(\lvert x \rvert \in \mathbb{R} \to y \in \mathbb{R} \to \lvert x \rvert \neq 0 \to \lvert x \rvert \cdot y = 1 \to y = \lvert x \rvert^{-1})$
\item $\varphi_{13} :\equiv \forall x \in \mathbb{R}\,(x \neq 0 \to x^{-1} \in \mathbb{R}) \to x \in \mathbb{R} \to x \neq 0 \to x^{-1} \in \mathbb{R}$
\item $\varphi_{14} :\equiv \forall w\,(x \cdot x^{-1} = w \to \lvert x \cdot x^{-1} \rvert = \lvert x \cdot x^{-1} \rvert \to \lvert x \cdot x^{-1} \rvert = \lvert w \rvert)$
\item $\varphi_{15} :\equiv (x \in \mathbb{R} \to x \neq 0 \to x^{-1} \in \mathbb{R}) \to x \in \mathbb{R} \to x \neq 0 \to \lvert x^{-1} \rvert \in \mathbb{R}$
\item $\varphi_{16} :\equiv \lvert x \cdot x^{-1} \rvert = \lvert x \rvert \cdot \lvert x^{-1} \rvert \to \lvert x \rvert \cdot \lvert x^{-1} \rvert = \lvert x \cdot x^{-1} \rvert$
\item $\varphi_{17} :\equiv (x \in \mathbb{R} \to x \neq 0 \to 0 \neq \lvert x \rvert) \to (\lvert x \rvert = 0 \to 0 = \lvert x \rvert) \to x \in \mathbb{R} \to x \neq 0 \to \lvert x \rvert \neq 0$
\item $\varphi_{18} :\equiv \lvert x^{-1} \rvert \in \mathbb{R} \to \lvert x \rvert \neq 0 \to \lvert x \rvert \cdot \lvert x^{-1} \rvert = 1 \to \lvert x^{-1} \rvert = \lvert x \rvert^{-1}$
\item $\varphi_{19} :\equiv \forall z\,(\lvert x \rvert \cdot \lvert x^{-1} \rvert = \lvert x \cdot x^{-1} \rvert \to \lvert x \cdot x^{-1} \rvert = z \to \lvert x \rvert \cdot \lvert x^{-1} \rvert = z)$
\item $\varphi_{20} :\equiv (x \in \mathbb{R} \to x \neq 0 \to \lvert x \rvert \cdot \lvert x^{-1} \rvert = 1) \to x \in \mathbb{R} \to x \neq 0 \to \lvert x^{-1} \rvert = \lvert x \rvert^{-1}$
\item $\varphi_{21} :\equiv \lvert x \rvert \cdot \lvert x^{-1} \rvert = \lvert x \cdot x^{-1} \rvert \to \lvert x \cdot x^{-1} \rvert = \lvert 1 \rvert \to \lvert x \rvert \cdot \lvert x^{-1} \rvert = \lvert 1 \rvert$
\item $\varphi_{22} :\equiv (x \in \mathbb{R} \to x \neq 0 \to x^{-1} \in \mathbb{R}) \to x \in \mathbb{R} \to x \neq 0 \to \lvert x \cdot x^{-1} \rvert = \lvert x \rvert \cdot \lvert x^{-1} \rvert$
\item $\varphi_{23} :\equiv \forall x\,\forall y\,(x \in \mathbb{R} \to y \in \mathbb{R} \to \lvert x \cdot y \rvert = \lvert x \rvert \cdot \lvert y \rvert) \to \forall y\,(x \in \mathbb{R} \to y \in \mathbb{R} \to \lvert x \cdot y \rvert = \lvert x \rvert \cdot \lvert y \rvert)$
\item $\varphi_{24} :\equiv \forall y\,(x \in \mathbb{R} \to y \in \mathbb{R} \to \lvert x \cdot y \rvert = \lvert x \rvert \cdot \lvert y \rvert) \to x \in \mathbb{R} \to x^{-1} \in \mathbb{R} \to \lvert x \cdot x^{-1} \rvert = \lvert x \rvert \cdot \lvert x^{-1} \rvert$
\item $\varphi_{25} :\equiv \lvert x \cdot x^{-1} \rvert = \lvert x \cdot x^{-1} \rvert \to \varphi_{11} \to x \in \mathbb{R} \to x \neq 0 \to \lvert x \cdot x^{-1} \rvert = \lvert 1 \rvert$
\item $\varphi_{26} :\equiv (x \in \mathbb{R} \to x \neq 0 \to \lvert x^{-1} \rvert \in \mathbb{R}) \to (x \in \mathbb{R} \to x \neq 0 \to \lvert x \rvert \neq 0) \to \varphi_{20}$
\item $\varphi_{27} :\equiv (x \in \mathbb{R} \to x \neq 0 \to \lvert x \rvert \cdot \lvert x^{-1} \rvert = \lvert 1 \rvert) \to \lvert 1 \rvert = 1 \to \varphi_{6} \to x \in \mathbb{R} \to x \neq 0 \to \lvert x \rvert \cdot \lvert x^{-1} \rvert = 1$
\item $\varphi_{28} :\equiv (x \in \mathbb{R} \to x \neq 0 \to \lvert x \cdot x^{-1} \rvert = \lvert 1 \rvert) \to \varphi_{21} \to x \in \mathbb{R} \to x \neq 0 \to \lvert x \rvert \cdot \lvert x^{-1} \rvert = \lvert 1 \rvert$
\item $\varphi_{29} :\equiv (x \in \mathbb{R} \to x^{-1} \in \mathbb{R} \to \lvert x \cdot x^{-1} \rvert = \lvert x \rvert \cdot \lvert x^{-1} \rvert) \to (x \in \mathbb{R} \to x \neq 0 \to x \in \mathbb{R}) \to \varphi_{22}$
\item $\varphi_{30} :\equiv (x \in \mathbb{R} \to x \neq 0 \to \lvert x \cdot x^{-1} \rvert = \lvert x \rvert \cdot \lvert x^{-1} \rvert) \to \varphi_{16} \to x \in \mathbb{R} \to x \neq 0 \to \lvert x \rvert \cdot \lvert x^{-1} \rvert = \lvert x \cdot x^{-1} \rvert$
\end{abbrevlist}
\end{fullproof}

% 自動生成 (prover/build.js): InvLt  ノード数 437
\begin{fullproof}{[E19]}{$\forall x\,\forall y\,(x \in \mathbb{R} \to y \in \mathbb{R} \to 0 < x \to x < y \to y^{-1} < x^{-1})$}{437}\phantomsection\label{tree:InvLt}
\noindent$\Pi_{1}$:
\begin{ptree}
\AXC{$z = w \to y^{-1} < z \to y^{-1} < w$\lbl{A6}}
\UIC{$\forall w\,(z = w \to y^{-1} < z \to y^{-1} < w)$}
\UIC{$\forall z\,\forall w\,(z = w \to y^{-1} < z \to y^{-1} < w)$}
\AXC{$\forall z\,\forall w\,(z = w \to y^{-1} < z \to y^{-1} < w) \to \varphi_{58}$\lbl{A2}}
\BIC{$\varphi_{58}$}
\AXC{$\varphi_{58} \to \varphi_{89}$\lbl{A2}}
\BIC{$\varphi_{89}$}
\end{ptree}
\noindent$\Pi_{2}$:
\begin{ptree}
\AXC{$z = w \to z < x^{-1} \cdot y^{-1} \cdot y \to w < x^{-1} \cdot y^{-1} \cdot y$\lbl{A6}}
\UIC{$\varphi_{34}$}
\UIC{$\forall z\,(\varphi_{34})$}
\AXC{$\forall z\,(\varphi_{34}) \to \varphi_{115}$\lbl{A2}}
\BIC{$\varphi_{115}$}
\AXC{$\varphi_{115} \to x^{-1} \cdot y^{-1} \cdot x = y^{-1} \to \varphi_{88}$\lbl{A2}}
\BIC{$x^{-1} \cdot y^{-1} \cdot x = y^{-1} \to \varphi_{88}$}
\end{ptree}
\noindent$\Pi_{3}$:
\begin{ptree}
\AXC{$\forall x\,\forall y\,\forall z\,(x < y \to y < z \to x < z)$\lbl{[O2]}}
\AXC{$\varphi_{35}$\lbl{A2}}
\BIC{$\forall y\,\forall z\,(0 < y \to y < z \to 0 < z)$}
\AXC{$\varphi_{20}$\lbl{A2}}
\BIC{$\forall z\,(0 < x \to x < z \to 0 < z)$}
\AXC{$\forall z\,(0 < x \to x < z \to 0 < z) \to 0 < x \to x < y \to 0 < y$\lbl{A2}}
\BIC{$0 < x \to x < y \to 0 < y$}
\end{ptree}
\noindent$\Pi_{4}$:
\begin{ptree}
\AXC{$x \in \mathbb{R} \to y \in \mathbb{R} \to 0 < x \to x < y \to 0 < x$\lbl{A1}}
\AXC{$\stackrel{\Pi_{3}}{\vdots}$}\noLine
\UIC{$0 < x \to x < y \to 0 < y$}
\AXC{$\varphi_{140}$\lbl{A1}}
\BIC{$(x \in \mathbb{R} \to y \in \mathbb{R} \to 0 < x \to x < y \to 0 < x) \to \varphi_{114}$}
\BIC{$\varphi_{114}$}
\end{ptree}
\noindent$\Pi_{5}$:
\begin{ptree}
\AXC{$x \in \mathbb{R} \to y \in \mathbb{R} \to 0 < x \to x < y \to x < y$\lbl{A1}}
\AXC{$\stackrel{\Pi_{4}}{\vdots}$}\noLine
\UIC{$\varphi_{114}$}
\BIC{$x \in \mathbb{R} \to y \in \mathbb{R} \to 0 < x \to x < y \to 0 < y$}
\end{ptree}
\noindent$\Pi_{6}$:
\begin{ptree}
\AXC{$x \in \mathbb{R} \to y \in \mathbb{R} \to 0 < x \to x < y \to y \in \mathbb{R}$\lbl{A1}}
\AXC{$\forall x \in \mathbb{R}\,(0 < x \to 0 < x^{-1})$\lbl{[F34]}}
\AXC{$\varphi_{51}$\lbl{A2}}
\BIC{$y \in \mathbb{R} \to 0 < y \to 0 < y^{-1}$}
\AXC{$\varphi_{142}$\lbl{A1}}
\BIC{$(x \in \mathbb{R} \to y \in \mathbb{R} \to 0 < x \to x < y \to y \in \mathbb{R}) \to \varphi_{121}$}
\BIC{$\varphi_{121}$}
\end{ptree}
\noindent$\Pi_{7}$:
\begin{ptree}
\AXC{$x \in \mathbb{R} \to y \in \mathbb{R} \to 0 < x \to x < y \to x \in \mathbb{R}$\lbl{A1}}
\AXC{$\forall x \in \mathbb{R}\,(0 < x \to 0 < x^{-1})$\lbl{[F34]}}
\AXC{$\varphi_{52}$\lbl{A2}}
\BIC{$x \in \mathbb{R} \to 0 < x \to 0 < x^{-1}$}
\AXC{$\varphi_{143}$\lbl{A1}}
\BIC{$(x \in \mathbb{R} \to y \in \mathbb{R} \to 0 < x \to x < y \to x \in \mathbb{R}) \to \varphi_{122}$}
\BIC{$\varphi_{122}$}
\end{ptree}
\noindent$\Pi_{8}$:
\begin{ptree}
\AXC{$x = x$\lbl{A5}}
\AXC{$x = y \to x = x \to y = x$\lbl{A6}}
\AXC{$(x = y \to x = x \to y = x) \to x = x \to x = y \to y = x$\lbl{A1}}
\BIC{$x = x \to x = y \to y = x$}
\BIC{$x = y \to y = x$}
\UIC{$\forall y\,(x = y \to y = x)$}
\UIC{$\forall x\,\forall y\,(x = y \to y = x)$}
\end{ptree}
\noindent$\Pi_{9}$:
\begin{ptree}
\AXC{$\stackrel{\Pi_{8}}{\vdots}$}\noLine
\UIC{$\forall x\,\forall y\,(x = y \to y = x)$}
\AXC{$\forall x\,\forall y\,(x = y \to y = x) \to \forall y\,(v_{524} = y \to y = v_{524})$\lbl{A2}}
\BIC{$\forall y\,(v_{524} = y \to y = v_{524})$}
\AXC{$\varphi_{1}$\lbl{A2}}
\BIC{$v_{524} = v_{525} \to v_{525} = v_{524}$}
\UIC{$\forall v_{525}\,(v_{524} = v_{525} \to v_{525} = v_{524})$}
\UIC{$\forall v_{524}\,\forall v_{525}\,(v_{524} = v_{525} \to v_{525} = v_{524})$}
\AXC{$\varphi_{19}$\lbl{A2}}
\BIC{$\forall v_{525}\,(y = v_{525} \to v_{525} = y)$}
\end{ptree}
\noindent$\Pi_{10}$:
\begin{ptree}
\AXC{$\forall x\,\forall y\,(x < y \to x \neq y)$\lbl{[F26]}}
\AXC{$\forall x\,\forall y\,(x < y \to x \neq y) \to \forall y\,(0 < y \to 0 \neq y)$\lbl{A2}}
\BIC{$\forall y\,(0 < y \to 0 \neq y)$}
\AXC{$\forall y\,(0 < y \to 0 \neq y) \to 0 < y \to 0 \neq y$\lbl{A2}}
\BIC{$0 < y \to 0 \neq y$}
\end{ptree}
\noindent$\Pi_{11}$:
\begin{ptree}
\AXC{$\stackrel{\Pi_{5}}{\vdots}$}\noLine
\UIC{$x \in \mathbb{R} \to y \in \mathbb{R} \to 0 < x \to x < y \to 0 < y$}
\AXC{$\stackrel{\Pi_{10}}{\vdots}$}\noLine
\UIC{$0 < y \to 0 \neq y$}
\AXC{$(0 < y \to 0 \neq y) \to \varphi_{112}$\lbl{A1}}
\BIC{$\varphi_{112}$}
\BIC{$x \in \mathbb{R} \to y \in \mathbb{R} \to 0 < x \to x < y \to 0 \neq y$}
\end{ptree}
\noindent$\Pi_{12}$:
\begin{ptree}
\AXC{$\stackrel{\Pi_{11}}{\vdots}$}\noLine
\UIC{$x \in \mathbb{R} \to y \in \mathbb{R} \to 0 < x \to x < y \to 0 \neq y$}
\AXC{$(x \in \mathbb{R} \to y \in \mathbb{R} \to 0 < x \to x < y \to 0 \neq y) \to \varphi_{21}$\lbl{A1}}
\BIC{$\varphi_{21}$}
\end{ptree}
\noindent$\Pi_{13}$:
\begin{ptree}
\AXC{$\stackrel{\Pi_{9}}{\vdots}$}\noLine
\UIC{$\forall v_{525}\,(y = v_{525} \to v_{525} = y)$}
\AXC{$\forall v_{525}\,(y = v_{525} \to v_{525} = y) \to y = 0 \to 0 = y$\lbl{A2}}
\BIC{$y = 0 \to 0 = y$}
\AXC{$\stackrel{\Pi_{12}}{\vdots}$}\noLine
\UIC{$\varphi_{21}$}
\BIC{$x \in \mathbb{R} \to y \in \mathbb{R} \to 0 < x \to x < y \to y \neq 0$}
\end{ptree}
\noindent$\Pi_{14}$:
\begin{ptree}
\AXC{$x \in \mathbb{R} \to y \in \mathbb{R} \to 0 < x \to x < y \to y \in \mathbb{R}$\lbl{A1}}
\AXC{$\forall x \in \mathbb{R}\,(x \neq 0 \to x^{-1} \in \mathbb{R})$\lbl{(R1)}}
\AXC{$\varphi_{67}$\lbl{A2}}
\BIC{$y \in \mathbb{R} \to y \neq 0 \to y^{-1} \in \mathbb{R}$}
\AXC{$\varphi_{146}$\lbl{A1}}
\BIC{$(x \in \mathbb{R} \to y \in \mathbb{R} \to 0 < x \to x < y \to y \in \mathbb{R}) \to \varphi_{125}$}
\BIC{$\varphi_{125}$}
\end{ptree}
\noindent$\Pi_{15}$:
\begin{ptree}
\AXC{$\stackrel{\Pi_{13}}{\vdots}$}\noLine
\UIC{$x \in \mathbb{R} \to y \in \mathbb{R} \to 0 < x \to x < y \to y \neq 0$}
\AXC{$\stackrel{\Pi_{14}}{\vdots}$}\noLine
\UIC{$\varphi_{125}$}
\BIC{$x \in \mathbb{R} \to y \in \mathbb{R} \to 0 < x \to x < y \to y^{-1} \in \mathbb{R}$}
\end{ptree}
\noindent$\Pi_{16}$:
\begin{ptree}
\AXC{$\stackrel{\Pi_{8}}{\vdots}$}\noLine
\UIC{$\forall x\,\forall y\,(x = y \to y = x)$}
\AXC{$\forall x\,\forall y\,(x = y \to y = x) \to \forall y\,(x = y \to y = x)$\lbl{A2}}
\BIC{$\forall y\,(x = y \to y = x)$}
\AXC{$\forall y\,(x = y \to y = x) \to x = 0 \to 0 = x$\lbl{A2}}
\BIC{$x = 0 \to 0 = x$}
\end{ptree}
\noindent$\Pi_{17}$:
\begin{ptree}
\AXC{$\forall x\,\forall y\,(x < y \to x \neq y)$\lbl{[F26]}}
\AXC{$\forall x\,\forall y\,(x < y \to x \neq y) \to \forall y\,(0 < y \to 0 \neq y)$\lbl{A2}}
\BIC{$\forall y\,(0 < y \to 0 \neq y)$}
\AXC{$\forall y\,(0 < y \to 0 \neq y) \to 0 < x \to 0 \neq x$\lbl{A2}}
\BIC{$0 < x \to 0 \neq x$}
\end{ptree}
\noindent$\Pi_{18}$:
\begin{ptree}
\AXC{$x \in \mathbb{R} \to y \in \mathbb{R} \to 0 < x \to x < y \to 0 < x$\lbl{A1}}
\AXC{$\stackrel{\Pi_{17}}{\vdots}$}\noLine
\UIC{$0 < x \to 0 \neq x$}
\AXC{$(0 < x \to 0 \neq x) \to \varphi_{113}$\lbl{A1}}
\BIC{$\varphi_{113}$}
\BIC{$x \in \mathbb{R} \to y \in \mathbb{R} \to 0 < x \to x < y \to 0 \neq x$}
\end{ptree}
\noindent$\Pi_{19}$:
\begin{ptree}
\AXC{$\stackrel{\Pi_{16}}{\vdots}$}\noLine
\UIC{$x = 0 \to 0 = x$}
\AXC{$\stackrel{\Pi_{18}}{\vdots}$}\noLine
\UIC{$x \in \mathbb{R} \to y \in \mathbb{R} \to 0 < x \to x < y \to 0 \neq x$}
\AXC{$(x \in \mathbb{R} \to y \in \mathbb{R} \to 0 < x \to x < y \to 0 \neq x) \to \varphi_{22}$\lbl{A1}}
\BIC{$\varphi_{22}$}
\BIC{$x \in \mathbb{R} \to y \in \mathbb{R} \to 0 < x \to x < y \to x \neq 0$}
\end{ptree}
\noindent$\Pi_{20}$:
\begin{ptree}
\AXC{$x \in \mathbb{R} \to y \in \mathbb{R} \to 0 < x \to x < y \to x \in \mathbb{R}$\lbl{A1}}
\AXC{$\forall x \in \mathbb{R}\,(x \neq 0 \to x^{-1} \in \mathbb{R})$\lbl{(R1)}}
\AXC{$\varphi_{68}$\lbl{A2}}
\BIC{$x \in \mathbb{R} \to x \neq 0 \to x^{-1} \in \mathbb{R}$}
\AXC{$\varphi_{147}$\lbl{A1}}
\BIC{$(x \in \mathbb{R} \to y \in \mathbb{R} \to 0 < x \to x < y \to x \in \mathbb{R}) \to \varphi_{126}$}
\BIC{$\varphi_{126}$}
\end{ptree}
\noindent$\Pi_{21}$:
\begin{ptree}
\AXC{$\stackrel{\Pi_{19}}{\vdots}$}\noLine
\UIC{$x \in \mathbb{R} \to y \in \mathbb{R} \to 0 < x \to x < y \to x \neq 0$}
\AXC{$\stackrel{\Pi_{20}}{\vdots}$}\noLine
\UIC{$\varphi_{126}$}
\BIC{$x \in \mathbb{R} \to y \in \mathbb{R} \to 0 < x \to x < y \to x^{-1} \in \mathbb{R}$}
\end{ptree}
\noindent$\Pi_{22}$:
\begin{ptree}
\AXC{$\forall x\,\forall y\,(x \in \mathbb{R} \to y \in \mathbb{R} \to 0 < x \to 0 < y \to 0 < x \cdot y)$\lbl{(R15)}}
\AXC{$\varphi_{136}$\lbl{A2}}
\BIC{$\varphi_{57}$}
\AXC{$\varphi_{57} \to \varphi_{92}$\lbl{A2}}
\BIC{$\varphi_{92}$}
\end{ptree}
\noindent$\Pi_{23}$:
\begin{ptree}
\AXC{$\stackrel{\Pi_{21}}{\vdots}$}\noLine
\UIC{$x \in \mathbb{R} \to y \in \mathbb{R} \to 0 < x \to x < y \to x^{-1} \in \mathbb{R}$}
\AXC{$\stackrel{\Pi_{22}}{\vdots}$}\noLine
\UIC{$\varphi_{92}$}
\AXC{$\varphi_{92} \to (x \in \mathbb{R} \to y \in \mathbb{R} \to 0 < x \to x < y \to x^{-1} \in \mathbb{R}) \to \varphi_{156}$\lbl{A1}}
\BIC{$(x \in \mathbb{R} \to y \in \mathbb{R} \to 0 < x \to x < y \to x^{-1} \in \mathbb{R}) \to \varphi_{156}$}
\BIC{$\varphi_{156}$}
\end{ptree}
\noindent$\Pi_{24}$:
\begin{ptree}
\AXC{$x \in \mathbb{R} \to y \in \mathbb{R} \to 0 < x \to x < y \to 0 < x$\lbl{A1}}
\AXC{$\stackrel{\Pi_{7}}{\vdots}$}\noLine
\UIC{$\varphi_{122}$}
\BIC{$x \in \mathbb{R} \to y \in \mathbb{R} \to 0 < x \to x < y \to 0 < x^{-1}$}
\AXC{$\stackrel{\Pi_{15}}{\vdots}$}\noLine
\UIC{$x \in \mathbb{R} \to y \in \mathbb{R} \to 0 < x \to x < y \to y^{-1} \in \mathbb{R}$}
\AXC{$\stackrel{\Pi_{23}}{\vdots}$}\noLine
\UIC{$\varphi_{156}$}
\BIC{$(x \in \mathbb{R} \to y \in \mathbb{R} \to 0 < x \to x < y \to 0 < x^{-1}) \to \varphi_{130}$}
\BIC{$\varphi_{130}$}
\end{ptree}
\noindent$\Pi_{25}$:
\begin{ptree}
\AXC{$\stackrel{\Pi_{5}}{\vdots}$}\noLine
\UIC{$x \in \mathbb{R} \to y \in \mathbb{R} \to 0 < x \to x < y \to 0 < y$}
\AXC{$\stackrel{\Pi_{6}}{\vdots}$}\noLine
\UIC{$\varphi_{121}$}
\BIC{$x \in \mathbb{R} \to y \in \mathbb{R} \to 0 < x \to x < y \to 0 < y^{-1}$}
\AXC{$\stackrel{\Pi_{24}}{\vdots}$}\noLine
\UIC{$\varphi_{130}$}
\BIC{$x \in \mathbb{R} \to y \in \mathbb{R} \to 0 < x \to x < y \to 0 < x^{-1} \cdot y^{-1}$}
\end{ptree}
\noindent$\Pi_{26}$:
\begin{ptree}
\AXC{$\forall x\,\forall y\,(x \in \mathbb{R} \to y \in \mathbb{R} \to x \cdot y \in \mathbb{R})$\lbl{(R1)}}
\AXC{$\varphi_{116}$\lbl{A2}}
\BIC{$\forall y\,(x^{-1} \in \mathbb{R} \to y \in \mathbb{R} \to x^{-1} \cdot y \in \mathbb{R})$}
\AXC{$\varphi_{129}$\lbl{A2}}
\BIC{$x^{-1} \in \mathbb{R} \to y^{-1} \in \mathbb{R} \to x^{-1} \cdot y^{-1} \in \mathbb{R}$}
\AXC{$\varphi_{160}$\lbl{A1}}
\BIC{$(x \in \mathbb{R} \to y \in \mathbb{R} \to 0 < x \to x < y \to x^{-1} \in \mathbb{R}) \to \varphi_{135}$}
\end{ptree}
\noindent$\Pi_{27}$:
\begin{ptree}
\AXC{$\stackrel{\Pi_{21}}{\vdots}$}\noLine
\UIC{$x \in \mathbb{R} \to y \in \mathbb{R} \to 0 < x \to x < y \to x^{-1} \in \mathbb{R}$}
\AXC{$\stackrel{\Pi_{26}}{\vdots}$}\noLine
\UIC{$(x \in \mathbb{R} \to y \in \mathbb{R} \to 0 < x \to x < y \to x^{-1} \in \mathbb{R}) \to \varphi_{135}$}
\BIC{$\varphi_{135}$}
\end{ptree}
\noindent$\Pi_{28}$:
\begin{ptree}
\AXC{$\stackrel{\Pi_{15}}{\vdots}$}\noLine
\UIC{$x \in \mathbb{R} \to y \in \mathbb{R} \to 0 < x \to x < y \to y^{-1} \in \mathbb{R}$}
\AXC{$\stackrel{\Pi_{27}}{\vdots}$}\noLine
\UIC{$\varphi_{135}$}
\BIC{$x \in \mathbb{R} \to y \in \mathbb{R} \to 0 < x \to x < y \to x^{-1} \cdot y^{-1} \in \mathbb{R}$}
\end{ptree}
\noindent$\Pi_{29}$:
\begin{ptree}
\AXC{$\forall x\,\forall y\,(\varphi_{69})$\lbl{[E8]}}
\AXC{$\forall x\,\forall y\,(\varphi_{69}) \to \forall y\,(\varphi_{69})$\lbl{A2}}
\BIC{$\forall y\,(\varphi_{69})$}
\AXC{$\forall y\,(\varphi_{69}) \to \varphi_{69}$\lbl{A2}}
\BIC{$\varphi_{69}$}
\AXC{$\varphi_{69} \to x \in \mathbb{R} \to y \in \mathbb{R} \to x^{-1} \cdot y^{-1} \in \mathbb{R} \to \varphi_{83}$\lbl{A2}}
\BIC{$x \in \mathbb{R} \to y \in \mathbb{R} \to x^{-1} \cdot y^{-1} \in \mathbb{R} \to \varphi_{83}$}
\end{ptree}
\noindent$\Pi_{30}$:
\begin{ptree}
\AXC{$\stackrel{\Pi_{29}}{\vdots}$}\noLine
\UIC{$x \in \mathbb{R} \to y \in \mathbb{R} \to x^{-1} \cdot y^{-1} \in \mathbb{R} \to \varphi_{83}$}
\AXC{$(x \in \mathbb{R} \to y \in \mathbb{R} \to x^{-1} \cdot y^{-1} \in \mathbb{R} \to \varphi_{83}) \to \varphi_{168}$\lbl{A1}}
\BIC{$\varphi_{168}$}
\end{ptree}
\noindent$\Pi_{31}$:
\begin{ptree}
\AXC{$x \in \mathbb{R} \to y \in \mathbb{R} \to 0 < x \to x < y \to y \in \mathbb{R}$\lbl{A1}}
\AXC{$x \in \mathbb{R} \to y \in \mathbb{R} \to 0 < x \to x < y \to x \in \mathbb{R}$\lbl{A1}}
\AXC{$\stackrel{\Pi_{30}}{\vdots}$}\noLine
\UIC{$\varphi_{168}$}
\BIC{$(x \in \mathbb{R} \to y \in \mathbb{R} \to 0 < x \to x < y \to y \in \mathbb{R}) \to \varphi_{163}$}
\BIC{$\varphi_{163}$}
\end{ptree}
\noindent$\Pi_{32}$:
\begin{ptree}
\AXC{$\stackrel{\Pi_{25}}{\vdots}$}\noLine
\UIC{$x \in \mathbb{R} \to y \in \mathbb{R} \to 0 < x \to x < y \to 0 < x^{-1} \cdot y^{-1}$}
\AXC{$\stackrel{\Pi_{28}}{\vdots}$}\noLine
\UIC{$x \in \mathbb{R} \to y \in \mathbb{R} \to 0 < x \to x < y \to x^{-1} \cdot y^{-1} \in \mathbb{R}$}
\AXC{$\stackrel{\Pi_{31}}{\vdots}$}\noLine
\UIC{$\varphi_{163}$}
\BIC{$\varphi_{148}$}
\BIC{$(x \in \mathbb{R} \to y \in \mathbb{R} \to 0 < x \to x < y \to x < y) \to x \in \mathbb{R} \to \varphi_{75}$}
\end{ptree}
\noindent$\Pi_{33}$:
\begin{ptree}
\AXC{$x \in \mathbb{R} \to y \in \mathbb{R} \to 0 < x \to x < y \to x < y$\lbl{A1}}
\AXC{$\stackrel{\Pi_{32}}{\vdots}$}\noLine
\UIC{$(x \in \mathbb{R} \to y \in \mathbb{R} \to 0 < x \to x < y \to x < y) \to x \in \mathbb{R} \to \varphi_{75}$}
\BIC{$x \in \mathbb{R} \to \varphi_{75}$}
\end{ptree}
\noindent$\Pi_{34}$:
\begin{ptree}
\AXC{$y = z \to x = y \to x = z$\lbl{A6}}
\AXC{$(y = z \to x = y \to x = z) \to x = y \to y = z \to x = z$\lbl{A1}}
\BIC{$x = y \to y = z \to x = z$}
\UIC{$\forall z\,(x = y \to y = z \to x = z)$}
\UIC{$\forall y\,\forall z\,(x = y \to y = z \to x = z)$}
\UIC{$\forall x\,\forall y\,\forall z\,(x = y \to y = z \to x = z)$}
\AXC{$\varphi_{12}$\lbl{A2}}
\BIC{$\forall y\,\forall z\,(v_{539} = y \to y = z \to v_{539} = z)$}
\AXC{$\varphi_{23}$\lbl{A2}}
\BIC{$\forall z\,(v_{539} = v_{540} \to v_{540} = z \to v_{539} = z)$}
\AXC{$\varphi_{36}$\lbl{A2}}
\BIC{$v_{539} = v_{540} \to v_{540} = v_{541} \to v_{539} = v_{541}$}
\UIC{$\forall v_{541}\,(v_{539} = v_{540} \to v_{540} = v_{541} \to v_{539} = v_{541})$}
\UIC{$\varphi_{3}$}
\UIC{$\forall v_{539}\,(\varphi_{3})$}
\end{ptree}
\noindent$\Pi_{35}$:
\begin{ptree}
\AXC{$\stackrel{\Pi_{34}}{\vdots}$}\noLine
\UIC{$\forall v_{539}\,(\varphi_{3})$}
\AXC{$\forall v_{539}\,(\varphi_{3}) \to \varphi_{76}$\lbl{A2}}
\BIC{$\varphi_{76}$}
\AXC{$\varphi_{76} \to \varphi_{95}$\lbl{A2}}
\BIC{$\varphi_{95}$}
\AXC{$\varphi_{95} \to \varphi_{93}$\lbl{A2}}
\BIC{$\varphi_{93}$}
\end{ptree}
\noindent$\Pi_{36}$:
\begin{ptree}
\AXC{$\forall x \in \mathbb{R}\,(x \cdot 1 = x)$\lbl{(R8)}}
\AXC{$\forall x \in \mathbb{R}\,(x \cdot 1 = x) \to y^{-1} \in \mathbb{R} \to y^{-1} \cdot 1 = y^{-1}$\lbl{A2}}
\BIC{$y^{-1} \in \mathbb{R} \to y^{-1} \cdot 1 = y^{-1}$}
\end{ptree}
\noindent$\Pi_{37}$:
\begin{ptree}
\AXC{$\stackrel{\Pi_{15}}{\vdots}$}\noLine
\UIC{$x \in \mathbb{R} \to y \in \mathbb{R} \to 0 < x \to x < y \to y^{-1} \in \mathbb{R}$}
\AXC{$\stackrel{\Pi_{36}}{\vdots}$}\noLine
\UIC{$y^{-1} \in \mathbb{R} \to y^{-1} \cdot 1 = y^{-1}$}
\AXC{$(y^{-1} \in \mathbb{R} \to y^{-1} \cdot 1 = y^{-1}) \to \varphi_{131}$\lbl{A1}}
\BIC{$\varphi_{131}$}
\BIC{$x \in \mathbb{R} \to y \in \mathbb{R} \to 0 < x \to x < y \to y^{-1} \cdot 1 = y^{-1}$}
\end{ptree}
\noindent$\Pi_{38}$:
\begin{ptree}
\AXC{$y = z \to x = y \to x = z$\lbl{A6}}
\AXC{$(y = z \to x = y \to x = z) \to x = y \to y = z \to x = z$\lbl{A1}}
\BIC{$x = y \to y = z \to x = z$}
\UIC{$\forall z\,(x = y \to y = z \to x = z)$}
\UIC{$\forall y\,\forall z\,(x = y \to y = z \to x = z)$}
\UIC{$\forall x\,\forall y\,\forall z\,(x = y \to y = z \to x = z)$}
\AXC{$\varphi_{13}$\lbl{A2}}
\BIC{$\forall y\,\forall z\,(v_{536} = y \to y = z \to v_{536} = z)$}
\AXC{$\varphi_{24}$\lbl{A2}}
\BIC{$\forall z\,(v_{536} = v_{537} \to v_{537} = z \to v_{536} = z)$}
\AXC{$\varphi_{37}$\lbl{A2}}
\BIC{$v_{536} = v_{537} \to v_{537} = v_{538} \to v_{536} = v_{538}$}
\UIC{$\forall v_{538}\,(v_{536} = v_{537} \to v_{537} = v_{538} \to v_{536} = v_{538})$}
\UIC{$\varphi_{4}$}
\UIC{$\forall v_{536}\,(\varphi_{4})$}
\end{ptree}
\noindent$\Pi_{39}$:
\begin{ptree}
\AXC{$\stackrel{\Pi_{38}}{\vdots}$}\noLine
\UIC{$\forall v_{536}\,(\varphi_{4})$}
\AXC{$\forall v_{536}\,(\varphi_{4}) \to \varphi_{77}$\lbl{A2}}
\BIC{$\varphi_{77}$}
\AXC{$\varphi_{77} \to \varphi_{96}$\lbl{A2}}
\BIC{$\varphi_{96}$}
\AXC{$\varphi_{96} \to \varphi_{103}$\lbl{A2}}
\BIC{$\varphi_{103}$}
\end{ptree}
\noindent$\Pi_{40}$:
\begin{ptree}
\AXC{$1 \in \mathbb{R}$\lbl{(R1)}}
\AXC{$\forall x\,\forall y\,(x \in \mathbb{R} \to y \in \mathbb{R} \to x \cdot y = y \cdot x)$\lbl{(R7)}}
\AXC{$\varphi_{102}$\lbl{A2}}
\BIC{$\forall y\,(1 \in \mathbb{R} \to y \in \mathbb{R} \to 1 \cdot y = y \cdot 1)$}
\AXC{$\varphi_{108}$\lbl{A2}}
\BIC{$1 \in \mathbb{R} \to y^{-1} \in \mathbb{R} \to 1 \cdot y^{-1} = y^{-1} \cdot 1$}
\AXC{$\varphi_{141}$\lbl{A1}}
\BIC{$1 \in \mathbb{R} \to \varphi_{134}$}
\BIC{$\varphi_{134}$}
\end{ptree}
\noindent$\Pi_{41}$:
\begin{ptree}
\AXC{$y = z \to x = y \to x = z$\lbl{A6}}
\AXC{$(y = z \to x = y \to x = z) \to x = y \to y = z \to x = z$\lbl{A1}}
\BIC{$x = y \to y = z \to x = z$}
\UIC{$\forall z\,(x = y \to y = z \to x = z)$}
\UIC{$\forall y\,\forall z\,(x = y \to y = z \to x = z)$}
\UIC{$\forall x\,\forall y\,\forall z\,(x = y \to y = z \to x = z)$}
\AXC{$\varphi_{14}$\lbl{A2}}
\BIC{$\forall y\,\forall z\,(v_{533} = y \to y = z \to v_{533} = z)$}
\AXC{$\varphi_{25}$\lbl{A2}}
\BIC{$\forall z\,(v_{533} = v_{534} \to v_{534} = z \to v_{533} = z)$}
\AXC{$\varphi_{38}$\lbl{A2}}
\BIC{$v_{533} = v_{534} \to v_{534} = v_{535} \to v_{533} = v_{535}$}
\UIC{$\forall v_{535}\,(v_{533} = v_{534} \to v_{534} = v_{535} \to v_{533} = v_{535})$}
\UIC{$\varphi_{5}$}
\UIC{$\forall v_{533}\,(\varphi_{5})$}
\end{ptree}
\noindent$\Pi_{42}$:
\begin{ptree}
\AXC{$\stackrel{\Pi_{41}}{\vdots}$}\noLine
\UIC{$\forall v_{533}\,(\varphi_{5})$}
\AXC{$\forall v_{533}\,(\varphi_{5}) \to \varphi_{78}$\lbl{A2}}
\BIC{$\varphi_{78}$}
\AXC{$\varphi_{78} \to \varphi_{117}$\lbl{A2}}
\BIC{$\varphi_{117}$}
\AXC{$\varphi_{117} \to x^{-1} \cdot y^{-1} \cdot x = x \cdot x^{-1} \cdot y^{-1} \to \varphi_{46}$\lbl{A2}}
\BIC{$x^{-1} \cdot y^{-1} \cdot x = x \cdot x^{-1} \cdot y^{-1} \to \varphi_{46}$}
\end{ptree}
\noindent$\Pi_{43}$:
\begin{ptree}
\AXC{$z = w \to \varphi_{47}$\lbl{A6}}
\UIC{$\forall w\,(z = w \to \varphi_{47})$}
\UIC{$\forall z\,\forall w\,(z = w \to \varphi_{47})$}
\AXC{$\forall z\,\forall w\,(z = w \to \varphi_{47}) \to \varphi_{105}$\lbl{A2}}
\BIC{$\varphi_{105}$}
\AXC{$\varphi_{105} \to x \cdot x^{-1} = 1 \to \varphi_{84}$\lbl{A2}}
\BIC{$x \cdot x^{-1} = 1 \to \varphi_{84}$}
\end{ptree}
\noindent$\Pi_{44}$:
\begin{ptree}
\AXC{$x \in \mathbb{R} \to y \in \mathbb{R} \to 0 < x \to x < y \to x \in \mathbb{R}$\lbl{A1}}
\AXC{$\forall x \in \mathbb{R}\,(x \neq 0 \to x \cdot x^{-1} = 1)$\lbl{(R9)}}
\AXC{$\varphi_{53}$\lbl{A2}}
\BIC{$x \in \mathbb{R} \to x \neq 0 \to x \cdot x^{-1} = 1$}
\AXC{$\varphi_{144}$\lbl{A1}}
\BIC{$(x \in \mathbb{R} \to y \in \mathbb{R} \to 0 < x \to x < y \to x \in \mathbb{R}) \to \varphi_{123}$}
\BIC{$\varphi_{123}$}
\end{ptree}
\noindent$\Pi_{45}$:
\begin{ptree}
\AXC{$\stackrel{\Pi_{19}}{\vdots}$}\noLine
\UIC{$x \in \mathbb{R} \to y \in \mathbb{R} \to 0 < x \to x < y \to x \neq 0$}
\AXC{$\stackrel{\Pi_{44}}{\vdots}$}\noLine
\UIC{$\varphi_{123}$}
\BIC{$x \in \mathbb{R} \to y \in \mathbb{R} \to 0 < x \to x < y \to x \cdot x^{-1} = 1$}
\AXC{$(x \in \mathbb{R} \to y \in \mathbb{R} \to 0 < x \to x < y \to x \cdot x^{-1} = 1) \to \varphi_{150}$\lbl{A1}}
\BIC{$\varphi_{150}$}
\end{ptree}
\noindent$\Pi_{46}$:
\begin{ptree}
\AXC{$\stackrel{\Pi_{43}}{\vdots}$}\noLine
\UIC{$x \cdot x^{-1} = 1 \to \varphi_{84}$}
\AXC{$x = x$\lbl{A5}}
\UIC{$\forall x\,(x = x)$}
\AXC{$\forall x\,(x = x) \to x \cdot x^{-1} \cdot y^{-1} = x \cdot x^{-1} \cdot y^{-1}$\lbl{A2}}
\BIC{$x \cdot x^{-1} \cdot y^{-1} = x \cdot x^{-1} \cdot y^{-1}$}
\AXC{$\stackrel{\Pi_{45}}{\vdots}$}\noLine
\UIC{$\varphi_{150}$}
\BIC{$(x \cdot x^{-1} = 1 \to \varphi_{84}) \to \varphi_{70}$}
\BIC{$\varphi_{70}$}
\end{ptree}
\noindent$\Pi_{47}$:
\begin{ptree}
\AXC{$y = z \to x = y \to x = z$\lbl{A6}}
\AXC{$(y = z \to x = y \to x = z) \to x = y \to y = z \to x = z$\lbl{A1}}
\BIC{$x = y \to y = z \to x = z$}
\UIC{$\forall z\,(x = y \to y = z \to x = z)$}
\UIC{$\forall y\,\forall z\,(x = y \to y = z \to x = z)$}
\UIC{$\forall x\,\forall y\,\forall z\,(x = y \to y = z \to x = z)$}
\AXC{$\varphi_{15}$\lbl{A2}}
\BIC{$\forall y\,\forall z\,(v_{530} = y \to y = z \to v_{530} = z)$}
\AXC{$\varphi_{26}$\lbl{A2}}
\BIC{$\forall z\,(v_{530} = v_{531} \to v_{531} = z \to v_{530} = z)$}
\AXC{$\varphi_{39}$\lbl{A2}}
\BIC{$v_{530} = v_{531} \to v_{531} = v_{532} \to v_{530} = v_{532}$}
\UIC{$\forall v_{532}\,(v_{530} = v_{531} \to v_{531} = v_{532} \to v_{530} = v_{532})$}
\UIC{$\varphi_{6}$}
\UIC{$\forall v_{530}\,(\varphi_{6})$}
\end{ptree}
\noindent$\Pi_{48}$:
\begin{ptree}
\AXC{$\stackrel{\Pi_{47}}{\vdots}$}\noLine
\UIC{$\forall v_{530}\,(\varphi_{6})$}
\AXC{$\forall v_{530}\,(\varphi_{6}) \to \varphi_{79}$\lbl{A2}}
\BIC{$\varphi_{79}$}
\AXC{$\varphi_{79} \to \varphi_{118}$\lbl{A2}}
\BIC{$\varphi_{118}$}
\AXC{$\varphi_{118} \to x^{-1} \cdot y^{-1} \cdot x = x \cdot (x^{-1} \cdot y^{-1}) \to \varphi_{99}$\lbl{A2}}
\BIC{$x^{-1} \cdot y^{-1} \cdot x = x \cdot (x^{-1} \cdot y^{-1}) \to \varphi_{99}$}
\end{ptree}
\noindent$\Pi_{49}$:
\begin{ptree}
\AXC{$\stackrel{\Pi_{8}}{\vdots}$}\noLine
\UIC{$\forall x\,\forall y\,(x = y \to y = x)$}
\AXC{$\forall x\,\forall y\,(x = y \to y = x) \to \forall y\,(v_{528} = y \to y = v_{528})$\lbl{A2}}
\BIC{$\forall y\,(v_{528} = y \to y = v_{528})$}
\AXC{$\varphi_{2}$\lbl{A2}}
\BIC{$v_{528} = v_{529} \to v_{529} = v_{528}$}
\UIC{$\forall v_{529}\,(v_{528} = v_{529} \to v_{529} = v_{528})$}
\UIC{$\forall v_{528}\,\forall v_{529}\,(v_{528} = v_{529} \to v_{529} = v_{528})$}
\end{ptree}
\noindent$\Pi_{50}$:
\begin{ptree}
\AXC{$\stackrel{\Pi_{49}}{\vdots}$}\noLine
\UIC{$\forall v_{528}\,\forall v_{529}\,(v_{528} = v_{529} \to v_{529} = v_{528})$}
\AXC{$\forall v_{528}\,\forall v_{529}\,(v_{528} = v_{529} \to v_{529} = v_{528}) \to \varphi_{30}$\lbl{A2}}
\BIC{$\varphi_{30}$}
\AXC{$\varphi_{30} \to \varphi_{98}$\lbl{A2}}
\BIC{$\varphi_{98}$}
\end{ptree}
\noindent$\Pi_{51}$:
\begin{ptree}
\AXC{$\forall x\,(\varphi_{44})$\lbl{(R6)}}
\AXC{$\forall x\,(\varphi_{44}) \to \varphi_{44}$\lbl{A2}}
\BIC{$\varphi_{44}$}
\AXC{$\varphi_{44} \to \varphi_{87}$\lbl{A2}}
\BIC{$\varphi_{87}$}
\AXC{$\varphi_{87} \to x \in \mathbb{R} \to \varphi_{85}$\lbl{A2}}
\BIC{$x \in \mathbb{R} \to \varphi_{85}$}
\AXC{$(x \in \mathbb{R} \to \varphi_{85}) \to \varphi_{161}$\lbl{A1}}
\BIC{$\varphi_{161}$}
\end{ptree}
\noindent$\Pi_{52}$:
\begin{ptree}
\AXC{$\stackrel{\Pi_{21}}{\vdots}$}\noLine
\UIC{$x \in \mathbb{R} \to y \in \mathbb{R} \to 0 < x \to x < y \to x^{-1} \in \mathbb{R}$}
\AXC{$x \in \mathbb{R} \to y \in \mathbb{R} \to 0 < x \to x < y \to x \in \mathbb{R}$\lbl{A1}}
\AXC{$\stackrel{\Pi_{51}}{\vdots}$}\noLine
\UIC{$\varphi_{161}$}
\BIC{$(x \in \mathbb{R} \to y \in \mathbb{R} \to 0 < x \to x < y \to x^{-1} \in \mathbb{R}) \to \varphi_{139}$}
\BIC{$\varphi_{139}$}
\end{ptree}
\noindent$\Pi_{53}$:
\begin{ptree}
\AXC{$\stackrel{\Pi_{50}}{\vdots}$}\noLine
\UIC{$\varphi_{98}$}
\AXC{$\stackrel{\Pi_{15}}{\vdots}$}\noLine
\UIC{$x \in \mathbb{R} \to y \in \mathbb{R} \to 0 < x \to x < y \to y^{-1} \in \mathbb{R}$}
\AXC{$\stackrel{\Pi_{52}}{\vdots}$}\noLine
\UIC{$\varphi_{139}$}
\BIC{$x \in \mathbb{R} \to \varphi_{59}$}
\AXC{$(x \in \mathbb{R} \to \varphi_{59}) \to \varphi_{98} \to x \in \mathbb{R} \to \varphi_{60}$\lbl{A1}}
\BIC{$\varphi_{98} \to x \in \mathbb{R} \to \varphi_{60}$}
\BIC{$x \in \mathbb{R} \to \varphi_{60}$}
\end{ptree}
\noindent$\Pi_{54}$:
\begin{ptree}
\AXC{$\forall x\,\forall y\,(x \in \mathbb{R} \to y \in \mathbb{R} \to x \cdot y = y \cdot x)$\lbl{(R7)}}
\AXC{$\varphi_{109}$\lbl{A2}}
\BIC{$\forall y\,(v_{526} \in \mathbb{R} \to y \in \mathbb{R} \to v_{526} \cdot y = y \cdot v_{526})$}
\AXC{$\varphi_{127}$\lbl{A2}}
\BIC{$v_{526} \in \mathbb{R} \to v_{527} \in \mathbb{R} \to v_{526} \cdot v_{527} = v_{527} \cdot v_{526}$}
\UIC{$\varphi_{32}$}
\UIC{$\forall v_{526}\,(\varphi_{32})$}
\AXC{$\forall v_{526}\,(\varphi_{32}) \to \varphi_{110}$\lbl{A2}}
\BIC{$\varphi_{110}$}
\AXC{$\varphi_{110} \to \varphi_{91}$\lbl{A2}}
\BIC{$\varphi_{91}$}
\end{ptree}
\noindent$\Pi_{55}$:
\begin{ptree}
\AXC{$x \in \mathbb{R} \to y \in \mathbb{R} \to 0 < x \to x < y \to x \in \mathbb{R}$\lbl{A1}}
\AXC{$\stackrel{\Pi_{28}}{\vdots}$}\noLine
\UIC{$x \in \mathbb{R} \to y \in \mathbb{R} \to 0 < x \to x < y \to x^{-1} \cdot y^{-1} \in \mathbb{R}$}
\AXC{$\stackrel{\Pi_{54}}{\vdots}$}\noLine
\UIC{$\varphi_{91}$}
\AXC{$\varphi_{91} \to \varphi_{149}$\lbl{A1}}
\BIC{$\varphi_{149}$}
\BIC{$\varphi_{137}$}
\BIC{$x \in \mathbb{R} \to \varphi_{61}$}
\end{ptree}
\noindent$\Pi_{56}$:
\begin{ptree}
\AXC{$\stackrel{\Pi_{48}}{\vdots}$}\noLine
\UIC{$x^{-1} \cdot y^{-1} \cdot x = x \cdot (x^{-1} \cdot y^{-1}) \to \varphi_{99}$}
\AXC{$\stackrel{\Pi_{53}}{\vdots}$}\noLine
\UIC{$x \in \mathbb{R} \to \varphi_{60}$}
\AXC{$\stackrel{\Pi_{55}}{\vdots}$}\noLine
\UIC{$x \in \mathbb{R} \to \varphi_{61}$}
\AXC{$(x \in \mathbb{R} \to \varphi_{61}) \to (x \in \mathbb{R} \to \varphi_{60}) \to \varphi_{152}$\lbl{A1}}
\BIC{$(x \in \mathbb{R} \to \varphi_{60}) \to \varphi_{152}$}
\BIC{$\varphi_{152}$}
\BIC{$x \in \mathbb{R} \to \varphi_{62}$}
\end{ptree}
\noindent$\Pi_{57}$:
\begin{ptree}
\AXC{$\stackrel{\Pi_{42}}{\vdots}$}\noLine
\UIC{$x^{-1} \cdot y^{-1} \cdot x = x \cdot x^{-1} \cdot y^{-1} \to \varphi_{46}$}
\AXC{$\stackrel{\Pi_{46}}{\vdots}$}\noLine
\UIC{$\varphi_{70}$}
\AXC{$\stackrel{\Pi_{56}}{\vdots}$}\noLine
\UIC{$x \in \mathbb{R} \to \varphi_{62}$}
\AXC{$\varphi_{166}$\lbl{A1}}
\BIC{$\varphi_{70} \to (x^{-1} \cdot y^{-1} \cdot x = x \cdot x^{-1} \cdot y^{-1} \to \varphi_{46}) \to \varphi_{71}$}
\BIC{$(x^{-1} \cdot y^{-1} \cdot x = x \cdot x^{-1} \cdot y^{-1} \to \varphi_{46}) \to \varphi_{71}$}
\BIC{$\varphi_{71}$}
\AXC{$\varphi_{71} \to \varphi_{158}$\lbl{A1}}
\BIC{$\varphi_{158}$}
\end{ptree}
\noindent$\Pi_{58}$:
\begin{ptree}
\AXC{$\stackrel{\Pi_{39}}{\vdots}$}\noLine
\UIC{$\varphi_{103}$}
\AXC{$\stackrel{\Pi_{15}}{\vdots}$}\noLine
\UIC{$x \in \mathbb{R} \to y \in \mathbb{R} \to 0 < x \to x < y \to y^{-1} \in \mathbb{R}$}
\AXC{$\stackrel{\Pi_{40}}{\vdots}$}\noLine
\UIC{$\varphi_{134}$}
\BIC{$x \in \mathbb{R} \to y \in \mathbb{R} \to 0 < x \to x < y \to 1 \cdot y^{-1} = y^{-1} \cdot 1$}
\AXC{$\stackrel{\Pi_{57}}{\vdots}$}\noLine
\UIC{$\varphi_{158}$}
\BIC{$\varphi_{103} \to \varphi_{72}$}
\BIC{$\varphi_{72}$}
\AXC{$\varphi_{72} \to \varphi_{153}$\lbl{A1}}
\BIC{$\varphi_{153}$}
\end{ptree}
\noindent$\Pi_{59}$:
\begin{ptree}
\AXC{$\stackrel{\Pi_{35}}{\vdots}$}\noLine
\UIC{$\varphi_{93}$}
\AXC{$\stackrel{\Pi_{37}}{\vdots}$}\noLine
\UIC{$x \in \mathbb{R} \to y \in \mathbb{R} \to 0 < x \to x < y \to y^{-1} \cdot 1 = y^{-1}$}
\AXC{$\stackrel{\Pi_{58}}{\vdots}$}\noLine
\UIC{$\varphi_{153}$}
\BIC{$\varphi_{93} \to \varphi_{54}$}
\BIC{$\varphi_{54}$}
\end{ptree}
\noindent$\Pi_{60}$:
\begin{ptree}
\AXC{$\stackrel{\Pi_{2}}{\vdots}$}\noLine
\UIC{$x^{-1} \cdot y^{-1} \cdot x = y^{-1} \to \varphi_{88}$}
\AXC{$\stackrel{\Pi_{33}}{\vdots}$}\noLine
\UIC{$x \in \mathbb{R} \to \varphi_{75}$}
\AXC{$\stackrel{\Pi_{59}}{\vdots}$}\noLine
\UIC{$\varphi_{54}$}
\AXC{$\varphi_{54} \to (x \in \mathbb{R} \to \varphi_{75}) \to (x^{-1} \cdot y^{-1} \cdot x = y^{-1} \to \varphi_{88}) \to \varphi_{66}$\lbl{A1}}
\BIC{$(x \in \mathbb{R} \to \varphi_{75}) \to (x^{-1} \cdot y^{-1} \cdot x = y^{-1} \to \varphi_{88}) \to \varphi_{66}$}
\BIC{$(x^{-1} \cdot y^{-1} \cdot x = y^{-1} \to \varphi_{88}) \to \varphi_{66}$}
\BIC{$\varphi_{66}$}
\end{ptree}
\noindent$\Pi_{61}$:
\begin{ptree}
\AXC{$y = z \to x = y \to x = z$\lbl{A6}}
\AXC{$(y = z \to x = y \to x = z) \to x = y \to y = z \to x = z$\lbl{A1}}
\BIC{$x = y \to y = z \to x = z$}
\UIC{$\forall z\,(x = y \to y = z \to x = z)$}
\UIC{$\forall y\,\forall z\,(x = y \to y = z \to x = z)$}
\UIC{$\forall x\,\forall y\,\forall z\,(x = y \to y = z \to x = z)$}
\AXC{$\varphi_{16}$\lbl{A2}}
\BIC{$\forall y\,\forall z\,(v_{550} = y \to y = z \to v_{550} = z)$}
\AXC{$\varphi_{27}$\lbl{A2}}
\BIC{$\forall z\,(v_{550} = v_{551} \to v_{551} = z \to v_{550} = z)$}
\AXC{$\varphi_{40}$\lbl{A2}}
\BIC{$v_{550} = v_{551} \to v_{551} = v_{552} \to v_{550} = v_{552}$}
\UIC{$\forall v_{552}\,(v_{550} = v_{551} \to v_{551} = v_{552} \to v_{550} = v_{552})$}
\UIC{$\varphi_{7}$}
\UIC{$\forall v_{550}\,(\varphi_{7})$}
\end{ptree}
\noindent$\Pi_{62}$:
\begin{ptree}
\AXC{$\stackrel{\Pi_{61}}{\vdots}$}\noLine
\UIC{$\forall v_{550}\,(\varphi_{7})$}
\AXC{$\forall v_{550}\,(\varphi_{7}) \to \varphi_{80}$\lbl{A2}}
\BIC{$\varphi_{80}$}
\AXC{$\varphi_{80} \to \varphi_{97}$\lbl{A2}}
\BIC{$\varphi_{97}$}
\AXC{$\varphi_{97} \to \varphi_{94}$\lbl{A2}}
\BIC{$\varphi_{94}$}
\end{ptree}
\noindent$\Pi_{63}$:
\begin{ptree}
\AXC{$\forall x \in \mathbb{R}\,(x \cdot 1 = x)$\lbl{(R8)}}
\AXC{$\forall x \in \mathbb{R}\,(x \cdot 1 = x) \to x^{-1} \in \mathbb{R} \to x^{-1} \cdot 1 = x^{-1}$\lbl{A2}}
\BIC{$x^{-1} \in \mathbb{R} \to x^{-1} \cdot 1 = x^{-1}$}
\end{ptree}
\noindent$\Pi_{64}$:
\begin{ptree}
\AXC{$\stackrel{\Pi_{21}}{\vdots}$}\noLine
\UIC{$x \in \mathbb{R} \to y \in \mathbb{R} \to 0 < x \to x < y \to x^{-1} \in \mathbb{R}$}
\AXC{$\stackrel{\Pi_{63}}{\vdots}$}\noLine
\UIC{$x^{-1} \in \mathbb{R} \to x^{-1} \cdot 1 = x^{-1}$}
\AXC{$(x^{-1} \in \mathbb{R} \to x^{-1} \cdot 1 = x^{-1}) \to \varphi_{132}$\lbl{A1}}
\BIC{$\varphi_{132}$}
\BIC{$x \in \mathbb{R} \to y \in \mathbb{R} \to 0 < x \to x < y \to x^{-1} \cdot 1 = x^{-1}$}
\end{ptree}
\noindent$\Pi_{65}$:
\begin{ptree}
\AXC{$y = z \to x = y \to x = z$\lbl{A6}}
\AXC{$(y = z \to x = y \to x = z) \to x = y \to y = z \to x = z$\lbl{A1}}
\BIC{$x = y \to y = z \to x = z$}
\UIC{$\forall z\,(x = y \to y = z \to x = z)$}
\UIC{$\forall y\,\forall z\,(x = y \to y = z \to x = z)$}
\UIC{$\forall x\,\forall y\,\forall z\,(x = y \to y = z \to x = z)$}
\AXC{$\varphi_{17}$\lbl{A2}}
\BIC{$\forall y\,\forall z\,(v_{547} = y \to y = z \to v_{547} = z)$}
\AXC{$\varphi_{28}$\lbl{A2}}
\BIC{$\forall z\,(v_{547} = v_{548} \to v_{548} = z \to v_{547} = z)$}
\AXC{$\varphi_{41}$\lbl{A2}}
\BIC{$v_{547} = v_{548} \to v_{548} = v_{549} \to v_{547} = v_{549}$}
\UIC{$\forall v_{549}\,(v_{547} = v_{548} \to v_{548} = v_{549} \to v_{547} = v_{549})$}
\UIC{$\varphi_{8}$}
\UIC{$\forall v_{547}\,(\varphi_{8})$}
\end{ptree}
\noindent$\Pi_{66}$:
\begin{ptree}
\AXC{$\stackrel{\Pi_{65}}{\vdots}$}\noLine
\UIC{$\forall v_{547}\,(\varphi_{8})$}
\AXC{$\forall v_{547}\,(\varphi_{8}) \to \varphi_{81}$\lbl{A2}}
\BIC{$\varphi_{81}$}
\AXC{$\varphi_{81} \to \varphi_{119}$\lbl{A2}}
\BIC{$\varphi_{119}$}
\AXC{$\varphi_{119} \to x^{-1} \cdot y^{-1} \cdot y = x^{-1} \cdot (y \cdot y^{-1}) \to \varphi_{48}$\lbl{A2}}
\BIC{$x^{-1} \cdot y^{-1} \cdot y = x^{-1} \cdot (y \cdot y^{-1}) \to \varphi_{48}$}
\end{ptree}
\noindent$\Pi_{67}$:
\begin{ptree}
\AXC{$z = w \to \varphi_{49}$\lbl{A6}}
\UIC{$\forall w\,(z = w \to \varphi_{49})$}
\UIC{$\forall z\,\forall w\,(z = w \to \varphi_{49})$}
\AXC{$\forall z\,\forall w\,(z = w \to \varphi_{49}) \to \varphi_{106}$\lbl{A2}}
\BIC{$\varphi_{106}$}
\AXC{$\varphi_{106} \to y \cdot y^{-1} = 1 \to \varphi_{86}$\lbl{A2}}
\BIC{$y \cdot y^{-1} = 1 \to \varphi_{86}$}
\end{ptree}
\noindent$\Pi_{68}$:
\begin{ptree}
\AXC{$x \in \mathbb{R} \to y \in \mathbb{R} \to 0 < x \to x < y \to y \in \mathbb{R}$\lbl{A1}}
\AXC{$\forall x \in \mathbb{R}\,(x \neq 0 \to x \cdot x^{-1} = 1)$\lbl{(R9)}}
\AXC{$\varphi_{55}$\lbl{A2}}
\BIC{$y \in \mathbb{R} \to y \neq 0 \to y \cdot y^{-1} = 1$}
\AXC{$\varphi_{145}$\lbl{A1}}
\BIC{$(x \in \mathbb{R} \to y \in \mathbb{R} \to 0 < x \to x < y \to y \in \mathbb{R}) \to \varphi_{124}$}
\BIC{$\varphi_{124}$}
\end{ptree}
\noindent$\Pi_{69}$:
\begin{ptree}
\AXC{$\stackrel{\Pi_{13}}{\vdots}$}\noLine
\UIC{$x \in \mathbb{R} \to y \in \mathbb{R} \to 0 < x \to x < y \to y \neq 0$}
\AXC{$\stackrel{\Pi_{68}}{\vdots}$}\noLine
\UIC{$\varphi_{124}$}
\BIC{$x \in \mathbb{R} \to y \in \mathbb{R} \to 0 < x \to x < y \to y \cdot y^{-1} = 1$}
\AXC{$(x \in \mathbb{R} \to y \in \mathbb{R} \to 0 < x \to x < y \to y \cdot y^{-1} = 1) \to \varphi_{151}$\lbl{A1}}
\BIC{$\varphi_{151}$}
\end{ptree}
\noindent$\Pi_{70}$:
\begin{ptree}
\AXC{$y = z \to x = y \to x = z$\lbl{A6}}
\AXC{$(y = z \to x = y \to x = z) \to x = y \to y = z \to x = z$\lbl{A1}}
\BIC{$x = y \to y = z \to x = z$}
\UIC{$\forall z\,(x = y \to y = z \to x = z)$}
\UIC{$\forall y\,\forall z\,(x = y \to y = z \to x = z)$}
\UIC{$\forall x\,\forall y\,\forall z\,(x = y \to y = z \to x = z)$}
\AXC{$\varphi_{18}$\lbl{A2}}
\BIC{$\forall y\,\forall z\,(v_{544} = y \to y = z \to v_{544} = z)$}
\AXC{$\varphi_{29}$\lbl{A2}}
\BIC{$\forall z\,(v_{544} = v_{545} \to v_{545} = z \to v_{544} = z)$}
\AXC{$\varphi_{42}$\lbl{A2}}
\BIC{$v_{544} = v_{545} \to v_{545} = v_{546} \to v_{544} = v_{546}$}
\UIC{$\forall v_{546}\,(v_{544} = v_{545} \to v_{545} = v_{546} \to v_{544} = v_{546})$}
\UIC{$\varphi_{9}$}
\UIC{$\forall v_{544}\,(\varphi_{9})$}
\end{ptree}
\noindent$\Pi_{71}$:
\begin{ptree}
\AXC{$z = w \to \varphi_{50}$\lbl{A6}}
\UIC{$\forall w\,(z = w \to \varphi_{50})$}
\UIC{$\forall z\,\forall w\,(z = w \to \varphi_{50})$}
\AXC{$\forall z\,\forall w\,(z = w \to \varphi_{50}) \to \varphi_{107}$\lbl{A2}}
\BIC{$\varphi_{107}$}
\AXC{$\varphi_{107} \to y^{-1} \cdot y = y \cdot y^{-1} \to \varphi_{101}$\lbl{A2}}
\BIC{$y^{-1} \cdot y = y \cdot y^{-1} \to \varphi_{101}$}
\end{ptree}
\noindent$\Pi_{72}$:
\begin{ptree}
\AXC{$\forall x\,\forall y\,(x \in \mathbb{R} \to y \in \mathbb{R} \to x \cdot y = y \cdot x)$\lbl{(R7)}}
\AXC{$\varphi_{111}$\lbl{A2}}
\BIC{$\forall y\,(v_{542} \in \mathbb{R} \to y \in \mathbb{R} \to v_{542} \cdot y = y \cdot v_{542})$}
\AXC{$\varphi_{128}$\lbl{A2}}
\BIC{$v_{542} \in \mathbb{R} \to v_{543} \in \mathbb{R} \to v_{542} \cdot v_{543} = v_{543} \cdot v_{542}$}
\UIC{$\varphi_{33}$}
\UIC{$\forall v_{542}\,(\varphi_{33})$}
\AXC{$\forall v_{542}\,(\varphi_{33}) \to \varphi_{45}$\lbl{A2}}
\BIC{$\varphi_{45}$}
\end{ptree}
\noindent$\Pi_{73}$:
\begin{ptree}
\AXC{$\stackrel{\Pi_{15}}{\vdots}$}\noLine
\UIC{$x \in \mathbb{R} \to y \in \mathbb{R} \to 0 < x \to x < y \to y^{-1} \in \mathbb{R}$}
\AXC{$\stackrel{\Pi_{72}}{\vdots}$}\noLine
\UIC{$\varphi_{45}$}
\AXC{$\varphi_{45} \to y^{-1} \in \mathbb{R} \to y \in \mathbb{R} \to y^{-1} \cdot y = y \cdot y^{-1}$\lbl{A2}}
\BIC{$y^{-1} \in \mathbb{R} \to y \in \mathbb{R} \to y^{-1} \cdot y = y \cdot y^{-1}$}
\AXC{$\varphi_{157}$\lbl{A1}}
\BIC{$(x \in \mathbb{R} \to y \in \mathbb{R} \to 0 < x \to x < y \to y^{-1} \in \mathbb{R}) \to \varphi_{133}$}
\BIC{$\varphi_{133}$}
\end{ptree}
\noindent$\Pi_{74}$:
\begin{ptree}
\AXC{$x = x$\lbl{A5}}
\UIC{$\forall x\,(x = x)$}
\AXC{$\varphi_{11}$\lbl{A2}}
\BIC{$x^{-1} \cdot (y^{-1} \cdot y) = x^{-1} \cdot (y^{-1} \cdot y)$}
\AXC{$x \in \mathbb{R} \to y \in \mathbb{R} \to 0 < x \to x < y \to y \in \mathbb{R}$\lbl{A1}}
\AXC{$\stackrel{\Pi_{73}}{\vdots}$}\noLine
\UIC{$\varphi_{133}$}
\BIC{$x \in \mathbb{R} \to y \in \mathbb{R} \to 0 < x \to x < y \to y^{-1} \cdot y = y \cdot y^{-1}$}
\AXC{$\varphi_{165}$\lbl{A1}}
\BIC{$\varphi_{159}$}
\BIC{$(y^{-1} \cdot y = y \cdot y^{-1} \to \varphi_{101}) \to x \in \mathbb{R} \to y \in \mathbb{R} \to \varphi_{31}$}
\end{ptree}
\noindent$\Pi_{75}$:
\begin{ptree}
\AXC{$\forall x\,(\varphi_{44})$\lbl{(R6)}}
\AXC{$\forall x\,(\varphi_{44}) \to \varphi_{90}$\lbl{A2}}
\BIC{$\varphi_{90}$}
\AXC{$\varphi_{90} \to \varphi_{104}$\lbl{A2}}
\BIC{$\varphi_{104}$}
\AXC{$\varphi_{104} \to x^{-1} \in \mathbb{R} \to \varphi_{63}$\lbl{A2}}
\BIC{$x^{-1} \in \mathbb{R} \to \varphi_{63}$}
\AXC{$(x^{-1} \in \mathbb{R} \to \varphi_{63}) \to \varphi_{162}$\lbl{A1}}
\BIC{$\varphi_{162}$}
\end{ptree}
\noindent$\Pi_{76}$:
\begin{ptree}
\AXC{$\stackrel{\Pi_{15}}{\vdots}$}\noLine
\UIC{$x \in \mathbb{R} \to y \in \mathbb{R} \to 0 < x \to x < y \to y^{-1} \in \mathbb{R}$}
\AXC{$\stackrel{\Pi_{21}}{\vdots}$}\noLine
\UIC{$x \in \mathbb{R} \to y \in \mathbb{R} \to 0 < x \to x < y \to x^{-1} \in \mathbb{R}$}
\AXC{$\stackrel{\Pi_{75}}{\vdots}$}\noLine
\UIC{$\varphi_{162}$}
\BIC{$(x \in \mathbb{R} \to y \in \mathbb{R} \to 0 < x \to x < y \to y^{-1} \in \mathbb{R}) \to \varphi_{138}$}
\BIC{$\varphi_{138}$}
\end{ptree}
\noindent$\Pi_{77}$:
\begin{ptree}
\AXC{$x \in \mathbb{R} \to y \in \mathbb{R} \to 0 < x \to x < y \to y \in \mathbb{R}$\lbl{A1}}
\AXC{$\stackrel{\Pi_{76}}{\vdots}$}\noLine
\UIC{$\varphi_{138}$}
\BIC{$x \in \mathbb{R} \to \varphi_{64}$}
\AXC{$(x \in \mathbb{R} \to \varphi_{64}) \to (x \in \mathbb{R} \to y \in \mathbb{R} \to \varphi_{31}) \to \varphi_{154}$\lbl{A1}}
\BIC{$(x \in \mathbb{R} \to y \in \mathbb{R} \to \varphi_{31}) \to \varphi_{154}$}
\end{ptree}
\noindent$\Pi_{78}$:
\begin{ptree}
\AXC{$\stackrel{\Pi_{71}}{\vdots}$}\noLine
\UIC{$y^{-1} \cdot y = y \cdot y^{-1} \to \varphi_{101}$}
\AXC{$\stackrel{\Pi_{74}}{\vdots}$}\noLine
\UIC{$(y^{-1} \cdot y = y \cdot y^{-1} \to \varphi_{101}) \to x \in \mathbb{R} \to y \in \mathbb{R} \to \varphi_{31}$}
\BIC{$x \in \mathbb{R} \to y \in \mathbb{R} \to \varphi_{31}$}
\AXC{$\stackrel{\Pi_{77}}{\vdots}$}\noLine
\UIC{$(x \in \mathbb{R} \to y \in \mathbb{R} \to \varphi_{31}) \to \varphi_{154}$}
\BIC{$\varphi_{154}$}
\end{ptree}
\noindent$\Pi_{79}$:
\begin{ptree}
\AXC{$\stackrel{\Pi_{70}}{\vdots}$}\noLine
\UIC{$\forall v_{544}\,(\varphi_{9})$}
\AXC{$\forall v_{544}\,(\varphi_{9}) \to \varphi_{82}$\lbl{A2}}
\BIC{$\varphi_{82}$}
\AXC{$\varphi_{82} \to \varphi_{120}$\lbl{A2}}
\BIC{$\varphi_{120}$}
\AXC{$\varphi_{120} \to x^{-1} \cdot y^{-1} \cdot y = x^{-1} \cdot (y^{-1} \cdot y) \to \varphi_{100}$\lbl{A2}}
\BIC{$x^{-1} \cdot y^{-1} \cdot y = x^{-1} \cdot (y^{-1} \cdot y) \to \varphi_{100}$}
\AXC{$\stackrel{\Pi_{78}}{\vdots}$}\noLine
\UIC{$\varphi_{154}$}
\BIC{$x \in \mathbb{R} \to \varphi_{65}$}
\end{ptree}
\noindent$\Pi_{80}$:
\begin{ptree}
\AXC{$\stackrel{\Pi_{67}}{\vdots}$}\noLine
\UIC{$y \cdot y^{-1} = 1 \to \varphi_{86}$}
\AXC{$x = x$\lbl{A5}}
\UIC{$\forall x\,(x = x)$}
\AXC{$\varphi_{10}$\lbl{A2}}
\BIC{$x^{-1} \cdot (y \cdot y^{-1}) = x^{-1} \cdot (y \cdot y^{-1})$}
\AXC{$\stackrel{\Pi_{69}}{\vdots}$}\noLine
\UIC{$\varphi_{151}$}
\BIC{$(y \cdot y^{-1} = 1 \to \varphi_{86}) \to \varphi_{73}$}
\BIC{$\varphi_{73}$}
\AXC{$\stackrel{\Pi_{79}}{\vdots}$}\noLine
\UIC{$x \in \mathbb{R} \to \varphi_{65}$}
\AXC{$\varphi_{167}$\lbl{A1}}
\BIC{$\varphi_{73} \to (x^{-1} \cdot y^{-1} \cdot y = x^{-1} \cdot (y \cdot y^{-1}) \to \varphi_{48}) \to \varphi_{74}$}
\BIC{$(x^{-1} \cdot y^{-1} \cdot y = x^{-1} \cdot (y \cdot y^{-1}) \to \varphi_{48}) \to \varphi_{74}$}
\end{ptree}
\noindent$\Pi_{81}$:
\begin{ptree}
\AXC{$\stackrel{\Pi_{66}}{\vdots}$}\noLine
\UIC{$x^{-1} \cdot y^{-1} \cdot y = x^{-1} \cdot (y \cdot y^{-1}) \to \varphi_{48}$}
\AXC{$\stackrel{\Pi_{80}}{\vdots}$}\noLine
\UIC{$(x^{-1} \cdot y^{-1} \cdot y = x^{-1} \cdot (y \cdot y^{-1}) \to \varphi_{48}) \to \varphi_{74}$}
\BIC{$\varphi_{74}$}
\AXC{$\varphi_{74} \to \varphi_{155}$\lbl{A1}}
\BIC{$\varphi_{155}$}
\end{ptree}
\noindent$\Pi_{82}$:
\begin{ptree}
\AXC{$\stackrel{\Pi_{60}}{\vdots}$}\noLine
\UIC{$\varphi_{66}$}
\AXC{$\stackrel{\Pi_{62}}{\vdots}$}\noLine
\UIC{$\varphi_{94}$}
\AXC{$\stackrel{\Pi_{64}}{\vdots}$}\noLine
\UIC{$x \in \mathbb{R} \to y \in \mathbb{R} \to 0 < x \to x < y \to x^{-1} \cdot 1 = x^{-1}$}
\AXC{$\stackrel{\Pi_{81}}{\vdots}$}\noLine
\UIC{$\varphi_{155}$}
\BIC{$\varphi_{94} \to \varphi_{56}$}
\BIC{$\varphi_{56}$}
\AXC{$\varphi_{164}$\lbl{A1}}
\BIC{$\varphi_{66} \to \varphi_{89} \to x \in \mathbb{R} \to y \in \mathbb{R} \to 0 < x \to x < y \to y^{-1} < x^{-1}$}
\BIC{$\varphi_{89} \to x \in \mathbb{R} \to y \in \mathbb{R} \to 0 < x \to x < y \to y^{-1} < x^{-1}$}
\end{ptree}
\noindent 主木:
\begin{ptree}
\AXC{$\stackrel{\Pi_{1}}{\vdots}$}\noLine
\UIC{$\varphi_{89}$}
\AXC{$\stackrel{\Pi_{82}}{\vdots}$}\noLine
\UIC{$\varphi_{89} \to x \in \mathbb{R} \to y \in \mathbb{R} \to 0 < x \to x < y \to y^{-1} < x^{-1}$}
\BIC{$x \in \mathbb{R} \to y \in \mathbb{R} \to 0 < x \to x < y \to y^{-1} < x^{-1}$}
\UIC{$\forall y\,(x \in \mathbb{R} \to y \in \mathbb{R} \to 0 < x \to x < y \to y^{-1} < x^{-1})$}
\UIC{$\varphi_{43}$}
\end{ptree}
\begin{abbrevlist}
\item $\varphi_{1} :\equiv \forall y\,(v_{524} = y \to y = v_{524}) \to v_{524} = v_{525} \to v_{525} = v_{524}$
\item $\varphi_{2} :\equiv \forall y\,(v_{528} = y \to y = v_{528}) \to v_{528} = v_{529} \to v_{529} = v_{528}$
\item $\varphi_{3} :\equiv \forall v_{540}\,\forall v_{541}\,(v_{539} = v_{540} \to v_{540} = v_{541} \to v_{539} = v_{541})$
\item $\varphi_{4} :\equiv \forall v_{537}\,\forall v_{538}\,(v_{536} = v_{537} \to v_{537} = v_{538} \to v_{536} = v_{538})$
\item $\varphi_{5} :\equiv \forall v_{534}\,\forall v_{535}\,(v_{533} = v_{534} \to v_{534} = v_{535} \to v_{533} = v_{535})$
\item $\varphi_{6} :\equiv \forall v_{531}\,\forall v_{532}\,(v_{530} = v_{531} \to v_{531} = v_{532} \to v_{530} = v_{532})$
\item $\varphi_{7} :\equiv \forall v_{551}\,\forall v_{552}\,(v_{550} = v_{551} \to v_{551} = v_{552} \to v_{550} = v_{552})$
\item $\varphi_{8} :\equiv \forall v_{548}\,\forall v_{549}\,(v_{547} = v_{548} \to v_{548} = v_{549} \to v_{547} = v_{549})$
\item $\varphi_{9} :\equiv \forall v_{545}\,\forall v_{546}\,(v_{544} = v_{545} \to v_{545} = v_{546} \to v_{544} = v_{546})$
\item $\varphi_{10} :\equiv \forall x\,(x = x) \to x^{-1} \cdot (y \cdot y^{-1}) = x^{-1} \cdot (y \cdot y^{-1})$
\item $\varphi_{11} :\equiv \forall x\,(x = x) \to x^{-1} \cdot (y^{-1} \cdot y) = x^{-1} \cdot (y^{-1} \cdot y)$
\item $\varphi_{12} :\equiv \forall x\,\forall y\,\forall z\,(x = y \to y = z \to x = z) \to \forall y\,\forall z\,(v_{539} = y \to y = z \to v_{539} = z)$
\item $\varphi_{13} :\equiv \forall x\,\forall y\,\forall z\,(x = y \to y = z \to x = z) \to \forall y\,\forall z\,(v_{536} = y \to y = z \to v_{536} = z)$
\item $\varphi_{14} :\equiv \forall x\,\forall y\,\forall z\,(x = y \to y = z \to x = z) \to \forall y\,\forall z\,(v_{533} = y \to y = z \to v_{533} = z)$
\item $\varphi_{15} :\equiv \forall x\,\forall y\,\forall z\,(x = y \to y = z \to x = z) \to \forall y\,\forall z\,(v_{530} = y \to y = z \to v_{530} = z)$
\item $\varphi_{16} :\equiv \forall x\,\forall y\,\forall z\,(x = y \to y = z \to x = z) \to \forall y\,\forall z\,(v_{550} = y \to y = z \to v_{550} = z)$
\item $\varphi_{17} :\equiv \forall x\,\forall y\,\forall z\,(x = y \to y = z \to x = z) \to \forall y\,\forall z\,(v_{547} = y \to y = z \to v_{547} = z)$
\item $\varphi_{18} :\equiv \forall x\,\forall y\,\forall z\,(x = y \to y = z \to x = z) \to \forall y\,\forall z\,(v_{544} = y \to y = z \to v_{544} = z)$
\item $\varphi_{19} :\equiv \forall v_{524}\,\forall v_{525}\,(v_{524} = v_{525} \to v_{525} = v_{524}) \to \forall v_{525}\,(y = v_{525} \to v_{525} = y)$
\item $\varphi_{20} :\equiv \forall y\,\forall z\,(0 < y \to y < z \to 0 < z) \to \forall z\,(0 < x \to x < z \to 0 < z)$
\item $\varphi_{21} :\equiv (y = 0 \to 0 = y) \to x \in \mathbb{R} \to y \in \mathbb{R} \to 0 < x \to x < y \to y \neq 0$
\item $\varphi_{22} :\equiv (x = 0 \to 0 = x) \to x \in \mathbb{R} \to y \in \mathbb{R} \to 0 < x \to x < y \to x \neq 0$
\item $\varphi_{23} :\equiv \forall y\,\forall z\,(v_{539} = y \to y = z \to v_{539} = z) \to \forall z\,(v_{539} = v_{540} \to v_{540} = z \to v_{539} = z)$
\item $\varphi_{24} :\equiv \forall y\,\forall z\,(v_{536} = y \to y = z \to v_{536} = z) \to \forall z\,(v_{536} = v_{537} \to v_{537} = z \to v_{536} = z)$
\item $\varphi_{25} :\equiv \forall y\,\forall z\,(v_{533} = y \to y = z \to v_{533} = z) \to \forall z\,(v_{533} = v_{534} \to v_{534} = z \to v_{533} = z)$
\item $\varphi_{26} :\equiv \forall y\,\forall z\,(v_{530} = y \to y = z \to v_{530} = z) \to \forall z\,(v_{530} = v_{531} \to v_{531} = z \to v_{530} = z)$
\item $\varphi_{27} :\equiv \forall y\,\forall z\,(v_{550} = y \to y = z \to v_{550} = z) \to \forall z\,(v_{550} = v_{551} \to v_{551} = z \to v_{550} = z)$
\item $\varphi_{28} :\equiv \forall y\,\forall z\,(v_{547} = y \to y = z \to v_{547} = z) \to \forall z\,(v_{547} = v_{548} \to v_{548} = z \to v_{547} = z)$
\item $\varphi_{29} :\equiv \forall y\,\forall z\,(v_{544} = y \to y = z \to v_{544} = z) \to \forall z\,(v_{544} = v_{545} \to v_{545} = z \to v_{544} = z)$
\item $\varphi_{30} :\equiv \forall v_{529}\,(x \cdot x^{-1} \cdot y^{-1} = v_{529} \to v_{529} = x \cdot x^{-1} \cdot y^{-1})$
\item $\varphi_{31} :\equiv 0 < x \to x < y \to x^{-1} \cdot (y^{-1} \cdot y) = x^{-1} \cdot (y \cdot y^{-1})$
\item $\varphi_{32} :\equiv \forall v_{527}\,(v_{526} \in \mathbb{R} \to v_{527} \in \mathbb{R} \to v_{526} \cdot v_{527} = v_{527} \cdot v_{526})$
\item $\varphi_{33} :\equiv \forall v_{543}\,(v_{542} \in \mathbb{R} \to v_{543} \in \mathbb{R} \to v_{542} \cdot v_{543} = v_{543} \cdot v_{542})$
\item $\varphi_{34} :\equiv \forall w\,(z = w \to z < x^{-1} \cdot y^{-1} \cdot y \to w < x^{-1} \cdot y^{-1} \cdot y)$
\item $\varphi_{35} :\equiv \forall x\,\forall y\,\forall z\,(x < y \to y < z \to x < z) \to \forall y\,\forall z\,(0 < y \to y < z \to 0 < z)$
\item $\varphi_{36} :\equiv \forall z\,(v_{539} = v_{540} \to v_{540} = z \to v_{539} = z) \to v_{539} = v_{540} \to v_{540} = v_{541} \to v_{539} = v_{541}$
\item $\varphi_{37} :\equiv \forall z\,(v_{536} = v_{537} \to v_{537} = z \to v_{536} = z) \to v_{536} = v_{537} \to v_{537} = v_{538} \to v_{536} = v_{538}$
\item $\varphi_{38} :\equiv \forall z\,(v_{533} = v_{534} \to v_{534} = z \to v_{533} = z) \to v_{533} = v_{534} \to v_{534} = v_{535} \to v_{533} = v_{535}$
\item $\varphi_{39} :\equiv \forall z\,(v_{530} = v_{531} \to v_{531} = z \to v_{530} = z) \to v_{530} = v_{531} \to v_{531} = v_{532} \to v_{530} = v_{532}$
\item $\varphi_{40} :\equiv \forall z\,(v_{550} = v_{551} \to v_{551} = z \to v_{550} = z) \to v_{550} = v_{551} \to v_{551} = v_{552} \to v_{550} = v_{552}$
\item $\varphi_{41} :\equiv \forall z\,(v_{547} = v_{548} \to v_{548} = z \to v_{547} = z) \to v_{547} = v_{548} \to v_{548} = v_{549} \to v_{547} = v_{549}$
\item $\varphi_{42} :\equiv \forall z\,(v_{544} = v_{545} \to v_{545} = z \to v_{544} = z) \to v_{544} = v_{545} \to v_{545} = v_{546} \to v_{544} = v_{546}$
\item $\varphi_{43} :\equiv \forall x\,\forall y\,(x \in \mathbb{R} \to y \in \mathbb{R} \to 0 < x \to x < y \to y^{-1} < x^{-1})$
\item $\varphi_{44} :\equiv \forall y\,\forall z\,(x \in \mathbb{R} \to y \in \mathbb{R} \to z \in \mathbb{R} \to x \cdot y \cdot z = x \cdot (y \cdot z))$
\item $\varphi_{45} :\equiv \forall v_{543}\,(y^{-1} \in \mathbb{R} \to v_{543} \in \mathbb{R} \to y^{-1} \cdot v_{543} = v_{543} \cdot y^{-1})$
\item $\varphi_{46} :\equiv x \cdot x^{-1} \cdot y^{-1} = 1 \cdot y^{-1} \to x^{-1} \cdot y^{-1} \cdot x = 1 \cdot y^{-1}$
\item $\varphi_{47} :\equiv x \cdot x^{-1} \cdot y^{-1} = z \cdot y^{-1} \to x \cdot x^{-1} \cdot y^{-1} = w \cdot y^{-1}$
\item $\varphi_{48} :\equiv x^{-1} \cdot (y \cdot y^{-1}) = x^{-1} \cdot 1 \to x^{-1} \cdot y^{-1} \cdot y = x^{-1} \cdot 1$
\item $\varphi_{49} :\equiv x^{-1} \cdot (y \cdot y^{-1}) = x^{-1} \cdot z \to x^{-1} \cdot (y \cdot y^{-1}) = x^{-1} \cdot w$
\item $\varphi_{50} :\equiv x^{-1} \cdot (y^{-1} \cdot y) = x^{-1} \cdot z \to x^{-1} \cdot (y^{-1} \cdot y) = x^{-1} \cdot w$
\item $\varphi_{51} :\equiv \forall x \in \mathbb{R}\,(0 < x \to 0 < x^{-1}) \to y \in \mathbb{R} \to 0 < y \to 0 < y^{-1}$
\item $\varphi_{52} :\equiv \forall x \in \mathbb{R}\,(0 < x \to 0 < x^{-1}) \to x \in \mathbb{R} \to 0 < x \to 0 < x^{-1}$
\item $\varphi_{53} :\equiv \forall x \in \mathbb{R}\,(x \neq 0 \to x \cdot x^{-1} = 1) \to x \in \mathbb{R} \to x \neq 0 \to x \cdot x^{-1} = 1$
\item $\varphi_{54} :\equiv x \in \mathbb{R} \to y \in \mathbb{R} \to 0 < x \to x < y \to x^{-1} \cdot y^{-1} \cdot x = y^{-1}$
\item $\varphi_{55} :\equiv \forall x \in \mathbb{R}\,(x \neq 0 \to x \cdot x^{-1} = 1) \to y \in \mathbb{R} \to y \neq 0 \to y \cdot y^{-1} = 1$
\item $\varphi_{56} :\equiv x \in \mathbb{R} \to y \in \mathbb{R} \to 0 < x \to x < y \to x^{-1} \cdot y^{-1} \cdot y = x^{-1}$
\item $\varphi_{57} :\equiv \forall y\,(x^{-1} \in \mathbb{R} \to y \in \mathbb{R} \to 0 < x^{-1} \to 0 < y \to 0 < x^{-1} \cdot y)$
\item $\varphi_{58} :\equiv \forall w\,(x^{-1} \cdot y^{-1} \cdot y = w \to y^{-1} < x^{-1} \cdot y^{-1} \cdot y \to y^{-1} < w)$
\item $\varphi_{59} :\equiv y \in \mathbb{R} \to 0 < x \to x < y \to x \cdot x^{-1} \cdot y^{-1} = x \cdot (x^{-1} \cdot y^{-1})$
\item $\varphi_{60} :\equiv y \in \mathbb{R} \to 0 < x \to x < y \to x \cdot (x^{-1} \cdot y^{-1}) = x \cdot x^{-1} \cdot y^{-1}$
\item $\varphi_{61} :\equiv y \in \mathbb{R} \to 0 < x \to x < y \to x^{-1} \cdot y^{-1} \cdot x = x \cdot (x^{-1} \cdot y^{-1})$
\item $\varphi_{62} :\equiv y \in \mathbb{R} \to 0 < x \to x < y \to x^{-1} \cdot y^{-1} \cdot x = x \cdot x^{-1} \cdot y^{-1}$
\item $\varphi_{63} :\equiv y^{-1} \in \mathbb{R} \to y \in \mathbb{R} \to x^{-1} \cdot y^{-1} \cdot y = x^{-1} \cdot (y^{-1} \cdot y)$
\item $\varphi_{64} :\equiv y \in \mathbb{R} \to 0 < x \to x < y \to x^{-1} \cdot y^{-1} \cdot y = x^{-1} \cdot (y^{-1} \cdot y)$
\item $\varphi_{65} :\equiv y \in \mathbb{R} \to 0 < x \to x < y \to x^{-1} \cdot y^{-1} \cdot y = x^{-1} \cdot (y \cdot y^{-1})$
\item $\varphi_{66} :\equiv x \in \mathbb{R} \to y \in \mathbb{R} \to 0 < x \to x < y \to y^{-1} < x^{-1} \cdot y^{-1} \cdot y$
\item $\varphi_{67} :\equiv \forall x \in \mathbb{R}\,(x \neq 0 \to x^{-1} \in \mathbb{R}) \to y \in \mathbb{R} \to y \neq 0 \to y^{-1} \in \mathbb{R}$
\item $\varphi_{68} :\equiv \forall x \in \mathbb{R}\,(x \neq 0 \to x^{-1} \in \mathbb{R}) \to x \in \mathbb{R} \to x \neq 0 \to x^{-1} \in \mathbb{R}$
\item $\varphi_{69} :\equiv \forall z\,(x \in \mathbb{R} \to y \in \mathbb{R} \to z \in \mathbb{R} \to 0 < z \to x < y \to z \cdot x < z \cdot y)$
\item $\varphi_{70} :\equiv x \in \mathbb{R} \to y \in \mathbb{R} \to 0 < x \to x < y \to x \cdot x^{-1} \cdot y^{-1} = 1 \cdot y^{-1}$
\item $\varphi_{71} :\equiv x \in \mathbb{R} \to y \in \mathbb{R} \to 0 < x \to x < y \to x^{-1} \cdot y^{-1} \cdot x = 1 \cdot y^{-1}$
\item $\varphi_{72} :\equiv x \in \mathbb{R} \to y \in \mathbb{R} \to 0 < x \to x < y \to x^{-1} \cdot y^{-1} \cdot x = y^{-1} \cdot 1$
\item $\varphi_{73} :\equiv x \in \mathbb{R} \to y \in \mathbb{R} \to 0 < x \to x < y \to x^{-1} \cdot (y \cdot y^{-1}) = x^{-1} \cdot 1$
\item $\varphi_{74} :\equiv x \in \mathbb{R} \to y \in \mathbb{R} \to 0 < x \to x < y \to x^{-1} \cdot y^{-1} \cdot y = x^{-1} \cdot 1$
\item $\varphi_{75} :\equiv y \in \mathbb{R} \to 0 < x \to x < y \to x^{-1} \cdot y^{-1} \cdot x < x^{-1} \cdot y^{-1} \cdot y$
\item $\varphi_{76} :\equiv \forall v_{540}\,\forall v_{541}\,(x^{-1} \cdot y^{-1} \cdot x = v_{540} \to v_{540} = v_{541} \to x^{-1} \cdot y^{-1} \cdot x = v_{541})$
\item $\varphi_{77} :\equiv \forall v_{537}\,\forall v_{538}\,(x^{-1} \cdot y^{-1} \cdot x = v_{537} \to v_{537} = v_{538} \to x^{-1} \cdot y^{-1} \cdot x = v_{538})$
\item $\varphi_{78} :\equiv \forall v_{534}\,\forall v_{535}\,(x^{-1} \cdot y^{-1} \cdot x = v_{534} \to v_{534} = v_{535} \to x^{-1} \cdot y^{-1} \cdot x = v_{535})$
\item $\varphi_{79} :\equiv \forall v_{531}\,\forall v_{532}\,(x^{-1} \cdot y^{-1} \cdot x = v_{531} \to v_{531} = v_{532} \to x^{-1} \cdot y^{-1} \cdot x = v_{532})$
\item $\varphi_{80} :\equiv \forall v_{551}\,\forall v_{552}\,(x^{-1} \cdot y^{-1} \cdot y = v_{551} \to v_{551} = v_{552} \to x^{-1} \cdot y^{-1} \cdot y = v_{552})$
\item $\varphi_{81} :\equiv \forall v_{548}\,\forall v_{549}\,(x^{-1} \cdot y^{-1} \cdot y = v_{548} \to v_{548} = v_{549} \to x^{-1} \cdot y^{-1} \cdot y = v_{549})$
\item $\varphi_{82} :\equiv \forall v_{545}\,\forall v_{546}\,(x^{-1} \cdot y^{-1} \cdot y = v_{545} \to v_{545} = v_{546} \to x^{-1} \cdot y^{-1} \cdot y = v_{546})$
\item $\varphi_{83} :\equiv 0 < x^{-1} \cdot y^{-1} \to x < y \to x^{-1} \cdot y^{-1} \cdot x < x^{-1} \cdot y^{-1} \cdot y$
\item $\varphi_{84} :\equiv x \cdot x^{-1} \cdot y^{-1} = x \cdot x^{-1} \cdot y^{-1} \to x \cdot x^{-1} \cdot y^{-1} = 1 \cdot y^{-1}$
\item $\varphi_{85} :\equiv x^{-1} \in \mathbb{R} \to y^{-1} \in \mathbb{R} \to x \cdot x^{-1} \cdot y^{-1} = x \cdot (x^{-1} \cdot y^{-1})$
\item $\varphi_{86} :\equiv x^{-1} \cdot (y \cdot y^{-1}) = x^{-1} \cdot (y \cdot y^{-1}) \to x^{-1} \cdot (y \cdot y^{-1}) = x^{-1} \cdot 1$
\item $\varphi_{87} :\equiv \forall z\,(x \in \mathbb{R} \to x^{-1} \in \mathbb{R} \to z \in \mathbb{R} \to x \cdot x^{-1} \cdot z = x \cdot (x^{-1} \cdot z))$
\item $\varphi_{88} :\equiv x^{-1} \cdot y^{-1} \cdot x < x^{-1} \cdot y^{-1} \cdot y \to y^{-1} < x^{-1} \cdot y^{-1} \cdot y$
\item $\varphi_{89} :\equiv x^{-1} \cdot y^{-1} \cdot y = x^{-1} \to y^{-1} < x^{-1} \cdot y^{-1} \cdot y \to y^{-1} < x^{-1}$
\item $\varphi_{90} :\equiv \forall y\,\forall z\,(x^{-1} \in \mathbb{R} \to y \in \mathbb{R} \to z \in \mathbb{R} \to x^{-1} \cdot y \cdot z = x^{-1} \cdot (y \cdot z))$
\item $\varphi_{91} :\equiv x^{-1} \cdot y^{-1} \in \mathbb{R} \to x \in \mathbb{R} \to x^{-1} \cdot y^{-1} \cdot x = x \cdot (x^{-1} \cdot y^{-1})$
\item $\varphi_{92} :\equiv x^{-1} \in \mathbb{R} \to y^{-1} \in \mathbb{R} \to 0 < x^{-1} \to 0 < y^{-1} \to 0 < x^{-1} \cdot y^{-1}$
\item $\varphi_{93} :\equiv x^{-1} \cdot y^{-1} \cdot x = y^{-1} \cdot 1 \to y^{-1} \cdot 1 = y^{-1} \to x^{-1} \cdot y^{-1} \cdot x = y^{-1}$
\item $\varphi_{94} :\equiv x^{-1} \cdot y^{-1} \cdot y = x^{-1} \cdot 1 \to x^{-1} \cdot 1 = x^{-1} \to x^{-1} \cdot y^{-1} \cdot y = x^{-1}$
\item $\varphi_{95} :\equiv \forall v_{541}\,(x^{-1} \cdot y^{-1} \cdot x = y^{-1} \cdot 1 \to y^{-1} \cdot 1 = v_{541} \to x^{-1} \cdot y^{-1} \cdot x = v_{541})$
\item $\varphi_{96} :\equiv \forall v_{538}\,(x^{-1} \cdot y^{-1} \cdot x = 1 \cdot y^{-1} \to 1 \cdot y^{-1} = v_{538} \to x^{-1} \cdot y^{-1} \cdot x = v_{538})$
\item $\varphi_{97} :\equiv \forall v_{552}\,(x^{-1} \cdot y^{-1} \cdot y = x^{-1} \cdot 1 \to x^{-1} \cdot 1 = v_{552} \to x^{-1} \cdot y^{-1} \cdot y = v_{552})$
\item $\varphi_{98} :\equiv x \cdot x^{-1} \cdot y^{-1} = x \cdot (x^{-1} \cdot y^{-1}) \to x \cdot (x^{-1} \cdot y^{-1}) = x \cdot x^{-1} \cdot y^{-1}$
\item $\varphi_{99} :\equiv x \cdot (x^{-1} \cdot y^{-1}) = x \cdot x^{-1} \cdot y^{-1} \to x^{-1} \cdot y^{-1} \cdot x = x \cdot x^{-1} \cdot y^{-1}$
\item $\varphi_{100} :\equiv x^{-1} \cdot (y^{-1} \cdot y) = x^{-1} \cdot (y \cdot y^{-1}) \to x^{-1} \cdot y^{-1} \cdot y = x^{-1} \cdot (y \cdot y^{-1})$
\item $\varphi_{101} :\equiv x^{-1} \cdot (y^{-1} \cdot y) = x^{-1} \cdot (y^{-1} \cdot y) \to x^{-1} \cdot (y^{-1} \cdot y) = x^{-1} \cdot (y \cdot y^{-1})$
\item $\varphi_{102} :\equiv \forall x\,\forall y\,(x \in \mathbb{R} \to y \in \mathbb{R} \to x \cdot y = y \cdot x) \to \forall y\,(1 \in \mathbb{R} \to y \in \mathbb{R} \to 1 \cdot y = y \cdot 1)$
\item $\varphi_{103} :\equiv x^{-1} \cdot y^{-1} \cdot x = 1 \cdot y^{-1} \to 1 \cdot y^{-1} = y^{-1} \cdot 1 \to x^{-1} \cdot y^{-1} \cdot x = y^{-1} \cdot 1$
\item $\varphi_{104} :\equiv \forall z\,(x^{-1} \in \mathbb{R} \to y^{-1} \in \mathbb{R} \to z \in \mathbb{R} \to x^{-1} \cdot y^{-1} \cdot z = x^{-1} \cdot (y^{-1} \cdot z))$
\item $\varphi_{105} :\equiv \forall w\,(x \cdot x^{-1} = w \to x \cdot x^{-1} \cdot y^{-1} = x \cdot x^{-1} \cdot y^{-1} \to x \cdot x^{-1} \cdot y^{-1} = w \cdot y^{-1})$
\item $\varphi_{106} :\equiv \forall w\,(y \cdot y^{-1} = w \to x^{-1} \cdot (y \cdot y^{-1}) = x^{-1} \cdot (y \cdot y^{-1}) \to x^{-1} \cdot (y \cdot y^{-1}) = x^{-1} \cdot w)$
\item $\varphi_{107} :\equiv \forall w\,(y^{-1} \cdot y = w \to x^{-1} \cdot (y^{-1} \cdot y) = x^{-1} \cdot (y^{-1} \cdot y) \to x^{-1} \cdot (y^{-1} \cdot y) = x^{-1} \cdot w)$
\item $\varphi_{108} :\equiv \forall y\,(1 \in \mathbb{R} \to y \in \mathbb{R} \to 1 \cdot y = y \cdot 1) \to 1 \in \mathbb{R} \to y^{-1} \in \mathbb{R} \to 1 \cdot y^{-1} = y^{-1} \cdot 1$
\item $\varphi_{109} :\equiv \forall x\,\forall y\,(x \in \mathbb{R} \to y \in \mathbb{R} \to x \cdot y = y \cdot x) \to \forall y\,(v_{526} \in \mathbb{R} \to y \in \mathbb{R} \to v_{526} \cdot y = y \cdot v_{526})$
\item $\varphi_{110} :\equiv \forall v_{527}\,(x^{-1} \cdot y^{-1} \in \mathbb{R} \to v_{527} \in \mathbb{R} \to x^{-1} \cdot y^{-1} \cdot v_{527} = v_{527} \cdot (x^{-1} \cdot y^{-1}))$
\item $\varphi_{111} :\equiv \forall x\,\forall y\,(x \in \mathbb{R} \to y \in \mathbb{R} \to x \cdot y = y \cdot x) \to \forall y\,(v_{542} \in \mathbb{R} \to y \in \mathbb{R} \to v_{542} \cdot y = y \cdot v_{542})$
\item $\varphi_{112} :\equiv (x \in \mathbb{R} \to y \in \mathbb{R} \to 0 < x \to x < y \to 0 < y) \to x \in \mathbb{R} \to y \in \mathbb{R} \to 0 < x \to x < y \to 0 \neq y$
\item $\varphi_{113} :\equiv (x \in \mathbb{R} \to y \in \mathbb{R} \to 0 < x \to x < y \to 0 < x) \to x \in \mathbb{R} \to y \in \mathbb{R} \to 0 < x \to x < y \to 0 \neq x$
\item $\varphi_{114} :\equiv (x \in \mathbb{R} \to y \in \mathbb{R} \to 0 < x \to x < y \to x < y) \to x \in \mathbb{R} \to y \in \mathbb{R} \to 0 < x \to x < y \to 0 < y$
\item $\varphi_{115} :\equiv \forall w\,(x^{-1} \cdot y^{-1} \cdot x = w \to x^{-1} \cdot y^{-1} \cdot x < x^{-1} \cdot y^{-1} \cdot y \to w < x^{-1} \cdot y^{-1} \cdot y)$
\item $\varphi_{116} :\equiv \forall x\,\forall y\,(x \in \mathbb{R} \to y \in \mathbb{R} \to x \cdot y \in \mathbb{R}) \to \forall y\,(x^{-1} \in \mathbb{R} \to y \in \mathbb{R} \to x^{-1} \cdot y \in \mathbb{R})$
\item $\varphi_{117} :\equiv \forall v_{535}\,(x^{-1} \cdot y^{-1} \cdot x = x \cdot x^{-1} \cdot y^{-1} \to x \cdot x^{-1} \cdot y^{-1} = v_{535} \to x^{-1} \cdot y^{-1} \cdot x = v_{535})$
\item $\varphi_{118} :\equiv \forall v_{532}\,(x^{-1} \cdot y^{-1} \cdot x = x \cdot (x^{-1} \cdot y^{-1}) \to x \cdot (x^{-1} \cdot y^{-1}) = v_{532} \to x^{-1} \cdot y^{-1} \cdot x = v_{532})$
\item $\varphi_{119} :\equiv \forall v_{549}\,(x^{-1} \cdot y^{-1} \cdot y = x^{-1} \cdot (y \cdot y^{-1}) \to x^{-1} \cdot (y \cdot y^{-1}) = v_{549} \to x^{-1} \cdot y^{-1} \cdot y = v_{549})$
\item $\varphi_{120} :\equiv \forall v_{546}\,(x^{-1} \cdot y^{-1} \cdot y = x^{-1} \cdot (y^{-1} \cdot y) \to x^{-1} \cdot (y^{-1} \cdot y) = v_{546} \to x^{-1} \cdot y^{-1} \cdot y = v_{546})$
\item $\varphi_{121} :\equiv (x \in \mathbb{R} \to y \in \mathbb{R} \to 0 < x \to x < y \to 0 < y) \to x \in \mathbb{R} \to y \in \mathbb{R} \to 0 < x \to x < y \to 0 < y^{-1}$
\item $\varphi_{122} :\equiv (x \in \mathbb{R} \to y \in \mathbb{R} \to 0 < x \to x < y \to 0 < x) \to x \in \mathbb{R} \to y \in \mathbb{R} \to 0 < x \to x < y \to 0 < x^{-1}$
\item $\varphi_{123} :\equiv (x \in \mathbb{R} \to y \in \mathbb{R} \to 0 < x \to x < y \to x \neq 0) \to x \in \mathbb{R} \to y \in \mathbb{R} \to 0 < x \to x < y \to x \cdot x^{-1} = 1$
\item $\varphi_{124} :\equiv (x \in \mathbb{R} \to y \in \mathbb{R} \to 0 < x \to x < y \to y \neq 0) \to x \in \mathbb{R} \to y \in \mathbb{R} \to 0 < x \to x < y \to y \cdot y^{-1} = 1$
\item $\varphi_{125} :\equiv (x \in \mathbb{R} \to y \in \mathbb{R} \to 0 < x \to x < y \to y \neq 0) \to x \in \mathbb{R} \to y \in \mathbb{R} \to 0 < x \to x < y \to y^{-1} \in \mathbb{R}$
\item $\varphi_{126} :\equiv (x \in \mathbb{R} \to y \in \mathbb{R} \to 0 < x \to x < y \to x \neq 0) \to x \in \mathbb{R} \to y \in \mathbb{R} \to 0 < x \to x < y \to x^{-1} \in \mathbb{R}$
\item $\varphi_{127} :\equiv \forall y\,(v_{526} \in \mathbb{R} \to y \in \mathbb{R} \to v_{526} \cdot y = y \cdot v_{526}) \to v_{526} \in \mathbb{R} \to v_{527} \in \mathbb{R} \to v_{526} \cdot v_{527} = v_{527} \cdot v_{526}$
\item $\varphi_{128} :\equiv \forall y\,(v_{542} \in \mathbb{R} \to y \in \mathbb{R} \to v_{542} \cdot y = y \cdot v_{542}) \to v_{542} \in \mathbb{R} \to v_{543} \in \mathbb{R} \to v_{542} \cdot v_{543} = v_{543} \cdot v_{542}$
\item $\varphi_{129} :\equiv \forall y\,(x^{-1} \in \mathbb{R} \to y \in \mathbb{R} \to x^{-1} \cdot y \in \mathbb{R}) \to x^{-1} \in \mathbb{R} \to y^{-1} \in \mathbb{R} \to x^{-1} \cdot y^{-1} \in \mathbb{R}$
\item $\varphi_{130} :\equiv (x \in \mathbb{R} \to y \in \mathbb{R} \to 0 < x \to x < y \to 0 < y^{-1}) \to x \in \mathbb{R} \to y \in \mathbb{R} \to 0 < x \to x < y \to 0 < x^{-1} \cdot y^{-1}$
\item $\varphi_{131} :\equiv (x \in \mathbb{R} \to y \in \mathbb{R} \to 0 < x \to x < y \to y^{-1} \in \mathbb{R}) \to x \in \mathbb{R} \to y \in \mathbb{R} \to 0 < x \to x < y \to y^{-1} \cdot 1 = y^{-1}$
\item $\varphi_{132} :\equiv (x \in \mathbb{R} \to y \in \mathbb{R} \to 0 < x \to x < y \to x^{-1} \in \mathbb{R}) \to x \in \mathbb{R} \to y \in \mathbb{R} \to 0 < x \to x < y \to x^{-1} \cdot 1 = x^{-1}$
\item $\varphi_{133} :\equiv (x \in \mathbb{R} \to y \in \mathbb{R} \to 0 < x \to x < y \to y \in \mathbb{R}) \to x \in \mathbb{R} \to y \in \mathbb{R} \to 0 < x \to x < y \to y^{-1} \cdot y = y \cdot y^{-1}$
\item $\varphi_{134} :\equiv (x \in \mathbb{R} \to y \in \mathbb{R} \to 0 < x \to x < y \to y^{-1} \in \mathbb{R}) \to x \in \mathbb{R} \to y \in \mathbb{R} \to 0 < x \to x < y \to 1 \cdot y^{-1} = y^{-1} \cdot 1$
\item $\varphi_{135} :\equiv (x \in \mathbb{R} \to y \in \mathbb{R} \to 0 < x \to x < y \to y^{-1} \in \mathbb{R}) \to x \in \mathbb{R} \to y \in \mathbb{R} \to 0 < x \to x < y \to x^{-1} \cdot y^{-1} \in \mathbb{R}$
\item $\varphi_{136} :\equiv \forall x\,\forall y\,(x \in \mathbb{R} \to y \in \mathbb{R} \to 0 < x \to 0 < y \to 0 < x \cdot y) \to \varphi_{57}$
\item $\varphi_{137} :\equiv (x \in \mathbb{R} \to y \in \mathbb{R} \to 0 < x \to x < y \to x \in \mathbb{R}) \to x \in \mathbb{R} \to \varphi_{61}$
\item $\varphi_{138} :\equiv (x \in \mathbb{R} \to y \in \mathbb{R} \to 0 < x \to x < y \to y \in \mathbb{R}) \to x \in \mathbb{R} \to \varphi_{64}$
\item $\varphi_{139} :\equiv (x \in \mathbb{R} \to y \in \mathbb{R} \to 0 < x \to x < y \to y^{-1} \in \mathbb{R}) \to x \in \mathbb{R} \to \varphi_{59}$
\item $\varphi_{140} :\equiv (0 < x \to x < y \to 0 < y) \to (x \in \mathbb{R} \to y \in \mathbb{R} \to 0 < x \to x < y \to 0 < x) \to \varphi_{114}$
\item $\varphi_{141} :\equiv (1 \in \mathbb{R} \to y^{-1} \in \mathbb{R} \to 1 \cdot y^{-1} = y^{-1} \cdot 1) \to 1 \in \mathbb{R} \to \varphi_{134}$
\item $\varphi_{142} :\equiv (y \in \mathbb{R} \to 0 < y \to 0 < y^{-1}) \to (x \in \mathbb{R} \to y \in \mathbb{R} \to 0 < x \to x < y \to y \in \mathbb{R}) \to \varphi_{121}$
\item $\varphi_{143} :\equiv (x \in \mathbb{R} \to 0 < x \to 0 < x^{-1}) \to (x \in \mathbb{R} \to y \in \mathbb{R} \to 0 < x \to x < y \to x \in \mathbb{R}) \to \varphi_{122}$
\item $\varphi_{144} :\equiv (x \in \mathbb{R} \to x \neq 0 \to x \cdot x^{-1} = 1) \to (x \in \mathbb{R} \to y \in \mathbb{R} \to 0 < x \to x < y \to x \in \mathbb{R}) \to \varphi_{123}$
\item $\varphi_{145} :\equiv (y \in \mathbb{R} \to y \neq 0 \to y \cdot y^{-1} = 1) \to (x \in \mathbb{R} \to y \in \mathbb{R} \to 0 < x \to x < y \to y \in \mathbb{R}) \to \varphi_{124}$
\item $\varphi_{146} :\equiv (y \in \mathbb{R} \to y \neq 0 \to y^{-1} \in \mathbb{R}) \to (x \in \mathbb{R} \to y \in \mathbb{R} \to 0 < x \to x < y \to y \in \mathbb{R}) \to \varphi_{125}$
\item $\varphi_{147} :\equiv (x \in \mathbb{R} \to x \neq 0 \to x^{-1} \in \mathbb{R}) \to (x \in \mathbb{R} \to y \in \mathbb{R} \to 0 < x \to x < y \to x \in \mathbb{R}) \to \varphi_{126}$
\item $\varphi_{148} :\equiv (x \in \mathbb{R} \to y \in \mathbb{R} \to 0 < x \to x < y \to 0 < x^{-1} \cdot y^{-1}) \to (x \in \mathbb{R} \to y \in \mathbb{R} \to 0 < x \to x < y \to x < y) \to x \in \mathbb{R} \to \varphi_{75}$
\item $\varphi_{149} :\equiv (x \in \mathbb{R} \to y \in \mathbb{R} \to 0 < x \to x < y \to x^{-1} \cdot y^{-1} \in \mathbb{R}) \to \varphi_{137}$
\item $\varphi_{150} :\equiv x \cdot x^{-1} \cdot y^{-1} = x \cdot x^{-1} \cdot y^{-1} \to (x \cdot x^{-1} = 1 \to \varphi_{84}) \to \varphi_{70}$
\item $\varphi_{151} :\equiv x^{-1} \cdot (y \cdot y^{-1}) = x^{-1} \cdot (y \cdot y^{-1}) \to (y \cdot y^{-1} = 1 \to \varphi_{86}) \to \varphi_{73}$
\item $\varphi_{152} :\equiv (x^{-1} \cdot y^{-1} \cdot x = x \cdot (x^{-1} \cdot y^{-1}) \to \varphi_{99}) \to x \in \mathbb{R} \to \varphi_{62}$
\item $\varphi_{153} :\equiv (x \in \mathbb{R} \to y \in \mathbb{R} \to 0 < x \to x < y \to y^{-1} \cdot 1 = y^{-1}) \to \varphi_{93} \to \varphi_{54}$
\item $\varphi_{154} :\equiv (x^{-1} \cdot y^{-1} \cdot y = x^{-1} \cdot (y^{-1} \cdot y) \to \varphi_{100}) \to x \in \mathbb{R} \to \varphi_{65}$
\item $\varphi_{155} :\equiv (x \in \mathbb{R} \to y \in \mathbb{R} \to 0 < x \to x < y \to x^{-1} \cdot 1 = x^{-1}) \to \varphi_{94} \to \varphi_{56}$
\item $\varphi_{156} :\equiv (x \in \mathbb{R} \to y \in \mathbb{R} \to 0 < x \to x < y \to y^{-1} \in \mathbb{R}) \to (x \in \mathbb{R} \to y \in \mathbb{R} \to 0 < x \to x < y \to 0 < x^{-1}) \to \varphi_{130}$
\item $\varphi_{157} :\equiv (y^{-1} \in \mathbb{R} \to y \in \mathbb{R} \to y^{-1} \cdot y = y \cdot y^{-1}) \to (x \in \mathbb{R} \to y \in \mathbb{R} \to 0 < x \to x < y \to y^{-1} \in \mathbb{R}) \to \varphi_{133}$
\item $\varphi_{158} :\equiv (x \in \mathbb{R} \to y \in \mathbb{R} \to 0 < x \to x < y \to 1 \cdot y^{-1} = y^{-1} \cdot 1) \to \varphi_{103} \to \varphi_{72}$
\item $\varphi_{159} :\equiv x^{-1} \cdot (y^{-1} \cdot y) = x^{-1} \cdot (y^{-1} \cdot y) \to (y^{-1} \cdot y = y \cdot y^{-1} \to \varphi_{101}) \to x \in \mathbb{R} \to y \in \mathbb{R} \to \varphi_{31}$
\item $\varphi_{160} :\equiv (x^{-1} \in \mathbb{R} \to y^{-1} \in \mathbb{R} \to x^{-1} \cdot y^{-1} \in \mathbb{R}) \to (x \in \mathbb{R} \to y \in \mathbb{R} \to 0 < x \to x < y \to x^{-1} \in \mathbb{R}) \to \varphi_{135}$
\item $\varphi_{161} :\equiv (x \in \mathbb{R} \to y \in \mathbb{R} \to 0 < x \to x < y \to x \in \mathbb{R}) \to (x \in \mathbb{R} \to y \in \mathbb{R} \to 0 < x \to x < y \to x^{-1} \in \mathbb{R}) \to \varphi_{139}$
\item $\varphi_{162} :\equiv (x \in \mathbb{R} \to y \in \mathbb{R} \to 0 < x \to x < y \to x^{-1} \in \mathbb{R}) \to (x \in \mathbb{R} \to y \in \mathbb{R} \to 0 < x \to x < y \to y^{-1} \in \mathbb{R}) \to \varphi_{138}$
\item $\varphi_{163} :\equiv (x \in \mathbb{R} \to y \in \mathbb{R} \to 0 < x \to x < y \to x^{-1} \cdot y^{-1} \in \mathbb{R}) \to \varphi_{148}$
\item $\varphi_{164} :\equiv \varphi_{56} \to \varphi_{66} \to \varphi_{89} \to x \in \mathbb{R} \to y \in \mathbb{R} \to 0 < x \to x < y \to y^{-1} < x^{-1}$
\item $\varphi_{165} :\equiv (x \in \mathbb{R} \to y \in \mathbb{R} \to 0 < x \to x < y \to y^{-1} \cdot y = y \cdot y^{-1}) \to \varphi_{159}$
\item $\varphi_{166} :\equiv (x \in \mathbb{R} \to \varphi_{62}) \to \varphi_{70} \to (x^{-1} \cdot y^{-1} \cdot x = x \cdot x^{-1} \cdot y^{-1} \to \varphi_{46}) \to \varphi_{71}$
\item $\varphi_{167} :\equiv (x \in \mathbb{R} \to \varphi_{65}) \to \varphi_{73} \to (x^{-1} \cdot y^{-1} \cdot y = x^{-1} \cdot (y \cdot y^{-1}) \to \varphi_{48}) \to \varphi_{74}$
\item $\varphi_{168} :\equiv (x \in \mathbb{R} \to y \in \mathbb{R} \to 0 < x \to x < y \to x \in \mathbb{R}) \to (x \in \mathbb{R} \to y \in \mathbb{R} \to 0 < x \to x < y \to y \in \mathbb{R}) \to \varphi_{163}$
\end{abbrevlist}
\end{fullproof}

% 自動生成 (prover/build.js): InvLe  ノード数 50
\begin{fullproof}{[E20]}{$\forall x\,\forall y\,(x \in \mathbb{R} \to y \in \mathbb{R} \to 0 < x \to x \leq y \to y^{-1} \leq x^{-1})$}{50}\phantomsection\label{tree:InvLe}
\noindent$\Pi_{1}$:
\begin{ptree}
\AXC{$x \in \mathbb{R} \to y \in \mathbb{R} \to 0 < x \to x \leq y \to x = y \to x = y$\lbl{A1}}
\AXC{$(x \in \mathbb{R} \to y \in \mathbb{R} \to 0 < x \to x \leq y \to x = y \to x = y) \to \varphi_{15}$\lbl{A1}}
\BIC{$\varphi_{15}$}
\end{ptree}
\noindent$\Pi_{2}$:
\begin{ptree}
\AXC{$x = x$\lbl{A5}}
\UIC{$\forall x\,(x = x)$}
\AXC{$\forall x\,(x = x) \to x^{-1} = x^{-1}$\lbl{A2}}
\BIC{$x^{-1} = x^{-1}$}
\AXC{$\varphi_{4}$\lbl{A1}}
\BIC{$x \in \mathbb{R} \to y \in \mathbb{R} \to 0 < x \to x \leq y \to x = y \to x^{-1} \leq x^{-1}$}
\AXC{$\stackrel{\Pi_{1}}{\vdots}$}\noLine
\UIC{$\varphi_{15}$}
\BIC{$(x = y \to x^{-1} < x^{-1} \to y^{-1} < x^{-1}) \to \varphi_{8}$}
\end{ptree}
\noindent$\Pi_{3}$:
\begin{ptree}
\AXC{$z = w \to z^{-1} < x^{-1} \to w^{-1} < x^{-1}$\lbl{A6}}
\UIC{$\forall w\,(z = w \to z^{-1} < x^{-1} \to w^{-1} < x^{-1})$}
\UIC{$\forall z\,\forall w\,(z = w \to z^{-1} < x^{-1} \to w^{-1} < x^{-1})$}
\AXC{$\varphi_{7}$\lbl{A2}}
\BIC{$\forall w\,(x = w \to x^{-1} < x^{-1} \to w^{-1} < x^{-1})$}
\AXC{$\varphi_{6}$\lbl{A2}}
\BIC{$x = y \to x^{-1} < x^{-1} \to y^{-1} < x^{-1}$}
\AXC{$\stackrel{\Pi_{2}}{\vdots}$}\noLine
\UIC{$(x = y \to x^{-1} < x^{-1} \to y^{-1} < x^{-1}) \to \varphi_{8}$}
\BIC{$\varphi_{8}$}
\end{ptree}
\noindent$\Pi_{4}$:
\begin{ptree}
\AXC{$\varphi_{2}$\lbl{[E19]}}
\AXC{$\varphi_{12}$\lbl{A2}}
\BIC{$\forall y\,(x \in \mathbb{R} \to y \in \mathbb{R} \to 0 < x \to x < y \to y^{-1} < x^{-1})$}
\AXC{$\varphi_{11}$\lbl{A2}}
\BIC{$x \in \mathbb{R} \to y \in \mathbb{R} \to 0 < x \to x < y \to y^{-1} < x^{-1}$}
\AXC{$\varphi_{17}$\lbl{A1}}
\BIC{$(x \in \mathbb{R} \to y \in \mathbb{R} \to 0 < x \to x \leq y \to x < y \to x \in \mathbb{R}) \to \varphi_{16}$}
\end{ptree}
\noindent$\Pi_{5}$:
\begin{ptree}
\AXC{$x \in \mathbb{R} \to y \in \mathbb{R} \to 0 < x \to x \leq y \to x < y \to x \in \mathbb{R}$\lbl{A1}}
\AXC{$\stackrel{\Pi_{4}}{\vdots}$}\noLine
\UIC{$(x \in \mathbb{R} \to y \in \mathbb{R} \to 0 < x \to x \leq y \to x < y \to x \in \mathbb{R}) \to \varphi_{16}$}
\BIC{$\varphi_{16}$}
\end{ptree}
\noindent$\Pi_{6}$:
\begin{ptree}
\AXC{$x \in \mathbb{R} \to y \in \mathbb{R} \to 0 < x \to x \leq y \to 0 < x$\lbl{A1}}
\AXC{$x \in \mathbb{R} \to y \in \mathbb{R} \to 0 < x \to x \leq y \to x < y \to y \in \mathbb{R}$\lbl{A1}}
\AXC{$\stackrel{\Pi_{5}}{\vdots}$}\noLine
\UIC{$\varphi_{16}$}
\BIC{$(x \in \mathbb{R} \to y \in \mathbb{R} \to 0 < x \to x \leq y \to 0 < x) \to \varphi_{10}$}
\BIC{$\varphi_{10}$}
\end{ptree}
\noindent$\Pi_{7}$:
\begin{ptree}
\AXC{$x \in \mathbb{R} \to y \in \mathbb{R} \to 0 < x \to x \leq y \to x \leq y$\lbl{A1}}
\AXC{$(x \in \mathbb{R} \to y \in \mathbb{R} \to 0 < x \to x \leq y \to x \leq y) \to \varphi_{14}$\lbl{A1}}
\BIC{$\varphi_{14}$}
\end{ptree}
\noindent$\Pi_{8}$:
\begin{ptree}
\AXC{$x \in \mathbb{R} \to y \in \mathbb{R} \to 0 < x \to x \leq y \to x < y \to x < y$\lbl{A1}}
\AXC{$\stackrel{\Pi_{6}}{\vdots}$}\noLine
\UIC{$\varphi_{10}$}
\BIC{$x \in \mathbb{R} \to y \in \mathbb{R} \to 0 < x \to x \leq y \to x < y \to y^{-1} < x^{-1}$}
\AXC{$\varphi_{13}$\lbl{A1}}
\BIC{$x \in \mathbb{R} \to y \in \mathbb{R} \to 0 < x \to x \leq y \to x < y \to y^{-1} \leq x^{-1}$}
\AXC{$\stackrel{\Pi_{7}}{\vdots}$}\noLine
\UIC{$\varphi_{14}$}
\BIC{$\varphi_{9}$}
\end{ptree}
\noindent 主木:
\begin{ptree}
\AXC{$z = w \to z^{-1} = x^{-1} \to w^{-1} = x^{-1}$\lbl{A6}}
\UIC{$\forall w\,(z = w \to z^{-1} = x^{-1} \to w^{-1} = x^{-1})$}
\UIC{$\forall z\,\forall w\,(z = w \to z^{-1} = x^{-1} \to w^{-1} = x^{-1})$}
\AXC{$\varphi_{5}$\lbl{A2}}
\BIC{$\forall w\,(x = w \to x^{-1} = x^{-1} \to w^{-1} = x^{-1})$}
\AXC{$\varphi_{3}$\lbl{A2}}
\BIC{$x = y \to x^{-1} = x^{-1} \to y^{-1} = x^{-1}$}
\AXC{$\stackrel{\Pi_{3}}{\vdots}$}\noLine
\UIC{$\varphi_{8}$}
\BIC{$x \in \mathbb{R} \to y \in \mathbb{R} \to 0 < x \to x \leq y \to x = y \to y^{-1} \leq x^{-1}$}
\AXC{$\stackrel{\Pi_{8}}{\vdots}$}\noLine
\UIC{$\varphi_{9}$}
\BIC{$x \in \mathbb{R} \to y \in \mathbb{R} \to 0 < x \to x \leq y \to y^{-1} \leq x^{-1}$}
\UIC{$\forall y\,(x \in \mathbb{R} \to y \in \mathbb{R} \to 0 < x \to x \leq y \to y^{-1} \leq x^{-1})$}
\UIC{$\varphi_{1}$}
\end{ptree}
\begin{abbrevlist}
\item $\varphi_{1} :\equiv \forall x\,\forall y\,(x \in \mathbb{R} \to y \in \mathbb{R} \to 0 < x \to x \leq y \to y^{-1} \leq x^{-1})$
\item $\varphi_{2} :\equiv \forall x\,\forall y\,(x \in \mathbb{R} \to y \in \mathbb{R} \to 0 < x \to x < y \to y^{-1} < x^{-1})$
\item $\varphi_{3} :\equiv \forall w\,(x = w \to x^{-1} = x^{-1} \to w^{-1} = x^{-1}) \to x = y \to x^{-1} = x^{-1} \to y^{-1} = x^{-1}$
\item $\varphi_{4} :\equiv x^{-1} = x^{-1} \to x \in \mathbb{R} \to y \in \mathbb{R} \to 0 < x \to x \leq y \to x = y \to x^{-1} \leq x^{-1}$
\item $\varphi_{5} :\equiv \forall z\,\forall w\,(z = w \to z^{-1} = x^{-1} \to w^{-1} = x^{-1}) \to \forall w\,(x = w \to x^{-1} = x^{-1} \to w^{-1} = x^{-1})$
\item $\varphi_{6} :\equiv \forall w\,(x = w \to x^{-1} < x^{-1} \to w^{-1} < x^{-1}) \to x = y \to x^{-1} < x^{-1} \to y^{-1} < x^{-1}$
\item $\varphi_{7} :\equiv \forall z\,\forall w\,(z = w \to z^{-1} < x^{-1} \to w^{-1} < x^{-1}) \to \forall w\,(x = w \to x^{-1} < x^{-1} \to w^{-1} < x^{-1})$
\item $\varphi_{8} :\equiv (x = y \to x^{-1} = x^{-1} \to y^{-1} = x^{-1}) \to x \in \mathbb{R} \to y \in \mathbb{R} \to 0 < x \to x \leq y \to x = y \to y^{-1} \leq x^{-1}$
\item $\varphi_{9} :\equiv (x \in \mathbb{R} \to y \in \mathbb{R} \to 0 < x \to x \leq y \to x = y \to y^{-1} \leq x^{-1}) \to x \in \mathbb{R} \to y \in \mathbb{R} \to 0 < x \to x \leq y \to y^{-1} \leq x^{-1}$
\item $\varphi_{10} :\equiv (x \in \mathbb{R} \to y \in \mathbb{R} \to 0 < x \to x \leq y \to x < y \to x < y) \to x \in \mathbb{R} \to y \in \mathbb{R} \to 0 < x \to x \leq y \to x < y \to y^{-1} < x^{-1}$
\item $\varphi_{11} :\equiv \forall y\,(x \in \mathbb{R} \to y \in \mathbb{R} \to 0 < x \to x < y \to y^{-1} < x^{-1}) \to x \in \mathbb{R} \to y \in \mathbb{R} \to 0 < x \to x < y \to y^{-1} < x^{-1}$
\item $\varphi_{12} :\equiv \varphi_{2} \to \forall y\,(x \in \mathbb{R} \to y \in \mathbb{R} \to 0 < x \to x < y \to y^{-1} < x^{-1})$
\item $\varphi_{13} :\equiv (x \in \mathbb{R} \to y \in \mathbb{R} \to 0 < x \to x \leq y \to x < y \to y^{-1} < x^{-1}) \to x \in \mathbb{R} \to y \in \mathbb{R} \to 0 < x \to x \leq y \to x < y \to y^{-1} \leq x^{-1}$
\item $\varphi_{14} :\equiv (x \in \mathbb{R} \to y \in \mathbb{R} \to 0 < x \to x \leq y \to x < y \to y^{-1} \leq x^{-1}) \to \varphi_{9}$
\item $\varphi_{15} :\equiv (x \in \mathbb{R} \to y \in \mathbb{R} \to 0 < x \to x \leq y \to x = y \to x^{-1} \leq x^{-1}) \to (x = y \to x^{-1} < x^{-1} \to y^{-1} < x^{-1}) \to \varphi_{8}$
\item $\varphi_{16} :\equiv (x \in \mathbb{R} \to y \in \mathbb{R} \to 0 < x \to x \leq y \to x < y \to y \in \mathbb{R}) \to (x \in \mathbb{R} \to y \in \mathbb{R} \to 0 < x \to x \leq y \to 0 < x) \to \varphi_{10}$
\item $\varphi_{17} :\equiv (x \in \mathbb{R} \to y \in \mathbb{R} \to 0 < x \to x < y \to y^{-1} < x^{-1}) \to (x \in \mathbb{R} \to y \in \mathbb{R} \to 0 < x \to x \leq y \to x < y \to x \in \mathbb{R}) \to \varphi_{16}$
\end{abbrevlist}
\end{fullproof}

% 自動生成 (prover/build.js): SubInv  ノード数 589
\begin{fullproof}{[E21]}{$\forall x\,\forall y\,(x \in \mathbb{R} \to y \in \mathbb{R} \to x \neq 0 \to y \neq 0 \to x^{-1} - y^{-1} = (y - x) \cdot (x^{-1} \cdot y^{-1}))$}{589}\phantomsection\label{tree:SubInv}
\noindent$\Pi_{1}$:
\begin{ptree}
\AXC{$y = z \to x = y \to x = z$\lbl{A6}}
\AXC{$(y = z \to x = y \to x = z) \to x = y \to y = z \to x = z$\lbl{A1}}
\BIC{$x = y \to y = z \to x = z$}
\UIC{$\forall z\,(x = y \to y = z \to x = z)$}
\UIC{$\forall y\,\forall z\,(x = y \to y = z \to x = z)$}
\UIC{$\forall x\,\forall y\,\forall z\,(x = y \to y = z \to x = z)$}
\AXC{$\varphi_{18}$\lbl{A2}}
\BIC{$\forall y\,\forall z\,(v_{608} = y \to y = z \to v_{608} = z)$}
\AXC{$\varphi_{36}$\lbl{A2}}
\BIC{$\forall z\,(v_{608} = v_{609} \to v_{609} = z \to v_{608} = z)$}
\AXC{$\varphi_{61}$\lbl{A2}}
\BIC{$v_{608} = v_{609} \to v_{609} = v_{610} \to v_{608} = v_{610}$}
\UIC{$\forall v_{610}\,(v_{608} = v_{609} \to v_{609} = v_{610} \to v_{608} = v_{610})$}
\UIC{$\varphi_{5}$}
\UIC{$\forall v_{608}\,(\varphi_{5})$}
\end{ptree}
\noindent$\Pi_{2}$:
\begin{ptree}
\AXC{$\stackrel{\Pi_{1}}{\vdots}$}\noLine
\UIC{$\forall v_{608}\,(\varphi_{5})$}
\AXC{$\forall v_{608}\,(\varphi_{5}) \to \varphi_{78}$\lbl{A2}}
\BIC{$\varphi_{78}$}
\AXC{$\varphi_{78} \to \varphi_{174}$\lbl{A2}}
\BIC{$\varphi_{174}$}
\AXC{$\varphi_{174} \to x^{-1} - y^{-1} = x^{-1} \cdot y^{-1} \cdot (y - x) \to \varphi_{165}$\lbl{A2}}
\BIC{$x^{-1} - y^{-1} = x^{-1} \cdot y^{-1} \cdot (y - x) \to \varphi_{165}$}
\end{ptree}
\noindent$\Pi_{3}$:
\begin{ptree}
\AXC{$\forall x\,\forall y\,(x \in \mathbb{R} \to y \in \mathbb{R} \to x - y \in \mathbb{R})$\lbl{[F1]}}
\AXC{$\varphi_{166}$\lbl{A2}}
\BIC{$\forall y\,(v_{594} \in \mathbb{R} \to y \in \mathbb{R} \to v_{594} - y \in \mathbb{R})$}
\AXC{$\varphi_{173}$\lbl{A2}}
\BIC{$v_{594} \in \mathbb{R} \to v_{595} \in \mathbb{R} \to v_{594} - v_{595} \in \mathbb{R}$}
\UIC{$\forall v_{595}\,(v_{594} \in \mathbb{R} \to v_{595} \in \mathbb{R} \to v_{594} - v_{595} \in \mathbb{R})$}
\UIC{$\varphi_{58}$}
\AXC{$\varphi_{58} \to \forall v_{595}\,(y \in \mathbb{R} \to v_{595} \in \mathbb{R} \to y - v_{595} \in \mathbb{R})$\lbl{A2}}
\BIC{$\forall v_{595}\,(y \in \mathbb{R} \to v_{595} \in \mathbb{R} \to y - v_{595} \in \mathbb{R})$}
\end{ptree}
\noindent$\Pi_{4}$:
\begin{ptree}
\AXC{$x \in \mathbb{R} \to y \in \mathbb{R} \to x \neq 0 \to y \neq 0 \to y \in \mathbb{R}$\lbl{A1}}
\AXC{$\stackrel{\Pi_{3}}{\vdots}$}\noLine
\UIC{$\forall v_{595}\,(y \in \mathbb{R} \to v_{595} \in \mathbb{R} \to y - v_{595} \in \mathbb{R})$}
\AXC{$\varphi_{156}$\lbl{A2}}
\BIC{$y \in \mathbb{R} \to x \in \mathbb{R} \to y - x \in \mathbb{R}$}
\AXC{$\varphi_{214}$\lbl{A1}}
\BIC{$(x \in \mathbb{R} \to y \in \mathbb{R} \to x \neq 0 \to y \neq 0 \to y \in \mathbb{R}) \to \varphi_{178}$}
\BIC{$\varphi_{178}$}
\end{ptree}
\noindent$\Pi_{5}$:
\begin{ptree}
\AXC{$x \in \mathbb{R} \to y \in \mathbb{R} \to x \neq 0 \to y \neq 0 \to y \in \mathbb{R}$\lbl{A1}}
\AXC{$\forall x \in \mathbb{R}\,(x \neq 0 \to x^{-1} \in \mathbb{R})$\lbl{(R1)}}
\AXC{$\varphi_{114}$\lbl{A2}}
\BIC{$y \in \mathbb{R} \to y \neq 0 \to y^{-1} \in \mathbb{R}$}
\AXC{$\varphi_{211}$\lbl{A1}}
\BIC{$(x \in \mathbb{R} \to y \in \mathbb{R} \to x \neq 0 \to y \neq 0 \to y \in \mathbb{R}) \to \varphi_{169}$}
\BIC{$\varphi_{169}$}
\end{ptree}
\noindent$\Pi_{6}$:
\begin{ptree}
\AXC{$x \in \mathbb{R} \to y \in \mathbb{R} \to x \neq 0 \to y \neq 0 \to y \neq 0$\lbl{A1}}
\AXC{$\stackrel{\Pi_{5}}{\vdots}$}\noLine
\UIC{$\varphi_{169}$}
\BIC{$x \in \mathbb{R} \to y \in \mathbb{R} \to x \neq 0 \to y \neq 0 \to y^{-1} \in \mathbb{R}$}
\end{ptree}
\noindent$\Pi_{7}$:
\begin{ptree}
\AXC{$x \in \mathbb{R} \to y \in \mathbb{R} \to x \neq 0 \to y \neq 0 \to x \in \mathbb{R}$\lbl{A1}}
\AXC{$\forall x \in \mathbb{R}\,(x \neq 0 \to x^{-1} \in \mathbb{R})$\lbl{(R1)}}
\AXC{$\varphi_{115}$\lbl{A2}}
\BIC{$x \in \mathbb{R} \to x \neq 0 \to x^{-1} \in \mathbb{R}$}
\AXC{$\varphi_{212}$\lbl{A1}}
\BIC{$(x \in \mathbb{R} \to y \in \mathbb{R} \to x \neq 0 \to y \neq 0 \to x \in \mathbb{R}) \to \varphi_{170}$}
\BIC{$\varphi_{170}$}
\end{ptree}
\noindent$\Pi_{8}$:
\begin{ptree}
\AXC{$x \in \mathbb{R} \to y \in \mathbb{R} \to x \neq 0 \to y \neq 0 \to x \neq 0$\lbl{A1}}
\AXC{$\stackrel{\Pi_{7}}{\vdots}$}\noLine
\UIC{$\varphi_{170}$}
\BIC{$x \in \mathbb{R} \to y \in \mathbb{R} \to x \neq 0 \to y \neq 0 \to x^{-1} \in \mathbb{R}$}
\end{ptree}
\noindent$\Pi_{9}$:
\begin{ptree}
\AXC{$\forall x\,\forall y\,(x \in \mathbb{R} \to y \in \mathbb{R} \to x \cdot y \in \mathbb{R})$\lbl{(R1)}}
\AXC{$\varphi_{171}$\lbl{A2}}
\BIC{$\forall y\,(x^{-1} \in \mathbb{R} \to y \in \mathbb{R} \to x^{-1} \cdot y \in \mathbb{R})$}
\AXC{$\varphi_{184}$\lbl{A2}}
\BIC{$x^{-1} \in \mathbb{R} \to y^{-1} \in \mathbb{R} \to x^{-1} \cdot y^{-1} \in \mathbb{R}$}
\AXC{$\varphi_{227}$\lbl{A1}}
\BIC{$(x \in \mathbb{R} \to y \in \mathbb{R} \to x \neq 0 \to y \neq 0 \to x^{-1} \in \mathbb{R}) \to \varphi_{188}$}
\end{ptree}
\noindent$\Pi_{10}$:
\begin{ptree}
\AXC{$\stackrel{\Pi_{8}}{\vdots}$}\noLine
\UIC{$x \in \mathbb{R} \to y \in \mathbb{R} \to x \neq 0 \to y \neq 0 \to x^{-1} \in \mathbb{R}$}
\AXC{$\stackrel{\Pi_{9}}{\vdots}$}\noLine
\UIC{$(x \in \mathbb{R} \to y \in \mathbb{R} \to x \neq 0 \to y \neq 0 \to x^{-1} \in \mathbb{R}) \to \varphi_{188}$}
\BIC{$\varphi_{188}$}
\end{ptree}
\noindent$\Pi_{11}$:
\begin{ptree}
\AXC{$\stackrel{\Pi_{6}}{\vdots}$}\noLine
\UIC{$x \in \mathbb{R} \to y \in \mathbb{R} \to x \neq 0 \to y \neq 0 \to y^{-1} \in \mathbb{R}$}
\AXC{$\stackrel{\Pi_{10}}{\vdots}$}\noLine
\UIC{$\varphi_{188}$}
\BIC{$x \in \mathbb{R} \to y \in \mathbb{R} \to x \neq 0 \to y \neq 0 \to x^{-1} \cdot y^{-1} \in \mathbb{R}$}
\end{ptree}
\noindent$\Pi_{12}$:
\begin{ptree}
\AXC{$\forall x\,\forall y\,(x \in \mathbb{R} \to y \in \mathbb{R} \to x \cdot y = y \cdot x)$\lbl{(R7)}}
\AXC{$\varphi_{160}$\lbl{A2}}
\BIC{$\forall y\,(v_{592} \in \mathbb{R} \to y \in \mathbb{R} \to v_{592} \cdot y = y \cdot v_{592})$}
\AXC{$\varphi_{179}$\lbl{A2}}
\BIC{$v_{592} \in \mathbb{R} \to v_{593} \in \mathbb{R} \to v_{592} \cdot v_{593} = v_{593} \cdot v_{592}$}
\UIC{$\varphi_{54}$}
\UIC{$\forall v_{592}\,(\varphi_{54})$}
\AXC{$\forall v_{592}\,(\varphi_{54}) \to \varphi_{161}$\lbl{A2}}
\BIC{$\varphi_{161}$}
\end{ptree}
\noindent$\Pi_{13}$:
\begin{ptree}
\AXC{$\stackrel{\Pi_{12}}{\vdots}$}\noLine
\UIC{$\varphi_{161}$}
\AXC{$\varphi_{161} \to x^{-1} \cdot y^{-1} \in \mathbb{R} \to \varphi_{97}$\lbl{A2}}
\BIC{$x^{-1} \cdot y^{-1} \in \mathbb{R} \to \varphi_{97}$}
\AXC{$(x^{-1} \cdot y^{-1} \in \mathbb{R} \to \varphi_{97}) \to \varphi_{221}$\lbl{A1}}
\BIC{$\varphi_{221}$}
\end{ptree}
\noindent$\Pi_{14}$:
\begin{ptree}
\AXC{$x \in \mathbb{R} \to y \in \mathbb{R} \to x \neq 0 \to y \neq 0 \to x \in \mathbb{R}$\lbl{A1}}
\AXC{$\stackrel{\Pi_{4}}{\vdots}$}\noLine
\UIC{$\varphi_{178}$}
\BIC{$x \in \mathbb{R} \to y \in \mathbb{R} \to x \neq 0 \to y \neq 0 \to y - x \in \mathbb{R}$}
\AXC{$\stackrel{\Pi_{11}}{\vdots}$}\noLine
\UIC{$x \in \mathbb{R} \to y \in \mathbb{R} \to x \neq 0 \to y \neq 0 \to x^{-1} \cdot y^{-1} \in \mathbb{R}$}
\AXC{$\stackrel{\Pi_{13}}{\vdots}$}\noLine
\UIC{$\varphi_{221}$}
\BIC{$\varphi_{199}$}
\BIC{$x \in \mathbb{R} \to y \in \mathbb{R} \to x \neq 0 \to \varphi_{48}$}
\end{ptree}
\noindent$\Pi_{15}$:
\begin{ptree}
\AXC{$y = z \to x = y \to x = z$\lbl{A6}}
\AXC{$(y = z \to x = y \to x = z) \to x = y \to y = z \to x = z$\lbl{A1}}
\BIC{$x = y \to y = z \to x = z$}
\UIC{$\forall z\,(x = y \to y = z \to x = z)$}
\UIC{$\forall y\,\forall z\,(x = y \to y = z \to x = z)$}
\UIC{$\forall x\,\forall y\,\forall z\,(x = y \to y = z \to x = z)$}
\AXC{$\varphi_{19}$\lbl{A2}}
\BIC{$\forall y\,\forall z\,(v_{605} = y \to y = z \to v_{605} = z)$}
\AXC{$\varphi_{37}$\lbl{A2}}
\BIC{$\forall z\,(v_{605} = v_{606} \to v_{606} = z \to v_{605} = z)$}
\AXC{$\varphi_{62}$\lbl{A2}}
\BIC{$v_{605} = v_{606} \to v_{606} = v_{607} \to v_{605} = v_{607}$}
\UIC{$\forall v_{607}\,(v_{605} = v_{606} \to v_{606} = v_{607} \to v_{605} = v_{607})$}
\UIC{$\varphi_{6}$}
\UIC{$\forall v_{605}\,(\varphi_{6})$}
\end{ptree}
\noindent$\Pi_{16}$:
\begin{ptree}
\AXC{$\varphi_{116}$\lbl{[F17]}}
\AXC{$\varphi_{116} \to \varphi_{134}$\lbl{A2}}
\BIC{$\varphi_{134}$}
\AXC{$\varphi_{134} \to \varphi_{147}$\lbl{A2}}
\BIC{$\varphi_{147}$}
\AXC{$\varphi_{147} \to v_{589} \in \mathbb{R} \to v_{590} \in \mathbb{R} \to \varphi_{60}$\lbl{A2}}
\BIC{$v_{589} \in \mathbb{R} \to v_{590} \in \mathbb{R} \to \varphi_{60}$}
\UIC{$\forall v_{591}\,(v_{589} \in \mathbb{R} \to v_{590} \in \mathbb{R} \to \varphi_{60})$}
\UIC{$\forall v_{590}\,\forall v_{591}\,(v_{589} \in \mathbb{R} \to v_{590} \in \mathbb{R} \to \varphi_{60})$}
\UIC{$\forall v_{589}\,\forall v_{590}\,\forall v_{591}\,(v_{589} \in \mathbb{R} \to v_{590} \in \mathbb{R} \to \varphi_{60})$}
\AXC{$\varphi_{231}$\lbl{A2}}
\BIC{$\varphi_{202}$}
\AXC{$\varphi_{202} \to \varphi_{195}$\lbl{A2}}
\BIC{$\varphi_{195}$}
\end{ptree}
\noindent$\Pi_{17}$:
\begin{ptree}
\AXC{$\stackrel{\Pi_{16}}{\vdots}$}\noLine
\UIC{$\varphi_{195}$}
\AXC{$\varphi_{195} \to x^{-1} \cdot y^{-1} \in \mathbb{R} \to y \in \mathbb{R} \to x \in \mathbb{R} \to \varphi_{109}$\lbl{A2}}
\BIC{$x^{-1} \cdot y^{-1} \in \mathbb{R} \to y \in \mathbb{R} \to x \in \mathbb{R} \to \varphi_{109}$}
\AXC{$(x^{-1} \cdot y^{-1} \in \mathbb{R} \to y \in \mathbb{R} \to x \in \mathbb{R} \to \varphi_{109}) \to \varphi_{232}$\lbl{A1}}
\BIC{$\varphi_{232}$}
\end{ptree}
\noindent$\Pi_{18}$:
\begin{ptree}
\AXC{$x \in \mathbb{R} \to y \in \mathbb{R} \to x \neq 0 \to y \neq 0 \to y \in \mathbb{R}$\lbl{A1}}
\AXC{$\stackrel{\Pi_{11}}{\vdots}$}\noLine
\UIC{$x \in \mathbb{R} \to y \in \mathbb{R} \to x \neq 0 \to y \neq 0 \to x^{-1} \cdot y^{-1} \in \mathbb{R}$}
\AXC{$\stackrel{\Pi_{17}}{\vdots}$}\noLine
\UIC{$\varphi_{232}$}
\BIC{$(x \in \mathbb{R} \to y \in \mathbb{R} \to x \neq 0 \to y \neq 0 \to y \in \mathbb{R}) \to \varphi_{204}$}
\BIC{$\varphi_{204}$}
\end{ptree}
\noindent$\Pi_{19}$:
\begin{ptree}
\AXC{$y = z \to x = y \to x = z$\lbl{A6}}
\AXC{$(y = z \to x = y \to x = z) \to x = y \to y = z \to x = z$\lbl{A1}}
\BIC{$x = y \to y = z \to x = z$}
\UIC{$\forall z\,(x = y \to y = z \to x = z)$}
\UIC{$\forall y\,\forall z\,(x = y \to y = z \to x = z)$}
\UIC{$\forall x\,\forall y\,\forall z\,(x = y \to y = z \to x = z)$}
\AXC{$\varphi_{20}$\lbl{A2}}
\BIC{$\forall y\,\forall z\,(v_{602} = y \to y = z \to v_{602} = z)$}
\AXC{$\varphi_{38}$\lbl{A2}}
\BIC{$\forall z\,(v_{602} = v_{603} \to v_{603} = z \to v_{602} = z)$}
\AXC{$\varphi_{63}$\lbl{A2}}
\BIC{$v_{602} = v_{603} \to v_{603} = v_{604} \to v_{602} = v_{604}$}
\UIC{$\forall v_{604}\,(v_{602} = v_{603} \to v_{603} = v_{604} \to v_{602} = v_{604})$}
\UIC{$\varphi_{7}$}
\UIC{$\forall v_{602}\,(\varphi_{7})$}
\end{ptree}
\noindent$\Pi_{20}$:
\begin{ptree}
\AXC{$\stackrel{\Pi_{19}}{\vdots}$}\noLine
\UIC{$\forall v_{602}\,(\varphi_{7})$}
\AXC{$\forall v_{602}\,(\varphi_{7}) \to \varphi_{80}$\lbl{A2}}
\BIC{$\varphi_{80}$}
\AXC{$\varphi_{80} \to \varphi_{197}$\lbl{A2}}
\BIC{$\varphi_{197}$}
\AXC{$\varphi_{197} \to \varphi_{223}$\lbl{A2}}
\BIC{$\varphi_{223}$}
\end{ptree}
\noindent$\Pi_{21}$:
\begin{ptree}
\AXC{$\forall x\,\forall y\,(x \in \mathbb{R} \to y \in \mathbb{R} \to x \cdot y = y \cdot x)$\lbl{(R7)}}
\AXC{$\varphi_{145}$\lbl{A2}}
\BIC{$\forall y\,(x \in \mathbb{R} \to y \in \mathbb{R} \to x \cdot y = y \cdot x)$}
\AXC{$\forall y\,(x \in \mathbb{R} \to y \in \mathbb{R} \to x \cdot y = y \cdot x) \to x \in \mathbb{R} \to \varphi_{98}$\lbl{A2}}
\BIC{$x \in \mathbb{R} \to \varphi_{98}$}
\end{ptree}
\noindent$\Pi_{22}$:
\begin{ptree}
\AXC{$x \in \mathbb{R} \to y \in \mathbb{R} \to x \neq 0 \to y \neq 0 \to x \in \mathbb{R}$\lbl{A1}}
\AXC{$\stackrel{\Pi_{21}}{\vdots}$}\noLine
\UIC{$x \in \mathbb{R} \to \varphi_{98}$}
\AXC{$\varphi_{233}$\lbl{A1}}
\BIC{$(x \in \mathbb{R} \to y \in \mathbb{R} \to x \neq 0 \to y \neq 0 \to x \in \mathbb{R}) \to \varphi_{200}$}
\BIC{$\varphi_{200}$}
\end{ptree}
\noindent$\Pi_{23}$:
\begin{ptree}
\AXC{$\stackrel{\Pi_{11}}{\vdots}$}\noLine
\UIC{$x \in \mathbb{R} \to y \in \mathbb{R} \to x \neq 0 \to y \neq 0 \to x^{-1} \cdot y^{-1} \in \mathbb{R}$}
\AXC{$\stackrel{\Pi_{22}}{\vdots}$}\noLine
\UIC{$\varphi_{200}$}
\BIC{$x \in \mathbb{R} \to \varphi_{87}$}
\AXC{$(x \in \mathbb{R} \to \varphi_{87}) \to \varphi_{150} \to \varphi_{241}$\lbl{A1}}
\BIC{$\varphi_{150} \to \varphi_{241}$}
\end{ptree}
\noindent$\Pi_{24}$:
\begin{ptree}
\AXC{$z = w \to \varphi_{111} \to \varphi_{110}$\lbl{A6}}
\UIC{$\forall w\,(z = w \to \varphi_{111} \to \varphi_{110})$}
\UIC{$\forall z\,\forall w\,(z = w \to \varphi_{111} \to \varphi_{110})$}
\AXC{$\varphi_{239}$\lbl{A2}}
\BIC{$\forall w\,(x \cdot (x^{-1} \cdot y^{-1}) = w \to \varphi_{150} \to \varphi_{110})$}
\AXC{$\varphi_{243}$\lbl{A2}}
\BIC{$x \cdot (x^{-1} \cdot y^{-1}) = x^{-1} \cdot y^{-1} \cdot x \to \varphi_{150} \to \varphi_{152}$}
\AXC{$x = x$\lbl{A5}}
\UIC{$\forall x\,(x = x)$}
\AXC{$\forall x\,(x = x) \to \varphi_{150}$\lbl{A2}}
\BIC{$\varphi_{150}$}
\AXC{$\stackrel{\Pi_{23}}{\vdots}$}\noLine
\UIC{$\varphi_{150} \to \varphi_{241}$}
\BIC{$\varphi_{241}$}
\BIC{$x \in \mathbb{R} \to y \in \mathbb{R} \to x \neq 0 \to y \neq 0 \to \varphi_{152}$}
\end{ptree}
\noindent$\Pi_{25}$:
\begin{ptree}
\AXC{$y = z \to x = y \to x = z$\lbl{A6}}
\AXC{$(y = z \to x = y \to x = z) \to x = y \to y = z \to x = z$\lbl{A1}}
\BIC{$x = y \to y = z \to x = z$}
\UIC{$\forall z\,(x = y \to y = z \to x = z)$}
\UIC{$\forall y\,\forall z\,(x = y \to y = z \to x = z)$}
\UIC{$\forall x\,\forall y\,\forall z\,(x = y \to y = z \to x = z)$}
\AXC{$\varphi_{21}$\lbl{A2}}
\BIC{$\forall y\,\forall z\,(v_{599} = y \to y = z \to v_{599} = z)$}
\AXC{$\varphi_{39}$\lbl{A2}}
\BIC{$\forall z\,(v_{599} = v_{600} \to v_{600} = z \to v_{599} = z)$}
\AXC{$\varphi_{64}$\lbl{A2}}
\BIC{$v_{599} = v_{600} \to v_{600} = v_{601} \to v_{599} = v_{601}$}
\UIC{$\forall v_{601}\,(v_{599} = v_{600} \to v_{600} = v_{601} \to v_{599} = v_{601})$}
\UIC{$\varphi_{8}$}
\UIC{$\forall v_{599}\,(\varphi_{8})$}
\end{ptree}
\noindent$\Pi_{26}$:
\begin{ptree}
\AXC{$\stackrel{\Pi_{25}}{\vdots}$}\noLine
\UIC{$\forall v_{599}\,(\varphi_{8})$}
\AXC{$\forall v_{599}\,(\varphi_{8}) \to \varphi_{81}$\lbl{A2}}
\BIC{$\varphi_{81}$}
\AXC{$\varphi_{81} \to \varphi_{198}$\lbl{A2}}
\BIC{$\varphi_{198}$}
\AXC{$\varphi_{198} \to \varphi_{57} \to \varphi_{203}$\lbl{A2}}
\BIC{$\varphi_{57} \to \varphi_{203}$}
\end{ptree}
\noindent$\Pi_{27}$:
\begin{ptree}
\AXC{$\forall x\,\forall y\,(x \in \mathbb{R} \to y \in \mathbb{R} \to x \cdot y = y \cdot x)$\lbl{(R7)}}
\AXC{$\varphi_{162}$\lbl{A2}}
\BIC{$\forall y\,(v_{587} \in \mathbb{R} \to y \in \mathbb{R} \to v_{587} \cdot y = y \cdot v_{587})$}
\AXC{$\varphi_{180}$\lbl{A2}}
\BIC{$v_{587} \in \mathbb{R} \to v_{588} \in \mathbb{R} \to v_{587} \cdot v_{588} = v_{588} \cdot v_{587}$}
\UIC{$\varphi_{55}$}
\UIC{$\forall v_{587}\,(\varphi_{55})$}
\AXC{$\varphi_{189}$\lbl{A2}}
\BIC{$\forall v_{588}\,(y \in \mathbb{R} \to v_{588} \in \mathbb{R} \to y \cdot v_{588} = v_{588} \cdot y)$}
\AXC{$\varphi_{192}$\lbl{A2}}
\BIC{$y \in \mathbb{R} \to \varphi_{99}$}
\AXC{$\varphi_{234}$\lbl{A1}}
\BIC{$(x \in \mathbb{R} \to y \in \mathbb{R} \to x \neq 0 \to y \neq 0 \to y \in \mathbb{R}) \to \varphi_{201}$}
\end{ptree}
\noindent$\Pi_{28}$:
\begin{ptree}
\AXC{$x \in \mathbb{R} \to y \in \mathbb{R} \to x \neq 0 \to y \neq 0 \to y \in \mathbb{R}$\lbl{A1}}
\AXC{$\stackrel{\Pi_{27}}{\vdots}$}\noLine
\UIC{$(x \in \mathbb{R} \to y \in \mathbb{R} \to x \neq 0 \to y \neq 0 \to y \in \mathbb{R}) \to \varphi_{201}$}
\BIC{$\varphi_{201}$}
\end{ptree}
\noindent$\Pi_{29}$:
\begin{ptree}
\AXC{$\stackrel{\Pi_{11}}{\vdots}$}\noLine
\UIC{$x \in \mathbb{R} \to y \in \mathbb{R} \to x \neq 0 \to y \neq 0 \to x^{-1} \cdot y^{-1} \in \mathbb{R}$}
\AXC{$\stackrel{\Pi_{28}}{\vdots}$}\noLine
\UIC{$\varphi_{201}$}
\BIC{$x \in \mathbb{R} \to \varphi_{88}$}
\AXC{$(x \in \mathbb{R} \to \varphi_{88}) \to \varphi_{151} \to \varphi_{242}$\lbl{A1}}
\BIC{$\varphi_{151} \to \varphi_{242}$}
\end{ptree}
\noindent$\Pi_{30}$:
\begin{ptree}
\AXC{$z = w \to \varphi_{113} \to \varphi_{112}$\lbl{A6}}
\UIC{$\forall w\,(z = w \to \varphi_{113} \to \varphi_{112})$}
\UIC{$\forall z\,\forall w\,(z = w \to \varphi_{113} \to \varphi_{112})$}
\AXC{$\varphi_{240}$\lbl{A2}}
\BIC{$\forall w\,(y \cdot (x^{-1} \cdot y^{-1}) = w \to \varphi_{151} \to \varphi_{112})$}
\AXC{$\varphi_{244}$\lbl{A2}}
\BIC{$y \cdot (x^{-1} \cdot y^{-1}) = x^{-1} \cdot y^{-1} \cdot y \to \varphi_{151} \to \varphi_{153}$}
\AXC{$x = x$\lbl{A5}}
\UIC{$\forall x\,(x = x)$}
\AXC{$\forall x\,(x = x) \to \varphi_{151}$\lbl{A2}}
\BIC{$\varphi_{151}$}
\AXC{$\stackrel{\Pi_{29}}{\vdots}$}\noLine
\UIC{$\varphi_{151} \to \varphi_{242}$}
\BIC{$\varphi_{242}$}
\BIC{$x \in \mathbb{R} \to y \in \mathbb{R} \to x \neq 0 \to y \neq 0 \to \varphi_{153}$}
\end{ptree}
\noindent$\Pi_{31}$:
\begin{ptree}
\AXC{$y = z \to x = y \to x = z$\lbl{A6}}
\AXC{$(y = z \to x = y \to x = z) \to x = y \to y = z \to x = z$\lbl{A1}}
\BIC{$x = y \to y = z \to x = z$}
\UIC{$\forall z\,(x = y \to y = z \to x = z)$}
\UIC{$\forall y\,\forall z\,(x = y \to y = z \to x = z)$}
\UIC{$\forall x\,\forall y\,\forall z\,(x = y \to y = z \to x = z)$}
\AXC{$\varphi_{22}$\lbl{A2}}
\BIC{$\forall y\,\forall z\,(v_{596} = y \to y = z \to v_{596} = z)$}
\AXC{$\varphi_{40}$\lbl{A2}}
\BIC{$\forall z\,(v_{596} = v_{597} \to v_{597} = z \to v_{596} = z)$}
\AXC{$\varphi_{65}$\lbl{A2}}
\BIC{$v_{596} = v_{597} \to v_{597} = v_{598} \to v_{596} = v_{598}$}
\UIC{$\forall v_{598}\,(v_{596} = v_{597} \to v_{597} = v_{598} \to v_{596} = v_{598})$}
\UIC{$\varphi_{9}$}
\UIC{$\forall v_{596}\,(\varphi_{9})$}
\end{ptree}
\noindent$\Pi_{32}$:
\begin{ptree}
\AXC{$\stackrel{\Pi_{31}}{\vdots}$}\noLine
\UIC{$\forall v_{596}\,(\varphi_{9})$}
\AXC{$\forall v_{596}\,(\varphi_{9}) \to \varphi_{82}$\lbl{A2}}
\BIC{$\varphi_{82}$}
\AXC{$\varphi_{82} \to \varphi_{183}$\lbl{A2}}
\BIC{$\varphi_{183}$}
\AXC{$\varphi_{183} \to x^{-1} - y^{-1} = y \cdot (x^{-1} \cdot y^{-1}) - y^{-1} \to \varphi_{119} \to \varphi_{57}$\lbl{A2}}
\BIC{$x^{-1} - y^{-1} = y \cdot (x^{-1} \cdot y^{-1}) - y^{-1} \to \varphi_{119} \to \varphi_{57}$}
\end{ptree}
\noindent$\Pi_{33}$:
\begin{ptree}
\AXC{$z = w \to \varphi_{185}$\lbl{A6}}
\UIC{$\forall w\,(z = w \to \varphi_{185})$}
\UIC{$\forall z\,\forall w\,(z = w \to \varphi_{185})$}
\AXC{$\forall z\,\forall w\,(z = w \to \varphi_{185}) \to \varphi_{191}$\lbl{A2}}
\BIC{$\varphi_{191}$}
\AXC{$\varphi_{191} \to y^{-1} = x \cdot (x^{-1} \cdot y^{-1}) \to \varphi_{56} \to \varphi_{119}$\lbl{A2}}
\BIC{$y^{-1} = x \cdot (x^{-1} \cdot y^{-1}) \to \varphi_{56} \to \varphi_{119}$}
\end{ptree}
\noindent$\Pi_{34}$:
\begin{ptree}
\AXC{$x = x$\lbl{A5}}
\AXC{$x = y \to x = x \to y = x$\lbl{A6}}
\AXC{$(x = y \to x = x \to y = x) \to x = x \to x = y \to y = x$\lbl{A1}}
\BIC{$x = x \to x = y \to y = x$}
\BIC{$x = y \to y = x$}
\UIC{$\forall y\,(x = y \to y = x)$}
\UIC{$\forall x\,\forall y\,(x = y \to y = x)$}
\end{ptree}
\noindent$\Pi_{35}$:
\begin{ptree}
\AXC{$\stackrel{\Pi_{34}}{\vdots}$}\noLine
\UIC{$\forall x\,\forall y\,(x = y \to y = x)$}
\AXC{$\forall x\,\forall y\,(x = y \to y = x) \to \forall y\,(v_{585} = y \to y = v_{585})$\lbl{A2}}
\BIC{$\forall y\,(v_{585} = y \to y = v_{585})$}
\AXC{$\varphi_{1}$\lbl{A2}}
\BIC{$v_{585} = v_{586} \to v_{586} = v_{585}$}
\UIC{$\forall v_{586}\,(v_{585} = v_{586} \to v_{586} = v_{585})$}
\UIC{$\forall v_{585}\,\forall v_{586}\,(v_{585} = v_{586} \to v_{586} = v_{585})$}
\end{ptree}
\noindent$\Pi_{36}$:
\begin{ptree}
\AXC{$\stackrel{\Pi_{35}}{\vdots}$}\noLine
\UIC{$\forall v_{585}\,\forall v_{586}\,(v_{585} = v_{586} \to v_{586} = v_{585})$}
\AXC{$\forall v_{585}\,\forall v_{586}\,(v_{585} = v_{586} \to v_{586} = v_{585}) \to \varphi_{49}$\lbl{A2}}
\BIC{$\varphi_{49}$}
\AXC{$\varphi_{49} \to \varphi_{33}$\lbl{A2}}
\BIC{$\varphi_{33}$}
\end{ptree}
\noindent$\Pi_{37}$:
\begin{ptree}
\AXC{$y = z \to x = y \to x = z$\lbl{A6}}
\AXC{$(y = z \to x = y \to x = z) \to x = y \to y = z \to x = z$\lbl{A1}}
\BIC{$x = y \to y = z \to x = z$}
\UIC{$\forall z\,(x = y \to y = z \to x = z)$}
\UIC{$\forall y\,\forall z\,(x = y \to y = z \to x = z)$}
\UIC{$\forall x\,\forall y\,\forall z\,(x = y \to y = z \to x = z)$}
\AXC{$\varphi_{23}$\lbl{A2}}
\BIC{$\forall y\,\forall z\,(v_{580} = y \to y = z \to v_{580} = z)$}
\AXC{$\varphi_{41}$\lbl{A2}}
\BIC{$\forall z\,(v_{580} = v_{581} \to v_{581} = z \to v_{580} = z)$}
\AXC{$\varphi_{66}$\lbl{A2}}
\BIC{$v_{580} = v_{581} \to v_{581} = v_{582} \to v_{580} = v_{582}$}
\UIC{$\forall v_{582}\,(v_{580} = v_{581} \to v_{581} = v_{582} \to v_{580} = v_{582})$}
\UIC{$\varphi_{10}$}
\UIC{$\forall v_{580}\,(\varphi_{10})$}
\end{ptree}
\noindent$\Pi_{38}$:
\begin{ptree}
\AXC{$\stackrel{\Pi_{37}}{\vdots}$}\noLine
\UIC{$\forall v_{580}\,(\varphi_{10})$}
\AXC{$\forall v_{580}\,(\varphi_{10}) \to \varphi_{120}$\lbl{A2}}
\BIC{$\varphi_{120}$}
\AXC{$\varphi_{120} \to \varphi_{137}$\lbl{A2}}
\BIC{$\varphi_{137}$}
\AXC{$\varphi_{137} \to \varphi_{135}$\lbl{A2}}
\BIC{$\varphi_{135}$}
\end{ptree}
\noindent$\Pi_{39}$:
\begin{ptree}
\AXC{$\forall x \in \mathbb{R}\,(x \cdot 1 = x)$\lbl{(R8)}}
\AXC{$\forall x \in \mathbb{R}\,(x \cdot 1 = x) \to y^{-1} \in \mathbb{R} \to y^{-1} \cdot 1 = y^{-1}$\lbl{A2}}
\BIC{$y^{-1} \in \mathbb{R} \to y^{-1} \cdot 1 = y^{-1}$}
\end{ptree}
\noindent$\Pi_{40}$:
\begin{ptree}
\AXC{$\stackrel{\Pi_{6}}{\vdots}$}\noLine
\UIC{$x \in \mathbb{R} \to y \in \mathbb{R} \to x \neq 0 \to y \neq 0 \to y^{-1} \in \mathbb{R}$}
\AXC{$\stackrel{\Pi_{39}}{\vdots}$}\noLine
\UIC{$y^{-1} \in \mathbb{R} \to y^{-1} \cdot 1 = y^{-1}$}
\AXC{$(y^{-1} \in \mathbb{R} \to y^{-1} \cdot 1 = y^{-1}) \to \varphi_{181}$\lbl{A1}}
\BIC{$\varphi_{181}$}
\BIC{$x \in \mathbb{R} \to y \in \mathbb{R} \to x \neq 0 \to y \neq 0 \to y^{-1} \cdot 1 = y^{-1}$}
\end{ptree}
\noindent$\Pi_{41}$:
\begin{ptree}
\AXC{$y = z \to x = y \to x = z$\lbl{A6}}
\AXC{$(y = z \to x = y \to x = z) \to x = y \to y = z \to x = z$\lbl{A1}}
\BIC{$x = y \to y = z \to x = z$}
\UIC{$\forall z\,(x = y \to y = z \to x = z)$}
\UIC{$\forall y\,\forall z\,(x = y \to y = z \to x = z)$}
\UIC{$\forall x\,\forall y\,\forall z\,(x = y \to y = z \to x = z)$}
\AXC{$\varphi_{24}$\lbl{A2}}
\BIC{$\forall y\,\forall z\,(v_{577} = y \to y = z \to v_{577} = z)$}
\AXC{$\varphi_{42}$\lbl{A2}}
\BIC{$\forall z\,(v_{577} = v_{578} \to v_{578} = z \to v_{577} = z)$}
\AXC{$\varphi_{67}$\lbl{A2}}
\BIC{$v_{577} = v_{578} \to v_{578} = v_{579} \to v_{577} = v_{579}$}
\UIC{$\forall v_{579}\,(v_{577} = v_{578} \to v_{578} = v_{579} \to v_{577} = v_{579})$}
\UIC{$\varphi_{11}$}
\UIC{$\forall v_{577}\,(\varphi_{11})$}
\end{ptree}
\noindent$\Pi_{42}$:
\begin{ptree}
\AXC{$1 \in \mathbb{R}$\lbl{(R1)}}
\AXC{$\forall x\,\forall y\,(x \in \mathbb{R} \to y \in \mathbb{R} \to x \cdot y = y \cdot x)$\lbl{(R7)}}
\AXC{$\varphi_{146}$\lbl{A2}}
\BIC{$\forall y\,(1 \in \mathbb{R} \to y \in \mathbb{R} \to 1 \cdot y = y \cdot 1)$}
\AXC{$\varphi_{163}$\lbl{A2}}
\BIC{$1 \in \mathbb{R} \to y^{-1} \in \mathbb{R} \to 1 \cdot y^{-1} = y^{-1} \cdot 1$}
\AXC{$\varphi_{209}$\lbl{A1}}
\BIC{$1 \in \mathbb{R} \to \varphi_{186}$}
\BIC{$\varphi_{186}$}
\end{ptree}
\noindent$\Pi_{43}$:
\begin{ptree}
\AXC{$y = z \to x = y \to x = z$\lbl{A6}}
\AXC{$(y = z \to x = y \to x = z) \to x = y \to y = z \to x = z$\lbl{A1}}
\BIC{$x = y \to y = z \to x = z$}
\UIC{$\forall z\,(x = y \to y = z \to x = z)$}
\UIC{$\forall y\,\forall z\,(x = y \to y = z \to x = z)$}
\UIC{$\forall x\,\forall y\,\forall z\,(x = y \to y = z \to x = z)$}
\AXC{$\varphi_{25}$\lbl{A2}}
\BIC{$\forall y\,\forall z\,(v_{574} = y \to y = z \to v_{574} = z)$}
\AXC{$\varphi_{43}$\lbl{A2}}
\BIC{$\forall z\,(v_{574} = v_{575} \to v_{575} = z \to v_{574} = z)$}
\AXC{$\varphi_{68}$\lbl{A2}}
\BIC{$v_{574} = v_{575} \to v_{575} = v_{576} \to v_{574} = v_{576}$}
\UIC{$\forall v_{576}\,(v_{574} = v_{575} \to v_{575} = v_{576} \to v_{574} = v_{576})$}
\UIC{$\varphi_{12}$}
\UIC{$\forall v_{574}\,(\varphi_{12})$}
\end{ptree}
\noindent$\Pi_{44}$:
\begin{ptree}
\AXC{$\stackrel{\Pi_{43}}{\vdots}$}\noLine
\UIC{$\forall v_{574}\,(\varphi_{12})$}
\AXC{$\forall v_{574}\,(\varphi_{12}) \to \varphi_{122}$\lbl{A2}}
\BIC{$\varphi_{122}$}
\AXC{$\varphi_{122} \to \varphi_{175}$\lbl{A2}}
\BIC{$\varphi_{175}$}
\AXC{$\varphi_{175} \to x \cdot (x^{-1} \cdot y^{-1}) = x \cdot x^{-1} \cdot y^{-1} \to \varphi_{83}$\lbl{A2}}
\BIC{$x \cdot (x^{-1} \cdot y^{-1}) = x \cdot x^{-1} \cdot y^{-1} \to \varphi_{83}$}
\end{ptree}
\noindent$\Pi_{45}$:
\begin{ptree}
\AXC{$z = w \to \varphi_{84}$\lbl{A6}}
\UIC{$\forall w\,(z = w \to \varphi_{84})$}
\UIC{$\forall z\,\forall w\,(z = w \to \varphi_{84})$}
\AXC{$\forall z\,\forall w\,(z = w \to \varphi_{84}) \to \varphi_{157}$\lbl{A2}}
\BIC{$\varphi_{157}$}
\AXC{$\varphi_{157} \to x \cdot x^{-1} = 1 \to \varphi_{127}$\lbl{A2}}
\BIC{$x \cdot x^{-1} = 1 \to \varphi_{127}$}
\end{ptree}
\noindent$\Pi_{46}$:
\begin{ptree}
\AXC{$x \in \mathbb{R} \to y \in \mathbb{R} \to x \neq 0 \to y \neq 0 \to x \in \mathbb{R}$\lbl{A1}}
\AXC{$\forall x \in \mathbb{R}\,(x \neq 0 \to x \cdot x^{-1} = 1)$\lbl{(R9)}}
\AXC{$\varphi_{100}$\lbl{A2}}
\BIC{$x \in \mathbb{R} \to x \neq 0 \to x \cdot x^{-1} = 1$}
\AXC{$\varphi_{206}$\lbl{A1}}
\BIC{$(x \in \mathbb{R} \to y \in \mathbb{R} \to x \neq 0 \to y \neq 0 \to x \in \mathbb{R}) \to \varphi_{167}$}
\BIC{$\varphi_{167}$}
\end{ptree}
\noindent$\Pi_{47}$:
\begin{ptree}
\AXC{$x \in \mathbb{R} \to y \in \mathbb{R} \to x \neq 0 \to y \neq 0 \to x \neq 0$\lbl{A1}}
\AXC{$\stackrel{\Pi_{46}}{\vdots}$}\noLine
\UIC{$\varphi_{167}$}
\BIC{$x \in \mathbb{R} \to y \in \mathbb{R} \to x \neq 0 \to y \neq 0 \to x \cdot x^{-1} = 1$}
\AXC{$(x \in \mathbb{R} \to y \in \mathbb{R} \to x \neq 0 \to y \neq 0 \to x \cdot x^{-1} = 1) \to \varphi_{217}$\lbl{A1}}
\BIC{$\varphi_{217}$}
\end{ptree}
\noindent$\Pi_{48}$:
\begin{ptree}
\AXC{$\stackrel{\Pi_{45}}{\vdots}$}\noLine
\UIC{$x \cdot x^{-1} = 1 \to \varphi_{127}$}
\AXC{$x = x$\lbl{A5}}
\UIC{$\forall x\,(x = x)$}
\AXC{$\forall x\,(x = x) \to x \cdot x^{-1} \cdot y^{-1} = x \cdot x^{-1} \cdot y^{-1}$\lbl{A2}}
\BIC{$x \cdot x^{-1} \cdot y^{-1} = x \cdot x^{-1} \cdot y^{-1}$}
\AXC{$\stackrel{\Pi_{47}}{\vdots}$}\noLine
\UIC{$\varphi_{217}$}
\BIC{$(x \cdot x^{-1} = 1 \to \varphi_{127}) \to \varphi_{102}$}
\BIC{$\varphi_{102}$}
\end{ptree}
\noindent$\Pi_{49}$:
\begin{ptree}
\AXC{$\stackrel{\Pi_{34}}{\vdots}$}\noLine
\UIC{$\forall x\,\forall y\,(x = y \to y = x)$}
\AXC{$\forall x\,\forall y\,(x = y \to y = x) \to \forall y\,(v_{572} = y \to y = v_{572})$\lbl{A2}}
\BIC{$\forall y\,(v_{572} = y \to y = v_{572})$}
\AXC{$\varphi_{2}$\lbl{A2}}
\BIC{$v_{572} = v_{573} \to v_{573} = v_{572}$}
\UIC{$\forall v_{573}\,(v_{572} = v_{573} \to v_{573} = v_{572})$}
\UIC{$\forall v_{572}\,\forall v_{573}\,(v_{572} = v_{573} \to v_{573} = v_{572})$}
\end{ptree}
\noindent$\Pi_{50}$:
\begin{ptree}
\AXC{$\stackrel{\Pi_{49}}{\vdots}$}\noLine
\UIC{$\forall v_{572}\,\forall v_{573}\,(v_{572} = v_{573} \to v_{573} = v_{572})$}
\AXC{$\forall v_{572}\,\forall v_{573}\,(v_{572} = v_{573} \to v_{573} = v_{572}) \to \varphi_{50}$\lbl{A2}}
\BIC{$\varphi_{50}$}
\AXC{$\varphi_{50} \to \varphi_{141}$\lbl{A2}}
\BIC{$\varphi_{141}$}
\end{ptree}
\noindent$\Pi_{51}$:
\begin{ptree}
\AXC{$\forall x\,(\varphi_{73})$\lbl{(R6)}}
\AXC{$\forall x\,(\varphi_{73}) \to \varphi_{73}$\lbl{A2}}
\BIC{$\varphi_{73}$}
\AXC{$\varphi_{73} \to \varphi_{131}$\lbl{A2}}
\BIC{$\varphi_{131}$}
\AXC{$\varphi_{131} \to x \in \mathbb{R} \to \varphi_{128}$\lbl{A2}}
\BIC{$x \in \mathbb{R} \to \varphi_{128}$}
\AXC{$(x \in \mathbb{R} \to \varphi_{128}) \to \varphi_{228}$\lbl{A1}}
\BIC{$\varphi_{228}$}
\end{ptree}
\noindent$\Pi_{52}$:
\begin{ptree}
\AXC{$\stackrel{\Pi_{8}}{\vdots}$}\noLine
\UIC{$x \in \mathbb{R} \to y \in \mathbb{R} \to x \neq 0 \to y \neq 0 \to x^{-1} \in \mathbb{R}$}
\AXC{$x \in \mathbb{R} \to y \in \mathbb{R} \to x \neq 0 \to y \neq 0 \to x \in \mathbb{R}$\lbl{A1}}
\AXC{$\stackrel{\Pi_{51}}{\vdots}$}\noLine
\UIC{$\varphi_{228}$}
\BIC{$(x \in \mathbb{R} \to y \in \mathbb{R} \to x \neq 0 \to y \neq 0 \to x^{-1} \in \mathbb{R}) \to \varphi_{193}$}
\BIC{$\varphi_{193}$}
\end{ptree}
\noindent$\Pi_{53}$:
\begin{ptree}
\AXC{$\stackrel{\Pi_{50}}{\vdots}$}\noLine
\UIC{$\varphi_{141}$}
\AXC{$\stackrel{\Pi_{6}}{\vdots}$}\noLine
\UIC{$x \in \mathbb{R} \to y \in \mathbb{R} \to x \neq 0 \to y \neq 0 \to y^{-1} \in \mathbb{R}$}
\AXC{$\stackrel{\Pi_{52}}{\vdots}$}\noLine
\UIC{$\varphi_{193}$}
\BIC{$x \in \mathbb{R} \to \varphi_{89}$}
\AXC{$(x \in \mathbb{R} \to \varphi_{89}) \to \varphi_{141} \to x \in \mathbb{R} \to \varphi_{90}$\lbl{A1}}
\BIC{$\varphi_{141} \to x \in \mathbb{R} \to \varphi_{90}$}
\BIC{$x \in \mathbb{R} \to \varphi_{90}$}
\AXC{$\varphi_{237}$\lbl{A1}}
\BIC{$\varphi_{102} \to (x \cdot (x^{-1} \cdot y^{-1}) = x \cdot x^{-1} \cdot y^{-1} \to \varphi_{83}) \to \varphi_{103}$}
\end{ptree}
\noindent$\Pi_{54}$:
\begin{ptree}
\AXC{$\stackrel{\Pi_{44}}{\vdots}$}\noLine
\UIC{$x \cdot (x^{-1} \cdot y^{-1}) = x \cdot x^{-1} \cdot y^{-1} \to \varphi_{83}$}
\AXC{$\stackrel{\Pi_{48}}{\vdots}$}\noLine
\UIC{$\varphi_{102}$}
\AXC{$\stackrel{\Pi_{53}}{\vdots}$}\noLine
\UIC{$\varphi_{102} \to (x \cdot (x^{-1} \cdot y^{-1}) = x \cdot x^{-1} \cdot y^{-1} \to \varphi_{83}) \to \varphi_{103}$}
\BIC{$(x \cdot (x^{-1} \cdot y^{-1}) = x \cdot x^{-1} \cdot y^{-1} \to \varphi_{83}) \to \varphi_{103}$}
\BIC{$\varphi_{103}$}
\end{ptree}
\noindent$\Pi_{55}$:
\begin{ptree}
\AXC{$\stackrel{\Pi_{6}}{\vdots}$}\noLine
\UIC{$x \in \mathbb{R} \to y \in \mathbb{R} \to x \neq 0 \to y \neq 0 \to y^{-1} \in \mathbb{R}$}
\AXC{$\stackrel{\Pi_{42}}{\vdots}$}\noLine
\UIC{$\varphi_{186}$}
\BIC{$x \in \mathbb{R} \to y \in \mathbb{R} \to x \neq 0 \to y \neq 0 \to 1 \cdot y^{-1} = y^{-1} \cdot 1$}
\AXC{$\stackrel{\Pi_{54}}{\vdots}$}\noLine
\UIC{$\varphi_{103}$}
\AXC{$\varphi_{103} \to \varphi_{225}$\lbl{A1}}
\BIC{$\varphi_{225}$}
\BIC{$\varphi_{154} \to \varphi_{104}$}
\end{ptree}
\noindent$\Pi_{56}$:
\begin{ptree}
\AXC{$\stackrel{\Pi_{41}}{\vdots}$}\noLine
\UIC{$\forall v_{577}\,(\varphi_{11})$}
\AXC{$\forall v_{577}\,(\varphi_{11}) \to \varphi_{121}$\lbl{A2}}
\BIC{$\varphi_{121}$}
\AXC{$\varphi_{121} \to \varphi_{138}$\lbl{A2}}
\BIC{$\varphi_{138}$}
\AXC{$\varphi_{138} \to \varphi_{154}$\lbl{A2}}
\BIC{$\varphi_{154}$}
\AXC{$\stackrel{\Pi_{55}}{\vdots}$}\noLine
\UIC{$\varphi_{154} \to \varphi_{104}$}
\BIC{$\varphi_{104}$}
\AXC{$\varphi_{104} \to \varphi_{218}$\lbl{A1}}
\BIC{$\varphi_{218}$}
\end{ptree}
\noindent$\Pi_{57}$:
\begin{ptree}
\AXC{$\stackrel{\Pi_{36}}{\vdots}$}\noLine
\UIC{$\varphi_{33}$}
\AXC{$\stackrel{\Pi_{38}}{\vdots}$}\noLine
\UIC{$\varphi_{135}$}
\AXC{$\stackrel{\Pi_{40}}{\vdots}$}\noLine
\UIC{$x \in \mathbb{R} \to y \in \mathbb{R} \to x \neq 0 \to y \neq 0 \to y^{-1} \cdot 1 = y^{-1}$}
\AXC{$\stackrel{\Pi_{56}}{\vdots}$}\noLine
\UIC{$\varphi_{218}$}
\BIC{$\varphi_{135} \to \varphi_{74}$}
\BIC{$\varphi_{74}$}
\AXC{$\varphi_{74} \to \varphi_{33} \to \varphi_{75}$\lbl{A1}}
\BIC{$\varphi_{33} \to \varphi_{75}$}
\BIC{$\varphi_{75}$}
\AXC{$\varphi_{75} \to \varphi_{56} \to \varphi_{235}$\lbl{A1}}
\BIC{$\varphi_{56} \to \varphi_{235}$}
\end{ptree}
\noindent$\Pi_{58}$:
\begin{ptree}
\AXC{$\stackrel{\Pi_{33}}{\vdots}$}\noLine
\UIC{$y^{-1} = x \cdot (x^{-1} \cdot y^{-1}) \to \varphi_{56} \to \varphi_{119}$}
\AXC{$x = x$\lbl{A5}}
\UIC{$\forall x\,(x = x)$}
\AXC{$\forall x\,(x = x) \to \varphi_{56}$\lbl{A2}}
\BIC{$\varphi_{56}$}
\AXC{$\stackrel{\Pi_{57}}{\vdots}$}\noLine
\UIC{$\varphi_{56} \to \varphi_{235}$}
\BIC{$\varphi_{235}$}
\BIC{$x \in \mathbb{R} \to y \in \mathbb{R} \to x \neq 0 \to y \neq 0 \to \varphi_{119}$}
\end{ptree}
\noindent$\Pi_{59}$:
\begin{ptree}
\AXC{$\stackrel{\Pi_{34}}{\vdots}$}\noLine
\UIC{$\forall x\,\forall y\,(x = y \to y = x)$}
\AXC{$\forall x\,\forall y\,(x = y \to y = x) \to \forall y\,(v_{583} = y \to y = v_{583})$\lbl{A2}}
\BIC{$\forall y\,(v_{583} = y \to y = v_{583})$}
\AXC{$\varphi_{3}$\lbl{A2}}
\BIC{$v_{583} = v_{584} \to v_{584} = v_{583}$}
\UIC{$\forall v_{584}\,(v_{583} = v_{584} \to v_{584} = v_{583})$}
\UIC{$\forall v_{583}\,\forall v_{584}\,(v_{583} = v_{584} \to v_{584} = v_{583})$}
\end{ptree}
\noindent$\Pi_{60}$:
\begin{ptree}
\AXC{$\stackrel{\Pi_{59}}{\vdots}$}\noLine
\UIC{$\forall v_{583}\,\forall v_{584}\,(v_{583} = v_{584} \to v_{584} = v_{583})$}
\AXC{$\forall v_{583}\,\forall v_{584}\,(v_{583} = v_{584} \to v_{584} = v_{583}) \to \varphi_{51}$\lbl{A2}}
\BIC{$\varphi_{51}$}
\AXC{$\varphi_{51} \to \varphi_{34}$\lbl{A2}}
\BIC{$\varphi_{34}$}
\end{ptree}
\noindent$\Pi_{61}$:
\begin{ptree}
\AXC{$y = z \to x = y \to x = z$\lbl{A6}}
\AXC{$(y = z \to x = y \to x = z) \to x = y \to y = z \to x = z$\lbl{A1}}
\BIC{$x = y \to y = z \to x = z$}
\UIC{$\forall z\,(x = y \to y = z \to x = z)$}
\UIC{$\forall y\,\forall z\,(x = y \to y = z \to x = z)$}
\UIC{$\forall x\,\forall y\,\forall z\,(x = y \to y = z \to x = z)$}
\AXC{$\varphi_{26}$\lbl{A2}}
\BIC{$\forall y\,\forall z\,(v_{569} = y \to y = z \to v_{569} = z)$}
\AXC{$\varphi_{44}$\lbl{A2}}
\BIC{$\forall z\,(v_{569} = v_{570} \to v_{570} = z \to v_{569} = z)$}
\AXC{$\varphi_{69}$\lbl{A2}}
\BIC{$v_{569} = v_{570} \to v_{570} = v_{571} \to v_{569} = v_{571}$}
\UIC{$\forall v_{571}\,(v_{569} = v_{570} \to v_{570} = v_{571} \to v_{569} = v_{571})$}
\UIC{$\varphi_{13}$}
\UIC{$\forall v_{569}\,(\varphi_{13})$}
\end{ptree}
\noindent$\Pi_{62}$:
\begin{ptree}
\AXC{$\stackrel{\Pi_{61}}{\vdots}$}\noLine
\UIC{$\forall v_{569}\,(\varphi_{13})$}
\AXC{$\forall v_{569}\,(\varphi_{13}) \to \varphi_{123}$\lbl{A2}}
\BIC{$\varphi_{123}$}
\AXC{$\varphi_{123} \to \varphi_{139}$\lbl{A2}}
\BIC{$\varphi_{139}$}
\AXC{$\varphi_{139} \to \varphi_{136}$\lbl{A2}}
\BIC{$\varphi_{136}$}
\end{ptree}
\noindent$\Pi_{63}$:
\begin{ptree}
\AXC{$\forall x \in \mathbb{R}\,(x \cdot 1 = x)$\lbl{(R8)}}
\AXC{$\forall x \in \mathbb{R}\,(x \cdot 1 = x) \to x^{-1} \in \mathbb{R} \to x^{-1} \cdot 1 = x^{-1}$\lbl{A2}}
\BIC{$x^{-1} \in \mathbb{R} \to x^{-1} \cdot 1 = x^{-1}$}
\end{ptree}
\noindent$\Pi_{64}$:
\begin{ptree}
\AXC{$\stackrel{\Pi_{8}}{\vdots}$}\noLine
\UIC{$x \in \mathbb{R} \to y \in \mathbb{R} \to x \neq 0 \to y \neq 0 \to x^{-1} \in \mathbb{R}$}
\AXC{$\stackrel{\Pi_{63}}{\vdots}$}\noLine
\UIC{$x^{-1} \in \mathbb{R} \to x^{-1} \cdot 1 = x^{-1}$}
\AXC{$(x^{-1} \in \mathbb{R} \to x^{-1} \cdot 1 = x^{-1}) \to \varphi_{182}$\lbl{A1}}
\BIC{$\varphi_{182}$}
\BIC{$x \in \mathbb{R} \to y \in \mathbb{R} \to x \neq 0 \to y \neq 0 \to x^{-1} \cdot 1 = x^{-1}$}
\end{ptree}
\noindent$\Pi_{65}$:
\begin{ptree}
\AXC{$y = z \to x = y \to x = z$\lbl{A6}}
\AXC{$(y = z \to x = y \to x = z) \to x = y \to y = z \to x = z$\lbl{A1}}
\BIC{$x = y \to y = z \to x = z$}
\UIC{$\forall z\,(x = y \to y = z \to x = z)$}
\UIC{$\forall y\,\forall z\,(x = y \to y = z \to x = z)$}
\UIC{$\forall x\,\forall y\,\forall z\,(x = y \to y = z \to x = z)$}
\AXC{$\varphi_{27}$\lbl{A2}}
\BIC{$\forall y\,\forall z\,(v_{566} = y \to y = z \to v_{566} = z)$}
\AXC{$\varphi_{45}$\lbl{A2}}
\BIC{$\forall z\,(v_{566} = v_{567} \to v_{567} = z \to v_{566} = z)$}
\AXC{$\varphi_{70}$\lbl{A2}}
\BIC{$v_{566} = v_{567} \to v_{567} = v_{568} \to v_{566} = v_{568}$}
\UIC{$\forall v_{568}\,(v_{566} = v_{567} \to v_{567} = v_{568} \to v_{566} = v_{568})$}
\UIC{$\varphi_{14}$}
\UIC{$\forall v_{566}\,(\varphi_{14})$}
\end{ptree}
\noindent$\Pi_{66}$:
\begin{ptree}
\AXC{$\stackrel{\Pi_{65}}{\vdots}$}\noLine
\UIC{$\forall v_{566}\,(\varphi_{14})$}
\AXC{$\forall v_{566}\,(\varphi_{14}) \to \varphi_{124}$\lbl{A2}}
\BIC{$\varphi_{124}$}
\AXC{$\varphi_{124} \to \varphi_{140}$\lbl{A2}}
\BIC{$\varphi_{140}$}
\AXC{$\varphi_{140} \to \varphi_{155}$\lbl{A2}}
\BIC{$\varphi_{155}$}
\end{ptree}
\noindent$\Pi_{67}$:
\begin{ptree}
\AXC{$1 \in \mathbb{R}$\lbl{(R1)}}
\AXC{$\forall x\,\forall y\,(x \in \mathbb{R} \to y \in \mathbb{R} \to x \cdot y = y \cdot x)$\lbl{(R7)}}
\AXC{$\varphi_{146}$\lbl{A2}}
\BIC{$\forall y\,(1 \in \mathbb{R} \to y \in \mathbb{R} \to 1 \cdot y = y \cdot 1)$}
\AXC{$\varphi_{164}$\lbl{A2}}
\BIC{$1 \in \mathbb{R} \to x^{-1} \in \mathbb{R} \to 1 \cdot x^{-1} = x^{-1} \cdot 1$}
\AXC{$\varphi_{210}$\lbl{A1}}
\BIC{$1 \in \mathbb{R} \to \varphi_{187}$}
\BIC{$\varphi_{187}$}
\end{ptree}
\noindent$\Pi_{68}$:
\begin{ptree}
\AXC{$y = z \to x = y \to x = z$\lbl{A6}}
\AXC{$(y = z \to x = y \to x = z) \to x = y \to y = z \to x = z$\lbl{A1}}
\BIC{$x = y \to y = z \to x = z$}
\UIC{$\forall z\,(x = y \to y = z \to x = z)$}
\UIC{$\forall y\,\forall z\,(x = y \to y = z \to x = z)$}
\UIC{$\forall x\,\forall y\,\forall z\,(x = y \to y = z \to x = z)$}
\AXC{$\varphi_{28}$\lbl{A2}}
\BIC{$\forall y\,\forall z\,(v_{563} = y \to y = z \to v_{563} = z)$}
\AXC{$\varphi_{46}$\lbl{A2}}
\BIC{$\forall z\,(v_{563} = v_{564} \to v_{564} = z \to v_{563} = z)$}
\AXC{$\varphi_{71}$\lbl{A2}}
\BIC{$v_{563} = v_{564} \to v_{564} = v_{565} \to v_{563} = v_{565}$}
\UIC{$\forall v_{565}\,(v_{563} = v_{564} \to v_{564} = v_{565} \to v_{563} = v_{565})$}
\UIC{$\varphi_{15}$}
\UIC{$\forall v_{563}\,(\varphi_{15})$}
\end{ptree}
\noindent$\Pi_{69}$:
\begin{ptree}
\AXC{$\stackrel{\Pi_{68}}{\vdots}$}\noLine
\UIC{$\forall v_{563}\,(\varphi_{15})$}
\AXC{$\forall v_{563}\,(\varphi_{15}) \to \varphi_{125}$\lbl{A2}}
\BIC{$\varphi_{125}$}
\AXC{$\varphi_{125} \to \varphi_{176}$\lbl{A2}}
\BIC{$\varphi_{176}$}
\AXC{$\varphi_{176} \to y \cdot (x^{-1} \cdot y^{-1}) = y \cdot y^{-1} \cdot x^{-1} \to \varphi_{85}$\lbl{A2}}
\BIC{$y \cdot (x^{-1} \cdot y^{-1}) = y \cdot y^{-1} \cdot x^{-1} \to \varphi_{85}$}
\end{ptree}
\noindent$\Pi_{70}$:
\begin{ptree}
\AXC{$z = w \to \varphi_{86}$\lbl{A6}}
\UIC{$\forall w\,(z = w \to \varphi_{86})$}
\UIC{$\forall z\,\forall w\,(z = w \to \varphi_{86})$}
\AXC{$\forall z\,\forall w\,(z = w \to \varphi_{86}) \to \varphi_{158}$\lbl{A2}}
\BIC{$\varphi_{158}$}
\AXC{$\varphi_{158} \to y \cdot y^{-1} = 1 \to \varphi_{129}$\lbl{A2}}
\BIC{$y \cdot y^{-1} = 1 \to \varphi_{129}$}
\end{ptree}
\noindent$\Pi_{71}$:
\begin{ptree}
\AXC{$x \in \mathbb{R} \to y \in \mathbb{R} \to x \neq 0 \to y \neq 0 \to y \in \mathbb{R}$\lbl{A1}}
\AXC{$\forall x \in \mathbb{R}\,(x \neq 0 \to x \cdot x^{-1} = 1)$\lbl{(R9)}}
\AXC{$\varphi_{101}$\lbl{A2}}
\BIC{$y \in \mathbb{R} \to y \neq 0 \to y \cdot y^{-1} = 1$}
\AXC{$\varphi_{207}$\lbl{A1}}
\BIC{$(x \in \mathbb{R} \to y \in \mathbb{R} \to x \neq 0 \to y \neq 0 \to y \in \mathbb{R}) \to \varphi_{168}$}
\BIC{$\varphi_{168}$}
\end{ptree}
\noindent$\Pi_{72}$:
\begin{ptree}
\AXC{$x \in \mathbb{R} \to y \in \mathbb{R} \to x \neq 0 \to y \neq 0 \to y \neq 0$\lbl{A1}}
\AXC{$\stackrel{\Pi_{71}}{\vdots}$}\noLine
\UIC{$\varphi_{168}$}
\BIC{$x \in \mathbb{R} \to y \in \mathbb{R} \to x \neq 0 \to y \neq 0 \to y \cdot y^{-1} = 1$}
\AXC{$(x \in \mathbb{R} \to y \in \mathbb{R} \to x \neq 0 \to y \neq 0 \to y \cdot y^{-1} = 1) \to \varphi_{219}$\lbl{A1}}
\BIC{$\varphi_{219}$}
\end{ptree}
\noindent$\Pi_{73}$:
\begin{ptree}
\AXC{$\stackrel{\Pi_{70}}{\vdots}$}\noLine
\UIC{$y \cdot y^{-1} = 1 \to \varphi_{129}$}
\AXC{$x = x$\lbl{A5}}
\UIC{$\forall x\,(x = x)$}
\AXC{$\forall x\,(x = x) \to y \cdot y^{-1} \cdot x^{-1} = y \cdot y^{-1} \cdot x^{-1}$\lbl{A2}}
\BIC{$y \cdot y^{-1} \cdot x^{-1} = y \cdot y^{-1} \cdot x^{-1}$}
\AXC{$\stackrel{\Pi_{72}}{\vdots}$}\noLine
\UIC{$\varphi_{219}$}
\BIC{$(y \cdot y^{-1} = 1 \to \varphi_{129}) \to \varphi_{105}$}
\BIC{$\varphi_{105}$}
\end{ptree}
\noindent$\Pi_{74}$:
\begin{ptree}
\AXC{$y = z \to x = y \to x = z$\lbl{A6}}
\AXC{$(y = z \to x = y \to x = z) \to x = y \to y = z \to x = z$\lbl{A1}}
\BIC{$x = y \to y = z \to x = z$}
\UIC{$\forall z\,(x = y \to y = z \to x = z)$}
\UIC{$\forall y\,\forall z\,(x = y \to y = z \to x = z)$}
\UIC{$\forall x\,\forall y\,\forall z\,(x = y \to y = z \to x = z)$}
\AXC{$\varphi_{29}$\lbl{A2}}
\BIC{$\forall y\,\forall z\,(v_{560} = y \to y = z \to v_{560} = z)$}
\AXC{$\varphi_{47}$\lbl{A2}}
\BIC{$\forall z\,(v_{560} = v_{561} \to v_{561} = z \to v_{560} = z)$}
\AXC{$\varphi_{72}$\lbl{A2}}
\BIC{$v_{560} = v_{561} \to v_{561} = v_{562} \to v_{560} = v_{562}$}
\UIC{$\forall v_{562}\,(v_{560} = v_{561} \to v_{561} = v_{562} \to v_{560} = v_{562})$}
\UIC{$\varphi_{16}$}
\UIC{$\forall v_{560}\,(\varphi_{16})$}
\end{ptree}
\noindent$\Pi_{75}$:
\begin{ptree}
\AXC{$\stackrel{\Pi_{74}}{\vdots}$}\noLine
\UIC{$\forall v_{560}\,(\varphi_{16})$}
\AXC{$\forall v_{560}\,(\varphi_{16}) \to \varphi_{126}$\lbl{A2}}
\BIC{$\varphi_{126}$}
\AXC{$\varphi_{126} \to \varphi_{177}$\lbl{A2}}
\BIC{$\varphi_{177}$}
\AXC{$\varphi_{177} \to y \cdot (x^{-1} \cdot y^{-1}) = y \cdot (y^{-1} \cdot x^{-1}) \to \varphi_{143}$\lbl{A2}}
\BIC{$y \cdot (x^{-1} \cdot y^{-1}) = y \cdot (y^{-1} \cdot x^{-1}) \to \varphi_{143}$}
\end{ptree}
\noindent$\Pi_{76}$:
\begin{ptree}
\AXC{$\stackrel{\Pi_{34}}{\vdots}$}\noLine
\UIC{$\forall x\,\forall y\,(x = y \to y = x)$}
\AXC{$\forall x\,\forall y\,(x = y \to y = x) \to \forall y\,(v_{558} = y \to y = v_{558})$\lbl{A2}}
\BIC{$\forall y\,(v_{558} = y \to y = v_{558})$}
\AXC{$\varphi_{4}$\lbl{A2}}
\BIC{$v_{558} = v_{559} \to v_{559} = v_{558}$}
\UIC{$\forall v_{559}\,(v_{558} = v_{559} \to v_{559} = v_{558})$}
\UIC{$\forall v_{558}\,\forall v_{559}\,(v_{558} = v_{559} \to v_{559} = v_{558})$}
\end{ptree}
\noindent$\Pi_{77}$:
\begin{ptree}
\AXC{$\stackrel{\Pi_{76}}{\vdots}$}\noLine
\UIC{$\forall v_{558}\,\forall v_{559}\,(v_{558} = v_{559} \to v_{559} = v_{558})$}
\AXC{$\forall v_{558}\,\forall v_{559}\,(v_{558} = v_{559} \to v_{559} = v_{558}) \to \varphi_{52}$\lbl{A2}}
\BIC{$\varphi_{52}$}
\AXC{$\varphi_{52} \to \varphi_{142}$\lbl{A2}}
\BIC{$\varphi_{142}$}
\end{ptree}
\noindent$\Pi_{78}$:
\begin{ptree}
\AXC{$\forall x\,(\varphi_{73})$\lbl{(R6)}}
\AXC{$\forall x\,(\varphi_{73}) \to \varphi_{117}$\lbl{A2}}
\BIC{$\varphi_{117}$}
\AXC{$\varphi_{117} \to \varphi_{132}$\lbl{A2}}
\BIC{$\varphi_{132}$}
\AXC{$\varphi_{132} \to v_{555} \in \mathbb{R} \to v_{556} \in \mathbb{R} \to \varphi_{30}$\lbl{A2}}
\BIC{$v_{555} \in \mathbb{R} \to v_{556} \in \mathbb{R} \to \varphi_{30}$}
\UIC{$\forall v_{557}\,(v_{555} \in \mathbb{R} \to v_{556} \in \mathbb{R} \to \varphi_{30})$}
\UIC{$\forall v_{556}\,\forall v_{557}\,(v_{555} \in \mathbb{R} \to v_{556} \in \mathbb{R} \to \varphi_{30})$}
\UIC{$\forall v_{555}\,\forall v_{556}\,\forall v_{557}\,(v_{555} \in \mathbb{R} \to v_{556} \in \mathbb{R} \to \varphi_{30})$}
\AXC{$\varphi_{213}$\lbl{A2}}
\BIC{$\varphi_{148}$}
\AXC{$\varphi_{148} \to \varphi_{149}$\lbl{A2}}
\BIC{$\varphi_{149}$}
\end{ptree}
\noindent$\Pi_{79}$:
\begin{ptree}
\AXC{$x \in \mathbb{R} \to y \in \mathbb{R} \to x \neq 0 \to y \neq 0 \to y \in \mathbb{R}$\lbl{A1}}
\AXC{$\stackrel{\Pi_{78}}{\vdots}$}\noLine
\UIC{$\varphi_{149}$}
\AXC{$\varphi_{149} \to y \in \mathbb{R} \to \varphi_{130}$\lbl{A2}}
\BIC{$y \in \mathbb{R} \to \varphi_{130}$}
\AXC{$(y \in \mathbb{R} \to \varphi_{130}) \to \varphi_{229}$\lbl{A1}}
\BIC{$\varphi_{229}$}
\BIC{$(x \in \mathbb{R} \to y \in \mathbb{R} \to x \neq 0 \to y \neq 0 \to y^{-1} \in \mathbb{R}) \to \varphi_{194}$}
\end{ptree}
\noindent$\Pi_{80}$:
\begin{ptree}
\AXC{$\stackrel{\Pi_{6}}{\vdots}$}\noLine
\UIC{$x \in \mathbb{R} \to y \in \mathbb{R} \to x \neq 0 \to y \neq 0 \to y^{-1} \in \mathbb{R}$}
\AXC{$\stackrel{\Pi_{79}}{\vdots}$}\noLine
\UIC{$(x \in \mathbb{R} \to y \in \mathbb{R} \to x \neq 0 \to y \neq 0 \to y^{-1} \in \mathbb{R}) \to \varphi_{194}$}
\BIC{$\varphi_{194}$}
\end{ptree}
\noindent$\Pi_{81}$:
\begin{ptree}
\AXC{$\stackrel{\Pi_{77}}{\vdots}$}\noLine
\UIC{$\varphi_{142}$}
\AXC{$\stackrel{\Pi_{8}}{\vdots}$}\noLine
\UIC{$x \in \mathbb{R} \to y \in \mathbb{R} \to x \neq 0 \to y \neq 0 \to x^{-1} \in \mathbb{R}$}
\AXC{$\stackrel{\Pi_{80}}{\vdots}$}\noLine
\UIC{$\varphi_{194}$}
\BIC{$x \in \mathbb{R} \to \varphi_{91}$}
\AXC{$(x \in \mathbb{R} \to \varphi_{91}) \to \varphi_{142} \to x \in \mathbb{R} \to \varphi_{92}$\lbl{A1}}
\BIC{$\varphi_{142} \to x \in \mathbb{R} \to \varphi_{92}$}
\BIC{$x \in \mathbb{R} \to \varphi_{92}$}
\end{ptree}
\noindent$\Pi_{82}$:
\begin{ptree}
\AXC{$\forall x\,\forall y\,(x \in \mathbb{R} \to y \in \mathbb{R} \to x \cdot y = y \cdot x)$\lbl{(R7)}}
\AXC{$\varphi_{172}$\lbl{A2}}
\BIC{$\forall y\,(x^{-1} \in \mathbb{R} \to y \in \mathbb{R} \to x^{-1} \cdot y = y \cdot x^{-1})$}
\AXC{$\varphi_{190}$\lbl{A2}}
\BIC{$x^{-1} \in \mathbb{R} \to y^{-1} \in \mathbb{R} \to x^{-1} \cdot y^{-1} = y^{-1} \cdot x^{-1}$}
\AXC{$\varphi_{230}$\lbl{A1}}
\BIC{$\varphi_{208}$}
\end{ptree}
\noindent$\Pi_{83}$:
\begin{ptree}
\AXC{$\stackrel{\Pi_{6}}{\vdots}$}\noLine
\UIC{$x \in \mathbb{R} \to y \in \mathbb{R} \to x \neq 0 \to y \neq 0 \to y^{-1} \in \mathbb{R}$}
\AXC{$\stackrel{\Pi_{8}}{\vdots}$}\noLine
\UIC{$x \in \mathbb{R} \to y \in \mathbb{R} \to x \neq 0 \to y \neq 0 \to x^{-1} \in \mathbb{R}$}
\AXC{$\stackrel{\Pi_{82}}{\vdots}$}\noLine
\UIC{$\varphi_{208}$}
\BIC{$(x \in \mathbb{R} \to y \in \mathbb{R} \to x \neq 0 \to y \neq 0 \to y^{-1} \in \mathbb{R}) \to \varphi_{106}$}
\BIC{$\varphi_{106}$}
\end{ptree}
\noindent$\Pi_{84}$:
\begin{ptree}
\AXC{$x = x$\lbl{A5}}
\UIC{$\forall x\,(x = x)$}
\AXC{$\varphi_{17}$\lbl{A2}}
\BIC{$y \cdot (x^{-1} \cdot y^{-1}) = y \cdot (x^{-1} \cdot y^{-1})$}
\AXC{$\stackrel{\Pi_{83}}{\vdots}$}\noLine
\UIC{$\varphi_{106}$}
\AXC{$\varphi_{106} \to y \cdot (x^{-1} \cdot y^{-1}) = y \cdot (x^{-1} \cdot y^{-1}) \to \varphi_{216}$\lbl{A1}}
\BIC{$y \cdot (x^{-1} \cdot y^{-1}) = y \cdot (x^{-1} \cdot y^{-1}) \to \varphi_{216}$}
\BIC{$\varphi_{216}$}
\end{ptree}
\noindent$\Pi_{85}$:
\begin{ptree}
\AXC{$z = w \to \varphi_{35}$\lbl{A6}}
\UIC{$\forall w\,(z = w \to \varphi_{35})$}
\UIC{$\forall z\,\forall w\,(z = w \to \varphi_{35})$}
\AXC{$\forall z\,\forall w\,(z = w \to \varphi_{35}) \to \varphi_{159}$\lbl{A2}}
\BIC{$\varphi_{159}$}
\AXC{$\varphi_{159} \to x^{-1} \cdot y^{-1} = y^{-1} \cdot x^{-1} \to \varphi_{144}$\lbl{A2}}
\BIC{$x^{-1} \cdot y^{-1} = y^{-1} \cdot x^{-1} \to \varphi_{144}$}
\AXC{$\stackrel{\Pi_{84}}{\vdots}$}\noLine
\UIC{$\varphi_{216}$}
\BIC{$x \in \mathbb{R} \to y \in \mathbb{R} \to \varphi_{31}$}
\end{ptree}
\noindent$\Pi_{86}$:
\begin{ptree}
\AXC{$\stackrel{\Pi_{81}}{\vdots}$}\noLine
\UIC{$x \in \mathbb{R} \to \varphi_{92}$}
\AXC{$\stackrel{\Pi_{85}}{\vdots}$}\noLine
\UIC{$x \in \mathbb{R} \to y \in \mathbb{R} \to \varphi_{31}$}
\AXC{$(x \in \mathbb{R} \to y \in \mathbb{R} \to \varphi_{31}) \to (x \in \mathbb{R} \to \varphi_{92}) \to \varphi_{222}$\lbl{A1}}
\BIC{$(x \in \mathbb{R} \to \varphi_{92}) \to \varphi_{222}$}
\BIC{$\varphi_{222}$}
\end{ptree}
\noindent$\Pi_{87}$:
\begin{ptree}
\AXC{$\stackrel{\Pi_{73}}{\vdots}$}\noLine
\UIC{$\varphi_{105}$}
\AXC{$\stackrel{\Pi_{75}}{\vdots}$}\noLine
\UIC{$y \cdot (x^{-1} \cdot y^{-1}) = y \cdot (y^{-1} \cdot x^{-1}) \to \varphi_{143}$}
\AXC{$\stackrel{\Pi_{86}}{\vdots}$}\noLine
\UIC{$\varphi_{222}$}
\BIC{$x \in \mathbb{R} \to \varphi_{93}$}
\AXC{$\varphi_{238}$\lbl{A1}}
\BIC{$\varphi_{105} \to (y \cdot (x^{-1} \cdot y^{-1}) = y \cdot y^{-1} \cdot x^{-1} \to \varphi_{85}) \to \varphi_{107}$}
\BIC{$(y \cdot (x^{-1} \cdot y^{-1}) = y \cdot y^{-1} \cdot x^{-1} \to \varphi_{85}) \to \varphi_{107}$}
\end{ptree}
\noindent$\Pi_{88}$:
\begin{ptree}
\AXC{$\stackrel{\Pi_{69}}{\vdots}$}\noLine
\UIC{$y \cdot (x^{-1} \cdot y^{-1}) = y \cdot y^{-1} \cdot x^{-1} \to \varphi_{85}$}
\AXC{$\stackrel{\Pi_{87}}{\vdots}$}\noLine
\UIC{$(y \cdot (x^{-1} \cdot y^{-1}) = y \cdot y^{-1} \cdot x^{-1} \to \varphi_{85}) \to \varphi_{107}$}
\BIC{$\varphi_{107}$}
\AXC{$\varphi_{107} \to \varphi_{226}$\lbl{A1}}
\BIC{$\varphi_{226}$}
\end{ptree}
\noindent$\Pi_{89}$:
\begin{ptree}
\AXC{$\stackrel{\Pi_{66}}{\vdots}$}\noLine
\UIC{$\varphi_{155}$}
\AXC{$\stackrel{\Pi_{8}}{\vdots}$}\noLine
\UIC{$x \in \mathbb{R} \to y \in \mathbb{R} \to x \neq 0 \to y \neq 0 \to x^{-1} \in \mathbb{R}$}
\AXC{$\stackrel{\Pi_{67}}{\vdots}$}\noLine
\UIC{$\varphi_{187}$}
\BIC{$x \in \mathbb{R} \to y \in \mathbb{R} \to x \neq 0 \to y \neq 0 \to 1 \cdot x^{-1} = x^{-1} \cdot 1$}
\AXC{$\stackrel{\Pi_{88}}{\vdots}$}\noLine
\UIC{$\varphi_{226}$}
\BIC{$\varphi_{155} \to \varphi_{108}$}
\BIC{$\varphi_{108}$}
\AXC{$\varphi_{108} \to \varphi_{220}$\lbl{A1}}
\BIC{$\varphi_{220}$}
\end{ptree}
\noindent$\Pi_{90}$:
\begin{ptree}
\AXC{$\stackrel{\Pi_{60}}{\vdots}$}\noLine
\UIC{$\varphi_{34}$}
\AXC{$\stackrel{\Pi_{62}}{\vdots}$}\noLine
\UIC{$\varphi_{136}$}
\AXC{$\stackrel{\Pi_{64}}{\vdots}$}\noLine
\UIC{$x \in \mathbb{R} \to y \in \mathbb{R} \to x \neq 0 \to y \neq 0 \to x^{-1} \cdot 1 = x^{-1}$}
\AXC{$\stackrel{\Pi_{89}}{\vdots}$}\noLine
\UIC{$\varphi_{220}$}
\BIC{$\varphi_{136} \to \varphi_{76}$}
\BIC{$\varphi_{76}$}
\AXC{$\varphi_{76} \to \varphi_{34} \to \varphi_{77}$\lbl{A1}}
\BIC{$\varphi_{34} \to \varphi_{77}$}
\BIC{$\varphi_{77}$}
\end{ptree}
\noindent$\Pi_{91}$:
\begin{ptree}
\AXC{$x = x$\lbl{A5}}
\UIC{$\forall x\,(x = x)$}
\AXC{$\forall x\,(x = x) \to x^{-1} - y^{-1} = x^{-1} - y^{-1}$\lbl{A2}}
\BIC{$x^{-1} - y^{-1} = x^{-1} - y^{-1}$}
\AXC{$\stackrel{\Pi_{90}}{\vdots}$}\noLine
\UIC{$\varphi_{77}$}
\AXC{$\varphi_{77} \to x^{-1} - y^{-1} = x^{-1} - y^{-1} \to \varphi_{215}$\lbl{A1}}
\BIC{$x^{-1} - y^{-1} = x^{-1} - y^{-1} \to \varphi_{215}$}
\BIC{$\varphi_{215}$}
\end{ptree}
\noindent$\Pi_{92}$:
\begin{ptree}
\AXC{$\varphi_{59}$\lbl{A6}}
\UIC{$\forall w\,(\varphi_{59})$}
\UIC{$\forall z\,\forall w\,(\varphi_{59})$}
\AXC{$\forall z\,\forall w\,(\varphi_{59}) \to \varphi_{118}$\lbl{A2}}
\BIC{$\varphi_{118}$}
\AXC{$\varphi_{118} \to x^{-1} = y \cdot (x^{-1} \cdot y^{-1}) \to \varphi_{133}$\lbl{A2}}
\BIC{$x^{-1} = y \cdot (x^{-1} \cdot y^{-1}) \to \varphi_{133}$}
\AXC{$\stackrel{\Pi_{91}}{\vdots}$}\noLine
\UIC{$\varphi_{215}$}
\BIC{$x \in \mathbb{R} \to y \in \mathbb{R} \to \varphi_{53}$}
\AXC{$\varphi_{248}$\lbl{A1}}
\BIC{$(x \in \mathbb{R} \to y \in \mathbb{R} \to x \neq 0 \to y \neq 0 \to \varphi_{119}) \to \varphi_{236}$}
\end{ptree}
\noindent$\Pi_{93}$:
\begin{ptree}
\AXC{$\stackrel{\Pi_{58}}{\vdots}$}\noLine
\UIC{$x \in \mathbb{R} \to y \in \mathbb{R} \to x \neq 0 \to y \neq 0 \to \varphi_{119}$}
\AXC{$\stackrel{\Pi_{92}}{\vdots}$}\noLine
\UIC{$(x \in \mathbb{R} \to y \in \mathbb{R} \to x \neq 0 \to y \neq 0 \to \varphi_{119}) \to \varphi_{236}$}
\BIC{$\varphi_{236}$}
\end{ptree}
\noindent$\Pi_{94}$:
\begin{ptree}
\AXC{$\stackrel{\Pi_{32}}{\vdots}$}\noLine
\UIC{$x^{-1} - y^{-1} = y \cdot (x^{-1} \cdot y^{-1}) - y^{-1} \to \varphi_{119} \to \varphi_{57}$}
\AXC{$\stackrel{\Pi_{93}}{\vdots}$}\noLine
\UIC{$\varphi_{236}$}
\BIC{$x \in \mathbb{R} \to y \in \mathbb{R} \to x \neq 0 \to y \neq 0 \to \varphi_{57}$}
\AXC{$(x \in \mathbb{R} \to y \in \mathbb{R} \to x \neq 0 \to y \neq 0 \to \varphi_{57}) \to \varphi_{245}$\lbl{A1}}
\BIC{$\varphi_{245}$}
\end{ptree}
\noindent$\Pi_{95}$:
\begin{ptree}
\AXC{$\stackrel{\Pi_{26}}{\vdots}$}\noLine
\UIC{$\varphi_{57} \to \varphi_{203}$}
\AXC{$\stackrel{\Pi_{30}}{\vdots}$}\noLine
\UIC{$x \in \mathbb{R} \to y \in \mathbb{R} \to x \neq 0 \to y \neq 0 \to \varphi_{153}$}
\AXC{$\stackrel{\Pi_{94}}{\vdots}$}\noLine
\UIC{$\varphi_{245}$}
\BIC{$(\varphi_{57} \to \varphi_{203}) \to x \in \mathbb{R} \to y \in \mathbb{R} \to x \neq 0 \to \varphi_{94}$}
\BIC{$x \in \mathbb{R} \to y \in \mathbb{R} \to x \neq 0 \to \varphi_{94}$}
\AXC{$(x \in \mathbb{R} \to y \in \mathbb{R} \to x \neq 0 \to \varphi_{94}) \to \varphi_{246}$\lbl{A1}}
\BIC{$\varphi_{246}$}
\end{ptree}
\noindent$\Pi_{96}$:
\begin{ptree}
\AXC{$\stackrel{\Pi_{20}}{\vdots}$}\noLine
\UIC{$\varphi_{223}$}
\AXC{$\stackrel{\Pi_{24}}{\vdots}$}\noLine
\UIC{$x \in \mathbb{R} \to y \in \mathbb{R} \to x \neq 0 \to y \neq 0 \to \varphi_{152}$}
\AXC{$\stackrel{\Pi_{95}}{\vdots}$}\noLine
\UIC{$\varphi_{246}$}
\BIC{$\varphi_{223} \to x \in \mathbb{R} \to y \in \mathbb{R} \to x \neq 0 \to \varphi_{95}$}
\BIC{$x \in \mathbb{R} \to y \in \mathbb{R} \to x \neq 0 \to \varphi_{95}$}
\AXC{$\varphi_{247}$\lbl{A1}}
\BIC{$(x \in \mathbb{R} \to y \in \mathbb{R} \to x \neq 0 \to y \neq 0 \to \varphi_{109}) \to \varphi_{205} \to x \in \mathbb{R} \to \varphi_{96}$}
\end{ptree}
\noindent$\Pi_{97}$:
\begin{ptree}
\AXC{$x \in \mathbb{R} \to y \in \mathbb{R} \to x \neq 0 \to y \neq 0 \to x \in \mathbb{R}$\lbl{A1}}
\AXC{$\stackrel{\Pi_{18}}{\vdots}$}\noLine
\UIC{$\varphi_{204}$}
\BIC{$x \in \mathbb{R} \to y \in \mathbb{R} \to x \neq 0 \to y \neq 0 \to \varphi_{109}$}
\AXC{$\stackrel{\Pi_{96}}{\vdots}$}\noLine
\UIC{$(x \in \mathbb{R} \to y \in \mathbb{R} \to x \neq 0 \to y \neq 0 \to \varphi_{109}) \to \varphi_{205} \to x \in \mathbb{R} \to \varphi_{96}$}
\BIC{$\varphi_{205} \to x \in \mathbb{R} \to \varphi_{96}$}
\end{ptree}
\noindent$\Pi_{98}$:
\begin{ptree}
\AXC{$\stackrel{\Pi_{15}}{\vdots}$}\noLine
\UIC{$\forall v_{605}\,(\varphi_{6})$}
\AXC{$\forall v_{605}\,(\varphi_{6}) \to \varphi_{79}$\lbl{A2}}
\BIC{$\varphi_{79}$}
\AXC{$\varphi_{79} \to \varphi_{196}$\lbl{A2}}
\BIC{$\varphi_{196}$}
\AXC{$\varphi_{196} \to \varphi_{205}$\lbl{A2}}
\BIC{$\varphi_{205}$}
\AXC{$\stackrel{\Pi_{97}}{\vdots}$}\noLine
\UIC{$\varphi_{205} \to x \in \mathbb{R} \to \varphi_{96}$}
\BIC{$x \in \mathbb{R} \to \varphi_{96}$}
\end{ptree}
\noindent$\Pi_{99}$:
\begin{ptree}
\AXC{$\stackrel{\Pi_{14}}{\vdots}$}\noLine
\UIC{$x \in \mathbb{R} \to y \in \mathbb{R} \to x \neq 0 \to \varphi_{48}$}
\AXC{$\stackrel{\Pi_{98}}{\vdots}$}\noLine
\UIC{$x \in \mathbb{R} \to \varphi_{96}$}
\AXC{$(x \in \mathbb{R} \to \varphi_{96}) \to (x \in \mathbb{R} \to y \in \mathbb{R} \to x \neq 0 \to \varphi_{48}) \to \varphi_{224}$\lbl{A1}}
\BIC{$(x \in \mathbb{R} \to y \in \mathbb{R} \to x \neq 0 \to \varphi_{48}) \to \varphi_{224}$}
\BIC{$\varphi_{224}$}
\end{ptree}
\noindent 主木:
\begin{ptree}
\AXC{$\stackrel{\Pi_{2}}{\vdots}$}\noLine
\UIC{$x^{-1} - y^{-1} = x^{-1} \cdot y^{-1} \cdot (y - x) \to \varphi_{165}$}
\AXC{$\stackrel{\Pi_{99}}{\vdots}$}\noLine
\UIC{$\varphi_{224}$}
\BIC{$x \in \mathbb{R} \to y \in \mathbb{R} \to \varphi_{32}$}
\UIC{$\forall y\,(x \in \mathbb{R} \to y \in \mathbb{R} \to \varphi_{32})$}
\UIC{$\forall x\,\forall y\,(x \in \mathbb{R} \to y \in \mathbb{R} \to \varphi_{32})$}
\end{ptree}
\begin{abbrevlist}
\item $\varphi_{1} :\equiv \forall y\,(v_{585} = y \to y = v_{585}) \to v_{585} = v_{586} \to v_{586} = v_{585}$
\item $\varphi_{2} :\equiv \forall y\,(v_{572} = y \to y = v_{572}) \to v_{572} = v_{573} \to v_{573} = v_{572}$
\item $\varphi_{3} :\equiv \forall y\,(v_{583} = y \to y = v_{583}) \to v_{583} = v_{584} \to v_{584} = v_{583}$
\item $\varphi_{4} :\equiv \forall y\,(v_{558} = y \to y = v_{558}) \to v_{558} = v_{559} \to v_{559} = v_{558}$
\item $\varphi_{5} :\equiv \forall v_{609}\,\forall v_{610}\,(v_{608} = v_{609} \to v_{609} = v_{610} \to v_{608} = v_{610})$
\item $\varphi_{6} :\equiv \forall v_{606}\,\forall v_{607}\,(v_{605} = v_{606} \to v_{606} = v_{607} \to v_{605} = v_{607})$
\item $\varphi_{7} :\equiv \forall v_{603}\,\forall v_{604}\,(v_{602} = v_{603} \to v_{603} = v_{604} \to v_{602} = v_{604})$
\item $\varphi_{8} :\equiv \forall v_{600}\,\forall v_{601}\,(v_{599} = v_{600} \to v_{600} = v_{601} \to v_{599} = v_{601})$
\item $\varphi_{9} :\equiv \forall v_{597}\,\forall v_{598}\,(v_{596} = v_{597} \to v_{597} = v_{598} \to v_{596} = v_{598})$
\item $\varphi_{10} :\equiv \forall v_{581}\,\forall v_{582}\,(v_{580} = v_{581} \to v_{581} = v_{582} \to v_{580} = v_{582})$
\item $\varphi_{11} :\equiv \forall v_{578}\,\forall v_{579}\,(v_{577} = v_{578} \to v_{578} = v_{579} \to v_{577} = v_{579})$
\item $\varphi_{12} :\equiv \forall v_{575}\,\forall v_{576}\,(v_{574} = v_{575} \to v_{575} = v_{576} \to v_{574} = v_{576})$
\item $\varphi_{13} :\equiv \forall v_{570}\,\forall v_{571}\,(v_{569} = v_{570} \to v_{570} = v_{571} \to v_{569} = v_{571})$
\item $\varphi_{14} :\equiv \forall v_{567}\,\forall v_{568}\,(v_{566} = v_{567} \to v_{567} = v_{568} \to v_{566} = v_{568})$
\item $\varphi_{15} :\equiv \forall v_{564}\,\forall v_{565}\,(v_{563} = v_{564} \to v_{564} = v_{565} \to v_{563} = v_{565})$
\item $\varphi_{16} :\equiv \forall v_{561}\,\forall v_{562}\,(v_{560} = v_{561} \to v_{561} = v_{562} \to v_{560} = v_{562})$
\item $\varphi_{17} :\equiv \forall x\,(x = x) \to y \cdot (x^{-1} \cdot y^{-1}) = y \cdot (x^{-1} \cdot y^{-1})$
\item $\varphi_{18} :\equiv \forall x\,\forall y\,\forall z\,(x = y \to y = z \to x = z) \to \forall y\,\forall z\,(v_{608} = y \to y = z \to v_{608} = z)$
\item $\varphi_{19} :\equiv \forall x\,\forall y\,\forall z\,(x = y \to y = z \to x = z) \to \forall y\,\forall z\,(v_{605} = y \to y = z \to v_{605} = z)$
\item $\varphi_{20} :\equiv \forall x\,\forall y\,\forall z\,(x = y \to y = z \to x = z) \to \forall y\,\forall z\,(v_{602} = y \to y = z \to v_{602} = z)$
\item $\varphi_{21} :\equiv \forall x\,\forall y\,\forall z\,(x = y \to y = z \to x = z) \to \forall y\,\forall z\,(v_{599} = y \to y = z \to v_{599} = z)$
\item $\varphi_{22} :\equiv \forall x\,\forall y\,\forall z\,(x = y \to y = z \to x = z) \to \forall y\,\forall z\,(v_{596} = y \to y = z \to v_{596} = z)$
\item $\varphi_{23} :\equiv \forall x\,\forall y\,\forall z\,(x = y \to y = z \to x = z) \to \forall y\,\forall z\,(v_{580} = y \to y = z \to v_{580} = z)$
\item $\varphi_{24} :\equiv \forall x\,\forall y\,\forall z\,(x = y \to y = z \to x = z) \to \forall y\,\forall z\,(v_{577} = y \to y = z \to v_{577} = z)$
\item $\varphi_{25} :\equiv \forall x\,\forall y\,\forall z\,(x = y \to y = z \to x = z) \to \forall y\,\forall z\,(v_{574} = y \to y = z \to v_{574} = z)$
\item $\varphi_{26} :\equiv \forall x\,\forall y\,\forall z\,(x = y \to y = z \to x = z) \to \forall y\,\forall z\,(v_{569} = y \to y = z \to v_{569} = z)$
\item $\varphi_{27} :\equiv \forall x\,\forall y\,\forall z\,(x = y \to y = z \to x = z) \to \forall y\,\forall z\,(v_{566} = y \to y = z \to v_{566} = z)$
\item $\varphi_{28} :\equiv \forall x\,\forall y\,\forall z\,(x = y \to y = z \to x = z) \to \forall y\,\forall z\,(v_{563} = y \to y = z \to v_{563} = z)$
\item $\varphi_{29} :\equiv \forall x\,\forall y\,\forall z\,(x = y \to y = z \to x = z) \to \forall y\,\forall z\,(v_{560} = y \to y = z \to v_{560} = z)$
\item $\varphi_{30} :\equiv v_{557} \in \mathbb{R} \to v_{555} \cdot v_{556} \cdot v_{557} = v_{555} \cdot (v_{556} \cdot v_{557})$
\item $\varphi_{31} :\equiv x \neq 0 \to y \neq 0 \to y \cdot (x^{-1} \cdot y^{-1}) = y \cdot (y^{-1} \cdot x^{-1})$
\item $\varphi_{32} :\equiv x \neq 0 \to y \neq 0 \to x^{-1} - y^{-1} = (y - x) \cdot (x^{-1} \cdot y^{-1})$
\item $\varphi_{33} :\equiv x \cdot (x^{-1} \cdot y^{-1}) = y^{-1} \to y^{-1} = x \cdot (x^{-1} \cdot y^{-1})$
\item $\varphi_{34} :\equiv y \cdot (x^{-1} \cdot y^{-1}) = x^{-1} \to x^{-1} = y \cdot (x^{-1} \cdot y^{-1})$
\item $\varphi_{35} :\equiv y \cdot (x^{-1} \cdot y^{-1}) = y \cdot z \to y \cdot (x^{-1} \cdot y^{-1}) = y \cdot w$
\item $\varphi_{36} :\equiv \forall y\,\forall z\,(v_{608} = y \to y = z \to v_{608} = z) \to \forall z\,(v_{608} = v_{609} \to v_{609} = z \to v_{608} = z)$
\item $\varphi_{37} :\equiv \forall y\,\forall z\,(v_{605} = y \to y = z \to v_{605} = z) \to \forall z\,(v_{605} = v_{606} \to v_{606} = z \to v_{605} = z)$
\item $\varphi_{38} :\equiv \forall y\,\forall z\,(v_{602} = y \to y = z \to v_{602} = z) \to \forall z\,(v_{602} = v_{603} \to v_{603} = z \to v_{602} = z)$
\item $\varphi_{39} :\equiv \forall y\,\forall z\,(v_{599} = y \to y = z \to v_{599} = z) \to \forall z\,(v_{599} = v_{600} \to v_{600} = z \to v_{599} = z)$
\item $\varphi_{40} :\equiv \forall y\,\forall z\,(v_{596} = y \to y = z \to v_{596} = z) \to \forall z\,(v_{596} = v_{597} \to v_{597} = z \to v_{596} = z)$
\item $\varphi_{41} :\equiv \forall y\,\forall z\,(v_{580} = y \to y = z \to v_{580} = z) \to \forall z\,(v_{580} = v_{581} \to v_{581} = z \to v_{580} = z)$
\item $\varphi_{42} :\equiv \forall y\,\forall z\,(v_{577} = y \to y = z \to v_{577} = z) \to \forall z\,(v_{577} = v_{578} \to v_{578} = z \to v_{577} = z)$
\item $\varphi_{43} :\equiv \forall y\,\forall z\,(v_{574} = y \to y = z \to v_{574} = z) \to \forall z\,(v_{574} = v_{575} \to v_{575} = z \to v_{574} = z)$
\item $\varphi_{44} :\equiv \forall y\,\forall z\,(v_{569} = y \to y = z \to v_{569} = z) \to \forall z\,(v_{569} = v_{570} \to v_{570} = z \to v_{569} = z)$
\item $\varphi_{45} :\equiv \forall y\,\forall z\,(v_{566} = y \to y = z \to v_{566} = z) \to \forall z\,(v_{566} = v_{567} \to v_{567} = z \to v_{566} = z)$
\item $\varphi_{46} :\equiv \forall y\,\forall z\,(v_{563} = y \to y = z \to v_{563} = z) \to \forall z\,(v_{563} = v_{564} \to v_{564} = z \to v_{563} = z)$
\item $\varphi_{47} :\equiv \forall y\,\forall z\,(v_{560} = y \to y = z \to v_{560} = z) \to \forall z\,(v_{560} = v_{561} \to v_{561} = z \to v_{560} = z)$
\item $\varphi_{48} :\equiv y \neq 0 \to x^{-1} \cdot y^{-1} \cdot (y - x) = (y - x) \cdot (x^{-1} \cdot y^{-1})$
\item $\varphi_{49} :\equiv \forall v_{586}\,(x \cdot (x^{-1} \cdot y^{-1}) = v_{586} \to v_{586} = x \cdot (x^{-1} \cdot y^{-1}))$
\item $\varphi_{50} :\equiv \forall v_{573}\,(x \cdot x^{-1} \cdot y^{-1} = v_{573} \to v_{573} = x \cdot x^{-1} \cdot y^{-1})$
\item $\varphi_{51} :\equiv \forall v_{584}\,(y \cdot (x^{-1} \cdot y^{-1}) = v_{584} \to v_{584} = y \cdot (x^{-1} \cdot y^{-1}))$
\item $\varphi_{52} :\equiv \forall v_{559}\,(y \cdot y^{-1} \cdot x^{-1} = v_{559} \to v_{559} = y \cdot y^{-1} \cdot x^{-1})$
\item $\varphi_{53} :\equiv x \neq 0 \to y \neq 0 \to x^{-1} - y^{-1} = y \cdot (x^{-1} \cdot y^{-1}) - y^{-1}$
\item $\varphi_{54} :\equiv \forall v_{593}\,(v_{592} \in \mathbb{R} \to v_{593} \in \mathbb{R} \to v_{592} \cdot v_{593} = v_{593} \cdot v_{592})$
\item $\varphi_{55} :\equiv \forall v_{588}\,(v_{587} \in \mathbb{R} \to v_{588} \in \mathbb{R} \to v_{587} \cdot v_{588} = v_{588} \cdot v_{587})$
\item $\varphi_{56} :\equiv y \cdot (x^{-1} \cdot y^{-1}) - y^{-1} = y \cdot (x^{-1} \cdot y^{-1}) - y^{-1}$
\item $\varphi_{57} :\equiv x^{-1} - y^{-1} = y \cdot (x^{-1} \cdot y^{-1}) - x \cdot (x^{-1} \cdot y^{-1})$
\item $\varphi_{58} :\equiv \forall v_{594}\,\forall v_{595}\,(v_{594} \in \mathbb{R} \to v_{595} \in \mathbb{R} \to v_{594} - v_{595} \in \mathbb{R})$
\item $\varphi_{59} :\equiv z = w \to x^{-1} - y^{-1} = z - y^{-1} \to x^{-1} - y^{-1} = w - y^{-1}$
\item $\varphi_{60} :\equiv v_{591} \in \mathbb{R} \to v_{589} \cdot v_{590} - v_{589} \cdot v_{591} = v_{589} \cdot (v_{590} - v_{591})$
\item $\varphi_{61} :\equiv \forall z\,(v_{608} = v_{609} \to v_{609} = z \to v_{608} = z) \to v_{608} = v_{609} \to v_{609} = v_{610} \to v_{608} = v_{610}$
\item $\varphi_{62} :\equiv \forall z\,(v_{605} = v_{606} \to v_{606} = z \to v_{605} = z) \to v_{605} = v_{606} \to v_{606} = v_{607} \to v_{605} = v_{607}$
\item $\varphi_{63} :\equiv \forall z\,(v_{602} = v_{603} \to v_{603} = z \to v_{602} = z) \to v_{602} = v_{603} \to v_{603} = v_{604} \to v_{602} = v_{604}$
\item $\varphi_{64} :\equiv \forall z\,(v_{599} = v_{600} \to v_{600} = z \to v_{599} = z) \to v_{599} = v_{600} \to v_{600} = v_{601} \to v_{599} = v_{601}$
\item $\varphi_{65} :\equiv \forall z\,(v_{596} = v_{597} \to v_{597} = z \to v_{596} = z) \to v_{596} = v_{597} \to v_{597} = v_{598} \to v_{596} = v_{598}$
\item $\varphi_{66} :\equiv \forall z\,(v_{580} = v_{581} \to v_{581} = z \to v_{580} = z) \to v_{580} = v_{581} \to v_{581} = v_{582} \to v_{580} = v_{582}$
\item $\varphi_{67} :\equiv \forall z\,(v_{577} = v_{578} \to v_{578} = z \to v_{577} = z) \to v_{577} = v_{578} \to v_{578} = v_{579} \to v_{577} = v_{579}$
\item $\varphi_{68} :\equiv \forall z\,(v_{574} = v_{575} \to v_{575} = z \to v_{574} = z) \to v_{574} = v_{575} \to v_{575} = v_{576} \to v_{574} = v_{576}$
\item $\varphi_{69} :\equiv \forall z\,(v_{569} = v_{570} \to v_{570} = z \to v_{569} = z) \to v_{569} = v_{570} \to v_{570} = v_{571} \to v_{569} = v_{571}$
\item $\varphi_{70} :\equiv \forall z\,(v_{566} = v_{567} \to v_{567} = z \to v_{566} = z) \to v_{566} = v_{567} \to v_{567} = v_{568} \to v_{566} = v_{568}$
\item $\varphi_{71} :\equiv \forall z\,(v_{563} = v_{564} \to v_{564} = z \to v_{563} = z) \to v_{563} = v_{564} \to v_{564} = v_{565} \to v_{563} = v_{565}$
\item $\varphi_{72} :\equiv \forall z\,(v_{560} = v_{561} \to v_{561} = z \to v_{560} = z) \to v_{560} = v_{561} \to v_{561} = v_{562} \to v_{560} = v_{562}$
\item $\varphi_{73} :\equiv \forall y\,\forall z\,(x \in \mathbb{R} \to y \in \mathbb{R} \to z \in \mathbb{R} \to x \cdot y \cdot z = x \cdot (y \cdot z))$
\item $\varphi_{74} :\equiv x \in \mathbb{R} \to y \in \mathbb{R} \to x \neq 0 \to y \neq 0 \to x \cdot (x^{-1} \cdot y^{-1}) = y^{-1}$
\item $\varphi_{75} :\equiv x \in \mathbb{R} \to y \in \mathbb{R} \to x \neq 0 \to y \neq 0 \to y^{-1} = x \cdot (x^{-1} \cdot y^{-1})$
\item $\varphi_{76} :\equiv x \in \mathbb{R} \to y \in \mathbb{R} \to x \neq 0 \to y \neq 0 \to y \cdot (x^{-1} \cdot y^{-1}) = x^{-1}$
\item $\varphi_{77} :\equiv x \in \mathbb{R} \to y \in \mathbb{R} \to x \neq 0 \to y \neq 0 \to x^{-1} = y \cdot (x^{-1} \cdot y^{-1})$
\item $\varphi_{78} :\equiv \forall v_{609}\,\forall v_{610}\,(x^{-1} - y^{-1} = v_{609} \to v_{609} = v_{610} \to x^{-1} - y^{-1} = v_{610})$
\item $\varphi_{79} :\equiv \forall v_{606}\,\forall v_{607}\,(x^{-1} - y^{-1} = v_{606} \to v_{606} = v_{607} \to x^{-1} - y^{-1} = v_{607})$
\item $\varphi_{80} :\equiv \forall v_{603}\,\forall v_{604}\,(x^{-1} - y^{-1} = v_{603} \to v_{603} = v_{604} \to x^{-1} - y^{-1} = v_{604})$
\item $\varphi_{81} :\equiv \forall v_{600}\,\forall v_{601}\,(x^{-1} - y^{-1} = v_{600} \to v_{600} = v_{601} \to x^{-1} - y^{-1} = v_{601})$
\item $\varphi_{82} :\equiv \forall v_{597}\,\forall v_{598}\,(x^{-1} - y^{-1} = v_{597} \to v_{597} = v_{598} \to x^{-1} - y^{-1} = v_{598})$
\item $\varphi_{83} :\equiv x \cdot x^{-1} \cdot y^{-1} = 1 \cdot y^{-1} \to x \cdot (x^{-1} \cdot y^{-1}) = 1 \cdot y^{-1}$
\item $\varphi_{84} :\equiv x \cdot x^{-1} \cdot y^{-1} = z \cdot y^{-1} \to x \cdot x^{-1} \cdot y^{-1} = w \cdot y^{-1}$
\item $\varphi_{85} :\equiv y \cdot y^{-1} \cdot x^{-1} = 1 \cdot x^{-1} \to y \cdot (x^{-1} \cdot y^{-1}) = 1 \cdot x^{-1}$
\item $\varphi_{86} :\equiv y \cdot y^{-1} \cdot x^{-1} = z \cdot x^{-1} \to y \cdot y^{-1} \cdot x^{-1} = w \cdot x^{-1}$
\item $\varphi_{87} :\equiv y \in \mathbb{R} \to x \neq 0 \to y \neq 0 \to x \cdot (x^{-1} \cdot y^{-1}) = x^{-1} \cdot y^{-1} \cdot x$
\item $\varphi_{88} :\equiv y \in \mathbb{R} \to x \neq 0 \to y \neq 0 \to y \cdot (x^{-1} \cdot y^{-1}) = x^{-1} \cdot y^{-1} \cdot y$
\item $\varphi_{89} :\equiv y \in \mathbb{R} \to x \neq 0 \to y \neq 0 \to x \cdot x^{-1} \cdot y^{-1} = x \cdot (x^{-1} \cdot y^{-1})$
\item $\varphi_{90} :\equiv y \in \mathbb{R} \to x \neq 0 \to y \neq 0 \to x \cdot (x^{-1} \cdot y^{-1}) = x \cdot x^{-1} \cdot y^{-1}$
\item $\varphi_{91} :\equiv y \in \mathbb{R} \to x \neq 0 \to y \neq 0 \to y \cdot y^{-1} \cdot x^{-1} = y \cdot (y^{-1} \cdot x^{-1})$
\item $\varphi_{92} :\equiv y \in \mathbb{R} \to x \neq 0 \to y \neq 0 \to y \cdot (y^{-1} \cdot x^{-1}) = y \cdot y^{-1} \cdot x^{-1}$
\item $\varphi_{93} :\equiv y \in \mathbb{R} \to x \neq 0 \to y \neq 0 \to y \cdot (x^{-1} \cdot y^{-1}) = y \cdot y^{-1} \cdot x^{-1}$
\item $\varphi_{94} :\equiv y \neq 0 \to x^{-1} - y^{-1} = x^{-1} \cdot y^{-1} \cdot y - x \cdot (x^{-1} \cdot y^{-1})$
\item $\varphi_{95} :\equiv y \neq 0 \to x^{-1} - y^{-1} = x^{-1} \cdot y^{-1} \cdot y - x^{-1} \cdot y^{-1} \cdot x$
\item $\varphi_{96} :\equiv y \in \mathbb{R} \to x \neq 0 \to y \neq 0 \to x^{-1} - y^{-1} = x^{-1} \cdot y^{-1} \cdot (y - x)$
\item $\varphi_{97} :\equiv y - x \in \mathbb{R} \to x^{-1} \cdot y^{-1} \cdot (y - x) = (y - x) \cdot (x^{-1} \cdot y^{-1})$
\item $\varphi_{98} :\equiv x^{-1} \cdot y^{-1} \in \mathbb{R} \to x \cdot (x^{-1} \cdot y^{-1}) = x^{-1} \cdot y^{-1} \cdot x$
\item $\varphi_{99} :\equiv x^{-1} \cdot y^{-1} \in \mathbb{R} \to y \cdot (x^{-1} \cdot y^{-1}) = x^{-1} \cdot y^{-1} \cdot y$
\item $\varphi_{100} :\equiv \forall x \in \mathbb{R}\,(x \neq 0 \to x \cdot x^{-1} = 1) \to x \in \mathbb{R} \to x \neq 0 \to x \cdot x^{-1} = 1$
\item $\varphi_{101} :\equiv \forall x \in \mathbb{R}\,(x \neq 0 \to x \cdot x^{-1} = 1) \to y \in \mathbb{R} \to y \neq 0 \to y \cdot y^{-1} = 1$
\item $\varphi_{102} :\equiv x \in \mathbb{R} \to y \in \mathbb{R} \to x \neq 0 \to y \neq 0 \to x \cdot x^{-1} \cdot y^{-1} = 1 \cdot y^{-1}$
\item $\varphi_{103} :\equiv x \in \mathbb{R} \to y \in \mathbb{R} \to x \neq 0 \to y \neq 0 \to x \cdot (x^{-1} \cdot y^{-1}) = 1 \cdot y^{-1}$
\item $\varphi_{104} :\equiv x \in \mathbb{R} \to y \in \mathbb{R} \to x \neq 0 \to y \neq 0 \to x \cdot (x^{-1} \cdot y^{-1}) = y^{-1} \cdot 1$
\item $\varphi_{105} :\equiv x \in \mathbb{R} \to y \in \mathbb{R} \to x \neq 0 \to y \neq 0 \to y \cdot y^{-1} \cdot x^{-1} = 1 \cdot x^{-1}$
\item $\varphi_{106} :\equiv x \in \mathbb{R} \to y \in \mathbb{R} \to x \neq 0 \to y \neq 0 \to x^{-1} \cdot y^{-1} = y^{-1} \cdot x^{-1}$
\item $\varphi_{107} :\equiv x \in \mathbb{R} \to y \in \mathbb{R} \to x \neq 0 \to y \neq 0 \to y \cdot (x^{-1} \cdot y^{-1}) = 1 \cdot x^{-1}$
\item $\varphi_{108} :\equiv x \in \mathbb{R} \to y \in \mathbb{R} \to x \neq 0 \to y \neq 0 \to y \cdot (x^{-1} \cdot y^{-1}) = x^{-1} \cdot 1$
\item $\varphi_{109} :\equiv x^{-1} \cdot y^{-1} \cdot y - x^{-1} \cdot y^{-1} \cdot x = x^{-1} \cdot y^{-1} \cdot (y - x)$
\item $\varphi_{110} :\equiv x^{-1} \cdot y^{-1} \cdot y - x \cdot (x^{-1} \cdot y^{-1}) = x^{-1} \cdot y^{-1} \cdot y - w$
\item $\varphi_{111} :\equiv x^{-1} \cdot y^{-1} \cdot y - x \cdot (x^{-1} \cdot y^{-1}) = x^{-1} \cdot y^{-1} \cdot y - z$
\item $\varphi_{112} :\equiv y \cdot (x^{-1} \cdot y^{-1}) - x \cdot (x^{-1} \cdot y^{-1}) = w - x \cdot (x^{-1} \cdot y^{-1})$
\item $\varphi_{113} :\equiv y \cdot (x^{-1} \cdot y^{-1}) - x \cdot (x^{-1} \cdot y^{-1}) = z - x \cdot (x^{-1} \cdot y^{-1})$
\item $\varphi_{114} :\equiv \forall x \in \mathbb{R}\,(x \neq 0 \to x^{-1} \in \mathbb{R}) \to y \in \mathbb{R} \to y \neq 0 \to y^{-1} \in \mathbb{R}$
\item $\varphi_{115} :\equiv \forall x \in \mathbb{R}\,(x \neq 0 \to x^{-1} \in \mathbb{R}) \to x \in \mathbb{R} \to x \neq 0 \to x^{-1} \in \mathbb{R}$
\item $\varphi_{116} :\equiv \forall c\,\forall x\,\forall a\,(c \in \mathbb{R} \to x \in \mathbb{R} \to a \in \mathbb{R} \to c \cdot x - c \cdot a = c \cdot (x - a))$
\item $\varphi_{117} :\equiv \forall y\,\forall z\,(v_{555} \in \mathbb{R} \to y \in \mathbb{R} \to z \in \mathbb{R} \to v_{555} \cdot y \cdot z = v_{555} \cdot (y \cdot z))$
\item $\varphi_{118} :\equiv \forall w\,(x^{-1} = w \to x^{-1} - y^{-1} = x^{-1} - y^{-1} \to x^{-1} - y^{-1} = w - y^{-1})$
\item $\varphi_{119} :\equiv y \cdot (x^{-1} \cdot y^{-1}) - y^{-1} = y \cdot (x^{-1} \cdot y^{-1}) - x \cdot (x^{-1} \cdot y^{-1})$
\item $\varphi_{120} :\equiv \forall v_{581}\,\forall v_{582}\,(x \cdot (x^{-1} \cdot y^{-1}) = v_{581} \to v_{581} = v_{582} \to x \cdot (x^{-1} \cdot y^{-1}) = v_{582})$
\item $\varphi_{121} :\equiv \forall v_{578}\,\forall v_{579}\,(x \cdot (x^{-1} \cdot y^{-1}) = v_{578} \to v_{578} = v_{579} \to x \cdot (x^{-1} \cdot y^{-1}) = v_{579})$
\item $\varphi_{122} :\equiv \forall v_{575}\,\forall v_{576}\,(x \cdot (x^{-1} \cdot y^{-1}) = v_{575} \to v_{575} = v_{576} \to x \cdot (x^{-1} \cdot y^{-1}) = v_{576})$
\item $\varphi_{123} :\equiv \forall v_{570}\,\forall v_{571}\,(y \cdot (x^{-1} \cdot y^{-1}) = v_{570} \to v_{570} = v_{571} \to y \cdot (x^{-1} \cdot y^{-1}) = v_{571})$
\item $\varphi_{124} :\equiv \forall v_{567}\,\forall v_{568}\,(y \cdot (x^{-1} \cdot y^{-1}) = v_{567} \to v_{567} = v_{568} \to y \cdot (x^{-1} \cdot y^{-1}) = v_{568})$
\item $\varphi_{125} :\equiv \forall v_{564}\,\forall v_{565}\,(y \cdot (x^{-1} \cdot y^{-1}) = v_{564} \to v_{564} = v_{565} \to y \cdot (x^{-1} \cdot y^{-1}) = v_{565})$
\item $\varphi_{126} :\equiv \forall v_{561}\,\forall v_{562}\,(y \cdot (x^{-1} \cdot y^{-1}) = v_{561} \to v_{561} = v_{562} \to y \cdot (x^{-1} \cdot y^{-1}) = v_{562})$
\item $\varphi_{127} :\equiv x \cdot x^{-1} \cdot y^{-1} = x \cdot x^{-1} \cdot y^{-1} \to x \cdot x^{-1} \cdot y^{-1} = 1 \cdot y^{-1}$
\item $\varphi_{128} :\equiv x^{-1} \in \mathbb{R} \to y^{-1} \in \mathbb{R} \to x \cdot x^{-1} \cdot y^{-1} = x \cdot (x^{-1} \cdot y^{-1})$
\item $\varphi_{129} :\equiv y \cdot y^{-1} \cdot x^{-1} = y \cdot y^{-1} \cdot x^{-1} \to y \cdot y^{-1} \cdot x^{-1} = 1 \cdot x^{-1}$
\item $\varphi_{130} :\equiv y^{-1} \in \mathbb{R} \to x^{-1} \in \mathbb{R} \to y \cdot y^{-1} \cdot x^{-1} = y \cdot (y^{-1} \cdot x^{-1})$
\item $\varphi_{131} :\equiv \forall z\,(x \in \mathbb{R} \to x^{-1} \in \mathbb{R} \to z \in \mathbb{R} \to x \cdot x^{-1} \cdot z = x \cdot (x^{-1} \cdot z))$
\item $\varphi_{132} :\equiv \forall z\,(v_{555} \in \mathbb{R} \to v_{556} \in \mathbb{R} \to z \in \mathbb{R} \to v_{555} \cdot v_{556} \cdot z = v_{555} \cdot (v_{556} \cdot z))$
\item $\varphi_{133} :\equiv x^{-1} - y^{-1} = x^{-1} - y^{-1} \to x^{-1} - y^{-1} = y \cdot (x^{-1} \cdot y^{-1}) - y^{-1}$
\item $\varphi_{134} :\equiv \forall x\,\forall a\,(v_{589} \in \mathbb{R} \to x \in \mathbb{R} \to a \in \mathbb{R} \to v_{589} \cdot x - v_{589} \cdot a = v_{589} \cdot (x - a))$
\item $\varphi_{135} :\equiv x \cdot (x^{-1} \cdot y^{-1}) = y^{-1} \cdot 1 \to y^{-1} \cdot 1 = y^{-1} \to x \cdot (x^{-1} \cdot y^{-1}) = y^{-1}$
\item $\varphi_{136} :\equiv y \cdot (x^{-1} \cdot y^{-1}) = x^{-1} \cdot 1 \to x^{-1} \cdot 1 = x^{-1} \to y \cdot (x^{-1} \cdot y^{-1}) = x^{-1}$
\item $\varphi_{137} :\equiv \forall v_{582}\,(x \cdot (x^{-1} \cdot y^{-1}) = y^{-1} \cdot 1 \to y^{-1} \cdot 1 = v_{582} \to x \cdot (x^{-1} \cdot y^{-1}) = v_{582})$
\item $\varphi_{138} :\equiv \forall v_{579}\,(x \cdot (x^{-1} \cdot y^{-1}) = 1 \cdot y^{-1} \to 1 \cdot y^{-1} = v_{579} \to x \cdot (x^{-1} \cdot y^{-1}) = v_{579})$
\item $\varphi_{139} :\equiv \forall v_{571}\,(y \cdot (x^{-1} \cdot y^{-1}) = x^{-1} \cdot 1 \to x^{-1} \cdot 1 = v_{571} \to y \cdot (x^{-1} \cdot y^{-1}) = v_{571})$
\item $\varphi_{140} :\equiv \forall v_{568}\,(y \cdot (x^{-1} \cdot y^{-1}) = 1 \cdot x^{-1} \to 1 \cdot x^{-1} = v_{568} \to y \cdot (x^{-1} \cdot y^{-1}) = v_{568})$
\item $\varphi_{141} :\equiv x \cdot x^{-1} \cdot y^{-1} = x \cdot (x^{-1} \cdot y^{-1}) \to x \cdot (x^{-1} \cdot y^{-1}) = x \cdot x^{-1} \cdot y^{-1}$
\item $\varphi_{142} :\equiv y \cdot y^{-1} \cdot x^{-1} = y \cdot (y^{-1} \cdot x^{-1}) \to y \cdot (y^{-1} \cdot x^{-1}) = y \cdot y^{-1} \cdot x^{-1}$
\item $\varphi_{143} :\equiv y \cdot (y^{-1} \cdot x^{-1}) = y \cdot y^{-1} \cdot x^{-1} \to y \cdot (x^{-1} \cdot y^{-1}) = y \cdot y^{-1} \cdot x^{-1}$
\item $\varphi_{144} :\equiv y \cdot (x^{-1} \cdot y^{-1}) = y \cdot (x^{-1} \cdot y^{-1}) \to y \cdot (x^{-1} \cdot y^{-1}) = y \cdot (y^{-1} \cdot x^{-1})$
\item $\varphi_{145} :\equiv \forall x\,\forall y\,(x \in \mathbb{R} \to y \in \mathbb{R} \to x \cdot y = y \cdot x) \to \forall y\,(x \in \mathbb{R} \to y \in \mathbb{R} \to x \cdot y = y \cdot x)$
\item $\varphi_{146} :\equiv \forall x\,\forall y\,(x \in \mathbb{R} \to y \in \mathbb{R} \to x \cdot y = y \cdot x) \to \forall y\,(1 \in \mathbb{R} \to y \in \mathbb{R} \to 1 \cdot y = y \cdot 1)$
\item $\varphi_{147} :\equiv \forall a\,(v_{589} \in \mathbb{R} \to v_{590} \in \mathbb{R} \to a \in \mathbb{R} \to v_{589} \cdot v_{590} - v_{589} \cdot a = v_{589} \cdot (v_{590} - a))$
\item $\varphi_{148} :\equiv \forall v_{556}\,\forall v_{557}\,(y \in \mathbb{R} \to v_{556} \in \mathbb{R} \to v_{557} \in \mathbb{R} \to y \cdot v_{556} \cdot v_{557} = y \cdot (v_{556} \cdot v_{557}))$
\item $\varphi_{149} :\equiv \forall v_{557}\,(y \in \mathbb{R} \to y^{-1} \in \mathbb{R} \to v_{557} \in \mathbb{R} \to y \cdot y^{-1} \cdot v_{557} = y \cdot (y^{-1} \cdot v_{557}))$
\item $\varphi_{150} :\equiv x^{-1} \cdot y^{-1} \cdot y - x \cdot (x^{-1} \cdot y^{-1}) = x^{-1} \cdot y^{-1} \cdot y - x \cdot (x^{-1} \cdot y^{-1})$
\item $\varphi_{151} :\equiv y \cdot (x^{-1} \cdot y^{-1}) - x \cdot (x^{-1} \cdot y^{-1}) = y \cdot (x^{-1} \cdot y^{-1}) - x \cdot (x^{-1} \cdot y^{-1})$
\item $\varphi_{152} :\equiv x^{-1} \cdot y^{-1} \cdot y - x \cdot (x^{-1} \cdot y^{-1}) = x^{-1} \cdot y^{-1} \cdot y - x^{-1} \cdot y^{-1} \cdot x$
\item $\varphi_{153} :\equiv y \cdot (x^{-1} \cdot y^{-1}) - x \cdot (x^{-1} \cdot y^{-1}) = x^{-1} \cdot y^{-1} \cdot y - x \cdot (x^{-1} \cdot y^{-1})$
\item $\varphi_{154} :\equiv x \cdot (x^{-1} \cdot y^{-1}) = 1 \cdot y^{-1} \to 1 \cdot y^{-1} = y^{-1} \cdot 1 \to x \cdot (x^{-1} \cdot y^{-1}) = y^{-1} \cdot 1$
\item $\varphi_{155} :\equiv y \cdot (x^{-1} \cdot y^{-1}) = 1 \cdot x^{-1} \to 1 \cdot x^{-1} = x^{-1} \cdot 1 \to y \cdot (x^{-1} \cdot y^{-1}) = x^{-1} \cdot 1$
\item $\varphi_{156} :\equiv \forall v_{595}\,(y \in \mathbb{R} \to v_{595} \in \mathbb{R} \to y - v_{595} \in \mathbb{R}) \to y \in \mathbb{R} \to x \in \mathbb{R} \to y - x \in \mathbb{R}$
\item $\varphi_{157} :\equiv \forall w\,(x \cdot x^{-1} = w \to x \cdot x^{-1} \cdot y^{-1} = x \cdot x^{-1} \cdot y^{-1} \to x \cdot x^{-1} \cdot y^{-1} = w \cdot y^{-1})$
\item $\varphi_{158} :\equiv \forall w\,(y \cdot y^{-1} = w \to y \cdot y^{-1} \cdot x^{-1} = y \cdot y^{-1} \cdot x^{-1} \to y \cdot y^{-1} \cdot x^{-1} = w \cdot x^{-1})$
\item $\varphi_{159} :\equiv \forall w\,(x^{-1} \cdot y^{-1} = w \to y \cdot (x^{-1} \cdot y^{-1}) = y \cdot (x^{-1} \cdot y^{-1}) \to y \cdot (x^{-1} \cdot y^{-1}) = y \cdot w)$
\item $\varphi_{160} :\equiv \forall x\,\forall y\,(x \in \mathbb{R} \to y \in \mathbb{R} \to x \cdot y = y \cdot x) \to \forall y\,(v_{592} \in \mathbb{R} \to y \in \mathbb{R} \to v_{592} \cdot y = y \cdot v_{592})$
\item $\varphi_{161} :\equiv \forall v_{593}\,(x^{-1} \cdot y^{-1} \in \mathbb{R} \to v_{593} \in \mathbb{R} \to x^{-1} \cdot y^{-1} \cdot v_{593} = v_{593} \cdot (x^{-1} \cdot y^{-1}))$
\item $\varphi_{162} :\equiv \forall x\,\forall y\,(x \in \mathbb{R} \to y \in \mathbb{R} \to x \cdot y = y \cdot x) \to \forall y\,(v_{587} \in \mathbb{R} \to y \in \mathbb{R} \to v_{587} \cdot y = y \cdot v_{587})$
\item $\varphi_{163} :\equiv \forall y\,(1 \in \mathbb{R} \to y \in \mathbb{R} \to 1 \cdot y = y \cdot 1) \to 1 \in \mathbb{R} \to y^{-1} \in \mathbb{R} \to 1 \cdot y^{-1} = y^{-1} \cdot 1$
\item $\varphi_{164} :\equiv \forall y\,(1 \in \mathbb{R} \to y \in \mathbb{R} \to 1 \cdot y = y \cdot 1) \to 1 \in \mathbb{R} \to x^{-1} \in \mathbb{R} \to 1 \cdot x^{-1} = x^{-1} \cdot 1$
\item $\varphi_{165} :\equiv x^{-1} \cdot y^{-1} \cdot (y - x) = (y - x) \cdot (x^{-1} \cdot y^{-1}) \to x^{-1} - y^{-1} = (y - x) \cdot (x^{-1} \cdot y^{-1})$
\item $\varphi_{166} :\equiv \forall x\,\forall y\,(x \in \mathbb{R} \to y \in \mathbb{R} \to x - y \in \mathbb{R}) \to \forall y\,(v_{594} \in \mathbb{R} \to y \in \mathbb{R} \to v_{594} - y \in \mathbb{R})$
\item $\varphi_{167} :\equiv (x \in \mathbb{R} \to y \in \mathbb{R} \to x \neq 0 \to y \neq 0 \to x \neq 0) \to x \in \mathbb{R} \to y \in \mathbb{R} \to x \neq 0 \to y \neq 0 \to x \cdot x^{-1} = 1$
\item $\varphi_{168} :\equiv (x \in \mathbb{R} \to y \in \mathbb{R} \to x \neq 0 \to y \neq 0 \to y \neq 0) \to x \in \mathbb{R} \to y \in \mathbb{R} \to x \neq 0 \to y \neq 0 \to y \cdot y^{-1} = 1$
\item $\varphi_{169} :\equiv (x \in \mathbb{R} \to y \in \mathbb{R} \to x \neq 0 \to y \neq 0 \to y \neq 0) \to x \in \mathbb{R} \to y \in \mathbb{R} \to x \neq 0 \to y \neq 0 \to y^{-1} \in \mathbb{R}$
\item $\varphi_{170} :\equiv (x \in \mathbb{R} \to y \in \mathbb{R} \to x \neq 0 \to y \neq 0 \to x \neq 0) \to x \in \mathbb{R} \to y \in \mathbb{R} \to x \neq 0 \to y \neq 0 \to x^{-1} \in \mathbb{R}$
\item $\varphi_{171} :\equiv \forall x\,\forall y\,(x \in \mathbb{R} \to y \in \mathbb{R} \to x \cdot y \in \mathbb{R}) \to \forall y\,(x^{-1} \in \mathbb{R} \to y \in \mathbb{R} \to x^{-1} \cdot y \in \mathbb{R})$
\item $\varphi_{172} :\equiv \forall x\,\forall y\,(x \in \mathbb{R} \to y \in \mathbb{R} \to x \cdot y = y \cdot x) \to \forall y\,(x^{-1} \in \mathbb{R} \to y \in \mathbb{R} \to x^{-1} \cdot y = y \cdot x^{-1})$
\item $\varphi_{173} :\equiv \forall y\,(v_{594} \in \mathbb{R} \to y \in \mathbb{R} \to v_{594} - y \in \mathbb{R}) \to v_{594} \in \mathbb{R} \to v_{595} \in \mathbb{R} \to v_{594} - v_{595} \in \mathbb{R}$
\item $\varphi_{174} :\equiv \forall v_{610}\,(x^{-1} - y^{-1} = x^{-1} \cdot y^{-1} \cdot (y - x) \to x^{-1} \cdot y^{-1} \cdot (y - x) = v_{610} \to x^{-1} - y^{-1} = v_{610})$
\item $\varphi_{175} :\equiv \forall v_{576}\,(x \cdot (x^{-1} \cdot y^{-1}) = x \cdot x^{-1} \cdot y^{-1} \to x \cdot x^{-1} \cdot y^{-1} = v_{576} \to x \cdot (x^{-1} \cdot y^{-1}) = v_{576})$
\item $\varphi_{176} :\equiv \forall v_{565}\,(y \cdot (x^{-1} \cdot y^{-1}) = y \cdot y^{-1} \cdot x^{-1} \to y \cdot y^{-1} \cdot x^{-1} = v_{565} \to y \cdot (x^{-1} \cdot y^{-1}) = v_{565})$
\item $\varphi_{177} :\equiv \forall v_{562}\,(y \cdot (x^{-1} \cdot y^{-1}) = y \cdot (y^{-1} \cdot x^{-1}) \to y \cdot (y^{-1} \cdot x^{-1}) = v_{562} \to y \cdot (x^{-1} \cdot y^{-1}) = v_{562})$
\item $\varphi_{178} :\equiv (x \in \mathbb{R} \to y \in \mathbb{R} \to x \neq 0 \to y \neq 0 \to x \in \mathbb{R}) \to x \in \mathbb{R} \to y \in \mathbb{R} \to x \neq 0 \to y \neq 0 \to y - x \in \mathbb{R}$
\item $\varphi_{179} :\equiv \forall y\,(v_{592} \in \mathbb{R} \to y \in \mathbb{R} \to v_{592} \cdot y = y \cdot v_{592}) \to v_{592} \in \mathbb{R} \to v_{593} \in \mathbb{R} \to v_{592} \cdot v_{593} = v_{593} \cdot v_{592}$
\item $\varphi_{180} :\equiv \forall y\,(v_{587} \in \mathbb{R} \to y \in \mathbb{R} \to v_{587} \cdot y = y \cdot v_{587}) \to v_{587} \in \mathbb{R} \to v_{588} \in \mathbb{R} \to v_{587} \cdot v_{588} = v_{588} \cdot v_{587}$
\item $\varphi_{181} :\equiv (x \in \mathbb{R} \to y \in \mathbb{R} \to x \neq 0 \to y \neq 0 \to y^{-1} \in \mathbb{R}) \to x \in \mathbb{R} \to y \in \mathbb{R} \to x \neq 0 \to y \neq 0 \to y^{-1} \cdot 1 = y^{-1}$
\item $\varphi_{182} :\equiv (x \in \mathbb{R} \to y \in \mathbb{R} \to x \neq 0 \to y \neq 0 \to x^{-1} \in \mathbb{R}) \to x \in \mathbb{R} \to y \in \mathbb{R} \to x \neq 0 \to y \neq 0 \to x^{-1} \cdot 1 = x^{-1}$
\item $\varphi_{183} :\equiv \forall v_{598}\,(x^{-1} - y^{-1} = y \cdot (x^{-1} \cdot y^{-1}) - y^{-1} \to y \cdot (x^{-1} \cdot y^{-1}) - y^{-1} = v_{598} \to x^{-1} - y^{-1} = v_{598})$
\item $\varphi_{184} :\equiv \forall y\,(x^{-1} \in \mathbb{R} \to y \in \mathbb{R} \to x^{-1} \cdot y \in \mathbb{R}) \to x^{-1} \in \mathbb{R} \to y^{-1} \in \mathbb{R} \to x^{-1} \cdot y^{-1} \in \mathbb{R}$
\item $\varphi_{185} :\equiv y \cdot (x^{-1} \cdot y^{-1}) - y^{-1} = y \cdot (x^{-1} \cdot y^{-1}) - z \to y \cdot (x^{-1} \cdot y^{-1}) - y^{-1} = y \cdot (x^{-1} \cdot y^{-1}) - w$
\item $\varphi_{186} :\equiv (x \in \mathbb{R} \to y \in \mathbb{R} \to x \neq 0 \to y \neq 0 \to y^{-1} \in \mathbb{R}) \to x \in \mathbb{R} \to y \in \mathbb{R} \to x \neq 0 \to y \neq 0 \to 1 \cdot y^{-1} = y^{-1} \cdot 1$
\item $\varphi_{187} :\equiv (x \in \mathbb{R} \to y \in \mathbb{R} \to x \neq 0 \to y \neq 0 \to x^{-1} \in \mathbb{R}) \to x \in \mathbb{R} \to y \in \mathbb{R} \to x \neq 0 \to y \neq 0 \to 1 \cdot x^{-1} = x^{-1} \cdot 1$
\item $\varphi_{188} :\equiv (x \in \mathbb{R} \to y \in \mathbb{R} \to x \neq 0 \to y \neq 0 \to y^{-1} \in \mathbb{R}) \to x \in \mathbb{R} \to y \in \mathbb{R} \to x \neq 0 \to y \neq 0 \to x^{-1} \cdot y^{-1} \in \mathbb{R}$
\item $\varphi_{189} :\equiv \forall v_{587}\,(\varphi_{55}) \to \forall v_{588}\,(y \in \mathbb{R} \to v_{588} \in \mathbb{R} \to y \cdot v_{588} = v_{588} \cdot y)$
\item $\varphi_{190} :\equiv \forall y\,(x^{-1} \in \mathbb{R} \to y \in \mathbb{R} \to x^{-1} \cdot y = y \cdot x^{-1}) \to x^{-1} \in \mathbb{R} \to y^{-1} \in \mathbb{R} \to x^{-1} \cdot y^{-1} = y^{-1} \cdot x^{-1}$
\item $\varphi_{191} :\equiv \forall w\,(y^{-1} = w \to \varphi_{56} \to y \cdot (x^{-1} \cdot y^{-1}) - y^{-1} = y \cdot (x^{-1} \cdot y^{-1}) - w)$
\item $\varphi_{192} :\equiv \forall v_{588}\,(y \in \mathbb{R} \to v_{588} \in \mathbb{R} \to y \cdot v_{588} = v_{588} \cdot y) \to y \in \mathbb{R} \to \varphi_{99}$
\item $\varphi_{193} :\equiv (x \in \mathbb{R} \to y \in \mathbb{R} \to x \neq 0 \to y \neq 0 \to y^{-1} \in \mathbb{R}) \to x \in \mathbb{R} \to \varphi_{89}$
\item $\varphi_{194} :\equiv (x \in \mathbb{R} \to y \in \mathbb{R} \to x \neq 0 \to y \neq 0 \to x^{-1} \in \mathbb{R}) \to x \in \mathbb{R} \to \varphi_{91}$
\item $\varphi_{195} :\equiv \forall v_{591}\,(x^{-1} \cdot y^{-1} \in \mathbb{R} \to y \in \mathbb{R} \to v_{591} \in \mathbb{R} \to x^{-1} \cdot y^{-1} \cdot y - x^{-1} \cdot y^{-1} \cdot v_{591} = x^{-1} \cdot y^{-1} \cdot (y - v_{591}))$
\item $\varphi_{196} :\equiv \forall v_{607}\,(x^{-1} - y^{-1} = x^{-1} \cdot y^{-1} \cdot y - x^{-1} \cdot y^{-1} \cdot x \to x^{-1} \cdot y^{-1} \cdot y - x^{-1} \cdot y^{-1} \cdot x = v_{607} \to x^{-1} - y^{-1} = v_{607})$
\item $\varphi_{197} :\equiv \forall v_{604}\,(x^{-1} - y^{-1} = x^{-1} \cdot y^{-1} \cdot y - x \cdot (x^{-1} \cdot y^{-1}) \to x^{-1} \cdot y^{-1} \cdot y - x \cdot (x^{-1} \cdot y^{-1}) = v_{604} \to x^{-1} - y^{-1} = v_{604})$
\item $\varphi_{198} :\equiv \forall v_{601}\,(\varphi_{57} \to y \cdot (x^{-1} \cdot y^{-1}) - x \cdot (x^{-1} \cdot y^{-1}) = v_{601} \to x^{-1} - y^{-1} = v_{601})$
\item $\varphi_{199} :\equiv (x \in \mathbb{R} \to y \in \mathbb{R} \to x \neq 0 \to y \neq 0 \to y - x \in \mathbb{R}) \to x \in \mathbb{R} \to y \in \mathbb{R} \to x \neq 0 \to \varphi_{48}$
\item $\varphi_{200} :\equiv (x \in \mathbb{R} \to y \in \mathbb{R} \to x \neq 0 \to y \neq 0 \to x^{-1} \cdot y^{-1} \in \mathbb{R}) \to x \in \mathbb{R} \to \varphi_{87}$
\item $\varphi_{201} :\equiv (x \in \mathbb{R} \to y \in \mathbb{R} \to x \neq 0 \to y \neq 0 \to x^{-1} \cdot y^{-1} \in \mathbb{R}) \to x \in \mathbb{R} \to \varphi_{88}$
\item $\varphi_{202} :\equiv \forall v_{590}\,\forall v_{591}\,(x^{-1} \cdot y^{-1} \in \mathbb{R} \to v_{590} \in \mathbb{R} \to v_{591} \in \mathbb{R} \to x^{-1} \cdot y^{-1} \cdot v_{590} - x^{-1} \cdot y^{-1} \cdot v_{591} = x^{-1} \cdot y^{-1} \cdot (v_{590} - v_{591}))$
\item $\varphi_{203} :\equiv \varphi_{153} \to x^{-1} - y^{-1} = x^{-1} \cdot y^{-1} \cdot y - x \cdot (x^{-1} \cdot y^{-1})$
\item $\varphi_{204} :\equiv (x \in \mathbb{R} \to y \in \mathbb{R} \to x \neq 0 \to y \neq 0 \to x \in \mathbb{R}) \to x \in \mathbb{R} \to y \in \mathbb{R} \to x \neq 0 \to y \neq 0 \to \varphi_{109}$
\item $\varphi_{205} :\equiv x^{-1} - y^{-1} = x^{-1} \cdot y^{-1} \cdot y - x^{-1} \cdot y^{-1} \cdot x \to \varphi_{109} \to x^{-1} - y^{-1} = x^{-1} \cdot y^{-1} \cdot (y - x)$
\item $\varphi_{206} :\equiv (x \in \mathbb{R} \to x \neq 0 \to x \cdot x^{-1} = 1) \to (x \in \mathbb{R} \to y \in \mathbb{R} \to x \neq 0 \to y \neq 0 \to x \in \mathbb{R}) \to \varphi_{167}$
\item $\varphi_{207} :\equiv (y \in \mathbb{R} \to y \neq 0 \to y \cdot y^{-1} = 1) \to (x \in \mathbb{R} \to y \in \mathbb{R} \to x \neq 0 \to y \neq 0 \to y \in \mathbb{R}) \to \varphi_{168}$
\item $\varphi_{208} :\equiv (x \in \mathbb{R} \to y \in \mathbb{R} \to x \neq 0 \to y \neq 0 \to x^{-1} \in \mathbb{R}) \to (x \in \mathbb{R} \to y \in \mathbb{R} \to x \neq 0 \to y \neq 0 \to y^{-1} \in \mathbb{R}) \to \varphi_{106}$
\item $\varphi_{209} :\equiv (1 \in \mathbb{R} \to y^{-1} \in \mathbb{R} \to 1 \cdot y^{-1} = y^{-1} \cdot 1) \to 1 \in \mathbb{R} \to \varphi_{186}$
\item $\varphi_{210} :\equiv (1 \in \mathbb{R} \to x^{-1} \in \mathbb{R} \to 1 \cdot x^{-1} = x^{-1} \cdot 1) \to 1 \in \mathbb{R} \to \varphi_{187}$
\item $\varphi_{211} :\equiv (y \in \mathbb{R} \to y \neq 0 \to y^{-1} \in \mathbb{R}) \to (x \in \mathbb{R} \to y \in \mathbb{R} \to x \neq 0 \to y \neq 0 \to y \in \mathbb{R}) \to \varphi_{169}$
\item $\varphi_{212} :\equiv (x \in \mathbb{R} \to x \neq 0 \to x^{-1} \in \mathbb{R}) \to (x \in \mathbb{R} \to y \in \mathbb{R} \to x \neq 0 \to y \neq 0 \to x \in \mathbb{R}) \to \varphi_{170}$
\item $\varphi_{213} :\equiv \forall v_{555}\,\forall v_{556}\,\forall v_{557}\,(v_{555} \in \mathbb{R} \to v_{556} \in \mathbb{R} \to \varphi_{30}) \to \varphi_{148}$
\item $\varphi_{214} :\equiv (y \in \mathbb{R} \to x \in \mathbb{R} \to y - x \in \mathbb{R}) \to (x \in \mathbb{R} \to y \in \mathbb{R} \to x \neq 0 \to y \neq 0 \to y \in \mathbb{R}) \to \varphi_{178}$
\item $\varphi_{215} :\equiv (x^{-1} = y \cdot (x^{-1} \cdot y^{-1}) \to \varphi_{133}) \to x \in \mathbb{R} \to y \in \mathbb{R} \to \varphi_{53}$
\item $\varphi_{216} :\equiv (x^{-1} \cdot y^{-1} = y^{-1} \cdot x^{-1} \to \varphi_{144}) \to x \in \mathbb{R} \to y \in \mathbb{R} \to \varphi_{31}$
\item $\varphi_{217} :\equiv x \cdot x^{-1} \cdot y^{-1} = x \cdot x^{-1} \cdot y^{-1} \to (x \cdot x^{-1} = 1 \to \varphi_{127}) \to \varphi_{102}$
\item $\varphi_{218} :\equiv (x \in \mathbb{R} \to y \in \mathbb{R} \to x \neq 0 \to y \neq 0 \to y^{-1} \cdot 1 = y^{-1}) \to \varphi_{135} \to \varphi_{74}$
\item $\varphi_{219} :\equiv y \cdot y^{-1} \cdot x^{-1} = y \cdot y^{-1} \cdot x^{-1} \to (y \cdot y^{-1} = 1 \to \varphi_{129}) \to \varphi_{105}$
\item $\varphi_{220} :\equiv (x \in \mathbb{R} \to y \in \mathbb{R} \to x \neq 0 \to y \neq 0 \to x^{-1} \cdot 1 = x^{-1}) \to \varphi_{136} \to \varphi_{76}$
\item $\varphi_{221} :\equiv (x \in \mathbb{R} \to y \in \mathbb{R} \to x \neq 0 \to y \neq 0 \to x^{-1} \cdot y^{-1} \in \mathbb{R}) \to \varphi_{199}$
\item $\varphi_{222} :\equiv (y \cdot (x^{-1} \cdot y^{-1}) = y \cdot (y^{-1} \cdot x^{-1}) \to \varphi_{143}) \to x \in \mathbb{R} \to \varphi_{93}$
\item $\varphi_{223} :\equiv x^{-1} - y^{-1} = x^{-1} \cdot y^{-1} \cdot y - x \cdot (x^{-1} \cdot y^{-1}) \to \varphi_{152} \to x^{-1} - y^{-1} = x^{-1} \cdot y^{-1} \cdot y - x^{-1} \cdot y^{-1} \cdot x$
\item $\varphi_{224} :\equiv (x^{-1} - y^{-1} = x^{-1} \cdot y^{-1} \cdot (y - x) \to \varphi_{165}) \to x \in \mathbb{R} \to y \in \mathbb{R} \to \varphi_{32}$
\item $\varphi_{225} :\equiv (x \in \mathbb{R} \to y \in \mathbb{R} \to x \neq 0 \to y \neq 0 \to 1 \cdot y^{-1} = y^{-1} \cdot 1) \to \varphi_{154} \to \varphi_{104}$
\item $\varphi_{226} :\equiv (x \in \mathbb{R} \to y \in \mathbb{R} \to x \neq 0 \to y \neq 0 \to 1 \cdot x^{-1} = x^{-1} \cdot 1) \to \varphi_{155} \to \varphi_{108}$
\item $\varphi_{227} :\equiv (x^{-1} \in \mathbb{R} \to y^{-1} \in \mathbb{R} \to x^{-1} \cdot y^{-1} \in \mathbb{R}) \to (x \in \mathbb{R} \to y \in \mathbb{R} \to x \neq 0 \to y \neq 0 \to x^{-1} \in \mathbb{R}) \to \varphi_{188}$
\item $\varphi_{228} :\equiv (x \in \mathbb{R} \to y \in \mathbb{R} \to x \neq 0 \to y \neq 0 \to x \in \mathbb{R}) \to (x \in \mathbb{R} \to y \in \mathbb{R} \to x \neq 0 \to y \neq 0 \to x^{-1} \in \mathbb{R}) \to \varphi_{193}$
\item $\varphi_{229} :\equiv (x \in \mathbb{R} \to y \in \mathbb{R} \to x \neq 0 \to y \neq 0 \to y \in \mathbb{R}) \to (x \in \mathbb{R} \to y \in \mathbb{R} \to x \neq 0 \to y \neq 0 \to y^{-1} \in \mathbb{R}) \to \varphi_{194}$
\item $\varphi_{230} :\equiv (x^{-1} \in \mathbb{R} \to y^{-1} \in \mathbb{R} \to x^{-1} \cdot y^{-1} = y^{-1} \cdot x^{-1}) \to \varphi_{208}$
\item $\varphi_{231} :\equiv \forall v_{589}\,\forall v_{590}\,\forall v_{591}\,(v_{589} \in \mathbb{R} \to v_{590} \in \mathbb{R} \to \varphi_{60}) \to \varphi_{202}$
\item $\varphi_{232} :\equiv (x \in \mathbb{R} \to y \in \mathbb{R} \to x \neq 0 \to y \neq 0 \to x^{-1} \cdot y^{-1} \in \mathbb{R}) \to (x \in \mathbb{R} \to y \in \mathbb{R} \to x \neq 0 \to y \neq 0 \to y \in \mathbb{R}) \to \varphi_{204}$
\item $\varphi_{233} :\equiv (x \in \mathbb{R} \to \varphi_{98}) \to (x \in \mathbb{R} \to y \in \mathbb{R} \to x \neq 0 \to y \neq 0 \to x \in \mathbb{R}) \to \varphi_{200}$
\item $\varphi_{234} :\equiv (y \in \mathbb{R} \to \varphi_{99}) \to (x \in \mathbb{R} \to y \in \mathbb{R} \to x \neq 0 \to y \neq 0 \to y \in \mathbb{R}) \to \varphi_{201}$
\item $\varphi_{235} :\equiv (y^{-1} = x \cdot (x^{-1} \cdot y^{-1}) \to \varphi_{56} \to \varphi_{119}) \to x \in \mathbb{R} \to y \in \mathbb{R} \to x \neq 0 \to y \neq 0 \to \varphi_{119}$
\item $\varphi_{236} :\equiv (x^{-1} - y^{-1} = y \cdot (x^{-1} \cdot y^{-1}) - y^{-1} \to \varphi_{119} \to \varphi_{57}) \to x \in \mathbb{R} \to y \in \mathbb{R} \to x \neq 0 \to y \neq 0 \to \varphi_{57}$
\item $\varphi_{237} :\equiv (x \in \mathbb{R} \to \varphi_{90}) \to \varphi_{102} \to (x \cdot (x^{-1} \cdot y^{-1}) = x \cdot x^{-1} \cdot y^{-1} \to \varphi_{83}) \to \varphi_{103}$
\item $\varphi_{238} :\equiv (x \in \mathbb{R} \to \varphi_{93}) \to \varphi_{105} \to (y \cdot (x^{-1} \cdot y^{-1}) = y \cdot y^{-1} \cdot x^{-1} \to \varphi_{85}) \to \varphi_{107}$
\item $\varphi_{239} :\equiv \forall z\,\forall w\,(z = w \to \varphi_{111} \to \varphi_{110}) \to \forall w\,(x \cdot (x^{-1} \cdot y^{-1}) = w \to \varphi_{150} \to \varphi_{110})$
\item $\varphi_{240} :\equiv \forall z\,\forall w\,(z = w \to \varphi_{113} \to \varphi_{112}) \to \forall w\,(y \cdot (x^{-1} \cdot y^{-1}) = w \to \varphi_{151} \to \varphi_{112})$
\item $\varphi_{241} :\equiv (x \cdot (x^{-1} \cdot y^{-1}) = x^{-1} \cdot y^{-1} \cdot x \to \varphi_{150} \to \varphi_{152}) \to x \in \mathbb{R} \to y \in \mathbb{R} \to x \neq 0 \to y \neq 0 \to \varphi_{152}$
\item $\varphi_{242} :\equiv (y \cdot (x^{-1} \cdot y^{-1}) = x^{-1} \cdot y^{-1} \cdot y \to \varphi_{151} \to \varphi_{153}) \to x \in \mathbb{R} \to y \in \mathbb{R} \to x \neq 0 \to y \neq 0 \to \varphi_{153}$
\item $\varphi_{243} :\equiv \forall w\,(x \cdot (x^{-1} \cdot y^{-1}) = w \to \varphi_{150} \to \varphi_{110}) \to x \cdot (x^{-1} \cdot y^{-1}) = x^{-1} \cdot y^{-1} \cdot x \to \varphi_{150} \to \varphi_{152}$
\item $\varphi_{244} :\equiv \forall w\,(y \cdot (x^{-1} \cdot y^{-1}) = w \to \varphi_{151} \to \varphi_{112}) \to y \cdot (x^{-1} \cdot y^{-1}) = x^{-1} \cdot y^{-1} \cdot y \to \varphi_{151} \to \varphi_{153}$
\item $\varphi_{245} :\equiv (x \in \mathbb{R} \to y \in \mathbb{R} \to x \neq 0 \to y \neq 0 \to \varphi_{153}) \to (\varphi_{57} \to \varphi_{203}) \to x \in \mathbb{R} \to y \in \mathbb{R} \to x \neq 0 \to \varphi_{94}$
\item $\varphi_{246} :\equiv (x \in \mathbb{R} \to y \in \mathbb{R} \to x \neq 0 \to y \neq 0 \to \varphi_{152}) \to \varphi_{223} \to x \in \mathbb{R} \to y \in \mathbb{R} \to x \neq 0 \to \varphi_{95}$
\item $\varphi_{247} :\equiv (x \in \mathbb{R} \to y \in \mathbb{R} \to x \neq 0 \to \varphi_{95}) \to (x \in \mathbb{R} \to y \in \mathbb{R} \to x \neq 0 \to y \neq 0 \to \varphi_{109}) \to \varphi_{205} \to x \in \mathbb{R} \to \varphi_{96}$
\item $\varphi_{248} :\equiv (x \in \mathbb{R} \to y \in \mathbb{R} \to \varphi_{53}) \to (x \in \mathbb{R} \to y \in \mathbb{R} \to x \neq 0 \to y \neq 0 \to \varphi_{119}) \to \varphi_{236}$
\end{abbrevlist}
\end{fullproof}

% 自動生成 (prover/build.js): InvMulL  ノード数 33
\begin{fullproof}{[E22]}{$\forall x \in \mathbb{R}\,(x \neq 0 \to x^{-1} \cdot x = 1)$}{33}\phantomsection\label{tree:InvMulL}
\noindent$\Pi_{1}$:
\begin{ptree}
\AXC{$y = z \to x = y \to x = z$\lbl{A6}}
\AXC{$(y = z \to x = y \to x = z) \to x = y \to y = z \to x = z$\lbl{A1}}
\BIC{$x = y \to y = z \to x = z$}
\UIC{$\forall z\,(x = y \to y = z \to x = z)$}
\UIC{$\forall y\,\forall z\,(x = y \to y = z \to x = z)$}
\UIC{$\forall x\,\forall y\,\forall z\,(x = y \to y = z \to x = z)$}
\AXC{$\varphi_{2}$\lbl{A2}}
\BIC{$\forall y\,\forall z\,(x^{-1} \cdot x = y \to y = z \to x^{-1} \cdot x = z)$}
\end{ptree}
\noindent$\Pi_{2}$:
\begin{ptree}
\AXC{$\stackrel{\Pi_{1}}{\vdots}$}\noLine
\UIC{$\forall y\,\forall z\,(x^{-1} \cdot x = y \to y = z \to x^{-1} \cdot x = z)$}
\AXC{$\forall y\,\forall z\,(x^{-1} \cdot x = y \to y = z \to x^{-1} \cdot x = z) \to \varphi_{1}$\lbl{A2}}
\BIC{$\varphi_{1}$}
\end{ptree}
\noindent$\Pi_{3}$:
\begin{ptree}
\AXC{$\forall x\,\forall y\,(x \in \mathbb{R} \to y \in \mathbb{R} \to x \cdot y = y \cdot x)$\lbl{(R7)}}
\AXC{$\varphi_{7}$\lbl{A2}}
\BIC{$\forall y\,(x^{-1} \in \mathbb{R} \to y \in \mathbb{R} \to x^{-1} \cdot y = y \cdot x^{-1})$}
\AXC{$\varphi_{8}$\lbl{A2}}
\BIC{$x^{-1} \in \mathbb{R} \to x \in \mathbb{R} \to x^{-1} \cdot x = x \cdot x^{-1}$}
\AXC{$\varphi_{9}$\lbl{A1}}
\BIC{$(x \in \mathbb{R} \to x \neq 0 \to x^{-1} \in \mathbb{R}) \to \varphi_{5}$}
\end{ptree}
\noindent$\Pi_{4}$:
\begin{ptree}
\AXC{$x \in \mathbb{R} \to x \neq 0 \to x \in \mathbb{R}$\lbl{A1}}
\AXC{$\forall x \in \mathbb{R}\,(x \neq 0 \to x^{-1} \in \mathbb{R})$\lbl{(R1)}}
\AXC{$\varphi_{4}$\lbl{A2}}
\BIC{$x \in \mathbb{R} \to x \neq 0 \to x^{-1} \in \mathbb{R}$}
\AXC{$\stackrel{\Pi_{3}}{\vdots}$}\noLine
\UIC{$(x \in \mathbb{R} \to x \neq 0 \to x^{-1} \in \mathbb{R}) \to \varphi_{5}$}
\BIC{$\varphi_{5}$}
\BIC{$x \in \mathbb{R} \to x \neq 0 \to x^{-1} \cdot x = x \cdot x^{-1}$}
\end{ptree}
\noindent$\Pi_{5}$:
\begin{ptree}
\AXC{$\forall x \in \mathbb{R}\,(x \neq 0 \to x \cdot x^{-1} = 1)$\lbl{(R9)}}
\AXC{$\varphi_{3}$\lbl{A2}}
\BIC{$x \in \mathbb{R} \to x \neq 0 \to x \cdot x^{-1} = 1$}
\AXC{$\stackrel{\Pi_{4}}{\vdots}$}\noLine
\UIC{$x \in \mathbb{R} \to x \neq 0 \to x^{-1} \cdot x = x \cdot x^{-1}$}
\AXC{$\varphi_{10}$\lbl{A1}}
\BIC{$(x \in \mathbb{R} \to x \neq 0 \to x \cdot x^{-1} = 1) \to \varphi_{6}$}
\BIC{$\varphi_{6}$}
\end{ptree}
\noindent 主木:
\begin{ptree}
\AXC{$\stackrel{\Pi_{2}}{\vdots}$}\noLine
\UIC{$\varphi_{1}$}
\AXC{$\varphi_{1} \to x^{-1} \cdot x = x \cdot x^{-1} \to x \cdot x^{-1} = 1 \to x^{-1} \cdot x = 1$\lbl{A2}}
\BIC{$x^{-1} \cdot x = x \cdot x^{-1} \to x \cdot x^{-1} = 1 \to x^{-1} \cdot x = 1$}
\AXC{$\stackrel{\Pi_{5}}{\vdots}$}\noLine
\UIC{$\varphi_{6}$}
\BIC{$x \in \mathbb{R} \to x \neq 0 \to x^{-1} \cdot x = 1$}
\UIC{$\forall x \in \mathbb{R}\,(x \neq 0 \to x^{-1} \cdot x = 1)$}
\end{ptree}
\begin{abbrevlist}
\item $\varphi_{1} :\equiv \forall z\,(x^{-1} \cdot x = x \cdot x^{-1} \to x \cdot x^{-1} = z \to x^{-1} \cdot x = z)$
\item $\varphi_{2} :\equiv \forall x\,\forall y\,\forall z\,(x = y \to y = z \to x = z) \to \forall y\,\forall z\,(x^{-1} \cdot x = y \to y = z \to x^{-1} \cdot x = z)$
\item $\varphi_{3} :\equiv \forall x \in \mathbb{R}\,(x \neq 0 \to x \cdot x^{-1} = 1) \to x \in \mathbb{R} \to x \neq 0 \to x \cdot x^{-1} = 1$
\item $\varphi_{4} :\equiv \forall x \in \mathbb{R}\,(x \neq 0 \to x^{-1} \in \mathbb{R}) \to x \in \mathbb{R} \to x \neq 0 \to x^{-1} \in \mathbb{R}$
\item $\varphi_{5} :\equiv (x \in \mathbb{R} \to x \neq 0 \to x \in \mathbb{R}) \to x \in \mathbb{R} \to x \neq 0 \to x^{-1} \cdot x = x \cdot x^{-1}$
\item $\varphi_{6} :\equiv (x^{-1} \cdot x = x \cdot x^{-1} \to x \cdot x^{-1} = 1 \to x^{-1} \cdot x = 1) \to x \in \mathbb{R} \to x \neq 0 \to x^{-1} \cdot x = 1$
\item $\varphi_{7} :\equiv \forall x\,\forall y\,(x \in \mathbb{R} \to y \in \mathbb{R} \to x \cdot y = y \cdot x) \to \forall y\,(x^{-1} \in \mathbb{R} \to y \in \mathbb{R} \to x^{-1} \cdot y = y \cdot x^{-1})$
\item $\varphi_{8} :\equiv \forall y\,(x^{-1} \in \mathbb{R} \to y \in \mathbb{R} \to x^{-1} \cdot y = y \cdot x^{-1}) \to x^{-1} \in \mathbb{R} \to x \in \mathbb{R} \to x^{-1} \cdot x = x \cdot x^{-1}$
\item $\varphi_{9} :\equiv (x^{-1} \in \mathbb{R} \to x \in \mathbb{R} \to x^{-1} \cdot x = x \cdot x^{-1}) \to (x \in \mathbb{R} \to x \neq 0 \to x^{-1} \in \mathbb{R}) \to \varphi_{5}$
\item $\varphi_{10} :\equiv (x \in \mathbb{R} \to x \neq 0 \to x^{-1} \cdot x = x \cdot x^{-1}) \to (x \in \mathbb{R} \to x \neq 0 \to x \cdot x^{-1} = 1) \to \varphi_{6}$
\end{abbrevlist}
\end{fullproof}

% 自動生成 (prover/build.js): DivMul  ノード数 99
\begin{fullproof}{[E23]}{$\forall x\,\forall y\,(x \in \mathbb{R} \to y \in \mathbb{R} \to y \neq 0 \to x \cdot y^{-1} \cdot y = x)$}{99}\phantomsection\label{tree:DivMul}
\noindent$\Pi_{1}$:
\begin{ptree}
\AXC{$y = z \to x = y \to x = z$\lbl{A6}}
\AXC{$(y = z \to x = y \to x = z) \to x = y \to y = z \to x = z$\lbl{A1}}
\BIC{$x = y \to y = z \to x = z$}
\UIC{$\forall z\,(x = y \to y = z \to x = z)$}
\UIC{$\forall y\,\forall z\,(x = y \to y = z \to x = z)$}
\UIC{$\forall x\,\forall y\,\forall z\,(x = y \to y = z \to x = z)$}
\AXC{$\varphi_{3}$\lbl{A2}}
\BIC{$\forall y\,\forall z\,(v_{614} = y \to y = z \to v_{614} = z)$}
\AXC{$\varphi_{7}$\lbl{A2}}
\BIC{$\forall z\,(v_{614} = v_{615} \to v_{615} = z \to v_{614} = z)$}
\AXC{$\varphi_{9}$\lbl{A2}}
\BIC{$v_{614} = v_{615} \to v_{615} = v_{616} \to v_{614} = v_{616}$}
\UIC{$\forall v_{616}\,(v_{614} = v_{615} \to v_{615} = v_{616} \to v_{614} = v_{616})$}
\UIC{$\varphi_{1}$}
\UIC{$\forall v_{614}\,(\varphi_{1})$}
\end{ptree}
\noindent$\Pi_{2}$:
\begin{ptree}
\AXC{$\stackrel{\Pi_{1}}{\vdots}$}\noLine
\UIC{$\forall v_{614}\,(\varphi_{1})$}
\AXC{$\forall v_{614}\,(\varphi_{1}) \to \varphi_{14}$\lbl{A2}}
\BIC{$\varphi_{14}$}
\AXC{$\varphi_{14} \to \varphi_{12}$\lbl{A2}}
\BIC{$\varphi_{12}$}
\AXC{$\varphi_{12} \to x \cdot y^{-1} \cdot y = x \cdot 1 \to x \cdot 1 = x \to x \cdot y^{-1} \cdot y = x$\lbl{A2}}
\BIC{$x \cdot y^{-1} \cdot y = x \cdot 1 \to x \cdot 1 = x \to x \cdot y^{-1} \cdot y = x$}
\end{ptree}
\noindent$\Pi_{3}$:
\begin{ptree}
\AXC{$\forall x \in \mathbb{R}\,(x \cdot 1 = x)$\lbl{(R8)}}
\AXC{$\forall x \in \mathbb{R}\,(x \cdot 1 = x) \to x \in \mathbb{R} \to x \cdot 1 = x$\lbl{A2}}
\BIC{$x \in \mathbb{R} \to x \cdot 1 = x$}
\AXC{$(x \in \mathbb{R} \to x \cdot 1 = x) \to \varphi_{23}$\lbl{A1}}
\BIC{$\varphi_{23}$}
\end{ptree}
\noindent$\Pi_{4}$:
\begin{ptree}
\AXC{$y = z \to x = y \to x = z$\lbl{A6}}
\AXC{$(y = z \to x = y \to x = z) \to x = y \to y = z \to x = z$\lbl{A1}}
\BIC{$x = y \to y = z \to x = z$}
\UIC{$\forall z\,(x = y \to y = z \to x = z)$}
\UIC{$\forall y\,\forall z\,(x = y \to y = z \to x = z)$}
\UIC{$\forall x\,\forall y\,\forall z\,(x = y \to y = z \to x = z)$}
\AXC{$\varphi_{4}$\lbl{A2}}
\BIC{$\forall y\,\forall z\,(v_{611} = y \to y = z \to v_{611} = z)$}
\AXC{$\varphi_{8}$\lbl{A2}}
\BIC{$\forall z\,(v_{611} = v_{612} \to v_{612} = z \to v_{611} = z)$}
\AXC{$\varphi_{10}$\lbl{A2}}
\BIC{$v_{611} = v_{612} \to v_{612} = v_{613} \to v_{611} = v_{613}$}
\UIC{$\forall v_{613}\,(v_{611} = v_{612} \to v_{612} = v_{613} \to v_{611} = v_{613})$}
\UIC{$\varphi_{2}$}
\UIC{$\forall v_{611}\,(\varphi_{2})$}
\end{ptree}
\noindent$\Pi_{5}$:
\begin{ptree}
\AXC{$\stackrel{\Pi_{4}}{\vdots}$}\noLine
\UIC{$\forall v_{611}\,(\varphi_{2})$}
\AXC{$\forall v_{611}\,(\varphi_{2}) \to \varphi_{15}$\lbl{A2}}
\BIC{$\varphi_{15}$}
\AXC{$\varphi_{15} \to \varphi_{24}$\lbl{A2}}
\BIC{$\varphi_{24}$}
\AXC{$\varphi_{24} \to \varphi_{21}$\lbl{A2}}
\BIC{$\varphi_{21}$}
\end{ptree}
\noindent$\Pi_{6}$:
\begin{ptree}
\AXC{$x = x$\lbl{A5}}
\UIC{$\forall x\,(x = x)$}
\AXC{$\forall x\,(x = x) \to x \cdot (y^{-1} \cdot y) = x \cdot (y^{-1} \cdot y)$\lbl{A2}}
\BIC{$x \cdot (y^{-1} \cdot y) = x \cdot (y^{-1} \cdot y)$}
\end{ptree}
\noindent$\Pi_{7}$:
\begin{ptree}
\AXC{$x \in \mathbb{R} \to y \in \mathbb{R} \to y \neq 0 \to y \in \mathbb{R}$\lbl{A1}}
\AXC{$\forall x \in \mathbb{R}\,(x \neq 0 \to x^{-1} \cdot x = 1)$\lbl{[E22]}}
\AXC{$\varphi_{16}$\lbl{A2}}
\BIC{$y \in \mathbb{R} \to y \neq 0 \to y^{-1} \cdot y = 1$}
\AXC{$\varphi_{29}$\lbl{A1}}
\BIC{$(x \in \mathbb{R} \to y \in \mathbb{R} \to y \neq 0 \to y \in \mathbb{R}) \to \varphi_{22}$}
\BIC{$\varphi_{22}$}
\end{ptree}
\noindent$\Pi_{8}$:
\begin{ptree}
\AXC{$\stackrel{\Pi_{6}}{\vdots}$}\noLine
\UIC{$x \cdot (y^{-1} \cdot y) = x \cdot (y^{-1} \cdot y)$}
\AXC{$x \in \mathbb{R} \to y \in \mathbb{R} \to y \neq 0 \to y \neq 0$\lbl{A1}}
\AXC{$\stackrel{\Pi_{7}}{\vdots}$}\noLine
\UIC{$\varphi_{22}$}
\BIC{$x \in \mathbb{R} \to y \in \mathbb{R} \to y \neq 0 \to y^{-1} \cdot y = 1$}
\AXC{$\varphi_{33}$\lbl{A1}}
\BIC{$x \cdot (y^{-1} \cdot y) = x \cdot (y^{-1} \cdot y) \to \varphi_{27}$}
\BIC{$\varphi_{27}$}
\end{ptree}
\noindent$\Pi_{9}$:
\begin{ptree}
\AXC{$x \in \mathbb{R} \to y \in \mathbb{R} \to y \neq 0 \to y \in \mathbb{R}$\lbl{A1}}
\AXC{$\forall x \in \mathbb{R}\,(x \neq 0 \to x^{-1} \in \mathbb{R})$\lbl{(R1)}}
\AXC{$\varphi_{17}$\lbl{A2}}
\BIC{$y \in \mathbb{R} \to y \neq 0 \to y^{-1} \in \mathbb{R}$}
\AXC{$\varphi_{30}$\lbl{A1}}
\BIC{$(x \in \mathbb{R} \to y \in \mathbb{R} \to y \neq 0 \to y \in \mathbb{R}) \to \varphi_{25}$}
\BIC{$\varphi_{25}$}
\end{ptree}
\noindent$\Pi_{10}$:
\begin{ptree}
\AXC{$\forall x\,(\varphi_{13})$\lbl{(R6)}}
\AXC{$\forall x\,(\varphi_{13}) \to \varphi_{13}$\lbl{A2}}
\BIC{$\varphi_{13}$}
\AXC{$\varphi_{13} \to \varphi_{20}$\lbl{A2}}
\BIC{$\varphi_{20}$}
\AXC{$\varphi_{20} \to \varphi_{18}$\lbl{A2}}
\BIC{$\varphi_{18}$}
\AXC{$\varphi_{18} \to (x \in \mathbb{R} \to y \in \mathbb{R} \to y \neq 0 \to x \in \mathbb{R}) \to \varphi_{28}$\lbl{A1}}
\BIC{$(x \in \mathbb{R} \to y \in \mathbb{R} \to y \neq 0 \to x \in \mathbb{R}) \to \varphi_{28}$}
\end{ptree}
\noindent$\Pi_{11}$:
\begin{ptree}
\AXC{$x \in \mathbb{R} \to y \in \mathbb{R} \to y \neq 0 \to x \in \mathbb{R}$\lbl{A1}}
\AXC{$\stackrel{\Pi_{10}}{\vdots}$}\noLine
\UIC{$(x \in \mathbb{R} \to y \in \mathbb{R} \to y \neq 0 \to x \in \mathbb{R}) \to \varphi_{28}$}
\BIC{$\varphi_{28}$}
\end{ptree}
\noindent$\Pi_{12}$:
\begin{ptree}
\AXC{$x \in \mathbb{R} \to y \in \mathbb{R} \to y \neq 0 \to y \in \mathbb{R}$\lbl{A1}}
\AXC{$x \in \mathbb{R} \to y \in \mathbb{R} \to y \neq 0 \to y \neq 0$\lbl{A1}}
\AXC{$\stackrel{\Pi_{9}}{\vdots}$}\noLine
\UIC{$\varphi_{25}$}
\BIC{$x \in \mathbb{R} \to y \in \mathbb{R} \to y \neq 0 \to y^{-1} \in \mathbb{R}$}
\AXC{$\stackrel{\Pi_{11}}{\vdots}$}\noLine
\UIC{$\varphi_{28}$}
\BIC{$(x \in \mathbb{R} \to y \in \mathbb{R} \to y \neq 0 \to y \in \mathbb{R}) \to \varphi_{11}$}
\BIC{$\varphi_{11}$}
\AXC{$\varphi_{11} \to \varphi_{31}$\lbl{A1}}
\BIC{$\varphi_{31}$}
\end{ptree}
\noindent$\Pi_{13}$:
\begin{ptree}
\AXC{$\varphi_{5}$\lbl{A6}}
\UIC{$\forall w\,(\varphi_{5})$}
\UIC{$\forall z\,\forall w\,(\varphi_{5})$}
\AXC{$\forall z\,\forall w\,(\varphi_{5}) \to \varphi_{19}$\lbl{A2}}
\BIC{$\varphi_{19}$}
\AXC{$\varphi_{19} \to y^{-1} \cdot y = 1 \to \varphi_{6}$\lbl{A2}}
\BIC{$y^{-1} \cdot y = 1 \to \varphi_{6}$}
\AXC{$\stackrel{\Pi_{8}}{\vdots}$}\noLine
\UIC{$\varphi_{27}$}
\BIC{$x \in \mathbb{R} \to y \in \mathbb{R} \to y \neq 0 \to x \cdot (y^{-1} \cdot y) = x \cdot 1$}
\AXC{$\stackrel{\Pi_{12}}{\vdots}$}\noLine
\UIC{$\varphi_{31}$}
\BIC{$\varphi_{21} \to x \in \mathbb{R} \to y \in \mathbb{R} \to y \neq 0 \to x \cdot y^{-1} \cdot y = x \cdot 1$}
\end{ptree}
\noindent$\Pi_{14}$:
\begin{ptree}
\AXC{$\stackrel{\Pi_{5}}{\vdots}$}\noLine
\UIC{$\varphi_{21}$}
\AXC{$\stackrel{\Pi_{13}}{\vdots}$}\noLine
\UIC{$\varphi_{21} \to x \in \mathbb{R} \to y \in \mathbb{R} \to y \neq 0 \to x \cdot y^{-1} \cdot y = x \cdot 1$}
\BIC{$x \in \mathbb{R} \to y \in \mathbb{R} \to y \neq 0 \to x \cdot y^{-1} \cdot y = x \cdot 1$}
\AXC{$\varphi_{32}$\lbl{A1}}
\BIC{$(x \in \mathbb{R} \to y \in \mathbb{R} \to y \neq 0 \to x \cdot 1 = x) \to \varphi_{26}$}
\end{ptree}
\noindent$\Pi_{15}$:
\begin{ptree}
\AXC{$x \in \mathbb{R} \to y \in \mathbb{R} \to y \neq 0 \to x \in \mathbb{R}$\lbl{A1}}
\AXC{$\stackrel{\Pi_{3}}{\vdots}$}\noLine
\UIC{$\varphi_{23}$}
\BIC{$x \in \mathbb{R} \to y \in \mathbb{R} \to y \neq 0 \to x \cdot 1 = x$}
\AXC{$\stackrel{\Pi_{14}}{\vdots}$}\noLine
\UIC{$(x \in \mathbb{R} \to y \in \mathbb{R} \to y \neq 0 \to x \cdot 1 = x) \to \varphi_{26}$}
\BIC{$\varphi_{26}$}
\end{ptree}
\noindent 主木:
\begin{ptree}
\AXC{$\stackrel{\Pi_{2}}{\vdots}$}\noLine
\UIC{$x \cdot y^{-1} \cdot y = x \cdot 1 \to x \cdot 1 = x \to x \cdot y^{-1} \cdot y = x$}
\AXC{$\stackrel{\Pi_{15}}{\vdots}$}\noLine
\UIC{$\varphi_{26}$}
\BIC{$x \in \mathbb{R} \to y \in \mathbb{R} \to y \neq 0 \to x \cdot y^{-1} \cdot y = x$}
\UIC{$\forall y\,(x \in \mathbb{R} \to y \in \mathbb{R} \to y \neq 0 \to x \cdot y^{-1} \cdot y = x)$}
\UIC{$\forall x\,\forall y\,(x \in \mathbb{R} \to y \in \mathbb{R} \to y \neq 0 \to x \cdot y^{-1} \cdot y = x)$}
\end{ptree}
\begin{abbrevlist}
\item $\varphi_{1} :\equiv \forall v_{615}\,\forall v_{616}\,(v_{614} = v_{615} \to v_{615} = v_{616} \to v_{614} = v_{616})$
\item $\varphi_{2} :\equiv \forall v_{612}\,\forall v_{613}\,(v_{611} = v_{612} \to v_{612} = v_{613} \to v_{611} = v_{613})$
\item $\varphi_{3} :\equiv \forall x\,\forall y\,\forall z\,(x = y \to y = z \to x = z) \to \forall y\,\forall z\,(v_{614} = y \to y = z \to v_{614} = z)$
\item $\varphi_{4} :\equiv \forall x\,\forall y\,\forall z\,(x = y \to y = z \to x = z) \to \forall y\,\forall z\,(v_{611} = y \to y = z \to v_{611} = z)$
\item $\varphi_{5} :\equiv z = w \to x \cdot (y^{-1} \cdot y) = x \cdot z \to x \cdot (y^{-1} \cdot y) = x \cdot w$
\item $\varphi_{6} :\equiv x \cdot (y^{-1} \cdot y) = x \cdot (y^{-1} \cdot y) \to x \cdot (y^{-1} \cdot y) = x \cdot 1$
\item $\varphi_{7} :\equiv \forall y\,\forall z\,(v_{614} = y \to y = z \to v_{614} = z) \to \forall z\,(v_{614} = v_{615} \to v_{615} = z \to v_{614} = z)$
\item $\varphi_{8} :\equiv \forall y\,\forall z\,(v_{611} = y \to y = z \to v_{611} = z) \to \forall z\,(v_{611} = v_{612} \to v_{612} = z \to v_{611} = z)$
\item $\varphi_{9} :\equiv \forall z\,(v_{614} = v_{615} \to v_{615} = z \to v_{614} = z) \to v_{614} = v_{615} \to v_{615} = v_{616} \to v_{614} = v_{616}$
\item $\varphi_{10} :\equiv \forall z\,(v_{611} = v_{612} \to v_{612} = z \to v_{611} = z) \to v_{611} = v_{612} \to v_{612} = v_{613} \to v_{611} = v_{613}$
\item $\varphi_{11} :\equiv x \in \mathbb{R} \to y \in \mathbb{R} \to y \neq 0 \to x \cdot y^{-1} \cdot y = x \cdot (y^{-1} \cdot y)$
\item $\varphi_{12} :\equiv \forall v_{616}\,(x \cdot y^{-1} \cdot y = x \cdot 1 \to x \cdot 1 = v_{616} \to x \cdot y^{-1} \cdot y = v_{616})$
\item $\varphi_{13} :\equiv \forall y\,\forall z\,(x \in \mathbb{R} \to y \in \mathbb{R} \to z \in \mathbb{R} \to x \cdot y \cdot z = x \cdot (y \cdot z))$
\item $\varphi_{14} :\equiv \forall v_{615}\,\forall v_{616}\,(x \cdot y^{-1} \cdot y = v_{615} \to v_{615} = v_{616} \to x \cdot y^{-1} \cdot y = v_{616})$
\item $\varphi_{15} :\equiv \forall v_{612}\,\forall v_{613}\,(x \cdot y^{-1} \cdot y = v_{612} \to v_{612} = v_{613} \to x \cdot y^{-1} \cdot y = v_{613})$
\item $\varphi_{16} :\equiv \forall x \in \mathbb{R}\,(x \neq 0 \to x^{-1} \cdot x = 1) \to y \in \mathbb{R} \to y \neq 0 \to y^{-1} \cdot y = 1$
\item $\varphi_{17} :\equiv \forall x \in \mathbb{R}\,(x \neq 0 \to x^{-1} \in \mathbb{R}) \to y \in \mathbb{R} \to y \neq 0 \to y^{-1} \in \mathbb{R}$
\item $\varphi_{18} :\equiv x \in \mathbb{R} \to y^{-1} \in \mathbb{R} \to y \in \mathbb{R} \to x \cdot y^{-1} \cdot y = x \cdot (y^{-1} \cdot y)$
\item $\varphi_{19} :\equiv \forall w\,(y^{-1} \cdot y = w \to x \cdot (y^{-1} \cdot y) = x \cdot (y^{-1} \cdot y) \to x \cdot (y^{-1} \cdot y) = x \cdot w)$
\item $\varphi_{20} :\equiv \forall z\,(x \in \mathbb{R} \to y^{-1} \in \mathbb{R} \to z \in \mathbb{R} \to x \cdot y^{-1} \cdot z = x \cdot (y^{-1} \cdot z))$
\item $\varphi_{21} :\equiv x \cdot y^{-1} \cdot y = x \cdot (y^{-1} \cdot y) \to x \cdot (y^{-1} \cdot y) = x \cdot 1 \to x \cdot y^{-1} \cdot y = x \cdot 1$
\item $\varphi_{22} :\equiv (x \in \mathbb{R} \to y \in \mathbb{R} \to y \neq 0 \to y \neq 0) \to x \in \mathbb{R} \to y \in \mathbb{R} \to y \neq 0 \to y^{-1} \cdot y = 1$
\item $\varphi_{23} :\equiv (x \in \mathbb{R} \to y \in \mathbb{R} \to y \neq 0 \to x \in \mathbb{R}) \to x \in \mathbb{R} \to y \in \mathbb{R} \to y \neq 0 \to x \cdot 1 = x$
\item $\varphi_{24} :\equiv \forall v_{613}\,(x \cdot y^{-1} \cdot y = x \cdot (y^{-1} \cdot y) \to x \cdot (y^{-1} \cdot y) = v_{613} \to x \cdot y^{-1} \cdot y = v_{613})$
\item $\varphi_{25} :\equiv (x \in \mathbb{R} \to y \in \mathbb{R} \to y \neq 0 \to y \neq 0) \to x \in \mathbb{R} \to y \in \mathbb{R} \to y \neq 0 \to y^{-1} \in \mathbb{R}$
\item $\varphi_{26} :\equiv (x \cdot y^{-1} \cdot y = x \cdot 1 \to x \cdot 1 = x \to x \cdot y^{-1} \cdot y = x) \to x \in \mathbb{R} \to y \in \mathbb{R} \to y \neq 0 \to x \cdot y^{-1} \cdot y = x$
\item $\varphi_{27} :\equiv (y^{-1} \cdot y = 1 \to \varphi_{6}) \to x \in \mathbb{R} \to y \in \mathbb{R} \to y \neq 0 \to x \cdot (y^{-1} \cdot y) = x \cdot 1$
\item $\varphi_{28} :\equiv (x \in \mathbb{R} \to y \in \mathbb{R} \to y \neq 0 \to y^{-1} \in \mathbb{R}) \to (x \in \mathbb{R} \to y \in \mathbb{R} \to y \neq 0 \to y \in \mathbb{R}) \to \varphi_{11}$
\item $\varphi_{29} :\equiv (y \in \mathbb{R} \to y \neq 0 \to y^{-1} \cdot y = 1) \to (x \in \mathbb{R} \to y \in \mathbb{R} \to y \neq 0 \to y \in \mathbb{R}) \to \varphi_{22}$
\item $\varphi_{30} :\equiv (y \in \mathbb{R} \to y \neq 0 \to y^{-1} \in \mathbb{R}) \to (x \in \mathbb{R} \to y \in \mathbb{R} \to y \neq 0 \to y \in \mathbb{R}) \to \varphi_{25}$
\item $\varphi_{31} :\equiv (x \in \mathbb{R} \to y \in \mathbb{R} \to y \neq 0 \to x \cdot (y^{-1} \cdot y) = x \cdot 1) \to \varphi_{21} \to x \in \mathbb{R} \to y \in \mathbb{R} \to y \neq 0 \to x \cdot y^{-1} \cdot y = x \cdot 1$
\item $\varphi_{32} :\equiv (x \in \mathbb{R} \to y \in \mathbb{R} \to y \neq 0 \to x \cdot y^{-1} \cdot y = x \cdot 1) \to (x \in \mathbb{R} \to y \in \mathbb{R} \to y \neq 0 \to x \cdot 1 = x) \to \varphi_{26}$
\item $\varphi_{33} :\equiv (x \in \mathbb{R} \to y \in \mathbb{R} \to y \neq 0 \to y^{-1} \cdot y = 1) \to x \cdot (y^{-1} \cdot y) = x \cdot (y^{-1} \cdot y) \to \varphi_{27}$
\end{abbrevlist}
\end{fullproof}

% 自動生成 (prover/build.js): LtDiv  ノード数 192
\begin{fullproof}{[E24]}{$\forall x\,\forall y\,\forall z\,(x \in \mathbb{R} \to y \in \mathbb{R} \to z \in \mathbb{R} \to 0 < y \to x < z \cdot y \to x \cdot y^{-1} < z)$}{192}\phantomsection\label{tree:LtDiv}
\noindent$\Pi_{1}$:
\begin{ptree}
\AXC{$z_{1} = w \to x \cdot y^{-1} < z_{1} \to x \cdot y^{-1} < w$\lbl{A6}}
\UIC{$\forall w\,(z_{1} = w \to x \cdot y^{-1} < z_{1} \to x \cdot y^{-1} < w)$}
\UIC{$\forall z_{1}\,\forall w\,(z_{1} = w \to x \cdot y^{-1} < z_{1} \to x \cdot y^{-1} < w)$}
\AXC{$\forall z_{1}\,\forall w\,(z_{1} = w \to x \cdot y^{-1} < z_{1} \to x \cdot y^{-1} < w) \to \varphi_{25}$\lbl{A2}}
\BIC{$\varphi_{25}$}
\AXC{$\varphi_{25} \to \varphi_{18}$\lbl{A2}}
\BIC{$\varphi_{18}$}
\end{ptree}
\noindent$\Pi_{2}$:
\begin{ptree}
\AXC{$\forall x \in \mathbb{R}\,(0 < x \to 0 < x^{-1})$\lbl{[F34]}}
\AXC{$\varphi_{21}$\lbl{A2}}
\BIC{$y \in \mathbb{R} \to 0 < y \to 0 < y^{-1}$}
\AXC{$(y \in \mathbb{R} \to 0 < y \to 0 < y^{-1}) \to \varphi_{62}$\lbl{A1}}
\BIC{$\varphi_{62}$}
\end{ptree}
\noindent$\Pi_{3}$:
\begin{ptree}
\AXC{$x \in \mathbb{R} \to y \in \mathbb{R} \to z \in \mathbb{R} \to 0 < y \to x < z \cdot y \to 0 < y$\lbl{A1}}
\AXC{$x \in \mathbb{R} \to y \in \mathbb{R} \to z \in \mathbb{R} \to 0 < y \to x < z \cdot y \to y \in \mathbb{R}$\lbl{A1}}
\AXC{$\stackrel{\Pi_{2}}{\vdots}$}\noLine
\UIC{$\varphi_{62}$}
\BIC{$\varphi_{52}$}
\BIC{$x \in \mathbb{R} \to y \in \mathbb{R} \to z \in \mathbb{R} \to 0 < y \to x < z \cdot y \to 0 < y^{-1}$}
\end{ptree}
\noindent$\Pi_{4}$:
\begin{ptree}
\AXC{$x = x$\lbl{A5}}
\AXC{$x = y \to x = x \to y = x$\lbl{A6}}
\AXC{$(x = y \to x = x \to y = x) \to x = x \to x = y \to y = x$\lbl{A1}}
\BIC{$x = x \to x = y \to y = x$}
\BIC{$x = y \to y = x$}
\UIC{$\forall y\,(x = y \to y = x)$}
\UIC{$\forall x\,\forall y\,(x = y \to y = x)$}
\end{ptree}
\noindent$\Pi_{5}$:
\begin{ptree}
\AXC{$\stackrel{\Pi_{4}}{\vdots}$}\noLine
\UIC{$\forall x\,\forall y\,(x = y \to y = x)$}
\AXC{$\forall x\,\forall y\,(x = y \to y = x) \to \forall y\,(v_{617} = y \to y = v_{617})$\lbl{A2}}
\BIC{$\forall y\,(v_{617} = y \to y = v_{617})$}
\AXC{$\varphi_{1}$\lbl{A2}}
\BIC{$v_{617} = v_{618} \to v_{618} = v_{617}$}
\UIC{$\forall v_{618}\,(v_{617} = v_{618} \to v_{618} = v_{617})$}
\UIC{$\forall v_{617}\,\forall v_{618}\,(v_{617} = v_{618} \to v_{618} = v_{617})$}
\AXC{$\varphi_{6}$\lbl{A2}}
\BIC{$\forall v_{618}\,(y = v_{618} \to v_{618} = y)$}
\end{ptree}
\noindent$\Pi_{6}$:
\begin{ptree}
\AXC{$\forall x\,\forall y\,(x < y \to x \neq y)$\lbl{[F26]}}
\AXC{$\forall x\,\forall y\,(x < y \to x \neq y) \to \forall y\,(0 < y \to 0 \neq y)$\lbl{A2}}
\BIC{$\forall y\,(0 < y \to 0 \neq y)$}
\AXC{$\forall y\,(0 < y \to 0 \neq y) \to 0 < y \to 0 \neq y$\lbl{A2}}
\BIC{$0 < y \to 0 \neq y$}
\end{ptree}
\noindent$\Pi_{7}$:
\begin{ptree}
\AXC{$x \in \mathbb{R} \to y \in \mathbb{R} \to z \in \mathbb{R} \to 0 < y \to x < z \cdot y \to 0 < y$\lbl{A1}}
\AXC{$\stackrel{\Pi_{6}}{\vdots}$}\noLine
\UIC{$0 < y \to 0 \neq y$}
\AXC{$(0 < y \to 0 \neq y) \to \varphi_{50}$\lbl{A1}}
\BIC{$\varphi_{50}$}
\BIC{$x \in \mathbb{R} \to y \in \mathbb{R} \to z \in \mathbb{R} \to 0 < y \to x < z \cdot y \to 0 \neq y$}
\AXC{$\varphi_{56}$\lbl{A1}}
\BIC{$\varphi_{34}$}
\end{ptree}
\noindent$\Pi_{8}$:
\begin{ptree}
\AXC{$\stackrel{\Pi_{5}}{\vdots}$}\noLine
\UIC{$\forall v_{618}\,(y = v_{618} \to v_{618} = y)$}
\AXC{$\forall v_{618}\,(y = v_{618} \to v_{618} = y) \to y = 0 \to 0 = y$\lbl{A2}}
\BIC{$y = 0 \to 0 = y$}
\AXC{$\stackrel{\Pi_{7}}{\vdots}$}\noLine
\UIC{$\varphi_{34}$}
\BIC{$x \in \mathbb{R} \to y \in \mathbb{R} \to z \in \mathbb{R} \to 0 < y \to x < z \cdot y \to y \neq 0$}
\end{ptree}
\noindent$\Pi_{9}$:
\begin{ptree}
\AXC{$\forall x \in \mathbb{R}\,(x \neq 0 \to x^{-1} \in \mathbb{R})$\lbl{(R1)}}
\AXC{$\varphi_{27}$\lbl{A2}}
\BIC{$y \in \mathbb{R} \to y \neq 0 \to y^{-1} \in \mathbb{R}$}
\AXC{$(y \in \mathbb{R} \to y \neq 0 \to y^{-1} \in \mathbb{R}) \to \varphi_{64}$\lbl{A1}}
\BIC{$\varphi_{64}$}
\end{ptree}
\noindent$\Pi_{10}$:
\begin{ptree}
\AXC{$\stackrel{\Pi_{8}}{\vdots}$}\noLine
\UIC{$x \in \mathbb{R} \to y \in \mathbb{R} \to z \in \mathbb{R} \to 0 < y \to x < z \cdot y \to y \neq 0$}
\AXC{$x \in \mathbb{R} \to y \in \mathbb{R} \to z \in \mathbb{R} \to 0 < y \to x < z \cdot y \to y \in \mathbb{R}$\lbl{A1}}
\AXC{$\stackrel{\Pi_{9}}{\vdots}$}\noLine
\UIC{$\varphi_{64}$}
\BIC{$\varphi_{55}$}
\BIC{$\varphi_{28}$}
\end{ptree}
\noindent$\Pi_{11}$:
\begin{ptree}
\AXC{$\forall x\,\forall y\,(x \in \mathbb{R} \to y \in \mathbb{R} \to x \cdot y \in \mathbb{R})$\lbl{(R1)}}
\AXC{$\varphi_{46}$\lbl{A2}}
\BIC{$\forall y\,(z \in \mathbb{R} \to y \in \mathbb{R} \to z \cdot y \in \mathbb{R})$}
\AXC{$\varphi_{43}$\lbl{A2}}
\BIC{$z \in \mathbb{R} \to y \in \mathbb{R} \to z \cdot y \in \mathbb{R}$}
\AXC{$(z \in \mathbb{R} \to y \in \mathbb{R} \to z \cdot y \in \mathbb{R}) \to \varphi_{65}$\lbl{A1}}
\BIC{$\varphi_{65}$}
\end{ptree}
\noindent$\Pi_{12}$:
\begin{ptree}
\AXC{$x \in \mathbb{R} \to y \in \mathbb{R} \to z \in \mathbb{R} \to 0 < y \to x < z \cdot y \to y \in \mathbb{R}$\lbl{A1}}
\AXC{$x \in \mathbb{R} \to y \in \mathbb{R} \to z \in \mathbb{R} \to 0 < y \to x < z \cdot y \to z \in \mathbb{R}$\lbl{A1}}
\AXC{$\stackrel{\Pi_{11}}{\vdots}$}\noLine
\UIC{$\varphi_{65}$}
\BIC{$\varphi_{57}$}
\BIC{$\varphi_{29}$}
\end{ptree}
\noindent$\Pi_{13}$:
\begin{ptree}
\AXC{$\varphi_{39}$\lbl{[F32]}}
\AXC{$\varphi_{39} \to \varphi_{44}$\lbl{A2}}
\BIC{$\varphi_{44}$}
\AXC{$\varphi_{44} \to \varphi_{47}$\lbl{A2}}
\BIC{$\varphi_{47}$}
\AXC{$\varphi_{47} \to v_{619} \in \mathbb{R} \to v_{620} \in \mathbb{R} \to \varphi_{20}$\lbl{A2}}
\BIC{$v_{619} \in \mathbb{R} \to v_{620} \in \mathbb{R} \to \varphi_{20}$}
\UIC{$\forall v_{621}\,(v_{619} \in \mathbb{R} \to v_{620} \in \mathbb{R} \to \varphi_{20})$}
\UIC{$\forall v_{620}\,\forall v_{621}\,(v_{619} \in \mathbb{R} \to v_{620} \in \mathbb{R} \to \varphi_{20})$}
\UIC{$\forall v_{619}\,\forall v_{620}\,\forall v_{621}\,(v_{619} \in \mathbb{R} \to v_{620} \in \mathbb{R} \to \varphi_{20})$}
\AXC{$\varphi_{66}$\lbl{A2}}
\BIC{$\varphi_{48}$}
\AXC{$\varphi_{48} \to \varphi_{49}$\lbl{A2}}
\BIC{$\varphi_{49}$}
\end{ptree}
\noindent$\Pi_{14}$:
\begin{ptree}
\AXC{$\stackrel{\Pi_{13}}{\vdots}$}\noLine
\UIC{$\varphi_{49}$}
\AXC{$\varphi_{49} \to x \in \mathbb{R} \to z \cdot y \in \mathbb{R} \to \varphi_{33}$\lbl{A2}}
\BIC{$x \in \mathbb{R} \to z \cdot y \in \mathbb{R} \to \varphi_{33}$}
\AXC{$(x \in \mathbb{R} \to z \cdot y \in \mathbb{R} \to \varphi_{33}) \to \varphi_{71}$\lbl{A1}}
\BIC{$\varphi_{71}$}
\end{ptree}
\noindent$\Pi_{15}$:
\begin{ptree}
\AXC{$\stackrel{\Pi_{10}}{\vdots}$}\noLine
\UIC{$\varphi_{28}$}
\AXC{$\stackrel{\Pi_{12}}{\vdots}$}\noLine
\UIC{$\varphi_{29}$}
\AXC{$x \in \mathbb{R} \to y \in \mathbb{R} \to z \in \mathbb{R} \to 0 < y \to x < z \cdot y \to x \in \mathbb{R}$\lbl{A1}}
\AXC{$\stackrel{\Pi_{14}}{\vdots}$}\noLine
\UIC{$\varphi_{71}$}
\BIC{$\varphi_{29} \to \varphi_{28} \to \varphi_{67}$}
\BIC{$\varphi_{28} \to \varphi_{67}$}
\BIC{$\varphi_{67}$}
\end{ptree}
\noindent$\Pi_{16}$:
\begin{ptree}
\AXC{$x \in \mathbb{R} \to y \in \mathbb{R} \to z \in \mathbb{R} \to 0 < y \to x < z \cdot y \to x < z \cdot y$\lbl{A1}}
\AXC{$\stackrel{\Pi_{3}}{\vdots}$}\noLine
\UIC{$x \in \mathbb{R} \to y \in \mathbb{R} \to z \in \mathbb{R} \to 0 < y \to x < z \cdot y \to 0 < y^{-1}$}
\AXC{$\stackrel{\Pi_{15}}{\vdots}$}\noLine
\UIC{$\varphi_{67}$}
\BIC{$\varphi_{58}$}
\BIC{$x \in \mathbb{R} \to \varphi_{38}$}
\end{ptree}
\noindent$\Pi_{17}$:
\begin{ptree}
\AXC{$y = z \to x = y \to x = z$\lbl{A6}}
\AXC{$(y = z \to x = y \to x = z) \to x = y \to y = z \to x = z$\lbl{A1}}
\BIC{$x = y \to y = z \to x = z$}
\UIC{$\forall z\,(x = y \to y = z \to x = z)$}
\UIC{$\forall y\,\forall z\,(x = y \to y = z \to x = z)$}
\UIC{$\forall x\,\forall y\,\forall z\,(x = y \to y = z \to x = z)$}
\AXC{$\varphi_{4}$\lbl{A2}}
\BIC{$\forall y\,\forall z\,(v_{628} = y \to y = z \to v_{628} = z)$}
\AXC{$\varphi_{10}$\lbl{A2}}
\BIC{$\forall z\,(v_{628} = v_{629} \to v_{629} = z \to v_{628} = z)$}
\AXC{$\varphi_{12}$\lbl{A2}}
\BIC{$v_{628} = v_{629} \to v_{629} = v_{630} \to v_{628} = v_{630}$}
\UIC{$\forall v_{630}\,(v_{628} = v_{629} \to v_{629} = v_{630} \to v_{628} = v_{630})$}
\UIC{$\varphi_{2}$}
\UIC{$\forall v_{628}\,(\varphi_{2})$}
\end{ptree}
\noindent$\Pi_{18}$:
\begin{ptree}
\AXC{$\stackrel{\Pi_{17}}{\vdots}$}\noLine
\UIC{$\forall v_{628}\,(\varphi_{2})$}
\AXC{$\forall v_{628}\,(\varphi_{2}) \to \varphi_{15}$\lbl{A2}}
\BIC{$\varphi_{15}$}
\AXC{$\varphi_{15} \to \varphi_{14}$\lbl{A2}}
\BIC{$\varphi_{14}$}
\AXC{$\varphi_{14} \to z \cdot y \cdot y^{-1} = z \cdot 1 \to z \cdot 1 = z \to z \cdot y \cdot y^{-1} = z$\lbl{A2}}
\BIC{$z \cdot y \cdot y^{-1} = z \cdot 1 \to z \cdot 1 = z \to z \cdot y \cdot y^{-1} = z$}
\end{ptree}
\noindent$\Pi_{19}$:
\begin{ptree}
\AXC{$\forall x \in \mathbb{R}\,(x \cdot 1 = x)$\lbl{(R8)}}
\AXC{$\forall x \in \mathbb{R}\,(x \cdot 1 = x) \to z \in \mathbb{R} \to z \cdot 1 = z$\lbl{A2}}
\BIC{$z \in \mathbb{R} \to z \cdot 1 = z$}
\AXC{$(z \in \mathbb{R} \to z \cdot 1 = z) \to \varphi_{54}$\lbl{A1}}
\BIC{$\varphi_{54}$}
\end{ptree}
\noindent$\Pi_{20}$:
\begin{ptree}
\AXC{$y = z \to x = y \to x = z$\lbl{A6}}
\AXC{$(y = z \to x = y \to x = z) \to x = y \to y = z \to x = z$\lbl{A1}}
\BIC{$x = y \to y = z \to x = z$}
\UIC{$\forall z\,(x = y \to y = z \to x = z)$}
\UIC{$\forall y\,\forall z\,(x = y \to y = z \to x = z)$}
\UIC{$\forall x\,\forall y\,\forall z\,(x = y \to y = z \to x = z)$}
\AXC{$\varphi_{5}$\lbl{A2}}
\BIC{$\forall y\,\forall z\,(v_{625} = y \to y = z \to v_{625} = z)$}
\AXC{$\varphi_{11}$\lbl{A2}}
\BIC{$\forall z\,(v_{625} = v_{626} \to v_{626} = z \to v_{625} = z)$}
\AXC{$\varphi_{13}$\lbl{A2}}
\BIC{$v_{625} = v_{626} \to v_{626} = v_{627} \to v_{625} = v_{627}$}
\UIC{$\forall v_{627}\,(v_{625} = v_{626} \to v_{626} = v_{627} \to v_{625} = v_{627})$}
\UIC{$\varphi_{3}$}
\UIC{$\forall v_{625}\,(\varphi_{3})$}
\end{ptree}
\noindent$\Pi_{21}$:
\begin{ptree}
\AXC{$\forall x \in \mathbb{R}\,(x \neq 0 \to x \cdot x^{-1} = 1)$\lbl{(R9)}}
\AXC{$\varphi_{22}$\lbl{A2}}
\BIC{$y \in \mathbb{R} \to y \neq 0 \to y \cdot y^{-1} = 1$}
\AXC{$(y \in \mathbb{R} \to y \neq 0 \to y \cdot y^{-1} = 1) \to \varphi_{63}$\lbl{A1}}
\BIC{$\varphi_{63}$}
\end{ptree}
\noindent$\Pi_{22}$:
\begin{ptree}
\AXC{$\stackrel{\Pi_{8}}{\vdots}$}\noLine
\UIC{$x \in \mathbb{R} \to y \in \mathbb{R} \to z \in \mathbb{R} \to 0 < y \to x < z \cdot y \to y \neq 0$}
\AXC{$x \in \mathbb{R} \to y \in \mathbb{R} \to z \in \mathbb{R} \to 0 < y \to x < z \cdot y \to y \in \mathbb{R}$\lbl{A1}}
\AXC{$\stackrel{\Pi_{21}}{\vdots}$}\noLine
\UIC{$\varphi_{63}$}
\BIC{$\varphi_{53}$}
\BIC{$\varphi_{26}$}
\end{ptree}
\noindent$\Pi_{23}$:
\begin{ptree}
\AXC{$x = x$\lbl{A5}}
\UIC{$\forall x\,(x = x)$}
\AXC{$\forall x\,(x = x) \to z \cdot (y \cdot y^{-1}) = z \cdot (y \cdot y^{-1})$\lbl{A2}}
\BIC{$z \cdot (y \cdot y^{-1}) = z \cdot (y \cdot y^{-1})$}
\AXC{$\stackrel{\Pi_{22}}{\vdots}$}\noLine
\UIC{$\varphi_{26}$}
\AXC{$\varphi_{26} \to \varphi_{59}$\lbl{A1}}
\BIC{$\varphi_{59}$}
\BIC{$(y \cdot y^{-1} = 1 \to \varphi_{9}) \to x \in \mathbb{R} \to \varphi_{23}$}
\end{ptree}
\noindent$\Pi_{24}$:
\begin{ptree}
\AXC{$\varphi_{7}$\lbl{A6}}
\UIC{$\forall w\,(\varphi_{7})$}
\UIC{$\forall z_{1}\,\forall w\,(\varphi_{7})$}
\AXC{$\forall z_{1}\,\forall w\,(\varphi_{7}) \to \varphi_{35}$\lbl{A2}}
\BIC{$\varphi_{35}$}
\AXC{$\varphi_{35} \to y \cdot y^{-1} = 1 \to \varphi_{9}$\lbl{A2}}
\BIC{$y \cdot y^{-1} = 1 \to \varphi_{9}$}
\AXC{$\stackrel{\Pi_{23}}{\vdots}$}\noLine
\UIC{$(y \cdot y^{-1} = 1 \to \varphi_{9}) \to x \in \mathbb{R} \to \varphi_{23}$}
\BIC{$x \in \mathbb{R} \to \varphi_{23}$}
\end{ptree}
\noindent$\Pi_{25}$:
\begin{ptree}
\AXC{$\varphi_{19}$\lbl{(R6)}}
\AXC{$\varphi_{19} \to \varphi_{30}$\lbl{A2}}
\BIC{$\varphi_{30}$}
\AXC{$\varphi_{30} \to \varphi_{40}$\lbl{A2}}
\BIC{$\varphi_{40}$}
\AXC{$\varphi_{40} \to v_{622} \in \mathbb{R} \to v_{623} \in \mathbb{R} \to \varphi_{8}$\lbl{A2}}
\BIC{$v_{622} \in \mathbb{R} \to v_{623} \in \mathbb{R} \to \varphi_{8}$}
\UIC{$\forall v_{624}\,(v_{622} \in \mathbb{R} \to v_{623} \in \mathbb{R} \to \varphi_{8})$}
\UIC{$\forall v_{623}\,\forall v_{624}\,(v_{622} \in \mathbb{R} \to v_{623} \in \mathbb{R} \to \varphi_{8})$}
\UIC{$\forall v_{622}\,\forall v_{623}\,\forall v_{624}\,(v_{622} \in \mathbb{R} \to v_{623} \in \mathbb{R} \to \varphi_{8})$}
\AXC{$\varphi_{60}$\lbl{A2}}
\BIC{$\varphi_{45}$}
\AXC{$\varphi_{45} \to \varphi_{31}$\lbl{A2}}
\BIC{$\varphi_{31}$}
\end{ptree}
\noindent$\Pi_{26}$:
\begin{ptree}
\AXC{$x \in \mathbb{R} \to y \in \mathbb{R} \to z \in \mathbb{R} \to 0 < y \to x < z \cdot y \to z \in \mathbb{R}$\lbl{A1}}
\AXC{$\stackrel{\Pi_{25}}{\vdots}$}\noLine
\UIC{$\varphi_{31}$}
\AXC{$\varphi_{31} \to \varphi_{32}$\lbl{A2}}
\BIC{$\varphi_{32}$}
\AXC{$\varphi_{32} \to \varphi_{69}$\lbl{A1}}
\BIC{$\varphi_{69}$}
\BIC{$\varphi_{68}$}
\end{ptree}
\noindent$\Pi_{27}$:
\begin{ptree}
\AXC{$\stackrel{\Pi_{24}}{\vdots}$}\noLine
\UIC{$x \in \mathbb{R} \to \varphi_{23}$}
\AXC{$\stackrel{\Pi_{10}}{\vdots}$}\noLine
\UIC{$\varphi_{28}$}
\AXC{$x \in \mathbb{R} \to y \in \mathbb{R} \to z \in \mathbb{R} \to 0 < y \to x < z \cdot y \to y \in \mathbb{R}$\lbl{A1}}
\AXC{$\stackrel{\Pi_{26}}{\vdots}$}\noLine
\UIC{$\varphi_{68}$}
\BIC{$\varphi_{28} \to x \in \mathbb{R} \to y \in \mathbb{R} \to \varphi_{17}$}
\BIC{$x \in \mathbb{R} \to y \in \mathbb{R} \to \varphi_{17}$}
\AXC{$\varphi_{70}$\lbl{A1}}
\BIC{$(x \in \mathbb{R} \to \varphi_{23}) \to \varphi_{41} \to x \in \mathbb{R} \to \varphi_{24}$}
\BIC{$\varphi_{41} \to x \in \mathbb{R} \to \varphi_{24}$}
\end{ptree}
\noindent$\Pi_{28}$:
\begin{ptree}
\AXC{$\stackrel{\Pi_{20}}{\vdots}$}\noLine
\UIC{$\forall v_{625}\,(\varphi_{3})$}
\AXC{$\forall v_{625}\,(\varphi_{3}) \to \varphi_{16}$\lbl{A2}}
\BIC{$\varphi_{16}$}
\AXC{$\varphi_{16} \to \varphi_{42}$\lbl{A2}}
\BIC{$\varphi_{42}$}
\AXC{$\varphi_{42} \to \varphi_{41}$\lbl{A2}}
\BIC{$\varphi_{41}$}
\AXC{$\stackrel{\Pi_{27}}{\vdots}$}\noLine
\UIC{$\varphi_{41} \to x \in \mathbb{R} \to \varphi_{24}$}
\BIC{$x \in \mathbb{R} \to \varphi_{24}$}
\end{ptree}
\noindent$\Pi_{29}$:
\begin{ptree}
\AXC{$x \in \mathbb{R} \to y \in \mathbb{R} \to z \in \mathbb{R} \to 0 < y \to x < z \cdot y \to z \in \mathbb{R}$\lbl{A1}}
\AXC{$\stackrel{\Pi_{19}}{\vdots}$}\noLine
\UIC{$\varphi_{54}$}
\BIC{$x \in \mathbb{R} \to y \in \mathbb{R} \to z \in \mathbb{R} \to 0 < y \to x < z \cdot y \to z \cdot 1 = z$}
\AXC{$\stackrel{\Pi_{28}}{\vdots}$}\noLine
\UIC{$x \in \mathbb{R} \to \varphi_{24}$}
\AXC{$(x \in \mathbb{R} \to \varphi_{24}) \to \varphi_{61}$\lbl{A1}}
\BIC{$\varphi_{61}$}
\BIC{$\varphi_{51}$}
\end{ptree}
\noindent$\Pi_{30}$:
\begin{ptree}
\AXC{$\stackrel{\Pi_{16}}{\vdots}$}\noLine
\UIC{$x \in \mathbb{R} \to \varphi_{38}$}
\AXC{$\stackrel{\Pi_{18}}{\vdots}$}\noLine
\UIC{$z \cdot y \cdot y^{-1} = z \cdot 1 \to z \cdot 1 = z \to z \cdot y \cdot y^{-1} = z$}
\AXC{$\stackrel{\Pi_{29}}{\vdots}$}\noLine
\UIC{$\varphi_{51}$}
\BIC{$\varphi_{37}$}
\AXC{$\varphi_{37} \to (x \in \mathbb{R} \to \varphi_{38}) \to \varphi_{18} \to \varphi_{36}$\lbl{A1}}
\BIC{$(x \in \mathbb{R} \to \varphi_{38}) \to \varphi_{18} \to \varphi_{36}$}
\BIC{$\varphi_{18} \to \varphi_{36}$}
\end{ptree}
\noindent 主木:
\begin{ptree}
\AXC{$\stackrel{\Pi_{1}}{\vdots}$}\noLine
\UIC{$\varphi_{18}$}
\AXC{$\stackrel{\Pi_{30}}{\vdots}$}\noLine
\UIC{$\varphi_{18} \to \varphi_{36}$}
\BIC{$\varphi_{36}$}
\UIC{$\forall z\,(\varphi_{36})$}
\UIC{$\forall y\,\forall z\,(\varphi_{36})$}
\UIC{$\forall x\,\forall y\,\forall z\,(\varphi_{36})$}
\end{ptree}
\begin{abbrevlist}
\item $\varphi_{1} :\equiv \forall y\,(v_{617} = y \to y = v_{617}) \to v_{617} = v_{618} \to v_{618} = v_{617}$
\item $\varphi_{2} :\equiv \forall v_{629}\,\forall v_{630}\,(v_{628} = v_{629} \to v_{629} = v_{630} \to v_{628} = v_{630})$
\item $\varphi_{3} :\equiv \forall v_{626}\,\forall v_{627}\,(v_{625} = v_{626} \to v_{626} = v_{627} \to v_{625} = v_{627})$
\item $\varphi_{4} :\equiv \forall x\,\forall y\,\forall z\,(x = y \to y = z \to x = z) \to \forall y\,\forall z\,(v_{628} = y \to y = z \to v_{628} = z)$
\item $\varphi_{5} :\equiv \forall x\,\forall y\,\forall z\,(x = y \to y = z \to x = z) \to \forall y\,\forall z\,(v_{625} = y \to y = z \to v_{625} = z)$
\item $\varphi_{6} :\equiv \forall v_{617}\,\forall v_{618}\,(v_{617} = v_{618} \to v_{618} = v_{617}) \to \forall v_{618}\,(y = v_{618} \to v_{618} = y)$
\item $\varphi_{7} :\equiv z_{1} = w \to z \cdot (y \cdot y^{-1}) = z \cdot z_{1} \to z \cdot (y \cdot y^{-1}) = z \cdot w$
\item $\varphi_{8} :\equiv v_{624} \in \mathbb{R} \to v_{622} \cdot v_{623} \cdot v_{624} = v_{622} \cdot (v_{623} \cdot v_{624})$
\item $\varphi_{9} :\equiv z \cdot (y \cdot y^{-1}) = z \cdot (y \cdot y^{-1}) \to z \cdot (y \cdot y^{-1}) = z \cdot 1$
\item $\varphi_{10} :\equiv \forall y\,\forall z\,(v_{628} = y \to y = z \to v_{628} = z) \to \forall z\,(v_{628} = v_{629} \to v_{629} = z \to v_{628} = z)$
\item $\varphi_{11} :\equiv \forall y\,\forall z\,(v_{625} = y \to y = z \to v_{625} = z) \to \forall z\,(v_{625} = v_{626} \to v_{626} = z \to v_{625} = z)$
\item $\varphi_{12} :\equiv \forall z\,(v_{628} = v_{629} \to v_{629} = z \to v_{628} = z) \to v_{628} = v_{629} \to v_{629} = v_{630} \to v_{628} = v_{630}$
\item $\varphi_{13} :\equiv \forall z\,(v_{625} = v_{626} \to v_{626} = z \to v_{625} = z) \to v_{625} = v_{626} \to v_{626} = v_{627} \to v_{625} = v_{627}$
\item $\varphi_{14} :\equiv \forall v_{630}\,(z \cdot y \cdot y^{-1} = z \cdot 1 \to z \cdot 1 = v_{630} \to z \cdot y \cdot y^{-1} = v_{630})$
\item $\varphi_{15} :\equiv \forall v_{629}\,\forall v_{630}\,(z \cdot y \cdot y^{-1} = v_{629} \to v_{629} = v_{630} \to z \cdot y \cdot y^{-1} = v_{630})$
\item $\varphi_{16} :\equiv \forall v_{626}\,\forall v_{627}\,(z \cdot y \cdot y^{-1} = v_{626} \to v_{626} = v_{627} \to z \cdot y \cdot y^{-1} = v_{627})$
\item $\varphi_{17} :\equiv z \in \mathbb{R} \to 0 < y \to x < z \cdot y \to z \cdot y \cdot y^{-1} = z \cdot (y \cdot y^{-1})$
\item $\varphi_{18} :\equiv z \cdot y \cdot y^{-1} = z \to x \cdot y^{-1} < z \cdot y \cdot y^{-1} \to x \cdot y^{-1} < z$
\item $\varphi_{19} :\equiv \forall x\,\forall y\,\forall z\,(x \in \mathbb{R} \to y \in \mathbb{R} \to z \in \mathbb{R} \to x \cdot y \cdot z = x \cdot (y \cdot z))$
\item $\varphi_{20} :\equiv v_{621} \in \mathbb{R} \to 0 < v_{621} \to v_{619} < v_{620} \to v_{619} \cdot v_{621} < v_{620} \cdot v_{621}$
\item $\varphi_{21} :\equiv \forall x \in \mathbb{R}\,(0 < x \to 0 < x^{-1}) \to y \in \mathbb{R} \to 0 < y \to 0 < y^{-1}$
\item $\varphi_{22} :\equiv \forall x \in \mathbb{R}\,(x \neq 0 \to x \cdot x^{-1} = 1) \to y \in \mathbb{R} \to y \neq 0 \to y \cdot y^{-1} = 1$
\item $\varphi_{23} :\equiv y \in \mathbb{R} \to z \in \mathbb{R} \to 0 < y \to x < z \cdot y \to z \cdot (y \cdot y^{-1}) = z \cdot 1$
\item $\varphi_{24} :\equiv y \in \mathbb{R} \to z \in \mathbb{R} \to 0 < y \to x < z \cdot y \to z \cdot y \cdot y^{-1} = z \cdot 1$
\item $\varphi_{25} :\equiv \forall w\,(z \cdot y \cdot y^{-1} = w \to x \cdot y^{-1} < z \cdot y \cdot y^{-1} \to x \cdot y^{-1} < w)$
\item $\varphi_{26} :\equiv x \in \mathbb{R} \to y \in \mathbb{R} \to z \in \mathbb{R} \to 0 < y \to x < z \cdot y \to y \cdot y^{-1} = 1$
\item $\varphi_{27} :\equiv \forall x \in \mathbb{R}\,(x \neq 0 \to x^{-1} \in \mathbb{R}) \to y \in \mathbb{R} \to y \neq 0 \to y^{-1} \in \mathbb{R}$
\item $\varphi_{28} :\equiv x \in \mathbb{R} \to y \in \mathbb{R} \to z \in \mathbb{R} \to 0 < y \to x < z \cdot y \to y^{-1} \in \mathbb{R}$
\item $\varphi_{29} :\equiv x \in \mathbb{R} \to y \in \mathbb{R} \to z \in \mathbb{R} \to 0 < y \to x < z \cdot y \to z \cdot y \in \mathbb{R}$
\item $\varphi_{30} :\equiv \forall y\,\forall z\,(v_{622} \in \mathbb{R} \to y \in \mathbb{R} \to z \in \mathbb{R} \to v_{622} \cdot y \cdot z = v_{622} \cdot (y \cdot z))$
\item $\varphi_{31} :\equiv \forall v_{624}\,(z \in \mathbb{R} \to y \in \mathbb{R} \to v_{624} \in \mathbb{R} \to z \cdot y \cdot v_{624} = z \cdot (y \cdot v_{624}))$
\item $\varphi_{32} :\equiv z \in \mathbb{R} \to y \in \mathbb{R} \to y^{-1} \in \mathbb{R} \to z \cdot y \cdot y^{-1} = z \cdot (y \cdot y^{-1})$
\item $\varphi_{33} :\equiv y^{-1} \in \mathbb{R} \to 0 < y^{-1} \to x < z \cdot y \to x \cdot y^{-1} < z \cdot y \cdot y^{-1}$
\item $\varphi_{34} :\equiv (y = 0 \to 0 = y) \to x \in \mathbb{R} \to y \in \mathbb{R} \to z \in \mathbb{R} \to 0 < y \to x < z \cdot y \to y \neq 0$
\item $\varphi_{35} :\equiv \forall w\,(y \cdot y^{-1} = w \to z \cdot (y \cdot y^{-1}) = z \cdot (y \cdot y^{-1}) \to z \cdot (y \cdot y^{-1}) = z \cdot w)$
\item $\varphi_{36} :\equiv x \in \mathbb{R} \to y \in \mathbb{R} \to z \in \mathbb{R} \to 0 < y \to x < z \cdot y \to x \cdot y^{-1} < z$
\item $\varphi_{37} :\equiv x \in \mathbb{R} \to y \in \mathbb{R} \to z \in \mathbb{R} \to 0 < y \to x < z \cdot y \to z \cdot y \cdot y^{-1} = z$
\item $\varphi_{38} :\equiv y \in \mathbb{R} \to z \in \mathbb{R} \to 0 < y \to x < z \cdot y \to x \cdot y^{-1} < z \cdot y \cdot y^{-1}$
\item $\varphi_{39} :\equiv \forall x\,\forall y\,\forall z\,(x \in \mathbb{R} \to y \in \mathbb{R} \to z \in \mathbb{R} \to 0 < z \to x < y \to x \cdot z < y \cdot z)$
\item $\varphi_{40} :\equiv \forall z\,(v_{622} \in \mathbb{R} \to v_{623} \in \mathbb{R} \to z \in \mathbb{R} \to v_{622} \cdot v_{623} \cdot z = v_{622} \cdot (v_{623} \cdot z))$
\item $\varphi_{41} :\equiv z \cdot y \cdot y^{-1} = z \cdot (y \cdot y^{-1}) \to z \cdot (y \cdot y^{-1}) = z \cdot 1 \to z \cdot y \cdot y^{-1} = z \cdot 1$
\item $\varphi_{42} :\equiv \forall v_{627}\,(z \cdot y \cdot y^{-1} = z \cdot (y \cdot y^{-1}) \to z \cdot (y \cdot y^{-1}) = v_{627} \to z \cdot y \cdot y^{-1} = v_{627})$
\item $\varphi_{43} :\equiv \forall y\,(z \in \mathbb{R} \to y \in \mathbb{R} \to z \cdot y \in \mathbb{R}) \to z \in \mathbb{R} \to y \in \mathbb{R} \to z \cdot y \in \mathbb{R}$
\item $\varphi_{44} :\equiv \forall y\,\forall z\,(v_{619} \in \mathbb{R} \to y \in \mathbb{R} \to z \in \mathbb{R} \to 0 < z \to v_{619} < y \to v_{619} \cdot z < y \cdot z)$
\item $\varphi_{45} :\equiv \forall v_{623}\,\forall v_{624}\,(z \in \mathbb{R} \to v_{623} \in \mathbb{R} \to v_{624} \in \mathbb{R} \to z \cdot v_{623} \cdot v_{624} = z \cdot (v_{623} \cdot v_{624}))$
\item $\varphi_{46} :\equiv \forall x\,\forall y\,(x \in \mathbb{R} \to y \in \mathbb{R} \to x \cdot y \in \mathbb{R}) \to \forall y\,(z \in \mathbb{R} \to y \in \mathbb{R} \to z \cdot y \in \mathbb{R})$
\item $\varphi_{47} :\equiv \forall z\,(v_{619} \in \mathbb{R} \to v_{620} \in \mathbb{R} \to z \in \mathbb{R} \to 0 < z \to v_{619} < v_{620} \to v_{619} \cdot z < v_{620} \cdot z)$
\item $\varphi_{48} :\equiv \forall v_{620}\,\forall v_{621}\,(x \in \mathbb{R} \to v_{620} \in \mathbb{R} \to v_{621} \in \mathbb{R} \to 0 < v_{621} \to x < v_{620} \to x \cdot v_{621} < v_{620} \cdot v_{621})$
\item $\varphi_{49} :\equiv \forall v_{621}\,(x \in \mathbb{R} \to z \cdot y \in \mathbb{R} \to v_{621} \in \mathbb{R} \to 0 < v_{621} \to x < z \cdot y \to x \cdot v_{621} < z \cdot y \cdot v_{621})$
\item $\varphi_{50} :\equiv (x \in \mathbb{R} \to y \in \mathbb{R} \to z \in \mathbb{R} \to 0 < y \to x < z \cdot y \to 0 < y) \to x \in \mathbb{R} \to y \in \mathbb{R} \to z \in \mathbb{R} \to 0 < y \to x < z \cdot y \to 0 \neq y$
\item $\varphi_{51} :\equiv (z \cdot y \cdot y^{-1} = z \cdot 1 \to z \cdot 1 = z \to z \cdot y \cdot y^{-1} = z) \to \varphi_{37}$
\item $\varphi_{52} :\equiv (x \in \mathbb{R} \to y \in \mathbb{R} \to z \in \mathbb{R} \to 0 < y \to x < z \cdot y \to 0 < y) \to x \in \mathbb{R} \to y \in \mathbb{R} \to z \in \mathbb{R} \to 0 < y \to x < z \cdot y \to 0 < y^{-1}$
\item $\varphi_{53} :\equiv (x \in \mathbb{R} \to y \in \mathbb{R} \to z \in \mathbb{R} \to 0 < y \to x < z \cdot y \to y \neq 0) \to \varphi_{26}$
\item $\varphi_{54} :\equiv (x \in \mathbb{R} \to y \in \mathbb{R} \to z \in \mathbb{R} \to 0 < y \to x < z \cdot y \to z \in \mathbb{R}) \to x \in \mathbb{R} \to y \in \mathbb{R} \to z \in \mathbb{R} \to 0 < y \to x < z \cdot y \to z \cdot 1 = z$
\item $\varphi_{55} :\equiv (x \in \mathbb{R} \to y \in \mathbb{R} \to z \in \mathbb{R} \to 0 < y \to x < z \cdot y \to y \neq 0) \to \varphi_{28}$
\item $\varphi_{56} :\equiv (x \in \mathbb{R} \to y \in \mathbb{R} \to z \in \mathbb{R} \to 0 < y \to x < z \cdot y \to 0 \neq y) \to \varphi_{34}$
\item $\varphi_{57} :\equiv (x \in \mathbb{R} \to y \in \mathbb{R} \to z \in \mathbb{R} \to 0 < y \to x < z \cdot y \to y \in \mathbb{R}) \to \varphi_{29}$
\item $\varphi_{58} :\equiv (x \in \mathbb{R} \to y \in \mathbb{R} \to z \in \mathbb{R} \to 0 < y \to x < z \cdot y \to x < z \cdot y) \to x \in \mathbb{R} \to \varphi_{38}$
\item $\varphi_{59} :\equiv z \cdot (y \cdot y^{-1}) = z \cdot (y \cdot y^{-1}) \to (y \cdot y^{-1} = 1 \to \varphi_{9}) \to x \in \mathbb{R} \to \varphi_{23}$
\item $\varphi_{60} :\equiv \forall v_{622}\,\forall v_{623}\,\forall v_{624}\,(v_{622} \in \mathbb{R} \to v_{623} \in \mathbb{R} \to \varphi_{8}) \to \varphi_{45}$
\item $\varphi_{61} :\equiv (x \in \mathbb{R} \to y \in \mathbb{R} \to z \in \mathbb{R} \to 0 < y \to x < z \cdot y \to z \cdot 1 = z) \to \varphi_{51}$
\item $\varphi_{62} :\equiv (x \in \mathbb{R} \to y \in \mathbb{R} \to z \in \mathbb{R} \to 0 < y \to x < z \cdot y \to y \in \mathbb{R}) \to \varphi_{52}$
\item $\varphi_{63} :\equiv (x \in \mathbb{R} \to y \in \mathbb{R} \to z \in \mathbb{R} \to 0 < y \to x < z \cdot y \to y \in \mathbb{R}) \to \varphi_{53}$
\item $\varphi_{64} :\equiv (x \in \mathbb{R} \to y \in \mathbb{R} \to z \in \mathbb{R} \to 0 < y \to x < z \cdot y \to y \in \mathbb{R}) \to \varphi_{55}$
\item $\varphi_{65} :\equiv (x \in \mathbb{R} \to y \in \mathbb{R} \to z \in \mathbb{R} \to 0 < y \to x < z \cdot y \to z \in \mathbb{R}) \to \varphi_{57}$
\item $\varphi_{66} :\equiv \forall v_{619}\,\forall v_{620}\,\forall v_{621}\,(v_{619} \in \mathbb{R} \to v_{620} \in \mathbb{R} \to \varphi_{20}) \to \varphi_{48}$
\item $\varphi_{67} :\equiv (x \in \mathbb{R} \to y \in \mathbb{R} \to z \in \mathbb{R} \to 0 < y \to x < z \cdot y \to 0 < y^{-1}) \to \varphi_{58}$
\item $\varphi_{68} :\equiv (x \in \mathbb{R} \to y \in \mathbb{R} \to z \in \mathbb{R} \to 0 < y \to x < z \cdot y \to y \in \mathbb{R}) \to \varphi_{28} \to x \in \mathbb{R} \to y \in \mathbb{R} \to \varphi_{17}$
\item $\varphi_{69} :\equiv (x \in \mathbb{R} \to y \in \mathbb{R} \to z \in \mathbb{R} \to 0 < y \to x < z \cdot y \to z \in \mathbb{R}) \to \varphi_{68}$
\item $\varphi_{70} :\equiv (x \in \mathbb{R} \to y \in \mathbb{R} \to \varphi_{17}) \to (x \in \mathbb{R} \to \varphi_{23}) \to \varphi_{41} \to x \in \mathbb{R} \to \varphi_{24}$
\item $\varphi_{71} :\equiv (x \in \mathbb{R} \to y \in \mathbb{R} \to z \in \mathbb{R} \to 0 < y \to x < z \cdot y \to x \in \mathbb{R}) \to \varphi_{29} \to \varphi_{28} \to \varphi_{67}$
\end{abbrevlist}
\end{fullproof}

% 自動生成 (prover/build.js): MulLtDiv  ノード数 183
\begin{fullproof}{[E25]}{$\forall x\,\forall y\,\forall z\,(x \in \mathbb{R} \to y \in \mathbb{R} \to z \in \mathbb{R} \to 0 < y \to x \cdot y < z \to x < z \cdot y^{-1})$}{183}\phantomsection\label{tree:MulLtDiv}
\noindent$\Pi_{1}$:
\begin{ptree}
\AXC{$z_{1} = w \to z_{1} < z \cdot y^{-1} \to w < z \cdot y^{-1}$\lbl{A6}}
\UIC{$\forall w\,(z_{1} = w \to z_{1} < z \cdot y^{-1} \to w < z \cdot y^{-1})$}
\UIC{$\forall z_{1}\,\forall w\,(z_{1} = w \to z_{1} < z \cdot y^{-1} \to w < z \cdot y^{-1})$}
\AXC{$\forall z_{1}\,\forall w\,(z_{1} = w \to z_{1} < z \cdot y^{-1} \to w < z \cdot y^{-1}) \to \varphi_{24}$\lbl{A2}}
\BIC{$\varphi_{24}$}
\AXC{$\varphi_{24} \to \varphi_{18}$\lbl{A2}}
\BIC{$\varphi_{18}$}
\end{ptree}
\noindent$\Pi_{2}$:
\begin{ptree}
\AXC{$\forall x \in \mathbb{R}\,(0 < x \to 0 < x^{-1})$\lbl{[F34]}}
\AXC{$\varphi_{20}$\lbl{A2}}
\BIC{$y \in \mathbb{R} \to 0 < y \to 0 < y^{-1}$}
\AXC{$(y \in \mathbb{R} \to 0 < y \to 0 < y^{-1}) \to \varphi_{56}$\lbl{A1}}
\BIC{$\varphi_{56}$}
\end{ptree}
\noindent$\Pi_{3}$:
\begin{ptree}
\AXC{$x \in \mathbb{R} \to y \in \mathbb{R} \to z \in \mathbb{R} \to 0 < y \to x \cdot y < z \to 0 < y$\lbl{A1}}
\AXC{$x \in \mathbb{R} \to y \in \mathbb{R} \to z \in \mathbb{R} \to 0 < y \to x \cdot y < z \to y \in \mathbb{R}$\lbl{A1}}
\AXC{$\stackrel{\Pi_{2}}{\vdots}$}\noLine
\UIC{$\varphi_{56}$}
\BIC{$\varphi_{47}$}
\BIC{$x \in \mathbb{R} \to y \in \mathbb{R} \to z \in \mathbb{R} \to 0 < y \to x \cdot y < z \to 0 < y^{-1}$}
\end{ptree}
\noindent$\Pi_{4}$:
\begin{ptree}
\AXC{$x = x$\lbl{A5}}
\AXC{$x = y \to x = x \to y = x$\lbl{A6}}
\AXC{$(x = y \to x = x \to y = x) \to x = x \to x = y \to y = x$\lbl{A1}}
\BIC{$x = x \to x = y \to y = x$}
\BIC{$x = y \to y = x$}
\UIC{$\forall y\,(x = y \to y = x)$}
\UIC{$\forall x\,\forall y\,(x = y \to y = x)$}
\end{ptree}
\noindent$\Pi_{5}$:
\begin{ptree}
\AXC{$\stackrel{\Pi_{4}}{\vdots}$}\noLine
\UIC{$\forall x\,\forall y\,(x = y \to y = x)$}
\AXC{$\forall x\,\forall y\,(x = y \to y = x) \to \forall y\,(v_{631} = y \to y = v_{631})$\lbl{A2}}
\BIC{$\forall y\,(v_{631} = y \to y = v_{631})$}
\AXC{$\varphi_{1}$\lbl{A2}}
\BIC{$v_{631} = v_{632} \to v_{632} = v_{631}$}
\UIC{$\forall v_{632}\,(v_{631} = v_{632} \to v_{632} = v_{631})$}
\UIC{$\forall v_{631}\,\forall v_{632}\,(v_{631} = v_{632} \to v_{632} = v_{631})$}
\AXC{$\varphi_{7}$\lbl{A2}}
\BIC{$\forall v_{632}\,(y = v_{632} \to v_{632} = y)$}
\end{ptree}
\noindent$\Pi_{6}$:
\begin{ptree}
\AXC{$\forall x\,\forall y\,(x < y \to x \neq y)$\lbl{[F26]}}
\AXC{$\forall x\,\forall y\,(x < y \to x \neq y) \to \forall y\,(0 < y \to 0 \neq y)$\lbl{A2}}
\BIC{$\forall y\,(0 < y \to 0 \neq y)$}
\AXC{$\forall y\,(0 < y \to 0 \neq y) \to 0 < y \to 0 \neq y$\lbl{A2}}
\BIC{$0 < y \to 0 \neq y$}
\end{ptree}
\noindent$\Pi_{7}$:
\begin{ptree}
\AXC{$x \in \mathbb{R} \to y \in \mathbb{R} \to z \in \mathbb{R} \to 0 < y \to x \cdot y < z \to 0 < y$\lbl{A1}}
\AXC{$\stackrel{\Pi_{6}}{\vdots}$}\noLine
\UIC{$0 < y \to 0 \neq y$}
\AXC{$(0 < y \to 0 \neq y) \to \varphi_{45}$\lbl{A1}}
\BIC{$\varphi_{45}$}
\BIC{$x \in \mathbb{R} \to y \in \mathbb{R} \to z \in \mathbb{R} \to 0 < y \to x \cdot y < z \to 0 \neq y$}
\AXC{$\varphi_{51}$\lbl{A1}}
\BIC{$\varphi_{31}$}
\end{ptree}
\noindent$\Pi_{8}$:
\begin{ptree}
\AXC{$\stackrel{\Pi_{5}}{\vdots}$}\noLine
\UIC{$\forall v_{632}\,(y = v_{632} \to v_{632} = y)$}
\AXC{$\forall v_{632}\,(y = v_{632} \to v_{632} = y) \to y = 0 \to 0 = y$\lbl{A2}}
\BIC{$y = 0 \to 0 = y$}
\AXC{$\stackrel{\Pi_{7}}{\vdots}$}\noLine
\UIC{$\varphi_{31}$}
\BIC{$x \in \mathbb{R} \to y \in \mathbb{R} \to z \in \mathbb{R} \to 0 < y \to x \cdot y < z \to y \neq 0$}
\end{ptree}
\noindent$\Pi_{9}$:
\begin{ptree}
\AXC{$\forall x \in \mathbb{R}\,(x \neq 0 \to x^{-1} \in \mathbb{R})$\lbl{(R1)}}
\AXC{$\varphi_{26}$\lbl{A2}}
\BIC{$y \in \mathbb{R} \to y \neq 0 \to y^{-1} \in \mathbb{R}$}
\AXC{$(y \in \mathbb{R} \to y \neq 0 \to y^{-1} \in \mathbb{R}) \to \varphi_{58}$\lbl{A1}}
\BIC{$\varphi_{58}$}
\end{ptree}
\noindent$\Pi_{10}$:
\begin{ptree}
\AXC{$\stackrel{\Pi_{8}}{\vdots}$}\noLine
\UIC{$x \in \mathbb{R} \to y \in \mathbb{R} \to z \in \mathbb{R} \to 0 < y \to x \cdot y < z \to y \neq 0$}
\AXC{$x \in \mathbb{R} \to y \in \mathbb{R} \to z \in \mathbb{R} \to 0 < y \to x \cdot y < z \to y \in \mathbb{R}$\lbl{A1}}
\AXC{$\stackrel{\Pi_{9}}{\vdots}$}\noLine
\UIC{$\varphi_{58}$}
\BIC{$\varphi_{50}$}
\BIC{$\varphi_{27}$}
\end{ptree}
\noindent$\Pi_{11}$:
\begin{ptree}
\AXC{$\forall x\,\forall y\,(x \in \mathbb{R} \to y \in \mathbb{R} \to x \cdot y \in \mathbb{R})$\lbl{(R1)}}
\AXC{$\varphi_{41}$\lbl{A2}}
\BIC{$\forall y\,(x \in \mathbb{R} \to y \in \mathbb{R} \to x \cdot y \in \mathbb{R})$}
\AXC{$\varphi_{39}$\lbl{A2}}
\BIC{$x \in \mathbb{R} \to y \in \mathbb{R} \to x \cdot y \in \mathbb{R}$}
\AXC{$(x \in \mathbb{R} \to y \in \mathbb{R} \to x \cdot y \in \mathbb{R}) \to \varphi_{59}$\lbl{A1}}
\BIC{$\varphi_{59}$}
\end{ptree}
\noindent$\Pi_{12}$:
\begin{ptree}
\AXC{$x \in \mathbb{R} \to y \in \mathbb{R} \to z \in \mathbb{R} \to 0 < y \to x \cdot y < z \to y \in \mathbb{R}$\lbl{A1}}
\AXC{$x \in \mathbb{R} \to y \in \mathbb{R} \to z \in \mathbb{R} \to 0 < y \to x \cdot y < z \to x \in \mathbb{R}$\lbl{A1}}
\AXC{$\stackrel{\Pi_{11}}{\vdots}$}\noLine
\UIC{$\varphi_{59}$}
\BIC{$\varphi_{52}$}
\BIC{$\varphi_{28}$}
\end{ptree}
\noindent$\Pi_{13}$:
\begin{ptree}
\AXC{$\varphi_{36}$\lbl{[F32]}}
\AXC{$\varphi_{36} \to \varphi_{40}$\lbl{A2}}
\BIC{$\varphi_{40}$}
\AXC{$\varphi_{40} \to \varphi_{42}$\lbl{A2}}
\BIC{$\varphi_{42}$}
\AXC{$\varphi_{42} \to v_{633} \in \mathbb{R} \to v_{634} \in \mathbb{R} \to \varphi_{19}$\lbl{A2}}
\BIC{$v_{633} \in \mathbb{R} \to v_{634} \in \mathbb{R} \to \varphi_{19}$}
\UIC{$\forall v_{635}\,(v_{633} \in \mathbb{R} \to v_{634} \in \mathbb{R} \to \varphi_{19})$}
\UIC{$\forall v_{634}\,\forall v_{635}\,(v_{633} \in \mathbb{R} \to v_{634} \in \mathbb{R} \to \varphi_{19})$}
\UIC{$\forall v_{633}\,\forall v_{634}\,\forall v_{635}\,(v_{633} \in \mathbb{R} \to v_{634} \in \mathbb{R} \to \varphi_{19})$}
\AXC{$\varphi_{60}$\lbl{A2}}
\BIC{$\varphi_{44}$}
\AXC{$\varphi_{44} \to \varphi_{43}$\lbl{A2}}
\BIC{$\varphi_{43}$}
\end{ptree}
\noindent$\Pi_{14}$:
\begin{ptree}
\AXC{$\stackrel{\Pi_{12}}{\vdots}$}\noLine
\UIC{$\varphi_{28}$}
\AXC{$\stackrel{\Pi_{13}}{\vdots}$}\noLine
\UIC{$\varphi_{43}$}
\AXC{$\varphi_{43} \to x \cdot y \in \mathbb{R} \to z \in \mathbb{R} \to \varphi_{30}$\lbl{A2}}
\BIC{$x \cdot y \in \mathbb{R} \to z \in \mathbb{R} \to \varphi_{30}$}
\AXC{$(x \cdot y \in \mathbb{R} \to z \in \mathbb{R} \to \varphi_{30}) \to \varphi_{28} \to \varphi_{65}$\lbl{A1}}
\BIC{$\varphi_{28} \to \varphi_{65}$}
\BIC{$\varphi_{65}$}
\end{ptree}
\noindent$\Pi_{15}$:
\begin{ptree}
\AXC{$\stackrel{\Pi_{10}}{\vdots}$}\noLine
\UIC{$\varphi_{27}$}
\AXC{$x \in \mathbb{R} \to y \in \mathbb{R} \to z \in \mathbb{R} \to 0 < y \to x \cdot y < z \to z \in \mathbb{R}$\lbl{A1}}
\AXC{$\stackrel{\Pi_{14}}{\vdots}$}\noLine
\UIC{$\varphi_{65}$}
\BIC{$\varphi_{27} \to \varphi_{61}$}
\BIC{$\varphi_{61}$}
\end{ptree}
\noindent$\Pi_{16}$:
\begin{ptree}
\AXC{$x \in \mathbb{R} \to y \in \mathbb{R} \to z \in \mathbb{R} \to 0 < y \to x \cdot y < z \to x \cdot y < z$\lbl{A1}}
\AXC{$\stackrel{\Pi_{3}}{\vdots}$}\noLine
\UIC{$x \in \mathbb{R} \to y \in \mathbb{R} \to z \in \mathbb{R} \to 0 < y \to x \cdot y < z \to 0 < y^{-1}$}
\AXC{$\stackrel{\Pi_{15}}{\vdots}$}\noLine
\UIC{$\varphi_{61}$}
\BIC{$\varphi_{53}$}
\BIC{$x \in \mathbb{R} \to \varphi_{35}$}
\end{ptree}
\noindent$\Pi_{17}$:
\begin{ptree}
\AXC{$y = z \to x = y \to x = z$\lbl{A6}}
\AXC{$(y = z \to x = y \to x = z) \to x = y \to y = z \to x = z$\lbl{A1}}
\BIC{$x = y \to y = z \to x = z$}
\UIC{$\forall z\,(x = y \to y = z \to x = z)$}
\UIC{$\forall y\,\forall z\,(x = y \to y = z \to x = z)$}
\UIC{$\forall x\,\forall y\,\forall z\,(x = y \to y = z \to x = z)$}
\AXC{$\varphi_{4}$\lbl{A2}}
\BIC{$\forall y\,\forall z\,(v_{639} = y \to y = z \to v_{639} = z)$}
\AXC{$\varphi_{9}$\lbl{A2}}
\BIC{$\forall z\,(v_{639} = v_{640} \to v_{640} = z \to v_{639} = z)$}
\AXC{$\varphi_{12}$\lbl{A2}}
\BIC{$v_{639} = v_{640} \to v_{640} = v_{641} \to v_{639} = v_{641}$}
\UIC{$\forall v_{641}\,(v_{639} = v_{640} \to v_{640} = v_{641} \to v_{639} = v_{641})$}
\UIC{$\varphi_{2}$}
\UIC{$\forall v_{639}\,(\varphi_{2})$}
\end{ptree}
\noindent$\Pi_{18}$:
\begin{ptree}
\AXC{$\stackrel{\Pi_{17}}{\vdots}$}\noLine
\UIC{$\forall v_{639}\,(\varphi_{2})$}
\AXC{$\forall v_{639}\,(\varphi_{2}) \to \varphi_{15}$\lbl{A2}}
\BIC{$\varphi_{15}$}
\AXC{$\varphi_{15} \to \varphi_{14}$\lbl{A2}}
\BIC{$\varphi_{14}$}
\AXC{$\varphi_{14} \to x \cdot y \cdot y^{-1} = x \cdot 1 \to x \cdot 1 = x \to x \cdot y \cdot y^{-1} = x$\lbl{A2}}
\BIC{$x \cdot y \cdot y^{-1} = x \cdot 1 \to x \cdot 1 = x \to x \cdot y \cdot y^{-1} = x$}
\end{ptree}
\noindent$\Pi_{19}$:
\begin{ptree}
\AXC{$\forall x \in \mathbb{R}\,(x \cdot 1 = x)$\lbl{(R8)}}
\AXC{$\forall x \in \mathbb{R}\,(x \cdot 1 = x) \to x \in \mathbb{R} \to x \cdot 1 = x$\lbl{A2}}
\BIC{$x \in \mathbb{R} \to x \cdot 1 = x$}
\AXC{$(x \in \mathbb{R} \to x \cdot 1 = x) \to \varphi_{49}$\lbl{A1}}
\BIC{$\varphi_{49}$}
\end{ptree}
\noindent$\Pi_{20}$:
\begin{ptree}
\AXC{$y = z \to x = y \to x = z$\lbl{A6}}
\AXC{$(y = z \to x = y \to x = z) \to x = y \to y = z \to x = z$\lbl{A1}}
\BIC{$x = y \to y = z \to x = z$}
\UIC{$\forall z\,(x = y \to y = z \to x = z)$}
\UIC{$\forall y\,\forall z\,(x = y \to y = z \to x = z)$}
\UIC{$\forall x\,\forall y\,\forall z\,(x = y \to y = z \to x = z)$}
\AXC{$\varphi_{5}$\lbl{A2}}
\BIC{$\forall y\,\forall z\,(v_{636} = y \to y = z \to v_{636} = z)$}
\AXC{$\varphi_{10}$\lbl{A2}}
\BIC{$\forall z\,(v_{636} = v_{637} \to v_{637} = z \to v_{636} = z)$}
\AXC{$\varphi_{13}$\lbl{A2}}
\BIC{$v_{636} = v_{637} \to v_{637} = v_{638} \to v_{636} = v_{638}$}
\UIC{$\forall v_{638}\,(v_{636} = v_{637} \to v_{637} = v_{638} \to v_{636} = v_{638})$}
\UIC{$\varphi_{3}$}
\UIC{$\forall v_{636}\,(\varphi_{3})$}
\end{ptree}
\noindent$\Pi_{21}$:
\begin{ptree}
\AXC{$\stackrel{\Pi_{20}}{\vdots}$}\noLine
\UIC{$\forall v_{636}\,(\varphi_{3})$}
\AXC{$\forall v_{636}\,(\varphi_{3}) \to \varphi_{16}$\lbl{A2}}
\BIC{$\varphi_{16}$}
\AXC{$\varphi_{16} \to \varphi_{38}$\lbl{A2}}
\BIC{$\varphi_{38}$}
\AXC{$\varphi_{38} \to \varphi_{37}$\lbl{A2}}
\BIC{$\varphi_{37}$}
\end{ptree}
\noindent$\Pi_{22}$:
\begin{ptree}
\AXC{$\forall x \in \mathbb{R}\,(x \neq 0 \to x \cdot x^{-1} = 1)$\lbl{(R9)}}
\AXC{$\varphi_{21}$\lbl{A2}}
\BIC{$y \in \mathbb{R} \to y \neq 0 \to y \cdot y^{-1} = 1$}
\AXC{$(y \in \mathbb{R} \to y \neq 0 \to y \cdot y^{-1} = 1) \to \varphi_{57}$\lbl{A1}}
\BIC{$\varphi_{57}$}
\end{ptree}
\noindent$\Pi_{23}$:
\begin{ptree}
\AXC{$\stackrel{\Pi_{8}}{\vdots}$}\noLine
\UIC{$x \in \mathbb{R} \to y \in \mathbb{R} \to z \in \mathbb{R} \to 0 < y \to x \cdot y < z \to y \neq 0$}
\AXC{$x \in \mathbb{R} \to y \in \mathbb{R} \to z \in \mathbb{R} \to 0 < y \to x \cdot y < z \to y \in \mathbb{R}$\lbl{A1}}
\AXC{$\stackrel{\Pi_{22}}{\vdots}$}\noLine
\UIC{$\varphi_{57}$}
\BIC{$\varphi_{48}$}
\BIC{$\varphi_{25}$}
\end{ptree}
\noindent$\Pi_{24}$:
\begin{ptree}
\AXC{$x = x$\lbl{A5}}
\UIC{$\forall x\,(x = x)$}
\AXC{$\forall x\,(x = x) \to x \cdot (y \cdot y^{-1}) = x \cdot (y \cdot y^{-1})$\lbl{A2}}
\BIC{$x \cdot (y \cdot y^{-1}) = x \cdot (y \cdot y^{-1})$}
\AXC{$\stackrel{\Pi_{23}}{\vdots}$}\noLine
\UIC{$\varphi_{25}$}
\AXC{$\varphi_{25} \to \varphi_{54}$\lbl{A1}}
\BIC{$\varphi_{54}$}
\BIC{$(y \cdot y^{-1} = 1 \to \varphi_{8}) \to x \in \mathbb{R} \to \varphi_{22}$}
\end{ptree}
\noindent$\Pi_{25}$:
\begin{ptree}
\AXC{$\varphi_{6}$\lbl{A6}}
\UIC{$\forall w\,(\varphi_{6})$}
\UIC{$\forall z\,\forall w\,(\varphi_{6})$}
\AXC{$\forall z\,\forall w\,(\varphi_{6}) \to \varphi_{32}$\lbl{A2}}
\BIC{$\varphi_{32}$}
\AXC{$\varphi_{32} \to y \cdot y^{-1} = 1 \to \varphi_{8}$\lbl{A2}}
\BIC{$y \cdot y^{-1} = 1 \to \varphi_{8}$}
\AXC{$\stackrel{\Pi_{24}}{\vdots}$}\noLine
\UIC{$(y \cdot y^{-1} = 1 \to \varphi_{8}) \to x \in \mathbb{R} \to \varphi_{22}$}
\BIC{$x \in \mathbb{R} \to \varphi_{22}$}
\end{ptree}
\noindent$\Pi_{26}$:
\begin{ptree}
\AXC{$\forall x\,\forall y\,(\varphi_{11})$\lbl{(R6)}}
\AXC{$\forall x\,\forall y\,(\varphi_{11}) \to \forall y\,(\varphi_{11})$\lbl{A2}}
\BIC{$\forall y\,(\varphi_{11})$}
\AXC{$\forall y\,(\varphi_{11}) \to \varphi_{11}$\lbl{A2}}
\BIC{$\varphi_{11}$}
\AXC{$\varphi_{11} \to \varphi_{29}$\lbl{A2}}
\BIC{$\varphi_{29}$}
\AXC{$\varphi_{29} \to \varphi_{63}$\lbl{A1}}
\BIC{$\varphi_{63}$}
\end{ptree}
\noindent$\Pi_{27}$:
\begin{ptree}
\AXC{$x \in \mathbb{R} \to y \in \mathbb{R} \to z \in \mathbb{R} \to 0 < y \to x \cdot y < z \to y \in \mathbb{R}$\lbl{A1}}
\AXC{$x \in \mathbb{R} \to y \in \mathbb{R} \to z \in \mathbb{R} \to 0 < y \to x \cdot y < z \to x \in \mathbb{R}$\lbl{A1}}
\AXC{$\stackrel{\Pi_{26}}{\vdots}$}\noLine
\UIC{$\varphi_{63}$}
\BIC{$\varphi_{62}$}
\BIC{$\varphi_{27} \to x \in \mathbb{R} \to y \in \mathbb{R} \to \varphi_{17}$}
\end{ptree}
\noindent$\Pi_{28}$:
\begin{ptree}
\AXC{$\stackrel{\Pi_{21}}{\vdots}$}\noLine
\UIC{$\varphi_{37}$}
\AXC{$\stackrel{\Pi_{25}}{\vdots}$}\noLine
\UIC{$x \in \mathbb{R} \to \varphi_{22}$}
\AXC{$\stackrel{\Pi_{10}}{\vdots}$}\noLine
\UIC{$\varphi_{27}$}
\AXC{$\stackrel{\Pi_{27}}{\vdots}$}\noLine
\UIC{$\varphi_{27} \to x \in \mathbb{R} \to y \in \mathbb{R} \to \varphi_{17}$}
\BIC{$x \in \mathbb{R} \to y \in \mathbb{R} \to \varphi_{17}$}
\AXC{$\varphi_{64}$\lbl{A1}}
\BIC{$(x \in \mathbb{R} \to \varphi_{22}) \to \varphi_{37} \to x \in \mathbb{R} \to \varphi_{23}$}
\BIC{$\varphi_{37} \to x \in \mathbb{R} \to \varphi_{23}$}
\BIC{$x \in \mathbb{R} \to \varphi_{23}$}
\AXC{$(x \in \mathbb{R} \to \varphi_{23}) \to \varphi_{55}$\lbl{A1}}
\BIC{$\varphi_{55}$}
\end{ptree}
\noindent$\Pi_{29}$:
\begin{ptree}
\AXC{$x \in \mathbb{R} \to y \in \mathbb{R} \to z \in \mathbb{R} \to 0 < y \to x \cdot y < z \to x \in \mathbb{R}$\lbl{A1}}
\AXC{$\stackrel{\Pi_{19}}{\vdots}$}\noLine
\UIC{$\varphi_{49}$}
\BIC{$x \in \mathbb{R} \to y \in \mathbb{R} \to z \in \mathbb{R} \to 0 < y \to x \cdot y < z \to x \cdot 1 = x$}
\AXC{$\stackrel{\Pi_{28}}{\vdots}$}\noLine
\UIC{$\varphi_{55}$}
\BIC{$\varphi_{46}$}
\end{ptree}
\noindent$\Pi_{30}$:
\begin{ptree}
\AXC{$\stackrel{\Pi_{16}}{\vdots}$}\noLine
\UIC{$x \in \mathbb{R} \to \varphi_{35}$}
\AXC{$\stackrel{\Pi_{18}}{\vdots}$}\noLine
\UIC{$x \cdot y \cdot y^{-1} = x \cdot 1 \to x \cdot 1 = x \to x \cdot y \cdot y^{-1} = x$}
\AXC{$\stackrel{\Pi_{29}}{\vdots}$}\noLine
\UIC{$\varphi_{46}$}
\BIC{$\varphi_{34}$}
\AXC{$\varphi_{34} \to (x \in \mathbb{R} \to \varphi_{35}) \to \varphi_{18} \to \varphi_{33}$\lbl{A1}}
\BIC{$(x \in \mathbb{R} \to \varphi_{35}) \to \varphi_{18} \to \varphi_{33}$}
\BIC{$\varphi_{18} \to \varphi_{33}$}
\end{ptree}
\noindent 主木:
\begin{ptree}
\AXC{$\stackrel{\Pi_{1}}{\vdots}$}\noLine
\UIC{$\varphi_{18}$}
\AXC{$\stackrel{\Pi_{30}}{\vdots}$}\noLine
\UIC{$\varphi_{18} \to \varphi_{33}$}
\BIC{$\varphi_{33}$}
\UIC{$\forall z\,(\varphi_{33})$}
\UIC{$\forall y\,\forall z\,(\varphi_{33})$}
\UIC{$\forall x\,\forall y\,\forall z\,(\varphi_{33})$}
\end{ptree}
\begin{abbrevlist}
\item $\varphi_{1} :\equiv \forall y\,(v_{631} = y \to y = v_{631}) \to v_{631} = v_{632} \to v_{632} = v_{631}$
\item $\varphi_{2} :\equiv \forall v_{640}\,\forall v_{641}\,(v_{639} = v_{640} \to v_{640} = v_{641} \to v_{639} = v_{641})$
\item $\varphi_{3} :\equiv \forall v_{637}\,\forall v_{638}\,(v_{636} = v_{637} \to v_{637} = v_{638} \to v_{636} = v_{638})$
\item $\varphi_{4} :\equiv \forall x\,\forall y\,\forall z\,(x = y \to y = z \to x = z) \to \forall y\,\forall z\,(v_{639} = y \to y = z \to v_{639} = z)$
\item $\varphi_{5} :\equiv \forall x\,\forall y\,\forall z\,(x = y \to y = z \to x = z) \to \forall y\,\forall z\,(v_{636} = y \to y = z \to v_{636} = z)$
\item $\varphi_{6} :\equiv z = w \to x \cdot (y \cdot y^{-1}) = x \cdot z \to x \cdot (y \cdot y^{-1}) = x \cdot w$
\item $\varphi_{7} :\equiv \forall v_{631}\,\forall v_{632}\,(v_{631} = v_{632} \to v_{632} = v_{631}) \to \forall v_{632}\,(y = v_{632} \to v_{632} = y)$
\item $\varphi_{8} :\equiv x \cdot (y \cdot y^{-1}) = x \cdot (y \cdot y^{-1}) \to x \cdot (y \cdot y^{-1}) = x \cdot 1$
\item $\varphi_{9} :\equiv \forall y\,\forall z\,(v_{639} = y \to y = z \to v_{639} = z) \to \forall z\,(v_{639} = v_{640} \to v_{640} = z \to v_{639} = z)$
\item $\varphi_{10} :\equiv \forall y\,\forall z\,(v_{636} = y \to y = z \to v_{636} = z) \to \forall z\,(v_{636} = v_{637} \to v_{637} = z \to v_{636} = z)$
\item $\varphi_{11} :\equiv \forall z\,(x \in \mathbb{R} \to y \in \mathbb{R} \to z \in \mathbb{R} \to x \cdot y \cdot z = x \cdot (y \cdot z))$
\item $\varphi_{12} :\equiv \forall z\,(v_{639} = v_{640} \to v_{640} = z \to v_{639} = z) \to v_{639} = v_{640} \to v_{640} = v_{641} \to v_{639} = v_{641}$
\item $\varphi_{13} :\equiv \forall z\,(v_{636} = v_{637} \to v_{637} = z \to v_{636} = z) \to v_{636} = v_{637} \to v_{637} = v_{638} \to v_{636} = v_{638}$
\item $\varphi_{14} :\equiv \forall v_{641}\,(x \cdot y \cdot y^{-1} = x \cdot 1 \to x \cdot 1 = v_{641} \to x \cdot y \cdot y^{-1} = v_{641})$
\item $\varphi_{15} :\equiv \forall v_{640}\,\forall v_{641}\,(x \cdot y \cdot y^{-1} = v_{640} \to v_{640} = v_{641} \to x \cdot y \cdot y^{-1} = v_{641})$
\item $\varphi_{16} :\equiv \forall v_{637}\,\forall v_{638}\,(x \cdot y \cdot y^{-1} = v_{637} \to v_{637} = v_{638} \to x \cdot y \cdot y^{-1} = v_{638})$
\item $\varphi_{17} :\equiv z \in \mathbb{R} \to 0 < y \to x \cdot y < z \to x \cdot y \cdot y^{-1} = x \cdot (y \cdot y^{-1})$
\item $\varphi_{18} :\equiv x \cdot y \cdot y^{-1} = x \to x \cdot y \cdot y^{-1} < z \cdot y^{-1} \to x < z \cdot y^{-1}$
\item $\varphi_{19} :\equiv v_{635} \in \mathbb{R} \to 0 < v_{635} \to v_{633} < v_{634} \to v_{633} \cdot v_{635} < v_{634} \cdot v_{635}$
\item $\varphi_{20} :\equiv \forall x \in \mathbb{R}\,(0 < x \to 0 < x^{-1}) \to y \in \mathbb{R} \to 0 < y \to 0 < y^{-1}$
\item $\varphi_{21} :\equiv \forall x \in \mathbb{R}\,(x \neq 0 \to x \cdot x^{-1} = 1) \to y \in \mathbb{R} \to y \neq 0 \to y \cdot y^{-1} = 1$
\item $\varphi_{22} :\equiv y \in \mathbb{R} \to z \in \mathbb{R} \to 0 < y \to x \cdot y < z \to x \cdot (y \cdot y^{-1}) = x \cdot 1$
\item $\varphi_{23} :\equiv y \in \mathbb{R} \to z \in \mathbb{R} \to 0 < y \to x \cdot y < z \to x \cdot y \cdot y^{-1} = x \cdot 1$
\item $\varphi_{24} :\equiv \forall w\,(x \cdot y \cdot y^{-1} = w \to x \cdot y \cdot y^{-1} < z \cdot y^{-1} \to w < z \cdot y^{-1})$
\item $\varphi_{25} :\equiv x \in \mathbb{R} \to y \in \mathbb{R} \to z \in \mathbb{R} \to 0 < y \to x \cdot y < z \to y \cdot y^{-1} = 1$
\item $\varphi_{26} :\equiv \forall x \in \mathbb{R}\,(x \neq 0 \to x^{-1} \in \mathbb{R}) \to y \in \mathbb{R} \to y \neq 0 \to y^{-1} \in \mathbb{R}$
\item $\varphi_{27} :\equiv x \in \mathbb{R} \to y \in \mathbb{R} \to z \in \mathbb{R} \to 0 < y \to x \cdot y < z \to y^{-1} \in \mathbb{R}$
\item $\varphi_{28} :\equiv x \in \mathbb{R} \to y \in \mathbb{R} \to z \in \mathbb{R} \to 0 < y \to x \cdot y < z \to x \cdot y \in \mathbb{R}$
\item $\varphi_{29} :\equiv x \in \mathbb{R} \to y \in \mathbb{R} \to y^{-1} \in \mathbb{R} \to x \cdot y \cdot y^{-1} = x \cdot (y \cdot y^{-1})$
\item $\varphi_{30} :\equiv y^{-1} \in \mathbb{R} \to 0 < y^{-1} \to x \cdot y < z \to x \cdot y \cdot y^{-1} < z \cdot y^{-1}$
\item $\varphi_{31} :\equiv (y = 0 \to 0 = y) \to x \in \mathbb{R} \to y \in \mathbb{R} \to z \in \mathbb{R} \to 0 < y \to x \cdot y < z \to y \neq 0$
\item $\varphi_{32} :\equiv \forall w\,(y \cdot y^{-1} = w \to x \cdot (y \cdot y^{-1}) = x \cdot (y \cdot y^{-1}) \to x \cdot (y \cdot y^{-1}) = x \cdot w)$
\item $\varphi_{33} :\equiv x \in \mathbb{R} \to y \in \mathbb{R} \to z \in \mathbb{R} \to 0 < y \to x \cdot y < z \to x < z \cdot y^{-1}$
\item $\varphi_{34} :\equiv x \in \mathbb{R} \to y \in \mathbb{R} \to z \in \mathbb{R} \to 0 < y \to x \cdot y < z \to x \cdot y \cdot y^{-1} = x$
\item $\varphi_{35} :\equiv y \in \mathbb{R} \to z \in \mathbb{R} \to 0 < y \to x \cdot y < z \to x \cdot y \cdot y^{-1} < z \cdot y^{-1}$
\item $\varphi_{36} :\equiv \forall x\,\forall y\,\forall z\,(x \in \mathbb{R} \to y \in \mathbb{R} \to z \in \mathbb{R} \to 0 < z \to x < y \to x \cdot z < y \cdot z)$
\item $\varphi_{37} :\equiv x \cdot y \cdot y^{-1} = x \cdot (y \cdot y^{-1}) \to x \cdot (y \cdot y^{-1}) = x \cdot 1 \to x \cdot y \cdot y^{-1} = x \cdot 1$
\item $\varphi_{38} :\equiv \forall v_{638}\,(x \cdot y \cdot y^{-1} = x \cdot (y \cdot y^{-1}) \to x \cdot (y \cdot y^{-1}) = v_{638} \to x \cdot y \cdot y^{-1} = v_{638})$
\item $\varphi_{39} :\equiv \forall y\,(x \in \mathbb{R} \to y \in \mathbb{R} \to x \cdot y \in \mathbb{R}) \to x \in \mathbb{R} \to y \in \mathbb{R} \to x \cdot y \in \mathbb{R}$
\item $\varphi_{40} :\equiv \forall y\,\forall z\,(v_{633} \in \mathbb{R} \to y \in \mathbb{R} \to z \in \mathbb{R} \to 0 < z \to v_{633} < y \to v_{633} \cdot z < y \cdot z)$
\item $\varphi_{41} :\equiv \forall x\,\forall y\,(x \in \mathbb{R} \to y \in \mathbb{R} \to x \cdot y \in \mathbb{R}) \to \forall y\,(x \in \mathbb{R} \to y \in \mathbb{R} \to x \cdot y \in \mathbb{R})$
\item $\varphi_{42} :\equiv \forall z\,(v_{633} \in \mathbb{R} \to v_{634} \in \mathbb{R} \to z \in \mathbb{R} \to 0 < z \to v_{633} < v_{634} \to v_{633} \cdot z < v_{634} \cdot z)$
\item $\varphi_{43} :\equiv \forall v_{635}\,(x \cdot y \in \mathbb{R} \to z \in \mathbb{R} \to v_{635} \in \mathbb{R} \to 0 < v_{635} \to x \cdot y < z \to x \cdot y \cdot v_{635} < z \cdot v_{635})$
\item $\varphi_{44} :\equiv \forall v_{634}\,\forall v_{635}\,(x \cdot y \in \mathbb{R} \to v_{634} \in \mathbb{R} \to v_{635} \in \mathbb{R} \to 0 < v_{635} \to x \cdot y < v_{634} \to x \cdot y \cdot v_{635} < v_{634} \cdot v_{635})$
\item $\varphi_{45} :\equiv (x \in \mathbb{R} \to y \in \mathbb{R} \to z \in \mathbb{R} \to 0 < y \to x \cdot y < z \to 0 < y) \to x \in \mathbb{R} \to y \in \mathbb{R} \to z \in \mathbb{R} \to 0 < y \to x \cdot y < z \to 0 \neq y$
\item $\varphi_{46} :\equiv (x \cdot y \cdot y^{-1} = x \cdot 1 \to x \cdot 1 = x \to x \cdot y \cdot y^{-1} = x) \to \varphi_{34}$
\item $\varphi_{47} :\equiv (x \in \mathbb{R} \to y \in \mathbb{R} \to z \in \mathbb{R} \to 0 < y \to x \cdot y < z \to 0 < y) \to x \in \mathbb{R} \to y \in \mathbb{R} \to z \in \mathbb{R} \to 0 < y \to x \cdot y < z \to 0 < y^{-1}$
\item $\varphi_{48} :\equiv (x \in \mathbb{R} \to y \in \mathbb{R} \to z \in \mathbb{R} \to 0 < y \to x \cdot y < z \to y \neq 0) \to \varphi_{25}$
\item $\varphi_{49} :\equiv (x \in \mathbb{R} \to y \in \mathbb{R} \to z \in \mathbb{R} \to 0 < y \to x \cdot y < z \to x \in \mathbb{R}) \to x \in \mathbb{R} \to y \in \mathbb{R} \to z \in \mathbb{R} \to 0 < y \to x \cdot y < z \to x \cdot 1 = x$
\item $\varphi_{50} :\equiv (x \in \mathbb{R} \to y \in \mathbb{R} \to z \in \mathbb{R} \to 0 < y \to x \cdot y < z \to y \neq 0) \to \varphi_{27}$
\item $\varphi_{51} :\equiv (x \in \mathbb{R} \to y \in \mathbb{R} \to z \in \mathbb{R} \to 0 < y \to x \cdot y < z \to 0 \neq y) \to \varphi_{31}$
\item $\varphi_{52} :\equiv (x \in \mathbb{R} \to y \in \mathbb{R} \to z \in \mathbb{R} \to 0 < y \to x \cdot y < z \to y \in \mathbb{R}) \to \varphi_{28}$
\item $\varphi_{53} :\equiv (x \in \mathbb{R} \to y \in \mathbb{R} \to z \in \mathbb{R} \to 0 < y \to x \cdot y < z \to x \cdot y < z) \to x \in \mathbb{R} \to \varphi_{35}$
\item $\varphi_{54} :\equiv x \cdot (y \cdot y^{-1}) = x \cdot (y \cdot y^{-1}) \to (y \cdot y^{-1} = 1 \to \varphi_{8}) \to x \in \mathbb{R} \to \varphi_{22}$
\item $\varphi_{55} :\equiv (x \in \mathbb{R} \to y \in \mathbb{R} \to z \in \mathbb{R} \to 0 < y \to x \cdot y < z \to x \cdot 1 = x) \to \varphi_{46}$
\item $\varphi_{56} :\equiv (x \in \mathbb{R} \to y \in \mathbb{R} \to z \in \mathbb{R} \to 0 < y \to x \cdot y < z \to y \in \mathbb{R}) \to \varphi_{47}$
\item $\varphi_{57} :\equiv (x \in \mathbb{R} \to y \in \mathbb{R} \to z \in \mathbb{R} \to 0 < y \to x \cdot y < z \to y \in \mathbb{R}) \to \varphi_{48}$
\item $\varphi_{58} :\equiv (x \in \mathbb{R} \to y \in \mathbb{R} \to z \in \mathbb{R} \to 0 < y \to x \cdot y < z \to y \in \mathbb{R}) \to \varphi_{50}$
\item $\varphi_{59} :\equiv (x \in \mathbb{R} \to y \in \mathbb{R} \to z \in \mathbb{R} \to 0 < y \to x \cdot y < z \to x \in \mathbb{R}) \to \varphi_{52}$
\item $\varphi_{60} :\equiv \forall v_{633}\,\forall v_{634}\,\forall v_{635}\,(v_{633} \in \mathbb{R} \to v_{634} \in \mathbb{R} \to \varphi_{19}) \to \varphi_{44}$
\item $\varphi_{61} :\equiv (x \in \mathbb{R} \to y \in \mathbb{R} \to z \in \mathbb{R} \to 0 < y \to x \cdot y < z \to 0 < y^{-1}) \to \varphi_{53}$
\item $\varphi_{62} :\equiv (x \in \mathbb{R} \to y \in \mathbb{R} \to z \in \mathbb{R} \to 0 < y \to x \cdot y < z \to y \in \mathbb{R}) \to \varphi_{27} \to x \in \mathbb{R} \to y \in \mathbb{R} \to \varphi_{17}$
\item $\varphi_{63} :\equiv (x \in \mathbb{R} \to y \in \mathbb{R} \to z \in \mathbb{R} \to 0 < y \to x \cdot y < z \to x \in \mathbb{R}) \to \varphi_{62}$
\item $\varphi_{64} :\equiv (x \in \mathbb{R} \to y \in \mathbb{R} \to \varphi_{17}) \to (x \in \mathbb{R} \to \varphi_{22}) \to \varphi_{37} \to x \in \mathbb{R} \to \varphi_{23}$
\item $\varphi_{65} :\equiv (x \in \mathbb{R} \to y \in \mathbb{R} \to z \in \mathbb{R} \to 0 < y \to x \cdot y < z \to z \in \mathbb{R}) \to \varphi_{27} \to \varphi_{61}$
\end{abbrevlist}
\end{fullproof}

% 自動生成 (prover/build.js): AbsLowerHalf  ノード数 204
\begin{fullproof}{[E26]}{$\forall x\,\forall \beta\,(x \in \mathbb{R} \to \beta \in \mathbb{R} \to \lvert x - \beta \rvert < \lvert \beta \rvert \cdot 2^{-1} \to \lvert \beta \rvert \cdot 2^{-1} < \lvert x \rvert)$}{204}\phantomsection\label{tree:AbsLowerHalf}
\noindent$\Pi_{1}$:
\begin{ptree}
\AXC{$x = x$\lbl{A5}}
\UIC{$\forall x\,(x = x)$}
\AXC{$\forall x\,(x = x) \to \lvert x \rvert = \lvert x \rvert$\lbl{A2}}
\BIC{$\lvert x \rvert = \lvert x \rvert$}
\AXC{$\lvert x \rvert = \lvert x \rvert \to \varphi_{8}$\lbl{A1}}
\BIC{$\varphi_{8}$}
\end{ptree}
\noindent$\Pi_{2}$:
\begin{ptree}
\AXC{$z = w \to z < \lvert \beta \rvert \cdot 2^{-1} \to w < \lvert \beta \rvert \cdot 2^{-1}$\lbl{A6}}
\UIC{$\forall w\,(z = w \to z < \lvert \beta \rvert \cdot 2^{-1} \to w < \lvert \beta \rvert \cdot 2^{-1})$}
\UIC{$\forall z\,\forall w\,(z = w \to z < \lvert \beta \rvert \cdot 2^{-1} \to w < \lvert \beta \rvert \cdot 2^{-1})$}
\AXC{$\varphi_{55}$\lbl{A2}}
\BIC{$\varphi_{10}$}
\AXC{$\varphi_{10} \to \varphi_{38}$\lbl{A2}}
\BIC{$\varphi_{38}$}
\end{ptree}
\noindent$\Pi_{3}$:
\begin{ptree}
\AXC{$\forall x\,\forall y\,(x \in \mathbb{R} \to y \in \mathbb{R} \to \lvert x - y \rvert = \lvert y - x \rvert)$\lbl{[A13]}}
\AXC{$\varphi_{47}$\lbl{A2}}
\BIC{$\forall y\,(x \in \mathbb{R} \to y \in \mathbb{R} \to \lvert x - y \rvert = \lvert y - x \rvert)$}
\AXC{$\varphi_{46}$\lbl{A2}}
\BIC{$x \in \mathbb{R} \to \beta \in \mathbb{R} \to \lvert x - \beta \rvert = \lvert \beta - x \rvert$}
\end{ptree}
\noindent$\Pi_{4}$:
\begin{ptree}
\AXC{$\stackrel{\Pi_{3}}{\vdots}$}\noLine
\UIC{$x \in \mathbb{R} \to \beta \in \mathbb{R} \to \lvert x - \beta \rvert = \lvert \beta - x \rvert$}
\AXC{$(x \in \mathbb{R} \to \beta \in \mathbb{R} \to \lvert x - \beta \rvert = \lvert \beta - x \rvert) \to \varphi_{74}$\lbl{A1}}
\BIC{$\varphi_{74}$}
\end{ptree}
\noindent$\Pi_{5}$:
\begin{ptree}
\AXC{$x \in \mathbb{R} \to \beta \in \mathbb{R} \to \lvert x - \beta \rvert < \lvert \beta \rvert \cdot 2^{-1} \to \beta \in \mathbb{R}$\lbl{A1}}
\AXC{$x \in \mathbb{R} \to \beta \in \mathbb{R} \to \lvert x - \beta \rvert < \lvert \beta \rvert \cdot 2^{-1} \to x \in \mathbb{R}$\lbl{A1}}
\AXC{$\stackrel{\Pi_{4}}{\vdots}$}\noLine
\UIC{$\varphi_{74}$}
\BIC{$\varphi_{63}$}
\BIC{$\varphi_{22}$}
\end{ptree}
\noindent$\Pi_{6}$:
\begin{ptree}
\AXC{$\stackrel{\Pi_{2}}{\vdots}$}\noLine
\UIC{$\varphi_{38}$}
\AXC{$x \in \mathbb{R} \to \varphi_{17}$\lbl{A1}}
\AXC{$\stackrel{\Pi_{5}}{\vdots}$}\noLine
\UIC{$\varphi_{22}$}
\AXC{$\varphi_{22} \to (x \in \mathbb{R} \to \varphi_{17}) \to \varphi_{38} \to x \in \mathbb{R} \to \varphi_{18}$\lbl{A1}}
\BIC{$(x \in \mathbb{R} \to \varphi_{17}) \to \varphi_{38} \to x \in \mathbb{R} \to \varphi_{18}$}
\BIC{$\varphi_{38} \to x \in \mathbb{R} \to \varphi_{18}$}
\BIC{$x \in \mathbb{R} \to \varphi_{18}$}
\end{ptree}
\noindent$\Pi_{7}$:
\begin{ptree}
\AXC{$\forall x \in \mathbb{R}\,(\lvert x \rvert \in \mathbb{R})$\lbl{[A5]}}
\AXC{$\forall x \in \mathbb{R}\,(\lvert x \rvert \in \mathbb{R}) \to x \in \mathbb{R} \to \lvert x \rvert \in \mathbb{R}$\lbl{A2}}
\BIC{$x \in \mathbb{R} \to \lvert x \rvert \in \mathbb{R}$}
\AXC{$(x \in \mathbb{R} \to \lvert x \rvert \in \mathbb{R}) \to \varphi_{59}$\lbl{A1}}
\BIC{$\varphi_{59}$}
\end{ptree}
\noindent$\Pi_{8}$:
\begin{ptree}
\AXC{$x \in \mathbb{R} \to \beta \in \mathbb{R} \to \lvert x - \beta \rvert < \lvert \beta \rvert \cdot 2^{-1} \to x \in \mathbb{R}$\lbl{A1}}
\AXC{$\stackrel{\Pi_{7}}{\vdots}$}\noLine
\UIC{$\varphi_{59}$}
\BIC{$x \in \mathbb{R} \to \beta \in \mathbb{R} \to \lvert x - \beta \rvert < \lvert \beta \rvert \cdot 2^{-1} \to \lvert x \rvert \in \mathbb{R}$}
\end{ptree}
\noindent$\Pi_{9}$:
\begin{ptree}
\AXC{$\forall x \in \mathbb{R}\,(\lvert x \rvert \in \mathbb{R})$\lbl{[A5]}}
\AXC{$\forall x \in \mathbb{R}\,(\lvert x \rvert \in \mathbb{R}) \to \beta \in \mathbb{R} \to \lvert \beta \rvert \in \mathbb{R}$\lbl{A2}}
\BIC{$\beta \in \mathbb{R} \to \lvert \beta \rvert \in \mathbb{R}$}
\AXC{$(\beta \in \mathbb{R} \to \lvert \beta \rvert \in \mathbb{R}) \to \varphi_{62}$\lbl{A1}}
\BIC{$\varphi_{62}$}
\end{ptree}
\noindent$\Pi_{10}$:
\begin{ptree}
\AXC{$x \in \mathbb{R} \to \beta \in \mathbb{R} \to \lvert x - \beta \rvert < \lvert \beta \rvert \cdot 2^{-1} \to \beta \in \mathbb{R}$\lbl{A1}}
\AXC{$\stackrel{\Pi_{9}}{\vdots}$}\noLine
\UIC{$\varphi_{62}$}
\BIC{$x \in \mathbb{R} \to \beta \in \mathbb{R} \to \lvert x - \beta \rvert < \lvert \beta \rvert \cdot 2^{-1} \to \lvert \beta \rvert \in \mathbb{R}$}
\end{ptree}
\noindent$\Pi_{11}$:
\begin{ptree}
\AXC{$\forall x\,\forall y\,(x \in \mathbb{R} \to y \in \mathbb{R} \to x \cdot y \in \mathbb{R})$\lbl{(R1)}}
\AXC{$\varphi_{45}$\lbl{A2}}
\BIC{$\forall y\,(\lvert \beta \rvert \in \mathbb{R} \to y \in \mathbb{R} \to \lvert \beta \rvert \cdot y \in \mathbb{R})$}
\AXC{$\varphi_{51}$\lbl{A2}}
\BIC{$\lvert \beta \rvert \in \mathbb{R} \to 2^{-1} \in \mathbb{R} \to \lvert \beta \rvert \cdot 2^{-1} \in \mathbb{R}$}
\AXC{$(\lvert \beta \rvert \in \mathbb{R} \to 2^{-1} \in \mathbb{R} \to \lvert \beta \rvert \cdot 2^{-1} \in \mathbb{R}) \to \varphi_{68}$\lbl{A1}}
\BIC{$\varphi_{68}$}
\end{ptree}
\noindent$\Pi_{12}$:
\begin{ptree}
\AXC{$2^{-1} \in \mathbb{R}$\lbl{[F31]}}
\AXC{$\stackrel{\Pi_{10}}{\vdots}$}\noLine
\UIC{$x \in \mathbb{R} \to \beta \in \mathbb{R} \to \lvert x - \beta \rvert < \lvert \beta \rvert \cdot 2^{-1} \to \lvert \beta \rvert \in \mathbb{R}$}
\AXC{$\stackrel{\Pi_{11}}{\vdots}$}\noLine
\UIC{$\varphi_{68}$}
\BIC{$2^{-1} \in \mathbb{R} \to \varphi_{24}$}
\BIC{$\varphi_{24}$}
\end{ptree}
\noindent$\Pi_{13}$:
\begin{ptree}
\AXC{$\forall x\,\forall y\,(x \in \mathbb{R} \to y \in \mathbb{R} \to x - y \in \mathbb{R})$\lbl{[F1]}}
\AXC{$\varphi_{37}$\lbl{A2}}
\BIC{$\forall y\,(\beta \in \mathbb{R} \to y \in \mathbb{R} \to \beta - y \in \mathbb{R})$}
\AXC{$\varphi_{31}$\lbl{A2}}
\BIC{$\beta \in \mathbb{R} \to x \in \mathbb{R} \to \beta - x \in \mathbb{R}$}
\AXC{$(\beta \in \mathbb{R} \to x \in \mathbb{R} \to \beta - x \in \mathbb{R}) \to \varphi_{73}$\lbl{A1}}
\BIC{$\varphi_{73}$}
\end{ptree}
\noindent$\Pi_{14}$:
\begin{ptree}
\AXC{$x \in \mathbb{R} \to \beta \in \mathbb{R} \to \lvert x - \beta \rvert < \lvert \beta \rvert \cdot 2^{-1} \to x \in \mathbb{R}$\lbl{A1}}
\AXC{$x \in \mathbb{R} \to \beta \in \mathbb{R} \to \lvert x - \beta \rvert < \lvert \beta \rvert \cdot 2^{-1} \to \beta \in \mathbb{R}$\lbl{A1}}
\AXC{$\stackrel{\Pi_{13}}{\vdots}$}\noLine
\UIC{$\varphi_{73}$}
\BIC{$\varphi_{61}$}
\BIC{$x \in \mathbb{R} \to \beta \in \mathbb{R} \to \lvert x - \beta \rvert < \lvert \beta \rvert \cdot 2^{-1} \to \beta - x \in \mathbb{R}$}
\end{ptree}
\noindent$\Pi_{15}$:
\begin{ptree}
\AXC{$\forall x \in \mathbb{R}\,(\lvert x \rvert \in \mathbb{R})$\lbl{[A5]}}
\AXC{$\forall x \in \mathbb{R}\,(\lvert x \rvert \in \mathbb{R}) \to \beta - x \in \mathbb{R} \to \lvert \beta - x \rvert \in \mathbb{R}$\lbl{A2}}
\BIC{$\beta - x \in \mathbb{R} \to \lvert \beta - x \rvert \in \mathbb{R}$}
\end{ptree}
\noindent$\Pi_{16}$:
\begin{ptree}
\AXC{$\stackrel{\Pi_{15}}{\vdots}$}\noLine
\UIC{$\beta - x \in \mathbb{R} \to \lvert \beta - x \rvert \in \mathbb{R}$}
\AXC{$(\beta - x \in \mathbb{R} \to \lvert \beta - x \rvert \in \mathbb{R}) \to \varphi_{64}$\lbl{A1}}
\BIC{$\varphi_{64}$}
\end{ptree}
\noindent$\Pi_{17}$:
\begin{ptree}
\AXC{$\varphi_{35}$\lbl{[F20]}}
\AXC{$\varphi_{35} \to \varphi_{52}$\lbl{A2}}
\BIC{$\varphi_{52}$}
\AXC{$\varphi_{52} \to \varphi_{65}$\lbl{A2}}
\BIC{$\varphi_{65}$}
\AXC{$\varphi_{65} \to \varphi_{67}$\lbl{A2}}
\BIC{$\varphi_{67}$}
\AXC{$\varphi_{67} \to \varphi_{70}$\lbl{A2}}
\BIC{$\varphi_{70}$}
\AXC{$\varphi_{70} \to \varphi_{19} \to \varphi_{24} \to \varphi_{84}$\lbl{A1}}
\BIC{$\varphi_{19} \to \varphi_{24} \to \varphi_{84}$}
\end{ptree}
\noindent$\Pi_{18}$:
\begin{ptree}
\AXC{$\stackrel{\Pi_{12}}{\vdots}$}\noLine
\UIC{$\varphi_{24}$}
\AXC{$\stackrel{\Pi_{14}}{\vdots}$}\noLine
\UIC{$x \in \mathbb{R} \to \beta \in \mathbb{R} \to \lvert x - \beta \rvert < \lvert \beta \rvert \cdot 2^{-1} \to \beta - x \in \mathbb{R}$}
\AXC{$\stackrel{\Pi_{16}}{\vdots}$}\noLine
\UIC{$\varphi_{64}$}
\BIC{$\varphi_{19}$}
\AXC{$\stackrel{\Pi_{17}}{\vdots}$}\noLine
\UIC{$\varphi_{19} \to \varphi_{24} \to \varphi_{84}$}
\BIC{$\varphi_{24} \to \varphi_{84}$}
\BIC{$\varphi_{84}$}
\end{ptree}
\noindent$\Pi_{19}$:
\begin{ptree}
\AXC{$\stackrel{\Pi_{8}}{\vdots}$}\noLine
\UIC{$x \in \mathbb{R} \to \beta \in \mathbb{R} \to \lvert x - \beta \rvert < \lvert \beta \rvert \cdot 2^{-1} \to \lvert x \rvert \in \mathbb{R}$}
\AXC{$\stackrel{\Pi_{8}}{\vdots}$}\noLine
\UIC{$x \in \mathbb{R} \to \beta \in \mathbb{R} \to \lvert x - \beta \rvert < \lvert \beta \rvert \cdot 2^{-1} \to \lvert x \rvert \in \mathbb{R}$}
\AXC{$\stackrel{\Pi_{18}}{\vdots}$}\noLine
\UIC{$\varphi_{84}$}
\BIC{$\varphi_{80}$}
\BIC{$(x \in \mathbb{R} \to \varphi_{18}) \to \varphi_{8} \to x \in \mathbb{R} \to \beta \in \mathbb{R} \to \varphi_{26}$}
\end{ptree}
\noindent$\Pi_{20}$:
\begin{ptree}
\AXC{$\forall a\,\forall b\,(a \in \mathbb{R} \to b \in \mathbb{R} \to \lvert a + b \rvert \leq \lvert a \rvert + \lvert b \rvert)$\lbl{命題1.2}}
\AXC{$\varphi_{54}$\lbl{A2}}
\BIC{$\varphi_{7}$}
\AXC{$\varphi_{58}$\lbl{A2}}
\BIC{$\beta - x \in \mathbb{R} \to x \in \mathbb{R} \to \lvert \beta - x + x \rvert \leq \lvert \beta - x \rvert + \lvert x \rvert$}
\AXC{$\varphi_{79}$\lbl{A1}}
\BIC{$\varphi_{75}$}
\end{ptree}
\noindent$\Pi_{21}$:
\begin{ptree}
\AXC{$x \in \mathbb{R} \to \beta \in \mathbb{R} \to \lvert x - \beta \rvert < \lvert \beta \rvert \cdot 2^{-1} \to x \in \mathbb{R}$\lbl{A1}}
\AXC{$\stackrel{\Pi_{14}}{\vdots}$}\noLine
\UIC{$x \in \mathbb{R} \to \beta \in \mathbb{R} \to \lvert x - \beta \rvert < \lvert \beta \rvert \cdot 2^{-1} \to \beta - x \in \mathbb{R}$}
\AXC{$\stackrel{\Pi_{20}}{\vdots}$}\noLine
\UIC{$\varphi_{75}$}
\BIC{$\varphi_{66}$}
\BIC{$x \in \mathbb{R} \to \beta \in \mathbb{R} \to \varphi_{9}$}
\end{ptree}
\noindent$\Pi_{22}$:
\begin{ptree}
\AXC{$\forall x\,\forall y\,(x \in \mathbb{R} \to y \in \mathbb{R} \to x - y + y = x)$\lbl{[F12]}}
\AXC{$\varphi_{33}$\lbl{A2}}
\BIC{$\forall y\,(\beta \in \mathbb{R} \to y \in \mathbb{R} \to \beta - y + y = \beta)$}
\AXC{$\varphi_{27}$\lbl{A2}}
\BIC{$\beta \in \mathbb{R} \to x \in \mathbb{R} \to \beta - x + x = \beta$}
\AXC{$(\beta \in \mathbb{R} \to x \in \mathbb{R} \to \beta - x + x = \beta) \to \varphi_{72}$\lbl{A1}}
\BIC{$\varphi_{72}$}
\end{ptree}
\noindent$\Pi_{23}$:
\begin{ptree}
\AXC{$x \in \mathbb{R} \to \beta \in \mathbb{R} \to \lvert x - \beta \rvert < \lvert \beta \rvert \cdot 2^{-1} \to x \in \mathbb{R}$\lbl{A1}}
\AXC{$x \in \mathbb{R} \to \beta \in \mathbb{R} \to \lvert x - \beta \rvert < \lvert \beta \rvert \cdot 2^{-1} \to \beta \in \mathbb{R}$\lbl{A1}}
\AXC{$\stackrel{\Pi_{22}}{\vdots}$}\noLine
\UIC{$\varphi_{72}$}
\BIC{$\varphi_{60}$}
\BIC{$\varphi_{11}$}
\end{ptree}
\noindent$\Pi_{24}$:
\begin{ptree}
\AXC{$\varphi_{4}$\lbl{A6}}
\UIC{$\forall w\,(\varphi_{4})$}
\UIC{$\forall z\,\forall w\,(\varphi_{4})$}
\AXC{$\forall z\,\forall w\,(\varphi_{4}) \to \varphi_{39}$\lbl{A2}}
\BIC{$\varphi_{39}$}
\AXC{$\varphi_{39} \to \beta - x + x = \beta \to \varphi_{13}$\lbl{A2}}
\BIC{$\beta - x + x = \beta \to \varphi_{13}$}
\AXC{$\stackrel{\Pi_{21}}{\vdots}$}\noLine
\UIC{$x \in \mathbb{R} \to \beta \in \mathbb{R} \to \varphi_{9}$}
\AXC{$\stackrel{\Pi_{23}}{\vdots}$}\noLine
\UIC{$\varphi_{11}$}
\AXC{$\varphi_{11} \to \varphi_{81}$\lbl{A1}}
\BIC{$\varphi_{81}$}
\BIC{$(\beta - x + x = \beta \to \varphi_{13}) \to (\beta - x + x = \beta \to \varphi_{6}) \to x \in \mathbb{R} \to \varphi_{14}$}
\BIC{$(\beta - x + x = \beta \to \varphi_{6}) \to x \in \mathbb{R} \to \varphi_{14}$}
\end{ptree}
\noindent$\Pi_{25}$:
\begin{ptree}
\AXC{$\varphi_{1}$\lbl{A6}}
\UIC{$\forall w\,(\varphi_{1})$}
\UIC{$\forall z\,\forall w\,(\varphi_{1})$}
\AXC{$\forall z\,\forall w\,(\varphi_{1}) \to \varphi_{32}$\lbl{A2}}
\BIC{$\varphi_{32}$}
\AXC{$\varphi_{32} \to \beta - x + x = \beta \to \varphi_{6}$\lbl{A2}}
\BIC{$\beta - x + x = \beta \to \varphi_{6}$}
\AXC{$\stackrel{\Pi_{24}}{\vdots}$}\noLine
\UIC{$(\beta - x + x = \beta \to \varphi_{6}) \to x \in \mathbb{R} \to \varphi_{14}$}
\BIC{$x \in \mathbb{R} \to \varphi_{14}$}
\end{ptree}
\noindent$\Pi_{26}$:
\begin{ptree}
\AXC{$\forall x\,\forall y\,\forall z\,(x \leq y \to y < z \to x < z)$\lbl{[O4]}}
\AXC{$\varphi_{2}$\lbl{A2}}
\BIC{$\forall y\,\forall z\,(\lvert \beta \rvert \leq y \to y < z \to \lvert \beta \rvert < z)$}
\AXC{$\forall y\,\forall z\,(\lvert \beta \rvert \leq y \to y < z \to \lvert \beta \rvert < z) \to \varphi_{12}$\lbl{A2}}
\BIC{$\varphi_{12}$}
\end{ptree}
\noindent$\Pi_{27}$:
\begin{ptree}
\AXC{$\stackrel{\Pi_{25}}{\vdots}$}\noLine
\UIC{$x \in \mathbb{R} \to \varphi_{14}$}
\AXC{$\stackrel{\Pi_{26}}{\vdots}$}\noLine
\UIC{$\varphi_{12}$}
\AXC{$\varphi_{12} \to \lvert \beta \rvert \leq \lvert \beta - x \rvert + \lvert x \rvert \to \varphi_{34}$\lbl{A2}}
\BIC{$\lvert \beta \rvert \leq \lvert \beta - x \rvert + \lvert x \rvert \to \varphi_{34}$}
\AXC{$\varphi_{83}$\lbl{A1}}
\BIC{$(x \in \mathbb{R} \to \varphi_{14}) \to (x \in \mathbb{R} \to \beta \in \mathbb{R} \to \varphi_{26}) \to x \in \mathbb{R} \to \varphi_{23}$}
\BIC{$(x \in \mathbb{R} \to \beta \in \mathbb{R} \to \varphi_{26}) \to x \in \mathbb{R} \to \varphi_{23}$}
\end{ptree}
\noindent$\Pi_{28}$:
\begin{ptree}
\AXC{$\stackrel{\Pi_{1}}{\vdots}$}\noLine
\UIC{$\varphi_{8}$}
\AXC{$\stackrel{\Pi_{6}}{\vdots}$}\noLine
\UIC{$x \in \mathbb{R} \to \varphi_{18}$}
\AXC{$\stackrel{\Pi_{19}}{\vdots}$}\noLine
\UIC{$(x \in \mathbb{R} \to \varphi_{18}) \to \varphi_{8} \to x \in \mathbb{R} \to \beta \in \mathbb{R} \to \varphi_{26}$}
\BIC{$\varphi_{8} \to x \in \mathbb{R} \to \beta \in \mathbb{R} \to \varphi_{26}$}
\BIC{$x \in \mathbb{R} \to \beta \in \mathbb{R} \to \varphi_{26}$}
\AXC{$\stackrel{\Pi_{27}}{\vdots}$}\noLine
\UIC{$(x \in \mathbb{R} \to \beta \in \mathbb{R} \to \varphi_{26}) \to x \in \mathbb{R} \to \varphi_{23}$}
\BIC{$x \in \mathbb{R} \to \varphi_{23}$}
\end{ptree}
\noindent$\Pi_{29}$:
\begin{ptree}
\AXC{$x = x$\lbl{A5}}
\AXC{$x = y \to x = x \to y = x$\lbl{A6}}
\AXC{$(x = y \to x = x \to y = x) \to x = x \to x = y \to y = x$\lbl{A1}}
\BIC{$x = x \to x = y \to y = x$}
\BIC{$x = y \to y = x$}
\UIC{$\forall y\,(x = y \to y = x)$}
\UIC{$\forall x\,\forall y\,(x = y \to y = x)$}
\end{ptree}
\noindent$\Pi_{30}$:
\begin{ptree}
\AXC{$\forall e \in \mathbb{R}\,(e \cdot 2^{-1} + e \cdot 2^{-1} = e)$\lbl{[F36]}}
\AXC{$\varphi_{49}$\lbl{A2}}
\BIC{$\lvert \beta \rvert \in \mathbb{R} \to \lvert \beta \rvert \cdot 2^{-1} + \lvert \beta \rvert \cdot 2^{-1} = \lvert \beta \rvert$}
\AXC{$(\lvert \beta \rvert \in \mathbb{R} \to \lvert \beta \rvert \cdot 2^{-1} + \lvert \beta \rvert \cdot 2^{-1} = \lvert \beta \rvert) \to \varphi_{69}$\lbl{A1}}
\BIC{$\varphi_{69}$}
\end{ptree}
\noindent$\Pi_{31}$:
\begin{ptree}
\AXC{$\stackrel{\Pi_{10}}{\vdots}$}\noLine
\UIC{$x \in \mathbb{R} \to \beta \in \mathbb{R} \to \lvert x - \beta \rvert < \lvert \beta \rvert \cdot 2^{-1} \to \lvert \beta \rvert \in \mathbb{R}$}
\AXC{$\stackrel{\Pi_{30}}{\vdots}$}\noLine
\UIC{$\varphi_{69}$}
\BIC{$x \in \mathbb{R} \to \beta \in \mathbb{R} \to \varphi_{15}$}
\AXC{$(x \in \mathbb{R} \to \beta \in \mathbb{R} \to \varphi_{15}) \to \varphi_{36} \to x \in \mathbb{R} \to \beta \in \mathbb{R} \to \varphi_{16}$\lbl{A1}}
\BIC{$\varphi_{36} \to x \in \mathbb{R} \to \beta \in \mathbb{R} \to \varphi_{16}$}
\end{ptree}
\noindent$\Pi_{32}$:
\begin{ptree}
\AXC{$\stackrel{\Pi_{29}}{\vdots}$}\noLine
\UIC{$\forall x\,\forall y\,(x = y \to y = x)$}
\AXC{$\forall x\,\forall y\,(x = y \to y = x) \to \varphi_{20}$\lbl{A2}}
\BIC{$\varphi_{20}$}
\AXC{$\varphi_{20} \to \varphi_{36}$\lbl{A2}}
\BIC{$\varphi_{36}$}
\AXC{$\stackrel{\Pi_{31}}{\vdots}$}\noLine
\UIC{$\varphi_{36} \to x \in \mathbb{R} \to \beta \in \mathbb{R} \to \varphi_{16}$}
\BIC{$x \in \mathbb{R} \to \beta \in \mathbb{R} \to \varphi_{16}$}
\end{ptree}
\noindent$\Pi_{33}$:
\begin{ptree}
\AXC{$\stackrel{\Pi_{28}}{\vdots}$}\noLine
\UIC{$x \in \mathbb{R} \to \varphi_{23}$}
\AXC{$\stackrel{\Pi_{32}}{\vdots}$}\noLine
\UIC{$x \in \mathbb{R} \to \beta \in \mathbb{R} \to \varphi_{16}$}
\AXC{$(x \in \mathbb{R} \to \beta \in \mathbb{R} \to \varphi_{16}) \to (x \in \mathbb{R} \to \varphi_{23}) \to \varphi_{76}$\lbl{A1}}
\BIC{$(x \in \mathbb{R} \to \varphi_{23}) \to \varphi_{76}$}
\BIC{$\varphi_{76}$}
\end{ptree}
\noindent$\Pi_{34}$:
\begin{ptree}
\AXC{$\varphi_{3}$\lbl{A6}}
\UIC{$\forall w\,(\varphi_{3})$}
\UIC{$\forall z\,\forall w\,(\varphi_{3})$}
\AXC{$\forall z\,\forall w\,(\varphi_{3}) \to \varphi_{21}$\lbl{A2}}
\BIC{$\varphi_{21}$}
\AXC{$\varphi_{21} \to \lvert \beta \rvert = \lvert \beta \rvert \cdot 2^{-1} + \lvert \beta \rvert \cdot 2^{-1} \to \varphi_{48}$\lbl{A2}}
\BIC{$\lvert \beta \rvert = \lvert \beta \rvert \cdot 2^{-1} + \lvert \beta \rvert \cdot 2^{-1} \to \varphi_{48}$}
\AXC{$\stackrel{\Pi_{33}}{\vdots}$}\noLine
\UIC{$\varphi_{76}$}
\BIC{$x \in \mathbb{R} \to \beta \in \mathbb{R} \to \varphi_{43}$}
\end{ptree}
\noindent$\Pi_{35}$:
\begin{ptree}
\AXC{$\forall x\,\forall y\,(x \in \mathbb{R} \to y \in \mathbb{R} \to x + y = y + x)$\lbl{(R3)}}
\AXC{$\forall x\,\forall y\,(x \in \mathbb{R} \to y \in \mathbb{R} \to x + y = y + x) \to \varphi_{28}$\lbl{A2}}
\BIC{$\varphi_{28}$}
\AXC{$\varphi_{28} \to \varphi_{41}$\lbl{A2}}
\BIC{$\varphi_{41}$}
\end{ptree}
\noindent$\Pi_{36}$:
\begin{ptree}
\AXC{$\stackrel{\Pi_{8}}{\vdots}$}\noLine
\UIC{$x \in \mathbb{R} \to \beta \in \mathbb{R} \to \lvert x - \beta \rvert < \lvert \beta \rvert \cdot 2^{-1} \to \lvert x \rvert \in \mathbb{R}$}
\AXC{$\stackrel{\Pi_{12}}{\vdots}$}\noLine
\UIC{$\varphi_{24}$}
\AXC{$\stackrel{\Pi_{35}}{\vdots}$}\noLine
\UIC{$\varphi_{41}$}
\AXC{$\varphi_{41} \to \varphi_{24} \to \varphi_{71}$\lbl{A1}}
\BIC{$\varphi_{24} \to \varphi_{71}$}
\BIC{$\varphi_{71}$}
\BIC{$x \in \mathbb{R} \to \beta \in \mathbb{R} \to \varphi_{30}$}
\end{ptree}
\noindent$\Pi_{37}$:
\begin{ptree}
\AXC{$\stackrel{\Pi_{34}}{\vdots}$}\noLine
\UIC{$x \in \mathbb{R} \to \beta \in \mathbb{R} \to \varphi_{43}$}
\AXC{$\stackrel{\Pi_{36}}{\vdots}$}\noLine
\UIC{$x \in \mathbb{R} \to \beta \in \mathbb{R} \to \varphi_{30}$}
\AXC{$(x \in \mathbb{R} \to \beta \in \mathbb{R} \to \varphi_{30}) \to (x \in \mathbb{R} \to \beta \in \mathbb{R} \to \varphi_{43}) \to \varphi_{78}$\lbl{A1}}
\BIC{$(x \in \mathbb{R} \to \beta \in \mathbb{R} \to \varphi_{43}) \to \varphi_{78}$}
\BIC{$\varphi_{78}$}
\end{ptree}
\noindent$\Pi_{38}$:
\begin{ptree}
\AXC{$z = w \to \varphi_{25}$\lbl{A6}}
\UIC{$\forall w\,(z = w \to \varphi_{25})$}
\UIC{$\forall z\,\forall w\,(z = w \to \varphi_{25})$}
\AXC{$\forall z\,\forall w\,(z = w \to \varphi_{25}) \to \varphi_{57}$\lbl{A2}}
\BIC{$\varphi_{57}$}
\AXC{$\varphi_{57} \to \lvert \beta \rvert \cdot 2^{-1} + \lvert x \rvert = \lvert x \rvert + \lvert \beta \rvert \cdot 2^{-1} \to \varphi_{56}$\lbl{A2}}
\BIC{$\lvert \beta \rvert \cdot 2^{-1} + \lvert x \rvert = \lvert x \rvert + \lvert \beta \rvert \cdot 2^{-1} \to \varphi_{56}$}
\AXC{$\stackrel{\Pi_{37}}{\vdots}$}\noLine
\UIC{$\varphi_{78}$}
\BIC{$x \in \mathbb{R} \to \beta \in \mathbb{R} \to \varphi_{44}$}
\end{ptree}
\noindent$\Pi_{39}$:
\begin{ptree}
\AXC{$\stackrel{\Pi_{12}}{\vdots}$}\noLine
\UIC{$\varphi_{24}$}
\AXC{$\varphi_{5}$\lbl{[E4]}}
\AXC{$\varphi_{5} \to \varphi_{50}$\lbl{A2}}
\BIC{$\varphi_{50}$}
\AXC{$\varphi_{50} \to \varphi_{53}$\lbl{A2}}
\BIC{$\varphi_{53}$}
\AXC{$\varphi_{77}$\lbl{A2}}
\BIC{$\lvert \beta \rvert \cdot 2^{-1} \in \mathbb{R} \to \lvert x \rvert \in \mathbb{R} \to \lvert \beta \rvert \cdot 2^{-1} \in \mathbb{R} \to \varphi_{40}$}
\AXC{$\varphi_{85}$\lbl{A1}}
\BIC{$\varphi_{24} \to \varphi_{82}$}
\BIC{$\varphi_{82}$}
\end{ptree}
\noindent 主木:
\begin{ptree}
\AXC{$\stackrel{\Pi_{38}}{\vdots}$}\noLine
\UIC{$x \in \mathbb{R} \to \beta \in \mathbb{R} \to \varphi_{44}$}
\AXC{$\stackrel{\Pi_{12}}{\vdots}$}\noLine
\UIC{$\varphi_{24}$}
\AXC{$\stackrel{\Pi_{8}}{\vdots}$}\noLine
\UIC{$x \in \mathbb{R} \to \beta \in \mathbb{R} \to \lvert x - \beta \rvert < \lvert \beta \rvert \cdot 2^{-1} \to \lvert x \rvert \in \mathbb{R}$}
\AXC{$\stackrel{\Pi_{39}}{\vdots}$}\noLine
\UIC{$\varphi_{82}$}
\BIC{$\varphi_{24} \to (x \in \mathbb{R} \to \beta \in \mathbb{R} \to \varphi_{44}) \to \varphi_{29}$}
\BIC{$(x \in \mathbb{R} \to \beta \in \mathbb{R} \to \varphi_{44}) \to \varphi_{29}$}
\BIC{$\varphi_{29}$}
\UIC{$\forall \beta\,(\varphi_{29})$}
\UIC{$\forall x\,\forall \beta\,(\varphi_{29})$}
\end{ptree}
\begin{abbrevlist}
\item $\varphi_{1} :\equiv z = w \to \lvert z \rvert = \lvert \beta - x \rvert + \lvert x \rvert \to \lvert w \rvert = \lvert \beta - x \rvert + \lvert x \rvert$
\item $\varphi_{2} :\equiv \forall x\,\forall y\,\forall z\,(x \leq y \to y < z \to x < z) \to \forall y\,\forall z\,(\lvert \beta \rvert \leq y \to y < z \to \lvert \beta \rvert < z)$
\item $\varphi_{3} :\equiv z = w \to z < \lvert \beta \rvert \cdot 2^{-1} + \lvert x \rvert \to w < \lvert \beta \rvert \cdot 2^{-1} + \lvert x \rvert$
\item $\varphi_{4} :\equiv z = w \to \lvert z \rvert < \lvert \beta - x \rvert + \lvert x \rvert \to \lvert w \rvert < \lvert \beta - x \rvert + \lvert x \rvert$
\item $\varphi_{5} :\equiv \forall x\,\forall y\,\forall z\,(x \in \mathbb{R} \to y \in \mathbb{R} \to z \in \mathbb{R} \to x + z < y + z \to x < y)$
\item $\varphi_{6} :\equiv \lvert \beta - x + x \rvert = \lvert \beta - x \rvert + \lvert x \rvert \to \lvert \beta \rvert = \lvert \beta - x \rvert + \lvert x \rvert$
\item $\varphi_{7} :\equiv \forall b\,(\beta - x \in \mathbb{R} \to b \in \mathbb{R} \to \lvert \beta - x + b \rvert \leq \lvert \beta - x \rvert + \lvert b \rvert)$
\item $\varphi_{8} :\equiv x \in \mathbb{R} \to \beta \in \mathbb{R} \to \lvert x - \beta \rvert < \lvert \beta \rvert \cdot 2^{-1} \to \lvert x \rvert \leq \lvert x \rvert$
\item $\varphi_{9} :\equiv \lvert x - \beta \rvert < \lvert \beta \rvert \cdot 2^{-1} \to \lvert \beta - x + x \rvert \leq \lvert \beta - x \rvert + \lvert x \rvert$
\item $\varphi_{10} :\equiv \forall w\,(\lvert x - \beta \rvert = w \to \lvert x - \beta \rvert < \lvert \beta \rvert \cdot 2^{-1} \to w < \lvert \beta \rvert \cdot 2^{-1})$
\item $\varphi_{11} :\equiv x \in \mathbb{R} \to \beta \in \mathbb{R} \to \lvert x - \beta \rvert < \lvert \beta \rvert \cdot 2^{-1} \to \beta - x + x = \beta$
\item $\varphi_{12} :\equiv \forall z\,(\lvert \beta \rvert \leq \lvert \beta - x \rvert + \lvert x \rvert \to \lvert \beta - x \rvert + \lvert x \rvert < z \to \lvert \beta \rvert < z)$
\item $\varphi_{13} :\equiv \lvert \beta - x + x \rvert < \lvert \beta - x \rvert + \lvert x \rvert \to \lvert \beta \rvert < \lvert \beta - x \rvert + \lvert x \rvert$
\item $\varphi_{14} :\equiv \beta \in \mathbb{R} \to \lvert x - \beta \rvert < \lvert \beta \rvert \cdot 2^{-1} \to \lvert \beta \rvert \leq \lvert \beta - x \rvert + \lvert x \rvert$
\item $\varphi_{15} :\equiv \lvert x - \beta \rvert < \lvert \beta \rvert \cdot 2^{-1} \to \lvert \beta \rvert \cdot 2^{-1} + \lvert \beta \rvert \cdot 2^{-1} = \lvert \beta \rvert$
\item $\varphi_{16} :\equiv \lvert x - \beta \rvert < \lvert \beta \rvert \cdot 2^{-1} \to \lvert \beta \rvert = \lvert \beta \rvert \cdot 2^{-1} + \lvert \beta \rvert \cdot 2^{-1}$
\item $\varphi_{17} :\equiv \beta \in \mathbb{R} \to \lvert x - \beta \rvert < \lvert \beta \rvert \cdot 2^{-1} \to \lvert x - \beta \rvert < \lvert \beta \rvert \cdot 2^{-1}$
\item $\varphi_{18} :\equiv \beta \in \mathbb{R} \to \lvert x - \beta \rvert < \lvert \beta \rvert \cdot 2^{-1} \to \lvert \beta - x \rvert < \lvert \beta \rvert \cdot 2^{-1}$
\item $\varphi_{19} :\equiv x \in \mathbb{R} \to \beta \in \mathbb{R} \to \lvert x - \beta \rvert < \lvert \beta \rvert \cdot 2^{-1} \to \lvert \beta - x \rvert \in \mathbb{R}$
\item $\varphi_{20} :\equiv \forall y\,(\lvert \beta \rvert \cdot 2^{-1} + \lvert \beta \rvert \cdot 2^{-1} = y \to y = \lvert \beta \rvert \cdot 2^{-1} + \lvert \beta \rvert \cdot 2^{-1})$
\item $\varphi_{21} :\equiv \forall w\,(\lvert \beta \rvert = w \to \lvert \beta \rvert < \lvert \beta \rvert \cdot 2^{-1} + \lvert x \rvert \to w < \lvert \beta \rvert \cdot 2^{-1} + \lvert x \rvert)$
\item $\varphi_{22} :\equiv x \in \mathbb{R} \to \beta \in \mathbb{R} \to \lvert x - \beta \rvert < \lvert \beta \rvert \cdot 2^{-1} \to \lvert x - \beta \rvert = \lvert \beta - x \rvert$
\item $\varphi_{23} :\equiv \beta \in \mathbb{R} \to \lvert x - \beta \rvert < \lvert \beta \rvert \cdot 2^{-1} \to \lvert \beta \rvert < \lvert \beta \rvert \cdot 2^{-1} + \lvert x \rvert$
\item $\varphi_{24} :\equiv x \in \mathbb{R} \to \beta \in \mathbb{R} \to \lvert x - \beta \rvert < \lvert \beta \rvert \cdot 2^{-1} \to \lvert \beta \rvert \cdot 2^{-1} \in \mathbb{R}$
\item $\varphi_{25} :\equiv \lvert \beta \rvert \cdot 2^{-1} + \lvert \beta \rvert \cdot 2^{-1} < z \to \lvert \beta \rvert \cdot 2^{-1} + \lvert \beta \rvert \cdot 2^{-1} < w$
\item $\varphi_{26} :\equiv \lvert x - \beta \rvert < \lvert \beta \rvert \cdot 2^{-1} \to \lvert \beta - x \rvert + \lvert x \rvert < \lvert \beta \rvert \cdot 2^{-1} + \lvert x \rvert$
\item $\varphi_{27} :\equiv \forall y\,(\beta \in \mathbb{R} \to y \in \mathbb{R} \to \beta - y + y = \beta) \to \beta \in \mathbb{R} \to x \in \mathbb{R} \to \beta - x + x = \beta$
\item $\varphi_{28} :\equiv \forall y\,(\lvert \beta \rvert \cdot 2^{-1} \in \mathbb{R} \to y \in \mathbb{R} \to \lvert \beta \rvert \cdot 2^{-1} + y = y + \lvert \beta \rvert \cdot 2^{-1})$
\item $\varphi_{29} :\equiv x \in \mathbb{R} \to \beta \in \mathbb{R} \to \lvert x - \beta \rvert < \lvert \beta \rvert \cdot 2^{-1} \to \lvert \beta \rvert \cdot 2^{-1} < \lvert x \rvert$
\item $\varphi_{30} :\equiv \lvert x - \beta \rvert < \lvert \beta \rvert \cdot 2^{-1} \to \lvert \beta \rvert \cdot 2^{-1} + \lvert x \rvert = \lvert x \rvert + \lvert \beta \rvert \cdot 2^{-1}$
\item $\varphi_{31} :\equiv \forall y\,(\beta \in \mathbb{R} \to y \in \mathbb{R} \to \beta - y \in \mathbb{R}) \to \beta \in \mathbb{R} \to x \in \mathbb{R} \to \beta - x \in \mathbb{R}$
\item $\varphi_{32} :\equiv \forall w\,(\beta - x + x = w \to \lvert \beta - x + x \rvert = \lvert \beta - x \rvert + \lvert x \rvert \to \lvert w \rvert = \lvert \beta - x \rvert + \lvert x \rvert)$
\item $\varphi_{33} :\equiv \forall x\,\forall y\,(x \in \mathbb{R} \to y \in \mathbb{R} \to x - y + y = x) \to \forall y\,(\beta \in \mathbb{R} \to y \in \mathbb{R} \to \beta - y + y = \beta)$
\item $\varphi_{34} :\equiv \lvert \beta - x \rvert + \lvert x \rvert < \lvert \beta \rvert \cdot 2^{-1} + \lvert x \rvert \to \lvert \beta \rvert < \lvert \beta \rvert \cdot 2^{-1} + \lvert x \rvert$
\item $\varphi_{35} :\equiv \forall x\,\forall y\,\forall u\,\forall v\,(x \in \mathbb{R} \to y \in \mathbb{R} \to u \in \mathbb{R} \to v \in \mathbb{R} \to x < y \to u \leq v \to x + u < y + v)$
\item $\varphi_{36} :\equiv \lvert \beta \rvert \cdot 2^{-1} + \lvert \beta \rvert \cdot 2^{-1} = \lvert \beta \rvert \to \lvert \beta \rvert = \lvert \beta \rvert \cdot 2^{-1} + \lvert \beta \rvert \cdot 2^{-1}$
\item $\varphi_{37} :\equiv \forall x\,\forall y\,(x \in \mathbb{R} \to y \in \mathbb{R} \to x - y \in \mathbb{R}) \to \forall y\,(\beta \in \mathbb{R} \to y \in \mathbb{R} \to \beta - y \in \mathbb{R})$
\item $\varphi_{38} :\equiv \lvert x - \beta \rvert = \lvert \beta - x \rvert \to \lvert x - \beta \rvert < \lvert \beta \rvert \cdot 2^{-1} \to \lvert \beta - x \rvert < \lvert \beta \rvert \cdot 2^{-1}$
\item $\varphi_{39} :\equiv \forall w\,(\beta - x + x = w \to \lvert \beta - x + x \rvert < \lvert \beta - x \rvert + \lvert x \rvert \to \lvert w \rvert < \lvert \beta - x \rvert + \lvert x \rvert)$
\item $\varphi_{40} :\equiv \lvert \beta \rvert \cdot 2^{-1} + \lvert \beta \rvert \cdot 2^{-1} < \lvert x \rvert + \lvert \beta \rvert \cdot 2^{-1} \to \lvert \beta \rvert \cdot 2^{-1} < \lvert x \rvert$
\item $\varphi_{41} :\equiv \lvert \beta \rvert \cdot 2^{-1} \in \mathbb{R} \to \lvert x \rvert \in \mathbb{R} \to \lvert \beta \rvert \cdot 2^{-1} + \lvert x \rvert = \lvert x \rvert + \lvert \beta \rvert \cdot 2^{-1}$
\item $\varphi_{42} :\equiv \lvert \beta - x \rvert < \lvert \beta \rvert \cdot 2^{-1} \to \lvert x \rvert \leq \lvert x \rvert \to \lvert \beta - x \rvert + \lvert x \rvert < \lvert \beta \rvert \cdot 2^{-1} + \lvert x \rvert$
\item $\varphi_{43} :\equiv \lvert x - \beta \rvert < \lvert \beta \rvert \cdot 2^{-1} \to \lvert \beta \rvert \cdot 2^{-1} + \lvert \beta \rvert \cdot 2^{-1} < \lvert \beta \rvert \cdot 2^{-1} + \lvert x \rvert$
\item $\varphi_{44} :\equiv \lvert x - \beta \rvert < \lvert \beta \rvert \cdot 2^{-1} \to \lvert \beta \rvert \cdot 2^{-1} + \lvert \beta \rvert \cdot 2^{-1} < \lvert x \rvert + \lvert \beta \rvert \cdot 2^{-1}$
\item $\varphi_{45} :\equiv \forall x\,\forall y\,(x \in \mathbb{R} \to y \in \mathbb{R} \to x \cdot y \in \mathbb{R}) \to \forall y\,(\lvert \beta \rvert \in \mathbb{R} \to y \in \mathbb{R} \to \lvert \beta \rvert \cdot y \in \mathbb{R})$
\item $\varphi_{46} :\equiv \forall y\,(x \in \mathbb{R} \to y \in \mathbb{R} \to \lvert x - y \rvert = \lvert y - x \rvert) \to x \in \mathbb{R} \to \beta \in \mathbb{R} \to \lvert x - \beta \rvert = \lvert \beta - x \rvert$
\item $\varphi_{47} :\equiv \forall x\,\forall y\,(x \in \mathbb{R} \to y \in \mathbb{R} \to \lvert x - y \rvert = \lvert y - x \rvert) \to \forall y\,(x \in \mathbb{R} \to y \in \mathbb{R} \to \lvert x - y \rvert = \lvert y - x \rvert)$
\item $\varphi_{48} :\equiv \lvert \beta \rvert < \lvert \beta \rvert \cdot 2^{-1} + \lvert x \rvert \to \lvert \beta \rvert \cdot 2^{-1} + \lvert \beta \rvert \cdot 2^{-1} < \lvert \beta \rvert \cdot 2^{-1} + \lvert x \rvert$
\item $\varphi_{49} :\equiv \forall e \in \mathbb{R}\,(e \cdot 2^{-1} + e \cdot 2^{-1} = e) \to \lvert \beta \rvert \in \mathbb{R} \to \lvert \beta \rvert \cdot 2^{-1} + \lvert \beta \rvert \cdot 2^{-1} = \lvert \beta \rvert$
\item $\varphi_{50} :\equiv \forall y\,\forall z\,(\lvert \beta \rvert \cdot 2^{-1} \in \mathbb{R} \to y \in \mathbb{R} \to z \in \mathbb{R} \to \lvert \beta \rvert \cdot 2^{-1} + z < y + z \to \lvert \beta \rvert \cdot 2^{-1} < y)$
\item $\varphi_{51} :\equiv \forall y\,(\lvert \beta \rvert \in \mathbb{R} \to y \in \mathbb{R} \to \lvert \beta \rvert \cdot y \in \mathbb{R}) \to \lvert \beta \rvert \in \mathbb{R} \to 2^{-1} \in \mathbb{R} \to \lvert \beta \rvert \cdot 2^{-1} \in \mathbb{R}$
\item $\varphi_{52} :\equiv \forall y\,\forall u\,\forall v\,(\lvert \beta - x \rvert \in \mathbb{R} \to y \in \mathbb{R} \to u \in \mathbb{R} \to v \in \mathbb{R} \to \lvert \beta - x \rvert < y \to u \leq v \to \lvert \beta - x \rvert + u < y + v)$
\item $\varphi_{53} :\equiv \forall z\,(\lvert \beta \rvert \cdot 2^{-1} \in \mathbb{R} \to \lvert x \rvert \in \mathbb{R} \to z \in \mathbb{R} \to \lvert \beta \rvert \cdot 2^{-1} + z < \lvert x \rvert + z \to \lvert \beta \rvert \cdot 2^{-1} < \lvert x \rvert)$
\item $\varphi_{54} :\equiv \forall a\,\forall b\,(a \in \mathbb{R} \to b \in \mathbb{R} \to \lvert a + b \rvert \leq \lvert a \rvert + \lvert b \rvert) \to \varphi_{7}$
\item $\varphi_{55} :\equiv \forall z\,\forall w\,(z = w \to z < \lvert \beta \rvert \cdot 2^{-1} \to w < \lvert \beta \rvert \cdot 2^{-1}) \to \varphi_{10}$
\item $\varphi_{56} :\equiv \lvert \beta \rvert \cdot 2^{-1} + \lvert \beta \rvert \cdot 2^{-1} < \lvert \beta \rvert \cdot 2^{-1} + \lvert x \rvert \to \lvert \beta \rvert \cdot 2^{-1} + \lvert \beta \rvert \cdot 2^{-1} < \lvert x \rvert + \lvert \beta \rvert \cdot 2^{-1}$
\item $\varphi_{57} :\equiv \forall w\,(\lvert \beta \rvert \cdot 2^{-1} + \lvert x \rvert = w \to \lvert \beta \rvert \cdot 2^{-1} + \lvert \beta \rvert \cdot 2^{-1} < \lvert \beta \rvert \cdot 2^{-1} + \lvert x \rvert \to \lvert \beta \rvert \cdot 2^{-1} + \lvert \beta \rvert \cdot 2^{-1} < w)$
\item $\varphi_{58} :\equiv \varphi_{7} \to \beta - x \in \mathbb{R} \to x \in \mathbb{R} \to \lvert \beta - x + x \rvert \leq \lvert \beta - x \rvert + \lvert x \rvert$
\item $\varphi_{59} :\equiv (x \in \mathbb{R} \to \beta \in \mathbb{R} \to \lvert x - \beta \rvert < \lvert \beta \rvert \cdot 2^{-1} \to x \in \mathbb{R}) \to x \in \mathbb{R} \to \beta \in \mathbb{R} \to \lvert x - \beta \rvert < \lvert \beta \rvert \cdot 2^{-1} \to \lvert x \rvert \in \mathbb{R}$
\item $\varphi_{60} :\equiv (x \in \mathbb{R} \to \beta \in \mathbb{R} \to \lvert x - \beta \rvert < \lvert \beta \rvert \cdot 2^{-1} \to x \in \mathbb{R}) \to \varphi_{11}$
\item $\varphi_{61} :\equiv (x \in \mathbb{R} \to \beta \in \mathbb{R} \to \lvert x - \beta \rvert < \lvert \beta \rvert \cdot 2^{-1} \to x \in \mathbb{R}) \to x \in \mathbb{R} \to \beta \in \mathbb{R} \to \lvert x - \beta \rvert < \lvert \beta \rvert \cdot 2^{-1} \to \beta - x \in \mathbb{R}$
\item $\varphi_{62} :\equiv (x \in \mathbb{R} \to \beta \in \mathbb{R} \to \lvert x - \beta \rvert < \lvert \beta \rvert \cdot 2^{-1} \to \beta \in \mathbb{R}) \to x \in \mathbb{R} \to \beta \in \mathbb{R} \to \lvert x - \beta \rvert < \lvert \beta \rvert \cdot 2^{-1} \to \lvert \beta \rvert \in \mathbb{R}$
\item $\varphi_{63} :\equiv (x \in \mathbb{R} \to \beta \in \mathbb{R} \to \lvert x - \beta \rvert < \lvert \beta \rvert \cdot 2^{-1} \to \beta \in \mathbb{R}) \to \varphi_{22}$
\item $\varphi_{64} :\equiv (x \in \mathbb{R} \to \beta \in \mathbb{R} \to \lvert x - \beta \rvert < \lvert \beta \rvert \cdot 2^{-1} \to \beta - x \in \mathbb{R}) \to \varphi_{19}$
\item $\varphi_{65} :\equiv \forall u\,\forall v\,(\lvert \beta - x \rvert \in \mathbb{R} \to \lvert \beta \rvert \cdot 2^{-1} \in \mathbb{R} \to u \in \mathbb{R} \to v \in \mathbb{R} \to \lvert \beta - x \rvert < \lvert \beta \rvert \cdot 2^{-1} \to u \leq v \to \lvert \beta - x \rvert + u < \lvert \beta \rvert \cdot 2^{-1} + v)$
\item $\varphi_{66} :\equiv (x \in \mathbb{R} \to \beta \in \mathbb{R} \to \lvert x - \beta \rvert < \lvert \beta \rvert \cdot 2^{-1} \to x \in \mathbb{R}) \to x \in \mathbb{R} \to \beta \in \mathbb{R} \to \varphi_{9}$
\item $\varphi_{67} :\equiv \forall v\,(\lvert \beta - x \rvert \in \mathbb{R} \to \lvert \beta \rvert \cdot 2^{-1} \in \mathbb{R} \to \lvert x \rvert \in \mathbb{R} \to v \in \mathbb{R} \to \lvert \beta - x \rvert < \lvert \beta \rvert \cdot 2^{-1} \to \lvert x \rvert \leq v \to \lvert \beta - x \rvert + \lvert x \rvert < \lvert \beta \rvert \cdot 2^{-1} + v)$
\item $\varphi_{68} :\equiv (x \in \mathbb{R} \to \beta \in \mathbb{R} \to \lvert x - \beta \rvert < \lvert \beta \rvert \cdot 2^{-1} \to \lvert \beta \rvert \in \mathbb{R}) \to 2^{-1} \in \mathbb{R} \to \varphi_{24}$
\item $\varphi_{69} :\equiv (x \in \mathbb{R} \to \beta \in \mathbb{R} \to \lvert x - \beta \rvert < \lvert \beta \rvert \cdot 2^{-1} \to \lvert \beta \rvert \in \mathbb{R}) \to x \in \mathbb{R} \to \beta \in \mathbb{R} \to \varphi_{15}$
\item $\varphi_{70} :\equiv \lvert \beta - x \rvert \in \mathbb{R} \to \lvert \beta \rvert \cdot 2^{-1} \in \mathbb{R} \to \lvert x \rvert \in \mathbb{R} \to \lvert x \rvert \in \mathbb{R} \to \varphi_{42}$
\item $\varphi_{71} :\equiv (x \in \mathbb{R} \to \beta \in \mathbb{R} \to \lvert x - \beta \rvert < \lvert \beta \rvert \cdot 2^{-1} \to \lvert x \rvert \in \mathbb{R}) \to x \in \mathbb{R} \to \beta \in \mathbb{R} \to \varphi_{30}$
\item $\varphi_{72} :\equiv (x \in \mathbb{R} \to \beta \in \mathbb{R} \to \lvert x - \beta \rvert < \lvert \beta \rvert \cdot 2^{-1} \to \beta \in \mathbb{R}) \to \varphi_{60}$
\item $\varphi_{73} :\equiv (x \in \mathbb{R} \to \beta \in \mathbb{R} \to \lvert x - \beta \rvert < \lvert \beta \rvert \cdot 2^{-1} \to \beta \in \mathbb{R}) \to \varphi_{61}$
\item $\varphi_{74} :\equiv (x \in \mathbb{R} \to \beta \in \mathbb{R} \to \lvert x - \beta \rvert < \lvert \beta \rvert \cdot 2^{-1} \to x \in \mathbb{R}) \to \varphi_{63}$
\item $\varphi_{75} :\equiv (x \in \mathbb{R} \to \beta \in \mathbb{R} \to \lvert x - \beta \rvert < \lvert \beta \rvert \cdot 2^{-1} \to \beta - x \in \mathbb{R}) \to \varphi_{66}$
\item $\varphi_{76} :\equiv (\lvert \beta \rvert = \lvert \beta \rvert \cdot 2^{-1} + \lvert \beta \rvert \cdot 2^{-1} \to \varphi_{48}) \to x \in \mathbb{R} \to \beta \in \mathbb{R} \to \varphi_{43}$
\item $\varphi_{77} :\equiv \varphi_{53} \to \lvert \beta \rvert \cdot 2^{-1} \in \mathbb{R} \to \lvert x \rvert \in \mathbb{R} \to \lvert \beta \rvert \cdot 2^{-1} \in \mathbb{R} \to \varphi_{40}$
\item $\varphi_{78} :\equiv (\lvert \beta \rvert \cdot 2^{-1} + \lvert x \rvert = \lvert x \rvert + \lvert \beta \rvert \cdot 2^{-1} \to \varphi_{56}) \to x \in \mathbb{R} \to \beta \in \mathbb{R} \to \varphi_{44}$
\item $\varphi_{79} :\equiv (\beta - x \in \mathbb{R} \to x \in \mathbb{R} \to \lvert \beta - x + x \rvert \leq \lvert \beta - x \rvert + \lvert x \rvert) \to \varphi_{75}$
\item $\varphi_{80} :\equiv (x \in \mathbb{R} \to \beta \in \mathbb{R} \to \lvert x - \beta \rvert < \lvert \beta \rvert \cdot 2^{-1} \to \lvert x \rvert \in \mathbb{R}) \to (x \in \mathbb{R} \to \varphi_{18}) \to \varphi_{8} \to x \in \mathbb{R} \to \beta \in \mathbb{R} \to \varphi_{26}$
\item $\varphi_{81} :\equiv (x \in \mathbb{R} \to \beta \in \mathbb{R} \to \varphi_{9}) \to (\beta - x + x = \beta \to \varphi_{13}) \to (\beta - x + x = \beta \to \varphi_{6}) \to x \in \mathbb{R} \to \varphi_{14}$
\item $\varphi_{82} :\equiv (x \in \mathbb{R} \to \beta \in \mathbb{R} \to \lvert x - \beta \rvert < \lvert \beta \rvert \cdot 2^{-1} \to \lvert x \rvert \in \mathbb{R}) \to \varphi_{24} \to (x \in \mathbb{R} \to \beta \in \mathbb{R} \to \varphi_{44}) \to \varphi_{29}$
\item $\varphi_{83} :\equiv (\lvert \beta \rvert \leq \lvert \beta - x \rvert + \lvert x \rvert \to \varphi_{34}) \to (x \in \mathbb{R} \to \varphi_{14}) \to (x \in \mathbb{R} \to \beta \in \mathbb{R} \to \varphi_{26}) \to x \in \mathbb{R} \to \varphi_{23}$
\item $\varphi_{84} :\equiv (x \in \mathbb{R} \to \beta \in \mathbb{R} \to \lvert x - \beta \rvert < \lvert \beta \rvert \cdot 2^{-1} \to \lvert x \rvert \in \mathbb{R}) \to \varphi_{80}$
\item $\varphi_{85} :\equiv (\lvert \beta \rvert \cdot 2^{-1} \in \mathbb{R} \to \lvert x \rvert \in \mathbb{R} \to \lvert \beta \rvert \cdot 2^{-1} \in \mathbb{R} \to \varphi_{40}) \to \varphi_{24} \to \varphi_{82}$
\end{abbrevlist}
\end{fullproof}

% 自動生成 (prover/build.js): RevTriangle  ノード数 658
\begin{fullproof}{[E27]}{$\forall x\,\forall y\,(x \in \mathbb{R} \to y \in \mathbb{R} \to \lvert \lvert x \rvert - \lvert y \rvert \rvert \leq \lvert x - y \rvert)$}{658}\phantomsection\label{tree:RevTriangle}
\noindent$\Pi_{1}$:
\begin{ptree}
\AXC{$z = w \to z = \lvert x - y \rvert \to w = \lvert x - y \rvert$\lbl{A6}}
\UIC{$\forall w\,(z = w \to z = \lvert x - y \rvert \to w = \lvert x - y \rvert)$}
\UIC{$\forall z\,\forall w\,(z = w \to z = \lvert x - y \rvert \to w = \lvert x - y \rvert)$}
\AXC{$\forall z\,\forall w\,(z = w \to z = \lvert x - y \rvert \to w = \lvert x - y \rvert) \to \varphi_{75}$\lbl{A2}}
\BIC{$\varphi_{75}$}
\AXC{$\varphi_{75} \to \varphi_{146}$\lbl{A2}}
\BIC{$\varphi_{146}$}
\end{ptree}
\noindent$\Pi_{2}$:
\begin{ptree}
\AXC{$z = w \to z < \lvert x - y \rvert \to w < \lvert x - y \rvert$\lbl{A6}}
\UIC{$\forall w\,(z = w \to z < \lvert x - y \rvert \to w < \lvert x - y \rvert)$}
\UIC{$\forall z\,\forall w\,(z = w \to z < \lvert x - y \rvert \to w < \lvert x - y \rvert)$}
\AXC{$\forall z\,\forall w\,(z = w \to z < \lvert x - y \rvert \to w < \lvert x - y \rvert) \to \varphi_{98}$\lbl{A2}}
\BIC{$\varphi_{98}$}
\AXC{$\varphi_{98} \to \varphi_{163}$\lbl{A2}}
\BIC{$\varphi_{163}$}
\end{ptree}
\noindent$\Pi_{3}$:
\begin{ptree}
\AXC{$z = w \to \lvert x \rvert + -\lvert y \rvert < z \to \lvert x \rvert + -\lvert y \rvert < w$\lbl{A6}}
\UIC{$\forall w\,(z = w \to \lvert x \rvert + -\lvert y \rvert < z \to \lvert x \rvert + -\lvert y \rvert < w)$}
\UIC{$\varphi_{41}$}
\AXC{$\varphi_{41} \to \varphi_{188}$\lbl{A2}}
\BIC{$\varphi_{188}$}
\AXC{$\varphi_{188} \to \lvert x - y \rvert + \lvert y \rvert + -\lvert y \rvert = \lvert x - y \rvert \to \varphi_{137}$\lbl{A2}}
\BIC{$\lvert x - y \rvert + \lvert y \rvert + -\lvert y \rvert = \lvert x - y \rvert \to \varphi_{137}$}
\end{ptree}
\noindent$\Pi_{4}$:
\begin{ptree}
\AXC{$\varphi_{51}$\lbl{A6}}
\UIC{$\forall w\,(\varphi_{51})$}
\UIC{$\forall z\,\forall w\,(\varphi_{51})$}
\AXC{$\forall z\,\forall w\,(\varphi_{51}) \to \varphi_{122}$\lbl{A2}}
\BIC{$\varphi_{122}$}
\AXC{$\varphi_{122} \to x - y + y = x \to \varphi_{71}$\lbl{A2}}
\BIC{$x - y + y = x \to \varphi_{71}$}
\end{ptree}
\noindent$\Pi_{5}$:
\begin{ptree}
\AXC{$\forall a\,\forall b\,(a \in \mathbb{R} \to b \in \mathbb{R} \to \lvert a + b \rvert \leq \lvert a \rvert + \lvert b \rvert)$\lbl{命題1.2}}
\AXC{$\varphi_{191}$\lbl{A2}}
\BIC{$\varphi_{55}$}
\AXC{$\varphi_{196}$\lbl{A2}}
\BIC{$x - y \in \mathbb{R} \to y \in \mathbb{R} \to \lvert x - y + y \rvert \leq \lvert x - y \rvert + \lvert y \rvert$}
\AXC{$\varphi_{221}$\lbl{A1}}
\BIC{$(x \in \mathbb{R} \to y \in \mathbb{R} \to x - y \in \mathbb{R}) \to \varphi_{140}$}
\end{ptree}
\noindent$\Pi_{6}$:
\begin{ptree}
\AXC{$\forall x\,\forall y\,(x \in \mathbb{R} \to y \in \mathbb{R} \to x - y \in \mathbb{R})$\lbl{[F1]}}
\AXC{$\varphi_{123}$\lbl{A2}}
\BIC{$\forall y\,(x \in \mathbb{R} \to y \in \mathbb{R} \to x - y \in \mathbb{R})$}
\AXC{$\varphi_{109}$\lbl{A2}}
\BIC{$x \in \mathbb{R} \to y \in \mathbb{R} \to x - y \in \mathbb{R}$}
\AXC{$\stackrel{\Pi_{5}}{\vdots}$}\noLine
\UIC{$(x \in \mathbb{R} \to y \in \mathbb{R} \to x - y \in \mathbb{R}) \to \varphi_{140}$}
\BIC{$\varphi_{140}$}
\end{ptree}
\noindent$\Pi_{7}$:
\begin{ptree}
\AXC{$\forall x\,\forall y\,(x \in \mathbb{R} \to y \in \mathbb{R} \to x - y + y = x)$\lbl{[F12]}}
\AXC{$\varphi_{114}$\lbl{A2}}
\BIC{$\forall y\,(x \in \mathbb{R} \to y \in \mathbb{R} \to x - y + y = x)$}
\AXC{$\varphi_{99}$\lbl{A2}}
\BIC{$x \in \mathbb{R} \to y \in \mathbb{R} \to x - y + y = x$}
\AXC{$(x \in \mathbb{R} \to y \in \mathbb{R} \to x - y + y = x) \to \varphi_{235}$\lbl{A1}}
\BIC{$\varphi_{235}$}
\end{ptree}
\noindent$\Pi_{8}$:
\begin{ptree}
\AXC{$\stackrel{\Pi_{4}}{\vdots}$}\noLine
\UIC{$x - y + y = x \to \varphi_{71}$}
\AXC{$x \in \mathbb{R} \to y \in \mathbb{R} \to y \in \mathbb{R}$\lbl{A1}}
\AXC{$\stackrel{\Pi_{6}}{\vdots}$}\noLine
\UIC{$\varphi_{140}$}
\BIC{$x \in \mathbb{R} \to y \in \mathbb{R} \to \lvert x - y + y \rvert \leq \lvert x - y \rvert + \lvert y \rvert$}
\AXC{$\stackrel{\Pi_{7}}{\vdots}$}\noLine
\UIC{$\varphi_{235}$}
\BIC{$(x - y + y = x \to \varphi_{71}) \to \varphi_{194}$}
\BIC{$\varphi_{194}$}
\end{ptree}
\noindent$\Pi_{9}$:
\begin{ptree}
\AXC{$\varphi_{30}$\lbl{A6}}
\UIC{$\forall w\,(\varphi_{30})$}
\UIC{$\forall z\,\forall w\,(\varphi_{30})$}
\AXC{$\forall z\,\forall w\,(\varphi_{30}) \to \varphi_{108}$\lbl{A2}}
\BIC{$\varphi_{108}$}
\AXC{$\varphi_{108} \to x - y + y = x \to \varphi_{48}$\lbl{A2}}
\BIC{$x - y + y = x \to \varphi_{48}$}
\AXC{$\stackrel{\Pi_{8}}{\vdots}$}\noLine
\UIC{$\varphi_{194}$}
\BIC{$x \in \mathbb{R} \to y \in \mathbb{R} \to \lvert x \rvert \leq \lvert x - y \rvert + \lvert y \rvert$}
\end{ptree}
\noindent$\Pi_{10}$:
\begin{ptree}
\AXC{$\forall x \in \mathbb{R}\,(\lvert x \rvert \in \mathbb{R})$\lbl{[A5]}}
\AXC{$\forall x \in \mathbb{R}\,(\lvert x \rvert \in \mathbb{R}) \to y \in \mathbb{R} \to \lvert y \rvert \in \mathbb{R}$\lbl{A2}}
\BIC{$y \in \mathbb{R} \to \lvert y \rvert \in \mathbb{R}$}
\AXC{$(y \in \mathbb{R} \to \lvert y \rvert \in \mathbb{R}) \to \varphi_{92}$\lbl{A1}}
\BIC{$\varphi_{92}$}
\end{ptree}
\noindent$\Pi_{11}$:
\begin{ptree}
\AXC{$x \in \mathbb{R} \to y \in \mathbb{R} \to y \in \mathbb{R}$\lbl{A1}}
\AXC{$\stackrel{\Pi_{10}}{\vdots}$}\noLine
\UIC{$\varphi_{92}$}
\BIC{$x \in \mathbb{R} \to y \in \mathbb{R} \to \lvert y \rvert \in \mathbb{R}$}
\end{ptree}
\noindent$\Pi_{12}$:
\begin{ptree}
\AXC{$\forall x \in \mathbb{R}\,(-x \in \mathbb{R})$\lbl{(R1)}}
\AXC{$\forall x \in \mathbb{R}\,(-x \in \mathbb{R}) \to \lvert y \rvert \in \mathbb{R} \to -\lvert y \rvert \in \mathbb{R}$\lbl{A2}}
\BIC{$\lvert y \rvert \in \mathbb{R} \to -\lvert y \rvert \in \mathbb{R}$}
\AXC{$(\lvert y \rvert \in \mathbb{R} \to -\lvert y \rvert \in \mathbb{R}) \to \varphi_{115}$\lbl{A1}}
\BIC{$\varphi_{115}$}
\end{ptree}
\noindent$\Pi_{13}$:
\begin{ptree}
\AXC{$\stackrel{\Pi_{11}}{\vdots}$}\noLine
\UIC{$x \in \mathbb{R} \to y \in \mathbb{R} \to \lvert y \rvert \in \mathbb{R}$}
\AXC{$\stackrel{\Pi_{12}}{\vdots}$}\noLine
\UIC{$\varphi_{115}$}
\BIC{$x \in \mathbb{R} \to y \in \mathbb{R} \to -\lvert y \rvert \in \mathbb{R}$}
\end{ptree}
\noindent$\Pi_{14}$:
\begin{ptree}
\AXC{$\forall x \in \mathbb{R}\,(\lvert x \rvert \in \mathbb{R})$\lbl{[A5]}}
\AXC{$\forall x \in \mathbb{R}\,(\lvert x \rvert \in \mathbb{R}) \to x - y \in \mathbb{R} \to \lvert x - y \rvert \in \mathbb{R}$\lbl{A2}}
\BIC{$x - y \in \mathbb{R} \to \lvert x - y \rvert \in \mathbb{R}$}
\end{ptree}
\noindent$\Pi_{15}$:
\begin{ptree}
\AXC{$\stackrel{\Pi_{14}}{\vdots}$}\noLine
\UIC{$x - y \in \mathbb{R} \to \lvert x - y \rvert \in \mathbb{R}$}
\AXC{$(x - y \in \mathbb{R} \to \lvert x - y \rvert \in \mathbb{R}) \to \varphi_{116}$\lbl{A1}}
\BIC{$\varphi_{116}$}
\end{ptree}
\noindent$\Pi_{16}$:
\begin{ptree}
\AXC{$\forall x\,\forall y\,(x \in \mathbb{R} \to y \in \mathbb{R} \to x - y \in \mathbb{R})$\lbl{[F1]}}
\AXC{$\varphi_{123}$\lbl{A2}}
\BIC{$\forall y\,(x \in \mathbb{R} \to y \in \mathbb{R} \to x - y \in \mathbb{R})$}
\AXC{$\varphi_{109}$\lbl{A2}}
\BIC{$x \in \mathbb{R} \to y \in \mathbb{R} \to x - y \in \mathbb{R}$}
\AXC{$\stackrel{\Pi_{15}}{\vdots}$}\noLine
\UIC{$\varphi_{116}$}
\BIC{$x \in \mathbb{R} \to y \in \mathbb{R} \to \lvert x - y \rvert \in \mathbb{R}$}
\end{ptree}
\noindent$\Pi_{17}$:
\begin{ptree}
\AXC{$\forall x\,\forall y\,(x \in \mathbb{R} \to y \in \mathbb{R} \to x + y \in \mathbb{R})$\lbl{(R1)}}
\AXC{$\varphi_{132}$\lbl{A2}}
\BIC{$\forall y\,(v_{646} \in \mathbb{R} \to y \in \mathbb{R} \to v_{646} + y \in \mathbb{R})$}
\AXC{$\varphi_{149}$\lbl{A2}}
\BIC{$v_{646} \in \mathbb{R} \to v_{647} \in \mathbb{R} \to v_{646} + v_{647} \in \mathbb{R}$}
\UIC{$\forall v_{647}\,(v_{646} \in \mathbb{R} \to v_{647} \in \mathbb{R} \to v_{646} + v_{647} \in \mathbb{R})$}
\UIC{$\varphi_{31}$}
\AXC{$\varphi_{31} \to \varphi_{58}$\lbl{A2}}
\BIC{$\varphi_{58}$}
\end{ptree}
\noindent$\Pi_{18}$:
\begin{ptree}
\AXC{$\stackrel{\Pi_{16}}{\vdots}$}\noLine
\UIC{$x \in \mathbb{R} \to y \in \mathbb{R} \to \lvert x - y \rvert \in \mathbb{R}$}
\AXC{$\stackrel{\Pi_{17}}{\vdots}$}\noLine
\UIC{$\varphi_{58}$}
\AXC{$\varphi_{58} \to \lvert x - y \rvert \in \mathbb{R} \to \lvert y \rvert \in \mathbb{R} \to \lvert x - y \rvert + \lvert y \rvert \in \mathbb{R}$\lbl{A2}}
\BIC{$\lvert x - y \rvert \in \mathbb{R} \to \lvert y \rvert \in \mathbb{R} \to \lvert x - y \rvert + \lvert y \rvert \in \mathbb{R}$}
\AXC{$\varphi_{223}$\lbl{A1}}
\BIC{$(x \in \mathbb{R} \to y \in \mathbb{R} \to \lvert x - y \rvert \in \mathbb{R}) \to \varphi_{133}$}
\BIC{$\varphi_{133}$}
\end{ptree}
\noindent$\Pi_{19}$:
\begin{ptree}
\AXC{$\forall x \in \mathbb{R}\,(\lvert x \rvert \in \mathbb{R})$\lbl{[A5]}}
\AXC{$\forall x \in \mathbb{R}\,(\lvert x \rvert \in \mathbb{R}) \to x \in \mathbb{R} \to \lvert x \rvert \in \mathbb{R}$\lbl{A2}}
\BIC{$x \in \mathbb{R} \to \lvert x \rvert \in \mathbb{R}$}
\AXC{$(x \in \mathbb{R} \to \lvert x \rvert \in \mathbb{R}) \to \varphi_{93}$\lbl{A1}}
\BIC{$\varphi_{93}$}
\end{ptree}
\noindent$\Pi_{20}$:
\begin{ptree}
\AXC{$x \in \mathbb{R} \to y \in \mathbb{R} \to x \in \mathbb{R}$\lbl{A1}}
\AXC{$\stackrel{\Pi_{19}}{\vdots}$}\noLine
\UIC{$\varphi_{93}$}
\BIC{$x \in \mathbb{R} \to y \in \mathbb{R} \to \lvert x \rvert \in \mathbb{R}$}
\end{ptree}
\noindent$\Pi_{21}$:
\begin{ptree}
\AXC{$\stackrel{\Pi_{20}}{\vdots}$}\noLine
\UIC{$x \in \mathbb{R} \to y \in \mathbb{R} \to \lvert x \rvert \in \mathbb{R}$}
\AXC{$\varphi_{42}$\lbl{[F19]}}
\AXC{$\varphi_{42} \to \varphi_{79}$\lbl{A2}}
\BIC{$\varphi_{79}$}
\AXC{$\varphi_{79} \to \varphi_{193}$\lbl{A2}}
\BIC{$\varphi_{193}$}
\AXC{$\varphi_{234}$\lbl{A2}}
\BIC{$\lvert x \rvert \in \mathbb{R} \to \lvert x - y \rvert + \lvert y \rvert \in \mathbb{R} \to -\lvert y \rvert \in \mathbb{R} \to \varphi_{112}$}
\AXC{$\varphi_{244}$\lbl{A1}}
\BIC{$\varphi_{231}$}
\BIC{$(x \in \mathbb{R} \to y \in \mathbb{R} \to \lvert x - y \rvert + \lvert y \rvert \in \mathbb{R}) \to \varphi_{213}$}
\end{ptree}
\noindent$\Pi_{22}$:
\begin{ptree}
\AXC{$\stackrel{\Pi_{11}}{\vdots}$}\noLine
\UIC{$x \in \mathbb{R} \to y \in \mathbb{R} \to \lvert y \rvert \in \mathbb{R}$}
\AXC{$\stackrel{\Pi_{18}}{\vdots}$}\noLine
\UIC{$\varphi_{133}$}
\BIC{$x \in \mathbb{R} \to y \in \mathbb{R} \to \lvert x - y \rvert + \lvert y \rvert \in \mathbb{R}$}
\AXC{$\stackrel{\Pi_{21}}{\vdots}$}\noLine
\UIC{$(x \in \mathbb{R} \to y \in \mathbb{R} \to \lvert x - y \rvert + \lvert y \rvert \in \mathbb{R}) \to \varphi_{213}$}
\BIC{$\varphi_{213}$}
\end{ptree}
\noindent$\Pi_{23}$:
\begin{ptree}
\AXC{$\stackrel{\Pi_{9}}{\vdots}$}\noLine
\UIC{$x \in \mathbb{R} \to y \in \mathbb{R} \to \lvert x \rvert \leq \lvert x - y \rvert + \lvert y \rvert$}
\AXC{$\stackrel{\Pi_{13}}{\vdots}$}\noLine
\UIC{$x \in \mathbb{R} \to y \in \mathbb{R} \to -\lvert y \rvert \in \mathbb{R}$}
\AXC{$\stackrel{\Pi_{22}}{\vdots}$}\noLine
\UIC{$\varphi_{213}$}
\BIC{$(x \in \mathbb{R} \to y \in \mathbb{R} \to \lvert x \rvert \leq \lvert x - y \rvert + \lvert y \rvert) \to \varphi_{100}$}
\BIC{$\varphi_{100}$}
\end{ptree}
\noindent$\Pi_{24}$:
\begin{ptree}
\AXC{$y = z \to x = y \to x = z$\lbl{A6}}
\AXC{$(y = z \to x = y \to x = z) \to x = y \to y = z \to x = z$\lbl{A1}}
\BIC{$x = y \to y = z \to x = z$}
\UIC{$\forall z\,(x = y \to y = z \to x = z)$}
\UIC{$\forall y\,\forall z\,(x = y \to y = z \to x = z)$}
\UIC{$\forall x\,\forall y\,\forall z\,(x = y \to y = z \to x = z)$}
\AXC{$\varphi_{8}$\lbl{A2}}
\BIC{$\forall y\,\forall z\,(v_{654} = y \to y = z \to v_{654} = z)$}
\AXC{$\varphi_{17}$\lbl{A2}}
\BIC{$\forall z\,(v_{654} = v_{655} \to v_{655} = z \to v_{654} = z)$}
\AXC{$\varphi_{36}$\lbl{A2}}
\BIC{$v_{654} = v_{655} \to v_{655} = v_{656} \to v_{654} = v_{656}$}
\UIC{$\forall v_{656}\,(v_{654} = v_{655} \to v_{655} = v_{656} \to v_{654} = v_{656})$}
\UIC{$\varphi_{4}$}
\UIC{$\forall v_{654}\,(\varphi_{4})$}
\end{ptree}
\noindent$\Pi_{25}$:
\begin{ptree}
\AXC{$\forall x \in \mathbb{R}\,(x + 0 = x)$\lbl{(R4)}}
\AXC{$\varphi_{76}$\lbl{A2}}
\BIC{$\lvert x - y \rvert \in \mathbb{R} \to \lvert x - y \rvert + 0 = \lvert x - y \rvert$}
\AXC{$(\lvert x - y \rvert \in \mathbb{R} \to \lvert x - y \rvert + 0 = \lvert x - y \rvert) \to \varphi_{141}$\lbl{A1}}
\BIC{$\varphi_{141}$}
\end{ptree}
\noindent$\Pi_{26}$:
\begin{ptree}
\AXC{$y = z \to x = y \to x = z$\lbl{A6}}
\AXC{$(y = z \to x = y \to x = z) \to x = y \to y = z \to x = z$\lbl{A1}}
\BIC{$x = y \to y = z \to x = z$}
\UIC{$\forall z\,(x = y \to y = z \to x = z)$}
\UIC{$\forall y\,\forall z\,(x = y \to y = z \to x = z)$}
\UIC{$\forall x\,\forall y\,\forall z\,(x = y \to y = z \to x = z)$}
\AXC{$\varphi_{9}$\lbl{A2}}
\BIC{$\forall y\,\forall z\,(v_{651} = y \to y = z \to v_{651} = z)$}
\AXC{$\varphi_{18}$\lbl{A2}}
\BIC{$\forall z\,(v_{651} = v_{652} \to v_{652} = z \to v_{651} = z)$}
\AXC{$\varphi_{37}$\lbl{A2}}
\BIC{$v_{651} = v_{652} \to v_{652} = v_{653} \to v_{651} = v_{653}$}
\UIC{$\forall v_{653}\,(v_{651} = v_{652} \to v_{652} = v_{653} \to v_{651} = v_{653})$}
\UIC{$\varphi_{5}$}
\UIC{$\forall v_{651}\,(\varphi_{5})$}
\end{ptree}
\noindent$\Pi_{27}$:
\begin{ptree}
\AXC{$z = w \to \varphi_{159}$\lbl{A6}}
\UIC{$\forall w\,(z = w \to \varphi_{159})$}
\UIC{$\forall z\,\forall w\,(z = w \to \varphi_{159})$}
\AXC{$\forall z\,\forall w\,(z = w \to \varphi_{159}) \to \varphi_{203}$\lbl{A2}}
\BIC{$\varphi_{203}$}
\AXC{$\varphi_{203} \to \varphi_{201}$\lbl{A2}}
\BIC{$\varphi_{201}$}
\end{ptree}
\noindent$\Pi_{28}$:
\begin{ptree}
\AXC{$\stackrel{\Pi_{11}}{\vdots}$}\noLine
\UIC{$x \in \mathbb{R} \to y \in \mathbb{R} \to \lvert y \rvert \in \mathbb{R}$}
\AXC{$\forall x \in \mathbb{R}\,(x + -x = 0)$\lbl{(R5)}}
\AXC{$\varphi_{60}$\lbl{A2}}
\BIC{$\lvert y \rvert \in \mathbb{R} \to \lvert y \rvert + -\lvert y \rvert = 0$}
\AXC{$(\lvert y \rvert \in \mathbb{R} \to \lvert y \rvert + -\lvert y \rvert = 0) \to \varphi_{120}$\lbl{A1}}
\BIC{$\varphi_{120}$}
\BIC{$x \in \mathbb{R} \to y \in \mathbb{R} \to \lvert y \rvert + -\lvert y \rvert = 0$}
\end{ptree}
\noindent$\Pi_{29}$:
\begin{ptree}
\AXC{$\stackrel{\Pi_{28}}{\vdots}$}\noLine
\UIC{$x \in \mathbb{R} \to y \in \mathbb{R} \to \lvert y \rvert + -\lvert y \rvert = 0$}
\AXC{$(x \in \mathbb{R} \to y \in \mathbb{R} \to \lvert y \rvert + -\lvert y \rvert = 0) \to \varphi_{59} \to \varphi_{201} \to \varphi_{80}$\lbl{A1}}
\BIC{$\varphi_{59} \to \varphi_{201} \to \varphi_{80}$}
\end{ptree}
\noindent$\Pi_{30}$:
\begin{ptree}
\AXC{$\stackrel{\Pi_{27}}{\vdots}$}\noLine
\UIC{$\varphi_{201}$}
\AXC{$x = x$\lbl{A5}}
\UIC{$\forall x\,(x = x)$}
\AXC{$\forall x\,(x = x) \to \varphi_{59}$\lbl{A2}}
\BIC{$\varphi_{59}$}
\AXC{$\stackrel{\Pi_{29}}{\vdots}$}\noLine
\UIC{$\varphi_{59} \to \varphi_{201} \to \varphi_{80}$}
\BIC{$\varphi_{201} \to \varphi_{80}$}
\BIC{$\varphi_{80}$}
\end{ptree}
\noindent$\Pi_{31}$:
\begin{ptree}
\AXC{$\varphi_{52}$\lbl{(R2)}}
\AXC{$\varphi_{52} \to \varphi_{73}$\lbl{A2}}
\BIC{$\varphi_{73}$}
\AXC{$\varphi_{73} \to \varphi_{96}$\lbl{A2}}
\BIC{$\varphi_{96}$}
\AXC{$\varphi_{96} \to v_{648} \in \mathbb{R} \to v_{649} \in \mathbb{R} \to \varphi_{13}$\lbl{A2}}
\BIC{$v_{648} \in \mathbb{R} \to v_{649} \in \mathbb{R} \to \varphi_{13}$}
\UIC{$\forall v_{650}\,(v_{648} \in \mathbb{R} \to v_{649} \in \mathbb{R} \to \varphi_{13})$}
\UIC{$\forall v_{649}\,\forall v_{650}\,(v_{648} \in \mathbb{R} \to v_{649} \in \mathbb{R} \to \varphi_{13})$}
\UIC{$\forall v_{648}\,\forall v_{649}\,\forall v_{650}\,(v_{648} \in \mathbb{R} \to v_{649} \in \mathbb{R} \to \varphi_{13})$}
\AXC{$\varphi_{226}$\lbl{A2}}
\BIC{$\varphi_{172}$}
\AXC{$\varphi_{172} \to \varphi_{173}$\lbl{A2}}
\BIC{$\varphi_{173}$}
\end{ptree}
\noindent$\Pi_{32}$:
\begin{ptree}
\AXC{$\stackrel{\Pi_{31}}{\vdots}$}\noLine
\UIC{$\varphi_{173}$}
\AXC{$\varphi_{173} \to \lvert x - y \rvert \in \mathbb{R} \to \lvert y \rvert \in \mathbb{R} \to -\lvert y \rvert \in \mathbb{R} \to \varphi_{65}$\lbl{A2}}
\BIC{$\lvert x - y \rvert \in \mathbb{R} \to \lvert y \rvert \in \mathbb{R} \to -\lvert y \rvert \in \mathbb{R} \to \varphi_{65}$}
\AXC{$(\lvert x - y \rvert \in \mathbb{R} \to \lvert y \rvert \in \mathbb{R} \to -\lvert y \rvert \in \mathbb{R} \to \varphi_{65}) \to \varphi_{229}$\lbl{A1}}
\BIC{$\varphi_{229}$}
\end{ptree}
\noindent$\Pi_{33}$:
\begin{ptree}
\AXC{$\stackrel{\Pi_{13}}{\vdots}$}\noLine
\UIC{$x \in \mathbb{R} \to y \in \mathbb{R} \to -\lvert y \rvert \in \mathbb{R}$}
\AXC{$\stackrel{\Pi_{11}}{\vdots}$}\noLine
\UIC{$x \in \mathbb{R} \to y \in \mathbb{R} \to \lvert y \rvert \in \mathbb{R}$}
\AXC{$\stackrel{\Pi_{16}}{\vdots}$}\noLine
\UIC{$x \in \mathbb{R} \to y \in \mathbb{R} \to \lvert x - y \rvert \in \mathbb{R}$}
\AXC{$\stackrel{\Pi_{32}}{\vdots}$}\noLine
\UIC{$\varphi_{229}$}
\BIC{$(x \in \mathbb{R} \to y \in \mathbb{R} \to \lvert y \rvert \in \mathbb{R}) \to \varphi_{198}$}
\BIC{$\varphi_{198}$}
\BIC{$x \in \mathbb{R} \to y \in \mathbb{R} \to \varphi_{65}$}
\end{ptree}
\noindent$\Pi_{34}$:
\begin{ptree}
\AXC{$\stackrel{\Pi_{30}}{\vdots}$}\noLine
\UIC{$\varphi_{80}$}
\AXC{$\stackrel{\Pi_{33}}{\vdots}$}\noLine
\UIC{$x \in \mathbb{R} \to y \in \mathbb{R} \to \varphi_{65}$}
\AXC{$(x \in \mathbb{R} \to y \in \mathbb{R} \to \varphi_{65}) \to \varphi_{80} \to (\varphi_{65} \to \varphi_{158}) \to \varphi_{81}$\lbl{A1}}
\BIC{$\varphi_{80} \to (\varphi_{65} \to \varphi_{158}) \to \varphi_{81}$}
\BIC{$(\varphi_{65} \to \varphi_{158}) \to \varphi_{81}$}
\end{ptree}
\noindent$\Pi_{35}$:
\begin{ptree}
\AXC{$\stackrel{\Pi_{26}}{\vdots}$}\noLine
\UIC{$\forall v_{651}\,(\varphi_{5})$}
\AXC{$\forall v_{651}\,(\varphi_{5}) \to \varphi_{154}$\lbl{A2}}
\BIC{$\varphi_{154}$}
\AXC{$\varphi_{154} \to \varphi_{209}$\lbl{A2}}
\BIC{$\varphi_{209}$}
\AXC{$\varphi_{209} \to \varphi_{65} \to \varphi_{158}$\lbl{A2}}
\BIC{$\varphi_{65} \to \varphi_{158}$}
\AXC{$\stackrel{\Pi_{34}}{\vdots}$}\noLine
\UIC{$(\varphi_{65} \to \varphi_{158}) \to \varphi_{81}$}
\BIC{$\varphi_{81}$}
\end{ptree}
\noindent$\Pi_{36}$:
\begin{ptree}
\AXC{$\stackrel{\Pi_{16}}{\vdots}$}\noLine
\UIC{$x \in \mathbb{R} \to y \in \mathbb{R} \to \lvert x - y \rvert \in \mathbb{R}$}
\AXC{$\stackrel{\Pi_{25}}{\vdots}$}\noLine
\UIC{$\varphi_{141}$}
\BIC{$x \in \mathbb{R} \to y \in \mathbb{R} \to \lvert x - y \rvert + 0 = \lvert x - y \rvert$}
\AXC{$\stackrel{\Pi_{35}}{\vdots}$}\noLine
\UIC{$\varphi_{81}$}
\AXC{$\varphi_{81} \to (x \in \mathbb{R} \to y \in \mathbb{R} \to \lvert x - y \rvert + 0 = \lvert x - y \rvert) \to \varphi_{218}$\lbl{A1}}
\BIC{$(x \in \mathbb{R} \to y \in \mathbb{R} \to \lvert x - y \rvert + 0 = \lvert x - y \rvert) \to \varphi_{218}$}
\BIC{$\varphi_{218}$}
\end{ptree}
\noindent$\Pi_{37}$:
\begin{ptree}
\AXC{$\stackrel{\Pi_{24}}{\vdots}$}\noLine
\UIC{$\forall v_{654}\,(\varphi_{4})$}
\AXC{$\forall v_{654}\,(\varphi_{4}) \to \varphi_{153}$\lbl{A2}}
\BIC{$\varphi_{153}$}
\AXC{$\varphi_{153} \to \varphi_{180}$\lbl{A2}}
\BIC{$\varphi_{180}$}
\AXC{$\varphi_{180} \to \lvert x - y \rvert + \lvert y \rvert + -\lvert y \rvert = \lvert x - y \rvert + 0 \to \varphi_{88}$\lbl{A2}}
\BIC{$\lvert x - y \rvert + \lvert y \rvert + -\lvert y \rvert = \lvert x - y \rvert + 0 \to \varphi_{88}$}
\AXC{$\stackrel{\Pi_{36}}{\vdots}$}\noLine
\UIC{$\varphi_{218}$}
\BIC{$x \in \mathbb{R} \to y \in \mathbb{R} \to \lvert x - y \rvert + \lvert y \rvert + -\lvert y \rvert = \lvert x - y \rvert$}
\end{ptree}
\noindent$\Pi_{38}$:
\begin{ptree}
\AXC{$\stackrel{\Pi_{3}}{\vdots}$}\noLine
\UIC{$\lvert x - y \rvert + \lvert y \rvert + -\lvert y \rvert = \lvert x - y \rvert \to \varphi_{137}$}
\AXC{$\stackrel{\Pi_{23}}{\vdots}$}\noLine
\UIC{$\varphi_{100}$}
\AXC{$\stackrel{\Pi_{37}}{\vdots}$}\noLine
\UIC{$x \in \mathbb{R} \to y \in \mathbb{R} \to \lvert x - y \rvert + \lvert y \rvert + -\lvert y \rvert = \lvert x - y \rvert$}
\AXC{$\varphi_{246}$\lbl{A1}}
\BIC{$\varphi_{100} \to (\lvert x - y \rvert + \lvert y \rvert + -\lvert y \rvert = \lvert x - y \rvert \to \varphi_{137}) \to \varphi_{216}$}
\BIC{$(\lvert x - y \rvert + \lvert y \rvert + -\lvert y \rvert = \lvert x - y \rvert \to \varphi_{137}) \to \varphi_{216}$}
\BIC{$\varphi_{216}$}
\end{ptree}
\noindent$\Pi_{39}$:
\begin{ptree}
\AXC{$z = w \to \lvert x \rvert + -\lvert y \rvert = z \to \lvert x \rvert + -\lvert y \rvert = w$\lbl{A6}}
\UIC{$\forall w\,(z = w \to \lvert x \rvert + -\lvert y \rvert = z \to \lvert x \rvert + -\lvert y \rvert = w)$}
\UIC{$\varphi_{22}$}
\AXC{$\varphi_{22} \to \varphi_{182}$\lbl{A2}}
\BIC{$\varphi_{182}$}
\AXC{$\varphi_{182} \to \lvert x - y \rvert + \lvert y \rvert + -\lvert y \rvert = \lvert x - y \rvert \to \varphi_{125}$\lbl{A2}}
\BIC{$\lvert x - y \rvert + \lvert y \rvert + -\lvert y \rvert = \lvert x - y \rvert \to \varphi_{125}$}
\AXC{$\stackrel{\Pi_{38}}{\vdots}$}\noLine
\UIC{$\varphi_{216}$}
\BIC{$x \in \mathbb{R} \to y \in \mathbb{R} \to \lvert x \rvert + -\lvert y \rvert \leq \lvert x - y \rvert$}
\end{ptree}
\noindent$\Pi_{40}$:
\begin{ptree}
\AXC{$x = x$\lbl{A5}}
\AXC{$x = y \to x = x \to y = x$\lbl{A6}}
\AXC{$(x = y \to x = x \to y = x) \to x = x \to x = y \to y = x$\lbl{A1}}
\BIC{$x = x \to x = y \to y = x$}
\BIC{$x = y \to y = x$}
\UIC{$\forall y\,(x = y \to y = x)$}
\UIC{$\forall x\,\forall y\,(x = y \to y = x)$}
\end{ptree}
\noindent$\Pi_{41}$:
\begin{ptree}
\AXC{$\stackrel{\Pi_{40}}{\vdots}$}\noLine
\UIC{$\forall x\,\forall y\,(x = y \to y = x)$}
\AXC{$\forall x\,\forall y\,(x = y \to y = x) \to \forall y\,(v_{657} = y \to y = v_{657})$\lbl{A2}}
\BIC{$\forall y\,(v_{657} = y \to y = v_{657})$}
\AXC{$\varphi_{1}$\lbl{A2}}
\BIC{$v_{657} = v_{658} \to v_{658} = v_{657}$}
\UIC{$\forall v_{658}\,(v_{657} = v_{658} \to v_{658} = v_{657})$}
\UIC{$\forall v_{657}\,\forall v_{658}\,(v_{657} = v_{658} \to v_{658} = v_{657})$}
\AXC{$\varphi_{106}$\lbl{A2}}
\BIC{$\forall v_{658}\,(\lvert x \rvert - \lvert y \rvert = v_{658} \to v_{658} = \lvert x \rvert - \lvert y \rvert)$}
\AXC{$\varphi_{186}$\lbl{A2}}
\BIC{$\varphi_{82}$}
\end{ptree}
\noindent$\Pi_{42}$:
\begin{ptree}
\AXC{$\stackrel{\Pi_{41}}{\vdots}$}\noLine
\UIC{$\varphi_{82}$}
\AXC{$\forall x\,\forall y\,(x - y = x + -y)$\lbl{Def$-$}}
\AXC{$\varphi_{15}$\lbl{A2}}
\BIC{$\forall y\,(\lvert x \rvert - y = \lvert x \rvert + -y)$}
\AXC{$\varphi_{56}$\lbl{A2}}
\BIC{$\lvert x \rvert - \lvert y \rvert = \lvert x \rvert + -\lvert y \rvert$}
\AXC{$\varphi_{206}$\lbl{A1}}
\BIC{$\varphi_{82} \to \lvert x \rvert + -\lvert y \rvert = \lvert x \rvert - \lvert y \rvert$}
\BIC{$\lvert x \rvert + -\lvert y \rvert = \lvert x \rvert - \lvert y \rvert$}
\AXC{$\lvert x \rvert + -\lvert y \rvert = \lvert x \rvert - \lvert y \rvert \to \varphi_{238}$\lbl{A1}}
\BIC{$\varphi_{238}$}
\end{ptree}
\noindent$\Pi_{43}$:
\begin{ptree}
\AXC{$z = w \to -\lvert x - y \rvert = z \to -\lvert x - y \rvert = w$\lbl{A6}}
\UIC{$\forall w\,(z = w \to -\lvert x - y \rvert = z \to -\lvert x - y \rvert = w)$}
\UIC{$\forall z\,\forall w\,(z = w \to -\lvert x - y \rvert = z \to -\lvert x - y \rvert = w)$}
\AXC{$\forall z\,\forall w\,(z = w \to -\lvert x - y \rvert = z \to -\lvert x - y \rvert = w) \to \varphi_{130}$\lbl{A2}}
\BIC{$\varphi_{130}$}
\end{ptree}
\noindent$\Pi_{44}$:
\begin{ptree}
\AXC{$z = w \to -\lvert x - y \rvert < z \to -\lvert x - y \rvert < w$\lbl{A6}}
\UIC{$\forall w\,(z = w \to -\lvert x - y \rvert < z \to -\lvert x - y \rvert < w)$}
\UIC{$\forall z\,\forall w\,(z = w \to -\lvert x - y \rvert < z \to -\lvert x - y \rvert < w)$}
\AXC{$\forall z\,\forall w\,(z = w \to -\lvert x - y \rvert < z \to -\lvert x - y \rvert < w) \to \varphi_{144}$\lbl{A2}}
\BIC{$\varphi_{144}$}
\end{ptree}
\noindent$\Pi_{45}$:
\begin{ptree}
\AXC{$z = w \to z = \lvert x - y \rvert \to w = \lvert x - y \rvert$\lbl{A6}}
\UIC{$\forall w\,(z = w \to z = \lvert x - y \rvert \to w = \lvert x - y \rvert)$}
\UIC{$\forall z\,\forall w\,(z = w \to z = \lvert x - y \rvert \to w = \lvert x - y \rvert)$}
\AXC{$\forall z\,\forall w\,(z = w \to z = \lvert x - y \rvert \to w = \lvert x - y \rvert) \to \varphi_{45}$\lbl{A2}}
\BIC{$\varphi_{45}$}
\AXC{$\varphi_{45} \to \varphi_{147}$\lbl{A2}}
\BIC{$\varphi_{147}$}
\end{ptree}
\noindent$\Pi_{46}$:
\begin{ptree}
\AXC{$z = w \to z < \lvert x - y \rvert \to w < \lvert x - y \rvert$\lbl{A6}}
\UIC{$\forall w\,(z = w \to z < \lvert x - y \rvert \to w < \lvert x - y \rvert)$}
\UIC{$\forall z\,\forall w\,(z = w \to z < \lvert x - y \rvert \to w < \lvert x - y \rvert)$}
\AXC{$\forall z\,\forall w\,(z = w \to z < \lvert x - y \rvert \to w < \lvert x - y \rvert) \to \varphi_{67}$\lbl{A2}}
\BIC{$\varphi_{67}$}
\AXC{$\varphi_{67} \to \varphi_{164}$\lbl{A2}}
\BIC{$\varphi_{164}$}
\end{ptree}
\noindent$\Pi_{47}$:
\begin{ptree}
\AXC{$z = w \to z = \lvert x - y \rvert \to w = \lvert x - y \rvert$\lbl{A6}}
\UIC{$\forall w\,(z = w \to z = \lvert x - y \rvert \to w = \lvert x - y \rvert)$}
\UIC{$\forall z\,\forall w\,(z = w \to z = \lvert x - y \rvert \to w = \lvert x - y \rvert)$}
\AXC{$\forall z\,\forall w\,(z = w \to z = \lvert x - y \rvert \to w = \lvert x - y \rvert) \to \varphi_{77}$\lbl{A2}}
\BIC{$\varphi_{77}$}
\AXC{$\varphi_{77} \to \varphi_{148}$\lbl{A2}}
\BIC{$\varphi_{148}$}
\end{ptree}
\noindent$\Pi_{48}$:
\begin{ptree}
\AXC{$z = w \to z < \lvert x - y \rvert \to w < \lvert x - y \rvert$\lbl{A6}}
\UIC{$\forall w\,(z = w \to z < \lvert x - y \rvert \to w < \lvert x - y \rvert)$}
\UIC{$\forall z\,\forall w\,(z = w \to z < \lvert x - y \rvert \to w < \lvert x - y \rvert)$}
\AXC{$\forall z\,\forall w\,(z = w \to z < \lvert x - y \rvert \to w < \lvert x - y \rvert) \to \varphi_{101}$\lbl{A2}}
\BIC{$\varphi_{101}$}
\AXC{$\varphi_{101} \to \varphi_{165}$\lbl{A2}}
\BIC{$\varphi_{165}$}
\end{ptree}
\noindent$\Pi_{49}$:
\begin{ptree}
\AXC{$z = w \to \lvert y \rvert + -\lvert x \rvert < z \to \lvert y \rvert + -\lvert x \rvert < w$\lbl{A6}}
\UIC{$\forall w\,(z = w \to \lvert y \rvert + -\lvert x \rvert < z \to \lvert y \rvert + -\lvert x \rvert < w)$}
\UIC{$\varphi_{43}$}
\AXC{$\varphi_{43} \to \varphi_{189}$\lbl{A2}}
\BIC{$\varphi_{189}$}
\AXC{$\varphi_{189} \to \lvert x - y \rvert + \lvert x \rvert + -\lvert x \rvert = \lvert x - y \rvert \to \varphi_{138}$\lbl{A2}}
\BIC{$\lvert x - y \rvert + \lvert x \rvert + -\lvert x \rvert = \lvert x - y \rvert \to \varphi_{138}$}
\end{ptree}
\noindent$\Pi_{50}$:
\begin{ptree}
\AXC{$z = w \to \lvert y \rvert = z + \lvert x \rvert \to \lvert y \rvert = w + \lvert x \rvert$\lbl{A6}}
\UIC{$\forall w\,(z = w \to \lvert y \rvert = z + \lvert x \rvert \to \lvert y \rvert = w + \lvert x \rvert)$}
\UIC{$\forall z\,\forall w\,(z = w \to \lvert y \rvert = z + \lvert x \rvert \to \lvert y \rvert = w + \lvert x \rvert)$}
\AXC{$\varphi_{174}$\lbl{A2}}
\BIC{$\varphi_{46}$}
\AXC{$\varphi_{46} \to \varphi_{103}$\lbl{A2}}
\BIC{$\varphi_{103}$}
\end{ptree}
\noindent$\Pi_{51}$:
\begin{ptree}
\AXC{$z = w \to \lvert y \rvert < z + \lvert x \rvert \to \lvert y \rvert < w + \lvert x \rvert$\lbl{A6}}
\UIC{$\forall w\,(z = w \to \lvert y \rvert < z + \lvert x \rvert \to \lvert y \rvert < w + \lvert x \rvert)$}
\UIC{$\forall z\,\forall w\,(z = w \to \lvert y \rvert < z + \lvert x \rvert \to \lvert y \rvert < w + \lvert x \rvert)$}
\AXC{$\varphi_{190}$\lbl{A2}}
\BIC{$\varphi_{68}$}
\AXC{$\varphi_{68} \to \varphi_{118}$\lbl{A2}}
\BIC{$\varphi_{118}$}
\end{ptree}
\noindent$\Pi_{52}$:
\begin{ptree}
\AXC{$\varphi_{53}$\lbl{A6}}
\UIC{$\forall w\,(\varphi_{53})$}
\UIC{$\forall z\,\forall w\,(\varphi_{53})$}
\AXC{$\forall z\,\forall w\,(\varphi_{53}) \to \varphi_{124}$\lbl{A2}}
\BIC{$\varphi_{124}$}
\AXC{$\varphi_{124} \to y - x + x = y \to \varphi_{72}$\lbl{A2}}
\BIC{$y - x + x = y \to \varphi_{72}$}
\end{ptree}
\noindent$\Pi_{53}$:
\begin{ptree}
\AXC{$\forall x\,\forall y\,(x \in \mathbb{R} \to y \in \mathbb{R} \to x - y \in \mathbb{R})$\lbl{[F1]}}
\AXC{$\varphi_{134}$\lbl{A2}}
\BIC{$\forall y\,(v_{295} \in \mathbb{R} \to y \in \mathbb{R} \to v_{295} - y \in \mathbb{R})$}
\AXC{$\varphi_{150}$\lbl{A2}}
\BIC{$v_{295} \in \mathbb{R} \to v_{296} \in \mathbb{R} \to v_{295} - v_{296} \in \mathbb{R}$}
\UIC{$\forall v_{296}\,(v_{295} \in \mathbb{R} \to v_{296} \in \mathbb{R} \to v_{295} - v_{296} \in \mathbb{R})$}
\UIC{$\varphi_{33}$}
\AXC{$\varphi_{33} \to \forall v_{296}\,(y \in \mathbb{R} \to v_{296} \in \mathbb{R} \to y - v_{296} \in \mathbb{R})$\lbl{A2}}
\BIC{$\forall v_{296}\,(y \in \mathbb{R} \to v_{296} \in \mathbb{R} \to y - v_{296} \in \mathbb{R})$}
\end{ptree}
\noindent$\Pi_{54}$:
\begin{ptree}
\AXC{$x \in \mathbb{R} \to y \in \mathbb{R} \to y \in \mathbb{R}$\lbl{A1}}
\AXC{$\stackrel{\Pi_{53}}{\vdots}$}\noLine
\UIC{$\forall v_{296}\,(y \in \mathbb{R} \to v_{296} \in \mathbb{R} \to y - v_{296} \in \mathbb{R})$}
\AXC{$\varphi_{129}$\lbl{A2}}
\BIC{$y \in \mathbb{R} \to x \in \mathbb{R} \to y - x \in \mathbb{R}$}
\AXC{$\varphi_{208}$\lbl{A1}}
\BIC{$(x \in \mathbb{R} \to y \in \mathbb{R} \to y \in \mathbb{R}) \to \varphi_{94}$}
\BIC{$\varphi_{94}$}
\end{ptree}
\noindent$\Pi_{55}$:
\begin{ptree}
\AXC{$\forall a\,\forall b\,(a \in \mathbb{R} \to b \in \mathbb{R} \to \lvert a + b \rvert \leq \lvert a \rvert + \lvert b \rvert)$\lbl{命題1.2}}
\AXC{$\varphi_{192}$\lbl{A2}}
\BIC{$\varphi_{57}$}
\AXC{$\varphi_{197}$\lbl{A2}}
\BIC{$y - x \in \mathbb{R} \to x \in \mathbb{R} \to \lvert y - x + x \rvert \leq \lvert y - x \rvert + \lvert x \rvert$}
\AXC{$\varphi_{222}$\lbl{A1}}
\BIC{$(x \in \mathbb{R} \to y \in \mathbb{R} \to y - x \in \mathbb{R}) \to \varphi_{142}$}
\end{ptree}
\noindent$\Pi_{56}$:
\begin{ptree}
\AXC{$x \in \mathbb{R} \to y \in \mathbb{R} \to x \in \mathbb{R}$\lbl{A1}}
\AXC{$x \in \mathbb{R} \to y \in \mathbb{R} \to x \in \mathbb{R}$\lbl{A1}}
\AXC{$\stackrel{\Pi_{54}}{\vdots}$}\noLine
\UIC{$\varphi_{94}$}
\BIC{$x \in \mathbb{R} \to y \in \mathbb{R} \to y - x \in \mathbb{R}$}
\AXC{$\stackrel{\Pi_{55}}{\vdots}$}\noLine
\UIC{$(x \in \mathbb{R} \to y \in \mathbb{R} \to y - x \in \mathbb{R}) \to \varphi_{142}$}
\BIC{$\varphi_{142}$}
\BIC{$x \in \mathbb{R} \to y \in \mathbb{R} \to \lvert y - x + x \rvert \leq \lvert y - x \rvert + \lvert x \rvert$}
\end{ptree}
\noindent$\Pi_{57}$:
\begin{ptree}
\AXC{$\forall x\,\forall y\,(x \in \mathbb{R} \to y \in \mathbb{R} \to x - y + y = x)$\lbl{[F12]}}
\AXC{$\varphi_{131}$\lbl{A2}}
\BIC{$\forall y\,(v_{642} \in \mathbb{R} \to y \in \mathbb{R} \to v_{642} - y + y = v_{642})$}
\AXC{$\varphi_{166}$\lbl{A2}}
\BIC{$v_{642} \in \mathbb{R} \to v_{643} \in \mathbb{R} \to v_{642} - v_{643} + v_{643} = v_{642}$}
\UIC{$\varphi_{27}$}
\UIC{$\forall v_{642}\,(\varphi_{27})$}
\AXC{$\varphi_{178}$\lbl{A2}}
\BIC{$\forall v_{643}\,(y \in \mathbb{R} \to v_{643} \in \mathbb{R} \to y - v_{643} + v_{643} = y)$}
\AXC{$\varphi_{128}$\lbl{A2}}
\BIC{$y \in \mathbb{R} \to x \in \mathbb{R} \to y - x + x = y$}
\AXC{$\varphi_{205}$\lbl{A1}}
\BIC{$(x \in \mathbb{R} \to y \in \mathbb{R} \to y \in \mathbb{R}) \to \varphi_{78}$}
\end{ptree}
\noindent$\Pi_{58}$:
\begin{ptree}
\AXC{$x \in \mathbb{R} \to y \in \mathbb{R} \to x \in \mathbb{R}$\lbl{A1}}
\AXC{$x \in \mathbb{R} \to y \in \mathbb{R} \to y \in \mathbb{R}$\lbl{A1}}
\AXC{$\stackrel{\Pi_{57}}{\vdots}$}\noLine
\UIC{$(x \in \mathbb{R} \to y \in \mathbb{R} \to y \in \mathbb{R}) \to \varphi_{78}$}
\BIC{$\varphi_{78}$}
\BIC{$x \in \mathbb{R} \to y \in \mathbb{R} \to y - x + x = y$}
\end{ptree}
\noindent$\Pi_{59}$:
\begin{ptree}
\AXC{$\stackrel{\Pi_{58}}{\vdots}$}\noLine
\UIC{$x \in \mathbb{R} \to y \in \mathbb{R} \to y - x + x = y$}
\AXC{$(x \in \mathbb{R} \to y \in \mathbb{R} \to y - x + x = y) \to \varphi_{236}$\lbl{A1}}
\BIC{$\varphi_{236}$}
\end{ptree}
\noindent$\Pi_{60}$:
\begin{ptree}
\AXC{$\stackrel{\Pi_{52}}{\vdots}$}\noLine
\UIC{$y - x + x = y \to \varphi_{72}$}
\AXC{$\stackrel{\Pi_{56}}{\vdots}$}\noLine
\UIC{$x \in \mathbb{R} \to y \in \mathbb{R} \to \lvert y - x + x \rvert \leq \lvert y - x \rvert + \lvert x \rvert$}
\AXC{$\stackrel{\Pi_{59}}{\vdots}$}\noLine
\UIC{$\varphi_{236}$}
\BIC{$(y - x + x = y \to \varphi_{72}) \to \varphi_{195}$}
\BIC{$\varphi_{195}$}
\end{ptree}
\noindent$\Pi_{61}$:
\begin{ptree}
\AXC{$\forall x\,\forall y\,(x \in \mathbb{R} \to y \in \mathbb{R} \to \lvert x - y \rvert = \lvert y - x \rvert)$\lbl{[A13]}}
\AXC{$\forall x\,\forall y\,(x \in \mathbb{R} \to y \in \mathbb{R} \to \lvert x - y \rvert = \lvert y - x \rvert) \to \varphi_{24}$\lbl{A2}}
\BIC{$\varphi_{24}$}
\AXC{$\varphi_{24} \to \varphi_{44}$\lbl{A2}}
\BIC{$\varphi_{44}$}
\UIC{$\forall v_{645}\,(\varphi_{44})$}
\UIC{$\forall v_{644}\,\forall v_{645}\,(\varphi_{44})$}
\end{ptree}
\noindent$\Pi_{62}$:
\begin{ptree}
\AXC{$\stackrel{\Pi_{61}}{\vdots}$}\noLine
\UIC{$\forall v_{644}\,\forall v_{645}\,(\varphi_{44})$}
\AXC{$\forall v_{644}\,\forall v_{645}\,(\varphi_{44}) \to \varphi_{29}$\lbl{A2}}
\BIC{$\varphi_{29}$}
\AXC{$\varphi_{29} \to y \in \mathbb{R} \to x \in \mathbb{R} \to \lvert y - x \rvert = \lvert x - y \rvert$\lbl{A2}}
\BIC{$y \in \mathbb{R} \to x \in \mathbb{R} \to \lvert y - x \rvert = \lvert x - y \rvert$}
\AXC{$\varphi_{211}$\lbl{A1}}
\BIC{$(x \in \mathbb{R} \to y \in \mathbb{R} \to y \in \mathbb{R}) \to \varphi_{111}$}
\end{ptree}
\noindent$\Pi_{63}$:
\begin{ptree}
\AXC{$x \in \mathbb{R} \to y \in \mathbb{R} \to x \in \mathbb{R}$\lbl{A1}}
\AXC{$x \in \mathbb{R} \to y \in \mathbb{R} \to y \in \mathbb{R}$\lbl{A1}}
\AXC{$\stackrel{\Pi_{62}}{\vdots}$}\noLine
\UIC{$(x \in \mathbb{R} \to y \in \mathbb{R} \to y \in \mathbb{R}) \to \varphi_{111}$}
\BIC{$\varphi_{111}$}
\BIC{$x \in \mathbb{R} \to y \in \mathbb{R} \to \lvert y - x \rvert = \lvert x - y \rvert$}
\end{ptree}
\noindent$\Pi_{64}$:
\begin{ptree}
\AXC{$\stackrel{\Pi_{63}}{\vdots}$}\noLine
\UIC{$x \in \mathbb{R} \to y \in \mathbb{R} \to \lvert y - x \rvert = \lvert x - y \rvert$}
\AXC{$(x \in \mathbb{R} \to y \in \mathbb{R} \to \lvert y - x \rvert = \lvert x - y \rvert) \to \varphi_{232}$\lbl{A1}}
\BIC{$\varphi_{232}$}
\end{ptree}
\noindent$\Pi_{65}$:
\begin{ptree}
\AXC{$\varphi_{32}$\lbl{A6}}
\UIC{$\forall w\,(\varphi_{32})$}
\UIC{$\forall z\,\forall w\,(\varphi_{32})$}
\AXC{$\forall z\,\forall w\,(\varphi_{32}) \to \varphi_{110}$\lbl{A2}}
\BIC{$\varphi_{110}$}
\AXC{$\varphi_{110} \to y - x + x = y \to \varphi_{49}$\lbl{A2}}
\BIC{$y - x + x = y \to \varphi_{49}$}
\AXC{$\stackrel{\Pi_{60}}{\vdots}$}\noLine
\UIC{$\varphi_{195}$}
\BIC{$x \in \mathbb{R} \to y \in \mathbb{R} \to \lvert y \rvert \leq \lvert y - x \rvert + \lvert x \rvert$}
\AXC{$\stackrel{\Pi_{64}}{\vdots}$}\noLine
\UIC{$\varphi_{232}$}
\BIC{$\varphi_{118} \to \varphi_{103} \to x \in \mathbb{R} \to y \in \mathbb{R} \to \lvert y \rvert \leq \lvert x - y \rvert + \lvert x \rvert$}
\end{ptree}
\noindent$\Pi_{66}$:
\begin{ptree}
\AXC{$\forall x \in \mathbb{R}\,(-x \in \mathbb{R})$\lbl{(R1)}}
\AXC{$\forall x \in \mathbb{R}\,(-x \in \mathbb{R}) \to \lvert x \rvert \in \mathbb{R} \to -\lvert x \rvert \in \mathbb{R}$\lbl{A2}}
\BIC{$\lvert x \rvert \in \mathbb{R} \to -\lvert x \rvert \in \mathbb{R}$}
\AXC{$(\lvert x \rvert \in \mathbb{R} \to -\lvert x \rvert \in \mathbb{R}) \to \varphi_{117}$\lbl{A1}}
\BIC{$\varphi_{117}$}
\end{ptree}
\noindent$\Pi_{67}$:
\begin{ptree}
\AXC{$\stackrel{\Pi_{20}}{\vdots}$}\noLine
\UIC{$x \in \mathbb{R} \to y \in \mathbb{R} \to \lvert x \rvert \in \mathbb{R}$}
\AXC{$\stackrel{\Pi_{66}}{\vdots}$}\noLine
\UIC{$\varphi_{117}$}
\BIC{$x \in \mathbb{R} \to y \in \mathbb{R} \to -\lvert x \rvert \in \mathbb{R}$}
\end{ptree}
\noindent$\Pi_{68}$:
\begin{ptree}
\AXC{$\forall x\,\forall y\,(x \in \mathbb{R} \to y \in \mathbb{R} \to x + y \in \mathbb{R})$\lbl{(R1)}}
\AXC{$\varphi_{135}$\lbl{A2}}
\BIC{$\forall y\,(v_{662} \in \mathbb{R} \to y \in \mathbb{R} \to v_{662} + y \in \mathbb{R})$}
\AXC{$\varphi_{151}$\lbl{A2}}
\BIC{$v_{662} \in \mathbb{R} \to v_{663} \in \mathbb{R} \to v_{662} + v_{663} \in \mathbb{R}$}
\UIC{$\forall v_{663}\,(v_{662} \in \mathbb{R} \to v_{663} \in \mathbb{R} \to v_{662} + v_{663} \in \mathbb{R})$}
\UIC{$\varphi_{34}$}
\AXC{$\varphi_{34} \to \varphi_{61}$\lbl{A2}}
\BIC{$\varphi_{61}$}
\end{ptree}
\noindent$\Pi_{69}$:
\begin{ptree}
\AXC{$\stackrel{\Pi_{16}}{\vdots}$}\noLine
\UIC{$x \in \mathbb{R} \to y \in \mathbb{R} \to \lvert x - y \rvert \in \mathbb{R}$}
\AXC{$\stackrel{\Pi_{68}}{\vdots}$}\noLine
\UIC{$\varphi_{61}$}
\AXC{$\varphi_{61} \to \lvert x - y \rvert \in \mathbb{R} \to \lvert x \rvert \in \mathbb{R} \to \lvert x - y \rvert + \lvert x \rvert \in \mathbb{R}$\lbl{A2}}
\BIC{$\lvert x - y \rvert \in \mathbb{R} \to \lvert x \rvert \in \mathbb{R} \to \lvert x - y \rvert + \lvert x \rvert \in \mathbb{R}$}
\AXC{$\varphi_{224}$\lbl{A1}}
\BIC{$(x \in \mathbb{R} \to y \in \mathbb{R} \to \lvert x - y \rvert \in \mathbb{R}) \to \varphi_{136}$}
\BIC{$\varphi_{136}$}
\end{ptree}
\noindent$\Pi_{70}$:
\begin{ptree}
\AXC{$\varphi_{42}$\lbl{[F19]}}
\AXC{$\varphi_{42} \to \varphi_{62}$\lbl{A2}}
\BIC{$\varphi_{62}$}
\AXC{$\varphi_{62} \to \varphi_{83}$\lbl{A2}}
\BIC{$\varphi_{83}$}
\AXC{$\varphi_{83} \to v_{659} \in \mathbb{R} \to \varphi_{50}$\lbl{A2}}
\BIC{$v_{659} \in \mathbb{R} \to \varphi_{50}$}
\UIC{$\forall v_{661}\,(v_{659} \in \mathbb{R} \to \varphi_{50})$}
\UIC{$\forall v_{660}\,\forall v_{661}\,(v_{659} \in \mathbb{R} \to \varphi_{50})$}
\UIC{$\forall v_{659}\,\forall v_{660}\,\forall v_{661}\,(v_{659} \in \mathbb{R} \to \varphi_{50})$}
\AXC{$\forall v_{659}\,\forall v_{660}\,\forall v_{661}\,(v_{659} \in \mathbb{R} \to \varphi_{50}) \to \varphi_{143}$\lbl{A2}}
\BIC{$\varphi_{143}$}
\end{ptree}
\noindent$\Pi_{71}$:
\begin{ptree}
\AXC{$\stackrel{\Pi_{11}}{\vdots}$}\noLine
\UIC{$x \in \mathbb{R} \to y \in \mathbb{R} \to \lvert y \rvert \in \mathbb{R}$}
\AXC{$\stackrel{\Pi_{70}}{\vdots}$}\noLine
\UIC{$\varphi_{143}$}
\AXC{$\varphi_{143} \to \varphi_{199}$\lbl{A2}}
\BIC{$\varphi_{199}$}
\AXC{$\varphi_{237}$\lbl{A2}}
\BIC{$\lvert y \rvert \in \mathbb{R} \to \lvert x - y \rvert + \lvert x \rvert \in \mathbb{R} \to -\lvert x \rvert \in \mathbb{R} \to \varphi_{113}$}
\AXC{$\varphi_{245}$\lbl{A1}}
\BIC{$\varphi_{233}$}
\BIC{$(x \in \mathbb{R} \to y \in \mathbb{R} \to \lvert x - y \rvert + \lvert x \rvert \in \mathbb{R}) \to \varphi_{214}$}
\end{ptree}
\noindent$\Pi_{72}$:
\begin{ptree}
\AXC{$\stackrel{\Pi_{20}}{\vdots}$}\noLine
\UIC{$x \in \mathbb{R} \to y \in \mathbb{R} \to \lvert x \rvert \in \mathbb{R}$}
\AXC{$\stackrel{\Pi_{69}}{\vdots}$}\noLine
\UIC{$\varphi_{136}$}
\BIC{$x \in \mathbb{R} \to y \in \mathbb{R} \to \lvert x - y \rvert + \lvert x \rvert \in \mathbb{R}$}
\AXC{$\stackrel{\Pi_{71}}{\vdots}$}\noLine
\UIC{$(x \in \mathbb{R} \to y \in \mathbb{R} \to \lvert x - y \rvert + \lvert x \rvert \in \mathbb{R}) \to \varphi_{214}$}
\BIC{$\varphi_{214}$}
\end{ptree}
\noindent$\Pi_{73}$:
\begin{ptree}
\AXC{$\stackrel{\Pi_{50}}{\vdots}$}\noLine
\UIC{$\varphi_{103}$}
\AXC{$\stackrel{\Pi_{51}}{\vdots}$}\noLine
\UIC{$\varphi_{118}$}
\AXC{$\stackrel{\Pi_{65}}{\vdots}$}\noLine
\UIC{$\varphi_{118} \to \varphi_{103} \to x \in \mathbb{R} \to y \in \mathbb{R} \to \lvert y \rvert \leq \lvert x - y \rvert + \lvert x \rvert$}
\BIC{$\varphi_{103} \to x \in \mathbb{R} \to y \in \mathbb{R} \to \lvert y \rvert \leq \lvert x - y \rvert + \lvert x \rvert$}
\BIC{$x \in \mathbb{R} \to y \in \mathbb{R} \to \lvert y \rvert \leq \lvert x - y \rvert + \lvert x \rvert$}
\AXC{$\stackrel{\Pi_{67}}{\vdots}$}\noLine
\UIC{$x \in \mathbb{R} \to y \in \mathbb{R} \to -\lvert x \rvert \in \mathbb{R}$}
\AXC{$\stackrel{\Pi_{72}}{\vdots}$}\noLine
\UIC{$\varphi_{214}$}
\BIC{$(x \in \mathbb{R} \to y \in \mathbb{R} \to \lvert y \rvert \leq \lvert x - y \rvert + \lvert x \rvert) \to \varphi_{102}$}
\BIC{$\varphi_{102}$}
\end{ptree}
\noindent$\Pi_{74}$:
\begin{ptree}
\AXC{$y = z \to x = y \to x = z$\lbl{A6}}
\AXC{$(y = z \to x = y \to x = z) \to x = y \to y = z \to x = z$\lbl{A1}}
\BIC{$x = y \to y = z \to x = z$}
\UIC{$\forall z\,(x = y \to y = z \to x = z)$}
\UIC{$\forall y\,\forall z\,(x = y \to y = z \to x = z)$}
\UIC{$\forall x\,\forall y\,\forall z\,(x = y \to y = z \to x = z)$}
\AXC{$\varphi_{10}$\lbl{A2}}
\BIC{$\forall y\,\forall z\,(v_{670} = y \to y = z \to v_{670} = z)$}
\AXC{$\varphi_{19}$\lbl{A2}}
\BIC{$\forall z\,(v_{670} = v_{671} \to v_{671} = z \to v_{670} = z)$}
\AXC{$\varphi_{38}$\lbl{A2}}
\BIC{$v_{670} = v_{671} \to v_{671} = v_{672} \to v_{670} = v_{672}$}
\UIC{$\forall v_{672}\,(v_{670} = v_{671} \to v_{671} = v_{672} \to v_{670} = v_{672})$}
\UIC{$\varphi_{6}$}
\UIC{$\forall v_{670}\,(\varphi_{6})$}
\end{ptree}
\noindent$\Pi_{75}$:
\begin{ptree}
\AXC{$\forall x \in \mathbb{R}\,(x + 0 = x)$\lbl{(R4)}}
\AXC{$\varphi_{76}$\lbl{A2}}
\BIC{$\lvert x - y \rvert \in \mathbb{R} \to \lvert x - y \rvert + 0 = \lvert x - y \rvert$}
\AXC{$(\lvert x - y \rvert \in \mathbb{R} \to \lvert x - y \rvert + 0 = \lvert x - y \rvert) \to \varphi_{141}$\lbl{A1}}
\BIC{$\varphi_{141}$}
\end{ptree}
\noindent$\Pi_{76}$:
\begin{ptree}
\AXC{$y = z \to x = y \to x = z$\lbl{A6}}
\AXC{$(y = z \to x = y \to x = z) \to x = y \to y = z \to x = z$\lbl{A1}}
\BIC{$x = y \to y = z \to x = z$}
\UIC{$\forall z\,(x = y \to y = z \to x = z)$}
\UIC{$\forall y\,\forall z\,(x = y \to y = z \to x = z)$}
\UIC{$\forall x\,\forall y\,\forall z\,(x = y \to y = z \to x = z)$}
\AXC{$\varphi_{11}$\lbl{A2}}
\BIC{$\forall y\,\forall z\,(v_{667} = y \to y = z \to v_{667} = z)$}
\AXC{$\varphi_{20}$\lbl{A2}}
\BIC{$\forall z\,(v_{667} = v_{668} \to v_{668} = z \to v_{667} = z)$}
\AXC{$\varphi_{39}$\lbl{A2}}
\BIC{$v_{667} = v_{668} \to v_{668} = v_{669} \to v_{667} = v_{669}$}
\UIC{$\forall v_{669}\,(v_{667} = v_{668} \to v_{668} = v_{669} \to v_{667} = v_{669})$}
\UIC{$\varphi_{7}$}
\UIC{$\forall v_{667}\,(\varphi_{7})$}
\end{ptree}
\noindent$\Pi_{77}$:
\begin{ptree}
\AXC{$z = w \to \varphi_{161}$\lbl{A6}}
\UIC{$\forall w\,(z = w \to \varphi_{161})$}
\UIC{$\forall z\,\forall w\,(z = w \to \varphi_{161})$}
\AXC{$\forall z\,\forall w\,(z = w \to \varphi_{161}) \to \varphi_{204}$\lbl{A2}}
\BIC{$\varphi_{204}$}
\AXC{$\varphi_{204} \to \varphi_{202}$\lbl{A2}}
\BIC{$\varphi_{202}$}
\end{ptree}
\noindent$\Pi_{78}$:
\begin{ptree}
\AXC{$\stackrel{\Pi_{20}}{\vdots}$}\noLine
\UIC{$x \in \mathbb{R} \to y \in \mathbb{R} \to \lvert x \rvert \in \mathbb{R}$}
\AXC{$\forall x \in \mathbb{R}\,(x + -x = 0)$\lbl{(R5)}}
\AXC{$\varphi_{64}$\lbl{A2}}
\BIC{$\lvert x \rvert \in \mathbb{R} \to \lvert x \rvert + -\lvert x \rvert = 0$}
\AXC{$(\lvert x \rvert \in \mathbb{R} \to \lvert x \rvert + -\lvert x \rvert = 0) \to \varphi_{121}$\lbl{A1}}
\BIC{$\varphi_{121}$}
\BIC{$x \in \mathbb{R} \to y \in \mathbb{R} \to \lvert x \rvert + -\lvert x \rvert = 0$}
\end{ptree}
\noindent$\Pi_{79}$:
\begin{ptree}
\AXC{$\stackrel{\Pi_{78}}{\vdots}$}\noLine
\UIC{$x \in \mathbb{R} \to y \in \mathbb{R} \to \lvert x \rvert + -\lvert x \rvert = 0$}
\AXC{$(x \in \mathbb{R} \to y \in \mathbb{R} \to \lvert x \rvert + -\lvert x \rvert = 0) \to \varphi_{63} \to \varphi_{202} \to \varphi_{84}$\lbl{A1}}
\BIC{$\varphi_{63} \to \varphi_{202} \to \varphi_{84}$}
\end{ptree}
\noindent$\Pi_{80}$:
\begin{ptree}
\AXC{$\stackrel{\Pi_{77}}{\vdots}$}\noLine
\UIC{$\varphi_{202}$}
\AXC{$x = x$\lbl{A5}}
\UIC{$\forall x\,(x = x)$}
\AXC{$\forall x\,(x = x) \to \varphi_{63}$\lbl{A2}}
\BIC{$\varphi_{63}$}
\AXC{$\stackrel{\Pi_{79}}{\vdots}$}\noLine
\UIC{$\varphi_{63} \to \varphi_{202} \to \varphi_{84}$}
\BIC{$\varphi_{202} \to \varphi_{84}$}
\BIC{$\varphi_{84}$}
\end{ptree}
\noindent$\Pi_{81}$:
\begin{ptree}
\AXC{$\varphi_{52}$\lbl{(R2)}}
\AXC{$\varphi_{52} \to \varphi_{74}$\lbl{A2}}
\BIC{$\varphi_{74}$}
\AXC{$\varphi_{74} \to \varphi_{97}$\lbl{A2}}
\BIC{$\varphi_{97}$}
\AXC{$\varphi_{97} \to v_{664} \in \mathbb{R} \to v_{665} \in \mathbb{R} \to \varphi_{14}$\lbl{A2}}
\BIC{$v_{664} \in \mathbb{R} \to v_{665} \in \mathbb{R} \to \varphi_{14}$}
\UIC{$\forall v_{666}\,(v_{664} \in \mathbb{R} \to v_{665} \in \mathbb{R} \to \varphi_{14})$}
\UIC{$\forall v_{665}\,\forall v_{666}\,(v_{664} \in \mathbb{R} \to v_{665} \in \mathbb{R} \to \varphi_{14})$}
\UIC{$\forall v_{664}\,\forall v_{665}\,\forall v_{666}\,(v_{664} \in \mathbb{R} \to v_{665} \in \mathbb{R} \to \varphi_{14})$}
\AXC{$\varphi_{227}$\lbl{A2}}
\BIC{$\varphi_{175}$}
\AXC{$\varphi_{175} \to \varphi_{176}$\lbl{A2}}
\BIC{$\varphi_{176}$}
\end{ptree}
\noindent$\Pi_{82}$:
\begin{ptree}
\AXC{$\stackrel{\Pi_{81}}{\vdots}$}\noLine
\UIC{$\varphi_{176}$}
\AXC{$\varphi_{176} \to \lvert x - y \rvert \in \mathbb{R} \to \lvert x \rvert \in \mathbb{R} \to -\lvert x \rvert \in \mathbb{R} \to \varphi_{66}$\lbl{A2}}
\BIC{$\lvert x - y \rvert \in \mathbb{R} \to \lvert x \rvert \in \mathbb{R} \to -\lvert x \rvert \in \mathbb{R} \to \varphi_{66}$}
\AXC{$(\lvert x - y \rvert \in \mathbb{R} \to \lvert x \rvert \in \mathbb{R} \to -\lvert x \rvert \in \mathbb{R} \to \varphi_{66}) \to \varphi_{230}$\lbl{A1}}
\BIC{$\varphi_{230}$}
\end{ptree}
\noindent$\Pi_{83}$:
\begin{ptree}
\AXC{$\stackrel{\Pi_{67}}{\vdots}$}\noLine
\UIC{$x \in \mathbb{R} \to y \in \mathbb{R} \to -\lvert x \rvert \in \mathbb{R}$}
\AXC{$\stackrel{\Pi_{20}}{\vdots}$}\noLine
\UIC{$x \in \mathbb{R} \to y \in \mathbb{R} \to \lvert x \rvert \in \mathbb{R}$}
\AXC{$\stackrel{\Pi_{16}}{\vdots}$}\noLine
\UIC{$x \in \mathbb{R} \to y \in \mathbb{R} \to \lvert x - y \rvert \in \mathbb{R}$}
\AXC{$\stackrel{\Pi_{82}}{\vdots}$}\noLine
\UIC{$\varphi_{230}$}
\BIC{$(x \in \mathbb{R} \to y \in \mathbb{R} \to \lvert x \rvert \in \mathbb{R}) \to \varphi_{200}$}
\BIC{$\varphi_{200}$}
\BIC{$x \in \mathbb{R} \to y \in \mathbb{R} \to \varphi_{66}$}
\end{ptree}
\noindent$\Pi_{84}$:
\begin{ptree}
\AXC{$\stackrel{\Pi_{80}}{\vdots}$}\noLine
\UIC{$\varphi_{84}$}
\AXC{$\stackrel{\Pi_{83}}{\vdots}$}\noLine
\UIC{$x \in \mathbb{R} \to y \in \mathbb{R} \to \varphi_{66}$}
\AXC{$(x \in \mathbb{R} \to y \in \mathbb{R} \to \varphi_{66}) \to \varphi_{84} \to (\varphi_{66} \to \varphi_{160}) \to \varphi_{85}$\lbl{A1}}
\BIC{$\varphi_{84} \to (\varphi_{66} \to \varphi_{160}) \to \varphi_{85}$}
\BIC{$(\varphi_{66} \to \varphi_{160}) \to \varphi_{85}$}
\end{ptree}
\noindent$\Pi_{85}$:
\begin{ptree}
\AXC{$\stackrel{\Pi_{76}}{\vdots}$}\noLine
\UIC{$\forall v_{667}\,(\varphi_{7})$}
\AXC{$\forall v_{667}\,(\varphi_{7}) \to \varphi_{156}$\lbl{A2}}
\BIC{$\varphi_{156}$}
\AXC{$\varphi_{156} \to \varphi_{210}$\lbl{A2}}
\BIC{$\varphi_{210}$}
\AXC{$\varphi_{210} \to \varphi_{66} \to \varphi_{160}$\lbl{A2}}
\BIC{$\varphi_{66} \to \varphi_{160}$}
\AXC{$\stackrel{\Pi_{84}}{\vdots}$}\noLine
\UIC{$(\varphi_{66} \to \varphi_{160}) \to \varphi_{85}$}
\BIC{$\varphi_{85}$}
\end{ptree}
\noindent$\Pi_{86}$:
\begin{ptree}
\AXC{$\stackrel{\Pi_{16}}{\vdots}$}\noLine
\UIC{$x \in \mathbb{R} \to y \in \mathbb{R} \to \lvert x - y \rvert \in \mathbb{R}$}
\AXC{$\stackrel{\Pi_{75}}{\vdots}$}\noLine
\UIC{$\varphi_{141}$}
\BIC{$x \in \mathbb{R} \to y \in \mathbb{R} \to \lvert x - y \rvert + 0 = \lvert x - y \rvert$}
\AXC{$\stackrel{\Pi_{85}}{\vdots}$}\noLine
\UIC{$\varphi_{85}$}
\AXC{$\varphi_{85} \to (x \in \mathbb{R} \to y \in \mathbb{R} \to \lvert x - y \rvert + 0 = \lvert x - y \rvert) \to \varphi_{219}$\lbl{A1}}
\BIC{$(x \in \mathbb{R} \to y \in \mathbb{R} \to \lvert x - y \rvert + 0 = \lvert x - y \rvert) \to \varphi_{219}$}
\BIC{$\varphi_{219}$}
\end{ptree}
\noindent$\Pi_{87}$:
\begin{ptree}
\AXC{$\stackrel{\Pi_{74}}{\vdots}$}\noLine
\UIC{$\forall v_{670}\,(\varphi_{6})$}
\AXC{$\forall v_{670}\,(\varphi_{6}) \to \varphi_{155}$\lbl{A2}}
\BIC{$\varphi_{155}$}
\AXC{$\varphi_{155} \to \varphi_{181}$\lbl{A2}}
\BIC{$\varphi_{181}$}
\AXC{$\varphi_{181} \to \lvert x - y \rvert + \lvert x \rvert + -\lvert x \rvert = \lvert x - y \rvert + 0 \to \varphi_{90}$\lbl{A2}}
\BIC{$\lvert x - y \rvert + \lvert x \rvert + -\lvert x \rvert = \lvert x - y \rvert + 0 \to \varphi_{90}$}
\AXC{$\stackrel{\Pi_{86}}{\vdots}$}\noLine
\UIC{$\varphi_{219}$}
\BIC{$x \in \mathbb{R} \to y \in \mathbb{R} \to \lvert x - y \rvert + \lvert x \rvert + -\lvert x \rvert = \lvert x - y \rvert$}
\end{ptree}
\noindent$\Pi_{88}$:
\begin{ptree}
\AXC{$\stackrel{\Pi_{49}}{\vdots}$}\noLine
\UIC{$\lvert x - y \rvert + \lvert x \rvert + -\lvert x \rvert = \lvert x - y \rvert \to \varphi_{138}$}
\AXC{$\stackrel{\Pi_{73}}{\vdots}$}\noLine
\UIC{$\varphi_{102}$}
\AXC{$\stackrel{\Pi_{87}}{\vdots}$}\noLine
\UIC{$x \in \mathbb{R} \to y \in \mathbb{R} \to \lvert x - y \rvert + \lvert x \rvert + -\lvert x \rvert = \lvert x - y \rvert$}
\AXC{$\varphi_{247}$\lbl{A1}}
\BIC{$\varphi_{102} \to (\lvert x - y \rvert + \lvert x \rvert + -\lvert x \rvert = \lvert x - y \rvert \to \varphi_{138}) \to \varphi_{217}$}
\BIC{$(\lvert x - y \rvert + \lvert x \rvert + -\lvert x \rvert = \lvert x - y \rvert \to \varphi_{138}) \to \varphi_{217}$}
\BIC{$\varphi_{217}$}
\end{ptree}
\noindent$\Pi_{89}$:
\begin{ptree}
\AXC{$z = w \to \lvert y \rvert + -\lvert x \rvert = z \to \lvert y \rvert + -\lvert x \rvert = w$\lbl{A6}}
\UIC{$\forall w\,(z = w \to \lvert y \rvert + -\lvert x \rvert = z \to \lvert y \rvert + -\lvert x \rvert = w)$}
\UIC{$\varphi_{23}$}
\AXC{$\varphi_{23} \to \varphi_{183}$\lbl{A2}}
\BIC{$\varphi_{183}$}
\AXC{$\varphi_{183} \to \lvert x - y \rvert + \lvert x \rvert + -\lvert x \rvert = \lvert x - y \rvert \to \varphi_{126}$\lbl{A2}}
\BIC{$\lvert x - y \rvert + \lvert x \rvert + -\lvert x \rvert = \lvert x - y \rvert \to \varphi_{126}$}
\AXC{$\stackrel{\Pi_{88}}{\vdots}$}\noLine
\UIC{$\varphi_{217}$}
\BIC{$x \in \mathbb{R} \to y \in \mathbb{R} \to \lvert y \rvert + -\lvert x \rvert \leq \lvert x - y \rvert$}
\end{ptree}
\noindent$\Pi_{90}$:
\begin{ptree}
\AXC{$\stackrel{\Pi_{40}}{\vdots}$}\noLine
\UIC{$\forall x\,\forall y\,(x = y \to y = x)$}
\AXC{$\forall x\,\forall y\,(x = y \to y = x) \to \forall y\,(v_{675} = y \to y = v_{675})$\lbl{A2}}
\BIC{$\forall y\,(v_{675} = y \to y = v_{675})$}
\AXC{$\varphi_{2}$\lbl{A2}}
\BIC{$v_{675} = v_{676} \to v_{676} = v_{675}$}
\UIC{$\forall v_{676}\,(v_{675} = v_{676} \to v_{676} = v_{675})$}
\UIC{$\forall v_{675}\,\forall v_{676}\,(v_{675} = v_{676} \to v_{676} = v_{675})$}
\AXC{$\varphi_{107}$\lbl{A2}}
\BIC{$\forall v_{676}\,(\lvert y \rvert - \lvert x \rvert = v_{676} \to v_{676} = \lvert y \rvert - \lvert x \rvert)$}
\AXC{$\varphi_{187}$\lbl{A2}}
\BIC{$\varphi_{86}$}
\end{ptree}
\noindent$\Pi_{91}$:
\begin{ptree}
\AXC{$\stackrel{\Pi_{90}}{\vdots}$}\noLine
\UIC{$\varphi_{86}$}
\AXC{$\forall x\,\forall y\,(x - y = x + -y)$\lbl{Def$-$}}
\AXC{$\varphi_{12}$\lbl{A2}}
\BIC{$\forall y\,(v_{673} - y = v_{673} + -y)$}
\AXC{$\varphi_{21}$\lbl{A2}}
\BIC{$v_{673} - v_{674} = v_{673} + -v_{674}$}
\UIC{$\forall v_{674}\,(v_{673} - v_{674} = v_{673} + -v_{674})$}
\UIC{$\forall v_{673}\,\forall v_{674}\,(v_{673} - v_{674} = v_{673} + -v_{674})$}
\AXC{$\varphi_{104}$\lbl{A2}}
\BIC{$\forall v_{674}\,(\lvert y \rvert - v_{674} = \lvert y \rvert + -v_{674})$}
\AXC{$\varphi_{95}$\lbl{A2}}
\BIC{$\lvert y \rvert - \lvert x \rvert = \lvert y \rvert + -\lvert x \rvert$}
\AXC{$\varphi_{207}$\lbl{A1}}
\BIC{$\varphi_{86} \to \lvert y \rvert + -\lvert x \rvert = \lvert y \rvert - \lvert x \rvert$}
\BIC{$\lvert y \rvert + -\lvert x \rvert = \lvert y \rvert - \lvert x \rvert$}
\end{ptree}
\noindent$\Pi_{92}$:
\begin{ptree}
\AXC{$\stackrel{\Pi_{89}}{\vdots}$}\noLine
\UIC{$x \in \mathbb{R} \to y \in \mathbb{R} \to \lvert y \rvert + -\lvert x \rvert \leq \lvert x - y \rvert$}
\AXC{$\stackrel{\Pi_{91}}{\vdots}$}\noLine
\UIC{$\lvert y \rvert + -\lvert x \rvert = \lvert y \rvert - \lvert x \rvert$}
\AXC{$\lvert y \rvert + -\lvert x \rvert = \lvert y \rvert - \lvert x \rvert \to \varphi_{239}$\lbl{A1}}
\BIC{$\varphi_{239}$}
\BIC{$\varphi_{165} \to \varphi_{148} \to x \in \mathbb{R} \to y \in \mathbb{R} \to \lvert y \rvert - \lvert x \rvert \leq \lvert x - y \rvert$}
\end{ptree}
\noindent$\Pi_{93}$:
\begin{ptree}
\AXC{$\stackrel{\Pi_{40}}{\vdots}$}\noLine
\UIC{$\forall x\,\forall y\,(x = y \to y = x)$}
\AXC{$\forall x\,\forall y\,(x = y \to y = x) \to \forall y\,(v_{679} = y \to y = v_{679})$\lbl{A2}}
\BIC{$\forall y\,(v_{679} = y \to y = v_{679})$}
\AXC{$\varphi_{3}$\lbl{A2}}
\BIC{$v_{679} = v_{680} \to v_{680} = v_{679}$}
\UIC{$\forall v_{680}\,(v_{679} = v_{680} \to v_{680} = v_{679})$}
\UIC{$\forall v_{679}\,\forall v_{680}\,(v_{679} = v_{680} \to v_{680} = v_{679})$}
\end{ptree}
\noindent$\Pi_{94}$:
\begin{ptree}
\AXC{$\stackrel{\Pi_{93}}{\vdots}$}\noLine
\UIC{$\forall v_{679}\,\forall v_{680}\,(v_{679} = v_{680} \to v_{680} = v_{679})$}
\AXC{$\forall v_{679}\,\forall v_{680}\,(v_{679} = v_{680} \to v_{680} = v_{679}) \to \varphi_{25}$\lbl{A2}}
\BIC{$\varphi_{25}$}
\AXC{$\varphi_{25} \to \varphi_{87}$\lbl{A2}}
\BIC{$\varphi_{87}$}
\end{ptree}
\noindent$\Pi_{95}$:
\begin{ptree}
\AXC{$\forall x\,\forall y\,(x \in \mathbb{R} \to y \in \mathbb{R} \to -(y - x) = x - y)$\lbl{[F14]}}
\AXC{$\varphi_{157}$\lbl{A2}}
\BIC{$\forall y\,(v_{677} \in \mathbb{R} \to y \in \mathbb{R} \to -(y - v_{677}) = v_{677} - y)$}
\AXC{$\varphi_{171}$\lbl{A2}}
\BIC{$\varphi_{26}$}
\UIC{$\forall v_{678}\,(\varphi_{26})$}
\UIC{$\forall v_{677}\,\forall v_{678}\,(\varphi_{26})$}
\AXC{$\forall v_{677}\,\forall v_{678}\,(\varphi_{26}) \to \varphi_{69}$\lbl{A2}}
\BIC{$\varphi_{69}$}
\AXC{$\varphi_{69} \to \varphi_{70}$\lbl{A2}}
\BIC{$\varphi_{70}$}
\end{ptree}
\noindent$\Pi_{96}$:
\begin{ptree}
\AXC{$\stackrel{\Pi_{20}}{\vdots}$}\noLine
\UIC{$x \in \mathbb{R} \to y \in \mathbb{R} \to \lvert x \rvert \in \mathbb{R}$}
\AXC{$\stackrel{\Pi_{11}}{\vdots}$}\noLine
\UIC{$x \in \mathbb{R} \to y \in \mathbb{R} \to \lvert y \rvert \in \mathbb{R}$}
\AXC{$\stackrel{\Pi_{95}}{\vdots}$}\noLine
\UIC{$\varphi_{70}$}
\AXC{$\varphi_{70} \to (x \in \mathbb{R} \to y \in \mathbb{R} \to \lvert y \rvert \in \mathbb{R}) \to \varphi_{152}$\lbl{A1}}
\BIC{$(x \in \mathbb{R} \to y \in \mathbb{R} \to \lvert y \rvert \in \mathbb{R}) \to \varphi_{152}$}
\BIC{$\varphi_{152}$}
\BIC{$x \in \mathbb{R} \to y \in \mathbb{R} \to -(\lvert x \rvert - \lvert y \rvert) = \lvert y \rvert - \lvert x \rvert$}
\end{ptree}
\noindent$\Pi_{97}$:
\begin{ptree}
\AXC{$\stackrel{\Pi_{94}}{\vdots}$}\noLine
\UIC{$\varphi_{87}$}
\AXC{$\stackrel{\Pi_{96}}{\vdots}$}\noLine
\UIC{$x \in \mathbb{R} \to y \in \mathbb{R} \to -(\lvert x \rvert - \lvert y \rvert) = \lvert y \rvert - \lvert x \rvert$}
\AXC{$\varphi_{220}$\lbl{A1}}
\BIC{$\varphi_{87} \to x \in \mathbb{R} \to y \in \mathbb{R} \to \lvert y \rvert - \lvert x \rvert = -(\lvert x \rvert - \lvert y \rvert)$}
\BIC{$x \in \mathbb{R} \to y \in \mathbb{R} \to \lvert y \rvert - \lvert x \rvert = -(\lvert x \rvert - \lvert y \rvert)$}
\AXC{$\varphi_{242}$\lbl{A1}}
\BIC{$\varphi_{240}$}
\end{ptree}
\noindent$\Pi_{98}$:
\begin{ptree}
\AXC{$\stackrel{\Pi_{45}}{\vdots}$}\noLine
\UIC{$\varphi_{147}$}
\AXC{$\stackrel{\Pi_{46}}{\vdots}$}\noLine
\UIC{$\varphi_{164}$}
\AXC{$\stackrel{\Pi_{47}}{\vdots}$}\noLine
\UIC{$\varphi_{148}$}
\AXC{$\stackrel{\Pi_{48}}{\vdots}$}\noLine
\UIC{$\varphi_{165}$}
\AXC{$\stackrel{\Pi_{92}}{\vdots}$}\noLine
\UIC{$\varphi_{165} \to \varphi_{148} \to x \in \mathbb{R} \to y \in \mathbb{R} \to \lvert y \rvert - \lvert x \rvert \leq \lvert x - y \rvert$}
\BIC{$\varphi_{148} \to x \in \mathbb{R} \to y \in \mathbb{R} \to \lvert y \rvert - \lvert x \rvert \leq \lvert x - y \rvert$}
\BIC{$x \in \mathbb{R} \to y \in \mathbb{R} \to \lvert y \rvert - \lvert x \rvert \leq \lvert x - y \rvert$}
\AXC{$\stackrel{\Pi_{97}}{\vdots}$}\noLine
\UIC{$\varphi_{240}$}
\BIC{$\varphi_{164} \to \varphi_{147} \to x \in \mathbb{R} \to y \in \mathbb{R} \to -(\lvert x \rvert - \lvert y \rvert) \leq \lvert x - y \rvert$}
\BIC{$\varphi_{147} \to x \in \mathbb{R} \to y \in \mathbb{R} \to -(\lvert x \rvert - \lvert y \rvert) \leq \lvert x - y \rvert$}
\BIC{$x \in \mathbb{R} \to y \in \mathbb{R} \to -(\lvert x \rvert - \lvert y \rvert) \leq \lvert x - y \rvert$}
\end{ptree}
\noindent$\Pi_{99}$:
\begin{ptree}
\AXC{$\stackrel{\Pi_{20}}{\vdots}$}\noLine
\UIC{$x \in \mathbb{R} \to y \in \mathbb{R} \to \lvert x \rvert \in \mathbb{R}$}
\AXC{$\forall x\,\forall y\,(x \in \mathbb{R} \to y \in \mathbb{R} \to x - y \in \mathbb{R})$\lbl{[F1]}}
\AXC{$\varphi_{145}$\lbl{A2}}
\BIC{$\forall y\,(\lvert x \rvert \in \mathbb{R} \to y \in \mathbb{R} \to \lvert x \rvert - y \in \mathbb{R})$}
\AXC{$\varphi_{169}$\lbl{A2}}
\BIC{$\lvert x \rvert \in \mathbb{R} \to \lvert y \rvert \in \mathbb{R} \to \lvert x \rvert - \lvert y \rvert \in \mathbb{R}$}
\AXC{$\varphi_{215}$\lbl{A1}}
\BIC{$(x \in \mathbb{R} \to y \in \mathbb{R} \to \lvert x \rvert \in \mathbb{R}) \to \varphi_{127}$}
\BIC{$\varphi_{127}$}
\end{ptree}
\noindent$\Pi_{100}$:
\begin{ptree}
\AXC{$\stackrel{\Pi_{11}}{\vdots}$}\noLine
\UIC{$x \in \mathbb{R} \to y \in \mathbb{R} \to \lvert y \rvert \in \mathbb{R}$}
\AXC{$\stackrel{\Pi_{99}}{\vdots}$}\noLine
\UIC{$\varphi_{127}$}
\BIC{$x \in \mathbb{R} \to y \in \mathbb{R} \to \lvert x \rvert - \lvert y \rvert \in \mathbb{R}$}
\end{ptree}
\noindent$\Pi_{101}$:
\begin{ptree}
\AXC{$\forall x \in \mathbb{R}\,(-x \in \mathbb{R})$\lbl{(R1)}}
\AXC{$\varphi_{119}$\lbl{A2}}
\BIC{$\lvert x \rvert - \lvert y \rvert \in \mathbb{R} \to -(\lvert x \rvert - \lvert y \rvert) \in \mathbb{R}$}
\AXC{$(\lvert x \rvert - \lvert y \rvert \in \mathbb{R} \to -(\lvert x \rvert - \lvert y \rvert) \in \mathbb{R}) \to \varphi_{162}$\lbl{A1}}
\BIC{$\varphi_{162}$}
\end{ptree}
\noindent$\Pi_{102}$:
\begin{ptree}
\AXC{$\forall x\,\forall y\,(x \in \mathbb{R} \to y \in \mathbb{R} \to x \leq y \to -y \leq -x)$\lbl{[F23]}}
\AXC{$\varphi_{167}$\lbl{A2}}
\BIC{$\forall y\,(v_{681} \in \mathbb{R} \to y \in \mathbb{R} \to v_{681} \leq y \to -y \leq -v_{681})$}
\AXC{$\varphi_{177}$\lbl{A2}}
\BIC{$\varphi_{28}$}
\UIC{$\forall v_{682}\,(\varphi_{28})$}
\UIC{$\forall v_{681}\,\forall v_{682}\,(\varphi_{28})$}
\AXC{$\forall v_{681}\,\forall v_{682}\,(\varphi_{28}) \to \varphi_{170}$\lbl{A2}}
\BIC{$\varphi_{170}$}
\end{ptree}
\noindent$\Pi_{103}$:
\begin{ptree}
\AXC{$\stackrel{\Pi_{102}}{\vdots}$}\noLine
\UIC{$\varphi_{170}$}
\AXC{$\varphi_{170} \to -(\lvert x \rvert - \lvert y \rvert) \in \mathbb{R} \to \lvert x - y \rvert \in \mathbb{R} \to \varphi_{91}$\lbl{A2}}
\BIC{$-(\lvert x \rvert - \lvert y \rvert) \in \mathbb{R} \to \lvert x - y \rvert \in \mathbb{R} \to \varphi_{91}$}
\AXC{$(-(\lvert x \rvert - \lvert y \rvert) \in \mathbb{R} \to \lvert x - y \rvert \in \mathbb{R} \to \varphi_{91}) \to \varphi_{228}$\lbl{A1}}
\BIC{$\varphi_{228}$}
\end{ptree}
\noindent$\Pi_{104}$:
\begin{ptree}
\AXC{$\stackrel{\Pi_{16}}{\vdots}$}\noLine
\UIC{$x \in \mathbb{R} \to y \in \mathbb{R} \to \lvert x - y \rvert \in \mathbb{R}$}
\AXC{$\stackrel{\Pi_{100}}{\vdots}$}\noLine
\UIC{$x \in \mathbb{R} \to y \in \mathbb{R} \to \lvert x \rvert - \lvert y \rvert \in \mathbb{R}$}
\AXC{$\stackrel{\Pi_{101}}{\vdots}$}\noLine
\UIC{$\varphi_{162}$}
\BIC{$x \in \mathbb{R} \to y \in \mathbb{R} \to -(\lvert x \rvert - \lvert y \rvert) \in \mathbb{R}$}
\AXC{$\stackrel{\Pi_{103}}{\vdots}$}\noLine
\UIC{$\varphi_{228}$}
\BIC{$(x \in \mathbb{R} \to y \in \mathbb{R} \to \lvert x - y \rvert \in \mathbb{R}) \to \varphi_{185}$}
\BIC{$\varphi_{185}$}
\end{ptree}
\noindent$\Pi_{105}$:
\begin{ptree}
\AXC{$\forall x \in \mathbb{R}\,(-(-x) = x)$\lbl{[F5]}}
\AXC{$\forall x \in \mathbb{R}\,(-(-x) = x) \to \varphi_{54}$\lbl{A2}}
\BIC{$\varphi_{54}$}
\AXC{$\varphi_{54} \to (x \in \mathbb{R} \to y \in \mathbb{R} \to \lvert x \rvert - \lvert y \rvert \in \mathbb{R}) \to \varphi_{47}$\lbl{A1}}
\BIC{$(x \in \mathbb{R} \to y \in \mathbb{R} \to \lvert x \rvert - \lvert y \rvert \in \mathbb{R}) \to \varphi_{47}$}
\end{ptree}
\noindent$\Pi_{106}$:
\begin{ptree}
\AXC{$\stackrel{\Pi_{100}}{\vdots}$}\noLine
\UIC{$x \in \mathbb{R} \to y \in \mathbb{R} \to \lvert x \rvert - \lvert y \rvert \in \mathbb{R}$}
\AXC{$\stackrel{\Pi_{105}}{\vdots}$}\noLine
\UIC{$(x \in \mathbb{R} \to y \in \mathbb{R} \to \lvert x \rvert - \lvert y \rvert \in \mathbb{R}) \to \varphi_{47}$}
\BIC{$\varphi_{47}$}
\AXC{$\varphi_{47} \to \varphi_{241}$\lbl{A1}}
\BIC{$\varphi_{241}$}
\end{ptree}
\noindent$\Pi_{107}$:
\begin{ptree}
\AXC{$\stackrel{\Pi_{98}}{\vdots}$}\noLine
\UIC{$x \in \mathbb{R} \to y \in \mathbb{R} \to -(\lvert x \rvert - \lvert y \rvert) \leq \lvert x - y \rvert$}
\AXC{$\stackrel{\Pi_{104}}{\vdots}$}\noLine
\UIC{$\varphi_{185}$}
\BIC{$x \in \mathbb{R} \to y \in \mathbb{R} \to -\lvert x - y \rvert \leq -(-(\lvert x \rvert - \lvert y \rvert))$}
\AXC{$\stackrel{\Pi_{106}}{\vdots}$}\noLine
\UIC{$\varphi_{241}$}
\BIC{$(-(-(\lvert x \rvert - \lvert y \rvert)) = \lvert x \rvert - \lvert y \rvert \to \varphi_{105}) \to \varphi_{212}$}
\end{ptree}
\noindent$\Pi_{108}$:
\begin{ptree}
\AXC{$\stackrel{\Pi_{44}}{\vdots}$}\noLine
\UIC{$\varphi_{144}$}
\AXC{$\varphi_{144} \to -(-(\lvert x \rvert - \lvert y \rvert)) = \lvert x \rvert - \lvert y \rvert \to \varphi_{105}$\lbl{A2}}
\BIC{$-(-(\lvert x \rvert - \lvert y \rvert)) = \lvert x \rvert - \lvert y \rvert \to \varphi_{105}$}
\AXC{$\stackrel{\Pi_{107}}{\vdots}$}\noLine
\UIC{$(-(-(\lvert x \rvert - \lvert y \rvert)) = \lvert x \rvert - \lvert y \rvert \to \varphi_{105}) \to \varphi_{212}$}
\BIC{$\varphi_{212}$}
\end{ptree}
\noindent$\Pi_{109}$:
\begin{ptree}
\AXC{$\forall x\,\forall y\,(x \in \mathbb{R} \to y \in \mathbb{R} \to -y \leq x \to x \leq y \to \lvert x \rvert \leq y)$\lbl{[A10]}}
\AXC{$\varphi_{179}$\lbl{A2}}
\BIC{$\varphi_{40}$}
\AXC{$\varphi_{40} \to v_{683} \in \mathbb{R} \to \varphi_{16}$\lbl{A2}}
\BIC{$v_{683} \in \mathbb{R} \to \varphi_{16}$}
\UIC{$\forall v_{684}\,(v_{683} \in \mathbb{R} \to \varphi_{16})$}
\UIC{$\forall v_{683}\,\forall v_{684}\,(v_{683} \in \mathbb{R} \to \varphi_{16})$}
\end{ptree}
\noindent$\Pi_{110}$:
\begin{ptree}
\AXC{$\stackrel{\Pi_{109}}{\vdots}$}\noLine
\UIC{$\forall v_{683}\,\forall v_{684}\,(v_{683} \in \mathbb{R} \to \varphi_{16})$}
\AXC{$\forall v_{683}\,\forall v_{684}\,(v_{683} \in \mathbb{R} \to \varphi_{16}) \to \varphi_{184}$\lbl{A2}}
\BIC{$\varphi_{184}$}
\AXC{$\varphi_{184} \to \lvert x \rvert - \lvert y \rvert \in \mathbb{R} \to \lvert x - y \rvert \in \mathbb{R} \to \varphi_{139}$\lbl{A2}}
\BIC{$\lvert x \rvert - \lvert y \rvert \in \mathbb{R} \to \lvert x - y \rvert \in \mathbb{R} \to \varphi_{139}$}
\end{ptree}
\noindent$\Pi_{111}$:
\begin{ptree}
\AXC{$\stackrel{\Pi_{100}}{\vdots}$}\noLine
\UIC{$x \in \mathbb{R} \to y \in \mathbb{R} \to \lvert x \rvert - \lvert y \rvert \in \mathbb{R}$}
\AXC{$\stackrel{\Pi_{110}}{\vdots}$}\noLine
\UIC{$\lvert x \rvert - \lvert y \rvert \in \mathbb{R} \to \lvert x - y \rvert \in \mathbb{R} \to \varphi_{139}$}
\AXC{$\varphi_{243}$\lbl{A1}}
\BIC{$(x \in \mathbb{R} \to y \in \mathbb{R} \to \lvert x \rvert - \lvert y \rvert \in \mathbb{R}) \to \varphi_{225}$}
\BIC{$\varphi_{225}$}
\end{ptree}
\noindent$\Pi_{112}$:
\begin{ptree}
\AXC{$\stackrel{\Pi_{43}}{\vdots}$}\noLine
\UIC{$\varphi_{130}$}
\AXC{$\varphi_{130} \to -(-(\lvert x \rvert - \lvert y \rvert)) = \lvert x \rvert - \lvert y \rvert \to \varphi_{89}$\lbl{A2}}
\BIC{$-(-(\lvert x \rvert - \lvert y \rvert)) = \lvert x \rvert - \lvert y \rvert \to \varphi_{89}$}
\AXC{$\stackrel{\Pi_{108}}{\vdots}$}\noLine
\UIC{$\varphi_{212}$}
\BIC{$x \in \mathbb{R} \to y \in \mathbb{R} \to -\lvert x - y \rvert \leq \lvert x \rvert - \lvert y \rvert$}
\AXC{$\stackrel{\Pi_{16}}{\vdots}$}\noLine
\UIC{$x \in \mathbb{R} \to y \in \mathbb{R} \to \lvert x - y \rvert \in \mathbb{R}$}
\AXC{$\stackrel{\Pi_{111}}{\vdots}$}\noLine
\UIC{$\varphi_{225}$}
\BIC{$(x \in \mathbb{R} \to y \in \mathbb{R} \to -\lvert x - y \rvert \leq \lvert x \rvert - \lvert y \rvert) \to \varphi_{168}$}
\BIC{$\varphi_{168}$}
\end{ptree}
\noindent 主木:
\begin{ptree}
\AXC{$\stackrel{\Pi_{1}}{\vdots}$}\noLine
\UIC{$\varphi_{146}$}
\AXC{$\stackrel{\Pi_{2}}{\vdots}$}\noLine
\UIC{$\varphi_{163}$}
\AXC{$\stackrel{\Pi_{39}}{\vdots}$}\noLine
\UIC{$x \in \mathbb{R} \to y \in \mathbb{R} \to \lvert x \rvert + -\lvert y \rvert \leq \lvert x - y \rvert$}
\AXC{$\stackrel{\Pi_{42}}{\vdots}$}\noLine
\UIC{$\varphi_{238}$}
\BIC{$\varphi_{163} \to \varphi_{146} \to x \in \mathbb{R} \to y \in \mathbb{R} \to \lvert x \rvert - \lvert y \rvert \leq \lvert x - y \rvert$}
\BIC{$\varphi_{146} \to x \in \mathbb{R} \to y \in \mathbb{R} \to \lvert x \rvert - \lvert y \rvert \leq \lvert x - y \rvert$}
\BIC{$x \in \mathbb{R} \to y \in \mathbb{R} \to \lvert x \rvert - \lvert y \rvert \leq \lvert x - y \rvert$}
\AXC{$\stackrel{\Pi_{112}}{\vdots}$}\noLine
\UIC{$\varphi_{168}$}
\BIC{$x \in \mathbb{R} \to y \in \mathbb{R} \to \lvert \lvert x \rvert - \lvert y \rvert \rvert \leq \lvert x - y \rvert$}
\UIC{$\forall y\,(x \in \mathbb{R} \to y \in \mathbb{R} \to \lvert \lvert x \rvert - \lvert y \rvert \rvert \leq \lvert x - y \rvert)$}
\UIC{$\varphi_{35}$}
\end{ptree}
\begin{abbrevlist}
\item $\varphi_{1} :\equiv \forall y\,(v_{657} = y \to y = v_{657}) \to v_{657} = v_{658} \to v_{658} = v_{657}$
\item $\varphi_{2} :\equiv \forall y\,(v_{675} = y \to y = v_{675}) \to v_{675} = v_{676} \to v_{676} = v_{675}$
\item $\varphi_{3} :\equiv \forall y\,(v_{679} = y \to y = v_{679}) \to v_{679} = v_{680} \to v_{680} = v_{679}$
\item $\varphi_{4} :\equiv \forall v_{655}\,\forall v_{656}\,(v_{654} = v_{655} \to v_{655} = v_{656} \to v_{654} = v_{656})$
\item $\varphi_{5} :\equiv \forall v_{652}\,\forall v_{653}\,(v_{651} = v_{652} \to v_{652} = v_{653} \to v_{651} = v_{653})$
\item $\varphi_{6} :\equiv \forall v_{671}\,\forall v_{672}\,(v_{670} = v_{671} \to v_{671} = v_{672} \to v_{670} = v_{672})$
\item $\varphi_{7} :\equiv \forall v_{668}\,\forall v_{669}\,(v_{667} = v_{668} \to v_{668} = v_{669} \to v_{667} = v_{669})$
\item $\varphi_{8} :\equiv \forall x\,\forall y\,\forall z\,(x = y \to y = z \to x = z) \to \forall y\,\forall z\,(v_{654} = y \to y = z \to v_{654} = z)$
\item $\varphi_{9} :\equiv \forall x\,\forall y\,\forall z\,(x = y \to y = z \to x = z) \to \forall y\,\forall z\,(v_{651} = y \to y = z \to v_{651} = z)$
\item $\varphi_{10} :\equiv \forall x\,\forall y\,\forall z\,(x = y \to y = z \to x = z) \to \forall y\,\forall z\,(v_{670} = y \to y = z \to v_{670} = z)$
\item $\varphi_{11} :\equiv \forall x\,\forall y\,\forall z\,(x = y \to y = z \to x = z) \to \forall y\,\forall z\,(v_{667} = y \to y = z \to v_{667} = z)$
\item $\varphi_{12} :\equiv \forall x\,\forall y\,(x - y = x + -y) \to \forall y\,(v_{673} - y = v_{673} + -y)$
\item $\varphi_{13} :\equiv v_{650} \in \mathbb{R} \to v_{648} + v_{649} + v_{650} = v_{648} + (v_{649} + v_{650})$
\item $\varphi_{14} :\equiv v_{666} \in \mathbb{R} \to v_{664} + v_{665} + v_{666} = v_{664} + (v_{665} + v_{666})$
\item $\varphi_{15} :\equiv \forall x\,\forall y\,(x - y = x + -y) \to \forall y\,(\lvert x \rvert - y = \lvert x \rvert + -y)$
\item $\varphi_{16} :\equiv v_{684} \in \mathbb{R} \to -v_{684} \leq v_{683} \to v_{683} \leq v_{684} \to \lvert v_{683} \rvert \leq v_{684}$
\item $\varphi_{17} :\equiv \forall y\,\forall z\,(v_{654} = y \to y = z \to v_{654} = z) \to \forall z\,(v_{654} = v_{655} \to v_{655} = z \to v_{654} = z)$
\item $\varphi_{18} :\equiv \forall y\,\forall z\,(v_{651} = y \to y = z \to v_{651} = z) \to \forall z\,(v_{651} = v_{652} \to v_{652} = z \to v_{651} = z)$
\item $\varphi_{19} :\equiv \forall y\,\forall z\,(v_{670} = y \to y = z \to v_{670} = z) \to \forall z\,(v_{670} = v_{671} \to v_{671} = z \to v_{670} = z)$
\item $\varphi_{20} :\equiv \forall y\,\forall z\,(v_{667} = y \to y = z \to v_{667} = z) \to \forall z\,(v_{667} = v_{668} \to v_{668} = z \to v_{667} = z)$
\item $\varphi_{21} :\equiv \forall y\,(v_{673} - y = v_{673} + -y) \to v_{673} - v_{674} = v_{673} + -v_{674}$
\item $\varphi_{22} :\equiv \forall z\,\forall w\,(z = w \to \lvert x \rvert + -\lvert y \rvert = z \to \lvert x \rvert + -\lvert y \rvert = w)$
\item $\varphi_{23} :\equiv \forall z\,\forall w\,(z = w \to \lvert y \rvert + -\lvert x \rvert = z \to \lvert y \rvert + -\lvert x \rvert = w)$
\item $\varphi_{24} :\equiv \forall y\,(v_{644} \in \mathbb{R} \to y \in \mathbb{R} \to \lvert v_{644} - y \rvert = \lvert y - v_{644} \rvert)$
\item $\varphi_{25} :\equiv \forall v_{680}\,(-(\lvert x \rvert - \lvert y \rvert) = v_{680} \to v_{680} = -(\lvert x \rvert - \lvert y \rvert))$
\item $\varphi_{26} :\equiv v_{677} \in \mathbb{R} \to v_{678} \in \mathbb{R} \to -(v_{678} - v_{677}) = v_{677} - v_{678}$
\item $\varphi_{27} :\equiv \forall v_{643}\,(v_{642} \in \mathbb{R} \to v_{643} \in \mathbb{R} \to v_{642} - v_{643} + v_{643} = v_{642})$
\item $\varphi_{28} :\equiv v_{681} \in \mathbb{R} \to v_{682} \in \mathbb{R} \to v_{681} \leq v_{682} \to -v_{682} \leq -v_{681}$
\item $\varphi_{29} :\equiv \forall v_{645}\,(y \in \mathbb{R} \to v_{645} \in \mathbb{R} \to \lvert y - v_{645} \rvert = \lvert v_{645} - y \rvert)$
\item $\varphi_{30} :\equiv z = w \to \lvert z \rvert = \lvert x - y \rvert + \lvert y \rvert \to \lvert w \rvert = \lvert x - y \rvert + \lvert y \rvert$
\item $\varphi_{31} :\equiv \forall v_{646}\,\forall v_{647}\,(v_{646} \in \mathbb{R} \to v_{647} \in \mathbb{R} \to v_{646} + v_{647} \in \mathbb{R})$
\item $\varphi_{32} :\equiv z = w \to \lvert z \rvert = \lvert y - x \rvert + \lvert x \rvert \to \lvert w \rvert = \lvert y - x \rvert + \lvert x \rvert$
\item $\varphi_{33} :\equiv \forall v_{295}\,\forall v_{296}\,(v_{295} \in \mathbb{R} \to v_{296} \in \mathbb{R} \to v_{295} - v_{296} \in \mathbb{R})$
\item $\varphi_{34} :\equiv \forall v_{662}\,\forall v_{663}\,(v_{662} \in \mathbb{R} \to v_{663} \in \mathbb{R} \to v_{662} + v_{663} \in \mathbb{R})$
\item $\varphi_{35} :\equiv \forall x\,\forall y\,(x \in \mathbb{R} \to y \in \mathbb{R} \to \lvert \lvert x \rvert - \lvert y \rvert \rvert \leq \lvert x - y \rvert)$
\item $\varphi_{36} :\equiv \forall z\,(v_{654} = v_{655} \to v_{655} = z \to v_{654} = z) \to v_{654} = v_{655} \to v_{655} = v_{656} \to v_{654} = v_{656}$
\item $\varphi_{37} :\equiv \forall z\,(v_{651} = v_{652} \to v_{652} = z \to v_{651} = z) \to v_{651} = v_{652} \to v_{652} = v_{653} \to v_{651} = v_{653}$
\item $\varphi_{38} :\equiv \forall z\,(v_{670} = v_{671} \to v_{671} = z \to v_{670} = z) \to v_{670} = v_{671} \to v_{671} = v_{672} \to v_{670} = v_{672}$
\item $\varphi_{39} :\equiv \forall z\,(v_{667} = v_{668} \to v_{668} = z \to v_{667} = z) \to v_{667} = v_{668} \to v_{668} = v_{669} \to v_{667} = v_{669}$
\item $\varphi_{40} :\equiv \forall y\,(v_{683} \in \mathbb{R} \to y \in \mathbb{R} \to -y \leq v_{683} \to v_{683} \leq y \to \lvert v_{683} \rvert \leq y)$
\item $\varphi_{41} :\equiv \forall z\,\forall w\,(z = w \to \lvert x \rvert + -\lvert y \rvert < z \to \lvert x \rvert + -\lvert y \rvert < w)$
\item $\varphi_{42} :\equiv \forall x\,\forall y\,\forall z\,(x \in \mathbb{R} \to y \in \mathbb{R} \to z \in \mathbb{R} \to x \leq y \to x + z \leq y + z)$
\item $\varphi_{43} :\equiv \forall z\,\forall w\,(z = w \to \lvert y \rvert + -\lvert x \rvert < z \to \lvert y \rvert + -\lvert x \rvert < w)$
\item $\varphi_{44} :\equiv v_{644} \in \mathbb{R} \to v_{645} \in \mathbb{R} \to \lvert v_{644} - v_{645} \rvert = \lvert v_{645} - v_{644} \rvert$
\item $\varphi_{45} :\equiv \forall w\,(\lvert y \rvert - \lvert x \rvert = w \to \lvert y \rvert - \lvert x \rvert = \lvert x - y \rvert \to w = \lvert x - y \rvert)$
\item $\varphi_{46} :\equiv \forall w\,(\lvert y - x \rvert = w \to \lvert y \rvert = \lvert y - x \rvert + \lvert x \rvert \to \lvert y \rvert = w + \lvert x \rvert)$
\item $\varphi_{47} :\equiv x \in \mathbb{R} \to y \in \mathbb{R} \to -(-(\lvert x \rvert - \lvert y \rvert)) = \lvert x \rvert - \lvert y \rvert$
\item $\varphi_{48} :\equiv \lvert x - y + y \rvert = \lvert x - y \rvert + \lvert y \rvert \to \lvert x \rvert = \lvert x - y \rvert + \lvert y \rvert$
\item $\varphi_{49} :\equiv \lvert y - x + x \rvert = \lvert y - x \rvert + \lvert x \rvert \to \lvert y \rvert = \lvert y - x \rvert + \lvert x \rvert$
\item $\varphi_{50} :\equiv v_{660} \in \mathbb{R} \to v_{661} \in \mathbb{R} \to v_{659} \leq v_{660} \to v_{659} + v_{661} \leq v_{660} + v_{661}$
\item $\varphi_{51} :\equiv z = w \to \lvert z \rvert < \lvert x - y \rvert + \lvert y \rvert \to \lvert w \rvert < \lvert x - y \rvert + \lvert y \rvert$
\item $\varphi_{52} :\equiv \forall x\,\forall y\,\forall z\,(x \in \mathbb{R} \to y \in \mathbb{R} \to z \in \mathbb{R} \to x + y + z = x + (y + z))$
\item $\varphi_{53} :\equiv z = w \to \lvert z \rvert < \lvert y - x \rvert + \lvert x \rvert \to \lvert w \rvert < \lvert y - x \rvert + \lvert x \rvert$
\item $\varphi_{54} :\equiv \lvert x \rvert - \lvert y \rvert \in \mathbb{R} \to -(-(\lvert x \rvert - \lvert y \rvert)) = \lvert x \rvert - \lvert y \rvert$
\item $\varphi_{55} :\equiv \forall b\,(x - y \in \mathbb{R} \to b \in \mathbb{R} \to \lvert x - y + b \rvert \leq \lvert x - y \rvert + \lvert b \rvert)$
\item $\varphi_{56} :\equiv \forall y\,(\lvert x \rvert - y = \lvert x \rvert + -y) \to \lvert x \rvert - \lvert y \rvert = \lvert x \rvert + -\lvert y \rvert$
\item $\varphi_{57} :\equiv \forall b\,(y - x \in \mathbb{R} \to b \in \mathbb{R} \to \lvert y - x + b \rvert \leq \lvert y - x \rvert + \lvert b \rvert)$
\item $\varphi_{58} :\equiv \forall v_{647}\,(\lvert x - y \rvert \in \mathbb{R} \to v_{647} \in \mathbb{R} \to \lvert x - y \rvert + v_{647} \in \mathbb{R})$
\item $\varphi_{59} :\equiv \lvert x - y \rvert + (\lvert y \rvert + -\lvert y \rvert) = \lvert x - y \rvert + (\lvert y \rvert + -\lvert y \rvert)$
\item $\varphi_{60} :\equiv \forall x \in \mathbb{R}\,(x + -x = 0) \to \lvert y \rvert \in \mathbb{R} \to \lvert y \rvert + -\lvert y \rvert = 0$
\item $\varphi_{61} :\equiv \forall v_{663}\,(\lvert x - y \rvert \in \mathbb{R} \to v_{663} \in \mathbb{R} \to \lvert x - y \rvert + v_{663} \in \mathbb{R})$
\item $\varphi_{62} :\equiv \forall y\,\forall z\,(v_{659} \in \mathbb{R} \to y \in \mathbb{R} \to z \in \mathbb{R} \to v_{659} \leq y \to v_{659} + z \leq y + z)$
\item $\varphi_{63} :\equiv \lvert x - y \rvert + (\lvert x \rvert + -\lvert x \rvert) = \lvert x - y \rvert + (\lvert x \rvert + -\lvert x \rvert)$
\item $\varphi_{64} :\equiv \forall x \in \mathbb{R}\,(x + -x = 0) \to \lvert x \rvert \in \mathbb{R} \to \lvert x \rvert + -\lvert x \rvert = 0$
\item $\varphi_{65} :\equiv \lvert x - y \rvert + \lvert y \rvert + -\lvert y \rvert = \lvert x - y \rvert + (\lvert y \rvert + -\lvert y \rvert)$
\item $\varphi_{66} :\equiv \lvert x - y \rvert + \lvert x \rvert + -\lvert x \rvert = \lvert x - y \rvert + (\lvert x \rvert + -\lvert x \rvert)$
\item $\varphi_{67} :\equiv \forall w\,(\lvert y \rvert - \lvert x \rvert = w \to \lvert y \rvert - \lvert x \rvert < \lvert x - y \rvert \to w < \lvert x - y \rvert)$
\item $\varphi_{68} :\equiv \forall w\,(\lvert y - x \rvert = w \to \lvert y \rvert < \lvert y - x \rvert + \lvert x \rvert \to \lvert y \rvert < w + \lvert x \rvert)$
\item $\varphi_{69} :\equiv \forall v_{678}\,(\lvert y \rvert \in \mathbb{R} \to v_{678} \in \mathbb{R} \to -(v_{678} - \lvert y \rvert) = \lvert y \rvert - v_{678})$
\item $\varphi_{70} :\equiv \lvert y \rvert \in \mathbb{R} \to \lvert x \rvert \in \mathbb{R} \to -(\lvert x \rvert - \lvert y \rvert) = \lvert y \rvert - \lvert x \rvert$
\item $\varphi_{71} :\equiv \lvert x - y + y \rvert < \lvert x - y \rvert + \lvert y \rvert \to \lvert x \rvert < \lvert x - y \rvert + \lvert y \rvert$
\item $\varphi_{72} :\equiv \lvert y - x + x \rvert < \lvert y - x \rvert + \lvert x \rvert \to \lvert y \rvert < \lvert y - x \rvert + \lvert x \rvert$
\item $\varphi_{73} :\equiv \forall y\,\forall z\,(v_{648} \in \mathbb{R} \to y \in \mathbb{R} \to z \in \mathbb{R} \to v_{648} + y + z = v_{648} + (y + z))$
\item $\varphi_{74} :\equiv \forall y\,\forall z\,(v_{664} \in \mathbb{R} \to y \in \mathbb{R} \to z \in \mathbb{R} \to v_{664} + y + z = v_{664} + (y + z))$
\item $\varphi_{75} :\equiv \forall w\,(\lvert x \rvert + -\lvert y \rvert = w \to \lvert x \rvert + -\lvert y \rvert = \lvert x - y \rvert \to w = \lvert x - y \rvert)$
\item $\varphi_{76} :\equiv \forall x \in \mathbb{R}\,(x + 0 = x) \to \lvert x - y \rvert \in \mathbb{R} \to \lvert x - y \rvert + 0 = \lvert x - y \rvert$
\item $\varphi_{77} :\equiv \forall w\,(\lvert y \rvert + -\lvert x \rvert = w \to \lvert y \rvert + -\lvert x \rvert = \lvert x - y \rvert \to w = \lvert x - y \rvert)$
\item $\varphi_{78} :\equiv (x \in \mathbb{R} \to y \in \mathbb{R} \to x \in \mathbb{R}) \to x \in \mathbb{R} \to y \in \mathbb{R} \to y - x + x = y$
\item $\varphi_{79} :\equiv \forall y\,\forall z\,(\lvert x \rvert \in \mathbb{R} \to y \in \mathbb{R} \to z \in \mathbb{R} \to \lvert x \rvert \leq y \to \lvert x \rvert + z \leq y + z)$
\item $\varphi_{80} :\equiv x \in \mathbb{R} \to y \in \mathbb{R} \to \lvert x - y \rvert + (\lvert y \rvert + -\lvert y \rvert) = \lvert x - y \rvert + 0$
\item $\varphi_{81} :\equiv x \in \mathbb{R} \to y \in \mathbb{R} \to \lvert x - y \rvert + \lvert y \rvert + -\lvert y \rvert = \lvert x - y \rvert + 0$
\item $\varphi_{82} :\equiv \lvert x \rvert - \lvert y \rvert = \lvert x \rvert + -\lvert y \rvert \to \lvert x \rvert + -\lvert y \rvert = \lvert x \rvert - \lvert y \rvert$
\item $\varphi_{83} :\equiv \forall z\,(v_{659} \in \mathbb{R} \to v_{660} \in \mathbb{R} \to z \in \mathbb{R} \to v_{659} \leq v_{660} \to v_{659} + z \leq v_{660} + z)$
\item $\varphi_{84} :\equiv x \in \mathbb{R} \to y \in \mathbb{R} \to \lvert x - y \rvert + (\lvert x \rvert + -\lvert x \rvert) = \lvert x - y \rvert + 0$
\item $\varphi_{85} :\equiv x \in \mathbb{R} \to y \in \mathbb{R} \to \lvert x - y \rvert + \lvert x \rvert + -\lvert x \rvert = \lvert x - y \rvert + 0$
\item $\varphi_{86} :\equiv \lvert y \rvert - \lvert x \rvert = \lvert y \rvert + -\lvert x \rvert \to \lvert y \rvert + -\lvert x \rvert = \lvert y \rvert - \lvert x \rvert$
\item $\varphi_{87} :\equiv -(\lvert x \rvert - \lvert y \rvert) = \lvert y \rvert - \lvert x \rvert \to \lvert y \rvert - \lvert x \rvert = -(\lvert x \rvert - \lvert y \rvert)$
\item $\varphi_{88} :\equiv \lvert x - y \rvert + 0 = \lvert x - y \rvert \to \lvert x - y \rvert + \lvert y \rvert + -\lvert y \rvert = \lvert x - y \rvert$
\item $\varphi_{89} :\equiv -\lvert x - y \rvert = -(-(\lvert x \rvert - \lvert y \rvert)) \to -\lvert x - y \rvert = \lvert x \rvert - \lvert y \rvert$
\item $\varphi_{90} :\equiv \lvert x - y \rvert + 0 = \lvert x - y \rvert \to \lvert x - y \rvert + \lvert x \rvert + -\lvert x \rvert = \lvert x - y \rvert$
\item $\varphi_{91} :\equiv -(\lvert x \rvert - \lvert y \rvert) \leq \lvert x - y \rvert \to -\lvert x - y \rvert \leq -(-(\lvert x \rvert - \lvert y \rvert))$
\item $\varphi_{92} :\equiv (x \in \mathbb{R} \to y \in \mathbb{R} \to y \in \mathbb{R}) \to x \in \mathbb{R} \to y \in \mathbb{R} \to \lvert y \rvert \in \mathbb{R}$
\item $\varphi_{93} :\equiv (x \in \mathbb{R} \to y \in \mathbb{R} \to x \in \mathbb{R}) \to x \in \mathbb{R} \to y \in \mathbb{R} \to \lvert x \rvert \in \mathbb{R}$
\item $\varphi_{94} :\equiv (x \in \mathbb{R} \to y \in \mathbb{R} \to x \in \mathbb{R}) \to x \in \mathbb{R} \to y \in \mathbb{R} \to y - x \in \mathbb{R}$
\item $\varphi_{95} :\equiv \forall v_{674}\,(\lvert y \rvert - v_{674} = \lvert y \rvert + -v_{674}) \to \lvert y \rvert - \lvert x \rvert = \lvert y \rvert + -\lvert x \rvert$
\item $\varphi_{96} :\equiv \forall z\,(v_{648} \in \mathbb{R} \to v_{649} \in \mathbb{R} \to z \in \mathbb{R} \to v_{648} + v_{649} + z = v_{648} + (v_{649} + z))$
\item $\varphi_{97} :\equiv \forall z\,(v_{664} \in \mathbb{R} \to v_{665} \in \mathbb{R} \to z \in \mathbb{R} \to v_{664} + v_{665} + z = v_{664} + (v_{665} + z))$
\item $\varphi_{98} :\equiv \forall w\,(\lvert x \rvert + -\lvert y \rvert = w \to \lvert x \rvert + -\lvert y \rvert < \lvert x - y \rvert \to w < \lvert x - y \rvert)$
\item $\varphi_{99} :\equiv \forall y\,(x \in \mathbb{R} \to y \in \mathbb{R} \to x - y + y = x) \to x \in \mathbb{R} \to y \in \mathbb{R} \to x - y + y = x$
\item $\varphi_{100} :\equiv x \in \mathbb{R} \to y \in \mathbb{R} \to \lvert x \rvert + -\lvert y \rvert \leq \lvert x - y \rvert + \lvert y \rvert + -\lvert y \rvert$
\item $\varphi_{101} :\equiv \forall w\,(\lvert y \rvert + -\lvert x \rvert = w \to \lvert y \rvert + -\lvert x \rvert < \lvert x - y \rvert \to w < \lvert x - y \rvert)$
\item $\varphi_{102} :\equiv x \in \mathbb{R} \to y \in \mathbb{R} \to \lvert y \rvert + -\lvert x \rvert \leq \lvert x - y \rvert + \lvert x \rvert + -\lvert x \rvert$
\item $\varphi_{103} :\equiv \lvert y - x \rvert = \lvert x - y \rvert \to \lvert y \rvert = \lvert y - x \rvert + \lvert x \rvert \to \lvert y \rvert = \lvert x - y \rvert + \lvert x \rvert$
\item $\varphi_{104} :\equiv \forall v_{673}\,\forall v_{674}\,(v_{673} - v_{674} = v_{673} + -v_{674}) \to \forall v_{674}\,(\lvert y \rvert - v_{674} = \lvert y \rvert + -v_{674})$
\item $\varphi_{105} :\equiv -\lvert x - y \rvert < -(-(\lvert x \rvert - \lvert y \rvert)) \to -\lvert x - y \rvert < \lvert x \rvert - \lvert y \rvert$
\item $\varphi_{106} :\equiv \forall v_{657}\,\forall v_{658}\,(v_{657} = v_{658} \to v_{658} = v_{657}) \to \forall v_{658}\,(\lvert x \rvert - \lvert y \rvert = v_{658} \to v_{658} = \lvert x \rvert - \lvert y \rvert)$
\item $\varphi_{107} :\equiv \forall v_{675}\,\forall v_{676}\,(v_{675} = v_{676} \to v_{676} = v_{675}) \to \forall v_{676}\,(\lvert y \rvert - \lvert x \rvert = v_{676} \to v_{676} = \lvert y \rvert - \lvert x \rvert)$
\item $\varphi_{108} :\equiv \forall w\,(x - y + y = w \to \lvert x - y + y \rvert = \lvert x - y \rvert + \lvert y \rvert \to \lvert w \rvert = \lvert x - y \rvert + \lvert y \rvert)$
\item $\varphi_{109} :\equiv \forall y\,(x \in \mathbb{R} \to y \in \mathbb{R} \to x - y \in \mathbb{R}) \to x \in \mathbb{R} \to y \in \mathbb{R} \to x - y \in \mathbb{R}$
\item $\varphi_{110} :\equiv \forall w\,(y - x + x = w \to \lvert y - x + x \rvert = \lvert y - x \rvert + \lvert x \rvert \to \lvert w \rvert = \lvert y - x \rvert + \lvert x \rvert)$
\item $\varphi_{111} :\equiv (x \in \mathbb{R} \to y \in \mathbb{R} \to x \in \mathbb{R}) \to x \in \mathbb{R} \to y \in \mathbb{R} \to \lvert y - x \rvert = \lvert x - y \rvert$
\item $\varphi_{112} :\equiv \lvert x \rvert \leq \lvert x - y \rvert + \lvert y \rvert \to \lvert x \rvert + -\lvert y \rvert \leq \lvert x - y \rvert + \lvert y \rvert + -\lvert y \rvert$
\item $\varphi_{113} :\equiv \lvert y \rvert \leq \lvert x - y \rvert + \lvert x \rvert \to \lvert y \rvert + -\lvert x \rvert \leq \lvert x - y \rvert + \lvert x \rvert + -\lvert x \rvert$
\item $\varphi_{114} :\equiv \forall x\,\forall y\,(x \in \mathbb{R} \to y \in \mathbb{R} \to x - y + y = x) \to \forall y\,(x \in \mathbb{R} \to y \in \mathbb{R} \to x - y + y = x)$
\item $\varphi_{115} :\equiv (x \in \mathbb{R} \to y \in \mathbb{R} \to \lvert y \rvert \in \mathbb{R}) \to x \in \mathbb{R} \to y \in \mathbb{R} \to -\lvert y \rvert \in \mathbb{R}$
\item $\varphi_{116} :\equiv (x \in \mathbb{R} \to y \in \mathbb{R} \to x - y \in \mathbb{R}) \to x \in \mathbb{R} \to y \in \mathbb{R} \to \lvert x - y \rvert \in \mathbb{R}$
\item $\varphi_{117} :\equiv (x \in \mathbb{R} \to y \in \mathbb{R} \to \lvert x \rvert \in \mathbb{R}) \to x \in \mathbb{R} \to y \in \mathbb{R} \to -\lvert x \rvert \in \mathbb{R}$
\item $\varphi_{118} :\equiv \lvert y - x \rvert = \lvert x - y \rvert \to \lvert y \rvert < \lvert y - x \rvert + \lvert x \rvert \to \lvert y \rvert < \lvert x - y \rvert + \lvert x \rvert$
\item $\varphi_{119} :\equiv \forall x \in \mathbb{R}\,(-x \in \mathbb{R}) \to \lvert x \rvert - \lvert y \rvert \in \mathbb{R} \to -(\lvert x \rvert - \lvert y \rvert) \in \mathbb{R}$
\item $\varphi_{120} :\equiv (x \in \mathbb{R} \to y \in \mathbb{R} \to \lvert y \rvert \in \mathbb{R}) \to x \in \mathbb{R} \to y \in \mathbb{R} \to \lvert y \rvert + -\lvert y \rvert = 0$
\item $\varphi_{121} :\equiv (x \in \mathbb{R} \to y \in \mathbb{R} \to \lvert x \rvert \in \mathbb{R}) \to x \in \mathbb{R} \to y \in \mathbb{R} \to \lvert x \rvert + -\lvert x \rvert = 0$
\item $\varphi_{122} :\equiv \forall w\,(x - y + y = w \to \lvert x - y + y \rvert < \lvert x - y \rvert + \lvert y \rvert \to \lvert w \rvert < \lvert x - y \rvert + \lvert y \rvert)$
\item $\varphi_{123} :\equiv \forall x\,\forall y\,(x \in \mathbb{R} \to y \in \mathbb{R} \to x - y \in \mathbb{R}) \to \forall y\,(x \in \mathbb{R} \to y \in \mathbb{R} \to x - y \in \mathbb{R})$
\item $\varphi_{124} :\equiv \forall w\,(y - x + x = w \to \lvert y - x + x \rvert < \lvert y - x \rvert + \lvert x \rvert \to \lvert w \rvert < \lvert y - x \rvert + \lvert x \rvert)$
\item $\varphi_{125} :\equiv \lvert x \rvert + -\lvert y \rvert = \lvert x - y \rvert + \lvert y \rvert + -\lvert y \rvert \to \lvert x \rvert + -\lvert y \rvert = \lvert x - y \rvert$
\item $\varphi_{126} :\equiv \lvert y \rvert + -\lvert x \rvert = \lvert x - y \rvert + \lvert x \rvert + -\lvert x \rvert \to \lvert y \rvert + -\lvert x \rvert = \lvert x - y \rvert$
\item $\varphi_{127} :\equiv (x \in \mathbb{R} \to y \in \mathbb{R} \to \lvert y \rvert \in \mathbb{R}) \to x \in \mathbb{R} \to y \in \mathbb{R} \to \lvert x \rvert - \lvert y \rvert \in \mathbb{R}$
\item $\varphi_{128} :\equiv \forall v_{643}\,(y \in \mathbb{R} \to v_{643} \in \mathbb{R} \to y - v_{643} + v_{643} = y) \to y \in \mathbb{R} \to x \in \mathbb{R} \to y - x + x = y$
\item $\varphi_{129} :\equiv \forall v_{296}\,(y \in \mathbb{R} \to v_{296} \in \mathbb{R} \to y - v_{296} \in \mathbb{R}) \to y \in \mathbb{R} \to x \in \mathbb{R} \to y - x \in \mathbb{R}$
\item $\varphi_{130} :\equiv \forall w\,(-(-(\lvert x \rvert - \lvert y \rvert)) = w \to -\lvert x - y \rvert = -(-(\lvert x \rvert - \lvert y \rvert)) \to -\lvert x - y \rvert = w)$
\item $\varphi_{131} :\equiv \forall x\,\forall y\,(x \in \mathbb{R} \to y \in \mathbb{R} \to x - y + y = x) \to \forall y\,(v_{642} \in \mathbb{R} \to y \in \mathbb{R} \to v_{642} - y + y = v_{642})$
\item $\varphi_{132} :\equiv \forall x\,\forall y\,(x \in \mathbb{R} \to y \in \mathbb{R} \to x + y \in \mathbb{R}) \to \forall y\,(v_{646} \in \mathbb{R} \to y \in \mathbb{R} \to v_{646} + y \in \mathbb{R})$
\item $\varphi_{133} :\equiv (x \in \mathbb{R} \to y \in \mathbb{R} \to \lvert y \rvert \in \mathbb{R}) \to x \in \mathbb{R} \to y \in \mathbb{R} \to \lvert x - y \rvert + \lvert y \rvert \in \mathbb{R}$
\item $\varphi_{134} :\equiv \forall x\,\forall y\,(x \in \mathbb{R} \to y \in \mathbb{R} \to x - y \in \mathbb{R}) \to \forall y\,(v_{295} \in \mathbb{R} \to y \in \mathbb{R} \to v_{295} - y \in \mathbb{R})$
\item $\varphi_{135} :\equiv \forall x\,\forall y\,(x \in \mathbb{R} \to y \in \mathbb{R} \to x + y \in \mathbb{R}) \to \forall y\,(v_{662} \in \mathbb{R} \to y \in \mathbb{R} \to v_{662} + y \in \mathbb{R})$
\item $\varphi_{136} :\equiv (x \in \mathbb{R} \to y \in \mathbb{R} \to \lvert x \rvert \in \mathbb{R}) \to x \in \mathbb{R} \to y \in \mathbb{R} \to \lvert x - y \rvert + \lvert x \rvert \in \mathbb{R}$
\item $\varphi_{137} :\equiv \lvert x \rvert + -\lvert y \rvert < \lvert x - y \rvert + \lvert y \rvert + -\lvert y \rvert \to \lvert x \rvert + -\lvert y \rvert < \lvert x - y \rvert$
\item $\varphi_{138} :\equiv \lvert y \rvert + -\lvert x \rvert < \lvert x - y \rvert + \lvert x \rvert + -\lvert x \rvert \to \lvert y \rvert + -\lvert x \rvert < \lvert x - y \rvert$
\item $\varphi_{139} :\equiv -\lvert x - y \rvert \leq \lvert x \rvert - \lvert y \rvert \to \lvert x \rvert - \lvert y \rvert \leq \lvert x - y \rvert \to \lvert \lvert x \rvert - \lvert y \rvert \rvert \leq \lvert x - y \rvert$
\item $\varphi_{140} :\equiv (x \in \mathbb{R} \to y \in \mathbb{R} \to y \in \mathbb{R}) \to x \in \mathbb{R} \to y \in \mathbb{R} \to \lvert x - y + y \rvert \leq \lvert x - y \rvert + \lvert y \rvert$
\item $\varphi_{141} :\equiv (x \in \mathbb{R} \to y \in \mathbb{R} \to \lvert x - y \rvert \in \mathbb{R}) \to x \in \mathbb{R} \to y \in \mathbb{R} \to \lvert x - y \rvert + 0 = \lvert x - y \rvert$
\item $\varphi_{142} :\equiv (x \in \mathbb{R} \to y \in \mathbb{R} \to x \in \mathbb{R}) \to x \in \mathbb{R} \to y \in \mathbb{R} \to \lvert y - x + x \rvert \leq \lvert y - x \rvert + \lvert x \rvert$
\item $\varphi_{143} :\equiv \forall v_{660}\,\forall v_{661}\,(\lvert y \rvert \in \mathbb{R} \to v_{660} \in \mathbb{R} \to v_{661} \in \mathbb{R} \to \lvert y \rvert \leq v_{660} \to \lvert y \rvert + v_{661} \leq v_{660} + v_{661})$
\item $\varphi_{144} :\equiv \forall w\,(-(-(\lvert x \rvert - \lvert y \rvert)) = w \to -\lvert x - y \rvert < -(-(\lvert x \rvert - \lvert y \rvert)) \to -\lvert x - y \rvert < w)$
\item $\varphi_{145} :\equiv \forall x\,\forall y\,(x \in \mathbb{R} \to y \in \mathbb{R} \to x - y \in \mathbb{R}) \to \forall y\,(\lvert x \rvert \in \mathbb{R} \to y \in \mathbb{R} \to \lvert x \rvert - y \in \mathbb{R})$
\item $\varphi_{146} :\equiv \lvert x \rvert + -\lvert y \rvert = \lvert x \rvert - \lvert y \rvert \to \lvert x \rvert + -\lvert y \rvert = \lvert x - y \rvert \to \lvert x \rvert - \lvert y \rvert = \lvert x - y \rvert$
\item $\varphi_{147} :\equiv \lvert y \rvert - \lvert x \rvert = -(\lvert x \rvert - \lvert y \rvert) \to \lvert y \rvert - \lvert x \rvert = \lvert x - y \rvert \to -(\lvert x \rvert - \lvert y \rvert) = \lvert x - y \rvert$
\item $\varphi_{148} :\equiv \lvert y \rvert + -\lvert x \rvert = \lvert y \rvert - \lvert x \rvert \to \lvert y \rvert + -\lvert x \rvert = \lvert x - y \rvert \to \lvert y \rvert - \lvert x \rvert = \lvert x - y \rvert$
\item $\varphi_{149} :\equiv \forall y\,(v_{646} \in \mathbb{R} \to y \in \mathbb{R} \to v_{646} + y \in \mathbb{R}) \to v_{646} \in \mathbb{R} \to v_{647} \in \mathbb{R} \to v_{646} + v_{647} \in \mathbb{R}$
\item $\varphi_{150} :\equiv \forall y\,(v_{295} \in \mathbb{R} \to y \in \mathbb{R} \to v_{295} - y \in \mathbb{R}) \to v_{295} \in \mathbb{R} \to v_{296} \in \mathbb{R} \to v_{295} - v_{296} \in \mathbb{R}$
\item $\varphi_{151} :\equiv \forall y\,(v_{662} \in \mathbb{R} \to y \in \mathbb{R} \to v_{662} + y \in \mathbb{R}) \to v_{662} \in \mathbb{R} \to v_{663} \in \mathbb{R} \to v_{662} + v_{663} \in \mathbb{R}$
\item $\varphi_{152} :\equiv (x \in \mathbb{R} \to y \in \mathbb{R} \to \lvert x \rvert \in \mathbb{R}) \to x \in \mathbb{R} \to y \in \mathbb{R} \to -(\lvert x \rvert - \lvert y \rvert) = \lvert y \rvert - \lvert x \rvert$
\item $\varphi_{153} :\equiv \forall v_{655}\,\forall v_{656}\,(\lvert x - y \rvert + \lvert y \rvert + -\lvert y \rvert = v_{655} \to v_{655} = v_{656} \to \lvert x - y \rvert + \lvert y \rvert + -\lvert y \rvert = v_{656})$
\item $\varphi_{154} :\equiv \forall v_{652}\,\forall v_{653}\,(\lvert x - y \rvert + \lvert y \rvert + -\lvert y \rvert = v_{652} \to v_{652} = v_{653} \to \lvert x - y \rvert + \lvert y \rvert + -\lvert y \rvert = v_{653})$
\item $\varphi_{155} :\equiv \forall v_{671}\,\forall v_{672}\,(\lvert x - y \rvert + \lvert x \rvert + -\lvert x \rvert = v_{671} \to v_{671} = v_{672} \to \lvert x - y \rvert + \lvert x \rvert + -\lvert x \rvert = v_{672})$
\item $\varphi_{156} :\equiv \forall v_{668}\,\forall v_{669}\,(\lvert x - y \rvert + \lvert x \rvert + -\lvert x \rvert = v_{668} \to v_{668} = v_{669} \to \lvert x - y \rvert + \lvert x \rvert + -\lvert x \rvert = v_{669})$
\item $\varphi_{157} :\equiv \forall x\,\forall y\,(x \in \mathbb{R} \to y \in \mathbb{R} \to -(y - x) = x - y) \to \forall y\,(v_{677} \in \mathbb{R} \to y \in \mathbb{R} \to -(y - v_{677}) = v_{677} - y)$
\item $\varphi_{158} :\equiv \lvert x - y \rvert + (\lvert y \rvert + -\lvert y \rvert) = \lvert x - y \rvert + 0 \to \lvert x - y \rvert + \lvert y \rvert + -\lvert y \rvert = \lvert x - y \rvert + 0$
\item $\varphi_{159} :\equiv \lvert x - y \rvert + (\lvert y \rvert + -\lvert y \rvert) = \lvert x - y \rvert + z \to \lvert x - y \rvert + (\lvert y \rvert + -\lvert y \rvert) = \lvert x - y \rvert + w$
\item $\varphi_{160} :\equiv \lvert x - y \rvert + (\lvert x \rvert + -\lvert x \rvert) = \lvert x - y \rvert + 0 \to \lvert x - y \rvert + \lvert x \rvert + -\lvert x \rvert = \lvert x - y \rvert + 0$
\item $\varphi_{161} :\equiv \lvert x - y \rvert + (\lvert x \rvert + -\lvert x \rvert) = \lvert x - y \rvert + z \to \lvert x - y \rvert + (\lvert x \rvert + -\lvert x \rvert) = \lvert x - y \rvert + w$
\item $\varphi_{162} :\equiv (x \in \mathbb{R} \to y \in \mathbb{R} \to \lvert x \rvert - \lvert y \rvert \in \mathbb{R}) \to x \in \mathbb{R} \to y \in \mathbb{R} \to -(\lvert x \rvert - \lvert y \rvert) \in \mathbb{R}$
\item $\varphi_{163} :\equiv \lvert x \rvert + -\lvert y \rvert = \lvert x \rvert - \lvert y \rvert \to \lvert x \rvert + -\lvert y \rvert < \lvert x - y \rvert \to \lvert x \rvert - \lvert y \rvert < \lvert x - y \rvert$
\item $\varphi_{164} :\equiv \lvert y \rvert - \lvert x \rvert = -(\lvert x \rvert - \lvert y \rvert) \to \lvert y \rvert - \lvert x \rvert < \lvert x - y \rvert \to -(\lvert x \rvert - \lvert y \rvert) < \lvert x - y \rvert$
\item $\varphi_{165} :\equiv \lvert y \rvert + -\lvert x \rvert = \lvert y \rvert - \lvert x \rvert \to \lvert y \rvert + -\lvert x \rvert < \lvert x - y \rvert \to \lvert y \rvert - \lvert x \rvert < \lvert x - y \rvert$
\item $\varphi_{166} :\equiv \forall y\,(v_{642} \in \mathbb{R} \to y \in \mathbb{R} \to v_{642} - y + y = v_{642}) \to v_{642} \in \mathbb{R} \to v_{643} \in \mathbb{R} \to v_{642} - v_{643} + v_{643} = v_{642}$
\item $\varphi_{167} :\equiv \forall x\,\forall y\,(x \in \mathbb{R} \to y \in \mathbb{R} \to x \leq y \to -y \leq -x) \to \forall y\,(v_{681} \in \mathbb{R} \to y \in \mathbb{R} \to v_{681} \leq y \to -y \leq -v_{681})$
\item $\varphi_{168} :\equiv (x \in \mathbb{R} \to y \in \mathbb{R} \to \lvert x \rvert - \lvert y \rvert \leq \lvert x - y \rvert) \to x \in \mathbb{R} \to y \in \mathbb{R} \to \lvert \lvert x \rvert - \lvert y \rvert \rvert \leq \lvert x - y \rvert$
\item $\varphi_{169} :\equiv \forall y\,(\lvert x \rvert \in \mathbb{R} \to y \in \mathbb{R} \to \lvert x \rvert - y \in \mathbb{R}) \to \lvert x \rvert \in \mathbb{R} \to \lvert y \rvert \in \mathbb{R} \to \lvert x \rvert - \lvert y \rvert \in \mathbb{R}$
\item $\varphi_{170} :\equiv \forall v_{682}\,(-(\lvert x \rvert - \lvert y \rvert) \in \mathbb{R} \to v_{682} \in \mathbb{R} \to -(\lvert x \rvert - \lvert y \rvert) \leq v_{682} \to -v_{682} \leq -(-(\lvert x \rvert - \lvert y \rvert)))$
\item $\varphi_{171} :\equiv \forall y\,(v_{677} \in \mathbb{R} \to y \in \mathbb{R} \to -(y - v_{677}) = v_{677} - y) \to \varphi_{26}$
\item $\varphi_{172} :\equiv \forall v_{649}\,\forall v_{650}\,(\lvert x - y \rvert \in \mathbb{R} \to v_{649} \in \mathbb{R} \to v_{650} \in \mathbb{R} \to \lvert x - y \rvert + v_{649} + v_{650} = \lvert x - y \rvert + (v_{649} + v_{650}))$
\item $\varphi_{173} :\equiv \forall v_{650}\,(\lvert x - y \rvert \in \mathbb{R} \to \lvert y \rvert \in \mathbb{R} \to v_{650} \in \mathbb{R} \to \lvert x - y \rvert + \lvert y \rvert + v_{650} = \lvert x - y \rvert + (\lvert y \rvert + v_{650}))$
\item $\varphi_{174} :\equiv \forall z\,\forall w\,(z = w \to \lvert y \rvert = z + \lvert x \rvert \to \lvert y \rvert = w + \lvert x \rvert) \to \varphi_{46}$
\item $\varphi_{175} :\equiv \forall v_{665}\,\forall v_{666}\,(\lvert x - y \rvert \in \mathbb{R} \to v_{665} \in \mathbb{R} \to v_{666} \in \mathbb{R} \to \lvert x - y \rvert + v_{665} + v_{666} = \lvert x - y \rvert + (v_{665} + v_{666}))$
\item $\varphi_{176} :\equiv \forall v_{666}\,(\lvert x - y \rvert \in \mathbb{R} \to \lvert x \rvert \in \mathbb{R} \to v_{666} \in \mathbb{R} \to \lvert x - y \rvert + \lvert x \rvert + v_{666} = \lvert x - y \rvert + (\lvert x \rvert + v_{666}))$
\item $\varphi_{177} :\equiv \forall y\,(v_{681} \in \mathbb{R} \to y \in \mathbb{R} \to v_{681} \leq y \to -y \leq -v_{681}) \to \varphi_{28}$
\item $\varphi_{178} :\equiv \forall v_{642}\,(\varphi_{27}) \to \forall v_{643}\,(y \in \mathbb{R} \to v_{643} \in \mathbb{R} \to y - v_{643} + v_{643} = y)$
\item $\varphi_{179} :\equiv \forall x\,\forall y\,(x \in \mathbb{R} \to y \in \mathbb{R} \to -y \leq x \to x \leq y \to \lvert x \rvert \leq y) \to \varphi_{40}$
\item $\varphi_{180} :\equiv \forall v_{656}\,(\lvert x - y \rvert + \lvert y \rvert + -\lvert y \rvert = \lvert x - y \rvert + 0 \to \lvert x - y \rvert + 0 = v_{656} \to \lvert x - y \rvert + \lvert y \rvert + -\lvert y \rvert = v_{656})$
\item $\varphi_{181} :\equiv \forall v_{672}\,(\lvert x - y \rvert + \lvert x \rvert + -\lvert x \rvert = \lvert x - y \rvert + 0 \to \lvert x - y \rvert + 0 = v_{672} \to \lvert x - y \rvert + \lvert x \rvert + -\lvert x \rvert = v_{672})$
\item $\varphi_{182} :\equiv \forall w\,(\lvert x - y \rvert + \lvert y \rvert + -\lvert y \rvert = w \to \lvert x \rvert + -\lvert y \rvert = \lvert x - y \rvert + \lvert y \rvert + -\lvert y \rvert \to \lvert x \rvert + -\lvert y \rvert = w)$
\item $\varphi_{183} :\equiv \forall w\,(\lvert x - y \rvert + \lvert x \rvert + -\lvert x \rvert = w \to \lvert y \rvert + -\lvert x \rvert = \lvert x - y \rvert + \lvert x \rvert + -\lvert x \rvert \to \lvert y \rvert + -\lvert x \rvert = w)$
\item $\varphi_{184} :\equiv \forall v_{684}\,(\lvert x \rvert - \lvert y \rvert \in \mathbb{R} \to v_{684} \in \mathbb{R} \to -v_{684} \leq \lvert x \rvert - \lvert y \rvert \to \lvert x \rvert - \lvert y \rvert \leq v_{684} \to \lvert \lvert x \rvert - \lvert y \rvert \rvert \leq v_{684})$
\item $\varphi_{185} :\equiv (x \in \mathbb{R} \to y \in \mathbb{R} \to -(\lvert x \rvert - \lvert y \rvert) \leq \lvert x - y \rvert) \to x \in \mathbb{R} \to y \in \mathbb{R} \to -\lvert x - y \rvert \leq -(-(\lvert x \rvert - \lvert y \rvert))$
\item $\varphi_{186} :\equiv \forall v_{658}\,(\lvert x \rvert - \lvert y \rvert = v_{658} \to v_{658} = \lvert x \rvert - \lvert y \rvert) \to \varphi_{82}$
\item $\varphi_{187} :\equiv \forall v_{676}\,(\lvert y \rvert - \lvert x \rvert = v_{676} \to v_{676} = \lvert y \rvert - \lvert x \rvert) \to \varphi_{86}$
\item $\varphi_{188} :\equiv \forall w\,(\lvert x - y \rvert + \lvert y \rvert + -\lvert y \rvert = w \to \lvert x \rvert + -\lvert y \rvert < \lvert x - y \rvert + \lvert y \rvert + -\lvert y \rvert \to \lvert x \rvert + -\lvert y \rvert < w)$
\item $\varphi_{189} :\equiv \forall w\,(\lvert x - y \rvert + \lvert x \rvert + -\lvert x \rvert = w \to \lvert y \rvert + -\lvert x \rvert < \lvert x - y \rvert + \lvert x \rvert + -\lvert x \rvert \to \lvert y \rvert + -\lvert x \rvert < w)$
\item $\varphi_{190} :\equiv \forall z\,\forall w\,(z = w \to \lvert y \rvert < z + \lvert x \rvert \to \lvert y \rvert < w + \lvert x \rvert) \to \varphi_{68}$
\item $\varphi_{191} :\equiv \forall a\,\forall b\,(a \in \mathbb{R} \to b \in \mathbb{R} \to \lvert a + b \rvert \leq \lvert a \rvert + \lvert b \rvert) \to \varphi_{55}$
\item $\varphi_{192} :\equiv \forall a\,\forall b\,(a \in \mathbb{R} \to b \in \mathbb{R} \to \lvert a + b \rvert \leq \lvert a \rvert + \lvert b \rvert) \to \varphi_{57}$
\item $\varphi_{193} :\equiv \forall z\,(\lvert x \rvert \in \mathbb{R} \to \lvert x - y \rvert + \lvert y \rvert \in \mathbb{R} \to z \in \mathbb{R} \to \lvert x \rvert \leq \lvert x - y \rvert + \lvert y \rvert \to \lvert x \rvert + z \leq \lvert x - y \rvert + \lvert y \rvert + z)$
\item $\varphi_{194} :\equiv (x - y + y = x \to \varphi_{48}) \to x \in \mathbb{R} \to y \in \mathbb{R} \to \lvert x \rvert \leq \lvert x - y \rvert + \lvert y \rvert$
\item $\varphi_{195} :\equiv (y - x + x = y \to \varphi_{49}) \to x \in \mathbb{R} \to y \in \mathbb{R} \to \lvert y \rvert \leq \lvert y - x \rvert + \lvert x \rvert$
\item $\varphi_{196} :\equiv \varphi_{55} \to x - y \in \mathbb{R} \to y \in \mathbb{R} \to \lvert x - y + y \rvert \leq \lvert x - y \rvert + \lvert y \rvert$
\item $\varphi_{197} :\equiv \varphi_{57} \to y - x \in \mathbb{R} \to x \in \mathbb{R} \to \lvert y - x + x \rvert \leq \lvert y - x \rvert + \lvert x \rvert$
\item $\varphi_{198} :\equiv (x \in \mathbb{R} \to y \in \mathbb{R} \to -\lvert y \rvert \in \mathbb{R}) \to x \in \mathbb{R} \to y \in \mathbb{R} \to \varphi_{65}$
\item $\varphi_{199} :\equiv \forall v_{661}\,(\lvert y \rvert \in \mathbb{R} \to \lvert x - y \rvert + \lvert x \rvert \in \mathbb{R} \to v_{661} \in \mathbb{R} \to \lvert y \rvert \leq \lvert x - y \rvert + \lvert x \rvert \to \lvert y \rvert + v_{661} \leq \lvert x - y \rvert + \lvert x \rvert + v_{661})$
\item $\varphi_{200} :\equiv (x \in \mathbb{R} \to y \in \mathbb{R} \to -\lvert x \rvert \in \mathbb{R}) \to x \in \mathbb{R} \to y \in \mathbb{R} \to \varphi_{66}$
\item $\varphi_{201} :\equiv \lvert y \rvert + -\lvert y \rvert = 0 \to \varphi_{59} \to \lvert x - y \rvert + (\lvert y \rvert + -\lvert y \rvert) = \lvert x - y \rvert + 0$
\item $\varphi_{202} :\equiv \lvert x \rvert + -\lvert x \rvert = 0 \to \varphi_{63} \to \lvert x - y \rvert + (\lvert x \rvert + -\lvert x \rvert) = \lvert x - y \rvert + 0$
\item $\varphi_{203} :\equiv \forall w\,(\lvert y \rvert + -\lvert y \rvert = w \to \varphi_{59} \to \lvert x - y \rvert + (\lvert y \rvert + -\lvert y \rvert) = \lvert x - y \rvert + w)$
\item $\varphi_{204} :\equiv \forall w\,(\lvert x \rvert + -\lvert x \rvert = w \to \varphi_{63} \to \lvert x - y \rvert + (\lvert x \rvert + -\lvert x \rvert) = \lvert x - y \rvert + w)$
\item $\varphi_{205} :\equiv (y \in \mathbb{R} \to x \in \mathbb{R} \to y - x + x = y) \to (x \in \mathbb{R} \to y \in \mathbb{R} \to y \in \mathbb{R}) \to \varphi_{78}$
\item $\varphi_{206} :\equiv \lvert x \rvert - \lvert y \rvert = \lvert x \rvert + -\lvert y \rvert \to \varphi_{82} \to \lvert x \rvert + -\lvert y \rvert = \lvert x \rvert - \lvert y \rvert$
\item $\varphi_{207} :\equiv \lvert y \rvert - \lvert x \rvert = \lvert y \rvert + -\lvert x \rvert \to \varphi_{86} \to \lvert y \rvert + -\lvert x \rvert = \lvert y \rvert - \lvert x \rvert$
\item $\varphi_{208} :\equiv (y \in \mathbb{R} \to x \in \mathbb{R} \to y - x \in \mathbb{R}) \to (x \in \mathbb{R} \to y \in \mathbb{R} \to y \in \mathbb{R}) \to \varphi_{94}$
\item $\varphi_{209} :\equiv \forall v_{653}\,(\varphi_{65} \to \lvert x - y \rvert + (\lvert y \rvert + -\lvert y \rvert) = v_{653} \to \lvert x - y \rvert + \lvert y \rvert + -\lvert y \rvert = v_{653})$
\item $\varphi_{210} :\equiv \forall v_{669}\,(\varphi_{66} \to \lvert x - y \rvert + (\lvert x \rvert + -\lvert x \rvert) = v_{669} \to \lvert x - y \rvert + \lvert x \rvert + -\lvert x \rvert = v_{669})$
\item $\varphi_{211} :\equiv (y \in \mathbb{R} \to x \in \mathbb{R} \to \lvert y - x \rvert = \lvert x - y \rvert) \to (x \in \mathbb{R} \to y \in \mathbb{R} \to y \in \mathbb{R}) \to \varphi_{111}$
\item $\varphi_{212} :\equiv (-(-(\lvert x \rvert - \lvert y \rvert)) = \lvert x \rvert - \lvert y \rvert \to \varphi_{89}) \to x \in \mathbb{R} \to y \in \mathbb{R} \to -\lvert x - y \rvert \leq \lvert x \rvert - \lvert y \rvert$
\item $\varphi_{213} :\equiv (x \in \mathbb{R} \to y \in \mathbb{R} \to -\lvert y \rvert \in \mathbb{R}) \to (x \in \mathbb{R} \to y \in \mathbb{R} \to \lvert x \rvert \leq \lvert x - y \rvert + \lvert y \rvert) \to \varphi_{100}$
\item $\varphi_{214} :\equiv (x \in \mathbb{R} \to y \in \mathbb{R} \to -\lvert x \rvert \in \mathbb{R}) \to (x \in \mathbb{R} \to y \in \mathbb{R} \to \lvert y \rvert \leq \lvert x - y \rvert + \lvert x \rvert) \to \varphi_{102}$
\item $\varphi_{215} :\equiv (\lvert x \rvert \in \mathbb{R} \to \lvert y \rvert \in \mathbb{R} \to \lvert x \rvert - \lvert y \rvert \in \mathbb{R}) \to (x \in \mathbb{R} \to y \in \mathbb{R} \to \lvert x \rvert \in \mathbb{R}) \to \varphi_{127}$
\item $\varphi_{216} :\equiv (\lvert x - y \rvert + \lvert y \rvert + -\lvert y \rvert = \lvert x - y \rvert \to \varphi_{125}) \to x \in \mathbb{R} \to y \in \mathbb{R} \to \lvert x \rvert + -\lvert y \rvert \leq \lvert x - y \rvert$
\item $\varphi_{217} :\equiv (\lvert x - y \rvert + \lvert x \rvert + -\lvert x \rvert = \lvert x - y \rvert \to \varphi_{126}) \to x \in \mathbb{R} \to y \in \mathbb{R} \to \lvert y \rvert + -\lvert x \rvert \leq \lvert x - y \rvert$
\item $\varphi_{218} :\equiv (\lvert x - y \rvert + \lvert y \rvert + -\lvert y \rvert = \lvert x - y \rvert + 0 \to \varphi_{88}) \to x \in \mathbb{R} \to y \in \mathbb{R} \to \lvert x - y \rvert + \lvert y \rvert + -\lvert y \rvert = \lvert x - y \rvert$
\item $\varphi_{219} :\equiv (\lvert x - y \rvert + \lvert x \rvert + -\lvert x \rvert = \lvert x - y \rvert + 0 \to \varphi_{90}) \to x \in \mathbb{R} \to y \in \mathbb{R} \to \lvert x - y \rvert + \lvert x \rvert + -\lvert x \rvert = \lvert x - y \rvert$
\item $\varphi_{220} :\equiv (x \in \mathbb{R} \to y \in \mathbb{R} \to -(\lvert x \rvert - \lvert y \rvert) = \lvert y \rvert - \lvert x \rvert) \to \varphi_{87} \to x \in \mathbb{R} \to y \in \mathbb{R} \to \lvert y \rvert - \lvert x \rvert = -(\lvert x \rvert - \lvert y \rvert)$
\item $\varphi_{221} :\equiv (x - y \in \mathbb{R} \to y \in \mathbb{R} \to \lvert x - y + y \rvert \leq \lvert x - y \rvert + \lvert y \rvert) \to (x \in \mathbb{R} \to y \in \mathbb{R} \to x - y \in \mathbb{R}) \to \varphi_{140}$
\item $\varphi_{222} :\equiv (y - x \in \mathbb{R} \to x \in \mathbb{R} \to \lvert y - x + x \rvert \leq \lvert y - x \rvert + \lvert x \rvert) \to (x \in \mathbb{R} \to y \in \mathbb{R} \to y - x \in \mathbb{R}) \to \varphi_{142}$
\item $\varphi_{223} :\equiv (\lvert x - y \rvert \in \mathbb{R} \to \lvert y \rvert \in \mathbb{R} \to \lvert x - y \rvert + \lvert y \rvert \in \mathbb{R}) \to (x \in \mathbb{R} \to y \in \mathbb{R} \to \lvert x - y \rvert \in \mathbb{R}) \to \varphi_{133}$
\item $\varphi_{224} :\equiv (\lvert x - y \rvert \in \mathbb{R} \to \lvert x \rvert \in \mathbb{R} \to \lvert x - y \rvert + \lvert x \rvert \in \mathbb{R}) \to (x \in \mathbb{R} \to y \in \mathbb{R} \to \lvert x - y \rvert \in \mathbb{R}) \to \varphi_{136}$
\item $\varphi_{225} :\equiv (x \in \mathbb{R} \to y \in \mathbb{R} \to \lvert x - y \rvert \in \mathbb{R}) \to (x \in \mathbb{R} \to y \in \mathbb{R} \to -\lvert x - y \rvert \leq \lvert x \rvert - \lvert y \rvert) \to \varphi_{168}$
\item $\varphi_{226} :\equiv \forall v_{648}\,\forall v_{649}\,\forall v_{650}\,(v_{648} \in \mathbb{R} \to v_{649} \in \mathbb{R} \to \varphi_{13}) \to \varphi_{172}$
\item $\varphi_{227} :\equiv \forall v_{664}\,\forall v_{665}\,\forall v_{666}\,(v_{664} \in \mathbb{R} \to v_{665} \in \mathbb{R} \to \varphi_{14}) \to \varphi_{175}$
\item $\varphi_{228} :\equiv (x \in \mathbb{R} \to y \in \mathbb{R} \to -(\lvert x \rvert - \lvert y \rvert) \in \mathbb{R}) \to (x \in \mathbb{R} \to y \in \mathbb{R} \to \lvert x - y \rvert \in \mathbb{R}) \to \varphi_{185}$
\item $\varphi_{229} :\equiv (x \in \mathbb{R} \to y \in \mathbb{R} \to \lvert x - y \rvert \in \mathbb{R}) \to (x \in \mathbb{R} \to y \in \mathbb{R} \to \lvert y \rvert \in \mathbb{R}) \to \varphi_{198}$
\item $\varphi_{230} :\equiv (x \in \mathbb{R} \to y \in \mathbb{R} \to \lvert x - y \rvert \in \mathbb{R}) \to (x \in \mathbb{R} \to y \in \mathbb{R} \to \lvert x \rvert \in \mathbb{R}) \to \varphi_{200}$
\item $\varphi_{231} :\equiv (x \in \mathbb{R} \to y \in \mathbb{R} \to \lvert x \rvert \in \mathbb{R}) \to (x \in \mathbb{R} \to y \in \mathbb{R} \to \lvert x - y \rvert + \lvert y \rvert \in \mathbb{R}) \to \varphi_{213}$
\item $\varphi_{232} :\equiv (x \in \mathbb{R} \to y \in \mathbb{R} \to \lvert y \rvert \leq \lvert y - x \rvert + \lvert x \rvert) \to \varphi_{118} \to \varphi_{103} \to x \in \mathbb{R} \to y \in \mathbb{R} \to \lvert y \rvert \leq \lvert x - y \rvert + \lvert x \rvert$
\item $\varphi_{233} :\equiv (x \in \mathbb{R} \to y \in \mathbb{R} \to \lvert y \rvert \in \mathbb{R}) \to (x \in \mathbb{R} \to y \in \mathbb{R} \to \lvert x - y \rvert + \lvert x \rvert \in \mathbb{R}) \to \varphi_{214}$
\item $\varphi_{234} :\equiv \varphi_{193} \to \lvert x \rvert \in \mathbb{R} \to \lvert x - y \rvert + \lvert y \rvert \in \mathbb{R} \to -\lvert y \rvert \in \mathbb{R} \to \varphi_{112}$
\item $\varphi_{235} :\equiv (x \in \mathbb{R} \to y \in \mathbb{R} \to \lvert x - y + y \rvert \leq \lvert x - y \rvert + \lvert y \rvert) \to (x - y + y = x \to \varphi_{71}) \to \varphi_{194}$
\item $\varphi_{236} :\equiv (x \in \mathbb{R} \to y \in \mathbb{R} \to \lvert y - x + x \rvert \leq \lvert y - x \rvert + \lvert x \rvert) \to (y - x + x = y \to \varphi_{72}) \to \varphi_{195}$
\item $\varphi_{237} :\equiv \varphi_{199} \to \lvert y \rvert \in \mathbb{R} \to \lvert x - y \rvert + \lvert x \rvert \in \mathbb{R} \to -\lvert x \rvert \in \mathbb{R} \to \varphi_{113}$
\item $\varphi_{238} :\equiv (x \in \mathbb{R} \to y \in \mathbb{R} \to \lvert x \rvert + -\lvert y \rvert \leq \lvert x - y \rvert) \to \varphi_{163} \to \varphi_{146} \to x \in \mathbb{R} \to y \in \mathbb{R} \to \lvert x \rvert - \lvert y \rvert \leq \lvert x - y \rvert$
\item $\varphi_{239} :\equiv (x \in \mathbb{R} \to y \in \mathbb{R} \to \lvert y \rvert + -\lvert x \rvert \leq \lvert x - y \rvert) \to \varphi_{165} \to \varphi_{148} \to x \in \mathbb{R} \to y \in \mathbb{R} \to \lvert y \rvert - \lvert x \rvert \leq \lvert x - y \rvert$
\item $\varphi_{240} :\equiv (x \in \mathbb{R} \to y \in \mathbb{R} \to \lvert y \rvert - \lvert x \rvert \leq \lvert x - y \rvert) \to \varphi_{164} \to \varphi_{147} \to x \in \mathbb{R} \to y \in \mathbb{R} \to -(\lvert x \rvert - \lvert y \rvert) \leq \lvert x - y \rvert$
\item $\varphi_{241} :\equiv (x \in \mathbb{R} \to y \in \mathbb{R} \to -\lvert x - y \rvert \leq -(-(\lvert x \rvert - \lvert y \rvert))) \to (-(-(\lvert x \rvert - \lvert y \rvert)) = \lvert x \rvert - \lvert y \rvert \to \varphi_{105}) \to \varphi_{212}$
\item $\varphi_{242} :\equiv (x \in \mathbb{R} \to y \in \mathbb{R} \to \lvert y \rvert - \lvert x \rvert = -(\lvert x \rvert - \lvert y \rvert)) \to \varphi_{240}$
\item $\varphi_{243} :\equiv (\lvert x \rvert - \lvert y \rvert \in \mathbb{R} \to \lvert x - y \rvert \in \mathbb{R} \to \varphi_{139}) \to (x \in \mathbb{R} \to y \in \mathbb{R} \to \lvert x \rvert - \lvert y \rvert \in \mathbb{R}) \to \varphi_{225}$
\item $\varphi_{244} :\equiv (\lvert x \rvert \in \mathbb{R} \to \lvert x - y \rvert + \lvert y \rvert \in \mathbb{R} \to -\lvert y \rvert \in \mathbb{R} \to \varphi_{112}) \to \varphi_{231}$
\item $\varphi_{245} :\equiv (\lvert y \rvert \in \mathbb{R} \to \lvert x - y \rvert + \lvert x \rvert \in \mathbb{R} \to -\lvert x \rvert \in \mathbb{R} \to \varphi_{113}) \to \varphi_{233}$
\item $\varphi_{246} :\equiv (x \in \mathbb{R} \to y \in \mathbb{R} \to \lvert x - y \rvert + \lvert y \rvert + -\lvert y \rvert = \lvert x - y \rvert) \to \varphi_{100} \to (\lvert x - y \rvert + \lvert y \rvert + -\lvert y \rvert = \lvert x - y \rvert \to \varphi_{137}) \to \varphi_{216}$
\item $\varphi_{247} :\equiv (x \in \mathbb{R} \to y \in \mathbb{R} \to \lvert x - y \rvert + \lvert x \rvert + -\lvert x \rvert = \lvert x - y \rvert) \to \varphi_{102} \to (\lvert x - y \rvert + \lvert x \rvert + -\lvert x \rvert = \lvert x - y \rvert \to \varphi_{138}) \to \varphi_{217}$
\end{abbrevlist}
\end{fullproof}

\section{自然数: ℕ の定義, 帰納法の構成手続き, ω の順序と ι の基本性質}
% 自動生成 (prover/build.js): NatSub  ノード数 16
\begin{fullproof}{[N7]}{$\mathbb{N} \subseteq \omega$}{16}\phantomsection\label{tree:NatSub}
\noindent$\Pi_{1}$:
\begin{ptree}
\AXC{$\forall n\,(n \in \mathbb{N} \leftrightarrow n \in \omega \land n \neq \emptyset)$\lbl{Def$\mathbb{N}$}}
\AXC{$\varphi_{2}$\lbl{A2}}
\BIC{$u \in \mathbb{N} \leftrightarrow u \in \omega \land u \neq \emptyset$}
\AXC{$u \in \mathbb{N} \to u \in \mathbb{N}$\lbl{A1}}
\AXC{$\varphi_{4}$\lbl{A1}}
\BIC{$(u \in \mathbb{N} \leftrightarrow u \in \omega \land u \neq \emptyset) \to u \in \mathbb{N} \to u \in \omega$}
\BIC{$u \in \mathbb{N} \to u \in \omega$}
\UIC{$\forall u \in \mathbb{N}\,(u \in \omega)$}
\end{ptree}
\noindent 主木:
\begin{ptree}
\AXC{$\stackrel{\Pi_{1}}{\vdots}$}\noLine
\UIC{$\forall u \in \mathbb{N}\,(u \in \omega)$}
\AXC{$\forall x\,\forall y\,(x \subseteq y \leftrightarrow \forall u \in x\,(u \in y))$\lbl{Def$\subseteq$}}
\AXC{$\varphi_{1}$\lbl{A2}}
\BIC{$\forall y\,(\mathbb{N} \subseteq y \leftrightarrow \forall u \in \mathbb{N}\,(u \in y))$}
\AXC{$\varphi_{3}$\lbl{A2}}
\BIC{$\mathbb{N} \subseteq \omega \leftrightarrow \forall u \in \mathbb{N}\,(u \in \omega)$}
\AXC{$\varphi_{5}$\lbl{A1}}
\BIC{$\forall u \in \mathbb{N}\,(u \in \omega) \to \mathbb{N} \subseteq \omega$}
\BIC{$\mathbb{N} \subseteq \omega$}
\end{ptree}
\begin{abbrevlist}
\item $\varphi_{1} :\equiv \forall x\,\forall y\,(x \subseteq y \leftrightarrow \forall u \in x\,(u \in y)) \to \forall y\,(\mathbb{N} \subseteq y \leftrightarrow \forall u \in \mathbb{N}\,(u \in y))$
\item $\varphi_{2} :\equiv \forall n\,(n \in \mathbb{N} \leftrightarrow n \in \omega \land n \neq \emptyset) \to (u \in \mathbb{N} \leftrightarrow u \in \omega \land u \neq \emptyset)$
\item $\varphi_{3} :\equiv \forall y\,(\mathbb{N} \subseteq y \leftrightarrow \forall u \in \mathbb{N}\,(u \in y)) \to (\mathbb{N} \subseteq \omega \leftrightarrow \forall u \in \mathbb{N}\,(u \in \omega))$
\item $\varphi_{4} :\equiv (u \in \mathbb{N} \to u \in \mathbb{N}) \to (u \in \mathbb{N} \leftrightarrow u \in \omega \land u \neq \emptyset) \to u \in \mathbb{N} \to u \in \omega$
\item $\varphi_{5} :\equiv (\mathbb{N} \subseteq \omega \leftrightarrow \forall u \in \mathbb{N}\,(u \in \omega)) \to \forall u \in \mathbb{N}\,(u \in \omega) \to \mathbb{N} \subseteq \omega$
\end{abbrevlist}
\end{fullproof}

% 自動生成 (prover/build.js): SNe0  ノード数 35
\begin{fullproof}{[N0]}{$\forall n\,(S(n) \neq \emptyset)$}{35}\phantomsection\label{tree:SNe0}
\noindent$\Pi_{1}$:
\begin{ptree}
\AXC{$\forall x\,\forall u\,(u \in S(x) \leftrightarrow u \in x \lor u = x)$\lbl{Def$S$}}
\AXC{$\varphi_{4}$\lbl{A2}}
\BIC{$\forall u\,(u \in S(n) \leftrightarrow u \in n \lor u = n)$}
\AXC{$\varphi_{1}$\lbl{A2}}
\BIC{$n \in S(n) \leftrightarrow n \in n \lor n = n$}
\AXC{$\varphi_{2}$\lbl{A1}}
\BIC{$n = n \to \neg S(n) \neq \emptyset \to n \in S(n)$}
\end{ptree}
\noindent$\Pi_{2}$:
\begin{ptree}
\AXC{$x = x$\lbl{A5}}
\UIC{$\forall x\,(x = x)$}
\AXC{$\forall x\,(x = x) \to n = n$\lbl{A2}}
\BIC{$n = n$}
\AXC{$\stackrel{\Pi_{1}}{\vdots}$}\noLine
\UIC{$n = n \to \neg S(n) \neq \emptyset \to n \in S(n)$}
\BIC{$\neg S(n) \neq \emptyset \to n \in S(n)$}
\end{ptree}
\noindent$\Pi_{3}$:
\begin{ptree}
\AXC{$\stackrel{\Pi_{2}}{\vdots}$}\noLine
\UIC{$\neg S(n) \neq \emptyset \to n \in S(n)$}
\AXC{$\neg S(n) \neq \emptyset \to \neg S(n) \neq \emptyset$\lbl{A1}}
\AXC{$\varphi_{5}$\lbl{A1}}
\BIC{$\neg S(n) \neq \emptyset \to S(n) = \emptyset$}
\AXC{$\varphi_{9}$\lbl{A1}}
\BIC{$(\neg S(n) \neq \emptyset \to n \in S(n)) \to \varphi_{7}$}
\BIC{$\varphi_{7}$}
\end{ptree}
\noindent$\Pi_{4}$:
\begin{ptree}
\AXC{$z = w \to n \in z \to n \in w$\lbl{A6}}
\UIC{$\forall w\,(z = w \to n \in z \to n \in w)$}
\UIC{$\forall z\,\forall w\,(z = w \to n \in z \to n \in w)$}
\AXC{$\varphi_{3}$\lbl{A2}}
\BIC{$\forall w\,(S(n) = w \to n \in S(n) \to n \in w)$}
\AXC{$\varphi_{6}$\lbl{A2}}
\BIC{$S(n) = \emptyset \to n \in S(n) \to n \in \emptyset$}
\AXC{$\stackrel{\Pi_{3}}{\vdots}$}\noLine
\UIC{$\varphi_{7}$}
\BIC{$\neg S(n) \neq \emptyset \to n \in \emptyset$}
\AXC{$\varphi_{8}$\lbl{A1}}
\BIC{$(\neg S(n) \neq \emptyset \to n \notin \emptyset) \to S(n) \neq \emptyset$}
\end{ptree}
\noindent 主木:
\begin{ptree}
\AXC{$\forall u\,(u \notin \emptyset)$\lbl{Def$\emptyset$}}
\AXC{$\forall u\,(u \notin \emptyset) \to n \notin \emptyset$\lbl{A2}}
\BIC{$n \notin \emptyset$}
\AXC{$n \notin \emptyset \to \neg S(n) \neq \emptyset \to n \notin \emptyset$\lbl{A1}}
\BIC{$\neg S(n) \neq \emptyset \to n \notin \emptyset$}
\AXC{$\stackrel{\Pi_{4}}{\vdots}$}\noLine
\UIC{$(\neg S(n) \neq \emptyset \to n \notin \emptyset) \to S(n) \neq \emptyset$}
\BIC{$S(n) \neq \emptyset$}
\UIC{$\forall n\,(S(n) \neq \emptyset)$}
\end{ptree}
\begin{abbrevlist}
\item $\varphi_{1} :\equiv \forall u\,(u \in S(n) \leftrightarrow u \in n \lor u = n) \to (n \in S(n) \leftrightarrow n \in n \lor n = n)$
\item $\varphi_{2} :\equiv (n \in S(n) \leftrightarrow n \in n \lor n = n) \to n = n \to \neg S(n) \neq \emptyset \to n \in S(n)$
\item $\varphi_{3} :\equiv \forall z\,\forall w\,(z = w \to n \in z \to n \in w) \to \forall w\,(S(n) = w \to n \in S(n) \to n \in w)$
\item $\varphi_{4} :\equiv \forall x\,\forall u\,(u \in S(x) \leftrightarrow u \in x \lor u = x) \to \forall u\,(u \in S(n) \leftrightarrow u \in n \lor u = n)$
\item $\varphi_{5} :\equiv (\neg S(n) \neq \emptyset \to \neg S(n) \neq \emptyset) \to \neg S(n) \neq \emptyset \to S(n) = \emptyset$
\item $\varphi_{6} :\equiv \forall w\,(S(n) = w \to n \in S(n) \to n \in w) \to S(n) = \emptyset \to n \in S(n) \to n \in \emptyset$
\item $\varphi_{7} :\equiv (S(n) = \emptyset \to n \in S(n) \to n \in \emptyset) \to \neg S(n) \neq \emptyset \to n \in \emptyset$
\item $\varphi_{8} :\equiv (\neg S(n) \neq \emptyset \to n \in \emptyset) \to (\neg S(n) \neq \emptyset \to n \notin \emptyset) \to S(n) \neq \emptyset$
\item $\varphi_{9} :\equiv (\neg S(n) \neq \emptyset \to S(n) = \emptyset) \to (\neg S(n) \neq \emptyset \to n \in S(n)) \to \varphi_{7}$
\end{abbrevlist}
\end{fullproof}

% 自動生成 (prover/build.js): NatSO  ノード数 24
\begin{fullproof}{[N2$_omega$]}{$\forall n \in \omega\,(S(n) \in \mathbb{N})$}{24}\phantomsection\label{tree:NatSO}
\noindent$\Pi_{1}$:
\begin{ptree}
\AXC{$\forall x\,(\mathrm{Ind}(x) \leftrightarrow \emptyset \in x \land \forall y \in x\,(S(y) \in x))$\lbl{DefInd}}
\AXC{$\varphi_{5}$\lbl{A2}}
\BIC{$\mathrm{Ind}(\omega) \leftrightarrow \emptyset \in \omega \land \forall y \in \omega\,(S(y) \in \omega)$}
\AXC{$\mathrm{Ind}(\omega)$\lbl{補題(ωは帰納的)}}
\AXC{$\mathrm{Ind}(\omega) \to \varphi_{4}$\lbl{A1}}
\BIC{$\varphi_{4}$}
\BIC{$n \in \omega \to \forall y \in \omega\,(S(y) \in \omega)$}
\end{ptree}
\noindent$\Pi_{2}$:
\begin{ptree}
\AXC{$\forall y \in \omega\,(S(y) \in \omega) \to n \in \omega \to S(n) \in \omega$\lbl{A2}}
\AXC{$\stackrel{\Pi_{1}}{\vdots}$}\noLine
\UIC{$n \in \omega \to \forall y \in \omega\,(S(y) \in \omega)$}
\AXC{$(n \in \omega \to \forall y \in \omega\,(S(y) \in \omega)) \to \varphi_{3}$\lbl{A1}}
\BIC{$\varphi_{3}$}
\BIC{$n \in \omega \to n \in \omega \to S(n) \in \omega$}
\end{ptree}
\noindent$\Pi_{3}$:
\begin{ptree}
\AXC{$n \in \omega \to n \in \omega$\lbl{A1}}
\AXC{$\stackrel{\Pi_{2}}{\vdots}$}\noLine
\UIC{$n \in \omega \to n \in \omega \to S(n) \in \omega$}
\AXC{$\varphi_{6}$\lbl{A1}}
\BIC{$(n \in \omega \to n \in \omega) \to S(n) \neq \emptyset \to \varphi_{1}$}
\BIC{$S(n) \neq \emptyset \to \varphi_{1}$}
\end{ptree}
\noindent$\Pi_{4}$:
\begin{ptree}
\AXC{$\forall n\,(S(n) \neq \emptyset)$\lbl{[N0]}}
\AXC{$\forall n\,(S(n) \neq \emptyset) \to S(n) \neq \emptyset$\lbl{A2}}
\BIC{$S(n) \neq \emptyset$}
\AXC{$\stackrel{\Pi_{3}}{\vdots}$}\noLine
\UIC{$S(n) \neq \emptyset \to \varphi_{1}$}
\BIC{$\varphi_{1}$}
\end{ptree}
\noindent 主木:
\begin{ptree}
\AXC{$\forall n\,(n \in \mathbb{N} \leftrightarrow n \in \omega \land n \neq \emptyset)$\lbl{Def$\mathbb{N}$}}
\AXC{$\varphi_{2}$\lbl{A2}}
\BIC{$S(n) \in \mathbb{N} \leftrightarrow S(n) \in \omega \land S(n) \neq \emptyset$}
\AXC{$\stackrel{\Pi_{4}}{\vdots}$}\noLine
\UIC{$\varphi_{1}$}
\BIC{$n \in \omega \to S(n) \in \mathbb{N}$}
\UIC{$\forall n \in \omega\,(S(n) \in \mathbb{N})$}
\end{ptree}
\begin{abbrevlist}
\item $\varphi_{1} :\equiv (S(n) \in \mathbb{N} \leftrightarrow S(n) \in \omega \land S(n) \neq \emptyset) \to n \in \omega \to S(n) \in \mathbb{N}$
\item $\varphi_{2} :\equiv \forall n\,(n \in \mathbb{N} \leftrightarrow n \in \omega \land n \neq \emptyset) \to (S(n) \in \mathbb{N} \leftrightarrow S(n) \in \omega \land S(n) \neq \emptyset)$
\item $\varphi_{3} :\equiv (\forall y \in \omega\,(S(y) \in \omega) \to n \in \omega \to S(n) \in \omega) \to n \in \omega \to n \in \omega \to S(n) \in \omega$
\item $\varphi_{4} :\equiv (\mathrm{Ind}(\omega) \leftrightarrow \emptyset \in \omega \land \forall y \in \omega\,(S(y) \in \omega)) \to n \in \omega \to \forall y \in \omega\,(S(y) \in \omega)$
\item $\varphi_{5} :\equiv \forall x\,(\mathrm{Ind}(x) \leftrightarrow \emptyset \in x \land \forall y \in x\,(S(y) \in x)) \to (\mathrm{Ind}(\omega) \leftrightarrow \emptyset \in \omega \land \forall y \in \omega\,(S(y) \in \omega))$
\item $\varphi_{6} :\equiv (n \in \omega \to n \in \omega \to S(n) \in \omega) \to (n \in \omega \to n \in \omega) \to S(n) \neq \emptyset \to \varphi_{1}$
\end{abbrevlist}
\end{fullproof}

% 自動生成 (prover/build.js): NatOne  ノード数 39
\begin{fullproof}{[N1]}{$1_{\omega} \in \mathbb{N}$}{39}\phantomsection\label{tree:NatOne}
\noindent$\Pi_{1}$:
\begin{ptree}
\AXC{$\forall x\,(\mathrm{Ind}(x) \leftrightarrow \emptyset \in x \land \forall y \in x\,(S(y) \in x))$\lbl{DefInd}}
\AXC{$\varphi_{8}$\lbl{A2}}
\BIC{$\mathrm{Ind}(\omega) \leftrightarrow \emptyset \in \omega \land \forall y \in \omega\,(S(y) \in \omega)$}
\AXC{$\mathrm{Ind}(\omega)$\lbl{補題(ωは帰納的)}}
\AXC{$\varphi_{6}$\lbl{A1}}
\BIC{$(\mathrm{Ind}(\omega) \leftrightarrow \emptyset \in \omega \land \forall y \in \omega\,(S(y) \in \omega)) \to \emptyset \in \omega$}
\BIC{$\emptyset \in \omega$}
\end{ptree}
\noindent$\Pi_{2}$:
\begin{ptree}
\AXC{$\forall n \in \omega\,(S(n) \in \mathbb{N})$\lbl{[N2$_omega$]}}
\AXC{$\forall n \in \omega\,(S(n) \in \mathbb{N}) \to \emptyset \in \omega \to S(\emptyset) \in \mathbb{N}$\lbl{A2}}
\BIC{$\emptyset \in \omega \to S(\emptyset) \in \mathbb{N}$}
\end{ptree}
\noindent$\Pi_{3}$:
\begin{ptree}
\AXC{$\stackrel{\Pi_{1}}{\vdots}$}\noLine
\UIC{$\emptyset \in \omega$}
\AXC{$\stackrel{\Pi_{2}}{\vdots}$}\noLine
\UIC{$\emptyset \in \omega \to S(\emptyset) \in \mathbb{N}$}
\AXC{$(\emptyset \in \omega \to S(\emptyset) \in \mathbb{N}) \to \emptyset \in \omega \to S(\emptyset) \in \mathbb{N}$\lbl{A1}}
\BIC{$\emptyset \in \omega \to S(\emptyset) \in \mathbb{N}$}
\BIC{$S(\emptyset) \in \mathbb{N}$}
\end{ptree}
\noindent$\Pi_{4}$:
\begin{ptree}
\AXC{$x = x$\lbl{A5}}
\AXC{$x = y \to x = x \to y = x$\lbl{A6}}
\AXC{$(x = y \to x = x \to y = x) \to x = x \to x = y \to y = x$\lbl{A1}}
\BIC{$x = x \to x = y \to y = x$}
\BIC{$x = y \to y = x$}
\UIC{$\forall y\,(x = y \to y = x)$}
\UIC{$\forall x\,\forall y\,(x = y \to y = x)$}
\AXC{$\varphi_{1}$\lbl{A2}}
\BIC{$\forall y\,(1_{\omega} = y \to y = 1_{\omega})$}
\AXC{$\varphi_{3}$\lbl{A2}}
\BIC{$1_{\omega} = S(\emptyset) \to S(\emptyset) = 1_{\omega}$}
\end{ptree}
\noindent$\Pi_{5}$:
\begin{ptree}
\AXC{$\stackrel{\Pi_{4}}{\vdots}$}\noLine
\UIC{$1_{\omega} = S(\emptyset) \to S(\emptyset) = 1_{\omega}$}
\AXC{$1_{\omega} = S(\emptyset)$\lbl{Def$1_\omega$}}
\AXC{$1_{\omega} = S(\emptyset) \to \varphi_{2}$\lbl{A1}}
\BIC{$\varphi_{2}$}
\BIC{$S(\emptyset) = 1_{\omega}$}
\end{ptree}
\noindent$\Pi_{6}$:
\begin{ptree}
\AXC{$\stackrel{\Pi_{3}}{\vdots}$}\noLine
\UIC{$S(\emptyset) \in \mathbb{N}$}
\AXC{$\stackrel{\Pi_{5}}{\vdots}$}\noLine
\UIC{$S(\emptyset) = 1_{\omega}$}
\AXC{$S(\emptyset) = 1_{\omega} \to S(\emptyset) \in \mathbb{N} \to \varphi_{4}$\lbl{A1}}
\BIC{$S(\emptyset) \in \mathbb{N} \to \varphi_{4}$}
\BIC{$\varphi_{4}$}
\end{ptree}
\noindent 主木:
\begin{ptree}
\AXC{$z = w \to z \in \mathbb{N} \to w \in \mathbb{N}$\lbl{A6}}
\UIC{$\forall w\,(z = w \to z \in \mathbb{N} \to w \in \mathbb{N})$}
\UIC{$\forall z\,\forall w\,(z = w \to z \in \mathbb{N} \to w \in \mathbb{N})$}
\AXC{$\varphi_{5}$\lbl{A2}}
\BIC{$\forall w\,(S(\emptyset) = w \to S(\emptyset) \in \mathbb{N} \to w \in \mathbb{N})$}
\AXC{$\varphi_{7}$\lbl{A2}}
\BIC{$S(\emptyset) = 1_{\omega} \to S(\emptyset) \in \mathbb{N} \to 1_{\omega} \in \mathbb{N}$}
\AXC{$\stackrel{\Pi_{6}}{\vdots}$}\noLine
\UIC{$\varphi_{4}$}
\BIC{$1_{\omega} \in \mathbb{N}$}
\end{ptree}
\begin{abbrevlist}
\item $\varphi_{1} :\equiv \forall x\,\forall y\,(x = y \to y = x) \to \forall y\,(1_{\omega} = y \to y = 1_{\omega})$
\item $\varphi_{2} :\equiv (1_{\omega} = S(\emptyset) \to S(\emptyset) = 1_{\omega}) \to S(\emptyset) = 1_{\omega}$
\item $\varphi_{3} :\equiv \forall y\,(1_{\omega} = y \to y = 1_{\omega}) \to 1_{\omega} = S(\emptyset) \to S(\emptyset) = 1_{\omega}$
\item $\varphi_{4} :\equiv (S(\emptyset) = 1_{\omega} \to S(\emptyset) \in \mathbb{N} \to 1_{\omega} \in \mathbb{N}) \to 1_{\omega} \in \mathbb{N}$
\item $\varphi_{5} :\equiv \forall z\,\forall w\,(z = w \to z \in \mathbb{N} \to w \in \mathbb{N}) \to \forall w\,(S(\emptyset) = w \to S(\emptyset) \in \mathbb{N} \to w \in \mathbb{N})$
\item $\varphi_{6} :\equiv \mathrm{Ind}(\omega) \to (\mathrm{Ind}(\omega) \leftrightarrow \emptyset \in \omega \land \forall y \in \omega\,(S(y) \in \omega)) \to \emptyset \in \omega$
\item $\varphi_{7} :\equiv \forall w\,(S(\emptyset) = w \to S(\emptyset) \in \mathbb{N} \to w \in \mathbb{N}) \to S(\emptyset) = 1_{\omega} \to S(\emptyset) \in \mathbb{N} \to 1_{\omega} \in \mathbb{N}$
\item $\varphi_{8} :\equiv \forall x\,(\mathrm{Ind}(x) \leftrightarrow \emptyset \in x \land \forall y \in x\,(S(y) \in x)) \to (\mathrm{Ind}(\omega) \leftrightarrow \emptyset \in \omega \land \forall y \in \omega\,(S(y) \in \omega))$
\end{abbrevlist}
\end{fullproof}

% 自動生成 (prover/build.js): NatS  ノード数 14
\begin{fullproof}{[N2]}{$\forall n \in \mathbb{N}\,(S(n) \in \mathbb{N})$}{14}\phantomsection\label{tree:NatS}
\noindent$\Pi_{1}$:
\begin{ptree}
\AXC{$\forall n\,(n \in \mathbb{N} \leftrightarrow n \in \omega \land n \neq \emptyset)$\lbl{Def$\mathbb{N}$}}
\AXC{$\varphi_{2}$\lbl{A2}}
\BIC{$n \in \mathbb{N} \leftrightarrow n \in \omega \land n \neq \emptyset$}
\AXC{$n \in \mathbb{N} \to n \in \mathbb{N}$\lbl{A1}}
\AXC{$\varphi_{3}$\lbl{A1}}
\BIC{$(n \in \mathbb{N} \leftrightarrow n \in \omega \land n \neq \emptyset) \to n \in \mathbb{N} \to n \in \omega$}
\BIC{$n \in \mathbb{N} \to n \in \omega$}
\end{ptree}
\noindent$\Pi_{2}$:
\begin{ptree}
\AXC{$\forall n \in \omega\,(S(n) \in \mathbb{N})$\lbl{[N2$_omega$]}}
\AXC{$\forall n \in \omega\,(S(n) \in \mathbb{N}) \to n \in \omega \to S(n) \in \mathbb{N}$\lbl{A2}}
\BIC{$n \in \omega \to S(n) \in \mathbb{N}$}
\AXC{$\varphi_{1}$\lbl{A1}}
\BIC{$(n \in \mathbb{N} \to n \in \omega) \to n \in \mathbb{N} \to S(n) \in \mathbb{N}$}
\end{ptree}
\noindent 主木:
\begin{ptree}
\AXC{$\stackrel{\Pi_{1}}{\vdots}$}\noLine
\UIC{$n \in \mathbb{N} \to n \in \omega$}
\AXC{$\stackrel{\Pi_{2}}{\vdots}$}\noLine
\UIC{$(n \in \mathbb{N} \to n \in \omega) \to n \in \mathbb{N} \to S(n) \in \mathbb{N}$}
\BIC{$n \in \mathbb{N} \to S(n) \in \mathbb{N}$}
\UIC{$\forall n \in \mathbb{N}\,(S(n) \in \mathbb{N})$}
\end{ptree}
\begin{abbrevlist}
\item $\varphi_{1} :\equiv (n \in \omega \to S(n) \in \mathbb{N}) \to (n \in \mathbb{N} \to n \in \omega) \to n \in \mathbb{N} \to S(n) \in \mathbb{N}$
\item $\varphi_{2} :\equiv \forall n\,(n \in \mathbb{N} \leftrightarrow n \in \omega \land n \neq \emptyset) \to (n \in \mathbb{N} \leftrightarrow n \in \omega \land n \neq \emptyset)$
\item $\varphi_{3} :\equiv (n \in \mathbb{N} \to n \in \mathbb{N}) \to (n \in \mathbb{N} \leftrightarrow n \in \omega \land n \neq \emptyset) \to n \in \mathbb{N} \to n \in \omega$
\end{abbrevlist}
\end{fullproof}

% 自動生成 (prover/build.js): NatOmega  ノード数 8
\begin{fullproof}{[N7']}{$\forall n \in \mathbb{N}\,(n \in \omega)$}{8}\phantomsection\label{tree:NatOmega}
\noindent 主木:
\begin{ptree}
\AXC{$\forall n\,(n \in \mathbb{N} \leftrightarrow n \in \omega \land n \neq \emptyset)$\lbl{Def$\mathbb{N}$}}
\AXC{$\varphi_{1}$\lbl{A2}}
\BIC{$n \in \mathbb{N} \leftrightarrow n \in \omega \land n \neq \emptyset$}
\AXC{$n \in \mathbb{N} \to n \in \mathbb{N}$\lbl{A1}}
\AXC{$\varphi_{2}$\lbl{A1}}
\BIC{$(n \in \mathbb{N} \leftrightarrow n \in \omega \land n \neq \emptyset) \to n \in \mathbb{N} \to n \in \omega$}
\BIC{$n \in \mathbb{N} \to n \in \omega$}
\UIC{$\forall n \in \mathbb{N}\,(n \in \omega)$}
\end{ptree}
\begin{abbrevlist}
\item $\varphi_{1} :\equiv \forall n\,(n \in \mathbb{N} \leftrightarrow n \in \omega \land n \neq \emptyset) \to (n \in \mathbb{N} \leftrightarrow n \in \omega \land n \neq \emptyset)$
\item $\varphi_{2} :\equiv (n \in \mathbb{N} \to n \in \mathbb{N}) \to (n \in \mathbb{N} \leftrightarrow n \in \omega \land n \neq \emptyset) \to n \in \mathbb{N} \to n \in \omega$
\end{abbrevlist}
\end{fullproof}

% 自動生成 (prover/build.js): LeS  ノード数 23
\begin{fullproof}{[N8]}{$\forall n\,(n \leq S(n))$}{23}\phantomsection\label{tree:LeS}
\noindent$\Pi_{1}$:
\begin{ptree}
\AXC{$\forall x\,\forall u\,(u \in S(x) \leftrightarrow u \in x \lor u = x)$\lbl{Def$S$}}
\AXC{$\varphi_{4}$\lbl{A2}}
\BIC{$\forall u\,(u \in S(n) \leftrightarrow u \in n \lor u = n)$}
\AXC{$\varphi_{1}$\lbl{A2}}
\BIC{$u \in S(n) \leftrightarrow u \in n \lor u = n$}
\AXC{$(u \in S(n) \leftrightarrow u \in n \lor u = n) \to u \in n \to u \in S(n)$\lbl{A1}}
\BIC{$u \in n \to u \in S(n)$}
\UIC{$\forall u \in n\,(u \in S(n))$}
\end{ptree}
\noindent$\Pi_{2}$:
\begin{ptree}
\AXC{$\forall x\,\forall y\,(x < y \leftrightarrow x \subseteq y \land x \neq y)$\lbl{Def$<$}}
\AXC{$\varphi_{2}$\lbl{A2}}
\BIC{$\forall y\,(n < y \leftrightarrow n \subseteq y \land n \neq y)$}
\AXC{$\varphi_{3}$\lbl{A2}}
\BIC{$n < S(n) \leftrightarrow n \subseteq S(n) \land n \neq S(n)$}
\AXC{$(n < S(n) \leftrightarrow n \subseteq S(n) \land n \neq S(n)) \to \varphi_{5}$\lbl{A1}}
\BIC{$\varphi_{5}$}
\end{ptree}
\noindent 主木:
\begin{ptree}
\AXC{$\stackrel{\Pi_{1}}{\vdots}$}\noLine
\UIC{$\forall u \in n\,(u \in S(n))$}
\AXC{$\forall x\,\forall y\,(x \subseteq y \leftrightarrow \forall u \in x\,(u \in y))$\lbl{Def$\subseteq$}}
\AXC{$\varphi_{6}$\lbl{A2}}
\BIC{$\forall y\,(n \subseteq y \leftrightarrow \forall u \in n\,(u \in y))$}
\AXC{$\varphi_{7}$\lbl{A2}}
\BIC{$n \subseteq S(n) \leftrightarrow \forall u \in n\,(u \in S(n))$}
\AXC{$\stackrel{\Pi_{2}}{\vdots}$}\noLine
\UIC{$\varphi_{5}$}
\BIC{$\forall u \in n\,(u \in S(n)) \to n \leq S(n)$}
\BIC{$n \leq S(n)$}
\UIC{$\forall n\,(n \leq S(n))$}
\end{ptree}
\begin{abbrevlist}
\item $\varphi_{1} :\equiv \forall u\,(u \in S(n) \leftrightarrow u \in n \lor u = n) \to (u \in S(n) \leftrightarrow u \in n \lor u = n)$
\item $\varphi_{2} :\equiv \forall x\,\forall y\,(x < y \leftrightarrow x \subseteq y \land x \neq y) \to \forall y\,(n < y \leftrightarrow n \subseteq y \land n \neq y)$
\item $\varphi_{3} :\equiv \forall y\,(n < y \leftrightarrow n \subseteq y \land n \neq y) \to (n < S(n) \leftrightarrow n \subseteq S(n) \land n \neq S(n))$
\item $\varphi_{4} :\equiv \forall x\,\forall u\,(u \in S(x) \leftrightarrow u \in x \lor u = x) \to \forall u\,(u \in S(n) \leftrightarrow u \in n \lor u = n)$
\item $\varphi_{5} :\equiv (n \subseteq S(n) \leftrightarrow \forall u \in n\,(u \in S(n))) \to \forall u \in n\,(u \in S(n)) \to n \leq S(n)$
\item $\varphi_{6} :\equiv \forall x\,\forall y\,(x \subseteq y \leftrightarrow \forall u \in x\,(u \in y)) \to \forall y\,(n \subseteq y \leftrightarrow \forall u \in n\,(u \in y))$
\item $\varphi_{7} :\equiv \forall y\,(n \subseteq y \leftrightarrow \forall u \in n\,(u \in y)) \to (n \subseteq S(n) \leftrightarrow \forall u \in n\,(u \in S(n)))$
\end{abbrevlist}
\end{fullproof}

% 自動生成 (prover/build.js): OneLe  ノード数 230
\begin{fullproof}{[N12]}{$\forall n \in \mathbb{N}\,(1_{\omega} \leq n)$}{230}\phantomsection\label{tree:OneLe}
\noindent$\Pi_{1}$:
\begin{ptree}
\AXC{$\forall n\,(n \in \mathbb{N} \leftrightarrow n \in \omega \land n \neq \emptyset)$\lbl{Def$\mathbb{N}$}}
\AXC{$\varphi_{41}$\lbl{A2}}
\BIC{$n \in \mathbb{N} \leftrightarrow n \in \omega \land n \neq \emptyset$}
\AXC{$n \in \mathbb{N} \to n \in \mathbb{N}$\lbl{A1}}
\AXC{$\varphi_{48}$\lbl{A1}}
\BIC{$(n \in \mathbb{N} \leftrightarrow n \in \omega \land n \neq \emptyset) \to n \in \mathbb{N} \to n \in \omega$}
\BIC{$n \in \mathbb{N} \to n \in \omega$}
\end{ptree}
\noindent$\Pi_{2}$:
\begin{ptree}
\AXC{$\varphi_{19}$\lbl{A4}}
\AXC{$\varphi_{11} \to \varphi_{23}$\lbl{A3}}
\AXC{$\forall n \in B_{1}\,(S(n) \in B_{1}) \to y \in B_{1} \to S(y) \in B_{1}$\lbl{A2}}
\UIC{$\varphi_{11}$}
\AXC{$\varphi_{11} \to (\varphi_{11} \to \varphi_{23}) \to \varphi_{68}$\lbl{A1}}
\BIC{$(\varphi_{11} \to \varphi_{23}) \to \varphi_{68}$}
\BIC{$\varphi_{68}$}
\BIC{$\forall n \in B_{1}\,(S(n) \in B_{1}) \to \forall y \in B_{1}\,(S(y) \in B_{1})$}
\end{ptree}
\noindent$\Pi_{3}$:
\begin{ptree}
\AXC{$\varphi_{20}$\lbl{A4}}
\AXC{$\varphi_{12} \to \varphi_{24}$\lbl{A3}}
\AXC{$\forall y \in B_{1}\,(S(y) \in B_{1}) \to n \in B_{1} \to S(n) \in B_{1}$\lbl{A2}}
\UIC{$\varphi_{12}$}
\AXC{$\varphi_{12} \to (\varphi_{12} \to \varphi_{24}) \to \varphi_{69}$\lbl{A1}}
\BIC{$(\varphi_{12} \to \varphi_{24}) \to \varphi_{69}$}
\BIC{$\varphi_{69}$}
\BIC{$\forall y \in B_{1}\,(S(y) \in B_{1}) \to \forall n \in B_{1}\,(S(n) \in B_{1})$}
\end{ptree}
\noindent$\Pi_{4}$:
\begin{ptree}
\AXC{$\varphi_{27}$\lbl{A1}}
\AXC{$\varphi_{52}$\lbl{A2}}
\AXC{$\varphi_{52} \to \varphi_{27} \to \varphi_{53}$\lbl{A1}}
\BIC{$\varphi_{27} \to \varphi_{53}$}
\BIC{$\varphi_{53}$}
\AXC{$\varphi_{53} \to \varphi_{30}$\lbl{A1}}
\BIC{$\varphi_{30}$}
\end{ptree}
\noindent$\Pi_{5}$:
\begin{ptree}
\AXC{$\varphi_{62}$\lbl{A1}}
\AXC{$\stackrel{\Pi_{4}}{\vdots}$}\noLine
\UIC{$\varphi_{30}$}
\AXC{$\forall n\,(n \in \mathbb{N} \leftrightarrow n \in \omega \land n \neq \emptyset)$\lbl{Def$\mathbb{N}$}}
\AXC{$\varphi_{41}$\lbl{A2}}
\BIC{$n \in \mathbb{N} \leftrightarrow n \in \omega \land n \neq \emptyset$}
\AXC{$(n \in \mathbb{N} \leftrightarrow n \in \omega \land n \neq \emptyset) \to \varphi_{30} \to \varphi_{62} \to \varphi_{59}$\lbl{A1}}
\BIC{$\varphi_{30} \to \varphi_{62} \to \varphi_{59}$}
\BIC{$\varphi_{62} \to \varphi_{59}$}
\BIC{$\varphi_{59}$}
\end{ptree}
\noindent$\Pi_{6}$:
\begin{ptree}
\AXC{$\stackrel{\Pi_{5}}{\vdots}$}\noLine
\UIC{$\varphi_{59}$}
\AXC{$\varphi_{27}$\lbl{A1}}
\AXC{$\varphi_{52}$\lbl{A2}}
\AXC{$\varphi_{52} \to \varphi_{27} \to \varphi_{53}$\lbl{A1}}
\BIC{$\varphi_{27} \to \varphi_{53}$}
\BIC{$\varphi_{53}$}
\AXC{$\varphi_{53} \to \varphi_{39}$\lbl{A1}}
\BIC{$\varphi_{39}$}
\AXC{$\varphi_{39} \to \varphi_{59} \to \varphi_{54}$\lbl{A1}}
\BIC{$\varphi_{59} \to \varphi_{54}$}
\BIC{$\varphi_{54}$}
\end{ptree}
\noindent$\Pi_{7}$:
\begin{ptree}
\AXC{$\stackrel{\Pi_{6}}{\vdots}$}\noLine
\UIC{$\varphi_{54}$}
\AXC{$\forall x\,\forall y\,\forall z\,(x \leq y \to y \leq z \to x \leq z)$\lbl{[O5]}}
\AXC{$\varphi_{13}$\lbl{A2}}
\BIC{$\forall y\,\forall z\,(1_{\omega} \leq y \to y \leq z \to 1_{\omega} \leq z)$}
\AXC{$\varphi_{14}$\lbl{A2}}
\BIC{$\forall z\,(1_{\omega} \leq n \to n \leq z \to 1_{\omega} \leq z)$}
\AXC{$\varphi_{15}$\lbl{A2}}
\BIC{$1_{\omega} \leq n \to n \leq S(n) \to 1_{\omega} \leq S(n)$}
\AXC{$\varphi_{79}$\lbl{A1}}
\BIC{$\varphi_{54} \to n \leq S(n) \to \varphi_{60}$}
\BIC{$n \leq S(n) \to \varphi_{60}$}
\end{ptree}
\noindent$\Pi_{8}$:
\begin{ptree}
\AXC{$z = w \to 1_{\omega} = z \to 1_{\omega} = w$\lbl{A6}}
\UIC{$\forall w\,(z = w \to 1_{\omega} = z \to 1_{\omega} = w)$}
\UIC{$\forall z\,\forall w\,(z = w \to 1_{\omega} = z \to 1_{\omega} = w)$}
\AXC{$\forall z\,\forall w\,(z = w \to 1_{\omega} = z \to 1_{\omega} = w) \to \varphi_{1}$\lbl{A2}}
\BIC{$\varphi_{1}$}
\AXC{$\varphi_{1} \to \varphi_{2}$\lbl{A2}}
\BIC{$\varphi_{2}$}
\end{ptree}
\noindent$\Pi_{9}$:
\begin{ptree}
\AXC{$z = w \to 1_{\omega} < z \to 1_{\omega} < w$\lbl{A6}}
\UIC{$\forall w\,(z = w \to 1_{\omega} < z \to 1_{\omega} < w)$}
\UIC{$\forall z\,\forall w\,(z = w \to 1_{\omega} < z \to 1_{\omega} < w)$}
\AXC{$\forall z\,\forall w\,(z = w \to 1_{\omega} < z \to 1_{\omega} < w) \to \varphi_{3}$\lbl{A2}}
\BIC{$\varphi_{3}$}
\AXC{$\varphi_{3} \to \varphi_{6}$\lbl{A2}}
\BIC{$\varphi_{6}$}
\end{ptree}
\noindent$\Pi_{10}$:
\begin{ptree}
\AXC{$x = x$\lbl{A5}}
\UIC{$\forall x\,(x = x)$}
\AXC{$\forall x\,(x = x) \to 1_{\omega} = 1_{\omega}$\lbl{A2}}
\BIC{$1_{\omega} = 1_{\omega}$}
\AXC{$1_{\omega} = 1_{\omega} \to 1_{\omega} \leq 1_{\omega}$\lbl{A1}}
\BIC{$1_{\omega} \leq 1_{\omega}$}
\end{ptree}
\noindent$\Pi_{11}$:
\begin{ptree}
\AXC{$x = x$\lbl{A5}}
\AXC{$x = y \to x = x \to y = x$\lbl{A6}}
\AXC{$(x = y \to x = x \to y = x) \to x = x \to x = y \to y = x$\lbl{A1}}
\BIC{$x = x \to x = y \to y = x$}
\BIC{$x = y \to y = x$}
\UIC{$\forall y\,(x = y \to y = x)$}
\UIC{$\forall x\,\forall y\,(x = y \to y = x)$}
\end{ptree}
\noindent$\Pi_{12}$:
\begin{ptree}
\AXC{$y = z \to x = y \to x = z$\lbl{A6}}
\AXC{$(y = z \to x = y \to x = z) \to x = y \to y = z \to x = z$\lbl{A1}}
\BIC{$x = y \to y = z \to x = z$}
\UIC{$\forall z\,(x = y \to y = z \to x = z)$}
\UIC{$\forall y\,\forall z\,(x = y \to y = z \to x = z)$}
\UIC{$\forall x\,\forall y\,\forall z\,(x = y \to y = z \to x = z)$}
\AXC{$\varphi_{16}$\lbl{A2}}
\BIC{$\forall y\,\forall z\,(S(n) = y \to y = z \to S(n) = z)$}
\AXC{$\varphi_{33}$\lbl{A2}}
\BIC{$\forall z\,(S(n) = S(\emptyset) \to S(\emptyset) = z \to S(n) = z)$}
\end{ptree}
\noindent$\Pi_{13}$:
\begin{ptree}
\AXC{$\stackrel{\Pi_{12}}{\vdots}$}\noLine
\UIC{$\forall z\,(S(n) = S(\emptyset) \to S(\emptyset) = z \to S(n) = z)$}
\AXC{$\forall z\,(S(n) = S(\emptyset) \to S(\emptyset) = z \to S(n) = z) \to \varphi_{8}$\lbl{A2}}
\BIC{$\varphi_{8}$}
\end{ptree}
\noindent$\Pi_{14}$:
\begin{ptree}
\AXC{$\stackrel{\Pi_{11}}{\vdots}$}\noLine
\UIC{$\forall x\,\forall y\,(x = y \to y = x)$}
\AXC{$\varphi_{5}$\lbl{A2}}
\BIC{$\forall y\,(1_{\omega} = y \to y = 1_{\omega})$}
\AXC{$\varphi_{21}$\lbl{A2}}
\BIC{$1_{\omega} = S(\emptyset) \to S(\emptyset) = 1_{\omega}$}
\AXC{$1_{\omega} = S(\emptyset)$\lbl{Def$1_\omega$}}
\AXC{$1_{\omega} = S(\emptyset) \to \varphi_{17}$\lbl{A1}}
\BIC{$\varphi_{17}$}
\BIC{$S(\emptyset) = 1_{\omega}$}
\end{ptree}
\noindent$\Pi_{15}$:
\begin{ptree}
\AXC{$z = w \to S(n) = S(z) \to S(n) = S(w)$\lbl{A6}}
\UIC{$\forall w\,(z = w \to S(n) = S(z) \to S(n) = S(w))$}
\UIC{$\forall z\,\forall w\,(z = w \to S(n) = S(z) \to S(n) = S(w))$}
\AXC{$\varphi_{26}$\lbl{A2}}
\BIC{$\forall w\,(n = w \to S(n) = S(n) \to S(n) = S(w))$}
\AXC{$\varphi_{34}$\lbl{A2}}
\BIC{$n = \emptyset \to S(n) = S(n) \to S(n) = S(\emptyset)$}
\end{ptree}
\noindent$\Pi_{16}$:
\begin{ptree}
\AXC{$x = x$\lbl{A5}}
\UIC{$\forall x\,(x = x)$}
\AXC{$\forall x\,(x = x) \to S(n) = S(n)$\lbl{A2}}
\BIC{$S(n) = S(n)$}
\AXC{$\varphi_{61}$\lbl{A1}}
\AXC{$\varphi_{61} \to \varphi_{75}$\lbl{A1}}
\BIC{$\varphi_{75}$}
\BIC{$(n = \emptyset \to S(n) = S(n) \to S(n) = S(\emptyset)) \to \varphi_{65}$}
\end{ptree}
\noindent$\Pi_{17}$:
\begin{ptree}
\AXC{$\stackrel{\Pi_{15}}{\vdots}$}\noLine
\UIC{$n = \emptyset \to S(n) = S(n) \to S(n) = S(\emptyset)$}
\AXC{$\stackrel{\Pi_{16}}{\vdots}$}\noLine
\UIC{$(n = \emptyset \to S(n) = S(n) \to S(n) = S(\emptyset)) \to \varphi_{65}$}
\BIC{$\varphi_{65}$}
\end{ptree}
\noindent$\Pi_{18}$:
\begin{ptree}
\AXC{$\stackrel{\Pi_{13}}{\vdots}$}\noLine
\UIC{$\varphi_{8}$}
\AXC{$\stackrel{\Pi_{14}}{\vdots}$}\noLine
\UIC{$S(\emptyset) = 1_{\omega}$}
\AXC{$\stackrel{\Pi_{17}}{\vdots}$}\noLine
\UIC{$\varphi_{65}$}
\AXC{$\varphi_{65} \to S(\emptyset) = 1_{\omega} \to \varphi_{8} \to \varphi_{55}$\lbl{A1}}
\BIC{$S(\emptyset) = 1_{\omega} \to \varphi_{8} \to \varphi_{55}$}
\BIC{$\varphi_{8} \to \varphi_{55}$}
\BIC{$\varphi_{55}$}
\end{ptree}
\noindent$\Pi_{19}$:
\begin{ptree}
\AXC{$\stackrel{\Pi_{11}}{\vdots}$}\noLine
\UIC{$\forall x\,\forall y\,(x = y \to y = x)$}
\AXC{$\varphi_{4}$\lbl{A2}}
\BIC{$\forall y\,(S(n) = y \to y = S(n))$}
\AXC{$\varphi_{7}$\lbl{A2}}
\BIC{$S(n) = 1_{\omega} \to 1_{\omega} = S(n)$}
\AXC{$\stackrel{\Pi_{18}}{\vdots}$}\noLine
\UIC{$\varphi_{55}$}
\AXC{$\varphi_{55} \to (S(n) = 1_{\omega} \to 1_{\omega} = S(n)) \to \varphi_{56}$\lbl{A1}}
\BIC{$(S(n) = 1_{\omega} \to 1_{\omega} = S(n)) \to \varphi_{56}$}
\BIC{$\varphi_{56}$}
\end{ptree}
\noindent$\Pi_{20}$:
\begin{ptree}
\AXC{$\stackrel{\Pi_{8}}{\vdots}$}\noLine
\UIC{$\varphi_{2}$}
\AXC{$\stackrel{\Pi_{9}}{\vdots}$}\noLine
\UIC{$\varphi_{6}$}
\AXC{$\stackrel{\Pi_{10}}{\vdots}$}\noLine
\UIC{$1_{\omega} \leq 1_{\omega}$}
\AXC{$\stackrel{\Pi_{19}}{\vdots}$}\noLine
\UIC{$\varphi_{56}$}
\AXC{$\varphi_{56} \to 1_{\omega} \leq 1_{\omega} \to \varphi_{6} \to \varphi_{2} \to \varphi_{57}$\lbl{A1}}
\BIC{$1_{\omega} \leq 1_{\omega} \to \varphi_{6} \to \varphi_{2} \to \varphi_{57}$}
\BIC{$\varphi_{6} \to \varphi_{2} \to \varphi_{57}$}
\BIC{$\varphi_{2} \to \varphi_{57}$}
\BIC{$\varphi_{57}$}
\AXC{$\varphi_{57} \to \varphi_{60} \to \varphi_{45}$\lbl{A1}}
\BIC{$\varphi_{60} \to \varphi_{45}$}
\end{ptree}
\noindent$\Pi_{21}$:
\begin{ptree}
\AXC{$\forall n\,(n \leq S(n))$\lbl{[N8]}}
\AXC{$\forall n\,(n \leq S(n)) \to n \leq S(n)$\lbl{A2}}
\BIC{$n \leq S(n)$}
\AXC{$\stackrel{\Pi_{7}}{\vdots}$}\noLine
\UIC{$n \leq S(n) \to \varphi_{60}$}
\BIC{$\varphi_{60}$}
\AXC{$\stackrel{\Pi_{20}}{\vdots}$}\noLine
\UIC{$\varphi_{60} \to \varphi_{45}$}
\BIC{$\varphi_{45}$}
\end{ptree}
\noindent$\Pi_{22}$:
\begin{ptree}
\AXC{$\forall x\,(\mathrm{Ind}(x) \leftrightarrow \emptyset \in x \land \forall y \in x\,(S(y) \in x))$\lbl{DefInd}}
\AXC{$\varphi_{72}$\lbl{A2}}
\BIC{$\mathrm{Ind}(\omega) \leftrightarrow \emptyset \in \omega \land \forall y \in \omega\,(S(y) \in \omega)$}
\AXC{$\mathrm{Ind}(\omega)$\lbl{補題(ωは帰納的)}}
\AXC{$\varphi_{77}$\lbl{A1}}
\BIC{$(\mathrm{Ind}(\omega) \leftrightarrow \emptyset \in \omega \land \forall y \in \omega\,(S(y) \in \omega)) \to \varphi_{63}$}
\BIC{$\varphi_{63}$}
\end{ptree}
\noindent$\Pi_{23}$:
\begin{ptree}
\AXC{$\stackrel{\Pi_{22}}{\vdots}$}\noLine
\UIC{$\varphi_{63}$}
\AXC{$\varphi_{63} \to \varphi_{51}$\lbl{A1}}
\BIC{$\varphi_{51}$}
\AXC{$\varphi_{51} \to (\forall y \in \omega\,(S(y) \in \omega) \to n \in \omega \to S(n) \in \omega) \to \varphi_{50}$\lbl{A1}}
\BIC{$(\forall y \in \omega\,(S(y) \in \omega) \to n \in \omega \to S(n) \in \omega) \to \varphi_{50}$}
\end{ptree}
\noindent$\Pi_{24}$:
\begin{ptree}
\AXC{$\forall y \in \omega\,(S(y) \in \omega) \to n \in \omega \to S(n) \in \omega$\lbl{A2}}
\AXC{$\stackrel{\Pi_{23}}{\vdots}$}\noLine
\UIC{$(\forall y \in \omega\,(S(y) \in \omega) \to n \in \omega \to S(n) \in \omega) \to \varphi_{50}$}
\BIC{$\varphi_{50}$}
\AXC{$\varphi_{50} \to \varphi_{30} \to \varphi_{35}$\lbl{A1}}
\BIC{$\varphi_{30} \to \varphi_{35}$}
\end{ptree}
\noindent$\Pi_{25}$:
\begin{ptree}
\AXC{$\stackrel{\Pi_{21}}{\vdots}$}\noLine
\UIC{$\varphi_{45}$}
\AXC{$\stackrel{\Pi_{4}}{\vdots}$}\noLine
\UIC{$\varphi_{30}$}
\AXC{$\stackrel{\Pi_{24}}{\vdots}$}\noLine
\UIC{$\varphi_{30} \to \varphi_{35}$}
\BIC{$\varphi_{35}$}
\AXC{$\forall n\,(n \in B_{1} \leftrightarrow n \in \omega \land (n \in \mathbb{N} \to 1_{\omega} \leq n)) \to \varphi_{18}$\lbl{A2}}
\AXC{$\varphi_{80}$\lbl{A1}}
\BIC{$\varphi_{35} \to \varphi_{45} \to \varphi_{31}$}
\BIC{$\varphi_{45} \to \varphi_{31}$}
\BIC{$\varphi_{31}$}
\UIC{$\forall n\,(\varphi_{31})$}
\end{ptree}
\noindent$\Pi_{26}$:
\begin{ptree}
\AXC{$\varphi_{58}$\lbl{A4}}
\AXC{$\forall n\,(\varphi_{31}) \to \varphi_{38}$\lbl{A3}}
\AXC{$\stackrel{\Pi_{25}}{\vdots}$}\noLine
\UIC{$\forall n\,(\varphi_{31})$}
\AXC{$\forall n\,(\varphi_{31}) \to (\forall n\,(\varphi_{31}) \to \varphi_{38}) \to \varphi_{58} \to \varphi_{36}$\lbl{A1}}
\BIC{$(\forall n\,(\varphi_{31}) \to \varphi_{38}) \to \varphi_{58} \to \varphi_{36}$}
\BIC{$\varphi_{58} \to \varphi_{36}$}
\BIC{$\varphi_{36}$}
\end{ptree}
\noindent$\Pi_{27}$:
\begin{ptree}
\AXC{$x = x$\lbl{A5}}
\UIC{$\forall x\,(x = x)$}
\AXC{$\forall x\,(x = x) \to \emptyset = \emptyset$\lbl{A2}}
\BIC{$\emptyset = \emptyset$}
\AXC{$\forall n\,(n \in \mathbb{N} \leftrightarrow n \in \omega \land n \neq \emptyset)$\lbl{Def$\mathbb{N}$}}
\AXC{$\varphi_{64}$\lbl{A2}}
\BIC{$\emptyset \in \mathbb{N} \leftrightarrow \emptyset \in \omega \land \emptyset \neq \emptyset$}
\AXC{$\varphi_{71}$\lbl{A1}}
\BIC{$\emptyset = \emptyset \to \emptyset \in \mathbb{N} \to 1_{\omega} \leq \emptyset$}
\BIC{$\emptyset \in \mathbb{N} \to 1_{\omega} \leq \emptyset$}
\end{ptree}
\noindent$\Pi_{28}$:
\begin{ptree}
\AXC{$\stackrel{\Pi_{27}}{\vdots}$}\noLine
\UIC{$\emptyset \in \mathbb{N} \to 1_{\omega} \leq \emptyset$}
\AXC{$\stackrel{\Pi_{22}}{\vdots}$}\noLine
\UIC{$\varphi_{63}$}
\AXC{$\varphi_{73}$\lbl{A2}}
\AXC{$\varphi_{73} \to \varphi_{63} \to (\emptyset \in \mathbb{N} \to 1_{\omega} \leq \emptyset) \to \varphi_{25}$\lbl{A1}}
\BIC{$\varphi_{63} \to (\emptyset \in \mathbb{N} \to 1_{\omega} \leq \emptyset) \to \varphi_{25}$}
\BIC{$(\emptyset \in \mathbb{N} \to 1_{\omega} \leq \emptyset) \to \varphi_{25}$}
\BIC{$\varphi_{25}$}
\end{ptree}
\noindent$\Pi_{29}$:
\begin{ptree}
\AXC{$\stackrel{\Pi_{26}}{\vdots}$}\noLine
\UIC{$\varphi_{36}$}
\AXC{$\stackrel{\Pi_{28}}{\vdots}$}\noLine
\UIC{$\varphi_{25}$}
\AXC{$\forall x\,(\mathrm{Ind}(x) \leftrightarrow \emptyset \in x \land \forall y \in x\,(S(y) \in x))$\lbl{DefInd}}
\AXC{$\varphi_{67}$\lbl{A2}}
\BIC{$\mathrm{Ind}(B_{1}) \leftrightarrow \emptyset \in B_{1} \land \forall y \in B_{1}\,(S(y) \in B_{1})$}
\AXC{$\varphi_{81}$\lbl{A1}}
\BIC{$\varphi_{25} \to \varphi_{36} \to \varphi_{76}$}
\BIC{$\varphi_{36} \to \varphi_{76}$}
\BIC{$\varphi_{76}$}
\end{ptree}
\noindent$\Pi_{30}$:
\begin{ptree}
\AXC{$\stackrel{\Pi_{2}}{\vdots}$}\noLine
\UIC{$\forall n \in B_{1}\,(S(n) \in B_{1}) \to \forall y \in B_{1}\,(S(y) \in B_{1})$}
\AXC{$\stackrel{\Pi_{3}}{\vdots}$}\noLine
\UIC{$\forall y \in B_{1}\,(S(y) \in B_{1}) \to \forall n \in B_{1}\,(S(n) \in B_{1})$}
\AXC{$\stackrel{\Pi_{29}}{\vdots}$}\noLine
\UIC{$\varphi_{76}$}
\BIC{$\varphi_{70}$}
\BIC{$\varphi_{22}$}
\end{ptree}
\noindent$\Pi_{31}$:
\begin{ptree}
\AXC{$\forall n\,(n \in \omega \leftrightarrow \forall x\,(\mathrm{Ind}(x) \to n \in x))$\lbl{Def$\omega$}}
\AXC{$\varphi_{32}$\lbl{A2}}
\BIC{$n \in \omega \leftrightarrow \forall x\,(\mathrm{Ind}(x) \to n \in x)$}
\AXC{$\varphi_{37}$\lbl{A1}}
\AXC{$\varphi_{37} \to (n \in \omega \leftrightarrow \forall x\,(\mathrm{Ind}(x) \to n \in x)) \to \varphi_{44}$\lbl{A1}}
\BIC{$(n \in \omega \leftrightarrow \forall x\,(\mathrm{Ind}(x) \to n \in x)) \to \varphi_{44}$}
\BIC{$\varphi_{44}$}
\end{ptree}
\noindent$\Pi_{32}$:
\begin{ptree}
\AXC{$\stackrel{\Pi_{30}}{\vdots}$}\noLine
\UIC{$\varphi_{22}$}
\AXC{$\forall x\,(\mathrm{Ind}(x) \to n \in x) \to \mathrm{Ind}(B_{1}) \to n \in B_{1}$\lbl{A2}}
\AXC{$\stackrel{\Pi_{31}}{\vdots}$}\noLine
\UIC{$\varphi_{44}$}
\AXC{$\varphi_{44} \to \varphi_{74}$\lbl{A1}}
\BIC{$\varphi_{74}$}
\BIC{$\varphi_{42}$}
\AXC{$\varphi_{42} \to \varphi_{22} \to \varphi_{52} \to \varphi_{43}$\lbl{A1}}
\BIC{$\varphi_{22} \to \varphi_{52} \to \varphi_{43}$}
\BIC{$\varphi_{52} \to \varphi_{43}$}
\end{ptree}
\noindent$\Pi_{33}$:
\begin{ptree}
\AXC{$\varphi_{58}$\lbl{A4}}
\AXC{$\forall n\,(\varphi_{43}) \to \varphi_{49}$\lbl{A3}}
\AXC{$\varphi_{52}$\lbl{A2}}
\AXC{$\stackrel{\Pi_{32}}{\vdots}$}\noLine
\UIC{$\varphi_{52} \to \varphi_{43}$}
\BIC{$\varphi_{43}$}
\UIC{$\forall n\,(\varphi_{43})$}
\AXC{$\forall n\,(\varphi_{43}) \to (\forall n\,(\varphi_{43}) \to \varphi_{49}) \to \varphi_{58} \to \varphi_{46}$\lbl{A1}}
\BIC{$(\forall n\,(\varphi_{43}) \to \varphi_{49}) \to \varphi_{58} \to \varphi_{46}$}
\BIC{$\varphi_{58} \to \varphi_{46}$}
\BIC{$\varphi_{46}$}
\AXC{$\varphi_{46} \to \varphi_{47}$\lbl{A1}}
\BIC{$\varphi_{47}$}
\UIC{$\forall B_{1}\,(\varphi_{47})$}
\end{ptree}
\noindent$\Pi_{34}$:
\begin{ptree}
\AXC{$\varphi_{40}$\lbl{A4}}
\AXC{$\varphi_{78}$\lbl{A3}}
\AXC{$\stackrel{\Pi_{33}}{\vdots}$}\noLine
\UIC{$\forall B_{1}\,(\varphi_{47})$}
\AXC{$\varphi_{9}$\lbl{Sep}}
\AXC{$\varphi_{66}$\lbl{A2}}
\BIC{$\exists B_{1}\,\forall n\,(n \in B_{1} \leftrightarrow n \in \omega \land (n \in \mathbb{N} \to 1_{\omega} \leq n))$}
\AXC{$\varphi_{82}$\lbl{A1}}
\BIC{$\forall B_{1}\,(\varphi_{47}) \to \varphi_{78} \to \varphi_{40} \to \forall n \in \omega\,(n \in \mathbb{N} \to 1_{\omega} \leq n)$}
\BIC{$\varphi_{78} \to \varphi_{40} \to \forall n \in \omega\,(n \in \mathbb{N} \to 1_{\omega} \leq n)$}
\BIC{$\varphi_{40} \to \forall n \in \omega\,(n \in \mathbb{N} \to 1_{\omega} \leq n)$}
\BIC{$\forall n \in \omega\,(n \in \mathbb{N} \to 1_{\omega} \leq n)$}
\AXC{$\varphi_{29}$\lbl{A2}}
\BIC{$n \in \omega \to n \in \mathbb{N} \to 1_{\omega} \leq n$}
\end{ptree}
\noindent 主木:
\begin{ptree}
\AXC{$n \in \mathbb{N} \to n \in \mathbb{N}$\lbl{A1}}
\AXC{$\stackrel{\Pi_{1}}{\vdots}$}\noLine
\UIC{$n \in \mathbb{N} \to n \in \omega$}
\AXC{$\stackrel{\Pi_{34}}{\vdots}$}\noLine
\UIC{$n \in \omega \to n \in \mathbb{N} \to 1_{\omega} \leq n$}
\AXC{$(n \in \omega \to n \in \mathbb{N} \to 1_{\omega} \leq n) \to \varphi_{28}$\lbl{A1}}
\BIC{$\varphi_{28}$}
\BIC{$(n \in \mathbb{N} \to n \in \mathbb{N}) \to n \in \mathbb{N} \to 1_{\omega} \leq n$}
\BIC{$n \in \mathbb{N} \to 1_{\omega} \leq n$}
\UIC{$\forall n \in \mathbb{N}\,(1_{\omega} \leq n)$}
\end{ptree}
\begin{abbrevlist}
\item $\varphi_{1} :\equiv \forall w\,(1_{\omega} = w \to 1_{\omega} = 1_{\omega} \to 1_{\omega} = w)$
\item $\varphi_{2} :\equiv 1_{\omega} = S(n) \to 1_{\omega} = 1_{\omega} \to 1_{\omega} = S(n)$
\item $\varphi_{3} :\equiv \forall w\,(1_{\omega} = w \to 1_{\omega} < 1_{\omega} \to 1_{\omega} < w)$
\item $\varphi_{4} :\equiv \forall x\,\forall y\,(x = y \to y = x) \to \forall y\,(S(n) = y \to y = S(n))$
\item $\varphi_{5} :\equiv \forall x\,\forall y\,(x = y \to y = x) \to \forall y\,(1_{\omega} = y \to y = 1_{\omega})$
\item $\varphi_{6} :\equiv 1_{\omega} = S(n) \to 1_{\omega} < 1_{\omega} \to 1_{\omega} < S(n)$
\item $\varphi_{7} :\equiv \forall y\,(S(n) = y \to y = S(n)) \to S(n) = 1_{\omega} \to 1_{\omega} = S(n)$
\item $\varphi_{8} :\equiv S(n) = S(\emptyset) \to S(\emptyset) = 1_{\omega} \to S(n) = 1_{\omega}$
\item $\varphi_{9} :\equiv \forall A_{2}\,\exists B_{1}\,\forall n\,(n \in B_{1} \leftrightarrow n \in A_{2} \land (n \in \mathbb{N} \to 1_{\omega} \leq n))$
\item $\varphi_{10} :\equiv \forall B_{1}\,(\neg \forall n\,(n \in B_{1} \leftrightarrow n \in \omega \land (n \in \mathbb{N} \to 1_{\omega} \leq n)))$
\item $\varphi_{11} :\equiv \forall y\,(\forall n \in B_{1}\,(S(n) \in B_{1}) \to y \in B_{1} \to S(y) \in B_{1})$
\item $\varphi_{12} :\equiv \forall n\,(\forall y \in B_{1}\,(S(y) \in B_{1}) \to n \in B_{1} \to S(n) \in B_{1})$
\item $\varphi_{13} :\equiv \forall x\,\forall y\,\forall z\,(x \leq y \to y \leq z \to x \leq z) \to \forall y\,\forall z\,(1_{\omega} \leq y \to y \leq z \to 1_{\omega} \leq z)$
\item $\varphi_{14} :\equiv \forall y\,\forall z\,(1_{\omega} \leq y \to y \leq z \to 1_{\omega} \leq z) \to \forall z\,(1_{\omega} \leq n \to n \leq z \to 1_{\omega} \leq z)$
\item $\varphi_{15} :\equiv \forall z\,(1_{\omega} \leq n \to n \leq z \to 1_{\omega} \leq z) \to 1_{\omega} \leq n \to n \leq S(n) \to 1_{\omega} \leq S(n)$
\item $\varphi_{16} :\equiv \forall x\,\forall y\,\forall z\,(x = y \to y = z \to x = z) \to \forall y\,\forall z\,(S(n) = y \to y = z \to S(n) = z)$
\item $\varphi_{17} :\equiv (1_{\omega} = S(\emptyset) \to S(\emptyset) = 1_{\omega}) \to S(\emptyset) = 1_{\omega}$
\item $\varphi_{18} :\equiv S(n) \in B_{1} \leftrightarrow S(n) \in \omega \land (S(n) \in \mathbb{N} \to 1_{\omega} \leq S(n))$
\item $\varphi_{19} :\equiv \forall n \in B_{1}\,(S(n) \in B_{1}) \to \forall y\,\forall n \in B_{1}\,(S(n) \in B_{1})$
\item $\varphi_{20} :\equiv \forall y \in B_{1}\,(S(y) \in B_{1}) \to \forall n\,\forall y \in B_{1}\,(S(y) \in B_{1})$
\item $\varphi_{21} :\equiv \forall y\,(1_{\omega} = y \to y = 1_{\omega}) \to 1_{\omega} = S(\emptyset) \to S(\emptyset) = 1_{\omega}$
\item $\varphi_{22} :\equiv \forall n\,(n \in B_{1} \leftrightarrow n \in \omega \land (n \in \mathbb{N} \to 1_{\omega} \leq n)) \to \mathrm{Ind}(B_{1})$
\item $\varphi_{23} :\equiv \forall y\,\forall n \in B_{1}\,(S(n) \in B_{1}) \to \forall y \in B_{1}\,(S(y) \in B_{1})$
\item $\varphi_{24} :\equiv \forall n\,\forall y \in B_{1}\,(S(y) \in B_{1}) \to \forall n \in B_{1}\,(S(n) \in B_{1})$
\item $\varphi_{25} :\equiv \forall n\,(n \in B_{1} \leftrightarrow n \in \omega \land (n \in \mathbb{N} \to 1_{\omega} \leq n)) \to \emptyset \in B_{1}$
\item $\varphi_{26} :\equiv \forall z\,\forall w\,(z = w \to S(n) = S(z) \to S(n) = S(w)) \to \forall w\,(n = w \to S(n) = S(n) \to S(n) = S(w))$
\item $\varphi_{27} :\equiv \forall n\,(n \in B_{1} \leftrightarrow n \in \omega \land (n \in \mathbb{N} \to 1_{\omega} \leq n)) \to n \in B_{1} \to n \in B_{1}$
\item $\varphi_{28} :\equiv (n \in \mathbb{N} \to n \in \omega) \to (n \in \mathbb{N} \to n \in \mathbb{N}) \to n \in \mathbb{N} \to 1_{\omega} \leq n$
\item $\varphi_{29} :\equiv \forall n \in \omega\,(n \in \mathbb{N} \to 1_{\omega} \leq n) \to n \in \omega \to n \in \mathbb{N} \to 1_{\omega} \leq n$
\item $\varphi_{30} :\equiv \forall n\,(n \in B_{1} \leftrightarrow n \in \omega \land (n \in \mathbb{N} \to 1_{\omega} \leq n)) \to n \in B_{1} \to n \in \omega$
\item $\varphi_{31} :\equiv \forall n\,(n \in B_{1} \leftrightarrow n \in \omega \land (n \in \mathbb{N} \to 1_{\omega} \leq n)) \to n \in B_{1} \to S(n) \in B_{1}$
\item $\varphi_{32} :\equiv \forall n\,(n \in \omega \leftrightarrow \forall x\,(\mathrm{Ind}(x) \to n \in x)) \to (n \in \omega \leftrightarrow \forall x\,(\mathrm{Ind}(x) \to n \in x))$
\item $\varphi_{33} :\equiv \forall y\,\forall z\,(S(n) = y \to y = z \to S(n) = z) \to \forall z\,(S(n) = S(\emptyset) \to S(\emptyset) = z \to S(n) = z)$
\item $\varphi_{34} :\equiv \forall w\,(n = w \to S(n) = S(n) \to S(n) = S(w)) \to n = \emptyset \to S(n) = S(n) \to S(n) = S(\emptyset)$
\item $\varphi_{35} :\equiv \forall n\,(n \in B_{1} \leftrightarrow n \in \omega \land (n \in \mathbb{N} \to 1_{\omega} \leq n)) \to n \in B_{1} \to S(n) \in \omega$
\item $\varphi_{36} :\equiv \forall n\,(n \in B_{1} \leftrightarrow n \in \omega \land (n \in \mathbb{N} \to 1_{\omega} \leq n)) \to \forall n \in B_{1}\,(S(n) \in B_{1})$
\item $\varphi_{37} :\equiv \forall n\,(n \in B_{1} \leftrightarrow n \in \omega \land (n \in \mathbb{N} \to 1_{\omega} \leq n)) \to n \in \omega \to n \in \omega$
\item $\varphi_{38} :\equiv \forall n\,\forall n\,(n \in B_{1} \leftrightarrow n \in \omega \land (n \in \mathbb{N} \to 1_{\omega} \leq n)) \to \forall n \in B_{1}\,(S(n) \in B_{1})$
\item $\varphi_{39} :\equiv \forall n\,(n \in B_{1} \leftrightarrow n \in \omega \land (n \in \mathbb{N} \to 1_{\omega} \leq n)) \to n \in B_{1} \to n \in \mathbb{N} \to 1_{\omega} \leq n$
\item $\varphi_{40} :\equiv \neg \forall n \in \omega\,(n \in \mathbb{N} \to 1_{\omega} \leq n) \to \forall B_{1}\,(\neg \forall n \in \omega\,(n \in \mathbb{N} \to 1_{\omega} \leq n))$
\item $\varphi_{41} :\equiv \forall n\,(n \in \mathbb{N} \leftrightarrow n \in \omega \land n \neq \emptyset) \to (n \in \mathbb{N} \leftrightarrow n \in \omega \land n \neq \emptyset)$
\item $\varphi_{42} :\equiv \forall n\,(n \in B_{1} \leftrightarrow n \in \omega \land (n \in \mathbb{N} \to 1_{\omega} \leq n)) \to n \in \omega \to \mathrm{Ind}(B_{1}) \to n \in B_{1}$
\item $\varphi_{43} :\equiv \forall n\,(n \in B_{1} \leftrightarrow n \in \omega \land (n \in \mathbb{N} \to 1_{\omega} \leq n)) \to n \in \omega \to n \in \mathbb{N} \to 1_{\omega} \leq n$
\item $\varphi_{44} :\equiv \forall n\,(n \in B_{1} \leftrightarrow n \in \omega \land (n \in \mathbb{N} \to 1_{\omega} \leq n)) \to n \in \omega \to \forall x\,(\mathrm{Ind}(x) \to n \in x)$
\item $\varphi_{45} :\equiv \forall n\,(n \in B_{1} \leftrightarrow n \in \omega \land (n \in \mathbb{N} \to 1_{\omega} \leq n)) \to n \in B_{1} \to S(n) \in \mathbb{N} \to 1_{\omega} \leq S(n)$
\item $\varphi_{46} :\equiv \forall n\,(n \in B_{1} \leftrightarrow n \in \omega \land (n \in \mathbb{N} \to 1_{\omega} \leq n)) \to \forall n \in \omega\,(n \in \mathbb{N} \to 1_{\omega} \leq n)$
\item $\varphi_{47} :\equiv \neg \forall n \in \omega\,(n \in \mathbb{N} \to 1_{\omega} \leq n) \to \neg \forall n\,(n \in B_{1} \leftrightarrow n \in \omega \land (n \in \mathbb{N} \to 1_{\omega} \leq n))$
\item $\varphi_{48} :\equiv (n \in \mathbb{N} \to n \in \mathbb{N}) \to (n \in \mathbb{N} \leftrightarrow n \in \omega \land n \neq \emptyset) \to n \in \mathbb{N} \to n \in \omega$
\item $\varphi_{49} :\equiv \forall n\,\forall n\,(n \in B_{1} \leftrightarrow n \in \omega \land (n \in \mathbb{N} \to 1_{\omega} \leq n)) \to \forall n \in \omega\,(n \in \mathbb{N} \to 1_{\omega} \leq n)$
\item $\varphi_{50} :\equiv \forall n\,(n \in B_{1} \leftrightarrow n \in \omega \land (n \in \mathbb{N} \to 1_{\omega} \leq n)) \to n \in B_{1} \to n \in \omega \to S(n) \in \omega$
\item $\varphi_{51} :\equiv \forall n\,(n \in B_{1} \leftrightarrow n \in \omega \land (n \in \mathbb{N} \to 1_{\omega} \leq n)) \to n \in B_{1} \to \forall y \in \omega\,(S(y) \in \omega)$
\item $\varphi_{52} :\equiv \forall n\,(n \in B_{1} \leftrightarrow n \in \omega \land (n \in \mathbb{N} \to 1_{\omega} \leq n)) \to (n \in B_{1} \leftrightarrow n \in \omega \land (n \in \mathbb{N} \to 1_{\omega} \leq n))$
\item $\varphi_{53} :\equiv \forall n\,(n \in B_{1} \leftrightarrow n \in \omega \land (n \in \mathbb{N} \to 1_{\omega} \leq n)) \to n \in B_{1} \to n \in \omega \land (n \in \mathbb{N} \to 1_{\omega} \leq n)$
\item $\varphi_{54} :\equiv \forall n\,(n \in B_{1} \leftrightarrow n \in \omega \land (n \in \mathbb{N} \to 1_{\omega} \leq n)) \to n \in B_{1} \to S(n) \in \mathbb{N} \to n \neq \emptyset \to 1_{\omega} \leq n$
\item $\varphi_{55} :\equiv \forall n\,(n \in B_{1} \leftrightarrow n \in \omega \land (n \in \mathbb{N} \to 1_{\omega} \leq n)) \to n \in B_{1} \to S(n) \in \mathbb{N} \to n = \emptyset \to S(n) = 1_{\omega}$
\item $\varphi_{56} :\equiv \forall n\,(n \in B_{1} \leftrightarrow n \in \omega \land (n \in \mathbb{N} \to 1_{\omega} \leq n)) \to n \in B_{1} \to S(n) \in \mathbb{N} \to n = \emptyset \to 1_{\omega} = S(n)$
\item $\varphi_{57} :\equiv \forall n\,(n \in B_{1} \leftrightarrow n \in \omega \land (n \in \mathbb{N} \to 1_{\omega} \leq n)) \to n \in B_{1} \to S(n) \in \mathbb{N} \to n = \emptyset \to 1_{\omega} \leq S(n)$
\item $\varphi_{58} :\equiv \forall n\,(n \in B_{1} \leftrightarrow n \in \omega \land (n \in \mathbb{N} \to 1_{\omega} \leq n)) \to \forall n\,\forall n\,(n \in B_{1} \leftrightarrow n \in \omega \land (n \in \mathbb{N} \to 1_{\omega} \leq n))$
\item $\varphi_{59} :\equiv \forall n\,(n \in B_{1} \leftrightarrow n \in \omega \land (n \in \mathbb{N} \to 1_{\omega} \leq n)) \to n \in B_{1} \to S(n) \in \mathbb{N} \to n \neq \emptyset \to n \in \mathbb{N}$
\item $\varphi_{60} :\equiv \forall n\,(n \in B_{1} \leftrightarrow n \in \omega \land (n \in \mathbb{N} \to 1_{\omega} \leq n)) \to n \in B_{1} \to S(n) \in \mathbb{N} \to n \neq \emptyset \to 1_{\omega} \leq S(n)$
\item $\varphi_{61} :\equiv \forall n\,(n \in B_{1} \leftrightarrow n \in \omega \land (n \in \mathbb{N} \to 1_{\omega} \leq n)) \to n \in B_{1} \to S(n) \in \mathbb{N} \to n = \emptyset \to n = \emptyset$
\item $\varphi_{62} :\equiv \forall n\,(n \in B_{1} \leftrightarrow n \in \omega \land (n \in \mathbb{N} \to 1_{\omega} \leq n)) \to n \in B_{1} \to S(n) \in \mathbb{N} \to n \neq \emptyset \to n \neq \emptyset$
\item $\varphi_{63} :\equiv \forall n\,(n \in B_{1} \leftrightarrow n \in \omega \land (n \in \mathbb{N} \to 1_{\omega} \leq n)) \to \emptyset \in \omega \land \forall y \in \omega\,(S(y) \in \omega)$
\item $\varphi_{64} :\equiv \forall n\,(n \in \mathbb{N} \leftrightarrow n \in \omega \land n \neq \emptyset) \to (\emptyset \in \mathbb{N} \leftrightarrow \emptyset \in \omega \land \emptyset \neq \emptyset)$
\item $\varphi_{65} :\equiv \forall n\,(n \in B_{1} \leftrightarrow n \in \omega \land (n \in \mathbb{N} \to 1_{\omega} \leq n)) \to n \in B_{1} \to S(n) \in \mathbb{N} \to n = \emptyset \to S(n) = S(\emptyset)$
\item $\varphi_{66} :\equiv \varphi_{9} \to \exists B_{1}\,\forall n\,(n \in B_{1} \leftrightarrow n \in \omega \land (n \in \mathbb{N} \to 1_{\omega} \leq n))$
\item $\varphi_{67} :\equiv \forall x\,(\mathrm{Ind}(x) \leftrightarrow \emptyset \in x \land \forall y \in x\,(S(y) \in x)) \to (\mathrm{Ind}(B_{1}) \leftrightarrow \emptyset \in B_{1} \land \forall y \in B_{1}\,(S(y) \in B_{1}))$
\item $\varphi_{68} :\equiv \varphi_{19} \to \forall n \in B_{1}\,(S(n) \in B_{1}) \to \forall y \in B_{1}\,(S(y) \in B_{1})$
\item $\varphi_{69} :\equiv \varphi_{20} \to \forall y \in B_{1}\,(S(y) \in B_{1}) \to \forall n \in B_{1}\,(S(n) \in B_{1})$
\item $\varphi_{70} :\equiv (\forall n \in B_{1}\,(S(n) \in B_{1}) \to \forall y \in B_{1}\,(S(y) \in B_{1})) \to \varphi_{22}$
\item $\varphi_{71} :\equiv (\emptyset \in \mathbb{N} \leftrightarrow \emptyset \in \omega \land \emptyset \neq \emptyset) \to \emptyset = \emptyset \to \emptyset \in \mathbb{N} \to 1_{\omega} \leq \emptyset$
\item $\varphi_{72} :\equiv \forall x\,(\mathrm{Ind}(x) \leftrightarrow \emptyset \in x \land \forall y \in x\,(S(y) \in x)) \to (\mathrm{Ind}(\omega) \leftrightarrow \emptyset \in \omega \land \forall y \in \omega\,(S(y) \in \omega))$
\item $\varphi_{73} :\equiv \forall n\,(n \in B_{1} \leftrightarrow n \in \omega \land (n \in \mathbb{N} \to 1_{\omega} \leq n)) \to (\emptyset \in B_{1} \leftrightarrow \emptyset \in \omega \land (\emptyset \in \mathbb{N} \to 1_{\omega} \leq \emptyset))$
\item $\varphi_{74} :\equiv (\forall x\,(\mathrm{Ind}(x) \to n \in x) \to \mathrm{Ind}(B_{1}) \to n \in B_{1}) \to \varphi_{42}$
\item $\varphi_{75} :\equiv S(n) = S(n) \to (n = \emptyset \to S(n) = S(n) \to S(n) = S(\emptyset)) \to \varphi_{65}$
\item $\varphi_{76} :\equiv (\forall y \in B_{1}\,(S(y) \in B_{1}) \to \forall n \in B_{1}\,(S(n) \in B_{1})) \to \varphi_{70}$
\item $\varphi_{77} :\equiv \mathrm{Ind}(\omega) \to (\mathrm{Ind}(\omega) \leftrightarrow \emptyset \in \omega \land \forall y \in \omega\,(S(y) \in \omega)) \to \varphi_{63}$
\item $\varphi_{78} :\equiv \forall B_{1}\,(\varphi_{47}) \to \forall B_{1}\,(\neg \forall n \in \omega\,(n \in \mathbb{N} \to 1_{\omega} \leq n)) \to \varphi_{10}$
\item $\varphi_{79} :\equiv (1_{\omega} \leq n \to n \leq S(n) \to 1_{\omega} \leq S(n)) \to \varphi_{54} \to n \leq S(n) \to \varphi_{60}$
\item $\varphi_{80} :\equiv (\forall n\,(n \in B_{1} \leftrightarrow n \in \omega \land (n \in \mathbb{N} \to 1_{\omega} \leq n)) \to \varphi_{18}) \to \varphi_{35} \to \varphi_{45} \to \varphi_{31}$
\item $\varphi_{81} :\equiv (\mathrm{Ind}(B_{1}) \leftrightarrow \emptyset \in B_{1} \land \forall y \in B_{1}\,(S(y) \in B_{1})) \to \varphi_{25} \to \varphi_{36} \to \varphi_{76}$
\item $\varphi_{82} :\equiv \exists B_{1}\,\forall n\,(n \in B_{1} \leftrightarrow n \in \omega \land (n \in \mathbb{N} \to 1_{\omega} \leq n)) \to \forall B_{1}\,(\varphi_{47}) \to \varphi_{78} \to \varphi_{40} \to \forall n \in \omega\,(n \in \mathbb{N} \to 1_{\omega} \leq n)$
\end{abbrevlist}
\end{fullproof}

% 自動生成 (prover/build.js): LtAddOne  ノード数 79
\begin{fullproof}{[N13]}{$\forall x \in \mathbb{R}\,(x < x + 1)$}{79}\phantomsection\label{tree:LtAddOne}
\noindent$\Pi_{1}$:
\begin{ptree}
\AXC{$z = w \to z < 1 + x \to w < 1 + x$\lbl{A6}}
\UIC{$\forall w\,(z = w \to z < 1 + x \to w < 1 + x)$}
\UIC{$\forall z\,\forall w\,(z = w \to z < 1 + x \to w < 1 + x)$}
\AXC{$\varphi_{12}$\lbl{A2}}
\BIC{$\forall w\,(0 + x = w \to 0 + x < 1 + x \to w < 1 + x)$}
\AXC{$\varphi_{15}$\lbl{A2}}
\BIC{$0 + x = x \to 0 + x < 1 + x \to x < 1 + x$}
\end{ptree}
\noindent$\Pi_{2}$:
\begin{ptree}
\AXC{$1 \in \mathbb{R}$\lbl{(R1)}}
\AXC{$0 \in \mathbb{R}$\lbl{(R1)}}
\AXC{$\varphi_{8}$\lbl{(R14)}}
\AXC{$\varphi_{8} \to \varphi_{6}$\lbl{A2}}
\BIC{$\varphi_{6}$}
\AXC{$\varphi_{21}$\lbl{A2}}
\BIC{$\forall z\,(0 \in \mathbb{R} \to 1 \in \mathbb{R} \to z \in \mathbb{R} \to 0 < 1 \to 0 + z < 1 + z)$}
\AXC{$\varphi_{20}$\lbl{A2}}
\BIC{$0 \in \mathbb{R} \to 1 \in \mathbb{R} \to x \in \mathbb{R} \to 0 < 1 \to 0 + x < 1 + x$}
\AXC{$\varphi_{22}$\lbl{A1}}
\BIC{$0 \in \mathbb{R} \to \varphi_{10}$}
\BIC{$\varphi_{10}$}
\BIC{$(x \in \mathbb{R} \to x \in \mathbb{R}) \to 0 < 1 \to x \in \mathbb{R} \to 0 + x < 1 + x$}
\end{ptree}
\noindent$\Pi_{3}$:
\begin{ptree}
\AXC{$0 < 1$\lbl{[F28]}}
\AXC{$x \in \mathbb{R} \to x \in \mathbb{R}$\lbl{A1}}
\AXC{$\stackrel{\Pi_{2}}{\vdots}$}\noLine
\UIC{$(x \in \mathbb{R} \to x \in \mathbb{R}) \to 0 < 1 \to x \in \mathbb{R} \to 0 + x < 1 + x$}
\BIC{$0 < 1 \to x \in \mathbb{R} \to 0 + x < 1 + x$}
\BIC{$x \in \mathbb{R} \to 0 + x < 1 + x$}
\end{ptree}
\noindent$\Pi_{4}$:
\begin{ptree}
\AXC{$y = z \to x = y \to x = z$\lbl{A6}}
\AXC{$(y = z \to x = y \to x = z) \to x = y \to y = z \to x = z$\lbl{A1}}
\BIC{$x = y \to y = z \to x = z$}
\UIC{$\forall z\,(x = y \to y = z \to x = z)$}
\UIC{$\forall y\,\forall z\,(x = y \to y = z \to x = z)$}
\UIC{$\forall x\,\forall y\,\forall z\,(x = y \to y = z \to x = z)$}
\AXC{$\varphi_{1}$\lbl{A2}}
\BIC{$\forall y\,\forall z\,(0 + x = y \to y = z \to 0 + x = z)$}
\AXC{$\varphi_{7}$\lbl{A2}}
\BIC{$\forall z\,(0 + x = x + 0 \to x + 0 = z \to 0 + x = z)$}
\AXC{$\varphi_{9}$\lbl{A2}}
\BIC{$0 + x = x + 0 \to x + 0 = x \to 0 + x = x$}
\end{ptree}
\noindent$\Pi_{5}$:
\begin{ptree}
\AXC{$0 \in \mathbb{R}$\lbl{(R1)}}
\AXC{$\forall x\,\forall y\,(x \in \mathbb{R} \to y \in \mathbb{R} \to x + y = y + x)$\lbl{(R3)}}
\AXC{$\varphi_{16}$\lbl{A2}}
\BIC{$\forall y\,(0 \in \mathbb{R} \to y \in \mathbb{R} \to 0 + y = y + 0)$}
\AXC{$\varphi_{13}$\lbl{A2}}
\BIC{$0 \in \mathbb{R} \to x \in \mathbb{R} \to 0 + x = x + 0$}
\AXC{$\varphi_{18}$\lbl{A1}}
\BIC{$0 \in \mathbb{R} \to (x \in \mathbb{R} \to x \in \mathbb{R}) \to x \in \mathbb{R} \to 0 + x = x + 0$}
\BIC{$(x \in \mathbb{R} \to x \in \mathbb{R}) \to x \in \mathbb{R} \to 0 + x = x + 0$}
\end{ptree}
\noindent$\Pi_{6}$:
\begin{ptree}
\AXC{$x \in \mathbb{R} \to x \in \mathbb{R}$\lbl{A1}}
\AXC{$\stackrel{\Pi_{5}}{\vdots}$}\noLine
\UIC{$(x \in \mathbb{R} \to x \in \mathbb{R}) \to x \in \mathbb{R} \to 0 + x = x + 0$}
\BIC{$x \in \mathbb{R} \to 0 + x = x + 0$}
\end{ptree}
\noindent$\Pi_{7}$:
\begin{ptree}
\AXC{$\stackrel{\Pi_{6}}{\vdots}$}\noLine
\UIC{$x \in \mathbb{R} \to 0 + x = x + 0$}
\AXC{$(x \in \mathbb{R} \to 0 + x = x + 0) \to (x \in \mathbb{R} \to x + 0 = x) \to \varphi_{2}$\lbl{A1}}
\BIC{$(x \in \mathbb{R} \to x + 0 = x) \to \varphi_{2}$}
\end{ptree}
\noindent$\Pi_{8}$:
\begin{ptree}
\AXC{$\forall x \in \mathbb{R}\,(x + 0 = x)$\lbl{(R4)}}
\AXC{$\forall x \in \mathbb{R}\,(x + 0 = x) \to x \in \mathbb{R} \to x + 0 = x$\lbl{A2}}
\BIC{$x \in \mathbb{R} \to x + 0 = x$}
\AXC{$\stackrel{\Pi_{7}}{\vdots}$}\noLine
\UIC{$(x \in \mathbb{R} \to x + 0 = x) \to \varphi_{2}$}
\BIC{$\varphi_{2}$}
\end{ptree}
\noindent$\Pi_{9}$:
\begin{ptree}
\AXC{$\stackrel{\Pi_{4}}{\vdots}$}\noLine
\UIC{$0 + x = x + 0 \to x + 0 = x \to 0 + x = x$}
\AXC{$\stackrel{\Pi_{8}}{\vdots}$}\noLine
\UIC{$\varphi_{2}$}
\BIC{$x \in \mathbb{R} \to 0 + x = x$}
\AXC{$(x \in \mathbb{R} \to 0 + x = x) \to (x \in \mathbb{R} \to 0 + x < 1 + x) \to \varphi_{4}$\lbl{A1}}
\BIC{$(x \in \mathbb{R} \to 0 + x < 1 + x) \to \varphi_{4}$}
\end{ptree}
\noindent$\Pi_{10}$:
\begin{ptree}
\AXC{$\stackrel{\Pi_{1}}{\vdots}$}\noLine
\UIC{$0 + x = x \to 0 + x < 1 + x \to x < 1 + x$}
\AXC{$\stackrel{\Pi_{3}}{\vdots}$}\noLine
\UIC{$x \in \mathbb{R} \to 0 + x < 1 + x$}
\AXC{$\stackrel{\Pi_{9}}{\vdots}$}\noLine
\UIC{$(x \in \mathbb{R} \to 0 + x < 1 + x) \to \varphi_{4}$}
\BIC{$\varphi_{4}$}
\BIC{$x \in \mathbb{R} \to x < 1 + x$}
\end{ptree}
\noindent$\Pi_{11}$:
\begin{ptree}
\AXC{$1 \in \mathbb{R}$\lbl{(R1)}}
\AXC{$\forall x\,\forall y\,(x \in \mathbb{R} \to y \in \mathbb{R} \to x + y = y + x)$\lbl{(R3)}}
\AXC{$\varphi_{17}$\lbl{A2}}
\BIC{$\forall y\,(1 \in \mathbb{R} \to y \in \mathbb{R} \to 1 + y = y + 1)$}
\AXC{$\varphi_{14}$\lbl{A2}}
\BIC{$1 \in \mathbb{R} \to x \in \mathbb{R} \to 1 + x = x + 1$}
\AXC{$\varphi_{19}$\lbl{A1}}
\BIC{$1 \in \mathbb{R} \to (x \in \mathbb{R} \to x \in \mathbb{R}) \to x \in \mathbb{R} \to 1 + x = x + 1$}
\BIC{$(x \in \mathbb{R} \to x \in \mathbb{R}) \to x \in \mathbb{R} \to 1 + x = x + 1$}
\end{ptree}
\noindent$\Pi_{12}$:
\begin{ptree}
\AXC{$x \in \mathbb{R} \to x \in \mathbb{R}$\lbl{A1}}
\AXC{$\stackrel{\Pi_{11}}{\vdots}$}\noLine
\UIC{$(x \in \mathbb{R} \to x \in \mathbb{R}) \to x \in \mathbb{R} \to 1 + x = x + 1$}
\BIC{$x \in \mathbb{R} \to 1 + x = x + 1$}
\end{ptree}
\noindent$\Pi_{13}$:
\begin{ptree}
\AXC{$\stackrel{\Pi_{10}}{\vdots}$}\noLine
\UIC{$x \in \mathbb{R} \to x < 1 + x$}
\AXC{$\stackrel{\Pi_{12}}{\vdots}$}\noLine
\UIC{$x \in \mathbb{R} \to 1 + x = x + 1$}
\AXC{$(x \in \mathbb{R} \to 1 + x = x + 1) \to (x \in \mathbb{R} \to x < 1 + x) \to \varphi_{5}$\lbl{A1}}
\BIC{$(x \in \mathbb{R} \to x < 1 + x) \to \varphi_{5}$}
\BIC{$\varphi_{5}$}
\end{ptree}
\noindent 主木:
\begin{ptree}
\AXC{$z = w \to x < z \to x < w$\lbl{A6}}
\UIC{$\forall w\,(z = w \to x < z \to x < w)$}
\UIC{$\forall z\,\forall w\,(z = w \to x < z \to x < w)$}
\AXC{$\varphi_{3}$\lbl{A2}}
\BIC{$\forall w\,(1 + x = w \to x < 1 + x \to x < w)$}
\AXC{$\varphi_{11}$\lbl{A2}}
\BIC{$1 + x = x + 1 \to x < 1 + x \to x < x + 1$}
\AXC{$\stackrel{\Pi_{13}}{\vdots}$}\noLine
\UIC{$\varphi_{5}$}
\BIC{$x \in \mathbb{R} \to x < x + 1$}
\UIC{$\forall x \in \mathbb{R}\,(x < x + 1)$}
\end{ptree}
\begin{abbrevlist}
\item $\varphi_{1} :\equiv \forall x\,\forall y\,\forall z\,(x = y \to y = z \to x = z) \to \forall y\,\forall z\,(0 + x = y \to y = z \to 0 + x = z)$
\item $\varphi_{2} :\equiv (0 + x = x + 0 \to x + 0 = x \to 0 + x = x) \to x \in \mathbb{R} \to 0 + x = x$
\item $\varphi_{3} :\equiv \forall z\,\forall w\,(z = w \to x < z \to x < w) \to \forall w\,(1 + x = w \to x < 1 + x \to x < w)$
\item $\varphi_{4} :\equiv (0 + x = x \to 0 + x < 1 + x \to x < 1 + x) \to x \in \mathbb{R} \to x < 1 + x$
\item $\varphi_{5} :\equiv (1 + x = x + 1 \to x < 1 + x \to x < x + 1) \to x \in \mathbb{R} \to x < x + 1$
\item $\varphi_{6} :\equiv \forall y\,\forall z\,(0 \in \mathbb{R} \to y \in \mathbb{R} \to z \in \mathbb{R} \to 0 < y \to 0 + z < y + z)$
\item $\varphi_{7} :\equiv \forall y\,\forall z\,(0 + x = y \to y = z \to 0 + x = z) \to \forall z\,(0 + x = x + 0 \to x + 0 = z \to 0 + x = z)$
\item $\varphi_{8} :\equiv \forall x\,\forall y\,\forall z\,(x \in \mathbb{R} \to y \in \mathbb{R} \to z \in \mathbb{R} \to x < y \to x + z < y + z)$
\item $\varphi_{9} :\equiv \forall z\,(0 + x = x + 0 \to x + 0 = z \to 0 + x = z) \to 0 + x = x + 0 \to x + 0 = x \to 0 + x = x$
\item $\varphi_{10} :\equiv 1 \in \mathbb{R} \to (x \in \mathbb{R} \to x \in \mathbb{R}) \to 0 < 1 \to x \in \mathbb{R} \to 0 + x < 1 + x$
\item $\varphi_{11} :\equiv \forall w\,(1 + x = w \to x < 1 + x \to x < w) \to 1 + x = x + 1 \to x < 1 + x \to x < x + 1$
\item $\varphi_{12} :\equiv \forall z\,\forall w\,(z = w \to z < 1 + x \to w < 1 + x) \to \forall w\,(0 + x = w \to 0 + x < 1 + x \to w < 1 + x)$
\item $\varphi_{13} :\equiv \forall y\,(0 \in \mathbb{R} \to y \in \mathbb{R} \to 0 + y = y + 0) \to 0 \in \mathbb{R} \to x \in \mathbb{R} \to 0 + x = x + 0$
\item $\varphi_{14} :\equiv \forall y\,(1 \in \mathbb{R} \to y \in \mathbb{R} \to 1 + y = y + 1) \to 1 \in \mathbb{R} \to x \in \mathbb{R} \to 1 + x = x + 1$
\item $\varphi_{15} :\equiv \forall w\,(0 + x = w \to 0 + x < 1 + x \to w < 1 + x) \to 0 + x = x \to 0 + x < 1 + x \to x < 1 + x$
\item $\varphi_{16} :\equiv \forall x\,\forall y\,(x \in \mathbb{R} \to y \in \mathbb{R} \to x + y = y + x) \to \forall y\,(0 \in \mathbb{R} \to y \in \mathbb{R} \to 0 + y = y + 0)$
\item $\varphi_{17} :\equiv \forall x\,\forall y\,(x \in \mathbb{R} \to y \in \mathbb{R} \to x + y = y + x) \to \forall y\,(1 \in \mathbb{R} \to y \in \mathbb{R} \to 1 + y = y + 1)$
\item $\varphi_{18} :\equiv (0 \in \mathbb{R} \to x \in \mathbb{R} \to 0 + x = x + 0) \to 0 \in \mathbb{R} \to (x \in \mathbb{R} \to x \in \mathbb{R}) \to x \in \mathbb{R} \to 0 + x = x + 0$
\item $\varphi_{19} :\equiv (1 \in \mathbb{R} \to x \in \mathbb{R} \to 1 + x = x + 1) \to 1 \in \mathbb{R} \to (x \in \mathbb{R} \to x \in \mathbb{R}) \to x \in \mathbb{R} \to 1 + x = x + 1$
\item $\varphi_{20} :\equiv \forall z\,(0 \in \mathbb{R} \to 1 \in \mathbb{R} \to z \in \mathbb{R} \to 0 < 1 \to 0 + z < 1 + z) \to 0 \in \mathbb{R} \to 1 \in \mathbb{R} \to x \in \mathbb{R} \to 0 < 1 \to 0 + x < 1 + x$
\item $\varphi_{21} :\equiv \varphi_{6} \to \forall z\,(0 \in \mathbb{R} \to 1 \in \mathbb{R} \to z \in \mathbb{R} \to 0 < 1 \to 0 + z < 1 + z)$
\item $\varphi_{22} :\equiv (0 \in \mathbb{R} \to 1 \in \mathbb{R} \to x \in \mathbb{R} \to 0 < 1 \to 0 + x < 1 + x) \to 0 \in \mathbb{R} \to \varphi_{10}$
\end{abbrevlist}
\end{fullproof}

% 自動生成 (prover/build.js): IotaGeOne  ノード数 271
\begin{fullproof}{[N14]}{$\forall n \in \mathbb{N}\,(1 \leq \underline{n})$}{271}\phantomsection\label{tree:IotaGeOne}
\noindent$\Pi_{1}$:
\begin{ptree}
\AXC{$\varphi_{12}$\lbl{A4}}
\AXC{$\varphi_{7} \to \varphi_{15}$\lbl{A3}}
\AXC{$\forall n \in B_{3}\,(S(n) \in B_{3}) \to y \in B_{3} \to S(y) \in B_{3}$\lbl{A2}}
\UIC{$\varphi_{7}$}
\AXC{$\varphi_{7} \to (\varphi_{7} \to \varphi_{15}) \to \varphi_{77}$\lbl{A1}}
\BIC{$(\varphi_{7} \to \varphi_{15}) \to \varphi_{77}$}
\BIC{$\varphi_{77}$}
\BIC{$\forall n \in B_{3}\,(S(n) \in B_{3}) \to \forall y \in B_{3}\,(S(y) \in B_{3})$}
\end{ptree}
\noindent$\Pi_{2}$:
\begin{ptree}
\AXC{$\varphi_{13}$\lbl{A4}}
\AXC{$\varphi_{8} \to \varphi_{16}$\lbl{A3}}
\AXC{$\forall y \in B_{3}\,(S(y) \in B_{3}) \to n \in B_{3} \to S(n) \in B_{3}$\lbl{A2}}
\UIC{$\varphi_{8}$}
\AXC{$\varphi_{8} \to (\varphi_{8} \to \varphi_{16}) \to \varphi_{78}$\lbl{A1}}
\BIC{$(\varphi_{8} \to \varphi_{16}) \to \varphi_{78}$}
\BIC{$\varphi_{78}$}
\BIC{$\forall y \in B_{3}\,(S(y) \in B_{3}) \to \forall n \in B_{3}\,(S(n) \in B_{3})$}
\end{ptree}
\noindent$\Pi_{3}$:
\begin{ptree}
\AXC{$\varphi_{24}$\lbl{A1}}
\AXC{$\varphi_{56}$\lbl{A2}}
\AXC{$\varphi_{56} \to \varphi_{24} \to \varphi_{57}$\lbl{A1}}
\BIC{$\varphi_{24} \to \varphi_{57}$}
\BIC{$\varphi_{57}$}
\AXC{$\varphi_{57} \to \varphi_{28}$\lbl{A1}}
\BIC{$\varphi_{28}$}
\end{ptree}
\noindent$\Pi_{4}$:
\begin{ptree}
\AXC{$\varphi_{65}$\lbl{A1}}
\AXC{$\stackrel{\Pi_{3}}{\vdots}$}\noLine
\UIC{$\varphi_{28}$}
\AXC{$\forall n\,(n \in \mathbb{N} \leftrightarrow n \in \omega \land n \neq \emptyset)$\lbl{Def$\mathbb{N}$}}
\AXC{$\varphi_{38}$\lbl{A2}}
\BIC{$n \in \mathbb{N} \leftrightarrow n \in \omega \land n \neq \emptyset$}
\AXC{$(n \in \mathbb{N} \leftrightarrow n \in \omega \land n \neq \emptyset) \to \varphi_{28} \to \varphi_{65} \to \varphi_{61}$\lbl{A1}}
\BIC{$\varphi_{28} \to \varphi_{65} \to \varphi_{61}$}
\BIC{$\varphi_{65} \to \varphi_{61}$}
\BIC{$\varphi_{61}$}
\end{ptree}
\noindent$\Pi_{5}$:
\begin{ptree}
\AXC{$\stackrel{\Pi_{4}}{\vdots}$}\noLine
\UIC{$\varphi_{61}$}
\AXC{$\forall n \in \mathbb{N}\,(\underline{n} \in \mathbb{R})$\lbl{[N4]}}
\AXC{$\forall n \in \mathbb{N}\,(\underline{n} \in \mathbb{R}) \to n \in \mathbb{N} \to \underline{n} \in \mathbb{R}$\lbl{A2}}
\BIC{$n \in \mathbb{N} \to \underline{n} \in \mathbb{R}$}
\AXC{$(n \in \mathbb{N} \to \underline{n} \in \mathbb{R}) \to \varphi_{61} \to \varphi_{73}$\lbl{A1}}
\BIC{$\varphi_{61} \to \varphi_{73}$}
\BIC{$\varphi_{73}$}
\end{ptree}
\noindent$\Pi_{6}$:
\begin{ptree}
\AXC{$\forall x \in \mathbb{R}\,(x < x + 1)$\lbl{[N13]}}
\AXC{$\forall x \in \mathbb{R}\,(x < x + 1) \to \underline{n} \in \mathbb{R} \to \underline{n} < \underline{n} + 1$\lbl{A2}}
\BIC{$\underline{n} \in \mathbb{R} \to \underline{n} < \underline{n} + 1$}
\AXC{$(\underline{n} \in \mathbb{R} \to \underline{n} < \underline{n} + 1) \to \varphi_{73} \to \varphi_{79}$\lbl{A1}}
\BIC{$\varphi_{73} \to \varphi_{79}$}
\end{ptree}
\noindent$\Pi_{7}$:
\begin{ptree}
\AXC{$\stackrel{\Pi_{4}}{\vdots}$}\noLine
\UIC{$\varphi_{61}$}
\AXC{$\varphi_{24}$\lbl{A1}}
\AXC{$\varphi_{56}$\lbl{A2}}
\AXC{$\varphi_{56} \to \varphi_{24} \to \varphi_{57}$\lbl{A1}}
\BIC{$\varphi_{24} \to \varphi_{57}$}
\BIC{$\varphi_{57}$}
\AXC{$\varphi_{57} \to \varphi_{41}$\lbl{A1}}
\BIC{$\varphi_{41}$}
\AXC{$\varphi_{41} \to \varphi_{61} \to \varphi_{62}$\lbl{A1}}
\BIC{$\varphi_{61} \to \varphi_{62}$}
\BIC{$\varphi_{62}$}
\end{ptree}
\noindent$\Pi_{8}$:
\begin{ptree}
\AXC{$\forall x\,\forall y\,\forall z\,(x \leq y \to y < z \to x < z)$\lbl{[O4]}}
\AXC{$\varphi_{20}$\lbl{A2}}
\BIC{$\forall y\,\forall z\,(1 \leq y \to y < z \to 1 < z)$}
\AXC{$\varphi_{22}$\lbl{A2}}
\BIC{$\forall z\,(1 \leq \underline{n} \to \underline{n} < z \to 1 < z)$}
\AXC{$\varphi_{54}$\lbl{A2}}
\BIC{$1 \leq \underline{n} \to \underline{n} < \underline{n} + 1 \to 1 < \underline{n} + 1$}
\end{ptree}
\noindent$\Pi_{9}$:
\begin{ptree}
\AXC{$\stackrel{\Pi_{7}}{\vdots}$}\noLine
\UIC{$\varphi_{62}$}
\AXC{$\stackrel{\Pi_{8}}{\vdots}$}\noLine
\UIC{$1 \leq \underline{n} \to \underline{n} < \underline{n} + 1 \to 1 < \underline{n} + 1$}
\AXC{$(1 \leq \underline{n} \to \underline{n} < \underline{n} + 1 \to 1 < \underline{n} + 1) \to \varphi_{62} \to \varphi_{79} \to \varphi_{75}$\lbl{A1}}
\BIC{$\varphi_{62} \to \varphi_{79} \to \varphi_{75}$}
\BIC{$\varphi_{79} \to \varphi_{75}$}
\end{ptree}
\noindent$\Pi_{10}$:
\begin{ptree}
\AXC{$x = x$\lbl{A5}}
\AXC{$x = y \to x = x \to y = x$\lbl{A6}}
\AXC{$(x = y \to x = x \to y = x) \to x = x \to x = y \to y = x$\lbl{A1}}
\BIC{$x = x \to x = y \to y = x$}
\BIC{$x = y \to y = x$}
\UIC{$\forall y\,(x = y \to y = x)$}
\UIC{$\forall x\,\forall y\,(x = y \to y = x)$}
\end{ptree}
\noindent$\Pi_{11}$:
\begin{ptree}
\AXC{$\stackrel{\Pi_{4}}{\vdots}$}\noLine
\UIC{$\varphi_{61}$}
\AXC{$\forall n \in \mathbb{N}\,(\underline{S(n)} = \underline{n} + 1)$\lbl{[N6]}}
\AXC{$\varphi_{33}$\lbl{A2}}
\BIC{$n \in \mathbb{N} \to \underline{S(n)} = \underline{n} + 1$}
\AXC{$(n \in \mathbb{N} \to \underline{S(n)} = \underline{n} + 1) \to \varphi_{61} \to \varphi_{80}$\lbl{A1}}
\BIC{$\varphi_{61} \to \varphi_{80}$}
\BIC{$\varphi_{80}$}
\end{ptree}
\noindent$\Pi_{12}$:
\begin{ptree}
\AXC{$\stackrel{\Pi_{10}}{\vdots}$}\noLine
\UIC{$\forall x\,\forall y\,(x = y \to y = x)$}
\AXC{$\varphi_{5}$\lbl{A2}}
\BIC{$\forall y\,(\underline{S(n)} = y \to y = \underline{S(n)})$}
\AXC{$\varphi_{40}$\lbl{A2}}
\BIC{$\underline{S(n)} = \underline{n} + 1 \to \underline{n} + 1 = \underline{S(n)}$}
\AXC{$\stackrel{\Pi_{11}}{\vdots}$}\noLine
\UIC{$\varphi_{80}$}
\AXC{$\varphi_{80} \to (\underline{S(n)} = \underline{n} + 1 \to \underline{n} + 1 = \underline{S(n)}) \to \varphi_{81}$\lbl{A1}}
\BIC{$(\underline{S(n)} = \underline{n} + 1 \to \underline{n} + 1 = \underline{S(n)}) \to \varphi_{81}$}
\BIC{$\varphi_{81}$}
\end{ptree}
\noindent$\Pi_{13}$:
\begin{ptree}
\AXC{$\stackrel{\Pi_{5}}{\vdots}$}\noLine
\UIC{$\varphi_{73}$}
\AXC{$\stackrel{\Pi_{6}}{\vdots}$}\noLine
\UIC{$\varphi_{73} \to \varphi_{79}$}
\BIC{$\varphi_{79}$}
\AXC{$\stackrel{\Pi_{9}}{\vdots}$}\noLine
\UIC{$\varphi_{79} \to \varphi_{75}$}
\BIC{$\varphi_{75}$}
\AXC{$\stackrel{\Pi_{12}}{\vdots}$}\noLine
\UIC{$\varphi_{81}$}
\AXC{$\varphi_{81} \to \varphi_{91}$\lbl{A1}}
\BIC{$\varphi_{91}$}
\BIC{$(\underline{n} + 1 = \underline{S(n)} \to 1 < \underline{n} + 1 \to 1 < \underline{S(n)}) \to \varphi_{71}$}
\end{ptree}
\noindent$\Pi_{14}$:
\begin{ptree}
\AXC{$z = w \to 1 < z \to 1 < w$\lbl{A6}}
\UIC{$\forall w\,(z = w \to 1 < z \to 1 < w)$}
\UIC{$\forall z\,\forall w\,(z = w \to 1 < z \to 1 < w)$}
\AXC{$\varphi_{32}$\lbl{A2}}
\BIC{$\forall w\,(\underline{n} + 1 = w \to 1 < \underline{n} + 1 \to 1 < w)$}
\AXC{$\varphi_{76}$\lbl{A2}}
\BIC{$\underline{n} + 1 = \underline{S(n)} \to 1 < \underline{n} + 1 \to 1 < \underline{S(n)}$}
\AXC{$\stackrel{\Pi_{13}}{\vdots}$}\noLine
\UIC{$(\underline{n} + 1 = \underline{S(n)} \to 1 < \underline{n} + 1 \to 1 < \underline{S(n)}) \to \varphi_{71}$}
\BIC{$\varphi_{71}$}
\end{ptree}
\noindent$\Pi_{15}$:
\begin{ptree}
\AXC{$\stackrel{\Pi_{10}}{\vdots}$}\noLine
\UIC{$\forall x\,\forall y\,(x = y \to y = x)$}
\AXC{$\varphi_{6}$\lbl{A2}}
\BIC{$\forall y\,(\underline{1_{\omega}} = y \to y = \underline{1_{\omega}})$}
\AXC{$\varphi_{21}$\lbl{A2}}
\BIC{$\underline{1_{\omega}} = 1 \to 1 = \underline{1_{\omega}}$}
\AXC{$\underline{1_{\omega}} = 1$\lbl{[N5]}}
\AXC{$\varphi_{17}$\lbl{A1}}
\BIC{$(\underline{1_{\omega}} = 1 \to 1 = \underline{1_{\omega}}) \to 1 = \underline{1_{\omega}}$}
\BIC{$1 = \underline{1_{\omega}}$}
\end{ptree}
\noindent$\Pi_{16}$:
\begin{ptree}
\AXC{$y = z \to x = y \to x = z$\lbl{A6}}
\AXC{$(y = z \to x = y \to x = z) \to x = y \to y = z \to x = z$\lbl{A1}}
\BIC{$x = y \to y = z \to x = z$}
\UIC{$\forall z\,(x = y \to y = z \to x = z)$}
\UIC{$\forall y\,\forall z\,(x = y \to y = z \to x = z)$}
\UIC{$\forall x\,\forall y\,\forall z\,(x = y \to y = z \to x = z)$}
\AXC{$\varphi_{9}$\lbl{A2}}
\BIC{$\forall y\,\forall z\,(S(n) = y \to y = z \to S(n) = z)$}
\AXC{$\varphi_{29}$\lbl{A2}}
\BIC{$\forall z\,(S(n) = S(\emptyset) \to S(\emptyset) = z \to S(n) = z)$}
\end{ptree}
\noindent$\Pi_{17}$:
\begin{ptree}
\AXC{$\stackrel{\Pi_{16}}{\vdots}$}\noLine
\UIC{$\forall z\,(S(n) = S(\emptyset) \to S(\emptyset) = z \to S(n) = z)$}
\AXC{$\forall z\,(S(n) = S(\emptyset) \to S(\emptyset) = z \to S(n) = z) \to \varphi_{4}$\lbl{A2}}
\BIC{$\varphi_{4}$}
\end{ptree}
\noindent$\Pi_{18}$:
\begin{ptree}
\AXC{$\stackrel{\Pi_{10}}{\vdots}$}\noLine
\UIC{$\forall x\,\forall y\,(x = y \to y = x)$}
\AXC{$\varphi_{2}$\lbl{A2}}
\BIC{$\forall y\,(1_{\omega} = y \to y = 1_{\omega})$}
\AXC{$\varphi_{14}$\lbl{A2}}
\BIC{$1_{\omega} = S(\emptyset) \to S(\emptyset) = 1_{\omega}$}
\AXC{$1_{\omega} = S(\emptyset)$\lbl{Def$1_\omega$}}
\AXC{$1_{\omega} = S(\emptyset) \to \varphi_{10}$\lbl{A1}}
\BIC{$\varphi_{10}$}
\BIC{$S(\emptyset) = 1_{\omega}$}
\end{ptree}
\noindent$\Pi_{19}$:
\begin{ptree}
\AXC{$z = w \to S(n) = S(z) \to S(n) = S(w)$\lbl{A6}}
\UIC{$\forall w\,(z = w \to S(n) = S(z) \to S(n) = S(w))$}
\UIC{$\forall z\,\forall w\,(z = w \to S(n) = S(z) \to S(n) = S(w))$}
\AXC{$\varphi_{23}$\lbl{A2}}
\BIC{$\forall w\,(n = w \to S(n) = S(n) \to S(n) = S(w))$}
\AXC{$\varphi_{30}$\lbl{A2}}
\BIC{$n = \emptyset \to S(n) = S(n) \to S(n) = S(\emptyset)$}
\end{ptree}
\noindent$\Pi_{20}$:
\begin{ptree}
\AXC{$x = x$\lbl{A5}}
\UIC{$\forall x\,(x = x)$}
\AXC{$\forall x\,(x = x) \to S(n) = S(n)$\lbl{A2}}
\BIC{$S(n) = S(n)$}
\AXC{$\varphi_{63}$\lbl{A1}}
\AXC{$\varphi_{63} \to \varphi_{87}$\lbl{A1}}
\BIC{$\varphi_{87}$}
\BIC{$(n = \emptyset \to S(n) = S(n) \to S(n) = S(\emptyset)) \to \varphi_{72}$}
\end{ptree}
\noindent$\Pi_{21}$:
\begin{ptree}
\AXC{$\stackrel{\Pi_{19}}{\vdots}$}\noLine
\UIC{$n = \emptyset \to S(n) = S(n) \to S(n) = S(\emptyset)$}
\AXC{$\stackrel{\Pi_{20}}{\vdots}$}\noLine
\UIC{$(n = \emptyset \to S(n) = S(n) \to S(n) = S(\emptyset)) \to \varphi_{72}$}
\BIC{$\varphi_{72}$}
\end{ptree}
\noindent$\Pi_{22}$:
\begin{ptree}
\AXC{$\stackrel{\Pi_{17}}{\vdots}$}\noLine
\UIC{$\varphi_{4}$}
\AXC{$\stackrel{\Pi_{18}}{\vdots}$}\noLine
\UIC{$S(\emptyset) = 1_{\omega}$}
\AXC{$\stackrel{\Pi_{21}}{\vdots}$}\noLine
\UIC{$\varphi_{72}$}
\AXC{$\varphi_{72} \to S(\emptyset) = 1_{\omega} \to \varphi_{4} \to \varphi_{58}$\lbl{A1}}
\BIC{$S(\emptyset) = 1_{\omega} \to \varphi_{4} \to \varphi_{58}$}
\BIC{$\varphi_{4} \to \varphi_{58}$}
\BIC{$\varphi_{58}$}
\end{ptree}
\noindent$\Pi_{23}$:
\begin{ptree}
\AXC{$\stackrel{\Pi_{10}}{\vdots}$}\noLine
\UIC{$\forall x\,\forall y\,(x = y \to y = x)$}
\AXC{$\varphi_{1}$\lbl{A2}}
\BIC{$\forall y\,(S(n) = y \to y = S(n))$}
\AXC{$\varphi_{3}$\lbl{A2}}
\BIC{$S(n) = 1_{\omega} \to 1_{\omega} = S(n)$}
\AXC{$\stackrel{\Pi_{22}}{\vdots}$}\noLine
\UIC{$\varphi_{58}$}
\AXC{$\varphi_{58} \to (S(n) = 1_{\omega} \to 1_{\omega} = S(n)) \to \varphi_{59}$\lbl{A1}}
\BIC{$(S(n) = 1_{\omega} \to 1_{\omega} = S(n)) \to \varphi_{59}$}
\BIC{$\varphi_{59}$}
\end{ptree}
\noindent$\Pi_{24}$:
\begin{ptree}
\AXC{$\stackrel{\Pi_{15}}{\vdots}$}\noLine
\UIC{$1 = \underline{1_{\omega}}$}
\AXC{$1 = \underline{1_{\omega}} \to 1 \leq \underline{1_{\omega}}$\lbl{A1}}
\BIC{$1 \leq \underline{1_{\omega}}$}
\AXC{$\stackrel{\Pi_{23}}{\vdots}$}\noLine
\UIC{$\varphi_{59}$}
\AXC{$\varphi_{59} \to 1 \leq \underline{1_{\omega}} \to \varphi_{90}$\lbl{A1}}
\BIC{$1 \leq \underline{1_{\omega}} \to \varphi_{90}$}
\BIC{$\varphi_{90}$}
\end{ptree}
\noindent$\Pi_{25}$:
\begin{ptree}
\AXC{$z = w \to 1 < \underline{z} \to 1 < \underline{w}$\lbl{A6}}
\UIC{$\forall w\,(z = w \to 1 < \underline{z} \to 1 < \underline{w})$}
\UIC{$\forall z\,\forall w\,(z = w \to 1 < \underline{z} \to 1 < \underline{w})$}
\AXC{$\varphi_{42}$\lbl{A2}}
\BIC{$\forall w\,(1_{\omega} = w \to 1 < \underline{1_{\omega}} \to 1 < \underline{w})$}
\AXC{$\varphi_{48}$\lbl{A2}}
\BIC{$1_{\omega} = S(n) \to 1 < \underline{1_{\omega}} \to 1 < \underline{S(n)}$}
\AXC{$\stackrel{\Pi_{24}}{\vdots}$}\noLine
\UIC{$\varphi_{90}$}
\BIC{$\varphi_{86}$}
\end{ptree}
\noindent$\Pi_{26}$:
\begin{ptree}
\AXC{$\stackrel{\Pi_{14}}{\vdots}$}\noLine
\UIC{$\varphi_{71}$}
\AXC{$\varphi_{71} \to \varphi_{68}$\lbl{A1}}
\BIC{$\varphi_{68}$}
\AXC{$z = w \to 1 = \underline{z} \to 1 = \underline{w}$\lbl{A6}}
\UIC{$\forall w\,(z = w \to 1 = \underline{z} \to 1 = \underline{w})$}
\UIC{$\forall z\,\forall w\,(z = w \to 1 = \underline{z} \to 1 = \underline{w})$}
\AXC{$\varphi_{26}$\lbl{A2}}
\BIC{$\forall w\,(1_{\omega} = w \to 1 = \underline{1_{\omega}} \to 1 = \underline{w})$}
\AXC{$\varphi_{37}$\lbl{A2}}
\BIC{$1_{\omega} = S(n) \to 1 = \underline{1_{\omega}} \to 1 = \underline{S(n)}$}
\AXC{$\stackrel{\Pi_{25}}{\vdots}$}\noLine
\UIC{$\varphi_{86}$}
\BIC{$\varphi_{66}$}
\AXC{$\varphi_{66} \to \varphi_{68} \to \varphi_{49}$\lbl{A1}}
\BIC{$\varphi_{68} \to \varphi_{49}$}
\BIC{$\varphi_{49}$}
\end{ptree}
\noindent$\Pi_{27}$:
\begin{ptree}
\AXC{$\forall x\,(\mathrm{Ind}(x) \leftrightarrow \emptyset \in x \land \forall y \in x\,(S(y) \in x))$\lbl{DefInd}}
\AXC{$\varphi_{83}$\lbl{A2}}
\BIC{$\mathrm{Ind}(\omega) \leftrightarrow \emptyset \in \omega \land \forall y \in \omega\,(S(y) \in \omega)$}
\AXC{$\mathrm{Ind}(\omega)$\lbl{補題(ωは帰納的)}}
\AXC{$\varphi_{89}$\lbl{A1}}
\BIC{$(\mathrm{Ind}(\omega) \leftrightarrow \emptyset \in \omega \land \forall y \in \omega\,(S(y) \in \omega)) \to \varphi_{69}$}
\BIC{$\varphi_{69}$}
\end{ptree}
\noindent$\Pi_{28}$:
\begin{ptree}
\AXC{$\stackrel{\Pi_{27}}{\vdots}$}\noLine
\UIC{$\varphi_{69}$}
\AXC{$\varphi_{69} \to \varphi_{55}$\lbl{A1}}
\BIC{$\varphi_{55}$}
\AXC{$\varphi_{55} \to (\forall y \in \omega\,(S(y) \in \omega) \to n \in \omega \to S(n) \in \omega) \to \varphi_{51}$\lbl{A1}}
\BIC{$(\forall y \in \omega\,(S(y) \in \omega) \to n \in \omega \to S(n) \in \omega) \to \varphi_{51}$}
\end{ptree}
\noindent$\Pi_{29}$:
\begin{ptree}
\AXC{$\forall y \in \omega\,(S(y) \in \omega) \to n \in \omega \to S(n) \in \omega$\lbl{A2}}
\AXC{$\stackrel{\Pi_{28}}{\vdots}$}\noLine
\UIC{$(\forall y \in \omega\,(S(y) \in \omega) \to n \in \omega \to S(n) \in \omega) \to \varphi_{51}$}
\BIC{$\varphi_{51}$}
\AXC{$\varphi_{51} \to \varphi_{28} \to \varphi_{34}$\lbl{A1}}
\BIC{$\varphi_{28} \to \varphi_{34}$}
\end{ptree}
\noindent$\Pi_{30}$:
\begin{ptree}
\AXC{$\stackrel{\Pi_{26}}{\vdots}$}\noLine
\UIC{$\varphi_{49}$}
\AXC{$\stackrel{\Pi_{3}}{\vdots}$}\noLine
\UIC{$\varphi_{28}$}
\AXC{$\stackrel{\Pi_{29}}{\vdots}$}\noLine
\UIC{$\varphi_{28} \to \varphi_{34}$}
\BIC{$\varphi_{34}$}
\AXC{$\forall n\,(n \in B_{3} \leftrightarrow n \in \omega \land (n \in \mathbb{N} \to 1 \leq \underline{n})) \to \varphi_{18}$\lbl{A2}}
\AXC{$\varphi_{93}$\lbl{A1}}
\BIC{$\varphi_{34} \to \varphi_{49} \to \varphi_{31}$}
\BIC{$\varphi_{49} \to \varphi_{31}$}
\BIC{$\varphi_{31}$}
\UIC{$\forall n\,(\varphi_{31})$}
\end{ptree}
\noindent$\Pi_{31}$:
\begin{ptree}
\AXC{$\varphi_{67}$\lbl{A4}}
\AXC{$\forall n\,(\varphi_{31}) \to \varphi_{39}$\lbl{A3}}
\AXC{$\stackrel{\Pi_{30}}{\vdots}$}\noLine
\UIC{$\forall n\,(\varphi_{31})$}
\AXC{$\forall n\,(\varphi_{31}) \to (\forall n\,(\varphi_{31}) \to \varphi_{39}) \to \varphi_{67} \to \varphi_{35}$\lbl{A1}}
\BIC{$(\forall n\,(\varphi_{31}) \to \varphi_{39}) \to \varphi_{67} \to \varphi_{35}$}
\BIC{$\varphi_{67} \to \varphi_{35}$}
\BIC{$\varphi_{35}$}
\end{ptree}
\noindent$\Pi_{32}$:
\begin{ptree}
\AXC{$\forall n\,(n \in \mathbb{N} \leftrightarrow n \in \omega \land n \neq \emptyset)$\lbl{Def$\mathbb{N}$}}
\AXC{$\varphi_{64}$\lbl{A2}}
\BIC{$\emptyset \in \mathbb{N} \leftrightarrow \emptyset \in \omega \land \emptyset \neq \emptyset$}
\AXC{$(\emptyset \in \mathbb{N} \leftrightarrow \emptyset \in \omega \land \emptyset \neq \emptyset) \to \emptyset = \emptyset \to \emptyset \in \mathbb{N} \to 1 \leq \underline{\emptyset}$\lbl{A1}}
\BIC{$\emptyset = \emptyset \to \emptyset \in \mathbb{N} \to 1 \leq \underline{\emptyset}$}
\end{ptree}
\noindent$\Pi_{33}$:
\begin{ptree}
\AXC{$x = x$\lbl{A5}}
\UIC{$\forall x\,(x = x)$}
\AXC{$\forall x\,(x = x) \to \emptyset = \emptyset$\lbl{A2}}
\BIC{$\emptyset = \emptyset$}
\AXC{$\stackrel{\Pi_{32}}{\vdots}$}\noLine
\UIC{$\emptyset = \emptyset \to \emptyset \in \mathbb{N} \to 1 \leq \underline{\emptyset}$}
\BIC{$\emptyset \in \mathbb{N} \to 1 \leq \underline{\emptyset}$}
\AXC{$\stackrel{\Pi_{27}}{\vdots}$}\noLine
\UIC{$\varphi_{69}$}
\AXC{$\varphi_{84}$\lbl{A2}}
\AXC{$\varphi_{84} \to \varphi_{69} \to \varphi_{70}$\lbl{A1}}
\BIC{$\varphi_{69} \to \varphi_{70}$}
\BIC{$\varphi_{70}$}
\BIC{$\forall n\,(n \in B_{3} \leftrightarrow n \in \omega \land (n \in \mathbb{N} \to 1 \leq \underline{n})) \to \emptyset \in B_{3}$}
\end{ptree}
\noindent$\Pi_{34}$:
\begin{ptree}
\AXC{$\forall x\,(\mathrm{Ind}(x) \leftrightarrow \emptyset \in x \land \forall y \in x\,(S(y) \in x))$\lbl{DefInd}}
\AXC{$\varphi_{74}$\lbl{A2}}
\BIC{$\mathrm{Ind}(B_{3}) \leftrightarrow \emptyset \in B_{3} \land \forall y \in B_{3}\,(S(y) \in B_{3})$}
\AXC{$\varphi_{94}$\lbl{A1}}
\BIC{$\varphi_{92}$}
\end{ptree}
\noindent$\Pi_{35}$:
\begin{ptree}
\AXC{$\stackrel{\Pi_{31}}{\vdots}$}\noLine
\UIC{$\varphi_{35}$}
\AXC{$\stackrel{\Pi_{33}}{\vdots}$}\noLine
\UIC{$\forall n\,(n \in B_{3} \leftrightarrow n \in \omega \land (n \in \mathbb{N} \to 1 \leq \underline{n})) \to \emptyset \in B_{3}$}
\AXC{$\stackrel{\Pi_{34}}{\vdots}$}\noLine
\UIC{$\varphi_{92}$}
\BIC{$\varphi_{35} \to \varphi_{88}$}
\BIC{$\varphi_{88}$}
\end{ptree}
\noindent$\Pi_{36}$:
\begin{ptree}
\AXC{$\stackrel{\Pi_{1}}{\vdots}$}\noLine
\UIC{$\forall n \in B_{3}\,(S(n) \in B_{3}) \to \forall y \in B_{3}\,(S(y) \in B_{3})$}
\AXC{$\stackrel{\Pi_{2}}{\vdots}$}\noLine
\UIC{$\forall y \in B_{3}\,(S(y) \in B_{3}) \to \forall n \in B_{3}\,(S(n) \in B_{3})$}
\AXC{$\stackrel{\Pi_{35}}{\vdots}$}\noLine
\UIC{$\varphi_{88}$}
\BIC{$\varphi_{82}$}
\BIC{$\varphi_{19}$}
\end{ptree}
\noindent$\Pi_{37}$:
\begin{ptree}
\AXC{$\forall n\,(n \in \omega \leftrightarrow \forall x\,(\mathrm{Ind}(x) \to n \in x))$\lbl{Def$\omega$}}
\AXC{$\varphi_{27}$\lbl{A2}}
\BIC{$n \in \omega \leftrightarrow \forall x\,(\mathrm{Ind}(x) \to n \in x)$}
\AXC{$\varphi_{36}$\lbl{A1}}
\AXC{$\varphi_{36} \to (n \in \omega \leftrightarrow \forall x\,(\mathrm{Ind}(x) \to n \in x)) \to \varphi_{45}$\lbl{A1}}
\BIC{$(n \in \omega \leftrightarrow \forall x\,(\mathrm{Ind}(x) \to n \in x)) \to \varphi_{45}$}
\BIC{$\varphi_{45}$}
\end{ptree}
\noindent$\Pi_{38}$:
\begin{ptree}
\AXC{$\stackrel{\Pi_{36}}{\vdots}$}\noLine
\UIC{$\varphi_{19}$}
\AXC{$\forall x\,(\mathrm{Ind}(x) \to n \in x) \to \mathrm{Ind}(B_{3}) \to n \in B_{3}$\lbl{A2}}
\AXC{$\stackrel{\Pi_{37}}{\vdots}$}\noLine
\UIC{$\varphi_{45}$}
\AXC{$\varphi_{45} \to \varphi_{85}$\lbl{A1}}
\BIC{$\varphi_{85}$}
\BIC{$\varphi_{43}$}
\AXC{$\varphi_{43} \to \varphi_{19} \to \varphi_{56} \to \varphi_{47}$\lbl{A1}}
\BIC{$\varphi_{19} \to \varphi_{56} \to \varphi_{47}$}
\BIC{$\varphi_{56} \to \varphi_{47}$}
\end{ptree}
\noindent$\Pi_{39}$:
\begin{ptree}
\AXC{$\varphi_{67}$\lbl{A4}}
\AXC{$\forall n\,(\varphi_{47}) \to \varphi_{53}$\lbl{A3}}
\AXC{$\varphi_{56}$\lbl{A2}}
\AXC{$\stackrel{\Pi_{38}}{\vdots}$}\noLine
\UIC{$\varphi_{56} \to \varphi_{47}$}
\BIC{$\varphi_{47}$}
\UIC{$\forall n\,(\varphi_{47})$}
\AXC{$\forall n\,(\varphi_{47}) \to (\forall n\,(\varphi_{47}) \to \varphi_{53}) \to \varphi_{67} \to \varphi_{50}$\lbl{A1}}
\BIC{$(\forall n\,(\varphi_{47}) \to \varphi_{53}) \to \varphi_{67} \to \varphi_{50}$}
\BIC{$\varphi_{67} \to \varphi_{50}$}
\BIC{$\varphi_{50}$}
\AXC{$\varphi_{50} \to \varphi_{52}$\lbl{A1}}
\BIC{$\varphi_{52}$}
\UIC{$\forall B_{3}\,(\varphi_{52})$}
\end{ptree}
\noindent$\Pi_{40}$:
\begin{ptree}
\AXC{$\varphi_{11}$\lbl{Sep}}
\AXC{$\varphi_{11} \to \exists B_{3}\,\forall n\,(n \in B_{3} \leftrightarrow n \in \omega \land (n \in \mathbb{N} \to 1 \leq \underline{n}))$\lbl{A2}}
\BIC{$\exists B_{3}\,\forall n\,(n \in B_{3} \leftrightarrow n \in \omega \land (n \in \mathbb{N} \to 1 \leq \underline{n}))$}
\AXC{$\exists B_{3}\,\forall n\,(n \in B_{3} \leftrightarrow n \in \omega \land (n \in \mathbb{N} \to 1 \leq \underline{n})) \to \varphi_{95}$\lbl{A1}}
\BIC{$\varphi_{95}$}
\end{ptree}
\noindent$\Pi_{41}$:
\begin{ptree}
\AXC{$\varphi_{44}$\lbl{A4}}
\AXC{$\forall B_{3}\,(\varphi_{52}) \to \varphi_{60}$\lbl{A3}}
\AXC{$\stackrel{\Pi_{39}}{\vdots}$}\noLine
\UIC{$\forall B_{3}\,(\varphi_{52})$}
\AXC{$\stackrel{\Pi_{40}}{\vdots}$}\noLine
\UIC{$\varphi_{95}$}
\BIC{$(\forall B_{3}\,(\varphi_{52}) \to \varphi_{60}) \to \varphi_{44} \to \forall n \in \omega\,(n \in \mathbb{N} \to 1 \leq \underline{n})$}
\BIC{$\varphi_{44} \to \forall n \in \omega\,(n \in \mathbb{N} \to 1 \leq \underline{n})$}
\BIC{$\forall n \in \omega\,(n \in \mathbb{N} \to 1 \leq \underline{n})$}
\end{ptree}
\noindent$\Pi_{42}$:
\begin{ptree}
\AXC{$\stackrel{\Pi_{41}}{\vdots}$}\noLine
\UIC{$\forall n \in \omega\,(n \in \mathbb{N} \to 1 \leq \underline{n})$}
\AXC{$\forall n \in \omega\,(n \in \mathbb{N} \to 1 \leq \underline{n}) \to n \in \omega \to n \in \mathbb{N} \to 1 \leq \underline{n}$\lbl{A2}}
\BIC{$n \in \omega \to n \in \mathbb{N} \to 1 \leq \underline{n}$}
\AXC{$(n \in \omega \to n \in \mathbb{N} \to 1 \leq \underline{n}) \to \varphi_{25}$\lbl{A1}}
\BIC{$\varphi_{25}$}
\end{ptree}
\noindent$\Pi_{43}$:
\begin{ptree}
\AXC{$\forall n\,(n \in \mathbb{N} \leftrightarrow n \in \omega \land n \neq \emptyset)$\lbl{Def$\mathbb{N}$}}
\AXC{$\varphi_{38}$\lbl{A2}}
\BIC{$n \in \mathbb{N} \leftrightarrow n \in \omega \land n \neq \emptyset$}
\AXC{$n \in \mathbb{N} \to n \in \mathbb{N}$\lbl{A1}}
\AXC{$\varphi_{46}$\lbl{A1}}
\BIC{$(n \in \mathbb{N} \leftrightarrow n \in \omega \land n \neq \emptyset) \to n \in \mathbb{N} \to n \in \omega$}
\BIC{$n \in \mathbb{N} \to n \in \omega$}
\AXC{$\stackrel{\Pi_{42}}{\vdots}$}\noLine
\UIC{$\varphi_{25}$}
\BIC{$(n \in \mathbb{N} \to n \in \mathbb{N}) \to n \in \mathbb{N} \to 1 \leq \underline{n}$}
\end{ptree}
\noindent 主木:
\begin{ptree}
\AXC{$n \in \mathbb{N} \to n \in \mathbb{N}$\lbl{A1}}
\AXC{$\stackrel{\Pi_{43}}{\vdots}$}\noLine
\UIC{$(n \in \mathbb{N} \to n \in \mathbb{N}) \to n \in \mathbb{N} \to 1 \leq \underline{n}$}
\BIC{$n \in \mathbb{N} \to 1 \leq \underline{n}$}
\UIC{$\forall n \in \mathbb{N}\,(1 \leq \underline{n})$}
\end{ptree}
\begin{abbrevlist}
\item $\varphi_{1} :\equiv \forall x\,\forall y\,(x = y \to y = x) \to \forall y\,(S(n) = y \to y = S(n))$
\item $\varphi_{2} :\equiv \forall x\,\forall y\,(x = y \to y = x) \to \forall y\,(1_{\omega} = y \to y = 1_{\omega})$
\item $\varphi_{3} :\equiv \forall y\,(S(n) = y \to y = S(n)) \to S(n) = 1_{\omega} \to 1_{\omega} = S(n)$
\item $\varphi_{4} :\equiv S(n) = S(\emptyset) \to S(\emptyset) = 1_{\omega} \to S(n) = 1_{\omega}$
\item $\varphi_{5} :\equiv \forall x\,\forall y\,(x = y \to y = x) \to \forall y\,(\underline{S(n)} = y \to y = \underline{S(n)})$
\item $\varphi_{6} :\equiv \forall x\,\forall y\,(x = y \to y = x) \to \forall y\,(\underline{1_{\omega}} = y \to y = \underline{1_{\omega}})$
\item $\varphi_{7} :\equiv \forall y\,(\forall n \in B_{3}\,(S(n) \in B_{3}) \to y \in B_{3} \to S(y) \in B_{3})$
\item $\varphi_{8} :\equiv \forall n\,(\forall y \in B_{3}\,(S(y) \in B_{3}) \to n \in B_{3} \to S(n) \in B_{3})$
\item $\varphi_{9} :\equiv \forall x\,\forall y\,\forall z\,(x = y \to y = z \to x = z) \to \forall y\,\forall z\,(S(n) = y \to y = z \to S(n) = z)$
\item $\varphi_{10} :\equiv (1_{\omega} = S(\emptyset) \to S(\emptyset) = 1_{\omega}) \to S(\emptyset) = 1_{\omega}$
\item $\varphi_{11} :\equiv \forall A_{4}\,\exists B_{3}\,\forall n\,(n \in B_{3} \leftrightarrow n \in A_{4} \land (n \in \mathbb{N} \to 1 \leq \underline{n}))$
\item $\varphi_{12} :\equiv \forall n \in B_{3}\,(S(n) \in B_{3}) \to \forall y\,\forall n \in B_{3}\,(S(n) \in B_{3})$
\item $\varphi_{13} :\equiv \forall y \in B_{3}\,(S(y) \in B_{3}) \to \forall n\,\forall y \in B_{3}\,(S(y) \in B_{3})$
\item $\varphi_{14} :\equiv \forall y\,(1_{\omega} = y \to y = 1_{\omega}) \to 1_{\omega} = S(\emptyset) \to S(\emptyset) = 1_{\omega}$
\item $\varphi_{15} :\equiv \forall y\,\forall n \in B_{3}\,(S(n) \in B_{3}) \to \forall y \in B_{3}\,(S(y) \in B_{3})$
\item $\varphi_{16} :\equiv \forall n\,\forall y \in B_{3}\,(S(y) \in B_{3}) \to \forall n \in B_{3}\,(S(n) \in B_{3})$
\item $\varphi_{17} :\equiv \underline{1_{\omega}} = 1 \to (\underline{1_{\omega}} = 1 \to 1 = \underline{1_{\omega}}) \to 1 = \underline{1_{\omega}}$
\item $\varphi_{18} :\equiv S(n) \in B_{3} \leftrightarrow S(n) \in \omega \land (S(n) \in \mathbb{N} \to 1 \leq \underline{S(n)})$
\item $\varphi_{19} :\equiv \forall n\,(n \in B_{3} \leftrightarrow n \in \omega \land (n \in \mathbb{N} \to 1 \leq \underline{n})) \to \mathrm{Ind}(B_{3})$
\item $\varphi_{20} :\equiv \forall x\,\forall y\,\forall z\,(x \leq y \to y < z \to x < z) \to \forall y\,\forall z\,(1 \leq y \to y < z \to 1 < z)$
\item $\varphi_{21} :\equiv \forall y\,(\underline{1_{\omega}} = y \to y = \underline{1_{\omega}}) \to \underline{1_{\omega}} = 1 \to 1 = \underline{1_{\omega}}$
\item $\varphi_{22} :\equiv \forall y\,\forall z\,(1 \leq y \to y < z \to 1 < z) \to \forall z\,(1 \leq \underline{n} \to \underline{n} < z \to 1 < z)$
\item $\varphi_{23} :\equiv \forall z\,\forall w\,(z = w \to S(n) = S(z) \to S(n) = S(w)) \to \forall w\,(n = w \to S(n) = S(n) \to S(n) = S(w))$
\item $\varphi_{24} :\equiv \forall n\,(n \in B_{3} \leftrightarrow n \in \omega \land (n \in \mathbb{N} \to 1 \leq \underline{n})) \to n \in B_{3} \to n \in B_{3}$
\item $\varphi_{25} :\equiv (n \in \mathbb{N} \to n \in \omega) \to (n \in \mathbb{N} \to n \in \mathbb{N}) \to n \in \mathbb{N} \to 1 \leq \underline{n}$
\item $\varphi_{26} :\equiv \forall z\,\forall w\,(z = w \to 1 = \underline{z} \to 1 = \underline{w}) \to \forall w\,(1_{\omega} = w \to 1 = \underline{1_{\omega}} \to 1 = \underline{w})$
\item $\varphi_{27} :\equiv \forall n\,(n \in \omega \leftrightarrow \forall x\,(\mathrm{Ind}(x) \to n \in x)) \to (n \in \omega \leftrightarrow \forall x\,(\mathrm{Ind}(x) \to n \in x))$
\item $\varphi_{28} :\equiv \forall n\,(n \in B_{3} \leftrightarrow n \in \omega \land (n \in \mathbb{N} \to 1 \leq \underline{n})) \to n \in B_{3} \to n \in \omega$
\item $\varphi_{29} :\equiv \forall y\,\forall z\,(S(n) = y \to y = z \to S(n) = z) \to \forall z\,(S(n) = S(\emptyset) \to S(\emptyset) = z \to S(n) = z)$
\item $\varphi_{30} :\equiv \forall w\,(n = w \to S(n) = S(n) \to S(n) = S(w)) \to n = \emptyset \to S(n) = S(n) \to S(n) = S(\emptyset)$
\item $\varphi_{31} :\equiv \forall n\,(n \in B_{3} \leftrightarrow n \in \omega \land (n \in \mathbb{N} \to 1 \leq \underline{n})) \to n \in B_{3} \to S(n) \in B_{3}$
\item $\varphi_{32} :\equiv \forall z\,\forall w\,(z = w \to 1 < z \to 1 < w) \to \forall w\,(\underline{n} + 1 = w \to 1 < \underline{n} + 1 \to 1 < w)$
\item $\varphi_{33} :\equiv \forall n \in \mathbb{N}\,(\underline{S(n)} = \underline{n} + 1) \to n \in \mathbb{N} \to \underline{S(n)} = \underline{n} + 1$
\item $\varphi_{34} :\equiv \forall n\,(n \in B_{3} \leftrightarrow n \in \omega \land (n \in \mathbb{N} \to 1 \leq \underline{n})) \to n \in B_{3} \to S(n) \in \omega$
\item $\varphi_{35} :\equiv \forall n\,(n \in B_{3} \leftrightarrow n \in \omega \land (n \in \mathbb{N} \to 1 \leq \underline{n})) \to \forall n \in B_{3}\,(S(n) \in B_{3})$
\item $\varphi_{36} :\equiv \forall n\,(n \in B_{3} \leftrightarrow n \in \omega \land (n \in \mathbb{N} \to 1 \leq \underline{n})) \to n \in \omega \to n \in \omega$
\item $\varphi_{37} :\equiv \forall w\,(1_{\omega} = w \to 1 = \underline{1_{\omega}} \to 1 = \underline{w}) \to 1_{\omega} = S(n) \to 1 = \underline{1_{\omega}} \to 1 = \underline{S(n)}$
\item $\varphi_{38} :\equiv \forall n\,(n \in \mathbb{N} \leftrightarrow n \in \omega \land n \neq \emptyset) \to (n \in \mathbb{N} \leftrightarrow n \in \omega \land n \neq \emptyset)$
\item $\varphi_{39} :\equiv \forall n\,\forall n\,(n \in B_{3} \leftrightarrow n \in \omega \land (n \in \mathbb{N} \to 1 \leq \underline{n})) \to \forall n \in B_{3}\,(S(n) \in B_{3})$
\item $\varphi_{40} :\equiv \forall y\,(\underline{S(n)} = y \to y = \underline{S(n)}) \to \underline{S(n)} = \underline{n} + 1 \to \underline{n} + 1 = \underline{S(n)}$
\item $\varphi_{41} :\equiv \forall n\,(n \in B_{3} \leftrightarrow n \in \omega \land (n \in \mathbb{N} \to 1 \leq \underline{n})) \to n \in B_{3} \to n \in \mathbb{N} \to 1 \leq \underline{n}$
\item $\varphi_{42} :\equiv \forall z\,\forall w\,(z = w \to 1 < \underline{z} \to 1 < \underline{w}) \to \forall w\,(1_{\omega} = w \to 1 < \underline{1_{\omega}} \to 1 < \underline{w})$
\item $\varphi_{43} :\equiv \forall n\,(n \in B_{3} \leftrightarrow n \in \omega \land (n \in \mathbb{N} \to 1 \leq \underline{n})) \to n \in \omega \to \mathrm{Ind}(B_{3}) \to n \in B_{3}$
\item $\varphi_{44} :\equiv \neg \forall n \in \omega\,(n \in \mathbb{N} \to 1 \leq \underline{n}) \to \forall B_{3}\,(\neg \forall n \in \omega\,(n \in \mathbb{N} \to 1 \leq \underline{n}))$
\item $\varphi_{45} :\equiv \forall n\,(n \in B_{3} \leftrightarrow n \in \omega \land (n \in \mathbb{N} \to 1 \leq \underline{n})) \to n \in \omega \to \forall x\,(\mathrm{Ind}(x) \to n \in x)$
\item $\varphi_{46} :\equiv (n \in \mathbb{N} \to n \in \mathbb{N}) \to (n \in \mathbb{N} \leftrightarrow n \in \omega \land n \neq \emptyset) \to n \in \mathbb{N} \to n \in \omega$
\item $\varphi_{47} :\equiv \forall n\,(n \in B_{3} \leftrightarrow n \in \omega \land (n \in \mathbb{N} \to 1 \leq \underline{n})) \to n \in \omega \to n \in \mathbb{N} \to 1 \leq \underline{n}$
\item $\varphi_{48} :\equiv \forall w\,(1_{\omega} = w \to 1 < \underline{1_{\omega}} \to 1 < \underline{w}) \to 1_{\omega} = S(n) \to 1 < \underline{1_{\omega}} \to 1 < \underline{S(n)}$
\item $\varphi_{49} :\equiv \forall n\,(n \in B_{3} \leftrightarrow n \in \omega \land (n \in \mathbb{N} \to 1 \leq \underline{n})) \to n \in B_{3} \to S(n) \in \mathbb{N} \to 1 \leq \underline{S(n)}$
\item $\varphi_{50} :\equiv \forall n\,(n \in B_{3} \leftrightarrow n \in \omega \land (n \in \mathbb{N} \to 1 \leq \underline{n})) \to \forall n \in \omega\,(n \in \mathbb{N} \to 1 \leq \underline{n})$
\item $\varphi_{51} :\equiv \forall n\,(n \in B_{3} \leftrightarrow n \in \omega \land (n \in \mathbb{N} \to 1 \leq \underline{n})) \to n \in B_{3} \to n \in \omega \to S(n) \in \omega$
\item $\varphi_{52} :\equiv \neg \forall n \in \omega\,(n \in \mathbb{N} \to 1 \leq \underline{n}) \to \neg \forall n\,(n \in B_{3} \leftrightarrow n \in \omega \land (n \in \mathbb{N} \to 1 \leq \underline{n}))$
\item $\varphi_{53} :\equiv \forall n\,\forall n\,(n \in B_{3} \leftrightarrow n \in \omega \land (n \in \mathbb{N} \to 1 \leq \underline{n})) \to \forall n \in \omega\,(n \in \mathbb{N} \to 1 \leq \underline{n})$
\item $\varphi_{54} :\equiv \forall z\,(1 \leq \underline{n} \to \underline{n} < z \to 1 < z) \to 1 \leq \underline{n} \to \underline{n} < \underline{n} + 1 \to 1 < \underline{n} + 1$
\item $\varphi_{55} :\equiv \forall n\,(n \in B_{3} \leftrightarrow n \in \omega \land (n \in \mathbb{N} \to 1 \leq \underline{n})) \to n \in B_{3} \to \forall y \in \omega\,(S(y) \in \omega)$
\item $\varphi_{56} :\equiv \forall n\,(n \in B_{3} \leftrightarrow n \in \omega \land (n \in \mathbb{N} \to 1 \leq \underline{n})) \to (n \in B_{3} \leftrightarrow n \in \omega \land (n \in \mathbb{N} \to 1 \leq \underline{n}))$
\item $\varphi_{57} :\equiv \forall n\,(n \in B_{3} \leftrightarrow n \in \omega \land (n \in \mathbb{N} \to 1 \leq \underline{n})) \to n \in B_{3} \to n \in \omega \land (n \in \mathbb{N} \to 1 \leq \underline{n})$
\item $\varphi_{58} :\equiv \forall n\,(n \in B_{3} \leftrightarrow n \in \omega \land (n \in \mathbb{N} \to 1 \leq \underline{n})) \to n \in B_{3} \to S(n) \in \mathbb{N} \to n = \emptyset \to S(n) = 1_{\omega}$
\item $\varphi_{59} :\equiv \forall n\,(n \in B_{3} \leftrightarrow n \in \omega \land (n \in \mathbb{N} \to 1 \leq \underline{n})) \to n \in B_{3} \to S(n) \in \mathbb{N} \to n = \emptyset \to 1_{\omega} = S(n)$
\item $\varphi_{60} :\equiv \forall B_{3}\,(\neg \forall n \in \omega\,(n \in \mathbb{N} \to 1 \leq \underline{n})) \to \forall B_{3}\,(\neg \forall n\,(n \in B_{3} \leftrightarrow n \in \omega \land (n \in \mathbb{N} \to 1 \leq \underline{n})))$
\item $\varphi_{61} :\equiv \forall n\,(n \in B_{3} \leftrightarrow n \in \omega \land (n \in \mathbb{N} \to 1 \leq \underline{n})) \to n \in B_{3} \to S(n) \in \mathbb{N} \to n \neq \emptyset \to n \in \mathbb{N}$
\item $\varphi_{62} :\equiv \forall n\,(n \in B_{3} \leftrightarrow n \in \omega \land (n \in \mathbb{N} \to 1 \leq \underline{n})) \to n \in B_{3} \to S(n) \in \mathbb{N} \to n \neq \emptyset \to 1 \leq \underline{n}$
\item $\varphi_{63} :\equiv \forall n\,(n \in B_{3} \leftrightarrow n \in \omega \land (n \in \mathbb{N} \to 1 \leq \underline{n})) \to n \in B_{3} \to S(n) \in \mathbb{N} \to n = \emptyset \to n = \emptyset$
\item $\varphi_{64} :\equiv \forall n\,(n \in \mathbb{N} \leftrightarrow n \in \omega \land n \neq \emptyset) \to (\emptyset \in \mathbb{N} \leftrightarrow \emptyset \in \omega \land \emptyset \neq \emptyset)$
\item $\varphi_{65} :\equiv \forall n\,(n \in B_{3} \leftrightarrow n \in \omega \land (n \in \mathbb{N} \to 1 \leq \underline{n})) \to n \in B_{3} \to S(n) \in \mathbb{N} \to n \neq \emptyset \to n \neq \emptyset$
\item $\varphi_{66} :\equiv \forall n\,(n \in B_{3} \leftrightarrow n \in \omega \land (n \in \mathbb{N} \to 1 \leq \underline{n})) \to n \in B_{3} \to S(n) \in \mathbb{N} \to n = \emptyset \to 1 \leq \underline{S(n)}$
\item $\varphi_{67} :\equiv \forall n\,(n \in B_{3} \leftrightarrow n \in \omega \land (n \in \mathbb{N} \to 1 \leq \underline{n})) \to \forall n\,\forall n\,(n \in B_{3} \leftrightarrow n \in \omega \land (n \in \mathbb{N} \to 1 \leq \underline{n}))$
\item $\varphi_{68} :\equiv \forall n\,(n \in B_{3} \leftrightarrow n \in \omega \land (n \in \mathbb{N} \to 1 \leq \underline{n})) \to n \in B_{3} \to S(n) \in \mathbb{N} \to n \neq \emptyset \to 1 \leq \underline{S(n)}$
\item $\varphi_{69} :\equiv \forall n\,(n \in B_{3} \leftrightarrow n \in \omega \land (n \in \mathbb{N} \to 1 \leq \underline{n})) \to \emptyset \in \omega \land \forall y \in \omega\,(S(y) \in \omega)$
\item $\varphi_{70} :\equiv (\emptyset \in \mathbb{N} \to 1 \leq \underline{\emptyset}) \to \forall n\,(n \in B_{3} \leftrightarrow n \in \omega \land (n \in \mathbb{N} \to 1 \leq \underline{n})) \to \emptyset \in B_{3}$
\item $\varphi_{71} :\equiv \forall n\,(n \in B_{3} \leftrightarrow n \in \omega \land (n \in \mathbb{N} \to 1 \leq \underline{n})) \to n \in B_{3} \to S(n) \in \mathbb{N} \to n \neq \emptyset \to 1 < \underline{S(n)}$
\item $\varphi_{72} :\equiv \forall n\,(n \in B_{3} \leftrightarrow n \in \omega \land (n \in \mathbb{N} \to 1 \leq \underline{n})) \to n \in B_{3} \to S(n) \in \mathbb{N} \to n = \emptyset \to S(n) = S(\emptyset)$
\item $\varphi_{73} :\equiv \forall n\,(n \in B_{3} \leftrightarrow n \in \omega \land (n \in \mathbb{N} \to 1 \leq \underline{n})) \to n \in B_{3} \to S(n) \in \mathbb{N} \to n \neq \emptyset \to \underline{n} \in \mathbb{R}$
\item $\varphi_{74} :\equiv \forall x\,(\mathrm{Ind}(x) \leftrightarrow \emptyset \in x \land \forall y \in x\,(S(y) \in x)) \to (\mathrm{Ind}(B_{3}) \leftrightarrow \emptyset \in B_{3} \land \forall y \in B_{3}\,(S(y) \in B_{3}))$
\item $\varphi_{75} :\equiv \forall n\,(n \in B_{3} \leftrightarrow n \in \omega \land (n \in \mathbb{N} \to 1 \leq \underline{n})) \to n \in B_{3} \to S(n) \in \mathbb{N} \to n \neq \emptyset \to 1 < \underline{n} + 1$
\item $\varphi_{76} :\equiv \forall w\,(\underline{n} + 1 = w \to 1 < \underline{n} + 1 \to 1 < w) \to \underline{n} + 1 = \underline{S(n)} \to 1 < \underline{n} + 1 \to 1 < \underline{S(n)}$
\item $\varphi_{77} :\equiv \varphi_{12} \to \forall n \in B_{3}\,(S(n) \in B_{3}) \to \forall y \in B_{3}\,(S(y) \in B_{3})$
\item $\varphi_{78} :\equiv \varphi_{13} \to \forall y \in B_{3}\,(S(y) \in B_{3}) \to \forall n \in B_{3}\,(S(n) \in B_{3})$
\item $\varphi_{79} :\equiv \forall n\,(n \in B_{3} \leftrightarrow n \in \omega \land (n \in \mathbb{N} \to 1 \leq \underline{n})) \to n \in B_{3} \to S(n) \in \mathbb{N} \to n \neq \emptyset \to \underline{n} < \underline{n} + 1$
\item $\varphi_{80} :\equiv \forall n\,(n \in B_{3} \leftrightarrow n \in \omega \land (n \in \mathbb{N} \to 1 \leq \underline{n})) \to n \in B_{3} \to S(n) \in \mathbb{N} \to n \neq \emptyset \to \underline{S(n)} = \underline{n} + 1$
\item $\varphi_{81} :\equiv \forall n\,(n \in B_{3} \leftrightarrow n \in \omega \land (n \in \mathbb{N} \to 1 \leq \underline{n})) \to n \in B_{3} \to S(n) \in \mathbb{N} \to n \neq \emptyset \to \underline{n} + 1 = \underline{S(n)}$
\item $\varphi_{82} :\equiv (\forall n \in B_{3}\,(S(n) \in B_{3}) \to \forall y \in B_{3}\,(S(y) \in B_{3})) \to \varphi_{19}$
\item $\varphi_{83} :\equiv \forall x\,(\mathrm{Ind}(x) \leftrightarrow \emptyset \in x \land \forall y \in x\,(S(y) \in x)) \to (\mathrm{Ind}(\omega) \leftrightarrow \emptyset \in \omega \land \forall y \in \omega\,(S(y) \in \omega))$
\item $\varphi_{84} :\equiv \forall n\,(n \in B_{3} \leftrightarrow n \in \omega \land (n \in \mathbb{N} \to 1 \leq \underline{n})) \to (\emptyset \in B_{3} \leftrightarrow \emptyset \in \omega \land (\emptyset \in \mathbb{N} \to 1 \leq \underline{\emptyset}))$
\item $\varphi_{85} :\equiv (\forall x\,(\mathrm{Ind}(x) \to n \in x) \to \mathrm{Ind}(B_{3}) \to n \in B_{3}) \to \varphi_{43}$
\item $\varphi_{86} :\equiv (1_{\omega} = S(n) \to 1 = \underline{1_{\omega}} \to 1 = \underline{S(n)}) \to \varphi_{66}$
\item $\varphi_{87} :\equiv S(n) = S(n) \to (n = \emptyset \to S(n) = S(n) \to S(n) = S(\emptyset)) \to \varphi_{72}$
\item $\varphi_{88} :\equiv (\forall y \in B_{3}\,(S(y) \in B_{3}) \to \forall n \in B_{3}\,(S(n) \in B_{3})) \to \varphi_{82}$
\item $\varphi_{89} :\equiv \mathrm{Ind}(\omega) \to (\mathrm{Ind}(\omega) \leftrightarrow \emptyset \in \omega \land \forall y \in \omega\,(S(y) \in \omega)) \to \varphi_{69}$
\item $\varphi_{90} :\equiv (1_{\omega} = S(n) \to 1 < \underline{1_{\omega}} \to 1 < \underline{S(n)}) \to \varphi_{86}$
\item $\varphi_{91} :\equiv \varphi_{75} \to (\underline{n} + 1 = \underline{S(n)} \to 1 < \underline{n} + 1 \to 1 < \underline{S(n)}) \to \varphi_{71}$
\item $\varphi_{92} :\equiv (\forall n\,(n \in B_{3} \leftrightarrow n \in \omega \land (n \in \mathbb{N} \to 1 \leq \underline{n})) \to \emptyset \in B_{3}) \to \varphi_{35} \to \varphi_{88}$
\item $\varphi_{93} :\equiv (\forall n\,(n \in B_{3} \leftrightarrow n \in \omega \land (n \in \mathbb{N} \to 1 \leq \underline{n})) \to \varphi_{18}) \to \varphi_{34} \to \varphi_{49} \to \varphi_{31}$
\item $\varphi_{94} :\equiv (\mathrm{Ind}(B_{3}) \leftrightarrow \emptyset \in B_{3} \land \forall y \in B_{3}\,(S(y) \in B_{3})) \to \varphi_{92}$
\item $\varphi_{95} :\equiv \forall B_{3}\,(\varphi_{52}) \to (\forall B_{3}\,(\varphi_{52}) \to \varphi_{60}) \to \varphi_{44} \to \forall n \in \omega\,(n \in \mathbb{N} \to 1 \leq \underline{n})$
\end{abbrevlist}
\end{fullproof}

% 自動生成 (prover/build.js): IotaPos  ノード数 16
\begin{fullproof}{[N15]}{$\forall n \in \mathbb{N}\,(0 < \underline{n})$}{16}\phantomsection\label{tree:IotaPos}
\noindent$\Pi_{1}$:
\begin{ptree}
\AXC{$\forall n \in \mathbb{N}\,(1 \leq \underline{n})$\lbl{[N14]}}
\AXC{$\forall n \in \mathbb{N}\,(1 \leq \underline{n}) \to n \in \mathbb{N} \to 1 \leq \underline{n}$\lbl{A2}}
\BIC{$n \in \mathbb{N} \to 1 \leq \underline{n}$}
\end{ptree}
\noindent$\Pi_{2}$:
\begin{ptree}
\AXC{$\forall x\,\forall y\,\forall z\,(x < y \to y \leq z \to x < z)$\lbl{[O3]}}
\AXC{$\varphi_{2}$\lbl{A2}}
\BIC{$\forall y\,\forall z\,(0 < y \to y \leq z \to 0 < z)$}
\AXC{$\varphi_{1}$\lbl{A2}}
\BIC{$\forall z\,(0 < 1 \to 1 \leq z \to 0 < z)$}
\AXC{$\forall z\,(0 < 1 \to 1 \leq z \to 0 < z) \to 0 < 1 \to 1 \leq \underline{n} \to 0 < \underline{n}$\lbl{A2}}
\BIC{$0 < 1 \to 1 \leq \underline{n} \to 0 < \underline{n}$}
\end{ptree}
\noindent 主木:
\begin{ptree}
\AXC{$\stackrel{\Pi_{1}}{\vdots}$}\noLine
\UIC{$n \in \mathbb{N} \to 1 \leq \underline{n}$}
\AXC{$0 < 1$\lbl{[F28]}}
\AXC{$\stackrel{\Pi_{2}}{\vdots}$}\noLine
\UIC{$0 < 1 \to 1 \leq \underline{n} \to 0 < \underline{n}$}
\AXC{$\varphi_{3}$\lbl{A1}}
\BIC{$0 < 1 \to (n \in \mathbb{N} \to 1 \leq \underline{n}) \to n \in \mathbb{N} \to 0 < \underline{n}$}
\BIC{$(n \in \mathbb{N} \to 1 \leq \underline{n}) \to n \in \mathbb{N} \to 0 < \underline{n}$}
\BIC{$n \in \mathbb{N} \to 0 < \underline{n}$}
\UIC{$\forall n \in \mathbb{N}\,(0 < \underline{n})$}
\end{ptree}
\begin{abbrevlist}
\item $\varphi_{1} :\equiv \forall y\,\forall z\,(0 < y \to y \leq z \to 0 < z) \to \forall z\,(0 < 1 \to 1 \leq z \to 0 < z)$
\item $\varphi_{2} :\equiv \forall x\,\forall y\,\forall z\,(x < y \to y \leq z \to x < z) \to \forall y\,\forall z\,(0 < y \to y \leq z \to 0 < z)$
\item $\varphi_{3} :\equiv (0 < 1 \to 1 \leq \underline{n} \to 0 < \underline{n}) \to 0 < 1 \to (n \in \mathbb{N} \to 1 \leq \underline{n}) \to n \in \mathbb{N} \to 0 < \underline{n}$
\end{abbrevlist}
\end{fullproof}

% 自動生成 (prover/build.js): IotaMono  ノード数 447
\begin{fullproof}{[N16]}{$\forall m\,\forall n\,(m \in \mathbb{N} \to n \in \mathbb{N} \to m \leq n \to \underline{m} \leq \underline{n})$}{447}\phantomsection\label{tree:IotaMono}
\noindent$\Pi_{1}$:
\begin{ptree}
\AXC{$\forall n\,(n \in \mathbb{N} \leftrightarrow n \in \omega \land n \neq \emptyset)$\lbl{Def$\mathbb{N}$}}
\AXC{$\varphi_{61}$\lbl{A2}}
\BIC{$n \in \mathbb{N} \leftrightarrow n \in \omega \land n \neq \emptyset$}
\AXC{$m \in \mathbb{N} \to n \in \mathbb{N} \to n \in \mathbb{N}$\lbl{A1}}
\AXC{$\varphi_{96}$\lbl{A1}}
\BIC{$(n \in \mathbb{N} \leftrightarrow n \in \omega \land n \neq \emptyset) \to m \in \mathbb{N} \to n \in \mathbb{N} \to n \in \omega$}
\BIC{$m \in \mathbb{N} \to n \in \mathbb{N} \to n \in \omega$}
\end{ptree}
\noindent$\Pi_{2}$:
\begin{ptree}
\AXC{$\varphi_{114}$\lbl{A2}}
\AXC{$m \in \mathbb{N} \to \varphi_{119}$\lbl{A1}}
\AXC{$(m \in \mathbb{N} \to \varphi_{119}) \to \varphi_{114} \to m \in \mathbb{N} \to \varphi_{114}$\lbl{A1}}
\BIC{$\varphi_{114} \to m \in \mathbb{N} \to \varphi_{114}$}
\BIC{$m \in \mathbb{N} \to \varphi_{114}$}
\end{ptree}
\noindent$\Pi_{3}$:
\begin{ptree}
\AXC{$x = x$\lbl{A5}}
\UIC{$\forall x\,(x = x)$}
\AXC{$\forall x\,(x = x) \to \underline{m} = \underline{m}$\lbl{A2}}
\BIC{$\underline{m} = \underline{m}$}
\AXC{$\underline{m} = \underline{m} \to m \in \mathbb{N} \to \varphi_{135}$\lbl{A1}}
\BIC{$m \in \mathbb{N} \to \varphi_{135}$}
\end{ptree}
\noindent$\Pi_{4}$:
\begin{ptree}
\AXC{$\stackrel{\Pi_{3}}{\vdots}$}\noLine
\UIC{$m \in \mathbb{N} \to \varphi_{135}$}
\AXC{$m \in \mathbb{N} \to \varphi_{123}$\lbl{A1}}
\AXC{$(m \in \mathbb{N} \to \varphi_{123}) \to (m \in \mathbb{N} \to \varphi_{135}) \to \varphi_{154}$\lbl{A1}}
\BIC{$(m \in \mathbb{N} \to \varphi_{135}) \to \varphi_{154}$}
\BIC{$\varphi_{154}$}
\end{ptree}
\noindent$\Pi_{5}$:
\begin{ptree}
\AXC{$z = w \to \underline{m} < \underline{z} \to \underline{m} < \underline{w}$\lbl{A6}}
\UIC{$\forall w\,(z = w \to \underline{m} < \underline{z} \to \underline{m} < \underline{w})$}
\UIC{$\forall z\,\forall w\,(z = w \to \underline{m} < \underline{z} \to \underline{m} < \underline{w})$}
\AXC{$\varphi_{82}$\lbl{A2}}
\BIC{$\forall w\,(m = w \to \underline{m} < \underline{m} \to \underline{m} < \underline{w})$}
\AXC{$\varphi_{83}$\lbl{A2}}
\BIC{$m = S(n) \to \underline{m} < \underline{m} \to \underline{m} < \underline{S(n)}$}
\AXC{$\stackrel{\Pi_{4}}{\vdots}$}\noLine
\UIC{$\varphi_{154}$}
\BIC{$(m = S(n) \to \underline{m} = \underline{m} \to \underline{m} = \underline{S(n)}) \to m \in \mathbb{N} \to \varphi_{138}$}
\end{ptree}
\noindent$\Pi_{6}$:
\begin{ptree}
\AXC{$z = w \to \underline{m} = \underline{z} \to \underline{m} = \underline{w}$\lbl{A6}}
\UIC{$\forall w\,(z = w \to \underline{m} = \underline{z} \to \underline{m} = \underline{w})$}
\UIC{$\forall z\,\forall w\,(z = w \to \underline{m} = \underline{z} \to \underline{m} = \underline{w})$}
\AXC{$\varphi_{71}$\lbl{A2}}
\BIC{$\forall w\,(m = w \to \underline{m} = \underline{m} \to \underline{m} = \underline{w})$}
\AXC{$\varphi_{72}$\lbl{A2}}
\BIC{$m = S(n) \to \underline{m} = \underline{m} \to \underline{m} = \underline{S(n)}$}
\AXC{$\stackrel{\Pi_{5}}{\vdots}$}\noLine
\UIC{$(m = S(n) \to \underline{m} = \underline{m} \to \underline{m} = \underline{S(n)}) \to m \in \mathbb{N} \to \varphi_{138}$}
\BIC{$m \in \mathbb{N} \to \varphi_{138}$}
\end{ptree}
\noindent$\Pi_{7}$:
\begin{ptree}
\AXC{$\varphi_{114}$\lbl{A2}}
\AXC{$m \in \mathbb{N} \to \varphi_{119}$\lbl{A1}}
\AXC{$(m \in \mathbb{N} \to \varphi_{119}) \to \varphi_{114} \to m \in \mathbb{N} \to \varphi_{114}$\lbl{A1}}
\BIC{$\varphi_{114} \to m \in \mathbb{N} \to \varphi_{114}$}
\BIC{$m \in \mathbb{N} \to \varphi_{114}$}
\AXC{$(m \in \mathbb{N} \to \varphi_{114}) \to (m \in \mathbb{N} \to \varphi_{64}) \to m \in \mathbb{N} \to \varphi_{115}$\lbl{A1}}
\BIC{$(m \in \mathbb{N} \to \varphi_{64}) \to m \in \mathbb{N} \to \varphi_{115}$}
\end{ptree}
\noindent$\Pi_{8}$:
\begin{ptree}
\AXC{$m \in \mathbb{N} \to \varphi_{64}$\lbl{A1}}
\AXC{$\stackrel{\Pi_{7}}{\vdots}$}\noLine
\UIC{$(m \in \mathbb{N} \to \varphi_{64}) \to m \in \mathbb{N} \to \varphi_{115}$}
\BIC{$m \in \mathbb{N} \to \varphi_{115}$}
\end{ptree}
\noindent$\Pi_{9}$:
\begin{ptree}
\AXC{$\stackrel{\Pi_{8}}{\vdots}$}\noLine
\UIC{$m \in \mathbb{N} \to \varphi_{115}$}
\AXC{$(m \in \mathbb{N} \to \varphi_{115}) \to m \in \mathbb{N} \to \varphi_{65}$\lbl{A1}}
\BIC{$m \in \mathbb{N} \to \varphi_{65}$}
\end{ptree}
\noindent$\Pi_{10}$:
\begin{ptree}
\AXC{$m \in \mathbb{N} \to \varphi_{108}$\lbl{A1}}
\AXC{$\stackrel{\Pi_{9}}{\vdots}$}\noLine
\UIC{$m \in \mathbb{N} \to \varphi_{65}$}
\AXC{$\forall n\,(n \in \mathbb{N} \leftrightarrow n \in \omega \land n \neq \emptyset)$\lbl{Def$\mathbb{N}$}}
\AXC{$\varphi_{61}$\lbl{A2}}
\BIC{$n \in \mathbb{N} \leftrightarrow n \in \omega \land n \neq \emptyset$}
\AXC{$\varphi_{164}$\lbl{A1}}
\BIC{$(m \in \mathbb{N} \to \varphi_{65}) \to (m \in \mathbb{N} \to \varphi_{108}) \to m \in \mathbb{N} \to \varphi_{102}$}
\BIC{$(m \in \mathbb{N} \to \varphi_{108}) \to m \in \mathbb{N} \to \varphi_{102}$}
\BIC{$m \in \mathbb{N} \to \varphi_{102}$}
\end{ptree}
\noindent$\Pi_{11}$:
\begin{ptree}
\AXC{$\forall n \in \mathbb{N}\,(\underline{n} \in \mathbb{R})$\lbl{[N4]}}
\AXC{$\forall n \in \mathbb{N}\,(\underline{n} \in \mathbb{R}) \to n \in \mathbb{N} \to \underline{n} \in \mathbb{R}$\lbl{A2}}
\BIC{$n \in \mathbb{N} \to \underline{n} \in \mathbb{R}$}
\end{ptree}
\noindent$\Pi_{12}$:
\begin{ptree}
\AXC{$\stackrel{\Pi_{10}}{\vdots}$}\noLine
\UIC{$m \in \mathbb{N} \to \varphi_{102}$}
\AXC{$\stackrel{\Pi_{11}}{\vdots}$}\noLine
\UIC{$n \in \mathbb{N} \to \underline{n} \in \mathbb{R}$}
\AXC{$(n \in \mathbb{N} \to \underline{n} \in \mathbb{R}) \to (m \in \mathbb{N} \to \varphi_{102}) \to m \in \mathbb{N} \to \varphi_{124}$\lbl{A1}}
\BIC{$(m \in \mathbb{N} \to \varphi_{102}) \to m \in \mathbb{N} \to \varphi_{124}$}
\BIC{$m \in \mathbb{N} \to \varphi_{124}$}
\end{ptree}
\noindent$\Pi_{13}$:
\begin{ptree}
\AXC{$\stackrel{\Pi_{12}}{\vdots}$}\noLine
\UIC{$m \in \mathbb{N} \to \varphi_{124}$}
\AXC{$\forall x \in \mathbb{R}\,(x < x + 1)$\lbl{[N13]}}
\AXC{$\forall x \in \mathbb{R}\,(x < x + 1) \to \underline{n} \in \mathbb{R} \to \underline{n} < \underline{n} + 1$\lbl{A2}}
\BIC{$\underline{n} \in \mathbb{R} \to \underline{n} < \underline{n} + 1$}
\AXC{$\varphi_{161}$\lbl{A1}}
\BIC{$(m \in \mathbb{N} \to \varphi_{124}) \to m \in \mathbb{N} \to \varphi_{131}$}
\BIC{$m \in \mathbb{N} \to \varphi_{131}$}
\end{ptree}
\noindent$\Pi_{14}$:
\begin{ptree}
\AXC{$x = x$\lbl{A5}}
\AXC{$x = y \to x = x \to y = x$\lbl{A6}}
\AXC{$(x = y \to x = x \to y = x) \to x = x \to x = y \to y = x$\lbl{A1}}
\BIC{$x = x \to x = y \to y = x$}
\BIC{$x = y \to y = x$}
\UIC{$\forall y\,(x = y \to y = x)$}
\UIC{$\forall x\,\forall y\,(x = y \to y = x)$}
\end{ptree}
\noindent$\Pi_{15}$:
\begin{ptree}
\AXC{$\stackrel{\Pi_{10}}{\vdots}$}\noLine
\UIC{$m \in \mathbb{N} \to \varphi_{102}$}
\AXC{$\forall n \in \mathbb{N}\,(\underline{S(n)} = \underline{n} + 1)$\lbl{[N6]}}
\AXC{$\varphi_{58}$\lbl{A2}}
\BIC{$n \in \mathbb{N} \to \underline{S(n)} = \underline{n} + 1$}
\AXC{$\varphi_{159}$\lbl{A1}}
\BIC{$(m \in \mathbb{N} \to \varphi_{102}) \to m \in \mathbb{N} \to \varphi_{132}$}
\BIC{$m \in \mathbb{N} \to \varphi_{132}$}
\AXC{$\varphi_{162}$\lbl{A1}}
\BIC{$(\underline{S(n)} = \underline{n} + 1 \to \underline{n} + 1 = \underline{S(n)}) \to m \in \mathbb{N} \to \varphi_{133}$}
\end{ptree}
\noindent$\Pi_{16}$:
\begin{ptree}
\AXC{$\stackrel{\Pi_{14}}{\vdots}$}\noLine
\UIC{$\forall x\,\forall y\,(x = y \to y = x)$}
\AXC{$\varphi_{8}$\lbl{A2}}
\BIC{$\forall y\,(\underline{S(n)} = y \to y = \underline{S(n)})$}
\AXC{$\varphi_{63}$\lbl{A2}}
\BIC{$\underline{S(n)} = \underline{n} + 1 \to \underline{n} + 1 = \underline{S(n)}$}
\AXC{$\stackrel{\Pi_{15}}{\vdots}$}\noLine
\UIC{$(\underline{S(n)} = \underline{n} + 1 \to \underline{n} + 1 = \underline{S(n)}) \to m \in \mathbb{N} \to \varphi_{133}$}
\BIC{$m \in \mathbb{N} \to \varphi_{133}$}
\end{ptree}
\noindent$\Pi_{17}$:
\begin{ptree}
\AXC{$\stackrel{\Pi_{13}}{\vdots}$}\noLine
\UIC{$m \in \mathbb{N} \to \varphi_{131}$}
\AXC{$\stackrel{\Pi_{16}}{\vdots}$}\noLine
\UIC{$m \in \mathbb{N} \to \varphi_{133}$}
\AXC{$(m \in \mathbb{N} \to \varphi_{133}) \to (m \in \mathbb{N} \to \varphi_{131}) \to \varphi_{151}$\lbl{A1}}
\BIC{$(m \in \mathbb{N} \to \varphi_{131}) \to \varphi_{151}$}
\BIC{$\varphi_{151}$}
\end{ptree}
\noindent$\Pi_{18}$:
\begin{ptree}
\AXC{$z = w \to \underline{n} < z \to \underline{n} < w$\lbl{A6}}
\UIC{$\forall w\,(z = w \to \underline{n} < z \to \underline{n} < w)$}
\UIC{$\forall z\,\forall w\,(z = w \to \underline{n} < z \to \underline{n} < w)$}
\AXC{$\varphi_{79}$\lbl{A2}}
\BIC{$\forall w\,(\underline{n} + 1 = w \to \underline{n} < \underline{n} + 1 \to \underline{n} < w)$}
\AXC{$\varphi_{128}$\lbl{A2}}
\BIC{$\underline{n} + 1 = \underline{S(n)} \to \underline{n} < \underline{n} + 1 \to \underline{n} < \underline{S(n)}$}
\AXC{$\stackrel{\Pi_{17}}{\vdots}$}\noLine
\UIC{$\varphi_{151}$}
\BIC{$m \in \mathbb{N} \to \varphi_{129}$}
\AXC{$(m \in \mathbb{N} \to \varphi_{129}) \to m \in \mathbb{N} \to \varphi_{126}$\lbl{A1}}
\BIC{$m \in \mathbb{N} \to \varphi_{126}$}
\end{ptree}
\noindent$\Pi_{19}$:
\begin{ptree}
\AXC{$\stackrel{\Pi_{10}}{\vdots}$}\noLine
\UIC{$m \in \mathbb{N} \to \varphi_{102}$}
\AXC{$\stackrel{\Pi_{8}}{\vdots}$}\noLine
\UIC{$m \in \mathbb{N} \to \varphi_{115}$}
\AXC{$(m \in \mathbb{N} \to \varphi_{115}) \to m \in \mathbb{N} \to \varphi_{94}$\lbl{A1}}
\BIC{$m \in \mathbb{N} \to \varphi_{94}$}
\AXC{$(m \in \mathbb{N} \to \varphi_{94}) \to (m \in \mathbb{N} \to \varphi_{102}) \to m \in \mathbb{N} \to \varphi_{120}$\lbl{A1}}
\BIC{$(m \in \mathbb{N} \to \varphi_{102}) \to m \in \mathbb{N} \to \varphi_{120}$}
\BIC{$m \in \mathbb{N} \to \varphi_{120}$}
\end{ptree}
\noindent$\Pi_{20}$:
\begin{ptree}
\AXC{$m \in \mathbb{N} \to \varphi_{116}$\lbl{A1}}
\AXC{$\stackrel{\Pi_{19}}{\vdots}$}\noLine
\UIC{$m \in \mathbb{N} \to \varphi_{120}$}
\AXC{$(m \in \mathbb{N} \to \varphi_{120}) \to (m \in \mathbb{N} \to \varphi_{116}) \to m \in \mathbb{N} \to \varphi_{130}$\lbl{A1}}
\BIC{$(m \in \mathbb{N} \to \varphi_{116}) \to m \in \mathbb{N} \to \varphi_{130}$}
\BIC{$m \in \mathbb{N} \to \varphi_{130}$}
\end{ptree}
\noindent$\Pi_{21}$:
\begin{ptree}
\AXC{$\forall x\,\forall y\,\forall z\,(x \leq y \to y \leq z \to x \leq z)$\lbl{[O5]}}
\AXC{$\varphi_{27}$\lbl{A2}}
\BIC{$\forall y\,\forall z\,(\underline{m} \leq y \to y \leq z \to \underline{m} \leq z)$}
\AXC{$\varphi_{59}$\lbl{A2}}
\BIC{$\forall z\,(\underline{m} \leq \underline{n} \to \underline{n} \leq z \to \underline{m} \leq z)$}
\AXC{$\varphi_{84}$\lbl{A2}}
\BIC{$\underline{m} \leq \underline{n} \to \underline{n} \leq \underline{S(n)} \to \underline{m} \leq \underline{S(n)}$}
\AXC{$\varphi_{170}$\lbl{A1}}
\BIC{$(m \in \mathbb{N} \to \varphi_{130}) \to (m \in \mathbb{N} \to \varphi_{126}) \to m \in \mathbb{N} \to \varphi_{136}$}
\end{ptree}
\noindent$\Pi_{22}$:
\begin{ptree}
\AXC{$\stackrel{\Pi_{18}}{\vdots}$}\noLine
\UIC{$m \in \mathbb{N} \to \varphi_{126}$}
\AXC{$\stackrel{\Pi_{20}}{\vdots}$}\noLine
\UIC{$m \in \mathbb{N} \to \varphi_{130}$}
\AXC{$\stackrel{\Pi_{21}}{\vdots}$}\noLine
\UIC{$(m \in \mathbb{N} \to \varphi_{130}) \to (m \in \mathbb{N} \to \varphi_{126}) \to m \in \mathbb{N} \to \varphi_{136}$}
\BIC{$(m \in \mathbb{N} \to \varphi_{126}) \to m \in \mathbb{N} \to \varphi_{136}$}
\BIC{$m \in \mathbb{N} \to \varphi_{136}$}
\end{ptree}
\noindent$\Pi_{23}$:
\begin{ptree}
\AXC{$\forall n \in \mathbb{N}\,(n \in \omega)$\lbl{[N7']}}
\AXC{$\forall n \in \mathbb{N}\,(n \in \omega) \to n \in \mathbb{N} \to n \in \omega$\lbl{A2}}
\BIC{$n \in \mathbb{N} \to n \in \omega$}
\AXC{$(n \in \mathbb{N} \to n \in \omega) \to (m \in \mathbb{N} \to \varphi_{102}) \to m \in \mathbb{N} \to \varphi_{117}$\lbl{A1}}
\BIC{$(m \in \mathbb{N} \to \varphi_{102}) \to m \in \mathbb{N} \to \varphi_{117}$}
\end{ptree}
\noindent$\Pi_{24}$:
\begin{ptree}
\AXC{$\forall n \in \mathbb{N}\,(n \in \omega)$\lbl{[N7']}}
\AXC{$\forall n \in \mathbb{N}\,(n \in \omega) \to m \in \mathbb{N} \to m \in \omega$\lbl{A2}}
\BIC{$m \in \mathbb{N} \to m \in \omega$}
\AXC{$(m \in \mathbb{N} \to m \in \omega) \to (m \in \mathbb{N} \to m \in \mathbb{N}) \to m \in \mathbb{N} \to \varphi_{118}$\lbl{A1}}
\BIC{$(m \in \mathbb{N} \to m \in \mathbb{N}) \to m \in \mathbb{N} \to \varphi_{118}$}
\end{ptree}
\noindent$\Pi_{25}$:
\begin{ptree}
\AXC{$\varphi_{19}$\lbl{[N9]}}
\AXC{$\varphi_{95}$\lbl{A2}}
\BIC{$\forall n\,(m \in \omega \to n \in \omega \to m \leq S(n) \to m \leq n \lor m = S(n))$}
\AXC{$\varphi_{89}$\lbl{A2}}
\BIC{$m \in \omega \to n \in \omega \to m \leq S(n) \to m \leq n \lor m = S(n)$}
\AXC{$(m \in \omega \to n \in \omega \to m \leq S(n) \to m \leq n \lor m = S(n)) \to \varphi_{171}$\lbl{A1}}
\BIC{$\varphi_{171}$}
\end{ptree}
\noindent$\Pi_{26}$:
\begin{ptree}
\AXC{$m \in \mathbb{N} \to m \in \mathbb{N}$\lbl{A1}}
\AXC{$\stackrel{\Pi_{24}}{\vdots}$}\noLine
\UIC{$(m \in \mathbb{N} \to m \in \mathbb{N}) \to m \in \mathbb{N} \to \varphi_{118}$}
\BIC{$m \in \mathbb{N} \to \varphi_{118}$}
\AXC{$\stackrel{\Pi_{25}}{\vdots}$}\noLine
\UIC{$\varphi_{171}$}
\BIC{$(m \in \mathbb{N} \to \varphi_{117}) \to (m \in \mathbb{N} \to \varphi_{111}) \to m \in \mathbb{N} \to \varphi_{121}$}
\end{ptree}
\noindent$\Pi_{27}$:
\begin{ptree}
\AXC{$m \in \mathbb{N} \to \varphi_{111}$\lbl{A1}}
\AXC{$\stackrel{\Pi_{10}}{\vdots}$}\noLine
\UIC{$m \in \mathbb{N} \to \varphi_{102}$}
\AXC{$\stackrel{\Pi_{23}}{\vdots}$}\noLine
\UIC{$(m \in \mathbb{N} \to \varphi_{102}) \to m \in \mathbb{N} \to \varphi_{117}$}
\BIC{$m \in \mathbb{N} \to \varphi_{117}$}
\AXC{$\stackrel{\Pi_{26}}{\vdots}$}\noLine
\UIC{$(m \in \mathbb{N} \to \varphi_{117}) \to (m \in \mathbb{N} \to \varphi_{111}) \to m \in \mathbb{N} \to \varphi_{121}$}
\BIC{$(m \in \mathbb{N} \to \varphi_{111}) \to m \in \mathbb{N} \to \varphi_{121}$}
\BIC{$m \in \mathbb{N} \to \varphi_{121}$}
\end{ptree}
\noindent$\Pi_{28}$:
\begin{ptree}
\AXC{$\stackrel{\Pi_{6}}{\vdots}$}\noLine
\UIC{$m \in \mathbb{N} \to \varphi_{138}$}
\AXC{$\stackrel{\Pi_{22}}{\vdots}$}\noLine
\UIC{$m \in \mathbb{N} \to \varphi_{136}$}
\AXC{$\stackrel{\Pi_{27}}{\vdots}$}\noLine
\UIC{$m \in \mathbb{N} \to \varphi_{121}$}
\AXC{$\varphi_{174}$\lbl{A1}}
\BIC{$(m \in \mathbb{N} \to \varphi_{136}) \to (m \in \mathbb{N} \to \varphi_{138}) \to m \in \mathbb{N} \to \varphi_{127}$}
\BIC{$(m \in \mathbb{N} \to \varphi_{138}) \to m \in \mathbb{N} \to \varphi_{127}$}
\BIC{$m \in \mathbb{N} \to \varphi_{127}$}
\end{ptree}
\noindent$\Pi_{29}$:
\begin{ptree}
\AXC{$y = z \to x = y \to x = z$\lbl{A6}}
\AXC{$(y = z \to x = y \to x = z) \to x = y \to y = z \to x = z$\lbl{A1}}
\BIC{$x = y \to y = z \to x = z$}
\UIC{$\forall z\,(x = y \to y = z \to x = z)$}
\UIC{$\forall y\,\forall z\,(x = y \to y = z \to x = z)$}
\UIC{$\forall x\,\forall y\,\forall z\,(x = y \to y = z \to x = z)$}
\AXC{$\varphi_{15}$\lbl{A2}}
\BIC{$\forall y\,\forall z\,(S(n) = y \to y = z \to S(n) = z)$}
\AXC{$\varphi_{55}$\lbl{A2}}
\BIC{$\forall z\,(S(n) = S(\emptyset) \to S(\emptyset) = z \to S(n) = z)$}
\end{ptree}
\noindent$\Pi_{30}$:
\begin{ptree}
\AXC{$\stackrel{\Pi_{29}}{\vdots}$}\noLine
\UIC{$\forall z\,(S(n) = S(\emptyset) \to S(\emptyset) = z \to S(n) = z)$}
\AXC{$\forall z\,(S(n) = S(\emptyset) \to S(\emptyset) = z \to S(n) = z) \to \varphi_{7}$\lbl{A2}}
\BIC{$\varphi_{7}$}
\end{ptree}
\noindent$\Pi_{31}$:
\begin{ptree}
\AXC{$\stackrel{\Pi_{14}}{\vdots}$}\noLine
\UIC{$\forall x\,\forall y\,(x = y \to y = x)$}
\AXC{$\varphi_{1}$\lbl{A2}}
\BIC{$\forall y\,(1_{\omega} = y \to y = 1_{\omega})$}
\AXC{$\varphi_{22}$\lbl{A2}}
\BIC{$1_{\omega} = S(\emptyset) \to S(\emptyset) = 1_{\omega}$}
\AXC{$1_{\omega} = S(\emptyset)$\lbl{Def$1_\omega$}}
\AXC{$1_{\omega} = S(\emptyset) \to \varphi_{16}$\lbl{A1}}
\BIC{$\varphi_{16}$}
\BIC{$S(\emptyset) = 1_{\omega}$}
\end{ptree}
\noindent$\Pi_{32}$:
\begin{ptree}
\AXC{$x = x$\lbl{A5}}
\UIC{$\forall x\,(x = x)$}
\AXC{$\forall x\,(x = x) \to S(n) = S(n)$\lbl{A2}}
\BIC{$S(n) = S(n)$}
\AXC{$m \in \mathbb{N} \to \varphi_{103}$\lbl{A1}}
\AXC{$(m \in \mathbb{N} \to \varphi_{103}) \to S(n) = S(n) \to \varphi_{143}$\lbl{A1}}
\BIC{$S(n) = S(n) \to \varphi_{143}$}
\BIC{$\varphi_{143}$}
\end{ptree}
\noindent$\Pi_{33}$:
\begin{ptree}
\AXC{$z = w \to S(n) = S(z) \to S(n) = S(w)$\lbl{A6}}
\UIC{$\forall w\,(z = w \to S(n) = S(z) \to S(n) = S(w))$}
\UIC{$\forall z\,\forall w\,(z = w \to S(n) = S(z) \to S(n) = S(w))$}
\AXC{$\varphi_{35}$\lbl{A2}}
\BIC{$\forall w\,(n = w \to S(n) = S(n) \to S(n) = S(w))$}
\AXC{$\varphi_{56}$\lbl{A2}}
\BIC{$n = \emptyset \to S(n) = S(n) \to S(n) = S(\emptyset)$}
\AXC{$\stackrel{\Pi_{32}}{\vdots}$}\noLine
\UIC{$\varphi_{143}$}
\BIC{$m \in \mathbb{N} \to \varphi_{112}$}
\end{ptree}
\noindent$\Pi_{34}$:
\begin{ptree}
\AXC{$\stackrel{\Pi_{30}}{\vdots}$}\noLine
\UIC{$\varphi_{7}$}
\AXC{$\stackrel{\Pi_{31}}{\vdots}$}\noLine
\UIC{$S(\emptyset) = 1_{\omega}$}
\AXC{$\stackrel{\Pi_{33}}{\vdots}$}\noLine
\UIC{$m \in \mathbb{N} \to \varphi_{112}$}
\AXC{$(m \in \mathbb{N} \to \varphi_{112}) \to S(\emptyset) = 1_{\omega} \to \varphi_{7} \to m \in \mathbb{N} \to \varphi_{100}$\lbl{A1}}
\BIC{$S(\emptyset) = 1_{\omega} \to \varphi_{7} \to m \in \mathbb{N} \to \varphi_{100}$}
\BIC{$\varphi_{7} \to m \in \mathbb{N} \to \varphi_{100}$}
\BIC{$m \in \mathbb{N} \to \varphi_{100}$}
\AXC{$\varphi_{156}$\lbl{A1}}
\BIC{$(S(n) = 1_{\omega} \to 1_{\omega} = S(n)) \to m \in \mathbb{N} \to \varphi_{101}$}
\end{ptree}
\noindent$\Pi_{35}$:
\begin{ptree}
\AXC{$\stackrel{\Pi_{14}}{\vdots}$}\noLine
\UIC{$\forall x\,\forall y\,(x = y \to y = x)$}
\AXC{$\varphi_{2}$\lbl{A2}}
\BIC{$\forall y\,(S(n) = y \to y = S(n))$}
\AXC{$\varphi_{4}$\lbl{A2}}
\BIC{$S(n) = 1_{\omega} \to 1_{\omega} = S(n)$}
\AXC{$\stackrel{\Pi_{34}}{\vdots}$}\noLine
\UIC{$(S(n) = 1_{\omega} \to 1_{\omega} = S(n)) \to m \in \mathbb{N} \to \varphi_{101}$}
\BIC{$m \in \mathbb{N} \to \varphi_{101}$}
\end{ptree}
\noindent$\Pi_{36}$:
\begin{ptree}
\AXC{$\stackrel{\Pi_{14}}{\vdots}$}\noLine
\UIC{$\forall x\,\forall y\,(x = y \to y = x)$}
\AXC{$\varphi_{1}$\lbl{A2}}
\BIC{$\forall y\,(1_{\omega} = y \to y = 1_{\omega})$}
\AXC{$\varphi_{3}$\lbl{A2}}
\BIC{$1_{\omega} = S(n) \to S(n) = 1_{\omega}$}
\AXC{$\stackrel{\Pi_{35}}{\vdots}$}\noLine
\UIC{$m \in \mathbb{N} \to \varphi_{101}$}
\AXC{$\varphi_{157}$\lbl{A1}}
\BIC{$(1_{\omega} = S(n) \to S(n) = 1_{\omega}) \to m \in \mathbb{N} \to \varphi_{100}$}
\BIC{$m \in \mathbb{N} \to \varphi_{100}$}
\end{ptree}
\noindent$\Pi_{37}$:
\begin{ptree}
\AXC{$\forall n \in \mathbb{N}\,(1_{\omega} \leq n)$\lbl{[N12]}}
\AXC{$\forall n \in \mathbb{N}\,(1_{\omega} \leq n) \to m \in \mathbb{N} \to 1_{\omega} \leq m$\lbl{A2}}
\BIC{$m \in \mathbb{N} \to 1_{\omega} \leq m$}
\AXC{$(m \in \mathbb{N} \to 1_{\omega} \leq m) \to \varphi_{11}$\lbl{A1}}
\BIC{$\varphi_{11}$}
\end{ptree}
\noindent$\Pi_{38}$:
\begin{ptree}
\AXC{$m \in \mathbb{N} \to m \leq 1_{\omega} \to m \leq 1_{\omega}$\lbl{A1}}
\AXC{$\forall x\,\forall y\,(x \leq y \to y \leq x \to x = y)$\lbl{[O6]}}
\AXC{$\varphi_{5}$\lbl{A2}}
\BIC{$\forall y\,(m \leq y \to y \leq m \to m = y)$}
\AXC{$\varphi_{6}$\lbl{A2}}
\BIC{$m \leq 1_{\omega} \to 1_{\omega} \leq m \to m = 1_{\omega}$}
\AXC{$\varphi_{106}$\lbl{A1}}
\BIC{$(m \in \mathbb{N} \to m \leq 1_{\omega} \to m \leq 1_{\omega}) \to \varphi_{25}$}
\BIC{$\varphi_{25}$}
\end{ptree}
\noindent$\Pi_{39}$:
\begin{ptree}
\AXC{$m \in \mathbb{N} \to m \in \mathbb{N}$\lbl{A1}}
\AXC{$\stackrel{\Pi_{37}}{\vdots}$}\noLine
\UIC{$\varphi_{11}$}
\BIC{$m \in \mathbb{N} \to m \leq 1_{\omega} \to 1_{\omega} \leq m$}
\AXC{$\stackrel{\Pi_{38}}{\vdots}$}\noLine
\UIC{$\varphi_{25}$}
\BIC{$m \in \mathbb{N} \to m \leq 1_{\omega} \to m = 1_{\omega}$}
\AXC{$(m \in \mathbb{N} \to m \leq 1_{\omega} \to m = 1_{\omega}) \to \varphi_{144}$\lbl{A1}}
\BIC{$\varphi_{144}$}
\end{ptree}
\noindent$\Pi_{40}$:
\begin{ptree}
\AXC{$x = x$\lbl{A5}}
\UIC{$\forall x\,(x = x)$}
\AXC{$\forall x\,(x = x) \to \underline{m} = \underline{m}$\lbl{A2}}
\BIC{$\underline{m} = \underline{m}$}
\AXC{$\underline{m} = \underline{m} \to m \in \mathbb{N} \to m \leq 1_{\omega} \to \underline{m} \leq \underline{m}$\lbl{A1}}
\BIC{$m \in \mathbb{N} \to m \leq 1_{\omega} \to \underline{m} \leq \underline{m}$}
\AXC{$\stackrel{\Pi_{39}}{\vdots}$}\noLine
\UIC{$\varphi_{144}$}
\BIC{$(m = 1_{\omega} \to \underline{m} < \underline{m} \to \underline{m} < \underline{1_{\omega}}) \to \varphi_{67}$}
\end{ptree}
\noindent$\Pi_{41}$:
\begin{ptree}
\AXC{$z = w \to \underline{m} < \underline{z} \to \underline{m} < \underline{w}$\lbl{A6}}
\UIC{$\forall w\,(z = w \to \underline{m} < \underline{z} \to \underline{m} < \underline{w})$}
\UIC{$\forall z\,\forall w\,(z = w \to \underline{m} < \underline{z} \to \underline{m} < \underline{w})$}
\AXC{$\varphi_{82}$\lbl{A2}}
\BIC{$\forall w\,(m = w \to \underline{m} < \underline{m} \to \underline{m} < \underline{w})$}
\AXC{$\varphi_{87}$\lbl{A2}}
\BIC{$m = 1_{\omega} \to \underline{m} < \underline{m} \to \underline{m} < \underline{1_{\omega}}$}
\AXC{$\stackrel{\Pi_{40}}{\vdots}$}\noLine
\UIC{$(m = 1_{\omega} \to \underline{m} < \underline{m} \to \underline{m} < \underline{1_{\omega}}) \to \varphi_{67}$}
\BIC{$\varphi_{67}$}
\end{ptree}
\noindent$\Pi_{42}$:
\begin{ptree}
\AXC{$z = w \to \underline{m} = \underline{z} \to \underline{m} = \underline{w}$\lbl{A6}}
\UIC{$\forall w\,(z = w \to \underline{m} = \underline{z} \to \underline{m} = \underline{w})$}
\UIC{$\forall z\,\forall w\,(z = w \to \underline{m} = \underline{z} \to \underline{m} = \underline{w})$}
\AXC{$\varphi_{71}$\lbl{A2}}
\BIC{$\forall w\,(m = w \to \underline{m} = \underline{m} \to \underline{m} = \underline{w})$}
\AXC{$\varphi_{73}$\lbl{A2}}
\BIC{$m = 1_{\omega} \to \underline{m} = \underline{m} \to \underline{m} = \underline{1_{\omega}}$}
\AXC{$\stackrel{\Pi_{41}}{\vdots}$}\noLine
\UIC{$\varphi_{67}$}
\BIC{$m \in \mathbb{N} \to m \leq 1_{\omega} \to \underline{m} \leq \underline{1_{\omega}}$}
\AXC{$\stackrel{\Pi_{35}}{\vdots}$}\noLine
\UIC{$m \in \mathbb{N} \to \varphi_{101}$}
\AXC{$\varphi_{177}$\lbl{A1}}
\BIC{$(m \in \mathbb{N} \to m \leq 1_{\omega} \to \underline{m} \leq \underline{1_{\omega}}) \to \varphi_{175}$}
\BIC{$\varphi_{175}$}
\end{ptree}
\noindent$\Pi_{43}$:
\begin{ptree}
\AXC{$z = w \to m < z \to m < w$\lbl{A6}}
\UIC{$\forall w\,(z = w \to m < z \to m < w)$}
\UIC{$\forall z\,\forall w\,(z = w \to m < z \to m < w)$}
\AXC{$\varphi_{29}$\lbl{A2}}
\BIC{$\forall w\,(1_{\omega} = w \to m < 1_{\omega} \to m < w)$}
\AXC{$\varphi_{37}$\lbl{A2}}
\BIC{$1_{\omega} = S(n) \to m < 1_{\omega} \to m < S(n)$}
\AXC{$\stackrel{\Pi_{42}}{\vdots}$}\noLine
\UIC{$\varphi_{175}$}
\BIC{$(1_{\omega} = S(n) \to m = 1_{\omega} \to m = S(n)) \to \varphi_{172}$}
\end{ptree}
\noindent$\Pi_{44}$:
\begin{ptree}
\AXC{$z = w \to m = z \to m = w$\lbl{A6}}
\UIC{$\forall w\,(z = w \to m = z \to m = w)$}
\UIC{$\forall z\,\forall w\,(z = w \to m = z \to m = w)$}
\AXC{$\varphi_{10}$\lbl{A2}}
\BIC{$\forall w\,(1_{\omega} = w \to m = 1_{\omega} \to m = w)$}
\AXC{$\varphi_{18}$\lbl{A2}}
\BIC{$1_{\omega} = S(n) \to m = 1_{\omega} \to m = S(n)$}
\AXC{$\stackrel{\Pi_{43}}{\vdots}$}\noLine
\UIC{$(1_{\omega} = S(n) \to m = 1_{\omega} \to m = S(n)) \to \varphi_{172}$}
\BIC{$\varphi_{172}$}
\end{ptree}
\noindent$\Pi_{45}$:
\begin{ptree}
\AXC{$z = w \to \underline{m} < \underline{z} \to \underline{m} < \underline{w}$\lbl{A6}}
\UIC{$\forall w\,(z = w \to \underline{m} < \underline{z} \to \underline{m} < \underline{w})$}
\UIC{$\forall z\,\forall w\,(z = w \to \underline{m} < \underline{z} \to \underline{m} < \underline{w})$}
\AXC{$\varphi_{92}$\lbl{A2}}
\BIC{$\forall w\,(1_{\omega} = w \to \underline{m} < \underline{1_{\omega}} \to \underline{m} < \underline{w})$}
\AXC{$\varphi_{105}$\lbl{A2}}
\BIC{$1_{\omega} = S(n) \to \underline{m} < \underline{1_{\omega}} \to \underline{m} < \underline{S(n)}$}
\AXC{$\stackrel{\Pi_{44}}{\vdots}$}\noLine
\UIC{$\varphi_{172}$}
\BIC{$\varphi_{169}$}
\end{ptree}
\noindent$\Pi_{46}$:
\begin{ptree}
\AXC{$z = w \to \underline{m} = \underline{z} \to \underline{m} = \underline{w}$\lbl{A6}}
\UIC{$\forall w\,(z = w \to \underline{m} = \underline{z} \to \underline{m} = \underline{w})$}
\UIC{$\forall z\,\forall w\,(z = w \to \underline{m} = \underline{z} \to \underline{m} = \underline{w})$}
\AXC{$\varphi_{75}$\lbl{A2}}
\BIC{$\forall w\,(1_{\omega} = w \to \underline{m} = \underline{1_{\omega}} \to \underline{m} = \underline{w})$}
\AXC{$\varphi_{86}$\lbl{A2}}
\BIC{$1_{\omega} = S(n) \to \underline{m} = \underline{1_{\omega}} \to \underline{m} = \underline{S(n)}$}
\AXC{$\stackrel{\Pi_{45}}{\vdots}$}\noLine
\UIC{$\varphi_{169}$}
\BIC{$\varphi_{165}$}
\end{ptree}
\noindent$\Pi_{47}$:
\begin{ptree}
\AXC{$z = w \to m < z \to m < w$\lbl{A6}}
\UIC{$\forall w\,(z = w \to m < z \to m < w)$}
\UIC{$\forall z\,\forall w\,(z = w \to m < z \to m < w)$}
\AXC{$\varphi_{28}$\lbl{A2}}
\BIC{$\forall w\,(S(n) = w \to m < S(n) \to m < w)$}
\AXC{$\varphi_{36}$\lbl{A2}}
\BIC{$S(n) = 1_{\omega} \to m < S(n) \to m < 1_{\omega}$}
\AXC{$\stackrel{\Pi_{46}}{\vdots}$}\noLine
\UIC{$\varphi_{165}$}
\BIC{$(S(n) = 1_{\omega} \to m = S(n) \to m = 1_{\omega}) \to \varphi_{163}$}
\end{ptree}
\noindent$\Pi_{48}$:
\begin{ptree}
\AXC{$z = w \to m = z \to m = w$\lbl{A6}}
\UIC{$\forall w\,(z = w \to m = z \to m = w)$}
\UIC{$\forall z\,\forall w\,(z = w \to m = z \to m = w)$}
\AXC{$\varphi_{9}$\lbl{A2}}
\BIC{$\forall w\,(S(n) = w \to m = S(n) \to m = w)$}
\AXC{$\varphi_{17}$\lbl{A2}}
\BIC{$S(n) = 1_{\omega} \to m = S(n) \to m = 1_{\omega}$}
\AXC{$\stackrel{\Pi_{47}}{\vdots}$}\noLine
\UIC{$(S(n) = 1_{\omega} \to m = S(n) \to m = 1_{\omega}) \to \varphi_{163}$}
\BIC{$\varphi_{163}$}
\end{ptree}
\noindent$\Pi_{49}$:
\begin{ptree}
\AXC{$z = w \to \underline{m} < \underline{z} \to \underline{m} < \underline{w}$\lbl{A6}}
\UIC{$\forall w\,(z = w \to \underline{m} < \underline{z} \to \underline{m} < \underline{w})$}
\UIC{$\forall z\,\forall w\,(z = w \to \underline{m} < \underline{z} \to \underline{m} < \underline{w})$}
\AXC{$\varphi_{91}$\lbl{A2}}
\BIC{$\forall w\,(S(n) = w \to \underline{m} < \underline{S(n)} \to \underline{m} < \underline{w})$}
\AXC{$\varphi_{104}$\lbl{A2}}
\BIC{$S(n) = 1_{\omega} \to \underline{m} < \underline{S(n)} \to \underline{m} < \underline{1_{\omega}}$}
\AXC{$\stackrel{\Pi_{48}}{\vdots}$}\noLine
\UIC{$\varphi_{163}$}
\BIC{$\varphi_{160}$}
\end{ptree}
\noindent$\Pi_{50}$:
\begin{ptree}
\AXC{$\stackrel{\Pi_{36}}{\vdots}$}\noLine
\UIC{$m \in \mathbb{N} \to \varphi_{100}$}
\AXC{$z = w \to \underline{m} = \underline{z} \to \underline{m} = \underline{w}$\lbl{A6}}
\UIC{$\forall w\,(z = w \to \underline{m} = \underline{z} \to \underline{m} = \underline{w})$}
\UIC{$\forall z\,\forall w\,(z = w \to \underline{m} = \underline{z} \to \underline{m} = \underline{w})$}
\AXC{$\varphi_{74}$\lbl{A2}}
\BIC{$\forall w\,(S(n) = w \to \underline{m} = \underline{S(n)} \to \underline{m} = \underline{w})$}
\AXC{$\varphi_{85}$\lbl{A2}}
\BIC{$S(n) = 1_{\omega} \to \underline{m} = \underline{S(n)} \to \underline{m} = \underline{1_{\omega}}$}
\AXC{$\stackrel{\Pi_{49}}{\vdots}$}\noLine
\UIC{$\varphi_{160}$}
\BIC{$(m \in \mathbb{N} \to \varphi_{100}) \to m \in \mathbb{N} \to \varphi_{125}$}
\BIC{$m \in \mathbb{N} \to \varphi_{125}$}
\end{ptree}
\noindent$\Pi_{51}$:
\begin{ptree}
\AXC{$\forall x\,(\mathrm{Ind}(x) \leftrightarrow \emptyset \in x \land \forall y \in x\,(S(y) \in x))$\lbl{DefInd}}
\AXC{$\varphi_{113}$\lbl{A2}}
\BIC{$\mathrm{Ind}(\omega) \leftrightarrow \emptyset \in \omega \land \forall y \in \omega\,(S(y) \in \omega)$}
\AXC{$\mathrm{Ind}(\omega)$\lbl{補題(ωは帰納的)}}
\AXC{$\mathrm{Ind}(\omega) \to \varphi_{148}$\lbl{A1}}
\BIC{$\varphi_{148}$}
\BIC{$m \in \mathbb{N} \to \varphi_{109}$}
\end{ptree}
\noindent$\Pi_{52}$:
\begin{ptree}
\AXC{$\stackrel{\Pi_{51}}{\vdots}$}\noLine
\UIC{$m \in \mathbb{N} \to \varphi_{109}$}
\AXC{$(m \in \mathbb{N} \to \varphi_{109}) \to m \in \mathbb{N} \to \varphi_{90}$\lbl{A1}}
\BIC{$m \in \mathbb{N} \to \varphi_{90}$}
\AXC{$(m \in \mathbb{N} \to \varphi_{90}) \to \varphi_{145}$\lbl{A1}}
\BIC{$\varphi_{145}$}
\end{ptree}
\noindent$\Pi_{53}$:
\begin{ptree}
\AXC{$\stackrel{\Pi_{9}}{\vdots}$}\noLine
\UIC{$m \in \mathbb{N} \to \varphi_{65}$}
\AXC{$\forall y \in \omega\,(S(y) \in \omega) \to n \in \omega \to S(n) \in \omega$\lbl{A2}}
\AXC{$\stackrel{\Pi_{52}}{\vdots}$}\noLine
\UIC{$\varphi_{145}$}
\BIC{$m \in \mathbb{N} \to \varphi_{88}$}
\AXC{$(m \in \mathbb{N} \to \varphi_{88}) \to (m \in \mathbb{N} \to \varphi_{65}) \to m \in \mathbb{N} \to \varphi_{68}$\lbl{A1}}
\BIC{$(m \in \mathbb{N} \to \varphi_{65}) \to m \in \mathbb{N} \to \varphi_{68}$}
\BIC{$m \in \mathbb{N} \to \varphi_{68}$}
\end{ptree}
\noindent$\Pi_{54}$:
\begin{ptree}
\AXC{$\stackrel{\Pi_{53}}{\vdots}$}\noLine
\UIC{$m \in \mathbb{N} \to \varphi_{68}$}
\AXC{$\varphi_{134}$\lbl{A2}}
\AXC{$m \in \mathbb{N} \to \varphi_{119}$\lbl{A1}}
\AXC{$(m \in \mathbb{N} \to \varphi_{119}) \to \varphi_{134} \to m \in \mathbb{N} \to \varphi_{134}$\lbl{A1}}
\BIC{$\varphi_{134} \to m \in \mathbb{N} \to \varphi_{134}$}
\BIC{$m \in \mathbb{N} \to \varphi_{134}$}
\AXC{$\varphi_{168}$\lbl{A1}}
\BIC{$(m \in \mathbb{N} \to \varphi_{68}) \to (m \in \mathbb{N} \to \varphi_{110}) \to m \in \mathbb{N} \to \varphi_{66}$}
\BIC{$(m \in \mathbb{N} \to \varphi_{110}) \to m \in \mathbb{N} \to \varphi_{66}$}
\end{ptree}
\noindent$\Pi_{55}$:
\begin{ptree}
\AXC{$\stackrel{\Pi_{28}}{\vdots}$}\noLine
\UIC{$m \in \mathbb{N} \to \varphi_{127}$}
\AXC{$\stackrel{\Pi_{50}}{\vdots}$}\noLine
\UIC{$m \in \mathbb{N} \to \varphi_{125}$}
\AXC{$(m \in \mathbb{N} \to \varphi_{125}) \to (m \in \mathbb{N} \to \varphi_{127}) \to m \in \mathbb{N} \to \varphi_{110}$\lbl{A1}}
\BIC{$(m \in \mathbb{N} \to \varphi_{127}) \to m \in \mathbb{N} \to \varphi_{110}$}
\BIC{$m \in \mathbb{N} \to \varphi_{110}$}
\AXC{$\stackrel{\Pi_{54}}{\vdots}$}\noLine
\UIC{$(m \in \mathbb{N} \to \varphi_{110}) \to m \in \mathbb{N} \to \varphi_{66}$}
\BIC{$m \in \mathbb{N} \to \varphi_{66}$}
\end{ptree}
\noindent$\Pi_{56}$:
\begin{ptree}
\AXC{$\varphi_{53} \to \forall n\,(\varphi_{53})$\lbl{A4}}
\AXC{$\varphi_{150}$\lbl{A3}}
\AXC{$\stackrel{\Pi_{55}}{\vdots}$}\noLine
\UIC{$m \in \mathbb{N} \to \varphi_{66}$}
\AXC{$(m \in \mathbb{N} \to \varphi_{66}) \to \varphi_{53} \to n \in B_{5} \to S(n) \in B_{5}$\lbl{A1}}
\BIC{$\varphi_{53} \to n \in B_{5} \to S(n) \in B_{5}$}
\UIC{$\forall n\,(\varphi_{53} \to n \in B_{5} \to S(n) \in B_{5})$}
\AXC{$\varphi_{173}$\lbl{A1}}
\BIC{$\varphi_{150} \to (\varphi_{53} \to \forall n\,(\varphi_{53})) \to m \in \mathbb{N} \to \varphi_{69}$}
\BIC{$(\varphi_{53} \to \forall n\,(\varphi_{53})) \to m \in \mathbb{N} \to \varphi_{69}$}
\BIC{$m \in \mathbb{N} \to \varphi_{69}$}
\end{ptree}
\noindent$\Pi_{57}$:
\begin{ptree}
\AXC{$x = x$\lbl{A5}}
\UIC{$\forall x\,(x = x)$}
\AXC{$\forall x\,(x = x) \to \emptyset = \emptyset$\lbl{A2}}
\BIC{$\emptyset = \emptyset$}
\AXC{$\forall n\,(n \in \mathbb{N} \leftrightarrow n \in \omega \land n \neq \emptyset)$\lbl{Def$\mathbb{N}$}}
\AXC{$\varphi_{80}$\lbl{A2}}
\BIC{$\emptyset \in \mathbb{N} \leftrightarrow \emptyset \in \omega \land \emptyset \neq \emptyset$}
\AXC{$\varphi_{139}$\lbl{A1}}
\BIC{$\emptyset = \emptyset \to m \in \mathbb{N} \to \emptyset \in \mathbb{N} \to m \leq \emptyset \to \underline{m} \leq \underline{\emptyset}$}
\BIC{$m \in \mathbb{N} \to \emptyset \in \mathbb{N} \to m \leq \emptyset \to \underline{m} \leq \underline{\emptyset}$}
\end{ptree}
\noindent$\Pi_{58}$:
\begin{ptree}
\AXC{$\stackrel{\Pi_{51}}{\vdots}$}\noLine
\UIC{$m \in \mathbb{N} \to \varphi_{109}$}
\AXC{$\varphi_{141}$\lbl{A2}}
\AXC{$m \in \mathbb{N} \to \varphi_{119}$\lbl{A1}}
\AXC{$(m \in \mathbb{N} \to \varphi_{119}) \to \varphi_{141} \to m \in \mathbb{N} \to \varphi_{141}$\lbl{A1}}
\BIC{$\varphi_{141} \to m \in \mathbb{N} \to \varphi_{141}$}
\BIC{$m \in \mathbb{N} \to \varphi_{141}$}
\AXC{$(m \in \mathbb{N} \to \varphi_{141}) \to \varphi_{158}$\lbl{A1}}
\BIC{$\varphi_{158}$}
\BIC{$(m \in \mathbb{N} \to \emptyset \in \mathbb{N} \to m \leq \emptyset \to \underline{m} \leq \underline{\emptyset}) \to m \in \mathbb{N} \to \varphi_{60}$}
\end{ptree}
\noindent$\Pi_{59}$:
\begin{ptree}
\AXC{$\stackrel{\Pi_{57}}{\vdots}$}\noLine
\UIC{$m \in \mathbb{N} \to \emptyset \in \mathbb{N} \to m \leq \emptyset \to \underline{m} \leq \underline{\emptyset}$}
\AXC{$\stackrel{\Pi_{58}}{\vdots}$}\noLine
\UIC{$(m \in \mathbb{N} \to \emptyset \in \mathbb{N} \to m \leq \emptyset \to \underline{m} \leq \underline{\emptyset}) \to m \in \mathbb{N} \to \varphi_{60}$}
\BIC{$m \in \mathbb{N} \to \varphi_{60}$}
\end{ptree}
\noindent$\Pi_{60}$:
\begin{ptree}
\AXC{$\stackrel{\Pi_{56}}{\vdots}$}\noLine
\UIC{$m \in \mathbb{N} \to \varphi_{69}$}
\AXC{$\stackrel{\Pi_{59}}{\vdots}$}\noLine
\UIC{$m \in \mathbb{N} \to \varphi_{60}$}
\AXC{$\forall x\,(\mathrm{Ind}(x) \leftrightarrow \emptyset \in x \land \forall y \in x\,(S(y) \in x))$\lbl{DefInd}}
\AXC{$\varphi_{93}$\lbl{A2}}
\BIC{$\mathrm{Ind}(B_{5}) \leftrightarrow \emptyset \in B_{5} \land \forall y \in B_{5}\,(S(y) \in B_{5})$}
\AXC{$\varphi_{166}$\lbl{A1}}
\BIC{$(m \in \mathbb{N} \to \varphi_{60}) \to (m \in \mathbb{N} \to \varphi_{69}) \to \varphi_{147}$}
\BIC{$(m \in \mathbb{N} \to \varphi_{69}) \to \varphi_{147}$}
\BIC{$\varphi_{147}$}
\end{ptree}
\noindent$\Pi_{61}$:
\begin{ptree}
\AXC{$\varphi_{21}$\lbl{A4}}
\AXC{$\varphi_{13} \to \varphi_{24}$\lbl{A3}}
\AXC{$\forall y \in B_{5}\,(S(y) \in B_{5}) \to n \in B_{5} \to S(n) \in B_{5}$\lbl{A2}}
\UIC{$\varphi_{13}$}
\AXC{$\varphi_{13} \to (\varphi_{13} \to \varphi_{24}) \to \varphi_{99}$\lbl{A1}}
\BIC{$(\varphi_{13} \to \varphi_{24}) \to \varphi_{99}$}
\BIC{$\varphi_{99}$}
\BIC{$\forall y \in B_{5}\,(S(y) \in B_{5}) \to \forall n \in B_{5}\,(S(n) \in B_{5})$}
\AXC{$\stackrel{\Pi_{60}}{\vdots}$}\noLine
\UIC{$\varphi_{147}$}
\BIC{$\varphi_{137}$}
\end{ptree}
\noindent$\Pi_{62}$:
\begin{ptree}
\AXC{$\varphi_{20}$\lbl{A4}}
\AXC{$\varphi_{12} \to \varphi_{23}$\lbl{A3}}
\AXC{$\forall n \in B_{5}\,(S(n) \in B_{5}) \to y \in B_{5} \to S(y) \in B_{5}$\lbl{A2}}
\UIC{$\varphi_{12}$}
\AXC{$\varphi_{12} \to (\varphi_{12} \to \varphi_{23}) \to \varphi_{98}$\lbl{A1}}
\BIC{$(\varphi_{12} \to \varphi_{23}) \to \varphi_{98}$}
\BIC{$\varphi_{98}$}
\BIC{$\forall n \in B_{5}\,(S(n) \in B_{5}) \to \forall y \in B_{5}\,(S(y) \in B_{5})$}
\AXC{$\stackrel{\Pi_{61}}{\vdots}$}\noLine
\UIC{$\varphi_{137}$}
\BIC{$m \in \mathbb{N} \to \varphi_{46}$}
\end{ptree}
\noindent$\Pi_{63}$:
\begin{ptree}
\AXC{$\forall n\,(n \in \omega \leftrightarrow \forall x\,(\mathrm{Ind}(x) \to n \in x))$\lbl{Def$\omega$}}
\AXC{$\varphi_{48}$\lbl{A2}}
\BIC{$n \in \omega \leftrightarrow \forall x\,(\mathrm{Ind}(x) \to n \in x)$}
\AXC{$m \in \mathbb{N} \to \varphi_{70}$\lbl{A1}}
\AXC{$\varphi_{153}$\lbl{A1}}
\BIC{$(n \in \omega \leftrightarrow \forall x\,(\mathrm{Ind}(x) \to n \in x)) \to m \in \mathbb{N} \to \varphi_{77}$}
\BIC{$m \in \mathbb{N} \to \varphi_{77}$}
\end{ptree}
\noindent$\Pi_{64}$:
\begin{ptree}
\AXC{$\forall x\,(\mathrm{Ind}(x) \to n \in x) \to \mathrm{Ind}(B_{5}) \to n \in B_{5}$\lbl{A2}}
\AXC{$\stackrel{\Pi_{63}}{\vdots}$}\noLine
\UIC{$m \in \mathbb{N} \to \varphi_{77}$}
\AXC{$(m \in \mathbb{N} \to \varphi_{77}) \to \varphi_{140}$\lbl{A1}}
\BIC{$\varphi_{140}$}
\BIC{$m \in \mathbb{N} \to \varphi_{76}$}
\AXC{$\varphi_{167}$\lbl{A1}}
\BIC{$(m \in \mathbb{N} \to \varphi_{46}) \to (m \in \mathbb{N} \to \varphi_{114}) \to m \in \mathbb{N} \to \varphi_{97}$}
\end{ptree}
\noindent$\Pi_{65}$:
\begin{ptree}
\AXC{$\stackrel{\Pi_{2}}{\vdots}$}\noLine
\UIC{$m \in \mathbb{N} \to \varphi_{114}$}
\AXC{$\stackrel{\Pi_{62}}{\vdots}$}\noLine
\UIC{$m \in \mathbb{N} \to \varphi_{46}$}
\AXC{$\stackrel{\Pi_{64}}{\vdots}$}\noLine
\UIC{$(m \in \mathbb{N} \to \varphi_{46}) \to (m \in \mathbb{N} \to \varphi_{114}) \to m \in \mathbb{N} \to \varphi_{97}$}
\BIC{$(m \in \mathbb{N} \to \varphi_{114}) \to m \in \mathbb{N} \to \varphi_{97}$}
\BIC{$m \in \mathbb{N} \to \varphi_{97}$}
\end{ptree}
\noindent$\Pi_{66}$:
\begin{ptree}
\AXC{$\varphi_{53} \to \forall n\,(\varphi_{53})$\lbl{A4}}
\AXC{$\varphi_{155}$\lbl{A3}}
\AXC{$\stackrel{\Pi_{65}}{\vdots}$}\noLine
\UIC{$m \in \mathbb{N} \to \varphi_{97}$}
\AXC{$(m \in \mathbb{N} \to \varphi_{97}) \to \varphi_{53} \to n \in \omega \to n \in \mathbb{N} \to m \leq n \to \underline{m} \leq \underline{n}$\lbl{A1}}
\BIC{$\varphi_{53} \to n \in \omega \to n \in \mathbb{N} \to m \leq n \to \underline{m} \leq \underline{n}$}
\UIC{$\forall n\,(\varphi_{53} \to n \in \omega \to n \in \mathbb{N} \to m \leq n \to \underline{m} \leq \underline{n})$}
\AXC{$\varphi_{176}$\lbl{A1}}
\BIC{$\varphi_{155} \to (\varphi_{53} \to \forall n\,(\varphi_{53})) \to m \in \mathbb{N} \to \varphi_{107}$}
\BIC{$(\varphi_{53} \to \forall n\,(\varphi_{53})) \to m \in \mathbb{N} \to \varphi_{107}$}
\BIC{$m \in \mathbb{N} \to \varphi_{107}$}
\end{ptree}
\noindent$\Pi_{67}$:
\begin{ptree}
\AXC{$\varphi_{49}$\lbl{Sep}}
\AXC{$\varphi_{49} \to \varphi_{38}$\lbl{A2}}
\BIC{$\varphi_{38}$}
\AXC{$\varphi_{38} \to \varphi_{33}$\lbl{A2}}
\BIC{$\varphi_{33}$}
\AXC{$\varphi_{33} \to m \in \mathbb{N} \to \varphi_{33}$\lbl{A1}}
\BIC{$m \in \mathbb{N} \to \varphi_{33}$}
\end{ptree}
\noindent$\Pi_{68}$:
\begin{ptree}
\AXC{$\stackrel{\Pi_{66}}{\vdots}$}\noLine
\UIC{$m \in \mathbb{N} \to \varphi_{107}$}
\AXC{$(m \in \mathbb{N} \to \varphi_{107}) \to \varphi_{122}$\lbl{A1}}
\BIC{$\varphi_{122}$}
\UIC{$\forall B_{5}\,(\varphi_{122})$}
\AXC{$\stackrel{\Pi_{67}}{\vdots}$}\noLine
\UIC{$m \in \mathbb{N} \to \varphi_{33}$}
\AXC{$(m \in \mathbb{N} \to \varphi_{33}) \to \forall B_{5}\,(\varphi_{122}) \to (\forall B_{5}\,(\varphi_{122}) \to \varphi_{41} \to \varphi_{39}) \to \varphi_{149}$\lbl{A1}}
\BIC{$\forall B_{5}\,(\varphi_{122}) \to (\forall B_{5}\,(\varphi_{122}) \to \varphi_{41} \to \varphi_{39}) \to \varphi_{149}$}
\BIC{$(\forall B_{5}\,(\varphi_{122}) \to \varphi_{41} \to \varphi_{39}) \to \varphi_{149}$}
\end{ptree}
\noindent$\Pi_{69}$:
\begin{ptree}
\AXC{$m \in \mathbb{N} \land \neg \forall n \in \omega\,(n \in \mathbb{N} \to m \leq n \to \underline{m} \leq \underline{n}) \to \varphi_{41}$\lbl{A4}}
\AXC{$\forall B_{5}\,(\varphi_{122}) \to \varphi_{41} \to \varphi_{39}$\lbl{A3}}
\AXC{$\stackrel{\Pi_{68}}{\vdots}$}\noLine
\UIC{$(\forall B_{5}\,(\varphi_{122}) \to \varphi_{41} \to \varphi_{39}) \to \varphi_{149}$}
\BIC{$\varphi_{149}$}
\BIC{$m \in \mathbb{N} \to \forall n \in \omega\,(n \in \mathbb{N} \to m \leq n \to \underline{m} \leq \underline{n})$}
\AXC{$\varphi_{152}$\lbl{A1}}
\BIC{$\varphi_{78} \to m \in \mathbb{N} \to n \in \omega \to n \in \mathbb{N} \to m \leq n \to \underline{m} \leq \underline{n}$}
\end{ptree}
\noindent$\Pi_{70}$:
\begin{ptree}
\AXC{$\stackrel{\Pi_{1}}{\vdots}$}\noLine
\UIC{$m \in \mathbb{N} \to n \in \mathbb{N} \to n \in \omega$}
\AXC{$\varphi_{78}$\lbl{A2}}
\AXC{$\stackrel{\Pi_{69}}{\vdots}$}\noLine
\UIC{$\varphi_{78} \to m \in \mathbb{N} \to n \in \omega \to n \in \mathbb{N} \to m \leq n \to \underline{m} \leq \underline{n}$}
\BIC{$m \in \mathbb{N} \to n \in \omega \to n \in \mathbb{N} \to m \leq n \to \underline{m} \leq \underline{n}$}
\AXC{$\varphi_{146}$\lbl{A1}}
\BIC{$(m \in \mathbb{N} \to n \in \mathbb{N} \to n \in \omega) \to \varphi_{62}$}
\BIC{$\varphi_{62}$}
\end{ptree}
\noindent$\Pi_{71}$:
\begin{ptree}
\AXC{$\varphi_{81}$\lbl{A3}}
\AXC{$m \in \mathbb{N} \to n \in \mathbb{N} \to n \in \mathbb{N}$\lbl{A1}}
\AXC{$\stackrel{\Pi_{70}}{\vdots}$}\noLine
\UIC{$\varphi_{62}$}
\BIC{$m \in \mathbb{N} \to n \in \mathbb{N} \to m \leq n \to \underline{m} \leq \underline{n}$}
\UIC{$\forall n\,(m \in \mathbb{N} \to n \in \mathbb{N} \to m \leq n \to \underline{m} \leq \underline{n})$}
\AXC{$\forall n\,(m \in \mathbb{N} \to n \in \mathbb{N} \to m \leq n \to \underline{m} \leq \underline{n}) \to \varphi_{81} \to \varphi_{57}$\lbl{A1}}
\BIC{$\varphi_{81} \to \varphi_{57}$}
\BIC{$\varphi_{57}$}
\end{ptree}
\noindent 主木:
\begin{ptree}
\AXC{$\varphi_{47}$\lbl{A2}}
\AXC{$m \in \mathbb{N} \to \forall n\,(m \in \mathbb{N})$\lbl{A4}}
\AXC{$\stackrel{\Pi_{71}}{\vdots}$}\noLine
\UIC{$\varphi_{57}$}
\BIC{$m \in \mathbb{N} \to \forall n \in \mathbb{N}\,(m \leq n \to \underline{m} \leq \underline{n})$}
\AXC{$\varphi_{142}$\lbl{A1}}
\BIC{$\varphi_{47} \to m \in \mathbb{N} \to n \in \mathbb{N} \to m \leq n \to \underline{m} \leq \underline{n}$}
\BIC{$m \in \mathbb{N} \to n \in \mathbb{N} \to m \leq n \to \underline{m} \leq \underline{n}$}
\UIC{$\forall n\,(m \in \mathbb{N} \to n \in \mathbb{N} \to m \leq n \to \underline{m} \leq \underline{n})$}
\UIC{$\forall m\,\forall n\,(m \in \mathbb{N} \to n \in \mathbb{N} \to m \leq n \to \underline{m} \leq \underline{n})$}
\end{ptree}
\begin{abbrevlist}
\item $\varphi_{1} :\equiv \forall x\,\forall y\,(x = y \to y = x) \to \forall y\,(1_{\omega} = y \to y = 1_{\omega})$
\item $\varphi_{2} :\equiv \forall x\,\forall y\,(x = y \to y = x) \to \forall y\,(S(n) = y \to y = S(n))$
\item $\varphi_{3} :\equiv \forall y\,(1_{\omega} = y \to y = 1_{\omega}) \to 1_{\omega} = S(n) \to S(n) = 1_{\omega}$
\item $\varphi_{4} :\equiv \forall y\,(S(n) = y \to y = S(n)) \to S(n) = 1_{\omega} \to 1_{\omega} = S(n)$
\item $\varphi_{5} :\equiv \forall x\,\forall y\,(x \leq y \to y \leq x \to x = y) \to \forall y\,(m \leq y \to y \leq m \to m = y)$
\item $\varphi_{6} :\equiv \forall y\,(m \leq y \to y \leq m \to m = y) \to m \leq 1_{\omega} \to 1_{\omega} \leq m \to m = 1_{\omega}$
\item $\varphi_{7} :\equiv S(n) = S(\emptyset) \to S(\emptyset) = 1_{\omega} \to S(n) = 1_{\omega}$
\item $\varphi_{8} :\equiv \forall x\,\forall y\,(x = y \to y = x) \to \forall y\,(\underline{S(n)} = y \to y = \underline{S(n)})$
\item $\varphi_{9} :\equiv \forall z\,\forall w\,(z = w \to m = z \to m = w) \to \forall w\,(S(n) = w \to m = S(n) \to m = w)$
\item $\varphi_{10} :\equiv \forall z\,\forall w\,(z = w \to m = z \to m = w) \to \forall w\,(1_{\omega} = w \to m = 1_{\omega} \to m = w)$
\item $\varphi_{11} :\equiv (m \in \mathbb{N} \to m \in \mathbb{N}) \to m \in \mathbb{N} \to m \leq 1_{\omega} \to 1_{\omega} \leq m$
\item $\varphi_{12} :\equiv \forall y\,(\forall n \in B_{5}\,(S(n) \in B_{5}) \to y \in B_{5} \to S(y) \in B_{5})$
\item $\varphi_{13} :\equiv \forall n\,(\forall y \in B_{5}\,(S(y) \in B_{5}) \to n \in B_{5} \to S(n) \in B_{5})$
\item $\varphi_{14} :\equiv S(n) \in \mathbb{N} \to n \neq \emptyset \to m \leq S(n) \to m = S(n) \to m = S(n)$
\item $\varphi_{15} :\equiv \forall x\,\forall y\,\forall z\,(x = y \to y = z \to x = z) \to \forall y\,\forall z\,(S(n) = y \to y = z \to S(n) = z)$
\item $\varphi_{16} :\equiv (1_{\omega} = S(\emptyset) \to S(\emptyset) = 1_{\omega}) \to S(\emptyset) = 1_{\omega}$
\item $\varphi_{17} :\equiv \forall w\,(S(n) = w \to m = S(n) \to m = w) \to S(n) = 1_{\omega} \to m = S(n) \to m = 1_{\omega}$
\item $\varphi_{18} :\equiv \forall w\,(1_{\omega} = w \to m = 1_{\omega} \to m = w) \to 1_{\omega} = S(n) \to m = 1_{\omega} \to m = S(n)$
\item $\varphi_{19} :\equiv \forall m\,\forall n\,(m \in \omega \to n \in \omega \to m \leq S(n) \to m \leq n \lor m = S(n))$
\item $\varphi_{20} :\equiv \forall n \in B_{5}\,(S(n) \in B_{5}) \to \forall y\,\forall n \in B_{5}\,(S(n) \in B_{5})$
\item $\varphi_{21} :\equiv \forall y \in B_{5}\,(S(y) \in B_{5}) \to \forall n\,\forall y \in B_{5}\,(S(y) \in B_{5})$
\item $\varphi_{22} :\equiv \forall y\,(1_{\omega} = y \to y = 1_{\omega}) \to 1_{\omega} = S(\emptyset) \to S(\emptyset) = 1_{\omega}$
\item $\varphi_{23} :\equiv \forall y\,\forall n \in B_{5}\,(S(n) \in B_{5}) \to \forall y \in B_{5}\,(S(y) \in B_{5})$
\item $\varphi_{24} :\equiv \forall n\,\forall y \in B_{5}\,(S(y) \in B_{5}) \to \forall n \in B_{5}\,(S(n) \in B_{5})$
\item $\varphi_{25} :\equiv (m \in \mathbb{N} \to m \leq 1_{\omega} \to 1_{\omega} \leq m) \to m \in \mathbb{N} \to m \leq 1_{\omega} \to m = 1_{\omega}$
\item $\varphi_{26} :\equiv n \in B_{5} \to S(n) \in \mathbb{N} \to n \neq \emptyset \to m \leq S(n) \to m \leq n \to m \leq n$
\item $\varphi_{27} :\equiv \forall x\,\forall y\,\forall z\,(x \leq y \to y \leq z \to x \leq z) \to \forall y\,\forall z\,(\underline{m} \leq y \to y \leq z \to \underline{m} \leq z)$
\item $\varphi_{28} :\equiv \forall z\,\forall w\,(z = w \to m < z \to m < w) \to \forall w\,(S(n) = w \to m < S(n) \to m < w)$
\item $\varphi_{29} :\equiv \forall z\,\forall w\,(z = w \to m < z \to m < w) \to \forall w\,(1_{\omega} = w \to m < 1_{\omega} \to m < w)$
\item $\varphi_{30} :\equiv n \in B_{5} \to S(n) \in \mathbb{N} \to n \neq \emptyset \to m \leq S(n) \to m \leq n \lor m = S(n)$
\item $\varphi_{31} :\equiv S(n) \in \mathbb{N} \to n \neq \emptyset \to m \leq S(n) \to m = S(n) \to \underline{m} \leq \underline{m}$
\item $\varphi_{32} :\equiv S(n) \in \mathbb{N} \to n \neq \emptyset \to m \leq S(n) \to m \leq n \to \underline{m} \leq \underline{S(n)}$
\item $\varphi_{33} :\equiv \exists B_{5}\,\forall n\,(n \in B_{5} \leftrightarrow n \in \omega \land (n \in \mathbb{N} \to m \leq n \to \underline{m} \leq \underline{n}))$
\item $\varphi_{34} :\equiv S(n) \in \mathbb{N} \to n \neq \emptyset \to m \leq S(n) \to m = S(n) \to \underline{m} \leq \underline{S(n)}$
\item $\varphi_{35} :\equiv \forall z\,\forall w\,(z = w \to S(n) = S(z) \to S(n) = S(w)) \to \forall w\,(n = w \to S(n) = S(n) \to S(n) = S(w))$
\item $\varphi_{36} :\equiv \forall w\,(S(n) = w \to m < S(n) \to m < w) \to S(n) = 1_{\omega} \to m < S(n) \to m < 1_{\omega}$
\item $\varphi_{37} :\equiv \forall w\,(1_{\omega} = w \to m < 1_{\omega} \to m < w) \to 1_{\omega} = S(n) \to m < 1_{\omega} \to m < S(n)$
\item $\varphi_{38} :\equiv \forall A_{6}\,\exists B_{5}\,\forall n\,(n \in B_{5} \leftrightarrow n \in A_{6} \land (n \in \mathbb{N} \to m \leq n \to \underline{m} \leq \underline{n}))$
\item $\varphi_{39} :\equiv \forall B_{5}\,(\neg \forall n\,(n \in B_{5} \leftrightarrow n \in \omega \land (n \in \mathbb{N} \to m \leq n \to \underline{m} \leq \underline{n})))$
\item $\varphi_{40} :\equiv n \in B_{5} \to S(n) \in \mathbb{N} \to n = \emptyset \to m \leq S(n) \to \underline{m} \leq \underline{S(n)}$
\item $\varphi_{41} :\equiv \forall B_{5}\,(m \in \mathbb{N} \land \neg \forall n \in \omega\,(n \in \mathbb{N} \to m \leq n \to \underline{m} \leq \underline{n}))$
\item $\varphi_{42} :\equiv n \in B_{5} \to S(n) \in \mathbb{N} \to n \neq \emptyset \to m \leq S(n) \to \underline{n} \leq \underline{S(n)}$
\item $\varphi_{43} :\equiv n \in B_{5} \to S(n) \in \mathbb{N} \to n \neq \emptyset \to m \leq S(n) \to \underline{m} \leq \underline{S(n)}$
\item $\varphi_{44} :\equiv n \in B_{5} \to S(n) \in \mathbb{N} \to n \neq \emptyset \to m \leq S(n) \to \underline{n} < \underline{S(n)}$
\item $\varphi_{45} :\equiv n \in B_{5} \to S(n) \in \mathbb{N} \to n \neq \emptyset \to m \leq S(n) \to m \leq n \to \underline{m} \leq \underline{n}$
\item $\varphi_{46} :\equiv \forall n\,(n \in B_{5} \leftrightarrow n \in \omega \land (n \in \mathbb{N} \to m \leq n \to \underline{m} \leq \underline{n})) \to \mathrm{Ind}(B_{5})$
\item $\varphi_{47} :\equiv \forall n \in \mathbb{N}\,(m \leq n \to \underline{m} \leq \underline{n}) \to n \in \mathbb{N} \to m \leq n \to \underline{m} \leq \underline{n}$
\item $\varphi_{48} :\equiv \forall n\,(n \in \omega \leftrightarrow \forall x\,(\mathrm{Ind}(x) \to n \in x)) \to (n \in \omega \leftrightarrow \forall x\,(\mathrm{Ind}(x) \to n \in x))$
\item $\varphi_{49} :\equiv \forall z_{7}\,\forall A_{6}\,\exists B_{5}\,\forall n\,(n \in B_{5} \leftrightarrow n \in A_{6} \land (n \in \mathbb{N} \to z_{7} \leq n \to \underline{z_{7}} \leq \underline{n}))$
\item $\varphi_{50} :\equiv n \in B_{5} \to S(n) \in \mathbb{N} \to n \neq \emptyset \to m \leq S(n) \to \underline{n} < \underline{n} + 1$
\item $\varphi_{51} :\equiv n \in B_{5} \to S(n) \in \mathbb{N} \to n \neq \emptyset \to m \leq S(n) \to \underline{S(n)} = \underline{n} + 1$
\item $\varphi_{52} :\equiv n \in B_{5} \to S(n) \in \mathbb{N} \to n \neq \emptyset \to m \leq S(n) \to \underline{n} + 1 = \underline{S(n)}$
\item $\varphi_{53} :\equiv m \in \mathbb{N} \land \forall n\,(n \in B_{5} \leftrightarrow n \in \omega \land (n \in \mathbb{N} \to m \leq n \to \underline{m} \leq \underline{n}))$
\item $\varphi_{54} :\equiv S(n) \in B_{5} \leftrightarrow S(n) \in \omega \land (S(n) \in \mathbb{N} \to m \leq S(n) \to \underline{m} \leq \underline{S(n)})$
\item $\varphi_{55} :\equiv \forall y\,\forall z\,(S(n) = y \to y = z \to S(n) = z) \to \forall z\,(S(n) = S(\emptyset) \to S(\emptyset) = z \to S(n) = z)$
\item $\varphi_{56} :\equiv \forall w\,(n = w \to S(n) = S(n) \to S(n) = S(w)) \to n = \emptyset \to S(n) = S(n) \to S(n) = S(\emptyset)$
\item $\varphi_{57} :\equiv (m \in \mathbb{N} \to \forall n\,(m \in \mathbb{N})) \to m \in \mathbb{N} \to \forall n \in \mathbb{N}\,(m \leq n \to \underline{m} \leq \underline{n})$
\item $\varphi_{58} :\equiv \forall n \in \mathbb{N}\,(\underline{S(n)} = \underline{n} + 1) \to n \in \mathbb{N} \to \underline{S(n)} = \underline{n} + 1$
\item $\varphi_{59} :\equiv \forall y\,\forall z\,(\underline{m} \leq y \to y \leq z \to \underline{m} \leq z) \to \forall z\,(\underline{m} \leq \underline{n} \to \underline{n} \leq z \to \underline{m} \leq z)$
\item $\varphi_{60} :\equiv \forall n\,(n \in B_{5} \leftrightarrow n \in \omega \land (n \in \mathbb{N} \to m \leq n \to \underline{m} \leq \underline{n})) \to \emptyset \in B_{5}$
\item $\varphi_{61} :\equiv \forall n\,(n \in \mathbb{N} \leftrightarrow n \in \omega \land n \neq \emptyset) \to (n \in \mathbb{N} \leftrightarrow n \in \omega \land n \neq \emptyset)$
\item $\varphi_{62} :\equiv (m \in \mathbb{N} \to n \in \mathbb{N} \to n \in \mathbb{N}) \to m \in \mathbb{N} \to n \in \mathbb{N} \to m \leq n \to \underline{m} \leq \underline{n}$
\item $\varphi_{63} :\equiv \forall y\,(\underline{S(n)} = y \to y = \underline{S(n)}) \to \underline{S(n)} = \underline{n} + 1 \to \underline{n} + 1 = \underline{S(n)}$
\item $\varphi_{64} :\equiv \forall n\,(n \in B_{5} \leftrightarrow n \in \omega \land (n \in \mathbb{N} \to m \leq n \to \underline{m} \leq \underline{n})) \to n \in B_{5} \to n \in B_{5}$
\item $\varphi_{65} :\equiv \forall n\,(n \in B_{5} \leftrightarrow n \in \omega \land (n \in \mathbb{N} \to m \leq n \to \underline{m} \leq \underline{n})) \to n \in B_{5} \to n \in \omega$
\item $\varphi_{66} :\equiv \forall n\,(n \in B_{5} \leftrightarrow n \in \omega \land (n \in \mathbb{N} \to m \leq n \to \underline{m} \leq \underline{n})) \to n \in B_{5} \to S(n) \in B_{5}$
\item $\varphi_{67} :\equiv (m = 1_{\omega} \to \underline{m} = \underline{m} \to \underline{m} = \underline{1_{\omega}}) \to m \in \mathbb{N} \to m \leq 1_{\omega} \to \underline{m} \leq \underline{1_{\omega}}$
\item $\varphi_{68} :\equiv \forall n\,(n \in B_{5} \leftrightarrow n \in \omega \land (n \in \mathbb{N} \to m \leq n \to \underline{m} \leq \underline{n})) \to n \in B_{5} \to S(n) \in \omega$
\item $\varphi_{69} :\equiv \forall n\,(n \in B_{5} \leftrightarrow n \in \omega \land (n \in \mathbb{N} \to m \leq n \to \underline{m} \leq \underline{n})) \to \forall n \in B_{5}\,(S(n) \in B_{5})$
\item $\varphi_{70} :\equiv \forall n\,(n \in B_{5} \leftrightarrow n \in \omega \land (n \in \mathbb{N} \to m \leq n \to \underline{m} \leq \underline{n})) \to n \in \omega \to n \in \omega$
\item $\varphi_{71} :\equiv \forall z\,\forall w\,(z = w \to \underline{m} = \underline{z} \to \underline{m} = \underline{w}) \to \forall w\,(m = w \to \underline{m} = \underline{m} \to \underline{m} = \underline{w})$
\item $\varphi_{72} :\equiv \forall w\,(m = w \to \underline{m} = \underline{m} \to \underline{m} = \underline{w}) \to m = S(n) \to \underline{m} = \underline{m} \to \underline{m} = \underline{S(n)}$
\item $\varphi_{73} :\equiv \forall w\,(m = w \to \underline{m} = \underline{m} \to \underline{m} = \underline{w}) \to m = 1_{\omega} \to \underline{m} = \underline{m} \to \underline{m} = \underline{1_{\omega}}$
\item $\varphi_{74} :\equiv \forall z\,\forall w\,(z = w \to \underline{m} = \underline{z} \to \underline{m} = \underline{w}) \to \forall w\,(S(n) = w \to \underline{m} = \underline{S(n)} \to \underline{m} = \underline{w})$
\item $\varphi_{75} :\equiv \forall z\,\forall w\,(z = w \to \underline{m} = \underline{z} \to \underline{m} = \underline{w}) \to \forall w\,(1_{\omega} = w \to \underline{m} = \underline{1_{\omega}} \to \underline{m} = \underline{w})$
\item $\varphi_{76} :\equiv \forall n\,(n \in B_{5} \leftrightarrow n \in \omega \land (n \in \mathbb{N} \to m \leq n \to \underline{m} \leq \underline{n})) \to n \in \omega \to \mathrm{Ind}(B_{5}) \to n \in B_{5}$
\item $\varphi_{77} :\equiv \forall n\,(n \in B_{5} \leftrightarrow n \in \omega \land (n \in \mathbb{N} \to m \leq n \to \underline{m} \leq \underline{n})) \to n \in \omega \to \forall x\,(\mathrm{Ind}(x) \to n \in x)$
\item $\varphi_{78} :\equiv \forall n \in \omega\,(n \in \mathbb{N} \to m \leq n \to \underline{m} \leq \underline{n}) \to n \in \omega \to n \in \mathbb{N} \to m \leq n \to \underline{m} \leq \underline{n}$
\item $\varphi_{79} :\equiv \forall z\,\forall w\,(z = w \to \underline{n} < z \to \underline{n} < w) \to \forall w\,(\underline{n} + 1 = w \to \underline{n} < \underline{n} + 1 \to \underline{n} < w)$
\item $\varphi_{80} :\equiv \forall n\,(n \in \mathbb{N} \leftrightarrow n \in \omega \land n \neq \emptyset) \to (\emptyset \in \mathbb{N} \leftrightarrow \emptyset \in \omega \land \emptyset \neq \emptyset)$
\item $\varphi_{81} :\equiv \forall n\,(m \in \mathbb{N} \to n \in \mathbb{N} \to m \leq n \to \underline{m} \leq \underline{n}) \to \forall n\,(m \in \mathbb{N}) \to \forall n \in \mathbb{N}\,(m \leq n \to \underline{m} \leq \underline{n})$
\item $\varphi_{82} :\equiv \forall z\,\forall w\,(z = w \to \underline{m} < \underline{z} \to \underline{m} < \underline{w}) \to \forall w\,(m = w \to \underline{m} < \underline{m} \to \underline{m} < \underline{w})$
\item $\varphi_{83} :\equiv \forall w\,(m = w \to \underline{m} < \underline{m} \to \underline{m} < \underline{w}) \to m = S(n) \to \underline{m} < \underline{m} \to \underline{m} < \underline{S(n)}$
\item $\varphi_{84} :\equiv \forall z\,(\underline{m} \leq \underline{n} \to \underline{n} \leq z \to \underline{m} \leq z) \to \underline{m} \leq \underline{n} \to \underline{n} \leq \underline{S(n)} \to \underline{m} \leq \underline{S(n)}$
\item $\varphi_{85} :\equiv \forall w\,(S(n) = w \to \underline{m} = \underline{S(n)} \to \underline{m} = \underline{w}) \to S(n) = 1_{\omega} \to \underline{m} = \underline{S(n)} \to \underline{m} = \underline{1_{\omega}}$
\item $\varphi_{86} :\equiv \forall w\,(1_{\omega} = w \to \underline{m} = \underline{1_{\omega}} \to \underline{m} = \underline{w}) \to 1_{\omega} = S(n) \to \underline{m} = \underline{1_{\omega}} \to \underline{m} = \underline{S(n)}$
\item $\varphi_{87} :\equiv \forall w\,(m = w \to \underline{m} < \underline{m} \to \underline{m} < \underline{w}) \to m = 1_{\omega} \to \underline{m} < \underline{m} \to \underline{m} < \underline{1_{\omega}}$
\item $\varphi_{88} :\equiv \forall n\,(n \in B_{5} \leftrightarrow n \in \omega \land (n \in \mathbb{N} \to m \leq n \to \underline{m} \leq \underline{n})) \to n \in B_{5} \to n \in \omega \to S(n) \in \omega$
\item $\varphi_{89} :\equiv \forall n\,(m \in \omega \to n \in \omega \to m \leq S(n) \to m \leq n \lor m = S(n)) \to m \in \omega \to n \in \omega \to m \leq S(n) \to m \leq n \lor m = S(n)$
\item $\varphi_{90} :\equiv \forall n\,(n \in B_{5} \leftrightarrow n \in \omega \land (n \in \mathbb{N} \to m \leq n \to \underline{m} \leq \underline{n})) \to n \in B_{5} \to \forall y \in \omega\,(S(y) \in \omega)$
\item $\varphi_{91} :\equiv \forall z\,\forall w\,(z = w \to \underline{m} < \underline{z} \to \underline{m} < \underline{w}) \to \forall w\,(S(n) = w \to \underline{m} < \underline{S(n)} \to \underline{m} < \underline{w})$
\item $\varphi_{92} :\equiv \forall z\,\forall w\,(z = w \to \underline{m} < \underline{z} \to \underline{m} < \underline{w}) \to \forall w\,(1_{\omega} = w \to \underline{m} < \underline{1_{\omega}} \to \underline{m} < \underline{w})$
\item $\varphi_{93} :\equiv \forall x\,(\mathrm{Ind}(x) \leftrightarrow \emptyset \in x \land \forall y \in x\,(S(y) \in x)) \to (\mathrm{Ind}(B_{5}) \leftrightarrow \emptyset \in B_{5} \land \forall y \in B_{5}\,(S(y) \in B_{5}))$
\item $\varphi_{94} :\equiv \forall n\,(n \in B_{5} \leftrightarrow n \in \omega \land (n \in \mathbb{N} \to m \leq n \to \underline{m} \leq \underline{n})) \to n \in B_{5} \to n \in \mathbb{N} \to m \leq n \to \underline{m} \leq \underline{n}$
\item $\varphi_{95} :\equiv \varphi_{19} \to \forall n\,(m \in \omega \to n \in \omega \to m \leq S(n) \to m \leq n \lor m = S(n))$
\item $\varphi_{96} :\equiv (m \in \mathbb{N} \to n \in \mathbb{N} \to n \in \mathbb{N}) \to (n \in \mathbb{N} \leftrightarrow n \in \omega \land n \neq \emptyset) \to m \in \mathbb{N} \to n \in \mathbb{N} \to n \in \omega$
\item $\varphi_{97} :\equiv \forall n\,(n \in B_{5} \leftrightarrow n \in \omega \land (n \in \mathbb{N} \to m \leq n \to \underline{m} \leq \underline{n})) \to n \in \omega \to n \in \mathbb{N} \to m \leq n \to \underline{m} \leq \underline{n}$
\item $\varphi_{98} :\equiv \varphi_{20} \to \forall n \in B_{5}\,(S(n) \in B_{5}) \to \forall y \in B_{5}\,(S(y) \in B_{5})$
\item $\varphi_{99} :\equiv \varphi_{21} \to \forall y \in B_{5}\,(S(y) \in B_{5}) \to \forall n \in B_{5}\,(S(n) \in B_{5})$
\item $\varphi_{100} :\equiv \forall n\,(n \in B_{5} \leftrightarrow n \in \omega \land (n \in \mathbb{N} \to m \leq n \to \underline{m} \leq \underline{n})) \to n \in B_{5} \to S(n) \in \mathbb{N} \to n = \emptyset \to S(n) = 1_{\omega}$
\item $\varphi_{101} :\equiv \forall n\,(n \in B_{5} \leftrightarrow n \in \omega \land (n \in \mathbb{N} \to m \leq n \to \underline{m} \leq \underline{n})) \to n \in B_{5} \to S(n) \in \mathbb{N} \to n = \emptyset \to 1_{\omega} = S(n)$
\item $\varphi_{102} :\equiv \forall n\,(n \in B_{5} \leftrightarrow n \in \omega \land (n \in \mathbb{N} \to m \leq n \to \underline{m} \leq \underline{n})) \to n \in B_{5} \to S(n) \in \mathbb{N} \to n \neq \emptyset \to n \in \mathbb{N}$
\item $\varphi_{103} :\equiv \forall n\,(n \in B_{5} \leftrightarrow n \in \omega \land (n \in \mathbb{N} \to m \leq n \to \underline{m} \leq \underline{n})) \to n \in B_{5} \to S(n) \in \mathbb{N} \to n = \emptyset \to n = \emptyset$
\item $\varphi_{104} :\equiv \forall w\,(S(n) = w \to \underline{m} < \underline{S(n)} \to \underline{m} < \underline{w}) \to S(n) = 1_{\omega} \to \underline{m} < \underline{S(n)} \to \underline{m} < \underline{1_{\omega}}$
\item $\varphi_{105} :\equiv \forall w\,(1_{\omega} = w \to \underline{m} < \underline{1_{\omega}} \to \underline{m} < \underline{w}) \to 1_{\omega} = S(n) \to \underline{m} < \underline{1_{\omega}} \to \underline{m} < \underline{S(n)}$
\item $\varphi_{106} :\equiv (m \leq 1_{\omega} \to 1_{\omega} \leq m \to m = 1_{\omega}) \to (m \in \mathbb{N} \to m \leq 1_{\omega} \to m \leq 1_{\omega}) \to \varphi_{25}$
\item $\varphi_{107} :\equiv \forall n\,(n \in B_{5} \leftrightarrow n \in \omega \land (n \in \mathbb{N} \to m \leq n \to \underline{m} \leq \underline{n})) \to \forall n \in \omega\,(n \in \mathbb{N} \to m \leq n \to \underline{m} \leq \underline{n})$
\item $\varphi_{108} :\equiv \forall n\,(n \in B_{5} \leftrightarrow n \in \omega \land (n \in \mathbb{N} \to m \leq n \to \underline{m} \leq \underline{n})) \to n \in B_{5} \to S(n) \in \mathbb{N} \to n \neq \emptyset \to n \neq \emptyset$
\item $\varphi_{109} :\equiv \forall n\,(n \in B_{5} \leftrightarrow n \in \omega \land (n \in \mathbb{N} \to m \leq n \to \underline{m} \leq \underline{n})) \to \emptyset \in \omega \land \forall y \in \omega\,(S(y) \in \omega)$
\item $\varphi_{110} :\equiv \forall n\,(n \in B_{5} \leftrightarrow n \in \omega \land (n \in \mathbb{N} \to m \leq n \to \underline{m} \leq \underline{n})) \to n \in B_{5} \to S(n) \in \mathbb{N} \to m \leq S(n) \to \underline{m} \leq \underline{S(n)}$
\item $\varphi_{111} :\equiv \forall n\,(n \in B_{5} \leftrightarrow n \in \omega \land (n \in \mathbb{N} \to m \leq n \to \underline{m} \leq \underline{n})) \to n \in B_{5} \to S(n) \in \mathbb{N} \to n \neq \emptyset \to m \leq S(n) \to m \leq S(n)$
\item $\varphi_{112} :\equiv \forall n\,(n \in B_{5} \leftrightarrow n \in \omega \land (n \in \mathbb{N} \to m \leq n \to \underline{m} \leq \underline{n})) \to n \in B_{5} \to S(n) \in \mathbb{N} \to n = \emptyset \to S(n) = S(\emptyset)$
\item $\varphi_{113} :\equiv \forall x\,(\mathrm{Ind}(x) \leftrightarrow \emptyset \in x \land \forall y \in x\,(S(y) \in x)) \to (\mathrm{Ind}(\omega) \leftrightarrow \emptyset \in \omega \land \forall y \in \omega\,(S(y) \in \omega))$
\item $\varphi_{114} :\equiv \forall n\,(n \in B_{5} \leftrightarrow n \in \omega \land (n \in \mathbb{N} \to m \leq n \to \underline{m} \leq \underline{n})) \to (n \in B_{5} \leftrightarrow n \in \omega \land (n \in \mathbb{N} \to m \leq n \to \underline{m} \leq \underline{n}))$
\item $\varphi_{115} :\equiv \forall n\,(n \in B_{5} \leftrightarrow n \in \omega \land (n \in \mathbb{N} \to m \leq n \to \underline{m} \leq \underline{n})) \to n \in B_{5} \to n \in \omega \land (n \in \mathbb{N} \to m \leq n \to \underline{m} \leq \underline{n})$
\item $\varphi_{116} :\equiv \forall n\,(n \in B_{5} \leftrightarrow n \in \omega \land (n \in \mathbb{N} \to m \leq n \to \underline{m} \leq \underline{n})) \to \varphi_{26}$
\item $\varphi_{117} :\equiv \forall n\,(n \in B_{5} \leftrightarrow n \in \omega \land (n \in \mathbb{N} \to m \leq n \to \underline{m} \leq \underline{n})) \to n \in B_{5} \to S(n) \in \mathbb{N} \to n \neq \emptyset \to m \leq S(n) \to n \in \omega$
\item $\varphi_{118} :\equiv \forall n\,(n \in B_{5} \leftrightarrow n \in \omega \land (n \in \mathbb{N} \to m \leq n \to \underline{m} \leq \underline{n})) \to n \in B_{5} \to S(n) \in \mathbb{N} \to n \neq \emptyset \to m \leq S(n) \to m \in \omega$
\item $\varphi_{119} :\equiv \forall n\,(n \in B_{5} \leftrightarrow n \in \omega \land (n \in \mathbb{N} \to m \leq n \to \underline{m} \leq \underline{n})) \to \forall n\,(n \in B_{5} \leftrightarrow n \in \omega \land (n \in \mathbb{N} \to m \leq n \to \underline{m} \leq \underline{n}))$
\item $\varphi_{120} :\equiv \forall n\,(n \in B_{5} \leftrightarrow n \in \omega \land (n \in \mathbb{N} \to m \leq n \to \underline{m} \leq \underline{n})) \to n \in B_{5} \to S(n) \in \mathbb{N} \to n \neq \emptyset \to m \leq n \to \underline{m} \leq \underline{n}$
\item $\varphi_{121} :\equiv \forall n\,(n \in B_{5} \leftrightarrow n \in \omega \land (n \in \mathbb{N} \to m \leq n \to \underline{m} \leq \underline{n})) \to \varphi_{30}$
\item $\varphi_{122} :\equiv m \in \mathbb{N} \land \neg \forall n \in \omega\,(n \in \mathbb{N} \to m \leq n \to \underline{m} \leq \underline{n}) \to \neg \forall n\,(n \in B_{5} \leftrightarrow n \in \omega \land (n \in \mathbb{N} \to m \leq n \to \underline{m} \leq \underline{n}))$
\item $\varphi_{123} :\equiv \forall n\,(n \in B_{5} \leftrightarrow n \in \omega \land (n \in \mathbb{N} \to m \leq n \to \underline{m} \leq \underline{n})) \to n \in B_{5} \to \varphi_{14}$
\item $\varphi_{124} :\equiv \forall n\,(n \in B_{5} \leftrightarrow n \in \omega \land (n \in \mathbb{N} \to m \leq n \to \underline{m} \leq \underline{n})) \to n \in B_{5} \to S(n) \in \mathbb{N} \to n \neq \emptyset \to m \leq S(n) \to \underline{n} \in \mathbb{R}$
\item $\varphi_{125} :\equiv \forall n\,(n \in B_{5} \leftrightarrow n \in \omega \land (n \in \mathbb{N} \to m \leq n \to \underline{m} \leq \underline{n})) \to \varphi_{40}$
\item $\varphi_{126} :\equiv \forall n\,(n \in B_{5} \leftrightarrow n \in \omega \land (n \in \mathbb{N} \to m \leq n \to \underline{m} \leq \underline{n})) \to \varphi_{42}$
\item $\varphi_{127} :\equiv \forall n\,(n \in B_{5} \leftrightarrow n \in \omega \land (n \in \mathbb{N} \to m \leq n \to \underline{m} \leq \underline{n})) \to \varphi_{43}$
\item $\varphi_{128} :\equiv \forall w\,(\underline{n} + 1 = w \to \underline{n} < \underline{n} + 1 \to \underline{n} < w) \to \underline{n} + 1 = \underline{S(n)} \to \underline{n} < \underline{n} + 1 \to \underline{n} < \underline{S(n)}$
\item $\varphi_{129} :\equiv \forall n\,(n \in B_{5} \leftrightarrow n \in \omega \land (n \in \mathbb{N} \to m \leq n \to \underline{m} \leq \underline{n})) \to \varphi_{44}$
\item $\varphi_{130} :\equiv \forall n\,(n \in B_{5} \leftrightarrow n \in \omega \land (n \in \mathbb{N} \to m \leq n \to \underline{m} \leq \underline{n})) \to \varphi_{45}$
\item $\varphi_{131} :\equiv \forall n\,(n \in B_{5} \leftrightarrow n \in \omega \land (n \in \mathbb{N} \to m \leq n \to \underline{m} \leq \underline{n})) \to \varphi_{50}$
\item $\varphi_{132} :\equiv \forall n\,(n \in B_{5} \leftrightarrow n \in \omega \land (n \in \mathbb{N} \to m \leq n \to \underline{m} \leq \underline{n})) \to \varphi_{51}$
\item $\varphi_{133} :\equiv \forall n\,(n \in B_{5} \leftrightarrow n \in \omega \land (n \in \mathbb{N} \to m \leq n \to \underline{m} \leq \underline{n})) \to \varphi_{52}$
\item $\varphi_{134} :\equiv \forall n\,(n \in B_{5} \leftrightarrow n \in \omega \land (n \in \mathbb{N} \to m \leq n \to \underline{m} \leq \underline{n})) \to \varphi_{54}$
\item $\varphi_{135} :\equiv \forall n\,(n \in B_{5} \leftrightarrow n \in \omega \land (n \in \mathbb{N} \to m \leq n \to \underline{m} \leq \underline{n})) \to n \in B_{5} \to \varphi_{31}$
\item $\varphi_{136} :\equiv \forall n\,(n \in B_{5} \leftrightarrow n \in \omega \land (n \in \mathbb{N} \to m \leq n \to \underline{m} \leq \underline{n})) \to n \in B_{5} \to \varphi_{32}$
\item $\varphi_{137} :\equiv (\forall n \in B_{5}\,(S(n) \in B_{5}) \to \forall y \in B_{5}\,(S(y) \in B_{5})) \to m \in \mathbb{N} \to \varphi_{46}$
\item $\varphi_{138} :\equiv \forall n\,(n \in B_{5} \leftrightarrow n \in \omega \land (n \in \mathbb{N} \to m \leq n \to \underline{m} \leq \underline{n})) \to n \in B_{5} \to \varphi_{34}$
\item $\varphi_{139} :\equiv (\emptyset \in \mathbb{N} \leftrightarrow \emptyset \in \omega \land \emptyset \neq \emptyset) \to \emptyset = \emptyset \to m \in \mathbb{N} \to \emptyset \in \mathbb{N} \to m \leq \emptyset \to \underline{m} \leq \underline{\emptyset}$
\item $\varphi_{140} :\equiv (\forall x\,(\mathrm{Ind}(x) \to n \in x) \to \mathrm{Ind}(B_{5}) \to n \in B_{5}) \to m \in \mathbb{N} \to \varphi_{76}$
\item $\varphi_{141} :\equiv \forall n\,(n \in B_{5} \leftrightarrow n \in \omega \land (n \in \mathbb{N} \to m \leq n \to \underline{m} \leq \underline{n})) \to (\emptyset \in B_{5} \leftrightarrow \emptyset \in \omega \land (\emptyset \in \mathbb{N} \to m \leq \emptyset \to \underline{m} \leq \underline{\emptyset}))$
\item $\varphi_{142} :\equiv (m \in \mathbb{N} \to \forall n \in \mathbb{N}\,(m \leq n \to \underline{m} \leq \underline{n})) \to \varphi_{47} \to m \in \mathbb{N} \to n \in \mathbb{N} \to m \leq n \to \underline{m} \leq \underline{n}$
\item $\varphi_{143} :\equiv (n = \emptyset \to S(n) = S(n) \to S(n) = S(\emptyset)) \to m \in \mathbb{N} \to \varphi_{112}$
\item $\varphi_{144} :\equiv (m \in \mathbb{N} \to m \leq 1_{\omega} \to \underline{m} \leq \underline{m}) \to (m = 1_{\omega} \to \underline{m} < \underline{m} \to \underline{m} < \underline{1_{\omega}}) \to \varphi_{67}$
\item $\varphi_{145} :\equiv (\forall y \in \omega\,(S(y) \in \omega) \to n \in \omega \to S(n) \in \omega) \to m \in \mathbb{N} \to \varphi_{88}$
\item $\varphi_{146} :\equiv (m \in \mathbb{N} \to n \in \omega \to n \in \mathbb{N} \to m \leq n \to \underline{m} \leq \underline{n}) \to (m \in \mathbb{N} \to n \in \mathbb{N} \to n \in \omega) \to \varphi_{62}$
\item $\varphi_{147} :\equiv (\forall y \in B_{5}\,(S(y) \in B_{5}) \to \forall n \in B_{5}\,(S(n) \in B_{5})) \to \varphi_{137}$
\item $\varphi_{148} :\equiv (\mathrm{Ind}(\omega) \leftrightarrow \emptyset \in \omega \land \forall y \in \omega\,(S(y) \in \omega)) \to m \in \mathbb{N} \to \varphi_{109}$
\item $\varphi_{149} :\equiv (m \in \mathbb{N} \land \neg \forall n \in \omega\,(n \in \mathbb{N} \to m \leq n \to \underline{m} \leq \underline{n}) \to \varphi_{41}) \to m \in \mathbb{N} \to \forall n \in \omega\,(n \in \mathbb{N} \to m \leq n \to \underline{m} \leq \underline{n})$
\item $\varphi_{150} :\equiv \forall n\,(\varphi_{53} \to n \in B_{5} \to S(n) \in B_{5}) \to \forall n\,(\varphi_{53}) \to \forall n \in B_{5}\,(S(n) \in B_{5})$
\item $\varphi_{151} :\equiv (\underline{n} + 1 = \underline{S(n)} \to \underline{n} < \underline{n} + 1 \to \underline{n} < \underline{S(n)}) \to m \in \mathbb{N} \to \varphi_{129}$
\item $\varphi_{152} :\equiv (m \in \mathbb{N} \to \forall n \in \omega\,(n \in \mathbb{N} \to m \leq n \to \underline{m} \leq \underline{n})) \to \varphi_{78} \to m \in \mathbb{N} \to n \in \omega \to n \in \mathbb{N} \to m \leq n \to \underline{m} \leq \underline{n}$
\item $\varphi_{153} :\equiv (m \in \mathbb{N} \to \varphi_{70}) \to (n \in \omega \leftrightarrow \forall x\,(\mathrm{Ind}(x) \to n \in x)) \to m \in \mathbb{N} \to \varphi_{77}$
\item $\varphi_{154} :\equiv (m = S(n) \to \underline{m} < \underline{m} \to \underline{m} < \underline{S(n)}) \to (m = S(n) \to \underline{m} = \underline{m} \to \underline{m} = \underline{S(n)}) \to m \in \mathbb{N} \to \varphi_{138}$
\item $\varphi_{155} :\equiv \forall n\,(\varphi_{53} \to n \in \omega \to n \in \mathbb{N} \to m \leq n \to \underline{m} \leq \underline{n}) \to \forall n\,(\varphi_{53}) \to \forall n \in \omega\,(n \in \mathbb{N} \to m \leq n \to \underline{m} \leq \underline{n})$
\item $\varphi_{156} :\equiv (m \in \mathbb{N} \to \varphi_{100}) \to (S(n) = 1_{\omega} \to 1_{\omega} = S(n)) \to m \in \mathbb{N} \to \varphi_{101}$
\item $\varphi_{157} :\equiv (m \in \mathbb{N} \to \varphi_{101}) \to (1_{\omega} = S(n) \to S(n) = 1_{\omega}) \to m \in \mathbb{N} \to \varphi_{100}$
\item $\varphi_{158} :\equiv (m \in \mathbb{N} \to \varphi_{109}) \to (m \in \mathbb{N} \to \emptyset \in \mathbb{N} \to m \leq \emptyset \to \underline{m} \leq \underline{\emptyset}) \to m \in \mathbb{N} \to \varphi_{60}$
\item $\varphi_{159} :\equiv (n \in \mathbb{N} \to \underline{S(n)} = \underline{n} + 1) \to (m \in \mathbb{N} \to \varphi_{102}) \to m \in \mathbb{N} \to \varphi_{132}$
\item $\varphi_{160} :\equiv (S(n) = 1_{\omega} \to \underline{m} = \underline{S(n)} \to \underline{m} = \underline{1_{\omega}}) \to (m \in \mathbb{N} \to \varphi_{100}) \to m \in \mathbb{N} \to \varphi_{125}$
\item $\varphi_{161} :\equiv (\underline{n} \in \mathbb{R} \to \underline{n} < \underline{n} + 1) \to (m \in \mathbb{N} \to \varphi_{124}) \to m \in \mathbb{N} \to \varphi_{131}$
\item $\varphi_{162} :\equiv (m \in \mathbb{N} \to \varphi_{132}) \to (\underline{S(n)} = \underline{n} + 1 \to \underline{n} + 1 = \underline{S(n)}) \to m \in \mathbb{N} \to \varphi_{133}$
\item $\varphi_{163} :\equiv (S(n) = 1_{\omega} \to \underline{m} < \underline{S(n)} \to \underline{m} < \underline{1_{\omega}}) \to \varphi_{160}$
\item $\varphi_{164} :\equiv (n \in \mathbb{N} \leftrightarrow n \in \omega \land n \neq \emptyset) \to (m \in \mathbb{N} \to \varphi_{65}) \to (m \in \mathbb{N} \to \varphi_{108}) \to m \in \mathbb{N} \to \varphi_{102}$
\item $\varphi_{165} :\equiv (S(n) = 1_{\omega} \to m < S(n) \to m < 1_{\omega}) \to (S(n) = 1_{\omega} \to m = S(n) \to m = 1_{\omega}) \to \varphi_{163}$
\item $\varphi_{166} :\equiv (\mathrm{Ind}(B_{5}) \leftrightarrow \emptyset \in B_{5} \land \forall y \in B_{5}\,(S(y) \in B_{5})) \to (m \in \mathbb{N} \to \varphi_{60}) \to (m \in \mathbb{N} \to \varphi_{69}) \to \varphi_{147}$
\item $\varphi_{167} :\equiv (m \in \mathbb{N} \to \varphi_{76}) \to (m \in \mathbb{N} \to \varphi_{46}) \to (m \in \mathbb{N} \to \varphi_{114}) \to m \in \mathbb{N} \to \varphi_{97}$
\item $\varphi_{168} :\equiv (m \in \mathbb{N} \to \varphi_{134}) \to (m \in \mathbb{N} \to \varphi_{68}) \to (m \in \mathbb{N} \to \varphi_{110}) \to m \in \mathbb{N} \to \varphi_{66}$
\item $\varphi_{169} :\equiv (1_{\omega} = S(n) \to \underline{m} = \underline{1_{\omega}} \to \underline{m} = \underline{S(n)}) \to \varphi_{165}$
\item $\varphi_{170} :\equiv (\underline{m} \leq \underline{n} \to \underline{n} \leq \underline{S(n)} \to \underline{m} \leq \underline{S(n)}) \to (m \in \mathbb{N} \to \varphi_{130}) \to (m \in \mathbb{N} \to \varphi_{126}) \to m \in \mathbb{N} \to \varphi_{136}$
\item $\varphi_{171} :\equiv (m \in \mathbb{N} \to \varphi_{118}) \to (m \in \mathbb{N} \to \varphi_{117}) \to (m \in \mathbb{N} \to \varphi_{111}) \to m \in \mathbb{N} \to \varphi_{121}$
\item $\varphi_{172} :\equiv (1_{\omega} = S(n) \to \underline{m} < \underline{1_{\omega}} \to \underline{m} < \underline{S(n)}) \to \varphi_{169}$
\item $\varphi_{173} :\equiv \forall n\,(\varphi_{53} \to n \in B_{5} \to S(n) \in B_{5}) \to \varphi_{150} \to (\varphi_{53} \to \forall n\,(\varphi_{53})) \to m \in \mathbb{N} \to \varphi_{69}$
\item $\varphi_{174} :\equiv (m \in \mathbb{N} \to \varphi_{121}) \to (m \in \mathbb{N} \to \varphi_{136}) \to (m \in \mathbb{N} \to \varphi_{138}) \to m \in \mathbb{N} \to \varphi_{127}$
\item $\varphi_{175} :\equiv (1_{\omega} = S(n) \to m < 1_{\omega} \to m < S(n)) \to (1_{\omega} = S(n) \to m = 1_{\omega} \to m = S(n)) \to \varphi_{172}$
\item $\varphi_{176} :\equiv \forall n\,(\varphi_{53} \to n \in \omega \to n \in \mathbb{N} \to m \leq n \to \underline{m} \leq \underline{n}) \to \varphi_{155} \to (\varphi_{53} \to \forall n\,(\varphi_{53})) \to m \in \mathbb{N} \to \varphi_{107}$
\item $\varphi_{177} :\equiv (m \in \mathbb{N} \to \varphi_{101}) \to (m \in \mathbb{N} \to m \leq 1_{\omega} \to \underline{m} \leq \underline{1_{\omega}}) \to \varphi_{175}$
\end{abbrevlist}
\end{fullproof}

\chapter{テキスト第1章 序論 の証明木}
本章の定理数 18, ノード数の合計 2716.

\section{絶対値 (テキスト第1章 1.1節)}
% 自動生成 (prover/build.js): AbsOfNonneg  ノード数 8
\begin{fullproof}{[A1]}{$\forall x\,(0 \leq x \to \lvert x \rvert = x)$}{8}\phantomsection\label{tree:AbsOfNonneg}
\noindent 主木:
\begin{ptree}
\AXC{$\forall x\,(0 \leq x \land \lvert x \rvert = x \lor \neg 0 \leq x \land \lvert x \rvert = -x)$\lbl{Def$\lvert\cdot\rvert$}}
\AXC{$\varphi_{2}$\lbl{A2}}
\BIC{$0 \leq x \land \lvert x \rvert = x \lor \neg 0 \leq x \land \lvert x \rvert = -x$}
\AXC{$0 \leq x \to 0 \leq x$\lbl{A1}}
\AXC{$(0 \leq x \to 0 \leq x) \to \varphi_{1}$\lbl{A1}}
\BIC{$\varphi_{1}$}
\BIC{$0 \leq x \to \lvert x \rvert = x$}
\UIC{$\forall x\,(0 \leq x \to \lvert x \rvert = x)$}
\end{ptree}
\begin{abbrevlist}
\item $\varphi_{1} :\equiv 0 \leq x \land \lvert x \rvert = x \lor \neg 0 \leq x \land \lvert x \rvert = -x \to 0 \leq x \to \lvert x \rvert = x$
\item $\varphi_{2} :\equiv \forall x\,(0 \leq x \land \lvert x \rvert = x \lor \neg 0 \leq x \land \lvert x \rvert = -x) \to 0 \leq x \land \lvert x \rvert = x \lor \neg 0 \leq x \land \lvert x \rvert = -x$
\end{abbrevlist}
\end{fullproof}

% 自動生成 (prover/build.js): NotLeLt  ノード数 17
\begin{fullproof}{[A2]}{$\forall x\,\forall y\,(x \in \mathbb{R} \to y \in \mathbb{R} \to \neg y \leq x \to x < y)$}{17}\phantomsection\label{tree:NotLeLt}
\noindent$\Pi_{1}$:
\begin{ptree}
\AXC{$\forall x\,\forall y\,(x \in \mathbb{R} \to y \in \mathbb{R} \to \neg x < y \to y \leq x)$\lbl{[O8]}}
\AXC{$\varphi_{3}$\lbl{A2}}
\BIC{$\forall y\,(x \in \mathbb{R} \to y \in \mathbb{R} \to \neg x < y \to y \leq x)$}
\AXC{$\varphi_{1}$\lbl{A2}}
\BIC{$x \in \mathbb{R} \to y \in \mathbb{R} \to \neg x < y \to y \leq x$}
\AXC{$\varphi_{5}$\lbl{A1}}
\BIC{$(x \in \mathbb{R} \to y \in \mathbb{R} \to \neg y \leq x \to x \in \mathbb{R}) \to \varphi_{4}$}
\end{ptree}
\noindent$\Pi_{2}$:
\begin{ptree}
\AXC{$x \in \mathbb{R} \to y \in \mathbb{R} \to \neg y \leq x \to y \in \mathbb{R}$\lbl{A1}}
\AXC{$x \in \mathbb{R} \to y \in \mathbb{R} \to \neg y \leq x \to x \in \mathbb{R}$\lbl{A1}}
\AXC{$\stackrel{\Pi_{1}}{\vdots}$}\noLine
\UIC{$(x \in \mathbb{R} \to y \in \mathbb{R} \to \neg y \leq x \to x \in \mathbb{R}) \to \varphi_{4}$}
\BIC{$\varphi_{4}$}
\BIC{$x \in \mathbb{R} \to y \in \mathbb{R} \to \neg y \leq x \to \neg x < y \to y \leq x$}
\end{ptree}
\noindent$\Pi_{3}$:
\begin{ptree}
\AXC{$x \in \mathbb{R} \to y \in \mathbb{R} \to \neg y \leq x \to \neg y \leq x$\lbl{A1}}
\AXC{$(x \in \mathbb{R} \to y \in \mathbb{R} \to \neg y \leq x \to \neg y \leq x) \to \varphi_{2}$\lbl{A1}}
\BIC{$\varphi_{2}$}
\end{ptree}
\noindent 主木:
\begin{ptree}
\AXC{$\stackrel{\Pi_{2}}{\vdots}$}\noLine
\UIC{$x \in \mathbb{R} \to y \in \mathbb{R} \to \neg y \leq x \to \neg x < y \to y \leq x$}
\AXC{$\stackrel{\Pi_{3}}{\vdots}$}\noLine
\UIC{$\varphi_{2}$}
\BIC{$x \in \mathbb{R} \to y \in \mathbb{R} \to \neg y \leq x \to x < y$}
\UIC{$\forall y\,(x \in \mathbb{R} \to y \in \mathbb{R} \to \neg y \leq x \to x < y)$}
\UIC{$\forall x\,\forall y\,(x \in \mathbb{R} \to y \in \mathbb{R} \to \neg y \leq x \to x < y)$}
\end{ptree}
\begin{abbrevlist}
\item $\varphi_{1} :\equiv \forall y\,(x \in \mathbb{R} \to y \in \mathbb{R} \to \neg x < y \to y \leq x) \to x \in \mathbb{R} \to y \in \mathbb{R} \to \neg x < y \to y \leq x$
\item $\varphi_{2} :\equiv (x \in \mathbb{R} \to y \in \mathbb{R} \to \neg y \leq x \to \neg x < y \to y \leq x) \to x \in \mathbb{R} \to y \in \mathbb{R} \to \neg y \leq x \to x < y$
\item $\varphi_{3} :\equiv \forall x\,\forall y\,(x \in \mathbb{R} \to y \in \mathbb{R} \to \neg x < y \to y \leq x) \to \forall y\,(x \in \mathbb{R} \to y \in \mathbb{R} \to \neg x < y \to y \leq x)$
\item $\varphi_{4} :\equiv (x \in \mathbb{R} \to y \in \mathbb{R} \to \neg y \leq x \to y \in \mathbb{R}) \to x \in \mathbb{R} \to y \in \mathbb{R} \to \neg y \leq x \to \neg x < y \to y \leq x$
\item $\varphi_{5} :\equiv (x \in \mathbb{R} \to y \in \mathbb{R} \to \neg x < y \to y \leq x) \to (x \in \mathbb{R} \to y \in \mathbb{R} \to \neg y \leq x \to x \in \mathbb{R}) \to \varphi_{4}$
\end{abbrevlist}
\end{fullproof}

% 自動生成 (prover/build.js): AbsOfNeg  ノード数 14
\begin{fullproof}{[A3]}{$\forall x\,(x < 0 \to \lvert x \rvert = -x)$}{14}\phantomsection\label{tree:AbsOfNeg}
\noindent$\Pi_{1}$:
\begin{ptree}
\AXC{$\forall x\,\forall y\,(x \leq y \to \neg y < x)$\lbl{[O7]}}
\AXC{$\forall x\,\forall y\,(x \leq y \to \neg y < x) \to \forall y\,(0 \leq y \to \neg y < 0)$\lbl{A2}}
\BIC{$\forall y\,(0 \leq y \to \neg y < 0)$}
\end{ptree}
\noindent$\Pi_{2}$:
\begin{ptree}
\AXC{$\forall x\,(0 \leq x \land \lvert x \rvert = x \lor \neg 0 \leq x \land \lvert x \rvert = -x)$\lbl{Def$\lvert\cdot\rvert$}}
\AXC{$\varphi_{2}$\lbl{A2}}
\BIC{$0 \leq x \land \lvert x \rvert = x \lor \neg 0 \leq x \land \lvert x \rvert = -x$}
\AXC{$x < 0 \to x < 0$\lbl{A1}}
\AXC{$(x < 0 \to x < 0) \to \varphi_{1}$\lbl{A1}}
\BIC{$\varphi_{1}$}
\BIC{$(0 \leq x \to \neg x < 0) \to x < 0 \to \lvert x \rvert = -x$}
\end{ptree}
\noindent 主木:
\begin{ptree}
\AXC{$\stackrel{\Pi_{1}}{\vdots}$}\noLine
\UIC{$\forall y\,(0 \leq y \to \neg y < 0)$}
\AXC{$\forall y\,(0 \leq y \to \neg y < 0) \to 0 \leq x \to \neg x < 0$\lbl{A2}}
\BIC{$0 \leq x \to \neg x < 0$}
\AXC{$\stackrel{\Pi_{2}}{\vdots}$}\noLine
\UIC{$(0 \leq x \to \neg x < 0) \to x < 0 \to \lvert x \rvert = -x$}
\BIC{$x < 0 \to \lvert x \rvert = -x$}
\UIC{$\forall x\,(x < 0 \to \lvert x \rvert = -x)$}
\end{ptree}
\begin{abbrevlist}
\item $\varphi_{1} :\equiv 0 \leq x \land \lvert x \rvert = x \lor \neg 0 \leq x \land \lvert x \rvert = -x \to (0 \leq x \to \neg x < 0) \to x < 0 \to \lvert x \rvert = -x$
\item $\varphi_{2} :\equiv \forall x\,(0 \leq x \land \lvert x \rvert = x \lor \neg 0 \leq x \land \lvert x \rvert = -x) \to 0 \leq x \land \lvert x \rvert = x \lor \neg 0 \leq x \land \lvert x \rvert = -x$
\end{abbrevlist}
\end{fullproof}

% 自動生成 (prover/build.js): AbsNonpos  ノード数 98
\begin{fullproof}{[A4]}{$\forall x \in \mathbb{R}\,(x \leq 0 \to \lvert x \rvert = -x)$}{98}\phantomsection\label{tree:AbsNonpos}
\noindent$\Pi_{1}$:
\begin{ptree}
\AXC{$y = z \to x = y \to x = z$\lbl{A6}}
\AXC{$(y = z \to x = y \to x = z) \to x = y \to y = z \to x = z$\lbl{A1}}
\BIC{$x = y \to y = z \to x = z$}
\UIC{$\forall z\,(x = y \to y = z \to x = z)$}
\UIC{$\forall y\,\forall z\,(x = y \to y = z \to x = z)$}
\UIC{$\forall x\,\forall y\,\forall z\,(x = y \to y = z \to x = z)$}
\AXC{$\varphi_{4}$\lbl{A2}}
\BIC{$\forall y\,\forall z\,(\lvert x \rvert = y \to y = z \to \lvert x \rvert = z)$}
\AXC{$\varphi_{7}$\lbl{A2}}
\BIC{$\forall z\,(\lvert x \rvert = x \to x = z \to \lvert x \rvert = z)$}
\AXC{$\varphi_{9}$\lbl{A2}}
\BIC{$\lvert x \rvert = x \to x = -x \to \lvert x \rvert = -x$}
\end{ptree}
\noindent$\Pi_{2}$:
\begin{ptree}
\AXC{$y = z \to x = y \to x = z$\lbl{A6}}
\AXC{$(y = z \to x = y \to x = z) \to x = y \to y = z \to x = z$\lbl{A1}}
\BIC{$x = y \to y = z \to x = z$}
\UIC{$\forall z\,(x = y \to y = z \to x = z)$}
\UIC{$\forall y\,\forall z\,(x = y \to y = z \to x = z)$}
\UIC{$\forall x\,\forall y\,\forall z\,(x = y \to y = z \to x = z)$}
\AXC{$\varphi_{2}$\lbl{A2}}
\BIC{$\forall y\,\forall z\,(x = y \to y = z \to x = z)$}
\AXC{$\varphi_{3}$\lbl{A2}}
\BIC{$\forall z\,(x = -0 \to -0 = z \to x = z)$}
\AXC{$\varphi_{10}$\lbl{A2}}
\BIC{$x = -0 \to -0 = -x \to x = -x$}
\end{ptree}
\noindent$\Pi_{3}$:
\begin{ptree}
\AXC{$x = x$\lbl{A5}}
\AXC{$x = y \to x = x \to y = x$\lbl{A6}}
\AXC{$(x = y \to x = x \to y = x) \to x = x \to x = y \to y = x$\lbl{A1}}
\BIC{$x = x \to x = y \to y = x$}
\BIC{$x = y \to y = x$}
\UIC{$\forall y\,(x = y \to y = x)$}
\UIC{$\forall x\,\forall y\,(x = y \to y = x)$}
\end{ptree}
\noindent$\Pi_{4}$:
\begin{ptree}
\AXC{$\stackrel{\Pi_{3}}{\vdots}$}\noLine
\UIC{$\forall x\,\forall y\,(x = y \to y = x)$}
\AXC{$\forall x\,\forall y\,(x = y \to y = x) \to \forall y\,(x = y \to y = x)$\lbl{A2}}
\BIC{$\forall y\,(x = y \to y = x)$}
\AXC{$\forall y\,(x = y \to y = x) \to x = 0 \to 0 = x$\lbl{A2}}
\BIC{$x = 0 \to 0 = x$}
\end{ptree}
\noindent$\Pi_{5}$:
\begin{ptree}
\AXC{$\stackrel{\Pi_{4}}{\vdots}$}\noLine
\UIC{$x = 0 \to 0 = x$}
\AXC{$x \in \mathbb{R} \to x \leq 0 \to x = 0 \to x = 0$\lbl{A1}}
\AXC{$\varphi_{11}$\lbl{A1}}
\BIC{$(x = 0 \to 0 = x) \to x \in \mathbb{R} \to x \leq 0 \to x = 0 \to 0 = x$}
\BIC{$x \in \mathbb{R} \to x \leq 0 \to x = 0 \to 0 = x$}
\AXC{$(x \in \mathbb{R} \to x \leq 0 \to x = 0 \to 0 = x) \to -0 = -0 \to \varphi_{15}$\lbl{A1}}
\BIC{$-0 = -0 \to \varphi_{15}$}
\end{ptree}
\noindent$\Pi_{6}$:
\begin{ptree}
\AXC{$x = x$\lbl{A5}}
\UIC{$\forall x\,(x = x)$}
\AXC{$\forall x\,(x = x) \to -0 = -0$\lbl{A2}}
\BIC{$-0 = -0$}
\AXC{$\stackrel{\Pi_{5}}{\vdots}$}\noLine
\UIC{$-0 = -0 \to \varphi_{15}$}
\BIC{$\varphi_{15}$}
\end{ptree}
\noindent$\Pi_{7}$:
\begin{ptree}
\AXC{$y = z \to x = y \to x = z$\lbl{A6}}
\AXC{$(y = z \to x = y \to x = z) \to x = y \to y = z \to x = z$\lbl{A1}}
\BIC{$x = y \to y = z \to x = z$}
\UIC{$\forall z\,(x = y \to y = z \to x = z)$}
\UIC{$\forall y\,\forall z\,(x = y \to y = z \to x = z)$}
\UIC{$\forall x\,\forall y\,\forall z\,(x = y \to y = z \to x = z)$}
\AXC{$\varphi_{2}$\lbl{A2}}
\BIC{$\forall y\,\forall z\,(x = y \to y = z \to x = z)$}
\AXC{$\varphi_{1}$\lbl{A2}}
\BIC{$\forall z\,(x = 0 \to 0 = z \to x = z)$}
\end{ptree}
\noindent$\Pi_{8}$:
\begin{ptree}
\AXC{$\stackrel{\Pi_{3}}{\vdots}$}\noLine
\UIC{$\forall x\,\forall y\,(x = y \to y = x)$}
\AXC{$\forall x\,\forall y\,(x = y \to y = x) \to \forall y\,(-0 = y \to y = -0)$\lbl{A2}}
\BIC{$\forall y\,(-0 = y \to y = -0)$}
\end{ptree}
\noindent$\Pi_{9}$:
\begin{ptree}
\AXC{$\stackrel{\Pi_{8}}{\vdots}$}\noLine
\UIC{$\forall y\,(-0 = y \to y = -0)$}
\AXC{$\forall y\,(-0 = y \to y = -0) \to -0 = 0 \to 0 = -0$\lbl{A2}}
\BIC{$-0 = 0 \to 0 = -0$}
\end{ptree}
\noindent$\Pi_{10}$:
\begin{ptree}
\AXC{$\stackrel{\Pi_{9}}{\vdots}$}\noLine
\UIC{$-0 = 0 \to 0 = -0$}
\AXC{$-0 = 0$\lbl{[F6]}}
\AXC{$-0 = 0 \to (-0 = 0 \to 0 = -0) \to 0 = -0$\lbl{A1}}
\BIC{$(-0 = 0 \to 0 = -0) \to 0 = -0$}
\BIC{$0 = -0$}
\end{ptree}
\noindent$\Pi_{11}$:
\begin{ptree}
\AXC{$\stackrel{\Pi_{10}}{\vdots}$}\noLine
\UIC{$0 = -0$}
\AXC{$x \in \mathbb{R} \to x \leq 0 \to x = 0 \to x = 0$\lbl{A1}}
\AXC{$(x \in \mathbb{R} \to x \leq 0 \to x = 0 \to x = 0) \to 0 = -0 \to \varphi_{6}$\lbl{A1}}
\BIC{$0 = -0 \to \varphi_{6}$}
\BIC{$\varphi_{6}$}
\end{ptree}
\noindent$\Pi_{12}$:
\begin{ptree}
\AXC{$\stackrel{\Pi_{7}}{\vdots}$}\noLine
\UIC{$\forall z\,(x = 0 \to 0 = z \to x = z)$}
\AXC{$\forall z\,(x = 0 \to 0 = z \to x = z) \to x = 0 \to 0 = -0 \to x = -0$\lbl{A2}}
\BIC{$x = 0 \to 0 = -0 \to x = -0$}
\AXC{$\stackrel{\Pi_{11}}{\vdots}$}\noLine
\UIC{$\varphi_{6}$}
\BIC{$x \in \mathbb{R} \to x \leq 0 \to x = 0 \to x = -0$}
\AXC{$\varphi_{18}$\lbl{A1}}
\BIC{$(x \in \mathbb{R} \to x \leq 0 \to x = 0 \to -0 = -x) \to \varphi_{12}$}
\end{ptree}
\noindent$\Pi_{13}$:
\begin{ptree}
\AXC{$z = w \to -0 = -z \to -0 = -w$\lbl{A6}}
\UIC{$\forall w\,(z = w \to -0 = -z \to -0 = -w)$}
\UIC{$\forall z\,\forall w\,(z = w \to -0 = -z \to -0 = -w)$}
\AXC{$\varphi_{17}$\lbl{A2}}
\BIC{$\forall w\,(0 = w \to -0 = -0 \to -0 = -w)$}
\AXC{$\varphi_{14}$\lbl{A2}}
\BIC{$0 = x \to -0 = -0 \to -0 = -x$}
\AXC{$\stackrel{\Pi_{6}}{\vdots}$}\noLine
\UIC{$\varphi_{15}$}
\BIC{$x \in \mathbb{R} \to x \leq 0 \to x = 0 \to -0 = -x$}
\AXC{$\stackrel{\Pi_{12}}{\vdots}$}\noLine
\UIC{$(x \in \mathbb{R} \to x \leq 0 \to x = 0 \to -0 = -x) \to \varphi_{12}$}
\BIC{$\varphi_{12}$}
\end{ptree}
\noindent$\Pi_{14}$:
\begin{ptree}
\AXC{$\stackrel{\Pi_{3}}{\vdots}$}\noLine
\UIC{$\forall x\,\forall y\,(x = y \to y = x)$}
\AXC{$\forall x\,\forall y\,(x = y \to y = x) \to \forall y\,(x = y \to y = x)$\lbl{A2}}
\BIC{$\forall y\,(x = y \to y = x)$}
\AXC{$\forall y\,(x = y \to y = x) \to x = 0 \to 0 = x$\lbl{A2}}
\BIC{$x = 0 \to 0 = x$}
\end{ptree}
\noindent$\Pi_{15}$:
\begin{ptree}
\AXC{$\forall x\,(0 \leq x \to \lvert x \rvert = x)$\lbl{[A1]}}
\AXC{$\forall x\,(0 \leq x \to \lvert x \rvert = x) \to 0 \leq x \to \lvert x \rvert = x$\lbl{A2}}
\BIC{$0 \leq x \to \lvert x \rvert = x$}
\AXC{$(0 \leq x \to \lvert x \rvert = x) \to \varphi_{8}$\lbl{A1}}
\BIC{$\varphi_{8}$}
\end{ptree}
\noindent$\Pi_{16}$:
\begin{ptree}
\AXC{$\stackrel{\Pi_{14}}{\vdots}$}\noLine
\UIC{$x = 0 \to 0 = x$}
\AXC{$x \in \mathbb{R} \to x \leq 0 \to x = 0 \to x = 0$\lbl{A1}}
\AXC{$\varphi_{11}$\lbl{A1}}
\BIC{$(x = 0 \to 0 = x) \to x \in \mathbb{R} \to x \leq 0 \to x = 0 \to 0 = x$}
\BIC{$x \in \mathbb{R} \to x \leq 0 \to x = 0 \to 0 = x$}
\AXC{$\varphi_{5}$\lbl{A1}}
\BIC{$x \in \mathbb{R} \to x \leq 0 \to x = 0 \to 0 \leq x$}
\AXC{$\stackrel{\Pi_{15}}{\vdots}$}\noLine
\UIC{$\varphi_{8}$}
\BIC{$x \in \mathbb{R} \to x \leq 0 \to x = 0 \to \lvert x \rvert = x$}
\AXC{$\varphi_{19}$\lbl{A1}}
\BIC{$(x \in \mathbb{R} \to x \leq 0 \to x = 0 \to x = -x) \to \varphi_{16}$}
\end{ptree}
\noindent$\Pi_{17}$:
\begin{ptree}
\AXC{$\stackrel{\Pi_{1}}{\vdots}$}\noLine
\UIC{$\lvert x \rvert = x \to x = -x \to \lvert x \rvert = -x$}
\AXC{$\stackrel{\Pi_{2}}{\vdots}$}\noLine
\UIC{$x = -0 \to -0 = -x \to x = -x$}
\AXC{$\stackrel{\Pi_{13}}{\vdots}$}\noLine
\UIC{$\varphi_{12}$}
\BIC{$x \in \mathbb{R} \to x \leq 0 \to x = 0 \to x = -x$}
\AXC{$\stackrel{\Pi_{16}}{\vdots}$}\noLine
\UIC{$(x \in \mathbb{R} \to x \leq 0 \to x = 0 \to x = -x) \to \varphi_{16}$}
\BIC{$\varphi_{16}$}
\BIC{$x \in \mathbb{R} \to x \leq 0 \to x = 0 \to \lvert x \rvert = -x$}
\end{ptree}
\noindent$\Pi_{18}$:
\begin{ptree}
\AXC{$x \in \mathbb{R} \to x \leq 0 \to x \leq 0$\lbl{A1}}
\AXC{$(x \in \mathbb{R} \to x \leq 0 \to x \leq 0) \to (x < 0 \to \lvert x \rvert = -x) \to \varphi_{13}$\lbl{A1}}
\BIC{$(x < 0 \to \lvert x \rvert = -x) \to \varphi_{13}$}
\end{ptree}
\noindent$\Pi_{19}$:
\begin{ptree}
\AXC{$\forall x\,(x < 0 \to \lvert x \rvert = -x)$\lbl{[A3]}}
\AXC{$\forall x\,(x < 0 \to \lvert x \rvert = -x) \to x < 0 \to \lvert x \rvert = -x$\lbl{A2}}
\BIC{$x < 0 \to \lvert x \rvert = -x$}
\AXC{$\stackrel{\Pi_{18}}{\vdots}$}\noLine
\UIC{$(x < 0 \to \lvert x \rvert = -x) \to \varphi_{13}$}
\BIC{$\varphi_{13}$}
\end{ptree}
\noindent 主木:
\begin{ptree}
\AXC{$\stackrel{\Pi_{17}}{\vdots}$}\noLine
\UIC{$x \in \mathbb{R} \to x \leq 0 \to x = 0 \to \lvert x \rvert = -x$}
\AXC{$\stackrel{\Pi_{19}}{\vdots}$}\noLine
\UIC{$\varphi_{13}$}
\BIC{$x \in \mathbb{R} \to x \leq 0 \to \lvert x \rvert = -x$}
\UIC{$\forall x \in \mathbb{R}\,(x \leq 0 \to \lvert x \rvert = -x)$}
\end{ptree}
\begin{abbrevlist}
\item $\varphi_{1} :\equiv \forall y\,\forall z\,(x = y \to y = z \to x = z) \to \forall z\,(x = 0 \to 0 = z \to x = z)$
\item $\varphi_{2} :\equiv \forall x\,\forall y\,\forall z\,(x = y \to y = z \to x = z) \to \forall y\,\forall z\,(x = y \to y = z \to x = z)$
\item $\varphi_{3} :\equiv \forall y\,\forall z\,(x = y \to y = z \to x = z) \to \forall z\,(x = -0 \to -0 = z \to x = z)$
\item $\varphi_{4} :\equiv \forall x\,\forall y\,\forall z\,(x = y \to y = z \to x = z) \to \forall y\,\forall z\,(\lvert x \rvert = y \to y = z \to \lvert x \rvert = z)$
\item $\varphi_{5} :\equiv (x \in \mathbb{R} \to x \leq 0 \to x = 0 \to 0 = x) \to x \in \mathbb{R} \to x \leq 0 \to x = 0 \to 0 \leq x$
\item $\varphi_{6} :\equiv (x = 0 \to 0 = -0 \to x = -0) \to x \in \mathbb{R} \to x \leq 0 \to x = 0 \to x = -0$
\item $\varphi_{7} :\equiv \forall y\,\forall z\,(\lvert x \rvert = y \to y = z \to \lvert x \rvert = z) \to \forall z\,(\lvert x \rvert = x \to x = z \to \lvert x \rvert = z)$
\item $\varphi_{8} :\equiv (x \in \mathbb{R} \to x \leq 0 \to x = 0 \to 0 \leq x) \to x \in \mathbb{R} \to x \leq 0 \to x = 0 \to \lvert x \rvert = x$
\item $\varphi_{9} :\equiv \forall z\,(\lvert x \rvert = x \to x = z \to \lvert x \rvert = z) \to \lvert x \rvert = x \to x = -x \to \lvert x \rvert = -x$
\item $\varphi_{10} :\equiv \forall z\,(x = -0 \to -0 = z \to x = z) \to x = -0 \to -0 = -x \to x = -x$
\item $\varphi_{11} :\equiv (x \in \mathbb{R} \to x \leq 0 \to x = 0 \to x = 0) \to (x = 0 \to 0 = x) \to x \in \mathbb{R} \to x \leq 0 \to x = 0 \to 0 = x$
\item $\varphi_{12} :\equiv (x = -0 \to -0 = -x \to x = -x) \to x \in \mathbb{R} \to x \leq 0 \to x = 0 \to x = -x$
\item $\varphi_{13} :\equiv (x \in \mathbb{R} \to x \leq 0 \to x = 0 \to \lvert x \rvert = -x) \to x \in \mathbb{R} \to x \leq 0 \to \lvert x \rvert = -x$
\item $\varphi_{14} :\equiv \forall w\,(0 = w \to -0 = -0 \to -0 = -w) \to 0 = x \to -0 = -0 \to -0 = -x$
\item $\varphi_{15} :\equiv (0 = x \to -0 = -0 \to -0 = -x) \to x \in \mathbb{R} \to x \leq 0 \to x = 0 \to -0 = -x$
\item $\varphi_{16} :\equiv (\lvert x \rvert = x \to x = -x \to \lvert x \rvert = -x) \to x \in \mathbb{R} \to x \leq 0 \to x = 0 \to \lvert x \rvert = -x$
\item $\varphi_{17} :\equiv \forall z\,\forall w\,(z = w \to -0 = -z \to -0 = -w) \to \forall w\,(0 = w \to -0 = -0 \to -0 = -w)$
\item $\varphi_{18} :\equiv (x \in \mathbb{R} \to x \leq 0 \to x = 0 \to x = -0) \to (x \in \mathbb{R} \to x \leq 0 \to x = 0 \to -0 = -x) \to \varphi_{12}$
\item $\varphi_{19} :\equiv (x \in \mathbb{R} \to x \leq 0 \to x = 0 \to \lvert x \rvert = x) \to (x \in \mathbb{R} \to x \leq 0 \to x = 0 \to x = -x) \to \varphi_{16}$
\end{abbrevlist}
\end{fullproof}

% 自動生成 (prover/build.js): AbsR  ノード数 65
\begin{fullproof}{[A5]}{$\forall x \in \mathbb{R}\,(\lvert x \rvert \in \mathbb{R})$}{65}\phantomsection\label{tree:AbsR}
\noindent$\Pi_{1}$:
\begin{ptree}
\AXC{$\forall x \in \mathbb{R}\,(-x \in \mathbb{R})$\lbl{(R1)}}
\AXC{$\forall x \in \mathbb{R}\,(-x \in \mathbb{R}) \to x \in \mathbb{R} \to -x \in \mathbb{R}$\lbl{A2}}
\BIC{$x \in \mathbb{R} \to -x \in \mathbb{R}$}
\AXC{$(x \in \mathbb{R} \to -x \in \mathbb{R}) \to \varphi_{11}$\lbl{A1}}
\BIC{$\varphi_{11}$}
\end{ptree}
\noindent$\Pi_{2}$:
\begin{ptree}
\AXC{$x = x$\lbl{A5}}
\AXC{$x = y \to x = x \to y = x$\lbl{A6}}
\AXC{$(x = y \to x = x \to y = x) \to x = x \to x = y \to y = x$\lbl{A1}}
\BIC{$x = x \to x = y \to y = x$}
\BIC{$x = y \to y = x$}
\UIC{$\forall y\,(x = y \to y = x)$}
\UIC{$\forall x\,\forall y\,(x = y \to y = x)$}
\end{ptree}
\noindent$\Pi_{3}$:
\begin{ptree}
\AXC{$\stackrel{\Pi_{2}}{\vdots}$}\noLine
\UIC{$\forall x\,\forall y\,(x = y \to y = x)$}
\AXC{$\forall x\,\forall y\,(x = y \to y = x) \to \forall y\,(\lvert x \rvert = y \to y = \lvert x \rvert)$\lbl{A2}}
\BIC{$\forall y\,(\lvert x \rvert = y \to y = \lvert x \rvert)$}
\end{ptree}
\noindent$\Pi_{4}$:
\begin{ptree}
\AXC{$x \in \mathbb{R} \to \neg 0 \leq x \land \lvert x \rvert = -x \to \neg 0 \leq x \land \lvert x \rvert = -x$\lbl{A1}}
\AXC{$\varphi_{13}$\lbl{A1}}
\BIC{$x \in \mathbb{R} \to \neg 0 \leq x \land \lvert x \rvert = -x \to \lvert x \rvert = -x$}
\AXC{$(x \in \mathbb{R} \to \neg 0 \leq x \land \lvert x \rvert = -x \to \lvert x \rvert = -x) \to \varphi_{6}$\lbl{A1}}
\BIC{$\varphi_{6}$}
\end{ptree}
\noindent$\Pi_{5}$:
\begin{ptree}
\AXC{$\stackrel{\Pi_{3}}{\vdots}$}\noLine
\UIC{$\forall y\,(\lvert x \rvert = y \to y = \lvert x \rvert)$}
\AXC{$\forall y\,(\lvert x \rvert = y \to y = \lvert x \rvert) \to \lvert x \rvert = -x \to -x = \lvert x \rvert$\lbl{A2}}
\BIC{$\lvert x \rvert = -x \to -x = \lvert x \rvert$}
\AXC{$\stackrel{\Pi_{4}}{\vdots}$}\noLine
\UIC{$\varphi_{6}$}
\BIC{$x \in \mathbb{R} \to \neg 0 \leq x \land \lvert x \rvert = -x \to -x = \lvert x \rvert$}
\AXC{$\varphi_{16}$\lbl{A1}}
\BIC{$(x \in \mathbb{R} \to \neg 0 \leq x \land \lvert x \rvert = -x \to -x \in \mathbb{R}) \to \varphi_{12}$}
\end{ptree}
\noindent$\Pi_{6}$:
\begin{ptree}
\AXC{$x \in \mathbb{R} \to \neg 0 \leq x \land \lvert x \rvert = -x \to x \in \mathbb{R}$\lbl{A1}}
\AXC{$\stackrel{\Pi_{1}}{\vdots}$}\noLine
\UIC{$\varphi_{11}$}
\BIC{$x \in \mathbb{R} \to \neg 0 \leq x \land \lvert x \rvert = -x \to -x \in \mathbb{R}$}
\AXC{$\stackrel{\Pi_{5}}{\vdots}$}\noLine
\UIC{$(x \in \mathbb{R} \to \neg 0 \leq x \land \lvert x \rvert = -x \to -x \in \mathbb{R}) \to \varphi_{12}$}
\BIC{$\varphi_{12}$}
\end{ptree}
\noindent$\Pi_{7}$:
\begin{ptree}
\AXC{$\stackrel{\Pi_{2}}{\vdots}$}\noLine
\UIC{$\forall x\,\forall y\,(x = y \to y = x)$}
\AXC{$\forall x\,\forall y\,(x = y \to y = x) \to \forall y\,(\lvert x \rvert = y \to y = \lvert x \rvert)$\lbl{A2}}
\BIC{$\forall y\,(\lvert x \rvert = y \to y = \lvert x \rvert)$}
\end{ptree}
\noindent$\Pi_{8}$:
\begin{ptree}
\AXC{$x \in \mathbb{R} \to 0 \leq x \land \lvert x \rvert = x \to 0 \leq x \land \lvert x \rvert = x$\lbl{A1}}
\AXC{$\varphi_{8}$\lbl{A1}}
\BIC{$x \in \mathbb{R} \to 0 \leq x \land \lvert x \rvert = x \to \lvert x \rvert = x$}
\AXC{$(x \in \mathbb{R} \to 0 \leq x \land \lvert x \rvert = x \to \lvert x \rvert = x) \to \varphi_{1}$\lbl{A1}}
\BIC{$\varphi_{1}$}
\end{ptree}
\noindent$\Pi_{9}$:
\begin{ptree}
\AXC{$\stackrel{\Pi_{7}}{\vdots}$}\noLine
\UIC{$\forall y\,(\lvert x \rvert = y \to y = \lvert x \rvert)$}
\AXC{$\forall y\,(\lvert x \rvert = y \to y = \lvert x \rvert) \to \lvert x \rvert = x \to x = \lvert x \rvert$\lbl{A2}}
\BIC{$\lvert x \rvert = x \to x = \lvert x \rvert$}
\AXC{$\stackrel{\Pi_{8}}{\vdots}$}\noLine
\UIC{$\varphi_{1}$}
\BIC{$x \in \mathbb{R} \to 0 \leq x \land \lvert x \rvert = x \to x = \lvert x \rvert$}
\AXC{$\varphi_{15}$\lbl{A1}}
\BIC{$(x \in \mathbb{R} \to 0 \leq x \land \lvert x \rvert = x \to x \in \mathbb{R}) \to \varphi_{9}$}
\end{ptree}
\noindent$\Pi_{10}$:
\begin{ptree}
\AXC{$x \in \mathbb{R} \to 0 \leq x \land \lvert x \rvert = x \to x \in \mathbb{R}$\lbl{A1}}
\AXC{$\stackrel{\Pi_{9}}{\vdots}$}\noLine
\UIC{$(x \in \mathbb{R} \to 0 \leq x \land \lvert x \rvert = x \to x \in \mathbb{R}) \to \varphi_{9}$}
\BIC{$\varphi_{9}$}
\end{ptree}
\noindent$\Pi_{11}$:
\begin{ptree}
\AXC{$\forall x\,(0 \leq x \land \lvert x \rvert = x \lor \neg 0 \leq x \land \lvert x \rvert = -x)$\lbl{Def$\lvert\cdot\rvert$}}
\AXC{$\varphi_{5}$\lbl{A2}}
\BIC{$0 \leq x \land \lvert x \rvert = x \lor \neg 0 \leq x \land \lvert x \rvert = -x$}
\AXC{$\varphi_{14}$\lbl{A1}}
\BIC{$(x \in \mathbb{R} \to 0 \leq x \land \lvert x \rvert = x \to \lvert x \rvert \in \mathbb{R}) \to \varphi_{7}$}
\end{ptree}
\noindent$\Pi_{12}$:
\begin{ptree}
\AXC{$z = w \to z \in \mathbb{R} \to w \in \mathbb{R}$\lbl{A6}}
\UIC{$\forall w\,(z = w \to z \in \mathbb{R} \to w \in \mathbb{R})$}
\UIC{$\forall z\,\forall w\,(z = w \to z \in \mathbb{R} \to w \in \mathbb{R})$}
\AXC{$\varphi_{2}$\lbl{A2}}
\BIC{$\forall w\,(x = w \to x \in \mathbb{R} \to w \in \mathbb{R})$}
\AXC{$\varphi_{3}$\lbl{A2}}
\BIC{$x = \lvert x \rvert \to x \in \mathbb{R} \to \lvert x \rvert \in \mathbb{R}$}
\AXC{$\stackrel{\Pi_{10}}{\vdots}$}\noLine
\UIC{$\varphi_{9}$}
\BIC{$x \in \mathbb{R} \to 0 \leq x \land \lvert x \rvert = x \to \lvert x \rvert \in \mathbb{R}$}
\AXC{$\stackrel{\Pi_{11}}{\vdots}$}\noLine
\UIC{$(x \in \mathbb{R} \to 0 \leq x \land \lvert x \rvert = x \to \lvert x \rvert \in \mathbb{R}) \to \varphi_{7}$}
\BIC{$\varphi_{7}$}
\end{ptree}
\noindent 主木:
\begin{ptree}
\AXC{$z = w \to z \in \mathbb{R} \to w \in \mathbb{R}$\lbl{A6}}
\UIC{$\forall w\,(z = w \to z \in \mathbb{R} \to w \in \mathbb{R})$}
\UIC{$\forall z\,\forall w\,(z = w \to z \in \mathbb{R} \to w \in \mathbb{R})$}
\AXC{$\varphi_{4}$\lbl{A2}}
\BIC{$\forall w\,(-x = w \to -x \in \mathbb{R} \to w \in \mathbb{R})$}
\AXC{$\varphi_{10}$\lbl{A2}}
\BIC{$-x = \lvert x \rvert \to -x \in \mathbb{R} \to \lvert x \rvert \in \mathbb{R}$}
\AXC{$\stackrel{\Pi_{6}}{\vdots}$}\noLine
\UIC{$\varphi_{12}$}
\BIC{$x \in \mathbb{R} \to \neg 0 \leq x \land \lvert x \rvert = -x \to \lvert x \rvert \in \mathbb{R}$}
\AXC{$\stackrel{\Pi_{12}}{\vdots}$}\noLine
\UIC{$\varphi_{7}$}
\BIC{$x \in \mathbb{R} \to \lvert x \rvert \in \mathbb{R}$}
\UIC{$\forall x \in \mathbb{R}\,(\lvert x \rvert \in \mathbb{R})$}
\end{ptree}
\begin{abbrevlist}
\item $\varphi_{1} :\equiv (\lvert x \rvert = x \to x = \lvert x \rvert) \to x \in \mathbb{R} \to 0 \leq x \land \lvert x \rvert = x \to x = \lvert x \rvert$
\item $\varphi_{2} :\equiv \forall z\,\forall w\,(z = w \to z \in \mathbb{R} \to w \in \mathbb{R}) \to \forall w\,(x = w \to x \in \mathbb{R} \to w \in \mathbb{R})$
\item $\varphi_{3} :\equiv \forall w\,(x = w \to x \in \mathbb{R} \to w \in \mathbb{R}) \to x = \lvert x \rvert \to x \in \mathbb{R} \to \lvert x \rvert \in \mathbb{R}$
\item $\varphi_{4} :\equiv \forall z\,\forall w\,(z = w \to z \in \mathbb{R} \to w \in \mathbb{R}) \to \forall w\,(-x = w \to -x \in \mathbb{R} \to w \in \mathbb{R})$
\item $\varphi_{5} :\equiv \forall x\,(0 \leq x \land \lvert x \rvert = x \lor \neg 0 \leq x \land \lvert x \rvert = -x) \to 0 \leq x \land \lvert x \rvert = x \lor \neg 0 \leq x \land \lvert x \rvert = -x$
\item $\varphi_{6} :\equiv (\lvert x \rvert = -x \to -x = \lvert x \rvert) \to x \in \mathbb{R} \to \neg 0 \leq x \land \lvert x \rvert = -x \to -x = \lvert x \rvert$
\item $\varphi_{7} :\equiv (x \in \mathbb{R} \to \neg 0 \leq x \land \lvert x \rvert = -x \to \lvert x \rvert \in \mathbb{R}) \to x \in \mathbb{R} \to \lvert x \rvert \in \mathbb{R}$
\item $\varphi_{8} :\equiv (x \in \mathbb{R} \to 0 \leq x \land \lvert x \rvert = x \to 0 \leq x \land \lvert x \rvert = x) \to x \in \mathbb{R} \to 0 \leq x \land \lvert x \rvert = x \to \lvert x \rvert = x$
\item $\varphi_{9} :\equiv (x = \lvert x \rvert \to x \in \mathbb{R} \to \lvert x \rvert \in \mathbb{R}) \to x \in \mathbb{R} \to 0 \leq x \land \lvert x \rvert = x \to \lvert x \rvert \in \mathbb{R}$
\item $\varphi_{10} :\equiv \forall w\,(-x = w \to -x \in \mathbb{R} \to w \in \mathbb{R}) \to -x = \lvert x \rvert \to -x \in \mathbb{R} \to \lvert x \rvert \in \mathbb{R}$
\item $\varphi_{11} :\equiv (x \in \mathbb{R} \to \neg 0 \leq x \land \lvert x \rvert = -x \to x \in \mathbb{R}) \to x \in \mathbb{R} \to \neg 0 \leq x \land \lvert x \rvert = -x \to -x \in \mathbb{R}$
\item $\varphi_{12} :\equiv (-x = \lvert x \rvert \to -x \in \mathbb{R} \to \lvert x \rvert \in \mathbb{R}) \to x \in \mathbb{R} \to \neg 0 \leq x \land \lvert x \rvert = -x \to \lvert x \rvert \in \mathbb{R}$
\item $\varphi_{13} :\equiv (x \in \mathbb{R} \to \neg 0 \leq x \land \lvert x \rvert = -x \to \neg 0 \leq x \land \lvert x \rvert = -x) \to x \in \mathbb{R} \to \neg 0 \leq x \land \lvert x \rvert = -x \to \lvert x \rvert = -x$
\item $\varphi_{14} :\equiv 0 \leq x \land \lvert x \rvert = x \lor \neg 0 \leq x \land \lvert x \rvert = -x \to (x \in \mathbb{R} \to 0 \leq x \land \lvert x \rvert = x \to \lvert x \rvert \in \mathbb{R}) \to \varphi_{7}$
\item $\varphi_{15} :\equiv (x \in \mathbb{R} \to 0 \leq x \land \lvert x \rvert = x \to x = \lvert x \rvert) \to (x \in \mathbb{R} \to 0 \leq x \land \lvert x \rvert = x \to x \in \mathbb{R}) \to \varphi_{9}$
\item $\varphi_{16} :\equiv (x \in \mathbb{R} \to \neg 0 \leq x \land \lvert x \rvert = -x \to -x = \lvert x \rvert) \to (x \in \mathbb{R} \to \neg 0 \leq x \land \lvert x \rvert = -x \to -x \in \mathbb{R}) \to \varphi_{12}$
\end{abbrevlist}
\end{fullproof}

% 自動生成 (prover/build.js): AbsNonneg  ノード数 106
\begin{fullproof}{[A6]}{$\forall x \in \mathbb{R}\,(0 \leq \lvert x \rvert)$}{106}\phantomsection\label{tree:AbsNonneg}
\noindent$\Pi_{1}$:
\begin{ptree}
\AXC{$x \in \mathbb{R} \to \neg 0 \leq x \to x \in \mathbb{R}$\lbl{A1}}
\AXC{$\forall x\,\forall y\,(x \in \mathbb{R} \to y \in \mathbb{R} \to \neg y \leq x \to x < y)$\lbl{[A2]}}
\AXC{$\varphi_{21}$\lbl{A2}}
\BIC{$\forall y\,(x \in \mathbb{R} \to y \in \mathbb{R} \to \neg y \leq x \to x < y)$}
\AXC{$\varphi_{19}$\lbl{A2}}
\BIC{$x \in \mathbb{R} \to 0 \in \mathbb{R} \to \neg 0 \leq x \to x < 0$}
\AXC{$\varphi_{25}$\lbl{A1}}
\BIC{$(x \in \mathbb{R} \to \neg 0 \leq x \to x \in \mathbb{R}) \to \varphi_{11}$}
\BIC{$\varphi_{11}$}
\end{ptree}
\noindent$\Pi_{2}$:
\begin{ptree}
\AXC{$x \in \mathbb{R} \to \neg 0 \leq x \to \neg 0 \leq x$\lbl{A1}}
\AXC{$0 \in \mathbb{R}$\lbl{(R1)}}
\AXC{$\stackrel{\Pi_{1}}{\vdots}$}\noLine
\UIC{$\varphi_{11}$}
\BIC{$(x \in \mathbb{R} \to \neg 0 \leq x \to \neg 0 \leq x) \to x \in \mathbb{R} \to \neg 0 \leq x \to x < 0$}
\BIC{$x \in \mathbb{R} \to \neg 0 \leq x \to x < 0$}
\end{ptree}
\noindent$\Pi_{3}$:
\begin{ptree}
\AXC{$x \in \mathbb{R} \to \neg 0 \leq x \to x \in \mathbb{R}$\lbl{A1}}
\AXC{$\forall x\,\forall y\,(x \in \mathbb{R} \to y \in \mathbb{R} \to x < y \to -y < -x)$\lbl{[F22]}}
\AXC{$\varphi_{23}$\lbl{A2}}
\BIC{$\forall y\,(x \in \mathbb{R} \to y \in \mathbb{R} \to x < y \to -y < -x)$}
\AXC{$\varphi_{22}$\lbl{A2}}
\BIC{$x \in \mathbb{R} \to 0 \in \mathbb{R} \to x < 0 \to -0 < -x$}
\AXC{$\varphi_{27}$\lbl{A1}}
\BIC{$(x \in \mathbb{R} \to \neg 0 \leq x \to x \in \mathbb{R}) \to 0 \in \mathbb{R} \to \varphi_{9}$}
\BIC{$0 \in \mathbb{R} \to \varphi_{9}$}
\end{ptree}
\noindent$\Pi_{4}$:
\begin{ptree}
\AXC{$\stackrel{\Pi_{2}}{\vdots}$}\noLine
\UIC{$x \in \mathbb{R} \to \neg 0 \leq x \to x < 0$}
\AXC{$0 \in \mathbb{R}$\lbl{(R1)}}
\AXC{$\stackrel{\Pi_{3}}{\vdots}$}\noLine
\UIC{$0 \in \mathbb{R} \to \varphi_{9}$}
\BIC{$\varphi_{9}$}
\BIC{$x \in \mathbb{R} \to \neg 0 \leq x \to -0 < -x$}
\end{ptree}
\noindent$\Pi_{5}$:
\begin{ptree}
\AXC{$\stackrel{\Pi_{4}}{\vdots}$}\noLine
\UIC{$x \in \mathbb{R} \to \neg 0 \leq x \to -0 < -x$}
\AXC{$-0 = 0$\lbl{[F6]}}
\AXC{$-0 = 0 \to (x \in \mathbb{R} \to \neg 0 \leq x \to -0 < -x) \to \varphi_{13}$\lbl{A1}}
\BIC{$(x \in \mathbb{R} \to \neg 0 \leq x \to -0 < -x) \to \varphi_{13}$}
\BIC{$\varphi_{13}$}
\end{ptree}
\noindent$\Pi_{6}$:
\begin{ptree}
\AXC{$x = x$\lbl{A5}}
\AXC{$x = y \to x = x \to y = x$\lbl{A6}}
\AXC{$(x = y \to x = x \to y = x) \to x = x \to x = y \to y = x$\lbl{A1}}
\BIC{$x = x \to x = y \to y = x$}
\BIC{$x = y \to y = x$}
\UIC{$\forall y\,(x = y \to y = x)$}
\UIC{$\forall x\,\forall y\,(x = y \to y = x)$}
\end{ptree}
\noindent$\Pi_{7}$:
\begin{ptree}
\AXC{$\stackrel{\Pi_{6}}{\vdots}$}\noLine
\UIC{$\forall x\,\forall y\,(x = y \to y = x)$}
\AXC{$\forall x\,\forall y\,(x = y \to y = x) \to \forall y\,(\lvert x \rvert = y \to y = \lvert x \rvert)$\lbl{A2}}
\BIC{$\forall y\,(\lvert x \rvert = y \to y = \lvert x \rvert)$}
\end{ptree}
\noindent$\Pi_{8}$:
\begin{ptree}
\AXC{$\stackrel{\Pi_{7}}{\vdots}$}\noLine
\UIC{$\forall y\,(\lvert x \rvert = y \to y = \lvert x \rvert)$}
\AXC{$\forall y\,(\lvert x \rvert = y \to y = \lvert x \rvert) \to \lvert x \rvert = -x \to -x = \lvert x \rvert$\lbl{A2}}
\BIC{$\lvert x \rvert = -x \to -x = \lvert x \rvert$}
\end{ptree}
\noindent$\Pi_{9}$:
\begin{ptree}
\AXC{$\forall x\,(x < 0 \to \lvert x \rvert = -x)$\lbl{[A3]}}
\AXC{$\forall x\,(x < 0 \to \lvert x \rvert = -x) \to x < 0 \to \lvert x \rvert = -x$\lbl{A2}}
\BIC{$x < 0 \to \lvert x \rvert = -x$}
\AXC{$(x < 0 \to \lvert x \rvert = -x) \to \varphi_{6}$\lbl{A1}}
\BIC{$\varphi_{6}$}
\end{ptree}
\noindent$\Pi_{10}$:
\begin{ptree}
\AXC{$\stackrel{\Pi_{8}}{\vdots}$}\noLine
\UIC{$\lvert x \rvert = -x \to -x = \lvert x \rvert$}
\AXC{$\stackrel{\Pi_{2}}{\vdots}$}\noLine
\UIC{$x \in \mathbb{R} \to \neg 0 \leq x \to x < 0$}
\AXC{$\stackrel{\Pi_{9}}{\vdots}$}\noLine
\UIC{$\varphi_{6}$}
\BIC{$x \in \mathbb{R} \to \neg 0 \leq x \to \lvert x \rvert = -x$}
\AXC{$(x \in \mathbb{R} \to \neg 0 \leq x \to \lvert x \rvert = -x) \to \varphi_{10}$\lbl{A1}}
\BIC{$\varphi_{10}$}
\BIC{$x \in \mathbb{R} \to \neg 0 \leq x \to -x = \lvert x \rvert$}
\end{ptree}
\noindent$\Pi_{11}$:
\begin{ptree}
\AXC{$z = w \to z < -x \to w < -x$\lbl{A6}}
\UIC{$\forall w\,(z = w \to z < -x \to w < -x)$}
\UIC{$\forall z\,\forall w\,(z = w \to z < -x \to w < -x)$}
\AXC{$\varphi_{18}$\lbl{A2}}
\BIC{$\forall w\,(-0 = w \to -0 < -x \to w < -x)$}
\AXC{$\varphi_{20}$\lbl{A2}}
\BIC{$-0 = 0 \to -0 < -x \to 0 < -x$}
\AXC{$\stackrel{\Pi_{5}}{\vdots}$}\noLine
\UIC{$\varphi_{13}$}
\BIC{$x \in \mathbb{R} \to \neg 0 \leq x \to 0 < -x$}
\AXC{$\stackrel{\Pi_{10}}{\vdots}$}\noLine
\UIC{$x \in \mathbb{R} \to \neg 0 \leq x \to -x = \lvert x \rvert$}
\AXC{$\varphi_{26}$\lbl{A1}}
\BIC{$(x \in \mathbb{R} \to \neg 0 \leq x \to 0 < -x) \to \varphi_{14}$}
\BIC{$\varphi_{14}$}
\end{ptree}
\noindent$\Pi_{12}$:
\begin{ptree}
\AXC{$z = w \to 0 < z \to 0 < w$\lbl{A6}}
\UIC{$\forall w\,(z = w \to 0 < z \to 0 < w)$}
\UIC{$\forall z\,\forall w\,(z = w \to 0 < z \to 0 < w)$}
\AXC{$\varphi_{8}$\lbl{A2}}
\BIC{$\forall w\,(-x = w \to 0 < -x \to 0 < w)$}
\AXC{$\varphi_{15}$\lbl{A2}}
\BIC{$-x = \lvert x \rvert \to 0 < -x \to 0 < \lvert x \rvert$}
\AXC{$\stackrel{\Pi_{11}}{\vdots}$}\noLine
\UIC{$\varphi_{14}$}
\BIC{$x \in \mathbb{R} \to \neg 0 \leq x \to 0 < \lvert x \rvert$}
\AXC{$\varphi_{7}$\lbl{A1}}
\BIC{$x \in \mathbb{R} \to \neg 0 \leq x \to 0 \leq \lvert x \rvert$}
\end{ptree}
\noindent$\Pi_{13}$:
\begin{ptree}
\AXC{$\stackrel{\Pi_{6}}{\vdots}$}\noLine
\UIC{$\forall x\,\forall y\,(x = y \to y = x)$}
\AXC{$\forall x\,\forall y\,(x = y \to y = x) \to \forall y\,(\lvert x \rvert = y \to y = \lvert x \rvert)$\lbl{A2}}
\BIC{$\forall y\,(\lvert x \rvert = y \to y = \lvert x \rvert)$}
\end{ptree}
\noindent$\Pi_{14}$:
\begin{ptree}
\AXC{$\stackrel{\Pi_{13}}{\vdots}$}\noLine
\UIC{$\forall y\,(\lvert x \rvert = y \to y = \lvert x \rvert)$}
\AXC{$\forall y\,(\lvert x \rvert = y \to y = \lvert x \rvert) \to \lvert x \rvert = x \to x = \lvert x \rvert$\lbl{A2}}
\BIC{$\lvert x \rvert = x \to x = \lvert x \rvert$}
\end{ptree}
\noindent$\Pi_{15}$:
\begin{ptree}
\AXC{$x \in \mathbb{R} \to 0 \leq x \to 0 \leq x$\lbl{A1}}
\AXC{$\forall x\,(0 \leq x \to \lvert x \rvert = x)$\lbl{[A1]}}
\AXC{$\forall x\,(0 \leq x \to \lvert x \rvert = x) \to 0 \leq x \to \lvert x \rvert = x$\lbl{A2}}
\BIC{$0 \leq x \to \lvert x \rvert = x$}
\AXC{$\varphi_{12}$\lbl{A1}}
\BIC{$(x \in \mathbb{R} \to 0 \leq x \to 0 \leq x) \to x \in \mathbb{R} \to 0 \leq x \to \lvert x \rvert = x$}
\BIC{$x \in \mathbb{R} \to 0 \leq x \to \lvert x \rvert = x$}
\end{ptree}
\noindent$\Pi_{16}$:
\begin{ptree}
\AXC{$\stackrel{\Pi_{14}}{\vdots}$}\noLine
\UIC{$\lvert x \rvert = x \to x = \lvert x \rvert$}
\AXC{$\stackrel{\Pi_{15}}{\vdots}$}\noLine
\UIC{$x \in \mathbb{R} \to 0 \leq x \to \lvert x \rvert = x$}
\AXC{$\varphi_{16}$\lbl{A1}}
\BIC{$(\lvert x \rvert = x \to x = \lvert x \rvert) \to x \in \mathbb{R} \to 0 \leq x \to x = \lvert x \rvert$}
\BIC{$x \in \mathbb{R} \to 0 \leq x \to x = \lvert x \rvert$}
\AXC{$(x \in \mathbb{R} \to 0 \leq x \to x = \lvert x \rvert) \to \varphi_{24}$\lbl{A1}}
\BIC{$\varphi_{24}$}
\end{ptree}
\noindent$\Pi_{17}$:
\begin{ptree}
\AXC{$z = w \to 0 < z \to 0 < w$\lbl{A6}}
\UIC{$\forall w\,(z = w \to 0 < z \to 0 < w)$}
\UIC{$\forall z\,\forall w\,(z = w \to 0 < z \to 0 < w)$}
\AXC{$\varphi_{3}$\lbl{A2}}
\BIC{$\forall w\,(x = w \to 0 < x \to 0 < w)$}
\AXC{$\varphi_{5}$\lbl{A2}}
\BIC{$x = \lvert x \rvert \to 0 < x \to 0 < \lvert x \rvert$}
\AXC{$x \in \mathbb{R} \to 0 \leq x \to 0 \leq x$\lbl{A1}}
\AXC{$\stackrel{\Pi_{16}}{\vdots}$}\noLine
\UIC{$\varphi_{24}$}
\BIC{$(x = \lvert x \rvert \to 0 < x \to 0 < \lvert x \rvert) \to \varphi_{4}$}
\BIC{$\varphi_{4}$}
\end{ptree}
\noindent 主木:
\begin{ptree}
\AXC{$\stackrel{\Pi_{12}}{\vdots}$}\noLine
\UIC{$x \in \mathbb{R} \to \neg 0 \leq x \to 0 \leq \lvert x \rvert$}
\AXC{$z = w \to 0 = z \to 0 = w$\lbl{A6}}
\UIC{$\forall w\,(z = w \to 0 = z \to 0 = w)$}
\UIC{$\forall z\,\forall w\,(z = w \to 0 = z \to 0 = w)$}
\AXC{$\varphi_{1}$\lbl{A2}}
\BIC{$\forall w\,(x = w \to 0 = x \to 0 = w)$}
\AXC{$\varphi_{2}$\lbl{A2}}
\BIC{$x = \lvert x \rvert \to 0 = x \to 0 = \lvert x \rvert$}
\AXC{$\stackrel{\Pi_{17}}{\vdots}$}\noLine
\UIC{$\varphi_{4}$}
\BIC{$x \in \mathbb{R} \to 0 \leq x \to 0 \leq \lvert x \rvert$}
\AXC{$\varphi_{17}$\lbl{A1}}
\BIC{$(x \in \mathbb{R} \to \neg 0 \leq x \to 0 \leq \lvert x \rvert) \to x \in \mathbb{R} \to 0 \leq \lvert x \rvert$}
\BIC{$x \in \mathbb{R} \to 0 \leq \lvert x \rvert$}
\UIC{$\forall x \in \mathbb{R}\,(0 \leq \lvert x \rvert)$}
\end{ptree}
\begin{abbrevlist}
\item $\varphi_{1} :\equiv \forall z\,\forall w\,(z = w \to 0 = z \to 0 = w) \to \forall w\,(x = w \to 0 = x \to 0 = w)$
\item $\varphi_{2} :\equiv \forall w\,(x = w \to 0 = x \to 0 = w) \to x = \lvert x \rvert \to 0 = x \to 0 = \lvert x \rvert$
\item $\varphi_{3} :\equiv \forall z\,\forall w\,(z = w \to 0 < z \to 0 < w) \to \forall w\,(x = w \to 0 < x \to 0 < w)$
\item $\varphi_{4} :\equiv (x = \lvert x \rvert \to 0 = x \to 0 = \lvert x \rvert) \to x \in \mathbb{R} \to 0 \leq x \to 0 \leq \lvert x \rvert$
\item $\varphi_{5} :\equiv \forall w\,(x = w \to 0 < x \to 0 < w) \to x = \lvert x \rvert \to 0 < x \to 0 < \lvert x \rvert$
\item $\varphi_{6} :\equiv (x \in \mathbb{R} \to \neg 0 \leq x \to x < 0) \to x \in \mathbb{R} \to \neg 0 \leq x \to \lvert x \rvert = -x$
\item $\varphi_{7} :\equiv (x \in \mathbb{R} \to \neg 0 \leq x \to 0 < \lvert x \rvert) \to x \in \mathbb{R} \to \neg 0 \leq x \to 0 \leq \lvert x \rvert$
\item $\varphi_{8} :\equiv \forall z\,\forall w\,(z = w \to 0 < z \to 0 < w) \to \forall w\,(-x = w \to 0 < -x \to 0 < w)$
\item $\varphi_{9} :\equiv (x \in \mathbb{R} \to \neg 0 \leq x \to x < 0) \to x \in \mathbb{R} \to \neg 0 \leq x \to -0 < -x$
\item $\varphi_{10} :\equiv (\lvert x \rvert = -x \to -x = \lvert x \rvert) \to x \in \mathbb{R} \to \neg 0 \leq x \to -x = \lvert x \rvert$
\item $\varphi_{11} :\equiv 0 \in \mathbb{R} \to (x \in \mathbb{R} \to \neg 0 \leq x \to \neg 0 \leq x) \to x \in \mathbb{R} \to \neg 0 \leq x \to x < 0$
\item $\varphi_{12} :\equiv (0 \leq x \to \lvert x \rvert = x) \to (x \in \mathbb{R} \to 0 \leq x \to 0 \leq x) \to x \in \mathbb{R} \to 0 \leq x \to \lvert x \rvert = x$
\item $\varphi_{13} :\equiv (-0 = 0 \to -0 < -x \to 0 < -x) \to x \in \mathbb{R} \to \neg 0 \leq x \to 0 < -x$
\item $\varphi_{14} :\equiv (-x = \lvert x \rvert \to 0 < -x \to 0 < \lvert x \rvert) \to x \in \mathbb{R} \to \neg 0 \leq x \to 0 < \lvert x \rvert$
\item $\varphi_{15} :\equiv \forall w\,(-x = w \to 0 < -x \to 0 < w) \to -x = \lvert x \rvert \to 0 < -x \to 0 < \lvert x \rvert$
\item $\varphi_{16} :\equiv (x \in \mathbb{R} \to 0 \leq x \to \lvert x \rvert = x) \to (\lvert x \rvert = x \to x = \lvert x \rvert) \to x \in \mathbb{R} \to 0 \leq x \to x = \lvert x \rvert$
\item $\varphi_{17} :\equiv (x \in \mathbb{R} \to 0 \leq x \to 0 \leq \lvert x \rvert) \to (x \in \mathbb{R} \to \neg 0 \leq x \to 0 \leq \lvert x \rvert) \to x \in \mathbb{R} \to 0 \leq \lvert x \rvert$
\item $\varphi_{18} :\equiv \forall z\,\forall w\,(z = w \to z < -x \to w < -x) \to \forall w\,(-0 = w \to -0 < -x \to w < -x)$
\item $\varphi_{19} :\equiv \forall y\,(x \in \mathbb{R} \to y \in \mathbb{R} \to \neg y \leq x \to x < y) \to x \in \mathbb{R} \to 0 \in \mathbb{R} \to \neg 0 \leq x \to x < 0$
\item $\varphi_{20} :\equiv \forall w\,(-0 = w \to -0 < -x \to w < -x) \to -0 = 0 \to -0 < -x \to 0 < -x$
\item $\varphi_{21} :\equiv \forall x\,\forall y\,(x \in \mathbb{R} \to y \in \mathbb{R} \to \neg y \leq x \to x < y) \to \forall y\,(x \in \mathbb{R} \to y \in \mathbb{R} \to \neg y \leq x \to x < y)$
\item $\varphi_{22} :\equiv \forall y\,(x \in \mathbb{R} \to y \in \mathbb{R} \to x < y \to -y < -x) \to x \in \mathbb{R} \to 0 \in \mathbb{R} \to x < 0 \to -0 < -x$
\item $\varphi_{23} :\equiv \forall x\,\forall y\,(x \in \mathbb{R} \to y \in \mathbb{R} \to x < y \to -y < -x) \to \forall y\,(x \in \mathbb{R} \to y \in \mathbb{R} \to x < y \to -y < -x)$
\item $\varphi_{24} :\equiv (x \in \mathbb{R} \to 0 \leq x \to 0 \leq x) \to (x = \lvert x \rvert \to 0 < x \to 0 < \lvert x \rvert) \to \varphi_{4}$
\item $\varphi_{25} :\equiv (x \in \mathbb{R} \to 0 \in \mathbb{R} \to \neg 0 \leq x \to x < 0) \to (x \in \mathbb{R} \to \neg 0 \leq x \to x \in \mathbb{R}) \to \varphi_{11}$
\item $\varphi_{26} :\equiv (x \in \mathbb{R} \to \neg 0 \leq x \to -x = \lvert x \rvert) \to (x \in \mathbb{R} \to \neg 0 \leq x \to 0 < -x) \to \varphi_{14}$
\item $\varphi_{27} :\equiv (x \in \mathbb{R} \to 0 \in \mathbb{R} \to x < 0 \to -0 < -x) \to (x \in \mathbb{R} \to \neg 0 \leq x \to x \in \mathbb{R}) \to 0 \in \mathbb{R} \to \varphi_{9}$
\end{abbrevlist}
\end{fullproof}

% 自動生成 (prover/build.js): AbsNeg  ノード数 216
\begin{fullproof}{[A7]}{$\forall x \in \mathbb{R}\,(\lvert -x \rvert = \lvert x \rvert)$}{216}\phantomsection\label{tree:AbsNeg}
\noindent$\Pi_{1}$:
\begin{ptree}
\AXC{$y = z \to x = y \to x = z$\lbl{A6}}
\AXC{$(y = z \to x = y \to x = z) \to x = y \to y = z \to x = z$\lbl{A1}}
\BIC{$x = y \to y = z \to x = z$}
\UIC{$\forall z\,(x = y \to y = z \to x = z)$}
\UIC{$\forall y\,\forall z\,(x = y \to y = z \to x = z)$}
\UIC{$\forall x\,\forall y\,\forall z\,(x = y \to y = z \to x = z)$}
\AXC{$\varphi_{7}$\lbl{A2}}
\BIC{$\forall y\,\forall z\,(\lvert -x \rvert = y \to y = z \to \lvert -x \rvert = z)$}
\AXC{$\varphi_{21}$\lbl{A2}}
\BIC{$\forall z\,(\lvert -x \rvert = x \to x = z \to \lvert -x \rvert = z)$}
\AXC{$\varphi_{27}$\lbl{A2}}
\BIC{$\lvert -x \rvert = x \to x = \lvert x \rvert \to \lvert -x \rvert = \lvert x \rvert$}
\end{ptree}
\noindent$\Pi_{2}$:
\begin{ptree}
\AXC{$x = x$\lbl{A5}}
\AXC{$x = y \to x = x \to y = x$\lbl{A6}}
\AXC{$(x = y \to x = x \to y = x) \to x = x \to x = y \to y = x$\lbl{A1}}
\BIC{$x = x \to x = y \to y = x$}
\BIC{$x = y \to y = x$}
\UIC{$\forall y\,(x = y \to y = x)$}
\UIC{$\forall x\,\forall y\,(x = y \to y = x)$}
\end{ptree}
\noindent$\Pi_{3}$:
\begin{ptree}
\AXC{$\stackrel{\Pi_{2}}{\vdots}$}\noLine
\UIC{$\forall x\,\forall y\,(x = y \to y = x)$}
\AXC{$\forall x\,\forall y\,(x = y \to y = x) \to \forall y\,(\lvert x \rvert = y \to y = \lvert x \rvert)$\lbl{A2}}
\BIC{$\forall y\,(\lvert x \rvert = y \to y = \lvert x \rvert)$}
\end{ptree}
\noindent$\Pi_{4}$:
\begin{ptree}
\AXC{$\stackrel{\Pi_{3}}{\vdots}$}\noLine
\UIC{$\forall y\,(\lvert x \rvert = y \to y = \lvert x \rvert)$}
\AXC{$\forall y\,(\lvert x \rvert = y \to y = \lvert x \rvert) \to \lvert x \rvert = x \to x = \lvert x \rvert$\lbl{A2}}
\BIC{$\lvert x \rvert = x \to x = \lvert x \rvert$}
\end{ptree}
\noindent$\Pi_{5}$:
\begin{ptree}
\AXC{$x \in \mathbb{R} \to 0 < x \to 0 < x$\lbl{A1}}
\AXC{$(x \in \mathbb{R} \to 0 < x \to 0 < x) \to x \in \mathbb{R} \to 0 < x \to 0 \leq x$\lbl{A1}}
\BIC{$x \in \mathbb{R} \to 0 < x \to 0 \leq x$}
\end{ptree}
\noindent$\Pi_{6}$:
\begin{ptree}
\AXC{$\forall x\,(0 \leq x \to \lvert x \rvert = x)$\lbl{[A1]}}
\AXC{$\forall x\,(0 \leq x \to \lvert x \rvert = x) \to 0 \leq x \to \lvert x \rvert = x$\lbl{A2}}
\BIC{$0 \leq x \to \lvert x \rvert = x$}
\AXC{$\varphi_{14}$\lbl{A1}}
\BIC{$(x \in \mathbb{R} \to 0 < x \to 0 \leq x) \to x \in \mathbb{R} \to 0 < x \to \lvert x \rvert = x$}
\end{ptree}
\noindent$\Pi_{7}$:
\begin{ptree}
\AXC{$\stackrel{\Pi_{5}}{\vdots}$}\noLine
\UIC{$x \in \mathbb{R} \to 0 < x \to 0 \leq x$}
\AXC{$\stackrel{\Pi_{6}}{\vdots}$}\noLine
\UIC{$(x \in \mathbb{R} \to 0 < x \to 0 \leq x) \to x \in \mathbb{R} \to 0 < x \to \lvert x \rvert = x$}
\BIC{$x \in \mathbb{R} \to 0 < x \to \lvert x \rvert = x$}
\AXC{$\varphi_{26}$\lbl{A1}}
\BIC{$(\lvert x \rvert = x \to x = \lvert x \rvert) \to x \in \mathbb{R} \to 0 < x \to x = \lvert x \rvert$}
\end{ptree}
\noindent$\Pi_{8}$:
\begin{ptree}
\AXC{$\stackrel{\Pi_{4}}{\vdots}$}\noLine
\UIC{$\lvert x \rvert = x \to x = \lvert x \rvert$}
\AXC{$\stackrel{\Pi_{7}}{\vdots}$}\noLine
\UIC{$(\lvert x \rvert = x \to x = \lvert x \rvert) \to x \in \mathbb{R} \to 0 < x \to x = \lvert x \rvert$}
\BIC{$x \in \mathbb{R} \to 0 < x \to x = \lvert x \rvert$}
\end{ptree}
\noindent$\Pi_{9}$:
\begin{ptree}
\AXC{$y = z \to x = y \to x = z$\lbl{A6}}
\AXC{$(y = z \to x = y \to x = z) \to x = y \to y = z \to x = z$\lbl{A1}}
\BIC{$x = y \to y = z \to x = z$}
\UIC{$\forall z\,(x = y \to y = z \to x = z)$}
\UIC{$\forall y\,\forall z\,(x = y \to y = z \to x = z)$}
\UIC{$\forall x\,\forall y\,\forall z\,(x = y \to y = z \to x = z)$}
\AXC{$\varphi_{7}$\lbl{A2}}
\BIC{$\forall y\,\forall z\,(\lvert -x \rvert = y \to y = z \to \lvert -x \rvert = z)$}
\AXC{$\varphi_{37}$\lbl{A2}}
\BIC{$\forall z\,(\lvert -x \rvert = -(-x) \to -(-x) = z \to \lvert -x \rvert = z)$}
\AXC{$\varphi_{46}$\lbl{A2}}
\BIC{$\lvert -x \rvert = -(-x) \to -(-x) = x \to \lvert -x \rvert = x$}
\end{ptree}
\noindent$\Pi_{10}$:
\begin{ptree}
\AXC{$\forall x \in \mathbb{R}\,(-(-x) = x)$\lbl{[F5]}}
\AXC{$\forall x \in \mathbb{R}\,(-(-x) = x) \to x \in \mathbb{R} \to -(-x) = x$\lbl{A2}}
\BIC{$x \in \mathbb{R} \to -(-x) = x$}
\AXC{$(x \in \mathbb{R} \to -(-x) = x) \to \varphi_{8}$\lbl{A1}}
\BIC{$\varphi_{8}$}
\end{ptree}
\noindent$\Pi_{11}$:
\begin{ptree}
\AXC{$\forall x\,\forall y\,(x \in \mathbb{R} \to y \in \mathbb{R} \to x < y \to -y < -x)$\lbl{[F22]}}
\AXC{$\varphi_{44}$\lbl{A2}}
\BIC{$\forall y\,(0 \in \mathbb{R} \to y \in \mathbb{R} \to 0 < y \to -y < -0)$}
\AXC{$\varphi_{41}$\lbl{A2}}
\BIC{$0 \in \mathbb{R} \to x \in \mathbb{R} \to 0 < x \to -x < -0$}
\end{ptree}
\noindent$\Pi_{12}$:
\begin{ptree}
\AXC{$0 \in \mathbb{R}$\lbl{(R1)}}
\AXC{$\stackrel{\Pi_{11}}{\vdots}$}\noLine
\UIC{$0 \in \mathbb{R} \to x \in \mathbb{R} \to 0 < x \to -x < -0$}
\AXC{$(0 \in \mathbb{R} \to x \in \mathbb{R} \to 0 < x \to -x < -0) \to 0 \in \mathbb{R} \to \varphi_{38}$\lbl{A1}}
\BIC{$0 \in \mathbb{R} \to \varphi_{38}$}
\BIC{$\varphi_{38}$}
\end{ptree}
\noindent$\Pi_{13}$:
\begin{ptree}
\AXC{$x \in \mathbb{R} \to 0 < x \to 0 < x$\lbl{A1}}
\AXC{$x \in \mathbb{R} \to 0 < x \to x \in \mathbb{R}$\lbl{A1}}
\AXC{$\stackrel{\Pi_{12}}{\vdots}$}\noLine
\UIC{$\varphi_{38}$}
\BIC{$(x \in \mathbb{R} \to 0 < x \to 0 < x) \to x \in \mathbb{R} \to 0 < x \to -x < -0$}
\BIC{$x \in \mathbb{R} \to 0 < x \to -x < -0$}
\end{ptree}
\noindent$\Pi_{14}$:
\begin{ptree}
\AXC{$\stackrel{\Pi_{13}}{\vdots}$}\noLine
\UIC{$x \in \mathbb{R} \to 0 < x \to -x < -0$}
\AXC{$-0 = 0$\lbl{[F6]}}
\AXC{$-0 = 0 \to (x \in \mathbb{R} \to 0 < x \to -x < -0) \to \varphi_{17}$\lbl{A1}}
\BIC{$(x \in \mathbb{R} \to 0 < x \to -x < -0) \to \varphi_{17}$}
\BIC{$\varphi_{17}$}
\end{ptree}
\noindent$\Pi_{15}$:
\begin{ptree}
\AXC{$z = w \to -x < z \to -x < w$\lbl{A6}}
\UIC{$\forall w\,(z = w \to -x < z \to -x < w)$}
\UIC{$\forall z\,\forall w\,(z = w \to -x < z \to -x < w)$}
\AXC{$\varphi_{24}$\lbl{A2}}
\BIC{$\forall w\,(-0 = w \to -x < -0 \to -x < w)$}
\AXC{$\varphi_{28}$\lbl{A2}}
\BIC{$-0 = 0 \to -x < -0 \to -x < 0$}
\AXC{$\stackrel{\Pi_{14}}{\vdots}$}\noLine
\UIC{$\varphi_{17}$}
\BIC{$x \in \mathbb{R} \to 0 < x \to -x < 0$}
\end{ptree}
\noindent$\Pi_{16}$:
\begin{ptree}
\AXC{$\forall x\,(x < 0 \to \lvert x \rvert = -x)$\lbl{[A3]}}
\AXC{$\forall x\,(x < 0 \to \lvert x \rvert = -x) \to -x < 0 \to \lvert -x \rvert = -(-x)$\lbl{A2}}
\BIC{$-x < 0 \to \lvert -x \rvert = -(-x)$}
\end{ptree}
\noindent$\Pi_{17}$:
\begin{ptree}
\AXC{$\stackrel{\Pi_{15}}{\vdots}$}\noLine
\UIC{$x \in \mathbb{R} \to 0 < x \to -x < 0$}
\AXC{$\stackrel{\Pi_{16}}{\vdots}$}\noLine
\UIC{$-x < 0 \to \lvert -x \rvert = -(-x)$}
\AXC{$(-x < 0 \to \lvert -x \rvert = -(-x)) \to \varphi_{20}$\lbl{A1}}
\BIC{$\varphi_{20}$}
\BIC{$x \in \mathbb{R} \to 0 < x \to \lvert -x \rvert = -(-x)$}
\AXC{$\varphi_{51}$\lbl{A1}}
\BIC{$(x \in \mathbb{R} \to 0 < x \to -(-x) = x) \to \varphi_{33}$}
\end{ptree}
\noindent$\Pi_{18}$:
\begin{ptree}
\AXC{$x \in \mathbb{R} \to 0 < x \to x \in \mathbb{R}$\lbl{A1}}
\AXC{$\stackrel{\Pi_{10}}{\vdots}$}\noLine
\UIC{$\varphi_{8}$}
\BIC{$x \in \mathbb{R} \to 0 < x \to -(-x) = x$}
\AXC{$\stackrel{\Pi_{17}}{\vdots}$}\noLine
\UIC{$(x \in \mathbb{R} \to 0 < x \to -(-x) = x) \to \varphi_{33}$}
\BIC{$\varphi_{33}$}
\end{ptree}
\noindent$\Pi_{19}$:
\begin{ptree}
\AXC{$\stackrel{\Pi_{8}}{\vdots}$}\noLine
\UIC{$x \in \mathbb{R} \to 0 < x \to x = \lvert x \rvert$}
\AXC{$\stackrel{\Pi_{9}}{\vdots}$}\noLine
\UIC{$\lvert -x \rvert = -(-x) \to -(-x) = x \to \lvert -x \rvert = x$}
\AXC{$\stackrel{\Pi_{18}}{\vdots}$}\noLine
\UIC{$\varphi_{33}$}
\BIC{$x \in \mathbb{R} \to 0 < x \to \lvert -x \rvert = x$}
\AXC{$\varphi_{49}$\lbl{A1}}
\BIC{$(x \in \mathbb{R} \to 0 < x \to x = \lvert x \rvert) \to \varphi_{29}$}
\BIC{$\varphi_{29}$}
\end{ptree}
\noindent$\Pi_{20}$:
\begin{ptree}
\AXC{$y = z \to x = y \to x = z$\lbl{A6}}
\AXC{$(y = z \to x = y \to x = z) \to x = y \to y = z \to x = z$\lbl{A1}}
\BIC{$x = y \to y = z \to x = z$}
\UIC{$\forall z\,(x = y \to y = z \to x = z)$}
\UIC{$\forall y\,\forall z\,(x = y \to y = z \to x = z)$}
\UIC{$\forall x\,\forall y\,\forall z\,(x = y \to y = z \to x = z)$}
\AXC{$\varphi_{2}$\lbl{A2}}
\BIC{$\forall y\,\forall z\,(-x = y \to y = z \to -x = z)$}
\AXC{$\varphi_{4}$\lbl{A2}}
\BIC{$\forall z\,(-x = 0 \to 0 = z \to -x = z)$}
\AXC{$\varphi_{1}$\lbl{A2}}
\BIC{$-x = 0 \to 0 = x \to -x = x$}
\end{ptree}
\noindent$\Pi_{21}$:
\begin{ptree}
\AXC{$\stackrel{\Pi_{2}}{\vdots}$}\noLine
\UIC{$\forall x\,\forall y\,(x = y \to y = x)$}
\AXC{$\forall x\,\forall y\,(x = y \to y = x) \to \forall y\,(x = y \to y = x)$\lbl{A2}}
\BIC{$\forall y\,(x = y \to y = x)$}
\AXC{$\forall y\,(x = y \to y = x) \to x = 0 \to 0 = x$\lbl{A2}}
\BIC{$x = 0 \to 0 = x$}
\end{ptree}
\noindent$\Pi_{22}$:
\begin{ptree}
\AXC{$y = z \to x = y \to x = z$\lbl{A6}}
\AXC{$(y = z \to x = y \to x = z) \to x = y \to y = z \to x = z$\lbl{A1}}
\BIC{$x = y \to y = z \to x = z$}
\UIC{$\forall z\,(x = y \to y = z \to x = z)$}
\UIC{$\forall y\,\forall z\,(x = y \to y = z \to x = z)$}
\UIC{$\forall x\,\forall y\,\forall z\,(x = y \to y = z \to x = z)$}
\AXC{$\varphi_{2}$\lbl{A2}}
\BIC{$\forall y\,\forall z\,(-x = y \to y = z \to -x = z)$}
\AXC{$\varphi_{10}$\lbl{A2}}
\BIC{$\forall z\,(-x = -0 \to -0 = z \to -x = z)$}
\AXC{$\varphi_{15}$\lbl{A2}}
\BIC{$-x = -0 \to -0 = 0 \to -x = 0$}
\end{ptree}
\noindent$\Pi_{23}$:
\begin{ptree}
\AXC{$x = x$\lbl{A5}}
\UIC{$\forall x\,(x = x)$}
\AXC{$\forall x\,(x = x) \to -x = -x$\lbl{A2}}
\BIC{$-x = -x$}
\AXC{$x \in \mathbb{R} \to x = 0 \to x = 0$\lbl{A1}}
\AXC{$(x \in \mathbb{R} \to x = 0 \to x = 0) \to -x = -x \to \varphi_{9}$\lbl{A1}}
\BIC{$-x = -x \to \varphi_{9}$}
\BIC{$\varphi_{9}$}
\end{ptree}
\noindent$\Pi_{24}$:
\begin{ptree}
\AXC{$z = w \to -x = -z \to -x = -w$\lbl{A6}}
\UIC{$\forall w\,(z = w \to -x = -z \to -x = -w)$}
\UIC{$\forall z\,\forall w\,(z = w \to -x = -z \to -x = -w)$}
\AXC{$\varphi_{22}$\lbl{A2}}
\BIC{$\forall w\,(x = w \to -x = -x \to -x = -w)$}
\AXC{$\varphi_{16}$\lbl{A2}}
\BIC{$x = 0 \to -x = -x \to -x = -0$}
\AXC{$\stackrel{\Pi_{23}}{\vdots}$}\noLine
\UIC{$\varphi_{9}$}
\BIC{$x \in \mathbb{R} \to x = 0 \to -x = -0$}
\AXC{$(x \in \mathbb{R} \to x = 0 \to -x = -0) \to -0 = 0 \to \varphi_{5}$\lbl{A1}}
\BIC{$-0 = 0 \to \varphi_{5}$}
\end{ptree}
\noindent$\Pi_{25}$:
\begin{ptree}
\AXC{$\stackrel{\Pi_{22}}{\vdots}$}\noLine
\UIC{$-x = -0 \to -0 = 0 \to -x = 0$}
\AXC{$-0 = 0$\lbl{[F6]}}
\AXC{$\stackrel{\Pi_{24}}{\vdots}$}\noLine
\UIC{$-0 = 0 \to \varphi_{5}$}
\BIC{$\varphi_{5}$}
\BIC{$x \in \mathbb{R} \to x = 0 \to -x = 0$}
\AXC{$(x \in \mathbb{R} \to x = 0 \to -x = 0) \to \varphi_{19}$\lbl{A1}}
\BIC{$\varphi_{19}$}
\end{ptree}
\noindent$\Pi_{26}$:
\begin{ptree}
\AXC{$\stackrel{\Pi_{20}}{\vdots}$}\noLine
\UIC{$-x = 0 \to 0 = x \to -x = x$}
\AXC{$\stackrel{\Pi_{21}}{\vdots}$}\noLine
\UIC{$x = 0 \to 0 = x$}
\AXC{$x \in \mathbb{R} \to x = 0 \to x = 0$\lbl{A1}}
\AXC{$\varphi_{3}$\lbl{A1}}
\BIC{$(x = 0 \to 0 = x) \to x \in \mathbb{R} \to x = 0 \to 0 = x$}
\BIC{$x \in \mathbb{R} \to x = 0 \to 0 = x$}
\AXC{$\stackrel{\Pi_{25}}{\vdots}$}\noLine
\UIC{$\varphi_{19}$}
\BIC{$(-x = 0 \to 0 = x \to -x = x) \to x \in \mathbb{R} \to x = 0 \to -x = x$}
\BIC{$x \in \mathbb{R} \to x = 0 \to -x = x$}
\end{ptree}
\noindent$\Pi_{27}$:
\begin{ptree}
\AXC{$\stackrel{\Pi_{26}}{\vdots}$}\noLine
\UIC{$x \in \mathbb{R} \to x = 0 \to -x = x$}
\AXC{$(x \in \mathbb{R} \to x = 0 \to -x = x) \to \lvert -x \rvert = \lvert -x \rvert \to \varphi_{34}$\lbl{A1}}
\BIC{$\lvert -x \rvert = \lvert -x \rvert \to \varphi_{34}$}
\end{ptree}
\noindent$\Pi_{28}$:
\begin{ptree}
\AXC{$x = x$\lbl{A5}}
\UIC{$\forall x\,(x = x)$}
\AXC{$\forall x\,(x = x) \to \lvert -x \rvert = \lvert -x \rvert$\lbl{A2}}
\BIC{$\lvert -x \rvert = \lvert -x \rvert$}
\AXC{$\stackrel{\Pi_{27}}{\vdots}$}\noLine
\UIC{$\lvert -x \rvert = \lvert -x \rvert \to \varphi_{34}$}
\BIC{$\varphi_{34}$}
\end{ptree}
\noindent$\Pi_{29}$:
\begin{ptree}
\AXC{$z = w \to \lvert -x \rvert = \lvert z \rvert \to \lvert -x \rvert = \lvert w \rvert$\lbl{A6}}
\UIC{$\forall w\,(z = w \to \lvert -x \rvert = \lvert z \rvert \to \lvert -x \rvert = \lvert w \rvert)$}
\UIC{$\forall z\,\forall w\,(z = w \to \lvert -x \rvert = \lvert z \rvert \to \lvert -x \rvert = \lvert w \rvert)$}
\AXC{$\varphi_{43}$\lbl{A2}}
\BIC{$\forall w\,(-x = w \to \lvert -x \rvert = \lvert -x \rvert \to \lvert -x \rvert = \lvert w \rvert)$}
\AXC{$\varphi_{47}$\lbl{A2}}
\BIC{$-x = x \to \lvert -x \rvert = \lvert -x \rvert \to \lvert -x \rvert = \lvert x \rvert$}
\AXC{$\stackrel{\Pi_{28}}{\vdots}$}\noLine
\UIC{$\varphi_{34}$}
\BIC{$x \in \mathbb{R} \to x = 0 \to \lvert -x \rvert = \lvert x \rvert$}
\end{ptree}
\noindent$\Pi_{30}$:
\begin{ptree}
\AXC{$y = z \to x = y \to x = z$\lbl{A6}}
\AXC{$(y = z \to x = y \to x = z) \to x = y \to y = z \to x = z$\lbl{A1}}
\BIC{$x = y \to y = z \to x = z$}
\UIC{$\forall z\,(x = y \to y = z \to x = z)$}
\UIC{$\forall y\,\forall z\,(x = y \to y = z \to x = z)$}
\UIC{$\forall x\,\forall y\,\forall z\,(x = y \to y = z \to x = z)$}
\AXC{$\varphi_{7}$\lbl{A2}}
\BIC{$\forall y\,\forall z\,(\lvert -x \rvert = y \to y = z \to \lvert -x \rvert = z)$}
\AXC{$\varphi_{32}$\lbl{A2}}
\BIC{$\forall z\,(\lvert -x \rvert = -x \to -x = z \to \lvert -x \rvert = z)$}
\AXC{$\varphi_{39}$\lbl{A2}}
\BIC{$\lvert -x \rvert = -x \to -x = \lvert x \rvert \to \lvert -x \rvert = \lvert x \rvert$}
\end{ptree}
\noindent$\Pi_{31}$:
\begin{ptree}
\AXC{$\stackrel{\Pi_{2}}{\vdots}$}\noLine
\UIC{$\forall x\,\forall y\,(x = y \to y = x)$}
\AXC{$\forall x\,\forall y\,(x = y \to y = x) \to \forall y\,(\lvert x \rvert = y \to y = \lvert x \rvert)$\lbl{A2}}
\BIC{$\forall y\,(\lvert x \rvert = y \to y = \lvert x \rvert)$}
\end{ptree}
\noindent$\Pi_{32}$:
\begin{ptree}
\AXC{$\stackrel{\Pi_{31}}{\vdots}$}\noLine
\UIC{$\forall y\,(\lvert x \rvert = y \to y = \lvert x \rvert)$}
\AXC{$\forall y\,(\lvert x \rvert = y \to y = \lvert x \rvert) \to \lvert x \rvert = -x \to -x = \lvert x \rvert$\lbl{A2}}
\BIC{$\lvert x \rvert = -x \to -x = \lvert x \rvert$}
\end{ptree}
\noindent$\Pi_{33}$:
\begin{ptree}
\AXC{$x \in \mathbb{R} \to x < 0 \to x < 0$\lbl{A1}}
\AXC{$\forall x\,(x < 0 \to \lvert x \rvert = -x)$\lbl{[A3]}}
\AXC{$\forall x\,(x < 0 \to \lvert x \rvert = -x) \to x < 0 \to \lvert x \rvert = -x$\lbl{A2}}
\BIC{$x < 0 \to \lvert x \rvert = -x$}
\AXC{$\varphi_{31}$\lbl{A1}}
\BIC{$(x \in \mathbb{R} \to x < 0 \to x < 0) \to x \in \mathbb{R} \to x < 0 \to \lvert x \rvert = -x$}
\BIC{$x \in \mathbb{R} \to x < 0 \to \lvert x \rvert = -x$}
\end{ptree}
\noindent$\Pi_{34}$:
\begin{ptree}
\AXC{$\stackrel{\Pi_{32}}{\vdots}$}\noLine
\UIC{$\lvert x \rvert = -x \to -x = \lvert x \rvert$}
\AXC{$\stackrel{\Pi_{33}}{\vdots}$}\noLine
\UIC{$x \in \mathbb{R} \to x < 0 \to \lvert x \rvert = -x$}
\AXC{$(x \in \mathbb{R} \to x < 0 \to \lvert x \rvert = -x) \to \varphi_{6}$\lbl{A1}}
\BIC{$\varphi_{6}$}
\BIC{$x \in \mathbb{R} \to x < 0 \to -x = \lvert x \rvert$}
\end{ptree}
\noindent$\Pi_{35}$:
\begin{ptree}
\AXC{$x \in \mathbb{R} \to x < 0 \to x \in \mathbb{R}$\lbl{A1}}
\AXC{$\forall x\,\forall y\,(x \in \mathbb{R} \to y \in \mathbb{R} \to x < y \to -y < -x)$\lbl{[F22]}}
\AXC{$\varphi_{45}$\lbl{A2}}
\BIC{$\forall y\,(x \in \mathbb{R} \to y \in \mathbb{R} \to x < y \to -y < -x)$}
\AXC{$\varphi_{42}$\lbl{A2}}
\BIC{$x \in \mathbb{R} \to 0 \in \mathbb{R} \to x < 0 \to -0 < -x$}
\AXC{$\varphi_{50}$\lbl{A1}}
\BIC{$(x \in \mathbb{R} \to x < 0 \to x \in \mathbb{R}) \to \varphi_{23}$}
\BIC{$\varphi_{23}$}
\end{ptree}
\noindent$\Pi_{36}$:
\begin{ptree}
\AXC{$x \in \mathbb{R} \to x < 0 \to x < 0$\lbl{A1}}
\AXC{$0 \in \mathbb{R}$\lbl{(R1)}}
\AXC{$\stackrel{\Pi_{35}}{\vdots}$}\noLine
\UIC{$\varphi_{23}$}
\BIC{$(x \in \mathbb{R} \to x < 0 \to x < 0) \to x \in \mathbb{R} \to x < 0 \to -0 < -x$}
\BIC{$x \in \mathbb{R} \to x < 0 \to -0 < -x$}
\end{ptree}
\noindent$\Pi_{37}$:
\begin{ptree}
\AXC{$\stackrel{\Pi_{36}}{\vdots}$}\noLine
\UIC{$x \in \mathbb{R} \to x < 0 \to -0 < -x$}
\AXC{$-0 = 0$\lbl{[F6]}}
\AXC{$-0 = 0 \to (x \in \mathbb{R} \to x < 0 \to -0 < -x) \to \varphi_{18}$\lbl{A1}}
\BIC{$(x \in \mathbb{R} \to x < 0 \to -0 < -x) \to \varphi_{18}$}
\BIC{$\varphi_{18}$}
\end{ptree}
\noindent$\Pi_{38}$:
\begin{ptree}
\AXC{$z = w \to z < -x \to w < -x$\lbl{A6}}
\UIC{$\forall w\,(z = w \to z < -x \to w < -x)$}
\UIC{$\forall z\,\forall w\,(z = w \to z < -x \to w < -x)$}
\AXC{$\varphi_{25}$\lbl{A2}}
\BIC{$\forall w\,(-0 = w \to -0 < -x \to w < -x)$}
\AXC{$\varphi_{30}$\lbl{A2}}
\BIC{$-0 = 0 \to -0 < -x \to 0 < -x$}
\AXC{$\stackrel{\Pi_{37}}{\vdots}$}\noLine
\UIC{$\varphi_{18}$}
\BIC{$x \in \mathbb{R} \to x < 0 \to 0 < -x$}
\end{ptree}
\noindent$\Pi_{39}$:
\begin{ptree}
\AXC{$\forall x\,(0 \leq x \to \lvert x \rvert = x)$\lbl{[A1]}}
\AXC{$\forall x\,(0 \leq x \to \lvert x \rvert = x) \to 0 \leq -x \to \lvert -x \rvert = -x$\lbl{A2}}
\BIC{$0 \leq -x \to \lvert -x \rvert = -x$}
\AXC{$(0 \leq -x \to \lvert -x \rvert = -x) \to \varphi_{11}$\lbl{A1}}
\BIC{$\varphi_{11}$}
\end{ptree}
\noindent$\Pi_{40}$:
\begin{ptree}
\AXC{$\stackrel{\Pi_{38}}{\vdots}$}\noLine
\UIC{$x \in \mathbb{R} \to x < 0 \to 0 < -x$}
\AXC{$(x \in \mathbb{R} \to x < 0 \to 0 < -x) \to x \in \mathbb{R} \to x < 0 \to 0 \leq -x$\lbl{A1}}
\BIC{$x \in \mathbb{R} \to x < 0 \to 0 \leq -x$}
\AXC{$\stackrel{\Pi_{39}}{\vdots}$}\noLine
\UIC{$\varphi_{11}$}
\BIC{$x \in \mathbb{R} \to x < 0 \to \lvert -x \rvert = -x$}
\AXC{$\varphi_{52}$\lbl{A1}}
\BIC{$(x \in \mathbb{R} \to x < 0 \to -x = \lvert x \rvert) \to \varphi_{35}$}
\end{ptree}
\noindent$\Pi_{41}$:
\begin{ptree}
\AXC{$\stackrel{\Pi_{30}}{\vdots}$}\noLine
\UIC{$\lvert -x \rvert = -x \to -x = \lvert x \rvert \to \lvert -x \rvert = \lvert x \rvert$}
\AXC{$\stackrel{\Pi_{34}}{\vdots}$}\noLine
\UIC{$x \in \mathbb{R} \to x < 0 \to -x = \lvert x \rvert$}
\AXC{$\stackrel{\Pi_{40}}{\vdots}$}\noLine
\UIC{$(x \in \mathbb{R} \to x < 0 \to -x = \lvert x \rvert) \to \varphi_{35}$}
\BIC{$\varphi_{35}$}
\BIC{$x \in \mathbb{R} \to x < 0 \to \lvert -x \rvert = \lvert x \rvert$}
\end{ptree}
\noindent$\Pi_{42}$:
\begin{ptree}
\AXC{$\forall x\,\forall y\,(x \in \mathbb{R} \to y \in \mathbb{R} \to x < y \lor x = y \lor y < x)$\lbl{(R13)}}
\AXC{$\varphi_{40}$\lbl{A2}}
\BIC{$\forall y\,(x \in \mathbb{R} \to y \in \mathbb{R} \to x < y \lor x = y \lor y < x)$}
\AXC{$\varphi_{36}$\lbl{A2}}
\BIC{$x \in \mathbb{R} \to 0 \in \mathbb{R} \to x < 0 \lor x = 0 \lor 0 < x$}
\end{ptree}
\noindent$\Pi_{43}$:
\begin{ptree}
\AXC{$x \in \mathbb{R} \to x \in \mathbb{R}$\lbl{A1}}
\AXC{$\stackrel{\Pi_{42}}{\vdots}$}\noLine
\UIC{$x \in \mathbb{R} \to 0 \in \mathbb{R} \to x < 0 \lor x = 0 \lor 0 < x$}
\AXC{$(x \in \mathbb{R} \to 0 \in \mathbb{R} \to x < 0 \lor x = 0 \lor 0 < x) \to \varphi_{12}$\lbl{A1}}
\BIC{$\varphi_{12}$}
\BIC{$0 \in \mathbb{R} \to x \in \mathbb{R} \to x < 0 \lor x = 0 \lor 0 < x$}
\end{ptree}
\noindent$\Pi_{44}$:
\begin{ptree}
\AXC{$0 \in \mathbb{R}$\lbl{(R1)}}
\AXC{$\stackrel{\Pi_{43}}{\vdots}$}\noLine
\UIC{$0 \in \mathbb{R} \to x \in \mathbb{R} \to x < 0 \lor x = 0 \lor 0 < x$}
\BIC{$x \in \mathbb{R} \to x < 0 \lor x = 0 \lor 0 < x$}
\AXC{$(x \in \mathbb{R} \to x < 0 \lor x = 0 \lor 0 < x) \to \varphi_{48}$\lbl{A1}}
\BIC{$\varphi_{48}$}
\end{ptree}
\noindent$\Pi_{45}$:
\begin{ptree}
\AXC{$\stackrel{\Pi_{29}}{\vdots}$}\noLine
\UIC{$x \in \mathbb{R} \to x = 0 \to \lvert -x \rvert = \lvert x \rvert$}
\AXC{$\stackrel{\Pi_{41}}{\vdots}$}\noLine
\UIC{$x \in \mathbb{R} \to x < 0 \to \lvert -x \rvert = \lvert x \rvert$}
\AXC{$\stackrel{\Pi_{44}}{\vdots}$}\noLine
\UIC{$\varphi_{48}$}
\BIC{$(x \in \mathbb{R} \to x = 0 \to \lvert -x \rvert = \lvert x \rvert) \to \varphi_{13}$}
\BIC{$\varphi_{13}$}
\end{ptree}
\noindent 主木:
\begin{ptree}
\AXC{$\stackrel{\Pi_{1}}{\vdots}$}\noLine
\UIC{$\lvert -x \rvert = x \to x = \lvert x \rvert \to \lvert -x \rvert = \lvert x \rvert$}
\AXC{$\stackrel{\Pi_{19}}{\vdots}$}\noLine
\UIC{$\varphi_{29}$}
\BIC{$x \in \mathbb{R} \to 0 < x \to \lvert -x \rvert = \lvert x \rvert$}
\AXC{$\stackrel{\Pi_{45}}{\vdots}$}\noLine
\UIC{$\varphi_{13}$}
\BIC{$x \in \mathbb{R} \to \lvert -x \rvert = \lvert x \rvert$}
\UIC{$\forall x \in \mathbb{R}\,(\lvert -x \rvert = \lvert x \rvert)$}
\end{ptree}
\begin{abbrevlist}
\item $\varphi_{1} :\equiv \forall z\,(-x = 0 \to 0 = z \to -x = z) \to -x = 0 \to 0 = x \to -x = x$
\item $\varphi_{2} :\equiv \forall x\,\forall y\,\forall z\,(x = y \to y = z \to x = z) \to \forall y\,\forall z\,(-x = y \to y = z \to -x = z)$
\item $\varphi_{3} :\equiv (x \in \mathbb{R} \to x = 0 \to x = 0) \to (x = 0 \to 0 = x) \to x \in \mathbb{R} \to x = 0 \to 0 = x$
\item $\varphi_{4} :\equiv \forall y\,\forall z\,(-x = y \to y = z \to -x = z) \to \forall z\,(-x = 0 \to 0 = z \to -x = z)$
\item $\varphi_{5} :\equiv (-x = -0 \to -0 = 0 \to -x = 0) \to x \in \mathbb{R} \to x = 0 \to -x = 0$
\item $\varphi_{6} :\equiv (\lvert x \rvert = -x \to -x = \lvert x \rvert) \to x \in \mathbb{R} \to x < 0 \to -x = \lvert x \rvert$
\item $\varphi_{7} :\equiv \forall x\,\forall y\,\forall z\,(x = y \to y = z \to x = z) \to \forall y\,\forall z\,(\lvert -x \rvert = y \to y = z \to \lvert -x \rvert = z)$
\item $\varphi_{8} :\equiv (x \in \mathbb{R} \to 0 < x \to x \in \mathbb{R}) \to x \in \mathbb{R} \to 0 < x \to -(-x) = x$
\item $\varphi_{9} :\equiv (x = 0 \to -x = -x \to -x = -0) \to x \in \mathbb{R} \to x = 0 \to -x = -0$
\item $\varphi_{10} :\equiv \forall y\,\forall z\,(-x = y \to y = z \to -x = z) \to \forall z\,(-x = -0 \to -0 = z \to -x = z)$
\item $\varphi_{11} :\equiv (x \in \mathbb{R} \to x < 0 \to 0 \leq -x) \to x \in \mathbb{R} \to x < 0 \to \lvert -x \rvert = -x$
\item $\varphi_{12} :\equiv (x \in \mathbb{R} \to x \in \mathbb{R}) \to 0 \in \mathbb{R} \to x \in \mathbb{R} \to x < 0 \lor x = 0 \lor 0 < x$
\item $\varphi_{13} :\equiv (x \in \mathbb{R} \to 0 < x \to \lvert -x \rvert = \lvert x \rvert) \to x \in \mathbb{R} \to \lvert -x \rvert = \lvert x \rvert$
\item $\varphi_{14} :\equiv (0 \leq x \to \lvert x \rvert = x) \to (x \in \mathbb{R} \to 0 < x \to 0 \leq x) \to x \in \mathbb{R} \to 0 < x \to \lvert x \rvert = x$
\item $\varphi_{15} :\equiv \forall z\,(-x = -0 \to -0 = z \to -x = z) \to -x = -0 \to -0 = 0 \to -x = 0$
\item $\varphi_{16} :\equiv \forall w\,(x = w \to -x = -x \to -x = -w) \to x = 0 \to -x = -x \to -x = -0$
\item $\varphi_{17} :\equiv (-0 = 0 \to -x < -0 \to -x < 0) \to x \in \mathbb{R} \to 0 < x \to -x < 0$
\item $\varphi_{18} :\equiv (-0 = 0 \to -0 < -x \to 0 < -x) \to x \in \mathbb{R} \to x < 0 \to 0 < -x$
\item $\varphi_{19} :\equiv (x \in \mathbb{R} \to x = 0 \to 0 = x) \to (-x = 0 \to 0 = x \to -x = x) \to x \in \mathbb{R} \to x = 0 \to -x = x$
\item $\varphi_{20} :\equiv (x \in \mathbb{R} \to 0 < x \to -x < 0) \to x \in \mathbb{R} \to 0 < x \to \lvert -x \rvert = -(-x)$
\item $\varphi_{21} :\equiv \forall y\,\forall z\,(\lvert -x \rvert = y \to y = z \to \lvert -x \rvert = z) \to \forall z\,(\lvert -x \rvert = x \to x = z \to \lvert -x \rvert = z)$
\item $\varphi_{22} :\equiv \forall z\,\forall w\,(z = w \to -x = -z \to -x = -w) \to \forall w\,(x = w \to -x = -x \to -x = -w)$
\item $\varphi_{23} :\equiv 0 \in \mathbb{R} \to (x \in \mathbb{R} \to x < 0 \to x < 0) \to x \in \mathbb{R} \to x < 0 \to -0 < -x$
\item $\varphi_{24} :\equiv \forall z\,\forall w\,(z = w \to -x < z \to -x < w) \to \forall w\,(-0 = w \to -x < -0 \to -x < w)$
\item $\varphi_{25} :\equiv \forall z\,\forall w\,(z = w \to z < -x \to w < -x) \to \forall w\,(-0 = w \to -0 < -x \to w < -x)$
\item $\varphi_{26} :\equiv (x \in \mathbb{R} \to 0 < x \to \lvert x \rvert = x) \to (\lvert x \rvert = x \to x = \lvert x \rvert) \to x \in \mathbb{R} \to 0 < x \to x = \lvert x \rvert$
\item $\varphi_{27} :\equiv \forall z\,(\lvert -x \rvert = x \to x = z \to \lvert -x \rvert = z) \to \lvert -x \rvert = x \to x = \lvert x \rvert \to \lvert -x \rvert = \lvert x \rvert$
\item $\varphi_{28} :\equiv \forall w\,(-0 = w \to -x < -0 \to -x < w) \to -0 = 0 \to -x < -0 \to -x < 0$
\item $\varphi_{29} :\equiv (\lvert -x \rvert = x \to x = \lvert x \rvert \to \lvert -x \rvert = \lvert x \rvert) \to x \in \mathbb{R} \to 0 < x \to \lvert -x \rvert = \lvert x \rvert$
\item $\varphi_{30} :\equiv \forall w\,(-0 = w \to -0 < -x \to w < -x) \to -0 = 0 \to -0 < -x \to 0 < -x$
\item $\varphi_{31} :\equiv (x < 0 \to \lvert x \rvert = -x) \to (x \in \mathbb{R} \to x < 0 \to x < 0) \to x \in \mathbb{R} \to x < 0 \to \lvert x \rvert = -x$
\item $\varphi_{32} :\equiv \forall y\,\forall z\,(\lvert -x \rvert = y \to y = z \to \lvert -x \rvert = z) \to \forall z\,(\lvert -x \rvert = -x \to -x = z \to \lvert -x \rvert = z)$
\item $\varphi_{33} :\equiv (\lvert -x \rvert = -(-x) \to -(-x) = x \to \lvert -x \rvert = x) \to x \in \mathbb{R} \to 0 < x \to \lvert -x \rvert = x$
\item $\varphi_{34} :\equiv (-x = x \to \lvert -x \rvert = \lvert -x \rvert \to \lvert -x \rvert = \lvert x \rvert) \to x \in \mathbb{R} \to x = 0 \to \lvert -x \rvert = \lvert x \rvert$
\item $\varphi_{35} :\equiv (\lvert -x \rvert = -x \to -x = \lvert x \rvert \to \lvert -x \rvert = \lvert x \rvert) \to x \in \mathbb{R} \to x < 0 \to \lvert -x \rvert = \lvert x \rvert$
\item $\varphi_{36} :\equiv \forall y\,(x \in \mathbb{R} \to y \in \mathbb{R} \to x < y \lor x = y \lor y < x) \to x \in \mathbb{R} \to 0 \in \mathbb{R} \to x < 0 \lor x = 0 \lor 0 < x$
\item $\varphi_{37} :\equiv \forall y\,\forall z\,(\lvert -x \rvert = y \to y = z \to \lvert -x \rvert = z) \to \forall z\,(\lvert -x \rvert = -(-x) \to -(-x) = z \to \lvert -x \rvert = z)$
\item $\varphi_{38} :\equiv (x \in \mathbb{R} \to 0 < x \to x \in \mathbb{R}) \to (x \in \mathbb{R} \to 0 < x \to 0 < x) \to x \in \mathbb{R} \to 0 < x \to -x < -0$
\item $\varphi_{39} :\equiv \forall z\,(\lvert -x \rvert = -x \to -x = z \to \lvert -x \rvert = z) \to \lvert -x \rvert = -x \to -x = \lvert x \rvert \to \lvert -x \rvert = \lvert x \rvert$
\item $\varphi_{40} :\equiv \forall x\,\forall y\,(x \in \mathbb{R} \to y \in \mathbb{R} \to x < y \lor x = y \lor y < x) \to \forall y\,(x \in \mathbb{R} \to y \in \mathbb{R} \to x < y \lor x = y \lor y < x)$
\item $\varphi_{41} :\equiv \forall y\,(0 \in \mathbb{R} \to y \in \mathbb{R} \to 0 < y \to -y < -0) \to 0 \in \mathbb{R} \to x \in \mathbb{R} \to 0 < x \to -x < -0$
\item $\varphi_{42} :\equiv \forall y\,(x \in \mathbb{R} \to y \in \mathbb{R} \to x < y \to -y < -x) \to x \in \mathbb{R} \to 0 \in \mathbb{R} \to x < 0 \to -0 < -x$
\item $\varphi_{43} :\equiv \forall z\,\forall w\,(z = w \to \lvert -x \rvert = \lvert z \rvert \to \lvert -x \rvert = \lvert w \rvert) \to \forall w\,(-x = w \to \lvert -x \rvert = \lvert -x \rvert \to \lvert -x \rvert = \lvert w \rvert)$
\item $\varphi_{44} :\equiv \forall x\,\forall y\,(x \in \mathbb{R} \to y \in \mathbb{R} \to x < y \to -y < -x) \to \forall y\,(0 \in \mathbb{R} \to y \in \mathbb{R} \to 0 < y \to -y < -0)$
\item $\varphi_{45} :\equiv \forall x\,\forall y\,(x \in \mathbb{R} \to y \in \mathbb{R} \to x < y \to -y < -x) \to \forall y\,(x \in \mathbb{R} \to y \in \mathbb{R} \to x < y \to -y < -x)$
\item $\varphi_{46} :\equiv \forall z\,(\lvert -x \rvert = -(-x) \to -(-x) = z \to \lvert -x \rvert = z) \to \lvert -x \rvert = -(-x) \to -(-x) = x \to \lvert -x \rvert = x$
\item $\varphi_{47} :\equiv \forall w\,(-x = w \to \lvert -x \rvert = \lvert -x \rvert \to \lvert -x \rvert = \lvert w \rvert) \to -x = x \to \lvert -x \rvert = \lvert -x \rvert \to \lvert -x \rvert = \lvert x \rvert$
\item $\varphi_{48} :\equiv (x \in \mathbb{R} \to x < 0 \to \lvert -x \rvert = \lvert x \rvert) \to (x \in \mathbb{R} \to x = 0 \to \lvert -x \rvert = \lvert x \rvert) \to \varphi_{13}$
\item $\varphi_{49} :\equiv (x \in \mathbb{R} \to 0 < x \to \lvert -x \rvert = x) \to (x \in \mathbb{R} \to 0 < x \to x = \lvert x \rvert) \to \varphi_{29}$
\item $\varphi_{50} :\equiv (x \in \mathbb{R} \to 0 \in \mathbb{R} \to x < 0 \to -0 < -x) \to (x \in \mathbb{R} \to x < 0 \to x \in \mathbb{R}) \to \varphi_{23}$
\item $\varphi_{51} :\equiv (x \in \mathbb{R} \to 0 < x \to \lvert -x \rvert = -(-x)) \to (x \in \mathbb{R} \to 0 < x \to -(-x) = x) \to \varphi_{33}$
\item $\varphi_{52} :\equiv (x \in \mathbb{R} \to x < 0 \to \lvert -x \rvert = -x) \to (x \in \mathbb{R} \to x < 0 \to -x = \lvert x \rvert) \to \varphi_{35}$
\end{abbrevlist}
\end{fullproof}

% 自動生成 (prover/build.js): AbsBound  ノード数 244
\begin{fullproof}{[A8]}{$\forall x \in \mathbb{R}\,(-\lvert x \rvert \leq x \land x \leq \lvert x \rvert)$}{244}\phantomsection\label{tree:AbsBound}
\noindent$\Pi_{1}$:
\begin{ptree}
\AXC{$x \in \mathbb{R} \to \neg 0 \leq x \to x \in \mathbb{R}$\lbl{A1}}
\AXC{$\forall x\,\forall y\,(x \in \mathbb{R} \to y \in \mathbb{R} \to \neg y \leq x \to x < y)$\lbl{[A2]}}
\AXC{$\varphi_{46}$\lbl{A2}}
\BIC{$\forall y\,(x \in \mathbb{R} \to y \in \mathbb{R} \to \neg y \leq x \to x < y)$}
\AXC{$\varphi_{40}$\lbl{A2}}
\BIC{$x \in \mathbb{R} \to 0 \in \mathbb{R} \to \neg 0 \leq x \to x < 0$}
\AXC{$\varphi_{59}$\lbl{A1}}
\BIC{$(x \in \mathbb{R} \to \neg 0 \leq x \to x \in \mathbb{R}) \to \varphi_{20}$}
\BIC{$\varphi_{20}$}
\end{ptree}
\noindent$\Pi_{2}$:
\begin{ptree}
\AXC{$x \in \mathbb{R} \to \neg 0 \leq x \to \neg 0 \leq x$\lbl{A1}}
\AXC{$0 \in \mathbb{R}$\lbl{(R1)}}
\AXC{$\stackrel{\Pi_{1}}{\vdots}$}\noLine
\UIC{$\varphi_{20}$}
\BIC{$(x \in \mathbb{R} \to \neg 0 \leq x \to \neg 0 \leq x) \to x \in \mathbb{R} \to \neg 0 \leq x \to x < 0$}
\BIC{$x \in \mathbb{R} \to \neg 0 \leq x \to x < 0$}
\end{ptree}
\noindent$\Pi_{3}$:
\begin{ptree}
\AXC{$x \in \mathbb{R} \to \neg 0 \leq x \to x \in \mathbb{R}$\lbl{A1}}
\AXC{$\forall x\,\forall y\,(x \in \mathbb{R} \to y \in \mathbb{R} \to x < y \to -y < -x)$\lbl{[F22]}}
\AXC{$\varphi_{54}$\lbl{A2}}
\BIC{$\forall y\,(x \in \mathbb{R} \to y \in \mathbb{R} \to x < y \to -y < -x)$}
\AXC{$\varphi_{51}$\lbl{A2}}
\BIC{$x \in \mathbb{R} \to 0 \in \mathbb{R} \to x < 0 \to -0 < -x$}
\AXC{$\varphi_{63}$\lbl{A1}}
\BIC{$(x \in \mathbb{R} \to \neg 0 \leq x \to x \in \mathbb{R}) \to 0 \in \mathbb{R} \to \varphi_{15}$}
\BIC{$0 \in \mathbb{R} \to \varphi_{15}$}
\end{ptree}
\noindent$\Pi_{4}$:
\begin{ptree}
\AXC{$\stackrel{\Pi_{2}}{\vdots}$}\noLine
\UIC{$x \in \mathbb{R} \to \neg 0 \leq x \to x < 0$}
\AXC{$0 \in \mathbb{R}$\lbl{(R1)}}
\AXC{$\stackrel{\Pi_{3}}{\vdots}$}\noLine
\UIC{$0 \in \mathbb{R} \to \varphi_{15}$}
\BIC{$\varphi_{15}$}
\BIC{$x \in \mathbb{R} \to \neg 0 \leq x \to -0 < -x$}
\end{ptree}
\noindent$\Pi_{5}$:
\begin{ptree}
\AXC{$\stackrel{\Pi_{4}}{\vdots}$}\noLine
\UIC{$x \in \mathbb{R} \to \neg 0 \leq x \to -0 < -x$}
\AXC{$-0 = 0$\lbl{[F6]}}
\AXC{$-0 = 0 \to (x \in \mathbb{R} \to \neg 0 \leq x \to -0 < -x) \to \varphi_{32}$\lbl{A1}}
\BIC{$(x \in \mathbb{R} \to \neg 0 \leq x \to -0 < -x) \to \varphi_{32}$}
\BIC{$\varphi_{32}$}
\end{ptree}
\noindent$\Pi_{6}$:
\begin{ptree}
\AXC{$z = w \to z < -x \to w < -x$\lbl{A6}}
\UIC{$\forall w\,(z = w \to z < -x \to w < -x)$}
\UIC{$\forall z\,\forall w\,(z = w \to z < -x \to w < -x)$}
\AXC{$\varphi_{39}$\lbl{A2}}
\BIC{$\forall w\,(-0 = w \to -0 < -x \to w < -x)$}
\AXC{$\varphi_{43}$\lbl{A2}}
\BIC{$-0 = 0 \to -0 < -x \to 0 < -x$}
\AXC{$\stackrel{\Pi_{5}}{\vdots}$}\noLine
\UIC{$\varphi_{32}$}
\BIC{$x \in \mathbb{R} \to \neg 0 \leq x \to 0 < -x$}
\end{ptree}
\noindent$\Pi_{7}$:
\begin{ptree}
\AXC{$\forall x\,\forall y\,\forall z\,(x < y \to y < z \to x < z)$\lbl{[O2]}}
\AXC{$\varphi_{21}$\lbl{A2}}
\BIC{$\forall y\,\forall z\,(x < y \to y < z \to x < z)$}
\AXC{$\varphi_{9}$\lbl{A2}}
\BIC{$\forall z\,(x < 0 \to 0 < z \to x < z)$}
\AXC{$\forall z\,(x < 0 \to 0 < z \to x < z) \to x < 0 \to 0 < -x \to x < -x$\lbl{A2}}
\BIC{$x < 0 \to 0 < -x \to x < -x$}
\end{ptree}
\noindent$\Pi_{8}$:
\begin{ptree}
\AXC{$\stackrel{\Pi_{6}}{\vdots}$}\noLine
\UIC{$x \in \mathbb{R} \to \neg 0 \leq x \to 0 < -x$}
\AXC{$\stackrel{\Pi_{2}}{\vdots}$}\noLine
\UIC{$x \in \mathbb{R} \to \neg 0 \leq x \to x < 0$}
\AXC{$\stackrel{\Pi_{7}}{\vdots}$}\noLine
\UIC{$x < 0 \to 0 < -x \to x < -x$}
\AXC{$\varphi_{58}$\lbl{A1}}
\BIC{$(x \in \mathbb{R} \to \neg 0 \leq x \to x < 0) \to \varphi_{16}$}
\BIC{$\varphi_{16}$}
\BIC{$x \in \mathbb{R} \to \neg 0 \leq x \to x < -x$}
\AXC{$\varphi_{10}$\lbl{A1}}
\BIC{$x \in \mathbb{R} \to \neg 0 \leq x \to x \leq -x$}
\end{ptree}
\noindent$\Pi_{9}$:
\begin{ptree}
\AXC{$x = x$\lbl{A5}}
\AXC{$x = y \to x = x \to y = x$\lbl{A6}}
\AXC{$(x = y \to x = x \to y = x) \to x = x \to x = y \to y = x$\lbl{A1}}
\BIC{$x = x \to x = y \to y = x$}
\BIC{$x = y \to y = x$}
\UIC{$\forall y\,(x = y \to y = x)$}
\UIC{$\forall x\,\forall y\,(x = y \to y = x)$}
\end{ptree}
\noindent$\Pi_{10}$:
\begin{ptree}
\AXC{$\stackrel{\Pi_{9}}{\vdots}$}\noLine
\UIC{$\forall x\,\forall y\,(x = y \to y = x)$}
\AXC{$\forall x\,\forall y\,(x = y \to y = x) \to \forall y\,(\lvert x \rvert = y \to y = \lvert x \rvert)$\lbl{A2}}
\BIC{$\forall y\,(\lvert x \rvert = y \to y = \lvert x \rvert)$}
\end{ptree}
\noindent$\Pi_{11}$:
\begin{ptree}
\AXC{$\stackrel{\Pi_{10}}{\vdots}$}\noLine
\UIC{$\forall y\,(\lvert x \rvert = y \to y = \lvert x \rvert)$}
\AXC{$\forall y\,(\lvert x \rvert = y \to y = \lvert x \rvert) \to \lvert x \rvert = -x \to -x = \lvert x \rvert$\lbl{A2}}
\BIC{$\lvert x \rvert = -x \to -x = \lvert x \rvert$}
\end{ptree}
\noindent$\Pi_{12}$:
\begin{ptree}
\AXC{$\forall x\,(x < 0 \to \lvert x \rvert = -x)$\lbl{[A3]}}
\AXC{$\forall x\,(x < 0 \to \lvert x \rvert = -x) \to x < 0 \to \lvert x \rvert = -x$\lbl{A2}}
\BIC{$x < 0 \to \lvert x \rvert = -x$}
\AXC{$(x < 0 \to \lvert x \rvert = -x) \to \varphi_{11}$\lbl{A1}}
\BIC{$\varphi_{11}$}
\end{ptree}
\noindent$\Pi_{13}$:
\begin{ptree}
\AXC{$\stackrel{\Pi_{2}}{\vdots}$}\noLine
\UIC{$x \in \mathbb{R} \to \neg 0 \leq x \to x < 0$}
\AXC{$\stackrel{\Pi_{12}}{\vdots}$}\noLine
\UIC{$\varphi_{11}$}
\BIC{$x \in \mathbb{R} \to \neg 0 \leq x \to \lvert x \rvert = -x$}
\end{ptree}
\noindent$\Pi_{14}$:
\begin{ptree}
\AXC{$\stackrel{\Pi_{11}}{\vdots}$}\noLine
\UIC{$\lvert x \rvert = -x \to -x = \lvert x \rvert$}
\AXC{$\stackrel{\Pi_{13}}{\vdots}$}\noLine
\UIC{$x \in \mathbb{R} \to \neg 0 \leq x \to \lvert x \rvert = -x$}
\AXC{$(x \in \mathbb{R} \to \neg 0 \leq x \to \lvert x \rvert = -x) \to \varphi_{17}$\lbl{A1}}
\BIC{$\varphi_{17}$}
\BIC{$x \in \mathbb{R} \to \neg 0 \leq x \to -x = \lvert x \rvert$}
\end{ptree}
\noindent$\Pi_{15}$:
\begin{ptree}
\AXC{$\stackrel{\Pi_{8}}{\vdots}$}\noLine
\UIC{$x \in \mathbb{R} \to \neg 0 \leq x \to x \leq -x$}
\AXC{$\stackrel{\Pi_{14}}{\vdots}$}\noLine
\UIC{$x \in \mathbb{R} \to \neg 0 \leq x \to -x = \lvert x \rvert$}
\AXC{$(x \in \mathbb{R} \to \neg 0 \leq x \to -x = \lvert x \rvert) \to \varphi_{60}$\lbl{A1}}
\BIC{$\varphi_{60}$}
\BIC{$(-x = \lvert x \rvert \to x < -x \to x < \lvert x \rvert) \to \varphi_{22}$}
\end{ptree}
\noindent$\Pi_{16}$:
\begin{ptree}
\AXC{$z = w \to x < z \to x < w$\lbl{A6}}
\UIC{$\forall w\,(z = w \to x < z \to x < w)$}
\UIC{$\forall z\,\forall w\,(z = w \to x < z \to x < w)$}
\AXC{$\varphi_{14}$\lbl{A2}}
\BIC{$\forall w\,(-x = w \to x < -x \to x < w)$}
\AXC{$\varphi_{35}$\lbl{A2}}
\BIC{$-x = \lvert x \rvert \to x < -x \to x < \lvert x \rvert$}
\AXC{$\stackrel{\Pi_{15}}{\vdots}$}\noLine
\UIC{$(-x = \lvert x \rvert \to x < -x \to x < \lvert x \rvert) \to \varphi_{22}$}
\BIC{$\varphi_{22}$}
\end{ptree}
\noindent$\Pi_{17}$:
\begin{ptree}
\AXC{$y = z \to x = y \to x = z$\lbl{A6}}
\AXC{$(y = z \to x = y \to x = z) \to x = y \to y = z \to x = z$\lbl{A1}}
\BIC{$x = y \to y = z \to x = z$}
\UIC{$\forall z\,(x = y \to y = z \to x = z)$}
\UIC{$\forall y\,\forall z\,(x = y \to y = z \to x = z)$}
\UIC{$\forall x\,\forall y\,\forall z\,(x = y \to y = z \to x = z)$}
\AXC{$\varphi_{25}$\lbl{A2}}
\BIC{$\forall y\,\forall z\,(-\lvert x \rvert = y \to y = z \to -\lvert x \rvert = z)$}
\AXC{$\varphi_{49}$\lbl{A2}}
\BIC{$\forall z\,(-\lvert x \rvert = -(-x) \to -(-x) = z \to -\lvert x \rvert = z)$}
\AXC{$\varphi_{56}$\lbl{A2}}
\BIC{$-\lvert x \rvert = -(-x) \to -(-x) = x \to -\lvert x \rvert = x$}
\end{ptree}
\noindent$\Pi_{18}$:
\begin{ptree}
\AXC{$\forall x \in \mathbb{R}\,(-(-x) = x)$\lbl{[F5]}}
\AXC{$\forall x \in \mathbb{R}\,(-(-x) = x) \to x \in \mathbb{R} \to -(-x) = x$\lbl{A2}}
\BIC{$x \in \mathbb{R} \to -(-x) = x$}
\AXC{$(x \in \mathbb{R} \to -(-x) = x) \to \varphi_{23}$\lbl{A1}}
\BIC{$\varphi_{23}$}
\end{ptree}
\noindent$\Pi_{19}$:
\begin{ptree}
\AXC{$\stackrel{\Pi_{13}}{\vdots}$}\noLine
\UIC{$x \in \mathbb{R} \to \neg 0 \leq x \to \lvert x \rvert = -x$}
\AXC{$(x \in \mathbb{R} \to \neg 0 \leq x \to \lvert x \rvert = -x) \to -\lvert x \rvert = -\lvert x \rvert \to \varphi_{52}$\lbl{A1}}
\BIC{$-\lvert x \rvert = -\lvert x \rvert \to \varphi_{52}$}
\end{ptree}
\noindent$\Pi_{20}$:
\begin{ptree}
\AXC{$x = x$\lbl{A5}}
\UIC{$\forall x\,(x = x)$}
\AXC{$\forall x\,(x = x) \to -\lvert x \rvert = -\lvert x \rvert$\lbl{A2}}
\BIC{$-\lvert x \rvert = -\lvert x \rvert$}
\AXC{$\stackrel{\Pi_{19}}{\vdots}$}\noLine
\UIC{$-\lvert x \rvert = -\lvert x \rvert \to \varphi_{52}$}
\BIC{$\varphi_{52}$}
\end{ptree}
\noindent$\Pi_{21}$:
\begin{ptree}
\AXC{$z = w \to -\lvert x \rvert = -z \to -\lvert x \rvert = -w$\lbl{A6}}
\UIC{$\forall w\,(z = w \to -\lvert x \rvert = -z \to -\lvert x \rvert = -w)$}
\UIC{$\forall z\,\forall w\,(z = w \to -\lvert x \rvert = -z \to -\lvert x \rvert = -w)$}
\AXC{$\varphi_{53}$\lbl{A2}}
\BIC{$\forall w\,(\lvert x \rvert = w \to -\lvert x \rvert = -\lvert x \rvert \to -\lvert x \rvert = -w)$}
\AXC{$\varphi_{57}$\lbl{A2}}
\BIC{$\lvert x \rvert = -x \to -\lvert x \rvert = -\lvert x \rvert \to -\lvert x \rvert = -(-x)$}
\AXC{$\stackrel{\Pi_{20}}{\vdots}$}\noLine
\UIC{$\varphi_{52}$}
\BIC{$x \in \mathbb{R} \to \neg 0 \leq x \to -\lvert x \rvert = -(-x)$}
\AXC{$\varphi_{64}$\lbl{A1}}
\BIC{$(x \in \mathbb{R} \to \neg 0 \leq x \to -(-x) = x) \to \varphi_{47}$}
\end{ptree}
\noindent$\Pi_{22}$:
\begin{ptree}
\AXC{$x \in \mathbb{R} \to \neg 0 \leq x \to x \in \mathbb{R}$\lbl{A1}}
\AXC{$\stackrel{\Pi_{18}}{\vdots}$}\noLine
\UIC{$\varphi_{23}$}
\BIC{$x \in \mathbb{R} \to \neg 0 \leq x \to -(-x) = x$}
\AXC{$\stackrel{\Pi_{21}}{\vdots}$}\noLine
\UIC{$(x \in \mathbb{R} \to \neg 0 \leq x \to -(-x) = x) \to \varphi_{47}$}
\BIC{$\varphi_{47}$}
\end{ptree}
\noindent$\Pi_{23}$:
\begin{ptree}
\AXC{$\stackrel{\Pi_{17}}{\vdots}$}\noLine
\UIC{$-\lvert x \rvert = -(-x) \to -(-x) = x \to -\lvert x \rvert = x$}
\AXC{$\stackrel{\Pi_{22}}{\vdots}$}\noLine
\UIC{$\varphi_{47}$}
\BIC{$x \in \mathbb{R} \to \neg 0 \leq x \to -\lvert x \rvert = x$}
\AXC{$\varphi_{26}$\lbl{A1}}
\BIC{$x \in \mathbb{R} \to \neg 0 \leq x \to -\lvert x \rvert \leq x$}
\AXC{$(x \in \mathbb{R} \to \neg 0 \leq x \to -\lvert x \rvert \leq x) \to \varphi_{33}$\lbl{A1}}
\BIC{$\varphi_{33}$}
\end{ptree}
\noindent$\Pi_{24}$:
\begin{ptree}
\AXC{$\stackrel{\Pi_{9}}{\vdots}$}\noLine
\UIC{$\forall x\,\forall y\,(x = y \to y = x)$}
\AXC{$\forall x\,\forall y\,(x = y \to y = x) \to \forall y\,(\lvert x \rvert = y \to y = \lvert x \rvert)$\lbl{A2}}
\BIC{$\forall y\,(\lvert x \rvert = y \to y = \lvert x \rvert)$}
\end{ptree}
\noindent$\Pi_{25}$:
\begin{ptree}
\AXC{$\stackrel{\Pi_{24}}{\vdots}$}\noLine
\UIC{$\forall y\,(\lvert x \rvert = y \to y = \lvert x \rvert)$}
\AXC{$\forall y\,(\lvert x \rvert = y \to y = \lvert x \rvert) \to \lvert x \rvert = x \to x = \lvert x \rvert$\lbl{A2}}
\BIC{$\lvert x \rvert = x \to x = \lvert x \rvert$}
\end{ptree}
\noindent$\Pi_{26}$:
\begin{ptree}
\AXC{$x \in \mathbb{R} \to 0 \leq x \to 0 \leq x$\lbl{A1}}
\AXC{$\forall x\,(0 \leq x \to \lvert x \rvert = x)$\lbl{[A1]}}
\AXC{$\forall x\,(0 \leq x \to \lvert x \rvert = x) \to 0 \leq x \to \lvert x \rvert = x$\lbl{A2}}
\BIC{$0 \leq x \to \lvert x \rvert = x$}
\AXC{$\varphi_{27}$\lbl{A1}}
\BIC{$(x \in \mathbb{R} \to 0 \leq x \to 0 \leq x) \to x \in \mathbb{R} \to 0 \leq x \to \lvert x \rvert = x$}
\BIC{$x \in \mathbb{R} \to 0 \leq x \to \lvert x \rvert = x$}
\end{ptree}
\noindent$\Pi_{27}$:
\begin{ptree}
\AXC{$\stackrel{\Pi_{25}}{\vdots}$}\noLine
\UIC{$\lvert x \rvert = x \to x = \lvert x \rvert$}
\AXC{$\stackrel{\Pi_{26}}{\vdots}$}\noLine
\UIC{$x \in \mathbb{R} \to 0 \leq x \to \lvert x \rvert = x$}
\AXC{$\varphi_{37}$\lbl{A1}}
\BIC{$(\lvert x \rvert = x \to x = \lvert x \rvert) \to x \in \mathbb{R} \to 0 \leq x \to x = \lvert x \rvert$}
\BIC{$x \in \mathbb{R} \to 0 \leq x \to x = \lvert x \rvert$}
\AXC{$(x \in \mathbb{R} \to 0 \leq x \to x = \lvert x \rvert) \to \varphi_{55}$\lbl{A1}}
\BIC{$\varphi_{55}$}
\end{ptree}
\noindent$\Pi_{28}$:
\begin{ptree}
\AXC{$x = x$\lbl{A5}}
\UIC{$\forall x\,(x = x)$}
\AXC{$\forall x\,(x = x) \to x = x$\lbl{A2}}
\BIC{$x = x$}
\AXC{$x = x \to x \in \mathbb{R} \to 0 \leq x \to x \leq x$\lbl{A1}}
\BIC{$x \in \mathbb{R} \to 0 \leq x \to x \leq x$}
\AXC{$\stackrel{\Pi_{27}}{\vdots}$}\noLine
\UIC{$\varphi_{55}$}
\BIC{$(x = \lvert x \rvert \to x < x \to x < \lvert x \rvert) \to \varphi_{5}$}
\end{ptree}
\noindent$\Pi_{29}$:
\begin{ptree}
\AXC{$z = w \to x < z \to x < w$\lbl{A6}}
\UIC{$\forall w\,(z = w \to x < z \to x < w)$}
\UIC{$\forall z\,\forall w\,(z = w \to x < z \to x < w)$}
\AXC{$\varphi_{4}$\lbl{A2}}
\BIC{$\forall w\,(x = w \to x < x \to x < w)$}
\AXC{$\varphi_{7}$\lbl{A2}}
\BIC{$x = \lvert x \rvert \to x < x \to x < \lvert x \rvert$}
\AXC{$\stackrel{\Pi_{28}}{\vdots}$}\noLine
\UIC{$(x = \lvert x \rvert \to x < x \to x < \lvert x \rvert) \to \varphi_{5}$}
\BIC{$\varphi_{5}$}
\end{ptree}
\noindent$\Pi_{30}$:
\begin{ptree}
\AXC{$\forall x\,\forall y\,(x \in \mathbb{R} \to y \in \mathbb{R} \to x \leq y \to -y \leq -x)$\lbl{[F23]}}
\AXC{$\varphi_{50}$\lbl{A2}}
\BIC{$\forall y\,(0 \in \mathbb{R} \to y \in \mathbb{R} \to 0 \leq y \to -y \leq -0)$}
\AXC{$\varphi_{48}$\lbl{A2}}
\BIC{$0 \in \mathbb{R} \to x \in \mathbb{R} \to 0 \leq x \to -x \leq -0$}
\end{ptree}
\noindent$\Pi_{31}$:
\begin{ptree}
\AXC{$0 \in \mathbb{R}$\lbl{(R1)}}
\AXC{$\stackrel{\Pi_{30}}{\vdots}$}\noLine
\UIC{$0 \in \mathbb{R} \to x \in \mathbb{R} \to 0 \leq x \to -x \leq -0$}
\AXC{$(0 \in \mathbb{R} \to x \in \mathbb{R} \to 0 \leq x \to -x \leq -0) \to 0 \in \mathbb{R} \to \varphi_{45}$\lbl{A1}}
\BIC{$0 \in \mathbb{R} \to \varphi_{45}$}
\BIC{$\varphi_{45}$}
\end{ptree}
\noindent$\Pi_{32}$:
\begin{ptree}
\AXC{$x \in \mathbb{R} \to 0 \leq x \to 0 \leq x$\lbl{A1}}
\AXC{$x \in \mathbb{R} \to 0 \leq x \to x \in \mathbb{R}$\lbl{A1}}
\AXC{$\stackrel{\Pi_{31}}{\vdots}$}\noLine
\UIC{$\varphi_{45}$}
\BIC{$(x \in \mathbb{R} \to 0 \leq x \to 0 \leq x) \to x \in \mathbb{R} \to 0 \leq x \to -x \leq -0$}
\BIC{$x \in \mathbb{R} \to 0 \leq x \to -x \leq -0$}
\AXC{$-0 = 0$\lbl{[F6]}}
\AXC{$-0 = 0 \to \varphi_{61}$\lbl{A1}}
\BIC{$\varphi_{61}$}
\BIC{$(-0 = 0 \to -x < -0 \to -x < 0) \to \varphi_{18}$}
\end{ptree}
\noindent$\Pi_{33}$:
\begin{ptree}
\AXC{$z = w \to -x < z \to -x < w$\lbl{A6}}
\UIC{$\forall w\,(z = w \to -x < z \to -x < w)$}
\UIC{$\forall z\,\forall w\,(z = w \to -x < z \to -x < w)$}
\AXC{$\varphi_{41}$\lbl{A2}}
\BIC{$\forall w\,(-0 = w \to -x < -0 \to -x < w)$}
\AXC{$\varphi_{44}$\lbl{A2}}
\BIC{$-0 = 0 \to -x < -0 \to -x < 0$}
\AXC{$\stackrel{\Pi_{32}}{\vdots}$}\noLine
\UIC{$(-0 = 0 \to -x < -0 \to -x < 0) \to \varphi_{18}$}
\BIC{$\varphi_{18}$}
\end{ptree}
\noindent$\Pi_{34}$:
\begin{ptree}
\AXC{$\forall x\,\forall y\,\forall z\,(x \leq y \to y \leq z \to x \leq z)$\lbl{[O5]}}
\AXC{$\varphi_{8}$\lbl{A2}}
\BIC{$\forall y\,\forall z\,(-x \leq y \to y \leq z \to -x \leq z)$}
\AXC{$\varphi_{13}$\lbl{A2}}
\BIC{$\forall z\,(-x \leq 0 \to 0 \leq z \to -x \leq z)$}
\AXC{$\varphi_{6}$\lbl{A2}}
\BIC{$-x \leq 0 \to 0 \leq x \to -x \leq x$}
\AXC{$(-x \leq 0 \to 0 \leq x \to -x \leq x) \to \varphi_{38}$\lbl{A1}}
\BIC{$\varphi_{38}$}
\end{ptree}
\noindent$\Pi_{35}$:
\begin{ptree}
\AXC{$x \in \mathbb{R} \to 0 \leq x \to 0 \leq x$\lbl{A1}}
\AXC{$z = w \to -x = z \to -x = w$\lbl{A6}}
\UIC{$\forall w\,(z = w \to -x = z \to -x = w)$}
\UIC{$\forall z\,\forall w\,(z = w \to -x = z \to -x = w)$}
\AXC{$\varphi_{29}$\lbl{A2}}
\BIC{$\forall w\,(-0 = w \to -x = -0 \to -x = w)$}
\AXC{$\varphi_{34}$\lbl{A2}}
\BIC{$-0 = 0 \to -x = -0 \to -x = 0$}
\AXC{$\stackrel{\Pi_{33}}{\vdots}$}\noLine
\UIC{$\varphi_{18}$}
\BIC{$x \in \mathbb{R} \to 0 \leq x \to -x \leq 0$}
\AXC{$\stackrel{\Pi_{34}}{\vdots}$}\noLine
\UIC{$\varphi_{38}$}
\BIC{$(x \in \mathbb{R} \to 0 \leq x \to 0 \leq x) \to x \in \mathbb{R} \to 0 \leq x \to -x \leq x$}
\BIC{$x \in \mathbb{R} \to 0 \leq x \to -x \leq x$}
\end{ptree}
\noindent$\Pi_{36}$:
\begin{ptree}
\AXC{$\stackrel{\Pi_{9}}{\vdots}$}\noLine
\UIC{$\forall x\,\forall y\,(x = y \to y = x)$}
\AXC{$\forall x\,\forall y\,(x = y \to y = x) \to \forall y\,(\lvert x \rvert = y \to y = \lvert x \rvert)$\lbl{A2}}
\BIC{$\forall y\,(\lvert x \rvert = y \to y = \lvert x \rvert)$}
\end{ptree}
\noindent$\Pi_{37}$:
\begin{ptree}
\AXC{$\stackrel{\Pi_{36}}{\vdots}$}\noLine
\UIC{$\forall y\,(\lvert x \rvert = y \to y = \lvert x \rvert)$}
\AXC{$\forall y\,(\lvert x \rvert = y \to y = \lvert x \rvert) \to \lvert x \rvert = x \to x = \lvert x \rvert$\lbl{A2}}
\BIC{$\lvert x \rvert = x \to x = \lvert x \rvert$}
\end{ptree}
\noindent$\Pi_{38}$:
\begin{ptree}
\AXC{$\stackrel{\Pi_{37}}{\vdots}$}\noLine
\UIC{$\lvert x \rvert = x \to x = \lvert x \rvert$}
\AXC{$\stackrel{\Pi_{26}}{\vdots}$}\noLine
\UIC{$x \in \mathbb{R} \to 0 \leq x \to \lvert x \rvert = x$}
\AXC{$\varphi_{37}$\lbl{A1}}
\BIC{$(\lvert x \rvert = x \to x = \lvert x \rvert) \to x \in \mathbb{R} \to 0 \leq x \to x = \lvert x \rvert$}
\BIC{$x \in \mathbb{R} \to 0 \leq x \to x = \lvert x \rvert$}
\AXC{$(x \in \mathbb{R} \to 0 \leq x \to x = \lvert x \rvert) \to \varphi_{62}$\lbl{A1}}
\BIC{$\varphi_{62}$}
\end{ptree}
\noindent$\Pi_{39}$:
\begin{ptree}
\AXC{$z = w \to -z < x \to -w < x$\lbl{A6}}
\UIC{$\forall w\,(z = w \to -z < x \to -w < x)$}
\UIC{$\forall z\,\forall w\,(z = w \to -z < x \to -w < x)$}
\AXC{$\varphi_{31}$\lbl{A2}}
\BIC{$\forall w\,(x = w \to -x < x \to -w < x)$}
\AXC{$\varphi_{36}$\lbl{A2}}
\BIC{$x = \lvert x \rvert \to -x < x \to -\lvert x \rvert < x$}
\AXC{$\stackrel{\Pi_{35}}{\vdots}$}\noLine
\UIC{$x \in \mathbb{R} \to 0 \leq x \to -x \leq x$}
\AXC{$\stackrel{\Pi_{38}}{\vdots}$}\noLine
\UIC{$\varphi_{62}$}
\BIC{$(x = \lvert x \rvert \to -x < x \to -\lvert x \rvert < x) \to \varphi_{28}$}
\BIC{$\varphi_{28}$}
\end{ptree}
\noindent$\Pi_{40}$:
\begin{ptree}
\AXC{$z = w \to -z = x \to -w = x$\lbl{A6}}
\UIC{$\forall w\,(z = w \to -z = x \to -w = x)$}
\UIC{$\forall z\,\forall w\,(z = w \to -z = x \to -w = x)$}
\AXC{$\varphi_{12}$\lbl{A2}}
\BIC{$\forall w\,(x = w \to -x = x \to -w = x)$}
\AXC{$\varphi_{24}$\lbl{A2}}
\BIC{$x = \lvert x \rvert \to -x = x \to -\lvert x \rvert = x$}
\AXC{$\stackrel{\Pi_{39}}{\vdots}$}\noLine
\UIC{$\varphi_{28}$}
\BIC{$x \in \mathbb{R} \to 0 \leq x \to -\lvert x \rvert \leq x$}
\AXC{$(x \in \mathbb{R} \to 0 \leq x \to -\lvert x \rvert \leq x) \to \varphi_{30}$\lbl{A1}}
\BIC{$\varphi_{30}$}
\end{ptree}
\noindent$\Pi_{41}$:
\begin{ptree}
\AXC{$z = w \to x = z \to x = w$\lbl{A6}}
\UIC{$\forall w\,(z = w \to x = z \to x = w)$}
\UIC{$\forall z\,\forall w\,(z = w \to x = z \to x = w)$}
\AXC{$\varphi_{1}$\lbl{A2}}
\BIC{$\forall w\,(x = w \to x = x \to x = w)$}
\AXC{$\varphi_{2}$\lbl{A2}}
\BIC{$x = \lvert x \rvert \to x = x \to x = \lvert x \rvert$}
\AXC{$\stackrel{\Pi_{29}}{\vdots}$}\noLine
\UIC{$\varphi_{5}$}
\BIC{$x \in \mathbb{R} \to 0 \leq x \to x \leq \lvert x \rvert$}
\AXC{$\stackrel{\Pi_{40}}{\vdots}$}\noLine
\UIC{$\varphi_{30}$}
\BIC{$x \in \mathbb{R} \to 0 \leq x \to -\lvert x \rvert \leq x \land x \leq \lvert x \rvert$}
\end{ptree}
\noindent$\Pi_{42}$:
\begin{ptree}
\AXC{$\stackrel{\Pi_{41}}{\vdots}$}\noLine
\UIC{$x \in \mathbb{R} \to 0 \leq x \to -\lvert x \rvert \leq x \land x \leq \lvert x \rvert$}
\AXC{$(x \in \mathbb{R} \to 0 \leq x \to -\lvert x \rvert \leq x \land x \leq \lvert x \rvert) \to \varphi_{42}$\lbl{A1}}
\BIC{$\varphi_{42}$}
\end{ptree}
\noindent 主木:
\begin{ptree}
\AXC{$z = w \to x = z \to x = w$\lbl{A6}}
\UIC{$\forall w\,(z = w \to x = z \to x = w)$}
\UIC{$\forall z\,\forall w\,(z = w \to x = z \to x = w)$}
\AXC{$\varphi_{3}$\lbl{A2}}
\BIC{$\forall w\,(-x = w \to x = -x \to x = w)$}
\AXC{$\varphi_{19}$\lbl{A2}}
\BIC{$-x = \lvert x \rvert \to x = -x \to x = \lvert x \rvert$}
\AXC{$\stackrel{\Pi_{16}}{\vdots}$}\noLine
\UIC{$\varphi_{22}$}
\BIC{$x \in \mathbb{R} \to \neg 0 \leq x \to x \leq \lvert x \rvert$}
\AXC{$\stackrel{\Pi_{23}}{\vdots}$}\noLine
\UIC{$\varphi_{33}$}
\BIC{$x \in \mathbb{R} \to \neg 0 \leq x \to -\lvert x \rvert \leq x \land x \leq \lvert x \rvert$}
\AXC{$\stackrel{\Pi_{42}}{\vdots}$}\noLine
\UIC{$\varphi_{42}$}
\BIC{$x \in \mathbb{R} \to -\lvert x \rvert \leq x \land x \leq \lvert x \rvert$}
\UIC{$\forall x \in \mathbb{R}\,(-\lvert x \rvert \leq x \land x \leq \lvert x \rvert)$}
\end{ptree}
\begin{abbrevlist}
\item $\varphi_{1} :\equiv \forall z\,\forall w\,(z = w \to x = z \to x = w) \to \forall w\,(x = w \to x = x \to x = w)$
\item $\varphi_{2} :\equiv \forall w\,(x = w \to x = x \to x = w) \to x = \lvert x \rvert \to x = x \to x = \lvert x \rvert$
\item $\varphi_{3} :\equiv \forall z\,\forall w\,(z = w \to x = z \to x = w) \to \forall w\,(-x = w \to x = -x \to x = w)$
\item $\varphi_{4} :\equiv \forall z\,\forall w\,(z = w \to x < z \to x < w) \to \forall w\,(x = w \to x < x \to x < w)$
\item $\varphi_{5} :\equiv (x = \lvert x \rvert \to x = x \to x = \lvert x \rvert) \to x \in \mathbb{R} \to 0 \leq x \to x \leq \lvert x \rvert$
\item $\varphi_{6} :\equiv \forall z\,(-x \leq 0 \to 0 \leq z \to -x \leq z) \to -x \leq 0 \to 0 \leq x \to -x \leq x$
\item $\varphi_{7} :\equiv \forall w\,(x = w \to x < x \to x < w) \to x = \lvert x \rvert \to x < x \to x < \lvert x \rvert$
\item $\varphi_{8} :\equiv \forall x\,\forall y\,\forall z\,(x \leq y \to y \leq z \to x \leq z) \to \forall y\,\forall z\,(-x \leq y \to y \leq z \to -x \leq z)$
\item $\varphi_{9} :\equiv \forall y\,\forall z\,(x < y \to y < z \to x < z) \to \forall z\,(x < 0 \to 0 < z \to x < z)$
\item $\varphi_{10} :\equiv (x \in \mathbb{R} \to \neg 0 \leq x \to x < -x) \to x \in \mathbb{R} \to \neg 0 \leq x \to x \leq -x$
\item $\varphi_{11} :\equiv (x \in \mathbb{R} \to \neg 0 \leq x \to x < 0) \to x \in \mathbb{R} \to \neg 0 \leq x \to \lvert x \rvert = -x$
\item $\varphi_{12} :\equiv \forall z\,\forall w\,(z = w \to -z = x \to -w = x) \to \forall w\,(x = w \to -x = x \to -w = x)$
\item $\varphi_{13} :\equiv \forall y\,\forall z\,(-x \leq y \to y \leq z \to -x \leq z) \to \forall z\,(-x \leq 0 \to 0 \leq z \to -x \leq z)$
\item $\varphi_{14} :\equiv \forall z\,\forall w\,(z = w \to x < z \to x < w) \to \forall w\,(-x = w \to x < -x \to x < w)$
\item $\varphi_{15} :\equiv (x \in \mathbb{R} \to \neg 0 \leq x \to x < 0) \to x \in \mathbb{R} \to \neg 0 \leq x \to -0 < -x$
\item $\varphi_{16} :\equiv (x \in \mathbb{R} \to \neg 0 \leq x \to 0 < -x) \to x \in \mathbb{R} \to \neg 0 \leq x \to x < -x$
\item $\varphi_{17} :\equiv (\lvert x \rvert = -x \to -x = \lvert x \rvert) \to x \in \mathbb{R} \to \neg 0 \leq x \to -x = \lvert x \rvert$
\item $\varphi_{18} :\equiv (-0 = 0 \to -x = -0 \to -x = 0) \to x \in \mathbb{R} \to 0 \leq x \to -x \leq 0$
\item $\varphi_{19} :\equiv \forall w\,(-x = w \to x = -x \to x = w) \to -x = \lvert x \rvert \to x = -x \to x = \lvert x \rvert$
\item $\varphi_{20} :\equiv 0 \in \mathbb{R} \to (x \in \mathbb{R} \to \neg 0 \leq x \to \neg 0 \leq x) \to x \in \mathbb{R} \to \neg 0 \leq x \to x < 0$
\item $\varphi_{21} :\equiv \forall x\,\forall y\,\forall z\,(x < y \to y < z \to x < z) \to \forall y\,\forall z\,(x < y \to y < z \to x < z)$
\item $\varphi_{22} :\equiv (-x = \lvert x \rvert \to x = -x \to x = \lvert x \rvert) \to x \in \mathbb{R} \to \neg 0 \leq x \to x \leq \lvert x \rvert$
\item $\varphi_{23} :\equiv (x \in \mathbb{R} \to \neg 0 \leq x \to x \in \mathbb{R}) \to x \in \mathbb{R} \to \neg 0 \leq x \to -(-x) = x$
\item $\varphi_{24} :\equiv \forall w\,(x = w \to -x = x \to -w = x) \to x = \lvert x \rvert \to -x = x \to -\lvert x \rvert = x$
\item $\varphi_{25} :\equiv \forall x\,\forall y\,\forall z\,(x = y \to y = z \to x = z) \to \forall y\,\forall z\,(-\lvert x \rvert = y \to y = z \to -\lvert x \rvert = z)$
\item $\varphi_{26} :\equiv (x \in \mathbb{R} \to \neg 0 \leq x \to -\lvert x \rvert = x) \to x \in \mathbb{R} \to \neg 0 \leq x \to -\lvert x \rvert \leq x$
\item $\varphi_{27} :\equiv (0 \leq x \to \lvert x \rvert = x) \to (x \in \mathbb{R} \to 0 \leq x \to 0 \leq x) \to x \in \mathbb{R} \to 0 \leq x \to \lvert x \rvert = x$
\item $\varphi_{28} :\equiv (x = \lvert x \rvert \to -x = x \to -\lvert x \rvert = x) \to x \in \mathbb{R} \to 0 \leq x \to -\lvert x \rvert \leq x$
\item $\varphi_{29} :\equiv \forall z\,\forall w\,(z = w \to -x = z \to -x = w) \to \forall w\,(-0 = w \to -x = -0 \to -x = w)$
\item $\varphi_{30} :\equiv (x \in \mathbb{R} \to 0 \leq x \to x \leq \lvert x \rvert) \to x \in \mathbb{R} \to 0 \leq x \to -\lvert x \rvert \leq x \land x \leq \lvert x \rvert$
\item $\varphi_{31} :\equiv \forall z\,\forall w\,(z = w \to -z < x \to -w < x) \to \forall w\,(x = w \to -x < x \to -w < x)$
\item $\varphi_{32} :\equiv (-0 = 0 \to -0 < -x \to 0 < -x) \to x \in \mathbb{R} \to \neg 0 \leq x \to 0 < -x$
\item $\varphi_{33} :\equiv (x \in \mathbb{R} \to \neg 0 \leq x \to x \leq \lvert x \rvert) \to x \in \mathbb{R} \to \neg 0 \leq x \to -\lvert x \rvert \leq x \land x \leq \lvert x \rvert$
\item $\varphi_{34} :\equiv \forall w\,(-0 = w \to -x = -0 \to -x = w) \to -0 = 0 \to -x = -0 \to -x = 0$
\item $\varphi_{35} :\equiv \forall w\,(-x = w \to x < -x \to x < w) \to -x = \lvert x \rvert \to x < -x \to x < \lvert x \rvert$
\item $\varphi_{36} :\equiv \forall w\,(x = w \to -x < x \to -w < x) \to x = \lvert x \rvert \to -x < x \to -\lvert x \rvert < x$
\item $\varphi_{37} :\equiv (x \in \mathbb{R} \to 0 \leq x \to \lvert x \rvert = x) \to (\lvert x \rvert = x \to x = \lvert x \rvert) \to x \in \mathbb{R} \to 0 \leq x \to x = \lvert x \rvert$
\item $\varphi_{38} :\equiv (x \in \mathbb{R} \to 0 \leq x \to -x \leq 0) \to (x \in \mathbb{R} \to 0 \leq x \to 0 \leq x) \to x \in \mathbb{R} \to 0 \leq x \to -x \leq x$
\item $\varphi_{39} :\equiv \forall z\,\forall w\,(z = w \to z < -x \to w < -x) \to \forall w\,(-0 = w \to -0 < -x \to w < -x)$
\item $\varphi_{40} :\equiv \forall y\,(x \in \mathbb{R} \to y \in \mathbb{R} \to \neg y \leq x \to x < y) \to x \in \mathbb{R} \to 0 \in \mathbb{R} \to \neg 0 \leq x \to x < 0$
\item $\varphi_{41} :\equiv \forall z\,\forall w\,(z = w \to -x < z \to -x < w) \to \forall w\,(-0 = w \to -x < -0 \to -x < w)$
\item $\varphi_{42} :\equiv (x \in \mathbb{R} \to \neg 0 \leq x \to -\lvert x \rvert \leq x \land x \leq \lvert x \rvert) \to x \in \mathbb{R} \to -\lvert x \rvert \leq x \land x \leq \lvert x \rvert$
\item $\varphi_{43} :\equiv \forall w\,(-0 = w \to -0 < -x \to w < -x) \to -0 = 0 \to -0 < -x \to 0 < -x$
\item $\varphi_{44} :\equiv \forall w\,(-0 = w \to -x < -0 \to -x < w) \to -0 = 0 \to -x < -0 \to -x < 0$
\item $\varphi_{45} :\equiv (x \in \mathbb{R} \to 0 \leq x \to x \in \mathbb{R}) \to (x \in \mathbb{R} \to 0 \leq x \to 0 \leq x) \to x \in \mathbb{R} \to 0 \leq x \to -x \leq -0$
\item $\varphi_{46} :\equiv \forall x\,\forall y\,(x \in \mathbb{R} \to y \in \mathbb{R} \to \neg y \leq x \to x < y) \to \forall y\,(x \in \mathbb{R} \to y \in \mathbb{R} \to \neg y \leq x \to x < y)$
\item $\varphi_{47} :\equiv (-\lvert x \rvert = -(-x) \to -(-x) = x \to -\lvert x \rvert = x) \to x \in \mathbb{R} \to \neg 0 \leq x \to -\lvert x \rvert = x$
\item $\varphi_{48} :\equiv \forall y\,(0 \in \mathbb{R} \to y \in \mathbb{R} \to 0 \leq y \to -y \leq -0) \to 0 \in \mathbb{R} \to x \in \mathbb{R} \to 0 \leq x \to -x \leq -0$
\item $\varphi_{49} :\equiv \forall y\,\forall z\,(-\lvert x \rvert = y \to y = z \to -\lvert x \rvert = z) \to \forall z\,(-\lvert x \rvert = -(-x) \to -(-x) = z \to -\lvert x \rvert = z)$
\item $\varphi_{50} :\equiv \forall x\,\forall y\,(x \in \mathbb{R} \to y \in \mathbb{R} \to x \leq y \to -y \leq -x) \to \forall y\,(0 \in \mathbb{R} \to y \in \mathbb{R} \to 0 \leq y \to -y \leq -0)$
\item $\varphi_{51} :\equiv \forall y\,(x \in \mathbb{R} \to y \in \mathbb{R} \to x < y \to -y < -x) \to x \in \mathbb{R} \to 0 \in \mathbb{R} \to x < 0 \to -0 < -x$
\item $\varphi_{52} :\equiv (\lvert x \rvert = -x \to -\lvert x \rvert = -\lvert x \rvert \to -\lvert x \rvert = -(-x)) \to x \in \mathbb{R} \to \neg 0 \leq x \to -\lvert x \rvert = -(-x)$
\item $\varphi_{53} :\equiv \forall z\,\forall w\,(z = w \to -\lvert x \rvert = -z \to -\lvert x \rvert = -w) \to \forall w\,(\lvert x \rvert = w \to -\lvert x \rvert = -\lvert x \rvert \to -\lvert x \rvert = -w)$
\item $\varphi_{54} :\equiv \forall x\,\forall y\,(x \in \mathbb{R} \to y \in \mathbb{R} \to x < y \to -y < -x) \to \forall y\,(x \in \mathbb{R} \to y \in \mathbb{R} \to x < y \to -y < -x)$
\item $\varphi_{55} :\equiv (x \in \mathbb{R} \to 0 \leq x \to x \leq x) \to (x = \lvert x \rvert \to x < x \to x < \lvert x \rvert) \to \varphi_{5}$
\item $\varphi_{56} :\equiv \forall z\,(-\lvert x \rvert = -(-x) \to -(-x) = z \to -\lvert x \rvert = z) \to -\lvert x \rvert = -(-x) \to -(-x) = x \to -\lvert x \rvert = x$
\item $\varphi_{57} :\equiv \forall w\,(\lvert x \rvert = w \to -\lvert x \rvert = -\lvert x \rvert \to -\lvert x \rvert = -w) \to \lvert x \rvert = -x \to -\lvert x \rvert = -\lvert x \rvert \to -\lvert x \rvert = -(-x)$
\item $\varphi_{58} :\equiv (x < 0 \to 0 < -x \to x < -x) \to (x \in \mathbb{R} \to \neg 0 \leq x \to x < 0) \to \varphi_{16}$
\item $\varphi_{59} :\equiv (x \in \mathbb{R} \to 0 \in \mathbb{R} \to \neg 0 \leq x \to x < 0) \to (x \in \mathbb{R} \to \neg 0 \leq x \to x \in \mathbb{R}) \to \varphi_{20}$
\item $\varphi_{60} :\equiv (x \in \mathbb{R} \to \neg 0 \leq x \to x \leq -x) \to (-x = \lvert x \rvert \to x < -x \to x < \lvert x \rvert) \to \varphi_{22}$
\item $\varphi_{61} :\equiv (x \in \mathbb{R} \to 0 \leq x \to -x \leq -0) \to (-0 = 0 \to -x < -0 \to -x < 0) \to \varphi_{18}$
\item $\varphi_{62} :\equiv (x \in \mathbb{R} \to 0 \leq x \to -x \leq x) \to (x = \lvert x \rvert \to -x < x \to -\lvert x \rvert < x) \to \varphi_{28}$
\item $\varphi_{63} :\equiv (x \in \mathbb{R} \to 0 \in \mathbb{R} \to x < 0 \to -0 < -x) \to (x \in \mathbb{R} \to \neg 0 \leq x \to x \in \mathbb{R}) \to 0 \in \mathbb{R} \to \varphi_{15}$
\item $\varphi_{64} :\equiv (x \in \mathbb{R} \to \neg 0 \leq x \to -\lvert x \rvert = -(-x)) \to (x \in \mathbb{R} \to \neg 0 \leq x \to -(-x) = x) \to \varphi_{47}$
\end{abbrevlist}
\end{fullproof}

% 自動生成 (prover/build.js): AbsLe  ノード数 66
\begin{fullproof}{[A9]}{$\forall x\,\forall y\,(x \in \mathbb{R} \to y \in \mathbb{R} \to \lvert x \rvert \leq y \to -y \leq x \land x \leq y)$}{66}\phantomsection\label{tree:AbsLe}
\noindent$\Pi_{1}$:
\begin{ptree}
\AXC{$\forall x \in \mathbb{R}\,(-\lvert x \rvert \leq x \land x \leq \lvert x \rvert)$\lbl{[A8]}}
\AXC{$\varphi_{9}$\lbl{A2}}
\BIC{$x \in \mathbb{R} \to -\lvert x \rvert \leq x \land x \leq \lvert x \rvert$}
\AXC{$(x \in \mathbb{R} \to -\lvert x \rvert \leq x \land x \leq \lvert x \rvert) \to \varphi_{18}$\lbl{A1}}
\BIC{$\varphi_{18}$}
\end{ptree}
\noindent$\Pi_{2}$:
\begin{ptree}
\AXC{$x \in \mathbb{R} \to y \in \mathbb{R} \to \lvert x \rvert \leq y \to x \in \mathbb{R}$\lbl{A1}}
\AXC{$\stackrel{\Pi_{1}}{\vdots}$}\noLine
\UIC{$\varphi_{18}$}
\BIC{$x \in \mathbb{R} \to y \in \mathbb{R} \to \lvert x \rvert \leq y \to -\lvert x \rvert \leq x \land x \leq \lvert x \rvert$}
\end{ptree}
\noindent$\Pi_{3}$:
\begin{ptree}
\AXC{$\forall x\,\forall y\,\forall z\,(x \leq y \to y \leq z \to x \leq z)$\lbl{[O5]}}
\AXC{$\varphi_{2}$\lbl{A2}}
\BIC{$\forall y\,\forall z\,(x \leq y \to y \leq z \to x \leq z)$}
\AXC{$\varphi_{4}$\lbl{A2}}
\BIC{$\forall z\,(x \leq \lvert x \rvert \to \lvert x \rvert \leq z \to x \leq z)$}
\AXC{$\varphi_{6}$\lbl{A2}}
\BIC{$x \leq \lvert x \rvert \to \lvert x \rvert \leq y \to x \leq y$}
\AXC{$\varphi_{22}$\lbl{A1}}
\BIC{$(x \in \mathbb{R} \to y \in \mathbb{R} \to \lvert x \rvert \leq y \to x \leq \lvert x \rvert) \to \varphi_{10}$}
\end{ptree}
\noindent$\Pi_{4}$:
\begin{ptree}
\AXC{$\stackrel{\Pi_{2}}{\vdots}$}\noLine
\UIC{$x \in \mathbb{R} \to y \in \mathbb{R} \to \lvert x \rvert \leq y \to -\lvert x \rvert \leq x \land x \leq \lvert x \rvert$}
\AXC{$\varphi_{16}$\lbl{A1}}
\BIC{$x \in \mathbb{R} \to y \in \mathbb{R} \to \lvert x \rvert \leq y \to x \leq \lvert x \rvert$}
\AXC{$\stackrel{\Pi_{3}}{\vdots}$}\noLine
\UIC{$(x \in \mathbb{R} \to y \in \mathbb{R} \to \lvert x \rvert \leq y \to x \leq \lvert x \rvert) \to \varphi_{10}$}
\BIC{$\varphi_{10}$}
\end{ptree}
\noindent$\Pi_{5}$:
\begin{ptree}
\AXC{$\forall x \in \mathbb{R}\,(\lvert x \rvert \in \mathbb{R})$\lbl{[A5]}}
\AXC{$\forall x \in \mathbb{R}\,(\lvert x \rvert \in \mathbb{R}) \to x \in \mathbb{R} \to \lvert x \rvert \in \mathbb{R}$\lbl{A2}}
\BIC{$x \in \mathbb{R} \to \lvert x \rvert \in \mathbb{R}$}
\AXC{$(x \in \mathbb{R} \to \lvert x \rvert \in \mathbb{R}) \to \varphi_{13}$\lbl{A1}}
\BIC{$\varphi_{13}$}
\end{ptree}
\noindent$\Pi_{6}$:
\begin{ptree}
\AXC{$\forall x\,\forall y\,(x \in \mathbb{R} \to y \in \mathbb{R} \to x \leq y \to -y \leq -x)$\lbl{[F23]}}
\AXC{$\varphi_{20}$\lbl{A2}}
\BIC{$\forall y\,(\lvert x \rvert \in \mathbb{R} \to y \in \mathbb{R} \to \lvert x \rvert \leq y \to -y \leq -\lvert x \rvert)$}
\AXC{$\varphi_{21}$\lbl{A2}}
\BIC{$\lvert x \rvert \in \mathbb{R} \to y \in \mathbb{R} \to \lvert x \rvert \leq y \to -y \leq -\lvert x \rvert$}
\end{ptree}
\noindent$\Pi_{7}$:
\begin{ptree}
\AXC{$\stackrel{\Pi_{6}}{\vdots}$}\noLine
\UIC{$\lvert x \rvert \in \mathbb{R} \to y \in \mathbb{R} \to \lvert x \rvert \leq y \to -y \leq -\lvert x \rvert$}
\AXC{$(\lvert x \rvert \in \mathbb{R} \to y \in \mathbb{R} \to \lvert x \rvert \leq y \to -y \leq -\lvert x \rvert) \to \varphi_{24}$\lbl{A1}}
\BIC{$\varphi_{24}$}
\end{ptree}
\noindent$\Pi_{8}$:
\begin{ptree}
\AXC{$x \in \mathbb{R} \to y \in \mathbb{R} \to \lvert x \rvert \leq y \to y \in \mathbb{R}$\lbl{A1}}
\AXC{$x \in \mathbb{R} \to y \in \mathbb{R} \to \lvert x \rvert \leq y \to x \in \mathbb{R}$\lbl{A1}}
\AXC{$\stackrel{\Pi_{5}}{\vdots}$}\noLine
\UIC{$\varphi_{13}$}
\BIC{$x \in \mathbb{R} \to y \in \mathbb{R} \to \lvert x \rvert \leq y \to \lvert x \rvert \in \mathbb{R}$}
\AXC{$\stackrel{\Pi_{7}}{\vdots}$}\noLine
\UIC{$\varphi_{24}$}
\BIC{$(x \in \mathbb{R} \to y \in \mathbb{R} \to \lvert x \rvert \leq y \to y \in \mathbb{R}) \to \varphi_{14}$}
\BIC{$\varphi_{14}$}
\end{ptree}
\noindent$\Pi_{9}$:
\begin{ptree}
\AXC{$\forall x\,\forall y\,\forall z\,(x \leq y \to y \leq z \to x \leq z)$\lbl{[O5]}}
\AXC{$\varphi_{5}$\lbl{A2}}
\BIC{$\forall y\,\forall z\,(v_{285} \leq y \to y \leq z \to v_{285} \leq z)$}
\AXC{$\varphi_{7}$\lbl{A2}}
\BIC{$\forall z\,(v_{285} \leq v_{286} \to v_{286} \leq z \to v_{285} \leq z)$}
\AXC{$\varphi_{8}$\lbl{A2}}
\BIC{$v_{285} \leq v_{286} \to v_{286} \leq v_{287} \to v_{285} \leq v_{287}$}
\UIC{$\forall v_{287}\,(v_{285} \leq v_{286} \to v_{286} \leq v_{287} \to v_{285} \leq v_{287})$}
\UIC{$\varphi_{1}$}
\UIC{$\forall v_{285}\,(\varphi_{1})$}
\AXC{$\forall v_{285}\,(\varphi_{1}) \to \varphi_{3}$\lbl{A2}}
\BIC{$\varphi_{3}$}
\AXC{$\varphi_{17}$\lbl{A2}}
\BIC{$\forall v_{287}\,(-y \leq -\lvert x \rvert \to -\lvert x \rvert \leq v_{287} \to -y \leq v_{287})$}
\end{ptree}
\noindent$\Pi_{10}$:
\begin{ptree}
\AXC{$\stackrel{\Pi_{9}}{\vdots}$}\noLine
\UIC{$\forall v_{287}\,(-y \leq -\lvert x \rvert \to -\lvert x \rvert \leq v_{287} \to -y \leq v_{287})$}
\AXC{$\varphi_{15}$\lbl{A2}}
\BIC{$-y \leq -\lvert x \rvert \to -\lvert x \rvert \leq x \to -y \leq x$}
\AXC{$\varphi_{23}$\lbl{A1}}
\BIC{$(x \in \mathbb{R} \to y \in \mathbb{R} \to \lvert x \rvert \leq y \to -y \leq -\lvert x \rvert) \to \varphi_{12}$}
\end{ptree}
\noindent$\Pi_{11}$:
\begin{ptree}
\AXC{$x \in \mathbb{R} \to y \in \mathbb{R} \to \lvert x \rvert \leq y \to \lvert x \rvert \leq y$\lbl{A1}}
\AXC{$\stackrel{\Pi_{8}}{\vdots}$}\noLine
\UIC{$\varphi_{14}$}
\BIC{$x \in \mathbb{R} \to y \in \mathbb{R} \to \lvert x \rvert \leq y \to -y \leq -\lvert x \rvert$}
\AXC{$\stackrel{\Pi_{10}}{\vdots}$}\noLine
\UIC{$(x \in \mathbb{R} \to y \in \mathbb{R} \to \lvert x \rvert \leq y \to -y \leq -\lvert x \rvert) \to \varphi_{12}$}
\BIC{$\varphi_{12}$}
\end{ptree}
\noindent$\Pi_{12}$:
\begin{ptree}
\AXC{$\stackrel{\Pi_{2}}{\vdots}$}\noLine
\UIC{$x \in \mathbb{R} \to y \in \mathbb{R} \to \lvert x \rvert \leq y \to -\lvert x \rvert \leq x \land x \leq \lvert x \rvert$}
\AXC{$\varphi_{19}$\lbl{A1}}
\BIC{$x \in \mathbb{R} \to y \in \mathbb{R} \to \lvert x \rvert \leq y \to -\lvert x \rvert \leq x$}
\AXC{$\stackrel{\Pi_{11}}{\vdots}$}\noLine
\UIC{$\varphi_{12}$}
\BIC{$x \in \mathbb{R} \to y \in \mathbb{R} \to \lvert x \rvert \leq y \to -y \leq x$}
\AXC{$(x \in \mathbb{R} \to y \in \mathbb{R} \to \lvert x \rvert \leq y \to -y \leq x) \to \varphi_{11}$\lbl{A1}}
\BIC{$\varphi_{11}$}
\end{ptree}
\noindent 主木:
\begin{ptree}
\AXC{$x \in \mathbb{R} \to y \in \mathbb{R} \to \lvert x \rvert \leq y \to \lvert x \rvert \leq y$\lbl{A1}}
\AXC{$\stackrel{\Pi_{4}}{\vdots}$}\noLine
\UIC{$\varphi_{10}$}
\BIC{$x \in \mathbb{R} \to y \in \mathbb{R} \to \lvert x \rvert \leq y \to x \leq y$}
\AXC{$\stackrel{\Pi_{12}}{\vdots}$}\noLine
\UIC{$\varphi_{11}$}
\BIC{$x \in \mathbb{R} \to y \in \mathbb{R} \to \lvert x \rvert \leq y \to -y \leq x \land x \leq y$}
\UIC{$\forall y\,(x \in \mathbb{R} \to y \in \mathbb{R} \to \lvert x \rvert \leq y \to -y \leq x \land x \leq y)$}
\UIC{$\forall x\,\forall y\,(x \in \mathbb{R} \to y \in \mathbb{R} \to \lvert x \rvert \leq y \to -y \leq x \land x \leq y)$}
\end{ptree}
\begin{abbrevlist}
\item $\varphi_{1} :\equiv \forall v_{286}\,\forall v_{287}\,(v_{285} \leq v_{286} \to v_{286} \leq v_{287} \to v_{285} \leq v_{287})$
\item $\varphi_{2} :\equiv \forall x\,\forall y\,\forall z\,(x \leq y \to y \leq z \to x \leq z) \to \forall y\,\forall z\,(x \leq y \to y \leq z \to x \leq z)$
\item $\varphi_{3} :\equiv \forall v_{286}\,\forall v_{287}\,(-y \leq v_{286} \to v_{286} \leq v_{287} \to -y \leq v_{287})$
\item $\varphi_{4} :\equiv \forall y\,\forall z\,(x \leq y \to y \leq z \to x \leq z) \to \forall z\,(x \leq \lvert x \rvert \to \lvert x \rvert \leq z \to x \leq z)$
\item $\varphi_{5} :\equiv \forall x\,\forall y\,\forall z\,(x \leq y \to y \leq z \to x \leq z) \to \forall y\,\forall z\,(v_{285} \leq y \to y \leq z \to v_{285} \leq z)$
\item $\varphi_{6} :\equiv \forall z\,(x \leq \lvert x \rvert \to \lvert x \rvert \leq z \to x \leq z) \to x \leq \lvert x \rvert \to \lvert x \rvert \leq y \to x \leq y$
\item $\varphi_{7} :\equiv \forall y\,\forall z\,(v_{285} \leq y \to y \leq z \to v_{285} \leq z) \to \forall z\,(v_{285} \leq v_{286} \to v_{286} \leq z \to v_{285} \leq z)$
\item $\varphi_{8} :\equiv \forall z\,(v_{285} \leq v_{286} \to v_{286} \leq z \to v_{285} \leq z) \to v_{285} \leq v_{286} \to v_{286} \leq v_{287} \to v_{285} \leq v_{287}$
\item $\varphi_{9} :\equiv \forall x \in \mathbb{R}\,(-\lvert x \rvert \leq x \land x \leq \lvert x \rvert) \to x \in \mathbb{R} \to -\lvert x \rvert \leq x \land x \leq \lvert x \rvert$
\item $\varphi_{10} :\equiv (x \in \mathbb{R} \to y \in \mathbb{R} \to \lvert x \rvert \leq y \to \lvert x \rvert \leq y) \to x \in \mathbb{R} \to y \in \mathbb{R} \to \lvert x \rvert \leq y \to x \leq y$
\item $\varphi_{11} :\equiv (x \in \mathbb{R} \to y \in \mathbb{R} \to \lvert x \rvert \leq y \to x \leq y) \to x \in \mathbb{R} \to y \in \mathbb{R} \to \lvert x \rvert \leq y \to -y \leq x \land x \leq y$
\item $\varphi_{12} :\equiv (x \in \mathbb{R} \to y \in \mathbb{R} \to \lvert x \rvert \leq y \to -\lvert x \rvert \leq x) \to x \in \mathbb{R} \to y \in \mathbb{R} \to \lvert x \rvert \leq y \to -y \leq x$
\item $\varphi_{13} :\equiv (x \in \mathbb{R} \to y \in \mathbb{R} \to \lvert x \rvert \leq y \to x \in \mathbb{R}) \to x \in \mathbb{R} \to y \in \mathbb{R} \to \lvert x \rvert \leq y \to \lvert x \rvert \in \mathbb{R}$
\item $\varphi_{14} :\equiv (x \in \mathbb{R} \to y \in \mathbb{R} \to \lvert x \rvert \leq y \to \lvert x \rvert \leq y) \to x \in \mathbb{R} \to y \in \mathbb{R} \to \lvert x \rvert \leq y \to -y \leq -\lvert x \rvert$
\item $\varphi_{15} :\equiv \forall v_{287}\,(-y \leq -\lvert x \rvert \to -\lvert x \rvert \leq v_{287} \to -y \leq v_{287}) \to -y \leq -\lvert x \rvert \to -\lvert x \rvert \leq x \to -y \leq x$
\item $\varphi_{16} :\equiv (x \in \mathbb{R} \to y \in \mathbb{R} \to \lvert x \rvert \leq y \to -\lvert x \rvert \leq x \land x \leq \lvert x \rvert) \to x \in \mathbb{R} \to y \in \mathbb{R} \to \lvert x \rvert \leq y \to x \leq \lvert x \rvert$
\item $\varphi_{17} :\equiv \varphi_{3} \to \forall v_{287}\,(-y \leq -\lvert x \rvert \to -\lvert x \rvert \leq v_{287} \to -y \leq v_{287})$
\item $\varphi_{18} :\equiv (x \in \mathbb{R} \to y \in \mathbb{R} \to \lvert x \rvert \leq y \to x \in \mathbb{R}) \to x \in \mathbb{R} \to y \in \mathbb{R} \to \lvert x \rvert \leq y \to -\lvert x \rvert \leq x \land x \leq \lvert x \rvert$
\item $\varphi_{19} :\equiv (x \in \mathbb{R} \to y \in \mathbb{R} \to \lvert x \rvert \leq y \to -\lvert x \rvert \leq x \land x \leq \lvert x \rvert) \to x \in \mathbb{R} \to y \in \mathbb{R} \to \lvert x \rvert \leq y \to -\lvert x \rvert \leq x$
\item $\varphi_{20} :\equiv \forall x\,\forall y\,(x \in \mathbb{R} \to y \in \mathbb{R} \to x \leq y \to -y \leq -x) \to \forall y\,(\lvert x \rvert \in \mathbb{R} \to y \in \mathbb{R} \to \lvert x \rvert \leq y \to -y \leq -\lvert x \rvert)$
\item $\varphi_{21} :\equiv \forall y\,(\lvert x \rvert \in \mathbb{R} \to y \in \mathbb{R} \to \lvert x \rvert \leq y \to -y \leq -\lvert x \rvert) \to \lvert x \rvert \in \mathbb{R} \to y \in \mathbb{R} \to \lvert x \rvert \leq y \to -y \leq -\lvert x \rvert$
\item $\varphi_{22} :\equiv (x \leq \lvert x \rvert \to \lvert x \rvert \leq y \to x \leq y) \to (x \in \mathbb{R} \to y \in \mathbb{R} \to \lvert x \rvert \leq y \to x \leq \lvert x \rvert) \to \varphi_{10}$
\item $\varphi_{23} :\equiv (-y \leq -\lvert x \rvert \to -\lvert x \rvert \leq x \to -y \leq x) \to (x \in \mathbb{R} \to y \in \mathbb{R} \to \lvert x \rvert \leq y \to -y \leq -\lvert x \rvert) \to \varphi_{12}$
\item $\varphi_{24} :\equiv (x \in \mathbb{R} \to y \in \mathbb{R} \to \lvert x \rvert \leq y \to \lvert x \rvert \in \mathbb{R}) \to (x \in \mathbb{R} \to y \in \mathbb{R} \to \lvert x \rvert \leq y \to y \in \mathbb{R}) \to \varphi_{14}$
\end{abbrevlist}
\end{fullproof}

% 自動生成 (prover/build.js): LeAbs  ノード数 142
\begin{fullproof}{[A10]}{$\forall x\,\forall y\,(x \in \mathbb{R} \to y \in \mathbb{R} \to -y \leq x \to x \leq y \to \lvert x \rvert \leq y)$}{142}\phantomsection\label{tree:LeAbs}
\noindent$\Pi_{1}$:
\begin{ptree}
\AXC{$\forall x \in \mathbb{R}\,(-x \in \mathbb{R})$\lbl{(R1)}}
\AXC{$\forall x \in \mathbb{R}\,(-x \in \mathbb{R}) \to y \in \mathbb{R} \to -y \in \mathbb{R}$\lbl{A2}}
\BIC{$y \in \mathbb{R} \to -y \in \mathbb{R}$}
\AXC{$(y \in \mathbb{R} \to -y \in \mathbb{R}) \to \varphi_{31}$\lbl{A1}}
\BIC{$\varphi_{31}$}
\end{ptree}
\noindent$\Pi_{2}$:
\begin{ptree}
\AXC{$\forall x\,\forall y\,(x \in \mathbb{R} \to y \in \mathbb{R} \to x \leq y \to -y \leq -x)$\lbl{[F23]}}
\AXC{$\varphi_{24}$\lbl{A2}}
\BIC{$\forall y\,(v_{288} \in \mathbb{R} \to y \in \mathbb{R} \to v_{288} \leq y \to -y \leq -v_{288})$}
\AXC{$\varphi_{29}$\lbl{A2}}
\BIC{$\varphi_{6}$}
\UIC{$\forall v_{289}\,(\varphi_{6})$}
\UIC{$\forall v_{288}\,\forall v_{289}\,(\varphi_{6})$}
\AXC{$\forall v_{288}\,\forall v_{289}\,(\varphi_{6}) \to \varphi_{11}$\lbl{A2}}
\BIC{$\varphi_{11}$}
\end{ptree}
\noindent$\Pi_{3}$:
\begin{ptree}
\AXC{$\stackrel{\Pi_{2}}{\vdots}$}\noLine
\UIC{$\varphi_{11}$}
\AXC{$\varphi_{11} \to -y \in \mathbb{R} \to x \in \mathbb{R} \to -y \leq x \to -x \leq -(-y)$\lbl{A2}}
\BIC{$-y \in \mathbb{R} \to x \in \mathbb{R} \to -y \leq x \to -x \leq -(-y)$}
\AXC{$(-y \in \mathbb{R} \to x \in \mathbb{R} \to -y \leq x \to -x \leq -(-y)) \to \varphi_{40}$\lbl{A1}}
\BIC{$\varphi_{40}$}
\end{ptree}
\noindent$\Pi_{4}$:
\begin{ptree}
\AXC{$x \in \mathbb{R} \to y \in \mathbb{R} \to -y \leq x \to x \leq y \to \neg 0 \leq x \to y \in \mathbb{R}$\lbl{A1}}
\AXC{$\stackrel{\Pi_{1}}{\vdots}$}\noLine
\UIC{$\varphi_{31}$}
\BIC{$x \in \mathbb{R} \to y \in \mathbb{R} \to -y \leq x \to x \leq y \to \neg 0 \leq x \to -y \in \mathbb{R}$}
\AXC{$\stackrel{\Pi_{3}}{\vdots}$}\noLine
\UIC{$\varphi_{40}$}
\BIC{$\varphi_{36}$}
\end{ptree}
\noindent$\Pi_{5}$:
\begin{ptree}
\AXC{$x \in \mathbb{R} \to y \in \mathbb{R} \to -y \leq x \to x \leq y \to -y \leq x$\lbl{A1}}
\AXC{$x \in \mathbb{R} \to y \in \mathbb{R} \to -y \leq x \to x \leq y \to \neg 0 \leq x \to x \in \mathbb{R}$\lbl{A1}}
\AXC{$\stackrel{\Pi_{4}}{\vdots}$}\noLine
\UIC{$\varphi_{36}$}
\BIC{$(x \in \mathbb{R} \to y \in \mathbb{R} \to -y \leq x \to x \leq y \to -y \leq x) \to \varphi_{10}$}
\BIC{$\varphi_{10}$}
\end{ptree}
\noindent$\Pi_{6}$:
\begin{ptree}
\AXC{$\forall x \in \mathbb{R}\,(-(-x) = x)$\lbl{[F5]}}
\AXC{$\forall x \in \mathbb{R}\,(-(-x) = x) \to y \in \mathbb{R} \to -(-y) = y$\lbl{A2}}
\BIC{$y \in \mathbb{R} \to -(-y) = y$}
\AXC{$(y \in \mathbb{R} \to -(-y) = y) \to \varphi_{30}$\lbl{A1}}
\BIC{$\varphi_{30}$}
\end{ptree}
\noindent$\Pi_{7}$:
\begin{ptree}
\AXC{$x \in \mathbb{R} \to y \in \mathbb{R} \to -y \leq x \to x \leq y \to \neg 0 \leq x \to y \in \mathbb{R}$\lbl{A1}}
\AXC{$\stackrel{\Pi_{6}}{\vdots}$}\noLine
\UIC{$\varphi_{30}$}
\BIC{$\varphi_{7}$}
\AXC{$\varphi_{7} \to \varphi_{10} \to (-(-y) = y \to -x < -(-y) \to -x < y) \to \varphi_{22}$\lbl{A1}}
\BIC{$\varphi_{10} \to (-(-y) = y \to -x < -(-y) \to -x < y) \to \varphi_{22}$}
\end{ptree}
\noindent$\Pi_{8}$:
\begin{ptree}
\AXC{$z = w \to -x < z \to -x < w$\lbl{A6}}
\UIC{$\forall w\,(z = w \to -x < z \to -x < w)$}
\UIC{$\forall z\,\forall w\,(z = w \to -x < z \to -x < w)$}
\AXC{$\varphi_{18}$\lbl{A2}}
\BIC{$\forall w\,(-(-y) = w \to -x < -(-y) \to -x < w)$}
\AXC{$\varphi_{23}$\lbl{A2}}
\BIC{$-(-y) = y \to -x < -(-y) \to -x < y$}
\AXC{$\stackrel{\Pi_{5}}{\vdots}$}\noLine
\UIC{$\varphi_{10}$}
\AXC{$\stackrel{\Pi_{7}}{\vdots}$}\noLine
\UIC{$\varphi_{10} \to (-(-y) = y \to -x < -(-y) \to -x < y) \to \varphi_{22}$}
\BIC{$(-(-y) = y \to -x < -(-y) \to -x < y) \to \varphi_{22}$}
\BIC{$\varphi_{22}$}
\end{ptree}
\noindent$\Pi_{9}$:
\begin{ptree}
\AXC{$x = x$\lbl{A5}}
\AXC{$x = y \to x = x \to y = x$\lbl{A6}}
\AXC{$(x = y \to x = x \to y = x) \to x = x \to x = y \to y = x$\lbl{A1}}
\BIC{$x = x \to x = y \to y = x$}
\BIC{$x = y \to y = x$}
\UIC{$\forall y\,(x = y \to y = x)$}
\UIC{$\forall x\,\forall y\,(x = y \to y = x)$}
\end{ptree}
\noindent$\Pi_{10}$:
\begin{ptree}
\AXC{$\stackrel{\Pi_{9}}{\vdots}$}\noLine
\UIC{$\forall x\,\forall y\,(x = y \to y = x)$}
\AXC{$\forall x\,\forall y\,(x = y \to y = x) \to \forall y\,(\lvert x \rvert = y \to y = \lvert x \rvert)$\lbl{A2}}
\BIC{$\forall y\,(\lvert x \rvert = y \to y = \lvert x \rvert)$}
\end{ptree}
\noindent$\Pi_{11}$:
\begin{ptree}
\AXC{$\forall x\,\forall y\,(x \in \mathbb{R} \to y \in \mathbb{R} \to \neg y \leq x \to x < y)$\lbl{[A2]}}
\AXC{$\varphi_{17}$\lbl{A2}}
\BIC{$\forall y\,(x \in \mathbb{R} \to y \in \mathbb{R} \to \neg y \leq x \to x < y)$}
\AXC{$\varphi_{15}$\lbl{A2}}
\BIC{$x \in \mathbb{R} \to 0 \in \mathbb{R} \to \neg 0 \leq x \to x < 0$}
\AXC{$(x \in \mathbb{R} \to 0 \in \mathbb{R} \to \neg 0 \leq x \to x < 0) \to \varphi_{37}$\lbl{A1}}
\BIC{$\varphi_{37}$}
\end{ptree}
\noindent$\Pi_{12}$:
\begin{ptree}
\AXC{$0 \in \mathbb{R}$\lbl{(R1)}}
\AXC{$x \in \mathbb{R} \to y \in \mathbb{R} \to -y \leq x \to x \leq y \to \neg 0 \leq x \to x \in \mathbb{R}$\lbl{A1}}
\AXC{$\stackrel{\Pi_{11}}{\vdots}$}\noLine
\UIC{$\varphi_{37}$}
\BIC{$0 \in \mathbb{R} \to \varphi_{27}$}
\BIC{$\varphi_{27}$}
\end{ptree}
\noindent$\Pi_{13}$:
\begin{ptree}
\AXC{$\forall x\,(x < 0 \to \lvert x \rvert = -x)$\lbl{[A3]}}
\AXC{$\forall x\,(x < 0 \to \lvert x \rvert = -x) \to x < 0 \to \lvert x \rvert = -x$\lbl{A2}}
\BIC{$x < 0 \to \lvert x \rvert = -x$}
\AXC{$(x < 0 \to \lvert x \rvert = -x) \to \varphi_{28}$\lbl{A1}}
\BIC{$\varphi_{28}$}
\end{ptree}
\noindent$\Pi_{14}$:
\begin{ptree}
\AXC{$x \in \mathbb{R} \to y \in \mathbb{R} \to -y \leq x \to x \leq y \to \neg 0 \leq x \to \neg 0 \leq x$\lbl{A1}}
\AXC{$\stackrel{\Pi_{12}}{\vdots}$}\noLine
\UIC{$\varphi_{27}$}
\BIC{$x \in \mathbb{R} \to y \in \mathbb{R} \to -y \leq x \to x \leq y \to \neg 0 \leq x \to x < 0$}
\AXC{$\stackrel{\Pi_{13}}{\vdots}$}\noLine
\UIC{$\varphi_{28}$}
\BIC{$x \in \mathbb{R} \to y \in \mathbb{R} \to -y \leq x \to x \leq y \to \neg 0 \leq x \to \lvert x \rvert = -x$}
\AXC{$\varphi_{33}$\lbl{A1}}
\BIC{$\varphi_{19}$}
\end{ptree}
\noindent$\Pi_{15}$:
\begin{ptree}
\AXC{$\stackrel{\Pi_{10}}{\vdots}$}\noLine
\UIC{$\forall y\,(\lvert x \rvert = y \to y = \lvert x \rvert)$}
\AXC{$\forall y\,(\lvert x \rvert = y \to y = \lvert x \rvert) \to \lvert x \rvert = -x \to -x = \lvert x \rvert$\lbl{A2}}
\BIC{$\lvert x \rvert = -x \to -x = \lvert x \rvert$}
\AXC{$\stackrel{\Pi_{14}}{\vdots}$}\noLine
\UIC{$\varphi_{19}$}
\BIC{$x \in \mathbb{R} \to y \in \mathbb{R} \to -y \leq x \to x \leq y \to \neg 0 \leq x \to -x = \lvert x \rvert$}
\AXC{$\varphi_{41}$\lbl{A1}}
\BIC{$\varphi_{38}$}
\end{ptree}
\noindent$\Pi_{16}$:
\begin{ptree}
\AXC{$z = w \to -x = z \to -x = w$\lbl{A6}}
\UIC{$\forall w\,(z = w \to -x = z \to -x = w)$}
\UIC{$\forall z\,\forall w\,(z = w \to -x = z \to -x = w)$}
\AXC{$\varphi_{14}$\lbl{A2}}
\BIC{$\forall w\,(-(-y) = w \to -x = -(-y) \to -x = w)$}
\AXC{$\varphi_{20}$\lbl{A2}}
\BIC{$-(-y) = y \to -x = -(-y) \to -x = y$}
\AXC{$\stackrel{\Pi_{8}}{\vdots}$}\noLine
\UIC{$\varphi_{22}$}
\BIC{$x \in \mathbb{R} \to y \in \mathbb{R} \to -y \leq x \to x \leq y \to \neg 0 \leq x \to -x \leq y$}
\AXC{$\stackrel{\Pi_{15}}{\vdots}$}\noLine
\UIC{$\varphi_{38}$}
\BIC{$(-x = \lvert x \rvert \to -x < y \to \lvert x \rvert < y) \to \varphi_{21}$}
\end{ptree}
\noindent$\Pi_{17}$:
\begin{ptree}
\AXC{$z = w \to z < y \to w < y$\lbl{A6}}
\UIC{$\forall w\,(z = w \to z < y \to w < y)$}
\UIC{$\forall z\,\forall w\,(z = w \to z < y \to w < y)$}
\AXC{$\varphi_{8}$\lbl{A2}}
\BIC{$\forall w\,(-x = w \to -x < y \to w < y)$}
\AXC{$\varphi_{12}$\lbl{A2}}
\BIC{$-x = \lvert x \rvert \to -x < y \to \lvert x \rvert < y$}
\AXC{$\stackrel{\Pi_{16}}{\vdots}$}\noLine
\UIC{$(-x = \lvert x \rvert \to -x < y \to \lvert x \rvert < y) \to \varphi_{21}$}
\BIC{$\varphi_{21}$}
\end{ptree}
\noindent$\Pi_{18}$:
\begin{ptree}
\AXC{$\stackrel{\Pi_{9}}{\vdots}$}\noLine
\UIC{$\forall x\,\forall y\,(x = y \to y = x)$}
\AXC{$\forall x\,\forall y\,(x = y \to y = x) \to \forall y\,(\lvert x \rvert = y \to y = \lvert x \rvert)$\lbl{A2}}
\BIC{$\forall y\,(\lvert x \rvert = y \to y = \lvert x \rvert)$}
\end{ptree}
\noindent$\Pi_{19}$:
\begin{ptree}
\AXC{$\stackrel{\Pi_{18}}{\vdots}$}\noLine
\UIC{$\forall y\,(\lvert x \rvert = y \to y = \lvert x \rvert)$}
\AXC{$\forall y\,(\lvert x \rvert = y \to y = \lvert x \rvert) \to \lvert x \rvert = x \to x = \lvert x \rvert$\lbl{A2}}
\BIC{$\lvert x \rvert = x \to x = \lvert x \rvert$}
\end{ptree}
\noindent$\Pi_{20}$:
\begin{ptree}
\AXC{$\forall x\,(0 \leq x \to \lvert x \rvert = x)$\lbl{[A1]}}
\AXC{$\forall x\,(0 \leq x \to \lvert x \rvert = x) \to 0 \leq x \to \lvert x \rvert = x$\lbl{A2}}
\BIC{$0 \leq x \to \lvert x \rvert = x$}
\AXC{$(0 \leq x \to \lvert x \rvert = x) \to \varphi_{25}$\lbl{A1}}
\BIC{$\varphi_{25}$}
\end{ptree}
\noindent$\Pi_{21}$:
\begin{ptree}
\AXC{$\stackrel{\Pi_{19}}{\vdots}$}\noLine
\UIC{$\lvert x \rvert = x \to x = \lvert x \rvert$}
\AXC{$x \in \mathbb{R} \to y \in \mathbb{R} \to -y \leq x \to x \leq y \to 0 \leq x \to 0 \leq x$\lbl{A1}}
\AXC{$\stackrel{\Pi_{20}}{\vdots}$}\noLine
\UIC{$\varphi_{25}$}
\BIC{$x \in \mathbb{R} \to y \in \mathbb{R} \to -y \leq x \to x \leq y \to 0 \leq x \to \lvert x \rvert = x$}
\AXC{$\varphi_{32}$\lbl{A1}}
\BIC{$\varphi_{13}$}
\BIC{$x \in \mathbb{R} \to y \in \mathbb{R} \to -y \leq x \to x \leq y \to 0 \leq x \to x = \lvert x \rvert$}
\AXC{$\varphi_{39}$\lbl{A1}}
\BIC{$\varphi_{34}$}
\end{ptree}
\noindent$\Pi_{22}$:
\begin{ptree}
\AXC{$z = w \to z < y \to w < y$\lbl{A6}}
\UIC{$\forall w\,(z = w \to z < y \to w < y)$}
\UIC{$\forall z\,\forall w\,(z = w \to z < y \to w < y)$}
\AXC{$\varphi_{4}$\lbl{A2}}
\BIC{$\forall w\,(x = w \to x < y \to w < y)$}
\AXC{$\varphi_{5}$\lbl{A2}}
\BIC{$x = \lvert x \rvert \to x < y \to \lvert x \rvert < y$}
\AXC{$x \in \mathbb{R} \to y \in \mathbb{R} \to -y \leq x \to x \leq y \to x \leq y$\lbl{A1}}
\AXC{$\stackrel{\Pi_{21}}{\vdots}$}\noLine
\UIC{$\varphi_{34}$}
\BIC{$(x = \lvert x \rvert \to x < y \to \lvert x \rvert < y) \to \varphi_{16}$}
\BIC{$\varphi_{16}$}
\end{ptree}
\noindent$\Pi_{23}$:
\begin{ptree}
\AXC{$z = w \to z = y \to w = y$\lbl{A6}}
\UIC{$\forall w\,(z = w \to z = y \to w = y)$}
\UIC{$\forall z\,\forall w\,(z = w \to z = y \to w = y)$}
\AXC{$\varphi_{1}$\lbl{A2}}
\BIC{$\forall w\,(x = w \to x = y \to w = y)$}
\AXC{$\varphi_{2}$\lbl{A2}}
\BIC{$x = \lvert x \rvert \to x = y \to \lvert x \rvert = y$}
\AXC{$\stackrel{\Pi_{22}}{\vdots}$}\noLine
\UIC{$\varphi_{16}$}
\BIC{$x \in \mathbb{R} \to y \in \mathbb{R} \to -y \leq x \to x \leq y \to 0 \leq x \to \lvert x \rvert \leq y$}
\AXC{$\varphi_{35}$\lbl{A1}}
\BIC{$\varphi_{26}$}
\end{ptree}
\noindent 主木:
\begin{ptree}
\AXC{$z = w \to z = y \to w = y$\lbl{A6}}
\UIC{$\forall w\,(z = w \to z = y \to w = y)$}
\UIC{$\forall z\,\forall w\,(z = w \to z = y \to w = y)$}
\AXC{$\varphi_{3}$\lbl{A2}}
\BIC{$\forall w\,(-x = w \to -x = y \to w = y)$}
\AXC{$\varphi_{9}$\lbl{A2}}
\BIC{$-x = \lvert x \rvert \to -x = y \to \lvert x \rvert = y$}
\AXC{$\stackrel{\Pi_{17}}{\vdots}$}\noLine
\UIC{$\varphi_{21}$}
\BIC{$x \in \mathbb{R} \to y \in \mathbb{R} \to -y \leq x \to x \leq y \to \neg 0 \leq x \to \lvert x \rvert \leq y$}
\AXC{$\stackrel{\Pi_{23}}{\vdots}$}\noLine
\UIC{$\varphi_{26}$}
\BIC{$x \in \mathbb{R} \to y \in \mathbb{R} \to -y \leq x \to x \leq y \to \lvert x \rvert \leq y$}
\UIC{$\forall y\,(x \in \mathbb{R} \to y \in \mathbb{R} \to -y \leq x \to x \leq y \to \lvert x \rvert \leq y)$}
\UIC{$\forall x\,\forall y\,(x \in \mathbb{R} \to y \in \mathbb{R} \to -y \leq x \to x \leq y \to \lvert x \rvert \leq y)$}
\end{ptree}
\begin{abbrevlist}
\item $\varphi_{1} :\equiv \forall z\,\forall w\,(z = w \to z = y \to w = y) \to \forall w\,(x = w \to x = y \to w = y)$
\item $\varphi_{2} :\equiv \forall w\,(x = w \to x = y \to w = y) \to x = \lvert x \rvert \to x = y \to \lvert x \rvert = y$
\item $\varphi_{3} :\equiv \forall z\,\forall w\,(z = w \to z = y \to w = y) \to \forall w\,(-x = w \to -x = y \to w = y)$
\item $\varphi_{4} :\equiv \forall z\,\forall w\,(z = w \to z < y \to w < y) \to \forall w\,(x = w \to x < y \to w < y)$
\item $\varphi_{5} :\equiv \forall w\,(x = w \to x < y \to w < y) \to x = \lvert x \rvert \to x < y \to \lvert x \rvert < y$
\item $\varphi_{6} :\equiv v_{288} \in \mathbb{R} \to v_{289} \in \mathbb{R} \to v_{288} \leq v_{289} \to -v_{289} \leq -v_{288}$
\item $\varphi_{7} :\equiv x \in \mathbb{R} \to y \in \mathbb{R} \to -y \leq x \to x \leq y \to \neg 0 \leq x \to -(-y) = y$
\item $\varphi_{8} :\equiv \forall z\,\forall w\,(z = w \to z < y \to w < y) \to \forall w\,(-x = w \to -x < y \to w < y)$
\item $\varphi_{9} :\equiv \forall w\,(-x = w \to -x = y \to w = y) \to -x = \lvert x \rvert \to -x = y \to \lvert x \rvert = y$
\item $\varphi_{10} :\equiv x \in \mathbb{R} \to y \in \mathbb{R} \to -y \leq x \to x \leq y \to \neg 0 \leq x \to -x \leq -(-y)$
\item $\varphi_{11} :\equiv \forall v_{289}\,(-y \in \mathbb{R} \to v_{289} \in \mathbb{R} \to -y \leq v_{289} \to -v_{289} \leq -(-y))$
\item $\varphi_{12} :\equiv \forall w\,(-x = w \to -x < y \to w < y) \to -x = \lvert x \rvert \to -x < y \to \lvert x \rvert < y$
\item $\varphi_{13} :\equiv (\lvert x \rvert = x \to x = \lvert x \rvert) \to x \in \mathbb{R} \to y \in \mathbb{R} \to -y \leq x \to x \leq y \to 0 \leq x \to x = \lvert x \rvert$
\item $\varphi_{14} :\equiv \forall z\,\forall w\,(z = w \to -x = z \to -x = w) \to \forall w\,(-(-y) = w \to -x = -(-y) \to -x = w)$
\item $\varphi_{15} :\equiv \forall y\,(x \in \mathbb{R} \to y \in \mathbb{R} \to \neg y \leq x \to x < y) \to x \in \mathbb{R} \to 0 \in \mathbb{R} \to \neg 0 \leq x \to x < 0$
\item $\varphi_{16} :\equiv (x = \lvert x \rvert \to x = y \to \lvert x \rvert = y) \to x \in \mathbb{R} \to y \in \mathbb{R} \to -y \leq x \to x \leq y \to 0 \leq x \to \lvert x \rvert \leq y$
\item $\varphi_{17} :\equiv \forall x\,\forall y\,(x \in \mathbb{R} \to y \in \mathbb{R} \to \neg y \leq x \to x < y) \to \forall y\,(x \in \mathbb{R} \to y \in \mathbb{R} \to \neg y \leq x \to x < y)$
\item $\varphi_{18} :\equiv \forall z\,\forall w\,(z = w \to -x < z \to -x < w) \to \forall w\,(-(-y) = w \to -x < -(-y) \to -x < w)$
\item $\varphi_{19} :\equiv (\lvert x \rvert = -x \to -x = \lvert x \rvert) \to x \in \mathbb{R} \to y \in \mathbb{R} \to -y \leq x \to x \leq y \to \neg 0 \leq x \to -x = \lvert x \rvert$
\item $\varphi_{20} :\equiv \forall w\,(-(-y) = w \to -x = -(-y) \to -x = w) \to -(-y) = y \to -x = -(-y) \to -x = y$
\item $\varphi_{21} :\equiv (-x = \lvert x \rvert \to -x = y \to \lvert x \rvert = y) \to x \in \mathbb{R} \to y \in \mathbb{R} \to -y \leq x \to x \leq y \to \neg 0 \leq x \to \lvert x \rvert \leq y$
\item $\varphi_{22} :\equiv (-(-y) = y \to -x = -(-y) \to -x = y) \to x \in \mathbb{R} \to y \in \mathbb{R} \to -y \leq x \to x \leq y \to \neg 0 \leq x \to -x \leq y$
\item $\varphi_{23} :\equiv \forall w\,(-(-y) = w \to -x < -(-y) \to -x < w) \to -(-y) = y \to -x < -(-y) \to -x < y$
\item $\varphi_{24} :\equiv \forall x\,\forall y\,(x \in \mathbb{R} \to y \in \mathbb{R} \to x \leq y \to -y \leq -x) \to \forall y\,(v_{288} \in \mathbb{R} \to y \in \mathbb{R} \to v_{288} \leq y \to -y \leq -v_{288})$
\item $\varphi_{25} :\equiv (x \in \mathbb{R} \to y \in \mathbb{R} \to -y \leq x \to x \leq y \to 0 \leq x \to 0 \leq x) \to x \in \mathbb{R} \to y \in \mathbb{R} \to -y \leq x \to x \leq y \to 0 \leq x \to \lvert x \rvert = x$
\item $\varphi_{26} :\equiv (x \in \mathbb{R} \to y \in \mathbb{R} \to -y \leq x \to x \leq y \to \neg 0 \leq x \to \lvert x \rvert \leq y) \to x \in \mathbb{R} \to y \in \mathbb{R} \to -y \leq x \to x \leq y \to \lvert x \rvert \leq y$
\item $\varphi_{27} :\equiv (x \in \mathbb{R} \to y \in \mathbb{R} \to -y \leq x \to x \leq y \to \neg 0 \leq x \to \neg 0 \leq x) \to x \in \mathbb{R} \to y \in \mathbb{R} \to -y \leq x \to x \leq y \to \neg 0 \leq x \to x < 0$
\item $\varphi_{28} :\equiv (x \in \mathbb{R} \to y \in \mathbb{R} \to -y \leq x \to x \leq y \to \neg 0 \leq x \to x < 0) \to x \in \mathbb{R} \to y \in \mathbb{R} \to -y \leq x \to x \leq y \to \neg 0 \leq x \to \lvert x \rvert = -x$
\item $\varphi_{29} :\equiv \forall y\,(v_{288} \in \mathbb{R} \to y \in \mathbb{R} \to v_{288} \leq y \to -y \leq -v_{288}) \to \varphi_{6}$
\item $\varphi_{30} :\equiv (x \in \mathbb{R} \to y \in \mathbb{R} \to -y \leq x \to x \leq y \to \neg 0 \leq x \to y \in \mathbb{R}) \to \varphi_{7}$
\item $\varphi_{31} :\equiv (x \in \mathbb{R} \to y \in \mathbb{R} \to -y \leq x \to x \leq y \to \neg 0 \leq x \to y \in \mathbb{R}) \to x \in \mathbb{R} \to y \in \mathbb{R} \to -y \leq x \to x \leq y \to \neg 0 \leq x \to -y \in \mathbb{R}$
\item $\varphi_{32} :\equiv (x \in \mathbb{R} \to y \in \mathbb{R} \to -y \leq x \to x \leq y \to 0 \leq x \to \lvert x \rvert = x) \to \varphi_{13}$
\item $\varphi_{33} :\equiv (x \in \mathbb{R} \to y \in \mathbb{R} \to -y \leq x \to x \leq y \to \neg 0 \leq x \to \lvert x \rvert = -x) \to \varphi_{19}$
\item $\varphi_{34} :\equiv (x \in \mathbb{R} \to y \in \mathbb{R} \to -y \leq x \to x \leq y \to x \leq y) \to (x = \lvert x \rvert \to x < y \to \lvert x \rvert < y) \to \varphi_{16}$
\item $\varphi_{35} :\equiv (x \in \mathbb{R} \to y \in \mathbb{R} \to -y \leq x \to x \leq y \to 0 \leq x \to \lvert x \rvert \leq y) \to \varphi_{26}$
\item $\varphi_{36} :\equiv (x \in \mathbb{R} \to y \in \mathbb{R} \to -y \leq x \to x \leq y \to \neg 0 \leq x \to x \in \mathbb{R}) \to (x \in \mathbb{R} \to y \in \mathbb{R} \to -y \leq x \to x \leq y \to -y \leq x) \to \varphi_{10}$
\item $\varphi_{37} :\equiv (x \in \mathbb{R} \to y \in \mathbb{R} \to -y \leq x \to x \leq y \to \neg 0 \leq x \to x \in \mathbb{R}) \to 0 \in \mathbb{R} \to \varphi_{27}$
\item $\varphi_{38} :\equiv (x \in \mathbb{R} \to y \in \mathbb{R} \to -y \leq x \to x \leq y \to \neg 0 \leq x \to -x \leq y) \to (-x = \lvert x \rvert \to -x < y \to \lvert x \rvert < y) \to \varphi_{21}$
\item $\varphi_{39} :\equiv (x \in \mathbb{R} \to y \in \mathbb{R} \to -y \leq x \to x \leq y \to 0 \leq x \to x = \lvert x \rvert) \to \varphi_{34}$
\item $\varphi_{40} :\equiv (x \in \mathbb{R} \to y \in \mathbb{R} \to -y \leq x \to x \leq y \to \neg 0 \leq x \to -y \in \mathbb{R}) \to \varphi_{36}$
\item $\varphi_{41} :\equiv (x \in \mathbb{R} \to y \in \mathbb{R} \to -y \leq x \to x \leq y \to \neg 0 \leq x \to -x = \lvert x \rvert) \to \varphi_{38}$
\end{abbrevlist}
\end{fullproof}

% 自動生成 (prover/build.js): AbsLt  ノード数 66
\begin{fullproof}{[A11]}{$\forall x\,\forall y\,(x \in \mathbb{R} \to y \in \mathbb{R} \to \lvert x \rvert < y \to -y < x \land x < y)$}{66}\phantomsection\label{tree:AbsLt}
\noindent$\Pi_{1}$:
\begin{ptree}
\AXC{$\forall x \in \mathbb{R}\,(-\lvert x \rvert \leq x \land x \leq \lvert x \rvert)$\lbl{[A8]}}
\AXC{$\varphi_{9}$\lbl{A2}}
\BIC{$x \in \mathbb{R} \to -\lvert x \rvert \leq x \land x \leq \lvert x \rvert$}
\AXC{$(x \in \mathbb{R} \to -\lvert x \rvert \leq x \land x \leq \lvert x \rvert) \to \varphi_{16}$\lbl{A1}}
\BIC{$\varphi_{16}$}
\end{ptree}
\noindent$\Pi_{2}$:
\begin{ptree}
\AXC{$x \in \mathbb{R} \to y \in \mathbb{R} \to \lvert x \rvert < y \to x \in \mathbb{R}$\lbl{A1}}
\AXC{$\stackrel{\Pi_{1}}{\vdots}$}\noLine
\UIC{$\varphi_{16}$}
\BIC{$x \in \mathbb{R} \to y \in \mathbb{R} \to \lvert x \rvert < y \to -\lvert x \rvert \leq x \land x \leq \lvert x \rvert$}
\end{ptree}
\noindent$\Pi_{3}$:
\begin{ptree}
\AXC{$\forall x\,\forall y\,\forall z\,(x \leq y \to y < z \to x < z)$\lbl{[O4]}}
\AXC{$\varphi_{3}$\lbl{A2}}
\BIC{$\forall y\,\forall z\,(x \leq y \to y < z \to x < z)$}
\AXC{$\varphi_{4}$\lbl{A2}}
\BIC{$\forall z\,(x \leq \lvert x \rvert \to \lvert x \rvert < z \to x < z)$}
\AXC{$\varphi_{6}$\lbl{A2}}
\BIC{$x \leq \lvert x \rvert \to \lvert x \rvert < y \to x < y$}
\AXC{$\varphi_{22}$\lbl{A1}}
\BIC{$(x \in \mathbb{R} \to y \in \mathbb{R} \to \lvert x \rvert < y \to x \leq \lvert x \rvert) \to \varphi_{10}$}
\end{ptree}
\noindent$\Pi_{4}$:
\begin{ptree}
\AXC{$\stackrel{\Pi_{2}}{\vdots}$}\noLine
\UIC{$x \in \mathbb{R} \to y \in \mathbb{R} \to \lvert x \rvert < y \to -\lvert x \rvert \leq x \land x \leq \lvert x \rvert$}
\AXC{$\varphi_{14}$\lbl{A1}}
\BIC{$x \in \mathbb{R} \to y \in \mathbb{R} \to \lvert x \rvert < y \to x \leq \lvert x \rvert$}
\AXC{$\stackrel{\Pi_{3}}{\vdots}$}\noLine
\UIC{$(x \in \mathbb{R} \to y \in \mathbb{R} \to \lvert x \rvert < y \to x \leq \lvert x \rvert) \to \varphi_{10}$}
\BIC{$\varphi_{10}$}
\end{ptree}
\noindent$\Pi_{5}$:
\begin{ptree}
\AXC{$\forall x \in \mathbb{R}\,(\lvert x \rvert \in \mathbb{R})$\lbl{[A5]}}
\AXC{$\forall x \in \mathbb{R}\,(\lvert x \rvert \in \mathbb{R}) \to x \in \mathbb{R} \to \lvert x \rvert \in \mathbb{R}$\lbl{A2}}
\BIC{$x \in \mathbb{R} \to \lvert x \rvert \in \mathbb{R}$}
\AXC{$(x \in \mathbb{R} \to \lvert x \rvert \in \mathbb{R}) \to \varphi_{12}$\lbl{A1}}
\BIC{$\varphi_{12}$}
\end{ptree}
\noindent$\Pi_{6}$:
\begin{ptree}
\AXC{$\forall x\,\forall y\,(x \in \mathbb{R} \to y \in \mathbb{R} \to x < y \to -y < -x)$\lbl{[F22]}}
\AXC{$\varphi_{20}$\lbl{A2}}
\BIC{$\forall y\,(\lvert x \rvert \in \mathbb{R} \to y \in \mathbb{R} \to \lvert x \rvert < y \to -y < -\lvert x \rvert)$}
\AXC{$\varphi_{21}$\lbl{A2}}
\BIC{$\lvert x \rvert \in \mathbb{R} \to y \in \mathbb{R} \to \lvert x \rvert < y \to -y < -\lvert x \rvert$}
\end{ptree}
\noindent$\Pi_{7}$:
\begin{ptree}
\AXC{$\stackrel{\Pi_{6}}{\vdots}$}\noLine
\UIC{$\lvert x \rvert \in \mathbb{R} \to y \in \mathbb{R} \to \lvert x \rvert < y \to -y < -\lvert x \rvert$}
\AXC{$(\lvert x \rvert \in \mathbb{R} \to y \in \mathbb{R} \to \lvert x \rvert < y \to -y < -\lvert x \rvert) \to \varphi_{24}$\lbl{A1}}
\BIC{$\varphi_{24}$}
\end{ptree}
\noindent$\Pi_{8}$:
\begin{ptree}
\AXC{$x \in \mathbb{R} \to y \in \mathbb{R} \to \lvert x \rvert < y \to y \in \mathbb{R}$\lbl{A1}}
\AXC{$x \in \mathbb{R} \to y \in \mathbb{R} \to \lvert x \rvert < y \to x \in \mathbb{R}$\lbl{A1}}
\AXC{$\stackrel{\Pi_{5}}{\vdots}$}\noLine
\UIC{$\varphi_{12}$}
\BIC{$x \in \mathbb{R} \to y \in \mathbb{R} \to \lvert x \rvert < y \to \lvert x \rvert \in \mathbb{R}$}
\AXC{$\stackrel{\Pi_{7}}{\vdots}$}\noLine
\UIC{$\varphi_{24}$}
\BIC{$(x \in \mathbb{R} \to y \in \mathbb{R} \to \lvert x \rvert < y \to y \in \mathbb{R}) \to \varphi_{15}$}
\BIC{$\varphi_{15}$}
\end{ptree}
\noindent$\Pi_{9}$:
\begin{ptree}
\AXC{$\forall x\,\forall y\,\forall z\,(x < y \to y \leq z \to x < z)$\lbl{[O3]}}
\AXC{$\varphi_{5}$\lbl{A2}}
\BIC{$\forall y\,\forall z\,(v_{290} < y \to y \leq z \to v_{290} < z)$}
\AXC{$\varphi_{7}$\lbl{A2}}
\BIC{$\forall z\,(v_{290} < v_{291} \to v_{291} \leq z \to v_{290} < z)$}
\AXC{$\varphi_{8}$\lbl{A2}}
\BIC{$v_{290} < v_{291} \to v_{291} \leq v_{292} \to v_{290} < v_{292}$}
\UIC{$\forall v_{292}\,(v_{290} < v_{291} \to v_{291} \leq v_{292} \to v_{290} < v_{292})$}
\UIC{$\varphi_{1}$}
\UIC{$\forall v_{290}\,(\varphi_{1})$}
\AXC{$\forall v_{290}\,(\varphi_{1}) \to \varphi_{2}$\lbl{A2}}
\BIC{$\varphi_{2}$}
\AXC{$\varphi_{19}$\lbl{A2}}
\BIC{$\forall v_{292}\,(-y < -\lvert x \rvert \to -\lvert x \rvert \leq v_{292} \to -y < v_{292})$}
\end{ptree}
\noindent$\Pi_{10}$:
\begin{ptree}
\AXC{$\stackrel{\Pi_{9}}{\vdots}$}\noLine
\UIC{$\forall v_{292}\,(-y < -\lvert x \rvert \to -\lvert x \rvert \leq v_{292} \to -y < v_{292})$}
\AXC{$\varphi_{17}$\lbl{A2}}
\BIC{$-y < -\lvert x \rvert \to -\lvert x \rvert \leq x \to -y < x$}
\AXC{$\varphi_{23}$\lbl{A1}}
\BIC{$(x \in \mathbb{R} \to y \in \mathbb{R} \to \lvert x \rvert < y \to -y < -\lvert x \rvert) \to \varphi_{11}$}
\end{ptree}
\noindent$\Pi_{11}$:
\begin{ptree}
\AXC{$x \in \mathbb{R} \to y \in \mathbb{R} \to \lvert x \rvert < y \to \lvert x \rvert < y$\lbl{A1}}
\AXC{$\stackrel{\Pi_{8}}{\vdots}$}\noLine
\UIC{$\varphi_{15}$}
\BIC{$x \in \mathbb{R} \to y \in \mathbb{R} \to \lvert x \rvert < y \to -y < -\lvert x \rvert$}
\AXC{$\stackrel{\Pi_{10}}{\vdots}$}\noLine
\UIC{$(x \in \mathbb{R} \to y \in \mathbb{R} \to \lvert x \rvert < y \to -y < -\lvert x \rvert) \to \varphi_{11}$}
\BIC{$\varphi_{11}$}
\end{ptree}
\noindent$\Pi_{12}$:
\begin{ptree}
\AXC{$\stackrel{\Pi_{2}}{\vdots}$}\noLine
\UIC{$x \in \mathbb{R} \to y \in \mathbb{R} \to \lvert x \rvert < y \to -\lvert x \rvert \leq x \land x \leq \lvert x \rvert$}
\AXC{$\varphi_{18}$\lbl{A1}}
\BIC{$x \in \mathbb{R} \to y \in \mathbb{R} \to \lvert x \rvert < y \to -\lvert x \rvert \leq x$}
\AXC{$\stackrel{\Pi_{11}}{\vdots}$}\noLine
\UIC{$\varphi_{11}$}
\BIC{$x \in \mathbb{R} \to y \in \mathbb{R} \to \lvert x \rvert < y \to -y < x$}
\AXC{$(x \in \mathbb{R} \to y \in \mathbb{R} \to \lvert x \rvert < y \to -y < x) \to \varphi_{13}$\lbl{A1}}
\BIC{$\varphi_{13}$}
\end{ptree}
\noindent 主木:
\begin{ptree}
\AXC{$x \in \mathbb{R} \to y \in \mathbb{R} \to \lvert x \rvert < y \to \lvert x \rvert < y$\lbl{A1}}
\AXC{$\stackrel{\Pi_{4}}{\vdots}$}\noLine
\UIC{$\varphi_{10}$}
\BIC{$x \in \mathbb{R} \to y \in \mathbb{R} \to \lvert x \rvert < y \to x < y$}
\AXC{$\stackrel{\Pi_{12}}{\vdots}$}\noLine
\UIC{$\varphi_{13}$}
\BIC{$x \in \mathbb{R} \to y \in \mathbb{R} \to \lvert x \rvert < y \to -y < x \land x < y$}
\UIC{$\forall y\,(x \in \mathbb{R} \to y \in \mathbb{R} \to \lvert x \rvert < y \to -y < x \land x < y)$}
\UIC{$\forall x\,\forall y\,(x \in \mathbb{R} \to y \in \mathbb{R} \to \lvert x \rvert < y \to -y < x \land x < y)$}
\end{ptree}
\begin{abbrevlist}
\item $\varphi_{1} :\equiv \forall v_{291}\,\forall v_{292}\,(v_{290} < v_{291} \to v_{291} \leq v_{292} \to v_{290} < v_{292})$
\item $\varphi_{2} :\equiv \forall v_{291}\,\forall v_{292}\,(-y < v_{291} \to v_{291} \leq v_{292} \to -y < v_{292})$
\item $\varphi_{3} :\equiv \forall x\,\forall y\,\forall z\,(x \leq y \to y < z \to x < z) \to \forall y\,\forall z\,(x \leq y \to y < z \to x < z)$
\item $\varphi_{4} :\equiv \forall y\,\forall z\,(x \leq y \to y < z \to x < z) \to \forall z\,(x \leq \lvert x \rvert \to \lvert x \rvert < z \to x < z)$
\item $\varphi_{5} :\equiv \forall x\,\forall y\,\forall z\,(x < y \to y \leq z \to x < z) \to \forall y\,\forall z\,(v_{290} < y \to y \leq z \to v_{290} < z)$
\item $\varphi_{6} :\equiv \forall z\,(x \leq \lvert x \rvert \to \lvert x \rvert < z \to x < z) \to x \leq \lvert x \rvert \to \lvert x \rvert < y \to x < y$
\item $\varphi_{7} :\equiv \forall y\,\forall z\,(v_{290} < y \to y \leq z \to v_{290} < z) \to \forall z\,(v_{290} < v_{291} \to v_{291} \leq z \to v_{290} < z)$
\item $\varphi_{8} :\equiv \forall z\,(v_{290} < v_{291} \to v_{291} \leq z \to v_{290} < z) \to v_{290} < v_{291} \to v_{291} \leq v_{292} \to v_{290} < v_{292}$
\item $\varphi_{9} :\equiv \forall x \in \mathbb{R}\,(-\lvert x \rvert \leq x \land x \leq \lvert x \rvert) \to x \in \mathbb{R} \to -\lvert x \rvert \leq x \land x \leq \lvert x \rvert$
\item $\varphi_{10} :\equiv (x \in \mathbb{R} \to y \in \mathbb{R} \to \lvert x \rvert < y \to \lvert x \rvert < y) \to x \in \mathbb{R} \to y \in \mathbb{R} \to \lvert x \rvert < y \to x < y$
\item $\varphi_{11} :\equiv (x \in \mathbb{R} \to y \in \mathbb{R} \to \lvert x \rvert < y \to -\lvert x \rvert \leq x) \to x \in \mathbb{R} \to y \in \mathbb{R} \to \lvert x \rvert < y \to -y < x$
\item $\varphi_{12} :\equiv (x \in \mathbb{R} \to y \in \mathbb{R} \to \lvert x \rvert < y \to x \in \mathbb{R}) \to x \in \mathbb{R} \to y \in \mathbb{R} \to \lvert x \rvert < y \to \lvert x \rvert \in \mathbb{R}$
\item $\varphi_{13} :\equiv (x \in \mathbb{R} \to y \in \mathbb{R} \to \lvert x \rvert < y \to x < y) \to x \in \mathbb{R} \to y \in \mathbb{R} \to \lvert x \rvert < y \to -y < x \land x < y$
\item $\varphi_{14} :\equiv (x \in \mathbb{R} \to y \in \mathbb{R} \to \lvert x \rvert < y \to -\lvert x \rvert \leq x \land x \leq \lvert x \rvert) \to x \in \mathbb{R} \to y \in \mathbb{R} \to \lvert x \rvert < y \to x \leq \lvert x \rvert$
\item $\varphi_{15} :\equiv (x \in \mathbb{R} \to y \in \mathbb{R} \to \lvert x \rvert < y \to \lvert x \rvert < y) \to x \in \mathbb{R} \to y \in \mathbb{R} \to \lvert x \rvert < y \to -y < -\lvert x \rvert$
\item $\varphi_{16} :\equiv (x \in \mathbb{R} \to y \in \mathbb{R} \to \lvert x \rvert < y \to x \in \mathbb{R}) \to x \in \mathbb{R} \to y \in \mathbb{R} \to \lvert x \rvert < y \to -\lvert x \rvert \leq x \land x \leq \lvert x \rvert$
\item $\varphi_{17} :\equiv \forall v_{292}\,(-y < -\lvert x \rvert \to -\lvert x \rvert \leq v_{292} \to -y < v_{292}) \to -y < -\lvert x \rvert \to -\lvert x \rvert \leq x \to -y < x$
\item $\varphi_{18} :\equiv (x \in \mathbb{R} \to y \in \mathbb{R} \to \lvert x \rvert < y \to -\lvert x \rvert \leq x \land x \leq \lvert x \rvert) \to x \in \mathbb{R} \to y \in \mathbb{R} \to \lvert x \rvert < y \to -\lvert x \rvert \leq x$
\item $\varphi_{19} :\equiv \varphi_{2} \to \forall v_{292}\,(-y < -\lvert x \rvert \to -\lvert x \rvert \leq v_{292} \to -y < v_{292})$
\item $\varphi_{20} :\equiv \forall x\,\forall y\,(x \in \mathbb{R} \to y \in \mathbb{R} \to x < y \to -y < -x) \to \forall y\,(\lvert x \rvert \in \mathbb{R} \to y \in \mathbb{R} \to \lvert x \rvert < y \to -y < -\lvert x \rvert)$
\item $\varphi_{21} :\equiv \forall y\,(\lvert x \rvert \in \mathbb{R} \to y \in \mathbb{R} \to \lvert x \rvert < y \to -y < -\lvert x \rvert) \to \lvert x \rvert \in \mathbb{R} \to y \in \mathbb{R} \to \lvert x \rvert < y \to -y < -\lvert x \rvert$
\item $\varphi_{22} :\equiv (x \leq \lvert x \rvert \to \lvert x \rvert < y \to x < y) \to (x \in \mathbb{R} \to y \in \mathbb{R} \to \lvert x \rvert < y \to x \leq \lvert x \rvert) \to \varphi_{10}$
\item $\varphi_{23} :\equiv (-y < -\lvert x \rvert \to -\lvert x \rvert \leq x \to -y < x) \to (x \in \mathbb{R} \to y \in \mathbb{R} \to \lvert x \rvert < y \to -y < -\lvert x \rvert) \to \varphi_{11}$
\item $\varphi_{24} :\equiv (x \in \mathbb{R} \to y \in \mathbb{R} \to \lvert x \rvert < y \to \lvert x \rvert \in \mathbb{R}) \to (x \in \mathbb{R} \to y \in \mathbb{R} \to \lvert x \rvert < y \to y \in \mathbb{R}) \to \varphi_{15}$
\end{abbrevlist}
\end{fullproof}

% 自動生成 (prover/build.js): LtAbs  ノード数 118
\begin{fullproof}{[A12]}{$\forall x\,\forall y\,(x \in \mathbb{R} \to y \in \mathbb{R} \to -y < x \to x < y \to \lvert x \rvert < y)$}{118}\phantomsection\label{tree:LtAbs}
\noindent$\Pi_{1}$:
\begin{ptree}
\AXC{$\forall x \in \mathbb{R}\,(-x \in \mathbb{R})$\lbl{(R1)}}
\AXC{$\forall x \in \mathbb{R}\,(-x \in \mathbb{R}) \to y \in \mathbb{R} \to -y \in \mathbb{R}$\lbl{A2}}
\BIC{$y \in \mathbb{R} \to -y \in \mathbb{R}$}
\AXC{$(y \in \mathbb{R} \to -y \in \mathbb{R}) \to \varphi_{25}$\lbl{A1}}
\BIC{$\varphi_{25}$}
\end{ptree}
\noindent$\Pi_{2}$:
\begin{ptree}
\AXC{$\forall x\,\forall y\,(x \in \mathbb{R} \to y \in \mathbb{R} \to x < y \to -y < -x)$\lbl{[F22]}}
\AXC{$\varphi_{18}$\lbl{A2}}
\BIC{$\forall y\,(v_{293} \in \mathbb{R} \to y \in \mathbb{R} \to v_{293} < y \to -y < -v_{293})$}
\AXC{$\varphi_{23}$\lbl{A2}}
\BIC{$\varphi_{4}$}
\UIC{$\forall v_{294}\,(\varphi_{4})$}
\UIC{$\forall v_{293}\,\forall v_{294}\,(\varphi_{4})$}
\AXC{$\forall v_{293}\,\forall v_{294}\,(\varphi_{4}) \to \varphi_{8}$\lbl{A2}}
\BIC{$\varphi_{8}$}
\end{ptree}
\noindent$\Pi_{3}$:
\begin{ptree}
\AXC{$\stackrel{\Pi_{2}}{\vdots}$}\noLine
\UIC{$\varphi_{8}$}
\AXC{$\varphi_{8} \to -y \in \mathbb{R} \to x \in \mathbb{R} \to -y < x \to -x < -(-y)$\lbl{A2}}
\BIC{$-y \in \mathbb{R} \to x \in \mathbb{R} \to -y < x \to -x < -(-y)$}
\AXC{$(-y \in \mathbb{R} \to x \in \mathbb{R} \to -y < x \to -x < -(-y)) \to \varphi_{34}$\lbl{A1}}
\BIC{$\varphi_{34}$}
\end{ptree}
\noindent$\Pi_{4}$:
\begin{ptree}
\AXC{$x \in \mathbb{R} \to y \in \mathbb{R} \to -y < x \to x < y \to \neg 0 \leq x \to y \in \mathbb{R}$\lbl{A1}}
\AXC{$\stackrel{\Pi_{1}}{\vdots}$}\noLine
\UIC{$\varphi_{25}$}
\BIC{$x \in \mathbb{R} \to y \in \mathbb{R} \to -y < x \to x < y \to \neg 0 \leq x \to -y \in \mathbb{R}$}
\AXC{$\stackrel{\Pi_{3}}{\vdots}$}\noLine
\UIC{$\varphi_{34}$}
\BIC{$\varphi_{31}$}
\end{ptree}
\noindent$\Pi_{5}$:
\begin{ptree}
\AXC{$x \in \mathbb{R} \to y \in \mathbb{R} \to -y < x \to x < y \to -y < x$\lbl{A1}}
\AXC{$x \in \mathbb{R} \to y \in \mathbb{R} \to -y < x \to x < y \to \neg 0 \leq x \to x \in \mathbb{R}$\lbl{A1}}
\AXC{$\stackrel{\Pi_{4}}{\vdots}$}\noLine
\UIC{$\varphi_{31}$}
\BIC{$(x \in \mathbb{R} \to y \in \mathbb{R} \to -y < x \to x < y \to -y < x) \to \varphi_{6}$}
\BIC{$\varphi_{6}$}
\end{ptree}
\noindent$\Pi_{6}$:
\begin{ptree}
\AXC{$\forall x \in \mathbb{R}\,(-(-x) = x)$\lbl{[F5]}}
\AXC{$\forall x \in \mathbb{R}\,(-(-x) = x) \to y \in \mathbb{R} \to -(-y) = y$\lbl{A2}}
\BIC{$y \in \mathbb{R} \to -(-y) = y$}
\AXC{$(y \in \mathbb{R} \to -(-y) = y) \to \varphi_{24}$\lbl{A1}}
\BIC{$\varphi_{24}$}
\end{ptree}
\noindent$\Pi_{7}$:
\begin{ptree}
\AXC{$\stackrel{\Pi_{5}}{\vdots}$}\noLine
\UIC{$\varphi_{6}$}
\AXC{$x \in \mathbb{R} \to y \in \mathbb{R} \to -y < x \to x < y \to \neg 0 \leq x \to y \in \mathbb{R}$\lbl{A1}}
\AXC{$\stackrel{\Pi_{6}}{\vdots}$}\noLine
\UIC{$\varphi_{24}$}
\BIC{$\varphi_{5}$}
\AXC{$\varphi_{5} \to \varphi_{6} \to \varphi_{17}$\lbl{A1}}
\BIC{$\varphi_{6} \to \varphi_{17}$}
\BIC{$\varphi_{17}$}
\end{ptree}
\noindent$\Pi_{8}$:
\begin{ptree}
\AXC{$x = x$\lbl{A5}}
\AXC{$x = y \to x = x \to y = x$\lbl{A6}}
\AXC{$(x = y \to x = x \to y = x) \to x = x \to x = y \to y = x$\lbl{A1}}
\BIC{$x = x \to x = y \to y = x$}
\BIC{$x = y \to y = x$}
\UIC{$\forall y\,(x = y \to y = x)$}
\UIC{$\forall x\,\forall y\,(x = y \to y = x)$}
\end{ptree}
\noindent$\Pi_{9}$:
\begin{ptree}
\AXC{$\stackrel{\Pi_{8}}{\vdots}$}\noLine
\UIC{$\forall x\,\forall y\,(x = y \to y = x)$}
\AXC{$\forall x\,\forall y\,(x = y \to y = x) \to \forall y\,(\lvert x \rvert = y \to y = \lvert x \rvert)$\lbl{A2}}
\BIC{$\forall y\,(\lvert x \rvert = y \to y = \lvert x \rvert)$}
\end{ptree}
\noindent$\Pi_{10}$:
\begin{ptree}
\AXC{$\forall x\,\forall y\,(x \in \mathbb{R} \to y \in \mathbb{R} \to \neg y \leq x \to x < y)$\lbl{[A2]}}
\AXC{$\varphi_{11}$\lbl{A2}}
\BIC{$\forall y\,(x \in \mathbb{R} \to y \in \mathbb{R} \to \neg y \leq x \to x < y)$}
\AXC{$\varphi_{9}$\lbl{A2}}
\BIC{$x \in \mathbb{R} \to 0 \in \mathbb{R} \to \neg 0 \leq x \to x < 0$}
\AXC{$(x \in \mathbb{R} \to 0 \in \mathbb{R} \to \neg 0 \leq x \to x < 0) \to \varphi_{30}$\lbl{A1}}
\BIC{$\varphi_{30}$}
\end{ptree}
\noindent$\Pi_{11}$:
\begin{ptree}
\AXC{$0 \in \mathbb{R}$\lbl{(R1)}}
\AXC{$x \in \mathbb{R} \to y \in \mathbb{R} \to -y < x \to x < y \to \neg 0 \leq x \to x \in \mathbb{R}$\lbl{A1}}
\AXC{$\stackrel{\Pi_{10}}{\vdots}$}\noLine
\UIC{$\varphi_{30}$}
\BIC{$0 \in \mathbb{R} \to \varphi_{20}$}
\BIC{$\varphi_{20}$}
\end{ptree}
\noindent$\Pi_{12}$:
\begin{ptree}
\AXC{$\forall x\,(x < 0 \to \lvert x \rvert = -x)$\lbl{[A3]}}
\AXC{$\forall x\,(x < 0 \to \lvert x \rvert = -x) \to x < 0 \to \lvert x \rvert = -x$\lbl{A2}}
\BIC{$x < 0 \to \lvert x \rvert = -x$}
\AXC{$(x < 0 \to \lvert x \rvert = -x) \to \varphi_{22}$\lbl{A1}}
\BIC{$\varphi_{22}$}
\end{ptree}
\noindent$\Pi_{13}$:
\begin{ptree}
\AXC{$x \in \mathbb{R} \to y \in \mathbb{R} \to -y < x \to x < y \to \neg 0 \leq x \to \neg 0 \leq x$\lbl{A1}}
\AXC{$\stackrel{\Pi_{11}}{\vdots}$}\noLine
\UIC{$\varphi_{20}$}
\BIC{$x \in \mathbb{R} \to y \in \mathbb{R} \to -y < x \to x < y \to \neg 0 \leq x \to x < 0$}
\AXC{$\stackrel{\Pi_{12}}{\vdots}$}\noLine
\UIC{$\varphi_{22}$}
\BIC{$x \in \mathbb{R} \to y \in \mathbb{R} \to -y < x \to x < y \to \neg 0 \leq x \to \lvert x \rvert = -x$}
\AXC{$\varphi_{27}$\lbl{A1}}
\BIC{$\varphi_{14}$}
\end{ptree}
\noindent$\Pi_{14}$:
\begin{ptree}
\AXC{$\stackrel{\Pi_{9}}{\vdots}$}\noLine
\UIC{$\forall y\,(\lvert x \rvert = y \to y = \lvert x \rvert)$}
\AXC{$\forall y\,(\lvert x \rvert = y \to y = \lvert x \rvert) \to \lvert x \rvert = -x \to -x = \lvert x \rvert$\lbl{A2}}
\BIC{$\lvert x \rvert = -x \to -x = \lvert x \rvert$}
\AXC{$\stackrel{\Pi_{13}}{\vdots}$}\noLine
\UIC{$\varphi_{14}$}
\BIC{$x \in \mathbb{R} \to y \in \mathbb{R} \to -y < x \to x < y \to \neg 0 \leq x \to -x = \lvert x \rvert$}
\AXC{$\varphi_{33}$\lbl{A1}}
\BIC{$\varphi_{28}$}
\end{ptree}
\noindent$\Pi_{15}$:
\begin{ptree}
\AXC{$z = w \to -x < z \to -x < w$\lbl{A6}}
\UIC{$\forall w\,(z = w \to -x < z \to -x < w)$}
\UIC{$\forall z\,\forall w\,(z = w \to -x < z \to -x < w)$}
\AXC{$\varphi_{12}$\lbl{A2}}
\BIC{$\forall w\,(-(-y) = w \to -x < -(-y) \to -x < w)$}
\AXC{$\varphi_{15}$\lbl{A2}}
\BIC{$-(-y) = y \to -x < -(-y) \to -x < y$}
\AXC{$\stackrel{\Pi_{7}}{\vdots}$}\noLine
\UIC{$\varphi_{17}$}
\BIC{$x \in \mathbb{R} \to y \in \mathbb{R} \to -y < x \to x < y \to \neg 0 \leq x \to -x < y$}
\AXC{$\stackrel{\Pi_{14}}{\vdots}$}\noLine
\UIC{$\varphi_{28}$}
\BIC{$\varphi_{16}$}
\end{ptree}
\noindent$\Pi_{16}$:
\begin{ptree}
\AXC{$\stackrel{\Pi_{8}}{\vdots}$}\noLine
\UIC{$\forall x\,\forall y\,(x = y \to y = x)$}
\AXC{$\forall x\,\forall y\,(x = y \to y = x) \to \forall y\,(\lvert x \rvert = y \to y = \lvert x \rvert)$\lbl{A2}}
\BIC{$\forall y\,(\lvert x \rvert = y \to y = \lvert x \rvert)$}
\end{ptree}
\noindent$\Pi_{17}$:
\begin{ptree}
\AXC{$\stackrel{\Pi_{16}}{\vdots}$}\noLine
\UIC{$\forall y\,(\lvert x \rvert = y \to y = \lvert x \rvert)$}
\AXC{$\forall y\,(\lvert x \rvert = y \to y = \lvert x \rvert) \to \lvert x \rvert = x \to x = \lvert x \rvert$\lbl{A2}}
\BIC{$\lvert x \rvert = x \to x = \lvert x \rvert$}
\end{ptree}
\noindent$\Pi_{18}$:
\begin{ptree}
\AXC{$\forall x\,(0 \leq x \to \lvert x \rvert = x)$\lbl{[A1]}}
\AXC{$\forall x\,(0 \leq x \to \lvert x \rvert = x) \to 0 \leq x \to \lvert x \rvert = x$\lbl{A2}}
\BIC{$0 \leq x \to \lvert x \rvert = x$}
\AXC{$(0 \leq x \to \lvert x \rvert = x) \to \varphi_{19}$\lbl{A1}}
\BIC{$\varphi_{19}$}
\end{ptree}
\noindent$\Pi_{19}$:
\begin{ptree}
\AXC{$\stackrel{\Pi_{17}}{\vdots}$}\noLine
\UIC{$\lvert x \rvert = x \to x = \lvert x \rvert$}
\AXC{$x \in \mathbb{R} \to y \in \mathbb{R} \to -y < x \to x < y \to 0 \leq x \to 0 \leq x$\lbl{A1}}
\AXC{$\stackrel{\Pi_{18}}{\vdots}$}\noLine
\UIC{$\varphi_{19}$}
\BIC{$x \in \mathbb{R} \to y \in \mathbb{R} \to -y < x \to x < y \to 0 \leq x \to \lvert x \rvert = x$}
\AXC{$\varphi_{26}$\lbl{A1}}
\BIC{$\varphi_{10}$}
\BIC{$x \in \mathbb{R} \to y \in \mathbb{R} \to -y < x \to x < y \to 0 \leq x \to x = \lvert x \rvert$}
\AXC{$\varphi_{32}$\lbl{A1}}
\BIC{$(x \in \mathbb{R} \to y \in \mathbb{R} \to -y < x \to x < y \to x < y) \to \varphi_{13}$}
\end{ptree}
\noindent$\Pi_{20}$:
\begin{ptree}
\AXC{$x \in \mathbb{R} \to y \in \mathbb{R} \to -y < x \to x < y \to x < y$\lbl{A1}}
\AXC{$\stackrel{\Pi_{19}}{\vdots}$}\noLine
\UIC{$(x \in \mathbb{R} \to y \in \mathbb{R} \to -y < x \to x < y \to x < y) \to \varphi_{13}$}
\BIC{$\varphi_{13}$}
\end{ptree}
\noindent$\Pi_{21}$:
\begin{ptree}
\AXC{$z = w \to z < y \to w < y$\lbl{A6}}
\UIC{$\forall w\,(z = w \to z < y \to w < y)$}
\UIC{$\forall z\,\forall w\,(z = w \to z < y \to w < y)$}
\AXC{$\varphi_{1}$\lbl{A2}}
\BIC{$\forall w\,(x = w \to x < y \to w < y)$}
\AXC{$\varphi_{2}$\lbl{A2}}
\BIC{$x = \lvert x \rvert \to x < y \to \lvert x \rvert < y$}
\AXC{$\stackrel{\Pi_{20}}{\vdots}$}\noLine
\UIC{$\varphi_{13}$}
\BIC{$x \in \mathbb{R} \to y \in \mathbb{R} \to -y < x \to x < y \to 0 \leq x \to \lvert x \rvert < y$}
\AXC{$\varphi_{29}$\lbl{A1}}
\BIC{$\varphi_{21}$}
\end{ptree}
\noindent 主木:
\begin{ptree}
\AXC{$z = w \to z < y \to w < y$\lbl{A6}}
\UIC{$\forall w\,(z = w \to z < y \to w < y)$}
\UIC{$\forall z\,\forall w\,(z = w \to z < y \to w < y)$}
\AXC{$\varphi_{3}$\lbl{A2}}
\BIC{$\forall w\,(-x = w \to -x < y \to w < y)$}
\AXC{$\varphi_{7}$\lbl{A2}}
\BIC{$-x = \lvert x \rvert \to -x < y \to \lvert x \rvert < y$}
\AXC{$\stackrel{\Pi_{15}}{\vdots}$}\noLine
\UIC{$\varphi_{16}$}
\BIC{$x \in \mathbb{R} \to y \in \mathbb{R} \to -y < x \to x < y \to \neg 0 \leq x \to \lvert x \rvert < y$}
\AXC{$\stackrel{\Pi_{21}}{\vdots}$}\noLine
\UIC{$\varphi_{21}$}
\BIC{$x \in \mathbb{R} \to y \in \mathbb{R} \to -y < x \to x < y \to \lvert x \rvert < y$}
\UIC{$\forall y\,(x \in \mathbb{R} \to y \in \mathbb{R} \to -y < x \to x < y \to \lvert x \rvert < y)$}
\UIC{$\forall x\,\forall y\,(x \in \mathbb{R} \to y \in \mathbb{R} \to -y < x \to x < y \to \lvert x \rvert < y)$}
\end{ptree}
\begin{abbrevlist}
\item $\varphi_{1} :\equiv \forall z\,\forall w\,(z = w \to z < y \to w < y) \to \forall w\,(x = w \to x < y \to w < y)$
\item $\varphi_{2} :\equiv \forall w\,(x = w \to x < y \to w < y) \to x = \lvert x \rvert \to x < y \to \lvert x \rvert < y$
\item $\varphi_{3} :\equiv \forall z\,\forall w\,(z = w \to z < y \to w < y) \to \forall w\,(-x = w \to -x < y \to w < y)$
\item $\varphi_{4} :\equiv v_{293} \in \mathbb{R} \to v_{294} \in \mathbb{R} \to v_{293} < v_{294} \to -v_{294} < -v_{293}$
\item $\varphi_{5} :\equiv x \in \mathbb{R} \to y \in \mathbb{R} \to -y < x \to x < y \to \neg 0 \leq x \to -(-y) = y$
\item $\varphi_{6} :\equiv x \in \mathbb{R} \to y \in \mathbb{R} \to -y < x \to x < y \to \neg 0 \leq x \to -x < -(-y)$
\item $\varphi_{7} :\equiv \forall w\,(-x = w \to -x < y \to w < y) \to -x = \lvert x \rvert \to -x < y \to \lvert x \rvert < y$
\item $\varphi_{8} :\equiv \forall v_{294}\,(-y \in \mathbb{R} \to v_{294} \in \mathbb{R} \to -y < v_{294} \to -v_{294} < -(-y))$
\item $\varphi_{9} :\equiv \forall y\,(x \in \mathbb{R} \to y \in \mathbb{R} \to \neg y \leq x \to x < y) \to x \in \mathbb{R} \to 0 \in \mathbb{R} \to \neg 0 \leq x \to x < 0$
\item $\varphi_{10} :\equiv (\lvert x \rvert = x \to x = \lvert x \rvert) \to x \in \mathbb{R} \to y \in \mathbb{R} \to -y < x \to x < y \to 0 \leq x \to x = \lvert x \rvert$
\item $\varphi_{11} :\equiv \forall x\,\forall y\,(x \in \mathbb{R} \to y \in \mathbb{R} \to \neg y \leq x \to x < y) \to \forall y\,(x \in \mathbb{R} \to y \in \mathbb{R} \to \neg y \leq x \to x < y)$
\item $\varphi_{12} :\equiv \forall z\,\forall w\,(z = w \to -x < z \to -x < w) \to \forall w\,(-(-y) = w \to -x < -(-y) \to -x < w)$
\item $\varphi_{13} :\equiv (x = \lvert x \rvert \to x < y \to \lvert x \rvert < y) \to x \in \mathbb{R} \to y \in \mathbb{R} \to -y < x \to x < y \to 0 \leq x \to \lvert x \rvert < y$
\item $\varphi_{14} :\equiv (\lvert x \rvert = -x \to -x = \lvert x \rvert) \to x \in \mathbb{R} \to y \in \mathbb{R} \to -y < x \to x < y \to \neg 0 \leq x \to -x = \lvert x \rvert$
\item $\varphi_{15} :\equiv \forall w\,(-(-y) = w \to -x < -(-y) \to -x < w) \to -(-y) = y \to -x < -(-y) \to -x < y$
\item $\varphi_{16} :\equiv (-x = \lvert x \rvert \to -x < y \to \lvert x \rvert < y) \to x \in \mathbb{R} \to y \in \mathbb{R} \to -y < x \to x < y \to \neg 0 \leq x \to \lvert x \rvert < y$
\item $\varphi_{17} :\equiv (-(-y) = y \to -x < -(-y) \to -x < y) \to x \in \mathbb{R} \to y \in \mathbb{R} \to -y < x \to x < y \to \neg 0 \leq x \to -x < y$
\item $\varphi_{18} :\equiv \forall x\,\forall y\,(x \in \mathbb{R} \to y \in \mathbb{R} \to x < y \to -y < -x) \to \forall y\,(v_{293} \in \mathbb{R} \to y \in \mathbb{R} \to v_{293} < y \to -y < -v_{293})$
\item $\varphi_{19} :\equiv (x \in \mathbb{R} \to y \in \mathbb{R} \to -y < x \to x < y \to 0 \leq x \to 0 \leq x) \to x \in \mathbb{R} \to y \in \mathbb{R} \to -y < x \to x < y \to 0 \leq x \to \lvert x \rvert = x$
\item $\varphi_{20} :\equiv (x \in \mathbb{R} \to y \in \mathbb{R} \to -y < x \to x < y \to \neg 0 \leq x \to \neg 0 \leq x) \to x \in \mathbb{R} \to y \in \mathbb{R} \to -y < x \to x < y \to \neg 0 \leq x \to x < 0$
\item $\varphi_{21} :\equiv (x \in \mathbb{R} \to y \in \mathbb{R} \to -y < x \to x < y \to \neg 0 \leq x \to \lvert x \rvert < y) \to x \in \mathbb{R} \to y \in \mathbb{R} \to -y < x \to x < y \to \lvert x \rvert < y$
\item $\varphi_{22} :\equiv (x \in \mathbb{R} \to y \in \mathbb{R} \to -y < x \to x < y \to \neg 0 \leq x \to x < 0) \to x \in \mathbb{R} \to y \in \mathbb{R} \to -y < x \to x < y \to \neg 0 \leq x \to \lvert x \rvert = -x$
\item $\varphi_{23} :\equiv \forall y\,(v_{293} \in \mathbb{R} \to y \in \mathbb{R} \to v_{293} < y \to -y < -v_{293}) \to \varphi_{4}$
\item $\varphi_{24} :\equiv (x \in \mathbb{R} \to y \in \mathbb{R} \to -y < x \to x < y \to \neg 0 \leq x \to y \in \mathbb{R}) \to \varphi_{5}$
\item $\varphi_{25} :\equiv (x \in \mathbb{R} \to y \in \mathbb{R} \to -y < x \to x < y \to \neg 0 \leq x \to y \in \mathbb{R}) \to x \in \mathbb{R} \to y \in \mathbb{R} \to -y < x \to x < y \to \neg 0 \leq x \to -y \in \mathbb{R}$
\item $\varphi_{26} :\equiv (x \in \mathbb{R} \to y \in \mathbb{R} \to -y < x \to x < y \to 0 \leq x \to \lvert x \rvert = x) \to \varphi_{10}$
\item $\varphi_{27} :\equiv (x \in \mathbb{R} \to y \in \mathbb{R} \to -y < x \to x < y \to \neg 0 \leq x \to \lvert x \rvert = -x) \to \varphi_{14}$
\item $\varphi_{28} :\equiv (x \in \mathbb{R} \to y \in \mathbb{R} \to -y < x \to x < y \to \neg 0 \leq x \to -x < y) \to \varphi_{16}$
\item $\varphi_{29} :\equiv (x \in \mathbb{R} \to y \in \mathbb{R} \to -y < x \to x < y \to 0 \leq x \to \lvert x \rvert < y) \to \varphi_{21}$
\item $\varphi_{30} :\equiv (x \in \mathbb{R} \to y \in \mathbb{R} \to -y < x \to x < y \to \neg 0 \leq x \to x \in \mathbb{R}) \to 0 \in \mathbb{R} \to \varphi_{20}$
\item $\varphi_{31} :\equiv (x \in \mathbb{R} \to y \in \mathbb{R} \to -y < x \to x < y \to \neg 0 \leq x \to x \in \mathbb{R}) \to (x \in \mathbb{R} \to y \in \mathbb{R} \to -y < x \to x < y \to -y < x) \to \varphi_{6}$
\item $\varphi_{32} :\equiv (x \in \mathbb{R} \to y \in \mathbb{R} \to -y < x \to x < y \to 0 \leq x \to x = \lvert x \rvert) \to (x \in \mathbb{R} \to y \in \mathbb{R} \to -y < x \to x < y \to x < y) \to \varphi_{13}$
\item $\varphi_{33} :\equiv (x \in \mathbb{R} \to y \in \mathbb{R} \to -y < x \to x < y \to \neg 0 \leq x \to -x = \lvert x \rvert) \to \varphi_{28}$
\item $\varphi_{34} :\equiv (x \in \mathbb{R} \to y \in \mathbb{R} \to -y < x \to x < y \to \neg 0 \leq x \to -y \in \mathbb{R}) \to \varphi_{31}$
\end{abbrevlist}
\end{fullproof}

% 自動生成 (prover/build.js): AbsSubComm  ノード数 41
\begin{fullproof}{[A13]}{$\forall x\,\forall y\,(x \in \mathbb{R} \to y \in \mathbb{R} \to \lvert x - y \rvert = \lvert y - x \rvert)$}{41}\phantomsection\label{tree:AbsSubComm}
\noindent$\Pi_{1}$:
\begin{ptree}
\AXC{$z = w \to \lvert z \rvert = \lvert y - x \rvert \to \lvert w \rvert = \lvert y - x \rvert$\lbl{A6}}
\UIC{$\forall w\,(z = w \to \lvert z \rvert = \lvert y - x \rvert \to \lvert w \rvert = \lvert y - x \rvert)$}
\UIC{$\forall z\,\forall w\,(z = w \to \lvert z \rvert = \lvert y - x \rvert \to \lvert w \rvert = \lvert y - x \rvert)$}
\AXC{$\varphi_{12}$\lbl{A2}}
\BIC{$\varphi_{2}$}
\AXC{$\varphi_{2} \to \varphi_{3}$\lbl{A2}}
\BIC{$\varphi_{3}$}
\end{ptree}
\noindent$\Pi_{2}$:
\begin{ptree}
\AXC{$\forall x\,\forall y\,(x \in \mathbb{R} \to y \in \mathbb{R} \to x - y \in \mathbb{R})$\lbl{[F1]}}
\AXC{$\varphi_{9}$\lbl{A2}}
\BIC{$\forall y\,(v_{295} \in \mathbb{R} \to y \in \mathbb{R} \to v_{295} - y \in \mathbb{R})$}
\AXC{$\varphi_{11}$\lbl{A2}}
\BIC{$v_{295} \in \mathbb{R} \to v_{296} \in \mathbb{R} \to v_{295} - v_{296} \in \mathbb{R}$}
\UIC{$\forall v_{296}\,(v_{295} \in \mathbb{R} \to v_{296} \in \mathbb{R} \to v_{295} - v_{296} \in \mathbb{R})$}
\UIC{$\varphi_{1}$}
\AXC{$\varphi_{1} \to \forall v_{296}\,(y \in \mathbb{R} \to v_{296} \in \mathbb{R} \to y - v_{296} \in \mathbb{R})$\lbl{A2}}
\BIC{$\forall v_{296}\,(y \in \mathbb{R} \to v_{296} \in \mathbb{R} \to y - v_{296} \in \mathbb{R})$}
\end{ptree}
\noindent$\Pi_{3}$:
\begin{ptree}
\AXC{$x \in \mathbb{R} \to y \in \mathbb{R} \to y \in \mathbb{R}$\lbl{A1}}
\AXC{$\stackrel{\Pi_{2}}{\vdots}$}\noLine
\UIC{$\forall v_{296}\,(y \in \mathbb{R} \to v_{296} \in \mathbb{R} \to y - v_{296} \in \mathbb{R})$}
\AXC{$\varphi_{7}$\lbl{A2}}
\BIC{$y \in \mathbb{R} \to x \in \mathbb{R} \to y - x \in \mathbb{R}$}
\AXC{$\varphi_{13}$\lbl{A1}}
\BIC{$(x \in \mathbb{R} \to y \in \mathbb{R} \to y \in \mathbb{R}) \to \varphi_{4}$}
\BIC{$\varphi_{4}$}
\end{ptree}
\noindent$\Pi_{4}$:
\begin{ptree}
\AXC{$\forall x \in \mathbb{R}\,(\lvert -x \rvert = \lvert x \rvert)$\lbl{[A7]}}
\AXC{$\varphi_{5}$\lbl{A2}}
\BIC{$y - x \in \mathbb{R} \to \lvert -(y - x) \rvert = \lvert y - x \rvert$}
\AXC{$(y - x \in \mathbb{R} \to \lvert -(y - x) \rvert = \lvert y - x \rvert) \to \varphi_{8}$\lbl{A1}}
\BIC{$\varphi_{8}$}
\end{ptree}
\noindent$\Pi_{5}$:
\begin{ptree}
\AXC{$\forall x\,\forall y\,(x \in \mathbb{R} \to y \in \mathbb{R} \to -(y - x) = x - y)$\lbl{[F14]}}
\AXC{$\varphi_{10}$\lbl{A2}}
\BIC{$\forall y\,(x \in \mathbb{R} \to y \in \mathbb{R} \to -(y - x) = x - y)$}
\AXC{$\varphi_{6}$\lbl{A2}}
\BIC{$x \in \mathbb{R} \to y \in \mathbb{R} \to -(y - x) = x - y$}
\AXC{$(x \in \mathbb{R} \to y \in \mathbb{R} \to -(y - x) = x - y) \to \varphi_{14}$\lbl{A1}}
\BIC{$\varphi_{14}$}
\end{ptree}
\noindent 主木:
\begin{ptree}
\AXC{$\stackrel{\Pi_{1}}{\vdots}$}\noLine
\UIC{$\varphi_{3}$}
\AXC{$x \in \mathbb{R} \to y \in \mathbb{R} \to x \in \mathbb{R}$\lbl{A1}}
\AXC{$\stackrel{\Pi_{3}}{\vdots}$}\noLine
\UIC{$\varphi_{4}$}
\BIC{$x \in \mathbb{R} \to y \in \mathbb{R} \to y - x \in \mathbb{R}$}
\AXC{$\stackrel{\Pi_{4}}{\vdots}$}\noLine
\UIC{$\varphi_{8}$}
\BIC{$x \in \mathbb{R} \to y \in \mathbb{R} \to \lvert -(y - x) \rvert = \lvert y - x \rvert$}
\AXC{$\stackrel{\Pi_{5}}{\vdots}$}\noLine
\UIC{$\varphi_{14}$}
\BIC{$\varphi_{3} \to x \in \mathbb{R} \to y \in \mathbb{R} \to \lvert x - y \rvert = \lvert y - x \rvert$}
\BIC{$x \in \mathbb{R} \to y \in \mathbb{R} \to \lvert x - y \rvert = \lvert y - x \rvert$}
\UIC{$\forall y\,(x \in \mathbb{R} \to y \in \mathbb{R} \to \lvert x - y \rvert = \lvert y - x \rvert)$}
\UIC{$\forall x\,\forall y\,(x \in \mathbb{R} \to y \in \mathbb{R} \to \lvert x - y \rvert = \lvert y - x \rvert)$}
\end{ptree}
\begin{abbrevlist}
\item $\varphi_{1} :\equiv \forall v_{295}\,\forall v_{296}\,(v_{295} \in \mathbb{R} \to v_{296} \in \mathbb{R} \to v_{295} - v_{296} \in \mathbb{R})$
\item $\varphi_{2} :\equiv \forall w\,(-(y - x) = w \to \lvert -(y - x) \rvert = \lvert y - x \rvert \to \lvert w \rvert = \lvert y - x \rvert)$
\item $\varphi_{3} :\equiv -(y - x) = x - y \to \lvert -(y - x) \rvert = \lvert y - x \rvert \to \lvert x - y \rvert = \lvert y - x \rvert$
\item $\varphi_{4} :\equiv (x \in \mathbb{R} \to y \in \mathbb{R} \to x \in \mathbb{R}) \to x \in \mathbb{R} \to y \in \mathbb{R} \to y - x \in \mathbb{R}$
\item $\varphi_{5} :\equiv \forall x \in \mathbb{R}\,(\lvert -x \rvert = \lvert x \rvert) \to y - x \in \mathbb{R} \to \lvert -(y - x) \rvert = \lvert y - x \rvert$
\item $\varphi_{6} :\equiv \forall y\,(x \in \mathbb{R} \to y \in \mathbb{R} \to -(y - x) = x - y) \to x \in \mathbb{R} \to y \in \mathbb{R} \to -(y - x) = x - y$
\item $\varphi_{7} :\equiv \forall v_{296}\,(y \in \mathbb{R} \to v_{296} \in \mathbb{R} \to y - v_{296} \in \mathbb{R}) \to y \in \mathbb{R} \to x \in \mathbb{R} \to y - x \in \mathbb{R}$
\item $\varphi_{8} :\equiv (x \in \mathbb{R} \to y \in \mathbb{R} \to y - x \in \mathbb{R}) \to x \in \mathbb{R} \to y \in \mathbb{R} \to \lvert -(y - x) \rvert = \lvert y - x \rvert$
\item $\varphi_{9} :\equiv \forall x\,\forall y\,(x \in \mathbb{R} \to y \in \mathbb{R} \to x - y \in \mathbb{R}) \to \forall y\,(v_{295} \in \mathbb{R} \to y \in \mathbb{R} \to v_{295} - y \in \mathbb{R})$
\item $\varphi_{10} :\equiv \forall x\,\forall y\,(x \in \mathbb{R} \to y \in \mathbb{R} \to -(y - x) = x - y) \to \forall y\,(x \in \mathbb{R} \to y \in \mathbb{R} \to -(y - x) = x - y)$
\item $\varphi_{11} :\equiv \forall y\,(v_{295} \in \mathbb{R} \to y \in \mathbb{R} \to v_{295} - y \in \mathbb{R}) \to v_{295} \in \mathbb{R} \to v_{296} \in \mathbb{R} \to v_{295} - v_{296} \in \mathbb{R}$
\item $\varphi_{12} :\equiv \forall z\,\forall w\,(z = w \to \lvert z \rvert = \lvert y - x \rvert \to \lvert w \rvert = \lvert y - x \rvert) \to \varphi_{2}$
\item $\varphi_{13} :\equiv (y \in \mathbb{R} \to x \in \mathbb{R} \to y - x \in \mathbb{R}) \to (x \in \mathbb{R} \to y \in \mathbb{R} \to y \in \mathbb{R}) \to \varphi_{4}$
\item $\varphi_{14} :\equiv (x \in \mathbb{R} \to y \in \mathbb{R} \to \lvert -(y - x) \rvert = \lvert y - x \rvert) \to \varphi_{3} \to x \in \mathbb{R} \to y \in \mathbb{R} \to \lvert x - y \rvert = \lvert y - x \rvert$
\end{abbrevlist}
\end{fullproof}

% 自動生成 (prover/build.js): Triangle  ノード数 155
\begin{fullproof}{命題1.2}{$\forall a\,\forall b\,(a \in \mathbb{R} \to b \in \mathbb{R} \to \lvert a + b \rvert \leq \lvert a \rvert + \lvert b \rvert)$}{155}\phantomsection\label{tree:Triangle}
\noindent$\Pi_{1}$:
\begin{ptree}
\AXC{$\forall x \in \mathbb{R}\,(-\lvert x \rvert \leq x \land x \leq \lvert x \rvert)$\lbl{[A8]}}
\AXC{$\varphi_{7}$\lbl{A2}}
\BIC{$b \in \mathbb{R} \to -\lvert b \rvert \leq b \land b \leq \lvert b \rvert$}
\AXC{$(b \in \mathbb{R} \to -\lvert b \rvert \leq b \land b \leq \lvert b \rvert) \to \varphi_{19}$\lbl{A1}}
\BIC{$\varphi_{19}$}
\end{ptree}
\noindent$\Pi_{2}$:
\begin{ptree}
\AXC{$a \in \mathbb{R} \to b \in \mathbb{R} \to b \in \mathbb{R}$\lbl{A1}}
\AXC{$\stackrel{\Pi_{1}}{\vdots}$}\noLine
\UIC{$\varphi_{19}$}
\BIC{$a \in \mathbb{R} \to b \in \mathbb{R} \to -\lvert b \rvert \leq b \land b \leq \lvert b \rvert$}
\end{ptree}
\noindent$\Pi_{3}$:
\begin{ptree}
\AXC{$\forall x \in \mathbb{R}\,(-\lvert x \rvert \leq x \land x \leq \lvert x \rvert)$\lbl{[A8]}}
\AXC{$\varphi_{8}$\lbl{A2}}
\BIC{$a \in \mathbb{R} \to -\lvert a \rvert \leq a \land a \leq \lvert a \rvert$}
\AXC{$(a \in \mathbb{R} \to -\lvert a \rvert \leq a \land a \leq \lvert a \rvert) \to \varphi_{20}$\lbl{A1}}
\BIC{$\varphi_{20}$}
\end{ptree}
\noindent$\Pi_{4}$:
\begin{ptree}
\AXC{$a \in \mathbb{R} \to b \in \mathbb{R} \to a \in \mathbb{R}$\lbl{A1}}
\AXC{$\stackrel{\Pi_{3}}{\vdots}$}\noLine
\UIC{$\varphi_{20}$}
\BIC{$a \in \mathbb{R} \to b \in \mathbb{R} \to -\lvert a \rvert \leq a \land a \leq \lvert a \rvert$}
\end{ptree}
\noindent$\Pi_{5}$:
\begin{ptree}
\AXC{$\forall x \in \mathbb{R}\,(\lvert x \rvert \in \mathbb{R})$\lbl{[A5]}}
\AXC{$\forall x \in \mathbb{R}\,(\lvert x \rvert \in \mathbb{R}) \to b \in \mathbb{R} \to \lvert b \rvert \in \mathbb{R}$\lbl{A2}}
\BIC{$b \in \mathbb{R} \to \lvert b \rvert \in \mathbb{R}$}
\AXC{$(b \in \mathbb{R} \to \lvert b \rvert \in \mathbb{R}) \to \varphi_{9}$\lbl{A1}}
\BIC{$\varphi_{9}$}
\end{ptree}
\noindent$\Pi_{6}$:
\begin{ptree}
\AXC{$a \in \mathbb{R} \to b \in \mathbb{R} \to b \in \mathbb{R}$\lbl{A1}}
\AXC{$\stackrel{\Pi_{5}}{\vdots}$}\noLine
\UIC{$\varphi_{9}$}
\BIC{$a \in \mathbb{R} \to b \in \mathbb{R} \to \lvert b \rvert \in \mathbb{R}$}
\end{ptree}
\noindent$\Pi_{7}$:
\begin{ptree}
\AXC{$\forall x \in \mathbb{R}\,(\lvert x \rvert \in \mathbb{R})$\lbl{[A5]}}
\AXC{$\forall x \in \mathbb{R}\,(\lvert x \rvert \in \mathbb{R}) \to a \in \mathbb{R} \to \lvert a \rvert \in \mathbb{R}$\lbl{A2}}
\BIC{$a \in \mathbb{R} \to \lvert a \rvert \in \mathbb{R}$}
\AXC{$(a \in \mathbb{R} \to \lvert a \rvert \in \mathbb{R}) \to \varphi_{10}$\lbl{A1}}
\BIC{$\varphi_{10}$}
\end{ptree}
\noindent$\Pi_{8}$:
\begin{ptree}
\AXC{$a \in \mathbb{R} \to b \in \mathbb{R} \to a \in \mathbb{R}$\lbl{A1}}
\AXC{$\stackrel{\Pi_{7}}{\vdots}$}\noLine
\UIC{$\varphi_{10}$}
\BIC{$a \in \mathbb{R} \to b \in \mathbb{R} \to \lvert a \rvert \in \mathbb{R}$}
\end{ptree}
\noindent$\Pi_{9}$:
\begin{ptree}
\AXC{$\varphi_{21}$\lbl{[F21]}}
\AXC{$\varphi_{21} \to \varphi_{14}$\lbl{A2}}
\BIC{$\varphi_{14}$}
\AXC{$\varphi_{14} \to \varphi_{29}$\lbl{A2}}
\BIC{$\varphi_{29}$}
\AXC{$\varphi_{29} \to \varphi_{27}$\lbl{A2}}
\BIC{$\varphi_{27}$}
\AXC{$\varphi_{27} \to a \in \mathbb{R} \to \lvert a \rvert \in \mathbb{R} \to \varphi_{4}$\lbl{A2}}
\BIC{$a \in \mathbb{R} \to \lvert a \rvert \in \mathbb{R} \to \varphi_{4}$}
\AXC{$\varphi_{53}$\lbl{A1}}
\BIC{$(a \in \mathbb{R} \to b \in \mathbb{R} \to a \in \mathbb{R}) \to \varphi_{48}$}
\end{ptree}
\noindent$\Pi_{10}$:
\begin{ptree}
\AXC{$\stackrel{\Pi_{8}}{\vdots}$}\noLine
\UIC{$a \in \mathbb{R} \to b \in \mathbb{R} \to \lvert a \rvert \in \mathbb{R}$}
\AXC{$a \in \mathbb{R} \to b \in \mathbb{R} \to a \in \mathbb{R}$\lbl{A1}}
\AXC{$\stackrel{\Pi_{9}}{\vdots}$}\noLine
\UIC{$(a \in \mathbb{R} \to b \in \mathbb{R} \to a \in \mathbb{R}) \to \varphi_{48}$}
\BIC{$\varphi_{48}$}
\BIC{$(a \in \mathbb{R} \to b \in \mathbb{R} \to b \in \mathbb{R}) \to \varphi_{43}$}
\end{ptree}
\noindent$\Pi_{11}$:
\begin{ptree}
\AXC{$\stackrel{\Pi_{6}}{\vdots}$}\noLine
\UIC{$a \in \mathbb{R} \to b \in \mathbb{R} \to \lvert b \rvert \in \mathbb{R}$}
\AXC{$a \in \mathbb{R} \to b \in \mathbb{R} \to b \in \mathbb{R}$\lbl{A1}}
\AXC{$\stackrel{\Pi_{10}}{\vdots}$}\noLine
\UIC{$(a \in \mathbb{R} \to b \in \mathbb{R} \to b \in \mathbb{R}) \to \varphi_{43}$}
\BIC{$\varphi_{43}$}
\BIC{$(a \in \mathbb{R} \to b \in \mathbb{R} \to a \leq \lvert a \rvert) \to \varphi_{15}$}
\end{ptree}
\noindent$\Pi_{12}$:
\begin{ptree}
\AXC{$\stackrel{\Pi_{4}}{\vdots}$}\noLine
\UIC{$a \in \mathbb{R} \to b \in \mathbb{R} \to -\lvert a \rvert \leq a \land a \leq \lvert a \rvert$}
\AXC{$\varphi_{17}$\lbl{A1}}
\BIC{$a \in \mathbb{R} \to b \in \mathbb{R} \to a \leq \lvert a \rvert$}
\AXC{$\stackrel{\Pi_{11}}{\vdots}$}\noLine
\UIC{$(a \in \mathbb{R} \to b \in \mathbb{R} \to a \leq \lvert a \rvert) \to \varphi_{15}$}
\BIC{$\varphi_{15}$}
\end{ptree}
\noindent$\Pi_{13}$:
\begin{ptree}
\AXC{$z = w \to z = a + b \to w = a + b$\lbl{A6}}
\UIC{$\forall w\,(z = w \to z = a + b \to w = a + b)$}
\UIC{$\forall z\,\forall w\,(z = w \to z = a + b \to w = a + b)$}
\AXC{$\forall z\,\forall w\,(z = w \to z = a + b \to w = a + b) \to \varphi_{5}$\lbl{A2}}
\BIC{$\varphi_{5}$}
\AXC{$\varphi_{5} \to \varphi_{35}$\lbl{A2}}
\BIC{$\varphi_{35}$}
\end{ptree}
\noindent$\Pi_{14}$:
\begin{ptree}
\AXC{$z = w \to z < a + b \to w < a + b$\lbl{A6}}
\UIC{$\forall w\,(z = w \to z < a + b \to w < a + b)$}
\UIC{$\forall z\,\forall w\,(z = w \to z < a + b \to w < a + b)$}
\AXC{$\forall z\,\forall w\,(z = w \to z < a + b \to w < a + b) \to \varphi_{11}$\lbl{A2}}
\BIC{$\varphi_{11}$}
\AXC{$\varphi_{11} \to \varphi_{39}$\lbl{A2}}
\BIC{$\varphi_{39}$}
\end{ptree}
\noindent$\Pi_{15}$:
\begin{ptree}
\AXC{$\forall x \in \mathbb{R}\,(-x \in \mathbb{R})$\lbl{(R1)}}
\AXC{$\forall x \in \mathbb{R}\,(-x \in \mathbb{R}) \to \lvert b \rvert \in \mathbb{R} \to -\lvert b \rvert \in \mathbb{R}$\lbl{A2}}
\BIC{$\lvert b \rvert \in \mathbb{R} \to -\lvert b \rvert \in \mathbb{R}$}
\AXC{$(\lvert b \rvert \in \mathbb{R} \to -\lvert b \rvert \in \mathbb{R}) \to \varphi_{22}$\lbl{A1}}
\BIC{$\varphi_{22}$}
\end{ptree}
\noindent$\Pi_{16}$:
\begin{ptree}
\AXC{$\forall x \in \mathbb{R}\,(-x \in \mathbb{R})$\lbl{(R1)}}
\AXC{$\forall x \in \mathbb{R}\,(-x \in \mathbb{R}) \to \lvert a \rvert \in \mathbb{R} \to -\lvert a \rvert \in \mathbb{R}$\lbl{A2}}
\BIC{$\lvert a \rvert \in \mathbb{R} \to -\lvert a \rvert \in \mathbb{R}$}
\AXC{$(\lvert a \rvert \in \mathbb{R} \to -\lvert a \rvert \in \mathbb{R}) \to \varphi_{23}$\lbl{A1}}
\BIC{$\varphi_{23}$}
\end{ptree}
\noindent$\Pi_{17}$:
\begin{ptree}
\AXC{$\varphi_{21}$\lbl{[F21]}}
\AXC{$\varphi_{21} \to \varphi_{36}$\lbl{A2}}
\BIC{$\varphi_{36}$}
\AXC{$\varphi_{36} \to \varphi_{34}$\lbl{A2}}
\BIC{$\varphi_{34}$}
\AXC{$\varphi_{34} \to \varphi_{41}$\lbl{A2}}
\BIC{$\varphi_{41}$}
\AXC{$\varphi_{41} \to -\lvert a \rvert \in \mathbb{R} \to a \in \mathbb{R} \to -\lvert b \rvert \in \mathbb{R} \to \varphi_{6}$\lbl{A2}}
\BIC{$-\lvert a \rvert \in \mathbb{R} \to a \in \mathbb{R} \to -\lvert b \rvert \in \mathbb{R} \to \varphi_{6}$}
\AXC{$\varphi_{55}$\lbl{A1}}
\BIC{$(a \in \mathbb{R} \to b \in \mathbb{R} \to -\lvert a \rvert \in \mathbb{R}) \to \varphi_{50}$}
\end{ptree}
\noindent$\Pi_{18}$:
\begin{ptree}
\AXC{$a \in \mathbb{R} \to b \in \mathbb{R} \to a \in \mathbb{R}$\lbl{A1}}
\AXC{$\stackrel{\Pi_{8}}{\vdots}$}\noLine
\UIC{$a \in \mathbb{R} \to b \in \mathbb{R} \to \lvert a \rvert \in \mathbb{R}$}
\AXC{$\stackrel{\Pi_{16}}{\vdots}$}\noLine
\UIC{$\varphi_{23}$}
\BIC{$a \in \mathbb{R} \to b \in \mathbb{R} \to -\lvert a \rvert \in \mathbb{R}$}
\AXC{$\stackrel{\Pi_{17}}{\vdots}$}\noLine
\UIC{$(a \in \mathbb{R} \to b \in \mathbb{R} \to -\lvert a \rvert \in \mathbb{R}) \to \varphi_{50}$}
\BIC{$\varphi_{50}$}
\BIC{$(a \in \mathbb{R} \to b \in \mathbb{R} \to -\lvert b \rvert \in \mathbb{R}) \to \varphi_{45}$}
\end{ptree}
\noindent$\Pi_{19}$:
\begin{ptree}
\AXC{$a \in \mathbb{R} \to b \in \mathbb{R} \to b \in \mathbb{R}$\lbl{A1}}
\AXC{$\stackrel{\Pi_{6}}{\vdots}$}\noLine
\UIC{$a \in \mathbb{R} \to b \in \mathbb{R} \to \lvert b \rvert \in \mathbb{R}$}
\AXC{$\stackrel{\Pi_{15}}{\vdots}$}\noLine
\UIC{$\varphi_{22}$}
\BIC{$a \in \mathbb{R} \to b \in \mathbb{R} \to -\lvert b \rvert \in \mathbb{R}$}
\AXC{$\stackrel{\Pi_{18}}{\vdots}$}\noLine
\UIC{$(a \in \mathbb{R} \to b \in \mathbb{R} \to -\lvert b \rvert \in \mathbb{R}) \to \varphi_{45}$}
\BIC{$\varphi_{45}$}
\BIC{$(a \in \mathbb{R} \to b \in \mathbb{R} \to -\lvert a \rvert \leq a) \to \varphi_{31}$}
\end{ptree}
\noindent$\Pi_{20}$:
\begin{ptree}
\AXC{$\stackrel{\Pi_{4}}{\vdots}$}\noLine
\UIC{$a \in \mathbb{R} \to b \in \mathbb{R} \to -\lvert a \rvert \leq a \land a \leq \lvert a \rvert$}
\AXC{$\varphi_{25}$\lbl{A1}}
\BIC{$a \in \mathbb{R} \to b \in \mathbb{R} \to -\lvert a \rvert \leq a$}
\AXC{$\stackrel{\Pi_{19}}{\vdots}$}\noLine
\UIC{$(a \in \mathbb{R} \to b \in \mathbb{R} \to -\lvert a \rvert \leq a) \to \varphi_{31}$}
\BIC{$\varphi_{31}$}
\end{ptree}
\noindent$\Pi_{21}$:
\begin{ptree}
\AXC{$x = x$\lbl{A5}}
\AXC{$x = y \to x = x \to y = x$\lbl{A6}}
\AXC{$(x = y \to x = x \to y = x) \to x = x \to x = y \to y = x$\lbl{A1}}
\BIC{$x = x \to x = y \to y = x$}
\BIC{$x = y \to y = x$}
\UIC{$\forall y\,(x = y \to y = x)$}
\UIC{$\forall x\,\forall y\,(x = y \to y = x)$}
\AXC{$\varphi_{1}$\lbl{A2}}
\BIC{$\forall y\,(-(\lvert a \rvert + \lvert b \rvert) = y \to y = -(\lvert a \rvert + \lvert b \rvert))$}
\AXC{$\varphi_{42}$\lbl{A2}}
\BIC{$\varphi_{30}$}
\end{ptree}
\noindent$\Pi_{22}$:
\begin{ptree}
\AXC{$\forall x\,\forall y\,(x \in \mathbb{R} \to y \in \mathbb{R} \to -(x + y) = -x + -y)$\lbl{[F7]}}
\AXC{$\forall x\,\forall y\,(x \in \mathbb{R} \to y \in \mathbb{R} \to -(x + y) = -x + -y) \to \varphi_{2}$\lbl{A2}}
\BIC{$\varphi_{2}$}
\end{ptree}
\noindent$\Pi_{23}$:
\begin{ptree}
\AXC{$\stackrel{\Pi_{8}}{\vdots}$}\noLine
\UIC{$a \in \mathbb{R} \to b \in \mathbb{R} \to \lvert a \rvert \in \mathbb{R}$}
\AXC{$\stackrel{\Pi_{22}}{\vdots}$}\noLine
\UIC{$\varphi_{2}$}
\AXC{$\varphi_{2} \to \varphi_{12}$\lbl{A2}}
\BIC{$\varphi_{12}$}
\AXC{$\varphi_{12} \to (a \in \mathbb{R} \to b \in \mathbb{R} \to \lvert a \rvert \in \mathbb{R}) \to \varphi_{37}$\lbl{A1}}
\BIC{$(a \in \mathbb{R} \to b \in \mathbb{R} \to \lvert a \rvert \in \mathbb{R}) \to \varphi_{37}$}
\BIC{$\varphi_{37}$}
\end{ptree}
\noindent$\Pi_{24}$:
\begin{ptree}
\AXC{$\stackrel{\Pi_{21}}{\vdots}$}\noLine
\UIC{$\varphi_{30}$}
\AXC{$\stackrel{\Pi_{6}}{\vdots}$}\noLine
\UIC{$a \in \mathbb{R} \to b \in \mathbb{R} \to \lvert b \rvert \in \mathbb{R}$}
\AXC{$\stackrel{\Pi_{23}}{\vdots}$}\noLine
\UIC{$\varphi_{37}$}
\BIC{$a \in \mathbb{R} \to b \in \mathbb{R} \to -(\lvert a \rvert + \lvert b \rvert) = -\lvert a \rvert + -\lvert b \rvert$}
\AXC{$\varphi_{49}$\lbl{A1}}
\BIC{$\varphi_{44}$}
\BIC{$a \in \mathbb{R} \to b \in \mathbb{R} \to -\lvert a \rvert + -\lvert b \rvert = -(\lvert a \rvert + \lvert b \rvert)$}
\AXC{$\varphi_{54}$\lbl{A1}}
\BIC{$\varphi_{51}$}
\end{ptree}
\noindent$\Pi_{25}$:
\begin{ptree}
\AXC{$\stackrel{\Pi_{13}}{\vdots}$}\noLine
\UIC{$\varphi_{35}$}
\AXC{$\stackrel{\Pi_{14}}{\vdots}$}\noLine
\UIC{$\varphi_{39}$}
\AXC{$\stackrel{\Pi_{2}}{\vdots}$}\noLine
\UIC{$a \in \mathbb{R} \to b \in \mathbb{R} \to -\lvert b \rvert \leq b \land b \leq \lvert b \rvert$}
\AXC{$\varphi_{24}$\lbl{A1}}
\BIC{$a \in \mathbb{R} \to b \in \mathbb{R} \to -\lvert b \rvert \leq b$}
\AXC{$\stackrel{\Pi_{20}}{\vdots}$}\noLine
\UIC{$\varphi_{31}$}
\BIC{$a \in \mathbb{R} \to b \in \mathbb{R} \to -\lvert a \rvert + -\lvert b \rvert \leq a + b$}
\AXC{$\stackrel{\Pi_{24}}{\vdots}$}\noLine
\UIC{$\varphi_{51}$}
\BIC{$\varphi_{39} \to \varphi_{35} \to a \in \mathbb{R} \to b \in \mathbb{R} \to -(\lvert a \rvert + \lvert b \rvert) \leq a + b$}
\BIC{$\varphi_{35} \to a \in \mathbb{R} \to b \in \mathbb{R} \to -(\lvert a \rvert + \lvert b \rvert) \leq a + b$}
\BIC{$a \in \mathbb{R} \to b \in \mathbb{R} \to -(\lvert a \rvert + \lvert b \rvert) \leq a + b$}
\end{ptree}
\noindent$\Pi_{26}$:
\begin{ptree}
\AXC{$\stackrel{\Pi_{8}}{\vdots}$}\noLine
\UIC{$a \in \mathbb{R} \to b \in \mathbb{R} \to \lvert a \rvert \in \mathbb{R}$}
\AXC{$\forall x\,\forall y\,(x \in \mathbb{R} \to y \in \mathbb{R} \to x + y \in \mathbb{R})$\lbl{(R1)}}
\AXC{$\varphi_{32}$\lbl{A2}}
\BIC{$\forall y\,(\lvert a \rvert \in \mathbb{R} \to y \in \mathbb{R} \to \lvert a \rvert + y \in \mathbb{R})$}
\AXC{$\varphi_{38}$\lbl{A2}}
\BIC{$\lvert a \rvert \in \mathbb{R} \to \lvert b \rvert \in \mathbb{R} \to \lvert a \rvert + \lvert b \rvert \in \mathbb{R}$}
\AXC{$\varphi_{46}$\lbl{A1}}
\BIC{$(a \in \mathbb{R} \to b \in \mathbb{R} \to \lvert a \rvert \in \mathbb{R}) \to \varphi_{28}$}
\BIC{$\varphi_{28}$}
\end{ptree}
\noindent$\Pi_{27}$:
\begin{ptree}
\AXC{$\forall x\,\forall y\,(x \in \mathbb{R} \to y \in \mathbb{R} \to -y \leq x \to x \leq y \to \lvert x \rvert \leq y)$\lbl{[A10]}}
\AXC{$\varphi_{40}$\lbl{A2}}
\BIC{$\varphi_{3}$}
\AXC{$\varphi_{3} \to a + b \in \mathbb{R} \to \lvert a \rvert + \lvert b \rvert \in \mathbb{R} \to \varphi_{13}$\lbl{A2}}
\BIC{$a + b \in \mathbb{R} \to \lvert a \rvert + \lvert b \rvert \in \mathbb{R} \to \varphi_{13}$}
\end{ptree}
\noindent$\Pi_{28}$:
\begin{ptree}
\AXC{$\forall x\,\forall y\,(x \in \mathbb{R} \to y \in \mathbb{R} \to x + y \in \mathbb{R})$\lbl{(R1)}}
\AXC{$\varphi_{26}$\lbl{A2}}
\BIC{$\forall y\,(a \in \mathbb{R} \to y \in \mathbb{R} \to a + y \in \mathbb{R})$}
\AXC{$\varphi_{18}$\lbl{A2}}
\BIC{$a \in \mathbb{R} \to b \in \mathbb{R} \to a + b \in \mathbb{R}$}
\AXC{$\stackrel{\Pi_{27}}{\vdots}$}\noLine
\UIC{$a + b \in \mathbb{R} \to \lvert a \rvert + \lvert b \rvert \in \mathbb{R} \to \varphi_{13}$}
\AXC{$\varphi_{52}$\lbl{A1}}
\BIC{$(a \in \mathbb{R} \to b \in \mathbb{R} \to a + b \in \mathbb{R}) \to \varphi_{47}$}
\BIC{$\varphi_{47}$}
\end{ptree}
\noindent$\Pi_{29}$:
\begin{ptree}
\AXC{$\stackrel{\Pi_{25}}{\vdots}$}\noLine
\UIC{$a \in \mathbb{R} \to b \in \mathbb{R} \to -(\lvert a \rvert + \lvert b \rvert) \leq a + b$}
\AXC{$\stackrel{\Pi_{6}}{\vdots}$}\noLine
\UIC{$a \in \mathbb{R} \to b \in \mathbb{R} \to \lvert b \rvert \in \mathbb{R}$}
\AXC{$\stackrel{\Pi_{26}}{\vdots}$}\noLine
\UIC{$\varphi_{28}$}
\BIC{$a \in \mathbb{R} \to b \in \mathbb{R} \to \lvert a \rvert + \lvert b \rvert \in \mathbb{R}$}
\AXC{$\stackrel{\Pi_{28}}{\vdots}$}\noLine
\UIC{$\varphi_{47}$}
\BIC{$(a \in \mathbb{R} \to b \in \mathbb{R} \to -(\lvert a \rvert + \lvert b \rvert) \leq a + b) \to \varphi_{33}$}
\BIC{$\varphi_{33}$}
\end{ptree}
\noindent 主木:
\begin{ptree}
\AXC{$\stackrel{\Pi_{2}}{\vdots}$}\noLine
\UIC{$a \in \mathbb{R} \to b \in \mathbb{R} \to -\lvert b \rvert \leq b \land b \leq \lvert b \rvert$}
\AXC{$\varphi_{16}$\lbl{A1}}
\BIC{$a \in \mathbb{R} \to b \in \mathbb{R} \to b \leq \lvert b \rvert$}
\AXC{$\stackrel{\Pi_{12}}{\vdots}$}\noLine
\UIC{$\varphi_{15}$}
\BIC{$a \in \mathbb{R} \to b \in \mathbb{R} \to a + b \leq \lvert a \rvert + \lvert b \rvert$}
\AXC{$\stackrel{\Pi_{29}}{\vdots}$}\noLine
\UIC{$\varphi_{33}$}
\BIC{$a \in \mathbb{R} \to b \in \mathbb{R} \to \lvert a + b \rvert \leq \lvert a \rvert + \lvert b \rvert$}
\UIC{$\forall b\,(a \in \mathbb{R} \to b \in \mathbb{R} \to \lvert a + b \rvert \leq \lvert a \rvert + \lvert b \rvert)$}
\UIC{$\forall a\,\forall b\,(a \in \mathbb{R} \to b \in \mathbb{R} \to \lvert a + b \rvert \leq \lvert a \rvert + \lvert b \rvert)$}
\end{ptree}
\begin{abbrevlist}
\item $\varphi_{1} :\equiv \forall x\,\forall y\,(x = y \to y = x) \to \forall y\,(-(\lvert a \rvert + \lvert b \rvert) = y \to y = -(\lvert a \rvert + \lvert b \rvert))$
\item $\varphi_{2} :\equiv \forall y\,(\lvert a \rvert \in \mathbb{R} \to y \in \mathbb{R} \to -(\lvert a \rvert + y) = -\lvert a \rvert + -y)$
\item $\varphi_{3} :\equiv \forall y\,(a + b \in \mathbb{R} \to y \in \mathbb{R} \to -y \leq a + b \to a + b \leq y \to \lvert a + b \rvert \leq y)$
\item $\varphi_{4} :\equiv b \in \mathbb{R} \to \lvert b \rvert \in \mathbb{R} \to a \leq \lvert a \rvert \to b \leq \lvert b \rvert \to a + b \leq \lvert a \rvert + \lvert b \rvert$
\item $\varphi_{5} :\equiv \forall w\,(-\lvert a \rvert + -\lvert b \rvert = w \to -\lvert a \rvert + -\lvert b \rvert = a + b \to w = a + b)$
\item $\varphi_{6} :\equiv b \in \mathbb{R} \to -\lvert a \rvert \leq a \to -\lvert b \rvert \leq b \to -\lvert a \rvert + -\lvert b \rvert \leq a + b$
\item $\varphi_{7} :\equiv \forall x \in \mathbb{R}\,(-\lvert x \rvert \leq x \land x \leq \lvert x \rvert) \to b \in \mathbb{R} \to -\lvert b \rvert \leq b \land b \leq \lvert b \rvert$
\item $\varphi_{8} :\equiv \forall x \in \mathbb{R}\,(-\lvert x \rvert \leq x \land x \leq \lvert x \rvert) \to a \in \mathbb{R} \to -\lvert a \rvert \leq a \land a \leq \lvert a \rvert$
\item $\varphi_{9} :\equiv (a \in \mathbb{R} \to b \in \mathbb{R} \to b \in \mathbb{R}) \to a \in \mathbb{R} \to b \in \mathbb{R} \to \lvert b \rvert \in \mathbb{R}$
\item $\varphi_{10} :\equiv (a \in \mathbb{R} \to b \in \mathbb{R} \to a \in \mathbb{R}) \to a \in \mathbb{R} \to b \in \mathbb{R} \to \lvert a \rvert \in \mathbb{R}$
\item $\varphi_{11} :\equiv \forall w\,(-\lvert a \rvert + -\lvert b \rvert = w \to -\lvert a \rvert + -\lvert b \rvert < a + b \to w < a + b)$
\item $\varphi_{12} :\equiv \lvert a \rvert \in \mathbb{R} \to \lvert b \rvert \in \mathbb{R} \to -(\lvert a \rvert + \lvert b \rvert) = -\lvert a \rvert + -\lvert b \rvert$
\item $\varphi_{13} :\equiv -(\lvert a \rvert + \lvert b \rvert) \leq a + b \to a + b \leq \lvert a \rvert + \lvert b \rvert \to \lvert a + b \rvert \leq \lvert a \rvert + \lvert b \rvert$
\item $\varphi_{14} :\equiv \forall y\,\forall u\,\forall v\,(a \in \mathbb{R} \to y \in \mathbb{R} \to u \in \mathbb{R} \to v \in \mathbb{R} \to a \leq y \to u \leq v \to a + u \leq y + v)$
\item $\varphi_{15} :\equiv (a \in \mathbb{R} \to b \in \mathbb{R} \to b \leq \lvert b \rvert) \to a \in \mathbb{R} \to b \in \mathbb{R} \to a + b \leq \lvert a \rvert + \lvert b \rvert$
\item $\varphi_{16} :\equiv (a \in \mathbb{R} \to b \in \mathbb{R} \to -\lvert b \rvert \leq b \land b \leq \lvert b \rvert) \to a \in \mathbb{R} \to b \in \mathbb{R} \to b \leq \lvert b \rvert$
\item $\varphi_{17} :\equiv (a \in \mathbb{R} \to b \in \mathbb{R} \to -\lvert a \rvert \leq a \land a \leq \lvert a \rvert) \to a \in \mathbb{R} \to b \in \mathbb{R} \to a \leq \lvert a \rvert$
\item $\varphi_{18} :\equiv \forall y\,(a \in \mathbb{R} \to y \in \mathbb{R} \to a + y \in \mathbb{R}) \to a \in \mathbb{R} \to b \in \mathbb{R} \to a + b \in \mathbb{R}$
\item $\varphi_{19} :\equiv (a \in \mathbb{R} \to b \in \mathbb{R} \to b \in \mathbb{R}) \to a \in \mathbb{R} \to b \in \mathbb{R} \to -\lvert b \rvert \leq b \land b \leq \lvert b \rvert$
\item $\varphi_{20} :\equiv (a \in \mathbb{R} \to b \in \mathbb{R} \to a \in \mathbb{R}) \to a \in \mathbb{R} \to b \in \mathbb{R} \to -\lvert a \rvert \leq a \land a \leq \lvert a \rvert$
\item $\varphi_{21} :\equiv \forall x\,\forall y\,\forall u\,\forall v\,(x \in \mathbb{R} \to y \in \mathbb{R} \to u \in \mathbb{R} \to v \in \mathbb{R} \to x \leq y \to u \leq v \to x + u \leq y + v)$
\item $\varphi_{22} :\equiv (a \in \mathbb{R} \to b \in \mathbb{R} \to \lvert b \rvert \in \mathbb{R}) \to a \in \mathbb{R} \to b \in \mathbb{R} \to -\lvert b \rvert \in \mathbb{R}$
\item $\varphi_{23} :\equiv (a \in \mathbb{R} \to b \in \mathbb{R} \to \lvert a \rvert \in \mathbb{R}) \to a \in \mathbb{R} \to b \in \mathbb{R} \to -\lvert a \rvert \in \mathbb{R}$
\item $\varphi_{24} :\equiv (a \in \mathbb{R} \to b \in \mathbb{R} \to -\lvert b \rvert \leq b \land b \leq \lvert b \rvert) \to a \in \mathbb{R} \to b \in \mathbb{R} \to -\lvert b \rvert \leq b$
\item $\varphi_{25} :\equiv (a \in \mathbb{R} \to b \in \mathbb{R} \to -\lvert a \rvert \leq a \land a \leq \lvert a \rvert) \to a \in \mathbb{R} \to b \in \mathbb{R} \to -\lvert a \rvert \leq a$
\item $\varphi_{26} :\equiv \forall x\,\forall y\,(x \in \mathbb{R} \to y \in \mathbb{R} \to x + y \in \mathbb{R}) \to \forall y\,(a \in \mathbb{R} \to y \in \mathbb{R} \to a + y \in \mathbb{R})$
\item $\varphi_{27} :\equiv \forall v\,(a \in \mathbb{R} \to \lvert a \rvert \in \mathbb{R} \to b \in \mathbb{R} \to v \in \mathbb{R} \to a \leq \lvert a \rvert \to b \leq v \to a + b \leq \lvert a \rvert + v)$
\item $\varphi_{28} :\equiv (a \in \mathbb{R} \to b \in \mathbb{R} \to \lvert b \rvert \in \mathbb{R}) \to a \in \mathbb{R} \to b \in \mathbb{R} \to \lvert a \rvert + \lvert b \rvert \in \mathbb{R}$
\item $\varphi_{29} :\equiv \forall u\,\forall v\,(a \in \mathbb{R} \to \lvert a \rvert \in \mathbb{R} \to u \in \mathbb{R} \to v \in \mathbb{R} \to a \leq \lvert a \rvert \to u \leq v \to a + u \leq \lvert a \rvert + v)$
\item $\varphi_{30} :\equiv -(\lvert a \rvert + \lvert b \rvert) = -\lvert a \rvert + -\lvert b \rvert \to -\lvert a \rvert + -\lvert b \rvert = -(\lvert a \rvert + \lvert b \rvert)$
\item $\varphi_{31} :\equiv (a \in \mathbb{R} \to b \in \mathbb{R} \to -\lvert b \rvert \leq b) \to a \in \mathbb{R} \to b \in \mathbb{R} \to -\lvert a \rvert + -\lvert b \rvert \leq a + b$
\item $\varphi_{32} :\equiv \forall x\,\forall y\,(x \in \mathbb{R} \to y \in \mathbb{R} \to x + y \in \mathbb{R}) \to \forall y\,(\lvert a \rvert \in \mathbb{R} \to y \in \mathbb{R} \to \lvert a \rvert + y \in \mathbb{R})$
\item $\varphi_{33} :\equiv (a \in \mathbb{R} \to b \in \mathbb{R} \to a + b \leq \lvert a \rvert + \lvert b \rvert) \to a \in \mathbb{R} \to b \in \mathbb{R} \to \lvert a + b \rvert \leq \lvert a \rvert + \lvert b \rvert$
\item $\varphi_{34} :\equiv \forall u\,\forall v\,(-\lvert a \rvert \in \mathbb{R} \to a \in \mathbb{R} \to u \in \mathbb{R} \to v \in \mathbb{R} \to -\lvert a \rvert \leq a \to u \leq v \to -\lvert a \rvert + u \leq a + v)$
\item $\varphi_{35} :\equiv -\lvert a \rvert + -\lvert b \rvert = -(\lvert a \rvert + \lvert b \rvert) \to -\lvert a \rvert + -\lvert b \rvert = a + b \to -(\lvert a \rvert + \lvert b \rvert) = a + b$
\item $\varphi_{36} :\equiv \forall y\,\forall u\,\forall v\,(-\lvert a \rvert \in \mathbb{R} \to y \in \mathbb{R} \to u \in \mathbb{R} \to v \in \mathbb{R} \to -\lvert a \rvert \leq y \to u \leq v \to -\lvert a \rvert + u \leq y + v)$
\item $\varphi_{37} :\equiv (a \in \mathbb{R} \to b \in \mathbb{R} \to \lvert b \rvert \in \mathbb{R}) \to a \in \mathbb{R} \to b \in \mathbb{R} \to -(\lvert a \rvert + \lvert b \rvert) = -\lvert a \rvert + -\lvert b \rvert$
\item $\varphi_{38} :\equiv \forall y\,(\lvert a \rvert \in \mathbb{R} \to y \in \mathbb{R} \to \lvert a \rvert + y \in \mathbb{R}) \to \lvert a \rvert \in \mathbb{R} \to \lvert b \rvert \in \mathbb{R} \to \lvert a \rvert + \lvert b \rvert \in \mathbb{R}$
\item $\varphi_{39} :\equiv -\lvert a \rvert + -\lvert b \rvert = -(\lvert a \rvert + \lvert b \rvert) \to -\lvert a \rvert + -\lvert b \rvert < a + b \to -(\lvert a \rvert + \lvert b \rvert) < a + b$
\item $\varphi_{40} :\equiv \forall x\,\forall y\,(x \in \mathbb{R} \to y \in \mathbb{R} \to -y \leq x \to x \leq y \to \lvert x \rvert \leq y) \to \varphi_{3}$
\item $\varphi_{41} :\equiv \forall v\,(-\lvert a \rvert \in \mathbb{R} \to a \in \mathbb{R} \to -\lvert b \rvert \in \mathbb{R} \to v \in \mathbb{R} \to -\lvert a \rvert \leq a \to -\lvert b \rvert \leq v \to -\lvert a \rvert + -\lvert b \rvert \leq a + v)$
\item $\varphi_{42} :\equiv \forall y\,(-(\lvert a \rvert + \lvert b \rvert) = y \to y = -(\lvert a \rvert + \lvert b \rvert)) \to \varphi_{30}$
\item $\varphi_{43} :\equiv (a \in \mathbb{R} \to b \in \mathbb{R} \to \lvert b \rvert \in \mathbb{R}) \to (a \in \mathbb{R} \to b \in \mathbb{R} \to a \leq \lvert a \rvert) \to \varphi_{15}$
\item $\varphi_{44} :\equiv \varphi_{30} \to a \in \mathbb{R} \to b \in \mathbb{R} \to -\lvert a \rvert + -\lvert b \rvert = -(\lvert a \rvert + \lvert b \rvert)$
\item $\varphi_{45} :\equiv (a \in \mathbb{R} \to b \in \mathbb{R} \to b \in \mathbb{R}) \to (a \in \mathbb{R} \to b \in \mathbb{R} \to -\lvert a \rvert \leq a) \to \varphi_{31}$
\item $\varphi_{46} :\equiv (\lvert a \rvert \in \mathbb{R} \to \lvert b \rvert \in \mathbb{R} \to \lvert a \rvert + \lvert b \rvert \in \mathbb{R}) \to (a \in \mathbb{R} \to b \in \mathbb{R} \to \lvert a \rvert \in \mathbb{R}) \to \varphi_{28}$
\item $\varphi_{47} :\equiv (a \in \mathbb{R} \to b \in \mathbb{R} \to \lvert a \rvert + \lvert b \rvert \in \mathbb{R}) \to (a \in \mathbb{R} \to b \in \mathbb{R} \to -(\lvert a \rvert + \lvert b \rvert) \leq a + b) \to \varphi_{33}$
\item $\varphi_{48} :\equiv (a \in \mathbb{R} \to b \in \mathbb{R} \to \lvert a \rvert \in \mathbb{R}) \to (a \in \mathbb{R} \to b \in \mathbb{R} \to b \in \mathbb{R}) \to \varphi_{43}$
\item $\varphi_{49} :\equiv (a \in \mathbb{R} \to b \in \mathbb{R} \to -(\lvert a \rvert + \lvert b \rvert) = -\lvert a \rvert + -\lvert b \rvert) \to \varphi_{44}$
\item $\varphi_{50} :\equiv (a \in \mathbb{R} \to b \in \mathbb{R} \to a \in \mathbb{R}) \to (a \in \mathbb{R} \to b \in \mathbb{R} \to -\lvert b \rvert \in \mathbb{R}) \to \varphi_{45}$
\item $\varphi_{51} :\equiv (a \in \mathbb{R} \to b \in \mathbb{R} \to -\lvert a \rvert + -\lvert b \rvert \leq a + b) \to \varphi_{39} \to \varphi_{35} \to a \in \mathbb{R} \to b \in \mathbb{R} \to -(\lvert a \rvert + \lvert b \rvert) \leq a + b$
\item $\varphi_{52} :\equiv (a + b \in \mathbb{R} \to \lvert a \rvert + \lvert b \rvert \in \mathbb{R} \to \varphi_{13}) \to (a \in \mathbb{R} \to b \in \mathbb{R} \to a + b \in \mathbb{R}) \to \varphi_{47}$
\item $\varphi_{53} :\equiv (a \in \mathbb{R} \to \lvert a \rvert \in \mathbb{R} \to \varphi_{4}) \to (a \in \mathbb{R} \to b \in \mathbb{R} \to a \in \mathbb{R}) \to \varphi_{48}$
\item $\varphi_{54} :\equiv (a \in \mathbb{R} \to b \in \mathbb{R} \to -\lvert a \rvert + -\lvert b \rvert = -(\lvert a \rvert + \lvert b \rvert)) \to \varphi_{51}$
\item $\varphi_{55} :\equiv (-\lvert a \rvert \in \mathbb{R} \to a \in \mathbb{R} \to -\lvert b \rvert \in \mathbb{R} \to \varphi_{6}) \to (a \in \mathbb{R} \to b \in \mathbb{R} \to -\lvert a \rvert \in \mathbb{R}) \to \varphi_{50}$
\end{abbrevlist}
\end{fullproof}

% 自動生成 (prover/build.js): AbsPos  ノード数 101
\begin{fullproof}{[A14]}{$\forall x \in \mathbb{R}\,(x \neq 0 \to 0 < \lvert x \rvert)$}{101}\phantomsection\label{tree:AbsPos}
\noindent$\Pi_{1}$:
\begin{ptree}
\AXC{$x = x$\lbl{A5}}
\AXC{$x = y \to x = x \to y = x$\lbl{A6}}
\AXC{$(x = y \to x = x \to y = x) \to x = x \to x = y \to y = x$\lbl{A1}}
\BIC{$x = x \to x = y \to y = x$}
\BIC{$x = y \to y = x$}
\UIC{$\forall y\,(x = y \to y = x)$}
\UIC{$\forall x\,\forall y\,(x = y \to y = x)$}
\end{ptree}
\noindent$\Pi_{2}$:
\begin{ptree}
\AXC{$\stackrel{\Pi_{1}}{\vdots}$}\noLine
\UIC{$\forall x\,\forall y\,(x = y \to y = x)$}
\AXC{$\forall x\,\forall y\,(x = y \to y = x) \to \forall y\,(\lvert x \rvert = y \to y = \lvert x \rvert)$\lbl{A2}}
\BIC{$\forall y\,(\lvert x \rvert = y \to y = \lvert x \rvert)$}
\end{ptree}
\noindent$\Pi_{3}$:
\begin{ptree}
\AXC{$\forall x\,(0 \leq x \to \lvert x \rvert = x)$\lbl{[A1]}}
\AXC{$\forall x\,(0 \leq x \to \lvert x \rvert = x) \to 0 \leq x \to \lvert x \rvert = x$\lbl{A2}}
\BIC{$0 \leq x \to \lvert x \rvert = x$}
\AXC{$(0 \leq x \to \lvert x \rvert = x) \to \varphi_{6}$\lbl{A1}}
\BIC{$\varphi_{6}$}
\end{ptree}
\noindent$\Pi_{4}$:
\begin{ptree}
\AXC{$x \in \mathbb{R} \to x \neq 0 \to 0 < x \to 0 < x$\lbl{A1}}
\AXC{$\varphi_{5}$\lbl{A1}}
\BIC{$x \in \mathbb{R} \to x \neq 0 \to 0 < x \to 0 \leq x$}
\AXC{$\stackrel{\Pi_{3}}{\vdots}$}\noLine
\UIC{$\varphi_{6}$}
\BIC{$x \in \mathbb{R} \to x \neq 0 \to 0 < x \to \lvert x \rvert = x$}
\AXC{$(x \in \mathbb{R} \to x \neq 0 \to 0 < x \to \lvert x \rvert = x) \to \varphi_{3}$\lbl{A1}}
\BIC{$\varphi_{3}$}
\end{ptree}
\noindent$\Pi_{5}$:
\begin{ptree}
\AXC{$\stackrel{\Pi_{2}}{\vdots}$}\noLine
\UIC{$\forall y\,(\lvert x \rvert = y \to y = \lvert x \rvert)$}
\AXC{$\forall y\,(\lvert x \rvert = y \to y = \lvert x \rvert) \to \lvert x \rvert = x \to x = \lvert x \rvert$\lbl{A2}}
\BIC{$\lvert x \rvert = x \to x = \lvert x \rvert$}
\AXC{$\stackrel{\Pi_{4}}{\vdots}$}\noLine
\UIC{$\varphi_{3}$}
\BIC{$x \in \mathbb{R} \to x \neq 0 \to 0 < x \to x = \lvert x \rvert$}
\AXC{$\varphi_{23}$\lbl{A1}}
\BIC{$(x \in \mathbb{R} \to x \neq 0 \to 0 < x \to 0 < x) \to \varphi_{7}$}
\end{ptree}
\noindent$\Pi_{6}$:
\begin{ptree}
\AXC{$x \in \mathbb{R} \to x \neq 0 \to 0 < x \to 0 < x$\lbl{A1}}
\AXC{$\stackrel{\Pi_{5}}{\vdots}$}\noLine
\UIC{$(x \in \mathbb{R} \to x \neq 0 \to 0 < x \to 0 < x) \to \varphi_{7}$}
\BIC{$\varphi_{7}$}
\end{ptree}
\noindent$\Pi_{7}$:
\begin{ptree}
\AXC{$\forall x\,\forall y\,(x \in \mathbb{R} \to y \in \mathbb{R} \to x < y \to -y < -x)$\lbl{[F22]}}
\AXC{$\varphi_{21}$\lbl{A2}}
\BIC{$\forall y\,(x \in \mathbb{R} \to y \in \mathbb{R} \to x < y \to -y < -x)$}
\AXC{$\varphi_{20}$\lbl{A2}}
\BIC{$x \in \mathbb{R} \to 0 \in \mathbb{R} \to x < 0 \to -0 < -x$}
\AXC{$\varphi_{25}$\lbl{A1}}
\BIC{$(x \in \mathbb{R} \to x \neq 0 \to x < 0 \to x \in \mathbb{R}) \to 0 \in \mathbb{R} \to \varphi_{12}$}
\end{ptree}
\noindent$\Pi_{8}$:
\begin{ptree}
\AXC{$0 \in \mathbb{R}$\lbl{(R1)}}
\AXC{$x \in \mathbb{R} \to x \neq 0 \to x < 0 \to x \in \mathbb{R}$\lbl{A1}}
\AXC{$\stackrel{\Pi_{7}}{\vdots}$}\noLine
\UIC{$(x \in \mathbb{R} \to x \neq 0 \to x < 0 \to x \in \mathbb{R}) \to 0 \in \mathbb{R} \to \varphi_{12}$}
\BIC{$0 \in \mathbb{R} \to \varphi_{12}$}
\BIC{$\varphi_{12}$}
\end{ptree}
\noindent$\Pi_{9}$:
\begin{ptree}
\AXC{$x \in \mathbb{R} \to x \neq 0 \to x < 0 \to x < 0$\lbl{A1}}
\AXC{$\stackrel{\Pi_{8}}{\vdots}$}\noLine
\UIC{$\varphi_{12}$}
\BIC{$x \in \mathbb{R} \to x \neq 0 \to x < 0 \to -0 < -x$}
\AXC{$-0 = 0$\lbl{[F6]}}
\AXC{$-0 = 0 \to (x \in \mathbb{R} \to x \neq 0 \to x < 0 \to -0 < -x) \to \varphi_{14}$\lbl{A1}}
\BIC{$(x \in \mathbb{R} \to x \neq 0 \to x < 0 \to -0 < -x) \to \varphi_{14}$}
\BIC{$\varphi_{14}$}
\end{ptree}
\noindent$\Pi_{10}$:
\begin{ptree}
\AXC{$\stackrel{\Pi_{1}}{\vdots}$}\noLine
\UIC{$\forall x\,\forall y\,(x = y \to y = x)$}
\AXC{$\forall x\,\forall y\,(x = y \to y = x) \to \forall y\,(\lvert x \rvert = y \to y = \lvert x \rvert)$\lbl{A2}}
\BIC{$\forall y\,(\lvert x \rvert = y \to y = \lvert x \rvert)$}
\end{ptree}
\noindent$\Pi_{11}$:
\begin{ptree}
\AXC{$\forall x\,(x < 0 \to \lvert x \rvert = -x)$\lbl{[A3]}}
\AXC{$\forall x\,(x < 0 \to \lvert x \rvert = -x) \to x < 0 \to \lvert x \rvert = -x$\lbl{A2}}
\BIC{$x < 0 \to \lvert x \rvert = -x$}
\AXC{$(x < 0 \to \lvert x \rvert = -x) \to \varphi_{11}$\lbl{A1}}
\BIC{$\varphi_{11}$}
\end{ptree}
\noindent$\Pi_{12}$:
\begin{ptree}
\AXC{$x \in \mathbb{R} \to x \neq 0 \to x < 0 \to x < 0$\lbl{A1}}
\AXC{$\stackrel{\Pi_{11}}{\vdots}$}\noLine
\UIC{$\varphi_{11}$}
\BIC{$x \in \mathbb{R} \to x \neq 0 \to x < 0 \to \lvert x \rvert = -x$}
\AXC{$(x \in \mathbb{R} \to x \neq 0 \to x < 0 \to \lvert x \rvert = -x) \to \varphi_{8}$\lbl{A1}}
\BIC{$\varphi_{8}$}
\end{ptree}
\noindent$\Pi_{13}$:
\begin{ptree}
\AXC{$\stackrel{\Pi_{10}}{\vdots}$}\noLine
\UIC{$\forall y\,(\lvert x \rvert = y \to y = \lvert x \rvert)$}
\AXC{$\forall y\,(\lvert x \rvert = y \to y = \lvert x \rvert) \to \lvert x \rvert = -x \to -x = \lvert x \rvert$\lbl{A2}}
\BIC{$\lvert x \rvert = -x \to -x = \lvert x \rvert$}
\AXC{$\stackrel{\Pi_{12}}{\vdots}$}\noLine
\UIC{$\varphi_{8}$}
\BIC{$x \in \mathbb{R} \to x \neq 0 \to x < 0 \to -x = \lvert x \rvert$}
\AXC{$\varphi_{24}$\lbl{A1}}
\BIC{$(x \in \mathbb{R} \to x \neq 0 \to x < 0 \to 0 < -x) \to \varphi_{15}$}
\end{ptree}
\noindent$\Pi_{14}$:
\begin{ptree}
\AXC{$z = w \to z < -x \to w < -x$\lbl{A6}}
\UIC{$\forall w\,(z = w \to z < -x \to w < -x)$}
\UIC{$\forall z\,\forall w\,(z = w \to z < -x \to w < -x)$}
\AXC{$\varphi_{13}$\lbl{A2}}
\BIC{$\forall w\,(-0 = w \to -0 < -x \to w < -x)$}
\AXC{$\varphi_{17}$\lbl{A2}}
\BIC{$-0 = 0 \to -0 < -x \to 0 < -x$}
\AXC{$\stackrel{\Pi_{9}}{\vdots}$}\noLine
\UIC{$\varphi_{14}$}
\BIC{$x \in \mathbb{R} \to x \neq 0 \to x < 0 \to 0 < -x$}
\AXC{$\stackrel{\Pi_{13}}{\vdots}$}\noLine
\UIC{$(x \in \mathbb{R} \to x \neq 0 \to x < 0 \to 0 < -x) \to \varphi_{15}$}
\BIC{$\varphi_{15}$}
\end{ptree}
\noindent$\Pi_{15}$:
\begin{ptree}
\AXC{$\forall x\,\forall y\,(x \in \mathbb{R} \to y \in \mathbb{R} \to x < y \lor x = y \lor y < x)$\lbl{(R13)}}
\AXC{$\varphi_{19}$\lbl{A2}}
\BIC{$\forall y\,(x \in \mathbb{R} \to y \in \mathbb{R} \to x < y \lor x = y \lor y < x)$}
\AXC{$\varphi_{18}$\lbl{A2}}
\BIC{$x \in \mathbb{R} \to 0 \in \mathbb{R} \to x < 0 \lor x = 0 \lor 0 < x$}
\end{ptree}
\noindent$\Pi_{16}$:
\begin{ptree}
\AXC{$x \in \mathbb{R} \to x \neq 0 \to x \in \mathbb{R}$\lbl{A1}}
\AXC{$\stackrel{\Pi_{15}}{\vdots}$}\noLine
\UIC{$x \in \mathbb{R} \to 0 \in \mathbb{R} \to x < 0 \lor x = 0 \lor 0 < x$}
\AXC{$(x \in \mathbb{R} \to 0 \in \mathbb{R} \to x < 0 \lor x = 0 \lor 0 < x) \to \varphi_{16}$\lbl{A1}}
\BIC{$\varphi_{16}$}
\BIC{$0 \in \mathbb{R} \to x \in \mathbb{R} \to x \neq 0 \to x < 0 \lor x = 0 \lor 0 < x$}
\end{ptree}
\noindent$\Pi_{17}$:
\begin{ptree}
\AXC{$0 \in \mathbb{R}$\lbl{(R1)}}
\AXC{$\stackrel{\Pi_{16}}{\vdots}$}\noLine
\UIC{$0 \in \mathbb{R} \to x \in \mathbb{R} \to x \neq 0 \to x < 0 \lor x = 0 \lor 0 < x$}
\BIC{$x \in \mathbb{R} \to x \neq 0 \to x < 0 \lor x = 0 \lor 0 < x$}
\AXC{$(x \in \mathbb{R} \to x \neq 0 \to x < 0 \lor x = 0 \lor 0 < x) \to \varphi_{22}$\lbl{A1}}
\BIC{$\varphi_{22}$}
\end{ptree}
\noindent$\Pi_{18}$:
\begin{ptree}
\AXC{$z = w \to 0 < z \to 0 < w$\lbl{A6}}
\UIC{$\forall w\,(z = w \to 0 < z \to 0 < w)$}
\UIC{$\forall z\,\forall w\,(z = w \to 0 < z \to 0 < w)$}
\AXC{$\varphi_{4}$\lbl{A2}}
\BIC{$\forall w\,(-x = w \to 0 < -x \to 0 < w)$}
\AXC{$\varphi_{10}$\lbl{A2}}
\BIC{$-x = \lvert x \rvert \to 0 < -x \to 0 < \lvert x \rvert$}
\AXC{$\stackrel{\Pi_{14}}{\vdots}$}\noLine
\UIC{$\varphi_{15}$}
\BIC{$x \in \mathbb{R} \to x \neq 0 \to x < 0 \to 0 < \lvert x \rvert$}
\AXC{$x \in \mathbb{R} \to x \neq 0 \to x \neq 0$\lbl{A1}}
\AXC{$\stackrel{\Pi_{17}}{\vdots}$}\noLine
\UIC{$\varphi_{22}$}
\BIC{$(x \in \mathbb{R} \to x \neq 0 \to x < 0 \to 0 < \lvert x \rvert) \to \varphi_{9}$}
\BIC{$\varphi_{9}$}
\end{ptree}
\noindent 主木:
\begin{ptree}
\AXC{$z = w \to 0 < z \to 0 < w$\lbl{A6}}
\UIC{$\forall w\,(z = w \to 0 < z \to 0 < w)$}
\UIC{$\forall z\,\forall w\,(z = w \to 0 < z \to 0 < w)$}
\AXC{$\varphi_{1}$\lbl{A2}}
\BIC{$\forall w\,(x = w \to 0 < x \to 0 < w)$}
\AXC{$\varphi_{2}$\lbl{A2}}
\BIC{$x = \lvert x \rvert \to 0 < x \to 0 < \lvert x \rvert$}
\AXC{$\stackrel{\Pi_{6}}{\vdots}$}\noLine
\UIC{$\varphi_{7}$}
\BIC{$x \in \mathbb{R} \to x \neq 0 \to 0 < x \to 0 < \lvert x \rvert$}
\AXC{$\stackrel{\Pi_{18}}{\vdots}$}\noLine
\UIC{$\varphi_{9}$}
\BIC{$x \in \mathbb{R} \to x \neq 0 \to 0 < \lvert x \rvert$}
\UIC{$\forall x \in \mathbb{R}\,(x \neq 0 \to 0 < \lvert x \rvert)$}
\end{ptree}
\begin{abbrevlist}
\item $\varphi_{1} :\equiv \forall z\,\forall w\,(z = w \to 0 < z \to 0 < w) \to \forall w\,(x = w \to 0 < x \to 0 < w)$
\item $\varphi_{2} :\equiv \forall w\,(x = w \to 0 < x \to 0 < w) \to x = \lvert x \rvert \to 0 < x \to 0 < \lvert x \rvert$
\item $\varphi_{3} :\equiv (\lvert x \rvert = x \to x = \lvert x \rvert) \to x \in \mathbb{R} \to x \neq 0 \to 0 < x \to x = \lvert x \rvert$
\item $\varphi_{4} :\equiv \forall z\,\forall w\,(z = w \to 0 < z \to 0 < w) \to \forall w\,(-x = w \to 0 < -x \to 0 < w)$
\item $\varphi_{5} :\equiv (x \in \mathbb{R} \to x \neq 0 \to 0 < x \to 0 < x) \to x \in \mathbb{R} \to x \neq 0 \to 0 < x \to 0 \leq x$
\item $\varphi_{6} :\equiv (x \in \mathbb{R} \to x \neq 0 \to 0 < x \to 0 \leq x) \to x \in \mathbb{R} \to x \neq 0 \to 0 < x \to \lvert x \rvert = x$
\item $\varphi_{7} :\equiv (x = \lvert x \rvert \to 0 < x \to 0 < \lvert x \rvert) \to x \in \mathbb{R} \to x \neq 0 \to 0 < x \to 0 < \lvert x \rvert$
\item $\varphi_{8} :\equiv (\lvert x \rvert = -x \to -x = \lvert x \rvert) \to x \in \mathbb{R} \to x \neq 0 \to x < 0 \to -x = \lvert x \rvert$
\item $\varphi_{9} :\equiv (x \in \mathbb{R} \to x \neq 0 \to 0 < x \to 0 < \lvert x \rvert) \to x \in \mathbb{R} \to x \neq 0 \to 0 < \lvert x \rvert$
\item $\varphi_{10} :\equiv \forall w\,(-x = w \to 0 < -x \to 0 < w) \to -x = \lvert x \rvert \to 0 < -x \to 0 < \lvert x \rvert$
\item $\varphi_{11} :\equiv (x \in \mathbb{R} \to x \neq 0 \to x < 0 \to x < 0) \to x \in \mathbb{R} \to x \neq 0 \to x < 0 \to \lvert x \rvert = -x$
\item $\varphi_{12} :\equiv (x \in \mathbb{R} \to x \neq 0 \to x < 0 \to x < 0) \to x \in \mathbb{R} \to x \neq 0 \to x < 0 \to -0 < -x$
\item $\varphi_{13} :\equiv \forall z\,\forall w\,(z = w \to z < -x \to w < -x) \to \forall w\,(-0 = w \to -0 < -x \to w < -x)$
\item $\varphi_{14} :\equiv (-0 = 0 \to -0 < -x \to 0 < -x) \to x \in \mathbb{R} \to x \neq 0 \to x < 0 \to 0 < -x$
\item $\varphi_{15} :\equiv (-x = \lvert x \rvert \to 0 < -x \to 0 < \lvert x \rvert) \to x \in \mathbb{R} \to x \neq 0 \to x < 0 \to 0 < \lvert x \rvert$
\item $\varphi_{16} :\equiv (x \in \mathbb{R} \to x \neq 0 \to x \in \mathbb{R}) \to 0 \in \mathbb{R} \to x \in \mathbb{R} \to x \neq 0 \to x < 0 \lor x = 0 \lor 0 < x$
\item $\varphi_{17} :\equiv \forall w\,(-0 = w \to -0 < -x \to w < -x) \to -0 = 0 \to -0 < -x \to 0 < -x$
\item $\varphi_{18} :\equiv \forall y\,(x \in \mathbb{R} \to y \in \mathbb{R} \to x < y \lor x = y \lor y < x) \to x \in \mathbb{R} \to 0 \in \mathbb{R} \to x < 0 \lor x = 0 \lor 0 < x$
\item $\varphi_{19} :\equiv \forall x\,\forall y\,(x \in \mathbb{R} \to y \in \mathbb{R} \to x < y \lor x = y \lor y < x) \to \forall y\,(x \in \mathbb{R} \to y \in \mathbb{R} \to x < y \lor x = y \lor y < x)$
\item $\varphi_{20} :\equiv \forall y\,(x \in \mathbb{R} \to y \in \mathbb{R} \to x < y \to -y < -x) \to x \in \mathbb{R} \to 0 \in \mathbb{R} \to x < 0 \to -0 < -x$
\item $\varphi_{21} :\equiv \forall x\,\forall y\,(x \in \mathbb{R} \to y \in \mathbb{R} \to x < y \to -y < -x) \to \forall y\,(x \in \mathbb{R} \to y \in \mathbb{R} \to x < y \to -y < -x)$
\item $\varphi_{22} :\equiv (x \in \mathbb{R} \to x \neq 0 \to x \neq 0) \to (x \in \mathbb{R} \to x \neq 0 \to x < 0 \to 0 < \lvert x \rvert) \to \varphi_{9}$
\item $\varphi_{23} :\equiv (x \in \mathbb{R} \to x \neq 0 \to 0 < x \to x = \lvert x \rvert) \to (x \in \mathbb{R} \to x \neq 0 \to 0 < x \to 0 < x) \to \varphi_{7}$
\item $\varphi_{24} :\equiv (x \in \mathbb{R} \to x \neq 0 \to x < 0 \to -x = \lvert x \rvert) \to (x \in \mathbb{R} \to x \neq 0 \to x < 0 \to 0 < -x) \to \varphi_{15}$
\item $\varphi_{25} :\equiv (x \in \mathbb{R} \to 0 \in \mathbb{R} \to x < 0 \to -0 < -x) \to (x \in \mathbb{R} \to x \neq 0 \to x < 0 \to x \in \mathbb{R}) \to 0 \in \mathbb{R} \to \varphi_{12}$
\end{abbrevlist}
\end{fullproof}

% 自動生成 (prover/build.js): SubNeZero  ノード数 107
\begin{fullproof}{[A15]}{$\forall x\,\forall y\,(x \in \mathbb{R} \to y \in \mathbb{R} \to x \neq y \to x - y \neq 0)$}{107}\phantomsection\label{tree:SubNeZero}
\noindent$\Pi_{1}$:
\begin{ptree}
\AXC{$y = z \to x = y \to x = z$\lbl{A6}}
\AXC{$(y = z \to x = y \to x = z) \to x = y \to y = z \to x = z$\lbl{A1}}
\BIC{$x = y \to y = z \to x = z$}
\UIC{$\forall z\,(x = y \to y = z \to x = z)$}
\UIC{$\forall y\,\forall z\,(x = y \to y = z \to x = z)$}
\UIC{$\forall x\,\forall y\,\forall z\,(x = y \to y = z \to x = z)$}
\AXC{$\varphi_{2}$\lbl{A2}}
\BIC{$\forall y\,\forall z\,(x = y \to y = z \to x = z)$}
\AXC{$\varphi_{3}$\lbl{A2}}
\BIC{$\forall z\,(x = y + 0 \to y + 0 = z \to x = z)$}
\AXC{$\varphi_{5}$\lbl{A2}}
\BIC{$x = y + 0 \to y + 0 = y \to x = y$}
\end{ptree}
\noindent$\Pi_{2}$:
\begin{ptree}
\AXC{$\forall x \in \mathbb{R}\,(x + 0 = x)$\lbl{(R4)}}
\AXC{$\forall x \in \mathbb{R}\,(x + 0 = x) \to y \in \mathbb{R} \to y + 0 = y$\lbl{A2}}
\BIC{$y \in \mathbb{R} \to y + 0 = y$}
\AXC{$(y \in \mathbb{R} \to y + 0 = y) \to \varphi_{21}$\lbl{A1}}
\BIC{$\varphi_{21}$}
\end{ptree}
\noindent$\Pi_{3}$:
\begin{ptree}
\AXC{$y = z \to x = y \to x = z$\lbl{A6}}
\AXC{$(y = z \to x = y \to x = z) \to x = y \to y = z \to x = z$\lbl{A1}}
\BIC{$x = y \to y = z \to x = z$}
\UIC{$\forall z\,(x = y \to y = z \to x = z)$}
\UIC{$\forall y\,\forall z\,(x = y \to y = z \to x = z)$}
\UIC{$\forall x\,\forall y\,\forall z\,(x = y \to y = z \to x = z)$}
\AXC{$\varphi_{2}$\lbl{A2}}
\BIC{$\forall y\,\forall z\,(x = y \to y = z \to x = z)$}
\AXC{$\varphi_{4}$\lbl{A2}}
\BIC{$\forall z\,(x = 0 + y \to 0 + y = z \to x = z)$}
\AXC{$\varphi_{8}$\lbl{A2}}
\BIC{$x = 0 + y \to 0 + y = y + 0 \to x = y + 0$}
\end{ptree}
\noindent$\Pi_{4}$:
\begin{ptree}
\AXC{$\forall x\,\forall y\,(x \in \mathbb{R} \to y \in \mathbb{R} \to x + y = y + x)$\lbl{(R3)}}
\AXC{$\varphi_{16}$\lbl{A2}}
\BIC{$\forall y\,(0 \in \mathbb{R} \to y \in \mathbb{R} \to 0 + y = y + 0)$}
\AXC{$\varphi_{12}$\lbl{A2}}
\BIC{$0 \in \mathbb{R} \to y \in \mathbb{R} \to 0 + y = y + 0$}
\AXC{$(0 \in \mathbb{R} \to y \in \mathbb{R} \to 0 + y = y + 0) \to 0 \in \mathbb{R} \to \varphi_{23}$\lbl{A1}}
\BIC{$0 \in \mathbb{R} \to \varphi_{23}$}
\end{ptree}
\noindent$\Pi_{5}$:
\begin{ptree}
\AXC{$y = z \to x = y \to x = z$\lbl{A6}}
\AXC{$(y = z \to x = y \to x = z) \to x = y \to y = z \to x = z$\lbl{A1}}
\BIC{$x = y \to y = z \to x = z$}
\UIC{$\forall z\,(x = y \to y = z \to x = z)$}
\UIC{$\forall y\,\forall z\,(x = y \to y = z \to x = z)$}
\UIC{$\forall x\,\forall y\,\forall z\,(x = y \to y = z \to x = z)$}
\AXC{$\varphi_{2}$\lbl{A2}}
\BIC{$\forall y\,\forall z\,(x = y \to y = z \to x = z)$}
\AXC{$\varphi_{7}$\lbl{A2}}
\BIC{$\forall z\,(x = x - y + y \to x - y + y = z \to x = z)$}
\AXC{$\varphi_{14}$\lbl{A2}}
\BIC{$x = x - y + y \to x - y + y = 0 + y \to x = 0 + y$}
\end{ptree}
\noindent$\Pi_{6}$:
\begin{ptree}
\AXC{$z = w \to x - y + y = z + y \to x - y + y = w + y$\lbl{A6}}
\UIC{$\forall w\,(z = w \to x - y + y = z + y \to x - y + y = w + y)$}
\UIC{$\forall z\,\forall w\,(z = w \to x - y + y = z + y \to x - y + y = w + y)$}
\AXC{$\varphi_{22}$\lbl{A2}}
\BIC{$\varphi_{6}$}
\AXC{$\varphi_{6} \to x - y = 0 \to x - y + y = x - y + y \to x - y + y = 0 + y$\lbl{A2}}
\BIC{$x - y = 0 \to x - y + y = x - y + y \to x - y + y = 0 + y$}
\end{ptree}
\noindent$\Pi_{7}$:
\begin{ptree}
\AXC{$x = x$\lbl{A5}}
\UIC{$\forall x\,(x = x)$}
\AXC{$\forall x\,(x = x) \to x - y + y = x - y + y$\lbl{A2}}
\BIC{$x - y + y = x - y + y$}
\AXC{$x \in \mathbb{R} \to y \in \mathbb{R} \to x \neq y \to x - y = 0 \to x - y = 0$\lbl{A1}}
\AXC{$\varphi_{30}$\lbl{A1}}
\BIC{$x - y + y = x - y + y \to \varphi_{25}$}
\BIC{$\varphi_{25}$}
\end{ptree}
\noindent$\Pi_{8}$:
\begin{ptree}
\AXC{$x = x$\lbl{A5}}
\AXC{$x = y \to x = x \to y = x$\lbl{A6}}
\AXC{$(x = y \to x = x \to y = x) \to x = x \to x = y \to y = x$\lbl{A1}}
\BIC{$x = x \to x = y \to y = x$}
\BIC{$x = y \to y = x$}
\UIC{$\forall y\,(x = y \to y = x)$}
\UIC{$\forall x\,\forall y\,(x = y \to y = x)$}
\end{ptree}
\noindent$\Pi_{9}$:
\begin{ptree}
\AXC{$\stackrel{\Pi_{8}}{\vdots}$}\noLine
\UIC{$\forall x\,\forall y\,(x = y \to y = x)$}
\AXC{$\forall x\,\forall y\,(x = y \to y = x) \to \forall y\,(v_{297} = y \to y = v_{297})$\lbl{A2}}
\BIC{$\forall y\,(v_{297} = y \to y = v_{297})$}
\AXC{$\varphi_{1}$\lbl{A2}}
\BIC{$v_{297} = v_{298} \to v_{298} = v_{297}$}
\UIC{$\forall v_{298}\,(v_{297} = v_{298} \to v_{298} = v_{297})$}
\UIC{$\forall v_{297}\,\forall v_{298}\,(v_{297} = v_{298} \to v_{298} = v_{297})$}
\AXC{$\varphi_{11}$\lbl{A2}}
\BIC{$\forall v_{298}\,(x - y + y = v_{298} \to v_{298} = x - y + y)$}
\AXC{$\varphi_{9}$\lbl{A2}}
\BIC{$x - y + y = x \to x = x - y + y$}
\end{ptree}
\noindent$\Pi_{10}$:
\begin{ptree}
\AXC{$\forall x\,\forall y\,(x \in \mathbb{R} \to y \in \mathbb{R} \to x - y + y = x)$\lbl{[F12]}}
\AXC{$\varphi_{17}$\lbl{A2}}
\BIC{$\forall y\,(x \in \mathbb{R} \to y \in \mathbb{R} \to x - y + y = x)$}
\AXC{$\varphi_{13}$\lbl{A2}}
\BIC{$x \in \mathbb{R} \to y \in \mathbb{R} \to x - y + y = x$}
\AXC{$\varphi_{31}$\lbl{A1}}
\BIC{$(x \in \mathbb{R} \to y \in \mathbb{R} \to x \neq y \to x - y = 0 \to x \in \mathbb{R}) \to \varphi_{24}$}
\end{ptree}
\noindent$\Pi_{11}$:
\begin{ptree}
\AXC{$x \in \mathbb{R} \to y \in \mathbb{R} \to x \neq y \to x - y = 0 \to x \in \mathbb{R}$\lbl{A1}}
\AXC{$\stackrel{\Pi_{10}}{\vdots}$}\noLine
\UIC{$(x \in \mathbb{R} \to y \in \mathbb{R} \to x \neq y \to x - y = 0 \to x \in \mathbb{R}) \to \varphi_{24}$}
\BIC{$\varphi_{24}$}
\end{ptree}
\noindent$\Pi_{12}$:
\begin{ptree}
\AXC{$\stackrel{\Pi_{9}}{\vdots}$}\noLine
\UIC{$x - y + y = x \to x = x - y + y$}
\AXC{$x \in \mathbb{R} \to y \in \mathbb{R} \to x \neq y \to x - y = 0 \to y \in \mathbb{R}$\lbl{A1}}
\AXC{$\stackrel{\Pi_{11}}{\vdots}$}\noLine
\UIC{$\varphi_{24}$}
\BIC{$x \in \mathbb{R} \to y \in \mathbb{R} \to x \neq y \to x - y = 0 \to x - y + y = x$}
\AXC{$\varphi_{26}$\lbl{A1}}
\BIC{$\varphi_{15}$}
\BIC{$x \in \mathbb{R} \to y \in \mathbb{R} \to x \neq y \to x - y = 0 \to x = x - y + y$}
\AXC{$\varphi_{32}$\lbl{A1}}
\BIC{$\varphi_{28}$}
\end{ptree}
\noindent$\Pi_{13}$:
\begin{ptree}
\AXC{$\stackrel{\Pi_{5}}{\vdots}$}\noLine
\UIC{$x = x - y + y \to x - y + y = 0 + y \to x = 0 + y$}
\AXC{$\stackrel{\Pi_{6}}{\vdots}$}\noLine
\UIC{$x - y = 0 \to x - y + y = x - y + y \to x - y + y = 0 + y$}
\AXC{$\stackrel{\Pi_{7}}{\vdots}$}\noLine
\UIC{$\varphi_{25}$}
\BIC{$x \in \mathbb{R} \to y \in \mathbb{R} \to x \neq y \to x - y = 0 \to x - y + y = 0 + y$}
\AXC{$\stackrel{\Pi_{12}}{\vdots}$}\noLine
\UIC{$\varphi_{28}$}
\BIC{$\varphi_{20}$}
\BIC{$x \in \mathbb{R} \to y \in \mathbb{R} \to x \neq y \to x - y = 0 \to x = 0 + y$}
\end{ptree}
\noindent$\Pi_{14}$:
\begin{ptree}
\AXC{$\stackrel{\Pi_{13}}{\vdots}$}\noLine
\UIC{$x \in \mathbb{R} \to y \in \mathbb{R} \to x \neq y \to x - y = 0 \to x = 0 + y$}
\AXC{$(x \in \mathbb{R} \to y \in \mathbb{R} \to x \neq y \to x - y = 0 \to x = 0 + y) \to \varphi_{27}$\lbl{A1}}
\BIC{$\varphi_{27}$}
\end{ptree}
\noindent$\Pi_{15}$:
\begin{ptree}
\AXC{$x \in \mathbb{R} \to y \in \mathbb{R} \to x \neq y \to x - y = 0 \to y \in \mathbb{R}$\lbl{A1}}
\AXC{$0 \in \mathbb{R}$\lbl{(R1)}}
\AXC{$\stackrel{\Pi_{4}}{\vdots}$}\noLine
\UIC{$0 \in \mathbb{R} \to \varphi_{23}$}
\BIC{$\varphi_{23}$}
\BIC{$x \in \mathbb{R} \to y \in \mathbb{R} \to x \neq y \to x - y = 0 \to 0 + y = y + 0$}
\AXC{$\stackrel{\Pi_{14}}{\vdots}$}\noLine
\UIC{$\varphi_{27}$}
\BIC{$\varphi_{18}$}
\end{ptree}
\noindent$\Pi_{16}$:
\begin{ptree}
\AXC{$x \in \mathbb{R} \to y \in \mathbb{R} \to x \neq y \to x - y = 0 \to y \in \mathbb{R}$\lbl{A1}}
\AXC{$\stackrel{\Pi_{2}}{\vdots}$}\noLine
\UIC{$\varphi_{21}$}
\BIC{$x \in \mathbb{R} \to y \in \mathbb{R} \to x \neq y \to x - y = 0 \to y + 0 = y$}
\AXC{$\stackrel{\Pi_{3}}{\vdots}$}\noLine
\UIC{$x = 0 + y \to 0 + y = y + 0 \to x = y + 0$}
\AXC{$\stackrel{\Pi_{15}}{\vdots}$}\noLine
\UIC{$\varphi_{18}$}
\BIC{$x \in \mathbb{R} \to y \in \mathbb{R} \to x \neq y \to x - y = 0 \to x = y + 0$}
\AXC{$\varphi_{29}$\lbl{A1}}
\BIC{$(x \in \mathbb{R} \to y \in \mathbb{R} \to x \neq y \to x - y = 0 \to y + 0 = y) \to \varphi_{10}$}
\BIC{$\varphi_{10}$}
\end{ptree}
\noindent 主木:
\begin{ptree}
\AXC{$\stackrel{\Pi_{1}}{\vdots}$}\noLine
\UIC{$x = y + 0 \to y + 0 = y \to x = y$}
\AXC{$\stackrel{\Pi_{16}}{\vdots}$}\noLine
\UIC{$\varphi_{10}$}
\BIC{$x \in \mathbb{R} \to y \in \mathbb{R} \to x \neq y \to x - y = 0 \to x = y$}
\AXC{$x \in \mathbb{R} \to y \in \mathbb{R} \to x \neq y \to x \neq y$\lbl{A1}}
\AXC{$(x \in \mathbb{R} \to y \in \mathbb{R} \to x \neq y \to x \neq y) \to \varphi_{19}$\lbl{A1}}
\BIC{$\varphi_{19}$}
\BIC{$x \in \mathbb{R} \to y \in \mathbb{R} \to x \neq y \to x - y \neq 0$}
\UIC{$\forall y\,(x \in \mathbb{R} \to y \in \mathbb{R} \to x \neq y \to x - y \neq 0)$}
\UIC{$\forall x\,\forall y\,(x \in \mathbb{R} \to y \in \mathbb{R} \to x \neq y \to x - y \neq 0)$}
\end{ptree}
\begin{abbrevlist}
\item $\varphi_{1} :\equiv \forall y\,(v_{297} = y \to y = v_{297}) \to v_{297} = v_{298} \to v_{298} = v_{297}$
\item $\varphi_{2} :\equiv \forall x\,\forall y\,\forall z\,(x = y \to y = z \to x = z) \to \forall y\,\forall z\,(x = y \to y = z \to x = z)$
\item $\varphi_{3} :\equiv \forall y\,\forall z\,(x = y \to y = z \to x = z) \to \forall z\,(x = y + 0 \to y + 0 = z \to x = z)$
\item $\varphi_{4} :\equiv \forall y\,\forall z\,(x = y \to y = z \to x = z) \to \forall z\,(x = 0 + y \to 0 + y = z \to x = z)$
\item $\varphi_{5} :\equiv \forall z\,(x = y + 0 \to y + 0 = z \to x = z) \to x = y + 0 \to y + 0 = y \to x = y$
\item $\varphi_{6} :\equiv \forall w\,(x - y = w \to x - y + y = x - y + y \to x - y + y = w + y)$
\item $\varphi_{7} :\equiv \forall y\,\forall z\,(x = y \to y = z \to x = z) \to \forall z\,(x = x - y + y \to x - y + y = z \to x = z)$
\item $\varphi_{8} :\equiv \forall z\,(x = 0 + y \to 0 + y = z \to x = z) \to x = 0 + y \to 0 + y = y + 0 \to x = y + 0$
\item $\varphi_{9} :\equiv \forall v_{298}\,(x - y + y = v_{298} \to v_{298} = x - y + y) \to x - y + y = x \to x = x - y + y$
\item $\varphi_{10} :\equiv (x = y + 0 \to y + 0 = y \to x = y) \to x \in \mathbb{R} \to y \in \mathbb{R} \to x \neq y \to x - y = 0 \to x = y$
\item $\varphi_{11} :\equiv \forall v_{297}\,\forall v_{298}\,(v_{297} = v_{298} \to v_{298} = v_{297}) \to \forall v_{298}\,(x - y + y = v_{298} \to v_{298} = x - y + y)$
\item $\varphi_{12} :\equiv \forall y\,(0 \in \mathbb{R} \to y \in \mathbb{R} \to 0 + y = y + 0) \to 0 \in \mathbb{R} \to y \in \mathbb{R} \to 0 + y = y + 0$
\item $\varphi_{13} :\equiv \forall y\,(x \in \mathbb{R} \to y \in \mathbb{R} \to x - y + y = x) \to x \in \mathbb{R} \to y \in \mathbb{R} \to x - y + y = x$
\item $\varphi_{14} :\equiv \forall z\,(x = x - y + y \to x - y + y = z \to x = z) \to x = x - y + y \to x - y + y = 0 + y \to x = 0 + y$
\item $\varphi_{15} :\equiv (x - y + y = x \to x = x - y + y) \to x \in \mathbb{R} \to y \in \mathbb{R} \to x \neq y \to x - y = 0 \to x = x - y + y$
\item $\varphi_{16} :\equiv \forall x\,\forall y\,(x \in \mathbb{R} \to y \in \mathbb{R} \to x + y = y + x) \to \forall y\,(0 \in \mathbb{R} \to y \in \mathbb{R} \to 0 + y = y + 0)$
\item $\varphi_{17} :\equiv \forall x\,\forall y\,(x \in \mathbb{R} \to y \in \mathbb{R} \to x - y + y = x) \to \forall y\,(x \in \mathbb{R} \to y \in \mathbb{R} \to x - y + y = x)$
\item $\varphi_{18} :\equiv (x = 0 + y \to 0 + y = y + 0 \to x = y + 0) \to x \in \mathbb{R} \to y \in \mathbb{R} \to x \neq y \to x - y = 0 \to x = y + 0$
\item $\varphi_{19} :\equiv (x \in \mathbb{R} \to y \in \mathbb{R} \to x \neq y \to x - y = 0 \to x = y) \to x \in \mathbb{R} \to y \in \mathbb{R} \to x \neq y \to x - y \neq 0$
\item $\varphi_{20} :\equiv (x = x - y + y \to x - y + y = 0 + y \to x = 0 + y) \to x \in \mathbb{R} \to y \in \mathbb{R} \to x \neq y \to x - y = 0 \to x = 0 + y$
\item $\varphi_{21} :\equiv (x \in \mathbb{R} \to y \in \mathbb{R} \to x \neq y \to x - y = 0 \to y \in \mathbb{R}) \to x \in \mathbb{R} \to y \in \mathbb{R} \to x \neq y \to x - y = 0 \to y + 0 = y$
\item $\varphi_{22} :\equiv \forall z\,\forall w\,(z = w \to x - y + y = z + y \to x - y + y = w + y) \to \varphi_{6}$
\item $\varphi_{23} :\equiv (x \in \mathbb{R} \to y \in \mathbb{R} \to x \neq y \to x - y = 0 \to y \in \mathbb{R}) \to x \in \mathbb{R} \to y \in \mathbb{R} \to x \neq y \to x - y = 0 \to 0 + y = y + 0$
\item $\varphi_{24} :\equiv (x \in \mathbb{R} \to y \in \mathbb{R} \to x \neq y \to x - y = 0 \to y \in \mathbb{R}) \to x \in \mathbb{R} \to y \in \mathbb{R} \to x \neq y \to x - y = 0 \to x - y + y = x$
\item $\varphi_{25} :\equiv (x - y = 0 \to x - y + y = x - y + y \to x - y + y = 0 + y) \to x \in \mathbb{R} \to y \in \mathbb{R} \to x \neq y \to x - y = 0 \to x - y + y = 0 + y$
\item $\varphi_{26} :\equiv (x \in \mathbb{R} \to y \in \mathbb{R} \to x \neq y \to x - y = 0 \to x - y + y = x) \to \varphi_{15}$
\item $\varphi_{27} :\equiv (x \in \mathbb{R} \to y \in \mathbb{R} \to x \neq y \to x - y = 0 \to 0 + y = y + 0) \to \varphi_{18}$
\item $\varphi_{28} :\equiv (x \in \mathbb{R} \to y \in \mathbb{R} \to x \neq y \to x - y = 0 \to x - y + y = 0 + y) \to \varphi_{20}$
\item $\varphi_{29} :\equiv (x \in \mathbb{R} \to y \in \mathbb{R} \to x \neq y \to x - y = 0 \to x = y + 0) \to (x \in \mathbb{R} \to y \in \mathbb{R} \to x \neq y \to x - y = 0 \to y + 0 = y) \to \varphi_{10}$
\item $\varphi_{30} :\equiv (x \in \mathbb{R} \to y \in \mathbb{R} \to x \neq y \to x - y = 0 \to x - y = 0) \to x - y + y = x - y + y \to \varphi_{25}$
\item $\varphi_{31} :\equiv (x \in \mathbb{R} \to y \in \mathbb{R} \to x - y + y = x) \to (x \in \mathbb{R} \to y \in \mathbb{R} \to x \neq y \to x - y = 0 \to x \in \mathbb{R}) \to \varphi_{24}$
\item $\varphi_{32} :\equiv (x \in \mathbb{R} \to y \in \mathbb{R} \to x \neq y \to x - y = 0 \to x = x - y + y) \to \varphi_{28}$
\end{abbrevlist}
\end{fullproof}

% 自動生成 (prover/build.js): MulNonneg  ノード数 82
\begin{fullproof}{[A16]}{$\forall x\,\forall y\,(x \in \mathbb{R} \to y \in \mathbb{R} \to 0 \leq x \to 0 \leq y \to 0 \leq x \cdot y)$}{82}\phantomsection\label{tree:MulNonneg}
\noindent$\Pi_{1}$:
\begin{ptree}
\AXC{$0 \in \mathbb{R}$\lbl{(R1)}}
\AXC{$\varphi_{9}$\lbl{[F33]}}
\AXC{$\varphi_{9} \to \varphi_{7}$\lbl{A2}}
\BIC{$\varphi_{7}$}
\AXC{$\varphi_{7} \to \varphi_{6}$\lbl{A2}}
\BIC{$\varphi_{6}$}
\AXC{$\varphi_{22}$\lbl{A2}}
\BIC{$0 \in \mathbb{R} \to x \in \mathbb{R} \to y \in \mathbb{R} \to 0 \leq y \to 0 \leq x \to 0 \cdot y \leq x \cdot y$}
\AXC{$\varphi_{27}$\lbl{A1}}
\BIC{$\varphi_{26}$}
\BIC{$(x \in \mathbb{R} \to y \in \mathbb{R} \to 0 \leq x \to 0 \leq y \to x \in \mathbb{R}) \to \varphi_{25}$}
\end{ptree}
\noindent$\Pi_{2}$:
\begin{ptree}
\AXC{$x \in \mathbb{R} \to y \in \mathbb{R} \to 0 \leq x \to 0 \leq y \to x \in \mathbb{R}$\lbl{A1}}
\AXC{$\stackrel{\Pi_{1}}{\vdots}$}\noLine
\UIC{$(x \in \mathbb{R} \to y \in \mathbb{R} \to 0 \leq x \to 0 \leq y \to x \in \mathbb{R}) \to \varphi_{25}$}
\BIC{$\varphi_{25}$}
\end{ptree}
\noindent$\Pi_{3}$:
\begin{ptree}
\AXC{$x \in \mathbb{R} \to y \in \mathbb{R} \to 0 \leq x \to 0 \leq y \to 0 \leq y$\lbl{A1}}
\AXC{$x \in \mathbb{R} \to y \in \mathbb{R} \to 0 \leq x \to 0 \leq y \to y \in \mathbb{R}$\lbl{A1}}
\AXC{$\stackrel{\Pi_{2}}{\vdots}$}\noLine
\UIC{$\varphi_{25}$}
\BIC{$(x \in \mathbb{R} \to y \in \mathbb{R} \to 0 \leq x \to 0 \leq y \to 0 \leq y) \to \varphi_{18}$}
\BIC{$\varphi_{18}$}
\end{ptree}
\noindent$\Pi_{4}$:
\begin{ptree}
\AXC{$y = z \to x = y \to x = z$\lbl{A6}}
\AXC{$(y = z \to x = y \to x = z) \to x = y \to y = z \to x = z$\lbl{A1}}
\BIC{$x = y \to y = z \to x = z$}
\UIC{$\forall z\,(x = y \to y = z \to x = z)$}
\UIC{$\forall y\,\forall z\,(x = y \to y = z \to x = z)$}
\UIC{$\forall x\,\forall y\,\forall z\,(x = y \to y = z \to x = z)$}
\AXC{$\varphi_{3}$\lbl{A2}}
\BIC{$\forall y\,\forall z\,(v_{299} = y \to y = z \to v_{299} = z)$}
\AXC{$\varphi_{4}$\lbl{A2}}
\BIC{$\forall z\,(v_{299} = v_{300} \to v_{300} = z \to v_{299} = z)$}
\AXC{$\varphi_{5}$\lbl{A2}}
\BIC{$v_{299} = v_{300} \to v_{300} = v_{301} \to v_{299} = v_{301}$}
\UIC{$\forall v_{301}\,(v_{299} = v_{300} \to v_{300} = v_{301} \to v_{299} = v_{301})$}
\UIC{$\varphi_{1}$}
\UIC{$\forall v_{299}\,(\varphi_{1})$}
\end{ptree}
\noindent$\Pi_{5}$:
\begin{ptree}
\AXC{$\forall x \in \mathbb{R}\,(x \cdot 0 = 0)$\lbl{[F8]}}
\AXC{$\forall x \in \mathbb{R}\,(x \cdot 0 = 0) \to y \in \mathbb{R} \to y \cdot 0 = 0$\lbl{A2}}
\BIC{$y \in \mathbb{R} \to y \cdot 0 = 0$}
\AXC{$(y \in \mathbb{R} \to y \cdot 0 = 0) \to \varphi_{19}$\lbl{A1}}
\BIC{$\varphi_{19}$}
\end{ptree}
\noindent$\Pi_{6}$:
\begin{ptree}
\AXC{$\forall x\,\forall y\,(x \in \mathbb{R} \to y \in \mathbb{R} \to x \cdot y = y \cdot x)$\lbl{(R7)}}
\AXC{$\varphi_{17}$\lbl{A2}}
\BIC{$\forall y\,(0 \in \mathbb{R} \to y \in \mathbb{R} \to 0 \cdot y = y \cdot 0)$}
\AXC{$\varphi_{12}$\lbl{A2}}
\BIC{$0 \in \mathbb{R} \to y \in \mathbb{R} \to 0 \cdot y = y \cdot 0$}
\AXC{$(0 \in \mathbb{R} \to y \in \mathbb{R} \to 0 \cdot y = y \cdot 0) \to 0 \in \mathbb{R} \to \varphi_{21}$\lbl{A1}}
\BIC{$0 \in \mathbb{R} \to \varphi_{21}$}
\end{ptree}
\noindent$\Pi_{7}$:
\begin{ptree}
\AXC{$x \in \mathbb{R} \to y \in \mathbb{R} \to 0 \leq x \to 0 \leq y \to y \in \mathbb{R}$\lbl{A1}}
\AXC{$0 \in \mathbb{R}$\lbl{(R1)}}
\AXC{$\stackrel{\Pi_{6}}{\vdots}$}\noLine
\UIC{$0 \in \mathbb{R} \to \varphi_{21}$}
\BIC{$\varphi_{21}$}
\BIC{$x \in \mathbb{R} \to y \in \mathbb{R} \to 0 \leq x \to 0 \leq y \to 0 \cdot y = y \cdot 0$}
\AXC{$\varphi_{24}$\lbl{A1}}
\BIC{$(x \in \mathbb{R} \to y \in \mathbb{R} \to 0 \leq x \to 0 \leq y \to y \cdot 0 = 0) \to \varphi_{13}$}
\end{ptree}
\noindent$\Pi_{8}$:
\begin{ptree}
\AXC{$x \in \mathbb{R} \to y \in \mathbb{R} \to 0 \leq x \to 0 \leq y \to y \in \mathbb{R}$\lbl{A1}}
\AXC{$\stackrel{\Pi_{5}}{\vdots}$}\noLine
\UIC{$\varphi_{19}$}
\BIC{$x \in \mathbb{R} \to y \in \mathbb{R} \to 0 \leq x \to 0 \leq y \to y \cdot 0 = 0$}
\AXC{$\stackrel{\Pi_{7}}{\vdots}$}\noLine
\UIC{$(x \in \mathbb{R} \to y \in \mathbb{R} \to 0 \leq x \to 0 \leq y \to y \cdot 0 = 0) \to \varphi_{13}$}
\BIC{$\varphi_{13}$}
\end{ptree}
\noindent$\Pi_{9}$:
\begin{ptree}
\AXC{$\stackrel{\Pi_{4}}{\vdots}$}\noLine
\UIC{$\forall v_{299}\,(\varphi_{1})$}
\AXC{$\forall v_{299}\,(\varphi_{1}) \to \varphi_{2}$\lbl{A2}}
\BIC{$\varphi_{2}$}
\AXC{$\varphi_{20}$\lbl{A2}}
\BIC{$\forall v_{301}\,(0 \cdot y = y \cdot 0 \to y \cdot 0 = v_{301} \to 0 \cdot y = v_{301})$}
\AXC{$\varphi_{15}$\lbl{A2}}
\BIC{$0 \cdot y = y \cdot 0 \to y \cdot 0 = 0 \to 0 \cdot y = 0$}
\AXC{$\stackrel{\Pi_{8}}{\vdots}$}\noLine
\UIC{$\varphi_{13}$}
\BIC{$x \in \mathbb{R} \to y \in \mathbb{R} \to 0 \leq x \to 0 \leq y \to 0 \cdot y = 0$}
\end{ptree}
\noindent$\Pi_{10}$:
\begin{ptree}
\AXC{$\stackrel{\Pi_{9}}{\vdots}$}\noLine
\UIC{$x \in \mathbb{R} \to y \in \mathbb{R} \to 0 \leq x \to 0 \leq y \to 0 \cdot y = 0$}
\AXC{$(x \in \mathbb{R} \to y \in \mathbb{R} \to 0 \leq x \to 0 \leq y \to 0 \cdot y = 0) \to \varphi_{23}$\lbl{A1}}
\BIC{$\varphi_{23}$}
\end{ptree}
\noindent$\Pi_{11}$:
\begin{ptree}
\AXC{$x \in \mathbb{R} \to y \in \mathbb{R} \to 0 \leq x \to 0 \leq y \to 0 \leq x$\lbl{A1}}
\AXC{$\stackrel{\Pi_{3}}{\vdots}$}\noLine
\UIC{$\varphi_{18}$}
\BIC{$x \in \mathbb{R} \to y \in \mathbb{R} \to 0 \leq x \to 0 \leq y \to 0 \cdot y \leq x \cdot y$}
\AXC{$\stackrel{\Pi_{10}}{\vdots}$}\noLine
\UIC{$\varphi_{23}$}
\BIC{$(0 \cdot y = 0 \to 0 \cdot y < x \cdot y \to 0 < x \cdot y) \to \varphi_{14}$}
\end{ptree}
\noindent$\Pi_{12}$:
\begin{ptree}
\AXC{$z = w \to z < x \cdot y \to w < x \cdot y$\lbl{A6}}
\UIC{$\forall w\,(z = w \to z < x \cdot y \to w < x \cdot y)$}
\UIC{$\forall z\,\forall w\,(z = w \to z < x \cdot y \to w < x \cdot y)$}
\AXC{$\varphi_{11}$\lbl{A2}}
\BIC{$\forall w\,(0 \cdot y = w \to 0 \cdot y < x \cdot y \to w < x \cdot y)$}
\AXC{$\varphi_{16}$\lbl{A2}}
\BIC{$0 \cdot y = 0 \to 0 \cdot y < x \cdot y \to 0 < x \cdot y$}
\AXC{$\stackrel{\Pi_{11}}{\vdots}$}\noLine
\UIC{$(0 \cdot y = 0 \to 0 \cdot y < x \cdot y \to 0 < x \cdot y) \to \varphi_{14}$}
\BIC{$\varphi_{14}$}
\end{ptree}
\noindent 主木:
\begin{ptree}
\AXC{$z = w \to z = x \cdot y \to w = x \cdot y$\lbl{A6}}
\UIC{$\forall w\,(z = w \to z = x \cdot y \to w = x \cdot y)$}
\UIC{$\forall z\,\forall w\,(z = w \to z = x \cdot y \to w = x \cdot y)$}
\AXC{$\varphi_{8}$\lbl{A2}}
\BIC{$\forall w\,(0 \cdot y = w \to 0 \cdot y = x \cdot y \to w = x \cdot y)$}
\AXC{$\varphi_{10}$\lbl{A2}}
\BIC{$0 \cdot y = 0 \to 0 \cdot y = x \cdot y \to 0 = x \cdot y$}
\AXC{$\stackrel{\Pi_{12}}{\vdots}$}\noLine
\UIC{$\varphi_{14}$}
\BIC{$x \in \mathbb{R} \to y \in \mathbb{R} \to 0 \leq x \to 0 \leq y \to 0 \leq x \cdot y$}
\UIC{$\forall y\,(x \in \mathbb{R} \to y \in \mathbb{R} \to 0 \leq x \to 0 \leq y \to 0 \leq x \cdot y)$}
\UIC{$\forall x\,\forall y\,(x \in \mathbb{R} \to y \in \mathbb{R} \to 0 \leq x \to 0 \leq y \to 0 \leq x \cdot y)$}
\end{ptree}
\begin{abbrevlist}
\item $\varphi_{1} :\equiv \forall v_{300}\,\forall v_{301}\,(v_{299} = v_{300} \to v_{300} = v_{301} \to v_{299} = v_{301})$
\item $\varphi_{2} :\equiv \forall v_{300}\,\forall v_{301}\,(0 \cdot y = v_{300} \to v_{300} = v_{301} \to 0 \cdot y = v_{301})$
\item $\varphi_{3} :\equiv \forall x\,\forall y\,\forall z\,(x = y \to y = z \to x = z) \to \forall y\,\forall z\,(v_{299} = y \to y = z \to v_{299} = z)$
\item $\varphi_{4} :\equiv \forall y\,\forall z\,(v_{299} = y \to y = z \to v_{299} = z) \to \forall z\,(v_{299} = v_{300} \to v_{300} = z \to v_{299} = z)$
\item $\varphi_{5} :\equiv \forall z\,(v_{299} = v_{300} \to v_{300} = z \to v_{299} = z) \to v_{299} = v_{300} \to v_{300} = v_{301} \to v_{299} = v_{301}$
\item $\varphi_{6} :\equiv \forall z\,(0 \in \mathbb{R} \to x \in \mathbb{R} \to z \in \mathbb{R} \to 0 \leq z \to 0 \leq x \to 0 \cdot z \leq x \cdot z)$
\item $\varphi_{7} :\equiv \forall y\,\forall z\,(0 \in \mathbb{R} \to y \in \mathbb{R} \to z \in \mathbb{R} \to 0 \leq z \to 0 \leq y \to 0 \cdot z \leq y \cdot z)$
\item $\varphi_{8} :\equiv \forall z\,\forall w\,(z = w \to z = x \cdot y \to w = x \cdot y) \to \forall w\,(0 \cdot y = w \to 0 \cdot y = x \cdot y \to w = x \cdot y)$
\item $\varphi_{9} :\equiv \forall x\,\forall y\,\forall z\,(x \in \mathbb{R} \to y \in \mathbb{R} \to z \in \mathbb{R} \to 0 \leq z \to x \leq y \to x \cdot z \leq y \cdot z)$
\item $\varphi_{10} :\equiv \forall w\,(0 \cdot y = w \to 0 \cdot y = x \cdot y \to w = x \cdot y) \to 0 \cdot y = 0 \to 0 \cdot y = x \cdot y \to 0 = x \cdot y$
\item $\varphi_{11} :\equiv \forall z\,\forall w\,(z = w \to z < x \cdot y \to w < x \cdot y) \to \forall w\,(0 \cdot y = w \to 0 \cdot y < x \cdot y \to w < x \cdot y)$
\item $\varphi_{12} :\equiv \forall y\,(0 \in \mathbb{R} \to y \in \mathbb{R} \to 0 \cdot y = y \cdot 0) \to 0 \in \mathbb{R} \to y \in \mathbb{R} \to 0 \cdot y = y \cdot 0$
\item $\varphi_{13} :\equiv (0 \cdot y = y \cdot 0 \to y \cdot 0 = 0 \to 0 \cdot y = 0) \to x \in \mathbb{R} \to y \in \mathbb{R} \to 0 \leq x \to 0 \leq y \to 0 \cdot y = 0$
\item $\varphi_{14} :\equiv (0 \cdot y = 0 \to 0 \cdot y = x \cdot y \to 0 = x \cdot y) \to x \in \mathbb{R} \to y \in \mathbb{R} \to 0 \leq x \to 0 \leq y \to 0 \leq x \cdot y$
\item $\varphi_{15} :\equiv \forall v_{301}\,(0 \cdot y = y \cdot 0 \to y \cdot 0 = v_{301} \to 0 \cdot y = v_{301}) \to 0 \cdot y = y \cdot 0 \to y \cdot 0 = 0 \to 0 \cdot y = 0$
\item $\varphi_{16} :\equiv \forall w\,(0 \cdot y = w \to 0 \cdot y < x \cdot y \to w < x \cdot y) \to 0 \cdot y = 0 \to 0 \cdot y < x \cdot y \to 0 < x \cdot y$
\item $\varphi_{17} :\equiv \forall x\,\forall y\,(x \in \mathbb{R} \to y \in \mathbb{R} \to x \cdot y = y \cdot x) \to \forall y\,(0 \in \mathbb{R} \to y \in \mathbb{R} \to 0 \cdot y = y \cdot 0)$
\item $\varphi_{18} :\equiv (x \in \mathbb{R} \to y \in \mathbb{R} \to 0 \leq x \to 0 \leq y \to 0 \leq x) \to x \in \mathbb{R} \to y \in \mathbb{R} \to 0 \leq x \to 0 \leq y \to 0 \cdot y \leq x \cdot y$
\item $\varphi_{19} :\equiv (x \in \mathbb{R} \to y \in \mathbb{R} \to 0 \leq x \to 0 \leq y \to y \in \mathbb{R}) \to x \in \mathbb{R} \to y \in \mathbb{R} \to 0 \leq x \to 0 \leq y \to y \cdot 0 = 0$
\item $\varphi_{20} :\equiv \varphi_{2} \to \forall v_{301}\,(0 \cdot y = y \cdot 0 \to y \cdot 0 = v_{301} \to 0 \cdot y = v_{301})$
\item $\varphi_{21} :\equiv (x \in \mathbb{R} \to y \in \mathbb{R} \to 0 \leq x \to 0 \leq y \to y \in \mathbb{R}) \to x \in \mathbb{R} \to y \in \mathbb{R} \to 0 \leq x \to 0 \leq y \to 0 \cdot y = y \cdot 0$
\item $\varphi_{22} :\equiv \varphi_{6} \to 0 \in \mathbb{R} \to x \in \mathbb{R} \to y \in \mathbb{R} \to 0 \leq y \to 0 \leq x \to 0 \cdot y \leq x \cdot y$
\item $\varphi_{23} :\equiv (x \in \mathbb{R} \to y \in \mathbb{R} \to 0 \leq x \to 0 \leq y \to 0 \cdot y \leq x \cdot y) \to (0 \cdot y = 0 \to 0 \cdot y < x \cdot y \to 0 < x \cdot y) \to \varphi_{14}$
\item $\varphi_{24} :\equiv (x \in \mathbb{R} \to y \in \mathbb{R} \to 0 \leq x \to 0 \leq y \to 0 \cdot y = y \cdot 0) \to (x \in \mathbb{R} \to y \in \mathbb{R} \to 0 \leq x \to 0 \leq y \to y \cdot 0 = 0) \to \varphi_{13}$
\item $\varphi_{25} :\equiv (x \in \mathbb{R} \to y \in \mathbb{R} \to 0 \leq x \to 0 \leq y \to y \in \mathbb{R}) \to (x \in \mathbb{R} \to y \in \mathbb{R} \to 0 \leq x \to 0 \leq y \to 0 \leq y) \to \varphi_{18}$
\item $\varphi_{26} :\equiv 0 \in \mathbb{R} \to (x \in \mathbb{R} \to y \in \mathbb{R} \to 0 \leq x \to 0 \leq y \to x \in \mathbb{R}) \to \varphi_{25}$
\item $\varphi_{27} :\equiv (0 \in \mathbb{R} \to x \in \mathbb{R} \to y \in \mathbb{R} \to 0 \leq y \to 0 \leq x \to 0 \cdot y \leq x \cdot y) \to \varphi_{26}$
\end{abbrevlist}
\end{fullproof}

% 自動生成 (prover/build.js): AbsMul  ノード数 1070
\begin{fullproof}{[A17]}{$\forall x\,\forall y\,(x \in \mathbb{R} \to y \in \mathbb{R} \to \lvert x \cdot y \rvert = \lvert x \rvert \cdot \lvert y \rvert)$}{1070}\phantomsection\label{tree:AbsMul}
\noindent$\Pi_{1}$:
\begin{ptree}
\AXC{$y = z \to x = y \to x = z$\lbl{A6}}
\AXC{$(y = z \to x = y \to x = z) \to x = y \to y = z \to x = z$\lbl{A1}}
\BIC{$x = y \to y = z \to x = z$}
\UIC{$\forall z\,(x = y \to y = z \to x = z)$}
\UIC{$\forall y\,\forall z\,(x = y \to y = z \to x = z)$}
\UIC{$\forall x\,\forall y\,\forall z\,(x = y \to y = z \to x = z)$}
\AXC{$\varphi_{19}$\lbl{A2}}
\BIC{$\forall y\,\forall z\,(v_{358} = y \to y = z \to v_{358} = z)$}
\AXC{$\varphi_{41}$\lbl{A2}}
\BIC{$\forall z\,(v_{358} = v_{359} \to v_{359} = z \to v_{358} = z)$}
\AXC{$\varphi_{68}$\lbl{A2}}
\BIC{$v_{358} = v_{359} \to v_{359} = v_{360} \to v_{358} = v_{360}$}
\UIC{$\forall v_{360}\,(v_{358} = v_{359} \to v_{359} = v_{360} \to v_{358} = v_{360})$}
\UIC{$\varphi_{8}$}
\UIC{$\forall v_{358}\,(\varphi_{8})$}
\end{ptree}
\noindent$\Pi_{2}$:
\begin{ptree}
\AXC{$x = x$\lbl{A5}}
\AXC{$x = y \to x = x \to y = x$\lbl{A6}}
\AXC{$(x = y \to x = x \to y = x) \to x = x \to x = y \to y = x$\lbl{A1}}
\BIC{$x = x \to x = y \to y = x$}
\BIC{$x = y \to y = x$}
\UIC{$\forall y\,(x = y \to y = x)$}
\UIC{$\forall x\,\forall y\,(x = y \to y = x)$}
\end{ptree}
\noindent$\Pi_{3}$:
\begin{ptree}
\AXC{$\stackrel{\Pi_{2}}{\vdots}$}\noLine
\UIC{$\forall x\,\forall y\,(x = y \to y = x)$}
\AXC{$\forall x\,\forall y\,(x = y \to y = x) \to \forall y\,(v_{356} = y \to y = v_{356})$\lbl{A2}}
\BIC{$\forall y\,(v_{356} = y \to y = v_{356})$}
\AXC{$\varphi_{1}$\lbl{A2}}
\BIC{$v_{356} = v_{357} \to v_{357} = v_{356}$}
\UIC{$\forall v_{357}\,(v_{356} = v_{357} \to v_{357} = v_{356})$}
\UIC{$\forall v_{356}\,\forall v_{357}\,(v_{356} = v_{357} \to v_{357} = v_{356})$}
\AXC{$\varphi_{145}$\lbl{A2}}
\BIC{$\forall v_{357}\,(\lvert x \rvert \cdot \lvert y \rvert = v_{357} \to v_{357} = \lvert x \rvert \cdot \lvert y \rvert)$}
\AXC{$\varphi_{246}$\lbl{A2}}
\BIC{$\lvert x \rvert \cdot \lvert y \rvert = -x \cdot -y \to -x \cdot -y = \lvert x \rvert \cdot \lvert y \rvert$}
\end{ptree}
\noindent$\Pi_{4}$:
\begin{ptree}
\AXC{$y = z \to x = y \to x = z$\lbl{A6}}
\AXC{$(y = z \to x = y \to x = z) \to x = y \to y = z \to x = z$\lbl{A1}}
\BIC{$x = y \to y = z \to x = z$}
\UIC{$\forall z\,(x = y \to y = z \to x = z)$}
\UIC{$\forall y\,\forall z\,(x = y \to y = z \to x = z)$}
\UIC{$\forall x\,\forall y\,\forall z\,(x = y \to y = z \to x = z)$}
\AXC{$\varphi_{20}$\lbl{A2}}
\BIC{$\forall y\,\forall z\,(v_{344} = y \to y = z \to v_{344} = z)$}
\AXC{$\varphi_{42}$\lbl{A2}}
\BIC{$\forall z\,(v_{344} = v_{345} \to v_{345} = z \to v_{344} = z)$}
\AXC{$\varphi_{69}$\lbl{A2}}
\BIC{$v_{344} = v_{345} \to v_{345} = v_{346} \to v_{344} = v_{346}$}
\UIC{$\forall v_{346}\,(v_{344} = v_{345} \to v_{345} = v_{346} \to v_{344} = v_{346})$}
\UIC{$\varphi_{9}$}
\UIC{$\forall v_{344}\,(\varphi_{9})$}
\end{ptree}
\noindent$\Pi_{5}$:
\begin{ptree}
\AXC{$\forall x\,\forall y\,(x \in \mathbb{R} \to y \in \mathbb{R} \to \neg y \leq x \to x < y)$\lbl{[A2]}}
\AXC{$\varphi_{179}$\lbl{A2}}
\BIC{$\forall y\,(v_{342} \in \mathbb{R} \to y \in \mathbb{R} \to \neg y \leq v_{342} \to v_{342} < y)$}
\AXC{$\varphi_{223}$\lbl{A2}}
\BIC{$v_{342} \in \mathbb{R} \to v_{343} \in \mathbb{R} \to \neg v_{343} \leq v_{342} \to v_{342} < v_{343}$}
\UIC{$\varphi_{61}$}
\UIC{$\forall v_{342}\,(\varphi_{61})$}
\AXC{$\varphi_{249}$\lbl{A2}}
\BIC{$\forall v_{343}\,(y \in \mathbb{R} \to v_{343} \in \mathbb{R} \to \neg v_{343} \leq y \to y < v_{343})$}
\AXC{$\varphi_{171}$\lbl{A2}}
\BIC{$y \in \mathbb{R} \to 0 \in \mathbb{R} \to \neg 0 \leq y \to y < 0$}
\end{ptree}
\noindent$\Pi_{6}$:
\begin{ptree}
\AXC{$\stackrel{\Pi_{5}}{\vdots}$}\noLine
\UIC{$y \in \mathbb{R} \to 0 \in \mathbb{R} \to \neg 0 \leq y \to y < 0$}
\AXC{$(y \in \mathbb{R} \to 0 \in \mathbb{R} \to \neg 0 \leq y \to y < 0) \to \varphi_{267}$\lbl{A1}}
\BIC{$\varphi_{267}$}
\end{ptree}
\noindent$\Pi_{7}$:
\begin{ptree}
\AXC{$x \in \mathbb{R} \to y \in \mathbb{R} \to \neg 0 \leq x \to \neg 0 \leq y \to \neg 0 \leq y$\lbl{A1}}
\AXC{$0 \in \mathbb{R}$\lbl{(R1)}}
\AXC{$x \in \mathbb{R} \to y \in \mathbb{R} \to \neg 0 \leq x \to \neg 0 \leq y \to y \in \mathbb{R}$\lbl{A1}}
\AXC{$\stackrel{\Pi_{6}}{\vdots}$}\noLine
\UIC{$\varphi_{267}$}
\BIC{$0 \in \mathbb{R} \to \varphi_{161}$}
\BIC{$\varphi_{161}$}
\BIC{$x \in \mathbb{R} \to y \in \mathbb{R} \to \neg 0 \leq x \to \neg 0 \leq y \to y < 0$}
\end{ptree}
\noindent$\Pi_{8}$:
\begin{ptree}
\AXC{$\forall x\,(x < 0 \to \lvert x \rvert = -x)$\lbl{[A3]}}
\AXC{$\forall x\,(x < 0 \to \lvert x \rvert = -x) \to y < 0 \to \lvert y \rvert = -y$\lbl{A2}}
\BIC{$y < 0 \to \lvert y \rvert = -y$}
\AXC{$(y < 0 \to \lvert y \rvert = -y) \to \varphi_{187}$\lbl{A1}}
\BIC{$\varphi_{187}$}
\end{ptree}
\noindent$\Pi_{9}$:
\begin{ptree}
\AXC{$\stackrel{\Pi_{7}}{\vdots}$}\noLine
\UIC{$x \in \mathbb{R} \to y \in \mathbb{R} \to \neg 0 \leq x \to \neg 0 \leq y \to y < 0$}
\AXC{$\stackrel{\Pi_{8}}{\vdots}$}\noLine
\UIC{$\varphi_{187}$}
\BIC{$x \in \mathbb{R} \to y \in \mathbb{R} \to \neg 0 \leq x \to \neg 0 \leq y \to \lvert y \rvert = -y$}
\AXC{$\varphi_{314}$\lbl{A1}}
\BIC{$-x \cdot \lvert y \rvert = -x \cdot \lvert y \rvert \to \varphi_{135} \to \varphi_{107}$}
\end{ptree}
\noindent$\Pi_{10}$:
\begin{ptree}
\AXC{$x = x$\lbl{A5}}
\UIC{$\forall x\,(x = x)$}
\AXC{$\forall x\,(x = x) \to -x \cdot \lvert y \rvert = -x \cdot \lvert y \rvert$\lbl{A2}}
\BIC{$-x \cdot \lvert y \rvert = -x \cdot \lvert y \rvert$}
\AXC{$\stackrel{\Pi_{9}}{\vdots}$}\noLine
\UIC{$-x \cdot \lvert y \rvert = -x \cdot \lvert y \rvert \to \varphi_{135} \to \varphi_{107}$}
\BIC{$\varphi_{135} \to \varphi_{107}$}
\end{ptree}
\noindent$\Pi_{11}$:
\begin{ptree}
\AXC{$z = w \to -x \cdot \lvert y \rvert = -x \cdot z \to -x \cdot \lvert y \rvert = -x \cdot w$\lbl{A6}}
\UIC{$\varphi_{94}$}
\UIC{$\forall z\,(\varphi_{94})$}
\AXC{$\forall z\,(\varphi_{94}) \to \varphi_{127}$\lbl{A2}}
\BIC{$\varphi_{127}$}
\AXC{$\varphi_{127} \to \varphi_{135}$\lbl{A2}}
\BIC{$\varphi_{135}$}
\AXC{$\stackrel{\Pi_{10}}{\vdots}$}\noLine
\UIC{$\varphi_{135} \to \varphi_{107}$}
\BIC{$\varphi_{107}$}
\end{ptree}
\noindent$\Pi_{12}$:
\begin{ptree}
\AXC{$\forall x\,\forall y\,(x \in \mathbb{R} \to y \in \mathbb{R} \to \neg y \leq x \to x < y)$\lbl{[A2]}}
\AXC{$\varphi_{156}$\lbl{A2}}
\BIC{$\forall y\,(x \in \mathbb{R} \to y \in \mathbb{R} \to \neg y \leq x \to x < y)$}
\AXC{$\varphi_{129}$\lbl{A2}}
\BIC{$x \in \mathbb{R} \to 0 \in \mathbb{R} \to \neg 0 \leq x \to x < 0$}
\AXC{$(x \in \mathbb{R} \to 0 \in \mathbb{R} \to \neg 0 \leq x \to x < 0) \to \varphi_{265}$\lbl{A1}}
\BIC{$\varphi_{265}$}
\end{ptree}
\noindent$\Pi_{13}$:
\begin{ptree}
\AXC{$x \in \mathbb{R} \to y \in \mathbb{R} \to \neg 0 \leq x \to \neg 0 \leq x$\lbl{A1}}
\AXC{$0 \in \mathbb{R}$\lbl{(R1)}}
\AXC{$x \in \mathbb{R} \to y \in \mathbb{R} \to \neg 0 \leq x \to \neg 0 \leq y \to x \in \mathbb{R}$\lbl{A1}}
\AXC{$\stackrel{\Pi_{12}}{\vdots}$}\noLine
\UIC{$\varphi_{265}$}
\BIC{$0 \in \mathbb{R} \to \varphi_{144}$}
\BIC{$\varphi_{144}$}
\BIC{$x \in \mathbb{R} \to y \in \mathbb{R} \to \neg 0 \leq x \to \neg 0 \leq y \to x < 0$}
\end{ptree}
\noindent$\Pi_{14}$:
\begin{ptree}
\AXC{$\forall x\,(x < 0 \to \lvert x \rvert = -x)$\lbl{[A3]}}
\AXC{$\forall x\,(x < 0 \to \lvert x \rvert = -x) \to x < 0 \to \lvert x \rvert = -x$\lbl{A2}}
\BIC{$x < 0 \to \lvert x \rvert = -x$}
\AXC{$(x < 0 \to \lvert x \rvert = -x) \to \varphi_{188}$\lbl{A1}}
\BIC{$\varphi_{188}$}
\end{ptree}
\noindent$\Pi_{15}$:
\begin{ptree}
\AXC{$\stackrel{\Pi_{13}}{\vdots}$}\noLine
\UIC{$x \in \mathbb{R} \to y \in \mathbb{R} \to \neg 0 \leq x \to \neg 0 \leq y \to x < 0$}
\AXC{$\stackrel{\Pi_{14}}{\vdots}$}\noLine
\UIC{$\varphi_{188}$}
\BIC{$x \in \mathbb{R} \to y \in \mathbb{R} \to \neg 0 \leq x \to \neg 0 \leq y \to \lvert x \rvert = -x$}
\AXC{$\varphi_{315}$\lbl{A1}}
\BIC{$\lvert x \rvert \cdot \lvert y \rvert = \lvert x \rvert \cdot \lvert y \rvert \to \varphi_{136} \to \varphi_{108}$}
\end{ptree}
\noindent$\Pi_{16}$:
\begin{ptree}
\AXC{$x = x$\lbl{A5}}
\UIC{$\forall x\,(x = x)$}
\AXC{$\forall x\,(x = x) \to \lvert x \rvert \cdot \lvert y \rvert = \lvert x \rvert \cdot \lvert y \rvert$\lbl{A2}}
\BIC{$\lvert x \rvert \cdot \lvert y \rvert = \lvert x \rvert \cdot \lvert y \rvert$}
\AXC{$\stackrel{\Pi_{15}}{\vdots}$}\noLine
\UIC{$\lvert x \rvert \cdot \lvert y \rvert = \lvert x \rvert \cdot \lvert y \rvert \to \varphi_{136} \to \varphi_{108}$}
\BIC{$\varphi_{136} \to \varphi_{108}$}
\end{ptree}
\noindent$\Pi_{17}$:
\begin{ptree}
\AXC{$\stackrel{\Pi_{11}}{\vdots}$}\noLine
\UIC{$\varphi_{107}$}
\AXC{$\varphi_{66}$\lbl{A6}}
\UIC{$\forall w\,(\varphi_{66})$}
\UIC{$\forall z\,\forall w\,(\varphi_{66})$}
\AXC{$\forall z\,\forall w\,(\varphi_{66}) \to \varphi_{128}$\lbl{A2}}
\BIC{$\varphi_{128}$}
\AXC{$\varphi_{128} \to \varphi_{136}$\lbl{A2}}
\BIC{$\varphi_{136}$}
\AXC{$\stackrel{\Pi_{16}}{\vdots}$}\noLine
\UIC{$\varphi_{136} \to \varphi_{108}$}
\BIC{$\varphi_{108}$}
\AXC{$\varphi_{108} \to \varphi_{107} \to \varphi_{201} \to \varphi_{109}$\lbl{A1}}
\BIC{$\varphi_{107} \to \varphi_{201} \to \varphi_{109}$}
\BIC{$\varphi_{201} \to \varphi_{109}$}
\end{ptree}
\noindent$\Pi_{18}$:
\begin{ptree}
\AXC{$\stackrel{\Pi_{4}}{\vdots}$}\noLine
\UIC{$\forall v_{344}\,(\varphi_{9})$}
\AXC{$\forall v_{344}\,(\varphi_{9}) \to \varphi_{97}$\lbl{A2}}
\BIC{$\varphi_{97}$}
\AXC{$\varphi_{97} \to \varphi_{146}$\lbl{A2}}
\BIC{$\varphi_{146}$}
\AXC{$\varphi_{146} \to \varphi_{201}$\lbl{A2}}
\BIC{$\varphi_{201}$}
\AXC{$\stackrel{\Pi_{17}}{\vdots}$}\noLine
\UIC{$\varphi_{201} \to \varphi_{109}$}
\BIC{$\varphi_{109}$}
\AXC{$\varphi_{109} \to \varphi_{262}$\lbl{A1}}
\BIC{$\varphi_{262}$}
\end{ptree}
\noindent$\Pi_{19}$:
\begin{ptree}
\AXC{$y = z \to x = y \to x = z$\lbl{A6}}
\AXC{$(y = z \to x = y \to x = z) \to x = y \to y = z \to x = z$\lbl{A1}}
\BIC{$x = y \to y = z \to x = z$}
\UIC{$\forall z\,(x = y \to y = z \to x = z)$}
\UIC{$\forall y\,\forall z\,(x = y \to y = z \to x = z)$}
\UIC{$\forall x\,\forall y\,\forall z\,(x = y \to y = z \to x = z)$}
\AXC{$\varphi_{21}$\lbl{A2}}
\BIC{$\forall y\,\forall z\,(v_{353} = y \to y = z \to v_{353} = z)$}
\AXC{$\varphi_{43}$\lbl{A2}}
\BIC{$\forall z\,(v_{353} = v_{354} \to v_{354} = z \to v_{353} = z)$}
\AXC{$\varphi_{70}$\lbl{A2}}
\BIC{$v_{353} = v_{354} \to v_{354} = v_{355} \to v_{353} = v_{355}$}
\UIC{$\forall v_{355}\,(v_{353} = v_{354} \to v_{354} = v_{355} \to v_{353} = v_{355})$}
\UIC{$\varphi_{10}$}
\UIC{$\forall v_{353}\,(\varphi_{10})$}
\end{ptree}
\noindent$\Pi_{20}$:
\begin{ptree}
\AXC{$\stackrel{\Pi_{19}}{\vdots}$}\noLine
\UIC{$\forall v_{353}\,(\varphi_{10})$}
\AXC{$\forall v_{353}\,(\varphi_{10}) \to \varphi_{35}$\lbl{A2}}
\BIC{$\varphi_{35}$}
\AXC{$\varphi_{35} \to \varphi_{30}$\lbl{A2}}
\BIC{$\varphi_{30}$}
\AXC{$\varphi_{30} \to \varphi_{114}$\lbl{A2}}
\BIC{$\varphi_{114}$}
\end{ptree}
\noindent$\Pi_{21}$:
\begin{ptree}
\AXC{$\stackrel{\Pi_{2}}{\vdots}$}\noLine
\UIC{$\forall x\,\forall y\,(x = y \to y = x)$}
\AXC{$\forall x\,\forall y\,(x = y \to y = x) \to \forall y\,(v_{351} = y \to y = v_{351})$\lbl{A2}}
\BIC{$\forall y\,(v_{351} = y \to y = v_{351})$}
\AXC{$\varphi_{2}$\lbl{A2}}
\BIC{$v_{351} = v_{352} \to v_{352} = v_{351}$}
\UIC{$\forall v_{352}\,(v_{351} = v_{352} \to v_{352} = v_{351})$}
\UIC{$\forall v_{351}\,\forall v_{352}\,(v_{351} = v_{352} \to v_{352} = v_{351})$}
\AXC{$\varphi_{147}$\lbl{A2}}
\BIC{$\forall v_{352}\,(-x \cdot -y = v_{352} \to v_{352} = -x \cdot -y)$}
\AXC{$\varphi_{193}$\lbl{A2}}
\BIC{$-x \cdot -y = x \cdot y \to x \cdot y = -x \cdot -y$}
\end{ptree}
\noindent$\Pi_{22}$:
\begin{ptree}
\AXC{$\forall x\,\forall y\,(x \in \mathbb{R} \to y \in \mathbb{R} \to -x \cdot -y = x \cdot y)$\lbl{[F11]}}
\AXC{$\varphi_{218}$\lbl{A2}}
\BIC{$\forall y\,(x \in \mathbb{R} \to y \in \mathbb{R} \to -x \cdot -y = x \cdot y)$}
\AXC{$\varphi_{197}$\lbl{A2}}
\BIC{$x \in \mathbb{R} \to y \in \mathbb{R} \to -x \cdot -y = x \cdot y$}
\AXC{$\varphi_{307}$\lbl{A1}}
\BIC{$(x \in \mathbb{R} \to y \in \mathbb{R} \to \neg 0 \leq x \to \neg 0 \leq y \to x \in \mathbb{R}) \to \varphi_{239}$}
\end{ptree}
\noindent$\Pi_{23}$:
\begin{ptree}
\AXC{$x \in \mathbb{R} \to y \in \mathbb{R} \to \neg 0 \leq x \to \neg 0 \leq y \to x \in \mathbb{R}$\lbl{A1}}
\AXC{$\stackrel{\Pi_{22}}{\vdots}$}\noLine
\UIC{$(x \in \mathbb{R} \to y \in \mathbb{R} \to \neg 0 \leq x \to \neg 0 \leq y \to x \in \mathbb{R}) \to \varphi_{239}$}
\BIC{$\varphi_{239}$}
\end{ptree}
\noindent$\Pi_{24}$:
\begin{ptree}
\AXC{$\stackrel{\Pi_{21}}{\vdots}$}\noLine
\UIC{$-x \cdot -y = x \cdot y \to x \cdot y = -x \cdot -y$}
\AXC{$x \in \mathbb{R} \to y \in \mathbb{R} \to \neg 0 \leq x \to \neg 0 \leq y \to y \in \mathbb{R}$\lbl{A1}}
\AXC{$\stackrel{\Pi_{23}}{\vdots}$}\noLine
\UIC{$\varphi_{239}$}
\BIC{$x \in \mathbb{R} \to y \in \mathbb{R} \to \neg 0 \leq x \to \neg 0 \leq y \to -x \cdot -y = x \cdot y$}
\AXC{$\varphi_{287}$\lbl{A1}}
\BIC{$\varphi_{238}$}
\BIC{$x \in \mathbb{R} \to y \in \mathbb{R} \to \neg 0 \leq x \to \neg 0 \leq y \to x \cdot y = -x \cdot -y$}
\end{ptree}
\noindent$\Pi_{25}$:
\begin{ptree}
\AXC{$\forall x\,\forall y\,(x \in \mathbb{R} \to y \in \mathbb{R} \to \neg y \leq x \to x < y)$\lbl{[A2]}}
\AXC{$\varphi_{180}$\lbl{A2}}
\BIC{$\forall y\,(v_{347} \in \mathbb{R} \to y \in \mathbb{R} \to \neg y \leq v_{347} \to v_{347} < y)$}
\AXC{$\varphi_{224}$\lbl{A2}}
\BIC{$v_{347} \in \mathbb{R} \to v_{348} \in \mathbb{R} \to \neg v_{348} \leq v_{347} \to v_{347} < v_{348}$}
\UIC{$\varphi_{62}$}
\UIC{$\forall v_{347}\,(\varphi_{62})$}
\AXC{$\varphi_{250}$\lbl{A2}}
\BIC{$\forall v_{348}\,(y \in \mathbb{R} \to v_{348} \in \mathbb{R} \to \neg v_{348} \leq y \to y < v_{348})$}
\AXC{$\varphi_{172}$\lbl{A2}}
\BIC{$y \in \mathbb{R} \to 0 \in \mathbb{R} \to \neg 0 \leq y \to y < 0$}
\end{ptree}
\noindent$\Pi_{26}$:
\begin{ptree}
\AXC{$\stackrel{\Pi_{25}}{\vdots}$}\noLine
\UIC{$y \in \mathbb{R} \to 0 \in \mathbb{R} \to \neg 0 \leq y \to y < 0$}
\AXC{$(y \in \mathbb{R} \to 0 \in \mathbb{R} \to \neg 0 \leq y \to y < 0) \to \varphi_{267}$\lbl{A1}}
\BIC{$\varphi_{267}$}
\end{ptree}
\noindent$\Pi_{27}$:
\begin{ptree}
\AXC{$x \in \mathbb{R} \to y \in \mathbb{R} \to \neg 0 \leq x \to \neg 0 \leq y \to \neg 0 \leq y$\lbl{A1}}
\AXC{$0 \in \mathbb{R}$\lbl{(R1)}}
\AXC{$x \in \mathbb{R} \to y \in \mathbb{R} \to \neg 0 \leq x \to \neg 0 \leq y \to y \in \mathbb{R}$\lbl{A1}}
\AXC{$\stackrel{\Pi_{26}}{\vdots}$}\noLine
\UIC{$\varphi_{267}$}
\BIC{$0 \in \mathbb{R} \to \varphi_{161}$}
\BIC{$\varphi_{161}$}
\BIC{$x \in \mathbb{R} \to y \in \mathbb{R} \to \neg 0 \leq x \to \neg 0 \leq y \to y < 0$}
\end{ptree}
\noindent$\Pi_{28}$:
\begin{ptree}
\AXC{$\forall x\,\forall y\,(x \in \mathbb{R} \to y \in \mathbb{R} \to x < y \to -y < -x)$\lbl{[F22]}}
\AXC{$\varphi_{247}$\lbl{A2}}
\BIC{$\forall y\,(v_{349} \in \mathbb{R} \to y \in \mathbb{R} \to v_{349} < y \to -y < -v_{349})$}
\AXC{$\varphi_{256}$\lbl{A2}}
\BIC{$\varphi_{98}$}
\UIC{$\forall v_{350}\,(\varphi_{98})$}
\UIC{$\forall v_{349}\,\forall v_{350}\,(\varphi_{98})$}
\AXC{$\forall v_{349}\,\forall v_{350}\,(\varphi_{98}) \to \varphi_{71}$\lbl{A2}}
\BIC{$\varphi_{71}$}
\end{ptree}
\noindent$\Pi_{29}$:
\begin{ptree}
\AXC{$\stackrel{\Pi_{28}}{\vdots}$}\noLine
\UIC{$\varphi_{71}$}
\AXC{$\varphi_{71} \to y \in \mathbb{R} \to 0 \in \mathbb{R} \to y < 0 \to -0 < -y$\lbl{A2}}
\BIC{$y \in \mathbb{R} \to 0 \in \mathbb{R} \to y < 0 \to -0 < -y$}
\AXC{$(y \in \mathbb{R} \to 0 \in \mathbb{R} \to y < 0 \to -0 < -y) \to \varphi_{278}$\lbl{A1}}
\BIC{$\varphi_{278}$}
\end{ptree}
\noindent$\Pi_{30}$:
\begin{ptree}
\AXC{$\stackrel{\Pi_{27}}{\vdots}$}\noLine
\UIC{$x \in \mathbb{R} \to y \in \mathbb{R} \to \neg 0 \leq x \to \neg 0 \leq y \to y < 0$}
\AXC{$0 \in \mathbb{R}$\lbl{(R1)}}
\AXC{$x \in \mathbb{R} \to y \in \mathbb{R} \to \neg 0 \leq x \to \neg 0 \leq y \to y \in \mathbb{R}$\lbl{A1}}
\AXC{$\stackrel{\Pi_{29}}{\vdots}$}\noLine
\UIC{$\varphi_{278}$}
\BIC{$0 \in \mathbb{R} \to \varphi_{202}$}
\BIC{$\varphi_{202}$}
\BIC{$x \in \mathbb{R} \to y \in \mathbb{R} \to \neg 0 \leq x \to \neg 0 \leq y \to -0 < -y$}
\end{ptree}
\noindent$\Pi_{31}$:
\begin{ptree}
\AXC{$\stackrel{\Pi_{30}}{\vdots}$}\noLine
\UIC{$x \in \mathbb{R} \to y \in \mathbb{R} \to \neg 0 \leq x \to \neg 0 \leq y \to -0 < -y$}
\AXC{$-0 = 0$\lbl{[F6]}}
\AXC{$\varphi_{272}$\lbl{A1}}
\BIC{$(x \in \mathbb{R} \to y \in \mathbb{R} \to \neg 0 \leq x \to \neg 0 \leq y \to -0 < -y) \to \varphi_{168}$}
\BIC{$\varphi_{168}$}
\end{ptree}
\noindent$\Pi_{32}$:
\begin{ptree}
\AXC{$\forall x\,\forall y\,(x \in \mathbb{R} \to y \in \mathbb{R} \to \neg y \leq x \to x < y)$\lbl{[A2]}}
\AXC{$\varphi_{156}$\lbl{A2}}
\BIC{$\forall y\,(x \in \mathbb{R} \to y \in \mathbb{R} \to \neg y \leq x \to x < y)$}
\AXC{$\varphi_{129}$\lbl{A2}}
\BIC{$x \in \mathbb{R} \to 0 \in \mathbb{R} \to \neg 0 \leq x \to x < 0$}
\AXC{$(x \in \mathbb{R} \to 0 \in \mathbb{R} \to \neg 0 \leq x \to x < 0) \to \varphi_{265}$\lbl{A1}}
\BIC{$\varphi_{265}$}
\end{ptree}
\noindent$\Pi_{33}$:
\begin{ptree}
\AXC{$x \in \mathbb{R} \to y \in \mathbb{R} \to \neg 0 \leq x \to \neg 0 \leq x$\lbl{A1}}
\AXC{$0 \in \mathbb{R}$\lbl{(R1)}}
\AXC{$x \in \mathbb{R} \to y \in \mathbb{R} \to \neg 0 \leq x \to \neg 0 \leq y \to x \in \mathbb{R}$\lbl{A1}}
\AXC{$\stackrel{\Pi_{32}}{\vdots}$}\noLine
\UIC{$\varphi_{265}$}
\BIC{$0 \in \mathbb{R} \to \varphi_{144}$}
\BIC{$\varphi_{144}$}
\BIC{$x \in \mathbb{R} \to y \in \mathbb{R} \to \neg 0 \leq x \to \neg 0 \leq y \to x < 0$}
\end{ptree}
\noindent$\Pi_{34}$:
\begin{ptree}
\AXC{$\forall x\,\forall y\,(x \in \mathbb{R} \to y \in \mathbb{R} \to x < y \to -y < -x)$\lbl{[F22]}}
\AXC{$\varphi_{232}$\lbl{A2}}
\BIC{$\forall y\,(x \in \mathbb{R} \to y \in \mathbb{R} \to x < y \to -y < -x)$}
\AXC{$\varphi_{211}$\lbl{A2}}
\BIC{$x \in \mathbb{R} \to 0 \in \mathbb{R} \to x < 0 \to -0 < -x$}
\AXC{$(x \in \mathbb{R} \to 0 \in \mathbb{R} \to x < 0 \to -0 < -x) \to \varphi_{279}$\lbl{A1}}
\BIC{$\varphi_{279}$}
\end{ptree}
\noindent$\Pi_{35}$:
\begin{ptree}
\AXC{$\stackrel{\Pi_{33}}{\vdots}$}\noLine
\UIC{$x \in \mathbb{R} \to y \in \mathbb{R} \to \neg 0 \leq x \to \neg 0 \leq y \to x < 0$}
\AXC{$0 \in \mathbb{R}$\lbl{(R1)}}
\AXC{$x \in \mathbb{R} \to y \in \mathbb{R} \to \neg 0 \leq x \to \neg 0 \leq y \to x \in \mathbb{R}$\lbl{A1}}
\AXC{$\stackrel{\Pi_{34}}{\vdots}$}\noLine
\UIC{$\varphi_{279}$}
\BIC{$0 \in \mathbb{R} \to \varphi_{203}$}
\BIC{$\varphi_{203}$}
\BIC{$x \in \mathbb{R} \to y \in \mathbb{R} \to \neg 0 \leq x \to \neg 0 \leq y \to -0 < -x$}
\end{ptree}
\noindent$\Pi_{36}$:
\begin{ptree}
\AXC{$\stackrel{\Pi_{35}}{\vdots}$}\noLine
\UIC{$x \in \mathbb{R} \to y \in \mathbb{R} \to \neg 0 \leq x \to \neg 0 \leq y \to -0 < -x$}
\AXC{$-0 = 0$\lbl{[F6]}}
\AXC{$\varphi_{273}$\lbl{A1}}
\BIC{$(x \in \mathbb{R} \to y \in \mathbb{R} \to \neg 0 \leq x \to \neg 0 \leq y \to -0 < -x) \to \varphi_{169}$}
\BIC{$\varphi_{169}$}
\end{ptree}
\noindent$\Pi_{37}$:
\begin{ptree}
\AXC{$z = w \to z < -x \to w < -x$\lbl{A6}}
\UIC{$\forall w\,(z = w \to z < -x \to w < -x)$}
\UIC{$\forall z\,\forall w\,(z = w \to z < -x \to w < -x)$}
\AXC{$\varphi_{132}$\lbl{A2}}
\BIC{$\forall w\,(-0 = w \to -0 < -x \to w < -x)$}
\AXC{$\varphi_{153}$\lbl{A2}}
\BIC{$-0 = 0 \to -0 < -x \to 0 < -x$}
\AXC{$\stackrel{\Pi_{36}}{\vdots}$}\noLine
\UIC{$\varphi_{169}$}
\BIC{$x \in \mathbb{R} \to y \in \mathbb{R} \to \neg 0 \leq x \to \neg 0 \leq y \to 0 < -x$}
\AXC{$\varphi_{190}$\lbl{A1}}
\BIC{$x \in \mathbb{R} \to y \in \mathbb{R} \to \neg 0 \leq x \to \neg 0 \leq y \to 0 \leq -x$}
\end{ptree}
\noindent$\Pi_{38}$:
\begin{ptree}
\AXC{$\forall x \in \mathbb{R}\,(-x \in \mathbb{R})$\lbl{(R1)}}
\AXC{$\forall x \in \mathbb{R}\,(-x \in \mathbb{R}) \to y \in \mathbb{R} \to -y \in \mathbb{R}$\lbl{A2}}
\BIC{$y \in \mathbb{R} \to -y \in \mathbb{R}$}
\AXC{$(y \in \mathbb{R} \to -y \in \mathbb{R}) \to \varphi_{212}$\lbl{A1}}
\BIC{$\varphi_{212}$}
\end{ptree}
\noindent$\Pi_{39}$:
\begin{ptree}
\AXC{$\forall x \in \mathbb{R}\,(-x \in \mathbb{R})$\lbl{(R1)}}
\AXC{$\forall x \in \mathbb{R}\,(-x \in \mathbb{R}) \to x \in \mathbb{R} \to -x \in \mathbb{R}$\lbl{A2}}
\BIC{$x \in \mathbb{R} \to -x \in \mathbb{R}$}
\AXC{$(x \in \mathbb{R} \to -x \in \mathbb{R}) \to \varphi_{213}$\lbl{A1}}
\BIC{$\varphi_{213}$}
\end{ptree}
\noindent$\Pi_{40}$:
\begin{ptree}
\AXC{$\forall x\,\forall y\,(x \in \mathbb{R} \to y \in \mathbb{R} \to 0 \leq x \to 0 \leq y \to 0 \leq x \cdot y)$\lbl{[A16]}}
\AXC{$\varphi_{245}$\lbl{A2}}
\BIC{$\forall y\,(-x \in \mathbb{R} \to y \in \mathbb{R} \to 0 \leq -x \to 0 \leq y \to 0 \leq -x \cdot y)$}
\AXC{$\varphi_{261}$\lbl{A2}}
\BIC{$-x \in \mathbb{R} \to -y \in \mathbb{R} \to 0 \leq -x \to 0 \leq -y \to 0 \leq -x \cdot -y$}
\AXC{$\varphi_{328}$\lbl{A1}}
\BIC{$(x \in \mathbb{R} \to y \in \mathbb{R} \to \neg 0 \leq x \to \neg 0 \leq y \to -x \in \mathbb{R}) \to \varphi_{303}$}
\end{ptree}
\noindent$\Pi_{41}$:
\begin{ptree}
\AXC{$x \in \mathbb{R} \to y \in \mathbb{R} \to \neg 0 \leq x \to \neg 0 \leq y \to x \in \mathbb{R}$\lbl{A1}}
\AXC{$\stackrel{\Pi_{39}}{\vdots}$}\noLine
\UIC{$\varphi_{213}$}
\BIC{$x \in \mathbb{R} \to y \in \mathbb{R} \to \neg 0 \leq x \to \neg 0 \leq y \to -x \in \mathbb{R}$}
\AXC{$\stackrel{\Pi_{40}}{\vdots}$}\noLine
\UIC{$(x \in \mathbb{R} \to y \in \mathbb{R} \to \neg 0 \leq x \to \neg 0 \leq y \to -x \in \mathbb{R}) \to \varphi_{303}$}
\BIC{$\varphi_{303}$}
\end{ptree}
\noindent$\Pi_{42}$:
\begin{ptree}
\AXC{$\stackrel{\Pi_{37}}{\vdots}$}\noLine
\UIC{$x \in \mathbb{R} \to y \in \mathbb{R} \to \neg 0 \leq x \to \neg 0 \leq y \to 0 \leq -x$}
\AXC{$x \in \mathbb{R} \to y \in \mathbb{R} \to \neg 0 \leq x \to \neg 0 \leq y \to y \in \mathbb{R}$\lbl{A1}}
\AXC{$\stackrel{\Pi_{38}}{\vdots}$}\noLine
\UIC{$\varphi_{212}$}
\BIC{$x \in \mathbb{R} \to y \in \mathbb{R} \to \neg 0 \leq x \to \neg 0 \leq y \to -y \in \mathbb{R}$}
\AXC{$\stackrel{\Pi_{41}}{\vdots}$}\noLine
\UIC{$\varphi_{303}$}
\BIC{$(x \in \mathbb{R} \to y \in \mathbb{R} \to \neg 0 \leq x \to \neg 0 \leq y \to 0 \leq -x) \to \varphi_{216}$}
\BIC{$\varphi_{216}$}
\end{ptree}
\noindent$\Pi_{43}$:
\begin{ptree}
\AXC{$z = w \to z < -y \to w < -y$\lbl{A6}}
\UIC{$\forall w\,(z = w \to z < -y \to w < -y)$}
\UIC{$\forall z\,\forall w\,(z = w \to z < -y \to w < -y)$}
\AXC{$\varphi_{131}$\lbl{A2}}
\BIC{$\forall w\,(-0 = w \to -0 < -y \to w < -y)$}
\AXC{$\varphi_{152}$\lbl{A2}}
\BIC{$-0 = 0 \to -0 < -y \to 0 < -y$}
\AXC{$\stackrel{\Pi_{31}}{\vdots}$}\noLine
\UIC{$\varphi_{168}$}
\BIC{$x \in \mathbb{R} \to y \in \mathbb{R} \to \neg 0 \leq x \to \neg 0 \leq y \to 0 < -y$}
\AXC{$\varphi_{189}$\lbl{A1}}
\BIC{$x \in \mathbb{R} \to y \in \mathbb{R} \to \neg 0 \leq x \to \neg 0 \leq y \to 0 \leq -y$}
\AXC{$\stackrel{\Pi_{42}}{\vdots}$}\noLine
\UIC{$\varphi_{216}$}
\BIC{$x \in \mathbb{R} \to y \in \mathbb{R} \to \neg 0 \leq x \to \neg 0 \leq y \to 0 \leq -x \cdot -y$}
\end{ptree}
\noindent$\Pi_{44}$:
\begin{ptree}
\AXC{$\forall x\,\forall y\,(x \in \mathbb{R} \to y \in \mathbb{R} \to -x \cdot -y = x \cdot y)$\lbl{[F11]}}
\AXC{$\varphi_{218}$\lbl{A2}}
\BIC{$\forall y\,(x \in \mathbb{R} \to y \in \mathbb{R} \to -x \cdot -y = x \cdot y)$}
\AXC{$\varphi_{197}$\lbl{A2}}
\BIC{$x \in \mathbb{R} \to y \in \mathbb{R} \to -x \cdot -y = x \cdot y$}
\AXC{$\varphi_{307}$\lbl{A1}}
\BIC{$(x \in \mathbb{R} \to y \in \mathbb{R} \to \neg 0 \leq x \to \neg 0 \leq y \to x \in \mathbb{R}) \to \varphi_{239}$}
\end{ptree}
\noindent$\Pi_{45}$:
\begin{ptree}
\AXC{$x \in \mathbb{R} \to y \in \mathbb{R} \to \neg 0 \leq x \to \neg 0 \leq y \to x \in \mathbb{R}$\lbl{A1}}
\AXC{$\stackrel{\Pi_{44}}{\vdots}$}\noLine
\UIC{$(x \in \mathbb{R} \to y \in \mathbb{R} \to \neg 0 \leq x \to \neg 0 \leq y \to x \in \mathbb{R}) \to \varphi_{239}$}
\BIC{$\varphi_{239}$}
\end{ptree}
\noindent$\Pi_{46}$:
\begin{ptree}
\AXC{$\stackrel{\Pi_{43}}{\vdots}$}\noLine
\UIC{$x \in \mathbb{R} \to y \in \mathbb{R} \to \neg 0 \leq x \to \neg 0 \leq y \to 0 \leq -x \cdot -y$}
\AXC{$x \in \mathbb{R} \to y \in \mathbb{R} \to \neg 0 \leq x \to \neg 0 \leq y \to y \in \mathbb{R}$\lbl{A1}}
\AXC{$\stackrel{\Pi_{45}}{\vdots}$}\noLine
\UIC{$\varphi_{239}$}
\BIC{$x \in \mathbb{R} \to y \in \mathbb{R} \to \neg 0 \leq x \to \neg 0 \leq y \to -x \cdot -y = x \cdot y$}
\AXC{$\varphi_{324}$\lbl{A1}}
\BIC{$\varphi_{309}$}
\BIC{$(-x \cdot -y = x \cdot y \to 0 < -x \cdot -y \to 0 < x \cdot y) \to \varphi_{206}$}
\end{ptree}
\noindent$\Pi_{47}$:
\begin{ptree}
\AXC{$z = w \to 0 < z \to 0 < w$\lbl{A6}}
\UIC{$\forall w\,(z = w \to 0 < z \to 0 < w)$}
\UIC{$\forall z\,\forall w\,(z = w \to 0 < z \to 0 < w)$}
\AXC{$\varphi_{130}$\lbl{A2}}
\BIC{$\forall w\,(-x \cdot -y = w \to 0 < -x \cdot -y \to 0 < w)$}
\AXC{$\varphi_{244}$\lbl{A2}}
\BIC{$-x \cdot -y = x \cdot y \to 0 < -x \cdot -y \to 0 < x \cdot y$}
\AXC{$\stackrel{\Pi_{46}}{\vdots}$}\noLine
\UIC{$(-x \cdot -y = x \cdot y \to 0 < -x \cdot -y \to 0 < x \cdot y) \to \varphi_{206}$}
\BIC{$\varphi_{206}$}
\end{ptree}
\noindent$\Pi_{48}$:
\begin{ptree}
\AXC{$\forall x\,(0 \leq x \to \lvert x \rvert = x)$\lbl{[A1]}}
\AXC{$\forall x\,(0 \leq x \to \lvert x \rvert = x) \to 0 \leq x \cdot y \to \lvert x \cdot y \rvert = x \cdot y$\lbl{A2}}
\BIC{$0 \leq x \cdot y \to \lvert x \cdot y \rvert = x \cdot y$}
\end{ptree}
\noindent$\Pi_{49}$:
\begin{ptree}
\AXC{$\stackrel{\Pi_{48}}{\vdots}$}\noLine
\UIC{$0 \leq x \cdot y \to \lvert x \cdot y \rvert = x \cdot y$}
\AXC{$(0 \leq x \cdot y \to \lvert x \cdot y \rvert = x \cdot y) \to \varphi_{217}$\lbl{A1}}
\BIC{$\varphi_{217}$}
\end{ptree}
\noindent$\Pi_{50}$:
\begin{ptree}
\AXC{$z = w \to 0 = z \to 0 = w$\lbl{A6}}
\UIC{$\forall w\,(z = w \to 0 = z \to 0 = w)$}
\UIC{$\forall z\,\forall w\,(z = w \to 0 = z \to 0 = w)$}
\AXC{$\varphi_{104}$\lbl{A2}}
\BIC{$\forall w\,(-x \cdot -y = w \to 0 = -x \cdot -y \to 0 = w)$}
\AXC{$\varphi_{205}$\lbl{A2}}
\BIC{$-x \cdot -y = x \cdot y \to 0 = -x \cdot -y \to 0 = x \cdot y$}
\AXC{$\stackrel{\Pi_{47}}{\vdots}$}\noLine
\UIC{$\varphi_{206}$}
\BIC{$x \in \mathbb{R} \to y \in \mathbb{R} \to \neg 0 \leq x \to \neg 0 \leq y \to 0 \leq x \cdot y$}
\AXC{$\stackrel{\Pi_{49}}{\vdots}$}\noLine
\UIC{$\varphi_{217}$}
\BIC{$x \in \mathbb{R} \to y \in \mathbb{R} \to \neg 0 \leq x \to \neg 0 \leq y \to \lvert x \cdot y \rvert = x \cdot y$}
\AXC{$\varphi_{319}$\lbl{A1}}
\BIC{$\varphi_{302}$}
\end{ptree}
\noindent$\Pi_{51}$:
\begin{ptree}
\AXC{$\stackrel{\Pi_{20}}{\vdots}$}\noLine
\UIC{$\varphi_{114}$}
\AXC{$\stackrel{\Pi_{24}}{\vdots}$}\noLine
\UIC{$x \in \mathbb{R} \to y \in \mathbb{R} \to \neg 0 \leq x \to \neg 0 \leq y \to x \cdot y = -x \cdot -y$}
\AXC{$\stackrel{\Pi_{50}}{\vdots}$}\noLine
\UIC{$\varphi_{302}$}
\BIC{$\varphi_{114} \to \varphi_{95}$}
\BIC{$\varphi_{95}$}
\AXC{$\varphi_{95} \to \varphi_{110} \to \varphi_{167} \to \varphi_{96}$\lbl{A1}}
\BIC{$\varphi_{110} \to \varphi_{167} \to \varphi_{96}$}
\end{ptree}
\noindent$\Pi_{52}$:
\begin{ptree}
\AXC{$\stackrel{\Pi_{3}}{\vdots}$}\noLine
\UIC{$\lvert x \rvert \cdot \lvert y \rvert = -x \cdot -y \to -x \cdot -y = \lvert x \rvert \cdot \lvert y \rvert$}
\AXC{$\stackrel{\Pi_{18}}{\vdots}$}\noLine
\UIC{$\varphi_{262}$}
\BIC{$\varphi_{110}$}
\AXC{$\stackrel{\Pi_{51}}{\vdots}$}\noLine
\UIC{$\varphi_{110} \to \varphi_{167} \to \varphi_{96}$}
\BIC{$\varphi_{167} \to \varphi_{96}$}
\end{ptree}
\noindent$\Pi_{53}$:
\begin{ptree}
\AXC{$\stackrel{\Pi_{1}}{\vdots}$}\noLine
\UIC{$\forall v_{358}\,(\varphi_{8})$}
\AXC{$\forall v_{358}\,(\varphi_{8}) \to \varphi_{34}$\lbl{A2}}
\BIC{$\varphi_{34}$}
\AXC{$\varphi_{34} \to \varphi_{122}$\lbl{A2}}
\BIC{$\varphi_{122}$}
\AXC{$\varphi_{122} \to \varphi_{167}$\lbl{A2}}
\BIC{$\varphi_{167}$}
\AXC{$\stackrel{\Pi_{52}}{\vdots}$}\noLine
\UIC{$\varphi_{167} \to \varphi_{96}$}
\BIC{$\varphi_{96}$}
\end{ptree}
\noindent$\Pi_{54}$:
\begin{ptree}
\AXC{$y = z \to x = y \to x = z$\lbl{A6}}
\AXC{$(y = z \to x = y \to x = z) \to x = y \to y = z \to x = z$\lbl{A1}}
\BIC{$x = y \to y = z \to x = z$}
\UIC{$\forall z\,(x = y \to y = z \to x = z)$}
\UIC{$\forall y\,\forall z\,(x = y \to y = z \to x = z)$}
\UIC{$\forall x\,\forall y\,\forall z\,(x = y \to y = z \to x = z)$}
\AXC{$\varphi_{22}$\lbl{A2}}
\BIC{$\forall y\,\forall z\,(v_{339} = y \to y = z \to v_{339} = z)$}
\AXC{$\varphi_{44}$\lbl{A2}}
\BIC{$\forall z\,(v_{339} = v_{340} \to v_{340} = z \to v_{339} = z)$}
\AXC{$\varphi_{72}$\lbl{A2}}
\BIC{$v_{339} = v_{340} \to v_{340} = v_{341} \to v_{339} = v_{341}$}
\UIC{$\forall v_{341}\,(v_{339} = v_{340} \to v_{340} = v_{341} \to v_{339} = v_{341})$}
\UIC{$\varphi_{11}$}
\UIC{$\forall v_{339}\,(\varphi_{11})$}
\end{ptree}
\noindent$\Pi_{55}$:
\begin{ptree}
\AXC{$\stackrel{\Pi_{2}}{\vdots}$}\noLine
\UIC{$\forall x\,\forall y\,(x = y \to y = x)$}
\AXC{$\forall x\,\forall y\,(x = y \to y = x) \to \forall y\,(v_{337} = y \to y = v_{337})$\lbl{A2}}
\BIC{$\forall y\,(v_{337} = y \to y = v_{337})$}
\AXC{$\varphi_{3}$\lbl{A2}}
\BIC{$v_{337} = v_{338} \to v_{338} = v_{337}$}
\UIC{$\forall v_{338}\,(v_{337} = v_{338} \to v_{338} = v_{337})$}
\UIC{$\forall v_{337}\,\forall v_{338}\,(v_{337} = v_{338} \to v_{338} = v_{337})$}
\AXC{$\varphi_{148}$\lbl{A2}}
\BIC{$\forall v_{338}\,(\lvert x \rvert \cdot \lvert y \rvert = v_{338} \to v_{338} = \lvert x \rvert \cdot \lvert y \rvert)$}
\AXC{$\varphi_{228}$\lbl{A2}}
\BIC{$\lvert x \rvert \cdot \lvert y \rvert = -x \cdot y \to -x \cdot y = \lvert x \rvert \cdot \lvert y \rvert$}
\end{ptree}
\noindent$\Pi_{56}$:
\begin{ptree}
\AXC{$y = z \to x = y \to x = z$\lbl{A6}}
\AXC{$(y = z \to x = y \to x = z) \to x = y \to y = z \to x = z$\lbl{A1}}
\BIC{$x = y \to y = z \to x = z$}
\UIC{$\forall z\,(x = y \to y = z \to x = z)$}
\UIC{$\forall y\,\forall z\,(x = y \to y = z \to x = z)$}
\UIC{$\forall x\,\forall y\,\forall z\,(x = y \to y = z \to x = z)$}
\AXC{$\varphi_{23}$\lbl{A2}}
\BIC{$\forall y\,\forall z\,(v_{329} = y \to y = z \to v_{329} = z)$}
\AXC{$\varphi_{45}$\lbl{A2}}
\BIC{$\forall z\,(v_{329} = v_{330} \to v_{330} = z \to v_{329} = z)$}
\AXC{$\varphi_{73}$\lbl{A2}}
\BIC{$v_{329} = v_{330} \to v_{330} = v_{331} \to v_{329} = v_{331}$}
\UIC{$\forall v_{331}\,(v_{329} = v_{330} \to v_{330} = v_{331} \to v_{329} = v_{331})$}
\UIC{$\varphi_{12}$}
\UIC{$\forall v_{329}\,(\varphi_{12})$}
\end{ptree}
\noindent$\Pi_{57}$:
\begin{ptree}
\AXC{$\stackrel{\Pi_{56}}{\vdots}$}\noLine
\UIC{$\forall v_{329}\,(\varphi_{12})$}
\AXC{$\forall v_{329}\,(\varphi_{12}) \to \varphi_{99}$\lbl{A2}}
\BIC{$\varphi_{99}$}
\AXC{$\varphi_{99} \to \varphi_{149}$\lbl{A2}}
\BIC{$\varphi_{149}$}
\AXC{$\varphi_{149} \to \varphi_{170}$\lbl{A2}}
\BIC{$\varphi_{170}$}
\end{ptree}
\noindent$\Pi_{58}$:
\begin{ptree}
\AXC{$\forall x\,(0 \leq x \to \lvert x \rvert = x)$\lbl{[A1]}}
\AXC{$\forall x\,(0 \leq x \to \lvert x \rvert = x) \to 0 \leq y \to \lvert y \rvert = y$\lbl{A2}}
\BIC{$0 \leq y \to \lvert y \rvert = y$}
\AXC{$(0 \leq y \to \lvert y \rvert = y) \to \varphi_{160}$\lbl{A1}}
\BIC{$\varphi_{160}$}
\end{ptree}
\noindent$\Pi_{59}$:
\begin{ptree}
\AXC{$x \in \mathbb{R} \to y \in \mathbb{R} \to \neg 0 \leq x \to 0 \leq y \to 0 \leq y$\lbl{A1}}
\AXC{$\stackrel{\Pi_{58}}{\vdots}$}\noLine
\UIC{$\varphi_{160}$}
\BIC{$x \in \mathbb{R} \to y \in \mathbb{R} \to \neg 0 \leq x \to 0 \leq y \to \lvert y \rvert = y$}
\AXC{$\varphi_{310}$\lbl{A1}}
\BIC{$-x \cdot \lvert y \rvert = -x \cdot \lvert y \rvert \to \varphi_{263}$}
\end{ptree}
\noindent$\Pi_{60}$:
\begin{ptree}
\AXC{$x = x$\lbl{A5}}
\UIC{$\forall x\,(x = x)$}
\AXC{$\forall x\,(x = x) \to -x \cdot \lvert y \rvert = -x \cdot \lvert y \rvert$\lbl{A2}}
\BIC{$-x \cdot \lvert y \rvert = -x \cdot \lvert y \rvert$}
\AXC{$\stackrel{\Pi_{59}}{\vdots}$}\noLine
\UIC{$-x \cdot \lvert y \rvert = -x \cdot \lvert y \rvert \to \varphi_{263}$}
\BIC{$\varphi_{263}$}
\end{ptree}
\noindent$\Pi_{61}$:
\begin{ptree}
\AXC{$z = w \to -x \cdot \lvert y \rvert = -x \cdot z \to -x \cdot \lvert y \rvert = -x \cdot w$\lbl{A6}}
\UIC{$\varphi_{94}$}
\UIC{$\forall z\,(\varphi_{94})$}
\AXC{$\forall z\,(\varphi_{94}) \to \varphi_{127}$\lbl{A2}}
\BIC{$\varphi_{127}$}
\AXC{$\varphi_{127} \to \varphi_{115}$\lbl{A2}}
\BIC{$\varphi_{115}$}
\AXC{$\stackrel{\Pi_{60}}{\vdots}$}\noLine
\UIC{$\varphi_{263}$}
\BIC{$x \in \mathbb{R} \to y \in \mathbb{R} \to \neg 0 \leq x \to 0 \leq y \to -x \cdot \lvert y \rvert = -x \cdot y$}
\end{ptree}
\noindent$\Pi_{62}$:
\begin{ptree}
\AXC{$\forall x\,\forall y\,(x \in \mathbb{R} \to y \in \mathbb{R} \to \neg y \leq x \to x < y)$\lbl{[A2]}}
\AXC{$\varphi_{156}$\lbl{A2}}
\BIC{$\forall y\,(x \in \mathbb{R} \to y \in \mathbb{R} \to \neg y \leq x \to x < y)$}
\AXC{$\varphi_{129}$\lbl{A2}}
\BIC{$x \in \mathbb{R} \to 0 \in \mathbb{R} \to \neg 0 \leq x \to x < 0$}
\AXC{$(x \in \mathbb{R} \to 0 \in \mathbb{R} \to \neg 0 \leq x \to x < 0) \to \varphi_{264}$\lbl{A1}}
\BIC{$\varphi_{264}$}
\end{ptree}
\noindent$\Pi_{63}$:
\begin{ptree}
\AXC{$x \in \mathbb{R} \to y \in \mathbb{R} \to \neg 0 \leq x \to \neg 0 \leq x$\lbl{A1}}
\AXC{$0 \in \mathbb{R}$\lbl{(R1)}}
\AXC{$x \in \mathbb{R} \to y \in \mathbb{R} \to \neg 0 \leq x \to 0 \leq y \to x \in \mathbb{R}$\lbl{A1}}
\AXC{$\stackrel{\Pi_{62}}{\vdots}$}\noLine
\UIC{$\varphi_{264}$}
\BIC{$0 \in \mathbb{R} \to \varphi_{138}$}
\BIC{$\varphi_{138}$}
\BIC{$x \in \mathbb{R} \to y \in \mathbb{R} \to \neg 0 \leq x \to 0 \leq y \to x < 0$}
\end{ptree}
\noindent$\Pi_{64}$:
\begin{ptree}
\AXC{$\forall x\,(x < 0 \to \lvert x \rvert = -x)$\lbl{[A3]}}
\AXC{$\forall x\,(x < 0 \to \lvert x \rvert = -x) \to x < 0 \to \lvert x \rvert = -x$\lbl{A2}}
\BIC{$x < 0 \to \lvert x \rvert = -x$}
\AXC{$(x < 0 \to \lvert x \rvert = -x) \to \varphi_{183}$\lbl{A1}}
\BIC{$\varphi_{183}$}
\end{ptree}
\noindent$\Pi_{65}$:
\begin{ptree}
\AXC{$\stackrel{\Pi_{63}}{\vdots}$}\noLine
\UIC{$x \in \mathbb{R} \to y \in \mathbb{R} \to \neg 0 \leq x \to 0 \leq y \to x < 0$}
\AXC{$\stackrel{\Pi_{64}}{\vdots}$}\noLine
\UIC{$\varphi_{183}$}
\BIC{$x \in \mathbb{R} \to y \in \mathbb{R} \to \neg 0 \leq x \to 0 \leq y \to \lvert x \rvert = -x$}
\AXC{$\varphi_{313}$\lbl{A1}}
\BIC{$\lvert x \rvert \cdot \lvert y \rvert = \lvert x \rvert \cdot \lvert y \rvert \to \varphi_{136} \to \varphi_{105}$}
\end{ptree}
\noindent$\Pi_{66}$:
\begin{ptree}
\AXC{$x = x$\lbl{A5}}
\UIC{$\forall x\,(x = x)$}
\AXC{$\forall x\,(x = x) \to \lvert x \rvert \cdot \lvert y \rvert = \lvert x \rvert \cdot \lvert y \rvert$\lbl{A2}}
\BIC{$\lvert x \rvert \cdot \lvert y \rvert = \lvert x \rvert \cdot \lvert y \rvert$}
\AXC{$\stackrel{\Pi_{65}}{\vdots}$}\noLine
\UIC{$\lvert x \rvert \cdot \lvert y \rvert = \lvert x \rvert \cdot \lvert y \rvert \to \varphi_{136} \to \varphi_{105}$}
\BIC{$\varphi_{136} \to \varphi_{105}$}
\end{ptree}
\noindent$\Pi_{67}$:
\begin{ptree}
\AXC{$\varphi_{66}$\lbl{A6}}
\UIC{$\forall w\,(\varphi_{66})$}
\UIC{$\forall z\,\forall w\,(\varphi_{66})$}
\AXC{$\forall z\,\forall w\,(\varphi_{66}) \to \varphi_{128}$\lbl{A2}}
\BIC{$\varphi_{128}$}
\AXC{$\varphi_{128} \to \varphi_{136}$\lbl{A2}}
\BIC{$\varphi_{136}$}
\AXC{$\stackrel{\Pi_{66}}{\vdots}$}\noLine
\UIC{$\varphi_{136} \to \varphi_{105}$}
\BIC{$\varphi_{105}$}
\AXC{$\varphi_{105} \to \varphi_{312}$\lbl{A1}}
\BIC{$\varphi_{312}$}
\end{ptree}
\noindent$\Pi_{68}$:
\begin{ptree}
\AXC{$\stackrel{\Pi_{57}}{\vdots}$}\noLine
\UIC{$\varphi_{170}$}
\AXC{$\stackrel{\Pi_{61}}{\vdots}$}\noLine
\UIC{$x \in \mathbb{R} \to y \in \mathbb{R} \to \neg 0 \leq x \to 0 \leq y \to -x \cdot \lvert y \rvert = -x \cdot y$}
\AXC{$\stackrel{\Pi_{67}}{\vdots}$}\noLine
\UIC{$\varphi_{312}$}
\BIC{$\varphi_{170} \to \varphi_{84}$}
\BIC{$\varphi_{84}$}
\AXC{$\varphi_{293}$\lbl{A1}}
\BIC{$(\lvert x \rvert \cdot \lvert y \rvert = -x \cdot y \to -x \cdot y = \lvert x \rvert \cdot \lvert y \rvert) \to \varphi_{85}$}
\end{ptree}
\noindent$\Pi_{69}$:
\begin{ptree}
\AXC{$\stackrel{\Pi_{55}}{\vdots}$}\noLine
\UIC{$\lvert x \rvert \cdot \lvert y \rvert = -x \cdot y \to -x \cdot y = \lvert x \rvert \cdot \lvert y \rvert$}
\AXC{$\stackrel{\Pi_{68}}{\vdots}$}\noLine
\UIC{$(\lvert x \rvert \cdot \lvert y \rvert = -x \cdot y \to -x \cdot y = \lvert x \rvert \cdot \lvert y \rvert) \to \varphi_{85}$}
\BIC{$\varphi_{85}$}
\end{ptree}
\noindent$\Pi_{70}$:
\begin{ptree}
\AXC{$y = z \to x = y \to x = z$\lbl{A6}}
\AXC{$(y = z \to x = y \to x = z) \to x = y \to y = z \to x = z$\lbl{A1}}
\BIC{$x = y \to y = z \to x = z$}
\UIC{$\forall z\,(x = y \to y = z \to x = z)$}
\UIC{$\forall y\,\forall z\,(x = y \to y = z \to x = z)$}
\UIC{$\forall x\,\forall y\,\forall z\,(x = y \to y = z \to x = z)$}
\AXC{$\varphi_{24}$\lbl{A2}}
\BIC{$\forall y\,\forall z\,(v_{334} = y \to y = z \to v_{334} = z)$}
\AXC{$\varphi_{46}$\lbl{A2}}
\BIC{$\forall z\,(v_{334} = v_{335} \to v_{335} = z \to v_{334} = z)$}
\AXC{$\varphi_{74}$\lbl{A2}}
\BIC{$v_{334} = v_{335} \to v_{335} = v_{336} \to v_{334} = v_{336}$}
\UIC{$\forall v_{336}\,(v_{334} = v_{335} \to v_{335} = v_{336} \to v_{334} = v_{336})$}
\UIC{$\varphi_{13}$}
\UIC{$\forall v_{334}\,(\varphi_{13})$}
\end{ptree}
\noindent$\Pi_{71}$:
\begin{ptree}
\AXC{$\stackrel{\Pi_{70}}{\vdots}$}\noLine
\UIC{$\forall v_{334}\,(\varphi_{13})$}
\AXC{$\forall v_{334}\,(\varphi_{13}) \to \varphi_{37}$\lbl{A2}}
\BIC{$\varphi_{37}$}
\AXC{$\varphi_{37} \to \varphi_{86}$\lbl{A2}}
\BIC{$\varphi_{86}$}
\AXC{$\varphi_{86} \to \varphi_{116}$\lbl{A2}}
\BIC{$\varphi_{116}$}
\end{ptree}
\noindent$\Pi_{72}$:
\begin{ptree}
\AXC{$\stackrel{\Pi_{2}}{\vdots}$}\noLine
\UIC{$\forall x\,\forall y\,(x = y \to y = x)$}
\AXC{$\forall x\,\forall y\,(x = y \to y = x) \to \forall y\,(v_{332} = y \to y = v_{332})$\lbl{A2}}
\BIC{$\forall y\,(v_{332} = y \to y = v_{332})$}
\AXC{$\varphi_{4}$\lbl{A2}}
\BIC{$v_{332} = v_{333} \to v_{333} = v_{332}$}
\UIC{$\forall v_{333}\,(v_{332} = v_{333} \to v_{333} = v_{332})$}
\UIC{$\forall v_{332}\,\forall v_{333}\,(v_{332} = v_{333} \to v_{333} = v_{332})$}
\AXC{$\varphi_{123}$\lbl{A2}}
\BIC{$\forall v_{333}\,(-x \cdot y = v_{333} \to v_{333} = -x \cdot y)$}
\AXC{$\varphi_{162}$\lbl{A2}}
\BIC{$-x \cdot y = -(x \cdot y) \to -(x \cdot y) = -x \cdot y$}
\end{ptree}
\noindent$\Pi_{73}$:
\begin{ptree}
\AXC{$\forall x\,\forall y\,(x \in \mathbb{R} \to y \in \mathbb{R} \to -x \cdot y = -(x \cdot y))$\lbl{[F9]}}
\AXC{$\varphi_{219}$\lbl{A2}}
\BIC{$\forall y\,(x \in \mathbb{R} \to y \in \mathbb{R} \to -x \cdot y = -(x \cdot y))$}
\AXC{$\varphi_{198}$\lbl{A2}}
\BIC{$x \in \mathbb{R} \to y \in \mathbb{R} \to -x \cdot y = -(x \cdot y)$}
\AXC{$\varphi_{304}$\lbl{A1}}
\BIC{$(x \in \mathbb{R} \to y \in \mathbb{R} \to \neg 0 \leq x \to 0 \leq y \to x \in \mathbb{R}) \to \varphi_{233}$}
\end{ptree}
\noindent$\Pi_{74}$:
\begin{ptree}
\AXC{$x \in \mathbb{R} \to y \in \mathbb{R} \to \neg 0 \leq x \to 0 \leq y \to x \in \mathbb{R}$\lbl{A1}}
\AXC{$\stackrel{\Pi_{73}}{\vdots}$}\noLine
\UIC{$(x \in \mathbb{R} \to y \in \mathbb{R} \to \neg 0 \leq x \to 0 \leq y \to x \in \mathbb{R}) \to \varphi_{233}$}
\BIC{$\varphi_{233}$}
\end{ptree}
\noindent$\Pi_{75}$:
\begin{ptree}
\AXC{$\stackrel{\Pi_{72}}{\vdots}$}\noLine
\UIC{$-x \cdot y = -(x \cdot y) \to -(x \cdot y) = -x \cdot y$}
\AXC{$x \in \mathbb{R} \to y \in \mathbb{R} \to \neg 0 \leq x \to 0 \leq y \to y \in \mathbb{R}$\lbl{A1}}
\AXC{$\stackrel{\Pi_{74}}{\vdots}$}\noLine
\UIC{$\varphi_{233}$}
\BIC{$x \in \mathbb{R} \to y \in \mathbb{R} \to \neg 0 \leq x \to 0 \leq y \to -x \cdot y = -(x \cdot y)$}
\AXC{$\varphi_{283}$\lbl{A1}}
\BIC{$\varphi_{234}$}
\BIC{$x \in \mathbb{R} \to y \in \mathbb{R} \to \neg 0 \leq x \to 0 \leq y \to -(x \cdot y) = -x \cdot y$}
\end{ptree}
\noindent$\Pi_{76}$:
\begin{ptree}
\AXC{$z = w \to z = -0 \to w = -0$\lbl{A6}}
\UIC{$\forall w\,(z = w \to z = -0 \to w = -0)$}
\UIC{$\forall z\,\forall w\,(z = w \to z = -0 \to w = -0)$}
\AXC{$\forall z\,\forall w\,(z = w \to z = -0 \to w = -0) \to \varphi_{32}$\lbl{A2}}
\BIC{$\varphi_{32}$}
\AXC{$\varphi_{32} \to \varphi_{67}$\lbl{A2}}
\BIC{$\varphi_{67}$}
\end{ptree}
\noindent$\Pi_{77}$:
\begin{ptree}
\AXC{$z = w \to z < -0 \to w < -0$\lbl{A6}}
\UIC{$\forall w\,(z = w \to z < -0 \to w < -0)$}
\UIC{$\forall z\,\forall w\,(z = w \to z < -0 \to w < -0)$}
\AXC{$\forall z\,\forall w\,(z = w \to z < -0 \to w < -0) \to \varphi_{65}$\lbl{A2}}
\BIC{$\varphi_{65}$}
\AXC{$\varphi_{65} \to \varphi_{103}$\lbl{A2}}
\BIC{$\varphi_{103}$}
\end{ptree}
\noindent$\Pi_{78}$:
\begin{ptree}
\AXC{$\forall x\,\forall y\,(x \in \mathbb{R} \to y \in \mathbb{R} \to \neg y \leq x \to x < y)$\lbl{[A2]}}
\AXC{$\varphi_{156}$\lbl{A2}}
\BIC{$\forall y\,(x \in \mathbb{R} \to y \in \mathbb{R} \to \neg y \leq x \to x < y)$}
\AXC{$\varphi_{129}$\lbl{A2}}
\BIC{$x \in \mathbb{R} \to 0 \in \mathbb{R} \to \neg 0 \leq x \to x < 0$}
\AXC{$(x \in \mathbb{R} \to 0 \in \mathbb{R} \to \neg 0 \leq x \to x < 0) \to \varphi_{264}$\lbl{A1}}
\BIC{$\varphi_{264}$}
\end{ptree}
\noindent$\Pi_{79}$:
\begin{ptree}
\AXC{$x \in \mathbb{R} \to y \in \mathbb{R} \to \neg 0 \leq x \to \neg 0 \leq x$\lbl{A1}}
\AXC{$0 \in \mathbb{R}$\lbl{(R1)}}
\AXC{$x \in \mathbb{R} \to y \in \mathbb{R} \to \neg 0 \leq x \to 0 \leq y \to x \in \mathbb{R}$\lbl{A1}}
\AXC{$\stackrel{\Pi_{78}}{\vdots}$}\noLine
\UIC{$\varphi_{264}$}
\BIC{$0 \in \mathbb{R} \to \varphi_{138}$}
\BIC{$\varphi_{138}$}
\BIC{$x \in \mathbb{R} \to y \in \mathbb{R} \to \neg 0 \leq x \to 0 \leq y \to x < 0$}
\end{ptree}
\noindent$\Pi_{80}$:
\begin{ptree}
\AXC{$\forall x\,\forall y\,(x \in \mathbb{R} \to y \in \mathbb{R} \to x < y \to -y < -x)$\lbl{[F22]}}
\AXC{$\varphi_{232}$\lbl{A2}}
\BIC{$\forall y\,(x \in \mathbb{R} \to y \in \mathbb{R} \to x < y \to -y < -x)$}
\AXC{$\varphi_{211}$\lbl{A2}}
\BIC{$x \in \mathbb{R} \to 0 \in \mathbb{R} \to x < 0 \to -0 < -x$}
\AXC{$(x \in \mathbb{R} \to 0 \in \mathbb{R} \to x < 0 \to -0 < -x) \to \varphi_{274}$\lbl{A1}}
\BIC{$\varphi_{274}$}
\end{ptree}
\noindent$\Pi_{81}$:
\begin{ptree}
\AXC{$\stackrel{\Pi_{79}}{\vdots}$}\noLine
\UIC{$x \in \mathbb{R} \to y \in \mathbb{R} \to \neg 0 \leq x \to 0 \leq y \to x < 0$}
\AXC{$0 \in \mathbb{R}$\lbl{(R1)}}
\AXC{$x \in \mathbb{R} \to y \in \mathbb{R} \to \neg 0 \leq x \to 0 \leq y \to x \in \mathbb{R}$\lbl{A1}}
\AXC{$\stackrel{\Pi_{80}}{\vdots}$}\noLine
\UIC{$\varphi_{274}$}
\BIC{$0 \in \mathbb{R} \to \varphi_{194}$}
\BIC{$\varphi_{194}$}
\BIC{$x \in \mathbb{R} \to y \in \mathbb{R} \to \neg 0 \leq x \to 0 \leq y \to -0 < -x$}
\end{ptree}
\noindent$\Pi_{82}$:
\begin{ptree}
\AXC{$\stackrel{\Pi_{81}}{\vdots}$}\noLine
\UIC{$x \in \mathbb{R} \to y \in \mathbb{R} \to \neg 0 \leq x \to 0 \leq y \to -0 < -x$}
\AXC{$-0 = 0$\lbl{[F6]}}
\AXC{$\varphi_{268}$\lbl{A1}}
\BIC{$(x \in \mathbb{R} \to y \in \mathbb{R} \to \neg 0 \leq x \to 0 \leq y \to -0 < -x) \to \varphi_{164}$}
\BIC{$\varphi_{164}$}
\end{ptree}
\noindent$\Pi_{83}$:
\begin{ptree}
\AXC{$z = w \to z < -x \to w < -x$\lbl{A6}}
\UIC{$\forall w\,(z = w \to z < -x \to w < -x)$}
\UIC{$\forall z\,\forall w\,(z = w \to z < -x \to w < -x)$}
\AXC{$\varphi_{132}$\lbl{A2}}
\BIC{$\forall w\,(-0 = w \to -0 < -x \to w < -x)$}
\AXC{$\varphi_{153}$\lbl{A2}}
\BIC{$-0 = 0 \to -0 < -x \to 0 < -x$}
\AXC{$\stackrel{\Pi_{82}}{\vdots}$}\noLine
\UIC{$\varphi_{164}$}
\BIC{$x \in \mathbb{R} \to y \in \mathbb{R} \to \neg 0 \leq x \to 0 \leq y \to 0 < -x$}
\AXC{$\varphi_{184}$\lbl{A1}}
\BIC{$x \in \mathbb{R} \to y \in \mathbb{R} \to \neg 0 \leq x \to 0 \leq y \to 0 \leq -x$}
\end{ptree}
\noindent$\Pi_{84}$:
\begin{ptree}
\AXC{$\forall x \in \mathbb{R}\,(-x \in \mathbb{R})$\lbl{(R1)}}
\AXC{$\forall x \in \mathbb{R}\,(-x \in \mathbb{R}) \to x \in \mathbb{R} \to -x \in \mathbb{R}$\lbl{A2}}
\BIC{$x \in \mathbb{R} \to -x \in \mathbb{R}$}
\AXC{$(x \in \mathbb{R} \to -x \in \mathbb{R}) \to \varphi_{207}$\lbl{A1}}
\BIC{$\varphi_{207}$}
\end{ptree}
\noindent$\Pi_{85}$:
\begin{ptree}
\AXC{$\forall x\,\forall y\,(x \in \mathbb{R} \to y \in \mathbb{R} \to 0 \leq x \to 0 \leq y \to 0 \leq x \cdot y)$\lbl{[A16]}}
\AXC{$\varphi_{245}$\lbl{A2}}
\BIC{$\forall y\,(-x \in \mathbb{R} \to y \in \mathbb{R} \to 0 \leq -x \to 0 \leq y \to 0 \leq -x \cdot y)$}
\AXC{$\varphi_{254}$\lbl{A2}}
\BIC{$-x \in \mathbb{R} \to y \in \mathbb{R} \to 0 \leq -x \to 0 \leq y \to 0 \leq -x \cdot y$}
\AXC{$\varphi_{326}$\lbl{A1}}
\BIC{$(x \in \mathbb{R} \to y \in \mathbb{R} \to \neg 0 \leq x \to 0 \leq y \to -x \in \mathbb{R}) \to \varphi_{294}$}
\end{ptree}
\noindent$\Pi_{86}$:
\begin{ptree}
\AXC{$x \in \mathbb{R} \to y \in \mathbb{R} \to \neg 0 \leq x \to 0 \leq y \to x \in \mathbb{R}$\lbl{A1}}
\AXC{$\stackrel{\Pi_{84}}{\vdots}$}\noLine
\UIC{$\varphi_{207}$}
\BIC{$x \in \mathbb{R} \to y \in \mathbb{R} \to \neg 0 \leq x \to 0 \leq y \to -x \in \mathbb{R}$}
\AXC{$\stackrel{\Pi_{85}}{\vdots}$}\noLine
\UIC{$(x \in \mathbb{R} \to y \in \mathbb{R} \to \neg 0 \leq x \to 0 \leq y \to -x \in \mathbb{R}) \to \varphi_{294}$}
\BIC{$\varphi_{294}$}
\end{ptree}
\noindent$\Pi_{87}$:
\begin{ptree}
\AXC{$\stackrel{\Pi_{83}}{\vdots}$}\noLine
\UIC{$x \in \mathbb{R} \to y \in \mathbb{R} \to \neg 0 \leq x \to 0 \leq y \to 0 \leq -x$}
\AXC{$x \in \mathbb{R} \to y \in \mathbb{R} \to \neg 0 \leq x \to 0 \leq y \to y \in \mathbb{R}$\lbl{A1}}
\AXC{$\stackrel{\Pi_{86}}{\vdots}$}\noLine
\UIC{$\varphi_{294}$}
\BIC{$(x \in \mathbb{R} \to y \in \mathbb{R} \to \neg 0 \leq x \to 0 \leq y \to 0 \leq -x) \to \varphi_{177}$}
\BIC{$\varphi_{177}$}
\end{ptree}
\noindent$\Pi_{88}$:
\begin{ptree}
\AXC{$\forall x\,\forall y\,(x \in \mathbb{R} \to y \in \mathbb{R} \to -x \cdot y = -(x \cdot y))$\lbl{[F9]}}
\AXC{$\varphi_{219}$\lbl{A2}}
\BIC{$\forall y\,(x \in \mathbb{R} \to y \in \mathbb{R} \to -x \cdot y = -(x \cdot y))$}
\AXC{$\varphi_{198}$\lbl{A2}}
\BIC{$x \in \mathbb{R} \to y \in \mathbb{R} \to -x \cdot y = -(x \cdot y)$}
\AXC{$\varphi_{304}$\lbl{A1}}
\BIC{$(x \in \mathbb{R} \to y \in \mathbb{R} \to \neg 0 \leq x \to 0 \leq y \to x \in \mathbb{R}) \to \varphi_{233}$}
\end{ptree}
\noindent$\Pi_{89}$:
\begin{ptree}
\AXC{$x \in \mathbb{R} \to y \in \mathbb{R} \to \neg 0 \leq x \to 0 \leq y \to x \in \mathbb{R}$\lbl{A1}}
\AXC{$\stackrel{\Pi_{88}}{\vdots}$}\noLine
\UIC{$(x \in \mathbb{R} \to y \in \mathbb{R} \to \neg 0 \leq x \to 0 \leq y \to x \in \mathbb{R}) \to \varphi_{233}$}
\BIC{$\varphi_{233}$}
\end{ptree}
\noindent$\Pi_{90}$:
\begin{ptree}
\AXC{$x \in \mathbb{R} \to y \in \mathbb{R} \to \neg 0 \leq x \to 0 \leq y \to y \in \mathbb{R}$\lbl{A1}}
\AXC{$\stackrel{\Pi_{89}}{\vdots}$}\noLine
\UIC{$\varphi_{233}$}
\BIC{$x \in \mathbb{R} \to y \in \mathbb{R} \to \neg 0 \leq x \to 0 \leq y \to -x \cdot y = -(x \cdot y)$}
\AXC{$\varphi_{322}$\lbl{A1}}
\BIC{$(x \in \mathbb{R} \to y \in \mathbb{R} \to \neg 0 \leq x \to 0 \leq y \to 0 \leq -x \cdot y) \to \varphi_{280}$}
\end{ptree}
\noindent$\Pi_{91}$:
\begin{ptree}
\AXC{$x \in \mathbb{R} \to y \in \mathbb{R} \to \neg 0 \leq x \to 0 \leq y \to 0 \leq y$\lbl{A1}}
\AXC{$\stackrel{\Pi_{87}}{\vdots}$}\noLine
\UIC{$\varphi_{177}$}
\BIC{$x \in \mathbb{R} \to y \in \mathbb{R} \to \neg 0 \leq x \to 0 \leq y \to 0 \leq -x \cdot y$}
\AXC{$\stackrel{\Pi_{90}}{\vdots}$}\noLine
\UIC{$(x \in \mathbb{R} \to y \in \mathbb{R} \to \neg 0 \leq x \to 0 \leq y \to 0 \leq -x \cdot y) \to \varphi_{280}$}
\BIC{$\varphi_{280}$}
\end{ptree}
\noindent$\Pi_{92}$:
\begin{ptree}
\AXC{$z = w \to 0 < z \to 0 < w$\lbl{A6}}
\UIC{$\forall w\,(z = w \to 0 < z \to 0 < w)$}
\UIC{$\forall z\,\forall w\,(z = w \to 0 < z \to 0 < w)$}
\AXC{$\varphi_{112}$\lbl{A2}}
\BIC{$\forall w\,(-x \cdot y = w \to 0 < -x \cdot y \to 0 < w)$}
\AXC{$\varphi_{214}$\lbl{A2}}
\BIC{$-x \cdot y = -(x \cdot y) \to 0 < -x \cdot y \to 0 < -(x \cdot y)$}
\AXC{$\stackrel{\Pi_{91}}{\vdots}$}\noLine
\UIC{$\varphi_{280}$}
\BIC{$\varphi_{220}$}
\end{ptree}
\noindent$\Pi_{93}$:
\begin{ptree}
\AXC{$\forall x\,\forall y\,(x \in \mathbb{R} \to y \in \mathbb{R} \to x \cdot y \in \mathbb{R})$\lbl{(R1)}}
\AXC{$\varphi_{165}$\lbl{A2}}
\BIC{$\forall y\,(x \in \mathbb{R} \to y \in \mathbb{R} \to x \cdot y \in \mathbb{R})$}
\AXC{$\varphi_{155}$\lbl{A2}}
\BIC{$x \in \mathbb{R} \to y \in \mathbb{R} \to x \cdot y \in \mathbb{R}$}
\AXC{$\varphi_{296}$\lbl{A1}}
\BIC{$(x \in \mathbb{R} \to y \in \mathbb{R} \to \neg 0 \leq x \to 0 \leq y \to x \in \mathbb{R}) \to \varphi_{208}$}
\end{ptree}
\noindent$\Pi_{94}$:
\begin{ptree}
\AXC{$x \in \mathbb{R} \to y \in \mathbb{R} \to \neg 0 \leq x \to 0 \leq y \to x \in \mathbb{R}$\lbl{A1}}
\AXC{$\stackrel{\Pi_{93}}{\vdots}$}\noLine
\UIC{$(x \in \mathbb{R} \to y \in \mathbb{R} \to \neg 0 \leq x \to 0 \leq y \to x \in \mathbb{R}) \to \varphi_{208}$}
\BIC{$\varphi_{208}$}
\end{ptree}
\noindent$\Pi_{95}$:
\begin{ptree}
\AXC{$x \in \mathbb{R} \to y \in \mathbb{R} \to \neg 0 \leq x \to 0 \leq y \to y \in \mathbb{R}$\lbl{A1}}
\AXC{$\stackrel{\Pi_{94}}{\vdots}$}\noLine
\UIC{$\varphi_{208}$}
\BIC{$x \in \mathbb{R} \to y \in \mathbb{R} \to \neg 0 \leq x \to 0 \leq y \to x \cdot y \in \mathbb{R}$}
\end{ptree}
\noindent$\Pi_{96}$:
\begin{ptree}
\AXC{$\forall x \in \mathbb{R}\,(-x \in \mathbb{R})$\lbl{(R1)}}
\AXC{$\forall x \in \mathbb{R}\,(-x \in \mathbb{R}) \to x \cdot y \in \mathbb{R} \to -(x \cdot y) \in \mathbb{R}$\lbl{A2}}
\BIC{$x \cdot y \in \mathbb{R} \to -(x \cdot y) \in \mathbb{R}$}
\end{ptree}
\noindent$\Pi_{97}$:
\begin{ptree}
\AXC{$\stackrel{\Pi_{95}}{\vdots}$}\noLine
\UIC{$x \in \mathbb{R} \to y \in \mathbb{R} \to \neg 0 \leq x \to 0 \leq y \to x \cdot y \in \mathbb{R}$}
\AXC{$\stackrel{\Pi_{96}}{\vdots}$}\noLine
\UIC{$x \cdot y \in \mathbb{R} \to -(x \cdot y) \in \mathbb{R}$}
\AXC{$(x \cdot y \in \mathbb{R} \to -(x \cdot y) \in \mathbb{R}) \to \varphi_{240}$\lbl{A1}}
\BIC{$\varphi_{240}$}
\BIC{$x \in \mathbb{R} \to y \in \mathbb{R} \to \neg 0 \leq x \to 0 \leq y \to -(x \cdot y) \in \mathbb{R}$}
\end{ptree}
\noindent$\Pi_{98}$:
\begin{ptree}
\AXC{$\forall x\,\forall y\,(x \in \mathbb{R} \to y \in \mathbb{R} \to x \leq y \to -y \leq -x)$\lbl{[F23]}}
\AXC{$\varphi_{191}$\lbl{A2}}
\BIC{$\forall y\,(0 \in \mathbb{R} \to y \in \mathbb{R} \to 0 \leq y \to -y \leq -0)$}
\AXC{$\forall y\,(0 \in \mathbb{R} \to y \in \mathbb{R} \to 0 \leq y \to -y \leq -0) \to 0 \in \mathbb{R} \to \varphi_{60}$\lbl{A2}}
\BIC{$0 \in \mathbb{R} \to \varphi_{60}$}
\end{ptree}
\noindent$\Pi_{99}$:
\begin{ptree}
\AXC{$\stackrel{\Pi_{97}}{\vdots}$}\noLine
\UIC{$x \in \mathbb{R} \to y \in \mathbb{R} \to \neg 0 \leq x \to 0 \leq y \to -(x \cdot y) \in \mathbb{R}$}
\AXC{$0 \in \mathbb{R}$\lbl{(R1)}}
\AXC{$\stackrel{\Pi_{98}}{\vdots}$}\noLine
\UIC{$0 \in \mathbb{R} \to \varphi_{60}$}
\AXC{$(0 \in \mathbb{R} \to \varphi_{60}) \to 0 \in \mathbb{R} \to \varphi_{284}$\lbl{A1}}
\BIC{$0 \in \mathbb{R} \to \varphi_{284}$}
\BIC{$\varphi_{284}$}
\BIC{$\varphi_{242}$}
\end{ptree}
\noindent$\Pi_{100}$:
\begin{ptree}
\AXC{$\stackrel{\Pi_{95}}{\vdots}$}\noLine
\UIC{$x \in \mathbb{R} \to y \in \mathbb{R} \to \neg 0 \leq x \to 0 \leq y \to x \cdot y \in \mathbb{R}$}
\AXC{$\forall x \in \mathbb{R}\,(-(-x) = x)$\lbl{[F5]}}
\AXC{$\varphi_{111}$\lbl{A2}}
\BIC{$x \cdot y \in \mathbb{R} \to -(-(x \cdot y)) = x \cdot y$}
\AXC{$\varphi_{269}$\lbl{A1}}
\BIC{$(x \in \mathbb{R} \to y \in \mathbb{R} \to \neg 0 \leq x \to 0 \leq y \to x \cdot y \in \mathbb{R}) \to \varphi_{54}$}
\BIC{$\varphi_{54}$}
\AXC{$\varphi_{54} \to \varphi_{316}$\lbl{A1}}
\BIC{$\varphi_{316}$}
\end{ptree}
\noindent$\Pi_{101}$:
\begin{ptree}
\AXC{$z = w \to 0 = z \to 0 = w$\lbl{A6}}
\UIC{$\forall w\,(z = w \to 0 = z \to 0 = w)$}
\UIC{$\forall z\,\forall w\,(z = w \to 0 = z \to 0 = w)$}
\AXC{$\varphi_{52}$\lbl{A2}}
\BIC{$\forall w\,(-x \cdot y = w \to 0 = -x \cdot y \to 0 = w)$}
\AXC{$\varphi_{173}$\lbl{A2}}
\BIC{$-x \cdot y = -(x \cdot y) \to 0 = -x \cdot y \to 0 = -(x \cdot y)$}
\AXC{$\stackrel{\Pi_{92}}{\vdots}$}\noLine
\UIC{$\varphi_{220}$}
\BIC{$x \in \mathbb{R} \to y \in \mathbb{R} \to \neg 0 \leq x \to 0 \leq y \to 0 \leq -(x \cdot y)$}
\AXC{$\stackrel{\Pi_{99}}{\vdots}$}\noLine
\UIC{$\varphi_{242}$}
\BIC{$\varphi_{53}$}
\AXC{$\stackrel{\Pi_{100}}{\vdots}$}\noLine
\UIC{$\varphi_{316}$}
\BIC{$\varphi_{103} \to \varphi_{67} \to x \in \mathbb{R} \to y \in \mathbb{R} \to \neg 0 \leq x \to 0 \leq y \to x \cdot y \leq -0$}
\end{ptree}
\noindent$\Pi_{102}$:
\begin{ptree}
\AXC{$\stackrel{\Pi_{76}}{\vdots}$}\noLine
\UIC{$\varphi_{67}$}
\AXC{$\stackrel{\Pi_{77}}{\vdots}$}\noLine
\UIC{$\varphi_{103}$}
\AXC{$\stackrel{\Pi_{101}}{\vdots}$}\noLine
\UIC{$\varphi_{103} \to \varphi_{67} \to x \in \mathbb{R} \to y \in \mathbb{R} \to \neg 0 \leq x \to 0 \leq y \to x \cdot y \leq -0$}
\BIC{$\varphi_{67} \to x \in \mathbb{R} \to y \in \mathbb{R} \to \neg 0 \leq x \to 0 \leq y \to x \cdot y \leq -0$}
\BIC{$x \in \mathbb{R} \to y \in \mathbb{R} \to \neg 0 \leq x \to 0 \leq y \to x \cdot y \leq -0$}
\AXC{$-0 = 0$\lbl{[F6]}}
\AXC{$-0 = 0 \to \varphi_{288}$\lbl{A1}}
\BIC{$\varphi_{288}$}
\BIC{$(-0 = 0 \to x \cdot y < -0 \to x \cdot y < 0) \to \varphi_{139}$}
\end{ptree}
\noindent$\Pi_{103}$:
\begin{ptree}
\AXC{$z = w \to x \cdot y < z \to x \cdot y < w$\lbl{A6}}
\UIC{$\forall w\,(z = w \to x \cdot y < z \to x \cdot y < w)$}
\UIC{$\forall z\,\forall w\,(z = w \to x \cdot y < z \to x \cdot y < w)$}
\AXC{$\varphi_{133}$\lbl{A2}}
\BIC{$\forall w\,(-0 = w \to x \cdot y < -0 \to x \cdot y < w)$}
\AXC{$\varphi_{154}$\lbl{A2}}
\BIC{$-0 = 0 \to x \cdot y < -0 \to x \cdot y < 0$}
\AXC{$\stackrel{\Pi_{102}}{\vdots}$}\noLine
\UIC{$(-0 = 0 \to x \cdot y < -0 \to x \cdot y < 0) \to \varphi_{139}$}
\BIC{$\varphi_{139}$}
\end{ptree}
\noindent$\Pi_{104}$:
\begin{ptree}
\AXC{$z = w \to x \cdot y = z \to x \cdot y = w$\lbl{A6}}
\UIC{$\forall w\,(z = w \to x \cdot y = z \to x \cdot y = w)$}
\UIC{$\forall z\,\forall w\,(z = w \to x \cdot y = z \to x \cdot y = w)$}
\AXC{$\varphi_{106}$\lbl{A2}}
\BIC{$\forall w\,(-0 = w \to x \cdot y = -0 \to x \cdot y = w)$}
\AXC{$\varphi_{121}$\lbl{A2}}
\BIC{$-0 = 0 \to x \cdot y = -0 \to x \cdot y = 0$}
\AXC{$\stackrel{\Pi_{103}}{\vdots}$}\noLine
\UIC{$\varphi_{139}$}
\BIC{$x \in \mathbb{R} \to y \in \mathbb{R} \to \neg 0 \leq x \to 0 \leq y \to x \cdot y \leq 0$}
\end{ptree}
\noindent$\Pi_{105}$:
\begin{ptree}
\AXC{$\forall x \in \mathbb{R}\,(x \leq 0 \to \lvert x \rvert = -x)$\lbl{[A4]}}
\AXC{$\varphi_{157}$\lbl{A2}}
\BIC{$x \cdot y \in \mathbb{R} \to x \cdot y \leq 0 \to \lvert x \cdot y \rvert = -(x \cdot y)$}
\AXC{$(x \cdot y \in \mathbb{R} \to x \cdot y \leq 0 \to \lvert x \cdot y \rvert = -(x \cdot y)) \to \varphi_{276}$\lbl{A1}}
\BIC{$\varphi_{276}$}
\end{ptree}
\noindent$\Pi_{106}$:
\begin{ptree}
\AXC{$\stackrel{\Pi_{104}}{\vdots}$}\noLine
\UIC{$x \in \mathbb{R} \to y \in \mathbb{R} \to \neg 0 \leq x \to 0 \leq y \to x \cdot y \leq 0$}
\AXC{$\stackrel{\Pi_{95}}{\vdots}$}\noLine
\UIC{$x \in \mathbb{R} \to y \in \mathbb{R} \to \neg 0 \leq x \to 0 \leq y \to x \cdot y \in \mathbb{R}$}
\AXC{$\stackrel{\Pi_{105}}{\vdots}$}\noLine
\UIC{$\varphi_{276}$}
\BIC{$(x \in \mathbb{R} \to y \in \mathbb{R} \to \neg 0 \leq x \to 0 \leq y \to x \cdot y \leq 0) \to \varphi_{55}$}
\BIC{$\varphi_{55}$}
\AXC{$\varphi_{55} \to \varphi_{298}$\lbl{A1}}
\BIC{$\varphi_{298}$}
\end{ptree}
\noindent$\Pi_{107}$:
\begin{ptree}
\AXC{$\stackrel{\Pi_{69}}{\vdots}$}\noLine
\UIC{$\varphi_{85}$}
\AXC{$\stackrel{\Pi_{71}}{\vdots}$}\noLine
\UIC{$\varphi_{116}$}
\AXC{$\stackrel{\Pi_{75}}{\vdots}$}\noLine
\UIC{$x \in \mathbb{R} \to y \in \mathbb{R} \to \neg 0 \leq x \to 0 \leq y \to -(x \cdot y) = -x \cdot y$}
\AXC{$\stackrel{\Pi_{106}}{\vdots}$}\noLine
\UIC{$\varphi_{298}$}
\BIC{$\varphi_{259}$}
\BIC{$x \in \mathbb{R} \to y \in \mathbb{R} \to \neg 0 \leq x \to 0 \leq y \to \lvert x \cdot y \rvert = -x \cdot y$}
\AXC{$\varphi_{320}$\lbl{A1}}
\BIC{$\varphi_{85} \to \varphi_{137} \to \varphi_{87}$}
\BIC{$\varphi_{137} \to \varphi_{87}$}
\end{ptree}
\noindent$\Pi_{108}$:
\begin{ptree}
\AXC{$\stackrel{\Pi_{54}}{\vdots}$}\noLine
\UIC{$\forall v_{339}\,(\varphi_{11})$}
\AXC{$\forall v_{339}\,(\varphi_{11}) \to \varphi_{36}$\lbl{A2}}
\BIC{$\varphi_{36}$}
\AXC{$\varphi_{36} \to \varphi_{83}$\lbl{A2}}
\BIC{$\varphi_{83}$}
\AXC{$\varphi_{83} \to \varphi_{137}$\lbl{A2}}
\BIC{$\varphi_{137}$}
\AXC{$\stackrel{\Pi_{107}}{\vdots}$}\noLine
\UIC{$\varphi_{137} \to \varphi_{87}$}
\BIC{$\varphi_{87}$}
\end{ptree}
\noindent$\Pi_{109}$:
\begin{ptree}
\AXC{$y = z \to x = y \to x = z$\lbl{A6}}
\AXC{$(y = z \to x = y \to x = z) \to x = y \to y = z \to x = z$\lbl{A1}}
\BIC{$x = y \to y = z \to x = z$}
\UIC{$\forall z\,(x = y \to y = z \to x = z)$}
\UIC{$\forall y\,\forall z\,(x = y \to y = z \to x = z)$}
\UIC{$\forall x\,\forall y\,\forall z\,(x = y \to y = z \to x = z)$}
\AXC{$\varphi_{25}$\lbl{A2}}
\BIC{$\forall y\,\forall z\,(v_{326} = y \to y = z \to v_{326} = z)$}
\AXC{$\varphi_{47}$\lbl{A2}}
\BIC{$\forall z\,(v_{326} = v_{327} \to v_{327} = z \to v_{326} = z)$}
\AXC{$\varphi_{75}$\lbl{A2}}
\BIC{$v_{326} = v_{327} \to v_{327} = v_{328} \to v_{326} = v_{328}$}
\UIC{$\forall v_{328}\,(v_{326} = v_{327} \to v_{327} = v_{328} \to v_{326} = v_{328})$}
\UIC{$\varphi_{14}$}
\UIC{$\forall v_{326}\,(\varphi_{14})$}
\end{ptree}
\noindent$\Pi_{110}$:
\begin{ptree}
\AXC{$\stackrel{\Pi_{2}}{\vdots}$}\noLine
\UIC{$\forall x\,\forall y\,(x = y \to y = x)$}
\AXC{$\forall x\,\forall y\,(x = y \to y = x) \to \forall y\,(v_{324} = y \to y = v_{324})$\lbl{A2}}
\BIC{$\forall y\,(v_{324} = y \to y = v_{324})$}
\AXC{$\varphi_{5}$\lbl{A2}}
\BIC{$v_{324} = v_{325} \to v_{325} = v_{324}$}
\UIC{$\forall v_{325}\,(v_{324} = v_{325} \to v_{325} = v_{324})$}
\UIC{$\forall v_{324}\,\forall v_{325}\,(v_{324} = v_{325} \to v_{325} = v_{324})$}
\AXC{$\varphi_{150}$\lbl{A2}}
\BIC{$\forall v_{325}\,(\lvert x \rvert \cdot \lvert y \rvert = v_{325} \to v_{325} = \lvert x \rvert \cdot \lvert y \rvert)$}
\AXC{$\varphi_{229}$\lbl{A2}}
\BIC{$\lvert x \rvert \cdot \lvert y \rvert = x \cdot -y \to x \cdot -y = \lvert x \rvert \cdot \lvert y \rvert$}
\end{ptree}
\noindent$\Pi_{111}$:
\begin{ptree}
\AXC{$y = z \to x = y \to x = z$\lbl{A6}}
\AXC{$(y = z \to x = y \to x = z) \to x = y \to y = z \to x = z$\lbl{A1}}
\BIC{$x = y \to y = z \to x = z$}
\UIC{$\forall z\,(x = y \to y = z \to x = z)$}
\UIC{$\forall y\,\forall z\,(x = y \to y = z \to x = z)$}
\UIC{$\forall x\,\forall y\,\forall z\,(x = y \to y = z \to x = z)$}
\AXC{$\varphi_{26}$\lbl{A2}}
\BIC{$\forall y\,\forall z\,(v_{312} = y \to y = z \to v_{312} = z)$}
\AXC{$\varphi_{48}$\lbl{A2}}
\BIC{$\forall z\,(v_{312} = v_{313} \to v_{313} = z \to v_{312} = z)$}
\AXC{$\varphi_{76}$\lbl{A2}}
\BIC{$v_{312} = v_{313} \to v_{313} = v_{314} \to v_{312} = v_{314}$}
\UIC{$\forall v_{314}\,(v_{312} = v_{313} \to v_{313} = v_{314} \to v_{312} = v_{314})$}
\UIC{$\varphi_{15}$}
\UIC{$\forall v_{312}\,(\varphi_{15})$}
\end{ptree}
\noindent$\Pi_{112}$:
\begin{ptree}
\AXC{$\forall x\,\forall y\,(x \in \mathbb{R} \to y \in \mathbb{R} \to \neg y \leq x \to x < y)$\lbl{[A2]}}
\AXC{$\varphi_{181}$\lbl{A2}}
\BIC{$\forall y\,(v_{310} \in \mathbb{R} \to y \in \mathbb{R} \to \neg y \leq v_{310} \to v_{310} < y)$}
\AXC{$\varphi_{225}$\lbl{A2}}
\BIC{$v_{310} \in \mathbb{R} \to v_{311} \in \mathbb{R} \to \neg v_{311} \leq v_{310} \to v_{310} < v_{311}$}
\UIC{$\varphi_{63}$}
\UIC{$\forall v_{310}\,(\varphi_{63})$}
\AXC{$\varphi_{252}$\lbl{A2}}
\BIC{$\forall v_{311}\,(y \in \mathbb{R} \to v_{311} \in \mathbb{R} \to \neg v_{311} \leq y \to y < v_{311})$}
\AXC{$\varphi_{174}$\lbl{A2}}
\BIC{$y \in \mathbb{R} \to 0 \in \mathbb{R} \to \neg 0 \leq y \to y < 0$}
\end{ptree}
\noindent$\Pi_{113}$:
\begin{ptree}
\AXC{$\stackrel{\Pi_{112}}{\vdots}$}\noLine
\UIC{$y \in \mathbb{R} \to 0 \in \mathbb{R} \to \neg 0 \leq y \to y < 0$}
\AXC{$(y \in \mathbb{R} \to 0 \in \mathbb{R} \to \neg 0 \leq y \to y < 0) \to \varphi_{266}$\lbl{A1}}
\BIC{$\varphi_{266}$}
\end{ptree}
\noindent$\Pi_{114}$:
\begin{ptree}
\AXC{$x \in \mathbb{R} \to y \in \mathbb{R} \to 0 \leq x \to \neg 0 \leq y \to \neg 0 \leq y$\lbl{A1}}
\AXC{$0 \in \mathbb{R}$\lbl{(R1)}}
\AXC{$x \in \mathbb{R} \to y \in \mathbb{R} \to 0 \leq x \to \neg 0 \leq y \to y \in \mathbb{R}$\lbl{A1}}
\AXC{$\stackrel{\Pi_{113}}{\vdots}$}\noLine
\UIC{$\varphi_{266}$}
\BIC{$0 \in \mathbb{R} \to \varphi_{159}$}
\BIC{$\varphi_{159}$}
\BIC{$x \in \mathbb{R} \to y \in \mathbb{R} \to 0 \leq x \to \neg 0 \leq y \to y < 0$}
\end{ptree}
\noindent$\Pi_{115}$:
\begin{ptree}
\AXC{$\forall x\,(x < 0 \to \lvert x \rvert = -x)$\lbl{[A3]}}
\AXC{$\forall x\,(x < 0 \to \lvert x \rvert = -x) \to y < 0 \to \lvert y \rvert = -y$\lbl{A2}}
\BIC{$y < 0 \to \lvert y \rvert = -y$}
\AXC{$(y < 0 \to \lvert y \rvert = -y) \to \varphi_{185}$\lbl{A1}}
\BIC{$\varphi_{185}$}
\end{ptree}
\noindent$\Pi_{116}$:
\begin{ptree}
\AXC{$x = x$\lbl{A5}}
\UIC{$\forall x\,(x = x)$}
\AXC{$\forall x\,(x = x) \to x \cdot \lvert y \rvert = x \cdot \lvert y \rvert$\lbl{A2}}
\BIC{$x \cdot \lvert y \rvert = x \cdot \lvert y \rvert$}
\AXC{$\stackrel{\Pi_{114}}{\vdots}$}\noLine
\UIC{$x \in \mathbb{R} \to y \in \mathbb{R} \to 0 \leq x \to \neg 0 \leq y \to y < 0$}
\AXC{$\stackrel{\Pi_{115}}{\vdots}$}\noLine
\UIC{$\varphi_{185}$}
\BIC{$x \in \mathbb{R} \to y \in \mathbb{R} \to 0 \leq x \to \neg 0 \leq y \to \lvert y \rvert = -y$}
\AXC{$\varphi_{301}$\lbl{A1}}
\BIC{$x \cdot \lvert y \rvert = x \cdot \lvert y \rvert \to \varphi_{255}$}
\BIC{$\varphi_{255}$}
\end{ptree}
\noindent$\Pi_{117}$:
\begin{ptree}
\AXC{$\forall x\,(0 \leq x \to \lvert x \rvert = x)$\lbl{[A1]}}
\AXC{$\forall x\,(0 \leq x \to \lvert x \rvert = x) \to 0 \leq x \to \lvert x \rvert = x$\lbl{A2}}
\BIC{$0 \leq x \to \lvert x \rvert = x$}
\AXC{$(0 \leq x \to \lvert x \rvert = x) \to \varphi_{142}$\lbl{A1}}
\BIC{$\varphi_{142}$}
\end{ptree}
\noindent$\Pi_{118}$:
\begin{ptree}
\AXC{$x = x$\lbl{A5}}
\UIC{$\forall x\,(x = x)$}
\AXC{$\forall x\,(x = x) \to \lvert x \rvert \cdot \lvert y \rvert = \lvert x \rvert \cdot \lvert y \rvert$\lbl{A2}}
\BIC{$\lvert x \rvert \cdot \lvert y \rvert = \lvert x \rvert \cdot \lvert y \rvert$}
\AXC{$x \in \mathbb{R} \to y \in \mathbb{R} \to 0 \leq x \to 0 \leq x$\lbl{A1}}
\AXC{$\stackrel{\Pi_{117}}{\vdots}$}\noLine
\UIC{$\varphi_{142}$}
\BIC{$x \in \mathbb{R} \to y \in \mathbb{R} \to 0 \leq x \to \neg 0 \leq y \to \lvert x \rvert = x$}
\AXC{$\varphi_{311}$\lbl{A1}}
\BIC{$\lvert x \rvert \cdot \lvert y \rvert = \lvert x \rvert \cdot \lvert y \rvert \to \varphi_{117} \to \varphi_{89}$}
\BIC{$\varphi_{117} \to \varphi_{89}$}
\end{ptree}
\noindent$\Pi_{119}$:
\begin{ptree}
\AXC{$\varphi_{66}$\lbl{A6}}
\UIC{$\forall w\,(\varphi_{66})$}
\UIC{$\forall z\,\forall w\,(\varphi_{66})$}
\AXC{$\forall z\,\forall w\,(\varphi_{66}) \to \varphi_{128}$\lbl{A2}}
\BIC{$\varphi_{128}$}
\AXC{$\varphi_{128} \to \varphi_{117}$\lbl{A2}}
\BIC{$\varphi_{117}$}
\AXC{$\stackrel{\Pi_{118}}{\vdots}$}\noLine
\UIC{$\varphi_{117} \to \varphi_{89}$}
\BIC{$\varphi_{89}$}
\AXC{$\varphi_{89} \to \varphi_{306}$\lbl{A1}}
\BIC{$\varphi_{306}$}
\end{ptree}
\noindent$\Pi_{120}$:
\begin{ptree}
\AXC{$z = w \to x \cdot \lvert y \rvert = x \cdot z \to x \cdot \lvert y \rvert = x \cdot w$\lbl{A6}}
\UIC{$\forall w\,(z = w \to x \cdot \lvert y \rvert = x \cdot z \to x \cdot \lvert y \rvert = x \cdot w)$}
\UIC{$\forall z\,\forall w\,(z = w \to x \cdot \lvert y \rvert = x \cdot z \to x \cdot \lvert y \rvert = x \cdot w)$}
\AXC{$\varphi_{227}$\lbl{A2}}
\BIC{$\varphi_{33}$}
\AXC{$\varphi_{251}$\lbl{A2}}
\BIC{$\lvert y \rvert = -y \to x \cdot \lvert y \rvert = x \cdot \lvert y \rvert \to x \cdot \lvert y \rvert = x \cdot -y$}
\AXC{$\stackrel{\Pi_{116}}{\vdots}$}\noLine
\UIC{$\varphi_{255}$}
\BIC{$x \in \mathbb{R} \to y \in \mathbb{R} \to 0 \leq x \to \neg 0 \leq y \to x \cdot \lvert y \rvert = x \cdot -y$}
\AXC{$\stackrel{\Pi_{119}}{\vdots}$}\noLine
\UIC{$\varphi_{306}$}
\BIC{$\varphi_{141} \to \varphi_{90}$}
\end{ptree}
\noindent$\Pi_{121}$:
\begin{ptree}
\AXC{$\stackrel{\Pi_{111}}{\vdots}$}\noLine
\UIC{$\forall v_{312}\,(\varphi_{15})$}
\AXC{$\forall v_{312}\,(\varphi_{15}) \to \varphi_{100}$\lbl{A2}}
\BIC{$\varphi_{100}$}
\AXC{$\varphi_{100} \to \varphi_{124}$\lbl{A2}}
\BIC{$\varphi_{124}$}
\AXC{$\varphi_{124} \to \varphi_{141}$\lbl{A2}}
\BIC{$\varphi_{141}$}
\AXC{$\stackrel{\Pi_{120}}{\vdots}$}\noLine
\UIC{$\varphi_{141} \to \varphi_{90}$}
\BIC{$\varphi_{90}$}
\AXC{$\varphi_{295}$\lbl{A1}}
\BIC{$(\lvert x \rvert \cdot \lvert y \rvert = x \cdot -y \to x \cdot -y = \lvert x \rvert \cdot \lvert y \rvert) \to \varphi_{91}$}
\end{ptree}
\noindent$\Pi_{122}$:
\begin{ptree}
\AXC{$\stackrel{\Pi_{110}}{\vdots}$}\noLine
\UIC{$\lvert x \rvert \cdot \lvert y \rvert = x \cdot -y \to x \cdot -y = \lvert x \rvert \cdot \lvert y \rvert$}
\AXC{$\stackrel{\Pi_{121}}{\vdots}$}\noLine
\UIC{$(\lvert x \rvert \cdot \lvert y \rvert = x \cdot -y \to x \cdot -y = \lvert x \rvert \cdot \lvert y \rvert) \to \varphi_{91}$}
\BIC{$\varphi_{91}$}
\end{ptree}
\noindent$\Pi_{123}$:
\begin{ptree}
\AXC{$y = z \to x = y \to x = z$\lbl{A6}}
\AXC{$(y = z \to x = y \to x = z) \to x = y \to y = z \to x = z$\lbl{A1}}
\BIC{$x = y \to y = z \to x = z$}
\UIC{$\forall z\,(x = y \to y = z \to x = z)$}
\UIC{$\forall y\,\forall z\,(x = y \to y = z \to x = z)$}
\UIC{$\forall x\,\forall y\,\forall z\,(x = y \to y = z \to x = z)$}
\AXC{$\varphi_{27}$\lbl{A2}}
\BIC{$\forall y\,\forall z\,(v_{321} = y \to y = z \to v_{321} = z)$}
\AXC{$\varphi_{49}$\lbl{A2}}
\BIC{$\forall z\,(v_{321} = v_{322} \to v_{322} = z \to v_{321} = z)$}
\AXC{$\varphi_{77}$\lbl{A2}}
\BIC{$v_{321} = v_{322} \to v_{322} = v_{323} \to v_{321} = v_{323}$}
\UIC{$\forall v_{323}\,(v_{321} = v_{322} \to v_{322} = v_{323} \to v_{321} = v_{323})$}
\UIC{$\varphi_{16}$}
\UIC{$\forall v_{321}\,(\varphi_{16})$}
\end{ptree}
\noindent$\Pi_{124}$:
\begin{ptree}
\AXC{$\stackrel{\Pi_{123}}{\vdots}$}\noLine
\UIC{$\forall v_{321}\,(\varphi_{16})$}
\AXC{$\forall v_{321}\,(\varphi_{16}) \to \varphi_{39}$\lbl{A2}}
\BIC{$\varphi_{39}$}
\AXC{$\varphi_{39} \to \varphi_{92}$\lbl{A2}}
\BIC{$\varphi_{92}$}
\AXC{$\varphi_{92} \to \varphi_{118}$\lbl{A2}}
\BIC{$\varphi_{118}$}
\end{ptree}
\noindent$\Pi_{125}$:
\begin{ptree}
\AXC{$\stackrel{\Pi_{2}}{\vdots}$}\noLine
\UIC{$\forall x\,\forall y\,(x = y \to y = x)$}
\AXC{$\forall x\,\forall y\,(x = y \to y = x) \to \forall y\,(v_{319} = y \to y = v_{319})$\lbl{A2}}
\BIC{$\forall y\,(v_{319} = y \to y = v_{319})$}
\AXC{$\varphi_{6}$\lbl{A2}}
\BIC{$v_{319} = v_{320} \to v_{320} = v_{319}$}
\UIC{$\forall v_{320}\,(v_{319} = v_{320} \to v_{320} = v_{319})$}
\UIC{$\forall v_{319}\,\forall v_{320}\,(v_{319} = v_{320} \to v_{320} = v_{319})$}
\AXC{$\varphi_{125}$\lbl{A2}}
\BIC{$\forall v_{320}\,(x \cdot -y = v_{320} \to v_{320} = x \cdot -y)$}
\AXC{$\varphi_{163}$\lbl{A2}}
\BIC{$x \cdot -y = -(x \cdot y) \to -(x \cdot y) = x \cdot -y$}
\end{ptree}
\noindent$\Pi_{126}$:
\begin{ptree}
\AXC{$\forall x\,\forall y\,(x \in \mathbb{R} \to y \in \mathbb{R} \to x \cdot -y = -(x \cdot y))$\lbl{[F10]}}
\AXC{$\varphi_{221}$\lbl{A2}}
\BIC{$\forall y\,(x \in \mathbb{R} \to y \in \mathbb{R} \to x \cdot -y = -(x \cdot y))$}
\AXC{$\varphi_{199}$\lbl{A2}}
\BIC{$x \in \mathbb{R} \to y \in \mathbb{R} \to x \cdot -y = -(x \cdot y)$}
\AXC{$\varphi_{305}$\lbl{A1}}
\BIC{$(x \in \mathbb{R} \to y \in \mathbb{R} \to 0 \leq x \to \neg 0 \leq y \to x \in \mathbb{R}) \to \varphi_{235}$}
\end{ptree}
\noindent$\Pi_{127}$:
\begin{ptree}
\AXC{$x \in \mathbb{R} \to y \in \mathbb{R} \to 0 \leq x \to \neg 0 \leq y \to x \in \mathbb{R}$\lbl{A1}}
\AXC{$\stackrel{\Pi_{126}}{\vdots}$}\noLine
\UIC{$(x \in \mathbb{R} \to y \in \mathbb{R} \to 0 \leq x \to \neg 0 \leq y \to x \in \mathbb{R}) \to \varphi_{235}$}
\BIC{$\varphi_{235}$}
\end{ptree}
\noindent$\Pi_{128}$:
\begin{ptree}
\AXC{$\stackrel{\Pi_{125}}{\vdots}$}\noLine
\UIC{$x \cdot -y = -(x \cdot y) \to -(x \cdot y) = x \cdot -y$}
\AXC{$x \in \mathbb{R} \to y \in \mathbb{R} \to 0 \leq x \to \neg 0 \leq y \to y \in \mathbb{R}$\lbl{A1}}
\AXC{$\stackrel{\Pi_{127}}{\vdots}$}\noLine
\UIC{$\varphi_{235}$}
\BIC{$x \in \mathbb{R} \to y \in \mathbb{R} \to 0 \leq x \to \neg 0 \leq y \to x \cdot -y = -(x \cdot y)$}
\AXC{$\varphi_{285}$\lbl{A1}}
\BIC{$\varphi_{236}$}
\BIC{$x \in \mathbb{R} \to y \in \mathbb{R} \to 0 \leq x \to \neg 0 \leq y \to -(x \cdot y) = x \cdot -y$}
\end{ptree}
\noindent$\Pi_{129}$:
\begin{ptree}
\AXC{$z = w \to z = -0 \to w = -0$\lbl{A6}}
\UIC{$\forall w\,(z = w \to z = -0 \to w = -0)$}
\UIC{$\forall z\,\forall w\,(z = w \to z = -0 \to w = -0)$}
\AXC{$\forall z\,\forall w\,(z = w \to z = -0 \to w = -0) \to \varphi_{32}$\lbl{A2}}
\BIC{$\varphi_{32}$}
\AXC{$\varphi_{32} \to \varphi_{67}$\lbl{A2}}
\BIC{$\varphi_{67}$}
\end{ptree}
\noindent$\Pi_{130}$:
\begin{ptree}
\AXC{$z = w \to z < -0 \to w < -0$\lbl{A6}}
\UIC{$\forall w\,(z = w \to z < -0 \to w < -0)$}
\UIC{$\forall z\,\forall w\,(z = w \to z < -0 \to w < -0)$}
\AXC{$\forall z\,\forall w\,(z = w \to z < -0 \to w < -0) \to \varphi_{65}$\lbl{A2}}
\BIC{$\varphi_{65}$}
\AXC{$\varphi_{65} \to \varphi_{103}$\lbl{A2}}
\BIC{$\varphi_{103}$}
\end{ptree}
\noindent$\Pi_{131}$:
\begin{ptree}
\AXC{$\forall x\,\forall y\,(x \in \mathbb{R} \to y \in \mathbb{R} \to \neg y \leq x \to x < y)$\lbl{[A2]}}
\AXC{$\varphi_{182}$\lbl{A2}}
\BIC{$\forall y\,(v_{315} \in \mathbb{R} \to y \in \mathbb{R} \to \neg y \leq v_{315} \to v_{315} < y)$}
\AXC{$\varphi_{226}$\lbl{A2}}
\BIC{$v_{315} \in \mathbb{R} \to v_{316} \in \mathbb{R} \to \neg v_{316} \leq v_{315} \to v_{315} < v_{316}$}
\UIC{$\varphi_{64}$}
\UIC{$\forall v_{315}\,(\varphi_{64})$}
\AXC{$\varphi_{253}$\lbl{A2}}
\BIC{$\forall v_{316}\,(y \in \mathbb{R} \to v_{316} \in \mathbb{R} \to \neg v_{316} \leq y \to y < v_{316})$}
\AXC{$\varphi_{176}$\lbl{A2}}
\BIC{$y \in \mathbb{R} \to 0 \in \mathbb{R} \to \neg 0 \leq y \to y < 0$}
\end{ptree}
\noindent$\Pi_{132}$:
\begin{ptree}
\AXC{$\stackrel{\Pi_{131}}{\vdots}$}\noLine
\UIC{$y \in \mathbb{R} \to 0 \in \mathbb{R} \to \neg 0 \leq y \to y < 0$}
\AXC{$(y \in \mathbb{R} \to 0 \in \mathbb{R} \to \neg 0 \leq y \to y < 0) \to \varphi_{266}$\lbl{A1}}
\BIC{$\varphi_{266}$}
\end{ptree}
\noindent$\Pi_{133}$:
\begin{ptree}
\AXC{$x \in \mathbb{R} \to y \in \mathbb{R} \to 0 \leq x \to \neg 0 \leq y \to \neg 0 \leq y$\lbl{A1}}
\AXC{$0 \in \mathbb{R}$\lbl{(R1)}}
\AXC{$x \in \mathbb{R} \to y \in \mathbb{R} \to 0 \leq x \to \neg 0 \leq y \to y \in \mathbb{R}$\lbl{A1}}
\AXC{$\stackrel{\Pi_{132}}{\vdots}$}\noLine
\UIC{$\varphi_{266}$}
\BIC{$0 \in \mathbb{R} \to \varphi_{159}$}
\BIC{$\varphi_{159}$}
\BIC{$x \in \mathbb{R} \to y \in \mathbb{R} \to 0 \leq x \to \neg 0 \leq y \to y < 0$}
\end{ptree}
\noindent$\Pi_{134}$:
\begin{ptree}
\AXC{$\forall x\,\forall y\,(x \in \mathbb{R} \to y \in \mathbb{R} \to x < y \to -y < -x)$\lbl{[F22]}}
\AXC{$\varphi_{248}$\lbl{A2}}
\BIC{$\forall y\,(v_{317} \in \mathbb{R} \to y \in \mathbb{R} \to v_{317} < y \to -y < -v_{317})$}
\AXC{$\varphi_{257}$\lbl{A2}}
\BIC{$\varphi_{101}$}
\UIC{$\forall v_{318}\,(\varphi_{101})$}
\UIC{$\forall v_{317}\,\forall v_{318}\,(\varphi_{101})$}
\AXC{$\forall v_{317}\,\forall v_{318}\,(\varphi_{101}) \to \varphi_{78}$\lbl{A2}}
\BIC{$\varphi_{78}$}
\end{ptree}
\noindent$\Pi_{135}$:
\begin{ptree}
\AXC{$\stackrel{\Pi_{134}}{\vdots}$}\noLine
\UIC{$\varphi_{78}$}
\AXC{$\varphi_{78} \to y \in \mathbb{R} \to 0 \in \mathbb{R} \to y < 0 \to -0 < -y$\lbl{A2}}
\BIC{$y \in \mathbb{R} \to 0 \in \mathbb{R} \to y < 0 \to -0 < -y$}
\AXC{$(y \in \mathbb{R} \to 0 \in \mathbb{R} \to y < 0 \to -0 < -y) \to \varphi_{275}$\lbl{A1}}
\BIC{$\varphi_{275}$}
\end{ptree}
\noindent$\Pi_{136}$:
\begin{ptree}
\AXC{$\stackrel{\Pi_{133}}{\vdots}$}\noLine
\UIC{$x \in \mathbb{R} \to y \in \mathbb{R} \to 0 \leq x \to \neg 0 \leq y \to y < 0$}
\AXC{$0 \in \mathbb{R}$\lbl{(R1)}}
\AXC{$x \in \mathbb{R} \to y \in \mathbb{R} \to 0 \leq x \to \neg 0 \leq y \to y \in \mathbb{R}$\lbl{A1}}
\AXC{$\stackrel{\Pi_{135}}{\vdots}$}\noLine
\UIC{$\varphi_{275}$}
\BIC{$0 \in \mathbb{R} \to \varphi_{195}$}
\BIC{$\varphi_{195}$}
\BIC{$x \in \mathbb{R} \to y \in \mathbb{R} \to 0 \leq x \to \neg 0 \leq y \to -0 < -y$}
\end{ptree}
\noindent$\Pi_{137}$:
\begin{ptree}
\AXC{$\stackrel{\Pi_{136}}{\vdots}$}\noLine
\UIC{$x \in \mathbb{R} \to y \in \mathbb{R} \to 0 \leq x \to \neg 0 \leq y \to -0 < -y$}
\AXC{$-0 = 0$\lbl{[F6]}}
\AXC{$\varphi_{270}$\lbl{A1}}
\BIC{$(x \in \mathbb{R} \to y \in \mathbb{R} \to 0 \leq x \to \neg 0 \leq y \to -0 < -y) \to \varphi_{166}$}
\BIC{$\varphi_{166}$}
\end{ptree}
\noindent$\Pi_{138}$:
\begin{ptree}
\AXC{$\forall x \in \mathbb{R}\,(-x \in \mathbb{R})$\lbl{(R1)}}
\AXC{$\forall x \in \mathbb{R}\,(-x \in \mathbb{R}) \to y \in \mathbb{R} \to -y \in \mathbb{R}$\lbl{A2}}
\BIC{$y \in \mathbb{R} \to -y \in \mathbb{R}$}
\AXC{$(y \in \mathbb{R} \to -y \in \mathbb{R}) \to \varphi_{209}$\lbl{A1}}
\BIC{$\varphi_{209}$}
\end{ptree}
\noindent$\Pi_{139}$:
\begin{ptree}
\AXC{$\forall x\,\forall y\,(x \in \mathbb{R} \to y \in \mathbb{R} \to 0 \leq x \to 0 \leq y \to 0 \leq x \cdot y)$\lbl{[A16]}}
\AXC{$\varphi_{200}$\lbl{A2}}
\BIC{$\forall y\,(x \in \mathbb{R} \to y \in \mathbb{R} \to 0 \leq x \to 0 \leq y \to 0 \leq x \cdot y)$}
\AXC{$\varphi_{230}$\lbl{A2}}
\BIC{$x \in \mathbb{R} \to -y \in \mathbb{R} \to 0 \leq x \to 0 \leq -y \to 0 \leq x \cdot -y$}
\AXC{$\varphi_{325}$\lbl{A1}}
\BIC{$(x \in \mathbb{R} \to y \in \mathbb{R} \to 0 \leq x \to \neg 0 \leq y \to x \in \mathbb{R}) \to \varphi_{291}$}
\end{ptree}
\noindent$\Pi_{140}$:
\begin{ptree}
\AXC{$x \in \mathbb{R} \to y \in \mathbb{R} \to 0 \leq x \to \neg 0 \leq y \to x \in \mathbb{R}$\lbl{A1}}
\AXC{$\stackrel{\Pi_{139}}{\vdots}$}\noLine
\UIC{$(x \in \mathbb{R} \to y \in \mathbb{R} \to 0 \leq x \to \neg 0 \leq y \to x \in \mathbb{R}) \to \varphi_{291}$}
\BIC{$\varphi_{291}$}
\end{ptree}
\noindent$\Pi_{141}$:
\begin{ptree}
\AXC{$x \in \mathbb{R} \to y \in \mathbb{R} \to 0 \leq x \to 0 \leq x$\lbl{A1}}
\AXC{$x \in \mathbb{R} \to y \in \mathbb{R} \to 0 \leq x \to \neg 0 \leq y \to y \in \mathbb{R}$\lbl{A1}}
\AXC{$\stackrel{\Pi_{138}}{\vdots}$}\noLine
\UIC{$\varphi_{209}$}
\BIC{$x \in \mathbb{R} \to y \in \mathbb{R} \to 0 \leq x \to \neg 0 \leq y \to -y \in \mathbb{R}$}
\AXC{$\stackrel{\Pi_{140}}{\vdots}$}\noLine
\UIC{$\varphi_{291}$}
\BIC{$(x \in \mathbb{R} \to y \in \mathbb{R} \to 0 \leq x \to 0 \leq x) \to \varphi_{192}$}
\BIC{$\varphi_{192}$}
\end{ptree}
\noindent$\Pi_{142}$:
\begin{ptree}
\AXC{$z = w \to z < -y \to w < -y$\lbl{A6}}
\UIC{$\forall w\,(z = w \to z < -y \to w < -y)$}
\UIC{$\forall z\,\forall w\,(z = w \to z < -y \to w < -y)$}
\AXC{$\varphi_{131}$\lbl{A2}}
\BIC{$\forall w\,(-0 = w \to -0 < -y \to w < -y)$}
\AXC{$\varphi_{152}$\lbl{A2}}
\BIC{$-0 = 0 \to -0 < -y \to 0 < -y$}
\AXC{$\stackrel{\Pi_{137}}{\vdots}$}\noLine
\UIC{$\varphi_{166}$}
\BIC{$x \in \mathbb{R} \to y \in \mathbb{R} \to 0 \leq x \to \neg 0 \leq y \to 0 < -y$}
\AXC{$\varphi_{186}$\lbl{A1}}
\BIC{$x \in \mathbb{R} \to y \in \mathbb{R} \to 0 \leq x \to \neg 0 \leq y \to 0 \leq -y$}
\AXC{$\stackrel{\Pi_{141}}{\vdots}$}\noLine
\UIC{$\varphi_{192}$}
\BIC{$x \in \mathbb{R} \to y \in \mathbb{R} \to 0 \leq x \to \neg 0 \leq y \to 0 \leq x \cdot -y$}
\end{ptree}
\noindent$\Pi_{143}$:
\begin{ptree}
\AXC{$\forall x\,\forall y\,(x \in \mathbb{R} \to y \in \mathbb{R} \to x \cdot -y = -(x \cdot y))$\lbl{[F10]}}
\AXC{$\varphi_{221}$\lbl{A2}}
\BIC{$\forall y\,(x \in \mathbb{R} \to y \in \mathbb{R} \to x \cdot -y = -(x \cdot y))$}
\AXC{$\varphi_{199}$\lbl{A2}}
\BIC{$x \in \mathbb{R} \to y \in \mathbb{R} \to x \cdot -y = -(x \cdot y)$}
\AXC{$\varphi_{305}$\lbl{A1}}
\BIC{$(x \in \mathbb{R} \to y \in \mathbb{R} \to 0 \leq x \to \neg 0 \leq y \to x \in \mathbb{R}) \to \varphi_{235}$}
\end{ptree}
\noindent$\Pi_{144}$:
\begin{ptree}
\AXC{$x \in \mathbb{R} \to y \in \mathbb{R} \to 0 \leq x \to \neg 0 \leq y \to x \in \mathbb{R}$\lbl{A1}}
\AXC{$\stackrel{\Pi_{143}}{\vdots}$}\noLine
\UIC{$(x \in \mathbb{R} \to y \in \mathbb{R} \to 0 \leq x \to \neg 0 \leq y \to x \in \mathbb{R}) \to \varphi_{235}$}
\BIC{$\varphi_{235}$}
\end{ptree}
\noindent$\Pi_{145}$:
\begin{ptree}
\AXC{$\stackrel{\Pi_{142}}{\vdots}$}\noLine
\UIC{$x \in \mathbb{R} \to y \in \mathbb{R} \to 0 \leq x \to \neg 0 \leq y \to 0 \leq x \cdot -y$}
\AXC{$x \in \mathbb{R} \to y \in \mathbb{R} \to 0 \leq x \to \neg 0 \leq y \to y \in \mathbb{R}$\lbl{A1}}
\AXC{$\stackrel{\Pi_{144}}{\vdots}$}\noLine
\UIC{$\varphi_{235}$}
\BIC{$x \in \mathbb{R} \to y \in \mathbb{R} \to 0 \leq x \to \neg 0 \leq y \to x \cdot -y = -(x \cdot y)$}
\AXC{$\varphi_{323}$\lbl{A1}}
\BIC{$(x \in \mathbb{R} \to y \in \mathbb{R} \to 0 \leq x \to \neg 0 \leq y \to 0 \leq x \cdot -y) \to \varphi_{281}$}
\BIC{$\varphi_{281}$}
\end{ptree}
\noindent$\Pi_{146}$:
\begin{ptree}
\AXC{$z = w \to 0 < z \to 0 < w$\lbl{A6}}
\UIC{$\forall w\,(z = w \to 0 < z \to 0 < w)$}
\UIC{$\forall z\,\forall w\,(z = w \to 0 < z \to 0 < w)$}
\AXC{$\varphi_{113}$\lbl{A2}}
\BIC{$\forall w\,(x \cdot -y = w \to 0 < x \cdot -y \to 0 < w)$}
\AXC{$\varphi_{215}$\lbl{A2}}
\BIC{$x \cdot -y = -(x \cdot y) \to 0 < x \cdot -y \to 0 < -(x \cdot y)$}
\AXC{$\stackrel{\Pi_{145}}{\vdots}$}\noLine
\UIC{$\varphi_{281}$}
\BIC{$\varphi_{222}$}
\end{ptree}
\noindent$\Pi_{147}$:
\begin{ptree}
\AXC{$\forall x\,\forall y\,(x \in \mathbb{R} \to y \in \mathbb{R} \to x \cdot y \in \mathbb{R})$\lbl{(R1)}}
\AXC{$\varphi_{165}$\lbl{A2}}
\BIC{$\forall y\,(x \in \mathbb{R} \to y \in \mathbb{R} \to x \cdot y \in \mathbb{R})$}
\AXC{$\varphi_{155}$\lbl{A2}}
\BIC{$x \in \mathbb{R} \to y \in \mathbb{R} \to x \cdot y \in \mathbb{R}$}
\AXC{$\varphi_{297}$\lbl{A1}}
\BIC{$(x \in \mathbb{R} \to y \in \mathbb{R} \to 0 \leq x \to \neg 0 \leq y \to x \in \mathbb{R}) \to \varphi_{210}$}
\end{ptree}
\noindent$\Pi_{148}$:
\begin{ptree}
\AXC{$x \in \mathbb{R} \to y \in \mathbb{R} \to 0 \leq x \to \neg 0 \leq y \to x \in \mathbb{R}$\lbl{A1}}
\AXC{$\stackrel{\Pi_{147}}{\vdots}$}\noLine
\UIC{$(x \in \mathbb{R} \to y \in \mathbb{R} \to 0 \leq x \to \neg 0 \leq y \to x \in \mathbb{R}) \to \varphi_{210}$}
\BIC{$\varphi_{210}$}
\end{ptree}
\noindent$\Pi_{149}$:
\begin{ptree}
\AXC{$x \in \mathbb{R} \to y \in \mathbb{R} \to 0 \leq x \to \neg 0 \leq y \to y \in \mathbb{R}$\lbl{A1}}
\AXC{$\stackrel{\Pi_{148}}{\vdots}$}\noLine
\UIC{$\varphi_{210}$}
\BIC{$x \in \mathbb{R} \to y \in \mathbb{R} \to 0 \leq x \to \neg 0 \leq y \to x \cdot y \in \mathbb{R}$}
\end{ptree}
\noindent$\Pi_{150}$:
\begin{ptree}
\AXC{$\forall x \in \mathbb{R}\,(-x \in \mathbb{R})$\lbl{(R1)}}
\AXC{$\forall x \in \mathbb{R}\,(-x \in \mathbb{R}) \to x \cdot y \in \mathbb{R} \to -(x \cdot y) \in \mathbb{R}$\lbl{A2}}
\BIC{$x \cdot y \in \mathbb{R} \to -(x \cdot y) \in \mathbb{R}$}
\end{ptree}
\noindent$\Pi_{151}$:
\begin{ptree}
\AXC{$\stackrel{\Pi_{149}}{\vdots}$}\noLine
\UIC{$x \in \mathbb{R} \to y \in \mathbb{R} \to 0 \leq x \to \neg 0 \leq y \to x \cdot y \in \mathbb{R}$}
\AXC{$\stackrel{\Pi_{150}}{\vdots}$}\noLine
\UIC{$x \cdot y \in \mathbb{R} \to -(x \cdot y) \in \mathbb{R}$}
\AXC{$(x \cdot y \in \mathbb{R} \to -(x \cdot y) \in \mathbb{R}) \to \varphi_{241}$\lbl{A1}}
\BIC{$\varphi_{241}$}
\BIC{$x \in \mathbb{R} \to y \in \mathbb{R} \to 0 \leq x \to \neg 0 \leq y \to -(x \cdot y) \in \mathbb{R}$}
\end{ptree}
\noindent$\Pi_{152}$:
\begin{ptree}
\AXC{$\forall x\,\forall y\,(x \in \mathbb{R} \to y \in \mathbb{R} \to x \leq y \to -y \leq -x)$\lbl{[F23]}}
\AXC{$\varphi_{191}$\lbl{A2}}
\BIC{$\forall y\,(0 \in \mathbb{R} \to y \in \mathbb{R} \to 0 \leq y \to -y \leq -0)$}
\AXC{$\forall y\,(0 \in \mathbb{R} \to y \in \mathbb{R} \to 0 \leq y \to -y \leq -0) \to 0 \in \mathbb{R} \to \varphi_{60}$\lbl{A2}}
\BIC{$0 \in \mathbb{R} \to \varphi_{60}$}
\end{ptree}
\noindent$\Pi_{153}$:
\begin{ptree}
\AXC{$\stackrel{\Pi_{151}}{\vdots}$}\noLine
\UIC{$x \in \mathbb{R} \to y \in \mathbb{R} \to 0 \leq x \to \neg 0 \leq y \to -(x \cdot y) \in \mathbb{R}$}
\AXC{$0 \in \mathbb{R}$\lbl{(R1)}}
\AXC{$\stackrel{\Pi_{152}}{\vdots}$}\noLine
\UIC{$0 \in \mathbb{R} \to \varphi_{60}$}
\AXC{$(0 \in \mathbb{R} \to \varphi_{60}) \to 0 \in \mathbb{R} \to \varphi_{286}$\lbl{A1}}
\BIC{$0 \in \mathbb{R} \to \varphi_{286}$}
\BIC{$\varphi_{286}$}
\BIC{$\varphi_{243}$}
\end{ptree}
\noindent$\Pi_{154}$:
\begin{ptree}
\AXC{$\stackrel{\Pi_{149}}{\vdots}$}\noLine
\UIC{$x \in \mathbb{R} \to y \in \mathbb{R} \to 0 \leq x \to \neg 0 \leq y \to x \cdot y \in \mathbb{R}$}
\AXC{$\forall x \in \mathbb{R}\,(-(-x) = x)$\lbl{[F5]}}
\AXC{$\varphi_{111}$\lbl{A2}}
\BIC{$x \cdot y \in \mathbb{R} \to -(-(x \cdot y)) = x \cdot y$}
\AXC{$\varphi_{271}$\lbl{A1}}
\BIC{$(x \in \mathbb{R} \to y \in \mathbb{R} \to 0 \leq x \to \neg 0 \leq y \to x \cdot y \in \mathbb{R}) \to \varphi_{58}$}
\BIC{$\varphi_{58}$}
\AXC{$\varphi_{58} \to \varphi_{317}$\lbl{A1}}
\BIC{$\varphi_{317}$}
\end{ptree}
\noindent$\Pi_{155}$:
\begin{ptree}
\AXC{$z = w \to 0 = z \to 0 = w$\lbl{A6}}
\UIC{$\forall w\,(z = w \to 0 = z \to 0 = w)$}
\UIC{$\forall z\,\forall w\,(z = w \to 0 = z \to 0 = w)$}
\AXC{$\varphi_{56}$\lbl{A2}}
\BIC{$\forall w\,(x \cdot -y = w \to 0 = x \cdot -y \to 0 = w)$}
\AXC{$\varphi_{175}$\lbl{A2}}
\BIC{$x \cdot -y = -(x \cdot y) \to 0 = x \cdot -y \to 0 = -(x \cdot y)$}
\AXC{$\stackrel{\Pi_{146}}{\vdots}$}\noLine
\UIC{$\varphi_{222}$}
\BIC{$x \in \mathbb{R} \to y \in \mathbb{R} \to 0 \leq x \to \neg 0 \leq y \to 0 \leq -(x \cdot y)$}
\AXC{$\stackrel{\Pi_{153}}{\vdots}$}\noLine
\UIC{$\varphi_{243}$}
\BIC{$\varphi_{57}$}
\AXC{$\stackrel{\Pi_{154}}{\vdots}$}\noLine
\UIC{$\varphi_{317}$}
\BIC{$\varphi_{103} \to \varphi_{67} \to x \in \mathbb{R} \to y \in \mathbb{R} \to 0 \leq x \to \neg 0 \leq y \to x \cdot y \leq -0$}
\end{ptree}
\noindent$\Pi_{156}$:
\begin{ptree}
\AXC{$\stackrel{\Pi_{129}}{\vdots}$}\noLine
\UIC{$\varphi_{67}$}
\AXC{$\stackrel{\Pi_{130}}{\vdots}$}\noLine
\UIC{$\varphi_{103}$}
\AXC{$\stackrel{\Pi_{155}}{\vdots}$}\noLine
\UIC{$\varphi_{103} \to \varphi_{67} \to x \in \mathbb{R} \to y \in \mathbb{R} \to 0 \leq x \to \neg 0 \leq y \to x \cdot y \leq -0$}
\BIC{$\varphi_{67} \to x \in \mathbb{R} \to y \in \mathbb{R} \to 0 \leq x \to \neg 0 \leq y \to x \cdot y \leq -0$}
\BIC{$x \in \mathbb{R} \to y \in \mathbb{R} \to 0 \leq x \to \neg 0 \leq y \to x \cdot y \leq -0$}
\AXC{$-0 = 0$\lbl{[F6]}}
\AXC{$-0 = 0 \to \varphi_{289}$\lbl{A1}}
\BIC{$\varphi_{289}$}
\BIC{$(-0 = 0 \to x \cdot y < -0 \to x \cdot y < 0) \to \varphi_{143}$}
\end{ptree}
\noindent$\Pi_{157}$:
\begin{ptree}
\AXC{$z = w \to x \cdot y < z \to x \cdot y < w$\lbl{A6}}
\UIC{$\forall w\,(z = w \to x \cdot y < z \to x \cdot y < w)$}
\UIC{$\forall z\,\forall w\,(z = w \to x \cdot y < z \to x \cdot y < w)$}
\AXC{$\varphi_{133}$\lbl{A2}}
\BIC{$\forall w\,(-0 = w \to x \cdot y < -0 \to x \cdot y < w)$}
\AXC{$\varphi_{154}$\lbl{A2}}
\BIC{$-0 = 0 \to x \cdot y < -0 \to x \cdot y < 0$}
\AXC{$\stackrel{\Pi_{156}}{\vdots}$}\noLine
\UIC{$(-0 = 0 \to x \cdot y < -0 \to x \cdot y < 0) \to \varphi_{143}$}
\BIC{$\varphi_{143}$}
\end{ptree}
\noindent$\Pi_{158}$:
\begin{ptree}
\AXC{$z = w \to x \cdot y = z \to x \cdot y = w$\lbl{A6}}
\UIC{$\forall w\,(z = w \to x \cdot y = z \to x \cdot y = w)$}
\UIC{$\forall z\,\forall w\,(z = w \to x \cdot y = z \to x \cdot y = w)$}
\AXC{$\varphi_{106}$\lbl{A2}}
\BIC{$\forall w\,(-0 = w \to x \cdot y = -0 \to x \cdot y = w)$}
\AXC{$\varphi_{121}$\lbl{A2}}
\BIC{$-0 = 0 \to x \cdot y = -0 \to x \cdot y = 0$}
\AXC{$\stackrel{\Pi_{157}}{\vdots}$}\noLine
\UIC{$\varphi_{143}$}
\BIC{$x \in \mathbb{R} \to y \in \mathbb{R} \to 0 \leq x \to \neg 0 \leq y \to x \cdot y \leq 0$}
\end{ptree}
\noindent$\Pi_{159}$:
\begin{ptree}
\AXC{$\forall x \in \mathbb{R}\,(x \leq 0 \to \lvert x \rvert = -x)$\lbl{[A4]}}
\AXC{$\varphi_{157}$\lbl{A2}}
\BIC{$x \cdot y \in \mathbb{R} \to x \cdot y \leq 0 \to \lvert x \cdot y \rvert = -(x \cdot y)$}
\AXC{$(x \cdot y \in \mathbb{R} \to x \cdot y \leq 0 \to \lvert x \cdot y \rvert = -(x \cdot y)) \to \varphi_{277}$\lbl{A1}}
\BIC{$\varphi_{277}$}
\end{ptree}
\noindent$\Pi_{160}$:
\begin{ptree}
\AXC{$\stackrel{\Pi_{158}}{\vdots}$}\noLine
\UIC{$x \in \mathbb{R} \to y \in \mathbb{R} \to 0 \leq x \to \neg 0 \leq y \to x \cdot y \leq 0$}
\AXC{$\stackrel{\Pi_{149}}{\vdots}$}\noLine
\UIC{$x \in \mathbb{R} \to y \in \mathbb{R} \to 0 \leq x \to \neg 0 \leq y \to x \cdot y \in \mathbb{R}$}
\AXC{$\stackrel{\Pi_{159}}{\vdots}$}\noLine
\UIC{$\varphi_{277}$}
\BIC{$(x \in \mathbb{R} \to y \in \mathbb{R} \to 0 \leq x \to \neg 0 \leq y \to x \cdot y \leq 0) \to \varphi_{59}$}
\BIC{$\varphi_{59}$}
\AXC{$\varphi_{59} \to \varphi_{299}$\lbl{A1}}
\BIC{$\varphi_{299}$}
\end{ptree}
\noindent$\Pi_{161}$:
\begin{ptree}
\AXC{$\stackrel{\Pi_{122}}{\vdots}$}\noLine
\UIC{$\varphi_{91}$}
\AXC{$\stackrel{\Pi_{124}}{\vdots}$}\noLine
\UIC{$\varphi_{118}$}
\AXC{$\stackrel{\Pi_{128}}{\vdots}$}\noLine
\UIC{$x \in \mathbb{R} \to y \in \mathbb{R} \to 0 \leq x \to \neg 0 \leq y \to -(x \cdot y) = x \cdot -y$}
\AXC{$\stackrel{\Pi_{160}}{\vdots}$}\noLine
\UIC{$\varphi_{299}$}
\BIC{$\varphi_{260}$}
\BIC{$x \in \mathbb{R} \to y \in \mathbb{R} \to 0 \leq x \to \neg 0 \leq y \to \lvert x \cdot y \rvert = x \cdot -y$}
\AXC{$\varphi_{321}$\lbl{A1}}
\BIC{$\varphi_{91} \to \varphi_{140} \to \varphi_{93}$}
\BIC{$\varphi_{140} \to \varphi_{93}$}
\end{ptree}
\noindent$\Pi_{162}$:
\begin{ptree}
\AXC{$\stackrel{\Pi_{109}}{\vdots}$}\noLine
\UIC{$\forall v_{326}\,(\varphi_{14})$}
\AXC{$\forall v_{326}\,(\varphi_{14}) \to \varphi_{38}$\lbl{A2}}
\BIC{$\varphi_{38}$}
\AXC{$\varphi_{38} \to \varphi_{88}$\lbl{A2}}
\BIC{$\varphi_{88}$}
\AXC{$\varphi_{88} \to \varphi_{140}$\lbl{A2}}
\BIC{$\varphi_{140}$}
\AXC{$\stackrel{\Pi_{161}}{\vdots}$}\noLine
\UIC{$\varphi_{140} \to \varphi_{93}$}
\BIC{$\varphi_{93}$}
\end{ptree}
\noindent$\Pi_{163}$:
\begin{ptree}
\AXC{$y = z \to x = y \to x = z$\lbl{A6}}
\AXC{$(y = z \to x = y \to x = z) \to x = y \to y = z \to x = z$\lbl{A1}}
\BIC{$x = y \to y = z \to x = z$}
\UIC{$\forall z\,(x = y \to y = z \to x = z)$}
\UIC{$\forall y\,\forall z\,(x = y \to y = z \to x = z)$}
\UIC{$\forall x\,\forall y\,\forall z\,(x = y \to y = z \to x = z)$}
\AXC{$\varphi_{28}$\lbl{A2}}
\BIC{$\forall y\,\forall z\,(v_{307} = y \to y = z \to v_{307} = z)$}
\AXC{$\varphi_{50}$\lbl{A2}}
\BIC{$\forall z\,(v_{307} = v_{308} \to v_{308} = z \to v_{307} = z)$}
\AXC{$\varphi_{79}$\lbl{A2}}
\BIC{$v_{307} = v_{308} \to v_{308} = v_{309} \to v_{307} = v_{309}$}
\UIC{$\forall v_{309}\,(v_{307} = v_{308} \to v_{308} = v_{309} \to v_{307} = v_{309})$}
\UIC{$\varphi_{17}$}
\UIC{$\forall v_{307}\,(\varphi_{17})$}
\end{ptree}
\noindent$\Pi_{164}$:
\begin{ptree}
\AXC{$\stackrel{\Pi_{2}}{\vdots}$}\noLine
\UIC{$\forall x\,\forall y\,(x = y \to y = x)$}
\AXC{$\forall x\,\forall y\,(x = y \to y = x) \to \forall y\,(v_{305} = y \to y = v_{305})$\lbl{A2}}
\BIC{$\forall y\,(v_{305} = y \to y = v_{305})$}
\AXC{$\varphi_{7}$\lbl{A2}}
\BIC{$v_{305} = v_{306} \to v_{306} = v_{305}$}
\UIC{$\forall v_{306}\,(v_{305} = v_{306} \to v_{306} = v_{305})$}
\UIC{$\forall v_{305}\,\forall v_{306}\,(v_{305} = v_{306} \to v_{306} = v_{305})$}
\AXC{$\varphi_{151}$\lbl{A2}}
\BIC{$\forall v_{306}\,(\lvert x \rvert \cdot \lvert y \rvert = v_{306} \to v_{306} = \lvert x \rvert \cdot \lvert y \rvert)$}
\AXC{$\varphi_{196}$\lbl{A2}}
\BIC{$\lvert x \rvert \cdot \lvert y \rvert = x \cdot y \to x \cdot y = \lvert x \rvert \cdot \lvert y \rvert$}
\end{ptree}
\noindent$\Pi_{165}$:
\begin{ptree}
\AXC{$y = z \to x = y \to x = z$\lbl{A6}}
\AXC{$(y = z \to x = y \to x = z) \to x = y \to y = z \to x = z$\lbl{A1}}
\BIC{$x = y \to y = z \to x = z$}
\UIC{$\forall z\,(x = y \to y = z \to x = z)$}
\UIC{$\forall y\,\forall z\,(x = y \to y = z \to x = z)$}
\UIC{$\forall x\,\forall y\,\forall z\,(x = y \to y = z \to x = z)$}
\AXC{$\varphi_{29}$\lbl{A2}}
\BIC{$\forall y\,\forall z\,(v_{302} = y \to y = z \to v_{302} = z)$}
\AXC{$\varphi_{51}$\lbl{A2}}
\BIC{$\forall z\,(v_{302} = v_{303} \to v_{303} = z \to v_{302} = z)$}
\AXC{$\varphi_{80}$\lbl{A2}}
\BIC{$v_{302} = v_{303} \to v_{303} = v_{304} \to v_{302} = v_{304}$}
\UIC{$\forall v_{304}\,(v_{302} = v_{303} \to v_{303} = v_{304} \to v_{302} = v_{304})$}
\UIC{$\varphi_{18}$}
\UIC{$\forall v_{302}\,(\varphi_{18})$}
\end{ptree}
\noindent$\Pi_{166}$:
\begin{ptree}
\AXC{$\stackrel{\Pi_{165}}{\vdots}$}\noLine
\UIC{$\forall v_{302}\,(\varphi_{18})$}
\AXC{$\forall v_{302}\,(\varphi_{18}) \to \varphi_{102}$\lbl{A2}}
\BIC{$\varphi_{102}$}
\AXC{$\varphi_{102} \to \varphi_{126}$\lbl{A2}}
\BIC{$\varphi_{126}$}
\AXC{$\varphi_{126} \to \varphi_{120}$\lbl{A2}}
\BIC{$\varphi_{120}$}
\end{ptree}
\noindent$\Pi_{167}$:
\begin{ptree}
\AXC{$z = w \to x \cdot \lvert y \rvert = x \cdot z \to x \cdot \lvert y \rvert = x \cdot w$\lbl{A6}}
\UIC{$\forall w\,(z = w \to x \cdot \lvert y \rvert = x \cdot z \to x \cdot \lvert y \rvert = x \cdot w)$}
\UIC{$\forall z\,\forall w\,(z = w \to x \cdot \lvert y \rvert = x \cdot z \to x \cdot \lvert y \rvert = x \cdot w)$}
\AXC{$\varphi_{227}$\lbl{A2}}
\BIC{$\varphi_{33}$}
\AXC{$\varphi_{33} \to \lvert y \rvert = y \to x \cdot \lvert y \rvert = x \cdot \lvert y \rvert \to x \cdot \lvert y \rvert = x \cdot y$\lbl{A2}}
\BIC{$\lvert y \rvert = y \to x \cdot \lvert y \rvert = x \cdot \lvert y \rvert \to x \cdot \lvert y \rvert = x \cdot y$}
\end{ptree}
\noindent$\Pi_{168}$:
\begin{ptree}
\AXC{$\forall x\,(0 \leq x \to \lvert x \rvert = x)$\lbl{[A1]}}
\AXC{$\forall x\,(0 \leq x \to \lvert x \rvert = x) \to 0 \leq y \to \lvert y \rvert = y$\lbl{A2}}
\BIC{$0 \leq y \to \lvert y \rvert = y$}
\AXC{$(0 \leq y \to \lvert y \rvert = y) \to \varphi_{158}$\lbl{A1}}
\BIC{$\varphi_{158}$}
\end{ptree}
\noindent$\Pi_{169}$:
\begin{ptree}
\AXC{$x = x$\lbl{A5}}
\UIC{$\forall x\,(x = x)$}
\AXC{$\forall x\,(x = x) \to x \cdot \lvert y \rvert = x \cdot \lvert y \rvert$\lbl{A2}}
\BIC{$x \cdot \lvert y \rvert = x \cdot \lvert y \rvert$}
\AXC{$x \in \mathbb{R} \to y \in \mathbb{R} \to 0 \leq x \to 0 \leq y \to 0 \leq y$\lbl{A1}}
\AXC{$\stackrel{\Pi_{168}}{\vdots}$}\noLine
\UIC{$\varphi_{158}$}
\BIC{$x \in \mathbb{R} \to y \in \mathbb{R} \to 0 \leq x \to 0 \leq y \to \lvert y \rvert = y$}
\AXC{$\varphi_{290}$\lbl{A1}}
\BIC{$x \cdot \lvert y \rvert = x \cdot \lvert y \rvert \to \varphi_{237}$}
\BIC{$\varphi_{237}$}
\end{ptree}
\noindent$\Pi_{170}$:
\begin{ptree}
\AXC{$\forall x\,(0 \leq x \to \lvert x \rvert = x)$\lbl{[A1]}}
\AXC{$\forall x\,(0 \leq x \to \lvert x \rvert = x) \to 0 \leq x \to \lvert x \rvert = x$\lbl{A2}}
\BIC{$0 \leq x \to \lvert x \rvert = x$}
\AXC{$(0 \leq x \to \lvert x \rvert = x) \to \varphi_{134}$\lbl{A1}}
\BIC{$\varphi_{134}$}
\end{ptree}
\noindent$\Pi_{171}$:
\begin{ptree}
\AXC{$x = x$\lbl{A5}}
\UIC{$\forall x\,(x = x)$}
\AXC{$\forall x\,(x = x) \to \lvert x \rvert \cdot \lvert y \rvert = \lvert x \rvert \cdot \lvert y \rvert$\lbl{A2}}
\BIC{$\lvert x \rvert \cdot \lvert y \rvert = \lvert x \rvert \cdot \lvert y \rvert$}
\AXC{$x \in \mathbb{R} \to y \in \mathbb{R} \to 0 \leq x \to 0 \leq x$\lbl{A1}}
\AXC{$\stackrel{\Pi_{170}}{\vdots}$}\noLine
\UIC{$\varphi_{134}$}
\BIC{$x \in \mathbb{R} \to y \in \mathbb{R} \to 0 \leq x \to 0 \leq y \to \lvert x \rvert = x$}
\AXC{$\varphi_{308}$\lbl{A1}}
\BIC{$\lvert x \rvert \cdot \lvert y \rvert = \lvert x \rvert \cdot \lvert y \rvert \to \varphi_{117} \to \varphi_{81}$}
\BIC{$\varphi_{117} \to \varphi_{81}$}
\end{ptree}
\noindent$\Pi_{172}$:
\begin{ptree}
\AXC{$\varphi_{66}$\lbl{A6}}
\UIC{$\forall w\,(\varphi_{66})$}
\UIC{$\forall z\,\forall w\,(\varphi_{66})$}
\AXC{$\forall z\,\forall w\,(\varphi_{66}) \to \varphi_{128}$\lbl{A2}}
\BIC{$\varphi_{128}$}
\AXC{$\varphi_{128} \to \varphi_{117}$\lbl{A2}}
\BIC{$\varphi_{117}$}
\AXC{$\stackrel{\Pi_{171}}{\vdots}$}\noLine
\UIC{$\varphi_{117} \to \varphi_{81}$}
\BIC{$\varphi_{81}$}
\AXC{$\varphi_{81} \to \varphi_{292}$\lbl{A1}}
\BIC{$\varphi_{292}$}
\end{ptree}
\noindent$\Pi_{173}$:
\begin{ptree}
\AXC{$\stackrel{\Pi_{166}}{\vdots}$}\noLine
\UIC{$\varphi_{120}$}
\AXC{$\stackrel{\Pi_{167}}{\vdots}$}\noLine
\UIC{$\lvert y \rvert = y \to x \cdot \lvert y \rvert = x \cdot \lvert y \rvert \to x \cdot \lvert y \rvert = x \cdot y$}
\AXC{$\stackrel{\Pi_{169}}{\vdots}$}\noLine
\UIC{$\varphi_{237}$}
\BIC{$x \in \mathbb{R} \to y \in \mathbb{R} \to 0 \leq x \to 0 \leq y \to x \cdot \lvert y \rvert = x \cdot y$}
\AXC{$\stackrel{\Pi_{172}}{\vdots}$}\noLine
\UIC{$\varphi_{292}$}
\BIC{$\varphi_{258}$}
\BIC{$x \in \mathbb{R} \to y \in \mathbb{R} \to 0 \leq x \to 0 \leq y \to \lvert x \rvert \cdot \lvert y \rvert = x \cdot y$}
\AXC{$\varphi_{282}$\lbl{A1}}
\BIC{$\varphi_{231}$}
\end{ptree}
\noindent$\Pi_{174}$:
\begin{ptree}
\AXC{$\forall x\,\forall y\,(x \in \mathbb{R} \to y \in \mathbb{R} \to 0 \leq x \to 0 \leq y \to 0 \leq x \cdot y)$\lbl{[A16]}}
\AXC{$\varphi_{200}$\lbl{A2}}
\BIC{$\forall y\,(x \in \mathbb{R} \to y \in \mathbb{R} \to 0 \leq x \to 0 \leq y \to 0 \leq x \cdot y)$}
\AXC{$\varphi_{178}$\lbl{A2}}
\BIC{$x \in \mathbb{R} \to y \in \mathbb{R} \to 0 \leq x \to 0 \leq y \to 0 \leq x \cdot y$}
\end{ptree}
\noindent$\Pi_{175}$:
\begin{ptree}
\AXC{$\forall x\,(0 \leq x \to \lvert x \rvert = x)$\lbl{[A1]}}
\AXC{$\forall x\,(0 \leq x \to \lvert x \rvert = x) \to 0 \leq x \cdot y \to \lvert x \cdot y \rvert = x \cdot y$\lbl{A2}}
\BIC{$0 \leq x \cdot y \to \lvert x \cdot y \rvert = x \cdot y$}
\end{ptree}
\noindent$\Pi_{176}$:
\begin{ptree}
\AXC{$\stackrel{\Pi_{174}}{\vdots}$}\noLine
\UIC{$x \in \mathbb{R} \to y \in \mathbb{R} \to 0 \leq x \to 0 \leq y \to 0 \leq x \cdot y$}
\AXC{$\stackrel{\Pi_{175}}{\vdots}$}\noLine
\UIC{$0 \leq x \cdot y \to \lvert x \cdot y \rvert = x \cdot y$}
\AXC{$(0 \leq x \cdot y \to \lvert x \cdot y \rvert = x \cdot y) \to \varphi_{204}$\lbl{A1}}
\BIC{$\varphi_{204}$}
\BIC{$x \in \mathbb{R} \to y \in \mathbb{R} \to 0 \leq x \to 0 \leq y \to \lvert x \cdot y \rvert = x \cdot y$}
\end{ptree}
\noindent$\Pi_{177}$:
\begin{ptree}
\AXC{$\stackrel{\Pi_{164}}{\vdots}$}\noLine
\UIC{$\lvert x \rvert \cdot \lvert y \rvert = x \cdot y \to x \cdot y = \lvert x \rvert \cdot \lvert y \rvert$}
\AXC{$\stackrel{\Pi_{173}}{\vdots}$}\noLine
\UIC{$\varphi_{231}$}
\BIC{$x \in \mathbb{R} \to y \in \mathbb{R} \to 0 \leq x \to 0 \leq y \to x \cdot y = \lvert x \rvert \cdot \lvert y \rvert$}
\AXC{$\stackrel{\Pi_{176}}{\vdots}$}\noLine
\UIC{$x \in \mathbb{R} \to y \in \mathbb{R} \to 0 \leq x \to 0 \leq y \to \lvert x \cdot y \rvert = x \cdot y$}
\AXC{$\varphi_{318}$\lbl{A1}}
\BIC{$\varphi_{300}$}
\BIC{$\varphi_{119} \to \varphi_{82}$}
\end{ptree}
\noindent$\Pi_{178}$:
\begin{ptree}
\AXC{$\stackrel{\Pi_{163}}{\vdots}$}\noLine
\UIC{$\forall v_{307}\,(\varphi_{17})$}
\AXC{$\forall v_{307}\,(\varphi_{17}) \to \varphi_{40}$\lbl{A2}}
\BIC{$\varphi_{40}$}
\AXC{$\varphi_{40} \to \varphi_{31}$\lbl{A2}}
\BIC{$\varphi_{31}$}
\AXC{$\varphi_{31} \to \varphi_{119}$\lbl{A2}}
\BIC{$\varphi_{119}$}
\AXC{$\stackrel{\Pi_{177}}{\vdots}$}\noLine
\UIC{$\varphi_{119} \to \varphi_{82}$}
\BIC{$\varphi_{82}$}
\AXC{$\varphi_{327}$\lbl{A1}}
\BIC{$\varphi_{93} \to \varphi_{87} \to \varphi_{96} \to x \in \mathbb{R} \to y \in \mathbb{R} \to \lvert x \cdot y \rvert = \lvert x \rvert \cdot \lvert y \rvert$}
\end{ptree}
\noindent 主木:
\begin{ptree}
\AXC{$\stackrel{\Pi_{53}}{\vdots}$}\noLine
\UIC{$\varphi_{96}$}
\AXC{$\stackrel{\Pi_{108}}{\vdots}$}\noLine
\UIC{$\varphi_{87}$}
\AXC{$\stackrel{\Pi_{162}}{\vdots}$}\noLine
\UIC{$\varphi_{93}$}
\AXC{$\stackrel{\Pi_{178}}{\vdots}$}\noLine
\UIC{$\varphi_{93} \to \varphi_{87} \to \varphi_{96} \to x \in \mathbb{R} \to y \in \mathbb{R} \to \lvert x \cdot y \rvert = \lvert x \rvert \cdot \lvert y \rvert$}
\BIC{$\varphi_{87} \to \varphi_{96} \to x \in \mathbb{R} \to y \in \mathbb{R} \to \lvert x \cdot y \rvert = \lvert x \rvert \cdot \lvert y \rvert$}
\BIC{$\varphi_{96} \to x \in \mathbb{R} \to y \in \mathbb{R} \to \lvert x \cdot y \rvert = \lvert x \rvert \cdot \lvert y \rvert$}
\BIC{$x \in \mathbb{R} \to y \in \mathbb{R} \to \lvert x \cdot y \rvert = \lvert x \rvert \cdot \lvert y \rvert$}
\UIC{$\forall y\,(x \in \mathbb{R} \to y \in \mathbb{R} \to \lvert x \cdot y \rvert = \lvert x \rvert \cdot \lvert y \rvert)$}
\UIC{$\forall x\,\forall y\,(x \in \mathbb{R} \to y \in \mathbb{R} \to \lvert x \cdot y \rvert = \lvert x \rvert \cdot \lvert y \rvert)$}
\end{ptree}
\begin{abbrevlist}
\item $\varphi_{1} :\equiv \forall y\,(v_{356} = y \to y = v_{356}) \to v_{356} = v_{357} \to v_{357} = v_{356}$
\item $\varphi_{2} :\equiv \forall y\,(v_{351} = y \to y = v_{351}) \to v_{351} = v_{352} \to v_{352} = v_{351}$
\item $\varphi_{3} :\equiv \forall y\,(v_{337} = y \to y = v_{337}) \to v_{337} = v_{338} \to v_{338} = v_{337}$
\item $\varphi_{4} :\equiv \forall y\,(v_{332} = y \to y = v_{332}) \to v_{332} = v_{333} \to v_{333} = v_{332}$
\item $\varphi_{5} :\equiv \forall y\,(v_{324} = y \to y = v_{324}) \to v_{324} = v_{325} \to v_{325} = v_{324}$
\item $\varphi_{6} :\equiv \forall y\,(v_{319} = y \to y = v_{319}) \to v_{319} = v_{320} \to v_{320} = v_{319}$
\item $\varphi_{7} :\equiv \forall y\,(v_{305} = y \to y = v_{305}) \to v_{305} = v_{306} \to v_{306} = v_{305}$
\item $\varphi_{8} :\equiv \forall v_{359}\,\forall v_{360}\,(v_{358} = v_{359} \to v_{359} = v_{360} \to v_{358} = v_{360})$
\item $\varphi_{9} :\equiv \forall v_{345}\,\forall v_{346}\,(v_{344} = v_{345} \to v_{345} = v_{346} \to v_{344} = v_{346})$
\item $\varphi_{10} :\equiv \forall v_{354}\,\forall v_{355}\,(v_{353} = v_{354} \to v_{354} = v_{355} \to v_{353} = v_{355})$
\item $\varphi_{11} :\equiv \forall v_{340}\,\forall v_{341}\,(v_{339} = v_{340} \to v_{340} = v_{341} \to v_{339} = v_{341})$
\item $\varphi_{12} :\equiv \forall v_{330}\,\forall v_{331}\,(v_{329} = v_{330} \to v_{330} = v_{331} \to v_{329} = v_{331})$
\item $\varphi_{13} :\equiv \forall v_{335}\,\forall v_{336}\,(v_{334} = v_{335} \to v_{335} = v_{336} \to v_{334} = v_{336})$
\item $\varphi_{14} :\equiv \forall v_{327}\,\forall v_{328}\,(v_{326} = v_{327} \to v_{327} = v_{328} \to v_{326} = v_{328})$
\item $\varphi_{15} :\equiv \forall v_{313}\,\forall v_{314}\,(v_{312} = v_{313} \to v_{313} = v_{314} \to v_{312} = v_{314})$
\item $\varphi_{16} :\equiv \forall v_{322}\,\forall v_{323}\,(v_{321} = v_{322} \to v_{322} = v_{323} \to v_{321} = v_{323})$
\item $\varphi_{17} :\equiv \forall v_{308}\,\forall v_{309}\,(v_{307} = v_{308} \to v_{308} = v_{309} \to v_{307} = v_{309})$
\item $\varphi_{18} :\equiv \forall v_{303}\,\forall v_{304}\,(v_{302} = v_{303} \to v_{303} = v_{304} \to v_{302} = v_{304})$
\item $\varphi_{19} :\equiv \forall x\,\forall y\,\forall z\,(x = y \to y = z \to x = z) \to \forall y\,\forall z\,(v_{358} = y \to y = z \to v_{358} = z)$
\item $\varphi_{20} :\equiv \forall x\,\forall y\,\forall z\,(x = y \to y = z \to x = z) \to \forall y\,\forall z\,(v_{344} = y \to y = z \to v_{344} = z)$
\item $\varphi_{21} :\equiv \forall x\,\forall y\,\forall z\,(x = y \to y = z \to x = z) \to \forall y\,\forall z\,(v_{353} = y \to y = z \to v_{353} = z)$
\item $\varphi_{22} :\equiv \forall x\,\forall y\,\forall z\,(x = y \to y = z \to x = z) \to \forall y\,\forall z\,(v_{339} = y \to y = z \to v_{339} = z)$
\item $\varphi_{23} :\equiv \forall x\,\forall y\,\forall z\,(x = y \to y = z \to x = z) \to \forall y\,\forall z\,(v_{329} = y \to y = z \to v_{329} = z)$
\item $\varphi_{24} :\equiv \forall x\,\forall y\,\forall z\,(x = y \to y = z \to x = z) \to \forall y\,\forall z\,(v_{334} = y \to y = z \to v_{334} = z)$
\item $\varphi_{25} :\equiv \forall x\,\forall y\,\forall z\,(x = y \to y = z \to x = z) \to \forall y\,\forall z\,(v_{326} = y \to y = z \to v_{326} = z)$
\item $\varphi_{26} :\equiv \forall x\,\forall y\,\forall z\,(x = y \to y = z \to x = z) \to \forall y\,\forall z\,(v_{312} = y \to y = z \to v_{312} = z)$
\item $\varphi_{27} :\equiv \forall x\,\forall y\,\forall z\,(x = y \to y = z \to x = z) \to \forall y\,\forall z\,(v_{321} = y \to y = z \to v_{321} = z)$
\item $\varphi_{28} :\equiv \forall x\,\forall y\,\forall z\,(x = y \to y = z \to x = z) \to \forall y\,\forall z\,(v_{307} = y \to y = z \to v_{307} = z)$
\item $\varphi_{29} :\equiv \forall x\,\forall y\,\forall z\,(x = y \to y = z \to x = z) \to \forall y\,\forall z\,(v_{302} = y \to y = z \to v_{302} = z)$
\item $\varphi_{30} :\equiv \forall v_{355}\,(\lvert x \cdot y \rvert = x \cdot y \to x \cdot y = v_{355} \to \lvert x \cdot y \rvert = v_{355})$
\item $\varphi_{31} :\equiv \forall v_{309}\,(\lvert x \cdot y \rvert = x \cdot y \to x \cdot y = v_{309} \to \lvert x \cdot y \rvert = v_{309})$
\item $\varphi_{32} :\equiv \forall w\,(-(-(x \cdot y)) = w \to -(-(x \cdot y)) = -0 \to w = -0)$
\item $\varphi_{33} :\equiv \forall w\,(\lvert y \rvert = w \to x \cdot \lvert y \rvert = x \cdot \lvert y \rvert \to x \cdot \lvert y \rvert = x \cdot w)$
\item $\varphi_{34} :\equiv \forall v_{359}\,\forall v_{360}\,(\lvert x \cdot y \rvert = v_{359} \to v_{359} = v_{360} \to \lvert x \cdot y \rvert = v_{360})$
\item $\varphi_{35} :\equiv \forall v_{354}\,\forall v_{355}\,(\lvert x \cdot y \rvert = v_{354} \to v_{354} = v_{355} \to \lvert x \cdot y \rvert = v_{355})$
\item $\varphi_{36} :\equiv \forall v_{340}\,\forall v_{341}\,(\lvert x \cdot y \rvert = v_{340} \to v_{340} = v_{341} \to \lvert x \cdot y \rvert = v_{341})$
\item $\varphi_{37} :\equiv \forall v_{335}\,\forall v_{336}\,(\lvert x \cdot y \rvert = v_{335} \to v_{335} = v_{336} \to \lvert x \cdot y \rvert = v_{336})$
\item $\varphi_{38} :\equiv \forall v_{327}\,\forall v_{328}\,(\lvert x \cdot y \rvert = v_{327} \to v_{327} = v_{328} \to \lvert x \cdot y \rvert = v_{328})$
\item $\varphi_{39} :\equiv \forall v_{322}\,\forall v_{323}\,(\lvert x \cdot y \rvert = v_{322} \to v_{322} = v_{323} \to \lvert x \cdot y \rvert = v_{323})$
\item $\varphi_{40} :\equiv \forall v_{308}\,\forall v_{309}\,(\lvert x \cdot y \rvert = v_{308} \to v_{308} = v_{309} \to \lvert x \cdot y \rvert = v_{309})$
\item $\varphi_{41} :\equiv \forall y\,\forall z\,(v_{358} = y \to y = z \to v_{358} = z) \to \forall z\,(v_{358} = v_{359} \to v_{359} = z \to v_{358} = z)$
\item $\varphi_{42} :\equiv \forall y\,\forall z\,(v_{344} = y \to y = z \to v_{344} = z) \to \forall z\,(v_{344} = v_{345} \to v_{345} = z \to v_{344} = z)$
\item $\varphi_{43} :\equiv \forall y\,\forall z\,(v_{353} = y \to y = z \to v_{353} = z) \to \forall z\,(v_{353} = v_{354} \to v_{354} = z \to v_{353} = z)$
\item $\varphi_{44} :\equiv \forall y\,\forall z\,(v_{339} = y \to y = z \to v_{339} = z) \to \forall z\,(v_{339} = v_{340} \to v_{340} = z \to v_{339} = z)$
\item $\varphi_{45} :\equiv \forall y\,\forall z\,(v_{329} = y \to y = z \to v_{329} = z) \to \forall z\,(v_{329} = v_{330} \to v_{330} = z \to v_{329} = z)$
\item $\varphi_{46} :\equiv \forall y\,\forall z\,(v_{334} = y \to y = z \to v_{334} = z) \to \forall z\,(v_{334} = v_{335} \to v_{335} = z \to v_{334} = z)$
\item $\varphi_{47} :\equiv \forall y\,\forall z\,(v_{326} = y \to y = z \to v_{326} = z) \to \forall z\,(v_{326} = v_{327} \to v_{327} = z \to v_{326} = z)$
\item $\varphi_{48} :\equiv \forall y\,\forall z\,(v_{312} = y \to y = z \to v_{312} = z) \to \forall z\,(v_{312} = v_{313} \to v_{313} = z \to v_{312} = z)$
\item $\varphi_{49} :\equiv \forall y\,\forall z\,(v_{321} = y \to y = z \to v_{321} = z) \to \forall z\,(v_{321} = v_{322} \to v_{322} = z \to v_{321} = z)$
\item $\varphi_{50} :\equiv \forall y\,\forall z\,(v_{307} = y \to y = z \to v_{307} = z) \to \forall z\,(v_{307} = v_{308} \to v_{308} = z \to v_{307} = z)$
\item $\varphi_{51} :\equiv \forall y\,\forall z\,(v_{302} = y \to y = z \to v_{302} = z) \to \forall z\,(v_{302} = v_{303} \to v_{303} = z \to v_{302} = z)$
\item $\varphi_{52} :\equiv \forall z\,\forall w\,(z = w \to 0 = z \to 0 = w) \to \forall w\,(-x \cdot y = w \to 0 = -x \cdot y \to 0 = w)$
\item $\varphi_{53} :\equiv x \in \mathbb{R} \to y \in \mathbb{R} \to \neg 0 \leq x \to 0 \leq y \to -(-(x \cdot y)) \leq -0$
\item $\varphi_{54} :\equiv x \in \mathbb{R} \to y \in \mathbb{R} \to \neg 0 \leq x \to 0 \leq y \to -(-(x \cdot y)) = x \cdot y$
\item $\varphi_{55} :\equiv x \in \mathbb{R} \to y \in \mathbb{R} \to \neg 0 \leq x \to 0 \leq y \to \lvert x \cdot y \rvert = -(x \cdot y)$
\item $\varphi_{56} :\equiv \forall z\,\forall w\,(z = w \to 0 = z \to 0 = w) \to \forall w\,(x \cdot -y = w \to 0 = x \cdot -y \to 0 = w)$
\item $\varphi_{57} :\equiv x \in \mathbb{R} \to y \in \mathbb{R} \to 0 \leq x \to \neg 0 \leq y \to -(-(x \cdot y)) \leq -0$
\item $\varphi_{58} :\equiv x \in \mathbb{R} \to y \in \mathbb{R} \to 0 \leq x \to \neg 0 \leq y \to -(-(x \cdot y)) = x \cdot y$
\item $\varphi_{59} :\equiv x \in \mathbb{R} \to y \in \mathbb{R} \to 0 \leq x \to \neg 0 \leq y \to \lvert x \cdot y \rvert = -(x \cdot y)$
\item $\varphi_{60} :\equiv -(x \cdot y) \in \mathbb{R} \to 0 \leq -(x \cdot y) \to -(-(x \cdot y)) \leq -0$
\item $\varphi_{61} :\equiv \forall v_{343}\,(v_{342} \in \mathbb{R} \to v_{343} \in \mathbb{R} \to \neg v_{343} \leq v_{342} \to v_{342} < v_{343})$
\item $\varphi_{62} :\equiv \forall v_{348}\,(v_{347} \in \mathbb{R} \to v_{348} \in \mathbb{R} \to \neg v_{348} \leq v_{347} \to v_{347} < v_{348})$
\item $\varphi_{63} :\equiv \forall v_{311}\,(v_{310} \in \mathbb{R} \to v_{311} \in \mathbb{R} \to \neg v_{311} \leq v_{310} \to v_{310} < v_{311})$
\item $\varphi_{64} :\equiv \forall v_{316}\,(v_{315} \in \mathbb{R} \to v_{316} \in \mathbb{R} \to \neg v_{316} \leq v_{315} \to v_{315} < v_{316})$
\item $\varphi_{65} :\equiv \forall w\,(-(-(x \cdot y)) = w \to -(-(x \cdot y)) < -0 \to w < -0)$
\item $\varphi_{66} :\equiv z = w \to \lvert x \rvert \cdot \lvert y \rvert = z \cdot \lvert y \rvert \to \lvert x \rvert \cdot \lvert y \rvert = w \cdot \lvert y \rvert$
\item $\varphi_{67} :\equiv -(-(x \cdot y)) = x \cdot y \to -(-(x \cdot y)) = -0 \to x \cdot y = -0$
\item $\varphi_{68} :\equiv \forall z\,(v_{358} = v_{359} \to v_{359} = z \to v_{358} = z) \to v_{358} = v_{359} \to v_{359} = v_{360} \to v_{358} = v_{360}$
\item $\varphi_{69} :\equiv \forall z\,(v_{344} = v_{345} \to v_{345} = z \to v_{344} = z) \to v_{344} = v_{345} \to v_{345} = v_{346} \to v_{344} = v_{346}$
\item $\varphi_{70} :\equiv \forall z\,(v_{353} = v_{354} \to v_{354} = z \to v_{353} = z) \to v_{353} = v_{354} \to v_{354} = v_{355} \to v_{353} = v_{355}$
\item $\varphi_{71} :\equiv \forall v_{350}\,(y \in \mathbb{R} \to v_{350} \in \mathbb{R} \to y < v_{350} \to -v_{350} < -y)$
\item $\varphi_{72} :\equiv \forall z\,(v_{339} = v_{340} \to v_{340} = z \to v_{339} = z) \to v_{339} = v_{340} \to v_{340} = v_{341} \to v_{339} = v_{341}$
\item $\varphi_{73} :\equiv \forall z\,(v_{329} = v_{330} \to v_{330} = z \to v_{329} = z) \to v_{329} = v_{330} \to v_{330} = v_{331} \to v_{329} = v_{331}$
\item $\varphi_{74} :\equiv \forall z\,(v_{334} = v_{335} \to v_{335} = z \to v_{334} = z) \to v_{334} = v_{335} \to v_{335} = v_{336} \to v_{334} = v_{336}$
\item $\varphi_{75} :\equiv \forall z\,(v_{326} = v_{327} \to v_{327} = z \to v_{326} = z) \to v_{326} = v_{327} \to v_{327} = v_{328} \to v_{326} = v_{328}$
\item $\varphi_{76} :\equiv \forall z\,(v_{312} = v_{313} \to v_{313} = z \to v_{312} = z) \to v_{312} = v_{313} \to v_{313} = v_{314} \to v_{312} = v_{314}$
\item $\varphi_{77} :\equiv \forall z\,(v_{321} = v_{322} \to v_{322} = z \to v_{321} = z) \to v_{321} = v_{322} \to v_{322} = v_{323} \to v_{321} = v_{323}$
\item $\varphi_{78} :\equiv \forall v_{318}\,(y \in \mathbb{R} \to v_{318} \in \mathbb{R} \to y < v_{318} \to -v_{318} < -y)$
\item $\varphi_{79} :\equiv \forall z\,(v_{307} = v_{308} \to v_{308} = z \to v_{307} = z) \to v_{307} = v_{308} \to v_{308} = v_{309} \to v_{307} = v_{309}$
\item $\varphi_{80} :\equiv \forall z\,(v_{302} = v_{303} \to v_{303} = z \to v_{302} = z) \to v_{302} = v_{303} \to v_{303} = v_{304} \to v_{302} = v_{304}$
\item $\varphi_{81} :\equiv x \in \mathbb{R} \to y \in \mathbb{R} \to 0 \leq x \to 0 \leq y \to \lvert x \rvert \cdot \lvert y \rvert = x \cdot \lvert y \rvert$
\item $\varphi_{82} :\equiv x \in \mathbb{R} \to y \in \mathbb{R} \to 0 \leq x \to 0 \leq y \to \lvert x \cdot y \rvert = \lvert x \rvert \cdot \lvert y \rvert$
\item $\varphi_{83} :\equiv \forall v_{341}\,(\lvert x \cdot y \rvert = -x \cdot y \to -x \cdot y = v_{341} \to \lvert x \cdot y \rvert = v_{341})$
\item $\varphi_{84} :\equiv x \in \mathbb{R} \to y \in \mathbb{R} \to \neg 0 \leq x \to 0 \leq y \to \lvert x \rvert \cdot \lvert y \rvert = -x \cdot y$
\item $\varphi_{85} :\equiv x \in \mathbb{R} \to y \in \mathbb{R} \to \neg 0 \leq x \to 0 \leq y \to -x \cdot y = \lvert x \rvert \cdot \lvert y \rvert$
\item $\varphi_{86} :\equiv \forall v_{336}\,(\lvert x \cdot y \rvert = -(x \cdot y) \to -(x \cdot y) = v_{336} \to \lvert x \cdot y \rvert = v_{336})$
\item $\varphi_{87} :\equiv x \in \mathbb{R} \to y \in \mathbb{R} \to \neg 0 \leq x \to 0 \leq y \to \lvert x \cdot y \rvert = \lvert x \rvert \cdot \lvert y \rvert$
\item $\varphi_{88} :\equiv \forall v_{328}\,(\lvert x \cdot y \rvert = x \cdot -y \to x \cdot -y = v_{328} \to \lvert x \cdot y \rvert = v_{328})$
\item $\varphi_{89} :\equiv x \in \mathbb{R} \to y \in \mathbb{R} \to 0 \leq x \to \neg 0 \leq y \to \lvert x \rvert \cdot \lvert y \rvert = x \cdot \lvert y \rvert$
\item $\varphi_{90} :\equiv x \in \mathbb{R} \to y \in \mathbb{R} \to 0 \leq x \to \neg 0 \leq y \to \lvert x \rvert \cdot \lvert y \rvert = x \cdot -y$
\item $\varphi_{91} :\equiv x \in \mathbb{R} \to y \in \mathbb{R} \to 0 \leq x \to \neg 0 \leq y \to x \cdot -y = \lvert x \rvert \cdot \lvert y \rvert$
\item $\varphi_{92} :\equiv \forall v_{323}\,(\lvert x \cdot y \rvert = -(x \cdot y) \to -(x \cdot y) = v_{323} \to \lvert x \cdot y \rvert = v_{323})$
\item $\varphi_{93} :\equiv x \in \mathbb{R} \to y \in \mathbb{R} \to 0 \leq x \to \neg 0 \leq y \to \lvert x \cdot y \rvert = \lvert x \rvert \cdot \lvert y \rvert$
\item $\varphi_{94} :\equiv \forall w\,(z = w \to -x \cdot \lvert y \rvert = -x \cdot z \to -x \cdot \lvert y \rvert = -x \cdot w)$
\item $\varphi_{95} :\equiv x \in \mathbb{R} \to y \in \mathbb{R} \to \neg 0 \leq x \to \neg 0 \leq y \to \lvert x \cdot y \rvert = -x \cdot -y$
\item $\varphi_{96} :\equiv x \in \mathbb{R} \to y \in \mathbb{R} \to \neg 0 \leq x \to \neg 0 \leq y \to \lvert x \cdot y \rvert = \lvert x \rvert \cdot \lvert y \rvert$
\item $\varphi_{97} :\equiv \forall v_{345}\,\forall v_{346}\,(\lvert x \rvert \cdot \lvert y \rvert = v_{345} \to v_{345} = v_{346} \to \lvert x \rvert \cdot \lvert y \rvert = v_{346})$
\item $\varphi_{98} :\equiv v_{349} \in \mathbb{R} \to v_{350} \in \mathbb{R} \to v_{349} < v_{350} \to -v_{350} < -v_{349}$
\item $\varphi_{99} :\equiv \forall v_{330}\,\forall v_{331}\,(\lvert x \rvert \cdot \lvert y \rvert = v_{330} \to v_{330} = v_{331} \to \lvert x \rvert \cdot \lvert y \rvert = v_{331})$
\item $\varphi_{100} :\equiv \forall v_{313}\,\forall v_{314}\,(\lvert x \rvert \cdot \lvert y \rvert = v_{313} \to v_{313} = v_{314} \to \lvert x \rvert \cdot \lvert y \rvert = v_{314})$
\item $\varphi_{101} :\equiv v_{317} \in \mathbb{R} \to v_{318} \in \mathbb{R} \to v_{317} < v_{318} \to -v_{318} < -v_{317}$
\item $\varphi_{102} :\equiv \forall v_{303}\,\forall v_{304}\,(\lvert x \rvert \cdot \lvert y \rvert = v_{303} \to v_{303} = v_{304} \to \lvert x \rvert \cdot \lvert y \rvert = v_{304})$
\item $\varphi_{103} :\equiv -(-(x \cdot y)) = x \cdot y \to -(-(x \cdot y)) < -0 \to x \cdot y < -0$
\item $\varphi_{104} :\equiv \forall z\,\forall w\,(z = w \to 0 = z \to 0 = w) \to \forall w\,(-x \cdot -y = w \to 0 = -x \cdot -y \to 0 = w)$
\item $\varphi_{105} :\equiv x \in \mathbb{R} \to y \in \mathbb{R} \to \neg 0 \leq x \to 0 \leq y \to \lvert x \rvert \cdot \lvert y \rvert = -x \cdot \lvert y \rvert$
\item $\varphi_{106} :\equiv \forall z\,\forall w\,(z = w \to x \cdot y = z \to x \cdot y = w) \to \forall w\,(-0 = w \to x \cdot y = -0 \to x \cdot y = w)$
\item $\varphi_{107} :\equiv x \in \mathbb{R} \to y \in \mathbb{R} \to \neg 0 \leq x \to \neg 0 \leq y \to -x \cdot \lvert y \rvert = -x \cdot -y$
\item $\varphi_{108} :\equiv x \in \mathbb{R} \to y \in \mathbb{R} \to \neg 0 \leq x \to \neg 0 \leq y \to \lvert x \rvert \cdot \lvert y \rvert = -x \cdot \lvert y \rvert$
\item $\varphi_{109} :\equiv x \in \mathbb{R} \to y \in \mathbb{R} \to \neg 0 \leq x \to \neg 0 \leq y \to \lvert x \rvert \cdot \lvert y \rvert = -x \cdot -y$
\item $\varphi_{110} :\equiv x \in \mathbb{R} \to y \in \mathbb{R} \to \neg 0 \leq x \to \neg 0 \leq y \to -x \cdot -y = \lvert x \rvert \cdot \lvert y \rvert$
\item $\varphi_{111} :\equiv \forall x \in \mathbb{R}\,(-(-x) = x) \to x \cdot y \in \mathbb{R} \to -(-(x \cdot y)) = x \cdot y$
\item $\varphi_{112} :\equiv \forall z\,\forall w\,(z = w \to 0 < z \to 0 < w) \to \forall w\,(-x \cdot y = w \to 0 < -x \cdot y \to 0 < w)$
\item $\varphi_{113} :\equiv \forall z\,\forall w\,(z = w \to 0 < z \to 0 < w) \to \forall w\,(x \cdot -y = w \to 0 < x \cdot -y \to 0 < w)$
\item $\varphi_{114} :\equiv \lvert x \cdot y \rvert = x \cdot y \to x \cdot y = -x \cdot -y \to \lvert x \cdot y \rvert = -x \cdot -y$
\item $\varphi_{115} :\equiv \lvert y \rvert = y \to -x \cdot \lvert y \rvert = -x \cdot \lvert y \rvert \to -x \cdot \lvert y \rvert = -x \cdot y$
\item $\varphi_{116} :\equiv \lvert x \cdot y \rvert = -(x \cdot y) \to -(x \cdot y) = -x \cdot y \to \lvert x \cdot y \rvert = -x \cdot y$
\item $\varphi_{117} :\equiv \lvert x \rvert = x \to \lvert x \rvert \cdot \lvert y \rvert = \lvert x \rvert \cdot \lvert y \rvert \to \lvert x \rvert \cdot \lvert y \rvert = x \cdot \lvert y \rvert$
\item $\varphi_{118} :\equiv \lvert x \cdot y \rvert = -(x \cdot y) \to -(x \cdot y) = x \cdot -y \to \lvert x \cdot y \rvert = x \cdot -y$
\item $\varphi_{119} :\equiv \lvert x \cdot y \rvert = x \cdot y \to x \cdot y = \lvert x \rvert \cdot \lvert y \rvert \to \lvert x \cdot y \rvert = \lvert x \rvert \cdot \lvert y \rvert$
\item $\varphi_{120} :\equiv \lvert x \rvert \cdot \lvert y \rvert = x \cdot \lvert y \rvert \to x \cdot \lvert y \rvert = x \cdot y \to \lvert x \rvert \cdot \lvert y \rvert = x \cdot y$
\item $\varphi_{121} :\equiv \forall w\,(-0 = w \to x \cdot y = -0 \to x \cdot y = w) \to -0 = 0 \to x \cdot y = -0 \to x \cdot y = 0$
\item $\varphi_{122} :\equiv \forall v_{360}\,(\lvert x \cdot y \rvert = -x \cdot -y \to -x \cdot -y = v_{360} \to \lvert x \cdot y \rvert = v_{360})$
\item $\varphi_{123} :\equiv \forall v_{332}\,\forall v_{333}\,(v_{332} = v_{333} \to v_{333} = v_{332}) \to \forall v_{333}\,(-x \cdot y = v_{333} \to v_{333} = -x \cdot y)$
\item $\varphi_{124} :\equiv \forall v_{314}\,(\lvert x \rvert \cdot \lvert y \rvert = x \cdot \lvert y \rvert \to x \cdot \lvert y \rvert = v_{314} \to \lvert x \rvert \cdot \lvert y \rvert = v_{314})$
\item $\varphi_{125} :\equiv \forall v_{319}\,\forall v_{320}\,(v_{319} = v_{320} \to v_{320} = v_{319}) \to \forall v_{320}\,(x \cdot -y = v_{320} \to v_{320} = x \cdot -y)$
\item $\varphi_{126} :\equiv \forall v_{304}\,(\lvert x \rvert \cdot \lvert y \rvert = x \cdot \lvert y \rvert \to x \cdot \lvert y \rvert = v_{304} \to \lvert x \rvert \cdot \lvert y \rvert = v_{304})$
\item $\varphi_{127} :\equiv \forall w\,(\lvert y \rvert = w \to -x \cdot \lvert y \rvert = -x \cdot \lvert y \rvert \to -x \cdot \lvert y \rvert = -x \cdot w)$
\item $\varphi_{128} :\equiv \forall w\,(\lvert x \rvert = w \to \lvert x \rvert \cdot \lvert y \rvert = \lvert x \rvert \cdot \lvert y \rvert \to \lvert x \rvert \cdot \lvert y \rvert = w \cdot \lvert y \rvert)$
\item $\varphi_{129} :\equiv \forall y\,(x \in \mathbb{R} \to y \in \mathbb{R} \to \neg y \leq x \to x < y) \to x \in \mathbb{R} \to 0 \in \mathbb{R} \to \neg 0 \leq x \to x < 0$
\item $\varphi_{130} :\equiv \forall z\,\forall w\,(z = w \to 0 < z \to 0 < w) \to \forall w\,(-x \cdot -y = w \to 0 < -x \cdot -y \to 0 < w)$
\item $\varphi_{131} :\equiv \forall z\,\forall w\,(z = w \to z < -y \to w < -y) \to \forall w\,(-0 = w \to -0 < -y \to w < -y)$
\item $\varphi_{132} :\equiv \forall z\,\forall w\,(z = w \to z < -x \to w < -x) \to \forall w\,(-0 = w \to -0 < -x \to w < -x)$
\item $\varphi_{133} :\equiv \forall z\,\forall w\,(z = w \to x \cdot y < z \to x \cdot y < w) \to \forall w\,(-0 = w \to x \cdot y < -0 \to x \cdot y < w)$
\item $\varphi_{134} :\equiv (x \in \mathbb{R} \to y \in \mathbb{R} \to 0 \leq x \to 0 \leq x) \to x \in \mathbb{R} \to y \in \mathbb{R} \to 0 \leq x \to 0 \leq y \to \lvert x \rvert = x$
\item $\varphi_{135} :\equiv \lvert y \rvert = -y \to -x \cdot \lvert y \rvert = -x \cdot \lvert y \rvert \to -x \cdot \lvert y \rvert = -x \cdot -y$
\item $\varphi_{136} :\equiv \lvert x \rvert = -x \to \lvert x \rvert \cdot \lvert y \rvert = \lvert x \rvert \cdot \lvert y \rvert \to \lvert x \rvert \cdot \lvert y \rvert = -x \cdot \lvert y \rvert$
\item $\varphi_{137} :\equiv \lvert x \cdot y \rvert = -x \cdot y \to -x \cdot y = \lvert x \rvert \cdot \lvert y \rvert \to \lvert x \cdot y \rvert = \lvert x \rvert \cdot \lvert y \rvert$
\item $\varphi_{138} :\equiv (x \in \mathbb{R} \to y \in \mathbb{R} \to \neg 0 \leq x \to \neg 0 \leq x) \to x \in \mathbb{R} \to y \in \mathbb{R} \to \neg 0 \leq x \to 0 \leq y \to x < 0$
\item $\varphi_{139} :\equiv (-0 = 0 \to x \cdot y = -0 \to x \cdot y = 0) \to x \in \mathbb{R} \to y \in \mathbb{R} \to \neg 0 \leq x \to 0 \leq y \to x \cdot y \leq 0$
\item $\varphi_{140} :\equiv \lvert x \cdot y \rvert = x \cdot -y \to x \cdot -y = \lvert x \rvert \cdot \lvert y \rvert \to \lvert x \cdot y \rvert = \lvert x \rvert \cdot \lvert y \rvert$
\item $\varphi_{141} :\equiv \lvert x \rvert \cdot \lvert y \rvert = x \cdot \lvert y \rvert \to x \cdot \lvert y \rvert = x \cdot -y \to \lvert x \rvert \cdot \lvert y \rvert = x \cdot -y$
\item $\varphi_{142} :\equiv (x \in \mathbb{R} \to y \in \mathbb{R} \to 0 \leq x \to 0 \leq x) \to x \in \mathbb{R} \to y \in \mathbb{R} \to 0 \leq x \to \neg 0 \leq y \to \lvert x \rvert = x$
\item $\varphi_{143} :\equiv (-0 = 0 \to x \cdot y = -0 \to x \cdot y = 0) \to x \in \mathbb{R} \to y \in \mathbb{R} \to 0 \leq x \to \neg 0 \leq y \to x \cdot y \leq 0$
\item $\varphi_{144} :\equiv (x \in \mathbb{R} \to y \in \mathbb{R} \to \neg 0 \leq x \to \neg 0 \leq x) \to x \in \mathbb{R} \to y \in \mathbb{R} \to \neg 0 \leq x \to \neg 0 \leq y \to x < 0$
\item $\varphi_{145} :\equiv \forall v_{356}\,\forall v_{357}\,(v_{356} = v_{357} \to v_{357} = v_{356}) \to \forall v_{357}\,(\lvert x \rvert \cdot \lvert y \rvert = v_{357} \to v_{357} = \lvert x \rvert \cdot \lvert y \rvert)$
\item $\varphi_{146} :\equiv \forall v_{346}\,(\lvert x \rvert \cdot \lvert y \rvert = -x \cdot \lvert y \rvert \to -x \cdot \lvert y \rvert = v_{346} \to \lvert x \rvert \cdot \lvert y \rvert = v_{346})$
\item $\varphi_{147} :\equiv \forall v_{351}\,\forall v_{352}\,(v_{351} = v_{352} \to v_{352} = v_{351}) \to \forall v_{352}\,(-x \cdot -y = v_{352} \to v_{352} = -x \cdot -y)$
\item $\varphi_{148} :\equiv \forall v_{337}\,\forall v_{338}\,(v_{337} = v_{338} \to v_{338} = v_{337}) \to \forall v_{338}\,(\lvert x \rvert \cdot \lvert y \rvert = v_{338} \to v_{338} = \lvert x \rvert \cdot \lvert y \rvert)$
\item $\varphi_{149} :\equiv \forall v_{331}\,(\lvert x \rvert \cdot \lvert y \rvert = -x \cdot \lvert y \rvert \to -x \cdot \lvert y \rvert = v_{331} \to \lvert x \rvert \cdot \lvert y \rvert = v_{331})$
\item $\varphi_{150} :\equiv \forall v_{324}\,\forall v_{325}\,(v_{324} = v_{325} \to v_{325} = v_{324}) \to \forall v_{325}\,(\lvert x \rvert \cdot \lvert y \rvert = v_{325} \to v_{325} = \lvert x \rvert \cdot \lvert y \rvert)$
\item $\varphi_{151} :\equiv \forall v_{305}\,\forall v_{306}\,(v_{305} = v_{306} \to v_{306} = v_{305}) \to \forall v_{306}\,(\lvert x \rvert \cdot \lvert y \rvert = v_{306} \to v_{306} = \lvert x \rvert \cdot \lvert y \rvert)$
\item $\varphi_{152} :\equiv \forall w\,(-0 = w \to -0 < -y \to w < -y) \to -0 = 0 \to -0 < -y \to 0 < -y$
\item $\varphi_{153} :\equiv \forall w\,(-0 = w \to -0 < -x \to w < -x) \to -0 = 0 \to -0 < -x \to 0 < -x$
\item $\varphi_{154} :\equiv \forall w\,(-0 = w \to x \cdot y < -0 \to x \cdot y < w) \to -0 = 0 \to x \cdot y < -0 \to x \cdot y < 0$
\item $\varphi_{155} :\equiv \forall y\,(x \in \mathbb{R} \to y \in \mathbb{R} \to x \cdot y \in \mathbb{R}) \to x \in \mathbb{R} \to y \in \mathbb{R} \to x \cdot y \in \mathbb{R}$
\item $\varphi_{156} :\equiv \forall x\,\forall y\,(x \in \mathbb{R} \to y \in \mathbb{R} \to \neg y \leq x \to x < y) \to \forall y\,(x \in \mathbb{R} \to y \in \mathbb{R} \to \neg y \leq x \to x < y)$
\item $\varphi_{157} :\equiv \forall x \in \mathbb{R}\,(x \leq 0 \to \lvert x \rvert = -x) \to x \cdot y \in \mathbb{R} \to x \cdot y \leq 0 \to \lvert x \cdot y \rvert = -(x \cdot y)$
\item $\varphi_{158} :\equiv (x \in \mathbb{R} \to y \in \mathbb{R} \to 0 \leq x \to 0 \leq y \to 0 \leq y) \to x \in \mathbb{R} \to y \in \mathbb{R} \to 0 \leq x \to 0 \leq y \to \lvert y \rvert = y$
\item $\varphi_{159} :\equiv (x \in \mathbb{R} \to y \in \mathbb{R} \to 0 \leq x \to \neg 0 \leq y \to \neg 0 \leq y) \to x \in \mathbb{R} \to y \in \mathbb{R} \to 0 \leq x \to \neg 0 \leq y \to y < 0$
\item $\varphi_{160} :\equiv (x \in \mathbb{R} \to y \in \mathbb{R} \to \neg 0 \leq x \to 0 \leq y \to 0 \leq y) \to x \in \mathbb{R} \to y \in \mathbb{R} \to \neg 0 \leq x \to 0 \leq y \to \lvert y \rvert = y$
\item $\varphi_{161} :\equiv (x \in \mathbb{R} \to y \in \mathbb{R} \to \neg 0 \leq x \to \neg 0 \leq y \to \neg 0 \leq y) \to x \in \mathbb{R} \to y \in \mathbb{R} \to \neg 0 \leq x \to \neg 0 \leq y \to y < 0$
\item $\varphi_{162} :\equiv \forall v_{333}\,(-x \cdot y = v_{333} \to v_{333} = -x \cdot y) \to -x \cdot y = -(x \cdot y) \to -(x \cdot y) = -x \cdot y$
\item $\varphi_{163} :\equiv \forall v_{320}\,(x \cdot -y = v_{320} \to v_{320} = x \cdot -y) \to x \cdot -y = -(x \cdot y) \to -(x \cdot y) = x \cdot -y$
\item $\varphi_{164} :\equiv (-0 = 0 \to -0 < -x \to 0 < -x) \to x \in \mathbb{R} \to y \in \mathbb{R} \to \neg 0 \leq x \to 0 \leq y \to 0 < -x$
\item $\varphi_{165} :\equiv \forall x\,\forall y\,(x \in \mathbb{R} \to y \in \mathbb{R} \to x \cdot y \in \mathbb{R}) \to \forall y\,(x \in \mathbb{R} \to y \in \mathbb{R} \to x \cdot y \in \mathbb{R})$
\item $\varphi_{166} :\equiv (-0 = 0 \to -0 < -y \to 0 < -y) \to x \in \mathbb{R} \to y \in \mathbb{R} \to 0 \leq x \to \neg 0 \leq y \to 0 < -y$
\item $\varphi_{167} :\equiv \lvert x \cdot y \rvert = -x \cdot -y \to -x \cdot -y = \lvert x \rvert \cdot \lvert y \rvert \to \lvert x \cdot y \rvert = \lvert x \rvert \cdot \lvert y \rvert$
\item $\varphi_{168} :\equiv (-0 = 0 \to -0 < -y \to 0 < -y) \to x \in \mathbb{R} \to y \in \mathbb{R} \to \neg 0 \leq x \to \neg 0 \leq y \to 0 < -y$
\item $\varphi_{169} :\equiv (-0 = 0 \to -0 < -x \to 0 < -x) \to x \in \mathbb{R} \to y \in \mathbb{R} \to \neg 0 \leq x \to \neg 0 \leq y \to 0 < -x$
\item $\varphi_{170} :\equiv \lvert x \rvert \cdot \lvert y \rvert = -x \cdot \lvert y \rvert \to -x \cdot \lvert y \rvert = -x \cdot y \to \lvert x \rvert \cdot \lvert y \rvert = -x \cdot y$
\item $\varphi_{171} :\equiv \forall v_{343}\,(y \in \mathbb{R} \to v_{343} \in \mathbb{R} \to \neg v_{343} \leq y \to y < v_{343}) \to y \in \mathbb{R} \to 0 \in \mathbb{R} \to \neg 0 \leq y \to y < 0$
\item $\varphi_{172} :\equiv \forall v_{348}\,(y \in \mathbb{R} \to v_{348} \in \mathbb{R} \to \neg v_{348} \leq y \to y < v_{348}) \to y \in \mathbb{R} \to 0 \in \mathbb{R} \to \neg 0 \leq y \to y < 0$
\item $\varphi_{173} :\equiv \forall w\,(-x \cdot y = w \to 0 = -x \cdot y \to 0 = w) \to -x \cdot y = -(x \cdot y) \to 0 = -x \cdot y \to 0 = -(x \cdot y)$
\item $\varphi_{174} :\equiv \forall v_{311}\,(y \in \mathbb{R} \to v_{311} \in \mathbb{R} \to \neg v_{311} \leq y \to y < v_{311}) \to y \in \mathbb{R} \to 0 \in \mathbb{R} \to \neg 0 \leq y \to y < 0$
\item $\varphi_{175} :\equiv \forall w\,(x \cdot -y = w \to 0 = x \cdot -y \to 0 = w) \to x \cdot -y = -(x \cdot y) \to 0 = x \cdot -y \to 0 = -(x \cdot y)$
\item $\varphi_{176} :\equiv \forall v_{316}\,(y \in \mathbb{R} \to v_{316} \in \mathbb{R} \to \neg v_{316} \leq y \to y < v_{316}) \to y \in \mathbb{R} \to 0 \in \mathbb{R} \to \neg 0 \leq y \to y < 0$
\item $\varphi_{177} :\equiv (x \in \mathbb{R} \to y \in \mathbb{R} \to \neg 0 \leq x \to 0 \leq y \to 0 \leq y) \to x \in \mathbb{R} \to y \in \mathbb{R} \to \neg 0 \leq x \to 0 \leq y \to 0 \leq -x \cdot y$
\item $\varphi_{178} :\equiv \forall y\,(x \in \mathbb{R} \to y \in \mathbb{R} \to 0 \leq x \to 0 \leq y \to 0 \leq x \cdot y) \to x \in \mathbb{R} \to y \in \mathbb{R} \to 0 \leq x \to 0 \leq y \to 0 \leq x \cdot y$
\item $\varphi_{179} :\equiv \forall x\,\forall y\,(x \in \mathbb{R} \to y \in \mathbb{R} \to \neg y \leq x \to x < y) \to \forall y\,(v_{342} \in \mathbb{R} \to y \in \mathbb{R} \to \neg y \leq v_{342} \to v_{342} < y)$
\item $\varphi_{180} :\equiv \forall x\,\forall y\,(x \in \mathbb{R} \to y \in \mathbb{R} \to \neg y \leq x \to x < y) \to \forall y\,(v_{347} \in \mathbb{R} \to y \in \mathbb{R} \to \neg y \leq v_{347} \to v_{347} < y)$
\item $\varphi_{181} :\equiv \forall x\,\forall y\,(x \in \mathbb{R} \to y \in \mathbb{R} \to \neg y \leq x \to x < y) \to \forall y\,(v_{310} \in \mathbb{R} \to y \in \mathbb{R} \to \neg y \leq v_{310} \to v_{310} < y)$
\item $\varphi_{182} :\equiv \forall x\,\forall y\,(x \in \mathbb{R} \to y \in \mathbb{R} \to \neg y \leq x \to x < y) \to \forall y\,(v_{315} \in \mathbb{R} \to y \in \mathbb{R} \to \neg y \leq v_{315} \to v_{315} < y)$
\item $\varphi_{183} :\equiv (x \in \mathbb{R} \to y \in \mathbb{R} \to \neg 0 \leq x \to 0 \leq y \to x < 0) \to x \in \mathbb{R} \to y \in \mathbb{R} \to \neg 0 \leq x \to 0 \leq y \to \lvert x \rvert = -x$
\item $\varphi_{184} :\equiv (x \in \mathbb{R} \to y \in \mathbb{R} \to \neg 0 \leq x \to 0 \leq y \to 0 < -x) \to x \in \mathbb{R} \to y \in \mathbb{R} \to \neg 0 \leq x \to 0 \leq y \to 0 \leq -x$
\item $\varphi_{185} :\equiv (x \in \mathbb{R} \to y \in \mathbb{R} \to 0 \leq x \to \neg 0 \leq y \to y < 0) \to x \in \mathbb{R} \to y \in \mathbb{R} \to 0 \leq x \to \neg 0 \leq y \to \lvert y \rvert = -y$
\item $\varphi_{186} :\equiv (x \in \mathbb{R} \to y \in \mathbb{R} \to 0 \leq x \to \neg 0 \leq y \to 0 < -y) \to x \in \mathbb{R} \to y \in \mathbb{R} \to 0 \leq x \to \neg 0 \leq y \to 0 \leq -y$
\item $\varphi_{187} :\equiv (x \in \mathbb{R} \to y \in \mathbb{R} \to \neg 0 \leq x \to \neg 0 \leq y \to y < 0) \to x \in \mathbb{R} \to y \in \mathbb{R} \to \neg 0 \leq x \to \neg 0 \leq y \to \lvert y \rvert = -y$
\item $\varphi_{188} :\equiv (x \in \mathbb{R} \to y \in \mathbb{R} \to \neg 0 \leq x \to \neg 0 \leq y \to x < 0) \to x \in \mathbb{R} \to y \in \mathbb{R} \to \neg 0 \leq x \to \neg 0 \leq y \to \lvert x \rvert = -x$
\item $\varphi_{189} :\equiv (x \in \mathbb{R} \to y \in \mathbb{R} \to \neg 0 \leq x \to \neg 0 \leq y \to 0 < -y) \to x \in \mathbb{R} \to y \in \mathbb{R} \to \neg 0 \leq x \to \neg 0 \leq y \to 0 \leq -y$
\item $\varphi_{190} :\equiv (x \in \mathbb{R} \to y \in \mathbb{R} \to \neg 0 \leq x \to \neg 0 \leq y \to 0 < -x) \to x \in \mathbb{R} \to y \in \mathbb{R} \to \neg 0 \leq x \to \neg 0 \leq y \to 0 \leq -x$
\item $\varphi_{191} :\equiv \forall x\,\forall y\,(x \in \mathbb{R} \to y \in \mathbb{R} \to x \leq y \to -y \leq -x) \to \forall y\,(0 \in \mathbb{R} \to y \in \mathbb{R} \to 0 \leq y \to -y \leq -0)$
\item $\varphi_{192} :\equiv (x \in \mathbb{R} \to y \in \mathbb{R} \to 0 \leq x \to \neg 0 \leq y \to 0 \leq -y) \to x \in \mathbb{R} \to y \in \mathbb{R} \to 0 \leq x \to \neg 0 \leq y \to 0 \leq x \cdot -y$
\item $\varphi_{193} :\equiv \forall v_{352}\,(-x \cdot -y = v_{352} \to v_{352} = -x \cdot -y) \to -x \cdot -y = x \cdot y \to x \cdot y = -x \cdot -y$
\item $\varphi_{194} :\equiv (x \in \mathbb{R} \to y \in \mathbb{R} \to \neg 0 \leq x \to 0 \leq y \to x < 0) \to x \in \mathbb{R} \to y \in \mathbb{R} \to \neg 0 \leq x \to 0 \leq y \to -0 < -x$
\item $\varphi_{195} :\equiv (x \in \mathbb{R} \to y \in \mathbb{R} \to 0 \leq x \to \neg 0 \leq y \to y < 0) \to x \in \mathbb{R} \to y \in \mathbb{R} \to 0 \leq x \to \neg 0 \leq y \to -0 < -y$
\item $\varphi_{196} :\equiv \forall v_{306}\,(\lvert x \rvert \cdot \lvert y \rvert = v_{306} \to v_{306} = \lvert x \rvert \cdot \lvert y \rvert) \to \lvert x \rvert \cdot \lvert y \rvert = x \cdot y \to x \cdot y = \lvert x \rvert \cdot \lvert y \rvert$
\item $\varphi_{197} :\equiv \forall y\,(x \in \mathbb{R} \to y \in \mathbb{R} \to -x \cdot -y = x \cdot y) \to x \in \mathbb{R} \to y \in \mathbb{R} \to -x \cdot -y = x \cdot y$
\item $\varphi_{198} :\equiv \forall y\,(x \in \mathbb{R} \to y \in \mathbb{R} \to -x \cdot y = -(x \cdot y)) \to x \in \mathbb{R} \to y \in \mathbb{R} \to -x \cdot y = -(x \cdot y)$
\item $\varphi_{199} :\equiv \forall y\,(x \in \mathbb{R} \to y \in \mathbb{R} \to x \cdot -y = -(x \cdot y)) \to x \in \mathbb{R} \to y \in \mathbb{R} \to x \cdot -y = -(x \cdot y)$
\item $\varphi_{200} :\equiv \forall x\,\forall y\,(x \in \mathbb{R} \to y \in \mathbb{R} \to 0 \leq x \to 0 \leq y \to 0 \leq x \cdot y) \to \forall y\,(x \in \mathbb{R} \to y \in \mathbb{R} \to 0 \leq x \to 0 \leq y \to 0 \leq x \cdot y)$
\item $\varphi_{201} :\equiv \lvert x \rvert \cdot \lvert y \rvert = -x \cdot \lvert y \rvert \to -x \cdot \lvert y \rvert = -x \cdot -y \to \lvert x \rvert \cdot \lvert y \rvert = -x \cdot -y$
\item $\varphi_{202} :\equiv (x \in \mathbb{R} \to y \in \mathbb{R} \to \neg 0 \leq x \to \neg 0 \leq y \to y < 0) \to x \in \mathbb{R} \to y \in \mathbb{R} \to \neg 0 \leq x \to \neg 0 \leq y \to -0 < -y$
\item $\varphi_{203} :\equiv (x \in \mathbb{R} \to y \in \mathbb{R} \to \neg 0 \leq x \to \neg 0 \leq y \to x < 0) \to x \in \mathbb{R} \to y \in \mathbb{R} \to \neg 0 \leq x \to \neg 0 \leq y \to -0 < -x$
\item $\varphi_{204} :\equiv (x \in \mathbb{R} \to y \in \mathbb{R} \to 0 \leq x \to 0 \leq y \to 0 \leq x \cdot y) \to x \in \mathbb{R} \to y \in \mathbb{R} \to 0 \leq x \to 0 \leq y \to \lvert x \cdot y \rvert = x \cdot y$
\item $\varphi_{205} :\equiv \forall w\,(-x \cdot -y = w \to 0 = -x \cdot -y \to 0 = w) \to -x \cdot -y = x \cdot y \to 0 = -x \cdot -y \to 0 = x \cdot y$
\item $\varphi_{206} :\equiv (-x \cdot -y = x \cdot y \to 0 = -x \cdot -y \to 0 = x \cdot y) \to x \in \mathbb{R} \to y \in \mathbb{R} \to \neg 0 \leq x \to \neg 0 \leq y \to 0 \leq x \cdot y$
\item $\varphi_{207} :\equiv (x \in \mathbb{R} \to y \in \mathbb{R} \to \neg 0 \leq x \to 0 \leq y \to x \in \mathbb{R}) \to x \in \mathbb{R} \to y \in \mathbb{R} \to \neg 0 \leq x \to 0 \leq y \to -x \in \mathbb{R}$
\item $\varphi_{208} :\equiv (x \in \mathbb{R} \to y \in \mathbb{R} \to \neg 0 \leq x \to 0 \leq y \to y \in \mathbb{R}) \to x \in \mathbb{R} \to y \in \mathbb{R} \to \neg 0 \leq x \to 0 \leq y \to x \cdot y \in \mathbb{R}$
\item $\varphi_{209} :\equiv (x \in \mathbb{R} \to y \in \mathbb{R} \to 0 \leq x \to \neg 0 \leq y \to y \in \mathbb{R}) \to x \in \mathbb{R} \to y \in \mathbb{R} \to 0 \leq x \to \neg 0 \leq y \to -y \in \mathbb{R}$
\item $\varphi_{210} :\equiv (x \in \mathbb{R} \to y \in \mathbb{R} \to 0 \leq x \to \neg 0 \leq y \to y \in \mathbb{R}) \to x \in \mathbb{R} \to y \in \mathbb{R} \to 0 \leq x \to \neg 0 \leq y \to x \cdot y \in \mathbb{R}$
\item $\varphi_{211} :\equiv \forall y\,(x \in \mathbb{R} \to y \in \mathbb{R} \to x < y \to -y < -x) \to x \in \mathbb{R} \to 0 \in \mathbb{R} \to x < 0 \to -0 < -x$
\item $\varphi_{212} :\equiv (x \in \mathbb{R} \to y \in \mathbb{R} \to \neg 0 \leq x \to \neg 0 \leq y \to y \in \mathbb{R}) \to x \in \mathbb{R} \to y \in \mathbb{R} \to \neg 0 \leq x \to \neg 0 \leq y \to -y \in \mathbb{R}$
\item $\varphi_{213} :\equiv (x \in \mathbb{R} \to y \in \mathbb{R} \to \neg 0 \leq x \to \neg 0 \leq y \to x \in \mathbb{R}) \to x \in \mathbb{R} \to y \in \mathbb{R} \to \neg 0 \leq x \to \neg 0 \leq y \to -x \in \mathbb{R}$
\item $\varphi_{214} :\equiv \forall w\,(-x \cdot y = w \to 0 < -x \cdot y \to 0 < w) \to -x \cdot y = -(x \cdot y) \to 0 < -x \cdot y \to 0 < -(x \cdot y)$
\item $\varphi_{215} :\equiv \forall w\,(x \cdot -y = w \to 0 < x \cdot -y \to 0 < w) \to x \cdot -y = -(x \cdot y) \to 0 < x \cdot -y \to 0 < -(x \cdot y)$
\item $\varphi_{216} :\equiv (x \in \mathbb{R} \to y \in \mathbb{R} \to \neg 0 \leq x \to \neg 0 \leq y \to 0 \leq -y) \to x \in \mathbb{R} \to y \in \mathbb{R} \to \neg 0 \leq x \to \neg 0 \leq y \to 0 \leq -x \cdot -y$
\item $\varphi_{217} :\equiv (x \in \mathbb{R} \to y \in \mathbb{R} \to \neg 0 \leq x \to \neg 0 \leq y \to 0 \leq x \cdot y) \to x \in \mathbb{R} \to y \in \mathbb{R} \to \neg 0 \leq x \to \neg 0 \leq y \to \lvert x \cdot y \rvert = x \cdot y$
\item $\varphi_{218} :\equiv \forall x\,\forall y\,(x \in \mathbb{R} \to y \in \mathbb{R} \to -x \cdot -y = x \cdot y) \to \forall y\,(x \in \mathbb{R} \to y \in \mathbb{R} \to -x \cdot -y = x \cdot y)$
\item $\varphi_{219} :\equiv \forall x\,\forall y\,(x \in \mathbb{R} \to y \in \mathbb{R} \to -x \cdot y = -(x \cdot y)) \to \forall y\,(x \in \mathbb{R} \to y \in \mathbb{R} \to -x \cdot y = -(x \cdot y))$
\item $\varphi_{220} :\equiv (-x \cdot y = -(x \cdot y) \to 0 = -x \cdot y \to 0 = -(x \cdot y)) \to x \in \mathbb{R} \to y \in \mathbb{R} \to \neg 0 \leq x \to 0 \leq y \to 0 \leq -(x \cdot y)$
\item $\varphi_{221} :\equiv \forall x\,\forall y\,(x \in \mathbb{R} \to y \in \mathbb{R} \to x \cdot -y = -(x \cdot y)) \to \forall y\,(x \in \mathbb{R} \to y \in \mathbb{R} \to x \cdot -y = -(x \cdot y))$
\item $\varphi_{222} :\equiv (x \cdot -y = -(x \cdot y) \to 0 = x \cdot -y \to 0 = -(x \cdot y)) \to x \in \mathbb{R} \to y \in \mathbb{R} \to 0 \leq x \to \neg 0 \leq y \to 0 \leq -(x \cdot y)$
\item $\varphi_{223} :\equiv \forall y\,(v_{342} \in \mathbb{R} \to y \in \mathbb{R} \to \neg y \leq v_{342} \to v_{342} < y) \to v_{342} \in \mathbb{R} \to v_{343} \in \mathbb{R} \to \neg v_{343} \leq v_{342} \to v_{342} < v_{343}$
\item $\varphi_{224} :\equiv \forall y\,(v_{347} \in \mathbb{R} \to y \in \mathbb{R} \to \neg y \leq v_{347} \to v_{347} < y) \to v_{347} \in \mathbb{R} \to v_{348} \in \mathbb{R} \to \neg v_{348} \leq v_{347} \to v_{347} < v_{348}$
\item $\varphi_{225} :\equiv \forall y\,(v_{310} \in \mathbb{R} \to y \in \mathbb{R} \to \neg y \leq v_{310} \to v_{310} < y) \to v_{310} \in \mathbb{R} \to v_{311} \in \mathbb{R} \to \neg v_{311} \leq v_{310} \to v_{310} < v_{311}$
\item $\varphi_{226} :\equiv \forall y\,(v_{315} \in \mathbb{R} \to y \in \mathbb{R} \to \neg y \leq v_{315} \to v_{315} < y) \to v_{315} \in \mathbb{R} \to v_{316} \in \mathbb{R} \to \neg v_{316} \leq v_{315} \to v_{315} < v_{316}$
\item $\varphi_{227} :\equiv \forall z\,\forall w\,(z = w \to x \cdot \lvert y \rvert = x \cdot z \to x \cdot \lvert y \rvert = x \cdot w) \to \varphi_{33}$
\item $\varphi_{228} :\equiv \forall v_{338}\,(\lvert x \rvert \cdot \lvert y \rvert = v_{338} \to v_{338} = \lvert x \rvert \cdot \lvert y \rvert) \to \lvert x \rvert \cdot \lvert y \rvert = -x \cdot y \to -x \cdot y = \lvert x \rvert \cdot \lvert y \rvert$
\item $\varphi_{229} :\equiv \forall v_{325}\,(\lvert x \rvert \cdot \lvert y \rvert = v_{325} \to v_{325} = \lvert x \rvert \cdot \lvert y \rvert) \to \lvert x \rvert \cdot \lvert y \rvert = x \cdot -y \to x \cdot -y = \lvert x \rvert \cdot \lvert y \rvert$
\item $\varphi_{230} :\equiv \forall y\,(x \in \mathbb{R} \to y \in \mathbb{R} \to 0 \leq x \to 0 \leq y \to 0 \leq x \cdot y) \to x \in \mathbb{R} \to -y \in \mathbb{R} \to 0 \leq x \to 0 \leq -y \to 0 \leq x \cdot -y$
\item $\varphi_{231} :\equiv (\lvert x \rvert \cdot \lvert y \rvert = x \cdot y \to x \cdot y = \lvert x \rvert \cdot \lvert y \rvert) \to x \in \mathbb{R} \to y \in \mathbb{R} \to 0 \leq x \to 0 \leq y \to x \cdot y = \lvert x \rvert \cdot \lvert y \rvert$
\item $\varphi_{232} :\equiv \forall x\,\forall y\,(x \in \mathbb{R} \to y \in \mathbb{R} \to x < y \to -y < -x) \to \forall y\,(x \in \mathbb{R} \to y \in \mathbb{R} \to x < y \to -y < -x)$
\item $\varphi_{233} :\equiv (x \in \mathbb{R} \to y \in \mathbb{R} \to \neg 0 \leq x \to 0 \leq y \to y \in \mathbb{R}) \to x \in \mathbb{R} \to y \in \mathbb{R} \to \neg 0 \leq x \to 0 \leq y \to -x \cdot y = -(x \cdot y)$
\item $\varphi_{234} :\equiv (-x \cdot y = -(x \cdot y) \to -(x \cdot y) = -x \cdot y) \to x \in \mathbb{R} \to y \in \mathbb{R} \to \neg 0 \leq x \to 0 \leq y \to -(x \cdot y) = -x \cdot y$
\item $\varphi_{235} :\equiv (x \in \mathbb{R} \to y \in \mathbb{R} \to 0 \leq x \to \neg 0 \leq y \to y \in \mathbb{R}) \to x \in \mathbb{R} \to y \in \mathbb{R} \to 0 \leq x \to \neg 0 \leq y \to x \cdot -y = -(x \cdot y)$
\item $\varphi_{236} :\equiv (x \cdot -y = -(x \cdot y) \to -(x \cdot y) = x \cdot -y) \to x \in \mathbb{R} \to y \in \mathbb{R} \to 0 \leq x \to \neg 0 \leq y \to -(x \cdot y) = x \cdot -y$
\item $\varphi_{237} :\equiv (\lvert y \rvert = y \to x \cdot \lvert y \rvert = x \cdot \lvert y \rvert \to x \cdot \lvert y \rvert = x \cdot y) \to x \in \mathbb{R} \to y \in \mathbb{R} \to 0 \leq x \to 0 \leq y \to x \cdot \lvert y \rvert = x \cdot y$
\item $\varphi_{238} :\equiv (-x \cdot -y = x \cdot y \to x \cdot y = -x \cdot -y) \to x \in \mathbb{R} \to y \in \mathbb{R} \to \neg 0 \leq x \to \neg 0 \leq y \to x \cdot y = -x \cdot -y$
\item $\varphi_{239} :\equiv (x \in \mathbb{R} \to y \in \mathbb{R} \to \neg 0 \leq x \to \neg 0 \leq y \to y \in \mathbb{R}) \to x \in \mathbb{R} \to y \in \mathbb{R} \to \neg 0 \leq x \to \neg 0 \leq y \to -x \cdot -y = x \cdot y$
\item $\varphi_{240} :\equiv (x \in \mathbb{R} \to y \in \mathbb{R} \to \neg 0 \leq x \to 0 \leq y \to x \cdot y \in \mathbb{R}) \to x \in \mathbb{R} \to y \in \mathbb{R} \to \neg 0 \leq x \to 0 \leq y \to -(x \cdot y) \in \mathbb{R}$
\item $\varphi_{241} :\equiv (x \in \mathbb{R} \to y \in \mathbb{R} \to 0 \leq x \to \neg 0 \leq y \to x \cdot y \in \mathbb{R}) \to x \in \mathbb{R} \to y \in \mathbb{R} \to 0 \leq x \to \neg 0 \leq y \to -(x \cdot y) \in \mathbb{R}$
\item $\varphi_{242} :\equiv (x \in \mathbb{R} \to y \in \mathbb{R} \to \neg 0 \leq x \to 0 \leq y \to 0 \leq -(x \cdot y)) \to \varphi_{53}$
\item $\varphi_{243} :\equiv (x \in \mathbb{R} \to y \in \mathbb{R} \to 0 \leq x \to \neg 0 \leq y \to 0 \leq -(x \cdot y)) \to \varphi_{57}$
\item $\varphi_{244} :\equiv \forall w\,(-x \cdot -y = w \to 0 < -x \cdot -y \to 0 < w) \to -x \cdot -y = x \cdot y \to 0 < -x \cdot -y \to 0 < x \cdot y$
\item $\varphi_{245} :\equiv \forall x\,\forall y\,(x \in \mathbb{R} \to y \in \mathbb{R} \to 0 \leq x \to 0 \leq y \to 0 \leq x \cdot y) \to \forall y\,(-x \in \mathbb{R} \to y \in \mathbb{R} \to 0 \leq -x \to 0 \leq y \to 0 \leq -x \cdot y)$
\item $\varphi_{246} :\equiv \forall v_{357}\,(\lvert x \rvert \cdot \lvert y \rvert = v_{357} \to v_{357} = \lvert x \rvert \cdot \lvert y \rvert) \to \lvert x \rvert \cdot \lvert y \rvert = -x \cdot -y \to -x \cdot -y = \lvert x \rvert \cdot \lvert y \rvert$
\item $\varphi_{247} :\equiv \forall x\,\forall y\,(x \in \mathbb{R} \to y \in \mathbb{R} \to x < y \to -y < -x) \to \forall y\,(v_{349} \in \mathbb{R} \to y \in \mathbb{R} \to v_{349} < y \to -y < -v_{349})$
\item $\varphi_{248} :\equiv \forall x\,\forall y\,(x \in \mathbb{R} \to y \in \mathbb{R} \to x < y \to -y < -x) \to \forall y\,(v_{317} \in \mathbb{R} \to y \in \mathbb{R} \to v_{317} < y \to -y < -v_{317})$
\item $\varphi_{249} :\equiv \forall v_{342}\,(\varphi_{61}) \to \forall v_{343}\,(y \in \mathbb{R} \to v_{343} \in \mathbb{R} \to \neg v_{343} \leq y \to y < v_{343})$
\item $\varphi_{250} :\equiv \forall v_{347}\,(\varphi_{62}) \to \forall v_{348}\,(y \in \mathbb{R} \to v_{348} \in \mathbb{R} \to \neg v_{348} \leq y \to y < v_{348})$
\item $\varphi_{251} :\equiv \varphi_{33} \to \lvert y \rvert = -y \to x \cdot \lvert y \rvert = x \cdot \lvert y \rvert \to x \cdot \lvert y \rvert = x \cdot -y$
\item $\varphi_{252} :\equiv \forall v_{310}\,(\varphi_{63}) \to \forall v_{311}\,(y \in \mathbb{R} \to v_{311} \in \mathbb{R} \to \neg v_{311} \leq y \to y < v_{311})$
\item $\varphi_{253} :\equiv \forall v_{315}\,(\varphi_{64}) \to \forall v_{316}\,(y \in \mathbb{R} \to v_{316} \in \mathbb{R} \to \neg v_{316} \leq y \to y < v_{316})$
\item $\varphi_{254} :\equiv \forall y\,(-x \in \mathbb{R} \to y \in \mathbb{R} \to 0 \leq -x \to 0 \leq y \to 0 \leq -x \cdot y) \to -x \in \mathbb{R} \to y \in \mathbb{R} \to 0 \leq -x \to 0 \leq y \to 0 \leq -x \cdot y$
\item $\varphi_{255} :\equiv (\lvert y \rvert = -y \to x \cdot \lvert y \rvert = x \cdot \lvert y \rvert \to x \cdot \lvert y \rvert = x \cdot -y) \to x \in \mathbb{R} \to y \in \mathbb{R} \to 0 \leq x \to \neg 0 \leq y \to x \cdot \lvert y \rvert = x \cdot -y$
\item $\varphi_{256} :\equiv \forall y\,(v_{349} \in \mathbb{R} \to y \in \mathbb{R} \to v_{349} < y \to -y < -v_{349}) \to \varphi_{98}$
\item $\varphi_{257} :\equiv \forall y\,(v_{317} \in \mathbb{R} \to y \in \mathbb{R} \to v_{317} < y \to -y < -v_{317}) \to \varphi_{101}$
\item $\varphi_{258} :\equiv \varphi_{120} \to x \in \mathbb{R} \to y \in \mathbb{R} \to 0 \leq x \to 0 \leq y \to \lvert x \rvert \cdot \lvert y \rvert = x \cdot y$
\item $\varphi_{259} :\equiv \varphi_{116} \to x \in \mathbb{R} \to y \in \mathbb{R} \to \neg 0 \leq x \to 0 \leq y \to \lvert x \cdot y \rvert = -x \cdot y$
\item $\varphi_{260} :\equiv \varphi_{118} \to x \in \mathbb{R} \to y \in \mathbb{R} \to 0 \leq x \to \neg 0 \leq y \to \lvert x \cdot y \rvert = x \cdot -y$
\item $\varphi_{261} :\equiv \forall y\,(-x \in \mathbb{R} \to y \in \mathbb{R} \to 0 \leq -x \to 0 \leq y \to 0 \leq -x \cdot y) \to -x \in \mathbb{R} \to -y \in \mathbb{R} \to 0 \leq -x \to 0 \leq -y \to 0 \leq -x \cdot -y$
\item $\varphi_{262} :\equiv (\lvert x \rvert \cdot \lvert y \rvert = -x \cdot -y \to -x \cdot -y = \lvert x \rvert \cdot \lvert y \rvert) \to \varphi_{110}$
\item $\varphi_{263} :\equiv \varphi_{115} \to x \in \mathbb{R} \to y \in \mathbb{R} \to \neg 0 \leq x \to 0 \leq y \to -x \cdot \lvert y \rvert = -x \cdot y$
\item $\varphi_{264} :\equiv (x \in \mathbb{R} \to y \in \mathbb{R} \to \neg 0 \leq x \to 0 \leq y \to x \in \mathbb{R}) \to 0 \in \mathbb{R} \to \varphi_{138}$
\item $\varphi_{265} :\equiv (x \in \mathbb{R} \to y \in \mathbb{R} \to \neg 0 \leq x \to \neg 0 \leq y \to x \in \mathbb{R}) \to 0 \in \mathbb{R} \to \varphi_{144}$
\item $\varphi_{266} :\equiv (x \in \mathbb{R} \to y \in \mathbb{R} \to 0 \leq x \to \neg 0 \leq y \to y \in \mathbb{R}) \to 0 \in \mathbb{R} \to \varphi_{159}$
\item $\varphi_{267} :\equiv (x \in \mathbb{R} \to y \in \mathbb{R} \to \neg 0 \leq x \to \neg 0 \leq y \to y \in \mathbb{R}) \to 0 \in \mathbb{R} \to \varphi_{161}$
\item $\varphi_{268} :\equiv -0 = 0 \to (x \in \mathbb{R} \to y \in \mathbb{R} \to \neg 0 \leq x \to 0 \leq y \to -0 < -x) \to \varphi_{164}$
\item $\varphi_{269} :\equiv (x \cdot y \in \mathbb{R} \to -(-(x \cdot y)) = x \cdot y) \to (x \in \mathbb{R} \to y \in \mathbb{R} \to \neg 0 \leq x \to 0 \leq y \to x \cdot y \in \mathbb{R}) \to \varphi_{54}$
\item $\varphi_{270} :\equiv -0 = 0 \to (x \in \mathbb{R} \to y \in \mathbb{R} \to 0 \leq x \to \neg 0 \leq y \to -0 < -y) \to \varphi_{166}$
\item $\varphi_{271} :\equiv (x \cdot y \in \mathbb{R} \to -(-(x \cdot y)) = x \cdot y) \to (x \in \mathbb{R} \to y \in \mathbb{R} \to 0 \leq x \to \neg 0 \leq y \to x \cdot y \in \mathbb{R}) \to \varphi_{58}$
\item $\varphi_{272} :\equiv -0 = 0 \to (x \in \mathbb{R} \to y \in \mathbb{R} \to \neg 0 \leq x \to \neg 0 \leq y \to -0 < -y) \to \varphi_{168}$
\item $\varphi_{273} :\equiv -0 = 0 \to (x \in \mathbb{R} \to y \in \mathbb{R} \to \neg 0 \leq x \to \neg 0 \leq y \to -0 < -x) \to \varphi_{169}$
\item $\varphi_{274} :\equiv (x \in \mathbb{R} \to y \in \mathbb{R} \to \neg 0 \leq x \to 0 \leq y \to x \in \mathbb{R}) \to 0 \in \mathbb{R} \to \varphi_{194}$
\item $\varphi_{275} :\equiv (x \in \mathbb{R} \to y \in \mathbb{R} \to 0 \leq x \to \neg 0 \leq y \to y \in \mathbb{R}) \to 0 \in \mathbb{R} \to \varphi_{195}$
\item $\varphi_{276} :\equiv (x \in \mathbb{R} \to y \in \mathbb{R} \to \neg 0 \leq x \to 0 \leq y \to x \cdot y \in \mathbb{R}) \to (x \in \mathbb{R} \to y \in \mathbb{R} \to \neg 0 \leq x \to 0 \leq y \to x \cdot y \leq 0) \to \varphi_{55}$
\item $\varphi_{277} :\equiv (x \in \mathbb{R} \to y \in \mathbb{R} \to 0 \leq x \to \neg 0 \leq y \to x \cdot y \in \mathbb{R}) \to (x \in \mathbb{R} \to y \in \mathbb{R} \to 0 \leq x \to \neg 0 \leq y \to x \cdot y \leq 0) \to \varphi_{59}$
\item $\varphi_{278} :\equiv (x \in \mathbb{R} \to y \in \mathbb{R} \to \neg 0 \leq x \to \neg 0 \leq y \to y \in \mathbb{R}) \to 0 \in \mathbb{R} \to \varphi_{202}$
\item $\varphi_{279} :\equiv (x \in \mathbb{R} \to y \in \mathbb{R} \to \neg 0 \leq x \to \neg 0 \leq y \to x \in \mathbb{R}) \to 0 \in \mathbb{R} \to \varphi_{203}$
\item $\varphi_{280} :\equiv (-x \cdot y = -(x \cdot y) \to 0 < -x \cdot y \to 0 < -(x \cdot y)) \to \varphi_{220}$
\item $\varphi_{281} :\equiv (x \cdot -y = -(x \cdot y) \to 0 < x \cdot -y \to 0 < -(x \cdot y)) \to \varphi_{222}$
\item $\varphi_{282} :\equiv (x \in \mathbb{R} \to y \in \mathbb{R} \to 0 \leq x \to 0 \leq y \to \lvert x \rvert \cdot \lvert y \rvert = x \cdot y) \to \varphi_{231}$
\item $\varphi_{283} :\equiv (x \in \mathbb{R} \to y \in \mathbb{R} \to \neg 0 \leq x \to 0 \leq y \to -x \cdot y = -(x \cdot y)) \to \varphi_{234}$
\item $\varphi_{284} :\equiv (x \in \mathbb{R} \to y \in \mathbb{R} \to \neg 0 \leq x \to 0 \leq y \to -(x \cdot y) \in \mathbb{R}) \to \varphi_{242}$
\item $\varphi_{285} :\equiv (x \in \mathbb{R} \to y \in \mathbb{R} \to 0 \leq x \to \neg 0 \leq y \to x \cdot -y = -(x \cdot y)) \to \varphi_{236}$
\item $\varphi_{286} :\equiv (x \in \mathbb{R} \to y \in \mathbb{R} \to 0 \leq x \to \neg 0 \leq y \to -(x \cdot y) \in \mathbb{R}) \to \varphi_{243}$
\item $\varphi_{287} :\equiv (x \in \mathbb{R} \to y \in \mathbb{R} \to \neg 0 \leq x \to \neg 0 \leq y \to -x \cdot -y = x \cdot y) \to \varphi_{238}$
\item $\varphi_{288} :\equiv (x \in \mathbb{R} \to y \in \mathbb{R} \to \neg 0 \leq x \to 0 \leq y \to x \cdot y \leq -0) \to (-0 = 0 \to x \cdot y < -0 \to x \cdot y < 0) \to \varphi_{139}$
\item $\varphi_{289} :\equiv (x \in \mathbb{R} \to y \in \mathbb{R} \to 0 \leq x \to \neg 0 \leq y \to x \cdot y \leq -0) \to (-0 = 0 \to x \cdot y < -0 \to x \cdot y < 0) \to \varphi_{143}$
\item $\varphi_{290} :\equiv (x \in \mathbb{R} \to y \in \mathbb{R} \to 0 \leq x \to 0 \leq y \to \lvert y \rvert = y) \to x \cdot \lvert y \rvert = x \cdot \lvert y \rvert \to \varphi_{237}$
\item $\varphi_{291} :\equiv (x \in \mathbb{R} \to y \in \mathbb{R} \to 0 \leq x \to \neg 0 \leq y \to -y \in \mathbb{R}) \to (x \in \mathbb{R} \to y \in \mathbb{R} \to 0 \leq x \to 0 \leq x) \to \varphi_{192}$
\item $\varphi_{292} :\equiv (x \in \mathbb{R} \to y \in \mathbb{R} \to 0 \leq x \to 0 \leq y \to x \cdot \lvert y \rvert = x \cdot y) \to \varphi_{258}$
\item $\varphi_{293} :\equiv \varphi_{84} \to (\lvert x \rvert \cdot \lvert y \rvert = -x \cdot y \to -x \cdot y = \lvert x \rvert \cdot \lvert y \rvert) \to \varphi_{85}$
\item $\varphi_{294} :\equiv (x \in \mathbb{R} \to y \in \mathbb{R} \to \neg 0 \leq x \to 0 \leq y \to y \in \mathbb{R}) \to (x \in \mathbb{R} \to y \in \mathbb{R} \to \neg 0 \leq x \to 0 \leq y \to 0 \leq -x) \to \varphi_{177}$
\item $\varphi_{295} :\equiv \varphi_{90} \to (\lvert x \rvert \cdot \lvert y \rvert = x \cdot -y \to x \cdot -y = \lvert x \rvert \cdot \lvert y \rvert) \to \varphi_{91}$
\item $\varphi_{296} :\equiv (x \in \mathbb{R} \to y \in \mathbb{R} \to x \cdot y \in \mathbb{R}) \to (x \in \mathbb{R} \to y \in \mathbb{R} \to \neg 0 \leq x \to 0 \leq y \to x \in \mathbb{R}) \to \varphi_{208}$
\item $\varphi_{297} :\equiv (x \in \mathbb{R} \to y \in \mathbb{R} \to x \cdot y \in \mathbb{R}) \to (x \in \mathbb{R} \to y \in \mathbb{R} \to 0 \leq x \to \neg 0 \leq y \to x \in \mathbb{R}) \to \varphi_{210}$
\item $\varphi_{298} :\equiv (x \in \mathbb{R} \to y \in \mathbb{R} \to \neg 0 \leq x \to 0 \leq y \to -(x \cdot y) = -x \cdot y) \to \varphi_{259}$
\item $\varphi_{299} :\equiv (x \in \mathbb{R} \to y \in \mathbb{R} \to 0 \leq x \to \neg 0 \leq y \to -(x \cdot y) = x \cdot -y) \to \varphi_{260}$
\item $\varphi_{300} :\equiv (x \in \mathbb{R} \to y \in \mathbb{R} \to 0 \leq x \to 0 \leq y \to x \cdot y = \lvert x \rvert \cdot \lvert y \rvert) \to \varphi_{119} \to \varphi_{82}$
\item $\varphi_{301} :\equiv (x \in \mathbb{R} \to y \in \mathbb{R} \to 0 \leq x \to \neg 0 \leq y \to \lvert y \rvert = -y) \to x \cdot \lvert y \rvert = x \cdot \lvert y \rvert \to \varphi_{255}$
\item $\varphi_{302} :\equiv (x \in \mathbb{R} \to y \in \mathbb{R} \to \neg 0 \leq x \to \neg 0 \leq y \to x \cdot y = -x \cdot -y) \to \varphi_{114} \to \varphi_{95}$
\item $\varphi_{303} :\equiv (x \in \mathbb{R} \to y \in \mathbb{R} \to \neg 0 \leq x \to \neg 0 \leq y \to -y \in \mathbb{R}) \to (x \in \mathbb{R} \to y \in \mathbb{R} \to \neg 0 \leq x \to \neg 0 \leq y \to 0 \leq -x) \to \varphi_{216}$
\item $\varphi_{304} :\equiv (x \in \mathbb{R} \to y \in \mathbb{R} \to -x \cdot y = -(x \cdot y)) \to (x \in \mathbb{R} \to y \in \mathbb{R} \to \neg 0 \leq x \to 0 \leq y \to x \in \mathbb{R}) \to \varphi_{233}$
\item $\varphi_{305} :\equiv (x \in \mathbb{R} \to y \in \mathbb{R} \to x \cdot -y = -(x \cdot y)) \to (x \in \mathbb{R} \to y \in \mathbb{R} \to 0 \leq x \to \neg 0 \leq y \to x \in \mathbb{R}) \to \varphi_{235}$
\item $\varphi_{306} :\equiv (x \in \mathbb{R} \to y \in \mathbb{R} \to 0 \leq x \to \neg 0 \leq y \to x \cdot \lvert y \rvert = x \cdot -y) \to \varphi_{141} \to \varphi_{90}$
\item $\varphi_{307} :\equiv (x \in \mathbb{R} \to y \in \mathbb{R} \to -x \cdot -y = x \cdot y) \to (x \in \mathbb{R} \to y \in \mathbb{R} \to \neg 0 \leq x \to \neg 0 \leq y \to x \in \mathbb{R}) \to \varphi_{239}$
\item $\varphi_{308} :\equiv (x \in \mathbb{R} \to y \in \mathbb{R} \to 0 \leq x \to 0 \leq y \to \lvert x \rvert = x) \to \lvert x \rvert \cdot \lvert y \rvert = \lvert x \rvert \cdot \lvert y \rvert \to \varphi_{117} \to \varphi_{81}$
\item $\varphi_{309} :\equiv (x \in \mathbb{R} \to y \in \mathbb{R} \to \neg 0 \leq x \to \neg 0 \leq y \to 0 \leq -x \cdot -y) \to (-x \cdot -y = x \cdot y \to 0 < -x \cdot -y \to 0 < x \cdot y) \to \varphi_{206}$
\item $\varphi_{310} :\equiv (x \in \mathbb{R} \to y \in \mathbb{R} \to \neg 0 \leq x \to 0 \leq y \to \lvert y \rvert = y) \to -x \cdot \lvert y \rvert = -x \cdot \lvert y \rvert \to \varphi_{263}$
\item $\varphi_{311} :\equiv (x \in \mathbb{R} \to y \in \mathbb{R} \to 0 \leq x \to \neg 0 \leq y \to \lvert x \rvert = x) \to \lvert x \rvert \cdot \lvert y \rvert = \lvert x \rvert \cdot \lvert y \rvert \to \varphi_{117} \to \varphi_{89}$
\item $\varphi_{312} :\equiv (x \in \mathbb{R} \to y \in \mathbb{R} \to \neg 0 \leq x \to 0 \leq y \to -x \cdot \lvert y \rvert = -x \cdot y) \to \varphi_{170} \to \varphi_{84}$
\item $\varphi_{313} :\equiv (x \in \mathbb{R} \to y \in \mathbb{R} \to \neg 0 \leq x \to 0 \leq y \to \lvert x \rvert = -x) \to \lvert x \rvert \cdot \lvert y \rvert = \lvert x \rvert \cdot \lvert y \rvert \to \varphi_{136} \to \varphi_{105}$
\item $\varphi_{314} :\equiv (x \in \mathbb{R} \to y \in \mathbb{R} \to \neg 0 \leq x \to \neg 0 \leq y \to \lvert y \rvert = -y) \to -x \cdot \lvert y \rvert = -x \cdot \lvert y \rvert \to \varphi_{135} \to \varphi_{107}$
\item $\varphi_{315} :\equiv (x \in \mathbb{R} \to y \in \mathbb{R} \to \neg 0 \leq x \to \neg 0 \leq y \to \lvert x \rvert = -x) \to \lvert x \rvert \cdot \lvert y \rvert = \lvert x \rvert \cdot \lvert y \rvert \to \varphi_{136} \to \varphi_{108}$
\item $\varphi_{316} :\equiv \varphi_{53} \to \varphi_{103} \to \varphi_{67} \to x \in \mathbb{R} \to y \in \mathbb{R} \to \neg 0 \leq x \to 0 \leq y \to x \cdot y \leq -0$
\item $\varphi_{317} :\equiv \varphi_{57} \to \varphi_{103} \to \varphi_{67} \to x \in \mathbb{R} \to y \in \mathbb{R} \to 0 \leq x \to \neg 0 \leq y \to x \cdot y \leq -0$
\item $\varphi_{318} :\equiv (x \in \mathbb{R} \to y \in \mathbb{R} \to 0 \leq x \to 0 \leq y \to \lvert x \cdot y \rvert = x \cdot y) \to \varphi_{300}$
\item $\varphi_{319} :\equiv (x \in \mathbb{R} \to y \in \mathbb{R} \to \neg 0 \leq x \to \neg 0 \leq y \to \lvert x \cdot y \rvert = x \cdot y) \to \varphi_{302}$
\item $\varphi_{320} :\equiv (x \in \mathbb{R} \to y \in \mathbb{R} \to \neg 0 \leq x \to 0 \leq y \to \lvert x \cdot y \rvert = -x \cdot y) \to \varphi_{85} \to \varphi_{137} \to \varphi_{87}$
\item $\varphi_{321} :\equiv (x \in \mathbb{R} \to y \in \mathbb{R} \to 0 \leq x \to \neg 0 \leq y \to \lvert x \cdot y \rvert = x \cdot -y) \to \varphi_{91} \to \varphi_{140} \to \varphi_{93}$
\item $\varphi_{322} :\equiv (x \in \mathbb{R} \to y \in \mathbb{R} \to \neg 0 \leq x \to 0 \leq y \to -x \cdot y = -(x \cdot y)) \to (x \in \mathbb{R} \to y \in \mathbb{R} \to \neg 0 \leq x \to 0 \leq y \to 0 \leq -x \cdot y) \to \varphi_{280}$
\item $\varphi_{323} :\equiv (x \in \mathbb{R} \to y \in \mathbb{R} \to 0 \leq x \to \neg 0 \leq y \to x \cdot -y = -(x \cdot y)) \to (x \in \mathbb{R} \to y \in \mathbb{R} \to 0 \leq x \to \neg 0 \leq y \to 0 \leq x \cdot -y) \to \varphi_{281}$
\item $\varphi_{324} :\equiv (x \in \mathbb{R} \to y \in \mathbb{R} \to \neg 0 \leq x \to \neg 0 \leq y \to -x \cdot -y = x \cdot y) \to \varphi_{309}$
\item $\varphi_{325} :\equiv (x \in \mathbb{R} \to -y \in \mathbb{R} \to 0 \leq x \to 0 \leq -y \to 0 \leq x \cdot -y) \to (x \in \mathbb{R} \to y \in \mathbb{R} \to 0 \leq x \to \neg 0 \leq y \to x \in \mathbb{R}) \to \varphi_{291}$
\item $\varphi_{326} :\equiv (-x \in \mathbb{R} \to y \in \mathbb{R} \to 0 \leq -x \to 0 \leq y \to 0 \leq -x \cdot y) \to (x \in \mathbb{R} \to y \in \mathbb{R} \to \neg 0 \leq x \to 0 \leq y \to -x \in \mathbb{R}) \to \varphi_{294}$
\item $\varphi_{327} :\equiv \varphi_{82} \to \varphi_{93} \to \varphi_{87} \to \varphi_{96} \to x \in \mathbb{R} \to y \in \mathbb{R} \to \lvert x \cdot y \rvert = \lvert x \rvert \cdot \lvert y \rvert$
\item $\varphi_{328} :\equiv (-x \in \mathbb{R} \to -y \in \mathbb{R} \to 0 \leq -x \to 0 \leq -y \to 0 \leq -x \cdot -y) \to (x \in \mathbb{R} \to y \in \mathbb{R} \to \neg 0 \leq x \to \neg 0 \leq y \to -x \in \mathbb{R}) \to \varphi_{303}$
\end{abbrevlist}
\end{fullproof}

\chapter{テキスト第2章 実数の集合 の証明木}
本章の定理数 10, ノード数の合計 1706.

\section{第2章: 上限・下限, 実数の連続性, アルキメデスの公理}
% 自動生成 (prover/build.js): NonemptyEx  ノード数 35
\begin{fullproof}{[B1]}{$\forall X\,(X \neq \emptyset \to \exists x\,(x \in X))$}{35}\phantomsection\label{tree:NonemptyEx}
\noindent$\Pi_{1}$:
\begin{ptree}
\AXC{$\forall x\,(x \notin X) \to z \notin X$\lbl{A2}}
\AXC{$X \neq \emptyset \to \neg \exists x\,(x \in X) \to \neg \exists x\,(x \in X)$\lbl{A1}}
\AXC{$\varphi_{3}$\lbl{A1}}
\BIC{$X \neq \emptyset \to \neg \exists x\,(x \in X) \to \forall x\,(x \notin X)$}
\AXC{$\varphi_{6}$\lbl{A1}}
\BIC{$(\forall x\,(x \notin X) \to z \notin X) \to X \neq \emptyset \to \neg \exists x\,(x \in X) \to z \notin X$}
\BIC{$X \neq \emptyset \to \neg \exists x\,(x \in X) \to z \notin X$}
\AXC{$\varphi_{8}$\lbl{A1}}
\BIC{$z \notin \emptyset \to X \neq \emptyset \to \neg \exists x\,(x \in X) \to (z \in X \leftrightarrow z \in \emptyset)$}
\end{ptree}
\noindent$\Pi_{2}$:
\begin{ptree}
\AXC{$\forall u\,(u \notin \emptyset)$\lbl{Def$\emptyset$}}
\AXC{$\forall u\,(u \notin \emptyset) \to z \notin \emptyset$\lbl{A2}}
\BIC{$z \notin \emptyset$}
\AXC{$\stackrel{\Pi_{1}}{\vdots}$}\noLine
\UIC{$z \notin \emptyset \to X \neq \emptyset \to \neg \exists x\,(x \in X) \to (z \in X \leftrightarrow z \in \emptyset)$}
\BIC{$X \neq \emptyset \to \neg \exists x\,(x \in X) \to (z \in X \leftrightarrow z \in \emptyset)$}
\AXC{$\varphi_{7}$\lbl{A1}}
\BIC{$X \neq \emptyset \land \neg \exists x\,(x \in X) \to (z \in X \leftrightarrow z \in \emptyset)$}
\UIC{$\forall z\,(X \neq \emptyset \land \neg \exists x\,(x \in X) \to (z \in X \leftrightarrow z \in \emptyset))$}
\end{ptree}
\noindent$\Pi_{3}$:
\begin{ptree}
\AXC{$\varphi_{10}$\lbl{A3}}
\AXC{$\stackrel{\Pi_{2}}{\vdots}$}\noLine
\UIC{$\forall z\,(X \neq \emptyset \land \neg \exists x\,(x \in X) \to (z \in X \leftrightarrow z \in \emptyset))$}
\AXC{$\forall z\,(X \neq \emptyset \land \neg \exists x\,(x \in X) \to (z \in X \leftrightarrow z \in \emptyset)) \to \varphi_{10} \to \varphi_{9}$\lbl{A1}}
\BIC{$\varphi_{10} \to \varphi_{9}$}
\BIC{$\varphi_{9}$}
\end{ptree}
\noindent$\Pi_{4}$:
\begin{ptree}
\AXC{$\forall x\,\forall y\,(\forall z\,(z \in x \leftrightarrow z \in y) \to x = y)$\lbl{Ext}}
\AXC{$\varphi_{1}$\lbl{A2}}
\BIC{$\forall y\,(\forall z\,(z \in X \leftrightarrow z \in y) \to X = y)$}
\AXC{$\varphi_{2}$\lbl{A2}}
\BIC{$\forall z\,(z \in X \leftrightarrow z \in \emptyset) \to X = \emptyset$}
\AXC{$(\forall z\,(z \in X \leftrightarrow z \in \emptyset) \to X = \emptyset) \to \varphi_{5}$\lbl{A1}}
\BIC{$\varphi_{5}$}
\end{ptree}
\noindent$\Pi_{5}$:
\begin{ptree}
\AXC{$X \neq \emptyset \land \neg \exists x\,(x \in X) \to \forall z\,(X \neq \emptyset \land \neg \exists x\,(x \in X))$\lbl{A4}}
\AXC{$\stackrel{\Pi_{3}}{\vdots}$}\noLine
\UIC{$\varphi_{9}$}
\BIC{$X \neq \emptyset \to \neg \exists x\,(x \in X) \to \forall z\,(z \in X \leftrightarrow z \in \emptyset)$}
\AXC{$\stackrel{\Pi_{4}}{\vdots}$}\noLine
\UIC{$\varphi_{5}$}
\BIC{$X \neq \emptyset \to \neg \exists x\,(x \in X) \to X = \emptyset$}
\end{ptree}
\noindent 主木:
\begin{ptree}
\AXC{$\stackrel{\Pi_{5}}{\vdots}$}\noLine
\UIC{$X \neq \emptyset \to \neg \exists x\,(x \in X) \to X = \emptyset$}
\AXC{$X \neq \emptyset \to X \neq \emptyset$\lbl{A1}}
\AXC{$\varphi_{4}$\lbl{A1}}
\BIC{$(X \neq \emptyset \to \neg \exists x\,(x \in X) \to X = \emptyset) \to X \neq \emptyset \to \exists x\,(x \in X)$}
\BIC{$X \neq \emptyset \to \exists x\,(x \in X)$}
\UIC{$\forall X\,(X \neq \emptyset \to \exists x\,(x \in X))$}
\end{ptree}
\begin{abbrevlist}
\item $\varphi_{1} :\equiv \forall x\,\forall y\,(\forall z\,(z \in x \leftrightarrow z \in y) \to x = y) \to \forall y\,(\forall z\,(z \in X \leftrightarrow z \in y) \to X = y)$
\item $\varphi_{2} :\equiv \forall y\,(\forall z\,(z \in X \leftrightarrow z \in y) \to X = y) \to \forall z\,(z \in X \leftrightarrow z \in \emptyset) \to X = \emptyset$
\item $\varphi_{3} :\equiv (X \neq \emptyset \to \neg \exists x\,(x \in X) \to \neg \exists x\,(x \in X)) \to X \neq \emptyset \to \neg \exists x\,(x \in X) \to \forall x\,(x \notin X)$
\item $\varphi_{4} :\equiv (X \neq \emptyset \to X \neq \emptyset) \to (X \neq \emptyset \to \neg \exists x\,(x \in X) \to X = \emptyset) \to X \neq \emptyset \to \exists x\,(x \in X)$
\item $\varphi_{5} :\equiv (X \neq \emptyset \to \neg \exists x\,(x \in X) \to \forall z\,(z \in X \leftrightarrow z \in \emptyset)) \to X \neq \emptyset \to \neg \exists x\,(x \in X) \to X = \emptyset$
\item $\varphi_{6} :\equiv (X \neq \emptyset \to \neg \exists x\,(x \in X) \to \forall x\,(x \notin X)) \to (\forall x\,(x \notin X) \to z \notin X) \to X \neq \emptyset \to \neg \exists x\,(x \in X) \to z \notin X$
\item $\varphi_{7} :\equiv (X \neq \emptyset \to \neg \exists x\,(x \in X) \to (z \in X \leftrightarrow z \in \emptyset)) \to X \neq \emptyset \land \neg \exists x\,(x \in X) \to (z \in X \leftrightarrow z \in \emptyset)$
\item $\varphi_{8} :\equiv (X \neq \emptyset \to \neg \exists x\,(x \in X) \to z \notin X) \to z \notin \emptyset \to X \neq \emptyset \to \neg \exists x\,(x \in X) \to (z \in X \leftrightarrow z \in \emptyset)$
\item $\varphi_{9} :\equiv (X \neq \emptyset \land \neg \exists x\,(x \in X) \to \forall z\,(X \neq \emptyset \land \neg \exists x\,(x \in X))) \to X \neq \emptyset \to \neg \exists x\,(x \in X) \to \forall z\,(z \in X \leftrightarrow z \in \emptyset)$
\item $\varphi_{10} :\equiv \forall z\,(X \neq \emptyset \land \neg \exists x\,(x \in X) \to (z \in X \leftrightarrow z \in \emptyset)) \to \forall z\,(X \neq \emptyset \land \neg \exists x\,(x \in X)) \to \forall z\,(z \in X \leftrightarrow z \in \emptyset)$
\end{abbrevlist}
\end{fullproof}

% 自動生成 (prover/build.js): ExNonempty  ノード数 21
\begin{fullproof}{[B2]}{$\forall X\,\forall x \in X\,(X \neq \emptyset)$}{21}\phantomsection\label{tree:ExNonempty}
\noindent$\Pi_{1}$:
\begin{ptree}
\AXC{$z = w \to x \in z \to x \in w$\lbl{A6}}
\UIC{$\forall w\,(z = w \to x \in z \to x \in w)$}
\UIC{$\forall z\,\forall w\,(z = w \to x \in z \to x \in w)$}
\AXC{$\varphi_{1}$\lbl{A2}}
\BIC{$\forall w\,(X = w \to x \in X \to x \in w)$}
\AXC{$\forall w\,(X = w \to x \in X \to x \in w) \to X = \emptyset \to x \in X \to x \in \emptyset$\lbl{A2}}
\BIC{$X = \emptyset \to x \in X \to x \in \emptyset$}
\end{ptree}
\noindent$\Pi_{2}$:
\begin{ptree}
\AXC{$\stackrel{\Pi_{1}}{\vdots}$}\noLine
\UIC{$X = \emptyset \to x \in X \to x \in \emptyset$}
\AXC{$x \in X \to x \in X$\lbl{A1}}
\AXC{$x \in X \to X = \emptyset \to X = \emptyset$\lbl{A1}}
\AXC{$(x \in X \to X = \emptyset \to X = \emptyset) \to \varphi_{2}$\lbl{A1}}
\BIC{$\varphi_{2}$}
\BIC{$(X = \emptyset \to x \in X \to x \in \emptyset) \to x \in X \to X = \emptyset \to x \in \emptyset$}
\BIC{$x \in X \to X = \emptyset \to x \in \emptyset$}
\end{ptree}
\noindent$\Pi_{3}$:
\begin{ptree}
\AXC{$\stackrel{\Pi_{2}}{\vdots}$}\noLine
\UIC{$x \in X \to X = \emptyset \to x \in \emptyset$}
\AXC{$(x \in X \to X = \emptyset \to x \in \emptyset) \to x \notin \emptyset \to x \in X \to X \neq \emptyset$\lbl{A1}}
\BIC{$x \notin \emptyset \to x \in X \to X \neq \emptyset$}
\end{ptree}
\noindent 主木:
\begin{ptree}
\AXC{$\forall u\,(u \notin \emptyset)$\lbl{Def$\emptyset$}}
\AXC{$\forall u\,(u \notin \emptyset) \to x \notin \emptyset$\lbl{A2}}
\BIC{$x \notin \emptyset$}
\AXC{$\stackrel{\Pi_{3}}{\vdots}$}\noLine
\UIC{$x \notin \emptyset \to x \in X \to X \neq \emptyset$}
\BIC{$x \in X \to X \neq \emptyset$}
\UIC{$\forall x \in X\,(X \neq \emptyset)$}
\UIC{$\forall X\,\forall x \in X\,(X \neq \emptyset)$}
\end{ptree}
\begin{abbrevlist}
\item $\varphi_{1} :\equiv \forall z\,\forall w\,(z = w \to x \in z \to x \in w) \to \forall w\,(X = w \to x \in X \to x \in w)$
\item $\varphi_{2} :\equiv (x \in X \to x \in X) \to (X = \emptyset \to x \in X \to x \in \emptyset) \to x \in X \to X = \emptyset \to x \in \emptyset$
\end{abbrevlist}
\end{fullproof}

% 自動生成 (prover/build.js): SupUnique  ノード数 50
\begin{fullproof}{[B3]}{$\forall s\,\forall t\,\forall X\,(\mathrm{IsSup}(s,X) \to \mathrm{IsSup}(t,X) \to s = t)$}{50}\phantomsection\label{tree:SupUnique}
\noindent$\Pi_{1}$:
\begin{ptree}
\AXC{$\forall s\,\forall X\,(\varphi_{5})$\lbl{DefSup}}
\AXC{$\forall s\,\forall X\,(\varphi_{5}) \to \forall X\,(\varphi_{3})$\lbl{A2}}
\BIC{$\forall X\,(\varphi_{3})$}
\AXC{$\forall X\,(\varphi_{3}) \to \varphi_{3}$\lbl{A2}}
\BIC{$\varphi_{3}$}
\end{ptree}
\noindent$\Pi_{2}$:
\begin{ptree}
\AXC{$\stackrel{\Pi_{1}}{\vdots}$}\noLine
\UIC{$\varphi_{3}$}
\AXC{$\varphi_{2}$\lbl{A1}}
\AXC{$\varphi_{2} \to \varphi_{3} \to \varphi_{7}$\lbl{A1}}
\BIC{$\varphi_{3} \to \varphi_{7}$}
\BIC{$\varphi_{7}$}
\AXC{$\varphi_{7} \to \varphi_{11}$\lbl{A1}}
\BIC{$\varphi_{11}$}
\end{ptree}
\noindent$\Pi_{3}$:
\begin{ptree}
\AXC{$\varphi_{2}$\lbl{A1}}
\AXC{$\mathrm{IsSup}(s,X) \to \mathrm{IsSup}(s,X)$\lbl{A1}}
\AXC{$(\mathrm{IsSup}(s,X) \to \mathrm{IsSup}(s,X)) \to \varphi_{2} \to \varphi_{14}$\lbl{A1}}
\BIC{$\varphi_{2} \to \varphi_{14}$}
\BIC{$\varphi_{14}$}
\end{ptree}
\noindent$\Pi_{4}$:
\begin{ptree}
\AXC{$\forall s\,\forall X\,(\varphi_{5})$\lbl{DefSup}}
\AXC{$\forall s\,\forall X\,(\varphi_{5}) \to \forall X\,(\varphi_{5})$\lbl{A2}}
\BIC{$\forall X\,(\varphi_{5})$}
\AXC{$\forall X\,(\varphi_{5}) \to \varphi_{5}$\lbl{A2}}
\BIC{$\varphi_{5}$}
\AXC{$\stackrel{\Pi_{3}}{\vdots}$}\noLine
\UIC{$\varphi_{14}$}
\BIC{$\varphi_{3} \to \varphi_{4} \to \mathrm{IsSup}(s,X) \to \mathrm{IsSup}(t,X) \to t \leq s$}
\end{ptree}
\noindent$\Pi_{5}$:
\begin{ptree}
\AXC{$\stackrel{\Pi_{1}}{\vdots}$}\noLine
\UIC{$\varphi_{3}$}
\AXC{$\stackrel{\Pi_{4}}{\vdots}$}\noLine
\UIC{$\varphi_{3} \to \varphi_{4} \to \mathrm{IsSup}(s,X) \to \mathrm{IsSup}(t,X) \to t \leq s$}
\BIC{$\varphi_{4} \to \mathrm{IsSup}(s,X) \to \mathrm{IsSup}(t,X) \to t \leq s$}
\end{ptree}
\noindent$\Pi_{6}$:
\begin{ptree}
\AXC{$\forall w\,(\mathrm{UB}(w,X) \to t \leq w) \to \mathrm{UB}(s,X) \to t \leq s$\lbl{A2}}
\AXC{$\stackrel{\Pi_{2}}{\vdots}$}\noLine
\UIC{$\varphi_{11}$}
\BIC{$\varphi_{4}$}
\AXC{$\stackrel{\Pi_{5}}{\vdots}$}\noLine
\UIC{$\varphi_{4} \to \mathrm{IsSup}(s,X) \to \mathrm{IsSup}(t,X) \to t \leq s$}
\BIC{$\mathrm{IsSup}(s,X) \to \mathrm{IsSup}(t,X) \to t \leq s$}
\end{ptree}
\noindent$\Pi_{7}$:
\begin{ptree}
\AXC{$\mathrm{IsSup}(s,X) \to \mathrm{IsSup}(s,X)$\lbl{A1}}
\AXC{$(\mathrm{IsSup}(s,X) \to \mathrm{IsSup}(s,X)) \to \varphi_{5} \to \varphi_{8}$\lbl{A1}}
\BIC{$\varphi_{5} \to \varphi_{8}$}
\end{ptree}
\noindent$\Pi_{8}$:
\begin{ptree}
\AXC{$\forall s\,\forall X\,(\varphi_{5})$\lbl{DefSup}}
\AXC{$\forall s\,\forall X\,(\varphi_{5}) \to \forall X\,(\varphi_{5})$\lbl{A2}}
\BIC{$\forall X\,(\varphi_{5})$}
\AXC{$\forall X\,(\varphi_{5}) \to \varphi_{5}$\lbl{A2}}
\BIC{$\varphi_{5}$}
\AXC{$\stackrel{\Pi_{7}}{\vdots}$}\noLine
\UIC{$\varphi_{5} \to \varphi_{8}$}
\BIC{$\varphi_{8}$}
\AXC{$\varphi_{8} \to \varphi_{12}$\lbl{A1}}
\BIC{$\varphi_{12}$}
\end{ptree}
\noindent$\Pi_{9}$:
\begin{ptree}
\AXC{$\varphi_{2}$\lbl{A1}}
\AXC{$\mathrm{IsSup}(s,X) \to \mathrm{IsSup}(s,X)$\lbl{A1}}
\AXC{$(\mathrm{IsSup}(s,X) \to \mathrm{IsSup}(s,X)) \to \varphi_{2} \to \varphi_{15}$\lbl{A1}}
\BIC{$\varphi_{2} \to \varphi_{15}$}
\BIC{$\varphi_{15}$}
\end{ptree}
\noindent$\Pi_{10}$:
\begin{ptree}
\AXC{$\stackrel{\Pi_{1}}{\vdots}$}\noLine
\UIC{$\varphi_{3}$}
\AXC{$\forall s\,\forall X\,(\varphi_{5})$\lbl{DefSup}}
\AXC{$\forall s\,\forall X\,(\varphi_{5}) \to \forall X\,(\varphi_{5})$\lbl{A2}}
\BIC{$\forall X\,(\varphi_{5})$}
\AXC{$\forall X\,(\varphi_{5}) \to \varphi_{5}$\lbl{A2}}
\BIC{$\varphi_{5}$}
\AXC{$\stackrel{\Pi_{9}}{\vdots}$}\noLine
\UIC{$\varphi_{15}$}
\BIC{$\varphi_{3} \to \varphi_{6} \to \mathrm{IsSup}(s,X) \to \mathrm{IsSup}(t,X) \to s \leq t$}
\BIC{$\varphi_{6} \to \mathrm{IsSup}(s,X) \to \mathrm{IsSup}(t,X) \to s \leq t$}
\end{ptree}
\noindent$\Pi_{11}$:
\begin{ptree}
\AXC{$\forall w\,(\mathrm{UB}(w,X) \to s \leq w) \to \mathrm{UB}(t,X) \to s \leq t$\lbl{A2}}
\AXC{$\stackrel{\Pi_{8}}{\vdots}$}\noLine
\UIC{$\varphi_{12}$}
\BIC{$\varphi_{6}$}
\AXC{$\stackrel{\Pi_{10}}{\vdots}$}\noLine
\UIC{$\varphi_{6} \to \mathrm{IsSup}(s,X) \to \mathrm{IsSup}(t,X) \to s \leq t$}
\BIC{$\mathrm{IsSup}(s,X) \to \mathrm{IsSup}(t,X) \to s \leq t$}
\end{ptree}
\noindent$\Pi_{12}$:
\begin{ptree}
\AXC{$\forall x\,\forall y\,(x \leq y \to y \leq x \to x = y)$\lbl{[O6]}}
\AXC{$\varphi_{9}$\lbl{A2}}
\BIC{$\forall y\,(s \leq y \to y \leq s \to s = y)$}
\AXC{$\forall y\,(s \leq y \to y \leq s \to s = y) \to s \leq t \to t \leq s \to s = t$\lbl{A2}}
\BIC{$s \leq t \to t \leq s \to s = t$}
\AXC{$\varphi_{13}$\lbl{A1}}
\BIC{$(\mathrm{IsSup}(s,X) \to \mathrm{IsSup}(t,X) \to s \leq t) \to \varphi_{10}$}
\end{ptree}
\noindent$\Pi_{13}$:
\begin{ptree}
\AXC{$\stackrel{\Pi_{11}}{\vdots}$}\noLine
\UIC{$\mathrm{IsSup}(s,X) \to \mathrm{IsSup}(t,X) \to s \leq t$}
\AXC{$\stackrel{\Pi_{12}}{\vdots}$}\noLine
\UIC{$(\mathrm{IsSup}(s,X) \to \mathrm{IsSup}(t,X) \to s \leq t) \to \varphi_{10}$}
\BIC{$\varphi_{10}$}
\end{ptree}
\noindent 主木:
\begin{ptree}
\AXC{$\stackrel{\Pi_{6}}{\vdots}$}\noLine
\UIC{$\mathrm{IsSup}(s,X) \to \mathrm{IsSup}(t,X) \to t \leq s$}
\AXC{$\stackrel{\Pi_{13}}{\vdots}$}\noLine
\UIC{$\varphi_{10}$}
\BIC{$\mathrm{IsSup}(s,X) \to \mathrm{IsSup}(t,X) \to s = t$}
\UIC{$\forall X\,(\mathrm{IsSup}(s,X) \to \mathrm{IsSup}(t,X) \to s = t)$}
\UIC{$\varphi_{1}$}
\UIC{$\forall s\,(\varphi_{1})$}
\end{ptree}
\begin{abbrevlist}
\item $\varphi_{1} :\equiv \forall t\,\forall X\,(\mathrm{IsSup}(s,X) \to \mathrm{IsSup}(t,X) \to s = t)$
\item $\varphi_{2} :\equiv \mathrm{IsSup}(s,X) \to \mathrm{IsSup}(t,X) \to \mathrm{IsSup}(t,X)$
\item $\varphi_{3} :\equiv \mathrm{IsSup}(t,X) \leftrightarrow \mathrm{UB}(t,X) \land \forall w\,(\mathrm{UB}(w,X) \to t \leq w)$
\item $\varphi_{4} :\equiv \mathrm{IsSup}(s,X) \to \mathrm{IsSup}(t,X) \to \mathrm{UB}(s,X) \to t \leq s$
\item $\varphi_{5} :\equiv \mathrm{IsSup}(s,X) \leftrightarrow \mathrm{UB}(s,X) \land \forall w\,(\mathrm{UB}(w,X) \to s \leq w)$
\item $\varphi_{6} :\equiv \mathrm{IsSup}(s,X) \to \mathrm{IsSup}(t,X) \to \mathrm{UB}(t,X) \to s \leq t$
\item $\varphi_{7} :\equiv \mathrm{IsSup}(s,X) \to \mathrm{IsSup}(t,X) \to \forall w\,(\mathrm{UB}(w,X) \to t \leq w)$
\item $\varphi_{8} :\equiv \mathrm{IsSup}(s,X) \to \mathrm{IsSup}(t,X) \to \forall w\,(\mathrm{UB}(w,X) \to s \leq w)$
\item $\varphi_{9} :\equiv \forall x\,\forall y\,(x \leq y \to y \leq x \to x = y) \to \forall y\,(s \leq y \to y \leq s \to s = y)$
\item $\varphi_{10} :\equiv (\mathrm{IsSup}(s,X) \to \mathrm{IsSup}(t,X) \to t \leq s) \to \mathrm{IsSup}(s,X) \to \mathrm{IsSup}(t,X) \to s = t$
\item $\varphi_{11} :\equiv (\forall w\,(\mathrm{UB}(w,X) \to t \leq w) \to \mathrm{UB}(s,X) \to t \leq s) \to \varphi_{4}$
\item $\varphi_{12} :\equiv (\forall w\,(\mathrm{UB}(w,X) \to s \leq w) \to \mathrm{UB}(t,X) \to s \leq t) \to \varphi_{6}$
\item $\varphi_{13} :\equiv (s \leq t \to t \leq s \to s = t) \to (\mathrm{IsSup}(s,X) \to \mathrm{IsSup}(t,X) \to s \leq t) \to \varphi_{10}$
\item $\varphi_{14} :\equiv \varphi_{5} \to \varphi_{3} \to \varphi_{4} \to \mathrm{IsSup}(s,X) \to \mathrm{IsSup}(t,X) \to t \leq s$
\item $\varphi_{15} :\equiv \varphi_{5} \to \varphi_{3} \to \varphi_{6} \to \mathrm{IsSup}(s,X) \to \mathrm{IsSup}(t,X) \to s \leq t$
\end{abbrevlist}
\end{fullproof}

% 自動生成 (prover/build.js): InfUnique  ノード数 50
\begin{fullproof}{[B4]}{$\forall s\,\forall t\,\forall X\,(\mathrm{IsInf}(s,X) \to \mathrm{IsInf}(t,X) \to s = t)$}{50}\phantomsection\label{tree:InfUnique}
\noindent$\Pi_{1}$:
\begin{ptree}
\AXC{$\mathrm{IsInf}(s,X) \to \mathrm{IsInf}(s,X)$\lbl{A1}}
\AXC{$(\mathrm{IsInf}(s,X) \to \mathrm{IsInf}(s,X)) \to \varphi_{3} \to \varphi_{7}$\lbl{A1}}
\BIC{$\varphi_{3} \to \varphi_{7}$}
\end{ptree}
\noindent$\Pi_{2}$:
\begin{ptree}
\AXC{$\forall s\,\forall X\,(\varphi_{3})$\lbl{DefInf}}
\AXC{$\forall s\,\forall X\,(\varphi_{3}) \to \forall X\,(\varphi_{3})$\lbl{A2}}
\BIC{$\forall X\,(\varphi_{3})$}
\AXC{$\forall X\,(\varphi_{3}) \to \varphi_{3}$\lbl{A2}}
\BIC{$\varphi_{3}$}
\AXC{$\stackrel{\Pi_{1}}{\vdots}$}\noLine
\UIC{$\varphi_{3} \to \varphi_{7}$}
\BIC{$\varphi_{7}$}
\AXC{$\varphi_{7} \to \varphi_{11}$\lbl{A1}}
\BIC{$\varphi_{11}$}
\end{ptree}
\noindent$\Pi_{3}$:
\begin{ptree}
\AXC{$\forall s\,\forall X\,(\varphi_{3})$\lbl{DefInf}}
\AXC{$\forall s\,\forall X\,(\varphi_{3}) \to \forall X\,(\varphi_{5})$\lbl{A2}}
\BIC{$\forall X\,(\varphi_{5})$}
\AXC{$\forall X\,(\varphi_{5}) \to \varphi_{5}$\lbl{A2}}
\BIC{$\varphi_{5}$}
\end{ptree}
\noindent$\Pi_{4}$:
\begin{ptree}
\AXC{$\varphi_{2}$\lbl{A1}}
\AXC{$\mathrm{IsInf}(s,X) \to \mathrm{IsInf}(s,X)$\lbl{A1}}
\AXC{$(\mathrm{IsInf}(s,X) \to \mathrm{IsInf}(s,X)) \to \varphi_{2} \to \varphi_{14}$\lbl{A1}}
\BIC{$\varphi_{2} \to \varphi_{14}$}
\BIC{$\varphi_{14}$}
\end{ptree}
\noindent$\Pi_{5}$:
\begin{ptree}
\AXC{$\forall s\,\forall X\,(\varphi_{3})$\lbl{DefInf}}
\AXC{$\forall s\,\forall X\,(\varphi_{3}) \to \forall X\,(\varphi_{3})$\lbl{A2}}
\BIC{$\forall X\,(\varphi_{3})$}
\AXC{$\forall X\,(\varphi_{3}) \to \varphi_{3}$\lbl{A2}}
\BIC{$\varphi_{3}$}
\AXC{$\stackrel{\Pi_{4}}{\vdots}$}\noLine
\UIC{$\varphi_{14}$}
\BIC{$\varphi_{5} \to \varphi_{4} \to \mathrm{IsInf}(s,X) \to \mathrm{IsInf}(t,X) \to t \leq s$}
\end{ptree}
\noindent$\Pi_{6}$:
\begin{ptree}
\AXC{$\stackrel{\Pi_{3}}{\vdots}$}\noLine
\UIC{$\varphi_{5}$}
\AXC{$\stackrel{\Pi_{5}}{\vdots}$}\noLine
\UIC{$\varphi_{5} \to \varphi_{4} \to \mathrm{IsInf}(s,X) \to \mathrm{IsInf}(t,X) \to t \leq s$}
\BIC{$\varphi_{4} \to \mathrm{IsInf}(s,X) \to \mathrm{IsInf}(t,X) \to t \leq s$}
\end{ptree}
\noindent$\Pi_{7}$:
\begin{ptree}
\AXC{$\forall w\,(\mathrm{LB}(w,X) \to w \leq s) \to \mathrm{LB}(t,X) \to t \leq s$\lbl{A2}}
\AXC{$\stackrel{\Pi_{2}}{\vdots}$}\noLine
\UIC{$\varphi_{11}$}
\BIC{$\varphi_{4}$}
\AXC{$\stackrel{\Pi_{6}}{\vdots}$}\noLine
\UIC{$\varphi_{4} \to \mathrm{IsInf}(s,X) \to \mathrm{IsInf}(t,X) \to t \leq s$}
\BIC{$\mathrm{IsInf}(s,X) \to \mathrm{IsInf}(t,X) \to t \leq s$}
\end{ptree}
\noindent$\Pi_{8}$:
\begin{ptree}
\AXC{$\forall w\,(\mathrm{LB}(w,X) \to w \leq t) \to \mathrm{LB}(s,X) \to s \leq t$\lbl{A2}}
\AXC{$\stackrel{\Pi_{3}}{\vdots}$}\noLine
\UIC{$\varphi_{5}$}
\AXC{$\varphi_{2}$\lbl{A1}}
\AXC{$\varphi_{2} \to \varphi_{5} \to \varphi_{8}$\lbl{A1}}
\BIC{$\varphi_{5} \to \varphi_{8}$}
\BIC{$\varphi_{8}$}
\AXC{$\varphi_{8} \to \varphi_{12}$\lbl{A1}}
\BIC{$\varphi_{12}$}
\BIC{$\varphi_{6}$}
\end{ptree}
\noindent$\Pi_{9}$:
\begin{ptree}
\AXC{$\varphi_{2}$\lbl{A1}}
\AXC{$\mathrm{IsInf}(s,X) \to \mathrm{IsInf}(s,X)$\lbl{A1}}
\AXC{$(\mathrm{IsInf}(s,X) \to \mathrm{IsInf}(s,X)) \to \varphi_{2} \to \varphi_{15}$\lbl{A1}}
\BIC{$\varphi_{2} \to \varphi_{15}$}
\BIC{$\varphi_{15}$}
\end{ptree}
\noindent$\Pi_{10}$:
\begin{ptree}
\AXC{$\stackrel{\Pi_{8}}{\vdots}$}\noLine
\UIC{$\varphi_{6}$}
\AXC{$\stackrel{\Pi_{3}}{\vdots}$}\noLine
\UIC{$\varphi_{5}$}
\AXC{$\forall s\,\forall X\,(\varphi_{3})$\lbl{DefInf}}
\AXC{$\forall s\,\forall X\,(\varphi_{3}) \to \forall X\,(\varphi_{3})$\lbl{A2}}
\BIC{$\forall X\,(\varphi_{3})$}
\AXC{$\forall X\,(\varphi_{3}) \to \varphi_{3}$\lbl{A2}}
\BIC{$\varphi_{3}$}
\AXC{$\stackrel{\Pi_{9}}{\vdots}$}\noLine
\UIC{$\varphi_{15}$}
\BIC{$\varphi_{5} \to \varphi_{6} \to \mathrm{IsInf}(s,X) \to \mathrm{IsInf}(t,X) \to s \leq t$}
\BIC{$\varphi_{6} \to \mathrm{IsInf}(s,X) \to \mathrm{IsInf}(t,X) \to s \leq t$}
\BIC{$\mathrm{IsInf}(s,X) \to \mathrm{IsInf}(t,X) \to s \leq t$}
\end{ptree}
\noindent$\Pi_{11}$:
\begin{ptree}
\AXC{$\forall x\,\forall y\,(x \leq y \to y \leq x \to x = y)$\lbl{[O6]}}
\AXC{$\varphi_{9}$\lbl{A2}}
\BIC{$\forall y\,(s \leq y \to y \leq s \to s = y)$}
\AXC{$\forall y\,(s \leq y \to y \leq s \to s = y) \to s \leq t \to t \leq s \to s = t$\lbl{A2}}
\BIC{$s \leq t \to t \leq s \to s = t$}
\AXC{$\varphi_{13}$\lbl{A1}}
\BIC{$(\mathrm{IsInf}(s,X) \to \mathrm{IsInf}(t,X) \to s \leq t) \to \varphi_{10}$}
\end{ptree}
\noindent$\Pi_{12}$:
\begin{ptree}
\AXC{$\stackrel{\Pi_{10}}{\vdots}$}\noLine
\UIC{$\mathrm{IsInf}(s,X) \to \mathrm{IsInf}(t,X) \to s \leq t$}
\AXC{$\stackrel{\Pi_{11}}{\vdots}$}\noLine
\UIC{$(\mathrm{IsInf}(s,X) \to \mathrm{IsInf}(t,X) \to s \leq t) \to \varphi_{10}$}
\BIC{$\varphi_{10}$}
\end{ptree}
\noindent 主木:
\begin{ptree}
\AXC{$\stackrel{\Pi_{7}}{\vdots}$}\noLine
\UIC{$\mathrm{IsInf}(s,X) \to \mathrm{IsInf}(t,X) \to t \leq s$}
\AXC{$\stackrel{\Pi_{12}}{\vdots}$}\noLine
\UIC{$\varphi_{10}$}
\BIC{$\mathrm{IsInf}(s,X) \to \mathrm{IsInf}(t,X) \to s = t$}
\UIC{$\forall X\,(\mathrm{IsInf}(s,X) \to \mathrm{IsInf}(t,X) \to s = t)$}
\UIC{$\varphi_{1}$}
\UIC{$\forall s\,(\varphi_{1})$}
\end{ptree}
\begin{abbrevlist}
\item $\varphi_{1} :\equiv \forall t\,\forall X\,(\mathrm{IsInf}(s,X) \to \mathrm{IsInf}(t,X) \to s = t)$
\item $\varphi_{2} :\equiv \mathrm{IsInf}(s,X) \to \mathrm{IsInf}(t,X) \to \mathrm{IsInf}(t,X)$
\item $\varphi_{3} :\equiv \mathrm{IsInf}(s,X) \leftrightarrow \mathrm{LB}(s,X) \land \forall w\,(\mathrm{LB}(w,X) \to w \leq s)$
\item $\varphi_{4} :\equiv \mathrm{IsInf}(s,X) \to \mathrm{IsInf}(t,X) \to \mathrm{LB}(t,X) \to t \leq s$
\item $\varphi_{5} :\equiv \mathrm{IsInf}(t,X) \leftrightarrow \mathrm{LB}(t,X) \land \forall w\,(\mathrm{LB}(w,X) \to w \leq t)$
\item $\varphi_{6} :\equiv \mathrm{IsInf}(s,X) \to \mathrm{IsInf}(t,X) \to \mathrm{LB}(s,X) \to s \leq t$
\item $\varphi_{7} :\equiv \mathrm{IsInf}(s,X) \to \mathrm{IsInf}(t,X) \to \forall w\,(\mathrm{LB}(w,X) \to w \leq s)$
\item $\varphi_{8} :\equiv \mathrm{IsInf}(s,X) \to \mathrm{IsInf}(t,X) \to \forall w\,(\mathrm{LB}(w,X) \to w \leq t)$
\item $\varphi_{9} :\equiv \forall x\,\forall y\,(x \leq y \to y \leq x \to x = y) \to \forall y\,(s \leq y \to y \leq s \to s = y)$
\item $\varphi_{10} :\equiv (\mathrm{IsInf}(s,X) \to \mathrm{IsInf}(t,X) \to t \leq s) \to \mathrm{IsInf}(s,X) \to \mathrm{IsInf}(t,X) \to s = t$
\item $\varphi_{11} :\equiv (\forall w\,(\mathrm{LB}(w,X) \to w \leq s) \to \mathrm{LB}(t,X) \to t \leq s) \to \varphi_{4}$
\item $\varphi_{12} :\equiv (\forall w\,(\mathrm{LB}(w,X) \to w \leq t) \to \mathrm{LB}(s,X) \to s \leq t) \to \varphi_{6}$
\item $\varphi_{13} :\equiv (s \leq t \to t \leq s \to s = t) \to (\mathrm{IsInf}(s,X) \to \mathrm{IsInf}(t,X) \to s \leq t) \to \varphi_{10}$
\item $\varphi_{14} :\equiv \varphi_{3} \to \varphi_{5} \to \varphi_{4} \to \mathrm{IsInf}(s,X) \to \mathrm{IsInf}(t,X) \to t \leq s$
\item $\varphi_{15} :\equiv \varphi_{3} \to \varphi_{5} \to \varphi_{6} \to \mathrm{IsInf}(s,X) \to \mathrm{IsInf}(t,X) \to s \leq t$
\end{abbrevlist}
\end{fullproof}

% 自動生成 (prover/build.js): SubElem  ノード数 21
\begin{fullproof}{[B5]}{$\forall X\,\forall Y\,\forall x\,(X \subseteq Y \to x \in X \to x \in Y)$}{21}\phantomsection\label{tree:SubElem}
\noindent$\Pi_{1}$:
\begin{ptree}
\AXC{$\forall x\,\forall y\,(x \subseteq y \leftrightarrow \forall u \in x\,(u \in y))$\lbl{Def$\subseteq$}}
\AXC{$\varphi_{2}$\lbl{A2}}
\BIC{$\forall y\,(X \subseteq y \leftrightarrow \forall u \in X\,(u \in y))$}
\AXC{$\varphi_{1}$\lbl{A2}}
\BIC{$X \subseteq Y \leftrightarrow \forall u \in X\,(u \in Y)$}
\AXC{$X \subseteq Y \to X \subseteq Y$\lbl{A1}}
\AXC{$(X \subseteq Y \to X \subseteq Y) \to \varphi_{3}$\lbl{A1}}
\BIC{$\varphi_{3}$}
\BIC{$X \subseteq Y \to x \in X \to \forall u \in X\,(u \in Y)$}
\end{ptree}
\noindent$\Pi_{2}$:
\begin{ptree}
\AXC{$\forall u \in X\,(u \in Y) \to x \in X \to x \in Y$\lbl{A2}}
\AXC{$\stackrel{\Pi_{1}}{\vdots}$}\noLine
\UIC{$X \subseteq Y \to x \in X \to \forall u \in X\,(u \in Y)$}
\AXC{$(X \subseteq Y \to x \in X \to \forall u \in X\,(u \in Y)) \to \varphi_{4}$\lbl{A1}}
\BIC{$\varphi_{4}$}
\BIC{$X \subseteq Y \to x \in X \to x \in X \to x \in Y$}
\end{ptree}
\noindent 主木:
\begin{ptree}
\AXC{$\stackrel{\Pi_{2}}{\vdots}$}\noLine
\UIC{$X \subseteq Y \to x \in X \to x \in X \to x \in Y$}
\AXC{$X \subseteq Y \to x \in X \to x \in X$\lbl{A1}}
\AXC{$X \subseteq Y \to X \subseteq Y$\lbl{A1}}
\AXC{$(X \subseteq Y \to X \subseteq Y) \to \varphi_{5}$\lbl{A1}}
\BIC{$\varphi_{5}$}
\BIC{$(X \subseteq Y \to x \in X \to x \in X \to x \in Y) \to X \subseteq Y \to x \in X \to x \in Y$}
\BIC{$X \subseteq Y \to x \in X \to x \in Y$}
\UIC{$\forall x\,(X \subseteq Y \to x \in X \to x \in Y)$}
\UIC{$\forall Y\,\forall x\,(X \subseteq Y \to x \in X \to x \in Y)$}
\UIC{$\forall X\,\forall Y\,\forall x\,(X \subseteq Y \to x \in X \to x \in Y)$}
\end{ptree}
\begin{abbrevlist}
\item $\varphi_{1} :\equiv \forall y\,(X \subseteq y \leftrightarrow \forall u \in X\,(u \in y)) \to (X \subseteq Y \leftrightarrow \forall u \in X\,(u \in Y))$
\item $\varphi_{2} :\equiv \forall x\,\forall y\,(x \subseteq y \leftrightarrow \forall u \in x\,(u \in y)) \to \forall y\,(X \subseteq y \leftrightarrow \forall u \in X\,(u \in y))$
\item $\varphi_{3} :\equiv (X \subseteq Y \leftrightarrow \forall u \in X\,(u \in Y)) \to X \subseteq Y \to x \in X \to \forall u \in X\,(u \in Y)$
\item $\varphi_{4} :\equiv (\forall u \in X\,(u \in Y) \to x \in X \to x \in Y) \to X \subseteq Y \to x \in X \to x \in X \to x \in Y$
\item $\varphi_{5} :\equiv (X \subseteq Y \to x \in X \to x \in X) \to (X \subseteq Y \to x \in X \to x \in X \to x \in Y) \to X \subseteq Y \to x \in X \to x \in Y$
\end{abbrevlist}
\end{fullproof}

% 自動生成 (prover/build.js): InfExists  ノード数 634
\begin{fullproof}{公理3.1(2)}{$\forall X\,(X \subseteq \mathbb{R} \to X \neq \emptyset \to \exists b\,(\mathrm{LB}(b,X)) \to \exists s\,(\mathrm{IsInf}(s,X)))$}{634}\phantomsection\label{tree:InfExists}
\noindent$\Pi_{1}$:
\begin{ptree}
\AXC{$\varphi_{101}$\lbl{A4}}
\AXC{$\forall w\,(\varphi_{80}) \to \varphi_{103}$\lbl{A3}}
\AXC{$\varphi_{80}$\lbl{A2}}
\UIC{$\forall w\,(\varphi_{80})$}
\AXC{$\forall w\,(\varphi_{80}) \to (\forall w\,(\varphi_{80}) \to \varphi_{103}) \to \varphi_{101} \to \varphi_{94}$\lbl{A1}}
\BIC{$(\forall w\,(\varphi_{80}) \to \varphi_{103}) \to \varphi_{101} \to \varphi_{94}$}
\BIC{$\varphi_{101} \to \varphi_{94}$}
\BIC{$\varphi_{94}$}
\end{ptree}
\noindent$\Pi_{2}$:
\begin{ptree}
\AXC{$\varphi_{102}$\lbl{A4}}
\AXC{$\forall a\,(\varphi_{81}) \to \varphi_{104}$\lbl{A3}}
\AXC{$\varphi_{81}$\lbl{A2}}
\UIC{$\forall a\,(\varphi_{81})$}
\AXC{$\forall a\,(\varphi_{81}) \to (\forall a\,(\varphi_{81}) \to \varphi_{104}) \to \varphi_{102} \to \varphi_{95}$\lbl{A1}}
\BIC{$(\forall a\,(\varphi_{81}) \to \varphi_{104}) \to \varphi_{102} \to \varphi_{95}$}
\BIC{$\varphi_{102} \to \varphi_{95}$}
\BIC{$\varphi_{95}$}
\end{ptree}
\noindent$\Pi_{3}$:
\begin{ptree}
\AXC{$\varphi_{129}$\lbl{A4}}
\AXC{$\forall x\,(\varphi_{113}) \to \varphi_{136}$\lbl{A3}}
\AXC{$\varphi_{113}$\lbl{A2}}
\UIC{$\forall x\,(\varphi_{113})$}
\AXC{$\forall x\,(\varphi_{113}) \to (\forall x\,(\varphi_{113}) \to \varphi_{136}) \to \varphi_{129} \to \varphi_{121}$\lbl{A1}}
\BIC{$(\forall x\,(\varphi_{113}) \to \varphi_{136}) \to \varphi_{129} \to \varphi_{121}$}
\BIC{$\varphi_{129} \to \varphi_{121}$}
\BIC{$\varphi_{121}$}
\end{ptree}
\noindent$\Pi_{4}$:
\begin{ptree}
\AXC{$\varphi_{130}$\lbl{A4}}
\AXC{$\forall y\,(\varphi_{114}) \to \varphi_{137}$\lbl{A3}}
\AXC{$\varphi_{114}$\lbl{A2}}
\UIC{$\forall y\,(\varphi_{114})$}
\AXC{$\forall y\,(\varphi_{114}) \to (\forall y\,(\varphi_{114}) \to \varphi_{137}) \to \varphi_{130} \to \varphi_{122}$\lbl{A1}}
\BIC{$(\forall y\,(\varphi_{114}) \to \varphi_{137}) \to \varphi_{130} \to \varphi_{122}$}
\BIC{$\varphi_{130} \to \varphi_{122}$}
\BIC{$\varphi_{122}$}
\end{ptree}
\noindent$\Pi_{5}$:
\begin{ptree}
\AXC{$\forall X\,\forall y\,(y \in \ominus X \leftrightarrow y \in \mathbb{R} \land -y \in X)$\lbl{Def$\ominus$}}
\AXC{$\varphi_{144}$\lbl{A2}}
\BIC{$\forall y\,(y \in \ominus X \leftrightarrow y \in \mathbb{R} \land -y \in X)$}
\AXC{$\varphi_{125}$\lbl{A2}}
\BIC{$y \in \ominus X \leftrightarrow y \in \mathbb{R} \land -y \in X$}
\AXC{$X \subseteq \mathbb{R} \to X \neq \emptyset \to \varphi_{99}$\lbl{A1}}
\AXC{$\varphi_{220}$\lbl{A1}}
\BIC{$(y \in \ominus X \leftrightarrow y \in \mathbb{R} \land -y \in X) \to X \subseteq \mathbb{R} \to X \neq \emptyset \to \varphi_{91}$}
\BIC{$X \subseteq \mathbb{R} \to X \neq \emptyset \to \varphi_{91}$}
\end{ptree}
\noindent$\Pi_{6}$:
\begin{ptree}
\AXC{$\forall a\,\forall X\,(\mathrm{LB}(a,X) \leftrightarrow a \in \mathbb{R} \land \forall x \in X\,(a \leq x))$\lbl{DefLB}}
\AXC{$\varphi_{117}$\lbl{A2}}
\BIC{$\forall X\,(\mathrm{LB}(a,X) \leftrightarrow a \in \mathbb{R} \land \forall x \in X\,(a \leq x))$}
\AXC{$\varphi_{111}$\lbl{A2}}
\BIC{$\mathrm{LB}(a,X) \leftrightarrow a \in \mathbb{R} \land \forall x \in X\,(a \leq x)$}
\end{ptree}
\noindent$\Pi_{7}$:
\begin{ptree}
\AXC{$\stackrel{\Pi_{6}}{\vdots}$}\noLine
\UIC{$\mathrm{LB}(a,X) \leftrightarrow a \in \mathbb{R} \land \forall x \in X\,(a \leq x)$}
\AXC{$X \subseteq \mathbb{R} \to X \neq \emptyset \to \varphi_{12}$\lbl{A1}}
\AXC{$(X \subseteq \mathbb{R} \to X \neq \emptyset \to \varphi_{12}) \to \varphi_{177}$\lbl{A1}}
\BIC{$\varphi_{177}$}
\BIC{$X \subseteq \mathbb{R} \to X \neq \emptyset \to \varphi_{45}$}
\end{ptree}
\noindent$\Pi_{8}$:
\begin{ptree}
\AXC{$\forall x \in X\,(a \leq x) \to -y \in X \to a \leq -y$\lbl{A2}}
\AXC{$\stackrel{\Pi_{7}}{\vdots}$}\noLine
\UIC{$X \subseteq \mathbb{R} \to X \neq \emptyset \to \varphi_{45}$}
\AXC{$(X \subseteq \mathbb{R} \to X \neq \emptyset \to \varphi_{45}) \to \varphi_{183}$\lbl{A1}}
\BIC{$\varphi_{183}$}
\BIC{$X \subseteq \mathbb{R} \to X \neq \emptyset \to \varphi_{68}$}
\AXC{$\varphi_{229}$\lbl{A1}}
\BIC{$(X \subseteq \mathbb{R} \to X \neq \emptyset \to \varphi_{91}) \to X \subseteq \mathbb{R} \to X \neq \emptyset \to \varphi_{72}$}
\end{ptree}
\noindent$\Pi_{9}$:
\begin{ptree}
\AXC{$\stackrel{\Pi_{5}}{\vdots}$}\noLine
\UIC{$X \subseteq \mathbb{R} \to X \neq \emptyset \to \varphi_{91}$}
\AXC{$\stackrel{\Pi_{8}}{\vdots}$}\noLine
\UIC{$(X \subseteq \mathbb{R} \to X \neq \emptyset \to \varphi_{91}) \to X \subseteq \mathbb{R} \to X \neq \emptyset \to \varphi_{72}$}
\BIC{$X \subseteq \mathbb{R} \to X \neq \emptyset \to \varphi_{72}$}
\end{ptree}
\noindent$\Pi_{10}$:
\begin{ptree}
\AXC{$\forall X\,\forall y\,(y \in \ominus X \leftrightarrow y \in \mathbb{R} \land -y \in X)$\lbl{Def$\ominus$}}
\AXC{$\varphi_{144}$\lbl{A2}}
\BIC{$\forall y\,(y \in \ominus X \leftrightarrow y \in \mathbb{R} \land -y \in X)$}
\AXC{$\varphi_{125}$\lbl{A2}}
\BIC{$y \in \ominus X \leftrightarrow y \in \mathbb{R} \land -y \in X$}
\AXC{$X \subseteq \mathbb{R} \to X \neq \emptyset \to \varphi_{99}$\lbl{A1}}
\AXC{$\varphi_{219}$\lbl{A1}}
\BIC{$(y \in \ominus X \leftrightarrow y \in \mathbb{R} \land -y \in X) \to X \subseteq \mathbb{R} \to X \neq \emptyset \to \varphi_{85}$}
\BIC{$X \subseteq \mathbb{R} \to X \neq \emptyset \to \varphi_{85}$}
\end{ptree}
\noindent$\Pi_{11}$:
\begin{ptree}
\AXC{$\stackrel{\Pi_{10}}{\vdots}$}\noLine
\UIC{$X \subseteq \mathbb{R} \to X \neq \emptyset \to \varphi_{85}$}
\AXC{$\forall x \in \mathbb{R}\,(-x \in \mathbb{R})$\lbl{(R1)}}
\AXC{$\forall x \in \mathbb{R}\,(-x \in \mathbb{R}) \to y \in \mathbb{R} \to -y \in \mathbb{R}$\lbl{A2}}
\BIC{$y \in \mathbb{R} \to -y \in \mathbb{R}$}
\AXC{$\varphi_{214}$\lbl{A1}}
\BIC{$(X \subseteq \mathbb{R} \to X \neq \emptyset \to \varphi_{85}) \to X \subseteq \mathbb{R} \to X \neq \emptyset \to \varphi_{100}$}
\BIC{$X \subseteq \mathbb{R} \to X \neq \emptyset \to \varphi_{100}$}
\end{ptree}
\noindent$\Pi_{12}$:
\begin{ptree}
\AXC{$\stackrel{\Pi_{6}}{\vdots}$}\noLine
\UIC{$\mathrm{LB}(a,X) \leftrightarrow a \in \mathbb{R} \land \forall x \in X\,(a \leq x)$}
\AXC{$X \subseteq \mathbb{R} \to X \neq \emptyset \to \varphi_{12}$\lbl{A1}}
\AXC{$(X \subseteq \mathbb{R} \to X \neq \emptyset \to \varphi_{12}) \to \varphi_{171}$\lbl{A1}}
\BIC{$\varphi_{171}$}
\BIC{$X \subseteq \mathbb{R} \to X \neq \emptyset \to \varphi_{24}$}
\end{ptree}
\noindent$\Pi_{13}$:
\begin{ptree}
\AXC{$\stackrel{\Pi_{12}}{\vdots}$}\noLine
\UIC{$X \subseteq \mathbb{R} \to X \neq \emptyset \to \varphi_{24}$}
\AXC{$\forall x\,\forall y\,(x \in \mathbb{R} \to y \in \mathbb{R} \to x \leq y \to -y \leq -x)$\lbl{[F23]}}
\AXC{$\varphi_{146}$\lbl{A2}}
\BIC{$\forall y\,(a \in \mathbb{R} \to y \in \mathbb{R} \to a \leq y \to -y \leq -a)$}
\AXC{$\varphi_{161}$\lbl{A2}}
\BIC{$a \in \mathbb{R} \to -y \in \mathbb{R} \to a \leq -y \to -(-y) \leq -a$}
\AXC{$\varphi_{239}$\lbl{A1}}
\BIC{$(X \subseteq \mathbb{R} \to X \neq \emptyset \to \varphi_{24}) \to (X \subseteq \mathbb{R} \to X \neq \emptyset \to \varphi_{100}) \to \varphi_{210}$}
\BIC{$(X \subseteq \mathbb{R} \to X \neq \emptyset \to \varphi_{100}) \to \varphi_{210}$}
\end{ptree}
\noindent$\Pi_{14}$:
\begin{ptree}
\AXC{$\forall x \in \mathbb{R}\,(-(-x) = x)$\lbl{[F5]}}
\AXC{$\forall x \in \mathbb{R}\,(-(-x) = x) \to y \in \mathbb{R} \to -(-y) = y$\lbl{A2}}
\BIC{$y \in \mathbb{R} \to -(-y) = y$}
\AXC{$(y \in \mathbb{R} \to -(-y) = y) \to \varphi_{208}$\lbl{A1}}
\BIC{$\varphi_{208}$}
\end{ptree}
\noindent$\Pi_{15}$:
\begin{ptree}
\AXC{$\stackrel{\Pi_{10}}{\vdots}$}\noLine
\UIC{$X \subseteq \mathbb{R} \to X \neq \emptyset \to \varphi_{85}$}
\AXC{$\stackrel{\Pi_{14}}{\vdots}$}\noLine
\UIC{$\varphi_{208}$}
\BIC{$X \subseteq \mathbb{R} \to X \neq \emptyset \to \exists b\,(\mathrm{LB}(b,X)) \to \varphi_{51}$}
\AXC{$(X \subseteq \mathbb{R} \to X \neq \emptyset \to \exists b\,(\mathrm{LB}(b,X)) \to \varphi_{51}) \to \varphi_{230}$\lbl{A1}}
\BIC{$\varphi_{230}$}
\end{ptree}
\noindent$\Pi_{16}$:
\begin{ptree}
\AXC{$\stackrel{\Pi_{9}}{\vdots}$}\noLine
\UIC{$X \subseteq \mathbb{R} \to X \neq \emptyset \to \varphi_{72}$}
\AXC{$\stackrel{\Pi_{11}}{\vdots}$}\noLine
\UIC{$X \subseteq \mathbb{R} \to X \neq \emptyset \to \varphi_{100}$}
\AXC{$\stackrel{\Pi_{13}}{\vdots}$}\noLine
\UIC{$(X \subseteq \mathbb{R} \to X \neq \emptyset \to \varphi_{100}) \to \varphi_{210}$}
\BIC{$\varphi_{210}$}
\BIC{$X \subseteq \mathbb{R} \to X \neq \emptyset \to \exists b\,(\mathrm{LB}(b,X)) \to \varphi_{69}$}
\AXC{$\stackrel{\Pi_{15}}{\vdots}$}\noLine
\UIC{$\varphi_{230}$}
\BIC{$(-(-y) = y \to -(-y) < -a \to y < -a) \to \varphi_{188}$}
\end{ptree}
\noindent$\Pi_{17}$:
\begin{ptree}
\AXC{$z = w \to z < -a \to w < -a$\lbl{A6}}
\UIC{$\forall w\,(z = w \to z < -a \to w < -a)$}
\UIC{$\forall z\,\forall w\,(z = w \to z < -a \to w < -a)$}
\AXC{$\varphi_{131}$\lbl{A2}}
\BIC{$\forall w\,(-(-y) = w \to -(-y) < -a \to w < -a)$}
\AXC{$\varphi_{157}$\lbl{A2}}
\BIC{$-(-y) = y \to -(-y) < -a \to y < -a$}
\AXC{$\stackrel{\Pi_{16}}{\vdots}$}\noLine
\UIC{$(-(-y) = y \to -(-y) < -a \to y < -a) \to \varphi_{188}$}
\BIC{$\varphi_{188}$}
\end{ptree}
\noindent$\Pi_{18}$:
\begin{ptree}
\AXC{$z = w \to z = -a \to w = -a$\lbl{A6}}
\UIC{$\forall w\,(z = w \to z = -a \to w = -a)$}
\UIC{$\forall z\,\forall w\,(z = w \to z = -a \to w = -a)$}
\AXC{$\varphi_{107}$\lbl{A2}}
\BIC{$\forall w\,(-(-y) = w \to -(-y) = -a \to w = -a)$}
\AXC{$\varphi_{141}$\lbl{A2}}
\BIC{$-(-y) = y \to -(-y) = -a \to y = -a$}
\AXC{$\stackrel{\Pi_{17}}{\vdots}$}\noLine
\UIC{$\varphi_{188}$}
\BIC{$X \subseteq \mathbb{R} \to X \neq \emptyset \to \varphi_{73}$}
\AXC{$\varphi_{202}$\lbl{A1}}
\BIC{$\varphi_{41} \land \mathrm{LB}(a,X) \to y \in \ominus X \to y \leq -a$}
\UIC{$\forall y\,(\varphi_{41} \land \mathrm{LB}(a,X) \to y \in \ominus X \to y \leq -a)$}
\end{ptree}
\noindent$\Pi_{19}$:
\begin{ptree}
\AXC{$\varphi_{209}$\lbl{A3}}
\AXC{$\stackrel{\Pi_{18}}{\vdots}$}\noLine
\UIC{$\forall y\,(\varphi_{41} \land \mathrm{LB}(a,X) \to y \in \ominus X \to y \leq -a)$}
\AXC{$\forall y\,(\varphi_{41} \land \mathrm{LB}(a,X) \to y \in \ominus X \to y \leq -a) \to \varphi_{209} \to \varphi_{217}$\lbl{A1}}
\BIC{$\varphi_{209} \to \varphi_{217}$}
\BIC{$\varphi_{217}$}
\end{ptree}
\noindent$\Pi_{20}$:
\begin{ptree}
\AXC{$\stackrel{\Pi_{12}}{\vdots}$}\noLine
\UIC{$X \subseteq \mathbb{R} \to X \neq \emptyset \to \varphi_{24}$}
\AXC{$\forall x \in \mathbb{R}\,(-x \in \mathbb{R})$\lbl{(R1)}}
\AXC{$\forall x \in \mathbb{R}\,(-x \in \mathbb{R}) \to a \in \mathbb{R} \to -a \in \mathbb{R}$\lbl{A2}}
\BIC{$a \in \mathbb{R} \to -a \in \mathbb{R}$}
\AXC{$\varphi_{205}$\lbl{A1}}
\BIC{$(X \subseteq \mathbb{R} \to X \neq \emptyset \to \varphi_{24}) \to X \subseteq \mathbb{R} \to X \neq \emptyset \to \varphi_{42}$}
\BIC{$X \subseteq \mathbb{R} \to X \neq \emptyset \to \varphi_{42}$}
\end{ptree}
\noindent$\Pi_{21}$:
\begin{ptree}
\AXC{$\stackrel{\Pi_{20}}{\vdots}$}\noLine
\UIC{$X \subseteq \mathbb{R} \to X \neq \emptyset \to \varphi_{42}$}
\AXC{$\forall a\,\forall X\,(\mathrm{UB}(a,X) \leftrightarrow a \in \mathbb{R} \land \forall x \in X\,(x \leq a))$\lbl{DefUB}}
\AXC{$\varphi_{151}$\lbl{A2}}
\BIC{$\forall X\,(\mathrm{UB}(-a,X) \leftrightarrow -a \in \mathbb{R} \land \forall x \in X\,(x \leq -a))$}
\AXC{$\varphi_{178}$\lbl{A2}}
\BIC{$\mathrm{UB}(-a,\ominus X) \leftrightarrow -a \in \mathbb{R} \land \forall x \in \ominus X\,(x \leq -a)$}
\AXC{$\varphi_{244}$\lbl{A1}}
\BIC{$(X \subseteq \mathbb{R} \to X \neq \emptyset \to \varphi_{42}) \to \varphi_{234}$}
\BIC{$\varphi_{234}$}
\end{ptree}
\noindent$\Pi_{22}$:
\begin{ptree}
\AXC{$\stackrel{\Pi_{3}}{\vdots}$}\noLine
\UIC{$\varphi_{121}$}
\AXC{$\stackrel{\Pi_{4}}{\vdots}$}\noLine
\UIC{$\varphi_{122}$}
\AXC{$\varphi_{41} \land \mathrm{LB}(a,X) \to \forall y\,(\varphi_{41} \land \mathrm{LB}(a,X))$\lbl{A4}}
\AXC{$\stackrel{\Pi_{19}}{\vdots}$}\noLine
\UIC{$\varphi_{217}$}
\BIC{$X \subseteq \mathbb{R} \to X \neq \emptyset \to \exists b\,(\mathrm{LB}(b,X)) \to \varphi_{46}$}
\AXC{$\stackrel{\Pi_{21}}{\vdots}$}\noLine
\UIC{$\varphi_{234}$}
\BIC{$\varphi_{122} \to \varphi_{121} \to X \subseteq \mathbb{R} \to X \neq \emptyset \to \varphi_{62}$}
\BIC{$\varphi_{121} \to X \subseteq \mathbb{R} \to X \neq \emptyset \to \varphi_{62}$}
\BIC{$X \subseteq \mathbb{R} \to X \neq \emptyset \to \varphi_{62}$}
\end{ptree}
\noindent$\Pi_{23}$:
\begin{ptree}
\AXC{$\forall s\,(\varphi_{2})$\lbl{DefSup}}
\AXC{$\forall s\,(\varphi_{2}) \to \varphi_{2}$\lbl{A2}}
\BIC{$\varphi_{2}$}
\AXC{$\varphi_{2} \to \varphi_{55}$\lbl{A2}}
\BIC{$\varphi_{55}$}
\AXC{$X \subseteq \mathbb{R} \to X \neq \emptyset \to \varphi_{15}$\lbl{A1}}
\AXC{$\varphi_{216}$\lbl{A1}}
\BIC{$\varphi_{55} \to X \subseteq \mathbb{R} \to X \neq \emptyset \to \exists b\,(\mathrm{LB}(b,X)) \to \varphi_{28}$}
\BIC{$X \subseteq \mathbb{R} \to X \neq \emptyset \to \exists b\,(\mathrm{LB}(b,X)) \to \varphi_{28}$}
\AXC{$\varphi_{221}$\lbl{A1}}
\BIC{$\varphi_{36} \to X \subseteq \mathbb{R} \to X \neq \emptyset \to \exists b\,(\mathrm{LB}(b,X)) \to \varphi_{52}$}
\end{ptree}
\noindent$\Pi_{24}$:
\begin{ptree}
\AXC{$\stackrel{\Pi_{22}}{\vdots}$}\noLine
\UIC{$X \subseteq \mathbb{R} \to X \neq \emptyset \to \varphi_{62}$}
\AXC{$\varphi_{36}$\lbl{A2}}
\AXC{$\stackrel{\Pi_{23}}{\vdots}$}\noLine
\UIC{$\varphi_{36} \to X \subseteq \mathbb{R} \to X \neq \emptyset \to \exists b\,(\mathrm{LB}(b,X)) \to \varphi_{52}$}
\BIC{$X \subseteq \mathbb{R} \to X \neq \emptyset \to \exists b\,(\mathrm{LB}(b,X)) \to \varphi_{52}$}
\AXC{$\varphi_{227}$\lbl{A1}}
\BIC{$(X \subseteq \mathbb{R} \to X \neq \emptyset \to \varphi_{62}) \to X \subseteq \mathbb{R} \to X \neq \emptyset \to \varphi_{16}$}
\BIC{$X \subseteq \mathbb{R} \to X \neq \emptyset \to \varphi_{16}$}
\end{ptree}
\noindent$\Pi_{25}$:
\begin{ptree}
\AXC{$\forall a\,\forall X\,(\mathrm{UB}(a,X) \leftrightarrow a \in \mathbb{R} \land \forall x \in X\,(x \leq a))$\lbl{DefUB}}
\AXC{$\varphi_{118}$\lbl{A2}}
\BIC{$\forall X\,(\mathrm{UB}(s,X) \leftrightarrow s \in \mathbb{R} \land \forall x \in X\,(x \leq s))$}
\AXC{$\varphi_{145}$\lbl{A2}}
\BIC{$\mathrm{UB}(s,\ominus X) \leftrightarrow s \in \mathbb{R} \land \forall x \in \ominus X\,(x \leq s)$}
\end{ptree}
\noindent$\Pi_{26}$:
\begin{ptree}
\AXC{$X \subseteq \mathbb{R} \to X \neq \emptyset \to \varphi_{15}$\lbl{A1}}
\AXC{$(X \subseteq \mathbb{R} \to X \neq \emptyset \to \varphi_{15}) \to \varphi_{55} \to X \subseteq \mathbb{R} \to X \neq \emptyset \to \varphi_{8}$\lbl{A1}}
\BIC{$\varphi_{55} \to X \subseteq \mathbb{R} \to X \neq \emptyset \to \varphi_{8}$}
\end{ptree}
\noindent$\Pi_{27}$:
\begin{ptree}
\AXC{$\forall s\,(\varphi_{2})$\lbl{DefSup}}
\AXC{$\forall s\,(\varphi_{2}) \to \varphi_{2}$\lbl{A2}}
\BIC{$\varphi_{2}$}
\AXC{$\varphi_{2} \to \varphi_{55}$\lbl{A2}}
\BIC{$\varphi_{55}$}
\AXC{$\stackrel{\Pi_{26}}{\vdots}$}\noLine
\UIC{$\varphi_{55} \to X \subseteq \mathbb{R} \to X \neq \emptyset \to \varphi_{8}$}
\BIC{$X \subseteq \mathbb{R} \to X \neq \emptyset \to \varphi_{8}$}
\end{ptree}
\noindent$\Pi_{28}$:
\begin{ptree}
\AXC{$\stackrel{\Pi_{25}}{\vdots}$}\noLine
\UIC{$\mathrm{UB}(s,\ominus X) \leftrightarrow s \in \mathbb{R} \land \forall x \in \ominus X\,(x \leq s)$}
\AXC{$\stackrel{\Pi_{27}}{\vdots}$}\noLine
\UIC{$X \subseteq \mathbb{R} \to X \neq \emptyset \to \varphi_{8}$}
\AXC{$(X \subseteq \mathbb{R} \to X \neq \emptyset \to \varphi_{8}) \to \varphi_{179}$\lbl{A1}}
\BIC{$\varphi_{179}$}
\BIC{$X \subseteq \mathbb{R} \to \varphi_{38}$}
\end{ptree}
\noindent$\Pi_{29}$:
\begin{ptree}
\AXC{$\stackrel{\Pi_{28}}{\vdots}$}\noLine
\UIC{$X \subseteq \mathbb{R} \to \varphi_{38}$}
\AXC{$\forall x\,\forall y\,(x \in \mathbb{R} \to y \in \mathbb{R} \to x \leq y \to -y \leq -x)$\lbl{[F23]}}
\AXC{$\varphi_{147}$\lbl{A2}}
\BIC{$\forall y\,(s \in \mathbb{R} \to y \in \mathbb{R} \to s \leq y \to -y \leq -s)$}
\AXC{$\varphi_{162}$\lbl{A2}}
\BIC{$s \in \mathbb{R} \to -a \in \mathbb{R} \to s \leq -a \to -(-a) \leq -s$}
\AXC{$\varphi_{236}$\lbl{A1}}
\BIC{$(X \subseteq \mathbb{R} \to \varphi_{38}) \to \varphi_{222}$}
\BIC{$\varphi_{222}$}
\end{ptree}
\noindent$\Pi_{30}$:
\begin{ptree}
\AXC{$\stackrel{\Pi_{12}}{\vdots}$}\noLine
\UIC{$X \subseteq \mathbb{R} \to X \neq \emptyset \to \varphi_{24}$}
\AXC{$\forall x \in \mathbb{R}\,(-(-x) = x)$\lbl{[F5]}}
\AXC{$\forall x \in \mathbb{R}\,(-(-x) = x) \to a \in \mathbb{R} \to -(-a) = a$\lbl{A2}}
\BIC{$a \in \mathbb{R} \to -(-a) = a$}
\AXC{$\varphi_{203}$\lbl{A1}}
\BIC{$(X \subseteq \mathbb{R} \to X \neq \emptyset \to \varphi_{24}) \to X \subseteq \mathbb{R} \to X \neq \emptyset \to \varphi_{32}$}
\BIC{$X \subseteq \mathbb{R} \to X \neq \emptyset \to \varphi_{32}$}
\end{ptree}
\noindent$\Pi_{31}$:
\begin{ptree}
\AXC{$\stackrel{\Pi_{24}}{\vdots}$}\noLine
\UIC{$X \subseteq \mathbb{R} \to X \neq \emptyset \to \varphi_{16}$}
\AXC{$\stackrel{\Pi_{20}}{\vdots}$}\noLine
\UIC{$X \subseteq \mathbb{R} \to X \neq \emptyset \to \varphi_{42}$}
\AXC{$\stackrel{\Pi_{29}}{\vdots}$}\noLine
\UIC{$\varphi_{222}$}
\BIC{$(X \subseteq \mathbb{R} \to X \neq \emptyset \to \varphi_{16}) \to X \subseteq \mathbb{R} \to X \neq \emptyset \to \varphi_{53}$}
\BIC{$X \subseteq \mathbb{R} \to X \neq \emptyset \to \varphi_{53}$}
\AXC{$\stackrel{\Pi_{30}}{\vdots}$}\noLine
\UIC{$X \subseteq \mathbb{R} \to X \neq \emptyset \to \varphi_{32}$}
\AXC{$(X \subseteq \mathbb{R} \to X \neq \emptyset \to \varphi_{32}) \to \varphi_{225}$\lbl{A1}}
\BIC{$\varphi_{225}$}
\BIC{$(-(-a) = a \to -(-a) < -s \to a < -s) \to \varphi_{173}$}
\end{ptree}
\noindent$\Pi_{32}$:
\begin{ptree}
\AXC{$z = w \to z < -s \to w < -s$\lbl{A6}}
\UIC{$\forall w\,(z = w \to z < -s \to w < -s)$}
\UIC{$\forall z\,\forall w\,(z = w \to z < -s \to w < -s)$}
\AXC{$\varphi_{128}$\lbl{A2}}
\BIC{$\forall w\,(-(-a) = w \to -(-a) < -s \to w < -s)$}
\AXC{$\varphi_{156}$\lbl{A2}}
\BIC{$-(-a) = a \to -(-a) < -s \to a < -s$}
\AXC{$\stackrel{\Pi_{31}}{\vdots}$}\noLine
\UIC{$(-(-a) = a \to -(-a) < -s \to a < -s) \to \varphi_{173}$}
\BIC{$\varphi_{173}$}
\end{ptree}
\noindent$\Pi_{33}$:
\begin{ptree}
\AXC{$z = w \to z = -s \to w = -s$\lbl{A6}}
\UIC{$\forall w\,(z = w \to z = -s \to w = -s)$}
\UIC{$\forall z\,\forall w\,(z = w \to z = -s \to w = -s)$}
\AXC{$\varphi_{106}$\lbl{A2}}
\BIC{$\forall w\,(-(-a) = w \to -(-a) = -s \to w = -s)$}
\AXC{$\varphi_{140}$\lbl{A2}}
\BIC{$-(-a) = a \to -(-a) = -s \to a = -s$}
\AXC{$\stackrel{\Pi_{32}}{\vdots}$}\noLine
\UIC{$\varphi_{173}$}
\BIC{$X \subseteq \mathbb{R} \to X \neq \emptyset \to \varphi_{17}$}
\AXC{$(X \subseteq \mathbb{R} \to X \neq \emptyset \to \varphi_{17}) \to \varphi_{41} \to \mathrm{LB}(a,X) \to a \leq -s$\lbl{A1}}
\BIC{$\varphi_{41} \to \mathrm{LB}(a,X) \to a \leq -s$}
\UIC{$\forall a\,(\varphi_{41} \to \mathrm{LB}(a,X) \to a \leq -s)$}
\end{ptree}
\noindent$\Pi_{34}$:
\begin{ptree}
\AXC{$\varphi_{41} \to \forall a\,(\varphi_{41})$\lbl{A4}}
\AXC{$\varphi_{193}$\lbl{A3}}
\AXC{$\stackrel{\Pi_{33}}{\vdots}$}\noLine
\UIC{$\forall a\,(\varphi_{41} \to \mathrm{LB}(a,X) \to a \leq -s)$}
\AXC{$\varphi_{240}$\lbl{A1}}
\BIC{$\varphi_{193} \to (\varphi_{41} \to \forall a\,(\varphi_{41})) \to X \subseteq \mathbb{R} \to X \neq \emptyset \to \varphi_{29}$}
\BIC{$(\varphi_{41} \to \forall a\,(\varphi_{41})) \to X \subseteq \mathbb{R} \to X \neq \emptyset \to \varphi_{29}$}
\BIC{$X \subseteq \mathbb{R} \to X \neq \emptyset \to \varphi_{29}$}
\end{ptree}
\noindent$\Pi_{35}$:
\begin{ptree}
\AXC{$x = x$\lbl{A5}}
\AXC{$x = y \to x = x \to y = x$\lbl{A6}}
\AXC{$(x = y \to x = x \to y = x) \to x = x \to x = y \to y = x$\lbl{A1}}
\BIC{$x = x \to x = y \to y = x$}
\BIC{$x = y \to y = x$}
\UIC{$\forall y\,(x = y \to y = x)$}
\UIC{$\forall x\,\forall y\,(x = y \to y = x)$}
\end{ptree}
\noindent$\Pi_{36}$:
\begin{ptree}
\AXC{$\forall X\,\forall Y\,\forall x\,(X \subseteq Y \to x \in X \to x \in Y)$\lbl{[B5]}}
\AXC{$\varphi_{86}$\lbl{A2}}
\BIC{$\forall Y\,\forall x\,(X \subseteq Y \to x \in X \to x \in Y)$}
\AXC{$\varphi_{87}$\lbl{A2}}
\BIC{$\forall x\,(X \subseteq \mathbb{R} \to x \in X \to x \in \mathbb{R})$}
\AXC{$\varphi_{88}$\lbl{A2}}
\BIC{$X \subseteq \mathbb{R} \to x \in X \to x \in \mathbb{R}$}
\AXC{$(X \subseteq \mathbb{R} \to x \in X \to x \in \mathbb{R}) \to \varphi_{200}$\lbl{A1}}
\BIC{$\varphi_{200}$}
\end{ptree}
\noindent$\Pi_{37}$:
\begin{ptree}
\AXC{$X \subseteq \mathbb{R} \to X \neq \emptyset \to \varphi_{13}$\lbl{A1}}
\AXC{$X \subseteq \mathbb{R} \to X \subseteq \mathbb{R}$\lbl{A1}}
\AXC{$\stackrel{\Pi_{36}}{\vdots}$}\noLine
\UIC{$\varphi_{200}$}
\BIC{$(X \subseteq \mathbb{R} \to X \neq \emptyset \to \varphi_{13}) \to X \subseteq \mathbb{R} \to X \neq \emptyset \to \varphi_{25}$}
\BIC{$X \subseteq \mathbb{R} \to X \neq \emptyset \to \varphi_{25}$}
\end{ptree}
\noindent$\Pi_{38}$:
\begin{ptree}
\AXC{$\stackrel{\Pi_{37}}{\vdots}$}\noLine
\UIC{$X \subseteq \mathbb{R} \to X \neq \emptyset \to \varphi_{25}$}
\AXC{$\forall x \in \mathbb{R}\,(-(-x) = x)$\lbl{[F5]}}
\AXC{$\forall x \in \mathbb{R}\,(-(-x) = x) \to x \in \mathbb{R} \to -(-x) = x$\lbl{A2}}
\BIC{$x \in \mathbb{R} \to -(-x) = x$}
\AXC{$\varphi_{204}$\lbl{A1}}
\BIC{$(X \subseteq \mathbb{R} \to X \neq \emptyset \to \varphi_{25}) \to X \subseteq \mathbb{R} \to X \neq \emptyset \to \varphi_{33}$}
\BIC{$X \subseteq \mathbb{R} \to X \neq \emptyset \to \varphi_{33}$}
\end{ptree}
\noindent$\Pi_{39}$:
\begin{ptree}
\AXC{$\stackrel{\Pi_{35}}{\vdots}$}\noLine
\UIC{$\forall x\,\forall y\,(x = y \to y = x)$}
\AXC{$\varphi_{10}$\lbl{A2}}
\BIC{$\forall y\,(-(-x) = y \to y = -(-x))$}
\AXC{$\varphi_{59}$\lbl{A2}}
\BIC{$-(-x) = x \to x = -(-x)$}
\AXC{$\stackrel{\Pi_{38}}{\vdots}$}\noLine
\UIC{$X \subseteq \mathbb{R} \to X \neq \emptyset \to \varphi_{33}$}
\AXC{$\varphi_{207}$\lbl{A1}}
\BIC{$(-(-x) = x \to x = -(-x)) \to X \subseteq \mathbb{R} \to X \neq \emptyset \to \varphi_{34}$}
\BIC{$X \subseteq \mathbb{R} \to X \neq \emptyset \to \varphi_{34}$}
\end{ptree}
\noindent$\Pi_{40}$:
\begin{ptree}
\AXC{$X \subseteq \mathbb{R} \to X \neq \emptyset \to \varphi_{13}$\lbl{A1}}
\AXC{$\stackrel{\Pi_{39}}{\vdots}$}\noLine
\UIC{$X \subseteq \mathbb{R} \to X \neq \emptyset \to \varphi_{34}$}
\AXC{$(X \subseteq \mathbb{R} \to X \neq \emptyset \to \varphi_{34}) \to (X \subseteq \mathbb{R} \to X \neq \emptyset \to \varphi_{13}) \to \varphi_{182}$\lbl{A1}}
\BIC{$(X \subseteq \mathbb{R} \to X \neq \emptyset \to \varphi_{13}) \to \varphi_{182}$}
\BIC{$\varphi_{182}$}
\end{ptree}
\noindent$\Pi_{41}$:
\begin{ptree}
\AXC{$\stackrel{\Pi_{37}}{\vdots}$}\noLine
\UIC{$X \subseteq \mathbb{R} \to X \neq \emptyset \to \varphi_{25}$}
\AXC{$\forall x \in \mathbb{R}\,(-x \in \mathbb{R})$\lbl{(R1)}}
\AXC{$\forall x \in \mathbb{R}\,(-x \in \mathbb{R}) \to x \in \mathbb{R} \to -x \in \mathbb{R}$\lbl{A2}}
\BIC{$x \in \mathbb{R} \to -x \in \mathbb{R}$}
\AXC{$\varphi_{206}$\lbl{A1}}
\BIC{$(X \subseteq \mathbb{R} \to X \neq \emptyset \to \varphi_{25}) \to X \subseteq \mathbb{R} \to X \neq \emptyset \to \varphi_{43}$}
\BIC{$X \subseteq \mathbb{R} \to X \neq \emptyset \to \varphi_{43}$}
\end{ptree}
\noindent$\Pi_{42}$:
\begin{ptree}
\AXC{$\forall X\,\forall y\,(y \in \ominus X \leftrightarrow y \in \mathbb{R} \land -y \in X)$\lbl{Def$\ominus$}}
\AXC{$\varphi_{144}$\lbl{A2}}
\BIC{$\forall y\,(y \in \ominus X \leftrightarrow y \in \mathbb{R} \land -y \in X)$}
\AXC{$\varphi_{155}$\lbl{A2}}
\BIC{$-x \in \ominus X \leftrightarrow -x \in \mathbb{R} \land -(-x) \in X$}
\AXC{$(-x \in \ominus X \leftrightarrow -x \in \mathbb{R} \land -(-x) \in X) \to \varphi_{224}$\lbl{A1}}
\BIC{$\varphi_{224}$}
\end{ptree}
\noindent$\Pi_{43}$:
\begin{ptree}
\AXC{$z = w \to z \in X \to w \in X$\lbl{A6}}
\UIC{$\forall w\,(z = w \to z \in X \to w \in X)$}
\UIC{$\forall z\,\forall w\,(z = w \to z \in X \to w \in X)$}
\AXC{$\varphi_{37}$\lbl{A2}}
\BIC{$\forall w\,(x = w \to x \in X \to w \in X)$}
\AXC{$\varphi_{82}$\lbl{A2}}
\BIC{$x = -(-x) \to x \in X \to -(-x) \in X$}
\AXC{$\stackrel{\Pi_{40}}{\vdots}$}\noLine
\UIC{$\varphi_{182}$}
\BIC{$X \subseteq \mathbb{R} \to X \neq \emptyset \to \varphi_{50}$}
\AXC{$\stackrel{\Pi_{41}}{\vdots}$}\noLine
\UIC{$X \subseteq \mathbb{R} \to X \neq \emptyset \to \varphi_{43}$}
\AXC{$\stackrel{\Pi_{42}}{\vdots}$}\noLine
\UIC{$\varphi_{224}$}
\BIC{$(X \subseteq \mathbb{R} \to X \neq \emptyset \to \varphi_{50}) \to X \subseteq \mathbb{R} \to X \neq \emptyset \to \varphi_{63}$}
\BIC{$X \subseteq \mathbb{R} \to X \neq \emptyset \to \varphi_{63}$}
\end{ptree}
\noindent$\Pi_{44}$:
\begin{ptree}
\AXC{$\varphi_{89}$\lbl{A4}}
\AXC{$\varphi_{76} \to \varphi_{92}$\lbl{A3}}
\AXC{$\forall y \in \ominus X\,(\neg y < s \to y = s) \to x \in \ominus X \to \neg x < s \to x = s$\lbl{A2}}
\UIC{$\varphi_{76}$}
\AXC{$\varphi_{76} \to (\varphi_{76} \to \varphi_{92}) \to \varphi_{89} \to \varphi_{77}$\lbl{A1}}
\BIC{$(\varphi_{76} \to \varphi_{92}) \to \varphi_{89} \to \varphi_{77}$}
\BIC{$\varphi_{89} \to \varphi_{77}$}
\BIC{$\varphi_{77}$}
\end{ptree}
\noindent$\Pi_{45}$:
\begin{ptree}
\AXC{$\varphi_{90}$\lbl{A4}}
\AXC{$\varphi_{78} \to \varphi_{93}$\lbl{A3}}
\AXC{$\forall x \in \ominus X\,(\neg x < s \to x = s) \to y \in \ominus X \to \neg y < s \to y = s$\lbl{A2}}
\UIC{$\varphi_{78}$}
\AXC{$\varphi_{78} \to (\varphi_{78} \to \varphi_{93}) \to \varphi_{90} \to \varphi_{79}$\lbl{A1}}
\BIC{$(\varphi_{78} \to \varphi_{93}) \to \varphi_{90} \to \varphi_{79}$}
\BIC{$\varphi_{90} \to \varphi_{79}$}
\BIC{$\varphi_{79}$}
\end{ptree}
\noindent$\Pi_{46}$:
\begin{ptree}
\AXC{$\stackrel{\Pi_{45}}{\vdots}$}\noLine
\UIC{$\varphi_{79}$}
\AXC{$\stackrel{\Pi_{25}}{\vdots}$}\noLine
\UIC{$\mathrm{UB}(s,\ominus X) \leftrightarrow s \in \mathbb{R} \land \forall x \in \ominus X\,(x \leq s)$}
\AXC{$\stackrel{\Pi_{27}}{\vdots}$}\noLine
\UIC{$X \subseteq \mathbb{R} \to X \neq \emptyset \to \varphi_{8}$}
\AXC{$(X \subseteq \mathbb{R} \to X \neq \emptyset \to \varphi_{8}) \to \varphi_{228}$\lbl{A1}}
\BIC{$\varphi_{228}$}
\BIC{$\varphi_{79} \to \varphi_{77} \to X \subseteq \mathbb{R} \to X \neq \emptyset \to \varphi_{39}$}
\BIC{$\varphi_{77} \to X \subseteq \mathbb{R} \to X \neq \emptyset \to \varphi_{39}$}
\end{ptree}
\noindent$\Pi_{47}$:
\begin{ptree}
\AXC{$\forall y \in \ominus X\,(y \leq s) \to -x \in \ominus X \to -x \leq s$\lbl{A2}}
\AXC{$\stackrel{\Pi_{44}}{\vdots}$}\noLine
\UIC{$\varphi_{77}$}
\AXC{$\stackrel{\Pi_{46}}{\vdots}$}\noLine
\UIC{$\varphi_{77} \to X \subseteq \mathbb{R} \to X \neq \emptyset \to \varphi_{39}$}
\BIC{$X \subseteq \mathbb{R} \to X \neq \emptyset \to \varphi_{39}$}
\AXC{$(X \subseteq \mathbb{R} \to X \neq \emptyset \to \varphi_{39}) \to \varphi_{189}$\lbl{A1}}
\BIC{$\varphi_{189}$}
\BIC{$X \subseteq \mathbb{R} \to X \neq \emptyset \to \varphi_{65}$}
\end{ptree}
\noindent$\Pi_{48}$:
\begin{ptree}
\AXC{$\stackrel{\Pi_{41}}{\vdots}$}\noLine
\UIC{$X \subseteq \mathbb{R} \to X \neq \emptyset \to \varphi_{43}$}
\AXC{$\forall x\,\forall y\,(x \in \mathbb{R} \to y \in \mathbb{R} \to x \leq y \to -y \leq -x)$\lbl{[F23]}}
\AXC{$\varphi_{169}$\lbl{A2}}
\BIC{$\forall y\,(-x \in \mathbb{R} \to y \in \mathbb{R} \to -x \leq y \to -y \leq -(-x))$}
\AXC{$\varphi_{172}$\lbl{A2}}
\BIC{$-x \in \mathbb{R} \to s \in \mathbb{R} \to -x \leq s \to -s \leq -(-x)$}
\AXC{$\varphi_{237}$\lbl{A1}}
\BIC{$(X \subseteq \mathbb{R} \to X \neq \emptyset \to \varphi_{43}) \to \varphi_{218}$}
\BIC{$\varphi_{218}$}
\end{ptree}
\noindent$\Pi_{49}$:
\begin{ptree}
\AXC{$\stackrel{\Pi_{43}}{\vdots}$}\noLine
\UIC{$X \subseteq \mathbb{R} \to X \neq \emptyset \to \varphi_{63}$}
\AXC{$\stackrel{\Pi_{47}}{\vdots}$}\noLine
\UIC{$X \subseteq \mathbb{R} \to X \neq \emptyset \to \varphi_{65}$}
\AXC{$\varphi_{223}$\lbl{A1}}
\BIC{$(X \subseteq \mathbb{R} \to X \neq \emptyset \to \varphi_{63}) \to X \subseteq \mathbb{R} \to X \neq \emptyset \to \varphi_{18}$}
\BIC{$X \subseteq \mathbb{R} \to X \neq \emptyset \to \varphi_{18}$}
\AXC{$\stackrel{\Pi_{28}}{\vdots}$}\noLine
\UIC{$X \subseteq \mathbb{R} \to \varphi_{38}$}
\AXC{$\stackrel{\Pi_{48}}{\vdots}$}\noLine
\UIC{$\varphi_{218}$}
\BIC{$(X \subseteq \mathbb{R} \to X \neq \emptyset \to \varphi_{18}) \to X \subseteq \mathbb{R} \to X \neq \emptyset \to \varphi_{54}$}
\BIC{$X \subseteq \mathbb{R} \to X \neq \emptyset \to \varphi_{54}$}
\end{ptree}
\noindent$\Pi_{50}$:
\begin{ptree}
\AXC{$\stackrel{\Pi_{37}}{\vdots}$}\noLine
\UIC{$X \subseteq \mathbb{R} \to X \neq \emptyset \to \varphi_{25}$}
\AXC{$\forall x \in \mathbb{R}\,(-(-x) = x)$\lbl{[F5]}}
\AXC{$\forall x \in \mathbb{R}\,(-(-x) = x) \to x \in \mathbb{R} \to -(-x) = x$\lbl{A2}}
\BIC{$x \in \mathbb{R} \to -(-x) = x$}
\AXC{$\varphi_{204}$\lbl{A1}}
\BIC{$(X \subseteq \mathbb{R} \to X \neq \emptyset \to \varphi_{25}) \to X \subseteq \mathbb{R} \to X \neq \emptyset \to \varphi_{33}$}
\BIC{$X \subseteq \mathbb{R} \to X \neq \emptyset \to \varphi_{33}$}
\end{ptree}
\noindent$\Pi_{51}$:
\begin{ptree}
\AXC{$\stackrel{\Pi_{49}}{\vdots}$}\noLine
\UIC{$X \subseteq \mathbb{R} \to X \neq \emptyset \to \varphi_{54}$}
\AXC{$\stackrel{\Pi_{50}}{\vdots}$}\noLine
\UIC{$X \subseteq \mathbb{R} \to X \neq \emptyset \to \varphi_{33}$}
\AXC{$(X \subseteq \mathbb{R} \to X \neq \emptyset \to \varphi_{33}) \to \varphi_{226}$\lbl{A1}}
\BIC{$\varphi_{226}$}
\BIC{$(-(-x) = x \to -s < -(-x) \to -s < x) \to \varphi_{174}$}
\end{ptree}
\noindent$\Pi_{52}$:
\begin{ptree}
\AXC{$z = w \to -s < z \to -s < w$\lbl{A6}}
\UIC{$\forall w\,(z = w \to -s < z \to -s < w)$}
\UIC{$\forall z\,\forall w\,(z = w \to -s < z \to -s < w)$}
\AXC{$\varphi_{132}$\lbl{A2}}
\BIC{$\forall w\,(-(-x) = w \to -s < -(-x) \to -s < w)$}
\AXC{$\varphi_{158}$\lbl{A2}}
\BIC{$-(-x) = x \to -s < -(-x) \to -s < x$}
\AXC{$\stackrel{\Pi_{51}}{\vdots}$}\noLine
\UIC{$(-(-x) = x \to -s < -(-x) \to -s < x) \to \varphi_{174}$}
\BIC{$\varphi_{174}$}
\end{ptree}
\noindent$\Pi_{53}$:
\begin{ptree}
\AXC{$z = w \to -s = z \to -s = w$\lbl{A6}}
\UIC{$\forall w\,(z = w \to -s = z \to -s = w)$}
\UIC{$\forall z\,\forall w\,(z = w \to -s = z \to -s = w)$}
\AXC{$\varphi_{108}$\lbl{A2}}
\BIC{$\forall w\,(-(-x) = w \to -s = -(-x) \to -s = w)$}
\AXC{$\varphi_{142}$\lbl{A2}}
\BIC{$-(-x) = x \to -s = -(-x) \to -s = x$}
\AXC{$\stackrel{\Pi_{52}}{\vdots}$}\noLine
\UIC{$\varphi_{174}$}
\BIC{$X \subseteq \mathbb{R} \to X \neq \emptyset \to \varphi_{19}$}
\AXC{$(X \subseteq \mathbb{R} \to X \neq \emptyset \to \varphi_{19}) \to \varphi_{41} \to x \in X \to -s \leq x$\lbl{A1}}
\BIC{$\varphi_{41} \to x \in X \to -s \leq x$}
\UIC{$\forall x\,(\varphi_{41} \to x \in X \to -s \leq x)$}
\end{ptree}
\noindent$\Pi_{54}$:
\begin{ptree}
\AXC{$\forall x \in \mathbb{R}\,(-x \in \mathbb{R})$\lbl{(R1)}}
\AXC{$\forall x \in \mathbb{R}\,(-x \in \mathbb{R}) \to s \in \mathbb{R} \to -s \in \mathbb{R}$\lbl{A2}}
\BIC{$s \in \mathbb{R} \to -s \in \mathbb{R}$}
\end{ptree}
\noindent$\Pi_{55}$:
\begin{ptree}
\AXC{$\forall a\,\forall X\,(\mathrm{LB}(a,X) \leftrightarrow a \in \mathbb{R} \land \forall x \in X\,(a \leq x))$\lbl{DefLB}}
\AXC{$\varphi_{152}$\lbl{A2}}
\BIC{$\forall X\,(\mathrm{LB}(-s,X) \leftrightarrow -s \in \mathbb{R} \land \forall x \in X\,(-s \leq x))$}
\AXC{$\varphi_{165}$\lbl{A2}}
\BIC{$\mathrm{LB}(-s,X) \leftrightarrow -s \in \mathbb{R} \land \forall x \in X\,(-s \leq x)$}
\end{ptree}
\noindent$\Pi_{56}$:
\begin{ptree}
\AXC{$\stackrel{\Pi_{55}}{\vdots}$}\noLine
\UIC{$\mathrm{LB}(-s,X) \leftrightarrow -s \in \mathbb{R} \land \forall x \in X\,(-s \leq x)$}
\AXC{$(\mathrm{LB}(-s,X) \leftrightarrow -s \in \mathbb{R} \land \forall x \in X\,(-s \leq x)) \to \varphi_{215}$\lbl{A1}}
\BIC{$\varphi_{215}$}
\end{ptree}
\noindent$\Pi_{57}$:
\begin{ptree}
\AXC{$\stackrel{\Pi_{28}}{\vdots}$}\noLine
\UIC{$X \subseteq \mathbb{R} \to \varphi_{38}$}
\AXC{$\stackrel{\Pi_{54}}{\vdots}$}\noLine
\UIC{$s \in \mathbb{R} \to -s \in \mathbb{R}$}
\AXC{$(s \in \mathbb{R} \to -s \in \mathbb{R}) \to (X \subseteq \mathbb{R} \to \varphi_{38}) \to X \subseteq \mathbb{R} \to \varphi_{56}$\lbl{A1}}
\BIC{$(X \subseteq \mathbb{R} \to \varphi_{38}) \to X \subseteq \mathbb{R} \to \varphi_{56}$}
\BIC{$X \subseteq \mathbb{R} \to \varphi_{56}$}
\AXC{$\stackrel{\Pi_{56}}{\vdots}$}\noLine
\UIC{$\varphi_{215}$}
\BIC{$(X \subseteq \mathbb{R} \to X \neq \emptyset \to \varphi_{30}) \to X \subseteq \mathbb{R} \to X \neq \emptyset \to \varphi_{5}$}
\end{ptree}
\noindent$\Pi_{58}$:
\begin{ptree}
\AXC{$\varphi_{41} \to \forall x\,(\varphi_{41})$\lbl{A4}}
\AXC{$\varphi_{194}$\lbl{A3}}
\AXC{$\stackrel{\Pi_{53}}{\vdots}$}\noLine
\UIC{$\forall x\,(\varphi_{41} \to x \in X \to -s \leq x)$}
\AXC{$\varphi_{241}$\lbl{A1}}
\BIC{$\varphi_{194} \to (\varphi_{41} \to \forall x\,(\varphi_{41})) \to X \subseteq \mathbb{R} \to X \neq \emptyset \to \varphi_{30}$}
\BIC{$(\varphi_{41} \to \forall x\,(\varphi_{41})) \to X \subseteq \mathbb{R} \to X \neq \emptyset \to \varphi_{30}$}
\BIC{$X \subseteq \mathbb{R} \to X \neq \emptyset \to \varphi_{30}$}
\AXC{$\stackrel{\Pi_{57}}{\vdots}$}\noLine
\UIC{$(X \subseteq \mathbb{R} \to X \neq \emptyset \to \varphi_{30}) \to X \subseteq \mathbb{R} \to X \neq \emptyset \to \varphi_{5}$}
\BIC{$X \subseteq \mathbb{R} \to X \neq \emptyset \to \varphi_{5}$}
\end{ptree}
\noindent$\Pi_{59}$:
\begin{ptree}
\AXC{$\stackrel{\Pi_{58}}{\vdots}$}\noLine
\UIC{$X \subseteq \mathbb{R} \to X \neq \emptyset \to \varphi_{5}$}
\AXC{$\varphi_{4}$\lbl{DefInf}}
\AXC{$\varphi_{4} \to \forall X\,(\varphi_{22})$\lbl{A2}}
\BIC{$\forall X\,(\varphi_{22})$}
\AXC{$\forall X\,(\varphi_{22}) \to \varphi_{22}$\lbl{A2}}
\BIC{$\varphi_{22}$}
\AXC{$\varphi_{22} \to \varphi_{235}$\lbl{A1}}
\BIC{$\varphi_{235}$}
\BIC{$(X \subseteq \mathbb{R} \to X \neq \emptyset \to \varphi_{29}) \to \varphi_{95} \to \varphi_{94} \to X \subseteq \mathbb{R} \to X \neq \emptyset \to \varphi_{9}$}
\end{ptree}
\noindent$\Pi_{60}$:
\begin{ptree}
\AXC{$\stackrel{\Pi_{1}}{\vdots}$}\noLine
\UIC{$\varphi_{94}$}
\AXC{$\stackrel{\Pi_{2}}{\vdots}$}\noLine
\UIC{$\varphi_{95}$}
\AXC{$\stackrel{\Pi_{34}}{\vdots}$}\noLine
\UIC{$X \subseteq \mathbb{R} \to X \neq \emptyset \to \varphi_{29}$}
\AXC{$\stackrel{\Pi_{59}}{\vdots}$}\noLine
\UIC{$(X \subseteq \mathbb{R} \to X \neq \emptyset \to \varphi_{29}) \to \varphi_{95} \to \varphi_{94} \to X \subseteq \mathbb{R} \to X \neq \emptyset \to \varphi_{9}$}
\BIC{$\varphi_{95} \to \varphi_{94} \to X \subseteq \mathbb{R} \to X \neq \emptyset \to \varphi_{9}$}
\BIC{$\varphi_{94} \to X \subseteq \mathbb{R} \to X \neq \emptyset \to \varphi_{9}$}
\BIC{$X \subseteq \mathbb{R} \to X \neq \emptyset \to \varphi_{9}$}
\end{ptree}
\noindent$\Pi_{61}$:
\begin{ptree}
\AXC{$\forall s\,(\neg \mathrm{IsInf}(s,X)) \to \neg \mathrm{IsInf}(-s,X)$\lbl{A2}}
\AXC{$\stackrel{\Pi_{60}}{\vdots}$}\noLine
\UIC{$X \subseteq \mathbb{R} \to X \neq \emptyset \to \varphi_{9}$}
\AXC{$(X \subseteq \mathbb{R} \to X \neq \emptyset \to \varphi_{9}) \to \varphi_{163}$\lbl{A1}}
\BIC{$\varphi_{163}$}
\BIC{$X \subseteq \mathbb{R} \to X \neq \emptyset \to \varphi_{6}$}
\end{ptree}
\noindent$\Pi_{62}$:
\begin{ptree}
\AXC{$\forall X\,\forall y\,(y \in \ominus X \leftrightarrow y \in \mathbb{R} \land -y \in X)$\lbl{Def$\ominus$}}
\AXC{$\varphi_{144}$\lbl{A2}}
\BIC{$\forall y\,(y \in \ominus X \leftrightarrow y \in \mathbb{R} \land -y \in X)$}
\AXC{$\varphi_{125}$\lbl{A2}}
\BIC{$y \in \ominus X \leftrightarrow y \in \mathbb{R} \land -y \in X$}
\AXC{$X \subseteq \mathbb{R} \to \varphi_{74}$\lbl{A1}}
\AXC{$\varphi_{213}$\lbl{A1}}
\BIC{$(y \in \ominus X \leftrightarrow y \in \mathbb{R} \land -y \in X) \to X \subseteq \mathbb{R} \to \varphi_{70}$}
\BIC{$X \subseteq \mathbb{R} \to \varphi_{70}$}
\end{ptree}
\noindent$\Pi_{63}$:
\begin{ptree}
\AXC{$\forall a\,\forall X\,(\mathrm{LB}(a,X) \leftrightarrow a \in \mathbb{R} \land \forall x \in X\,(a \leq x))$\lbl{DefLB}}
\AXC{$\varphi_{119}$\lbl{A2}}
\BIC{$\forall X\,(\mathrm{LB}(b,X) \leftrightarrow b \in \mathbb{R} \land \forall x \in X\,(b \leq x))$}
\AXC{$\varphi_{112}$\lbl{A2}}
\BIC{$\mathrm{LB}(b,X) \leftrightarrow b \in \mathbb{R} \land \forall x \in X\,(b \leq x)$}
\end{ptree}
\noindent$\Pi_{64}$:
\begin{ptree}
\AXC{$\stackrel{\Pi_{63}}{\vdots}$}\noLine
\UIC{$\mathrm{LB}(b,X) \leftrightarrow b \in \mathbb{R} \land \forall x \in X\,(b \leq x)$}
\AXC{$X \subseteq \mathbb{R} \to \varphi_{3}$\lbl{A1}}
\AXC{$(X \subseteq \mathbb{R} \to \varphi_{3}) \to \varphi_{160}$\lbl{A1}}
\BIC{$\varphi_{160}$}
\BIC{$X \subseteq \mathbb{R} \to \varphi_{21}$}
\AXC{$\varphi_{199}$\lbl{A1}}
\BIC{$(\forall x \in X\,(b \leq x) \to -y \in X \to b \leq -y) \to X \subseteq \mathbb{R} \to \varphi_{48}$}
\end{ptree}
\noindent$\Pi_{65}$:
\begin{ptree}
\AXC{$\forall x \in X\,(b \leq x) \to -y \in X \to b \leq -y$\lbl{A2}}
\AXC{$\stackrel{\Pi_{64}}{\vdots}$}\noLine
\UIC{$(\forall x \in X\,(b \leq x) \to -y \in X \to b \leq -y) \to X \subseteq \mathbb{R} \to \varphi_{48}$}
\BIC{$X \subseteq \mathbb{R} \to \varphi_{48}$}
\end{ptree}
\noindent$\Pi_{66}$:
\begin{ptree}
\AXC{$\stackrel{\Pi_{62}}{\vdots}$}\noLine
\UIC{$X \subseteq \mathbb{R} \to \varphi_{70}$}
\AXC{$\stackrel{\Pi_{65}}{\vdots}$}\noLine
\UIC{$X \subseteq \mathbb{R} \to \varphi_{48}$}
\AXC{$(X \subseteq \mathbb{R} \to \varphi_{48}) \to (X \subseteq \mathbb{R} \to \varphi_{70}) \to X \subseteq \mathbb{R} \to \varphi_{57}$\lbl{A1}}
\BIC{$(X \subseteq \mathbb{R} \to \varphi_{70}) \to X \subseteq \mathbb{R} \to \varphi_{57}$}
\BIC{$X \subseteq \mathbb{R} \to \varphi_{57}$}
\end{ptree}
\noindent$\Pi_{67}$:
\begin{ptree}
\AXC{$\forall X\,\forall y\,(y \in \ominus X \leftrightarrow y \in \mathbb{R} \land -y \in X)$\lbl{Def$\ominus$}}
\AXC{$\varphi_{144}$\lbl{A2}}
\BIC{$\forall y\,(y \in \ominus X \leftrightarrow y \in \mathbb{R} \land -y \in X)$}
\AXC{$\varphi_{125}$\lbl{A2}}
\BIC{$y \in \ominus X \leftrightarrow y \in \mathbb{R} \land -y \in X$}
\AXC{$X \subseteq \mathbb{R} \to \varphi_{74}$\lbl{A1}}
\AXC{$\varphi_{212}$\lbl{A1}}
\BIC{$(y \in \ominus X \leftrightarrow y \in \mathbb{R} \land -y \in X) \to X \subseteq \mathbb{R} \to \varphi_{66}$}
\BIC{$X \subseteq \mathbb{R} \to \varphi_{66}$}
\end{ptree}
\noindent$\Pi_{68}$:
\begin{ptree}
\AXC{$\forall x \in \mathbb{R}\,(-x \in \mathbb{R})$\lbl{(R1)}}
\AXC{$\forall x \in \mathbb{R}\,(-x \in \mathbb{R}) \to y \in \mathbb{R} \to -y \in \mathbb{R}$\lbl{A2}}
\BIC{$y \in \mathbb{R} \to -y \in \mathbb{R}$}
\end{ptree}
\noindent$\Pi_{69}$:
\begin{ptree}
\AXC{$\stackrel{\Pi_{67}}{\vdots}$}\noLine
\UIC{$X \subseteq \mathbb{R} \to \varphi_{66}$}
\AXC{$\stackrel{\Pi_{68}}{\vdots}$}\noLine
\UIC{$y \in \mathbb{R} \to -y \in \mathbb{R}$}
\AXC{$(y \in \mathbb{R} \to -y \in \mathbb{R}) \to (X \subseteq \mathbb{R} \to \varphi_{66}) \to X \subseteq \mathbb{R} \to \varphi_{75}$\lbl{A1}}
\BIC{$(X \subseteq \mathbb{R} \to \varphi_{66}) \to X \subseteq \mathbb{R} \to \varphi_{75}$}
\BIC{$X \subseteq \mathbb{R} \to \varphi_{75}$}
\end{ptree}
\noindent$\Pi_{70}$:
\begin{ptree}
\AXC{$\stackrel{\Pi_{63}}{\vdots}$}\noLine
\UIC{$\mathrm{LB}(b,X) \leftrightarrow b \in \mathbb{R} \land \forall x \in X\,(b \leq x)$}
\AXC{$X \subseteq \mathbb{R} \to \varphi_{3}$\lbl{A1}}
\AXC{$\varphi_{191}$\lbl{A1}}
\BIC{$(\mathrm{LB}(b,X) \leftrightarrow b \in \mathbb{R} \land \forall x \in X\,(b \leq x)) \to \varphi_{49}$}
\BIC{$\varphi_{49}$}
\end{ptree}
\noindent$\Pi_{71}$:
\begin{ptree}
\AXC{$\stackrel{\Pi_{70}}{\vdots}$}\noLine
\UIC{$\varphi_{49}$}
\AXC{$\forall x\,\forall y\,(x \in \mathbb{R} \to y \in \mathbb{R} \to x \leq y \to -y \leq -x)$\lbl{[F23]}}
\AXC{$\varphi_{148}$\lbl{A2}}
\BIC{$\forall y\,(b \in \mathbb{R} \to y \in \mathbb{R} \to b \leq y \to -y \leq -b)$}
\AXC{$\varphi_{164}$\lbl{A2}}
\BIC{$b \in \mathbb{R} \to -y \in \mathbb{R} \to b \leq -y \to -(-y) \leq -b$}
\AXC{$\varphi_{233}$\lbl{A1}}
\BIC{$\varphi_{49} \to (X \subseteq \mathbb{R} \to \varphi_{75}) \to (X \subseteq \mathbb{R} \to \varphi_{57}) \to X \subseteq \mathbb{R} \to X \neq \emptyset \to \varphi_{44}$}
\BIC{$(X \subseteq \mathbb{R} \to \varphi_{75}) \to (X \subseteq \mathbb{R} \to \varphi_{57}) \to X \subseteq \mathbb{R} \to X \neq \emptyset \to \varphi_{44}$}
\end{ptree}
\noindent$\Pi_{72}$:
\begin{ptree}
\AXC{$\stackrel{\Pi_{66}}{\vdots}$}\noLine
\UIC{$X \subseteq \mathbb{R} \to \varphi_{57}$}
\AXC{$\stackrel{\Pi_{69}}{\vdots}$}\noLine
\UIC{$X \subseteq \mathbb{R} \to \varphi_{75}$}
\AXC{$\stackrel{\Pi_{71}}{\vdots}$}\noLine
\UIC{$(X \subseteq \mathbb{R} \to \varphi_{75}) \to (X \subseteq \mathbb{R} \to \varphi_{57}) \to X \subseteq \mathbb{R} \to X \neq \emptyset \to \varphi_{44}$}
\BIC{$(X \subseteq \mathbb{R} \to \varphi_{57}) \to X \subseteq \mathbb{R} \to X \neq \emptyset \to \varphi_{44}$}
\BIC{$X \subseteq \mathbb{R} \to X \neq \emptyset \to \varphi_{44}$}
\end{ptree}
\noindent$\Pi_{73}$:
\begin{ptree}
\AXC{$\stackrel{\Pi_{67}}{\vdots}$}\noLine
\UIC{$X \subseteq \mathbb{R} \to \varphi_{66}$}
\AXC{$\forall x \in \mathbb{R}\,(-(-x) = x)$\lbl{[F5]}}
\AXC{$\forall x \in \mathbb{R}\,(-(-x) = x) \to y \in \mathbb{R} \to -(-y) = y$\lbl{A2}}
\BIC{$y \in \mathbb{R} \to -(-y) = y$}
\AXC{$\varphi_{201}$\lbl{A1}}
\BIC{$(X \subseteq \mathbb{R} \to \varphi_{66}) \to X \subseteq \mathbb{R} \to X \neq \emptyset \to \varphi_{26}$}
\BIC{$X \subseteq \mathbb{R} \to X \neq \emptyset \to \varphi_{26}$}
\end{ptree}
\noindent$\Pi_{74}$:
\begin{ptree}
\AXC{$\stackrel{\Pi_{72}}{\vdots}$}\noLine
\UIC{$X \subseteq \mathbb{R} \to X \neq \emptyset \to \varphi_{44}$}
\AXC{$\stackrel{\Pi_{73}}{\vdots}$}\noLine
\UIC{$X \subseteq \mathbb{R} \to X \neq \emptyset \to \varphi_{26}$}
\AXC{$(X \subseteq \mathbb{R} \to X \neq \emptyset \to \varphi_{26}) \to (X \subseteq \mathbb{R} \to X \neq \emptyset \to \varphi_{44}) \to \varphi_{196}$\lbl{A1}}
\BIC{$(X \subseteq \mathbb{R} \to X \neq \emptyset \to \varphi_{44}) \to \varphi_{196}$}
\BIC{$\varphi_{196}$}
\end{ptree}
\noindent$\Pi_{75}$:
\begin{ptree}
\AXC{$z = w \to z < -b \to w < -b$\lbl{A6}}
\UIC{$\forall w\,(z = w \to z < -b \to w < -b)$}
\UIC{$\forall z\,\forall w\,(z = w \to z < -b \to w < -b)$}
\AXC{$\varphi_{135}$\lbl{A2}}
\BIC{$\forall w\,(-(-y) = w \to -(-y) < -b \to w < -b)$}
\AXC{$\varphi_{159}$\lbl{A2}}
\BIC{$-(-y) = y \to -(-y) < -b \to y < -b$}
\AXC{$\stackrel{\Pi_{74}}{\vdots}$}\noLine
\UIC{$\varphi_{196}$}
\BIC{$(-(-y) = y \to -(-y) = -b \to y = -b) \to X \subseteq \mathbb{R} \to \varphi_{58}$}
\end{ptree}
\noindent$\Pi_{76}$:
\begin{ptree}
\AXC{$z = w \to z = -b \to w = -b$\lbl{A6}}
\UIC{$\forall w\,(z = w \to z = -b \to w = -b)$}
\UIC{$\forall z\,\forall w\,(z = w \to z = -b \to w = -b)$}
\AXC{$\varphi_{109}$\lbl{A2}}
\BIC{$\forall w\,(-(-y) = w \to -(-y) = -b \to w = -b)$}
\AXC{$\varphi_{143}$\lbl{A2}}
\BIC{$-(-y) = y \to -(-y) = -b \to y = -b$}
\AXC{$\stackrel{\Pi_{75}}{\vdots}$}\noLine
\UIC{$(-(-y) = y \to -(-y) = -b \to y = -b) \to X \subseteq \mathbb{R} \to \varphi_{58}$}
\BIC{$X \subseteq \mathbb{R} \to \varphi_{58}$}
\end{ptree}
\noindent$\Pi_{77}$:
\begin{ptree}
\AXC{$\stackrel{\Pi_{76}}{\vdots}$}\noLine
\UIC{$X \subseteq \mathbb{R} \to \varphi_{58}$}
\AXC{$(X \subseteq \mathbb{R} \to \varphi_{58}) \to \varphi_{98}$\lbl{A1}}
\BIC{$\varphi_{98}$}
\UIC{$\forall y\,(\varphi_{98})$}
\AXC{$\forall y\,(\varphi_{98}) \to (\forall y\,(\varphi_{98}) \to \varphi_{14} \to \forall y \in \ominus X\,(y \leq -b)) \to \varphi_{197}$\lbl{A1}}
\BIC{$(\forall y\,(\varphi_{98}) \to \varphi_{14} \to \forall y \in \ominus X\,(y \leq -b)) \to \varphi_{197}$}
\end{ptree}
\noindent$\Pi_{78}$:
\begin{ptree}
\AXC{$\forall y\,(\varphi_{98}) \to \varphi_{14} \to \forall y \in \ominus X\,(y \leq -b)$\lbl{A3}}
\AXC{$\stackrel{\Pi_{77}}{\vdots}$}\noLine
\UIC{$(\forall y\,(\varphi_{98}) \to \varphi_{14} \to \forall y \in \ominus X\,(y \leq -b)) \to \varphi_{197}$}
\BIC{$\varphi_{197}$}
\end{ptree}
\noindent$\Pi_{79}$:
\begin{ptree}
\AXC{$\stackrel{\Pi_{70}}{\vdots}$}\noLine
\UIC{$\varphi_{49}$}
\AXC{$\forall x \in \mathbb{R}\,(-x \in \mathbb{R})$\lbl{(R1)}}
\AXC{$\forall x \in \mathbb{R}\,(-x \in \mathbb{R}) \to b \in \mathbb{R} \to -b \in \mathbb{R}$\lbl{A2}}
\BIC{$b \in \mathbb{R} \to -b \in \mathbb{R}$}
\AXC{$(b \in \mathbb{R} \to -b \in \mathbb{R}) \to \varphi_{49} \to \varphi_{67}$\lbl{A1}}
\BIC{$\varphi_{49} \to \varphi_{67}$}
\BIC{$\varphi_{67}$}
\end{ptree}
\noindent$\Pi_{80}$:
\begin{ptree}
\AXC{$\stackrel{\Pi_{79}}{\vdots}$}\noLine
\UIC{$\varphi_{67}$}
\AXC{$\forall a\,\forall X\,(\mathrm{UB}(a,X) \leftrightarrow a \in \mathbb{R} \land \forall x \in X\,(x \leq a))$\lbl{DefUB}}
\AXC{$\varphi_{153}$\lbl{A2}}
\BIC{$\forall X\,(\mathrm{UB}(-b,X) \leftrightarrow -b \in \mathbb{R} \land \forall x \in X\,(x \leq -b))$}
\AXC{$\varphi_{180}$\lbl{A2}}
\BIC{$\mathrm{UB}(-b,\ominus X) \leftrightarrow -b \in \mathbb{R} \land \forall x \in \ominus X\,(x \leq -b)$}
\AXC{$\varphi_{243}$\lbl{A1}}
\BIC{$\varphi_{67} \to (X \subseteq \mathbb{R} \to X \neq \emptyset \to \varphi_{20}) \to \varphi_{124} \to \varphi_{123} \to X \subseteq \mathbb{R} \to \varphi_{35}$}
\BIC{$(X \subseteq \mathbb{R} \to X \neq \emptyset \to \varphi_{20}) \to \varphi_{124} \to \varphi_{123} \to X \subseteq \mathbb{R} \to \varphi_{35}$}
\end{ptree}
\noindent$\Pi_{81}$:
\begin{ptree}
\AXC{$X \subseteq \mathbb{R} \land X \neq \emptyset \land \exists b\,(\mathrm{LB}(b,X)) \land \mathrm{LB}(b,X) \to \varphi_{14}$\lbl{A4}}
\AXC{$\stackrel{\Pi_{78}}{\vdots}$}\noLine
\UIC{$\varphi_{197}$}
\BIC{$X \subseteq \mathbb{R} \to X \neq \emptyset \to \varphi_{20}$}
\AXC{$\stackrel{\Pi_{80}}{\vdots}$}\noLine
\UIC{$(X \subseteq \mathbb{R} \to X \neq \emptyset \to \varphi_{20}) \to \varphi_{124} \to \varphi_{123} \to X \subseteq \mathbb{R} \to \varphi_{35}$}
\BIC{$\varphi_{124} \to \varphi_{123} \to X \subseteq \mathbb{R} \to \varphi_{35}$}
\end{ptree}
\noindent$\Pi_{82}$:
\begin{ptree}
\AXC{$\varphi_{134}$\lbl{A4}}
\AXC{$\forall y\,(\varphi_{116}) \to \varphi_{139}$\lbl{A3}}
\AXC{$\varphi_{116}$\lbl{A2}}
\UIC{$\forall y\,(\varphi_{116})$}
\AXC{$\forall y\,(\varphi_{116}) \to (\forall y\,(\varphi_{116}) \to \varphi_{139}) \to \varphi_{134} \to \varphi_{124}$\lbl{A1}}
\BIC{$(\forall y\,(\varphi_{116}) \to \varphi_{139}) \to \varphi_{134} \to \varphi_{124}$}
\BIC{$\varphi_{134} \to \varphi_{124}$}
\BIC{$\varphi_{124}$}
\AXC{$\stackrel{\Pi_{81}}{\vdots}$}\noLine
\UIC{$\varphi_{124} \to \varphi_{123} \to X \subseteq \mathbb{R} \to \varphi_{35}$}
\BIC{$\varphi_{123} \to X \subseteq \mathbb{R} \to \varphi_{35}$}
\end{ptree}
\noindent$\Pi_{83}$:
\begin{ptree}
\AXC{$\varphi_{133}$\lbl{A4}}
\AXC{$\forall x\,(\varphi_{115}) \to \varphi_{138}$\lbl{A3}}
\AXC{$\varphi_{115}$\lbl{A2}}
\UIC{$\forall x\,(\varphi_{115})$}
\AXC{$\forall x\,(\varphi_{115}) \to (\forall x\,(\varphi_{115}) \to \varphi_{138}) \to \varphi_{133} \to \varphi_{123}$\lbl{A1}}
\BIC{$(\forall x\,(\varphi_{115}) \to \varphi_{138}) \to \varphi_{133} \to \varphi_{123}$}
\BIC{$\varphi_{133} \to \varphi_{123}$}
\BIC{$\varphi_{123}$}
\AXC{$\stackrel{\Pi_{82}}{\vdots}$}\noLine
\UIC{$\varphi_{123} \to X \subseteq \mathbb{R} \to \varphi_{35}$}
\BIC{$X \subseteq \mathbb{R} \to \varphi_{35}$}
\end{ptree}
\noindent$\Pi_{84}$:
\begin{ptree}
\AXC{$\forall b\,(\neg \mathrm{UB}(b,\ominus X)) \to \neg \mathrm{UB}(-b,\ominus X)$\lbl{A2}}
\AXC{$\stackrel{\Pi_{83}}{\vdots}$}\noLine
\UIC{$X \subseteq \mathbb{R} \to \varphi_{35}$}
\AXC{$(X \subseteq \mathbb{R} \to \varphi_{35}) \to \varphi_{166}$\lbl{A1}}
\BIC{$\varphi_{166}$}
\BIC{$X \subseteq \mathbb{R} \to X \neq \emptyset \to \varphi_{1}$}
\end{ptree}
\noindent$\Pi_{85}$:
\begin{ptree}
\AXC{$\stackrel{\Pi_{84}}{\vdots}$}\noLine
\UIC{$X \subseteq \mathbb{R} \to X \neq \emptyset \to \varphi_{1}$}
\AXC{$(X \subseteq \mathbb{R} \to X \neq \emptyset \to \varphi_{1}) \to \varphi_{47} \to \neg \mathrm{LB}(b,X)$\lbl{A1}}
\BIC{$\varphi_{47} \to \neg \mathrm{LB}(b,X)$}
\UIC{$\forall b\,(\varphi_{47} \to \neg \mathrm{LB}(b,X))$}
\end{ptree}
\noindent$\Pi_{86}$:
\begin{ptree}
\AXC{$\varphi_{190}$\lbl{A3}}
\AXC{$\stackrel{\Pi_{85}}{\vdots}$}\noLine
\UIC{$\forall b\,(\varphi_{47} \to \neg \mathrm{LB}(b,X))$}
\AXC{$\varphi_{11}$\lbl{A1}}
\AXC{$\varphi_{11} \to \forall b\,(\varphi_{47} \to \neg \mathrm{LB}(b,X)) \to \varphi_{190} \to (\varphi_{47} \to \forall b\,(\varphi_{47})) \to \varphi_{40}$\lbl{A1}}
\BIC{$\forall b\,(\varphi_{47} \to \neg \mathrm{LB}(b,X)) \to \varphi_{190} \to (\varphi_{47} \to \forall b\,(\varphi_{47})) \to \varphi_{40}$}
\BIC{$\varphi_{190} \to (\varphi_{47} \to \forall b\,(\varphi_{47})) \to \varphi_{40}$}
\BIC{$(\varphi_{47} \to \forall b\,(\varphi_{47})) \to \varphi_{40}$}
\end{ptree}
\noindent$\Pi_{87}$:
\begin{ptree}
\AXC{$\stackrel{\Pi_{35}}{\vdots}$}\noLine
\UIC{$\forall x\,\forall y\,(x = y \to y = x)$}
\AXC{$\forall x\,\forall y\,(x = y \to y = x) \to \forall y\,(\ominus X = y \to y = \ominus X)$\lbl{A2}}
\BIC{$\forall y\,(\ominus X = y \to y = \ominus X)$}
\end{ptree}
\noindent$\Pi_{88}$:
\begin{ptree}
\AXC{$\stackrel{\Pi_{35}}{\vdots}$}\noLine
\UIC{$\forall x\,\forall y\,(x = y \to y = x)$}
\AXC{$\forall x\,\forall y\,(x = y \to y = x) \to \forall y\,(\emptyset = y \to y = \emptyset)$\lbl{A2}}
\BIC{$\forall y\,(\emptyset = y \to y = \emptyset)$}
\AXC{$\forall y\,(\emptyset = y \to y = \emptyset) \to \emptyset = \ominus X \to \ominus X = \emptyset$\lbl{A2}}
\BIC{$\emptyset = \ominus X \to \ominus X = \emptyset$}
\end{ptree}
\noindent$\Pi_{89}$:
\begin{ptree}
\AXC{$\forall X\,\forall Y\,\forall x\,(X \subseteq Y \to x \in X \to x \in Y)$\lbl{[B5]}}
\AXC{$\varphi_{86}$\lbl{A2}}
\BIC{$\forall Y\,\forall x\,(X \subseteq Y \to x \in X \to x \in Y)$}
\AXC{$\varphi_{87}$\lbl{A2}}
\BIC{$\forall x\,(X \subseteq \mathbb{R} \to x \in X \to x \in \mathbb{R})$}
\AXC{$\varphi_{88}$\lbl{A2}}
\BIC{$X \subseteq \mathbb{R} \to x \in X \to x \in \mathbb{R}$}
\end{ptree}
\noindent$\Pi_{90}$:
\begin{ptree}
\AXC{$X \subseteq \mathbb{R} \to X \subseteq \mathbb{R}$\lbl{A1}}
\AXC{$\stackrel{\Pi_{89}}{\vdots}$}\noLine
\UIC{$X \subseteq \mathbb{R} \to x \in X \to x \in \mathbb{R}$}
\AXC{$(X \subseteq \mathbb{R} \to x \in X \to x \in \mathbb{R}) \to (X \subseteq \mathbb{R} \to X \subseteq \mathbb{R}) \to \varphi_{170}$\lbl{A1}}
\BIC{$(X \subseteq \mathbb{R} \to X \subseteq \mathbb{R}) \to \varphi_{170}$}
\BIC{$\varphi_{170}$}
\end{ptree}
\noindent$\Pi_{91}$:
\begin{ptree}
\AXC{$X \subseteq \mathbb{R} \to X \neq \emptyset \to \exists b\,(\mathrm{LB}(b,X)) \to x \in X \to x \in X$\lbl{A1}}
\AXC{$\stackrel{\Pi_{90}}{\vdots}$}\noLine
\UIC{$\varphi_{170}$}
\BIC{$X \subseteq \mathbb{R} \to X \neq \emptyset \to \exists b\,(\mathrm{LB}(b,X)) \to x \in X \to x \in \mathbb{R}$}
\end{ptree}
\noindent$\Pi_{92}$:
\begin{ptree}
\AXC{$\forall x \in \mathbb{R}\,(-(-x) = x)$\lbl{[F5]}}
\AXC{$\forall x \in \mathbb{R}\,(-(-x) = x) \to x \in \mathbb{R} \to -(-x) = x$\lbl{A2}}
\BIC{$x \in \mathbb{R} \to -(-x) = x$}
\AXC{$(x \in \mathbb{R} \to -(-x) = x) \to \varphi_{176}$\lbl{A1}}
\BIC{$\varphi_{176}$}
\end{ptree}
\noindent$\Pi_{93}$:
\begin{ptree}
\AXC{$\stackrel{\Pi_{91}}{\vdots}$}\noLine
\UIC{$X \subseteq \mathbb{R} \to X \neq \emptyset \to \exists b\,(\mathrm{LB}(b,X)) \to x \in X \to x \in \mathbb{R}$}
\AXC{$\stackrel{\Pi_{92}}{\vdots}$}\noLine
\UIC{$\varphi_{176}$}
\BIC{$\varphi_{60}$}
\AXC{$\varphi_{60} \to (-(-x) = x \to x = -(-x)) \to \varphi_{61}$\lbl{A1}}
\BIC{$(-(-x) = x \to x = -(-x)) \to \varphi_{61}$}
\end{ptree}
\noindent$\Pi_{94}$:
\begin{ptree}
\AXC{$\stackrel{\Pi_{35}}{\vdots}$}\noLine
\UIC{$\forall x\,\forall y\,(x = y \to y = x)$}
\AXC{$\varphi_{10}$\lbl{A2}}
\BIC{$\forall y\,(-(-x) = y \to y = -(-x))$}
\AXC{$\varphi_{59}$\lbl{A2}}
\BIC{$-(-x) = x \to x = -(-x)$}
\AXC{$\stackrel{\Pi_{93}}{\vdots}$}\noLine
\UIC{$(-(-x) = x \to x = -(-x)) \to \varphi_{61}$}
\BIC{$\varphi_{61}$}
\AXC{$\varphi_{61} \to \varphi_{195}$\lbl{A1}}
\BIC{$\varphi_{195}$}
\end{ptree}
\noindent$\Pi_{95}$:
\begin{ptree}
\AXC{$z = w \to z \in X \to w \in X$\lbl{A6}}
\UIC{$\forall w\,(z = w \to z \in X \to w \in X)$}
\UIC{$\forall z\,\forall w\,(z = w \to z \in X \to w \in X)$}
\AXC{$\varphi_{37}$\lbl{A2}}
\BIC{$\forall w\,(x = w \to x \in X \to w \in X)$}
\AXC{$\varphi_{82}$\lbl{A2}}
\BIC{$x = -(-x) \to x \in X \to -(-x) \in X$}
\AXC{$X \subseteq \mathbb{R} \to X \neq \emptyset \to \exists b\,(\mathrm{LB}(b,X)) \to x \in X \to x \in X$\lbl{A1}}
\AXC{$\stackrel{\Pi_{94}}{\vdots}$}\noLine
\UIC{$\varphi_{195}$}
\BIC{$(x = -(-x) \to x \in X \to -(-x) \in X) \to \varphi_{71}$}
\BIC{$\varphi_{71}$}
\end{ptree}
\noindent$\Pi_{96}$:
\begin{ptree}
\AXC{$\forall x \in \mathbb{R}\,(-x \in \mathbb{R})$\lbl{(R1)}}
\AXC{$\forall x \in \mathbb{R}\,(-x \in \mathbb{R}) \to x \in \mathbb{R} \to -x \in \mathbb{R}$\lbl{A2}}
\BIC{$x \in \mathbb{R} \to -x \in \mathbb{R}$}
\AXC{$(x \in \mathbb{R} \to -x \in \mathbb{R}) \to \varphi_{181}$\lbl{A1}}
\BIC{$\varphi_{181}$}
\end{ptree}
\noindent$\Pi_{97}$:
\begin{ptree}
\AXC{$\forall X\,\forall y\,(y \in \ominus X \leftrightarrow y \in \mathbb{R} \land -y \in X)$\lbl{Def$\ominus$}}
\AXC{$\varphi_{144}$\lbl{A2}}
\BIC{$\forall y\,(y \in \ominus X \leftrightarrow y \in \mathbb{R} \land -y \in X)$}
\AXC{$\varphi_{155}$\lbl{A2}}
\BIC{$-x \in \ominus X \leftrightarrow -x \in \mathbb{R} \land -(-x) \in X$}
\AXC{$(-x \in \ominus X \leftrightarrow -x \in \mathbb{R} \land -(-x) \in X) \to \varphi_{211}$\lbl{A1}}
\BIC{$\varphi_{211}$}
\end{ptree}
\noindent$\Pi_{98}$:
\begin{ptree}
\AXC{$\stackrel{\Pi_{95}}{\vdots}$}\noLine
\UIC{$\varphi_{71}$}
\AXC{$\stackrel{\Pi_{91}}{\vdots}$}\noLine
\UIC{$X \subseteq \mathbb{R} \to X \neq \emptyset \to \exists b\,(\mathrm{LB}(b,X)) \to x \in X \to x \in \mathbb{R}$}
\AXC{$\stackrel{\Pi_{96}}{\vdots}$}\noLine
\UIC{$\varphi_{181}$}
\BIC{$X \subseteq \mathbb{R} \to X \neq \emptyset \to \exists b\,(\mathrm{LB}(b,X)) \to x \in X \to -x \in \mathbb{R}$}
\AXC{$\stackrel{\Pi_{97}}{\vdots}$}\noLine
\UIC{$\varphi_{211}$}
\BIC{$\varphi_{185}$}
\BIC{$X \subseteq \mathbb{R} \to X \neq \emptyset \to \exists b\,(\mathrm{LB}(b,X)) \to x \in X \to -x \in \ominus X$}
\end{ptree}
\noindent$\Pi_{99}$:
\begin{ptree}
\AXC{$\forall X\,\forall x \in X\,(X \neq \emptyset)$\lbl{[B2]}}
\AXC{$\forall X\,\forall x \in X\,(X \neq \emptyset) \to \forall x \in \ominus X\,(\ominus X \neq \emptyset)$\lbl{A2}}
\BIC{$\forall x \in \ominus X\,(\ominus X \neq \emptyset)$}
\AXC{$\forall x \in \ominus X\,(\ominus X \neq \emptyset) \to -x \in \ominus X \to \ominus X \neq \emptyset$\lbl{A2}}
\BIC{$-x \in \ominus X \to \ominus X \neq \emptyset$}
\end{ptree}
\noindent$\Pi_{100}$:
\begin{ptree}
\AXC{$\stackrel{\Pi_{98}}{\vdots}$}\noLine
\UIC{$X \subseteq \mathbb{R} \to X \neq \emptyset \to \exists b\,(\mathrm{LB}(b,X)) \to x \in X \to -x \in \ominus X$}
\AXC{$\stackrel{\Pi_{99}}{\vdots}$}\noLine
\UIC{$-x \in \ominus X \to \ominus X \neq \emptyset$}
\AXC{$(-x \in \ominus X \to \ominus X \neq \emptyset) \to \varphi_{186}$\lbl{A1}}
\BIC{$\varphi_{186}$}
\BIC{$X \subseteq \mathbb{R} \to X \neq \emptyset \to \exists b\,(\mathrm{LB}(b,X)) \to x \in X \to \ominus X \neq \emptyset$}
\end{ptree}
\noindent$\Pi_{101}$:
\begin{ptree}
\AXC{$\stackrel{\Pi_{100}}{\vdots}$}\noLine
\UIC{$X \subseteq \mathbb{R} \to X \neq \emptyset \to \exists b\,(\mathrm{LB}(b,X)) \to x \in X \to \ominus X \neq \emptyset$}
\AXC{$\varphi_{187}$\lbl{A1}}
\BIC{$X \subseteq \mathbb{R} \land X \neq \emptyset \land \exists b\,(\mathrm{LB}(b,X)) \land \neg \ominus X \neq \emptyset \to x \notin X$}
\UIC{$\varphi_{96}$}
\end{ptree}
\noindent$\Pi_{102}$:
\begin{ptree}
\AXC{$\forall X\,(X \neq \emptyset \to \exists x\,(x \in X))$\lbl{[B1]}}
\AXC{$\forall X\,(X \neq \emptyset \to \exists x\,(x \in X)) \to X \neq \emptyset \to \exists x\,(x \in X)$\lbl{A2}}
\BIC{$X \neq \emptyset \to \exists x\,(x \in X)$}
\AXC{$(X \neq \emptyset \to \exists x\,(x \in X)) \to \varphi_{127}$\lbl{A1}}
\BIC{$\varphi_{127}$}
\end{ptree}
\noindent$\Pi_{103}$:
\begin{ptree}
\AXC{$\varphi_{175}$\lbl{A4}}
\AXC{$\varphi_{96} \to \varphi_{105}$\lbl{A3}}
\AXC{$\stackrel{\Pi_{101}}{\vdots}$}\noLine
\UIC{$\varphi_{96}$}
\AXC{$X \subseteq \mathbb{R} \to X \neq \emptyset \to X \neq \emptyset$\lbl{A1}}
\AXC{$\stackrel{\Pi_{102}}{\vdots}$}\noLine
\UIC{$\varphi_{127}$}
\BIC{$X \subseteq \mathbb{R} \to X \neq \emptyset \to \exists b\,(\mathrm{LB}(b,X)) \to \exists x\,(x \in X)$}
\AXC{$\varphi_{242}$\lbl{A1}}
\BIC{$\varphi_{96} \to \varphi_{232}$}
\BIC{$\varphi_{232}$}
\BIC{$\varphi_{175} \to X \subseteq \mathbb{R} \to X \neq \emptyset \to \exists b\,(\mathrm{LB}(b,X)) \to \ominus X \neq \emptyset$}
\BIC{$X \subseteq \mathbb{R} \to X \neq \emptyset \to \exists b\,(\mathrm{LB}(b,X)) \to \ominus X \neq \emptyset$}
\end{ptree}
\noindent$\Pi_{104}$:
\begin{ptree}
\AXC{$\stackrel{\Pi_{103}}{\vdots}$}\noLine
\UIC{$X \subseteq \mathbb{R} \to X \neq \emptyset \to \exists b\,(\mathrm{LB}(b,X)) \to \ominus X \neq \emptyset$}
\AXC{$(X \subseteq \mathbb{R} \to X \neq \emptyset \to \exists b\,(\mathrm{LB}(b,X)) \to \ominus X \neq \emptyset) \to \varphi_{149}$\lbl{A1}}
\BIC{$\varphi_{149}$}
\end{ptree}
\noindent$\Pi_{105}$:
\begin{ptree}
\AXC{$\stackrel{\Pi_{88}}{\vdots}$}\noLine
\UIC{$\emptyset = \ominus X \to \ominus X = \emptyset$}
\AXC{$\stackrel{\Pi_{104}}{\vdots}$}\noLine
\UIC{$\varphi_{149}$}
\BIC{$X \subseteq \mathbb{R} \to X \neq \emptyset \to \exists b\,(\mathrm{LB}(b,X)) \to \emptyset \neq \ominus X$}
\AXC{$(X \subseteq \mathbb{R} \to X \neq \emptyset \to \exists b\,(\mathrm{LB}(b,X)) \to \emptyset \neq \ominus X) \to \varphi_{150}$\lbl{A1}}
\BIC{$\varphi_{150}$}
\end{ptree}
\noindent$\Pi_{106}$:
\begin{ptree}
\AXC{$\forall X\,\forall y\,(y \in \ominus X \leftrightarrow y \in \mathbb{R} \land -y \in X)$\lbl{Def$\ominus$}}
\AXC{$\varphi_{144}$\lbl{A2}}
\BIC{$\forall y\,(y \in \ominus X \leftrightarrow y \in \mathbb{R} \land -y \in X)$}
\AXC{$\varphi_{126}$\lbl{A2}}
\BIC{$u \in \ominus X \leftrightarrow u \in \mathbb{R} \land -u \in X$}
\AXC{$\varphi_{167}$\lbl{A1}}
\BIC{$X \subseteq \mathbb{R} \to X \neq \emptyset \to \exists b\,(\mathrm{LB}(b,X)) \to u \in \ominus X \to u \in \mathbb{R}$}
\AXC{$\varphi_{184}$\lbl{A1}}
\BIC{$X \subseteq \mathbb{R} \land X \neq \emptyset \land \exists b\,(\mathrm{LB}(b,X)) \to u \in \ominus X \to u \in \mathbb{R}$}
\UIC{$\varphi_{83}$}
\AXC{$\varphi_{83} \to (\varphi_{83} \to \varphi_{97}) \to \varphi_{110} \to \varphi_{84}$\lbl{A1}}
\BIC{$(\varphi_{83} \to \varphi_{97}) \to \varphi_{110} \to \varphi_{84}$}
\end{ptree}
\noindent$\Pi_{107}$:
\begin{ptree}
\AXC{$\varphi_{110}$\lbl{A4}}
\AXC{$\varphi_{83} \to \varphi_{97}$\lbl{A3}}
\AXC{$\stackrel{\Pi_{106}}{\vdots}$}\noLine
\UIC{$(\varphi_{83} \to \varphi_{97}) \to \varphi_{110} \to \varphi_{84}$}
\BIC{$\varphi_{110} \to \varphi_{84}$}
\BIC{$\varphi_{84}$}
\AXC{$\varphi_{84} \to \varphi_{168}$\lbl{A1}}
\BIC{$\varphi_{168}$}
\end{ptree}
\noindent$\Pi_{108}$:
\begin{ptree}
\AXC{$\forall x\,\forall y\,(x \subseteq y \leftrightarrow \forall u \in x\,(u \in y))$\lbl{Def$\subseteq$}}
\AXC{$\varphi_{120}$\lbl{A2}}
\BIC{$\forall y\,(\ominus X \subseteq y \leftrightarrow \forall u \in \ominus X\,(u \in y))$}
\AXC{$\varphi_{154}$\lbl{A2}}
\BIC{$\ominus X \subseteq \mathbb{R} \leftrightarrow \forall u \in \ominus X\,(u \in \mathbb{R})$}
\AXC{$\stackrel{\Pi_{107}}{\vdots}$}\noLine
\UIC{$\varphi_{168}$}
\BIC{$X \subseteq \mathbb{R} \to X \neq \emptyset \to \exists b\,(\mathrm{LB}(b,X)) \to \ominus X \subseteq \mathbb{R}$}
\end{ptree}
\noindent$\Pi_{109}$:
\begin{ptree}
\AXC{$\stackrel{\Pi_{108}}{\vdots}$}\noLine
\UIC{$X \subseteq \mathbb{R} \to X \neq \emptyset \to \exists b\,(\mathrm{LB}(b,X)) \to \ominus X \subseteq \mathbb{R}$}
\AXC{$\varphi_{31}$\lbl{(R16)}}
\AXC{$\varphi_{31} \to \ominus X \subseteq \mathbb{R} \to \varphi_{64}$\lbl{A2}}
\BIC{$\ominus X \subseteq \mathbb{R} \to \varphi_{64}$}
\AXC{$\varphi_{231}$\lbl{A1}}
\BIC{$(X \subseteq \mathbb{R} \to X \neq \emptyset \to \exists b\,(\mathrm{LB}(b,X)) \to \ominus X \subseteq \mathbb{R}) \to \varphi_{198}$}
\BIC{$\varphi_{198}$}
\end{ptree}
\noindent$\Pi_{110}$:
\begin{ptree}
\AXC{$\stackrel{\Pi_{87}}{\vdots}$}\noLine
\UIC{$\forall y\,(\ominus X = y \to y = \ominus X)$}
\AXC{$\forall y\,(\ominus X = y \to y = \ominus X) \to \ominus X = \emptyset \to \emptyset = \ominus X$\lbl{A2}}
\BIC{$\ominus X = \emptyset \to \emptyset = \ominus X$}
\AXC{$\stackrel{\Pi_{105}}{\vdots}$}\noLine
\UIC{$\varphi_{150}$}
\BIC{$X \subseteq \mathbb{R} \to X \neq \emptyset \to \exists b\,(\mathrm{LB}(b,X)) \to \ominus X \neq \emptyset$}
\AXC{$\stackrel{\Pi_{109}}{\vdots}$}\noLine
\UIC{$\varphi_{198}$}
\BIC{$\varphi_{40} \to X \subseteq \mathbb{R} \to \varphi_{7}$}
\end{ptree}
\noindent$\Pi_{111}$:
\begin{ptree}
\AXC{$\varphi_{47} \to \forall b\,(\varphi_{47})$\lbl{A4}}
\AXC{$\stackrel{\Pi_{86}}{\vdots}$}\noLine
\UIC{$(\varphi_{47} \to \forall b\,(\varphi_{47})) \to \varphi_{40}$}
\BIC{$\varphi_{40}$}
\AXC{$\stackrel{\Pi_{110}}{\vdots}$}\noLine
\UIC{$\varphi_{40} \to X \subseteq \mathbb{R} \to \varphi_{7}$}
\BIC{$X \subseteq \mathbb{R} \to \varphi_{7}$}
\AXC{$(X \subseteq \mathbb{R} \to \varphi_{7}) \to \varphi_{238}$\lbl{A1}}
\BIC{$\varphi_{238}$}
\end{ptree}
\noindent 主木:
\begin{ptree}
\AXC{$\varphi_{27} \to \forall s\,(\varphi_{27})$\lbl{A4}}
\AXC{$\varphi_{192}$\lbl{A3}}
\AXC{$\stackrel{\Pi_{61}}{\vdots}$}\noLine
\UIC{$X \subseteq \mathbb{R} \to X \neq \emptyset \to \varphi_{6}$}
\AXC{$(X \subseteq \mathbb{R} \to X \neq \emptyset \to \varphi_{6}) \to \varphi_{27} \to \neg \mathrm{IsSup}(s,\ominus X)$\lbl{A1}}
\BIC{$\varphi_{27} \to \neg \mathrm{IsSup}(s,\ominus X)$}
\UIC{$\forall s\,(\varphi_{27} \to \neg \mathrm{IsSup}(s,\ominus X))$}
\AXC{$\stackrel{\Pi_{111}}{\vdots}$}\noLine
\UIC{$\varphi_{238}$}
\BIC{$\varphi_{192} \to (\varphi_{27} \to \forall s\,(\varphi_{27})) \to \varphi_{23}$}
\BIC{$(\varphi_{27} \to \forall s\,(\varphi_{27})) \to \varphi_{23}$}
\BIC{$\varphi_{23}$}
\UIC{$\forall X\,(\varphi_{23})$}
\end{ptree}
\begin{abbrevlist}
\item $\varphi_{1} :\equiv \exists b\,(\mathrm{LB}(b,X)) \to \mathrm{LB}(b,X) \to \exists b\,(\mathrm{UB}(b,\ominus X))$
\item $\varphi_{2} :\equiv \forall X\,(\mathrm{IsSup}(s,X) \leftrightarrow \mathrm{UB}(s,X) \land \forall w\,(\mathrm{UB}(w,X) \to s \leq w))$
\item $\varphi_{3} :\equiv X \neq \emptyset \to \exists b\,(\mathrm{LB}(b,X)) \to \mathrm{LB}(b,X) \to \mathrm{LB}(b,X)$
\item $\varphi_{4} :\equiv \forall s\,\forall X\,(\mathrm{IsInf}(s,X) \leftrightarrow \mathrm{LB}(s,X) \land \forall w\,(\mathrm{LB}(w,X) \to w \leq s))$
\item $\varphi_{5} :\equiv \exists b\,(\mathrm{LB}(b,X)) \to \mathrm{IsSup}(s,\ominus X) \to \mathrm{LB}(-s,X)$
\item $\varphi_{6} :\equiv \exists b\,(\mathrm{LB}(b,X)) \to \mathrm{IsSup}(s,\ominus X) \to \exists s\,(\mathrm{IsInf}(s,X))$
\item $\varphi_{7} :\equiv X \neq \emptyset \to \exists b\,(\mathrm{LB}(b,X)) \to \exists s\,(\mathrm{IsSup}(s,\ominus X))$
\item $\varphi_{8} :\equiv \exists b\,(\mathrm{LB}(b,X)) \to \mathrm{IsSup}(s,\ominus X) \to \mathrm{UB}(s,\ominus X)$
\item $\varphi_{9} :\equiv \exists b\,(\mathrm{LB}(b,X)) \to \mathrm{IsSup}(s,\ominus X) \to \mathrm{IsInf}(-s,X)$
\item $\varphi_{10} :\equiv \forall x\,\forall y\,(x = y \to y = x) \to \forall y\,(-(-x) = y \to y = -(-x))$
\item $\varphi_{11} :\equiv X \subseteq \mathbb{R} \to X \neq \emptyset \to \exists b\,(\mathrm{LB}(b,X)) \to \exists b\,(\mathrm{LB}(b,X))$
\item $\varphi_{12} :\equiv \exists b\,(\mathrm{LB}(b,X)) \to \mathrm{IsSup}(s,\ominus X) \to \mathrm{LB}(a,X) \to \mathrm{LB}(a,X)$
\item $\varphi_{13} :\equiv \exists b\,(\mathrm{LB}(b,X)) \to \mathrm{IsSup}(s,\ominus X) \to x \in X \to x \in X$
\item $\varphi_{14} :\equiv \forall y\,(X \subseteq \mathbb{R} \land X \neq \emptyset \land \exists b\,(\mathrm{LB}(b,X)) \land \mathrm{LB}(b,X))$
\item $\varphi_{15} :\equiv \exists b\,(\mathrm{LB}(b,X)) \to \mathrm{IsSup}(s,\ominus X) \to \mathrm{IsSup}(s,\ominus X)$
\item $\varphi_{16} :\equiv \exists b\,(\mathrm{LB}(b,X)) \to \mathrm{IsSup}(s,\ominus X) \to \mathrm{LB}(a,X) \to s \leq -a$
\item $\varphi_{17} :\equiv \exists b\,(\mathrm{LB}(b,X)) \to \mathrm{IsSup}(s,\ominus X) \to \mathrm{LB}(a,X) \to a \leq -s$
\item $\varphi_{18} :\equiv \exists b\,(\mathrm{LB}(b,X)) \to \mathrm{IsSup}(s,\ominus X) \to x \in X \to -x \leq s$
\item $\varphi_{19} :\equiv \exists b\,(\mathrm{LB}(b,X)) \to \mathrm{IsSup}(s,\ominus X) \to x \in X \to -s \leq x$
\item $\varphi_{20} :\equiv \exists b\,(\mathrm{LB}(b,X)) \to \mathrm{LB}(b,X) \to \forall y \in \ominus X\,(y \leq -b)$
\item $\varphi_{21} :\equiv X \neq \emptyset \to \exists b\,(\mathrm{LB}(b,X)) \to \mathrm{LB}(b,X) \to \forall x \in X\,(b \leq x)$
\item $\varphi_{22} :\equiv \mathrm{IsInf}(-s,X) \leftrightarrow \mathrm{LB}(-s,X) \land \forall w\,(\mathrm{LB}(w,X) \to w \leq -s)$
\item $\varphi_{23} :\equiv X \subseteq \mathbb{R} \to X \neq \emptyset \to \exists b\,(\mathrm{LB}(b,X)) \to \exists s\,(\mathrm{IsInf}(s,X))$
\item $\varphi_{24} :\equiv \exists b\,(\mathrm{LB}(b,X)) \to \mathrm{IsSup}(s,\ominus X) \to \mathrm{LB}(a,X) \to a \in \mathbb{R}$
\item $\varphi_{25} :\equiv \exists b\,(\mathrm{LB}(b,X)) \to \mathrm{IsSup}(s,\ominus X) \to x \in X \to x \in \mathbb{R}$
\item $\varphi_{26} :\equiv \exists b\,(\mathrm{LB}(b,X)) \to \mathrm{LB}(b,X) \to y \in \ominus X \to -(-y) = y$
\item $\varphi_{27} :\equiv X \subseteq \mathbb{R} \land X \neq \emptyset \land \exists b\,(\mathrm{LB}(b,X)) \land \neg \exists s\,(\mathrm{IsInf}(s,X))$
\item $\varphi_{28} :\equiv \mathrm{IsSup}(s,\ominus X) \to \mathrm{LB}(a,X) \to \forall w\,(\mathrm{UB}(w,\ominus X) \to s \leq w)$
\item $\varphi_{29} :\equiv \exists b\,(\mathrm{LB}(b,X)) \to \mathrm{IsSup}(s,\ominus X) \to \forall a\,(\mathrm{LB}(a,X) \to a \leq -s)$
\item $\varphi_{30} :\equiv \exists b\,(\mathrm{LB}(b,X)) \to \mathrm{IsSup}(s,\ominus X) \to \forall x \in X\,(-s \leq x)$
\item $\varphi_{31} :\equiv \forall X\,(X \subseteq \mathbb{R} \to X \neq \emptyset \to \exists b\,(\mathrm{UB}(b,X)) \to \exists s\,(\mathrm{IsSup}(s,X)))$
\item $\varphi_{32} :\equiv \exists b\,(\mathrm{LB}(b,X)) \to \mathrm{IsSup}(s,\ominus X) \to \mathrm{LB}(a,X) \to -(-a) = a$
\item $\varphi_{33} :\equiv \exists b\,(\mathrm{LB}(b,X)) \to \mathrm{IsSup}(s,\ominus X) \to x \in X \to -(-x) = x$
\item $\varphi_{34} :\equiv \exists b\,(\mathrm{LB}(b,X)) \to \mathrm{IsSup}(s,\ominus X) \to x \in X \to x = -(-x)$
\item $\varphi_{35} :\equiv X \neq \emptyset \to \exists b\,(\mathrm{LB}(b,X)) \to \mathrm{LB}(b,X) \to \mathrm{UB}(-b,\ominus X)$
\item $\varphi_{36} :\equiv \forall w\,(\mathrm{UB}(w,\ominus X) \to s \leq w) \to \mathrm{UB}(-a,\ominus X) \to s \leq -a$
\item $\varphi_{37} :\equiv \forall z\,\forall w\,(z = w \to z \in X \to w \in X) \to \forall w\,(x = w \to x \in X \to w \in X)$
\item $\varphi_{38} :\equiv X \neq \emptyset \to \exists b\,(\mathrm{LB}(b,X)) \to \mathrm{IsSup}(s,\ominus X) \to s \in \mathbb{R}$
\item $\varphi_{39} :\equiv \exists b\,(\mathrm{LB}(b,X)) \to \mathrm{IsSup}(s,\ominus X) \to \forall y \in \ominus X\,(y \leq s)$
\item $\varphi_{40} :\equiv X \subseteq \mathbb{R} \to X \neq \emptyset \to \exists b\,(\mathrm{LB}(b,X)) \to \exists b\,(\mathrm{UB}(b,\ominus X))$
\item $\varphi_{41} :\equiv X \subseteq \mathbb{R} \land X \neq \emptyset \land \exists b\,(\mathrm{LB}(b,X)) \land \mathrm{IsSup}(s,\ominus X)$
\item $\varphi_{42} :\equiv \exists b\,(\mathrm{LB}(b,X)) \to \mathrm{IsSup}(s,\ominus X) \to \mathrm{LB}(a,X) \to -a \in \mathbb{R}$
\item $\varphi_{43} :\equiv \exists b\,(\mathrm{LB}(b,X)) \to \mathrm{IsSup}(s,\ominus X) \to x \in X \to -x \in \mathbb{R}$
\item $\varphi_{44} :\equiv \exists b\,(\mathrm{LB}(b,X)) \to \mathrm{LB}(b,X) \to y \in \ominus X \to -(-y) \leq -b$
\item $\varphi_{45} :\equiv \exists b\,(\mathrm{LB}(b,X)) \to \mathrm{IsSup}(s,\ominus X) \to \mathrm{LB}(a,X) \to \forall x \in X\,(a \leq x)$
\item $\varphi_{46} :\equiv \mathrm{IsSup}(s,\ominus X) \to \mathrm{LB}(a,X) \to \forall y \in \ominus X\,(y \leq -a)$
\item $\varphi_{47} :\equiv X \subseteq \mathbb{R} \land X \neq \emptyset \land \exists b\,(\mathrm{LB}(b,X)) \land \neg \exists b\,(\mathrm{UB}(b,\ominus X))$
\item $\varphi_{48} :\equiv X \neq \emptyset \to \exists b\,(\mathrm{LB}(b,X)) \to \mathrm{LB}(b,X) \to -y \in X \to b \leq -y$
\item $\varphi_{49} :\equiv X \subseteq \mathbb{R} \to X \neq \emptyset \to \exists b\,(\mathrm{LB}(b,X)) \to \mathrm{LB}(b,X) \to b \in \mathbb{R}$
\item $\varphi_{50} :\equiv \exists b\,(\mathrm{LB}(b,X)) \to \mathrm{IsSup}(s,\ominus X) \to x \in X \to -(-x) \in X$
\item $\varphi_{51} :\equiv \mathrm{IsSup}(s,\ominus X) \to \mathrm{LB}(a,X) \to y \in \ominus X \to -(-y) = y$
\item $\varphi_{52} :\equiv \mathrm{IsSup}(s,\ominus X) \to \mathrm{LB}(a,X) \to \mathrm{UB}(-a,\ominus X) \to s \leq -a$
\item $\varphi_{53} :\equiv \exists b\,(\mathrm{LB}(b,X)) \to \mathrm{IsSup}(s,\ominus X) \to \mathrm{LB}(a,X) \to -(-a) \leq -s$
\item $\varphi_{54} :\equiv \exists b\,(\mathrm{LB}(b,X)) \to \mathrm{IsSup}(s,\ominus X) \to x \in X \to -s \leq -(-x)$
\item $\varphi_{55} :\equiv \mathrm{IsSup}(s,\ominus X) \leftrightarrow \mathrm{UB}(s,\ominus X) \land \forall w\,(\mathrm{UB}(w,\ominus X) \to s \leq w)$
\item $\varphi_{56} :\equiv X \neq \emptyset \to \exists b\,(\mathrm{LB}(b,X)) \to \mathrm{IsSup}(s,\ominus X) \to -s \in \mathbb{R}$
\item $\varphi_{57} :\equiv X \neq \emptyset \to \exists b\,(\mathrm{LB}(b,X)) \to \mathrm{LB}(b,X) \to y \in \ominus X \to b \leq -y$
\item $\varphi_{58} :\equiv X \neq \emptyset \to \exists b\,(\mathrm{LB}(b,X)) \to \mathrm{LB}(b,X) \to y \in \ominus X \to y \leq -b$
\item $\varphi_{59} :\equiv \forall y\,(-(-x) = y \to y = -(-x)) \to -(-x) = x \to x = -(-x)$
\item $\varphi_{60} :\equiv X \subseteq \mathbb{R} \to X \neq \emptyset \to \exists b\,(\mathrm{LB}(b,X)) \to x \in X \to -(-x) = x$
\item $\varphi_{61} :\equiv X \subseteq \mathbb{R} \to X \neq \emptyset \to \exists b\,(\mathrm{LB}(b,X)) \to x \in X \to x = -(-x)$
\item $\varphi_{62} :\equiv \exists b\,(\mathrm{LB}(b,X)) \to \mathrm{IsSup}(s,\ominus X) \to \mathrm{LB}(a,X) \to \mathrm{UB}(-a,\ominus X)$
\item $\varphi_{63} :\equiv \exists b\,(\mathrm{LB}(b,X)) \to \mathrm{IsSup}(s,\ominus X) \to x \in X \to -x \in \ominus X$
\item $\varphi_{64} :\equiv \ominus X \neq \emptyset \to \exists b\,(\mathrm{UB}(b,\ominus X)) \to \exists s\,(\mathrm{IsSup}(s,\ominus X))$
\item $\varphi_{65} :\equiv \exists b\,(\mathrm{LB}(b,X)) \to \mathrm{IsSup}(s,\ominus X) \to -x \in \ominus X \to -x \leq s$
\item $\varphi_{66} :\equiv X \neq \emptyset \to \exists b\,(\mathrm{LB}(b,X)) \to \mathrm{LB}(b,X) \to y \in \ominus X \to y \in \mathbb{R}$
\item $\varphi_{67} :\equiv X \subseteq \mathbb{R} \to X \neq \emptyset \to \exists b\,(\mathrm{LB}(b,X)) \to \mathrm{LB}(b,X) \to -b \in \mathbb{R}$
\item $\varphi_{68} :\equiv \exists b\,(\mathrm{LB}(b,X)) \to \mathrm{IsSup}(s,\ominus X) \to \mathrm{LB}(a,X) \to -y \in X \to a \leq -y$
\item $\varphi_{69} :\equiv \mathrm{IsSup}(s,\ominus X) \to \mathrm{LB}(a,X) \to y \in \ominus X \to -(-y) \leq -a$
\item $\varphi_{70} :\equiv X \neq \emptyset \to \exists b\,(\mathrm{LB}(b,X)) \to \mathrm{LB}(b,X) \to y \in \ominus X \to -y \in X$
\item $\varphi_{71} :\equiv X \subseteq \mathbb{R} \to X \neq \emptyset \to \exists b\,(\mathrm{LB}(b,X)) \to x \in X \to -(-x) \in X$
\item $\varphi_{72} :\equiv \exists b\,(\mathrm{LB}(b,X)) \to \mathrm{IsSup}(s,\ominus X) \to \mathrm{LB}(a,X) \to y \in \ominus X \to a \leq -y$
\item $\varphi_{73} :\equiv \exists b\,(\mathrm{LB}(b,X)) \to \mathrm{IsSup}(s,\ominus X) \to \mathrm{LB}(a,X) \to y \in \ominus X \to y \leq -a$
\item $\varphi_{74} :\equiv X \neq \emptyset \to \exists b\,(\mathrm{LB}(b,X)) \to \mathrm{LB}(b,X) \to y \in \ominus X \to y \in \ominus X$
\item $\varphi_{75} :\equiv X \neq \emptyset \to \exists b\,(\mathrm{LB}(b,X)) \to \mathrm{LB}(b,X) \to y \in \ominus X \to -y \in \mathbb{R}$
\item $\varphi_{76} :\equiv \forall x\,(\forall y \in \ominus X\,(\neg y < s \to y = s) \to x \in \ominus X \to \neg x < s \to x = s)$
\item $\varphi_{77} :\equiv \forall y \in \ominus X\,(\neg y < s \to y = s) \to \forall x \in \ominus X\,(\neg x < s \to x = s)$
\item $\varphi_{78} :\equiv \forall y\,(\forall x \in \ominus X\,(\neg x < s \to x = s) \to y \in \ominus X \to \neg y < s \to y = s)$
\item $\varphi_{79} :\equiv \forall x \in \ominus X\,(\neg x < s \to x = s) \to \forall y \in \ominus X\,(\neg y < s \to y = s)$
\item $\varphi_{80} :\equiv \forall a\,(\mathrm{LB}(a,X) \to \neg a < -s \to a = -s) \to \mathrm{LB}(w,X) \to \neg w < -s \to w = -s$
\item $\varphi_{81} :\equiv \forall w\,(\mathrm{LB}(w,X) \to \neg w < -s \to w = -s) \to \mathrm{LB}(a,X) \to \neg a < -s \to a = -s$
\item $\varphi_{82} :\equiv \forall w\,(x = w \to x \in X \to w \in X) \to x = -(-x) \to x \in X \to -(-x) \in X$
\item $\varphi_{83} :\equiv \forall u\,(X \subseteq \mathbb{R} \land X \neq \emptyset \land \exists b\,(\mathrm{LB}(b,X)) \to u \in \ominus X \to u \in \mathbb{R})$
\item $\varphi_{84} :\equiv X \subseteq \mathbb{R} \to X \neq \emptyset \to \exists b\,(\mathrm{LB}(b,X)) \to \forall u \in \ominus X\,(u \in \mathbb{R})$
\item $\varphi_{85} :\equiv \exists b\,(\mathrm{LB}(b,X)) \to \mathrm{IsSup}(s,\ominus X) \to \mathrm{LB}(a,X) \to y \in \ominus X \to y \in \mathbb{R}$
\item $\varphi_{86} :\equiv \forall X\,\forall Y\,\forall x\,(X \subseteq Y \to x \in X \to x \in Y) \to \forall Y\,\forall x\,(X \subseteq Y \to x \in X \to x \in Y)$
\item $\varphi_{87} :\equiv \forall Y\,\forall x\,(X \subseteq Y \to x \in X \to x \in Y) \to \forall x\,(X \subseteq \mathbb{R} \to x \in X \to x \in \mathbb{R})$
\item $\varphi_{88} :\equiv \forall x\,(X \subseteq \mathbb{R} \to x \in X \to x \in \mathbb{R}) \to X \subseteq \mathbb{R} \to x \in X \to x \in \mathbb{R}$
\item $\varphi_{89} :\equiv \forall y \in \ominus X\,(\neg y < s \to y = s) \to \forall x\,\forall y \in \ominus X\,(\neg y < s \to y = s)$
\item $\varphi_{90} :\equiv \forall x \in \ominus X\,(\neg x < s \to x = s) \to \forall y\,\forall x \in \ominus X\,(\neg x < s \to x = s)$
\item $\varphi_{91} :\equiv \exists b\,(\mathrm{LB}(b,X)) \to \mathrm{IsSup}(s,\ominus X) \to \mathrm{LB}(a,X) \to y \in \ominus X \to -y \in X$
\item $\varphi_{92} :\equiv \forall x\,\forall y \in \ominus X\,(\neg y < s \to y = s) \to \forall x \in \ominus X\,(\neg x < s \to x = s)$
\item $\varphi_{93} :\equiv \forall y\,\forall x \in \ominus X\,(\neg x < s \to x = s) \to \forall y \in \ominus X\,(\neg y < s \to y = s)$
\item $\varphi_{94} :\equiv \forall a\,(\mathrm{LB}(a,X) \to \neg a < -s \to a = -s) \to \forall w\,(\mathrm{LB}(w,X) \to \neg w < -s \to w = -s)$
\item $\varphi_{95} :\equiv \forall w\,(\mathrm{LB}(w,X) \to \neg w < -s \to w = -s) \to \forall a\,(\mathrm{LB}(a,X) \to \neg a < -s \to a = -s)$
\item $\varphi_{96} :\equiv \forall x\,(X \subseteq \mathbb{R} \land X \neq \emptyset \land \exists b\,(\mathrm{LB}(b,X)) \land \neg \ominus X \neq \emptyset \to x \notin X)$
\item $\varphi_{97} :\equiv \forall u\,(X \subseteq \mathbb{R} \land X \neq \emptyset \land \exists b\,(\mathrm{LB}(b,X))) \to \forall u \in \ominus X\,(u \in \mathbb{R})$
\item $\varphi_{98} :\equiv X \subseteq \mathbb{R} \land X \neq \emptyset \land \exists b\,(\mathrm{LB}(b,X)) \land \mathrm{LB}(b,X) \to y \in \ominus X \to y \leq -b$
\item $\varphi_{99} :\equiv \exists b\,(\mathrm{LB}(b,X)) \to \mathrm{IsSup}(s,\ominus X) \to \mathrm{LB}(a,X) \to y \in \ominus X \to y \in \ominus X$
\item $\varphi_{100} :\equiv \exists b\,(\mathrm{LB}(b,X)) \to \mathrm{IsSup}(s,\ominus X) \to \mathrm{LB}(a,X) \to y \in \ominus X \to -y \in \mathbb{R}$
\item $\varphi_{101} :\equiv \forall a\,(\mathrm{LB}(a,X) \to \neg a < -s \to a = -s) \to \forall w\,\forall a\,(\mathrm{LB}(a,X) \to \neg a < -s \to a = -s)$
\item $\varphi_{102} :\equiv \forall w\,(\mathrm{LB}(w,X) \to \neg w < -s \to w = -s) \to \forall a\,\forall w\,(\mathrm{LB}(w,X) \to \neg w < -s \to w = -s)$
\item $\varphi_{103} :\equiv \forall w\,\forall a\,(\mathrm{LB}(a,X) \to \neg a < -s \to a = -s) \to \forall w\,(\mathrm{LB}(w,X) \to \neg w < -s \to w = -s)$
\item $\varphi_{104} :\equiv \forall a\,\forall w\,(\mathrm{LB}(w,X) \to \neg w < -s \to w = -s) \to \forall a\,(\mathrm{LB}(a,X) \to \neg a < -s \to a = -s)$
\item $\varphi_{105} :\equiv \forall x\,(X \subseteq \mathbb{R} \land X \neq \emptyset \land \exists b\,(\mathrm{LB}(b,X)) \land \neg \ominus X \neq \emptyset) \to \forall x\,(x \notin X)$
\item $\varphi_{106} :\equiv \forall z\,\forall w\,(z = w \to z = -s \to w = -s) \to \forall w\,(-(-a) = w \to -(-a) = -s \to w = -s)$
\item $\varphi_{107} :\equiv \forall z\,\forall w\,(z = w \to z = -a \to w = -a) \to \forall w\,(-(-y) = w \to -(-y) = -a \to w = -a)$
\item $\varphi_{108} :\equiv \forall z\,\forall w\,(z = w \to -s = z \to -s = w) \to \forall w\,(-(-x) = w \to -s = -(-x) \to -s = w)$
\item $\varphi_{109} :\equiv \forall z\,\forall w\,(z = w \to z = -b \to w = -b) \to \forall w\,(-(-y) = w \to -(-y) = -b \to w = -b)$
\item $\varphi_{110} :\equiv X \subseteq \mathbb{R} \land X \neq \emptyset \land \exists b\,(\mathrm{LB}(b,X)) \to \forall u\,(X \subseteq \mathbb{R} \land X \neq \emptyset \land \exists b\,(\mathrm{LB}(b,X)))$
\item $\varphi_{111} :\equiv \forall X\,(\mathrm{LB}(a,X) \leftrightarrow a \in \mathbb{R} \land \forall x \in X\,(a \leq x)) \to (\mathrm{LB}(a,X) \leftrightarrow a \in \mathbb{R} \land \forall x \in X\,(a \leq x))$
\item $\varphi_{112} :\equiv \forall X\,(\mathrm{LB}(b,X) \leftrightarrow b \in \mathbb{R} \land \forall x \in X\,(b \leq x)) \to (\mathrm{LB}(b,X) \leftrightarrow b \in \mathbb{R} \land \forall x \in X\,(b \leq x))$
\item $\varphi_{113} :\equiv \forall y \in \ominus X\,(\neg y < -a \to y = -a) \to x \in \ominus X \to \neg x < -a \to x = -a$
\item $\varphi_{114} :\equiv \forall x \in \ominus X\,(\neg x < -a \to x = -a) \to y \in \ominus X \to \neg y < -a \to y = -a$
\item $\varphi_{115} :\equiv \forall y \in \ominus X\,(\neg y < -b \to y = -b) \to x \in \ominus X \to \neg x < -b \to x = -b$
\item $\varphi_{116} :\equiv \forall x \in \ominus X\,(\neg x < -b \to x = -b) \to y \in \ominus X \to \neg y < -b \to y = -b$
\item $\varphi_{117} :\equiv \forall a\,\forall X\,(\mathrm{LB}(a,X) \leftrightarrow a \in \mathbb{R} \land \forall x \in X\,(a \leq x)) \to \forall X\,(\mathrm{LB}(a,X) \leftrightarrow a \in \mathbb{R} \land \forall x \in X\,(a \leq x))$
\item $\varphi_{118} :\equiv \forall a\,\forall X\,(\mathrm{UB}(a,X) \leftrightarrow a \in \mathbb{R} \land \forall x \in X\,(x \leq a)) \to \forall X\,(\mathrm{UB}(s,X) \leftrightarrow s \in \mathbb{R} \land \forall x \in X\,(x \leq s))$
\item $\varphi_{119} :\equiv \forall a\,\forall X\,(\mathrm{LB}(a,X) \leftrightarrow a \in \mathbb{R} \land \forall x \in X\,(a \leq x)) \to \forall X\,(\mathrm{LB}(b,X) \leftrightarrow b \in \mathbb{R} \land \forall x \in X\,(b \leq x))$
\item $\varphi_{120} :\equiv \forall x\,\forall y\,(x \subseteq y \leftrightarrow \forall u \in x\,(u \in y)) \to \forall y\,(\ominus X \subseteq y \leftrightarrow \forall u \in \ominus X\,(u \in y))$
\item $\varphi_{121} :\equiv \forall y \in \ominus X\,(\neg y < -a \to y = -a) \to \forall x \in \ominus X\,(\neg x < -a \to x = -a)$
\item $\varphi_{122} :\equiv \forall x \in \ominus X\,(\neg x < -a \to x = -a) \to \forall y \in \ominus X\,(\neg y < -a \to y = -a)$
\item $\varphi_{123} :\equiv \forall y \in \ominus X\,(\neg y < -b \to y = -b) \to \forall x \in \ominus X\,(\neg x < -b \to x = -b)$
\item $\varphi_{124} :\equiv \forall x \in \ominus X\,(\neg x < -b \to x = -b) \to \forall y \in \ominus X\,(\neg y < -b \to y = -b)$
\item $\varphi_{125} :\equiv \forall y\,(y \in \ominus X \leftrightarrow y \in \mathbb{R} \land -y \in X) \to (y \in \ominus X \leftrightarrow y \in \mathbb{R} \land -y \in X)$
\item $\varphi_{126} :\equiv \forall y\,(y \in \ominus X \leftrightarrow y \in \mathbb{R} \land -y \in X) \to (u \in \ominus X \leftrightarrow u \in \mathbb{R} \land -u \in X)$
\item $\varphi_{127} :\equiv (X \subseteq \mathbb{R} \to X \neq \emptyset \to X \neq \emptyset) \to X \subseteq \mathbb{R} \to X \neq \emptyset \to \exists b\,(\mathrm{LB}(b,X)) \to \exists x\,(x \in X)$
\item $\varphi_{128} :\equiv \forall z\,\forall w\,(z = w \to z < -s \to w < -s) \to \forall w\,(-(-a) = w \to -(-a) < -s \to w < -s)$
\item $\varphi_{129} :\equiv \forall y \in \ominus X\,(\neg y < -a \to y = -a) \to \forall x\,\forall y \in \ominus X\,(\neg y < -a \to y = -a)$
\item $\varphi_{130} :\equiv \forall x \in \ominus X\,(\neg x < -a \to x = -a) \to \forall y\,\forall x \in \ominus X\,(\neg x < -a \to x = -a)$
\item $\varphi_{131} :\equiv \forall z\,\forall w\,(z = w \to z < -a \to w < -a) \to \forall w\,(-(-y) = w \to -(-y) < -a \to w < -a)$
\item $\varphi_{132} :\equiv \forall z\,\forall w\,(z = w \to -s < z \to -s < w) \to \forall w\,(-(-x) = w \to -s < -(-x) \to -s < w)$
\item $\varphi_{133} :\equiv \forall y \in \ominus X\,(\neg y < -b \to y = -b) \to \forall x\,\forall y \in \ominus X\,(\neg y < -b \to y = -b)$
\item $\varphi_{134} :\equiv \forall x \in \ominus X\,(\neg x < -b \to x = -b) \to \forall y\,\forall x \in \ominus X\,(\neg x < -b \to x = -b)$
\item $\varphi_{135} :\equiv \forall z\,\forall w\,(z = w \to z < -b \to w < -b) \to \forall w\,(-(-y) = w \to -(-y) < -b \to w < -b)$
\item $\varphi_{136} :\equiv \forall x\,\forall y \in \ominus X\,(\neg y < -a \to y = -a) \to \forall x \in \ominus X\,(\neg x < -a \to x = -a)$
\item $\varphi_{137} :\equiv \forall y\,\forall x \in \ominus X\,(\neg x < -a \to x = -a) \to \forall y \in \ominus X\,(\neg y < -a \to y = -a)$
\item $\varphi_{138} :\equiv \forall x\,\forall y \in \ominus X\,(\neg y < -b \to y = -b) \to \forall x \in \ominus X\,(\neg x < -b \to x = -b)$
\item $\varphi_{139} :\equiv \forall y\,\forall x \in \ominus X\,(\neg x < -b \to x = -b) \to \forall y \in \ominus X\,(\neg y < -b \to y = -b)$
\item $\varphi_{140} :\equiv \forall w\,(-(-a) = w \to -(-a) = -s \to w = -s) \to -(-a) = a \to -(-a) = -s \to a = -s$
\item $\varphi_{141} :\equiv \forall w\,(-(-y) = w \to -(-y) = -a \to w = -a) \to -(-y) = y \to -(-y) = -a \to y = -a$
\item $\varphi_{142} :\equiv \forall w\,(-(-x) = w \to -s = -(-x) \to -s = w) \to -(-x) = x \to -s = -(-x) \to -s = x$
\item $\varphi_{143} :\equiv \forall w\,(-(-y) = w \to -(-y) = -b \to w = -b) \to -(-y) = y \to -(-y) = -b \to y = -b$
\item $\varphi_{144} :\equiv \forall X\,\forall y\,(y \in \ominus X \leftrightarrow y \in \mathbb{R} \land -y \in X) \to \forall y\,(y \in \ominus X \leftrightarrow y \in \mathbb{R} \land -y \in X)$
\item $\varphi_{145} :\equiv \forall X\,(\mathrm{UB}(s,X) \leftrightarrow s \in \mathbb{R} \land \forall x \in X\,(x \leq s)) \to (\mathrm{UB}(s,\ominus X) \leftrightarrow s \in \mathbb{R} \land \forall x \in \ominus X\,(x \leq s))$
\item $\varphi_{146} :\equiv \forall x\,\forall y\,(x \in \mathbb{R} \to y \in \mathbb{R} \to x \leq y \to -y \leq -x) \to \forall y\,(a \in \mathbb{R} \to y \in \mathbb{R} \to a \leq y \to -y \leq -a)$
\item $\varphi_{147} :\equiv \forall x\,\forall y\,(x \in \mathbb{R} \to y \in \mathbb{R} \to x \leq y \to -y \leq -x) \to \forall y\,(s \in \mathbb{R} \to y \in \mathbb{R} \to s \leq y \to -y \leq -s)$
\item $\varphi_{148} :\equiv \forall x\,\forall y\,(x \in \mathbb{R} \to y \in \mathbb{R} \to x \leq y \to -y \leq -x) \to \forall y\,(b \in \mathbb{R} \to y \in \mathbb{R} \to b \leq y \to -y \leq -b)$
\item $\varphi_{149} :\equiv (\emptyset = \ominus X \to \ominus X = \emptyset) \to X \subseteq \mathbb{R} \to X \neq \emptyset \to \exists b\,(\mathrm{LB}(b,X)) \to \emptyset \neq \ominus X$
\item $\varphi_{150} :\equiv (\ominus X = \emptyset \to \emptyset = \ominus X) \to X \subseteq \mathbb{R} \to X \neq \emptyset \to \exists b\,(\mathrm{LB}(b,X)) \to \ominus X \neq \emptyset$
\item $\varphi_{151} :\equiv \forall a\,\forall X\,(\mathrm{UB}(a,X) \leftrightarrow a \in \mathbb{R} \land \forall x \in X\,(x \leq a)) \to \forall X\,(\mathrm{UB}(-a,X) \leftrightarrow -a \in \mathbb{R} \land \forall x \in X\,(x \leq -a))$
\item $\varphi_{152} :\equiv \forall a\,\forall X\,(\mathrm{LB}(a,X) \leftrightarrow a \in \mathbb{R} \land \forall x \in X\,(a \leq x)) \to \forall X\,(\mathrm{LB}(-s,X) \leftrightarrow -s \in \mathbb{R} \land \forall x \in X\,(-s \leq x))$
\item $\varphi_{153} :\equiv \forall a\,\forall X\,(\mathrm{UB}(a,X) \leftrightarrow a \in \mathbb{R} \land \forall x \in X\,(x \leq a)) \to \forall X\,(\mathrm{UB}(-b,X) \leftrightarrow -b \in \mathbb{R} \land \forall x \in X\,(x \leq -b))$
\item $\varphi_{154} :\equiv \forall y\,(\ominus X \subseteq y \leftrightarrow \forall u \in \ominus X\,(u \in y)) \to (\ominus X \subseteq \mathbb{R} \leftrightarrow \forall u \in \ominus X\,(u \in \mathbb{R}))$
\item $\varphi_{155} :\equiv \forall y\,(y \in \ominus X \leftrightarrow y \in \mathbb{R} \land -y \in X) \to (-x \in \ominus X \leftrightarrow -x \in \mathbb{R} \land -(-x) \in X)$
\item $\varphi_{156} :\equiv \forall w\,(-(-a) = w \to -(-a) < -s \to w < -s) \to -(-a) = a \to -(-a) < -s \to a < -s$
\item $\varphi_{157} :\equiv \forall w\,(-(-y) = w \to -(-y) < -a \to w < -a) \to -(-y) = y \to -(-y) < -a \to y < -a$
\item $\varphi_{158} :\equiv \forall w\,(-(-x) = w \to -s < -(-x) \to -s < w) \to -(-x) = x \to -s < -(-x) \to -s < x$
\item $\varphi_{159} :\equiv \forall w\,(-(-y) = w \to -(-y) < -b \to w < -b) \to -(-y) = y \to -(-y) < -b \to y < -b$
\item $\varphi_{160} :\equiv (\mathrm{LB}(b,X) \leftrightarrow b \in \mathbb{R} \land \forall x \in X\,(b \leq x)) \to X \subseteq \mathbb{R} \to \varphi_{21}$
\item $\varphi_{161} :\equiv \forall y\,(a \in \mathbb{R} \to y \in \mathbb{R} \to a \leq y \to -y \leq -a) \to a \in \mathbb{R} \to -y \in \mathbb{R} \to a \leq -y \to -(-y) \leq -a$
\item $\varphi_{162} :\equiv \forall y\,(s \in \mathbb{R} \to y \in \mathbb{R} \to s \leq y \to -y \leq -s) \to s \in \mathbb{R} \to -a \in \mathbb{R} \to s \leq -a \to -(-a) \leq -s$
\item $\varphi_{163} :\equiv (\forall s\,(\neg \mathrm{IsInf}(s,X)) \to \neg \mathrm{IsInf}(-s,X)) \to X \subseteq \mathbb{R} \to X \neq \emptyset \to \varphi_{6}$
\item $\varphi_{164} :\equiv \forall y\,(b \in \mathbb{R} \to y \in \mathbb{R} \to b \leq y \to -y \leq -b) \to b \in \mathbb{R} \to -y \in \mathbb{R} \to b \leq -y \to -(-y) \leq -b$
\item $\varphi_{165} :\equiv \forall X\,(\mathrm{LB}(-s,X) \leftrightarrow -s \in \mathbb{R} \land \forall x \in X\,(-s \leq x)) \to (\mathrm{LB}(-s,X) \leftrightarrow -s \in \mathbb{R} \land \forall x \in X\,(-s \leq x))$
\item $\varphi_{166} :\equiv (\forall b\,(\neg \mathrm{UB}(b,\ominus X)) \to \neg \mathrm{UB}(-b,\ominus X)) \to X \subseteq \mathbb{R} \to X \neq \emptyset \to \varphi_{1}$
\item $\varphi_{167} :\equiv (u \in \ominus X \leftrightarrow u \in \mathbb{R} \land -u \in X) \to X \subseteq \mathbb{R} \to X \neq \emptyset \to \exists b\,(\mathrm{LB}(b,X)) \to u \in \ominus X \to u \in \mathbb{R}$
\item $\varphi_{168} :\equiv (\ominus X \subseteq \mathbb{R} \leftrightarrow \forall u \in \ominus X\,(u \in \mathbb{R})) \to X \subseteq \mathbb{R} \to X \neq \emptyset \to \exists b\,(\mathrm{LB}(b,X)) \to \ominus X \subseteq \mathbb{R}$
\item $\varphi_{169} :\equiv \forall x\,\forall y\,(x \in \mathbb{R} \to y \in \mathbb{R} \to x \leq y \to -y \leq -x) \to \forall y\,(-x \in \mathbb{R} \to y \in \mathbb{R} \to -x \leq y \to -y \leq -(-x))$
\item $\varphi_{170} :\equiv (X \subseteq \mathbb{R} \to X \neq \emptyset \to \exists b\,(\mathrm{LB}(b,X)) \to x \in X \to x \in X) \to X \subseteq \mathbb{R} \to X \neq \emptyset \to \exists b\,(\mathrm{LB}(b,X)) \to x \in X \to x \in \mathbb{R}$
\item $\varphi_{171} :\equiv (\mathrm{LB}(a,X) \leftrightarrow a \in \mathbb{R} \land \forall x \in X\,(a \leq x)) \to X \subseteq \mathbb{R} \to X \neq \emptyset \to \varphi_{24}$
\item $\varphi_{172} :\equiv \forall y\,(-x \in \mathbb{R} \to y \in \mathbb{R} \to -x \leq y \to -y \leq -(-x)) \to -x \in \mathbb{R} \to s \in \mathbb{R} \to -x \leq s \to -s \leq -(-x)$
\item $\varphi_{173} :\equiv (-(-a) = a \to -(-a) = -s \to a = -s) \to X \subseteq \mathbb{R} \to X \neq \emptyset \to \varphi_{17}$
\item $\varphi_{174} :\equiv (-(-x) = x \to -s = -(-x) \to -s = x) \to X \subseteq \mathbb{R} \to X \neq \emptyset \to \varphi_{19}$
\item $\varphi_{175} :\equiv X \subseteq \mathbb{R} \land X \neq \emptyset \land \exists b\,(\mathrm{LB}(b,X)) \land \neg \ominus X \neq \emptyset \to \forall x\,(X \subseteq \mathbb{R} \land X \neq \emptyset \land \exists b\,(\mathrm{LB}(b,X)) \land \neg \ominus X \neq \emptyset)$
\item $\varphi_{176} :\equiv (X \subseteq \mathbb{R} \to X \neq \emptyset \to \exists b\,(\mathrm{LB}(b,X)) \to x \in X \to x \in \mathbb{R}) \to \varphi_{60}$
\item $\varphi_{177} :\equiv (\mathrm{LB}(a,X) \leftrightarrow a \in \mathbb{R} \land \forall x \in X\,(a \leq x)) \to X \subseteq \mathbb{R} \to X \neq \emptyset \to \varphi_{45}$
\item $\varphi_{178} :\equiv \forall X\,(\mathrm{UB}(-a,X) \leftrightarrow -a \in \mathbb{R} \land \forall x \in X\,(x \leq -a)) \to (\mathrm{UB}(-a,\ominus X) \leftrightarrow -a \in \mathbb{R} \land \forall x \in \ominus X\,(x \leq -a))$
\item $\varphi_{179} :\equiv (\mathrm{UB}(s,\ominus X) \leftrightarrow s \in \mathbb{R} \land \forall x \in \ominus X\,(x \leq s)) \to X \subseteq \mathbb{R} \to \varphi_{38}$
\item $\varphi_{180} :\equiv \forall X\,(\mathrm{UB}(-b,X) \leftrightarrow -b \in \mathbb{R} \land \forall x \in X\,(x \leq -b)) \to (\mathrm{UB}(-b,\ominus X) \leftrightarrow -b \in \mathbb{R} \land \forall x \in \ominus X\,(x \leq -b))$
\item $\varphi_{181} :\equiv (X \subseteq \mathbb{R} \to X \neq \emptyset \to \exists b\,(\mathrm{LB}(b,X)) \to x \in X \to x \in \mathbb{R}) \to X \subseteq \mathbb{R} \to X \neq \emptyset \to \exists b\,(\mathrm{LB}(b,X)) \to x \in X \to -x \in \mathbb{R}$
\item $\varphi_{182} :\equiv (x = -(-x) \to x \in X \to -(-x) \in X) \to X \subseteq \mathbb{R} \to X \neq \emptyset \to \varphi_{50}$
\item $\varphi_{183} :\equiv (\forall x \in X\,(a \leq x) \to -y \in X \to a \leq -y) \to X \subseteq \mathbb{R} \to X \neq \emptyset \to \varphi_{68}$
\item $\varphi_{184} :\equiv (X \subseteq \mathbb{R} \to X \neq \emptyset \to \exists b\,(\mathrm{LB}(b,X)) \to u \in \ominus X \to u \in \mathbb{R}) \to X \subseteq \mathbb{R} \land X \neq \emptyset \land \exists b\,(\mathrm{LB}(b,X)) \to u \in \ominus X \to u \in \mathbb{R}$
\item $\varphi_{185} :\equiv \varphi_{71} \to X \subseteq \mathbb{R} \to X \neq \emptyset \to \exists b\,(\mathrm{LB}(b,X)) \to x \in X \to -x \in \ominus X$
\item $\varphi_{186} :\equiv (X \subseteq \mathbb{R} \to X \neq \emptyset \to \exists b\,(\mathrm{LB}(b,X)) \to x \in X \to -x \in \ominus X) \to X \subseteq \mathbb{R} \to X \neq \emptyset \to \exists b\,(\mathrm{LB}(b,X)) \to x \in X \to \ominus X \neq \emptyset$
\item $\varphi_{187} :\equiv (X \subseteq \mathbb{R} \to X \neq \emptyset \to \exists b\,(\mathrm{LB}(b,X)) \to x \in X \to \ominus X \neq \emptyset) \to X \subseteq \mathbb{R} \land X \neq \emptyset \land \exists b\,(\mathrm{LB}(b,X)) \land \neg \ominus X \neq \emptyset \to x \notin X$
\item $\varphi_{188} :\equiv (-(-y) = y \to -(-y) = -a \to y = -a) \to X \subseteq \mathbb{R} \to X \neq \emptyset \to \varphi_{73}$
\item $\varphi_{189} :\equiv (\forall y \in \ominus X\,(y \leq s) \to -x \in \ominus X \to -x \leq s) \to X \subseteq \mathbb{R} \to X \neq \emptyset \to \varphi_{65}$
\item $\varphi_{190} :\equiv \forall b\,(\varphi_{47} \to \neg \mathrm{LB}(b,X)) \to \forall b\,(\varphi_{47}) \to \forall b\,(\neg \mathrm{LB}(b,X))$
\item $\varphi_{191} :\equiv (X \subseteq \mathbb{R} \to \varphi_{3}) \to (\mathrm{LB}(b,X) \leftrightarrow b \in \mathbb{R} \land \forall x \in X\,(b \leq x)) \to \varphi_{49}$
\item $\varphi_{192} :\equiv \forall s\,(\varphi_{27} \to \neg \mathrm{IsSup}(s,\ominus X)) \to \forall s\,(\varphi_{27}) \to \forall s\,(\neg \mathrm{IsSup}(s,\ominus X))$
\item $\varphi_{193} :\equiv \forall a\,(\varphi_{41} \to \mathrm{LB}(a,X) \to a \leq -s) \to \forall a\,(\varphi_{41}) \to \forall a\,(\mathrm{LB}(a,X) \to a \leq -s)$
\item $\varphi_{194} :\equiv \forall x\,(\varphi_{41} \to x \in X \to -s \leq x) \to \forall x\,(\varphi_{41}) \to \forall x \in X\,(-s \leq x)$
\item $\varphi_{195} :\equiv (X \subseteq \mathbb{R} \to X \neq \emptyset \to \exists b\,(\mathrm{LB}(b,X)) \to x \in X \to x \in X) \to (x = -(-x) \to x \in X \to -(-x) \in X) \to \varphi_{71}$
\item $\varphi_{196} :\equiv (-(-y) = y \to -(-y) < -b \to y < -b) \to (-(-y) = y \to -(-y) = -b \to y = -b) \to X \subseteq \mathbb{R} \to \varphi_{58}$
\item $\varphi_{197} :\equiv (X \subseteq \mathbb{R} \land X \neq \emptyset \land \exists b\,(\mathrm{LB}(b,X)) \land \mathrm{LB}(b,X) \to \varphi_{14}) \to X \subseteq \mathbb{R} \to X \neq \emptyset \to \varphi_{20}$
\item $\varphi_{198} :\equiv (X \subseteq \mathbb{R} \to X \neq \emptyset \to \exists b\,(\mathrm{LB}(b,X)) \to \ominus X \neq \emptyset) \to \varphi_{40} \to X \subseteq \mathbb{R} \to \varphi_{7}$
\item $\varphi_{199} :\equiv (X \subseteq \mathbb{R} \to \varphi_{21}) \to (\forall x \in X\,(b \leq x) \to -y \in X \to b \leq -y) \to X \subseteq \mathbb{R} \to \varphi_{48}$
\item $\varphi_{200} :\equiv (X \subseteq \mathbb{R} \to X \subseteq \mathbb{R}) \to (X \subseteq \mathbb{R} \to X \neq \emptyset \to \varphi_{13}) \to X \subseteq \mathbb{R} \to X \neq \emptyset \to \varphi_{25}$
\item $\varphi_{201} :\equiv (y \in \mathbb{R} \to -(-y) = y) \to (X \subseteq \mathbb{R} \to \varphi_{66}) \to X \subseteq \mathbb{R} \to X \neq \emptyset \to \varphi_{26}$
\item $\varphi_{202} :\equiv (X \subseteq \mathbb{R} \to X \neq \emptyset \to \varphi_{73}) \to \varphi_{41} \land \mathrm{LB}(a,X) \to y \in \ominus X \to y \leq -a$
\item $\varphi_{203} :\equiv (a \in \mathbb{R} \to -(-a) = a) \to (X \subseteq \mathbb{R} \to X \neq \emptyset \to \varphi_{24}) \to X \subseteq \mathbb{R} \to X \neq \emptyset \to \varphi_{32}$
\item $\varphi_{204} :\equiv (x \in \mathbb{R} \to -(-x) = x) \to (X \subseteq \mathbb{R} \to X \neq \emptyset \to \varphi_{25}) \to X \subseteq \mathbb{R} \to X \neq \emptyset \to \varphi_{33}$
\item $\varphi_{205} :\equiv (a \in \mathbb{R} \to -a \in \mathbb{R}) \to (X \subseteq \mathbb{R} \to X \neq \emptyset \to \varphi_{24}) \to X \subseteq \mathbb{R} \to X \neq \emptyset \to \varphi_{42}$
\item $\varphi_{206} :\equiv (x \in \mathbb{R} \to -x \in \mathbb{R}) \to (X \subseteq \mathbb{R} \to X \neq \emptyset \to \varphi_{25}) \to X \subseteq \mathbb{R} \to X \neq \emptyset \to \varphi_{43}$
\item $\varphi_{207} :\equiv (X \subseteq \mathbb{R} \to X \neq \emptyset \to \varphi_{33}) \to (-(-x) = x \to x = -(-x)) \to X \subseteq \mathbb{R} \to X \neq \emptyset \to \varphi_{34}$
\item $\varphi_{208} :\equiv (X \subseteq \mathbb{R} \to X \neq \emptyset \to \varphi_{85}) \to X \subseteq \mathbb{R} \to X \neq \emptyset \to \exists b\,(\mathrm{LB}(b,X)) \to \varphi_{51}$
\item $\varphi_{209} :\equiv \forall y\,(\varphi_{41} \land \mathrm{LB}(a,X) \to y \in \ominus X \to y \leq -a) \to \forall y\,(\varphi_{41} \land \mathrm{LB}(a,X)) \to \forall y \in \ominus X\,(y \leq -a)$
\item $\varphi_{210} :\equiv (X \subseteq \mathbb{R} \to X \neq \emptyset \to \varphi_{72}) \to X \subseteq \mathbb{R} \to X \neq \emptyset \to \exists b\,(\mathrm{LB}(b,X)) \to \varphi_{69}$
\item $\varphi_{211} :\equiv (X \subseteq \mathbb{R} \to X \neq \emptyset \to \exists b\,(\mathrm{LB}(b,X)) \to x \in X \to -x \in \mathbb{R}) \to \varphi_{185}$
\item $\varphi_{212} :\equiv (X \subseteq \mathbb{R} \to \varphi_{74}) \to (y \in \ominus X \leftrightarrow y \in \mathbb{R} \land -y \in X) \to X \subseteq \mathbb{R} \to \varphi_{66}$
\item $\varphi_{213} :\equiv (X \subseteq \mathbb{R} \to \varphi_{74}) \to (y \in \ominus X \leftrightarrow y \in \mathbb{R} \land -y \in X) \to X \subseteq \mathbb{R} \to \varphi_{70}$
\item $\varphi_{214} :\equiv (y \in \mathbb{R} \to -y \in \mathbb{R}) \to (X \subseteq \mathbb{R} \to X \neq \emptyset \to \varphi_{85}) \to X \subseteq \mathbb{R} \to X \neq \emptyset \to \varphi_{100}$
\item $\varphi_{215} :\equiv (X \subseteq \mathbb{R} \to \varphi_{56}) \to (X \subseteq \mathbb{R} \to X \neq \emptyset \to \varphi_{30}) \to X \subseteq \mathbb{R} \to X \neq \emptyset \to \varphi_{5}$
\item $\varphi_{216} :\equiv (X \subseteq \mathbb{R} \to X \neq \emptyset \to \varphi_{15}) \to \varphi_{55} \to X \subseteq \mathbb{R} \to X \neq \emptyset \to \exists b\,(\mathrm{LB}(b,X)) \to \varphi_{28}$
\item $\varphi_{217} :\equiv (\varphi_{41} \land \mathrm{LB}(a,X) \to \forall y\,(\varphi_{41} \land \mathrm{LB}(a,X))) \to X \subseteq \mathbb{R} \to X \neq \emptyset \to \exists b\,(\mathrm{LB}(b,X)) \to \varphi_{46}$
\item $\varphi_{218} :\equiv (X \subseteq \mathbb{R} \to \varphi_{38}) \to (X \subseteq \mathbb{R} \to X \neq \emptyset \to \varphi_{18}) \to X \subseteq \mathbb{R} \to X \neq \emptyset \to \varphi_{54}$
\item $\varphi_{219} :\equiv (X \subseteq \mathbb{R} \to X \neq \emptyset \to \varphi_{99}) \to (y \in \ominus X \leftrightarrow y \in \mathbb{R} \land -y \in X) \to X \subseteq \mathbb{R} \to X \neq \emptyset \to \varphi_{85}$
\item $\varphi_{220} :\equiv (X \subseteq \mathbb{R} \to X \neq \emptyset \to \varphi_{99}) \to (y \in \ominus X \leftrightarrow y \in \mathbb{R} \land -y \in X) \to X \subseteq \mathbb{R} \to X \neq \emptyset \to \varphi_{91}$
\item $\varphi_{221} :\equiv (X \subseteq \mathbb{R} \to X \neq \emptyset \to \exists b\,(\mathrm{LB}(b,X)) \to \varphi_{28}) \to \varphi_{36} \to X \subseteq \mathbb{R} \to X \neq \emptyset \to \exists b\,(\mathrm{LB}(b,X)) \to \varphi_{52}$
\item $\varphi_{222} :\equiv (X \subseteq \mathbb{R} \to X \neq \emptyset \to \varphi_{42}) \to (X \subseteq \mathbb{R} \to X \neq \emptyset \to \varphi_{16}) \to X \subseteq \mathbb{R} \to X \neq \emptyset \to \varphi_{53}$
\item $\varphi_{223} :\equiv (X \subseteq \mathbb{R} \to X \neq \emptyset \to \varphi_{65}) \to (X \subseteq \mathbb{R} \to X \neq \emptyset \to \varphi_{63}) \to X \subseteq \mathbb{R} \to X \neq \emptyset \to \varphi_{18}$
\item $\varphi_{224} :\equiv (X \subseteq \mathbb{R} \to X \neq \emptyset \to \varphi_{43}) \to (X \subseteq \mathbb{R} \to X \neq \emptyset \to \varphi_{50}) \to X \subseteq \mathbb{R} \to X \neq \emptyset \to \varphi_{63}$
\item $\varphi_{225} :\equiv (X \subseteq \mathbb{R} \to X \neq \emptyset \to \varphi_{53}) \to (-(-a) = a \to -(-a) < -s \to a < -s) \to \varphi_{173}$
\item $\varphi_{226} :\equiv (X \subseteq \mathbb{R} \to X \neq \emptyset \to \varphi_{54}) \to (-(-x) = x \to -s < -(-x) \to -s < x) \to \varphi_{174}$
\item $\varphi_{227} :\equiv (X \subseteq \mathbb{R} \to X \neq \emptyset \to \exists b\,(\mathrm{LB}(b,X)) \to \varphi_{52}) \to (X \subseteq \mathbb{R} \to X \neq \emptyset \to \varphi_{62}) \to X \subseteq \mathbb{R} \to X \neq \emptyset \to \varphi_{16}$
\item $\varphi_{228} :\equiv (\mathrm{UB}(s,\ominus X) \leftrightarrow s \in \mathbb{R} \land \forall x \in \ominus X\,(x \leq s)) \to \varphi_{79} \to \varphi_{77} \to X \subseteq \mathbb{R} \to X \neq \emptyset \to \varphi_{39}$
\item $\varphi_{229} :\equiv (X \subseteq \mathbb{R} \to X \neq \emptyset \to \varphi_{68}) \to (X \subseteq \mathbb{R} \to X \neq \emptyset \to \varphi_{91}) \to X \subseteq \mathbb{R} \to X \neq \emptyset \to \varphi_{72}$
\item $\varphi_{230} :\equiv (X \subseteq \mathbb{R} \to X \neq \emptyset \to \exists b\,(\mathrm{LB}(b,X)) \to \varphi_{69}) \to (-(-y) = y \to -(-y) < -a \to y < -a) \to \varphi_{188}$
\item $\varphi_{231} :\equiv (\ominus X \subseteq \mathbb{R} \to \varphi_{64}) \to (X \subseteq \mathbb{R} \to X \neq \emptyset \to \exists b\,(\mathrm{LB}(b,X)) \to \ominus X \subseteq \mathbb{R}) \to \varphi_{198}$
\item $\varphi_{232} :\equiv (\varphi_{96} \to \varphi_{105}) \to \varphi_{175} \to X \subseteq \mathbb{R} \to X \neq \emptyset \to \exists b\,(\mathrm{LB}(b,X)) \to \ominus X \neq \emptyset$
\item $\varphi_{233} :\equiv (b \in \mathbb{R} \to -y \in \mathbb{R} \to b \leq -y \to -(-y) \leq -b) \to \varphi_{49} \to (X \subseteq \mathbb{R} \to \varphi_{75}) \to (X \subseteq \mathbb{R} \to \varphi_{57}) \to X \subseteq \mathbb{R} \to X \neq \emptyset \to \varphi_{44}$
\item $\varphi_{234} :\equiv (X \subseteq \mathbb{R} \to X \neq \emptyset \to \exists b\,(\mathrm{LB}(b,X)) \to \varphi_{46}) \to \varphi_{122} \to \varphi_{121} \to X \subseteq \mathbb{R} \to X \neq \emptyset \to \varphi_{62}$
\item $\varphi_{235} :\equiv (X \subseteq \mathbb{R} \to X \neq \emptyset \to \varphi_{5}) \to (X \subseteq \mathbb{R} \to X \neq \emptyset \to \varphi_{29}) \to \varphi_{95} \to \varphi_{94} \to X \subseteq \mathbb{R} \to X \neq \emptyset \to \varphi_{9}$
\item $\varphi_{236} :\equiv (s \in \mathbb{R} \to -a \in \mathbb{R} \to s \leq -a \to -(-a) \leq -s) \to (X \subseteq \mathbb{R} \to \varphi_{38}) \to \varphi_{222}$
\item $\varphi_{237} :\equiv (-x \in \mathbb{R} \to s \in \mathbb{R} \to -x \leq s \to -s \leq -(-x)) \to (X \subseteq \mathbb{R} \to X \neq \emptyset \to \varphi_{43}) \to \varphi_{218}$
\item $\varphi_{238} :\equiv \forall s\,(\varphi_{27} \to \neg \mathrm{IsSup}(s,\ominus X)) \to \varphi_{192} \to (\varphi_{27} \to \forall s\,(\varphi_{27})) \to \varphi_{23}$
\item $\varphi_{239} :\equiv (a \in \mathbb{R} \to -y \in \mathbb{R} \to a \leq -y \to -(-y) \leq -a) \to (X \subseteq \mathbb{R} \to X \neq \emptyset \to \varphi_{24}) \to (X \subseteq \mathbb{R} \to X \neq \emptyset \to \varphi_{100}) \to \varphi_{210}$
\item $\varphi_{240} :\equiv \forall a\,(\varphi_{41} \to \mathrm{LB}(a,X) \to a \leq -s) \to \varphi_{193} \to (\varphi_{41} \to \forall a\,(\varphi_{41})) \to X \subseteq \mathbb{R} \to X \neq \emptyset \to \varphi_{29}$
\item $\varphi_{241} :\equiv \forall x\,(\varphi_{41} \to x \in X \to -s \leq x) \to \varphi_{194} \to (\varphi_{41} \to \forall x\,(\varphi_{41})) \to X \subseteq \mathbb{R} \to X \neq \emptyset \to \varphi_{30}$
\item $\varphi_{242} :\equiv (X \subseteq \mathbb{R} \to X \neq \emptyset \to \exists b\,(\mathrm{LB}(b,X)) \to \exists x\,(x \in X)) \to \varphi_{96} \to \varphi_{232}$
\item $\varphi_{243} :\equiv (\mathrm{UB}(-b,\ominus X) \leftrightarrow -b \in \mathbb{R} \land \forall x \in \ominus X\,(x \leq -b)) \to \varphi_{67} \to (X \subseteq \mathbb{R} \to X \neq \emptyset \to \varphi_{20}) \to \varphi_{124} \to \varphi_{123} \to X \subseteq \mathbb{R} \to \varphi_{35}$
\item $\varphi_{244} :\equiv (\mathrm{UB}(-a,\ominus X) \leftrightarrow -a \in \mathbb{R} \land \forall x \in \ominus X\,(x \leq -a)) \to (X \subseteq \mathbb{R} \to X \neq \emptyset \to \varphi_{42}) \to \varphi_{234}$
\end{abbrevlist}
\end{fullproof}

% 自動生成 (prover/build.js): SubSubSelf  ノード数 146
\begin{fullproof}{[B6]}{$\forall x\,\forall y\,(x \in \mathbb{R} \to y \in \mathbb{R} \to x - (x - y) = y)$}{146}\phantomsection\label{tree:SubSubSelf}
\noindent$\Pi_{1}$:
\begin{ptree}
\AXC{$y = z \to x = y \to x = z$\lbl{A6}}
\AXC{$(y = z \to x = y \to x = z) \to x = y \to y = z \to x = z$\lbl{A1}}
\BIC{$x = y \to y = z \to x = z$}
\UIC{$\forall z\,(x = y \to y = z \to x = z)$}
\UIC{$\forall y\,\forall z\,(x = y \to y = z \to x = z)$}
\UIC{$\forall x\,\forall y\,\forall z\,(x = y \to y = z \to x = z)$}
\AXC{$\varphi_{4}$\lbl{A2}}
\BIC{$\forall y\,\forall z\,(v_{371} = y \to y = z \to v_{371} = z)$}
\AXC{$\varphi_{14}$\lbl{A2}}
\BIC{$\forall z\,(v_{371} = v_{372} \to v_{372} = z \to v_{371} = z)$}
\AXC{$\varphi_{21}$\lbl{A2}}
\BIC{$v_{371} = v_{372} \to v_{372} = v_{373} \to v_{371} = v_{373}$}
\UIC{$\forall v_{373}\,(v_{371} = v_{372} \to v_{372} = v_{373} \to v_{371} = v_{373})$}
\UIC{$\varphi_{1}$}
\UIC{$\forall v_{371}\,(\varphi_{1})$}
\end{ptree}
\noindent$\Pi_{2}$:
\begin{ptree}
\AXC{$\stackrel{\Pi_{1}}{\vdots}$}\noLine
\UIC{$\forall v_{371}\,(\varphi_{1})$}
\AXC{$\forall v_{371}\,(\varphi_{1}) \to \varphi_{10}$\lbl{A2}}
\BIC{$\varphi_{10}$}
\AXC{$\varphi_{10} \to \varphi_{24}$\lbl{A2}}
\BIC{$\varphi_{24}$}
\AXC{$\varphi_{24} \to \varphi_{8}$\lbl{A2}}
\BIC{$\varphi_{8}$}
\end{ptree}
\noindent$\Pi_{3}$:
\begin{ptree}
\AXC{$\forall x\,\forall y\,(x \in \mathbb{R} \to y \in \mathbb{R} \to x - y + y = x)$\lbl{[F12]}}
\AXC{$\varphi_{36}$\lbl{A2}}
\BIC{$\forall y\,(v_{363} \in \mathbb{R} \to y \in \mathbb{R} \to v_{363} - y + y = v_{363})$}
\AXC{$\varphi_{41}$\lbl{A2}}
\BIC{$v_{363} \in \mathbb{R} \to v_{364} \in \mathbb{R} \to v_{363} - v_{364} + v_{364} = v_{363}$}
\UIC{$\varphi_{18}$}
\UIC{$\forall v_{363}\,(\varphi_{18})$}
\AXC{$\varphi_{43}$\lbl{A2}}
\BIC{$\forall v_{364}\,(y \in \mathbb{R} \to v_{364} \in \mathbb{R} \to y - v_{364} + v_{364} = y)$}
\AXC{$\varphi_{34}$\lbl{A2}}
\BIC{$y \in \mathbb{R} \to x \in \mathbb{R} \to y - x + x = y$}
\AXC{$\varphi_{47}$\lbl{A1}}
\BIC{$(x \in \mathbb{R} \to y \in \mathbb{R} \to y \in \mathbb{R}) \to \varphi_{28}$}
\end{ptree}
\noindent$\Pi_{4}$:
\begin{ptree}
\AXC{$x \in \mathbb{R} \to y \in \mathbb{R} \to x \in \mathbb{R}$\lbl{A1}}
\AXC{$x \in \mathbb{R} \to y \in \mathbb{R} \to y \in \mathbb{R}$\lbl{A1}}
\AXC{$\stackrel{\Pi_{3}}{\vdots}$}\noLine
\UIC{$(x \in \mathbb{R} \to y \in \mathbb{R} \to y \in \mathbb{R}) \to \varphi_{28}$}
\BIC{$\varphi_{28}$}
\BIC{$x \in \mathbb{R} \to y \in \mathbb{R} \to y - x + x = y$}
\end{ptree}
\noindent$\Pi_{5}$:
\begin{ptree}
\AXC{$y = z \to x = y \to x = z$\lbl{A6}}
\AXC{$(y = z \to x = y \to x = z) \to x = y \to y = z \to x = z$\lbl{A1}}
\BIC{$x = y \to y = z \to x = z$}
\UIC{$\forall z\,(x = y \to y = z \to x = z)$}
\UIC{$\forall y\,\forall z\,(x = y \to y = z \to x = z)$}
\UIC{$\forall x\,\forall y\,\forall z\,(x = y \to y = z \to x = z)$}
\AXC{$\varphi_{5}$\lbl{A2}}
\BIC{$\forall y\,\forall z\,(v_{368} = y \to y = z \to v_{368} = z)$}
\AXC{$\varphi_{15}$\lbl{A2}}
\BIC{$\forall z\,(v_{368} = v_{369} \to v_{369} = z \to v_{368} = z)$}
\AXC{$\varphi_{22}$\lbl{A2}}
\BIC{$v_{368} = v_{369} \to v_{369} = v_{370} \to v_{368} = v_{370}$}
\UIC{$\forall v_{370}\,(v_{368} = v_{369} \to v_{369} = v_{370} \to v_{368} = v_{370})$}
\UIC{$\varphi_{2}$}
\UIC{$\forall v_{368}\,(\varphi_{2})$}
\end{ptree}
\noindent$\Pi_{6}$:
\begin{ptree}
\AXC{$\stackrel{\Pi_{5}}{\vdots}$}\noLine
\UIC{$\forall v_{368}\,(\varphi_{2})$}
\AXC{$\forall v_{368}\,(\varphi_{2}) \to \varphi_{11}$\lbl{A2}}
\BIC{$\varphi_{11}$}
\AXC{$\varphi_{11} \to \varphi_{25}$\lbl{A2}}
\BIC{$\varphi_{25}$}
\AXC{$\varphi_{25} \to \varphi_{26}$\lbl{A2}}
\BIC{$\varphi_{26}$}
\end{ptree}
\noindent$\Pi_{7}$:
\begin{ptree}
\AXC{$\forall x\,\forall y\,(x \in \mathbb{R} \to y \in \mathbb{R} \to x - y \in \mathbb{R})$\lbl{[F1]}}
\AXC{$\varphi_{38}$\lbl{A2}}
\BIC{$\forall y\,(v_{295} \in \mathbb{R} \to y \in \mathbb{R} \to v_{295} - y \in \mathbb{R})$}
\AXC{$\varphi_{39}$\lbl{A2}}
\BIC{$v_{295} \in \mathbb{R} \to v_{296} \in \mathbb{R} \to v_{295} - v_{296} \in \mathbb{R}$}
\UIC{$\forall v_{296}\,(v_{295} \in \mathbb{R} \to v_{296} \in \mathbb{R} \to v_{295} - v_{296} \in \mathbb{R})$}
\UIC{$\varphi_{20}$}
\AXC{$\varphi_{20} \to \forall v_{296}\,(y \in \mathbb{R} \to v_{296} \in \mathbb{R} \to y - v_{296} \in \mathbb{R})$\lbl{A2}}
\BIC{$\forall v_{296}\,(y \in \mathbb{R} \to v_{296} \in \mathbb{R} \to y - v_{296} \in \mathbb{R})$}
\end{ptree}
\noindent$\Pi_{8}$:
\begin{ptree}
\AXC{$x \in \mathbb{R} \to y \in \mathbb{R} \to y \in \mathbb{R}$\lbl{A1}}
\AXC{$\stackrel{\Pi_{7}}{\vdots}$}\noLine
\UIC{$\forall v_{296}\,(y \in \mathbb{R} \to v_{296} \in \mathbb{R} \to y - v_{296} \in \mathbb{R})$}
\AXC{$\varphi_{35}$\lbl{A2}}
\BIC{$y \in \mathbb{R} \to x \in \mathbb{R} \to y - x \in \mathbb{R}$}
\AXC{$\varphi_{48}$\lbl{A1}}
\BIC{$(x \in \mathbb{R} \to y \in \mathbb{R} \to y \in \mathbb{R}) \to \varphi_{30}$}
\BIC{$\varphi_{30}$}
\end{ptree}
\noindent$\Pi_{9}$:
\begin{ptree}
\AXC{$x \in \mathbb{R} \to y \in \mathbb{R} \to x \in \mathbb{R}$\lbl{A1}}
\AXC{$\forall x\,\forall y\,(x \in \mathbb{R} \to y \in \mathbb{R} \to x + y = y + x)$\lbl{(R3)}}
\AXC{$\varphi_{32}$\lbl{A2}}
\BIC{$\forall y\,(x \in \mathbb{R} \to y \in \mathbb{R} \to x + y = y + x)$}
\AXC{$\varphi_{37}$\lbl{A2}}
\BIC{$x \in \mathbb{R} \to y - x \in \mathbb{R} \to x + (y - x) = y - x + x$}
\AXC{$\varphi_{51}$\lbl{A1}}
\BIC{$(x \in \mathbb{R} \to y \in \mathbb{R} \to x \in \mathbb{R}) \to \varphi_{33}$}
\BIC{$\varphi_{33}$}
\end{ptree}
\noindent$\Pi_{10}$:
\begin{ptree}
\AXC{$y = z \to x = y \to x = z$\lbl{A6}}
\AXC{$(y = z \to x = y \to x = z) \to x = y \to y = z \to x = z$\lbl{A1}}
\BIC{$x = y \to y = z \to x = z$}
\UIC{$\forall z\,(x = y \to y = z \to x = z)$}
\UIC{$\forall y\,\forall z\,(x = y \to y = z \to x = z)$}
\UIC{$\forall x\,\forall y\,\forall z\,(x = y \to y = z \to x = z)$}
\AXC{$\varphi_{6}$\lbl{A2}}
\BIC{$\forall y\,\forall z\,(v_{365} = y \to y = z \to v_{365} = z)$}
\AXC{$\varphi_{16}$\lbl{A2}}
\BIC{$\forall z\,(v_{365} = v_{366} \to v_{366} = z \to v_{365} = z)$}
\AXC{$\varphi_{23}$\lbl{A2}}
\BIC{$v_{365} = v_{366} \to v_{366} = v_{367} \to v_{365} = v_{367}$}
\UIC{$\forall v_{367}\,(v_{365} = v_{366} \to v_{366} = v_{367} \to v_{365} = v_{367})$}
\UIC{$\varphi_{3}$}
\UIC{$\forall v_{365}\,(\varphi_{3})$}
\end{ptree}
\noindent$\Pi_{11}$:
\begin{ptree}
\AXC{$\stackrel{\Pi_{10}}{\vdots}$}\noLine
\UIC{$\forall v_{365}\,(\varphi_{3})$}
\AXC{$\forall v_{365}\,(\varphi_{3}) \to \varphi_{12}$\lbl{A2}}
\BIC{$\varphi_{12}$}
\AXC{$\varphi_{12} \to \varphi_{27}$\lbl{A2}}
\BIC{$\varphi_{27}$}
\AXC{$\varphi_{27} \to x - (x - y) = x + -(x - y) \to \varphi_{7}$\lbl{A2}}
\BIC{$x - (x - y) = x + -(x - y) \to \varphi_{7}$}
\end{ptree}
\noindent$\Pi_{12}$:
\begin{ptree}
\AXC{$\forall x\,\forall y\,(x \in \mathbb{R} \to y \in \mathbb{R} \to -(y - x) = x - y)$\lbl{[F14]}}
\AXC{$\varphi_{40}$\lbl{A2}}
\BIC{$\forall y\,(v_{361} \in \mathbb{R} \to y \in \mathbb{R} \to -(y - v_{361}) = v_{361} - y)$}
\AXC{$\varphi_{42}$\lbl{A2}}
\BIC{$\varphi_{17}$}
\UIC{$\forall v_{362}\,(\varphi_{17})$}
\UIC{$\forall v_{361}\,\forall v_{362}\,(\varphi_{17})$}
\AXC{$\forall v_{361}\,\forall v_{362}\,(\varphi_{17}) \to \varphi_{13}$\lbl{A2}}
\BIC{$\varphi_{13}$}
\end{ptree}
\noindent$\Pi_{13}$:
\begin{ptree}
\AXC{$x \in \mathbb{R} \to y \in \mathbb{R} \to y \in \mathbb{R}$\lbl{A1}}
\AXC{$\stackrel{\Pi_{12}}{\vdots}$}\noLine
\UIC{$\varphi_{13}$}
\AXC{$\varphi_{13} \to y \in \mathbb{R} \to x \in \mathbb{R} \to -(x - y) = y - x$\lbl{A2}}
\BIC{$y \in \mathbb{R} \to x \in \mathbb{R} \to -(x - y) = y - x$}
\AXC{$\varphi_{49}$\lbl{A1}}
\BIC{$(x \in \mathbb{R} \to y \in \mathbb{R} \to y \in \mathbb{R}) \to \varphi_{31}$}
\BIC{$\varphi_{31}$}
\end{ptree}
\noindent$\Pi_{14}$:
\begin{ptree}
\AXC{$x = x$\lbl{A5}}
\UIC{$\forall x\,(x = x)$}
\AXC{$\forall x\,(x = x) \to x + -(x - y) = x + -(x - y)$\lbl{A2}}
\BIC{$x + -(x - y) = x + -(x - y)$}
\AXC{$x \in \mathbb{R} \to y \in \mathbb{R} \to x \in \mathbb{R}$\lbl{A1}}
\AXC{$\stackrel{\Pi_{13}}{\vdots}$}\noLine
\UIC{$\varphi_{31}$}
\BIC{$x \in \mathbb{R} \to y \in \mathbb{R} \to -(x - y) = y - x$}
\AXC{$\varphi_{52}$\lbl{A1}}
\BIC{$x + -(x - y) = x + -(x - y) \to \varphi_{46}$}
\BIC{$\varphi_{46}$}
\end{ptree}
\noindent$\Pi_{15}$:
\begin{ptree}
\AXC{$z = w \to x + -(x - y) = x + z \to x + -(x - y) = x + w$\lbl{A6}}
\UIC{$\varphi_{9}$}
\UIC{$\forall z\,(\varphi_{9})$}
\AXC{$\forall z\,(\varphi_{9}) \to \varphi_{29}$\lbl{A2}}
\BIC{$\varphi_{29}$}
\AXC{$\varphi_{29} \to -(x - y) = y - x \to \varphi_{19}$\lbl{A2}}
\BIC{$-(x - y) = y - x \to \varphi_{19}$}
\AXC{$\stackrel{\Pi_{14}}{\vdots}$}\noLine
\UIC{$\varphi_{46}$}
\BIC{$x \in \mathbb{R} \to y \in \mathbb{R} \to x + -(x - y) = x + (y - x)$}
\end{ptree}
\noindent$\Pi_{16}$:
\begin{ptree}
\AXC{$\forall x\,\forall y\,(x - y = x + -y)$\lbl{Def$-$}}
\AXC{$\forall x\,\forall y\,(x - y = x + -y) \to \forall y\,(x - y = x + -y)$\lbl{A2}}
\BIC{$\forall y\,(x - y = x + -y)$}
\end{ptree}
\noindent$\Pi_{17}$:
\begin{ptree}
\AXC{$\stackrel{\Pi_{16}}{\vdots}$}\noLine
\UIC{$\forall y\,(x - y = x + -y)$}
\AXC{$\forall y\,(x - y = x + -y) \to x - (x - y) = x + -(x - y)$\lbl{A2}}
\BIC{$x - (x - y) = x + -(x - y)$}
\AXC{$\varphi_{53}$\lbl{A1}}
\BIC{$(x \in \mathbb{R} \to y \in \mathbb{R} \to x + -(x - y) = x + (y - x)) \to \varphi_{44}$}
\end{ptree}
\noindent$\Pi_{18}$:
\begin{ptree}
\AXC{$\stackrel{\Pi_{15}}{\vdots}$}\noLine
\UIC{$x \in \mathbb{R} \to y \in \mathbb{R} \to x + -(x - y) = x + (y - x)$}
\AXC{$\stackrel{\Pi_{17}}{\vdots}$}\noLine
\UIC{$(x \in \mathbb{R} \to y \in \mathbb{R} \to x + -(x - y) = x + (y - x)) \to \varphi_{44}$}
\BIC{$\varphi_{44}$}
\end{ptree}
\noindent$\Pi_{19}$:
\begin{ptree}
\AXC{$\stackrel{\Pi_{11}}{\vdots}$}\noLine
\UIC{$x - (x - y) = x + -(x - y) \to \varphi_{7}$}
\AXC{$\stackrel{\Pi_{18}}{\vdots}$}\noLine
\UIC{$\varphi_{44}$}
\BIC{$x \in \mathbb{R} \to y \in \mathbb{R} \to x - (x - y) = x + (y - x)$}
\AXC{$(x \in \mathbb{R} \to y \in \mathbb{R} \to x - (x - y) = x + (y - x)) \to \varphi_{50}$\lbl{A1}}
\BIC{$\varphi_{50}$}
\end{ptree}
\noindent$\Pi_{20}$:
\begin{ptree}
\AXC{$\stackrel{\Pi_{6}}{\vdots}$}\noLine
\UIC{$\varphi_{26}$}
\AXC{$x \in \mathbb{R} \to y \in \mathbb{R} \to x \in \mathbb{R}$\lbl{A1}}
\AXC{$\stackrel{\Pi_{8}}{\vdots}$}\noLine
\UIC{$\varphi_{30}$}
\BIC{$x \in \mathbb{R} \to y \in \mathbb{R} \to y - x \in \mathbb{R}$}
\AXC{$\stackrel{\Pi_{9}}{\vdots}$}\noLine
\UIC{$\varphi_{33}$}
\BIC{$x \in \mathbb{R} \to y \in \mathbb{R} \to x + (y - x) = y - x + x$}
\AXC{$\stackrel{\Pi_{19}}{\vdots}$}\noLine
\UIC{$\varphi_{50}$}
\BIC{$\varphi_{26} \to x \in \mathbb{R} \to y \in \mathbb{R} \to x - (x - y) = y - x + x$}
\BIC{$x \in \mathbb{R} \to y \in \mathbb{R} \to x - (x - y) = y - x + x$}
\end{ptree}
\noindent$\Pi_{21}$:
\begin{ptree}
\AXC{$\stackrel{\Pi_{20}}{\vdots}$}\noLine
\UIC{$x \in \mathbb{R} \to y \in \mathbb{R} \to x - (x - y) = y - x + x$}
\AXC{$(x \in \mathbb{R} \to y \in \mathbb{R} \to x - (x - y) = y - x + x) \to \varphi_{45}$\lbl{A1}}
\BIC{$\varphi_{45}$}
\end{ptree}
\noindent 主木:
\begin{ptree}
\AXC{$\stackrel{\Pi_{2}}{\vdots}$}\noLine
\UIC{$\varphi_{8}$}
\AXC{$\stackrel{\Pi_{4}}{\vdots}$}\noLine
\UIC{$x \in \mathbb{R} \to y \in \mathbb{R} \to y - x + x = y$}
\AXC{$\stackrel{\Pi_{21}}{\vdots}$}\noLine
\UIC{$\varphi_{45}$}
\BIC{$\varphi_{8} \to x \in \mathbb{R} \to y \in \mathbb{R} \to x - (x - y) = y$}
\BIC{$x \in \mathbb{R} \to y \in \mathbb{R} \to x - (x - y) = y$}
\UIC{$\forall y\,(x \in \mathbb{R} \to y \in \mathbb{R} \to x - (x - y) = y)$}
\UIC{$\forall x\,\forall y\,(x \in \mathbb{R} \to y \in \mathbb{R} \to x - (x - y) = y)$}
\end{ptree}
\begin{abbrevlist}
\item $\varphi_{1} :\equiv \forall v_{372}\,\forall v_{373}\,(v_{371} = v_{372} \to v_{372} = v_{373} \to v_{371} = v_{373})$
\item $\varphi_{2} :\equiv \forall v_{369}\,\forall v_{370}\,(v_{368} = v_{369} \to v_{369} = v_{370} \to v_{368} = v_{370})$
\item $\varphi_{3} :\equiv \forall v_{366}\,\forall v_{367}\,(v_{365} = v_{366} \to v_{366} = v_{367} \to v_{365} = v_{367})$
\item $\varphi_{4} :\equiv \forall x\,\forall y\,\forall z\,(x = y \to y = z \to x = z) \to \forall y\,\forall z\,(v_{371} = y \to y = z \to v_{371} = z)$
\item $\varphi_{5} :\equiv \forall x\,\forall y\,\forall z\,(x = y \to y = z \to x = z) \to \forall y\,\forall z\,(v_{368} = y \to y = z \to v_{368} = z)$
\item $\varphi_{6} :\equiv \forall x\,\forall y\,\forall z\,(x = y \to y = z \to x = z) \to \forall y\,\forall z\,(v_{365} = y \to y = z \to v_{365} = z)$
\item $\varphi_{7} :\equiv x + -(x - y) = x + (y - x) \to x - (x - y) = x + (y - x)$
\item $\varphi_{8} :\equiv x - (x - y) = y - x + x \to y - x + x = y \to x - (x - y) = y$
\item $\varphi_{9} :\equiv \forall w\,(z = w \to x + -(x - y) = x + z \to x + -(x - y) = x + w)$
\item $\varphi_{10} :\equiv \forall v_{372}\,\forall v_{373}\,(x - (x - y) = v_{372} \to v_{372} = v_{373} \to x - (x - y) = v_{373})$
\item $\varphi_{11} :\equiv \forall v_{369}\,\forall v_{370}\,(x - (x - y) = v_{369} \to v_{369} = v_{370} \to x - (x - y) = v_{370})$
\item $\varphi_{12} :\equiv \forall v_{366}\,\forall v_{367}\,(x - (x - y) = v_{366} \to v_{366} = v_{367} \to x - (x - y) = v_{367})$
\item $\varphi_{13} :\equiv \forall v_{362}\,(y \in \mathbb{R} \to v_{362} \in \mathbb{R} \to -(v_{362} - y) = y - v_{362})$
\item $\varphi_{14} :\equiv \forall y\,\forall z\,(v_{371} = y \to y = z \to v_{371} = z) \to \forall z\,(v_{371} = v_{372} \to v_{372} = z \to v_{371} = z)$
\item $\varphi_{15} :\equiv \forall y\,\forall z\,(v_{368} = y \to y = z \to v_{368} = z) \to \forall z\,(v_{368} = v_{369} \to v_{369} = z \to v_{368} = z)$
\item $\varphi_{16} :\equiv \forall y\,\forall z\,(v_{365} = y \to y = z \to v_{365} = z) \to \forall z\,(v_{365} = v_{366} \to v_{366} = z \to v_{365} = z)$
\item $\varphi_{17} :\equiv v_{361} \in \mathbb{R} \to v_{362} \in \mathbb{R} \to -(v_{362} - v_{361}) = v_{361} - v_{362}$
\item $\varphi_{18} :\equiv \forall v_{364}\,(v_{363} \in \mathbb{R} \to v_{364} \in \mathbb{R} \to v_{363} - v_{364} + v_{364} = v_{363})$
\item $\varphi_{19} :\equiv x + -(x - y) = x + -(x - y) \to x + -(x - y) = x + (y - x)$
\item $\varphi_{20} :\equiv \forall v_{295}\,\forall v_{296}\,(v_{295} \in \mathbb{R} \to v_{296} \in \mathbb{R} \to v_{295} - v_{296} \in \mathbb{R})$
\item $\varphi_{21} :\equiv \forall z\,(v_{371} = v_{372} \to v_{372} = z \to v_{371} = z) \to v_{371} = v_{372} \to v_{372} = v_{373} \to v_{371} = v_{373}$
\item $\varphi_{22} :\equiv \forall z\,(v_{368} = v_{369} \to v_{369} = z \to v_{368} = z) \to v_{368} = v_{369} \to v_{369} = v_{370} \to v_{368} = v_{370}$
\item $\varphi_{23} :\equiv \forall z\,(v_{365} = v_{366} \to v_{366} = z \to v_{365} = z) \to v_{365} = v_{366} \to v_{366} = v_{367} \to v_{365} = v_{367}$
\item $\varphi_{24} :\equiv \forall v_{373}\,(x - (x - y) = y - x + x \to y - x + x = v_{373} \to x - (x - y) = v_{373})$
\item $\varphi_{25} :\equiv \forall v_{370}\,(x - (x - y) = x + (y - x) \to x + (y - x) = v_{370} \to x - (x - y) = v_{370})$
\item $\varphi_{26} :\equiv x - (x - y) = x + (y - x) \to x + (y - x) = y - x + x \to x - (x - y) = y - x + x$
\item $\varphi_{27} :\equiv \forall v_{367}\,(x - (x - y) = x + -(x - y) \to x + -(x - y) = v_{367} \to x - (x - y) = v_{367})$
\item $\varphi_{28} :\equiv (x \in \mathbb{R} \to y \in \mathbb{R} \to x \in \mathbb{R}) \to x \in \mathbb{R} \to y \in \mathbb{R} \to y - x + x = y$
\item $\varphi_{29} :\equiv \forall w\,(-(x - y) = w \to x + -(x - y) = x + -(x - y) \to x + -(x - y) = x + w)$
\item $\varphi_{30} :\equiv (x \in \mathbb{R} \to y \in \mathbb{R} \to x \in \mathbb{R}) \to x \in \mathbb{R} \to y \in \mathbb{R} \to y - x \in \mathbb{R}$
\item $\varphi_{31} :\equiv (x \in \mathbb{R} \to y \in \mathbb{R} \to x \in \mathbb{R}) \to x \in \mathbb{R} \to y \in \mathbb{R} \to -(x - y) = y - x$
\item $\varphi_{32} :\equiv \forall x\,\forall y\,(x \in \mathbb{R} \to y \in \mathbb{R} \to x + y = y + x) \to \forall y\,(x \in \mathbb{R} \to y \in \mathbb{R} \to x + y = y + x)$
\item $\varphi_{33} :\equiv (x \in \mathbb{R} \to y \in \mathbb{R} \to y - x \in \mathbb{R}) \to x \in \mathbb{R} \to y \in \mathbb{R} \to x + (y - x) = y - x + x$
\item $\varphi_{34} :\equiv \forall v_{364}\,(y \in \mathbb{R} \to v_{364} \in \mathbb{R} \to y - v_{364} + v_{364} = y) \to y \in \mathbb{R} \to x \in \mathbb{R} \to y - x + x = y$
\item $\varphi_{35} :\equiv \forall v_{296}\,(y \in \mathbb{R} \to v_{296} \in \mathbb{R} \to y - v_{296} \in \mathbb{R}) \to y \in \mathbb{R} \to x \in \mathbb{R} \to y - x \in \mathbb{R}$
\item $\varphi_{36} :\equiv \forall x\,\forall y\,(x \in \mathbb{R} \to y \in \mathbb{R} \to x - y + y = x) \to \forall y\,(v_{363} \in \mathbb{R} \to y \in \mathbb{R} \to v_{363} - y + y = v_{363})$
\item $\varphi_{37} :\equiv \forall y\,(x \in \mathbb{R} \to y \in \mathbb{R} \to x + y = y + x) \to x \in \mathbb{R} \to y - x \in \mathbb{R} \to x + (y - x) = y - x + x$
\item $\varphi_{38} :\equiv \forall x\,\forall y\,(x \in \mathbb{R} \to y \in \mathbb{R} \to x - y \in \mathbb{R}) \to \forall y\,(v_{295} \in \mathbb{R} \to y \in \mathbb{R} \to v_{295} - y \in \mathbb{R})$
\item $\varphi_{39} :\equiv \forall y\,(v_{295} \in \mathbb{R} \to y \in \mathbb{R} \to v_{295} - y \in \mathbb{R}) \to v_{295} \in \mathbb{R} \to v_{296} \in \mathbb{R} \to v_{295} - v_{296} \in \mathbb{R}$
\item $\varphi_{40} :\equiv \forall x\,\forall y\,(x \in \mathbb{R} \to y \in \mathbb{R} \to -(y - x) = x - y) \to \forall y\,(v_{361} \in \mathbb{R} \to y \in \mathbb{R} \to -(y - v_{361}) = v_{361} - y)$
\item $\varphi_{41} :\equiv \forall y\,(v_{363} \in \mathbb{R} \to y \in \mathbb{R} \to v_{363} - y + y = v_{363}) \to v_{363} \in \mathbb{R} \to v_{364} \in \mathbb{R} \to v_{363} - v_{364} + v_{364} = v_{363}$
\item $\varphi_{42} :\equiv \forall y\,(v_{361} \in \mathbb{R} \to y \in \mathbb{R} \to -(y - v_{361}) = v_{361} - y) \to \varphi_{17}$
\item $\varphi_{43} :\equiv \forall v_{363}\,(\varphi_{18}) \to \forall v_{364}\,(y \in \mathbb{R} \to v_{364} \in \mathbb{R} \to y - v_{364} + v_{364} = y)$
\item $\varphi_{44} :\equiv (x - (x - y) = x + -(x - y) \to \varphi_{7}) \to x \in \mathbb{R} \to y \in \mathbb{R} \to x - (x - y) = x + (y - x)$
\item $\varphi_{45} :\equiv (x \in \mathbb{R} \to y \in \mathbb{R} \to y - x + x = y) \to \varphi_{8} \to x \in \mathbb{R} \to y \in \mathbb{R} \to x - (x - y) = y$
\item $\varphi_{46} :\equiv (-(x - y) = y - x \to \varphi_{19}) \to x \in \mathbb{R} \to y \in \mathbb{R} \to x + -(x - y) = x + (y - x)$
\item $\varphi_{47} :\equiv (y \in \mathbb{R} \to x \in \mathbb{R} \to y - x + x = y) \to (x \in \mathbb{R} \to y \in \mathbb{R} \to y \in \mathbb{R}) \to \varphi_{28}$
\item $\varphi_{48} :\equiv (y \in \mathbb{R} \to x \in \mathbb{R} \to y - x \in \mathbb{R}) \to (x \in \mathbb{R} \to y \in \mathbb{R} \to y \in \mathbb{R}) \to \varphi_{30}$
\item $\varphi_{49} :\equiv (y \in \mathbb{R} \to x \in \mathbb{R} \to -(x - y) = y - x) \to (x \in \mathbb{R} \to y \in \mathbb{R} \to y \in \mathbb{R}) \to \varphi_{31}$
\item $\varphi_{50} :\equiv (x \in \mathbb{R} \to y \in \mathbb{R} \to x + (y - x) = y - x + x) \to \varphi_{26} \to x \in \mathbb{R} \to y \in \mathbb{R} \to x - (x - y) = y - x + x$
\item $\varphi_{51} :\equiv (x \in \mathbb{R} \to y - x \in \mathbb{R} \to x + (y - x) = y - x + x) \to (x \in \mathbb{R} \to y \in \mathbb{R} \to x \in \mathbb{R}) \to \varphi_{33}$
\item $\varphi_{52} :\equiv (x \in \mathbb{R} \to y \in \mathbb{R} \to -(x - y) = y - x) \to x + -(x - y) = x + -(x - y) \to \varphi_{46}$
\item $\varphi_{53} :\equiv x - (x - y) = x + -(x - y) \to (x \in \mathbb{R} \to y \in \mathbb{R} \to x + -(x - y) = x + (y - x)) \to \varphi_{44}$
\end{abbrevlist}
\end{fullproof}

% 自動生成 (prover/build.js): LtSubAdd  ノード数 52
\begin{fullproof}{[B7]}{$\forall x\,\forall a\,\forall y\,(x \in \mathbb{R} \to a \in \mathbb{R} \to y \in \mathbb{R} \to x - a < y \to x < y + a)$}{52}\phantomsection\label{tree:LtSubAdd}
\noindent$\Pi_{1}$:
\begin{ptree}
\AXC{$\forall x\,\forall y\,(x \in \mathbb{R} \to y \in \mathbb{R} \to x - y \in \mathbb{R})$\lbl{[F1]}}
\AXC{$\varphi_{9}$\lbl{A2}}
\BIC{$\forall y\,(x \in \mathbb{R} \to y \in \mathbb{R} \to x - y \in \mathbb{R})$}
\AXC{$\varphi_{6}$\lbl{A2}}
\BIC{$x \in \mathbb{R} \to a \in \mathbb{R} \to x - a \in \mathbb{R}$}
\AXC{$\varphi_{17}$\lbl{A1}}
\BIC{$(x \in \mathbb{R} \to a \in \mathbb{R} \to y \in \mathbb{R} \to x - a < y \to x \in \mathbb{R}) \to \varphi_{13}$}
\end{ptree}
\noindent$\Pi_{2}$:
\begin{ptree}
\AXC{$x \in \mathbb{R} \to a \in \mathbb{R} \to y \in \mathbb{R} \to x - a < y \to x \in \mathbb{R}$\lbl{A1}}
\AXC{$\stackrel{\Pi_{1}}{\vdots}$}\noLine
\UIC{$(x \in \mathbb{R} \to a \in \mathbb{R} \to y \in \mathbb{R} \to x - a < y \to x \in \mathbb{R}) \to \varphi_{13}$}
\BIC{$\varphi_{13}$}
\end{ptree}
\noindent$\Pi_{3}$:
\begin{ptree}
\AXC{$\varphi_{2}$\lbl{(R14)}}
\AXC{$\varphi_{2} \to \forall y\,(\varphi_{4})$\lbl{A2}}
\BIC{$\forall y\,(\varphi_{4})$}
\AXC{$\forall y\,(\varphi_{4}) \to \varphi_{4}$\lbl{A2}}
\BIC{$\varphi_{4}$}
\AXC{$\varphi_{4} \to \varphi_{3}$\lbl{A2}}
\BIC{$\varphi_{3}$}
\AXC{$\varphi_{3} \to \varphi_{20}$\lbl{A1}}
\BIC{$\varphi_{20}$}
\end{ptree}
\noindent$\Pi_{4}$:
\begin{ptree}
\AXC{$x \in \mathbb{R} \to a \in \mathbb{R} \to y \in \mathbb{R} \to x - a < y \to y \in \mathbb{R}$\lbl{A1}}
\AXC{$x \in \mathbb{R} \to a \in \mathbb{R} \to y \in \mathbb{R} \to x - a < y \to a \in \mathbb{R}$\lbl{A1}}
\AXC{$\stackrel{\Pi_{2}}{\vdots}$}\noLine
\UIC{$\varphi_{13}$}
\BIC{$x \in \mathbb{R} \to a \in \mathbb{R} \to y \in \mathbb{R} \to x - a < y \to x - a \in \mathbb{R}$}
\AXC{$\stackrel{\Pi_{3}}{\vdots}$}\noLine
\UIC{$\varphi_{20}$}
\BIC{$\varphi_{19}$}
\BIC{$(x \in \mathbb{R} \to a \in \mathbb{R} \to y \in \mathbb{R} \to x - a < y \to a \in \mathbb{R}) \to \varphi_{14}$}
\end{ptree}
\noindent$\Pi_{5}$:
\begin{ptree}
\AXC{$x \in \mathbb{R} \to a \in \mathbb{R} \to y \in \mathbb{R} \to x - a < y \to a \in \mathbb{R}$\lbl{A1}}
\AXC{$\stackrel{\Pi_{4}}{\vdots}$}\noLine
\UIC{$(x \in \mathbb{R} \to a \in \mathbb{R} \to y \in \mathbb{R} \to x - a < y \to a \in \mathbb{R}) \to \varphi_{14}$}
\BIC{$\varphi_{14}$}
\end{ptree}
\noindent$\Pi_{6}$:
\begin{ptree}
\AXC{$\forall x\,\forall y\,(x \in \mathbb{R} \to y \in \mathbb{R} \to x - y + y = x)$\lbl{[F12]}}
\AXC{$\varphi_{7}$\lbl{A2}}
\BIC{$\forall y\,(x \in \mathbb{R} \to y \in \mathbb{R} \to x - y + y = x)$}
\AXC{$\varphi_{5}$\lbl{A2}}
\BIC{$x \in \mathbb{R} \to a \in \mathbb{R} \to x - a + a = x$}
\AXC{$\varphi_{16}$\lbl{A1}}
\BIC{$(x \in \mathbb{R} \to a \in \mathbb{R} \to y \in \mathbb{R} \to x - a < y \to x \in \mathbb{R}) \to \varphi_{12}$}
\end{ptree}
\noindent$\Pi_{7}$:
\begin{ptree}
\AXC{$x \in \mathbb{R} \to a \in \mathbb{R} \to y \in \mathbb{R} \to x - a < y \to x \in \mathbb{R}$\lbl{A1}}
\AXC{$\stackrel{\Pi_{6}}{\vdots}$}\noLine
\UIC{$(x \in \mathbb{R} \to a \in \mathbb{R} \to y \in \mathbb{R} \to x - a < y \to x \in \mathbb{R}) \to \varphi_{12}$}
\BIC{$\varphi_{12}$}
\end{ptree}
\noindent$\Pi_{8}$:
\begin{ptree}
\AXC{$x \in \mathbb{R} \to a \in \mathbb{R} \to y \in \mathbb{R} \to x - a < y \to a \in \mathbb{R}$\lbl{A1}}
\AXC{$\stackrel{\Pi_{7}}{\vdots}$}\noLine
\UIC{$\varphi_{12}$}
\BIC{$x \in \mathbb{R} \to a \in \mathbb{R} \to y \in \mathbb{R} \to x - a < y \to x - a + a = x$}
\AXC{$\varphi_{18}$\lbl{A1}}
\BIC{$\varphi_{15}$}
\end{ptree}
\noindent$\Pi_{9}$:
\begin{ptree}
\AXC{$x \in \mathbb{R} \to a \in \mathbb{R} \to y \in \mathbb{R} \to x - a < y \to x - a < y$\lbl{A1}}
\AXC{$\stackrel{\Pi_{5}}{\vdots}$}\noLine
\UIC{$\varphi_{14}$}
\BIC{$x \in \mathbb{R} \to a \in \mathbb{R} \to y \in \mathbb{R} \to x - a < y \to x - a + a < y + a$}
\AXC{$\stackrel{\Pi_{8}}{\vdots}$}\noLine
\UIC{$\varphi_{15}$}
\BIC{$\varphi_{11}$}
\end{ptree}
\noindent 主木:
\begin{ptree}
\AXC{$z = w \to z < y + a \to w < y + a$\lbl{A6}}
\UIC{$\forall w\,(z = w \to z < y + a \to w < y + a)$}
\UIC{$\forall z\,\forall w\,(z = w \to z < y + a \to w < y + a)$}
\AXC{$\varphi_{8}$\lbl{A2}}
\BIC{$\forall w\,(x - a + a = w \to x - a + a < y + a \to w < y + a)$}
\AXC{$\varphi_{10}$\lbl{A2}}
\BIC{$x - a + a = x \to x - a + a < y + a \to x < y + a$}
\AXC{$\stackrel{\Pi_{9}}{\vdots}$}\noLine
\UIC{$\varphi_{11}$}
\BIC{$x \in \mathbb{R} \to a \in \mathbb{R} \to y \in \mathbb{R} \to x - a < y \to x < y + a$}
\UIC{$\forall y\,(x \in \mathbb{R} \to a \in \mathbb{R} \to y \in \mathbb{R} \to x - a < y \to x < y + a)$}
\UIC{$\varphi_{1}$}
\UIC{$\forall x\,(\varphi_{1})$}
\end{ptree}
\begin{abbrevlist}
\item $\varphi_{1} :\equiv \forall a\,\forall y\,(x \in \mathbb{R} \to a \in \mathbb{R} \to y \in \mathbb{R} \to x - a < y \to x < y + a)$
\item $\varphi_{2} :\equiv \forall x\,\forall y\,\forall z\,(x \in \mathbb{R} \to y \in \mathbb{R} \to z \in \mathbb{R} \to x < y \to x + z < y + z)$
\item $\varphi_{3} :\equiv x - a \in \mathbb{R} \to y \in \mathbb{R} \to a \in \mathbb{R} \to x - a < y \to x - a + a < y + a$
\item $\varphi_{4} :\equiv \forall z\,(x - a \in \mathbb{R} \to y \in \mathbb{R} \to z \in \mathbb{R} \to x - a < y \to x - a + z < y + z)$
\item $\varphi_{5} :\equiv \forall y\,(x \in \mathbb{R} \to y \in \mathbb{R} \to x - y + y = x) \to x \in \mathbb{R} \to a \in \mathbb{R} \to x - a + a = x$
\item $\varphi_{6} :\equiv \forall y\,(x \in \mathbb{R} \to y \in \mathbb{R} \to x - y \in \mathbb{R}) \to x \in \mathbb{R} \to a \in \mathbb{R} \to x - a \in \mathbb{R}$
\item $\varphi_{7} :\equiv \forall x\,\forall y\,(x \in \mathbb{R} \to y \in \mathbb{R} \to x - y + y = x) \to \forall y\,(x \in \mathbb{R} \to y \in \mathbb{R} \to x - y + y = x)$
\item $\varphi_{8} :\equiv \forall z\,\forall w\,(z = w \to z < y + a \to w < y + a) \to \forall w\,(x - a + a = w \to x - a + a < y + a \to w < y + a)$
\item $\varphi_{9} :\equiv \forall x\,\forall y\,(x \in \mathbb{R} \to y \in \mathbb{R} \to x - y \in \mathbb{R}) \to \forall y\,(x \in \mathbb{R} \to y \in \mathbb{R} \to x - y \in \mathbb{R})$
\item $\varphi_{10} :\equiv \forall w\,(x - a + a = w \to x - a + a < y + a \to w < y + a) \to x - a + a = x \to x - a + a < y + a \to x < y + a$
\item $\varphi_{11} :\equiv (x - a + a = x \to x - a + a < y + a \to x < y + a) \to x \in \mathbb{R} \to a \in \mathbb{R} \to y \in \mathbb{R} \to x - a < y \to x < y + a$
\item $\varphi_{12} :\equiv (x \in \mathbb{R} \to a \in \mathbb{R} \to y \in \mathbb{R} \to x - a < y \to a \in \mathbb{R}) \to x \in \mathbb{R} \to a \in \mathbb{R} \to y \in \mathbb{R} \to x - a < y \to x - a + a = x$
\item $\varphi_{13} :\equiv (x \in \mathbb{R} \to a \in \mathbb{R} \to y \in \mathbb{R} \to x - a < y \to a \in \mathbb{R}) \to x \in \mathbb{R} \to a \in \mathbb{R} \to y \in \mathbb{R} \to x - a < y \to x - a \in \mathbb{R}$
\item $\varphi_{14} :\equiv (x \in \mathbb{R} \to a \in \mathbb{R} \to y \in \mathbb{R} \to x - a < y \to x - a < y) \to x \in \mathbb{R} \to a \in \mathbb{R} \to y \in \mathbb{R} \to x - a < y \to x - a + a < y + a$
\item $\varphi_{15} :\equiv (x \in \mathbb{R} \to a \in \mathbb{R} \to y \in \mathbb{R} \to x - a < y \to x - a + a < y + a) \to \varphi_{11}$
\item $\varphi_{16} :\equiv (x \in \mathbb{R} \to a \in \mathbb{R} \to x - a + a = x) \to (x \in \mathbb{R} \to a \in \mathbb{R} \to y \in \mathbb{R} \to x - a < y \to x \in \mathbb{R}) \to \varphi_{12}$
\item $\varphi_{17} :\equiv (x \in \mathbb{R} \to a \in \mathbb{R} \to x - a \in \mathbb{R}) \to (x \in \mathbb{R} \to a \in \mathbb{R} \to y \in \mathbb{R} \to x - a < y \to x \in \mathbb{R}) \to \varphi_{13}$
\item $\varphi_{18} :\equiv (x \in \mathbb{R} \to a \in \mathbb{R} \to y \in \mathbb{R} \to x - a < y \to x - a + a = x) \to \varphi_{15}$
\item $\varphi_{19} :\equiv (x \in \mathbb{R} \to a \in \mathbb{R} \to y \in \mathbb{R} \to x - a < y \to y \in \mathbb{R}) \to (x \in \mathbb{R} \to a \in \mathbb{R} \to y \in \mathbb{R} \to x - a < y \to a \in \mathbb{R}) \to \varphi_{14}$
\item $\varphi_{20} :\equiv (x \in \mathbb{R} \to a \in \mathbb{R} \to y \in \mathbb{R} \to x - a < y \to x - a \in \mathbb{R}) \to \varphi_{19}$
\end{abbrevlist}
\end{fullproof}

% 自動生成 (prover/build.js): IotaSMul  ノード数 153
\begin{fullproof}{[B8]}{$\forall m\,\forall a\,(m \in \mathbb{N} \to a \in \mathbb{R} \to \underline{S(m)} \cdot a = \underline{m} \cdot a + a)$}{153}\phantomsection\label{tree:IotaSMul}
\noindent$\Pi_{1}$:
\begin{ptree}
\AXC{$y = z \to x = y \to x = z$\lbl{A6}}
\AXC{$(y = z \to x = y \to x = z) \to x = y \to y = z \to x = z$\lbl{A1}}
\BIC{$x = y \to y = z \to x = z$}
\UIC{$\forall z\,(x = y \to y = z \to x = z)$}
\UIC{$\forall y\,\forall z\,(x = y \to y = z \to x = z)$}
\UIC{$\forall x\,\forall y\,\forall z\,(x = y \to y = z \to x = z)$}
\AXC{$\varphi_{8}$\lbl{A2}}
\BIC{$\forall y\,\forall z\,(\underline{S(m)} \cdot a = y \to y = z \to \underline{S(m)} \cdot a = z)$}
\AXC{$\varphi_{38}$\lbl{A2}}
\BIC{$\varphi_{21}$}
\AXC{$\varphi_{21} \to \varphi_{35}$\lbl{A2}}
\BIC{$\varphi_{35}$}
\end{ptree}
\noindent$\Pi_{2}$:
\begin{ptree}
\AXC{$\forall x \in \mathbb{R}\,(x \cdot 1 = x)$\lbl{(R8)}}
\AXC{$\forall x \in \mathbb{R}\,(x \cdot 1 = x) \to a \in \mathbb{R} \to a \cdot 1 = a$\lbl{A2}}
\BIC{$a \in \mathbb{R} \to a \cdot 1 = a$}
\AXC{$(a \in \mathbb{R} \to a \cdot 1 = a) \to \varphi_{3}$\lbl{A1}}
\BIC{$\varphi_{3}$}
\end{ptree}
\noindent$\Pi_{3}$:
\begin{ptree}
\AXC{$x = x$\lbl{A5}}
\UIC{$\forall x\,(x = x)$}
\AXC{$\forall x\,(x = x) \to \underline{m} \cdot a + a \cdot 1 = \underline{m} \cdot a + a \cdot 1$\lbl{A2}}
\BIC{$\underline{m} \cdot a + a \cdot 1 = \underline{m} \cdot a + a \cdot 1$}
\AXC{$m \in \mathbb{N} \to a \in \mathbb{R} \to a \in \mathbb{R}$\lbl{A1}}
\AXC{$\stackrel{\Pi_{2}}{\vdots}$}\noLine
\UIC{$\varphi_{3}$}
\BIC{$m \in \mathbb{N} \to a \in \mathbb{R} \to a \cdot 1 = a$}
\AXC{$\varphi_{53}$\lbl{A1}}
\BIC{$\underline{m} \cdot a + a \cdot 1 = \underline{m} \cdot a + a \cdot 1 \to \varphi_{43}$}
\BIC{$\varphi_{43}$}
\end{ptree}
\noindent$\Pi_{4}$:
\begin{ptree}
\AXC{$y = z \to x = y \to x = z$\lbl{A6}}
\AXC{$(y = z \to x = y \to x = z) \to x = y \to y = z \to x = z$\lbl{A1}}
\BIC{$x = y \to y = z \to x = z$}
\UIC{$\forall z\,(x = y \to y = z \to x = z)$}
\UIC{$\forall y\,\forall z\,(x = y \to y = z \to x = z)$}
\UIC{$\forall x\,\forall y\,\forall z\,(x = y \to y = z \to x = z)$}
\AXC{$\varphi_{8}$\lbl{A2}}
\BIC{$\forall y\,\forall z\,(\underline{S(m)} \cdot a = y \to y = z \to \underline{S(m)} \cdot a = z)$}
\AXC{$\varphi_{39}$\lbl{A2}}
\BIC{$\varphi_{22}$}
\end{ptree}
\noindent$\Pi_{5}$:
\begin{ptree}
\AXC{$\forall n \in \mathbb{N}\,(\underline{n} \in \mathbb{R})$\lbl{[N4]}}
\AXC{$\forall n \in \mathbb{N}\,(\underline{n} \in \mathbb{R}) \to m \in \mathbb{N} \to \underline{m} \in \mathbb{R}$\lbl{A2}}
\BIC{$m \in \mathbb{N} \to \underline{m} \in \mathbb{R}$}
\AXC{$(m \in \mathbb{N} \to \underline{m} \in \mathbb{R}) \to \varphi_{12}$\lbl{A1}}
\BIC{$\varphi_{12}$}
\end{ptree}
\noindent$\Pi_{6}$:
\begin{ptree}
\AXC{$m \in \mathbb{N} \to a \in \mathbb{R} \to m \in \mathbb{N}$\lbl{A1}}
\AXC{$\stackrel{\Pi_{5}}{\vdots}$}\noLine
\UIC{$\varphi_{12}$}
\BIC{$m \in \mathbb{N} \to a \in \mathbb{R} \to \underline{m} \in \mathbb{R}$}
\end{ptree}
\noindent$\Pi_{7}$:
\begin{ptree}
\AXC{$m \in \mathbb{N} \to a \in \mathbb{R} \to a \in \mathbb{R}$\lbl{A1}}
\AXC{$\forall x\,\forall y\,(x \in \mathbb{R} \to y \in \mathbb{R} \to x \cdot y = y \cdot x)$\lbl{(R7)}}
\AXC{$\varphi_{25}$\lbl{A2}}
\BIC{$\forall y\,(a \in \mathbb{R} \to y \in \mathbb{R} \to a \cdot y = y \cdot a)$}
\AXC{$\varphi_{31}$\lbl{A2}}
\BIC{$a \in \mathbb{R} \to \underline{m} \in \mathbb{R} \to a \cdot \underline{m} = \underline{m} \cdot a$}
\AXC{$\varphi_{46}$\lbl{A1}}
\BIC{$(m \in \mathbb{N} \to a \in \mathbb{R} \to a \in \mathbb{R}) \to \varphi_{27}$}
\BIC{$\varphi_{27}$}
\end{ptree}
\noindent$\Pi_{8}$:
\begin{ptree}
\AXC{$x = x$\lbl{A5}}
\UIC{$\forall x\,(x = x)$}
\AXC{$\forall x\,(x = x) \to a \cdot \underline{m} + a \cdot 1 = a \cdot \underline{m} + a \cdot 1$\lbl{A2}}
\BIC{$a \cdot \underline{m} + a \cdot 1 = a \cdot \underline{m} + a \cdot 1$}
\AXC{$\stackrel{\Pi_{6}}{\vdots}$}\noLine
\UIC{$m \in \mathbb{N} \to a \in \mathbb{R} \to \underline{m} \in \mathbb{R}$}
\AXC{$\stackrel{\Pi_{7}}{\vdots}$}\noLine
\UIC{$\varphi_{27}$}
\BIC{$m \in \mathbb{N} \to a \in \mathbb{R} \to a \cdot \underline{m} = \underline{m} \cdot a$}
\AXC{$\varphi_{56}$\lbl{A1}}
\BIC{$a \cdot \underline{m} + a \cdot 1 = a \cdot \underline{m} + a \cdot 1 \to \varphi_{48}$}
\BIC{$\varphi_{48}$}
\end{ptree}
\noindent$\Pi_{9}$:
\begin{ptree}
\AXC{$y = z \to x = y \to x = z$\lbl{A6}}
\AXC{$(y = z \to x = y \to x = z) \to x = y \to y = z \to x = z$\lbl{A1}}
\BIC{$x = y \to y = z \to x = z$}
\UIC{$\forall z\,(x = y \to y = z \to x = z)$}
\UIC{$\forall y\,\forall z\,(x = y \to y = z \to x = z)$}
\UIC{$\forall x\,\forall y\,\forall z\,(x = y \to y = z \to x = z)$}
\AXC{$\varphi_{8}$\lbl{A2}}
\BIC{$\forall y\,\forall z\,(\underline{S(m)} \cdot a = y \to y = z \to \underline{S(m)} \cdot a = z)$}
\AXC{$\varphi_{36}$\lbl{A2}}
\BIC{$\varphi_{13}$}
\end{ptree}
\noindent$\Pi_{10}$:
\begin{ptree}
\AXC{$m \in \mathbb{N} \to a \in \mathbb{R} \to a \in \mathbb{R}$\lbl{A1}}
\AXC{$\varphi_{9}$\lbl{(R10)}}
\AXC{$\varphi_{9} \to \varphi_{5}$\lbl{A2}}
\BIC{$\varphi_{5}$}
\AXC{$\varphi_{5} \to \varphi_{23}$\lbl{A2}}
\BIC{$\varphi_{23}$}
\AXC{$\varphi_{23} \to \varphi_{18}$\lbl{A2}}
\BIC{$\varphi_{18}$}
\AXC{$\varphi_{18} \to \varphi_{42}$\lbl{A1}}
\BIC{$\varphi_{42}$}
\BIC{$(m \in \mathbb{N} \to a \in \mathbb{R} \to \underline{m} \in \mathbb{R}) \to \varphi_{14}$}
\end{ptree}
\noindent$\Pi_{11}$:
\begin{ptree}
\AXC{$y = z \to x = y \to x = z$\lbl{A6}}
\AXC{$(y = z \to x = y \to x = z) \to x = y \to y = z \to x = z$\lbl{A1}}
\BIC{$x = y \to y = z \to x = z$}
\UIC{$\forall z\,(x = y \to y = z \to x = z)$}
\UIC{$\forall y\,\forall z\,(x = y \to y = z \to x = z)$}
\UIC{$\forall x\,\forall y\,\forall z\,(x = y \to y = z \to x = z)$}
\AXC{$\varphi_{8}$\lbl{A2}}
\BIC{$\forall y\,\forall z\,(\underline{S(m)} \cdot a = y \to y = z \to \underline{S(m)} \cdot a = z)$}
\AXC{$\varphi_{37}$\lbl{A2}}
\BIC{$\varphi_{15}$}
\end{ptree}
\noindent$\Pi_{12}$:
\begin{ptree}
\AXC{$\forall x\,\forall y\,(x \in \mathbb{R} \to y \in \mathbb{R} \to x + y \in \mathbb{R})$\lbl{(R1)}}
\AXC{$\varphi_{33}$\lbl{A2}}
\BIC{$\forall y\,(\underline{m} \in \mathbb{R} \to y \in \mathbb{R} \to \underline{m} + y \in \mathbb{R})$}
\AXC{$\varphi_{34}$\lbl{A2}}
\BIC{$\underline{m} \in \mathbb{R} \to 1 \in \mathbb{R} \to \underline{m} + 1 \in \mathbb{R}$}
\AXC{$(\underline{m} \in \mathbb{R} \to 1 \in \mathbb{R} \to \underline{m} + 1 \in \mathbb{R}) \to \varphi_{32}$\lbl{A1}}
\BIC{$\varphi_{32}$}
\end{ptree}
\noindent$\Pi_{13}$:
\begin{ptree}
\AXC{$\forall x\,\forall y\,(x \in \mathbb{R} \to y \in \mathbb{R} \to x \cdot y = y \cdot x)$\lbl{(R7)}}
\AXC{$\forall x\,\forall y\,(x \in \mathbb{R} \to y \in \mathbb{R} \to x \cdot y = y \cdot x) \to \varphi_{7}$\lbl{A2}}
\BIC{$\varphi_{7}$}
\AXC{$\varphi_{40}$\lbl{A2}}
\BIC{$\underline{m} + 1 \in \mathbb{R} \to a \in \mathbb{R} \to (\underline{m} + 1) \cdot a = a \cdot (\underline{m} + 1)$}
\AXC{$\varphi_{49}$\lbl{A1}}
\BIC{$(m \in \mathbb{N} \to a \in \mathbb{R} \to \underline{m} + 1 \in \mathbb{R}) \to \varphi_{28}$}
\end{ptree}
\noindent$\Pi_{14}$:
\begin{ptree}
\AXC{$1 \in \mathbb{R}$\lbl{(R1)}}
\AXC{$\stackrel{\Pi_{6}}{\vdots}$}\noLine
\UIC{$m \in \mathbb{N} \to a \in \mathbb{R} \to \underline{m} \in \mathbb{R}$}
\AXC{$\stackrel{\Pi_{12}}{\vdots}$}\noLine
\UIC{$\varphi_{32}$}
\BIC{$1 \in \mathbb{R} \to m \in \mathbb{N} \to a \in \mathbb{R} \to \underline{m} + 1 \in \mathbb{R}$}
\BIC{$m \in \mathbb{N} \to a \in \mathbb{R} \to \underline{m} + 1 \in \mathbb{R}$}
\AXC{$\stackrel{\Pi_{13}}{\vdots}$}\noLine
\UIC{$(m \in \mathbb{N} \to a \in \mathbb{R} \to \underline{m} + 1 \in \mathbb{R}) \to \varphi_{28}$}
\BIC{$\varphi_{28}$}
\end{ptree}
\noindent$\Pi_{15}$:
\begin{ptree}
\AXC{$m \in \mathbb{N} \to a \in \mathbb{R} \to m \in \mathbb{N}$\lbl{A1}}
\AXC{$\forall n \in \mathbb{N}\,(\underline{S(n)} = \underline{n} + 1)$\lbl{[N6]}}
\AXC{$\varphi_{10}$\lbl{A2}}
\BIC{$m \in \mathbb{N} \to \underline{S(m)} = \underline{m} + 1$}
\AXC{$(m \in \mathbb{N} \to \underline{S(m)} = \underline{m} + 1) \to \varphi_{19}$\lbl{A1}}
\BIC{$\varphi_{19}$}
\BIC{$m \in \mathbb{N} \to a \in \mathbb{R} \to \underline{S(m)} = \underline{m} + 1$}
\end{ptree}
\noindent$\Pi_{16}$:
\begin{ptree}
\AXC{$x = x$\lbl{A5}}
\UIC{$\forall x\,(x = x)$}
\AXC{$\forall x\,(x = x) \to \underline{S(m)} \cdot a = \underline{S(m)} \cdot a$\lbl{A2}}
\BIC{$\underline{S(m)} \cdot a = \underline{S(m)} \cdot a$}
\AXC{$\stackrel{\Pi_{15}}{\vdots}$}\noLine
\UIC{$m \in \mathbb{N} \to a \in \mathbb{R} \to \underline{S(m)} = \underline{m} + 1$}
\AXC{$\varphi_{51}$\lbl{A1}}
\BIC{$\underline{S(m)} \cdot a = \underline{S(m)} \cdot a \to \varphi_{41}$}
\BIC{$\varphi_{41}$}
\end{ptree}
\noindent$\Pi_{17}$:
\begin{ptree}
\AXC{$z = w \to \underline{S(m)} \cdot a = z \cdot a \to \underline{S(m)} \cdot a = w \cdot a$\lbl{A6}}
\UIC{$\varphi_{1}$}
\UIC{$\forall z\,(\varphi_{1})$}
\AXC{$\forall z\,(\varphi_{1}) \to \varphi_{6}$\lbl{A2}}
\BIC{$\varphi_{6}$}
\AXC{$\varphi_{6} \to \underline{S(m)} = \underline{m} + 1 \to \varphi_{2}$\lbl{A2}}
\BIC{$\underline{S(m)} = \underline{m} + 1 \to \varphi_{2}$}
\AXC{$\stackrel{\Pi_{16}}{\vdots}$}\noLine
\UIC{$\varphi_{41}$}
\BIC{$m \in \mathbb{N} \to a \in \mathbb{R} \to \underline{S(m)} \cdot a = (\underline{m} + 1) \cdot a$}
\end{ptree}
\noindent$\Pi_{18}$:
\begin{ptree}
\AXC{$\stackrel{\Pi_{17}}{\vdots}$}\noLine
\UIC{$m \in \mathbb{N} \to a \in \mathbb{R} \to \underline{S(m)} \cdot a = (\underline{m} + 1) \cdot a$}
\AXC{$(m \in \mathbb{N} \to a \in \mathbb{R} \to \underline{S(m)} \cdot a = (\underline{m} + 1) \cdot a) \to \varphi_{50}$\lbl{A1}}
\BIC{$\varphi_{50}$}
\end{ptree}
\noindent$\Pi_{19}$:
\begin{ptree}
\AXC{$\stackrel{\Pi_{11}}{\vdots}$}\noLine
\UIC{$\varphi_{15}$}
\AXC{$\varphi_{15} \to \underline{S(m)} \cdot a = (\underline{m} + 1) \cdot a \to \varphi_{4}$\lbl{A2}}
\BIC{$\underline{S(m)} \cdot a = (\underline{m} + 1) \cdot a \to \varphi_{4}$}
\AXC{$m \in \mathbb{N} \to a \in \mathbb{R} \to a \in \mathbb{R}$\lbl{A1}}
\AXC{$\stackrel{\Pi_{14}}{\vdots}$}\noLine
\UIC{$\varphi_{28}$}
\BIC{$m \in \mathbb{N} \to a \in \mathbb{R} \to (\underline{m} + 1) \cdot a = a \cdot (\underline{m} + 1)$}
\AXC{$\stackrel{\Pi_{18}}{\vdots}$}\noLine
\UIC{$\varphi_{50}$}
\BIC{$\varphi_{44}$}
\BIC{$m \in \mathbb{N} \to a \in \mathbb{R} \to \underline{S(m)} \cdot a = a \cdot (\underline{m} + 1)$}
\end{ptree}
\noindent$\Pi_{20}$:
\begin{ptree}
\AXC{$\stackrel{\Pi_{19}}{\vdots}$}\noLine
\UIC{$m \in \mathbb{N} \to a \in \mathbb{R} \to \underline{S(m)} \cdot a = a \cdot (\underline{m} + 1)$}
\AXC{$(m \in \mathbb{N} \to a \in \mathbb{R} \to \underline{S(m)} \cdot a = a \cdot (\underline{m} + 1)) \to \varphi_{54}$\lbl{A1}}
\BIC{$\varphi_{54}$}
\end{ptree}
\noindent$\Pi_{21}$:
\begin{ptree}
\AXC{$1 \in \mathbb{R}$\lbl{(R1)}}
\AXC{$\stackrel{\Pi_{6}}{\vdots}$}\noLine
\UIC{$m \in \mathbb{N} \to a \in \mathbb{R} \to \underline{m} \in \mathbb{R}$}
\AXC{$\stackrel{\Pi_{10}}{\vdots}$}\noLine
\UIC{$(m \in \mathbb{N} \to a \in \mathbb{R} \to \underline{m} \in \mathbb{R}) \to \varphi_{14}$}
\BIC{$\varphi_{14}$}
\BIC{$m \in \mathbb{N} \to a \in \mathbb{R} \to a \cdot (\underline{m} + 1) = a \cdot \underline{m} + a \cdot 1$}
\AXC{$\stackrel{\Pi_{20}}{\vdots}$}\noLine
\UIC{$\varphi_{54}$}
\BIC{$\varphi_{45}$}
\end{ptree}
\noindent$\Pi_{22}$:
\begin{ptree}
\AXC{$\stackrel{\Pi_{9}}{\vdots}$}\noLine
\UIC{$\varphi_{13}$}
\AXC{$\varphi_{13} \to \underline{S(m)} \cdot a = a \cdot (\underline{m} + 1) \to \varphi_{16}$\lbl{A2}}
\BIC{$\underline{S(m)} \cdot a = a \cdot (\underline{m} + 1) \to \varphi_{16}$}
\AXC{$\stackrel{\Pi_{21}}{\vdots}$}\noLine
\UIC{$\varphi_{45}$}
\BIC{$m \in \mathbb{N} \to a \in \mathbb{R} \to \underline{S(m)} \cdot a = a \cdot \underline{m} + a \cdot 1$}
\AXC{$(m \in \mathbb{N} \to a \in \mathbb{R} \to \underline{S(m)} \cdot a = a \cdot \underline{m} + a \cdot 1) \to \varphi_{55}$\lbl{A1}}
\BIC{$\varphi_{55}$}
\end{ptree}
\noindent$\Pi_{23}$:
\begin{ptree}
\AXC{$\varphi_{11}$\lbl{A6}}
\UIC{$\forall w\,(\varphi_{11})$}
\UIC{$\forall z\,\forall w\,(\varphi_{11})$}
\AXC{$\forall z\,\forall w\,(\varphi_{11}) \to \varphi_{30}$\lbl{A2}}
\BIC{$\varphi_{30}$}
\AXC{$\varphi_{30} \to a \cdot \underline{m} = \underline{m} \cdot a \to \varphi_{26}$\lbl{A2}}
\BIC{$a \cdot \underline{m} = \underline{m} \cdot a \to \varphi_{26}$}
\AXC{$\stackrel{\Pi_{8}}{\vdots}$}\noLine
\UIC{$\varphi_{48}$}
\BIC{$m \in \mathbb{N} \to a \in \mathbb{R} \to a \cdot \underline{m} + a \cdot 1 = \underline{m} \cdot a + a \cdot 1$}
\AXC{$\stackrel{\Pi_{22}}{\vdots}$}\noLine
\UIC{$\varphi_{55}$}
\BIC{$\varphi_{47}$}
\end{ptree}
\noindent$\Pi_{24}$:
\begin{ptree}
\AXC{$\stackrel{\Pi_{4}}{\vdots}$}\noLine
\UIC{$\varphi_{22}$}
\AXC{$\varphi_{22} \to \underline{S(m)} \cdot a = a \cdot \underline{m} + a \cdot 1 \to \varphi_{20}$\lbl{A2}}
\BIC{$\underline{S(m)} \cdot a = a \cdot \underline{m} + a \cdot 1 \to \varphi_{20}$}
\AXC{$\stackrel{\Pi_{23}}{\vdots}$}\noLine
\UIC{$\varphi_{47}$}
\BIC{$m \in \mathbb{N} \to a \in \mathbb{R} \to \underline{S(m)} \cdot a = \underline{m} \cdot a + a \cdot 1$}
\AXC{$(m \in \mathbb{N} \to a \in \mathbb{R} \to \underline{S(m)} \cdot a = \underline{m} \cdot a + a \cdot 1) \to \varphi_{52}$\lbl{A1}}
\BIC{$\varphi_{52}$}
\end{ptree}
\noindent 主木:
\begin{ptree}
\AXC{$\stackrel{\Pi_{1}}{\vdots}$}\noLine
\UIC{$\varphi_{35}$}
\AXC{$z = w \to \varphi_{17}$\lbl{A6}}
\UIC{$\forall w\,(z = w \to \varphi_{17})$}
\UIC{$\forall z\,\forall w\,(z = w \to \varphi_{17})$}
\AXC{$\forall z\,\forall w\,(z = w \to \varphi_{17}) \to \varphi_{29}$\lbl{A2}}
\BIC{$\varphi_{29}$}
\AXC{$\varphi_{29} \to a \cdot 1 = a \to \varphi_{24}$\lbl{A2}}
\BIC{$a \cdot 1 = a \to \varphi_{24}$}
\AXC{$\stackrel{\Pi_{3}}{\vdots}$}\noLine
\UIC{$\varphi_{43}$}
\BIC{$m \in \mathbb{N} \to a \in \mathbb{R} \to \underline{m} \cdot a + a \cdot 1 = \underline{m} \cdot a + a$}
\AXC{$\stackrel{\Pi_{24}}{\vdots}$}\noLine
\UIC{$\varphi_{52}$}
\BIC{$\varphi_{35} \to m \in \mathbb{N} \to a \in \mathbb{R} \to \underline{S(m)} \cdot a = \underline{m} \cdot a + a$}
\BIC{$m \in \mathbb{N} \to a \in \mathbb{R} \to \underline{S(m)} \cdot a = \underline{m} \cdot a + a$}
\UIC{$\forall a\,(m \in \mathbb{N} \to a \in \mathbb{R} \to \underline{S(m)} \cdot a = \underline{m} \cdot a + a)$}
\UIC{$\forall m\,\forall a\,(m \in \mathbb{N} \to a \in \mathbb{R} \to \underline{S(m)} \cdot a = \underline{m} \cdot a + a)$}
\end{ptree}
\begin{abbrevlist}
\item $\varphi_{1} :\equiv \forall w\,(z = w \to \underline{S(m)} \cdot a = z \cdot a \to \underline{S(m)} \cdot a = w \cdot a)$
\item $\varphi_{2} :\equiv \underline{S(m)} \cdot a = \underline{S(m)} \cdot a \to \underline{S(m)} \cdot a = (\underline{m} + 1) \cdot a$
\item $\varphi_{3} :\equiv (m \in \mathbb{N} \to a \in \mathbb{R} \to a \in \mathbb{R}) \to m \in \mathbb{N} \to a \in \mathbb{R} \to a \cdot 1 = a$
\item $\varphi_{4} :\equiv (\underline{m} + 1) \cdot a = a \cdot (\underline{m} + 1) \to \underline{S(m)} \cdot a = a \cdot (\underline{m} + 1)$
\item $\varphi_{5} :\equiv \forall y\,\forall z\,(a \in \mathbb{R} \to y \in \mathbb{R} \to z \in \mathbb{R} \to a \cdot (y + z) = a \cdot y + a \cdot z)$
\item $\varphi_{6} :\equiv \forall w\,(\underline{S(m)} = w \to \underline{S(m)} \cdot a = \underline{S(m)} \cdot a \to \underline{S(m)} \cdot a = w \cdot a)$
\item $\varphi_{7} :\equiv \forall y\,(\underline{m} + 1 \in \mathbb{R} \to y \in \mathbb{R} \to (\underline{m} + 1) \cdot y = y \cdot (\underline{m} + 1))$
\item $\varphi_{8} :\equiv \forall x\,\forall y\,\forall z\,(x = y \to y = z \to x = z) \to \forall y\,\forall z\,(\underline{S(m)} \cdot a = y \to y = z \to \underline{S(m)} \cdot a = z)$
\item $\varphi_{9} :\equiv \forall x\,\forall y\,\forall z\,(x \in \mathbb{R} \to y \in \mathbb{R} \to z \in \mathbb{R} \to x \cdot (y + z) = x \cdot y + x \cdot z)$
\item $\varphi_{10} :\equiv \forall n \in \mathbb{N}\,(\underline{S(n)} = \underline{n} + 1) \to m \in \mathbb{N} \to \underline{S(m)} = \underline{m} + 1$
\item $\varphi_{11} :\equiv z = w \to a \cdot \underline{m} + a \cdot 1 = z + a \cdot 1 \to a \cdot \underline{m} + a \cdot 1 = w + a \cdot 1$
\item $\varphi_{12} :\equiv (m \in \mathbb{N} \to a \in \mathbb{R} \to m \in \mathbb{N}) \to m \in \mathbb{N} \to a \in \mathbb{R} \to \underline{m} \in \mathbb{R}$
\item $\varphi_{13} :\equiv \forall z\,(\underline{S(m)} \cdot a = a \cdot (\underline{m} + 1) \to a \cdot (\underline{m} + 1) = z \to \underline{S(m)} \cdot a = z)$
\item $\varphi_{14} :\equiv 1 \in \mathbb{R} \to m \in \mathbb{N} \to a \in \mathbb{R} \to a \cdot (\underline{m} + 1) = a \cdot \underline{m} + a \cdot 1$
\item $\varphi_{15} :\equiv \forall z\,(\underline{S(m)} \cdot a = (\underline{m} + 1) \cdot a \to (\underline{m} + 1) \cdot a = z \to \underline{S(m)} \cdot a = z)$
\item $\varphi_{16} :\equiv a \cdot (\underline{m} + 1) = a \cdot \underline{m} + a \cdot 1 \to \underline{S(m)} \cdot a = a \cdot \underline{m} + a \cdot 1$
\item $\varphi_{17} :\equiv \underline{m} \cdot a + a \cdot 1 = \underline{m} \cdot a + z \to \underline{m} \cdot a + a \cdot 1 = \underline{m} \cdot a + w$
\item $\varphi_{18} :\equiv a \in \mathbb{R} \to \underline{m} \in \mathbb{R} \to 1 \in \mathbb{R} \to a \cdot (\underline{m} + 1) = a \cdot \underline{m} + a \cdot 1$
\item $\varphi_{19} :\equiv (m \in \mathbb{N} \to a \in \mathbb{R} \to m \in \mathbb{N}) \to m \in \mathbb{N} \to a \in \mathbb{R} \to \underline{S(m)} = \underline{m} + 1$
\item $\varphi_{20} :\equiv a \cdot \underline{m} + a \cdot 1 = \underline{m} \cdot a + a \cdot 1 \to \underline{S(m)} \cdot a = \underline{m} \cdot a + a \cdot 1$
\item $\varphi_{21} :\equiv \forall z\,(\underline{S(m)} \cdot a = \underline{m} \cdot a + a \cdot 1 \to \underline{m} \cdot a + a \cdot 1 = z \to \underline{S(m)} \cdot a = z)$
\item $\varphi_{22} :\equiv \forall z\,(\underline{S(m)} \cdot a = a \cdot \underline{m} + a \cdot 1 \to a \cdot \underline{m} + a \cdot 1 = z \to \underline{S(m)} \cdot a = z)$
\item $\varphi_{23} :\equiv \forall z\,(a \in \mathbb{R} \to \underline{m} \in \mathbb{R} \to z \in \mathbb{R} \to a \cdot (\underline{m} + z) = a \cdot \underline{m} + a \cdot z)$
\item $\varphi_{24} :\equiv \underline{m} \cdot a + a \cdot 1 = \underline{m} \cdot a + a \cdot 1 \to \underline{m} \cdot a + a \cdot 1 = \underline{m} \cdot a + a$
\item $\varphi_{25} :\equiv \forall x\,\forall y\,(x \in \mathbb{R} \to y \in \mathbb{R} \to x \cdot y = y \cdot x) \to \forall y\,(a \in \mathbb{R} \to y \in \mathbb{R} \to a \cdot y = y \cdot a)$
\item $\varphi_{26} :\equiv a \cdot \underline{m} + a \cdot 1 = a \cdot \underline{m} + a \cdot 1 \to a \cdot \underline{m} + a \cdot 1 = \underline{m} \cdot a + a \cdot 1$
\item $\varphi_{27} :\equiv (m \in \mathbb{N} \to a \in \mathbb{R} \to \underline{m} \in \mathbb{R}) \to m \in \mathbb{N} \to a \in \mathbb{R} \to a \cdot \underline{m} = \underline{m} \cdot a$
\item $\varphi_{28} :\equiv (m \in \mathbb{N} \to a \in \mathbb{R} \to a \in \mathbb{R}) \to m \in \mathbb{N} \to a \in \mathbb{R} \to (\underline{m} + 1) \cdot a = a \cdot (\underline{m} + 1)$
\item $\varphi_{29} :\equiv \forall w\,(a \cdot 1 = w \to \underline{m} \cdot a + a \cdot 1 = \underline{m} \cdot a + a \cdot 1 \to \underline{m} \cdot a + a \cdot 1 = \underline{m} \cdot a + w)$
\item $\varphi_{30} :\equiv \forall w\,(a \cdot \underline{m} = w \to a \cdot \underline{m} + a \cdot 1 = a \cdot \underline{m} + a \cdot 1 \to a \cdot \underline{m} + a \cdot 1 = w + a \cdot 1)$
\item $\varphi_{31} :\equiv \forall y\,(a \in \mathbb{R} \to y \in \mathbb{R} \to a \cdot y = y \cdot a) \to a \in \mathbb{R} \to \underline{m} \in \mathbb{R} \to a \cdot \underline{m} = \underline{m} \cdot a$
\item $\varphi_{32} :\equiv (m \in \mathbb{N} \to a \in \mathbb{R} \to \underline{m} \in \mathbb{R}) \to 1 \in \mathbb{R} \to m \in \mathbb{N} \to a \in \mathbb{R} \to \underline{m} + 1 \in \mathbb{R}$
\item $\varphi_{33} :\equiv \forall x\,\forall y\,(x \in \mathbb{R} \to y \in \mathbb{R} \to x + y \in \mathbb{R}) \to \forall y\,(\underline{m} \in \mathbb{R} \to y \in \mathbb{R} \to \underline{m} + y \in \mathbb{R})$
\item $\varphi_{34} :\equiv \forall y\,(\underline{m} \in \mathbb{R} \to y \in \mathbb{R} \to \underline{m} + y \in \mathbb{R}) \to \underline{m} \in \mathbb{R} \to 1 \in \mathbb{R} \to \underline{m} + 1 \in \mathbb{R}$
\item $\varphi_{35} :\equiv \underline{S(m)} \cdot a = \underline{m} \cdot a + a \cdot 1 \to \underline{m} \cdot a + a \cdot 1 = \underline{m} \cdot a + a \to \underline{S(m)} \cdot a = \underline{m} \cdot a + a$
\item $\varphi_{36} :\equiv \forall y\,\forall z\,(\underline{S(m)} \cdot a = y \to y = z \to \underline{S(m)} \cdot a = z) \to \varphi_{13}$
\item $\varphi_{37} :\equiv \forall y\,\forall z\,(\underline{S(m)} \cdot a = y \to y = z \to \underline{S(m)} \cdot a = z) \to \varphi_{15}$
\item $\varphi_{38} :\equiv \forall y\,\forall z\,(\underline{S(m)} \cdot a = y \to y = z \to \underline{S(m)} \cdot a = z) \to \varphi_{21}$
\item $\varphi_{39} :\equiv \forall y\,\forall z\,(\underline{S(m)} \cdot a = y \to y = z \to \underline{S(m)} \cdot a = z) \to \varphi_{22}$
\item $\varphi_{40} :\equiv \varphi_{7} \to \underline{m} + 1 \in \mathbb{R} \to a \in \mathbb{R} \to (\underline{m} + 1) \cdot a = a \cdot (\underline{m} + 1)$
\item $\varphi_{41} :\equiv (\underline{S(m)} = \underline{m} + 1 \to \varphi_{2}) \to m \in \mathbb{N} \to a \in \mathbb{R} \to \underline{S(m)} \cdot a = (\underline{m} + 1) \cdot a$
\item $\varphi_{42} :\equiv (m \in \mathbb{N} \to a \in \mathbb{R} \to a \in \mathbb{R}) \to (m \in \mathbb{N} \to a \in \mathbb{R} \to \underline{m} \in \mathbb{R}) \to \varphi_{14}$
\item $\varphi_{43} :\equiv (a \cdot 1 = a \to \varphi_{24}) \to m \in \mathbb{N} \to a \in \mathbb{R} \to \underline{m} \cdot a + a \cdot 1 = \underline{m} \cdot a + a$
\item $\varphi_{44} :\equiv (\underline{S(m)} \cdot a = (\underline{m} + 1) \cdot a \to \varphi_{4}) \to m \in \mathbb{N} \to a \in \mathbb{R} \to \underline{S(m)} \cdot a = a \cdot (\underline{m} + 1)$
\item $\varphi_{45} :\equiv (\underline{S(m)} \cdot a = a \cdot (\underline{m} + 1) \to \varphi_{16}) \to m \in \mathbb{N} \to a \in \mathbb{R} \to \underline{S(m)} \cdot a = a \cdot \underline{m} + a \cdot 1$
\item $\varphi_{46} :\equiv (a \in \mathbb{R} \to \underline{m} \in \mathbb{R} \to a \cdot \underline{m} = \underline{m} \cdot a) \to (m \in \mathbb{N} \to a \in \mathbb{R} \to a \in \mathbb{R}) \to \varphi_{27}$
\item $\varphi_{47} :\equiv (\underline{S(m)} \cdot a = a \cdot \underline{m} + a \cdot 1 \to \varphi_{20}) \to m \in \mathbb{N} \to a \in \mathbb{R} \to \underline{S(m)} \cdot a = \underline{m} \cdot a + a \cdot 1$
\item $\varphi_{48} :\equiv (a \cdot \underline{m} = \underline{m} \cdot a \to \varphi_{26}) \to m \in \mathbb{N} \to a \in \mathbb{R} \to a \cdot \underline{m} + a \cdot 1 = \underline{m} \cdot a + a \cdot 1$
\item $\varphi_{49} :\equiv (\underline{m} + 1 \in \mathbb{R} \to a \in \mathbb{R} \to (\underline{m} + 1) \cdot a = a \cdot (\underline{m} + 1)) \to (m \in \mathbb{N} \to a \in \mathbb{R} \to \underline{m} + 1 \in \mathbb{R}) \to \varphi_{28}$
\item $\varphi_{50} :\equiv (m \in \mathbb{N} \to a \in \mathbb{R} \to (\underline{m} + 1) \cdot a = a \cdot (\underline{m} + 1)) \to \varphi_{44}$
\item $\varphi_{51} :\equiv (m \in \mathbb{N} \to a \in \mathbb{R} \to \underline{S(m)} = \underline{m} + 1) \to \underline{S(m)} \cdot a = \underline{S(m)} \cdot a \to \varphi_{41}$
\item $\varphi_{52} :\equiv (m \in \mathbb{N} \to a \in \mathbb{R} \to \underline{m} \cdot a + a \cdot 1 = \underline{m} \cdot a + a) \to \varphi_{35} \to m \in \mathbb{N} \to a \in \mathbb{R} \to \underline{S(m)} \cdot a = \underline{m} \cdot a + a$
\item $\varphi_{53} :\equiv (m \in \mathbb{N} \to a \in \mathbb{R} \to a \cdot 1 = a) \to \underline{m} \cdot a + a \cdot 1 = \underline{m} \cdot a + a \cdot 1 \to \varphi_{43}$
\item $\varphi_{54} :\equiv (m \in \mathbb{N} \to a \in \mathbb{R} \to a \cdot (\underline{m} + 1) = a \cdot \underline{m} + a \cdot 1) \to \varphi_{45}$
\item $\varphi_{55} :\equiv (m \in \mathbb{N} \to a \in \mathbb{R} \to a \cdot \underline{m} + a \cdot 1 = \underline{m} \cdot a + a \cdot 1) \to \varphi_{47}$
\item $\varphi_{56} :\equiv (m \in \mathbb{N} \to a \in \mathbb{R} \to a \cdot \underline{m} = \underline{m} \cdot a) \to a \cdot \underline{m} + a \cdot 1 = a \cdot \underline{m} + a \cdot 1 \to \varphi_{48}$
\end{abbrevlist}
\end{fullproof}

% 自動生成 (prover/build.js): Archimedes  ノード数 544
\begin{fullproof}{命題3.2}{$\forall a\,\forall b\,(a \in \mathbb{R} \to b \in \mathbb{R} \to 0 < a \to 0 < b \to \exists n \in \mathbb{N}\,(b < \underline{n} \cdot a))$}{544}\phantomsection\label{tree:Archimedes}
\noindent$\Pi_{1}$:
\begin{ptree}
\AXC{$\forall x\,(\neg x < x)$\lbl{[O1]}}
\AXC{$\forall x\,(\neg x < x) \to \neg 0 < 0$\lbl{A2}}
\BIC{$\neg 0 < 0$}
\AXC{$\neg 0 < 0 \to a \in \mathbb{R} \to \varphi_{69}$\lbl{A1}}
\BIC{$a \in \mathbb{R} \to \varphi_{69}$}
\end{ptree}
\noindent$\Pi_{2}$:
\begin{ptree}
\AXC{$\forall n \in \mathbb{N}\,(S(n) \in \mathbb{N})$\lbl{[N2]}}
\AXC{$\forall n \in \mathbb{N}\,(S(n) \in \mathbb{N}) \to m \in \mathbb{N} \to S(m) \in \mathbb{N}$\lbl{A2}}
\BIC{$m \in \mathbb{N} \to S(m) \in \mathbb{N}$}
\AXC{$\varphi_{251}$\lbl{A1}}
\BIC{$\varphi_{203} \to a \in \mathbb{R} \to b \in \mathbb{R} \to 0 < a \to \varphi_{166}$}
\end{ptree}
\noindent$\Pi_{3}$:
\begin{ptree}
\AXC{$a \in \mathbb{R} \to b \in \mathbb{R} \to 0 < a \to \varphi_{187}$\lbl{A1}}
\AXC{$(a \in \mathbb{R} \to b \in \mathbb{R} \to 0 < a \to \varphi_{187}) \to \varphi_{203}$\lbl{A1}}
\BIC{$\varphi_{203}$}
\AXC{$\stackrel{\Pi_{2}}{\vdots}$}\noLine
\UIC{$\varphi_{203} \to a \in \mathbb{R} \to b \in \mathbb{R} \to 0 < a \to \varphi_{166}$}
\BIC{$a \in \mathbb{R} \to b \in \mathbb{R} \to 0 < a \to \varphi_{166}$}
\end{ptree}
\noindent$\Pi_{4}$:
\begin{ptree}
\AXC{$\stackrel{\Pi_{3}}{\vdots}$}\noLine
\UIC{$a \in \mathbb{R} \to b \in \mathbb{R} \to 0 < a \to \varphi_{166}$}
\AXC{$(a \in \mathbb{R} \to b \in \mathbb{R} \to 0 < a \to \varphi_{166}) \to \varphi_{223}$\lbl{A1}}
\BIC{$\varphi_{223}$}
\end{ptree}
\noindent$\Pi_{5}$:
\begin{ptree}
\AXC{$\forall n \in \mathbb{N}\,(\underline{n} \in \mathbb{R})$\lbl{[N4]}}
\AXC{$\forall n \in \mathbb{N}\,(\underline{n} \in \mathbb{R}) \to S(m) \in \mathbb{N} \to \underline{S(m)} \in \mathbb{R}$\lbl{A2}}
\BIC{$S(m) \in \mathbb{N} \to \underline{S(m)} \in \mathbb{R}$}
\end{ptree}
\noindent$\Pi_{6}$:
\begin{ptree}
\AXC{$\forall x\,\forall y\,(x \in \mathbb{R} \to y \in \mathbb{R} \to x \cdot y \in \mathbb{R})$\lbl{(R1)}}
\AXC{$\varphi_{149}$\lbl{A2}}
\BIC{$\forall y\,(\underline{S(m)} \in \mathbb{R} \to y \in \mathbb{R} \to \underline{S(m)} \cdot y \in \mathbb{R})$}
\AXC{$\varphi_{157}$\lbl{A2}}
\BIC{$\underline{S(m)} \in \mathbb{R} \to a \in \mathbb{R} \to \underline{S(m)} \cdot a \in \mathbb{R}$}
\AXC{$(\underline{S(m)} \in \mathbb{R} \to a \in \mathbb{R} \to \underline{S(m)} \cdot a \in \mathbb{R}) \to \varphi_{267}$\lbl{A1}}
\BIC{$\varphi_{267}$}
\end{ptree}
\noindent$\Pi_{7}$:
\begin{ptree}
\AXC{$\stackrel{\Pi_{3}}{\vdots}$}\noLine
\UIC{$a \in \mathbb{R} \to b \in \mathbb{R} \to 0 < a \to \varphi_{166}$}
\AXC{$\stackrel{\Pi_{5}}{\vdots}$}\noLine
\UIC{$S(m) \in \mathbb{N} \to \underline{S(m)} \in \mathbb{R}$}
\AXC{$(S(m) \in \mathbb{N} \to \underline{S(m)} \in \mathbb{R}) \to \varphi_{246}$\lbl{A1}}
\BIC{$\varphi_{246}$}
\BIC{$a \in \mathbb{R} \to b \in \mathbb{R} \to 0 < a \to \varphi_{178}$}
\AXC{$\stackrel{\Pi_{6}}{\vdots}$}\noLine
\UIC{$\varphi_{267}$}
\BIC{$\varphi_{204} \to a \in \mathbb{R} \to b \in \mathbb{R} \to 0 < a \to \varphi_{186}$}
\end{ptree}
\noindent$\Pi_{8}$:
\begin{ptree}
\AXC{$\forall a\,\forall y\,(y \in X_{a} \leftrightarrow y \in \mathbb{R} \land \exists n \in \mathbb{N}\,(y = \underline{n} \cdot a))$\lbl{Def$X_a$}}
\AXC{$\varphi_{163}$\lbl{A2}}
\BIC{$\forall y\,(y \in X_{a} \leftrightarrow y \in \mathbb{R} \land \exists n \in \mathbb{N}\,(y = \underline{n} \cdot a))$}
\AXC{$\forall y\,(y \in X_{a} \leftrightarrow y \in \mathbb{R} \land \exists n \in \mathbb{N}\,(y = \underline{n} \cdot a)) \to \varphi_{141}$\lbl{A2}}
\BIC{$\varphi_{141}$}
\end{ptree}
\noindent$\Pi_{9}$:
\begin{ptree}
\AXC{$\varphi_{204}$\lbl{A1}}
\AXC{$\stackrel{\Pi_{7}}{\vdots}$}\noLine
\UIC{$\varphi_{204} \to a \in \mathbb{R} \to b \in \mathbb{R} \to 0 < a \to \varphi_{186}$}
\BIC{$a \in \mathbb{R} \to b \in \mathbb{R} \to 0 < a \to \varphi_{186}$}
\AXC{$\stackrel{\Pi_{8}}{\vdots}$}\noLine
\UIC{$\varphi_{141}$}
\AXC{$\varphi_{141} \to (a \in \mathbb{R} \to b \in \mathbb{R} \to 0 < a \to \varphi_{186}) \to \varphi_{255}$\lbl{A1}}
\BIC{$(a \in \mathbb{R} \to b \in \mathbb{R} \to 0 < a \to \varphi_{186}) \to \varphi_{255}$}
\BIC{$\varphi_{255}$}
\end{ptree}
\noindent$\Pi_{10}$:
\begin{ptree}
\AXC{$\varphi_{128}$\lbl{A2}}
\AXC{$x = x$\lbl{A5}}
\UIC{$\forall x\,(x = x)$}
\AXC{$\forall x\,(x = x) \to \underline{S(m)} \cdot a = \underline{S(m)} \cdot a$\lbl{A2}}
\BIC{$\underline{S(m)} \cdot a = \underline{S(m)} \cdot a$}
\AXC{$\stackrel{\Pi_{4}}{\vdots}$}\noLine
\UIC{$\varphi_{223}$}
\BIC{$a \in \mathbb{R} \to b \in \mathbb{R} \to 0 < a \to 0 < b \to \varphi_{198}$}
\AXC{$\varphi_{263}$\lbl{A1}}
\BIC{$\varphi_{128} \to a \in \mathbb{R} \to b \in \mathbb{R} \to 0 < a \to 0 < b \to \varphi_{194}$}
\BIC{$a \in \mathbb{R} \to b \in \mathbb{R} \to 0 < a \to 0 < b \to \varphi_{194}$}
\AXC{$\stackrel{\Pi_{9}}{\vdots}$}\noLine
\UIC{$\varphi_{255}$}
\BIC{$a \in \mathbb{R} \to b \in \mathbb{R} \to 0 < a \to \varphi_{189}$}
\end{ptree}
\noindent$\Pi_{11}$:
\begin{ptree}
\AXC{$\forall s\,(\varphi_{1})$\lbl{DefSup}}
\AXC{$\forall s\,(\varphi_{1}) \to \varphi_{1}$\lbl{A2}}
\BIC{$\varphi_{1}$}
\AXC{$\varphi_{1} \to \varphi_{14}$\lbl{A2}}
\BIC{$\varphi_{14}$}
\AXC{$a \in \mathbb{R} \to b \in \mathbb{R} \to 0 < a \to 0 < b \to \varphi_{64}$\lbl{A1}}
\AXC{$\varphi_{238}$\lbl{A1}}
\BIC{$\varphi_{14} \to a \in \mathbb{R} \to b \in \mathbb{R} \to 0 < a \to 0 < b \to \varphi_{52}$}
\BIC{$a \in \mathbb{R} \to b \in \mathbb{R} \to 0 < a \to 0 < b \to \varphi_{52}$}
\end{ptree}
\noindent$\Pi_{12}$:
\begin{ptree}
\AXC{$\stackrel{\Pi_{11}}{\vdots}$}\noLine
\UIC{$a \in \mathbb{R} \to b \in \mathbb{R} \to 0 < a \to 0 < b \to \varphi_{52}$}
\AXC{$(a \in \mathbb{R} \to b \in \mathbb{R} \to 0 < a \to 0 < b \to \varphi_{52}) \to \varphi_{222}$\lbl{A1}}
\BIC{$\varphi_{222}$}
\end{ptree}
\noindent$\Pi_{13}$:
\begin{ptree}
\AXC{$\varphi_{67}$\lbl{A2}}
\AXC{$\forall a\,\forall X\,(\mathrm{UB}(a,X) \leftrightarrow a \in \mathbb{R} \land \forall x \in X\,(x \leq a))$\lbl{DefUB}}
\AXC{$\varphi_{116}$\lbl{A2}}
\BIC{$\forall X\,(\mathrm{UB}(s,X) \leftrightarrow s \in \mathbb{R} \land \forall x \in X\,(x \leq s))$}
\AXC{$\varphi_{136}$\lbl{A2}}
\BIC{$\mathrm{UB}(s,X_{a}) \leftrightarrow s \in \mathbb{R} \land \forall x \in X_{a}\,(x \leq s)$}
\AXC{$\stackrel{\Pi_{12}}{\vdots}$}\noLine
\UIC{$\varphi_{222}$}
\BIC{$a \in \mathbb{R} \to b \in \mathbb{R} \to 0 < a \to \varphi_{185}$}
\AXC{$\varphi_{261}$\lbl{A1}}
\BIC{$\varphi_{67} \to a \in \mathbb{R} \to b \in \mathbb{R} \to 0 < a \to 0 < b \to \varphi_{199}$}
\BIC{$a \in \mathbb{R} \to b \in \mathbb{R} \to 0 < a \to 0 < b \to \varphi_{199}$}
\end{ptree}
\noindent$\Pi_{14}$:
\begin{ptree}
\AXC{$\stackrel{\Pi_{10}}{\vdots}$}\noLine
\UIC{$a \in \mathbb{R} \to b \in \mathbb{R} \to 0 < a \to \varphi_{189}$}
\AXC{$\stackrel{\Pi_{13}}{\vdots}$}\noLine
\UIC{$a \in \mathbb{R} \to b \in \mathbb{R} \to 0 < a \to 0 < b \to \varphi_{199}$}
\AXC{$(a \in \mathbb{R} \to b \in \mathbb{R} \to 0 < a \to 0 < b \to \varphi_{199}) \to \varphi_{249}$\lbl{A1}}
\BIC{$\varphi_{249}$}
\BIC{$a \in \mathbb{R} \to b \in \mathbb{R} \to 0 < a \to \varphi_{175}$}
\end{ptree}
\noindent$\Pi_{15}$:
\begin{ptree}
\AXC{$z = w \to s < z \to s < w$\lbl{A6}}
\UIC{$\forall w\,(z = w \to s < z \to s < w)$}
\UIC{$\forall z\,\forall w\,(z = w \to s < z \to s < w)$}
\AXC{$\varphi_{109}$\lbl{A2}}
\BIC{$\forall w\,(\underline{m} \cdot a + a = w \to s < \underline{m} \cdot a + a \to s < w)$}
\AXC{$\forall w\,(\underline{m} \cdot a + a = w \to s < \underline{m} \cdot a + a \to s < w) \to \varphi_{90}$\lbl{A2}}
\BIC{$\varphi_{90}$}
\end{ptree}
\noindent$\Pi_{16}$:
\begin{ptree}
\AXC{$\varphi_{191}$\lbl{A1}}
\AXC{$\varphi_{191} \to a \in \mathbb{R} \to b \in \mathbb{R} \to 0 < a \to 0 < b \to \varphi_{139}$\lbl{A1}}
\BIC{$a \in \mathbb{R} \to b \in \mathbb{R} \to 0 < a \to 0 < b \to \varphi_{139}$}
\AXC{$(a \in \mathbb{R} \to b \in \mathbb{R} \to 0 < a \to 0 < b \to \varphi_{139}) \to \varphi_{210}$\lbl{A1}}
\BIC{$\varphi_{210}$}
\end{ptree}
\noindent$\Pi_{17}$:
\begin{ptree}
\AXC{$\forall a\,\forall y\,(y \in X_{a} \leftrightarrow y \in \mathbb{R} \land \exists n \in \mathbb{N}\,(y = \underline{n} \cdot a))$\lbl{Def$X_a$}}
\AXC{$\varphi_{163}$\lbl{A2}}
\BIC{$\forall y\,(y \in X_{a} \leftrightarrow y \in \mathbb{R} \land \exists n \in \mathbb{N}\,(y = \underline{n} \cdot a))$}
\AXC{$\varphi_{158}$\lbl{A2}}
\BIC{$y \in X_{a} \leftrightarrow y \in \mathbb{R} \land \exists n \in \mathbb{N}\,(y = \underline{n} \cdot a)$}
\AXC{$\stackrel{\Pi_{16}}{\vdots}$}\noLine
\UIC{$\varphi_{210}$}
\BIC{$a \in \mathbb{R} \to b \in \mathbb{R} \to 0 < a \to 0 < b \to \varphi_{131}$}
\end{ptree}
\noindent$\Pi_{18}$:
\begin{ptree}
\AXC{$\stackrel{\Pi_{11}}{\vdots}$}\noLine
\UIC{$a \in \mathbb{R} \to b \in \mathbb{R} \to 0 < a \to 0 < b \to \varphi_{52}$}
\AXC{$(a \in \mathbb{R} \to b \in \mathbb{R} \to 0 < a \to 0 < b \to \varphi_{52}) \to \varphi_{202}$\lbl{A1}}
\BIC{$\varphi_{202}$}
\end{ptree}
\noindent$\Pi_{19}$:
\begin{ptree}
\AXC{$\forall a\,\forall X\,(\mathrm{UB}(a,X) \leftrightarrow a \in \mathbb{R} \land \forall x \in X\,(x \leq a))$\lbl{DefUB}}
\AXC{$\varphi_{116}$\lbl{A2}}
\BIC{$\forall X\,(\mathrm{UB}(s,X) \leftrightarrow s \in \mathbb{R} \land \forall x \in X\,(x \leq s))$}
\AXC{$\varphi_{136}$\lbl{A2}}
\BIC{$\mathrm{UB}(s,X_{a}) \leftrightarrow s \in \mathbb{R} \land \forall x \in X_{a}\,(x \leq s)$}
\AXC{$\stackrel{\Pi_{18}}{\vdots}$}\noLine
\UIC{$\varphi_{202}$}
\BIC{$a \in \mathbb{R} \to b \in \mathbb{R} \to 0 < a \to 0 < b \to \varphi_{36}$}
\end{ptree}
\noindent$\Pi_{20}$:
\begin{ptree}
\AXC{$\forall x\,\forall y\,(x \in \mathbb{R} \to y \in \mathbb{R} \to x - y \in \mathbb{R})$\lbl{[F1]}}
\AXC{$\varphi_{129}$\lbl{A2}}
\BIC{$\forall y\,(s \in \mathbb{R} \to y \in \mathbb{R} \to s - y \in \mathbb{R})$}
\AXC{$\varphi_{113}$\lbl{A2}}
\BIC{$s \in \mathbb{R} \to a \in \mathbb{R} \to s - a \in \mathbb{R}$}
\AXC{$(s \in \mathbb{R} \to a \in \mathbb{R} \to s - a \in \mathbb{R}) \to \varphi_{256}$\lbl{A1}}
\BIC{$\varphi_{256}$}
\end{ptree}
\noindent$\Pi_{21}$:
\begin{ptree}
\AXC{$a \in \mathbb{R} \to b \in \mathbb{R} \to 0 < a \to 0 < b \to \varphi_{132}$\lbl{A1}}
\AXC{$\stackrel{\Pi_{19}}{\vdots}$}\noLine
\UIC{$a \in \mathbb{R} \to b \in \mathbb{R} \to 0 < a \to 0 < b \to \varphi_{36}$}
\AXC{$\stackrel{\Pi_{20}}{\vdots}$}\noLine
\UIC{$\varphi_{256}$}
\BIC{$(a \in \mathbb{R} \to b \in \mathbb{R} \to 0 < a \to 0 < b \to \varphi_{132}) \to \varphi_{181}$}
\BIC{$\varphi_{181}$}
\end{ptree}
\noindent$\Pi_{22}$:
\begin{ptree}
\AXC{$\stackrel{\Pi_{21}}{\vdots}$}\noLine
\UIC{$\varphi_{181}$}
\AXC{$\forall x\,\forall y\,(x \in \mathbb{R} \to y \in \mathbb{R} \to \neg y \leq x \to x < y)$\lbl{[A2]}}
\AXC{$\varphi_{146}$\lbl{A2}}
\BIC{$\forall y\,(s - a \in \mathbb{R} \to y \in \mathbb{R} \to \neg y \leq s - a \to s - a < y)$}
\AXC{$\varphi_{153}$\lbl{A2}}
\BIC{$s - a \in \mathbb{R} \to y \in \mathbb{R} \to \neg y \leq s - a \to s - a < y$}
\AXC{$\varphi_{274}$\lbl{A1}}
\BIC{$\varphi_{181} \to (a \in \mathbb{R} \to b \in \mathbb{R} \to 0 < a \to 0 < b \to \varphi_{131}) \to \varphi_{172} \to \varphi_{173}$}
\BIC{$(a \in \mathbb{R} \to b \in \mathbb{R} \to 0 < a \to 0 < b \to \varphi_{131}) \to \varphi_{172} \to \varphi_{173}$}
\end{ptree}
\noindent$\Pi_{23}$:
\begin{ptree}
\AXC{$\varphi_{191}$\lbl{A1}}
\AXC{$\varphi_{191} \to \varphi_{172}$\lbl{A1}}
\BIC{$\varphi_{172}$}
\AXC{$\stackrel{\Pi_{17}}{\vdots}$}\noLine
\UIC{$a \in \mathbb{R} \to b \in \mathbb{R} \to 0 < a \to 0 < b \to \varphi_{131}$}
\AXC{$\stackrel{\Pi_{22}}{\vdots}$}\noLine
\UIC{$(a \in \mathbb{R} \to b \in \mathbb{R} \to 0 < a \to 0 < b \to \varphi_{131}) \to \varphi_{172} \to \varphi_{173}$}
\BIC{$\varphi_{172} \to \varphi_{173}$}
\BIC{$\varphi_{173}$}
\end{ptree}
\noindent$\Pi_{24}$:
\begin{ptree}
\AXC{$\stackrel{\Pi_{19}}{\vdots}$}\noLine
\UIC{$a \in \mathbb{R} \to b \in \mathbb{R} \to 0 < a \to 0 < b \to \varphi_{36}$}
\AXC{$\varphi_{63}$\lbl{[B7]}}
\AXC{$\varphi_{63} \to \varphi_{50}$\lbl{A2}}
\BIC{$\varphi_{50}$}
\AXC{$\varphi_{174}$\lbl{A2}}
\BIC{$\forall y\,(s \in \mathbb{R} \to a \in \mathbb{R} \to y \in \mathbb{R} \to s - a < y \to s < y + a)$}
\AXC{$\varphi_{164}$\lbl{A2}}
\BIC{$s \in \mathbb{R} \to a \in \mathbb{R} \to y \in \mathbb{R} \to s - a < y \to s < y + a$}
\AXC{$\varphi_{277}$\lbl{A1}}
\BIC{$(a \in \mathbb{R} \to b \in \mathbb{R} \to 0 < a \to 0 < b \to \varphi_{36}) \to \varphi_{204} \to \varphi_{259}$}
\BIC{$\varphi_{204} \to \varphi_{259}$}
\end{ptree}
\noindent$\Pi_{25}$:
\begin{ptree}
\AXC{$\stackrel{\Pi_{23}}{\vdots}$}\noLine
\UIC{$\varphi_{173}$}
\AXC{$\stackrel{\Pi_{17}}{\vdots}$}\noLine
\UIC{$a \in \mathbb{R} \to b \in \mathbb{R} \to 0 < a \to 0 < b \to \varphi_{131}$}
\AXC{$\varphi_{204}$\lbl{A1}}
\AXC{$\stackrel{\Pi_{24}}{\vdots}$}\noLine
\UIC{$\varphi_{204} \to \varphi_{259}$}
\BIC{$\varphi_{259}$}
\BIC{$\varphi_{173} \to a \in \mathbb{R} \to b \in \mathbb{R} \to 0 < a \to \varphi_{165}$}
\BIC{$a \in \mathbb{R} \to b \in \mathbb{R} \to 0 < a \to \varphi_{165}$}
\end{ptree}
\noindent$\Pi_{26}$:
\begin{ptree}
\AXC{$\stackrel{\Pi_{25}}{\vdots}$}\noLine
\UIC{$a \in \mathbb{R} \to b \in \mathbb{R} \to 0 < a \to \varphi_{165}$}
\AXC{$a \in \mathbb{R} \to b \in \mathbb{R} \to 0 < a \to \varphi_{187}$\lbl{A1}}
\AXC{$\varphi_{248}$\lbl{A1}}
\BIC{$a \in \mathbb{R} \to b \in \mathbb{R} \to 0 < a \to \varphi_{170}$}
\AXC{$\varphi_{269}$\lbl{A1}}
\BIC{$(a \in \mathbb{R} \to b \in \mathbb{R} \to 0 < a \to \varphi_{165}) \to \varphi_{221}$}
\BIC{$\varphi_{221}$}
\end{ptree}
\noindent$\Pi_{27}$:
\begin{ptree}
\AXC{$x = x$\lbl{A5}}
\AXC{$x = y \to x = x \to y = x$\lbl{A6}}
\AXC{$(x = y \to x = x \to y = x) \to x = x \to x = y \to y = x$\lbl{A1}}
\BIC{$x = x \to x = y \to y = x$}
\BIC{$x = y \to y = x$}
\UIC{$\forall y\,(x = y \to y = x)$}
\UIC{$\forall x\,\forall y\,(x = y \to y = x)$}
\end{ptree}
\noindent$\Pi_{28}$:
\begin{ptree}
\AXC{$\forall m\,\forall a\,(m \in \mathbb{N} \to a \in \mathbb{R} \to \underline{S(m)} \cdot a = \underline{m} \cdot a + a)$\lbl{[B8]}}
\AXC{$\varphi_{161}$\lbl{A2}}
\BIC{$\forall a\,(m \in \mathbb{N} \to a \in \mathbb{R} \to \underline{S(m)} \cdot a = \underline{m} \cdot a + a)$}
\AXC{$\varphi_{156}$\lbl{A2}}
\BIC{$m \in \mathbb{N} \to a \in \mathbb{R} \to \underline{S(m)} \cdot a = \underline{m} \cdot a + a$}
\AXC{$\varphi_{271}$\lbl{A1}}
\BIC{$\varphi_{203} \to \varphi_{204} \to a \in \mathbb{R} \to b \in \mathbb{R} \to 0 < a \to \varphi_{192}$}
\end{ptree}
\noindent$\Pi_{29}$:
\begin{ptree}
\AXC{$a \in \mathbb{R} \to b \in \mathbb{R} \to 0 < a \to \varphi_{187}$\lbl{A1}}
\AXC{$(a \in \mathbb{R} \to b \in \mathbb{R} \to 0 < a \to \varphi_{187}) \to \varphi_{203}$\lbl{A1}}
\BIC{$\varphi_{203}$}
\AXC{$\stackrel{\Pi_{28}}{\vdots}$}\noLine
\UIC{$\varphi_{203} \to \varphi_{204} \to a \in \mathbb{R} \to b \in \mathbb{R} \to 0 < a \to \varphi_{192}$}
\BIC{$\varphi_{204} \to a \in \mathbb{R} \to b \in \mathbb{R} \to 0 < a \to \varphi_{192}$}
\end{ptree}
\noindent$\Pi_{30}$:
\begin{ptree}
\AXC{$\varphi_{204}$\lbl{A1}}
\AXC{$\stackrel{\Pi_{29}}{\vdots}$}\noLine
\UIC{$\varphi_{204} \to a \in \mathbb{R} \to b \in \mathbb{R} \to 0 < a \to \varphi_{192}$}
\BIC{$a \in \mathbb{R} \to b \in \mathbb{R} \to 0 < a \to \varphi_{192}$}
\AXC{$(a \in \mathbb{R} \to b \in \mathbb{R} \to 0 < a \to \varphi_{192}) \to \varphi_{226}$\lbl{A1}}
\BIC{$\varphi_{226}$}
\end{ptree}
\noindent$\Pi_{31}$:
\begin{ptree}
\AXC{$\stackrel{\Pi_{27}}{\vdots}$}\noLine
\UIC{$\forall x\,\forall y\,(x = y \to y = x)$}
\AXC{$\varphi_{10}$\lbl{A2}}
\BIC{$\forall y\,(\underline{S(m)} \cdot a = y \to y = \underline{S(m)} \cdot a)$}
\AXC{$\varphi_{150}$\lbl{A2}}
\BIC{$\underline{S(m)} \cdot a = \underline{m} \cdot a + a \to \underline{m} \cdot a + a = \underline{S(m)} \cdot a$}
\AXC{$\stackrel{\Pi_{30}}{\vdots}$}\noLine
\UIC{$\varphi_{226}$}
\BIC{$a \in \mathbb{R} \to b \in \mathbb{R} \to 0 < a \to \varphi_{193}$}
\AXC{$(a \in \mathbb{R} \to b \in \mathbb{R} \to 0 < a \to \varphi_{193}) \to \varphi_{258}$\lbl{A1}}
\BIC{$\varphi_{258}$}
\end{ptree}
\noindent$\Pi_{32}$:
\begin{ptree}
\AXC{$\stackrel{\Pi_{15}}{\vdots}$}\noLine
\UIC{$\varphi_{90}$}
\AXC{$z = w \to s < z + a \to s < w + a$\lbl{A6}}
\UIC{$\forall w\,(z = w \to s < z + a \to s < w + a)$}
\UIC{$\forall z\,\forall w\,(z = w \to s < z + a \to s < w + a)$}
\AXC{$\varphi_{68}$\lbl{A2}}
\BIC{$\forall w\,(y = w \to s < y + a \to s < w + a)$}
\AXC{$\varphi_{117}$\lbl{A2}}
\BIC{$y = \underline{m} \cdot a \to s < y + a \to s < \underline{m} \cdot a + a$}
\AXC{$\stackrel{\Pi_{26}}{\vdots}$}\noLine
\UIC{$\varphi_{221}$}
\BIC{$a \in \mathbb{R} \to b \in \mathbb{R} \to 0 < a \to \varphi_{184}$}
\AXC{$\stackrel{\Pi_{31}}{\vdots}$}\noLine
\UIC{$\varphi_{258}$}
\BIC{$\varphi_{90} \to a \in \mathbb{R} \to b \in \mathbb{R} \to 0 < a \to \varphi_{180}$}
\BIC{$a \in \mathbb{R} \to b \in \mathbb{R} \to 0 < a \to \varphi_{180}$}
\end{ptree}
\noindent$\Pi_{33}$:
\begin{ptree}
\AXC{$\forall x\,\forall y\,\forall z\,(x < y \to y \leq z \to x < z)$\lbl{[O3]}}
\AXC{$\varphi_{15}$\lbl{A2}}
\BIC{$\forall y\,\forall z\,(s < y \to y \leq z \to s < z)$}
\AXC{$\varphi_{93}$\lbl{A2}}
\BIC{$\forall z\,(s < \underline{S(m)} \cdot a \to \underline{S(m)} \cdot a \leq z \to s < z)$}
\AXC{$\varphi_{145}$\lbl{A2}}
\BIC{$s < \underline{S(m)} \cdot a \to \underline{S(m)} \cdot a \leq s \to s < s$}
\end{ptree}
\noindent$\Pi_{34}$:
\begin{ptree}
\AXC{$\stackrel{\Pi_{33}}{\vdots}$}\noLine
\UIC{$s < \underline{S(m)} \cdot a \to \underline{S(m)} \cdot a \leq s \to s < s$}
\AXC{$(s < \underline{S(m)} \cdot a \to \underline{S(m)} \cdot a \leq s \to s < s) \to \varphi_{264}$\lbl{A1}}
\BIC{$\varphi_{264}$}
\end{ptree}
\noindent$\Pi_{35}$:
\begin{ptree}
\AXC{$\stackrel{\Pi_{14}}{\vdots}$}\noLine
\UIC{$a \in \mathbb{R} \to b \in \mathbb{R} \to 0 < a \to \varphi_{175}$}
\AXC{$\stackrel{\Pi_{32}}{\vdots}$}\noLine
\UIC{$a \in \mathbb{R} \to b \in \mathbb{R} \to 0 < a \to \varphi_{180}$}
\AXC{$\stackrel{\Pi_{34}}{\vdots}$}\noLine
\UIC{$\varphi_{264}$}
\BIC{$(a \in \mathbb{R} \to b \in \mathbb{R} \to 0 < a \to \varphi_{175}) \to \varphi_{200}$}
\BIC{$\varphi_{200}$}
\AXC{$\varphi_{200} \to \neg s < s \to \varphi_{201}$\lbl{A1}}
\BIC{$\neg s < s \to \varphi_{201}$}
\end{ptree}
\noindent$\Pi_{36}$:
\begin{ptree}
\AXC{$\forall x\,(\neg x < x)$\lbl{[O1]}}
\AXC{$\forall x\,(\neg x < x) \to \neg s < s$\lbl{A2}}
\BIC{$\neg s < s$}
\AXC{$\stackrel{\Pi_{35}}{\vdots}$}\noLine
\UIC{$\neg s < s \to \varphi_{201}$}
\BIC{$\varphi_{201}$}
\AXC{$\varphi_{201} \to \varphi_{168} \to \neg (m \in \mathbb{N} \land y = \underline{m} \cdot a)$\lbl{A1}}
\BIC{$\varphi_{168} \to \neg (m \in \mathbb{N} \land y = \underline{m} \cdot a)$}
\UIC{$\forall m\,(\varphi_{168} \to \neg (m \in \mathbb{N} \land y = \underline{m} \cdot a))$}
\end{ptree}
\noindent$\Pi_{37}$:
\begin{ptree}
\AXC{$\forall a\,\forall y\,(y \in X_{a} \leftrightarrow y \in \mathbb{R} \land \exists n \in \mathbb{N}\,(y = \underline{n} \cdot a))$\lbl{Def$X_a$}}
\AXC{$\varphi_{163}$\lbl{A2}}
\BIC{$\forall y\,(y \in X_{a} \leftrightarrow y \in \mathbb{R} \land \exists n \in \mathbb{N}\,(y = \underline{n} \cdot a))$}
\AXC{$\varphi_{158}$\lbl{A2}}
\BIC{$y \in X_{a} \leftrightarrow y \in \mathbb{R} \land \exists n \in \mathbb{N}\,(y = \underline{n} \cdot a)$}
\end{ptree}
\noindent$\Pi_{38}$:
\begin{ptree}
\AXC{$\stackrel{\Pi_{37}}{\vdots}$}\noLine
\UIC{$y \in X_{a} \leftrightarrow y \in \mathbb{R} \land \exists n \in \mathbb{N}\,(y = \underline{n} \cdot a)$}
\AXC{$\varphi_{191}$\lbl{A1}}
\AXC{$\varphi_{191} \to a \in \mathbb{R} \to b \in \mathbb{R} \to 0 < a \to 0 < b \to \varphi_{139}$\lbl{A1}}
\BIC{$a \in \mathbb{R} \to b \in \mathbb{R} \to 0 < a \to 0 < b \to \varphi_{139}$}
\AXC{$\varphi_{250}$\lbl{A1}}
\BIC{$(y \in X_{a} \leftrightarrow y \in \mathbb{R} \land \exists n \in \mathbb{N}\,(y = \underline{n} \cdot a)) \to \varphi_{195}$}
\BIC{$\varphi_{195}$}
\end{ptree}
\noindent$\Pi_{39}$:
\begin{ptree}
\AXC{$\varphi_{13}$\lbl{A2}}
\UIC{$\forall n\,(\varphi_{13})$}
\AXC{$\varphi_{168} \to \forall m\,(\varphi_{168})$\lbl{A4}}
\AXC{$\varphi_{244}$\lbl{A3}}
\AXC{$\stackrel{\Pi_{36}}{\vdots}$}\noLine
\UIC{$\forall m\,(\varphi_{168} \to \neg (m \in \mathbb{N} \land y = \underline{m} \cdot a))$}
\AXC{$\stackrel{\Pi_{38}}{\vdots}$}\noLine
\UIC{$\varphi_{195}$}
\AXC{$\varphi_{195} \to \varphi_{279}$\lbl{A1}}
\BIC{$\varphi_{279}$}
\BIC{$\varphi_{244} \to (\varphi_{168} \to \forall m\,(\varphi_{168})) \to \forall n\,(\varphi_{13}) \to \varphi_{243}$}
\BIC{$(\varphi_{168} \to \forall m\,(\varphi_{168})) \to \forall n\,(\varphi_{13}) \to \varphi_{243}$}
\BIC{$\forall n\,(\varphi_{13}) \to \varphi_{243}$}
\BIC{$\varphi_{243}$}
\end{ptree}
\noindent$\Pi_{40}$:
\begin{ptree}
\AXC{$\varphi_{31}$\lbl{A4}}
\AXC{$\forall n\,(\varphi_{13}) \to \varphi_{35}$\lbl{A3}}
\AXC{$\stackrel{\Pi_{39}}{\vdots}$}\noLine
\UIC{$\varphi_{243}$}
\BIC{$\varphi_{31} \to a \in \mathbb{R} \to b \in \mathbb{R} \to 0 < a \to 0 < b \to \varphi_{123}$}
\BIC{$a \in \mathbb{R} \to b \in \mathbb{R} \to 0 < a \to 0 < b \to \varphi_{123}$}
\AXC{$(a \in \mathbb{R} \to b \in \mathbb{R} \to 0 < a \to 0 < b \to \varphi_{123}) \to \varphi_{169}$\lbl{A1}}
\BIC{$\varphi_{169}$}
\UIC{$\forall y\,(\varphi_{169})$}
\end{ptree}
\noindent$\Pi_{41}$:
\begin{ptree}
\AXC{$\stackrel{\Pi_{19}}{\vdots}$}\noLine
\UIC{$a \in \mathbb{R} \to b \in \mathbb{R} \to 0 < a \to 0 < b \to \varphi_{36}$}
\AXC{$\forall x\,\forall y\,(x \in \mathbb{R} \to y \in \mathbb{R} \to x - y \in \mathbb{R})$\lbl{[F1]}}
\AXC{$\varphi_{129}$\lbl{A2}}
\BIC{$\forall y\,(s \in \mathbb{R} \to y \in \mathbb{R} \to s - y \in \mathbb{R})$}
\AXC{$\varphi_{113}$\lbl{A2}}
\BIC{$s \in \mathbb{R} \to a \in \mathbb{R} \to s - a \in \mathbb{R}$}
\AXC{$\varphi_{254}$\lbl{A1}}
\BIC{$(a \in \mathbb{R} \to b \in \mathbb{R} \to 0 < a \to 0 < b \to \varphi_{36}) \to \varphi_{217}$}
\BIC{$\varphi_{217}$}
\end{ptree}
\noindent$\Pi_{42}$:
\begin{ptree}
\AXC{$a \in \mathbb{R} \to b \in \mathbb{R} \to 0 < a \to 0 < b \to \varphi_{37}$\lbl{A1}}
\AXC{$\stackrel{\Pi_{41}}{\vdots}$}\noLine
\UIC{$\varphi_{217}$}
\BIC{$a \in \mathbb{R} \to b \in \mathbb{R} \to 0 < a \to 0 < b \to \varphi_{53}$}
\end{ptree}
\noindent$\Pi_{43}$:
\begin{ptree}
\AXC{$\forall a\,\forall X\,(\mathrm{UB}(a,X) \leftrightarrow a \in \mathbb{R} \land \forall x \in X\,(x \leq a))$\lbl{DefUB}}
\AXC{$\forall a\,\forall X\,(\mathrm{UB}(a,X) \leftrightarrow a \in \mathbb{R} \land \forall x \in X\,(x \leq a)) \to \varphi_{7}$\lbl{A2}}
\BIC{$\varphi_{7}$}
\AXC{$\varphi_{162}$\lbl{A2}}
\BIC{$\mathrm{UB}(s - a,X_{a}) \leftrightarrow s - a \in \mathbb{R} \land \forall x \in X_{a}\,(x \leq s - a)$}
\end{ptree}
\noindent$\Pi_{44}$:
\begin{ptree}
\AXC{$\forall x\,\forall y\,(x \in \mathbb{R} \to y \in \mathbb{R} \to x - (x - y) = y)$\lbl{[B6]}}
\AXC{$\varphi_{118}$\lbl{A2}}
\BIC{$\forall y\,(s \in \mathbb{R} \to y \in \mathbb{R} \to s - (s - y) = y)$}
\AXC{$\varphi_{98}$\lbl{A2}}
\BIC{$s \in \mathbb{R} \to a \in \mathbb{R} \to s - (s - a) = a$}
\AXC{$\varphi_{253}$\lbl{A1}}
\BIC{$(a \in \mathbb{R} \to b \in \mathbb{R} \to 0 < a \to 0 < b \to \varphi_{36}) \to \varphi_{214}$}
\end{ptree}
\noindent$\Pi_{45}$:
\begin{ptree}
\AXC{$a \in \mathbb{R} \to b \in \mathbb{R} \to 0 < a \to 0 < b \to \varphi_{37}$\lbl{A1}}
\AXC{$\stackrel{\Pi_{19}}{\vdots}$}\noLine
\UIC{$a \in \mathbb{R} \to b \in \mathbb{R} \to 0 < a \to 0 < b \to \varphi_{36}$}
\AXC{$\stackrel{\Pi_{44}}{\vdots}$}\noLine
\UIC{$(a \in \mathbb{R} \to b \in \mathbb{R} \to 0 < a \to 0 < b \to \varphi_{36}) \to \varphi_{214}$}
\BIC{$\varphi_{214}$}
\BIC{$a \in \mathbb{R} \to b \in \mathbb{R} \to 0 < a \to 0 < b \to \varphi_{47}$}
\end{ptree}
\noindent$\Pi_{46}$:
\begin{ptree}
\AXC{$\stackrel{\Pi_{45}}{\vdots}$}\noLine
\UIC{$a \in \mathbb{R} \to b \in \mathbb{R} \to 0 < a \to 0 < b \to \varphi_{47}$}
\AXC{$(a \in \mathbb{R} \to b \in \mathbb{R} \to 0 < a \to 0 < b \to \varphi_{47}) \to \varphi_{182}$\lbl{A1}}
\BIC{$\varphi_{182}$}
\end{ptree}
\noindent$\Pi_{47}$:
\begin{ptree}
\AXC{$a \in \mathbb{R} \to b \in \mathbb{R} \to 0 < a \to 0 < b \to 0 < a$\lbl{A1}}
\AXC{$\stackrel{\Pi_{27}}{\vdots}$}\noLine
\UIC{$\forall x\,\forall y\,(x = y \to y = x)$}
\AXC{$\varphi_{3}$\lbl{A2}}
\BIC{$\forall y\,(s - (s - a) = y \to y = s - (s - a))$}
\AXC{$\varphi_{23}$\lbl{A2}}
\BIC{$s - (s - a) = a \to a = s - (s - a)$}
\AXC{$\stackrel{\Pi_{46}}{\vdots}$}\noLine
\UIC{$\varphi_{182}$}
\BIC{$a \in \mathbb{R} \to b \in \mathbb{R} \to 0 < a \to 0 < b \to \varphi_{48}$}
\AXC{$\varphi_{242}$\lbl{A1}}
\BIC{$(a \in \mathbb{R} \to b \in \mathbb{R} \to 0 < a \to 0 < b \to 0 < a) \to \varphi_{196}$}
\BIC{$\varphi_{196}$}
\end{ptree}
\noindent$\Pi_{48}$:
\begin{ptree}
\AXC{$\forall x\,\forall y\,(x \in \mathbb{R} \to y \in \mathbb{R} \to 0 < y - x \to x < y)$\lbl{[F25]}}
\AXC{$\forall x\,\forall y\,(x \in \mathbb{R} \to y \in \mathbb{R} \to 0 < y - x \to x < y) \to \varphi_{24}$\lbl{A2}}
\BIC{$\varphi_{24}$}
\end{ptree}
\noindent$\Pi_{49}$:
\begin{ptree}
\AXC{$\stackrel{\Pi_{42}}{\vdots}$}\noLine
\UIC{$a \in \mathbb{R} \to b \in \mathbb{R} \to 0 < a \to 0 < b \to \varphi_{53}$}
\AXC{$\stackrel{\Pi_{48}}{\vdots}$}\noLine
\UIC{$\varphi_{24}$}
\AXC{$\varphi_{24} \to s - a \in \mathbb{R} \to s \in \mathbb{R} \to 0 < s - (s - a) \to s - a < s$\lbl{A2}}
\BIC{$s - a \in \mathbb{R} \to s \in \mathbb{R} \to 0 < s - (s - a) \to s - a < s$}
\AXC{$\varphi_{265}$\lbl{A1}}
\BIC{$\varphi_{260}$}
\BIC{$(a \in \mathbb{R} \to b \in \mathbb{R} \to 0 < a \to 0 < b \to \varphi_{36}) \to \varphi_{219}$}
\end{ptree}
\noindent$\Pi_{50}$:
\begin{ptree}
\AXC{$\stackrel{\Pi_{19}}{\vdots}$}\noLine
\UIC{$a \in \mathbb{R} \to b \in \mathbb{R} \to 0 < a \to 0 < b \to \varphi_{36}$}
\AXC{$\stackrel{\Pi_{49}}{\vdots}$}\noLine
\UIC{$(a \in \mathbb{R} \to b \in \mathbb{R} \to 0 < a \to 0 < b \to \varphi_{36}) \to \varphi_{219}$}
\BIC{$\varphi_{219}$}
\end{ptree}
\noindent$\Pi_{51}$:
\begin{ptree}
\AXC{$z = w \to 0 < z \to 0 < w$\lbl{A6}}
\UIC{$\forall w\,(z = w \to 0 < z \to 0 < w)$}
\UIC{$\forall z\,\forall w\,(z = w \to 0 < z \to 0 < w)$}
\AXC{$\varphi_{4}$\lbl{A2}}
\BIC{$\forall w\,(a = w \to 0 < a \to 0 < w)$}
\AXC{$\varphi_{39}$\lbl{A2}}
\BIC{$a = s - (s - a) \to 0 < a \to 0 < s - (s - a)$}
\AXC{$\stackrel{\Pi_{47}}{\vdots}$}\noLine
\UIC{$\varphi_{196}$}
\BIC{$a \in \mathbb{R} \to b \in \mathbb{R} \to 0 < a \to 0 < b \to \varphi_{59}$}
\AXC{$\stackrel{\Pi_{50}}{\vdots}$}\noLine
\UIC{$\varphi_{219}$}
\BIC{$a \in \mathbb{R} \to b \in \mathbb{R} \to 0 < a \to 0 < b \to \varphi_{38}$}
\end{ptree}
\noindent$\Pi_{52}$:
\begin{ptree}
\AXC{$\forall x\,\forall y\,(x \leq y \to \neg y < x)$\lbl{[O7]}}
\AXC{$\forall x\,\forall y\,(x \leq y \to \neg y < x) \to \forall y\,(s \leq y \to \neg y < s)$\lbl{A2}}
\BIC{$\forall y\,(s \leq y \to \neg y < s)$}
\end{ptree}
\noindent$\Pi_{53}$:
\begin{ptree}
\AXC{$a \in \mathbb{R} \to b \in \mathbb{R} \to 0 < a \to 0 < b \to \varphi_{64}$\lbl{A1}}
\AXC{$(a \in \mathbb{R} \to b \in \mathbb{R} \to 0 < a \to 0 < b \to \varphi_{64}) \to \varphi_{14} \to \varphi_{183}$\lbl{A1}}
\BIC{$\varphi_{14} \to \varphi_{183}$}
\end{ptree}
\noindent$\Pi_{54}$:
\begin{ptree}
\AXC{$\varphi_{6}$\lbl{A2}}
\AXC{$\forall s\,(\varphi_{1})$\lbl{DefSup}}
\AXC{$\forall s\,(\varphi_{1}) \to \varphi_{1}$\lbl{A2}}
\BIC{$\varphi_{1}$}
\AXC{$\varphi_{1} \to \varphi_{14}$\lbl{A2}}
\BIC{$\varphi_{14}$}
\AXC{$\stackrel{\Pi_{53}}{\vdots}$}\noLine
\UIC{$\varphi_{14} \to \varphi_{183}$}
\BIC{$\varphi_{183}$}
\AXC{$\varphi_{183} \to \varphi_{6} \to \varphi_{190}$\lbl{A1}}
\BIC{$\varphi_{6} \to \varphi_{190}$}
\BIC{$\varphi_{190}$}
\end{ptree}
\noindent$\Pi_{55}$:
\begin{ptree}
\AXC{$\varphi_{176}$\lbl{A1}}
\AXC{$\stackrel{\Pi_{54}}{\vdots}$}\noLine
\UIC{$\varphi_{190}$}
\AXC{$\varphi_{190} \to \varphi_{176} \to a \in \mathbb{R} \to b \in \mathbb{R} \to 0 < a \to 0 < b \to \varphi_{111}$\lbl{A1}}
\BIC{$\varphi_{176} \to a \in \mathbb{R} \to b \in \mathbb{R} \to 0 < a \to 0 < b \to \varphi_{111}$}
\BIC{$a \in \mathbb{R} \to b \in \mathbb{R} \to 0 < a \to 0 < b \to \varphi_{111}$}
\AXC{$(a \in \mathbb{R} \to b \in \mathbb{R} \to 0 < a \to 0 < b \to \varphi_{111}) \to \varphi_{231}$\lbl{A1}}
\BIC{$\varphi_{231}$}
\end{ptree}
\noindent$\Pi_{56}$:
\begin{ptree}
\AXC{$\stackrel{\Pi_{52}}{\vdots}$}\noLine
\UIC{$\forall y\,(s \leq y \to \neg y < s)$}
\AXC{$\forall y\,(s \leq y \to \neg y < s) \to s \leq s - a \to \neg s - a < s$\lbl{A2}}
\BIC{$s \leq s - a \to \neg s - a < s$}
\AXC{$\stackrel{\Pi_{55}}{\vdots}$}\noLine
\UIC{$\varphi_{231}$}
\BIC{$(a \in \mathbb{R} \to b \in \mathbb{R} \to 0 < a \to 0 < b \to \varphi_{38}) \to \varphi_{177}$}
\end{ptree}
\noindent$\Pi_{57}$:
\begin{ptree}
\AXC{$\stackrel{\Pi_{51}}{\vdots}$}\noLine
\UIC{$a \in \mathbb{R} \to b \in \mathbb{R} \to 0 < a \to 0 < b \to \varphi_{38}$}
\AXC{$\stackrel{\Pi_{56}}{\vdots}$}\noLine
\UIC{$(a \in \mathbb{R} \to b \in \mathbb{R} \to 0 < a \to 0 < b \to \varphi_{38}) \to \varphi_{177}$}
\BIC{$\varphi_{177}$}
\end{ptree}
\noindent$\Pi_{58}$:
\begin{ptree}
\AXC{$\stackrel{\Pi_{57}}{\vdots}$}\noLine
\UIC{$\varphi_{177}$}
\AXC{$\varphi_{177} \to a \in \mathbb{R} \to b \in \mathbb{R} \to 0 < a \to 0 < b \to \varphi_{77}$\lbl{A1}}
\BIC{$a \in \mathbb{R} \to b \in \mathbb{R} \to 0 < a \to 0 < b \to \varphi_{77}$}
\AXC{$(a \in \mathbb{R} \to b \in \mathbb{R} \to 0 < a \to 0 < b \to \varphi_{77}) \to \varphi_{241}$\lbl{A1}}
\BIC{$\varphi_{241}$}
\end{ptree}
\noindent$\Pi_{59}$:
\begin{ptree}
\AXC{$\stackrel{\Pi_{42}}{\vdots}$}\noLine
\UIC{$a \in \mathbb{R} \to b \in \mathbb{R} \to 0 < a \to 0 < b \to \varphi_{53}$}
\AXC{$\stackrel{\Pi_{43}}{\vdots}$}\noLine
\UIC{$\mathrm{UB}(s - a,X_{a}) \leftrightarrow s - a \in \mathbb{R} \land \forall x \in X_{a}\,(x \leq s - a)$}
\AXC{$\stackrel{\Pi_{58}}{\vdots}$}\noLine
\UIC{$\varphi_{241}$}
\BIC{$\varphi_{224}$}
\BIC{$a \in \mathbb{R} \to b \in \mathbb{R} \to 0 < a \to 0 < b \to \varphi_{103}$}
\AXC{$\varphi_{266}$\lbl{A1}}
\BIC{$\varphi_{18} \to (\varphi_{18} \to \varphi_{27}) \to \varphi_{32} \to a \in \mathbb{R} \to b \in \mathbb{R} \to 0 < a \to 0 < b \to \varphi_{104}$}
\end{ptree}
\noindent$\Pi_{60}$:
\begin{ptree}
\AXC{$\varphi_{18} \to \varphi_{27}$\lbl{A3}}
\AXC{$\varphi_{16}$\lbl{A2}}
\AXC{$\varphi_{155}$\lbl{A1}}
\BIC{$\forall x\,(\neg \neg (x \in X_{a} \to x \leq s - a)) \to x \in X_{a} \to x \leq s - a$}
\UIC{$\varphi_{18}$}
\AXC{$\stackrel{\Pi_{59}}{\vdots}$}\noLine
\UIC{$\varphi_{18} \to (\varphi_{18} \to \varphi_{27}) \to \varphi_{32} \to a \in \mathbb{R} \to b \in \mathbb{R} \to 0 < a \to 0 < b \to \varphi_{104}$}
\BIC{$(\varphi_{18} \to \varphi_{27}) \to \varphi_{32} \to a \in \mathbb{R} \to b \in \mathbb{R} \to 0 < a \to 0 < b \to \varphi_{104}$}
\BIC{$\varphi_{32} \to a \in \mathbb{R} \to b \in \mathbb{R} \to 0 < a \to 0 < b \to \varphi_{104}$}
\end{ptree}
\noindent$\Pi_{61}$:
\begin{ptree}
\AXC{$\varphi_{213}$\lbl{A4}}
\AXC{$\forall y\,(\varphi_{169}) \to \varphi_{179}$\lbl{A3}}
\AXC{$\stackrel{\Pi_{40}}{\vdots}$}\noLine
\UIC{$\forall y\,(\varphi_{169})$}
\AXC{$\varphi_{32}$\lbl{A4}}
\AXC{$\stackrel{\Pi_{60}}{\vdots}$}\noLine
\UIC{$\varphi_{32} \to a \in \mathbb{R} \to b \in \mathbb{R} \to 0 < a \to 0 < b \to \varphi_{104}$}
\BIC{$a \in \mathbb{R} \to b \in \mathbb{R} \to 0 < a \to 0 < b \to \varphi_{104}$}
\AXC{$\varphi_{278}$\lbl{A1}}
\BIC{$\forall y\,(\varphi_{169}) \to (\forall y\,(\varphi_{169}) \to \varphi_{179}) \to \varphi_{213} \to \varphi_{252}$}
\BIC{$(\forall y\,(\varphi_{169}) \to \varphi_{179}) \to \varphi_{213} \to \varphi_{252}$}
\BIC{$\varphi_{213} \to \varphi_{252}$}
\BIC{$\varphi_{252}$}
\end{ptree}
\noindent$\Pi_{62}$:
\begin{ptree}
\AXC{$\varphi_{55}$\lbl{A4}}
\AXC{$\varphi_{42} \to \varphi_{61}$\lbl{A3}}
\AXC{$\forall y \in X_{a}\,(\neg y < b \to y = b) \to x \in X_{a} \to \neg x < b \to x = b$\lbl{A2}}
\UIC{$\varphi_{42}$}
\AXC{$\varphi_{42} \to (\varphi_{42} \to \varphi_{61}) \to \varphi_{55} \to \varphi_{43}$\lbl{A1}}
\BIC{$(\varphi_{42} \to \varphi_{61}) \to \varphi_{55} \to \varphi_{43}$}
\BIC{$\varphi_{55} \to \varphi_{43}$}
\BIC{$\varphi_{43}$}
\end{ptree}
\noindent$\Pi_{63}$:
\begin{ptree}
\AXC{$\varphi_{56}$\lbl{A4}}
\AXC{$\varphi_{44} \to \varphi_{62}$\lbl{A3}}
\AXC{$\forall x \in X_{a}\,(\neg x < b \to x = b) \to y \in X_{a} \to \neg y < b \to y = b$\lbl{A2}}
\UIC{$\varphi_{44}$}
\AXC{$\varphi_{44} \to (\varphi_{44} \to \varphi_{62}) \to \varphi_{56} \to \varphi_{45}$\lbl{A1}}
\BIC{$(\varphi_{44} \to \varphi_{62}) \to \varphi_{56} \to \varphi_{45}$}
\BIC{$\varphi_{56} \to \varphi_{45}$}
\BIC{$\varphi_{45}$}
\end{ptree}
\noindent$\Pi_{64}$:
\begin{ptree}
\AXC{$\varphi_{33}$\lbl{A2}}
\AXC{$a \in \mathbb{R} \to b \in \mathbb{R} \to 0 < a \to \varphi_{85}$\lbl{A1}}
\AXC{$\varphi_{218}$\lbl{A1}}
\BIC{$a \in \mathbb{R} \to b \in \mathbb{R} \to 0 < a \to 0 < b \to \varphi_{49}$}
\AXC{$\varphi_{236}$\lbl{A1}}
\BIC{$\varphi_{33} \to a \in \mathbb{R} \to b \in \mathbb{R} \to 0 < a \to \varphi_{74}$}
\BIC{$a \in \mathbb{R} \to b \in \mathbb{R} \to 0 < a \to \varphi_{74}$}
\AXC{$(a \in \mathbb{R} \to b \in \mathbb{R} \to 0 < a \to \varphi_{74}) \to \varphi_{232}$\lbl{A1}}
\BIC{$\varphi_{232}$}
\end{ptree}
\noindent$\Pi_{65}$:
\begin{ptree}
\AXC{$\forall n \in \mathbb{N}\,(\underline{n} \in \mathbb{R})$\lbl{[N4]}}
\AXC{$\forall n \in \mathbb{N}\,(\underline{n} \in \mathbb{R}) \to n \in \mathbb{N} \to \underline{n} \in \mathbb{R}$\lbl{A2}}
\BIC{$n \in \mathbb{N} \to \underline{n} \in \mathbb{R}$}
\AXC{$(n \in \mathbb{N} \to \underline{n} \in \mathbb{R}) \to \varphi_{230}$\lbl{A1}}
\BIC{$\varphi_{230}$}
\end{ptree}
\noindent$\Pi_{66}$:
\begin{ptree}
\AXC{$\forall x\,\forall y\,(x \in \mathbb{R} \to y \in \mathbb{R} \to x \cdot y \in \mathbb{R})$\lbl{(R1)}}
\AXC{$\varphi_{147}$\lbl{A2}}
\BIC{$\forall y\,(\underline{n} \in \mathbb{R} \to y \in \mathbb{R} \to \underline{n} \cdot y \in \mathbb{R})$}
\AXC{$\varphi_{151}$\lbl{A2}}
\BIC{$\underline{n} \in \mathbb{R} \to a \in \mathbb{R} \to \underline{n} \cdot a \in \mathbb{R}$}
\AXC{$\varphi_{262}$\lbl{A1}}
\BIC{$(a \in \mathbb{R} \to b \in \mathbb{R} \to 0 < a \to 0 < b \to \varphi_{138}) \to \varphi_{233}$}
\end{ptree}
\noindent$\Pi_{67}$:
\begin{ptree}
\AXC{$a \in \mathbb{R} \to b \in \mathbb{R} \to 0 < a \to 0 < b \to \varphi_{148}$\lbl{A1}}
\AXC{$\varphi_{235}$\lbl{A1}}
\BIC{$a \in \mathbb{R} \to b \in \mathbb{R} \to 0 < a \to 0 < b \to \varphi_{124}$}
\AXC{$\stackrel{\Pi_{65}}{\vdots}$}\noLine
\UIC{$\varphi_{230}$}
\BIC{$a \in \mathbb{R} \to b \in \mathbb{R} \to 0 < a \to 0 < b \to \varphi_{138}$}
\AXC{$\stackrel{\Pi_{66}}{\vdots}$}\noLine
\UIC{$(a \in \mathbb{R} \to b \in \mathbb{R} \to 0 < a \to 0 < b \to \varphi_{138}) \to \varphi_{233}$}
\BIC{$\varphi_{233}$}
\end{ptree}
\noindent$\Pi_{68}$:
\begin{ptree}
\AXC{$\forall x\,\forall y\,(x \in \mathbb{R} \to y \in \mathbb{R} \to \neg x < y \to y \leq x)$\lbl{[O8]}}
\AXC{$\varphi_{119}$\lbl{A2}}
\BIC{$\forall y\,(b \in \mathbb{R} \to y \in \mathbb{R} \to \neg b < y \to y \leq b)$}
\AXC{$\varphi_{154}$\lbl{A2}}
\BIC{$b \in \mathbb{R} \to \underline{n} \cdot a \in \mathbb{R} \to \neg b < \underline{n} \cdot a \to \underline{n} \cdot a \leq b$}
\AXC{$\varphi_{275}$\lbl{A1}}
\BIC{$\varphi_{268}$}
\end{ptree}
\noindent$\Pi_{69}$:
\begin{ptree}
\AXC{$a \in \mathbb{R} \to b \in \mathbb{R} \to 0 < a \to 0 < b \to \varphi_{126}$\lbl{A1}}
\AXC{$\stackrel{\Pi_{67}}{\vdots}$}\noLine
\UIC{$\varphi_{233}$}
\BIC{$a \in \mathbb{R} \to b \in \mathbb{R} \to 0 < a \to 0 < b \to \varphi_{144}$}
\AXC{$a \in \mathbb{R} \to b \in \mathbb{R} \to 0 < a \to 0 < b \to \varphi_{127}$\lbl{A1}}
\AXC{$\stackrel{\Pi_{68}}{\vdots}$}\noLine
\UIC{$\varphi_{268}$}
\BIC{$(a \in \mathbb{R} \to b \in \mathbb{R} \to 0 < a \to 0 < b \to \varphi_{144}) \to \varphi_{234}$}
\BIC{$\varphi_{234}$}
\end{ptree}
\noindent$\Pi_{70}$:
\begin{ptree}
\AXC{$\stackrel{\Pi_{27}}{\vdots}$}\noLine
\UIC{$\forall x\,\forall y\,(x = y \to y = x)$}
\AXC{$\forall x\,\forall y\,(x = y \to y = x) \to \forall y\,(v_{374} = y \to y = v_{374})$\lbl{A2}}
\BIC{$\forall y\,(v_{374} = y \to y = v_{374})$}
\AXC{$\varphi_{2}$\lbl{A2}}
\BIC{$v_{374} = v_{375} \to v_{375} = v_{374}$}
\UIC{$\forall v_{375}\,(v_{374} = v_{375} \to v_{375} = v_{374})$}
\UIC{$\forall v_{374}\,\forall v_{375}\,(v_{374} = v_{375} \to v_{375} = v_{374})$}
\AXC{$\varphi_{8}$\lbl{A2}}
\BIC{$\forall v_{375}\,(y = v_{375} \to v_{375} = y)$}
\AXC{$\varphi_{5}$\lbl{A2}}
\BIC{$y = \underline{n} \cdot a \to \underline{n} \cdot a = y$}
\end{ptree}
\noindent$\Pi_{71}$:
\begin{ptree}
\AXC{$a \in \mathbb{R} \to b \in \mathbb{R} \to 0 < a \to 0 < b \to \varphi_{148}$\lbl{A1}}
\AXC{$\varphi_{237}$\lbl{A1}}
\BIC{$a \in \mathbb{R} \to b \in \mathbb{R} \to 0 < a \to 0 < b \to \varphi_{134}$}
\AXC{$(a \in \mathbb{R} \to b \in \mathbb{R} \to 0 < a \to 0 < b \to \varphi_{134}) \to \varphi_{207}$\lbl{A1}}
\BIC{$\varphi_{207}$}
\end{ptree}
\noindent$\Pi_{72}$:
\begin{ptree}
\AXC{$\stackrel{\Pi_{70}}{\vdots}$}\noLine
\UIC{$y = \underline{n} \cdot a \to \underline{n} \cdot a = y$}
\AXC{$\stackrel{\Pi_{71}}{\vdots}$}\noLine
\UIC{$\varphi_{207}$}
\BIC{$a \in \mathbb{R} \to b \in \mathbb{R} \to 0 < a \to 0 < b \to \varphi_{135}$}
\AXC{$(a \in \mathbb{R} \to b \in \mathbb{R} \to 0 < a \to 0 < b \to \varphi_{135}) \to \varphi_{247}$\lbl{A1}}
\BIC{$\varphi_{247}$}
\end{ptree}
\noindent$\Pi_{73}$:
\begin{ptree}
\AXC{$a \in \mathbb{R} \to b \in \mathbb{R} \to 0 < a \to 0 < b \to \varphi_{148}$\lbl{A1}}
\AXC{$\varphi_{235}$\lbl{A1}}
\BIC{$a \in \mathbb{R} \to b \in \mathbb{R} \to 0 < a \to 0 < b \to \varphi_{124}$}
\AXC{$\stackrel{\Pi_{64}}{\vdots}$}\noLine
\UIC{$\varphi_{232}$}
\BIC{$a \in \mathbb{R} \to b \in \mathbb{R} \to 0 < a \to 0 < b \to \varphi_{140}$}
\AXC{$\stackrel{\Pi_{69}}{\vdots}$}\noLine
\UIC{$\varphi_{234}$}
\BIC{$a \in \mathbb{R} \to b \in \mathbb{R} \to 0 < a \to 0 < b \to \varphi_{133}$}
\AXC{$\stackrel{\Pi_{72}}{\vdots}$}\noLine
\UIC{$\varphi_{247}$}
\BIC{$(\underline{n} \cdot a = y \to \underline{n} \cdot a < b \to y < b) \to \varphi_{205}$}
\end{ptree}
\noindent$\Pi_{74}$:
\begin{ptree}
\AXC{$z = w \to z < b \to w < b$\lbl{A6}}
\UIC{$\forall w\,(z = w \to z < b \to w < b)$}
\UIC{$\forall z\,\forall w\,(z = w \to z < b \to w < b)$}
\AXC{$\varphi_{75}$\lbl{A2}}
\BIC{$\forall w\,(\underline{n} \cdot a = w \to \underline{n} \cdot a < b \to w < b)$}
\AXC{$\varphi_{125}$\lbl{A2}}
\BIC{$\underline{n} \cdot a = y \to \underline{n} \cdot a < b \to y < b$}
\AXC{$\stackrel{\Pi_{73}}{\vdots}$}\noLine
\UIC{$(\underline{n} \cdot a = y \to \underline{n} \cdot a < b \to y < b) \to \varphi_{205}$}
\BIC{$\varphi_{205}$}
\end{ptree}
\noindent$\Pi_{75}$:
\begin{ptree}
\AXC{$z = w \to z = b \to w = b$\lbl{A6}}
\UIC{$\forall w\,(z = w \to z = b \to w = b)$}
\UIC{$\forall z\,\forall w\,(z = w \to z = b \to w = b)$}
\AXC{$\varphi_{25}$\lbl{A2}}
\BIC{$\forall w\,(\underline{n} \cdot a = w \to \underline{n} \cdot a = b \to w = b)$}
\AXC{$\varphi_{88}$\lbl{A2}}
\BIC{$\underline{n} \cdot a = y \to \underline{n} \cdot a = b \to y = b$}
\AXC{$\stackrel{\Pi_{74}}{\vdots}$}\noLine
\UIC{$\varphi_{205}$}
\BIC{$a \in \mathbb{R} \to b \in \mathbb{R} \to 0 < a \to 0 < b \to \varphi_{105}$}
\AXC{$\varphi_{227}$\lbl{A1}}
\BIC{$\varphi_{80} \land y \in X_{a} \land \neg y \leq b \to \neg (n \in \mathbb{N} \land y = \underline{n} \cdot a)$}
\UIC{$\forall n\,(\varphi_{80} \land y \in X_{a} \land \neg y \leq b \to \neg (n \in \mathbb{N} \land y = \underline{n} \cdot a))$}
\end{ptree}
\noindent$\Pi_{76}$:
\begin{ptree}
\AXC{$\forall a\,\forall y\,(y \in X_{a} \leftrightarrow y \in \mathbb{R} \land \exists n \in \mathbb{N}\,(y = \underline{n} \cdot a))$\lbl{Def$X_a$}}
\AXC{$\varphi_{163}$\lbl{A2}}
\BIC{$\forall y\,(y \in X_{a} \leftrightarrow y \in \mathbb{R} \land \exists n \in \mathbb{N}\,(y = \underline{n} \cdot a))$}
\AXC{$\varphi_{158}$\lbl{A2}}
\BIC{$y \in X_{a} \leftrightarrow y \in \mathbb{R} \land \exists n \in \mathbb{N}\,(y = \underline{n} \cdot a)$}
\end{ptree}
\noindent$\Pi_{77}$:
\begin{ptree}
\AXC{$\stackrel{\Pi_{76}}{\vdots}$}\noLine
\UIC{$y \in X_{a} \leftrightarrow y \in \mathbb{R} \land \exists n \in \mathbb{N}\,(y = \underline{n} \cdot a)$}
\AXC{$a \in \mathbb{R} \to b \in \mathbb{R} \to \varphi_{106}$\lbl{A1}}
\AXC{$(a \in \mathbb{R} \to b \in \mathbb{R} \to \varphi_{106}) \to \varphi_{208}$\lbl{A1}}
\BIC{$\varphi_{208}$}
\BIC{$a \in \mathbb{R} \to b \in \mathbb{R} \to 0 < a \to 0 < b \to \varphi_{95}$}
\end{ptree}
\noindent$\Pi_{78}$:
\begin{ptree}
\AXC{$\stackrel{\Pi_{77}}{\vdots}$}\noLine
\UIC{$a \in \mathbb{R} \to b \in \mathbb{R} \to 0 < a \to 0 < b \to \varphi_{95}$}
\AXC{$(a \in \mathbb{R} \to b \in \mathbb{R} \to 0 < a \to 0 < b \to \varphi_{95}) \to \varphi_{276}$\lbl{A1}}
\BIC{$\varphi_{276}$}
\end{ptree}
\noindent$\Pi_{79}$:
\begin{ptree}
\AXC{$\varphi_{229}$\lbl{A3}}
\AXC{$\stackrel{\Pi_{75}}{\vdots}$}\noLine
\UIC{$\forall n\,(\varphi_{80} \land y \in X_{a} \land \neg y \leq b \to \neg (n \in \mathbb{N} \land y = \underline{n} \cdot a))$}
\AXC{$\stackrel{\Pi_{78}}{\vdots}$}\noLine
\UIC{$\varphi_{276}$}
\BIC{$\varphi_{229} \to \varphi_{240}$}
\BIC{$\varphi_{240}$}
\end{ptree}
\noindent$\Pi_{80}$:
\begin{ptree}
\AXC{$\varphi_{80} \land y \in X_{a} \land \neg y \leq b \to \forall n\,(\varphi_{80} \land y \in X_{a} \land \neg y \leq b)$\lbl{A4}}
\AXC{$\stackrel{\Pi_{79}}{\vdots}$}\noLine
\UIC{$\varphi_{240}$}
\BIC{$a \in \mathbb{R} \to b \in \mathbb{R} \to \varphi_{70}$}
\AXC{$(a \in \mathbb{R} \to b \in \mathbb{R} \to \varphi_{70}) \to \varphi_{80} \to y \in X_{a} \to y \leq b$\lbl{A1}}
\BIC{$\varphi_{80} \to y \in X_{a} \to y \leq b$}
\UIC{$\forall y\,(\varphi_{80} \to y \in X_{a} \to y \leq b)$}
\AXC{$\varphi_{270}$\lbl{A1}}
\BIC{$\varphi_{211} \to (\varphi_{80} \to \forall y\,(\varphi_{80})) \to a \in \mathbb{R} \to b \in \mathbb{R} \to \varphi_{81}$}
\end{ptree}
\noindent$\Pi_{81}$:
\begin{ptree}
\AXC{$\varphi_{80} \to \forall y\,(\varphi_{80})$\lbl{A4}}
\AXC{$\varphi_{211}$\lbl{A3}}
\AXC{$\stackrel{\Pi_{80}}{\vdots}$}\noLine
\UIC{$\varphi_{211} \to (\varphi_{80} \to \forall y\,(\varphi_{80})) \to a \in \mathbb{R} \to b \in \mathbb{R} \to \varphi_{81}$}
\BIC{$(\varphi_{80} \to \forall y\,(\varphi_{80})) \to a \in \mathbb{R} \to b \in \mathbb{R} \to \varphi_{81}$}
\BIC{$a \in \mathbb{R} \to b \in \mathbb{R} \to \varphi_{81}$}
\end{ptree}
\noindent$\Pi_{82}$:
\begin{ptree}
\AXC{$\forall a\,\forall X\,(\mathrm{UB}(a,X) \leftrightarrow a \in \mathbb{R} \land \forall x \in X\,(x \leq a))$\lbl{DefUB}}
\AXC{$\varphi_{120}$\lbl{A2}}
\BIC{$\forall X\,(\mathrm{UB}(b,X) \leftrightarrow b \in \mathbb{R} \land \forall x \in X\,(x \leq b))$}
\AXC{$\varphi_{137}$\lbl{A2}}
\BIC{$\mathrm{UB}(b,X_{a}) \leftrightarrow b \in \mathbb{R} \land \forall x \in X_{a}\,(x \leq b)$}
\AXC{$\varphi_{257}$\lbl{A1}}
\BIC{$(a \in \mathbb{R} \to b \in \mathbb{R} \to \varphi_{81}) \to \varphi_{45} \to \varphi_{43} \to a \in \mathbb{R} \to b \in \mathbb{R} \to \varphi_{41}$}
\end{ptree}
\noindent$\Pi_{83}$:
\begin{ptree}
\AXC{$\stackrel{\Pi_{62}}{\vdots}$}\noLine
\UIC{$\varphi_{43}$}
\AXC{$\stackrel{\Pi_{63}}{\vdots}$}\noLine
\UIC{$\varphi_{45}$}
\AXC{$\stackrel{\Pi_{81}}{\vdots}$}\noLine
\UIC{$a \in \mathbb{R} \to b \in \mathbb{R} \to \varphi_{81}$}
\AXC{$\stackrel{\Pi_{82}}{\vdots}$}\noLine
\UIC{$(a \in \mathbb{R} \to b \in \mathbb{R} \to \varphi_{81}) \to \varphi_{45} \to \varphi_{43} \to a \in \mathbb{R} \to b \in \mathbb{R} \to \varphi_{41}$}
\BIC{$\varphi_{45} \to \varphi_{43} \to a \in \mathbb{R} \to b \in \mathbb{R} \to \varphi_{41}$}
\BIC{$\varphi_{43} \to a \in \mathbb{R} \to b \in \mathbb{R} \to \varphi_{41}$}
\BIC{$a \in \mathbb{R} \to b \in \mathbb{R} \to \varphi_{41}$}
\end{ptree}
\noindent$\Pi_{84}$:
\begin{ptree}
\AXC{$\forall b\,(\neg \mathrm{UB}(b,X_{a})) \to \neg \mathrm{UB}(b,X_{a})$\lbl{A2}}
\AXC{$\stackrel{\Pi_{83}}{\vdots}$}\noLine
\UIC{$a \in \mathbb{R} \to b \in \mathbb{R} \to \varphi_{41}$}
\AXC{$(a \in \mathbb{R} \to b \in \mathbb{R} \to \varphi_{41}) \to \varphi_{167}$\lbl{A1}}
\BIC{$\varphi_{167}$}
\BIC{$a \in \mathbb{R} \to b \in \mathbb{R} \to 0 < a \to \varphi_{22}$}
\end{ptree}
\noindent$\Pi_{85}$:
\begin{ptree}
\AXC{$x = x$\lbl{A5}}
\UIC{$\forall x\,(x = x)$}
\AXC{$\forall x\,(x = x) \to \underline{1_{\omega}} \cdot a = \underline{1_{\omega}} \cdot a$\lbl{A2}}
\BIC{$\underline{1_{\omega}} \cdot a = \underline{1_{\omega}} \cdot a$}
\AXC{$1_{\omega} \in \mathbb{N}$\lbl{[N1]}}
\AXC{$1_{\omega} \in \mathbb{N} \to \varphi_{197}$\lbl{A1}}
\BIC{$\varphi_{197}$}
\BIC{$a \in \mathbb{R} \to b \in \mathbb{R} \to 0 < a \to 0 < b \to \varphi_{96}$}
\end{ptree}
\noindent$\Pi_{86}$:
\begin{ptree}
\AXC{$\forall x\,\forall y\,(x \in \mathbb{R} \to y \in \mathbb{R} \to x \cdot y \in \mathbb{R})$\lbl{(R1)}}
\AXC{$\varphi_{152}$\lbl{A2}}
\BIC{$\forall y\,(\underline{1_{\omega}} \in \mathbb{R} \to y \in \mathbb{R} \to \underline{1_{\omega}} \cdot y \in \mathbb{R})$}
\AXC{$\varphi_{159}$\lbl{A2}}
\BIC{$\underline{1_{\omega}} \in \mathbb{R} \to a \in \mathbb{R} \to \underline{1_{\omega}} \cdot a \in \mathbb{R}$}
\end{ptree}
\noindent$\Pi_{87}$:
\begin{ptree}
\AXC{$\stackrel{\Pi_{86}}{\vdots}$}\noLine
\UIC{$\underline{1_{\omega}} \in \mathbb{R} \to a \in \mathbb{R} \to \underline{1_{\omega}} \cdot a \in \mathbb{R}$}
\AXC{$(\underline{1_{\omega}} \in \mathbb{R} \to a \in \mathbb{R} \to \underline{1_{\omega}} \cdot a \in \mathbb{R}) \to \varphi_{239}$\lbl{A1}}
\BIC{$\varphi_{239}$}
\end{ptree}
\noindent$\Pi_{88}$:
\begin{ptree}
\AXC{$a \in \mathbb{R} \to \varphi_{82}$\lbl{A1}}
\AXC{$1_{\omega} \in \mathbb{N}$\lbl{[N1]}}
\AXC{$\forall n \in \mathbb{N}\,(\underline{n} \in \mathbb{R})$\lbl{[N4]}}
\AXC{$\varphi_{40}$\lbl{A2}}
\BIC{$1_{\omega} \in \mathbb{N} \to \underline{1_{\omega}} \in \mathbb{R}$}
\AXC{$\varphi_{188}$\lbl{A1}}
\BIC{$1_{\omega} \in \mathbb{N} \to a \in \mathbb{R} \to b \in \mathbb{R} \to \varphi_{65}$}
\BIC{$a \in \mathbb{R} \to b \in \mathbb{R} \to \varphi_{65}$}
\AXC{$\stackrel{\Pi_{87}}{\vdots}$}\noLine
\UIC{$\varphi_{239}$}
\BIC{$(a \in \mathbb{R} \to \varphi_{82}) \to a \in \mathbb{R} \to b \in \mathbb{R} \to \varphi_{86}$}
\BIC{$a \in \mathbb{R} \to b \in \mathbb{R} \to \varphi_{86}$}
\end{ptree}
\noindent$\Pi_{89}$:
\begin{ptree}
\AXC{$\forall a\,\forall y\,(y \in X_{a} \leftrightarrow y \in \mathbb{R} \land \exists n \in \mathbb{N}\,(y = \underline{n} \cdot a))$\lbl{Def$X_a$}}
\AXC{$\varphi_{163}$\lbl{A2}}
\BIC{$\forall y\,(y \in X_{a} \leftrightarrow y \in \mathbb{R} \land \exists n \in \mathbb{N}\,(y = \underline{n} \cdot a))$}
\AXC{$\forall y\,(y \in X_{a} \leftrightarrow y \in \mathbb{R} \land \exists n \in \mathbb{N}\,(y = \underline{n} \cdot a)) \to \varphi_{142}$\lbl{A2}}
\BIC{$\varphi_{142}$}
\end{ptree}
\noindent$\Pi_{90}$:
\begin{ptree}
\AXC{$\stackrel{\Pi_{88}}{\vdots}$}\noLine
\UIC{$a \in \mathbb{R} \to b \in \mathbb{R} \to \varphi_{86}$}
\AXC{$\stackrel{\Pi_{89}}{\vdots}$}\noLine
\UIC{$\varphi_{142}$}
\AXC{$\varphi_{142} \to (a \in \mathbb{R} \to b \in \mathbb{R} \to \varphi_{86}) \to \varphi_{220}$\lbl{A1}}
\BIC{$(a \in \mathbb{R} \to b \in \mathbb{R} \to \varphi_{86}) \to \varphi_{220}$}
\BIC{$\varphi_{220}$}
\end{ptree}
\noindent$\Pi_{91}$:
\begin{ptree}
\AXC{$\forall X\,\forall x \in X\,(X \neq \emptyset)$\lbl{[B2]}}
\AXC{$\forall X\,\forall x \in X\,(X \neq \emptyset) \to \forall x \in X_{a}\,(X_{a} \neq \emptyset)$\lbl{A2}}
\BIC{$\forall x \in X_{a}\,(X_{a} \neq \emptyset)$}
\end{ptree}
\noindent$\Pi_{92}$:
\begin{ptree}
\AXC{$\stackrel{\Pi_{91}}{\vdots}$}\noLine
\UIC{$\forall x \in X_{a}\,(X_{a} \neq \emptyset)$}
\AXC{$\forall x \in X_{a}\,(X_{a} \neq \emptyset) \to \underline{1_{\omega}} \cdot a \in X_{a} \to X_{a} \neq \emptyset$\lbl{A2}}
\BIC{$\underline{1_{\omega}} \cdot a \in X_{a} \to X_{a} \neq \emptyset$}
\AXC{$\varphi_{225}$\lbl{A1}}
\BIC{$(a \in \mathbb{R} \to b \in \mathbb{R} \to \varphi_{97}) \to a \in \mathbb{R} \to \varphi_{107}$}
\end{ptree}
\noindent$\Pi_{93}$:
\begin{ptree}
\AXC{$\varphi_{130}$\lbl{A2}}
\AXC{$\stackrel{\Pi_{85}}{\vdots}$}\noLine
\UIC{$a \in \mathbb{R} \to b \in \mathbb{R} \to 0 < a \to 0 < b \to \varphi_{96}$}
\AXC{$\varphi_{245}$\lbl{A1}}
\BIC{$\varphi_{130} \to a \in \mathbb{R} \to b \in \mathbb{R} \to 0 < a \to 0 < b \to \varphi_{87}$}
\BIC{$a \in \mathbb{R} \to b \in \mathbb{R} \to 0 < a \to 0 < b \to \varphi_{87}$}
\AXC{$\stackrel{\Pi_{90}}{\vdots}$}\noLine
\UIC{$\varphi_{220}$}
\BIC{$a \in \mathbb{R} \to b \in \mathbb{R} \to \varphi_{97}$}
\AXC{$\stackrel{\Pi_{92}}{\vdots}$}\noLine
\UIC{$(a \in \mathbb{R} \to b \in \mathbb{R} \to \varphi_{97}) \to a \in \mathbb{R} \to \varphi_{107}$}
\BIC{$a \in \mathbb{R} \to \varphi_{107}$}
\end{ptree}
\noindent$\Pi_{94}$:
\begin{ptree}
\AXC{$\forall a\,\forall y\,(y \in X_{a} \leftrightarrow y \in \mathbb{R} \land \exists n \in \mathbb{N}\,(y = \underline{n} \cdot a))$\lbl{Def$X_a$}}
\AXC{$\varphi_{163}$\lbl{A2}}
\BIC{$\forall y\,(y \in X_{a} \leftrightarrow y \in \mathbb{R} \land \exists n \in \mathbb{N}\,(y = \underline{n} \cdot a))$}
\AXC{$\varphi_{160}$\lbl{A2}}
\BIC{$u \in X_{a} \leftrightarrow u \in \mathbb{R} \land \exists n \in \mathbb{N}\,(u = \underline{n} \cdot a)$}
\AXC{$\varphi_{206}$\lbl{A1}}
\BIC{$a \in \mathbb{R} \to b \in \mathbb{R} \to \varphi_{92}$}
\end{ptree}
\noindent$\Pi_{95}$:
\begin{ptree}
\AXC{$\varphi_{80} \to \forall u\,(\varphi_{80})$\lbl{A4}}
\AXC{$\varphi_{215}$\lbl{A3}}
\AXC{$\stackrel{\Pi_{94}}{\vdots}$}\noLine
\UIC{$a \in \mathbb{R} \to b \in \mathbb{R} \to \varphi_{92}$}
\AXC{$(a \in \mathbb{R} \to b \in \mathbb{R} \to \varphi_{92}) \to \varphi_{80} \to u \in X_{a} \to u \in \mathbb{R}$\lbl{A1}}
\BIC{$\varphi_{80} \to u \in X_{a} \to u \in \mathbb{R}$}
\UIC{$\forall u\,(\varphi_{80} \to u \in X_{a} \to u \in \mathbb{R})$}
\AXC{$\varphi_{273}$\lbl{A1}}
\BIC{$\varphi_{215} \to (\varphi_{80} \to \forall u\,(\varphi_{80})) \to a \in \mathbb{R} \to b \in \mathbb{R} \to \varphi_{102}$}
\BIC{$(\varphi_{80} \to \forall u\,(\varphi_{80})) \to a \in \mathbb{R} \to b \in \mathbb{R} \to \varphi_{102}$}
\BIC{$a \in \mathbb{R} \to b \in \mathbb{R} \to \varphi_{102}$}
\end{ptree}
\noindent$\Pi_{96}$:
\begin{ptree}
\AXC{$\forall x\,\forall y\,(x \subseteq y \leftrightarrow \forall u \in x\,(u \in y))$\lbl{Def$\subseteq$}}
\AXC{$\varphi_{114}$\lbl{A2}}
\BIC{$\forall y\,(X_{a} \subseteq y \leftrightarrow \forall u \in X_{a}\,(u \in y))$}
\AXC{$\varphi_{143}$\lbl{A2}}
\BIC{$X_{a} \subseteq \mathbb{R} \leftrightarrow \forall u \in X_{a}\,(u \in \mathbb{R})$}
\AXC{$\stackrel{\Pi_{95}}{\vdots}$}\noLine
\UIC{$a \in \mathbb{R} \to b \in \mathbb{R} \to \varphi_{102}$}
\AXC{$\varphi_{228}$\lbl{A1}}
\BIC{$(X_{a} \subseteq \mathbb{R} \leftrightarrow \forall u \in X_{a}\,(u \in \mathbb{R})) \to a \in \mathbb{R} \to \varphi_{108}$}
\BIC{$a \in \mathbb{R} \to \varphi_{108}$}
\end{ptree}
\noindent$\Pi_{97}$:
\begin{ptree}
\AXC{$\stackrel{\Pi_{96}}{\vdots}$}\noLine
\UIC{$a \in \mathbb{R} \to \varphi_{108}$}
\AXC{$\varphi_{9}$\lbl{(R16)}}
\AXC{$\varphi_{9} \to X_{a} \subseteq \mathbb{R} \to \varphi_{20}$\lbl{A2}}
\BIC{$X_{a} \subseteq \mathbb{R} \to \varphi_{20}$}
\AXC{$(X_{a} \subseteq \mathbb{R} \to \varphi_{20}) \to (a \in \mathbb{R} \to \varphi_{108}) \to (a \in \mathbb{R} \to \varphi_{107}) \to \varphi_{209}$\lbl{A1}}
\BIC{$(a \in \mathbb{R} \to \varphi_{108}) \to (a \in \mathbb{R} \to \varphi_{107}) \to \varphi_{209}$}
\BIC{$(a \in \mathbb{R} \to \varphi_{107}) \to \varphi_{209}$}
\end{ptree}
\noindent$\Pi_{98}$:
\begin{ptree}
\AXC{$\stackrel{\Pi_{84}}{\vdots}$}\noLine
\UIC{$a \in \mathbb{R} \to b \in \mathbb{R} \to 0 < a \to \varphi_{22}$}
\AXC{$\stackrel{\Pi_{93}}{\vdots}$}\noLine
\UIC{$a \in \mathbb{R} \to \varphi_{107}$}
\AXC{$\stackrel{\Pi_{97}}{\vdots}$}\noLine
\UIC{$(a \in \mathbb{R} \to \varphi_{107}) \to \varphi_{209}$}
\BIC{$\varphi_{209}$}
\BIC{$a \in \mathbb{R} \to b \in \mathbb{R} \to 0 < a \to \varphi_{29}$}
\AXC{$(a \in \mathbb{R} \to b \in \mathbb{R} \to 0 < a \to \varphi_{29}) \to \varphi_{272}$\lbl{A1}}
\BIC{$\varphi_{272}$}
\end{ptree}
\noindent$\Pi_{99}$:
\begin{ptree}
\AXC{$\varphi_{216}$\lbl{A3}}
\AXC{$\varphi_{30}$\lbl{A4}}
\AXC{$\forall x\,(\varphi_{12}) \to \varphi_{34}$\lbl{A3}}
\AXC{$\varphi_{12}$\lbl{A2}}
\UIC{$\forall x\,(\varphi_{12})$}
\AXC{$\stackrel{\Pi_{61}}{\vdots}$}\noLine
\UIC{$\varphi_{252}$}
\BIC{$(\forall x\,(\varphi_{12}) \to \varphi_{34}) \to \varphi_{30} \to a \in \mathbb{R} \to b \in \mathbb{R} \to 0 < a \to \varphi_{54}$}
\BIC{$\varphi_{30} \to a \in \mathbb{R} \to b \in \mathbb{R} \to 0 < a \to \varphi_{54}$}
\BIC{$a \in \mathbb{R} \to b \in \mathbb{R} \to 0 < a \to \varphi_{54}$}
\AXC{$\varphi_{212}$\lbl{A1}}
\BIC{$\varphi_{80} \land \neg 0 < 0 \to \neg \mathrm{IsSup}(s,X_{a})$}
\UIC{$\forall s\,(\varphi_{80} \land \neg 0 < 0 \to \neg \mathrm{IsSup}(s,X_{a}))$}
\AXC{$\stackrel{\Pi_{98}}{\vdots}$}\noLine
\UIC{$\varphi_{272}$}
\BIC{$\varphi_{216} \to (\varphi_{80} \land \neg 0 < 0 \to \forall s\,(\varphi_{80} \land \neg 0 < 0)) \to a \in \mathbb{R} \to \varphi_{66}$}
\BIC{$(\varphi_{80} \land \neg 0 < 0 \to \forall s\,(\varphi_{80} \land \neg 0 < 0)) \to a \in \mathbb{R} \to \varphi_{66}$}
\end{ptree}
\noindent$\Pi_{100}$:
\begin{ptree}
\AXC{$\varphi_{80} \land \neg 0 < 0 \to \forall s\,(\varphi_{80} \land \neg 0 < 0)$\lbl{A4}}
\AXC{$\stackrel{\Pi_{99}}{\vdots}$}\noLine
\UIC{$(\varphi_{80} \land \neg 0 < 0 \to \forall s\,(\varphi_{80} \land \neg 0 < 0)) \to a \in \mathbb{R} \to \varphi_{66}$}
\BIC{$a \in \mathbb{R} \to \varphi_{66}$}
\end{ptree}
\noindent 主木:
\begin{ptree}
\AXC{$\stackrel{\Pi_{1}}{\vdots}$}\noLine
\UIC{$a \in \mathbb{R} \to \varphi_{69}$}
\AXC{$\stackrel{\Pi_{100}}{\vdots}$}\noLine
\UIC{$a \in \mathbb{R} \to \varphi_{66}$}
\AXC{$(a \in \mathbb{R} \to \varphi_{66}) \to (a \in \mathbb{R} \to \varphi_{69}) \to \varphi_{76}$\lbl{A1}}
\BIC{$(a \in \mathbb{R} \to \varphi_{69}) \to \varphi_{76}$}
\BIC{$\varphi_{76}$}
\UIC{$\forall b\,(\varphi_{76})$}
\UIC{$\forall a\,\forall b\,(\varphi_{76})$}
\end{ptree}
\begin{abbrevlist}
\item $\varphi_{1} :\equiv \forall X\,(\mathrm{IsSup}(s,X) \leftrightarrow \mathrm{UB}(s,X) \land \forall w\,(\mathrm{UB}(w,X) \to s \leq w))$
\item $\varphi_{2} :\equiv \forall y\,(v_{374} = y \to y = v_{374}) \to v_{374} = v_{375} \to v_{375} = v_{374}$
\item $\varphi_{3} :\equiv \forall x\,\forall y\,(x = y \to y = x) \to \forall y\,(s - (s - a) = y \to y = s - (s - a))$
\item $\varphi_{4} :\equiv \forall z\,\forall w\,(z = w \to 0 < z \to 0 < w) \to \forall w\,(a = w \to 0 < a \to 0 < w)$
\item $\varphi_{5} :\equiv \forall v_{375}\,(y = v_{375} \to v_{375} = y) \to y = \underline{n} \cdot a \to \underline{n} \cdot a = y$
\item $\varphi_{6} :\equiv \forall w\,(\mathrm{UB}(w,X_{a}) \to s \leq w) \to \mathrm{UB}(s - a,X_{a}) \to s \leq s - a$
\item $\varphi_{7} :\equiv \forall X\,(\mathrm{UB}(s - a,X) \leftrightarrow s - a \in \mathbb{R} \land \forall x \in X\,(x \leq s - a))$
\item $\varphi_{8} :\equiv \forall v_{374}\,\forall v_{375}\,(v_{374} = v_{375} \to v_{375} = v_{374}) \to \forall v_{375}\,(y = v_{375} \to v_{375} = y)$
\item $\varphi_{9} :\equiv \forall X\,(X \subseteq \mathbb{R} \to X \neq \emptyset \to \exists b\,(\mathrm{UB}(b,X)) \to \exists s\,(\mathrm{IsSup}(s,X)))$
\item $\varphi_{10} :\equiv \forall x\,\forall y\,(x = y \to y = x) \to \forall y\,(\underline{S(m)} \cdot a = y \to y = \underline{S(m)} \cdot a)$
\item $\varphi_{11} :\equiv \mathrm{IsSup}(s,X_{a}) \to \neg (y \in X_{a} \to y \leq s - a) \to \neg y \leq s - a$
\item $\varphi_{12} :\equiv \forall y\,(\neg \neg (y \in X_{a} \to y \leq s - a)) \to \neg \neg (x \in X_{a} \to x \leq s - a)$
\item $\varphi_{13} :\equiv \forall m\,(\neg (m \in \mathbb{N} \land y = \underline{m} \cdot a)) \to \neg (n \in \mathbb{N} \land y = \underline{n} \cdot a)$
\item $\varphi_{14} :\equiv \mathrm{IsSup}(s,X_{a}) \leftrightarrow \mathrm{UB}(s,X_{a}) \land \forall w\,(\mathrm{UB}(w,X_{a}) \to s \leq w)$
\item $\varphi_{15} :\equiv \forall x\,\forall y\,\forall z\,(x < y \to y \leq z \to x < z) \to \forall y\,\forall z\,(s < y \to y \leq z \to s < z)$
\item $\varphi_{16} :\equiv \forall x\,(\neg \neg (x \in X_{a} \to x \leq s - a)) \to \neg \neg (x \in X_{a} \to x \leq s - a)$
\item $\varphi_{17} :\equiv \mathrm{IsSup}(s,X_{a}) \to \neg (y \in X_{a} \to y \leq s - a) \to s - a < y$
\item $\varphi_{18} :\equiv \forall x\,(\forall x\,(\neg \neg (x \in X_{a} \to x \leq s - a)) \to x \in X_{a} \to x \leq s - a)$
\item $\varphi_{19} :\equiv \mathrm{IsSup}(s,X_{a}) \to \mathrm{UB}(s - a,X_{a}) \to \mathrm{UB}(s - a,X_{a})$
\item $\varphi_{20} :\equiv X_{a} \neq \emptyset \to \exists b\,(\mathrm{UB}(b,X_{a})) \to \exists s\,(\mathrm{IsSup}(s,X_{a}))$
\item $\varphi_{21} :\equiv \mathrm{IsSup}(s,X_{a}) \to \mathrm{UB}(s - a,X_{a}) \to \neg \mathrm{UB}(s - a,X_{a})$
\item $\varphi_{22} :\equiv 0 < b \to \neg \exists n \in \mathbb{N}\,(b < \underline{n} \cdot a) \to \exists b\,(\mathrm{UB}(b,X_{a}))$
\item $\varphi_{23} :\equiv \forall y\,(s - (s - a) = y \to y = s - (s - a)) \to s - (s - a) = a \to a = s - (s - a)$
\item $\varphi_{24} :\equiv \forall y\,(s - a \in \mathbb{R} \to y \in \mathbb{R} \to 0 < y - (s - a) \to s - a < y)$
\item $\varphi_{25} :\equiv \forall z\,\forall w\,(z = w \to z = b \to w = b) \to \forall w\,(\underline{n} \cdot a = w \to \underline{n} \cdot a = b \to w = b)$
\item $\varphi_{26} :\equiv \mathrm{IsSup}(s,X_{a}) \to \neg (y \in X_{a} \to y \leq s - a) \to s - a \in \mathbb{R}$
\item $\varphi_{27} :\equiv \forall x\,\forall x\,(\neg \neg (x \in X_{a} \to x \leq s - a)) \to \forall x \in X_{a}\,(x \leq s - a)$
\item $\varphi_{28} :\equiv \mathrm{IsSup}(s,X_{a}) \to \mathrm{UB}(s - a,X_{a}) \to \forall w\,(\mathrm{UB}(w,X_{a}) \to s \leq w)$
\item $\varphi_{29} :\equiv 0 < b \to \neg \exists n \in \mathbb{N}\,(b < \underline{n} \cdot a) \to \exists s\,(\mathrm{IsSup}(s,X_{a}))$
\item $\varphi_{30} :\equiv \forall y\,(\neg \neg (y \in X_{a} \to y \leq s - a)) \to \forall x\,\forall y\,(\neg \neg (y \in X_{a} \to y \leq s - a))$
\item $\varphi_{31} :\equiv \forall m\,(\neg (m \in \mathbb{N} \land y = \underline{m} \cdot a)) \to \forall n\,\forall m\,(\neg (m \in \mathbb{N} \land y = \underline{m} \cdot a))$
\item $\varphi_{32} :\equiv \forall x\,(\neg \neg (x \in X_{a} \to x \leq s - a)) \to \forall x\,\forall x\,(\neg \neg (x \in X_{a} \to x \leq s - a))$
\item $\varphi_{33} :\equiv \forall n\,(\neg (n \in \mathbb{N} \land b < \underline{n} \cdot a)) \to \neg (n \in \mathbb{N} \land b < \underline{n} \cdot a)$
\item $\varphi_{34} :\equiv \forall x\,\forall y\,(\neg \neg (y \in X_{a} \to y \leq s - a)) \to \forall x\,(\neg \neg (x \in X_{a} \to x \leq s - a))$
\item $\varphi_{35} :\equiv \forall n\,\forall m\,(\neg (m \in \mathbb{N} \land y = \underline{m} \cdot a)) \to \forall n\,(\neg (n \in \mathbb{N} \land y = \underline{n} \cdot a))$
\item $\varphi_{36} :\equiv \neg \exists n \in \mathbb{N}\,(b < \underline{n} \cdot a) \to \mathrm{IsSup}(s,X_{a}) \to s \in \mathbb{R}$
\item $\varphi_{37} :\equiv \neg \exists n \in \mathbb{N}\,(b < \underline{n} \cdot a) \to \mathrm{IsSup}(s,X_{a}) \to a \in \mathbb{R}$
\item $\varphi_{38} :\equiv \neg \exists n \in \mathbb{N}\,(b < \underline{n} \cdot a) \to \mathrm{IsSup}(s,X_{a}) \to s - a < s$
\item $\varphi_{39} :\equiv \forall w\,(a = w \to 0 < a \to 0 < w) \to a = s - (s - a) \to 0 < a \to 0 < s - (s - a)$
\item $\varphi_{40} :\equiv \forall n \in \mathbb{N}\,(\underline{n} \in \mathbb{R}) \to 1_{\omega} \in \mathbb{N} \to \underline{1_{\omega}} \in \mathbb{R}$
\item $\varphi_{41} :\equiv 0 < a \to 0 < b \to \neg \exists n \in \mathbb{N}\,(b < \underline{n} \cdot a) \to \mathrm{UB}(b,X_{a})$
\item $\varphi_{42} :\equiv \forall x\,(\forall y \in X_{a}\,(\neg y < b \to y = b) \to x \in X_{a} \to \neg x < b \to x = b)$
\item $\varphi_{43} :\equiv \forall y \in X_{a}\,(\neg y < b \to y = b) \to \forall x \in X_{a}\,(\neg x < b \to x = b)$
\item $\varphi_{44} :\equiv \forall y\,(\forall x \in X_{a}\,(\neg x < b \to x = b) \to y \in X_{a} \to \neg y < b \to y = b)$
\item $\varphi_{45} :\equiv \forall x \in X_{a}\,(\neg x < b \to x = b) \to \forall y \in X_{a}\,(\neg y < b \to y = b)$
\item $\varphi_{46} :\equiv \neg (y \in X_{a} \to y \leq s - a) \to m \in \mathbb{N} \land y = \underline{m} \cdot a \to s < y + a$
\item $\varphi_{47} :\equiv \neg \exists n \in \mathbb{N}\,(b < \underline{n} \cdot a) \to \mathrm{IsSup}(s,X_{a}) \to s - (s - a) = a$
\item $\varphi_{48} :\equiv \neg \exists n \in \mathbb{N}\,(b < \underline{n} \cdot a) \to \mathrm{IsSup}(s,X_{a}) \to a = s - (s - a)$
\item $\varphi_{49} :\equiv \neg \exists n \in \mathbb{N}\,(b < \underline{n} \cdot a) \to \forall n\,(\neg (n \in \mathbb{N} \land b < \underline{n} \cdot a))$
\item $\varphi_{50} :\equiv \forall a\,\forall y\,(s \in \mathbb{R} \to a \in \mathbb{R} \to y \in \mathbb{R} \to s - a < y \to s < y + a)$
\item $\varphi_{51} :\equiv \neg (y \in X_{a} \to y \leq s - a) \to m \in \mathbb{N} \land y = \underline{m} \cdot a \to S(m) \in \mathbb{N}$
\item $\varphi_{52} :\equiv \neg \exists n \in \mathbb{N}\,(b < \underline{n} \cdot a) \to \mathrm{IsSup}(s,X_{a}) \to \mathrm{UB}(s,X_{a})$
\item $\varphi_{53} :\equiv \neg \exists n \in \mathbb{N}\,(b < \underline{n} \cdot a) \to \mathrm{IsSup}(s,X_{a}) \to s - a \in \mathbb{R}$
\item $\varphi_{54} :\equiv 0 < b \to \neg \exists n \in \mathbb{N}\,(b < \underline{n} \cdot a) \to \mathrm{IsSup}(s,X_{a}) \to 0 < 0$
\item $\varphi_{55} :\equiv \forall y \in X_{a}\,(\neg y < b \to y = b) \to \forall x\,\forall y \in X_{a}\,(\neg y < b \to y = b)$
\item $\varphi_{56} :\equiv \forall x \in X_{a}\,(\neg x < b \to x = b) \to \forall y\,\forall x \in X_{a}\,(\neg x < b \to x = b)$
\item $\varphi_{57} :\equiv m \in \mathbb{N} \land y = \underline{m} \cdot a \to \exists n \in \mathbb{N}\,(\underline{S(m)} \cdot a = \underline{n} \cdot a)$
\item $\varphi_{58} :\equiv \neg (y \in X_{a} \to y \leq s - a) \to m \in \mathbb{N} \land y = \underline{m} \cdot a \to y = \underline{m} \cdot a$
\item $\varphi_{59} :\equiv \neg \exists n \in \mathbb{N}\,(b < \underline{n} \cdot a) \to \mathrm{IsSup}(s,X_{a}) \to 0 < s - (s - a)$
\item $\varphi_{60} :\equiv \mathrm{IsSup}(s,X_{a}) \to \mathrm{UB}(s - a,X_{a}) \to \mathrm{UB}(s - a,X_{a}) \to s \leq s - a$
\item $\varphi_{61} :\equiv \forall x\,\forall y \in X_{a}\,(\neg y < b \to y = b) \to \forall x \in X_{a}\,(\neg x < b \to x = b)$
\item $\varphi_{62} :\equiv \forall y\,\forall x \in X_{a}\,(\neg x < b \to x = b) \to \forall y \in X_{a}\,(\neg y < b \to y = b)$
\item $\varphi_{63} :\equiv \forall x\,\forall a\,\forall y\,(x \in \mathbb{R} \to a \in \mathbb{R} \to y \in \mathbb{R} \to x - a < y \to x < y + a)$
\item $\varphi_{64} :\equiv \neg \exists n \in \mathbb{N}\,(b < \underline{n} \cdot a) \to \mathrm{IsSup}(s,X_{a}) \to \mathrm{IsSup}(s,X_{a})$
\item $\varphi_{65} :\equiv 0 < a \to 0 < b \to \neg \exists n \in \mathbb{N}\,(b < \underline{n} \cdot a) \to \underline{1_{\omega}} \in \mathbb{R}$
\item $\varphi_{66} :\equiv b \in \mathbb{R} \to 0 < a \to 0 < b \to \neg \exists n \in \mathbb{N}\,(b < \underline{n} \cdot a) \to 0 < 0$
\item $\varphi_{67} :\equiv \forall x \in X_{a}\,(x \leq s) \to \underline{S(m)} \cdot a \in X_{a} \to \underline{S(m)} \cdot a \leq s$
\item $\varphi_{68} :\equiv \forall z\,\forall w\,(z = w \to s < z + a \to s < w + a) \to \forall w\,(y = w \to s < y + a \to s < w + a)$
\item $\varphi_{69} :\equiv b \in \mathbb{R} \to 0 < a \to 0 < b \to \neg \exists n \in \mathbb{N}\,(b < \underline{n} \cdot a) \to \neg 0 < 0$
\item $\varphi_{70} :\equiv 0 < a \to 0 < b \to \neg \exists n \in \mathbb{N}\,(b < \underline{n} \cdot a) \to y \in X_{a} \to y \leq b$
\item $\varphi_{71} :\equiv m \in \mathbb{N} \land y = \underline{m} \cdot a \to S(m) \in \mathbb{N} \land \underline{S(m)} \cdot a = \underline{S(m)} \cdot a$
\item $\varphi_{72} :\equiv \neg (y \in X_{a} \to y \leq s - a) \to m \in \mathbb{N} \land y = \underline{m} \cdot a \to \underline{S(m)} \cdot a \leq s$
\item $\varphi_{73} :\equiv \mathrm{IsSup}(s,X_{a}) \to \neg (y \in X_{a} \to y \leq s - a) \to \neg (y \in X_{a} \to y \leq s - a)$
\item $\varphi_{74} :\equiv 0 < b \to \neg \exists n \in \mathbb{N}\,(b < \underline{n} \cdot a) \to \neg (n \in \mathbb{N} \land b < \underline{n} \cdot a)$
\item $\varphi_{75} :\equiv \forall z\,\forall w\,(z = w \to z < b \to w < b) \to \forall w\,(\underline{n} \cdot a = w \to \underline{n} \cdot a < b \to w < b)$
\item $\varphi_{76} :\equiv a \in \mathbb{R} \to b \in \mathbb{R} \to 0 < a \to 0 < b \to \exists n \in \mathbb{N}\,(b < \underline{n} \cdot a)$
\item $\varphi_{77} :\equiv \neg \exists n \in \mathbb{N}\,(b < \underline{n} \cdot a) \to \mathrm{IsSup}(s,X_{a}) \to \neg \mathrm{UB}(s - a,X_{a})$
\item $\varphi_{78} :\equiv \neg (y \in X_{a} \to y \leq s - a) \to m \in \mathbb{N} \land y = \underline{m} \cdot a \to \underline{S(m)} \in \mathbb{R}$
\item $\varphi_{79} :\equiv m \in \mathbb{N} \land y = \underline{m} \cdot a \to \underline{S(m)} \cdot a \in X_{a} \to \underline{S(m)} \cdot a \leq s$
\item $\varphi_{80} :\equiv a \in \mathbb{R} \land b \in \mathbb{R} \land 0 < a \land 0 < b \land \neg \exists n \in \mathbb{N}\,(b < \underline{n} \cdot a)$
\item $\varphi_{81} :\equiv 0 < a \to 0 < b \to \neg \exists n \in \mathbb{N}\,(b < \underline{n} \cdot a) \to \forall y \in X_{a}\,(y \leq b)$
\item $\varphi_{82} :\equiv b \in \mathbb{R} \to 0 < a \to 0 < b \to \neg \exists n \in \mathbb{N}\,(b < \underline{n} \cdot a) \to a \in \mathbb{R}$
\item $\varphi_{83} :\equiv \neg (y \in X_{a} \to y \leq s - a) \to m \in \mathbb{N} \land y = \underline{m} \cdot a \to s < \underline{S(m)} \cdot a$
\item $\varphi_{84} :\equiv \mathrm{IsSup}(s,X_{a}) \to \neg (y \in X_{a} \to y \leq s - a) \to \exists n \in \mathbb{N}\,(y = \underline{n} \cdot a)$
\item $\varphi_{85} :\equiv 0 < b \to \neg \exists n \in \mathbb{N}\,(b < \underline{n} \cdot a) \to \neg \exists n \in \mathbb{N}\,(b < \underline{n} \cdot a)$
\item $\varphi_{86} :\equiv 0 < a \to 0 < b \to \neg \exists n \in \mathbb{N}\,(b < \underline{n} \cdot a) \to \underline{1_{\omega}} \cdot a \in \mathbb{R}$
\item $\varphi_{87} :\equiv \neg \exists n \in \mathbb{N}\,(b < \underline{n} \cdot a) \to \exists n \in \mathbb{N}\,(\underline{1_{\omega}} \cdot a = \underline{n} \cdot a)$
\item $\varphi_{88} :\equiv \forall w\,(\underline{n} \cdot a = w \to \underline{n} \cdot a = b \to w = b) \to \underline{n} \cdot a = y \to \underline{n} \cdot a = b \to y = b$
\item $\varphi_{89} :\equiv \neg (y \in X_{a} \to y \leq s - a) \to m \in \mathbb{N} \land y = \underline{m} \cdot a \to s < \underline{m} \cdot a + a$
\item $\varphi_{90} :\equiv \underline{m} \cdot a + a = \underline{S(m)} \cdot a \to s < \underline{m} \cdot a + a \to s < \underline{S(m)} \cdot a$
\item $\varphi_{91} :\equiv \neg (y \in X_{a} \to y \leq s - a) \to m \in \mathbb{N} \land y = \underline{m} \cdot a \to \forall x \in X_{a}\,(x \leq s)$
\item $\varphi_{92} :\equiv 0 < a \to 0 < b \to \neg \exists n \in \mathbb{N}\,(b < \underline{n} \cdot a) \to u \in X_{a} \to u \in \mathbb{R}$
\item $\varphi_{93} :\equiv \forall y\,\forall z\,(s < y \to y \leq z \to s < z) \to \forall z\,(s < \underline{S(m)} \cdot a \to \underline{S(m)} \cdot a \leq z \to s < z)$
\item $\varphi_{94} :\equiv \neg (y \in X_{a} \to y \leq s - a) \to m \in \mathbb{N} \land y = \underline{m} \cdot a \to \underline{S(m)} \cdot a \in \mathbb{R}$
\item $\varphi_{95} :\equiv \neg \exists n \in \mathbb{N}\,(b < \underline{n} \cdot a) \to y \in X_{a} \to \exists n \in \mathbb{N}\,(y = \underline{n} \cdot a)$
\item $\varphi_{96} :\equiv \neg \exists n \in \mathbb{N}\,(b < \underline{n} \cdot a) \to 1_{\omega} \in \mathbb{N} \land \underline{1_{\omega}} \cdot a = \underline{1_{\omega}} \cdot a$
\item $\varphi_{97} :\equiv 0 < a \to 0 < b \to \neg \exists n \in \mathbb{N}\,(b < \underline{n} \cdot a) \to \underline{1_{\omega}} \cdot a \in X_{a}$
\item $\varphi_{98} :\equiv \forall y\,(s \in \mathbb{R} \to y \in \mathbb{R} \to s - (s - y) = y) \to s \in \mathbb{R} \to a \in \mathbb{R} \to s - (s - a) = a$
\item $\varphi_{99} :\equiv \neg (y \in X_{a} \to y \leq s - a) \to m \in \mathbb{N} \land y = \underline{m} \cdot a \to m \in \mathbb{N} \land y = \underline{m} \cdot a$
\item $\varphi_{100} :\equiv \mathrm{IsSup}(s,X_{a}) \to \neg (y \in X_{a} \to y \leq s - a) \to m \in \mathbb{N} \land y = \underline{m} \cdot a \to s < s$
\item $\varphi_{101} :\equiv \mathrm{IsSup}(s,X_{a}) \to \neg (y \in X_{a} \to y \leq s - a) \to m \in \mathbb{N} \land y = \underline{m} \cdot a \to 0 < 0$
\item $\varphi_{102} :\equiv 0 < a \to 0 < b \to \neg \exists n \in \mathbb{N}\,(b < \underline{n} \cdot a) \to \forall u \in X_{a}\,(u \in \mathbb{R})$
\item $\varphi_{103} :\equiv \neg \exists n \in \mathbb{N}\,(b < \underline{n} \cdot a) \to \mathrm{IsSup}(s,X_{a}) \to \neg \forall x \in X_{a}\,(x \leq s - a)$
\item $\varphi_{104} :\equiv \neg \exists n \in \mathbb{N}\,(b < \underline{n} \cdot a) \to \mathrm{IsSup}(s,X_{a}) \to \exists x\,(\neg (x \in X_{a} \to x \leq s - a))$
\item $\varphi_{105} :\equiv \neg \exists n \in \mathbb{N}\,(b < \underline{n} \cdot a) \to y \in X_{a} \to n \in \mathbb{N} \land y = \underline{n} \cdot a \to y \leq b$
\item $\varphi_{106} :\equiv 0 < a \to 0 < b \to \neg \exists n \in \mathbb{N}\,(b < \underline{n} \cdot a) \to y \in X_{a} \to y \in X_{a}$
\item $\varphi_{107} :\equiv b \in \mathbb{R} \to 0 < a \to 0 < b \to \neg \exists n \in \mathbb{N}\,(b < \underline{n} \cdot a) \to X_{a} \neq \emptyset$
\item $\varphi_{108} :\equiv b \in \mathbb{R} \to 0 < a \to 0 < b \to \neg \exists n \in \mathbb{N}\,(b < \underline{n} \cdot a) \to X_{a} \subseteq \mathbb{R}$
\item $\varphi_{109} :\equiv \forall z\,\forall w\,(z = w \to s < z \to s < w) \to \forall w\,(\underline{m} \cdot a + a = w \to s < \underline{m} \cdot a + a \to s < w)$
\item $\varphi_{110} :\equiv \neg (y \in X_{a} \to y \leq s - a) \to m \in \mathbb{N} \land y = \underline{m} \cdot a \to \underline{S(m)} \cdot a \in X_{a}$
\item $\varphi_{111} :\equiv \neg \exists n \in \mathbb{N}\,(b < \underline{n} \cdot a) \to \mathrm{IsSup}(s,X_{a}) \to \mathrm{UB}(s - a,X_{a}) \to s \leq s - a$
\item $\varphi_{112} :\equiv \mathrm{IsSup}(s,X_{a}) \to \neg (y \in X_{a} \to y \leq s - a) \to m \in \mathbb{N} \land y = \underline{m} \cdot a \to m \in \mathbb{N}$
\item $\varphi_{113} :\equiv \forall y\,(s \in \mathbb{R} \to y \in \mathbb{R} \to s - y \in \mathbb{R}) \to s \in \mathbb{R} \to a \in \mathbb{R} \to s - a \in \mathbb{R}$
\item $\varphi_{114} :\equiv \forall x\,\forall y\,(x \subseteq y \leftrightarrow \forall u \in x\,(u \in y)) \to \forall y\,(X_{a} \subseteq y \leftrightarrow \forall u \in X_{a}\,(u \in y))$
\item $\varphi_{115} :\equiv \mathrm{IsSup}(s,X_{a}) \to \neg (y \in X_{a} \to y \leq s - a) \to m \in \mathbb{N} \land y = \underline{m} \cdot a \to a \in \mathbb{R}$
\item $\varphi_{116} :\equiv \forall a\,\forall X\,(\mathrm{UB}(a,X) \leftrightarrow a \in \mathbb{R} \land \forall x \in X\,(x \leq a)) \to \forall X\,(\mathrm{UB}(s,X) \leftrightarrow s \in \mathbb{R} \land \forall x \in X\,(x \leq s))$
\item $\varphi_{117} :\equiv \forall w\,(y = w \to s < y + a \to s < w + a) \to y = \underline{m} \cdot a \to s < y + a \to s < \underline{m} \cdot a + a$
\item $\varphi_{118} :\equiv \forall x\,\forall y\,(x \in \mathbb{R} \to y \in \mathbb{R} \to x - (x - y) = y) \to \forall y\,(s \in \mathbb{R} \to y \in \mathbb{R} \to s - (s - y) = y)$
\item $\varphi_{119} :\equiv \forall x\,\forall y\,(x \in \mathbb{R} \to y \in \mathbb{R} \to \neg x < y \to y \leq x) \to \forall y\,(b \in \mathbb{R} \to y \in \mathbb{R} \to \neg b < y \to y \leq b)$
\item $\varphi_{120} :\equiv \forall a\,\forall X\,(\mathrm{UB}(a,X) \leftrightarrow a \in \mathbb{R} \land \forall x \in X\,(x \leq a)) \to \forall X\,(\mathrm{UB}(b,X) \leftrightarrow b \in \mathbb{R} \land \forall x \in X\,(x \leq b))$
\item $\varphi_{121} :\equiv \neg (y \in X_{a} \to y \leq s - a) \to m \in \mathbb{N} \land y = \underline{m} \cdot a \to \underline{S(m)} \cdot a = \underline{m} \cdot a + a$
\item $\varphi_{122} :\equiv \neg (y \in X_{a} \to y \leq s - a) \to m \in \mathbb{N} \land y = \underline{m} \cdot a \to \underline{m} \cdot a + a = \underline{S(m)} \cdot a$
\item $\varphi_{123} :\equiv \neg \exists n \in \mathbb{N}\,(b < \underline{n} \cdot a) \to \mathrm{IsSup}(s,X_{a}) \to \neg (y \in X_{a} \to y \leq s - a) \to 0 < 0$
\item $\varphi_{124} :\equiv \neg \exists n \in \mathbb{N}\,(b < \underline{n} \cdot a) \to y \in X_{a} \to n \in \mathbb{N} \land y = \underline{n} \cdot a \to n \in \mathbb{N}$
\item $\varphi_{125} :\equiv \forall w\,(\underline{n} \cdot a = w \to \underline{n} \cdot a < b \to w < b) \to \underline{n} \cdot a = y \to \underline{n} \cdot a < b \to y < b$
\item $\varphi_{126} :\equiv \neg \exists n \in \mathbb{N}\,(b < \underline{n} \cdot a) \to y \in X_{a} \to n \in \mathbb{N} \land y = \underline{n} \cdot a \to a \in \mathbb{R}$
\item $\varphi_{127} :\equiv \neg \exists n \in \mathbb{N}\,(b < \underline{n} \cdot a) \to y \in X_{a} \to n \in \mathbb{N} \land y = \underline{n} \cdot a \to b \in \mathbb{R}$
\item $\varphi_{128} :\equiv \forall n\,(\neg (n \in \mathbb{N} \land \underline{S(m)} \cdot a = \underline{n} \cdot a)) \to \neg (S(m) \in \mathbb{N} \land \underline{S(m)} \cdot a = \underline{S(m)} \cdot a)$
\item $\varphi_{129} :\equiv \forall x\,\forall y\,(x \in \mathbb{R} \to y \in \mathbb{R} \to x - y \in \mathbb{R}) \to \forall y\,(s \in \mathbb{R} \to y \in \mathbb{R} \to s - y \in \mathbb{R})$
\item $\varphi_{130} :\equiv \forall n\,(\neg (n \in \mathbb{N} \land \underline{1_{\omega}} \cdot a = \underline{n} \cdot a)) \to \neg (1_{\omega} \in \mathbb{N} \land \underline{1_{\omega}} \cdot a = \underline{1_{\omega}} \cdot a)$
\item $\varphi_{131} :\equiv \neg \exists n \in \mathbb{N}\,(b < \underline{n} \cdot a) \to \mathrm{IsSup}(s,X_{a}) \to \neg (y \in X_{a} \to y \leq s - a) \to y \in \mathbb{R}$
\item $\varphi_{132} :\equiv \neg \exists n \in \mathbb{N}\,(b < \underline{n} \cdot a) \to \mathrm{IsSup}(s,X_{a}) \to \neg (y \in X_{a} \to y \leq s - a) \to a \in \mathbb{R}$
\item $\varphi_{133} :\equiv \neg \exists n \in \mathbb{N}\,(b < \underline{n} \cdot a) \to y \in X_{a} \to n \in \mathbb{N} \land y = \underline{n} \cdot a \to \underline{n} \cdot a \leq b$
\item $\varphi_{134} :\equiv \neg \exists n \in \mathbb{N}\,(b < \underline{n} \cdot a) \to y \in X_{a} \to n \in \mathbb{N} \land y = \underline{n} \cdot a \to y = \underline{n} \cdot a$
\item $\varphi_{135} :\equiv \neg \exists n \in \mathbb{N}\,(b < \underline{n} \cdot a) \to y \in X_{a} \to n \in \mathbb{N} \land y = \underline{n} \cdot a \to \underline{n} \cdot a = y$
\item $\varphi_{136} :\equiv \forall X\,(\mathrm{UB}(s,X) \leftrightarrow s \in \mathbb{R} \land \forall x \in X\,(x \leq s)) \to (\mathrm{UB}(s,X_{a}) \leftrightarrow s \in \mathbb{R} \land \forall x \in X_{a}\,(x \leq s))$
\item $\varphi_{137} :\equiv \forall X\,(\mathrm{UB}(b,X) \leftrightarrow b \in \mathbb{R} \land \forall x \in X\,(x \leq b)) \to (\mathrm{UB}(b,X_{a}) \leftrightarrow b \in \mathbb{R} \land \forall x \in X_{a}\,(x \leq b))$
\item $\varphi_{138} :\equiv \neg \exists n \in \mathbb{N}\,(b < \underline{n} \cdot a) \to y \in X_{a} \to n \in \mathbb{N} \land y = \underline{n} \cdot a \to \underline{n} \in \mathbb{R}$
\item $\varphi_{139} :\equiv \neg \exists n \in \mathbb{N}\,(b < \underline{n} \cdot a) \to \mathrm{IsSup}(s,X_{a}) \to \neg (y \in X_{a} \to y \leq s - a) \to y \in X_{a}$
\item $\varphi_{140} :\equiv \neg \exists n \in \mathbb{N}\,(b < \underline{n} \cdot a) \to y \in X_{a} \to n \in \mathbb{N} \land y = \underline{n} \cdot a \to \neg b < \underline{n} \cdot a$
\item $\varphi_{141} :\equiv \underline{S(m)} \cdot a \in X_{a} \leftrightarrow \underline{S(m)} \cdot a \in \mathbb{R} \land \exists n \in \mathbb{N}\,(\underline{S(m)} \cdot a = \underline{n} \cdot a)$
\item $\varphi_{142} :\equiv \underline{1_{\omega}} \cdot a \in X_{a} \leftrightarrow \underline{1_{\omega}} \cdot a \in \mathbb{R} \land \exists n \in \mathbb{N}\,(\underline{1_{\omega}} \cdot a = \underline{n} \cdot a)$
\item $\varphi_{143} :\equiv \forall y\,(X_{a} \subseteq y \leftrightarrow \forall u \in X_{a}\,(u \in y)) \to (X_{a} \subseteq \mathbb{R} \leftrightarrow \forall u \in X_{a}\,(u \in \mathbb{R}))$
\item $\varphi_{144} :\equiv \neg \exists n \in \mathbb{N}\,(b < \underline{n} \cdot a) \to y \in X_{a} \to n \in \mathbb{N} \land y = \underline{n} \cdot a \to \underline{n} \cdot a \in \mathbb{R}$
\item $\varphi_{145} :\equiv \forall z\,(s < \underline{S(m)} \cdot a \to \underline{S(m)} \cdot a \leq z \to s < z) \to s < \underline{S(m)} \cdot a \to \underline{S(m)} \cdot a \leq s \to s < s$
\item $\varphi_{146} :\equiv \forall x\,\forall y\,(x \in \mathbb{R} \to y \in \mathbb{R} \to \neg y \leq x \to x < y) \to \forall y\,(s - a \in \mathbb{R} \to y \in \mathbb{R} \to \neg y \leq s - a \to s - a < y)$
\item $\varphi_{147} :\equiv \forall x\,\forall y\,(x \in \mathbb{R} \to y \in \mathbb{R} \to x \cdot y \in \mathbb{R}) \to \forall y\,(\underline{n} \in \mathbb{R} \to y \in \mathbb{R} \to \underline{n} \cdot y \in \mathbb{R})$
\item $\varphi_{148} :\equiv \neg \exists n \in \mathbb{N}\,(b < \underline{n} \cdot a) \to y \in X_{a} \to n \in \mathbb{N} \land y = \underline{n} \cdot a \to n \in \mathbb{N} \land y = \underline{n} \cdot a$
\item $\varphi_{149} :\equiv \forall x\,\forall y\,(x \in \mathbb{R} \to y \in \mathbb{R} \to x \cdot y \in \mathbb{R}) \to \forall y\,(\underline{S(m)} \in \mathbb{R} \to y \in \mathbb{R} \to \underline{S(m)} \cdot y \in \mathbb{R})$
\item $\varphi_{150} :\equiv \forall y\,(\underline{S(m)} \cdot a = y \to y = \underline{S(m)} \cdot a) \to \underline{S(m)} \cdot a = \underline{m} \cdot a + a \to \underline{m} \cdot a + a = \underline{S(m)} \cdot a$
\item $\varphi_{151} :\equiv \forall y\,(\underline{n} \in \mathbb{R} \to y \in \mathbb{R} \to \underline{n} \cdot y \in \mathbb{R}) \to \underline{n} \in \mathbb{R} \to a \in \mathbb{R} \to \underline{n} \cdot a \in \mathbb{R}$
\item $\varphi_{152} :\equiv \forall x\,\forall y\,(x \in \mathbb{R} \to y \in \mathbb{R} \to x \cdot y \in \mathbb{R}) \to \forall y\,(\underline{1_{\omega}} \in \mathbb{R} \to y \in \mathbb{R} \to \underline{1_{\omega}} \cdot y \in \mathbb{R})$
\item $\varphi_{153} :\equiv \forall y\,(s - a \in \mathbb{R} \to y \in \mathbb{R} \to \neg y \leq s - a \to s - a < y) \to s - a \in \mathbb{R} \to y \in \mathbb{R} \to \neg y \leq s - a \to s - a < y$
\item $\varphi_{154} :\equiv \forall y\,(b \in \mathbb{R} \to y \in \mathbb{R} \to \neg b < y \to y \leq b) \to b \in \mathbb{R} \to \underline{n} \cdot a \in \mathbb{R} \to \neg b < \underline{n} \cdot a \to \underline{n} \cdot a \leq b$
\item $\varphi_{155} :\equiv \varphi_{16} \to \forall x\,(\neg \neg (x \in X_{a} \to x \leq s - a)) \to x \in X_{a} \to x \leq s - a$
\item $\varphi_{156} :\equiv \forall a\,(m \in \mathbb{N} \to a \in \mathbb{R} \to \underline{S(m)} \cdot a = \underline{m} \cdot a + a) \to m \in \mathbb{N} \to a \in \mathbb{R} \to \underline{S(m)} \cdot a = \underline{m} \cdot a + a$
\item $\varphi_{157} :\equiv \forall y\,(\underline{S(m)} \in \mathbb{R} \to y \in \mathbb{R} \to \underline{S(m)} \cdot y \in \mathbb{R}) \to \underline{S(m)} \in \mathbb{R} \to a \in \mathbb{R} \to \underline{S(m)} \cdot a \in \mathbb{R}$
\item $\varphi_{158} :\equiv \forall y\,(y \in X_{a} \leftrightarrow y \in \mathbb{R} \land \exists n \in \mathbb{N}\,(y = \underline{n} \cdot a)) \to (y \in X_{a} \leftrightarrow y \in \mathbb{R} \land \exists n \in \mathbb{N}\,(y = \underline{n} \cdot a))$
\item $\varphi_{159} :\equiv \forall y\,(\underline{1_{\omega}} \in \mathbb{R} \to y \in \mathbb{R} \to \underline{1_{\omega}} \cdot y \in \mathbb{R}) \to \underline{1_{\omega}} \in \mathbb{R} \to a \in \mathbb{R} \to \underline{1_{\omega}} \cdot a \in \mathbb{R}$
\item $\varphi_{160} :\equiv \forall y\,(y \in X_{a} \leftrightarrow y \in \mathbb{R} \land \exists n \in \mathbb{N}\,(y = \underline{n} \cdot a)) \to (u \in X_{a} \leftrightarrow u \in \mathbb{R} \land \exists n \in \mathbb{N}\,(u = \underline{n} \cdot a))$
\item $\varphi_{161} :\equiv \forall m\,\forall a\,(m \in \mathbb{N} \to a \in \mathbb{R} \to \underline{S(m)} \cdot a = \underline{m} \cdot a + a) \to \forall a\,(m \in \mathbb{N} \to a \in \mathbb{R} \to \underline{S(m)} \cdot a = \underline{m} \cdot a + a)$
\item $\varphi_{162} :\equiv \varphi_{7} \to (\mathrm{UB}(s - a,X_{a}) \leftrightarrow s - a \in \mathbb{R} \land \forall x \in X_{a}\,(x \leq s - a))$
\item $\varphi_{163} :\equiv \forall a\,\forall y\,(y \in X_{a} \leftrightarrow y \in \mathbb{R} \land \exists n \in \mathbb{N}\,(y = \underline{n} \cdot a)) \to \forall y\,(y \in X_{a} \leftrightarrow y \in \mathbb{R} \land \exists n \in \mathbb{N}\,(y = \underline{n} \cdot a))$
\item $\varphi_{164} :\equiv \forall y\,(s \in \mathbb{R} \to a \in \mathbb{R} \to y \in \mathbb{R} \to s - a < y \to s < y + a) \to s \in \mathbb{R} \to a \in \mathbb{R} \to y \in \mathbb{R} \to s - a < y \to s < y + a$
\item $\varphi_{165} :\equiv 0 < b \to \neg \exists n \in \mathbb{N}\,(b < \underline{n} \cdot a) \to \mathrm{IsSup}(s,X_{a}) \to \varphi_{46}$
\item $\varphi_{166} :\equiv 0 < b \to \neg \exists n \in \mathbb{N}\,(b < \underline{n} \cdot a) \to \mathrm{IsSup}(s,X_{a}) \to \varphi_{51}$
\item $\varphi_{167} :\equiv (\forall b\,(\neg \mathrm{UB}(b,X_{a})) \to \neg \mathrm{UB}(b,X_{a})) \to a \in \mathbb{R} \to b \in \mathbb{R} \to 0 < a \to \varphi_{22}$
\item $\varphi_{168} :\equiv \varphi_{80} \land \mathrm{IsSup}(s,X_{a}) \land \neg (y \in X_{a} \to y \leq s - a) \land \neg 0 < 0$
\item $\varphi_{169} :\equiv \varphi_{80} \land \mathrm{IsSup}(s,X_{a}) \land \neg 0 < 0 \to \neg \neg (y \in X_{a} \to y \leq s - a)$
\item $\varphi_{170} :\equiv 0 < b \to \neg \exists n \in \mathbb{N}\,(b < \underline{n} \cdot a) \to \mathrm{IsSup}(s,X_{a}) \to \varphi_{58}$
\item $\varphi_{171} :\equiv \forall n\,(\varphi_{80} \land y \in X_{a} \land \neg y \leq b) \to \forall n\,(\neg (n \in \mathbb{N} \land y = \underline{n} \cdot a))$
\item $\varphi_{172} :\equiv a \in \mathbb{R} \to b \in \mathbb{R} \to 0 < a \to 0 < b \to \neg \exists n \in \mathbb{N}\,(b < \underline{n} \cdot a) \to \varphi_{11}$
\item $\varphi_{173} :\equiv a \in \mathbb{R} \to b \in \mathbb{R} \to 0 < a \to 0 < b \to \neg \exists n \in \mathbb{N}\,(b < \underline{n} \cdot a) \to \varphi_{17}$
\item $\varphi_{174} :\equiv \varphi_{50} \to \forall y\,(s \in \mathbb{R} \to a \in \mathbb{R} \to y \in \mathbb{R} \to s - a < y \to s < y + a)$
\item $\varphi_{175} :\equiv 0 < b \to \neg \exists n \in \mathbb{N}\,(b < \underline{n} \cdot a) \to \mathrm{IsSup}(s,X_{a}) \to \varphi_{72}$
\item $\varphi_{176} :\equiv a \in \mathbb{R} \to b \in \mathbb{R} \to 0 < a \to 0 < b \to \neg \exists n \in \mathbb{N}\,(b < \underline{n} \cdot a) \to \varphi_{19}$
\item $\varphi_{177} :\equiv a \in \mathbb{R} \to b \in \mathbb{R} \to 0 < a \to 0 < b \to \neg \exists n \in \mathbb{N}\,(b < \underline{n} \cdot a) \to \varphi_{21}$
\item $\varphi_{178} :\equiv 0 < b \to \neg \exists n \in \mathbb{N}\,(b < \underline{n} \cdot a) \to \mathrm{IsSup}(s,X_{a}) \to \varphi_{78}$
\item $\varphi_{179} :\equiv \forall y\,(\varphi_{80} \land \mathrm{IsSup}(s,X_{a}) \land \neg 0 < 0) \to \forall y\,(\neg \neg (y \in X_{a} \to y \leq s - a))$
\item $\varphi_{180} :\equiv 0 < b \to \neg \exists n \in \mathbb{N}\,(b < \underline{n} \cdot a) \to \mathrm{IsSup}(s,X_{a}) \to \varphi_{83}$
\item $\varphi_{181} :\equiv a \in \mathbb{R} \to b \in \mathbb{R} \to 0 < a \to 0 < b \to \neg \exists n \in \mathbb{N}\,(b < \underline{n} \cdot a) \to \varphi_{26}$
\item $\varphi_{182} :\equiv (s - (s - a) = a \to a = s - (s - a)) \to a \in \mathbb{R} \to b \in \mathbb{R} \to 0 < a \to 0 < b \to \varphi_{48}$
\item $\varphi_{183} :\equiv a \in \mathbb{R} \to b \in \mathbb{R} \to 0 < a \to 0 < b \to \neg \exists n \in \mathbb{N}\,(b < \underline{n} \cdot a) \to \varphi_{28}$
\item $\varphi_{184} :\equiv 0 < b \to \neg \exists n \in \mathbb{N}\,(b < \underline{n} \cdot a) \to \mathrm{IsSup}(s,X_{a}) \to \varphi_{89}$
\item $\varphi_{185} :\equiv 0 < b \to \neg \exists n \in \mathbb{N}\,(b < \underline{n} \cdot a) \to \mathrm{IsSup}(s,X_{a}) \to \varphi_{91}$
\item $\varphi_{186} :\equiv 0 < b \to \neg \exists n \in \mathbb{N}\,(b < \underline{n} \cdot a) \to \mathrm{IsSup}(s,X_{a}) \to \varphi_{94}$
\item $\varphi_{187} :\equiv 0 < b \to \neg \exists n \in \mathbb{N}\,(b < \underline{n} \cdot a) \to \mathrm{IsSup}(s,X_{a}) \to \varphi_{99}$
\item $\varphi_{188} :\equiv (1_{\omega} \in \mathbb{N} \to \underline{1_{\omega}} \in \mathbb{R}) \to 1_{\omega} \in \mathbb{N} \to a \in \mathbb{R} \to b \in \mathbb{R} \to \varphi_{65}$
\item $\varphi_{189} :\equiv 0 < b \to \neg \exists n \in \mathbb{N}\,(b < \underline{n} \cdot a) \to \mathrm{IsSup}(s,X_{a}) \to \varphi_{110}$
\item $\varphi_{190} :\equiv a \in \mathbb{R} \to b \in \mathbb{R} \to 0 < a \to 0 < b \to \neg \exists n \in \mathbb{N}\,(b < \underline{n} \cdot a) \to \varphi_{60}$
\item $\varphi_{191} :\equiv a \in \mathbb{R} \to b \in \mathbb{R} \to 0 < a \to 0 < b \to \neg \exists n \in \mathbb{N}\,(b < \underline{n} \cdot a) \to \varphi_{73}$
\item $\varphi_{192} :\equiv 0 < b \to \neg \exists n \in \mathbb{N}\,(b < \underline{n} \cdot a) \to \mathrm{IsSup}(s,X_{a}) \to \varphi_{121}$
\item $\varphi_{193} :\equiv 0 < b \to \neg \exists n \in \mathbb{N}\,(b < \underline{n} \cdot a) \to \mathrm{IsSup}(s,X_{a}) \to \varphi_{122}$
\item $\varphi_{194} :\equiv \neg \exists n \in \mathbb{N}\,(b < \underline{n} \cdot a) \to \mathrm{IsSup}(s,X_{a}) \to \neg (y \in X_{a} \to y \leq s - a) \to \varphi_{57}$
\item $\varphi_{195} :\equiv a \in \mathbb{R} \to b \in \mathbb{R} \to 0 < a \to 0 < b \to \neg \exists n \in \mathbb{N}\,(b < \underline{n} \cdot a) \to \varphi_{84}$
\item $\varphi_{196} :\equiv (a = s - (s - a) \to 0 < a \to 0 < s - (s - a)) \to a \in \mathbb{R} \to b \in \mathbb{R} \to 0 < a \to 0 < b \to \varphi_{59}$
\item $\varphi_{197} :\equiv \underline{1_{\omega}} \cdot a = \underline{1_{\omega}} \cdot a \to a \in \mathbb{R} \to b \in \mathbb{R} \to 0 < a \to 0 < b \to \varphi_{96}$
\item $\varphi_{198} :\equiv \neg \exists n \in \mathbb{N}\,(b < \underline{n} \cdot a) \to \mathrm{IsSup}(s,X_{a}) \to \neg (y \in X_{a} \to y \leq s - a) \to \varphi_{71}$
\item $\varphi_{199} :\equiv \neg \exists n \in \mathbb{N}\,(b < \underline{n} \cdot a) \to \mathrm{IsSup}(s,X_{a}) \to \neg (y \in X_{a} \to y \leq s - a) \to \varphi_{79}$
\item $\varphi_{200} :\equiv a \in \mathbb{R} \to b \in \mathbb{R} \to 0 < a \to 0 < b \to \neg \exists n \in \mathbb{N}\,(b < \underline{n} \cdot a) \to \varphi_{100}$
\item $\varphi_{201} :\equiv a \in \mathbb{R} \to b \in \mathbb{R} \to 0 < a \to 0 < b \to \neg \exists n \in \mathbb{N}\,(b < \underline{n} \cdot a) \to \varphi_{101}$
\item $\varphi_{202} :\equiv (\mathrm{UB}(s,X_{a}) \leftrightarrow s \in \mathbb{R} \land \forall x \in X_{a}\,(x \leq s)) \to a \in \mathbb{R} \to b \in \mathbb{R} \to 0 < a \to 0 < b \to \varphi_{36}$
\item $\varphi_{203} :\equiv a \in \mathbb{R} \to b \in \mathbb{R} \to 0 < a \to 0 < b \to \neg \exists n \in \mathbb{N}\,(b < \underline{n} \cdot a) \to \varphi_{112}$
\item $\varphi_{204} :\equiv a \in \mathbb{R} \to b \in \mathbb{R} \to 0 < a \to 0 < b \to \neg \exists n \in \mathbb{N}\,(b < \underline{n} \cdot a) \to \varphi_{115}$
\item $\varphi_{205} :\equiv (\underline{n} \cdot a = y \to \underline{n} \cdot a = b \to y = b) \to a \in \mathbb{R} \to b \in \mathbb{R} \to 0 < a \to 0 < b \to \varphi_{105}$
\item $\varphi_{206} :\equiv (u \in X_{a} \leftrightarrow u \in \mathbb{R} \land \exists n \in \mathbb{N}\,(u = \underline{n} \cdot a)) \to a \in \mathbb{R} \to b \in \mathbb{R} \to \varphi_{92}$
\item $\varphi_{207} :\equiv (y = \underline{n} \cdot a \to \underline{n} \cdot a = y) \to a \in \mathbb{R} \to b \in \mathbb{R} \to 0 < a \to 0 < b \to \varphi_{135}$
\item $\varphi_{208} :\equiv (y \in X_{a} \leftrightarrow y \in \mathbb{R} \land \exists n \in \mathbb{N}\,(y = \underline{n} \cdot a)) \to a \in \mathbb{R} \to b \in \mathbb{R} \to 0 < a \to 0 < b \to \varphi_{95}$
\item $\varphi_{209} :\equiv (a \in \mathbb{R} \to b \in \mathbb{R} \to 0 < a \to \varphi_{22}) \to a \in \mathbb{R} \to b \in \mathbb{R} \to 0 < a \to \varphi_{29}$
\item $\varphi_{210} :\equiv (y \in X_{a} \leftrightarrow y \in \mathbb{R} \land \exists n \in \mathbb{N}\,(y = \underline{n} \cdot a)) \to a \in \mathbb{R} \to b \in \mathbb{R} \to 0 < a \to 0 < b \to \varphi_{131}$
\item $\varphi_{211} :\equiv \forall y\,(\varphi_{80} \to y \in X_{a} \to y \leq b) \to \forall y\,(\varphi_{80}) \to \forall y \in X_{a}\,(y \leq b)$
\item $\varphi_{212} :\equiv (a \in \mathbb{R} \to b \in \mathbb{R} \to 0 < a \to \varphi_{54}) \to \varphi_{80} \land \neg 0 < 0 \to \neg \mathrm{IsSup}(s,X_{a})$
\item $\varphi_{213} :\equiv \varphi_{80} \land \mathrm{IsSup}(s,X_{a}) \land \neg 0 < 0 \to \forall y\,(\varphi_{80} \land \mathrm{IsSup}(s,X_{a}) \land \neg 0 < 0)$
\item $\varphi_{214} :\equiv (a \in \mathbb{R} \to b \in \mathbb{R} \to 0 < a \to 0 < b \to \varphi_{37}) \to a \in \mathbb{R} \to b \in \mathbb{R} \to 0 < a \to 0 < b \to \varphi_{47}$
\item $\varphi_{215} :\equiv \forall u\,(\varphi_{80} \to u \in X_{a} \to u \in \mathbb{R}) \to \forall u\,(\varphi_{80}) \to \forall u \in X_{a}\,(u \in \mathbb{R})$
\item $\varphi_{216} :\equiv \forall s\,(\varphi_{80} \land \neg 0 < 0 \to \neg \mathrm{IsSup}(s,X_{a})) \to \forall s\,(\varphi_{80} \land \neg 0 < 0) \to \forall s\,(\neg \mathrm{IsSup}(s,X_{a}))$
\item $\varphi_{217} :\equiv (a \in \mathbb{R} \to b \in \mathbb{R} \to 0 < a \to 0 < b \to \varphi_{37}) \to a \in \mathbb{R} \to b \in \mathbb{R} \to 0 < a \to 0 < b \to \varphi_{53}$
\item $\varphi_{218} :\equiv (a \in \mathbb{R} \to b \in \mathbb{R} \to 0 < a \to \varphi_{85}) \to a \in \mathbb{R} \to b \in \mathbb{R} \to 0 < a \to 0 < b \to \varphi_{49}$
\item $\varphi_{219} :\equiv (a \in \mathbb{R} \to b \in \mathbb{R} \to 0 < a \to 0 < b \to \varphi_{59}) \to a \in \mathbb{R} \to b \in \mathbb{R} \to 0 < a \to 0 < b \to \varphi_{38}$
\item $\varphi_{220} :\equiv (a \in \mathbb{R} \to b \in \mathbb{R} \to 0 < a \to 0 < b \to \varphi_{87}) \to a \in \mathbb{R} \to b \in \mathbb{R} \to \varphi_{97}$
\item $\varphi_{221} :\equiv (y = \underline{m} \cdot a \to s < y + a \to s < \underline{m} \cdot a + a) \to a \in \mathbb{R} \to b \in \mathbb{R} \to 0 < a \to \varphi_{184}$
\item $\varphi_{222} :\equiv (\mathrm{UB}(s,X_{a}) \leftrightarrow s \in \mathbb{R} \land \forall x \in X_{a}\,(x \leq s)) \to a \in \mathbb{R} \to b \in \mathbb{R} \to 0 < a \to \varphi_{185}$
\item $\varphi_{223} :\equiv \underline{S(m)} \cdot a = \underline{S(m)} \cdot a \to a \in \mathbb{R} \to b \in \mathbb{R} \to 0 < a \to 0 < b \to \varphi_{198}$
\item $\varphi_{224} :\equiv (a \in \mathbb{R} \to b \in \mathbb{R} \to 0 < a \to 0 < b \to \varphi_{53}) \to a \in \mathbb{R} \to b \in \mathbb{R} \to 0 < a \to 0 < b \to \varphi_{103}$
\item $\varphi_{225} :\equiv (\underline{1_{\omega}} \cdot a \in X_{a} \to X_{a} \neq \emptyset) \to (a \in \mathbb{R} \to b \in \mathbb{R} \to \varphi_{97}) \to a \in \mathbb{R} \to \varphi_{107}$
\item $\varphi_{226} :\equiv (\underline{S(m)} \cdot a = \underline{m} \cdot a + a \to \underline{m} \cdot a + a = \underline{S(m)} \cdot a) \to a \in \mathbb{R} \to b \in \mathbb{R} \to 0 < a \to \varphi_{193}$
\item $\varphi_{227} :\equiv (a \in \mathbb{R} \to b \in \mathbb{R} \to 0 < a \to 0 < b \to \varphi_{105}) \to \varphi_{80} \land y \in X_{a} \land \neg y \leq b \to \neg (n \in \mathbb{N} \land y = \underline{n} \cdot a)$
\item $\varphi_{228} :\equiv (a \in \mathbb{R} \to b \in \mathbb{R} \to \varphi_{102}) \to (X_{a} \subseteq \mathbb{R} \leftrightarrow \forall u \in X_{a}\,(u \in \mathbb{R})) \to a \in \mathbb{R} \to \varphi_{108}$
\item $\varphi_{229} :\equiv \forall n\,(\varphi_{80} \land y \in X_{a} \land \neg y \leq b \to \neg (n \in \mathbb{N} \land y = \underline{n} \cdot a)) \to \varphi_{171}$
\item $\varphi_{230} :\equiv (a \in \mathbb{R} \to b \in \mathbb{R} \to 0 < a \to 0 < b \to \varphi_{124}) \to a \in \mathbb{R} \to b \in \mathbb{R} \to 0 < a \to 0 < b \to \varphi_{138}$
\item $\varphi_{231} :\equiv (s \leq s - a \to \neg s - a < s) \to (a \in \mathbb{R} \to b \in \mathbb{R} \to 0 < a \to 0 < b \to \varphi_{38}) \to \varphi_{177}$
\item $\varphi_{232} :\equiv (a \in \mathbb{R} \to b \in \mathbb{R} \to 0 < a \to 0 < b \to \varphi_{124}) \to a \in \mathbb{R} \to b \in \mathbb{R} \to 0 < a \to 0 < b \to \varphi_{140}$
\item $\varphi_{233} :\equiv (a \in \mathbb{R} \to b \in \mathbb{R} \to 0 < a \to 0 < b \to \varphi_{126}) \to a \in \mathbb{R} \to b \in \mathbb{R} \to 0 < a \to 0 < b \to \varphi_{144}$
\item $\varphi_{234} :\equiv (a \in \mathbb{R} \to b \in \mathbb{R} \to 0 < a \to 0 < b \to \varphi_{140}) \to a \in \mathbb{R} \to b \in \mathbb{R} \to 0 < a \to 0 < b \to \varphi_{133}$
\item $\varphi_{235} :\equiv (a \in \mathbb{R} \to b \in \mathbb{R} \to 0 < a \to 0 < b \to \varphi_{148}) \to a \in \mathbb{R} \to b \in \mathbb{R} \to 0 < a \to 0 < b \to \varphi_{124}$
\item $\varphi_{236} :\equiv (a \in \mathbb{R} \to b \in \mathbb{R} \to 0 < a \to 0 < b \to \varphi_{49}) \to \varphi_{33} \to a \in \mathbb{R} \to b \in \mathbb{R} \to 0 < a \to \varphi_{74}$
\item $\varphi_{237} :\equiv (a \in \mathbb{R} \to b \in \mathbb{R} \to 0 < a \to 0 < b \to \varphi_{148}) \to a \in \mathbb{R} \to b \in \mathbb{R} \to 0 < a \to 0 < b \to \varphi_{134}$
\item $\varphi_{238} :\equiv (a \in \mathbb{R} \to b \in \mathbb{R} \to 0 < a \to 0 < b \to \varphi_{64}) \to \varphi_{14} \to a \in \mathbb{R} \to b \in \mathbb{R} \to 0 < a \to 0 < b \to \varphi_{52}$
\item $\varphi_{239} :\equiv (a \in \mathbb{R} \to b \in \mathbb{R} \to \varphi_{65}) \to (a \in \mathbb{R} \to \varphi_{82}) \to a \in \mathbb{R} \to b \in \mathbb{R} \to \varphi_{86}$
\item $\varphi_{240} :\equiv (\varphi_{80} \land y \in X_{a} \land \neg y \leq b \to \forall n\,(\varphi_{80} \land y \in X_{a} \land \neg y \leq b)) \to a \in \mathbb{R} \to b \in \mathbb{R} \to \varphi_{70}$
\item $\varphi_{241} :\equiv (\mathrm{UB}(s - a,X_{a}) \leftrightarrow s - a \in \mathbb{R} \land \forall x \in X_{a}\,(x \leq s - a)) \to \varphi_{224}$
\item $\varphi_{242} :\equiv (a \in \mathbb{R} \to b \in \mathbb{R} \to 0 < a \to 0 < b \to \varphi_{48}) \to (a \in \mathbb{R} \to b \in \mathbb{R} \to 0 < a \to 0 < b \to 0 < a) \to \varphi_{196}$
\item $\varphi_{243} :\equiv (\forall n\,(\varphi_{13}) \to \varphi_{35}) \to \varphi_{31} \to a \in \mathbb{R} \to b \in \mathbb{R} \to 0 < a \to 0 < b \to \varphi_{123}$
\item $\varphi_{244} :\equiv \forall m\,(\varphi_{168} \to \neg (m \in \mathbb{N} \land y = \underline{m} \cdot a)) \to \forall m\,(\varphi_{168}) \to \forall m\,(\neg (m \in \mathbb{N} \land y = \underline{m} \cdot a))$
\item $\varphi_{245} :\equiv (a \in \mathbb{R} \to b \in \mathbb{R} \to 0 < a \to 0 < b \to \varphi_{96}) \to \varphi_{130} \to a \in \mathbb{R} \to b \in \mathbb{R} \to 0 < a \to 0 < b \to \varphi_{87}$
\item $\varphi_{246} :\equiv (a \in \mathbb{R} \to b \in \mathbb{R} \to 0 < a \to \varphi_{166}) \to a \in \mathbb{R} \to b \in \mathbb{R} \to 0 < a \to \varphi_{178}$
\item $\varphi_{247} :\equiv (a \in \mathbb{R} \to b \in \mathbb{R} \to 0 < a \to 0 < b \to \varphi_{133}) \to (\underline{n} \cdot a = y \to \underline{n} \cdot a < b \to y < b) \to \varphi_{205}$
\item $\varphi_{248} :\equiv (a \in \mathbb{R} \to b \in \mathbb{R} \to 0 < a \to \varphi_{187}) \to a \in \mathbb{R} \to b \in \mathbb{R} \to 0 < a \to \varphi_{170}$
\item $\varphi_{249} :\equiv (a \in \mathbb{R} \to b \in \mathbb{R} \to 0 < a \to \varphi_{189}) \to a \in \mathbb{R} \to b \in \mathbb{R} \to 0 < a \to \varphi_{175}$
\item $\varphi_{250} :\equiv (a \in \mathbb{R} \to b \in \mathbb{R} \to 0 < a \to 0 < b \to \varphi_{139}) \to (y \in X_{a} \leftrightarrow y \in \mathbb{R} \land \exists n \in \mathbb{N}\,(y = \underline{n} \cdot a)) \to \varphi_{195}$
\item $\varphi_{251} :\equiv (m \in \mathbb{N} \to S(m) \in \mathbb{N}) \to \varphi_{203} \to a \in \mathbb{R} \to b \in \mathbb{R} \to 0 < a \to \varphi_{166}$
\item $\varphi_{252} :\equiv \forall x\,(\varphi_{12}) \to (\forall x\,(\varphi_{12}) \to \varphi_{34}) \to \varphi_{30} \to a \in \mathbb{R} \to b \in \mathbb{R} \to 0 < a \to \varphi_{54}$
\item $\varphi_{253} :\equiv (s \in \mathbb{R} \to a \in \mathbb{R} \to s - (s - a) = a) \to (a \in \mathbb{R} \to b \in \mathbb{R} \to 0 < a \to 0 < b \to \varphi_{36}) \to \varphi_{214}$
\item $\varphi_{254} :\equiv (s \in \mathbb{R} \to a \in \mathbb{R} \to s - a \in \mathbb{R}) \to (a \in \mathbb{R} \to b \in \mathbb{R} \to 0 < a \to 0 < b \to \varphi_{36}) \to \varphi_{217}$
\item $\varphi_{255} :\equiv (a \in \mathbb{R} \to b \in \mathbb{R} \to 0 < a \to 0 < b \to \varphi_{194}) \to a \in \mathbb{R} \to b \in \mathbb{R} \to 0 < a \to \varphi_{189}$
\item $\varphi_{256} :\equiv (a \in \mathbb{R} \to b \in \mathbb{R} \to 0 < a \to 0 < b \to \varphi_{36}) \to (a \in \mathbb{R} \to b \in \mathbb{R} \to 0 < a \to 0 < b \to \varphi_{132}) \to \varphi_{181}$
\item $\varphi_{257} :\equiv (\mathrm{UB}(b,X_{a}) \leftrightarrow b \in \mathbb{R} \land \forall x \in X_{a}\,(x \leq b)) \to (a \in \mathbb{R} \to b \in \mathbb{R} \to \varphi_{81}) \to \varphi_{45} \to \varphi_{43} \to a \in \mathbb{R} \to b \in \mathbb{R} \to \varphi_{41}$
\item $\varphi_{258} :\equiv (a \in \mathbb{R} \to b \in \mathbb{R} \to 0 < a \to \varphi_{184}) \to \varphi_{90} \to a \in \mathbb{R} \to b \in \mathbb{R} \to 0 < a \to \varphi_{180}$
\item $\varphi_{259} :\equiv (a \in \mathbb{R} \to b \in \mathbb{R} \to 0 < a \to 0 < b \to \varphi_{131}) \to \varphi_{173} \to a \in \mathbb{R} \to b \in \mathbb{R} \to 0 < a \to \varphi_{165}$
\item $\varphi_{260} :\equiv (a \in \mathbb{R} \to b \in \mathbb{R} \to 0 < a \to 0 < b \to \varphi_{53}) \to (a \in \mathbb{R} \to b \in \mathbb{R} \to 0 < a \to 0 < b \to \varphi_{36}) \to \varphi_{219}$
\item $\varphi_{261} :\equiv (a \in \mathbb{R} \to b \in \mathbb{R} \to 0 < a \to \varphi_{185}) \to \varphi_{67} \to a \in \mathbb{R} \to b \in \mathbb{R} \to 0 < a \to 0 < b \to \varphi_{199}$
\item $\varphi_{262} :\equiv (\underline{n} \in \mathbb{R} \to a \in \mathbb{R} \to \underline{n} \cdot a \in \mathbb{R}) \to (a \in \mathbb{R} \to b \in \mathbb{R} \to 0 < a \to 0 < b \to \varphi_{138}) \to \varphi_{233}$
\item $\varphi_{263} :\equiv (a \in \mathbb{R} \to b \in \mathbb{R} \to 0 < a \to 0 < b \to \varphi_{198}) \to \varphi_{128} \to a \in \mathbb{R} \to b \in \mathbb{R} \to 0 < a \to 0 < b \to \varphi_{194}$
\item $\varphi_{264} :\equiv (a \in \mathbb{R} \to b \in \mathbb{R} \to 0 < a \to \varphi_{180}) \to (a \in \mathbb{R} \to b \in \mathbb{R} \to 0 < a \to \varphi_{175}) \to \varphi_{200}$
\item $\varphi_{265} :\equiv (s - a \in \mathbb{R} \to s \in \mathbb{R} \to 0 < s - (s - a) \to s - a < s) \to \varphi_{260}$
\item $\varphi_{266} :\equiv (a \in \mathbb{R} \to b \in \mathbb{R} \to 0 < a \to 0 < b \to \varphi_{103}) \to \varphi_{18} \to (\varphi_{18} \to \varphi_{27}) \to \varphi_{32} \to a \in \mathbb{R} \to b \in \mathbb{R} \to 0 < a \to 0 < b \to \varphi_{104}$
\item $\varphi_{267} :\equiv (a \in \mathbb{R} \to b \in \mathbb{R} \to 0 < a \to \varphi_{178}) \to \varphi_{204} \to a \in \mathbb{R} \to b \in \mathbb{R} \to 0 < a \to \varphi_{186}$
\item $\varphi_{268} :\equiv (a \in \mathbb{R} \to b \in \mathbb{R} \to 0 < a \to 0 < b \to \varphi_{127}) \to (a \in \mathbb{R} \to b \in \mathbb{R} \to 0 < a \to 0 < b \to \varphi_{144}) \to \varphi_{234}$
\item $\varphi_{269} :\equiv (a \in \mathbb{R} \to b \in \mathbb{R} \to 0 < a \to \varphi_{170}) \to (a \in \mathbb{R} \to b \in \mathbb{R} \to 0 < a \to \varphi_{165}) \to \varphi_{221}$
\item $\varphi_{270} :\equiv \forall y\,(\varphi_{80} \to y \in X_{a} \to y \leq b) \to \varphi_{211} \to (\varphi_{80} \to \forall y\,(\varphi_{80})) \to a \in \mathbb{R} \to b \in \mathbb{R} \to \varphi_{81}$
\item $\varphi_{271} :\equiv (m \in \mathbb{N} \to a \in \mathbb{R} \to \underline{S(m)} \cdot a = \underline{m} \cdot a + a) \to \varphi_{203} \to \varphi_{204} \to a \in \mathbb{R} \to b \in \mathbb{R} \to 0 < a \to \varphi_{192}$
\item $\varphi_{272} :\equiv \forall s\,(\varphi_{80} \land \neg 0 < 0 \to \neg \mathrm{IsSup}(s,X_{a})) \to \varphi_{216} \to (\varphi_{80} \land \neg 0 < 0 \to \forall s\,(\varphi_{80} \land \neg 0 < 0)) \to a \in \mathbb{R} \to \varphi_{66}$
\item $\varphi_{273} :\equiv \forall u\,(\varphi_{80} \to u \in X_{a} \to u \in \mathbb{R}) \to \varphi_{215} \to (\varphi_{80} \to \forall u\,(\varphi_{80})) \to a \in \mathbb{R} \to b \in \mathbb{R} \to \varphi_{102}$
\item $\varphi_{274} :\equiv (s - a \in \mathbb{R} \to y \in \mathbb{R} \to \neg y \leq s - a \to s - a < y) \to \varphi_{181} \to (a \in \mathbb{R} \to b \in \mathbb{R} \to 0 < a \to 0 < b \to \varphi_{131}) \to \varphi_{172} \to \varphi_{173}$
\item $\varphi_{275} :\equiv (b \in \mathbb{R} \to \underline{n} \cdot a \in \mathbb{R} \to \neg b < \underline{n} \cdot a \to \underline{n} \cdot a \leq b) \to \varphi_{268}$
\item $\varphi_{276} :\equiv \forall n\,(\varphi_{80} \land y \in X_{a} \land \neg y \leq b \to \neg (n \in \mathbb{N} \land y = \underline{n} \cdot a)) \to \varphi_{229} \to \varphi_{240}$
\item $\varphi_{277} :\equiv (s \in \mathbb{R} \to a \in \mathbb{R} \to y \in \mathbb{R} \to s - a < y \to s < y + a) \to (a \in \mathbb{R} \to b \in \mathbb{R} \to 0 < a \to 0 < b \to \varphi_{36}) \to \varphi_{204} \to \varphi_{259}$
\item $\varphi_{278} :\equiv (a \in \mathbb{R} \to b \in \mathbb{R} \to 0 < a \to 0 < b \to \varphi_{104}) \to \forall y\,(\varphi_{169}) \to (\forall y\,(\varphi_{169}) \to \varphi_{179}) \to \varphi_{213} \to \varphi_{252}$
\item $\varphi_{279} :\equiv \forall m\,(\varphi_{168} \to \neg (m \in \mathbb{N} \land y = \underline{m} \cdot a)) \to \varphi_{244} \to (\varphi_{168} \to \forall m\,(\varphi_{168})) \to \forall n\,(\varphi_{13}) \to \varphi_{243}$
\end{abbrevlist}
\end{fullproof}

\chapter{テキスト第3章 数列の極限 の証明木}
本章の定理数 84, ノード数の合計 24131.

\section{第3章: 数列の極限}
% 自動生成 (prover/build.js): ConvSeq  ノード数 11
\begin{fullproof}{[C1]}{$\forall a\,\forall \alpha\,(\mathrm{Conv}(a,\alpha) \to \mathrm{Seq}(a))$}{11}\phantomsection\label{tree:ConvSeq}
\noindent$\Pi_{1}$:
\begin{ptree}
\AXC{$\forall a\,\forall \alpha\,(\varphi_{1})$\lbl{DefConv}}
\AXC{$\forall a\,\forall \alpha\,(\varphi_{1}) \to \forall \alpha\,(\varphi_{1})$\lbl{A2}}
\BIC{$\forall \alpha\,(\varphi_{1})$}
\AXC{$\forall \alpha\,(\varphi_{1}) \to \varphi_{1}$\lbl{A2}}
\BIC{$\varphi_{1}$}
\end{ptree}
\noindent 主木:
\begin{ptree}
\AXC{$\stackrel{\Pi_{1}}{\vdots}$}\noLine
\UIC{$\varphi_{1}$}
\AXC{$\mathrm{Conv}(a,\alpha) \to \mathrm{Conv}(a,\alpha)$\lbl{A1}}
\AXC{$\varphi_{2}$\lbl{A1}}
\BIC{$\varphi_{1} \to \mathrm{Conv}(a,\alpha) \to \mathrm{Seq}(a)$}
\BIC{$\mathrm{Conv}(a,\alpha) \to \mathrm{Seq}(a)$}
\UIC{$\forall \alpha\,(\mathrm{Conv}(a,\alpha) \to \mathrm{Seq}(a))$}
\UIC{$\forall a\,\forall \alpha\,(\mathrm{Conv}(a,\alpha) \to \mathrm{Seq}(a))$}
\end{ptree}
\begin{abbrevlist}
\item $\varphi_{1} :\equiv \mathrm{Conv}(a,\alpha) \leftrightarrow \mathrm{Seq}(a) \land \alpha \in \mathbb{R} \land \forall \varepsilon \in \mathbb{R}\,(0 < \varepsilon \to \exists N \in \mathbb{N}\,\forall n \in \mathbb{N}\,(N \leq n \to \lvert a_{n} - \alpha \rvert < \varepsilon))$
\item $\varphi_{2} :\equiv (\mathrm{Conv}(a,\alpha) \to \mathrm{Conv}(a,\alpha)) \to \varphi_{1} \to \mathrm{Conv}(a,\alpha) \to \mathrm{Seq}(a)$
\end{abbrevlist}
\end{fullproof}

% 自動生成 (prover/build.js): ConvR  ノード数 11
\begin{fullproof}{[C2]}{$\forall a\,\forall \alpha\,(\mathrm{Conv}(a,\alpha) \to \alpha \in \mathbb{R})$}{11}\phantomsection\label{tree:ConvR}
\noindent$\Pi_{1}$:
\begin{ptree}
\AXC{$\forall a\,\forall \alpha\,(\varphi_{1})$\lbl{DefConv}}
\AXC{$\forall a\,\forall \alpha\,(\varphi_{1}) \to \forall \alpha\,(\varphi_{1})$\lbl{A2}}
\BIC{$\forall \alpha\,(\varphi_{1})$}
\AXC{$\forall \alpha\,(\varphi_{1}) \to \varphi_{1}$\lbl{A2}}
\BIC{$\varphi_{1}$}
\end{ptree}
\noindent 主木:
\begin{ptree}
\AXC{$\stackrel{\Pi_{1}}{\vdots}$}\noLine
\UIC{$\varphi_{1}$}
\AXC{$\mathrm{Conv}(a,\alpha) \to \mathrm{Conv}(a,\alpha)$\lbl{A1}}
\AXC{$\varphi_{2}$\lbl{A1}}
\BIC{$\varphi_{1} \to \mathrm{Conv}(a,\alpha) \to \alpha \in \mathbb{R}$}
\BIC{$\mathrm{Conv}(a,\alpha) \to \alpha \in \mathbb{R}$}
\UIC{$\forall \alpha\,(\mathrm{Conv}(a,\alpha) \to \alpha \in \mathbb{R})$}
\UIC{$\forall a\,\forall \alpha\,(\mathrm{Conv}(a,\alpha) \to \alpha \in \mathbb{R})$}
\end{ptree}
\begin{abbrevlist}
\item $\varphi_{1} :\equiv \mathrm{Conv}(a,\alpha) \leftrightarrow \mathrm{Seq}(a) \land \alpha \in \mathbb{R} \land \forall \varepsilon \in \mathbb{R}\,(0 < \varepsilon \to \exists N \in \mathbb{N}\,\forall n \in \mathbb{N}\,(N \leq n \to \lvert a_{n} - \alpha \rvert < \varepsilon))$
\item $\varphi_{2} :\equiv (\mathrm{Conv}(a,\alpha) \to \mathrm{Conv}(a,\alpha)) \to \varphi_{1} \to \mathrm{Conv}(a,\alpha) \to \alpha \in \mathbb{R}$
\end{abbrevlist}
\end{fullproof}

% 自動生成 (prover/build.js): ConvE  ノード数 22
\begin{fullproof}{[C3]}{$\forall a\,\forall \alpha\,\forall \varepsilon\,(\mathrm{Conv}(a,\alpha) \to \varepsilon \in \mathbb{R} \to 0 < \varepsilon \to \exists N \in \mathbb{N}\,\forall n \in \mathbb{N}\,(N \leq n \to \lvert a_{n} - \alpha \rvert < \varepsilon))$}{22}\phantomsection\label{tree:ConvE}
\noindent$\Pi_{1}$:
\begin{ptree}
\AXC{$\mathrm{Conv}(a,\alpha) \to \mathrm{Conv}(a,\alpha)$\lbl{A1}}
\AXC{$(\mathrm{Conv}(a,\alpha) \to \mathrm{Conv}(a,\alpha)) \to \varphi_{4}$\lbl{A1}}
\BIC{$\varphi_{4}$}
\end{ptree}
\noindent$\Pi_{2}$:
\begin{ptree}
\AXC{$\varphi_{2} \to \varphi_{1}$\lbl{A2}}
\AXC{$\forall a\,\forall \alpha\,(\mathrm{Conv}(a,\alpha) \leftrightarrow \mathrm{Seq}(a) \land \alpha \in \mathbb{R} \land \varphi_{2})$\lbl{DefConv}}
\AXC{$\varphi_{6}$\lbl{A2}}
\BIC{$\forall \alpha\,(\mathrm{Conv}(a,\alpha) \leftrightarrow \mathrm{Seq}(a) \land \alpha \in \mathbb{R} \land \varphi_{2})$}
\AXC{$\varphi_{5}$\lbl{A2}}
\BIC{$\mathrm{Conv}(a,\alpha) \leftrightarrow \mathrm{Seq}(a) \land \alpha \in \mathbb{R} \land \varphi_{2}$}
\AXC{$\stackrel{\Pi_{1}}{\vdots}$}\noLine
\UIC{$\varphi_{4}$}
\BIC{$\mathrm{Conv}(a,\alpha) \to \varepsilon \in \mathbb{R} \to 0 < \varepsilon \to \varphi_{2}$}
\AXC{$\varphi_{8}$\lbl{A1}}
\BIC{$(\varphi_{2} \to \varphi_{1}) \to \mathrm{Conv}(a,\alpha) \to \varepsilon \in \mathbb{R} \to 0 < \varepsilon \to \varphi_{1}$}
\BIC{$\mathrm{Conv}(a,\alpha) \to \varepsilon \in \mathbb{R} \to 0 < \varepsilon \to \varphi_{1}$}
\end{ptree}
\noindent$\Pi_{3}$:
\begin{ptree}
\AXC{$\mathrm{Conv}(a,\alpha) \to \varepsilon \in \mathbb{R} \to \varepsilon \in \mathbb{R}$\lbl{A1}}
\AXC{$\stackrel{\Pi_{2}}{\vdots}$}\noLine
\UIC{$\mathrm{Conv}(a,\alpha) \to \varepsilon \in \mathbb{R} \to 0 < \varepsilon \to \varphi_{1}$}
\AXC{$\varphi_{7}$\lbl{A1}}
\BIC{$(\mathrm{Conv}(a,\alpha) \to \varepsilon \in \mathbb{R} \to \varepsilon \in \mathbb{R}) \to \varphi_{3}$}
\BIC{$\varphi_{3}$}
\end{ptree}
\noindent 主木:
\begin{ptree}
\AXC{$\mathrm{Conv}(a,\alpha) \to \varepsilon \in \mathbb{R} \to 0 < \varepsilon \to 0 < \varepsilon$\lbl{A1}}
\AXC{$\stackrel{\Pi_{3}}{\vdots}$}\noLine
\UIC{$\varphi_{3}$}
\BIC{$\mathrm{Conv}(a,\alpha) \to \varphi_{1}$}
\UIC{$\forall \varepsilon\,(\mathrm{Conv}(a,\alpha) \to \varphi_{1})$}
\UIC{$\forall \alpha\,\forall \varepsilon\,(\mathrm{Conv}(a,\alpha) \to \varphi_{1})$}
\UIC{$\forall a\,\forall \alpha\,\forall \varepsilon\,(\mathrm{Conv}(a,\alpha) \to \varphi_{1})$}
\end{ptree}
\begin{abbrevlist}
\item $\varphi_{1} :\equiv \varepsilon \in \mathbb{R} \to 0 < \varepsilon \to \exists N \in \mathbb{N}\,\forall n \in \mathbb{N}\,(N \leq n \to \lvert a_{n} - \alpha \rvert < \varepsilon)$
\item $\varphi_{2} :\equiv \forall \varepsilon \in \mathbb{R}\,(0 < \varepsilon \to \exists N \in \mathbb{N}\,\forall n \in \mathbb{N}\,(N \leq n \to \lvert a_{n} - \alpha \rvert < \varepsilon))$
\item $\varphi_{3} :\equiv (\mathrm{Conv}(a,\alpha) \to \varepsilon \in \mathbb{R} \to 0 < \varepsilon \to 0 < \varepsilon) \to \mathrm{Conv}(a,\alpha) \to \varphi_{1}$
\item $\varphi_{4} :\equiv (\mathrm{Conv}(a,\alpha) \leftrightarrow \mathrm{Seq}(a) \land \alpha \in \mathbb{R} \land \varphi_{2}) \to \mathrm{Conv}(a,\alpha) \to \varepsilon \in \mathbb{R} \to 0 < \varepsilon \to \varphi_{2}$
\item $\varphi_{5} :\equiv \forall \alpha\,(\mathrm{Conv}(a,\alpha) \leftrightarrow \mathrm{Seq}(a) \land \alpha \in \mathbb{R} \land \varphi_{2}) \to (\mathrm{Conv}(a,\alpha) \leftrightarrow \mathrm{Seq}(a) \land \alpha \in \mathbb{R} \land \varphi_{2})$
\item $\varphi_{6} :\equiv \forall a\,\forall \alpha\,(\mathrm{Conv}(a,\alpha) \leftrightarrow \mathrm{Seq}(a) \land \alpha \in \mathbb{R} \land \varphi_{2}) \to \forall \alpha\,(\mathrm{Conv}(a,\alpha) \leftrightarrow \mathrm{Seq}(a) \land \alpha \in \mathbb{R} \land \varphi_{2})$
\item $\varphi_{7} :\equiv (\mathrm{Conv}(a,\alpha) \to \varepsilon \in \mathbb{R} \to 0 < \varepsilon \to \varphi_{1}) \to (\mathrm{Conv}(a,\alpha) \to \varepsilon \in \mathbb{R} \to \varepsilon \in \mathbb{R}) \to \varphi_{3}$
\item $\varphi_{8} :\equiv (\mathrm{Conv}(a,\alpha) \to \varepsilon \in \mathbb{R} \to 0 < \varepsilon \to \varphi_{2}) \to (\varphi_{2} \to \varphi_{1}) \to \mathrm{Conv}(a,\alpha) \to \varepsilon \in \mathbb{R} \to 0 < \varepsilon \to \varphi_{1}$
\end{abbrevlist}
\end{fullproof}

% 自動生成 (prover/build.js): SubCupL  ノード数 20
\begin{fullproof}{[C4]}{$\forall m\,\forall n\,(m \subseteq m \cup n)$}{20}\phantomsection\label{tree:SubCupL}
\noindent$\Pi_{1}$:
\begin{ptree}
\AXC{$\forall x\,\forall y\,\forall u\,(u \in x \cup y \leftrightarrow u \in x \lor u \in y)$\lbl{Def$\cup$}}
\AXC{$\varphi_{5}$\lbl{A2}}
\BIC{$\forall y\,\forall u\,(u \in m \cup y \leftrightarrow u \in m \lor u \in y)$}
\AXC{$\varphi_{3}$\lbl{A2}}
\BIC{$\forall u\,(u \in m \cup n \leftrightarrow u \in m \lor u \in n)$}
\AXC{$\varphi_{2}$\lbl{A2}}
\BIC{$u \in m \cup n \leftrightarrow u \in m \lor u \in n$}
\end{ptree}
\noindent$\Pi_{2}$:
\begin{ptree}
\AXC{$\stackrel{\Pi_{1}}{\vdots}$}\noLine
\UIC{$u \in m \cup n \leftrightarrow u \in m \lor u \in n$}
\AXC{$(u \in m \cup n \leftrightarrow u \in m \lor u \in n) \to u \in m \to u \in m \cup n$\lbl{A1}}
\BIC{$u \in m \to u \in m \cup n$}
\UIC{$\forall u \in m\,(u \in m \cup n)$}
\AXC{$\varphi_{6}$\lbl{A1}}
\BIC{$(m \subseteq m \cup n \leftrightarrow \forall u \in m\,(u \in m \cup n)) \to m \subseteq m \cup n$}
\end{ptree}
\noindent 主木:
\begin{ptree}
\AXC{$\forall x\,\forall y\,(x \subseteq y \leftrightarrow \forall u \in x\,(u \in y))$\lbl{Def$\subseteq$}}
\AXC{$\varphi_{1}$\lbl{A2}}
\BIC{$\forall y\,(m \subseteq y \leftrightarrow \forall u \in m\,(u \in y))$}
\AXC{$\varphi_{4}$\lbl{A2}}
\BIC{$m \subseteq m \cup n \leftrightarrow \forall u \in m\,(u \in m \cup n)$}
\AXC{$\stackrel{\Pi_{2}}{\vdots}$}\noLine
\UIC{$(m \subseteq m \cup n \leftrightarrow \forall u \in m\,(u \in m \cup n)) \to m \subseteq m \cup n$}
\BIC{$m \subseteq m \cup n$}
\UIC{$\forall n\,(m \subseteq m \cup n)$}
\UIC{$\forall m\,\forall n\,(m \subseteq m \cup n)$}
\end{ptree}
\begin{abbrevlist}
\item $\varphi_{1} :\equiv \forall x\,\forall y\,(x \subseteq y \leftrightarrow \forall u \in x\,(u \in y)) \to \forall y\,(m \subseteq y \leftrightarrow \forall u \in m\,(u \in y))$
\item $\varphi_{2} :\equiv \forall u\,(u \in m \cup n \leftrightarrow u \in m \lor u \in n) \to (u \in m \cup n \leftrightarrow u \in m \lor u \in n)$
\item $\varphi_{3} :\equiv \forall y\,\forall u\,(u \in m \cup y \leftrightarrow u \in m \lor u \in y) \to \forall u\,(u \in m \cup n \leftrightarrow u \in m \lor u \in n)$
\item $\varphi_{4} :\equiv \forall y\,(m \subseteq y \leftrightarrow \forall u \in m\,(u \in y)) \to (m \subseteq m \cup n \leftrightarrow \forall u \in m\,(u \in m \cup n))$
\item $\varphi_{5} :\equiv \forall x\,\forall y\,\forall u\,(u \in x \cup y \leftrightarrow u \in x \lor u \in y) \to \forall y\,\forall u\,(u \in m \cup y \leftrightarrow u \in m \lor u \in y)$
\item $\varphi_{6} :\equiv \forall u \in m\,(u \in m \cup n) \to (m \subseteq m \cup n \leftrightarrow \forall u \in m\,(u \in m \cup n)) \to m \subseteq m \cup n$
\end{abbrevlist}
\end{fullproof}

% 自動生成 (prover/build.js): SubCupR  ノード数 20
\begin{fullproof}{[C5]}{$\forall m\,\forall n\,(n \subseteq m \cup n)$}{20}\phantomsection\label{tree:SubCupR}
\noindent$\Pi_{1}$:
\begin{ptree}
\AXC{$\forall x\,\forall y\,\forall u\,(u \in x \cup y \leftrightarrow u \in x \lor u \in y)$\lbl{Def$\cup$}}
\AXC{$\varphi_{5}$\lbl{A2}}
\BIC{$\forall y\,\forall u\,(u \in m \cup y \leftrightarrow u \in m \lor u \in y)$}
\AXC{$\varphi_{3}$\lbl{A2}}
\BIC{$\forall u\,(u \in m \cup n \leftrightarrow u \in m \lor u \in n)$}
\AXC{$\varphi_{2}$\lbl{A2}}
\BIC{$u \in m \cup n \leftrightarrow u \in m \lor u \in n$}
\end{ptree}
\noindent$\Pi_{2}$:
\begin{ptree}
\AXC{$\stackrel{\Pi_{1}}{\vdots}$}\noLine
\UIC{$u \in m \cup n \leftrightarrow u \in m \lor u \in n$}
\AXC{$(u \in m \cup n \leftrightarrow u \in m \lor u \in n) \to u \in n \to u \in m \cup n$\lbl{A1}}
\BIC{$u \in n \to u \in m \cup n$}
\UIC{$\forall u \in n\,(u \in m \cup n)$}
\AXC{$\varphi_{6}$\lbl{A1}}
\BIC{$(n \subseteq m \cup n \leftrightarrow \forall u \in n\,(u \in m \cup n)) \to n \subseteq m \cup n$}
\end{ptree}
\noindent 主木:
\begin{ptree}
\AXC{$\forall x\,\forall y\,(x \subseteq y \leftrightarrow \forall u \in x\,(u \in y))$\lbl{Def$\subseteq$}}
\AXC{$\varphi_{1}$\lbl{A2}}
\BIC{$\forall y\,(n \subseteq y \leftrightarrow \forall u \in n\,(u \in y))$}
\AXC{$\varphi_{4}$\lbl{A2}}
\BIC{$n \subseteq m \cup n \leftrightarrow \forall u \in n\,(u \in m \cup n)$}
\AXC{$\stackrel{\Pi_{2}}{\vdots}$}\noLine
\UIC{$(n \subseteq m \cup n \leftrightarrow \forall u \in n\,(u \in m \cup n)) \to n \subseteq m \cup n$}
\BIC{$n \subseteq m \cup n$}
\UIC{$\forall n\,(n \subseteq m \cup n)$}
\UIC{$\forall m\,\forall n\,(n \subseteq m \cup n)$}
\end{ptree}
\begin{abbrevlist}
\item $\varphi_{1} :\equiv \forall x\,\forall y\,(x \subseteq y \leftrightarrow \forall u \in x\,(u \in y)) \to \forall y\,(n \subseteq y \leftrightarrow \forall u \in n\,(u \in y))$
\item $\varphi_{2} :\equiv \forall u\,(u \in m \cup n \leftrightarrow u \in m \lor u \in n) \to (u \in m \cup n \leftrightarrow u \in m \lor u \in n)$
\item $\varphi_{3} :\equiv \forall y\,\forall u\,(u \in m \cup y \leftrightarrow u \in m \lor u \in y) \to \forall u\,(u \in m \cup n \leftrightarrow u \in m \lor u \in n)$
\item $\varphi_{4} :\equiv \forall y\,(n \subseteq y \leftrightarrow \forall u \in n\,(u \in y)) \to (n \subseteq m \cup n \leftrightarrow \forall u \in n\,(u \in m \cup n))$
\item $\varphi_{5} :\equiv \forall x\,\forall y\,\forall u\,(u \in x \cup y \leftrightarrow u \in x \lor u \in y) \to \forall y\,\forall u\,(u \in m \cup y \leftrightarrow u \in m \lor u \in y)$
\item $\varphi_{6} :\equiv \forall u \in n\,(u \in m \cup n) \to (n \subseteq m \cup n \leftrightarrow \forall u \in n\,(u \in m \cup n)) \to n \subseteq m \cup n$
\end{abbrevlist}
\end{fullproof}

% 自動生成 (prover/build.js): SubLe  ノード数 11
\begin{fullproof}{[C6]}{$\forall x\,\forall y\,(x \subseteq y \to x \leq y)$}{11}\phantomsection\label{tree:SubLe}
\noindent$\Pi_{1}$:
\begin{ptree}
\AXC{$\forall x\,\forall y\,(x < y \leftrightarrow x \subseteq y \land x \neq y)$\lbl{Def$<$}}
\AXC{$\varphi_{1}$\lbl{A2}}
\BIC{$\forall y\,(x < y \leftrightarrow x \subseteq y \land x \neq y)$}
\AXC{$\forall y\,(x < y \leftrightarrow x \subseteq y \land x \neq y) \to (x < y \leftrightarrow x \subseteq y \land x \neq y)$\lbl{A2}}
\BIC{$x < y \leftrightarrow x \subseteq y \land x \neq y$}
\end{ptree}
\noindent 主木:
\begin{ptree}
\AXC{$\stackrel{\Pi_{1}}{\vdots}$}\noLine
\UIC{$x < y \leftrightarrow x \subseteq y \land x \neq y$}
\AXC{$x \subseteq y \to x \subseteq y$\lbl{A1}}
\AXC{$\varphi_{2}$\lbl{A1}}
\BIC{$(x < y \leftrightarrow x \subseteq y \land x \neq y) \to x \subseteq y \to x \leq y$}
\BIC{$x \subseteq y \to x \leq y$}
\UIC{$\forall y\,(x \subseteq y \to x \leq y)$}
\UIC{$\forall x\,\forall y\,(x \subseteq y \to x \leq y)$}
\end{ptree}
\begin{abbrevlist}
\item $\varphi_{1} :\equiv \forall x\,\forall y\,(x < y \leftrightarrow x \subseteq y \land x \neq y) \to \forall y\,(x < y \leftrightarrow x \subseteq y \land x \neq y)$
\item $\varphi_{2} :\equiv (x \subseteq y \to x \subseteq y) \to (x < y \leftrightarrow x \subseteq y \land x \neq y) \to x \subseteq y \to x \leq y$
\end{abbrevlist}
\end{fullproof}

% 自動生成 (prover/build.js): LeCupL  ノード数 15
\begin{fullproof}{[C7]}{$\forall m\,\forall n\,(m \leq m \cup n)$}{15}\phantomsection\label{tree:LeCupL}
\noindent$\Pi_{1}$:
\begin{ptree}
\AXC{$\forall m\,\forall n\,(m \subseteq m \cup n)$\lbl{[C4]}}
\AXC{$\forall m\,\forall n\,(m \subseteq m \cup n) \to \forall n\,(m \subseteq m \cup n)$\lbl{A2}}
\BIC{$\forall n\,(m \subseteq m \cup n)$}
\AXC{$\forall n\,(m \subseteq m \cup n) \to m \subseteq m \cup n$\lbl{A2}}
\BIC{$m \subseteq m \cup n$}
\end{ptree}
\noindent$\Pi_{2}$:
\begin{ptree}
\AXC{$\forall x\,\forall y\,(x \subseteq y \to x \leq y)$\lbl{[C6]}}
\AXC{$\forall x\,\forall y\,(x \subseteq y \to x \leq y) \to \forall y\,(m \subseteq y \to m \leq y)$\lbl{A2}}
\BIC{$\forall y\,(m \subseteq y \to m \leq y)$}
\end{ptree}
\noindent$\Pi_{3}$:
\begin{ptree}
\AXC{$\stackrel{\Pi_{2}}{\vdots}$}\noLine
\UIC{$\forall y\,(m \subseteq y \to m \leq y)$}
\AXC{$\forall y\,(m \subseteq y \to m \leq y) \to m \subseteq m \cup n \to m \leq m \cup n$\lbl{A2}}
\BIC{$m \subseteq m \cup n \to m \leq m \cup n$}
\AXC{$(m \subseteq m \cup n \to m \leq m \cup n) \to m \subseteq m \cup n \to m \leq m \cup n$\lbl{A1}}
\BIC{$m \subseteq m \cup n \to m \leq m \cup n$}
\end{ptree}
\noindent 主木:
\begin{ptree}
\AXC{$\stackrel{\Pi_{1}}{\vdots}$}\noLine
\UIC{$m \subseteq m \cup n$}
\AXC{$\stackrel{\Pi_{3}}{\vdots}$}\noLine
\UIC{$m \subseteq m \cup n \to m \leq m \cup n$}
\BIC{$m \leq m \cup n$}
\UIC{$\forall n\,(m \leq m \cup n)$}
\UIC{$\forall m\,\forall n\,(m \leq m \cup n)$}
\end{ptree}
\end{fullproof}

% 自動生成 (prover/build.js): LeCupR  ノード数 15
\begin{fullproof}{[C8]}{$\forall m\,\forall n\,(n \leq m \cup n)$}{15}\phantomsection\label{tree:LeCupR}
\noindent$\Pi_{1}$:
\begin{ptree}
\AXC{$\forall m\,\forall n\,(n \subseteq m \cup n)$\lbl{[C5]}}
\AXC{$\forall m\,\forall n\,(n \subseteq m \cup n) \to \forall n\,(n \subseteq m \cup n)$\lbl{A2}}
\BIC{$\forall n\,(n \subseteq m \cup n)$}
\AXC{$\forall n\,(n \subseteq m \cup n) \to n \subseteq m \cup n$\lbl{A2}}
\BIC{$n \subseteq m \cup n$}
\end{ptree}
\noindent$\Pi_{2}$:
\begin{ptree}
\AXC{$\forall x\,\forall y\,(x \subseteq y \to x \leq y)$\lbl{[C6]}}
\AXC{$\forall x\,\forall y\,(x \subseteq y \to x \leq y) \to \forall y\,(n \subseteq y \to n \leq y)$\lbl{A2}}
\BIC{$\forall y\,(n \subseteq y \to n \leq y)$}
\end{ptree}
\noindent$\Pi_{3}$:
\begin{ptree}
\AXC{$\stackrel{\Pi_{2}}{\vdots}$}\noLine
\UIC{$\forall y\,(n \subseteq y \to n \leq y)$}
\AXC{$\forall y\,(n \subseteq y \to n \leq y) \to n \subseteq m \cup n \to n \leq m \cup n$\lbl{A2}}
\BIC{$n \subseteq m \cup n \to n \leq m \cup n$}
\AXC{$(n \subseteq m \cup n \to n \leq m \cup n) \to n \subseteq m \cup n \to n \leq m \cup n$\lbl{A1}}
\BIC{$n \subseteq m \cup n \to n \leq m \cup n$}
\end{ptree}
\noindent 主木:
\begin{ptree}
\AXC{$\stackrel{\Pi_{1}}{\vdots}$}\noLine
\UIC{$n \subseteq m \cup n$}
\AXC{$\stackrel{\Pi_{3}}{\vdots}$}\noLine
\UIC{$n \subseteq m \cup n \to n \leq m \cup n$}
\BIC{$n \leq m \cup n$}
\UIC{$\forall n\,(n \leq m \cup n)$}
\UIC{$\forall m\,\forall n\,(n \leq m \cup n)$}
\end{ptree}
\end{fullproof}

% 自動生成 (prover/build.js): Tail  ノード数 157
\begin{fullproof}{[C9]}{$\forall a\,\forall \alpha\,\forall b\,\forall \beta\,\forall \varepsilon\,(\mathrm{Conv}(a,\alpha) \to \mathrm{Conv}(b,\beta) \to \varepsilon \in \mathbb{R} \to 0 < \varepsilon \to \exists N \in \mathbb{N}\,\forall n \in \mathbb{N}\,(N \leq n \to \lvert a_{n} - \alpha \rvert < \varepsilon \land \lvert b_{n} - \beta \rvert < \varepsilon))$}{157}\phantomsection\label{tree:Tail}
\noindent$\Pi_{1}$:
\begin{ptree}
\AXC{$\forall m\,\forall n\,(n \leq m \cup n)$\lbl{[C8]}}
\AXC{$\forall m\,\forall n\,(n \leq m \cup n) \to \forall n\,(n \leq N_{1} \cup n)$\lbl{A2}}
\BIC{$\forall n\,(n \leq N_{1} \cup n)$}
\AXC{$\forall n\,(n \leq N_{1} \cup n) \to N_{2} \leq N_{1} \cup N_{2}$\lbl{A2}}
\BIC{$N_{2} \leq N_{1} \cup N_{2}$}
\end{ptree}
\noindent$\Pi_{2}$:
\begin{ptree}
\AXC{$\stackrel{\Pi_{1}}{\vdots}$}\noLine
\UIC{$N_{2} \leq N_{1} \cup N_{2}$}
\AXC{$\forall x\,\forall y\,\forall z\,(x \leq y \to y \leq z \to x \leq z)$\lbl{[O5]}}
\AXC{$\varphi_{1}$\lbl{A2}}
\BIC{$\forall y\,\forall z\,(N_{2} \leq y \to y \leq z \to N_{2} \leq z)$}
\AXC{$\varphi_{6}$\lbl{A2}}
\BIC{$\forall z\,(N_{2} \leq N_{1} \cup N_{2} \to N_{1} \cup N_{2} \leq z \to N_{2} \leq z)$}
\AXC{$\varphi_{16}$\lbl{A2}}
\BIC{$N_{2} \leq N_{1} \cup N_{2} \to N_{1} \cup N_{2} \leq n \to N_{2} \leq n$}
\AXC{$\varphi_{88}$\lbl{A1}}
\BIC{$N_{2} \leq N_{1} \cup N_{2} \to \varphi_{81}$}
\BIC{$\varphi_{81}$}
\end{ptree}
\noindent$\Pi_{3}$:
\begin{ptree}
\AXC{$\varphi_{29}$\lbl{A2}}
\AXC{$\mathrm{Conv}(a,\alpha) \to \mathrm{Conv}(b,\beta) \to \varepsilon \in \mathbb{R} \to \varphi_{55}$\lbl{A1}}
\AXC{$\varphi_{91}$\lbl{A1}}
\BIC{$\mathrm{Conv}(a,\alpha) \to \mathrm{Conv}(b,\beta) \to \varepsilon \in \mathbb{R} \to 0 < \varepsilon \to \varphi_{58}$}
\AXC{$\varphi_{94}$\lbl{A1}}
\BIC{$\varphi_{76}$}
\BIC{$\mathrm{Conv}(a,\alpha) \to \mathrm{Conv}(b,\beta) \to \varepsilon \in \mathbb{R} \to 0 < \varepsilon \to \varphi_{56}$}
\AXC{$\varphi_{98}$\lbl{A1}}
\BIC{$(\mathrm{Conv}(a,\alpha) \to \mathrm{Conv}(b,\beta) \to \varepsilon \in \mathbb{R} \to \varphi_{46}) \to \varphi_{85}$}
\end{ptree}
\noindent$\Pi_{4}$:
\begin{ptree}
\AXC{$\mathrm{Conv}(a,\alpha) \to \mathrm{Conv}(b,\beta) \to \varepsilon \in \mathbb{R} \to \varphi_{46}$\lbl{A1}}
\AXC{$\stackrel{\Pi_{3}}{\vdots}$}\noLine
\UIC{$(\mathrm{Conv}(a,\alpha) \to \mathrm{Conv}(b,\beta) \to \varepsilon \in \mathbb{R} \to \varphi_{46}) \to \varphi_{85}$}
\BIC{$\varphi_{85}$}
\end{ptree}
\noindent$\Pi_{5}$:
\begin{ptree}
\AXC{$\forall m\,\forall n\,(m \leq m \cup n)$\lbl{[C7]}}
\AXC{$\forall m\,\forall n\,(m \leq m \cup n) \to \forall n\,(N_{1} \leq N_{1} \cup n)$\lbl{A2}}
\BIC{$\forall n\,(N_{1} \leq N_{1} \cup n)$}
\AXC{$\forall n\,(N_{1} \leq N_{1} \cup n) \to N_{1} \leq N_{1} \cup N_{2}$\lbl{A2}}
\BIC{$N_{1} \leq N_{1} \cup N_{2}$}
\end{ptree}
\noindent$\Pi_{6}$:
\begin{ptree}
\AXC{$\stackrel{\Pi_{5}}{\vdots}$}\noLine
\UIC{$N_{1} \leq N_{1} \cup N_{2}$}
\AXC{$\forall x\,\forall y\,\forall z\,(x \leq y \to y \leq z \to x \leq z)$\lbl{[O5]}}
\AXC{$\varphi_{2}$\lbl{A2}}
\BIC{$\forall y\,\forall z\,(N_{1} \leq y \to y \leq z \to N_{1} \leq z)$}
\AXC{$\varphi_{7}$\lbl{A2}}
\BIC{$\forall z\,(N_{1} \leq N_{1} \cup N_{2} \to N_{1} \cup N_{2} \leq z \to N_{1} \leq z)$}
\AXC{$\varphi_{17}$\lbl{A2}}
\BIC{$N_{1} \leq N_{1} \cup N_{2} \to N_{1} \cup N_{2} \leq n \to N_{1} \leq n$}
\AXC{$\varphi_{89}$\lbl{A1}}
\BIC{$N_{1} \leq N_{1} \cup N_{2} \to \varphi_{82}$}
\BIC{$\varphi_{82}$}
\end{ptree}
\noindent$\Pi_{7}$:
\begin{ptree}
\AXC{$\varphi_{30}$\lbl{A2}}
\AXC{$\mathrm{Conv}(a,\alpha) \to \mathrm{Conv}(b,\beta) \to \varepsilon \in \mathbb{R} \to 0 < \varepsilon \to \varphi_{36}$\lbl{A1}}
\AXC{$\varphi_{84}$\lbl{A1}}
\BIC{$\mathrm{Conv}(a,\alpha) \to \mathrm{Conv}(b,\beta) \to \varepsilon \in \mathbb{R} \to 0 < \varepsilon \to \varphi_{59}$}
\AXC{$\varphi_{95}$\lbl{A1}}
\BIC{$\varphi_{77}$}
\BIC{$\mathrm{Conv}(a,\alpha) \to \mathrm{Conv}(b,\beta) \to \varepsilon \in \mathbb{R} \to 0 < \varepsilon \to \varphi_{57}$}
\AXC{$\varphi_{99}$\lbl{A1}}
\BIC{$(\mathrm{Conv}(a,\alpha) \to \mathrm{Conv}(b,\beta) \to \varepsilon \in \mathbb{R} \to \varphi_{46}) \to \varphi_{86}$}
\end{ptree}
\noindent$\Pi_{8}$:
\begin{ptree}
\AXC{$\mathrm{Conv}(a,\alpha) \to \mathrm{Conv}(b,\beta) \to \varepsilon \in \mathbb{R} \to \varphi_{46}$\lbl{A1}}
\AXC{$\stackrel{\Pi_{7}}{\vdots}$}\noLine
\UIC{$(\mathrm{Conv}(a,\alpha) \to \mathrm{Conv}(b,\beta) \to \varepsilon \in \mathbb{R} \to \varphi_{46}) \to \varphi_{86}$}
\BIC{$\varphi_{86}$}
\end{ptree}
\noindent$\Pi_{9}$:
\begin{ptree}
\AXC{$\mathrm{Conv}(a,\alpha) \to \mathrm{Conv}(b,\beta) \to \varepsilon \in \mathbb{R} \to \varphi_{52}$\lbl{A1}}
\AXC{$\stackrel{\Pi_{6}}{\vdots}$}\noLine
\UIC{$\varphi_{82}$}
\BIC{$\mathrm{Conv}(a,\alpha) \to \mathrm{Conv}(b,\beta) \to \varepsilon \in \mathbb{R} \to \varphi_{49}$}
\AXC{$\stackrel{\Pi_{8}}{\vdots}$}\noLine
\UIC{$\varphi_{86}$}
\BIC{$\mathrm{Conv}(a,\alpha) \to \mathrm{Conv}(b,\beta) \to \varepsilon \in \mathbb{R} \to \varphi_{54}$}
\end{ptree}
\noindent$\Pi_{10}$:
\begin{ptree}
\AXC{$\stackrel{\Pi_{9}}{\vdots}$}\noLine
\UIC{$\mathrm{Conv}(a,\alpha) \to \mathrm{Conv}(b,\beta) \to \varepsilon \in \mathbb{R} \to \varphi_{54}$}
\AXC{$(\mathrm{Conv}(a,\alpha) \to \mathrm{Conv}(b,\beta) \to \varepsilon \in \mathbb{R} \to \varphi_{54}) \to \varphi_{90}$\lbl{A1}}
\BIC{$\varphi_{90}$}
\end{ptree}
\noindent$\Pi_{11}$:
\begin{ptree}
\AXC{$\mathrm{Conv}(a,\alpha) \to \mathrm{Conv}(b,\beta) \to \varepsilon \in \mathbb{R} \to \varphi_{52}$\lbl{A1}}
\AXC{$\stackrel{\Pi_{2}}{\vdots}$}\noLine
\UIC{$\varphi_{81}$}
\BIC{$\mathrm{Conv}(a,\alpha) \to \mathrm{Conv}(b,\beta) \to \varepsilon \in \mathbb{R} \to \varphi_{48}$}
\AXC{$\stackrel{\Pi_{4}}{\vdots}$}\noLine
\UIC{$\varphi_{85}$}
\BIC{$\mathrm{Conv}(a,\alpha) \to \mathrm{Conv}(b,\beta) \to \varepsilon \in \mathbb{R} \to \varphi_{53}$}
\AXC{$\stackrel{\Pi_{10}}{\vdots}$}\noLine
\UIC{$\varphi_{90}$}
\BIC{$\mathrm{Conv}(a,\alpha) \to \mathrm{Conv}(b,\beta) \to \varepsilon \in \mathbb{R} \to \varphi_{65}$}
\end{ptree}
\noindent$\Pi_{12}$:
\begin{ptree}
\AXC{$\stackrel{\Pi_{11}}{\vdots}$}\noLine
\UIC{$\mathrm{Conv}(a,\alpha) \to \mathrm{Conv}(b,\beta) \to \varepsilon \in \mathbb{R} \to \varphi_{65}$}
\AXC{$(\mathrm{Conv}(a,\alpha) \to \mathrm{Conv}(b,\beta) \to \varepsilon \in \mathbb{R} \to \varphi_{65}) \to \varphi_{70}$\lbl{A1}}
\BIC{$\varphi_{70}$}
\UIC{$\forall n\,(\varphi_{70})$}
\end{ptree}
\noindent$\Pi_{13}$:
\begin{ptree}
\AXC{$\varphi_{51} \to \forall n\,(\varphi_{51})$\lbl{A4}}
\AXC{$\forall n\,(\varphi_{70}) \to \forall n\,(\varphi_{51}) \to \varphi_{19}$\lbl{A3}}
\AXC{$\stackrel{\Pi_{12}}{\vdots}$}\noLine
\UIC{$\forall n\,(\varphi_{70})$}
\AXC{$\forall n\,(\varphi_{70}) \to (\forall n\,(\varphi_{70}) \to \forall n\,(\varphi_{51}) \to \varphi_{19}) \to \varphi_{96}$\lbl{A1}}
\BIC{$(\forall n\,(\varphi_{70}) \to \forall n\,(\varphi_{51}) \to \varphi_{19}) \to \varphi_{96}$}
\BIC{$\varphi_{96}$}
\BIC{$\mathrm{Conv}(a,\alpha) \to \mathrm{Conv}(b,\beta) \to \varepsilon \in \mathbb{R} \to 0 < \varepsilon \to \varphi_{62}$}
\end{ptree}
\noindent$\Pi_{14}$:
\begin{ptree}
\AXC{$\forall m\,\forall n\,(m \in \mathbb{N} \to n \in \mathbb{N} \to m \cup n \in \mathbb{N})$\lbl{[N3]}}
\AXC{$\varphi_{26}$\lbl{A2}}
\BIC{$\forall n\,(N_{1} \in \mathbb{N} \to n \in \mathbb{N} \to N_{1} \cup n \in \mathbb{N})$}
\AXC{$\varphi_{24}$\lbl{A2}}
\BIC{$N_{1} \in \mathbb{N} \to N_{2} \in \mathbb{N} \to N_{1} \cup N_{2} \in \mathbb{N}$}
\AXC{$\varphi_{97}$\lbl{A1}}
\BIC{$(\mathrm{Conv}(a,\alpha) \to \mathrm{Conv}(b,\beta) \to \varepsilon \in \mathbb{R} \to \varphi_{44}) \to \varphi_{80}$}
\end{ptree}
\noindent$\Pi_{15}$:
\begin{ptree}
\AXC{$\mathrm{Conv}(a,\alpha) \to \mathrm{Conv}(b,\beta) \to \varepsilon \in \mathbb{R} \to 0 < \varepsilon \to \varphi_{36}$\lbl{A1}}
\AXC{$\varphi_{79}$\lbl{A1}}
\BIC{$\mathrm{Conv}(a,\alpha) \to \mathrm{Conv}(b,\beta) \to \varepsilon \in \mathbb{R} \to \varphi_{44}$}
\AXC{$\stackrel{\Pi_{14}}{\vdots}$}\noLine
\UIC{$(\mathrm{Conv}(a,\alpha) \to \mathrm{Conv}(b,\beta) \to \varepsilon \in \mathbb{R} \to \varphi_{44}) \to \varphi_{80}$}
\BIC{$\varphi_{80}$}
\end{ptree}
\noindent$\Pi_{16}$:
\begin{ptree}
\AXC{$\mathrm{Conv}(a,\alpha) \to \mathrm{Conv}(b,\beta) \to \varepsilon \in \mathbb{R} \to \varphi_{55}$\lbl{A1}}
\AXC{$\varphi_{83}$\lbl{A1}}
\BIC{$\mathrm{Conv}(a,\alpha) \to \mathrm{Conv}(b,\beta) \to \varepsilon \in \mathbb{R} \to \varphi_{43}$}
\AXC{$\stackrel{\Pi_{15}}{\vdots}$}\noLine
\UIC{$\varphi_{80}$}
\BIC{$\mathrm{Conv}(a,\alpha) \to \mathrm{Conv}(b,\beta) \to \varepsilon \in \mathbb{R} \to \varphi_{45}$}
\end{ptree}
\noindent$\Pi_{17}$:
\begin{ptree}
\AXC{$\stackrel{\Pi_{16}}{\vdots}$}\noLine
\UIC{$\mathrm{Conv}(a,\alpha) \to \mathrm{Conv}(b,\beta) \to \varepsilon \in \mathbb{R} \to \varphi_{45}$}
\AXC{$(\mathrm{Conv}(a,\alpha) \to \mathrm{Conv}(b,\beta) \to \varepsilon \in \mathbb{R} \to \varphi_{45}) \to \varphi_{93}$\lbl{A1}}
\BIC{$\varphi_{93}$}
\end{ptree}
\noindent$\Pi_{18}$:
\begin{ptree}
\AXC{$\stackrel{\Pi_{13}}{\vdots}$}\noLine
\UIC{$\mathrm{Conv}(a,\alpha) \to \mathrm{Conv}(b,\beta) \to \varepsilon \in \mathbb{R} \to 0 < \varepsilon \to \varphi_{62}$}
\AXC{$\stackrel{\Pi_{17}}{\vdots}$}\noLine
\UIC{$\varphi_{93}$}
\BIC{$\mathrm{Conv}(a,\alpha) \to \mathrm{Conv}(b,\beta) \to \varepsilon \in \mathbb{R} \to \varphi_{69}$}
\AXC{$(\mathrm{Conv}(a,\alpha) \to \mathrm{Conv}(b,\beta) \to \varepsilon \in \mathbb{R} \to \varphi_{69}) \to \varphi_{87}$\lbl{A1}}
\BIC{$\varphi_{87}$}
\end{ptree}
\noindent$\Pi_{19}$:
\begin{ptree}
\AXC{$\forall a\,\forall \alpha\,\forall \varepsilon\,(\mathrm{Conv}(a,\alpha) \to \varphi_{22})$\lbl{[C3]}}
\AXC{$\forall a\,\forall \alpha\,\forall \varepsilon\,(\mathrm{Conv}(a,\alpha) \to \varphi_{22}) \to \varphi_{32}$\lbl{A2}}
\BIC{$\varphi_{32}$}
\AXC{$\varphi_{32} \to \forall \varepsilon\,(\mathrm{Conv}(b,\beta) \to \varphi_{21})$\lbl{A2}}
\BIC{$\forall \varepsilon\,(\mathrm{Conv}(b,\beta) \to \varphi_{21})$}
\end{ptree}
\noindent$\Pi_{20}$:
\begin{ptree}
\AXC{$\stackrel{\Pi_{19}}{\vdots}$}\noLine
\UIC{$\forall \varepsilon\,(\mathrm{Conv}(b,\beta) \to \varphi_{21})$}
\AXC{$\forall \varepsilon\,(\mathrm{Conv}(b,\beta) \to \varphi_{21}) \to \mathrm{Conv}(b,\beta) \to \varphi_{21}$\lbl{A2}}
\BIC{$\mathrm{Conv}(b,\beta) \to \varphi_{21}$}
\AXC{$(\mathrm{Conv}(b,\beta) \to \varphi_{21}) \to \varphi_{71}$\lbl{A1}}
\BIC{$\varphi_{71}$}
\end{ptree}
\noindent$\Pi_{21}$:
\begin{ptree}
\AXC{$\mathrm{Conv}(a,\alpha) \to \mathrm{Conv}(b,\beta) \to \varepsilon \in \mathbb{R} \to \varepsilon \in \mathbb{R}$\lbl{A1}}
\AXC{$\mathrm{Conv}(a,\alpha) \to \mathrm{Conv}(b,\beta) \to \mathrm{Conv}(b,\beta)$\lbl{A1}}
\AXC{$\stackrel{\Pi_{20}}{\vdots}$}\noLine
\UIC{$\varphi_{71}$}
\BIC{$\varphi_{67}$}
\BIC{$\varphi_{3} \to \mathrm{Conv}(a,\alpha) \to \mathrm{Conv}(b,\beta) \to \varphi_{21}$}
\end{ptree}
\noindent$\Pi_{22}$:
\begin{ptree}
\AXC{$\varphi_{3}$\lbl{A1}}
\AXC{$\stackrel{\Pi_{21}}{\vdots}$}\noLine
\UIC{$\varphi_{3} \to \mathrm{Conv}(a,\alpha) \to \mathrm{Conv}(b,\beta) \to \varphi_{21}$}
\BIC{$\mathrm{Conv}(a,\alpha) \to \mathrm{Conv}(b,\beta) \to \varphi_{21}$}
\AXC{$\varphi_{74}$\lbl{A1}}
\BIC{$\mathrm{Conv}(a,\alpha) \to \mathrm{Conv}(b,\beta) \to \varepsilon \in \mathbb{R} \to 0 < \varepsilon \to \varphi_{37}$}
\AXC{$\varphi_{103}$\lbl{A1}}
\BIC{$\forall N_{2}\,(\varphi_{72}) \to \varphi_{101}$}
\end{ptree}
\noindent$\Pi_{23}$:
\begin{ptree}
\AXC{$\forall N_{2}\,(\varphi_{72}) \to \forall N_{2}\,(\varphi_{60}) \to \varphi_{9}$\lbl{A3}}
\AXC{$\varphi_{25} \to \neg (N_{1} \cup N_{2} \in \mathbb{N} \land \varphi_{19})$\lbl{A2}}
\AXC{$\stackrel{\Pi_{18}}{\vdots}$}\noLine
\UIC{$\varphi_{87}$}
\BIC{$\mathrm{Conv}(a,\alpha) \to \mathrm{Conv}(b,\beta) \to \varepsilon \in \mathbb{R} \to 0 < \varepsilon \to \varphi_{64}$}
\AXC{$\varphi_{92}$\lbl{A1}}
\BIC{$\varphi_{72}$}
\UIC{$\forall N_{2}\,(\varphi_{72})$}
\AXC{$\stackrel{\Pi_{22}}{\vdots}$}\noLine
\UIC{$\forall N_{2}\,(\varphi_{72}) \to \varphi_{101}$}
\BIC{$\varphi_{101}$}
\BIC{$(\varphi_{60} \to \forall N_{2}\,(\varphi_{60})) \to \varphi_{39} \to (\varphi_{39} \to \forall N\,(\varphi_{9}) \to \varphi_{5}) \to \varphi_{75}$}
\end{ptree}
\noindent$\Pi_{24}$:
\begin{ptree}
\AXC{$\varphi_{60} \to \forall N_{2}\,(\varphi_{60})$\lbl{A4}}
\AXC{$\stackrel{\Pi_{23}}{\vdots}$}\noLine
\UIC{$(\varphi_{60} \to \forall N_{2}\,(\varphi_{60})) \to \varphi_{39} \to (\varphi_{39} \to \forall N\,(\varphi_{9}) \to \varphi_{5}) \to \varphi_{75}$}
\BIC{$\varphi_{39} \to (\varphi_{39} \to \forall N\,(\varphi_{9}) \to \varphi_{5}) \to \varphi_{75}$}
\end{ptree}
\noindent$\Pi_{25}$:
\begin{ptree}
\AXC{$\varphi_{39} \to \forall N\,(\varphi_{9}) \to \varphi_{5}$\lbl{A3}}
\AXC{$\varphi_{9} \to \neg (N \in \mathbb{N} \land \forall n \in \mathbb{N}\,(N \leq n \to \lvert b_{n} - \beta \rvert < \varepsilon))$\lbl{A2}}
\UIC{$\varphi_{39}$}
\AXC{$\stackrel{\Pi_{24}}{\vdots}$}\noLine
\UIC{$\varphi_{39} \to (\varphi_{39} \to \forall N\,(\varphi_{9}) \to \varphi_{5}) \to \varphi_{75}$}
\BIC{$(\varphi_{39} \to \forall N\,(\varphi_{9}) \to \varphi_{5}) \to \varphi_{75}$}
\BIC{$\varphi_{75}$}
\end{ptree}
\noindent$\Pi_{26}$:
\begin{ptree}
\AXC{$\varphi_{9} \to \forall N\,(\varphi_{9})$\lbl{A4}}
\AXC{$\stackrel{\Pi_{25}}{\vdots}$}\noLine
\UIC{$\varphi_{75}$}
\BIC{$\mathrm{Conv}(a,\alpha) \to \mathrm{Conv}(b,\beta) \to \varepsilon \in \mathbb{R} \to \varphi_{47}$}
\AXC{$(\mathrm{Conv}(a,\alpha) \to \mathrm{Conv}(b,\beta) \to \varepsilon \in \mathbb{R} \to \varphi_{47}) \to \varphi_{61}$\lbl{A1}}
\BIC{$\varphi_{61}$}
\UIC{$\forall N_{1}\,(\varphi_{61})$}
\end{ptree}
\noindent$\Pi_{27}$:
\begin{ptree}
\AXC{$\forall a\,\forall \alpha\,\forall \varepsilon\,(\mathrm{Conv}(a,\alpha) \to \varphi_{22})$\lbl{[C3]}}
\AXC{$\varphi_{66}$\lbl{A2}}
\BIC{$\forall \alpha\,\forall \varepsilon\,(\mathrm{Conv}(a,\alpha) \to \varphi_{22})$}
\AXC{$\varphi_{63}$\lbl{A2}}
\BIC{$\forall \varepsilon\,(\mathrm{Conv}(a,\alpha) \to \varphi_{22})$}
\AXC{$\forall \varepsilon\,(\mathrm{Conv}(a,\alpha) \to \varphi_{22}) \to \mathrm{Conv}(a,\alpha) \to \varphi_{22}$\lbl{A2}}
\BIC{$\mathrm{Conv}(a,\alpha) \to \varphi_{22}$}
\end{ptree}
\noindent$\Pi_{28}$:
\begin{ptree}
\AXC{$\mathrm{Conv}(a,\alpha) \to \mathrm{Conv}(a,\alpha)$\lbl{A1}}
\AXC{$\stackrel{\Pi_{27}}{\vdots}$}\noLine
\UIC{$\mathrm{Conv}(a,\alpha) \to \varphi_{22}$}
\AXC{$\varphi_{78}$\lbl{A1}}
\BIC{$(\mathrm{Conv}(a,\alpha) \to \mathrm{Conv}(a,\alpha)) \to \varphi_{68}$}
\BIC{$\varphi_{68}$}
\end{ptree}
\noindent$\Pi_{29}$:
\begin{ptree}
\AXC{$\stackrel{\Pi_{26}}{\vdots}$}\noLine
\UIC{$\forall N_{1}\,(\varphi_{61})$}
\AXC{$\varphi_{3}$\lbl{A1}}
\AXC{$\mathrm{Conv}(a,\alpha) \to \mathrm{Conv}(b,\beta) \to \varepsilon \in \mathbb{R} \to \varepsilon \in \mathbb{R}$\lbl{A1}}
\AXC{$\stackrel{\Pi_{28}}{\vdots}$}\noLine
\UIC{$\varphi_{68}$}
\BIC{$\varphi_{3} \to \mathrm{Conv}(a,\alpha) \to \mathrm{Conv}(b,\beta) \to \varphi_{22}$}
\BIC{$\mathrm{Conv}(a,\alpha) \to \mathrm{Conv}(b,\beta) \to \varphi_{22}$}
\AXC{$\varphi_{102}$\lbl{A1}}
\BIC{$\forall N_{1}\,(\varphi_{61}) \to \varphi_{100}$}
\BIC{$\varphi_{100}$}
\end{ptree}
\noindent$\Pi_{30}$:
\begin{ptree}
\AXC{$\varphi_{8} \to \neg (N \in \mathbb{N} \land \forall n \in \mathbb{N}\,(N \leq n \to \lvert a_{n} - \alpha \rvert < \varepsilon))$\lbl{A2}}
\UIC{$\varphi_{38}$}
\AXC{$\varphi_{40} \to \forall N_{1}\,(\varphi_{40})$\lbl{A4}}
\AXC{$\forall N_{1}\,(\varphi_{61}) \to \forall N_{1}\,(\varphi_{40}) \to \varphi_{8}$\lbl{A3}}
\AXC{$\stackrel{\Pi_{29}}{\vdots}$}\noLine
\UIC{$\varphi_{100}$}
\BIC{$(\varphi_{40} \to \forall N_{1}\,(\varphi_{40})) \to \varphi_{38} \to (\varphi_{38} \to \forall N\,(\varphi_{8}) \to \varphi_{4}) \to \varphi_{73}$}
\BIC{$\varphi_{38} \to (\varphi_{38} \to \forall N\,(\varphi_{8}) \to \varphi_{4}) \to \varphi_{73}$}
\BIC{$(\varphi_{38} \to \forall N\,(\varphi_{8}) \to \varphi_{4}) \to \varphi_{73}$}
\end{ptree}
\noindent 主木:
\begin{ptree}
\AXC{$\varphi_{8} \to \forall N\,(\varphi_{8})$\lbl{A4}}
\AXC{$\varphi_{38} \to \forall N\,(\varphi_{8}) \to \varphi_{4}$\lbl{A3}}
\AXC{$\stackrel{\Pi_{30}}{\vdots}$}\noLine
\UIC{$(\varphi_{38} \to \forall N\,(\varphi_{8}) \to \varphi_{4}) \to \varphi_{73}$}
\BIC{$\varphi_{73}$}
\BIC{$\mathrm{Conv}(a,\alpha) \to \mathrm{Conv}(b,\beta) \to \varepsilon \in \mathbb{R} \to 0 < \varepsilon \to \varphi_{23}$}
\UIC{$\varphi_{42}$}
\UIC{$\forall \beta\,(\varphi_{42})$}
\UIC{$\forall b\,\forall \beta\,(\varphi_{42})$}
\UIC{$\forall \alpha\,\forall b\,\forall \beta\,(\varphi_{42})$}
\UIC{$\forall a\,\forall \alpha\,\forall b\,\forall \beta\,(\varphi_{42})$}
\end{ptree}
\begin{abbrevlist}
\item $\varphi_{1} :\equiv \forall x\,\forall y\,\forall z\,(x \leq y \to y \leq z \to x \leq z) \to \forall y\,\forall z\,(N_{2} \leq y \to y \leq z \to N_{2} \leq z)$
\item $\varphi_{2} :\equiv \forall x\,\forall y\,\forall z\,(x \leq y \to y \leq z \to x \leq z) \to \forall y\,\forall z\,(N_{1} \leq y \to y \leq z \to N_{1} \leq z)$
\item $\varphi_{3} :\equiv \mathrm{Conv}(a,\alpha) \to \mathrm{Conv}(b,\beta) \to \varepsilon \in \mathbb{R} \to 0 < \varepsilon \to 0 < \varepsilon$
\item $\varphi_{4} :\equiv \forall N\,(\neg (N \in \mathbb{N} \land \forall n \in \mathbb{N}\,(N \leq n \to \lvert a_{n} - \alpha \rvert < \varepsilon)))$
\item $\varphi_{5} :\equiv \forall N\,(\neg (N \in \mathbb{N} \land \forall n \in \mathbb{N}\,(N \leq n \to \lvert b_{n} - \beta \rvert < \varepsilon)))$
\item $\varphi_{6} :\equiv \forall y\,\forall z\,(N_{2} \leq y \to y \leq z \to N_{2} \leq z) \to \forall z\,(N_{2} \leq N_{1} \cup N_{2} \to N_{1} \cup N_{2} \leq z \to N_{2} \leq z)$
\item $\varphi_{7} :\equiv \forall y\,\forall z\,(N_{1} \leq y \to y \leq z \to N_{1} \leq z) \to \forall z\,(N_{1} \leq N_{1} \cup N_{2} \to N_{1} \cup N_{2} \leq z \to N_{1} \leq z)$
\item $\varphi_{8} :\equiv \forall N_{1}\,(\neg (N_{1} \in \mathbb{N} \land \forall n \in \mathbb{N}\,(N_{1} \leq n \to \lvert a_{n} - \alpha \rvert < \varepsilon)))$
\item $\varphi_{9} :\equiv \forall N_{2}\,(\neg (N_{2} \in \mathbb{N} \land \forall n \in \mathbb{N}\,(N_{2} \leq n \to \lvert b_{n} - \beta \rvert < \varepsilon)))$
\item $\varphi_{10} :\equiv n \in \mathbb{N} \to N_{1} \cup N_{2} \leq n \to n \in \mathbb{N} \to N_{2} \leq n \to \lvert b_{n} - \beta \rvert < \varepsilon$
\item $\varphi_{11} :\equiv n \in \mathbb{N} \to N_{1} \cup N_{2} \leq n \to n \in \mathbb{N} \to N_{1} \leq n \to \lvert a_{n} - \alpha \rvert < \varepsilon$
\item $\varphi_{12} :\equiv N_{2} \in \mathbb{N} \land \forall n \in \mathbb{N}\,(N_{2} \leq n \to \lvert b_{n} - \beta \rvert < \varepsilon) \to N_{2} \in \mathbb{N}$
\item $\varphi_{13} :\equiv N_{2} \in \mathbb{N} \land \forall n \in \mathbb{N}\,(N_{2} \leq n \to \lvert b_{n} - \beta \rvert < \varepsilon) \to N_{1} \in \mathbb{N}$
\item $\varphi_{14} :\equiv n \in \mathbb{N} \to N_{1} \cup N_{2} \leq n \to \forall n \in \mathbb{N}\,(N_{2} \leq n \to \lvert b_{n} - \beta \rvert < \varepsilon)$
\item $\varphi_{15} :\equiv n \in \mathbb{N} \to N_{1} \cup N_{2} \leq n \to \forall n \in \mathbb{N}\,(N_{1} \leq n \to \lvert a_{n} - \alpha \rvert < \varepsilon)$
\item $\varphi_{16} :\equiv \forall z\,(N_{2} \leq N_{1} \cup N_{2} \to N_{1} \cup N_{2} \leq z \to N_{2} \leq z) \to N_{2} \leq N_{1} \cup N_{2} \to N_{1} \cup N_{2} \leq n \to N_{2} \leq n$
\item $\varphi_{17} :\equiv \forall z\,(N_{1} \leq N_{1} \cup N_{2} \to N_{1} \cup N_{2} \leq z \to N_{1} \leq z) \to N_{1} \leq N_{1} \cup N_{2} \to N_{1} \cup N_{2} \leq n \to N_{1} \leq n$
\item $\varphi_{18} :\equiv N_{2} \in \mathbb{N} \land \forall n \in \mathbb{N}\,(N_{2} \leq n \to \lvert b_{n} - \beta \rvert < \varepsilon) \to N_{1} \cup N_{2} \in \mathbb{N}$
\item $\varphi_{19} :\equiv \forall n \in \mathbb{N}\,(N_{1} \cup N_{2} \leq n \to \lvert a_{n} - \alpha \rvert < \varepsilon \land \lvert b_{n} - \beta \rvert < \varepsilon)$
\item $\varphi_{20} :\equiv N_{2} \in \mathbb{N} \land \forall n \in \mathbb{N}\,(N_{2} \leq n \to \lvert b_{n} - \beta \rvert < \varepsilon) \to n \in \mathbb{N} \to n \in \mathbb{N}$
\item $\varphi_{21} :\equiv \varepsilon \in \mathbb{R} \to 0 < \varepsilon \to \exists N \in \mathbb{N}\,\forall n \in \mathbb{N}\,(N \leq n \to \lvert b_{n} - \beta \rvert < \varepsilon)$
\item $\varphi_{22} :\equiv \varepsilon \in \mathbb{R} \to 0 < \varepsilon \to \exists N \in \mathbb{N}\,\forall n \in \mathbb{N}\,(N \leq n \to \lvert a_{n} - \alpha \rvert < \varepsilon)$
\item $\varphi_{23} :\equiv \exists N \in \mathbb{N}\,\forall n \in \mathbb{N}\,(N \leq n \to \lvert a_{n} - \alpha \rvert < \varepsilon \land \lvert b_{n} - \beta \rvert < \varepsilon)$
\item $\varphi_{24} :\equiv \forall n\,(N_{1} \in \mathbb{N} \to n \in \mathbb{N} \to N_{1} \cup n \in \mathbb{N}) \to N_{1} \in \mathbb{N} \to N_{2} \in \mathbb{N} \to N_{1} \cup N_{2} \in \mathbb{N}$
\item $\varphi_{25} :\equiv \forall N\,(\neg (N \in \mathbb{N} \land \forall n \in \mathbb{N}\,(N \leq n \to \lvert a_{n} - \alpha \rvert < \varepsilon \land \lvert b_{n} - \beta \rvert < \varepsilon)))$
\item $\varphi_{26} :\equiv \forall m\,\forall n\,(m \in \mathbb{N} \to n \in \mathbb{N} \to m \cup n \in \mathbb{N}) \to \forall n\,(N_{1} \in \mathbb{N} \to n \in \mathbb{N} \to N_{1} \cup n \in \mathbb{N})$
\item $\varphi_{27} :\equiv N_{2} \in \mathbb{N} \land \forall n \in \mathbb{N}\,(N_{2} \leq n \to \lvert b_{n} - \beta \rvert < \varepsilon) \to n \in \mathbb{N} \to N_{1} \cup N_{2} \leq n \to N_{2} \leq n$
\item $\varphi_{28} :\equiv N_{2} \in \mathbb{N} \land \forall n \in \mathbb{N}\,(N_{2} \leq n \to \lvert b_{n} - \beta \rvert < \varepsilon) \to n \in \mathbb{N} \to N_{1} \cup N_{2} \leq n \to N_{1} \leq n$
\item $\varphi_{29} :\equiv \forall n \in \mathbb{N}\,(N_{2} \leq n \to \lvert b_{n} - \beta \rvert < \varepsilon) \to n \in \mathbb{N} \to N_{2} \leq n \to \lvert b_{n} - \beta \rvert < \varepsilon$
\item $\varphi_{30} :\equiv \forall n \in \mathbb{N}\,(N_{1} \leq n \to \lvert a_{n} - \alpha \rvert < \varepsilon) \to n \in \mathbb{N} \to N_{1} \leq n \to \lvert a_{n} - \alpha \rvert < \varepsilon$
\item $\varphi_{31} :\equiv N_{2} \in \mathbb{N} \land \forall n \in \mathbb{N}\,(N_{2} \leq n \to \lvert b_{n} - \beta \rvert < \varepsilon) \to n \in \mathbb{N} \to N_{1} \cup N_{2} \leq n \to N_{1} \cup N_{2} \leq n$
\item $\varphi_{32} :\equiv \forall \alpha\,\forall \varepsilon\,(\mathrm{Conv}(b,\alpha) \to \varepsilon \in \mathbb{R} \to 0 < \varepsilon \to \exists N \in \mathbb{N}\,\forall n \in \mathbb{N}\,(N \leq n \to \lvert b_{n} - \alpha \rvert < \varepsilon))$
\item $\varphi_{33} :\equiv N_{2} \in \mathbb{N} \land \forall n \in \mathbb{N}\,(N_{2} \leq n \to \lvert b_{n} - \beta \rvert < \varepsilon) \to n \in \mathbb{N} \to N_{1} \cup N_{2} \leq n \to \lvert b_{n} - \beta \rvert < \varepsilon$
\item $\varphi_{34} :\equiv N_{2} \in \mathbb{N} \land \forall n \in \mathbb{N}\,(N_{2} \leq n \to \lvert b_{n} - \beta \rvert < \varepsilon) \to n \in \mathbb{N} \to N_{1} \cup N_{2} \leq n \to \lvert a_{n} - \alpha \rvert < \varepsilon$
\item $\varphi_{35} :\equiv N_{2} \in \mathbb{N} \land \forall n \in \mathbb{N}\,(N_{2} \leq n \to \lvert b_{n} - \beta \rvert < \varepsilon) \to N_{2} \in \mathbb{N} \land \forall n \in \mathbb{N}\,(N_{2} \leq n \to \lvert b_{n} - \beta \rvert < \varepsilon)$
\item $\varphi_{36} :\equiv N_{1} \in \mathbb{N} \land \forall n \in \mathbb{N}\,(N_{1} \leq n \to \lvert a_{n} - \alpha \rvert < \varepsilon) \to N_{1} \in \mathbb{N} \land \forall n \in \mathbb{N}\,(N_{1} \leq n \to \lvert a_{n} - \alpha \rvert < \varepsilon)$
\item $\varphi_{37} :\equiv N_{1} \in \mathbb{N} \land \forall n \in \mathbb{N}\,(N_{1} \leq n \to \lvert a_{n} - \alpha \rvert < \varepsilon) \to \exists N \in \mathbb{N}\,\forall n \in \mathbb{N}\,(N \leq n \to \lvert b_{n} - \beta \rvert < \varepsilon)$
\item $\varphi_{38} :\equiv \forall N\,(\varphi_{8} \to \neg (N \in \mathbb{N} \land \forall n \in \mathbb{N}\,(N \leq n \to \lvert a_{n} - \alpha \rvert < \varepsilon)))$
\item $\varphi_{39} :\equiv \forall N\,(\varphi_{9} \to \neg (N \in \mathbb{N} \land \forall n \in \mathbb{N}\,(N \leq n \to \lvert b_{n} - \beta \rvert < \varepsilon)))$
\item $\varphi_{40} :\equiv \mathrm{Conv}(a,\alpha) \land \mathrm{Conv}(b,\beta) \land \varepsilon \in \mathbb{R} \land 0 < \varepsilon \land \neg \varphi_{23}$
\item $\varphi_{41} :\equiv N_{2} \in \mathbb{N} \land \forall n \in \mathbb{N}\,(N_{2} \leq n \to \lvert b_{n} - \beta \rvert < \varepsilon) \to n \in \mathbb{N} \to N_{1} \cup N_{2} \leq n \to \lvert a_{n} - \alpha \rvert < \varepsilon \land \lvert b_{n} - \beta \rvert < \varepsilon$
\item $\varphi_{42} :\equiv \forall \varepsilon\,(\mathrm{Conv}(a,\alpha) \to \mathrm{Conv}(b,\beta) \to \varepsilon \in \mathbb{R} \to 0 < \varepsilon \to \varphi_{23})$
\item $\varphi_{43} :\equiv 0 < \varepsilon \to N_{1} \in \mathbb{N} \land \forall n \in \mathbb{N}\,(N_{1} \leq n \to \lvert a_{n} - \alpha \rvert < \varepsilon) \to \varphi_{12}$
\item $\varphi_{44} :\equiv 0 < \varepsilon \to N_{1} \in \mathbb{N} \land \forall n \in \mathbb{N}\,(N_{1} \leq n \to \lvert a_{n} - \alpha \rvert < \varepsilon) \to \varphi_{13}$
\item $\varphi_{45} :\equiv 0 < \varepsilon \to N_{1} \in \mathbb{N} \land \forall n \in \mathbb{N}\,(N_{1} \leq n \to \lvert a_{n} - \alpha \rvert < \varepsilon) \to \varphi_{18}$
\item $\varphi_{46} :\equiv 0 < \varepsilon \to N_{1} \in \mathbb{N} \land \forall n \in \mathbb{N}\,(N_{1} \leq n \to \lvert a_{n} - \alpha \rvert < \varepsilon) \to \varphi_{20}$
\item $\varphi_{47} :\equiv 0 < \varepsilon \to N_{1} \in \mathbb{N} \land \forall n \in \mathbb{N}\,(N_{1} \leq n \to \lvert a_{n} - \alpha \rvert < \varepsilon) \to \varphi_{23}$
\item $\varphi_{48} :\equiv 0 < \varepsilon \to N_{1} \in \mathbb{N} \land \forall n \in \mathbb{N}\,(N_{1} \leq n \to \lvert a_{n} - \alpha \rvert < \varepsilon) \to \varphi_{27}$
\item $\varphi_{49} :\equiv 0 < \varepsilon \to N_{1} \in \mathbb{N} \land \forall n \in \mathbb{N}\,(N_{1} \leq n \to \lvert a_{n} - \alpha \rvert < \varepsilon) \to \varphi_{28}$
\item $\varphi_{50} :\equiv N_{2} \in \mathbb{N} \land \forall n \in \mathbb{N}\,(N_{2} \leq n \to \lvert b_{n} - \beta \rvert < \varepsilon) \to N_{1} \cup N_{2} \in \mathbb{N} \land \varphi_{19}$
\item $\varphi_{51} :\equiv \mathrm{Conv}(a,\alpha) \land \mathrm{Conv}(b,\beta) \land \varepsilon \in \mathbb{R} \land 0 < \varepsilon \land (N_{1} \in \mathbb{N} \land \forall n \in \mathbb{N}\,(N_{1} \leq n \to \lvert a_{n} - \alpha \rvert < \varepsilon)) \land (N_{2} \in \mathbb{N} \land \forall n \in \mathbb{N}\,(N_{2} \leq n \to \lvert b_{n} - \beta \rvert < \varepsilon))$
\item $\varphi_{52} :\equiv 0 < \varepsilon \to N_{1} \in \mathbb{N} \land \forall n \in \mathbb{N}\,(N_{1} \leq n \to \lvert a_{n} - \alpha \rvert < \varepsilon) \to \varphi_{31}$
\item $\varphi_{53} :\equiv 0 < \varepsilon \to N_{1} \in \mathbb{N} \land \forall n \in \mathbb{N}\,(N_{1} \leq n \to \lvert a_{n} - \alpha \rvert < \varepsilon) \to \varphi_{33}$
\item $\varphi_{54} :\equiv 0 < \varepsilon \to N_{1} \in \mathbb{N} \land \forall n \in \mathbb{N}\,(N_{1} \leq n \to \lvert a_{n} - \alpha \rvert < \varepsilon) \to \varphi_{34}$
\item $\varphi_{55} :\equiv 0 < \varepsilon \to N_{1} \in \mathbb{N} \land \forall n \in \mathbb{N}\,(N_{1} \leq n \to \lvert a_{n} - \alpha \rvert < \varepsilon) \to \varphi_{35}$
\item $\varphi_{56} :\equiv N_{1} \in \mathbb{N} \land \forall n \in \mathbb{N}\,(N_{1} \leq n \to \lvert a_{n} - \alpha \rvert < \varepsilon) \to N_{2} \in \mathbb{N} \land \forall n \in \mathbb{N}\,(N_{2} \leq n \to \lvert b_{n} - \beta \rvert < \varepsilon) \to \varphi_{10}$
\item $\varphi_{57} :\equiv N_{1} \in \mathbb{N} \land \forall n \in \mathbb{N}\,(N_{1} \leq n \to \lvert a_{n} - \alpha \rvert < \varepsilon) \to N_{2} \in \mathbb{N} \land \forall n \in \mathbb{N}\,(N_{2} \leq n \to \lvert b_{n} - \beta \rvert < \varepsilon) \to \varphi_{11}$
\item $\varphi_{58} :\equiv N_{1} \in \mathbb{N} \land \forall n \in \mathbb{N}\,(N_{1} \leq n \to \lvert a_{n} - \alpha \rvert < \varepsilon) \to N_{2} \in \mathbb{N} \land \forall n \in \mathbb{N}\,(N_{2} \leq n \to \lvert b_{n} - \beta \rvert < \varepsilon) \to \varphi_{14}$
\item $\varphi_{59} :\equiv N_{1} \in \mathbb{N} \land \forall n \in \mathbb{N}\,(N_{1} \leq n \to \lvert a_{n} - \alpha \rvert < \varepsilon) \to N_{2} \in \mathbb{N} \land \forall n \in \mathbb{N}\,(N_{2} \leq n \to \lvert b_{n} - \beta \rvert < \varepsilon) \to \varphi_{15}$
\item $\varphi_{60} :\equiv \mathrm{Conv}(a,\alpha) \land \mathrm{Conv}(b,\beta) \land \varepsilon \in \mathbb{R} \land 0 < \varepsilon \land (N_{1} \in \mathbb{N} \land \forall n \in \mathbb{N}\,(N_{1} \leq n \to \lvert a_{n} - \alpha \rvert < \varepsilon)) \land \neg \varphi_{23}$
\item $\varphi_{61} :\equiv \varphi_{40} \to \neg (N_{1} \in \mathbb{N} \land \forall n \in \mathbb{N}\,(N_{1} \leq n \to \lvert a_{n} - \alpha \rvert < \varepsilon))$
\item $\varphi_{62} :\equiv N_{1} \in \mathbb{N} \land \forall n \in \mathbb{N}\,(N_{1} \leq n \to \lvert a_{n} - \alpha \rvert < \varepsilon) \to N_{2} \in \mathbb{N} \land \forall n \in \mathbb{N}\,(N_{2} \leq n \to \lvert b_{n} - \beta \rvert < \varepsilon) \to \varphi_{19}$
\item $\varphi_{63} :\equiv \forall \alpha\,\forall \varepsilon\,(\mathrm{Conv}(a,\alpha) \to \varphi_{22}) \to \forall \varepsilon\,(\mathrm{Conv}(a,\alpha) \to \varphi_{22})$
\item $\varphi_{64} :\equiv N_{1} \in \mathbb{N} \land \forall n \in \mathbb{N}\,(N_{1} \leq n \to \lvert a_{n} - \alpha \rvert < \varepsilon) \to N_{2} \in \mathbb{N} \land \forall n \in \mathbb{N}\,(N_{2} \leq n \to \lvert b_{n} - \beta \rvert < \varepsilon) \to \varphi_{23}$
\item $\varphi_{65} :\equiv 0 < \varepsilon \to N_{1} \in \mathbb{N} \land \forall n \in \mathbb{N}\,(N_{1} \leq n \to \lvert a_{n} - \alpha \rvert < \varepsilon) \to \varphi_{41}$
\item $\varphi_{66} :\equiv \forall a\,\forall \alpha\,\forall \varepsilon\,(\mathrm{Conv}(a,\alpha) \to \varphi_{22}) \to \forall \alpha\,\forall \varepsilon\,(\mathrm{Conv}(a,\alpha) \to \varphi_{22})$
\item $\varphi_{67} :\equiv (\mathrm{Conv}(a,\alpha) \to \mathrm{Conv}(b,\beta) \to \varepsilon \in \mathbb{R} \to \varepsilon \in \mathbb{R}) \to \varphi_{3} \to \mathrm{Conv}(a,\alpha) \to \mathrm{Conv}(b,\beta) \to \varphi_{21}$
\item $\varphi_{68} :\equiv (\mathrm{Conv}(a,\alpha) \to \mathrm{Conv}(b,\beta) \to \varepsilon \in \mathbb{R} \to \varepsilon \in \mathbb{R}) \to \varphi_{3} \to \mathrm{Conv}(a,\alpha) \to \mathrm{Conv}(b,\beta) \to \varphi_{22}$
\item $\varphi_{69} :\equiv 0 < \varepsilon \to N_{1} \in \mathbb{N} \land \forall n \in \mathbb{N}\,(N_{1} \leq n \to \lvert a_{n} - \alpha \rvert < \varepsilon) \to \varphi_{50}$
\item $\varphi_{70} :\equiv \varphi_{51} \to n \in \mathbb{N} \to N_{1} \cup N_{2} \leq n \to \lvert a_{n} - \alpha \rvert < \varepsilon \land \lvert b_{n} - \beta \rvert < \varepsilon$
\item $\varphi_{71} :\equiv (\mathrm{Conv}(a,\alpha) \to \mathrm{Conv}(b,\beta) \to \mathrm{Conv}(b,\beta)) \to \varphi_{67}$
\item $\varphi_{72} :\equiv \varphi_{60} \to \neg (N_{2} \in \mathbb{N} \land \forall n \in \mathbb{N}\,(N_{2} \leq n \to \lvert b_{n} - \beta \rvert < \varepsilon))$
\item $\varphi_{73} :\equiv (\varphi_{8} \to \forall N\,(\varphi_{8})) \to \mathrm{Conv}(a,\alpha) \to \mathrm{Conv}(b,\beta) \to \varepsilon \in \mathbb{R} \to 0 < \varepsilon \to \varphi_{23}$
\item $\varphi_{74} :\equiv (\mathrm{Conv}(a,\alpha) \to \mathrm{Conv}(b,\beta) \to \varphi_{21}) \to \mathrm{Conv}(a,\alpha) \to \mathrm{Conv}(b,\beta) \to \varepsilon \in \mathbb{R} \to 0 < \varepsilon \to \varphi_{37}$
\item $\varphi_{75} :\equiv (\varphi_{9} \to \forall N\,(\varphi_{9})) \to \mathrm{Conv}(a,\alpha) \to \mathrm{Conv}(b,\beta) \to \varepsilon \in \mathbb{R} \to \varphi_{47}$
\item $\varphi_{76} :\equiv \varphi_{29} \to \mathrm{Conv}(a,\alpha) \to \mathrm{Conv}(b,\beta) \to \varepsilon \in \mathbb{R} \to 0 < \varepsilon \to \varphi_{56}$
\item $\varphi_{77} :\equiv \varphi_{30} \to \mathrm{Conv}(a,\alpha) \to \mathrm{Conv}(b,\beta) \to \varepsilon \in \mathbb{R} \to 0 < \varepsilon \to \varphi_{57}$
\item $\varphi_{78} :\equiv (\mathrm{Conv}(a,\alpha) \to \varphi_{22}) \to (\mathrm{Conv}(a,\alpha) \to \mathrm{Conv}(a,\alpha)) \to \varphi_{68}$
\item $\varphi_{79} :\equiv (\mathrm{Conv}(a,\alpha) \to \mathrm{Conv}(b,\beta) \to \varepsilon \in \mathbb{R} \to 0 < \varepsilon \to \varphi_{36}) \to \mathrm{Conv}(a,\alpha) \to \mathrm{Conv}(b,\beta) \to \varepsilon \in \mathbb{R} \to \varphi_{44}$
\item $\varphi_{80} :\equiv (\mathrm{Conv}(a,\alpha) \to \mathrm{Conv}(b,\beta) \to \varepsilon \in \mathbb{R} \to \varphi_{43}) \to \mathrm{Conv}(a,\alpha) \to \mathrm{Conv}(b,\beta) \to \varepsilon \in \mathbb{R} \to \varphi_{45}$
\item $\varphi_{81} :\equiv (\mathrm{Conv}(a,\alpha) \to \mathrm{Conv}(b,\beta) \to \varepsilon \in \mathbb{R} \to \varphi_{52}) \to \mathrm{Conv}(a,\alpha) \to \mathrm{Conv}(b,\beta) \to \varepsilon \in \mathbb{R} \to \varphi_{48}$
\item $\varphi_{82} :\equiv (\mathrm{Conv}(a,\alpha) \to \mathrm{Conv}(b,\beta) \to \varepsilon \in \mathbb{R} \to \varphi_{52}) \to \mathrm{Conv}(a,\alpha) \to \mathrm{Conv}(b,\beta) \to \varepsilon \in \mathbb{R} \to \varphi_{49}$
\item $\varphi_{83} :\equiv (\mathrm{Conv}(a,\alpha) \to \mathrm{Conv}(b,\beta) \to \varepsilon \in \mathbb{R} \to \varphi_{55}) \to \mathrm{Conv}(a,\alpha) \to \mathrm{Conv}(b,\beta) \to \varepsilon \in \mathbb{R} \to \varphi_{43}$
\item $\varphi_{84} :\equiv (\mathrm{Conv}(a,\alpha) \to \mathrm{Conv}(b,\beta) \to \varepsilon \in \mathbb{R} \to 0 < \varepsilon \to \varphi_{36}) \to \mathrm{Conv}(a,\alpha) \to \mathrm{Conv}(b,\beta) \to \varepsilon \in \mathbb{R} \to 0 < \varepsilon \to \varphi_{59}$
\item $\varphi_{85} :\equiv (\mathrm{Conv}(a,\alpha) \to \mathrm{Conv}(b,\beta) \to \varepsilon \in \mathbb{R} \to \varphi_{48}) \to \mathrm{Conv}(a,\alpha) \to \mathrm{Conv}(b,\beta) \to \varepsilon \in \mathbb{R} \to \varphi_{53}$
\item $\varphi_{86} :\equiv (\mathrm{Conv}(a,\alpha) \to \mathrm{Conv}(b,\beta) \to \varepsilon \in \mathbb{R} \to \varphi_{49}) \to \mathrm{Conv}(a,\alpha) \to \mathrm{Conv}(b,\beta) \to \varepsilon \in \mathbb{R} \to \varphi_{54}$
\item $\varphi_{87} :\equiv (\varphi_{25} \to \neg (N_{1} \cup N_{2} \in \mathbb{N} \land \varphi_{19})) \to \mathrm{Conv}(a,\alpha) \to \mathrm{Conv}(b,\beta) \to \varepsilon \in \mathbb{R} \to 0 < \varepsilon \to \varphi_{64}$
\item $\varphi_{88} :\equiv (N_{2} \leq N_{1} \cup N_{2} \to N_{1} \cup N_{2} \leq n \to N_{2} \leq n) \to N_{2} \leq N_{1} \cup N_{2} \to \varphi_{81}$
\item $\varphi_{89} :\equiv (N_{1} \leq N_{1} \cup N_{2} \to N_{1} \cup N_{2} \leq n \to N_{1} \leq n) \to N_{1} \leq N_{1} \cup N_{2} \to \varphi_{82}$
\item $\varphi_{90} :\equiv (\mathrm{Conv}(a,\alpha) \to \mathrm{Conv}(b,\beta) \to \varepsilon \in \mathbb{R} \to \varphi_{53}) \to \mathrm{Conv}(a,\alpha) \to \mathrm{Conv}(b,\beta) \to \varepsilon \in \mathbb{R} \to \varphi_{65}$
\item $\varphi_{91} :\equiv (\mathrm{Conv}(a,\alpha) \to \mathrm{Conv}(b,\beta) \to \varepsilon \in \mathbb{R} \to \varphi_{55}) \to \mathrm{Conv}(a,\alpha) \to \mathrm{Conv}(b,\beta) \to \varepsilon \in \mathbb{R} \to 0 < \varepsilon \to \varphi_{58}$
\item $\varphi_{92} :\equiv (\mathrm{Conv}(a,\alpha) \to \mathrm{Conv}(b,\beta) \to \varepsilon \in \mathbb{R} \to 0 < \varepsilon \to \varphi_{64}) \to \varphi_{72}$
\item $\varphi_{93} :\equiv (\mathrm{Conv}(a,\alpha) \to \mathrm{Conv}(b,\beta) \to \varepsilon \in \mathbb{R} \to 0 < \varepsilon \to \varphi_{62}) \to \mathrm{Conv}(a,\alpha) \to \mathrm{Conv}(b,\beta) \to \varepsilon \in \mathbb{R} \to \varphi_{69}$
\item $\varphi_{94} :\equiv (\mathrm{Conv}(a,\alpha) \to \mathrm{Conv}(b,\beta) \to \varepsilon \in \mathbb{R} \to 0 < \varepsilon \to \varphi_{58}) \to \varphi_{76}$
\item $\varphi_{95} :\equiv (\mathrm{Conv}(a,\alpha) \to \mathrm{Conv}(b,\beta) \to \varepsilon \in \mathbb{R} \to 0 < \varepsilon \to \varphi_{59}) \to \varphi_{77}$
\item $\varphi_{96} :\equiv (\varphi_{51} \to \forall n\,(\varphi_{51})) \to \mathrm{Conv}(a,\alpha) \to \mathrm{Conv}(b,\beta) \to \varepsilon \in \mathbb{R} \to 0 < \varepsilon \to \varphi_{62}$
\item $\varphi_{97} :\equiv (N_{1} \in \mathbb{N} \to N_{2} \in \mathbb{N} \to N_{1} \cup N_{2} \in \mathbb{N}) \to (\mathrm{Conv}(a,\alpha) \to \mathrm{Conv}(b,\beta) \to \varepsilon \in \mathbb{R} \to \varphi_{44}) \to \varphi_{80}$
\item $\varphi_{98} :\equiv (\mathrm{Conv}(a,\alpha) \to \mathrm{Conv}(b,\beta) \to \varepsilon \in \mathbb{R} \to 0 < \varepsilon \to \varphi_{56}) \to (\mathrm{Conv}(a,\alpha) \to \mathrm{Conv}(b,\beta) \to \varepsilon \in \mathbb{R} \to \varphi_{46}) \to \varphi_{85}$
\item $\varphi_{99} :\equiv (\mathrm{Conv}(a,\alpha) \to \mathrm{Conv}(b,\beta) \to \varepsilon \in \mathbb{R} \to 0 < \varepsilon \to \varphi_{57}) \to (\mathrm{Conv}(a,\alpha) \to \mathrm{Conv}(b,\beta) \to \varepsilon \in \mathbb{R} \to \varphi_{46}) \to \varphi_{86}$
\item $\varphi_{100} :\equiv (\forall N_{1}\,(\varphi_{61}) \to \forall N_{1}\,(\varphi_{40}) \to \varphi_{8}) \to (\varphi_{40} \to \forall N_{1}\,(\varphi_{40})) \to \varphi_{38} \to (\varphi_{38} \to \forall N\,(\varphi_{8}) \to \varphi_{4}) \to \varphi_{73}$
\item $\varphi_{101} :\equiv (\forall N_{2}\,(\varphi_{72}) \to \forall N_{2}\,(\varphi_{60}) \to \varphi_{9}) \to (\varphi_{60} \to \forall N_{2}\,(\varphi_{60})) \to \varphi_{39} \to (\varphi_{39} \to \forall N\,(\varphi_{9}) \to \varphi_{5}) \to \varphi_{75}$
\item $\varphi_{102} :\equiv (\mathrm{Conv}(a,\alpha) \to \mathrm{Conv}(b,\beta) \to \varphi_{22}) \to \forall N_{1}\,(\varphi_{61}) \to \varphi_{100}$
\item $\varphi_{103} :\equiv (\mathrm{Conv}(a,\alpha) \to \mathrm{Conv}(b,\beta) \to \varepsilon \in \mathbb{R} \to 0 < \varepsilon \to \varphi_{37}) \to \forall N_{2}\,(\varphi_{72}) \to \varphi_{101}$
\end{abbrevlist}
\end{fullproof}

% 自動生成 (prover/build.js): EpsEq  ノード数 56
\begin{fullproof}{命題1.1}{$\forall a\,\forall b\,(a \in \mathbb{R} \to b \in \mathbb{R} \to \forall \varepsilon \in \mathbb{R}\,(0 < \varepsilon \to \lvert a - b \rvert < \varepsilon) \to a = b)$}{56}\phantomsection\label{tree:EpsEq}
\noindent$\Pi_{1}$:
\begin{ptree}
\AXC{$\forall x\,(\neg x < x)$\lbl{[O1]}}
\AXC{$\forall x\,(\neg x < x) \to \neg \lvert a - b \rvert < \lvert a - b \rvert$\lbl{A2}}
\BIC{$\neg \lvert a - b \rvert < \lvert a - b \rvert$}
\AXC{$\neg \lvert a - b \rvert < \lvert a - b \rvert \to a \in \mathbb{R} \to b \in \mathbb{R} \to \varphi_{9}$\lbl{A1}}
\BIC{$a \in \mathbb{R} \to b \in \mathbb{R} \to \varphi_{9}$}
\end{ptree}
\noindent$\Pi_{2}$:
\begin{ptree}
\AXC{$\forall x\,\forall y\,(x \in \mathbb{R} \to y \in \mathbb{R} \to x \neq y \to x - y \neq 0)$\lbl{[A15]}}
\AXC{$\varphi_{15}$\lbl{A2}}
\BIC{$\forall y\,(a \in \mathbb{R} \to y \in \mathbb{R} \to a \neq y \to a - y \neq 0)$}
\AXC{$\varphi_{13}$\lbl{A2}}
\BIC{$a \in \mathbb{R} \to b \in \mathbb{R} \to a \neq b \to a - b \neq 0$}
\AXC{$(a \in \mathbb{R} \to b \in \mathbb{R} \to a \neq b \to a - b \neq 0) \to \varphi_{22}$\lbl{A1}}
\BIC{$\varphi_{22}$}
\end{ptree}
\noindent$\Pi_{3}$:
\begin{ptree}
\AXC{$a \in \mathbb{R} \to \varphi_{1}$\lbl{A1}}
\AXC{$a \in \mathbb{R} \to \varphi_{5}$\lbl{A1}}
\AXC{$a \in \mathbb{R} \to \varphi_{6}$\lbl{A1}}
\AXC{$\stackrel{\Pi_{2}}{\vdots}$}\noLine
\UIC{$\varphi_{22}$}
\BIC{$(a \in \mathbb{R} \to \varphi_{5}) \to (a \in \mathbb{R} \to \varphi_{1}) \to a \in \mathbb{R} \to \varphi_{4}$}
\BIC{$(a \in \mathbb{R} \to \varphi_{1}) \to a \in \mathbb{R} \to \varphi_{4}$}
\BIC{$a \in \mathbb{R} \to \varphi_{4}$}
\end{ptree}
\noindent$\Pi_{4}$:
\begin{ptree}
\AXC{$a \in \mathbb{R} \to \varphi_{6}$\lbl{A1}}
\AXC{$\forall x\,\forall y\,(x \in \mathbb{R} \to y \in \mathbb{R} \to x - y \in \mathbb{R})$\lbl{[F1]}}
\AXC{$\varphi_{14}$\lbl{A2}}
\BIC{$\forall y\,(a \in \mathbb{R} \to y \in \mathbb{R} \to a - y \in \mathbb{R})$}
\AXC{$\varphi_{11}$\lbl{A2}}
\BIC{$a \in \mathbb{R} \to b \in \mathbb{R} \to a - b \in \mathbb{R}$}
\AXC{$\varphi_{20}$\lbl{A1}}
\BIC{$(a \in \mathbb{R} \to \varphi_{6}) \to (a \in \mathbb{R} \to \varphi_{5}) \to a \in \mathbb{R} \to \varphi_{10}$}
\BIC{$(a \in \mathbb{R} \to \varphi_{5}) \to a \in \mathbb{R} \to \varphi_{10}$}
\end{ptree}
\noindent$\Pi_{5}$:
\begin{ptree}
\AXC{$a \in \mathbb{R} \to \varphi_{5}$\lbl{A1}}
\AXC{$\stackrel{\Pi_{4}}{\vdots}$}\noLine
\UIC{$(a \in \mathbb{R} \to \varphi_{5}) \to a \in \mathbb{R} \to \varphi_{10}$}
\BIC{$a \in \mathbb{R} \to \varphi_{10}$}
\end{ptree}
\noindent$\Pi_{6}$:
\begin{ptree}
\AXC{$\stackrel{\Pi_{3}}{\vdots}$}\noLine
\UIC{$a \in \mathbb{R} \to \varphi_{4}$}
\AXC{$\stackrel{\Pi_{5}}{\vdots}$}\noLine
\UIC{$a \in \mathbb{R} \to \varphi_{10}$}
\AXC{$\forall x \in \mathbb{R}\,(x \neq 0 \to 0 < \lvert x \rvert)$\lbl{[A14]}}
\AXC{$\varphi_{7}$\lbl{A2}}
\BIC{$a - b \in \mathbb{R} \to a - b \neq 0 \to 0 < \lvert a - b \rvert$}
\AXC{$\varphi_{21}$\lbl{A1}}
\BIC{$(a \in \mathbb{R} \to \varphi_{10}) \to (a \in \mathbb{R} \to \varphi_{4}) \to a \in \mathbb{R} \to \varphi_{12}$}
\BIC{$(a \in \mathbb{R} \to \varphi_{4}) \to a \in \mathbb{R} \to \varphi_{12}$}
\BIC{$a \in \mathbb{R} \to \varphi_{12}$}
\end{ptree}
\noindent$\Pi_{7}$:
\begin{ptree}
\AXC{$\stackrel{\Pi_{5}}{\vdots}$}\noLine
\UIC{$a \in \mathbb{R} \to \varphi_{10}$}
\AXC{$\forall x \in \mathbb{R}\,(\lvert x \rvert \in \mathbb{R})$\lbl{[A5]}}
\AXC{$\forall x \in \mathbb{R}\,(\lvert x \rvert \in \mathbb{R}) \to a - b \in \mathbb{R} \to \lvert a - b \rvert \in \mathbb{R}$\lbl{A2}}
\BIC{$a - b \in \mathbb{R} \to \lvert a - b \rvert \in \mathbb{R}$}
\AXC{$\varphi_{18}$\lbl{A1}}
\BIC{$(a \in \mathbb{R} \to \varphi_{10}) \to a \in \mathbb{R} \to b \in \mathbb{R} \to \varphi_{2}$}
\BIC{$a \in \mathbb{R} \to b \in \mathbb{R} \to \varphi_{2}$}
\end{ptree}
\noindent$\Pi_{8}$:
\begin{ptree}
\AXC{$\varphi_{17}$\lbl{A2}}
\AXC{$a \in \mathbb{R} \to b \in \mathbb{R} \to \varphi_{16}$\lbl{A1}}
\AXC{$(a \in \mathbb{R} \to b \in \mathbb{R} \to \varphi_{16}) \to \varphi_{17} \to a \in \mathbb{R} \to b \in \mathbb{R} \to \varphi_{17}$\lbl{A1}}
\BIC{$\varphi_{17} \to a \in \mathbb{R} \to b \in \mathbb{R} \to \varphi_{17}$}
\BIC{$a \in \mathbb{R} \to b \in \mathbb{R} \to \varphi_{17}$}
\end{ptree}
\noindent$\Pi_{9}$:
\begin{ptree}
\AXC{$\stackrel{\Pi_{6}}{\vdots}$}\noLine
\UIC{$a \in \mathbb{R} \to \varphi_{12}$}
\AXC{$\stackrel{\Pi_{7}}{\vdots}$}\noLine
\UIC{$a \in \mathbb{R} \to b \in \mathbb{R} \to \varphi_{2}$}
\AXC{$\stackrel{\Pi_{8}}{\vdots}$}\noLine
\UIC{$a \in \mathbb{R} \to b \in \mathbb{R} \to \varphi_{17}$}
\AXC{$(a \in \mathbb{R} \to b \in \mathbb{R} \to \varphi_{17}) \to \varphi_{19}$\lbl{A1}}
\BIC{$\varphi_{19}$}
\BIC{$(a \in \mathbb{R} \to \varphi_{12}) \to a \in \mathbb{R} \to b \in \mathbb{R} \to \varphi_{8}$}
\BIC{$a \in \mathbb{R} \to b \in \mathbb{R} \to \varphi_{8}$}
\end{ptree}
\noindent 主木:
\begin{ptree}
\AXC{$\stackrel{\Pi_{1}}{\vdots}$}\noLine
\UIC{$a \in \mathbb{R} \to b \in \mathbb{R} \to \varphi_{9}$}
\AXC{$\stackrel{\Pi_{9}}{\vdots}$}\noLine
\UIC{$a \in \mathbb{R} \to b \in \mathbb{R} \to \varphi_{8}$}
\AXC{$(a \in \mathbb{R} \to b \in \mathbb{R} \to \varphi_{8}) \to (a \in \mathbb{R} \to b \in \mathbb{R} \to \varphi_{9}) \to \varphi_{3}$\lbl{A1}}
\BIC{$(a \in \mathbb{R} \to b \in \mathbb{R} \to \varphi_{9}) \to \varphi_{3}$}
\BIC{$\varphi_{3}$}
\UIC{$\forall b\,(\varphi_{3})$}
\UIC{$\forall a\,\forall b\,(\varphi_{3})$}
\end{ptree}
\begin{abbrevlist}
\item $\varphi_{1} :\equiv b \in \mathbb{R} \to \forall \varepsilon \in \mathbb{R}\,(0 < \varepsilon \to \lvert a - b \rvert < \varepsilon) \to a \neq b \to a \neq b$
\item $\varphi_{2} :\equiv \forall \varepsilon \in \mathbb{R}\,(0 < \varepsilon \to \lvert a - b \rvert < \varepsilon) \to a \neq b \to \lvert a - b \rvert \in \mathbb{R}$
\item $\varphi_{3} :\equiv a \in \mathbb{R} \to b \in \mathbb{R} \to \forall \varepsilon \in \mathbb{R}\,(0 < \varepsilon \to \lvert a - b \rvert < \varepsilon) \to a = b$
\item $\varphi_{4} :\equiv b \in \mathbb{R} \to \forall \varepsilon \in \mathbb{R}\,(0 < \varepsilon \to \lvert a - b \rvert < \varepsilon) \to a \neq b \to a - b \neq 0$
\item $\varphi_{5} :\equiv b \in \mathbb{R} \to \forall \varepsilon \in \mathbb{R}\,(0 < \varepsilon \to \lvert a - b \rvert < \varepsilon) \to a \neq b \to b \in \mathbb{R}$
\item $\varphi_{6} :\equiv b \in \mathbb{R} \to \forall \varepsilon \in \mathbb{R}\,(0 < \varepsilon \to \lvert a - b \rvert < \varepsilon) \to a \neq b \to a \in \mathbb{R}$
\item $\varphi_{7} :\equiv \forall x \in \mathbb{R}\,(x \neq 0 \to 0 < \lvert x \rvert) \to a - b \in \mathbb{R} \to a - b \neq 0 \to 0 < \lvert a - b \rvert$
\item $\varphi_{8} :\equiv \forall \varepsilon \in \mathbb{R}\,(0 < \varepsilon \to \lvert a - b \rvert < \varepsilon) \to a \neq b \to \lvert a - b \rvert < \lvert a - b \rvert$
\item $\varphi_{9} :\equiv \forall \varepsilon \in \mathbb{R}\,(0 < \varepsilon \to \lvert a - b \rvert < \varepsilon) \to a \neq b \to \neg \lvert a - b \rvert < \lvert a - b \rvert$
\item $\varphi_{10} :\equiv b \in \mathbb{R} \to \forall \varepsilon \in \mathbb{R}\,(0 < \varepsilon \to \lvert a - b \rvert < \varepsilon) \to a \neq b \to a - b \in \mathbb{R}$
\item $\varphi_{11} :\equiv \forall y\,(a \in \mathbb{R} \to y \in \mathbb{R} \to a - y \in \mathbb{R}) \to a \in \mathbb{R} \to b \in \mathbb{R} \to a - b \in \mathbb{R}$
\item $\varphi_{12} :\equiv b \in \mathbb{R} \to \forall \varepsilon \in \mathbb{R}\,(0 < \varepsilon \to \lvert a - b \rvert < \varepsilon) \to a \neq b \to 0 < \lvert a - b \rvert$
\item $\varphi_{13} :\equiv \forall y\,(a \in \mathbb{R} \to y \in \mathbb{R} \to a \neq y \to a - y \neq 0) \to a \in \mathbb{R} \to b \in \mathbb{R} \to a \neq b \to a - b \neq 0$
\item $\varphi_{14} :\equiv \forall x\,\forall y\,(x \in \mathbb{R} \to y \in \mathbb{R} \to x - y \in \mathbb{R}) \to \forall y\,(a \in \mathbb{R} \to y \in \mathbb{R} \to a - y \in \mathbb{R})$
\item $\varphi_{15} :\equiv \forall x\,\forall y\,(x \in \mathbb{R} \to y \in \mathbb{R} \to x \neq y \to x - y \neq 0) \to \forall y\,(a \in \mathbb{R} \to y \in \mathbb{R} \to a \neq y \to a - y \neq 0)$
\item $\varphi_{16} :\equiv \forall \varepsilon \in \mathbb{R}\,(0 < \varepsilon \to \lvert a - b \rvert < \varepsilon) \to \forall \varepsilon \in \mathbb{R}\,(0 < \varepsilon \to \lvert a - b \rvert < \varepsilon)$
\item $\varphi_{17} :\equiv \forall \varepsilon \in \mathbb{R}\,(0 < \varepsilon \to \lvert a - b \rvert < \varepsilon) \to \lvert a - b \rvert \in \mathbb{R} \to 0 < \lvert a - b \rvert \to \lvert a - b \rvert < \lvert a - b \rvert$
\item $\varphi_{18} :\equiv (a - b \in \mathbb{R} \to \lvert a - b \rvert \in \mathbb{R}) \to (a \in \mathbb{R} \to \varphi_{10}) \to a \in \mathbb{R} \to b \in \mathbb{R} \to \varphi_{2}$
\item $\varphi_{19} :\equiv (a \in \mathbb{R} \to b \in \mathbb{R} \to \varphi_{2}) \to (a \in \mathbb{R} \to \varphi_{12}) \to a \in \mathbb{R} \to b \in \mathbb{R} \to \varphi_{8}$
\item $\varphi_{20} :\equiv (a \in \mathbb{R} \to b \in \mathbb{R} \to a - b \in \mathbb{R}) \to (a \in \mathbb{R} \to \varphi_{6}) \to (a \in \mathbb{R} \to \varphi_{5}) \to a \in \mathbb{R} \to \varphi_{10}$
\item $\varphi_{21} :\equiv (a - b \in \mathbb{R} \to a - b \neq 0 \to 0 < \lvert a - b \rvert) \to (a \in \mathbb{R} \to \varphi_{10}) \to (a \in \mathbb{R} \to \varphi_{4}) \to a \in \mathbb{R} \to \varphi_{12}$
\item $\varphi_{22} :\equiv (a \in \mathbb{R} \to \varphi_{6}) \to (a \in \mathbb{R} \to \varphi_{5}) \to (a \in \mathbb{R} \to \varphi_{1}) \to a \in \mathbb{R} \to \varphi_{4}$
\end{abbrevlist}
\end{fullproof}

% 自動生成 (prover/build.js): ErrTwo  ノード数 151
\begin{fullproof}{[C10]}{$\forall x\,\forall \alpha\,\forall \beta\,\forall \varepsilon\,(x \in \mathbb{R} \to \alpha \in \mathbb{R} \to \beta \in \mathbb{R} \to \varepsilon \in \mathbb{R} \to \lvert x - \alpha \rvert < \varepsilon \to \lvert x - \beta \rvert < \varepsilon \to \lvert \alpha - \beta \rvert < \varepsilon + \varepsilon)$}{151}\phantomsection\label{tree:ErrTwo}
\noindent$\Pi_{1}$:
\begin{ptree}
\AXC{$\forall x\,\forall y\,(x \in \mathbb{R} \to y \in \mathbb{R} \to \lvert x - y \rvert = \lvert y - x \rvert)$\lbl{[A13]}}
\AXC{$\varphi_{40}$\lbl{A2}}
\BIC{$\forall y\,(x \in \mathbb{R} \to y \in \mathbb{R} \to \lvert x - y \rvert = \lvert y - x \rvert)$}
\AXC{$\varphi_{39}$\lbl{A2}}
\BIC{$x \in \mathbb{R} \to \alpha \in \mathbb{R} \to \lvert x - \alpha \rvert = \lvert \alpha - x \rvert$}
\end{ptree}
\noindent$\Pi_{2}$:
\begin{ptree}
\AXC{$x \in \mathbb{R} \to \varphi_{31}$\lbl{A1}}
\AXC{$\stackrel{\Pi_{1}}{\vdots}$}\noLine
\UIC{$x \in \mathbb{R} \to \alpha \in \mathbb{R} \to \lvert x - \alpha \rvert = \lvert \alpha - x \rvert$}
\AXC{$(x \in \mathbb{R} \to \alpha \in \mathbb{R} \to \lvert x - \alpha \rvert = \lvert \alpha - x \rvert) \to \varphi_{58}$\lbl{A1}}
\BIC{$\varphi_{58}$}
\BIC{$(x \in \mathbb{R} \to \varphi_{33}) \to x \in \mathbb{R} \to \alpha \in \mathbb{R} \to \beta \in \mathbb{R} \to \varphi_{11}$}
\end{ptree}
\noindent$\Pi_{3}$:
\begin{ptree}
\AXC{$x \in \mathbb{R} \to \alpha \in \mathbb{R} \to \varphi_{24}$\lbl{A1}}
\AXC{$x \in \mathbb{R} \to \varphi_{33}$\lbl{A1}}
\AXC{$\stackrel{\Pi_{2}}{\vdots}$}\noLine
\UIC{$(x \in \mathbb{R} \to \varphi_{33}) \to x \in \mathbb{R} \to \alpha \in \mathbb{R} \to \beta \in \mathbb{R} \to \varphi_{11}$}
\BIC{$x \in \mathbb{R} \to \alpha \in \mathbb{R} \to \beta \in \mathbb{R} \to \varphi_{11}$}
\AXC{$\varphi_{62}$\lbl{A1}}
\BIC{$(x \in \mathbb{R} \to \alpha \in \mathbb{R} \to \varphi_{24}) \to \varphi_{50}$}
\BIC{$\varphi_{50}$}
\end{ptree}
\noindent$\Pi_{4}$:
\begin{ptree}
\AXC{$z = w \to z < \varepsilon \to w < \varepsilon$\lbl{A6}}
\UIC{$\forall w\,(z = w \to z < \varepsilon \to w < \varepsilon)$}
\UIC{$\forall z\,\forall w\,(z = w \to z < \varepsilon \to w < \varepsilon)$}
\AXC{$\varphi_{9}$\lbl{A2}}
\BIC{$\forall w\,(\lvert x - \alpha \rvert = w \to \lvert x - \alpha \rvert < \varepsilon \to w < \varepsilon)$}
\AXC{$\varphi_{43}$\lbl{A2}}
\BIC{$\lvert x - \alpha \rvert = \lvert \alpha - x \rvert \to \lvert x - \alpha \rvert < \varepsilon \to \lvert \alpha - x \rvert < \varepsilon$}
\AXC{$\stackrel{\Pi_{3}}{\vdots}$}\noLine
\UIC{$\varphi_{50}$}
\BIC{$x \in \mathbb{R} \to \alpha \in \mathbb{R} \to \varphi_{25}$}
\end{ptree}
\noindent$\Pi_{5}$:
\begin{ptree}
\AXC{$x \in \mathbb{R} \to \varphi_{31}$\lbl{A1}}
\AXC{$\forall x\,\forall y\,(x \in \mathbb{R} \to y \in \mathbb{R} \to x - y \in \mathbb{R})$\lbl{[F1]}}
\AXC{$\varphi_{26}$\lbl{A2}}
\BIC{$\forall y\,(x \in \mathbb{R} \to y \in \mathbb{R} \to x - y \in \mathbb{R})$}
\AXC{$\varphi_{20}$\lbl{A2}}
\BIC{$x \in \mathbb{R} \to \beta \in \mathbb{R} \to x - \beta \in \mathbb{R}$}
\AXC{$\varphi_{60}$\lbl{A1}}
\BIC{$(x \in \mathbb{R} \to \varphi_{31}) \to (x \in \mathbb{R} \to \varphi_{34}) \to x \in \mathbb{R} \to \alpha \in \mathbb{R} \to \varphi_{18}$}
\BIC{$(x \in \mathbb{R} \to \varphi_{34}) \to x \in \mathbb{R} \to \alpha \in \mathbb{R} \to \varphi_{18}$}
\end{ptree}
\noindent$\Pi_{6}$:
\begin{ptree}
\AXC{$x \in \mathbb{R} \to \varphi_{34}$\lbl{A1}}
\AXC{$\stackrel{\Pi_{5}}{\vdots}$}\noLine
\UIC{$(x \in \mathbb{R} \to \varphi_{34}) \to x \in \mathbb{R} \to \alpha \in \mathbb{R} \to \varphi_{18}$}
\BIC{$x \in \mathbb{R} \to \alpha \in \mathbb{R} \to \varphi_{18}$}
\end{ptree}
\noindent$\Pi_{7}$:
\begin{ptree}
\AXC{$\forall x \in \mathbb{R}\,(\lvert x \rvert \in \mathbb{R})$\lbl{[A5]}}
\AXC{$\forall x \in \mathbb{R}\,(\lvert x \rvert \in \mathbb{R}) \to x - \beta \in \mathbb{R} \to \lvert x - \beta \rvert \in \mathbb{R}$\lbl{A2}}
\BIC{$x - \beta \in \mathbb{R} \to \lvert x - \beta \rvert \in \mathbb{R}$}
\AXC{$\varphi_{55}$\lbl{A1}}
\BIC{$(x \in \mathbb{R} \to \alpha \in \mathbb{R} \to \varphi_{18}) \to x \in \mathbb{R} \to \alpha \in \mathbb{R} \to \varphi_{27}$}
\end{ptree}
\noindent$\Pi_{8}$:
\begin{ptree}
\AXC{$x \in \mathbb{R} \to \varphi_{33}$\lbl{A1}}
\AXC{$\forall x\,\forall y\,(x \in \mathbb{R} \to y \in \mathbb{R} \to x - y \in \mathbb{R})$\lbl{[F1]}}
\AXC{$\varphi_{30}$\lbl{A2}}
\BIC{$\forall y\,(\alpha \in \mathbb{R} \to y \in \mathbb{R} \to \alpha - y \in \mathbb{R})$}
\AXC{$\varphi_{22}$\lbl{A2}}
\BIC{$\alpha \in \mathbb{R} \to x \in \mathbb{R} \to \alpha - x \in \mathbb{R}$}
\AXC{$\varphi_{61}$\lbl{A1}}
\BIC{$(x \in \mathbb{R} \to \varphi_{33}) \to (x \in \mathbb{R} \to \varphi_{31}) \to x \in \mathbb{R} \to \alpha \in \mathbb{R} \to \varphi_{19}$}
\BIC{$(x \in \mathbb{R} \to \varphi_{31}) \to x \in \mathbb{R} \to \alpha \in \mathbb{R} \to \varphi_{19}$}
\end{ptree}
\noindent$\Pi_{9}$:
\begin{ptree}
\AXC{$x \in \mathbb{R} \to \varphi_{31}$\lbl{A1}}
\AXC{$\stackrel{\Pi_{8}}{\vdots}$}\noLine
\UIC{$(x \in \mathbb{R} \to \varphi_{31}) \to x \in \mathbb{R} \to \alpha \in \mathbb{R} \to \varphi_{19}$}
\BIC{$x \in \mathbb{R} \to \alpha \in \mathbb{R} \to \varphi_{19}$}
\end{ptree}
\noindent$\Pi_{10}$:
\begin{ptree}
\AXC{$\forall x \in \mathbb{R}\,(\lvert x \rvert \in \mathbb{R})$\lbl{[A5]}}
\AXC{$\forall x \in \mathbb{R}\,(\lvert x \rvert \in \mathbb{R}) \to \alpha - x \in \mathbb{R} \to \lvert \alpha - x \rvert \in \mathbb{R}$\lbl{A2}}
\BIC{$\alpha - x \in \mathbb{R} \to \lvert \alpha - x \rvert \in \mathbb{R}$}
\AXC{$\varphi_{56}$\lbl{A1}}
\BIC{$(x \in \mathbb{R} \to \alpha \in \mathbb{R} \to \varphi_{19}) \to x \in \mathbb{R} \to \alpha \in \mathbb{R} \to \varphi_{28}$}
\end{ptree}
\noindent$\Pi_{11}$:
\begin{ptree}
\AXC{$\varphi_{29}$\lbl{[F20]}}
\AXC{$\varphi_{29} \to \varphi_{45}$\lbl{A2}}
\BIC{$\varphi_{45}$}
\AXC{$\varphi_{45} \to \varphi_{46}$\lbl{A2}}
\BIC{$\varphi_{46}$}
\AXC{$\varphi_{46} \to \varphi_{48}$\lbl{A2}}
\BIC{$\varphi_{48}$}
\AXC{$\varphi_{57}$\lbl{A2}}
\BIC{$\lvert \alpha - x \rvert \in \mathbb{R} \to \varepsilon \in \mathbb{R} \to \lvert x - \beta \rvert \in \mathbb{R} \to \varepsilon \in \mathbb{R} \to \varphi_{8}$}
\AXC{$\varphi_{66}$\lbl{A1}}
\BIC{$(x \in \mathbb{R} \to \alpha \in \mathbb{R} \to \varphi_{28}) \to (x \in \mathbb{R} \to \varphi_{35}) \to \varphi_{65}$}
\end{ptree}
\noindent$\Pi_{12}$:
\begin{ptree}
\AXC{$\stackrel{\Pi_{9}}{\vdots}$}\noLine
\UIC{$x \in \mathbb{R} \to \alpha \in \mathbb{R} \to \varphi_{19}$}
\AXC{$\stackrel{\Pi_{10}}{\vdots}$}\noLine
\UIC{$(x \in \mathbb{R} \to \alpha \in \mathbb{R} \to \varphi_{19}) \to x \in \mathbb{R} \to \alpha \in \mathbb{R} \to \varphi_{28}$}
\BIC{$x \in \mathbb{R} \to \alpha \in \mathbb{R} \to \varphi_{28}$}
\AXC{$\stackrel{\Pi_{11}}{\vdots}$}\noLine
\UIC{$(x \in \mathbb{R} \to \alpha \in \mathbb{R} \to \varphi_{28}) \to (x \in \mathbb{R} \to \varphi_{35}) \to \varphi_{65}$}
\BIC{$(x \in \mathbb{R} \to \varphi_{35}) \to \varphi_{65}$}
\end{ptree}
\noindent$\Pi_{13}$:
\begin{ptree}
\AXC{$\stackrel{\Pi_{6}}{\vdots}$}\noLine
\UIC{$x \in \mathbb{R} \to \alpha \in \mathbb{R} \to \varphi_{18}$}
\AXC{$\stackrel{\Pi_{7}}{\vdots}$}\noLine
\UIC{$(x \in \mathbb{R} \to \alpha \in \mathbb{R} \to \varphi_{18}) \to x \in \mathbb{R} \to \alpha \in \mathbb{R} \to \varphi_{27}$}
\BIC{$x \in \mathbb{R} \to \alpha \in \mathbb{R} \to \varphi_{27}$}
\AXC{$x \in \mathbb{R} \to \varphi_{35}$\lbl{A1}}
\AXC{$\stackrel{\Pi_{12}}{\vdots}$}\noLine
\UIC{$(x \in \mathbb{R} \to \varphi_{35}) \to \varphi_{65}$}
\BIC{$\varphi_{65}$}
\BIC{$(x \in \mathbb{R} \to \varphi_{35}) \to (x \in \mathbb{R} \to \alpha \in \mathbb{R} \to \varphi_{25}) \to \varphi_{51}$}
\end{ptree}
\noindent$\Pi_{14}$:
\begin{ptree}
\AXC{$\stackrel{\Pi_{4}}{\vdots}$}\noLine
\UIC{$x \in \mathbb{R} \to \alpha \in \mathbb{R} \to \varphi_{25}$}
\AXC{$x \in \mathbb{R} \to \varphi_{35}$\lbl{A1}}
\AXC{$\stackrel{\Pi_{13}}{\vdots}$}\noLine
\UIC{$(x \in \mathbb{R} \to \varphi_{35}) \to (x \in \mathbb{R} \to \alpha \in \mathbb{R} \to \varphi_{25}) \to \varphi_{51}$}
\BIC{$(x \in \mathbb{R} \to \alpha \in \mathbb{R} \to \varphi_{25}) \to \varphi_{51}$}
\BIC{$\varphi_{51}$}
\end{ptree}
\noindent$\Pi_{15}$:
\begin{ptree}
\AXC{$z = w \to \varphi_{4}$\lbl{A6}}
\UIC{$\forall w\,(z = w \to \varphi_{4})$}
\UIC{$\forall z\,\forall w\,(z = w \to \varphi_{4})$}
\AXC{$\forall z\,\forall w\,(z = w \to \varphi_{4}) \to \varphi_{41}$\lbl{A2}}
\BIC{$\varphi_{41}$}
\AXC{$\varphi_{41} \to \alpha - x + (x - \beta) = \alpha - \beta \to \varphi_{32}$\lbl{A2}}
\BIC{$\alpha - x + (x - \beta) = \alpha - \beta \to \varphi_{32}$}
\end{ptree}
\noindent$\Pi_{16}$:
\begin{ptree}
\AXC{$z = w \to \varphi_{7}$\lbl{A6}}
\UIC{$\forall w\,(z = w \to \varphi_{7})$}
\UIC{$\forall z\,\forall w\,(z = w \to \varphi_{7})$}
\AXC{$\forall z\,\forall w\,(z = w \to \varphi_{7}) \to \varphi_{44}$\lbl{A2}}
\BIC{$\varphi_{44}$}
\AXC{$\varphi_{44} \to \alpha - x + (x - \beta) = \alpha - \beta \to \varphi_{38}$\lbl{A2}}
\BIC{$\alpha - x + (x - \beta) = \alpha - \beta \to \varphi_{38}$}
\end{ptree}
\noindent$\Pi_{17}$:
\begin{ptree}
\AXC{$\forall a\,\forall b\,(a \in \mathbb{R} \to b \in \mathbb{R} \to \lvert a + b \rvert \leq \lvert a \rvert + \lvert b \rvert)$\lbl{命題1.2}}
\AXC{$\varphi_{47}$\lbl{A2}}
\BIC{$\varphi_{6}$}
\AXC{$\varphi_{6} \to \varphi_{21}$\lbl{A2}}
\BIC{$\varphi_{21}$}
\AXC{$\varphi_{21} \to (x \in \mathbb{R} \to \alpha \in \mathbb{R} \to \varphi_{19}) \to \varphi_{52}$\lbl{A1}}
\BIC{$(x \in \mathbb{R} \to \alpha \in \mathbb{R} \to \varphi_{19}) \to \varphi_{52}$}
\end{ptree}
\noindent$\Pi_{18}$:
\begin{ptree}
\AXC{$x = x$\lbl{A5}}
\AXC{$x = y \to x = x \to y = x$\lbl{A6}}
\AXC{$(x = y \to x = x \to y = x) \to x = x \to x = y \to y = x$\lbl{A1}}
\BIC{$x = x \to x = y \to y = x$}
\BIC{$x = y \to y = x$}
\UIC{$\forall y\,(x = y \to y = x)$}
\UIC{$\forall x\,\forall y\,(x = y \to y = x)$}
\end{ptree}
\noindent$\Pi_{19}$:
\begin{ptree}
\AXC{$\stackrel{\Pi_{18}}{\vdots}$}\noLine
\UIC{$\forall x\,\forall y\,(x = y \to y = x)$}
\AXC{$\forall x\,\forall y\,(x = y \to y = x) \to \forall y\,(\alpha - \beta = y \to y = \alpha - \beta)$\lbl{A2}}
\BIC{$\forall y\,(\alpha - \beta = y \to y = \alpha - \beta)$}
\AXC{$\varphi_{12}$\lbl{A2}}
\BIC{$\alpha - \beta = \alpha - x + (x - \beta) \to \alpha - x + (x - \beta) = \alpha - \beta$}
\end{ptree}
\noindent$\Pi_{20}$:
\begin{ptree}
\AXC{$\varphi_{1}$\lbl{[F16]}}
\AXC{$\varphi_{1} \to \varphi_{2}$\lbl{A2}}
\BIC{$\varphi_{2}$}
\AXC{$\varphi_{2} \to \varphi_{3}$\lbl{A2}}
\BIC{$\varphi_{3}$}
\AXC{$\varphi_{3} \to \alpha \in \mathbb{R} \to \beta \in \mathbb{R} \to x \in \mathbb{R} \to \alpha - \beta = \alpha - x + (x - \beta)$\lbl{A2}}
\BIC{$\alpha \in \mathbb{R} \to \beta \in \mathbb{R} \to x \in \mathbb{R} \to \alpha - \beta = \alpha - x + (x - \beta)$}
\AXC{$\varphi_{63}$\lbl{A1}}
\BIC{$(x \in \mathbb{R} \to \varphi_{33}) \to \varphi_{59}$}
\end{ptree}
\noindent$\Pi_{21}$:
\begin{ptree}
\AXC{$x \in \mathbb{R} \to \varphi_{31}$\lbl{A1}}
\AXC{$x \in \mathbb{R} \to \varphi_{34}$\lbl{A1}}
\AXC{$x \in \mathbb{R} \to \varphi_{33}$\lbl{A1}}
\AXC{$\stackrel{\Pi_{20}}{\vdots}$}\noLine
\UIC{$(x \in \mathbb{R} \to \varphi_{33}) \to \varphi_{59}$}
\BIC{$\varphi_{59}$}
\BIC{$(x \in \mathbb{R} \to \varphi_{31}) \to x \in \mathbb{R} \to \alpha \in \mathbb{R} \to \beta \in \mathbb{R} \to \varphi_{14}$}
\BIC{$x \in \mathbb{R} \to \alpha \in \mathbb{R} \to \beta \in \mathbb{R} \to \varphi_{14}$}
\end{ptree}
\noindent$\Pi_{22}$:
\begin{ptree}
\AXC{$\stackrel{\Pi_{21}}{\vdots}$}\noLine
\UIC{$x \in \mathbb{R} \to \alpha \in \mathbb{R} \to \beta \in \mathbb{R} \to \varphi_{14}$}
\AXC{$(x \in \mathbb{R} \to \alpha \in \mathbb{R} \to \beta \in \mathbb{R} \to \varphi_{14}) \to \varphi_{49}$\lbl{A1}}
\BIC{$\varphi_{49}$}
\end{ptree}
\noindent$\Pi_{23}$:
\begin{ptree}
\AXC{$\stackrel{\Pi_{19}}{\vdots}$}\noLine
\UIC{$\alpha - \beta = \alpha - x + (x - \beta) \to \alpha - x + (x - \beta) = \alpha - \beta$}
\AXC{$\stackrel{\Pi_{22}}{\vdots}$}\noLine
\UIC{$\varphi_{49}$}
\BIC{$x \in \mathbb{R} \to \alpha \in \mathbb{R} \to \beta \in \mathbb{R} \to \varphi_{15}$}
\AXC{$(x \in \mathbb{R} \to \alpha \in \mathbb{R} \to \beta \in \mathbb{R} \to \varphi_{15}) \to \varphi_{64}$\lbl{A1}}
\BIC{$\varphi_{64}$}
\end{ptree}
\noindent$\Pi_{24}$:
\begin{ptree}
\AXC{$\stackrel{\Pi_{6}}{\vdots}$}\noLine
\UIC{$x \in \mathbb{R} \to \alpha \in \mathbb{R} \to \varphi_{18}$}
\AXC{$\stackrel{\Pi_{9}}{\vdots}$}\noLine
\UIC{$x \in \mathbb{R} \to \alpha \in \mathbb{R} \to \varphi_{19}$}
\AXC{$\stackrel{\Pi_{17}}{\vdots}$}\noLine
\UIC{$(x \in \mathbb{R} \to \alpha \in \mathbb{R} \to \varphi_{19}) \to \varphi_{52}$}
\BIC{$\varphi_{52}$}
\BIC{$x \in \mathbb{R} \to \alpha \in \mathbb{R} \to \beta \in \mathbb{R} \to \varepsilon \in \mathbb{R} \to \lvert x - \alpha \rvert < \varepsilon \to \varphi_{5}$}
\AXC{$\stackrel{\Pi_{23}}{\vdots}$}\noLine
\UIC{$\varphi_{64}$}
\BIC{$(\alpha - x + (x - \beta) = \alpha - \beta \to \varphi_{38}) \to \varphi_{53}$}
\end{ptree}
\noindent$\Pi_{25}$:
\begin{ptree}
\AXC{$\stackrel{\Pi_{15}}{\vdots}$}\noLine
\UIC{$\alpha - x + (x - \beta) = \alpha - \beta \to \varphi_{32}$}
\AXC{$\stackrel{\Pi_{16}}{\vdots}$}\noLine
\UIC{$\alpha - x + (x - \beta) = \alpha - \beta \to \varphi_{38}$}
\AXC{$\stackrel{\Pi_{24}}{\vdots}$}\noLine
\UIC{$(\alpha - x + (x - \beta) = \alpha - \beta \to \varphi_{38}) \to \varphi_{53}$}
\BIC{$\varphi_{53}$}
\BIC{$x \in \mathbb{R} \to \alpha \in \mathbb{R} \to \beta \in \mathbb{R} \to \varepsilon \in \mathbb{R} \to \varphi_{16}$}
\end{ptree}
\noindent$\Pi_{26}$:
\begin{ptree}
\AXC{$\forall x\,\forall y\,\forall z\,(x \leq y \to y < z \to x < z)$\lbl{[O4]}}
\AXC{$\varphi_{10}$\lbl{A2}}
\BIC{$\forall y\,\forall z\,(\lvert \alpha - \beta \rvert \leq y \to y < z \to \lvert \alpha - \beta \rvert < z)$}
\AXC{$\forall y\,\forall z\,(\lvert \alpha - \beta \rvert \leq y \to y < z \to \lvert \alpha - \beta \rvert < z) \to \varphi_{37}$\lbl{A2}}
\BIC{$\varphi_{37}$}
\AXC{$\varphi_{37} \to \varphi_{42}$\lbl{A2}}
\BIC{$\varphi_{42}$}
\end{ptree}
\noindent$\Pi_{27}$:
\begin{ptree}
\AXC{$\stackrel{\Pi_{25}}{\vdots}$}\noLine
\UIC{$x \in \mathbb{R} \to \alpha \in \mathbb{R} \to \beta \in \mathbb{R} \to \varepsilon \in \mathbb{R} \to \varphi_{16}$}
\AXC{$\stackrel{\Pi_{26}}{\vdots}$}\noLine
\UIC{$\varphi_{42}$}
\AXC{$\varphi_{42} \to (x \in \mathbb{R} \to \alpha \in \mathbb{R} \to \beta \in \mathbb{R} \to \varepsilon \in \mathbb{R} \to \varphi_{16}) \to \varphi_{54}$\lbl{A1}}
\BIC{$(x \in \mathbb{R} \to \alpha \in \mathbb{R} \to \beta \in \mathbb{R} \to \varepsilon \in \mathbb{R} \to \varphi_{16}) \to \varphi_{54}$}
\BIC{$\varphi_{54}$}
\end{ptree}
\noindent 主木:
\begin{ptree}
\AXC{$x \in \mathbb{R} \to \alpha \in \mathbb{R} \to \varphi_{23}$\lbl{A1}}
\AXC{$(x \in \mathbb{R} \to \alpha \in \mathbb{R} \to \varphi_{23}) \to x \in \mathbb{R} \to \alpha \in \mathbb{R} \to \varphi_{17}$\lbl{A1}}
\BIC{$x \in \mathbb{R} \to \alpha \in \mathbb{R} \to \varphi_{17}$}
\AXC{$\stackrel{\Pi_{14}}{\vdots}$}\noLine
\UIC{$\varphi_{51}$}
\BIC{$x \in \mathbb{R} \to \alpha \in \mathbb{R} \to \beta \in \mathbb{R} \to \varepsilon \in \mathbb{R} \to \varphi_{13}$}
\AXC{$\stackrel{\Pi_{27}}{\vdots}$}\noLine
\UIC{$\varphi_{54}$}
\BIC{$x \in \mathbb{R} \to \alpha \in \mathbb{R} \to \varphi_{36}$}
\UIC{$\forall \varepsilon\,(x \in \mathbb{R} \to \alpha \in \mathbb{R} \to \varphi_{36})$}
\UIC{$\forall \beta\,\forall \varepsilon\,(x \in \mathbb{R} \to \alpha \in \mathbb{R} \to \varphi_{36})$}
\UIC{$\forall \alpha\,\forall \beta\,\forall \varepsilon\,(x \in \mathbb{R} \to \alpha \in \mathbb{R} \to \varphi_{36})$}
\UIC{$\forall x\,\forall \alpha\,\forall \beta\,\forall \varepsilon\,(x \in \mathbb{R} \to \alpha \in \mathbb{R} \to \varphi_{36})$}
\end{ptree}
\begin{abbrevlist}
\item $\varphi_{1} :\equiv \forall a\,\forall b\,\forall x\,(a \in \mathbb{R} \to b \in \mathbb{R} \to x \in \mathbb{R} \to a - b = a - x + (x - b))$
\item $\varphi_{2} :\equiv \forall b\,\forall x\,(\alpha \in \mathbb{R} \to b \in \mathbb{R} \to x \in \mathbb{R} \to \alpha - b = \alpha - x + (x - b))$
\item $\varphi_{3} :\equiv \forall x\,(\alpha \in \mathbb{R} \to \beta \in \mathbb{R} \to x \in \mathbb{R} \to \alpha - \beta = \alpha - x + (x - \beta))$
\item $\varphi_{4} :\equiv \lvert z \rvert = \lvert \alpha - x \rvert + \lvert x - \beta \rvert \to \lvert w \rvert = \lvert \alpha - x \rvert + \lvert x - \beta \rvert$
\item $\varphi_{5} :\equiv \lvert x - \beta \rvert < \varepsilon \to \lvert \alpha - x + (x - \beta) \rvert \leq \lvert \alpha - x \rvert + \lvert x - \beta \rvert$
\item $\varphi_{6} :\equiv \forall b\,(\alpha - x \in \mathbb{R} \to b \in \mathbb{R} \to \lvert \alpha - x + b \rvert \leq \lvert \alpha - x \rvert + \lvert b \rvert)$
\item $\varphi_{7} :\equiv \lvert z \rvert < \lvert \alpha - x \rvert + \lvert x - \beta \rvert \to \lvert w \rvert < \lvert \alpha - x \rvert + \lvert x - \beta \rvert$
\item $\varphi_{8} :\equiv \lvert \alpha - x \rvert < \varepsilon \to \lvert x - \beta \rvert \leq \varepsilon \to \lvert \alpha - x \rvert + \lvert x - \beta \rvert < \varepsilon + \varepsilon$
\item $\varphi_{9} :\equiv \forall z\,\forall w\,(z = w \to z < \varepsilon \to w < \varepsilon) \to \forall w\,(\lvert x - \alpha \rvert = w \to \lvert x - \alpha \rvert < \varepsilon \to w < \varepsilon)$
\item $\varphi_{10} :\equiv \forall x\,\forall y\,\forall z\,(x \leq y \to y < z \to x < z) \to \forall y\,\forall z\,(\lvert \alpha - \beta \rvert \leq y \to y < z \to \lvert \alpha - \beta \rvert < z)$
\item $\varphi_{11} :\equiv \varepsilon \in \mathbb{R} \to \lvert x - \alpha \rvert < \varepsilon \to \lvert x - \beta \rvert < \varepsilon \to \lvert x - \alpha \rvert = \lvert \alpha - x \rvert$
\item $\varphi_{12} :\equiv \forall y\,(\alpha - \beta = y \to y = \alpha - \beta) \to \alpha - \beta = \alpha - x + (x - \beta) \to \alpha - x + (x - \beta) = \alpha - \beta$
\item $\varphi_{13} :\equiv \lvert x - \alpha \rvert < \varepsilon \to \lvert x - \beta \rvert < \varepsilon \to \lvert \alpha - x \rvert + \lvert x - \beta \rvert < \varepsilon + \varepsilon$
\item $\varphi_{14} :\equiv \varepsilon \in \mathbb{R} \to \lvert x - \alpha \rvert < \varepsilon \to \lvert x - \beta \rvert < \varepsilon \to \alpha - \beta = \alpha - x + (x - \beta)$
\item $\varphi_{15} :\equiv \varepsilon \in \mathbb{R} \to \lvert x - \alpha \rvert < \varepsilon \to \lvert x - \beta \rvert < \varepsilon \to \alpha - x + (x - \beta) = \alpha - \beta$
\item $\varphi_{16} :\equiv \lvert x - \alpha \rvert < \varepsilon \to \lvert x - \beta \rvert < \varepsilon \to \lvert \alpha - \beta \rvert \leq \lvert \alpha - x \rvert + \lvert x - \beta \rvert$
\item $\varphi_{17} :\equiv \beta \in \mathbb{R} \to \varepsilon \in \mathbb{R} \to \lvert x - \alpha \rvert < \varepsilon \to \lvert x - \beta \rvert < \varepsilon \to \lvert x - \beta \rvert \leq \varepsilon$
\item $\varphi_{18} :\equiv \beta \in \mathbb{R} \to \varepsilon \in \mathbb{R} \to \lvert x - \alpha \rvert < \varepsilon \to \lvert x - \beta \rvert < \varepsilon \to x - \beta \in \mathbb{R}$
\item $\varphi_{19} :\equiv \beta \in \mathbb{R} \to \varepsilon \in \mathbb{R} \to \lvert x - \alpha \rvert < \varepsilon \to \lvert x - \beta \rvert < \varepsilon \to \alpha - x \in \mathbb{R}$
\item $\varphi_{20} :\equiv \forall y\,(x \in \mathbb{R} \to y \in \mathbb{R} \to x - y \in \mathbb{R}) \to x \in \mathbb{R} \to \beta \in \mathbb{R} \to x - \beta \in \mathbb{R}$
\item $\varphi_{21} :\equiv \alpha - x \in \mathbb{R} \to x - \beta \in \mathbb{R} \to \lvert \alpha - x + (x - \beta) \rvert \leq \lvert \alpha - x \rvert + \lvert x - \beta \rvert$
\item $\varphi_{22} :\equiv \forall y\,(\alpha \in \mathbb{R} \to y \in \mathbb{R} \to \alpha - y \in \mathbb{R}) \to \alpha \in \mathbb{R} \to x \in \mathbb{R} \to \alpha - x \in \mathbb{R}$
\item $\varphi_{23} :\equiv \beta \in \mathbb{R} \to \varepsilon \in \mathbb{R} \to \lvert x - \alpha \rvert < \varepsilon \to \lvert x - \beta \rvert < \varepsilon \to \lvert x - \beta \rvert < \varepsilon$
\item $\varphi_{24} :\equiv \beta \in \mathbb{R} \to \varepsilon \in \mathbb{R} \to \lvert x - \alpha \rvert < \varepsilon \to \lvert x - \beta \rvert < \varepsilon \to \lvert x - \alpha \rvert < \varepsilon$
\item $\varphi_{25} :\equiv \beta \in \mathbb{R} \to \varepsilon \in \mathbb{R} \to \lvert x - \alpha \rvert < \varepsilon \to \lvert x - \beta \rvert < \varepsilon \to \lvert \alpha - x \rvert < \varepsilon$
\item $\varphi_{26} :\equiv \forall x\,\forall y\,(x \in \mathbb{R} \to y \in \mathbb{R} \to x - y \in \mathbb{R}) \to \forall y\,(x \in \mathbb{R} \to y \in \mathbb{R} \to x - y \in \mathbb{R})$
\item $\varphi_{27} :\equiv \beta \in \mathbb{R} \to \varepsilon \in \mathbb{R} \to \lvert x - \alpha \rvert < \varepsilon \to \lvert x - \beta \rvert < \varepsilon \to \lvert x - \beta \rvert \in \mathbb{R}$
\item $\varphi_{28} :\equiv \beta \in \mathbb{R} \to \varepsilon \in \mathbb{R} \to \lvert x - \alpha \rvert < \varepsilon \to \lvert x - \beta \rvert < \varepsilon \to \lvert \alpha - x \rvert \in \mathbb{R}$
\item $\varphi_{29} :\equiv \forall x\,\forall y\,\forall u\,\forall v\,(x \in \mathbb{R} \to y \in \mathbb{R} \to u \in \mathbb{R} \to v \in \mathbb{R} \to x < y \to u \leq v \to x + u < y + v)$
\item $\varphi_{30} :\equiv \forall x\,\forall y\,(x \in \mathbb{R} \to y \in \mathbb{R} \to x - y \in \mathbb{R}) \to \forall y\,(\alpha \in \mathbb{R} \to y \in \mathbb{R} \to \alpha - y \in \mathbb{R})$
\item $\varphi_{31} :\equiv \alpha \in \mathbb{R} \to \beta \in \mathbb{R} \to \varepsilon \in \mathbb{R} \to \lvert x - \alpha \rvert < \varepsilon \to \lvert x - \beta \rvert < \varepsilon \to x \in \mathbb{R}$
\item $\varphi_{32} :\equiv \lvert \alpha - x + (x - \beta) \rvert = \lvert \alpha - x \rvert + \lvert x - \beta \rvert \to \lvert \alpha - \beta \rvert = \lvert \alpha - x \rvert + \lvert x - \beta \rvert$
\item $\varphi_{33} :\equiv \alpha \in \mathbb{R} \to \beta \in \mathbb{R} \to \varepsilon \in \mathbb{R} \to \lvert x - \alpha \rvert < \varepsilon \to \lvert x - \beta \rvert < \varepsilon \to \alpha \in \mathbb{R}$
\item $\varphi_{34} :\equiv \alpha \in \mathbb{R} \to \beta \in \mathbb{R} \to \varepsilon \in \mathbb{R} \to \lvert x - \alpha \rvert < \varepsilon \to \lvert x - \beta \rvert < \varepsilon \to \beta \in \mathbb{R}$
\item $\varphi_{35} :\equiv \alpha \in \mathbb{R} \to \beta \in \mathbb{R} \to \varepsilon \in \mathbb{R} \to \lvert x - \alpha \rvert < \varepsilon \to \lvert x - \beta \rvert < \varepsilon \to \varepsilon \in \mathbb{R}$
\item $\varphi_{36} :\equiv \beta \in \mathbb{R} \to \varepsilon \in \mathbb{R} \to \lvert x - \alpha \rvert < \varepsilon \to \lvert x - \beta \rvert < \varepsilon \to \lvert \alpha - \beta \rvert < \varepsilon + \varepsilon$
\item $\varphi_{37} :\equiv \forall z\,(\lvert \alpha - \beta \rvert \leq \lvert \alpha - x \rvert + \lvert x - \beta \rvert \to \lvert \alpha - x \rvert + \lvert x - \beta \rvert < z \to \lvert \alpha - \beta \rvert < z)$
\item $\varphi_{38} :\equiv \lvert \alpha - x + (x - \beta) \rvert < \lvert \alpha - x \rvert + \lvert x - \beta \rvert \to \lvert \alpha - \beta \rvert < \lvert \alpha - x \rvert + \lvert x - \beta \rvert$
\item $\varphi_{39} :\equiv \forall y\,(x \in \mathbb{R} \to y \in \mathbb{R} \to \lvert x - y \rvert = \lvert y - x \rvert) \to x \in \mathbb{R} \to \alpha \in \mathbb{R} \to \lvert x - \alpha \rvert = \lvert \alpha - x \rvert$
\item $\varphi_{40} :\equiv \forall x\,\forall y\,(x \in \mathbb{R} \to y \in \mathbb{R} \to \lvert x - y \rvert = \lvert y - x \rvert) \to \forall y\,(x \in \mathbb{R} \to y \in \mathbb{R} \to \lvert x - y \rvert = \lvert y - x \rvert)$
\item $\varphi_{41} :\equiv \forall w\,(\alpha - x + (x - \beta) = w \to \lvert \alpha - x + (x - \beta) \rvert = \lvert \alpha - x \rvert + \lvert x - \beta \rvert \to \lvert w \rvert = \lvert \alpha - x \rvert + \lvert x - \beta \rvert)$
\item $\varphi_{42} :\equiv \lvert \alpha - \beta \rvert \leq \lvert \alpha - x \rvert + \lvert x - \beta \rvert \to \lvert \alpha - x \rvert + \lvert x - \beta \rvert < \varepsilon + \varepsilon \to \lvert \alpha - \beta \rvert < \varepsilon + \varepsilon$
\item $\varphi_{43} :\equiv \forall w\,(\lvert x - \alpha \rvert = w \to \lvert x - \alpha \rvert < \varepsilon \to w < \varepsilon) \to \lvert x - \alpha \rvert = \lvert \alpha - x \rvert \to \lvert x - \alpha \rvert < \varepsilon \to \lvert \alpha - x \rvert < \varepsilon$
\item $\varphi_{44} :\equiv \forall w\,(\alpha - x + (x - \beta) = w \to \lvert \alpha - x + (x - \beta) \rvert < \lvert \alpha - x \rvert + \lvert x - \beta \rvert \to \lvert w \rvert < \lvert \alpha - x \rvert + \lvert x - \beta \rvert)$
\item $\varphi_{45} :\equiv \forall y\,\forall u\,\forall v\,(\lvert \alpha - x \rvert \in \mathbb{R} \to y \in \mathbb{R} \to u \in \mathbb{R} \to v \in \mathbb{R} \to \lvert \alpha - x \rvert < y \to u \leq v \to \lvert \alpha - x \rvert + u < y + v)$
\item $\varphi_{46} :\equiv \forall u\,\forall v\,(\lvert \alpha - x \rvert \in \mathbb{R} \to \varepsilon \in \mathbb{R} \to u \in \mathbb{R} \to v \in \mathbb{R} \to \lvert \alpha - x \rvert < \varepsilon \to u \leq v \to \lvert \alpha - x \rvert + u < \varepsilon + v)$
\item $\varphi_{47} :\equiv \forall a\,\forall b\,(a \in \mathbb{R} \to b \in \mathbb{R} \to \lvert a + b \rvert \leq \lvert a \rvert + \lvert b \rvert) \to \varphi_{6}$
\item $\varphi_{48} :\equiv \forall v\,(\lvert \alpha - x \rvert \in \mathbb{R} \to \varepsilon \in \mathbb{R} \to \lvert x - \beta \rvert \in \mathbb{R} \to v \in \mathbb{R} \to \lvert \alpha - x \rvert < \varepsilon \to \lvert x - \beta \rvert \leq v \to \lvert \alpha - x \rvert + \lvert x - \beta \rvert < \varepsilon + v)$
\item $\varphi_{49} :\equiv (\alpha - \beta = \alpha - x + (x - \beta) \to \alpha - x + (x - \beta) = \alpha - \beta) \to x \in \mathbb{R} \to \alpha \in \mathbb{R} \to \beta \in \mathbb{R} \to \varphi_{15}$
\item $\varphi_{50} :\equiv (\lvert x - \alpha \rvert = \lvert \alpha - x \rvert \to \lvert x - \alpha \rvert < \varepsilon \to \lvert \alpha - x \rvert < \varepsilon) \to x \in \mathbb{R} \to \alpha \in \mathbb{R} \to \varphi_{25}$
\item $\varphi_{51} :\equiv (x \in \mathbb{R} \to \alpha \in \mathbb{R} \to \varphi_{17}) \to x \in \mathbb{R} \to \alpha \in \mathbb{R} \to \beta \in \mathbb{R} \to \varepsilon \in \mathbb{R} \to \varphi_{13}$
\item $\varphi_{52} :\equiv (x \in \mathbb{R} \to \alpha \in \mathbb{R} \to \varphi_{18}) \to x \in \mathbb{R} \to \alpha \in \mathbb{R} \to \beta \in \mathbb{R} \to \varepsilon \in \mathbb{R} \to \lvert x - \alpha \rvert < \varepsilon \to \varphi_{5}$
\item $\varphi_{53} :\equiv (\alpha - x + (x - \beta) = \alpha - \beta \to \varphi_{32}) \to x \in \mathbb{R} \to \alpha \in \mathbb{R} \to \beta \in \mathbb{R} \to \varepsilon \in \mathbb{R} \to \varphi_{16}$
\item $\varphi_{54} :\equiv (x \in \mathbb{R} \to \alpha \in \mathbb{R} \to \beta \in \mathbb{R} \to \varepsilon \in \mathbb{R} \to \varphi_{13}) \to x \in \mathbb{R} \to \alpha \in \mathbb{R} \to \varphi_{36}$
\item $\varphi_{55} :\equiv (x - \beta \in \mathbb{R} \to \lvert x - \beta \rvert \in \mathbb{R}) \to (x \in \mathbb{R} \to \alpha \in \mathbb{R} \to \varphi_{18}) \to x \in \mathbb{R} \to \alpha \in \mathbb{R} \to \varphi_{27}$
\item $\varphi_{56} :\equiv (\alpha - x \in \mathbb{R} \to \lvert \alpha - x \rvert \in \mathbb{R}) \to (x \in \mathbb{R} \to \alpha \in \mathbb{R} \to \varphi_{19}) \to x \in \mathbb{R} \to \alpha \in \mathbb{R} \to \varphi_{28}$
\item $\varphi_{57} :\equiv \varphi_{48} \to \lvert \alpha - x \rvert \in \mathbb{R} \to \varepsilon \in \mathbb{R} \to \lvert x - \beta \rvert \in \mathbb{R} \to \varepsilon \in \mathbb{R} \to \varphi_{8}$
\item $\varphi_{58} :\equiv (x \in \mathbb{R} \to \varphi_{31}) \to (x \in \mathbb{R} \to \varphi_{33}) \to x \in \mathbb{R} \to \alpha \in \mathbb{R} \to \beta \in \mathbb{R} \to \varphi_{11}$
\item $\varphi_{59} :\equiv (x \in \mathbb{R} \to \varphi_{34}) \to (x \in \mathbb{R} \to \varphi_{31}) \to x \in \mathbb{R} \to \alpha \in \mathbb{R} \to \beta \in \mathbb{R} \to \varphi_{14}$
\item $\varphi_{60} :\equiv (x \in \mathbb{R} \to \beta \in \mathbb{R} \to x - \beta \in \mathbb{R}) \to (x \in \mathbb{R} \to \varphi_{31}) \to (x \in \mathbb{R} \to \varphi_{34}) \to x \in \mathbb{R} \to \alpha \in \mathbb{R} \to \varphi_{18}$
\item $\varphi_{61} :\equiv (\alpha \in \mathbb{R} \to x \in \mathbb{R} \to \alpha - x \in \mathbb{R}) \to (x \in \mathbb{R} \to \varphi_{33}) \to (x \in \mathbb{R} \to \varphi_{31}) \to x \in \mathbb{R} \to \alpha \in \mathbb{R} \to \varphi_{19}$
\item $\varphi_{62} :\equiv (x \in \mathbb{R} \to \alpha \in \mathbb{R} \to \beta \in \mathbb{R} \to \varphi_{11}) \to (x \in \mathbb{R} \to \alpha \in \mathbb{R} \to \varphi_{24}) \to \varphi_{50}$
\item $\varphi_{63} :\equiv (\alpha \in \mathbb{R} \to \beta \in \mathbb{R} \to x \in \mathbb{R} \to \alpha - \beta = \alpha - x + (x - \beta)) \to (x \in \mathbb{R} \to \varphi_{33}) \to \varphi_{59}$
\item $\varphi_{64} :\equiv (x \in \mathbb{R} \to \alpha \in \mathbb{R} \to \beta \in \mathbb{R} \to \varepsilon \in \mathbb{R} \to \lvert x - \alpha \rvert < \varepsilon \to \varphi_{5}) \to (\alpha - x + (x - \beta) = \alpha - \beta \to \varphi_{38}) \to \varphi_{53}$
\item $\varphi_{65} :\equiv (x \in \mathbb{R} \to \alpha \in \mathbb{R} \to \varphi_{27}) \to (x \in \mathbb{R} \to \varphi_{35}) \to (x \in \mathbb{R} \to \alpha \in \mathbb{R} \to \varphi_{25}) \to \varphi_{51}$
\item $\varphi_{66} :\equiv (\lvert \alpha - x \rvert \in \mathbb{R} \to \varepsilon \in \mathbb{R} \to \lvert x - \beta \rvert \in \mathbb{R} \to \varepsilon \in \mathbb{R} \to \varphi_{8}) \to (x \in \mathbb{R} \to \alpha \in \mathbb{R} \to \varphi_{28}) \to (x \in \mathbb{R} \to \varphi_{35}) \to \varphi_{65}$
\end{abbrevlist}
\end{fullproof}

% 自動生成 (prover/build.js): LimUnique  ノード数 209
\begin{fullproof}{命題1.3}{$\forall a\,\forall \alpha\,\forall \beta\,(\mathrm{Conv}(a,\alpha) \to \mathrm{Conv}(a,\beta) \to \alpha = \beta)$}{209}\phantomsection\label{tree:LimUnique}
\noindent$\Pi_{1}$:
\begin{ptree}
\AXC{$z = w \to \lvert \alpha - \beta \rvert < z \to \lvert \alpha - \beta \rvert < w$\lbl{A6}}
\UIC{$\forall w\,(z = w \to \lvert \alpha - \beta \rvert < z \to \lvert \alpha - \beta \rvert < w)$}
\UIC{$\forall z\,\forall w\,(z = w \to \lvert \alpha - \beta \rvert < z \to \lvert \alpha - \beta \rvert < w)$}
\AXC{$\forall z\,\forall w\,(z = w \to \lvert \alpha - \beta \rvert < z \to \lvert \alpha - \beta \rvert < w) \to \varphi_{38}$\lbl{A2}}
\BIC{$\varphi_{38}$}
\end{ptree}
\noindent$\Pi_{2}$:
\begin{ptree}
\AXC{$\varphi_{31} \to N \in \mathbb{N} \to N \leq N \to \varphi_{24}$\lbl{A2}}
\AXC{$\varphi_{67}$\lbl{A1}}
\AXC{$\varphi_{67} \to \varphi_{66}$\lbl{A1}}
\BIC{$\varphi_{66}$}
\AXC{$\varphi_{66} \to (\varphi_{31} \to N \in \mathbb{N} \to N \leq N \to \varphi_{24}) \to \mathrm{Conv}(a,\alpha) \to \varphi_{65}$\lbl{A1}}
\BIC{$(\varphi_{31} \to N \in \mathbb{N} \to N \leq N \to \varphi_{24}) \to \mathrm{Conv}(a,\alpha) \to \varphi_{65}$}
\BIC{$\mathrm{Conv}(a,\alpha) \to \varphi_{65}$}
\end{ptree}
\noindent$\Pi_{3}$:
\begin{ptree}
\AXC{$\varphi_{67}$\lbl{A1}}
\AXC{$\varphi_{67} \to \varphi_{48}$\lbl{A1}}
\BIC{$\varphi_{48}$}
\AXC{$\stackrel{\Pi_{2}}{\vdots}$}\noLine
\UIC{$\mathrm{Conv}(a,\alpha) \to \varphi_{65}$}
\AXC{$(\mathrm{Conv}(a,\alpha) \to \varphi_{65}) \to \varphi_{48} \to N = N \to \varphi_{63}$\lbl{A1}}
\BIC{$\varphi_{48} \to N = N \to \varphi_{63}$}
\BIC{$N = N \to \varphi_{63}$}
\end{ptree}
\noindent$\Pi_{4}$:
\begin{ptree}
\AXC{$x = x$\lbl{A5}}
\UIC{$\forall x\,(x = x)$}
\AXC{$\forall x\,(x = x) \to N = N$\lbl{A2}}
\BIC{$N = N$}
\AXC{$\stackrel{\Pi_{3}}{\vdots}$}\noLine
\UIC{$N = N \to \varphi_{63}$}
\BIC{$\varphi_{63}$}
\end{ptree}
\noindent$\Pi_{5}$:
\begin{ptree}
\AXC{$\forall a\,\forall \alpha\,(\mathrm{Conv}(a,\alpha) \to \alpha \in \mathbb{R})$\lbl{[C2]}}
\AXC{$\varphi_{11}$\lbl{A2}}
\BIC{$\forall \alpha\,(\mathrm{Conv}(a,\alpha) \to \alpha \in \mathbb{R})$}
\AXC{$\varphi_{7}$\lbl{A2}}
\BIC{$\mathrm{Conv}(a,\beta) \to \beta \in \mathbb{R}$}
\AXC{$(\mathrm{Conv}(a,\beta) \to \beta \in \mathbb{R}) \to \varphi_{15}$\lbl{A1}}
\BIC{$\varphi_{15}$}
\end{ptree}
\noindent$\Pi_{6}$:
\begin{ptree}
\AXC{$\mathrm{Conv}(a,\alpha) \to \mathrm{Conv}(a,\beta) \to \mathrm{Conv}(a,\beta)$\lbl{A1}}
\AXC{$\stackrel{\Pi_{5}}{\vdots}$}\noLine
\UIC{$\varphi_{15}$}
\BIC{$\mathrm{Conv}(a,\alpha) \to \mathrm{Conv}(a,\beta) \to \beta \in \mathbb{R}$}
\end{ptree}
\noindent$\Pi_{7}$:
\begin{ptree}
\AXC{$\forall a\,\forall \alpha\,(\mathrm{Conv}(a,\alpha) \to \alpha \in \mathbb{R})$\lbl{[C2]}}
\AXC{$\varphi_{11}$\lbl{A2}}
\BIC{$\forall \alpha\,(\mathrm{Conv}(a,\alpha) \to \alpha \in \mathbb{R})$}
\AXC{$\varphi_{8}$\lbl{A2}}
\BIC{$\mathrm{Conv}(a,\alpha) \to \alpha \in \mathbb{R}$}
\AXC{$(\mathrm{Conv}(a,\alpha) \to \alpha \in \mathbb{R}) \to \varphi_{12}$\lbl{A1}}
\BIC{$\varphi_{12}$}
\end{ptree}
\noindent$\Pi_{8}$:
\begin{ptree}
\AXC{$\mathrm{Conv}(a,\alpha) \to \mathrm{Conv}(a,\alpha)$\lbl{A1}}
\AXC{$\stackrel{\Pi_{7}}{\vdots}$}\noLine
\UIC{$\varphi_{12}$}
\BIC{$\mathrm{Conv}(a,\alpha) \to \mathrm{Conv}(a,\beta) \to \alpha \in \mathbb{R}$}
\end{ptree}
\noindent$\Pi_{9}$:
\begin{ptree}
\AXC{$\forall x\,\forall y\,(x \in \mathbb{R} \to y \in \mathbb{R} \to x - y \in \mathbb{R})$\lbl{[F1]}}
\AXC{$\varphi_{25}$\lbl{A2}}
\BIC{$\forall y\,(\alpha \in \mathbb{R} \to y \in \mathbb{R} \to \alpha - y \in \mathbb{R})$}
\AXC{$\varphi_{23}$\lbl{A2}}
\BIC{$\alpha \in \mathbb{R} \to \beta \in \mathbb{R} \to \alpha - \beta \in \mathbb{R}$}
\AXC{$(\alpha \in \mathbb{R} \to \beta \in \mathbb{R} \to \alpha - \beta \in \mathbb{R}) \to \varphi_{64}$\lbl{A1}}
\BIC{$\varphi_{64}$}
\end{ptree}
\noindent$\Pi_{10}$:
\begin{ptree}
\AXC{$\stackrel{\Pi_{6}}{\vdots}$}\noLine
\UIC{$\mathrm{Conv}(a,\alpha) \to \mathrm{Conv}(a,\beta) \to \beta \in \mathbb{R}$}
\AXC{$\stackrel{\Pi_{8}}{\vdots}$}\noLine
\UIC{$\mathrm{Conv}(a,\alpha) \to \mathrm{Conv}(a,\beta) \to \alpha \in \mathbb{R}$}
\AXC{$\stackrel{\Pi_{9}}{\vdots}$}\noLine
\UIC{$\varphi_{64}$}
\BIC{$(\mathrm{Conv}(a,\alpha) \to \mathrm{Conv}(a,\beta) \to \beta \in \mathbb{R}) \to \varphi_{50}$}
\BIC{$\varphi_{50}$}
\end{ptree}
\noindent$\Pi_{11}$:
\begin{ptree}
\AXC{$\forall x \in \mathbb{R}\,(\lvert x \rvert \in \mathbb{R})$\lbl{[A5]}}
\AXC{$\forall x \in \mathbb{R}\,(\lvert x \rvert \in \mathbb{R}) \to \alpha - \beta \in \mathbb{R} \to \lvert \alpha - \beta \rvert \in \mathbb{R}$\lbl{A2}}
\BIC{$\alpha - \beta \in \mathbb{R} \to \lvert \alpha - \beta \rvert \in \mathbb{R}$}
\end{ptree}
\noindent$\Pi_{12}$:
\begin{ptree}
\AXC{$\stackrel{\Pi_{10}}{\vdots}$}\noLine
\UIC{$\varphi_{50}$}
\AXC{$\stackrel{\Pi_{11}}{\vdots}$}\noLine
\UIC{$\alpha - \beta \in \mathbb{R} \to \lvert \alpha - \beta \rvert \in \mathbb{R}$}
\AXC{$(\alpha - \beta \in \mathbb{R} \to \lvert \alpha - \beta \rvert \in \mathbb{R}) \to \varphi_{50} \to \varphi_{54}$\lbl{A1}}
\BIC{$\varphi_{50} \to \varphi_{54}$}
\BIC{$\varphi_{54}$}
\end{ptree}
\noindent$\Pi_{13}$:
\begin{ptree}
\AXC{$2^{-1} \in \mathbb{R}$\lbl{[F31]}}
\AXC{$\stackrel{\Pi_{12}}{\vdots}$}\noLine
\UIC{$\varphi_{54}$}
\AXC{$\forall x\,\forall y\,(x \in \mathbb{R} \to y \in \mathbb{R} \to x \cdot y \in \mathbb{R})$\lbl{(R1)}}
\AXC{$\varphi_{33}$\lbl{A2}}
\BIC{$\forall y\,(\lvert \alpha - \beta \rvert \in \mathbb{R} \to y \in \mathbb{R} \to \lvert \alpha - \beta \rvert \cdot y \in \mathbb{R})$}
\AXC{$\varphi_{39}$\lbl{A2}}
\BIC{$\lvert \alpha - \beta \rvert \in \mathbb{R} \to 2^{-1} \in \mathbb{R} \to \lvert \alpha - \beta \rvert \cdot 2^{-1} \in \mathbb{R}$}
\AXC{$\varphi_{72}$\lbl{A1}}
\BIC{$\varphi_{54} \to 2^{-1} \in \mathbb{R} \to \mathrm{Conv}(a,\alpha) \to \varphi_{53}$}
\BIC{$2^{-1} \in \mathbb{R} \to \mathrm{Conv}(a,\alpha) \to \varphi_{53}$}
\BIC{$\mathrm{Conv}(a,\alpha) \to \varphi_{53}$}
\end{ptree}
\noindent$\Pi_{14}$:
\begin{ptree}
\AXC{$\forall a\,\forall \alpha\,(\mathrm{Conv}(a,\alpha) \to \mathrm{Seq}(a))$\lbl{[C1]}}
\AXC{$\varphi_{5}$\lbl{A2}}
\BIC{$\forall \alpha\,(\mathrm{Conv}(a,\alpha) \to \mathrm{Seq}(a))$}
\AXC{$\varphi_{2}$\lbl{A2}}
\BIC{$\mathrm{Conv}(a,\alpha) \to \mathrm{Seq}(a)$}
\AXC{$(\mathrm{Conv}(a,\alpha) \to \mathrm{Seq}(a)) \to \varphi_{10}$\lbl{A1}}
\BIC{$\varphi_{10}$}
\end{ptree}
\noindent$\Pi_{15}$:
\begin{ptree}
\AXC{$\forall a\,\forall n\,(\mathrm{Seq}(a) \to n \in \mathbb{N} \to a_{n} \in \mathbb{R})$\lbl{[S1]}}
\AXC{$\varphi_{18}$\lbl{A2}}
\BIC{$\forall n\,(\mathrm{Seq}(a) \to n \in \mathbb{N} \to a_{n} \in \mathbb{R})$}
\AXC{$\varphi_{17}$\lbl{A2}}
\BIC{$\mathrm{Seq}(a) \to N \in \mathbb{N} \to a_{N} \in \mathbb{R}$}
\AXC{$(\mathrm{Seq}(a) \to N \in \mathbb{N} \to a_{N} \in \mathbb{R}) \to \varphi_{70}$\lbl{A1}}
\BIC{$\varphi_{70}$}
\end{ptree}
\noindent$\Pi_{16}$:
\begin{ptree}
\AXC{$\varphi_{67}$\lbl{A1}}
\AXC{$\varphi_{67} \to \varphi_{48}$\lbl{A1}}
\BIC{$\varphi_{48}$}
\AXC{$\mathrm{Conv}(a,\alpha) \to \mathrm{Conv}(a,\alpha)$\lbl{A1}}
\AXC{$\stackrel{\Pi_{14}}{\vdots}$}\noLine
\UIC{$\varphi_{10}$}
\BIC{$\mathrm{Conv}(a,\alpha) \to \mathrm{Conv}(a,\beta) \to \mathrm{Seq}(a)$}
\AXC{$\stackrel{\Pi_{15}}{\vdots}$}\noLine
\UIC{$\varphi_{70}$}
\BIC{$\varphi_{48} \to \varphi_{51}$}
\BIC{$\varphi_{51}$}
\end{ptree}
\noindent$\Pi_{17}$:
\begin{ptree}
\AXC{$\varphi_{37}$\lbl{[C10]}}
\AXC{$\varphi_{37} \to \forall \alpha\,\forall \beta\,(\varphi_{41})$\lbl{A2}}
\BIC{$\forall \alpha\,\forall \beta\,(\varphi_{41})$}
\AXC{$\forall \alpha\,\forall \beta\,(\varphi_{41}) \to \forall \beta\,(\varphi_{41})$\lbl{A2}}
\BIC{$\forall \beta\,(\varphi_{41})$}
\AXC{$\forall \beta\,(\varphi_{41}) \to \varphi_{41}$\lbl{A2}}
\BIC{$\varphi_{41}$}
\end{ptree}
\noindent$\Pi_{18}$:
\begin{ptree}
\AXC{$\stackrel{\Pi_{8}}{\vdots}$}\noLine
\UIC{$\mathrm{Conv}(a,\alpha) \to \mathrm{Conv}(a,\beta) \to \alpha \in \mathbb{R}$}
\AXC{$\stackrel{\Pi_{16}}{\vdots}$}\noLine
\UIC{$\varphi_{51}$}
\AXC{$\stackrel{\Pi_{17}}{\vdots}$}\noLine
\UIC{$\varphi_{41}$}
\AXC{$\varphi_{41} \to a_{N} \in \mathbb{R} \to \alpha \in \mathbb{R} \to \beta \in \mathbb{R} \to \varphi_{60}$\lbl{A2}}
\BIC{$a_{N} \in \mathbb{R} \to \alpha \in \mathbb{R} \to \beta \in \mathbb{R} \to \varphi_{60}$}
\AXC{$\varphi_{77}$\lbl{A1}}
\BIC{$\varphi_{51} \to (\mathrm{Conv}(a,\alpha) \to \mathrm{Conv}(a,\beta) \to \alpha \in \mathbb{R}) \to \varphi_{76}$}
\BIC{$(\mathrm{Conv}(a,\alpha) \to \mathrm{Conv}(a,\beta) \to \alpha \in \mathbb{R}) \to \varphi_{76}$}
\BIC{$\varphi_{76}$}
\end{ptree}
\noindent$\Pi_{19}$:
\begin{ptree}
\AXC{$\stackrel{\Pi_{13}}{\vdots}$}\noLine
\UIC{$\mathrm{Conv}(a,\alpha) \to \varphi_{53}$}
\AXC{$\stackrel{\Pi_{6}}{\vdots}$}\noLine
\UIC{$\mathrm{Conv}(a,\alpha) \to \mathrm{Conv}(a,\beta) \to \beta \in \mathbb{R}$}
\AXC{$\stackrel{\Pi_{18}}{\vdots}$}\noLine
\UIC{$\varphi_{76}$}
\BIC{$(\mathrm{Conv}(a,\alpha) \to \varphi_{53}) \to (\mathrm{Conv}(a,\alpha) \to \varphi_{59}) \to \varphi_{71}$}
\BIC{$(\mathrm{Conv}(a,\alpha) \to \varphi_{59}) \to \varphi_{71}$}
\end{ptree}
\noindent$\Pi_{20}$:
\begin{ptree}
\AXC{$\stackrel{\Pi_{4}}{\vdots}$}\noLine
\UIC{$\varphi_{63}$}
\AXC{$\varphi_{63} \to \mathrm{Conv}(a,\alpha) \to \varphi_{58}$\lbl{A1}}
\BIC{$\mathrm{Conv}(a,\alpha) \to \varphi_{58}$}
\AXC{$\stackrel{\Pi_{4}}{\vdots}$}\noLine
\UIC{$\varphi_{63}$}
\AXC{$\varphi_{63} \to \mathrm{Conv}(a,\alpha) \to \varphi_{59}$\lbl{A1}}
\BIC{$\mathrm{Conv}(a,\alpha) \to \varphi_{59}$}
\AXC{$\stackrel{\Pi_{19}}{\vdots}$}\noLine
\UIC{$(\mathrm{Conv}(a,\alpha) \to \varphi_{59}) \to \varphi_{71}$}
\BIC{$\varphi_{71}$}
\BIC{$\mathrm{Conv}(a,\alpha) \to \mathrm{Conv}(a,\beta) \to \alpha \neq \beta \to \varphi_{57}$}
\end{ptree}
\noindent$\Pi_{21}$:
\begin{ptree}
\AXC{$\forall e \in \mathbb{R}\,(e \cdot 2^{-1} + e \cdot 2^{-1} = e)$\lbl{[F36]}}
\AXC{$\forall e \in \mathbb{R}\,(e \cdot 2^{-1} + e \cdot 2^{-1} = e) \to \varphi_{19}$\lbl{A2}}
\BIC{$\varphi_{19}$}
\end{ptree}
\noindent$\Pi_{22}$:
\begin{ptree}
\AXC{$\stackrel{\Pi_{12}}{\vdots}$}\noLine
\UIC{$\varphi_{54}$}
\AXC{$\stackrel{\Pi_{21}}{\vdots}$}\noLine
\UIC{$\varphi_{19}$}
\AXC{$\varphi_{19} \to \varphi_{54} \to \mathrm{Conv}(a,\alpha) \to \mathrm{Conv}(a,\beta) \to \alpha \neq \beta \to \varphi_{55}$\lbl{A1}}
\BIC{$\varphi_{54} \to \mathrm{Conv}(a,\alpha) \to \mathrm{Conv}(a,\beta) \to \alpha \neq \beta \to \varphi_{55}$}
\BIC{$\mathrm{Conv}(a,\alpha) \to \mathrm{Conv}(a,\beta) \to \alpha \neq \beta \to \varphi_{55}$}
\AXC{$\varphi_{75}$\lbl{A1}}
\BIC{$(\mathrm{Conv}(a,\alpha) \to \mathrm{Conv}(a,\beta) \to \alpha \neq \beta \to \varphi_{57}) \to \varphi_{69}$}
\end{ptree}
\noindent$\Pi_{23}$:
\begin{ptree}
\AXC{$\stackrel{\Pi_{20}}{\vdots}$}\noLine
\UIC{$\mathrm{Conv}(a,\alpha) \to \mathrm{Conv}(a,\beta) \to \alpha \neq \beta \to \varphi_{57}$}
\AXC{$\stackrel{\Pi_{22}}{\vdots}$}\noLine
\UIC{$(\mathrm{Conv}(a,\alpha) \to \mathrm{Conv}(a,\beta) \to \alpha \neq \beta \to \varphi_{57}) \to \varphi_{69}$}
\BIC{$\varphi_{69}$}
\end{ptree}
\noindent$\Pi_{24}$:
\begin{ptree}
\AXC{$\stackrel{\Pi_{1}}{\vdots}$}\noLine
\UIC{$\varphi_{38}$}
\AXC{$\varphi_{38} \to \lvert \alpha - \beta \rvert \cdot 2^{-1} + \lvert \alpha - \beta \rvert \cdot 2^{-1} = \lvert \alpha - \beta \rvert \to \varphi_{28}$\lbl{A2}}
\BIC{$\lvert \alpha - \beta \rvert \cdot 2^{-1} + \lvert \alpha - \beta \rvert \cdot 2^{-1} = \lvert \alpha - \beta \rvert \to \varphi_{28}$}
\AXC{$\stackrel{\Pi_{23}}{\vdots}$}\noLine
\UIC{$\varphi_{69}$}
\BIC{$\mathrm{Conv}(a,\alpha) \to \varphi_{52}$}
\end{ptree}
\noindent$\Pi_{25}$:
\begin{ptree}
\AXC{$\stackrel{\Pi_{24}}{\vdots}$}\noLine
\UIC{$\mathrm{Conv}(a,\alpha) \to \varphi_{52}$}
\AXC{$(\mathrm{Conv}(a,\alpha) \to \varphi_{52}) \to \neg \lvert \alpha - \beta \rvert < \lvert \alpha - \beta \rvert \to \varphi_{46}$\lbl{A1}}
\BIC{$\neg \lvert \alpha - \beta \rvert < \lvert \alpha - \beta \rvert \to \varphi_{46}$}
\end{ptree}
\noindent$\Pi_{26}$:
\begin{ptree}
\AXC{$\forall x\,(\neg x < x)$\lbl{[O1]}}
\AXC{$\forall x\,(\neg x < x) \to \neg \lvert \alpha - \beta \rvert < \lvert \alpha - \beta \rvert$\lbl{A2}}
\BIC{$\neg \lvert \alpha - \beta \rvert < \lvert \alpha - \beta \rvert$}
\AXC{$\stackrel{\Pi_{25}}{\vdots}$}\noLine
\UIC{$\neg \lvert \alpha - \beta \rvert < \lvert \alpha - \beta \rvert \to \varphi_{46}$}
\BIC{$\varphi_{46}$}
\AXC{$\varphi_{46} \to \varphi_{47}$\lbl{A1}}
\BIC{$\varphi_{47}$}
\UIC{$\forall N\,(\varphi_{47})$}
\end{ptree}
\noindent$\Pi_{27}$:
\begin{ptree}
\AXC{$\forall x\,\forall y\,(x \in \mathbb{R} \to y \in \mathbb{R} \to x \neq y \to x - y \neq 0)$\lbl{[A15]}}
\AXC{$\varphi_{26}$\lbl{A2}}
\BIC{$\forall y\,(\alpha \in \mathbb{R} \to y \in \mathbb{R} \to \alpha \neq y \to \alpha - y \neq 0)$}
\AXC{$\varphi_{27}$\lbl{A2}}
\BIC{$\alpha \in \mathbb{R} \to \beta \in \mathbb{R} \to \alpha \neq \beta \to \alpha - \beta \neq 0$}
\AXC{$(\alpha \in \mathbb{R} \to \beta \in \mathbb{R} \to \alpha \neq \beta \to \alpha - \beta \neq 0) \to \varphi_{45}$\lbl{A1}}
\BIC{$\varphi_{45}$}
\end{ptree}
\noindent$\Pi_{28}$:
\begin{ptree}
\AXC{$\stackrel{\Pi_{6}}{\vdots}$}\noLine
\UIC{$\mathrm{Conv}(a,\alpha) \to \mathrm{Conv}(a,\beta) \to \beta \in \mathbb{R}$}
\AXC{$\stackrel{\Pi_{8}}{\vdots}$}\noLine
\UIC{$\mathrm{Conv}(a,\alpha) \to \mathrm{Conv}(a,\beta) \to \alpha \in \mathbb{R}$}
\AXC{$\stackrel{\Pi_{27}}{\vdots}$}\noLine
\UIC{$\varphi_{45}$}
\BIC{$(\mathrm{Conv}(a,\alpha) \to \mathrm{Conv}(a,\beta) \to \beta \in \mathbb{R}) \to \varphi_{21}$}
\BIC{$\varphi_{21}$}
\end{ptree}
\noindent$\Pi_{29}$:
\begin{ptree}
\AXC{$\forall x\,\forall y\,(x \in \mathbb{R} \to y \in \mathbb{R} \to x - y \in \mathbb{R})$\lbl{[F1]}}
\AXC{$\varphi_{25}$\lbl{A2}}
\BIC{$\forall y\,(\alpha \in \mathbb{R} \to y \in \mathbb{R} \to \alpha - y \in \mathbb{R})$}
\AXC{$\varphi_{23}$\lbl{A2}}
\BIC{$\alpha \in \mathbb{R} \to \beta \in \mathbb{R} \to \alpha - \beta \in \mathbb{R}$}
\AXC{$\varphi_{49}$\lbl{A1}}
\BIC{$(\mathrm{Conv}(a,\alpha) \to \mathrm{Conv}(a,\beta) \to \alpha \in \mathbb{R}) \to \varphi_{22}$}
\end{ptree}
\noindent$\Pi_{30}$:
\begin{ptree}
\AXC{$\stackrel{\Pi_{8}}{\vdots}$}\noLine
\UIC{$\mathrm{Conv}(a,\alpha) \to \mathrm{Conv}(a,\beta) \to \alpha \in \mathbb{R}$}
\AXC{$\stackrel{\Pi_{29}}{\vdots}$}\noLine
\UIC{$(\mathrm{Conv}(a,\alpha) \to \mathrm{Conv}(a,\beta) \to \alpha \in \mathbb{R}) \to \varphi_{22}$}
\BIC{$\varphi_{22}$}
\end{ptree}
\noindent$\Pi_{31}$:
\begin{ptree}
\AXC{$\stackrel{\Pi_{6}}{\vdots}$}\noLine
\UIC{$\mathrm{Conv}(a,\alpha) \to \mathrm{Conv}(a,\beta) \to \beta \in \mathbb{R}$}
\AXC{$\stackrel{\Pi_{30}}{\vdots}$}\noLine
\UIC{$\varphi_{22}$}
\BIC{$\mathrm{Conv}(a,\alpha) \to \mathrm{Conv}(a,\beta) \to \alpha \neq \beta \to \alpha - \beta \in \mathbb{R}$}
\end{ptree}
\noindent$\Pi_{32}$:
\begin{ptree}
\AXC{$\forall x \in \mathbb{R}\,(x \neq 0 \to 0 < \lvert x \rvert)$\lbl{[A14]}}
\AXC{$\varphi_{20}$\lbl{A2}}
\BIC{$\alpha - \beta \in \mathbb{R} \to \alpha - \beta \neq 0 \to 0 < \lvert \alpha - \beta \rvert$}
\AXC{$(\alpha - \beta \in \mathbb{R} \to \alpha - \beta \neq 0 \to 0 < \lvert \alpha - \beta \rvert) \to \varphi_{43}$\lbl{A1}}
\BIC{$\varphi_{43}$}
\end{ptree}
\noindent$\Pi_{33}$:
\begin{ptree}
\AXC{$\stackrel{\Pi_{31}}{\vdots}$}\noLine
\UIC{$\mathrm{Conv}(a,\alpha) \to \mathrm{Conv}(a,\beta) \to \alpha \neq \beta \to \alpha - \beta \in \mathbb{R}$}
\AXC{$\stackrel{\Pi_{32}}{\vdots}$}\noLine
\UIC{$\varphi_{43}$}
\BIC{$\varphi_{29}$}
\end{ptree}
\noindent$\Pi_{34}$:
\begin{ptree}
\AXC{$\forall x \in \mathbb{R}\,(\lvert x \rvert \in \mathbb{R})$\lbl{[A5]}}
\AXC{$\forall x \in \mathbb{R}\,(\lvert x \rvert \in \mathbb{R}) \to \alpha - \beta \in \mathbb{R} \to \lvert \alpha - \beta \rvert \in \mathbb{R}$\lbl{A2}}
\BIC{$\alpha - \beta \in \mathbb{R} \to \lvert \alpha - \beta \rvert \in \mathbb{R}$}
\end{ptree}
\noindent$\Pi_{35}$:
\begin{ptree}
\AXC{$\stackrel{\Pi_{34}}{\vdots}$}\noLine
\UIC{$\alpha - \beta \in \mathbb{R} \to \lvert \alpha - \beta \rvert \in \mathbb{R}$}
\AXC{$(\alpha - \beta \in \mathbb{R} \to \lvert \alpha - \beta \rvert \in \mathbb{R}) \to \varphi_{30}$\lbl{A1}}
\BIC{$\varphi_{30}$}
\end{ptree}
\noindent$\Pi_{36}$:
\begin{ptree}
\AXC{$\stackrel{\Pi_{31}}{\vdots}$}\noLine
\UIC{$\mathrm{Conv}(a,\alpha) \to \mathrm{Conv}(a,\beta) \to \alpha \neq \beta \to \alpha - \beta \in \mathbb{R}$}
\AXC{$\stackrel{\Pi_{35}}{\vdots}$}\noLine
\UIC{$\varphi_{30}$}
\BIC{$\varphi_{9}$}
\end{ptree}
\noindent$\Pi_{37}$:
\begin{ptree}
\AXC{$\stackrel{\Pi_{36}}{\vdots}$}\noLine
\UIC{$\varphi_{9}$}
\AXC{$\forall e \in \mathbb{R}\,(0 < e \to 0 < e \cdot 2^{-1})$\lbl{[F35]}}
\AXC{$\varphi_{34}$\lbl{A2}}
\BIC{$\lvert \alpha - \beta \rvert \in \mathbb{R} \to 0 < \lvert \alpha - \beta \rvert \to 0 < \lvert \alpha - \beta \rvert \cdot 2^{-1}$}
\AXC{$\varphi_{62}$\lbl{A1}}
\BIC{$\varphi_{9} \to \varphi_{6} \to \varphi_{13}$}
\BIC{$\varphi_{6} \to \varphi_{13}$}
\end{ptree}
\noindent$\Pi_{38}$:
\begin{ptree}
\AXC{$\mathrm{Conv}(a,\alpha) \to \mathrm{Conv}(a,\beta) \to \alpha \neq \beta \to \alpha \neq \beta$\lbl{A1}}
\AXC{$\stackrel{\Pi_{28}}{\vdots}$}\noLine
\UIC{$\varphi_{21}$}
\BIC{$\mathrm{Conv}(a,\alpha) \to \mathrm{Conv}(a,\beta) \to \alpha \neq \beta \to \alpha - \beta \neq 0$}
\AXC{$\stackrel{\Pi_{33}}{\vdots}$}\noLine
\UIC{$\varphi_{29}$}
\BIC{$\varphi_{6}$}
\AXC{$\stackrel{\Pi_{37}}{\vdots}$}\noLine
\UIC{$\varphi_{6} \to \varphi_{13}$}
\BIC{$\varphi_{13}$}
\end{ptree}
\noindent$\Pi_{39}$:
\begin{ptree}
\AXC{$2^{-1} \in \mathbb{R}$\lbl{[F31]}}
\AXC{$\stackrel{\Pi_{36}}{\vdots}$}\noLine
\UIC{$\varphi_{9}$}
\AXC{$\forall x\,\forall y\,(x \in \mathbb{R} \to y \in \mathbb{R} \to x \cdot y \in \mathbb{R})$\lbl{(R1)}}
\AXC{$\varphi_{33}$\lbl{A2}}
\BIC{$\forall y\,(\lvert \alpha - \beta \rvert \in \mathbb{R} \to y \in \mathbb{R} \to \lvert \alpha - \beta \rvert \cdot y \in \mathbb{R})$}
\AXC{$\varphi_{39}$\lbl{A2}}
\BIC{$\lvert \alpha - \beta \rvert \in \mathbb{R} \to 2^{-1} \in \mathbb{R} \to \lvert \alpha - \beta \rvert \cdot 2^{-1} \in \mathbb{R}$}
\AXC{$\varphi_{61}$\lbl{A1}}
\BIC{$\varphi_{9} \to 2^{-1} \in \mathbb{R} \to \varphi_{14}$}
\BIC{$2^{-1} \in \mathbb{R} \to \varphi_{14}$}
\BIC{$\varphi_{14}$}
\end{ptree}
\noindent$\Pi_{40}$:
\begin{ptree}
\AXC{$\forall a\,\forall \alpha\,(\varphi_{42})$\lbl{[C9]}}
\AXC{$\forall a\,\forall \alpha\,(\varphi_{42}) \to \forall \alpha\,(\varphi_{42})$\lbl{A2}}
\BIC{$\forall \alpha\,(\varphi_{42})$}
\AXC{$\forall \alpha\,(\varphi_{42}) \to \varphi_{42}$\lbl{A2}}
\BIC{$\varphi_{42}$}
\AXC{$\varphi_{42} \to \forall \beta\,(\varphi_{40})$\lbl{A2}}
\BIC{$\forall \beta\,(\varphi_{40})$}
\AXC{$\forall \beta\,(\varphi_{40}) \to \varphi_{40}$\lbl{A2}}
\BIC{$\varphi_{40}$}
\end{ptree}
\noindent$\Pi_{41}$:
\begin{ptree}
\AXC{$\mathrm{Conv}(a,\alpha) \to \mathrm{Conv}(a,\alpha)$\lbl{A1}}
\AXC{$\stackrel{\Pi_{40}}{\vdots}$}\noLine
\UIC{$\varphi_{40}$}
\AXC{$\varphi_{40} \to \mathrm{Conv}(a,\alpha) \to \mathrm{Conv}(a,\beta) \to \varphi_{56}$\lbl{A2}}
\BIC{$\mathrm{Conv}(a,\alpha) \to \mathrm{Conv}(a,\beta) \to \varphi_{56}$}
\AXC{$\varphi_{73}$\lbl{A1}}
\BIC{$(\mathrm{Conv}(a,\alpha) \to \mathrm{Conv}(a,\alpha)) \to \varphi_{68}$}
\BIC{$\varphi_{68}$}
\end{ptree}
\noindent$\Pi_{42}$:
\begin{ptree}
\AXC{$\stackrel{\Pi_{26}}{\vdots}$}\noLine
\UIC{$\forall N\,(\varphi_{47})$}
\AXC{$\stackrel{\Pi_{38}}{\vdots}$}\noLine
\UIC{$\varphi_{13}$}
\AXC{$\stackrel{\Pi_{39}}{\vdots}$}\noLine
\UIC{$\varphi_{14}$}
\AXC{$\mathrm{Conv}(a,\alpha) \to \mathrm{Conv}(a,\beta) \to \mathrm{Conv}(a,\beta)$\lbl{A1}}
\AXC{$\stackrel{\Pi_{41}}{\vdots}$}\noLine
\UIC{$\varphi_{68}$}
\BIC{$\varphi_{14} \to \varphi_{13} \to \varphi_{44}$}
\BIC{$\varphi_{13} \to \varphi_{44}$}
\BIC{$\varphi_{44}$}
\AXC{$\varphi_{44} \to \varphi_{74}$\lbl{A1}}
\BIC{$\varphi_{74}$}
\BIC{$(\forall N\,(\varphi_{47}) \to \varphi_{3} \to \forall N\,(\neg (N \in \mathbb{N} \land \varphi_{31}))) \to \varphi_{36}$}
\end{ptree}
\noindent$\Pi_{43}$:
\begin{ptree}
\AXC{$\forall N\,(\varphi_{47}) \to \varphi_{3} \to \forall N\,(\neg (N \in \mathbb{N} \land \varphi_{31}))$\lbl{A3}}
\AXC{$\stackrel{\Pi_{42}}{\vdots}$}\noLine
\UIC{$(\forall N\,(\varphi_{47}) \to \varphi_{3} \to \forall N\,(\neg (N \in \mathbb{N} \land \varphi_{31}))) \to \varphi_{36}$}
\BIC{$\varphi_{36}$}
\end{ptree}
\noindent$\Pi_{44}$:
\begin{ptree}
\AXC{$\mathrm{Conv}(a,\alpha) \land \mathrm{Conv}(a,\beta) \land \alpha \neq \beta \land \neg 0 < 0 \to \varphi_{3}$\lbl{A4}}
\AXC{$\stackrel{\Pi_{43}}{\vdots}$}\noLine
\UIC{$\varphi_{36}$}
\BIC{$\mathrm{Conv}(a,\alpha) \to \mathrm{Conv}(a,\beta) \to \alpha \neq \beta \to 0 < 0$}
\AXC{$\varphi_{35}$\lbl{A1}}
\BIC{$\varphi_{16}$}
\end{ptree}
\noindent 主木:
\begin{ptree}
\AXC{$\forall x\,(\neg x < x)$\lbl{[O1]}}
\AXC{$\forall x\,(\neg x < x) \to \neg 0 < 0$\lbl{A2}}
\BIC{$\neg 0 < 0$}
\AXC{$\varphi_{4}$\lbl{A1}}
\BIC{$\mathrm{Conv}(a,\alpha) \to \mathrm{Conv}(a,\beta) \to \alpha \neq \beta \to \neg 0 < 0$}
\AXC{$\stackrel{\Pi_{44}}{\vdots}$}\noLine
\UIC{$\varphi_{16}$}
\BIC{$\mathrm{Conv}(a,\alpha) \to \mathrm{Conv}(a,\beta) \to \alpha = \beta$}
\UIC{$\forall \beta\,(\mathrm{Conv}(a,\alpha) \to \mathrm{Conv}(a,\beta) \to \alpha = \beta)$}
\UIC{$\forall \alpha\,\forall \beta\,(\mathrm{Conv}(a,\alpha) \to \mathrm{Conv}(a,\beta) \to \alpha = \beta)$}
\UIC{$\varphi_{1}$}
\end{ptree}
\begin{abbrevlist}
\item $\varphi_{1} :\equiv \forall a\,\forall \alpha\,\forall \beta\,(\mathrm{Conv}(a,\alpha) \to \mathrm{Conv}(a,\beta) \to \alpha = \beta)$
\item $\varphi_{2} :\equiv \forall \alpha\,(\mathrm{Conv}(a,\alpha) \to \mathrm{Seq}(a)) \to \mathrm{Conv}(a,\alpha) \to \mathrm{Seq}(a)$
\item $\varphi_{3} :\equiv \forall N\,(\mathrm{Conv}(a,\alpha) \land \mathrm{Conv}(a,\beta) \land \alpha \neq \beta \land \neg 0 < 0)$
\item $\varphi_{4} :\equiv \neg 0 < 0 \to \mathrm{Conv}(a,\alpha) \to \mathrm{Conv}(a,\beta) \to \alpha \neq \beta \to \neg 0 < 0$
\item $\varphi_{5} :\equiv \forall a\,\forall \alpha\,(\mathrm{Conv}(a,\alpha) \to \mathrm{Seq}(a)) \to \forall \alpha\,(\mathrm{Conv}(a,\alpha) \to \mathrm{Seq}(a))$
\item $\varphi_{6} :\equiv \mathrm{Conv}(a,\alpha) \to \mathrm{Conv}(a,\beta) \to \alpha \neq \beta \to 0 < \lvert \alpha - \beta \rvert$
\item $\varphi_{7} :\equiv \forall \alpha\,(\mathrm{Conv}(a,\alpha) \to \alpha \in \mathbb{R}) \to \mathrm{Conv}(a,\beta) \to \beta \in \mathbb{R}$
\item $\varphi_{8} :\equiv \forall \alpha\,(\mathrm{Conv}(a,\alpha) \to \alpha \in \mathbb{R}) \to \mathrm{Conv}(a,\alpha) \to \alpha \in \mathbb{R}$
\item $\varphi_{9} :\equiv \mathrm{Conv}(a,\alpha) \to \mathrm{Conv}(a,\beta) \to \alpha \neq \beta \to \lvert \alpha - \beta \rvert \in \mathbb{R}$
\item $\varphi_{10} :\equiv (\mathrm{Conv}(a,\alpha) \to \mathrm{Conv}(a,\alpha)) \to \mathrm{Conv}(a,\alpha) \to \mathrm{Conv}(a,\beta) \to \mathrm{Seq}(a)$
\item $\varphi_{11} :\equiv \forall a\,\forall \alpha\,(\mathrm{Conv}(a,\alpha) \to \alpha \in \mathbb{R}) \to \forall \alpha\,(\mathrm{Conv}(a,\alpha) \to \alpha \in \mathbb{R})$
\item $\varphi_{12} :\equiv (\mathrm{Conv}(a,\alpha) \to \mathrm{Conv}(a,\alpha)) \to \mathrm{Conv}(a,\alpha) \to \mathrm{Conv}(a,\beta) \to \alpha \in \mathbb{R}$
\item $\varphi_{13} :\equiv \mathrm{Conv}(a,\alpha) \to \mathrm{Conv}(a,\beta) \to \alpha \neq \beta \to 0 < \lvert \alpha - \beta \rvert \cdot 2^{-1}$
\item $\varphi_{14} :\equiv \mathrm{Conv}(a,\alpha) \to \mathrm{Conv}(a,\beta) \to \alpha \neq \beta \to \lvert \alpha - \beta \rvert \cdot 2^{-1} \in \mathbb{R}$
\item $\varphi_{15} :\equiv (\mathrm{Conv}(a,\alpha) \to \mathrm{Conv}(a,\beta) \to \mathrm{Conv}(a,\beta)) \to \mathrm{Conv}(a,\alpha) \to \mathrm{Conv}(a,\beta) \to \beta \in \mathbb{R}$
\item $\varphi_{16} :\equiv (\mathrm{Conv}(a,\alpha) \to \mathrm{Conv}(a,\beta) \to \alpha \neq \beta \to \neg 0 < 0) \to \mathrm{Conv}(a,\alpha) \to \mathrm{Conv}(a,\beta) \to \alpha = \beta$
\item $\varphi_{17} :\equiv \forall n\,(\mathrm{Seq}(a) \to n \in \mathbb{N} \to a_{n} \in \mathbb{R}) \to \mathrm{Seq}(a) \to N \in \mathbb{N} \to a_{N} \in \mathbb{R}$
\item $\varphi_{18} :\equiv \forall a\,\forall n\,(\mathrm{Seq}(a) \to n \in \mathbb{N} \to a_{n} \in \mathbb{R}) \to \forall n\,(\mathrm{Seq}(a) \to n \in \mathbb{N} \to a_{n} \in \mathbb{R})$
\item $\varphi_{19} :\equiv \lvert \alpha - \beta \rvert \in \mathbb{R} \to \lvert \alpha - \beta \rvert \cdot 2^{-1} + \lvert \alpha - \beta \rvert \cdot 2^{-1} = \lvert \alpha - \beta \rvert$
\item $\varphi_{20} :\equiv \forall x \in \mathbb{R}\,(x \neq 0 \to 0 < \lvert x \rvert) \to \alpha - \beta \in \mathbb{R} \to \alpha - \beta \neq 0 \to 0 < \lvert \alpha - \beta \rvert$
\item $\varphi_{21} :\equiv (\mathrm{Conv}(a,\alpha) \to \mathrm{Conv}(a,\beta) \to \alpha \neq \beta \to \alpha \neq \beta) \to \mathrm{Conv}(a,\alpha) \to \mathrm{Conv}(a,\beta) \to \alpha \neq \beta \to \alpha - \beta \neq 0$
\item $\varphi_{22} :\equiv (\mathrm{Conv}(a,\alpha) \to \mathrm{Conv}(a,\beta) \to \beta \in \mathbb{R}) \to \mathrm{Conv}(a,\alpha) \to \mathrm{Conv}(a,\beta) \to \alpha \neq \beta \to \alpha - \beta \in \mathbb{R}$
\item $\varphi_{23} :\equiv \forall y\,(\alpha \in \mathbb{R} \to y \in \mathbb{R} \to \alpha - y \in \mathbb{R}) \to \alpha \in \mathbb{R} \to \beta \in \mathbb{R} \to \alpha - \beta \in \mathbb{R}$
\item $\varphi_{24} :\equiv \lvert a_{N} - \alpha \rvert < \lvert \alpha - \beta \rvert \cdot 2^{-1} \land \lvert a_{N} - \beta \rvert < \lvert \alpha - \beta \rvert \cdot 2^{-1}$
\item $\varphi_{25} :\equiv \forall x\,\forall y\,(x \in \mathbb{R} \to y \in \mathbb{R} \to x - y \in \mathbb{R}) \to \forall y\,(\alpha \in \mathbb{R} \to y \in \mathbb{R} \to \alpha - y \in \mathbb{R})$
\item $\varphi_{26} :\equiv \forall x\,\forall y\,(x \in \mathbb{R} \to y \in \mathbb{R} \to x \neq y \to x - y \neq 0) \to \forall y\,(\alpha \in \mathbb{R} \to y \in \mathbb{R} \to \alpha \neq y \to \alpha - y \neq 0)$
\item $\varphi_{27} :\equiv \forall y\,(\alpha \in \mathbb{R} \to y \in \mathbb{R} \to \alpha \neq y \to \alpha - y \neq 0) \to \alpha \in \mathbb{R} \to \beta \in \mathbb{R} \to \alpha \neq \beta \to \alpha - \beta \neq 0$
\item $\varphi_{28} :\equiv \lvert \alpha - \beta \rvert < \lvert \alpha - \beta \rvert \cdot 2^{-1} + \lvert \alpha - \beta \rvert \cdot 2^{-1} \to \lvert \alpha - \beta \rvert < \lvert \alpha - \beta \rvert$
\item $\varphi_{29} :\equiv (\mathrm{Conv}(a,\alpha) \to \mathrm{Conv}(a,\beta) \to \alpha \neq \beta \to \alpha - \beta \neq 0) \to \varphi_{6}$
\item $\varphi_{30} :\equiv (\mathrm{Conv}(a,\alpha) \to \mathrm{Conv}(a,\beta) \to \alpha \neq \beta \to \alpha - \beta \in \mathbb{R}) \to \varphi_{9}$
\item $\varphi_{31} :\equiv \forall n \in \mathbb{N}\,(N \leq n \to \lvert a_{n} - \alpha \rvert < \lvert \alpha - \beta \rvert \cdot 2^{-1} \land \lvert a_{n} - \beta \rvert < \lvert \alpha - \beta \rvert \cdot 2^{-1})$
\item $\varphi_{32} :\equiv \lvert a_{N} - \beta \rvert < \lvert \alpha - \beta \rvert \cdot 2^{-1} \to \lvert \alpha - \beta \rvert < \lvert \alpha - \beta \rvert \cdot 2^{-1} + \lvert \alpha - \beta \rvert \cdot 2^{-1}$
\item $\varphi_{33} :\equiv \forall x\,\forall y\,(x \in \mathbb{R} \to y \in \mathbb{R} \to x \cdot y \in \mathbb{R}) \to \forall y\,(\lvert \alpha - \beta \rvert \in \mathbb{R} \to y \in \mathbb{R} \to \lvert \alpha - \beta \rvert \cdot y \in \mathbb{R})$
\item $\varphi_{34} :\equiv \forall e \in \mathbb{R}\,(0 < e \to 0 < e \cdot 2^{-1}) \to \lvert \alpha - \beta \rvert \in \mathbb{R} \to 0 < \lvert \alpha - \beta \rvert \to 0 < \lvert \alpha - \beta \rvert \cdot 2^{-1}$
\item $\varphi_{35} :\equiv (\mathrm{Conv}(a,\alpha) \to \mathrm{Conv}(a,\beta) \to \alpha \neq \beta \to 0 < 0) \to \varphi_{16}$
\item $\varphi_{36} :\equiv (\mathrm{Conv}(a,\alpha) \land \mathrm{Conv}(a,\beta) \land \alpha \neq \beta \land \neg 0 < 0 \to \varphi_{3}) \to \mathrm{Conv}(a,\alpha) \to \mathrm{Conv}(a,\beta) \to \alpha \neq \beta \to 0 < 0$
\item $\varphi_{37} :\equiv \forall x\,\forall \alpha\,\forall \beta\,\forall \varepsilon\,(x \in \mathbb{R} \to \alpha \in \mathbb{R} \to \beta \in \mathbb{R} \to \varepsilon \in \mathbb{R} \to \lvert x - \alpha \rvert < \varepsilon \to \lvert x - \beta \rvert < \varepsilon \to \lvert \alpha - \beta \rvert < \varepsilon + \varepsilon)$
\item $\varphi_{38} :\equiv \forall w\,(\lvert \alpha - \beta \rvert \cdot 2^{-1} + \lvert \alpha - \beta \rvert \cdot 2^{-1} = w \to \lvert \alpha - \beta \rvert < \lvert \alpha - \beta \rvert \cdot 2^{-1} + \lvert \alpha - \beta \rvert \cdot 2^{-1} \to \lvert \alpha - \beta \rvert < w)$
\item $\varphi_{39} :\equiv \forall y\,(\lvert \alpha - \beta \rvert \in \mathbb{R} \to y \in \mathbb{R} \to \lvert \alpha - \beta \rvert \cdot y \in \mathbb{R}) \to \lvert \alpha - \beta \rvert \in \mathbb{R} \to 2^{-1} \in \mathbb{R} \to \lvert \alpha - \beta \rvert \cdot 2^{-1} \in \mathbb{R}$
\item $\varphi_{40} :\equiv \forall \varepsilon\,(\mathrm{Conv}(a,\alpha) \to \mathrm{Conv}(a,\beta) \to \varepsilon \in \mathbb{R} \to 0 < \varepsilon \to \exists N \in \mathbb{N}\,\forall n \in \mathbb{N}\,(N \leq n \to \lvert a_{n} - \alpha \rvert < \varepsilon \land \lvert a_{n} - \beta \rvert < \varepsilon))$
\item $\varphi_{41} :\equiv \forall \varepsilon\,(a_{N} \in \mathbb{R} \to \alpha \in \mathbb{R} \to \beta \in \mathbb{R} \to \varepsilon \in \mathbb{R} \to \lvert a_{N} - \alpha \rvert < \varepsilon \to \lvert a_{N} - \beta \rvert < \varepsilon \to \lvert \alpha - \beta \rvert < \varepsilon + \varepsilon)$
\item $\varphi_{42} :\equiv \forall b\,\forall \beta\,\forall \varepsilon\,(\mathrm{Conv}(a,\alpha) \to \mathrm{Conv}(b,\beta) \to \varepsilon \in \mathbb{R} \to 0 < \varepsilon \to \exists N \in \mathbb{N}\,\forall n \in \mathbb{N}\,(N \leq n \to \lvert a_{n} - \alpha \rvert < \varepsilon \land \lvert b_{n} - \beta \rvert < \varepsilon))$
\item $\varphi_{43} :\equiv (\mathrm{Conv}(a,\alpha) \to \mathrm{Conv}(a,\beta) \to \alpha \neq \beta \to \alpha - \beta \in \mathbb{R}) \to \varphi_{29}$
\item $\varphi_{44} :\equiv \mathrm{Conv}(a,\alpha) \to \mathrm{Conv}(a,\beta) \to \alpha \neq \beta \to \exists N \in \mathbb{N}\,(\varphi_{31})$
\item $\varphi_{45} :\equiv (\mathrm{Conv}(a,\alpha) \to \mathrm{Conv}(a,\beta) \to \alpha \in \mathbb{R}) \to (\mathrm{Conv}(a,\alpha) \to \mathrm{Conv}(a,\beta) \to \beta \in \mathbb{R}) \to \varphi_{21}$
\item $\varphi_{46} :\equiv \mathrm{Conv}(a,\alpha) \to \mathrm{Conv}(a,\beta) \to \alpha \neq \beta \to N \in \mathbb{N} \land \varphi_{31} \to 0 < 0$
\item $\varphi_{47} :\equiv \mathrm{Conv}(a,\alpha) \land \mathrm{Conv}(a,\beta) \land \alpha \neq \beta \land \neg 0 < 0 \to \neg (N \in \mathbb{N} \land \varphi_{31})$
\item $\varphi_{48} :\equiv \mathrm{Conv}(a,\alpha) \to \mathrm{Conv}(a,\beta) \to \alpha \neq \beta \to N \in \mathbb{N} \land \varphi_{31} \to N \in \mathbb{N}$
\item $\varphi_{49} :\equiv (\alpha \in \mathbb{R} \to \beta \in \mathbb{R} \to \alpha - \beta \in \mathbb{R}) \to (\mathrm{Conv}(a,\alpha) \to \mathrm{Conv}(a,\beta) \to \alpha \in \mathbb{R}) \to \varphi_{22}$
\item $\varphi_{50} :\equiv \mathrm{Conv}(a,\alpha) \to \mathrm{Conv}(a,\beta) \to \alpha \neq \beta \to N \in \mathbb{N} \land \varphi_{31} \to \alpha - \beta \in \mathbb{R}$
\item $\varphi_{51} :\equiv \mathrm{Conv}(a,\alpha) \to \mathrm{Conv}(a,\beta) \to \alpha \neq \beta \to N \in \mathbb{N} \land \varphi_{31} \to a_{N} \in \mathbb{R}$
\item $\varphi_{52} :\equiv \mathrm{Conv}(a,\beta) \to \alpha \neq \beta \to N \in \mathbb{N} \land \varphi_{31} \to \lvert \alpha - \beta \rvert < \lvert \alpha - \beta \rvert$
\item $\varphi_{53} :\equiv \mathrm{Conv}(a,\beta) \to \alpha \neq \beta \to N \in \mathbb{N} \land \varphi_{31} \to \lvert \alpha - \beta \rvert \cdot 2^{-1} \in \mathbb{R}$
\item $\varphi_{54} :\equiv \mathrm{Conv}(a,\alpha) \to \mathrm{Conv}(a,\beta) \to \alpha \neq \beta \to N \in \mathbb{N} \land \varphi_{31} \to \lvert \alpha - \beta \rvert \in \mathbb{R}$
\item $\varphi_{55} :\equiv N \in \mathbb{N} \land \varphi_{31} \to \lvert \alpha - \beta \rvert \cdot 2^{-1} + \lvert \alpha - \beta \rvert \cdot 2^{-1} = \lvert \alpha - \beta \rvert$
\item $\varphi_{56} :\equiv \lvert \alpha - \beta \rvert \cdot 2^{-1} \in \mathbb{R} \to 0 < \lvert \alpha - \beta \rvert \cdot 2^{-1} \to \exists N \in \mathbb{N}\,(\varphi_{31})$
\item $\varphi_{57} :\equiv N \in \mathbb{N} \land \varphi_{31} \to \lvert \alpha - \beta \rvert < \lvert \alpha - \beta \rvert \cdot 2^{-1} + \lvert \alpha - \beta \rvert \cdot 2^{-1}$
\item $\varphi_{58} :\equiv \mathrm{Conv}(a,\beta) \to \alpha \neq \beta \to N \in \mathbb{N} \land \varphi_{31} \to \lvert a_{N} - \beta \rvert < \lvert \alpha - \beta \rvert \cdot 2^{-1}$
\item $\varphi_{59} :\equiv \mathrm{Conv}(a,\beta) \to \alpha \neq \beta \to N \in \mathbb{N} \land \varphi_{31} \to \lvert a_{N} - \alpha \rvert < \lvert \alpha - \beta \rvert \cdot 2^{-1}$
\item $\varphi_{60} :\equiv \lvert \alpha - \beta \rvert \cdot 2^{-1} \in \mathbb{R} \to \lvert a_{N} - \alpha \rvert < \lvert \alpha - \beta \rvert \cdot 2^{-1} \to \varphi_{32}$
\item $\varphi_{61} :\equiv (\lvert \alpha - \beta \rvert \in \mathbb{R} \to 2^{-1} \in \mathbb{R} \to \lvert \alpha - \beta \rvert \cdot 2^{-1} \in \mathbb{R}) \to \varphi_{9} \to 2^{-1} \in \mathbb{R} \to \varphi_{14}$
\item $\varphi_{62} :\equiv (\lvert \alpha - \beta \rvert \in \mathbb{R} \to 0 < \lvert \alpha - \beta \rvert \to 0 < \lvert \alpha - \beta \rvert \cdot 2^{-1}) \to \varphi_{9} \to \varphi_{6} \to \varphi_{13}$
\item $\varphi_{63} :\equiv \mathrm{Conv}(a,\alpha) \to \mathrm{Conv}(a,\beta) \to \alpha \neq \beta \to N \in \mathbb{N} \land \varphi_{31} \to \varphi_{24}$
\item $\varphi_{64} :\equiv (\mathrm{Conv}(a,\alpha) \to \mathrm{Conv}(a,\beta) \to \alpha \in \mathbb{R}) \to (\mathrm{Conv}(a,\alpha) \to \mathrm{Conv}(a,\beta) \to \beta \in \mathbb{R}) \to \varphi_{50}$
\item $\varphi_{65} :\equiv \mathrm{Conv}(a,\beta) \to \alpha \neq \beta \to N \in \mathbb{N} \land \varphi_{31} \to N \in \mathbb{N} \to N \leq N \to \varphi_{24}$
\item $\varphi_{66} :\equiv \mathrm{Conv}(a,\alpha) \to \mathrm{Conv}(a,\beta) \to \alpha \neq \beta \to N \in \mathbb{N} \land \varphi_{31} \to \varphi_{31}$
\item $\varphi_{67} :\equiv \mathrm{Conv}(a,\alpha) \to \mathrm{Conv}(a,\beta) \to \alpha \neq \beta \to N \in \mathbb{N} \land \varphi_{31} \to N \in \mathbb{N} \land \varphi_{31}$
\item $\varphi_{68} :\equiv (\mathrm{Conv}(a,\alpha) \to \mathrm{Conv}(a,\beta) \to \mathrm{Conv}(a,\beta)) \to \varphi_{14} \to \varphi_{13} \to \varphi_{44}$
\item $\varphi_{69} :\equiv (\lvert \alpha - \beta \rvert \cdot 2^{-1} + \lvert \alpha - \beta \rvert \cdot 2^{-1} = \lvert \alpha - \beta \rvert \to \varphi_{28}) \to \mathrm{Conv}(a,\alpha) \to \varphi_{52}$
\item $\varphi_{70} :\equiv (\mathrm{Conv}(a,\alpha) \to \mathrm{Conv}(a,\beta) \to \mathrm{Seq}(a)) \to \varphi_{48} \to \varphi_{51}$
\item $\varphi_{71} :\equiv (\mathrm{Conv}(a,\alpha) \to \varphi_{58}) \to \mathrm{Conv}(a,\alpha) \to \mathrm{Conv}(a,\beta) \to \alpha \neq \beta \to \varphi_{57}$
\item $\varphi_{72} :\equiv (\lvert \alpha - \beta \rvert \in \mathbb{R} \to 2^{-1} \in \mathbb{R} \to \lvert \alpha - \beta \rvert \cdot 2^{-1} \in \mathbb{R}) \to \varphi_{54} \to 2^{-1} \in \mathbb{R} \to \mathrm{Conv}(a,\alpha) \to \varphi_{53}$
\item $\varphi_{73} :\equiv (\mathrm{Conv}(a,\alpha) \to \mathrm{Conv}(a,\beta) \to \varphi_{56}) \to (\mathrm{Conv}(a,\alpha) \to \mathrm{Conv}(a,\alpha)) \to \varphi_{68}$
\item $\varphi_{74} :\equiv \forall N\,(\varphi_{47}) \to (\forall N\,(\varphi_{47}) \to \varphi_{3} \to \forall N\,(\neg (N \in \mathbb{N} \land \varphi_{31}))) \to \varphi_{36}$
\item $\varphi_{75} :\equiv (\mathrm{Conv}(a,\alpha) \to \mathrm{Conv}(a,\beta) \to \alpha \neq \beta \to \varphi_{55}) \to (\mathrm{Conv}(a,\alpha) \to \mathrm{Conv}(a,\beta) \to \alpha \neq \beta \to \varphi_{57}) \to \varphi_{69}$
\item $\varphi_{76} :\equiv (\mathrm{Conv}(a,\alpha) \to \mathrm{Conv}(a,\beta) \to \beta \in \mathbb{R}) \to (\mathrm{Conv}(a,\alpha) \to \varphi_{53}) \to (\mathrm{Conv}(a,\alpha) \to \varphi_{59}) \to \varphi_{71}$
\item $\varphi_{77} :\equiv (a_{N} \in \mathbb{R} \to \alpha \in \mathbb{R} \to \beta \in \mathbb{R} \to \varphi_{60}) \to \varphi_{51} \to (\mathrm{Conv}(a,\alpha) \to \mathrm{Conv}(a,\beta) \to \alpha \in \mathbb{R}) \to \varphi_{76}$
\end{abbrevlist}
\end{fullproof}

% 自動生成 (prover/build.js): ErrSum  ノード数 155
\begin{fullproof}{[C11]}{$\forall x\,\forall y\,\forall \alpha\,\forall \beta\,\forall h\,(x \in \mathbb{R} \to y \in \mathbb{R} \to \alpha \in \mathbb{R} \to \beta \in \mathbb{R} \to h \in \mathbb{R} \to \lvert x - \alpha \rvert < h \to \lvert y - \beta \rvert < h \to \lvert x + y - (\alpha + \beta) \rvert < h + h)$}{155}\phantomsection\label{tree:ErrSum}
\noindent$\Pi_{1}$:
\begin{ptree}
\AXC{$\forall x\,\forall y\,(x \in \mathbb{R} \to y \in \mathbb{R} \to x - y \in \mathbb{R})$\lbl{[F1]}}
\AXC{$\varphi_{41}$\lbl{A2}}
\BIC{$\forall y\,(v_{378} \in \mathbb{R} \to y \in \mathbb{R} \to v_{378} - y \in \mathbb{R})$}
\AXC{$\varphi_{42}$\lbl{A2}}
\BIC{$v_{378} \in \mathbb{R} \to v_{379} \in \mathbb{R} \to v_{378} - v_{379} \in \mathbb{R}$}
\UIC{$\forall v_{379}\,(v_{378} \in \mathbb{R} \to v_{379} \in \mathbb{R} \to v_{378} - v_{379} \in \mathbb{R})$}
\UIC{$\varphi_{4}$}
\AXC{$\varphi_{4} \to \forall v_{379}\,(y \in \mathbb{R} \to v_{379} \in \mathbb{R} \to y - v_{379} \in \mathbb{R})$\lbl{A2}}
\BIC{$\forall v_{379}\,(y \in \mathbb{R} \to v_{379} \in \mathbb{R} \to y - v_{379} \in \mathbb{R})$}
\end{ptree}
\noindent$\Pi_{2}$:
\begin{ptree}
\AXC{$\stackrel{\Pi_{1}}{\vdots}$}\noLine
\UIC{$\forall v_{379}\,(y \in \mathbb{R} \to v_{379} \in \mathbb{R} \to y - v_{379} \in \mathbb{R})$}
\AXC{$\varphi_{40}$\lbl{A2}}
\BIC{$y \in \mathbb{R} \to \beta \in \mathbb{R} \to y - \beta \in \mathbb{R}$}
\AXC{$(y \in \mathbb{R} \to \beta \in \mathbb{R} \to y - \beta \in \mathbb{R}) \to \varphi_{62}$\lbl{A1}}
\BIC{$\varphi_{62}$}
\end{ptree}
\noindent$\Pi_{3}$:
\begin{ptree}
\AXC{$x \in \mathbb{R} \to y \in \mathbb{R} \to \varphi_{36}$\lbl{A1}}
\AXC{$x \in \mathbb{R} \to y \in \mathbb{R} \to \varphi_{34}$\lbl{A1}}
\AXC{$\stackrel{\Pi_{2}}{\vdots}$}\noLine
\UIC{$\varphi_{62}$}
\BIC{$(x \in \mathbb{R} \to y \in \mathbb{R} \to \varphi_{36}) \to x \in \mathbb{R} \to y \in \mathbb{R} \to \alpha \in \mathbb{R} \to \varphi_{22}$}
\BIC{$x \in \mathbb{R} \to y \in \mathbb{R} \to \alpha \in \mathbb{R} \to \varphi_{22}$}
\end{ptree}
\noindent$\Pi_{4}$:
\begin{ptree}
\AXC{$\forall x \in \mathbb{R}\,(\lvert x \rvert \in \mathbb{R})$\lbl{[A5]}}
\AXC{$\forall x \in \mathbb{R}\,(\lvert x \rvert \in \mathbb{R}) \to y - \beta \in \mathbb{R} \to \lvert y - \beta \rvert \in \mathbb{R}$\lbl{A2}}
\BIC{$y - \beta \in \mathbb{R} \to \lvert y - \beta \rvert \in \mathbb{R}$}
\end{ptree}
\noindent$\Pi_{5}$:
\begin{ptree}
\AXC{$\stackrel{\Pi_{3}}{\vdots}$}\noLine
\UIC{$x \in \mathbb{R} \to y \in \mathbb{R} \to \alpha \in \mathbb{R} \to \varphi_{22}$}
\AXC{$\stackrel{\Pi_{4}}{\vdots}$}\noLine
\UIC{$y - \beta \in \mathbb{R} \to \lvert y - \beta \rvert \in \mathbb{R}$}
\AXC{$(y - \beta \in \mathbb{R} \to \lvert y - \beta \rvert \in \mathbb{R}) \to \varphi_{56}$\lbl{A1}}
\BIC{$\varphi_{56}$}
\BIC{$x \in \mathbb{R} \to y \in \mathbb{R} \to \alpha \in \mathbb{R} \to \varphi_{26}$}
\end{ptree}
\noindent$\Pi_{6}$:
\begin{ptree}
\AXC{$\forall x\,\forall y\,(x \in \mathbb{R} \to y \in \mathbb{R} \to x - y \in \mathbb{R})$\lbl{[F1]}}
\AXC{$\varphi_{38}$\lbl{A2}}
\BIC{$\forall y\,(x \in \mathbb{R} \to y \in \mathbb{R} \to x - y \in \mathbb{R})$}
\AXC{$\varphi_{29}$\lbl{A2}}
\BIC{$x \in \mathbb{R} \to \alpha \in \mathbb{R} \to x - \alpha \in \mathbb{R}$}
\AXC{$(x \in \mathbb{R} \to \alpha \in \mathbb{R} \to x - \alpha \in \mathbb{R}) \to \varphi_{63}$\lbl{A1}}
\BIC{$\varphi_{63}$}
\end{ptree}
\noindent$\Pi_{7}$:
\begin{ptree}
\AXC{$x \in \mathbb{R} \to y \in \mathbb{R} \to \varphi_{37}$\lbl{A1}}
\AXC{$x \in \mathbb{R} \to y \in \mathbb{R} \to \varphi_{35}$\lbl{A1}}
\AXC{$\stackrel{\Pi_{6}}{\vdots}$}\noLine
\UIC{$\varphi_{63}$}
\BIC{$(x \in \mathbb{R} \to y \in \mathbb{R} \to \varphi_{37}) \to x \in \mathbb{R} \to y \in \mathbb{R} \to \alpha \in \mathbb{R} \to \varphi_{23}$}
\BIC{$x \in \mathbb{R} \to y \in \mathbb{R} \to \alpha \in \mathbb{R} \to \varphi_{23}$}
\end{ptree}
\noindent$\Pi_{8}$:
\begin{ptree}
\AXC{$\forall x \in \mathbb{R}\,(\lvert x \rvert \in \mathbb{R})$\lbl{[A5]}}
\AXC{$\forall x \in \mathbb{R}\,(\lvert x \rvert \in \mathbb{R}) \to x - \alpha \in \mathbb{R} \to \lvert x - \alpha \rvert \in \mathbb{R}$\lbl{A2}}
\BIC{$x - \alpha \in \mathbb{R} \to \lvert x - \alpha \rvert \in \mathbb{R}$}
\end{ptree}
\noindent$\Pi_{9}$:
\begin{ptree}
\AXC{$\stackrel{\Pi_{7}}{\vdots}$}\noLine
\UIC{$x \in \mathbb{R} \to y \in \mathbb{R} \to \alpha \in \mathbb{R} \to \varphi_{23}$}
\AXC{$\stackrel{\Pi_{8}}{\vdots}$}\noLine
\UIC{$x - \alpha \in \mathbb{R} \to \lvert x - \alpha \rvert \in \mathbb{R}$}
\AXC{$(x - \alpha \in \mathbb{R} \to \lvert x - \alpha \rvert \in \mathbb{R}) \to \varphi_{57}$\lbl{A1}}
\BIC{$\varphi_{57}$}
\BIC{$x \in \mathbb{R} \to y \in \mathbb{R} \to \alpha \in \mathbb{R} \to \varphi_{27}$}
\end{ptree}
\noindent$\Pi_{10}$:
\begin{ptree}
\AXC{$\stackrel{\Pi_{9}}{\vdots}$}\noLine
\UIC{$x \in \mathbb{R} \to y \in \mathbb{R} \to \alpha \in \mathbb{R} \to \varphi_{27}$}
\AXC{$\varphi_{39}$\lbl{[F20]}}
\AXC{$\varphi_{39} \to \varphi_{48}$\lbl{A2}}
\BIC{$\varphi_{48}$}
\AXC{$\varphi_{48} \to \varphi_{46}$\lbl{A2}}
\BIC{$\varphi_{46}$}
\AXC{$\varphi_{46} \to \varphi_{52}$\lbl{A2}}
\BIC{$\varphi_{52}$}
\AXC{$\varphi_{61}$\lbl{A2}}
\BIC{$\lvert x - \alpha \rvert \in \mathbb{R} \to h \in \mathbb{R} \to \lvert y - \beta \rvert \in \mathbb{R} \to h \in \mathbb{R} \to \varphi_{13}$}
\AXC{$\varphi_{72}$\lbl{A1}}
\BIC{$(x \in \mathbb{R} \to y \in \mathbb{R} \to \alpha \in \mathbb{R} \to \varphi_{27}) \to \varphi_{70}$}
\BIC{$\varphi_{70}$}
\end{ptree}
\noindent$\Pi_{11}$:
\begin{ptree}
\AXC{$x \in \mathbb{R} \to y \in \mathbb{R} \to \varphi_{33}$\lbl{A1}}
\AXC{$\stackrel{\Pi_{5}}{\vdots}$}\noLine
\UIC{$x \in \mathbb{R} \to y \in \mathbb{R} \to \alpha \in \mathbb{R} \to \varphi_{26}$}
\AXC{$x \in \mathbb{R} \to y \in \mathbb{R} \to \varphi_{33}$\lbl{A1}}
\AXC{$\stackrel{\Pi_{10}}{\vdots}$}\noLine
\UIC{$\varphi_{70}$}
\BIC{$(x \in \mathbb{R} \to y \in \mathbb{R} \to \alpha \in \mathbb{R} \to \varphi_{26}) \to \varphi_{67}$}
\BIC{$\varphi_{67}$}
\BIC{$(x \in \mathbb{R} \to y \in \mathbb{R} \to \alpha \in \mathbb{R} \to \varphi_{25}) \to \varphi_{59}$}
\end{ptree}
\noindent$\Pi_{12}$:
\begin{ptree}
\AXC{$x \in \mathbb{R} \to y \in \mathbb{R} \to \alpha \in \mathbb{R} \to \varphi_{24}$\lbl{A1}}
\AXC{$\varphi_{55}$\lbl{A1}}
\BIC{$x \in \mathbb{R} \to y \in \mathbb{R} \to \alpha \in \mathbb{R} \to \varphi_{21}$}
\AXC{$x \in \mathbb{R} \to y \in \mathbb{R} \to \alpha \in \mathbb{R} \to \varphi_{25}$\lbl{A1}}
\AXC{$\stackrel{\Pi_{11}}{\vdots}$}\noLine
\UIC{$(x \in \mathbb{R} \to y \in \mathbb{R} \to \alpha \in \mathbb{R} \to \varphi_{25}) \to \varphi_{59}$}
\BIC{$\varphi_{59}$}
\BIC{$x \in \mathbb{R} \to y \in \mathbb{R} \to \alpha \in \mathbb{R} \to \beta \in \mathbb{R} \to h \in \mathbb{R} \to \varphi_{18}$}
\end{ptree}
\noindent$\Pi_{13}$:
\begin{ptree}
\AXC{$z = w \to \varphi_{9}$\lbl{A6}}
\UIC{$\forall w\,(z = w \to \varphi_{9})$}
\UIC{$\forall z\,\forall w\,(z = w \to \varphi_{9})$}
\AXC{$\forall z\,\forall w\,(z = w \to \varphi_{9}) \to \varphi_{45}$\lbl{A2}}
\BIC{$\varphi_{45}$}
\AXC{$\varphi_{45} \to x - \alpha + (y - \beta) = x + y - (\alpha + \beta) \to \varphi_{43}$\lbl{A2}}
\BIC{$x - \alpha + (y - \beta) = x + y - (\alpha + \beta) \to \varphi_{43}$}
\end{ptree}
\noindent$\Pi_{14}$:
\begin{ptree}
\AXC{$z = w \to \varphi_{20}$\lbl{A6}}
\UIC{$\forall w\,(z = w \to \varphi_{20})$}
\UIC{$\forall z\,\forall w\,(z = w \to \varphi_{20})$}
\AXC{$\forall z\,\forall w\,(z = w \to \varphi_{20}) \to \varphi_{47}$\lbl{A2}}
\BIC{$\varphi_{47}$}
\AXC{$\varphi_{47} \to x - \alpha + (y - \beta) = x + y - (\alpha + \beta) \to \varphi_{44}$\lbl{A2}}
\BIC{$x - \alpha + (y - \beta) = x + y - (\alpha + \beta) \to \varphi_{44}$}
\end{ptree}
\noindent$\Pi_{15}$:
\begin{ptree}
\AXC{$\forall a\,\forall b\,(a \in \mathbb{R} \to b \in \mathbb{R} \to \lvert a + b \rvert \leq \lvert a \rvert + \lvert b \rvert)$\lbl{命題1.2}}
\AXC{$\varphi_{50}$\lbl{A2}}
\BIC{$\varphi_{12}$}
\AXC{$\varphi_{12} \to \varphi_{30}$\lbl{A2}}
\BIC{$\varphi_{30}$}
\AXC{$\varphi_{30} \to \varphi_{65}$\lbl{A1}}
\BIC{$\varphi_{65}$}
\end{ptree}
\noindent$\Pi_{16}$:
\begin{ptree}
\AXC{$\stackrel{\Pi_{3}}{\vdots}$}\noLine
\UIC{$x \in \mathbb{R} \to y \in \mathbb{R} \to \alpha \in \mathbb{R} \to \varphi_{22}$}
\AXC{$\stackrel{\Pi_{7}}{\vdots}$}\noLine
\UIC{$x \in \mathbb{R} \to y \in \mathbb{R} \to \alpha \in \mathbb{R} \to \varphi_{23}$}
\AXC{$\stackrel{\Pi_{15}}{\vdots}$}\noLine
\UIC{$\varphi_{65}$}
\BIC{$(x \in \mathbb{R} \to y \in \mathbb{R} \to \alpha \in \mathbb{R} \to \varphi_{22}) \to \varphi_{53}$}
\BIC{$\varphi_{53}$}
\end{ptree}
\noindent$\Pi_{17}$:
\begin{ptree}
\AXC{$x = x$\lbl{A5}}
\AXC{$x = y \to x = x \to y = x$\lbl{A6}}
\AXC{$(x = y \to x = x \to y = x) \to x = x \to x = y \to y = x$\lbl{A1}}
\BIC{$x = x \to x = y \to y = x$}
\BIC{$x = y \to y = x$}
\UIC{$\forall y\,(x = y \to y = x)$}
\UIC{$\forall x\,\forall y\,(x = y \to y = x)$}
\end{ptree}
\noindent$\Pi_{18}$:
\begin{ptree}
\AXC{$\stackrel{\Pi_{17}}{\vdots}$}\noLine
\UIC{$\forall x\,\forall y\,(x = y \to y = x)$}
\AXC{$\forall x\,\forall y\,(x = y \to y = x) \to \forall y\,(v_{376} = y \to y = v_{376})$\lbl{A2}}
\BIC{$\forall y\,(v_{376} = y \to y = v_{376})$}
\AXC{$\varphi_{1}$\lbl{A2}}
\BIC{$v_{376} = v_{377} \to v_{377} = v_{376}$}
\UIC{$\forall v_{377}\,(v_{376} = v_{377} \to v_{377} = v_{376})$}
\UIC{$\forall v_{376}\,\forall v_{377}\,(v_{376} = v_{377} \to v_{377} = v_{376})$}
\end{ptree}
\noindent$\Pi_{19}$:
\begin{ptree}
\AXC{$\stackrel{\Pi_{18}}{\vdots}$}\noLine
\UIC{$\forall v_{376}\,\forall v_{377}\,(v_{376} = v_{377} \to v_{377} = v_{376})$}
\AXC{$\forall v_{376}\,\forall v_{377}\,(v_{376} = v_{377} \to v_{377} = v_{376}) \to \varphi_{3}$\lbl{A2}}
\BIC{$\varphi_{3}$}
\AXC{$\varphi_{3} \to \varphi_{14}$\lbl{A2}}
\BIC{$\varphi_{14}$}
\end{ptree}
\noindent$\Pi_{20}$:
\begin{ptree}
\AXC{$\forall x\,\forall y\,(\varphi_{31})$\lbl{[F15]}}
\AXC{$\forall x\,\forall y\,(\varphi_{31}) \to \forall y\,(\varphi_{31})$\lbl{A2}}
\BIC{$\forall y\,(\varphi_{31})$}
\AXC{$\forall y\,(\varphi_{31}) \to \varphi_{31}$\lbl{A2}}
\BIC{$\varphi_{31}$}
\AXC{$\varphi_{31} \to \varphi_{32}$\lbl{A2}}
\BIC{$\varphi_{32}$}
\AXC{$\varphi_{32} \to x \in \mathbb{R} \to \varphi_{19}$\lbl{A2}}
\BIC{$x \in \mathbb{R} \to \varphi_{19}$}
\end{ptree}
\noindent$\Pi_{21}$:
\begin{ptree}
\AXC{$x \in \mathbb{R} \to y \in \mathbb{R} \to \varphi_{34}$\lbl{A1}}
\AXC{$x \in \mathbb{R} \to y \in \mathbb{R} \to \varphi_{35}$\lbl{A1}}
\AXC{$\stackrel{\Pi_{20}}{\vdots}$}\noLine
\UIC{$x \in \mathbb{R} \to \varphi_{19}$}
\AXC{$(x \in \mathbb{R} \to \varphi_{19}) \to \varphi_{69}$\lbl{A1}}
\BIC{$\varphi_{69}$}
\BIC{$(x \in \mathbb{R} \to y \in \mathbb{R} \to \varphi_{34}) \to (x \in \mathbb{R} \to y \in \mathbb{R} \to \varphi_{37}) \to \varphi_{58}$}
\BIC{$(x \in \mathbb{R} \to y \in \mathbb{R} \to \varphi_{37}) \to \varphi_{58}$}
\end{ptree}
\noindent$\Pi_{22}$:
\begin{ptree}
\AXC{$\stackrel{\Pi_{19}}{\vdots}$}\noLine
\UIC{$\varphi_{14}$}
\AXC{$x \in \mathbb{R} \to y \in \mathbb{R} \to \varphi_{36}$\lbl{A1}}
\AXC{$x \in \mathbb{R} \to y \in \mathbb{R} \to \varphi_{37}$\lbl{A1}}
\AXC{$\stackrel{\Pi_{21}}{\vdots}$}\noLine
\UIC{$(x \in \mathbb{R} \to y \in \mathbb{R} \to \varphi_{37}) \to \varphi_{58}$}
\BIC{$\varphi_{58}$}
\BIC{$x \in \mathbb{R} \to y \in \mathbb{R} \to \alpha \in \mathbb{R} \to \beta \in \mathbb{R} \to h \in \mathbb{R} \to \varphi_{15}$}
\AXC{$\varphi_{64}$\lbl{A1}}
\BIC{$\varphi_{14} \to x \in \mathbb{R} \to y \in \mathbb{R} \to \alpha \in \mathbb{R} \to \beta \in \mathbb{R} \to h \in \mathbb{R} \to \varphi_{16}$}
\BIC{$x \in \mathbb{R} \to y \in \mathbb{R} \to \alpha \in \mathbb{R} \to \beta \in \mathbb{R} \to h \in \mathbb{R} \to \varphi_{16}$}
\end{ptree}
\noindent$\Pi_{23}$:
\begin{ptree}
\AXC{$\stackrel{\Pi_{14}}{\vdots}$}\noLine
\UIC{$x - \alpha + (y - \beta) = x + y - (\alpha + \beta) \to \varphi_{44}$}
\AXC{$\stackrel{\Pi_{16}}{\vdots}$}\noLine
\UIC{$\varphi_{53}$}
\AXC{$\stackrel{\Pi_{22}}{\vdots}$}\noLine
\UIC{$x \in \mathbb{R} \to y \in \mathbb{R} \to \alpha \in \mathbb{R} \to \beta \in \mathbb{R} \to h \in \mathbb{R} \to \varphi_{16}$}
\AXC{$\varphi_{71}$\lbl{A1}}
\BIC{$\varphi_{53} \to \varphi_{66}$}
\BIC{$\varphi_{66}$}
\BIC{$(x - \alpha + (y - \beta) = x + y - (\alpha + \beta) \to \varphi_{43}) \to \varphi_{54}$}
\end{ptree}
\noindent$\Pi_{24}$:
\begin{ptree}
\AXC{$\stackrel{\Pi_{13}}{\vdots}$}\noLine
\UIC{$x - \alpha + (y - \beta) = x + y - (\alpha + \beta) \to \varphi_{43}$}
\AXC{$\stackrel{\Pi_{23}}{\vdots}$}\noLine
\UIC{$(x - \alpha + (y - \beta) = x + y - (\alpha + \beta) \to \varphi_{43}) \to \varphi_{54}$}
\BIC{$\varphi_{54}$}
\end{ptree}
\noindent$\Pi_{25}$:
\begin{ptree}
\AXC{$\forall x\,\forall y\,\forall z\,(x \leq y \to y < z \to x < z)$\lbl{[O4]}}
\AXC{$\varphi_{5}$\lbl{A2}}
\BIC{$\forall y\,\forall z\,(v_{380} \leq y \to y < z \to v_{380} < z)$}
\AXC{$\varphi_{6}$\lbl{A2}}
\BIC{$\forall z\,(v_{380} \leq v_{381} \to v_{381} < z \to v_{380} < z)$}
\AXC{$\varphi_{17}$\lbl{A2}}
\BIC{$v_{380} \leq v_{381} \to v_{381} < v_{382} \to v_{380} < v_{382}$}
\UIC{$\forall v_{382}\,(v_{380} \leq v_{381} \to v_{381} < v_{382} \to v_{380} < v_{382})$}
\UIC{$\varphi_{2}$}
\UIC{$\forall v_{380}\,(\varphi_{2})$}
\AXC{$\forall v_{380}\,(\varphi_{2}) \to \varphi_{28}$\lbl{A2}}
\BIC{$\varphi_{28}$}
\AXC{$\varphi_{28} \to \varphi_{49}$\lbl{A2}}
\BIC{$\varphi_{49}$}
\end{ptree}
\noindent$\Pi_{26}$:
\begin{ptree}
\AXC{$\stackrel{\Pi_{24}}{\vdots}$}\noLine
\UIC{$\varphi_{54}$}
\AXC{$\stackrel{\Pi_{25}}{\vdots}$}\noLine
\UIC{$\varphi_{49}$}
\AXC{$\varphi_{49} \to \lvert x + y - (\alpha + \beta) \rvert \leq \lvert x - \alpha \rvert + \lvert y - \beta \rvert \to \varphi_{11}$\lbl{A2}}
\BIC{$\lvert x + y - (\alpha + \beta) \rvert \leq \lvert x - \alpha \rvert + \lvert y - \beta \rvert \to \varphi_{11}$}
\AXC{$\varphi_{68}$\lbl{A1}}
\BIC{$\varphi_{54} \to \varphi_{60}$}
\BIC{$\varphi_{60}$}
\end{ptree}
\noindent 主木:
\begin{ptree}
\AXC{$\stackrel{\Pi_{12}}{\vdots}$}\noLine
\UIC{$x \in \mathbb{R} \to y \in \mathbb{R} \to \alpha \in \mathbb{R} \to \beta \in \mathbb{R} \to h \in \mathbb{R} \to \varphi_{18}$}
\AXC{$\stackrel{\Pi_{26}}{\vdots}$}\noLine
\UIC{$\varphi_{60}$}
\BIC{$x \in \mathbb{R} \to y \in \mathbb{R} \to \alpha \in \mathbb{R} \to \beta \in \mathbb{R} \to h \in \mathbb{R} \to \varphi_{10}$}
\UIC{$\forall h\,(x \in \mathbb{R} \to y \in \mathbb{R} \to \alpha \in \mathbb{R} \to \beta \in \mathbb{R} \to h \in \mathbb{R} \to \varphi_{10})$}
\UIC{$\varphi_{51}$}
\UIC{$\forall \alpha\,(\varphi_{51})$}
\UIC{$\forall y\,\forall \alpha\,(\varphi_{51})$}
\UIC{$\forall x\,\forall y\,\forall \alpha\,(\varphi_{51})$}
\end{ptree}
\begin{abbrevlist}
\item $\varphi_{1} :\equiv \forall y\,(v_{376} = y \to y = v_{376}) \to v_{376} = v_{377} \to v_{377} = v_{376}$
\item $\varphi_{2} :\equiv \forall v_{381}\,\forall v_{382}\,(v_{380} \leq v_{381} \to v_{381} < v_{382} \to v_{380} < v_{382})$
\item $\varphi_{3} :\equiv \forall v_{377}\,(x + y - (\alpha + \beta) = v_{377} \to v_{377} = x + y - (\alpha + \beta))$
\item $\varphi_{4} :\equiv \forall v_{378}\,\forall v_{379}\,(v_{378} \in \mathbb{R} \to v_{379} \in \mathbb{R} \to v_{378} - v_{379} \in \mathbb{R})$
\item $\varphi_{5} :\equiv \forall x\,\forall y\,\forall z\,(x \leq y \to y < z \to x < z) \to \forall y\,\forall z\,(v_{380} \leq y \to y < z \to v_{380} < z)$
\item $\varphi_{6} :\equiv \forall y\,\forall z\,(v_{380} \leq y \to y < z \to v_{380} < z) \to \forall z\,(v_{380} \leq v_{381} \to v_{381} < z \to v_{380} < z)$
\item $\varphi_{7} :\equiv \lvert y - \beta \rvert < h \to \lvert x - \alpha + (y - \beta) \rvert \leq \lvert x - \alpha \rvert + \lvert y - \beta \rvert$
\item $\varphi_{8} :\equiv \lvert y - \beta \rvert < h \to \lvert x + y - (\alpha + \beta) \rvert \leq \lvert x - \alpha \rvert + \lvert y - \beta \rvert$
\item $\varphi_{9} :\equiv \lvert z \rvert = \lvert x - \alpha \rvert + \lvert y - \beta \rvert \to \lvert w \rvert = \lvert x - \alpha \rvert + \lvert y - \beta \rvert$
\item $\varphi_{10} :\equiv \lvert x - \alpha \rvert < h \to \lvert y - \beta \rvert < h \to \lvert x + y - (\alpha + \beta) \rvert < h + h$
\item $\varphi_{11} :\equiv \lvert x - \alpha \rvert + \lvert y - \beta \rvert < h + h \to \lvert x + y - (\alpha + \beta) \rvert < h + h$
\item $\varphi_{12} :\equiv \forall b\,(x - \alpha \in \mathbb{R} \to b \in \mathbb{R} \to \lvert x - \alpha + b \rvert \leq \lvert x - \alpha \rvert + \lvert b \rvert)$
\item $\varphi_{13} :\equiv \lvert x - \alpha \rvert < h \to \lvert y - \beta \rvert \leq h \to \lvert x - \alpha \rvert + \lvert y - \beta \rvert < h + h$
\item $\varphi_{14} :\equiv x + y - (\alpha + \beta) = x - \alpha + (y - \beta) \to x - \alpha + (y - \beta) = x + y - (\alpha + \beta)$
\item $\varphi_{15} :\equiv \lvert x - \alpha \rvert < h \to \lvert y - \beta \rvert < h \to x + y - (\alpha + \beta) = x - \alpha + (y - \beta)$
\item $\varphi_{16} :\equiv \lvert x - \alpha \rvert < h \to \lvert y - \beta \rvert < h \to x - \alpha + (y - \beta) = x + y - (\alpha + \beta)$
\item $\varphi_{17} :\equiv \forall z\,(v_{380} \leq v_{381} \to v_{381} < z \to v_{380} < z) \to v_{380} \leq v_{381} \to v_{381} < v_{382} \to v_{380} < v_{382}$
\item $\varphi_{18} :\equiv \lvert x - \alpha \rvert < h \to \lvert y - \beta \rvert < h \to \lvert x - \alpha \rvert + \lvert y - \beta \rvert < h + h$
\item $\varphi_{19} :\equiv y \in \mathbb{R} \to \alpha \in \mathbb{R} \to \beta \in \mathbb{R} \to x + y - (\alpha + \beta) = x - \alpha + (y - \beta)$
\item $\varphi_{20} :\equiv \lvert z \rvert < \lvert x - \alpha \rvert + \lvert y - \beta \rvert \to \lvert w \rvert < \lvert x - \alpha \rvert + \lvert y - \beta \rvert$
\item $\varphi_{21} :\equiv \beta \in \mathbb{R} \to h \in \mathbb{R} \to \lvert x - \alpha \rvert < h \to \lvert y - \beta \rvert < h \to \lvert y - \beta \rvert \leq h$
\item $\varphi_{22} :\equiv \beta \in \mathbb{R} \to h \in \mathbb{R} \to \lvert x - \alpha \rvert < h \to \lvert y - \beta \rvert < h \to y - \beta \in \mathbb{R}$
\item $\varphi_{23} :\equiv \beta \in \mathbb{R} \to h \in \mathbb{R} \to \lvert x - \alpha \rvert < h \to \lvert y - \beta \rvert < h \to x - \alpha \in \mathbb{R}$
\item $\varphi_{24} :\equiv \beta \in \mathbb{R} \to h \in \mathbb{R} \to \lvert x - \alpha \rvert < h \to \lvert y - \beta \rvert < h \to \lvert y - \beta \rvert < h$
\item $\varphi_{25} :\equiv \beta \in \mathbb{R} \to h \in \mathbb{R} \to \lvert x - \alpha \rvert < h \to \lvert y - \beta \rvert < h \to \lvert x - \alpha \rvert < h$
\item $\varphi_{26} :\equiv \beta \in \mathbb{R} \to h \in \mathbb{R} \to \lvert x - \alpha \rvert < h \to \lvert y - \beta \rvert < h \to \lvert y - \beta \rvert \in \mathbb{R}$
\item $\varphi_{27} :\equiv \beta \in \mathbb{R} \to h \in \mathbb{R} \to \lvert x - \alpha \rvert < h \to \lvert y - \beta \rvert < h \to \lvert x - \alpha \rvert \in \mathbb{R}$
\item $\varphi_{28} :\equiv \forall v_{381}\,\forall v_{382}\,(\lvert x + y - (\alpha + \beta) \rvert \leq v_{381} \to v_{381} < v_{382} \to \lvert x + y - (\alpha + \beta) \rvert < v_{382})$
\item $\varphi_{29} :\equiv \forall y\,(x \in \mathbb{R} \to y \in \mathbb{R} \to x - y \in \mathbb{R}) \to x \in \mathbb{R} \to \alpha \in \mathbb{R} \to x - \alpha \in \mathbb{R}$
\item $\varphi_{30} :\equiv x - \alpha \in \mathbb{R} \to y - \beta \in \mathbb{R} \to \lvert x - \alpha + (y - \beta) \rvert \leq \lvert x - \alpha \rvert + \lvert y - \beta \rvert$
\item $\varphi_{31} :\equiv \forall a\,\forall b\,(x \in \mathbb{R} \to y \in \mathbb{R} \to a \in \mathbb{R} \to b \in \mathbb{R} \to x + y - (a + b) = x - a + (y - b))$
\item $\varphi_{32} :\equiv \forall b\,(x \in \mathbb{R} \to y \in \mathbb{R} \to \alpha \in \mathbb{R} \to b \in \mathbb{R} \to x + y - (\alpha + b) = x - \alpha + (y - b))$
\item $\varphi_{33} :\equiv \alpha \in \mathbb{R} \to \beta \in \mathbb{R} \to h \in \mathbb{R} \to \lvert x - \alpha \rvert < h \to \lvert y - \beta \rvert < h \to h \in \mathbb{R}$
\item $\varphi_{34} :\equiv \alpha \in \mathbb{R} \to \beta \in \mathbb{R} \to h \in \mathbb{R} \to \lvert x - \alpha \rvert < h \to \lvert y - \beta \rvert < h \to y \in \mathbb{R}$
\item $\varphi_{35} :\equiv \alpha \in \mathbb{R} \to \beta \in \mathbb{R} \to h \in \mathbb{R} \to \lvert x - \alpha \rvert < h \to \lvert y - \beta \rvert < h \to x \in \mathbb{R}$
\item $\varphi_{36} :\equiv \alpha \in \mathbb{R} \to \beta \in \mathbb{R} \to h \in \mathbb{R} \to \lvert x - \alpha \rvert < h \to \lvert y - \beta \rvert < h \to \beta \in \mathbb{R}$
\item $\varphi_{37} :\equiv \alpha \in \mathbb{R} \to \beta \in \mathbb{R} \to h \in \mathbb{R} \to \lvert x - \alpha \rvert < h \to \lvert y - \beta \rvert < h \to \alpha \in \mathbb{R}$
\item $\varphi_{38} :\equiv \forall x\,\forall y\,(x \in \mathbb{R} \to y \in \mathbb{R} \to x - y \in \mathbb{R}) \to \forall y\,(x \in \mathbb{R} \to y \in \mathbb{R} \to x - y \in \mathbb{R})$
\item $\varphi_{39} :\equiv \forall x\,\forall y\,\forall u\,\forall v\,(x \in \mathbb{R} \to y \in \mathbb{R} \to u \in \mathbb{R} \to v \in \mathbb{R} \to x < y \to u \leq v \to x + u < y + v)$
\item $\varphi_{40} :\equiv \forall v_{379}\,(y \in \mathbb{R} \to v_{379} \in \mathbb{R} \to y - v_{379} \in \mathbb{R}) \to y \in \mathbb{R} \to \beta \in \mathbb{R} \to y - \beta \in \mathbb{R}$
\item $\varphi_{41} :\equiv \forall x\,\forall y\,(x \in \mathbb{R} \to y \in \mathbb{R} \to x - y \in \mathbb{R}) \to \forall y\,(v_{378} \in \mathbb{R} \to y \in \mathbb{R} \to v_{378} - y \in \mathbb{R})$
\item $\varphi_{42} :\equiv \forall y\,(v_{378} \in \mathbb{R} \to y \in \mathbb{R} \to v_{378} - y \in \mathbb{R}) \to v_{378} \in \mathbb{R} \to v_{379} \in \mathbb{R} \to v_{378} - v_{379} \in \mathbb{R}$
\item $\varphi_{43} :\equiv \lvert x - \alpha + (y - \beta) \rvert = \lvert x - \alpha \rvert + \lvert y - \beta \rvert \to \lvert x + y - (\alpha + \beta) \rvert = \lvert x - \alpha \rvert + \lvert y - \beta \rvert$
\item $\varphi_{44} :\equiv \lvert x - \alpha + (y - \beta) \rvert < \lvert x - \alpha \rvert + \lvert y - \beta \rvert \to \lvert x + y - (\alpha + \beta) \rvert < \lvert x - \alpha \rvert + \lvert y - \beta \rvert$
\item $\varphi_{45} :\equiv \forall w\,(x - \alpha + (y - \beta) = w \to \lvert x - \alpha + (y - \beta) \rvert = \lvert x - \alpha \rvert + \lvert y - \beta \rvert \to \lvert w \rvert = \lvert x - \alpha \rvert + \lvert y - \beta \rvert)$
\item $\varphi_{46} :\equiv \forall u\,\forall v\,(\lvert x - \alpha \rvert \in \mathbb{R} \to h \in \mathbb{R} \to u \in \mathbb{R} \to v \in \mathbb{R} \to \lvert x - \alpha \rvert < h \to u \leq v \to \lvert x - \alpha \rvert + u < h + v)$
\item $\varphi_{47} :\equiv \forall w\,(x - \alpha + (y - \beta) = w \to \lvert x - \alpha + (y - \beta) \rvert < \lvert x - \alpha \rvert + \lvert y - \beta \rvert \to \lvert w \rvert < \lvert x - \alpha \rvert + \lvert y - \beta \rvert)$
\item $\varphi_{48} :\equiv \forall y\,\forall u\,\forall v\,(\lvert x - \alpha \rvert \in \mathbb{R} \to y \in \mathbb{R} \to u \in \mathbb{R} \to v \in \mathbb{R} \to \lvert x - \alpha \rvert < y \to u \leq v \to \lvert x - \alpha \rvert + u < y + v)$
\item $\varphi_{49} :\equiv \forall v_{382}\,(\lvert x + y - (\alpha + \beta) \rvert \leq \lvert x - \alpha \rvert + \lvert y - \beta \rvert \to \lvert x - \alpha \rvert + \lvert y - \beta \rvert < v_{382} \to \lvert x + y - (\alpha + \beta) \rvert < v_{382})$
\item $\varphi_{50} :\equiv \forall a\,\forall b\,(a \in \mathbb{R} \to b \in \mathbb{R} \to \lvert a + b \rvert \leq \lvert a \rvert + \lvert b \rvert) \to \varphi_{12}$
\item $\varphi_{51} :\equiv \forall \beta\,\forall h\,(x \in \mathbb{R} \to y \in \mathbb{R} \to \alpha \in \mathbb{R} \to \beta \in \mathbb{R} \to h \in \mathbb{R} \to \varphi_{10})$
\item $\varphi_{52} :\equiv \forall v\,(\lvert x - \alpha \rvert \in \mathbb{R} \to h \in \mathbb{R} \to \lvert y - \beta \rvert \in \mathbb{R} \to v \in \mathbb{R} \to \lvert x - \alpha \rvert < h \to \lvert y - \beta \rvert \leq v \to \lvert x - \alpha \rvert + \lvert y - \beta \rvert < h + v)$
\item $\varphi_{53} :\equiv x \in \mathbb{R} \to y \in \mathbb{R} \to \alpha \in \mathbb{R} \to \beta \in \mathbb{R} \to h \in \mathbb{R} \to \lvert x - \alpha \rvert < h \to \varphi_{7}$
\item $\varphi_{54} :\equiv x \in \mathbb{R} \to y \in \mathbb{R} \to \alpha \in \mathbb{R} \to \beta \in \mathbb{R} \to h \in \mathbb{R} \to \lvert x - \alpha \rvert < h \to \varphi_{8}$
\item $\varphi_{55} :\equiv (x \in \mathbb{R} \to y \in \mathbb{R} \to \alpha \in \mathbb{R} \to \varphi_{24}) \to x \in \mathbb{R} \to y \in \mathbb{R} \to \alpha \in \mathbb{R} \to \varphi_{21}$
\item $\varphi_{56} :\equiv (x \in \mathbb{R} \to y \in \mathbb{R} \to \alpha \in \mathbb{R} \to \varphi_{22}) \to x \in \mathbb{R} \to y \in \mathbb{R} \to \alpha \in \mathbb{R} \to \varphi_{26}$
\item $\varphi_{57} :\equiv (x \in \mathbb{R} \to y \in \mathbb{R} \to \alpha \in \mathbb{R} \to \varphi_{23}) \to x \in \mathbb{R} \to y \in \mathbb{R} \to \alpha \in \mathbb{R} \to \varphi_{27}$
\item $\varphi_{58} :\equiv (x \in \mathbb{R} \to y \in \mathbb{R} \to \varphi_{36}) \to x \in \mathbb{R} \to y \in \mathbb{R} \to \alpha \in \mathbb{R} \to \beta \in \mathbb{R} \to h \in \mathbb{R} \to \varphi_{15}$
\item $\varphi_{59} :\equiv (x \in \mathbb{R} \to y \in \mathbb{R} \to \alpha \in \mathbb{R} \to \varphi_{21}) \to x \in \mathbb{R} \to y \in \mathbb{R} \to \alpha \in \mathbb{R} \to \beta \in \mathbb{R} \to h \in \mathbb{R} \to \varphi_{18}$
\item $\varphi_{60} :\equiv (x \in \mathbb{R} \to y \in \mathbb{R} \to \alpha \in \mathbb{R} \to \beta \in \mathbb{R} \to h \in \mathbb{R} \to \varphi_{18}) \to x \in \mathbb{R} \to y \in \mathbb{R} \to \alpha \in \mathbb{R} \to \beta \in \mathbb{R} \to h \in \mathbb{R} \to \varphi_{10}$
\item $\varphi_{61} :\equiv \varphi_{52} \to \lvert x - \alpha \rvert \in \mathbb{R} \to h \in \mathbb{R} \to \lvert y - \beta \rvert \in \mathbb{R} \to h \in \mathbb{R} \to \varphi_{13}$
\item $\varphi_{62} :\equiv (x \in \mathbb{R} \to y \in \mathbb{R} \to \varphi_{34}) \to (x \in \mathbb{R} \to y \in \mathbb{R} \to \varphi_{36}) \to x \in \mathbb{R} \to y \in \mathbb{R} \to \alpha \in \mathbb{R} \to \varphi_{22}$
\item $\varphi_{63} :\equiv (x \in \mathbb{R} \to y \in \mathbb{R} \to \varphi_{35}) \to (x \in \mathbb{R} \to y \in \mathbb{R} \to \varphi_{37}) \to x \in \mathbb{R} \to y \in \mathbb{R} \to \alpha \in \mathbb{R} \to \varphi_{23}$
\item $\varphi_{64} :\equiv (x \in \mathbb{R} \to y \in \mathbb{R} \to \alpha \in \mathbb{R} \to \beta \in \mathbb{R} \to h \in \mathbb{R} \to \varphi_{15}) \to \varphi_{14} \to x \in \mathbb{R} \to y \in \mathbb{R} \to \alpha \in \mathbb{R} \to \beta \in \mathbb{R} \to h \in \mathbb{R} \to \varphi_{16}$
\item $\varphi_{65} :\equiv (x \in \mathbb{R} \to y \in \mathbb{R} \to \alpha \in \mathbb{R} \to \varphi_{23}) \to (x \in \mathbb{R} \to y \in \mathbb{R} \to \alpha \in \mathbb{R} \to \varphi_{22}) \to \varphi_{53}$
\item $\varphi_{66} :\equiv (x - \alpha + (y - \beta) = x + y - (\alpha + \beta) \to \varphi_{44}) \to (x - \alpha + (y - \beta) = x + y - (\alpha + \beta) \to \varphi_{43}) \to \varphi_{54}$
\item $\varphi_{67} :\equiv (x \in \mathbb{R} \to y \in \mathbb{R} \to \varphi_{33}) \to (x \in \mathbb{R} \to y \in \mathbb{R} \to \alpha \in \mathbb{R} \to \varphi_{25}) \to \varphi_{59}$
\item $\varphi_{68} :\equiv (\lvert x + y - (\alpha + \beta) \rvert \leq \lvert x - \alpha \rvert + \lvert y - \beta \rvert \to \varphi_{11}) \to \varphi_{54} \to \varphi_{60}$
\item $\varphi_{69} :\equiv (x \in \mathbb{R} \to y \in \mathbb{R} \to \varphi_{35}) \to (x \in \mathbb{R} \to y \in \mathbb{R} \to \varphi_{34}) \to (x \in \mathbb{R} \to y \in \mathbb{R} \to \varphi_{37}) \to \varphi_{58}$
\item $\varphi_{70} :\equiv (x \in \mathbb{R} \to y \in \mathbb{R} \to \varphi_{33}) \to (x \in \mathbb{R} \to y \in \mathbb{R} \to \alpha \in \mathbb{R} \to \varphi_{26}) \to \varphi_{67}$
\item $\varphi_{71} :\equiv (x \in \mathbb{R} \to y \in \mathbb{R} \to \alpha \in \mathbb{R} \to \beta \in \mathbb{R} \to h \in \mathbb{R} \to \varphi_{16}) \to \varphi_{53} \to \varphi_{66}$
\item $\varphi_{72} :\equiv (\lvert x - \alpha \rvert \in \mathbb{R} \to h \in \mathbb{R} \to \lvert y - \beta \rvert \in \mathbb{R} \to h \in \mathbb{R} \to \varphi_{13}) \to (x \in \mathbb{R} \to y \in \mathbb{R} \to \alpha \in \mathbb{R} \to \varphi_{27}) \to \varphi_{70}$
\end{abbrevlist}
\end{fullproof}

% 自動生成 (prover/build.js): LimSum  ノード数 252
\begin{fullproof}{定理2.1(1)}{$\forall a\,\forall \alpha\,\forall b\,\forall \beta\,(\mathrm{Conv}(a,\alpha) \to \mathrm{Conv}(b,\beta) \to \mathrm{Conv}(a \oplus b,\alpha + \beta))$}{252}\phantomsection\label{tree:LimSum}
\noindent$\Pi_{1}$:
\begin{ptree}
\AXC{$\varphi_{39} \to \varphi_{35}$\lbl{A2}}
\AXC{$\mathrm{Conv}(a,\alpha) \to \varphi_{88}$\lbl{A1}}
\AXC{$(\mathrm{Conv}(a,\alpha) \to \varphi_{88}) \to \varphi_{89}$\lbl{A1}}
\BIC{$\varphi_{89}$}
\AXC{$\varphi_{89} \to (\varphi_{39} \to \varphi_{35}) \to \varphi_{87}$\lbl{A1}}
\BIC{$(\varphi_{39} \to \varphi_{35}) \to \varphi_{87}$}
\BIC{$\varphi_{87}$}
\end{ptree}
\noindent$\Pi_{2}$:
\begin{ptree}
\AXC{$\forall x\,\forall y\,(x \in \mathbb{R} \to y \in \mathbb{R} \to x \cdot y \in \mathbb{R})$\lbl{(R1)}}
\AXC{$\varphi_{40}$\lbl{A2}}
\BIC{$\forall y\,(\varepsilon \in \mathbb{R} \to y \in \mathbb{R} \to \varepsilon \cdot y \in \mathbb{R})$}
\AXC{$\varphi_{42}$\lbl{A2}}
\BIC{$\varepsilon \in \mathbb{R} \to 2^{-1} \in \mathbb{R} \to \varepsilon \cdot 2^{-1} \in \mathbb{R}$}
\AXC{$(\varepsilon \in \mathbb{R} \to 2^{-1} \in \mathbb{R} \to \varepsilon \cdot 2^{-1} \in \mathbb{R}) \to \varphi_{93}$\lbl{A1}}
\BIC{$\varphi_{93}$}
\end{ptree}
\noindent$\Pi_{3}$:
\begin{ptree}
\AXC{$2^{-1} \in \mathbb{R}$\lbl{[F31]}}
\AXC{$\mathrm{Conv}(a,\alpha) \to \mathrm{Conv}(b,\beta) \to \varepsilon \in \mathbb{R} \to \varepsilon \in \mathbb{R}$\lbl{A1}}
\AXC{$\stackrel{\Pi_{2}}{\vdots}$}\noLine
\UIC{$\varphi_{93}$}
\BIC{$2^{-1} \in \mathbb{R} \to \mathrm{Conv}(a,\alpha) \to \mathrm{Conv}(b,\beta) \to \varphi_{73}$}
\BIC{$\mathrm{Conv}(a,\alpha) \to \mathrm{Conv}(b,\beta) \to \varphi_{73}$}
\end{ptree}
\noindent$\Pi_{4}$:
\begin{ptree}
\AXC{$\forall a\,\forall \alpha\,(\mathrm{Conv}(a,\alpha) \to \alpha \in \mathbb{R})$\lbl{[C2]}}
\AXC{$\varphi_{12}$\lbl{A2}}
\BIC{$\forall \alpha\,(\mathrm{Conv}(b,\alpha) \to \alpha \in \mathbb{R})$}
\AXC{$\varphi_{7}$\lbl{A2}}
\BIC{$\mathrm{Conv}(b,\beta) \to \beta \in \mathbb{R}$}
\AXC{$(\mathrm{Conv}(b,\beta) \to \beta \in \mathbb{R}) \to \varphi_{20}$\lbl{A1}}
\BIC{$\varphi_{20}$}
\end{ptree}
\noindent$\Pi_{5}$:
\begin{ptree}
\AXC{$\mathrm{Conv}(a,\alpha) \to \mathrm{Conv}(b,\beta) \to \mathrm{Conv}(b,\beta)$\lbl{A1}}
\AXC{$\stackrel{\Pi_{4}}{\vdots}$}\noLine
\UIC{$\varphi_{20}$}
\BIC{$\mathrm{Conv}(a,\alpha) \to \mathrm{Conv}(b,\beta) \to \beta \in \mathbb{R}$}
\end{ptree}
\noindent$\Pi_{6}$:
\begin{ptree}
\AXC{$\forall a\,\forall \alpha\,(\mathrm{Conv}(a,\alpha) \to \alpha \in \mathbb{R})$\lbl{[C2]}}
\AXC{$\varphi_{13}$\lbl{A2}}
\BIC{$\forall \alpha\,(\mathrm{Conv}(a,\alpha) \to \alpha \in \mathbb{R})$}
\AXC{$\varphi_{8}$\lbl{A2}}
\BIC{$\mathrm{Conv}(a,\alpha) \to \alpha \in \mathbb{R}$}
\AXC{$(\mathrm{Conv}(a,\alpha) \to \alpha \in \mathbb{R}) \to \varphi_{15}$\lbl{A1}}
\BIC{$\varphi_{15}$}
\end{ptree}
\noindent$\Pi_{7}$:
\begin{ptree}
\AXC{$\mathrm{Conv}(a,\alpha) \to \mathrm{Conv}(a,\alpha)$\lbl{A1}}
\AXC{$\stackrel{\Pi_{6}}{\vdots}$}\noLine
\UIC{$\varphi_{15}$}
\BIC{$\mathrm{Conv}(a,\alpha) \to \mathrm{Conv}(b,\beta) \to \alpha \in \mathbb{R}$}
\end{ptree}
\noindent$\Pi_{8}$:
\begin{ptree}
\AXC{$\forall a\,\forall \alpha\,(\mathrm{Conv}(a,\alpha) \to \mathrm{Seq}(a))$\lbl{[C1]}}
\AXC{$\varphi_{4}$\lbl{A2}}
\BIC{$\forall \alpha\,(\mathrm{Conv}(b,\alpha) \to \mathrm{Seq}(b))$}
\AXC{$\varphi_{1}$\lbl{A2}}
\BIC{$\mathrm{Conv}(b,\beta) \to \mathrm{Seq}(b)$}
\AXC{$(\mathrm{Conv}(b,\beta) \to \mathrm{Seq}(b)) \to \varphi_{17}$\lbl{A1}}
\BIC{$\varphi_{17}$}
\end{ptree}
\noindent$\Pi_{9}$:
\begin{ptree}
\AXC{$\mathrm{Conv}(a,\alpha) \to \mathrm{Conv}(b,\beta) \to \mathrm{Conv}(b,\beta)$\lbl{A1}}
\AXC{$\stackrel{\Pi_{8}}{\vdots}$}\noLine
\UIC{$\varphi_{17}$}
\BIC{$\mathrm{Conv}(a,\alpha) \to \mathrm{Conv}(b,\beta) \to \mathrm{Seq}(b)$}
\end{ptree}
\noindent$\Pi_{10}$:
\begin{ptree}
\AXC{$\forall a\,\forall n\,(\mathrm{Seq}(a) \to n \in \mathbb{N} \to a_{n} \in \mathbb{R})$\lbl{[S1]}}
\AXC{$\varphi_{30}$\lbl{A2}}
\BIC{$\forall n\,(\mathrm{Seq}(b) \to n \in \mathbb{N} \to b_{n} \in \mathbb{R})$}
\AXC{$\varphi_{25}$\lbl{A2}}
\BIC{$\mathrm{Seq}(b) \to n \in \mathbb{N} \to b_{n} \in \mathbb{R}$}
\AXC{$(\mathrm{Seq}(b) \to n \in \mathbb{N} \to b_{n} \in \mathbb{R}) \to \varphi_{101}$\lbl{A1}}
\BIC{$\varphi_{101}$}
\end{ptree}
\noindent$\Pi_{11}$:
\begin{ptree}
\AXC{$\mathrm{Conv}(a,\alpha) \to \varphi_{64}$\lbl{A1}}
\AXC{$\stackrel{\Pi_{9}}{\vdots}$}\noLine
\UIC{$\mathrm{Conv}(a,\alpha) \to \mathrm{Conv}(b,\beta) \to \mathrm{Seq}(b)$}
\AXC{$\stackrel{\Pi_{10}}{\vdots}$}\noLine
\UIC{$\varphi_{101}$}
\BIC{$(\mathrm{Conv}(a,\alpha) \to \varphi_{64}) \to \mathrm{Conv}(a,\alpha) \to \varphi_{77}$}
\BIC{$\mathrm{Conv}(a,\alpha) \to \varphi_{77}$}
\end{ptree}
\noindent$\Pi_{12}$:
\begin{ptree}
\AXC{$\forall a\,\forall \alpha\,(\mathrm{Conv}(a,\alpha) \to \mathrm{Seq}(a))$\lbl{[C1]}}
\AXC{$\varphi_{5}$\lbl{A2}}
\BIC{$\forall \alpha\,(\mathrm{Conv}(a,\alpha) \to \mathrm{Seq}(a))$}
\AXC{$\varphi_{2}$\lbl{A2}}
\BIC{$\mathrm{Conv}(a,\alpha) \to \mathrm{Seq}(a)$}
\AXC{$(\mathrm{Conv}(a,\alpha) \to \mathrm{Seq}(a)) \to \varphi_{9}$\lbl{A1}}
\BIC{$\varphi_{9}$}
\end{ptree}
\noindent$\Pi_{13}$:
\begin{ptree}
\AXC{$\mathrm{Conv}(a,\alpha) \to \mathrm{Conv}(a,\alpha)$\lbl{A1}}
\AXC{$\stackrel{\Pi_{12}}{\vdots}$}\noLine
\UIC{$\varphi_{9}$}
\BIC{$\mathrm{Conv}(a,\alpha) \to \mathrm{Conv}(b,\beta) \to \mathrm{Seq}(a)$}
\end{ptree}
\noindent$\Pi_{14}$:
\begin{ptree}
\AXC{$\forall a\,\forall n\,(\mathrm{Seq}(a) \to n \in \mathbb{N} \to a_{n} \in \mathbb{R})$\lbl{[S1]}}
\AXC{$\varphi_{31}$\lbl{A2}}
\BIC{$\forall n\,(\mathrm{Seq}(a) \to n \in \mathbb{N} \to a_{n} \in \mathbb{R})$}
\AXC{$\varphi_{26}$\lbl{A2}}
\BIC{$\mathrm{Seq}(a) \to n \in \mathbb{N} \to a_{n} \in \mathbb{R}$}
\AXC{$(\mathrm{Seq}(a) \to n \in \mathbb{N} \to a_{n} \in \mathbb{R}) \to \varphi_{102}$\lbl{A1}}
\BIC{$\varphi_{102}$}
\end{ptree}
\noindent$\Pi_{15}$:
\begin{ptree}
\AXC{$\mathrm{Conv}(a,\alpha) \to \varphi_{64}$\lbl{A1}}
\AXC{$\stackrel{\Pi_{13}}{\vdots}$}\noLine
\UIC{$\mathrm{Conv}(a,\alpha) \to \mathrm{Conv}(b,\beta) \to \mathrm{Seq}(a)$}
\AXC{$\stackrel{\Pi_{14}}{\vdots}$}\noLine
\UIC{$\varphi_{102}$}
\BIC{$(\mathrm{Conv}(a,\alpha) \to \varphi_{64}) \to \mathrm{Conv}(a,\alpha) \to \varphi_{78}$}
\BIC{$\mathrm{Conv}(a,\alpha) \to \varphi_{78}$}
\end{ptree}
\noindent$\Pi_{16}$:
\begin{ptree}
\AXC{$\varphi_{57}$\lbl{[C11]}}
\AXC{$\varphi_{57} \to \varphi_{70}$\lbl{A2}}
\BIC{$\varphi_{70}$}
\AXC{$\varphi_{70} \to \forall \alpha\,\forall \beta\,(\varphi_{76})$\lbl{A2}}
\BIC{$\forall \alpha\,\forall \beta\,(\varphi_{76})$}
\AXC{$\forall \alpha\,\forall \beta\,(\varphi_{76}) \to \forall \beta\,(\varphi_{76})$\lbl{A2}}
\BIC{$\forall \beta\,(\varphi_{76})$}
\AXC{$\forall \beta\,(\varphi_{76}) \to \varphi_{76}$\lbl{A2}}
\BIC{$\varphi_{76}$}
\end{ptree}
\noindent$\Pi_{17}$:
\begin{ptree}
\AXC{$\stackrel{\Pi_{11}}{\vdots}$}\noLine
\UIC{$\mathrm{Conv}(a,\alpha) \to \varphi_{77}$}
\AXC{$\stackrel{\Pi_{15}}{\vdots}$}\noLine
\UIC{$\mathrm{Conv}(a,\alpha) \to \varphi_{78}$}
\AXC{$\stackrel{\Pi_{16}}{\vdots}$}\noLine
\UIC{$\varphi_{76}$}
\AXC{$\varphi_{76} \to a_{n} \in \mathbb{R} \to b_{n} \in \mathbb{R} \to \varphi_{80}$\lbl{A2}}
\BIC{$a_{n} \in \mathbb{R} \to b_{n} \in \mathbb{R} \to \varphi_{80}$}
\AXC{$\varphi_{113}$\lbl{A1}}
\BIC{$(\mathrm{Conv}(a,\alpha) \to \varphi_{78}) \to \varphi_{112}$}
\BIC{$\varphi_{112}$}
\BIC{$(\mathrm{Conv}(a,\alpha) \to \mathrm{Conv}(b,\beta) \to \alpha \in \mathbb{R}) \to \varphi_{111}$}
\end{ptree}
\noindent$\Pi_{18}$:
\begin{ptree}
\AXC{$\stackrel{\Pi_{7}}{\vdots}$}\noLine
\UIC{$\mathrm{Conv}(a,\alpha) \to \mathrm{Conv}(b,\beta) \to \alpha \in \mathbb{R}$}
\AXC{$\stackrel{\Pi_{17}}{\vdots}$}\noLine
\UIC{$(\mathrm{Conv}(a,\alpha) \to \mathrm{Conv}(b,\beta) \to \alpha \in \mathbb{R}) \to \varphi_{111}$}
\BIC{$\varphi_{111}$}
\end{ptree}
\noindent$\Pi_{19}$:
\begin{ptree}
\AXC{$\stackrel{\Pi_{3}}{\vdots}$}\noLine
\UIC{$\mathrm{Conv}(a,\alpha) \to \mathrm{Conv}(b,\beta) \to \varphi_{73}$}
\AXC{$\stackrel{\Pi_{5}}{\vdots}$}\noLine
\UIC{$\mathrm{Conv}(a,\alpha) \to \mathrm{Conv}(b,\beta) \to \beta \in \mathbb{R}$}
\AXC{$\stackrel{\Pi_{18}}{\vdots}$}\noLine
\UIC{$\varphi_{111}$}
\BIC{$(\mathrm{Conv}(a,\alpha) \to \mathrm{Conv}(b,\beta) \to \varphi_{73}) \to \varphi_{109}$}
\BIC{$\varphi_{109}$}
\end{ptree}
\noindent$\Pi_{20}$:
\begin{ptree}
\AXC{$\stackrel{\Pi_{1}}{\vdots}$}\noLine
\UIC{$\varphi_{87}$}
\AXC{$\varphi_{87} \to \mathrm{Conv}(a,\alpha) \to \mathrm{Conv}(b,\beta) \to \varepsilon \in \mathbb{R} \to \varphi_{72}$\lbl{A1}}
\BIC{$\mathrm{Conv}(a,\alpha) \to \mathrm{Conv}(b,\beta) \to \varepsilon \in \mathbb{R} \to \varphi_{72}$}
\AXC{$\stackrel{\Pi_{19}}{\vdots}$}\noLine
\UIC{$\varphi_{109}$}
\BIC{$(\mathrm{Conv}(a,\alpha) \to \mathrm{Conv}(b,\beta) \to \varepsilon \in \mathbb{R} \to \varphi_{71}) \to \varphi_{86}$}
\end{ptree}
\noindent$\Pi_{21}$:
\begin{ptree}
\AXC{$\stackrel{\Pi_{1}}{\vdots}$}\noLine
\UIC{$\varphi_{87}$}
\AXC{$\varphi_{87} \to \mathrm{Conv}(a,\alpha) \to \mathrm{Conv}(b,\beta) \to \varepsilon \in \mathbb{R} \to \varphi_{71}$\lbl{A1}}
\BIC{$\mathrm{Conv}(a,\alpha) \to \mathrm{Conv}(b,\beta) \to \varepsilon \in \mathbb{R} \to \varphi_{71}$}
\AXC{$\stackrel{\Pi_{20}}{\vdots}$}\noLine
\UIC{$(\mathrm{Conv}(a,\alpha) \to \mathrm{Conv}(b,\beta) \to \varepsilon \in \mathbb{R} \to \varphi_{71}) \to \varphi_{86}$}
\BIC{$\varphi_{86}$}
\end{ptree}
\noindent$\Pi_{22}$:
\begin{ptree}
\AXC{$\forall e \in \mathbb{R}\,(e \cdot 2^{-1} + e \cdot 2^{-1} = e)$\lbl{[F36]}}
\AXC{$\varphi_{37}$\lbl{A2}}
\BIC{$\varepsilon \in \mathbb{R} \to \varepsilon \cdot 2^{-1} + \varepsilon \cdot 2^{-1} = \varepsilon$}
\AXC{$(\varepsilon \in \mathbb{R} \to \varepsilon \cdot 2^{-1} + \varepsilon \cdot 2^{-1} = \varepsilon) \to \varphi_{92}$\lbl{A1}}
\BIC{$\varphi_{92}$}
\end{ptree}
\noindent$\Pi_{23}$:
\begin{ptree}
\AXC{$\stackrel{\Pi_{21}}{\vdots}$}\noLine
\UIC{$\varphi_{86}$}
\AXC{$\mathrm{Conv}(a,\alpha) \to \mathrm{Conv}(b,\beta) \to \varepsilon \in \mathbb{R} \to \varepsilon \in \mathbb{R}$\lbl{A1}}
\AXC{$\stackrel{\Pi_{22}}{\vdots}$}\noLine
\UIC{$\varphi_{92}$}
\BIC{$\mathrm{Conv}(a,\alpha) \to \mathrm{Conv}(b,\beta) \to \varepsilon \in \mathbb{R} \to \varphi_{69}$}
\AXC{$\varphi_{110}$\lbl{A1}}
\BIC{$\varphi_{86} \to \varphi_{98}$}
\BIC{$\varphi_{98}$}
\end{ptree}
\noindent$\Pi_{24}$:
\begin{ptree}
\AXC{$x = x$\lbl{A5}}
\AXC{$x = y \to x = x \to y = x$\lbl{A6}}
\AXC{$(x = y \to x = x \to y = x) \to x = x \to x = y \to y = x$\lbl{A1}}
\BIC{$x = x \to x = y \to y = x$}
\BIC{$x = y \to y = x$}
\UIC{$\forall y\,(x = y \to y = x)$}
\UIC{$\forall x\,\forall y\,(x = y \to y = x)$}
\AXC{$\varphi_{14}$\lbl{A2}}
\BIC{$\forall y\,(a \oplus b(n) = y \to y = a \oplus b(n))$}
\AXC{$\varphi_{46}$\lbl{A2}}
\BIC{$a \oplus b(n) = a_{n} + b_{n} \to a_{n} + b_{n} = a \oplus b(n)$}
\end{ptree}
\noindent$\Pi_{25}$:
\begin{ptree}
\AXC{$\forall a\,\forall b\,(\mathrm{Seq}(a) \to \varphi_{24})$\lbl{[S2]}}
\AXC{$\forall a\,\forall b\,(\mathrm{Seq}(a) \to \varphi_{24}) \to \forall b\,(\mathrm{Seq}(a) \to \varphi_{24})$\lbl{A2}}
\BIC{$\forall b\,(\mathrm{Seq}(a) \to \varphi_{24})$}
\end{ptree}
\noindent$\Pi_{26}$:
\begin{ptree}
\AXC{$\stackrel{\Pi_{25}}{\vdots}$}\noLine
\UIC{$\forall b\,(\mathrm{Seq}(a) \to \varphi_{24})$}
\AXC{$\forall b\,(\mathrm{Seq}(a) \to \varphi_{24}) \to \mathrm{Seq}(a) \to \varphi_{24}$\lbl{A2}}
\BIC{$\mathrm{Seq}(a) \to \varphi_{24}$}
\AXC{$\varphi_{90}$\lbl{A1}}
\BIC{$(\mathrm{Conv}(a,\alpha) \to \mathrm{Conv}(b,\beta) \to \mathrm{Seq}(a)) \to \varphi_{48}$}
\end{ptree}
\noindent$\Pi_{27}$:
\begin{ptree}
\AXC{$\stackrel{\Pi_{13}}{\vdots}$}\noLine
\UIC{$\mathrm{Conv}(a,\alpha) \to \mathrm{Conv}(b,\beta) \to \mathrm{Seq}(a)$}
\AXC{$\stackrel{\Pi_{26}}{\vdots}$}\noLine
\UIC{$(\mathrm{Conv}(a,\alpha) \to \mathrm{Conv}(b,\beta) \to \mathrm{Seq}(a)) \to \varphi_{48}$}
\BIC{$\varphi_{48}$}
\end{ptree}
\noindent$\Pi_{28}$:
\begin{ptree}
\AXC{$\stackrel{\Pi_{9}}{\vdots}$}\noLine
\UIC{$\mathrm{Conv}(a,\alpha) \to \mathrm{Conv}(b,\beta) \to \mathrm{Seq}(b)$}
\AXC{$\stackrel{\Pi_{27}}{\vdots}$}\noLine
\UIC{$\varphi_{48}$}
\BIC{$\mathrm{Conv}(a,\alpha) \to \varphi_{28}$}
\end{ptree}
\noindent$\Pi_{29}$:
\begin{ptree}
\AXC{$\varphi_{43}$\lbl{A2}}
\AXC{$\stackrel{\Pi_{28}}{\vdots}$}\noLine
\UIC{$\mathrm{Conv}(a,\alpha) \to \varphi_{28}$}
\AXC{$(\mathrm{Conv}(a,\alpha) \to \varphi_{28}) \to \varphi_{23}$\lbl{A1}}
\BIC{$\varphi_{23}$}
\AXC{$\varphi_{23} \to \varphi_{43} \to \varphi_{21}$\lbl{A1}}
\BIC{$\varphi_{43} \to \varphi_{21}$}
\BIC{$\varphi_{21}$}
\AXC{$\varphi_{21} \to \varphi_{100}$\lbl{A1}}
\BIC{$\varphi_{100}$}
\end{ptree}
\noindent$\Pi_{30}$:
\begin{ptree}
\AXC{$\mathrm{Conv}(a,\alpha) \to \varphi_{64}$\lbl{A1}}
\AXC{$\stackrel{\Pi_{29}}{\vdots}$}\noLine
\UIC{$\varphi_{100}$}
\BIC{$\mathrm{Conv}(a,\alpha) \to \mathrm{Conv}(b,\beta) \to \varepsilon \in \mathbb{R} \to \varphi_{74}$}
\AXC{$(\mathrm{Conv}(a,\alpha) \to \mathrm{Conv}(b,\beta) \to \varepsilon \in \mathbb{R} \to \varphi_{74}) \to \varphi_{95}$\lbl{A1}}
\BIC{$\varphi_{95}$}
\end{ptree}
\noindent$\Pi_{31}$:
\begin{ptree}
\AXC{$\stackrel{\Pi_{24}}{\vdots}$}\noLine
\UIC{$a \oplus b(n) = a_{n} + b_{n} \to a_{n} + b_{n} = a \oplus b(n)$}
\AXC{$\stackrel{\Pi_{30}}{\vdots}$}\noLine
\UIC{$\varphi_{95}$}
\BIC{$\mathrm{Conv}(a,\alpha) \to \mathrm{Conv}(b,\beta) \to \varepsilon \in \mathbb{R} \to \varphi_{75}$}
\AXC{$(\mathrm{Conv}(a,\alpha) \to \mathrm{Conv}(b,\beta) \to \varepsilon \in \mathbb{R} \to \varphi_{75}) \to \varphi_{108}$\lbl{A1}}
\BIC{$\varphi_{108}$}
\end{ptree}
\noindent$\Pi_{32}$:
\begin{ptree}
\AXC{$z = w \to \varphi_{29}$\lbl{A6}}
\UIC{$\forall w\,(z = w \to \varphi_{29})$}
\UIC{$\forall z\,\forall w\,(z = w \to \varphi_{29})$}
\AXC{$\forall z\,\forall w\,(z = w \to \varphi_{29}) \to \varphi_{56}$\lbl{A2}}
\BIC{$\varphi_{56}$}
\AXC{$\varphi_{56} \to \varepsilon \cdot 2^{-1} + \varepsilon \cdot 2^{-1} = \varepsilon \to \varphi_{45}$\lbl{A2}}
\BIC{$\varepsilon \cdot 2^{-1} + \varepsilon \cdot 2^{-1} = \varepsilon \to \varphi_{45}$}
\AXC{$\stackrel{\Pi_{23}}{\vdots}$}\noLine
\UIC{$\varphi_{98}$}
\BIC{$\mathrm{Conv}(a,\alpha) \to \mathrm{Conv}(b,\beta) \to \varepsilon \in \mathbb{R} \to 0 < \varepsilon \to \varphi_{66}$}
\AXC{$\stackrel{\Pi_{31}}{\vdots}$}\noLine
\UIC{$\varphi_{108}$}
\BIC{$\varphi_{97}$}
\end{ptree}
\noindent$\Pi_{33}$:
\begin{ptree}
\AXC{$z = w \to \lvert z - (\alpha + \beta) \rvert < \varepsilon \to \lvert w - (\alpha + \beta) \rvert < \varepsilon$\lbl{A6}}
\UIC{$\varphi_{11}$}
\UIC{$\forall z\,(\varphi_{11})$}
\AXC{$\forall z\,(\varphi_{11}) \to \varphi_{41}$\lbl{A2}}
\BIC{$\varphi_{41}$}
\AXC{$\varphi_{41} \to a_{n} + b_{n} = a \oplus b(n) \to \varphi_{33}$\lbl{A2}}
\BIC{$a_{n} + b_{n} = a \oplus b(n) \to \varphi_{33}$}
\AXC{$\stackrel{\Pi_{32}}{\vdots}$}\noLine
\UIC{$\varphi_{97}$}
\BIC{$\mathrm{Conv}(a,\alpha) \to \mathrm{Conv}(b,\beta) \to \varepsilon \in \mathbb{R} \to 0 < \varepsilon \to \varphi_{61}$}
\end{ptree}
\noindent$\Pi_{34}$:
\begin{ptree}
\AXC{$\varphi_{81} \to \varphi_{82}$\lbl{A3}}
\AXC{$\stackrel{\Pi_{33}}{\vdots}$}\noLine
\UIC{$\mathrm{Conv}(a,\alpha) \to \mathrm{Conv}(b,\beta) \to \varepsilon \in \mathbb{R} \to 0 < \varepsilon \to \varphi_{61}$}
\AXC{$\varphi_{103}$\lbl{A1}}
\BIC{$\varphi_{59} \to n \in \mathbb{N} \to N \leq n \to \lvert a \oplus b(n) - (\alpha + \beta) \rvert < \varepsilon$}
\UIC{$\varphi_{81}$}
\AXC{$\varphi_{81} \to (\varphi_{81} \to \varphi_{82}) \to \varphi_{106}$\lbl{A1}}
\BIC{$(\varphi_{81} \to \varphi_{82}) \to \varphi_{106}$}
\BIC{$\varphi_{106}$}
\end{ptree}
\noindent$\Pi_{35}$:
\begin{ptree}
\AXC{$\mathrm{Conv}(a,\alpha) \to \varphi_{88}$\lbl{A1}}
\AXC{$(\mathrm{Conv}(a,\alpha) \to \varphi_{88}) \to \varphi_{65}$\lbl{A1}}
\BIC{$\varphi_{65}$}
\AXC{$\varphi_{65} \to \varphi_{104}$\lbl{A1}}
\BIC{$\varphi_{104}$}
\end{ptree}
\noindent$\Pi_{36}$:
\begin{ptree}
\AXC{$\forall N\,(\neg \varphi_{19}) \to \neg \varphi_{19}$\lbl{A2}}
\AXC{$\varphi_{59} \to \forall n\,(\varphi_{59})$\lbl{A4}}
\AXC{$\stackrel{\Pi_{34}}{\vdots}$}\noLine
\UIC{$\varphi_{106}$}
\BIC{$\mathrm{Conv}(a,\alpha) \to \mathrm{Conv}(b,\beta) \to \varepsilon \in \mathbb{R} \to 0 < \varepsilon \to \varphi_{67}$}
\AXC{$\stackrel{\Pi_{35}}{\vdots}$}\noLine
\UIC{$\varphi_{104}$}
\BIC{$\varphi_{84}$}
\AXC{$\varphi_{84} \to (\forall N\,(\neg \varphi_{19}) \to \neg \varphi_{19}) \to \varphi_{85}$\lbl{A1}}
\BIC{$(\forall N\,(\neg \varphi_{19}) \to \neg \varphi_{19}) \to \varphi_{85}$}
\BIC{$\varphi_{85}$}
\end{ptree}
\noindent$\Pi_{37}$:
\begin{ptree}
\AXC{$\forall e \in \mathbb{R}\,(0 < e \to 0 < e \cdot 2^{-1})$\lbl{[F35]}}
\AXC{$\varphi_{32}$\lbl{A2}}
\BIC{$\varepsilon \in \mathbb{R} \to 0 < \varepsilon \to 0 < \varepsilon \cdot 2^{-1}$}
\AXC{$(\varepsilon \in \mathbb{R} \to 0 < \varepsilon \to 0 < \varepsilon \cdot 2^{-1}) \to \varphi_{79}$\lbl{A1}}
\BIC{$\varphi_{79}$}
\end{ptree}
\noindent$\Pi_{38}$:
\begin{ptree}
\AXC{$\varphi_{10}$\lbl{A1}}
\AXC{$\mathrm{Conv}(a,\alpha) \to \mathrm{Conv}(b,\beta) \to \varepsilon \in \mathbb{R} \to \varepsilon \in \mathbb{R}$\lbl{A1}}
\AXC{$\stackrel{\Pi_{37}}{\vdots}$}\noLine
\UIC{$\varphi_{79}$}
\BIC{$\varphi_{10} \to \varphi_{16}$}
\BIC{$\varphi_{16}$}
\end{ptree}
\noindent$\Pi_{39}$:
\begin{ptree}
\AXC{$\forall x\,\forall y\,(x \in \mathbb{R} \to y \in \mathbb{R} \to x \cdot y \in \mathbb{R})$\lbl{(R1)}}
\AXC{$\varphi_{40}$\lbl{A2}}
\BIC{$\forall y\,(\varepsilon \in \mathbb{R} \to y \in \mathbb{R} \to \varepsilon \cdot y \in \mathbb{R})$}
\AXC{$\varphi_{42}$\lbl{A2}}
\BIC{$\varepsilon \in \mathbb{R} \to 2^{-1} \in \mathbb{R} \to \varepsilon \cdot 2^{-1} \in \mathbb{R}$}
\AXC{$(\varepsilon \in \mathbb{R} \to 2^{-1} \in \mathbb{R} \to \varepsilon \cdot 2^{-1} \in \mathbb{R}) \to \varphi_{54}$\lbl{A1}}
\BIC{$\varphi_{54}$}
\end{ptree}
\noindent$\Pi_{40}$:
\begin{ptree}
\AXC{$2^{-1} \in \mathbb{R}$\lbl{[F31]}}
\AXC{$\mathrm{Conv}(a,\alpha) \to \mathrm{Conv}(b,\beta) \to \varepsilon \in \mathbb{R} \to \varepsilon \in \mathbb{R}$\lbl{A1}}
\AXC{$\stackrel{\Pi_{39}}{\vdots}$}\noLine
\UIC{$\varphi_{54}$}
\BIC{$2^{-1} \in \mathbb{R} \to \varphi_{18}$}
\BIC{$\varphi_{18}$}
\end{ptree}
\noindent$\Pi_{41}$:
\begin{ptree}
\AXC{$\forall a\,\forall \alpha\,\forall b\,\forall \beta\,(\varphi_{55})$\lbl{[C9]}}
\AXC{$\forall a\,\forall \alpha\,\forall b\,\forall \beta\,(\varphi_{55}) \to \forall \alpha\,\forall b\,\forall \beta\,(\varphi_{55})$\lbl{A2}}
\BIC{$\forall \alpha\,\forall b\,\forall \beta\,(\varphi_{55})$}
\AXC{$\forall \alpha\,\forall b\,\forall \beta\,(\varphi_{55}) \to \forall b\,\forall \beta\,(\varphi_{55})$\lbl{A2}}
\BIC{$\forall b\,\forall \beta\,(\varphi_{55})$}
\end{ptree}
\noindent$\Pi_{42}$:
\begin{ptree}
\AXC{$\stackrel{\Pi_{41}}{\vdots}$}\noLine
\UIC{$\forall b\,\forall \beta\,(\varphi_{55})$}
\AXC{$\forall b\,\forall \beta\,(\varphi_{55}) \to \forall \beta\,(\varphi_{55})$\lbl{A2}}
\BIC{$\forall \beta\,(\varphi_{55})$}
\AXC{$\forall \beta\,(\varphi_{55}) \to \varphi_{55}$\lbl{A2}}
\BIC{$\varphi_{55}$}
\AXC{$\varphi_{55} \to \mathrm{Conv}(a,\alpha) \to \varphi_{62}$\lbl{A2}}
\BIC{$\mathrm{Conv}(a,\alpha) \to \varphi_{62}$}
\AXC{$\varphi_{107}$\lbl{A1}}
\BIC{$(\mathrm{Conv}(a,\alpha) \to \mathrm{Conv}(a,\alpha)) \to \varphi_{96}$}
\end{ptree}
\noindent$\Pi_{43}$:
\begin{ptree}
\AXC{$\mathrm{Conv}(a,\alpha) \to \mathrm{Conv}(a,\alpha)$\lbl{A1}}
\AXC{$\stackrel{\Pi_{42}}{\vdots}$}\noLine
\UIC{$(\mathrm{Conv}(a,\alpha) \to \mathrm{Conv}(a,\alpha)) \to \varphi_{96}$}
\BIC{$\varphi_{96}$}
\end{ptree}
\noindent$\Pi_{44}$:
\begin{ptree}
\AXC{$\stackrel{\Pi_{38}}{\vdots}$}\noLine
\UIC{$\varphi_{16}$}
\AXC{$\stackrel{\Pi_{40}}{\vdots}$}\noLine
\UIC{$\varphi_{18}$}
\AXC{$\mathrm{Conv}(a,\alpha) \to \mathrm{Conv}(b,\beta) \to \mathrm{Conv}(b,\beta)$\lbl{A1}}
\AXC{$\stackrel{\Pi_{43}}{\vdots}$}\noLine
\UIC{$\varphi_{96}$}
\BIC{$\varphi_{18} \to \varphi_{16} \to \varphi_{60}$}
\BIC{$\varphi_{16} \to \varphi_{60}$}
\BIC{$\varphi_{60}$}
\end{ptree}
\noindent$\Pi_{45}$:
\begin{ptree}
\AXC{$\varphi_{105}$\lbl{A3}}
\AXC{$\stackrel{\Pi_{36}}{\vdots}$}\noLine
\UIC{$\varphi_{85}$}
\AXC{$\varphi_{85} \to \varphi_{47} \to \neg (N \in \mathbb{N} \land \varphi_{39})$\lbl{A1}}
\BIC{$\varphi_{47} \to \neg (N \in \mathbb{N} \land \varphi_{39})$}
\UIC{$\forall N\,(\varphi_{47} \to \neg (N \in \mathbb{N} \land \varphi_{39}))$}
\AXC{$\stackrel{\Pi_{44}}{\vdots}$}\noLine
\UIC{$\varphi_{60}$}
\AXC{$\varphi_{60} \to \forall N\,(\varphi_{47} \to \neg (N \in \mathbb{N} \land \varphi_{39})) \to \varphi_{105} \to \varphi_{99}$\lbl{A1}}
\BIC{$\forall N\,(\varphi_{47} \to \neg (N \in \mathbb{N} \land \varphi_{39})) \to \varphi_{105} \to \varphi_{99}$}
\BIC{$\varphi_{105} \to \varphi_{99}$}
\BIC{$\varphi_{99}$}
\end{ptree}
\noindent$\Pi_{46}$:
\begin{ptree}
\AXC{$\varphi_{6}$\lbl{A4}}
\AXC{$\varphi_{49} \to \varphi_{52}$\lbl{A3}}
\AXC{$\varphi_{47} \to \forall N\,(\varphi_{47})$\lbl{A4}}
\AXC{$\stackrel{\Pi_{45}}{\vdots}$}\noLine
\UIC{$\varphi_{99}$}
\BIC{$\mathrm{Conv}(a,\alpha) \to \mathrm{Conv}(b,\beta) \to \varepsilon \in \mathbb{R} \to 0 < \varepsilon \to \varphi_{22}$}
\AXC{$\varphi_{91}$\lbl{A1}}
\BIC{$\mathrm{Conv}(a,\alpha) \land \mathrm{Conv}(b,\beta) \to \varepsilon \in \mathbb{R} \to 0 < \varepsilon \to \varphi_{22}$}
\UIC{$\varphi_{49}$}
\AXC{$\varphi_{49} \to (\varphi_{49} \to \varphi_{52}) \to \varphi_{6} \to \varphi_{50}$\lbl{A1}}
\BIC{$(\varphi_{49} \to \varphi_{52}) \to \varphi_{6} \to \varphi_{50}$}
\BIC{$\varphi_{6} \to \varphi_{50}$}
\BIC{$\varphi_{50}$}
\end{ptree}
\noindent$\Pi_{47}$:
\begin{ptree}
\AXC{$\forall x\,\forall y\,(x \in \mathbb{R} \to y \in \mathbb{R} \to x + y \in \mathbb{R})$\lbl{(R1)}}
\AXC{$\varphi_{38}$\lbl{A2}}
\BIC{$\forall y\,(\alpha \in \mathbb{R} \to y \in \mathbb{R} \to \alpha + y \in \mathbb{R})$}
\AXC{$\varphi_{36}$\lbl{A2}}
\BIC{$\alpha \in \mathbb{R} \to \beta \in \mathbb{R} \to \alpha + \beta \in \mathbb{R}$}
\AXC{$\varphi_{63}$\lbl{A1}}
\BIC{$(\mathrm{Conv}(a,\alpha) \to \mathrm{Conv}(b,\beta) \to \alpha \in \mathbb{R}) \to \varphi_{27}$}
\end{ptree}
\noindent$\Pi_{48}$:
\begin{ptree}
\AXC{$\stackrel{\Pi_{7}}{\vdots}$}\noLine
\UIC{$\mathrm{Conv}(a,\alpha) \to \mathrm{Conv}(b,\beta) \to \alpha \in \mathbb{R}$}
\AXC{$\stackrel{\Pi_{47}}{\vdots}$}\noLine
\UIC{$(\mathrm{Conv}(a,\alpha) \to \mathrm{Conv}(b,\beta) \to \alpha \in \mathbb{R}) \to \varphi_{27}$}
\BIC{$\varphi_{27}$}
\end{ptree}
\noindent$\Pi_{49}$:
\begin{ptree}
\AXC{$\stackrel{\Pi_{28}}{\vdots}$}\noLine
\UIC{$\mathrm{Conv}(a,\alpha) \to \varphi_{28}$}
\AXC{$\varphi_{53}$\lbl{A1}}
\BIC{$\mathrm{Conv}(a,\alpha) \to \mathrm{Conv}(b,\beta) \to \mathrm{Seq}(a \oplus b)$}
\AXC{$\varphi_{51}$\lbl{DefConv}}
\AXC{$\varphi_{51} \to \varphi_{58}$\lbl{A2}}
\BIC{$\varphi_{58}$}
\AXC{$\varphi_{58} \to \varphi_{68}$\lbl{A2}}
\BIC{$\varphi_{68}$}
\AXC{$\varphi_{68} \to \varphi_{94}$\lbl{A1}}
\BIC{$\varphi_{94}$}
\BIC{$\varphi_{83}$}
\end{ptree}
\noindent 主木:
\begin{ptree}
\AXC{$\stackrel{\Pi_{46}}{\vdots}$}\noLine
\UIC{$\varphi_{50}$}
\AXC{$\stackrel{\Pi_{5}}{\vdots}$}\noLine
\UIC{$\mathrm{Conv}(a,\alpha) \to \mathrm{Conv}(b,\beta) \to \beta \in \mathbb{R}$}
\AXC{$\stackrel{\Pi_{48}}{\vdots}$}\noLine
\UIC{$\varphi_{27}$}
\BIC{$\mathrm{Conv}(a,\alpha) \to \mathrm{Conv}(b,\beta) \to \alpha + \beta \in \mathbb{R}$}
\AXC{$\stackrel{\Pi_{49}}{\vdots}$}\noLine
\UIC{$\varphi_{83}$}
\BIC{$\varphi_{50} \to \varphi_{3}$}
\BIC{$\varphi_{3}$}
\UIC{$\forall \beta\,(\varphi_{3})$}
\UIC{$\forall b\,\forall \beta\,(\varphi_{3})$}
\UIC{$\forall \alpha\,\forall b\,\forall \beta\,(\varphi_{3})$}
\UIC{$\forall a\,\forall \alpha\,\forall b\,\forall \beta\,(\varphi_{3})$}
\end{ptree}
\begin{abbrevlist}
\item $\varphi_{1} :\equiv \forall \alpha\,(\mathrm{Conv}(b,\alpha) \to \mathrm{Seq}(b)) \to \mathrm{Conv}(b,\beta) \to \mathrm{Seq}(b)$
\item $\varphi_{2} :\equiv \forall \alpha\,(\mathrm{Conv}(a,\alpha) \to \mathrm{Seq}(a)) \to \mathrm{Conv}(a,\alpha) \to \mathrm{Seq}(a)$
\item $\varphi_{3} :\equiv \mathrm{Conv}(a,\alpha) \to \mathrm{Conv}(b,\beta) \to \mathrm{Conv}(a \oplus b,\alpha + \beta)$
\item $\varphi_{4} :\equiv \forall a\,\forall \alpha\,(\mathrm{Conv}(a,\alpha) \to \mathrm{Seq}(a)) \to \forall \alpha\,(\mathrm{Conv}(b,\alpha) \to \mathrm{Seq}(b))$
\item $\varphi_{5} :\equiv \forall a\,\forall \alpha\,(\mathrm{Conv}(a,\alpha) \to \mathrm{Seq}(a)) \to \forall \alpha\,(\mathrm{Conv}(a,\alpha) \to \mathrm{Seq}(a))$
\item $\varphi_{6} :\equiv \mathrm{Conv}(a,\alpha) \land \mathrm{Conv}(b,\beta) \to \forall \varepsilon\,(\mathrm{Conv}(a,\alpha) \land \mathrm{Conv}(b,\beta))$
\item $\varphi_{7} :\equiv \forall \alpha\,(\mathrm{Conv}(b,\alpha) \to \alpha \in \mathbb{R}) \to \mathrm{Conv}(b,\beta) \to \beta \in \mathbb{R}$
\item $\varphi_{8} :\equiv \forall \alpha\,(\mathrm{Conv}(a,\alpha) \to \alpha \in \mathbb{R}) \to \mathrm{Conv}(a,\alpha) \to \alpha \in \mathbb{R}$
\item $\varphi_{9} :\equiv (\mathrm{Conv}(a,\alpha) \to \mathrm{Conv}(a,\alpha)) \to \mathrm{Conv}(a,\alpha) \to \mathrm{Conv}(b,\beta) \to \mathrm{Seq}(a)$
\item $\varphi_{10} :\equiv \mathrm{Conv}(a,\alpha) \to \mathrm{Conv}(b,\beta) \to \varepsilon \in \mathbb{R} \to 0 < \varepsilon \to 0 < \varepsilon$
\item $\varphi_{11} :\equiv \forall w\,(z = w \to \lvert z - (\alpha + \beta) \rvert < \varepsilon \to \lvert w - (\alpha + \beta) \rvert < \varepsilon)$
\item $\varphi_{12} :\equiv \forall a\,\forall \alpha\,(\mathrm{Conv}(a,\alpha) \to \alpha \in \mathbb{R}) \to \forall \alpha\,(\mathrm{Conv}(b,\alpha) \to \alpha \in \mathbb{R})$
\item $\varphi_{13} :\equiv \forall a\,\forall \alpha\,(\mathrm{Conv}(a,\alpha) \to \alpha \in \mathbb{R}) \to \forall \alpha\,(\mathrm{Conv}(a,\alpha) \to \alpha \in \mathbb{R})$
\item $\varphi_{14} :\equiv \forall x\,\forall y\,(x = y \to y = x) \to \forall y\,(a \oplus b(n) = y \to y = a \oplus b(n))$
\item $\varphi_{15} :\equiv (\mathrm{Conv}(a,\alpha) \to \mathrm{Conv}(a,\alpha)) \to \mathrm{Conv}(a,\alpha) \to \mathrm{Conv}(b,\beta) \to \alpha \in \mathbb{R}$
\item $\varphi_{16} :\equiv \mathrm{Conv}(a,\alpha) \to \mathrm{Conv}(b,\beta) \to \varepsilon \in \mathbb{R} \to 0 < \varepsilon \to 0 < \varepsilon \cdot 2^{-1}$
\item $\varphi_{17} :\equiv (\mathrm{Conv}(a,\alpha) \to \mathrm{Conv}(b,\beta) \to \mathrm{Conv}(b,\beta)) \to \mathrm{Conv}(a,\alpha) \to \mathrm{Conv}(b,\beta) \to \mathrm{Seq}(b)$
\item $\varphi_{18} :\equiv \mathrm{Conv}(a,\alpha) \to \mathrm{Conv}(b,\beta) \to \varepsilon \in \mathbb{R} \to 0 < \varepsilon \to \varepsilon \cdot 2^{-1} \in \mathbb{R}$
\item $\varphi_{19} :\equiv N \in \mathbb{N} \land \forall n \in \mathbb{N}\,(N \leq n \to \lvert a \oplus b(n) - (\alpha + \beta) \rvert < \varepsilon)$
\item $\varphi_{20} :\equiv (\mathrm{Conv}(a,\alpha) \to \mathrm{Conv}(b,\beta) \to \mathrm{Conv}(b,\beta)) \to \mathrm{Conv}(a,\alpha) \to \mathrm{Conv}(b,\beta) \to \beta \in \mathbb{R}$
\item $\varphi_{21} :\equiv \mathrm{Conv}(a,\alpha) \to \mathrm{Conv}(b,\beta) \to n \in \mathbb{N} \to a \oplus b(n) = a_{n} + b_{n}$
\item $\varphi_{22} :\equiv \exists N \in \mathbb{N}\,\forall n \in \mathbb{N}\,(N \leq n \to \lvert a \oplus b(n) - (\alpha + \beta) \rvert < \varepsilon)$
\item $\varphi_{23} :\equiv \mathrm{Conv}(a,\alpha) \to \mathrm{Conv}(b,\beta) \to \forall n \in \mathbb{N}\,(a \oplus b(n) = a_{n} + b_{n})$
\item $\varphi_{24} :\equiv \mathrm{Seq}(b) \to \mathrm{Seq}(a \oplus b) \land \forall n \in \mathbb{N}\,(a \oplus b(n) = a_{n} + b_{n})$
\item $\varphi_{25} :\equiv \forall n\,(\mathrm{Seq}(b) \to n \in \mathbb{N} \to b_{n} \in \mathbb{R}) \to \mathrm{Seq}(b) \to n \in \mathbb{N} \to b_{n} \in \mathbb{R}$
\item $\varphi_{26} :\equiv \forall n\,(\mathrm{Seq}(a) \to n \in \mathbb{N} \to a_{n} \in \mathbb{R}) \to \mathrm{Seq}(a) \to n \in \mathbb{N} \to a_{n} \in \mathbb{R}$
\item $\varphi_{27} :\equiv (\mathrm{Conv}(a,\alpha) \to \mathrm{Conv}(b,\beta) \to \beta \in \mathbb{R}) \to \mathrm{Conv}(a,\alpha) \to \mathrm{Conv}(b,\beta) \to \alpha + \beta \in \mathbb{R}$
\item $\varphi_{28} :\equiv \mathrm{Conv}(b,\beta) \to \mathrm{Seq}(a \oplus b) \land \forall n \in \mathbb{N}\,(a \oplus b(n) = a_{n} + b_{n})$
\item $\varphi_{29} :\equiv \lvert a_{n} + b_{n} - (\alpha + \beta) \rvert < z \to \lvert a_{n} + b_{n} - (\alpha + \beta) \rvert < w$
\item $\varphi_{30} :\equiv \forall a\,\forall n\,(\mathrm{Seq}(a) \to n \in \mathbb{N} \to a_{n} \in \mathbb{R}) \to \forall n\,(\mathrm{Seq}(b) \to n \in \mathbb{N} \to b_{n} \in \mathbb{R})$
\item $\varphi_{31} :\equiv \forall a\,\forall n\,(\mathrm{Seq}(a) \to n \in \mathbb{N} \to a_{n} \in \mathbb{R}) \to \forall n\,(\mathrm{Seq}(a) \to n \in \mathbb{N} \to a_{n} \in \mathbb{R})$
\item $\varphi_{32} :\equiv \forall e \in \mathbb{R}\,(0 < e \to 0 < e \cdot 2^{-1}) \to \varepsilon \in \mathbb{R} \to 0 < \varepsilon \to 0 < \varepsilon \cdot 2^{-1}$
\item $\varphi_{33} :\equiv \lvert a_{n} + b_{n} - (\alpha + \beta) \rvert < \varepsilon \to \lvert a \oplus b(n) - (\alpha + \beta) \rvert < \varepsilon$
\item $\varphi_{34} :\equiv n \in \mathbb{N} \to N \leq n \to \lvert a_{n} + b_{n} - (\alpha + \beta) \rvert < \varepsilon \cdot 2^{-1} + \varepsilon \cdot 2^{-1}$
\item $\varphi_{35} :\equiv n \in \mathbb{N} \to N \leq n \to \lvert a_{n} - \alpha \rvert < \varepsilon \cdot 2^{-1} \land \lvert b_{n} - \beta \rvert < \varepsilon \cdot 2^{-1}$
\item $\varphi_{36} :\equiv \forall y\,(\alpha \in \mathbb{R} \to y \in \mathbb{R} \to \alpha + y \in \mathbb{R}) \to \alpha \in \mathbb{R} \to \beta \in \mathbb{R} \to \alpha + \beta \in \mathbb{R}$
\item $\varphi_{37} :\equiv \forall e \in \mathbb{R}\,(e \cdot 2^{-1} + e \cdot 2^{-1} = e) \to \varepsilon \in \mathbb{R} \to \varepsilon \cdot 2^{-1} + \varepsilon \cdot 2^{-1} = \varepsilon$
\item $\varphi_{38} :\equiv \forall x\,\forall y\,(x \in \mathbb{R} \to y \in \mathbb{R} \to x + y \in \mathbb{R}) \to \forall y\,(\alpha \in \mathbb{R} \to y \in \mathbb{R} \to \alpha + y \in \mathbb{R})$
\item $\varphi_{39} :\equiv \forall n \in \mathbb{N}\,(N \leq n \to \lvert a_{n} - \alpha \rvert < \varepsilon \cdot 2^{-1} \land \lvert b_{n} - \beta \rvert < \varepsilon \cdot 2^{-1})$
\item $\varphi_{40} :\equiv \forall x\,\forall y\,(x \in \mathbb{R} \to y \in \mathbb{R} \to x \cdot y \in \mathbb{R}) \to \forall y\,(\varepsilon \in \mathbb{R} \to y \in \mathbb{R} \to \varepsilon \cdot y \in \mathbb{R})$
\item $\varphi_{41} :\equiv \forall w\,(a_{n} + b_{n} = w \to \lvert a_{n} + b_{n} - (\alpha + \beta) \rvert < \varepsilon \to \lvert w - (\alpha + \beta) \rvert < \varepsilon)$
\item $\varphi_{42} :\equiv \forall y\,(\varepsilon \in \mathbb{R} \to y \in \mathbb{R} \to \varepsilon \cdot y \in \mathbb{R}) \to \varepsilon \in \mathbb{R} \to 2^{-1} \in \mathbb{R} \to \varepsilon \cdot 2^{-1} \in \mathbb{R}$
\item $\varphi_{43} :\equiv \forall n \in \mathbb{N}\,(a \oplus b(n) = a_{n} + b_{n}) \to n \in \mathbb{N} \to a \oplus b(n) = a_{n} + b_{n}$
\item $\varphi_{44} :\equiv \lvert b_{n} - \beta \rvert < \varepsilon \cdot 2^{-1} \to \lvert a_{n} + b_{n} - (\alpha + \beta) \rvert < \varepsilon \cdot 2^{-1} + \varepsilon \cdot 2^{-1}$
\item $\varphi_{45} :\equiv \lvert a_{n} + b_{n} - (\alpha + \beta) \rvert < \varepsilon \cdot 2^{-1} + \varepsilon \cdot 2^{-1} \to \lvert a_{n} + b_{n} - (\alpha + \beta) \rvert < \varepsilon$
\item $\varphi_{46} :\equiv \forall y\,(a \oplus b(n) = y \to y = a \oplus b(n)) \to a \oplus b(n) = a_{n} + b_{n} \to a_{n} + b_{n} = a \oplus b(n)$
\item $\varphi_{47} :\equiv \mathrm{Conv}(a,\alpha) \land \mathrm{Conv}(b,\beta) \land \varepsilon \in \mathbb{R} \land 0 < \varepsilon \land \neg \varphi_{22}$
\item $\varphi_{48} :\equiv (\mathrm{Conv}(a,\alpha) \to \mathrm{Conv}(b,\beta) \to \mathrm{Seq}(b)) \to \mathrm{Conv}(a,\alpha) \to \varphi_{28}$
\item $\varphi_{49} :\equiv \forall \varepsilon\,(\mathrm{Conv}(a,\alpha) \land \mathrm{Conv}(b,\beta) \to \varepsilon \in \mathbb{R} \to 0 < \varepsilon \to \varphi_{22})$
\item $\varphi_{50} :\equiv \mathrm{Conv}(a,\alpha) \to \mathrm{Conv}(b,\beta) \to \forall \varepsilon \in \mathbb{R}\,(0 < \varepsilon \to \varphi_{22})$
\item $\varphi_{51} :\equiv \forall a\,\forall \alpha\,(\mathrm{Conv}(a,\alpha) \leftrightarrow \mathrm{Seq}(a) \land \alpha \in \mathbb{R} \land \forall \varepsilon \in \mathbb{R}\,(0 < \varepsilon \to \exists N \in \mathbb{N}\,\forall n \in \mathbb{N}\,(N \leq n \to \lvert a_{n} - \alpha \rvert < \varepsilon)))$
\item $\varphi_{52} :\equiv \forall \varepsilon\,(\mathrm{Conv}(a,\alpha) \land \mathrm{Conv}(b,\beta)) \to \forall \varepsilon \in \mathbb{R}\,(0 < \varepsilon \to \varphi_{22})$
\item $\varphi_{53} :\equiv (\mathrm{Conv}(a,\alpha) \to \varphi_{28}) \to \mathrm{Conv}(a,\alpha) \to \mathrm{Conv}(b,\beta) \to \mathrm{Seq}(a \oplus b)$
\item $\varphi_{54} :\equiv (\mathrm{Conv}(a,\alpha) \to \mathrm{Conv}(b,\beta) \to \varepsilon \in \mathbb{R} \to \varepsilon \in \mathbb{R}) \to 2^{-1} \in \mathbb{R} \to \varphi_{18}$
\item $\varphi_{55} :\equiv \forall \varepsilon\,(\mathrm{Conv}(a,\alpha) \to \mathrm{Conv}(b,\beta) \to \varepsilon \in \mathbb{R} \to 0 < \varepsilon \to \exists N \in \mathbb{N}\,\forall n \in \mathbb{N}\,(N \leq n \to \lvert a_{n} - \alpha \rvert < \varepsilon \land \lvert b_{n} - \beta \rvert < \varepsilon))$
\item $\varphi_{56} :\equiv \forall w\,(\varepsilon \cdot 2^{-1} + \varepsilon \cdot 2^{-1} = w \to \lvert a_{n} + b_{n} - (\alpha + \beta) \rvert < \varepsilon \cdot 2^{-1} + \varepsilon \cdot 2^{-1} \to \lvert a_{n} + b_{n} - (\alpha + \beta) \rvert < w)$
\item $\varphi_{57} :\equiv \forall x\,\forall y\,\forall \alpha\,\forall \beta\,\forall h\,(x \in \mathbb{R} \to y \in \mathbb{R} \to \alpha \in \mathbb{R} \to \beta \in \mathbb{R} \to h \in \mathbb{R} \to \lvert x - \alpha \rvert < h \to \lvert y - \beta \rvert < h \to \lvert x + y - (\alpha + \beta) \rvert < h + h)$
\item $\varphi_{58} :\equiv \forall \alpha\,(\mathrm{Conv}(a \oplus b,\alpha) \leftrightarrow \mathrm{Seq}(a \oplus b) \land \alpha \in \mathbb{R} \land \forall \varepsilon \in \mathbb{R}\,(0 < \varepsilon \to \exists N \in \mathbb{N}\,\forall n \in \mathbb{N}\,(N \leq n \to \lvert a \oplus b(n) - \alpha \rvert < \varepsilon)))$
\item $\varphi_{59} :\equiv \mathrm{Conv}(a,\alpha) \land \mathrm{Conv}(b,\beta) \land \varepsilon \in \mathbb{R} \land 0 < \varepsilon \land (N \in \mathbb{N} \land \varphi_{39})$
\item $\varphi_{60} :\equiv \mathrm{Conv}(a,\alpha) \to \mathrm{Conv}(b,\beta) \to \varepsilon \in \mathbb{R} \to 0 < \varepsilon \to \exists N \in \mathbb{N}\,(\varphi_{39})$
\item $\varphi_{61} :\equiv N \in \mathbb{N} \land \varphi_{39} \to n \in \mathbb{N} \to N \leq n \to \lvert a \oplus b(n) - (\alpha + \beta) \rvert < \varepsilon$
\item $\varphi_{62} :\equiv \mathrm{Conv}(b,\beta) \to \varepsilon \cdot 2^{-1} \in \mathbb{R} \to 0 < \varepsilon \cdot 2^{-1} \to \exists N \in \mathbb{N}\,(\varphi_{39})$
\item $\varphi_{63} :\equiv (\alpha \in \mathbb{R} \to \beta \in \mathbb{R} \to \alpha + \beta \in \mathbb{R}) \to (\mathrm{Conv}(a,\alpha) \to \mathrm{Conv}(b,\beta) \to \alpha \in \mathbb{R}) \to \varphi_{27}$
\item $\varphi_{64} :\equiv \mathrm{Conv}(b,\beta) \to \varepsilon \in \mathbb{R} \to 0 < \varepsilon \to N \in \mathbb{N} \land \varphi_{39} \to n \in \mathbb{N} \to n \in \mathbb{N}$
\item $\varphi_{65} :\equiv \mathrm{Conv}(a,\alpha) \to \mathrm{Conv}(b,\beta) \to \varepsilon \in \mathbb{R} \to 0 < \varepsilon \to N \in \mathbb{N} \land \varphi_{39} \to N \in \mathbb{N}$
\item $\varphi_{66} :\equiv N \in \mathbb{N} \land \varphi_{39} \to n \in \mathbb{N} \to N \leq n \to \lvert a_{n} + b_{n} - (\alpha + \beta) \rvert < \varepsilon$
\item $\varphi_{67} :\equiv N \in \mathbb{N} \land \varphi_{39} \to \forall n \in \mathbb{N}\,(N \leq n \to \lvert a \oplus b(n) - (\alpha + \beta) \rvert < \varepsilon)$
\item $\varphi_{68} :\equiv \mathrm{Conv}(a \oplus b,\alpha + \beta) \leftrightarrow \mathrm{Seq}(a \oplus b) \land \alpha + \beta \in \mathbb{R} \land \forall \varepsilon \in \mathbb{R}\,(0 < \varepsilon \to \varphi_{22})$
\item $\varphi_{69} :\equiv 0 < \varepsilon \to N \in \mathbb{N} \land \varphi_{39} \to n \in \mathbb{N} \to N \leq n \to \varepsilon \cdot 2^{-1} + \varepsilon \cdot 2^{-1} = \varepsilon$
\item $\varphi_{70} :\equiv \forall y\,\forall \alpha\,\forall \beta\,\forall h\,(a_{n} \in \mathbb{R} \to y \in \mathbb{R} \to \alpha \in \mathbb{R} \to \beta \in \mathbb{R} \to h \in \mathbb{R} \to \lvert a_{n} - \alpha \rvert < h \to \lvert y - \beta \rvert < h \to \lvert a_{n} + y - (\alpha + \beta) \rvert < h + h)$
\item $\varphi_{71} :\equiv 0 < \varepsilon \to N \in \mathbb{N} \land \varphi_{39} \to n \in \mathbb{N} \to N \leq n \to \lvert b_{n} - \beta \rvert < \varepsilon \cdot 2^{-1}$
\item $\varphi_{72} :\equiv 0 < \varepsilon \to N \in \mathbb{N} \land \varphi_{39} \to n \in \mathbb{N} \to N \leq n \to \lvert a_{n} - \alpha \rvert < \varepsilon \cdot 2^{-1}$
\item $\varphi_{73} :\equiv \varepsilon \in \mathbb{R} \to 0 < \varepsilon \to N \in \mathbb{N} \land \varphi_{39} \to n \in \mathbb{N} \to N \leq n \to \varepsilon \cdot 2^{-1} \in \mathbb{R}$
\item $\varphi_{74} :\equiv 0 < \varepsilon \to N \in \mathbb{N} \land \varphi_{39} \to n \in \mathbb{N} \to N \leq n \to a \oplus b(n) = a_{n} + b_{n}$
\item $\varphi_{75} :\equiv 0 < \varepsilon \to N \in \mathbb{N} \land \varphi_{39} \to n \in \mathbb{N} \to N \leq n \to a_{n} + b_{n} = a \oplus b(n)$
\item $\varphi_{76} :\equiv \forall h\,(a_{n} \in \mathbb{R} \to b_{n} \in \mathbb{R} \to \alpha \in \mathbb{R} \to \beta \in \mathbb{R} \to h \in \mathbb{R} \to \lvert a_{n} - \alpha \rvert < h \to \lvert b_{n} - \beta \rvert < h \to \lvert a_{n} + b_{n} - (\alpha + \beta) \rvert < h + h)$
\item $\varphi_{77} :\equiv \mathrm{Conv}(b,\beta) \to \varepsilon \in \mathbb{R} \to 0 < \varepsilon \to N \in \mathbb{N} \land \varphi_{39} \to n \in \mathbb{N} \to N \leq n \to b_{n} \in \mathbb{R}$
\item $\varphi_{78} :\equiv \mathrm{Conv}(b,\beta) \to \varepsilon \in \mathbb{R} \to 0 < \varepsilon \to N \in \mathbb{N} \land \varphi_{39} \to n \in \mathbb{N} \to N \leq n \to a_{n} \in \mathbb{R}$
\item $\varphi_{79} :\equiv (\mathrm{Conv}(a,\alpha) \to \mathrm{Conv}(b,\beta) \to \varepsilon \in \mathbb{R} \to \varepsilon \in \mathbb{R}) \to \varphi_{10} \to \varphi_{16}$
\item $\varphi_{80} :\equiv \alpha \in \mathbb{R} \to \beta \in \mathbb{R} \to \varepsilon \cdot 2^{-1} \in \mathbb{R} \to \lvert a_{n} - \alpha \rvert < \varepsilon \cdot 2^{-1} \to \varphi_{44}$
\item $\varphi_{81} :\equiv \forall n\,(\varphi_{59} \to n \in \mathbb{N} \to N \leq n \to \lvert a \oplus b(n) - (\alpha + \beta) \rvert < \varepsilon)$
\item $\varphi_{82} :\equiv \forall n\,(\varphi_{59}) \to \forall n \in \mathbb{N}\,(N \leq n \to \lvert a \oplus b(n) - (\alpha + \beta) \rvert < \varepsilon)$
\item $\varphi_{83} :\equiv (\mathrm{Conv}(a,\alpha) \to \mathrm{Conv}(b,\beta) \to \alpha + \beta \in \mathbb{R}) \to \varphi_{50} \to \varphi_{3}$
\item $\varphi_{84} :\equiv \mathrm{Conv}(a,\alpha) \to \mathrm{Conv}(b,\beta) \to \varepsilon \in \mathbb{R} \to 0 < \varepsilon \to N \in \mathbb{N} \land \varphi_{39} \to \varphi_{19}$
\item $\varphi_{85} :\equiv \mathrm{Conv}(a,\alpha) \to \mathrm{Conv}(b,\beta) \to \varepsilon \in \mathbb{R} \to 0 < \varepsilon \to N \in \mathbb{N} \land \varphi_{39} \to \varphi_{22}$
\item $\varphi_{86} :\equiv \mathrm{Conv}(a,\alpha) \to \mathrm{Conv}(b,\beta) \to \varepsilon \in \mathbb{R} \to 0 < \varepsilon \to N \in \mathbb{N} \land \varphi_{39} \to \varphi_{34}$
\item $\varphi_{87} :\equiv \mathrm{Conv}(a,\alpha) \to \mathrm{Conv}(b,\beta) \to \varepsilon \in \mathbb{R} \to 0 < \varepsilon \to N \in \mathbb{N} \land \varphi_{39} \to \varphi_{35}$
\item $\varphi_{88} :\equiv \mathrm{Conv}(b,\beta) \to \varepsilon \in \mathbb{R} \to 0 < \varepsilon \to N \in \mathbb{N} \land \varphi_{39} \to N \in \mathbb{N} \land \varphi_{39}$
\item $\varphi_{89} :\equiv \mathrm{Conv}(a,\alpha) \to \mathrm{Conv}(b,\beta) \to \varepsilon \in \mathbb{R} \to 0 < \varepsilon \to N \in \mathbb{N} \land \varphi_{39} \to \varphi_{39}$
\item $\varphi_{90} :\equiv (\mathrm{Seq}(a) \to \varphi_{24}) \to (\mathrm{Conv}(a,\alpha) \to \mathrm{Conv}(b,\beta) \to \mathrm{Seq}(a)) \to \varphi_{48}$
\item $\varphi_{91} :\equiv (\mathrm{Conv}(a,\alpha) \to \mathrm{Conv}(b,\beta) \to \varepsilon \in \mathbb{R} \to 0 < \varepsilon \to \varphi_{22}) \to \mathrm{Conv}(a,\alpha) \land \mathrm{Conv}(b,\beta) \to \varepsilon \in \mathbb{R} \to 0 < \varepsilon \to \varphi_{22}$
\item $\varphi_{92} :\equiv (\mathrm{Conv}(a,\alpha) \to \mathrm{Conv}(b,\beta) \to \varepsilon \in \mathbb{R} \to \varepsilon \in \mathbb{R}) \to \mathrm{Conv}(a,\alpha) \to \mathrm{Conv}(b,\beta) \to \varepsilon \in \mathbb{R} \to \varphi_{69}$
\item $\varphi_{93} :\equiv (\mathrm{Conv}(a,\alpha) \to \mathrm{Conv}(b,\beta) \to \varepsilon \in \mathbb{R} \to \varepsilon \in \mathbb{R}) \to 2^{-1} \in \mathbb{R} \to \mathrm{Conv}(a,\alpha) \to \mathrm{Conv}(b,\beta) \to \varphi_{73}$
\item $\varphi_{94} :\equiv (\mathrm{Conv}(a,\alpha) \to \mathrm{Conv}(b,\beta) \to \mathrm{Seq}(a \oplus b)) \to \varphi_{83}$
\item $\varphi_{95} :\equiv (a \oplus b(n) = a_{n} + b_{n} \to a_{n} + b_{n} = a \oplus b(n)) \to \mathrm{Conv}(a,\alpha) \to \mathrm{Conv}(b,\beta) \to \varepsilon \in \mathbb{R} \to \varphi_{75}$
\item $\varphi_{96} :\equiv (\mathrm{Conv}(a,\alpha) \to \mathrm{Conv}(b,\beta) \to \mathrm{Conv}(b,\beta)) \to \varphi_{18} \to \varphi_{16} \to \varphi_{60}$
\item $\varphi_{97} :\equiv (a_{n} + b_{n} = a \oplus b(n) \to \varphi_{33}) \to \mathrm{Conv}(a,\alpha) \to \mathrm{Conv}(b,\beta) \to \varepsilon \in \mathbb{R} \to 0 < \varepsilon \to \varphi_{61}$
\item $\varphi_{98} :\equiv (\varepsilon \cdot 2^{-1} + \varepsilon \cdot 2^{-1} = \varepsilon \to \varphi_{45}) \to \mathrm{Conv}(a,\alpha) \to \mathrm{Conv}(b,\beta) \to \varepsilon \in \mathbb{R} \to 0 < \varepsilon \to \varphi_{66}$
\item $\varphi_{99} :\equiv (\varphi_{47} \to \forall N\,(\varphi_{47})) \to \mathrm{Conv}(a,\alpha) \to \mathrm{Conv}(b,\beta) \to \varepsilon \in \mathbb{R} \to 0 < \varepsilon \to \varphi_{22}$
\item $\varphi_{100} :\equiv (\mathrm{Conv}(a,\alpha) \to \varphi_{64}) \to \mathrm{Conv}(a,\alpha) \to \mathrm{Conv}(b,\beta) \to \varepsilon \in \mathbb{R} \to \varphi_{74}$
\item $\varphi_{101} :\equiv (\mathrm{Conv}(a,\alpha) \to \mathrm{Conv}(b,\beta) \to \mathrm{Seq}(b)) \to (\mathrm{Conv}(a,\alpha) \to \varphi_{64}) \to \mathrm{Conv}(a,\alpha) \to \varphi_{77}$
\item $\varphi_{102} :\equiv (\mathrm{Conv}(a,\alpha) \to \mathrm{Conv}(b,\beta) \to \mathrm{Seq}(a)) \to (\mathrm{Conv}(a,\alpha) \to \varphi_{64}) \to \mathrm{Conv}(a,\alpha) \to \varphi_{78}$
\item $\varphi_{103} :\equiv (\mathrm{Conv}(a,\alpha) \to \mathrm{Conv}(b,\beta) \to \varepsilon \in \mathbb{R} \to 0 < \varepsilon \to \varphi_{61}) \to \varphi_{59} \to n \in \mathbb{N} \to N \leq n \to \lvert a \oplus b(n) - (\alpha + \beta) \rvert < \varepsilon$
\item $\varphi_{104} :\equiv (\mathrm{Conv}(a,\alpha) \to \mathrm{Conv}(b,\beta) \to \varepsilon \in \mathbb{R} \to 0 < \varepsilon \to \varphi_{67}) \to \varphi_{84}$
\item $\varphi_{105} :\equiv \forall N\,(\varphi_{47} \to \neg (N \in \mathbb{N} \land \varphi_{39})) \to \forall N\,(\varphi_{47}) \to \forall N\,(\neg (N \in \mathbb{N} \land \varphi_{39}))$
\item $\varphi_{106} :\equiv (\varphi_{59} \to \forall n\,(\varphi_{59})) \to \mathrm{Conv}(a,\alpha) \to \mathrm{Conv}(b,\beta) \to \varepsilon \in \mathbb{R} \to 0 < \varepsilon \to \varphi_{67}$
\item $\varphi_{107} :\equiv (\mathrm{Conv}(a,\alpha) \to \varphi_{62}) \to (\mathrm{Conv}(a,\alpha) \to \mathrm{Conv}(a,\alpha)) \to \varphi_{96}$
\item $\varphi_{108} :\equiv (\mathrm{Conv}(a,\alpha) \to \mathrm{Conv}(b,\beta) \to \varepsilon \in \mathbb{R} \to 0 < \varepsilon \to \varphi_{66}) \to \varphi_{97}$
\item $\varphi_{109} :\equiv (\mathrm{Conv}(a,\alpha) \to \mathrm{Conv}(b,\beta) \to \varepsilon \in \mathbb{R} \to \varphi_{72}) \to (\mathrm{Conv}(a,\alpha) \to \mathrm{Conv}(b,\beta) \to \varepsilon \in \mathbb{R} \to \varphi_{71}) \to \varphi_{86}$
\item $\varphi_{110} :\equiv (\mathrm{Conv}(a,\alpha) \to \mathrm{Conv}(b,\beta) \to \varepsilon \in \mathbb{R} \to \varphi_{69}) \to \varphi_{86} \to \varphi_{98}$
\item $\varphi_{111} :\equiv (\mathrm{Conv}(a,\alpha) \to \mathrm{Conv}(b,\beta) \to \beta \in \mathbb{R}) \to (\mathrm{Conv}(a,\alpha) \to \mathrm{Conv}(b,\beta) \to \varphi_{73}) \to \varphi_{109}$
\item $\varphi_{112} :\equiv (\mathrm{Conv}(a,\alpha) \to \varphi_{77}) \to (\mathrm{Conv}(a,\alpha) \to \mathrm{Conv}(b,\beta) \to \alpha \in \mathbb{R}) \to \varphi_{111}$
\item $\varphi_{113} :\equiv (a_{n} \in \mathbb{R} \to b_{n} \in \mathbb{R} \to \varphi_{80}) \to (\mathrm{Conv}(a,\alpha) \to \varphi_{78}) \to \varphi_{112}$
\end{abbrevlist}
\end{fullproof}

% 自動生成 (prover/build.js): LimScalar  ノード数 720
\begin{fullproof}{定理2.1(2)}{$\forall c\,\forall a\,\forall \alpha\,(c \in \mathbb{R} \to \mathrm{Conv}(a,\alpha) \to \mathrm{Conv}(c \odot a,c \cdot \alpha))$}{720}\phantomsection\label{tree:LimScalar}
\noindent$\Pi_{1}$:
\begin{ptree}
\AXC{$z = w \to z < \varepsilon \to w < \varepsilon$\lbl{A6}}
\UIC{$\forall w\,(z = w \to z < \varepsilon \to w < \varepsilon)$}
\UIC{$\forall z\,\forall w\,(z = w \to z < \varepsilon \to w < \varepsilon)$}
\AXC{$\forall z\,\forall w\,(z = w \to z < \varepsilon \to w < \varepsilon) \to \varphi_{88}$\lbl{A2}}
\BIC{$\varphi_{88}$}
\AXC{$\varphi_{88} \to \varphi_{186}$\lbl{A2}}
\BIC{$\varphi_{186}$}
\end{ptree}
\noindent$\Pi_{2}$:
\begin{ptree}
\AXC{$z = w \to \lvert a_{n} - \alpha \rvert \cdot \lvert c \rvert < z \to \lvert a_{n} - \alpha \rvert \cdot \lvert c \rvert < w$\lbl{A6}}
\UIC{$\varphi_{73}$}
\UIC{$\forall z\,(\varphi_{73})$}
\AXC{$\forall z\,(\varphi_{73}) \to \varphi_{178}$\lbl{A2}}
\BIC{$\varphi_{178}$}
\AXC{$\varphi_{178} \to \varepsilon \cdot \lvert c \rvert^{-1} \cdot \lvert c \rvert = \varepsilon \to \varphi_{119}$\lbl{A2}}
\BIC{$\varepsilon \cdot \lvert c \rvert^{-1} \cdot \lvert c \rvert = \varepsilon \to \varphi_{119}$}
\end{ptree}
\noindent$\Pi_{3}$:
\begin{ptree}
\AXC{$\varphi_{154}$\lbl{A2}}
\AXC{$c \in \mathbb{R} \to \mathrm{Conv}(a,\alpha) \to \varepsilon \in \mathbb{R} \to 0 < \varepsilon \to c \neq 0 \to \varphi_{202}$\lbl{A1}}
\AXC{$\varphi_{286}$\lbl{A1}}
\BIC{$c \in \mathbb{R} \to \mathrm{Conv}(a,\alpha) \to \varepsilon \in \mathbb{R} \to 0 < \varepsilon \to c \neq 0 \to \varphi_{190}$}
\AXC{$\varphi_{305}$\lbl{A1}}
\BIC{$\varphi_{256}$}
\BIC{$c \in \mathbb{R} \to \mathrm{Conv}(a,\alpha) \to \varepsilon \in \mathbb{R} \to 0 < \varepsilon \to c \neq 0 \to \varphi_{185}$}
\end{ptree}
\noindent$\Pi_{4}$:
\begin{ptree}
\AXC{$\varphi_{15}$\lbl{A1}}
\AXC{$\varphi_{30}$\lbl{A1}}
\AXC{$\forall x \in \mathbb{R}\,(x \neq 0 \to 0 < \lvert x \rvert)$\lbl{[A14]}}
\AXC{$\varphi_{25}$\lbl{A2}}
\BIC{$c \in \mathbb{R} \to c \neq 0 \to 0 < \lvert c \rvert$}
\AXC{$(c \in \mathbb{R} \to c \neq 0 \to 0 < \lvert c \rvert) \to \varphi_{30} \to \varphi_{15} \to \varphi_{34}$\lbl{A1}}
\BIC{$\varphi_{30} \to \varphi_{15} \to \varphi_{34}$}
\BIC{$\varphi_{15} \to \varphi_{34}$}
\BIC{$\varphi_{34}$}
\end{ptree}
\noindent$\Pi_{5}$:
\begin{ptree}
\AXC{$\forall x \in \mathbb{R}\,(\lvert x \rvert \in \mathbb{R})$\lbl{[A5]}}
\AXC{$\forall x \in \mathbb{R}\,(\lvert x \rvert \in \mathbb{R}) \to c \in \mathbb{R} \to \lvert c \rvert \in \mathbb{R}$\lbl{A2}}
\BIC{$c \in \mathbb{R} \to \lvert c \rvert \in \mathbb{R}$}
\AXC{$(c \in \mathbb{R} \to \lvert c \rvert \in \mathbb{R}) \to \varphi_{261}$\lbl{A1}}
\BIC{$\varphi_{261}$}
\end{ptree}
\noindent$\Pi_{6}$:
\begin{ptree}
\AXC{$c \in \mathbb{R} \to \mathrm{Conv}(a,\alpha) \to \varepsilon \in \mathbb{R} \to 0 < \varepsilon \to c \neq 0 \to \varphi_{125}$\lbl{A1}}
\AXC{$\stackrel{\Pi_{5}}{\vdots}$}\noLine
\UIC{$\varphi_{261}$}
\BIC{$c \in \mathbb{R} \to \mathrm{Conv}(a,\alpha) \to \varepsilon \in \mathbb{R} \to 0 < \varepsilon \to c \neq 0 \to \varphi_{132}$}
\end{ptree}
\noindent$\Pi_{7}$:
\begin{ptree}
\AXC{$x = x$\lbl{A5}}
\AXC{$x = y \to x = x \to y = x$\lbl{A6}}
\AXC{$(x = y \to x = x \to y = x) \to x = x \to x = y \to y = x$\lbl{A1}}
\BIC{$x = x \to x = y \to y = x$}
\BIC{$x = y \to y = x$}
\UIC{$\forall y\,(x = y \to y = x)$}
\UIC{$\forall x\,\forall y\,(x = y \to y = x)$}
\end{ptree}
\noindent$\Pi_{8}$:
\begin{ptree}
\AXC{$\stackrel{\Pi_{7}}{\vdots}$}\noLine
\UIC{$\forall x\,\forall y\,(x = y \to y = x)$}
\AXC{$\forall x\,\forall y\,(x = y \to y = x) \to \forall y\,(\lvert c \rvert = y \to y = \lvert c \rvert)$\lbl{A2}}
\BIC{$\forall y\,(\lvert c \rvert = y \to y = \lvert c \rvert)$}
\end{ptree}
\noindent$\Pi_{9}$:
\begin{ptree}
\AXC{$\stackrel{\Pi_{8}}{\vdots}$}\noLine
\UIC{$\forall y\,(\lvert c \rvert = y \to y = \lvert c \rvert)$}
\AXC{$\forall y\,(\lvert c \rvert = y \to y = \lvert c \rvert) \to \lvert c \rvert = 0 \to 0 = \lvert c \rvert$\lbl{A2}}
\BIC{$\lvert c \rvert = 0 \to 0 = \lvert c \rvert$}
\end{ptree}
\noindent$\Pi_{10}$:
\begin{ptree}
\AXC{$\forall x\,\forall y\,(x < y \to x \neq y)$\lbl{[F26]}}
\AXC{$\forall x\,\forall y\,(x < y \to x \neq y) \to \forall y\,(0 < y \to 0 \neq y)$\lbl{A2}}
\BIC{$\forall y\,(0 < y \to 0 \neq y)$}
\end{ptree}
\noindent$\Pi_{11}$:
\begin{ptree}
\AXC{$\stackrel{\Pi_{4}}{\vdots}$}\noLine
\UIC{$\varphi_{34}$}
\AXC{$\stackrel{\Pi_{10}}{\vdots}$}\noLine
\UIC{$\forall y\,(0 < y \to 0 \neq y)$}
\AXC{$\forall y\,(0 < y \to 0 \neq y) \to 0 < \lvert c \rvert \to 0 \neq \lvert c \rvert$\lbl{A2}}
\BIC{$0 < \lvert c \rvert \to 0 \neq \lvert c \rvert$}
\AXC{$(0 < \lvert c \rvert \to 0 \neq \lvert c \rvert) \to \varphi_{34} \to \varphi_{26}$\lbl{A1}}
\BIC{$\varphi_{34} \to \varphi_{26}$}
\BIC{$\varphi_{26}$}
\end{ptree}
\noindent$\Pi_{12}$:
\begin{ptree}
\AXC{$\stackrel{\Pi_{9}}{\vdots}$}\noLine
\UIC{$\lvert c \rvert = 0 \to 0 = \lvert c \rvert$}
\AXC{$\stackrel{\Pi_{11}}{\vdots}$}\noLine
\UIC{$\varphi_{26}$}
\AXC{$\varphi_{26} \to (\lvert c \rvert = 0 \to 0 = \lvert c \rvert) \to \varphi_{27}$\lbl{A1}}
\BIC{$(\lvert c \rvert = 0 \to 0 = \lvert c \rvert) \to \varphi_{27}$}
\BIC{$\varphi_{27}$}
\end{ptree}
\noindent$\Pi_{13}$:
\begin{ptree}
\AXC{$\forall x \in \mathbb{R}\,(x \neq 0 \to x^{-1} \in \mathbb{R})$\lbl{(R1)}}
\AXC{$\varphi_{104}$\lbl{A2}}
\BIC{$\lvert c \rvert \in \mathbb{R} \to \lvert c \rvert \neq 0 \to \lvert c \rvert^{-1} \in \mathbb{R}$}
\AXC{$(\lvert c \rvert \in \mathbb{R} \to \lvert c \rvert \neq 0 \to \lvert c \rvert^{-1} \in \mathbb{R}) \to \varphi_{291}$\lbl{A1}}
\BIC{$\varphi_{291}$}
\end{ptree}
\noindent$\Pi_{14}$:
\begin{ptree}
\AXC{$\stackrel{\Pi_{12}}{\vdots}$}\noLine
\UIC{$\varphi_{27}$}
\AXC{$\stackrel{\Pi_{6}}{\vdots}$}\noLine
\UIC{$c \in \mathbb{R} \to \mathrm{Conv}(a,\alpha) \to \varepsilon \in \mathbb{R} \to 0 < \varepsilon \to c \neq 0 \to \varphi_{132}$}
\AXC{$\stackrel{\Pi_{13}}{\vdots}$}\noLine
\UIC{$\varphi_{291}$}
\BIC{$\varphi_{237}$}
\BIC{$c \in \mathbb{R} \to \mathrm{Conv}(a,\alpha) \to \varepsilon \in \mathbb{R} \to 0 < \varepsilon \to c \neq 0 \to \varphi_{136}$}
\end{ptree}
\noindent$\Pi_{15}$:
\begin{ptree}
\AXC{$\forall x\,\forall y\,(x \in \mathbb{R} \to y \in \mathbb{R} \to x \cdot y \in \mathbb{R})$\lbl{(R1)}}
\AXC{$\varphi_{116}$\lbl{A2}}
\BIC{$\forall y\,(\varepsilon \in \mathbb{R} \to y \in \mathbb{R} \to \varepsilon \cdot y \in \mathbb{R})$}
\AXC{$\varphi_{137}$\lbl{A2}}
\BIC{$\varepsilon \in \mathbb{R} \to \lvert c \rvert^{-1} \in \mathbb{R} \to \varepsilon \cdot \lvert c \rvert^{-1} \in \mathbb{R}$}
\AXC{$\varphi_{302}$\lbl{A1}}
\BIC{$(c \in \mathbb{R} \to \mathrm{Conv}(a,\alpha) \to \varepsilon \in \mathbb{R} \to \varepsilon \in \mathbb{R}) \to \varphi_{265}$}
\end{ptree}
\noindent$\Pi_{16}$:
\begin{ptree}
\AXC{$c \in \mathbb{R} \to \mathrm{Conv}(a,\alpha) \to \varepsilon \in \mathbb{R} \to \varepsilon \in \mathbb{R}$\lbl{A1}}
\AXC{$\stackrel{\Pi_{15}}{\vdots}$}\noLine
\UIC{$(c \in \mathbb{R} \to \mathrm{Conv}(a,\alpha) \to \varepsilon \in \mathbb{R} \to \varepsilon \in \mathbb{R}) \to \varphi_{265}$}
\BIC{$\varphi_{265}$}
\end{ptree}
\noindent$\Pi_{17}$:
\begin{ptree}
\AXC{$\stackrel{\Pi_{14}}{\vdots}$}\noLine
\UIC{$c \in \mathbb{R} \to \mathrm{Conv}(a,\alpha) \to \varepsilon \in \mathbb{R} \to 0 < \varepsilon \to c \neq 0 \to \varphi_{136}$}
\AXC{$\stackrel{\Pi_{16}}{\vdots}$}\noLine
\UIC{$\varphi_{265}$}
\BIC{$c \in \mathbb{R} \to \mathrm{Conv}(a,\alpha) \to \varepsilon \in \mathbb{R} \to 0 < \varepsilon \to c \neq 0 \to \varphi_{152}$}
\end{ptree}
\noindent$\Pi_{18}$:
\begin{ptree}
\AXC{$\forall a\,\forall \alpha\,(\mathrm{Conv}(a,\alpha) \to \alpha \in \mathbb{R})$\lbl{[C2]}}
\AXC{$\varphi_{13}$\lbl{A2}}
\BIC{$\forall \alpha\,(\mathrm{Conv}(a,\alpha) \to \alpha \in \mathbb{R})$}
\AXC{$\varphi_{5}$\lbl{A2}}
\BIC{$\mathrm{Conv}(a,\alpha) \to \alpha \in \mathbb{R}$}
\AXC{$(\mathrm{Conv}(a,\alpha) \to \alpha \in \mathbb{R}) \to \varphi_{61}$\lbl{A1}}
\BIC{$\varphi_{61}$}
\end{ptree}
\noindent$\Pi_{19}$:
\begin{ptree}
\AXC{$c \in \mathbb{R} \to \mathrm{Conv}(a,\alpha) \to \mathrm{Conv}(a,\alpha)$\lbl{A1}}
\AXC{$\stackrel{\Pi_{18}}{\vdots}$}\noLine
\UIC{$\varphi_{61}$}
\BIC{$c \in \mathbb{R} \to \mathrm{Conv}(a,\alpha) \to \alpha \in \mathbb{R}$}
\end{ptree}
\noindent$\Pi_{20}$:
\begin{ptree}
\AXC{$\forall a\,\forall \alpha\,(\mathrm{Conv}(a,\alpha) \to \mathrm{Seq}(a))$\lbl{[C1]}}
\AXC{$\varphi_{2}$\lbl{A2}}
\BIC{$\forall \alpha\,(\mathrm{Conv}(a,\alpha) \to \mathrm{Seq}(a))$}
\AXC{$\varphi_{1}$\lbl{A2}}
\BIC{$\mathrm{Conv}(a,\alpha) \to \mathrm{Seq}(a)$}
\AXC{$(\mathrm{Conv}(a,\alpha) \to \mathrm{Seq}(a)) \to \varphi_{36}$\lbl{A1}}
\BIC{$\varphi_{36}$}
\end{ptree}
\noindent$\Pi_{21}$:
\begin{ptree}
\AXC{$c \in \mathbb{R} \to \mathrm{Conv}(a,\alpha) \to \mathrm{Conv}(a,\alpha)$\lbl{A1}}
\AXC{$\stackrel{\Pi_{20}}{\vdots}$}\noLine
\UIC{$\varphi_{36}$}
\BIC{$c \in \mathbb{R} \to \mathrm{Conv}(a,\alpha) \to \mathrm{Seq}(a)$}
\end{ptree}
\noindent$\Pi_{22}$:
\begin{ptree}
\AXC{$\forall a\,\forall n\,(\mathrm{Seq}(a) \to n \in \mathbb{N} \to a_{n} \in \mathbb{R})$\lbl{[S1]}}
\AXC{$\varphi_{99}$\lbl{A2}}
\BIC{$\forall n\,(\mathrm{Seq}(a) \to n \in \mathbb{N} \to a_{n} \in \mathbb{R})$}
\AXC{$\varphi_{80}$\lbl{A2}}
\BIC{$\mathrm{Seq}(a) \to n \in \mathbb{N} \to a_{n} \in \mathbb{R}$}
\AXC{$\varphi_{284}$\lbl{A1}}
\BIC{$(c \in \mathbb{R} \to \mathrm{Conv}(a,\alpha) \to \mathrm{Seq}(a)) \to \varphi_{259}$}
\end{ptree}
\noindent$\Pi_{23}$:
\begin{ptree}
\AXC{$\stackrel{\Pi_{21}}{\vdots}$}\noLine
\UIC{$c \in \mathbb{R} \to \mathrm{Conv}(a,\alpha) \to \mathrm{Seq}(a)$}
\AXC{$\stackrel{\Pi_{22}}{\vdots}$}\noLine
\UIC{$(c \in \mathbb{R} \to \mathrm{Conv}(a,\alpha) \to \mathrm{Seq}(a)) \to \varphi_{259}$}
\BIC{$\varphi_{259}$}
\end{ptree}
\noindent$\Pi_{24}$:
\begin{ptree}
\AXC{$c \in \mathbb{R} \to \mathrm{Conv}(a,\alpha) \to \varepsilon \in \mathbb{R} \to 0 < \varepsilon \to c \neq 0 \to \varphi_{117}$\lbl{A1}}
\AXC{$\stackrel{\Pi_{23}}{\vdots}$}\noLine
\UIC{$\varphi_{259}$}
\BIC{$c \in \mathbb{R} \to \mathrm{Conv}(a,\alpha) \to \varepsilon \in \mathbb{R} \to 0 < \varepsilon \to c \neq 0 \to \varphi_{134}$}
\end{ptree}
\noindent$\Pi_{25}$:
\begin{ptree}
\AXC{$\forall x\,\forall y\,(x \in \mathbb{R} \to y \in \mathbb{R} \to x - y \in \mathbb{R})$\lbl{[F1]}}
\AXC{$\varphi_{130}$\lbl{A2}}
\BIC{$\forall y\,(a_{n} \in \mathbb{R} \to y \in \mathbb{R} \to a_{n} - y \in \mathbb{R})$}
\AXC{$\varphi_{143}$\lbl{A2}}
\BIC{$a_{n} \in \mathbb{R} \to \alpha \in \mathbb{R} \to a_{n} - \alpha \in \mathbb{R}$}
\AXC{$(a_{n} \in \mathbb{R} \to \alpha \in \mathbb{R} \to a_{n} - \alpha \in \mathbb{R}) \to \varphi_{279}$\lbl{A1}}
\BIC{$\varphi_{279}$}
\end{ptree}
\noindent$\Pi_{26}$:
\begin{ptree}
\AXC{$\stackrel{\Pi_{19}}{\vdots}$}\noLine
\UIC{$c \in \mathbb{R} \to \mathrm{Conv}(a,\alpha) \to \alpha \in \mathbb{R}$}
\AXC{$\stackrel{\Pi_{24}}{\vdots}$}\noLine
\UIC{$c \in \mathbb{R} \to \mathrm{Conv}(a,\alpha) \to \varepsilon \in \mathbb{R} \to 0 < \varepsilon \to c \neq 0 \to \varphi_{134}$}
\AXC{$\stackrel{\Pi_{25}}{\vdots}$}\noLine
\UIC{$\varphi_{279}$}
\BIC{$\varphi_{230}$}
\BIC{$c \in \mathbb{R} \to \mathrm{Conv}(a,\alpha) \to \varepsilon \in \mathbb{R} \to 0 < \varepsilon \to c \neq 0 \to \varphi_{145}$}
\end{ptree}
\noindent$\Pi_{27}$:
\begin{ptree}
\AXC{$\forall x \in \mathbb{R}\,(\lvert x \rvert \in \mathbb{R})$\lbl{[A5]}}
\AXC{$\varphi_{100}$\lbl{A2}}
\BIC{$a_{n} - \alpha \in \mathbb{R} \to \lvert a_{n} - \alpha \rvert \in \mathbb{R}$}
\AXC{$(a_{n} - \alpha \in \mathbb{R} \to \lvert a_{n} - \alpha \rvert \in \mathbb{R}) \to \varphi_{267}$\lbl{A1}}
\BIC{$\varphi_{267}$}
\end{ptree}
\noindent$\Pi_{28}$:
\begin{ptree}
\AXC{$\stackrel{\Pi_{26}}{\vdots}$}\noLine
\UIC{$c \in \mathbb{R} \to \mathrm{Conv}(a,\alpha) \to \varepsilon \in \mathbb{R} \to 0 < \varepsilon \to c \neq 0 \to \varphi_{145}$}
\AXC{$\stackrel{\Pi_{27}}{\vdots}$}\noLine
\UIC{$\varphi_{267}$}
\BIC{$c \in \mathbb{R} \to \mathrm{Conv}(a,\alpha) \to \varepsilon \in \mathbb{R} \to 0 < \varepsilon \to c \neq 0 \to \varphi_{153}$}
\end{ptree}
\noindent$\Pi_{29}$:
\begin{ptree}
\AXC{$\varphi_{81}$\lbl{[F32]}}
\AXC{$\varphi_{81} \to \varphi_{175}$\lbl{A2}}
\BIC{$\varphi_{175}$}
\AXC{$\varphi_{175} \to \varphi_{219}$\lbl{A2}}
\BIC{$\varphi_{219}$}
\AXC{$\varphi_{219} \to \varphi_{223}$\lbl{A2}}
\BIC{$\varphi_{223}$}
\AXC{$\varphi_{223} \to \varphi_{335}$\lbl{A1}}
\BIC{$\varphi_{335}$}
\end{ptree}
\noindent$\Pi_{30}$:
\begin{ptree}
\AXC{$\stackrel{\Pi_{17}}{\vdots}$}\noLine
\UIC{$c \in \mathbb{R} \to \mathrm{Conv}(a,\alpha) \to \varepsilon \in \mathbb{R} \to 0 < \varepsilon \to c \neq 0 \to \varphi_{152}$}
\AXC{$\stackrel{\Pi_{28}}{\vdots}$}\noLine
\UIC{$c \in \mathbb{R} \to \mathrm{Conv}(a,\alpha) \to \varepsilon \in \mathbb{R} \to 0 < \varepsilon \to c \neq 0 \to \varphi_{153}$}
\AXC{$\stackrel{\Pi_{29}}{\vdots}$}\noLine
\UIC{$\varphi_{335}$}
\BIC{$\varphi_{334}$}
\BIC{$\varphi_{328}$}
\end{ptree}
\noindent$\Pi_{31}$:
\begin{ptree}
\AXC{$\stackrel{\Pi_{4}}{\vdots}$}\noLine
\UIC{$\varphi_{34}$}
\AXC{$\stackrel{\Pi_{6}}{\vdots}$}\noLine
\UIC{$c \in \mathbb{R} \to \mathrm{Conv}(a,\alpha) \to \varepsilon \in \mathbb{R} \to 0 < \varepsilon \to c \neq 0 \to \varphi_{132}$}
\AXC{$\stackrel{\Pi_{30}}{\vdots}$}\noLine
\UIC{$\varphi_{328}$}
\BIC{$\varphi_{34} \to \varphi_{288}$}
\BIC{$\varphi_{288}$}
\end{ptree}
\noindent$\Pi_{32}$:
\begin{ptree}
\AXC{$y = z \to x = y \to x = z$\lbl{A6}}
\AXC{$(y = z \to x = y \to x = z) \to x = y \to y = z \to x = z$\lbl{A1}}
\BIC{$x = y \to y = z \to x = z$}
\UIC{$\forall z\,(x = y \to y = z \to x = z)$}
\UIC{$\forall y\,\forall z\,(x = y \to y = z \to x = z)$}
\UIC{$\forall x\,\forall y\,\forall z\,(x = y \to y = z \to x = z)$}
\AXC{$\forall x\,\forall y\,\forall z\,(x = y \to y = z \to x = z) \to \varphi_{51}$\lbl{A2}}
\BIC{$\varphi_{51}$}
\end{ptree}
\noindent$\Pi_{33}$:
\begin{ptree}
\AXC{$\forall x \in \mathbb{R}\,(x \cdot 1 = x)$\lbl{(R8)}}
\AXC{$\forall x \in \mathbb{R}\,(x \cdot 1 = x) \to \varepsilon \in \mathbb{R} \to \varepsilon \cdot 1 = \varepsilon$\lbl{A2}}
\BIC{$\varepsilon \in \mathbb{R} \to \varepsilon \cdot 1 = \varepsilon$}
\AXC{$(\varepsilon \in \mathbb{R} \to \varepsilon \cdot 1 = \varepsilon) \to \varphi_{231}$\lbl{A1}}
\BIC{$\varphi_{231}$}
\end{ptree}
\noindent$\Pi_{34}$:
\begin{ptree}
\AXC{$y = z \to x = y \to x = z$\lbl{A6}}
\AXC{$(y = z \to x = y \to x = z) \to x = y \to y = z \to x = z$\lbl{A1}}
\BIC{$x = y \to y = z \to x = z$}
\UIC{$\forall z\,(x = y \to y = z \to x = z)$}
\UIC{$\forall y\,\forall z\,(x = y \to y = z \to x = z)$}
\UIC{$\forall x\,\forall y\,\forall z\,(x = y \to y = z \to x = z)$}
\AXC{$\forall x\,\forall y\,\forall z\,(x = y \to y = z \to x = z) \to \varphi_{51}$\lbl{A2}}
\BIC{$\varphi_{51}$}
\end{ptree}
\noindent$\Pi_{35}$:
\begin{ptree}
\AXC{$\forall x \in \mathbb{R}\,(x \neq 0 \to x \cdot x^{-1} = 1)$\lbl{(R9)}}
\AXC{$\varphi_{108}$\lbl{A2}}
\BIC{$\lvert c \rvert \in \mathbb{R} \to \lvert c \rvert \neq 0 \to \lvert c \rvert \cdot \lvert c \rvert^{-1} = 1$}
\AXC{$(\lvert c \rvert \in \mathbb{R} \to \lvert c \rvert \neq 0 \to \lvert c \rvert \cdot \lvert c \rvert^{-1} = 1) \to \varphi_{292}$\lbl{A1}}
\BIC{$\varphi_{292}$}
\end{ptree}
\noindent$\Pi_{36}$:
\begin{ptree}
\AXC{$\stackrel{\Pi_{12}}{\vdots}$}\noLine
\UIC{$\varphi_{27}$}
\AXC{$\stackrel{\Pi_{6}}{\vdots}$}\noLine
\UIC{$c \in \mathbb{R} \to \mathrm{Conv}(a,\alpha) \to \varepsilon \in \mathbb{R} \to 0 < \varepsilon \to c \neq 0 \to \varphi_{132}$}
\AXC{$\stackrel{\Pi_{35}}{\vdots}$}\noLine
\UIC{$\varphi_{292}$}
\BIC{$\varphi_{238}$}
\BIC{$c \in \mathbb{R} \to \mathrm{Conv}(a,\alpha) \to \varepsilon \in \mathbb{R} \to 0 < \varepsilon \to c \neq 0 \to \varphi_{146}$}
\AXC{$\varphi_{313}$\lbl{A1}}
\BIC{$\varepsilon \cdot (\lvert c \rvert \cdot \lvert c \rvert^{-1}) = \varepsilon \cdot (\lvert c \rvert \cdot \lvert c \rvert^{-1}) \to \varphi_{255}$}
\end{ptree}
\noindent$\Pi_{37}$:
\begin{ptree}
\AXC{$x = x$\lbl{A5}}
\UIC{$\forall x\,(x = x)$}
\AXC{$\varphi_{19}$\lbl{A2}}
\BIC{$\varepsilon \cdot (\lvert c \rvert \cdot \lvert c \rvert^{-1}) = \varepsilon \cdot (\lvert c \rvert \cdot \lvert c \rvert^{-1})$}
\AXC{$\stackrel{\Pi_{36}}{\vdots}$}\noLine
\UIC{$\varepsilon \cdot (\lvert c \rvert \cdot \lvert c \rvert^{-1}) = \varepsilon \cdot (\lvert c \rvert \cdot \lvert c \rvert^{-1}) \to \varphi_{255}$}
\BIC{$\varphi_{255}$}
\end{ptree}
\noindent$\Pi_{38}$:
\begin{ptree}
\AXC{$y = z \to x = y \to x = z$\lbl{A6}}
\AXC{$(y = z \to x = y \to x = z) \to x = y \to y = z \to x = z$\lbl{A1}}
\BIC{$x = y \to y = z \to x = z$}
\UIC{$\forall z\,(x = y \to y = z \to x = z)$}
\UIC{$\forall y\,\forall z\,(x = y \to y = z \to x = z)$}
\UIC{$\forall x\,\forall y\,\forall z\,(x = y \to y = z \to x = z)$}
\AXC{$\forall x\,\forall y\,\forall z\,(x = y \to y = z \to x = z) \to \varphi_{51}$\lbl{A2}}
\BIC{$\varphi_{51}$}
\end{ptree}
\noindent$\Pi_{39}$:
\begin{ptree}
\AXC{$\forall x\,\forall y\,(x \in \mathbb{R} \to y \in \mathbb{R} \to x \cdot y = y \cdot x)$\lbl{(R7)}}
\AXC{$\varphi_{148}$\lbl{A2}}
\BIC{$\forall y\,(\lvert c \rvert^{-1} \in \mathbb{R} \to y \in \mathbb{R} \to \lvert c \rvert^{-1} \cdot y = y \cdot \lvert c \rvert^{-1})$}
\AXC{$\varphi_{210}$\lbl{A2}}
\BIC{$\lvert c \rvert^{-1} \in \mathbb{R} \to \lvert c \rvert \in \mathbb{R} \to \lvert c \rvert^{-1} \cdot \lvert c \rvert = \lvert c \rvert \cdot \lvert c \rvert^{-1}$}
\AXC{$\varphi_{323}$\lbl{A1}}
\BIC{$\varphi_{309}$}
\end{ptree}
\noindent$\Pi_{40}$:
\begin{ptree}
\AXC{$\stackrel{\Pi_{6}}{\vdots}$}\noLine
\UIC{$c \in \mathbb{R} \to \mathrm{Conv}(a,\alpha) \to \varepsilon \in \mathbb{R} \to 0 < \varepsilon \to c \neq 0 \to \varphi_{132}$}
\AXC{$\stackrel{\Pi_{14}}{\vdots}$}\noLine
\UIC{$c \in \mathbb{R} \to \mathrm{Conv}(a,\alpha) \to \varepsilon \in \mathbb{R} \to 0 < \varepsilon \to c \neq 0 \to \varphi_{136}$}
\AXC{$\stackrel{\Pi_{39}}{\vdots}$}\noLine
\UIC{$\varphi_{309}$}
\BIC{$\varphi_{270}$}
\BIC{$c \in \mathbb{R} \to \mathrm{Conv}(a,\alpha) \to \varepsilon \in \mathbb{R} \to 0 < \varepsilon \to c \neq 0 \to \varphi_{189}$}
\end{ptree}
\noindent$\Pi_{41}$:
\begin{ptree}
\AXC{$x = x$\lbl{A5}}
\UIC{$\forall x\,(x = x)$}
\AXC{$\varphi_{20}$\lbl{A2}}
\BIC{$\varepsilon \cdot (\lvert c \rvert^{-1} \cdot \lvert c \rvert) = \varepsilon \cdot (\lvert c \rvert^{-1} \cdot \lvert c \rvert)$}
\AXC{$\stackrel{\Pi_{40}}{\vdots}$}\noLine
\UIC{$c \in \mathbb{R} \to \mathrm{Conv}(a,\alpha) \to \varepsilon \in \mathbb{R} \to 0 < \varepsilon \to c \neq 0 \to \varphi_{189}$}
\AXC{$\varphi_{324}$\lbl{A1}}
\BIC{$\varepsilon \cdot (\lvert c \rvert^{-1} \cdot \lvert c \rvert) = \varepsilon \cdot (\lvert c \rvert^{-1} \cdot \lvert c \rvert) \to \varphi_{273}$}
\BIC{$\varphi_{273}$}
\end{ptree}
\noindent$\Pi_{42}$:
\begin{ptree}
\AXC{$z = w \to \varphi_{70}$\lbl{A6}}
\UIC{$\forall w\,(z = w \to \varphi_{70})$}
\UIC{$\forall z\,\forall w\,(z = w \to \varphi_{70})$}
\AXC{$\forall z\,\forall w\,(z = w \to \varphi_{70}) \to \varphi_{161}$\lbl{A2}}
\BIC{$\varphi_{161}$}
\AXC{$\varphi_{161} \to \lvert c \rvert^{-1} \cdot \lvert c \rvert = \lvert c \rvert \cdot \lvert c \rvert^{-1} \to \varphi_{142}$\lbl{A2}}
\BIC{$\lvert c \rvert^{-1} \cdot \lvert c \rvert = \lvert c \rvert \cdot \lvert c \rvert^{-1} \to \varphi_{142}$}
\AXC{$\stackrel{\Pi_{41}}{\vdots}$}\noLine
\UIC{$\varphi_{273}$}
\BIC{$c \in \mathbb{R} \to \mathrm{Conv}(a,\alpha) \to \varepsilon \in \mathbb{R} \to 0 < \varepsilon \to \varphi_{211}$}
\end{ptree}
\noindent$\Pi_{43}$:
\begin{ptree}
\AXC{$\varphi_{37}$\lbl{(R6)}}
\AXC{$\varphi_{37} \to \varphi_{47}$\lbl{A2}}
\BIC{$\varphi_{47}$}
\AXC{$\varphi_{47} \to \varphi_{120}$\lbl{A2}}
\BIC{$\varphi_{120}$}
\AXC{$\varphi_{120} \to \varepsilon \in \mathbb{R} \to \varphi_{118}$\lbl{A2}}
\BIC{$\varepsilon \in \mathbb{R} \to \varphi_{118}$}
\AXC{$\varphi_{329}$\lbl{A1}}
\BIC{$(c \in \mathbb{R} \to \mathrm{Conv}(a,\alpha) \to \varepsilon \in \mathbb{R} \to \varepsilon \in \mathbb{R}) \to \varphi_{314}$}
\end{ptree}
\noindent$\Pi_{44}$:
\begin{ptree}
\AXC{$c \in \mathbb{R} \to \mathrm{Conv}(a,\alpha) \to \varepsilon \in \mathbb{R} \to \varepsilon \in \mathbb{R}$\lbl{A1}}
\AXC{$\stackrel{\Pi_{43}}{\vdots}$}\noLine
\UIC{$(c \in \mathbb{R} \to \mathrm{Conv}(a,\alpha) \to \varepsilon \in \mathbb{R} \to \varepsilon \in \mathbb{R}) \to \varphi_{314}$}
\BIC{$\varphi_{314}$}
\end{ptree}
\noindent$\Pi_{45}$:
\begin{ptree}
\AXC{$\stackrel{\Pi_{6}}{\vdots}$}\noLine
\UIC{$c \in \mathbb{R} \to \mathrm{Conv}(a,\alpha) \to \varepsilon \in \mathbb{R} \to 0 < \varepsilon \to c \neq 0 \to \varphi_{132}$}
\AXC{$\stackrel{\Pi_{14}}{\vdots}$}\noLine
\UIC{$c \in \mathbb{R} \to \mathrm{Conv}(a,\alpha) \to \varepsilon \in \mathbb{R} \to 0 < \varepsilon \to c \neq 0 \to \varphi_{136}$}
\AXC{$\stackrel{\Pi_{44}}{\vdots}$}\noLine
\UIC{$\varphi_{314}$}
\BIC{$\varphi_{275}$}
\BIC{$c \in \mathbb{R} \to \mathrm{Conv}(a,\alpha) \to \varepsilon \in \mathbb{R} \to 0 < \varepsilon \to c \neq 0 \to \varphi_{203}$}
\end{ptree}
\noindent$\Pi_{46}$:
\begin{ptree}
\AXC{$\stackrel{\Pi_{42}}{\vdots}$}\noLine
\UIC{$c \in \mathbb{R} \to \mathrm{Conv}(a,\alpha) \to \varepsilon \in \mathbb{R} \to 0 < \varepsilon \to \varphi_{211}$}
\AXC{$\stackrel{\Pi_{45}}{\vdots}$}\noLine
\UIC{$c \in \mathbb{R} \to \mathrm{Conv}(a,\alpha) \to \varepsilon \in \mathbb{R} \to 0 < \varepsilon \to c \neq 0 \to \varphi_{203}$}
\AXC{$\varphi_{333}$\lbl{A1}}
\BIC{$(c \in \mathbb{R} \to \mathrm{Conv}(a,\alpha) \to \varepsilon \in \mathbb{R} \to 0 < \varepsilon \to \varphi_{211}) \to \varphi_{280}$}
\BIC{$\varphi_{280}$}
\end{ptree}
\noindent$\Pi_{47}$:
\begin{ptree}
\AXC{$\stackrel{\Pi_{38}}{\vdots}$}\noLine
\UIC{$\varphi_{51}$}
\AXC{$\varphi_{51} \to \varphi_{160}$\lbl{A2}}
\BIC{$\varphi_{160}$}
\AXC{$\varphi_{160} \to \varepsilon \cdot \lvert c \rvert^{-1} \cdot \lvert c \rvert = \varepsilon \cdot (\lvert c \rvert^{-1} \cdot \lvert c \rvert) \to \varphi_{141}$\lbl{A2}}
\BIC{$\varepsilon \cdot \lvert c \rvert^{-1} \cdot \lvert c \rvert = \varepsilon \cdot (\lvert c \rvert^{-1} \cdot \lvert c \rvert) \to \varphi_{141}$}
\AXC{$\stackrel{\Pi_{46}}{\vdots}$}\noLine
\UIC{$\varphi_{280}$}
\BIC{$c \in \mathbb{R} \to \mathrm{Conv}(a,\alpha) \to \varepsilon \in \mathbb{R} \to 0 < \varepsilon \to c \neq 0 \to \varphi_{204}$}
\AXC{$\varphi_{331}$\lbl{A1}}
\BIC{$\varphi_{307}$}
\end{ptree}
\noindent$\Pi_{48}$:
\begin{ptree}
\AXC{$z = w \to \varphi_{69}$\lbl{A6}}
\UIC{$\forall w\,(z = w \to \varphi_{69})$}
\UIC{$\forall z\,\forall w\,(z = w \to \varphi_{69})$}
\AXC{$\forall z\,\forall w\,(z = w \to \varphi_{69}) \to \varphi_{159}$\lbl{A2}}
\BIC{$\varphi_{159}$}
\AXC{$\varphi_{159} \to \lvert c \rvert \cdot \lvert c \rvert^{-1} = 1 \to \varphi_{115}$\lbl{A2}}
\BIC{$\lvert c \rvert \cdot \lvert c \rvert^{-1} = 1 \to \varphi_{115}$}
\AXC{$\stackrel{\Pi_{37}}{\vdots}$}\noLine
\UIC{$\varphi_{255}$}
\BIC{$c \in \mathbb{R} \to \mathrm{Conv}(a,\alpha) \to \varepsilon \in \mathbb{R} \to 0 < \varepsilon \to c \neq 0 \to \varphi_{173}$}
\AXC{$\stackrel{\Pi_{47}}{\vdots}$}\noLine
\UIC{$\varphi_{307}$}
\BIC{$\varphi_{260}$}
\end{ptree}
\noindent$\Pi_{49}$:
\begin{ptree}
\AXC{$\stackrel{\Pi_{34}}{\vdots}$}\noLine
\UIC{$\varphi_{51}$}
\AXC{$\varphi_{51} \to \varphi_{158}$\lbl{A2}}
\BIC{$\varphi_{158}$}
\AXC{$\varphi_{158} \to \varphi_{187}$\lbl{A2}}
\BIC{$\varphi_{187}$}
\AXC{$\stackrel{\Pi_{48}}{\vdots}$}\noLine
\UIC{$\varphi_{260}$}
\BIC{$c \in \mathbb{R} \to \mathrm{Conv}(a,\alpha) \to \varepsilon \in \mathbb{R} \to 0 < \varepsilon \to c \neq 0 \to \varphi_{174}$}
\AXC{$\varphi_{327}$\lbl{A1}}
\BIC{$\varphi_{295}$}
\end{ptree}
\noindent$\Pi_{50}$:
\begin{ptree}
\AXC{$\stackrel{\Pi_{32}}{\vdots}$}\noLine
\UIC{$\varphi_{51}$}
\AXC{$\varphi_{51} \to \varphi_{93}$\lbl{A2}}
\BIC{$\varphi_{93}$}
\AXC{$\varphi_{93} \to \varphi_{96}$\lbl{A2}}
\BIC{$\varphi_{96}$}
\AXC{$c \in \mathbb{R} \to \mathrm{Conv}(a,\alpha) \to \varepsilon \in \mathbb{R} \to \varepsilon \in \mathbb{R}$\lbl{A1}}
\AXC{$\stackrel{\Pi_{33}}{\vdots}$}\noLine
\UIC{$\varphi_{231}$}
\BIC{$c \in \mathbb{R} \to \mathrm{Conv}(a,\alpha) \to \varepsilon \in \mathbb{R} \to 0 < \varepsilon \to c \neq 0 \to \varphi_{129}$}
\AXC{$\stackrel{\Pi_{49}}{\vdots}$}\noLine
\UIC{$\varphi_{295}$}
\BIC{$\varphi_{247}$}
\BIC{$c \in \mathbb{R} \to \mathrm{Conv}(a,\alpha) \to \varepsilon \in \mathbb{R} \to 0 < \varepsilon \to c \neq 0 \to \varphi_{164}$}
\end{ptree}
\noindent$\Pi_{51}$:
\begin{ptree}
\AXC{$\stackrel{\Pi_{50}}{\vdots}$}\noLine
\UIC{$c \in \mathbb{R} \to \mathrm{Conv}(a,\alpha) \to \varepsilon \in \mathbb{R} \to 0 < \varepsilon \to c \neq 0 \to \varphi_{164}$}
\AXC{$\varphi_{330}$\lbl{A1}}
\BIC{$\varphi_{310}$}
\end{ptree}
\noindent$\Pi_{52}$:
\begin{ptree}
\AXC{$\stackrel{\Pi_{2}}{\vdots}$}\noLine
\UIC{$\varepsilon \cdot \lvert c \rvert^{-1} \cdot \lvert c \rvert = \varepsilon \to \varphi_{119}$}
\AXC{$\stackrel{\Pi_{3}}{\vdots}$}\noLine
\UIC{$c \in \mathbb{R} \to \mathrm{Conv}(a,\alpha) \to \varepsilon \in \mathbb{R} \to 0 < \varepsilon \to c \neq 0 \to \varphi_{185}$}
\AXC{$\stackrel{\Pi_{31}}{\vdots}$}\noLine
\UIC{$\varphi_{288}$}
\BIC{$c \in \mathbb{R} \to \mathrm{Conv}(a,\alpha) \to \varepsilon \in \mathbb{R} \to 0 < \varepsilon \to c \neq 0 \to \varphi_{207}$}
\AXC{$\stackrel{\Pi_{51}}{\vdots}$}\noLine
\UIC{$\varphi_{310}$}
\BIC{$\varphi_{257}$}
\BIC{$c \in \mathbb{R} \to \mathrm{Conv}(a,\alpha) \to \varepsilon \in \mathbb{R} \to 0 < \varepsilon \to c \neq 0 \to \varphi_{171}$}
\end{ptree}
\noindent$\Pi_{53}$:
\begin{ptree}
\AXC{$\forall x\,\forall y\,(x \in \mathbb{R} \to y \in \mathbb{R} \to x \cdot y = y \cdot x)$\lbl{(R7)}}
\AXC{$\forall x\,\forall y\,(x \in \mathbb{R} \to y \in \mathbb{R} \to x \cdot y = y \cdot x) \to \varphi_{112}$\lbl{A2}}
\BIC{$\varphi_{112}$}
\AXC{$\varphi_{112} \to \varphi_{131}$\lbl{A2}}
\BIC{$\varphi_{131}$}
\end{ptree}
\noindent$\Pi_{54}$:
\begin{ptree}
\AXC{$\stackrel{\Pi_{28}}{\vdots}$}\noLine
\UIC{$c \in \mathbb{R} \to \mathrm{Conv}(a,\alpha) \to \varepsilon \in \mathbb{R} \to 0 < \varepsilon \to c \neq 0 \to \varphi_{153}$}
\AXC{$\stackrel{\Pi_{53}}{\vdots}$}\noLine
\UIC{$\varphi_{131}$}
\AXC{$\varphi_{131} \to \varphi_{317}$\lbl{A1}}
\BIC{$\varphi_{317}$}
\BIC{$\varphi_{276}$}
\end{ptree}
\noindent$\Pi_{55}$:
\begin{ptree}
\AXC{$\stackrel{\Pi_{6}}{\vdots}$}\noLine
\UIC{$c \in \mathbb{R} \to \mathrm{Conv}(a,\alpha) \to \varepsilon \in \mathbb{R} \to 0 < \varepsilon \to c \neq 0 \to \varphi_{132}$}
\AXC{$\stackrel{\Pi_{54}}{\vdots}$}\noLine
\UIC{$\varphi_{276}$}
\BIC{$c \in \mathbb{R} \to \mathrm{Conv}(a,\alpha) \to \varepsilon \in \mathbb{R} \to 0 < \varepsilon \to c \neq 0 \to \varphi_{206}$}
\AXC{$\varphi_{332}$\lbl{A1}}
\BIC{$\varphi_{306}$}
\end{ptree}
\noindent$\Pi_{56}$:
\begin{ptree}
\AXC{$\stackrel{\Pi_{7}}{\vdots}$}\noLine
\UIC{$\forall x\,\forall y\,(x = y \to y = x)$}
\AXC{$\forall x\,\forall y\,(x = y \to y = x) \to \varphi_{48}$\lbl{A2}}
\BIC{$\varphi_{48}$}
\AXC{$\varphi_{48} \to \varphi_{163}$\lbl{A2}}
\BIC{$\varphi_{163}$}
\end{ptree}
\noindent$\Pi_{57}$:
\begin{ptree}
\AXC{$y = z \to x = y \to x = z$\lbl{A6}}
\AXC{$(y = z \to x = y \to x = z) \to x = y \to y = z \to x = z$\lbl{A1}}
\BIC{$x = y \to y = z \to x = z$}
\UIC{$\forall z\,(x = y \to y = z \to x = z)$}
\UIC{$\forall y\,\forall z\,(x = y \to y = z \to x = z)$}
\UIC{$\forall x\,\forall y\,\forall z\,(x = y \to y = z \to x = z)$}
\AXC{$\forall x\,\forall y\,\forall z\,(x = y \to y = z \to x = z) \to \varphi_{86}$\lbl{A2}}
\BIC{$\varphi_{86}$}
\end{ptree}
\noindent$\Pi_{58}$:
\begin{ptree}
\AXC{$\stackrel{\Pi_{57}}{\vdots}$}\noLine
\UIC{$\varphi_{86}$}
\AXC{$\varphi_{86} \to \varphi_{162}$\lbl{A2}}
\BIC{$\varphi_{162}$}
\AXC{$\varphi_{162} \to \lvert c \odot a(n) - c \cdot \alpha \rvert = \lvert c \cdot (a_{n} - \alpha) \rvert \to \varphi_{147}$\lbl{A2}}
\BIC{$\lvert c \odot a(n) - c \cdot \alpha \rvert = \lvert c \cdot (a_{n} - \alpha) \rvert \to \varphi_{147}$}
\end{ptree}
\noindent$\Pi_{59}$:
\begin{ptree}
\AXC{$\forall x\,\forall y\,(x \in \mathbb{R} \to y \in \mathbb{R} \to \lvert x \cdot y \rvert = \lvert x \rvert \cdot \lvert y \rvert)$\lbl{[A17]}}
\AXC{$\varphi_{149}$\lbl{A2}}
\BIC{$\forall y\,(c \in \mathbb{R} \to y \in \mathbb{R} \to \lvert c \cdot y \rvert = \lvert c \rvert \cdot \lvert y \rvert)$}
\AXC{$\forall y\,(c \in \mathbb{R} \to y \in \mathbb{R} \to \lvert c \cdot y \rvert = \lvert c \rvert \cdot \lvert y \rvert) \to \varphi_{109}$\lbl{A2}}
\BIC{$\varphi_{109}$}
\end{ptree}
\noindent$\Pi_{60}$:
\begin{ptree}
\AXC{$c \in \mathbb{R} \to \mathrm{Conv}(a,\alpha) \to \varepsilon \in \mathbb{R} \to 0 < \varepsilon \to c \neq 0 \to \varphi_{125}$\lbl{A1}}
\AXC{$\stackrel{\Pi_{59}}{\vdots}$}\noLine
\UIC{$\varphi_{109}$}
\AXC{$\varphi_{109} \to \varphi_{311}$\lbl{A1}}
\BIC{$\varphi_{311}$}
\BIC{$\varphi_{277}$}
\end{ptree}
\noindent$\Pi_{61}$:
\begin{ptree}
\AXC{$y = z \to x = y \to x = z$\lbl{A6}}
\AXC{$(y = z \to x = y \to x = z) \to x = y \to y = z \to x = z$\lbl{A1}}
\BIC{$x = y \to y = z \to x = z$}
\UIC{$\forall z\,(x = y \to y = z \to x = z)$}
\UIC{$\forall y\,\forall z\,(x = y \to y = z \to x = z)$}
\UIC{$\forall x\,\forall y\,\forall z\,(x = y \to y = z \to x = z)$}
\AXC{$\forall x\,\forall y\,\forall z\,(x = y \to y = z \to x = z) \to \varphi_{86}$\lbl{A2}}
\BIC{$\varphi_{86}$}
\end{ptree}
\noindent$\Pi_{62}$:
\begin{ptree}
\AXC{$\varphi_{62}$\lbl{[F17]}}
\AXC{$\varphi_{62} \to \varphi_{91}$\lbl{A2}}
\BIC{$\varphi_{91}$}
\AXC{$\varphi_{91} \to \varphi_{105}$\lbl{A2}}
\BIC{$\varphi_{105}$}
\AXC{$\varphi_{105} \to v_{386} \in \mathbb{R} \to v_{387} \in \mathbb{R} \to \varphi_{17}$\lbl{A2}}
\BIC{$v_{386} \in \mathbb{R} \to v_{387} \in \mathbb{R} \to \varphi_{17}$}
\UIC{$\forall v_{388}\,(v_{386} \in \mathbb{R} \to v_{387} \in \mathbb{R} \to \varphi_{17})$}
\UIC{$\forall v_{387}\,\forall v_{388}\,(v_{386} \in \mathbb{R} \to v_{387} \in \mathbb{R} \to \varphi_{17})$}
\UIC{$\forall v_{386}\,\forall v_{387}\,\forall v_{388}\,(v_{386} \in \mathbb{R} \to v_{387} \in \mathbb{R} \to \varphi_{17})$}
\AXC{$\varphi_{227}$\lbl{A2}}
\BIC{$\varphi_{113}$}
\AXC{$\varphi_{113} \to \varphi_{121}$\lbl{A2}}
\BIC{$\varphi_{121}$}
\end{ptree}
\noindent$\Pi_{63}$:
\begin{ptree}
\AXC{$c \in \mathbb{R} \to \mathrm{Conv}(a,\alpha) \to \varepsilon \in \mathbb{R} \to 0 < \varepsilon \to c \neq 0 \to \varphi_{125}$\lbl{A1}}
\AXC{$\stackrel{\Pi_{62}}{\vdots}$}\noLine
\UIC{$\varphi_{121}$}
\AXC{$\varphi_{121} \to \varphi_{106}$\lbl{A2}}
\BIC{$\varphi_{106}$}
\AXC{$\varphi_{106} \to \varphi_{319}$\lbl{A1}}
\BIC{$\varphi_{319}$}
\BIC{$\varphi_{287}$}
\end{ptree}
\noindent$\Pi_{64}$:
\begin{ptree}
\AXC{$\stackrel{\Pi_{19}}{\vdots}$}\noLine
\UIC{$c \in \mathbb{R} \to \mathrm{Conv}(a,\alpha) \to \alpha \in \mathbb{R}$}
\AXC{$\stackrel{\Pi_{24}}{\vdots}$}\noLine
\UIC{$c \in \mathbb{R} \to \mathrm{Conv}(a,\alpha) \to \varepsilon \in \mathbb{R} \to 0 < \varepsilon \to c \neq 0 \to \varphi_{134}$}
\AXC{$\stackrel{\Pi_{63}}{\vdots}$}\noLine
\UIC{$\varphi_{287}$}
\BIC{$\varphi_{233}$}
\BIC{$c \in \mathbb{R} \to \mathrm{Conv}(a,\alpha) \to \varepsilon \in \mathbb{R} \to 0 < \varepsilon \to c \neq 0 \to \varphi_{191}$}
\AXC{$\varphi_{325}$\lbl{A1}}
\BIC{$\lvert c \cdot a_{n} - c \cdot \alpha \rvert = \lvert c \cdot a_{n} - c \cdot \alpha \rvert \to \varphi_{271}$}
\end{ptree}
\noindent$\Pi_{65}$:
\begin{ptree}
\AXC{$x = x$\lbl{A5}}
\UIC{$\forall x\,(x = x)$}
\AXC{$\forall x\,(x = x) \to \lvert c \cdot a_{n} - c \cdot \alpha \rvert = \lvert c \cdot a_{n} - c \cdot \alpha \rvert$\lbl{A2}}
\BIC{$\lvert c \cdot a_{n} - c \cdot \alpha \rvert = \lvert c \cdot a_{n} - c \cdot \alpha \rvert$}
\AXC{$\stackrel{\Pi_{64}}{\vdots}$}\noLine
\UIC{$\lvert c \cdot a_{n} - c \cdot \alpha \rvert = \lvert c \cdot a_{n} - c \cdot \alpha \rvert \to \varphi_{271}$}
\BIC{$\varphi_{271}$}
\end{ptree}
\noindent$\Pi_{66}$:
\begin{ptree}
\AXC{$\forall c\,\forall a\,(c \in \mathbb{R} \to \varphi_{40})$\lbl{[S3]}}
\AXC{$\forall c\,\forall a\,(c \in \mathbb{R} \to \varphi_{40}) \to \forall a\,(c \in \mathbb{R} \to \varphi_{40})$\lbl{A2}}
\BIC{$\forall a\,(c \in \mathbb{R} \to \varphi_{40})$}
\AXC{$\forall a\,(c \in \mathbb{R} \to \varphi_{40}) \to c \in \mathbb{R} \to \varphi_{40}$\lbl{A2}}
\BIC{$c \in \mathbb{R} \to \varphi_{40}$}
\end{ptree}
\noindent$\Pi_{67}$:
\begin{ptree}
\AXC{$\stackrel{\Pi_{21}}{\vdots}$}\noLine
\UIC{$c \in \mathbb{R} \to \mathrm{Conv}(a,\alpha) \to \mathrm{Seq}(a)$}
\AXC{$c \in \mathbb{R} \to c \in \mathbb{R}$\lbl{A1}}
\AXC{$\stackrel{\Pi_{66}}{\vdots}$}\noLine
\UIC{$c \in \mathbb{R} \to \varphi_{40}$}
\AXC{$(c \in \mathbb{R} \to \varphi_{40}) \to \varphi_{198}$\lbl{A1}}
\BIC{$\varphi_{198}$}
\BIC{$(c \in \mathbb{R} \to \mathrm{Conv}(a,\alpha) \to \mathrm{Seq}(a)) \to c \in \mathbb{R} \to \varphi_{54}$}
\BIC{$c \in \mathbb{R} \to \varphi_{54}$}
\end{ptree}
\noindent$\Pi_{68}$:
\begin{ptree}
\AXC{$\stackrel{\Pi_{67}}{\vdots}$}\noLine
\UIC{$c \in \mathbb{R} \to \varphi_{54}$}
\AXC{$(c \in \mathbb{R} \to \varphi_{54}) \to \varphi_{38}$\lbl{A1}}
\BIC{$\varphi_{38}$}
\end{ptree}
\noindent$\Pi_{69}$:
\begin{ptree}
\AXC{$c \in \mathbb{R} \to \mathrm{Conv}(a,\alpha) \to \varepsilon \in \mathbb{R} \to 0 < \varepsilon \to c \neq 0 \to \varphi_{117}$\lbl{A1}}
\AXC{$\varphi_{101}$\lbl{A2}}
\AXC{$\stackrel{\Pi_{68}}{\vdots}$}\noLine
\UIC{$\varphi_{38}$}
\AXC{$\varphi_{38} \to \varphi_{101} \to \varphi_{28}$\lbl{A1}}
\BIC{$\varphi_{101} \to \varphi_{28}$}
\BIC{$\varphi_{28}$}
\AXC{$\varphi_{28} \to \varphi_{264}$\lbl{A1}}
\BIC{$\varphi_{264}$}
\BIC{$c \in \mathbb{R} \to \mathrm{Conv}(a,\alpha) \to \varepsilon \in \mathbb{R} \to 0 < \varepsilon \to c \neq 0 \to \varphi_{157}$}
\end{ptree}
\noindent$\Pi_{70}$:
\begin{ptree}
\AXC{$x = x$\lbl{A5}}
\UIC{$\forall x\,(x = x)$}
\AXC{$\varphi_{49}$\lbl{A2}}
\BIC{$\lvert c \odot a(n) - c \cdot \alpha \rvert = \lvert c \odot a(n) - c \cdot \alpha \rvert$}
\AXC{$\stackrel{\Pi_{69}}{\vdots}$}\noLine
\UIC{$c \in \mathbb{R} \to \mathrm{Conv}(a,\alpha) \to \varepsilon \in \mathbb{R} \to 0 < \varepsilon \to c \neq 0 \to \varphi_{157}$}
\AXC{$\varphi_{326}$\lbl{A1}}
\BIC{$\lvert c \odot a(n) - c \cdot \alpha \rvert = \lvert c \odot a(n) - c \cdot \alpha \rvert \to \varphi_{274}$}
\BIC{$\varphi_{274}$}
\end{ptree}
\noindent$\Pi_{71}$:
\begin{ptree}
\AXC{$z = w \to \varphi_{123}$\lbl{A6}}
\UIC{$\forall w\,(z = w \to \varphi_{123})$}
\UIC{$\forall z\,\forall w\,(z = w \to \varphi_{123})$}
\AXC{$\forall z\,\forall w\,(z = w \to \varphi_{123}) \to \varphi_{194}$\lbl{A2}}
\BIC{$\varphi_{194}$}
\AXC{$\varphi_{194} \to c \odot a(n) = c \cdot a_{n} \to \varphi_{169}$\lbl{A2}}
\BIC{$c \odot a(n) = c \cdot a_{n} \to \varphi_{169}$}
\AXC{$\stackrel{\Pi_{70}}{\vdots}$}\noLine
\UIC{$\varphi_{274}$}
\BIC{$c \in \mathbb{R} \to \mathrm{Conv}(a,\alpha) \to \varepsilon \in \mathbb{R} \to 0 < \varepsilon \to \varphi_{214}$}
\end{ptree}
\noindent$\Pi_{72}$:
\begin{ptree}
\AXC{$\stackrel{\Pi_{71}}{\vdots}$}\noLine
\UIC{$c \in \mathbb{R} \to \mathrm{Conv}(a,\alpha) \to \varepsilon \in \mathbb{R} \to 0 < \varepsilon \to \varphi_{214}$}
\AXC{$(c \in \mathbb{R} \to \mathrm{Conv}(a,\alpha) \to \varepsilon \in \mathbb{R} \to 0 < \varepsilon \to \varphi_{214}) \to \varphi_{321}$\lbl{A1}}
\BIC{$\varphi_{321}$}
\end{ptree}
\noindent$\Pi_{73}$:
\begin{ptree}
\AXC{$\varphi_{77}$\lbl{A6}}
\UIC{$\forall w\,(\varphi_{77})$}
\UIC{$\forall z\,\forall w\,(\varphi_{77})$}
\AXC{$\forall z\,\forall w\,(\varphi_{77}) \to \varphi_{166}$\lbl{A2}}
\BIC{$\varphi_{166}$}
\AXC{$\varphi_{166} \to c \cdot a_{n} - c \cdot \alpha = c \cdot (a_{n} - \alpha) \to \varphi_{138}$\lbl{A2}}
\BIC{$c \cdot a_{n} - c \cdot \alpha = c \cdot (a_{n} - \alpha) \to \varphi_{138}$}
\AXC{$\stackrel{\Pi_{65}}{\vdots}$}\noLine
\UIC{$\varphi_{271}$}
\BIC{$c \in \mathbb{R} \to \mathrm{Conv}(a,\alpha) \to \varepsilon \in \mathbb{R} \to 0 < \varepsilon \to c \neq 0 \to \varphi_{200}$}
\AXC{$\stackrel{\Pi_{72}}{\vdots}$}\noLine
\UIC{$\varphi_{321}$}
\BIC{$\varphi_{281}$}
\end{ptree}
\noindent$\Pi_{74}$:
\begin{ptree}
\AXC{$\stackrel{\Pi_{61}}{\vdots}$}\noLine
\UIC{$\varphi_{86}$}
\AXC{$\varphi_{86} \to \varphi_{181}$\lbl{A2}}
\BIC{$\varphi_{181}$}
\AXC{$\varphi_{181} \to \lvert c \odot a(n) - c \cdot \alpha \rvert = \lvert c \cdot a_{n} - c \cdot \alpha \rvert \to \varphi_{135}$\lbl{A2}}
\BIC{$\lvert c \odot a(n) - c \cdot \alpha \rvert = \lvert c \cdot a_{n} - c \cdot \alpha \rvert \to \varphi_{135}$}
\AXC{$\stackrel{\Pi_{73}}{\vdots}$}\noLine
\UIC{$\varphi_{281}$}
\BIC{$c \in \mathbb{R} \to \mathrm{Conv}(a,\alpha) \to \varepsilon \in \mathbb{R} \to 0 < \varepsilon \to \varphi_{212}$}
\end{ptree}
\noindent$\Pi_{75}$:
\begin{ptree}
\AXC{$\stackrel{\Pi_{74}}{\vdots}$}\noLine
\UIC{$c \in \mathbb{R} \to \mathrm{Conv}(a,\alpha) \to \varepsilon \in \mathbb{R} \to 0 < \varepsilon \to \varphi_{212}$}
\AXC{$(c \in \mathbb{R} \to \mathrm{Conv}(a,\alpha) \to \varepsilon \in \mathbb{R} \to 0 < \varepsilon \to \varphi_{212}) \to \varphi_{322}$\lbl{A1}}
\BIC{$\varphi_{322}$}
\end{ptree}
\noindent$\Pi_{76}$:
\begin{ptree}
\AXC{$\stackrel{\Pi_{58}}{\vdots}$}\noLine
\UIC{$\lvert c \odot a(n) - c \cdot \alpha \rvert = \lvert c \cdot (a_{n} - \alpha) \rvert \to \varphi_{147}$}
\AXC{$\stackrel{\Pi_{26}}{\vdots}$}\noLine
\UIC{$c \in \mathbb{R} \to \mathrm{Conv}(a,\alpha) \to \varepsilon \in \mathbb{R} \to 0 < \varepsilon \to c \neq 0 \to \varphi_{145}$}
\AXC{$\stackrel{\Pi_{60}}{\vdots}$}\noLine
\UIC{$\varphi_{277}$}
\BIC{$c \in \mathbb{R} \to \mathrm{Conv}(a,\alpha) \to \varepsilon \in \mathbb{R} \to 0 < \varepsilon \to c \neq 0 \to \varphi_{199}$}
\AXC{$\stackrel{\Pi_{75}}{\vdots}$}\noLine
\UIC{$\varphi_{322}$}
\BIC{$\varphi_{283}$}
\BIC{$c \in \mathbb{R} \to \mathrm{Conv}(a,\alpha) \to \varepsilon \in \mathbb{R} \to 0 < \varepsilon \to \varphi_{215}$}
\end{ptree}
\noindent$\Pi_{77}$:
\begin{ptree}
\AXC{$\stackrel{\Pi_{56}}{\vdots}$}\noLine
\UIC{$\varphi_{163}$}
\AXC{$\stackrel{\Pi_{76}}{\vdots}$}\noLine
\UIC{$c \in \mathbb{R} \to \mathrm{Conv}(a,\alpha) \to \varepsilon \in \mathbb{R} \to 0 < \varepsilon \to \varphi_{215}$}
\AXC{$\varphi_{315}$\lbl{A1}}
\BIC{$\varphi_{163} \to c \in \mathbb{R} \to \mathrm{Conv}(a,\alpha) \to \varepsilon \in \mathbb{R} \to 0 < \varepsilon \to \varphi_{216}$}
\BIC{$c \in \mathbb{R} \to \mathrm{Conv}(a,\alpha) \to \varepsilon \in \mathbb{R} \to 0 < \varepsilon \to \varphi_{216}$}
\end{ptree}
\noindent$\Pi_{78}$:
\begin{ptree}
\AXC{$\stackrel{\Pi_{77}}{\vdots}$}\noLine
\UIC{$c \in \mathbb{R} \to \mathrm{Conv}(a,\alpha) \to \varepsilon \in \mathbb{R} \to 0 < \varepsilon \to \varphi_{216}$}
\AXC{$(c \in \mathbb{R} \to \mathrm{Conv}(a,\alpha) \to \varepsilon \in \mathbb{R} \to 0 < \varepsilon \to \varphi_{216}) \to \varphi_{308}$\lbl{A1}}
\BIC{$\varphi_{308}$}
\end{ptree}
\noindent$\Pi_{79}$:
\begin{ptree}
\AXC{$\stackrel{\Pi_{1}}{\vdots}$}\noLine
\UIC{$\varphi_{186}$}
\AXC{$\stackrel{\Pi_{52}}{\vdots}$}\noLine
\UIC{$c \in \mathbb{R} \to \mathrm{Conv}(a,\alpha) \to \varepsilon \in \mathbb{R} \to 0 < \varepsilon \to c \neq 0 \to \varphi_{171}$}
\AXC{$\stackrel{\Pi_{55}}{\vdots}$}\noLine
\UIC{$\varphi_{306}$}
\BIC{$\varphi_{258}$}
\BIC{$c \in \mathbb{R} \to \mathrm{Conv}(a,\alpha) \to \varepsilon \in \mathbb{R} \to 0 < \varepsilon \to c \neq 0 \to \varphi_{172}$}
\AXC{$\stackrel{\Pi_{78}}{\vdots}$}\noLine
\UIC{$\varphi_{308}$}
\BIC{$\varphi_{263}$}
\end{ptree}
\noindent$\Pi_{80}$:
\begin{ptree}
\AXC{$z = w \to z < \varepsilon \to w < \varepsilon$\lbl{A6}}
\UIC{$\forall w\,(z = w \to z < \varepsilon \to w < \varepsilon)$}
\UIC{$\forall z\,\forall w\,(z = w \to z < \varepsilon \to w < \varepsilon)$}
\AXC{$\forall z\,\forall w\,(z = w \to z < \varepsilon \to w < \varepsilon) \to \varphi_{87}$\lbl{A2}}
\BIC{$\varphi_{87}$}
\AXC{$\varphi_{87} \to \varphi_{196}$\lbl{A2}}
\BIC{$\varphi_{196}$}
\AXC{$\stackrel{\Pi_{79}}{\vdots}$}\noLine
\UIC{$\varphi_{263}$}
\BIC{$c \in \mathbb{R} \to \mathrm{Conv}(a,\alpha) \to \varepsilon \in \mathbb{R} \to 0 < \varepsilon \to c \neq 0 \to \varphi_{179}$}
\AXC{$\varphi_{278}$\lbl{A1}}
\BIC{$\varphi_{177} \to n \in \mathbb{N} \to N \leq n \to \lvert c \odot a(n) - c \cdot \alpha \rvert < \varepsilon$}
\UIC{$\forall n\,(\varphi_{177} \to n \in \mathbb{N} \to N \leq n \to \lvert c \odot a(n) - c \cdot \alpha \rvert < \varepsilon)$}
\end{ptree}
\noindent$\Pi_{81}$:
\begin{ptree}
\AXC{$\varphi_{177} \to \forall n\,(\varphi_{177})$\lbl{A4}}
\AXC{$\varphi_{282}$\lbl{A3}}
\AXC{$\stackrel{\Pi_{80}}{\vdots}$}\noLine
\UIC{$\forall n\,(\varphi_{177} \to n \in \mathbb{N} \to N \leq n \to \lvert c \odot a(n) - c \cdot \alpha \rvert < \varepsilon)$}
\AXC{$\varphi_{336}$\lbl{A1}}
\BIC{$\varphi_{282} \to \varphi_{300}$}
\BIC{$\varphi_{300}$}
\BIC{$c \in \mathbb{R} \to \mathrm{Conv}(a,\alpha) \to \varepsilon \in \mathbb{R} \to 0 < \varepsilon \to c \neq 0 \to \varphi_{188}$}
\end{ptree}
\noindent$\Pi_{82}$:
\begin{ptree}
\AXC{$c \in \mathbb{R} \to \mathrm{Conv}(a,\alpha) \to \varepsilon \in \mathbb{R} \to 0 < \varepsilon \to c \neq 0 \to \varphi_{202}$\lbl{A1}}
\AXC{$\varphi_{266}$\lbl{A1}}
\BIC{$c \in \mathbb{R} \to \mathrm{Conv}(a,\alpha) \to \varepsilon \in \mathbb{R} \to 0 < \varepsilon \to c \neq 0 \to \varphi_{94}$}
\AXC{$\varphi_{312}$\lbl{A1}}
\BIC{$\varphi_{285}$}
\end{ptree}
\noindent$\Pi_{83}$:
\begin{ptree}
\AXC{$\forall N\,(\varphi_{45}) \to \varphi_{45}$\lbl{A2}}
\AXC{$\stackrel{\Pi_{81}}{\vdots}$}\noLine
\UIC{$c \in \mathbb{R} \to \mathrm{Conv}(a,\alpha) \to \varepsilon \in \mathbb{R} \to 0 < \varepsilon \to c \neq 0 \to \varphi_{188}$}
\AXC{$\stackrel{\Pi_{82}}{\vdots}$}\noLine
\UIC{$\varphi_{285}$}
\BIC{$c \in \mathbb{R} \to \mathrm{Conv}(a,\alpha) \to \varepsilon \in \mathbb{R} \to 0 < \varepsilon \to c \neq 0 \to \varphi_{201}$}
\AXC{$\varphi_{318}$\lbl{A1}}
\BIC{$\varphi_{269}$}
\BIC{$c \in \mathbb{R} \to \mathrm{Conv}(a,\alpha) \to \varepsilon \in \mathbb{R} \to 0 < \varepsilon \to c \neq 0 \to \varphi_{205}$}
\AXC{$\varphi_{293}$\lbl{A1}}
\BIC{$\varphi_{180} \to \varphi_{56}$}
\UIC{$\forall N\,(\varphi_{180} \to \varphi_{56})$}
\end{ptree}
\noindent$\Pi_{84}$:
\begin{ptree}
\AXC{$\forall x \in \mathbb{R}\,(\lvert x \rvert \in \mathbb{R})$\lbl{[A5]}}
\AXC{$\forall x \in \mathbb{R}\,(\lvert x \rvert \in \mathbb{R}) \to c \in \mathbb{R} \to \lvert c \rvert \in \mathbb{R}$\lbl{A2}}
\BIC{$c \in \mathbb{R} \to \lvert c \rvert \in \mathbb{R}$}
\AXC{$(c \in \mathbb{R} \to \lvert c \rvert \in \mathbb{R}) \to \varphi_{30} \to \varphi_{50}$\lbl{A1}}
\BIC{$\varphi_{30} \to \varphi_{50}$}
\end{ptree}
\noindent$\Pi_{85}$:
\begin{ptree}
\AXC{$\varphi_{30}$\lbl{A1}}
\AXC{$\stackrel{\Pi_{84}}{\vdots}$}\noLine
\UIC{$\varphi_{30} \to \varphi_{50}$}
\BIC{$\varphi_{50}$}
\end{ptree}
\noindent$\Pi_{86}$:
\begin{ptree}
\AXC{$\stackrel{\Pi_{4}}{\vdots}$}\noLine
\UIC{$\varphi_{34}$}
\AXC{$\stackrel{\Pi_{85}}{\vdots}$}\noLine
\UIC{$\varphi_{50}$}
\AXC{$\forall x \in \mathbb{R}\,(0 < x \to 0 < x^{-1})$\lbl{[F34]}}
\AXC{$\varphi_{97}$\lbl{A2}}
\BIC{$\lvert c \rvert \in \mathbb{R} \to 0 < \lvert c \rvert \to 0 < \lvert c \rvert^{-1}$}
\AXC{$(\lvert c \rvert \in \mathbb{R} \to 0 < \lvert c \rvert \to 0 < \lvert c \rvert^{-1}) \to \varphi_{50} \to \varphi_{34} \to \varphi_{52}$\lbl{A1}}
\BIC{$\varphi_{50} \to \varphi_{34} \to \varphi_{52}$}
\BIC{$\varphi_{34} \to \varphi_{52}$}
\BIC{$\varphi_{52}$}
\end{ptree}
\noindent$\Pi_{87}$:
\begin{ptree}
\AXC{$\stackrel{\Pi_{85}}{\vdots}$}\noLine
\UIC{$\varphi_{50}$}
\AXC{$\forall x \in \mathbb{R}\,(x \neq 0 \to x^{-1} \in \mathbb{R})$\lbl{(R1)}}
\AXC{$\varphi_{104}$\lbl{A2}}
\BIC{$\lvert c \rvert \in \mathbb{R} \to \lvert c \rvert \neq 0 \to \lvert c \rvert^{-1} \in \mathbb{R}$}
\AXC{$(\lvert c \rvert \in \mathbb{R} \to \lvert c \rvert \neq 0 \to \lvert c \rvert^{-1} \in \mathbb{R}) \to \varphi_{50} \to \varphi_{27} \to \varphi_{72}$\lbl{A1}}
\BIC{$\varphi_{50} \to \varphi_{27} \to \varphi_{72}$}
\BIC{$\varphi_{27} \to \varphi_{72}$}
\end{ptree}
\noindent$\Pi_{88}$:
\begin{ptree}
\AXC{$\stackrel{\Pi_{12}}{\vdots}$}\noLine
\UIC{$\varphi_{27}$}
\AXC{$\stackrel{\Pi_{87}}{\vdots}$}\noLine
\UIC{$\varphi_{27} \to \varphi_{72}$}
\BIC{$\varphi_{72}$}
\end{ptree}
\noindent$\Pi_{89}$:
\begin{ptree}
\AXC{$\forall x\,\forall y\,(x \in \mathbb{R} \to y \in \mathbb{R} \to 0 < x \to 0 < y \to 0 < x \cdot y)$\lbl{(R15)}}
\AXC{$\varphi_{165}$\lbl{A2}}
\BIC{$\forall y\,(\varepsilon \in \mathbb{R} \to y \in \mathbb{R} \to 0 < \varepsilon \to 0 < y \to 0 < \varepsilon \cdot y)$}
\AXC{$\varphi_{209}$\lbl{A2}}
\BIC{$\varepsilon \in \mathbb{R} \to \lvert c \rvert^{-1} \in \mathbb{R} \to 0 < \varepsilon \to 0 < \lvert c \rvert^{-1} \to 0 < \varepsilon \cdot \lvert c \rvert^{-1}$}
\AXC{$\varphi_{298}$\lbl{A1}}
\BIC{$(c \in \mathbb{R} \to \mathrm{Conv}(a,\alpha) \to \varepsilon \in \mathbb{R} \to \varepsilon \in \mathbb{R}) \to \varphi_{72} \to \varphi_{232}$}
\end{ptree}
\noindent$\Pi_{90}$:
\begin{ptree}
\AXC{$\stackrel{\Pi_{88}}{\vdots}$}\noLine
\UIC{$\varphi_{72}$}
\AXC{$c \in \mathbb{R} \to \mathrm{Conv}(a,\alpha) \to \varepsilon \in \mathbb{R} \to \varepsilon \in \mathbb{R}$\lbl{A1}}
\AXC{$\stackrel{\Pi_{89}}{\vdots}$}\noLine
\UIC{$(c \in \mathbb{R} \to \mathrm{Conv}(a,\alpha) \to \varepsilon \in \mathbb{R} \to \varepsilon \in \mathbb{R}) \to \varphi_{72} \to \varphi_{232}$}
\BIC{$\varphi_{72} \to \varphi_{232}$}
\BIC{$\varphi_{232}$}
\end{ptree}
\noindent$\Pi_{91}$:
\begin{ptree}
\AXC{$\stackrel{\Pi_{86}}{\vdots}$}\noLine
\UIC{$\varphi_{52}$}
\AXC{$c \in \mathbb{R} \to \mathrm{Conv}(a,\alpha) \to \varepsilon \in \mathbb{R} \to 0 < \varepsilon \to 0 < \varepsilon$\lbl{A1}}
\AXC{$\stackrel{\Pi_{90}}{\vdots}$}\noLine
\UIC{$\varphi_{232}$}
\BIC{$\varphi_{52} \to \varphi_{82}$}
\BIC{$\varphi_{82}$}
\end{ptree}
\noindent$\Pi_{92}$:
\begin{ptree}
\AXC{$\stackrel{\Pi_{88}}{\vdots}$}\noLine
\UIC{$\varphi_{72}$}
\AXC{$\varphi_{39}$\lbl{A1}}
\AXC{$\forall x\,\forall y\,(x \in \mathbb{R} \to y \in \mathbb{R} \to x \cdot y \in \mathbb{R})$\lbl{(R1)}}
\AXC{$\varphi_{116}$\lbl{A2}}
\BIC{$\forall y\,(\varepsilon \in \mathbb{R} \to y \in \mathbb{R} \to \varepsilon \cdot y \in \mathbb{R})$}
\AXC{$\varphi_{137}$\lbl{A2}}
\BIC{$\varepsilon \in \mathbb{R} \to \lvert c \rvert^{-1} \in \mathbb{R} \to \varepsilon \cdot \lvert c \rvert^{-1} \in \mathbb{R}$}
\AXC{$\varphi_{252}$\lbl{A1}}
\BIC{$\varphi_{39} \to \varphi_{72} \to c \in \mathbb{R} \to \varphi_{46}$}
\BIC{$\varphi_{72} \to c \in \mathbb{R} \to \varphi_{46}$}
\BIC{$c \in \mathbb{R} \to \varphi_{46}$}
\end{ptree}
\noindent$\Pi_{93}$:
\begin{ptree}
\AXC{$\forall a\,\forall \alpha\,(\varphi_{127})$\lbl{[C3]}}
\AXC{$\forall a\,\forall \alpha\,(\varphi_{127}) \to \forall \alpha\,(\varphi_{127})$\lbl{A2}}
\BIC{$\forall \alpha\,(\varphi_{127})$}
\AXC{$\forall \alpha\,(\varphi_{127}) \to \varphi_{127}$\lbl{A2}}
\BIC{$\varphi_{127}$}
\AXC{$\varphi_{127} \to \mathrm{Conv}(a,\alpha) \to \varepsilon \cdot \lvert c \rvert^{-1} \in \mathbb{R} \to \varphi_{124}$\lbl{A2}}
\BIC{$\mathrm{Conv}(a,\alpha) \to \varepsilon \cdot \lvert c \rvert^{-1} \in \mathbb{R} \to \varphi_{124}$}
\end{ptree}
\noindent$\Pi_{94}$:
\begin{ptree}
\AXC{$\stackrel{\Pi_{92}}{\vdots}$}\noLine
\UIC{$c \in \mathbb{R} \to \varphi_{46}$}
\AXC{$c \in \mathbb{R} \to \mathrm{Conv}(a,\alpha) \to \mathrm{Conv}(a,\alpha)$\lbl{A1}}
\AXC{$\stackrel{\Pi_{93}}{\vdots}$}\noLine
\UIC{$\mathrm{Conv}(a,\alpha) \to \varepsilon \cdot \lvert c \rvert^{-1} \in \mathbb{R} \to \varphi_{124}$}
\AXC{$\varphi_{304}$\lbl{A1}}
\BIC{$(c \in \mathbb{R} \to \mathrm{Conv}(a,\alpha) \to \mathrm{Conv}(a,\alpha)) \to \varphi_{253}$}
\BIC{$\varphi_{253}$}
\BIC{$\varphi_{82} \to c \in \mathbb{R} \to \mathrm{Conv}(a,\alpha) \to \varepsilon \in \mathbb{R} \to 0 < \varepsilon \to \varphi_{90}$}
\end{ptree}
\noindent$\Pi_{95}$:
\begin{ptree}
\AXC{$\stackrel{\Pi_{91}}{\vdots}$}\noLine
\UIC{$\varphi_{82}$}
\AXC{$\stackrel{\Pi_{94}}{\vdots}$}\noLine
\UIC{$\varphi_{82} \to c \in \mathbb{R} \to \mathrm{Conv}(a,\alpha) \to \varepsilon \in \mathbb{R} \to 0 < \varepsilon \to \varphi_{90}$}
\BIC{$c \in \mathbb{R} \to \mathrm{Conv}(a,\alpha) \to \varepsilon \in \mathbb{R} \to 0 < \varepsilon \to \varphi_{90}$}
\AXC{$(c \in \mathbb{R} \to \mathrm{Conv}(a,\alpha) \to \varepsilon \in \mathbb{R} \to 0 < \varepsilon \to \varphi_{90}) \to \varphi_{337}$\lbl{A1}}
\BIC{$\varphi_{337}$}
\end{ptree}
\noindent$\Pi_{96}$:
\begin{ptree}
\AXC{$\varphi_{180} \to \forall N\,(\varphi_{180})$\lbl{A4}}
\AXC{$\forall N\,(\varphi_{180} \to \varphi_{56}) \to \forall N\,(\varphi_{180}) \to \forall N\,(\varphi_{56})$\lbl{A3}}
\AXC{$\stackrel{\Pi_{83}}{\vdots}$}\noLine
\UIC{$\forall N\,(\varphi_{180} \to \varphi_{56})$}
\AXC{$\stackrel{\Pi_{95}}{\vdots}$}\noLine
\UIC{$\varphi_{337}$}
\BIC{$(\forall N\,(\varphi_{180} \to \varphi_{56}) \to \forall N\,(\varphi_{180}) \to \forall N\,(\varphi_{56})) \to \varphi_{289}$}
\BIC{$\varphi_{289}$}
\BIC{$c \in \mathbb{R} \to \mathrm{Conv}(a,\alpha) \to \varepsilon \in \mathbb{R} \to 0 < \varepsilon \to \varphi_{85}$}
\end{ptree}
\noindent$\Pi_{97}$:
\begin{ptree}
\AXC{$z = w \to z < \varepsilon \to w < \varepsilon$\lbl{A6}}
\UIC{$\forall w\,(z = w \to z < \varepsilon \to w < \varepsilon)$}
\UIC{$\forall z\,\forall w\,(z = w \to z < \varepsilon \to w < \varepsilon)$}
\AXC{$\varphi_{14}$\lbl{A2}}
\BIC{$\forall w\,(0 = w \to 0 < \varepsilon \to w < \varepsilon)$}
\AXC{$\forall w\,(0 = w \to 0 < \varepsilon \to w < \varepsilon) \to \varphi_{98}$\lbl{A2}}
\BIC{$\varphi_{98}$}
\end{ptree}
\noindent$\Pi_{98}$:
\begin{ptree}
\AXC{$y = z \to x = y \to x = z$\lbl{A6}}
\AXC{$(y = z \to x = y \to x = z) \to x = y \to y = z \to x = z$\lbl{A1}}
\BIC{$x = y \to y = z \to x = z$}
\UIC{$\forall z\,(x = y \to y = z \to x = z)$}
\UIC{$\forall y\,\forall z\,(x = y \to y = z \to x = z)$}
\UIC{$\forall x\,\forall y\,\forall z\,(x = y \to y = z \to x = z)$}
\AXC{$\forall x\,\forall y\,\forall z\,(x = y \to y = z \to x = z) \to \varphi_{86}$\lbl{A2}}
\BIC{$\varphi_{86}$}
\end{ptree}
\noindent$\Pi_{99}$:
\begin{ptree}
\AXC{$y = z \to x = y \to x = z$\lbl{A6}}
\AXC{$(y = z \to x = y \to x = z) \to x = y \to y = z \to x = z$\lbl{A1}}
\BIC{$x = y \to y = z \to x = z$}
\UIC{$\forall z\,(x = y \to y = z \to x = z)$}
\UIC{$\forall y\,\forall z\,(x = y \to y = z \to x = z)$}
\UIC{$\forall x\,\forall y\,\forall z\,(x = y \to y = z \to x = z)$}
\AXC{$\forall x\,\forall y\,\forall z\,(x = y \to y = z \to x = z) \to \varphi_{68}$\lbl{A2}}
\BIC{$\varphi_{68}$}
\end{ptree}
\noindent$\Pi_{100}$:
\begin{ptree}
\AXC{$\forall a\,\forall n\,(\mathrm{Seq}(a) \to n \in \mathbb{N} \to a_{n} \in \mathbb{R})$\lbl{[S1]}}
\AXC{$\varphi_{99}$\lbl{A2}}
\BIC{$\forall n\,(\mathrm{Seq}(a) \to n \in \mathbb{N} \to a_{n} \in \mathbb{R})$}
\AXC{$\varphi_{80}$\lbl{A2}}
\BIC{$\mathrm{Seq}(a) \to n \in \mathbb{N} \to a_{n} \in \mathbb{R}$}
\AXC{$(\mathrm{Seq}(a) \to n \in \mathbb{N} \to a_{n} \in \mathbb{R}) \to \varphi_{226}$\lbl{A1}}
\BIC{$\varphi_{226}$}
\end{ptree}
\noindent$\Pi_{101}$:
\begin{ptree}
\AXC{$c \in \mathbb{R} \to \varphi_{21}$\lbl{A1}}
\AXC{$\stackrel{\Pi_{21}}{\vdots}$}\noLine
\UIC{$c \in \mathbb{R} \to \mathrm{Conv}(a,\alpha) \to \mathrm{Seq}(a)$}
\AXC{$\stackrel{\Pi_{100}}{\vdots}$}\noLine
\UIC{$\varphi_{226}$}
\BIC{$(c \in \mathbb{R} \to \varphi_{21}) \to c \in \mathbb{R} \to \mathrm{Conv}(a,\alpha) \to \varphi_{35}$}
\BIC{$c \in \mathbb{R} \to \mathrm{Conv}(a,\alpha) \to \varphi_{35}$}
\end{ptree}
\noindent$\Pi_{102}$:
\begin{ptree}
\AXC{$\stackrel{\Pi_{101}}{\vdots}$}\noLine
\UIC{$c \in \mathbb{R} \to \mathrm{Conv}(a,\alpha) \to \varphi_{35}$}
\AXC{$\forall x\,\forall y\,(x \in \mathbb{R} \to y \in \mathbb{R} \to x - y \in \mathbb{R})$\lbl{[F1]}}
\AXC{$\varphi_{130}$\lbl{A2}}
\BIC{$\forall y\,(a_{n} \in \mathbb{R} \to y \in \mathbb{R} \to a_{n} - y \in \mathbb{R})$}
\AXC{$\varphi_{143}$\lbl{A2}}
\BIC{$a_{n} \in \mathbb{R} \to \alpha \in \mathbb{R} \to a_{n} - \alpha \in \mathbb{R}$}
\AXC{$\varphi_{251}$\lbl{A1}}
\BIC{$(c \in \mathbb{R} \to \mathrm{Conv}(a,\alpha) \to \varphi_{35}) \to \varphi_{195}$}
\BIC{$\varphi_{195}$}
\end{ptree}
\noindent$\Pi_{103}$:
\begin{ptree}
\AXC{$\stackrel{\Pi_{19}}{\vdots}$}\noLine
\UIC{$c \in \mathbb{R} \to \mathrm{Conv}(a,\alpha) \to \alpha \in \mathbb{R}$}
\AXC{$\stackrel{\Pi_{102}}{\vdots}$}\noLine
\UIC{$\varphi_{195}$}
\BIC{$c \in \mathbb{R} \to \mathrm{Conv}(a,\alpha) \to \varphi_{65}$}
\end{ptree}
\noindent$\Pi_{104}$:
\begin{ptree}
\AXC{$\stackrel{\Pi_{103}}{\vdots}$}\noLine
\UIC{$c \in \mathbb{R} \to \mathrm{Conv}(a,\alpha) \to \varphi_{65}$}
\AXC{$\forall x \in \mathbb{R}\,(\lvert x \rvert \in \mathbb{R})$\lbl{[A5]}}
\AXC{$\varphi_{100}$\lbl{A2}}
\BIC{$a_{n} - \alpha \in \mathbb{R} \to \lvert a_{n} - \alpha \rvert \in \mathbb{R}$}
\AXC{$(a_{n} - \alpha \in \mathbb{R} \to \lvert a_{n} - \alpha \rvert \in \mathbb{R}) \to \varphi_{225}$\lbl{A1}}
\BIC{$\varphi_{225}$}
\BIC{$c \in \mathbb{R} \to \mathrm{Conv}(a,\alpha) \to \varphi_{78}$}
\end{ptree}
\noindent$\Pi_{105}$:
\begin{ptree}
\AXC{$\stackrel{\Pi_{104}}{\vdots}$}\noLine
\UIC{$c \in \mathbb{R} \to \mathrm{Conv}(a,\alpha) \to \varphi_{78}$}
\AXC{$\forall x \in \mathbb{R}\,(x \cdot 0 = 0)$\lbl{[F8]}}
\AXC{$\varphi_{89}$\lbl{A2}}
\BIC{$\lvert a_{n} - \alpha \rvert \in \mathbb{R} \to \lvert a_{n} - \alpha \rvert \cdot 0 = 0$}
\AXC{$(\lvert a_{n} - \alpha \rvert \in \mathbb{R} \to \lvert a_{n} - \alpha \rvert \cdot 0 = 0) \to \varphi_{228}$\lbl{A1}}
\BIC{$\varphi_{228}$}
\BIC{$c \in \mathbb{R} \to \mathrm{Conv}(a,\alpha) \to \varepsilon \in \mathbb{R} \to \varphi_{16}$}
\end{ptree}
\noindent$\Pi_{106}$:
\begin{ptree}
\AXC{$y = z \to x = y \to x = z$\lbl{A6}}
\AXC{$(y = z \to x = y \to x = z) \to x = y \to y = z \to x = z$\lbl{A1}}
\BIC{$x = y \to y = z \to x = z$}
\UIC{$\forall z\,(x = y \to y = z \to x = z)$}
\UIC{$\forall y\,\forall z\,(x = y \to y = z \to x = z)$}
\UIC{$\forall x\,\forall y\,\forall z\,(x = y \to y = z \to x = z)$}
\AXC{$\forall x\,\forall y\,\forall z\,(x = y \to y = z \to x = z) \to \varphi_{68}$\lbl{A2}}
\BIC{$\varphi_{68}$}
\end{ptree}
\noindent$\Pi_{107}$:
\begin{ptree}
\AXC{$\forall x\,\forall y\,(x \in \mathbb{R} \to y \in \mathbb{R} \to x \cdot y = y \cdot x)$\lbl{(R7)}}
\AXC{$\varphi_{102}$\lbl{A2}}
\BIC{$\forall y\,(0 \in \mathbb{R} \to y \in \mathbb{R} \to 0 \cdot y = y \cdot 0)$}
\AXC{$\forall y\,(0 \in \mathbb{R} \to y \in \mathbb{R} \to 0 \cdot y = y \cdot 0) \to \varphi_{110}$\lbl{A2}}
\BIC{$\varphi_{110}$}
\end{ptree}
\noindent$\Pi_{108}$:
\begin{ptree}
\AXC{$\stackrel{\Pi_{104}}{\vdots}$}\noLine
\UIC{$c \in \mathbb{R} \to \mathrm{Conv}(a,\alpha) \to \varphi_{78}$}
\AXC{$0 \in \mathbb{R}$\lbl{(R1)}}
\AXC{$\stackrel{\Pi_{107}}{\vdots}$}\noLine
\UIC{$\varphi_{110}$}
\AXC{$\varphi_{110} \to 0 \in \mathbb{R} \to \varphi_{234}$\lbl{A1}}
\BIC{$0 \in \mathbb{R} \to \varphi_{234}$}
\BIC{$\varphi_{234}$}
\BIC{$c \in \mathbb{R} \to \mathrm{Conv}(a,\alpha) \to \varepsilon \in \mathbb{R} \to 0 < \varepsilon \to \varphi_{66}$}
\end{ptree}
\noindent$\Pi_{109}$:
\begin{ptree}
\AXC{$y = z \to x = y \to x = z$\lbl{A6}}
\AXC{$(y = z \to x = y \to x = z) \to x = y \to y = z \to x = z$\lbl{A1}}
\BIC{$x = y \to y = z \to x = z$}
\UIC{$\forall z\,(x = y \to y = z \to x = z)$}
\UIC{$\forall y\,\forall z\,(x = y \to y = z \to x = z)$}
\UIC{$\forall x\,\forall y\,\forall z\,(x = y \to y = z \to x = z)$}
\AXC{$\varphi_{6}$\lbl{A2}}
\BIC{$\forall y\,\forall z\,(\lvert c \rvert = y \to y = z \to \lvert c \rvert = z)$}
\AXC{$\varphi_{44}$\lbl{A2}}
\BIC{$\forall z\,(\lvert c \rvert = \lvert 0 \rvert \to \lvert 0 \rvert = z \to \lvert c \rvert = z)$}
\AXC{$\varphi_{63}$\lbl{A2}}
\BIC{$\lvert c \rvert = \lvert 0 \rvert \to \lvert 0 \rvert = 0 \to \lvert c \rvert = 0$}
\end{ptree}
\noindent$\Pi_{110}$:
\begin{ptree}
\AXC{$\forall x\,(0 \leq x \to \lvert x \rvert = x)$\lbl{[A1]}}
\AXC{$\forall x\,(0 \leq x \to \lvert x \rvert = x) \to 0 \leq 0 \to \lvert 0 \rvert = 0$\lbl{A2}}
\BIC{$0 \leq 0 \to \lvert 0 \rvert = 0$}
\AXC{$\varphi_{222}$\lbl{A1}}
\BIC{$(c \in \mathbb{R} \to \varphi_{31}) \to c \in \mathbb{R} \to \mathrm{Conv}(a,\alpha) \to \varphi_{9}$}
\end{ptree}
\noindent$\Pi_{111}$:
\begin{ptree}
\AXC{$x = x$\lbl{A5}}
\UIC{$\forall x\,(x = x)$}
\AXC{$\forall x\,(x = x) \to 0 = 0$\lbl{A2}}
\BIC{$0 = 0$}
\AXC{$0 = 0 \to c \in \mathbb{R} \to \varphi_{31}$\lbl{A1}}
\BIC{$c \in \mathbb{R} \to \varphi_{31}$}
\AXC{$\stackrel{\Pi_{110}}{\vdots}$}\noLine
\UIC{$(c \in \mathbb{R} \to \varphi_{31}) \to c \in \mathbb{R} \to \mathrm{Conv}(a,\alpha) \to \varphi_{9}$}
\BIC{$c \in \mathbb{R} \to \mathrm{Conv}(a,\alpha) \to \varphi_{9}$}
\end{ptree}
\noindent$\Pi_{112}$:
\begin{ptree}
\AXC{$x = x$\lbl{A5}}
\UIC{$\forall x\,(x = x)$}
\AXC{$\forall x\,(x = x) \to \lvert c \rvert = \lvert c \rvert$\lbl{A2}}
\BIC{$\lvert c \rvert = \lvert c \rvert$}
\AXC{$\varphi_{8}$\lbl{A1}}
\AXC{$\varphi_{8} \to \lvert c \rvert = \lvert c \rvert \to \varphi_{168}$\lbl{A1}}
\BIC{$\lvert c \rvert = \lvert c \rvert \to \varphi_{168}$}
\BIC{$\varphi_{168}$}
\end{ptree}
\noindent$\Pi_{113}$:
\begin{ptree}
\AXC{$z = w \to \lvert c \rvert = \lvert z \rvert \to \lvert c \rvert = \lvert w \rvert$\lbl{A6}}
\UIC{$\forall w\,(z = w \to \lvert c \rvert = \lvert z \rvert \to \lvert c \rvert = \lvert w \rvert)$}
\UIC{$\forall z\,\forall w\,(z = w \to \lvert c \rvert = \lvert z \rvert \to \lvert c \rvert = \lvert w \rvert)$}
\AXC{$\varphi_{83}$\lbl{A2}}
\BIC{$\forall w\,(c = w \to \lvert c \rvert = \lvert c \rvert \to \lvert c \rvert = \lvert w \rvert)$}
\AXC{$\varphi_{64}$\lbl{A2}}
\BIC{$c = 0 \to \lvert c \rvert = \lvert c \rvert \to \lvert c \rvert = \lvert 0 \rvert$}
\AXC{$\stackrel{\Pi_{112}}{\vdots}$}\noLine
\UIC{$\varphi_{168}$}
\BIC{$c \in \mathbb{R} \to \mathrm{Conv}(a,\alpha) \to \varphi_{22}$}
\AXC{$\varphi_{254}$\lbl{A1}}
\BIC{$(c \in \mathbb{R} \to \mathrm{Conv}(a,\alpha) \to \varphi_{9}) \to \varphi_{156}$}
\end{ptree}
\noindent$\Pi_{114}$:
\begin{ptree}
\AXC{$\stackrel{\Pi_{109}}{\vdots}$}\noLine
\UIC{$\lvert c \rvert = \lvert 0 \rvert \to \lvert 0 \rvert = 0 \to \lvert c \rvert = 0$}
\AXC{$\stackrel{\Pi_{111}}{\vdots}$}\noLine
\UIC{$c \in \mathbb{R} \to \mathrm{Conv}(a,\alpha) \to \varphi_{9}$}
\AXC{$\stackrel{\Pi_{113}}{\vdots}$}\noLine
\UIC{$(c \in \mathbb{R} \to \mathrm{Conv}(a,\alpha) \to \varphi_{9}) \to \varphi_{156}$}
\BIC{$\varphi_{156}$}
\BIC{$c \in \mathbb{R} \to \mathrm{Conv}(a,\alpha) \to \varphi_{10}$}
\end{ptree}
\noindent$\Pi_{115}$:
\begin{ptree}
\AXC{$x = x$\lbl{A5}}
\UIC{$\forall x\,(x = x)$}
\AXC{$\forall x\,(x = x) \to \lvert c \rvert \cdot \lvert a_{n} - \alpha \rvert = \lvert c \rvert \cdot \lvert a_{n} - \alpha \rvert$\lbl{A2}}
\BIC{$\lvert c \rvert \cdot \lvert a_{n} - \alpha \rvert = \lvert c \rvert \cdot \lvert a_{n} - \alpha \rvert$}
\AXC{$\stackrel{\Pi_{114}}{\vdots}$}\noLine
\UIC{$c \in \mathbb{R} \to \mathrm{Conv}(a,\alpha) \to \varphi_{10}$}
\AXC{$\varphi_{290}$\lbl{A1}}
\BIC{$\lvert c \rvert \cdot \lvert a_{n} - \alpha \rvert = \lvert c \rvert \cdot \lvert a_{n} - \alpha \rvert \to \varphi_{242}$}
\BIC{$\varphi_{242}$}
\end{ptree}
\noindent$\Pi_{116}$:
\begin{ptree}
\AXC{$z = w \to \varphi_{133}$\lbl{A6}}
\UIC{$\forall w\,(z = w \to \varphi_{133})$}
\UIC{$\forall z\,\forall w\,(z = w \to \varphi_{133})$}
\AXC{$\forall z\,\forall w\,(z = w \to \varphi_{133}) \to \varphi_{167}$\lbl{A2}}
\BIC{$\varphi_{167}$}
\AXC{$\varphi_{167} \to \lvert c \rvert = 0 \to \varphi_{139}$\lbl{A2}}
\BIC{$\lvert c \rvert = 0 \to \varphi_{139}$}
\AXC{$\stackrel{\Pi_{115}}{\vdots}$}\noLine
\UIC{$\varphi_{242}$}
\BIC{$c \in \mathbb{R} \to \mathrm{Conv}(a,\alpha) \to \varepsilon \in \mathbb{R} \to 0 < \varepsilon \to c = 0 \to \varphi_{57}$}
\AXC{$\varphi_{316}$\lbl{A1}}
\BIC{$(c \in \mathbb{R} \to \mathrm{Conv}(a,\alpha) \to \varepsilon \in \mathbb{R} \to 0 < \varepsilon \to \varphi_{66}) \to \varphi_{249}$}
\end{ptree}
\noindent$\Pi_{117}$:
\begin{ptree}
\AXC{$\stackrel{\Pi_{108}}{\vdots}$}\noLine
\UIC{$c \in \mathbb{R} \to \mathrm{Conv}(a,\alpha) \to \varepsilon \in \mathbb{R} \to 0 < \varepsilon \to \varphi_{66}$}
\AXC{$\stackrel{\Pi_{116}}{\vdots}$}\noLine
\UIC{$(c \in \mathbb{R} \to \mathrm{Conv}(a,\alpha) \to \varepsilon \in \mathbb{R} \to 0 < \varepsilon \to \varphi_{66}) \to \varphi_{249}$}
\BIC{$\varphi_{249}$}
\end{ptree}
\noindent$\Pi_{118}$:
\begin{ptree}
\AXC{$\stackrel{\Pi_{106}}{\vdots}$}\noLine
\UIC{$\varphi_{68}$}
\AXC{$\varphi_{68} \to \varphi_{151}$\lbl{A2}}
\BIC{$\varphi_{151}$}
\AXC{$\varphi_{151} \to \lvert c \rvert \cdot \lvert a_{n} - \alpha \rvert = 0 \cdot \lvert a_{n} - \alpha \rvert \to \varphi_{128}$\lbl{A2}}
\BIC{$\lvert c \rvert \cdot \lvert a_{n} - \alpha \rvert = 0 \cdot \lvert a_{n} - \alpha \rvert \to \varphi_{128}$}
\AXC{$\stackrel{\Pi_{117}}{\vdots}$}\noLine
\UIC{$\varphi_{249}$}
\BIC{$c \in \mathbb{R} \to \mathrm{Conv}(a,\alpha) \to \varepsilon \in \mathbb{R} \to 0 < \varepsilon \to c = 0 \to \varphi_{58}$}
\AXC{$\varphi_{301}$\lbl{A1}}
\BIC{$\varphi_{262}$}
\end{ptree}
\noindent$\Pi_{119}$:
\begin{ptree}
\AXC{$y = z \to x = y \to x = z$\lbl{A6}}
\AXC{$(y = z \to x = y \to x = z) \to x = y \to y = z \to x = z$\lbl{A1}}
\BIC{$x = y \to y = z \to x = z$}
\UIC{$\forall z\,(x = y \to y = z \to x = z)$}
\UIC{$\forall y\,\forall z\,(x = y \to y = z \to x = z)$}
\UIC{$\forall x\,\forall y\,\forall z\,(x = y \to y = z \to x = z)$}
\AXC{$\forall x\,\forall y\,\forall z\,(x = y \to y = z \to x = z) \to \varphi_{86}$\lbl{A2}}
\BIC{$\varphi_{86}$}
\end{ptree}
\noindent$\Pi_{120}$:
\begin{ptree}
\AXC{$\stackrel{\Pi_{119}}{\vdots}$}\noLine
\UIC{$\varphi_{86}$}
\AXC{$\varphi_{86} \to \varphi_{162}$\lbl{A2}}
\BIC{$\varphi_{162}$}
\AXC{$\varphi_{162} \to \lvert c \odot a(n) - c \cdot \alpha \rvert = \lvert c \cdot (a_{n} - \alpha) \rvert \to \varphi_{147}$\lbl{A2}}
\BIC{$\lvert c \odot a(n) - c \cdot \alpha \rvert = \lvert c \cdot (a_{n} - \alpha) \rvert \to \varphi_{147}$}
\end{ptree}
\noindent$\Pi_{121}$:
\begin{ptree}
\AXC{$\forall x\,\forall y\,(x \in \mathbb{R} \to y \in \mathbb{R} \to \lvert x \cdot y \rvert = \lvert x \rvert \cdot \lvert y \rvert)$\lbl{[A17]}}
\AXC{$\varphi_{149}$\lbl{A2}}
\BIC{$\forall y\,(c \in \mathbb{R} \to y \in \mathbb{R} \to \lvert c \cdot y \rvert = \lvert c \rvert \cdot \lvert y \rvert)$}
\AXC{$\forall y\,(c \in \mathbb{R} \to y \in \mathbb{R} \to \lvert c \cdot y \rvert = \lvert c \rvert \cdot \lvert y \rvert) \to \varphi_{109}$\lbl{A2}}
\BIC{$\varphi_{109}$}
\end{ptree}
\noindent$\Pi_{122}$:
\begin{ptree}
\AXC{$y = z \to x = y \to x = z$\lbl{A6}}
\AXC{$(y = z \to x = y \to x = z) \to x = y \to y = z \to x = z$\lbl{A1}}
\BIC{$x = y \to y = z \to x = z$}
\UIC{$\forall z\,(x = y \to y = z \to x = z)$}
\UIC{$\forall y\,\forall z\,(x = y \to y = z \to x = z)$}
\UIC{$\forall x\,\forall y\,\forall z\,(x = y \to y = z \to x = z)$}
\AXC{$\forall x\,\forall y\,\forall z\,(x = y \to y = z \to x = z) \to \varphi_{86}$\lbl{A2}}
\BIC{$\varphi_{86}$}
\end{ptree}
\noindent$\Pi_{123}$:
\begin{ptree}
\AXC{$\stackrel{\Pi_{122}}{\vdots}$}\noLine
\UIC{$\varphi_{86}$}
\AXC{$\varphi_{86} \to \varphi_{181}$\lbl{A2}}
\BIC{$\varphi_{181}$}
\AXC{$\varphi_{181} \to \lvert c \odot a(n) - c \cdot \alpha \rvert = \lvert c \cdot a_{n} - c \cdot \alpha \rvert \to \varphi_{135}$\lbl{A2}}
\BIC{$\lvert c \odot a(n) - c \cdot \alpha \rvert = \lvert c \cdot a_{n} - c \cdot \alpha \rvert \to \varphi_{135}$}
\end{ptree}
\noindent$\Pi_{124}$:
\begin{ptree}
\AXC{$x = x$\lbl{A5}}
\UIC{$\forall x\,(x = x)$}
\AXC{$\forall x\,(x = x) \to \lvert c \cdot a_{n} - c \cdot \alpha \rvert = \lvert c \cdot a_{n} - c \cdot \alpha \rvert$\lbl{A2}}
\BIC{$\lvert c \cdot a_{n} - c \cdot \alpha \rvert = \lvert c \cdot a_{n} - c \cdot \alpha \rvert$}
\end{ptree}
\noindent$\Pi_{125}$:
\begin{ptree}
\AXC{$\varphi_{62}$\lbl{[F17]}}
\AXC{$\varphi_{62} \to \varphi_{92}$\lbl{A2}}
\BIC{$\varphi_{92}$}
\AXC{$\varphi_{92} \to \varphi_{107}$\lbl{A2}}
\BIC{$\varphi_{107}$}
\AXC{$\varphi_{107} \to v_{383} \in \mathbb{R} \to v_{384} \in \mathbb{R} \to \varphi_{18}$\lbl{A2}}
\BIC{$v_{383} \in \mathbb{R} \to v_{384} \in \mathbb{R} \to \varphi_{18}$}
\UIC{$\forall v_{385}\,(v_{383} \in \mathbb{R} \to v_{384} \in \mathbb{R} \to \varphi_{18})$}
\UIC{$\forall v_{384}\,\forall v_{385}\,(v_{383} \in \mathbb{R} \to v_{384} \in \mathbb{R} \to \varphi_{18})$}
\UIC{$\forall v_{383}\,\forall v_{384}\,\forall v_{385}\,(v_{383} \in \mathbb{R} \to v_{384} \in \mathbb{R} \to \varphi_{18})$}
\AXC{$\varphi_{229}$\lbl{A2}}
\BIC{$\varphi_{114}$}
\AXC{$\varphi_{114} \to \varphi_{122}$\lbl{A2}}
\BIC{$\varphi_{122}$}
\end{ptree}
\noindent$\Pi_{126}$:
\begin{ptree}
\AXC{$\stackrel{\Pi_{101}}{\vdots}$}\noLine
\UIC{$c \in \mathbb{R} \to \mathrm{Conv}(a,\alpha) \to \varphi_{35}$}
\AXC{$c \in \mathbb{R} \to \varphi_{67}$\lbl{A1}}
\AXC{$\stackrel{\Pi_{125}}{\vdots}$}\noLine
\UIC{$\varphi_{122}$}
\AXC{$\varphi_{122} \to \varphi_{106}$\lbl{A2}}
\BIC{$\varphi_{106}$}
\AXC{$\varphi_{296}$\lbl{A1}}
\BIC{$(c \in \mathbb{R} \to \varphi_{67}) \to (c \in \mathbb{R} \to \mathrm{Conv}(a,\alpha) \to \varphi_{35}) \to \varphi_{213}$}
\BIC{$(c \in \mathbb{R} \to \mathrm{Conv}(a,\alpha) \to \varphi_{35}) \to \varphi_{213}$}
\BIC{$\varphi_{213}$}
\end{ptree}
\noindent$\Pi_{127}$:
\begin{ptree}
\AXC{$\stackrel{\Pi_{124}}{\vdots}$}\noLine
\UIC{$\lvert c \cdot a_{n} - c \cdot \alpha \rvert = \lvert c \cdot a_{n} - c \cdot \alpha \rvert$}
\AXC{$\stackrel{\Pi_{19}}{\vdots}$}\noLine
\UIC{$c \in \mathbb{R} \to \mathrm{Conv}(a,\alpha) \to \alpha \in \mathbb{R}$}
\AXC{$\stackrel{\Pi_{126}}{\vdots}$}\noLine
\UIC{$\varphi_{213}$}
\BIC{$c \in \mathbb{R} \to \mathrm{Conv}(a,\alpha) \to \varepsilon \in \mathbb{R} \to 0 < \varepsilon \to \varphi_{41}$}
\AXC{$\varphi_{303}$\lbl{A1}}
\BIC{$\lvert c \cdot a_{n} - c \cdot \alpha \rvert = \lvert c \cdot a_{n} - c \cdot \alpha \rvert \to \varphi_{250}$}
\BIC{$\varphi_{250}$}
\end{ptree}
\noindent$\Pi_{128}$:
\begin{ptree}
\AXC{$z = w \to \varphi_{123}$\lbl{A6}}
\UIC{$\forall w\,(z = w \to \varphi_{123})$}
\UIC{$\forall z\,\forall w\,(z = w \to \varphi_{123})$}
\AXC{$\forall z\,\forall w\,(z = w \to \varphi_{123}) \to \varphi_{194}$\lbl{A2}}
\BIC{$\varphi_{194}$}
\AXC{$\varphi_{194} \to c \odot a(n) = c \cdot a_{n} \to \varphi_{169}$\lbl{A2}}
\BIC{$c \odot a(n) = c \cdot a_{n} \to \varphi_{169}$}
\end{ptree}
\noindent$\Pi_{129}$:
\begin{ptree}
\AXC{$c \in \mathbb{R} \to \varphi_{21}$\lbl{A1}}
\AXC{$\varphi_{101}$\lbl{A2}}
\AXC{$\stackrel{\Pi_{68}}{\vdots}$}\noLine
\UIC{$\varphi_{38}$}
\AXC{$\varphi_{38} \to \varphi_{101} \to \varphi_{28}$\lbl{A1}}
\BIC{$\varphi_{101} \to \varphi_{28}$}
\BIC{$\varphi_{28}$}
\AXC{$\varphi_{241}$\lbl{A1}}
\BIC{$(c \in \mathbb{R} \to \varphi_{21}) \to c \in \mathbb{R} \to \mathrm{Conv}(a,\alpha) \to \varepsilon \in \mathbb{R} \to \varphi_{29}$}
\BIC{$c \in \mathbb{R} \to \mathrm{Conv}(a,\alpha) \to \varepsilon \in \mathbb{R} \to \varphi_{29}$}
\end{ptree}
\noindent$\Pi_{130}$:
\begin{ptree}
\AXC{$\stackrel{\Pi_{129}}{\vdots}$}\noLine
\UIC{$c \in \mathbb{R} \to \mathrm{Conv}(a,\alpha) \to \varepsilon \in \mathbb{R} \to \varphi_{29}$}
\AXC{$(c \in \mathbb{R} \to \mathrm{Conv}(a,\alpha) \to \varepsilon \in \mathbb{R} \to \varphi_{29}) \to \varphi_{272}$\lbl{A1}}
\BIC{$\varphi_{272}$}
\end{ptree}
\noindent$\Pi_{131}$:
\begin{ptree}
\AXC{$\stackrel{\Pi_{128}}{\vdots}$}\noLine
\UIC{$c \odot a(n) = c \cdot a_{n} \to \varphi_{169}$}
\AXC{$x = x$\lbl{A5}}
\UIC{$\forall x\,(x = x)$}
\AXC{$\varphi_{49}$\lbl{A2}}
\BIC{$\lvert c \odot a(n) - c \cdot \alpha \rvert = \lvert c \odot a(n) - c \cdot \alpha \rvert$}
\AXC{$\stackrel{\Pi_{130}}{\vdots}$}\noLine
\UIC{$\varphi_{272}$}
\BIC{$(c \odot a(n) = c \cdot a_{n} \to \varphi_{169}) \to \varphi_{192}$}
\BIC{$\varphi_{192}$}
\AXC{$\varphi_{192} \to \varphi_{297}$\lbl{A1}}
\BIC{$\varphi_{297}$}
\end{ptree}
\noindent$\Pi_{132}$:
\begin{ptree}
\AXC{$\varphi_{77}$\lbl{A6}}
\UIC{$\forall w\,(\varphi_{77})$}
\UIC{$\forall z\,\forall w\,(\varphi_{77})$}
\AXC{$\forall z\,\forall w\,(\varphi_{77}) \to \varphi_{166}$\lbl{A2}}
\BIC{$\varphi_{166}$}
\AXC{$\varphi_{166} \to c \cdot a_{n} - c \cdot \alpha = c \cdot (a_{n} - \alpha) \to \varphi_{138}$\lbl{A2}}
\BIC{$c \cdot a_{n} - c \cdot \alpha = c \cdot (a_{n} - \alpha) \to \varphi_{138}$}
\AXC{$\stackrel{\Pi_{127}}{\vdots}$}\noLine
\UIC{$\varphi_{250}$}
\BIC{$c \in \mathbb{R} \to \mathrm{Conv}(a,\alpha) \to \varepsilon \in \mathbb{R} \to 0 < \varepsilon \to c = 0 \to \varphi_{60}$}
\AXC{$\stackrel{\Pi_{131}}{\vdots}$}\noLine
\UIC{$\varphi_{297}$}
\BIC{$(\lvert c \odot a(n) - c \cdot \alpha \rvert = \lvert c \cdot a_{n} - c \cdot \alpha \rvert \to \varphi_{135}) \to \varphi_{183}$}
\end{ptree}
\noindent$\Pi_{133}$:
\begin{ptree}
\AXC{$\stackrel{\Pi_{123}}{\vdots}$}\noLine
\UIC{$\lvert c \odot a(n) - c \cdot \alpha \rvert = \lvert c \cdot a_{n} - c \cdot \alpha \rvert \to \varphi_{135}$}
\AXC{$\stackrel{\Pi_{132}}{\vdots}$}\noLine
\UIC{$(\lvert c \odot a(n) - c \cdot \alpha \rvert = \lvert c \cdot a_{n} - c \cdot \alpha \rvert \to \varphi_{135}) \to \varphi_{183}$}
\BIC{$\varphi_{183}$}
\AXC{$\varphi_{183} \to \varphi_{299}$\lbl{A1}}
\BIC{$\varphi_{299}$}
\end{ptree}
\noindent$\Pi_{134}$:
\begin{ptree}
\AXC{$\stackrel{\Pi_{103}}{\vdots}$}\noLine
\UIC{$c \in \mathbb{R} \to \mathrm{Conv}(a,\alpha) \to \varphi_{65}$}
\AXC{$c \in \mathbb{R} \to \varphi_{67}$\lbl{A1}}
\AXC{$\stackrel{\Pi_{121}}{\vdots}$}\noLine
\UIC{$\varphi_{109}$}
\AXC{$\varphi_{109} \to (c \in \mathbb{R} \to \varphi_{67}) \to \varphi_{235}$\lbl{A1}}
\BIC{$(c \in \mathbb{R} \to \varphi_{67}) \to \varphi_{235}$}
\BIC{$\varphi_{235}$}
\BIC{$c \in \mathbb{R} \to \mathrm{Conv}(a,\alpha) \to \varepsilon \in \mathbb{R} \to 0 < \varepsilon \to c = 0 \to \varphi_{59}$}
\AXC{$\stackrel{\Pi_{133}}{\vdots}$}\noLine
\UIC{$\varphi_{299}$}
\BIC{$(\lvert c \odot a(n) - c \cdot \alpha \rvert = \lvert c \cdot (a_{n} - \alpha) \rvert \to \varphi_{147}) \to \varphi_{193}$}
\end{ptree}
\noindent$\Pi_{135}$:
\begin{ptree}
\AXC{$\stackrel{\Pi_{120}}{\vdots}$}\noLine
\UIC{$\lvert c \odot a(n) - c \cdot \alpha \rvert = \lvert c \cdot (a_{n} - \alpha) \rvert \to \varphi_{147}$}
\AXC{$\stackrel{\Pi_{134}}{\vdots}$}\noLine
\UIC{$(\lvert c \odot a(n) - c \cdot \alpha \rvert = \lvert c \cdot (a_{n} - \alpha) \rvert \to \varphi_{147}) \to \varphi_{193}$}
\BIC{$\varphi_{193}$}
\AXC{$\varphi_{193} \to \varphi_{268}$\lbl{A1}}
\BIC{$\varphi_{268}$}
\end{ptree}
\noindent$\Pi_{136}$:
\begin{ptree}
\AXC{$\stackrel{\Pi_{99}}{\vdots}$}\noLine
\UIC{$\varphi_{68}$}
\AXC{$\varphi_{68} \to \varphi_{150}$\lbl{A2}}
\BIC{$\varphi_{150}$}
\AXC{$\varphi_{150} \to \varphi_{140}$\lbl{A2}}
\BIC{$\varphi_{140}$}
\AXC{$\stackrel{\Pi_{105}}{\vdots}$}\noLine
\UIC{$c \in \mathbb{R} \to \mathrm{Conv}(a,\alpha) \to \varepsilon \in \mathbb{R} \to \varphi_{16}$}
\AXC{$\stackrel{\Pi_{118}}{\vdots}$}\noLine
\UIC{$\varphi_{262}$}
\BIC{$\varphi_{140} \to c \in \mathbb{R} \to \mathrm{Conv}(a,\alpha) \to \varepsilon \in \mathbb{R} \to \varphi_{32}$}
\BIC{$c \in \mathbb{R} \to \mathrm{Conv}(a,\alpha) \to \varepsilon \in \mathbb{R} \to \varphi_{32}$}
\AXC{$\stackrel{\Pi_{135}}{\vdots}$}\noLine
\UIC{$\varphi_{268}$}
\BIC{$\varphi_{176} \to c \in \mathbb{R} \to \mathrm{Conv}(a,\alpha) \to \varepsilon \in \mathbb{R} \to 0 < \varepsilon \to \varphi_{11}$}
\end{ptree}
\noindent$\Pi_{137}$:
\begin{ptree}
\AXC{$\stackrel{\Pi_{98}}{\vdots}$}\noLine
\UIC{$\varphi_{86}$}
\AXC{$\varphi_{86} \to \varphi_{182}$\lbl{A2}}
\BIC{$\varphi_{182}$}
\AXC{$\varphi_{182} \to \varphi_{176}$\lbl{A2}}
\BIC{$\varphi_{176}$}
\AXC{$\stackrel{\Pi_{136}}{\vdots}$}\noLine
\UIC{$\varphi_{176} \to c \in \mathbb{R} \to \mathrm{Conv}(a,\alpha) \to \varepsilon \in \mathbb{R} \to 0 < \varepsilon \to \varphi_{11}$}
\BIC{$c \in \mathbb{R} \to \mathrm{Conv}(a,\alpha) \to \varepsilon \in \mathbb{R} \to 0 < \varepsilon \to \varphi_{11}$}
\end{ptree}
\noindent$\Pi_{138}$:
\begin{ptree}
\AXC{$\stackrel{\Pi_{137}}{\vdots}$}\noLine
\UIC{$c \in \mathbb{R} \to \mathrm{Conv}(a,\alpha) \to \varepsilon \in \mathbb{R} \to 0 < \varepsilon \to \varphi_{11}$}
\AXC{$(c \in \mathbb{R} \to \mathrm{Conv}(a,\alpha) \to \varepsilon \in \mathbb{R} \to 0 < \varepsilon \to \varphi_{11}) \to \varphi_{221}$\lbl{A1}}
\BIC{$\varphi_{221}$}
\end{ptree}
\noindent$\Pi_{139}$:
\begin{ptree}
\AXC{$\stackrel{\Pi_{7}}{\vdots}$}\noLine
\UIC{$\forall x\,\forall y\,(x = y \to y = x)$}
\AXC{$\forall x\,\forall y\,(x = y \to y = x) \to \varphi_{48}$\lbl{A2}}
\BIC{$\varphi_{48}$}
\AXC{$\varphi_{197}$\lbl{A2}}
\BIC{$\lvert c \odot a(n) - c \cdot \alpha \rvert = 0 \to 0 = \lvert c \odot a(n) - c \cdot \alpha \rvert$}
\AXC{$\stackrel{\Pi_{138}}{\vdots}$}\noLine
\UIC{$\varphi_{221}$}
\BIC{$c \in \mathbb{R} \to \mathrm{Conv}(a,\alpha) \to \varepsilon \in \mathbb{R} \to 0 < \varepsilon \to \varphi_{12}$}
\end{ptree}
\noindent$\Pi_{140}$:
\begin{ptree}
\AXC{$\stackrel{\Pi_{139}}{\vdots}$}\noLine
\UIC{$c \in \mathbb{R} \to \mathrm{Conv}(a,\alpha) \to \varepsilon \in \mathbb{R} \to 0 < \varepsilon \to \varphi_{12}$}
\AXC{$(c \in \mathbb{R} \to \mathrm{Conv}(a,\alpha) \to \varepsilon \in \mathbb{R} \to 0 < \varepsilon \to \varphi_{12}) \to \varphi_{244}$\lbl{A1}}
\BIC{$\varphi_{244}$}
\end{ptree}
\noindent$\Pi_{141}$:
\begin{ptree}
\AXC{$\stackrel{\Pi_{97}}{\vdots}$}\noLine
\UIC{$\varphi_{98}$}
\AXC{$c \in \mathbb{R} \to \mathrm{Conv}(a,\alpha) \to \varepsilon \in \mathbb{R} \to 0 < \varepsilon \to 0 < \varepsilon$\lbl{A1}}
\AXC{$\stackrel{\Pi_{140}}{\vdots}$}\noLine
\UIC{$\varphi_{244}$}
\BIC{$\varphi_{98} \to c \in \mathbb{R} \to \mathrm{Conv}(a,\alpha) \to \varepsilon \in \mathbb{R} \to 0 < \varepsilon \to \varphi_{24}$}
\BIC{$c \in \mathbb{R} \to \mathrm{Conv}(a,\alpha) \to \varepsilon \in \mathbb{R} \to 0 < \varepsilon \to \varphi_{24}$}
\end{ptree}
\noindent$\Pi_{142}$:
\begin{ptree}
\AXC{$\stackrel{\Pi_{141}}{\vdots}$}\noLine
\UIC{$c \in \mathbb{R} \to \mathrm{Conv}(a,\alpha) \to \varepsilon \in \mathbb{R} \to 0 < \varepsilon \to \varphi_{24}$}
\AXC{$(c \in \mathbb{R} \to \mathrm{Conv}(a,\alpha) \to \varepsilon \in \mathbb{R} \to 0 < \varepsilon \to \varphi_{24}) \to \varphi_{144}$\lbl{A1}}
\BIC{$\varphi_{144}$}
\UIC{$\forall n\,(\varphi_{144})$}
\end{ptree}
\noindent$\Pi_{143}$:
\begin{ptree}
\AXC{$\forall n\,(\varphi_{144}) \to \varphi_{155}$\lbl{A3}}
\AXC{$\stackrel{\Pi_{142}}{\vdots}$}\noLine
\UIC{$\forall n\,(\varphi_{144})$}
\AXC{$\forall n\,(\varphi_{144}) \to (\forall n\,(\varphi_{144}) \to \varphi_{155}) \to \varphi_{236}$\lbl{A1}}
\BIC{$(\forall n\,(\varphi_{144}) \to \varphi_{155}) \to \varphi_{236}$}
\BIC{$\varphi_{236}$}
\end{ptree}
\noindent$\Pi_{144}$:
\begin{ptree}
\AXC{$c \in \mathbb{R} \land \mathrm{Conv}(a,\alpha) \land \varepsilon \in \mathbb{R} \land 0 < \varepsilon \land c = 0 \to \varphi_{7}$\lbl{A4}}
\AXC{$\stackrel{\Pi_{143}}{\vdots}$}\noLine
\UIC{$\varphi_{236}$}
\BIC{$c \in \mathbb{R} \to \mathrm{Conv}(a,\alpha) \to \varepsilon \in \mathbb{R} \to 0 < \varepsilon \to \varphi_{33}$}
\AXC{$1_{\omega} \in \mathbb{N}$\lbl{[N1]}}
\AXC{$1_{\omega} \in \mathbb{N} \to \varphi_{240}$\lbl{A1}}
\BIC{$\varphi_{240}$}
\BIC{$c \in \mathbb{R} \to \mathrm{Conv}(a,\alpha) \to \varepsilon \in \mathbb{R} \to 0 < \varepsilon \to c = 0 \to \varphi_{71}$}
\end{ptree}
\noindent$\Pi_{145}$:
\begin{ptree}
\AXC{$\forall N\,(\varphi_{45}) \to \neg \varphi_{71}$\lbl{A2}}
\AXC{$\stackrel{\Pi_{144}}{\vdots}$}\noLine
\UIC{$c \in \mathbb{R} \to \mathrm{Conv}(a,\alpha) \to \varepsilon \in \mathbb{R} \to 0 < \varepsilon \to c = 0 \to \varphi_{71}$}
\AXC{$\varphi_{294}$\lbl{A1}}
\BIC{$\varphi_{248}$}
\BIC{$c \in \mathbb{R} \to \mathrm{Conv}(a,\alpha) \to \varepsilon \in \mathbb{R} \to 0 < \varepsilon \to \varphi_{79}$}
\end{ptree}
\noindent$\Pi_{146}$:
\begin{ptree}
\AXC{$\stackrel{\Pi_{145}}{\vdots}$}\noLine
\UIC{$c \in \mathbb{R} \to \mathrm{Conv}(a,\alpha) \to \varepsilon \in \mathbb{R} \to 0 < \varepsilon \to \varphi_{79}$}
\AXC{$(c \in \mathbb{R} \to \mathrm{Conv}(a,\alpha) \to \varepsilon \in \mathbb{R} \to 0 < \varepsilon \to \varphi_{79}) \to \varphi_{243}$\lbl{A1}}
\BIC{$\varphi_{243}$}
\end{ptree}
\noindent$\Pi_{147}$:
\begin{ptree}
\AXC{$\varphi_{4}$\lbl{A4}}
\AXC{$\varphi_{245}$\lbl{A3}}
\AXC{$\stackrel{\Pi_{96}}{\vdots}$}\noLine
\UIC{$c \in \mathbb{R} \to \mathrm{Conv}(a,\alpha) \to \varepsilon \in \mathbb{R} \to 0 < \varepsilon \to \varphi_{85}$}
\AXC{$\stackrel{\Pi_{146}}{\vdots}$}\noLine
\UIC{$\varphi_{243}$}
\BIC{$c \in \mathbb{R} \to \mathrm{Conv}(a,\alpha) \to \varepsilon \in \mathbb{R} \to \varphi_{95}$}
\AXC{$\varphi_{239}$\lbl{A1}}
\BIC{$c \in \mathbb{R} \land \mathrm{Conv}(a,\alpha) \to \varepsilon \in \mathbb{R} \to \varphi_{95}$}
\UIC{$\forall \varepsilon\,(c \in \mathbb{R} \land \mathrm{Conv}(a,\alpha) \to \varepsilon \in \mathbb{R} \to \varphi_{95})$}
\AXC{$\varphi_{320}$\lbl{A1}}
\BIC{$\varphi_{245} \to \varphi_{4} \to c \in \mathbb{R} \to \mathrm{Conv}(a,\alpha) \to \forall \varepsilon \in \mathbb{R}\,(\varphi_{95})$}
\BIC{$\varphi_{4} \to c \in \mathbb{R} \to \mathrm{Conv}(a,\alpha) \to \forall \varepsilon \in \mathbb{R}\,(\varphi_{95})$}
\BIC{$c \in \mathbb{R} \to \mathrm{Conv}(a,\alpha) \to \forall \varepsilon \in \mathbb{R}\,(\varphi_{95})$}
\end{ptree}
\noindent$\Pi_{148}$:
\begin{ptree}
\AXC{$c \in \mathbb{R} \to \mathrm{Conv}(a,\alpha) \to c \in \mathbb{R}$\lbl{A1}}
\AXC{$\forall x\,\forall y\,(x \in \mathbb{R} \to y \in \mathbb{R} \to x \cdot y \in \mathbb{R})$\lbl{(R1)}}
\AXC{$\varphi_{111}$\lbl{A2}}
\BIC{$\forall y\,(c \in \mathbb{R} \to y \in \mathbb{R} \to c \cdot y \in \mathbb{R})$}
\AXC{$\varphi_{103}$\lbl{A2}}
\BIC{$c \in \mathbb{R} \to \alpha \in \mathbb{R} \to c \cdot \alpha \in \mathbb{R}$}
\AXC{$\varphi_{217}$\lbl{A1}}
\BIC{$(c \in \mathbb{R} \to \mathrm{Conv}(a,\alpha) \to c \in \mathbb{R}) \to \varphi_{84}$}
\BIC{$\varphi_{84}$}
\end{ptree}
\noindent$\Pi_{149}$:
\begin{ptree}
\AXC{$\stackrel{\Pi_{67}}{\vdots}$}\noLine
\UIC{$c \in \mathbb{R} \to \varphi_{54}$}
\AXC{$\varphi_{170}$\lbl{A1}}
\BIC{$c \in \mathbb{R} \to \mathrm{Conv}(a,\alpha) \to \mathrm{Seq}(c \odot a)$}
\AXC{$\varphi_{184}$\lbl{DefConv}}
\AXC{$\varphi_{184} \to \varphi_{208}$\lbl{A2}}
\BIC{$\varphi_{208}$}
\AXC{$\varphi_{208} \to \varphi_{218}$\lbl{A2}}
\BIC{$\varphi_{218}$}
\AXC{$\varphi_{218} \to \varphi_{246}$\lbl{A1}}
\BIC{$\varphi_{246}$}
\BIC{$(c \in \mathbb{R} \to \mathrm{Conv}(a,\alpha) \to c \cdot \alpha \in \mathbb{R}) \to \varphi_{220}$}
\end{ptree}
\noindent$\Pi_{150}$:
\begin{ptree}
\AXC{$\stackrel{\Pi_{19}}{\vdots}$}\noLine
\UIC{$c \in \mathbb{R} \to \mathrm{Conv}(a,\alpha) \to \alpha \in \mathbb{R}$}
\AXC{$\stackrel{\Pi_{148}}{\vdots}$}\noLine
\UIC{$\varphi_{84}$}
\BIC{$c \in \mathbb{R} \to \mathrm{Conv}(a,\alpha) \to c \cdot \alpha \in \mathbb{R}$}
\AXC{$\stackrel{\Pi_{149}}{\vdots}$}\noLine
\UIC{$(c \in \mathbb{R} \to \mathrm{Conv}(a,\alpha) \to c \cdot \alpha \in \mathbb{R}) \to \varphi_{220}$}
\BIC{$\varphi_{220}$}
\end{ptree}
\noindent 主木:
\begin{ptree}
\AXC{$\stackrel{\Pi_{147}}{\vdots}$}\noLine
\UIC{$c \in \mathbb{R} \to \mathrm{Conv}(a,\alpha) \to \forall \varepsilon \in \mathbb{R}\,(\varphi_{95})$}
\AXC{$\stackrel{\Pi_{150}}{\vdots}$}\noLine
\UIC{$\varphi_{220}$}
\BIC{$c \in \mathbb{R} \to \mathrm{Conv}(a,\alpha) \to \mathrm{Conv}(c \odot a,c \cdot \alpha)$}
\UIC{$\varphi_{3}$}
\UIC{$\forall a\,(\varphi_{3})$}
\UIC{$\forall c\,\forall a\,(\varphi_{3})$}
\end{ptree}
\begin{abbrevlist}
\item $\varphi_{1} :\equiv \forall \alpha\,(\mathrm{Conv}(a,\alpha) \to \mathrm{Seq}(a)) \to \mathrm{Conv}(a,\alpha) \to \mathrm{Seq}(a)$
\item $\varphi_{2} :\equiv \forall a\,\forall \alpha\,(\mathrm{Conv}(a,\alpha) \to \mathrm{Seq}(a)) \to \forall \alpha\,(\mathrm{Conv}(a,\alpha) \to \mathrm{Seq}(a))$
\item $\varphi_{3} :\equiv \forall \alpha\,(c \in \mathbb{R} \to \mathrm{Conv}(a,\alpha) \to \mathrm{Conv}(c \odot a,c \cdot \alpha))$
\item $\varphi_{4} :\equiv c \in \mathbb{R} \land \mathrm{Conv}(a,\alpha) \to \forall \varepsilon\,(c \in \mathbb{R} \land \mathrm{Conv}(a,\alpha))$
\item $\varphi_{5} :\equiv \forall \alpha\,(\mathrm{Conv}(a,\alpha) \to \alpha \in \mathbb{R}) \to \mathrm{Conv}(a,\alpha) \to \alpha \in \mathbb{R}$
\item $\varphi_{6} :\equiv \forall x\,\forall y\,\forall z\,(x = y \to y = z \to x = z) \to \forall y\,\forall z\,(\lvert c \rvert = y \to y = z \to \lvert c \rvert = z)$
\item $\varphi_{7} :\equiv \forall n\,(c \in \mathbb{R} \land \mathrm{Conv}(a,\alpha) \land \varepsilon \in \mathbb{R} \land 0 < \varepsilon \land c = 0)$
\item $\varphi_{8} :\equiv c \in \mathbb{R} \to \mathrm{Conv}(a,\alpha) \to \varepsilon \in \mathbb{R} \to 0 < \varepsilon \to c = 0 \to c = 0$
\item $\varphi_{9} :\equiv \varepsilon \in \mathbb{R} \to 0 < \varepsilon \to c = 0 \to n \in \mathbb{N} \to 1_{\omega} \leq n \to \lvert 0 \rvert = 0$
\item $\varphi_{10} :\equiv \varepsilon \in \mathbb{R} \to 0 < \varepsilon \to c = 0 \to n \in \mathbb{N} \to 1_{\omega} \leq n \to \lvert c \rvert = 0$
\item $\varphi_{11} :\equiv c = 0 \to n \in \mathbb{N} \to 1_{\omega} \leq n \to \lvert c \odot a(n) - c \cdot \alpha \rvert = 0$
\item $\varphi_{12} :\equiv c = 0 \to n \in \mathbb{N} \to 1_{\omega} \leq n \to 0 = \lvert c \odot a(n) - c \cdot \alpha \rvert$
\item $\varphi_{13} :\equiv \forall a\,\forall \alpha\,(\mathrm{Conv}(a,\alpha) \to \alpha \in \mathbb{R}) \to \forall \alpha\,(\mathrm{Conv}(a,\alpha) \to \alpha \in \mathbb{R})$
\item $\varphi_{14} :\equiv \forall z\,\forall w\,(z = w \to z < \varepsilon \to w < \varepsilon) \to \forall w\,(0 = w \to 0 < \varepsilon \to w < \varepsilon)$
\item $\varphi_{15} :\equiv c \in \mathbb{R} \to \mathrm{Conv}(a,\alpha) \to \varepsilon \in \mathbb{R} \to 0 < \varepsilon \to c \neq 0 \to c \neq 0$
\item $\varphi_{16} :\equiv 0 < \varepsilon \to c = 0 \to n \in \mathbb{N} \to 1_{\omega} \leq n \to \lvert a_{n} - \alpha \rvert \cdot 0 = 0$
\item $\varphi_{17} :\equiv v_{388} \in \mathbb{R} \to v_{386} \cdot v_{387} - v_{386} \cdot v_{388} = v_{386} \cdot (v_{387} - v_{388})$
\item $\varphi_{18} :\equiv v_{385} \in \mathbb{R} \to v_{383} \cdot v_{384} - v_{383} \cdot v_{385} = v_{383} \cdot (v_{384} - v_{385})$
\item $\varphi_{19} :\equiv \forall x\,(x = x) \to \varepsilon \cdot (\lvert c \rvert \cdot \lvert c \rvert^{-1}) = \varepsilon \cdot (\lvert c \rvert \cdot \lvert c \rvert^{-1})$
\item $\varphi_{20} :\equiv \forall x\,(x = x) \to \varepsilon \cdot (\lvert c \rvert^{-1} \cdot \lvert c \rvert) = \varepsilon \cdot (\lvert c \rvert^{-1} \cdot \lvert c \rvert)$
\item $\varphi_{21} :\equiv \mathrm{Conv}(a,\alpha) \to \varepsilon \in \mathbb{R} \to 0 < \varepsilon \to c = 0 \to n \in \mathbb{N} \to n \in \mathbb{N}$
\item $\varphi_{22} :\equiv \varepsilon \in \mathbb{R} \to 0 < \varepsilon \to c = 0 \to n \in \mathbb{N} \to 1_{\omega} \leq n \to \lvert c \rvert = \lvert 0 \rvert$
\item $\varphi_{23} :\equiv 1_{\omega} \leq n \to \lvert c \odot a(n) - c \cdot \alpha \rvert = \lvert c \cdot (a_{n} - \alpha) \rvert$
\item $\varphi_{24} :\equiv c = 0 \to n \in \mathbb{N} \to 1_{\omega} \leq n \to \lvert c \odot a(n) - c \cdot \alpha \rvert < \varepsilon$
\item $\varphi_{25} :\equiv \forall x \in \mathbb{R}\,(x \neq 0 \to 0 < \lvert x \rvert) \to c \in \mathbb{R} \to c \neq 0 \to 0 < \lvert c \rvert$
\item $\varphi_{26} :\equiv c \in \mathbb{R} \to \mathrm{Conv}(a,\alpha) \to \varepsilon \in \mathbb{R} \to 0 < \varepsilon \to c \neq 0 \to 0 \neq \lvert c \rvert$
\item $\varphi_{27} :\equiv c \in \mathbb{R} \to \mathrm{Conv}(a,\alpha) \to \varepsilon \in \mathbb{R} \to 0 < \varepsilon \to c \neq 0 \to \lvert c \rvert \neq 0$
\item $\varphi_{28} :\equiv c \in \mathbb{R} \to \mathrm{Conv}(a,\alpha) \to n \in \mathbb{N} \to c \odot a(n) = c \cdot a_{n}$
\item $\varphi_{29} :\equiv 0 < \varepsilon \to c = 0 \to n \in \mathbb{N} \to 1_{\omega} \leq n \to c \odot a(n) = c \cdot a_{n}$
\item $\varphi_{30} :\equiv c \in \mathbb{R} \to \mathrm{Conv}(a,\alpha) \to \varepsilon \in \mathbb{R} \to 0 < \varepsilon \to c \neq 0 \to c \in \mathbb{R}$
\item $\varphi_{31} :\equiv \mathrm{Conv}(a,\alpha) \to \varepsilon \in \mathbb{R} \to 0 < \varepsilon \to c = 0 \to n \in \mathbb{N} \to 1_{\omega} \leq n \to 0 \leq 0$
\item $\varphi_{32} :\equiv 0 < \varepsilon \to c = 0 \to n \in \mathbb{N} \to 1_{\omega} \leq n \to \lvert c \rvert \cdot \lvert a_{n} - \alpha \rvert = 0$
\item $\varphi_{33} :\equiv c = 0 \to \forall n \in \mathbb{N}\,(1_{\omega} \leq n \to \lvert c \odot a(n) - c \cdot \alpha \rvert < \varepsilon)$
\item $\varphi_{34} :\equiv c \in \mathbb{R} \to \mathrm{Conv}(a,\alpha) \to \varepsilon \in \mathbb{R} \to 0 < \varepsilon \to c \neq 0 \to 0 < \lvert c \rvert$
\item $\varphi_{35} :\equiv \varepsilon \in \mathbb{R} \to 0 < \varepsilon \to c = 0 \to n \in \mathbb{N} \to 1_{\omega} \leq n \to a_{n} \in \mathbb{R}$
\item $\varphi_{36} :\equiv (c \in \mathbb{R} \to \mathrm{Conv}(a,\alpha) \to \mathrm{Conv}(a,\alpha)) \to c \in \mathbb{R} \to \mathrm{Conv}(a,\alpha) \to \mathrm{Seq}(a)$
\item $\varphi_{37} :\equiv \forall x\,\forall y\,\forall z\,(x \in \mathbb{R} \to y \in \mathbb{R} \to z \in \mathbb{R} \to x \cdot y \cdot z = x \cdot (y \cdot z))$
\item $\varphi_{38} :\equiv c \in \mathbb{R} \to \mathrm{Conv}(a,\alpha) \to \forall n \in \mathbb{N}\,(c \odot a(n) = c \cdot a_{n})$
\item $\varphi_{39} :\equiv c \in \mathbb{R} \to \mathrm{Conv}(a,\alpha) \to \varepsilon \in \mathbb{R} \to 0 < \varepsilon \to c \neq 0 \to \varepsilon \in \mathbb{R}$
\item $\varphi_{40} :\equiv \mathrm{Seq}(a) \to \mathrm{Seq}(c \odot a) \land \forall n \in \mathbb{N}\,(c \odot a(n) = c \cdot a_{n})$
\item $\varphi_{41} :\equiv c = 0 \to n \in \mathbb{N} \to 1_{\omega} \leq n \to c \cdot a_{n} - c \cdot \alpha = c \cdot (a_{n} - \alpha)$
\item $\varphi_{42} :\equiv 1_{\omega} \leq n \to \lvert c \odot a(n) - c \cdot \alpha \rvert = \lvert c \cdot a_{n} - c \cdot \alpha \rvert$
\item $\varphi_{43} :\equiv 1_{\omega} \leq n \to \lvert c \odot a(n) - c \cdot \alpha \rvert = \lvert c \rvert \cdot \lvert a_{n} - \alpha \rvert$
\item $\varphi_{44} :\equiv \forall y\,\forall z\,(\lvert c \rvert = y \to y = z \to \lvert c \rvert = z) \to \forall z\,(\lvert c \rvert = \lvert 0 \rvert \to \lvert 0 \rvert = z \to \lvert c \rvert = z)$
\item $\varphi_{45} :\equiv \neg (N \in \mathbb{N} \land \forall n \in \mathbb{N}\,(N \leq n \to \lvert c \odot a(n) - c \cdot \alpha \rvert < \varepsilon))$
\item $\varphi_{46} :\equiv \mathrm{Conv}(a,\alpha) \to \varepsilon \in \mathbb{R} \to 0 < \varepsilon \to c \neq 0 \to \varepsilon \cdot \lvert c \rvert^{-1} \in \mathbb{R}$
\item $\varphi_{47} :\equiv \forall y\,\forall z\,(\varepsilon \in \mathbb{R} \to y \in \mathbb{R} \to z \in \mathbb{R} \to \varepsilon \cdot y \cdot z = \varepsilon \cdot (y \cdot z))$
\item $\varphi_{48} :\equiv \forall y\,(\lvert c \odot a(n) - c \cdot \alpha \rvert = y \to y = \lvert c \odot a(n) - c \cdot \alpha \rvert)$
\item $\varphi_{49} :\equiv \forall x\,(x = x) \to \lvert c \odot a(n) - c \cdot \alpha \rvert = \lvert c \odot a(n) - c \cdot \alpha \rvert$
\item $\varphi_{50} :\equiv c \in \mathbb{R} \to \mathrm{Conv}(a,\alpha) \to \varepsilon \in \mathbb{R} \to 0 < \varepsilon \to c \neq 0 \to \lvert c \rvert \in \mathbb{R}$
\item $\varphi_{51} :\equiv \forall y\,\forall z\,(\varepsilon \cdot \lvert c \rvert^{-1} \cdot \lvert c \rvert = y \to y = z \to \varepsilon \cdot \lvert c \rvert^{-1} \cdot \lvert c \rvert = z)$
\item $\varphi_{52} :\equiv c \in \mathbb{R} \to \mathrm{Conv}(a,\alpha) \to \varepsilon \in \mathbb{R} \to 0 < \varepsilon \to c \neq 0 \to 0 < \lvert c \rvert^{-1}$
\item $\varphi_{53} :\equiv n \in \mathbb{N} \to N \leq n \to \varepsilon \cdot (\lvert c \rvert^{-1} \cdot \lvert c \rvert) = \varepsilon \cdot (\lvert c \rvert \cdot \lvert c \rvert^{-1})$
\item $\varphi_{54} :\equiv \mathrm{Conv}(a,\alpha) \to \mathrm{Seq}(c \odot a) \land \forall n \in \mathbb{N}\,(c \odot a(n) = c \cdot a_{n})$
\item $\varphi_{55} :\equiv n \in \mathbb{N} \to N \leq n \to \lvert c \odot a(n) - c \cdot \alpha \rvert = \lvert c \cdot (a_{n} - \alpha) \rvert$
\item $\varphi_{56} :\equiv \neg (N \in \mathbb{N} \land \forall n \in \mathbb{N}\,(N \leq n \to \lvert a_{n} - \alpha \rvert < \varepsilon \cdot \lvert c \rvert^{-1}))$
\item $\varphi_{57} :\equiv n \in \mathbb{N} \to 1_{\omega} \leq n \to \lvert c \rvert \cdot \lvert a_{n} - \alpha \rvert = 0 \cdot \lvert a_{n} - \alpha \rvert$
\item $\varphi_{58} :\equiv n \in \mathbb{N} \to 1_{\omega} \leq n \to \lvert c \rvert \cdot \lvert a_{n} - \alpha \rvert = \lvert a_{n} - \alpha \rvert \cdot 0$
\item $\varphi_{59} :\equiv n \in \mathbb{N} \to 1_{\omega} \leq n \to \lvert c \cdot (a_{n} - \alpha) \rvert = \lvert c \rvert \cdot \lvert a_{n} - \alpha \rvert$
\item $\varphi_{60} :\equiv n \in \mathbb{N} \to 1_{\omega} \leq n \to \lvert c \cdot a_{n} - c \cdot \alpha \rvert = \lvert c \cdot (a_{n} - \alpha) \rvert$
\item $\varphi_{61} :\equiv (c \in \mathbb{R} \to \mathrm{Conv}(a,\alpha) \to \mathrm{Conv}(a,\alpha)) \to c \in \mathbb{R} \to \mathrm{Conv}(a,\alpha) \to \alpha \in \mathbb{R}$
\item $\varphi_{62} :\equiv \forall c\,\forall x\,\forall a\,(c \in \mathbb{R} \to x \in \mathbb{R} \to a \in \mathbb{R} \to c \cdot x - c \cdot a = c \cdot (x - a))$
\item $\varphi_{63} :\equiv \forall z\,(\lvert c \rvert = \lvert 0 \rvert \to \lvert 0 \rvert = z \to \lvert c \rvert = z) \to \lvert c \rvert = \lvert 0 \rvert \to \lvert 0 \rvert = 0 \to \lvert c \rvert = 0$
\item $\varphi_{64} :\equiv \forall w\,(c = w \to \lvert c \rvert = \lvert c \rvert \to \lvert c \rvert = \lvert w \rvert) \to c = 0 \to \lvert c \rvert = \lvert c \rvert \to \lvert c \rvert = \lvert 0 \rvert$
\item $\varphi_{65} :\equiv \varepsilon \in \mathbb{R} \to 0 < \varepsilon \to c = 0 \to n \in \mathbb{N} \to 1_{\omega} \leq n \to a_{n} - \alpha \in \mathbb{R}$
\item $\varphi_{66} :\equiv c = 0 \to n \in \mathbb{N} \to 1_{\omega} \leq n \to 0 \cdot \lvert a_{n} - \alpha \rvert = \lvert a_{n} - \alpha \rvert \cdot 0$
\item $\varphi_{67} :\equiv \mathrm{Conv}(a,\alpha) \to \varepsilon \in \mathbb{R} \to 0 < \varepsilon \to c = 0 \to n \in \mathbb{N} \to 1_{\omega} \leq n \to c \in \mathbb{R}$
\item $\varphi_{68} :\equiv \forall y\,\forall z\,(\lvert c \rvert \cdot \lvert a_{n} - \alpha \rvert = y \to y = z \to \lvert c \rvert \cdot \lvert a_{n} - \alpha \rvert = z)$
\item $\varphi_{69} :\equiv \varepsilon \cdot (\lvert c \rvert \cdot \lvert c \rvert^{-1}) = \varepsilon \cdot z \to \varepsilon \cdot (\lvert c \rvert \cdot \lvert c \rvert^{-1}) = \varepsilon \cdot w$
\item $\varphi_{70} :\equiv \varepsilon \cdot (\lvert c \rvert^{-1} \cdot \lvert c \rvert) = \varepsilon \cdot z \to \varepsilon \cdot (\lvert c \rvert^{-1} \cdot \lvert c \rvert) = \varepsilon \cdot w$
\item $\varphi_{71} :\equiv 1_{\omega} \in \mathbb{N} \land \forall n \in \mathbb{N}\,(1_{\omega} \leq n \to \lvert c \odot a(n) - c \cdot \alpha \rvert < \varepsilon)$
\item $\varphi_{72} :\equiv c \in \mathbb{R} \to \mathrm{Conv}(a,\alpha) \to \varepsilon \in \mathbb{R} \to 0 < \varepsilon \to c \neq 0 \to \lvert c \rvert^{-1} \in \mathbb{R}$
\item $\varphi_{73} :\equiv \forall w\,(z = w \to \lvert a_{n} - \alpha \rvert \cdot \lvert c \rvert < z \to \lvert a_{n} - \alpha \rvert \cdot \lvert c \rvert < w)$
\item $\varphi_{74} :\equiv n \in \mathbb{N} \to N \leq n \to \lvert c \odot a(n) - c \cdot \alpha \rvert = \lvert c \cdot a_{n} - c \cdot \alpha \rvert$
\item $\varphi_{75} :\equiv n \in \mathbb{N} \to N \leq n \to \lvert c \odot a(n) - c \cdot \alpha \rvert = \lvert c \rvert \cdot \lvert a_{n} - \alpha \rvert$
\item $\varphi_{76} :\equiv n \in \mathbb{N} \to N \leq n \to \lvert c \rvert \cdot \lvert a_{n} - \alpha \rvert = \lvert c \odot a(n) - c \cdot \alpha \rvert$
\item $\varphi_{77} :\equiv z = w \to \lvert c \cdot a_{n} - c \cdot \alpha \rvert = \lvert z \rvert \to \lvert c \cdot a_{n} - c \cdot \alpha \rvert = \lvert w \rvert$
\item $\varphi_{78} :\equiv \varepsilon \in \mathbb{R} \to 0 < \varepsilon \to c = 0 \to n \in \mathbb{N} \to 1_{\omega} \leq n \to \lvert a_{n} - \alpha \rvert \in \mathbb{R}$
\item $\varphi_{79} :\equiv c = 0 \to \exists N \in \mathbb{N}\,\forall n \in \mathbb{N}\,(N \leq n \to \lvert c \odot a(n) - c \cdot \alpha \rvert < \varepsilon)$
\item $\varphi_{80} :\equiv \forall n\,(\mathrm{Seq}(a) \to n \in \mathbb{N} \to a_{n} \in \mathbb{R}) \to \mathrm{Seq}(a) \to n \in \mathbb{N} \to a_{n} \in \mathbb{R}$
\item $\varphi_{81} :\equiv \forall x\,\forall y\,\forall z\,(x \in \mathbb{R} \to y \in \mathbb{R} \to z \in \mathbb{R} \to 0 < z \to x < y \to x \cdot z < y \cdot z)$
\item $\varphi_{82} :\equiv c \in \mathbb{R} \to \mathrm{Conv}(a,\alpha) \to \varepsilon \in \mathbb{R} \to 0 < \varepsilon \to c \neq 0 \to 0 < \varepsilon \cdot \lvert c \rvert^{-1}$
\item $\varphi_{83} :\equiv \forall z\,\forall w\,(z = w \to \lvert c \rvert = \lvert z \rvert \to \lvert c \rvert = \lvert w \rvert) \to \forall w\,(c = w \to \lvert c \rvert = \lvert c \rvert \to \lvert c \rvert = \lvert w \rvert)$
\item $\varphi_{84} :\equiv (c \in \mathbb{R} \to \mathrm{Conv}(a,\alpha) \to \alpha \in \mathbb{R}) \to c \in \mathbb{R} \to \mathrm{Conv}(a,\alpha) \to c \cdot \alpha \in \mathbb{R}$
\item $\varphi_{85} :\equiv c \neq 0 \to \exists N \in \mathbb{N}\,\forall n \in \mathbb{N}\,(N \leq n \to \lvert c \odot a(n) - c \cdot \alpha \rvert < \varepsilon)$
\item $\varphi_{86} :\equiv \forall y\,\forall z\,(\lvert c \odot a(n) - c \cdot \alpha \rvert = y \to y = z \to \lvert c \odot a(n) - c \cdot \alpha \rvert = z)$
\item $\varphi_{87} :\equiv \forall w\,(\lvert c \rvert \cdot \lvert a_{n} - \alpha \rvert = w \to \lvert c \rvert \cdot \lvert a_{n} - \alpha \rvert < \varepsilon \to w < \varepsilon)$
\item $\varphi_{88} :\equiv \forall w\,(\lvert a_{n} - \alpha \rvert \cdot \lvert c \rvert = w \to \lvert a_{n} - \alpha \rvert \cdot \lvert c \rvert < \varepsilon \to w < \varepsilon)$
\item $\varphi_{89} :\equiv \forall x \in \mathbb{R}\,(x \cdot 0 = 0) \to \lvert a_{n} - \alpha \rvert \in \mathbb{R} \to \lvert a_{n} - \alpha \rvert \cdot 0 = 0$
\item $\varphi_{90} :\equiv c \neq 0 \to \exists N \in \mathbb{N}\,\forall n \in \mathbb{N}\,(N \leq n \to \lvert a_{n} - \alpha \rvert < \varepsilon \cdot \lvert c \rvert^{-1})$
\item $\varphi_{91} :\equiv \forall x\,\forall a\,(v_{386} \in \mathbb{R} \to x \in \mathbb{R} \to a \in \mathbb{R} \to v_{386} \cdot x - v_{386} \cdot a = v_{386} \cdot (x - a))$
\item $\varphi_{92} :\equiv \forall x\,\forall a\,(v_{383} \in \mathbb{R} \to x \in \mathbb{R} \to a \in \mathbb{R} \to v_{383} \cdot x - v_{383} \cdot a = v_{383} \cdot (x - a))$
\item $\varphi_{93} :\equiv \forall z\,(\varepsilon \cdot \lvert c \rvert^{-1} \cdot \lvert c \rvert = \varepsilon \cdot 1 \to \varepsilon \cdot 1 = z \to \varepsilon \cdot \lvert c \rvert^{-1} \cdot \lvert c \rvert = z)$
\item $\varphi_{94} :\equiv N \in \mathbb{N} \land \forall n \in \mathbb{N}\,(N \leq n \to \lvert a_{n} - \alpha \rvert < \varepsilon \cdot \lvert c \rvert^{-1}) \to N \in \mathbb{N}$
\item $\varphi_{95} :\equiv 0 < \varepsilon \to \exists N \in \mathbb{N}\,\forall n \in \mathbb{N}\,(N \leq n \to \lvert c \odot a(n) - c \cdot \alpha \rvert < \varepsilon)$
\item $\varphi_{96} :\equiv \varepsilon \cdot \lvert c \rvert^{-1} \cdot \lvert c \rvert = \varepsilon \cdot 1 \to \varepsilon \cdot 1 = \varepsilon \to \varepsilon \cdot \lvert c \rvert^{-1} \cdot \lvert c \rvert = \varepsilon$
\item $\varphi_{97} :\equiv \forall x \in \mathbb{R}\,(0 < x \to 0 < x^{-1}) \to \lvert c \rvert \in \mathbb{R} \to 0 < \lvert c \rvert \to 0 < \lvert c \rvert^{-1}$
\item $\varphi_{98} :\equiv 0 = \lvert c \odot a(n) - c \cdot \alpha \rvert \to 0 < \varepsilon \to \lvert c \odot a(n) - c \cdot \alpha \rvert < \varepsilon$
\item $\varphi_{99} :\equiv \forall a\,\forall n\,(\mathrm{Seq}(a) \to n \in \mathbb{N} \to a_{n} \in \mathbb{R}) \to \forall n\,(\mathrm{Seq}(a) \to n \in \mathbb{N} \to a_{n} \in \mathbb{R})$
\item $\varphi_{100} :\equiv \forall x \in \mathbb{R}\,(\lvert x \rvert \in \mathbb{R}) \to a_{n} - \alpha \in \mathbb{R} \to \lvert a_{n} - \alpha \rvert \in \mathbb{R}$
\item $\varphi_{101} :\equiv \forall n \in \mathbb{N}\,(c \odot a(n) = c \cdot a_{n}) \to n \in \mathbb{N} \to c \odot a(n) = c \cdot a_{n}$
\item $\varphi_{102} :\equiv \forall x\,\forall y\,(x \in \mathbb{R} \to y \in \mathbb{R} \to x \cdot y = y \cdot x) \to \forall y\,(0 \in \mathbb{R} \to y \in \mathbb{R} \to 0 \cdot y = y \cdot 0)$
\item $\varphi_{103} :\equiv \forall y\,(c \in \mathbb{R} \to y \in \mathbb{R} \to c \cdot y \in \mathbb{R}) \to c \in \mathbb{R} \to \alpha \in \mathbb{R} \to c \cdot \alpha \in \mathbb{R}$
\item $\varphi_{104} :\equiv \forall x \in \mathbb{R}\,(x \neq 0 \to x^{-1} \in \mathbb{R}) \to \lvert c \rvert \in \mathbb{R} \to \lvert c \rvert \neq 0 \to \lvert c \rvert^{-1} \in \mathbb{R}$
\item $\varphi_{105} :\equiv \forall a\,(v_{386} \in \mathbb{R} \to v_{387} \in \mathbb{R} \to a \in \mathbb{R} \to v_{386} \cdot v_{387} - v_{386} \cdot a = v_{386} \cdot (v_{387} - a))$
\item $\varphi_{106} :\equiv c \in \mathbb{R} \to a_{n} \in \mathbb{R} \to \alpha \in \mathbb{R} \to c \cdot a_{n} - c \cdot \alpha = c \cdot (a_{n} - \alpha)$
\item $\varphi_{107} :\equiv \forall a\,(v_{383} \in \mathbb{R} \to v_{384} \in \mathbb{R} \to a \in \mathbb{R} \to v_{383} \cdot v_{384} - v_{383} \cdot a = v_{383} \cdot (v_{384} - a))$
\item $\varphi_{108} :\equiv \forall x \in \mathbb{R}\,(x \neq 0 \to x \cdot x^{-1} = 1) \to \lvert c \rvert \in \mathbb{R} \to \lvert c \rvert \neq 0 \to \lvert c \rvert \cdot \lvert c \rvert^{-1} = 1$
\item $\varphi_{109} :\equiv c \in \mathbb{R} \to a_{n} - \alpha \in \mathbb{R} \to \lvert c \cdot (a_{n} - \alpha) \rvert = \lvert c \rvert \cdot \lvert a_{n} - \alpha \rvert$
\item $\varphi_{110} :\equiv 0 \in \mathbb{R} \to \lvert a_{n} - \alpha \rvert \in \mathbb{R} \to 0 \cdot \lvert a_{n} - \alpha \rvert = \lvert a_{n} - \alpha \rvert \cdot 0$
\item $\varphi_{111} :\equiv \forall x\,\forall y\,(x \in \mathbb{R} \to y \in \mathbb{R} \to x \cdot y \in \mathbb{R}) \to \forall y\,(c \in \mathbb{R} \to y \in \mathbb{R} \to c \cdot y \in \mathbb{R})$
\item $\varphi_{112} :\equiv \forall y\,(\lvert a_{n} - \alpha \rvert \in \mathbb{R} \to y \in \mathbb{R} \to \lvert a_{n} - \alpha \rvert \cdot y = y \cdot \lvert a_{n} - \alpha \rvert)$
\item $\varphi_{113} :\equiv \forall v_{387}\,\forall v_{388}\,(c \in \mathbb{R} \to v_{387} \in \mathbb{R} \to v_{388} \in \mathbb{R} \to c \cdot v_{387} - c \cdot v_{388} = c \cdot (v_{387} - v_{388}))$
\item $\varphi_{114} :\equiv \forall v_{384}\,\forall v_{385}\,(c \in \mathbb{R} \to v_{384} \in \mathbb{R} \to v_{385} \in \mathbb{R} \to c \cdot v_{384} - c \cdot v_{385} = c \cdot (v_{384} - v_{385}))$
\item $\varphi_{115} :\equiv \varepsilon \cdot (\lvert c \rvert \cdot \lvert c \rvert^{-1}) = \varepsilon \cdot (\lvert c \rvert \cdot \lvert c \rvert^{-1}) \to \varepsilon \cdot (\lvert c \rvert \cdot \lvert c \rvert^{-1}) = \varepsilon \cdot 1$
\item $\varphi_{116} :\equiv \forall x\,\forall y\,(x \in \mathbb{R} \to y \in \mathbb{R} \to x \cdot y \in \mathbb{R}) \to \forall y\,(\varepsilon \in \mathbb{R} \to y \in \mathbb{R} \to \varepsilon \cdot y \in \mathbb{R})$
\item $\varphi_{117} :\equiv N \in \mathbb{N} \land \forall n \in \mathbb{N}\,(N \leq n \to \lvert a_{n} - \alpha \rvert < \varepsilon \cdot \lvert c \rvert^{-1}) \to n \in \mathbb{N} \to n \in \mathbb{N}$
\item $\varphi_{118} :\equiv \lvert c \rvert^{-1} \in \mathbb{R} \to \lvert c \rvert \in \mathbb{R} \to \varepsilon \cdot \lvert c \rvert^{-1} \cdot \lvert c \rvert = \varepsilon \cdot (\lvert c \rvert^{-1} \cdot \lvert c \rvert)$
\item $\varphi_{119} :\equiv \lvert a_{n} - \alpha \rvert \cdot \lvert c \rvert < \varepsilon \cdot \lvert c \rvert^{-1} \cdot \lvert c \rvert \to \lvert a_{n} - \alpha \rvert \cdot \lvert c \rvert < \varepsilon$
\item $\varphi_{120} :\equiv \forall z\,(\varepsilon \in \mathbb{R} \to \lvert c \rvert^{-1} \in \mathbb{R} \to z \in \mathbb{R} \to \varepsilon \cdot \lvert c \rvert^{-1} \cdot z = \varepsilon \cdot (\lvert c \rvert^{-1} \cdot z))$
\item $\varphi_{121} :\equiv \forall v_{388}\,(c \in \mathbb{R} \to a_{n} \in \mathbb{R} \to v_{388} \in \mathbb{R} \to c \cdot a_{n} - c \cdot v_{388} = c \cdot (a_{n} - v_{388}))$
\item $\varphi_{122} :\equiv \forall v_{385}\,(c \in \mathbb{R} \to a_{n} \in \mathbb{R} \to v_{385} \in \mathbb{R} \to c \cdot a_{n} - c \cdot v_{385} = c \cdot (a_{n} - v_{385}))$
\item $\varphi_{123} :\equiv \lvert c \odot a(n) - c \cdot \alpha \rvert = \lvert z - c \cdot \alpha \rvert \to \lvert c \odot a(n) - c \cdot \alpha \rvert = \lvert w - c \cdot \alpha \rvert$
\item $\varphi_{124} :\equiv 0 < \varepsilon \cdot \lvert c \rvert^{-1} \to \exists N \in \mathbb{N}\,\forall n \in \mathbb{N}\,(N \leq n \to \lvert a_{n} - \alpha \rvert < \varepsilon \cdot \lvert c \rvert^{-1})$
\item $\varphi_{125} :\equiv N \in \mathbb{N} \land \forall n \in \mathbb{N}\,(N \leq n \to \lvert a_{n} - \alpha \rvert < \varepsilon \cdot \lvert c \rvert^{-1}) \to n \in \mathbb{N} \to N \leq n \to c \in \mathbb{R}$
\item $\varphi_{126} :\equiv \lvert a_{n} - \alpha \rvert < \varepsilon \cdot \lvert c \rvert^{-1} \to \lvert a_{n} - \alpha \rvert \cdot \lvert c \rvert < \varepsilon \cdot \lvert c \rvert^{-1} \cdot \lvert c \rvert$
\item $\varphi_{127} :\equiv \forall \varepsilon\,(\mathrm{Conv}(a,\alpha) \to \varepsilon \in \mathbb{R} \to 0 < \varepsilon \to \exists N \in \mathbb{N}\,\forall n \in \mathbb{N}\,(N \leq n \to \lvert a_{n} - \alpha \rvert < \varepsilon))$
\item $\varphi_{128} :\equiv 0 \cdot \lvert a_{n} - \alpha \rvert = \lvert a_{n} - \alpha \rvert \cdot 0 \to \lvert c \rvert \cdot \lvert a_{n} - \alpha \rvert = \lvert a_{n} - \alpha \rvert \cdot 0$
\item $\varphi_{129} :\equiv N \in \mathbb{N} \land \forall n \in \mathbb{N}\,(N \leq n \to \lvert a_{n} - \alpha \rvert < \varepsilon \cdot \lvert c \rvert^{-1}) \to n \in \mathbb{N} \to N \leq n \to \varepsilon \cdot 1 = \varepsilon$
\item $\varphi_{130} :\equiv \forall x\,\forall y\,(x \in \mathbb{R} \to y \in \mathbb{R} \to x - y \in \mathbb{R}) \to \forall y\,(a_{n} \in \mathbb{R} \to y \in \mathbb{R} \to a_{n} - y \in \mathbb{R})$
\item $\varphi_{131} :\equiv \lvert a_{n} - \alpha \rvert \in \mathbb{R} \to \lvert c \rvert \in \mathbb{R} \to \lvert a_{n} - \alpha \rvert \cdot \lvert c \rvert = \lvert c \rvert \cdot \lvert a_{n} - \alpha \rvert$
\item $\varphi_{132} :\equiv N \in \mathbb{N} \land \forall n \in \mathbb{N}\,(N \leq n \to \lvert a_{n} - \alpha \rvert < \varepsilon \cdot \lvert c \rvert^{-1}) \to n \in \mathbb{N} \to N \leq n \to \lvert c \rvert \in \mathbb{R}$
\item $\varphi_{133} :\equiv \lvert c \rvert \cdot \lvert a_{n} - \alpha \rvert = z \cdot \lvert a_{n} - \alpha \rvert \to \lvert c \rvert \cdot \lvert a_{n} - \alpha \rvert = w \cdot \lvert a_{n} - \alpha \rvert$
\item $\varphi_{134} :\equiv N \in \mathbb{N} \land \forall n \in \mathbb{N}\,(N \leq n \to \lvert a_{n} - \alpha \rvert < \varepsilon \cdot \lvert c \rvert^{-1}) \to n \in \mathbb{N} \to N \leq n \to a_{n} \in \mathbb{R}$
\item $\varphi_{135} :\equiv \lvert c \cdot a_{n} - c \cdot \alpha \rvert = \lvert c \cdot (a_{n} - \alpha) \rvert \to \lvert c \odot a(n) - c \cdot \alpha \rvert = \lvert c \cdot (a_{n} - \alpha) \rvert$
\item $\varphi_{136} :\equiv N \in \mathbb{N} \land \forall n \in \mathbb{N}\,(N \leq n \to \lvert a_{n} - \alpha \rvert < \varepsilon \cdot \lvert c \rvert^{-1}) \to n \in \mathbb{N} \to N \leq n \to \lvert c \rvert^{-1} \in \mathbb{R}$
\item $\varphi_{137} :\equiv \forall y\,(\varepsilon \in \mathbb{R} \to y \in \mathbb{R} \to \varepsilon \cdot y \in \mathbb{R}) \to \varepsilon \in \mathbb{R} \to \lvert c \rvert^{-1} \in \mathbb{R} \to \varepsilon \cdot \lvert c \rvert^{-1} \in \mathbb{R}$
\item $\varphi_{138} :\equiv \lvert c \cdot a_{n} - c \cdot \alpha \rvert = \lvert c \cdot a_{n} - c \cdot \alpha \rvert \to \lvert c \cdot a_{n} - c \cdot \alpha \rvert = \lvert c \cdot (a_{n} - \alpha) \rvert$
\item $\varphi_{139} :\equiv \lvert c \rvert \cdot \lvert a_{n} - \alpha \rvert = \lvert c \rvert \cdot \lvert a_{n} - \alpha \rvert \to \lvert c \rvert \cdot \lvert a_{n} - \alpha \rvert = 0 \cdot \lvert a_{n} - \alpha \rvert$
\item $\varphi_{140} :\equiv \lvert c \rvert \cdot \lvert a_{n} - \alpha \rvert = \lvert a_{n} - \alpha \rvert \cdot 0 \to \lvert a_{n} - \alpha \rvert \cdot 0 = 0 \to \lvert c \rvert \cdot \lvert a_{n} - \alpha \rvert = 0$
\item $\varphi_{141} :\equiv \varepsilon \cdot (\lvert c \rvert^{-1} \cdot \lvert c \rvert) = \varepsilon \cdot (\lvert c \rvert \cdot \lvert c \rvert^{-1}) \to \varepsilon \cdot \lvert c \rvert^{-1} \cdot \lvert c \rvert = \varepsilon \cdot (\lvert c \rvert \cdot \lvert c \rvert^{-1})$
\item $\varphi_{142} :\equiv \varepsilon \cdot (\lvert c \rvert^{-1} \cdot \lvert c \rvert) = \varepsilon \cdot (\lvert c \rvert^{-1} \cdot \lvert c \rvert) \to \varepsilon \cdot (\lvert c \rvert^{-1} \cdot \lvert c \rvert) = \varepsilon \cdot (\lvert c \rvert \cdot \lvert c \rvert^{-1})$
\item $\varphi_{143} :\equiv \forall y\,(a_{n} \in \mathbb{R} \to y \in \mathbb{R} \to a_{n} - y \in \mathbb{R}) \to a_{n} \in \mathbb{R} \to \alpha \in \mathbb{R} \to a_{n} - \alpha \in \mathbb{R}$
\item $\varphi_{144} :\equiv c \in \mathbb{R} \land \mathrm{Conv}(a,\alpha) \land \varepsilon \in \mathbb{R} \land 0 < \varepsilon \land c = 0 \to n \in \mathbb{N} \to 1_{\omega} \leq n \to \lvert c \odot a(n) - c \cdot \alpha \rvert < \varepsilon$
\item $\varphi_{145} :\equiv N \in \mathbb{N} \land \forall n \in \mathbb{N}\,(N \leq n \to \lvert a_{n} - \alpha \rvert < \varepsilon \cdot \lvert c \rvert^{-1}) \to n \in \mathbb{N} \to N \leq n \to a_{n} - \alpha \in \mathbb{R}$
\item $\varphi_{146} :\equiv N \in \mathbb{N} \land \forall n \in \mathbb{N}\,(N \leq n \to \lvert a_{n} - \alpha \rvert < \varepsilon \cdot \lvert c \rvert^{-1}) \to n \in \mathbb{N} \to N \leq n \to \lvert c \rvert \cdot \lvert c \rvert^{-1} = 1$
\item $\varphi_{147} :\equiv \lvert c \cdot (a_{n} - \alpha) \rvert = \lvert c \rvert \cdot \lvert a_{n} - \alpha \rvert \to \lvert c \odot a(n) - c \cdot \alpha \rvert = \lvert c \rvert \cdot \lvert a_{n} - \alpha \rvert$
\item $\varphi_{148} :\equiv \forall x\,\forall y\,(x \in \mathbb{R} \to y \in \mathbb{R} \to x \cdot y = y \cdot x) \to \forall y\,(\lvert c \rvert^{-1} \in \mathbb{R} \to y \in \mathbb{R} \to \lvert c \rvert^{-1} \cdot y = y \cdot \lvert c \rvert^{-1})$
\item $\varphi_{149} :\equiv \forall x\,\forall y\,(x \in \mathbb{R} \to y \in \mathbb{R} \to \lvert x \cdot y \rvert = \lvert x \rvert \cdot \lvert y \rvert) \to \forall y\,(c \in \mathbb{R} \to y \in \mathbb{R} \to \lvert c \cdot y \rvert = \lvert c \rvert \cdot \lvert y \rvert)$
\item $\varphi_{150} :\equiv \forall z\,(\lvert c \rvert \cdot \lvert a_{n} - \alpha \rvert = \lvert a_{n} - \alpha \rvert \cdot 0 \to \lvert a_{n} - \alpha \rvert \cdot 0 = z \to \lvert c \rvert \cdot \lvert a_{n} - \alpha \rvert = z)$
\item $\varphi_{151} :\equiv \forall z\,(\lvert c \rvert \cdot \lvert a_{n} - \alpha \rvert = 0 \cdot \lvert a_{n} - \alpha \rvert \to 0 \cdot \lvert a_{n} - \alpha \rvert = z \to \lvert c \rvert \cdot \lvert a_{n} - \alpha \rvert = z)$
\item $\varphi_{152} :\equiv N \in \mathbb{N} \land \forall n \in \mathbb{N}\,(N \leq n \to \lvert a_{n} - \alpha \rvert < \varepsilon \cdot \lvert c \rvert^{-1}) \to n \in \mathbb{N} \to N \leq n \to \varepsilon \cdot \lvert c \rvert^{-1} \in \mathbb{R}$
\item $\varphi_{153} :\equiv N \in \mathbb{N} \land \forall n \in \mathbb{N}\,(N \leq n \to \lvert a_{n} - \alpha \rvert < \varepsilon \cdot \lvert c \rvert^{-1}) \to n \in \mathbb{N} \to N \leq n \to \lvert a_{n} - \alpha \rvert \in \mathbb{R}$
\item $\varphi_{154} :\equiv \forall n \in \mathbb{N}\,(N \leq n \to \lvert a_{n} - \alpha \rvert < \varepsilon \cdot \lvert c \rvert^{-1}) \to n \in \mathbb{N} \to N \leq n \to \lvert a_{n} - \alpha \rvert < \varepsilon \cdot \lvert c \rvert^{-1}$
\item $\varphi_{155} :\equiv \varphi_{7} \to \forall n \in \mathbb{N}\,(1_{\omega} \leq n \to \lvert c \odot a(n) - c \cdot \alpha \rvert < \varepsilon)$
\item $\varphi_{156} :\equiv (\lvert c \rvert = \lvert 0 \rvert \to \lvert 0 \rvert = 0 \to \lvert c \rvert = 0) \to c \in \mathbb{R} \to \mathrm{Conv}(a,\alpha) \to \varphi_{10}$
\item $\varphi_{157} :\equiv N \in \mathbb{N} \land \forall n \in \mathbb{N}\,(N \leq n \to \lvert a_{n} - \alpha \rvert < \varepsilon \cdot \lvert c \rvert^{-1}) \to n \in \mathbb{N} \to N \leq n \to c \odot a(n) = c \cdot a_{n}$
\item $\varphi_{158} :\equiv \forall z\,(\varepsilon \cdot \lvert c \rvert^{-1} \cdot \lvert c \rvert = \varepsilon \cdot (\lvert c \rvert \cdot \lvert c \rvert^{-1}) \to \varepsilon \cdot (\lvert c \rvert \cdot \lvert c \rvert^{-1}) = z \to \varepsilon \cdot \lvert c \rvert^{-1} \cdot \lvert c \rvert = z)$
\item $\varphi_{159} :\equiv \forall w\,(\lvert c \rvert \cdot \lvert c \rvert^{-1} = w \to \varepsilon \cdot (\lvert c \rvert \cdot \lvert c \rvert^{-1}) = \varepsilon \cdot (\lvert c \rvert \cdot \lvert c \rvert^{-1}) \to \varepsilon \cdot (\lvert c \rvert \cdot \lvert c \rvert^{-1}) = \varepsilon \cdot w)$
\item $\varphi_{160} :\equiv \forall z\,(\varepsilon \cdot \lvert c \rvert^{-1} \cdot \lvert c \rvert = \varepsilon \cdot (\lvert c \rvert^{-1} \cdot \lvert c \rvert) \to \varepsilon \cdot (\lvert c \rvert^{-1} \cdot \lvert c \rvert) = z \to \varepsilon \cdot \lvert c \rvert^{-1} \cdot \lvert c \rvert = z)$
\item $\varphi_{161} :\equiv \forall w\,(\lvert c \rvert^{-1} \cdot \lvert c \rvert = w \to \varepsilon \cdot (\lvert c \rvert^{-1} \cdot \lvert c \rvert) = \varepsilon \cdot (\lvert c \rvert^{-1} \cdot \lvert c \rvert) \to \varepsilon \cdot (\lvert c \rvert^{-1} \cdot \lvert c \rvert) = \varepsilon \cdot w)$
\item $\varphi_{162} :\equiv \forall z\,(\lvert c \odot a(n) - c \cdot \alpha \rvert = \lvert c \cdot (a_{n} - \alpha) \rvert \to \lvert c \cdot (a_{n} - \alpha) \rvert = z \to \lvert c \odot a(n) - c \cdot \alpha \rvert = z)$
\item $\varphi_{163} :\equiv \lvert c \odot a(n) - c \cdot \alpha \rvert = \lvert c \rvert \cdot \lvert a_{n} - \alpha \rvert \to \lvert c \rvert \cdot \lvert a_{n} - \alpha \rvert = \lvert c \odot a(n) - c \cdot \alpha \rvert$
\item $\varphi_{164} :\equiv N \in \mathbb{N} \land \forall n \in \mathbb{N}\,(N \leq n \to \lvert a_{n} - \alpha \rvert < \varepsilon \cdot \lvert c \rvert^{-1}) \to n \in \mathbb{N} \to N \leq n \to \varepsilon \cdot \lvert c \rvert^{-1} \cdot \lvert c \rvert = \varepsilon$
\item $\varphi_{165} :\equiv \forall x\,\forall y\,(x \in \mathbb{R} \to y \in \mathbb{R} \to 0 < x \to 0 < y \to 0 < x \cdot y) \to \forall y\,(\varepsilon \in \mathbb{R} \to y \in \mathbb{R} \to 0 < \varepsilon \to 0 < y \to 0 < \varepsilon \cdot y)$
\item $\varphi_{166} :\equiv \forall w\,(c \cdot a_{n} - c \cdot \alpha = w \to \lvert c \cdot a_{n} - c \cdot \alpha \rvert = \lvert c \cdot a_{n} - c \cdot \alpha \rvert \to \lvert c \cdot a_{n} - c \cdot \alpha \rvert = \lvert w \rvert)$
\item $\varphi_{167} :\equiv \forall w\,(\lvert c \rvert = w \to \lvert c \rvert \cdot \lvert a_{n} - \alpha \rvert = \lvert c \rvert \cdot \lvert a_{n} - \alpha \rvert \to \lvert c \rvert \cdot \lvert a_{n} - \alpha \rvert = w \cdot \lvert a_{n} - \alpha \rvert)$
\item $\varphi_{168} :\equiv (c = 0 \to \lvert c \rvert = \lvert c \rvert \to \lvert c \rvert = \lvert 0 \rvert) \to c \in \mathbb{R} \to \mathrm{Conv}(a,\alpha) \to \varphi_{22}$
\item $\varphi_{169} :\equiv \lvert c \odot a(n) - c \cdot \alpha \rvert = \lvert c \odot a(n) - c \cdot \alpha \rvert \to \lvert c \odot a(n) - c \cdot \alpha \rvert = \lvert c \cdot a_{n} - c \cdot \alpha \rvert$
\item $\varphi_{170} :\equiv (c \in \mathbb{R} \to \varphi_{54}) \to c \in \mathbb{R} \to \mathrm{Conv}(a,\alpha) \to \mathrm{Seq}(c \odot a)$
\item $\varphi_{171} :\equiv N \in \mathbb{N} \land \forall n \in \mathbb{N}\,(N \leq n \to \lvert a_{n} - \alpha \rvert < \varepsilon \cdot \lvert c \rvert^{-1}) \to n \in \mathbb{N} \to N \leq n \to \lvert a_{n} - \alpha \rvert \cdot \lvert c \rvert < \varepsilon$
\item $\varphi_{172} :\equiv N \in \mathbb{N} \land \forall n \in \mathbb{N}\,(N \leq n \to \lvert a_{n} - \alpha \rvert < \varepsilon \cdot \lvert c \rvert^{-1}) \to n \in \mathbb{N} \to N \leq n \to \lvert c \rvert \cdot \lvert a_{n} - \alpha \rvert < \varepsilon$
\item $\varphi_{173} :\equiv N \in \mathbb{N} \land \forall n \in \mathbb{N}\,(N \leq n \to \lvert a_{n} - \alpha \rvert < \varepsilon \cdot \lvert c \rvert^{-1}) \to n \in \mathbb{N} \to N \leq n \to \varepsilon \cdot (\lvert c \rvert \cdot \lvert c \rvert^{-1}) = \varepsilon \cdot 1$
\item $\varphi_{174} :\equiv N \in \mathbb{N} \land \forall n \in \mathbb{N}\,(N \leq n \to \lvert a_{n} - \alpha \rvert < \varepsilon \cdot \lvert c \rvert^{-1}) \to n \in \mathbb{N} \to N \leq n \to \varepsilon \cdot \lvert c \rvert^{-1} \cdot \lvert c \rvert = \varepsilon \cdot 1$
\item $\varphi_{175} :\equiv \forall y\,\forall z\,(\lvert a_{n} - \alpha \rvert \in \mathbb{R} \to y \in \mathbb{R} \to z \in \mathbb{R} \to 0 < z \to \lvert a_{n} - \alpha \rvert < y \to \lvert a_{n} - \alpha \rvert \cdot z < y \cdot z)$
\item $\varphi_{176} :\equiv \lvert c \odot a(n) - c \cdot \alpha \rvert = \lvert c \rvert \cdot \lvert a_{n} - \alpha \rvert \to \lvert c \rvert \cdot \lvert a_{n} - \alpha \rvert = 0 \to \lvert c \odot a(n) - c \cdot \alpha \rvert = 0$
\item $\varphi_{177} :\equiv c \in \mathbb{R} \land \mathrm{Conv}(a,\alpha) \land \varepsilon \in \mathbb{R} \land 0 < \varepsilon \land c \neq 0 \land (N \in \mathbb{N} \land \forall n \in \mathbb{N}\,(N \leq n \to \lvert a_{n} - \alpha \rvert < \varepsilon \cdot \lvert c \rvert^{-1}))$
\item $\varphi_{178} :\equiv \forall w\,(\varepsilon \cdot \lvert c \rvert^{-1} \cdot \lvert c \rvert = w \to \lvert a_{n} - \alpha \rvert \cdot \lvert c \rvert < \varepsilon \cdot \lvert c \rvert^{-1} \cdot \lvert c \rvert \to \lvert a_{n} - \alpha \rvert \cdot \lvert c \rvert < w)$
\item $\varphi_{179} :\equiv N \in \mathbb{N} \land \forall n \in \mathbb{N}\,(N \leq n \to \lvert a_{n} - \alpha \rvert < \varepsilon \cdot \lvert c \rvert^{-1}) \to n \in \mathbb{N} \to N \leq n \to \lvert c \odot a(n) - c \cdot \alpha \rvert < \varepsilon$
\item $\varphi_{180} :\equiv c \in \mathbb{R} \land \mathrm{Conv}(a,\alpha) \land \varepsilon \in \mathbb{R} \land 0 < \varepsilon \land c \neq 0 \land \neg \exists N \in \mathbb{N}\,\forall n \in \mathbb{N}\,(N \leq n \to \lvert c \odot a(n) - c \cdot \alpha \rvert < \varepsilon)$
\item $\varphi_{181} :\equiv \forall z\,(\lvert c \odot a(n) - c \cdot \alpha \rvert = \lvert c \cdot a_{n} - c \cdot \alpha \rvert \to \lvert c \cdot a_{n} - c \cdot \alpha \rvert = z \to \lvert c \odot a(n) - c \cdot \alpha \rvert = z)$
\item $\varphi_{182} :\equiv \forall z\,(\lvert c \odot a(n) - c \cdot \alpha \rvert = \lvert c \rvert \cdot \lvert a_{n} - \alpha \rvert \to \lvert c \rvert \cdot \lvert a_{n} - \alpha \rvert = z \to \lvert c \odot a(n) - c \cdot \alpha \rvert = z)$
\item $\varphi_{183} :\equiv c \in \mathbb{R} \to \mathrm{Conv}(a,\alpha) \to \varepsilon \in \mathbb{R} \to 0 < \varepsilon \to c = 0 \to n \in \mathbb{N} \to \varphi_{23}$
\item $\varphi_{184} :\equiv \forall a\,\forall \alpha\,(\mathrm{Conv}(a,\alpha) \leftrightarrow \mathrm{Seq}(a) \land \alpha \in \mathbb{R} \land \forall \varepsilon \in \mathbb{R}\,(0 < \varepsilon \to \exists N \in \mathbb{N}\,\forall n \in \mathbb{N}\,(N \leq n \to \lvert a_{n} - \alpha \rvert < \varepsilon)))$
\item $\varphi_{185} :\equiv N \in \mathbb{N} \land \forall n \in \mathbb{N}\,(N \leq n \to \lvert a_{n} - \alpha \rvert < \varepsilon \cdot \lvert c \rvert^{-1}) \to n \in \mathbb{N} \to N \leq n \to \lvert a_{n} - \alpha \rvert < \varepsilon \cdot \lvert c \rvert^{-1}$
\item $\varphi_{186} :\equiv \lvert a_{n} - \alpha \rvert \cdot \lvert c \rvert = \lvert c \rvert \cdot \lvert a_{n} - \alpha \rvert \to \lvert a_{n} - \alpha \rvert \cdot \lvert c \rvert < \varepsilon \to \lvert c \rvert \cdot \lvert a_{n} - \alpha \rvert < \varepsilon$
\item $\varphi_{187} :\equiv \varepsilon \cdot \lvert c \rvert^{-1} \cdot \lvert c \rvert = \varepsilon \cdot (\lvert c \rvert \cdot \lvert c \rvert^{-1}) \to \varepsilon \cdot (\lvert c \rvert \cdot \lvert c \rvert^{-1}) = \varepsilon \cdot 1 \to \varepsilon \cdot \lvert c \rvert^{-1} \cdot \lvert c \rvert = \varepsilon \cdot 1$
\item $\varphi_{188} :\equiv N \in \mathbb{N} \land \forall n \in \mathbb{N}\,(N \leq n \to \lvert a_{n} - \alpha \rvert < \varepsilon \cdot \lvert c \rvert^{-1}) \to \forall n \in \mathbb{N}\,(N \leq n \to \lvert c \odot a(n) - c \cdot \alpha \rvert < \varepsilon)$
\item $\varphi_{189} :\equiv N \in \mathbb{N} \land \forall n \in \mathbb{N}\,(N \leq n \to \lvert a_{n} - \alpha \rvert < \varepsilon \cdot \lvert c \rvert^{-1}) \to n \in \mathbb{N} \to N \leq n \to \lvert c \rvert^{-1} \cdot \lvert c \rvert = \lvert c \rvert \cdot \lvert c \rvert^{-1}$
\item $\varphi_{190} :\equiv N \in \mathbb{N} \land \forall n \in \mathbb{N}\,(N \leq n \to \lvert a_{n} - \alpha \rvert < \varepsilon \cdot \lvert c \rvert^{-1}) \to \forall n \in \mathbb{N}\,(N \leq n \to \lvert a_{n} - \alpha \rvert < \varepsilon \cdot \lvert c \rvert^{-1})$
\item $\varphi_{191} :\equiv N \in \mathbb{N} \land \forall n \in \mathbb{N}\,(N \leq n \to \lvert a_{n} - \alpha \rvert < \varepsilon \cdot \lvert c \rvert^{-1}) \to n \in \mathbb{N} \to N \leq n \to c \cdot a_{n} - c \cdot \alpha = c \cdot (a_{n} - \alpha)$
\item $\varphi_{192} :\equiv c \in \mathbb{R} \to \mathrm{Conv}(a,\alpha) \to \varepsilon \in \mathbb{R} \to 0 < \varepsilon \to c = 0 \to n \in \mathbb{N} \to \varphi_{42}$
\item $\varphi_{193} :\equiv c \in \mathbb{R} \to \mathrm{Conv}(a,\alpha) \to \varepsilon \in \mathbb{R} \to 0 < \varepsilon \to c = 0 \to n \in \mathbb{N} \to \varphi_{43}$
\item $\varphi_{194} :\equiv \forall w\,(c \odot a(n) = w \to \lvert c \odot a(n) - c \cdot \alpha \rvert = \lvert c \odot a(n) - c \cdot \alpha \rvert \to \lvert c \odot a(n) - c \cdot \alpha \rvert = \lvert w - c \cdot \alpha \rvert)$
\item $\varphi_{195} :\equiv (c \in \mathbb{R} \to \mathrm{Conv}(a,\alpha) \to \alpha \in \mathbb{R}) \to c \in \mathbb{R} \to \mathrm{Conv}(a,\alpha) \to \varphi_{65}$
\item $\varphi_{196} :\equiv \lvert c \rvert \cdot \lvert a_{n} - \alpha \rvert = \lvert c \odot a(n) - c \cdot \alpha \rvert \to \lvert c \rvert \cdot \lvert a_{n} - \alpha \rvert < \varepsilon \to \lvert c \odot a(n) - c \cdot \alpha \rvert < \varepsilon$
\item $\varphi_{197} :\equiv \varphi_{48} \to \lvert c \odot a(n) - c \cdot \alpha \rvert = 0 \to 0 = \lvert c \odot a(n) - c \cdot \alpha \rvert$
\item $\varphi_{198} :\equiv (c \in \mathbb{R} \to c \in \mathbb{R}) \to (c \in \mathbb{R} \to \mathrm{Conv}(a,\alpha) \to \mathrm{Seq}(a)) \to c \in \mathbb{R} \to \varphi_{54}$
\item $\varphi_{199} :\equiv N \in \mathbb{N} \land \forall n \in \mathbb{N}\,(N \leq n \to \lvert a_{n} - \alpha \rvert < \varepsilon \cdot \lvert c \rvert^{-1}) \to n \in \mathbb{N} \to N \leq n \to \lvert c \cdot (a_{n} - \alpha) \rvert = \lvert c \rvert \cdot \lvert a_{n} - \alpha \rvert$
\item $\varphi_{200} :\equiv N \in \mathbb{N} \land \forall n \in \mathbb{N}\,(N \leq n \to \lvert a_{n} - \alpha \rvert < \varepsilon \cdot \lvert c \rvert^{-1}) \to n \in \mathbb{N} \to N \leq n \to \lvert c \cdot a_{n} - c \cdot \alpha \rvert = \lvert c \cdot (a_{n} - \alpha) \rvert$
\item $\varphi_{201} :\equiv N \in \mathbb{N} \land \forall n \in \mathbb{N}\,(N \leq n \to \lvert a_{n} - \alpha \rvert < \varepsilon \cdot \lvert c \rvert^{-1}) \to N \in \mathbb{N} \land \forall n \in \mathbb{N}\,(N \leq n \to \lvert c \odot a(n) - c \cdot \alpha \rvert < \varepsilon)$
\item $\varphi_{202} :\equiv N \in \mathbb{N} \land \forall n \in \mathbb{N}\,(N \leq n \to \lvert a_{n} - \alpha \rvert < \varepsilon \cdot \lvert c \rvert^{-1}) \to N \in \mathbb{N} \land \forall n \in \mathbb{N}\,(N \leq n \to \lvert a_{n} - \alpha \rvert < \varepsilon \cdot \lvert c \rvert^{-1})$
\item $\varphi_{203} :\equiv N \in \mathbb{N} \land \forall n \in \mathbb{N}\,(N \leq n \to \lvert a_{n} - \alpha \rvert < \varepsilon \cdot \lvert c \rvert^{-1}) \to n \in \mathbb{N} \to N \leq n \to \varepsilon \cdot \lvert c \rvert^{-1} \cdot \lvert c \rvert = \varepsilon \cdot (\lvert c \rvert^{-1} \cdot \lvert c \rvert)$
\item $\varphi_{204} :\equiv N \in \mathbb{N} \land \forall n \in \mathbb{N}\,(N \leq n \to \lvert a_{n} - \alpha \rvert < \varepsilon \cdot \lvert c \rvert^{-1}) \to n \in \mathbb{N} \to N \leq n \to \varepsilon \cdot \lvert c \rvert^{-1} \cdot \lvert c \rvert = \varepsilon \cdot (\lvert c \rvert \cdot \lvert c \rvert^{-1})$
\item $\varphi_{205} :\equiv N \in \mathbb{N} \land \forall n \in \mathbb{N}\,(N \leq n \to \lvert a_{n} - \alpha \rvert < \varepsilon \cdot \lvert c \rvert^{-1}) \to \exists N \in \mathbb{N}\,\forall n \in \mathbb{N}\,(N \leq n \to \lvert c \odot a(n) - c \cdot \alpha \rvert < \varepsilon)$
\item $\varphi_{206} :\equiv N \in \mathbb{N} \land \forall n \in \mathbb{N}\,(N \leq n \to \lvert a_{n} - \alpha \rvert < \varepsilon \cdot \lvert c \rvert^{-1}) \to n \in \mathbb{N} \to N \leq n \to \lvert a_{n} - \alpha \rvert \cdot \lvert c \rvert = \lvert c \rvert \cdot \lvert a_{n} - \alpha \rvert$
\item $\varphi_{207} :\equiv N \in \mathbb{N} \land \forall n \in \mathbb{N}\,(N \leq n \to \lvert a_{n} - \alpha \rvert < \varepsilon \cdot \lvert c \rvert^{-1}) \to n \in \mathbb{N} \to N \leq n \to \lvert a_{n} - \alpha \rvert \cdot \lvert c \rvert < \varepsilon \cdot \lvert c \rvert^{-1} \cdot \lvert c \rvert$
\item $\varphi_{208} :\equiv \forall \alpha\,(\mathrm{Conv}(c \odot a,\alpha) \leftrightarrow \mathrm{Seq}(c \odot a) \land \alpha \in \mathbb{R} \land \forall \varepsilon \in \mathbb{R}\,(0 < \varepsilon \to \exists N \in \mathbb{N}\,\forall n \in \mathbb{N}\,(N \leq n \to \lvert c \odot a(n) - \alpha \rvert < \varepsilon)))$
\item $\varphi_{209} :\equiv \forall y\,(\varepsilon \in \mathbb{R} \to y \in \mathbb{R} \to 0 < \varepsilon \to 0 < y \to 0 < \varepsilon \cdot y) \to \varepsilon \in \mathbb{R} \to \lvert c \rvert^{-1} \in \mathbb{R} \to 0 < \varepsilon \to 0 < \lvert c \rvert^{-1} \to 0 < \varepsilon \cdot \lvert c \rvert^{-1}$
\item $\varphi_{210} :\equiv \forall y\,(\lvert c \rvert^{-1} \in \mathbb{R} \to y \in \mathbb{R} \to \lvert c \rvert^{-1} \cdot y = y \cdot \lvert c \rvert^{-1}) \to \lvert c \rvert^{-1} \in \mathbb{R} \to \lvert c \rvert \in \mathbb{R} \to \lvert c \rvert^{-1} \cdot \lvert c \rvert = \lvert c \rvert \cdot \lvert c \rvert^{-1}$
\item $\varphi_{211} :\equiv c \neq 0 \to N \in \mathbb{N} \land \forall n \in \mathbb{N}\,(N \leq n \to \lvert a_{n} - \alpha \rvert < \varepsilon \cdot \lvert c \rvert^{-1}) \to \varphi_{53}$
\item $\varphi_{212} :\equiv c \neq 0 \to N \in \mathbb{N} \land \forall n \in \mathbb{N}\,(N \leq n \to \lvert a_{n} - \alpha \rvert < \varepsilon \cdot \lvert c \rvert^{-1}) \to \varphi_{55}$
\item $\varphi_{213} :\equiv (c \in \mathbb{R} \to \mathrm{Conv}(a,\alpha) \to \alpha \in \mathbb{R}) \to c \in \mathbb{R} \to \mathrm{Conv}(a,\alpha) \to \varepsilon \in \mathbb{R} \to 0 < \varepsilon \to \varphi_{41}$
\item $\varphi_{214} :\equiv c \neq 0 \to N \in \mathbb{N} \land \forall n \in \mathbb{N}\,(N \leq n \to \lvert a_{n} - \alpha \rvert < \varepsilon \cdot \lvert c \rvert^{-1}) \to \varphi_{74}$
\item $\varphi_{215} :\equiv c \neq 0 \to N \in \mathbb{N} \land \forall n \in \mathbb{N}\,(N \leq n \to \lvert a_{n} - \alpha \rvert < \varepsilon \cdot \lvert c \rvert^{-1}) \to \varphi_{75}$
\item $\varphi_{216} :\equiv c \neq 0 \to N \in \mathbb{N} \land \forall n \in \mathbb{N}\,(N \leq n \to \lvert a_{n} - \alpha \rvert < \varepsilon \cdot \lvert c \rvert^{-1}) \to \varphi_{76}$
\item $\varphi_{217} :\equiv (c \in \mathbb{R} \to \alpha \in \mathbb{R} \to c \cdot \alpha \in \mathbb{R}) \to (c \in \mathbb{R} \to \mathrm{Conv}(a,\alpha) \to c \in \mathbb{R}) \to \varphi_{84}$
\item $\varphi_{218} :\equiv \mathrm{Conv}(c \odot a,c \cdot \alpha) \leftrightarrow \mathrm{Seq}(c \odot a) \land c \cdot \alpha \in \mathbb{R} \land \forall \varepsilon \in \mathbb{R}\,(\varphi_{95})$
\item $\varphi_{219} :\equiv \forall z\,(\lvert a_{n} - \alpha \rvert \in \mathbb{R} \to \varepsilon \cdot \lvert c \rvert^{-1} \in \mathbb{R} \to z \in \mathbb{R} \to 0 < z \to \lvert a_{n} - \alpha \rvert < \varepsilon \cdot \lvert c \rvert^{-1} \to \lvert a_{n} - \alpha \rvert \cdot z < \varepsilon \cdot \lvert c \rvert^{-1} \cdot z)$
\item $\varphi_{220} :\equiv (c \in \mathbb{R} \to \mathrm{Conv}(a,\alpha) \to \forall \varepsilon \in \mathbb{R}\,(\varphi_{95})) \to c \in \mathbb{R} \to \mathrm{Conv}(a,\alpha) \to \mathrm{Conv}(c \odot a,c \cdot \alpha)$
\item $\varphi_{221} :\equiv (\lvert c \odot a(n) - c \cdot \alpha \rvert = 0 \to 0 = \lvert c \odot a(n) - c \cdot \alpha \rvert) \to c \in \mathbb{R} \to \mathrm{Conv}(a,\alpha) \to \varepsilon \in \mathbb{R} \to 0 < \varepsilon \to \varphi_{12}$
\item $\varphi_{222} :\equiv (0 \leq 0 \to \lvert 0 \rvert = 0) \to (c \in \mathbb{R} \to \varphi_{31}) \to c \in \mathbb{R} \to \mathrm{Conv}(a,\alpha) \to \varphi_{9}$
\item $\varphi_{223} :\equiv \lvert a_{n} - \alpha \rvert \in \mathbb{R} \to \varepsilon \cdot \lvert c \rvert^{-1} \in \mathbb{R} \to \lvert c \rvert \in \mathbb{R} \to 0 < \lvert c \rvert \to \varphi_{126}$
\item $\varphi_{224} :\equiv \forall n\,(\varphi_{177}) \to \forall n \in \mathbb{N}\,(N \leq n \to \lvert c \odot a(n) - c \cdot \alpha \rvert < \varepsilon)$
\item $\varphi_{225} :\equiv (c \in \mathbb{R} \to \mathrm{Conv}(a,\alpha) \to \varphi_{65}) \to c \in \mathbb{R} \to \mathrm{Conv}(a,\alpha) \to \varphi_{78}$
\item $\varphi_{226} :\equiv (c \in \mathbb{R} \to \mathrm{Conv}(a,\alpha) \to \mathrm{Seq}(a)) \to (c \in \mathbb{R} \to \varphi_{21}) \to c \in \mathbb{R} \to \mathrm{Conv}(a,\alpha) \to \varphi_{35}$
\item $\varphi_{227} :\equiv \forall v_{386}\,\forall v_{387}\,\forall v_{388}\,(v_{386} \in \mathbb{R} \to v_{387} \in \mathbb{R} \to \varphi_{17}) \to \varphi_{113}$
\item $\varphi_{228} :\equiv (c \in \mathbb{R} \to \mathrm{Conv}(a,\alpha) \to \varphi_{78}) \to c \in \mathbb{R} \to \mathrm{Conv}(a,\alpha) \to \varepsilon \in \mathbb{R} \to \varphi_{16}$
\item $\varphi_{229} :\equiv \forall v_{383}\,\forall v_{384}\,\forall v_{385}\,(v_{383} \in \mathbb{R} \to v_{384} \in \mathbb{R} \to \varphi_{18}) \to \varphi_{114}$
\item $\varphi_{230} :\equiv (c \in \mathbb{R} \to \mathrm{Conv}(a,\alpha) \to \alpha \in \mathbb{R}) \to c \in \mathbb{R} \to \mathrm{Conv}(a,\alpha) \to \varepsilon \in \mathbb{R} \to 0 < \varepsilon \to c \neq 0 \to \varphi_{145}$
\item $\varphi_{231} :\equiv (c \in \mathbb{R} \to \mathrm{Conv}(a,\alpha) \to \varepsilon \in \mathbb{R} \to \varepsilon \in \mathbb{R}) \to c \in \mathbb{R} \to \mathrm{Conv}(a,\alpha) \to \varepsilon \in \mathbb{R} \to 0 < \varepsilon \to c \neq 0 \to \varphi_{129}$
\item $\varphi_{232} :\equiv (c \in \mathbb{R} \to \mathrm{Conv}(a,\alpha) \to \varepsilon \in \mathbb{R} \to 0 < \varepsilon \to 0 < \varepsilon) \to \varphi_{52} \to \varphi_{82}$
\item $\varphi_{233} :\equiv (c \in \mathbb{R} \to \mathrm{Conv}(a,\alpha) \to \alpha \in \mathbb{R}) \to c \in \mathbb{R} \to \mathrm{Conv}(a,\alpha) \to \varepsilon \in \mathbb{R} \to 0 < \varepsilon \to c \neq 0 \to \varphi_{191}$
\item $\varphi_{234} :\equiv (c \in \mathbb{R} \to \mathrm{Conv}(a,\alpha) \to \varphi_{78}) \to c \in \mathbb{R} \to \mathrm{Conv}(a,\alpha) \to \varepsilon \in \mathbb{R} \to 0 < \varepsilon \to \varphi_{66}$
\item $\varphi_{235} :\equiv (c \in \mathbb{R} \to \mathrm{Conv}(a,\alpha) \to \varphi_{65}) \to c \in \mathbb{R} \to \mathrm{Conv}(a,\alpha) \to \varepsilon \in \mathbb{R} \to 0 < \varepsilon \to c = 0 \to \varphi_{59}$
\item $\varphi_{236} :\equiv (c \in \mathbb{R} \land \mathrm{Conv}(a,\alpha) \land \varepsilon \in \mathbb{R} \land 0 < \varepsilon \land c = 0 \to \varphi_{7}) \to c \in \mathbb{R} \to \mathrm{Conv}(a,\alpha) \to \varepsilon \in \mathbb{R} \to 0 < \varepsilon \to \varphi_{33}$
\item $\varphi_{237} :\equiv \varphi_{27} \to c \in \mathbb{R} \to \mathrm{Conv}(a,\alpha) \to \varepsilon \in \mathbb{R} \to 0 < \varepsilon \to c \neq 0 \to \varphi_{136}$
\item $\varphi_{238} :\equiv \varphi_{27} \to c \in \mathbb{R} \to \mathrm{Conv}(a,\alpha) \to \varepsilon \in \mathbb{R} \to 0 < \varepsilon \to c \neq 0 \to \varphi_{146}$
\item $\varphi_{239} :\equiv (c \in \mathbb{R} \to \mathrm{Conv}(a,\alpha) \to \varepsilon \in \mathbb{R} \to \varphi_{95}) \to c \in \mathbb{R} \land \mathrm{Conv}(a,\alpha) \to \varepsilon \in \mathbb{R} \to \varphi_{95}$
\item $\varphi_{240} :\equiv (c \in \mathbb{R} \to \mathrm{Conv}(a,\alpha) \to \varepsilon \in \mathbb{R} \to 0 < \varepsilon \to \varphi_{33}) \to c \in \mathbb{R} \to \mathrm{Conv}(a,\alpha) \to \varepsilon \in \mathbb{R} \to 0 < \varepsilon \to c = 0 \to \varphi_{71}$
\item $\varphi_{241} :\equiv \varphi_{28} \to (c \in \mathbb{R} \to \varphi_{21}) \to c \in \mathbb{R} \to \mathrm{Conv}(a,\alpha) \to \varepsilon \in \mathbb{R} \to \varphi_{29}$
\item $\varphi_{242} :\equiv (\lvert c \rvert = 0 \to \varphi_{139}) \to c \in \mathbb{R} \to \mathrm{Conv}(a,\alpha) \to \varepsilon \in \mathbb{R} \to 0 < \varepsilon \to c = 0 \to \varphi_{57}$
\item $\varphi_{243} :\equiv (c \in \mathbb{R} \to \mathrm{Conv}(a,\alpha) \to \varepsilon \in \mathbb{R} \to 0 < \varepsilon \to \varphi_{85}) \to c \in \mathbb{R} \to \mathrm{Conv}(a,\alpha) \to \varepsilon \in \mathbb{R} \to \varphi_{95}$
\item $\varphi_{244} :\equiv (c \in \mathbb{R} \to \mathrm{Conv}(a,\alpha) \to \varepsilon \in \mathbb{R} \to 0 < \varepsilon \to 0 < \varepsilon) \to \varphi_{98} \to c \in \mathbb{R} \to \mathrm{Conv}(a,\alpha) \to \varepsilon \in \mathbb{R} \to 0 < \varepsilon \to \varphi_{24}$
\item $\varphi_{245} :\equiv \forall \varepsilon\,(c \in \mathbb{R} \land \mathrm{Conv}(a,\alpha) \to \varepsilon \in \mathbb{R} \to \varphi_{95}) \to \forall \varepsilon\,(c \in \mathbb{R} \land \mathrm{Conv}(a,\alpha)) \to \forall \varepsilon \in \mathbb{R}\,(\varphi_{95})$
\item $\varphi_{246} :\equiv (c \in \mathbb{R} \to \mathrm{Conv}(a,\alpha) \to \mathrm{Seq}(c \odot a)) \to (c \in \mathbb{R} \to \mathrm{Conv}(a,\alpha) \to c \cdot \alpha \in \mathbb{R}) \to \varphi_{220}$
\item $\varphi_{247} :\equiv \varphi_{96} \to c \in \mathbb{R} \to \mathrm{Conv}(a,\alpha) \to \varepsilon \in \mathbb{R} \to 0 < \varepsilon \to c \neq 0 \to \varphi_{164}$
\item $\varphi_{248} :\equiv (\forall N\,(\varphi_{45}) \to \neg \varphi_{71}) \to c \in \mathbb{R} \to \mathrm{Conv}(a,\alpha) \to \varepsilon \in \mathbb{R} \to 0 < \varepsilon \to \varphi_{79}$
\item $\varphi_{249} :\equiv (\lvert c \rvert \cdot \lvert a_{n} - \alpha \rvert = 0 \cdot \lvert a_{n} - \alpha \rvert \to \varphi_{128}) \to c \in \mathbb{R} \to \mathrm{Conv}(a,\alpha) \to \varepsilon \in \mathbb{R} \to 0 < \varepsilon \to c = 0 \to \varphi_{58}$
\item $\varphi_{250} :\equiv (c \cdot a_{n} - c \cdot \alpha = c \cdot (a_{n} - \alpha) \to \varphi_{138}) \to c \in \mathbb{R} \to \mathrm{Conv}(a,\alpha) \to \varepsilon \in \mathbb{R} \to 0 < \varepsilon \to c = 0 \to \varphi_{60}$
\item $\varphi_{251} :\equiv (a_{n} \in \mathbb{R} \to \alpha \in \mathbb{R} \to a_{n} - \alpha \in \mathbb{R}) \to (c \in \mathbb{R} \to \mathrm{Conv}(a,\alpha) \to \varphi_{35}) \to \varphi_{195}$
\item $\varphi_{252} :\equiv (\varepsilon \in \mathbb{R} \to \lvert c \rvert^{-1} \in \mathbb{R} \to \varepsilon \cdot \lvert c \rvert^{-1} \in \mathbb{R}) \to \varphi_{39} \to \varphi_{72} \to c \in \mathbb{R} \to \varphi_{46}$
\item $\varphi_{253} :\equiv (c \in \mathbb{R} \to \varphi_{46}) \to \varphi_{82} \to c \in \mathbb{R} \to \mathrm{Conv}(a,\alpha) \to \varepsilon \in \mathbb{R} \to 0 < \varepsilon \to \varphi_{90}$
\item $\varphi_{254} :\equiv (c \in \mathbb{R} \to \mathrm{Conv}(a,\alpha) \to \varphi_{22}) \to (c \in \mathbb{R} \to \mathrm{Conv}(a,\alpha) \to \varphi_{9}) \to \varphi_{156}$
\item $\varphi_{255} :\equiv (\lvert c \rvert \cdot \lvert c \rvert^{-1} = 1 \to \varphi_{115}) \to c \in \mathbb{R} \to \mathrm{Conv}(a,\alpha) \to \varepsilon \in \mathbb{R} \to 0 < \varepsilon \to c \neq 0 \to \varphi_{173}$
\item $\varphi_{256} :\equiv \varphi_{154} \to c \in \mathbb{R} \to \mathrm{Conv}(a,\alpha) \to \varepsilon \in \mathbb{R} \to 0 < \varepsilon \to c \neq 0 \to \varphi_{185}$
\item $\varphi_{257} :\equiv (\varepsilon \cdot \lvert c \rvert^{-1} \cdot \lvert c \rvert = \varepsilon \to \varphi_{119}) \to c \in \mathbb{R} \to \mathrm{Conv}(a,\alpha) \to \varepsilon \in \mathbb{R} \to 0 < \varepsilon \to c \neq 0 \to \varphi_{171}$
\item $\varphi_{258} :\equiv \varphi_{186} \to c \in \mathbb{R} \to \mathrm{Conv}(a,\alpha) \to \varepsilon \in \mathbb{R} \to 0 < \varepsilon \to c \neq 0 \to \varphi_{172}$
\item $\varphi_{259} :\equiv (c \in \mathbb{R} \to \mathrm{Conv}(a,\alpha) \to \varepsilon \in \mathbb{R} \to 0 < \varepsilon \to c \neq 0 \to \varphi_{117}) \to c \in \mathbb{R} \to \mathrm{Conv}(a,\alpha) \to \varepsilon \in \mathbb{R} \to 0 < \varepsilon \to c \neq 0 \to \varphi_{134}$
\item $\varphi_{260} :\equiv \varphi_{187} \to c \in \mathbb{R} \to \mathrm{Conv}(a,\alpha) \to \varepsilon \in \mathbb{R} \to 0 < \varepsilon \to c \neq 0 \to \varphi_{174}$
\item $\varphi_{261} :\equiv (c \in \mathbb{R} \to \mathrm{Conv}(a,\alpha) \to \varepsilon \in \mathbb{R} \to 0 < \varepsilon \to c \neq 0 \to \varphi_{125}) \to c \in \mathbb{R} \to \mathrm{Conv}(a,\alpha) \to \varepsilon \in \mathbb{R} \to 0 < \varepsilon \to c \neq 0 \to \varphi_{132}$
\item $\varphi_{262} :\equiv (c \in \mathbb{R} \to \mathrm{Conv}(a,\alpha) \to \varepsilon \in \mathbb{R} \to \varphi_{16}) \to \varphi_{140} \to c \in \mathbb{R} \to \mathrm{Conv}(a,\alpha) \to \varepsilon \in \mathbb{R} \to \varphi_{32}$
\item $\varphi_{263} :\equiv \varphi_{196} \to c \in \mathbb{R} \to \mathrm{Conv}(a,\alpha) \to \varepsilon \in \mathbb{R} \to 0 < \varepsilon \to c \neq 0 \to \varphi_{179}$
\item $\varphi_{264} :\equiv (c \in \mathbb{R} \to \mathrm{Conv}(a,\alpha) \to \varepsilon \in \mathbb{R} \to 0 < \varepsilon \to c \neq 0 \to \varphi_{117}) \to c \in \mathbb{R} \to \mathrm{Conv}(a,\alpha) \to \varepsilon \in \mathbb{R} \to 0 < \varepsilon \to c \neq 0 \to \varphi_{157}$
\item $\varphi_{265} :\equiv (c \in \mathbb{R} \to \mathrm{Conv}(a,\alpha) \to \varepsilon \in \mathbb{R} \to 0 < \varepsilon \to c \neq 0 \to \varphi_{136}) \to c \in \mathbb{R} \to \mathrm{Conv}(a,\alpha) \to \varepsilon \in \mathbb{R} \to 0 < \varepsilon \to c \neq 0 \to \varphi_{152}$
\item $\varphi_{266} :\equiv (c \in \mathbb{R} \to \mathrm{Conv}(a,\alpha) \to \varepsilon \in \mathbb{R} \to 0 < \varepsilon \to c \neq 0 \to \varphi_{202}) \to c \in \mathbb{R} \to \mathrm{Conv}(a,\alpha) \to \varepsilon \in \mathbb{R} \to 0 < \varepsilon \to c \neq 0 \to \varphi_{94}$
\item $\varphi_{267} :\equiv (c \in \mathbb{R} \to \mathrm{Conv}(a,\alpha) \to \varepsilon \in \mathbb{R} \to 0 < \varepsilon \to c \neq 0 \to \varphi_{145}) \to c \in \mathbb{R} \to \mathrm{Conv}(a,\alpha) \to \varepsilon \in \mathbb{R} \to 0 < \varepsilon \to c \neq 0 \to \varphi_{153}$
\item $\varphi_{268} :\equiv (c \in \mathbb{R} \to \mathrm{Conv}(a,\alpha) \to \varepsilon \in \mathbb{R} \to \varphi_{32}) \to \varphi_{176} \to c \in \mathbb{R} \to \mathrm{Conv}(a,\alpha) \to \varepsilon \in \mathbb{R} \to 0 < \varepsilon \to \varphi_{11}$
\item $\varphi_{269} :\equiv (\forall N\,(\varphi_{45}) \to \varphi_{45}) \to c \in \mathbb{R} \to \mathrm{Conv}(a,\alpha) \to \varepsilon \in \mathbb{R} \to 0 < \varepsilon \to c \neq 0 \to \varphi_{205}$
\item $\varphi_{270} :\equiv (c \in \mathbb{R} \to \mathrm{Conv}(a,\alpha) \to \varepsilon \in \mathbb{R} \to 0 < \varepsilon \to c \neq 0 \to \varphi_{132}) \to c \in \mathbb{R} \to \mathrm{Conv}(a,\alpha) \to \varepsilon \in \mathbb{R} \to 0 < \varepsilon \to c \neq 0 \to \varphi_{189}$
\item $\varphi_{271} :\equiv (c \cdot a_{n} - c \cdot \alpha = c \cdot (a_{n} - \alpha) \to \varphi_{138}) \to c \in \mathbb{R} \to \mathrm{Conv}(a,\alpha) \to \varepsilon \in \mathbb{R} \to 0 < \varepsilon \to c \neq 0 \to \varphi_{200}$
\item $\varphi_{272} :\equiv \lvert c \odot a(n) - c \cdot \alpha \rvert = \lvert c \odot a(n) - c \cdot \alpha \rvert \to (c \odot a(n) = c \cdot a_{n} \to \varphi_{169}) \to \varphi_{192}$
\item $\varphi_{273} :\equiv (\lvert c \rvert^{-1} \cdot \lvert c \rvert = \lvert c \rvert \cdot \lvert c \rvert^{-1} \to \varphi_{142}) \to c \in \mathbb{R} \to \mathrm{Conv}(a,\alpha) \to \varepsilon \in \mathbb{R} \to 0 < \varepsilon \to \varphi_{211}$
\item $\varphi_{274} :\equiv (c \odot a(n) = c \cdot a_{n} \to \varphi_{169}) \to c \in \mathbb{R} \to \mathrm{Conv}(a,\alpha) \to \varepsilon \in \mathbb{R} \to 0 < \varepsilon \to \varphi_{214}$
\item $\varphi_{275} :\equiv (c \in \mathbb{R} \to \mathrm{Conv}(a,\alpha) \to \varepsilon \in \mathbb{R} \to 0 < \varepsilon \to c \neq 0 \to \varphi_{132}) \to c \in \mathbb{R} \to \mathrm{Conv}(a,\alpha) \to \varepsilon \in \mathbb{R} \to 0 < \varepsilon \to c \neq 0 \to \varphi_{203}$
\item $\varphi_{276} :\equiv (c \in \mathbb{R} \to \mathrm{Conv}(a,\alpha) \to \varepsilon \in \mathbb{R} \to 0 < \varepsilon \to c \neq 0 \to \varphi_{132}) \to c \in \mathbb{R} \to \mathrm{Conv}(a,\alpha) \to \varepsilon \in \mathbb{R} \to 0 < \varepsilon \to c \neq 0 \to \varphi_{206}$
\item $\varphi_{277} :\equiv (c \in \mathbb{R} \to \mathrm{Conv}(a,\alpha) \to \varepsilon \in \mathbb{R} \to 0 < \varepsilon \to c \neq 0 \to \varphi_{145}) \to c \in \mathbb{R} \to \mathrm{Conv}(a,\alpha) \to \varepsilon \in \mathbb{R} \to 0 < \varepsilon \to c \neq 0 \to \varphi_{199}$
\item $\varphi_{278} :\equiv (c \in \mathbb{R} \to \mathrm{Conv}(a,\alpha) \to \varepsilon \in \mathbb{R} \to 0 < \varepsilon \to c \neq 0 \to \varphi_{179}) \to \varphi_{177} \to n \in \mathbb{N} \to N \leq n \to \lvert c \odot a(n) - c \cdot \alpha \rvert < \varepsilon$
\item $\varphi_{279} :\equiv (c \in \mathbb{R} \to \mathrm{Conv}(a,\alpha) \to \varepsilon \in \mathbb{R} \to 0 < \varepsilon \to c \neq 0 \to \varphi_{134}) \to \varphi_{230}$
\item $\varphi_{280} :\equiv (\varepsilon \cdot \lvert c \rvert^{-1} \cdot \lvert c \rvert = \varepsilon \cdot (\lvert c \rvert^{-1} \cdot \lvert c \rvert) \to \varphi_{141}) \to c \in \mathbb{R} \to \mathrm{Conv}(a,\alpha) \to \varepsilon \in \mathbb{R} \to 0 < \varepsilon \to c \neq 0 \to \varphi_{204}$
\item $\varphi_{281} :\equiv (\lvert c \odot a(n) - c \cdot \alpha \rvert = \lvert c \cdot a_{n} - c \cdot \alpha \rvert \to \varphi_{135}) \to c \in \mathbb{R} \to \mathrm{Conv}(a,\alpha) \to \varepsilon \in \mathbb{R} \to 0 < \varepsilon \to \varphi_{212}$
\item $\varphi_{282} :\equiv \forall n\,(\varphi_{177} \to n \in \mathbb{N} \to N \leq n \to \lvert c \odot a(n) - c \cdot \alpha \rvert < \varepsilon) \to \varphi_{224}$
\item $\varphi_{283} :\equiv (\lvert c \odot a(n) - c \cdot \alpha \rvert = \lvert c \cdot (a_{n} - \alpha) \rvert \to \varphi_{147}) \to c \in \mathbb{R} \to \mathrm{Conv}(a,\alpha) \to \varepsilon \in \mathbb{R} \to 0 < \varepsilon \to \varphi_{215}$
\item $\varphi_{284} :\equiv (\mathrm{Seq}(a) \to n \in \mathbb{N} \to a_{n} \in \mathbb{R}) \to (c \in \mathbb{R} \to \mathrm{Conv}(a,\alpha) \to \mathrm{Seq}(a)) \to \varphi_{259}$
\item $\varphi_{285} :\equiv (c \in \mathbb{R} \to \mathrm{Conv}(a,\alpha) \to \varepsilon \in \mathbb{R} \to 0 < \varepsilon \to c \neq 0 \to \varphi_{188}) \to c \in \mathbb{R} \to \mathrm{Conv}(a,\alpha) \to \varepsilon \in \mathbb{R} \to 0 < \varepsilon \to c \neq 0 \to \varphi_{201}$
\item $\varphi_{286} :\equiv (c \in \mathbb{R} \to \mathrm{Conv}(a,\alpha) \to \varepsilon \in \mathbb{R} \to 0 < \varepsilon \to c \neq 0 \to \varphi_{202}) \to c \in \mathbb{R} \to \mathrm{Conv}(a,\alpha) \to \varepsilon \in \mathbb{R} \to 0 < \varepsilon \to c \neq 0 \to \varphi_{190}$
\item $\varphi_{287} :\equiv (c \in \mathbb{R} \to \mathrm{Conv}(a,\alpha) \to \varepsilon \in \mathbb{R} \to 0 < \varepsilon \to c \neq 0 \to \varphi_{134}) \to \varphi_{233}$
\item $\varphi_{288} :\equiv (c \in \mathbb{R} \to \mathrm{Conv}(a,\alpha) \to \varepsilon \in \mathbb{R} \to 0 < \varepsilon \to c \neq 0 \to \varphi_{185}) \to c \in \mathbb{R} \to \mathrm{Conv}(a,\alpha) \to \varepsilon \in \mathbb{R} \to 0 < \varepsilon \to c \neq 0 \to \varphi_{207}$
\item $\varphi_{289} :\equiv (\varphi_{180} \to \forall N\,(\varphi_{180})) \to c \in \mathbb{R} \to \mathrm{Conv}(a,\alpha) \to \varepsilon \in \mathbb{R} \to 0 < \varepsilon \to \varphi_{85}$
\item $\varphi_{290} :\equiv (c \in \mathbb{R} \to \mathrm{Conv}(a,\alpha) \to \varphi_{10}) \to \lvert c \rvert \cdot \lvert a_{n} - \alpha \rvert = \lvert c \rvert \cdot \lvert a_{n} - \alpha \rvert \to \varphi_{242}$
\item $\varphi_{291} :\equiv (c \in \mathbb{R} \to \mathrm{Conv}(a,\alpha) \to \varepsilon \in \mathbb{R} \to 0 < \varepsilon \to c \neq 0 \to \varphi_{132}) \to \varphi_{237}$
\item $\varphi_{292} :\equiv (c \in \mathbb{R} \to \mathrm{Conv}(a,\alpha) \to \varepsilon \in \mathbb{R} \to 0 < \varepsilon \to c \neq 0 \to \varphi_{132}) \to \varphi_{238}$
\item $\varphi_{293} :\equiv (c \in \mathbb{R} \to \mathrm{Conv}(a,\alpha) \to \varepsilon \in \mathbb{R} \to 0 < \varepsilon \to c \neq 0 \to \varphi_{205}) \to \varphi_{180} \to \varphi_{56}$
\item $\varphi_{294} :\equiv (c \in \mathbb{R} \to \mathrm{Conv}(a,\alpha) \to \varepsilon \in \mathbb{R} \to 0 < \varepsilon \to c = 0 \to \varphi_{71}) \to \varphi_{248}$
\item $\varphi_{295} :\equiv (c \in \mathbb{R} \to \mathrm{Conv}(a,\alpha) \to \varepsilon \in \mathbb{R} \to 0 < \varepsilon \to c \neq 0 \to \varphi_{129}) \to \varphi_{247}$
\item $\varphi_{296} :\equiv \varphi_{106} \to (c \in \mathbb{R} \to \varphi_{67}) \to (c \in \mathbb{R} \to \mathrm{Conv}(a,\alpha) \to \varphi_{35}) \to \varphi_{213}$
\item $\varphi_{297} :\equiv (c \in \mathbb{R} \to \mathrm{Conv}(a,\alpha) \to \varepsilon \in \mathbb{R} \to 0 < \varepsilon \to c = 0 \to \varphi_{60}) \to (\lvert c \odot a(n) - c \cdot \alpha \rvert = \lvert c \cdot a_{n} - c \cdot \alpha \rvert \to \varphi_{135}) \to \varphi_{183}$
\item $\varphi_{298} :\equiv (\varepsilon \in \mathbb{R} \to \lvert c \rvert^{-1} \in \mathbb{R} \to 0 < \varepsilon \to 0 < \lvert c \rvert^{-1} \to 0 < \varepsilon \cdot \lvert c \rvert^{-1}) \to (c \in \mathbb{R} \to \mathrm{Conv}(a,\alpha) \to \varepsilon \in \mathbb{R} \to \varepsilon \in \mathbb{R}) \to \varphi_{72} \to \varphi_{232}$
\item $\varphi_{299} :\equiv (c \in \mathbb{R} \to \mathrm{Conv}(a,\alpha) \to \varepsilon \in \mathbb{R} \to 0 < \varepsilon \to c = 0 \to \varphi_{59}) \to (\lvert c \odot a(n) - c \cdot \alpha \rvert = \lvert c \cdot (a_{n} - \alpha) \rvert \to \varphi_{147}) \to \varphi_{193}$
\item $\varphi_{300} :\equiv (\varphi_{177} \to \forall n\,(\varphi_{177})) \to c \in \mathbb{R} \to \mathrm{Conv}(a,\alpha) \to \varepsilon \in \mathbb{R} \to 0 < \varepsilon \to c \neq 0 \to \varphi_{188}$
\item $\varphi_{301} :\equiv (c \in \mathbb{R} \to \mathrm{Conv}(a,\alpha) \to \varepsilon \in \mathbb{R} \to 0 < \varepsilon \to c = 0 \to \varphi_{58}) \to \varphi_{262}$
\item $\varphi_{302} :\equiv (\varepsilon \in \mathbb{R} \to \lvert c \rvert^{-1} \in \mathbb{R} \to \varepsilon \cdot \lvert c \rvert^{-1} \in \mathbb{R}) \to (c \in \mathbb{R} \to \mathrm{Conv}(a,\alpha) \to \varepsilon \in \mathbb{R} \to \varepsilon \in \mathbb{R}) \to \varphi_{265}$
\item $\varphi_{303} :\equiv (c \in \mathbb{R} \to \mathrm{Conv}(a,\alpha) \to \varepsilon \in \mathbb{R} \to 0 < \varepsilon \to \varphi_{41}) \to \lvert c \cdot a_{n} - c \cdot \alpha \rvert = \lvert c \cdot a_{n} - c \cdot \alpha \rvert \to \varphi_{250}$
\item $\varphi_{304} :\equiv (\mathrm{Conv}(a,\alpha) \to \varepsilon \cdot \lvert c \rvert^{-1} \in \mathbb{R} \to \varphi_{124}) \to (c \in \mathbb{R} \to \mathrm{Conv}(a,\alpha) \to \mathrm{Conv}(a,\alpha)) \to \varphi_{253}$
\item $\varphi_{305} :\equiv (c \in \mathbb{R} \to \mathrm{Conv}(a,\alpha) \to \varepsilon \in \mathbb{R} \to 0 < \varepsilon \to c \neq 0 \to \varphi_{190}) \to \varphi_{256}$
\item $\varphi_{306} :\equiv (c \in \mathbb{R} \to \mathrm{Conv}(a,\alpha) \to \varepsilon \in \mathbb{R} \to 0 < \varepsilon \to c \neq 0 \to \varphi_{171}) \to \varphi_{258}$
\item $\varphi_{307} :\equiv (c \in \mathbb{R} \to \mathrm{Conv}(a,\alpha) \to \varepsilon \in \mathbb{R} \to 0 < \varepsilon \to c \neq 0 \to \varphi_{173}) \to \varphi_{260}$
\item $\varphi_{308} :\equiv (c \in \mathbb{R} \to \mathrm{Conv}(a,\alpha) \to \varepsilon \in \mathbb{R} \to 0 < \varepsilon \to c \neq 0 \to \varphi_{172}) \to \varphi_{263}$
\item $\varphi_{309} :\equiv (c \in \mathbb{R} \to \mathrm{Conv}(a,\alpha) \to \varepsilon \in \mathbb{R} \to 0 < \varepsilon \to c \neq 0 \to \varphi_{136}) \to \varphi_{270}$
\item $\varphi_{310} :\equiv (c \in \mathbb{R} \to \mathrm{Conv}(a,\alpha) \to \varepsilon \in \mathbb{R} \to 0 < \varepsilon \to c \neq 0 \to \varphi_{207}) \to \varphi_{257}$
\item $\varphi_{311} :\equiv (c \in \mathbb{R} \to \mathrm{Conv}(a,\alpha) \to \varepsilon \in \mathbb{R} \to 0 < \varepsilon \to c \neq 0 \to \varphi_{125}) \to \varphi_{277}$
\item $\varphi_{312} :\equiv (c \in \mathbb{R} \to \mathrm{Conv}(a,\alpha) \to \varepsilon \in \mathbb{R} \to 0 < \varepsilon \to c \neq 0 \to \varphi_{94}) \to \varphi_{285}$
\item $\varphi_{313} :\equiv (c \in \mathbb{R} \to \mathrm{Conv}(a,\alpha) \to \varepsilon \in \mathbb{R} \to 0 < \varepsilon \to c \neq 0 \to \varphi_{146}) \to \varepsilon \cdot (\lvert c \rvert \cdot \lvert c \rvert^{-1}) = \varepsilon \cdot (\lvert c \rvert \cdot \lvert c \rvert^{-1}) \to \varphi_{255}$
\item $\varphi_{314} :\equiv (c \in \mathbb{R} \to \mathrm{Conv}(a,\alpha) \to \varepsilon \in \mathbb{R} \to 0 < \varepsilon \to c \neq 0 \to \varphi_{136}) \to \varphi_{275}$
\item $\varphi_{315} :\equiv (c \in \mathbb{R} \to \mathrm{Conv}(a,\alpha) \to \varepsilon \in \mathbb{R} \to 0 < \varepsilon \to \varphi_{215}) \to \varphi_{163} \to c \in \mathbb{R} \to \mathrm{Conv}(a,\alpha) \to \varepsilon \in \mathbb{R} \to 0 < \varepsilon \to \varphi_{216}$
\item $\varphi_{316} :\equiv (c \in \mathbb{R} \to \mathrm{Conv}(a,\alpha) \to \varepsilon \in \mathbb{R} \to 0 < \varepsilon \to c = 0 \to \varphi_{57}) \to (c \in \mathbb{R} \to \mathrm{Conv}(a,\alpha) \to \varepsilon \in \mathbb{R} \to 0 < \varepsilon \to \varphi_{66}) \to \varphi_{249}$
\item $\varphi_{317} :\equiv (c \in \mathbb{R} \to \mathrm{Conv}(a,\alpha) \to \varepsilon \in \mathbb{R} \to 0 < \varepsilon \to c \neq 0 \to \varphi_{153}) \to \varphi_{276}$
\item $\varphi_{318} :\equiv (c \in \mathbb{R} \to \mathrm{Conv}(a,\alpha) \to \varepsilon \in \mathbb{R} \to 0 < \varepsilon \to c \neq 0 \to \varphi_{201}) \to \varphi_{269}$
\item $\varphi_{319} :\equiv (c \in \mathbb{R} \to \mathrm{Conv}(a,\alpha) \to \varepsilon \in \mathbb{R} \to 0 < \varepsilon \to c \neq 0 \to \varphi_{125}) \to \varphi_{287}$
\item $\varphi_{320} :\equiv \forall \varepsilon\,(c \in \mathbb{R} \land \mathrm{Conv}(a,\alpha) \to \varepsilon \in \mathbb{R} \to \varphi_{95}) \to \varphi_{245} \to \varphi_{4} \to c \in \mathbb{R} \to \mathrm{Conv}(a,\alpha) \to \forall \varepsilon \in \mathbb{R}\,(\varphi_{95})$
\item $\varphi_{321} :\equiv (c \in \mathbb{R} \to \mathrm{Conv}(a,\alpha) \to \varepsilon \in \mathbb{R} \to 0 < \varepsilon \to c \neq 0 \to \varphi_{200}) \to \varphi_{281}$
\item $\varphi_{322} :\equiv (c \in \mathbb{R} \to \mathrm{Conv}(a,\alpha) \to \varepsilon \in \mathbb{R} \to 0 < \varepsilon \to c \neq 0 \to \varphi_{199}) \to \varphi_{283}$
\item $\varphi_{323} :\equiv (\lvert c \rvert^{-1} \in \mathbb{R} \to \lvert c \rvert \in \mathbb{R} \to \lvert c \rvert^{-1} \cdot \lvert c \rvert = \lvert c \rvert \cdot \lvert c \rvert^{-1}) \to \varphi_{309}$
\item $\varphi_{324} :\equiv (c \in \mathbb{R} \to \mathrm{Conv}(a,\alpha) \to \varepsilon \in \mathbb{R} \to 0 < \varepsilon \to c \neq 0 \to \varphi_{189}) \to \varepsilon \cdot (\lvert c \rvert^{-1} \cdot \lvert c \rvert) = \varepsilon \cdot (\lvert c \rvert^{-1} \cdot \lvert c \rvert) \to \varphi_{273}$
\item $\varphi_{325} :\equiv (c \in \mathbb{R} \to \mathrm{Conv}(a,\alpha) \to \varepsilon \in \mathbb{R} \to 0 < \varepsilon \to c \neq 0 \to \varphi_{191}) \to \lvert c \cdot a_{n} - c \cdot \alpha \rvert = \lvert c \cdot a_{n} - c \cdot \alpha \rvert \to \varphi_{271}$
\item $\varphi_{326} :\equiv (c \in \mathbb{R} \to \mathrm{Conv}(a,\alpha) \to \varepsilon \in \mathbb{R} \to 0 < \varepsilon \to c \neq 0 \to \varphi_{157}) \to \lvert c \odot a(n) - c \cdot \alpha \rvert = \lvert c \odot a(n) - c \cdot \alpha \rvert \to \varphi_{274}$
\item $\varphi_{327} :\equiv (c \in \mathbb{R} \to \mathrm{Conv}(a,\alpha) \to \varepsilon \in \mathbb{R} \to 0 < \varepsilon \to c \neq 0 \to \varphi_{174}) \to \varphi_{295}$
\item $\varphi_{328} :\equiv (c \in \mathbb{R} \to \mathrm{Conv}(a,\alpha) \to \varepsilon \in \mathbb{R} \to 0 < \varepsilon \to c \neq 0 \to \varphi_{132}) \to \varphi_{34} \to \varphi_{288}$
\item $\varphi_{329} :\equiv (\varepsilon \in \mathbb{R} \to \varphi_{118}) \to (c \in \mathbb{R} \to \mathrm{Conv}(a,\alpha) \to \varepsilon \in \mathbb{R} \to \varepsilon \in \mathbb{R}) \to \varphi_{314}$
\item $\varphi_{330} :\equiv (c \in \mathbb{R} \to \mathrm{Conv}(a,\alpha) \to \varepsilon \in \mathbb{R} \to 0 < \varepsilon \to c \neq 0 \to \varphi_{164}) \to \varphi_{310}$
\item $\varphi_{331} :\equiv (c \in \mathbb{R} \to \mathrm{Conv}(a,\alpha) \to \varepsilon \in \mathbb{R} \to 0 < \varepsilon \to c \neq 0 \to \varphi_{204}) \to \varphi_{307}$
\item $\varphi_{332} :\equiv (c \in \mathbb{R} \to \mathrm{Conv}(a,\alpha) \to \varepsilon \in \mathbb{R} \to 0 < \varepsilon \to c \neq 0 \to \varphi_{206}) \to \varphi_{306}$
\item $\varphi_{333} :\equiv (c \in \mathbb{R} \to \mathrm{Conv}(a,\alpha) \to \varepsilon \in \mathbb{R} \to 0 < \varepsilon \to c \neq 0 \to \varphi_{203}) \to (c \in \mathbb{R} \to \mathrm{Conv}(a,\alpha) \to \varepsilon \in \mathbb{R} \to 0 < \varepsilon \to \varphi_{211}) \to \varphi_{280}$
\item $\varphi_{334} :\equiv (c \in \mathbb{R} \to \mathrm{Conv}(a,\alpha) \to \varepsilon \in \mathbb{R} \to 0 < \varepsilon \to c \neq 0 \to \varphi_{152}) \to \varphi_{328}$
\item $\varphi_{335} :\equiv (c \in \mathbb{R} \to \mathrm{Conv}(a,\alpha) \to \varepsilon \in \mathbb{R} \to 0 < \varepsilon \to c \neq 0 \to \varphi_{153}) \to \varphi_{334}$
\item $\varphi_{336} :\equiv \forall n\,(\varphi_{177} \to n \in \mathbb{N} \to N \leq n \to \lvert c \odot a(n) - c \cdot \alpha \rvert < \varepsilon) \to \varphi_{282} \to \varphi_{300}$
\item $\varphi_{337} :\equiv \forall N\,(\varphi_{180} \to \varphi_{56}) \to (\forall N\,(\varphi_{180} \to \varphi_{56}) \to \forall N\,(\varphi_{180}) \to \forall N\,(\varphi_{56})) \to \varphi_{289}$
\end{abbrevlist}
\end{fullproof}

% 自動生成 (prover/build.js): ErrSqueeze  ノード数 321
\begin{fullproof}{[C12]}{$\forall x\,\forall y\,\forall z\,\forall \alpha\,\forall \varepsilon\,(x \in \mathbb{R} \to y \in \mathbb{R} \to z \in \mathbb{R} \to \alpha \in \mathbb{R} \to \varepsilon \in \mathbb{R} \to x \leq z \to z \leq y \to \lvert x - \alpha \rvert < \varepsilon \to \lvert y - \alpha \rvert < \varepsilon \to \lvert z - \alpha \rvert < \varepsilon)$}{321}\phantomsection\label{tree:ErrSqueeze}
\noindent$\Pi_{1}$:
\begin{ptree}
\AXC{$\forall x\,\forall y\,(x \in \mathbb{R} \to y \in \mathbb{R} \to x - y \in \mathbb{R})$\lbl{[F1]}}
\AXC{$\varphi_{71}$\lbl{A2}}
\BIC{$\forall y\,(v_{391} \in \mathbb{R} \to y \in \mathbb{R} \to v_{391} - y \in \mathbb{R})$}
\AXC{$\varphi_{75}$\lbl{A2}}
\BIC{$v_{391} \in \mathbb{R} \to v_{392} \in \mathbb{R} \to v_{391} - v_{392} \in \mathbb{R}$}
\UIC{$\forall v_{392}\,(v_{391} \in \mathbb{R} \to v_{392} \in \mathbb{R} \to v_{391} - v_{392} \in \mathbb{R})$}
\UIC{$\varphi_{6}$}
\AXC{$\varphi_{6} \to \forall v_{392}\,(y \in \mathbb{R} \to v_{392} \in \mathbb{R} \to y - v_{392} \in \mathbb{R})$\lbl{A2}}
\BIC{$\forall v_{392}\,(y \in \mathbb{R} \to v_{392} \in \mathbb{R} \to y - v_{392} \in \mathbb{R})$}
\end{ptree}
\noindent$\Pi_{2}$:
\begin{ptree}
\AXC{$x \in \mathbb{R} \to y \in \mathbb{R} \to z \in \mathbb{R} \to \alpha \in \mathbb{R} \to \varphi_{36}$\lbl{A1}}
\AXC{$\stackrel{\Pi_{1}}{\vdots}$}\noLine
\UIC{$\forall v_{392}\,(y \in \mathbb{R} \to v_{392} \in \mathbb{R} \to y - v_{392} \in \mathbb{R})$}
\AXC{$\varphi_{70}$\lbl{A2}}
\BIC{$y \in \mathbb{R} \to \alpha \in \mathbb{R} \to y - \alpha \in \mathbb{R}$}
\AXC{$\varphi_{118}$\lbl{A1}}
\BIC{$(x \in \mathbb{R} \to y \in \mathbb{R} \to z \in \mathbb{R} \to \alpha \in \mathbb{R} \to \varphi_{36}) \to \varphi_{101}$}
\BIC{$\varphi_{101}$}
\end{ptree}
\noindent$\Pi_{3}$:
\begin{ptree}
\AXC{$\forall x\,\forall y\,(x \in \mathbb{R} \to y \in \mathbb{R} \to \lvert x \rvert < y \to -y < x \land x < y)$\lbl{[A11]}}
\AXC{$\varphi_{89}$\lbl{A2}}
\BIC{$\varphi_{25}$}
\AXC{$\varphi_{25} \to v_{389} \in \mathbb{R} \to \varphi_{9}$\lbl{A2}}
\BIC{$v_{389} \in \mathbb{R} \to \varphi_{9}$}
\UIC{$\forall v_{390}\,(v_{389} \in \mathbb{R} \to \varphi_{9})$}
\UIC{$\forall v_{389}\,\forall v_{390}\,(v_{389} \in \mathbb{R} \to \varphi_{9})$}
\end{ptree}
\noindent$\Pi_{4}$:
\begin{ptree}
\AXC{$\stackrel{\Pi_{3}}{\vdots}$}\noLine
\UIC{$\forall v_{389}\,\forall v_{390}\,(v_{389} \in \mathbb{R} \to \varphi_{9})$}
\AXC{$\forall v_{389}\,\forall v_{390}\,(v_{389} \in \mathbb{R} \to \varphi_{9}) \to \varphi_{72}$\lbl{A2}}
\BIC{$\varphi_{72}$}
\AXC{$\varphi_{72} \to \varphi_{64}$\lbl{A2}}
\BIC{$\varphi_{64}$}
\AXC{$\varphi_{64} \to \varphi_{128}$\lbl{A1}}
\BIC{$\varphi_{128}$}
\end{ptree}
\noindent$\Pi_{5}$:
\begin{ptree}
\AXC{$x \in \mathbb{R} \to y \in \mathbb{R} \to z \in \mathbb{R} \to \alpha \in \mathbb{R} \to \varphi_{42}$\lbl{A1}}
\AXC{$x \in \mathbb{R} \to y \in \mathbb{R} \to z \in \mathbb{R} \to \alpha \in \mathbb{R} \to \varphi_{41}$\lbl{A1}}
\AXC{$\stackrel{\Pi_{2}}{\vdots}$}\noLine
\UIC{$\varphi_{101}$}
\BIC{$x \in \mathbb{R} \to y \in \mathbb{R} \to z \in \mathbb{R} \to \alpha \in \mathbb{R} \to \varphi_{49}$}
\AXC{$\stackrel{\Pi_{4}}{\vdots}$}\noLine
\UIC{$\varphi_{128}$}
\BIC{$(x \in \mathbb{R} \to y \in \mathbb{R} \to z \in \mathbb{R} \to \alpha \in \mathbb{R} \to \varphi_{42}) \to \varphi_{113}$}
\BIC{$\varphi_{113}$}
\end{ptree}
\noindent$\Pi_{6}$:
\begin{ptree}
\AXC{$\forall x \in \mathbb{R}\,(-x \in \mathbb{R})$\lbl{(R1)}}
\AXC{$\forall x \in \mathbb{R}\,(-x \in \mathbb{R}) \to \alpha \in \mathbb{R} \to -\alpha \in \mathbb{R}$\lbl{A2}}
\BIC{$\alpha \in \mathbb{R} \to -\alpha \in \mathbb{R}$}
\AXC{$(\alpha \in \mathbb{R} \to -\alpha \in \mathbb{R}) \to \varphi_{102}$\lbl{A1}}
\BIC{$\varphi_{102}$}
\end{ptree}
\noindent$\Pi_{7}$:
\begin{ptree}
\AXC{$x \in \mathbb{R} \to y \in \mathbb{R} \to z \in \mathbb{R} \to \alpha \in \mathbb{R} \to \varphi_{41}$\lbl{A1}}
\AXC{$\stackrel{\Pi_{6}}{\vdots}$}\noLine
\UIC{$\varphi_{102}$}
\BIC{$x \in \mathbb{R} \to y \in \mathbb{R} \to z \in \mathbb{R} \to \alpha \in \mathbb{R} \to \varphi_{50}$}
\end{ptree}
\noindent$\Pi_{8}$:
\begin{ptree}
\AXC{$\varphi_{8}$\lbl{[F19]}}
\AXC{$\varphi_{8} \to \varphi_{18}$\lbl{A2}}
\BIC{$\varphi_{18}$}
\AXC{$\varphi_{18} \to \varphi_{33}$\lbl{A2}}
\BIC{$\varphi_{33}$}
\AXC{$\varphi_{33} \to v_{396} \in \mathbb{R} \to \varphi_{11}$\lbl{A2}}
\BIC{$v_{396} \in \mathbb{R} \to \varphi_{11}$}
\UIC{$\forall v_{398}\,(v_{396} \in \mathbb{R} \to \varphi_{11})$}
\UIC{$\forall v_{397}\,\forall v_{398}\,(v_{396} \in \mathbb{R} \to \varphi_{11})$}
\UIC{$\forall v_{396}\,\forall v_{397}\,\forall v_{398}\,(v_{396} \in \mathbb{R} \to \varphi_{11})$}
\AXC{$\forall v_{396}\,\forall v_{397}\,\forall v_{398}\,(v_{396} \in \mathbb{R} \to \varphi_{11}) \to \varphi_{53}$\lbl{A2}}
\BIC{$\varphi_{53}$}
\end{ptree}
\noindent$\Pi_{9}$:
\begin{ptree}
\AXC{$x \in \mathbb{R} \to y \in \mathbb{R} \to z \in \mathbb{R} \to \alpha \in \mathbb{R} \to \varphi_{37}$\lbl{A1}}
\AXC{$\stackrel{\Pi_{8}}{\vdots}$}\noLine
\UIC{$\varphi_{53}$}
\AXC{$\varphi_{53} \to \varphi_{19}$\lbl{A2}}
\BIC{$\varphi_{19}$}
\AXC{$\varphi_{92}$\lbl{A2}}
\BIC{$z \in \mathbb{R} \to y \in \mathbb{R} \to -\alpha \in \mathbb{R} \to z \leq y \to z + -\alpha \leq y + -\alpha$}
\AXC{$\varphi_{130}$\lbl{A1}}
\BIC{$(x \in \mathbb{R} \to y \in \mathbb{R} \to z \in \mathbb{R} \to \alpha \in \mathbb{R} \to \varphi_{37}) \to \varphi_{125}$}
\BIC{$\varphi_{125}$}
\end{ptree}
\noindent$\Pi_{10}$:
\begin{ptree}
\AXC{$\stackrel{\Pi_{7}}{\vdots}$}\noLine
\UIC{$x \in \mathbb{R} \to y \in \mathbb{R} \to z \in \mathbb{R} \to \alpha \in \mathbb{R} \to \varphi_{50}$}
\AXC{$x \in \mathbb{R} \to y \in \mathbb{R} \to z \in \mathbb{R} \to \alpha \in \mathbb{R} \to \varphi_{36}$\lbl{A1}}
\AXC{$\stackrel{\Pi_{9}}{\vdots}$}\noLine
\UIC{$\varphi_{125}$}
\BIC{$(x \in \mathbb{R} \to y \in \mathbb{R} \to z \in \mathbb{R} \to \alpha \in \mathbb{R} \to \varphi_{50}) \to \varphi_{105}$}
\BIC{$\varphi_{105}$}
\end{ptree}
\noindent$\Pi_{11}$:
\begin{ptree}
\AXC{$x = x$\lbl{A5}}
\AXC{$x = y \to x = x \to y = x$\lbl{A6}}
\AXC{$(x = y \to x = x \to y = x) \to x = x \to x = y \to y = x$\lbl{A1}}
\BIC{$x = x \to x = y \to y = x$}
\BIC{$x = y \to y = x$}
\UIC{$\forall y\,(x = y \to y = x)$}
\UIC{$\forall x\,\forall y\,(x = y \to y = x)$}
\end{ptree}
\noindent$\Pi_{12}$:
\begin{ptree}
\AXC{$\stackrel{\Pi_{11}}{\vdots}$}\noLine
\UIC{$\forall x\,\forall y\,(x = y \to y = x)$}
\AXC{$\forall x\,\forall y\,(x = y \to y = x) \to \forall y\,(z - \alpha = y \to y = z - \alpha)$\lbl{A2}}
\BIC{$\forall y\,(z - \alpha = y \to y = z - \alpha)$}
\AXC{$\varphi_{13}$\lbl{A2}}
\BIC{$z - \alpha = z + -\alpha \to z + -\alpha = z - \alpha$}
\end{ptree}
\noindent$\Pi_{13}$:
\begin{ptree}
\AXC{$\forall x\,\forall y\,(x - y = x + -y)$\lbl{Def$-$}}
\AXC{$\forall x\,\forall y\,(x - y = x + -y) \to \forall y\,(z - y = z + -y)$\lbl{A2}}
\BIC{$\forall y\,(z - y = z + -y)$}
\end{ptree}
\noindent$\Pi_{14}$:
\begin{ptree}
\AXC{$\stackrel{\Pi_{12}}{\vdots}$}\noLine
\UIC{$z - \alpha = z + -\alpha \to z + -\alpha = z - \alpha$}
\AXC{$\stackrel{\Pi_{13}}{\vdots}$}\noLine
\UIC{$\forall y\,(z - y = z + -y)$}
\AXC{$\forall y\,(z - y = z + -y) \to z - \alpha = z + -\alpha$\lbl{A2}}
\BIC{$z - \alpha = z + -\alpha$}
\AXC{$\varphi_{55}$\lbl{A1}}
\BIC{$(z - \alpha = z + -\alpha \to z + -\alpha = z - \alpha) \to z + -\alpha = z - \alpha$}
\BIC{$z + -\alpha = z - \alpha$}
\end{ptree}
\noindent$\Pi_{15}$:
\begin{ptree}
\AXC{$x \in \mathbb{R} \to y \in \mathbb{R} \to z \in \mathbb{R} \to \alpha \in \mathbb{R} \to \varphi_{23}$\lbl{A1}}
\AXC{$\stackrel{\Pi_{10}}{\vdots}$}\noLine
\UIC{$\varphi_{105}$}
\BIC{$x \in \mathbb{R} \to y \in \mathbb{R} \to z \in \mathbb{R} \to \alpha \in \mathbb{R} \to \varepsilon \in \mathbb{R} \to \varphi_{38}$}
\AXC{$\stackrel{\Pi_{14}}{\vdots}$}\noLine
\UIC{$z + -\alpha = z - \alpha$}
\AXC{$\varphi_{121}$\lbl{A1}}
\BIC{$(x \in \mathbb{R} \to y \in \mathbb{R} \to z \in \mathbb{R} \to \alpha \in \mathbb{R} \to \varepsilon \in \mathbb{R} \to \varphi_{38}) \to \varphi_{110}$}
\BIC{$\varphi_{110}$}
\end{ptree}
\noindent$\Pi_{16}$:
\begin{ptree}
\AXC{$z_{1} = w \to z_{1} < y + -\alpha \to w < y + -\alpha$\lbl{A6}}
\UIC{$\forall w\,(z_{1} = w \to z_{1} < y + -\alpha \to w < y + -\alpha)$}
\UIC{$\forall z_{1}\,\forall w\,(z_{1} = w \to z_{1} < y + -\alpha \to w < y + -\alpha)$}
\AXC{$\varphi_{83}$\lbl{A2}}
\BIC{$\forall w\,(z + -\alpha = w \to z + -\alpha < y + -\alpha \to w < y + -\alpha)$}
\AXC{$\varphi_{90}$\lbl{A2}}
\BIC{$z + -\alpha = z - \alpha \to z + -\alpha < y + -\alpha \to z - \alpha < y + -\alpha$}
\AXC{$\stackrel{\Pi_{15}}{\vdots}$}\noLine
\UIC{$\varphi_{110}$}
\BIC{$\varphi_{99}$}
\end{ptree}
\noindent$\Pi_{17}$:
\begin{ptree}
\AXC{$\stackrel{\Pi_{11}}{\vdots}$}\noLine
\UIC{$\forall x\,\forall y\,(x = y \to y = x)$}
\AXC{$\forall x\,\forall y\,(x = y \to y = x) \to \forall y\,(v_{401} = y \to y = v_{401})$\lbl{A2}}
\BIC{$\forall y\,(v_{401} = y \to y = v_{401})$}
\AXC{$\varphi_{1}$\lbl{A2}}
\BIC{$v_{401} = v_{402} \to v_{402} = v_{401}$}
\UIC{$\forall v_{402}\,(v_{401} = v_{402} \to v_{402} = v_{401})$}
\UIC{$\forall v_{401}\,\forall v_{402}\,(v_{401} = v_{402} \to v_{402} = v_{401})$}
\AXC{$\varphi_{10}$\lbl{A2}}
\BIC{$\forall v_{402}\,(y - \alpha = v_{402} \to v_{402} = y - \alpha)$}
\AXC{$\varphi_{34}$\lbl{A2}}
\BIC{$y - \alpha = y + -\alpha \to y + -\alpha = y - \alpha$}
\end{ptree}
\noindent$\Pi_{18}$:
\begin{ptree}
\AXC{$\forall x\,\forall y\,(x - y = x + -y)$\lbl{Def$-$}}
\AXC{$\varphi_{3}$\lbl{A2}}
\BIC{$\forall y\,(v_{399} - y = v_{399} + -y)$}
\AXC{$\varphi_{5}$\lbl{A2}}
\BIC{$v_{399} - v_{400} = v_{399} + -v_{400}$}
\UIC{$\forall v_{400}\,(v_{399} - v_{400} = v_{399} + -v_{400})$}
\UIC{$\forall v_{399}\,\forall v_{400}\,(v_{399} - v_{400} = v_{399} + -v_{400})$}
\AXC{$\varphi_{26}$\lbl{A2}}
\BIC{$\forall v_{400}\,(y - v_{400} = y + -v_{400})$}
\AXC{$\forall v_{400}\,(y - v_{400} = y + -v_{400}) \to y - \alpha = y + -\alpha$\lbl{A2}}
\BIC{$y - \alpha = y + -\alpha$}
\end{ptree}
\noindent$\Pi_{19}$:
\begin{ptree}
\AXC{$\stackrel{\Pi_{17}}{\vdots}$}\noLine
\UIC{$y - \alpha = y + -\alpha \to y + -\alpha = y - \alpha$}
\AXC{$\stackrel{\Pi_{18}}{\vdots}$}\noLine
\UIC{$y - \alpha = y + -\alpha$}
\AXC{$\varphi_{56}$\lbl{A1}}
\BIC{$(y - \alpha = y + -\alpha \to y + -\alpha = y - \alpha) \to y + -\alpha = y - \alpha$}
\BIC{$y + -\alpha = y - \alpha$}
\AXC{$y + -\alpha = y - \alpha \to \varphi_{116}$\lbl{A1}}
\BIC{$\varphi_{116}$}
\end{ptree}
\noindent$\Pi_{20}$:
\begin{ptree}
\AXC{$z_{1} = w \to z_{1} = y + -\alpha \to w = y + -\alpha$\lbl{A6}}
\UIC{$\forall w\,(z_{1} = w \to z_{1} = y + -\alpha \to w = y + -\alpha)$}
\UIC{$\forall z_{1}\,\forall w\,(z_{1} = w \to z_{1} = y + -\alpha \to w = y + -\alpha)$}
\AXC{$\varphi_{77}$\lbl{A2}}
\BIC{$\forall w\,(z + -\alpha = w \to z + -\alpha = y + -\alpha \to w = y + -\alpha)$}
\AXC{$\varphi_{87}$\lbl{A2}}
\BIC{$z + -\alpha = z - \alpha \to z + -\alpha = y + -\alpha \to z - \alpha = y + -\alpha$}
\AXC{$\stackrel{\Pi_{16}}{\vdots}$}\noLine
\UIC{$\varphi_{99}$}
\BIC{$x \in \mathbb{R} \to y \in \mathbb{R} \to z \in \mathbb{R} \to \alpha \in \mathbb{R} \to \varepsilon \in \mathbb{R} \to \varphi_{30}$}
\AXC{$\stackrel{\Pi_{19}}{\vdots}$}\noLine
\UIC{$\varphi_{116}$}
\BIC{$(y + -\alpha = y - \alpha \to z - \alpha < y + -\alpha \to z - \alpha < y - \alpha) \to \varphi_{97}$}
\end{ptree}
\noindent$\Pi_{21}$:
\begin{ptree}
\AXC{$z_{1} = w \to z - \alpha < z_{1} \to z - \alpha < w$\lbl{A6}}
\UIC{$\forall w\,(z_{1} = w \to z - \alpha < z_{1} \to z - \alpha < w)$}
\UIC{$\forall z_{1}\,\forall w\,(z_{1} = w \to z - \alpha < z_{1} \to z - \alpha < w)$}
\AXC{$\varphi_{73}$\lbl{A2}}
\BIC{$\forall w\,(y + -\alpha = w \to z - \alpha < y + -\alpha \to z - \alpha < w)$}
\AXC{$\varphi_{85}$\lbl{A2}}
\BIC{$y + -\alpha = y - \alpha \to z - \alpha < y + -\alpha \to z - \alpha < y - \alpha$}
\AXC{$\stackrel{\Pi_{20}}{\vdots}$}\noLine
\UIC{$(y + -\alpha = y - \alpha \to z - \alpha < y + -\alpha \to z - \alpha < y - \alpha) \to \varphi_{97}$}
\BIC{$\varphi_{97}$}
\end{ptree}
\noindent$\Pi_{22}$:
\begin{ptree}
\AXC{$z_{1} = w \to z - \alpha = z_{1} \to z - \alpha = w$\lbl{A6}}
\UIC{$\forall w\,(z_{1} = w \to z - \alpha = z_{1} \to z - \alpha = w)$}
\UIC{$\forall z_{1}\,\forall w\,(z_{1} = w \to z - \alpha = z_{1} \to z - \alpha = w)$}
\AXC{$\varphi_{62}$\lbl{A2}}
\BIC{$\forall w\,(y + -\alpha = w \to z - \alpha = y + -\alpha \to z - \alpha = w)$}
\AXC{$\varphi_{80}$\lbl{A2}}
\BIC{$y + -\alpha = y - \alpha \to z - \alpha = y + -\alpha \to z - \alpha = y - \alpha$}
\AXC{$\stackrel{\Pi_{21}}{\vdots}$}\noLine
\UIC{$\varphi_{97}$}
\BIC{$x \in \mathbb{R} \to y \in \mathbb{R} \to z \in \mathbb{R} \to \alpha \in \mathbb{R} \to \varepsilon \in \mathbb{R} \to \varphi_{15}$}
\end{ptree}
\noindent$\Pi_{23}$:
\begin{ptree}
\AXC{$\forall x\,\forall y\,\forall z\,(x \leq y \to y < z \to x < z)$\lbl{[O4]}}
\AXC{$\varphi_{7}$\lbl{A2}}
\BIC{$\forall y\,\forall z\,(v_{403} \leq y \to y < z \to v_{403} < z)$}
\AXC{$\varphi_{17}$\lbl{A2}}
\BIC{$\forall z\,(v_{403} \leq v_{404} \to v_{404} < z \to v_{403} < z)$}
\AXC{$\varphi_{32}$\lbl{A2}}
\BIC{$v_{403} \leq v_{404} \to v_{404} < v_{405} \to v_{403} < v_{405}$}
\UIC{$\forall v_{405}\,(v_{403} \leq v_{404} \to v_{404} < v_{405} \to v_{403} < v_{405})$}
\UIC{$\varphi_{2}$}
\UIC{$\forall v_{403}\,(\varphi_{2})$}
\AXC{$\forall v_{403}\,(\varphi_{2}) \to \varphi_{4}$\lbl{A2}}
\BIC{$\varphi_{4}$}
\AXC{$\varphi_{82}$\lbl{A2}}
\BIC{$\forall v_{405}\,(z - \alpha \leq y - \alpha \to y - \alpha < v_{405} \to z - \alpha < v_{405})$}
\end{ptree}
\noindent$\Pi_{24}$:
\begin{ptree}
\AXC{$\stackrel{\Pi_{23}}{\vdots}$}\noLine
\UIC{$\forall v_{405}\,(z - \alpha \leq y - \alpha \to y - \alpha < v_{405} \to z - \alpha < v_{405})$}
\AXC{$\varphi_{76}$\lbl{A2}}
\BIC{$z - \alpha \leq y - \alpha \to y - \alpha < \varepsilon \to z - \alpha < \varepsilon$}
\AXC{$\varphi_{122}$\lbl{A1}}
\BIC{$(x \in \mathbb{R} \to y \in \mathbb{R} \to z \in \mathbb{R} \to \alpha \in \mathbb{R} \to \varepsilon \in \mathbb{R} \to \varphi_{15}) \to \varphi_{106}$}
\end{ptree}
\noindent$\Pi_{25}$:
\begin{ptree}
\AXC{$\stackrel{\Pi_{22}}{\vdots}$}\noLine
\UIC{$x \in \mathbb{R} \to y \in \mathbb{R} \to z \in \mathbb{R} \to \alpha \in \mathbb{R} \to \varepsilon \in \mathbb{R} \to \varphi_{15}$}
\AXC{$\stackrel{\Pi_{24}}{\vdots}$}\noLine
\UIC{$(x \in \mathbb{R} \to y \in \mathbb{R} \to z \in \mathbb{R} \to \alpha \in \mathbb{R} \to \varepsilon \in \mathbb{R} \to \varphi_{15}) \to \varphi_{106}$}
\BIC{$\varphi_{106}$}
\end{ptree}
\noindent$\Pi_{26}$:
\begin{ptree}
\AXC{$\varphi_{8}$\lbl{[F19]}}
\AXC{$\varphi_{8} \to \varphi_{20}$\lbl{A2}}
\BIC{$\varphi_{20}$}
\AXC{$\varphi_{20} \to \varphi_{35}$\lbl{A2}}
\BIC{$\varphi_{35}$}
\AXC{$\varphi_{35} \to v_{393} \in \mathbb{R} \to \varphi_{12}$\lbl{A2}}
\BIC{$v_{393} \in \mathbb{R} \to \varphi_{12}$}
\UIC{$\forall v_{395}\,(v_{393} \in \mathbb{R} \to \varphi_{12})$}
\UIC{$\forall v_{394}\,\forall v_{395}\,(v_{393} \in \mathbb{R} \to \varphi_{12})$}
\UIC{$\forall v_{393}\,\forall v_{394}\,\forall v_{395}\,(v_{393} \in \mathbb{R} \to \varphi_{12})$}
\AXC{$\forall v_{393}\,\forall v_{394}\,\forall v_{395}\,(v_{393} \in \mathbb{R} \to \varphi_{12}) \to \varphi_{54}$\lbl{A2}}
\BIC{$\varphi_{54}$}
\end{ptree}
\noindent$\Pi_{27}$:
\begin{ptree}
\AXC{$x \in \mathbb{R} \to y \in \mathbb{R} \to z \in \mathbb{R} \to \alpha \in \mathbb{R} \to \varphi_{39}$\lbl{A1}}
\AXC{$\stackrel{\Pi_{26}}{\vdots}$}\noLine
\UIC{$\varphi_{54}$}
\AXC{$\varphi_{54} \to \varphi_{21}$\lbl{A2}}
\BIC{$\varphi_{21}$}
\AXC{$\varphi_{93}$\lbl{A2}}
\BIC{$x \in \mathbb{R} \to z \in \mathbb{R} \to -\alpha \in \mathbb{R} \to x \leq z \to x + -\alpha \leq z + -\alpha$}
\AXC{$\varphi_{131}$\lbl{A1}}
\BIC{$(x \in \mathbb{R} \to y \in \mathbb{R} \to z \in \mathbb{R} \to \alpha \in \mathbb{R} \to \varphi_{39}) \to \varphi_{126}$}
\BIC{$\varphi_{126}$}
\end{ptree}
\noindent$\Pi_{28}$:
\begin{ptree}
\AXC{$\stackrel{\Pi_{7}}{\vdots}$}\noLine
\UIC{$x \in \mathbb{R} \to y \in \mathbb{R} \to z \in \mathbb{R} \to \alpha \in \mathbb{R} \to \varphi_{50}$}
\AXC{$x \in \mathbb{R} \to y \in \mathbb{R} \to z \in \mathbb{R} \to \alpha \in \mathbb{R} \to \varphi_{37}$\lbl{A1}}
\AXC{$\stackrel{\Pi_{27}}{\vdots}$}\noLine
\UIC{$\varphi_{126}$}
\BIC{$(x \in \mathbb{R} \to y \in \mathbb{R} \to z \in \mathbb{R} \to \alpha \in \mathbb{R} \to \varphi_{50}) \to \varphi_{107}$}
\BIC{$\varphi_{107}$}
\end{ptree}
\noindent$\Pi_{29}$:
\begin{ptree}
\AXC{$\stackrel{\Pi_{11}}{\vdots}$}\noLine
\UIC{$\forall x\,\forall y\,(x = y \to y = x)$}
\AXC{$\forall x\,\forall y\,(x = y \to y = x) \to \forall y\,(x - \alpha = y \to y = x - \alpha)$\lbl{A2}}
\BIC{$\forall y\,(x - \alpha = y \to y = x - \alpha)$}
\AXC{$\varphi_{14}$\lbl{A2}}
\BIC{$x - \alpha = x + -\alpha \to x + -\alpha = x - \alpha$}
\end{ptree}
\noindent$\Pi_{30}$:
\begin{ptree}
\AXC{$\forall x\,\forall y\,(x - y = x + -y)$\lbl{Def$-$}}
\AXC{$\forall x\,\forall y\,(x - y = x + -y) \to \forall y\,(x - y = x + -y)$\lbl{A2}}
\BIC{$\forall y\,(x - y = x + -y)$}
\end{ptree}
\noindent$\Pi_{31}$:
\begin{ptree}
\AXC{$\stackrel{\Pi_{29}}{\vdots}$}\noLine
\UIC{$x - \alpha = x + -\alpha \to x + -\alpha = x - \alpha$}
\AXC{$\stackrel{\Pi_{30}}{\vdots}$}\noLine
\UIC{$\forall y\,(x - y = x + -y)$}
\AXC{$\forall y\,(x - y = x + -y) \to x - \alpha = x + -\alpha$\lbl{A2}}
\BIC{$x - \alpha = x + -\alpha$}
\AXC{$\varphi_{57}$\lbl{A1}}
\BIC{$(x - \alpha = x + -\alpha \to x + -\alpha = x - \alpha) \to x + -\alpha = x - \alpha$}
\BIC{$x + -\alpha = x - \alpha$}
\end{ptree}
\noindent$\Pi_{32}$:
\begin{ptree}
\AXC{$x \in \mathbb{R} \to y \in \mathbb{R} \to z \in \mathbb{R} \to \alpha \in \mathbb{R} \to \varphi_{24}$\lbl{A1}}
\AXC{$\stackrel{\Pi_{28}}{\vdots}$}\noLine
\UIC{$\varphi_{107}$}
\BIC{$x \in \mathbb{R} \to y \in \mathbb{R} \to z \in \mathbb{R} \to \alpha \in \mathbb{R} \to \varepsilon \in \mathbb{R} \to \varphi_{40}$}
\AXC{$\stackrel{\Pi_{31}}{\vdots}$}\noLine
\UIC{$x + -\alpha = x - \alpha$}
\AXC{$\varphi_{123}$\lbl{A1}}
\BIC{$(x \in \mathbb{R} \to y \in \mathbb{R} \to z \in \mathbb{R} \to \alpha \in \mathbb{R} \to \varepsilon \in \mathbb{R} \to \varphi_{40}) \to \varphi_{111}$}
\BIC{$\varphi_{111}$}
\end{ptree}
\noindent$\Pi_{33}$:
\begin{ptree}
\AXC{$z_{1} = w \to z_{1} < z + -\alpha \to w < z + -\alpha$\lbl{A6}}
\UIC{$\forall w\,(z_{1} = w \to z_{1} < z + -\alpha \to w < z + -\alpha)$}
\UIC{$\forall z_{1}\,\forall w\,(z_{1} = w \to z_{1} < z + -\alpha \to w < z + -\alpha)$}
\AXC{$\varphi_{84}$\lbl{A2}}
\BIC{$\forall w\,(x + -\alpha = w \to x + -\alpha < z + -\alpha \to w < z + -\alpha)$}
\AXC{$\varphi_{91}$\lbl{A2}}
\BIC{$x + -\alpha = x - \alpha \to x + -\alpha < z + -\alpha \to x - \alpha < z + -\alpha$}
\AXC{$\stackrel{\Pi_{32}}{\vdots}$}\noLine
\UIC{$\varphi_{111}$}
\BIC{$\varphi_{100}$}
\end{ptree}
\noindent$\Pi_{34}$:
\begin{ptree}
\AXC{$z_{1} = w \to z_{1} = z + -\alpha \to w = z + -\alpha$\lbl{A6}}
\UIC{$\forall w\,(z_{1} = w \to z_{1} = z + -\alpha \to w = z + -\alpha)$}
\UIC{$\forall z_{1}\,\forall w\,(z_{1} = w \to z_{1} = z + -\alpha \to w = z + -\alpha)$}
\AXC{$\varphi_{78}$\lbl{A2}}
\BIC{$\forall w\,(x + -\alpha = w \to x + -\alpha = z + -\alpha \to w = z + -\alpha)$}
\AXC{$\varphi_{88}$\lbl{A2}}
\BIC{$x + -\alpha = x - \alpha \to x + -\alpha = z + -\alpha \to x - \alpha = z + -\alpha$}
\AXC{$\stackrel{\Pi_{33}}{\vdots}$}\noLine
\UIC{$\varphi_{100}$}
\BIC{$x \in \mathbb{R} \to y \in \mathbb{R} \to z \in \mathbb{R} \to \alpha \in \mathbb{R} \to \varepsilon \in \mathbb{R} \to \varphi_{31}$}
\end{ptree}
\noindent$\Pi_{35}$:
\begin{ptree}
\AXC{$\stackrel{\Pi_{11}}{\vdots}$}\noLine
\UIC{$\forall x\,\forall y\,(x = y \to y = x)$}
\AXC{$\forall x\,\forall y\,(x = y \to y = x) \to \forall y\,(z - \alpha = y \to y = z - \alpha)$\lbl{A2}}
\BIC{$\forall y\,(z - \alpha = y \to y = z - \alpha)$}
\AXC{$\varphi_{13}$\lbl{A2}}
\BIC{$z - \alpha = z + -\alpha \to z + -\alpha = z - \alpha$}
\end{ptree}
\noindent$\Pi_{36}$:
\begin{ptree}
\AXC{$\forall x\,\forall y\,(x - y = x + -y)$\lbl{Def$-$}}
\AXC{$\forall x\,\forall y\,(x - y = x + -y) \to \forall y\,(z - y = z + -y)$\lbl{A2}}
\BIC{$\forall y\,(z - y = z + -y)$}
\end{ptree}
\noindent$\Pi_{37}$:
\begin{ptree}
\AXC{$\stackrel{\Pi_{35}}{\vdots}$}\noLine
\UIC{$z - \alpha = z + -\alpha \to z + -\alpha = z - \alpha$}
\AXC{$\stackrel{\Pi_{36}}{\vdots}$}\noLine
\UIC{$\forall y\,(z - y = z + -y)$}
\AXC{$\forall y\,(z - y = z + -y) \to z - \alpha = z + -\alpha$\lbl{A2}}
\BIC{$z - \alpha = z + -\alpha$}
\AXC{$\varphi_{55}$\lbl{A1}}
\BIC{$(z - \alpha = z + -\alpha \to z + -\alpha = z - \alpha) \to z + -\alpha = z - \alpha$}
\BIC{$z + -\alpha = z - \alpha$}
\end{ptree}
\noindent$\Pi_{38}$:
\begin{ptree}
\AXC{$\stackrel{\Pi_{34}}{\vdots}$}\noLine
\UIC{$x \in \mathbb{R} \to y \in \mathbb{R} \to z \in \mathbb{R} \to \alpha \in \mathbb{R} \to \varepsilon \in \mathbb{R} \to \varphi_{31}$}
\AXC{$\stackrel{\Pi_{37}}{\vdots}$}\noLine
\UIC{$z + -\alpha = z - \alpha$}
\AXC{$z + -\alpha = z - \alpha \to \varphi_{117}$\lbl{A1}}
\BIC{$\varphi_{117}$}
\BIC{$(z + -\alpha = z - \alpha \to x - \alpha < z + -\alpha \to x - \alpha < z - \alpha) \to \varphi_{98}$}
\end{ptree}
\noindent$\Pi_{39}$:
\begin{ptree}
\AXC{$z_{1} = w \to x - \alpha < z_{1} \to x - \alpha < w$\lbl{A6}}
\UIC{$\forall w\,(z_{1} = w \to x - \alpha < z_{1} \to x - \alpha < w)$}
\UIC{$\forall z_{1}\,\forall w\,(z_{1} = w \to x - \alpha < z_{1} \to x - \alpha < w)$}
\AXC{$\varphi_{74}$\lbl{A2}}
\BIC{$\forall w\,(z + -\alpha = w \to x - \alpha < z + -\alpha \to x - \alpha < w)$}
\AXC{$\varphi_{86}$\lbl{A2}}
\BIC{$z + -\alpha = z - \alpha \to x - \alpha < z + -\alpha \to x - \alpha < z - \alpha$}
\AXC{$\stackrel{\Pi_{38}}{\vdots}$}\noLine
\UIC{$(z + -\alpha = z - \alpha \to x - \alpha < z + -\alpha \to x - \alpha < z - \alpha) \to \varphi_{98}$}
\BIC{$\varphi_{98}$}
\end{ptree}
\noindent$\Pi_{40}$:
\begin{ptree}
\AXC{$\forall x\,\forall y\,(x \in \mathbb{R} \to y \in \mathbb{R} \to x - y \in \mathbb{R})$\lbl{[F1]}}
\AXC{$\varphi_{67}$\lbl{A2}}
\BIC{$\forall y\,(x \in \mathbb{R} \to y \in \mathbb{R} \to x - y \in \mathbb{R})$}
\AXC{$\varphi_{58}$\lbl{A2}}
\BIC{$x \in \mathbb{R} \to \alpha \in \mathbb{R} \to x - \alpha \in \mathbb{R}$}
\AXC{$\varphi_{119}$\lbl{A1}}
\BIC{$(x \in \mathbb{R} \to y \in \mathbb{R} \to z \in \mathbb{R} \to \alpha \in \mathbb{R} \to \varphi_{39}) \to \varphi_{103}$}
\end{ptree}
\noindent$\Pi_{41}$:
\begin{ptree}
\AXC{$x \in \mathbb{R} \to y \in \mathbb{R} \to z \in \mathbb{R} \to \alpha \in \mathbb{R} \to \varphi_{39}$\lbl{A1}}
\AXC{$\stackrel{\Pi_{40}}{\vdots}$}\noLine
\UIC{$(x \in \mathbb{R} \to y \in \mathbb{R} \to z \in \mathbb{R} \to \alpha \in \mathbb{R} \to \varphi_{39}) \to \varphi_{103}$}
\BIC{$\varphi_{103}$}
\end{ptree}
\noindent$\Pi_{42}$:
\begin{ptree}
\AXC{$\forall x\,\forall y\,(x \in \mathbb{R} \to y \in \mathbb{R} \to \lvert x \rvert < y \to -y < x \land x < y)$\lbl{[A11]}}
\AXC{$\varphi_{95}$\lbl{A2}}
\BIC{$\varphi_{47}$}
\AXC{$\varphi_{47} \to \varphi_{65}$\lbl{A2}}
\BIC{$\varphi_{65}$}
\AXC{$\varphi_{65} \to \varphi_{129}$\lbl{A1}}
\BIC{$\varphi_{129}$}
\end{ptree}
\noindent$\Pi_{43}$:
\begin{ptree}
\AXC{$x \in \mathbb{R} \to y \in \mathbb{R} \to z \in \mathbb{R} \to \alpha \in \mathbb{R} \to \varphi_{42}$\lbl{A1}}
\AXC{$x \in \mathbb{R} \to y \in \mathbb{R} \to z \in \mathbb{R} \to \alpha \in \mathbb{R} \to \varphi_{41}$\lbl{A1}}
\AXC{$\stackrel{\Pi_{41}}{\vdots}$}\noLine
\UIC{$\varphi_{103}$}
\BIC{$x \in \mathbb{R} \to y \in \mathbb{R} \to z \in \mathbb{R} \to \alpha \in \mathbb{R} \to \varphi_{51}$}
\AXC{$\stackrel{\Pi_{42}}{\vdots}$}\noLine
\UIC{$\varphi_{129}$}
\BIC{$(x \in \mathbb{R} \to y \in \mathbb{R} \to z \in \mathbb{R} \to \alpha \in \mathbb{R} \to \varphi_{42}) \to \varphi_{114}$}
\BIC{$\varphi_{114}$}
\end{ptree}
\noindent$\Pi_{44}$:
\begin{ptree}
\AXC{$\forall x\,\forall y\,\forall z\,(x < y \to y \leq z \to x < z)$\lbl{[O3]}}
\AXC{$\varphi_{22}$\lbl{A2}}
\BIC{$\forall y\,\forall z\,(-\varepsilon < y \to y \leq z \to -\varepsilon < z)$}
\AXC{$\varphi_{68}$\lbl{A2}}
\BIC{$\forall z\,(-\varepsilon < x - \alpha \to x - \alpha \leq z \to -\varepsilon < z)$}
\AXC{$\varphi_{79}$\lbl{A2}}
\BIC{$-\varepsilon < x - \alpha \to x - \alpha \leq z - \alpha \to -\varepsilon < z - \alpha$}
\AXC{$\varphi_{124}$\lbl{A1}}
\BIC{$(x \in \mathbb{R} \to y \in \mathbb{R} \to z \in \mathbb{R} \to \alpha \in \mathbb{R} \to \varphi_{60}) \to \varphi_{109}$}
\end{ptree}
\noindent$\Pi_{45}$:
\begin{ptree}
\AXC{$x \in \mathbb{R} \to y \in \mathbb{R} \to z \in \mathbb{R} \to \alpha \in \mathbb{R} \to \varepsilon \in \mathbb{R} \to \varphi_{28}$\lbl{A1}}
\AXC{$\stackrel{\Pi_{43}}{\vdots}$}\noLine
\UIC{$\varphi_{114}$}
\BIC{$x \in \mathbb{R} \to y \in \mathbb{R} \to z \in \mathbb{R} \to \alpha \in \mathbb{R} \to \varepsilon \in \mathbb{R} \to x \leq z \to \varphi_{46}$}
\AXC{$\varphi_{115}$\lbl{A1}}
\BIC{$x \in \mathbb{R} \to y \in \mathbb{R} \to z \in \mathbb{R} \to \alpha \in \mathbb{R} \to \varphi_{60}$}
\AXC{$\stackrel{\Pi_{44}}{\vdots}$}\noLine
\UIC{$(x \in \mathbb{R} \to y \in \mathbb{R} \to z \in \mathbb{R} \to \alpha \in \mathbb{R} \to \varphi_{60}) \to \varphi_{109}$}
\BIC{$\varphi_{109}$}
\end{ptree}
\noindent$\Pi_{46}$:
\begin{ptree}
\AXC{$z_{1} = w \to x - \alpha = z_{1} \to x - \alpha = w$\lbl{A6}}
\UIC{$\forall w\,(z_{1} = w \to x - \alpha = z_{1} \to x - \alpha = w)$}
\UIC{$\forall z_{1}\,\forall w\,(z_{1} = w \to x - \alpha = z_{1} \to x - \alpha = w)$}
\AXC{$\varphi_{63}$\lbl{A2}}
\BIC{$\forall w\,(z + -\alpha = w \to x - \alpha = z + -\alpha \to x - \alpha = w)$}
\AXC{$\varphi_{81}$\lbl{A2}}
\BIC{$z + -\alpha = z - \alpha \to x - \alpha = z + -\alpha \to x - \alpha = z - \alpha$}
\AXC{$\stackrel{\Pi_{39}}{\vdots}$}\noLine
\UIC{$\varphi_{98}$}
\BIC{$x \in \mathbb{R} \to y \in \mathbb{R} \to z \in \mathbb{R} \to \alpha \in \mathbb{R} \to \varepsilon \in \mathbb{R} \to \varphi_{16}$}
\AXC{$\stackrel{\Pi_{45}}{\vdots}$}\noLine
\UIC{$\varphi_{109}$}
\BIC{$x \in \mathbb{R} \to y \in \mathbb{R} \to z \in \mathbb{R} \to \alpha \in \mathbb{R} \to \varphi_{61}$}
\end{ptree}
\noindent$\Pi_{47}$:
\begin{ptree}
\AXC{$\forall x\,\forall y\,(x \in \mathbb{R} \to y \in \mathbb{R} \to x - y \in \mathbb{R})$\lbl{[F1]}}
\AXC{$\varphi_{69}$\lbl{A2}}
\BIC{$\forall y\,(z \in \mathbb{R} \to y \in \mathbb{R} \to z - y \in \mathbb{R})$}
\AXC{$\varphi_{59}$\lbl{A2}}
\BIC{$z \in \mathbb{R} \to \alpha \in \mathbb{R} \to z - \alpha \in \mathbb{R}$}
\AXC{$\varphi_{120}$\lbl{A1}}
\BIC{$(x \in \mathbb{R} \to y \in \mathbb{R} \to z \in \mathbb{R} \to \alpha \in \mathbb{R} \to \varphi_{37}) \to \varphi_{104}$}
\end{ptree}
\noindent$\Pi_{48}$:
\begin{ptree}
\AXC{$x \in \mathbb{R} \to y \in \mathbb{R} \to z \in \mathbb{R} \to \alpha \in \mathbb{R} \to \varphi_{37}$\lbl{A1}}
\AXC{$\stackrel{\Pi_{47}}{\vdots}$}\noLine
\UIC{$(x \in \mathbb{R} \to y \in \mathbb{R} \to z \in \mathbb{R} \to \alpha \in \mathbb{R} \to \varphi_{37}) \to \varphi_{104}$}
\BIC{$\varphi_{104}$}
\end{ptree}
\noindent$\Pi_{49}$:
\begin{ptree}
\AXC{$\forall x\,\forall y\,(x \in \mathbb{R} \to y \in \mathbb{R} \to -y < x \to x < y \to \lvert x \rvert < y)$\lbl{[A12]}}
\AXC{$\varphi_{96}$\lbl{A2}}
\BIC{$\varphi_{48}$}
\AXC{$\varphi_{48} \to \varphi_{66}$\lbl{A2}}
\BIC{$\varphi_{66}$}
\end{ptree}
\noindent$\Pi_{50}$:
\begin{ptree}
\AXC{$x \in \mathbb{R} \to y \in \mathbb{R} \to z \in \mathbb{R} \to \alpha \in \mathbb{R} \to \varphi_{41}$\lbl{A1}}
\AXC{$\stackrel{\Pi_{48}}{\vdots}$}\noLine
\UIC{$\varphi_{104}$}
\BIC{$x \in \mathbb{R} \to y \in \mathbb{R} \to z \in \mathbb{R} \to \alpha \in \mathbb{R} \to \varphi_{52}$}
\AXC{$\stackrel{\Pi_{49}}{\vdots}$}\noLine
\UIC{$\varphi_{66}$}
\AXC{$\varphi_{66} \to (x \in \mathbb{R} \to y \in \mathbb{R} \to z \in \mathbb{R} \to \alpha \in \mathbb{R} \to \varphi_{52}) \to \varphi_{127}$\lbl{A1}}
\BIC{$(x \in \mathbb{R} \to y \in \mathbb{R} \to z \in \mathbb{R} \to \alpha \in \mathbb{R} \to \varphi_{52}) \to \varphi_{127}$}
\BIC{$\varphi_{127}$}
\end{ptree}
\noindent$\Pi_{51}$:
\begin{ptree}
\AXC{$\stackrel{\Pi_{46}}{\vdots}$}\noLine
\UIC{$x \in \mathbb{R} \to y \in \mathbb{R} \to z \in \mathbb{R} \to \alpha \in \mathbb{R} \to \varphi_{61}$}
\AXC{$x \in \mathbb{R} \to y \in \mathbb{R} \to z \in \mathbb{R} \to \alpha \in \mathbb{R} \to \varphi_{42}$\lbl{A1}}
\AXC{$\stackrel{\Pi_{50}}{\vdots}$}\noLine
\UIC{$\varphi_{127}$}
\BIC{$(x \in \mathbb{R} \to y \in \mathbb{R} \to z \in \mathbb{R} \to \alpha \in \mathbb{R} \to \varphi_{61}) \to \varphi_{108}$}
\BIC{$\varphi_{108}$}
\end{ptree}
\noindent 主木:
\begin{ptree}
\AXC{$x \in \mathbb{R} \to y \in \mathbb{R} \to z \in \mathbb{R} \to \alpha \in \mathbb{R} \to \varepsilon \in \mathbb{R} \to \varphi_{27}$\lbl{A1}}
\AXC{$\stackrel{\Pi_{5}}{\vdots}$}\noLine
\UIC{$\varphi_{113}$}
\BIC{$x \in \mathbb{R} \to y \in \mathbb{R} \to z \in \mathbb{R} \to \alpha \in \mathbb{R} \to \varepsilon \in \mathbb{R} \to x \leq z \to \varphi_{45}$}
\AXC{$\varphi_{112}$\lbl{A1}}
\BIC{$x \in \mathbb{R} \to y \in \mathbb{R} \to z \in \mathbb{R} \to \alpha \in \mathbb{R} \to \varphi_{43}$}
\AXC{$\stackrel{\Pi_{25}}{\vdots}$}\noLine
\UIC{$\varphi_{106}$}
\BIC{$x \in \mathbb{R} \to y \in \mathbb{R} \to z \in \mathbb{R} \to \alpha \in \mathbb{R} \to \varphi_{44}$}
\AXC{$\stackrel{\Pi_{51}}{\vdots}$}\noLine
\UIC{$\varphi_{108}$}
\BIC{$x \in \mathbb{R} \to y \in \mathbb{R} \to z \in \mathbb{R} \to \alpha \in \mathbb{R} \to \varepsilon \in \mathbb{R} \to \varphi_{29}$}
\UIC{$\forall \varepsilon\,(x \in \mathbb{R} \to y \in \mathbb{R} \to z \in \mathbb{R} \to \alpha \in \mathbb{R} \to \varepsilon \in \mathbb{R} \to \varphi_{29})$}
\UIC{$\varphi_{94}$}
\UIC{$\forall z\,(\varphi_{94})$}
\UIC{$\forall y\,\forall z\,(\varphi_{94})$}
\UIC{$\forall x\,\forall y\,\forall z\,(\varphi_{94})$}
\end{ptree}
\begin{abbrevlist}
\item $\varphi_{1} :\equiv \forall y\,(v_{401} = y \to y = v_{401}) \to v_{401} = v_{402} \to v_{402} = v_{401}$
\item $\varphi_{2} :\equiv \forall v_{404}\,\forall v_{405}\,(v_{403} \leq v_{404} \to v_{404} < v_{405} \to v_{403} < v_{405})$
\item $\varphi_{3} :\equiv \forall x\,\forall y\,(x - y = x + -y) \to \forall y\,(v_{399} - y = v_{399} + -y)$
\item $\varphi_{4} :\equiv \forall v_{404}\,\forall v_{405}\,(z - \alpha \leq v_{404} \to v_{404} < v_{405} \to z - \alpha < v_{405})$
\item $\varphi_{5} :\equiv \forall y\,(v_{399} - y = v_{399} + -y) \to v_{399} - v_{400} = v_{399} + -v_{400}$
\item $\varphi_{6} :\equiv \forall v_{391}\,\forall v_{392}\,(v_{391} \in \mathbb{R} \to v_{392} \in \mathbb{R} \to v_{391} - v_{392} \in \mathbb{R})$
\item $\varphi_{7} :\equiv \forall x\,\forall y\,\forall z\,(x \leq y \to y < z \to x < z) \to \forall y\,\forall z\,(v_{403} \leq y \to y < z \to v_{403} < z)$
\item $\varphi_{8} :\equiv \forall x\,\forall y\,\forall z\,(x \in \mathbb{R} \to y \in \mathbb{R} \to z \in \mathbb{R} \to x \leq y \to x + z \leq y + z)$
\item $\varphi_{9} :\equiv v_{390} \in \mathbb{R} \to \lvert v_{389} \rvert < v_{390} \to -v_{390} < v_{389} \land v_{389} < v_{390}$
\item $\varphi_{10} :\equiv \forall v_{401}\,\forall v_{402}\,(v_{401} = v_{402} \to v_{402} = v_{401}) \to \forall v_{402}\,(y - \alpha = v_{402} \to v_{402} = y - \alpha)$
\item $\varphi_{11} :\equiv v_{397} \in \mathbb{R} \to v_{398} \in \mathbb{R} \to v_{396} \leq v_{397} \to v_{396} + v_{398} \leq v_{397} + v_{398}$
\item $\varphi_{12} :\equiv v_{394} \in \mathbb{R} \to v_{395} \in \mathbb{R} \to v_{393} \leq v_{394} \to v_{393} + v_{395} \leq v_{394} + v_{395}$
\item $\varphi_{13} :\equiv \forall y\,(z - \alpha = y \to y = z - \alpha) \to z - \alpha = z + -\alpha \to z + -\alpha = z - \alpha$
\item $\varphi_{14} :\equiv \forall y\,(x - \alpha = y \to y = x - \alpha) \to x - \alpha = x + -\alpha \to x + -\alpha = x - \alpha$
\item $\varphi_{15} :\equiv x \leq z \to z \leq y \to \lvert x - \alpha \rvert < \varepsilon \to \lvert y - \alpha \rvert < \varepsilon \to z - \alpha \leq y - \alpha$
\item $\varphi_{16} :\equiv x \leq z \to z \leq y \to \lvert x - \alpha \rvert < \varepsilon \to \lvert y - \alpha \rvert < \varepsilon \to x - \alpha \leq z - \alpha$
\item $\varphi_{17} :\equiv \forall y\,\forall z\,(v_{403} \leq y \to y < z \to v_{403} < z) \to \forall z\,(v_{403} \leq v_{404} \to v_{404} < z \to v_{403} < z)$
\item $\varphi_{18} :\equiv \forall y\,\forall z\,(v_{396} \in \mathbb{R} \to y \in \mathbb{R} \to z \in \mathbb{R} \to v_{396} \leq y \to v_{396} + z \leq y + z)$
\item $\varphi_{19} :\equiv \forall v_{398}\,(z \in \mathbb{R} \to y \in \mathbb{R} \to v_{398} \in \mathbb{R} \to z \leq y \to z + v_{398} \leq y + v_{398})$
\item $\varphi_{20} :\equiv \forall y\,\forall z\,(v_{393} \in \mathbb{R} \to y \in \mathbb{R} \to z \in \mathbb{R} \to v_{393} \leq y \to v_{393} + z \leq y + z)$
\item $\varphi_{21} :\equiv \forall v_{395}\,(x \in \mathbb{R} \to z \in \mathbb{R} \to v_{395} \in \mathbb{R} \to x \leq z \to x + v_{395} \leq z + v_{395})$
\item $\varphi_{22} :\equiv \forall x\,\forall y\,\forall z\,(x < y \to y \leq z \to x < z) \to \forall y\,\forall z\,(-\varepsilon < y \to y \leq z \to -\varepsilon < z)$
\item $\varphi_{23} :\equiv \varepsilon \in \mathbb{R} \to x \leq z \to z \leq y \to \lvert x - \alpha \rvert < \varepsilon \to \lvert y - \alpha \rvert < \varepsilon \to z \leq y$
\item $\varphi_{24} :\equiv \varepsilon \in \mathbb{R} \to x \leq z \to z \leq y \to \lvert x - \alpha \rvert < \varepsilon \to \lvert y - \alpha \rvert < \varepsilon \to x \leq z$
\item $\varphi_{25} :\equiv \forall y\,(v_{389} \in \mathbb{R} \to y \in \mathbb{R} \to \lvert v_{389} \rvert < y \to -y < v_{389} \land v_{389} < y)$
\item $\varphi_{26} :\equiv \forall v_{399}\,\forall v_{400}\,(v_{399} - v_{400} = v_{399} + -v_{400}) \to \forall v_{400}\,(y - v_{400} = y + -v_{400})$
\item $\varphi_{27} :\equiv x \leq z \to z \leq y \to \lvert x - \alpha \rvert < \varepsilon \to \lvert y - \alpha \rvert < \varepsilon \to \lvert y - \alpha \rvert < \varepsilon$
\item $\varphi_{28} :\equiv x \leq z \to z \leq y \to \lvert x - \alpha \rvert < \varepsilon \to \lvert y - \alpha \rvert < \varepsilon \to \lvert x - \alpha \rvert < \varepsilon$
\item $\varphi_{29} :\equiv x \leq z \to z \leq y \to \lvert x - \alpha \rvert < \varepsilon \to \lvert y - \alpha \rvert < \varepsilon \to \lvert z - \alpha \rvert < \varepsilon$
\item $\varphi_{30} :\equiv x \leq z \to z \leq y \to \lvert x - \alpha \rvert < \varepsilon \to \lvert y - \alpha \rvert < \varepsilon \to z - \alpha \leq y + -\alpha$
\item $\varphi_{31} :\equiv x \leq z \to z \leq y \to \lvert x - \alpha \rvert < \varepsilon \to \lvert y - \alpha \rvert < \varepsilon \to x - \alpha \leq z + -\alpha$
\item $\varphi_{32} :\equiv \forall z\,(v_{403} \leq v_{404} \to v_{404} < z \to v_{403} < z) \to v_{403} \leq v_{404} \to v_{404} < v_{405} \to v_{403} < v_{405}$
\item $\varphi_{33} :\equiv \forall z\,(v_{396} \in \mathbb{R} \to v_{397} \in \mathbb{R} \to z \in \mathbb{R} \to v_{396} \leq v_{397} \to v_{396} + z \leq v_{397} + z)$
\item $\varphi_{34} :\equiv \forall v_{402}\,(y - \alpha = v_{402} \to v_{402} = y - \alpha) \to y - \alpha = y + -\alpha \to y + -\alpha = y - \alpha$
\item $\varphi_{35} :\equiv \forall z\,(v_{393} \in \mathbb{R} \to v_{394} \in \mathbb{R} \to z \in \mathbb{R} \to v_{393} \leq v_{394} \to v_{393} + z \leq v_{394} + z)$
\item $\varphi_{36} :\equiv \varepsilon \in \mathbb{R} \to x \leq z \to z \leq y \to \lvert x - \alpha \rvert < \varepsilon \to \lvert y - \alpha \rvert < \varepsilon \to y \in \mathbb{R}$
\item $\varphi_{37} :\equiv \varepsilon \in \mathbb{R} \to x \leq z \to z \leq y \to \lvert x - \alpha \rvert < \varepsilon \to \lvert y - \alpha \rvert < \varepsilon \to z \in \mathbb{R}$
\item $\varphi_{38} :\equiv x \leq z \to z \leq y \to \lvert x - \alpha \rvert < \varepsilon \to \lvert y - \alpha \rvert < \varepsilon \to z + -\alpha \leq y + -\alpha$
\item $\varphi_{39} :\equiv \varepsilon \in \mathbb{R} \to x \leq z \to z \leq y \to \lvert x - \alpha \rvert < \varepsilon \to \lvert y - \alpha \rvert < \varepsilon \to x \in \mathbb{R}$
\item $\varphi_{40} :\equiv x \leq z \to z \leq y \to \lvert x - \alpha \rvert < \varepsilon \to \lvert y - \alpha \rvert < \varepsilon \to x + -\alpha \leq z + -\alpha$
\item $\varphi_{41} :\equiv \varepsilon \in \mathbb{R} \to x \leq z \to z \leq y \to \lvert x - \alpha \rvert < \varepsilon \to \lvert y - \alpha \rvert < \varepsilon \to \alpha \in \mathbb{R}$
\item $\varphi_{42} :\equiv \varepsilon \in \mathbb{R} \to x \leq z \to z \leq y \to \lvert x - \alpha \rvert < \varepsilon \to \lvert y - \alpha \rvert < \varepsilon \to \varepsilon \in \mathbb{R}$
\item $\varphi_{43} :\equiv \varepsilon \in \mathbb{R} \to x \leq z \to z \leq y \to \lvert x - \alpha \rvert < \varepsilon \to \lvert y - \alpha \rvert < \varepsilon \to y - \alpha < \varepsilon$
\item $\varphi_{44} :\equiv \varepsilon \in \mathbb{R} \to x \leq z \to z \leq y \to \lvert x - \alpha \rvert < \varepsilon \to \lvert y - \alpha \rvert < \varepsilon \to z - \alpha < \varepsilon$
\item $\varphi_{45} :\equiv z \leq y \to \lvert x - \alpha \rvert < \varepsilon \to \lvert y - \alpha \rvert < \varepsilon \to -\varepsilon < y - \alpha \land y - \alpha < \varepsilon$
\item $\varphi_{46} :\equiv z \leq y \to \lvert x - \alpha \rvert < \varepsilon \to \lvert y - \alpha \rvert < \varepsilon \to -\varepsilon < x - \alpha \land x - \alpha < \varepsilon$
\item $\varphi_{47} :\equiv \forall y\,(x - \alpha \in \mathbb{R} \to y \in \mathbb{R} \to \lvert x - \alpha \rvert < y \to -y < x - \alpha \land x - \alpha < y)$
\item $\varphi_{48} :\equiv \forall y\,(z - \alpha \in \mathbb{R} \to y \in \mathbb{R} \to -y < z - \alpha \to z - \alpha < y \to \lvert z - \alpha \rvert < y)$
\item $\varphi_{49} :\equiv \varepsilon \in \mathbb{R} \to x \leq z \to z \leq y \to \lvert x - \alpha \rvert < \varepsilon \to \lvert y - \alpha \rvert < \varepsilon \to y - \alpha \in \mathbb{R}$
\item $\varphi_{50} :\equiv \varepsilon \in \mathbb{R} \to x \leq z \to z \leq y \to \lvert x - \alpha \rvert < \varepsilon \to \lvert y - \alpha \rvert < \varepsilon \to -\alpha \in \mathbb{R}$
\item $\varphi_{51} :\equiv \varepsilon \in \mathbb{R} \to x \leq z \to z \leq y \to \lvert x - \alpha \rvert < \varepsilon \to \lvert y - \alpha \rvert < \varepsilon \to x - \alpha \in \mathbb{R}$
\item $\varphi_{52} :\equiv \varepsilon \in \mathbb{R} \to x \leq z \to z \leq y \to \lvert x - \alpha \rvert < \varepsilon \to \lvert y - \alpha \rvert < \varepsilon \to z - \alpha \in \mathbb{R}$
\item $\varphi_{53} :\equiv \forall v_{397}\,\forall v_{398}\,(z \in \mathbb{R} \to v_{397} \in \mathbb{R} \to v_{398} \in \mathbb{R} \to z \leq v_{397} \to z + v_{398} \leq v_{397} + v_{398})$
\item $\varphi_{54} :\equiv \forall v_{394}\,\forall v_{395}\,(x \in \mathbb{R} \to v_{394} \in \mathbb{R} \to v_{395} \in \mathbb{R} \to x \leq v_{394} \to x + v_{395} \leq v_{394} + v_{395})$
\item $\varphi_{55} :\equiv z - \alpha = z + -\alpha \to (z - \alpha = z + -\alpha \to z + -\alpha = z - \alpha) \to z + -\alpha = z - \alpha$
\item $\varphi_{56} :\equiv y - \alpha = y + -\alpha \to (y - \alpha = y + -\alpha \to y + -\alpha = y - \alpha) \to y + -\alpha = y - \alpha$
\item $\varphi_{57} :\equiv x - \alpha = x + -\alpha \to (x - \alpha = x + -\alpha \to x + -\alpha = x - \alpha) \to x + -\alpha = x - \alpha$
\item $\varphi_{58} :\equiv \forall y\,(x \in \mathbb{R} \to y \in \mathbb{R} \to x - y \in \mathbb{R}) \to x \in \mathbb{R} \to \alpha \in \mathbb{R} \to x - \alpha \in \mathbb{R}$
\item $\varphi_{59} :\equiv \forall y\,(z \in \mathbb{R} \to y \in \mathbb{R} \to z - y \in \mathbb{R}) \to z \in \mathbb{R} \to \alpha \in \mathbb{R} \to z - \alpha \in \mathbb{R}$
\item $\varphi_{60} :\equiv \varepsilon \in \mathbb{R} \to x \leq z \to z \leq y \to \lvert x - \alpha \rvert < \varepsilon \to \lvert y - \alpha \rvert < \varepsilon \to -\varepsilon < x - \alpha$
\item $\varphi_{61} :\equiv \varepsilon \in \mathbb{R} \to x \leq z \to z \leq y \to \lvert x - \alpha \rvert < \varepsilon \to \lvert y - \alpha \rvert < \varepsilon \to -\varepsilon < z - \alpha$
\item $\varphi_{62} :\equiv \forall z_{1}\,\forall w\,(z_{1} = w \to z - \alpha = z_{1} \to z - \alpha = w) \to \forall w\,(y + -\alpha = w \to z - \alpha = y + -\alpha \to z - \alpha = w)$
\item $\varphi_{63} :\equiv \forall z_{1}\,\forall w\,(z_{1} = w \to x - \alpha = z_{1} \to x - \alpha = w) \to \forall w\,(z + -\alpha = w \to x - \alpha = z + -\alpha \to x - \alpha = w)$
\item $\varphi_{64} :\equiv y - \alpha \in \mathbb{R} \to \varepsilon \in \mathbb{R} \to \lvert y - \alpha \rvert < \varepsilon \to -\varepsilon < y - \alpha \land y - \alpha < \varepsilon$
\item $\varphi_{65} :\equiv x - \alpha \in \mathbb{R} \to \varepsilon \in \mathbb{R} \to \lvert x - \alpha \rvert < \varepsilon \to -\varepsilon < x - \alpha \land x - \alpha < \varepsilon$
\item $\varphi_{66} :\equiv z - \alpha \in \mathbb{R} \to \varepsilon \in \mathbb{R} \to -\varepsilon < z - \alpha \to z - \alpha < \varepsilon \to \lvert z - \alpha \rvert < \varepsilon$
\item $\varphi_{67} :\equiv \forall x\,\forall y\,(x \in \mathbb{R} \to y \in \mathbb{R} \to x - y \in \mathbb{R}) \to \forall y\,(x \in \mathbb{R} \to y \in \mathbb{R} \to x - y \in \mathbb{R})$
\item $\varphi_{68} :\equiv \forall y\,\forall z\,(-\varepsilon < y \to y \leq z \to -\varepsilon < z) \to \forall z\,(-\varepsilon < x - \alpha \to x - \alpha \leq z \to -\varepsilon < z)$
\item $\varphi_{69} :\equiv \forall x\,\forall y\,(x \in \mathbb{R} \to y \in \mathbb{R} \to x - y \in \mathbb{R}) \to \forall y\,(z \in \mathbb{R} \to y \in \mathbb{R} \to z - y \in \mathbb{R})$
\item $\varphi_{70} :\equiv \forall v_{392}\,(y \in \mathbb{R} \to v_{392} \in \mathbb{R} \to y - v_{392} \in \mathbb{R}) \to y \in \mathbb{R} \to \alpha \in \mathbb{R} \to y - \alpha \in \mathbb{R}$
\item $\varphi_{71} :\equiv \forall x\,\forall y\,(x \in \mathbb{R} \to y \in \mathbb{R} \to x - y \in \mathbb{R}) \to \forall y\,(v_{391} \in \mathbb{R} \to y \in \mathbb{R} \to v_{391} - y \in \mathbb{R})$
\item $\varphi_{72} :\equiv \forall v_{390}\,(y - \alpha \in \mathbb{R} \to v_{390} \in \mathbb{R} \to \lvert y - \alpha \rvert < v_{390} \to -v_{390} < y - \alpha \land y - \alpha < v_{390})$
\item $\varphi_{73} :\equiv \forall z_{1}\,\forall w\,(z_{1} = w \to z - \alpha < z_{1} \to z - \alpha < w) \to \forall w\,(y + -\alpha = w \to z - \alpha < y + -\alpha \to z - \alpha < w)$
\item $\varphi_{74} :\equiv \forall z_{1}\,\forall w\,(z_{1} = w \to x - \alpha < z_{1} \to x - \alpha < w) \to \forall w\,(z + -\alpha = w \to x - \alpha < z + -\alpha \to x - \alpha < w)$
\item $\varphi_{75} :\equiv \forall y\,(v_{391} \in \mathbb{R} \to y \in \mathbb{R} \to v_{391} - y \in \mathbb{R}) \to v_{391} \in \mathbb{R} \to v_{392} \in \mathbb{R} \to v_{391} - v_{392} \in \mathbb{R}$
\item $\varphi_{76} :\equiv \forall v_{405}\,(z - \alpha \leq y - \alpha \to y - \alpha < v_{405} \to z - \alpha < v_{405}) \to z - \alpha \leq y - \alpha \to y - \alpha < \varepsilon \to z - \alpha < \varepsilon$
\item $\varphi_{77} :\equiv \forall z_{1}\,\forall w\,(z_{1} = w \to z_{1} = y + -\alpha \to w = y + -\alpha) \to \forall w\,(z + -\alpha = w \to z + -\alpha = y + -\alpha \to w = y + -\alpha)$
\item $\varphi_{78} :\equiv \forall z_{1}\,\forall w\,(z_{1} = w \to z_{1} = z + -\alpha \to w = z + -\alpha) \to \forall w\,(x + -\alpha = w \to x + -\alpha = z + -\alpha \to w = z + -\alpha)$
\item $\varphi_{79} :\equiv \forall z\,(-\varepsilon < x - \alpha \to x - \alpha \leq z \to -\varepsilon < z) \to -\varepsilon < x - \alpha \to x - \alpha \leq z - \alpha \to -\varepsilon < z - \alpha$
\item $\varphi_{80} :\equiv \forall w\,(y + -\alpha = w \to z - \alpha = y + -\alpha \to z - \alpha = w) \to y + -\alpha = y - \alpha \to z - \alpha = y + -\alpha \to z - \alpha = y - \alpha$
\item $\varphi_{81} :\equiv \forall w\,(z + -\alpha = w \to x - \alpha = z + -\alpha \to x - \alpha = w) \to z + -\alpha = z - \alpha \to x - \alpha = z + -\alpha \to x - \alpha = z - \alpha$
\item $\varphi_{82} :\equiv \varphi_{4} \to \forall v_{405}\,(z - \alpha \leq y - \alpha \to y - \alpha < v_{405} \to z - \alpha < v_{405})$
\item $\varphi_{83} :\equiv \forall z_{1}\,\forall w\,(z_{1} = w \to z_{1} < y + -\alpha \to w < y + -\alpha) \to \forall w\,(z + -\alpha = w \to z + -\alpha < y + -\alpha \to w < y + -\alpha)$
\item $\varphi_{84} :\equiv \forall z_{1}\,\forall w\,(z_{1} = w \to z_{1} < z + -\alpha \to w < z + -\alpha) \to \forall w\,(x + -\alpha = w \to x + -\alpha < z + -\alpha \to w < z + -\alpha)$
\item $\varphi_{85} :\equiv \forall w\,(y + -\alpha = w \to z - \alpha < y + -\alpha \to z - \alpha < w) \to y + -\alpha = y - \alpha \to z - \alpha < y + -\alpha \to z - \alpha < y - \alpha$
\item $\varphi_{86} :\equiv \forall w\,(z + -\alpha = w \to x - \alpha < z + -\alpha \to x - \alpha < w) \to z + -\alpha = z - \alpha \to x - \alpha < z + -\alpha \to x - \alpha < z - \alpha$
\item $\varphi_{87} :\equiv \forall w\,(z + -\alpha = w \to z + -\alpha = y + -\alpha \to w = y + -\alpha) \to z + -\alpha = z - \alpha \to z + -\alpha = y + -\alpha \to z - \alpha = y + -\alpha$
\item $\varphi_{88} :\equiv \forall w\,(x + -\alpha = w \to x + -\alpha = z + -\alpha \to w = z + -\alpha) \to x + -\alpha = x - \alpha \to x + -\alpha = z + -\alpha \to x - \alpha = z + -\alpha$
\item $\varphi_{89} :\equiv \forall x\,\forall y\,(x \in \mathbb{R} \to y \in \mathbb{R} \to \lvert x \rvert < y \to -y < x \land x < y) \to \varphi_{25}$
\item $\varphi_{90} :\equiv \forall w\,(z + -\alpha = w \to z + -\alpha < y + -\alpha \to w < y + -\alpha) \to z + -\alpha = z - \alpha \to z + -\alpha < y + -\alpha \to z - \alpha < y + -\alpha$
\item $\varphi_{91} :\equiv \forall w\,(x + -\alpha = w \to x + -\alpha < z + -\alpha \to w < z + -\alpha) \to x + -\alpha = x - \alpha \to x + -\alpha < z + -\alpha \to x - \alpha < z + -\alpha$
\item $\varphi_{92} :\equiv \varphi_{19} \to z \in \mathbb{R} \to y \in \mathbb{R} \to -\alpha \in \mathbb{R} \to z \leq y \to z + -\alpha \leq y + -\alpha$
\item $\varphi_{93} :\equiv \varphi_{21} \to x \in \mathbb{R} \to z \in \mathbb{R} \to -\alpha \in \mathbb{R} \to x \leq z \to x + -\alpha \leq z + -\alpha$
\item $\varphi_{94} :\equiv \forall \alpha\,\forall \varepsilon\,(x \in \mathbb{R} \to y \in \mathbb{R} \to z \in \mathbb{R} \to \alpha \in \mathbb{R} \to \varepsilon \in \mathbb{R} \to \varphi_{29})$
\item $\varphi_{95} :\equiv \forall x\,\forall y\,(x \in \mathbb{R} \to y \in \mathbb{R} \to \lvert x \rvert < y \to -y < x \land x < y) \to \varphi_{47}$
\item $\varphi_{96} :\equiv \forall x\,\forall y\,(x \in \mathbb{R} \to y \in \mathbb{R} \to -y < x \to x < y \to \lvert x \rvert < y) \to \varphi_{48}$
\item $\varphi_{97} :\equiv (y + -\alpha = y - \alpha \to z - \alpha = y + -\alpha \to z - \alpha = y - \alpha) \to x \in \mathbb{R} \to y \in \mathbb{R} \to z \in \mathbb{R} \to \alpha \in \mathbb{R} \to \varepsilon \in \mathbb{R} \to \varphi_{15}$
\item $\varphi_{98} :\equiv (z + -\alpha = z - \alpha \to x - \alpha = z + -\alpha \to x - \alpha = z - \alpha) \to x \in \mathbb{R} \to y \in \mathbb{R} \to z \in \mathbb{R} \to \alpha \in \mathbb{R} \to \varepsilon \in \mathbb{R} \to \varphi_{16}$
\item $\varphi_{99} :\equiv (z + -\alpha = z - \alpha \to z + -\alpha = y + -\alpha \to z - \alpha = y + -\alpha) \to x \in \mathbb{R} \to y \in \mathbb{R} \to z \in \mathbb{R} \to \alpha \in \mathbb{R} \to \varepsilon \in \mathbb{R} \to \varphi_{30}$
\item $\varphi_{100} :\equiv (x + -\alpha = x - \alpha \to x + -\alpha = z + -\alpha \to x - \alpha = z + -\alpha) \to x \in \mathbb{R} \to y \in \mathbb{R} \to z \in \mathbb{R} \to \alpha \in \mathbb{R} \to \varepsilon \in \mathbb{R} \to \varphi_{31}$
\item $\varphi_{101} :\equiv (x \in \mathbb{R} \to y \in \mathbb{R} \to z \in \mathbb{R} \to \alpha \in \mathbb{R} \to \varphi_{41}) \to x \in \mathbb{R} \to y \in \mathbb{R} \to z \in \mathbb{R} \to \alpha \in \mathbb{R} \to \varphi_{49}$
\item $\varphi_{102} :\equiv (x \in \mathbb{R} \to y \in \mathbb{R} \to z \in \mathbb{R} \to \alpha \in \mathbb{R} \to \varphi_{41}) \to x \in \mathbb{R} \to y \in \mathbb{R} \to z \in \mathbb{R} \to \alpha \in \mathbb{R} \to \varphi_{50}$
\item $\varphi_{103} :\equiv (x \in \mathbb{R} \to y \in \mathbb{R} \to z \in \mathbb{R} \to \alpha \in \mathbb{R} \to \varphi_{41}) \to x \in \mathbb{R} \to y \in \mathbb{R} \to z \in \mathbb{R} \to \alpha \in \mathbb{R} \to \varphi_{51}$
\item $\varphi_{104} :\equiv (x \in \mathbb{R} \to y \in \mathbb{R} \to z \in \mathbb{R} \to \alpha \in \mathbb{R} \to \varphi_{41}) \to x \in \mathbb{R} \to y \in \mathbb{R} \to z \in \mathbb{R} \to \alpha \in \mathbb{R} \to \varphi_{52}$
\item $\varphi_{105} :\equiv (x \in \mathbb{R} \to y \in \mathbb{R} \to z \in \mathbb{R} \to \alpha \in \mathbb{R} \to \varphi_{23}) \to x \in \mathbb{R} \to y \in \mathbb{R} \to z \in \mathbb{R} \to \alpha \in \mathbb{R} \to \varepsilon \in \mathbb{R} \to \varphi_{38}$
\item $\varphi_{106} :\equiv (x \in \mathbb{R} \to y \in \mathbb{R} \to z \in \mathbb{R} \to \alpha \in \mathbb{R} \to \varphi_{43}) \to x \in \mathbb{R} \to y \in \mathbb{R} \to z \in \mathbb{R} \to \alpha \in \mathbb{R} \to \varphi_{44}$
\item $\varphi_{107} :\equiv (x \in \mathbb{R} \to y \in \mathbb{R} \to z \in \mathbb{R} \to \alpha \in \mathbb{R} \to \varphi_{24}) \to x \in \mathbb{R} \to y \in \mathbb{R} \to z \in \mathbb{R} \to \alpha \in \mathbb{R} \to \varepsilon \in \mathbb{R} \to \varphi_{40}$
\item $\varphi_{108} :\equiv (x \in \mathbb{R} \to y \in \mathbb{R} \to z \in \mathbb{R} \to \alpha \in \mathbb{R} \to \varphi_{44}) \to x \in \mathbb{R} \to y \in \mathbb{R} \to z \in \mathbb{R} \to \alpha \in \mathbb{R} \to \varepsilon \in \mathbb{R} \to \varphi_{29}$
\item $\varphi_{109} :\equiv (x \in \mathbb{R} \to y \in \mathbb{R} \to z \in \mathbb{R} \to \alpha \in \mathbb{R} \to \varepsilon \in \mathbb{R} \to \varphi_{16}) \to x \in \mathbb{R} \to y \in \mathbb{R} \to z \in \mathbb{R} \to \alpha \in \mathbb{R} \to \varphi_{61}$
\item $\varphi_{110} :\equiv (z + -\alpha = z - \alpha \to z + -\alpha < y + -\alpha \to z - \alpha < y + -\alpha) \to \varphi_{99}$
\item $\varphi_{111} :\equiv (x + -\alpha = x - \alpha \to x + -\alpha < z + -\alpha \to x - \alpha < z + -\alpha) \to \varphi_{100}$
\item $\varphi_{112} :\equiv (x \in \mathbb{R} \to y \in \mathbb{R} \to z \in \mathbb{R} \to \alpha \in \mathbb{R} \to \varepsilon \in \mathbb{R} \to x \leq z \to \varphi_{45}) \to x \in \mathbb{R} \to y \in \mathbb{R} \to z \in \mathbb{R} \to \alpha \in \mathbb{R} \to \varphi_{43}$
\item $\varphi_{113} :\equiv (x \in \mathbb{R} \to y \in \mathbb{R} \to z \in \mathbb{R} \to \alpha \in \mathbb{R} \to \varepsilon \in \mathbb{R} \to \varphi_{27}) \to x \in \mathbb{R} \to y \in \mathbb{R} \to z \in \mathbb{R} \to \alpha \in \mathbb{R} \to \varepsilon \in \mathbb{R} \to x \leq z \to \varphi_{45}$
\item $\varphi_{114} :\equiv (x \in \mathbb{R} \to y \in \mathbb{R} \to z \in \mathbb{R} \to \alpha \in \mathbb{R} \to \varepsilon \in \mathbb{R} \to \varphi_{28}) \to x \in \mathbb{R} \to y \in \mathbb{R} \to z \in \mathbb{R} \to \alpha \in \mathbb{R} \to \varepsilon \in \mathbb{R} \to x \leq z \to \varphi_{46}$
\item $\varphi_{115} :\equiv (x \in \mathbb{R} \to y \in \mathbb{R} \to z \in \mathbb{R} \to \alpha \in \mathbb{R} \to \varepsilon \in \mathbb{R} \to x \leq z \to \varphi_{46}) \to x \in \mathbb{R} \to y \in \mathbb{R} \to z \in \mathbb{R} \to \alpha \in \mathbb{R} \to \varphi_{60}$
\item $\varphi_{116} :\equiv (x \in \mathbb{R} \to y \in \mathbb{R} \to z \in \mathbb{R} \to \alpha \in \mathbb{R} \to \varepsilon \in \mathbb{R} \to \varphi_{30}) \to (y + -\alpha = y - \alpha \to z - \alpha < y + -\alpha \to z - \alpha < y - \alpha) \to \varphi_{97}$
\item $\varphi_{117} :\equiv (x \in \mathbb{R} \to y \in \mathbb{R} \to z \in \mathbb{R} \to \alpha \in \mathbb{R} \to \varepsilon \in \mathbb{R} \to \varphi_{31}) \to (z + -\alpha = z - \alpha \to x - \alpha < z + -\alpha \to x - \alpha < z - \alpha) \to \varphi_{98}$
\item $\varphi_{118} :\equiv (y \in \mathbb{R} \to \alpha \in \mathbb{R} \to y - \alpha \in \mathbb{R}) \to (x \in \mathbb{R} \to y \in \mathbb{R} \to z \in \mathbb{R} \to \alpha \in \mathbb{R} \to \varphi_{36}) \to \varphi_{101}$
\item $\varphi_{119} :\equiv (x \in \mathbb{R} \to \alpha \in \mathbb{R} \to x - \alpha \in \mathbb{R}) \to (x \in \mathbb{R} \to y \in \mathbb{R} \to z \in \mathbb{R} \to \alpha \in \mathbb{R} \to \varphi_{39}) \to \varphi_{103}$
\item $\varphi_{120} :\equiv (z \in \mathbb{R} \to \alpha \in \mathbb{R} \to z - \alpha \in \mathbb{R}) \to (x \in \mathbb{R} \to y \in \mathbb{R} \to z \in \mathbb{R} \to \alpha \in \mathbb{R} \to \varphi_{37}) \to \varphi_{104}$
\item $\varphi_{121} :\equiv z + -\alpha = z - \alpha \to (x \in \mathbb{R} \to y \in \mathbb{R} \to z \in \mathbb{R} \to \alpha \in \mathbb{R} \to \varepsilon \in \mathbb{R} \to \varphi_{38}) \to \varphi_{110}$
\item $\varphi_{122} :\equiv (z - \alpha \leq y - \alpha \to y - \alpha < \varepsilon \to z - \alpha < \varepsilon) \to (x \in \mathbb{R} \to y \in \mathbb{R} \to z \in \mathbb{R} \to \alpha \in \mathbb{R} \to \varepsilon \in \mathbb{R} \to \varphi_{15}) \to \varphi_{106}$
\item $\varphi_{123} :\equiv x + -\alpha = x - \alpha \to (x \in \mathbb{R} \to y \in \mathbb{R} \to z \in \mathbb{R} \to \alpha \in \mathbb{R} \to \varepsilon \in \mathbb{R} \to \varphi_{40}) \to \varphi_{111}$
\item $\varphi_{124} :\equiv (-\varepsilon < x - \alpha \to x - \alpha \leq z - \alpha \to -\varepsilon < z - \alpha) \to (x \in \mathbb{R} \to y \in \mathbb{R} \to z \in \mathbb{R} \to \alpha \in \mathbb{R} \to \varphi_{60}) \to \varphi_{109}$
\item $\varphi_{125} :\equiv (x \in \mathbb{R} \to y \in \mathbb{R} \to z \in \mathbb{R} \to \alpha \in \mathbb{R} \to \varphi_{36}) \to (x \in \mathbb{R} \to y \in \mathbb{R} \to z \in \mathbb{R} \to \alpha \in \mathbb{R} \to \varphi_{50}) \to \varphi_{105}$
\item $\varphi_{126} :\equiv (x \in \mathbb{R} \to y \in \mathbb{R} \to z \in \mathbb{R} \to \alpha \in \mathbb{R} \to \varphi_{37}) \to (x \in \mathbb{R} \to y \in \mathbb{R} \to z \in \mathbb{R} \to \alpha \in \mathbb{R} \to \varphi_{50}) \to \varphi_{107}$
\item $\varphi_{127} :\equiv (x \in \mathbb{R} \to y \in \mathbb{R} \to z \in \mathbb{R} \to \alpha \in \mathbb{R} \to \varphi_{42}) \to (x \in \mathbb{R} \to y \in \mathbb{R} \to z \in \mathbb{R} \to \alpha \in \mathbb{R} \to \varphi_{61}) \to \varphi_{108}$
\item $\varphi_{128} :\equiv (x \in \mathbb{R} \to y \in \mathbb{R} \to z \in \mathbb{R} \to \alpha \in \mathbb{R} \to \varphi_{49}) \to (x \in \mathbb{R} \to y \in \mathbb{R} \to z \in \mathbb{R} \to \alpha \in \mathbb{R} \to \varphi_{42}) \to \varphi_{113}$
\item $\varphi_{129} :\equiv (x \in \mathbb{R} \to y \in \mathbb{R} \to z \in \mathbb{R} \to \alpha \in \mathbb{R} \to \varphi_{51}) \to (x \in \mathbb{R} \to y \in \mathbb{R} \to z \in \mathbb{R} \to \alpha \in \mathbb{R} \to \varphi_{42}) \to \varphi_{114}$
\item $\varphi_{130} :\equiv (z \in \mathbb{R} \to y \in \mathbb{R} \to -\alpha \in \mathbb{R} \to z \leq y \to z + -\alpha \leq y + -\alpha) \to (x \in \mathbb{R} \to y \in \mathbb{R} \to z \in \mathbb{R} \to \alpha \in \mathbb{R} \to \varphi_{37}) \to \varphi_{125}$
\item $\varphi_{131} :\equiv (x \in \mathbb{R} \to z \in \mathbb{R} \to -\alpha \in \mathbb{R} \to x \leq z \to x + -\alpha \leq z + -\alpha) \to (x \in \mathbb{R} \to y \in \mathbb{R} \to z \in \mathbb{R} \to \alpha \in \mathbb{R} \to \varphi_{39}) \to \varphi_{126}$
\end{abbrevlist}
\end{fullproof}

% 自動生成 (prover/build.js): Squeeze  ノード数 170
\begin{fullproof}{定理2.3}{$\forall a\,\forall b\,\forall c\,\forall \alpha\,(\mathrm{Seq}(c) \to \forall n \in \mathbb{N}\,(a_{n} \leq c_{n} \land c_{n} \leq b_{n}) \to \mathrm{Conv}(a,\alpha) \to \mathrm{Conv}(b,\alpha) \to \mathrm{Conv}(c,\alpha))$}{170}\phantomsection\label{tree:Squeeze}
\noindent$\Pi_{1}$:
\begin{ptree}
\AXC{$\varphi_{32}$\lbl{A2}}
\AXC{$\mathrm{Seq}(c) \to \forall n \in \mathbb{N}\,(a_{n} \leq c_{n} \land c_{n} \leq b_{n}) \to \varphi_{61}$\lbl{A1}}
\AXC{$\varphi_{90}$\lbl{A1}}
\BIC{$\mathrm{Seq}(c) \to \varphi_{78}$}
\AXC{$(\mathrm{Seq}(c) \to \varphi_{78}) \to \varphi_{32} \to \mathrm{Seq}(c) \to \varphi_{77}$\lbl{A1}}
\BIC{$\varphi_{32} \to \mathrm{Seq}(c) \to \varphi_{77}$}
\BIC{$\mathrm{Seq}(c) \to \varphi_{77}$}
\end{ptree}
\noindent$\Pi_{2}$:
\begin{ptree}
\AXC{$\varphi_{26}$\lbl{A2}}
\AXC{$\mathrm{Seq}(c) \to \varphi_{28}$\lbl{A1}}
\AXC{$(\mathrm{Seq}(c) \to \varphi_{28}) \to \varphi_{26} \to \mathrm{Seq}(c) \to \varphi_{26}$\lbl{A1}}
\BIC{$\varphi_{26} \to \mathrm{Seq}(c) \to \varphi_{26}$}
\BIC{$\mathrm{Seq}(c) \to \varphi_{26}$}
\AXC{$(\mathrm{Seq}(c) \to \varphi_{26}) \to \varphi_{86}$\lbl{A1}}
\BIC{$\varphi_{86}$}
\end{ptree}
\noindent$\Pi_{3}$:
\begin{ptree}
\AXC{$\mathrm{Seq}(c) \to \forall n \in \mathbb{N}\,(a_{n} \leq c_{n} \land c_{n} \leq b_{n}) \to \varphi_{36}$\lbl{A1}}
\AXC{$\stackrel{\Pi_{2}}{\vdots}$}\noLine
\UIC{$\varphi_{86}$}
\BIC{$\mathrm{Seq}(c) \to \varphi_{70}$}
\end{ptree}
\noindent$\Pi_{4}$:
\begin{ptree}
\AXC{$\mathrm{Seq}(c) \to \varphi_{9}$\lbl{A1}}
\AXC{$\forall a\,\forall \alpha\,(\mathrm{Conv}(a,\alpha) \to \alpha \in \mathbb{R})$\lbl{[C2]}}
\AXC{$\varphi_{7}$\lbl{A2}}
\BIC{$\forall \alpha\,(\mathrm{Conv}(a,\alpha) \to \alpha \in \mathbb{R})$}
\AXC{$\varphi_{5}$\lbl{A2}}
\BIC{$\mathrm{Conv}(a,\alpha) \to \alpha \in \mathbb{R}$}
\AXC{$\varphi_{59}$\lbl{A1}}
\BIC{$(\mathrm{Seq}(c) \to \varphi_{9}) \to \mathrm{Seq}(c) \to \varphi_{24}$}
\BIC{$\mathrm{Seq}(c) \to \varphi_{24}$}
\end{ptree}
\noindent$\Pi_{5}$:
\begin{ptree}
\AXC{$\mathrm{Seq}(c) \to \mathrm{Seq}(c)$\lbl{A1}}
\AXC{$\forall a\,\forall n\,(\mathrm{Seq}(a) \to n \in \mathbb{N} \to a_{n} \in \mathbb{R})$\lbl{[S1]}}
\AXC{$\varphi_{16}$\lbl{A2}}
\BIC{$\forall n\,(\mathrm{Seq}(c) \to n \in \mathbb{N} \to c_{n} \in \mathbb{R})$}
\AXC{$\varphi_{10}$\lbl{A2}}
\BIC{$\mathrm{Seq}(c) \to n \in \mathbb{N} \to c_{n} \in \mathbb{R}$}
\AXC{$\varphi_{89}$\lbl{A1}}
\BIC{$(\mathrm{Seq}(c) \to \mathrm{Seq}(c)) \to \varphi_{83}$}
\BIC{$\varphi_{83}$}
\end{ptree}
\noindent$\Pi_{6}$:
\begin{ptree}
\AXC{$\mathrm{Seq}(c) \to \varphi_{21}$\lbl{A1}}
\AXC{$\forall a\,\forall \alpha\,(\mathrm{Conv}(a,\alpha) \to \mathrm{Seq}(a))$\lbl{[C1]}}
\AXC{$\varphi_{3}$\lbl{A2}}
\BIC{$\forall \alpha\,(\mathrm{Conv}(b,\alpha) \to \mathrm{Seq}(b))$}
\AXC{$\varphi_{1}$\lbl{A2}}
\BIC{$\mathrm{Conv}(b,\alpha) \to \mathrm{Seq}(b)$}
\AXC{$\varphi_{60}$\lbl{A1}}
\BIC{$(\mathrm{Seq}(c) \to \varphi_{21}) \to \mathrm{Seq}(c) \to \varphi_{14}$}
\BIC{$\mathrm{Seq}(c) \to \varphi_{14}$}
\end{ptree}
\noindent$\Pi_{7}$:
\begin{ptree}
\AXC{$\stackrel{\Pi_{6}}{\vdots}$}\noLine
\UIC{$\mathrm{Seq}(c) \to \varphi_{14}$}
\AXC{$\forall a\,\forall n\,(\mathrm{Seq}(a) \to n \in \mathbb{N} \to a_{n} \in \mathbb{R})$\lbl{[S1]}}
\AXC{$\varphi_{17}$\lbl{A2}}
\BIC{$\forall n\,(\mathrm{Seq}(b) \to n \in \mathbb{N} \to b_{n} \in \mathbb{R})$}
\AXC{$\varphi_{11}$\lbl{A2}}
\BIC{$\mathrm{Seq}(b) \to n \in \mathbb{N} \to b_{n} \in \mathbb{R}$}
\AXC{$\varphi_{91}$\lbl{A1}}
\BIC{$(\mathrm{Seq}(c) \to \varphi_{14}) \to \varphi_{84}$}
\BIC{$\varphi_{84}$}
\end{ptree}
\noindent$\Pi_{8}$:
\begin{ptree}
\AXC{$\mathrm{Seq}(c) \to \varphi_{9}$\lbl{A1}}
\AXC{$\forall a\,\forall \alpha\,(\mathrm{Conv}(a,\alpha) \to \mathrm{Seq}(a))$\lbl{[C1]}}
\AXC{$\varphi_{4}$\lbl{A2}}
\BIC{$\forall \alpha\,(\mathrm{Conv}(a,\alpha) \to \mathrm{Seq}(a))$}
\AXC{$\varphi_{2}$\lbl{A2}}
\BIC{$\mathrm{Conv}(a,\alpha) \to \mathrm{Seq}(a)$}
\AXC{$\varphi_{57}$\lbl{A1}}
\BIC{$(\mathrm{Seq}(c) \to \varphi_{9}) \to \mathrm{Seq}(c) \to \varphi_{15}$}
\BIC{$\mathrm{Seq}(c) \to \varphi_{15}$}
\end{ptree}
\noindent$\Pi_{9}$:
\begin{ptree}
\AXC{$\stackrel{\Pi_{8}}{\vdots}$}\noLine
\UIC{$\mathrm{Seq}(c) \to \varphi_{15}$}
\AXC{$\forall a\,\forall n\,(\mathrm{Seq}(a) \to n \in \mathbb{N} \to a_{n} \in \mathbb{R})$\lbl{[S1]}}
\AXC{$\varphi_{18}$\lbl{A2}}
\BIC{$\forall n\,(\mathrm{Seq}(a) \to n \in \mathbb{N} \to a_{n} \in \mathbb{R})$}
\AXC{$\varphi_{12}$\lbl{A2}}
\BIC{$\mathrm{Seq}(a) \to n \in \mathbb{N} \to a_{n} \in \mathbb{R}$}
\AXC{$\varphi_{92}$\lbl{A1}}
\BIC{$(\mathrm{Seq}(c) \to \varphi_{15}) \to \varphi_{85}$}
\BIC{$\varphi_{85}$}
\end{ptree}
\noindent$\Pi_{10}$:
\begin{ptree}
\AXC{$\varphi_{35}$\lbl{[C12]}}
\AXC{$\varphi_{35} \to \varphi_{52}$\lbl{A2}}
\BIC{$\varphi_{52}$}
\AXC{$\varphi_{52} \to \varphi_{55}$\lbl{A2}}
\BIC{$\varphi_{55}$}
\AXC{$\varphi_{55} \to \forall \alpha\,\forall \varepsilon\,(\varphi_{58})$\lbl{A2}}
\BIC{$\forall \alpha\,\forall \varepsilon\,(\varphi_{58})$}
\AXC{$\forall \alpha\,\forall \varepsilon\,(\varphi_{58}) \to \forall \varepsilon\,(\varphi_{58})$\lbl{A2}}
\BIC{$\forall \varepsilon\,(\varphi_{58})$}
\AXC{$\forall \varepsilon\,(\varphi_{58}) \to \varphi_{58}$\lbl{A2}}
\BIC{$\varphi_{58}$}
\end{ptree}
\noindent$\Pi_{11}$:
\begin{ptree}
\AXC{$\mathrm{Seq}(c) \to \forall n \in \mathbb{N}\,(a_{n} \leq c_{n} \land c_{n} \leq b_{n}) \to \varphi_{36}$\lbl{A1}}
\AXC{$\stackrel{\Pi_{9}}{\vdots}$}\noLine
\UIC{$\varphi_{85}$}
\BIC{$\varphi_{66}$}
\AXC{$\stackrel{\Pi_{10}}{\vdots}$}\noLine
\UIC{$\varphi_{58}$}
\AXC{$\varphi_{58} \to \varphi_{66} \to \varphi_{65} \to \varphi_{64} \to \varphi_{99}$\lbl{A1}}
\BIC{$\varphi_{66} \to \varphi_{65} \to \varphi_{64} \to \varphi_{99}$}
\BIC{$\varphi_{65} \to \varphi_{64} \to \varphi_{99}$}
\end{ptree}
\noindent$\Pi_{12}$:
\begin{ptree}
\AXC{$\mathrm{Seq}(c) \to \forall n \in \mathbb{N}\,(a_{n} \leq c_{n} \land c_{n} \leq b_{n}) \to \varphi_{36}$\lbl{A1}}
\AXC{$\stackrel{\Pi_{7}}{\vdots}$}\noLine
\UIC{$\varphi_{84}$}
\BIC{$\varphi_{65}$}
\AXC{$\stackrel{\Pi_{11}}{\vdots}$}\noLine
\UIC{$\varphi_{65} \to \varphi_{64} \to \varphi_{99}$}
\BIC{$\varphi_{64} \to \varphi_{99}$}
\end{ptree}
\noindent$\Pi_{13}$:
\begin{ptree}
\AXC{$\mathrm{Seq}(c) \to \varphi_{29}$\lbl{A1}}
\AXC{$\stackrel{\Pi_{4}}{\vdots}$}\noLine
\UIC{$\mathrm{Seq}(c) \to \varphi_{24}$}
\AXC{$\mathrm{Seq}(c) \to \forall n \in \mathbb{N}\,(a_{n} \leq c_{n} \land c_{n} \leq b_{n}) \to \varphi_{36}$\lbl{A1}}
\AXC{$\stackrel{\Pi_{5}}{\vdots}$}\noLine
\UIC{$\varphi_{83}$}
\BIC{$\varphi_{64}$}
\AXC{$\stackrel{\Pi_{12}}{\vdots}$}\noLine
\UIC{$\varphi_{64} \to \varphi_{99}$}
\BIC{$\varphi_{99}$}
\BIC{$(\mathrm{Seq}(c) \to \varphi_{29}) \to \varphi_{63} \to \varphi_{62} \to \varphi_{68} \to \varphi_{67} \to \varphi_{69}$}
\BIC{$\varphi_{63} \to \varphi_{62} \to \varphi_{68} \to \varphi_{67} \to \varphi_{69}$}
\end{ptree}
\noindent$\Pi_{14}$:
\begin{ptree}
\AXC{$\stackrel{\Pi_{3}}{\vdots}$}\noLine
\UIC{$\mathrm{Seq}(c) \to \varphi_{70}$}
\AXC{$(\mathrm{Seq}(c) \to \varphi_{70}) \to \varphi_{63}$\lbl{A1}}
\BIC{$\varphi_{63}$}
\AXC{$\stackrel{\Pi_{13}}{\vdots}$}\noLine
\UIC{$\varphi_{63} \to \varphi_{62} \to \varphi_{68} \to \varphi_{67} \to \varphi_{69}$}
\BIC{$\varphi_{62} \to \varphi_{68} \to \varphi_{67} \to \varphi_{69}$}
\end{ptree}
\noindent$\Pi_{15}$:
\begin{ptree}
\AXC{$\stackrel{\Pi_{1}}{\vdots}$}\noLine
\UIC{$\mathrm{Seq}(c) \to \varphi_{77}$}
\AXC{$(\mathrm{Seq}(c) \to \varphi_{77}) \to \varphi_{68}$\lbl{A1}}
\BIC{$\varphi_{68}$}
\AXC{$\stackrel{\Pi_{3}}{\vdots}$}\noLine
\UIC{$\mathrm{Seq}(c) \to \varphi_{70}$}
\AXC{$(\mathrm{Seq}(c) \to \varphi_{70}) \to \varphi_{62}$\lbl{A1}}
\BIC{$\varphi_{62}$}
\AXC{$\stackrel{\Pi_{14}}{\vdots}$}\noLine
\UIC{$\varphi_{62} \to \varphi_{68} \to \varphi_{67} \to \varphi_{69}$}
\BIC{$\varphi_{68} \to \varphi_{67} \to \varphi_{69}$}
\BIC{$\varphi_{67} \to \varphi_{69}$}
\end{ptree}
\noindent$\Pi_{16}$:
\begin{ptree}
\AXC{$\varphi_{71} \to \varphi_{73}$\lbl{A3}}
\AXC{$\stackrel{\Pi_{1}}{\vdots}$}\noLine
\UIC{$\mathrm{Seq}(c) \to \varphi_{77}$}
\AXC{$(\mathrm{Seq}(c) \to \varphi_{77}) \to \varphi_{67}$\lbl{A1}}
\BIC{$\varphi_{67}$}
\AXC{$\stackrel{\Pi_{15}}{\vdots}$}\noLine
\UIC{$\varphi_{67} \to \varphi_{69}$}
\BIC{$\varphi_{69}$}
\AXC{$\varphi_{87}$\lbl{A1}}
\BIC{$\varphi_{13} \land \varepsilon \in \mathbb{R} \land 0 < \varepsilon \land \varphi_{19} \to n \in \mathbb{N} \to N \leq n \to \lvert c_{n} - \alpha \rvert < \varepsilon$}
\UIC{$\varphi_{71}$}
\AXC{$\varphi_{71} \to (\varphi_{71} \to \varphi_{73}) \to \varphi_{94}$\lbl{A1}}
\BIC{$(\varphi_{71} \to \varphi_{73}) \to \varphi_{94}$}
\BIC{$\varphi_{94}$}
\end{ptree}
\noindent$\Pi_{17}$:
\begin{ptree}
\AXC{$\mathrm{Seq}(c) \to \forall n \in \mathbb{N}\,(a_{n} \leq c_{n} \land c_{n} \leq b_{n}) \to \varphi_{61}$\lbl{A1}}
\AXC{$\varphi_{88}$\lbl{A1}}
\BIC{$\mathrm{Seq}(c) \to \forall n \in \mathbb{N}\,(a_{n} \leq c_{n} \land c_{n} \leq b_{n}) \to \varphi_{34}$}
\AXC{$\varphi_{95}$\lbl{A1}}
\BIC{$\varphi_{72} \to \mathrm{Seq}(c) \to \varphi_{74}$}
\end{ptree}
\noindent$\Pi_{18}$:
\begin{ptree}
\AXC{$\varphi_{13} \land \varepsilon \in \mathbb{R} \land 0 < \varepsilon \land \varphi_{19} \to \forall n\,(\varphi_{13} \land \varepsilon \in \mathbb{R} \land 0 < \varepsilon \land \varphi_{19})$\lbl{A4}}
\AXC{$\stackrel{\Pi_{16}}{\vdots}$}\noLine
\UIC{$\varphi_{94}$}
\BIC{$\varphi_{72}$}
\AXC{$\stackrel{\Pi_{17}}{\vdots}$}\noLine
\UIC{$\varphi_{72} \to \mathrm{Seq}(c) \to \varphi_{74}$}
\BIC{$\mathrm{Seq}(c) \to \varphi_{74}$}
\AXC{$(\mathrm{Seq}(c) \to \varphi_{74}) \to \varphi_{81}$\lbl{A1}}
\BIC{$\varphi_{81}$}
\end{ptree}
\noindent$\Pi_{19}$:
\begin{ptree}
\AXC{$\varphi_{8} \to \neg (N \in \mathbb{N} \land \forall n \in \mathbb{N}\,(N \leq n \to \lvert c_{n} - \alpha \rvert < \varepsilon))$\lbl{A2}}
\AXC{$\stackrel{\Pi_{18}}{\vdots}$}\noLine
\UIC{$\varphi_{81}$}
\BIC{$\mathrm{Seq}(c) \to \varphi_{75}$}
\AXC{$(\mathrm{Seq}(c) \to \varphi_{75}) \to \varphi_{53} \to \neg \varphi_{19}$\lbl{A1}}
\BIC{$\varphi_{53} \to \neg \varphi_{19}$}
\UIC{$\forall N\,(\varphi_{53} \to \neg \varphi_{19})$}
\end{ptree}
\noindent$\Pi_{20}$:
\begin{ptree}
\AXC{$\forall a\,\forall \alpha\,\forall b\,(\varphi_{33})$\lbl{[C9]}}
\AXC{$\forall a\,\forall \alpha\,\forall b\,(\varphi_{33}) \to \forall \alpha\,\forall b\,(\varphi_{33})$\lbl{A2}}
\BIC{$\forall \alpha\,\forall b\,(\varphi_{33})$}
\AXC{$\forall \alpha\,\forall b\,(\varphi_{33}) \to \forall b\,(\varphi_{33})$\lbl{A2}}
\BIC{$\forall b\,(\varphi_{33})$}
\AXC{$\forall b\,(\varphi_{33}) \to \varphi_{33}$\lbl{A2}}
\BIC{$\varphi_{33}$}
\end{ptree}
\noindent$\Pi_{21}$:
\begin{ptree}
\AXC{$\mathrm{Seq}(c) \to \varphi_{21}$\lbl{A1}}
\AXC{$\mathrm{Seq}(c) \to \varphi_{9}$\lbl{A1}}
\AXC{$\stackrel{\Pi_{20}}{\vdots}$}\noLine
\UIC{$\varphi_{33}$}
\AXC{$\varphi_{33} \to \varphi_{31}$\lbl{A2}}
\BIC{$\varphi_{31}$}
\AXC{$\varphi_{79}$\lbl{A2}}
\BIC{$\mathrm{Conv}(a,\alpha) \to \mathrm{Conv}(b,\alpha) \to \varepsilon \in \mathbb{R} \to 0 < \varepsilon \to \varphi_{22}$}
\AXC{$\varphi_{97}$\lbl{A1}}
\BIC{$\varphi_{93}$}
\BIC{$(\mathrm{Seq}(c) \to \varphi_{21}) \to (\mathrm{Seq}(c) \to \varphi_{29}) \to \varphi_{80}$}
\BIC{$(\mathrm{Seq}(c) \to \varphi_{29}) \to \varphi_{80}$}
\end{ptree}
\noindent$\Pi_{22}$:
\begin{ptree}
\AXC{$\mathrm{Seq}(c) \to \forall n \in \mathbb{N}\,(a_{n} \leq c_{n} \land c_{n} \leq b_{n}) \to \varphi_{6}$\lbl{A1}}
\AXC{$\mathrm{Seq}(c) \to \varphi_{29}$\lbl{A1}}
\AXC{$\stackrel{\Pi_{21}}{\vdots}$}\noLine
\UIC{$(\mathrm{Seq}(c) \to \varphi_{29}) \to \varphi_{80}$}
\BIC{$\varphi_{80}$}
\BIC{$\mathrm{Seq}(c) \to \varphi_{56}$}
\end{ptree}
\noindent$\Pi_{23}$:
\begin{ptree}
\AXC{$\stackrel{\Pi_{19}}{\vdots}$}\noLine
\UIC{$\forall N\,(\varphi_{53} \to \neg \varphi_{19})$}
\AXC{$\stackrel{\Pi_{22}}{\vdots}$}\noLine
\UIC{$\mathrm{Seq}(c) \to \varphi_{56}$}
\AXC{$(\mathrm{Seq}(c) \to \varphi_{56}) \to \forall N\,(\varphi_{53} \to \neg \varphi_{19}) \to \varphi_{98}$\lbl{A1}}
\BIC{$\forall N\,(\varphi_{53} \to \neg \varphi_{19}) \to \varphi_{98}$}
\BIC{$\varphi_{98}$}
\end{ptree}
\noindent$\Pi_{24}$:
\begin{ptree}
\AXC{$\varphi_{53} \to \forall N\,(\varphi_{53})$\lbl{A4}}
\AXC{$\forall N\,(\varphi_{53} \to \neg \varphi_{19}) \to \forall N\,(\varphi_{53}) \to \forall N\,(\neg \varphi_{19})$\lbl{A3}}
\AXC{$\stackrel{\Pi_{23}}{\vdots}$}\noLine
\UIC{$\varphi_{98}$}
\BIC{$(\varphi_{53} \to \forall N\,(\varphi_{53})) \to \mathrm{Seq}(c) \to \varphi_{42}$}
\BIC{$\mathrm{Seq}(c) \to \varphi_{42}$}
\AXC{$(\mathrm{Seq}(c) \to \varphi_{42}) \to \varphi_{13} \to \varphi_{20}$\lbl{A1}}
\BIC{$\varphi_{13} \to \varphi_{20}$}
\UIC{$\forall \varepsilon\,(\varphi_{13} \to \varphi_{20})$}
\end{ptree}
\noindent$\Pi_{25}$:
\begin{ptree}
\AXC{$\varphi_{13} \to \forall \varepsilon\,(\varphi_{13})$\lbl{A4}}
\AXC{$\forall \varepsilon\,(\varphi_{13} \to \varphi_{20}) \to \forall \varepsilon\,(\varphi_{13}) \to \varphi_{25}$\lbl{A3}}
\AXC{$\stackrel{\Pi_{24}}{\vdots}$}\noLine
\UIC{$\forall \varepsilon\,(\varphi_{13} \to \varphi_{20})$}
\AXC{$\forall \varepsilon\,(\varphi_{13} \to \varphi_{20}) \to \varphi_{96}$\lbl{A1}}
\BIC{$\varphi_{96}$}
\BIC{$(\varphi_{13} \to \forall \varepsilon\,(\varphi_{13})) \to \mathrm{Seq}(c) \to \varphi_{47}$}
\BIC{$\mathrm{Seq}(c) \to \varphi_{47}$}
\end{ptree}
\noindent$\Pi_{26}$:
\begin{ptree}
\AXC{$\varphi_{30}$\lbl{DefConv}}
\AXC{$\varphi_{30} \to \forall \alpha\,(\mathrm{Conv}(c,\alpha) \leftrightarrow \mathrm{Seq}(c) \land \alpha \in \mathbb{R} \land \varphi_{25})$\lbl{A2}}
\BIC{$\forall \alpha\,(\mathrm{Conv}(c,\alpha) \leftrightarrow \mathrm{Seq}(c) \land \alpha \in \mathbb{R} \land \varphi_{25})$}
\AXC{$\varphi_{76}$\lbl{A2}}
\BIC{$\mathrm{Conv}(c,\alpha) \leftrightarrow \mathrm{Seq}(c) \land \alpha \in \mathbb{R} \land \varphi_{25}$}
\end{ptree}
\noindent$\Pi_{27}$:
\begin{ptree}
\AXC{$\stackrel{\Pi_{26}}{\vdots}$}\noLine
\UIC{$\mathrm{Conv}(c,\alpha) \leftrightarrow \mathrm{Seq}(c) \land \alpha \in \mathbb{R} \land \varphi_{25}$}
\AXC{$(\mathrm{Conv}(c,\alpha) \leftrightarrow \mathrm{Seq}(c) \land \alpha \in \mathbb{R} \land \varphi_{25}) \to \varphi_{82}$\lbl{A1}}
\BIC{$\varphi_{82}$}
\end{ptree}
\noindent 主木:
\begin{ptree}
\AXC{$\stackrel{\Pi_{25}}{\vdots}$}\noLine
\UIC{$\mathrm{Seq}(c) \to \varphi_{47}$}
\AXC{$\stackrel{\Pi_{4}}{\vdots}$}\noLine
\UIC{$\mathrm{Seq}(c) \to \varphi_{24}$}
\AXC{$\mathrm{Seq}(c) \to \mathrm{Seq}(c)$\lbl{A1}}
\AXC{$\stackrel{\Pi_{27}}{\vdots}$}\noLine
\UIC{$\varphi_{82}$}
\BIC{$(\mathrm{Seq}(c) \to \varphi_{24}) \to (\mathrm{Seq}(c) \to \varphi_{47}) \to \mathrm{Seq}(c) \to \varphi_{23}$}
\BIC{$(\mathrm{Seq}(c) \to \varphi_{47}) \to \mathrm{Seq}(c) \to \varphi_{23}$}
\BIC{$\mathrm{Seq}(c) \to \varphi_{23}$}
\UIC{$\forall \alpha\,(\mathrm{Seq}(c) \to \varphi_{23})$}
\UIC{$\forall c\,\forall \alpha\,(\mathrm{Seq}(c) \to \varphi_{23})$}
\UIC{$\forall b\,\forall c\,\forall \alpha\,(\mathrm{Seq}(c) \to \varphi_{23})$}
\UIC{$\forall a\,\forall b\,\forall c\,\forall \alpha\,(\mathrm{Seq}(c) \to \varphi_{23})$}
\end{ptree}
\begin{abbrevlist}
\item $\varphi_{1} :\equiv \forall \alpha\,(\mathrm{Conv}(b,\alpha) \to \mathrm{Seq}(b)) \to \mathrm{Conv}(b,\alpha) \to \mathrm{Seq}(b)$
\item $\varphi_{2} :\equiv \forall \alpha\,(\mathrm{Conv}(a,\alpha) \to \mathrm{Seq}(a)) \to \mathrm{Conv}(a,\alpha) \to \mathrm{Seq}(a)$
\item $\varphi_{3} :\equiv \forall a\,\forall \alpha\,(\mathrm{Conv}(a,\alpha) \to \mathrm{Seq}(a)) \to \forall \alpha\,(\mathrm{Conv}(b,\alpha) \to \mathrm{Seq}(b))$
\item $\varphi_{4} :\equiv \forall a\,\forall \alpha\,(\mathrm{Conv}(a,\alpha) \to \mathrm{Seq}(a)) \to \forall \alpha\,(\mathrm{Conv}(a,\alpha) \to \mathrm{Seq}(a))$
\item $\varphi_{5} :\equiv \forall \alpha\,(\mathrm{Conv}(a,\alpha) \to \alpha \in \mathbb{R}) \to \mathrm{Conv}(a,\alpha) \to \alpha \in \mathbb{R}$
\item $\varphi_{6} :\equiv \mathrm{Conv}(a,\alpha) \to \mathrm{Conv}(b,\alpha) \to \varepsilon \in \mathbb{R} \to 0 < \varepsilon \to 0 < \varepsilon$
\item $\varphi_{7} :\equiv \forall a\,\forall \alpha\,(\mathrm{Conv}(a,\alpha) \to \alpha \in \mathbb{R}) \to \forall \alpha\,(\mathrm{Conv}(a,\alpha) \to \alpha \in \mathbb{R})$
\item $\varphi_{8} :\equiv \forall N\,(\neg (N \in \mathbb{N} \land \forall n \in \mathbb{N}\,(N \leq n \to \lvert c_{n} - \alpha \rvert < \varepsilon)))$
\item $\varphi_{9} :\equiv \forall n \in \mathbb{N}\,(a_{n} \leq c_{n} \land c_{n} \leq b_{n}) \to \mathrm{Conv}(a,\alpha) \to \mathrm{Conv}(a,\alpha)$
\item $\varphi_{10} :\equiv \forall n\,(\mathrm{Seq}(c) \to n \in \mathbb{N} \to c_{n} \in \mathbb{R}) \to \mathrm{Seq}(c) \to n \in \mathbb{N} \to c_{n} \in \mathbb{R}$
\item $\varphi_{11} :\equiv \forall n\,(\mathrm{Seq}(b) \to n \in \mathbb{N} \to b_{n} \in \mathbb{R}) \to \mathrm{Seq}(b) \to n \in \mathbb{N} \to b_{n} \in \mathbb{R}$
\item $\varphi_{12} :\equiv \forall n\,(\mathrm{Seq}(a) \to n \in \mathbb{N} \to a_{n} \in \mathbb{R}) \to \mathrm{Seq}(a) \to n \in \mathbb{N} \to a_{n} \in \mathbb{R}$
\item $\varphi_{13} :\equiv \mathrm{Seq}(c) \land \forall n \in \mathbb{N}\,(a_{n} \leq c_{n} \land c_{n} \leq b_{n}) \land \mathrm{Conv}(a,\alpha) \land \mathrm{Conv}(b,\alpha)$
\item $\varphi_{14} :\equiv \forall n \in \mathbb{N}\,(a_{n} \leq c_{n} \land c_{n} \leq b_{n}) \to \mathrm{Conv}(a,\alpha) \to \mathrm{Conv}(b,\alpha) \to \mathrm{Seq}(b)$
\item $\varphi_{15} :\equiv \forall n \in \mathbb{N}\,(a_{n} \leq c_{n} \land c_{n} \leq b_{n}) \to \mathrm{Conv}(a,\alpha) \to \mathrm{Conv}(b,\alpha) \to \mathrm{Seq}(a)$
\item $\varphi_{16} :\equiv \forall a\,\forall n\,(\mathrm{Seq}(a) \to n \in \mathbb{N} \to a_{n} \in \mathbb{R}) \to \forall n\,(\mathrm{Seq}(c) \to n \in \mathbb{N} \to c_{n} \in \mathbb{R})$
\item $\varphi_{17} :\equiv \forall a\,\forall n\,(\mathrm{Seq}(a) \to n \in \mathbb{N} \to a_{n} \in \mathbb{R}) \to \forall n\,(\mathrm{Seq}(b) \to n \in \mathbb{N} \to b_{n} \in \mathbb{R})$
\item $\varphi_{18} :\equiv \forall a\,\forall n\,(\mathrm{Seq}(a) \to n \in \mathbb{N} \to a_{n} \in \mathbb{R}) \to \forall n\,(\mathrm{Seq}(a) \to n \in \mathbb{N} \to a_{n} \in \mathbb{R})$
\item $\varphi_{19} :\equiv N \in \mathbb{N} \land \forall n \in \mathbb{N}\,(N \leq n \to \lvert a_{n} - \alpha \rvert < \varepsilon \land \lvert b_{n} - \alpha \rvert < \varepsilon)$
\item $\varphi_{20} :\equiv \varepsilon \in \mathbb{R} \to 0 < \varepsilon \to \exists N \in \mathbb{N}\,\forall n \in \mathbb{N}\,(N \leq n \to \lvert c_{n} - \alpha \rvert < \varepsilon)$
\item $\varphi_{21} :\equiv \forall n \in \mathbb{N}\,(a_{n} \leq c_{n} \land c_{n} \leq b_{n}) \to \mathrm{Conv}(a,\alpha) \to \mathrm{Conv}(b,\alpha) \to \mathrm{Conv}(b,\alpha)$
\item $\varphi_{22} :\equiv \exists N \in \mathbb{N}\,\forall n \in \mathbb{N}\,(N \leq n \to \lvert a_{n} - \alpha \rvert < \varepsilon \land \lvert b_{n} - \alpha \rvert < \varepsilon)$
\item $\varphi_{23} :\equiv \forall n \in \mathbb{N}\,(a_{n} \leq c_{n} \land c_{n} \leq b_{n}) \to \mathrm{Conv}(a,\alpha) \to \mathrm{Conv}(b,\alpha) \to \mathrm{Conv}(c,\alpha)$
\item $\varphi_{24} :\equiv \forall n \in \mathbb{N}\,(a_{n} \leq c_{n} \land c_{n} \leq b_{n}) \to \mathrm{Conv}(a,\alpha) \to \mathrm{Conv}(b,\alpha) \to \alpha \in \mathbb{R}$
\item $\varphi_{25} :\equiv \forall \varepsilon \in \mathbb{R}\,(0 < \varepsilon \to \exists N \in \mathbb{N}\,\forall n \in \mathbb{N}\,(N \leq n \to \lvert c_{n} - \alpha \rvert < \varepsilon))$
\item $\varphi_{26} :\equiv \forall n \in \mathbb{N}\,(a_{n} \leq c_{n} \land c_{n} \leq b_{n}) \to n \in \mathbb{N} \to a_{n} \leq c_{n} \land c_{n} \leq b_{n}$
\item $\varphi_{27} :\equiv c_{n} \leq b_{n} \to \lvert a_{n} - \alpha \rvert < \varepsilon \to \lvert b_{n} - \alpha \rvert < \varepsilon \to \lvert c_{n} - \alpha \rvert < \varepsilon$
\item $\varphi_{28} :\equiv \forall n \in \mathbb{N}\,(a_{n} \leq c_{n} \land c_{n} \leq b_{n}) \to \forall n \in \mathbb{N}\,(a_{n} \leq c_{n} \land c_{n} \leq b_{n})$
\item $\varphi_{29} :\equiv \forall n \in \mathbb{N}\,(a_{n} \leq c_{n} \land c_{n} \leq b_{n}) \to \mathrm{Conv}(a,\alpha) \to \mathrm{Conv}(b,\alpha) \to \varepsilon \in \mathbb{R} \to \varepsilon \in \mathbb{R}$
\item $\varphi_{30} :\equiv \forall a\,\forall \alpha\,(\mathrm{Conv}(a,\alpha) \leftrightarrow \mathrm{Seq}(a) \land \alpha \in \mathbb{R} \land \forall \varepsilon \in \mathbb{R}\,(0 < \varepsilon \to \exists N \in \mathbb{N}\,\forall n \in \mathbb{N}\,(N \leq n \to \lvert a_{n} - \alpha \rvert < \varepsilon)))$
\item $\varphi_{31} :\equiv \forall \varepsilon\,(\mathrm{Conv}(a,\alpha) \to \mathrm{Conv}(b,\alpha) \to \varepsilon \in \mathbb{R} \to 0 < \varepsilon \to \varphi_{22})$
\item $\varphi_{32} :\equiv \forall n \in \mathbb{N}\,(N \leq n \to \lvert a_{n} - \alpha \rvert < \varepsilon \land \lvert b_{n} - \alpha \rvert < \varepsilon) \to n \in \mathbb{N} \to N \leq n \to \lvert a_{n} - \alpha \rvert < \varepsilon \land \lvert b_{n} - \alpha \rvert < \varepsilon$
\item $\varphi_{33} :\equiv \forall \beta\,\forall \varepsilon\,(\mathrm{Conv}(a,\alpha) \to \mathrm{Conv}(b,\beta) \to \varepsilon \in \mathbb{R} \to 0 < \varepsilon \to \exists N \in \mathbb{N}\,\forall n \in \mathbb{N}\,(N \leq n \to \lvert a_{n} - \alpha \rvert < \varepsilon \land \lvert b_{n} - \beta \rvert < \varepsilon))$
\item $\varphi_{34} :\equiv \mathrm{Conv}(a,\alpha) \to \mathrm{Conv}(b,\alpha) \to \varepsilon \in \mathbb{R} \to 0 < \varepsilon \to \varphi_{19} \to N \in \mathbb{N}$
\item $\varphi_{35} :\equiv \forall x\,\forall y\,\forall z\,\forall \alpha\,\forall \varepsilon\,(x \in \mathbb{R} \to y \in \mathbb{R} \to z \in \mathbb{R} \to \alpha \in \mathbb{R} \to \varepsilon \in \mathbb{R} \to x \leq z \to z \leq y \to \lvert x - \alpha \rvert < \varepsilon \to \lvert y - \alpha \rvert < \varepsilon \to \lvert z - \alpha \rvert < \varepsilon)$
\item $\varphi_{36} :\equiv \mathrm{Conv}(a,\alpha) \to \mathrm{Conv}(b,\alpha) \to \varepsilon \in \mathbb{R} \to 0 < \varepsilon \to \varphi_{19} \to n \in \mathbb{N} \to n \in \mathbb{N}$
\item $\varphi_{37} :\equiv \mathrm{Conv}(b,\alpha) \to \varepsilon \in \mathbb{R} \to 0 < \varepsilon \to \varphi_{19} \to n \in \mathbb{N} \to N \leq n \to c_{n} \leq b_{n}$
\item $\varphi_{38} :\equiv \mathrm{Conv}(b,\alpha) \to \varepsilon \in \mathbb{R} \to 0 < \varepsilon \to \varphi_{19} \to n \in \mathbb{N} \to N \leq n \to a_{n} \leq c_{n}$
\item $\varphi_{39} :\equiv \mathrm{Conv}(b,\alpha) \to \varepsilon \in \mathbb{R} \to 0 < \varepsilon \to \varphi_{19} \to n \in \mathbb{N} \to N \leq n \to c_{n} \in \mathbb{R}$
\item $\varphi_{40} :\equiv \mathrm{Conv}(b,\alpha) \to \varepsilon \in \mathbb{R} \to 0 < \varepsilon \to \varphi_{19} \to n \in \mathbb{N} \to N \leq n \to b_{n} \in \mathbb{R}$
\item $\varphi_{41} :\equiv \mathrm{Conv}(b,\alpha) \to \varepsilon \in \mathbb{R} \to 0 < \varepsilon \to \varphi_{19} \to n \in \mathbb{N} \to N \leq n \to a_{n} \in \mathbb{R}$
\item $\varphi_{42} :\equiv \forall n \in \mathbb{N}\,(a_{n} \leq c_{n} \land c_{n} \leq b_{n}) \to \mathrm{Conv}(a,\alpha) \to \mathrm{Conv}(b,\alpha) \to \varphi_{20}$
\item $\varphi_{43} :\equiv \varepsilon \in \mathbb{R} \to 0 < \varepsilon \to \varphi_{19} \to n \in \mathbb{N} \to N \leq n \to a_{n} \leq c_{n} \land c_{n} \leq b_{n}$
\item $\varphi_{44} :\equiv \mathrm{Conv}(b,\alpha) \to \varepsilon \in \mathbb{R} \to 0 < \varepsilon \to \varphi_{19} \to n \in \mathbb{N} \to N \leq n \to \lvert b_{n} - \alpha \rvert < \varepsilon$
\item $\varphi_{45} :\equiv \mathrm{Conv}(b,\alpha) \to \varepsilon \in \mathbb{R} \to 0 < \varepsilon \to \varphi_{19} \to n \in \mathbb{N} \to N \leq n \to \lvert a_{n} - \alpha \rvert < \varepsilon$
\item $\varphi_{46} :\equiv \mathrm{Conv}(b,\alpha) \to \varepsilon \in \mathbb{R} \to 0 < \varepsilon \to \varphi_{19} \to n \in \mathbb{N} \to N \leq n \to \lvert c_{n} - \alpha \rvert < \varepsilon$
\item $\varphi_{47} :\equiv \forall n \in \mathbb{N}\,(a_{n} \leq c_{n} \land c_{n} \leq b_{n}) \to \mathrm{Conv}(a,\alpha) \to \mathrm{Conv}(b,\alpha) \to \varphi_{25}$
\item $\varphi_{48} :\equiv 0 < \varepsilon \to \varphi_{19} \to n \in \mathbb{N} \to N \leq n \to \lvert a_{n} - \alpha \rvert < \varepsilon \land \lvert b_{n} - \alpha \rvert < \varepsilon$
\item $\varphi_{49} :\equiv \varepsilon \in \mathbb{R} \to 0 < \varepsilon \to \varphi_{19} \to N \in \mathbb{N} \land \forall n \in \mathbb{N}\,(N \leq n \to \lvert c_{n} - \alpha \rvert < \varepsilon)$
\item $\varphi_{50} :\equiv \mathrm{Conv}(b,\alpha) \to \varepsilon \in \mathbb{R} \to 0 < \varepsilon \to \varphi_{19} \to \forall n \in \mathbb{N}\,(N \leq n \to \lvert c_{n} - \alpha \rvert < \varepsilon)$
\item $\varphi_{51} :\equiv 0 < \varepsilon \to \varphi_{19} \to \forall n \in \mathbb{N}\,(N \leq n \to \lvert a_{n} - \alpha \rvert < \varepsilon \land \lvert b_{n} - \alpha \rvert < \varepsilon)$
\item $\varphi_{52} :\equiv \forall y\,\forall z\,\forall \alpha\,\forall \varepsilon\,(a_{n} \in \mathbb{R} \to y \in \mathbb{R} \to z \in \mathbb{R} \to \alpha \in \mathbb{R} \to \varepsilon \in \mathbb{R} \to a_{n} \leq z \to z \leq y \to \lvert a_{n} - \alpha \rvert < \varepsilon \to \lvert y - \alpha \rvert < \varepsilon \to \lvert z - \alpha \rvert < \varepsilon)$
\item $\varphi_{53} :\equiv \varphi_{13} \land \varepsilon \in \mathbb{R} \land 0 < \varepsilon \land \neg \exists N \in \mathbb{N}\,\forall n \in \mathbb{N}\,(N \leq n \to \lvert c_{n} - \alpha \rvert < \varepsilon)$
\item $\varphi_{54} :\equiv \varepsilon \in \mathbb{R} \to 0 < \varepsilon \to \varphi_{19} \to \exists N \in \mathbb{N}\,\forall n \in \mathbb{N}\,(N \leq n \to \lvert c_{n} - \alpha \rvert < \varepsilon)$
\item $\varphi_{55} :\equiv \forall z\,\forall \alpha\,\forall \varepsilon\,(a_{n} \in \mathbb{R} \to b_{n} \in \mathbb{R} \to z \in \mathbb{R} \to \alpha \in \mathbb{R} \to \varepsilon \in \mathbb{R} \to a_{n} \leq z \to z \leq b_{n} \to \lvert a_{n} - \alpha \rvert < \varepsilon \to \lvert b_{n} - \alpha \rvert < \varepsilon \to \lvert z - \alpha \rvert < \varepsilon)$
\item $\varphi_{56} :\equiv \forall n \in \mathbb{N}\,(a_{n} \leq c_{n} \land c_{n} \leq b_{n}) \to \mathrm{Conv}(a,\alpha) \to \mathrm{Conv}(b,\alpha) \to \varepsilon \in \mathbb{R} \to 0 < \varepsilon \to \varphi_{22}$
\item $\varphi_{57} :\equiv (\mathrm{Conv}(a,\alpha) \to \mathrm{Seq}(a)) \to (\mathrm{Seq}(c) \to \varphi_{9}) \to \mathrm{Seq}(c) \to \varphi_{15}$
\item $\varphi_{58} :\equiv a_{n} \in \mathbb{R} \to b_{n} \in \mathbb{R} \to c_{n} \in \mathbb{R} \to \alpha \in \mathbb{R} \to \varepsilon \in \mathbb{R} \to a_{n} \leq c_{n} \to \varphi_{27}$
\item $\varphi_{59} :\equiv (\mathrm{Conv}(a,\alpha) \to \alpha \in \mathbb{R}) \to (\mathrm{Seq}(c) \to \varphi_{9}) \to \mathrm{Seq}(c) \to \varphi_{24}$
\item $\varphi_{60} :\equiv (\mathrm{Conv}(b,\alpha) \to \mathrm{Seq}(b)) \to (\mathrm{Seq}(c) \to \varphi_{21}) \to \mathrm{Seq}(c) \to \varphi_{14}$
\item $\varphi_{61} :\equiv \mathrm{Conv}(a,\alpha) \to \mathrm{Conv}(b,\alpha) \to \varepsilon \in \mathbb{R} \to 0 < \varepsilon \to \varphi_{19} \to \varphi_{19}$
\item $\varphi_{62} :\equiv \mathrm{Seq}(c) \to \forall n \in \mathbb{N}\,(a_{n} \leq c_{n} \land c_{n} \leq b_{n}) \to \mathrm{Conv}(a,\alpha) \to \varphi_{37}$
\item $\varphi_{63} :\equiv \mathrm{Seq}(c) \to \forall n \in \mathbb{N}\,(a_{n} \leq c_{n} \land c_{n} \leq b_{n}) \to \mathrm{Conv}(a,\alpha) \to \varphi_{38}$
\item $\varphi_{64} :\equiv \mathrm{Seq}(c) \to \forall n \in \mathbb{N}\,(a_{n} \leq c_{n} \land c_{n} \leq b_{n}) \to \mathrm{Conv}(a,\alpha) \to \varphi_{39}$
\item $\varphi_{65} :\equiv \mathrm{Seq}(c) \to \forall n \in \mathbb{N}\,(a_{n} \leq c_{n} \land c_{n} \leq b_{n}) \to \mathrm{Conv}(a,\alpha) \to \varphi_{40}$
\item $\varphi_{66} :\equiv \mathrm{Seq}(c) \to \forall n \in \mathbb{N}\,(a_{n} \leq c_{n} \land c_{n} \leq b_{n}) \to \mathrm{Conv}(a,\alpha) \to \varphi_{41}$
\item $\varphi_{67} :\equiv \mathrm{Seq}(c) \to \forall n \in \mathbb{N}\,(a_{n} \leq c_{n} \land c_{n} \leq b_{n}) \to \mathrm{Conv}(a,\alpha) \to \varphi_{44}$
\item $\varphi_{68} :\equiv \mathrm{Seq}(c) \to \forall n \in \mathbb{N}\,(a_{n} \leq c_{n} \land c_{n} \leq b_{n}) \to \mathrm{Conv}(a,\alpha) \to \varphi_{45}$
\item $\varphi_{69} :\equiv \mathrm{Seq}(c) \to \forall n \in \mathbb{N}\,(a_{n} \leq c_{n} \land c_{n} \leq b_{n}) \to \mathrm{Conv}(a,\alpha) \to \varphi_{46}$
\item $\varphi_{70} :\equiv \forall n \in \mathbb{N}\,(a_{n} \leq c_{n} \land c_{n} \leq b_{n}) \to \mathrm{Conv}(a,\alpha) \to \mathrm{Conv}(b,\alpha) \to \varphi_{43}$
\item $\varphi_{71} :\equiv \forall n\,(\varphi_{13} \land \varepsilon \in \mathbb{R} \land 0 < \varepsilon \land \varphi_{19} \to n \in \mathbb{N} \to N \leq n \to \lvert c_{n} - \alpha \rvert < \varepsilon)$
\item $\varphi_{72} :\equiv \mathrm{Seq}(c) \to \forall n \in \mathbb{N}\,(a_{n} \leq c_{n} \land c_{n} \leq b_{n}) \to \mathrm{Conv}(a,\alpha) \to \varphi_{50}$
\item $\varphi_{73} :\equiv \forall n\,(\varphi_{13} \land \varepsilon \in \mathbb{R} \land 0 < \varepsilon \land \varphi_{19}) \to \forall n \in \mathbb{N}\,(N \leq n \to \lvert c_{n} - \alpha \rvert < \varepsilon)$
\item $\varphi_{74} :\equiv \forall n \in \mathbb{N}\,(a_{n} \leq c_{n} \land c_{n} \leq b_{n}) \to \mathrm{Conv}(a,\alpha) \to \mathrm{Conv}(b,\alpha) \to \varphi_{49}$
\item $\varphi_{75} :\equiv \forall n \in \mathbb{N}\,(a_{n} \leq c_{n} \land c_{n} \leq b_{n}) \to \mathrm{Conv}(a,\alpha) \to \mathrm{Conv}(b,\alpha) \to \varphi_{54}$
\item $\varphi_{76} :\equiv \forall \alpha\,(\mathrm{Conv}(c,\alpha) \leftrightarrow \mathrm{Seq}(c) \land \alpha \in \mathbb{R} \land \varphi_{25}) \to (\mathrm{Conv}(c,\alpha) \leftrightarrow \mathrm{Seq}(c) \land \alpha \in \mathbb{R} \land \varphi_{25})$
\item $\varphi_{77} :\equiv \forall n \in \mathbb{N}\,(a_{n} \leq c_{n} \land c_{n} \leq b_{n}) \to \mathrm{Conv}(a,\alpha) \to \mathrm{Conv}(b,\alpha) \to \varepsilon \in \mathbb{R} \to \varphi_{48}$
\item $\varphi_{78} :\equiv \forall n \in \mathbb{N}\,(a_{n} \leq c_{n} \land c_{n} \leq b_{n}) \to \mathrm{Conv}(a,\alpha) \to \mathrm{Conv}(b,\alpha) \to \varepsilon \in \mathbb{R} \to \varphi_{51}$
\item $\varphi_{79} :\equiv \varphi_{31} \to \mathrm{Conv}(a,\alpha) \to \mathrm{Conv}(b,\alpha) \to \varepsilon \in \mathbb{R} \to 0 < \varepsilon \to \varphi_{22}$
\item $\varphi_{80} :\equiv (\mathrm{Seq}(c) \to \forall n \in \mathbb{N}\,(a_{n} \leq c_{n} \land c_{n} \leq b_{n}) \to \varphi_{6}) \to \mathrm{Seq}(c) \to \varphi_{56}$
\item $\varphi_{81} :\equiv (\varphi_{8} \to \neg (N \in \mathbb{N} \land \forall n \in \mathbb{N}\,(N \leq n \to \lvert c_{n} - \alpha \rvert < \varepsilon))) \to \mathrm{Seq}(c) \to \varphi_{75}$
\item $\varphi_{82} :\equiv (\mathrm{Seq}(c) \to \mathrm{Seq}(c)) \to (\mathrm{Seq}(c) \to \varphi_{24}) \to (\mathrm{Seq}(c) \to \varphi_{47}) \to \mathrm{Seq}(c) \to \varphi_{23}$
\item $\varphi_{83} :\equiv (\mathrm{Seq}(c) \to \forall n \in \mathbb{N}\,(a_{n} \leq c_{n} \land c_{n} \leq b_{n}) \to \varphi_{36}) \to \varphi_{64}$
\item $\varphi_{84} :\equiv (\mathrm{Seq}(c) \to \forall n \in \mathbb{N}\,(a_{n} \leq c_{n} \land c_{n} \leq b_{n}) \to \varphi_{36}) \to \varphi_{65}$
\item $\varphi_{85} :\equiv (\mathrm{Seq}(c) \to \forall n \in \mathbb{N}\,(a_{n} \leq c_{n} \land c_{n} \leq b_{n}) \to \varphi_{36}) \to \varphi_{66}$
\item $\varphi_{86} :\equiv (\mathrm{Seq}(c) \to \forall n \in \mathbb{N}\,(a_{n} \leq c_{n} \land c_{n} \leq b_{n}) \to \varphi_{36}) \to \mathrm{Seq}(c) \to \varphi_{70}$
\item $\varphi_{87} :\equiv \varphi_{69} \to \varphi_{13} \land \varepsilon \in \mathbb{R} \land 0 < \varepsilon \land \varphi_{19} \to n \in \mathbb{N} \to N \leq n \to \lvert c_{n} - \alpha \rvert < \varepsilon$
\item $\varphi_{88} :\equiv (\mathrm{Seq}(c) \to \forall n \in \mathbb{N}\,(a_{n} \leq c_{n} \land c_{n} \leq b_{n}) \to \varphi_{61}) \to \mathrm{Seq}(c) \to \forall n \in \mathbb{N}\,(a_{n} \leq c_{n} \land c_{n} \leq b_{n}) \to \varphi_{34}$
\item $\varphi_{89} :\equiv (\mathrm{Seq}(c) \to n \in \mathbb{N} \to c_{n} \in \mathbb{R}) \to (\mathrm{Seq}(c) \to \mathrm{Seq}(c)) \to \varphi_{83}$
\item $\varphi_{90} :\equiv (\mathrm{Seq}(c) \to \forall n \in \mathbb{N}\,(a_{n} \leq c_{n} \land c_{n} \leq b_{n}) \to \varphi_{61}) \to \mathrm{Seq}(c) \to \varphi_{78}$
\item $\varphi_{91} :\equiv (\mathrm{Seq}(b) \to n \in \mathbb{N} \to b_{n} \in \mathbb{R}) \to (\mathrm{Seq}(c) \to \varphi_{14}) \to \varphi_{84}$
\item $\varphi_{92} :\equiv (\mathrm{Seq}(a) \to n \in \mathbb{N} \to a_{n} \in \mathbb{R}) \to (\mathrm{Seq}(c) \to \varphi_{15}) \to \varphi_{85}$
\item $\varphi_{93} :\equiv (\mathrm{Seq}(c) \to \varphi_{9}) \to (\mathrm{Seq}(c) \to \varphi_{21}) \to (\mathrm{Seq}(c) \to \varphi_{29}) \to \varphi_{80}$
\item $\varphi_{94} :\equiv (\varphi_{13} \land \varepsilon \in \mathbb{R} \land 0 < \varepsilon \land \varphi_{19} \to \forall n\,(\varphi_{13} \land \varepsilon \in \mathbb{R} \land 0 < \varepsilon \land \varphi_{19})) \to \varphi_{72}$
\item $\varphi_{95} :\equiv (\mathrm{Seq}(c) \to \forall n \in \mathbb{N}\,(a_{n} \leq c_{n} \land c_{n} \leq b_{n}) \to \varphi_{34}) \to \varphi_{72} \to \mathrm{Seq}(c) \to \varphi_{74}$
\item $\varphi_{96} :\equiv (\forall \varepsilon\,(\varphi_{13} \to \varphi_{20}) \to \forall \varepsilon\,(\varphi_{13}) \to \varphi_{25}) \to (\varphi_{13} \to \forall \varepsilon\,(\varphi_{13})) \to \mathrm{Seq}(c) \to \varphi_{47}$
\item $\varphi_{97} :\equiv (\mathrm{Conv}(a,\alpha) \to \mathrm{Conv}(b,\alpha) \to \varepsilon \in \mathbb{R} \to 0 < \varepsilon \to \varphi_{22}) \to \varphi_{93}$
\item $\varphi_{98} :\equiv (\forall N\,(\varphi_{53} \to \neg \varphi_{19}) \to \forall N\,(\varphi_{53}) \to \forall N\,(\neg \varphi_{19})) \to (\varphi_{53} \to \forall N\,(\varphi_{53})) \to \mathrm{Seq}(c) \to \varphi_{42}$
\item $\varphi_{99} :\equiv (\mathrm{Seq}(c) \to \varphi_{24}) \to (\mathrm{Seq}(c) \to \varphi_{29}) \to \varphi_{63} \to \varphi_{62} \to \varphi_{68} \to \varphi_{67} \to \varphi_{69}$
\end{abbrevlist}
\end{fullproof}

% 自動生成 (prover/build.js): LimLe  ノード数 568
\begin{fullproof}{定理2.5}{$\forall a\,\forall b\,\forall \alpha\,\forall \beta\,(\mathrm{Conv}(a,\alpha) \to \mathrm{Conv}(b,\beta) \to \forall n \in \mathbb{N}\,(a_{n} \leq b_{n}) \to \alpha \leq \beta)$}{568}\phantomsection\label{tree:LimLe}
\noindent$\Pi_{1}$:
\begin{ptree}
\AXC{$z = w \to \alpha - \beta < z \to \alpha - \beta < w$\lbl{A6}}
\UIC{$\forall w\,(z = w \to \alpha - \beta < z \to \alpha - \beta < w)$}
\UIC{$\forall z\,\forall w\,(z = w \to \alpha - \beta < z \to \alpha - \beta < w)$}
\AXC{$\forall z\,\forall w\,(z = w \to \alpha - \beta < z \to \alpha - \beta < w) \to \varphi_{93}$\lbl{A2}}
\BIC{$\varphi_{93}$}
\end{ptree}
\noindent$\Pi_{2}$:
\begin{ptree}
\AXC{$z = w \to z = (\alpha - \beta) \cdot 2^{-1} \to w = (\alpha - \beta) \cdot 2^{-1}$\lbl{A6}}
\UIC{$\varphi_{3}$}
\UIC{$\forall z\,(\varphi_{3})$}
\AXC{$\forall z\,(\varphi_{3}) \to \varphi_{105}$\lbl{A2}}
\BIC{$\varphi_{105}$}
\AXC{$\varphi_{105} \to a_{N} - b_{N} + (b_{N} - \beta) = a_{N} - \beta \to \varphi_{60}$\lbl{A2}}
\BIC{$a_{N} - b_{N} + (b_{N} - \beta) = a_{N} - \beta \to \varphi_{60}$}
\end{ptree}
\noindent$\Pi_{3}$:
\begin{ptree}
\AXC{$z = w \to z < (\alpha - \beta) \cdot 2^{-1} \to w < (\alpha - \beta) \cdot 2^{-1}$\lbl{A6}}
\UIC{$\varphi_{9}$}
\UIC{$\forall z\,(\varphi_{9})$}
\AXC{$\forall z\,(\varphi_{9}) \to \varphi_{110}$\lbl{A2}}
\BIC{$\varphi_{110}$}
\AXC{$\varphi_{110} \to a_{N} - b_{N} + (b_{N} - \beta) = a_{N} - \beta \to \varphi_{70}$\lbl{A2}}
\BIC{$a_{N} - b_{N} + (b_{N} - \beta) = a_{N} - \beta \to \varphi_{70}$}
\end{ptree}
\noindent$\Pi_{4}$:
\begin{ptree}
\AXC{$\varphi_{50}$\lbl{A6}}
\UIC{$\forall w\,(\varphi_{50})$}
\UIC{$\forall z\,\forall w\,(\varphi_{50})$}
\AXC{$\forall z\,\forall w\,(\varphi_{50}) \to \varphi_{111}$\lbl{A2}}
\BIC{$\varphi_{111}$}
\AXC{$\varphi_{111} \to 0 + (\alpha - \beta) \cdot 2^{-1} = (\alpha - \beta) \cdot 2^{-1} \to \varphi_{103}$\lbl{A2}}
\BIC{$0 + (\alpha - \beta) \cdot 2^{-1} = (\alpha - \beta) \cdot 2^{-1} \to \varphi_{103}$}
\end{ptree}
\noindent$\Pi_{5}$:
\begin{ptree}
\AXC{$\mathrm{Conv}(a,\alpha) \to \varphi_{195}$\lbl{A1}}
\AXC{$(\mathrm{Conv}(a,\alpha) \to \varphi_{195}) \to \mathrm{Conv}(a,\alpha) \to \varphi_{143}$\lbl{A1}}
\BIC{$\mathrm{Conv}(a,\alpha) \to \varphi_{143}$}
\end{ptree}
\noindent$\Pi_{6}$:
\begin{ptree}
\AXC{$\mathrm{Conv}(a,\alpha) \to \varphi_{195}$\lbl{A1}}
\AXC{$(\mathrm{Conv}(a,\alpha) \to \varphi_{195}) \to \mathrm{Conv}(a,\alpha) \to \varphi_{193}$\lbl{A1}}
\BIC{$\mathrm{Conv}(a,\alpha) \to \varphi_{193}$}
\AXC{$(\mathrm{Conv}(a,\alpha) \to \varphi_{193}) \to \varphi_{228}$\lbl{A1}}
\BIC{$\varphi_{228}$}
\end{ptree}
\noindent$\Pi_{7}$:
\begin{ptree}
\AXC{$\stackrel{\Pi_{5}}{\vdots}$}\noLine
\UIC{$\mathrm{Conv}(a,\alpha) \to \varphi_{143}$}
\AXC{$\varphi_{78} \to N \in \mathbb{N} \to N \leq N \to \varphi_{45}$\lbl{A2}}
\AXC{$\stackrel{\Pi_{6}}{\vdots}$}\noLine
\UIC{$\varphi_{228}$}
\BIC{$\mathrm{Conv}(a,\alpha) \to \mathrm{Conv}(b,\beta) \to \varphi_{191}$}
\AXC{$\varphi_{246}$\lbl{A1}}
\BIC{$(\mathrm{Conv}(a,\alpha) \to \varphi_{143}) \to N = N \to \mathrm{Conv}(a,\alpha) \to \varphi_{190}$}
\BIC{$N = N \to \mathrm{Conv}(a,\alpha) \to \varphi_{190}$}
\end{ptree}
\noindent$\Pi_{8}$:
\begin{ptree}
\AXC{$x = x$\lbl{A5}}
\UIC{$\forall x\,(x = x)$}
\AXC{$\forall x\,(x = x) \to N = N$\lbl{A2}}
\BIC{$N = N$}
\AXC{$\stackrel{\Pi_{7}}{\vdots}$}\noLine
\UIC{$N = N \to \mathrm{Conv}(a,\alpha) \to \varphi_{190}$}
\BIC{$\mathrm{Conv}(a,\alpha) \to \varphi_{190}$}
\end{ptree}
\noindent$\Pi_{9}$:
\begin{ptree}
\AXC{$\forall a\,\forall \alpha\,(\mathrm{Conv}(a,\alpha) \to \alpha \in \mathbb{R})$\lbl{[C2]}}
\AXC{$\varphi_{10}$\lbl{A2}}
\BIC{$\forall \alpha\,(\mathrm{Conv}(b,\alpha) \to \alpha \in \mathbb{R})$}
\AXC{$\varphi_{6}$\lbl{A2}}
\BIC{$\mathrm{Conv}(b,\beta) \to \beta \in \mathbb{R}$}
\AXC{$(\mathrm{Conv}(b,\beta) \to \beta \in \mathbb{R}) \to \varphi_{86}$\lbl{A1}}
\BIC{$\varphi_{86}$}
\end{ptree}
\noindent$\Pi_{10}$:
\begin{ptree}
\AXC{$\mathrm{Conv}(a,\alpha) \to \mathrm{Conv}(b,\beta) \to \mathrm{Conv}(b,\beta)$\lbl{A1}}
\AXC{$\stackrel{\Pi_{9}}{\vdots}$}\noLine
\UIC{$\varphi_{86}$}
\BIC{$\varphi_{23}$}
\end{ptree}
\noindent$\Pi_{11}$:
\begin{ptree}
\AXC{$\mathrm{Conv}(a,\alpha) \to \mathrm{Conv}(a,\alpha)$\lbl{A1}}
\AXC{$\forall a\,\forall \alpha\,(\mathrm{Conv}(a,\alpha) \to \alpha \in \mathbb{R})$\lbl{[C2]}}
\AXC{$\varphi_{11}$\lbl{A2}}
\BIC{$\forall \alpha\,(\mathrm{Conv}(a,\alpha) \to \alpha \in \mathbb{R})$}
\AXC{$\varphi_{7}$\lbl{A2}}
\BIC{$\mathrm{Conv}(a,\alpha) \to \alpha \in \mathbb{R}$}
\AXC{$\varphi_{104}$\lbl{A1}}
\BIC{$(\mathrm{Conv}(a,\alpha) \to \mathrm{Conv}(a,\alpha)) \to \varphi_{24}$}
\BIC{$\varphi_{24}$}
\end{ptree}
\noindent$\Pi_{12}$:
\begin{ptree}
\AXC{$\stackrel{\Pi_{10}}{\vdots}$}\noLine
\UIC{$\varphi_{23}$}
\AXC{$\stackrel{\Pi_{11}}{\vdots}$}\noLine
\UIC{$\varphi_{24}$}
\AXC{$\forall x\,\forall y\,(x \in \mathbb{R} \to y \in \mathbb{R} \to x - y \in \mathbb{R})$\lbl{[F1]}}
\AXC{$\varphi_{67}$\lbl{A2}}
\BIC{$\forall y\,(\alpha \in \mathbb{R} \to y \in \mathbb{R} \to \alpha - y \in \mathbb{R})$}
\AXC{$\varphi_{61}$\lbl{A2}}
\BIC{$\alpha \in \mathbb{R} \to \beta \in \mathbb{R} \to \alpha - \beta \in \mathbb{R}$}
\AXC{$\varphi_{203}$\lbl{A1}}
\BIC{$\varphi_{24} \to \varphi_{23} \to \mathrm{Conv}(a,\alpha) \to \varphi_{155}$}
\BIC{$\varphi_{23} \to \mathrm{Conv}(a,\alpha) \to \varphi_{155}$}
\BIC{$\mathrm{Conv}(a,\alpha) \to \varphi_{155}$}
\end{ptree}
\noindent$\Pi_{13}$:
\begin{ptree}
\AXC{$\forall x\,\forall y\,(x \in \mathbb{R} \to y \in \mathbb{R} \to x \cdot y \in \mathbb{R})$\lbl{(R1)}}
\AXC{$\varphi_{82}$\lbl{A2}}
\BIC{$\forall y\,(\alpha - \beta \in \mathbb{R} \to y \in \mathbb{R} \to (\alpha - \beta) \cdot y \in \mathbb{R})$}
\AXC{$\varphi_{106}$\lbl{A2}}
\BIC{$\alpha - \beta \in \mathbb{R} \to 2^{-1} \in \mathbb{R} \to (\alpha - \beta) \cdot 2^{-1} \in \mathbb{R}$}
\end{ptree}
\noindent$\Pi_{14}$:
\begin{ptree}
\AXC{$\stackrel{\Pi_{13}}{\vdots}$}\noLine
\UIC{$\alpha - \beta \in \mathbb{R} \to 2^{-1} \in \mathbb{R} \to (\alpha - \beta) \cdot 2^{-1} \in \mathbb{R}$}
\AXC{$(\alpha - \beta \in \mathbb{R} \to 2^{-1} \in \mathbb{R} \to (\alpha - \beta) \cdot 2^{-1} \in \mathbb{R}) \to \varphi_{214}$\lbl{A1}}
\BIC{$\varphi_{214}$}
\end{ptree}
\noindent$\Pi_{15}$:
\begin{ptree}
\AXC{$2^{-1} \in \mathbb{R}$\lbl{[F31]}}
\AXC{$\stackrel{\Pi_{12}}{\vdots}$}\noLine
\UIC{$\mathrm{Conv}(a,\alpha) \to \varphi_{155}$}
\AXC{$\stackrel{\Pi_{14}}{\vdots}$}\noLine
\UIC{$\varphi_{214}$}
\BIC{$2^{-1} \in \mathbb{R} \to \mathrm{Conv}(a,\alpha) \to \mathrm{Conv}(b,\beta) \to \varphi_{151}$}
\BIC{$\mathrm{Conv}(a,\alpha) \to \mathrm{Conv}(b,\beta) \to \varphi_{151}$}
\end{ptree}
\noindent$\Pi_{16}$:
\begin{ptree}
\AXC{$\forall a\,\forall \alpha\,(\mathrm{Conv}(a,\alpha) \to \mathrm{Seq}(a))$\lbl{[C1]}}
\AXC{$\varphi_{4}$\lbl{A2}}
\BIC{$\forall \alpha\,(\mathrm{Conv}(b,\alpha) \to \mathrm{Seq}(b))$}
\AXC{$\varphi_{1}$\lbl{A2}}
\BIC{$\mathrm{Conv}(b,\beta) \to \mathrm{Seq}(b)$}
\AXC{$(\mathrm{Conv}(b,\beta) \to \mathrm{Seq}(b)) \to \varphi_{75}$\lbl{A1}}
\BIC{$\varphi_{75}$}
\end{ptree}
\noindent$\Pi_{17}$:
\begin{ptree}
\AXC{$\forall a\,\forall n\,(\mathrm{Seq}(a) \to n \in \mathbb{N} \to a_{n} \in \mathbb{R})$\lbl{[S1]}}
\AXC{$\varphi_{48}$\lbl{A2}}
\BIC{$\forall n\,(\mathrm{Seq}(b) \to n \in \mathbb{N} \to b_{n} \in \mathbb{R})$}
\AXC{$\varphi_{40}$\lbl{A2}}
\BIC{$\mathrm{Seq}(b) \to N \in \mathbb{N} \to b_{N} \in \mathbb{R}$}
\AXC{$\varphi_{232}$\lbl{A1}}
\BIC{$\varphi_{14} \to (\mathrm{Conv}(a,\alpha) \to \varphi_{143}) \to \mathrm{Conv}(a,\alpha) \to \varphi_{156}$}
\end{ptree}
\noindent$\Pi_{18}$:
\begin{ptree}
\AXC{$\mathrm{Conv}(a,\alpha) \to \mathrm{Conv}(b,\beta) \to \mathrm{Conv}(b,\beta)$\lbl{A1}}
\AXC{$\stackrel{\Pi_{16}}{\vdots}$}\noLine
\UIC{$\varphi_{75}$}
\BIC{$\varphi_{14}$}
\AXC{$\stackrel{\Pi_{17}}{\vdots}$}\noLine
\UIC{$\varphi_{14} \to (\mathrm{Conv}(a,\alpha) \to \varphi_{143}) \to \mathrm{Conv}(a,\alpha) \to \varphi_{156}$}
\BIC{$(\mathrm{Conv}(a,\alpha) \to \varphi_{143}) \to \mathrm{Conv}(a,\alpha) \to \varphi_{156}$}
\end{ptree}
\noindent$\Pi_{19}$:
\begin{ptree}
\AXC{$\stackrel{\Pi_{5}}{\vdots}$}\noLine
\UIC{$\mathrm{Conv}(a,\alpha) \to \varphi_{143}$}
\AXC{$\stackrel{\Pi_{18}}{\vdots}$}\noLine
\UIC{$(\mathrm{Conv}(a,\alpha) \to \varphi_{143}) \to \mathrm{Conv}(a,\alpha) \to \varphi_{156}$}
\BIC{$\mathrm{Conv}(a,\alpha) \to \varphi_{156}$}
\end{ptree}
\noindent$\Pi_{20}$:
\begin{ptree}
\AXC{$\forall x\,\forall y\,(x \in \mathbb{R} \to y \in \mathbb{R} \to x - y \in \mathbb{R})$\lbl{[F1]}}
\AXC{$\varphi_{83}$\lbl{A2}}
\BIC{$\forall y\,(b_{N} \in \mathbb{R} \to y \in \mathbb{R} \to b_{N} - y \in \mathbb{R})$}
\AXC{$\varphi_{97}$\lbl{A2}}
\BIC{$b_{N} \in \mathbb{R} \to \beta \in \mathbb{R} \to b_{N} - \beta \in \mathbb{R}$}
\AXC{$(b_{N} \in \mathbb{R} \to \beta \in \mathbb{R} \to b_{N} - \beta \in \mathbb{R}) \to \varphi_{220}$\lbl{A1}}
\BIC{$\varphi_{220}$}
\end{ptree}
\noindent$\Pi_{21}$:
\begin{ptree}
\AXC{$\stackrel{\Pi_{10}}{\vdots}$}\noLine
\UIC{$\varphi_{23}$}
\AXC{$\stackrel{\Pi_{19}}{\vdots}$}\noLine
\UIC{$\mathrm{Conv}(a,\alpha) \to \varphi_{156}$}
\AXC{$\stackrel{\Pi_{20}}{\vdots}$}\noLine
\UIC{$\varphi_{220}$}
\BIC{$\varphi_{23} \to \mathrm{Conv}(a,\alpha) \to \mathrm{Conv}(b,\beta) \to \varphi_{146}$}
\BIC{$\mathrm{Conv}(a,\alpha) \to \mathrm{Conv}(b,\beta) \to \varphi_{146}$}
\end{ptree}
\noindent$\Pi_{22}$:
\begin{ptree}
\AXC{$\forall x\,\forall y\,(x \in \mathbb{R} \to y \in \mathbb{R} \to \lvert x \rvert < y \to -y < x \land x < y)$\lbl{[A11]}}
\AXC{$\varphi_{138}$\lbl{A2}}
\BIC{$\varphi_{90}$}
\AXC{$\varphi_{90} \to b_{N} - \beta \in \mathbb{R} \to \varphi_{119}$\lbl{A2}}
\BIC{$b_{N} - \beta \in \mathbb{R} \to \varphi_{119}$}
\end{ptree}
\noindent$\Pi_{23}$:
\begin{ptree}
\AXC{$\stackrel{\Pi_{21}}{\vdots}$}\noLine
\UIC{$\mathrm{Conv}(a,\alpha) \to \mathrm{Conv}(b,\beta) \to \varphi_{146}$}
\AXC{$\stackrel{\Pi_{22}}{\vdots}$}\noLine
\UIC{$b_{N} - \beta \in \mathbb{R} \to \varphi_{119}$}
\AXC{$(b_{N} - \beta \in \mathbb{R} \to \varphi_{119}) \to \varphi_{249}$\lbl{A1}}
\BIC{$\varphi_{249}$}
\BIC{$(\mathrm{Conv}(a,\alpha) \to \mathrm{Conv}(b,\beta) \to \varphi_{151}) \to \varphi_{225}$}
\end{ptree}
\noindent$\Pi_{24}$:
\begin{ptree}
\AXC{$\stackrel{\Pi_{15}}{\vdots}$}\noLine
\UIC{$\mathrm{Conv}(a,\alpha) \to \mathrm{Conv}(b,\beta) \to \varphi_{151}$}
\AXC{$\stackrel{\Pi_{23}}{\vdots}$}\noLine
\UIC{$(\mathrm{Conv}(a,\alpha) \to \mathrm{Conv}(b,\beta) \to \varphi_{151}) \to \varphi_{225}$}
\BIC{$\varphi_{225}$}
\end{ptree}
\noindent$\Pi_{25}$:
\begin{ptree}
\AXC{$\stackrel{\Pi_{8}}{\vdots}$}\noLine
\UIC{$\mathrm{Conv}(a,\alpha) \to \varphi_{190}$}
\AXC{$\varphi_{230}$\lbl{A1}}
\BIC{$\mathrm{Conv}(a,\alpha) \to \mathrm{Conv}(b,\beta) \to \varphi_{166}$}
\AXC{$\stackrel{\Pi_{24}}{\vdots}$}\noLine
\UIC{$\varphi_{225}$}
\BIC{$\mathrm{Conv}(a,\alpha) \to \varphi_{188}$}
\AXC{$\varphi_{223}$\lbl{A1}}
\BIC{$\mathrm{Conv}(a,\alpha) \to \mathrm{Conv}(b,\beta) \to \varphi_{164}$}
\AXC{$\varphi_{216}$\lbl{A1}}
\BIC{$\mathrm{Conv}(a,\alpha) \to \mathrm{Conv}(b,\beta) \to \varphi_{160}$}
\end{ptree}
\noindent$\Pi_{26}$:
\begin{ptree}
\AXC{$z = w \to a_{N} + -b_{N} = z \to a_{N} + -b_{N} = w$\lbl{A6}}
\UIC{$\forall w\,(z = w \to a_{N} + -b_{N} = z \to a_{N} + -b_{N} = w)$}
\UIC{$\forall z\,\forall w\,(z = w \to a_{N} + -b_{N} = z \to a_{N} + -b_{N} = w)$}
\AXC{$\forall z\,\forall w\,(z = w \to a_{N} + -b_{N} = z \to a_{N} + -b_{N} = w) \to \varphi_{88}$\lbl{A2}}
\BIC{$\varphi_{88}$}
\AXC{$\varphi_{88} \to \varphi_{87}$\lbl{A2}}
\BIC{$\varphi_{87}$}
\end{ptree}
\noindent$\Pi_{27}$:
\begin{ptree}
\AXC{$z = w \to a_{N} + -b_{N} < z \to a_{N} + -b_{N} < w$\lbl{A6}}
\UIC{$\forall w\,(z = w \to a_{N} + -b_{N} < z \to a_{N} + -b_{N} < w)$}
\UIC{$\forall z\,\forall w\,(z = w \to a_{N} + -b_{N} < z \to a_{N} + -b_{N} < w)$}
\AXC{$\forall z\,\forall w\,(z = w \to a_{N} + -b_{N} < z \to a_{N} + -b_{N} < w) \to \varphi_{100}$\lbl{A2}}
\BIC{$\varphi_{100}$}
\AXC{$\varphi_{100} \to \varphi_{96}$\lbl{A2}}
\BIC{$\varphi_{96}$}
\end{ptree}
\noindent$\Pi_{28}$:
\begin{ptree}
\AXC{$\forall n \in \mathbb{N}\,(a_{n} \leq b_{n}) \to N \in \mathbb{N} \to a_{N} \leq b_{N}$\lbl{A2}}
\AXC{$\mathrm{Conv}(a,\alpha) \to \varphi_{37}$\lbl{A1}}
\AXC{$(\mathrm{Conv}(a,\alpha) \to \varphi_{37}) \to \varphi_{121}$\lbl{A1}}
\BIC{$\varphi_{121}$}
\BIC{$\mathrm{Conv}(a,\alpha) \to \varphi_{31}$}
\AXC{$\varphi_{224}$\lbl{A1}}
\BIC{$(\mathrm{Conv}(a,\alpha) \to \varphi_{143}) \to \mathrm{Conv}(a,\alpha) \to \varphi_{157}$}
\end{ptree}
\noindent$\Pi_{29}$:
\begin{ptree}
\AXC{$\stackrel{\Pi_{5}}{\vdots}$}\noLine
\UIC{$\mathrm{Conv}(a,\alpha) \to \varphi_{143}$}
\AXC{$\stackrel{\Pi_{28}}{\vdots}$}\noLine
\UIC{$(\mathrm{Conv}(a,\alpha) \to \varphi_{143}) \to \mathrm{Conv}(a,\alpha) \to \varphi_{157}$}
\BIC{$\mathrm{Conv}(a,\alpha) \to \varphi_{157}$}
\end{ptree}
\noindent$\Pi_{30}$:
\begin{ptree}
\AXC{$\stackrel{\Pi_{19}}{\vdots}$}\noLine
\UIC{$\mathrm{Conv}(a,\alpha) \to \varphi_{156}$}
\AXC{$\forall x \in \mathbb{R}\,(-x \in \mathbb{R})$\lbl{(R1)}}
\AXC{$\forall x \in \mathbb{R}\,(-x \in \mathbb{R}) \to b_{N} \in \mathbb{R} \to -b_{N} \in \mathbb{R}$\lbl{A2}}
\BIC{$b_{N} \in \mathbb{R} \to -b_{N} \in \mathbb{R}$}
\AXC{$\varphi_{218}$\lbl{A1}}
\BIC{$(\mathrm{Conv}(a,\alpha) \to \varphi_{156}) \to \mathrm{Conv}(a,\alpha) \to \varphi_{163}$}
\BIC{$\mathrm{Conv}(a,\alpha) \to \varphi_{163}$}
\end{ptree}
\noindent$\Pi_{31}$:
\begin{ptree}
\AXC{$\mathrm{Conv}(a,\alpha) \to \mathrm{Conv}(a,\alpha)$\lbl{A1}}
\AXC{$\forall a\,\forall \alpha\,(\mathrm{Conv}(a,\alpha) \to \mathrm{Seq}(a))$\lbl{[C1]}}
\AXC{$\varphi_{5}$\lbl{A2}}
\BIC{$\forall \alpha\,(\mathrm{Conv}(a,\alpha) \to \mathrm{Seq}(a))$}
\AXC{$\varphi_{2}$\lbl{A2}}
\BIC{$\mathrm{Conv}(a,\alpha) \to \mathrm{Seq}(a)$}
\AXC{$\varphi_{91}$\lbl{A1}}
\BIC{$(\mathrm{Conv}(a,\alpha) \to \mathrm{Conv}(a,\alpha)) \to \varphi_{15}$}
\BIC{$\varphi_{15}$}
\end{ptree}
\noindent$\Pi_{32}$:
\begin{ptree}
\AXC{$\stackrel{\Pi_{5}}{\vdots}$}\noLine
\UIC{$\mathrm{Conv}(a,\alpha) \to \varphi_{143}$}
\AXC{$\stackrel{\Pi_{31}}{\vdots}$}\noLine
\UIC{$\varphi_{15}$}
\AXC{$\forall a\,\forall n\,(\mathrm{Seq}(a) \to n \in \mathbb{N} \to a_{n} \in \mathbb{R})$\lbl{[S1]}}
\AXC{$\varphi_{49}$\lbl{A2}}
\BIC{$\forall n\,(\mathrm{Seq}(a) \to n \in \mathbb{N} \to a_{n} \in \mathbb{R})$}
\AXC{$\varphi_{41}$\lbl{A2}}
\BIC{$\mathrm{Seq}(a) \to N \in \mathbb{N} \to a_{N} \in \mathbb{R}$}
\AXC{$\varphi_{233}$\lbl{A1}}
\BIC{$\varphi_{15} \to (\mathrm{Conv}(a,\alpha) \to \varphi_{143}) \to \mathrm{Conv}(a,\alpha) \to \varphi_{158}$}
\BIC{$(\mathrm{Conv}(a,\alpha) \to \varphi_{143}) \to \mathrm{Conv}(a,\alpha) \to \varphi_{158}$}
\BIC{$\mathrm{Conv}(a,\alpha) \to \varphi_{158}$}
\end{ptree}
\noindent$\Pi_{33}$:
\begin{ptree}
\AXC{$\varphi_{16}$\lbl{[F19]}}
\AXC{$\varphi_{16} \to \varphi_{43}$\lbl{A2}}
\BIC{$\varphi_{43}$}
\AXC{$\varphi_{43} \to \varphi_{72}$\lbl{A2}}
\BIC{$\varphi_{72}$}
\AXC{$\varphi_{72} \to a_{N} \in \mathbb{R} \to \varphi_{95}$\lbl{A2}}
\BIC{$a_{N} \in \mathbb{R} \to \varphi_{95}$}
\AXC{$(a_{N} \in \mathbb{R} \to \varphi_{95}) \to \varphi_{255}$\lbl{A1}}
\BIC{$\varphi_{255}$}
\end{ptree}
\noindent$\Pi_{34}$:
\begin{ptree}
\AXC{$\stackrel{\Pi_{30}}{\vdots}$}\noLine
\UIC{$\mathrm{Conv}(a,\alpha) \to \varphi_{163}$}
\AXC{$\stackrel{\Pi_{19}}{\vdots}$}\noLine
\UIC{$\mathrm{Conv}(a,\alpha) \to \varphi_{156}$}
\AXC{$\stackrel{\Pi_{32}}{\vdots}$}\noLine
\UIC{$\mathrm{Conv}(a,\alpha) \to \varphi_{158}$}
\AXC{$\stackrel{\Pi_{33}}{\vdots}$}\noLine
\UIC{$\varphi_{255}$}
\BIC{$(\mathrm{Conv}(a,\alpha) \to \varphi_{156}) \to (\mathrm{Conv}(a,\alpha) \to \varphi_{163}) \to \varphi_{215}$}
\BIC{$(\mathrm{Conv}(a,\alpha) \to \varphi_{163}) \to \varphi_{215}$}
\BIC{$\varphi_{215}$}
\end{ptree}
\noindent$\Pi_{35}$:
\begin{ptree}
\AXC{$\forall x \in \mathbb{R}\,(x + -x = 0)$\lbl{(R5)}}
\AXC{$\forall x \in \mathbb{R}\,(x + -x = 0) \to b_{N} \in \mathbb{R} \to b_{N} + -b_{N} = 0$\lbl{A2}}
\BIC{$b_{N} \in \mathbb{R} \to b_{N} + -b_{N} = 0$}
\AXC{$(b_{N} \in \mathbb{R} \to b_{N} + -b_{N} = 0) \to \varphi_{208}$\lbl{A1}}
\BIC{$\varphi_{208}$}
\end{ptree}
\noindent$\Pi_{36}$:
\begin{ptree}
\AXC{$\stackrel{\Pi_{19}}{\vdots}$}\noLine
\UIC{$\mathrm{Conv}(a,\alpha) \to \varphi_{156}$}
\AXC{$\stackrel{\Pi_{35}}{\vdots}$}\noLine
\UIC{$\varphi_{208}$}
\BIC{$\mathrm{Conv}(a,\alpha) \to \mathrm{Conv}(b,\beta) \to \varphi_{152}$}
\AXC{$(\mathrm{Conv}(a,\alpha) \to \mathrm{Conv}(b,\beta) \to \varphi_{152}) \to \varphi_{242}$\lbl{A1}}
\BIC{$\varphi_{242}$}
\end{ptree}
\noindent$\Pi_{37}$:
\begin{ptree}
\AXC{$x = x$\lbl{A5}}
\AXC{$x = y \to x = x \to y = x$\lbl{A6}}
\AXC{$(x = y \to x = x \to y = x) \to x = x \to x = y \to y = x$\lbl{A1}}
\BIC{$x = x \to x = y \to y = x$}
\BIC{$x = y \to y = x$}
\UIC{$\forall y\,(x = y \to y = x)$}
\UIC{$\forall x\,\forall y\,(x = y \to y = x)$}
\end{ptree}
\noindent$\Pi_{38}$:
\begin{ptree}
\AXC{$\forall x\,\forall y\,(x - y = x + -y)$\lbl{Def$-$}}
\AXC{$\forall x\,\forall y\,(x - y = x + -y) \to \forall y\,(a_{N} - y = a_{N} + -y)$\lbl{A2}}
\BIC{$\forall y\,(a_{N} - y = a_{N} + -y)$}
\end{ptree}
\noindent$\Pi_{39}$:
\begin{ptree}
\AXC{$\stackrel{\Pi_{38}}{\vdots}$}\noLine
\UIC{$\forall y\,(a_{N} - y = a_{N} + -y)$}
\AXC{$\forall y\,(a_{N} - y = a_{N} + -y) \to a_{N} - b_{N} = a_{N} + -b_{N}$\lbl{A2}}
\BIC{$a_{N} - b_{N} = a_{N} + -b_{N}$}
\AXC{$a_{N} - b_{N} = a_{N} + -b_{N} \to \varphi_{115}$\lbl{A1}}
\BIC{$\varphi_{115}$}
\end{ptree}
\noindent$\Pi_{40}$:
\begin{ptree}
\AXC{$\stackrel{\Pi_{37}}{\vdots}$}\noLine
\UIC{$\forall x\,\forall y\,(x = y \to y = x)$}
\AXC{$\varphi_{17}$\lbl{A2}}
\BIC{$\forall y\,(a_{N} - b_{N} = y \to y = a_{N} - b_{N})$}
\AXC{$\varphi_{116}$\lbl{A2}}
\BIC{$a_{N} - b_{N} = a_{N} + -b_{N} \to a_{N} + -b_{N} = a_{N} - b_{N}$}
\AXC{$\stackrel{\Pi_{39}}{\vdots}$}\noLine
\UIC{$\varphi_{115}$}
\BIC{$a_{N} + -b_{N} = a_{N} - b_{N}$}
\AXC{$\varphi_{243}$\lbl{A1}}
\BIC{$(\mathrm{Conv}(a,\alpha) \to \mathrm{Conv}(b,\beta) \to \varphi_{153}) \to \varphi_{205}$}
\end{ptree}
\noindent$\Pi_{41}$:
\begin{ptree}
\AXC{$\stackrel{\Pi_{26}}{\vdots}$}\noLine
\UIC{$\varphi_{87}$}
\AXC{$\stackrel{\Pi_{27}}{\vdots}$}\noLine
\UIC{$\varphi_{96}$}
\AXC{$\stackrel{\Pi_{29}}{\vdots}$}\noLine
\UIC{$\mathrm{Conv}(a,\alpha) \to \varphi_{157}$}
\AXC{$\stackrel{\Pi_{34}}{\vdots}$}\noLine
\UIC{$\varphi_{215}$}
\BIC{$\mathrm{Conv}(a,\alpha) \to \mathrm{Conv}(b,\beta) \to \varphi_{174}$}
\AXC{$\stackrel{\Pi_{36}}{\vdots}$}\noLine
\UIC{$\varphi_{242}$}
\BIC{$\varphi_{96} \to \varphi_{87} \to \mathrm{Conv}(a,\alpha) \to \mathrm{Conv}(b,\beta) \to \varphi_{153}$}
\BIC{$\varphi_{87} \to \mathrm{Conv}(a,\alpha) \to \mathrm{Conv}(b,\beta) \to \varphi_{153}$}
\BIC{$\mathrm{Conv}(a,\alpha) \to \mathrm{Conv}(b,\beta) \to \varphi_{153}$}
\AXC{$\stackrel{\Pi_{40}}{\vdots}$}\noLine
\UIC{$(\mathrm{Conv}(a,\alpha) \to \mathrm{Conv}(b,\beta) \to \varphi_{153}) \to \varphi_{205}$}
\BIC{$\varphi_{205}$}
\end{ptree}
\noindent$\Pi_{42}$:
\begin{ptree}
\AXC{$z = w \to z < 0 \to w < 0$\lbl{A6}}
\UIC{$\forall w\,(z = w \to z < 0 \to w < 0)$}
\UIC{$\forall z\,\forall w\,(z = w \to z < 0 \to w < 0)$}
\AXC{$\varphi_{71}$\lbl{A2}}
\BIC{$\forall w\,(a_{N} + -b_{N} = w \to a_{N} + -b_{N} < 0 \to w < 0)$}
\AXC{$\varphi_{135}$\lbl{A2}}
\BIC{$a_{N} + -b_{N} = a_{N} - b_{N} \to a_{N} + -b_{N} < 0 \to a_{N} - b_{N} < 0$}
\AXC{$\stackrel{\Pi_{41}}{\vdots}$}\noLine
\UIC{$\varphi_{205}$}
\BIC{$\varphi_{196}$}
\end{ptree}
\noindent$\Pi_{43}$:
\begin{ptree}
\AXC{$\stackrel{\Pi_{32}}{\vdots}$}\noLine
\UIC{$\mathrm{Conv}(a,\alpha) \to \varphi_{158}$}
\AXC{$\forall x\,\forall y\,(x \in \mathbb{R} \to y \in \mathbb{R} \to x - y \in \mathbb{R})$\lbl{[F1]}}
\AXC{$\varphi_{84}$\lbl{A2}}
\BIC{$\forall y\,(a_{N} \in \mathbb{R} \to y \in \mathbb{R} \to a_{N} - y \in \mathbb{R})$}
\AXC{$\varphi_{108}$\lbl{A2}}
\BIC{$a_{N} \in \mathbb{R} \to b_{N} \in \mathbb{R} \to a_{N} - b_{N} \in \mathbb{R}$}
\AXC{$\varphi_{245}$\lbl{A1}}
\BIC{$(\mathrm{Conv}(a,\alpha) \to \varphi_{158}) \to \varphi_{209}$}
\BIC{$\varphi_{209}$}
\end{ptree}
\noindent$\Pi_{44}$:
\begin{ptree}
\AXC{$\varphi_{51}$\lbl{[F21]}}
\AXC{$\varphi_{51} \to \varphi_{114}$\lbl{A2}}
\BIC{$\varphi_{114}$}
\AXC{$\varphi_{114} \to \varphi_{112}$\lbl{A2}}
\BIC{$\varphi_{112}$}
\AXC{$\varphi_{112} \to \varphi_{128}$\lbl{A2}}
\BIC{$\varphi_{128}$}
\AXC{$\varphi_{128} \to a_{N} - b_{N} \in \mathbb{R} \to \varphi_{150}$\lbl{A2}}
\BIC{$a_{N} - b_{N} \in \mathbb{R} \to \varphi_{150}$}
\AXC{$\varphi_{258}$\lbl{A1}}
\BIC{$(\mathrm{Conv}(a,\alpha) \to \mathrm{Conv}(b,\beta) \to \varphi_{154}) \to 0 \in \mathbb{R} \to \varphi_{256}$}
\end{ptree}
\noindent$\Pi_{45}$:
\begin{ptree}
\AXC{$0 \in \mathbb{R}$\lbl{(R1)}}
\AXC{$\stackrel{\Pi_{19}}{\vdots}$}\noLine
\UIC{$\mathrm{Conv}(a,\alpha) \to \varphi_{156}$}
\AXC{$\stackrel{\Pi_{43}}{\vdots}$}\noLine
\UIC{$\varphi_{209}$}
\BIC{$\mathrm{Conv}(a,\alpha) \to \mathrm{Conv}(b,\beta) \to \varphi_{154}$}
\AXC{$\stackrel{\Pi_{44}}{\vdots}$}\noLine
\UIC{$(\mathrm{Conv}(a,\alpha) \to \mathrm{Conv}(b,\beta) \to \varphi_{154}) \to 0 \in \mathbb{R} \to \varphi_{256}$}
\BIC{$0 \in \mathbb{R} \to \varphi_{256}$}
\BIC{$\varphi_{256}$}
\end{ptree}
\noindent$\Pi_{46}$:
\begin{ptree}
\AXC{$\stackrel{\Pi_{15}}{\vdots}$}\noLine
\UIC{$\mathrm{Conv}(a,\alpha) \to \mathrm{Conv}(b,\beta) \to \varphi_{151}$}
\AXC{$\stackrel{\Pi_{21}}{\vdots}$}\noLine
\UIC{$\mathrm{Conv}(a,\alpha) \to \mathrm{Conv}(b,\beta) \to \varphi_{146}$}
\AXC{$\stackrel{\Pi_{45}}{\vdots}$}\noLine
\UIC{$\varphi_{256}$}
\BIC{$(\mathrm{Conv}(a,\alpha) \to \mathrm{Conv}(b,\beta) \to \varphi_{151}) \to \varphi_{244}$}
\BIC{$\varphi_{244}$}
\end{ptree}
\noindent$\Pi_{47}$:
\begin{ptree}
\AXC{$z = w \to z = 0 \to w = 0$\lbl{A6}}
\UIC{$\forall w\,(z = w \to z = 0 \to w = 0)$}
\UIC{$\forall z\,\forall w\,(z = w \to z = 0 \to w = 0)$}
\AXC{$\varphi_{55}$\lbl{A2}}
\BIC{$\forall w\,(a_{N} + -b_{N} = w \to a_{N} + -b_{N} = 0 \to w = 0)$}
\AXC{$\varphi_{124}$\lbl{A2}}
\BIC{$a_{N} + -b_{N} = a_{N} - b_{N} \to a_{N} + -b_{N} = 0 \to a_{N} - b_{N} = 0$}
\AXC{$\stackrel{\Pi_{42}}{\vdots}$}\noLine
\UIC{$\varphi_{196}$}
\BIC{$\mathrm{Conv}(a,\alpha) \to \mathrm{Conv}(b,\beta) \to \varphi_{144}$}
\AXC{$\stackrel{\Pi_{46}}{\vdots}$}\noLine
\UIC{$\varphi_{244}$}
\BIC{$(\mathrm{Conv}(a,\alpha) \to \mathrm{Conv}(b,\beta) \to \varphi_{160}) \to \varphi_{186}$}
\end{ptree}
\noindent$\Pi_{48}$:
\begin{ptree}
\AXC{$y = z \to x = y \to x = z$\lbl{A6}}
\AXC{$(y = z \to x = y \to x = z) \to x = y \to y = z \to x = z$\lbl{A1}}
\BIC{$x = y \to y = z \to x = z$}
\UIC{$\forall z\,(x = y \to y = z \to x = z)$}
\UIC{$\forall y\,\forall z\,(x = y \to y = z \to x = z)$}
\UIC{$\forall x\,\forall y\,\forall z\,(x = y \to y = z \to x = z)$}
\AXC{$\forall x\,\forall y\,\forall z\,(x = y \to y = z \to x = z) \to \varphi_{12}$\lbl{A2}}
\BIC{$\varphi_{12}$}
\end{ptree}
\noindent$\Pi_{49}$:
\begin{ptree}
\AXC{$\stackrel{\Pi_{48}}{\vdots}$}\noLine
\UIC{$\varphi_{12}$}
\AXC{$\varphi_{12} \to \varphi_{76}$\lbl{A2}}
\BIC{$\varphi_{76}$}
\AXC{$\varphi_{76} \to 0 + (\alpha - \beta) \cdot 2^{-1} = (\alpha - \beta) \cdot 2^{-1} + 0 \to \varphi_{47}$\lbl{A2}}
\BIC{$0 + (\alpha - \beta) \cdot 2^{-1} = (\alpha - \beta) \cdot 2^{-1} + 0 \to \varphi_{47}$}
\end{ptree}
\noindent$\Pi_{50}$:
\begin{ptree}
\AXC{$\forall x \in \mathbb{R}\,(x + 0 = x)$\lbl{(R4)}}
\AXC{$\forall x \in \mathbb{R}\,(x + 0 = x) \to \varphi_{25}$\lbl{A2}}
\BIC{$\varphi_{25}$}
\AXC{$\varphi_{25} \to (\mathrm{Conv}(a,\alpha) \to \mathrm{Conv}(b,\beta) \to \varphi_{151}) \to \varphi_{177}$\lbl{A1}}
\BIC{$(\mathrm{Conv}(a,\alpha) \to \mathrm{Conv}(b,\beta) \to \varphi_{151}) \to \varphi_{177}$}
\end{ptree}
\noindent$\Pi_{51}$:
\begin{ptree}
\AXC{$\forall x\,\forall y\,(x \in \mathbb{R} \to y \in \mathbb{R} \to x + y = y + x)$\lbl{(R3)}}
\AXC{$\varphi_{56}$\lbl{A2}}
\BIC{$\forall y\,(0 \in \mathbb{R} \to y \in \mathbb{R} \to 0 + y = y + 0)$}
\AXC{$\forall y\,(0 \in \mathbb{R} \to y \in \mathbb{R} \to 0 + y = y + 0) \to 0 \in \mathbb{R} \to \varphi_{32}$\lbl{A2}}
\BIC{$0 \in \mathbb{R} \to \varphi_{32}$}
\end{ptree}
\noindent$\Pi_{52}$:
\begin{ptree}
\AXC{$\stackrel{\Pi_{15}}{\vdots}$}\noLine
\UIC{$\mathrm{Conv}(a,\alpha) \to \mathrm{Conv}(b,\beta) \to \varphi_{151}$}
\AXC{$0 \in \mathbb{R}$\lbl{(R1)}}
\AXC{$\stackrel{\Pi_{51}}{\vdots}$}\noLine
\UIC{$0 \in \mathbb{R} \to \varphi_{32}$}
\AXC{$\varphi_{236}$\lbl{A1}}
\BIC{$0 \in \mathbb{R} \to (\mathrm{Conv}(a,\alpha) \to \mathrm{Conv}(b,\beta) \to \varphi_{151}) \to \varphi_{180}$}
\BIC{$(\mathrm{Conv}(a,\alpha) \to \mathrm{Conv}(b,\beta) \to \varphi_{151}) \to \varphi_{180}$}
\BIC{$\varphi_{180}$}
\end{ptree}
\noindent$\Pi_{53}$:
\begin{ptree}
\AXC{$\stackrel{\Pi_{15}}{\vdots}$}\noLine
\UIC{$\mathrm{Conv}(a,\alpha) \to \mathrm{Conv}(b,\beta) \to \varphi_{151}$}
\AXC{$\stackrel{\Pi_{50}}{\vdots}$}\noLine
\UIC{$(\mathrm{Conv}(a,\alpha) \to \mathrm{Conv}(b,\beta) \to \varphi_{151}) \to \varphi_{177}$}
\BIC{$\varphi_{177}$}
\AXC{$\stackrel{\Pi_{52}}{\vdots}$}\noLine
\UIC{$\varphi_{180}$}
\AXC{$\varphi_{180} \to \varphi_{239}$\lbl{A1}}
\BIC{$\varphi_{239}$}
\BIC{$(0 + (\alpha - \beta) \cdot 2^{-1} = (\alpha - \beta) \cdot 2^{-1} + 0 \to \varphi_{47}) \to \varphi_{178}$}
\end{ptree}
\noindent$\Pi_{54}$:
\begin{ptree}
\AXC{$\stackrel{\Pi_{49}}{\vdots}$}\noLine
\UIC{$0 + (\alpha - \beta) \cdot 2^{-1} = (\alpha - \beta) \cdot 2^{-1} + 0 \to \varphi_{47}$}
\AXC{$\stackrel{\Pi_{53}}{\vdots}$}\noLine
\UIC{$(0 + (\alpha - \beta) \cdot 2^{-1} = (\alpha - \beta) \cdot 2^{-1} + 0 \to \varphi_{47}) \to \varphi_{178}$}
\BIC{$\varphi_{178}$}
\AXC{$\varphi_{178} \to \varphi_{186} \to \varphi_{237}$\lbl{A1}}
\BIC{$\varphi_{186} \to \varphi_{237}$}
\end{ptree}
\noindent$\Pi_{55}$:
\begin{ptree}
\AXC{$\stackrel{\Pi_{25}}{\vdots}$}\noLine
\UIC{$\mathrm{Conv}(a,\alpha) \to \mathrm{Conv}(b,\beta) \to \varphi_{160}$}
\AXC{$\stackrel{\Pi_{47}}{\vdots}$}\noLine
\UIC{$(\mathrm{Conv}(a,\alpha) \to \mathrm{Conv}(b,\beta) \to \varphi_{160}) \to \varphi_{186}$}
\BIC{$\varphi_{186}$}
\AXC{$\stackrel{\Pi_{54}}{\vdots}$}\noLine
\UIC{$\varphi_{186} \to \varphi_{237}$}
\BIC{$\varphi_{237}$}
\end{ptree}
\noindent$\Pi_{56}$:
\begin{ptree}
\AXC{$\varphi_{63}$\lbl{A6}}
\UIC{$\forall w\,(\varphi_{63})$}
\UIC{$\forall z\,\forall w\,(\varphi_{63})$}
\AXC{$\forall z\,\forall w\,(\varphi_{63}) \to \varphi_{113}$\lbl{A2}}
\BIC{$\varphi_{113}$}
\AXC{$\varphi_{113} \to 0 + (\alpha - \beta) \cdot 2^{-1} = (\alpha - \beta) \cdot 2^{-1} \to \varphi_{107}$\lbl{A2}}
\BIC{$0 + (\alpha - \beta) \cdot 2^{-1} = (\alpha - \beta) \cdot 2^{-1} \to \varphi_{107}$}
\AXC{$\stackrel{\Pi_{55}}{\vdots}$}\noLine
\UIC{$\varphi_{237}$}
\BIC{$(0 + (\alpha - \beta) \cdot 2^{-1} = (\alpha - \beta) \cdot 2^{-1} \to \varphi_{103}) \to \varphi_{185}$}
\end{ptree}
\noindent$\Pi_{57}$:
\begin{ptree}
\AXC{$\stackrel{\Pi_{4}}{\vdots}$}\noLine
\UIC{$0 + (\alpha - \beta) \cdot 2^{-1} = (\alpha - \beta) \cdot 2^{-1} \to \varphi_{103}$}
\AXC{$\stackrel{\Pi_{56}}{\vdots}$}\noLine
\UIC{$(0 + (\alpha - \beta) \cdot 2^{-1} = (\alpha - \beta) \cdot 2^{-1} \to \varphi_{103}) \to \varphi_{185}$}
\BIC{$\varphi_{185}$}
\end{ptree}
\noindent$\Pi_{58}$:
\begin{ptree}
\AXC{$\stackrel{\Pi_{37}}{\vdots}$}\noLine
\UIC{$\forall x\,\forall y\,(x = y \to y = x)$}
\AXC{$\varphi_{8}$\lbl{A2}}
\BIC{$\forall y\,(a_{N} - \beta = y \to y = a_{N} - \beta)$}
\AXC{$\forall y\,(a_{N} - \beta = y \to y = a_{N} - \beta) \to \varphi_{85}$\lbl{A2}}
\BIC{$\varphi_{85}$}
\end{ptree}
\noindent$\Pi_{59}$:
\begin{ptree}
\AXC{$\varphi_{20}$\lbl{[F16]}}
\AXC{$\varphi_{20} \to \varphi_{52}$\lbl{A2}}
\BIC{$\varphi_{52}$}
\AXC{$\varphi_{52} \to \varphi_{53}$\lbl{A2}}
\BIC{$\varphi_{53}$}
\AXC{$\varphi_{53} \to \varphi_{79}$\lbl{A2}}
\BIC{$\varphi_{79}$}
\AXC{$\varphi_{79} \to (\mathrm{Conv}(a,\alpha) \to \varphi_{158}) \to \varphi_{23} \to \varphi_{217}$\lbl{A1}}
\BIC{$(\mathrm{Conv}(a,\alpha) \to \varphi_{158}) \to \varphi_{23} \to \varphi_{217}$}
\end{ptree}
\noindent$\Pi_{60}$:
\begin{ptree}
\AXC{$\stackrel{\Pi_{19}}{\vdots}$}\noLine
\UIC{$\mathrm{Conv}(a,\alpha) \to \varphi_{156}$}
\AXC{$\stackrel{\Pi_{10}}{\vdots}$}\noLine
\UIC{$\varphi_{23}$}
\AXC{$\stackrel{\Pi_{32}}{\vdots}$}\noLine
\UIC{$\mathrm{Conv}(a,\alpha) \to \varphi_{158}$}
\AXC{$\stackrel{\Pi_{59}}{\vdots}$}\noLine
\UIC{$(\mathrm{Conv}(a,\alpha) \to \varphi_{158}) \to \varphi_{23} \to \varphi_{217}$}
\BIC{$\varphi_{23} \to \varphi_{217}$}
\BIC{$\varphi_{217}$}
\BIC{$\mathrm{Conv}(a,\alpha) \to \mathrm{Conv}(b,\beta) \to \varphi_{175}$}
\AXC{$\varphi_{238}$\lbl{A1}}
\BIC{$\varphi_{85} \to \mathrm{Conv}(a,\alpha) \to \mathrm{Conv}(b,\beta) \to \varphi_{176}$}
\end{ptree}
\noindent$\Pi_{61}$:
\begin{ptree}
\AXC{$\stackrel{\Pi_{3}}{\vdots}$}\noLine
\UIC{$a_{N} - b_{N} + (b_{N} - \beta) = a_{N} - \beta \to \varphi_{70}$}
\AXC{$\stackrel{\Pi_{57}}{\vdots}$}\noLine
\UIC{$\varphi_{185}$}
\AXC{$\stackrel{\Pi_{58}}{\vdots}$}\noLine
\UIC{$\varphi_{85}$}
\AXC{$\stackrel{\Pi_{60}}{\vdots}$}\noLine
\UIC{$\varphi_{85} \to \mathrm{Conv}(a,\alpha) \to \mathrm{Conv}(b,\beta) \to \varphi_{176}$}
\BIC{$\mathrm{Conv}(a,\alpha) \to \mathrm{Conv}(b,\beta) \to \varphi_{176}$}
\AXC{$\varphi_{252}$\lbl{A1}}
\BIC{$\varphi_{185} \to (a_{N} - b_{N} + (b_{N} - \beta) = a_{N} - \beta \to \varphi_{70}) \to \varphi_{200}$}
\BIC{$(a_{N} - b_{N} + (b_{N} - \beta) = a_{N} - \beta \to \varphi_{70}) \to \varphi_{200}$}
\BIC{$\varphi_{200}$}
\end{ptree}
\noindent$\Pi_{62}$:
\begin{ptree}
\AXC{$z = w \to -(a_{N} - \alpha) < z \to -(a_{N} - \alpha) < w$\lbl{A6}}
\UIC{$\forall w\,(z = w \to -(a_{N} - \alpha) < z \to -(a_{N} - \alpha) < w)$}
\UIC{$\forall z\,\forall w\,(z = w \to -(a_{N} - \alpha) < z \to -(a_{N} - \alpha) < w)$}
\AXC{$\varphi_{132}$\lbl{A2}}
\BIC{$\varphi_{89}$}
\AXC{$\varphi_{89} \to -(-((\alpha - \beta) \cdot 2^{-1})) = (\alpha - \beta) \cdot 2^{-1} \to \varphi_{65}$\lbl{A2}}
\BIC{$-(-((\alpha - \beta) \cdot 2^{-1})) = (\alpha - \beta) \cdot 2^{-1} \to \varphi_{65}$}
\end{ptree}
\noindent$\Pi_{63}$:
\begin{ptree}
\AXC{$\forall x\,\forall y\,(x \in \mathbb{R} \to y \in \mathbb{R} \to x - y \in \mathbb{R})$\lbl{[F1]}}
\AXC{$\varphi_{84}$\lbl{A2}}
\BIC{$\forall y\,(a_{N} \in \mathbb{R} \to y \in \mathbb{R} \to a_{N} - y \in \mathbb{R})$}
\AXC{$\varphi_{98}$\lbl{A2}}
\BIC{$a_{N} \in \mathbb{R} \to \alpha \in \mathbb{R} \to a_{N} - \alpha \in \mathbb{R}$}
\AXC{$(a_{N} \in \mathbb{R} \to \alpha \in \mathbb{R} \to a_{N} - \alpha \in \mathbb{R}) \to \varphi_{221}$\lbl{A1}}
\BIC{$\varphi_{221}$}
\end{ptree}
\noindent$\Pi_{64}$:
\begin{ptree}
\AXC{$\stackrel{\Pi_{11}}{\vdots}$}\noLine
\UIC{$\varphi_{24}$}
\AXC{$\stackrel{\Pi_{32}}{\vdots}$}\noLine
\UIC{$\mathrm{Conv}(a,\alpha) \to \varphi_{158}$}
\AXC{$\stackrel{\Pi_{63}}{\vdots}$}\noLine
\UIC{$\varphi_{221}$}
\BIC{$\varphi_{24} \to \mathrm{Conv}(a,\alpha) \to \mathrm{Conv}(b,\beta) \to \varphi_{147}$}
\BIC{$\mathrm{Conv}(a,\alpha) \to \mathrm{Conv}(b,\beta) \to \varphi_{147}$}
\end{ptree}
\noindent$\Pi_{65}$:
\begin{ptree}
\AXC{$\forall x\,\forall y\,(x \in \mathbb{R} \to y \in \mathbb{R} \to \lvert x \rvert < y \to -y < x \land x < y)$\lbl{[A11]}}
\AXC{$\varphi_{139}$\lbl{A2}}
\BIC{$\varphi_{92}$}
\AXC{$\varphi_{92} \to a_{N} - \alpha \in \mathbb{R} \to \varphi_{120}$\lbl{A2}}
\BIC{$a_{N} - \alpha \in \mathbb{R} \to \varphi_{120}$}
\end{ptree}
\noindent$\Pi_{66}$:
\begin{ptree}
\AXC{$\stackrel{\Pi_{64}}{\vdots}$}\noLine
\UIC{$\mathrm{Conv}(a,\alpha) \to \mathrm{Conv}(b,\beta) \to \varphi_{147}$}
\AXC{$\stackrel{\Pi_{65}}{\vdots}$}\noLine
\UIC{$a_{N} - \alpha \in \mathbb{R} \to \varphi_{120}$}
\AXC{$(a_{N} - \alpha \in \mathbb{R} \to \varphi_{120}) \to \varphi_{250}$\lbl{A1}}
\BIC{$\varphi_{250}$}
\BIC{$(\mathrm{Conv}(a,\alpha) \to \mathrm{Conv}(b,\beta) \to \varphi_{151}) \to \varphi_{226}$}
\end{ptree}
\noindent$\Pi_{67}$:
\begin{ptree}
\AXC{$\stackrel{\Pi_{15}}{\vdots}$}\noLine
\UIC{$\mathrm{Conv}(a,\alpha) \to \mathrm{Conv}(b,\beta) \to \varphi_{151}$}
\AXC{$\stackrel{\Pi_{66}}{\vdots}$}\noLine
\UIC{$(\mathrm{Conv}(a,\alpha) \to \mathrm{Conv}(b,\beta) \to \varphi_{151}) \to \varphi_{226}$}
\BIC{$\varphi_{226}$}
\end{ptree}
\noindent$\Pi_{68}$:
\begin{ptree}
\AXC{$\stackrel{\Pi_{15}}{\vdots}$}\noLine
\UIC{$\mathrm{Conv}(a,\alpha) \to \mathrm{Conv}(b,\beta) \to \varphi_{151}$}
\AXC{$\forall x \in \mathbb{R}\,(-x \in \mathbb{R})$\lbl{(R1)}}
\AXC{$\varphi_{68}$\lbl{A2}}
\BIC{$(\alpha - \beta) \cdot 2^{-1} \in \mathbb{R} \to -((\alpha - \beta) \cdot 2^{-1}) \in \mathbb{R}$}
\AXC{$\varphi_{229}$\lbl{A1}}
\BIC{$\varphi_{213}$}
\BIC{$\mathrm{Conv}(a,\alpha) \to \mathrm{Conv}(b,\beta) \to \varphi_{159}$}
\end{ptree}
\noindent$\Pi_{69}$:
\begin{ptree}
\AXC{$\forall x\,\forall y\,(x \in \mathbb{R} \to y \in \mathbb{R} \to x < y \to -y < -x)$\lbl{[F22]}}
\AXC{$\forall x\,\forall y\,(x \in \mathbb{R} \to y \in \mathbb{R} \to x < y \to -y < -x) \to \varphi_{101}$\lbl{A2}}
\BIC{$\varphi_{101}$}
\end{ptree}
\noindent$\Pi_{70}$:
\begin{ptree}
\AXC{$\stackrel{\Pi_{68}}{\vdots}$}\noLine
\UIC{$\mathrm{Conv}(a,\alpha) \to \mathrm{Conv}(b,\beta) \to \varphi_{159}$}
\AXC{$\stackrel{\Pi_{69}}{\vdots}$}\noLine
\UIC{$\varphi_{101}$}
\AXC{$\varphi_{101} \to -((\alpha - \beta) \cdot 2^{-1}) \in \mathbb{R} \to a_{N} - \alpha \in \mathbb{R} \to \varphi_{66}$\lbl{A2}}
\BIC{$-((\alpha - \beta) \cdot 2^{-1}) \in \mathbb{R} \to a_{N} - \alpha \in \mathbb{R} \to \varphi_{66}$}
\AXC{$\varphi_{254}$\lbl{A1}}
\BIC{$(\mathrm{Conv}(a,\alpha) \to \mathrm{Conv}(b,\beta) \to \varphi_{159}) \to \varphi_{241}$}
\BIC{$\varphi_{241}$}
\end{ptree}
\noindent$\Pi_{71}$:
\begin{ptree}
\AXC{$\stackrel{\Pi_{8}}{\vdots}$}\noLine
\UIC{$\mathrm{Conv}(a,\alpha) \to \varphi_{190}$}
\AXC{$\varphi_{231}$\lbl{A1}}
\BIC{$\mathrm{Conv}(a,\alpha) \to \mathrm{Conv}(b,\beta) \to \varphi_{167}$}
\AXC{$\stackrel{\Pi_{67}}{\vdots}$}\noLine
\UIC{$\varphi_{226}$}
\BIC{$\mathrm{Conv}(a,\alpha) \to \varphi_{189}$}
\AXC{$\varphi_{227}$\lbl{A1}}
\BIC{$\mathrm{Conv}(a,\alpha) \to \mathrm{Conv}(b,\beta) \to \varphi_{168}$}
\AXC{$\stackrel{\Pi_{64}}{\vdots}$}\noLine
\UIC{$\mathrm{Conv}(a,\alpha) \to \mathrm{Conv}(b,\beta) \to \varphi_{147}$}
\AXC{$\stackrel{\Pi_{70}}{\vdots}$}\noLine
\UIC{$\varphi_{241}$}
\BIC{$(\mathrm{Conv}(a,\alpha) \to \mathrm{Conv}(b,\beta) \to \varphi_{168}) \to \varphi_{183}$}
\BIC{$\varphi_{183}$}
\end{ptree}
\noindent$\Pi_{72}$:
\begin{ptree}
\AXC{$\forall x \in \mathbb{R}\,(-(-x) = x)$\lbl{[F5]}}
\AXC{$\forall x \in \mathbb{R}\,(-(-x) = x) \to \varphi_{30}$\lbl{A2}}
\BIC{$\varphi_{30}$}
\AXC{$\varphi_{30} \to (\mathrm{Conv}(a,\alpha) \to \mathrm{Conv}(b,\beta) \to \varphi_{151}) \to \varphi_{181}$\lbl{A1}}
\BIC{$(\mathrm{Conv}(a,\alpha) \to \mathrm{Conv}(b,\beta) \to \varphi_{151}) \to \varphi_{181}$}
\end{ptree}
\noindent$\Pi_{73}$:
\begin{ptree}
\AXC{$\stackrel{\Pi_{71}}{\vdots}$}\noLine
\UIC{$\varphi_{183}$}
\AXC{$\stackrel{\Pi_{15}}{\vdots}$}\noLine
\UIC{$\mathrm{Conv}(a,\alpha) \to \mathrm{Conv}(b,\beta) \to \varphi_{151}$}
\AXC{$\stackrel{\Pi_{72}}{\vdots}$}\noLine
\UIC{$(\mathrm{Conv}(a,\alpha) \to \mathrm{Conv}(b,\beta) \to \varphi_{151}) \to \varphi_{181}$}
\BIC{$\varphi_{181}$}
\AXC{$\varphi_{181} \to \varphi_{183} \to \varphi_{201}$\lbl{A1}}
\BIC{$\varphi_{183} \to \varphi_{201}$}
\BIC{$\varphi_{201}$}
\end{ptree}
\noindent$\Pi_{74}$:
\begin{ptree}
\AXC{$\forall x\,\forall y\,(x \in \mathbb{R} \to y \in \mathbb{R} \to -(y - x) = x - y)$\lbl{[F14]}}
\AXC{$\varphi_{73}$\lbl{A2}}
\BIC{$\forall y\,(\alpha \in \mathbb{R} \to y \in \mathbb{R} \to -(y - \alpha) = \alpha - y)$}
\AXC{$\varphi_{102}$\lbl{A2}}
\BIC{$\alpha \in \mathbb{R} \to a_{N} \in \mathbb{R} \to -(a_{N} - \alpha) = \alpha - a_{N}$}
\end{ptree}
\noindent$\Pi_{75}$:
\begin{ptree}
\AXC{$\stackrel{\Pi_{11}}{\vdots}$}\noLine
\UIC{$\varphi_{24}$}
\AXC{$\stackrel{\Pi_{74}}{\vdots}$}\noLine
\UIC{$\alpha \in \mathbb{R} \to a_{N} \in \mathbb{R} \to -(a_{N} - \alpha) = \alpha - a_{N}$}
\AXC{$(\alpha \in \mathbb{R} \to a_{N} \in \mathbb{R} \to -(a_{N} - \alpha) = \alpha - a_{N}) \to \varphi_{24} \to \varphi_{211}$\lbl{A1}}
\BIC{$\varphi_{24} \to \varphi_{211}$}
\BIC{$\varphi_{211}$}
\end{ptree}
\noindent$\Pi_{76}$:
\begin{ptree}
\AXC{$\stackrel{\Pi_{62}}{\vdots}$}\noLine
\UIC{$-(-((\alpha - \beta) \cdot 2^{-1})) = (\alpha - \beta) \cdot 2^{-1} \to \varphi_{65}$}
\AXC{$\stackrel{\Pi_{73}}{\vdots}$}\noLine
\UIC{$\varphi_{201}$}
\BIC{$\mathrm{Conv}(a,\alpha) \to \mathrm{Conv}(b,\beta) \to \varphi_{169}$}
\AXC{$\stackrel{\Pi_{32}}{\vdots}$}\noLine
\UIC{$\mathrm{Conv}(a,\alpha) \to \varphi_{158}$}
\AXC{$\stackrel{\Pi_{75}}{\vdots}$}\noLine
\UIC{$\varphi_{211}$}
\BIC{$\mathrm{Conv}(a,\alpha) \to \mathrm{Conv}(b,\beta) \to \varphi_{162}$}
\AXC{$\varphi_{247}$\lbl{A1}}
\BIC{$(\mathrm{Conv}(a,\alpha) \to \mathrm{Conv}(b,\beta) \to \varphi_{169}) \to \varphi_{198}$}
\BIC{$\varphi_{198}$}
\end{ptree}
\noindent$\Pi_{77}$:
\begin{ptree}
\AXC{$z = w \to z < (\alpha - \beta) \cdot 2^{-1} \to w < (\alpha - \beta) \cdot 2^{-1}$\lbl{A6}}
\UIC{$\varphi_{9}$}
\UIC{$\forall z\,(\varphi_{9})$}
\AXC{$\forall z\,(\varphi_{9}) \to \varphi_{62}$\lbl{A2}}
\BIC{$\varphi_{62}$}
\AXC{$\varphi_{62} \to -(a_{N} - \alpha) = \alpha - a_{N} \to \varphi_{38}$\lbl{A2}}
\BIC{$-(a_{N} - \alpha) = \alpha - a_{N} \to \varphi_{38}$}
\AXC{$\stackrel{\Pi_{76}}{\vdots}$}\noLine
\UIC{$\varphi_{198}$}
\BIC{$\mathrm{Conv}(a,\alpha) \to \mathrm{Conv}(b,\beta) \to \varphi_{165}$}
\end{ptree}
\noindent$\Pi_{78}$:
\begin{ptree}
\AXC{$\forall x\,\forall y\,(x \in \mathbb{R} \to y \in \mathbb{R} \to x - y \in \mathbb{R})$\lbl{[F1]}}
\AXC{$\varphi_{84}$\lbl{A2}}
\BIC{$\forall y\,(a_{N} \in \mathbb{R} \to y \in \mathbb{R} \to a_{N} - y \in \mathbb{R})$}
\AXC{$\varphi_{99}$\lbl{A2}}
\BIC{$a_{N} \in \mathbb{R} \to \beta \in \mathbb{R} \to a_{N} - \beta \in \mathbb{R}$}
\AXC{$(a_{N} \in \mathbb{R} \to \beta \in \mathbb{R} \to a_{N} - \beta \in \mathbb{R}) \to \varphi_{222}$\lbl{A1}}
\BIC{$\varphi_{222}$}
\end{ptree}
\noindent$\Pi_{79}$:
\begin{ptree}
\AXC{$\forall x\,\forall y\,(x \in \mathbb{R} \to y \in \mathbb{R} \to x - y \in \mathbb{R})$\lbl{[F1]}}
\AXC{$\varphi_{67}$\lbl{A2}}
\BIC{$\forall y\,(\alpha \in \mathbb{R} \to y \in \mathbb{R} \to \alpha - y \in \mathbb{R})$}
\AXC{$\varphi_{77}$\lbl{A2}}
\BIC{$\alpha \in \mathbb{R} \to a_{N} \in \mathbb{R} \to \alpha - a_{N} \in \mathbb{R}$}
\AXC{$(\alpha \in \mathbb{R} \to a_{N} \in \mathbb{R} \to \alpha - a_{N} \in \mathbb{R}) \to \varphi_{24} \to \varphi_{206}$\lbl{A1}}
\BIC{$\varphi_{24} \to \varphi_{206}$}
\end{ptree}
\noindent$\Pi_{80}$:
\begin{ptree}
\AXC{$\varphi_{64}$\lbl{[F20]}}
\AXC{$\varphi_{64} \to \varphi_{109}$\lbl{A2}}
\BIC{$\varphi_{109}$}
\AXC{$\varphi_{109} \to \varphi_{130}$\lbl{A2}}
\BIC{$\varphi_{130}$}
\AXC{$\varphi_{130} \to \varphi_{173}$\lbl{A2}}
\BIC{$\varphi_{173}$}
\AXC{$\varphi_{173} \to \alpha - a_{N} \in \mathbb{R} \to \varphi_{179}$\lbl{A2}}
\BIC{$\alpha - a_{N} \in \mathbb{R} \to \varphi_{179}$}
\AXC{$\varphi_{259}$\lbl{A1}}
\BIC{$(\mathrm{Conv}(a,\alpha) \to \mathrm{Conv}(b,\beta) \to \varphi_{149}) \to \varphi_{257}$}
\end{ptree}
\noindent$\Pi_{81}$:
\begin{ptree}
\AXC{$\stackrel{\Pi_{32}}{\vdots}$}\noLine
\UIC{$\mathrm{Conv}(a,\alpha) \to \varphi_{158}$}
\AXC{$\stackrel{\Pi_{11}}{\vdots}$}\noLine
\UIC{$\varphi_{24}$}
\AXC{$\stackrel{\Pi_{79}}{\vdots}$}\noLine
\UIC{$\varphi_{24} \to \varphi_{206}$}
\BIC{$\varphi_{206}$}
\BIC{$\mathrm{Conv}(a,\alpha) \to \mathrm{Conv}(b,\beta) \to \varphi_{149}$}
\AXC{$\stackrel{\Pi_{80}}{\vdots}$}\noLine
\UIC{$(\mathrm{Conv}(a,\alpha) \to \mathrm{Conv}(b,\beta) \to \varphi_{149}) \to \varphi_{257}$}
\BIC{$\varphi_{257}$}
\end{ptree}
\noindent$\Pi_{82}$:
\begin{ptree}
\AXC{$\stackrel{\Pi_{10}}{\vdots}$}\noLine
\UIC{$\varphi_{23}$}
\AXC{$\stackrel{\Pi_{32}}{\vdots}$}\noLine
\UIC{$\mathrm{Conv}(a,\alpha) \to \varphi_{158}$}
\AXC{$\stackrel{\Pi_{78}}{\vdots}$}\noLine
\UIC{$\varphi_{222}$}
\BIC{$\varphi_{23} \to \mathrm{Conv}(a,\alpha) \to \mathrm{Conv}(b,\beta) \to \varphi_{148}$}
\BIC{$\mathrm{Conv}(a,\alpha) \to \mathrm{Conv}(b,\beta) \to \varphi_{148}$}
\AXC{$\stackrel{\Pi_{15}}{\vdots}$}\noLine
\UIC{$\mathrm{Conv}(a,\alpha) \to \mathrm{Conv}(b,\beta) \to \varphi_{151}$}
\AXC{$\stackrel{\Pi_{81}}{\vdots}$}\noLine
\UIC{$\varphi_{257}$}
\BIC{$(\mathrm{Conv}(a,\alpha) \to \mathrm{Conv}(b,\beta) \to \varphi_{148}) \to \varphi_{251}$}
\BIC{$\varphi_{251}$}
\end{ptree}
\noindent$\Pi_{83}$:
\begin{ptree}
\AXC{$\stackrel{\Pi_{77}}{\vdots}$}\noLine
\UIC{$\mathrm{Conv}(a,\alpha) \to \mathrm{Conv}(b,\beta) \to \varphi_{165}$}
\AXC{$\stackrel{\Pi_{15}}{\vdots}$}\noLine
\UIC{$\mathrm{Conv}(a,\alpha) \to \mathrm{Conv}(b,\beta) \to \varphi_{151}$}
\AXC{$\stackrel{\Pi_{82}}{\vdots}$}\noLine
\UIC{$\varphi_{251}$}
\BIC{$(\mathrm{Conv}(a,\alpha) \to \mathrm{Conv}(b,\beta) \to \varphi_{165}) \to \varphi_{219}$}
\BIC{$\varphi_{219}$}
\end{ptree}
\noindent$\Pi_{84}$:
\begin{ptree}
\AXC{$\stackrel{\Pi_{37}}{\vdots}$}\noLine
\UIC{$\forall x\,\forall y\,(x = y \to y = x)$}
\AXC{$\forall x\,\forall y\,(x = y \to y = x) \to \forall y\,(\alpha - \beta = y \to y = \alpha - \beta)$\lbl{A2}}
\BIC{$\forall y\,(\alpha - \beta = y \to y = \alpha - \beta)$}
\AXC{$\forall y\,(\alpha - \beta = y \to y = \alpha - \beta) \to \varphi_{44}$\lbl{A2}}
\BIC{$\varphi_{44}$}
\end{ptree}
\noindent$\Pi_{85}$:
\begin{ptree}
\AXC{$\stackrel{\Pi_{10}}{\vdots}$}\noLine
\UIC{$\varphi_{23}$}
\AXC{$\stackrel{\Pi_{11}}{\vdots}$}\noLine
\UIC{$\varphi_{24}$}
\AXC{$\varphi_{20}$\lbl{[F16]}}
\AXC{$\varphi_{20} \to \varphi_{21}$\lbl{A2}}
\BIC{$\varphi_{21}$}
\AXC{$\varphi_{21} \to \varphi_{22}$\lbl{A2}}
\BIC{$\varphi_{22}$}
\AXC{$\varphi_{22} \to \varphi_{54}$\lbl{A2}}
\BIC{$\varphi_{54}$}
\AXC{$\varphi_{54} \to \varphi_{24} \to \varphi_{23} \to \varphi_{212}$\lbl{A1}}
\BIC{$\varphi_{24} \to \varphi_{23} \to \varphi_{212}$}
\BIC{$\varphi_{23} \to \varphi_{212}$}
\BIC{$\varphi_{212}$}
\end{ptree}
\noindent$\Pi_{86}$:
\begin{ptree}
\AXC{$\stackrel{\Pi_{84}}{\vdots}$}\noLine
\UIC{$\varphi_{44}$}
\AXC{$\stackrel{\Pi_{32}}{\vdots}$}\noLine
\UIC{$\mathrm{Conv}(a,\alpha) \to \varphi_{158}$}
\AXC{$\stackrel{\Pi_{85}}{\vdots}$}\noLine
\UIC{$\varphi_{212}$}
\BIC{$\mathrm{Conv}(a,\alpha) \to \mathrm{Conv}(b,\beta) \to \varphi_{170}$}
\AXC{$\varphi_{235}$\lbl{A1}}
\BIC{$\varphi_{44} \to \mathrm{Conv}(a,\alpha) \to \mathrm{Conv}(b,\beta) \to \varphi_{171}$}
\BIC{$\mathrm{Conv}(a,\alpha) \to \mathrm{Conv}(b,\beta) \to \varphi_{171}$}
\AXC{$\varphi_{248}$\lbl{A1}}
\BIC{$(\mathrm{Conv}(a,\alpha) \to \varphi_{187}) \to \varphi_{127} \to \varphi_{184}$}
\end{ptree}
\noindent$\Pi_{87}$:
\begin{ptree}
\AXC{$\stackrel{\Pi_{2}}{\vdots}$}\noLine
\UIC{$a_{N} - b_{N} + (b_{N} - \beta) = a_{N} - \beta \to \varphi_{60}$}
\AXC{$\stackrel{\Pi_{61}}{\vdots}$}\noLine
\UIC{$\varphi_{200}$}
\BIC{$\mathrm{Conv}(a,\alpha) \to \mathrm{Conv}(b,\beta) \to \varphi_{161}$}
\AXC{$\stackrel{\Pi_{83}}{\vdots}$}\noLine
\UIC{$\varphi_{219}$}
\BIC{$\mathrm{Conv}(a,\alpha) \to \varphi_{187}$}
\AXC{$\stackrel{\Pi_{86}}{\vdots}$}\noLine
\UIC{$(\mathrm{Conv}(a,\alpha) \to \varphi_{187}) \to \varphi_{127} \to \varphi_{184}$}
\BIC{$\varphi_{127} \to \varphi_{184}$}
\end{ptree}
\noindent$\Pi_{88}$:
\begin{ptree}
\AXC{$\stackrel{\Pi_{12}}{\vdots}$}\noLine
\UIC{$\mathrm{Conv}(a,\alpha) \to \varphi_{155}$}
\AXC{$\forall e \in \mathbb{R}\,(e \cdot 2^{-1} + e \cdot 2^{-1} = e)$\lbl{[F36]}}
\AXC{$\varphi_{94}$\lbl{A2}}
\BIC{$\alpha - \beta \in \mathbb{R} \to (\alpha - \beta) \cdot 2^{-1} + (\alpha - \beta) \cdot 2^{-1} = \alpha - \beta$}
\AXC{$\varphi_{234}$\lbl{A1}}
\BIC{$(\mathrm{Conv}(a,\alpha) \to \varphi_{155}) \to \varphi_{182}$}
\BIC{$\varphi_{182}$}
\AXC{$\varphi_{182} \to \varphi_{184} \to \varphi_{197}$\lbl{A1}}
\BIC{$\varphi_{184} \to \varphi_{197}$}
\end{ptree}
\noindent$\Pi_{89}$:
\begin{ptree}
\AXC{$z = w \to \varphi_{59}$\lbl{A6}}
\UIC{$\forall w\,(z = w \to \varphi_{59})$}
\UIC{$\forall z\,\forall w\,(z = w \to \varphi_{59})$}
\AXC{$\forall z\,\forall w\,(z = w \to \varphi_{59}) \to \varphi_{122}$\lbl{A2}}
\BIC{$\varphi_{122}$}
\AXC{$\varphi_{122} \to \varphi_{127}$\lbl{A2}}
\BIC{$\varphi_{127}$}
\AXC{$\stackrel{\Pi_{87}}{\vdots}$}\noLine
\UIC{$\varphi_{127} \to \varphi_{184}$}
\BIC{$\varphi_{184}$}
\AXC{$\stackrel{\Pi_{88}}{\vdots}$}\noLine
\UIC{$\varphi_{184} \to \varphi_{197}$}
\BIC{$\varphi_{197}$}
\end{ptree}
\noindent$\Pi_{90}$:
\begin{ptree}
\AXC{$\stackrel{\Pi_{1}}{\vdots}$}\noLine
\UIC{$\varphi_{93}$}
\AXC{$\varphi_{93} \to (\alpha - \beta) \cdot 2^{-1} + (\alpha - \beta) \cdot 2^{-1} = \alpha - \beta \to \varphi_{27}$\lbl{A2}}
\BIC{$(\alpha - \beta) \cdot 2^{-1} + (\alpha - \beta) \cdot 2^{-1} = \alpha - \beta \to \varphi_{27}$}
\AXC{$\stackrel{\Pi_{89}}{\vdots}$}\noLine
\UIC{$\varphi_{197}$}
\BIC{$\mathrm{Conv}(a,\alpha) \to \mathrm{Conv}(b,\beta) \to \varphi_{142}$}
\AXC{$\varphi_{210}$\lbl{A1}}
\BIC{$\neg \alpha - \beta < \alpha - \beta \to \mathrm{Conv}(a,\alpha) \to \varphi_{141}$}
\end{ptree}
\noindent$\Pi_{91}$:
\begin{ptree}
\AXC{$\forall x\,(\neg x < x)$\lbl{[O1]}}
\AXC{$\forall x\,(\neg x < x) \to \neg \alpha - \beta < \alpha - \beta$\lbl{A2}}
\BIC{$\neg \alpha - \beta < \alpha - \beta$}
\AXC{$\stackrel{\Pi_{90}}{\vdots}$}\noLine
\UIC{$\neg \alpha - \beta < \alpha - \beta \to \mathrm{Conv}(a,\alpha) \to \varphi_{141}$}
\BIC{$\mathrm{Conv}(a,\alpha) \to \varphi_{141}$}
\end{ptree}
\noindent$\Pi_{92}$:
\begin{ptree}
\AXC{$\stackrel{\Pi_{91}}{\vdots}$}\noLine
\UIC{$\mathrm{Conv}(a,\alpha) \to \varphi_{141}$}
\AXC{$(\mathrm{Conv}(a,\alpha) \to \varphi_{141}) \to \varphi_{34} \to \neg (N \in \mathbb{N} \land \varphi_{78})$\lbl{A1}}
\BIC{$\varphi_{34} \to \neg (N \in \mathbb{N} \land \varphi_{78})$}
\UIC{$\forall N\,(\varphi_{34} \to \neg (N \in \mathbb{N} \land \varphi_{78}))$}
\end{ptree}
\noindent$\Pi_{93}$:
\begin{ptree}
\AXC{$\stackrel{\Pi_{11}}{\vdots}$}\noLine
\UIC{$\varphi_{24}$}
\AXC{$\stackrel{\Pi_{10}}{\vdots}$}\noLine
\UIC{$\varphi_{23}$}
\AXC{$\forall x\,\forall y\,(x \in \mathbb{R} \to y \in \mathbb{R} \to x < y \to 0 < y - x)$\lbl{[F24]}}
\AXC{$\varphi_{80}$\lbl{A2}}
\BIC{$\forall y\,(\beta \in \mathbb{R} \to y \in \mathbb{R} \to \beta < y \to 0 < y - \beta)$}
\AXC{$\varphi_{81}$\lbl{A2}}
\BIC{$\beta \in \mathbb{R} \to \alpha \in \mathbb{R} \to \beta < \alpha \to 0 < \alpha - \beta$}
\AXC{$\varphi_{199}$\lbl{A1}}
\BIC{$\varphi_{23} \to \varphi_{24} \to \varphi_{35} \to \mathrm{Conv}(a,\alpha) \to \varphi_{26}$}
\BIC{$\varphi_{24} \to \varphi_{35} \to \mathrm{Conv}(a,\alpha) \to \varphi_{26}$}
\BIC{$\varphi_{35} \to \mathrm{Conv}(a,\alpha) \to \varphi_{26}$}
\end{ptree}
\noindent$\Pi_{94}$:
\begin{ptree}
\AXC{$\stackrel{\Pi_{10}}{\vdots}$}\noLine
\UIC{$\varphi_{23}$}
\AXC{$\stackrel{\Pi_{11}}{\vdots}$}\noLine
\UIC{$\varphi_{24}$}
\AXC{$\forall x\,\forall y\,(x \in \mathbb{R} \to y \in \mathbb{R} \to x - y \in \mathbb{R})$\lbl{[F1]}}
\AXC{$\varphi_{67}$\lbl{A2}}
\BIC{$\forall y\,(\alpha \in \mathbb{R} \to y \in \mathbb{R} \to \alpha - y \in \mathbb{R})$}
\AXC{$\varphi_{61}$\lbl{A2}}
\BIC{$\alpha \in \mathbb{R} \to \beta \in \mathbb{R} \to \alpha - \beta \in \mathbb{R}$}
\AXC{$\varphi_{194}$\lbl{A1}}
\BIC{$\varphi_{24} \to \varphi_{23} \to \mathrm{Conv}(a,\alpha) \to \varphi_{29}$}
\BIC{$\varphi_{23} \to \mathrm{Conv}(a,\alpha) \to \varphi_{29}$}
\BIC{$\mathrm{Conv}(a,\alpha) \to \varphi_{29}$}
\end{ptree}
\noindent$\Pi_{95}$:
\begin{ptree}
\AXC{$\forall e \in \mathbb{R}\,(0 < e \to 0 < e \cdot 2^{-1})$\lbl{[F35]}}
\AXC{$\varphi_{74}$\lbl{A2}}
\BIC{$\alpha - \beta \in \mathbb{R} \to 0 < \alpha - \beta \to 0 < (\alpha - \beta) \cdot 2^{-1}$}
\AXC{$(\alpha - \beta \in \mathbb{R} \to 0 < \alpha - \beta \to 0 < (\alpha - \beta) \cdot 2^{-1}) \to \varphi_{192}$\lbl{A1}}
\BIC{$\varphi_{192}$}
\end{ptree}
\noindent$\Pi_{96}$:
\begin{ptree}
\AXC{$\varphi_{35}$\lbl{A1}}
\AXC{$\stackrel{\Pi_{93}}{\vdots}$}\noLine
\UIC{$\varphi_{35} \to \mathrm{Conv}(a,\alpha) \to \varphi_{26}$}
\BIC{$\mathrm{Conv}(a,\alpha) \to \varphi_{26}$}
\AXC{$\stackrel{\Pi_{94}}{\vdots}$}\noLine
\UIC{$\mathrm{Conv}(a,\alpha) \to \varphi_{29}$}
\AXC{$\stackrel{\Pi_{95}}{\vdots}$}\noLine
\UIC{$\varphi_{192}$}
\BIC{$(\mathrm{Conv}(a,\alpha) \to \varphi_{26}) \to \mathrm{Conv}(a,\alpha) \to \varphi_{42}$}
\BIC{$\mathrm{Conv}(a,\alpha) \to \varphi_{42}$}
\end{ptree}
\noindent$\Pi_{97}$:
\begin{ptree}
\AXC{$\forall x\,\forall y\,(x \in \mathbb{R} \to y \in \mathbb{R} \to x \cdot y \in \mathbb{R})$\lbl{(R1)}}
\AXC{$\varphi_{82}$\lbl{A2}}
\BIC{$\forall y\,(\alpha - \beta \in \mathbb{R} \to y \in \mathbb{R} \to (\alpha - \beta) \cdot y \in \mathbb{R})$}
\AXC{$\varphi_{106}$\lbl{A2}}
\BIC{$\alpha - \beta \in \mathbb{R} \to 2^{-1} \in \mathbb{R} \to (\alpha - \beta) \cdot 2^{-1} \in \mathbb{R}$}
\end{ptree}
\noindent$\Pi_{98}$:
\begin{ptree}
\AXC{$\stackrel{\Pi_{97}}{\vdots}$}\noLine
\UIC{$\alpha - \beta \in \mathbb{R} \to 2^{-1} \in \mathbb{R} \to (\alpha - \beta) \cdot 2^{-1} \in \mathbb{R}$}
\AXC{$(\alpha - \beta \in \mathbb{R} \to 2^{-1} \in \mathbb{R} \to (\alpha - \beta) \cdot 2^{-1} \in \mathbb{R}) \to \varphi_{172}$\lbl{A1}}
\BIC{$\varphi_{172}$}
\end{ptree}
\noindent$\Pi_{99}$:
\begin{ptree}
\AXC{$\forall a\,\forall \alpha\,\forall b\,\forall \beta\,(\varphi_{117})$\lbl{[C9]}}
\AXC{$\forall a\,\forall \alpha\,\forall b\,\forall \beta\,(\varphi_{117}) \to \forall \alpha\,\forall b\,\forall \beta\,(\varphi_{117})$\lbl{A2}}
\BIC{$\forall \alpha\,\forall b\,\forall \beta\,(\varphi_{117})$}
\AXC{$\forall \alpha\,\forall b\,\forall \beta\,(\varphi_{117}) \to \forall b\,\forall \beta\,(\varphi_{117})$\lbl{A2}}
\BIC{$\forall b\,\forall \beta\,(\varphi_{117})$}
\end{ptree}
\noindent$\Pi_{100}$:
\begin{ptree}
\AXC{$\stackrel{\Pi_{99}}{\vdots}$}\noLine
\UIC{$\forall b\,\forall \beta\,(\varphi_{117})$}
\AXC{$\forall b\,\forall \beta\,(\varphi_{117}) \to \forall \beta\,(\varphi_{117})$\lbl{A2}}
\BIC{$\forall \beta\,(\varphi_{117})$}
\AXC{$\forall \beta\,(\varphi_{117}) \to \varphi_{117}$\lbl{A2}}
\BIC{$\varphi_{117}$}
\AXC{$\varphi_{117} \to \mathrm{Conv}(a,\alpha) \to \mathrm{Conv}(b,\beta) \to \varphi_{129}$\lbl{A2}}
\BIC{$\mathrm{Conv}(a,\alpha) \to \mathrm{Conv}(b,\beta) \to \varphi_{129}$}
\end{ptree}
\noindent$\Pi_{101}$:
\begin{ptree}
\AXC{$\mathrm{Conv}(a,\alpha) \to \mathrm{Conv}(a,\alpha)$\lbl{A1}}
\AXC{$\stackrel{\Pi_{100}}{\vdots}$}\noLine
\UIC{$\mathrm{Conv}(a,\alpha) \to \mathrm{Conv}(b,\beta) \to \varphi_{129}$}
\AXC{$\varphi_{240}$\lbl{A1}}
\BIC{$(\mathrm{Conv}(a,\alpha) \to \mathrm{Conv}(a,\alpha)) \to \varphi_{207}$}
\BIC{$\varphi_{207}$}
\end{ptree}
\noindent$\Pi_{102}$:
\begin{ptree}
\AXC{$\mathrm{Conv}(a,\alpha) \to \mathrm{Conv}(b,\beta) \to \mathrm{Conv}(b,\beta)$\lbl{A1}}
\AXC{$\stackrel{\Pi_{101}}{\vdots}$}\noLine
\UIC{$\varphi_{207}$}
\BIC{$\varphi_{202}$}
\end{ptree}
\noindent$\Pi_{103}$:
\begin{ptree}
\AXC{$\stackrel{\Pi_{96}}{\vdots}$}\noLine
\UIC{$\mathrm{Conv}(a,\alpha) \to \varphi_{42}$}
\AXC{$2^{-1} \in \mathbb{R}$\lbl{[F31]}}
\AXC{$\stackrel{\Pi_{94}}{\vdots}$}\noLine
\UIC{$\mathrm{Conv}(a,\alpha) \to \varphi_{29}$}
\AXC{$\stackrel{\Pi_{98}}{\vdots}$}\noLine
\UIC{$\varphi_{172}$}
\BIC{$2^{-1} \in \mathbb{R} \to \mathrm{Conv}(a,\alpha) \to \varphi_{46}$}
\BIC{$\mathrm{Conv}(a,\alpha) \to \varphi_{46}$}
\AXC{$\stackrel{\Pi_{102}}{\vdots}$}\noLine
\UIC{$\varphi_{202}$}
\BIC{$(\mathrm{Conv}(a,\alpha) \to \varphi_{42}) \to \mathrm{Conv}(a,\alpha) \to \varphi_{133}$}
\BIC{$\mathrm{Conv}(a,\alpha) \to \varphi_{133}$}
\AXC{$\varphi_{253}$\lbl{A1}}
\BIC{$\forall N\,(\varphi_{34} \to \neg (N \in \mathbb{N} \land \varphi_{78})) \to \varphi_{204} \to (\varphi_{34} \to \forall N\,(\varphi_{34})) \to \varphi_{33}$}
\end{ptree}
\noindent$\Pi_{104}$:
\begin{ptree}
\AXC{$\varphi_{34} \to \forall N\,(\varphi_{34})$\lbl{A4}}
\AXC{$\varphi_{204}$\lbl{A3}}
\AXC{$\stackrel{\Pi_{92}}{\vdots}$}\noLine
\UIC{$\forall N\,(\varphi_{34} \to \neg (N \in \mathbb{N} \land \varphi_{78}))$}
\AXC{$\stackrel{\Pi_{103}}{\vdots}$}\noLine
\UIC{$\forall N\,(\varphi_{34} \to \neg (N \in \mathbb{N} \land \varphi_{78})) \to \varphi_{204} \to (\varphi_{34} \to \forall N\,(\varphi_{34})) \to \varphi_{33}$}
\BIC{$\varphi_{204} \to (\varphi_{34} \to \forall N\,(\varphi_{34})) \to \varphi_{33}$}
\BIC{$(\varphi_{34} \to \forall N\,(\varphi_{34})) \to \varphi_{33}$}
\BIC{$\varphi_{33}$}
\end{ptree}
\noindent$\Pi_{105}$:
\begin{ptree}
\AXC{$\forall x\,(\neg x < x)$\lbl{[O1]}}
\AXC{$\forall x\,(\neg x < x) \to \neg 0 < 0$\lbl{A2}}
\BIC{$\neg 0 < 0$}
\AXC{$\stackrel{\Pi_{104}}{\vdots}$}\noLine
\UIC{$\varphi_{33}$}
\AXC{$\varphi_{33} \to \neg 0 < 0 \to \mathrm{Conv}(a,\alpha) \to \varphi_{18}$\lbl{A1}}
\BIC{$\neg 0 < 0 \to \mathrm{Conv}(a,\alpha) \to \varphi_{18}$}
\BIC{$\mathrm{Conv}(a,\alpha) \to \varphi_{18}$}
\end{ptree}
\noindent$\Pi_{106}$:
\begin{ptree}
\AXC{$\forall x\,\forall y\,(x \in \mathbb{R} \to y \in \mathbb{R} \to \neg x < y \to y \leq x)$\lbl{[O8]}}
\AXC{$\varphi_{57}$\lbl{A2}}
\BIC{$\forall y\,(\beta \in \mathbb{R} \to y \in \mathbb{R} \to \neg \beta < y \to y \leq \beta)$}
\AXC{$\varphi_{58}$\lbl{A2}}
\BIC{$\beta \in \mathbb{R} \to \alpha \in \mathbb{R} \to \neg \beta < \alpha \to \alpha \leq \beta$}
\end{ptree}
\noindent$\Pi_{107}$:
\begin{ptree}
\AXC{$\stackrel{\Pi_{11}}{\vdots}$}\noLine
\UIC{$\varphi_{24}$}
\AXC{$\stackrel{\Pi_{10}}{\vdots}$}\noLine
\UIC{$\varphi_{23}$}
\AXC{$\stackrel{\Pi_{106}}{\vdots}$}\noLine
\UIC{$\beta \in \mathbb{R} \to \alpha \in \mathbb{R} \to \neg \beta < \alpha \to \alpha \leq \beta$}
\AXC{$(\beta \in \mathbb{R} \to \alpha \in \mathbb{R} \to \neg \beta < \alpha \to \alpha \leq \beta) \to \varphi_{23} \to \varphi_{24} \to \varphi_{19} \to \varphi_{13}$\lbl{A1}}
\BIC{$\varphi_{23} \to \varphi_{24} \to \varphi_{19} \to \varphi_{13}$}
\BIC{$\varphi_{24} \to \varphi_{19} \to \varphi_{13}$}
\BIC{$\varphi_{19} \to \varphi_{13}$}
\end{ptree}
\noindent 主木:
\begin{ptree}
\AXC{$\stackrel{\Pi_{105}}{\vdots}$}\noLine
\UIC{$\mathrm{Conv}(a,\alpha) \to \varphi_{18}$}
\AXC{$(\mathrm{Conv}(a,\alpha) \to \varphi_{18}) \to \varphi_{19}$\lbl{A1}}
\BIC{$\varphi_{19}$}
\AXC{$\stackrel{\Pi_{107}}{\vdots}$}\noLine
\UIC{$\varphi_{19} \to \varphi_{13}$}
\BIC{$\varphi_{13}$}
\UIC{$\forall \beta\,(\varphi_{13})$}
\UIC{$\forall \alpha\,\forall \beta\,(\varphi_{13})$}
\UIC{$\forall b\,\forall \alpha\,\forall \beta\,(\varphi_{13})$}
\UIC{$\forall a\,\forall b\,\forall \alpha\,\forall \beta\,(\varphi_{13})$}
\end{ptree}
\begin{abbrevlist}
\item $\varphi_{1} :\equiv \forall \alpha\,(\mathrm{Conv}(b,\alpha) \to \mathrm{Seq}(b)) \to \mathrm{Conv}(b,\beta) \to \mathrm{Seq}(b)$
\item $\varphi_{2} :\equiv \forall \alpha\,(\mathrm{Conv}(a,\alpha) \to \mathrm{Seq}(a)) \to \mathrm{Conv}(a,\alpha) \to \mathrm{Seq}(a)$
\item $\varphi_{3} :\equiv \forall w\,(z = w \to z = (\alpha - \beta) \cdot 2^{-1} \to w = (\alpha - \beta) \cdot 2^{-1})$
\item $\varphi_{4} :\equiv \forall a\,\forall \alpha\,(\mathrm{Conv}(a,\alpha) \to \mathrm{Seq}(a)) \to \forall \alpha\,(\mathrm{Conv}(b,\alpha) \to \mathrm{Seq}(b))$
\item $\varphi_{5} :\equiv \forall a\,\forall \alpha\,(\mathrm{Conv}(a,\alpha) \to \mathrm{Seq}(a)) \to \forall \alpha\,(\mathrm{Conv}(a,\alpha) \to \mathrm{Seq}(a))$
\item $\varphi_{6} :\equiv \forall \alpha\,(\mathrm{Conv}(b,\alpha) \to \alpha \in \mathbb{R}) \to \mathrm{Conv}(b,\beta) \to \beta \in \mathbb{R}$
\item $\varphi_{7} :\equiv \forall \alpha\,(\mathrm{Conv}(a,\alpha) \to \alpha \in \mathbb{R}) \to \mathrm{Conv}(a,\alpha) \to \alpha \in \mathbb{R}$
\item $\varphi_{8} :\equiv \forall x\,\forall y\,(x = y \to y = x) \to \forall y\,(a_{N} - \beta = y \to y = a_{N} - \beta)$
\item $\varphi_{9} :\equiv \forall w\,(z = w \to z < (\alpha - \beta) \cdot 2^{-1} \to w < (\alpha - \beta) \cdot 2^{-1})$
\item $\varphi_{10} :\equiv \forall a\,\forall \alpha\,(\mathrm{Conv}(a,\alpha) \to \alpha \in \mathbb{R}) \to \forall \alpha\,(\mathrm{Conv}(b,\alpha) \to \alpha \in \mathbb{R})$
\item $\varphi_{11} :\equiv \forall a\,\forall \alpha\,(\mathrm{Conv}(a,\alpha) \to \alpha \in \mathbb{R}) \to \forall \alpha\,(\mathrm{Conv}(a,\alpha) \to \alpha \in \mathbb{R})$
\item $\varphi_{12} :\equiv \forall y\,\forall z\,(0 + (\alpha - \beta) \cdot 2^{-1} = y \to y = z \to 0 + (\alpha - \beta) \cdot 2^{-1} = z)$
\item $\varphi_{13} :\equiv \mathrm{Conv}(a,\alpha) \to \mathrm{Conv}(b,\beta) \to \forall n \in \mathbb{N}\,(a_{n} \leq b_{n}) \to \alpha \leq \beta$
\item $\varphi_{14} :\equiv \mathrm{Conv}(a,\alpha) \to \mathrm{Conv}(b,\beta) \to \forall n \in \mathbb{N}\,(a_{n} \leq b_{n}) \to \mathrm{Seq}(b)$
\item $\varphi_{15} :\equiv \mathrm{Conv}(a,\alpha) \to \mathrm{Conv}(b,\beta) \to \forall n \in \mathbb{N}\,(a_{n} \leq b_{n}) \to \mathrm{Seq}(a)$
\item $\varphi_{16} :\equiv \forall x\,\forall y\,\forall z\,(x \in \mathbb{R} \to y \in \mathbb{R} \to z \in \mathbb{R} \to x \leq y \to x + z \leq y + z)$
\item $\varphi_{17} :\equiv \forall x\,\forall y\,(x = y \to y = x) \to \forall y\,(a_{N} - b_{N} = y \to y = a_{N} - b_{N})$
\item $\varphi_{18} :\equiv \mathrm{Conv}(b,\beta) \to \forall n \in \mathbb{N}\,(a_{n} \leq b_{n}) \to \beta < \alpha \to \neg \beta < \alpha$
\item $\varphi_{19} :\equiv \mathrm{Conv}(a,\alpha) \to \mathrm{Conv}(b,\beta) \to \forall n \in \mathbb{N}\,(a_{n} \leq b_{n}) \to \neg \beta < \alpha$
\item $\varphi_{20} :\equiv \forall a\,\forall b\,\forall x\,(a \in \mathbb{R} \to b \in \mathbb{R} \to x \in \mathbb{R} \to a - b = a - x + (x - b))$
\item $\varphi_{21} :\equiv \forall b\,\forall x\,(\alpha \in \mathbb{R} \to b \in \mathbb{R} \to x \in \mathbb{R} \to \alpha - b = \alpha - x + (x - b))$
\item $\varphi_{22} :\equiv \forall x\,(\alpha \in \mathbb{R} \to \beta \in \mathbb{R} \to x \in \mathbb{R} \to \alpha - \beta = \alpha - x + (x - \beta))$
\item $\varphi_{23} :\equiv \mathrm{Conv}(a,\alpha) \to \mathrm{Conv}(b,\beta) \to \forall n \in \mathbb{N}\,(a_{n} \leq b_{n}) \to \beta \in \mathbb{R}$
\item $\varphi_{24} :\equiv \mathrm{Conv}(a,\alpha) \to \mathrm{Conv}(b,\beta) \to \forall n \in \mathbb{N}\,(a_{n} \leq b_{n}) \to \alpha \in \mathbb{R}$
\item $\varphi_{25} :\equiv (\alpha - \beta) \cdot 2^{-1} \in \mathbb{R} \to (\alpha - \beta) \cdot 2^{-1} + 0 = (\alpha - \beta) \cdot 2^{-1}$
\item $\varphi_{26} :\equiv \mathrm{Conv}(b,\beta) \to \forall n \in \mathbb{N}\,(a_{n} \leq b_{n}) \to \beta < \alpha \to 0 < \alpha - \beta$
\item $\varphi_{27} :\equiv \alpha - \beta < (\alpha - \beta) \cdot 2^{-1} + (\alpha - \beta) \cdot 2^{-1} \to \alpha - \beta < \alpha - \beta$
\item $\varphi_{28} :\equiv \alpha - a_{N} + (a_{N} - \beta) < (\alpha - \beta) \cdot 2^{-1} + (\alpha - \beta) \cdot 2^{-1}$
\item $\varphi_{29} :\equiv \mathrm{Conv}(b,\beta) \to \forall n \in \mathbb{N}\,(a_{n} \leq b_{n}) \to \beta < \alpha \to \alpha - \beta \in \mathbb{R}$
\item $\varphi_{30} :\equiv (\alpha - \beta) \cdot 2^{-1} \in \mathbb{R} \to -(-((\alpha - \beta) \cdot 2^{-1})) = (\alpha - \beta) \cdot 2^{-1}$
\item $\varphi_{31} :\equiv \mathrm{Conv}(b,\beta) \to \forall n \in \mathbb{N}\,(a_{n} \leq b_{n}) \to N \in \mathbb{N} \to a_{N} \leq b_{N}$
\item $\varphi_{32} :\equiv (\alpha - \beta) \cdot 2^{-1} \in \mathbb{R} \to 0 + (\alpha - \beta) \cdot 2^{-1} = (\alpha - \beta) \cdot 2^{-1} + 0$
\item $\varphi_{33} :\equiv \mathrm{Conv}(a,\alpha) \to \mathrm{Conv}(b,\beta) \to \forall n \in \mathbb{N}\,(a_{n} \leq b_{n}) \to \beta < \alpha \to 0 < 0$
\item $\varphi_{34} :\equiv \mathrm{Conv}(a,\alpha) \land \mathrm{Conv}(b,\beta) \land \forall n \in \mathbb{N}\,(a_{n} \leq b_{n}) \land \beta < \alpha \land \neg 0 < 0$
\item $\varphi_{35} :\equiv \mathrm{Conv}(a,\alpha) \to \mathrm{Conv}(b,\beta) \to \forall n \in \mathbb{N}\,(a_{n} \leq b_{n}) \to \beta < \alpha \to \beta < \alpha$
\item $\varphi_{36} :\equiv -((\alpha - \beta) \cdot 2^{-1}) < b_{N} - \beta \land b_{N} - \beta < (\alpha - \beta) \cdot 2^{-1}$
\item $\varphi_{37} :\equiv \mathrm{Conv}(b,\beta) \to \forall n \in \mathbb{N}\,(a_{n} \leq b_{n}) \to \forall n \in \mathbb{N}\,(a_{n} \leq b_{n})$
\item $\varphi_{38} :\equiv -(a_{N} - \alpha) < (\alpha - \beta) \cdot 2^{-1} \to \alpha - a_{N} < (\alpha - \beta) \cdot 2^{-1}$
\item $\varphi_{39} :\equiv -((\alpha - \beta) \cdot 2^{-1}) < a_{N} - \alpha \land a_{N} - \alpha < (\alpha - \beta) \cdot 2^{-1}$
\item $\varphi_{40} :\equiv \forall n\,(\mathrm{Seq}(b) \to n \in \mathbb{N} \to b_{n} \in \mathbb{R}) \to \mathrm{Seq}(b) \to N \in \mathbb{N} \to b_{N} \in \mathbb{R}$
\item $\varphi_{41} :\equiv \forall n\,(\mathrm{Seq}(a) \to n \in \mathbb{N} \to a_{n} \in \mathbb{R}) \to \mathrm{Seq}(a) \to N \in \mathbb{N} \to a_{N} \in \mathbb{R}$
\item $\varphi_{42} :\equiv \mathrm{Conv}(b,\beta) \to \forall n \in \mathbb{N}\,(a_{n} \leq b_{n}) \to \beta < \alpha \to 0 < (\alpha - \beta) \cdot 2^{-1}$
\item $\varphi_{43} :\equiv \forall y\,\forall z\,(a_{N} \in \mathbb{R} \to y \in \mathbb{R} \to z \in \mathbb{R} \to a_{N} \leq y \to a_{N} + z \leq y + z)$
\item $\varphi_{44} :\equiv \alpha - \beta = \alpha - a_{N} + (a_{N} - \beta) \to \alpha - a_{N} + (a_{N} - \beta) = \alpha - \beta$
\item $\varphi_{45} :\equiv \lvert a_{N} - \alpha \rvert < (\alpha - \beta) \cdot 2^{-1} \land \lvert b_{N} - \beta \rvert < (\alpha - \beta) \cdot 2^{-1}$
\item $\varphi_{46} :\equiv \mathrm{Conv}(b,\beta) \to \forall n \in \mathbb{N}\,(a_{n} \leq b_{n}) \to \beta < \alpha \to (\alpha - \beta) \cdot 2^{-1} \in \mathbb{R}$
\item $\varphi_{47} :\equiv (\alpha - \beta) \cdot 2^{-1} + 0 = (\alpha - \beta) \cdot 2^{-1} \to 0 + (\alpha - \beta) \cdot 2^{-1} = (\alpha - \beta) \cdot 2^{-1}$
\item $\varphi_{48} :\equiv \forall a\,\forall n\,(\mathrm{Seq}(a) \to n \in \mathbb{N} \to a_{n} \in \mathbb{R}) \to \forall n\,(\mathrm{Seq}(b) \to n \in \mathbb{N} \to b_{n} \in \mathbb{R})$
\item $\varphi_{49} :\equiv \forall a\,\forall n\,(\mathrm{Seq}(a) \to n \in \mathbb{N} \to a_{n} \in \mathbb{R}) \to \forall n\,(\mathrm{Seq}(a) \to n \in \mathbb{N} \to a_{n} \in \mathbb{R})$
\item $\varphi_{50} :\equiv z = w \to a_{N} - b_{N} + (b_{N} - \beta) = z \to a_{N} - b_{N} + (b_{N} - \beta) = w$
\item $\varphi_{51} :\equiv \forall x\,\forall y\,\forall u\,\forall v\,(x \in \mathbb{R} \to y \in \mathbb{R} \to u \in \mathbb{R} \to v \in \mathbb{R} \to x \leq y \to u \leq v \to x + u \leq y + v)$
\item $\varphi_{52} :\equiv \forall b\,\forall x\,(a_{N} \in \mathbb{R} \to b \in \mathbb{R} \to x \in \mathbb{R} \to a_{N} - b = a_{N} - x + (x - b))$
\item $\varphi_{53} :\equiv \forall x\,(a_{N} \in \mathbb{R} \to \beta \in \mathbb{R} \to x \in \mathbb{R} \to a_{N} - \beta = a_{N} - x + (x - \beta))$
\item $\varphi_{54} :\equiv \alpha \in \mathbb{R} \to \beta \in \mathbb{R} \to a_{N} \in \mathbb{R} \to \alpha - \beta = \alpha - a_{N} + (a_{N} - \beta)$
\item $\varphi_{55} :\equiv \forall z\,\forall w\,(z = w \to z = 0 \to w = 0) \to \forall w\,(a_{N} + -b_{N} = w \to a_{N} + -b_{N} = 0 \to w = 0)$
\item $\varphi_{56} :\equiv \forall x\,\forall y\,(x \in \mathbb{R} \to y \in \mathbb{R} \to x + y = y + x) \to \forall y\,(0 \in \mathbb{R} \to y \in \mathbb{R} \to 0 + y = y + 0)$
\item $\varphi_{57} :\equiv \forall x\,\forall y\,(x \in \mathbb{R} \to y \in \mathbb{R} \to \neg x < y \to y \leq x) \to \forall y\,(\beta \in \mathbb{R} \to y \in \mathbb{R} \to \neg \beta < y \to y \leq \beta)$
\item $\varphi_{58} :\equiv \forall y\,(\beta \in \mathbb{R} \to y \in \mathbb{R} \to \neg \beta < y \to y \leq \beta) \to \beta \in \mathbb{R} \to \alpha \in \mathbb{R} \to \neg \beta < \alpha \to \alpha \leq \beta$
\item $\varphi_{59} :\equiv z < (\alpha - \beta) \cdot 2^{-1} + (\alpha - \beta) \cdot 2^{-1} \to w < (\alpha - \beta) \cdot 2^{-1} + (\alpha - \beta) \cdot 2^{-1}$
\item $\varphi_{60} :\equiv a_{N} - b_{N} + (b_{N} - \beta) = (\alpha - \beta) \cdot 2^{-1} \to a_{N} - \beta = (\alpha - \beta) \cdot 2^{-1}$
\item $\varphi_{61} :\equiv \forall y\,(\alpha \in \mathbb{R} \to y \in \mathbb{R} \to \alpha - y \in \mathbb{R}) \to \alpha \in \mathbb{R} \to \beta \in \mathbb{R} \to \alpha - \beta \in \mathbb{R}$
\item $\varphi_{62} :\equiv \forall w\,(-(a_{N} - \alpha) = w \to -(a_{N} - \alpha) < (\alpha - \beta) \cdot 2^{-1} \to w < (\alpha - \beta) \cdot 2^{-1})$
\item $\varphi_{63} :\equiv z = w \to a_{N} - b_{N} + (b_{N} - \beta) < z \to a_{N} - b_{N} + (b_{N} - \beta) < w$
\item $\varphi_{64} :\equiv \forall x\,\forall y\,\forall u\,\forall v\,(x \in \mathbb{R} \to y \in \mathbb{R} \to u \in \mathbb{R} \to v \in \mathbb{R} \to x < y \to u \leq v \to x + u < y + v)$
\item $\varphi_{65} :\equiv -(a_{N} - \alpha) < -(-((\alpha - \beta) \cdot 2^{-1})) \to -(a_{N} - \alpha) < (\alpha - \beta) \cdot 2^{-1}$
\item $\varphi_{66} :\equiv -((\alpha - \beta) \cdot 2^{-1}) < a_{N} - \alpha \to -(a_{N} - \alpha) < -(-((\alpha - \beta) \cdot 2^{-1}))$
\item $\varphi_{67} :\equiv \forall x\,\forall y\,(x \in \mathbb{R} \to y \in \mathbb{R} \to x - y \in \mathbb{R}) \to \forall y\,(\alpha \in \mathbb{R} \to y \in \mathbb{R} \to \alpha - y \in \mathbb{R})$
\item $\varphi_{68} :\equiv \forall x \in \mathbb{R}\,(-x \in \mathbb{R}) \to (\alpha - \beta) \cdot 2^{-1} \in \mathbb{R} \to -((\alpha - \beta) \cdot 2^{-1}) \in \mathbb{R}$
\item $\varphi_{69} :\equiv b_{N} - \beta \leq (\alpha - \beta) \cdot 2^{-1} \to a_{N} - b_{N} + (b_{N} - \beta) \leq 0 + (\alpha - \beta) \cdot 2^{-1}$
\item $\varphi_{70} :\equiv a_{N} - b_{N} + (b_{N} - \beta) < (\alpha - \beta) \cdot 2^{-1} \to a_{N} - \beta < (\alpha - \beta) \cdot 2^{-1}$
\item $\varphi_{71} :\equiv \forall z\,\forall w\,(z = w \to z < 0 \to w < 0) \to \forall w\,(a_{N} + -b_{N} = w \to a_{N} + -b_{N} < 0 \to w < 0)$
\item $\varphi_{72} :\equiv \forall z\,(a_{N} \in \mathbb{R} \to b_{N} \in \mathbb{R} \to z \in \mathbb{R} \to a_{N} \leq b_{N} \to a_{N} + z \leq b_{N} + z)$
\item $\varphi_{73} :\equiv \forall x\,\forall y\,(x \in \mathbb{R} \to y \in \mathbb{R} \to -(y - x) = x - y) \to \forall y\,(\alpha \in \mathbb{R} \to y \in \mathbb{R} \to -(y - \alpha) = \alpha - y)$
\item $\varphi_{74} :\equiv \forall e \in \mathbb{R}\,(0 < e \to 0 < e \cdot 2^{-1}) \to \alpha - \beta \in \mathbb{R} \to 0 < \alpha - \beta \to 0 < (\alpha - \beta) \cdot 2^{-1}$
\item $\varphi_{75} :\equiv (\mathrm{Conv}(a,\alpha) \to \mathrm{Conv}(b,\beta) \to \mathrm{Conv}(b,\beta)) \to \varphi_{14}$
\item $\varphi_{76} :\equiv \forall z\,(0 + (\alpha - \beta) \cdot 2^{-1} = (\alpha - \beta) \cdot 2^{-1} + 0 \to (\alpha - \beta) \cdot 2^{-1} + 0 = z \to 0 + (\alpha - \beta) \cdot 2^{-1} = z)$
\item $\varphi_{77} :\equiv \forall y\,(\alpha \in \mathbb{R} \to y \in \mathbb{R} \to \alpha - y \in \mathbb{R}) \to \alpha \in \mathbb{R} \to a_{N} \in \mathbb{R} \to \alpha - a_{N} \in \mathbb{R}$
\item $\varphi_{78} :\equiv \forall n \in \mathbb{N}\,(N \leq n \to \lvert a_{n} - \alpha \rvert < (\alpha - \beta) \cdot 2^{-1} \land \lvert b_{n} - \beta \rvert < (\alpha - \beta) \cdot 2^{-1})$
\item $\varphi_{79} :\equiv a_{N} \in \mathbb{R} \to \beta \in \mathbb{R} \to b_{N} \in \mathbb{R} \to a_{N} - \beta = a_{N} - b_{N} + (b_{N} - \beta)$
\item $\varphi_{80} :\equiv \forall x\,\forall y\,(x \in \mathbb{R} \to y \in \mathbb{R} \to x < y \to 0 < y - x) \to \forall y\,(\beta \in \mathbb{R} \to y \in \mathbb{R} \to \beta < y \to 0 < y - \beta)$
\item $\varphi_{81} :\equiv \forall y\,(\beta \in \mathbb{R} \to y \in \mathbb{R} \to \beta < y \to 0 < y - \beta) \to \beta \in \mathbb{R} \to \alpha \in \mathbb{R} \to \beta < \alpha \to 0 < \alpha - \beta$
\item $\varphi_{82} :\equiv \forall x\,\forall y\,(x \in \mathbb{R} \to y \in \mathbb{R} \to x \cdot y \in \mathbb{R}) \to \forall y\,(\alpha - \beta \in \mathbb{R} \to y \in \mathbb{R} \to (\alpha - \beta) \cdot y \in \mathbb{R})$
\item $\varphi_{83} :\equiv \forall x\,\forall y\,(x \in \mathbb{R} \to y \in \mathbb{R} \to x - y \in \mathbb{R}) \to \forall y\,(b_{N} \in \mathbb{R} \to y \in \mathbb{R} \to b_{N} - y \in \mathbb{R})$
\item $\varphi_{84} :\equiv \forall x\,\forall y\,(x \in \mathbb{R} \to y \in \mathbb{R} \to x - y \in \mathbb{R}) \to \forall y\,(a_{N} \in \mathbb{R} \to y \in \mathbb{R} \to a_{N} - y \in \mathbb{R})$
\item $\varphi_{85} :\equiv a_{N} - \beta = a_{N} - b_{N} + (b_{N} - \beta) \to a_{N} - b_{N} + (b_{N} - \beta) = a_{N} - \beta$
\item $\varphi_{86} :\equiv (\mathrm{Conv}(a,\alpha) \to \mathrm{Conv}(b,\beta) \to \mathrm{Conv}(b,\beta)) \to \varphi_{23}$
\item $\varphi_{87} :\equiv b_{N} + -b_{N} = 0 \to a_{N} + -b_{N} = b_{N} + -b_{N} \to a_{N} + -b_{N} = 0$
\item $\varphi_{88} :\equiv \forall w\,(b_{N} + -b_{N} = w \to a_{N} + -b_{N} = b_{N} + -b_{N} \to a_{N} + -b_{N} = w)$
\item $\varphi_{89} :\equiv \forall w\,(-(-((\alpha - \beta) \cdot 2^{-1})) = w \to -(a_{N} - \alpha) < -(-((\alpha - \beta) \cdot 2^{-1})) \to -(a_{N} - \alpha) < w)$
\item $\varphi_{90} :\equiv \forall y\,(b_{N} - \beta \in \mathbb{R} \to y \in \mathbb{R} \to \lvert b_{N} - \beta \rvert < y \to -y < b_{N} - \beta \land b_{N} - \beta < y)$
\item $\varphi_{91} :\equiv (\mathrm{Conv}(a,\alpha) \to \mathrm{Seq}(a)) \to (\mathrm{Conv}(a,\alpha) \to \mathrm{Conv}(a,\alpha)) \to \varphi_{15}$
\item $\varphi_{92} :\equiv \forall y\,(a_{N} - \alpha \in \mathbb{R} \to y \in \mathbb{R} \to \lvert a_{N} - \alpha \rvert < y \to -y < a_{N} - \alpha \land a_{N} - \alpha < y)$
\item $\varphi_{93} :\equiv \forall w\,((\alpha - \beta) \cdot 2^{-1} + (\alpha - \beta) \cdot 2^{-1} = w \to \alpha - \beta < (\alpha - \beta) \cdot 2^{-1} + (\alpha - \beta) \cdot 2^{-1} \to \alpha - \beta < w)$
\item $\varphi_{94} :\equiv \forall e \in \mathbb{R}\,(e \cdot 2^{-1} + e \cdot 2^{-1} = e) \to \alpha - \beta \in \mathbb{R} \to (\alpha - \beta) \cdot 2^{-1} + (\alpha - \beta) \cdot 2^{-1} = \alpha - \beta$
\item $\varphi_{95} :\equiv b_{N} \in \mathbb{R} \to -b_{N} \in \mathbb{R} \to a_{N} \leq b_{N} \to a_{N} + -b_{N} \leq b_{N} + -b_{N}$
\item $\varphi_{96} :\equiv b_{N} + -b_{N} = 0 \to a_{N} + -b_{N} < b_{N} + -b_{N} \to a_{N} + -b_{N} < 0$
\item $\varphi_{97} :\equiv \forall y\,(b_{N} \in \mathbb{R} \to y \in \mathbb{R} \to b_{N} - y \in \mathbb{R}) \to b_{N} \in \mathbb{R} \to \beta \in \mathbb{R} \to b_{N} - \beta \in \mathbb{R}$
\item $\varphi_{98} :\equiv \forall y\,(a_{N} \in \mathbb{R} \to y \in \mathbb{R} \to a_{N} - y \in \mathbb{R}) \to a_{N} \in \mathbb{R} \to \alpha \in \mathbb{R} \to a_{N} - \alpha \in \mathbb{R}$
\item $\varphi_{99} :\equiv \forall y\,(a_{N} \in \mathbb{R} \to y \in \mathbb{R} \to a_{N} - y \in \mathbb{R}) \to a_{N} \in \mathbb{R} \to \beta \in \mathbb{R} \to a_{N} - \beta \in \mathbb{R}$
\item $\varphi_{100} :\equiv \forall w\,(b_{N} + -b_{N} = w \to a_{N} + -b_{N} < b_{N} + -b_{N} \to a_{N} + -b_{N} < w)$
\item $\varphi_{101} :\equiv \forall y\,(-((\alpha - \beta) \cdot 2^{-1}) \in \mathbb{R} \to y \in \mathbb{R} \to -((\alpha - \beta) \cdot 2^{-1}) < y \to -y < -(-((\alpha - \beta) \cdot 2^{-1})))$
\item $\varphi_{102} :\equiv \forall y\,(\alpha \in \mathbb{R} \to y \in \mathbb{R} \to -(y - \alpha) = \alpha - y) \to \alpha \in \mathbb{R} \to a_{N} \in \mathbb{R} \to -(a_{N} - \alpha) = \alpha - a_{N}$
\item $\varphi_{103} :\equiv a_{N} - b_{N} + (b_{N} - \beta) = 0 + (\alpha - \beta) \cdot 2^{-1} \to a_{N} - b_{N} + (b_{N} - \beta) = (\alpha - \beta) \cdot 2^{-1}$
\item $\varphi_{104} :\equiv (\mathrm{Conv}(a,\alpha) \to \alpha \in \mathbb{R}) \to (\mathrm{Conv}(a,\alpha) \to \mathrm{Conv}(a,\alpha)) \to \varphi_{24}$
\item $\varphi_{105} :\equiv \forall w\,(a_{N} - b_{N} + (b_{N} - \beta) = w \to a_{N} - b_{N} + (b_{N} - \beta) = (\alpha - \beta) \cdot 2^{-1} \to w = (\alpha - \beta) \cdot 2^{-1})$
\item $\varphi_{106} :\equiv \forall y\,(\alpha - \beta \in \mathbb{R} \to y \in \mathbb{R} \to (\alpha - \beta) \cdot y \in \mathbb{R}) \to \alpha - \beta \in \mathbb{R} \to 2^{-1} \in \mathbb{R} \to (\alpha - \beta) \cdot 2^{-1} \in \mathbb{R}$
\item $\varphi_{107} :\equiv a_{N} - b_{N} + (b_{N} - \beta) < 0 + (\alpha - \beta) \cdot 2^{-1} \to a_{N} - b_{N} + (b_{N} - \beta) < (\alpha - \beta) \cdot 2^{-1}$
\item $\varphi_{108} :\equiv \forall y\,(a_{N} \in \mathbb{R} \to y \in \mathbb{R} \to a_{N} - y \in \mathbb{R}) \to a_{N} \in \mathbb{R} \to b_{N} \in \mathbb{R} \to a_{N} - b_{N} \in \mathbb{R}$
\item $\varphi_{109} :\equiv \forall y\,\forall u\,\forall v\,(\alpha - a_{N} \in \mathbb{R} \to y \in \mathbb{R} \to u \in \mathbb{R} \to v \in \mathbb{R} \to \alpha - a_{N} < y \to u \leq v \to \alpha - a_{N} + u < y + v)$
\item $\varphi_{110} :\equiv \forall w\,(a_{N} - b_{N} + (b_{N} - \beta) = w \to a_{N} - b_{N} + (b_{N} - \beta) < (\alpha - \beta) \cdot 2^{-1} \to w < (\alpha - \beta) \cdot 2^{-1})$
\item $\varphi_{111} :\equiv \forall w\,(0 + (\alpha - \beta) \cdot 2^{-1} = w \to a_{N} - b_{N} + (b_{N} - \beta) = 0 + (\alpha - \beta) \cdot 2^{-1} \to a_{N} - b_{N} + (b_{N} - \beta) = w)$
\item $\varphi_{112} :\equiv \forall u\,\forall v\,(a_{N} - b_{N} \in \mathbb{R} \to 0 \in \mathbb{R} \to u \in \mathbb{R} \to v \in \mathbb{R} \to a_{N} - b_{N} \leq 0 \to u \leq v \to a_{N} - b_{N} + u \leq 0 + v)$
\item $\varphi_{113} :\equiv \forall w\,(0 + (\alpha - \beta) \cdot 2^{-1} = w \to a_{N} - b_{N} + (b_{N} - \beta) < 0 + (\alpha - \beta) \cdot 2^{-1} \to a_{N} - b_{N} + (b_{N} - \beta) < w)$
\item $\varphi_{114} :\equiv \forall y\,\forall u\,\forall v\,(a_{N} - b_{N} \in \mathbb{R} \to y \in \mathbb{R} \to u \in \mathbb{R} \to v \in \mathbb{R} \to a_{N} - b_{N} \leq y \to u \leq v \to a_{N} - b_{N} + u \leq y + v)$
\item $\varphi_{115} :\equiv (a_{N} - b_{N} = a_{N} + -b_{N} \to a_{N} + -b_{N} = a_{N} - b_{N}) \to a_{N} + -b_{N} = a_{N} - b_{N}$
\item $\varphi_{116} :\equiv \forall y\,(a_{N} - b_{N} = y \to y = a_{N} - b_{N}) \to a_{N} - b_{N} = a_{N} + -b_{N} \to a_{N} + -b_{N} = a_{N} - b_{N}$
\item $\varphi_{117} :\equiv \forall \varepsilon\,(\mathrm{Conv}(a,\alpha) \to \mathrm{Conv}(b,\beta) \to \varepsilon \in \mathbb{R} \to 0 < \varepsilon \to \exists N \in \mathbb{N}\,\forall n \in \mathbb{N}\,(N \leq n \to \lvert a_{n} - \alpha \rvert < \varepsilon \land \lvert b_{n} - \beta \rvert < \varepsilon))$
\item $\varphi_{118} :\equiv \alpha - a_{N} < (\alpha - \beta) \cdot 2^{-1} \to a_{N} - \beta \leq (\alpha - \beta) \cdot 2^{-1} \to \varphi_{28}$
\item $\varphi_{119} :\equiv (\alpha - \beta) \cdot 2^{-1} \in \mathbb{R} \to \lvert b_{N} - \beta \rvert < (\alpha - \beta) \cdot 2^{-1} \to \varphi_{36}$
\item $\varphi_{120} :\equiv (\alpha - \beta) \cdot 2^{-1} \in \mathbb{R} \to \lvert a_{N} - \alpha \rvert < (\alpha - \beta) \cdot 2^{-1} \to \varphi_{39}$
\item $\varphi_{121} :\equiv (\forall n \in \mathbb{N}\,(a_{n} \leq b_{n}) \to N \in \mathbb{N} \to a_{N} \leq b_{N}) \to \mathrm{Conv}(a,\alpha) \to \varphi_{31}$
\item $\varphi_{122} :\equiv \forall w\,(\alpha - a_{N} + (a_{N} - \beta) = w \to \varphi_{28} \to w < (\alpha - \beta) \cdot 2^{-1} + (\alpha - \beta) \cdot 2^{-1})$
\item $\varphi_{123} :\equiv N \in \mathbb{N} \land \varphi_{78} \to -(-((\alpha - \beta) \cdot 2^{-1})) = (\alpha - \beta) \cdot 2^{-1}$
\item $\varphi_{124} :\equiv \forall w\,(a_{N} + -b_{N} = w \to a_{N} + -b_{N} = 0 \to w = 0) \to a_{N} + -b_{N} = a_{N} - b_{N} \to a_{N} + -b_{N} = 0 \to a_{N} - b_{N} = 0$
\item $\varphi_{125} :\equiv \beta < \alpha \to N \in \mathbb{N} \land \varphi_{78} \to (\alpha - \beta) \cdot 2^{-1} + 0 = (\alpha - \beta) \cdot 2^{-1}$
\item $\varphi_{126} :\equiv \beta < \alpha \to N \in \mathbb{N} \land \varphi_{78} \to 0 + (\alpha - \beta) \cdot 2^{-1} = (\alpha - \beta) \cdot 2^{-1}$
\item $\varphi_{127} :\equiv \alpha - a_{N} + (a_{N} - \beta) = \alpha - \beta \to \varphi_{28} \to \alpha - \beta < (\alpha - \beta) \cdot 2^{-1} + (\alpha - \beta) \cdot 2^{-1}$
\item $\varphi_{128} :\equiv \forall v\,(a_{N} - b_{N} \in \mathbb{R} \to 0 \in \mathbb{R} \to b_{N} - \beta \in \mathbb{R} \to v \in \mathbb{R} \to a_{N} - b_{N} \leq 0 \to b_{N} - \beta \leq v \to a_{N} - b_{N} + (b_{N} - \beta) \leq 0 + v)$
\item $\varphi_{129} :\equiv (\alpha - \beta) \cdot 2^{-1} \in \mathbb{R} \to 0 < (\alpha - \beta) \cdot 2^{-1} \to \exists N \in \mathbb{N}\,(\varphi_{78})$
\item $\varphi_{130} :\equiv \forall u\,\forall v\,(\alpha - a_{N} \in \mathbb{R} \to (\alpha - \beta) \cdot 2^{-1} \in \mathbb{R} \to u \in \mathbb{R} \to v \in \mathbb{R} \to \alpha - a_{N} < (\alpha - \beta) \cdot 2^{-1} \to u \leq v \to \alpha - a_{N} + u < (\alpha - \beta) \cdot 2^{-1} + v)$
\item $\varphi_{131} :\equiv \beta < \alpha \to N \in \mathbb{N} \land \varphi_{78} \to 0 + (\alpha - \beta) \cdot 2^{-1} = (\alpha - \beta) \cdot 2^{-1} + 0$
\item $\varphi_{132} :\equiv \forall z\,\forall w\,(z = w \to -(a_{N} - \alpha) < z \to -(a_{N} - \alpha) < w) \to \varphi_{89}$
\item $\varphi_{133} :\equiv \mathrm{Conv}(b,\beta) \to \forall n \in \mathbb{N}\,(a_{n} \leq b_{n}) \to \beta < \alpha \to \exists N \in \mathbb{N}\,(\varphi_{78})$
\item $\varphi_{134} :\equiv \beta < \alpha \to N \in \mathbb{N} \land \varphi_{78} \to (\alpha - \beta) \cdot 2^{-1} + (\alpha - \beta) \cdot 2^{-1} = \alpha - \beta$
\item $\varphi_{135} :\equiv \forall w\,(a_{N} + -b_{N} = w \to a_{N} + -b_{N} < 0 \to w < 0) \to a_{N} + -b_{N} = a_{N} - b_{N} \to a_{N} + -b_{N} < 0 \to a_{N} - b_{N} < 0$
\item $\varphi_{136} :\equiv N \in \mathbb{N} \land \varphi_{78} \to a_{N} - b_{N} + (b_{N} - \beta) \leq 0 + (\alpha - \beta) \cdot 2^{-1}$
\item $\varphi_{137} :\equiv \beta < \alpha \to N \in \mathbb{N} \land \varphi_{78} \to -(a_{N} - \alpha) < -(-((\alpha - \beta) \cdot 2^{-1}))$
\item $\varphi_{138} :\equiv \forall x\,\forall y\,(x \in \mathbb{R} \to y \in \mathbb{R} \to \lvert x \rvert < y \to -y < x \land x < y) \to \varphi_{90}$
\item $\varphi_{139} :\equiv \forall x\,\forall y\,(x \in \mathbb{R} \to y \in \mathbb{R} \to \lvert x \rvert < y \to -y < x \land x < y) \to \varphi_{92}$
\item $\varphi_{140} :\equiv \beta < \alpha \to N \in \mathbb{N} \land \varphi_{78} \to \alpha - \beta < (\alpha - \beta) \cdot 2^{-1} + (\alpha - \beta) \cdot 2^{-1}$
\item $\varphi_{141} :\equiv \mathrm{Conv}(b,\beta) \to \forall n \in \mathbb{N}\,(a_{n} \leq b_{n}) \to \beta < \alpha \to N \in \mathbb{N} \land \varphi_{78} \to 0 < 0$
\item $\varphi_{142} :\equiv \forall n \in \mathbb{N}\,(a_{n} \leq b_{n}) \to \beta < \alpha \to N \in \mathbb{N} \land \varphi_{78} \to \alpha - \beta < \alpha - \beta$
\item $\varphi_{143} :\equiv \mathrm{Conv}(b,\beta) \to \forall n \in \mathbb{N}\,(a_{n} \leq b_{n}) \to \beta < \alpha \to N \in \mathbb{N} \land \varphi_{78} \to N \in \mathbb{N}$
\item $\varphi_{144} :\equiv \forall n \in \mathbb{N}\,(a_{n} \leq b_{n}) \to \beta < \alpha \to N \in \mathbb{N} \land \varphi_{78} \to a_{N} - b_{N} \leq 0$
\item $\varphi_{145} :\equiv \beta < \alpha \to N \in \mathbb{N} \land \varphi_{78} \to a_{N} - b_{N} + (b_{N} - \beta) \leq (\alpha - \beta) \cdot 2^{-1}$
\item $\varphi_{146} :\equiv \forall n \in \mathbb{N}\,(a_{n} \leq b_{n}) \to \beta < \alpha \to N \in \mathbb{N} \land \varphi_{78} \to b_{N} - \beta \in \mathbb{R}$
\item $\varphi_{147} :\equiv \forall n \in \mathbb{N}\,(a_{n} \leq b_{n}) \to \beta < \alpha \to N \in \mathbb{N} \land \varphi_{78} \to a_{N} - \alpha \in \mathbb{R}$
\item $\varphi_{148} :\equiv \forall n \in \mathbb{N}\,(a_{n} \leq b_{n}) \to \beta < \alpha \to N \in \mathbb{N} \land \varphi_{78} \to a_{N} - \beta \in \mathbb{R}$
\item $\varphi_{149} :\equiv \forall n \in \mathbb{N}\,(a_{n} \leq b_{n}) \to \beta < \alpha \to N \in \mathbb{N} \land \varphi_{78} \to \alpha - a_{N} \in \mathbb{R}$
\item $\varphi_{150} :\equiv 0 \in \mathbb{R} \to b_{N} - \beta \in \mathbb{R} \to (\alpha - \beta) \cdot 2^{-1} \in \mathbb{R} \to a_{N} - b_{N} \leq 0 \to \varphi_{69}$
\item $\varphi_{151} :\equiv \forall n \in \mathbb{N}\,(a_{n} \leq b_{n}) \to \beta < \alpha \to N \in \mathbb{N} \land \varphi_{78} \to (\alpha - \beta) \cdot 2^{-1} \in \mathbb{R}$
\item $\varphi_{152} :\equiv \forall n \in \mathbb{N}\,(a_{n} \leq b_{n}) \to \beta < \alpha \to N \in \mathbb{N} \land \varphi_{78} \to b_{N} + -b_{N} = 0$
\item $\varphi_{153} :\equiv \forall n \in \mathbb{N}\,(a_{n} \leq b_{n}) \to \beta < \alpha \to N \in \mathbb{N} \land \varphi_{78} \to a_{N} + -b_{N} \leq 0$
\item $\varphi_{154} :\equiv \forall n \in \mathbb{N}\,(a_{n} \leq b_{n}) \to \beta < \alpha \to N \in \mathbb{N} \land \varphi_{78} \to a_{N} - b_{N} \in \mathbb{R}$
\item $\varphi_{155} :\equiv \mathrm{Conv}(b,\beta) \to \forall n \in \mathbb{N}\,(a_{n} \leq b_{n}) \to \beta < \alpha \to N \in \mathbb{N} \land \varphi_{78} \to \alpha - \beta \in \mathbb{R}$
\item $\varphi_{156} :\equiv \mathrm{Conv}(b,\beta) \to \forall n \in \mathbb{N}\,(a_{n} \leq b_{n}) \to \beta < \alpha \to N \in \mathbb{N} \land \varphi_{78} \to b_{N} \in \mathbb{R}$
\item $\varphi_{157} :\equiv \mathrm{Conv}(b,\beta) \to \forall n \in \mathbb{N}\,(a_{n} \leq b_{n}) \to \beta < \alpha \to N \in \mathbb{N} \land \varphi_{78} \to a_{N} \leq b_{N}$
\item $\varphi_{158} :\equiv \mathrm{Conv}(b,\beta) \to \forall n \in \mathbb{N}\,(a_{n} \leq b_{n}) \to \beta < \alpha \to N \in \mathbb{N} \land \varphi_{78} \to a_{N} \in \mathbb{R}$
\item $\varphi_{159} :\equiv \forall n \in \mathbb{N}\,(a_{n} \leq b_{n}) \to \beta < \alpha \to N \in \mathbb{N} \land \varphi_{78} \to -((\alpha - \beta) \cdot 2^{-1}) \in \mathbb{R}$
\item $\varphi_{160} :\equiv \forall n \in \mathbb{N}\,(a_{n} \leq b_{n}) \to \beta < \alpha \to N \in \mathbb{N} \land \varphi_{78} \to b_{N} - \beta \leq (\alpha - \beta) \cdot 2^{-1}$
\item $\varphi_{161} :\equiv \forall n \in \mathbb{N}\,(a_{n} \leq b_{n}) \to \beta < \alpha \to N \in \mathbb{N} \land \varphi_{78} \to a_{N} - \beta \leq (\alpha - \beta) \cdot 2^{-1}$
\item $\varphi_{162} :\equiv \forall n \in \mathbb{N}\,(a_{n} \leq b_{n}) \to \beta < \alpha \to N \in \mathbb{N} \land \varphi_{78} \to -(a_{N} - \alpha) = \alpha - a_{N}$
\item $\varphi_{163} :\equiv \mathrm{Conv}(b,\beta) \to \forall n \in \mathbb{N}\,(a_{n} \leq b_{n}) \to \beta < \alpha \to N \in \mathbb{N} \land \varphi_{78} \to -b_{N} \in \mathbb{R}$
\item $\varphi_{164} :\equiv \forall n \in \mathbb{N}\,(a_{n} \leq b_{n}) \to \beta < \alpha \to N \in \mathbb{N} \land \varphi_{78} \to b_{N} - \beta < (\alpha - \beta) \cdot 2^{-1}$
\item $\varphi_{165} :\equiv \forall n \in \mathbb{N}\,(a_{n} \leq b_{n}) \to \beta < \alpha \to N \in \mathbb{N} \land \varphi_{78} \to \alpha - a_{N} < (\alpha - \beta) \cdot 2^{-1}$
\item $\varphi_{166} :\equiv \forall n \in \mathbb{N}\,(a_{n} \leq b_{n}) \to \beta < \alpha \to N \in \mathbb{N} \land \varphi_{78} \to \lvert b_{N} - \beta \rvert < (\alpha - \beta) \cdot 2^{-1}$
\item $\varphi_{167} :\equiv \forall n \in \mathbb{N}\,(a_{n} \leq b_{n}) \to \beta < \alpha \to N \in \mathbb{N} \land \varphi_{78} \to \lvert a_{N} - \alpha \rvert < (\alpha - \beta) \cdot 2^{-1}$
\item $\varphi_{168} :\equiv \forall n \in \mathbb{N}\,(a_{n} \leq b_{n}) \to \beta < \alpha \to N \in \mathbb{N} \land \varphi_{78} \to -((\alpha - \beta) \cdot 2^{-1}) < a_{N} - \alpha$
\item $\varphi_{169} :\equiv \forall n \in \mathbb{N}\,(a_{n} \leq b_{n}) \to \beta < \alpha \to N \in \mathbb{N} \land \varphi_{78} \to -(a_{N} - \alpha) < (\alpha - \beta) \cdot 2^{-1}$
\item $\varphi_{170} :\equiv \forall n \in \mathbb{N}\,(a_{n} \leq b_{n}) \to \beta < \alpha \to N \in \mathbb{N} \land \varphi_{78} \to \alpha - \beta = \alpha - a_{N} + (a_{N} - \beta)$
\item $\varphi_{171} :\equiv \forall n \in \mathbb{N}\,(a_{n} \leq b_{n}) \to \beta < \alpha \to N \in \mathbb{N} \land \varphi_{78} \to \alpha - a_{N} + (a_{N} - \beta) = \alpha - \beta$
\item $\varphi_{172} :\equiv (\mathrm{Conv}(a,\alpha) \to \varphi_{29}) \to 2^{-1} \in \mathbb{R} \to \mathrm{Conv}(a,\alpha) \to \varphi_{46}$
\item $\varphi_{173} :\equiv \forall v\,(\alpha - a_{N} \in \mathbb{R} \to (\alpha - \beta) \cdot 2^{-1} \in \mathbb{R} \to a_{N} - \beta \in \mathbb{R} \to v \in \mathbb{R} \to \alpha - a_{N} < (\alpha - \beta) \cdot 2^{-1} \to a_{N} - \beta \leq v \to \alpha - a_{N} + (a_{N} - \beta) < (\alpha - \beta) \cdot 2^{-1} + v)$
\item $\varphi_{174} :\equiv \forall n \in \mathbb{N}\,(a_{n} \leq b_{n}) \to \beta < \alpha \to N \in \mathbb{N} \land \varphi_{78} \to a_{N} + -b_{N} \leq b_{N} + -b_{N}$
\item $\varphi_{175} :\equiv \forall n \in \mathbb{N}\,(a_{n} \leq b_{n}) \to \beta < \alpha \to N \in \mathbb{N} \land \varphi_{78} \to a_{N} - \beta = a_{N} - b_{N} + (b_{N} - \beta)$
\item $\varphi_{176} :\equiv \forall n \in \mathbb{N}\,(a_{n} \leq b_{n}) \to \beta < \alpha \to N \in \mathbb{N} \land \varphi_{78} \to a_{N} - b_{N} + (b_{N} - \beta) = a_{N} - \beta$
\item $\varphi_{177} :\equiv \mathrm{Conv}(a,\alpha) \to \mathrm{Conv}(b,\beta) \to \forall n \in \mathbb{N}\,(a_{n} \leq b_{n}) \to \varphi_{125}$
\item $\varphi_{178} :\equiv \mathrm{Conv}(a,\alpha) \to \mathrm{Conv}(b,\beta) \to \forall n \in \mathbb{N}\,(a_{n} \leq b_{n}) \to \varphi_{126}$
\item $\varphi_{179} :\equiv (\alpha - \beta) \cdot 2^{-1} \in \mathbb{R} \to a_{N} - \beta \in \mathbb{R} \to (\alpha - \beta) \cdot 2^{-1} \in \mathbb{R} \to \varphi_{118}$
\item $\varphi_{180} :\equiv \mathrm{Conv}(a,\alpha) \to \mathrm{Conv}(b,\beta) \to \forall n \in \mathbb{N}\,(a_{n} \leq b_{n}) \to \varphi_{131}$
\item $\varphi_{181} :\equiv \mathrm{Conv}(a,\alpha) \to \mathrm{Conv}(b,\beta) \to \forall n \in \mathbb{N}\,(a_{n} \leq b_{n}) \to \beta < \alpha \to \varphi_{123}$
\item $\varphi_{182} :\equiv \mathrm{Conv}(a,\alpha) \to \mathrm{Conv}(b,\beta) \to \forall n \in \mathbb{N}\,(a_{n} \leq b_{n}) \to \varphi_{134}$
\item $\varphi_{183} :\equiv \mathrm{Conv}(a,\alpha) \to \mathrm{Conv}(b,\beta) \to \forall n \in \mathbb{N}\,(a_{n} \leq b_{n}) \to \varphi_{137}$
\item $\varphi_{184} :\equiv \mathrm{Conv}(a,\alpha) \to \mathrm{Conv}(b,\beta) \to \forall n \in \mathbb{N}\,(a_{n} \leq b_{n}) \to \varphi_{140}$
\item $\varphi_{185} :\equiv \mathrm{Conv}(a,\alpha) \to \mathrm{Conv}(b,\beta) \to \forall n \in \mathbb{N}\,(a_{n} \leq b_{n}) \to \varphi_{145}$
\item $\varphi_{186} :\equiv \mathrm{Conv}(a,\alpha) \to \mathrm{Conv}(b,\beta) \to \forall n \in \mathbb{N}\,(a_{n} \leq b_{n}) \to \beta < \alpha \to \varphi_{136}$
\item $\varphi_{187} :\equiv \mathrm{Conv}(b,\beta) \to \forall n \in \mathbb{N}\,(a_{n} \leq b_{n}) \to \beta < \alpha \to N \in \mathbb{N} \land \varphi_{78} \to \varphi_{28}$
\item $\varphi_{188} :\equiv \mathrm{Conv}(b,\beta) \to \forall n \in \mathbb{N}\,(a_{n} \leq b_{n}) \to \beta < \alpha \to N \in \mathbb{N} \land \varphi_{78} \to \varphi_{36}$
\item $\varphi_{189} :\equiv \mathrm{Conv}(b,\beta) \to \forall n \in \mathbb{N}\,(a_{n} \leq b_{n}) \to \beta < \alpha \to N \in \mathbb{N} \land \varphi_{78} \to \varphi_{39}$
\item $\varphi_{190} :\equiv \mathrm{Conv}(b,\beta) \to \forall n \in \mathbb{N}\,(a_{n} \leq b_{n}) \to \beta < \alpha \to N \in \mathbb{N} \land \varphi_{78} \to \varphi_{45}$
\item $\varphi_{191} :\equiv \forall n \in \mathbb{N}\,(a_{n} \leq b_{n}) \to \beta < \alpha \to N \in \mathbb{N} \land \varphi_{78} \to N \in \mathbb{N} \to N \leq N \to \varphi_{45}$
\item $\varphi_{192} :\equiv (\mathrm{Conv}(a,\alpha) \to \varphi_{29}) \to (\mathrm{Conv}(a,\alpha) \to \varphi_{26}) \to \mathrm{Conv}(a,\alpha) \to \varphi_{42}$
\item $\varphi_{193} :\equiv \mathrm{Conv}(b,\beta) \to \forall n \in \mathbb{N}\,(a_{n} \leq b_{n}) \to \beta < \alpha \to N \in \mathbb{N} \land \varphi_{78} \to \varphi_{78}$
\item $\varphi_{194} :\equiv (\alpha \in \mathbb{R} \to \beta \in \mathbb{R} \to \alpha - \beta \in \mathbb{R}) \to \varphi_{24} \to \varphi_{23} \to \mathrm{Conv}(a,\alpha) \to \varphi_{29}$
\item $\varphi_{195} :\equiv \mathrm{Conv}(b,\beta) \to \forall n \in \mathbb{N}\,(a_{n} \leq b_{n}) \to \beta < \alpha \to N \in \mathbb{N} \land \varphi_{78} \to N \in \mathbb{N} \land \varphi_{78}$
\item $\varphi_{196} :\equiv (a_{N} + -b_{N} = a_{N} - b_{N} \to a_{N} + -b_{N} = 0 \to a_{N} - b_{N} = 0) \to \mathrm{Conv}(a,\alpha) \to \mathrm{Conv}(b,\beta) \to \varphi_{144}$
\item $\varphi_{197} :\equiv ((\alpha - \beta) \cdot 2^{-1} + (\alpha - \beta) \cdot 2^{-1} = \alpha - \beta \to \varphi_{27}) \to \mathrm{Conv}(a,\alpha) \to \mathrm{Conv}(b,\beta) \to \varphi_{142}$
\item $\varphi_{198} :\equiv (-(a_{N} - \alpha) = \alpha - a_{N} \to \varphi_{38}) \to \mathrm{Conv}(a,\alpha) \to \mathrm{Conv}(b,\beta) \to \varphi_{165}$
\item $\varphi_{199} :\equiv (\beta \in \mathbb{R} \to \alpha \in \mathbb{R} \to \beta < \alpha \to 0 < \alpha - \beta) \to \varphi_{23} \to \varphi_{24} \to \varphi_{35} \to \mathrm{Conv}(a,\alpha) \to \varphi_{26}$
\item $\varphi_{200} :\equiv (a_{N} - b_{N} + (b_{N} - \beta) = a_{N} - \beta \to \varphi_{60}) \to \mathrm{Conv}(a,\alpha) \to \mathrm{Conv}(b,\beta) \to \varphi_{161}$
\item $\varphi_{201} :\equiv (-(-((\alpha - \beta) \cdot 2^{-1})) = (\alpha - \beta) \cdot 2^{-1} \to \varphi_{65}) \to \mathrm{Conv}(a,\alpha) \to \mathrm{Conv}(b,\beta) \to \varphi_{169}$
\item $\varphi_{202} :\equiv (\mathrm{Conv}(a,\alpha) \to \varphi_{46}) \to (\mathrm{Conv}(a,\alpha) \to \varphi_{42}) \to \mathrm{Conv}(a,\alpha) \to \varphi_{133}$
\item $\varphi_{203} :\equiv (\alpha \in \mathbb{R} \to \beta \in \mathbb{R} \to \alpha - \beta \in \mathbb{R}) \to \varphi_{24} \to \varphi_{23} \to \mathrm{Conv}(a,\alpha) \to \varphi_{155}$
\item $\varphi_{204} :\equiv \forall N\,(\varphi_{34} \to \neg (N \in \mathbb{N} \land \varphi_{78})) \to \forall N\,(\varphi_{34}) \to \forall N\,(\neg (N \in \mathbb{N} \land \varphi_{78}))$
\item $\varphi_{205} :\equiv (a_{N} + -b_{N} = a_{N} - b_{N} \to a_{N} + -b_{N} < 0 \to a_{N} - b_{N} < 0) \to \varphi_{196}$
\item $\varphi_{206} :\equiv (\mathrm{Conv}(a,\alpha) \to \varphi_{158}) \to \mathrm{Conv}(a,\alpha) \to \mathrm{Conv}(b,\beta) \to \varphi_{149}$
\item $\varphi_{207} :\equiv (\mathrm{Conv}(a,\alpha) \to \mathrm{Conv}(b,\beta) \to \mathrm{Conv}(b,\beta)) \to \varphi_{202}$
\item $\varphi_{208} :\equiv (\mathrm{Conv}(a,\alpha) \to \varphi_{156}) \to \mathrm{Conv}(a,\alpha) \to \mathrm{Conv}(b,\beta) \to \varphi_{152}$
\item $\varphi_{209} :\equiv (\mathrm{Conv}(a,\alpha) \to \varphi_{156}) \to \mathrm{Conv}(a,\alpha) \to \mathrm{Conv}(b,\beta) \to \varphi_{154}$
\item $\varphi_{210} :\equiv (\mathrm{Conv}(a,\alpha) \to \mathrm{Conv}(b,\beta) \to \varphi_{142}) \to \neg \alpha - \beta < \alpha - \beta \to \mathrm{Conv}(a,\alpha) \to \varphi_{141}$
\item $\varphi_{211} :\equiv (\mathrm{Conv}(a,\alpha) \to \varphi_{158}) \to \mathrm{Conv}(a,\alpha) \to \mathrm{Conv}(b,\beta) \to \varphi_{162}$
\item $\varphi_{212} :\equiv (\mathrm{Conv}(a,\alpha) \to \varphi_{158}) \to \mathrm{Conv}(a,\alpha) \to \mathrm{Conv}(b,\beta) \to \varphi_{170}$
\item $\varphi_{213} :\equiv (\mathrm{Conv}(a,\alpha) \to \mathrm{Conv}(b,\beta) \to \varphi_{151}) \to \mathrm{Conv}(a,\alpha) \to \mathrm{Conv}(b,\beta) \to \varphi_{159}$
\item $\varphi_{214} :\equiv (\mathrm{Conv}(a,\alpha) \to \varphi_{155}) \to 2^{-1} \in \mathbb{R} \to \mathrm{Conv}(a,\alpha) \to \mathrm{Conv}(b,\beta) \to \varphi_{151}$
\item $\varphi_{215} :\equiv (\mathrm{Conv}(a,\alpha) \to \varphi_{157}) \to \mathrm{Conv}(a,\alpha) \to \mathrm{Conv}(b,\beta) \to \varphi_{174}$
\item $\varphi_{216} :\equiv (\mathrm{Conv}(a,\alpha) \to \mathrm{Conv}(b,\beta) \to \varphi_{164}) \to \mathrm{Conv}(a,\alpha) \to \mathrm{Conv}(b,\beta) \to \varphi_{160}$
\item $\varphi_{217} :\equiv (\mathrm{Conv}(a,\alpha) \to \varphi_{156}) \to \mathrm{Conv}(a,\alpha) \to \mathrm{Conv}(b,\beta) \to \varphi_{175}$
\item $\varphi_{218} :\equiv (b_{N} \in \mathbb{R} \to -b_{N} \in \mathbb{R}) \to (\mathrm{Conv}(a,\alpha) \to \varphi_{156}) \to \mathrm{Conv}(a,\alpha) \to \varphi_{163}$
\item $\varphi_{219} :\equiv (\mathrm{Conv}(a,\alpha) \to \mathrm{Conv}(b,\beta) \to \varphi_{161}) \to \mathrm{Conv}(a,\alpha) \to \varphi_{187}$
\item $\varphi_{220} :\equiv (\mathrm{Conv}(a,\alpha) \to \varphi_{156}) \to \varphi_{23} \to \mathrm{Conv}(a,\alpha) \to \mathrm{Conv}(b,\beta) \to \varphi_{146}$
\item $\varphi_{221} :\equiv (\mathrm{Conv}(a,\alpha) \to \varphi_{158}) \to \varphi_{24} \to \mathrm{Conv}(a,\alpha) \to \mathrm{Conv}(b,\beta) \to \varphi_{147}$
\item $\varphi_{222} :\equiv (\mathrm{Conv}(a,\alpha) \to \varphi_{158}) \to \varphi_{23} \to \mathrm{Conv}(a,\alpha) \to \mathrm{Conv}(b,\beta) \to \varphi_{148}$
\item $\varphi_{223} :\equiv (\mathrm{Conv}(a,\alpha) \to \varphi_{188}) \to \mathrm{Conv}(a,\alpha) \to \mathrm{Conv}(b,\beta) \to \varphi_{164}$
\item $\varphi_{224} :\equiv (\mathrm{Conv}(a,\alpha) \to \varphi_{31}) \to (\mathrm{Conv}(a,\alpha) \to \varphi_{143}) \to \mathrm{Conv}(a,\alpha) \to \varphi_{157}$
\item $\varphi_{225} :\equiv (\mathrm{Conv}(a,\alpha) \to \mathrm{Conv}(b,\beta) \to \varphi_{166}) \to \mathrm{Conv}(a,\alpha) \to \varphi_{188}$
\item $\varphi_{226} :\equiv (\mathrm{Conv}(a,\alpha) \to \mathrm{Conv}(b,\beta) \to \varphi_{167}) \to \mathrm{Conv}(a,\alpha) \to \varphi_{189}$
\item $\varphi_{227} :\equiv (\mathrm{Conv}(a,\alpha) \to \varphi_{189}) \to \mathrm{Conv}(a,\alpha) \to \mathrm{Conv}(b,\beta) \to \varphi_{168}$
\item $\varphi_{228} :\equiv (\varphi_{78} \to N \in \mathbb{N} \to N \leq N \to \varphi_{45}) \to \mathrm{Conv}(a,\alpha) \to \mathrm{Conv}(b,\beta) \to \varphi_{191}$
\item $\varphi_{229} :\equiv ((\alpha - \beta) \cdot 2^{-1} \in \mathbb{R} \to -((\alpha - \beta) \cdot 2^{-1}) \in \mathbb{R}) \to \varphi_{213}$
\item $\varphi_{230} :\equiv (\mathrm{Conv}(a,\alpha) \to \varphi_{190}) \to \mathrm{Conv}(a,\alpha) \to \mathrm{Conv}(b,\beta) \to \varphi_{166}$
\item $\varphi_{231} :\equiv (\mathrm{Conv}(a,\alpha) \to \varphi_{190}) \to \mathrm{Conv}(a,\alpha) \to \mathrm{Conv}(b,\beta) \to \varphi_{167}$
\item $\varphi_{232} :\equiv (\mathrm{Seq}(b) \to N \in \mathbb{N} \to b_{N} \in \mathbb{R}) \to \varphi_{14} \to (\mathrm{Conv}(a,\alpha) \to \varphi_{143}) \to \mathrm{Conv}(a,\alpha) \to \varphi_{156}$
\item $\varphi_{233} :\equiv (\mathrm{Seq}(a) \to N \in \mathbb{N} \to a_{N} \in \mathbb{R}) \to \varphi_{15} \to (\mathrm{Conv}(a,\alpha) \to \varphi_{143}) \to \mathrm{Conv}(a,\alpha) \to \varphi_{158}$
\item $\varphi_{234} :\equiv (\alpha - \beta \in \mathbb{R} \to (\alpha - \beta) \cdot 2^{-1} + (\alpha - \beta) \cdot 2^{-1} = \alpha - \beta) \to (\mathrm{Conv}(a,\alpha) \to \varphi_{155}) \to \varphi_{182}$
\item $\varphi_{235} :\equiv (\mathrm{Conv}(a,\alpha) \to \mathrm{Conv}(b,\beta) \to \varphi_{170}) \to \varphi_{44} \to \mathrm{Conv}(a,\alpha) \to \mathrm{Conv}(b,\beta) \to \varphi_{171}$
\item $\varphi_{236} :\equiv (0 \in \mathbb{R} \to \varphi_{32}) \to 0 \in \mathbb{R} \to (\mathrm{Conv}(a,\alpha) \to \mathrm{Conv}(b,\beta) \to \varphi_{151}) \to \varphi_{180}$
\item $\varphi_{237} :\equiv (0 + (\alpha - \beta) \cdot 2^{-1} = (\alpha - \beta) \cdot 2^{-1} \to \varphi_{107}) \to (0 + (\alpha - \beta) \cdot 2^{-1} = (\alpha - \beta) \cdot 2^{-1} \to \varphi_{103}) \to \varphi_{185}$
\item $\varphi_{238} :\equiv (\mathrm{Conv}(a,\alpha) \to \mathrm{Conv}(b,\beta) \to \varphi_{175}) \to \varphi_{85} \to \mathrm{Conv}(a,\alpha) \to \mathrm{Conv}(b,\beta) \to \varphi_{176}$
\item $\varphi_{239} :\equiv \varphi_{177} \to (0 + (\alpha - \beta) \cdot 2^{-1} = (\alpha - \beta) \cdot 2^{-1} + 0 \to \varphi_{47}) \to \varphi_{178}$
\item $\varphi_{240} :\equiv (\mathrm{Conv}(a,\alpha) \to \mathrm{Conv}(b,\beta) \to \varphi_{129}) \to (\mathrm{Conv}(a,\alpha) \to \mathrm{Conv}(a,\alpha)) \to \varphi_{207}$
\item $\varphi_{241} :\equiv (\mathrm{Conv}(a,\alpha) \to \mathrm{Conv}(b,\beta) \to \varphi_{147}) \to (\mathrm{Conv}(a,\alpha) \to \mathrm{Conv}(b,\beta) \to \varphi_{168}) \to \varphi_{183}$
\item $\varphi_{242} :\equiv (\mathrm{Conv}(a,\alpha) \to \mathrm{Conv}(b,\beta) \to \varphi_{174}) \to \varphi_{96} \to \varphi_{87} \to \mathrm{Conv}(a,\alpha) \to \mathrm{Conv}(b,\beta) \to \varphi_{153}$
\item $\varphi_{243} :\equiv a_{N} + -b_{N} = a_{N} - b_{N} \to (\mathrm{Conv}(a,\alpha) \to \mathrm{Conv}(b,\beta) \to \varphi_{153}) \to \varphi_{205}$
\item $\varphi_{244} :\equiv (\mathrm{Conv}(a,\alpha) \to \mathrm{Conv}(b,\beta) \to \varphi_{144}) \to (\mathrm{Conv}(a,\alpha) \to \mathrm{Conv}(b,\beta) \to \varphi_{160}) \to \varphi_{186}$
\item $\varphi_{245} :\equiv (a_{N} \in \mathbb{R} \to b_{N} \in \mathbb{R} \to a_{N} - b_{N} \in \mathbb{R}) \to (\mathrm{Conv}(a,\alpha) \to \varphi_{158}) \to \varphi_{209}$
\item $\varphi_{246} :\equiv (\mathrm{Conv}(a,\alpha) \to \mathrm{Conv}(b,\beta) \to \varphi_{191}) \to (\mathrm{Conv}(a,\alpha) \to \varphi_{143}) \to N = N \to \mathrm{Conv}(a,\alpha) \to \varphi_{190}$
\item $\varphi_{247} :\equiv (\mathrm{Conv}(a,\alpha) \to \mathrm{Conv}(b,\beta) \to \varphi_{162}) \to (\mathrm{Conv}(a,\alpha) \to \mathrm{Conv}(b,\beta) \to \varphi_{169}) \to \varphi_{198}$
\item $\varphi_{248} :\equiv (\mathrm{Conv}(a,\alpha) \to \mathrm{Conv}(b,\beta) \to \varphi_{171}) \to (\mathrm{Conv}(a,\alpha) \to \varphi_{187}) \to \varphi_{127} \to \varphi_{184}$
\item $\varphi_{249} :\equiv (\mathrm{Conv}(a,\alpha) \to \mathrm{Conv}(b,\beta) \to \varphi_{146}) \to (\mathrm{Conv}(a,\alpha) \to \mathrm{Conv}(b,\beta) \to \varphi_{151}) \to \varphi_{225}$
\item $\varphi_{250} :\equiv (\mathrm{Conv}(a,\alpha) \to \mathrm{Conv}(b,\beta) \to \varphi_{147}) \to (\mathrm{Conv}(a,\alpha) \to \mathrm{Conv}(b,\beta) \to \varphi_{151}) \to \varphi_{226}$
\item $\varphi_{251} :\equiv (\mathrm{Conv}(a,\alpha) \to \mathrm{Conv}(b,\beta) \to \varphi_{151}) \to (\mathrm{Conv}(a,\alpha) \to \mathrm{Conv}(b,\beta) \to \varphi_{165}) \to \varphi_{219}$
\item $\varphi_{252} :\equiv (\mathrm{Conv}(a,\alpha) \to \mathrm{Conv}(b,\beta) \to \varphi_{176}) \to \varphi_{185} \to (a_{N} - b_{N} + (b_{N} - \beta) = a_{N} - \beta \to \varphi_{70}) \to \varphi_{200}$
\item $\varphi_{253} :\equiv (\mathrm{Conv}(a,\alpha) \to \varphi_{133}) \to \forall N\,(\varphi_{34} \to \neg (N \in \mathbb{N} \land \varphi_{78})) \to \varphi_{204} \to (\varphi_{34} \to \forall N\,(\varphi_{34})) \to \varphi_{33}$
\item $\varphi_{254} :\equiv (-((\alpha - \beta) \cdot 2^{-1}) \in \mathbb{R} \to a_{N} - \alpha \in \mathbb{R} \to \varphi_{66}) \to (\mathrm{Conv}(a,\alpha) \to \mathrm{Conv}(b,\beta) \to \varphi_{159}) \to \varphi_{241}$
\item $\varphi_{255} :\equiv (\mathrm{Conv}(a,\alpha) \to \varphi_{158}) \to (\mathrm{Conv}(a,\alpha) \to \varphi_{156}) \to (\mathrm{Conv}(a,\alpha) \to \varphi_{163}) \to \varphi_{215}$
\item $\varphi_{256} :\equiv (\mathrm{Conv}(a,\alpha) \to \mathrm{Conv}(b,\beta) \to \varphi_{146}) \to (\mathrm{Conv}(a,\alpha) \to \mathrm{Conv}(b,\beta) \to \varphi_{151}) \to \varphi_{244}$
\item $\varphi_{257} :\equiv (\mathrm{Conv}(a,\alpha) \to \mathrm{Conv}(b,\beta) \to \varphi_{151}) \to (\mathrm{Conv}(a,\alpha) \to \mathrm{Conv}(b,\beta) \to \varphi_{148}) \to \varphi_{251}$
\item $\varphi_{258} :\equiv (a_{N} - b_{N} \in \mathbb{R} \to \varphi_{150}) \to (\mathrm{Conv}(a,\alpha) \to \mathrm{Conv}(b,\beta) \to \varphi_{154}) \to 0 \in \mathbb{R} \to \varphi_{256}$
\item $\varphi_{259} :\equiv (\alpha - a_{N} \in \mathbb{R} \to \varphi_{179}) \to (\mathrm{Conv}(a,\alpha) \to \mathrm{Conv}(b,\beta) \to \varphi_{149}) \to \varphi_{257}$
\end{abbrevlist}
\end{fullproof}

\section{第3章 2節(続き): 有界性, 積・商の極限, 絶対値, ∞ への発散, 部分列}
% 自動生成 (prover/build.js): ConvExt  ノード数 106
\begin{fullproof}{[C13]}{$\forall a\,\forall b\,\forall \alpha\,(\mathrm{Seq}(b) \to \forall n \in \mathbb{N}\,(a_{n} = b_{n}) \to \mathrm{Conv}(a,\alpha) \to \mathrm{Conv}(b,\alpha))$}{106}\phantomsection\label{tree:ConvExt}
\noindent$\Pi_{1}$:
\begin{ptree}
\AXC{$\varphi_{18}$\lbl{A2}}
\AXC{$\mathrm{Seq}(b) \to \varphi_{46}$\lbl{A1}}
\AXC{$(\mathrm{Seq}(b) \to \varphi_{46}) \to \mathrm{Seq}(b) \to \varphi_{50}$\lbl{A1}}
\BIC{$\mathrm{Seq}(b) \to \varphi_{50}$}
\AXC{$(\mathrm{Seq}(b) \to \varphi_{50}) \to \varphi_{18} \to \mathrm{Seq}(b) \to \varphi_{48}$\lbl{A1}}
\BIC{$\varphi_{18} \to \mathrm{Seq}(b) \to \varphi_{48}$}
\BIC{$\mathrm{Seq}(b) \to \varphi_{48}$}
\end{ptree}
\noindent$\Pi_{2}$:
\begin{ptree}
\AXC{$\mathrm{Seq}(b) \to \varphi_{40}$\lbl{A1}}
\AXC{$\mathrm{Seq}(b) \to \varphi_{41}$\lbl{A1}}
\AXC{$\stackrel{\Pi_{1}}{\vdots}$}\noLine
\UIC{$\mathrm{Seq}(b) \to \varphi_{48}$}
\AXC{$\varphi_{62}$\lbl{A1}}
\BIC{$(\mathrm{Seq}(b) \to \varphi_{41}) \to (\mathrm{Seq}(b) \to \varphi_{40}) \to \mathrm{Seq}(b) \to \varphi_{43}$}
\BIC{$(\mathrm{Seq}(b) \to \varphi_{40}) \to \mathrm{Seq}(b) \to \varphi_{43}$}
\BIC{$\mathrm{Seq}(b) \to \varphi_{43}$}
\end{ptree}
\noindent$\Pi_{3}$:
\begin{ptree}
\AXC{$\mathrm{Seq}(b) \to \varphi_{41}$\lbl{A1}}
\AXC{$\forall n \in \mathbb{N}\,(a_{n} = b_{n}) \to n \in \mathbb{N} \to a_{n} = b_{n}$\lbl{A2}}
\AXC{$\varphi_{10}$\lbl{A1}}
\AXC{$\varphi_{10} \to \varphi_{31}$\lbl{A1}}
\BIC{$\varphi_{31}$}
\BIC{$\mathrm{Seq}(b) \to \forall n \in \mathbb{N}\,(a_{n} = b_{n}) \to n \in \mathbb{N} \to a_{n} = b_{n}$}
\AXC{$\varphi_{60}$\lbl{A1}}
\BIC{$(\mathrm{Seq}(b) \to \varphi_{41}) \to \mathrm{Seq}(b) \to \varphi_{42}$}
\BIC{$\mathrm{Seq}(b) \to \varphi_{42}$}
\end{ptree}
\noindent$\Pi_{4}$:
\begin{ptree}
\AXC{$\stackrel{\Pi_{2}}{\vdots}$}\noLine
\UIC{$\mathrm{Seq}(b) \to \varphi_{43}$}
\AXC{$\stackrel{\Pi_{3}}{\vdots}$}\noLine
\UIC{$\mathrm{Seq}(b) \to \varphi_{42}$}
\AXC{$(\mathrm{Seq}(b) \to \varphi_{42}) \to (\mathrm{Seq}(b) \to \varphi_{43}) \to \varphi_{55}$\lbl{A1}}
\BIC{$(\mathrm{Seq}(b) \to \varphi_{43}) \to \varphi_{55}$}
\BIC{$\varphi_{55}$}
\end{ptree}
\noindent$\Pi_{5}$:
\begin{ptree}
\AXC{$z = w \to \lvert z - \alpha \rvert < \varepsilon \to \lvert w - \alpha \rvert < \varepsilon$\lbl{A6}}
\UIC{$\forall w\,(z = w \to \lvert z - \alpha \rvert < \varepsilon \to \lvert w - \alpha \rvert < \varepsilon)$}
\UIC{$\forall z\,\forall w\,(z = w \to \lvert z - \alpha \rvert < \varepsilon \to \lvert w - \alpha \rvert < \varepsilon)$}
\AXC{$\varphi_{24}$\lbl{A2}}
\BIC{$\forall w\,(a_{n} = w \to \lvert a_{n} - \alpha \rvert < \varepsilon \to \lvert w - \alpha \rvert < \varepsilon)$}
\AXC{$\varphi_{32}$\lbl{A2}}
\BIC{$a_{n} = b_{n} \to \lvert a_{n} - \alpha \rvert < \varepsilon \to \lvert b_{n} - \alpha \rvert < \varepsilon$}
\AXC{$\stackrel{\Pi_{4}}{\vdots}$}\noLine
\UIC{$\varphi_{55}$}
\BIC{$\mathrm{Seq}(b) \to \varphi_{44}$}
\end{ptree}
\noindent$\Pi_{6}$:
\begin{ptree}
\AXC{$\varphi_{57}$\lbl{A3}}
\AXC{$\stackrel{\Pi_{5}}{\vdots}$}\noLine
\UIC{$\mathrm{Seq}(b) \to \varphi_{44}$}
\AXC{$(\mathrm{Seq}(b) \to \varphi_{44}) \to \varphi_{33} \to n \in \mathbb{N} \to N \leq n \to \lvert b_{n} - \alpha \rvert < \varepsilon$\lbl{A1}}
\BIC{$\varphi_{33} \to n \in \mathbb{N} \to N \leq n \to \lvert b_{n} - \alpha \rvert < \varepsilon$}
\UIC{$\forall n\,(\varphi_{33} \to n \in \mathbb{N} \to N \leq n \to \lvert b_{n} - \alpha \rvert < \varepsilon)$}
\AXC{$\varphi_{63}$\lbl{A1}}
\BIC{$\varphi_{57} \to (\varphi_{33} \to \forall n\,(\varphi_{33})) \to \mathrm{Seq}(b) \to \varphi_{45}$}
\BIC{$(\varphi_{33} \to \forall n\,(\varphi_{33})) \to \mathrm{Seq}(b) \to \varphi_{45}$}
\end{ptree}
\noindent$\Pi_{7}$:
\begin{ptree}
\AXC{$\mathrm{Seq}(b) \to \varphi_{46}$\lbl{A1}}
\AXC{$(\mathrm{Seq}(b) \to \varphi_{46}) \to \varphi_{39}$\lbl{A1}}
\BIC{$\varphi_{39}$}
\AXC{$\varphi_{39} \to (\mathrm{Seq}(b) \to \varphi_{45}) \to \mathrm{Seq}(b) \to \varphi_{47}$\lbl{A1}}
\BIC{$(\mathrm{Seq}(b) \to \varphi_{45}) \to \mathrm{Seq}(b) \to \varphi_{47}$}
\end{ptree}
\noindent$\Pi_{8}$:
\begin{ptree}
\AXC{$\varphi_{33} \to \forall n\,(\varphi_{33})$\lbl{A4}}
\AXC{$\stackrel{\Pi_{6}}{\vdots}$}\noLine
\UIC{$(\varphi_{33} \to \forall n\,(\varphi_{33})) \to \mathrm{Seq}(b) \to \varphi_{45}$}
\BIC{$\mathrm{Seq}(b) \to \varphi_{45}$}
\AXC{$\stackrel{\Pi_{7}}{\vdots}$}\noLine
\UIC{$(\mathrm{Seq}(b) \to \varphi_{45}) \to \mathrm{Seq}(b) \to \varphi_{47}$}
\BIC{$\mathrm{Seq}(b) \to \varphi_{47}$}
\end{ptree}
\noindent$\Pi_{9}$:
\begin{ptree}
\AXC{$\varphi_{7} \to \neg (N \in \mathbb{N} \land \forall n \in \mathbb{N}\,(N \leq n \to \lvert b_{n} - \alpha \rvert < \varepsilon))$\lbl{A2}}
\AXC{$\stackrel{\Pi_{8}}{\vdots}$}\noLine
\UIC{$\mathrm{Seq}(b) \to \varphi_{47}$}
\AXC{$(\mathrm{Seq}(b) \to \varphi_{47}) \to \varphi_{56}$\lbl{A1}}
\BIC{$\varphi_{56}$}
\BIC{$\mathrm{Seq}(b) \to \varphi_{49}$}
\AXC{$\varphi_{59}$\lbl{A1}}
\BIC{$\varphi_{37} \to \neg (N \in \mathbb{N} \land \forall n \in \mathbb{N}\,(N \leq n \to \lvert a_{n} - \alpha \rvert < \varepsilon))$}
\UIC{$\varphi_{51}$}
\end{ptree}
\noindent$\Pi_{10}$:
\begin{ptree}
\AXC{$\forall a\,\forall \alpha\,\forall \varepsilon\,(\mathrm{Conv}(a,\alpha) \to \varphi_{15})$\lbl{[C3]}}
\AXC{$\varphi_{53}$\lbl{A2}}
\BIC{$\forall \alpha\,\forall \varepsilon\,(\mathrm{Conv}(a,\alpha) \to \varphi_{15})$}
\AXC{$\varphi_{52}$\lbl{A2}}
\BIC{$\forall \varepsilon\,(\mathrm{Conv}(a,\alpha) \to \varphi_{15})$}
\AXC{$\forall \varepsilon\,(\mathrm{Conv}(a,\alpha) \to \varphi_{15}) \to \mathrm{Conv}(a,\alpha) \to \varphi_{15}$\lbl{A2}}
\BIC{$\mathrm{Conv}(a,\alpha) \to \varphi_{15}$}
\end{ptree}
\noindent$\Pi_{11}$:
\begin{ptree}
\AXC{$\mathrm{Seq}(b) \to \varphi_{9}$\lbl{A1}}
\AXC{$\mathrm{Seq}(b) \to \varphi_{2}$\lbl{A1}}
\AXC{$\stackrel{\Pi_{10}}{\vdots}$}\noLine
\UIC{$\mathrm{Conv}(a,\alpha) \to \varphi_{15}$}
\AXC{$(\mathrm{Conv}(a,\alpha) \to \varphi_{15}) \to \varphi_{58}$\lbl{A1}}
\BIC{$\varphi_{58}$}
\BIC{$(\mathrm{Seq}(b) \to \varphi_{9}) \to (\mathrm{Seq}(b) \to \varphi_{12}) \to \varphi_{34}$}
\BIC{$(\mathrm{Seq}(b) \to \varphi_{12}) \to \varphi_{34}$}
\end{ptree}
\noindent$\Pi_{12}$:
\begin{ptree}
\AXC{$\stackrel{\Pi_{9}}{\vdots}$}\noLine
\UIC{$\varphi_{51}$}
\AXC{$\mathrm{Seq}(b) \to \varphi_{12}$\lbl{A1}}
\AXC{$\stackrel{\Pi_{11}}{\vdots}$}\noLine
\UIC{$(\mathrm{Seq}(b) \to \varphi_{12}) \to \varphi_{34}$}
\BIC{$\varphi_{34}$}
\AXC{$\varphi_{34} \to \varphi_{51} \to (\varphi_{51} \to \forall N\,(\varphi_{37}) \to \varphi_{6}) \to (\varphi_{37} \to \forall N\,(\varphi_{37})) \to \varphi_{35}$\lbl{A1}}
\BIC{$\varphi_{51} \to (\varphi_{51} \to \forall N\,(\varphi_{37}) \to \varphi_{6}) \to (\varphi_{37} \to \forall N\,(\varphi_{37})) \to \varphi_{35}$}
\BIC{$(\varphi_{51} \to \forall N\,(\varphi_{37}) \to \varphi_{6}) \to (\varphi_{37} \to \forall N\,(\varphi_{37})) \to \varphi_{35}$}
\end{ptree}
\noindent$\Pi_{13}$:
\begin{ptree}
\AXC{$\varphi_{37} \to \forall N\,(\varphi_{37})$\lbl{A4}}
\AXC{$\varphi_{51} \to \forall N\,(\varphi_{37}) \to \varphi_{6}$\lbl{A3}}
\AXC{$\stackrel{\Pi_{12}}{\vdots}$}\noLine
\UIC{$(\varphi_{51} \to \forall N\,(\varphi_{37}) \to \varphi_{6}) \to (\varphi_{37} \to \forall N\,(\varphi_{37})) \to \varphi_{35}$}
\BIC{$(\varphi_{37} \to \forall N\,(\varphi_{37})) \to \varphi_{35}$}
\BIC{$\varphi_{35}$}
\AXC{$\varphi_{35} \to \varphi_{36}$\lbl{A1}}
\BIC{$\varphi_{36}$}
\UIC{$\forall \varepsilon\,(\varphi_{36})$}
\end{ptree}
\noindent$\Pi_{14}$:
\begin{ptree}
\AXC{$\varphi_{23}$\lbl{A4}}
\AXC{$\forall \varepsilon\,(\varphi_{36}) \to \varphi_{4} \to \varphi_{17}$\lbl{A3}}
\AXC{$\stackrel{\Pi_{13}}{\vdots}$}\noLine
\UIC{$\forall \varepsilon\,(\varphi_{36})$}
\AXC{$\forall \varepsilon\,(\varphi_{36}) \to (\forall \varepsilon\,(\varphi_{36}) \to \varphi_{4} \to \varphi_{17}) \to \varphi_{23} \to \varphi_{38}$\lbl{A1}}
\BIC{$(\forall \varepsilon\,(\varphi_{36}) \to \varphi_{4} \to \varphi_{17}) \to \varphi_{23} \to \varphi_{38}$}
\BIC{$\varphi_{23} \to \varphi_{38}$}
\BIC{$\varphi_{38}$}
\end{ptree}
\noindent$\Pi_{15}$:
\begin{ptree}
\AXC{$\forall a\,\forall \alpha\,(\mathrm{Conv}(a,\alpha) \to \alpha \in \mathbb{R})$\lbl{[C2]}}
\AXC{$\varphi_{5}$\lbl{A2}}
\BIC{$\forall \alpha\,(\mathrm{Conv}(a,\alpha) \to \alpha \in \mathbb{R})$}
\AXC{$\varphi_{1}$\lbl{A2}}
\BIC{$\mathrm{Conv}(a,\alpha) \to \alpha \in \mathbb{R}$}
\AXC{$(\mathrm{Conv}(a,\alpha) \to \alpha \in \mathbb{R}) \to (\mathrm{Seq}(b) \to \varphi_{2}) \to \varphi_{8}$\lbl{A1}}
\BIC{$(\mathrm{Seq}(b) \to \varphi_{2}) \to \varphi_{8}$}
\end{ptree}
\noindent$\Pi_{16}$:
\begin{ptree}
\AXC{$\varphi_{30}$\lbl{DefConv}}
\AXC{$\varphi_{30} \to \forall \alpha\,(\mathrm{Conv}(b,\alpha) \leftrightarrow \mathrm{Seq}(b) \land \alpha \in \mathbb{R} \land \varphi_{17})$\lbl{A2}}
\BIC{$\forall \alpha\,(\mathrm{Conv}(b,\alpha) \leftrightarrow \mathrm{Seq}(b) \land \alpha \in \mathbb{R} \land \varphi_{17})$}
\AXC{$\varphi_{54}$\lbl{A2}}
\BIC{$\mathrm{Conv}(b,\alpha) \leftrightarrow \mathrm{Seq}(b) \land \alpha \in \mathbb{R} \land \varphi_{17}$}
\AXC{$\varphi_{61}$\lbl{A1}}
\BIC{$(\mathrm{Seq}(b) \to \mathrm{Seq}(b)) \to \varphi_{8} \to \varphi_{38} \to \mathrm{Seq}(b) \to \varphi_{3}$}
\end{ptree}
\noindent$\Pi_{17}$:
\begin{ptree}
\AXC{$\mathrm{Seq}(b) \to \mathrm{Seq}(b)$\lbl{A1}}
\AXC{$\stackrel{\Pi_{16}}{\vdots}$}\noLine
\UIC{$(\mathrm{Seq}(b) \to \mathrm{Seq}(b)) \to \varphi_{8} \to \varphi_{38} \to \mathrm{Seq}(b) \to \varphi_{3}$}
\BIC{$\varphi_{8} \to \varphi_{38} \to \mathrm{Seq}(b) \to \varphi_{3}$}
\end{ptree}
\noindent 主木:
\begin{ptree}
\AXC{$\stackrel{\Pi_{14}}{\vdots}$}\noLine
\UIC{$\varphi_{38}$}
\AXC{$\mathrm{Seq}(b) \to \varphi_{2}$\lbl{A1}}
\AXC{$\stackrel{\Pi_{15}}{\vdots}$}\noLine
\UIC{$(\mathrm{Seq}(b) \to \varphi_{2}) \to \varphi_{8}$}
\BIC{$\varphi_{8}$}
\AXC{$\stackrel{\Pi_{17}}{\vdots}$}\noLine
\UIC{$\varphi_{8} \to \varphi_{38} \to \mathrm{Seq}(b) \to \varphi_{3}$}
\BIC{$\varphi_{38} \to \mathrm{Seq}(b) \to \varphi_{3}$}
\BIC{$\mathrm{Seq}(b) \to \varphi_{3}$}
\UIC{$\forall \alpha\,(\mathrm{Seq}(b) \to \varphi_{3})$}
\UIC{$\forall b\,\forall \alpha\,(\mathrm{Seq}(b) \to \varphi_{3})$}
\UIC{$\forall a\,\forall b\,\forall \alpha\,(\mathrm{Seq}(b) \to \varphi_{3})$}
\end{ptree}
\begin{abbrevlist}
\item $\varphi_{1} :\equiv \forall \alpha\,(\mathrm{Conv}(a,\alpha) \to \alpha \in \mathbb{R}) \to \mathrm{Conv}(a,\alpha) \to \alpha \in \mathbb{R}$
\item $\varphi_{2} :\equiv \forall n \in \mathbb{N}\,(a_{n} = b_{n}) \to \mathrm{Conv}(a,\alpha) \to \mathrm{Conv}(a,\alpha)$
\item $\varphi_{3} :\equiv \forall n \in \mathbb{N}\,(a_{n} = b_{n}) \to \mathrm{Conv}(a,\alpha) \to \mathrm{Conv}(b,\alpha)$
\item $\varphi_{4} :\equiv \forall \varepsilon\,(\mathrm{Seq}(b) \land \forall n \in \mathbb{N}\,(a_{n} = b_{n}) \land \mathrm{Conv}(a,\alpha))$
\item $\varphi_{5} :\equiv \forall a\,\forall \alpha\,(\mathrm{Conv}(a,\alpha) \to \alpha \in \mathbb{R}) \to \forall \alpha\,(\mathrm{Conv}(a,\alpha) \to \alpha \in \mathbb{R})$
\item $\varphi_{6} :\equiv \forall N\,(\neg (N \in \mathbb{N} \land \forall n \in \mathbb{N}\,(N \leq n \to \lvert a_{n} - \alpha \rvert < \varepsilon)))$
\item $\varphi_{7} :\equiv \forall N\,(\neg (N \in \mathbb{N} \land \forall n \in \mathbb{N}\,(N \leq n \to \lvert b_{n} - \alpha \rvert < \varepsilon)))$
\item $\varphi_{8} :\equiv \mathrm{Seq}(b) \to \forall n \in \mathbb{N}\,(a_{n} = b_{n}) \to \mathrm{Conv}(a,\alpha) \to \alpha \in \mathbb{R}$
\item $\varphi_{9} :\equiv \forall n \in \mathbb{N}\,(a_{n} = b_{n}) \to \mathrm{Conv}(a,\alpha) \to \varepsilon \in \mathbb{R} \to \varepsilon \in \mathbb{R}$
\item $\varphi_{10} :\equiv \mathrm{Seq}(b) \to \forall n \in \mathbb{N}\,(a_{n} = b_{n}) \to \forall n \in \mathbb{N}\,(a_{n} = b_{n})$
\item $\varphi_{11} :\equiv 0 < \varepsilon \to N \in \mathbb{N} \land \forall n \in \mathbb{N}\,(N \leq n \to \lvert a_{n} - \alpha \rvert < \varepsilon) \to N \in \mathbb{N}$
\item $\varphi_{12} :\equiv \forall n \in \mathbb{N}\,(a_{n} = b_{n}) \to \mathrm{Conv}(a,\alpha) \to \varepsilon \in \mathbb{R} \to 0 < \varepsilon \to 0 < \varepsilon$
\item $\varphi_{13} :\equiv N \in \mathbb{N} \land \forall n \in \mathbb{N}\,(N \leq n \to \lvert a_{n} - \alpha \rvert < \varepsilon) \to n \in \mathbb{N} \to N \leq n \to N \leq n$
\item $\varphi_{14} :\equiv N \in \mathbb{N} \land \forall n \in \mathbb{N}\,(N \leq n \to \lvert a_{n} - \alpha \rvert < \varepsilon) \to n \in \mathbb{N} \to n \in \mathbb{N}$
\item $\varphi_{15} :\equiv \varepsilon \in \mathbb{R} \to 0 < \varepsilon \to \exists N \in \mathbb{N}\,\forall n \in \mathbb{N}\,(N \leq n \to \lvert a_{n} - \alpha \rvert < \varepsilon)$
\item $\varphi_{16} :\equiv \varepsilon \in \mathbb{R} \to 0 < \varepsilon \to \exists N \in \mathbb{N}\,\forall n \in \mathbb{N}\,(N \leq n \to \lvert b_{n} - \alpha \rvert < \varepsilon)$
\item $\varphi_{17} :\equiv \forall \varepsilon \in \mathbb{R}\,(0 < \varepsilon \to \exists N \in \mathbb{N}\,\forall n \in \mathbb{N}\,(N \leq n \to \lvert b_{n} - \alpha \rvert < \varepsilon))$
\item $\varphi_{18} :\equiv \forall n \in \mathbb{N}\,(N \leq n \to \lvert a_{n} - \alpha \rvert < \varepsilon) \to n \in \mathbb{N} \to N \leq n \to \lvert a_{n} - \alpha \rvert < \varepsilon$
\item $\varphi_{19} :\equiv N \in \mathbb{N} \land \forall n \in \mathbb{N}\,(N \leq n \to \lvert a_{n} - \alpha \rvert < \varepsilon) \to n \in \mathbb{N} \to N \leq n \to a_{n} = b_{n}$
\item $\varphi_{20} :\equiv N \in \mathbb{N} \land \forall n \in \mathbb{N}\,(N \leq n \to \lvert a_{n} - \alpha \rvert < \varepsilon) \to n \in \mathbb{N} \to N \leq n \to \lvert a_{n} - \alpha \rvert < \varepsilon$
\item $\varphi_{21} :\equiv N \in \mathbb{N} \land \forall n \in \mathbb{N}\,(N \leq n \to \lvert a_{n} - \alpha \rvert < \varepsilon) \to n \in \mathbb{N} \to N \leq n \to \lvert b_{n} - \alpha \rvert < \varepsilon$
\item $\varphi_{22} :\equiv N \in \mathbb{N} \land \forall n \in \mathbb{N}\,(N \leq n \to \lvert a_{n} - \alpha \rvert < \varepsilon) \to \forall n \in \mathbb{N}\,(N \leq n \to \lvert b_{n} - \alpha \rvert < \varepsilon)$
\item $\varphi_{23} :\equiv \mathrm{Seq}(b) \land \forall n \in \mathbb{N}\,(a_{n} = b_{n}) \land \mathrm{Conv}(a,\alpha) \to \varphi_{4}$
\item $\varphi_{24} :\equiv \forall z\,\forall w\,(z = w \to \lvert z - \alpha \rvert < \varepsilon \to \lvert w - \alpha \rvert < \varepsilon) \to \forall w\,(a_{n} = w \to \lvert a_{n} - \alpha \rvert < \varepsilon \to \lvert w - \alpha \rvert < \varepsilon)$
\item $\varphi_{25} :\equiv N \in \mathbb{N} \land \forall n \in \mathbb{N}\,(N \leq n \to \lvert a_{n} - \alpha \rvert < \varepsilon) \to N \in \mathbb{N} \land \forall n \in \mathbb{N}\,(N \leq n \to \lvert a_{n} - \alpha \rvert < \varepsilon)$
\item $\varphi_{26} :\equiv N \in \mathbb{N} \land \forall n \in \mathbb{N}\,(N \leq n \to \lvert a_{n} - \alpha \rvert < \varepsilon) \to N \in \mathbb{N} \land \forall n \in \mathbb{N}\,(N \leq n \to \lvert b_{n} - \alpha \rvert < \varepsilon)$
\item $\varphi_{27} :\equiv N \in \mathbb{N} \land \forall n \in \mathbb{N}\,(N \leq n \to \lvert a_{n} - \alpha \rvert < \varepsilon) \to n \in \mathbb{N} \to N \leq n \to n \in \mathbb{N} \to N \leq n \to \lvert a_{n} - \alpha \rvert < \varepsilon$
\item $\varphi_{28} :\equiv N \in \mathbb{N} \land \forall n \in \mathbb{N}\,(N \leq n \to \lvert a_{n} - \alpha \rvert < \varepsilon) \to \exists N \in \mathbb{N}\,\forall n \in \mathbb{N}\,(N \leq n \to \lvert b_{n} - \alpha \rvert < \varepsilon)$
\item $\varphi_{29} :\equiv N \in \mathbb{N} \land \forall n \in \mathbb{N}\,(N \leq n \to \lvert a_{n} - \alpha \rvert < \varepsilon) \to n \in \mathbb{N} \to N \leq n \to \forall n \in \mathbb{N}\,(N \leq n \to \lvert a_{n} - \alpha \rvert < \varepsilon)$
\item $\varphi_{30} :\equiv \forall a\,\forall \alpha\,(\mathrm{Conv}(a,\alpha) \leftrightarrow \mathrm{Seq}(a) \land \alpha \in \mathbb{R} \land \forall \varepsilon \in \mathbb{R}\,(0 < \varepsilon \to \exists N \in \mathbb{N}\,\forall n \in \mathbb{N}\,(N \leq n \to \lvert a_{n} - \alpha \rvert < \varepsilon)))$
\item $\varphi_{31} :\equiv (\forall n \in \mathbb{N}\,(a_{n} = b_{n}) \to n \in \mathbb{N} \to a_{n} = b_{n}) \to \mathrm{Seq}(b) \to \forall n \in \mathbb{N}\,(a_{n} = b_{n}) \to n \in \mathbb{N} \to a_{n} = b_{n}$
\item $\varphi_{32} :\equiv \forall w\,(a_{n} = w \to \lvert a_{n} - \alpha \rvert < \varepsilon \to \lvert w - \alpha \rvert < \varepsilon) \to a_{n} = b_{n} \to \lvert a_{n} - \alpha \rvert < \varepsilon \to \lvert b_{n} - \alpha \rvert < \varepsilon$
\item $\varphi_{33} :\equiv \mathrm{Seq}(b) \land \forall n \in \mathbb{N}\,(a_{n} = b_{n}) \land \mathrm{Conv}(a,\alpha) \land \varepsilon \in \mathbb{R} \land 0 < \varepsilon \land (N \in \mathbb{N} \land \forall n \in \mathbb{N}\,(N \leq n \to \lvert a_{n} - \alpha \rvert < \varepsilon))$
\item $\varphi_{34} :\equiv \mathrm{Seq}(b) \to \forall n \in \mathbb{N}\,(a_{n} = b_{n}) \to \mathrm{Conv}(a,\alpha) \to \varphi_{15}$
\item $\varphi_{35} :\equiv \mathrm{Seq}(b) \to \forall n \in \mathbb{N}\,(a_{n} = b_{n}) \to \mathrm{Conv}(a,\alpha) \to \varphi_{16}$
\item $\varphi_{36} :\equiv \mathrm{Seq}(b) \land \forall n \in \mathbb{N}\,(a_{n} = b_{n}) \land \mathrm{Conv}(a,\alpha) \to \varphi_{16}$
\item $\varphi_{37} :\equiv \mathrm{Seq}(b) \land \forall n \in \mathbb{N}\,(a_{n} = b_{n}) \land \mathrm{Conv}(a,\alpha) \land \varepsilon \in \mathbb{R} \land 0 < \varepsilon \land \neg \exists N \in \mathbb{N}\,\forall n \in \mathbb{N}\,(N \leq n \to \lvert b_{n} - \alpha \rvert < \varepsilon)$
\item $\varphi_{38} :\equiv \mathrm{Seq}(b) \to \forall n \in \mathbb{N}\,(a_{n} = b_{n}) \to \mathrm{Conv}(a,\alpha) \to \varphi_{17}$
\item $\varphi_{39} :\equiv \mathrm{Seq}(b) \to \forall n \in \mathbb{N}\,(a_{n} = b_{n}) \to \mathrm{Conv}(a,\alpha) \to \varepsilon \in \mathbb{R} \to \varphi_{11}$
\item $\varphi_{40} :\equiv \forall n \in \mathbb{N}\,(a_{n} = b_{n}) \to \mathrm{Conv}(a,\alpha) \to \varepsilon \in \mathbb{R} \to 0 < \varepsilon \to \varphi_{13}$
\item $\varphi_{41} :\equiv \forall n \in \mathbb{N}\,(a_{n} = b_{n}) \to \mathrm{Conv}(a,\alpha) \to \varepsilon \in \mathbb{R} \to 0 < \varepsilon \to \varphi_{14}$
\item $\varphi_{42} :\equiv \forall n \in \mathbb{N}\,(a_{n} = b_{n}) \to \mathrm{Conv}(a,\alpha) \to \varepsilon \in \mathbb{R} \to 0 < \varepsilon \to \varphi_{19}$
\item $\varphi_{43} :\equiv \forall n \in \mathbb{N}\,(a_{n} = b_{n}) \to \mathrm{Conv}(a,\alpha) \to \varepsilon \in \mathbb{R} \to 0 < \varepsilon \to \varphi_{20}$
\item $\varphi_{44} :\equiv \forall n \in \mathbb{N}\,(a_{n} = b_{n}) \to \mathrm{Conv}(a,\alpha) \to \varepsilon \in \mathbb{R} \to 0 < \varepsilon \to \varphi_{21}$
\item $\varphi_{45} :\equiv \forall n \in \mathbb{N}\,(a_{n} = b_{n}) \to \mathrm{Conv}(a,\alpha) \to \varepsilon \in \mathbb{R} \to 0 < \varepsilon \to \varphi_{22}$
\item $\varphi_{46} :\equiv \forall n \in \mathbb{N}\,(a_{n} = b_{n}) \to \mathrm{Conv}(a,\alpha) \to \varepsilon \in \mathbb{R} \to 0 < \varepsilon \to \varphi_{25}$
\item $\varphi_{47} :\equiv \forall n \in \mathbb{N}\,(a_{n} = b_{n}) \to \mathrm{Conv}(a,\alpha) \to \varepsilon \in \mathbb{R} \to 0 < \varepsilon \to \varphi_{26}$
\item $\varphi_{48} :\equiv \forall n \in \mathbb{N}\,(a_{n} = b_{n}) \to \mathrm{Conv}(a,\alpha) \to \varepsilon \in \mathbb{R} \to 0 < \varepsilon \to \varphi_{27}$
\item $\varphi_{49} :\equiv \forall n \in \mathbb{N}\,(a_{n} = b_{n}) \to \mathrm{Conv}(a,\alpha) \to \varepsilon \in \mathbb{R} \to 0 < \varepsilon \to \varphi_{28}$
\item $\varphi_{50} :\equiv \forall n \in \mathbb{N}\,(a_{n} = b_{n}) \to \mathrm{Conv}(a,\alpha) \to \varepsilon \in \mathbb{R} \to 0 < \varepsilon \to \varphi_{29}$
\item $\varphi_{51} :\equiv \forall N\,(\varphi_{37} \to \neg (N \in \mathbb{N} \land \forall n \in \mathbb{N}\,(N \leq n \to \lvert a_{n} - \alpha \rvert < \varepsilon)))$
\item $\varphi_{52} :\equiv \forall \alpha\,\forall \varepsilon\,(\mathrm{Conv}(a,\alpha) \to \varphi_{15}) \to \forall \varepsilon\,(\mathrm{Conv}(a,\alpha) \to \varphi_{15})$
\item $\varphi_{53} :\equiv \forall a\,\forall \alpha\,\forall \varepsilon\,(\mathrm{Conv}(a,\alpha) \to \varphi_{15}) \to \forall \alpha\,\forall \varepsilon\,(\mathrm{Conv}(a,\alpha) \to \varphi_{15})$
\item $\varphi_{54} :\equiv \forall \alpha\,(\mathrm{Conv}(b,\alpha) \leftrightarrow \mathrm{Seq}(b) \land \alpha \in \mathbb{R} \land \varphi_{17}) \to (\mathrm{Conv}(b,\alpha) \leftrightarrow \mathrm{Seq}(b) \land \alpha \in \mathbb{R} \land \varphi_{17})$
\item $\varphi_{55} :\equiv (a_{n} = b_{n} \to \lvert a_{n} - \alpha \rvert < \varepsilon \to \lvert b_{n} - \alpha \rvert < \varepsilon) \to \mathrm{Seq}(b) \to \varphi_{44}$
\item $\varphi_{56} :\equiv (\varphi_{7} \to \neg (N \in \mathbb{N} \land \forall n \in \mathbb{N}\,(N \leq n \to \lvert b_{n} - \alpha \rvert < \varepsilon))) \to \mathrm{Seq}(b) \to \varphi_{49}$
\item $\varphi_{57} :\equiv \forall n\,(\varphi_{33} \to n \in \mathbb{N} \to N \leq n \to \lvert b_{n} - \alpha \rvert < \varepsilon) \to \forall n\,(\varphi_{33}) \to \forall n \in \mathbb{N}\,(N \leq n \to \lvert b_{n} - \alpha \rvert < \varepsilon)$
\item $\varphi_{58} :\equiv (\mathrm{Seq}(b) \to \varphi_{2}) \to (\mathrm{Seq}(b) \to \varphi_{9}) \to (\mathrm{Seq}(b) \to \varphi_{12}) \to \varphi_{34}$
\item $\varphi_{59} :\equiv (\mathrm{Seq}(b) \to \varphi_{49}) \to \varphi_{37} \to \neg (N \in \mathbb{N} \land \forall n \in \mathbb{N}\,(N \leq n \to \lvert a_{n} - \alpha \rvert < \varepsilon))$
\item $\varphi_{60} :\equiv (\mathrm{Seq}(b) \to \forall n \in \mathbb{N}\,(a_{n} = b_{n}) \to n \in \mathbb{N} \to a_{n} = b_{n}) \to (\mathrm{Seq}(b) \to \varphi_{41}) \to \mathrm{Seq}(b) \to \varphi_{42}$
\item $\varphi_{61} :\equiv (\mathrm{Conv}(b,\alpha) \leftrightarrow \mathrm{Seq}(b) \land \alpha \in \mathbb{R} \land \varphi_{17}) \to (\mathrm{Seq}(b) \to \mathrm{Seq}(b)) \to \varphi_{8} \to \varphi_{38} \to \mathrm{Seq}(b) \to \varphi_{3}$
\item $\varphi_{62} :\equiv (\mathrm{Seq}(b) \to \varphi_{48}) \to (\mathrm{Seq}(b) \to \varphi_{41}) \to (\mathrm{Seq}(b) \to \varphi_{40}) \to \mathrm{Seq}(b) \to \varphi_{43}$
\item $\varphi_{63} :\equiv \forall n\,(\varphi_{33} \to n \in \mathbb{N} \to N \leq n \to \lvert b_{n} - \alpha \rvert < \varepsilon) \to \varphi_{57} \to (\varphi_{33} \to \forall n\,(\varphi_{33})) \to \mathrm{Seq}(b) \to \varphi_{45}$
\end{abbrevlist}
\end{fullproof}

% 自動生成 (prover/build.js): FinBound  ノード数 636
\begin{fullproof}{[C14]}{$\forall a\,(\mathrm{Seq}(a) \to \forall N \in \mathbb{N}\,\exists M\,(M \in \mathbb{R} \land 0 \leq M \land \forall n \in \mathbb{N}\,(n \leq N \to \lvert a_{n} \rvert \leq M)))$}{636}\phantomsection\label{tree:FinBound}
\noindent$\Pi_{1}$:
\begin{ptree}
\AXC{$\forall N\,(N \in B_{8} \leftrightarrow N \in \omega \land \varphi_{51}) \to (N \in B_{8} \leftrightarrow N \in \omega \land \varphi_{51})$\lbl{A2}}
\AXC{$\varphi_{175}$\lbl{A1}}
\AXC{$\varphi_{175} \to \varphi_{248}$\lbl{A1}}
\BIC{$\varphi_{248}$}
\BIC{$\varphi_{171}$}
\end{ptree}
\noindent$\Pi_{2}$:
\begin{ptree}
\AXC{$\varphi_{19}$\lbl{A4}}
\AXC{$\varphi_{12} \to \varphi_{22}$\lbl{A3}}
\AXC{$\forall N \in B_{8}\,(S(N) \in B_{8}) \to y \in B_{8} \to S(y) \in B_{8}$\lbl{A2}}
\UIC{$\varphi_{12}$}
\AXC{$\varphi_{12} \to (\varphi_{12} \to \varphi_{22}) \to \varphi_{119}$\lbl{A1}}
\BIC{$(\varphi_{12} \to \varphi_{22}) \to \varphi_{119}$}
\BIC{$\varphi_{119}$}
\BIC{$\forall N \in B_{8}\,(S(N) \in B_{8}) \to \forall y \in B_{8}\,(S(y) \in B_{8})$}
\end{ptree}
\noindent$\Pi_{3}$:
\begin{ptree}
\AXC{$\varphi_{20}$\lbl{A4}}
\AXC{$\varphi_{13} \to \varphi_{23}$\lbl{A3}}
\AXC{$\forall y \in B_{8}\,(S(y) \in B_{8}) \to N \in B_{8} \to S(N) \in B_{8}$\lbl{A2}}
\UIC{$\varphi_{13}$}
\AXC{$\varphi_{13} \to (\varphi_{13} \to \varphi_{23}) \to \varphi_{120}$\lbl{A1}}
\BIC{$(\varphi_{13} \to \varphi_{23}) \to \varphi_{120}$}
\BIC{$\varphi_{120}$}
\BIC{$\forall y \in B_{8}\,(S(y) \in B_{8}) \to \forall N \in B_{8}\,(S(N) \in B_{8})$}
\end{ptree}
\noindent$\Pi_{4}$:
\begin{ptree}
\AXC{$\varphi_{132}$\lbl{A1}}
\AXC{$\forall N\,(N \in B_{8} \leftrightarrow N \in \omega \land \varphi_{51}) \to (N \in B_{8} \leftrightarrow N \in \omega \land \varphi_{51})$\lbl{A2}}
\AXC{$\varphi_{175}$\lbl{A1}}
\AXC{$\varphi_{175} \to \varphi_{248}$\lbl{A1}}
\BIC{$\varphi_{248}$}
\BIC{$\varphi_{171}$}
\AXC{$\varphi_{171} \to \varphi_{132} \to \varphi_{172}$\lbl{A1}}
\BIC{$\varphi_{132} \to \varphi_{172}$}
\BIC{$\varphi_{172}$}
\end{ptree}
\noindent$\Pi_{5}$:
\begin{ptree}
\AXC{$\stackrel{\Pi_{4}}{\vdots}$}\noLine
\UIC{$\varphi_{172}$}
\AXC{$\varphi_{242}$\lbl{A1}}
\BIC{$\mathrm{Seq}(a) \to \forall N\,(N \in B_{8} \leftrightarrow N \in \omega \land \varphi_{51}) \to N \in B_{8} \to N \in \omega$}
\end{ptree}
\noindent$\Pi_{6}$:
\begin{ptree}
\AXC{$\forall n\,(n \in \mathbb{N} \leftrightarrow n \in \omega \land n \neq \emptyset)$\lbl{Def$\mathbb{N}$}}
\AXC{$\varphi_{70}$\lbl{A2}}
\BIC{$N \in \mathbb{N} \leftrightarrow N \in \omega \land N \neq \emptyset$}
\AXC{$(N \in \mathbb{N} \leftrightarrow N \in \omega \land N \neq \emptyset) \to \varphi_{254}$\lbl{A1}}
\BIC{$\varphi_{254}$}
\end{ptree}
\noindent$\Pi_{7}$:
\begin{ptree}
\AXC{$\mathrm{Seq}(a) \to \varphi_{155}$\lbl{A1}}
\AXC{$\stackrel{\Pi_{5}}{\vdots}$}\noLine
\UIC{$\mathrm{Seq}(a) \to \forall N\,(N \in B_{8} \leftrightarrow N \in \omega \land \varphi_{51}) \to N \in B_{8} \to N \in \omega$}
\AXC{$\stackrel{\Pi_{6}}{\vdots}$}\noLine
\UIC{$\varphi_{254}$}
\BIC{$(\mathrm{Seq}(a) \to \varphi_{155}) \to \mathrm{Seq}(a) \to \varphi_{153}$}
\BIC{$\mathrm{Seq}(a) \to \varphi_{153}$}
\end{ptree}
\noindent$\Pi_{8}$:
\begin{ptree}
\AXC{$\stackrel{\Pi_{7}}{\vdots}$}\noLine
\UIC{$\mathrm{Seq}(a) \to \varphi_{153}$}
\AXC{$\forall n \in \mathbb{N}\,(S(n) \in \mathbb{N})$\lbl{[N2]}}
\AXC{$\forall n \in \mathbb{N}\,(S(n) \in \mathbb{N}) \to N \in \mathbb{N} \to S(N) \in \mathbb{N}$\lbl{A2}}
\BIC{$N \in \mathbb{N} \to S(N) \in \mathbb{N}$}
\AXC{$\varphi_{249}$\lbl{A1}}
\BIC{$(\mathrm{Seq}(a) \to \varphi_{153}) \to \mathrm{Seq}(a) \to \varphi_{182}$}
\BIC{$\mathrm{Seq}(a) \to \varphi_{182}$}
\end{ptree}
\noindent$\Pi_{9}$:
\begin{ptree}
\AXC{$\forall a\,\forall n\,(\mathrm{Seq}(a) \to n \in \mathbb{N} \to a_{n} \in \mathbb{R})$\lbl{[S1]}}
\AXC{$\varphi_{85}$\lbl{A2}}
\BIC{$\forall n\,(\mathrm{Seq}(a) \to n \in \mathbb{N} \to a_{n} \in \mathbb{R})$}
\AXC{$\varphi_{86}$\lbl{A2}}
\BIC{$\mathrm{Seq}(a) \to S(N) \in \mathbb{N} \to a_{S(N)} \in \mathbb{R}$}
\AXC{$(\mathrm{Seq}(a) \to S(N) \in \mathbb{N} \to a_{S(N)} \in \mathbb{R}) \to \varphi_{266}$\lbl{A1}}
\BIC{$\varphi_{266}$}
\end{ptree}
\noindent$\Pi_{10}$:
\begin{ptree}
\AXC{$\stackrel{\Pi_{8}}{\vdots}$}\noLine
\UIC{$\mathrm{Seq}(a) \to \varphi_{182}$}
\AXC{$\mathrm{Seq}(a) \to \mathrm{Seq}(a)$\lbl{A1}}
\AXC{$\stackrel{\Pi_{9}}{\vdots}$}\noLine
\UIC{$\varphi_{266}$}
\BIC{$(\mathrm{Seq}(a) \to \varphi_{182}) \to \mathrm{Seq}(a) \to \varphi_{209}$}
\BIC{$\mathrm{Seq}(a) \to \varphi_{209}$}
\AXC{$\forall x \in \mathbb{R}\,(\lvert x \rvert \in \mathbb{R})$\lbl{[A5]}}
\AXC{$\varphi_{72}$\lbl{A2}}
\BIC{$a_{S(N)} \in \mathbb{R} \to \lvert a_{S(N)} \rvert \in \mathbb{R}$}
\AXC{$\varphi_{276}$\lbl{A1}}
\BIC{$(\mathrm{Seq}(a) \to \varphi_{209}) \to \mathrm{Seq}(a) \to \varphi_{213}$}
\BIC{$\mathrm{Seq}(a) \to \varphi_{213}$}
\end{ptree}
\noindent$\Pi_{11}$:
\begin{ptree}
\AXC{$\forall x\,\forall y\,(x \in \mathbb{R} \to y \in \mathbb{R} \to 0 \leq y \to x \leq y + x)$\lbl{[E2]}}
\AXC{$\forall x\,\forall y\,(x \in \mathbb{R} \to y \in \mathbb{R} \to 0 \leq y \to x \leq y + x) \to \varphi_{87}$\lbl{A2}}
\BIC{$\varphi_{87}$}
\AXC{$\varphi_{87} \to \varphi_{79}$\lbl{A2}}
\BIC{$\varphi_{79}$}
\end{ptree}
\noindent$\Pi_{12}$:
\begin{ptree}
\AXC{$\stackrel{\Pi_{10}}{\vdots}$}\noLine
\UIC{$\mathrm{Seq}(a) \to \varphi_{213}$}
\AXC{$\stackrel{\Pi_{11}}{\vdots}$}\noLine
\UIC{$\varphi_{79}$}
\AXC{$\varphi_{79} \to (\mathrm{Seq}(a) \to \varphi_{213}) \to (\mathrm{Seq}(a) \to \varphi_{180}) \to \varphi_{181} \to \varphi_{228}$\lbl{A1}}
\BIC{$(\mathrm{Seq}(a) \to \varphi_{213}) \to (\mathrm{Seq}(a) \to \varphi_{180}) \to \varphi_{181} \to \varphi_{228}$}
\BIC{$(\mathrm{Seq}(a) \to \varphi_{180}) \to \varphi_{181} \to \varphi_{228}$}
\end{ptree}
\noindent$\Pi_{13}$:
\begin{ptree}
\AXC{$\mathrm{Seq}(a) \to \varphi_{212}$\lbl{A1}}
\AXC{$(\mathrm{Seq}(a) \to \varphi_{212}) \to \mathrm{Seq}(a) \to \varphi_{180}$\lbl{A1}}
\BIC{$\mathrm{Seq}(a) \to \varphi_{180}$}
\AXC{$\stackrel{\Pi_{12}}{\vdots}$}\noLine
\UIC{$(\mathrm{Seq}(a) \to \varphi_{180}) \to \varphi_{181} \to \varphi_{228}$}
\BIC{$\varphi_{181} \to \varphi_{228}$}
\end{ptree}
\noindent$\Pi_{14}$:
\begin{ptree}
\AXC{$\mathrm{Seq}(a) \to \varphi_{212}$\lbl{A1}}
\AXC{$(\mathrm{Seq}(a) \to \varphi_{212}) \to \varphi_{181}$\lbl{A1}}
\BIC{$\varphi_{181}$}
\AXC{$\stackrel{\Pi_{13}}{\vdots}$}\noLine
\UIC{$\varphi_{181} \to \varphi_{228}$}
\BIC{$\varphi_{228}$}
\end{ptree}
\noindent$\Pi_{15}$:
\begin{ptree}
\AXC{$x = x$\lbl{A5}}
\AXC{$x = y \to x = x \to y = x$\lbl{A6}}
\AXC{$(x = y \to x = x \to y = x) \to x = x \to x = y \to y = x$\lbl{A1}}
\BIC{$x = x \to x = y \to y = x$}
\BIC{$x = y \to y = x$}
\UIC{$\forall y\,(x = y \to y = x)$}
\UIC{$\forall x\,\forall y\,(x = y \to y = x)$}
\end{ptree}
\noindent$\Pi_{16}$:
\begin{ptree}
\AXC{$\stackrel{\Pi_{15}}{\vdots}$}\noLine
\UIC{$\forall x\,\forall y\,(x = y \to y = x)$}
\AXC{$\forall x\,\forall y\,(x = y \to y = x) \to \forall y\,(n = y \to y = n)$\lbl{A2}}
\BIC{$\forall y\,(n = y \to y = n)$}
\end{ptree}
\noindent$\Pi_{17}$:
\begin{ptree}
\AXC{$\stackrel{\Pi_{16}}{\vdots}$}\noLine
\UIC{$\forall y\,(n = y \to y = n)$}
\AXC{$\forall y\,(n = y \to y = n) \to n = S(N) \to S(N) = n$\lbl{A2}}
\BIC{$n = S(N) \to S(N) = n$}
\AXC{$\mathrm{Seq}(a) \to \varphi_{204}$\lbl{A1}}
\AXC{$\varphi_{268}$\lbl{A1}}
\BIC{$(n = S(N) \to S(N) = n) \to \mathrm{Seq}(a) \to \varphi_{205}$}
\BIC{$\mathrm{Seq}(a) \to \varphi_{205}$}
\end{ptree}
\noindent$\Pi_{18}$:
\begin{ptree}
\AXC{$\stackrel{\Pi_{14}}{\vdots}$}\noLine
\UIC{$\varphi_{228}$}
\AXC{$\stackrel{\Pi_{17}}{\vdots}$}\noLine
\UIC{$\mathrm{Seq}(a) \to \varphi_{205}$}
\AXC{$(\mathrm{Seq}(a) \to \varphi_{205}) \to \varphi_{228} \to \varphi_{97} \to \varphi_{77} \to \varphi_{227}$\lbl{A1}}
\BIC{$\varphi_{228} \to \varphi_{97} \to \varphi_{77} \to \varphi_{227}$}
\BIC{$\varphi_{97} \to \varphi_{77} \to \varphi_{227}$}
\end{ptree}
\noindent$\Pi_{19}$:
\begin{ptree}
\AXC{$\varphi_{78}$\lbl{A6}}
\UIC{$\forall w\,(\varphi_{78})$}
\UIC{$\forall z\,\forall w\,(\varphi_{78})$}
\AXC{$\forall z\,\forall w\,(\varphi_{78}) \to \varphi_{100}$\lbl{A2}}
\BIC{$\varphi_{100}$}
\AXC{$\varphi_{100} \to \varphi_{97}$\lbl{A2}}
\BIC{$\varphi_{97}$}
\AXC{$\stackrel{\Pi_{18}}{\vdots}$}\noLine
\UIC{$\varphi_{97} \to \varphi_{77} \to \varphi_{227}$}
\BIC{$\varphi_{77} \to \varphi_{227}$}
\end{ptree}
\noindent$\Pi_{20}$:
\begin{ptree}
\AXC{$\forall a\,\forall n\,(\mathrm{Seq}(a) \to n \in \mathbb{N} \to a_{n} \in \mathbb{R})$\lbl{[S1]}}
\AXC{$\varphi_{85}$\lbl{A2}}
\BIC{$\forall n\,(\mathrm{Seq}(a) \to n \in \mathbb{N} \to a_{n} \in \mathbb{R})$}
\AXC{$\varphi_{86}$\lbl{A2}}
\BIC{$\mathrm{Seq}(a) \to S(N) \in \mathbb{N} \to a_{S(N)} \in \mathbb{R}$}
\AXC{$(\mathrm{Seq}(a) \to S(N) \in \mathbb{N} \to a_{S(N)} \in \mathbb{R}) \to \varphi_{260}$\lbl{A1}}
\BIC{$\varphi_{260}$}
\end{ptree}
\noindent$\Pi_{21}$:
\begin{ptree}
\AXC{$\stackrel{\Pi_{8}}{\vdots}$}\noLine
\UIC{$\mathrm{Seq}(a) \to \varphi_{182}$}
\AXC{$\mathrm{Seq}(a) \to \mathrm{Seq}(a)$\lbl{A1}}
\AXC{$\stackrel{\Pi_{20}}{\vdots}$}\noLine
\UIC{$\varphi_{260}$}
\BIC{$(\mathrm{Seq}(a) \to \varphi_{182}) \to \mathrm{Seq}(a) \to \varphi_{189}$}
\BIC{$\mathrm{Seq}(a) \to \varphi_{189}$}
\end{ptree}
\noindent$\Pi_{22}$:
\begin{ptree}
\AXC{$\stackrel{\Pi_{21}}{\vdots}$}\noLine
\UIC{$\mathrm{Seq}(a) \to \varphi_{189}$}
\AXC{$\forall x \in \mathbb{R}\,(0 \leq \lvert x \rvert)$\lbl{[A6]}}
\AXC{$\forall x \in \mathbb{R}\,(0 \leq \lvert x \rvert) \to a_{S(N)} \in \mathbb{R} \to 0 \leq \lvert a_{S(N)} \rvert$\lbl{A2}}
\BIC{$a_{S(N)} \in \mathbb{R} \to 0 \leq \lvert a_{S(N)} \rvert$}
\AXC{$\varphi_{267}$\lbl{A1}}
\BIC{$(\mathrm{Seq}(a) \to \varphi_{189}) \to \mathrm{Seq}(a) \to \varphi_{186}$}
\BIC{$\mathrm{Seq}(a) \to \varphi_{186}$}
\end{ptree}
\noindent$\Pi_{23}$:
\begin{ptree}
\AXC{$\forall a\,\forall n\,(\mathrm{Seq}(a) \to n \in \mathbb{N} \to a_{n} \in \mathbb{R})$\lbl{[S1]}}
\AXC{$\varphi_{85}$\lbl{A2}}
\BIC{$\forall n\,(\mathrm{Seq}(a) \to n \in \mathbb{N} \to a_{n} \in \mathbb{R})$}
\AXC{$\varphi_{86}$\lbl{A2}}
\BIC{$\mathrm{Seq}(a) \to S(N) \in \mathbb{N} \to a_{S(N)} \in \mathbb{R}$}
\AXC{$(\mathrm{Seq}(a) \to S(N) \in \mathbb{N} \to a_{S(N)} \in \mathbb{R}) \to \varphi_{265}$\lbl{A1}}
\BIC{$\varphi_{265}$}
\end{ptree}
\noindent$\Pi_{24}$:
\begin{ptree}
\AXC{$\stackrel{\Pi_{8}}{\vdots}$}\noLine
\UIC{$\mathrm{Seq}(a) \to \varphi_{182}$}
\AXC{$\mathrm{Seq}(a) \to \mathrm{Seq}(a)$\lbl{A1}}
\AXC{$\stackrel{\Pi_{23}}{\vdots}$}\noLine
\UIC{$\varphi_{265}$}
\BIC{$(\mathrm{Seq}(a) \to \varphi_{182}) \to \mathrm{Seq}(a) \to \varphi_{208}$}
\BIC{$\mathrm{Seq}(a) \to \varphi_{208}$}
\AXC{$\forall x \in \mathbb{R}\,(\lvert x \rvert \in \mathbb{R})$\lbl{[A5]}}
\AXC{$\varphi_{72}$\lbl{A2}}
\BIC{$a_{S(N)} \in \mathbb{R} \to \lvert a_{S(N)} \rvert \in \mathbb{R}$}
\AXC{$\varphi_{275}$\lbl{A1}}
\BIC{$(\mathrm{Seq}(a) \to \varphi_{208}) \to \mathrm{Seq}(a) \to \varphi_{211}$}
\BIC{$\mathrm{Seq}(a) \to \varphi_{211}$}
\end{ptree}
\noindent$\Pi_{25}$:
\begin{ptree}
\AXC{$\forall x\,\forall y\,(x \in \mathbb{R} \to y \in \mathbb{R} \to 0 \leq y \to x \leq x + y)$\lbl{[E1]}}
\AXC{$\varphi_{99}$\lbl{A2}}
\BIC{$\forall y\,(M \in \mathbb{R} \to y \in \mathbb{R} \to 0 \leq y \to M \leq M + y)$}
\AXC{$\forall y\,(M \in \mathbb{R} \to y \in \mathbb{R} \to 0 \leq y \to M \leq M + y) \to \varphi_{80}$\lbl{A2}}
\BIC{$\varphi_{80}$}
\end{ptree}
\noindent$\Pi_{26}$:
\begin{ptree}
\AXC{$\stackrel{\Pi_{24}}{\vdots}$}\noLine
\UIC{$\mathrm{Seq}(a) \to \varphi_{211}$}
\AXC{$\mathrm{Seq}(a) \to \varphi_{212}$\lbl{A1}}
\AXC{$(\mathrm{Seq}(a) \to \varphi_{212}) \to \mathrm{Seq}(a) \to \varphi_{180}$\lbl{A1}}
\BIC{$\mathrm{Seq}(a) \to \varphi_{180}$}
\AXC{$\stackrel{\Pi_{25}}{\vdots}$}\noLine
\UIC{$\varphi_{80}$}
\AXC{$\varphi_{80} \to \varphi_{289}$\lbl{A1}}
\BIC{$\varphi_{289}$}
\BIC{$(\mathrm{Seq}(a) \to \varphi_{211}) \to (\mathrm{Seq}(a) \to \varphi_{186}) \to \mathrm{Seq}(a) \to \varphi_{210}$}
\BIC{$(\mathrm{Seq}(a) \to \varphi_{186}) \to \mathrm{Seq}(a) \to \varphi_{210}$}
\end{ptree}
\noindent$\Pi_{27}$:
\begin{ptree}
\AXC{$\mathrm{Seq}(a) \to \varphi_{191}$\lbl{A1}}
\AXC{$\varphi_{54}$\lbl{A2}}
\AXC{$\mathrm{Seq}(a) \to \varphi_{212}$\lbl{A1}}
\AXC{$(\mathrm{Seq}(a) \to \varphi_{212}) \to \varphi_{225}$\lbl{A1}}
\BIC{$\varphi_{225}$}
\AXC{$\varphi_{225} \to \varphi_{54} \to \varphi_{221}$\lbl{A1}}
\BIC{$\varphi_{54} \to \varphi_{221}$}
\BIC{$\varphi_{221}$}
\AXC{$\varphi_{290}$\lbl{A1}}
\BIC{$(\mathrm{Seq}(a) \to \varphi_{191}) \to (\mathrm{Seq}(a) \to \varphi_{197}) \to \mathrm{Seq}(a) \to \varphi_{206}$}
\BIC{$(\mathrm{Seq}(a) \to \varphi_{197}) \to \mathrm{Seq}(a) \to \varphi_{206}$}
\end{ptree}
\noindent$\Pi_{28}$:
\begin{ptree}
\AXC{$\forall x\,\forall y\,\forall z\,(x \leq y \to y \leq z \to x \leq z)$\lbl{[O5]}}
\AXC{$\varphi_{55}$\lbl{A2}}
\BIC{$\forall y\,\forall z\,(\lvert a_{n} \rvert \leq y \to y \leq z \to \lvert a_{n} \rvert \leq z)$}
\AXC{$\varphi_{93}$\lbl{A2}}
\BIC{$\forall z\,(\lvert a_{n} \rvert \leq M \to M \leq z \to \lvert a_{n} \rvert \leq z)$}
\AXC{$\forall z\,(\lvert a_{n} \rvert \leq M \to M \leq z \to \lvert a_{n} \rvert \leq z) \to \varphi_{68}$\lbl{A2}}
\BIC{$\varphi_{68}$}
\end{ptree}
\noindent$\Pi_{29}$:
\begin{ptree}
\AXC{$\mathrm{Seq}(a) \to \varphi_{197}$\lbl{A1}}
\AXC{$\stackrel{\Pi_{27}}{\vdots}$}\noLine
\UIC{$(\mathrm{Seq}(a) \to \varphi_{197}) \to \mathrm{Seq}(a) \to \varphi_{206}$}
\BIC{$\mathrm{Seq}(a) \to \varphi_{206}$}
\AXC{$\stackrel{\Pi_{28}}{\vdots}$}\noLine
\UIC{$\varphi_{68}$}
\AXC{$\varphi_{284}$\lbl{A1}}
\BIC{$(\mathrm{Seq}(a) \to \varphi_{206}) \to (\mathrm{Seq}(a) \to \varphi_{210}) \to \mathrm{Seq}(a) \to \varphi_{219}$}
\BIC{$(\mathrm{Seq}(a) \to \varphi_{210}) \to \mathrm{Seq}(a) \to \varphi_{219}$}
\end{ptree}
\noindent$\Pi_{30}$:
\begin{ptree}
\AXC{$\stackrel{\Pi_{22}}{\vdots}$}\noLine
\UIC{$\mathrm{Seq}(a) \to \varphi_{186}$}
\AXC{$\stackrel{\Pi_{26}}{\vdots}$}\noLine
\UIC{$(\mathrm{Seq}(a) \to \varphi_{186}) \to \mathrm{Seq}(a) \to \varphi_{210}$}
\BIC{$\mathrm{Seq}(a) \to \varphi_{210}$}
\AXC{$\stackrel{\Pi_{29}}{\vdots}$}\noLine
\UIC{$(\mathrm{Seq}(a) \to \varphi_{210}) \to \mathrm{Seq}(a) \to \varphi_{219}$}
\BIC{$\mathrm{Seq}(a) \to \varphi_{219}$}
\end{ptree}
\noindent$\Pi_{31}$:
\begin{ptree}
\AXC{$\forall n \in \mathbb{N}\,(n \in \omega)$\lbl{[N7']}}
\AXC{$\forall n \in \mathbb{N}\,(n \in \omega) \to N \in \mathbb{N} \to N \in \omega$\lbl{A2}}
\BIC{$N \in \mathbb{N} \to N \in \omega$}
\AXC{$(N \in \mathbb{N} \to N \in \omega) \to (\mathrm{Seq}(a) \to \varphi_{153}) \to \mathrm{Seq}(a) \to \varphi_{200}$\lbl{A1}}
\BIC{$(\mathrm{Seq}(a) \to \varphi_{153}) \to \mathrm{Seq}(a) \to \varphi_{200}$}
\end{ptree}
\noindent$\Pi_{32}$:
\begin{ptree}
\AXC{$\forall n \in \mathbb{N}\,(n \in \omega)$\lbl{[N7']}}
\AXC{$\forall n \in \mathbb{N}\,(n \in \omega) \to n \in \mathbb{N} \to n \in \omega$\lbl{A2}}
\BIC{$n \in \mathbb{N} \to n \in \omega$}
\AXC{$(n \in \mathbb{N} \to n \in \omega) \to (\mathrm{Seq}(a) \to \varphi_{191}) \to \mathrm{Seq}(a) \to \varphi_{201}$\lbl{A1}}
\BIC{$(\mathrm{Seq}(a) \to \varphi_{191}) \to \mathrm{Seq}(a) \to \varphi_{201}$}
\end{ptree}
\noindent$\Pi_{33}$:
\begin{ptree}
\AXC{$\varphi_{18}$\lbl{[N9]}}
\AXC{$\varphi_{18} \to \varphi_{38}$\lbl{A2}}
\BIC{$\varphi_{38}$}
\AXC{$\varphi_{38} \to v_{685} \in \omega \to \varphi_{24}$\lbl{A2}}
\BIC{$v_{685} \in \omega \to \varphi_{24}$}
\UIC{$\forall v_{686}\,(v_{685} \in \omega \to \varphi_{24})$}
\UIC{$\forall v_{685}\,\forall v_{686}\,(v_{685} \in \omega \to \varphi_{24})$}
\AXC{$\forall v_{685}\,\forall v_{686}\,(v_{685} \in \omega \to \varphi_{24}) \to \varphi_{45}$\lbl{A2}}
\BIC{$\varphi_{45}$}
\end{ptree}
\noindent$\Pi_{34}$:
\begin{ptree}
\AXC{$\stackrel{\Pi_{33}}{\vdots}$}\noLine
\UIC{$\varphi_{45}$}
\AXC{$\varphi_{45} \to n \in \omega \to N \in \omega \to n \leq S(N) \to n \leq N \lor n = S(N)$\lbl{A2}}
\BIC{$n \in \omega \to N \in \omega \to n \leq S(N) \to n \leq N \lor n = S(N)$}
\AXC{$(n \in \omega \to N \in \omega \to n \leq S(N) \to n \leq N \lor n = S(N)) \to \varphi_{288}$\lbl{A1}}
\BIC{$\varphi_{288}$}
\end{ptree}
\noindent$\Pi_{35}$:
\begin{ptree}
\AXC{$\mathrm{Seq}(a) \to \varphi_{191}$\lbl{A1}}
\AXC{$\stackrel{\Pi_{32}}{\vdots}$}\noLine
\UIC{$(\mathrm{Seq}(a) \to \varphi_{191}) \to \mathrm{Seq}(a) \to \varphi_{201}$}
\BIC{$\mathrm{Seq}(a) \to \varphi_{201}$}
\AXC{$\stackrel{\Pi_{34}}{\vdots}$}\noLine
\UIC{$\varphi_{288}$}
\BIC{$(\mathrm{Seq}(a) \to \varphi_{200}) \to (\mathrm{Seq}(a) \to \varphi_{195}) \to \mathrm{Seq}(a) \to \varphi_{202}$}
\end{ptree}
\noindent$\Pi_{36}$:
\begin{ptree}
\AXC{$\stackrel{\Pi_{7}}{\vdots}$}\noLine
\UIC{$\mathrm{Seq}(a) \to \varphi_{153}$}
\AXC{$\stackrel{\Pi_{31}}{\vdots}$}\noLine
\UIC{$(\mathrm{Seq}(a) \to \varphi_{153}) \to \mathrm{Seq}(a) \to \varphi_{200}$}
\BIC{$\mathrm{Seq}(a) \to \varphi_{200}$}
\AXC{$\stackrel{\Pi_{35}}{\vdots}$}\noLine
\UIC{$(\mathrm{Seq}(a) \to \varphi_{200}) \to (\mathrm{Seq}(a) \to \varphi_{195}) \to \mathrm{Seq}(a) \to \varphi_{202}$}
\BIC{$(\mathrm{Seq}(a) \to \varphi_{195}) \to \mathrm{Seq}(a) \to \varphi_{202}$}
\end{ptree}
\noindent$\Pi_{37}$:
\begin{ptree}
\AXC{$\stackrel{\Pi_{30}}{\vdots}$}\noLine
\UIC{$\mathrm{Seq}(a) \to \varphi_{219}$}
\AXC{$\mathrm{Seq}(a) \to \varphi_{195}$\lbl{A1}}
\AXC{$\stackrel{\Pi_{36}}{\vdots}$}\noLine
\UIC{$(\mathrm{Seq}(a) \to \varphi_{195}) \to \mathrm{Seq}(a) \to \varphi_{202}$}
\BIC{$\mathrm{Seq}(a) \to \varphi_{202}$}
\AXC{$\varphi_{291}$\lbl{A1}}
\BIC{$(\mathrm{Seq}(a) \to \varphi_{219}) \to \varphi_{227} \to \mathrm{Seq}(a) \to \varphi_{214}$}
\BIC{$\varphi_{227} \to \mathrm{Seq}(a) \to \varphi_{214}$}
\end{ptree}
\noindent$\Pi_{38}$:
\begin{ptree}
\AXC{$\varphi_{67}$\lbl{A6}}
\UIC{$\forall w\,(\varphi_{67})$}
\UIC{$\forall z\,\forall w\,(\varphi_{67})$}
\AXC{$\forall z\,\forall w\,(\varphi_{67}) \to \varphi_{84}$\lbl{A2}}
\BIC{$\varphi_{84}$}
\AXC{$\varphi_{84} \to \varphi_{77}$\lbl{A2}}
\BIC{$\varphi_{77}$}
\AXC{$\stackrel{\Pi_{19}}{\vdots}$}\noLine
\UIC{$\varphi_{77} \to \varphi_{227}$}
\BIC{$\varphi_{227}$}
\AXC{$\stackrel{\Pi_{37}}{\vdots}$}\noLine
\UIC{$\varphi_{227} \to \mathrm{Seq}(a) \to \varphi_{214}$}
\BIC{$\mathrm{Seq}(a) \to \varphi_{214}$}
\AXC{$\varphi_{272}$\lbl{A1}}
\BIC{$\varphi_{173} \to n \in \mathbb{N} \to n \leq S(N) \to \lvert a_{n} \rvert \leq M + \lvert a_{S(N)} \rvert$}
\UIC{$\varphi_{220}$}
\end{ptree}
\noindent$\Pi_{39}$:
\begin{ptree}
\AXC{$\varphi_{173} \to \forall n\,(\varphi_{173})$\lbl{A4}}
\AXC{$\varphi_{220} \to \varphi_{224}$\lbl{A3}}
\AXC{$\stackrel{\Pi_{38}}{\vdots}$}\noLine
\UIC{$\varphi_{220}$}
\AXC{$\varphi_{220} \to (\varphi_{220} \to \varphi_{224}) \to (\varphi_{173} \to \forall n\,(\varphi_{173})) \to \mathrm{Seq}(a) \to \varphi_{215}$\lbl{A1}}
\BIC{$(\varphi_{220} \to \varphi_{224}) \to (\varphi_{173} \to \forall n\,(\varphi_{173})) \to \mathrm{Seq}(a) \to \varphi_{215}$}
\BIC{$(\varphi_{173} \to \forall n\,(\varphi_{173})) \to \mathrm{Seq}(a) \to \varphi_{215}$}
\BIC{$\mathrm{Seq}(a) \to \varphi_{215}$}
\end{ptree}
\noindent$\Pi_{40}$:
\begin{ptree}
\AXC{$\stackrel{\Pi_{21}}{\vdots}$}\noLine
\UIC{$\mathrm{Seq}(a) \to \varphi_{189}$}
\AXC{$\forall x \in \mathbb{R}\,(\lvert x \rvert \in \mathbb{R})$\lbl{[A5]}}
\AXC{$\varphi_{72}$\lbl{A2}}
\BIC{$a_{S(N)} \in \mathbb{R} \to \lvert a_{S(N)} \rvert \in \mathbb{R}$}
\AXC{$\varphi_{269}$\lbl{A1}}
\BIC{$(\mathrm{Seq}(a) \to \varphi_{189}) \to \mathrm{Seq}(a) \to \varphi_{194}$}
\BIC{$\mathrm{Seq}(a) \to \varphi_{194}$}
\end{ptree}
\noindent$\Pi_{41}$:
\begin{ptree}
\AXC{$\forall x\,\forall y\,(x \in \mathbb{R} \to y \in \mathbb{R} \to 0 \leq x \to 0 \leq y \to 0 \leq x + y)$\lbl{[E6]}}
\AXC{$\varphi_{107}$\lbl{A2}}
\BIC{$\forall y\,(M \in \mathbb{R} \to y \in \mathbb{R} \to 0 \leq M \to 0 \leq y \to 0 \leq M + y)$}
\AXC{$\varphi_{146}$\lbl{A2}}
\BIC{$M \in \mathbb{R} \to \varphi_{66}$}
\AXC{$(M \in \mathbb{R} \to \varphi_{66}) \to (\mathrm{Seq}(a) \to \varphi_{180}) \to \varphi_{285}$\lbl{A1}}
\BIC{$(\mathrm{Seq}(a) \to \varphi_{180}) \to \varphi_{285}$}
\end{ptree}
\noindent$\Pi_{42}$:
\begin{ptree}
\AXC{$\stackrel{\Pi_{40}}{\vdots}$}\noLine
\UIC{$\mathrm{Seq}(a) \to \varphi_{194}$}
\AXC{$\mathrm{Seq}(a) \to \varphi_{212}$\lbl{A1}}
\AXC{$(\mathrm{Seq}(a) \to \varphi_{212}) \to \mathrm{Seq}(a) \to \varphi_{180}$\lbl{A1}}
\BIC{$\mathrm{Seq}(a) \to \varphi_{180}$}
\AXC{$\stackrel{\Pi_{41}}{\vdots}$}\noLine
\UIC{$(\mathrm{Seq}(a) \to \varphi_{180}) \to \varphi_{285}$}
\BIC{$\varphi_{285}$}
\BIC{$\varphi_{181} \to (\mathrm{Seq}(a) \to \varphi_{186}) \to \mathrm{Seq}(a) \to \varphi_{193}$}
\end{ptree}
\noindent$\Pi_{43}$:
\begin{ptree}
\AXC{$\stackrel{\Pi_{22}}{\vdots}$}\noLine
\UIC{$\mathrm{Seq}(a) \to \varphi_{186}$}
\AXC{$\mathrm{Seq}(a) \to \varphi_{212}$\lbl{A1}}
\AXC{$(\mathrm{Seq}(a) \to \varphi_{212}) \to \varphi_{181}$\lbl{A1}}
\BIC{$\varphi_{181}$}
\AXC{$\stackrel{\Pi_{42}}{\vdots}$}\noLine
\UIC{$\varphi_{181} \to (\mathrm{Seq}(a) \to \varphi_{186}) \to \mathrm{Seq}(a) \to \varphi_{193}$}
\BIC{$(\mathrm{Seq}(a) \to \varphi_{186}) \to \mathrm{Seq}(a) \to \varphi_{193}$}
\BIC{$\mathrm{Seq}(a) \to \varphi_{193}$}
\end{ptree}
\noindent$\Pi_{44}$:
\begin{ptree}
\AXC{$\forall x\,\forall y\,(x \in \mathbb{R} \to y \in \mathbb{R} \to x + y \in \mathbb{R})$\lbl{(R1)}}
\AXC{$\varphi_{101}$\lbl{A2}}
\BIC{$\forall y\,(M \in \mathbb{R} \to y \in \mathbb{R} \to M + y \in \mathbb{R})$}
\AXC{$\varphi_{126}$\lbl{A2}}
\BIC{$M \in \mathbb{R} \to \lvert a_{S(N)} \rvert \in \mathbb{R} \to M + \lvert a_{S(N)} \rvert \in \mathbb{R}$}
\AXC{$\varphi_{283}$\lbl{A1}}
\BIC{$(\mathrm{Seq}(a) \to \varphi_{180}) \to (\mathrm{Seq}(a) \to \varphi_{194}) \to \mathrm{Seq}(a) \to \varphi_{198}$}
\end{ptree}
\noindent$\Pi_{45}$:
\begin{ptree}
\AXC{$\mathrm{Seq}(a) \to \varphi_{212}$\lbl{A1}}
\AXC{$(\mathrm{Seq}(a) \to \varphi_{212}) \to \mathrm{Seq}(a) \to \varphi_{180}$\lbl{A1}}
\BIC{$\mathrm{Seq}(a) \to \varphi_{180}$}
\AXC{$\stackrel{\Pi_{44}}{\vdots}$}\noLine
\UIC{$(\mathrm{Seq}(a) \to \varphi_{180}) \to (\mathrm{Seq}(a) \to \varphi_{194}) \to \mathrm{Seq}(a) \to \varphi_{198}$}
\BIC{$(\mathrm{Seq}(a) \to \varphi_{194}) \to \mathrm{Seq}(a) \to \varphi_{198}$}
\end{ptree}
\noindent$\Pi_{46}$:
\begin{ptree}
\AXC{$\stackrel{\Pi_{39}}{\vdots}$}\noLine
\UIC{$\mathrm{Seq}(a) \to \varphi_{215}$}
\AXC{$\stackrel{\Pi_{43}}{\vdots}$}\noLine
\UIC{$\mathrm{Seq}(a) \to \varphi_{193}$}
\AXC{$\stackrel{\Pi_{40}}{\vdots}$}\noLine
\UIC{$\mathrm{Seq}(a) \to \varphi_{194}$}
\AXC{$\stackrel{\Pi_{45}}{\vdots}$}\noLine
\UIC{$(\mathrm{Seq}(a) \to \varphi_{194}) \to \mathrm{Seq}(a) \to \varphi_{198}$}
\BIC{$\mathrm{Seq}(a) \to \varphi_{198}$}
\AXC{$\varphi_{292}$\lbl{A1}}
\BIC{$(\mathrm{Seq}(a) \to \varphi_{193}) \to (\mathrm{Seq}(a) \to \varphi_{215}) \to \varphi_{238}$}
\BIC{$(\mathrm{Seq}(a) \to \varphi_{215}) \to \varphi_{238}$}
\BIC{$\varphi_{238}$}
\end{ptree}
\noindent$\Pi_{47}$:
\begin{ptree}
\AXC{$\forall M\,(\varphi_{115}) \to \neg \varphi_{127}$\lbl{A2}}
\AXC{$\stackrel{\Pi_{46}}{\vdots}$}\noLine
\UIC{$\varphi_{238}$}
\AXC{$\varphi_{238} \to (\forall M\,(\varphi_{115}) \to \neg \varphi_{127}) \to \mathrm{Seq}(a) \to \varphi_{216}$\lbl{A1}}
\BIC{$(\forall M\,(\varphi_{115}) \to \neg \varphi_{127}) \to \mathrm{Seq}(a) \to \varphi_{216}$}
\BIC{$\mathrm{Seq}(a) \to \varphi_{216}$}
\AXC{$(\mathrm{Seq}(a) \to \varphi_{216}) \to \varphi_{226}$\lbl{A1}}
\BIC{$\varphi_{226}$}
\UIC{$\forall M\,(\varphi_{226})$}
\end{ptree}
\noindent$\Pi_{48}$:
\begin{ptree}
\AXC{$\stackrel{\Pi_{7}}{\vdots}$}\noLine
\UIC{$\mathrm{Seq}(a) \to \varphi_{153}$}
\AXC{$\stackrel{\Pi_{4}}{\vdots}$}\noLine
\UIC{$\varphi_{172}$}
\AXC{$\varphi_{172} \to \mathrm{Seq}(a) \to \forall N\,(N \in B_{8} \leftrightarrow N \in \omega \land \varphi_{51}) \to N \in B_{8} \to \varphi_{51}$\lbl{A1}}
\BIC{$\mathrm{Seq}(a) \to \forall N\,(N \in B_{8} \leftrightarrow N \in \omega \land \varphi_{51}) \to N \in B_{8} \to \varphi_{51}$}
\AXC{$\varphi_{274}$\lbl{A1}}
\BIC{$(\mathrm{Seq}(a) \to \varphi_{153}) \to \varphi_{176}$}
\BIC{$\varphi_{176}$}
\end{ptree}
\noindent$\Pi_{49}$:
\begin{ptree}
\AXC{$\forall M\,(\varphi_{226}) \to \forall M\,(\varphi_{183}) \to \varphi_{27}$\lbl{A3}}
\AXC{$\stackrel{\Pi_{47}}{\vdots}$}\noLine
\UIC{$\forall M\,(\varphi_{226})$}
\AXC{$\stackrel{\Pi_{48}}{\vdots}$}\noLine
\UIC{$\varphi_{176}$}
\AXC{$\varphi_{176} \to \forall M\,(\varphi_{226}) \to (\forall M\,(\varphi_{226}) \to \forall M\,(\varphi_{183}) \to \varphi_{27}) \to \varphi_{278}$\lbl{A1}}
\BIC{$\forall M\,(\varphi_{226}) \to (\forall M\,(\varphi_{226}) \to \forall M\,(\varphi_{183}) \to \varphi_{27}) \to \varphi_{278}$}
\BIC{$(\forall M\,(\varphi_{226}) \to \forall M\,(\varphi_{183}) \to \varphi_{27}) \to \varphi_{278}$}
\BIC{$\varphi_{278}$}
\end{ptree}
\noindent$\Pi_{50}$:
\begin{ptree}
\AXC{$z = w \to n = z \to n = w$\lbl{A6}}
\UIC{$\forall w\,(z = w \to n = z \to n = w)$}
\UIC{$\forall z\,\forall w\,(z = w \to n = z \to n = w)$}
\AXC{$\varphi_{10}$\lbl{A2}}
\BIC{$\forall w\,(1_{\omega} = w \to n = 1_{\omega} \to n = w)$}
\AXC{$\varphi_{15}$\lbl{A2}}
\BIC{$1_{\omega} = S(N) \to n = 1_{\omega} \to n = S(N)$}
\AXC{$\varphi_{118}$\lbl{A1}}
\BIC{$(S(N) = 1_{\omega} \to n = S(N) \to n = 1_{\omega}) \to \varphi_{11}$}
\end{ptree}
\noindent$\Pi_{51}$:
\begin{ptree}
\AXC{$z = w \to n = z \to n = w$\lbl{A6}}
\UIC{$\forall w\,(z = w \to n = z \to n = w)$}
\UIC{$\forall z\,\forall w\,(z = w \to n = z \to n = w)$}
\AXC{$\varphi_{9}$\lbl{A2}}
\BIC{$\forall w\,(S(N) = w \to n = S(N) \to n = w)$}
\AXC{$\varphi_{14}$\lbl{A2}}
\BIC{$S(N) = 1_{\omega} \to n = S(N) \to n = 1_{\omega}$}
\AXC{$\stackrel{\Pi_{50}}{\vdots}$}\noLine
\UIC{$(S(N) = 1_{\omega} \to n = S(N) \to n = 1_{\omega}) \to \varphi_{11}$}
\BIC{$\varphi_{11}$}
\end{ptree}
\noindent$\Pi_{52}$:
\begin{ptree}
\AXC{$z = w \to n < z \to n < w$\lbl{A6}}
\UIC{$\forall w\,(z = w \to n < z \to n < w)$}
\UIC{$\forall z\,\forall w\,(z = w \to n < z \to n < w)$}
\AXC{$\varphi_{32}$\lbl{A2}}
\BIC{$\forall w\,(1_{\omega} = w \to n < 1_{\omega} \to n < w)$}
\AXC{$\varphi_{41}$\lbl{A2}}
\BIC{$1_{\omega} = S(N) \to n < 1_{\omega} \to n < S(N)$}
\AXC{$\varphi_{133}$\lbl{A1}}
\BIC{$(S(N) = 1_{\omega} \to n < S(N) \to n < 1_{\omega}) \to \varphi_{25}$}
\end{ptree}
\noindent$\Pi_{53}$:
\begin{ptree}
\AXC{$z = w \to n < z \to n < w$\lbl{A6}}
\UIC{$\forall w\,(z = w \to n < z \to n < w)$}
\UIC{$\forall z\,\forall w\,(z = w \to n < z \to n < w)$}
\AXC{$\varphi_{31}$\lbl{A2}}
\BIC{$\forall w\,(S(N) = w \to n < S(N) \to n < w)$}
\AXC{$\varphi_{40}$\lbl{A2}}
\BIC{$S(N) = 1_{\omega} \to n < S(N) \to n < 1_{\omega}$}
\AXC{$\stackrel{\Pi_{52}}{\vdots}$}\noLine
\UIC{$(S(N) = 1_{\omega} \to n < S(N) \to n < 1_{\omega}) \to \varphi_{25}$}
\BIC{$\varphi_{25}$}
\end{ptree}
\noindent$\Pi_{54}$:
\begin{ptree}
\AXC{$\stackrel{\Pi_{15}}{\vdots}$}\noLine
\UIC{$\forall x\,\forall y\,(x = y \to y = x)$}
\AXC{$\varphi_{1}$\lbl{A2}}
\BIC{$\forall y\,(1_{\omega} = y \to y = 1_{\omega})$}
\AXC{$\varphi_{4}$\lbl{A2}}
\BIC{$1_{\omega} = S(N) \to S(N) = 1_{\omega}$}
\AXC{$\stackrel{\Pi_{53}}{\vdots}$}\noLine
\UIC{$\varphi_{25}$}
\BIC{$1_{\omega} = S(N) \to (n < 1_{\omega} \leftrightarrow n < S(N))$}
\AXC{$\varphi_{53}$\lbl{A1}}
\BIC{$1_{\omega} = S(N) \to (\neg n < 1_{\omega} \leftrightarrow \neg n < S(N))$}
\AXC{$\varphi_{131}$\lbl{A1}}
\BIC{$\varphi_{74}$}
\end{ptree}
\noindent$\Pi_{55}$:
\begin{ptree}
\AXC{$\varphi_{91}$\lbl{A1}}
\AXC{$\stackrel{\Pi_{15}}{\vdots}$}\noLine
\UIC{$\forall x\,\forall y\,(x = y \to y = x)$}
\AXC{$\varphi_{1}$\lbl{A2}}
\BIC{$\forall y\,(1_{\omega} = y \to y = 1_{\omega})$}
\AXC{$\varphi_{4}$\lbl{A2}}
\BIC{$1_{\omega} = S(N) \to S(N) = 1_{\omega}$}
\AXC{$\stackrel{\Pi_{51}}{\vdots}$}\noLine
\UIC{$\varphi_{11}$}
\BIC{$1_{\omega} = S(N) \to (n = 1_{\omega} \leftrightarrow n = S(N))$}
\AXC{$\stackrel{\Pi_{54}}{\vdots}$}\noLine
\UIC{$\varphi_{74}$}
\BIC{$1_{\omega} = S(N) \to \varphi_{3}$}
\AXC{$\varphi_{230}$\lbl{A1}}
\BIC{$\varphi_{91} \to 1_{\omega} = S(N) \to (\varphi_{28} \leftrightarrow \varphi_{29})$}
\BIC{$1_{\omega} = S(N) \to (\varphi_{28} \leftrightarrow \varphi_{29})$}
\end{ptree}
\noindent$\Pi_{56}$:
\begin{ptree}
\AXC{$1_{\omega} = S(N) \to (n \in \mathbb{N} \leftrightarrow n \in \mathbb{N})$\lbl{A1}}
\AXC{$(1_{\omega} = S(N) \to (n \in \mathbb{N} \leftrightarrow n \in \mathbb{N})) \to \varphi_{231}$\lbl{A1}}
\BIC{$\varphi_{231}$}
\end{ptree}
\noindent$\Pi_{57}$:
\begin{ptree}
\AXC{$\stackrel{\Pi_{55}}{\vdots}$}\noLine
\UIC{$1_{\omega} = S(N) \to (\varphi_{28} \leftrightarrow \varphi_{29})$}
\AXC{$\stackrel{\Pi_{56}}{\vdots}$}\noLine
\UIC{$\varphi_{231}$}
\BIC{$1_{\omega} = S(N) \to (n \in \mathbb{N} \to \varphi_{28} \leftrightarrow n \in \mathbb{N} \to \varphi_{29})$}
\end{ptree}
\noindent$\Pi_{58}$:
\begin{ptree}
\AXC{$\varphi_{239}$\lbl{A3}}
\AXC{$\stackrel{\Pi_{57}}{\vdots}$}\noLine
\UIC{$1_{\omega} = S(N) \to (n \in \mathbb{N} \to \varphi_{28} \leftrightarrow n \in \mathbb{N} \to \varphi_{29})$}
\AXC{$\varphi_{235}$\lbl{A1}}
\BIC{$1_{\omega} = S(N) \to (n \in \mathbb{N} \to \varphi_{29}) \to n \in \mathbb{N} \to \varphi_{28}$}
\UIC{$\forall n\,(1_{\omega} = S(N) \to (n \in \mathbb{N} \to \varphi_{29}) \to n \in \mathbb{N} \to \varphi_{28})$}
\AXC{$\varphi_{279}$\lbl{A1}}
\BIC{$\varphi_{239} \to \varphi_{164}$}
\BIC{$\varphi_{164}$}
\end{ptree}
\noindent$\Pi_{59}$:
\begin{ptree}
\AXC{$\varphi_{240}$\lbl{A3}}
\AXC{$\stackrel{\Pi_{57}}{\vdots}$}\noLine
\UIC{$1_{\omega} = S(N) \to (n \in \mathbb{N} \to \varphi_{28} \leftrightarrow n \in \mathbb{N} \to \varphi_{29})$}
\AXC{$\varphi_{236}$\lbl{A1}}
\BIC{$1_{\omega} = S(N) \to (n \in \mathbb{N} \to \varphi_{28}) \to n \in \mathbb{N} \to \varphi_{29}$}
\UIC{$\forall n\,(1_{\omega} = S(N) \to (n \in \mathbb{N} \to \varphi_{28}) \to n \in \mathbb{N} \to \varphi_{29})$}
\AXC{$\varphi_{280}$\lbl{A1}}
\BIC{$\varphi_{240} \to \varphi_{165}$}
\BIC{$\varphi_{165}$}
\end{ptree}
\noindent$\Pi_{60}$:
\begin{ptree}
\AXC{$1_{\omega} = S(N) \to \forall n\,(1_{\omega} = S(N))$\lbl{A4}}
\AXC{$\stackrel{\Pi_{59}}{\vdots}$}\noLine
\UIC{$\varphi_{165}$}
\BIC{$1_{\omega} = S(N) \to \forall n\,((n \in \mathbb{N} \to \varphi_{28}) \to n \in \mathbb{N} \to \varphi_{29})$}
\AXC{$\varphi_{293}$\lbl{A1}}
\BIC{$\varphi_{287}$}
\end{ptree}
\noindent$\Pi_{61}$:
\begin{ptree}
\AXC{$\varphi_{232}$\lbl{A3}}
\AXC{$\varphi_{233}$\lbl{A3}}
\AXC{$1_{\omega} = S(N) \to \forall n\,(1_{\omega} = S(N))$\lbl{A4}}
\AXC{$\stackrel{\Pi_{58}}{\vdots}$}\noLine
\UIC{$\varphi_{164}$}
\BIC{$1_{\omega} = S(N) \to \forall n\,((n \in \mathbb{N} \to \varphi_{29}) \to n \in \mathbb{N} \to \varphi_{28})$}
\AXC{$\stackrel{\Pi_{60}}{\vdots}$}\noLine
\UIC{$\varphi_{287}$}
\BIC{$\varphi_{233} \to \varphi_{259}$}
\BIC{$\varphi_{259}$}
\BIC{$1_{\omega} = S(N) \to (\forall n \in \mathbb{N}\,(\varphi_{28}) \leftrightarrow \forall n \in \mathbb{N}\,(\varphi_{29}))$}
\AXC{$\varphi_{241}$\lbl{A1}}
\BIC{$\varphi_{160}$}
\end{ptree}
\noindent$\Pi_{62}$:
\begin{ptree}
\AXC{$\stackrel{\Pi_{61}}{\vdots}$}\noLine
\UIC{$\varphi_{160}$}
\AXC{$1_{\omega} = S(N) \to \varphi_{30}$\lbl{A1}}
\AXC{$(1_{\omega} = S(N) \to \varphi_{30}) \to \varphi_{160} \to 1_{\omega} = S(N) \to \varphi_{184}$\lbl{A1}}
\BIC{$\varphi_{160} \to 1_{\omega} = S(N) \to \varphi_{184}$}
\BIC{$1_{\omega} = S(N) \to \varphi_{184}$}
\end{ptree}
\noindent$\Pi_{63}$:
\begin{ptree}
\AXC{$\stackrel{\Pi_{62}}{\vdots}$}\noLine
\UIC{$1_{\omega} = S(N) \to \varphi_{184}$}
\AXC{$(1_{\omega} = S(N) \to \varphi_{184}) \to 1_{\omega} = S(N) \to \varphi_{187}$\lbl{A1}}
\BIC{$1_{\omega} = S(N) \to \varphi_{187}$}
\AXC{$\varphi_{256}$\lbl{A1}}
\BIC{$1_{\omega} = S(N) \to (\varphi_{114} \leftrightarrow \varphi_{115})$}
\end{ptree}
\noindent$\Pi_{64}$:
\begin{ptree}
\AXC{$\varphi_{261}$\lbl{A3}}
\AXC{$\stackrel{\Pi_{63}}{\vdots}$}\noLine
\UIC{$1_{\omega} = S(N) \to (\varphi_{114} \leftrightarrow \varphi_{115})$}
\AXC{$\varphi_{257}$\lbl{A1}}
\BIC{$1_{\omega} = S(N) \to \varphi_{115} \to \varphi_{114}$}
\UIC{$\forall M\,(1_{\omega} = S(N) \to \varphi_{115} \to \varphi_{114})$}
\AXC{$\forall M\,(1_{\omega} = S(N) \to \varphi_{115} \to \varphi_{114}) \to \varphi_{261} \to \varphi_{217}$\lbl{A1}}
\BIC{$\varphi_{261} \to \varphi_{217}$}
\BIC{$\varphi_{217}$}
\end{ptree}
\noindent$\Pi_{65}$:
\begin{ptree}
\AXC{$\varphi_{262}$\lbl{A3}}
\AXC{$\stackrel{\Pi_{63}}{\vdots}$}\noLine
\UIC{$1_{\omega} = S(N) \to (\varphi_{114} \leftrightarrow \varphi_{115})$}
\AXC{$\varphi_{258}$\lbl{A1}}
\BIC{$1_{\omega} = S(N) \to \varphi_{114} \to \varphi_{115}$}
\UIC{$\forall M\,(1_{\omega} = S(N) \to \varphi_{114} \to \varphi_{115})$}
\AXC{$\forall M\,(1_{\omega} = S(N) \to \varphi_{114} \to \varphi_{115}) \to \varphi_{262} \to \varphi_{218}$\lbl{A1}}
\BIC{$\varphi_{262} \to \varphi_{218}$}
\BIC{$\varphi_{218}$}
\end{ptree}
\noindent$\Pi_{66}$:
\begin{ptree}
\AXC{$1_{\omega} = S(N) \to \forall M\,(1_{\omega} = S(N))$\lbl{A4}}
\AXC{$\stackrel{\Pi_{65}}{\vdots}$}\noLine
\UIC{$\varphi_{218}$}
\BIC{$1_{\omega} = S(N) \to \forall M\,(\varphi_{114} \to \varphi_{115})$}
\AXC{$(1_{\omega} = S(N) \to \forall M\,(\varphi_{114} \to \varphi_{115})) \to \varphi_{294}$\lbl{A1}}
\BIC{$\varphi_{294}$}
\end{ptree}
\noindent$\Pi_{67}$:
\begin{ptree}
\AXC{$\forall M\,(\varphi_{114} \to \varphi_{115}) \to \forall M\,(\varphi_{114}) \to \forall M\,(\varphi_{115})$\lbl{A3}}
\AXC{$1_{\omega} = S(N) \to \forall M\,(1_{\omega} = S(N))$\lbl{A4}}
\AXC{$\stackrel{\Pi_{64}}{\vdots}$}\noLine
\UIC{$\varphi_{217}$}
\BIC{$1_{\omega} = S(N) \to \forall M\,(\varphi_{115} \to \varphi_{114})$}
\AXC{$\stackrel{\Pi_{66}}{\vdots}$}\noLine
\UIC{$\varphi_{294}$}
\BIC{$(\forall M\,(\varphi_{114} \to \varphi_{115}) \to \forall M\,(\varphi_{114}) \to \forall M\,(\varphi_{115})) \to \varphi_{281}$}
\BIC{$\varphi_{281}$}
\end{ptree}
\noindent$\Pi_{68}$:
\begin{ptree}
\AXC{$\forall M\,(\varphi_{115} \to \varphi_{114}) \to \forall M\,(\varphi_{115}) \to \forall M\,(\varphi_{114})$\lbl{A3}}
\AXC{$\stackrel{\Pi_{67}}{\vdots}$}\noLine
\UIC{$\varphi_{281}$}
\BIC{$1_{\omega} = S(N) \to (\forall M\,(\varphi_{114}) \leftrightarrow \forall M\,(\varphi_{115}))$}
\AXC{$\varphi_{263}$\lbl{A1}}
\BIC{$1_{\omega} = S(N) \to (\neg \forall M\,(\varphi_{114}) \leftrightarrow \neg \forall M\,(\varphi_{115}))$}
\end{ptree}
\noindent$\Pi_{69}$:
\begin{ptree}
\AXC{$z = w \to \lvert a_{n} \rvert = \lvert a_{z} \rvert \to \lvert a_{n} \rvert = \lvert a_{w} \rvert$\lbl{A6}}
\UIC{$\forall w\,(z = w \to \lvert a_{n} \rvert = \lvert a_{z} \rvert \to \lvert a_{n} \rvert = \lvert a_{w} \rvert)$}
\UIC{$\forall z\,\forall w\,(z = w \to \lvert a_{n} \rvert = \lvert a_{z} \rvert \to \lvert a_{n} \rvert = \lvert a_{w} \rvert)$}
\AXC{$\varphi_{143}$\lbl{A2}}
\BIC{$\forall w\,(n = w \to \lvert a_{n} \rvert = \lvert a_{n} \rvert \to \lvert a_{n} \rvert = \lvert a_{w} \rvert)$}
\AXC{$\forall w\,(n = w \to \lvert a_{n} \rvert = \lvert a_{n} \rvert \to \lvert a_{n} \rvert = \lvert a_{w} \rvert) \to \varphi_{49}$\lbl{A2}}
\BIC{$\varphi_{49}$}
\end{ptree}
\noindent$\Pi_{70}$:
\begin{ptree}
\AXC{$z = w \to \lvert a_{n} \rvert < \lvert a_{z} \rvert \to \lvert a_{n} \rvert < \lvert a_{w} \rvert$\lbl{A6}}
\UIC{$\forall w\,(z = w \to \lvert a_{n} \rvert < \lvert a_{z} \rvert \to \lvert a_{n} \rvert < \lvert a_{w} \rvert)$}
\UIC{$\forall z\,\forall w\,(z = w \to \lvert a_{n} \rvert < \lvert a_{z} \rvert \to \lvert a_{n} \rvert < \lvert a_{w} \rvert)$}
\AXC{$\varphi_{152}$\lbl{A2}}
\BIC{$\forall w\,(n = w \to \lvert a_{n} \rvert < \lvert a_{n} \rvert \to \lvert a_{n} \rvert < \lvert a_{w} \rvert)$}
\AXC{$\forall w\,(n = w \to \lvert a_{n} \rvert < \lvert a_{n} \rvert \to \lvert a_{n} \rvert < \lvert a_{w} \rvert) \to \varphi_{61}$\lbl{A2}}
\BIC{$\varphi_{61}$}
\end{ptree}
\noindent$\Pi_{71}$:
\begin{ptree}
\AXC{$\forall n \in \mathbb{N}\,(1_{\omega} \leq n)$\lbl{[N12]}}
\AXC{$\forall n \in \mathbb{N}\,(1_{\omega} \leq n) \to n \in \mathbb{N} \to 1_{\omega} \leq n$\lbl{A2}}
\BIC{$n \in \mathbb{N} \to 1_{\omega} \leq n$}
\AXC{$(n \in \mathbb{N} \to 1_{\omega} \leq n) \to \varphi_{50}$\lbl{A1}}
\BIC{$\varphi_{50}$}
\end{ptree}
\noindent$\Pi_{72}$:
\begin{ptree}
\AXC{$\forall x\,\forall y\,(x \leq y \to y \leq x \to x = y)$\lbl{[O6]}}
\AXC{$\varphi_{5}$\lbl{A2}}
\BIC{$\forall y\,(n \leq y \to y \leq n \to n = y)$}
\AXC{$\varphi_{7}$\lbl{A2}}
\BIC{$n \leq 1_{\omega} \to 1_{\omega} \leq n \to n = 1_{\omega}$}
\AXC{$(n \leq 1_{\omega} \to 1_{\omega} \leq n \to n = 1_{\omega}) \to \varphi_{123}$\lbl{A1}}
\BIC{$\varphi_{123}$}
\end{ptree}
\noindent$\Pi_{73}$:
\begin{ptree}
\AXC{$\mathrm{Seq}(a) \to n \in \mathbb{N} \to n \in \mathbb{N}$\lbl{A1}}
\AXC{$\stackrel{\Pi_{71}}{\vdots}$}\noLine
\UIC{$\varphi_{50}$}
\BIC{$\mathrm{Seq}(a) \to n \in \mathbb{N} \to n \leq 1_{\omega} \to 1_{\omega} \leq n$}
\AXC{$\mathrm{Seq}(a) \to n \in \mathbb{N} \to n \leq 1_{\omega} \to n \leq 1_{\omega}$\lbl{A1}}
\AXC{$\stackrel{\Pi_{72}}{\vdots}$}\noLine
\UIC{$\varphi_{123}$}
\BIC{$\varphi_{62}$}
\BIC{$\mathrm{Seq}(a) \to n \in \mathbb{N} \to n \leq 1_{\omega} \to n = 1_{\omega}$}
\AXC{$\varphi_{234}$\lbl{A1}}
\BIC{$\varphi_{222}$}
\end{ptree}
\noindent$\Pi_{74}$:
\begin{ptree}
\AXC{$\stackrel{\Pi_{69}}{\vdots}$}\noLine
\UIC{$\varphi_{49}$}
\AXC{$\stackrel{\Pi_{70}}{\vdots}$}\noLine
\UIC{$\varphi_{61}$}
\AXC{$x = x$\lbl{A5}}
\UIC{$\forall x\,(x = x)$}
\AXC{$\forall x\,(x = x) \to \lvert a_{n} \rvert = \lvert a_{n} \rvert$\lbl{A2}}
\BIC{$\lvert a_{n} \rvert = \lvert a_{n} \rvert$}
\AXC{$\varphi_{82}$\lbl{A1}}
\BIC{$\mathrm{Seq}(a) \to n \in \mathbb{N} \to n \leq 1_{\omega} \to \lvert a_{n} \rvert \leq \lvert a_{n} \rvert$}
\AXC{$\stackrel{\Pi_{73}}{\vdots}$}\noLine
\UIC{$\varphi_{222}$}
\BIC{$\varphi_{61} \to \varphi_{49} \to \varphi_{33}$}
\BIC{$\varphi_{49} \to \varphi_{33}$}
\BIC{$\varphi_{33}$}
\UIC{$\forall n\,(\varphi_{33})$}
\AXC{$\forall n\,(\varphi_{33}) \to \varphi_{188}$\lbl{A1}}
\BIC{$\varphi_{188}$}
\end{ptree}
\noindent$\Pi_{75}$:
\begin{ptree}
\AXC{$\mathrm{Seq}(a) \to \forall n\,(\mathrm{Seq}(a))$\lbl{A4}}
\AXC{$\forall n\,(\varphi_{33}) \to \varphi_{44}$\lbl{A3}}
\AXC{$\stackrel{\Pi_{74}}{\vdots}$}\noLine
\UIC{$\varphi_{188}$}
\BIC{$(\mathrm{Seq}(a) \to \forall n\,(\mathrm{Seq}(a))) \to \varphi_{36}$}
\BIC{$\varphi_{36}$}
\end{ptree}
\noindent$\Pi_{76}$:
\begin{ptree}
\AXC{$\forall a\,\forall n\,(\mathrm{Seq}(a) \to n \in \mathbb{N} \to a_{n} \in \mathbb{R})$\lbl{[S1]}}
\AXC{$\varphi_{85}$\lbl{A2}}
\BIC{$\forall n\,(\mathrm{Seq}(a) \to n \in \mathbb{N} \to a_{n} \in \mathbb{R})$}
\AXC{$\varphi_{88}$\lbl{A2}}
\BIC{$\mathrm{Seq}(a) \to 1_{\omega} \in \mathbb{N} \to a_{1_{\omega}} \in \mathbb{R}$}
\end{ptree}
\noindent$\Pi_{77}$:
\begin{ptree}
\AXC{$\stackrel{\Pi_{76}}{\vdots}$}\noLine
\UIC{$\mathrm{Seq}(a) \to 1_{\omega} \in \mathbb{N} \to a_{1_{\omega}} \in \mathbb{R}$}
\AXC{$(\mathrm{Seq}(a) \to 1_{\omega} \in \mathbb{N} \to a_{1_{\omega}} \in \mathbb{R}) \to \varphi_{35}$\lbl{A1}}
\BIC{$\varphi_{35}$}
\end{ptree}
\noindent$\Pi_{78}$:
\begin{ptree}
\AXC{$1_{\omega} \in \mathbb{N}$\lbl{[N1]}}
\AXC{$\mathrm{Seq}(a) \to \mathrm{Seq}(a)$\lbl{A1}}
\AXC{$\stackrel{\Pi_{77}}{\vdots}$}\noLine
\UIC{$\varphi_{35}$}
\BIC{$1_{\omega} \in \mathbb{N} \to \mathrm{Seq}(a) \to a_{1_{\omega}} \in \mathbb{R}$}
\BIC{$\mathrm{Seq}(a) \to a_{1_{\omega}} \in \mathbb{R}$}
\end{ptree}
\noindent$\Pi_{79}$:
\begin{ptree}
\AXC{$\forall x \in \mathbb{R}\,(0 \leq \lvert x \rvert)$\lbl{[A6]}}
\AXC{$\varphi_{46}$\lbl{A2}}
\BIC{$a_{1_{\omega}} \in \mathbb{R} \to 0 \leq \lvert a_{1_{\omega}} \rvert$}
\AXC{$(a_{1_{\omega}} \in \mathbb{R} \to 0 \leq \lvert a_{1_{\omega}} \rvert) \to \varphi_{26}$\lbl{A1}}
\BIC{$\varphi_{26}$}
\end{ptree}
\noindent$\Pi_{80}$:
\begin{ptree}
\AXC{$\forall x \in \mathbb{R}\,(\lvert x \rvert \in \mathbb{R})$\lbl{[A5]}}
\AXC{$\varphi_{73}$\lbl{A2}}
\BIC{$a_{1_{\omega}} \in \mathbb{R} \to \lvert a_{1_{\omega}} \rvert \in \mathbb{R}$}
\AXC{$(a_{1_{\omega}} \in \mathbb{R} \to \lvert a_{1_{\omega}} \rvert \in \mathbb{R}) \to \varphi_{42}$\lbl{A1}}
\BIC{$\varphi_{42}$}
\end{ptree}
\noindent$\Pi_{81}$:
\begin{ptree}
\AXC{$\stackrel{\Pi_{78}}{\vdots}$}\noLine
\UIC{$\mathrm{Seq}(a) \to a_{1_{\omega}} \in \mathbb{R}$}
\AXC{$\stackrel{\Pi_{80}}{\vdots}$}\noLine
\UIC{$\varphi_{42}$}
\BIC{$\mathrm{Seq}(a) \to \lvert a_{1_{\omega}} \rvert \in \mathbb{R}$}
\AXC{$(\mathrm{Seq}(a) \to \lvert a_{1_{\omega}} \rvert \in \mathbb{R}) \to \varphi_{177}$\lbl{A1}}
\BIC{$\varphi_{177}$}
\end{ptree}
\noindent$\Pi_{82}$:
\begin{ptree}
\AXC{$\forall M\,(\varphi_{114}) \to \neg \varphi_{106}$\lbl{A2}}
\AXC{$\stackrel{\Pi_{75}}{\vdots}$}\noLine
\UIC{$\varphi_{36}$}
\AXC{$\stackrel{\Pi_{78}}{\vdots}$}\noLine
\UIC{$\mathrm{Seq}(a) \to a_{1_{\omega}} \in \mathbb{R}$}
\AXC{$\stackrel{\Pi_{79}}{\vdots}$}\noLine
\UIC{$\varphi_{26}$}
\BIC{$\mathrm{Seq}(a) \to 0 \leq \lvert a_{1_{\omega}} \rvert$}
\AXC{$\stackrel{\Pi_{81}}{\vdots}$}\noLine
\UIC{$\varphi_{177}$}
\BIC{$\varphi_{36} \to \mathrm{Seq}(a) \to \varphi_{106}$}
\BIC{$\mathrm{Seq}(a) \to \varphi_{106}$}
\AXC{$(\mathrm{Seq}(a) \to \varphi_{106}) \to \varphi_{207}$\lbl{A1}}
\BIC{$\varphi_{207}$}
\BIC{$\mathrm{Seq}(a) \to \exists M\,(M \in \mathbb{R} \land 0 \leq M \land \forall n \in \mathbb{N}\,(\varphi_{28}))$}
\end{ptree}
\noindent$\Pi_{83}$:
\begin{ptree}
\AXC{$y = z \to x = y \to x = z$\lbl{A6}}
\AXC{$(y = z \to x = y \to x = z) \to x = y \to y = z \to x = z$\lbl{A1}}
\BIC{$x = y \to y = z \to x = z$}
\UIC{$\forall z\,(x = y \to y = z \to x = z)$}
\UIC{$\forall y\,\forall z\,(x = y \to y = z \to x = z)$}
\UIC{$\forall x\,\forall y\,\forall z\,(x = y \to y = z \to x = z)$}
\AXC{$\varphi_{16}$\lbl{A2}}
\BIC{$\forall y\,\forall z\,(S(N) = y \to y = z \to S(N) = z)$}
\AXC{$\varphi_{59}$\lbl{A2}}
\BIC{$\forall z\,(S(N) = S(\emptyset) \to S(\emptyset) = z \to S(N) = z)$}
\end{ptree}
\noindent$\Pi_{84}$:
\begin{ptree}
\AXC{$\stackrel{\Pi_{83}}{\vdots}$}\noLine
\UIC{$\forall z\,(S(N) = S(\emptyset) \to S(\emptyset) = z \to S(N) = z)$}
\AXC{$\forall z\,(S(N) = S(\emptyset) \to S(\emptyset) = z \to S(N) = z) \to \varphi_{8}$\lbl{A2}}
\BIC{$\varphi_{8}$}
\end{ptree}
\noindent$\Pi_{85}$:
\begin{ptree}
\AXC{$\stackrel{\Pi_{15}}{\vdots}$}\noLine
\UIC{$\forall x\,\forall y\,(x = y \to y = x)$}
\AXC{$\varphi_{1}$\lbl{A2}}
\BIC{$\forall y\,(1_{\omega} = y \to y = 1_{\omega})$}
\AXC{$\varphi_{21}$\lbl{A2}}
\BIC{$1_{\omega} = S(\emptyset) \to S(\emptyset) = 1_{\omega}$}
\AXC{$1_{\omega} = S(\emptyset)$\lbl{Def$1_\omega$}}
\AXC{$1_{\omega} = S(\emptyset) \to \varphi_{17}$\lbl{A1}}
\BIC{$\varphi_{17}$}
\BIC{$S(\emptyset) = 1_{\omega}$}
\end{ptree}
\noindent$\Pi_{86}$:
\begin{ptree}
\AXC{$x = x$\lbl{A5}}
\UIC{$\forall x\,(x = x)$}
\AXC{$\forall x\,(x = x) \to S(N) = S(N)$\lbl{A2}}
\BIC{$S(N) = S(N)$}
\AXC{$\mathrm{Seq}(a) \to \varphi_{154}$\lbl{A1}}
\AXC{$(\mathrm{Seq}(a) \to \varphi_{154}) \to S(N) = S(N) \to \varphi_{179}$\lbl{A1}}
\BIC{$S(N) = S(N) \to \varphi_{179}$}
\BIC{$\varphi_{179}$}
\end{ptree}
\noindent$\Pi_{87}$:
\begin{ptree}
\AXC{$z = w \to S(N) = S(z) \to S(N) = S(w)$\lbl{A6}}
\UIC{$\forall w\,(z = w \to S(N) = S(z) \to S(N) = S(w))$}
\UIC{$\forall z\,\forall w\,(z = w \to S(N) = S(z) \to S(N) = S(w))$}
\AXC{$\varphi_{43}$\lbl{A2}}
\BIC{$\forall w\,(N = w \to S(N) = S(N) \to S(N) = S(w))$}
\AXC{$\varphi_{60}$\lbl{A2}}
\BIC{$N = \emptyset \to S(N) = S(N) \to S(N) = S(\emptyset)$}
\AXC{$\stackrel{\Pi_{86}}{\vdots}$}\noLine
\UIC{$\varphi_{179}$}
\BIC{$\mathrm{Seq}(a) \to \varphi_{157}$}
\AXC{$\varphi_{246}$\lbl{A1}}
\BIC{$S(\emptyset) = 1_{\omega} \to \varphi_{8} \to \mathrm{Seq}(a) \to \varphi_{150}$}
\end{ptree}
\noindent$\Pi_{88}$:
\begin{ptree}
\AXC{$\stackrel{\Pi_{84}}{\vdots}$}\noLine
\UIC{$\varphi_{8}$}
\AXC{$\stackrel{\Pi_{85}}{\vdots}$}\noLine
\UIC{$S(\emptyset) = 1_{\omega}$}
\AXC{$\stackrel{\Pi_{87}}{\vdots}$}\noLine
\UIC{$S(\emptyset) = 1_{\omega} \to \varphi_{8} \to \mathrm{Seq}(a) \to \varphi_{150}$}
\BIC{$\varphi_{8} \to \mathrm{Seq}(a) \to \varphi_{150}$}
\BIC{$\mathrm{Seq}(a) \to \varphi_{150}$}
\AXC{$(\mathrm{Seq}(a) \to \varphi_{150}) \to \varphi_{162}$\lbl{A1}}
\BIC{$\varphi_{162}$}
\end{ptree}
\noindent$\Pi_{89}$:
\begin{ptree}
\AXC{$\stackrel{\Pi_{15}}{\vdots}$}\noLine
\UIC{$\forall x\,\forall y\,(x = y \to y = x)$}
\AXC{$\varphi_{2}$\lbl{A2}}
\BIC{$\forall y\,(S(N) = y \to y = S(N))$}
\AXC{$\varphi_{6}$\lbl{A2}}
\BIC{$S(N) = 1_{\omega} \to 1_{\omega} = S(N)$}
\AXC{$\stackrel{\Pi_{88}}{\vdots}$}\noLine
\UIC{$\varphi_{162}$}
\BIC{$\mathrm{Seq}(a) \to \varphi_{151}$}
\AXC{$(\mathrm{Seq}(a) \to \varphi_{151}) \to \varphi_{271}$\lbl{A1}}
\BIC{$\varphi_{271}$}
\end{ptree}
\noindent$\Pi_{90}$:
\begin{ptree}
\AXC{$\stackrel{\Pi_{68}}{\vdots}$}\noLine
\UIC{$1_{\omega} = S(N) \to (\neg \forall M\,(\varphi_{114}) \leftrightarrow \neg \forall M\,(\varphi_{115}))$}
\AXC{$\stackrel{\Pi_{82}}{\vdots}$}\noLine
\UIC{$\mathrm{Seq}(a) \to \exists M\,(M \in \mathbb{R} \land 0 \leq M \land \forall n \in \mathbb{N}\,(\varphi_{28}))$}
\AXC{$\stackrel{\Pi_{89}}{\vdots}$}\noLine
\UIC{$\varphi_{271}$}
\BIC{$\varphi_{253}$}
\BIC{$\mathrm{Seq}(a) \to \forall N\,(N \in B_{8} \leftrightarrow N \in \omega \land \varphi_{51}) \to N \in B_{8} \to \varphi_{92}$}
\end{ptree}
\noindent$\Pi_{91}$:
\begin{ptree}
\AXC{$\varphi_{183} \to \forall M\,(\varphi_{183})$\lbl{A4}}
\AXC{$\stackrel{\Pi_{49}}{\vdots}$}\noLine
\UIC{$\varphi_{278}$}
\BIC{$\mathrm{Seq}(a) \to \forall N\,(N \in B_{8} \leftrightarrow N \in \omega \land \varphi_{51}) \to N \in B_{8} \to \varphi_{96}$}
\AXC{$\stackrel{\Pi_{90}}{\vdots}$}\noLine
\UIC{$\mathrm{Seq}(a) \to \forall N\,(N \in B_{8} \leftrightarrow N \in \omega \land \varphi_{51}) \to N \in B_{8} \to \varphi_{92}$}
\AXC{$\varphi_{277}$\lbl{A1}}
\BIC{$\varphi_{250}$}
\BIC{$\mathrm{Seq}(a) \to \forall N\,(N \in B_{8} \leftrightarrow N \in \omega \land \varphi_{51}) \to \varphi_{83}$}
\end{ptree}
\noindent$\Pi_{92}$:
\begin{ptree}
\AXC{$\forall x\,(\mathrm{Ind}(x) \leftrightarrow \emptyset \in x \land \forall y \in x\,(S(y) \in x))$\lbl{DefInd}}
\AXC{$\varphi_{129}$\lbl{A2}}
\BIC{$\mathrm{Ind}(\omega) \leftrightarrow \emptyset \in \omega \land \forall y \in \omega\,(S(y) \in \omega)$}
\AXC{$\mathrm{Ind}(\omega)$\lbl{補題(ωは帰納的)}}
\AXC{$\mathrm{Ind}(\omega) \to \varphi_{199}$\lbl{A1}}
\BIC{$\varphi_{199}$}
\BIC{$\mathrm{Seq}(a) \to \varphi_{156}$}
\end{ptree}
\noindent$\Pi_{93}$:
\begin{ptree}
\AXC{$\forall y \in \omega\,(S(y) \in \omega) \to N \in \omega \to S(N) \in \omega$\lbl{A2}}
\AXC{$\stackrel{\Pi_{92}}{\vdots}$}\noLine
\UIC{$\mathrm{Seq}(a) \to \varphi_{156}$}
\AXC{$(\mathrm{Seq}(a) \to \varphi_{156}) \to \mathrm{Seq}(a) \to \varphi_{147}$\lbl{A1}}
\BIC{$\mathrm{Seq}(a) \to \varphi_{147}$}
\AXC{$\varphi_{243}$\lbl{A1}}
\BIC{$(\forall y \in \omega\,(S(y) \in \omega) \to N \in \omega \to S(N) \in \omega) \to \varphi_{149}$}
\BIC{$\varphi_{149}$}
\end{ptree}
\noindent$\Pi_{94}$:
\begin{ptree}
\AXC{$\stackrel{\Pi_{5}}{\vdots}$}\noLine
\UIC{$\mathrm{Seq}(a) \to \forall N\,(N \in B_{8} \leftrightarrow N \in \omega \land \varphi_{51}) \to N \in B_{8} \to N \in \omega$}
\AXC{$\stackrel{\Pi_{93}}{\vdots}$}\noLine
\UIC{$\varphi_{149}$}
\AXC{$\varphi_{149} \to \varphi_{223}$\lbl{A1}}
\BIC{$\varphi_{223}$}
\BIC{$\varphi_{139}$}
\end{ptree}
\noindent$\Pi_{95}$:
\begin{ptree}
\AXC{$\stackrel{\Pi_{94}}{\vdots}$}\noLine
\UIC{$\varphi_{139}$}
\AXC{$\forall N\,(N \in B_{8} \leftrightarrow N \in \omega \land \varphi_{51}) \to \varphi_{110}$\lbl{A2}}
\AXC{$\varphi_{175}$\lbl{A1}}
\AXC{$\varphi_{175} \to \varphi_{252}$\lbl{A1}}
\BIC{$\varphi_{252}$}
\BIC{$\mathrm{Seq}(a) \to \forall N\,(N \in B_{8} \leftrightarrow N \in \omega \land \varphi_{51}) \to \varphi_{110}$}
\AXC{$\varphi_{282}$\lbl{A1}}
\BIC{$\varphi_{139} \to (\mathrm{Seq}(a) \to \forall N\,(N \in B_{8} \leftrightarrow N \in \omega \land \varphi_{51}) \to \varphi_{83}) \to \varphi_{134}$}
\BIC{$(\mathrm{Seq}(a) \to \forall N\,(N \in B_{8} \leftrightarrow N \in \omega \land \varphi_{51}) \to \varphi_{83}) \to \varphi_{134}$}
\end{ptree}
\noindent$\Pi_{96}$:
\begin{ptree}
\AXC{$\stackrel{\Pi_{91}}{\vdots}$}\noLine
\UIC{$\mathrm{Seq}(a) \to \forall N\,(N \in B_{8} \leftrightarrow N \in \omega \land \varphi_{51}) \to \varphi_{83}$}
\AXC{$\stackrel{\Pi_{95}}{\vdots}$}\noLine
\UIC{$(\mathrm{Seq}(a) \to \forall N\,(N \in B_{8} \leftrightarrow N \in \omega \land \varphi_{51}) \to \varphi_{83}) \to \varphi_{134}$}
\BIC{$\varphi_{134}$}
\AXC{$\varphi_{134} \to \varphi_{135}$\lbl{A1}}
\BIC{$\varphi_{135}$}
\UIC{$\forall N\,(\varphi_{135})$}
\end{ptree}
\noindent$\Pi_{97}$:
\begin{ptree}
\AXC{$\varphi_{190}$\lbl{A4}}
\AXC{$\forall N\,(\varphi_{135}) \to \varphi_{144}$\lbl{A3}}
\AXC{$\stackrel{\Pi_{96}}{\vdots}$}\noLine
\UIC{$\forall N\,(\varphi_{135})$}
\AXC{$\forall N\,(\varphi_{135}) \to (\forall N\,(\varphi_{135}) \to \varphi_{144}) \to \varphi_{190} \to \varphi_{140}$\lbl{A1}}
\BIC{$(\forall N\,(\varphi_{135}) \to \varphi_{144}) \to \varphi_{190} \to \varphi_{140}$}
\BIC{$\varphi_{190} \to \varphi_{140}$}
\BIC{$\varphi_{140}$}
\end{ptree}
\noindent$\Pi_{98}$:
\begin{ptree}
\AXC{$\forall n\,(n \in \mathbb{N} \leftrightarrow n \in \omega \land n \neq \emptyset)$\lbl{Def$\mathbb{N}$}}
\AXC{$\varphi_{105}$\lbl{A2}}
\BIC{$\emptyset \in \mathbb{N} \leftrightarrow \emptyset \in \omega \land \emptyset \neq \emptyset$}
\AXC{$(\emptyset \in \mathbb{N} \leftrightarrow \emptyset \in \omega \land \emptyset \neq \emptyset) \to \emptyset = \emptyset \to \mathrm{Seq}(a) \to \varphi_{76}$\lbl{A1}}
\BIC{$\emptyset = \emptyset \to \mathrm{Seq}(a) \to \varphi_{76}$}
\end{ptree}
\noindent$\Pi_{99}$:
\begin{ptree}
\AXC{$\forall N\,(N \in B_{8} \leftrightarrow N \in \omega \land \varphi_{51}) \to (\emptyset \in B_{8} \leftrightarrow \emptyset \in \omega \land \varphi_{76})$\lbl{A2}}
\AXC{$\varphi_{175}$\lbl{A1}}
\AXC{$\varphi_{175} \to \varphi_{255}$\lbl{A1}}
\BIC{$\varphi_{255}$}
\BIC{$\varphi_{203}$}
\AXC{$\varphi_{203} \to (\mathrm{Seq}(a) \to \varphi_{156}) \to \varphi_{192}$\lbl{A1}}
\BIC{$(\mathrm{Seq}(a) \to \varphi_{156}) \to \varphi_{192}$}
\end{ptree}
\noindent$\Pi_{100}$:
\begin{ptree}
\AXC{$x = x$\lbl{A5}}
\UIC{$\forall x\,(x = x)$}
\AXC{$\forall x\,(x = x) \to \emptyset = \emptyset$\lbl{A2}}
\BIC{$\emptyset = \emptyset$}
\AXC{$\stackrel{\Pi_{98}}{\vdots}$}\noLine
\UIC{$\emptyset = \emptyset \to \mathrm{Seq}(a) \to \varphi_{76}$}
\BIC{$\mathrm{Seq}(a) \to \varphi_{76}$}
\AXC{$\stackrel{\Pi_{92}}{\vdots}$}\noLine
\UIC{$\mathrm{Seq}(a) \to \varphi_{156}$}
\AXC{$\stackrel{\Pi_{99}}{\vdots}$}\noLine
\UIC{$(\mathrm{Seq}(a) \to \varphi_{156}) \to \varphi_{192}$}
\BIC{$\varphi_{192}$}
\BIC{$\mathrm{Seq}(a) \to \forall N\,(N \in B_{8} \leftrightarrow N \in \omega \land \varphi_{51}) \to \emptyset \in B_{8}$}
\end{ptree}
\noindent$\Pi_{101}$:
\begin{ptree}
\AXC{$\forall x\,(\mathrm{Ind}(x) \leftrightarrow \emptyset \in x \land \forall y \in x\,(S(y) \in x))$\lbl{DefInd}}
\AXC{$\varphi_{111}$\lbl{A2}}
\BIC{$\mathrm{Ind}(B_{8}) \leftrightarrow \emptyset \in B_{8} \land \forall y \in B_{8}\,(S(y) \in B_{8})$}
\AXC{$\varphi_{273}$\lbl{A1}}
\BIC{$\varphi_{264}$}
\end{ptree}
\noindent$\Pi_{102}$:
\begin{ptree}
\AXC{$\stackrel{\Pi_{97}}{\vdots}$}\noLine
\UIC{$\varphi_{140}$}
\AXC{$\stackrel{\Pi_{100}}{\vdots}$}\noLine
\UIC{$\mathrm{Seq}(a) \to \forall N\,(N \in B_{8} \leftrightarrow N \in \omega \land \varphi_{51}) \to \emptyset \in B_{8}$}
\AXC{$\stackrel{\Pi_{101}}{\vdots}$}\noLine
\UIC{$\varphi_{264}$}
\BIC{$\varphi_{140} \to \varphi_{196}$}
\BIC{$\varphi_{196}$}
\end{ptree}
\noindent$\Pi_{103}$:
\begin{ptree}
\AXC{$\stackrel{\Pi_{2}}{\vdots}$}\noLine
\UIC{$\forall N \in B_{8}\,(S(N) \in B_{8}) \to \forall y \in B_{8}\,(S(y) \in B_{8})$}
\AXC{$\stackrel{\Pi_{3}}{\vdots}$}\noLine
\UIC{$\forall y \in B_{8}\,(S(y) \in B_{8}) \to \forall N \in B_{8}\,(S(N) \in B_{8})$}
\AXC{$\stackrel{\Pi_{102}}{\vdots}$}\noLine
\UIC{$\varphi_{196}$}
\BIC{$\varphi_{161}$}
\BIC{$\mathrm{Seq}(a) \to \forall N\,(N \in B_{8} \leftrightarrow N \in \omega \land \varphi_{51}) \to \mathrm{Ind}(B_{8})$}
\end{ptree}
\noindent$\Pi_{104}$:
\begin{ptree}
\AXC{$\mathrm{Seq}(a) \to \forall N\,(N \in B_{8} \leftrightarrow N \in \omega \land \varphi_{51}) \to N \in \omega \to N \in \omega$\lbl{A1}}
\AXC{$\varphi_{237}$\lbl{A1}}
\BIC{$(N \in \omega \leftrightarrow \forall x\,(\mathrm{Ind}(x) \to N \in x)) \to \mathrm{Seq}(a) \to \varphi_{138}$}
\end{ptree}
\noindent$\Pi_{105}$:
\begin{ptree}
\AXC{$\forall n\,(n \in \omega \leftrightarrow \forall x\,(\mathrm{Ind}(x) \to n \in x))$\lbl{Def$\omega$}}
\AXC{$\varphi_{56}$\lbl{A2}}
\BIC{$N \in \omega \leftrightarrow \forall x\,(\mathrm{Ind}(x) \to N \in x)$}
\AXC{$\stackrel{\Pi_{104}}{\vdots}$}\noLine
\UIC{$(N \in \omega \leftrightarrow \forall x\,(\mathrm{Ind}(x) \to N \in x)) \to \mathrm{Seq}(a) \to \varphi_{138}$}
\BIC{$\mathrm{Seq}(a) \to \varphi_{138}$}
\end{ptree}
\noindent$\Pi_{106}$:
\begin{ptree}
\AXC{$\forall x\,(\mathrm{Ind}(x) \to N \in x) \to \mathrm{Ind}(B_{8}) \to N \in B_{8}$\lbl{A2}}
\AXC{$\stackrel{\Pi_{105}}{\vdots}$}\noLine
\UIC{$\mathrm{Seq}(a) \to \varphi_{138}$}
\AXC{$(\mathrm{Seq}(a) \to \varphi_{138}) \to \varphi_{163}$\lbl{A1}}
\BIC{$\varphi_{163}$}
\BIC{$\mathrm{Seq}(a) \to \varphi_{137}$}
\end{ptree}
\noindent$\Pi_{107}$:
\begin{ptree}
\AXC{$\stackrel{\Pi_{103}}{\vdots}$}\noLine
\UIC{$\mathrm{Seq}(a) \to \forall N\,(N \in B_{8} \leftrightarrow N \in \omega \land \varphi_{51}) \to \mathrm{Ind}(B_{8})$}
\AXC{$\stackrel{\Pi_{106}}{\vdots}$}\noLine
\UIC{$\mathrm{Seq}(a) \to \varphi_{137}$}
\AXC{$(\mathrm{Seq}(a) \to \varphi_{137}) \to \varphi_{270}$\lbl{A1}}
\BIC{$\varphi_{270}$}
\BIC{$\varphi_{171} \to \mathrm{Seq}(a) \to \forall N\,(N \in B_{8} \leftrightarrow N \in \omega \land \varphi_{51}) \to N \in \omega \to \varphi_{51}$}
\end{ptree}
\noindent$\Pi_{108}$:
\begin{ptree}
\AXC{$\varphi_{167} \to \varphi_{169}$\lbl{A3}}
\AXC{$\stackrel{\Pi_{1}}{\vdots}$}\noLine
\UIC{$\varphi_{171}$}
\AXC{$\stackrel{\Pi_{107}}{\vdots}$}\noLine
\UIC{$\varphi_{171} \to \mathrm{Seq}(a) \to \forall N\,(N \in B_{8} \leftrightarrow N \in \omega \land \varphi_{51}) \to N \in \omega \to \varphi_{51}$}
\BIC{$\mathrm{Seq}(a) \to \forall N\,(N \in B_{8} \leftrightarrow N \in \omega \land \varphi_{51}) \to N \in \omega \to \varphi_{51}$}
\AXC{$\varphi_{245}$\lbl{A1}}
\BIC{$\mathrm{Seq}(a) \land \forall N\,(N \in B_{8} \leftrightarrow N \in \omega \land \varphi_{51}) \to N \in \omega \to \varphi_{51}$}
\UIC{$\varphi_{167}$}
\AXC{$\varphi_{167} \to (\varphi_{167} \to \varphi_{169}) \to \varphi_{251}$\lbl{A1}}
\BIC{$(\varphi_{167} \to \varphi_{169}) \to \varphi_{251}$}
\BIC{$\varphi_{251}$}
\end{ptree}
\noindent$\Pi_{109}$:
\begin{ptree}
\AXC{$\varphi_{116}$\lbl{Sep}}
\AXC{$\varphi_{116} \to \forall A_{9}\,\exists B_{8}\,\forall N\,(N \in B_{8} \leftrightarrow N \in A_{9} \land \varphi_{51})$\lbl{A2}}
\BIC{$\forall A_{9}\,\exists B_{8}\,\forall N\,(N \in B_{8} \leftrightarrow N \in A_{9} \land \varphi_{51})$}
\AXC{$\varphi_{178}$\lbl{A2}}
\BIC{$\exists B_{8}\,\forall N\,(N \in B_{8} \leftrightarrow N \in \omega \land \varphi_{51})$}
\AXC{$\varphi_{185}$\lbl{A1}}
\BIC{$\mathrm{Seq}(a) \to \exists B_{8}\,\forall N\,(N \in B_{8} \leftrightarrow N \in \omega \land \varphi_{51})$}
\end{ptree}
\noindent$\Pi_{110}$:
\begin{ptree}
\AXC{$\stackrel{\Pi_{109}}{\vdots}$}\noLine
\UIC{$\mathrm{Seq}(a) \to \exists B_{8}\,\forall N\,(N \in B_{8} \leftrightarrow N \in \omega \land \varphi_{51})$}
\AXC{$(\mathrm{Seq}(a) \to \exists B_{8}\,\forall N\,(N \in B_{8} \leftrightarrow N \in \omega \land \varphi_{51})) \to \varphi_{286}$\lbl{A1}}
\BIC{$\varphi_{286}$}
\end{ptree}
\noindent$\Pi_{111}$:
\begin{ptree}
\AXC{$\varphi_{170}$\lbl{A4}}
\AXC{$\forall B_{8}\,(\varphi_{168}) \to \varphi_{174}$\lbl{A3}}
\AXC{$\varphi_{190}$\lbl{A4}}
\AXC{$\stackrel{\Pi_{108}}{\vdots}$}\noLine
\UIC{$\varphi_{251}$}
\BIC{$\mathrm{Seq}(a) \to \forall N\,(N \in B_{8} \leftrightarrow N \in \omega \land \varphi_{51}) \to \forall N \in \omega\,(\varphi_{51})$}
\AXC{$\varphi_{247}$\lbl{A1}}
\BIC{$\varphi_{168}$}
\UIC{$\forall B_{8}\,(\varphi_{168})$}
\AXC{$\stackrel{\Pi_{110}}{\vdots}$}\noLine
\UIC{$\varphi_{286}$}
\BIC{$(\forall B_{8}\,(\varphi_{168}) \to \varphi_{174}) \to \varphi_{170} \to \mathrm{Seq}(a) \to \forall N \in \omega\,(\varphi_{51})$}
\BIC{$\varphi_{170} \to \mathrm{Seq}(a) \to \forall N \in \omega\,(\varphi_{51})$}
\BIC{$\mathrm{Seq}(a) \to \forall N \in \omega\,(\varphi_{51})$}
\AXC{$\varphi_{244}$\lbl{A1}}
\BIC{$(\forall N \in \omega\,(\varphi_{51}) \to N \in \omega \to \varphi_{51}) \to \mathrm{Seq}(a) \to N \in \omega \to \varphi_{51}$}
\end{ptree}
\noindent$\Pi_{112}$:
\begin{ptree}
\AXC{$\forall N \in \omega\,(\varphi_{51}) \to N \in \omega \to \varphi_{51}$\lbl{A2}}
\AXC{$\stackrel{\Pi_{111}}{\vdots}$}\noLine
\UIC{$(\forall N \in \omega\,(\varphi_{51}) \to N \in \omega \to \varphi_{51}) \to \mathrm{Seq}(a) \to N \in \omega \to \varphi_{51}$}
\BIC{$\mathrm{Seq}(a) \to N \in \omega \to \varphi_{51}$}
\AXC{$(\mathrm{Seq}(a) \to N \in \omega \to \varphi_{51}) \to \varphi_{148}$\lbl{A1}}
\BIC{$\varphi_{148}$}
\end{ptree}
\noindent$\Pi_{113}$:
\begin{ptree}
\AXC{$\forall n\,(n \in \mathbb{N} \leftrightarrow n \in \omega \land n \neq \emptyset)$\lbl{Def$\mathbb{N}$}}
\AXC{$\varphi_{70}$\lbl{A2}}
\BIC{$N \in \mathbb{N} \leftrightarrow N \in \omega \land N \neq \emptyset$}
\AXC{$\mathrm{Seq}(a) \to N \in \mathbb{N} \to N \in \mathbb{N}$\lbl{A1}}
\AXC{$\varphi_{108}$\lbl{A1}}
\BIC{$(N \in \mathbb{N} \leftrightarrow N \in \omega \land N \neq \emptyset) \to \mathrm{Seq}(a) \to N \in \mathbb{N} \to N \in \omega$}
\BIC{$\mathrm{Seq}(a) \to N \in \mathbb{N} \to N \in \omega$}
\AXC{$\stackrel{\Pi_{112}}{\vdots}$}\noLine
\UIC{$\varphi_{148}$}
\BIC{$(\mathrm{Seq}(a) \to N \in \mathbb{N} \to N \in \mathbb{N}) \to \mathrm{Seq}(a) \to \varphi_{51}$}
\end{ptree}
\noindent$\Pi_{114}$:
\begin{ptree}
\AXC{$\mathrm{Seq}(a) \to N \in \mathbb{N} \to N \in \mathbb{N}$\lbl{A1}}
\AXC{$\stackrel{\Pi_{113}}{\vdots}$}\noLine
\UIC{$(\mathrm{Seq}(a) \to N \in \mathbb{N} \to N \in \mathbb{N}) \to \mathrm{Seq}(a) \to \varphi_{51}$}
\BIC{$\mathrm{Seq}(a) \to \varphi_{51}$}
\UIC{$\forall N\,(\mathrm{Seq}(a) \to \varphi_{51})$}
\AXC{$\forall N\,(\mathrm{Seq}(a) \to \varphi_{51}) \to \varphi_{229}$\lbl{A1}}
\BIC{$\varphi_{229}$}
\end{ptree}
\noindent 主木:
\begin{ptree}
\AXC{$\mathrm{Seq}(a) \to \forall N\,(\mathrm{Seq}(a))$\lbl{A4}}
\AXC{$\forall N\,(\mathrm{Seq}(a) \to \varphi_{51}) \to \forall N\,(\mathrm{Seq}(a)) \to \varphi_{58}$\lbl{A3}}
\AXC{$\stackrel{\Pi_{114}}{\vdots}$}\noLine
\UIC{$\varphi_{229}$}
\BIC{$(\mathrm{Seq}(a) \to \forall N\,(\mathrm{Seq}(a))) \to \mathrm{Seq}(a) \to \varphi_{58}$}
\BIC{$\mathrm{Seq}(a) \to \varphi_{58}$}
\UIC{$\forall a\,(\mathrm{Seq}(a) \to \varphi_{58})$}
\end{ptree}
\begin{abbrevlist}
\item $\varphi_{1} :\equiv \forall x\,\forall y\,(x = y \to y = x) \to \forall y\,(1_{\omega} = y \to y = 1_{\omega})$
\item $\varphi_{2} :\equiv \forall x\,\forall y\,(x = y \to y = x) \to \forall y\,(S(N) = y \to y = S(N))$
\item $\varphi_{3} :\equiv \neg n < 1_{\omega} \to n = 1_{\omega} \leftrightarrow \neg n < S(N) \to n = S(N)$
\item $\varphi_{4} :\equiv \forall y\,(1_{\omega} = y \to y = 1_{\omega}) \to 1_{\omega} = S(N) \to S(N) = 1_{\omega}$
\item $\varphi_{5} :\equiv \forall x\,\forall y\,(x \leq y \to y \leq x \to x = y) \to \forall y\,(n \leq y \to y \leq n \to n = y)$
\item $\varphi_{6} :\equiv \forall y\,(S(N) = y \to y = S(N)) \to S(N) = 1_{\omega} \to 1_{\omega} = S(N)$
\item $\varphi_{7} :\equiv \forall y\,(n \leq y \to y \leq n \to n = y) \to n \leq 1_{\omega} \to 1_{\omega} \leq n \to n = 1_{\omega}$
\item $\varphi_{8} :\equiv S(N) = S(\emptyset) \to S(\emptyset) = 1_{\omega} \to S(N) = 1_{\omega}$
\item $\varphi_{9} :\equiv \forall z\,\forall w\,(z = w \to n = z \to n = w) \to \forall w\,(S(N) = w \to n = S(N) \to n = w)$
\item $\varphi_{10} :\equiv \forall z\,\forall w\,(z = w \to n = z \to n = w) \to \forall w\,(1_{\omega} = w \to n = 1_{\omega} \to n = w)$
\item $\varphi_{11} :\equiv (1_{\omega} = S(N) \to S(N) = 1_{\omega}) \to 1_{\omega} = S(N) \to (n = 1_{\omega} \leftrightarrow n = S(N))$
\item $\varphi_{12} :\equiv \forall y\,(\forall N \in B_{8}\,(S(N) \in B_{8}) \to y \in B_{8} \to S(y) \in B_{8})$
\item $\varphi_{13} :\equiv \forall N\,(\forall y \in B_{8}\,(S(y) \in B_{8}) \to N \in B_{8} \to S(N) \in B_{8})$
\item $\varphi_{14} :\equiv \forall w\,(S(N) = w \to n = S(N) \to n = w) \to S(N) = 1_{\omega} \to n = S(N) \to n = 1_{\omega}$
\item $\varphi_{15} :\equiv \forall w\,(1_{\omega} = w \to n = 1_{\omega} \to n = w) \to 1_{\omega} = S(N) \to n = 1_{\omega} \to n = S(N)$
\item $\varphi_{16} :\equiv \forall x\,\forall y\,\forall z\,(x = y \to y = z \to x = z) \to \forall y\,\forall z\,(S(N) = y \to y = z \to S(N) = z)$
\item $\varphi_{17} :\equiv (1_{\omega} = S(\emptyset) \to S(\emptyset) = 1_{\omega}) \to S(\emptyset) = 1_{\omega}$
\item $\varphi_{18} :\equiv \forall m\,\forall n\,(m \in \omega \to n \in \omega \to m \leq S(n) \to m \leq n \lor m = S(n))$
\item $\varphi_{19} :\equiv \forall N \in B_{8}\,(S(N) \in B_{8}) \to \forall y\,\forall N \in B_{8}\,(S(N) \in B_{8})$
\item $\varphi_{20} :\equiv \forall y \in B_{8}\,(S(y) \in B_{8}) \to \forall N\,\forall y \in B_{8}\,(S(y) \in B_{8})$
\item $\varphi_{21} :\equiv \forall y\,(1_{\omega} = y \to y = 1_{\omega}) \to 1_{\omega} = S(\emptyset) \to S(\emptyset) = 1_{\omega}$
\item $\varphi_{22} :\equiv \forall y\,\forall N \in B_{8}\,(S(N) \in B_{8}) \to \forall y \in B_{8}\,(S(y) \in B_{8})$
\item $\varphi_{23} :\equiv \forall N\,\forall y \in B_{8}\,(S(y) \in B_{8}) \to \forall N \in B_{8}\,(S(N) \in B_{8})$
\item $\varphi_{24} :\equiv v_{686} \in \omega \to v_{685} \leq S(v_{686}) \to v_{685} \leq v_{686} \lor v_{685} = S(v_{686})$
\item $\varphi_{25} :\equiv (1_{\omega} = S(N) \to S(N) = 1_{\omega}) \to 1_{\omega} = S(N) \to (n < 1_{\omega} \leftrightarrow n < S(N))$
\item $\varphi_{26} :\equiv (\mathrm{Seq}(a) \to a_{1_{\omega}} \in \mathbb{R}) \to \mathrm{Seq}(a) \to 0 \leq \lvert a_{1_{\omega}} \rvert$
\item $\varphi_{27} :\equiv \forall M\,(\neg (M \in \mathbb{R} \land 0 \leq M \land \forall n \in \mathbb{N}\,(n \leq N \to \lvert a_{n} \rvert \leq M)))$
\item $\varphi_{28} :\equiv (\neg n < 1_{\omega} \to n = 1_{\omega}) \to \neg \lvert a_{n} \rvert < M \to \lvert a_{n} \rvert = M$
\item $\varphi_{29} :\equiv (\neg n < S(N) \to n = S(N)) \to \neg \lvert a_{n} \rvert < M \to \lvert a_{n} \rvert = M$
\item $\varphi_{30} :\equiv \neg (M \in \mathbb{R} \to \neg (\neg 0 < M \to 0 = M)) \leftrightarrow \neg (M \in \mathbb{R} \to \neg (\neg 0 < M \to 0 = M))$
\item $\varphi_{31} :\equiv \forall z\,\forall w\,(z = w \to n < z \to n < w) \to \forall w\,(S(N) = w \to n < S(N) \to n < w)$
\item $\varphi_{32} :\equiv \forall z\,\forall w\,(z = w \to n < z \to n < w) \to \forall w\,(1_{\omega} = w \to n < 1_{\omega} \to n < w)$
\item $\varphi_{33} :\equiv \mathrm{Seq}(a) \to n \in \mathbb{N} \to n \leq 1_{\omega} \to \lvert a_{n} \rvert \leq \lvert a_{1_{\omega}} \rvert$
\item $\varphi_{34} :\equiv n \in \mathbb{N} \to n \leq S(N) \to n \leq N \to n \in \mathbb{N} \to n \leq N \to \lvert a_{n} \rvert \leq M$
\item $\varphi_{35} :\equiv (\mathrm{Seq}(a) \to \mathrm{Seq}(a)) \to 1_{\omega} \in \mathbb{N} \to \mathrm{Seq}(a) \to a_{1_{\omega}} \in \mathbb{R}$
\item $\varphi_{36} :\equiv \mathrm{Seq}(a) \to \forall n \in \mathbb{N}\,(n \leq 1_{\omega} \to \lvert a_{n} \rvert \leq \lvert a_{1_{\omega}} \rvert)$
\item $\varphi_{37} :\equiv n \in \mathbb{N} \to n \leq S(N) \to n \leq N \to \forall n \in \mathbb{N}\,(n \leq N \to \lvert a_{n} \rvert \leq M)$
\item $\varphi_{38} :\equiv \forall n\,(v_{685} \in \omega \to n \in \omega \to v_{685} \leq S(n) \to v_{685} \leq n \lor v_{685} = S(n))$
\item $\varphi_{39} :\equiv n \in \mathbb{N} \to n \leq S(N) \to n = S(N) \to \lvert a_{n} \rvert \leq M + \lvert a_{S(N)} \rvert$
\item $\varphi_{40} :\equiv \forall w\,(S(N) = w \to n < S(N) \to n < w) \to S(N) = 1_{\omega} \to n < S(N) \to n < 1_{\omega}$
\item $\varphi_{41} :\equiv \forall w\,(1_{\omega} = w \to n < 1_{\omega} \to n < w) \to 1_{\omega} = S(N) \to n < 1_{\omega} \to n < S(N)$
\item $\varphi_{42} :\equiv (\mathrm{Seq}(a) \to a_{1_{\omega}} \in \mathbb{R}) \to \mathrm{Seq}(a) \to \lvert a_{1_{\omega}} \rvert \in \mathbb{R}$
\item $\varphi_{43} :\equiv \forall z\,\forall w\,(z = w \to S(N) = S(z) \to S(N) = S(w)) \to \forall w\,(N = w \to S(N) = S(N) \to S(N) = S(w))$
\item $\varphi_{44} :\equiv \forall n\,(\mathrm{Seq}(a)) \to \forall n \in \mathbb{N}\,(n \leq 1_{\omega} \to \lvert a_{n} \rvert \leq \lvert a_{1_{\omega}} \rvert)$
\item $\varphi_{45} :\equiv \forall v_{686}\,(n \in \omega \to v_{686} \in \omega \to n \leq S(v_{686}) \to n \leq v_{686} \lor n = S(v_{686}))$
\item $\varphi_{46} :\equiv \forall x \in \mathbb{R}\,(0 \leq \lvert x \rvert) \to a_{1_{\omega}} \in \mathbb{R} \to 0 \leq \lvert a_{1_{\omega}} \rvert$
\item $\varphi_{47} :\equiv M \in \mathbb{R} \land 0 \leq M \land \forall n \in \mathbb{N}\,(n \leq N \to \lvert a_{n} \rvert \leq M) \to M \in \mathbb{R}$
\item $\varphi_{48} :\equiv n \in \mathbb{N} \to n \leq S(N) \to n = S(N) \to \lvert a_{S(N)} \rvert \leq M + \lvert a_{S(N)} \rvert$
\item $\varphi_{49} :\equiv n = 1_{\omega} \to \lvert a_{n} \rvert = \lvert a_{n} \rvert \to \lvert a_{n} \rvert = \lvert a_{1_{\omega}} \rvert$
\item $\varphi_{50} :\equiv (\mathrm{Seq}(a) \to n \in \mathbb{N} \to n \in \mathbb{N}) \to \mathrm{Seq}(a) \to n \in \mathbb{N} \to n \leq 1_{\omega} \to 1_{\omega} \leq n$
\item $\varphi_{51} :\equiv N \in \mathbb{N} \to \exists M\,(M \in \mathbb{R} \land 0 \leq M \land \forall n \in \mathbb{N}\,(n \leq N \to \lvert a_{n} \rvert \leq M))$
\item $\varphi_{52} :\equiv M \in \mathbb{R} \land 0 \leq M \land \forall n \in \mathbb{N}\,(n \leq N \to \lvert a_{n} \rvert \leq M) \to S(N) \in \mathbb{N}$
\item $\varphi_{53} :\equiv (1_{\omega} = S(N) \to (n < 1_{\omega} \leftrightarrow n < S(N))) \to 1_{\omega} = S(N) \to (\neg n < 1_{\omega} \leftrightarrow \neg n < S(N))$
\item $\varphi_{54} :\equiv \forall n \in \mathbb{N}\,(n \leq N \to \lvert a_{n} \rvert \leq M) \to n \in \mathbb{N} \to n \leq N \to \lvert a_{n} \rvert \leq M$
\item $\varphi_{55} :\equiv \forall x\,\forall y\,\forall z\,(x \leq y \to y \leq z \to x \leq z) \to \forall y\,\forall z\,(\lvert a_{n} \rvert \leq y \to y \leq z \to \lvert a_{n} \rvert \leq z)$
\item $\varphi_{56} :\equiv \forall n\,(n \in \omega \leftrightarrow \forall x\,(\mathrm{Ind}(x) \to n \in x)) \to (N \in \omega \leftrightarrow \forall x\,(\mathrm{Ind}(x) \to N \in x))$
\item $\varphi_{57} :\equiv N \neq \emptyset \to \exists M\,(M \in \mathbb{R} \land 0 \leq M \land \forall n \in \mathbb{N}\,(n \leq N \to \lvert a_{n} \rvert \leq M))$
\item $\varphi_{58} :\equiv \forall N \in \mathbb{N}\,\exists M\,(M \in \mathbb{R} \land 0 \leq M \land \forall n \in \mathbb{N}\,(n \leq N \to \lvert a_{n} \rvert \leq M))$
\item $\varphi_{59} :\equiv \forall y\,\forall z\,(S(N) = y \to y = z \to S(N) = z) \to \forall z\,(S(N) = S(\emptyset) \to S(\emptyset) = z \to S(N) = z)$
\item $\varphi_{60} :\equiv \forall w\,(N = w \to S(N) = S(N) \to S(N) = S(w)) \to N = \emptyset \to S(N) = S(N) \to S(N) = S(\emptyset)$
\item $\varphi_{61} :\equiv n = 1_{\omega} \to \lvert a_{n} \rvert < \lvert a_{n} \rvert \to \lvert a_{n} \rvert < \lvert a_{1_{\omega}} \rvert$
\item $\varphi_{62} :\equiv (\mathrm{Seq}(a) \to n \in \mathbb{N} \to n \leq 1_{\omega} \to 1_{\omega} \leq n) \to \mathrm{Seq}(a) \to n \in \mathbb{N} \to n \leq 1_{\omega} \to n = 1_{\omega}$
\item $\varphi_{63} :\equiv N \neq \emptyset \to M \in \mathbb{R} \land 0 \leq M \land \forall n \in \mathbb{N}\,(n \leq N \to \lvert a_{n} \rvert \leq M) \to 0 \leq M$
\item $\varphi_{64} :\equiv M \in \mathbb{R} \land 0 \leq M \land \forall n \in \mathbb{N}\,(n \leq N \to \lvert a_{n} \rvert \leq M) \to 0 \leq \lvert a_{S(N)} \rvert$
\item $\varphi_{65} :\equiv M \in \mathbb{R} \land 0 \leq M \land \forall n \in \mathbb{N}\,(n \leq N \to \lvert a_{n} \rvert \leq M) \to a_{S(N)} \in \mathbb{R}$
\item $\varphi_{66} :\equiv \lvert a_{S(N)} \rvert \in \mathbb{R} \to 0 \leq M \to 0 \leq \lvert a_{S(N)} \rvert \to 0 \leq M + \lvert a_{S(N)} \rvert$
\item $\varphi_{67} :\equiv z = w \to \lvert a_{z} \rvert = M + \lvert a_{S(N)} \rvert \to \lvert a_{w} \rvert = M + \lvert a_{S(N)} \rvert$
\item $\varphi_{68} :\equiv \lvert a_{n} \rvert \leq M \to M \leq M + \lvert a_{S(N)} \rvert \to \lvert a_{n} \rvert \leq M + \lvert a_{S(N)} \rvert$
\item $\varphi_{69} :\equiv M \in \mathbb{R} \land 0 \leq M \land \forall n \in \mathbb{N}\,(n \leq N \to \lvert a_{n} \rvert \leq M) \to n \in \mathbb{N} \to n \in \mathbb{N}$
\item $\varphi_{70} :\equiv \forall n\,(n \in \mathbb{N} \leftrightarrow n \in \omega \land n \neq \emptyset) \to (N \in \mathbb{N} \leftrightarrow N \in \omega \land N \neq \emptyset)$
\item $\varphi_{71} :\equiv M \in \mathbb{R} \land 0 \leq M \land \forall n \in \mathbb{N}\,(n \leq N \to \lvert a_{n} \rvert \leq M) \to 0 \leq M + \lvert a_{S(N)} \rvert$
\item $\varphi_{72} :\equiv \forall x \in \mathbb{R}\,(\lvert x \rvert \in \mathbb{R}) \to a_{S(N)} \in \mathbb{R} \to \lvert a_{S(N)} \rvert \in \mathbb{R}$
\item $\varphi_{73} :\equiv \forall x \in \mathbb{R}\,(\lvert x \rvert \in \mathbb{R}) \to a_{1_{\omega}} \in \mathbb{R} \to \lvert a_{1_{\omega}} \rvert \in \mathbb{R}$
\item $\varphi_{74} :\equiv (1_{\omega} = S(N) \to (n = 1_{\omega} \leftrightarrow n = S(N))) \to 1_{\omega} = S(N) \to \varphi_{3}$
\item $\varphi_{75} :\equiv M \in \mathbb{R} \land 0 \leq M \land \forall n \in \mathbb{N}\,(n \leq N \to \lvert a_{n} \rvert \leq M) \to \lvert a_{S(N)} \rvert \in \mathbb{R}$
\item $\varphi_{76} :\equiv \emptyset \in \mathbb{N} \to \exists M\,(M \in \mathbb{R} \land 0 \leq M \land \forall n \in \mathbb{N}\,(n \leq \emptyset \to \lvert a_{n} \rvert \leq M))$
\item $\varphi_{77} :\equiv S(N) = n \to \lvert a_{S(N)} \rvert = M + \lvert a_{S(N)} \rvert \to \lvert a_{n} \rvert = M + \lvert a_{S(N)} \rvert$
\item $\varphi_{78} :\equiv z = w \to \lvert a_{z} \rvert < M + \lvert a_{S(N)} \rvert \to \lvert a_{w} \rvert < M + \lvert a_{S(N)} \rvert$
\item $\varphi_{79} :\equiv \lvert a_{S(N)} \rvert \in \mathbb{R} \to M \in \mathbb{R} \to 0 \leq M \to \lvert a_{S(N)} \rvert \leq M + \lvert a_{S(N)} \rvert$
\item $\varphi_{80} :\equiv M \in \mathbb{R} \to \lvert a_{S(N)} \rvert \in \mathbb{R} \to 0 \leq \lvert a_{S(N)} \rvert \to M \leq M + \lvert a_{S(N)} \rvert$
\item $\varphi_{81} :\equiv M \in \mathbb{R} \land 0 \leq M \land \forall n \in \mathbb{N}\,(n \leq N \to \lvert a_{n} \rvert \leq M) \to n \in \mathbb{N} \to n \leq S(N) \to n \leq S(N)$
\item $\varphi_{82} :\equiv \lvert a_{n} \rvert = \lvert a_{n} \rvert \to \mathrm{Seq}(a) \to n \in \mathbb{N} \to n \leq 1_{\omega} \to \lvert a_{n} \rvert \leq \lvert a_{n} \rvert$
\item $\varphi_{83} :\equiv N \in B_{8} \to S(N) \in \mathbb{N} \to \exists M\,(M \in \mathbb{R} \land 0 \leq M \land \forall n \in \mathbb{N}\,(\varphi_{29}))$
\item $\varphi_{84} :\equiv \forall w\,(S(N) = w \to \lvert a_{S(N)} \rvert = M + \lvert a_{S(N)} \rvert \to \lvert a_{w} \rvert = M + \lvert a_{S(N)} \rvert)$
\item $\varphi_{85} :\equiv \forall a\,\forall n\,(\mathrm{Seq}(a) \to n \in \mathbb{N} \to a_{n} \in \mathbb{R}) \to \forall n\,(\mathrm{Seq}(a) \to n \in \mathbb{N} \to a_{n} \in \mathbb{R})$
\item $\varphi_{86} :\equiv \forall n\,(\mathrm{Seq}(a) \to n \in \mathbb{N} \to a_{n} \in \mathbb{R}) \to \mathrm{Seq}(a) \to S(N) \in \mathbb{N} \to a_{S(N)} \in \mathbb{R}$
\item $\varphi_{87} :\equiv \forall y\,(\lvert a_{S(N)} \rvert \in \mathbb{R} \to y \in \mathbb{R} \to 0 \leq y \to \lvert a_{S(N)} \rvert \leq y + \lvert a_{S(N)} \rvert)$
\item $\varphi_{88} :\equiv \forall n\,(\mathrm{Seq}(a) \to n \in \mathbb{N} \to a_{n} \in \mathbb{R}) \to \mathrm{Seq}(a) \to 1_{\omega} \in \mathbb{N} \to a_{1_{\omega}} \in \mathbb{R}$
\item $\varphi_{89} :\equiv M \in \mathbb{R} \land 0 \leq M \land \forall n \in \mathbb{N}\,(n \leq N \to \lvert a_{n} \rvert \leq M) \to n \in \mathbb{N} \to n \leq S(N) \to n \leq N \to n \leq N$
\item $\varphi_{90} :\equiv M \in \mathbb{R} \land 0 \leq M \land \forall n \in \mathbb{N}\,(n \leq N \to \lvert a_{n} \rvert \leq M) \to M + \lvert a_{S(N)} \rvert \in \mathbb{R}$
\item $\varphi_{91} :\equiv 1_{\omega} = S(N) \to (\neg \lvert a_{n} \rvert < M \to \lvert a_{n} \rvert = M \leftrightarrow \neg \lvert a_{n} \rvert < M \to \lvert a_{n} \rvert = M)$
\item $\varphi_{92} :\equiv S(N) \in \mathbb{N} \to N = \emptyset \to \exists M\,(M \in \mathbb{R} \land 0 \leq M \land \forall n \in \mathbb{N}\,(\varphi_{29}))$
\item $\varphi_{93} :\equiv \forall y\,\forall z\,(\lvert a_{n} \rvert \leq y \to y \leq z \to \lvert a_{n} \rvert \leq z) \to \forall z\,(\lvert a_{n} \rvert \leq M \to M \leq z \to \lvert a_{n} \rvert \leq z)$
\item $\varphi_{94} :\equiv M \in \mathbb{R} \land 0 \leq M \land \forall n \in \mathbb{N}\,(n \leq N \to \lvert a_{n} \rvert \leq M) \to n \in \mathbb{N} \to n \leq S(N) \to N \in \omega$
\item $\varphi_{95} :\equiv M \in \mathbb{R} \land 0 \leq M \land \forall n \in \mathbb{N}\,(n \leq N \to \lvert a_{n} \rvert \leq M) \to n \in \mathbb{N} \to n \leq S(N) \to n \in \omega$
\item $\varphi_{96} :\equiv S(N) \in \mathbb{N} \to N \neq \emptyset \to \exists M\,(M \in \mathbb{R} \land 0 \leq M \land \forall n \in \mathbb{N}\,(\varphi_{29}))$
\item $\varphi_{97} :\equiv S(N) = n \to \lvert a_{S(N)} \rvert < M + \lvert a_{S(N)} \rvert \to \lvert a_{n} \rvert < M + \lvert a_{S(N)} \rvert$
\item $\varphi_{98} :\equiv M \in \mathbb{R} \land 0 \leq M \land \forall n \in \mathbb{N}\,(n \leq N \to \lvert a_{n} \rvert \leq M) \to n \in \mathbb{N} \to n \leq S(N) \to n \leq N \lor n = S(N)$
\item $\varphi_{99} :\equiv \forall x\,\forall y\,(x \in \mathbb{R} \to y \in \mathbb{R} \to 0 \leq y \to x \leq x + y) \to \forall y\,(M \in \mathbb{R} \to y \in \mathbb{R} \to 0 \leq y \to M \leq M + y)$
\item $\varphi_{100} :\equiv \forall w\,(S(N) = w \to \lvert a_{S(N)} \rvert < M + \lvert a_{S(N)} \rvert \to \lvert a_{w} \rvert < M + \lvert a_{S(N)} \rvert)$
\item $\varphi_{101} :\equiv \forall x\,\forall y\,(x \in \mathbb{R} \to y \in \mathbb{R} \to x + y \in \mathbb{R}) \to \forall y\,(M \in \mathbb{R} \to y \in \mathbb{R} \to M + y \in \mathbb{R})$
\item $\varphi_{102} :\equiv M \in \mathbb{R} \land 0 \leq M \land \forall n \in \mathbb{N}\,(n \leq N \to \lvert a_{n} \rvert \leq M) \to n \in \mathbb{N} \to n \leq S(N) \to n = S(N) \to n = S(N)$
\item $\varphi_{103} :\equiv M \in \mathbb{R} \land 0 \leq M \land \forall n \in \mathbb{N}\,(n \leq N \to \lvert a_{n} \rvert \leq M) \to n \in \mathbb{N} \to n \leq S(N) \to n = S(N) \to S(N) = n$
\item $\varphi_{104} :\equiv M \in \mathbb{R} \land 0 \leq M \land \forall n \in \mathbb{N}\,(n \leq N \to \lvert a_{n} \rvert \leq M) \to n \in \mathbb{N} \to n \leq S(N) \to n \leq N \to \lvert a_{n} \rvert \leq M$
\item $\varphi_{105} :\equiv \forall n\,(n \in \mathbb{N} \leftrightarrow n \in \omega \land n \neq \emptyset) \to (\emptyset \in \mathbb{N} \leftrightarrow \emptyset \in \omega \land \emptyset \neq \emptyset)$
\item $\varphi_{106} :\equiv \lvert a_{1_{\omega}} \rvert \in \mathbb{R} \land 0 \leq \lvert a_{1_{\omega}} \rvert \land \forall n \in \mathbb{N}\,(n \leq 1_{\omega} \to \lvert a_{n} \rvert \leq \lvert a_{1_{\omega}} \rvert)$
\item $\varphi_{107} :\equiv \forall x\,\forall y\,(x \in \mathbb{R} \to y \in \mathbb{R} \to 0 \leq x \to 0 \leq y \to 0 \leq x + y) \to \forall y\,(M \in \mathbb{R} \to y \in \mathbb{R} \to 0 \leq M \to 0 \leq y \to 0 \leq M + y)$
\item $\varphi_{108} :\equiv (\mathrm{Seq}(a) \to N \in \mathbb{N} \to N \in \mathbb{N}) \to (N \in \mathbb{N} \leftrightarrow N \in \omega \land N \neq \emptyset) \to \mathrm{Seq}(a) \to N \in \mathbb{N} \to N \in \omega$
\item $\varphi_{109} :\equiv M \in \mathbb{R} \land 0 \leq M \land \forall n \in \mathbb{N}\,(n \leq N \to \lvert a_{n} \rvert \leq M) \to n \in \mathbb{N} \to n \leq S(N) \to n \leq N \to a_{S(N)} \in \mathbb{R}$
\item $\varphi_{110} :\equiv S(N) \in B_{8} \leftrightarrow S(N) \in \omega \land (S(N) \in \mathbb{N} \to \exists M\,(M \in \mathbb{R} \land 0 \leq M \land \forall n \in \mathbb{N}\,(\varphi_{29})))$
\item $\varphi_{111} :\equiv \forall x\,(\mathrm{Ind}(x) \leftrightarrow \emptyset \in x \land \forall y \in x\,(S(y) \in x)) \to (\mathrm{Ind}(B_{8}) \leftrightarrow \emptyset \in B_{8} \land \forall y \in B_{8}\,(S(y) \in B_{8}))$
\item $\varphi_{112} :\equiv M \in \mathbb{R} \land 0 \leq M \land \forall n \in \mathbb{N}\,(n \leq N \to \lvert a_{n} \rvert \leq M) \to n \in \mathbb{N} \to n \leq S(N) \to n = S(N) \to a_{S(N)} \in \mathbb{R}$
\item $\varphi_{113} :\equiv M \in \mathbb{R} \land 0 \leq M \land \forall n \in \mathbb{N}\,(n \leq N \to \lvert a_{n} \rvert \leq M) \to n \in \mathbb{N} \to n \leq S(N) \to n \leq N \to M \leq M + \lvert a_{S(N)} \rvert$
\item $\varphi_{114} :\equiv \neg \neg (\neg (M \in \mathbb{R} \to \neg (\neg 0 < M \to 0 = M)) \to \neg \forall n \in \mathbb{N}\,(\varphi_{28}))$
\item $\varphi_{115} :\equiv \neg \neg (\neg (M \in \mathbb{R} \to \neg (\neg 0 < M \to 0 = M)) \to \neg \forall n \in \mathbb{N}\,(\varphi_{29}))$
\item $\varphi_{116} :\equiv \forall z_{10}\,\forall A_{9}\,\exists B_{8}\,\forall N\,(N \in B_{8} \leftrightarrow N \in A_{9} \land (N \in \mathbb{N} \to \exists M\,(M \in \mathbb{R} \land 0 \leq M \land \forall n \in \mathbb{N}\,(n \leq N \to \lvert z_{10}(n) \rvert \leq M))))$
\item $\varphi_{117} :\equiv M \in \mathbb{R} \land 0 \leq M \land \forall n \in \mathbb{N}\,(n \leq N \to \lvert a_{n} \rvert \leq M) \to n \in \mathbb{N} \to n \leq S(N) \to n \leq N \to \lvert a_{S(N)} \rvert \in \mathbb{R}$
\item $\varphi_{118} :\equiv (1_{\omega} = S(N) \to n = 1_{\omega} \to n = S(N)) \to (S(N) = 1_{\omega} \to n = S(N) \to n = 1_{\omega}) \to \varphi_{11}$
\item $\varphi_{119} :\equiv \varphi_{19} \to \forall N \in B_{8}\,(S(N) \in B_{8}) \to \forall y \in B_{8}\,(S(y) \in B_{8})$
\item $\varphi_{120} :\equiv \varphi_{20} \to \forall y \in B_{8}\,(S(y) \in B_{8}) \to \forall N \in B_{8}\,(S(N) \in B_{8})$
\item $\varphi_{121} :\equiv M \in \mathbb{R} \land 0 \leq M \land \forall n \in \mathbb{N}\,(n \leq N \to \lvert a_{n} \rvert \leq M) \to M \in \mathbb{R} \land 0 \leq M \land \forall n \in \mathbb{N}\,(n \leq N \to \lvert a_{n} \rvert \leq M)$
\item $\varphi_{122} :\equiv M \in \mathbb{R} \land 0 \leq M \land \forall n \in \mathbb{N}\,(n \leq N \to \lvert a_{n} \rvert \leq M) \to n \in \mathbb{N} \to n \leq S(N) \to n = S(N) \to \lvert a_{S(N)} \rvert \in \mathbb{R}$
\item $\varphi_{123} :\equiv (\mathrm{Seq}(a) \to n \in \mathbb{N} \to n \leq 1_{\omega} \to n \leq 1_{\omega}) \to \varphi_{62}$
\item $\varphi_{124} :\equiv M \in \mathbb{R} \land 0 \leq M \land \forall n \in \mathbb{N}\,(n \leq N \to \lvert a_{n} \rvert \leq M) \to n \in \mathbb{N} \to n \leq S(N) \to \lvert a_{n} \rvert \leq M + \lvert a_{S(N)} \rvert$
\item $\varphi_{125} :\equiv M \in \mathbb{R} \land 0 \leq M \land \forall n \in \mathbb{N}\,(n \leq N \to \lvert a_{n} \rvert \leq M) \to \forall n \in \mathbb{N}\,(n \leq S(N) \to \lvert a_{n} \rvert \leq M + \lvert a_{S(N)} \rvert)$
\item $\varphi_{126} :\equiv \forall y\,(M \in \mathbb{R} \to y \in \mathbb{R} \to M + y \in \mathbb{R}) \to M \in \mathbb{R} \to \lvert a_{S(N)} \rvert \in \mathbb{R} \to M + \lvert a_{S(N)} \rvert \in \mathbb{R}$
\item $\varphi_{127} :\equiv M + \lvert a_{S(N)} \rvert \in \mathbb{R} \land 0 \leq M + \lvert a_{S(N)} \rvert \land \forall n \in \mathbb{N}\,(n \leq S(N) \to \lvert a_{n} \rvert \leq M + \lvert a_{S(N)} \rvert)$
\item $\varphi_{128} :\equiv M \in \mathbb{R} \land 0 \leq M \land \forall n \in \mathbb{N}\,(n \leq N \to \lvert a_{n} \rvert \leq M) \to \exists M\,(M \in \mathbb{R} \land 0 \leq M \land \forall n \in \mathbb{N}\,(\varphi_{29}))$
\item $\varphi_{129} :\equiv \forall x\,(\mathrm{Ind}(x) \leftrightarrow \emptyset \in x \land \forall y \in x\,(S(y) \in x)) \to (\mathrm{Ind}(\omega) \leftrightarrow \emptyset \in \omega \land \forall y \in \omega\,(S(y) \in \omega))$
\item $\varphi_{130} :\equiv M \in \mathbb{R} \land 0 \leq M \land \forall n \in \mathbb{N}\,(n \leq N \to \lvert a_{n} \rvert \leq M) \to n \in \mathbb{N} \to n \leq S(N) \to n \leq N \to \lvert a_{n} \rvert \leq M + \lvert a_{S(N)} \rvert$
\item $\varphi_{131} :\equiv (1_{\omega} = S(N) \to (\neg n < 1_{\omega} \leftrightarrow \neg n < S(N))) \to \varphi_{74}$
\item $\varphi_{132} :\equiv \mathrm{Seq}(a) \to \forall N\,(N \in B_{8} \leftrightarrow N \in \omega \land \varphi_{51}) \to N \in B_{8} \to N \in B_{8}$
\item $\varphi_{133} :\equiv (1_{\omega} = S(N) \to n < 1_{\omega} \to n < S(N)) \to (S(N) = 1_{\omega} \to n < S(N) \to n < 1_{\omega}) \to \varphi_{25}$
\item $\varphi_{134} :\equiv \mathrm{Seq}(a) \to \forall N\,(N \in B_{8} \leftrightarrow N \in \omega \land \varphi_{51}) \to N \in B_{8} \to S(N) \in B_{8}$
\item $\varphi_{135} :\equiv \mathrm{Seq}(a) \land \forall N\,(N \in B_{8} \leftrightarrow N \in \omega \land \varphi_{51}) \to N \in B_{8} \to S(N) \in B_{8}$
\item $\varphi_{136} :\equiv N \neq \emptyset \to M \in \mathbb{R} \land 0 \leq M \land \forall n \in \mathbb{N}\,(n \leq N \to \lvert a_{n} \rvert \leq M) \to \varphi_{34}$
\item $\varphi_{137} :\equiv \forall N\,(N \in B_{8} \leftrightarrow N \in \omega \land \varphi_{51}) \to N \in \omega \to \mathrm{Ind}(B_{8}) \to N \in B_{8}$
\item $\varphi_{138} :\equiv \forall N\,(N \in B_{8} \leftrightarrow N \in \omega \land \varphi_{51}) \to N \in \omega \to \forall x\,(\mathrm{Ind}(x) \to N \in x)$
\item $\varphi_{139} :\equiv \mathrm{Seq}(a) \to \forall N\,(N \in B_{8} \leftrightarrow N \in \omega \land \varphi_{51}) \to N \in B_{8} \to S(N) \in \omega$
\item $\varphi_{140} :\equiv \mathrm{Seq}(a) \to \forall N\,(N \in B_{8} \leftrightarrow N \in \omega \land \varphi_{51}) \to \forall N \in B_{8}\,(S(N) \in B_{8})$
\item $\varphi_{141} :\equiv N \neq \emptyset \to M \in \mathbb{R} \land 0 \leq M \land \forall n \in \mathbb{N}\,(n \leq N \to \lvert a_{n} \rvert \leq M) \to \varphi_{37}$
\item $\varphi_{142} :\equiv N \neq \emptyset \to M \in \mathbb{R} \land 0 \leq M \land \forall n \in \mathbb{N}\,(n \leq N \to \lvert a_{n} \rvert \leq M) \to \varphi_{39}$
\item $\varphi_{143} :\equiv \forall z\,\forall w\,(z = w \to \lvert a_{n} \rvert = \lvert a_{z} \rvert \to \lvert a_{n} \rvert = \lvert a_{w} \rvert) \to \forall w\,(n = w \to \lvert a_{n} \rvert = \lvert a_{n} \rvert \to \lvert a_{n} \rvert = \lvert a_{w} \rvert)$
\item $\varphi_{144} :\equiv \forall N\,(\mathrm{Seq}(a) \land \forall N\,(N \in B_{8} \leftrightarrow N \in \omega \land \varphi_{51})) \to \forall N \in B_{8}\,(S(N) \in B_{8})$
\item $\varphi_{145} :\equiv N \neq \emptyset \to M \in \mathbb{R} \land 0 \leq M \land \forall n \in \mathbb{N}\,(n \leq N \to \lvert a_{n} \rvert \leq M) \to \varphi_{48}$
\item $\varphi_{146} :\equiv \forall y\,(M \in \mathbb{R} \to y \in \mathbb{R} \to 0 \leq M \to 0 \leq y \to 0 \leq M + y) \to M \in \mathbb{R} \to \varphi_{66}$
\item $\varphi_{147} :\equiv \forall N\,(N \in B_{8} \leftrightarrow N \in \omega \land \varphi_{51}) \to N \in B_{8} \to \forall y \in \omega\,(S(y) \in \omega)$
\item $\varphi_{148} :\equiv (\mathrm{Seq}(a) \to N \in \mathbb{N} \to N \in \omega) \to (\mathrm{Seq}(a) \to N \in \mathbb{N} \to N \in \mathbb{N}) \to \mathrm{Seq}(a) \to \varphi_{51}$
\item $\varphi_{149} :\equiv \mathrm{Seq}(a) \to \forall N\,(N \in B_{8} \leftrightarrow N \in \omega \land \varphi_{51}) \to N \in B_{8} \to N \in \omega \to S(N) \in \omega$
\item $\varphi_{150} :\equiv \forall N\,(N \in B_{8} \leftrightarrow N \in \omega \land \varphi_{51}) \to N \in B_{8} \to S(N) \in \mathbb{N} \to N = \emptyset \to S(N) = 1_{\omega}$
\item $\varphi_{151} :\equiv \forall N\,(N \in B_{8} \leftrightarrow N \in \omega \land \varphi_{51}) \to N \in B_{8} \to S(N) \in \mathbb{N} \to N = \emptyset \to 1_{\omega} = S(N)$
\item $\varphi_{152} :\equiv \forall z\,\forall w\,(z = w \to \lvert a_{n} \rvert < \lvert a_{z} \rvert \to \lvert a_{n} \rvert < \lvert a_{w} \rvert) \to \forall w\,(n = w \to \lvert a_{n} \rvert < \lvert a_{n} \rvert \to \lvert a_{n} \rvert < \lvert a_{w} \rvert)$
\item $\varphi_{153} :\equiv \forall N\,(N \in B_{8} \leftrightarrow N \in \omega \land \varphi_{51}) \to N \in B_{8} \to S(N) \in \mathbb{N} \to N \neq \emptyset \to N \in \mathbb{N}$
\item $\varphi_{154} :\equiv \forall N\,(N \in B_{8} \leftrightarrow N \in \omega \land \varphi_{51}) \to N \in B_{8} \to S(N) \in \mathbb{N} \to N = \emptyset \to N = \emptyset$
\item $\varphi_{155} :\equiv \forall N\,(N \in B_{8} \leftrightarrow N \in \omega \land \varphi_{51}) \to N \in B_{8} \to S(N) \in \mathbb{N} \to N \neq \emptyset \to N \neq \emptyset$
\item $\varphi_{156} :\equiv \forall N\,(N \in B_{8} \leftrightarrow N \in \omega \land \varphi_{51}) \to \emptyset \in \omega \land \forall y \in \omega\,(S(y) \in \omega)$
\item $\varphi_{157} :\equiv \forall N\,(N \in B_{8} \leftrightarrow N \in \omega \land \varphi_{51}) \to N \in B_{8} \to S(N) \in \mathbb{N} \to N = \emptyset \to S(N) = S(\emptyset)$
\item $\varphi_{158} :\equiv \forall n\,(1_{\omega} = S(N)) \to \forall n\,((n \in \mathbb{N} \to \varphi_{29}) \to n \in \mathbb{N} \to \varphi_{28})$
\item $\varphi_{159} :\equiv \forall n\,(1_{\omega} = S(N)) \to \forall n\,((n \in \mathbb{N} \to \varphi_{28}) \to n \in \mathbb{N} \to \varphi_{29})$
\item $\varphi_{160} :\equiv 1_{\omega} = S(N) \to (\neg \forall n \in \mathbb{N}\,(\varphi_{28}) \leftrightarrow \neg \forall n \in \mathbb{N}\,(\varphi_{29}))$
\item $\varphi_{161} :\equiv (\forall N \in B_{8}\,(S(N) \in B_{8}) \to \forall y \in B_{8}\,(S(y) \in B_{8})) \to \mathrm{Seq}(a) \to \forall N\,(N \in B_{8} \leftrightarrow N \in \omega \land \varphi_{51}) \to \mathrm{Ind}(B_{8})$
\item $\varphi_{162} :\equiv (S(N) = 1_{\omega} \to 1_{\omega} = S(N)) \to \mathrm{Seq}(a) \to \varphi_{151}$
\item $\varphi_{163} :\equiv (\forall x\,(\mathrm{Ind}(x) \to N \in x) \to \mathrm{Ind}(B_{8}) \to N \in B_{8}) \to \mathrm{Seq}(a) \to \varphi_{137}$
\item $\varphi_{164} :\equiv (1_{\omega} = S(N) \to \forall n\,(1_{\omega} = S(N))) \to 1_{\omega} = S(N) \to \forall n\,((n \in \mathbb{N} \to \varphi_{29}) \to n \in \mathbb{N} \to \varphi_{28})$
\item $\varphi_{165} :\equiv (1_{\omega} = S(N) \to \forall n\,(1_{\omega} = S(N))) \to 1_{\omega} = S(N) \to \forall n\,((n \in \mathbb{N} \to \varphi_{28}) \to n \in \mathbb{N} \to \varphi_{29})$
\item $\varphi_{166} :\equiv N \neq \emptyset \to M \in \mathbb{R} \land 0 \leq M \land \forall n \in \mathbb{N}\,(n \leq N \to \lvert a_{n} \rvert \leq M) \to \varphi_{127}$
\item $\varphi_{167} :\equiv \forall N\,(\mathrm{Seq}(a) \land \forall N\,(N \in B_{8} \leftrightarrow N \in \omega \land \varphi_{51}) \to N \in \omega \to \varphi_{51})$
\item $\varphi_{168} :\equiv \mathrm{Seq}(a) \land \neg \forall N \in \omega\,(\varphi_{51}) \to \neg \forall N\,(N \in B_{8} \leftrightarrow N \in \omega \land \varphi_{51})$
\item $\varphi_{169} :\equiv \forall N\,(\mathrm{Seq}(a) \land \forall N\,(N \in B_{8} \leftrightarrow N \in \omega \land \varphi_{51})) \to \forall N \in \omega\,(\varphi_{51})$
\item $\varphi_{170} :\equiv \mathrm{Seq}(a) \land \neg \forall N \in \omega\,(\varphi_{51}) \to \forall B_{8}\,(\mathrm{Seq}(a) \land \neg \forall N \in \omega\,(\varphi_{51}))$
\item $\varphi_{171} :\equiv \mathrm{Seq}(a) \to \forall N\,(N \in B_{8} \leftrightarrow N \in \omega \land \varphi_{51}) \to (N \in B_{8} \leftrightarrow N \in \omega \land \varphi_{51})$
\item $\varphi_{172} :\equiv \mathrm{Seq}(a) \to \forall N\,(N \in B_{8} \leftrightarrow N \in \omega \land \varphi_{51}) \to N \in B_{8} \to N \in \omega \land \varphi_{51}$
\item $\varphi_{173} :\equiv \mathrm{Seq}(a) \land \forall N\,(N \in B_{8} \leftrightarrow N \in \omega \land \varphi_{51}) \land N \in B_{8} \land S(N) \in \mathbb{N} \land N \neq \emptyset \land (M \in \mathbb{R} \land 0 \leq M \land \forall n \in \mathbb{N}\,(n \leq N \to \lvert a_{n} \rvert \leq M))$
\item $\varphi_{174} :\equiv \forall B_{8}\,(\mathrm{Seq}(a) \land \neg \forall N \in \omega\,(\varphi_{51})) \to \forall B_{8}\,(\neg \forall N\,(N \in B_{8} \leftrightarrow N \in \omega \land \varphi_{51}))$
\item $\varphi_{175} :\equiv \mathrm{Seq}(a) \to \forall N\,(N \in B_{8} \leftrightarrow N \in \omega \land \varphi_{51}) \to \forall N\,(N \in B_{8} \leftrightarrow N \in \omega \land \varphi_{51})$
\item $\varphi_{176} :\equiv \mathrm{Seq}(a) \to \forall N\,(N \in B_{8} \leftrightarrow N \in \omega \land \varphi_{51}) \to N \in B_{8} \to S(N) \in \mathbb{N} \to \varphi_{57}$
\item $\varphi_{177} :\equiv (\mathrm{Seq}(a) \to 0 \leq \lvert a_{1_{\omega}} \rvert) \to \varphi_{36} \to \mathrm{Seq}(a) \to \varphi_{106}$
\item $\varphi_{178} :\equiv \forall A_{9}\,\exists B_{8}\,\forall N\,(N \in B_{8} \leftrightarrow N \in A_{9} \land \varphi_{51}) \to \exists B_{8}\,\forall N\,(N \in B_{8} \leftrightarrow N \in \omega \land \varphi_{51})$
\item $\varphi_{179} :\equiv (N = \emptyset \to S(N) = S(N) \to S(N) = S(\emptyset)) \to \mathrm{Seq}(a) \to \varphi_{157}$
\item $\varphi_{180} :\equiv \forall N\,(N \in B_{8} \leftrightarrow N \in \omega \land \varphi_{51}) \to N \in B_{8} \to S(N) \in \mathbb{N} \to N \neq \emptyset \to \varphi_{47}$
\item $\varphi_{181} :\equiv \mathrm{Seq}(a) \to \forall N\,(N \in B_{8} \leftrightarrow N \in \omega \land \varphi_{51}) \to N \in B_{8} \to S(N) \in \mathbb{N} \to \varphi_{63}$
\item $\varphi_{182} :\equiv \forall N\,(N \in B_{8} \leftrightarrow N \in \omega \land \varphi_{51}) \to N \in B_{8} \to S(N) \in \mathbb{N} \to N \neq \emptyset \to \varphi_{52}$
\item $\varphi_{183} :\equiv \mathrm{Seq}(a) \land \forall N\,(N \in B_{8} \leftrightarrow N \in \omega \land \varphi_{51}) \land N \in B_{8} \land S(N) \in \mathbb{N} \land N \neq \emptyset \land \neg \exists M\,(M \in \mathbb{R} \land 0 \leq M \land \forall n \in \mathbb{N}\,(\varphi_{29}))$
\item $\varphi_{184} :\equiv \neg (M \in \mathbb{R} \to \neg (\neg 0 < M \to 0 = M)) \to \neg \forall n \in \mathbb{N}\,(\varphi_{28}) \leftrightarrow \neg (M \in \mathbb{R} \to \neg (\neg 0 < M \to 0 = M)) \to \neg \forall n \in \mathbb{N}\,(\varphi_{29})$
\item $\varphi_{185} :\equiv \exists B_{8}\,\forall N\,(N \in B_{8} \leftrightarrow N \in \omega \land \varphi_{51}) \to \mathrm{Seq}(a) \to \exists B_{8}\,\forall N\,(N \in B_{8} \leftrightarrow N \in \omega \land \varphi_{51})$
\item $\varphi_{186} :\equiv \forall N\,(N \in B_{8} \leftrightarrow N \in \omega \land \varphi_{51}) \to N \in B_{8} \to S(N) \in \mathbb{N} \to N \neq \emptyset \to \varphi_{64}$
\item $\varphi_{187} :\equiv \neg (\neg (M \in \mathbb{R} \to \neg (\neg 0 < M \to 0 = M)) \to \neg \forall n \in \mathbb{N}\,(\varphi_{28})) \leftrightarrow \neg (\neg (M \in \mathbb{R} \to \neg (\neg 0 < M \to 0 = M)) \to \neg \forall n \in \mathbb{N}\,(\varphi_{29}))$
\item $\varphi_{188} :\equiv (\forall n\,(\varphi_{33}) \to \varphi_{44}) \to (\mathrm{Seq}(a) \to \forall n\,(\mathrm{Seq}(a))) \to \varphi_{36}$
\item $\varphi_{189} :\equiv \forall N\,(N \in B_{8} \leftrightarrow N \in \omega \land \varphi_{51}) \to N \in B_{8} \to S(N) \in \mathbb{N} \to N \neq \emptyset \to \varphi_{65}$
\item $\varphi_{190} :\equiv \mathrm{Seq}(a) \land \forall N\,(N \in B_{8} \leftrightarrow N \in \omega \land \varphi_{51}) \to \forall N\,(\mathrm{Seq}(a) \land \forall N\,(N \in B_{8} \leftrightarrow N \in \omega \land \varphi_{51}))$
\item $\varphi_{191} :\equiv \forall N\,(N \in B_{8} \leftrightarrow N \in \omega \land \varphi_{51}) \to N \in B_{8} \to S(N) \in \mathbb{N} \to N \neq \emptyset \to \varphi_{69}$
\item $\varphi_{192} :\equiv (\mathrm{Seq}(a) \to \varphi_{76}) \to \mathrm{Seq}(a) \to \forall N\,(N \in B_{8} \leftrightarrow N \in \omega \land \varphi_{51}) \to \emptyset \in B_{8}$
\item $\varphi_{193} :\equiv \forall N\,(N \in B_{8} \leftrightarrow N \in \omega \land \varphi_{51}) \to N \in B_{8} \to S(N) \in \mathbb{N} \to N \neq \emptyset \to \varphi_{71}$
\item $\varphi_{194} :\equiv \forall N\,(N \in B_{8} \leftrightarrow N \in \omega \land \varphi_{51}) \to N \in B_{8} \to S(N) \in \mathbb{N} \to N \neq \emptyset \to \varphi_{75}$
\item $\varphi_{195} :\equiv \forall N\,(N \in B_{8} \leftrightarrow N \in \omega \land \varphi_{51}) \to N \in B_{8} \to S(N) \in \mathbb{N} \to N \neq \emptyset \to \varphi_{81}$
\item $\varphi_{196} :\equiv (\forall y \in B_{8}\,(S(y) \in B_{8}) \to \forall N \in B_{8}\,(S(N) \in B_{8})) \to \varphi_{161}$
\item $\varphi_{197} :\equiv \forall N\,(N \in B_{8} \leftrightarrow N \in \omega \land \varphi_{51}) \to N \in B_{8} \to S(N) \in \mathbb{N} \to N \neq \emptyset \to \varphi_{89}$
\item $\varphi_{198} :\equiv \forall N\,(N \in B_{8} \leftrightarrow N \in \omega \land \varphi_{51}) \to N \in B_{8} \to S(N) \in \mathbb{N} \to N \neq \emptyset \to \varphi_{90}$
\item $\varphi_{199} :\equiv (\mathrm{Ind}(\omega) \leftrightarrow \emptyset \in \omega \land \forall y \in \omega\,(S(y) \in \omega)) \to \mathrm{Seq}(a) \to \varphi_{156}$
\item $\varphi_{200} :\equiv \forall N\,(N \in B_{8} \leftrightarrow N \in \omega \land \varphi_{51}) \to N \in B_{8} \to S(N) \in \mathbb{N} \to N \neq \emptyset \to \varphi_{94}$
\item $\varphi_{201} :\equiv \forall N\,(N \in B_{8} \leftrightarrow N \in \omega \land \varphi_{51}) \to N \in B_{8} \to S(N) \in \mathbb{N} \to N \neq \emptyset \to \varphi_{95}$
\item $\varphi_{202} :\equiv \forall N\,(N \in B_{8} \leftrightarrow N \in \omega \land \varphi_{51}) \to N \in B_{8} \to S(N) \in \mathbb{N} \to N \neq \emptyset \to \varphi_{98}$
\item $\varphi_{203} :\equiv \mathrm{Seq}(a) \to \forall N\,(N \in B_{8} \leftrightarrow N \in \omega \land \varphi_{51}) \to (\emptyset \in B_{8} \leftrightarrow \emptyset \in \omega \land \varphi_{76})$
\item $\varphi_{204} :\equiv \forall N\,(N \in B_{8} \leftrightarrow N \in \omega \land \varphi_{51}) \to N \in B_{8} \to S(N) \in \mathbb{N} \to N \neq \emptyset \to \varphi_{102}$
\item $\varphi_{205} :\equiv \forall N\,(N \in B_{8} \leftrightarrow N \in \omega \land \varphi_{51}) \to N \in B_{8} \to S(N) \in \mathbb{N} \to N \neq \emptyset \to \varphi_{103}$
\item $\varphi_{206} :\equiv \forall N\,(N \in B_{8} \leftrightarrow N \in \omega \land \varphi_{51}) \to N \in B_{8} \to S(N) \in \mathbb{N} \to N \neq \emptyset \to \varphi_{104}$
\item $\varphi_{207} :\equiv (\forall M\,(\varphi_{114}) \to \neg \varphi_{106}) \to \mathrm{Seq}(a) \to \exists M\,(M \in \mathbb{R} \land 0 \leq M \land \forall n \in \mathbb{N}\,(\varphi_{28}))$
\item $\varphi_{208} :\equiv \forall N\,(N \in B_{8} \leftrightarrow N \in \omega \land \varphi_{51}) \to N \in B_{8} \to S(N) \in \mathbb{N} \to N \neq \emptyset \to \varphi_{109}$
\item $\varphi_{209} :\equiv \forall N\,(N \in B_{8} \leftrightarrow N \in \omega \land \varphi_{51}) \to N \in B_{8} \to S(N) \in \mathbb{N} \to N \neq \emptyset \to \varphi_{112}$
\item $\varphi_{210} :\equiv \forall N\,(N \in B_{8} \leftrightarrow N \in \omega \land \varphi_{51}) \to N \in B_{8} \to S(N) \in \mathbb{N} \to N \neq \emptyset \to \varphi_{113}$
\item $\varphi_{211} :\equiv \forall N\,(N \in B_{8} \leftrightarrow N \in \omega \land \varphi_{51}) \to N \in B_{8} \to S(N) \in \mathbb{N} \to N \neq \emptyset \to \varphi_{117}$
\item $\varphi_{212} :\equiv \forall N\,(N \in B_{8} \leftrightarrow N \in \omega \land \varphi_{51}) \to N \in B_{8} \to S(N) \in \mathbb{N} \to N \neq \emptyset \to \varphi_{121}$
\item $\varphi_{213} :\equiv \forall N\,(N \in B_{8} \leftrightarrow N \in \omega \land \varphi_{51}) \to N \in B_{8} \to S(N) \in \mathbb{N} \to N \neq \emptyset \to \varphi_{122}$
\item $\varphi_{214} :\equiv \forall N\,(N \in B_{8} \leftrightarrow N \in \omega \land \varphi_{51}) \to N \in B_{8} \to S(N) \in \mathbb{N} \to N \neq \emptyset \to \varphi_{124}$
\item $\varphi_{215} :\equiv \forall N\,(N \in B_{8} \leftrightarrow N \in \omega \land \varphi_{51}) \to N \in B_{8} \to S(N) \in \mathbb{N} \to N \neq \emptyset \to \varphi_{125}$
\item $\varphi_{216} :\equiv \forall N\,(N \in B_{8} \leftrightarrow N \in \omega \land \varphi_{51}) \to N \in B_{8} \to S(N) \in \mathbb{N} \to N \neq \emptyset \to \varphi_{128}$
\item $\varphi_{217} :\equiv (1_{\omega} = S(N) \to \forall M\,(1_{\omega} = S(N))) \to 1_{\omega} = S(N) \to \forall M\,(\varphi_{115} \to \varphi_{114})$
\item $\varphi_{218} :\equiv (1_{\omega} = S(N) \to \forall M\,(1_{\omega} = S(N))) \to 1_{\omega} = S(N) \to \forall M\,(\varphi_{114} \to \varphi_{115})$
\item $\varphi_{219} :\equiv \forall N\,(N \in B_{8} \leftrightarrow N \in \omega \land \varphi_{51}) \to N \in B_{8} \to S(N) \in \mathbb{N} \to N \neq \emptyset \to \varphi_{130}$
\item $\varphi_{220} :\equiv \forall n\,(\varphi_{173} \to n \in \mathbb{N} \to n \leq S(N) \to \lvert a_{n} \rvert \leq M + \lvert a_{S(N)} \rvert)$
\item $\varphi_{221} :\equiv \mathrm{Seq}(a) \to \forall N\,(N \in B_{8} \leftrightarrow N \in \omega \land \varphi_{51}) \to N \in B_{8} \to S(N) \in \mathbb{N} \to \varphi_{136}$
\item $\varphi_{222} :\equiv (\mathrm{Seq}(a) \to n \in \mathbb{N} \to n \leq 1_{\omega} \to \lvert a_{n} \rvert \leq \lvert a_{n} \rvert) \to \varphi_{61} \to \varphi_{49} \to \varphi_{33}$
\item $\varphi_{223} :\equiv (\mathrm{Seq}(a) \to \forall N\,(N \in B_{8} \leftrightarrow N \in \omega \land \varphi_{51}) \to N \in B_{8} \to N \in \omega) \to \varphi_{139}$
\item $\varphi_{224} :\equiv \forall n\,(\varphi_{173}) \to \forall n \in \mathbb{N}\,(n \leq S(N) \to \lvert a_{n} \rvert \leq M + \lvert a_{S(N)} \rvert)$
\item $\varphi_{225} :\equiv \mathrm{Seq}(a) \to \forall N\,(N \in B_{8} \leftrightarrow N \in \omega \land \varphi_{51}) \to N \in B_{8} \to S(N) \in \mathbb{N} \to \varphi_{141}$
\item $\varphi_{226} :\equiv \varphi_{183} \to \neg (M \in \mathbb{R} \land 0 \leq M \land \forall n \in \mathbb{N}\,(n \leq N \to \lvert a_{n} \rvert \leq M))$
\item $\varphi_{227} :\equiv \mathrm{Seq}(a) \to \forall N\,(N \in B_{8} \leftrightarrow N \in \omega \land \varphi_{51}) \to N \in B_{8} \to S(N) \in \mathbb{N} \to \varphi_{142}$
\item $\varphi_{228} :\equiv \mathrm{Seq}(a) \to \forall N\,(N \in B_{8} \leftrightarrow N \in \omega \land \varphi_{51}) \to N \in B_{8} \to S(N) \in \mathbb{N} \to \varphi_{145}$
\item $\varphi_{229} :\equiv (\forall N\,(\mathrm{Seq}(a) \to \varphi_{51}) \to \forall N\,(\mathrm{Seq}(a)) \to \varphi_{58}) \to (\mathrm{Seq}(a) \to \forall N\,(\mathrm{Seq}(a))) \to \mathrm{Seq}(a) \to \varphi_{58}$
\item $\varphi_{230} :\equiv (1_{\omega} = S(N) \to \varphi_{3}) \to \varphi_{91} \to 1_{\omega} = S(N) \to (\varphi_{28} \leftrightarrow \varphi_{29})$
\item $\varphi_{231} :\equiv (1_{\omega} = S(N) \to (\varphi_{28} \leftrightarrow \varphi_{29})) \to 1_{\omega} = S(N) \to (n \in \mathbb{N} \to \varphi_{28} \leftrightarrow n \in \mathbb{N} \to \varphi_{29})$
\item $\varphi_{232} :\equiv \forall n\,((n \in \mathbb{N} \to \varphi_{29}) \to n \in \mathbb{N} \to \varphi_{28}) \to \forall n \in \mathbb{N}\,(\varphi_{29}) \to \forall n \in \mathbb{N}\,(\varphi_{28})$
\item $\varphi_{233} :\equiv \forall n\,((n \in \mathbb{N} \to \varphi_{28}) \to n \in \mathbb{N} \to \varphi_{29}) \to \forall n \in \mathbb{N}\,(\varphi_{28}) \to \forall n \in \mathbb{N}\,(\varphi_{29})$
\item $\varphi_{234} :\equiv (\mathrm{Seq}(a) \to n \in \mathbb{N} \to n \leq 1_{\omega} \to n = 1_{\omega}) \to \varphi_{222}$
\item $\varphi_{235} :\equiv (1_{\omega} = S(N) \to (n \in \mathbb{N} \to \varphi_{28} \leftrightarrow n \in \mathbb{N} \to \varphi_{29})) \to 1_{\omega} = S(N) \to (n \in \mathbb{N} \to \varphi_{29}) \to n \in \mathbb{N} \to \varphi_{28}$
\item $\varphi_{236} :\equiv (1_{\omega} = S(N) \to (n \in \mathbb{N} \to \varphi_{28} \leftrightarrow n \in \mathbb{N} \to \varphi_{29})) \to 1_{\omega} = S(N) \to (n \in \mathbb{N} \to \varphi_{28}) \to n \in \mathbb{N} \to \varphi_{29}$
\item $\varphi_{237} :\equiv (\mathrm{Seq}(a) \to \forall N\,(N \in B_{8} \leftrightarrow N \in \omega \land \varphi_{51}) \to N \in \omega \to N \in \omega) \to (N \in \omega \leftrightarrow \forall x\,(\mathrm{Ind}(x) \to N \in x)) \to \mathrm{Seq}(a) \to \varphi_{138}$
\item $\varphi_{238} :\equiv \mathrm{Seq}(a) \to \forall N\,(N \in B_{8} \leftrightarrow N \in \omega \land \varphi_{51}) \to N \in B_{8} \to S(N) \in \mathbb{N} \to \varphi_{166}$
\item $\varphi_{239} :\equiv \forall n\,(1_{\omega} = S(N) \to (n \in \mathbb{N} \to \varphi_{29}) \to n \in \mathbb{N} \to \varphi_{28}) \to \varphi_{158}$
\item $\varphi_{240} :\equiv \forall n\,(1_{\omega} = S(N) \to (n \in \mathbb{N} \to \varphi_{28}) \to n \in \mathbb{N} \to \varphi_{29}) \to \varphi_{159}$
\item $\varphi_{241} :\equiv (1_{\omega} = S(N) \to (\forall n \in \mathbb{N}\,(\varphi_{28}) \leftrightarrow \forall n \in \mathbb{N}\,(\varphi_{29}))) \to \varphi_{160}$
\item $\varphi_{242} :\equiv \varphi_{172} \to \mathrm{Seq}(a) \to \forall N\,(N \in B_{8} \leftrightarrow N \in \omega \land \varphi_{51}) \to N \in B_{8} \to N \in \omega$
\item $\varphi_{243} :\equiv (\mathrm{Seq}(a) \to \varphi_{147}) \to (\forall y \in \omega\,(S(y) \in \omega) \to N \in \omega \to S(N) \in \omega) \to \varphi_{149}$
\item $\varphi_{244} :\equiv (\mathrm{Seq}(a) \to \forall N \in \omega\,(\varphi_{51})) \to (\forall N \in \omega\,(\varphi_{51}) \to N \in \omega \to \varphi_{51}) \to \mathrm{Seq}(a) \to N \in \omega \to \varphi_{51}$
\item $\varphi_{245} :\equiv (\mathrm{Seq}(a) \to \forall N\,(N \in B_{8} \leftrightarrow N \in \omega \land \varphi_{51}) \to N \in \omega \to \varphi_{51}) \to \mathrm{Seq}(a) \land \forall N\,(N \in B_{8} \leftrightarrow N \in \omega \land \varphi_{51}) \to N \in \omega \to \varphi_{51}$
\item $\varphi_{246} :\equiv (\mathrm{Seq}(a) \to \varphi_{157}) \to S(\emptyset) = 1_{\omega} \to \varphi_{8} \to \mathrm{Seq}(a) \to \varphi_{150}$
\item $\varphi_{247} :\equiv (\mathrm{Seq}(a) \to \forall N\,(N \in B_{8} \leftrightarrow N \in \omega \land \varphi_{51}) \to \forall N \in \omega\,(\varphi_{51})) \to \varphi_{168}$
\item $\varphi_{248} :\equiv (\forall N\,(N \in B_{8} \leftrightarrow N \in \omega \land \varphi_{51}) \to (N \in B_{8} \leftrightarrow N \in \omega \land \varphi_{51})) \to \varphi_{171}$
\item $\varphi_{249} :\equiv (N \in \mathbb{N} \to S(N) \in \mathbb{N}) \to (\mathrm{Seq}(a) \to \varphi_{153}) \to \mathrm{Seq}(a) \to \varphi_{182}$
\item $\varphi_{250} :\equiv (\mathrm{Seq}(a) \to \forall N\,(N \in B_{8} \leftrightarrow N \in \omega \land \varphi_{51}) \to N \in B_{8} \to \varphi_{96}) \to \mathrm{Seq}(a) \to \forall N\,(N \in B_{8} \leftrightarrow N \in \omega \land \varphi_{51}) \to \varphi_{83}$
\item $\varphi_{251} :\equiv \varphi_{190} \to \mathrm{Seq}(a) \to \forall N\,(N \in B_{8} \leftrightarrow N \in \omega \land \varphi_{51}) \to \forall N \in \omega\,(\varphi_{51})$
\item $\varphi_{252} :\equiv (\forall N\,(N \in B_{8} \leftrightarrow N \in \omega \land \varphi_{51}) \to \varphi_{110}) \to \mathrm{Seq}(a) \to \forall N\,(N \in B_{8} \leftrightarrow N \in \omega \land \varphi_{51}) \to \varphi_{110}$
\item $\varphi_{253} :\equiv (1_{\omega} = S(N) \to (\neg \forall M\,(\varphi_{114}) \leftrightarrow \neg \forall M\,(\varphi_{115}))) \to \mathrm{Seq}(a) \to \forall N\,(N \in B_{8} \leftrightarrow N \in \omega \land \varphi_{51}) \to N \in B_{8} \to \varphi_{92}$
\item $\varphi_{254} :\equiv (\mathrm{Seq}(a) \to \forall N\,(N \in B_{8} \leftrightarrow N \in \omega \land \varphi_{51}) \to N \in B_{8} \to N \in \omega) \to (\mathrm{Seq}(a) \to \varphi_{155}) \to \mathrm{Seq}(a) \to \varphi_{153}$
\item $\varphi_{255} :\equiv (\forall N\,(N \in B_{8} \leftrightarrow N \in \omega \land \varphi_{51}) \to (\emptyset \in B_{8} \leftrightarrow \emptyset \in \omega \land \varphi_{76})) \to \varphi_{203}$
\item $\varphi_{256} :\equiv (1_{\omega} = S(N) \to \varphi_{187}) \to 1_{\omega} = S(N) \to (\varphi_{114} \leftrightarrow \varphi_{115})$
\item $\varphi_{257} :\equiv (1_{\omega} = S(N) \to (\varphi_{114} \leftrightarrow \varphi_{115})) \to 1_{\omega} = S(N) \to \varphi_{115} \to \varphi_{114}$
\item $\varphi_{258} :\equiv (1_{\omega} = S(N) \to (\varphi_{114} \leftrightarrow \varphi_{115})) \to 1_{\omega} = S(N) \to \varphi_{114} \to \varphi_{115}$
\item $\varphi_{259} :\equiv \varphi_{232} \to 1_{\omega} = S(N) \to (\forall n \in \mathbb{N}\,(\varphi_{28}) \leftrightarrow \forall n \in \mathbb{N}\,(\varphi_{29}))$
\item $\varphi_{260} :\equiv (\mathrm{Seq}(a) \to \mathrm{Seq}(a)) \to (\mathrm{Seq}(a) \to \varphi_{182}) \to \mathrm{Seq}(a) \to \varphi_{189}$
\item $\varphi_{261} :\equiv \forall M\,(1_{\omega} = S(N) \to \varphi_{115} \to \varphi_{114}) \to \forall M\,(1_{\omega} = S(N)) \to \forall M\,(\varphi_{115} \to \varphi_{114})$
\item $\varphi_{262} :\equiv \forall M\,(1_{\omega} = S(N) \to \varphi_{114} \to \varphi_{115}) \to \forall M\,(1_{\omega} = S(N)) \to \forall M\,(\varphi_{114} \to \varphi_{115})$
\item $\varphi_{263} :\equiv (1_{\omega} = S(N) \to (\forall M\,(\varphi_{114}) \leftrightarrow \forall M\,(\varphi_{115}))) \to 1_{\omega} = S(N) \to (\neg \forall M\,(\varphi_{114}) \leftrightarrow \neg \forall M\,(\varphi_{115}))$
\item $\varphi_{264} :\equiv (\mathrm{Seq}(a) \to \forall N\,(N \in B_{8} \leftrightarrow N \in \omega \land \varphi_{51}) \to \emptyset \in B_{8}) \to \varphi_{140} \to \varphi_{196}$
\item $\varphi_{265} :\equiv (\mathrm{Seq}(a) \to \mathrm{Seq}(a)) \to (\mathrm{Seq}(a) \to \varphi_{182}) \to \mathrm{Seq}(a) \to \varphi_{208}$
\item $\varphi_{266} :\equiv (\mathrm{Seq}(a) \to \mathrm{Seq}(a)) \to (\mathrm{Seq}(a) \to \varphi_{182}) \to \mathrm{Seq}(a) \to \varphi_{209}$
\item $\varphi_{267} :\equiv (a_{S(N)} \in \mathbb{R} \to 0 \leq \lvert a_{S(N)} \rvert) \to (\mathrm{Seq}(a) \to \varphi_{189}) \to \mathrm{Seq}(a) \to \varphi_{186}$
\item $\varphi_{268} :\equiv (\mathrm{Seq}(a) \to \varphi_{204}) \to (n = S(N) \to S(N) = n) \to \mathrm{Seq}(a) \to \varphi_{205}$
\item $\varphi_{269} :\equiv (a_{S(N)} \in \mathbb{R} \to \lvert a_{S(N)} \rvert \in \mathbb{R}) \to (\mathrm{Seq}(a) \to \varphi_{189}) \to \mathrm{Seq}(a) \to \varphi_{194}$
\item $\varphi_{270} :\equiv (\mathrm{Seq}(a) \to \forall N\,(N \in B_{8} \leftrightarrow N \in \omega \land \varphi_{51}) \to \mathrm{Ind}(B_{8})) \to \varphi_{171} \to \mathrm{Seq}(a) \to \forall N\,(N \in B_{8} \leftrightarrow N \in \omega \land \varphi_{51}) \to N \in \omega \to \varphi_{51}$
\item $\varphi_{271} :\equiv (\mathrm{Seq}(a) \to \exists M\,(M \in \mathbb{R} \land 0 \leq M \land \forall n \in \mathbb{N}\,(\varphi_{28}))) \to \varphi_{253}$
\item $\varphi_{272} :\equiv (\mathrm{Seq}(a) \to \varphi_{214}) \to \varphi_{173} \to n \in \mathbb{N} \to n \leq S(N) \to \lvert a_{n} \rvert \leq M + \lvert a_{S(N)} \rvert$
\item $\varphi_{273} :\equiv (\mathrm{Ind}(B_{8}) \leftrightarrow \emptyset \in B_{8} \land \forall y \in B_{8}\,(S(y) \in B_{8})) \to \varphi_{264}$
\item $\varphi_{274} :\equiv (\mathrm{Seq}(a) \to \forall N\,(N \in B_{8} \leftrightarrow N \in \omega \land \varphi_{51}) \to N \in B_{8} \to \varphi_{51}) \to (\mathrm{Seq}(a) \to \varphi_{153}) \to \varphi_{176}$
\item $\varphi_{275} :\equiv (a_{S(N)} \in \mathbb{R} \to \lvert a_{S(N)} \rvert \in \mathbb{R}) \to (\mathrm{Seq}(a) \to \varphi_{208}) \to \mathrm{Seq}(a) \to \varphi_{211}$
\item $\varphi_{276} :\equiv (a_{S(N)} \in \mathbb{R} \to \lvert a_{S(N)} \rvert \in \mathbb{R}) \to (\mathrm{Seq}(a) \to \varphi_{209}) \to \mathrm{Seq}(a) \to \varphi_{213}$
\item $\varphi_{277} :\equiv (\mathrm{Seq}(a) \to \forall N\,(N \in B_{8} \leftrightarrow N \in \omega \land \varphi_{51}) \to N \in B_{8} \to \varphi_{92}) \to \varphi_{250}$
\item $\varphi_{278} :\equiv (\varphi_{183} \to \forall M\,(\varphi_{183})) \to \mathrm{Seq}(a) \to \forall N\,(N \in B_{8} \leftrightarrow N \in \omega \land \varphi_{51}) \to N \in B_{8} \to \varphi_{96}$
\item $\varphi_{279} :\equiv \forall n\,(1_{\omega} = S(N) \to (n \in \mathbb{N} \to \varphi_{29}) \to n \in \mathbb{N} \to \varphi_{28}) \to \varphi_{239} \to \varphi_{164}$
\item $\varphi_{280} :\equiv \forall n\,(1_{\omega} = S(N) \to (n \in \mathbb{N} \to \varphi_{28}) \to n \in \mathbb{N} \to \varphi_{29}) \to \varphi_{240} \to \varphi_{165}$
\item $\varphi_{281} :\equiv (\forall M\,(\varphi_{115} \to \varphi_{114}) \to \forall M\,(\varphi_{115}) \to \forall M\,(\varphi_{114})) \to 1_{\omega} = S(N) \to (\forall M\,(\varphi_{114}) \leftrightarrow \forall M\,(\varphi_{115}))$
\item $\varphi_{282} :\equiv (\mathrm{Seq}(a) \to \forall N\,(N \in B_{8} \leftrightarrow N \in \omega \land \varphi_{51}) \to \varphi_{110}) \to \varphi_{139} \to (\mathrm{Seq}(a) \to \forall N\,(N \in B_{8} \leftrightarrow N \in \omega \land \varphi_{51}) \to \varphi_{83}) \to \varphi_{134}$
\item $\varphi_{283} :\equiv (M \in \mathbb{R} \to \lvert a_{S(N)} \rvert \in \mathbb{R} \to M + \lvert a_{S(N)} \rvert \in \mathbb{R}) \to (\mathrm{Seq}(a) \to \varphi_{180}) \to (\mathrm{Seq}(a) \to \varphi_{194}) \to \mathrm{Seq}(a) \to \varphi_{198}$
\item $\varphi_{284} :\equiv \varphi_{68} \to (\mathrm{Seq}(a) \to \varphi_{206}) \to (\mathrm{Seq}(a) \to \varphi_{210}) \to \mathrm{Seq}(a) \to \varphi_{219}$
\item $\varphi_{285} :\equiv (\mathrm{Seq}(a) \to \varphi_{194}) \to \varphi_{181} \to (\mathrm{Seq}(a) \to \varphi_{186}) \to \mathrm{Seq}(a) \to \varphi_{193}$
\item $\varphi_{286} :\equiv \forall B_{8}\,(\varphi_{168}) \to (\forall B_{8}\,(\varphi_{168}) \to \varphi_{174}) \to \varphi_{170} \to \mathrm{Seq}(a) \to \forall N \in \omega\,(\varphi_{51})$
\item $\varphi_{287} :\equiv (1_{\omega} = S(N) \to \forall n\,((n \in \mathbb{N} \to \varphi_{29}) \to n \in \mathbb{N} \to \varphi_{28})) \to \varphi_{233} \to \varphi_{259}$
\item $\varphi_{288} :\equiv (\mathrm{Seq}(a) \to \varphi_{201}) \to (\mathrm{Seq}(a) \to \varphi_{200}) \to (\mathrm{Seq}(a) \to \varphi_{195}) \to \mathrm{Seq}(a) \to \varphi_{202}$
\item $\varphi_{289} :\equiv (\mathrm{Seq}(a) \to \varphi_{180}) \to (\mathrm{Seq}(a) \to \varphi_{211}) \to (\mathrm{Seq}(a) \to \varphi_{186}) \to \mathrm{Seq}(a) \to \varphi_{210}$
\item $\varphi_{290} :\equiv \varphi_{221} \to (\mathrm{Seq}(a) \to \varphi_{191}) \to (\mathrm{Seq}(a) \to \varphi_{197}) \to \mathrm{Seq}(a) \to \varphi_{206}$
\item $\varphi_{291} :\equiv (\mathrm{Seq}(a) \to \varphi_{202}) \to (\mathrm{Seq}(a) \to \varphi_{219}) \to \varphi_{227} \to \mathrm{Seq}(a) \to \varphi_{214}$
\item $\varphi_{292} :\equiv (\mathrm{Seq}(a) \to \varphi_{198}) \to (\mathrm{Seq}(a) \to \varphi_{193}) \to (\mathrm{Seq}(a) \to \varphi_{215}) \to \varphi_{238}$
\item $\varphi_{293} :\equiv (1_{\omega} = S(N) \to \forall n\,((n \in \mathbb{N} \to \varphi_{28}) \to n \in \mathbb{N} \to \varphi_{29})) \to \varphi_{287}$
\item $\varphi_{294} :\equiv (1_{\omega} = S(N) \to \forall M\,(\varphi_{115} \to \varphi_{114})) \to (\forall M\,(\varphi_{114} \to \varphi_{115}) \to \forall M\,(\varphi_{114}) \to \forall M\,(\varphi_{115})) \to \varphi_{281}$
\end{abbrevlist}
\end{fullproof}

% 自動生成 (prover/build.js): ConvBdd  ノード数 356
\begin{fullproof}{定理2.11}{$\forall a\,\forall \alpha\,(\mathrm{Conv}(a,\alpha) \to \mathrm{SeqBdd}(a))$}{356}\phantomsection\label{tree:ConvBdd}
\noindent$\Pi_{1}$:
\begin{ptree}
\AXC{$\forall x \in \mathbb{R}\,(0 \leq \lvert x \rvert)$\lbl{[A6]}}
\AXC{$\forall x \in \mathbb{R}\,(0 \leq \lvert x \rvert) \to \alpha \in \mathbb{R} \to 0 \leq \lvert \alpha \rvert$\lbl{A2}}
\BIC{$\alpha \in \mathbb{R} \to 0 \leq \lvert \alpha \rvert$}
\AXC{$\varphi_{151}$\lbl{A1}}
\BIC{$(\mathrm{Conv}(a,\alpha) \to \alpha \in \mathbb{R}) \to \varphi_{105}$}
\end{ptree}
\noindent$\Pi_{2}$:
\begin{ptree}
\AXC{$\forall a\,\forall \alpha\,(\mathrm{Conv}(a,\alpha) \to \alpha \in \mathbb{R})$\lbl{[C2]}}
\AXC{$\varphi_{7}$\lbl{A2}}
\BIC{$\forall \alpha\,(\mathrm{Conv}(a,\alpha) \to \alpha \in \mathbb{R})$}
\AXC{$\varphi_{4}$\lbl{A2}}
\BIC{$\mathrm{Conv}(a,\alpha) \to \alpha \in \mathbb{R}$}
\AXC{$\stackrel{\Pi_{1}}{\vdots}$}\noLine
\UIC{$(\mathrm{Conv}(a,\alpha) \to \alpha \in \mathbb{R}) \to \varphi_{105}$}
\BIC{$\varphi_{105}$}
\end{ptree}
\noindent$\Pi_{3}$:
\begin{ptree}
\AXC{$0 < 1$\lbl{[F28]}}
\AXC{$0 < 1 \to \mathrm{Conv}(a,\alpha) \to 0 < 1$\lbl{A1}}
\BIC{$\mathrm{Conv}(a,\alpha) \to 0 < 1$}
\AXC{$(\mathrm{Conv}(a,\alpha) \to 0 < 1) \to \mathrm{Conv}(a,\alpha) \to \varphi_{94}$\lbl{A1}}
\BIC{$\mathrm{Conv}(a,\alpha) \to \varphi_{94}$}
\end{ptree}
\noindent$\Pi_{4}$:
\begin{ptree}
\AXC{$\forall x \in \mathbb{R}\,(\lvert x \rvert \in \mathbb{R})$\lbl{[A5]}}
\AXC{$\forall x \in \mathbb{R}\,(\lvert x \rvert \in \mathbb{R}) \to \alpha \in \mathbb{R} \to \lvert \alpha \rvert \in \mathbb{R}$\lbl{A2}}
\BIC{$\alpha \in \mathbb{R} \to \lvert \alpha \rvert \in \mathbb{R}$}
\AXC{$\varphi_{157}$\lbl{A1}}
\BIC{$(\mathrm{Conv}(a,\alpha) \to \alpha \in \mathbb{R}) \to \varphi_{114}$}
\end{ptree}
\noindent$\Pi_{5}$:
\begin{ptree}
\AXC{$\forall a\,\forall \alpha\,(\mathrm{Conv}(a,\alpha) \to \alpha \in \mathbb{R})$\lbl{[C2]}}
\AXC{$\varphi_{7}$\lbl{A2}}
\BIC{$\forall \alpha\,(\mathrm{Conv}(a,\alpha) \to \alpha \in \mathbb{R})$}
\AXC{$\varphi_{4}$\lbl{A2}}
\BIC{$\mathrm{Conv}(a,\alpha) \to \alpha \in \mathbb{R}$}
\AXC{$\stackrel{\Pi_{4}}{\vdots}$}\noLine
\UIC{$(\mathrm{Conv}(a,\alpha) \to \alpha \in \mathbb{R}) \to \varphi_{114}$}
\BIC{$\varphi_{114}$}
\end{ptree}
\noindent$\Pi_{6}$:
\begin{ptree}
\AXC{$\stackrel{\Pi_{5}}{\vdots}$}\noLine
\UIC{$\varphi_{114}$}
\AXC{$1 \in \mathbb{R}$\lbl{(R1)}}
\AXC{$\forall x\,\forall y\,(x \in \mathbb{R} \to y \in \mathbb{R} \to 0 \leq x \to 0 \leq y \to 0 \leq x + y)$\lbl{[E6]}}
\AXC{$\varphi_{69}$\lbl{A2}}
\BIC{$\forall y\,(1 \in \mathbb{R} \to y \in \mathbb{R} \to 0 \leq 1 \to 0 \leq y \to 0 \leq 1 + y)$}
\AXC{$\varphi_{88}$\lbl{A2}}
\BIC{$1 \in \mathbb{R} \to \lvert \alpha \rvert \in \mathbb{R} \to 0 \leq 1 \to 0 \leq \lvert \alpha \rvert \to 0 \leq 1 + \lvert \alpha \rvert$}
\AXC{$\varphi_{173}$\lbl{A1}}
\BIC{$1 \in \mathbb{R} \to \varphi_{114} \to (\mathrm{Conv}(a,\alpha) \to \varphi_{94}) \to \varphi_{105} \to \varphi_{111}$}
\BIC{$\varphi_{114} \to (\mathrm{Conv}(a,\alpha) \to \varphi_{94}) \to \varphi_{105} \to \varphi_{111}$}
\BIC{$(\mathrm{Conv}(a,\alpha) \to \varphi_{94}) \to \varphi_{105} \to \varphi_{111}$}
\end{ptree}
\noindent$\Pi_{7}$:
\begin{ptree}
\AXC{$\forall x \in \mathbb{R}\,(\lvert x \rvert \in \mathbb{R})$\lbl{[A5]}}
\AXC{$\forall x \in \mathbb{R}\,(\lvert x \rvert \in \mathbb{R}) \to \alpha \in \mathbb{R} \to \lvert \alpha \rvert \in \mathbb{R}$\lbl{A2}}
\BIC{$\alpha \in \mathbb{R} \to \lvert \alpha \rvert \in \mathbb{R}$}
\AXC{$\varphi_{159}$\lbl{A1}}
\BIC{$(\mathrm{Conv}(a,\alpha) \to \alpha \in \mathbb{R}) \to \varphi_{130}$}
\end{ptree}
\noindent$\Pi_{8}$:
\begin{ptree}
\AXC{$\forall a\,\forall \alpha\,(\mathrm{Conv}(a,\alpha) \to \alpha \in \mathbb{R})$\lbl{[C2]}}
\AXC{$\varphi_{7}$\lbl{A2}}
\BIC{$\forall \alpha\,(\mathrm{Conv}(a,\alpha) \to \alpha \in \mathbb{R})$}
\AXC{$\varphi_{4}$\lbl{A2}}
\BIC{$\mathrm{Conv}(a,\alpha) \to \alpha \in \mathbb{R}$}
\AXC{$\stackrel{\Pi_{7}}{\vdots}$}\noLine
\UIC{$(\mathrm{Conv}(a,\alpha) \to \alpha \in \mathbb{R}) \to \varphi_{130}$}
\BIC{$\varphi_{130}$}
\end{ptree}
\noindent$\Pi_{9}$:
\begin{ptree}
\AXC{$\stackrel{\Pi_{8}}{\vdots}$}\noLine
\UIC{$\varphi_{130}$}
\AXC{$1 \in \mathbb{R}$\lbl{(R1)}}
\AXC{$\forall x\,\forall y\,(x \in \mathbb{R} \to y \in \mathbb{R} \to x + y \in \mathbb{R})$\lbl{(R1)}}
\AXC{$\varphi_{53}$\lbl{A2}}
\BIC{$\forall y\,(1 \in \mathbb{R} \to y \in \mathbb{R} \to 1 + y \in \mathbb{R})$}
\AXC{$\varphi_{63}$\lbl{A2}}
\BIC{$1 \in \mathbb{R} \to \lvert \alpha \rvert \in \mathbb{R} \to 1 + \lvert \alpha \rvert \in \mathbb{R}$}
\AXC{$\varphi_{167}$\lbl{A1}}
\BIC{$1 \in \mathbb{R} \to \varphi_{130} \to \varphi_{134}$}
\BIC{$\varphi_{130} \to \varphi_{134}$}
\BIC{$\varphi_{134}$}
\end{ptree}
\noindent$\Pi_{10}$:
\begin{ptree}
\AXC{$\forall x\,\forall y\,(x \in \mathbb{R} \to y \in \mathbb{R} \to 0 \leq y \to x \leq x + y)$\lbl{[E1]}}
\AXC{$\varphi_{48}$\lbl{A2}}
\BIC{$\forall y\,(M \in \mathbb{R} \to y \in \mathbb{R} \to 0 \leq y \to M \leq M + y)$}
\AXC{$\forall y\,(M \in \mathbb{R} \to y \in \mathbb{R} \to 0 \leq y \to M \leq M + y) \to \varphi_{18}$\lbl{A2}}
\BIC{$\varphi_{18}$}
\end{ptree}
\noindent$\Pi_{11}$:
\begin{ptree}
\AXC{$\stackrel{\Pi_{9}}{\vdots}$}\noLine
\UIC{$\varphi_{134}$}
\AXC{$\varphi_{148}$\lbl{A1}}
\AXC{$\varphi_{148} \to \varphi_{107}$\lbl{A1}}
\BIC{$\varphi_{107}$}
\AXC{$\stackrel{\Pi_{10}}{\vdots}$}\noLine
\UIC{$\varphi_{18}$}
\AXC{$\varphi_{18} \to \varphi_{107} \to \varphi_{134} \to \varphi_{111} \to \varphi_{132}$\lbl{A1}}
\BIC{$\varphi_{107} \to \varphi_{134} \to \varphi_{111} \to \varphi_{132}$}
\BIC{$\varphi_{134} \to \varphi_{111} \to \varphi_{132}$}
\BIC{$\varphi_{111} \to \varphi_{132}$}
\end{ptree}
\noindent$\Pi_{12}$:
\begin{ptree}
\AXC{$\stackrel{\Pi_{2}}{\vdots}$}\noLine
\UIC{$\varphi_{105}$}
\AXC{$\stackrel{\Pi_{3}}{\vdots}$}\noLine
\UIC{$\mathrm{Conv}(a,\alpha) \to \varphi_{94}$}
\AXC{$\stackrel{\Pi_{6}}{\vdots}$}\noLine
\UIC{$(\mathrm{Conv}(a,\alpha) \to \varphi_{94}) \to \varphi_{105} \to \varphi_{111}$}
\BIC{$\varphi_{105} \to \varphi_{111}$}
\BIC{$\varphi_{111}$}
\AXC{$\stackrel{\Pi_{11}}{\vdots}$}\noLine
\UIC{$\varphi_{111} \to \varphi_{132}$}
\BIC{$\varphi_{132}$}
\end{ptree}
\noindent$\Pi_{13}$:
\begin{ptree}
\AXC{$\varphi_{122}$\lbl{A1}}
\AXC{$\varphi_{115}$\lbl{A1}}
\AXC{$\varphi_{15}$\lbl{A2}}
\AXC{$\varphi_{148}$\lbl{A1}}
\AXC{$\varphi_{148} \to \varphi_{149}$\lbl{A1}}
\BIC{$\varphi_{149}$}
\AXC{$\varphi_{149} \to \varphi_{15} \to \varphi_{147}$\lbl{A1}}
\BIC{$\varphi_{15} \to \varphi_{147}$}
\BIC{$\varphi_{147}$}
\AXC{$\varphi_{147} \to \varphi_{115} \to \varphi_{122} \to \varphi_{129}$\lbl{A1}}
\BIC{$\varphi_{115} \to \varphi_{122} \to \varphi_{129}$}
\BIC{$\varphi_{122} \to \varphi_{129}$}
\BIC{$\varphi_{129}$}
\end{ptree}
\noindent$\Pi_{14}$:
\begin{ptree}
\AXC{$\forall x\,\forall y\,\forall z\,(x \leq y \to y \leq z \to x \leq z)$\lbl{[O5]}}
\AXC{$\varphi_{16}$\lbl{A2}}
\BIC{$\forall y\,\forall z\,(\lvert a_{n} \rvert \leq y \to y \leq z \to \lvert a_{n} \rvert \leq z)$}
\AXC{$\varphi_{44}$\lbl{A2}}
\BIC{$\forall z\,(\lvert a_{n} \rvert \leq M \to M \leq z \to \lvert a_{n} \rvert \leq z)$}
\AXC{$\forall z\,(\lvert a_{n} \rvert \leq M \to M \leq z \to \lvert a_{n} \rvert \leq z) \to \varphi_{14}$\lbl{A2}}
\BIC{$\varphi_{14}$}
\end{ptree}
\noindent$\Pi_{15}$:
\begin{ptree}
\AXC{$\forall x \in \mathbb{R}\,(\lvert x \rvert \in \mathbb{R})$\lbl{[A5]}}
\AXC{$\forall x \in \mathbb{R}\,(\lvert x \rvert \in \mathbb{R}) \to \alpha \in \mathbb{R} \to \lvert \alpha \rvert \in \mathbb{R}$\lbl{A2}}
\BIC{$\alpha \in \mathbb{R} \to \lvert \alpha \rvert \in \mathbb{R}$}
\AXC{$\varphi_{158}$\lbl{A1}}
\BIC{$(\mathrm{Conv}(a,\alpha) \to \alpha \in \mathbb{R}) \to \varphi_{127}$}
\end{ptree}
\noindent$\Pi_{16}$:
\begin{ptree}
\AXC{$\forall a\,\forall \alpha\,(\mathrm{Conv}(a,\alpha) \to \alpha \in \mathbb{R})$\lbl{[C2]}}
\AXC{$\varphi_{7}$\lbl{A2}}
\BIC{$\forall \alpha\,(\mathrm{Conv}(a,\alpha) \to \alpha \in \mathbb{R})$}
\AXC{$\varphi_{4}$\lbl{A2}}
\BIC{$\mathrm{Conv}(a,\alpha) \to \alpha \in \mathbb{R}$}
\AXC{$\stackrel{\Pi_{15}}{\vdots}$}\noLine
\UIC{$(\mathrm{Conv}(a,\alpha) \to \alpha \in \mathbb{R}) \to \varphi_{127}$}
\BIC{$\varphi_{127}$}
\end{ptree}
\noindent$\Pi_{17}$:
\begin{ptree}
\AXC{$\stackrel{\Pi_{16}}{\vdots}$}\noLine
\UIC{$\varphi_{127}$}
\AXC{$1 \in \mathbb{R}$\lbl{(R1)}}
\AXC{$\forall x\,\forall y\,(x \in \mathbb{R} \to y \in \mathbb{R} \to x + y \in \mathbb{R})$\lbl{(R1)}}
\AXC{$\varphi_{53}$\lbl{A2}}
\BIC{$\forall y\,(1 \in \mathbb{R} \to y \in \mathbb{R} \to 1 + y \in \mathbb{R})$}
\AXC{$\varphi_{63}$\lbl{A2}}
\BIC{$1 \in \mathbb{R} \to \lvert \alpha \rvert \in \mathbb{R} \to 1 + \lvert \alpha \rvert \in \mathbb{R}$}
\AXC{$\varphi_{166}$\lbl{A1}}
\BIC{$1 \in \mathbb{R} \to \varphi_{127} \to \varphi_{131}$}
\BIC{$\varphi_{127} \to \varphi_{131}$}
\BIC{$\varphi_{131}$}
\end{ptree}
\noindent$\Pi_{18}$:
\begin{ptree}
\AXC{$\forall x\,\forall y\,(x \in \mathbb{R} \to y \in \mathbb{R} \to 0 \leq y \to x \leq y + x)$\lbl{[E2]}}
\AXC{$\forall x\,\forall y\,(x \in \mathbb{R} \to y \in \mathbb{R} \to 0 \leq y \to x \leq y + x) \to \varphi_{24}$\lbl{A2}}
\BIC{$\varphi_{24}$}
\AXC{$\varphi_{24} \to \varphi_{19}$\lbl{A2}}
\BIC{$\varphi_{19}$}
\end{ptree}
\noindent$\Pi_{19}$:
\begin{ptree}
\AXC{$\varphi_{148}$\lbl{A1}}
\AXC{$\varphi_{148} \to \varphi_{107}$\lbl{A1}}
\BIC{$\varphi_{107}$}
\AXC{$\stackrel{\Pi_{17}}{\vdots}$}\noLine
\UIC{$\varphi_{131}$}
\AXC{$\stackrel{\Pi_{18}}{\vdots}$}\noLine
\UIC{$\varphi_{19}$}
\AXC{$\varphi_{19} \to \varphi_{131} \to \varphi_{107} \to (\mathrm{Conv}(a,\alpha) \to \varphi_{95}) \to \varphi_{143}$\lbl{A1}}
\BIC{$\varphi_{131} \to \varphi_{107} \to (\mathrm{Conv}(a,\alpha) \to \varphi_{95}) \to \varphi_{143}$}
\BIC{$\varphi_{107} \to (\mathrm{Conv}(a,\alpha) \to \varphi_{95}) \to \varphi_{143}$}
\BIC{$(\mathrm{Conv}(a,\alpha) \to \varphi_{95}) \to \varphi_{143}$}
\end{ptree}
\noindent$\Pi_{20}$:
\begin{ptree}
\AXC{$\varphi_{45}$\lbl{A2}}
\AXC{$\mathrm{Conv}(a,\alpha) \to \varphi_{89}$\lbl{A1}}
\AXC{$(\mathrm{Conv}(a,\alpha) \to \varphi_{89}) \to \varphi_{154}$\lbl{A1}}
\BIC{$\varphi_{154}$}
\AXC{$\varphi_{154} \to \varphi_{45} \to \varphi_{150}$\lbl{A1}}
\BIC{$\varphi_{45} \to \varphi_{150}$}
\BIC{$\varphi_{150}$}
\AXC{$\varphi_{150} \to \varphi_{115} \to \varphi_{116} \to \varphi_{135}$\lbl{A1}}
\BIC{$\varphi_{115} \to \varphi_{116} \to \varphi_{135}$}
\end{ptree}
\noindent$\Pi_{21}$:
\begin{ptree}
\AXC{$\forall a\,\forall n\,(\mathrm{Seq}(a) \to n \in \mathbb{N} \to a_{n} \in \mathbb{R})$\lbl{[S1]}}
\AXC{$\varphi_{41}$\lbl{A2}}
\BIC{$\forall n\,(\mathrm{Seq}(a) \to n \in \mathbb{N} \to a_{n} \in \mathbb{R})$}
\AXC{$\varphi_{30}$\lbl{A2}}
\BIC{$\mathrm{Seq}(a) \to n \in \mathbb{N} \to a_{n} \in \mathbb{R}$}
\AXC{$\varphi_{165}$\lbl{A1}}
\BIC{$(\mathrm{Conv}(a,\alpha) \to \mathrm{Seq}(a)) \to \varphi_{115} \to \varphi_{128}$}
\end{ptree}
\noindent$\Pi_{22}$:
\begin{ptree}
\AXC{$\forall a\,\forall \alpha\,(\mathrm{Conv}(a,\alpha) \to \mathrm{Seq}(a))$\lbl{[C1]}}
\AXC{$\varphi_{3}$\lbl{A2}}
\BIC{$\forall \alpha\,(\mathrm{Conv}(a,\alpha) \to \mathrm{Seq}(a))$}
\AXC{$\varphi_{1}$\lbl{A2}}
\BIC{$\mathrm{Conv}(a,\alpha) \to \mathrm{Seq}(a)$}
\AXC{$\stackrel{\Pi_{21}}{\vdots}$}\noLine
\UIC{$(\mathrm{Conv}(a,\alpha) \to \mathrm{Seq}(a)) \to \varphi_{115} \to \varphi_{128}$}
\BIC{$\varphi_{115} \to \varphi_{128}$}
\end{ptree}
\noindent$\Pi_{23}$:
\begin{ptree}
\AXC{$\varphi_{115}$\lbl{A1}}
\AXC{$\stackrel{\Pi_{22}}{\vdots}$}\noLine
\UIC{$\varphi_{115} \to \varphi_{128}$}
\BIC{$\varphi_{128}$}
\end{ptree}
\noindent$\Pi_{24}$:
\begin{ptree}
\AXC{$\stackrel{\Pi_{23}}{\vdots}$}\noLine
\UIC{$\varphi_{128}$}
\AXC{$\forall x\,\forall y\,(x \in \mathbb{R} \to y \in \mathbb{R} \to x - y \in \mathbb{R})$\lbl{[F1]}}
\AXC{$\varphi_{79}$\lbl{A2}}
\BIC{$\forall y\,(a_{n} \in \mathbb{R} \to y \in \mathbb{R} \to a_{n} - y \in \mathbb{R})$}
\AXC{$\varphi_{90}$\lbl{A2}}
\BIC{$a_{n} \in \mathbb{R} \to \alpha \in \mathbb{R} \to a_{n} - \alpha \in \mathbb{R}$}
\AXC{$\varphi_{169}$\lbl{A1}}
\BIC{$\varphi_{128} \to (\mathrm{Conv}(a,\alpha) \to \alpha \in \mathbb{R}) \to \varphi_{133}$}
\BIC{$(\mathrm{Conv}(a,\alpha) \to \alpha \in \mathbb{R}) \to \varphi_{133}$}
\end{ptree}
\noindent$\Pi_{25}$:
\begin{ptree}
\AXC{$\forall a\,\forall \alpha\,(\mathrm{Conv}(a,\alpha) \to \alpha \in \mathbb{R})$\lbl{[C2]}}
\AXC{$\varphi_{7}$\lbl{A2}}
\BIC{$\forall \alpha\,(\mathrm{Conv}(a,\alpha) \to \alpha \in \mathbb{R})$}
\AXC{$\varphi_{4}$\lbl{A2}}
\BIC{$\mathrm{Conv}(a,\alpha) \to \alpha \in \mathbb{R}$}
\AXC{$\stackrel{\Pi_{24}}{\vdots}$}\noLine
\UIC{$(\mathrm{Conv}(a,\alpha) \to \alpha \in \mathbb{R}) \to \varphi_{133}$}
\BIC{$\varphi_{133}$}
\end{ptree}
\noindent$\Pi_{26}$:
\begin{ptree}
\AXC{$\stackrel{\Pi_{25}}{\vdots}$}\noLine
\UIC{$\varphi_{133}$}
\AXC{$\forall x \in \mathbb{R}\,(\lvert x \rvert \in \mathbb{R})$\lbl{[A5]}}
\AXC{$\varphi_{42}$\lbl{A2}}
\BIC{$a_{n} - \alpha \in \mathbb{R} \to \lvert a_{n} - \alpha \rvert \in \mathbb{R}$}
\AXC{$(a_{n} - \alpha \in \mathbb{R} \to \lvert a_{n} - \alpha \rvert \in \mathbb{R}) \to \varphi_{133} \to \varphi_{138}$\lbl{A1}}
\BIC{$\varphi_{133} \to \varphi_{138}$}
\BIC{$\varphi_{138}$}
\end{ptree}
\noindent$\Pi_{27}$:
\begin{ptree}
\AXC{$1 \in \mathbb{R}$\lbl{(R1)}}
\AXC{$\stackrel{\Pi_{26}}{\vdots}$}\noLine
\UIC{$\varphi_{138}$}
\AXC{$\varphi_{56}$\lbl{[F20]}}
\AXC{$\varphi_{56} \to \varphi_{112}$\lbl{A2}}
\BIC{$\varphi_{112}$}
\AXC{$\varphi_{112} \to \varphi_{109}$\lbl{A2}}
\BIC{$\varphi_{109}$}
\AXC{$\varphi_{109} \to \varphi_{123}$\lbl{A2}}
\BIC{$\varphi_{123}$}
\AXC{$\varphi_{162}$\lbl{A2}}
\BIC{$\lvert a_{n} - \alpha \rvert \in \mathbb{R} \to 1 \in \mathbb{R} \to \lvert \alpha \rvert \in \mathbb{R} \to \lvert \alpha \rvert \in \mathbb{R} \to \varphi_{47}$}
\AXC{$\varphi_{176}$\lbl{A1}}
\BIC{$\varphi_{138} \to 1 \in \mathbb{R} \to \varphi_{127} \to \varphi_{127} \to \varphi_{135} \to \varphi_{126} \to \varphi_{155}$}
\BIC{$1 \in \mathbb{R} \to \varphi_{127} \to \varphi_{127} \to \varphi_{135} \to \varphi_{126} \to \varphi_{155}$}
\BIC{$\varphi_{127} \to \varphi_{127} \to \varphi_{135} \to \varphi_{126} \to \varphi_{155}$}
\end{ptree}
\noindent$\Pi_{28}$:
\begin{ptree}
\AXC{$\varphi_{116}$\lbl{A1}}
\AXC{$\varphi_{115}$\lbl{A1}}
\AXC{$\stackrel{\Pi_{20}}{\vdots}$}\noLine
\UIC{$\varphi_{115} \to \varphi_{116} \to \varphi_{135}$}
\BIC{$\varphi_{116} \to \varphi_{135}$}
\BIC{$\varphi_{135}$}
\AXC{$\stackrel{\Pi_{16}}{\vdots}$}\noLine
\UIC{$\varphi_{127}$}
\AXC{$\stackrel{\Pi_{16}}{\vdots}$}\noLine
\UIC{$\varphi_{127}$}
\AXC{$\stackrel{\Pi_{27}}{\vdots}$}\noLine
\UIC{$\varphi_{127} \to \varphi_{127} \to \varphi_{135} \to \varphi_{126} \to \varphi_{155}$}
\BIC{$\varphi_{127} \to \varphi_{135} \to \varphi_{126} \to \varphi_{155}$}
\BIC{$\varphi_{135} \to \varphi_{126} \to \varphi_{155}$}
\BIC{$\varphi_{126} \to \varphi_{155}$}
\end{ptree}
\noindent$\Pi_{29}$:
\begin{ptree}
\AXC{$x = x$\lbl{A5}}
\UIC{$\forall x\,(x = x)$}
\AXC{$\forall x\,(x = x) \to \lvert \alpha \rvert = \lvert \alpha \rvert$\lbl{A2}}
\BIC{$\lvert \alpha \rvert = \lvert \alpha \rvert$}
\AXC{$\lvert \alpha \rvert = \lvert \alpha \rvert \to \varphi_{126}$\lbl{A1}}
\BIC{$\varphi_{126}$}
\AXC{$\stackrel{\Pi_{28}}{\vdots}$}\noLine
\UIC{$\varphi_{126} \to \varphi_{155}$}
\BIC{$\varphi_{155}$}
\end{ptree}
\noindent$\Pi_{30}$:
\begin{ptree}
\AXC{$\varphi_{43}$\lbl{A6}}
\UIC{$\forall w\,(\varphi_{43})$}
\UIC{$\forall z\,\forall w\,(\varphi_{43})$}
\AXC{$\forall z\,\forall w\,(\varphi_{43}) \to \varphi_{102}$\lbl{A2}}
\BIC{$\varphi_{102}$}
\AXC{$\varphi_{102} \to a_{n} - \alpha + \alpha = a_{n} \to \varphi_{80}$\lbl{A2}}
\BIC{$a_{n} - \alpha + \alpha = a_{n} \to \varphi_{80}$}
\end{ptree}
\noindent$\Pi_{31}$:
\begin{ptree}
\AXC{$\forall a\,\forall b\,(a \in \mathbb{R} \to b \in \mathbb{R} \to \lvert a + b \rvert \leq \lvert a \rvert + \lvert b \rvert)$\lbl{命題1.2}}
\AXC{$\varphi_{124}$\lbl{A2}}
\BIC{$\varphi_{58}$}
\AXC{$\varphi_{58} \to \varphi_{59}$\lbl{A2}}
\BIC{$\varphi_{59}$}
\end{ptree}
\noindent$\Pi_{32}$:
\begin{ptree}
\AXC{$\stackrel{\Pi_{25}}{\vdots}$}\noLine
\UIC{$\varphi_{133}$}
\AXC{$\stackrel{\Pi_{31}}{\vdots}$}\noLine
\UIC{$\varphi_{59}$}
\AXC{$\varphi_{59} \to \varphi_{133} \to (\mathrm{Conv}(a,\alpha) \to \alpha \in \mathbb{R}) \to \varphi_{156}$\lbl{A1}}
\BIC{$\varphi_{133} \to (\mathrm{Conv}(a,\alpha) \to \alpha \in \mathbb{R}) \to \varphi_{156}$}
\BIC{$(\mathrm{Conv}(a,\alpha) \to \alpha \in \mathbb{R}) \to \varphi_{156}$}
\end{ptree}
\noindent$\Pi_{33}$:
\begin{ptree}
\AXC{$\stackrel{\Pi_{23}}{\vdots}$}\noLine
\UIC{$\varphi_{128}$}
\AXC{$\forall x\,\forall y\,(x \in \mathbb{R} \to y \in \mathbb{R} \to x - y + y = x)$\lbl{[F12]}}
\AXC{$\varphi_{83}$\lbl{A2}}
\BIC{$\forall y\,(a_{n} \in \mathbb{R} \to y \in \mathbb{R} \to a_{n} - y + y = a_{n})$}
\AXC{$\varphi_{101}$\lbl{A2}}
\BIC{$a_{n} \in \mathbb{R} \to \alpha \in \mathbb{R} \to a_{n} - \alpha + \alpha = a_{n}$}
\AXC{$\varphi_{170}$\lbl{A1}}
\BIC{$\varphi_{128} \to (\mathrm{Conv}(a,\alpha) \to \alpha \in \mathbb{R}) \to \varphi_{141}$}
\BIC{$(\mathrm{Conv}(a,\alpha) \to \alpha \in \mathbb{R}) \to \varphi_{141}$}
\end{ptree}
\noindent$\Pi_{34}$:
\begin{ptree}
\AXC{$\forall a\,\forall \alpha\,(\mathrm{Conv}(a,\alpha) \to \alpha \in \mathbb{R})$\lbl{[C2]}}
\AXC{$\varphi_{7}$\lbl{A2}}
\BIC{$\forall \alpha\,(\mathrm{Conv}(a,\alpha) \to \alpha \in \mathbb{R})$}
\AXC{$\varphi_{4}$\lbl{A2}}
\BIC{$\mathrm{Conv}(a,\alpha) \to \alpha \in \mathbb{R}$}
\AXC{$\stackrel{\Pi_{33}}{\vdots}$}\noLine
\UIC{$(\mathrm{Conv}(a,\alpha) \to \alpha \in \mathbb{R}) \to \varphi_{141}$}
\BIC{$\varphi_{141}$}
\AXC{$\varphi_{141} \to \varphi_{174}$\lbl{A1}}
\BIC{$\varphi_{174}$}
\end{ptree}
\noindent$\Pi_{35}$:
\begin{ptree}
\AXC{$\forall a\,\forall \alpha\,(\mathrm{Conv}(a,\alpha) \to \alpha \in \mathbb{R})$\lbl{[C2]}}
\AXC{$\varphi_{7}$\lbl{A2}}
\BIC{$\forall \alpha\,(\mathrm{Conv}(a,\alpha) \to \alpha \in \mathbb{R})$}
\AXC{$\varphi_{4}$\lbl{A2}}
\BIC{$\mathrm{Conv}(a,\alpha) \to \alpha \in \mathbb{R}$}
\AXC{$\stackrel{\Pi_{32}}{\vdots}$}\noLine
\UIC{$(\mathrm{Conv}(a,\alpha) \to \alpha \in \mathbb{R}) \to \varphi_{156}$}
\BIC{$\varphi_{156}$}
\AXC{$\stackrel{\Pi_{34}}{\vdots}$}\noLine
\UIC{$\varphi_{174}$}
\BIC{$(a_{n} - \alpha + \alpha = a_{n} \to \varphi_{80}) \to (a_{n} - \alpha + \alpha = a_{n} \to \varphi_{68}) \to \varphi_{152}$}
\end{ptree}
\noindent$\Pi_{36}$:
\begin{ptree}
\AXC{$\stackrel{\Pi_{30}}{\vdots}$}\noLine
\UIC{$a_{n} - \alpha + \alpha = a_{n} \to \varphi_{80}$}
\AXC{$\stackrel{\Pi_{35}}{\vdots}$}\noLine
\UIC{$(a_{n} - \alpha + \alpha = a_{n} \to \varphi_{80}) \to (a_{n} - \alpha + \alpha = a_{n} \to \varphi_{68}) \to \varphi_{152}$}
\BIC{$(a_{n} - \alpha + \alpha = a_{n} \to \varphi_{68}) \to \varphi_{152}$}
\end{ptree}
\noindent$\Pi_{37}$:
\begin{ptree}
\AXC{$\varphi_{34}$\lbl{A6}}
\UIC{$\forall w\,(\varphi_{34})$}
\UIC{$\forall z\,\forall w\,(\varphi_{34})$}
\AXC{$\forall z\,\forall w\,(\varphi_{34}) \to \varphi_{98}$\lbl{A2}}
\BIC{$\varphi_{98}$}
\AXC{$\varphi_{98} \to a_{n} - \alpha + \alpha = a_{n} \to \varphi_{68}$\lbl{A2}}
\BIC{$a_{n} - \alpha + \alpha = a_{n} \to \varphi_{68}$}
\AXC{$\stackrel{\Pi_{36}}{\vdots}$}\noLine
\UIC{$(a_{n} - \alpha + \alpha = a_{n} \to \varphi_{68}) \to \varphi_{152}$}
\BIC{$\varphi_{152}$}
\end{ptree}
\noindent$\Pi_{38}$:
\begin{ptree}
\AXC{$\forall x\,\forall y\,\forall z\,(x \leq y \to y < z \to x < z)$\lbl{[O4]}}
\AXC{$\varphi_{36}$\lbl{A2}}
\BIC{$\forall y\,\forall z\,(\lvert a_{n} \rvert \leq y \to y < z \to \lvert a_{n} \rvert < z)$}
\AXC{$\forall y\,\forall z\,(\lvert a_{n} \rvert \leq y \to y < z \to \lvert a_{n} \rvert < z) \to \varphi_{75}$\lbl{A2}}
\BIC{$\varphi_{75}$}
\AXC{$\varphi_{75} \to \varphi_{99}$\lbl{A2}}
\BIC{$\varphi_{99}$}
\end{ptree}
\noindent$\Pi_{39}$:
\begin{ptree}
\AXC{$\forall x\,\forall y\,\forall z\,(x \leq y \to y \leq z \to x \leq z)$\lbl{[O5]}}
\AXC{$\varphi_{16}$\lbl{A2}}
\BIC{$\forall y\,\forall z\,(\lvert a_{n} \rvert \leq y \to y \leq z \to \lvert a_{n} \rvert \leq z)$}
\AXC{$\varphi_{84}$\lbl{A2}}
\BIC{$\forall z\,(\lvert a_{n} \rvert \leq 1 + \lvert \alpha \rvert \to 1 + \lvert \alpha \rvert \leq z \to \lvert a_{n} \rvert \leq z)$}
\AXC{$\varphi_{125}$\lbl{A2}}
\BIC{$\varphi_{51}$}
\AXC{$\varphi_{51} \to \varphi_{137} \to \varphi_{143} \to \varphi_{145}$\lbl{A1}}
\BIC{$\varphi_{137} \to \varphi_{143} \to \varphi_{145}$}
\end{ptree}
\noindent$\Pi_{40}$:
\begin{ptree}
\AXC{$\stackrel{\Pi_{29}}{\vdots}$}\noLine
\UIC{$\varphi_{155}$}
\AXC{$\stackrel{\Pi_{37}}{\vdots}$}\noLine
\UIC{$\varphi_{152}$}
\AXC{$\stackrel{\Pi_{38}}{\vdots}$}\noLine
\UIC{$\varphi_{99}$}
\AXC{$\varphi_{99} \to \varphi_{152} \to \varphi_{155} \to \varphi_{142}$\lbl{A1}}
\BIC{$\varphi_{152} \to \varphi_{155} \to \varphi_{142}$}
\BIC{$\varphi_{155} \to \varphi_{142}$}
\BIC{$\varphi_{142}$}
\AXC{$\varphi_{142} \to \varphi_{137}$\lbl{A1}}
\BIC{$\varphi_{137}$}
\AXC{$\stackrel{\Pi_{39}}{\vdots}$}\noLine
\UIC{$\varphi_{137} \to \varphi_{143} \to \varphi_{145}$}
\BIC{$\varphi_{143} \to \varphi_{145}$}
\end{ptree}
\noindent$\Pi_{41}$:
\begin{ptree}
\AXC{$\varphi_{148}$\lbl{A1}}
\AXC{$\varphi_{148} \to \mathrm{Conv}(a,\alpha) \to \varphi_{95}$\lbl{A1}}
\BIC{$\mathrm{Conv}(a,\alpha) \to \varphi_{95}$}
\AXC{$\stackrel{\Pi_{19}}{\vdots}$}\noLine
\UIC{$(\mathrm{Conv}(a,\alpha) \to \varphi_{95}) \to \varphi_{143}$}
\BIC{$\varphi_{143}$}
\AXC{$\stackrel{\Pi_{40}}{\vdots}$}\noLine
\UIC{$\varphi_{143} \to \varphi_{145}$}
\BIC{$\varphi_{145}$}
\end{ptree}
\noindent$\Pi_{42}$:
\begin{ptree}
\AXC{$\varphi_{115}$\lbl{A1}}
\AXC{$\forall n \in \mathbb{N}\,(n \in \omega)$\lbl{[N7']}}
\AXC{$\forall n \in \mathbb{N}\,(n \in \omega) \to n \in \mathbb{N} \to n \in \omega$\lbl{A2}}
\BIC{$n \in \mathbb{N} \to n \in \omega$}
\AXC{$(n \in \mathbb{N} \to n \in \omega) \to \varphi_{115} \to \varphi_{118}$\lbl{A1}}
\BIC{$\varphi_{115} \to \varphi_{118}$}
\BIC{$\varphi_{118}$}
\end{ptree}
\noindent$\Pi_{43}$:
\begin{ptree}
\AXC{$\forall n \in \mathbb{N}\,(n \in \omega)$\lbl{[N7']}}
\AXC{$\forall n \in \mathbb{N}\,(n \in \omega) \to N \in \mathbb{N} \to N \in \omega$\lbl{A2}}
\BIC{$N \in \mathbb{N} \to N \in \omega$}
\AXC{$(N \in \mathbb{N} \to N \in \omega) \to \varphi_{31} \to \varphi_{119}$\lbl{A1}}
\BIC{$\varphi_{31} \to \varphi_{119}$}
\end{ptree}
\noindent$\Pi_{44}$:
\begin{ptree}
\AXC{$\forall m\,\forall n\,(m \in \omega \to n \in \omega \to m \leq n \lor n < m)$\lbl{[N11]}}
\AXC{$\varphi_{49}$\lbl{A2}}
\BIC{$\forall n\,(N \in \omega \to n \in \omega \to N \leq n \lor n < N)$}
\AXC{$\varphi_{40}$\lbl{A2}}
\BIC{$N \in \omega \to n \in \omega \to N \leq n \lor n < N$}
\AXC{$(N \in \omega \to n \in \omega \to N \leq n \lor n < N) \to \varphi_{119} \to \varphi_{118} \to \varphi_{120}$\lbl{A1}}
\BIC{$\varphi_{119} \to \varphi_{118} \to \varphi_{120}$}
\end{ptree}
\noindent$\Pi_{45}$:
\begin{ptree}
\AXC{$\stackrel{\Pi_{42}}{\vdots}$}\noLine
\UIC{$\varphi_{118}$}
\AXC{$\mathrm{Conv}(a,\alpha) \to \varphi_{89}$\lbl{A1}}
\AXC{$(\mathrm{Conv}(a,\alpha) \to \varphi_{89}) \to \varphi_{31}$\lbl{A1}}
\BIC{$\varphi_{31}$}
\AXC{$\stackrel{\Pi_{43}}{\vdots}$}\noLine
\UIC{$\varphi_{31} \to \varphi_{119}$}
\BIC{$\varphi_{119}$}
\AXC{$\stackrel{\Pi_{44}}{\vdots}$}\noLine
\UIC{$\varphi_{119} \to \varphi_{118} \to \varphi_{120}$}
\BIC{$\varphi_{118} \to \varphi_{120}$}
\BIC{$\varphi_{120}$}
\end{ptree}
\noindent$\Pi_{46}$:
\begin{ptree}
\AXC{$\stackrel{\Pi_{12}}{\vdots}$}\noLine
\UIC{$\varphi_{132}$}
\AXC{$\stackrel{\Pi_{13}}{\vdots}$}\noLine
\UIC{$\varphi_{129}$}
\AXC{$\stackrel{\Pi_{14}}{\vdots}$}\noLine
\UIC{$\varphi_{14}$}
\AXC{$\varphi_{14} \to \varphi_{129} \to \varphi_{132} \to \varphi_{146}$\lbl{A1}}
\BIC{$\varphi_{129} \to \varphi_{132} \to \varphi_{146}$}
\BIC{$\varphi_{132} \to \varphi_{146}$}
\BIC{$\varphi_{146}$}
\AXC{$\stackrel{\Pi_{41}}{\vdots}$}\noLine
\UIC{$\varphi_{145}$}
\AXC{$\stackrel{\Pi_{45}}{\vdots}$}\noLine
\UIC{$\varphi_{120}$}
\AXC{$\varphi_{120} \to \varphi_{145} \to \varphi_{146} \to \varphi_{136}$\lbl{A1}}
\BIC{$\varphi_{145} \to \varphi_{146} \to \varphi_{136}$}
\BIC{$\varphi_{146} \to \varphi_{136}$}
\BIC{$\varphi_{136}$}
\end{ptree}
\noindent$\Pi_{47}$:
\begin{ptree}
\AXC{$\varphi_{100} \to \forall n\,(\varphi_{100})$\lbl{A4}}
\AXC{$\varphi_{163}$\lbl{A3}}
\AXC{$\stackrel{\Pi_{46}}{\vdots}$}\noLine
\UIC{$\varphi_{136}$}
\AXC{$\varphi_{136} \to \varphi_{100} \to n \in \mathbb{N} \to \lvert a_{n} \rvert \leq M + (1 + \lvert \alpha \rvert)$\lbl{A1}}
\BIC{$\varphi_{100} \to n \in \mathbb{N} \to \lvert a_{n} \rvert \leq M + (1 + \lvert \alpha \rvert)$}
\UIC{$\forall n\,(\varphi_{100} \to n \in \mathbb{N} \to \lvert a_{n} \rvert \leq M + (1 + \lvert \alpha \rvert))$}
\AXC{$\varphi_{175}$\lbl{A1}}
\BIC{$\varphi_{163} \to (\varphi_{100} \to \forall n\,(\varphi_{100})) \to \varphi_{139}$}
\BIC{$(\varphi_{100} \to \forall n\,(\varphi_{100})) \to \varphi_{139}$}
\BIC{$\varphi_{139}$}
\end{ptree}
\noindent$\Pi_{48}$:
\begin{ptree}
\AXC{$\stackrel{\Pi_{5}}{\vdots}$}\noLine
\UIC{$\varphi_{114}$}
\AXC{$1 \in \mathbb{R}$\lbl{(R1)}}
\AXC{$\forall x\,\forall y\,(x \in \mathbb{R} \to y \in \mathbb{R} \to x + y \in \mathbb{R})$\lbl{(R1)}}
\AXC{$\varphi_{53}$\lbl{A2}}
\BIC{$\forall y\,(1 \in \mathbb{R} \to y \in \mathbb{R} \to 1 + y \in \mathbb{R})$}
\AXC{$\varphi_{63}$\lbl{A2}}
\BIC{$1 \in \mathbb{R} \to \lvert \alpha \rvert \in \mathbb{R} \to 1 + \lvert \alpha \rvert \in \mathbb{R}$}
\AXC{$\varphi_{164}$\lbl{A1}}
\BIC{$1 \in \mathbb{R} \to \varphi_{114} \to \varphi_{117}$}
\BIC{$\varphi_{114} \to \varphi_{117}$}
\BIC{$\varphi_{117}$}
\end{ptree}
\noindent$\Pi_{49}$:
\begin{ptree}
\AXC{$\varphi_{148}$\lbl{A1}}
\AXC{$\varphi_{148} \to \varphi_{107}$\lbl{A1}}
\BIC{$\varphi_{107}$}
\AXC{$\forall x\,\forall y\,(x \in \mathbb{R} \to y \in \mathbb{R} \to x + y \in \mathbb{R})$\lbl{(R1)}}
\AXC{$\varphi_{54}$\lbl{A2}}
\BIC{$\forall y\,(M \in \mathbb{R} \to y \in \mathbb{R} \to M + y \in \mathbb{R})$}
\AXC{$\varphi_{81}$\lbl{A2}}
\BIC{$M \in \mathbb{R} \to 1 + \lvert \alpha \rvert \in \mathbb{R} \to M + (1 + \lvert \alpha \rvert) \in \mathbb{R}$}
\AXC{$\varphi_{171}$\lbl{A1}}
\BIC{$\varphi_{107} \to \varphi_{117} \to \varphi_{121}$}
\BIC{$\varphi_{117} \to \varphi_{121}$}
\end{ptree}
\noindent$\Pi_{50}$:
\begin{ptree}
\AXC{$\stackrel{\Pi_{47}}{\vdots}$}\noLine
\UIC{$\varphi_{139}$}
\AXC{$\stackrel{\Pi_{48}}{\vdots}$}\noLine
\UIC{$\varphi_{117}$}
\AXC{$\stackrel{\Pi_{49}}{\vdots}$}\noLine
\UIC{$\varphi_{117} \to \varphi_{121}$}
\BIC{$\varphi_{121}$}
\AXC{$\varphi_{121} \to \varphi_{139} \to \mathrm{Conv}(a,\alpha) \to \varphi_{153}$\lbl{A1}}
\BIC{$\varphi_{139} \to \mathrm{Conv}(a,\alpha) \to \varphi_{153}$}
\BIC{$\mathrm{Conv}(a,\alpha) \to \varphi_{153}$}
\AXC{$(\mathrm{Conv}(a,\alpha) \to \varphi_{153}) \to \varphi_{161}$\lbl{A1}}
\BIC{$\varphi_{161}$}
\end{ptree}
\noindent$\Pi_{51}$:
\begin{ptree}
\AXC{$\forall a\,(\varphi_{8})$\lbl{DefSeqBdd}}
\AXC{$\forall a\,(\varphi_{8}) \to \varphi_{8}$\lbl{A2}}
\BIC{$\varphi_{8}$}
\AXC{$\varphi_{8} \to (\mathrm{Conv}(a,\alpha) \to \mathrm{Seq}(a)) \to \varphi_{140} \to \varphi_{106}$\lbl{A1}}
\BIC{$(\mathrm{Conv}(a,\alpha) \to \mathrm{Seq}(a)) \to \varphi_{140} \to \varphi_{106}$}
\end{ptree}
\noindent$\Pi_{52}$:
\begin{ptree}
\AXC{$\forall a\,\forall \alpha\,(\mathrm{Conv}(a,\alpha) \to \mathrm{Seq}(a))$\lbl{[C1]}}
\AXC{$\varphi_{3}$\lbl{A2}}
\BIC{$\forall \alpha\,(\mathrm{Conv}(a,\alpha) \to \mathrm{Seq}(a))$}
\AXC{$\varphi_{1}$\lbl{A2}}
\BIC{$\mathrm{Conv}(a,\alpha) \to \mathrm{Seq}(a)$}
\AXC{$\stackrel{\Pi_{51}}{\vdots}$}\noLine
\UIC{$(\mathrm{Conv}(a,\alpha) \to \mathrm{Seq}(a)) \to \varphi_{140} \to \varphi_{106}$}
\BIC{$\varphi_{140} \to \varphi_{106}$}
\end{ptree}
\noindent$\Pi_{53}$:
\begin{ptree}
\AXC{$\forall M\,(\neg (M \in \mathbb{R} \land \forall n \in \mathbb{N}\,(\lvert a_{n} \rvert \leq M))) \to \neg \varphi_{23}$\lbl{A2}}
\AXC{$\stackrel{\Pi_{50}}{\vdots}$}\noLine
\UIC{$\varphi_{161}$}
\BIC{$\varphi_{140}$}
\AXC{$\stackrel{\Pi_{52}}{\vdots}$}\noLine
\UIC{$\varphi_{140} \to \varphi_{106}$}
\BIC{$\varphi_{106}$}
\AXC{$\varphi_{106} \to \varphi_{108}$\lbl{A1}}
\BIC{$\varphi_{108}$}
\UIC{$\forall M\,(\varphi_{108})$}
\end{ptree}
\noindent$\Pi_{54}$:
\begin{ptree}
\AXC{$\forall a\,(\mathrm{Seq}(a) \to \varphi_{21})$\lbl{[C14]}}
\AXC{$\forall a\,(\mathrm{Seq}(a) \to \varphi_{21}) \to \mathrm{Seq}(a) \to \varphi_{21}$\lbl{A2}}
\BIC{$\mathrm{Seq}(a) \to \varphi_{21}$}
\AXC{$\varphi_{160}$\lbl{A1}}
\BIC{$(\mathrm{Conv}(a,\alpha) \to \mathrm{Seq}(a)) \to \varphi_{113}$}
\end{ptree}
\noindent$\Pi_{55}$:
\begin{ptree}
\AXC{$\forall a\,\forall \alpha\,(\mathrm{Conv}(a,\alpha) \to \mathrm{Seq}(a))$\lbl{[C1]}}
\AXC{$\varphi_{3}$\lbl{A2}}
\BIC{$\forall \alpha\,(\mathrm{Conv}(a,\alpha) \to \mathrm{Seq}(a))$}
\AXC{$\varphi_{1}$\lbl{A2}}
\BIC{$\mathrm{Conv}(a,\alpha) \to \mathrm{Seq}(a)$}
\AXC{$\stackrel{\Pi_{54}}{\vdots}$}\noLine
\UIC{$(\mathrm{Conv}(a,\alpha) \to \mathrm{Seq}(a)) \to \varphi_{113}$}
\BIC{$\varphi_{113}$}
\end{ptree}
\noindent$\Pi_{56}$:
\begin{ptree}
\AXC{$\varphi_{21} \to \varphi_{13}$\lbl{A2}}
\AXC{$\stackrel{\Pi_{55}}{\vdots}$}\noLine
\UIC{$\varphi_{113}$}
\AXC{$\varphi_{113} \to (\varphi_{21} \to \varphi_{13}) \to \varphi_{110}$\lbl{A1}}
\BIC{$(\varphi_{21} \to \varphi_{13}) \to \varphi_{110}$}
\BIC{$\varphi_{110}$}
\AXC{$\varphi_{110} \to \varphi_{31} \to \mathrm{Conv}(a,\alpha) \to \varphi_{92}$\lbl{A1}}
\BIC{$\varphi_{31} \to \mathrm{Conv}(a,\alpha) \to \varphi_{92}$}
\end{ptree}
\noindent$\Pi_{57}$:
\begin{ptree}
\AXC{$\mathrm{Conv}(a,\alpha) \to \varphi_{89}$\lbl{A1}}
\AXC{$(\mathrm{Conv}(a,\alpha) \to \varphi_{89}) \to \varphi_{31}$\lbl{A1}}
\BIC{$\varphi_{31}$}
\AXC{$\stackrel{\Pi_{56}}{\vdots}$}\noLine
\UIC{$\varphi_{31} \to \mathrm{Conv}(a,\alpha) \to \varphi_{92}$}
\BIC{$\mathrm{Conv}(a,\alpha) \to \varphi_{92}$}
\end{ptree}
\noindent$\Pi_{58}$:
\begin{ptree}
\AXC{$\forall M\,(\varphi_{108}) \to \forall M\,(\varphi_{35}) \to \varphi_{5}$\lbl{A3}}
\AXC{$\stackrel{\Pi_{53}}{\vdots}$}\noLine
\UIC{$\forall M\,(\varphi_{108})$}
\AXC{$\stackrel{\Pi_{57}}{\vdots}$}\noLine
\UIC{$\mathrm{Conv}(a,\alpha) \to \varphi_{92}$}
\AXC{$(\mathrm{Conv}(a,\alpha) \to \varphi_{92}) \to \forall M\,(\varphi_{108}) \to \varphi_{172}$\lbl{A1}}
\BIC{$\forall M\,(\varphi_{108}) \to \varphi_{172}$}
\BIC{$\varphi_{172}$}
\BIC{$(\varphi_{35} \to \forall M\,(\varphi_{35})) \to \mathrm{Conv}(a,\alpha) \to \varphi_{12}$}
\end{ptree}
\noindent$\Pi_{59}$:
\begin{ptree}
\AXC{$\varphi_{35} \to \forall M\,(\varphi_{35})$\lbl{A4}}
\AXC{$\stackrel{\Pi_{58}}{\vdots}$}\noLine
\UIC{$(\varphi_{35} \to \forall M\,(\varphi_{35})) \to \mathrm{Conv}(a,\alpha) \to \varphi_{12}$}
\BIC{$\mathrm{Conv}(a,\alpha) \to \varphi_{12}$}
\AXC{$(\mathrm{Conv}(a,\alpha) \to \varphi_{12}) \to \varphi_{37}$\lbl{A1}}
\BIC{$\varphi_{37}$}
\UIC{$\forall N\,(\varphi_{37})$}
\end{ptree}
\noindent$\Pi_{60}$:
\begin{ptree}
\AXC{$\forall a\,\forall \alpha\,(\varphi_{76})$\lbl{[C3]}}
\AXC{$\forall a\,\forall \alpha\,(\varphi_{76}) \to \forall \alpha\,(\varphi_{76})$\lbl{A2}}
\BIC{$\forall \alpha\,(\varphi_{76})$}
\AXC{$\forall \alpha\,(\varphi_{76}) \to \varphi_{76}$\lbl{A2}}
\BIC{$\varphi_{76}$}
\AXC{$\varphi_{76} \to \mathrm{Conv}(a,\alpha) \to \varphi_{33}$\lbl{A2}}
\BIC{$\mathrm{Conv}(a,\alpha) \to \varphi_{33}$}
\end{ptree}
\noindent$\Pi_{61}$:
\begin{ptree}
\AXC{$1 \in \mathbb{R}$\lbl{(R1)}}
\AXC{$\mathrm{Conv}(a,\alpha) \to \mathrm{Conv}(a,\alpha)$\lbl{A1}}
\AXC{$\stackrel{\Pi_{60}}{\vdots}$}\noLine
\UIC{$\mathrm{Conv}(a,\alpha) \to \varphi_{33}$}
\AXC{$(\mathrm{Conv}(a,\alpha) \to \varphi_{33}) \to \varphi_{104}$\lbl{A1}}
\BIC{$\varphi_{104}$}
\BIC{$1 \in \mathbb{R} \to (\mathrm{Conv}(a,\alpha) \to 0 < 1) \to \varphi_{20}$}
\BIC{$(\mathrm{Conv}(a,\alpha) \to 0 < 1) \to \varphi_{20}$}
\end{ptree}
\noindent$\Pi_{62}$:
\begin{ptree}
\AXC{$0 < 1$\lbl{[F28]}}
\AXC{$0 < 1 \to \mathrm{Conv}(a,\alpha) \to 0 < 1$\lbl{A1}}
\BIC{$\mathrm{Conv}(a,\alpha) \to 0 < 1$}
\AXC{$\stackrel{\Pi_{61}}{\vdots}$}\noLine
\UIC{$(\mathrm{Conv}(a,\alpha) \to 0 < 1) \to \varphi_{20}$}
\BIC{$\varphi_{20}$}
\AXC{$\varphi_{168}$\lbl{A1}}
\BIC{$\forall N\,(\varphi_{37}) \to \varphi_{144} \to \varphi_{2} \to \mathrm{Conv}(a,\alpha) \to \mathrm{SeqBdd}(a)$}
\end{ptree}
\noindent 主木:
\begin{ptree}
\AXC{$\varphi_{2}$\lbl{A4}}
\AXC{$\varphi_{144}$\lbl{A3}}
\AXC{$\stackrel{\Pi_{59}}{\vdots}$}\noLine
\UIC{$\forall N\,(\varphi_{37})$}
\AXC{$\stackrel{\Pi_{62}}{\vdots}$}\noLine
\UIC{$\forall N\,(\varphi_{37}) \to \varphi_{144} \to \varphi_{2} \to \mathrm{Conv}(a,\alpha) \to \mathrm{SeqBdd}(a)$}
\BIC{$\varphi_{144} \to \varphi_{2} \to \mathrm{Conv}(a,\alpha) \to \mathrm{SeqBdd}(a)$}
\BIC{$\varphi_{2} \to \mathrm{Conv}(a,\alpha) \to \mathrm{SeqBdd}(a)$}
\BIC{$\mathrm{Conv}(a,\alpha) \to \mathrm{SeqBdd}(a)$}
\UIC{$\forall \alpha\,(\mathrm{Conv}(a,\alpha) \to \mathrm{SeqBdd}(a))$}
\UIC{$\forall a\,\forall \alpha\,(\mathrm{Conv}(a,\alpha) \to \mathrm{SeqBdd}(a))$}
\end{ptree}
\begin{abbrevlist}
\item $\varphi_{1} :\equiv \forall \alpha\,(\mathrm{Conv}(a,\alpha) \to \mathrm{Seq}(a)) \to \mathrm{Conv}(a,\alpha) \to \mathrm{Seq}(a)$
\item $\varphi_{2} :\equiv \mathrm{Conv}(a,\alpha) \land \neg \mathrm{SeqBdd}(a) \to \forall N\,(\mathrm{Conv}(a,\alpha) \land \neg \mathrm{SeqBdd}(a))$
\item $\varphi_{3} :\equiv \forall a\,\forall \alpha\,(\mathrm{Conv}(a,\alpha) \to \mathrm{Seq}(a)) \to \forall \alpha\,(\mathrm{Conv}(a,\alpha) \to \mathrm{Seq}(a))$
\item $\varphi_{4} :\equiv \forall \alpha\,(\mathrm{Conv}(a,\alpha) \to \alpha \in \mathbb{R}) \to \mathrm{Conv}(a,\alpha) \to \alpha \in \mathbb{R}$
\item $\varphi_{5} :\equiv \forall M\,(\neg (M \in \mathbb{R} \land 0 \leq M \land \forall n \in \mathbb{N}\,(n \leq N \to \lvert a_{n} \rvert \leq M)))$
\item $\varphi_{6} :\equiv \forall N\,(\neg (N \in \mathbb{N} \land \forall n \in \mathbb{N}\,(N \leq n \to \lvert a_{n} - \alpha \rvert < 1)))$
\item $\varphi_{7} :\equiv \forall a\,\forall \alpha\,(\mathrm{Conv}(a,\alpha) \to \alpha \in \mathbb{R}) \to \forall \alpha\,(\mathrm{Conv}(a,\alpha) \to \alpha \in \mathbb{R})$
\item $\varphi_{8} :\equiv \mathrm{SeqBdd}(a) \leftrightarrow \mathrm{Seq}(a) \land \exists M \in \mathbb{R}\,\forall n \in \mathbb{N}\,(\lvert a_{n} \rvert \leq M)$
\item $\varphi_{9} :\equiv M \in \mathbb{R} \land 0 \leq M \land \forall n \in \mathbb{N}\,(n \leq N \to \lvert a_{n} \rvert \leq M) \to 0 \leq \lvert \alpha \rvert$
\item $\varphi_{10} :\equiv M \in \mathbb{R} \land 0 \leq M \land \forall n \in \mathbb{N}\,(n \leq N \to \lvert a_{n} \rvert \leq M) \to \mathrm{SeqBdd}(a)$
\item $\varphi_{11} :\equiv M \in \mathbb{R} \land 0 \leq M \land \forall n \in \mathbb{N}\,(n \leq N \to \lvert a_{n} \rvert \leq M) \to M \in \mathbb{R}$
\item $\varphi_{12} :\equiv N \in \mathbb{N} \land \forall n \in \mathbb{N}\,(N \leq n \to \lvert a_{n} - \alpha \rvert < 1) \to \mathrm{SeqBdd}(a)$
\item $\varphi_{13} :\equiv N \in \mathbb{N} \to \exists M\,(M \in \mathbb{R} \land 0 \leq M \land \forall n \in \mathbb{N}\,(n \leq N \to \lvert a_{n} \rvert \leq M))$
\item $\varphi_{14} :\equiv \lvert a_{n} \rvert \leq M \to M \leq M + (1 + \lvert \alpha \rvert) \to \lvert a_{n} \rvert \leq M + (1 + \lvert \alpha \rvert)$
\item $\varphi_{15} :\equiv \forall n \in \mathbb{N}\,(n \leq N \to \lvert a_{n} \rvert \leq M) \to n \in \mathbb{N} \to n \leq N \to \lvert a_{n} \rvert \leq M$
\item $\varphi_{16} :\equiv \forall x\,\forall y\,\forall z\,(x \leq y \to y \leq z \to x \leq z) \to \forall y\,\forall z\,(\lvert a_{n} \rvert \leq y \to y \leq z \to \lvert a_{n} \rvert \leq z)$
\item $\varphi_{17} :\equiv M \in \mathbb{R} \land 0 \leq M \land \forall n \in \mathbb{N}\,(n \leq N \to \lvert a_{n} \rvert \leq M) \to 0 \leq 1 + \lvert \alpha \rvert$
\item $\varphi_{18} :\equiv M \in \mathbb{R} \to 1 + \lvert \alpha \rvert \in \mathbb{R} \to 0 \leq 1 + \lvert \alpha \rvert \to M \leq M + (1 + \lvert \alpha \rvert)$
\item $\varphi_{19} :\equiv 1 + \lvert \alpha \rvert \in \mathbb{R} \to M \in \mathbb{R} \to 0 \leq M \to 1 + \lvert \alpha \rvert \leq M + (1 + \lvert \alpha \rvert)$
\item $\varphi_{20} :\equiv \mathrm{Conv}(a,\alpha) \to \exists N \in \mathbb{N}\,\forall n \in \mathbb{N}\,(N \leq n \to \lvert a_{n} - \alpha \rvert < 1)$
\item $\varphi_{21} :\equiv \forall N \in \mathbb{N}\,\exists M\,(M \in \mathbb{R} \land 0 \leq M \land \forall n \in \mathbb{N}\,(n \leq N \to \lvert a_{n} \rvert \leq M))$
\item $\varphi_{22} :\equiv M \in \mathbb{R} \land 0 \leq M \land \forall n \in \mathbb{N}\,(n \leq N \to \lvert a_{n} \rvert \leq M) \to \lvert \alpha \rvert \in \mathbb{R}$
\item $\varphi_{23} :\equiv M + (1 + \lvert \alpha \rvert) \in \mathbb{R} \land \forall n \in \mathbb{N}\,(\lvert a_{n} \rvert \leq M + (1 + \lvert \alpha \rvert))$
\item $\varphi_{24} :\equiv \forall y\,(1 + \lvert \alpha \rvert \in \mathbb{R} \to y \in \mathbb{R} \to 0 \leq y \to 1 + \lvert \alpha \rvert \leq y + (1 + \lvert \alpha \rvert))$
\item $\varphi_{25} :\equiv M \in \mathbb{R} \land 0 \leq M \land \forall n \in \mathbb{N}\,(n \leq N \to \lvert a_{n} \rvert \leq M) \to n \in \mathbb{N} \to n \in \mathbb{N}$
\item $\varphi_{26} :\equiv M \in \mathbb{R} \land 0 \leq M \land \forall n \in \mathbb{N}\,(n \leq N \to \lvert a_{n} \rvert \leq M) \to n \in \mathbb{N} \to N \leq n \to N \leq n$
\item $\varphi_{27} :\equiv M \in \mathbb{R} \land 0 \leq M \land \forall n \in \mathbb{N}\,(n \leq N \to \lvert a_{n} \rvert \leq M) \to 1 + \lvert \alpha \rvert \in \mathbb{R}$
\item $\varphi_{28} :\equiv M \in \mathbb{R} \land 0 \leq M \land \forall n \in \mathbb{N}\,(n \leq N \to \lvert a_{n} \rvert \leq M) \to n \in \mathbb{N} \to n \in \omega$
\item $\varphi_{29} :\equiv M \in \mathbb{R} \land 0 \leq M \land \forall n \in \mathbb{N}\,(n \leq N \to \lvert a_{n} \rvert \leq M) \to n \in \mathbb{N} \to N \in \omega$
\item $\varphi_{30} :\equiv \forall n\,(\mathrm{Seq}(a) \to n \in \mathbb{N} \to a_{n} \in \mathbb{R}) \to \mathrm{Seq}(a) \to n \in \mathbb{N} \to a_{n} \in \mathbb{R}$
\item $\varphi_{31} :\equiv \mathrm{Conv}(a,\alpha) \to N \in \mathbb{N} \land \forall n \in \mathbb{N}\,(N \leq n \to \lvert a_{n} - \alpha \rvert < 1) \to N \in \mathbb{N}$
\item $\varphi_{32} :\equiv M \in \mathbb{R} \land 0 \leq M \land \forall n \in \mathbb{N}\,(n \leq N \to \lvert a_{n} \rvert \leq M) \to n \in \mathbb{N} \to N \leq n \lor n < N$
\item $\varphi_{33} :\equiv 1 \in \mathbb{R} \to 0 < 1 \to \exists N \in \mathbb{N}\,\forall n \in \mathbb{N}\,(N \leq n \to \lvert a_{n} - \alpha \rvert < 1)$
\item $\varphi_{34} :\equiv z = w \to \lvert z \rvert = \lvert a_{n} - \alpha \rvert + \lvert \alpha \rvert \to \lvert w \rvert = \lvert a_{n} - \alpha \rvert + \lvert \alpha \rvert$
\item $\varphi_{35} :\equiv \mathrm{Conv}(a,\alpha) \land (N \in \mathbb{N} \land \forall n \in \mathbb{N}\,(N \leq n \to \lvert a_{n} - \alpha \rvert < 1)) \land \neg \mathrm{SeqBdd}(a)$
\item $\varphi_{36} :\equiv \forall x\,\forall y\,\forall z\,(x \leq y \to y < z \to x < z) \to \forall y\,\forall z\,(\lvert a_{n} \rvert \leq y \to y < z \to \lvert a_{n} \rvert < z)$
\item $\varphi_{37} :\equiv \mathrm{Conv}(a,\alpha) \land \neg \mathrm{SeqBdd}(a) \to \neg (N \in \mathbb{N} \land \forall n \in \mathbb{N}\,(N \leq n \to \lvert a_{n} - \alpha \rvert < 1))$
\item $\varphi_{38} :\equiv M \in \mathbb{R} \land 0 \leq M \land \forall n \in \mathbb{N}\,(n \leq N \to \lvert a_{n} \rvert \leq M) \to M + (1 + \lvert \alpha \rvert) \in \mathbb{R}$
\item $\varphi_{39} :\equiv M \in \mathbb{R} \land 0 \leq M \land \forall n \in \mathbb{N}\,(n \leq N \to \lvert a_{n} \rvert \leq M) \to n \in \mathbb{N} \to n < N \to n < N$
\item $\varphi_{40} :\equiv \forall n\,(N \in \omega \to n \in \omega \to N \leq n \lor n < N) \to N \in \omega \to n \in \omega \to N \leq n \lor n < N$
\item $\varphi_{41} :\equiv \forall a\,\forall n\,(\mathrm{Seq}(a) \to n \in \mathbb{N} \to a_{n} \in \mathbb{R}) \to \forall n\,(\mathrm{Seq}(a) \to n \in \mathbb{N} \to a_{n} \in \mathbb{R})$
\item $\varphi_{42} :\equiv \forall x \in \mathbb{R}\,(\lvert x \rvert \in \mathbb{R}) \to a_{n} - \alpha \in \mathbb{R} \to \lvert a_{n} - \alpha \rvert \in \mathbb{R}$
\item $\varphi_{43} :\equiv z = w \to \lvert z \rvert < \lvert a_{n} - \alpha \rvert + \lvert \alpha \rvert \to \lvert w \rvert < \lvert a_{n} - \alpha \rvert + \lvert \alpha \rvert$
\item $\varphi_{44} :\equiv \forall y\,\forall z\,(\lvert a_{n} \rvert \leq y \to y \leq z \to \lvert a_{n} \rvert \leq z) \to \forall z\,(\lvert a_{n} \rvert \leq M \to M \leq z \to \lvert a_{n} \rvert \leq z)$
\item $\varphi_{45} :\equiv \forall n \in \mathbb{N}\,(N \leq n \to \lvert a_{n} - \alpha \rvert < 1) \to n \in \mathbb{N} \to N \leq n \to \lvert a_{n} - \alpha \rvert < 1$
\item $\varphi_{46} :\equiv M \in \mathbb{R} \land 0 \leq M \land \forall n \in \mathbb{N}\,(n \leq N \to \lvert a_{n} \rvert \leq M) \to n \in \mathbb{N} \to N \leq n \to \lvert \alpha \rvert \leq \lvert \alpha \rvert$
\item $\varphi_{47} :\equiv \lvert a_{n} - \alpha \rvert < 1 \to \lvert \alpha \rvert \leq \lvert \alpha \rvert \to \lvert a_{n} - \alpha \rvert + \lvert \alpha \rvert < 1 + \lvert \alpha \rvert$
\item $\varphi_{48} :\equiv \forall x\,\forall y\,(x \in \mathbb{R} \to y \in \mathbb{R} \to 0 \leq y \to x \leq x + y) \to \forall y\,(M \in \mathbb{R} \to y \in \mathbb{R} \to 0 \leq y \to M \leq M + y)$
\item $\varphi_{49} :\equiv \forall m\,\forall n\,(m \in \omega \to n \in \omega \to m \leq n \lor n < m) \to \forall n\,(N \in \omega \to n \in \omega \to N \leq n \lor n < N)$
\item $\varphi_{50} :\equiv M \in \mathbb{R} \land 0 \leq M \land \forall n \in \mathbb{N}\,(n \leq N \to \lvert a_{n} \rvert \leq M) \to n \in \mathbb{N} \to N \leq n \to \lvert \alpha \rvert \in \mathbb{R}$
\item $\varphi_{51} :\equiv \lvert a_{n} \rvert \leq 1 + \lvert \alpha \rvert \to 1 + \lvert \alpha \rvert \leq M + (1 + \lvert \alpha \rvert) \to \lvert a_{n} \rvert \leq M + (1 + \lvert \alpha \rvert)$
\item $\varphi_{52} :\equiv M \in \mathbb{R} \land 0 \leq M \land \forall n \in \mathbb{N}\,(n \leq N \to \lvert a_{n} \rvert \leq M) \to n \in \mathbb{N} \to N \leq n \to a_{n} \in \mathbb{R}$
\item $\varphi_{53} :\equiv \forall x\,\forall y\,(x \in \mathbb{R} \to y \in \mathbb{R} \to x + y \in \mathbb{R}) \to \forall y\,(1 \in \mathbb{R} \to y \in \mathbb{R} \to 1 + y \in \mathbb{R})$
\item $\varphi_{54} :\equiv \forall x\,\forall y\,(x \in \mathbb{R} \to y \in \mathbb{R} \to x + y \in \mathbb{R}) \to \forall y\,(M \in \mathbb{R} \to y \in \mathbb{R} \to M + y \in \mathbb{R})$
\item $\varphi_{55} :\equiv M \in \mathbb{R} \land 0 \leq M \land \forall n \in \mathbb{N}\,(n \leq N \to \lvert a_{n} \rvert \leq M) \to n \in \mathbb{N} \to n < N \to \lvert a_{n} \rvert \leq M$
\item $\varphi_{56} :\equiv \forall x\,\forall y\,\forall u\,\forall v\,(x \in \mathbb{R} \to y \in \mathbb{R} \to u \in \mathbb{R} \to v \in \mathbb{R} \to x < y \to u \leq v \to x + u < y + v)$
\item $\varphi_{57} :\equiv M \in \mathbb{R} \land 0 \leq M \land \forall n \in \mathbb{N}\,(n \leq N \to \lvert a_{n} \rvert \leq M) \to n \in \mathbb{N} \to n < N \to \lvert \alpha \rvert \in \mathbb{R}$
\item $\varphi_{58} :\equiv \forall b\,(a_{n} - \alpha \in \mathbb{R} \to b \in \mathbb{R} \to \lvert a_{n} - \alpha + b \rvert \leq \lvert a_{n} - \alpha \rvert + \lvert b \rvert)$
\item $\varphi_{59} :\equiv a_{n} - \alpha \in \mathbb{R} \to \alpha \in \mathbb{R} \to \lvert a_{n} - \alpha + \alpha \rvert \leq \lvert a_{n} - \alpha \rvert + \lvert \alpha \rvert$
\item $\varphi_{60} :\equiv M \in \mathbb{R} \land 0 \leq M \land \forall n \in \mathbb{N}\,(n \leq N \to \lvert a_{n} \rvert \leq M) \to n \in \mathbb{N} \to N \leq n \to 1 + \lvert \alpha \rvert \in \mathbb{R}$
\item $\varphi_{61} :\equiv M \in \mathbb{R} \land 0 \leq M \land \forall n \in \mathbb{N}\,(n \leq N \to \lvert a_{n} \rvert \leq M) \to n \in \mathbb{N} \to n < N \to M \leq M + (1 + \lvert \alpha \rvert)$
\item $\varphi_{62} :\equiv M \in \mathbb{R} \land 0 \leq M \land \forall n \in \mathbb{N}\,(n \leq N \to \lvert a_{n} \rvert \leq M) \to n \in \mathbb{N} \to N \leq n \to a_{n} - \alpha \in \mathbb{R}$
\item $\varphi_{63} :\equiv \forall y\,(1 \in \mathbb{R} \to y \in \mathbb{R} \to 1 + y \in \mathbb{R}) \to 1 \in \mathbb{R} \to \lvert \alpha \rvert \in \mathbb{R} \to 1 + \lvert \alpha \rvert \in \mathbb{R}$
\item $\varphi_{64} :\equiv M \in \mathbb{R} \land 0 \leq M \land \forall n \in \mathbb{N}\,(n \leq N \to \lvert a_{n} \rvert \leq M) \to n \in \mathbb{N} \to n < N \to 1 + \lvert \alpha \rvert \in \mathbb{R}$
\item $\varphi_{65} :\equiv M \in \mathbb{R} \land 0 \leq M \land \forall n \in \mathbb{N}\,(n \leq N \to \lvert a_{n} \rvert \leq M) \to n \in \mathbb{N} \to N \leq n \to \lvert a_{n} - \alpha \rvert < 1$
\item $\varphi_{66} :\equiv M \in \mathbb{R} \land 0 \leq M \land \forall n \in \mathbb{N}\,(n \leq N \to \lvert a_{n} \rvert \leq M) \to n \in \mathbb{N} \to \lvert a_{n} \rvert \leq M + (1 + \lvert \alpha \rvert)$
\item $\varphi_{67} :\equiv M \in \mathbb{R} \land 0 \leq M \land \forall n \in \mathbb{N}\,(n \leq N \to \lvert a_{n} \rvert \leq M) \to n \in \mathbb{N} \to N \leq n \to \lvert a_{n} \rvert \leq 1 + \lvert \alpha \rvert$
\item $\varphi_{68} :\equiv \lvert a_{n} - \alpha + \alpha \rvert = \lvert a_{n} - \alpha \rvert + \lvert \alpha \rvert \to \lvert a_{n} \rvert = \lvert a_{n} - \alpha \rvert + \lvert \alpha \rvert$
\item $\varphi_{69} :\equiv \forall x\,\forall y\,(x \in \mathbb{R} \to y \in \mathbb{R} \to 0 \leq x \to 0 \leq y \to 0 \leq x + y) \to \forall y\,(1 \in \mathbb{R} \to y \in \mathbb{R} \to 0 \leq 1 \to 0 \leq y \to 0 \leq 1 + y)$
\item $\varphi_{70} :\equiv M \in \mathbb{R} \land 0 \leq M \land \forall n \in \mathbb{N}\,(n \leq N \to \lvert a_{n} \rvert \leq M) \to n \in \mathbb{N} \to N \leq n \to \lvert a_{n} - \alpha \rvert \in \mathbb{R}$
\item $\varphi_{71} :\equiv M \in \mathbb{R} \land 0 \leq M \land \forall n \in \mathbb{N}\,(n \leq N \to \lvert a_{n} \rvert \leq M) \to \forall n \in \mathbb{N}\,(\lvert a_{n} \rvert \leq M + (1 + \lvert \alpha \rvert))$
\item $\varphi_{72} :\equiv M \in \mathbb{R} \land 0 \leq M \land \forall n \in \mathbb{N}\,(n \leq N \to \lvert a_{n} \rvert \leq M) \to \exists M \in \mathbb{R}\,\forall n \in \mathbb{N}\,(\lvert a_{n} \rvert \leq M)$
\item $\varphi_{73} :\equiv M \in \mathbb{R} \land 0 \leq M \land \forall n \in \mathbb{N}\,(n \leq N \to \lvert a_{n} \rvert \leq M) \to n \in \mathbb{N} \to N \leq n \to a_{n} - \alpha + \alpha = a_{n}$
\item $\varphi_{74} :\equiv M \in \mathbb{R} \land 0 \leq M \land \forall n \in \mathbb{N}\,(n \leq N \to \lvert a_{n} \rvert \leq M) \to n \in \mathbb{N} \to N \leq n \to \lvert a_{n} \rvert < 1 + \lvert \alpha \rvert$
\item $\varphi_{75} :\equiv \forall z\,(\lvert a_{n} \rvert \leq \lvert a_{n} - \alpha \rvert + \lvert \alpha \rvert \to \lvert a_{n} - \alpha \rvert + \lvert \alpha \rvert < z \to \lvert a_{n} \rvert < z)$
\item $\varphi_{76} :\equiv \forall \varepsilon\,(\mathrm{Conv}(a,\alpha) \to \varepsilon \in \mathbb{R} \to 0 < \varepsilon \to \exists N \in \mathbb{N}\,\forall n \in \mathbb{N}\,(N \leq n \to \lvert a_{n} - \alpha \rvert < \varepsilon))$
\item $\varphi_{77} :\equiv M \in \mathbb{R} \land 0 \leq M \land \forall n \in \mathbb{N}\,(n \leq N \to \lvert a_{n} \rvert \leq M) \to n \in \mathbb{N} \to N \leq n \to 1 + \lvert \alpha \rvert \leq M + (1 + \lvert \alpha \rvert)$
\item $\varphi_{78} :\equiv M \in \mathbb{R} \land 0 \leq M \land \forall n \in \mathbb{N}\,(n \leq N \to \lvert a_{n} \rvert \leq M) \to n \in \mathbb{N} \to N \leq n \to \lvert a_{n} \rvert \leq M + (1 + \lvert \alpha \rvert)$
\item $\varphi_{79} :\equiv \forall x\,\forall y\,(x \in \mathbb{R} \to y \in \mathbb{R} \to x - y \in \mathbb{R}) \to \forall y\,(a_{n} \in \mathbb{R} \to y \in \mathbb{R} \to a_{n} - y \in \mathbb{R})$
\item $\varphi_{80} :\equiv \lvert a_{n} - \alpha + \alpha \rvert < \lvert a_{n} - \alpha \rvert + \lvert \alpha \rvert \to \lvert a_{n} \rvert < \lvert a_{n} - \alpha \rvert + \lvert \alpha \rvert$
\item $\varphi_{81} :\equiv \forall y\,(M \in \mathbb{R} \to y \in \mathbb{R} \to M + y \in \mathbb{R}) \to M \in \mathbb{R} \to 1 + \lvert \alpha \rvert \in \mathbb{R} \to M + (1 + \lvert \alpha \rvert) \in \mathbb{R}$
\item $\varphi_{82} :\equiv M \in \mathbb{R} \land 0 \leq M \land \forall n \in \mathbb{N}\,(n \leq N \to \lvert a_{n} \rvert \leq M) \to n \in \mathbb{N} \to n < N \to \lvert a_{n} \rvert \leq M + (1 + \lvert \alpha \rvert)$
\item $\varphi_{83} :\equiv \forall x\,\forall y\,(x \in \mathbb{R} \to y \in \mathbb{R} \to x - y + y = x) \to \forall y\,(a_{n} \in \mathbb{R} \to y \in \mathbb{R} \to a_{n} - y + y = a_{n})$
\item $\varphi_{84} :\equiv \forall y\,\forall z\,(\lvert a_{n} \rvert \leq y \to y \leq z \to \lvert a_{n} \rvert \leq z) \to \forall z\,(\lvert a_{n} \rvert \leq 1 + \lvert \alpha \rvert \to 1 + \lvert \alpha \rvert \leq z \to \lvert a_{n} \rvert \leq z)$
\item $\varphi_{85} :\equiv M \in \mathbb{R} \land 0 \leq M \land \forall n \in \mathbb{N}\,(n \leq N \to \lvert a_{n} \rvert \leq M) \to n \in \mathbb{N} \to n < N \to n \in \mathbb{N} \to n \leq N \to \lvert a_{n} \rvert \leq M$
\item $\varphi_{86} :\equiv M \in \mathbb{R} \land 0 \leq M \land \forall n \in \mathbb{N}\,(n \leq N \to \lvert a_{n} \rvert \leq M) \to M \in \mathbb{R} \land 0 \leq M \land \forall n \in \mathbb{N}\,(n \leq N \to \lvert a_{n} \rvert \leq M)$
\item $\varphi_{87} :\equiv M \in \mathbb{R} \land 0 \leq M \land \forall n \in \mathbb{N}\,(n \leq N \to \lvert a_{n} \rvert \leq M) \to n \in \mathbb{N} \to n < N \to \forall n \in \mathbb{N}\,(n \leq N \to \lvert a_{n} \rvert \leq M)$
\item $\varphi_{88} :\equiv \forall y\,(1 \in \mathbb{R} \to y \in \mathbb{R} \to 0 \leq 1 \to 0 \leq y \to 0 \leq 1 + y) \to 1 \in \mathbb{R} \to \lvert \alpha \rvert \in \mathbb{R} \to 0 \leq 1 \to 0 \leq \lvert \alpha \rvert \to 0 \leq 1 + \lvert \alpha \rvert$
\item $\varphi_{89} :\equiv N \in \mathbb{N} \land \forall n \in \mathbb{N}\,(N \leq n \to \lvert a_{n} - \alpha \rvert < 1) \to N \in \mathbb{N} \land \forall n \in \mathbb{N}\,(N \leq n \to \lvert a_{n} - \alpha \rvert < 1)$
\item $\varphi_{90} :\equiv \forall y\,(a_{n} \in \mathbb{R} \to y \in \mathbb{R} \to a_{n} - y \in \mathbb{R}) \to a_{n} \in \mathbb{R} \to \alpha \in \mathbb{R} \to a_{n} - \alpha \in \mathbb{R}$
\item $\varphi_{91} :\equiv M \in \mathbb{R} \land 0 \leq M \land \forall n \in \mathbb{N}\,(n \leq N \to \lvert a_{n} \rvert \leq M) \to n \in \mathbb{N} \to N \leq n \to n \in \mathbb{N} \to N \leq n \to \lvert a_{n} - \alpha \rvert < 1$
\item $\varphi_{92} :\equiv N \in \mathbb{N} \land \forall n \in \mathbb{N}\,(N \leq n \to \lvert a_{n} - \alpha \rvert < 1) \to \exists M\,(M \in \mathbb{R} \land 0 \leq M \land \forall n \in \mathbb{N}\,(n \leq N \to \lvert a_{n} \rvert \leq M))$
\item $\varphi_{93} :\equiv M \in \mathbb{R} \land 0 \leq M \land \forall n \in \mathbb{N}\,(n \leq N \to \lvert a_{n} \rvert \leq M) \to n \in \mathbb{N} \to N \leq n \to \lvert a_{n} \rvert \leq \lvert a_{n} - \alpha \rvert + \lvert \alpha \rvert$
\item $\varphi_{94} :\equiv N \in \mathbb{N} \land \forall n \in \mathbb{N}\,(N \leq n \to \lvert a_{n} - \alpha \rvert < 1) \to M \in \mathbb{R} \land 0 \leq M \land \forall n \in \mathbb{N}\,(n \leq N \to \lvert a_{n} \rvert \leq M) \to 0 \leq 1$
\item $\varphi_{95} :\equiv N \in \mathbb{N} \land \forall n \in \mathbb{N}\,(N \leq n \to \lvert a_{n} - \alpha \rvert < 1) \to M \in \mathbb{R} \land 0 \leq M \land \forall n \in \mathbb{N}\,(n \leq N \to \lvert a_{n} \rvert \leq M) \to 0 \leq M$
\item $\varphi_{96} :\equiv M \in \mathbb{R} \land 0 \leq M \land \forall n \in \mathbb{N}\,(n \leq N \to \lvert a_{n} \rvert \leq M) \to n \in \mathbb{N} \to N \leq n \to \forall n \in \mathbb{N}\,(N \leq n \to \lvert a_{n} - \alpha \rvert < 1)$
\item $\varphi_{97} :\equiv M \in \mathbb{R} \land 0 \leq M \land \forall n \in \mathbb{N}\,(n \leq N \to \lvert a_{n} \rvert \leq M) \to n \in \mathbb{N} \to N \leq n \to \lvert a_{n} - \alpha \rvert + \lvert \alpha \rvert < 1 + \lvert \alpha \rvert$
\item $\varphi_{98} :\equiv \forall w\,(a_{n} - \alpha + \alpha = w \to \lvert a_{n} - \alpha + \alpha \rvert = \lvert a_{n} - \alpha \rvert + \lvert \alpha \rvert \to \lvert w \rvert = \lvert a_{n} - \alpha \rvert + \lvert \alpha \rvert)$
\item $\varphi_{99} :\equiv \lvert a_{n} \rvert \leq \lvert a_{n} - \alpha \rvert + \lvert \alpha \rvert \to \lvert a_{n} - \alpha \rvert + \lvert \alpha \rvert < 1 + \lvert \alpha \rvert \to \lvert a_{n} \rvert < 1 + \lvert \alpha \rvert$
\item $\varphi_{100} :\equiv \mathrm{Conv}(a,\alpha) \land (N \in \mathbb{N} \land \forall n \in \mathbb{N}\,(N \leq n \to \lvert a_{n} - \alpha \rvert < 1)) \land (M \in \mathbb{R} \land 0 \leq M \land \forall n \in \mathbb{N}\,(n \leq N \to \lvert a_{n} \rvert \leq M))$
\item $\varphi_{101} :\equiv \forall y\,(a_{n} \in \mathbb{R} \to y \in \mathbb{R} \to a_{n} - y + y = a_{n}) \to a_{n} \in \mathbb{R} \to \alpha \in \mathbb{R} \to a_{n} - \alpha + \alpha = a_{n}$
\item $\varphi_{102} :\equiv \forall w\,(a_{n} - \alpha + \alpha = w \to \lvert a_{n} - \alpha + \alpha \rvert < \lvert a_{n} - \alpha \rvert + \lvert \alpha \rvert \to \lvert w \rvert < \lvert a_{n} - \alpha \rvert + \lvert \alpha \rvert)$
\item $\varphi_{103} :\equiv M \in \mathbb{R} \land 0 \leq M \land \forall n \in \mathbb{N}\,(n \leq N \to \lvert a_{n} \rvert \leq M) \to n \in \mathbb{N} \to N \leq n \to \lvert a_{n} - \alpha + \alpha \rvert \leq \lvert a_{n} - \alpha \rvert + \lvert \alpha \rvert$
\item $\varphi_{104} :\equiv (\mathrm{Conv}(a,\alpha) \to \mathrm{Conv}(a,\alpha)) \to 1 \in \mathbb{R} \to (\mathrm{Conv}(a,\alpha) \to 0 < 1) \to \varphi_{20}$
\item $\varphi_{105} :\equiv \mathrm{Conv}(a,\alpha) \to N \in \mathbb{N} \land \forall n \in \mathbb{N}\,(N \leq n \to \lvert a_{n} - \alpha \rvert < 1) \to \varphi_{9}$
\item $\varphi_{106} :\equiv \mathrm{Conv}(a,\alpha) \to N \in \mathbb{N} \land \forall n \in \mathbb{N}\,(N \leq n \to \lvert a_{n} - \alpha \rvert < 1) \to \varphi_{10}$
\item $\varphi_{107} :\equiv \mathrm{Conv}(a,\alpha) \to N \in \mathbb{N} \land \forall n \in \mathbb{N}\,(N \leq n \to \lvert a_{n} - \alpha \rvert < 1) \to \varphi_{11}$
\item $\varphi_{108} :\equiv \varphi_{35} \to \neg (M \in \mathbb{R} \land 0 \leq M \land \forall n \in \mathbb{N}\,(n \leq N \to \lvert a_{n} \rvert \leq M))$
\item $\varphi_{109} :\equiv \forall u\,\forall v\,(\lvert a_{n} - \alpha \rvert \in \mathbb{R} \to 1 \in \mathbb{R} \to u \in \mathbb{R} \to v \in \mathbb{R} \to \lvert a_{n} - \alpha \rvert < 1 \to u \leq v \to \lvert a_{n} - \alpha \rvert + u < 1 + v)$
\item $\varphi_{110} :\equiv \mathrm{Conv}(a,\alpha) \to N \in \mathbb{N} \land \forall n \in \mathbb{N}\,(N \leq n \to \lvert a_{n} - \alpha \rvert < 1) \to \varphi_{13}$
\item $\varphi_{111} :\equiv \mathrm{Conv}(a,\alpha) \to N \in \mathbb{N} \land \forall n \in \mathbb{N}\,(N \leq n \to \lvert a_{n} - \alpha \rvert < 1) \to \varphi_{17}$
\item $\varphi_{112} :\equiv \forall y\,\forall u\,\forall v\,(\lvert a_{n} - \alpha \rvert \in \mathbb{R} \to y \in \mathbb{R} \to u \in \mathbb{R} \to v \in \mathbb{R} \to \lvert a_{n} - \alpha \rvert < y \to u \leq v \to \lvert a_{n} - \alpha \rvert + u < y + v)$
\item $\varphi_{113} :\equiv \mathrm{Conv}(a,\alpha) \to N \in \mathbb{N} \land \forall n \in \mathbb{N}\,(N \leq n \to \lvert a_{n} - \alpha \rvert < 1) \to \varphi_{21}$
\item $\varphi_{114} :\equiv \mathrm{Conv}(a,\alpha) \to N \in \mathbb{N} \land \forall n \in \mathbb{N}\,(N \leq n \to \lvert a_{n} - \alpha \rvert < 1) \to \varphi_{22}$
\item $\varphi_{115} :\equiv \mathrm{Conv}(a,\alpha) \to N \in \mathbb{N} \land \forall n \in \mathbb{N}\,(N \leq n \to \lvert a_{n} - \alpha \rvert < 1) \to \varphi_{25}$
\item $\varphi_{116} :\equiv \mathrm{Conv}(a,\alpha) \to N \in \mathbb{N} \land \forall n \in \mathbb{N}\,(N \leq n \to \lvert a_{n} - \alpha \rvert < 1) \to \varphi_{26}$
\item $\varphi_{117} :\equiv \mathrm{Conv}(a,\alpha) \to N \in \mathbb{N} \land \forall n \in \mathbb{N}\,(N \leq n \to \lvert a_{n} - \alpha \rvert < 1) \to \varphi_{27}$
\item $\varphi_{118} :\equiv \mathrm{Conv}(a,\alpha) \to N \in \mathbb{N} \land \forall n \in \mathbb{N}\,(N \leq n \to \lvert a_{n} - \alpha \rvert < 1) \to \varphi_{28}$
\item $\varphi_{119} :\equiv \mathrm{Conv}(a,\alpha) \to N \in \mathbb{N} \land \forall n \in \mathbb{N}\,(N \leq n \to \lvert a_{n} - \alpha \rvert < 1) \to \varphi_{29}$
\item $\varphi_{120} :\equiv \mathrm{Conv}(a,\alpha) \to N \in \mathbb{N} \land \forall n \in \mathbb{N}\,(N \leq n \to \lvert a_{n} - \alpha \rvert < 1) \to \varphi_{32}$
\item $\varphi_{121} :\equiv \mathrm{Conv}(a,\alpha) \to N \in \mathbb{N} \land \forall n \in \mathbb{N}\,(N \leq n \to \lvert a_{n} - \alpha \rvert < 1) \to \varphi_{38}$
\item $\varphi_{122} :\equiv \mathrm{Conv}(a,\alpha) \to N \in \mathbb{N} \land \forall n \in \mathbb{N}\,(N \leq n \to \lvert a_{n} - \alpha \rvert < 1) \to \varphi_{39}$
\item $\varphi_{123} :\equiv \forall v\,(\lvert a_{n} - \alpha \rvert \in \mathbb{R} \to 1 \in \mathbb{R} \to \lvert \alpha \rvert \in \mathbb{R} \to v \in \mathbb{R} \to \lvert a_{n} - \alpha \rvert < 1 \to \lvert \alpha \rvert \leq v \to \lvert a_{n} - \alpha \rvert + \lvert \alpha \rvert < 1 + v)$
\item $\varphi_{124} :\equiv \forall a\,\forall b\,(a \in \mathbb{R} \to b \in \mathbb{R} \to \lvert a + b \rvert \leq \lvert a \rvert + \lvert b \rvert) \to \varphi_{58}$
\item $\varphi_{125} :\equiv \forall z\,(\lvert a_{n} \rvert \leq 1 + \lvert \alpha \rvert \to 1 + \lvert \alpha \rvert \leq z \to \lvert a_{n} \rvert \leq z) \to \varphi_{51}$
\item $\varphi_{126} :\equiv \mathrm{Conv}(a,\alpha) \to N \in \mathbb{N} \land \forall n \in \mathbb{N}\,(N \leq n \to \lvert a_{n} - \alpha \rvert < 1) \to \varphi_{46}$
\item $\varphi_{127} :\equiv \mathrm{Conv}(a,\alpha) \to N \in \mathbb{N} \land \forall n \in \mathbb{N}\,(N \leq n \to \lvert a_{n} - \alpha \rvert < 1) \to \varphi_{50}$
\item $\varphi_{128} :\equiv \mathrm{Conv}(a,\alpha) \to N \in \mathbb{N} \land \forall n \in \mathbb{N}\,(N \leq n \to \lvert a_{n} - \alpha \rvert < 1) \to \varphi_{52}$
\item $\varphi_{129} :\equiv \mathrm{Conv}(a,\alpha) \to N \in \mathbb{N} \land \forall n \in \mathbb{N}\,(N \leq n \to \lvert a_{n} - \alpha \rvert < 1) \to \varphi_{55}$
\item $\varphi_{130} :\equiv \mathrm{Conv}(a,\alpha) \to N \in \mathbb{N} \land \forall n \in \mathbb{N}\,(N \leq n \to \lvert a_{n} - \alpha \rvert < 1) \to \varphi_{57}$
\item $\varphi_{131} :\equiv \mathrm{Conv}(a,\alpha) \to N \in \mathbb{N} \land \forall n \in \mathbb{N}\,(N \leq n \to \lvert a_{n} - \alpha \rvert < 1) \to \varphi_{60}$
\item $\varphi_{132} :\equiv \mathrm{Conv}(a,\alpha) \to N \in \mathbb{N} \land \forall n \in \mathbb{N}\,(N \leq n \to \lvert a_{n} - \alpha \rvert < 1) \to \varphi_{61}$
\item $\varphi_{133} :\equiv \mathrm{Conv}(a,\alpha) \to N \in \mathbb{N} \land \forall n \in \mathbb{N}\,(N \leq n \to \lvert a_{n} - \alpha \rvert < 1) \to \varphi_{62}$
\item $\varphi_{134} :\equiv \mathrm{Conv}(a,\alpha) \to N \in \mathbb{N} \land \forall n \in \mathbb{N}\,(N \leq n \to \lvert a_{n} - \alpha \rvert < 1) \to \varphi_{64}$
\item $\varphi_{135} :\equiv \mathrm{Conv}(a,\alpha) \to N \in \mathbb{N} \land \forall n \in \mathbb{N}\,(N \leq n \to \lvert a_{n} - \alpha \rvert < 1) \to \varphi_{65}$
\item $\varphi_{136} :\equiv \mathrm{Conv}(a,\alpha) \to N \in \mathbb{N} \land \forall n \in \mathbb{N}\,(N \leq n \to \lvert a_{n} - \alpha \rvert < 1) \to \varphi_{66}$
\item $\varphi_{137} :\equiv \mathrm{Conv}(a,\alpha) \to N \in \mathbb{N} \land \forall n \in \mathbb{N}\,(N \leq n \to \lvert a_{n} - \alpha \rvert < 1) \to \varphi_{67}$
\item $\varphi_{138} :\equiv \mathrm{Conv}(a,\alpha) \to N \in \mathbb{N} \land \forall n \in \mathbb{N}\,(N \leq n \to \lvert a_{n} - \alpha \rvert < 1) \to \varphi_{70}$
\item $\varphi_{139} :\equiv \mathrm{Conv}(a,\alpha) \to N \in \mathbb{N} \land \forall n \in \mathbb{N}\,(N \leq n \to \lvert a_{n} - \alpha \rvert < 1) \to \varphi_{71}$
\item $\varphi_{140} :\equiv \mathrm{Conv}(a,\alpha) \to N \in \mathbb{N} \land \forall n \in \mathbb{N}\,(N \leq n \to \lvert a_{n} - \alpha \rvert < 1) \to \varphi_{72}$
\item $\varphi_{141} :\equiv \mathrm{Conv}(a,\alpha) \to N \in \mathbb{N} \land \forall n \in \mathbb{N}\,(N \leq n \to \lvert a_{n} - \alpha \rvert < 1) \to \varphi_{73}$
\item $\varphi_{142} :\equiv \mathrm{Conv}(a,\alpha) \to N \in \mathbb{N} \land \forall n \in \mathbb{N}\,(N \leq n \to \lvert a_{n} - \alpha \rvert < 1) \to \varphi_{74}$
\item $\varphi_{143} :\equiv \mathrm{Conv}(a,\alpha) \to N \in \mathbb{N} \land \forall n \in \mathbb{N}\,(N \leq n \to \lvert a_{n} - \alpha \rvert < 1) \to \varphi_{77}$
\item $\varphi_{144} :\equiv \forall N\,(\varphi_{37}) \to \forall N\,(\mathrm{Conv}(a,\alpha) \land \neg \mathrm{SeqBdd}(a)) \to \varphi_{6}$
\item $\varphi_{145} :\equiv \mathrm{Conv}(a,\alpha) \to N \in \mathbb{N} \land \forall n \in \mathbb{N}\,(N \leq n \to \lvert a_{n} - \alpha \rvert < 1) \to \varphi_{78}$
\item $\varphi_{146} :\equiv \mathrm{Conv}(a,\alpha) \to N \in \mathbb{N} \land \forall n \in \mathbb{N}\,(N \leq n \to \lvert a_{n} - \alpha \rvert < 1) \to \varphi_{82}$
\item $\varphi_{147} :\equiv \mathrm{Conv}(a,\alpha) \to N \in \mathbb{N} \land \forall n \in \mathbb{N}\,(N \leq n \to \lvert a_{n} - \alpha \rvert < 1) \to \varphi_{85}$
\item $\varphi_{148} :\equiv \mathrm{Conv}(a,\alpha) \to N \in \mathbb{N} \land \forall n \in \mathbb{N}\,(N \leq n \to \lvert a_{n} - \alpha \rvert < 1) \to \varphi_{86}$
\item $\varphi_{149} :\equiv \mathrm{Conv}(a,\alpha) \to N \in \mathbb{N} \land \forall n \in \mathbb{N}\,(N \leq n \to \lvert a_{n} - \alpha \rvert < 1) \to \varphi_{87}$
\item $\varphi_{150} :\equiv \mathrm{Conv}(a,\alpha) \to N \in \mathbb{N} \land \forall n \in \mathbb{N}\,(N \leq n \to \lvert a_{n} - \alpha \rvert < 1) \to \varphi_{91}$
\item $\varphi_{151} :\equiv (\alpha \in \mathbb{R} \to 0 \leq \lvert \alpha \rvert) \to (\mathrm{Conv}(a,\alpha) \to \alpha \in \mathbb{R}) \to \varphi_{105}$
\item $\varphi_{152} :\equiv \mathrm{Conv}(a,\alpha) \to N \in \mathbb{N} \land \forall n \in \mathbb{N}\,(N \leq n \to \lvert a_{n} - \alpha \rvert < 1) \to \varphi_{93}$
\item $\varphi_{153} :\equiv N \in \mathbb{N} \land \forall n \in \mathbb{N}\,(N \leq n \to \lvert a_{n} - \alpha \rvert < 1) \to M \in \mathbb{R} \land 0 \leq M \land \forall n \in \mathbb{N}\,(n \leq N \to \lvert a_{n} \rvert \leq M) \to \varphi_{23}$
\item $\varphi_{154} :\equiv \mathrm{Conv}(a,\alpha) \to N \in \mathbb{N} \land \forall n \in \mathbb{N}\,(N \leq n \to \lvert a_{n} - \alpha \rvert < 1) \to \varphi_{96}$
\item $\varphi_{155} :\equiv \mathrm{Conv}(a,\alpha) \to N \in \mathbb{N} \land \forall n \in \mathbb{N}\,(N \leq n \to \lvert a_{n} - \alpha \rvert < 1) \to \varphi_{97}$
\item $\varphi_{156} :\equiv \mathrm{Conv}(a,\alpha) \to N \in \mathbb{N} \land \forall n \in \mathbb{N}\,(N \leq n \to \lvert a_{n} - \alpha \rvert < 1) \to \varphi_{103}$
\item $\varphi_{157} :\equiv (\alpha \in \mathbb{R} \to \lvert \alpha \rvert \in \mathbb{R}) \to (\mathrm{Conv}(a,\alpha) \to \alpha \in \mathbb{R}) \to \varphi_{114}$
\item $\varphi_{158} :\equiv (\alpha \in \mathbb{R} \to \lvert \alpha \rvert \in \mathbb{R}) \to (\mathrm{Conv}(a,\alpha) \to \alpha \in \mathbb{R}) \to \varphi_{127}$
\item $\varphi_{159} :\equiv (\alpha \in \mathbb{R} \to \lvert \alpha \rvert \in \mathbb{R}) \to (\mathrm{Conv}(a,\alpha) \to \alpha \in \mathbb{R}) \to \varphi_{130}$
\item $\varphi_{160} :\equiv (\mathrm{Seq}(a) \to \varphi_{21}) \to (\mathrm{Conv}(a,\alpha) \to \mathrm{Seq}(a)) \to \varphi_{113}$
\item $\varphi_{161} :\equiv (\forall M\,(\neg (M \in \mathbb{R} \land \forall n \in \mathbb{N}\,(\lvert a_{n} \rvert \leq M))) \to \neg \varphi_{23}) \to \varphi_{140}$
\item $\varphi_{162} :\equiv \varphi_{123} \to \lvert a_{n} - \alpha \rvert \in \mathbb{R} \to 1 \in \mathbb{R} \to \lvert \alpha \rvert \in \mathbb{R} \to \lvert \alpha \rvert \in \mathbb{R} \to \varphi_{47}$
\item $\varphi_{163} :\equiv \forall n\,(\varphi_{100} \to n \in \mathbb{N} \to \lvert a_{n} \rvert \leq M + (1 + \lvert \alpha \rvert)) \to \forall n\,(\varphi_{100}) \to \forall n \in \mathbb{N}\,(\lvert a_{n} \rvert \leq M + (1 + \lvert \alpha \rvert))$
\item $\varphi_{164} :\equiv (1 \in \mathbb{R} \to \lvert \alpha \rvert \in \mathbb{R} \to 1 + \lvert \alpha \rvert \in \mathbb{R}) \to 1 \in \mathbb{R} \to \varphi_{114} \to \varphi_{117}$
\item $\varphi_{165} :\equiv (\mathrm{Seq}(a) \to n \in \mathbb{N} \to a_{n} \in \mathbb{R}) \to (\mathrm{Conv}(a,\alpha) \to \mathrm{Seq}(a)) \to \varphi_{115} \to \varphi_{128}$
\item $\varphi_{166} :\equiv (1 \in \mathbb{R} \to \lvert \alpha \rvert \in \mathbb{R} \to 1 + \lvert \alpha \rvert \in \mathbb{R}) \to 1 \in \mathbb{R} \to \varphi_{127} \to \varphi_{131}$
\item $\varphi_{167} :\equiv (1 \in \mathbb{R} \to \lvert \alpha \rvert \in \mathbb{R} \to 1 + \lvert \alpha \rvert \in \mathbb{R}) \to 1 \in \mathbb{R} \to \varphi_{130} \to \varphi_{134}$
\item $\varphi_{168} :\equiv \varphi_{20} \to \forall N\,(\varphi_{37}) \to \varphi_{144} \to \varphi_{2} \to \mathrm{Conv}(a,\alpha) \to \mathrm{SeqBdd}(a)$
\item $\varphi_{169} :\equiv (a_{n} \in \mathbb{R} \to \alpha \in \mathbb{R} \to a_{n} - \alpha \in \mathbb{R}) \to \varphi_{128} \to (\mathrm{Conv}(a,\alpha) \to \alpha \in \mathbb{R}) \to \varphi_{133}$
\item $\varphi_{170} :\equiv (a_{n} \in \mathbb{R} \to \alpha \in \mathbb{R} \to a_{n} - \alpha + \alpha = a_{n}) \to \varphi_{128} \to (\mathrm{Conv}(a,\alpha) \to \alpha \in \mathbb{R}) \to \varphi_{141}$
\item $\varphi_{171} :\equiv (M \in \mathbb{R} \to 1 + \lvert \alpha \rvert \in \mathbb{R} \to M + (1 + \lvert \alpha \rvert) \in \mathbb{R}) \to \varphi_{107} \to \varphi_{117} \to \varphi_{121}$
\item $\varphi_{172} :\equiv (\forall M\,(\varphi_{108}) \to \forall M\,(\varphi_{35}) \to \varphi_{5}) \to (\varphi_{35} \to \forall M\,(\varphi_{35})) \to \mathrm{Conv}(a,\alpha) \to \varphi_{12}$
\item $\varphi_{173} :\equiv (1 \in \mathbb{R} \to \lvert \alpha \rvert \in \mathbb{R} \to 0 \leq 1 \to 0 \leq \lvert \alpha \rvert \to 0 \leq 1 + \lvert \alpha \rvert) \to 1 \in \mathbb{R} \to \varphi_{114} \to (\mathrm{Conv}(a,\alpha) \to \varphi_{94}) \to \varphi_{105} \to \varphi_{111}$
\item $\varphi_{174} :\equiv \varphi_{156} \to (a_{n} - \alpha + \alpha = a_{n} \to \varphi_{80}) \to (a_{n} - \alpha + \alpha = a_{n} \to \varphi_{68}) \to \varphi_{152}$
\item $\varphi_{175} :\equiv \forall n\,(\varphi_{100} \to n \in \mathbb{N} \to \lvert a_{n} \rvert \leq M + (1 + \lvert \alpha \rvert)) \to \varphi_{163} \to (\varphi_{100} \to \forall n\,(\varphi_{100})) \to \varphi_{139}$
\item $\varphi_{176} :\equiv (\lvert a_{n} - \alpha \rvert \in \mathbb{R} \to 1 \in \mathbb{R} \to \lvert \alpha \rvert \in \mathbb{R} \to \lvert \alpha \rvert \in \mathbb{R} \to \varphi_{47}) \to \varphi_{138} \to 1 \in \mathbb{R} \to \varphi_{127} \to \varphi_{127} \to \varphi_{135} \to \varphi_{126} \to \varphi_{155}$
\end{abbrevlist}
\end{fullproof}

% 自動生成 (prover/build.js): ErrProd  ノード数 516
\begin{fullproof}{[C15]}{$\forall x\,\forall y\,\forall \alpha\,\forall \beta\,\forall M\,\forall h\,(x \in \mathbb{R} \to y \in \mathbb{R} \to \alpha \in \mathbb{R} \to \beta \in \mathbb{R} \to M \in \mathbb{R} \to h \in \mathbb{R} \to 0 \leq M \to \lvert x \rvert \leq M \to \lvert x - \alpha \rvert < h \to \lvert y - \beta \rvert < h \to \lvert x \cdot y - \alpha \cdot \beta \rvert \leq (M + \lvert \beta \rvert) \cdot h)$}{516}\phantomsection\label{tree:ErrProd}
\noindent$\Pi_{1}$:
\begin{ptree}
\AXC{$z = w \to \varphi_{117}$\lbl{A6}}
\UIC{$\forall w\,(z = w \to \varphi_{117})$}
\UIC{$\forall z\,\forall w\,(z = w \to \varphi_{117})$}
\AXC{$\forall z\,\forall w\,(z = w \to \varphi_{117}) \to \varphi_{180}$\lbl{A2}}
\BIC{$\varphi_{180}$}
\AXC{$\varphi_{180} \to M \cdot h + \lvert \beta \rvert \cdot h = (M + \lvert \beta \rvert) \cdot h \to \varphi_{142}$\lbl{A2}}
\BIC{$M \cdot h + \lvert \beta \rvert \cdot h = (M + \lvert \beta \rvert) \cdot h \to \varphi_{142}$}
\end{ptree}
\noindent$\Pi_{2}$:
\begin{ptree}
\AXC{$z = w \to \varphi_{123}$\lbl{A6}}
\UIC{$\forall w\,(z = w \to \varphi_{123})$}
\UIC{$\forall z\,\forall w\,(z = w \to \varphi_{123})$}
\AXC{$\forall z\,\forall w\,(z = w \to \varphi_{123}) \to \varphi_{195}$\lbl{A2}}
\BIC{$\varphi_{195}$}
\AXC{$\varphi_{195} \to M \cdot h + \lvert \beta \rvert \cdot h = (M + \lvert \beta \rvert) \cdot h \to \varphi_{155}$\lbl{A2}}
\BIC{$M \cdot h + \lvert \beta \rvert \cdot h = (M + \lvert \beta \rvert) \cdot h \to \varphi_{155}$}
\end{ptree}
\noindent$\Pi_{3}$:
\begin{ptree}
\AXC{$\forall x\,\forall y\,(x \in \mathbb{R} \to y \in \mathbb{R} \to x - y \in \mathbb{R})$\lbl{[F1]}}
\AXC{$\varphi_{92}$\lbl{A2}}
\BIC{$\forall y\,(x \in \mathbb{R} \to y \in \mathbb{R} \to x - y \in \mathbb{R})$}
\AXC{$\varphi_{81}$\lbl{A2}}
\BIC{$x \in \mathbb{R} \to \alpha \in \mathbb{R} \to x - \alpha \in \mathbb{R}$}
\AXC{$(x \in \mathbb{R} \to \alpha \in \mathbb{R} \to x - \alpha \in \mathbb{R}) \to \varphi_{229}$\lbl{A1}}
\BIC{$\varphi_{229}$}
\end{ptree}
\noindent$\Pi_{4}$:
\begin{ptree}
\AXC{$x \in \mathbb{R} \to y \in \mathbb{R} \to \alpha \in \mathbb{R} \to \beta \in \mathbb{R} \to M \in \mathbb{R} \to \varphi_{62}$\lbl{A1}}
\AXC{$x \in \mathbb{R} \to y \in \mathbb{R} \to \alpha \in \mathbb{R} \to \beta \in \mathbb{R} \to M \in \mathbb{R} \to \varphi_{52}$\lbl{A1}}
\AXC{$\stackrel{\Pi_{3}}{\vdots}$}\noLine
\UIC{$\varphi_{229}$}
\BIC{$(x \in \mathbb{R} \to y \in \mathbb{R} \to \alpha \in \mathbb{R} \to \beta \in \mathbb{R} \to M \in \mathbb{R} \to \varphi_{62}) \to \varphi_{171}$}
\BIC{$\varphi_{171}$}
\end{ptree}
\noindent$\Pi_{5}$:
\begin{ptree}
\AXC{$\forall x \in \mathbb{R}\,(0 \leq \lvert x \rvert)$\lbl{[A6]}}
\AXC{$\forall x \in \mathbb{R}\,(0 \leq \lvert x \rvert) \to x - \alpha \in \mathbb{R} \to 0 \leq \lvert x - \alpha \rvert$\lbl{A2}}
\BIC{$x - \alpha \in \mathbb{R} \to 0 \leq \lvert x - \alpha \rvert$}
\AXC{$(x - \alpha \in \mathbb{R} \to 0 \leq \lvert x - \alpha \rvert) \to \varphi_{171} \to \varphi_{159}$\lbl{A1}}
\BIC{$\varphi_{171} \to \varphi_{159}$}
\end{ptree}
\noindent$\Pi_{6}$:
\begin{ptree}
\AXC{$\forall x \in \mathbb{R}\,(0 \leq \lvert x \rvert)$\lbl{[A6]}}
\AXC{$\forall x \in \mathbb{R}\,(0 \leq \lvert x \rvert) \to \beta \in \mathbb{R} \to 0 \leq \lvert \beta \rvert$\lbl{A2}}
\BIC{$\beta \in \mathbb{R} \to 0 \leq \lvert \beta \rvert$}
\AXC{$\varphi_{216}$\lbl{A1}}
\BIC{$(x \in \mathbb{R} \to y \in \mathbb{R} \to \alpha \in \mathbb{R} \to \beta \in \mathbb{R} \to M \in \mathbb{R} \to \varphi_{63}) \to \varphi_{150}$}
\end{ptree}
\noindent$\Pi_{7}$:
\begin{ptree}
\AXC{$x \in \mathbb{R} \to y \in \mathbb{R} \to \alpha \in \mathbb{R} \to \beta \in \mathbb{R} \to M \in \mathbb{R} \to \varphi_{63}$\lbl{A1}}
\AXC{$\stackrel{\Pi_{6}}{\vdots}$}\noLine
\UIC{$(x \in \mathbb{R} \to y \in \mathbb{R} \to \alpha \in \mathbb{R} \to \beta \in \mathbb{R} \to M \in \mathbb{R} \to \varphi_{63}) \to \varphi_{150}$}
\BIC{$\varphi_{150}$}
\end{ptree}
\noindent$\Pi_{8}$:
\begin{ptree}
\AXC{$\forall x \in \mathbb{R}\,(\lvert x \rvert \in \mathbb{R})$\lbl{[A5]}}
\AXC{$\forall x \in \mathbb{R}\,(\lvert x \rvert \in \mathbb{R}) \to x - \alpha \in \mathbb{R} \to \lvert x - \alpha \rvert \in \mathbb{R}$\lbl{A2}}
\BIC{$x - \alpha \in \mathbb{R} \to \lvert x - \alpha \rvert \in \mathbb{R}$}
\end{ptree}
\noindent$\Pi_{9}$:
\begin{ptree}
\AXC{$\stackrel{\Pi_{4}}{\vdots}$}\noLine
\UIC{$\varphi_{171}$}
\AXC{$\stackrel{\Pi_{8}}{\vdots}$}\noLine
\UIC{$x - \alpha \in \mathbb{R} \to \lvert x - \alpha \rvert \in \mathbb{R}$}
\AXC{$(x - \alpha \in \mathbb{R} \to \lvert x - \alpha \rvert \in \mathbb{R}) \to \varphi_{171} \to \varphi_{187}$\lbl{A1}}
\BIC{$\varphi_{171} \to \varphi_{187}$}
\BIC{$\varphi_{187}$}
\end{ptree}
\noindent$\Pi_{10}$:
\begin{ptree}
\AXC{$\forall x \in \mathbb{R}\,(\lvert x \rvert \in \mathbb{R})$\lbl{[A5]}}
\AXC{$\forall x \in \mathbb{R}\,(\lvert x \rvert \in \mathbb{R}) \to \beta \in \mathbb{R} \to \lvert \beta \rvert \in \mathbb{R}$\lbl{A2}}
\BIC{$\beta \in \mathbb{R} \to \lvert \beta \rvert \in \mathbb{R}$}
\AXC{$\varphi_{220}$\lbl{A1}}
\BIC{$(x \in \mathbb{R} \to y \in \mathbb{R} \to \alpha \in \mathbb{R} \to \beta \in \mathbb{R} \to M \in \mathbb{R} \to \varphi_{63}) \to \varphi_{172}$}
\end{ptree}
\noindent$\Pi_{11}$:
\begin{ptree}
\AXC{$x \in \mathbb{R} \to y \in \mathbb{R} \to \alpha \in \mathbb{R} \to \beta \in \mathbb{R} \to M \in \mathbb{R} \to \varphi_{63}$\lbl{A1}}
\AXC{$\stackrel{\Pi_{10}}{\vdots}$}\noLine
\UIC{$(x \in \mathbb{R} \to y \in \mathbb{R} \to \alpha \in \mathbb{R} \to \beta \in \mathbb{R} \to M \in \mathbb{R} \to \varphi_{63}) \to \varphi_{172}$}
\BIC{$\varphi_{172}$}
\end{ptree}
\noindent$\Pi_{12}$:
\begin{ptree}
\AXC{$\varphi_{105}$\lbl{[E12]}}
\AXC{$\varphi_{105} \to \varphi_{133}$\lbl{A2}}
\BIC{$\varphi_{133}$}
\AXC{$\varphi_{133} \to \varphi_{138}$\lbl{A2}}
\BIC{$\varphi_{138}$}
\AXC{$\varphi_{138} \to \varphi_{206}$\lbl{A2}}
\BIC{$\varphi_{206}$}
\AXC{$\varphi_{206} \to \lvert \beta \rvert \in \mathbb{R} \to \varphi_{182}$\lbl{A2}}
\BIC{$\lvert \beta \rvert \in \mathbb{R} \to \varphi_{182}$}
\end{ptree}
\noindent$\Pi_{13}$:
\begin{ptree}
\AXC{$\stackrel{\Pi_{9}}{\vdots}$}\noLine
\UIC{$\varphi_{187}$}
\AXC{$\stackrel{\Pi_{11}}{\vdots}$}\noLine
\UIC{$\varphi_{172}$}
\AXC{$\stackrel{\Pi_{11}}{\vdots}$}\noLine
\UIC{$\varphi_{172}$}
\AXC{$\stackrel{\Pi_{12}}{\vdots}$}\noLine
\UIC{$\lvert \beta \rvert \in \mathbb{R} \to \varphi_{182}$}
\AXC{$(\lvert \beta \rvert \in \mathbb{R} \to \varphi_{182}) \to \varphi_{172} \to \varphi_{172} \to \varphi_{187} \to \varphi_{254}$\lbl{A1}}
\BIC{$\varphi_{172} \to \varphi_{172} \to \varphi_{187} \to \varphi_{254}$}
\BIC{$\varphi_{172} \to \varphi_{187} \to \varphi_{254}$}
\BIC{$\varphi_{187} \to \varphi_{254}$}
\BIC{$\varphi_{254}$}
\end{ptree}
\noindent$\Pi_{14}$:
\begin{ptree}
\AXC{$\stackrel{\Pi_{4}}{\vdots}$}\noLine
\UIC{$\varphi_{171}$}
\AXC{$\stackrel{\Pi_{5}}{\vdots}$}\noLine
\UIC{$\varphi_{171} \to \varphi_{159}$}
\BIC{$\varphi_{159}$}
\AXC{$\stackrel{\Pi_{7}}{\vdots}$}\noLine
\UIC{$\varphi_{150}$}
\AXC{$x \in \mathbb{R} \to y \in \mathbb{R} \to \alpha \in \mathbb{R} \to \beta \in \mathbb{R} \to M \in \mathbb{R} \to \varphi_{53}$\lbl{A1}}
\AXC{$\stackrel{\Pi_{13}}{\vdots}$}\noLine
\UIC{$\varphi_{254}$}
\BIC{$\varphi_{150} \to \varphi_{159} \to \varphi_{167} \to \varphi_{158} \to x \in \mathbb{R} \to y \in \mathbb{R} \to \varphi_{153}$}
\BIC{$\varphi_{159} \to \varphi_{167} \to \varphi_{158} \to x \in \mathbb{R} \to y \in \mathbb{R} \to \varphi_{153}$}
\BIC{$\varphi_{167} \to \varphi_{158} \to x \in \mathbb{R} \to y \in \mathbb{R} \to \varphi_{153}$}
\end{ptree}
\noindent$\Pi_{15}$:
\begin{ptree}
\AXC{$x = x$\lbl{A5}}
\UIC{$\forall x\,(x = x)$}
\AXC{$\forall x\,(x = x) \to \lvert \beta \rvert = \lvert \beta \rvert$\lbl{A2}}
\BIC{$\lvert \beta \rvert = \lvert \beta \rvert$}
\AXC{$\lvert \beta \rvert = \lvert \beta \rvert \to \varphi_{167}$\lbl{A1}}
\BIC{$\varphi_{167}$}
\AXC{$\stackrel{\Pi_{14}}{\vdots}$}\noLine
\UIC{$\varphi_{167} \to \varphi_{158} \to x \in \mathbb{R} \to y \in \mathbb{R} \to \varphi_{153}$}
\BIC{$\varphi_{158} \to x \in \mathbb{R} \to y \in \mathbb{R} \to \varphi_{153}$}
\end{ptree}
\noindent$\Pi_{16}$:
\begin{ptree}
\AXC{$\forall x\,\forall y\,(x \in \mathbb{R} \to y \in \mathbb{R} \to x - y \in \mathbb{R})$\lbl{[F1]}}
\AXC{$\varphi_{102}$\lbl{A2}}
\BIC{$\forall y\,(v_{689} \in \mathbb{R} \to y \in \mathbb{R} \to v_{689} - y \in \mathbb{R})$}
\AXC{$\varphi_{119}$\lbl{A2}}
\BIC{$v_{689} \in \mathbb{R} \to v_{690} \in \mathbb{R} \to v_{689} - v_{690} \in \mathbb{R}$}
\UIC{$\forall v_{690}\,(v_{689} \in \mathbb{R} \to v_{690} \in \mathbb{R} \to v_{689} - v_{690} \in \mathbb{R})$}
\UIC{$\varphi_{8}$}
\AXC{$\varphi_{8} \to \forall v_{690}\,(y \in \mathbb{R} \to v_{690} \in \mathbb{R} \to y - v_{690} \in \mathbb{R})$\lbl{A2}}
\BIC{$\forall v_{690}\,(y \in \mathbb{R} \to v_{690} \in \mathbb{R} \to y - v_{690} \in \mathbb{R})$}
\end{ptree}
\noindent$\Pi_{17}$:
\begin{ptree}
\AXC{$\stackrel{\Pi_{16}}{\vdots}$}\noLine
\UIC{$\forall v_{690}\,(y \in \mathbb{R} \to v_{690} \in \mathbb{R} \to y - v_{690} \in \mathbb{R})$}
\AXC{$\varphi_{100}$\lbl{A2}}
\BIC{$y \in \mathbb{R} \to \beta \in \mathbb{R} \to y - \beta \in \mathbb{R}$}
\AXC{$(y \in \mathbb{R} \to \beta \in \mathbb{R} \to y - \beta \in \mathbb{R}) \to \varphi_{230}$\lbl{A1}}
\BIC{$\varphi_{230}$}
\end{ptree}
\noindent$\Pi_{18}$:
\begin{ptree}
\AXC{$x \in \mathbb{R} \to y \in \mathbb{R} \to \alpha \in \mathbb{R} \to \beta \in \mathbb{R} \to M \in \mathbb{R} \to \varphi_{63}$\lbl{A1}}
\AXC{$x \in \mathbb{R} \to y \in \mathbb{R} \to \alpha \in \mathbb{R} \to \beta \in \mathbb{R} \to M \in \mathbb{R} \to \varphi_{54}$\lbl{A1}}
\AXC{$\stackrel{\Pi_{17}}{\vdots}$}\noLine
\UIC{$\varphi_{230}$}
\BIC{$(x \in \mathbb{R} \to y \in \mathbb{R} \to \alpha \in \mathbb{R} \to \beta \in \mathbb{R} \to M \in \mathbb{R} \to \varphi_{63}) \to \varphi_{173}$}
\BIC{$\varphi_{173}$}
\end{ptree}
\noindent$\Pi_{19}$:
\begin{ptree}
\AXC{$\forall x \in \mathbb{R}\,(0 \leq \lvert x \rvert)$\lbl{[A6]}}
\AXC{$\forall x \in \mathbb{R}\,(0 \leq \lvert x \rvert) \to y - \beta \in \mathbb{R} \to 0 \leq \lvert y - \beta \rvert$\lbl{A2}}
\BIC{$y - \beta \in \mathbb{R} \to 0 \leq \lvert y - \beta \rvert$}
\AXC{$(y - \beta \in \mathbb{R} \to 0 \leq \lvert y - \beta \rvert) \to \varphi_{173} \to \varphi_{161}$\lbl{A1}}
\BIC{$\varphi_{173} \to \varphi_{161}$}
\end{ptree}
\noindent$\Pi_{20}$:
\begin{ptree}
\AXC{$\forall x \in \mathbb{R}\,(0 \leq \lvert x \rvert)$\lbl{[A6]}}
\AXC{$\forall x \in \mathbb{R}\,(0 \leq \lvert x \rvert) \to x \in \mathbb{R} \to 0 \leq \lvert x \rvert$\lbl{A2}}
\BIC{$x \in \mathbb{R} \to 0 \leq \lvert x \rvert$}
\AXC{$\varphi_{215}$\lbl{A1}}
\BIC{$(x \in \mathbb{R} \to y \in \mathbb{R} \to \alpha \in \mathbb{R} \to \beta \in \mathbb{R} \to M \in \mathbb{R} \to \varphi_{52}) \to \varphi_{147}$}
\end{ptree}
\noindent$\Pi_{21}$:
\begin{ptree}
\AXC{$x \in \mathbb{R} \to y \in \mathbb{R} \to \alpha \in \mathbb{R} \to \beta \in \mathbb{R} \to M \in \mathbb{R} \to \varphi_{52}$\lbl{A1}}
\AXC{$\stackrel{\Pi_{20}}{\vdots}$}\noLine
\UIC{$(x \in \mathbb{R} \to y \in \mathbb{R} \to \alpha \in \mathbb{R} \to \beta \in \mathbb{R} \to M \in \mathbb{R} \to \varphi_{52}) \to \varphi_{147}$}
\BIC{$\varphi_{147}$}
\end{ptree}
\noindent$\Pi_{22}$:
\begin{ptree}
\AXC{$\forall x \in \mathbb{R}\,(\lvert x \rvert \in \mathbb{R})$\lbl{[A5]}}
\AXC{$\forall x \in \mathbb{R}\,(\lvert x \rvert \in \mathbb{R}) \to y - \beta \in \mathbb{R} \to \lvert y - \beta \rvert \in \mathbb{R}$\lbl{A2}}
\BIC{$y - \beta \in \mathbb{R} \to \lvert y - \beta \rvert \in \mathbb{R}$}
\end{ptree}
\noindent$\Pi_{23}$:
\begin{ptree}
\AXC{$\stackrel{\Pi_{18}}{\vdots}$}\noLine
\UIC{$\varphi_{173}$}
\AXC{$\stackrel{\Pi_{22}}{\vdots}$}\noLine
\UIC{$y - \beta \in \mathbb{R} \to \lvert y - \beta \rvert \in \mathbb{R}$}
\AXC{$(y - \beta \in \mathbb{R} \to \lvert y - \beta \rvert \in \mathbb{R}) \to \varphi_{173} \to \varphi_{188}$\lbl{A1}}
\BIC{$\varphi_{173} \to \varphi_{188}$}
\BIC{$\varphi_{188}$}
\end{ptree}
\noindent$\Pi_{24}$:
\begin{ptree}
\AXC{$\forall x \in \mathbb{R}\,(\lvert x \rvert \in \mathbb{R})$\lbl{[A5]}}
\AXC{$\forall x \in \mathbb{R}\,(\lvert x \rvert \in \mathbb{R}) \to x \in \mathbb{R} \to \lvert x \rvert \in \mathbb{R}$\lbl{A2}}
\BIC{$x \in \mathbb{R} \to \lvert x \rvert \in \mathbb{R}$}
\AXC{$\varphi_{218}$\lbl{A1}}
\BIC{$(x \in \mathbb{R} \to y \in \mathbb{R} \to \alpha \in \mathbb{R} \to \beta \in \mathbb{R} \to M \in \mathbb{R} \to \varphi_{52}) \to \varphi_{168}$}
\end{ptree}
\noindent$\Pi_{25}$:
\begin{ptree}
\AXC{$x \in \mathbb{R} \to y \in \mathbb{R} \to \alpha \in \mathbb{R} \to \beta \in \mathbb{R} \to M \in \mathbb{R} \to \varphi_{52}$\lbl{A1}}
\AXC{$\stackrel{\Pi_{24}}{\vdots}$}\noLine
\UIC{$(x \in \mathbb{R} \to y \in \mathbb{R} \to \alpha \in \mathbb{R} \to \beta \in \mathbb{R} \to M \in \mathbb{R} \to \varphi_{52}) \to \varphi_{168}$}
\BIC{$\varphi_{168}$}
\end{ptree}
\noindent$\Pi_{26}$:
\begin{ptree}
\AXC{$\stackrel{\Pi_{25}}{\vdots}$}\noLine
\UIC{$\varphi_{168}$}
\AXC{$\varphi_{105}$\lbl{[E12]}}
\AXC{$\varphi_{105} \to \varphi_{129}$\lbl{A2}}
\BIC{$\varphi_{129}$}
\AXC{$\varphi_{129} \to \varphi_{128}$\lbl{A2}}
\BIC{$\varphi_{128}$}
\AXC{$\varphi_{128} \to \varphi_{181}$\lbl{A2}}
\BIC{$\varphi_{181}$}
\AXC{$\varphi_{181} \to \varphi_{174}$\lbl{A2}}
\BIC{$\varphi_{174}$}
\AXC{$\varphi_{174} \to \varphi_{168} \to \varphi_{255}$\lbl{A1}}
\BIC{$\varphi_{168} \to \varphi_{255}$}
\BIC{$\varphi_{255}$}
\end{ptree}
\noindent$\Pi_{27}$:
\begin{ptree}
\AXC{$x \in \mathbb{R} \to y \in \mathbb{R} \to \alpha \in \mathbb{R} \to \beta \in \mathbb{R} \to M \in \mathbb{R} \to \varphi_{53}$\lbl{A1}}
\AXC{$\stackrel{\Pi_{23}}{\vdots}$}\noLine
\UIC{$\varphi_{188}$}
\AXC{$x \in \mathbb{R} \to y \in \mathbb{R} \to \alpha \in \mathbb{R} \to \beta \in \mathbb{R} \to M \in \mathbb{R} \to \varphi_{55}$\lbl{A1}}
\AXC{$\stackrel{\Pi_{26}}{\vdots}$}\noLine
\UIC{$\varphi_{255}$}
\BIC{$\varphi_{188} \to \varphi_{253}$}
\BIC{$\varphi_{253}$}
\BIC{$\varphi_{147} \to \varphi_{161} \to \varphi_{146} \to \varphi_{160} \to \varphi_{199}$}
\end{ptree}
\noindent$\Pi_{28}$:
\begin{ptree}
\AXC{$\varphi_{178}$\lbl{A1}}
\AXC{$\varphi_{178} \to \varphi_{160}$\lbl{A1}}
\BIC{$\varphi_{160}$}
\AXC{$\varphi_{146}$\lbl{A1}}
\AXC{$\stackrel{\Pi_{18}}{\vdots}$}\noLine
\UIC{$\varphi_{173}$}
\AXC{$\stackrel{\Pi_{19}}{\vdots}$}\noLine
\UIC{$\varphi_{173} \to \varphi_{161}$}
\BIC{$\varphi_{161}$}
\AXC{$\stackrel{\Pi_{21}}{\vdots}$}\noLine
\UIC{$\varphi_{147}$}
\AXC{$\stackrel{\Pi_{27}}{\vdots}$}\noLine
\UIC{$\varphi_{147} \to \varphi_{161} \to \varphi_{146} \to \varphi_{160} \to \varphi_{199}$}
\BIC{$\varphi_{161} \to \varphi_{146} \to \varphi_{160} \to \varphi_{199}$}
\BIC{$\varphi_{146} \to \varphi_{160} \to \varphi_{199}$}
\BIC{$\varphi_{160} \to \varphi_{199}$}
\BIC{$\varphi_{199}$}
\end{ptree}
\noindent$\Pi_{29}$:
\begin{ptree}
\AXC{$\forall x\,\forall y\,(x \in \mathbb{R} \to y \in \mathbb{R} \to x \cdot y \in \mathbb{R})$\lbl{(R1)}}
\AXC{$\varphi_{120}$\lbl{A2}}
\BIC{$\forall y\,(\lvert \beta \rvert \in \mathbb{R} \to y \in \mathbb{R} \to \lvert \beta \rvert \cdot y \in \mathbb{R})$}
\AXC{$\varphi_{127}$\lbl{A2}}
\BIC{$\lvert \beta \rvert \in \mathbb{R} \to h \in \mathbb{R} \to \lvert \beta \rvert \cdot h \in \mathbb{R}$}
\AXC{$(\lvert \beta \rvert \in \mathbb{R} \to h \in \mathbb{R} \to \lvert \beta \rvert \cdot h \in \mathbb{R}) \to \varphi_{231}$\lbl{A1}}
\BIC{$\varphi_{231}$}
\end{ptree}
\noindent$\Pi_{30}$:
\begin{ptree}
\AXC{$x \in \mathbb{R} \to y \in \mathbb{R} \to \alpha \in \mathbb{R} \to \beta \in \mathbb{R} \to M \in \mathbb{R} \to \varphi_{53}$\lbl{A1}}
\AXC{$\stackrel{\Pi_{11}}{\vdots}$}\noLine
\UIC{$\varphi_{172}$}
\AXC{$\stackrel{\Pi_{29}}{\vdots}$}\noLine
\UIC{$\varphi_{231}$}
\BIC{$(x \in \mathbb{R} \to y \in \mathbb{R} \to \alpha \in \mathbb{R} \to \beta \in \mathbb{R} \to M \in \mathbb{R} \to \varphi_{53}) \to \varphi_{189}$}
\BIC{$\varphi_{189}$}
\end{ptree}
\noindent$\Pi_{31}$:
\begin{ptree}
\AXC{$\stackrel{\Pi_{9}}{\vdots}$}\noLine
\UIC{$\varphi_{187}$}
\AXC{$\stackrel{\Pi_{11}}{\vdots}$}\noLine
\UIC{$\varphi_{172}$}
\AXC{$\forall x\,\forall y\,(x \in \mathbb{R} \to y \in \mathbb{R} \to x \cdot y \in \mathbb{R})$\lbl{(R1)}}
\AXC{$\varphi_{120}$\lbl{A2}}
\BIC{$\forall y\,(\lvert \beta \rvert \in \mathbb{R} \to y \in \mathbb{R} \to \lvert \beta \rvert \cdot y \in \mathbb{R})$}
\AXC{$\varphi_{136}$\lbl{A2}}
\BIC{$\lvert \beta \rvert \in \mathbb{R} \to \lvert x - \alpha \rvert \in \mathbb{R} \to \lvert \beta \rvert \cdot \lvert x - \alpha \rvert \in \mathbb{R}$}
\AXC{$\varphi_{243}$\lbl{A1}}
\BIC{$\varphi_{172} \to \varphi_{187} \to \varphi_{203}$}
\BIC{$\varphi_{187} \to \varphi_{203}$}
\BIC{$\varphi_{203}$}
\end{ptree}
\noindent$\Pi_{32}$:
\begin{ptree}
\AXC{$\forall x\,\forall y\,(x \in \mathbb{R} \to y \in \mathbb{R} \to x \cdot y \in \mathbb{R})$\lbl{(R1)}}
\AXC{$\varphi_{93}$\lbl{A2}}
\BIC{$\forall y\,(M \in \mathbb{R} \to y \in \mathbb{R} \to M \cdot y \in \mathbb{R})$}
\AXC{$\varphi_{78}$\lbl{A2}}
\BIC{$M \in \mathbb{R} \to h \in \mathbb{R} \to M \cdot h \in \mathbb{R}$}
\AXC{$(M \in \mathbb{R} \to h \in \mathbb{R} \to M \cdot h \in \mathbb{R}) \to \varphi_{228}$\lbl{A1}}
\BIC{$\varphi_{228}$}
\end{ptree}
\noindent$\Pi_{33}$:
\begin{ptree}
\AXC{$x \in \mathbb{R} \to y \in \mathbb{R} \to \alpha \in \mathbb{R} \to \beta \in \mathbb{R} \to M \in \mathbb{R} \to \varphi_{53}$\lbl{A1}}
\AXC{$x \in \mathbb{R} \to y \in \mathbb{R} \to \alpha \in \mathbb{R} \to \beta \in \mathbb{R} \to M \in \mathbb{R} \to \varphi_{55}$\lbl{A1}}
\AXC{$\stackrel{\Pi_{32}}{\vdots}$}\noLine
\UIC{$\varphi_{228}$}
\BIC{$(x \in \mathbb{R} \to y \in \mathbb{R} \to \alpha \in \mathbb{R} \to \beta \in \mathbb{R} \to M \in \mathbb{R} \to \varphi_{53}) \to \varphi_{169}$}
\BIC{$\varphi_{169}$}
\end{ptree}
\noindent$\Pi_{34}$:
\begin{ptree}
\AXC{$\stackrel{\Pi_{23}}{\vdots}$}\noLine
\UIC{$\varphi_{188}$}
\AXC{$\stackrel{\Pi_{25}}{\vdots}$}\noLine
\UIC{$\varphi_{168}$}
\AXC{$\forall x\,\forall y\,(x \in \mathbb{R} \to y \in \mathbb{R} \to x \cdot y \in \mathbb{R})$\lbl{(R1)}}
\AXC{$\varphi_{115}$\lbl{A2}}
\BIC{$\forall y\,(\lvert x \rvert \in \mathbb{R} \to y \in \mathbb{R} \to \lvert x \rvert \cdot y \in \mathbb{R})$}
\AXC{$\varphi_{135}$\lbl{A2}}
\BIC{$\lvert x \rvert \in \mathbb{R} \to \lvert y - \beta \rvert \in \mathbb{R} \to \lvert x \rvert \cdot \lvert y - \beta \rvert \in \mathbb{R}$}
\AXC{$\varphi_{242}$\lbl{A1}}
\BIC{$\varphi_{168} \to \varphi_{188} \to \varphi_{202}$}
\BIC{$\varphi_{188} \to \varphi_{202}$}
\BIC{$\varphi_{202}$}
\end{ptree}
\noindent$\Pi_{35}$:
\begin{ptree}
\AXC{$\varphi_{80}$\lbl{[F21]}}
\AXC{$\varphi_{80} \to \varphi_{96}$\lbl{A2}}
\BIC{$\varphi_{96}$}
\AXC{$\varphi_{96} \to \varphi_{106}$\lbl{A2}}
\BIC{$\varphi_{106}$}
\AXC{$\varphi_{106} \to \varphi_{121}$\lbl{A2}}
\BIC{$\varphi_{121}$}
\AXC{$\varphi_{211}$\lbl{A2}}
\BIC{$v_{691} \in \mathbb{R} \to v_{692} \in \mathbb{R} \to v_{693} \in \mathbb{R} \to v_{694} \in \mathbb{R} \to \varphi_{3}$}
\UIC{$\varphi_{130}$}
\UIC{$\forall v_{693}\,(\varphi_{130})$}
\UIC{$\forall v_{692}\,\forall v_{693}\,(\varphi_{130})$}
\UIC{$\forall v_{691}\,\forall v_{692}\,\forall v_{693}\,(\varphi_{130})$}
\AXC{$\forall v_{691}\,\forall v_{692}\,\forall v_{693}\,(\varphi_{130}) \to \varphi_{207}$\lbl{A2}}
\BIC{$\varphi_{207}$}
\end{ptree}
\noindent$\Pi_{36}$:
\begin{ptree}
\AXC{$\stackrel{\Pi_{34}}{\vdots}$}\noLine
\UIC{$\varphi_{202}$}
\AXC{$\stackrel{\Pi_{35}}{\vdots}$}\noLine
\UIC{$\varphi_{207}$}
\AXC{$\varphi_{207} \to \varphi_{208}$\lbl{A2}}
\BIC{$\varphi_{208}$}
\AXC{$\varphi_{208} \to \varphi_{212}$\lbl{A2}}
\BIC{$\varphi_{212}$}
\AXC{$\varphi_{212} \to \lvert x \rvert \cdot \lvert y - \beta \rvert \in \mathbb{R} \to M \cdot h \in \mathbb{R} \to \varphi_{210}$\lbl{A2}}
\BIC{$\lvert x \rvert \cdot \lvert y - \beta \rvert \in \mathbb{R} \to M \cdot h \in \mathbb{R} \to \varphi_{210}$}
\AXC{$\varphi_{256}$\lbl{A1}}
\BIC{$\varphi_{202} \to \varphi_{169} \to \varphi_{203} \to \varphi_{189} \to \varphi_{241}$}
\BIC{$\varphi_{169} \to \varphi_{203} \to \varphi_{189} \to \varphi_{241}$}
\end{ptree}
\noindent$\Pi_{37}$:
\begin{ptree}
\AXC{$\stackrel{\Pi_{28}}{\vdots}$}\noLine
\UIC{$\varphi_{199}$}
\AXC{$\stackrel{\Pi_{30}}{\vdots}$}\noLine
\UIC{$\varphi_{189}$}
\AXC{$\stackrel{\Pi_{31}}{\vdots}$}\noLine
\UIC{$\varphi_{203}$}
\AXC{$\stackrel{\Pi_{33}}{\vdots}$}\noLine
\UIC{$\varphi_{169}$}
\AXC{$\stackrel{\Pi_{36}}{\vdots}$}\noLine
\UIC{$\varphi_{169} \to \varphi_{203} \to \varphi_{189} \to \varphi_{241}$}
\BIC{$\varphi_{203} \to \varphi_{189} \to \varphi_{241}$}
\BIC{$\varphi_{189} \to \varphi_{241}$}
\BIC{$\varphi_{241}$}
\BIC{$(x \in \mathbb{R} \to y \in \mathbb{R} \to \varphi_{153}) \to x \in \mathbb{R} \to y \in \mathbb{R} \to \alpha \in \mathbb{R} \to \varphi_{204}$}
\end{ptree}
\noindent$\Pi_{38}$:
\begin{ptree}
\AXC{$\varphi_{177}$\lbl{A1}}
\AXC{$\varphi_{177} \to \varphi_{158}$\lbl{A1}}
\BIC{$\varphi_{158}$}
\AXC{$\stackrel{\Pi_{15}}{\vdots}$}\noLine
\UIC{$\varphi_{158} \to x \in \mathbb{R} \to y \in \mathbb{R} \to \varphi_{153}$}
\BIC{$x \in \mathbb{R} \to y \in \mathbb{R} \to \varphi_{153}$}
\AXC{$\stackrel{\Pi_{37}}{\vdots}$}\noLine
\UIC{$(x \in \mathbb{R} \to y \in \mathbb{R} \to \varphi_{153}) \to x \in \mathbb{R} \to y \in \mathbb{R} \to \alpha \in \mathbb{R} \to \varphi_{204}$}
\BIC{$x \in \mathbb{R} \to y \in \mathbb{R} \to \alpha \in \mathbb{R} \to \varphi_{204}$}
\end{ptree}
\noindent$\Pi_{39}$:
\begin{ptree}
\AXC{$y = z \to x = y \to x = z$\lbl{A6}}
\AXC{$(y = z \to x = y \to x = z) \to x = y \to y = z \to x = z$\lbl{A1}}
\BIC{$x = y \to y = z \to x = z$}
\UIC{$\forall z\,(x = y \to y = z \to x = z)$}
\UIC{$\forall y\,\forall z\,(x = y \to y = z \to x = z)$}
\UIC{$\forall x\,\forall y\,\forall z\,(x = y \to y = z \to x = z)$}
\AXC{$\forall x\,\forall y\,\forall z\,(x = y \to y = z \to x = z) \to \varphi_{7}$\lbl{A2}}
\BIC{$\varphi_{7}$}
\end{ptree}
\noindent$\Pi_{40}$:
\begin{ptree}
\AXC{$\forall x\,\forall y\,(x \in \mathbb{R} \to y \in \mathbb{R} \to x + y \in \mathbb{R})$\lbl{(R1)}}
\AXC{$\varphi_{94}$\lbl{A2}}
\BIC{$\forall y\,(M \in \mathbb{R} \to y \in \mathbb{R} \to M + y \in \mathbb{R})$}
\AXC{$\varphi_{103}$\lbl{A2}}
\BIC{$M \in \mathbb{R} \to \lvert \beta \rvert \in \mathbb{R} \to M + \lvert \beta \rvert \in \mathbb{R}$}
\AXC{$(M \in \mathbb{R} \to \lvert \beta \rvert \in \mathbb{R} \to M + \lvert \beta \rvert \in \mathbb{R}) \to \varphi_{232}$\lbl{A1}}
\BIC{$\varphi_{232}$}
\end{ptree}
\noindent$\Pi_{41}$:
\begin{ptree}
\AXC{$\stackrel{\Pi_{11}}{\vdots}$}\noLine
\UIC{$\varphi_{172}$}
\AXC{$x \in \mathbb{R} \to y \in \mathbb{R} \to \alpha \in \mathbb{R} \to \beta \in \mathbb{R} \to M \in \mathbb{R} \to \varphi_{55}$\lbl{A1}}
\AXC{$\stackrel{\Pi_{40}}{\vdots}$}\noLine
\UIC{$\varphi_{232}$}
\BIC{$\varphi_{172} \to \varphi_{190}$}
\BIC{$\varphi_{190}$}
\end{ptree}
\noindent$\Pi_{42}$:
\begin{ptree}
\AXC{$\forall x\,\forall y\,(x \in \mathbb{R} \to y \in \mathbb{R} \to x \cdot y = y \cdot x)$\lbl{(R7)}}
\AXC{$\varphi_{82}$\lbl{A2}}
\BIC{$\forall y\,(h \in \mathbb{R} \to y \in \mathbb{R} \to h \cdot y = y \cdot h)$}
\AXC{$\forall y\,(h \in \mathbb{R} \to y \in \mathbb{R} \to h \cdot y = y \cdot h) \to \varphi_{23}$\lbl{A2}}
\BIC{$\varphi_{23}$}
\AXC{$\varphi_{23} \to \varphi_{237}$\lbl{A1}}
\BIC{$\varphi_{237}$}
\end{ptree}
\noindent$\Pi_{43}$:
\begin{ptree}
\AXC{$y = z \to x = y \to x = z$\lbl{A6}}
\AXC{$(y = z \to x = y \to x = z) \to x = y \to y = z \to x = z$\lbl{A1}}
\BIC{$x = y \to y = z \to x = z$}
\UIC{$\forall z\,(x = y \to y = z \to x = z)$}
\UIC{$\forall y\,\forall z\,(x = y \to y = z \to x = z)$}
\UIC{$\forall x\,\forall y\,\forall z\,(x = y \to y = z \to x = z)$}
\AXC{$\forall x\,\forall y\,\forall z\,(x = y \to y = z \to x = z) \to \varphi_{7}$\lbl{A2}}
\BIC{$\varphi_{7}$}
\end{ptree}
\noindent$\Pi_{44}$:
\begin{ptree}
\AXC{$x = x$\lbl{A5}}
\AXC{$x = y \to x = x \to y = x$\lbl{A6}}
\AXC{$(x = y \to x = x \to y = x) \to x = x \to x = y \to y = x$\lbl{A1}}
\BIC{$x = x \to x = y \to y = x$}
\BIC{$x = y \to y = x$}
\UIC{$\forall y\,(x = y \to y = x)$}
\UIC{$\forall x\,\forall y\,(x = y \to y = x)$}
\end{ptree}
\noindent$\Pi_{45}$:
\begin{ptree}
\AXC{$\stackrel{\Pi_{44}}{\vdots}$}\noLine
\UIC{$\forall x\,\forall y\,(x = y \to y = x)$}
\AXC{$\varphi_{6}$\lbl{A2}}
\BIC{$\forall y\,(h \cdot (M + \lvert \beta \rvert) = y \to y = h \cdot (M + \lvert \beta \rvert))$}
\AXC{$\forall y\,(h \cdot (M + \lvert \beta \rvert) = y \to y = h \cdot (M + \lvert \beta \rvert)) \to \varphi_{67}$\lbl{A2}}
\BIC{$\varphi_{67}$}
\end{ptree}
\noindent$\Pi_{46}$:
\begin{ptree}
\AXC{$\varphi_{42}$\lbl{(R10)}}
\AXC{$\varphi_{42} \to \varphi_{30}$\lbl{A2}}
\BIC{$\varphi_{30}$}
\AXC{$\varphi_{30} \to \varphi_{15}$\lbl{A2}}
\BIC{$\varphi_{15}$}
\AXC{$\varphi_{15} \to h \in \mathbb{R} \to \varphi_{28}$\lbl{A2}}
\BIC{$h \in \mathbb{R} \to \varphi_{28}$}
\AXC{$\varphi_{247}$\lbl{A1}}
\BIC{$(x \in \mathbb{R} \to y \in \mathbb{R} \to \alpha \in \mathbb{R} \to \beta \in \mathbb{R} \to M \in \mathbb{R} \to \varphi_{53}) \to \varphi_{238}$}
\end{ptree}
\noindent$\Pi_{47}$:
\begin{ptree}
\AXC{$x \in \mathbb{R} \to y \in \mathbb{R} \to \alpha \in \mathbb{R} \to \beta \in \mathbb{R} \to M \in \mathbb{R} \to \varphi_{53}$\lbl{A1}}
\AXC{$\stackrel{\Pi_{46}}{\vdots}$}\noLine
\UIC{$(x \in \mathbb{R} \to y \in \mathbb{R} \to \alpha \in \mathbb{R} \to \beta \in \mathbb{R} \to M \in \mathbb{R} \to \varphi_{53}) \to \varphi_{238}$}
\BIC{$\varphi_{238}$}
\end{ptree}
\noindent$\Pi_{48}$:
\begin{ptree}
\AXC{$\stackrel{\Pi_{11}}{\vdots}$}\noLine
\UIC{$\varphi_{172}$}
\AXC{$x \in \mathbb{R} \to y \in \mathbb{R} \to \alpha \in \mathbb{R} \to \beta \in \mathbb{R} \to M \in \mathbb{R} \to \varphi_{55}$\lbl{A1}}
\AXC{$\stackrel{\Pi_{47}}{\vdots}$}\noLine
\UIC{$\varphi_{238}$}
\BIC{$\varphi_{172} \to x \in \mathbb{R} \to y \in \mathbb{R} \to \varphi_{162}$}
\BIC{$x \in \mathbb{R} \to y \in \mathbb{R} \to \varphi_{162}$}
\end{ptree}
\noindent$\Pi_{49}$:
\begin{ptree}
\AXC{$\stackrel{\Pi_{45}}{\vdots}$}\noLine
\UIC{$\varphi_{67}$}
\AXC{$\stackrel{\Pi_{48}}{\vdots}$}\noLine
\UIC{$x \in \mathbb{R} \to y \in \mathbb{R} \to \varphi_{162}$}
\AXC{$(x \in \mathbb{R} \to y \in \mathbb{R} \to \varphi_{162}) \to \varphi_{67} \to x \in \mathbb{R} \to y \in \mathbb{R} \to \varphi_{163}$\lbl{A1}}
\BIC{$\varphi_{67} \to x \in \mathbb{R} \to y \in \mathbb{R} \to \varphi_{163}$}
\BIC{$x \in \mathbb{R} \to y \in \mathbb{R} \to \varphi_{163}$}
\end{ptree}
\noindent$\Pi_{50}$:
\begin{ptree}
\AXC{$y = z \to x = y \to x = z$\lbl{A6}}
\AXC{$(y = z \to x = y \to x = z) \to x = y \to y = z \to x = z$\lbl{A1}}
\BIC{$x = y \to y = z \to x = z$}
\UIC{$\forall z\,(x = y \to y = z \to x = z)$}
\UIC{$\forall y\,\forall z\,(x = y \to y = z \to x = z)$}
\UIC{$\forall x\,\forall y\,\forall z\,(x = y \to y = z \to x = z)$}
\AXC{$\forall x\,\forall y\,\forall z\,(x = y \to y = z \to x = z) \to \varphi_{7}$\lbl{A2}}
\BIC{$\varphi_{7}$}
\end{ptree}
\noindent$\Pi_{51}$:
\begin{ptree}
\AXC{$\stackrel{\Pi_{50}}{\vdots}$}\noLine
\UIC{$\varphi_{7}$}
\AXC{$\varphi_{7} \to \varphi_{111}$\lbl{A2}}
\BIC{$\varphi_{111}$}
\AXC{$\varphi_{111} \to M \cdot h + \lvert \beta \rvert \cdot h = h \cdot M + \lvert \beta \rvert \cdot h \to \varphi_{88}$\lbl{A2}}
\BIC{$M \cdot h + \lvert \beta \rvert \cdot h = h \cdot M + \lvert \beta \rvert \cdot h \to \varphi_{88}$}
\end{ptree}
\noindent$\Pi_{52}$:
\begin{ptree}
\AXC{$z = w \to \varphi_{22}$\lbl{A6}}
\UIC{$\forall w\,(z = w \to \varphi_{22})$}
\UIC{$\forall z\,\forall w\,(z = w \to \varphi_{22})$}
\AXC{$\forall z\,\forall w\,(z = w \to \varphi_{22}) \to \varphi_{112}$\lbl{A2}}
\BIC{$\varphi_{112}$}
\AXC{$\varphi_{112} \to \lvert \beta \rvert \cdot h = h \cdot \lvert \beta \rvert \to \varphi_{89}$\lbl{A2}}
\BIC{$\lvert \beta \rvert \cdot h = h \cdot \lvert \beta \rvert \to \varphi_{89}$}
\end{ptree}
\noindent$\Pi_{53}$:
\begin{ptree}
\AXC{$\forall x\,\forall y\,(x \in \mathbb{R} \to y \in \mathbb{R} \to x \cdot y = y \cdot x)$\lbl{(R7)}}
\AXC{$\varphi_{122}$\lbl{A2}}
\BIC{$\forall y\,(\lvert \beta \rvert \in \mathbb{R} \to y \in \mathbb{R} \to \lvert \beta \rvert \cdot y = y \cdot \lvert \beta \rvert)$}
\AXC{$\varphi_{131}$\lbl{A2}}
\BIC{$\lvert \beta \rvert \in \mathbb{R} \to h \in \mathbb{R} \to \lvert \beta \rvert \cdot h = h \cdot \lvert \beta \rvert$}
\AXC{$(\lvert \beta \rvert \in \mathbb{R} \to h \in \mathbb{R} \to \lvert \beta \rvert \cdot h = h \cdot \lvert \beta \rvert) \to \varphi_{235}$\lbl{A1}}
\BIC{$\varphi_{235}$}
\end{ptree}
\noindent$\Pi_{54}$:
\begin{ptree}
\AXC{$x \in \mathbb{R} \to y \in \mathbb{R} \to \alpha \in \mathbb{R} \to \beta \in \mathbb{R} \to M \in \mathbb{R} \to \varphi_{53}$\lbl{A1}}
\AXC{$\stackrel{\Pi_{11}}{\vdots}$}\noLine
\UIC{$\varphi_{172}$}
\AXC{$\stackrel{\Pi_{53}}{\vdots}$}\noLine
\UIC{$\varphi_{235}$}
\BIC{$(x \in \mathbb{R} \to y \in \mathbb{R} \to \alpha \in \mathbb{R} \to \beta \in \mathbb{R} \to M \in \mathbb{R} \to \varphi_{53}) \to \varphi_{196}$}
\BIC{$\varphi_{196}$}
\AXC{$\varphi_{196} \to \varphi_{222}$\lbl{A1}}
\BIC{$\varphi_{222}$}
\end{ptree}
\noindent$\Pi_{55}$:
\begin{ptree}
\AXC{$\stackrel{\Pi_{52}}{\vdots}$}\noLine
\UIC{$\lvert \beta \rvert \cdot h = h \cdot \lvert \beta \rvert \to \varphi_{89}$}
\AXC{$x = x$\lbl{A5}}
\UIC{$\forall x\,(x = x)$}
\AXC{$\forall x\,(x = x) \to h \cdot M + \lvert \beta \rvert \cdot h = h \cdot M + \lvert \beta \rvert \cdot h$\lbl{A2}}
\BIC{$h \cdot M + \lvert \beta \rvert \cdot h = h \cdot M + \lvert \beta \rvert \cdot h$}
\AXC{$\stackrel{\Pi_{54}}{\vdots}$}\noLine
\UIC{$\varphi_{222}$}
\BIC{$(\lvert \beta \rvert \cdot h = h \cdot \lvert \beta \rvert \to \varphi_{89}) \to x \in \mathbb{R} \to y \in \mathbb{R} \to \varphi_{183}$}
\BIC{$x \in \mathbb{R} \to y \in \mathbb{R} \to \varphi_{183}$}
\end{ptree}
\noindent$\Pi_{56}$:
\begin{ptree}
\AXC{$z = w \to \varphi_{71}$\lbl{A6}}
\UIC{$\forall w\,(z = w \to \varphi_{71})$}
\UIC{$\forall z\,\forall w\,(z = w \to \varphi_{71})$}
\AXC{$\forall z\,\forall w\,(z = w \to \varphi_{71}) \to \varphi_{113}$\lbl{A2}}
\BIC{$\varphi_{113}$}
\AXC{$\varphi_{113} \to M \cdot h = h \cdot M \to \varphi_{90}$\lbl{A2}}
\BIC{$M \cdot h = h \cdot M \to \varphi_{90}$}
\end{ptree}
\noindent$\Pi_{57}$:
\begin{ptree}
\AXC{$\forall x\,\forall y\,(x \in \mathbb{R} \to y \in \mathbb{R} \to x \cdot y = y \cdot x)$\lbl{(R7)}}
\AXC{$\varphi_{83}$\lbl{A2}}
\BIC{$\forall y\,(M \in \mathbb{R} \to y \in \mathbb{R} \to M \cdot y = y \cdot M)$}
\AXC{$\varphi_{73}$\lbl{A2}}
\BIC{$M \in \mathbb{R} \to h \in \mathbb{R} \to M \cdot h = h \cdot M$}
\AXC{$(M \in \mathbb{R} \to h \in \mathbb{R} \to M \cdot h = h \cdot M) \to \varphi_{227}$\lbl{A1}}
\BIC{$\varphi_{227}$}
\end{ptree}
\noindent$\Pi_{58}$:
\begin{ptree}
\AXC{$x \in \mathbb{R} \to y \in \mathbb{R} \to \alpha \in \mathbb{R} \to \beta \in \mathbb{R} \to M \in \mathbb{R} \to \varphi_{53}$\lbl{A1}}
\AXC{$x \in \mathbb{R} \to y \in \mathbb{R} \to \alpha \in \mathbb{R} \to \beta \in \mathbb{R} \to M \in \mathbb{R} \to \varphi_{55}$\lbl{A1}}
\AXC{$\stackrel{\Pi_{57}}{\vdots}$}\noLine
\UIC{$\varphi_{227}$}
\BIC{$(x \in \mathbb{R} \to y \in \mathbb{R} \to \alpha \in \mathbb{R} \to \beta \in \mathbb{R} \to M \in \mathbb{R} \to \varphi_{53}) \to \varphi_{156}$}
\BIC{$\varphi_{156}$}
\end{ptree}
\noindent$\Pi_{59}$:
\begin{ptree}
\AXC{$x = x$\lbl{A5}}
\UIC{$\forall x\,(x = x)$}
\AXC{$\forall x\,(x = x) \to M \cdot h + \lvert \beta \rvert \cdot h = M \cdot h + \lvert \beta \rvert \cdot h$\lbl{A2}}
\BIC{$M \cdot h + \lvert \beta \rvert \cdot h = M \cdot h + \lvert \beta \rvert \cdot h$}
\AXC{$\stackrel{\Pi_{58}}{\vdots}$}\noLine
\UIC{$\varphi_{156}$}
\AXC{$\varphi_{156} \to \varphi_{221}$\lbl{A1}}
\BIC{$\varphi_{221}$}
\BIC{$(M \cdot h = h \cdot M \to \varphi_{90}) \to x \in \mathbb{R} \to y \in \mathbb{R} \to \varphi_{184}$}
\end{ptree}
\noindent$\Pi_{60}$:
\begin{ptree}
\AXC{$\stackrel{\Pi_{56}}{\vdots}$}\noLine
\UIC{$M \cdot h = h \cdot M \to \varphi_{90}$}
\AXC{$\stackrel{\Pi_{59}}{\vdots}$}\noLine
\UIC{$(M \cdot h = h \cdot M \to \varphi_{90}) \to x \in \mathbb{R} \to y \in \mathbb{R} \to \varphi_{184}$}
\BIC{$x \in \mathbb{R} \to y \in \mathbb{R} \to \varphi_{184}$}
\AXC{$(x \in \mathbb{R} \to y \in \mathbb{R} \to \varphi_{184}) \to (x \in \mathbb{R} \to y \in \mathbb{R} \to \varphi_{183}) \to \varphi_{217}$\lbl{A1}}
\BIC{$(x \in \mathbb{R} \to y \in \mathbb{R} \to \varphi_{183}) \to \varphi_{217}$}
\end{ptree}
\noindent$\Pi_{61}$:
\begin{ptree}
\AXC{$\stackrel{\Pi_{51}}{\vdots}$}\noLine
\UIC{$M \cdot h + \lvert \beta \rvert \cdot h = h \cdot M + \lvert \beta \rvert \cdot h \to \varphi_{88}$}
\AXC{$\stackrel{\Pi_{55}}{\vdots}$}\noLine
\UIC{$x \in \mathbb{R} \to y \in \mathbb{R} \to \varphi_{183}$}
\AXC{$\stackrel{\Pi_{60}}{\vdots}$}\noLine
\UIC{$(x \in \mathbb{R} \to y \in \mathbb{R} \to \varphi_{183}) \to \varphi_{217}$}
\BIC{$\varphi_{217}$}
\BIC{$x \in \mathbb{R} \to y \in \mathbb{R} \to \varphi_{185}$}
\end{ptree}
\noindent$\Pi_{62}$:
\begin{ptree}
\AXC{$\stackrel{\Pi_{49}}{\vdots}$}\noLine
\UIC{$x \in \mathbb{R} \to y \in \mathbb{R} \to \varphi_{163}$}
\AXC{$\stackrel{\Pi_{61}}{\vdots}$}\noLine
\UIC{$x \in \mathbb{R} \to y \in \mathbb{R} \to \varphi_{185}$}
\AXC{$(x \in \mathbb{R} \to y \in \mathbb{R} \to \varphi_{185}) \to (x \in \mathbb{R} \to y \in \mathbb{R} \to \varphi_{163}) \to \varphi_{214}$\lbl{A1}}
\BIC{$(x \in \mathbb{R} \to y \in \mathbb{R} \to \varphi_{163}) \to \varphi_{214}$}
\BIC{$\varphi_{214}$}
\end{ptree}
\noindent$\Pi_{63}$:
\begin{ptree}
\AXC{$\stackrel{\Pi_{43}}{\vdots}$}\noLine
\UIC{$\varphi_{7}$}
\AXC{$\varphi_{7} \to \varphi_{110}$\lbl{A2}}
\BIC{$\varphi_{110}$}
\AXC{$\varphi_{110} \to M \cdot h + \lvert \beta \rvert \cdot h = h \cdot M + h \cdot \lvert \beta \rvert \to \varphi_{70}$\lbl{A2}}
\BIC{$M \cdot h + \lvert \beta \rvert \cdot h = h \cdot M + h \cdot \lvert \beta \rvert \to \varphi_{70}$}
\AXC{$\stackrel{\Pi_{62}}{\vdots}$}\noLine
\UIC{$\varphi_{214}$}
\BIC{$x \in \mathbb{R} \to y \in \mathbb{R} \to \varphi_{164}$}
\end{ptree}
\noindent$\Pi_{64}$:
\begin{ptree}
\AXC{$\stackrel{\Pi_{63}}{\vdots}$}\noLine
\UIC{$x \in \mathbb{R} \to y \in \mathbb{R} \to \varphi_{164}$}
\AXC{$(x \in \mathbb{R} \to y \in \mathbb{R} \to \varphi_{164}) \to (x \in \mathbb{R} \to y \in \mathbb{R} \to \varphi_{151}) \to \varphi_{213}$\lbl{A1}}
\BIC{$(x \in \mathbb{R} \to y \in \mathbb{R} \to \varphi_{151}) \to \varphi_{213}$}
\end{ptree}
\noindent$\Pi_{65}$:
\begin{ptree}
\AXC{$\stackrel{\Pi_{41}}{\vdots}$}\noLine
\UIC{$\varphi_{190}$}
\AXC{$x \in \mathbb{R} \to y \in \mathbb{R} \to \alpha \in \mathbb{R} \to \beta \in \mathbb{R} \to M \in \mathbb{R} \to \varphi_{53}$\lbl{A1}}
\AXC{$\stackrel{\Pi_{42}}{\vdots}$}\noLine
\UIC{$\varphi_{237}$}
\BIC{$\varphi_{190} \to x \in \mathbb{R} \to y \in \mathbb{R} \to \varphi_{151}$}
\BIC{$x \in \mathbb{R} \to y \in \mathbb{R} \to \varphi_{151}$}
\AXC{$\stackrel{\Pi_{64}}{\vdots}$}\noLine
\UIC{$(x \in \mathbb{R} \to y \in \mathbb{R} \to \varphi_{151}) \to \varphi_{213}$}
\BIC{$\varphi_{213}$}
\end{ptree}
\noindent$\Pi_{66}$:
\begin{ptree}
\AXC{$\stackrel{\Pi_{39}}{\vdots}$}\noLine
\UIC{$\varphi_{7}$}
\AXC{$\varphi_{7} \to \varphi_{86}$\lbl{A2}}
\BIC{$\varphi_{86}$}
\AXC{$\varphi_{86} \to M \cdot h + \lvert \beta \rvert \cdot h = h \cdot (M + \lvert \beta \rvert) \to \varphi_{50}$\lbl{A2}}
\BIC{$M \cdot h + \lvert \beta \rvert \cdot h = h \cdot (M + \lvert \beta \rvert) \to \varphi_{50}$}
\AXC{$\stackrel{\Pi_{65}}{\vdots}$}\noLine
\UIC{$\varphi_{213}$}
\BIC{$x \in \mathbb{R} \to y \in \mathbb{R} \to \varphi_{165}$}
\AXC{$(x \in \mathbb{R} \to y \in \mathbb{R} \to \varphi_{165}) \to \varphi_{249}$\lbl{A1}}
\BIC{$\varphi_{249}$}
\end{ptree}
\noindent$\Pi_{67}$:
\begin{ptree}
\AXC{$\stackrel{\Pi_{2}}{\vdots}$}\noLine
\UIC{$M \cdot h + \lvert \beta \rvert \cdot h = (M + \lvert \beta \rvert) \cdot h \to \varphi_{155}$}
\AXC{$\stackrel{\Pi_{38}}{\vdots}$}\noLine
\UIC{$x \in \mathbb{R} \to y \in \mathbb{R} \to \alpha \in \mathbb{R} \to \varphi_{204}$}
\AXC{$\stackrel{\Pi_{66}}{\vdots}$}\noLine
\UIC{$\varphi_{249}$}
\BIC{$(M \cdot h + \lvert \beta \rvert \cdot h = (M + \lvert \beta \rvert) \cdot h \to \varphi_{155}) \to \varphi_{224}$}
\BIC{$\varphi_{224}$}
\end{ptree}
\noindent$\Pi_{68}$:
\begin{ptree}
\AXC{$z = w \to \varphi_{118}$\lbl{A6}}
\UIC{$\forall w\,(z = w \to \varphi_{118})$}
\UIC{$\forall z\,\forall w\,(z = w \to \varphi_{118})$}
\AXC{$\forall z\,\forall w\,(z = w \to \varphi_{118}) \to \varphi_{154}$\lbl{A2}}
\BIC{$\varphi_{154}$}
\AXC{$\varphi_{154} \to \lvert \beta \cdot (x - \alpha) \rvert = \lvert \beta \rvert \cdot \lvert x - \alpha \rvert \to \varphi_{144}$\lbl{A2}}
\BIC{$\lvert \beta \cdot (x - \alpha) \rvert = \lvert \beta \rvert \cdot \lvert x - \alpha \rvert \to \varphi_{144}$}
\end{ptree}
\noindent$\Pi_{69}$:
\begin{ptree}
\AXC{$z = w \to \varphi_{125}$\lbl{A6}}
\UIC{$\forall w\,(z = w \to \varphi_{125})$}
\UIC{$\forall z\,\forall w\,(z = w \to \varphi_{125})$}
\AXC{$\forall z\,\forall w\,(z = w \to \varphi_{125}) \to \varphi_{179}$\lbl{A2}}
\BIC{$\varphi_{179}$}
\AXC{$\varphi_{179} \to \lvert \beta \cdot (x - \alpha) \rvert = \lvert \beta \rvert \cdot \lvert x - \alpha \rvert \to \varphi_{157}$\lbl{A2}}
\BIC{$\lvert \beta \cdot (x - \alpha) \rvert = \lvert \beta \rvert \cdot \lvert x - \alpha \rvert \to \varphi_{157}$}
\end{ptree}
\noindent$\Pi_{70}$:
\begin{ptree}
\AXC{$z = w \to \varphi_{114}$\lbl{A6}}
\UIC{$\forall w\,(z = w \to \varphi_{114})$}
\UIC{$\forall z\,\forall w\,(z = w \to \varphi_{114})$}
\AXC{$\forall z\,\forall w\,(z = w \to \varphi_{114}) \to \varphi_{143}$\lbl{A2}}
\BIC{$\varphi_{143}$}
\AXC{$\varphi_{143} \to \lvert x \cdot (y - \beta) \rvert = \lvert x \rvert \cdot \lvert y - \beta \rvert \to \varphi_{140}$\lbl{A2}}
\BIC{$\lvert x \cdot (y - \beta) \rvert = \lvert x \rvert \cdot \lvert y - \beta \rvert \to \varphi_{140}$}
\end{ptree}
\noindent$\Pi_{71}$:
\begin{ptree}
\AXC{$z = w \to \varphi_{97}$\lbl{A6}}
\UIC{$\forall w\,(z = w \to \varphi_{97})$}
\UIC{$\forall z\,\forall w\,(z = w \to \varphi_{97})$}
\AXC{$\forall z\,\forall w\,(z = w \to \varphi_{97}) \to \varphi_{193}$\lbl{A2}}
\BIC{$\varphi_{193}$}
\AXC{$\varphi_{193} \to x \cdot (y - \beta) + \beta \cdot (x - \alpha) = x \cdot y - \alpha \cdot \beta \to \varphi_{141}$\lbl{A2}}
\BIC{$x \cdot (y - \beta) + \beta \cdot (x - \alpha) = x \cdot y - \alpha \cdot \beta \to \varphi_{141}$}
\end{ptree}
\noindent$\Pi_{72}$:
\begin{ptree}
\AXC{$z = w \to \varphi_{109}$\lbl{A6}}
\UIC{$\forall w\,(z = w \to \varphi_{109})$}
\UIC{$\forall z\,\forall w\,(z = w \to \varphi_{109})$}
\AXC{$\forall z\,\forall w\,(z = w \to \varphi_{109}) \to \varphi_{200}$\lbl{A2}}
\BIC{$\varphi_{200}$}
\AXC{$\varphi_{200} \to x \cdot (y - \beta) + \beta \cdot (x - \alpha) = x \cdot y - \alpha \cdot \beta \to \varphi_{152}$\lbl{A2}}
\BIC{$x \cdot (y - \beta) + \beta \cdot (x - \alpha) = x \cdot y - \alpha \cdot \beta \to \varphi_{152}$}
\end{ptree}
\noindent$\Pi_{73}$:
\begin{ptree}
\AXC{$\forall x\,\forall y\,(x \in \mathbb{R} \to y \in \mathbb{R} \to x \cdot y \in \mathbb{R})$\lbl{(R1)}}
\AXC{$\varphi_{98}$\lbl{A2}}
\BIC{$\forall y\,(\beta \in \mathbb{R} \to y \in \mathbb{R} \to \beta \cdot y \in \mathbb{R})$}
\AXC{$\varphi_{116}$\lbl{A2}}
\BIC{$\beta \in \mathbb{R} \to x - \alpha \in \mathbb{R} \to \beta \cdot (x - \alpha) \in \mathbb{R}$}
\AXC{$(\beta \in \mathbb{R} \to x - \alpha \in \mathbb{R} \to \beta \cdot (x - \alpha) \in \mathbb{R}) \to \varphi_{234}$\lbl{A1}}
\BIC{$\varphi_{234}$}
\end{ptree}
\noindent$\Pi_{74}$:
\begin{ptree}
\AXC{$\stackrel{\Pi_{4}}{\vdots}$}\noLine
\UIC{$\varphi_{171}$}
\AXC{$x \in \mathbb{R} \to y \in \mathbb{R} \to \alpha \in \mathbb{R} \to \beta \in \mathbb{R} \to M \in \mathbb{R} \to \varphi_{63}$\lbl{A1}}
\AXC{$\stackrel{\Pi_{73}}{\vdots}$}\noLine
\UIC{$\varphi_{234}$}
\BIC{$\varphi_{171} \to \varphi_{192}$}
\BIC{$\varphi_{192}$}
\end{ptree}
\noindent$\Pi_{75}$:
\begin{ptree}
\AXC{$\forall x\,\forall y\,(x \in \mathbb{R} \to y \in \mathbb{R} \to x \cdot y \in \mathbb{R})$\lbl{(R1)}}
\AXC{$\varphi_{95}$\lbl{A2}}
\BIC{$\forall y\,(x \in \mathbb{R} \to y \in \mathbb{R} \to x \cdot y \in \mathbb{R})$}
\AXC{$\varphi_{104}$\lbl{A2}}
\BIC{$x \in \mathbb{R} \to y - \beta \in \mathbb{R} \to x \cdot (y - \beta) \in \mathbb{R}$}
\AXC{$(x \in \mathbb{R} \to y - \beta \in \mathbb{R} \to x \cdot (y - \beta) \in \mathbb{R}) \to \varphi_{233}$\lbl{A1}}
\BIC{$\varphi_{233}$}
\end{ptree}
\noindent$\Pi_{76}$:
\begin{ptree}
\AXC{$\stackrel{\Pi_{18}}{\vdots}$}\noLine
\UIC{$\varphi_{173}$}
\AXC{$x \in \mathbb{R} \to y \in \mathbb{R} \to \alpha \in \mathbb{R} \to \beta \in \mathbb{R} \to M \in \mathbb{R} \to \varphi_{52}$\lbl{A1}}
\AXC{$\stackrel{\Pi_{75}}{\vdots}$}\noLine
\UIC{$\varphi_{233}$}
\BIC{$\varphi_{173} \to \varphi_{191}$}
\BIC{$\varphi_{191}$}
\end{ptree}
\noindent$\Pi_{77}$:
\begin{ptree}
\AXC{$\forall a\,\forall b\,(a \in \mathbb{R} \to b \in \mathbb{R} \to \lvert a + b \rvert \leq \lvert a \rvert + \lvert b \rvert)$\lbl{命題1.2}}
\AXC{$\varphi_{175}$\lbl{A2}}
\BIC{$\varphi_{84}$}
\AXC{$\varphi_{84} \to x \cdot (y - \beta) \in \mathbb{R} \to \beta \cdot (x - \alpha) \in \mathbb{R} \to \varphi_{39}$\lbl{A2}}
\BIC{$x \cdot (y - \beta) \in \mathbb{R} \to \beta \cdot (x - \alpha) \in \mathbb{R} \to \varphi_{39}$}
\end{ptree}
\noindent$\Pi_{78}$:
\begin{ptree}
\AXC{$\stackrel{\Pi_{74}}{\vdots}$}\noLine
\UIC{$\varphi_{192}$}
\AXC{$\stackrel{\Pi_{76}}{\vdots}$}\noLine
\UIC{$\varphi_{191}$}
\AXC{$\stackrel{\Pi_{77}}{\vdots}$}\noLine
\UIC{$x \cdot (y - \beta) \in \mathbb{R} \to \beta \cdot (x - \alpha) \in \mathbb{R} \to \varphi_{39}$}
\AXC{$\varphi_{246}$\lbl{A1}}
\BIC{$\varphi_{191} \to \varphi_{192} \to x \in \mathbb{R} \to y \in \mathbb{R} \to \alpha \in \mathbb{R} \to \beta \in \mathbb{R} \to M \in \mathbb{R} \to \varphi_{145}$}
\BIC{$\varphi_{192} \to x \in \mathbb{R} \to y \in \mathbb{R} \to \alpha \in \mathbb{R} \to \beta \in \mathbb{R} \to M \in \mathbb{R} \to \varphi_{145}$}
\BIC{$x \in \mathbb{R} \to y \in \mathbb{R} \to \alpha \in \mathbb{R} \to \beta \in \mathbb{R} \to M \in \mathbb{R} \to \varphi_{145}$}
\end{ptree}
\noindent$\Pi_{79}$:
\begin{ptree}
\AXC{$\stackrel{\Pi_{44}}{\vdots}$}\noLine
\UIC{$\forall x\,\forall y\,(x = y \to y = x)$}
\AXC{$\forall x\,\forall y\,(x = y \to y = x) \to \forall y\,(v_{687} = y \to y = v_{687})$\lbl{A2}}
\BIC{$\forall y\,(v_{687} = y \to y = v_{687})$}
\AXC{$\varphi_{1}$\lbl{A2}}
\BIC{$v_{687} = v_{688} \to v_{688} = v_{687}$}
\UIC{$\forall v_{688}\,(v_{687} = v_{688} \to v_{688} = v_{687})$}
\UIC{$\forall v_{687}\,\forall v_{688}\,(v_{687} = v_{688} \to v_{688} = v_{687})$}
\AXC{$\varphi_{85}$\lbl{A2}}
\BIC{$\forall v_{688}\,(x \cdot y - \alpha \cdot \beta = v_{688} \to v_{688} = x \cdot y - \alpha \cdot \beta)$}
\AXC{$\varphi_{186}$\lbl{A2}}
\BIC{$\varphi_{101}$}
\end{ptree}
\noindent$\Pi_{80}$:
\begin{ptree}
\AXC{$\forall x\,\forall y\,(\varphi_{107})$\lbl{[F18]}}
\AXC{$\forall x\,\forall y\,(\varphi_{107}) \to \forall y\,(\varphi_{107})$\lbl{A2}}
\BIC{$\forall y\,(\varphi_{107})$}
\AXC{$\forall y\,(\varphi_{107}) \to \varphi_{107}$\lbl{A2}}
\BIC{$\varphi_{107}$}
\AXC{$\varphi_{107} \to \varphi_{108}$\lbl{A2}}
\BIC{$\varphi_{108}$}
\AXC{$\varphi_{108} \to x \in \mathbb{R} \to y \in \mathbb{R} \to \varphi_{45}$\lbl{A2}}
\BIC{$x \in \mathbb{R} \to y \in \mathbb{R} \to \varphi_{45}$}
\end{ptree}
\noindent$\Pi_{81}$:
\begin{ptree}
\AXC{$x \in \mathbb{R} \to y \in \mathbb{R} \to \alpha \in \mathbb{R} \to \beta \in \mathbb{R} \to M \in \mathbb{R} \to \varphi_{52}$\lbl{A1}}
\AXC{$\stackrel{\Pi_{80}}{\vdots}$}\noLine
\UIC{$x \in \mathbb{R} \to y \in \mathbb{R} \to \varphi_{45}$}
\AXC{$\varphi_{252}$\lbl{A1}}
\BIC{$(x \in \mathbb{R} \to y \in \mathbb{R} \to \alpha \in \mathbb{R} \to \beta \in \mathbb{R} \to M \in \mathbb{R} \to \varphi_{52}) \to \varphi_{245}$}
\BIC{$\varphi_{245}$}
\end{ptree}
\noindent$\Pi_{82}$:
\begin{ptree}
\AXC{$x \in \mathbb{R} \to y \in \mathbb{R} \to \alpha \in \mathbb{R} \to \beta \in \mathbb{R} \to M \in \mathbb{R} \to \varphi_{62}$\lbl{A1}}
\AXC{$x \in \mathbb{R} \to y \in \mathbb{R} \to \alpha \in \mathbb{R} \to \beta \in \mathbb{R} \to M \in \mathbb{R} \to \varphi_{54}$\lbl{A1}}
\AXC{$\stackrel{\Pi_{81}}{\vdots}$}\noLine
\UIC{$\varphi_{245}$}
\BIC{$(x \in \mathbb{R} \to y \in \mathbb{R} \to \alpha \in \mathbb{R} \to \beta \in \mathbb{R} \to M \in \mathbb{R} \to \varphi_{62}) \to \varphi_{219}$}
\BIC{$\varphi_{219}$}
\end{ptree}
\noindent$\Pi_{83}$:
\begin{ptree}
\AXC{$\stackrel{\Pi_{79}}{\vdots}$}\noLine
\UIC{$\varphi_{101}$}
\AXC{$x \in \mathbb{R} \to y \in \mathbb{R} \to \alpha \in \mathbb{R} \to \beta \in \mathbb{R} \to M \in \mathbb{R} \to \varphi_{63}$\lbl{A1}}
\AXC{$\stackrel{\Pi_{82}}{\vdots}$}\noLine
\UIC{$\varphi_{219}$}
\BIC{$x \in \mathbb{R} \to y \in \mathbb{R} \to \alpha \in \mathbb{R} \to \varphi_{148}$}
\AXC{$\varphi_{236}$\lbl{A1}}
\BIC{$\varphi_{101} \to x \in \mathbb{R} \to y \in \mathbb{R} \to \alpha \in \mathbb{R} \to \varphi_{149}$}
\BIC{$x \in \mathbb{R} \to y \in \mathbb{R} \to \alpha \in \mathbb{R} \to \varphi_{149}$}
\end{ptree}
\noindent$\Pi_{84}$:
\begin{ptree}
\AXC{$\stackrel{\Pi_{78}}{\vdots}$}\noLine
\UIC{$x \in \mathbb{R} \to y \in \mathbb{R} \to \alpha \in \mathbb{R} \to \beta \in \mathbb{R} \to M \in \mathbb{R} \to \varphi_{145}$}
\AXC{$\stackrel{\Pi_{83}}{\vdots}$}\noLine
\UIC{$x \in \mathbb{R} \to y \in \mathbb{R} \to \alpha \in \mathbb{R} \to \varphi_{149}$}
\AXC{$(x \in \mathbb{R} \to y \in \mathbb{R} \to \alpha \in \mathbb{R} \to \varphi_{149}) \to \varphi_{250}$\lbl{A1}}
\BIC{$\varphi_{250}$}
\BIC{$(x \cdot (y - \beta) + \beta \cdot (x - \alpha) = x \cdot y - \alpha \cdot \beta \to \varphi_{152}) \to \varphi_{225}$}
\end{ptree}
\noindent$\Pi_{85}$:
\begin{ptree}
\AXC{$\stackrel{\Pi_{72}}{\vdots}$}\noLine
\UIC{$x \cdot (y - \beta) + \beta \cdot (x - \alpha) = x \cdot y - \alpha \cdot \beta \to \varphi_{152}$}
\AXC{$\stackrel{\Pi_{84}}{\vdots}$}\noLine
\UIC{$(x \cdot (y - \beta) + \beta \cdot (x - \alpha) = x \cdot y - \alpha \cdot \beta \to \varphi_{152}) \to \varphi_{225}$}
\BIC{$\varphi_{225}$}
\end{ptree}
\noindent$\Pi_{86}$:
\begin{ptree}
\AXC{$\forall x\,\forall y\,(x \in \mathbb{R} \to y \in \mathbb{R} \to \lvert x \cdot y \rvert = \lvert x \rvert \cdot \lvert y \rvert)$\lbl{[A17]}}
\AXC{$\varphi_{132}$\lbl{A2}}
\BIC{$\forall y\,(x \in \mathbb{R} \to y \in \mathbb{R} \to \lvert x \cdot y \rvert = \lvert x \rvert \cdot \lvert y \rvert)$}
\AXC{$\forall y\,(x \in \mathbb{R} \to y \in \mathbb{R} \to \lvert x \cdot y \rvert = \lvert x \rvert \cdot \lvert y \rvert) \to \varphi_{24}$\lbl{A2}}
\BIC{$\varphi_{24}$}
\end{ptree}
\noindent$\Pi_{87}$:
\begin{ptree}
\AXC{$\stackrel{\Pi_{18}}{\vdots}$}\noLine
\UIC{$\varphi_{173}$}
\AXC{$x \in \mathbb{R} \to y \in \mathbb{R} \to \alpha \in \mathbb{R} \to \beta \in \mathbb{R} \to M \in \mathbb{R} \to \varphi_{52}$\lbl{A1}}
\AXC{$\stackrel{\Pi_{86}}{\vdots}$}\noLine
\UIC{$\varphi_{24}$}
\AXC{$\varphi_{24} \to \varphi_{239}$\lbl{A1}}
\BIC{$\varphi_{239}$}
\BIC{$\varphi_{173} \to x \in \mathbb{R} \to y \in \mathbb{R} \to \varphi_{166}$}
\BIC{$x \in \mathbb{R} \to y \in \mathbb{R} \to \varphi_{166}$}
\AXC{$(x \in \mathbb{R} \to y \in \mathbb{R} \to \varphi_{166}) \to \varphi_{248}$\lbl{A1}}
\BIC{$\varphi_{248}$}
\end{ptree}
\noindent$\Pi_{88}$:
\begin{ptree}
\AXC{$\stackrel{\Pi_{70}}{\vdots}$}\noLine
\UIC{$\lvert x \cdot (y - \beta) \rvert = \lvert x \rvert \cdot \lvert y - \beta \rvert \to \varphi_{140}$}
\AXC{$\stackrel{\Pi_{71}}{\vdots}$}\noLine
\UIC{$x \cdot (y - \beta) + \beta \cdot (x - \alpha) = x \cdot y - \alpha \cdot \beta \to \varphi_{141}$}
\AXC{$\stackrel{\Pi_{85}}{\vdots}$}\noLine
\UIC{$\varphi_{225}$}
\BIC{$x \in \mathbb{R} \to y \in \mathbb{R} \to \alpha \in \mathbb{R} \to \varphi_{194}$}
\AXC{$\stackrel{\Pi_{87}}{\vdots}$}\noLine
\UIC{$\varphi_{248}$}
\BIC{$(\lvert x \cdot (y - \beta) \rvert = \lvert x \rvert \cdot \lvert y - \beta \rvert \to \varphi_{140}) \to \varphi_{223}$}
\BIC{$\varphi_{223}$}
\end{ptree}
\noindent$\Pi_{89}$:
\begin{ptree}
\AXC{$\forall x\,\forall y\,(x \in \mathbb{R} \to y \in \mathbb{R} \to \lvert x \cdot y \rvert = \lvert x \rvert \cdot \lvert y \rvert)$\lbl{[A17]}}
\AXC{$\varphi_{134}$\lbl{A2}}
\BIC{$\forall y\,(\beta \in \mathbb{R} \to y \in \mathbb{R} \to \lvert \beta \cdot y \rvert = \lvert \beta \rvert \cdot \lvert y \rvert)$}
\AXC{$\forall y\,(\beta \in \mathbb{R} \to y \in \mathbb{R} \to \lvert \beta \cdot y \rvert = \lvert \beta \rvert \cdot \lvert y \rvert) \to \varphi_{40}$\lbl{A2}}
\BIC{$\varphi_{40}$}
\end{ptree}
\noindent$\Pi_{90}$:
\begin{ptree}
\AXC{$\stackrel{\Pi_{4}}{\vdots}$}\noLine
\UIC{$\varphi_{171}$}
\AXC{$x \in \mathbb{R} \to y \in \mathbb{R} \to \alpha \in \mathbb{R} \to \beta \in \mathbb{R} \to M \in \mathbb{R} \to \varphi_{63}$\lbl{A1}}
\AXC{$\stackrel{\Pi_{89}}{\vdots}$}\noLine
\UIC{$\varphi_{40}$}
\AXC{$\varphi_{40} \to \varphi_{240}$\lbl{A1}}
\BIC{$\varphi_{240}$}
\BIC{$\varphi_{171} \to x \in \mathbb{R} \to y \in \mathbb{R} \to \varphi_{176}$}
\BIC{$x \in \mathbb{R} \to y \in \mathbb{R} \to \varphi_{176}$}
\AXC{$(x \in \mathbb{R} \to y \in \mathbb{R} \to \varphi_{176}) \to \varphi_{251}$\lbl{A1}}
\BIC{$\varphi_{251}$}
\end{ptree}
\noindent$\Pi_{91}$:
\begin{ptree}
\AXC{$z = w \to \varphi_{99}$\lbl{A6}}
\UIC{$\forall w\,(z = w \to \varphi_{99})$}
\UIC{$\forall z\,\forall w\,(z = w \to \varphi_{99})$}
\AXC{$\forall z\,\forall w\,(z = w \to \varphi_{99}) \to \varphi_{139}$\lbl{A2}}
\BIC{$\varphi_{139}$}
\AXC{$\varphi_{139} \to \lvert x \cdot (y - \beta) \rvert = \lvert x \rvert \cdot \lvert y - \beta \rvert \to \varphi_{137}$\lbl{A2}}
\BIC{$\lvert x \cdot (y - \beta) \rvert = \lvert x \rvert \cdot \lvert y - \beta \rvert \to \varphi_{137}$}
\AXC{$\stackrel{\Pi_{88}}{\vdots}$}\noLine
\UIC{$\varphi_{223}$}
\BIC{$x \in \mathbb{R} \to y \in \mathbb{R} \to \alpha \in \mathbb{R} \to \varphi_{198}$}
\AXC{$\stackrel{\Pi_{90}}{\vdots}$}\noLine
\UIC{$\varphi_{251}$}
\BIC{$(\lvert \beta \cdot (x - \alpha) \rvert = \lvert \beta \rvert \cdot \lvert x - \alpha \rvert \to \varphi_{157}) \to \varphi_{226}$}
\end{ptree}
\noindent$\Pi_{92}$:
\begin{ptree}
\AXC{$\stackrel{\Pi_{69}}{\vdots}$}\noLine
\UIC{$\lvert \beta \cdot (x - \alpha) \rvert = \lvert \beta \rvert \cdot \lvert x - \alpha \rvert \to \varphi_{157}$}
\AXC{$\stackrel{\Pi_{91}}{\vdots}$}\noLine
\UIC{$(\lvert \beta \cdot (x - \alpha) \rvert = \lvert \beta \rvert \cdot \lvert x - \alpha \rvert \to \varphi_{157}) \to \varphi_{226}$}
\BIC{$\varphi_{226}$}
\end{ptree}
\noindent$\Pi_{93}$:
\begin{ptree}
\AXC{$\forall x\,\forall y\,\forall z\,(x \leq y \to y \leq z \to x \leq z)$\lbl{[O5]}}
\AXC{$\varphi_{4}$\lbl{A2}}
\BIC{$\forall y\,\forall z\,(v_{695} \leq y \to y \leq z \to v_{695} \leq z)$}
\AXC{$\varphi_{5}$\lbl{A2}}
\BIC{$\forall z\,(v_{695} \leq v_{696} \to v_{696} \leq z \to v_{695} \leq z)$}
\AXC{$\varphi_{11}$\lbl{A2}}
\BIC{$v_{695} \leq v_{696} \to v_{696} \leq v_{697} \to v_{695} \leq v_{697}$}
\UIC{$\forall v_{697}\,(v_{695} \leq v_{696} \to v_{696} \leq v_{697} \to v_{695} \leq v_{697})$}
\UIC{$\varphi_{2}$}
\UIC{$\forall v_{695}\,(\varphi_{2})$}
\AXC{$\forall v_{695}\,(\varphi_{2}) \to \varphi_{68}$\lbl{A2}}
\BIC{$\varphi_{68}$}
\AXC{$\varphi_{68} \to \varphi_{201}$\lbl{A2}}
\BIC{$\varphi_{201}$}
\end{ptree}
\noindent$\Pi_{94}$:
\begin{ptree}
\AXC{$\stackrel{\Pi_{68}}{\vdots}$}\noLine
\UIC{$\lvert \beta \cdot (x - \alpha) \rvert = \lvert \beta \rvert \cdot \lvert x - \alpha \rvert \to \varphi_{144}$}
\AXC{$\stackrel{\Pi_{92}}{\vdots}$}\noLine
\UIC{$\varphi_{226}$}
\BIC{$x \in \mathbb{R} \to y \in \mathbb{R} \to \alpha \in \mathbb{R} \to \varphi_{205}$}
\AXC{$\stackrel{\Pi_{93}}{\vdots}$}\noLine
\UIC{$\varphi_{201}$}
\AXC{$\varphi_{201} \to \varphi_{209}$\lbl{A2}}
\BIC{$\varphi_{209}$}
\AXC{$\varphi_{209} \to \varphi_{244}$\lbl{A1}}
\BIC{$\varphi_{244}$}
\BIC{$(x \in \mathbb{R} \to y \in \mathbb{R} \to \alpha \in \mathbb{R} \to \varphi_{197}) \to x \in \mathbb{R} \to y \in \mathbb{R} \to \varphi_{170}$}
\end{ptree}
\noindent 主木:
\begin{ptree}
\AXC{$\stackrel{\Pi_{1}}{\vdots}$}\noLine
\UIC{$M \cdot h + \lvert \beta \rvert \cdot h = (M + \lvert \beta \rvert) \cdot h \to \varphi_{142}$}
\AXC{$\stackrel{\Pi_{67}}{\vdots}$}\noLine
\UIC{$\varphi_{224}$}
\BIC{$x \in \mathbb{R} \to y \in \mathbb{R} \to \alpha \in \mathbb{R} \to \varphi_{197}$}
\AXC{$\stackrel{\Pi_{94}}{\vdots}$}\noLine
\UIC{$(x \in \mathbb{R} \to y \in \mathbb{R} \to \alpha \in \mathbb{R} \to \varphi_{197}) \to x \in \mathbb{R} \to y \in \mathbb{R} \to \varphi_{170}$}
\BIC{$x \in \mathbb{R} \to y \in \mathbb{R} \to \varphi_{170}$}
\UIC{$\forall h\,(x \in \mathbb{R} \to y \in \mathbb{R} \to \varphi_{170})$}
\UIC{$\forall M\,\forall h\,(x \in \mathbb{R} \to y \in \mathbb{R} \to \varphi_{170})$}
\UIC{$\forall \beta\,\forall M\,\forall h\,(x \in \mathbb{R} \to y \in \mathbb{R} \to \varphi_{170})$}
\UIC{$\forall \alpha\,\forall \beta\,\forall M\,\forall h\,(x \in \mathbb{R} \to y \in \mathbb{R} \to \varphi_{170})$}
\UIC{$\forall y\,\forall \alpha\,\forall \beta\,\forall M\,\forall h\,(x \in \mathbb{R} \to y \in \mathbb{R} \to \varphi_{170})$}
\UIC{$\forall x\,\forall y\,\forall \alpha\,\forall \beta\,\forall M\,\forall h\,(x \in \mathbb{R} \to y \in \mathbb{R} \to \varphi_{170})$}
\end{ptree}
\begin{abbrevlist}
\item $\varphi_{1} :\equiv \forall y\,(v_{687} = y \to y = v_{687}) \to v_{687} = v_{688} \to v_{688} = v_{687}$
\item $\varphi_{2} :\equiv \forall v_{696}\,\forall v_{697}\,(v_{695} \leq v_{696} \to v_{696} \leq v_{697} \to v_{695} \leq v_{697})$
\item $\varphi_{3} :\equiv v_{691} \leq v_{692} \to v_{693} \leq v_{694} \to v_{691} + v_{693} \leq v_{692} + v_{694}$
\item $\varphi_{4} :\equiv \forall x\,\forall y\,\forall z\,(x \leq y \to y \leq z \to x \leq z) \to \forall y\,\forall z\,(v_{695} \leq y \to y \leq z \to v_{695} \leq z)$
\item $\varphi_{5} :\equiv \forall y\,\forall z\,(v_{695} \leq y \to y \leq z \to v_{695} \leq z) \to \forall z\,(v_{695} \leq v_{696} \to v_{696} \leq z \to v_{695} \leq z)$
\item $\varphi_{6} :\equiv \forall x\,\forall y\,(x = y \to y = x) \to \forall y\,(h \cdot (M + \lvert \beta \rvert) = y \to y = h \cdot (M + \lvert \beta \rvert))$
\item $\varphi_{7} :\equiv \forall y\,\forall z\,(M \cdot h + \lvert \beta \rvert \cdot h = y \to y = z \to M \cdot h + \lvert \beta \rvert \cdot h = z)$
\item $\varphi_{8} :\equiv \forall v_{689}\,\forall v_{690}\,(v_{689} \in \mathbb{R} \to v_{690} \in \mathbb{R} \to v_{689} - v_{690} \in \mathbb{R})$
\item $\varphi_{9} :\equiv 0 \leq M \to \lvert x \rvert \leq M \to \lvert x - \alpha \rvert < h \to \lvert y - \beta \rvert < h \to \lvert x \rvert \leq M$
\item $\varphi_{10} :\equiv 0 \leq M \to \lvert x \rvert \leq M \to \lvert x - \alpha \rvert < h \to \lvert y - \beta \rvert < h \to 0 \leq \lvert x \rvert$
\item $\varphi_{11} :\equiv \forall z\,(v_{695} \leq v_{696} \to v_{696} \leq z \to v_{695} \leq z) \to v_{695} \leq v_{696} \to v_{696} \leq v_{697} \to v_{695} \leq v_{697}$
\item $\varphi_{12} :\equiv 0 \leq M \to \lvert x \rvert \leq M \to \lvert x - \alpha \rvert < h \to \lvert y - \beta \rvert < h \to 0 \leq \lvert \beta \rvert$
\item $\varphi_{13} :\equiv \lvert y - \beta \rvert < h \to x \cdot y - \alpha \cdot \beta = x \cdot (y - \beta) + \beta \cdot (x - \alpha)$
\item $\varphi_{14} :\equiv \lvert y - \beta \rvert < h \to x \cdot (y - \beta) + \beta \cdot (x - \alpha) = x \cdot y - \alpha \cdot \beta$
\item $\varphi_{15} :\equiv \forall z\,(h \in \mathbb{R} \to M \in \mathbb{R} \to z \in \mathbb{R} \to h \cdot (M + z) = h \cdot M + h \cdot z)$
\item $\varphi_{16} :\equiv 0 \leq M \to \lvert x \rvert \leq M \to \lvert x - \alpha \rvert < h \to \lvert y - \beta \rvert < h \to M \cdot h = h \cdot M$
\item $\varphi_{17} :\equiv 0 \leq M \to \lvert x \rvert \leq M \to \lvert x - \alpha \rvert < h \to \lvert y - \beta \rvert < h \to \lvert x - \alpha \rvert \leq h$
\item $\varphi_{18} :\equiv 0 \leq M \to \lvert x \rvert \leq M \to \lvert x - \alpha \rvert < h \to \lvert y - \beta \rvert < h \to 0 \leq \lvert x - \alpha \rvert$
\item $\varphi_{19} :\equiv \lvert \beta \rvert \leq \lvert \beta \rvert \to \lvert x - \alpha \rvert \leq h \to \lvert \beta \rvert \cdot \lvert x - \alpha \rvert \leq \lvert \beta \rvert \cdot h$
\item $\varphi_{20} :\equiv 0 \leq M \to \lvert x \rvert \leq M \to \lvert x - \alpha \rvert < h \to \lvert y - \beta \rvert < h \to \lvert y - \beta \rvert \leq h$
\item $\varphi_{21} :\equiv 0 \leq M \to \lvert x \rvert \leq M \to \lvert x - \alpha \rvert < h \to \lvert y - \beta \rvert < h \to 0 \leq \lvert y - \beta \rvert$
\item $\varphi_{22} :\equiv h \cdot M + \lvert \beta \rvert \cdot h = h \cdot M + z \to h \cdot M + \lvert \beta \rvert \cdot h = h \cdot M + w$
\item $\varphi_{23} :\equiv h \in \mathbb{R} \to M + \lvert \beta \rvert \in \mathbb{R} \to h \cdot (M + \lvert \beta \rvert) = (M + \lvert \beta \rvert) \cdot h$
\item $\varphi_{24} :\equiv x \in \mathbb{R} \to y - \beta \in \mathbb{R} \to \lvert x \cdot (y - \beta) \rvert = \lvert x \rvert \cdot \lvert y - \beta \rvert$
\item $\varphi_{25} :\equiv 0 \leq M \to \lvert x \rvert \leq M \to \lvert x - \alpha \rvert < h \to \lvert y - \beta \rvert < h \to \lvert \beta \rvert \leq \lvert \beta \rvert$
\item $\varphi_{26} :\equiv 0 \leq M \to \lvert x \rvert \leq M \to \lvert x - \alpha \rvert < h \to \lvert y - \beta \rvert < h \to \lvert x \rvert \in \mathbb{R}$
\item $\varphi_{27} :\equiv 0 \leq M \to \lvert x \rvert \leq M \to \lvert x - \alpha \rvert < h \to \lvert y - \beta \rvert < h \to M \cdot h \in \mathbb{R}$
\item $\varphi_{28} :\equiv M \in \mathbb{R} \to \lvert \beta \rvert \in \mathbb{R} \to h \cdot (M + \lvert \beta \rvert) = h \cdot M + h \cdot \lvert \beta \rvert$
\item $\varphi_{29} :\equiv \lvert x \rvert \leq M \to \lvert x - \alpha \rvert < h \to \lvert y - \beta \rvert < h \to x \cdot (y - \beta) \in \mathbb{R}$
\item $\varphi_{30} :\equiv \forall y\,\forall z\,(h \in \mathbb{R} \to y \in \mathbb{R} \to z \in \mathbb{R} \to h \cdot (y + z) = h \cdot y + h \cdot z)$
\item $\varphi_{31} :\equiv 0 \leq M \to \lvert x \rvert \leq M \to \lvert x - \alpha \rvert < h \to \lvert y - \beta \rvert < h \to x - \alpha \in \mathbb{R}$
\item $\varphi_{32} :\equiv 0 \leq M \to \lvert x \rvert \leq M \to \lvert x - \alpha \rvert < h \to \lvert y - \beta \rvert < h \to \lvert \beta \rvert \in \mathbb{R}$
\item $\varphi_{33} :\equiv 0 \leq M \to \lvert x \rvert \leq M \to \lvert x - \alpha \rvert < h \to \lvert y - \beta \rvert < h \to y - \beta \in \mathbb{R}$
\item $\varphi_{34} :\equiv \lvert x \rvert \leq M \to \lvert x - \alpha \rvert < h \to \lvert y - \beta \rvert < h \to \beta \cdot (x - \alpha) \in \mathbb{R}$
\item $\varphi_{35} :\equiv 0 \leq M \to \lvert x \rvert \leq M \to \lvert x - \alpha \rvert < h \to \lvert y - \beta \rvert < h \to \lvert x - \alpha \rvert < h$
\item $\varphi_{36} :\equiv 0 \leq M \to \lvert x \rvert \leq M \to \lvert x - \alpha \rvert < h \to \lvert y - \beta \rvert < h \to \lvert y - \beta \rvert < h$
\item $\varphi_{37} :\equiv 0 \leq \lvert y - \beta \rvert \to \lvert x \rvert \leq M \to \lvert y - \beta \rvert \leq h \to \lvert x \rvert \cdot \lvert y - \beta \rvert \leq M \cdot h$
\item $\varphi_{38} :\equiv \lvert x \cdot (y - \beta) + \beta \cdot (x - \alpha) \rvert = \lvert x \cdot (y - \beta) \rvert + \lvert \beta \cdot (x - \alpha) \rvert$
\item $\varphi_{39} :\equiv \lvert x \cdot (y - \beta) + \beta \cdot (x - \alpha) \rvert \leq \lvert x \cdot (y - \beta) \rvert + \lvert \beta \cdot (x - \alpha) \rvert$
\item $\varphi_{40} :\equiv \beta \in \mathbb{R} \to x - \alpha \in \mathbb{R} \to \lvert \beta \cdot (x - \alpha) \rvert = \lvert \beta \rvert \cdot \lvert x - \alpha \rvert$
\item $\varphi_{41} :\equiv \lvert x - \alpha \rvert < h \to \lvert y - \beta \rvert < h \to h \cdot (M + \lvert \beta \rvert) = (M + \lvert \beta \rvert) \cdot h$
\item $\varphi_{42} :\equiv \forall x\,\forall y\,\forall z\,(x \in \mathbb{R} \to y \in \mathbb{R} \to z \in \mathbb{R} \to x \cdot (y + z) = x \cdot y + x \cdot z)$
\item $\varphi_{43} :\equiv \lvert x - \alpha \rvert < h \to \lvert y - \beta \rvert < h \to \lvert \beta \rvert \cdot \lvert x - \alpha \rvert \leq \lvert \beta \rvert \cdot h$
\item $\varphi_{44} :\equiv \lvert x \rvert \leq M \to \lvert x - \alpha \rvert < h \to \lvert y - \beta \rvert < h \to \lvert \beta \rvert \cdot h = h \cdot \lvert \beta \rvert$
\item $\varphi_{45} :\equiv \alpha \in \mathbb{R} \to \beta \in \mathbb{R} \to x \cdot y - \alpha \cdot \beta = x \cdot (y - \beta) + \beta \cdot (x - \alpha)$
\item $\varphi_{46} :\equiv \lvert x \cdot (y - \beta) + \beta \cdot (x - \alpha) \rvert < \lvert x \cdot (y - \beta) \rvert + \lvert \beta \cdot (x - \alpha) \rvert$
\item $\varphi_{47} :\equiv 0 \leq M \to \lvert x \rvert \leq M \to \lvert x - \alpha \rvert < h \to \lvert y - \beta \rvert < h \to \lvert x - \alpha \rvert \in \mathbb{R}$
\item $\varphi_{48} :\equiv 0 \leq M \to \lvert x \rvert \leq M \to \lvert x - \alpha \rvert < h \to \lvert y - \beta \rvert < h \to \lvert y - \beta \rvert \in \mathbb{R}$
\item $\varphi_{49} :\equiv 0 \leq M \to \lvert x \rvert \leq M \to \lvert x - \alpha \rvert < h \to \lvert y - \beta \rvert < h \to \lvert \beta \rvert \cdot h \in \mathbb{R}$
\item $\varphi_{50} :\equiv h \cdot (M + \lvert \beta \rvert) = (M + \lvert \beta \rvert) \cdot h \to M \cdot h + \lvert \beta \rvert \cdot h = (M + \lvert \beta \rvert) \cdot h$
\item $\varphi_{51} :\equiv 0 \leq M \to \lvert x \rvert \leq M \to \lvert x - \alpha \rvert < h \to \lvert y - \beta \rvert < h \to M + \lvert \beta \rvert \in \mathbb{R}$
\item $\varphi_{52} :\equiv h \in \mathbb{R} \to 0 \leq M \to \lvert x \rvert \leq M \to \lvert x - \alpha \rvert < h \to \lvert y - \beta \rvert < h \to x \in \mathbb{R}$
\item $\varphi_{53} :\equiv h \in \mathbb{R} \to 0 \leq M \to \lvert x \rvert \leq M \to \lvert x - \alpha \rvert < h \to \lvert y - \beta \rvert < h \to h \in \mathbb{R}$
\item $\varphi_{54} :\equiv h \in \mathbb{R} \to 0 \leq M \to \lvert x \rvert \leq M \to \lvert x - \alpha \rvert < h \to \lvert y - \beta \rvert < h \to y \in \mathbb{R}$
\item $\varphi_{55} :\equiv h \in \mathbb{R} \to 0 \leq M \to \lvert x \rvert \leq M \to \lvert x - \alpha \rvert < h \to \lvert y - \beta \rvert < h \to M \in \mathbb{R}$
\item $\varphi_{56} :\equiv \lvert x \rvert \leq M \to \lvert x - \alpha \rvert < h \to \lvert y - \beta \rvert < h \to \lvert x \rvert \cdot \lvert y - \beta \rvert \leq M \cdot h$
\item $\varphi_{57} :\equiv \lvert x - \alpha \rvert < h \to \lvert y - \beta \rvert < h \to h \cdot (M + \lvert \beta \rvert) = h \cdot M + h \cdot \lvert \beta \rvert$
\item $\varphi_{58} :\equiv \lvert x - \alpha \rvert < h \to \lvert y - \beta \rvert < h \to h \cdot M + h \cdot \lvert \beta \rvert = h \cdot (M + \lvert \beta \rvert)$
\item $\varphi_{59} :\equiv \lvert x - \alpha \rvert < h \to \lvert y - \beta \rvert < h \to M \cdot h + \lvert \beta \rvert \cdot h = h \cdot (M + \lvert \beta \rvert)$
\item $\varphi_{60} :\equiv \lvert x - \alpha \rvert < h \to \lvert y - \beta \rvert < h \to M \cdot h + \lvert \beta \rvert \cdot h = (M + \lvert \beta \rvert) \cdot h$
\item $\varphi_{61} :\equiv \lvert x - \alpha \rvert < h \to \lvert y - \beta \rvert < h \to \lvert x \cdot (y - \beta) \rvert = \lvert x \rvert \cdot \lvert y - \beta \rvert$
\item $\varphi_{62} :\equiv h \in \mathbb{R} \to 0 \leq M \to \lvert x \rvert \leq M \to \lvert x - \alpha \rvert < h \to \lvert y - \beta \rvert < h \to \alpha \in \mathbb{R}$
\item $\varphi_{63} :\equiv h \in \mathbb{R} \to 0 \leq M \to \lvert x \rvert \leq M \to \lvert x - \alpha \rvert < h \to \lvert y - \beta \rvert < h \to \beta \in \mathbb{R}$
\item $\varphi_{64} :\equiv \lvert x - \alpha \rvert < h \to \lvert y - \beta \rvert < h \to \lvert x \cdot y - \alpha \cdot \beta \rvert \leq (M + \lvert \beta \rvert) \cdot h$
\item $\varphi_{65} :\equiv \lvert x \rvert \leq M \to \lvert x - \alpha \rvert < h \to \lvert y - \beta \rvert < h \to \lvert x \rvert \cdot \lvert y - \beta \rvert \in \mathbb{R}$
\item $\varphi_{66} :\equiv \lvert x - \alpha \rvert < h \to \lvert y - \beta \rvert < h \to \lvert \beta \cdot (x - \alpha) \rvert = \lvert \beta \rvert \cdot \lvert x - \alpha \rvert$
\item $\varphi_{67} :\equiv h \cdot (M + \lvert \beta \rvert) = h \cdot M + h \cdot \lvert \beta \rvert \to h \cdot M + h \cdot \lvert \beta \rvert = h \cdot (M + \lvert \beta \rvert)$
\item $\varphi_{68} :\equiv \forall v_{696}\,\forall v_{697}\,(\lvert x \cdot y - \alpha \cdot \beta \rvert \leq v_{696} \to v_{696} \leq v_{697} \to \lvert x \cdot y - \alpha \cdot \beta \rvert \leq v_{697})$
\item $\varphi_{69} :\equiv \lvert x \rvert \leq M \to \lvert x - \alpha \rvert < h \to \lvert y - \beta \rvert < h \to \lvert \beta \rvert \cdot \lvert x - \alpha \rvert \in \mathbb{R}$
\item $\varphi_{70} :\equiv h \cdot M + h \cdot \lvert \beta \rvert = h \cdot (M + \lvert \beta \rvert) \to M \cdot h + \lvert \beta \rvert \cdot h = h \cdot (M + \lvert \beta \rvert)$
\item $\varphi_{71} :\equiv M \cdot h + \lvert \beta \rvert \cdot h = z + \lvert \beta \rvert \cdot h \to M \cdot h + \lvert \beta \rvert \cdot h = w + \lvert \beta \rvert \cdot h$
\item $\varphi_{72} :\equiv \lvert y - \beta \rvert < h \to \lvert x \cdot y - \alpha \cdot \beta \rvert \leq \lvert x \cdot (y - \beta) \rvert + \lvert \beta \cdot (x - \alpha) \rvert$
\item $\varphi_{73} :\equiv \forall y\,(M \in \mathbb{R} \to y \in \mathbb{R} \to M \cdot y = y \cdot M) \to M \in \mathbb{R} \to h \in \mathbb{R} \to M \cdot h = h \cdot M$
\item $\varphi_{74} :\equiv \lvert x - \alpha \rvert < h \to \lvert y - \beta \rvert < h \to h \cdot M + \lvert \beta \rvert \cdot h = h \cdot M + h \cdot \lvert \beta \rvert$
\item $\varphi_{75} :\equiv \lvert x - \alpha \rvert < h \to \lvert y - \beta \rvert < h \to M \cdot h + \lvert \beta \rvert \cdot h = h \cdot M + \lvert \beta \rvert \cdot h$
\item $\varphi_{76} :\equiv \lvert x - \alpha \rvert < h \to \lvert y - \beta \rvert < h \to M \cdot h + \lvert \beta \rvert \cdot h = h \cdot M + h \cdot \lvert \beta \rvert$
\item $\varphi_{77} :\equiv \lvert y - \beta \rvert < h \to \lvert x \rvert \cdot \lvert y - \beta \rvert + \lvert \beta \rvert \cdot \lvert x - \alpha \rvert \leq (M + \lvert \beta \rvert) \cdot h$
\item $\varphi_{78} :\equiv \forall y\,(M \in \mathbb{R} \to y \in \mathbb{R} \to M \cdot y \in \mathbb{R}) \to M \in \mathbb{R} \to h \in \mathbb{R} \to M \cdot h \in \mathbb{R}$
\item $\varphi_{79} :\equiv \lvert y - \beta \rvert < h \to \lvert x \cdot y - \alpha \cdot \beta \rvert \leq \lvert x \rvert \cdot \lvert y - \beta \rvert + \lvert \beta \cdot (x - \alpha) \rvert$
\item $\varphi_{80} :\equiv \forall x\,\forall y\,\forall u\,\forall v\,(x \in \mathbb{R} \to y \in \mathbb{R} \to u \in \mathbb{R} \to v \in \mathbb{R} \to x \leq y \to u \leq v \to x + u \leq y + v)$
\item $\varphi_{81} :\equiv \forall y\,(x \in \mathbb{R} \to y \in \mathbb{R} \to x - y \in \mathbb{R}) \to x \in \mathbb{R} \to \alpha \in \mathbb{R} \to x - \alpha \in \mathbb{R}$
\item $\varphi_{82} :\equiv \forall x\,\forall y\,(x \in \mathbb{R} \to y \in \mathbb{R} \to x \cdot y = y \cdot x) \to \forall y\,(h \in \mathbb{R} \to y \in \mathbb{R} \to h \cdot y = y \cdot h)$
\item $\varphi_{83} :\equiv \forall x\,\forall y\,(x \in \mathbb{R} \to y \in \mathbb{R} \to x \cdot y = y \cdot x) \to \forall y\,(M \in \mathbb{R} \to y \in \mathbb{R} \to M \cdot y = y \cdot M)$
\item $\varphi_{84} :\equiv \forall b\,(x \cdot (y - \beta) \in \mathbb{R} \to b \in \mathbb{R} \to \lvert x \cdot (y - \beta) + b \rvert \leq \lvert x \cdot (y - \beta) \rvert + \lvert b \rvert)$
\item $\varphi_{85} :\equiv \forall v_{687}\,\forall v_{688}\,(v_{687} = v_{688} \to v_{688} = v_{687}) \to \forall v_{688}\,(x \cdot y - \alpha \cdot \beta = v_{688} \to v_{688} = x \cdot y - \alpha \cdot \beta)$
\item $\varphi_{86} :\equiv \forall z\,(M \cdot h + \lvert \beta \rvert \cdot h = h \cdot (M + \lvert \beta \rvert) \to h \cdot (M + \lvert \beta \rvert) = z \to M \cdot h + \lvert \beta \rvert \cdot h = z)$
\item $\varphi_{87} :\equiv \lvert y - \beta \rvert < h \to \lvert x \rvert \cdot \lvert y - \beta \rvert + \lvert \beta \rvert \cdot \lvert x - \alpha \rvert \leq M \cdot h + \lvert \beta \rvert \cdot h$
\item $\varphi_{88} :\equiv h \cdot M + \lvert \beta \rvert \cdot h = h \cdot M + h \cdot \lvert \beta \rvert \to M \cdot h + \lvert \beta \rvert \cdot h = h \cdot M + h \cdot \lvert \beta \rvert$
\item $\varphi_{89} :\equiv h \cdot M + \lvert \beta \rvert \cdot h = h \cdot M + \lvert \beta \rvert \cdot h \to h \cdot M + \lvert \beta \rvert \cdot h = h \cdot M + h \cdot \lvert \beta \rvert$
\item $\varphi_{90} :\equiv M \cdot h + \lvert \beta \rvert \cdot h = M \cdot h + \lvert \beta \rvert \cdot h \to M \cdot h + \lvert \beta \rvert \cdot h = h \cdot M + \lvert \beta \rvert \cdot h$
\item $\varphi_{91} :\equiv \lvert y - \beta \rvert < h \to \lvert x \cdot y - \alpha \cdot \beta \rvert \leq \lvert x \rvert \cdot \lvert y - \beta \rvert + \lvert \beta \rvert \cdot \lvert x - \alpha \rvert$
\item $\varphi_{92} :\equiv \forall x\,\forall y\,(x \in \mathbb{R} \to y \in \mathbb{R} \to x - y \in \mathbb{R}) \to \forall y\,(x \in \mathbb{R} \to y \in \mathbb{R} \to x - y \in \mathbb{R})$
\item $\varphi_{93} :\equiv \forall x\,\forall y\,(x \in \mathbb{R} \to y \in \mathbb{R} \to x \cdot y \in \mathbb{R}) \to \forall y\,(M \in \mathbb{R} \to y \in \mathbb{R} \to M \cdot y \in \mathbb{R})$
\item $\varphi_{94} :\equiv \forall x\,\forall y\,(x \in \mathbb{R} \to y \in \mathbb{R} \to x + y \in \mathbb{R}) \to \forall y\,(M \in \mathbb{R} \to y \in \mathbb{R} \to M + y \in \mathbb{R})$
\item $\varphi_{95} :\equiv \forall x\,\forall y\,(x \in \mathbb{R} \to y \in \mathbb{R} \to x \cdot y \in \mathbb{R}) \to \forall y\,(x \in \mathbb{R} \to y \in \mathbb{R} \to x \cdot y \in \mathbb{R})$
\item $\varphi_{96} :\equiv \forall y\,\forall u\,\forall v\,(v_{691} \in \mathbb{R} \to y \in \mathbb{R} \to u \in \mathbb{R} \to v \in \mathbb{R} \to v_{691} \leq y \to u \leq v \to v_{691} + u \leq y + v)$
\item $\varphi_{97} :\equiv \lvert z \rvert = \lvert x \cdot (y - \beta) \rvert + \lvert \beta \cdot (x - \alpha) \rvert \to \lvert w \rvert = \lvert x \cdot (y - \beta) \rvert + \lvert \beta \cdot (x - \alpha) \rvert$
\item $\varphi_{98} :\equiv \forall x\,\forall y\,(x \in \mathbb{R} \to y \in \mathbb{R} \to x \cdot y \in \mathbb{R}) \to \forall y\,(\beta \in \mathbb{R} \to y \in \mathbb{R} \to \beta \cdot y \in \mathbb{R})$
\item $\varphi_{99} :\equiv \lvert x \cdot y - \alpha \cdot \beta \rvert = z + \lvert \beta \cdot (x - \alpha) \rvert \to \lvert x \cdot y - \alpha \cdot \beta \rvert = w + \lvert \beta \cdot (x - \alpha) \rvert$
\item $\varphi_{100} :\equiv \forall v_{690}\,(y \in \mathbb{R} \to v_{690} \in \mathbb{R} \to y - v_{690} \in \mathbb{R}) \to y \in \mathbb{R} \to \beta \in \mathbb{R} \to y - \beta \in \mathbb{R}$
\item $\varphi_{101} :\equiv x \cdot y - \alpha \cdot \beta = x \cdot (y - \beta) + \beta \cdot (x - \alpha) \to x \cdot (y - \beta) + \beta \cdot (x - \alpha) = x \cdot y - \alpha \cdot \beta$
\item $\varphi_{102} :\equiv \forall x\,\forall y\,(x \in \mathbb{R} \to y \in \mathbb{R} \to x - y \in \mathbb{R}) \to \forall y\,(v_{689} \in \mathbb{R} \to y \in \mathbb{R} \to v_{689} - y \in \mathbb{R})$
\item $\varphi_{103} :\equiv \forall y\,(M \in \mathbb{R} \to y \in \mathbb{R} \to M + y \in \mathbb{R}) \to M \in \mathbb{R} \to \lvert \beta \rvert \in \mathbb{R} \to M + \lvert \beta \rvert \in \mathbb{R}$
\item $\varphi_{104} :\equiv \forall y\,(x \in \mathbb{R} \to y \in \mathbb{R} \to x \cdot y \in \mathbb{R}) \to x \in \mathbb{R} \to y - \beta \in \mathbb{R} \to x \cdot (y - \beta) \in \mathbb{R}$
\item $\varphi_{105} :\equiv \forall x\,\forall y\,\forall u\,\forall v\,(x \in \mathbb{R} \to y \in \mathbb{R} \to u \in \mathbb{R} \to v \in \mathbb{R} \to 0 \leq x \to 0 \leq u \to x \leq y \to u \leq v \to x \cdot u \leq y \cdot v)$
\item $\varphi_{106} :\equiv \forall u\,\forall v\,(v_{691} \in \mathbb{R} \to v_{692} \in \mathbb{R} \to u \in \mathbb{R} \to v \in \mathbb{R} \to v_{691} \leq v_{692} \to u \leq v \to v_{691} + u \leq v_{692} + v)$
\item $\varphi_{107} :\equiv \forall a\,\forall b\,(x \in \mathbb{R} \to y \in \mathbb{R} \to a \in \mathbb{R} \to b \in \mathbb{R} \to x \cdot y - a \cdot b = x \cdot (y - b) + b \cdot (x - a))$
\item $\varphi_{108} :\equiv \forall b\,(x \in \mathbb{R} \to y \in \mathbb{R} \to \alpha \in \mathbb{R} \to b \in \mathbb{R} \to x \cdot y - \alpha \cdot b = x \cdot (y - b) + b \cdot (x - \alpha))$
\item $\varphi_{109} :\equiv \lvert z \rvert < \lvert x \cdot (y - \beta) \rvert + \lvert \beta \cdot (x - \alpha) \rvert \to \lvert w \rvert < \lvert x \cdot (y - \beta) \rvert + \lvert \beta \cdot (x - \alpha) \rvert$
\item $\varphi_{110} :\equiv \forall z\,(M \cdot h + \lvert \beta \rvert \cdot h = h \cdot M + h \cdot \lvert \beta \rvert \to h \cdot M + h \cdot \lvert \beta \rvert = z \to M \cdot h + \lvert \beta \rvert \cdot h = z)$
\item $\varphi_{111} :\equiv \forall z\,(M \cdot h + \lvert \beta \rvert \cdot h = h \cdot M + \lvert \beta \rvert \cdot h \to h \cdot M + \lvert \beta \rvert \cdot h = z \to M \cdot h + \lvert \beta \rvert \cdot h = z)$
\item $\varphi_{112} :\equiv \forall w\,(\lvert \beta \rvert \cdot h = w \to h \cdot M + \lvert \beta \rvert \cdot h = h \cdot M + \lvert \beta \rvert \cdot h \to h \cdot M + \lvert \beta \rvert \cdot h = h \cdot M + w)$
\item $\varphi_{113} :\equiv \forall w\,(M \cdot h = w \to M \cdot h + \lvert \beta \rvert \cdot h = M \cdot h + \lvert \beta \rvert \cdot h \to M \cdot h + \lvert \beta \rvert \cdot h = w + \lvert \beta \rvert \cdot h)$
\item $\varphi_{114} :\equiv \lvert x \cdot y - \alpha \cdot \beta \rvert < z + \lvert \beta \cdot (x - \alpha) \rvert \to \lvert x \cdot y - \alpha \cdot \beta \rvert < w + \lvert \beta \cdot (x - \alpha) \rvert$
\item $\varphi_{115} :\equiv \forall x\,\forall y\,(x \in \mathbb{R} \to y \in \mathbb{R} \to x \cdot y \in \mathbb{R}) \to \forall y\,(\lvert x \rvert \in \mathbb{R} \to y \in \mathbb{R} \to \lvert x \rvert \cdot y \in \mathbb{R})$
\item $\varphi_{116} :\equiv \forall y\,(\beta \in \mathbb{R} \to y \in \mathbb{R} \to \beta \cdot y \in \mathbb{R}) \to \beta \in \mathbb{R} \to x - \alpha \in \mathbb{R} \to \beta \cdot (x - \alpha) \in \mathbb{R}$
\item $\varphi_{117} :\equiv \lvert x \rvert \cdot \lvert y - \beta \rvert + \lvert \beta \rvert \cdot \lvert x - \alpha \rvert = z \to \lvert x \rvert \cdot \lvert y - \beta \rvert + \lvert \beta \rvert \cdot \lvert x - \alpha \rvert = w$
\item $\varphi_{118} :\equiv \lvert x \cdot y - \alpha \cdot \beta \rvert = \lvert x \rvert \cdot \lvert y - \beta \rvert + z \to \lvert x \cdot y - \alpha \cdot \beta \rvert = \lvert x \rvert \cdot \lvert y - \beta \rvert + w$
\item $\varphi_{119} :\equiv \forall y\,(v_{689} \in \mathbb{R} \to y \in \mathbb{R} \to v_{689} - y \in \mathbb{R}) \to v_{689} \in \mathbb{R} \to v_{690} \in \mathbb{R} \to v_{689} - v_{690} \in \mathbb{R}$
\item $\varphi_{120} :\equiv \forall x\,\forall y\,(x \in \mathbb{R} \to y \in \mathbb{R} \to x \cdot y \in \mathbb{R}) \to \forall y\,(\lvert \beta \rvert \in \mathbb{R} \to y \in \mathbb{R} \to \lvert \beta \rvert \cdot y \in \mathbb{R})$
\item $\varphi_{121} :\equiv \forall v\,(v_{691} \in \mathbb{R} \to v_{692} \in \mathbb{R} \to v_{693} \in \mathbb{R} \to v \in \mathbb{R} \to v_{691} \leq v_{692} \to v_{693} \leq v \to v_{691} + v_{693} \leq v_{692} + v)$
\item $\varphi_{122} :\equiv \forall x\,\forall y\,(x \in \mathbb{R} \to y \in \mathbb{R} \to x \cdot y = y \cdot x) \to \forall y\,(\lvert \beta \rvert \in \mathbb{R} \to y \in \mathbb{R} \to \lvert \beta \rvert \cdot y = y \cdot \lvert \beta \rvert)$
\item $\varphi_{123} :\equiv \lvert x \rvert \cdot \lvert y - \beta \rvert + \lvert \beta \rvert \cdot \lvert x - \alpha \rvert < z \to \lvert x \rvert \cdot \lvert y - \beta \rvert + \lvert \beta \rvert \cdot \lvert x - \alpha \rvert < w$
\item $\varphi_{124} :\equiv \lvert \beta \rvert \cdot \lvert x - \alpha \rvert \leq \lvert \beta \rvert \cdot h \to \lvert x \rvert \cdot \lvert y - \beta \rvert + \lvert \beta \rvert \cdot \lvert x - \alpha \rvert \leq M \cdot h + \lvert \beta \rvert \cdot h$
\item $\varphi_{125} :\equiv \lvert x \cdot y - \alpha \cdot \beta \rvert < \lvert x \rvert \cdot \lvert y - \beta \rvert + z \to \lvert x \cdot y - \alpha \cdot \beta \rvert < \lvert x \rvert \cdot \lvert y - \beta \rvert + w$
\item $\varphi_{126} :\equiv \lvert x \rvert \cdot \lvert y - \beta \rvert + \lvert \beta \rvert \cdot \lvert x - \alpha \rvert \leq (M + \lvert \beta \rvert) \cdot h \to \lvert x \cdot y - \alpha \cdot \beta \rvert \leq (M + \lvert \beta \rvert) \cdot h$
\item $\varphi_{127} :\equiv \forall y\,(\lvert \beta \rvert \in \mathbb{R} \to y \in \mathbb{R} \to \lvert \beta \rvert \cdot y \in \mathbb{R}) \to \lvert \beta \rvert \in \mathbb{R} \to h \in \mathbb{R} \to \lvert \beta \rvert \cdot h \in \mathbb{R}$
\item $\varphi_{128} :\equiv \forall u\,\forall v\,(\lvert x \rvert \in \mathbb{R} \to M \in \mathbb{R} \to u \in \mathbb{R} \to v \in \mathbb{R} \to 0 \leq \lvert x \rvert \to 0 \leq u \to \lvert x \rvert \leq M \to u \leq v \to \lvert x \rvert \cdot u \leq M \cdot v)$
\item $\varphi_{129} :\equiv \forall y\,\forall u\,\forall v\,(\lvert x \rvert \in \mathbb{R} \to y \in \mathbb{R} \to u \in \mathbb{R} \to v \in \mathbb{R} \to 0 \leq \lvert x \rvert \to 0 \leq u \to \lvert x \rvert \leq y \to u \leq v \to \lvert x \rvert \cdot u \leq y \cdot v)$
\item $\varphi_{130} :\equiv \forall v_{694}\,(v_{691} \in \mathbb{R} \to v_{692} \in \mathbb{R} \to v_{693} \in \mathbb{R} \to v_{694} \in \mathbb{R} \to \varphi_{3})$
\item $\varphi_{131} :\equiv \forall y\,(\lvert \beta \rvert \in \mathbb{R} \to y \in \mathbb{R} \to \lvert \beta \rvert \cdot y = y \cdot \lvert \beta \rvert) \to \lvert \beta \rvert \in \mathbb{R} \to h \in \mathbb{R} \to \lvert \beta \rvert \cdot h = h \cdot \lvert \beta \rvert$
\item $\varphi_{132} :\equiv \forall x\,\forall y\,(x \in \mathbb{R} \to y \in \mathbb{R} \to \lvert x \cdot y \rvert = \lvert x \rvert \cdot \lvert y \rvert) \to \forall y\,(x \in \mathbb{R} \to y \in \mathbb{R} \to \lvert x \cdot y \rvert = \lvert x \rvert \cdot \lvert y \rvert)$
\item $\varphi_{133} :\equiv \forall y\,\forall u\,\forall v\,(\lvert \beta \rvert \in \mathbb{R} \to y \in \mathbb{R} \to u \in \mathbb{R} \to v \in \mathbb{R} \to 0 \leq \lvert \beta \rvert \to 0 \leq u \to \lvert \beta \rvert \leq y \to u \leq v \to \lvert \beta \rvert \cdot u \leq y \cdot v)$
\item $\varphi_{134} :\equiv \forall x\,\forall y\,(x \in \mathbb{R} \to y \in \mathbb{R} \to \lvert x \cdot y \rvert = \lvert x \rvert \cdot \lvert y \rvert) \to \forall y\,(\beta \in \mathbb{R} \to y \in \mathbb{R} \to \lvert \beta \cdot y \rvert = \lvert \beta \rvert \cdot \lvert y \rvert)$
\item $\varphi_{135} :\equiv \forall y\,(\lvert x \rvert \in \mathbb{R} \to y \in \mathbb{R} \to \lvert x \rvert \cdot y \in \mathbb{R}) \to \lvert x \rvert \in \mathbb{R} \to \lvert y - \beta \rvert \in \mathbb{R} \to \lvert x \rvert \cdot \lvert y - \beta \rvert \in \mathbb{R}$
\item $\varphi_{136} :\equiv \forall y\,(\lvert \beta \rvert \in \mathbb{R} \to y \in \mathbb{R} \to \lvert \beta \rvert \cdot y \in \mathbb{R}) \to \lvert \beta \rvert \in \mathbb{R} \to \lvert x - \alpha \rvert \in \mathbb{R} \to \lvert \beta \rvert \cdot \lvert x - \alpha \rvert \in \mathbb{R}$
\item $\varphi_{137} :\equiv \lvert x \cdot y - \alpha \cdot \beta \rvert = \lvert x \cdot (y - \beta) \rvert + \lvert \beta \cdot (x - \alpha) \rvert \to \lvert x \cdot y - \alpha \cdot \beta \rvert = \lvert x \rvert \cdot \lvert y - \beta \rvert + \lvert \beta \cdot (x - \alpha) \rvert$
\item $\varphi_{138} :\equiv \forall u\,\forall v\,(\lvert \beta \rvert \in \mathbb{R} \to \lvert \beta \rvert \in \mathbb{R} \to u \in \mathbb{R} \to v \in \mathbb{R} \to 0 \leq \lvert \beta \rvert \to 0 \leq u \to \lvert \beta \rvert \leq \lvert \beta \rvert \to u \leq v \to \lvert \beta \rvert \cdot u \leq \lvert \beta \rvert \cdot v)$
\item $\varphi_{139} :\equiv \forall w\,(\lvert x \cdot (y - \beta) \rvert = w \to \lvert x \cdot y - \alpha \cdot \beta \rvert = \lvert x \cdot (y - \beta) \rvert + \lvert \beta \cdot (x - \alpha) \rvert \to \lvert x \cdot y - \alpha \cdot \beta \rvert = w + \lvert \beta \cdot (x - \alpha) \rvert)$
\item $\varphi_{140} :\equiv \lvert x \cdot y - \alpha \cdot \beta \rvert < \lvert x \cdot (y - \beta) \rvert + \lvert \beta \cdot (x - \alpha) \rvert \to \lvert x \cdot y - \alpha \cdot \beta \rvert < \lvert x \rvert \cdot \lvert y - \beta \rvert + \lvert \beta \cdot (x - \alpha) \rvert$
\item $\varphi_{141} :\equiv \varphi_{38} \to \lvert x \cdot y - \alpha \cdot \beta \rvert = \lvert x \cdot (y - \beta) \rvert + \lvert \beta \cdot (x - \alpha) \rvert$
\item $\varphi_{142} :\equiv \lvert x \rvert \cdot \lvert y - \beta \rvert + \lvert \beta \rvert \cdot \lvert x - \alpha \rvert = M \cdot h + \lvert \beta \rvert \cdot h \to \lvert x \rvert \cdot \lvert y - \beta \rvert + \lvert \beta \rvert \cdot \lvert x - \alpha \rvert = (M + \lvert \beta \rvert) \cdot h$
\item $\varphi_{143} :\equiv \forall w\,(\lvert x \cdot (y - \beta) \rvert = w \to \lvert x \cdot y - \alpha \cdot \beta \rvert < \lvert x \cdot (y - \beta) \rvert + \lvert \beta \cdot (x - \alpha) \rvert \to \lvert x \cdot y - \alpha \cdot \beta \rvert < w + \lvert \beta \cdot (x - \alpha) \rvert)$
\item $\varphi_{144} :\equiv \lvert x \cdot y - \alpha \cdot \beta \rvert = \lvert x \rvert \cdot \lvert y - \beta \rvert + \lvert \beta \cdot (x - \alpha) \rvert \to \lvert x \cdot y - \alpha \cdot \beta \rvert = \lvert x \rvert \cdot \lvert y - \beta \rvert + \lvert \beta \rvert \cdot \lvert x - \alpha \rvert$
\item $\varphi_{145} :\equiv h \in \mathbb{R} \to 0 \leq M \to \lvert x \rvert \leq M \to \lvert x - \alpha \rvert < h \to \lvert y - \beta \rvert < h \to \varphi_{39}$
\item $\varphi_{146} :\equiv x \in \mathbb{R} \to y \in \mathbb{R} \to \alpha \in \mathbb{R} \to \beta \in \mathbb{R} \to M \in \mathbb{R} \to h \in \mathbb{R} \to \varphi_{9}$
\item $\varphi_{147} :\equiv x \in \mathbb{R} \to y \in \mathbb{R} \to \alpha \in \mathbb{R} \to \beta \in \mathbb{R} \to M \in \mathbb{R} \to h \in \mathbb{R} \to \varphi_{10}$
\item $\varphi_{148} :\equiv \beta \in \mathbb{R} \to M \in \mathbb{R} \to h \in \mathbb{R} \to 0 \leq M \to \lvert x \rvert \leq M \to \lvert x - \alpha \rvert < h \to \varphi_{13}$
\item $\varphi_{149} :\equiv \beta \in \mathbb{R} \to M \in \mathbb{R} \to h \in \mathbb{R} \to 0 \leq M \to \lvert x \rvert \leq M \to \lvert x - \alpha \rvert < h \to \varphi_{14}$
\item $\varphi_{150} :\equiv x \in \mathbb{R} \to y \in \mathbb{R} \to \alpha \in \mathbb{R} \to \beta \in \mathbb{R} \to M \in \mathbb{R} \to h \in \mathbb{R} \to \varphi_{12}$
\item $\varphi_{151} :\equiv \alpha \in \mathbb{R} \to \beta \in \mathbb{R} \to M \in \mathbb{R} \to h \in \mathbb{R} \to 0 \leq M \to \lvert x \rvert \leq M \to \varphi_{41}$
\item $\varphi_{152} :\equiv \varphi_{46} \to \lvert x \cdot y - \alpha \cdot \beta \rvert < \lvert x \cdot (y - \beta) \rvert + \lvert \beta \cdot (x - \alpha) \rvert$
\item $\varphi_{153} :\equiv \alpha \in \mathbb{R} \to \beta \in \mathbb{R} \to M \in \mathbb{R} \to h \in \mathbb{R} \to 0 \leq M \to \lvert x \rvert \leq M \to \varphi_{43}$
\item $\varphi_{154} :\equiv \forall w\,(\lvert \beta \cdot (x - \alpha) \rvert = w \to \lvert x \cdot y - \alpha \cdot \beta \rvert = \lvert x \rvert \cdot \lvert y - \beta \rvert + \lvert \beta \cdot (x - \alpha) \rvert \to \lvert x \cdot y - \alpha \cdot \beta \rvert = \lvert x \rvert \cdot \lvert y - \beta \rvert + w)$
\item $\varphi_{155} :\equiv \lvert x \rvert \cdot \lvert y - \beta \rvert + \lvert \beta \rvert \cdot \lvert x - \alpha \rvert < M \cdot h + \lvert \beta \rvert \cdot h \to \lvert x \rvert \cdot \lvert y - \beta \rvert + \lvert \beta \rvert \cdot \lvert x - \alpha \rvert < (M + \lvert \beta \rvert) \cdot h$
\item $\varphi_{156} :\equiv x \in \mathbb{R} \to y \in \mathbb{R} \to \alpha \in \mathbb{R} \to \beta \in \mathbb{R} \to M \in \mathbb{R} \to h \in \mathbb{R} \to \varphi_{16}$
\item $\varphi_{157} :\equiv \lvert x \cdot y - \alpha \cdot \beta \rvert < \lvert x \rvert \cdot \lvert y - \beta \rvert + \lvert \beta \cdot (x - \alpha) \rvert \to \lvert x \cdot y - \alpha \cdot \beta \rvert < \lvert x \rvert \cdot \lvert y - \beta \rvert + \lvert \beta \rvert \cdot \lvert x - \alpha \rvert$
\item $\varphi_{158} :\equiv x \in \mathbb{R} \to y \in \mathbb{R} \to \alpha \in \mathbb{R} \to \beta \in \mathbb{R} \to M \in \mathbb{R} \to h \in \mathbb{R} \to \varphi_{17}$
\item $\varphi_{159} :\equiv x \in \mathbb{R} \to y \in \mathbb{R} \to \alpha \in \mathbb{R} \to \beta \in \mathbb{R} \to M \in \mathbb{R} \to h \in \mathbb{R} \to \varphi_{18}$
\item $\varphi_{160} :\equiv x \in \mathbb{R} \to y \in \mathbb{R} \to \alpha \in \mathbb{R} \to \beta \in \mathbb{R} \to M \in \mathbb{R} \to h \in \mathbb{R} \to \varphi_{20}$
\item $\varphi_{161} :\equiv x \in \mathbb{R} \to y \in \mathbb{R} \to \alpha \in \mathbb{R} \to \beta \in \mathbb{R} \to M \in \mathbb{R} \to h \in \mathbb{R} \to \varphi_{21}$
\item $\varphi_{162} :\equiv \alpha \in \mathbb{R} \to \beta \in \mathbb{R} \to M \in \mathbb{R} \to h \in \mathbb{R} \to 0 \leq M \to \lvert x \rvert \leq M \to \varphi_{57}$
\item $\varphi_{163} :\equiv \alpha \in \mathbb{R} \to \beta \in \mathbb{R} \to M \in \mathbb{R} \to h \in \mathbb{R} \to 0 \leq M \to \lvert x \rvert \leq M \to \varphi_{58}$
\item $\varphi_{164} :\equiv \alpha \in \mathbb{R} \to \beta \in \mathbb{R} \to M \in \mathbb{R} \to h \in \mathbb{R} \to 0 \leq M \to \lvert x \rvert \leq M \to \varphi_{59}$
\item $\varphi_{165} :\equiv \alpha \in \mathbb{R} \to \beta \in \mathbb{R} \to M \in \mathbb{R} \to h \in \mathbb{R} \to 0 \leq M \to \lvert x \rvert \leq M \to \varphi_{60}$
\item $\varphi_{166} :\equiv \alpha \in \mathbb{R} \to \beta \in \mathbb{R} \to M \in \mathbb{R} \to h \in \mathbb{R} \to 0 \leq M \to \lvert x \rvert \leq M \to \varphi_{61}$
\item $\varphi_{167} :\equiv x \in \mathbb{R} \to y \in \mathbb{R} \to \alpha \in \mathbb{R} \to \beta \in \mathbb{R} \to M \in \mathbb{R} \to h \in \mathbb{R} \to \varphi_{25}$
\item $\varphi_{168} :\equiv x \in \mathbb{R} \to y \in \mathbb{R} \to \alpha \in \mathbb{R} \to \beta \in \mathbb{R} \to M \in \mathbb{R} \to h \in \mathbb{R} \to \varphi_{26}$
\item $\varphi_{169} :\equiv x \in \mathbb{R} \to y \in \mathbb{R} \to \alpha \in \mathbb{R} \to \beta \in \mathbb{R} \to M \in \mathbb{R} \to h \in \mathbb{R} \to \varphi_{27}$
\item $\varphi_{170} :\equiv \alpha \in \mathbb{R} \to \beta \in \mathbb{R} \to M \in \mathbb{R} \to h \in \mathbb{R} \to 0 \leq M \to \lvert x \rvert \leq M \to \varphi_{64}$
\item $\varphi_{171} :\equiv x \in \mathbb{R} \to y \in \mathbb{R} \to \alpha \in \mathbb{R} \to \beta \in \mathbb{R} \to M \in \mathbb{R} \to h \in \mathbb{R} \to \varphi_{31}$
\item $\varphi_{172} :\equiv x \in \mathbb{R} \to y \in \mathbb{R} \to \alpha \in \mathbb{R} \to \beta \in \mathbb{R} \to M \in \mathbb{R} \to h \in \mathbb{R} \to \varphi_{32}$
\item $\varphi_{173} :\equiv x \in \mathbb{R} \to y \in \mathbb{R} \to \alpha \in \mathbb{R} \to \beta \in \mathbb{R} \to M \in \mathbb{R} \to h \in \mathbb{R} \to \varphi_{33}$
\item $\varphi_{174} :\equiv \lvert x \rvert \in \mathbb{R} \to M \in \mathbb{R} \to \lvert y - \beta \rvert \in \mathbb{R} \to h \in \mathbb{R} \to 0 \leq \lvert x \rvert \to \varphi_{37}$
\item $\varphi_{175} :\equiv \forall a\,\forall b\,(a \in \mathbb{R} \to b \in \mathbb{R} \to \lvert a + b \rvert \leq \lvert a \rvert + \lvert b \rvert) \to \varphi_{84}$
\item $\varphi_{176} :\equiv \alpha \in \mathbb{R} \to \beta \in \mathbb{R} \to M \in \mathbb{R} \to h \in \mathbb{R} \to 0 \leq M \to \lvert x \rvert \leq M \to \varphi_{66}$
\item $\varphi_{177} :\equiv x \in \mathbb{R} \to y \in \mathbb{R} \to \alpha \in \mathbb{R} \to \beta \in \mathbb{R} \to M \in \mathbb{R} \to h \in \mathbb{R} \to \varphi_{35}$
\item $\varphi_{178} :\equiv x \in \mathbb{R} \to y \in \mathbb{R} \to \alpha \in \mathbb{R} \to \beta \in \mathbb{R} \to M \in \mathbb{R} \to h \in \mathbb{R} \to \varphi_{36}$
\item $\varphi_{179} :\equiv \forall w\,(\lvert \beta \cdot (x - \alpha) \rvert = w \to \lvert x \cdot y - \alpha \cdot \beta \rvert < \lvert x \rvert \cdot \lvert y - \beta \rvert + \lvert \beta \cdot (x - \alpha) \rvert \to \lvert x \cdot y - \alpha \cdot \beta \rvert < \lvert x \rvert \cdot \lvert y - \beta \rvert + w)$
\item $\varphi_{180} :\equiv \forall w\,(M \cdot h + \lvert \beta \rvert \cdot h = w \to \lvert x \rvert \cdot \lvert y - \beta \rvert + \lvert \beta \rvert \cdot \lvert x - \alpha \rvert = M \cdot h + \lvert \beta \rvert \cdot h \to \lvert x \rvert \cdot \lvert y - \beta \rvert + \lvert \beta \rvert \cdot \lvert x - \alpha \rvert = w)$
\item $\varphi_{181} :\equiv \forall v\,(\lvert x \rvert \in \mathbb{R} \to M \in \mathbb{R} \to \lvert y - \beta \rvert \in \mathbb{R} \to v \in \mathbb{R} \to 0 \leq \lvert x \rvert \to 0 \leq \lvert y - \beta \rvert \to \lvert x \rvert \leq M \to \lvert y - \beta \rvert \leq v \to \lvert x \rvert \cdot \lvert y - \beta \rvert \leq M \cdot v)$
\item $\varphi_{182} :\equiv \lvert \beta \rvert \in \mathbb{R} \to \lvert x - \alpha \rvert \in \mathbb{R} \to h \in \mathbb{R} \to 0 \leq \lvert \beta \rvert \to 0 \leq \lvert x - \alpha \rvert \to \varphi_{19}$
\item $\varphi_{183} :\equiv \alpha \in \mathbb{R} \to \beta \in \mathbb{R} \to M \in \mathbb{R} \to h \in \mathbb{R} \to 0 \leq M \to \lvert x \rvert \leq M \to \varphi_{74}$
\item $\varphi_{184} :\equiv \alpha \in \mathbb{R} \to \beta \in \mathbb{R} \to M \in \mathbb{R} \to h \in \mathbb{R} \to 0 \leq M \to \lvert x \rvert \leq M \to \varphi_{75}$
\item $\varphi_{185} :\equiv \alpha \in \mathbb{R} \to \beta \in \mathbb{R} \to M \in \mathbb{R} \to h \in \mathbb{R} \to 0 \leq M \to \lvert x \rvert \leq M \to \varphi_{76}$
\item $\varphi_{186} :\equiv \forall v_{688}\,(x \cdot y - \alpha \cdot \beta = v_{688} \to v_{688} = x \cdot y - \alpha \cdot \beta) \to \varphi_{101}$
\item $\varphi_{187} :\equiv x \in \mathbb{R} \to y \in \mathbb{R} \to \alpha \in \mathbb{R} \to \beta \in \mathbb{R} \to M \in \mathbb{R} \to h \in \mathbb{R} \to \varphi_{47}$
\item $\varphi_{188} :\equiv x \in \mathbb{R} \to y \in \mathbb{R} \to \alpha \in \mathbb{R} \to \beta \in \mathbb{R} \to M \in \mathbb{R} \to h \in \mathbb{R} \to \varphi_{48}$
\item $\varphi_{189} :\equiv x \in \mathbb{R} \to y \in \mathbb{R} \to \alpha \in \mathbb{R} \to \beta \in \mathbb{R} \to M \in \mathbb{R} \to h \in \mathbb{R} \to \varphi_{49}$
\item $\varphi_{190} :\equiv x \in \mathbb{R} \to y \in \mathbb{R} \to \alpha \in \mathbb{R} \to \beta \in \mathbb{R} \to M \in \mathbb{R} \to h \in \mathbb{R} \to \varphi_{51}$
\item $\varphi_{191} :\equiv x \in \mathbb{R} \to y \in \mathbb{R} \to \alpha \in \mathbb{R} \to \beta \in \mathbb{R} \to M \in \mathbb{R} \to h \in \mathbb{R} \to 0 \leq M \to \varphi_{29}$
\item $\varphi_{192} :\equiv x \in \mathbb{R} \to y \in \mathbb{R} \to \alpha \in \mathbb{R} \to \beta \in \mathbb{R} \to M \in \mathbb{R} \to h \in \mathbb{R} \to 0 \leq M \to \varphi_{34}$
\item $\varphi_{193} :\equiv \forall w\,(x \cdot (y - \beta) + \beta \cdot (x - \alpha) = w \to \varphi_{38} \to \lvert w \rvert = \lvert x \cdot (y - \beta) \rvert + \lvert \beta \cdot (x - \alpha) \rvert)$
\item $\varphi_{194} :\equiv \beta \in \mathbb{R} \to M \in \mathbb{R} \to h \in \mathbb{R} \to 0 \leq M \to \lvert x \rvert \leq M \to \lvert x - \alpha \rvert < h \to \varphi_{72}$
\item $\varphi_{195} :\equiv \forall w\,(M \cdot h + \lvert \beta \rvert \cdot h = w \to \lvert x \rvert \cdot \lvert y - \beta \rvert + \lvert \beta \rvert \cdot \lvert x - \alpha \rvert < M \cdot h + \lvert \beta \rvert \cdot h \to \lvert x \rvert \cdot \lvert y - \beta \rvert + \lvert \beta \rvert \cdot \lvert x - \alpha \rvert < w)$
\item $\varphi_{196} :\equiv x \in \mathbb{R} \to y \in \mathbb{R} \to \alpha \in \mathbb{R} \to \beta \in \mathbb{R} \to M \in \mathbb{R} \to h \in \mathbb{R} \to 0 \leq M \to \varphi_{44}$
\item $\varphi_{197} :\equiv \beta \in \mathbb{R} \to M \in \mathbb{R} \to h \in \mathbb{R} \to 0 \leq M \to \lvert x \rvert \leq M \to \lvert x - \alpha \rvert < h \to \varphi_{77}$
\item $\varphi_{198} :\equiv \beta \in \mathbb{R} \to M \in \mathbb{R} \to h \in \mathbb{R} \to 0 \leq M \to \lvert x \rvert \leq M \to \lvert x - \alpha \rvert < h \to \varphi_{79}$
\item $\varphi_{199} :\equiv x \in \mathbb{R} \to y \in \mathbb{R} \to \alpha \in \mathbb{R} \to \beta \in \mathbb{R} \to M \in \mathbb{R} \to h \in \mathbb{R} \to 0 \leq M \to \varphi_{56}$
\item $\varphi_{200} :\equiv \forall w\,(x \cdot (y - \beta) + \beta \cdot (x - \alpha) = w \to \varphi_{46} \to \lvert w \rvert < \lvert x \cdot (y - \beta) \rvert + \lvert \beta \cdot (x - \alpha) \rvert)$
\item $\varphi_{201} :\equiv \forall v_{697}\,(\lvert x \cdot y - \alpha \cdot \beta \rvert \leq \lvert x \rvert \cdot \lvert y - \beta \rvert + \lvert \beta \rvert \cdot \lvert x - \alpha \rvert \to \lvert x \rvert \cdot \lvert y - \beta \rvert + \lvert \beta \rvert \cdot \lvert x - \alpha \rvert \leq v_{697} \to \lvert x \cdot y - \alpha \cdot \beta \rvert \leq v_{697})$
\item $\varphi_{202} :\equiv x \in \mathbb{R} \to y \in \mathbb{R} \to \alpha \in \mathbb{R} \to \beta \in \mathbb{R} \to M \in \mathbb{R} \to h \in \mathbb{R} \to 0 \leq M \to \varphi_{65}$
\item $\varphi_{203} :\equiv x \in \mathbb{R} \to y \in \mathbb{R} \to \alpha \in \mathbb{R} \to \beta \in \mathbb{R} \to M \in \mathbb{R} \to h \in \mathbb{R} \to 0 \leq M \to \varphi_{69}$
\item $\varphi_{204} :\equiv \beta \in \mathbb{R} \to M \in \mathbb{R} \to h \in \mathbb{R} \to 0 \leq M \to \lvert x \rvert \leq M \to \lvert x - \alpha \rvert < h \to \varphi_{87}$
\item $\varphi_{205} :\equiv \beta \in \mathbb{R} \to M \in \mathbb{R} \to h \in \mathbb{R} \to 0 \leq M \to \lvert x \rvert \leq M \to \lvert x - \alpha \rvert < h \to \varphi_{91}$
\item $\varphi_{206} :\equiv \forall v\,(\lvert \beta \rvert \in \mathbb{R} \to \lvert \beta \rvert \in \mathbb{R} \to \lvert x - \alpha \rvert \in \mathbb{R} \to v \in \mathbb{R} \to 0 \leq \lvert \beta \rvert \to 0 \leq \lvert x - \alpha \rvert \to \lvert \beta \rvert \leq \lvert \beta \rvert \to \lvert x - \alpha \rvert \leq v \to \lvert \beta \rvert \cdot \lvert x - \alpha \rvert \leq \lvert \beta \rvert \cdot v)$
\item $\varphi_{207} :\equiv \forall v_{692}\,\forall v_{693}\,\forall v_{694}\,(\lvert x \rvert \cdot \lvert y - \beta \rvert \in \mathbb{R} \to v_{692} \in \mathbb{R} \to v_{693} \in \mathbb{R} \to v_{694} \in \mathbb{R} \to \lvert x \rvert \cdot \lvert y - \beta \rvert \leq v_{692} \to v_{693} \leq v_{694} \to \lvert x \rvert \cdot \lvert y - \beta \rvert + v_{693} \leq v_{692} + v_{694})$
\item $\varphi_{208} :\equiv \forall v_{693}\,\forall v_{694}\,(\lvert x \rvert \cdot \lvert y - \beta \rvert \in \mathbb{R} \to M \cdot h \in \mathbb{R} \to v_{693} \in \mathbb{R} \to v_{694} \in \mathbb{R} \to \lvert x \rvert \cdot \lvert y - \beta \rvert \leq M \cdot h \to v_{693} \leq v_{694} \to \lvert x \rvert \cdot \lvert y - \beta \rvert + v_{693} \leq M \cdot h + v_{694})$
\item $\varphi_{209} :\equiv \lvert x \cdot y - \alpha \cdot \beta \rvert \leq \lvert x \rvert \cdot \lvert y - \beta \rvert + \lvert \beta \rvert \cdot \lvert x - \alpha \rvert \to \varphi_{126}$
\item $\varphi_{210} :\equiv \lvert \beta \rvert \cdot \lvert x - \alpha \rvert \in \mathbb{R} \to \lvert \beta \rvert \cdot h \in \mathbb{R} \to \lvert x \rvert \cdot \lvert y - \beta \rvert \leq M \cdot h \to \varphi_{124}$
\item $\varphi_{211} :\equiv \varphi_{121} \to v_{691} \in \mathbb{R} \to v_{692} \in \mathbb{R} \to v_{693} \in \mathbb{R} \to v_{694} \in \mathbb{R} \to \varphi_{3}$
\item $\varphi_{212} :\equiv \forall v_{694}\,(\lvert x \rvert \cdot \lvert y - \beta \rvert \in \mathbb{R} \to M \cdot h \in \mathbb{R} \to \lvert \beta \rvert \cdot \lvert x - \alpha \rvert \in \mathbb{R} \to v_{694} \in \mathbb{R} \to \lvert x \rvert \cdot \lvert y - \beta \rvert \leq M \cdot h \to \lvert \beta \rvert \cdot \lvert x - \alpha \rvert \leq v_{694} \to \lvert x \rvert \cdot \lvert y - \beta \rvert + \lvert \beta \rvert \cdot \lvert x - \alpha \rvert \leq M \cdot h + v_{694})$
\item $\varphi_{213} :\equiv (M \cdot h + \lvert \beta \rvert \cdot h = h \cdot (M + \lvert \beta \rvert) \to \varphi_{50}) \to x \in \mathbb{R} \to y \in \mathbb{R} \to \varphi_{165}$
\item $\varphi_{214} :\equiv (M \cdot h + \lvert \beta \rvert \cdot h = h \cdot M + h \cdot \lvert \beta \rvert \to \varphi_{70}) \to x \in \mathbb{R} \to y \in \mathbb{R} \to \varphi_{164}$
\item $\varphi_{215} :\equiv (x \in \mathbb{R} \to 0 \leq \lvert x \rvert) \to (x \in \mathbb{R} \to y \in \mathbb{R} \to \alpha \in \mathbb{R} \to \beta \in \mathbb{R} \to M \in \mathbb{R} \to \varphi_{52}) \to \varphi_{147}$
\item $\varphi_{216} :\equiv (\beta \in \mathbb{R} \to 0 \leq \lvert \beta \rvert) \to (x \in \mathbb{R} \to y \in \mathbb{R} \to \alpha \in \mathbb{R} \to \beta \in \mathbb{R} \to M \in \mathbb{R} \to \varphi_{63}) \to \varphi_{150}$
\item $\varphi_{217} :\equiv (M \cdot h + \lvert \beta \rvert \cdot h = h \cdot M + \lvert \beta \rvert \cdot h \to \varphi_{88}) \to x \in \mathbb{R} \to y \in \mathbb{R} \to \varphi_{185}$
\item $\varphi_{218} :\equiv (x \in \mathbb{R} \to \lvert x \rvert \in \mathbb{R}) \to (x \in \mathbb{R} \to y \in \mathbb{R} \to \alpha \in \mathbb{R} \to \beta \in \mathbb{R} \to M \in \mathbb{R} \to \varphi_{52}) \to \varphi_{168}$
\item $\varphi_{219} :\equiv (x \in \mathbb{R} \to y \in \mathbb{R} \to \alpha \in \mathbb{R} \to \beta \in \mathbb{R} \to M \in \mathbb{R} \to \varphi_{63}) \to x \in \mathbb{R} \to y \in \mathbb{R} \to \alpha \in \mathbb{R} \to \varphi_{148}$
\item $\varphi_{220} :\equiv (\beta \in \mathbb{R} \to \lvert \beta \rvert \in \mathbb{R}) \to (x \in \mathbb{R} \to y \in \mathbb{R} \to \alpha \in \mathbb{R} \to \beta \in \mathbb{R} \to M \in \mathbb{R} \to \varphi_{63}) \to \varphi_{172}$
\item $\varphi_{221} :\equiv M \cdot h + \lvert \beta \rvert \cdot h = M \cdot h + \lvert \beta \rvert \cdot h \to (M \cdot h = h \cdot M \to \varphi_{90}) \to x \in \mathbb{R} \to y \in \mathbb{R} \to \varphi_{184}$
\item $\varphi_{222} :\equiv h \cdot M + \lvert \beta \rvert \cdot h = h \cdot M + \lvert \beta \rvert \cdot h \to (\lvert \beta \rvert \cdot h = h \cdot \lvert \beta \rvert \to \varphi_{89}) \to x \in \mathbb{R} \to y \in \mathbb{R} \to \varphi_{183}$
\item $\varphi_{223} :\equiv (\lvert x \cdot (y - \beta) \rvert = \lvert x \rvert \cdot \lvert y - \beta \rvert \to \varphi_{137}) \to x \in \mathbb{R} \to y \in \mathbb{R} \to \alpha \in \mathbb{R} \to \varphi_{198}$
\item $\varphi_{224} :\equiv (M \cdot h + \lvert \beta \rvert \cdot h = (M + \lvert \beta \rvert) \cdot h \to \varphi_{142}) \to x \in \mathbb{R} \to y \in \mathbb{R} \to \alpha \in \mathbb{R} \to \varphi_{197}$
\item $\varphi_{225} :\equiv (x \cdot (y - \beta) + \beta \cdot (x - \alpha) = x \cdot y - \alpha \cdot \beta \to \varphi_{141}) \to x \in \mathbb{R} \to y \in \mathbb{R} \to \alpha \in \mathbb{R} \to \varphi_{194}$
\item $\varphi_{226} :\equiv (\lvert \beta \cdot (x - \alpha) \rvert = \lvert \beta \rvert \cdot \lvert x - \alpha \rvert \to \varphi_{144}) \to x \in \mathbb{R} \to y \in \mathbb{R} \to \alpha \in \mathbb{R} \to \varphi_{205}$
\item $\varphi_{227} :\equiv (x \in \mathbb{R} \to y \in \mathbb{R} \to \alpha \in \mathbb{R} \to \beta \in \mathbb{R} \to M \in \mathbb{R} \to \varphi_{55}) \to (x \in \mathbb{R} \to y \in \mathbb{R} \to \alpha \in \mathbb{R} \to \beta \in \mathbb{R} \to M \in \mathbb{R} \to \varphi_{53}) \to \varphi_{156}$
\item $\varphi_{228} :\equiv (x \in \mathbb{R} \to y \in \mathbb{R} \to \alpha \in \mathbb{R} \to \beta \in \mathbb{R} \to M \in \mathbb{R} \to \varphi_{55}) \to (x \in \mathbb{R} \to y \in \mathbb{R} \to \alpha \in \mathbb{R} \to \beta \in \mathbb{R} \to M \in \mathbb{R} \to \varphi_{53}) \to \varphi_{169}$
\item $\varphi_{229} :\equiv (x \in \mathbb{R} \to y \in \mathbb{R} \to \alpha \in \mathbb{R} \to \beta \in \mathbb{R} \to M \in \mathbb{R} \to \varphi_{52}) \to (x \in \mathbb{R} \to y \in \mathbb{R} \to \alpha \in \mathbb{R} \to \beta \in \mathbb{R} \to M \in \mathbb{R} \to \varphi_{62}) \to \varphi_{171}$
\item $\varphi_{230} :\equiv (x \in \mathbb{R} \to y \in \mathbb{R} \to \alpha \in \mathbb{R} \to \beta \in \mathbb{R} \to M \in \mathbb{R} \to \varphi_{54}) \to (x \in \mathbb{R} \to y \in \mathbb{R} \to \alpha \in \mathbb{R} \to \beta \in \mathbb{R} \to M \in \mathbb{R} \to \varphi_{63}) \to \varphi_{173}$
\item $\varphi_{231} :\equiv \varphi_{172} \to (x \in \mathbb{R} \to y \in \mathbb{R} \to \alpha \in \mathbb{R} \to \beta \in \mathbb{R} \to M \in \mathbb{R} \to \varphi_{53}) \to \varphi_{189}$
\item $\varphi_{232} :\equiv (x \in \mathbb{R} \to y \in \mathbb{R} \to \alpha \in \mathbb{R} \to \beta \in \mathbb{R} \to M \in \mathbb{R} \to \varphi_{55}) \to \varphi_{172} \to \varphi_{190}$
\item $\varphi_{233} :\equiv (x \in \mathbb{R} \to y \in \mathbb{R} \to \alpha \in \mathbb{R} \to \beta \in \mathbb{R} \to M \in \mathbb{R} \to \varphi_{52}) \to \varphi_{173} \to \varphi_{191}$
\item $\varphi_{234} :\equiv (x \in \mathbb{R} \to y \in \mathbb{R} \to \alpha \in \mathbb{R} \to \beta \in \mathbb{R} \to M \in \mathbb{R} \to \varphi_{63}) \to \varphi_{171} \to \varphi_{192}$
\item $\varphi_{235} :\equiv \varphi_{172} \to (x \in \mathbb{R} \to y \in \mathbb{R} \to \alpha \in \mathbb{R} \to \beta \in \mathbb{R} \to M \in \mathbb{R} \to \varphi_{53}) \to \varphi_{196}$
\item $\varphi_{236} :\equiv (x \in \mathbb{R} \to y \in \mathbb{R} \to \alpha \in \mathbb{R} \to \varphi_{148}) \to \varphi_{101} \to x \in \mathbb{R} \to y \in \mathbb{R} \to \alpha \in \mathbb{R} \to \varphi_{149}$
\item $\varphi_{237} :\equiv (x \in \mathbb{R} \to y \in \mathbb{R} \to \alpha \in \mathbb{R} \to \beta \in \mathbb{R} \to M \in \mathbb{R} \to \varphi_{53}) \to \varphi_{190} \to x \in \mathbb{R} \to y \in \mathbb{R} \to \varphi_{151}$
\item $\varphi_{238} :\equiv (x \in \mathbb{R} \to y \in \mathbb{R} \to \alpha \in \mathbb{R} \to \beta \in \mathbb{R} \to M \in \mathbb{R} \to \varphi_{55}) \to \varphi_{172} \to x \in \mathbb{R} \to y \in \mathbb{R} \to \varphi_{162}$
\item $\varphi_{239} :\equiv (x \in \mathbb{R} \to y \in \mathbb{R} \to \alpha \in \mathbb{R} \to \beta \in \mathbb{R} \to M \in \mathbb{R} \to \varphi_{52}) \to \varphi_{173} \to x \in \mathbb{R} \to y \in \mathbb{R} \to \varphi_{166}$
\item $\varphi_{240} :\equiv (x \in \mathbb{R} \to y \in \mathbb{R} \to \alpha \in \mathbb{R} \to \beta \in \mathbb{R} \to M \in \mathbb{R} \to \varphi_{63}) \to \varphi_{171} \to x \in \mathbb{R} \to y \in \mathbb{R} \to \varphi_{176}$
\item $\varphi_{241} :\equiv \varphi_{199} \to (x \in \mathbb{R} \to y \in \mathbb{R} \to \varphi_{153}) \to x \in \mathbb{R} \to y \in \mathbb{R} \to \alpha \in \mathbb{R} \to \varphi_{204}$
\item $\varphi_{242} :\equiv (\lvert x \rvert \in \mathbb{R} \to \lvert y - \beta \rvert \in \mathbb{R} \to \lvert x \rvert \cdot \lvert y - \beta \rvert \in \mathbb{R}) \to \varphi_{168} \to \varphi_{188} \to \varphi_{202}$
\item $\varphi_{243} :\equiv (\lvert \beta \rvert \in \mathbb{R} \to \lvert x - \alpha \rvert \in \mathbb{R} \to \lvert \beta \rvert \cdot \lvert x - \alpha \rvert \in \mathbb{R}) \to \varphi_{172} \to \varphi_{187} \to \varphi_{203}$
\item $\varphi_{244} :\equiv (x \in \mathbb{R} \to y \in \mathbb{R} \to \alpha \in \mathbb{R} \to \varphi_{205}) \to (x \in \mathbb{R} \to y \in \mathbb{R} \to \alpha \in \mathbb{R} \to \varphi_{197}) \to x \in \mathbb{R} \to y \in \mathbb{R} \to \varphi_{170}$
\item $\varphi_{245} :\equiv (x \in \mathbb{R} \to y \in \mathbb{R} \to \alpha \in \mathbb{R} \to \beta \in \mathbb{R} \to M \in \mathbb{R} \to \varphi_{54}) \to (x \in \mathbb{R} \to y \in \mathbb{R} \to \alpha \in \mathbb{R} \to \beta \in \mathbb{R} \to M \in \mathbb{R} \to \varphi_{62}) \to \varphi_{219}$
\item $\varphi_{246} :\equiv (x \cdot (y - \beta) \in \mathbb{R} \to \beta \cdot (x - \alpha) \in \mathbb{R} \to \varphi_{39}) \to \varphi_{191} \to \varphi_{192} \to x \in \mathbb{R} \to y \in \mathbb{R} \to \alpha \in \mathbb{R} \to \beta \in \mathbb{R} \to M \in \mathbb{R} \to \varphi_{145}$
\item $\varphi_{247} :\equiv (h \in \mathbb{R} \to \varphi_{28}) \to (x \in \mathbb{R} \to y \in \mathbb{R} \to \alpha \in \mathbb{R} \to \beta \in \mathbb{R} \to M \in \mathbb{R} \to \varphi_{53}) \to \varphi_{238}$
\item $\varphi_{248} :\equiv (x \in \mathbb{R} \to y \in \mathbb{R} \to \alpha \in \mathbb{R} \to \varphi_{194}) \to (\lvert x \cdot (y - \beta) \rvert = \lvert x \rvert \cdot \lvert y - \beta \rvert \to \varphi_{140}) \to \varphi_{223}$
\item $\varphi_{249} :\equiv (x \in \mathbb{R} \to y \in \mathbb{R} \to \alpha \in \mathbb{R} \to \varphi_{204}) \to (M \cdot h + \lvert \beta \rvert \cdot h = (M + \lvert \beta \rvert) \cdot h \to \varphi_{155}) \to \varphi_{224}$
\item $\varphi_{250} :\equiv (x \in \mathbb{R} \to y \in \mathbb{R} \to \alpha \in \mathbb{R} \to \beta \in \mathbb{R} \to M \in \mathbb{R} \to \varphi_{145}) \to (x \cdot (y - \beta) + \beta \cdot (x - \alpha) = x \cdot y - \alpha \cdot \beta \to \varphi_{152}) \to \varphi_{225}$
\item $\varphi_{251} :\equiv (x \in \mathbb{R} \to y \in \mathbb{R} \to \alpha \in \mathbb{R} \to \varphi_{198}) \to (\lvert \beta \cdot (x - \alpha) \rvert = \lvert \beta \rvert \cdot \lvert x - \alpha \rvert \to \varphi_{157}) \to \varphi_{226}$
\item $\varphi_{252} :\equiv (x \in \mathbb{R} \to y \in \mathbb{R} \to \varphi_{45}) \to (x \in \mathbb{R} \to y \in \mathbb{R} \to \alpha \in \mathbb{R} \to \beta \in \mathbb{R} \to M \in \mathbb{R} \to \varphi_{52}) \to \varphi_{245}$
\item $\varphi_{253} :\equiv (x \in \mathbb{R} \to y \in \mathbb{R} \to \alpha \in \mathbb{R} \to \beta \in \mathbb{R} \to M \in \mathbb{R} \to \varphi_{53}) \to \varphi_{147} \to \varphi_{161} \to \varphi_{146} \to \varphi_{160} \to \varphi_{199}$
\item $\varphi_{254} :\equiv (x \in \mathbb{R} \to y \in \mathbb{R} \to \alpha \in \mathbb{R} \to \beta \in \mathbb{R} \to M \in \mathbb{R} \to \varphi_{53}) \to \varphi_{150} \to \varphi_{159} \to \varphi_{167} \to \varphi_{158} \to x \in \mathbb{R} \to y \in \mathbb{R} \to \varphi_{153}$
\item $\varphi_{255} :\equiv (x \in \mathbb{R} \to y \in \mathbb{R} \to \alpha \in \mathbb{R} \to \beta \in \mathbb{R} \to M \in \mathbb{R} \to \varphi_{55}) \to \varphi_{188} \to \varphi_{253}$
\item $\varphi_{256} :\equiv (\lvert x \rvert \cdot \lvert y - \beta \rvert \in \mathbb{R} \to M \cdot h \in \mathbb{R} \to \varphi_{210}) \to \varphi_{202} \to \varphi_{169} \to \varphi_{203} \to \varphi_{189} \to \varphi_{241}$
\end{abbrevlist}
\end{fullproof}

% 自動生成 (prover/build.js): LimProd  ノード数 529
\begin{fullproof}{定理2.7}{$\forall a\,\forall \alpha\,\forall b\,\forall \beta\,(\mathrm{Conv}(a,\alpha) \to \mathrm{Conv}(b,\beta) \to \mathrm{Conv}(a \otimes b,\alpha \cdot \beta))$}{529}\phantomsection\label{tree:LimProd}
\noindent$\Pi_{1}$:
\begin{ptree}
\AXC{$\forall a\,\forall \alpha\,(\mathrm{Conv}(a,\alpha) \to \alpha \in \mathbb{R})$\lbl{[C2]}}
\AXC{$\varphi_{17}$\lbl{A2}}
\BIC{$\forall \alpha\,(\mathrm{Conv}(b,\alpha) \to \alpha \in \mathbb{R})$}
\AXC{$\varphi_{10}$\lbl{A2}}
\BIC{$\mathrm{Conv}(b,\beta) \to \beta \in \mathbb{R}$}
\AXC{$(\mathrm{Conv}(b,\beta) \to \beta \in \mathbb{R}) \to \varphi_{41}$\lbl{A1}}
\BIC{$\varphi_{41}$}
\end{ptree}
\noindent$\Pi_{2}$:
\begin{ptree}
\AXC{$\mathrm{Conv}(a,\alpha) \to \mathrm{Conv}(b,\beta) \to \mathrm{Conv}(b,\beta)$\lbl{A1}}
\AXC{$\stackrel{\Pi_{1}}{\vdots}$}\noLine
\UIC{$\varphi_{41}$}
\BIC{$\mathrm{Conv}(a,\alpha) \to \mathrm{Conv}(b,\beta) \to \beta \in \mathbb{R}$}
\end{ptree}
\noindent$\Pi_{3}$:
\begin{ptree}
\AXC{$\forall x \in \mathbb{R}\,(\lvert x \rvert \in \mathbb{R})$\lbl{[A5]}}
\AXC{$\forall x \in \mathbb{R}\,(\lvert x \rvert \in \mathbb{R}) \to \beta \in \mathbb{R} \to \lvert \beta \rvert \in \mathbb{R}$\lbl{A2}}
\BIC{$\beta \in \mathbb{R} \to \lvert \beta \rvert \in \mathbb{R}$}
\AXC{$(\beta \in \mathbb{R} \to \lvert \beta \rvert \in \mathbb{R}) \to \varphi_{223}$\lbl{A1}}
\BIC{$\varphi_{223}$}
\end{ptree}
\noindent$\Pi_{4}$:
\begin{ptree}
\AXC{$\forall x\,\forall y\,(x \in \mathbb{R} \to y \in \mathbb{R} \to x + y \in \mathbb{R})$\lbl{(R1)}}
\AXC{$\varphi_{77}$\lbl{A2}}
\BIC{$\forall y\,(M \in \mathbb{R} \to y \in \mathbb{R} \to M + y \in \mathbb{R})$}
\AXC{$\varphi_{87}$\lbl{A2}}
\BIC{$M \in \mathbb{R} \to \lvert \beta \rvert \in \mathbb{R} \to M + \lvert \beta \rvert \in \mathbb{R}$}
\AXC{$(M \in \mathbb{R} \to \lvert \beta \rvert \in \mathbb{R} \to M + \lvert \beta \rvert \in \mathbb{R}) \to \varphi_{247}$\lbl{A1}}
\BIC{$\varphi_{247}$}
\end{ptree}
\noindent$\Pi_{5}$:
\begin{ptree}
\AXC{$\mathrm{Conv}(a,\alpha) \to \mathrm{Conv}(b,\beta) \to \varphi_{73}$\lbl{A1}}
\AXC{$\varphi_{144}$\lbl{A1}}
\BIC{$\mathrm{Conv}(a,\alpha) \to \varphi_{29}$}
\AXC{$\stackrel{\Pi_{4}}{\vdots}$}\noLine
\UIC{$\varphi_{247}$}
\BIC{$(\mathrm{Conv}(a,\alpha) \to \varphi_{185}) \to \mathrm{Conv}(a,\alpha) \to \varphi_{191}$}
\end{ptree}
\noindent$\Pi_{6}$:
\begin{ptree}
\AXC{$\stackrel{\Pi_{2}}{\vdots}$}\noLine
\UIC{$\mathrm{Conv}(a,\alpha) \to \mathrm{Conv}(b,\beta) \to \beta \in \mathbb{R}$}
\AXC{$\stackrel{\Pi_{3}}{\vdots}$}\noLine
\UIC{$\varphi_{223}$}
\BIC{$\mathrm{Conv}(a,\alpha) \to \varphi_{185}$}
\AXC{$\stackrel{\Pi_{5}}{\vdots}$}\noLine
\UIC{$(\mathrm{Conv}(a,\alpha) \to \varphi_{185}) \to \mathrm{Conv}(a,\alpha) \to \varphi_{191}$}
\BIC{$\mathrm{Conv}(a,\alpha) \to \varphi_{191}$}
\end{ptree}
\noindent$\Pi_{7}$:
\begin{ptree}
\AXC{$\stackrel{\Pi_{6}}{\vdots}$}\noLine
\UIC{$\mathrm{Conv}(a,\alpha) \to \varphi_{191}$}
\AXC{$\forall x \in \mathbb{R}\,(x < x + 1)$\lbl{[N13]}}
\AXC{$\varphi_{67}$\lbl{A2}}
\BIC{$M + \lvert \beta \rvert \in \mathbb{R} \to M + \lvert \beta \rvert < M + \lvert \beta \rvert + 1$}
\AXC{$\varphi_{245}$\lbl{A1}}
\BIC{$(\mathrm{Conv}(a,\alpha) \to \varphi_{191}) \to \mathrm{Conv}(a,\alpha) \to \varphi_{200}$}
\BIC{$\mathrm{Conv}(a,\alpha) \to \varphi_{200}$}
\end{ptree}
\noindent$\Pi_{8}$:
\begin{ptree}
\AXC{$\forall x \in \mathbb{R}\,(\lvert x \rvert \in \mathbb{R})$\lbl{[A5]}}
\AXC{$\forall x \in \mathbb{R}\,(\lvert x \rvert \in \mathbb{R}) \to \beta \in \mathbb{R} \to \lvert \beta \rvert \in \mathbb{R}$\lbl{A2}}
\BIC{$\beta \in \mathbb{R} \to \lvert \beta \rvert \in \mathbb{R}$}
\AXC{$(\beta \in \mathbb{R} \to \lvert \beta \rvert \in \mathbb{R}) \to \varphi_{104}$\lbl{A1}}
\BIC{$\varphi_{104}$}
\end{ptree}
\noindent$\Pi_{9}$:
\begin{ptree}
\AXC{$\stackrel{\Pi_{2}}{\vdots}$}\noLine
\UIC{$\mathrm{Conv}(a,\alpha) \to \mathrm{Conv}(b,\beta) \to \beta \in \mathbb{R}$}
\AXC{$\stackrel{\Pi_{8}}{\vdots}$}\noLine
\UIC{$\varphi_{104}$}
\BIC{$\mathrm{Conv}(a,\alpha) \to \varphi_{38}$}
\end{ptree}
\noindent$\Pi_{10}$:
\begin{ptree}
\AXC{$\forall x\,\forall y\,(x \in \mathbb{R} \to y \in \mathbb{R} \to x + y \in \mathbb{R})$\lbl{(R1)}}
\AXC{$\varphi_{77}$\lbl{A2}}
\BIC{$\forall y\,(M \in \mathbb{R} \to y \in \mathbb{R} \to M + y \in \mathbb{R})$}
\AXC{$\varphi_{87}$\lbl{A2}}
\BIC{$M \in \mathbb{R} \to \lvert \beta \rvert \in \mathbb{R} \to M + \lvert \beta \rvert \in \mathbb{R}$}
\AXC{$(M \in \mathbb{R} \to \lvert \beta \rvert \in \mathbb{R} \to M + \lvert \beta \rvert \in \mathbb{R}) \to \varphi_{193}$\lbl{A1}}
\BIC{$\varphi_{193}$}
\end{ptree}
\noindent$\Pi_{11}$:
\begin{ptree}
\AXC{$\stackrel{\Pi_{9}}{\vdots}$}\noLine
\UIC{$\mathrm{Conv}(a,\alpha) \to \varphi_{38}$}
\AXC{$\mathrm{Conv}(a,\alpha) \to \mathrm{Conv}(b,\beta) \to \varphi_{73}$\lbl{A1}}
\AXC{$\varphi_{144}$\lbl{A1}}
\BIC{$\mathrm{Conv}(a,\alpha) \to \varphi_{29}$}
\AXC{$\stackrel{\Pi_{10}}{\vdots}$}\noLine
\UIC{$\varphi_{193}$}
\BIC{$(\mathrm{Conv}(a,\alpha) \to \varphi_{38}) \to \mathrm{Conv}(a,\alpha) \to \varphi_{54}$}
\BIC{$\mathrm{Conv}(a,\alpha) \to \varphi_{54}$}
\end{ptree}
\noindent$\Pi_{12}$:
\begin{ptree}
\AXC{$\stackrel{\Pi_{11}}{\vdots}$}\noLine
\UIC{$\mathrm{Conv}(a,\alpha) \to \varphi_{54}$}
\AXC{$\forall x \in \mathbb{R}\,(x < x + 1)$\lbl{[N13]}}
\AXC{$\varphi_{67}$\lbl{A2}}
\BIC{$M + \lvert \beta \rvert \in \mathbb{R} \to M + \lvert \beta \rvert < M + \lvert \beta \rvert + 1$}
\AXC{$(M + \lvert \beta \rvert \in \mathbb{R} \to M + \lvert \beta \rvert < M + \lvert \beta \rvert + 1) \to \varphi_{147}$\lbl{A1}}
\BIC{$\varphi_{147}$}
\BIC{$\mathrm{Conv}(a,\alpha) \to \mathrm{Conv}(b,\beta) \to \varphi_{49}$}
\end{ptree}
\noindent$\Pi_{13}$:
\begin{ptree}
\AXC{$\forall x \in \mathbb{R}\,(0 \leq \lvert x \rvert)$\lbl{[A6]}}
\AXC{$\forall x \in \mathbb{R}\,(0 \leq \lvert x \rvert) \to \beta \in \mathbb{R} \to 0 \leq \lvert \beta \rvert$\lbl{A2}}
\BIC{$\beta \in \mathbb{R} \to 0 \leq \lvert \beta \rvert$}
\AXC{$(\beta \in \mathbb{R} \to 0 \leq \lvert \beta \rvert) \to \varphi_{103}$\lbl{A1}}
\BIC{$\varphi_{103}$}
\end{ptree}
\noindent$\Pi_{14}$:
\begin{ptree}
\AXC{$1_{\omega} \in \mathbb{N}$\lbl{[N1]}}
\AXC{$\varphi_{24}$\lbl{A2}}
\AXC{$\mathrm{Conv}(a,\alpha) \to \mathrm{Conv}(b,\beta) \to \varphi_{73}$\lbl{A1}}
\AXC{$\varphi_{164}$\lbl{A1}}
\BIC{$\mathrm{Conv}(a,\alpha) \to \mathrm{Conv}(b,\beta) \to \varphi_{44}$}
\AXC{$\varphi_{199}$\lbl{A1}}
\BIC{$\varphi_{24} \to \mathrm{Conv}(a,\alpha) \to \mathrm{Conv}(b,\beta) \to \varphi_{50}$}
\BIC{$\mathrm{Conv}(a,\alpha) \to \mathrm{Conv}(b,\beta) \to \varphi_{50}$}
\AXC{$\varphi_{156}$\lbl{A1}}
\BIC{$1_{\omega} \in \mathbb{N} \to \mathrm{Conv}(a,\alpha) \to \varphi_{46}$}
\BIC{$\mathrm{Conv}(a,\alpha) \to \varphi_{46}$}
\end{ptree}
\noindent$\Pi_{15}$:
\begin{ptree}
\AXC{$\forall a\,\forall \alpha\,(\mathrm{Conv}(a,\alpha) \to \mathrm{Seq}(a))$\lbl{[C1]}}
\AXC{$\varphi_{6}$\lbl{A2}}
\BIC{$\forall \alpha\,(\mathrm{Conv}(a,\alpha) \to \mathrm{Seq}(a))$}
\AXC{$\varphi_{1}$\lbl{A2}}
\BIC{$\mathrm{Conv}(a,\alpha) \to \mathrm{Seq}(a)$}
\AXC{$(\mathrm{Conv}(a,\alpha) \to \mathrm{Seq}(a)) \to \varphi_{12}$\lbl{A1}}
\BIC{$\varphi_{12}$}
\end{ptree}
\noindent$\Pi_{16}$:
\begin{ptree}
\AXC{$\mathrm{Conv}(a,\alpha) \to \mathrm{Conv}(a,\alpha)$\lbl{A1}}
\AXC{$\stackrel{\Pi_{15}}{\vdots}$}\noLine
\UIC{$\varphi_{12}$}
\BIC{$\mathrm{Conv}(a,\alpha) \to \mathrm{Conv}(b,\beta) \to \mathrm{Seq}(a)$}
\end{ptree}
\noindent$\Pi_{17}$:
\begin{ptree}
\AXC{$\forall a\,\forall n\,(\mathrm{Seq}(a) \to n \in \mathbb{N} \to a_{n} \in \mathbb{R})$\lbl{[S1]}}
\AXC{$\varphi_{68}$\lbl{A2}}
\BIC{$\forall n\,(\mathrm{Seq}(a) \to n \in \mathbb{N} \to a_{n} \in \mathbb{R})$}
\AXC{$\varphi_{69}$\lbl{A2}}
\BIC{$\mathrm{Seq}(a) \to 1_{\omega} \in \mathbb{N} \to a_{1_{\omega}} \in \mathbb{R}$}
\end{ptree}
\noindent$\Pi_{18}$:
\begin{ptree}
\AXC{$\stackrel{\Pi_{17}}{\vdots}$}\noLine
\UIC{$\mathrm{Seq}(a) \to 1_{\omega} \in \mathbb{N} \to a_{1_{\omega}} \in \mathbb{R}$}
\AXC{$(\mathrm{Seq}(a) \to 1_{\omega} \in \mathbb{N} \to a_{1_{\omega}} \in \mathbb{R}) \to \varphi_{113}$\lbl{A1}}
\BIC{$\varphi_{113}$}
\end{ptree}
\noindent$\Pi_{19}$:
\begin{ptree}
\AXC{$1_{\omega} \in \mathbb{N}$\lbl{[N1]}}
\AXC{$\stackrel{\Pi_{16}}{\vdots}$}\noLine
\UIC{$\mathrm{Conv}(a,\alpha) \to \mathrm{Conv}(b,\beta) \to \mathrm{Seq}(a)$}
\AXC{$\stackrel{\Pi_{18}}{\vdots}$}\noLine
\UIC{$\varphi_{113}$}
\BIC{$1_{\omega} \in \mathbb{N} \to \mathrm{Conv}(a,\alpha) \to \varphi_{51}$}
\BIC{$\mathrm{Conv}(a,\alpha) \to \varphi_{51}$}
\end{ptree}
\noindent$\Pi_{20}$:
\begin{ptree}
\AXC{$\stackrel{\Pi_{19}}{\vdots}$}\noLine
\UIC{$\mathrm{Conv}(a,\alpha) \to \varphi_{51}$}
\AXC{$\forall x \in \mathbb{R}\,(0 \leq \lvert x \rvert)$\lbl{[A6]}}
\AXC{$\varphi_{26}$\lbl{A2}}
\BIC{$a_{1_{\omega}} \in \mathbb{R} \to 0 \leq \lvert a_{1_{\omega}} \rvert$}
\AXC{$\varphi_{169}$\lbl{A1}}
\BIC{$(\mathrm{Conv}(a,\alpha) \to \varphi_{51}) \to \mathrm{Conv}(a,\alpha) \to \varphi_{47}$}
\BIC{$\mathrm{Conv}(a,\alpha) \to \varphi_{47}$}
\end{ptree}
\noindent$\Pi_{21}$:
\begin{ptree}
\AXC{$\forall x\,\forall y\,\forall z\,(x \leq y \to y \leq z \to x \leq z)$\lbl{[O5]}}
\AXC{$\varphi_{3}$\lbl{A2}}
\BIC{$\forall y\,\forall z\,(0 \leq y \to y \leq z \to 0 \leq z)$}
\AXC{$\forall y\,\forall z\,(0 \leq y \to y \leq z \to 0 \leq z) \to \varphi_{4}$\lbl{A2}}
\BIC{$\varphi_{4}$}
\end{ptree}
\noindent$\Pi_{22}$:
\begin{ptree}
\AXC{$\stackrel{\Pi_{20}}{\vdots}$}\noLine
\UIC{$\mathrm{Conv}(a,\alpha) \to \varphi_{47}$}
\AXC{$\stackrel{\Pi_{21}}{\vdots}$}\noLine
\UIC{$\varphi_{4}$}
\AXC{$\varphi_{4} \to 0 \leq \lvert a_{1_{\omega}} \rvert \to \lvert a_{1_{\omega}} \rvert \leq M \to 0 \leq M$\lbl{A2}}
\BIC{$0 \leq \lvert a_{1_{\omega}} \rvert \to \lvert a_{1_{\omega}} \rvert \leq M \to 0 \leq M$}
\AXC{$\varphi_{221}$\lbl{A1}}
\BIC{$\varphi_{183}$}
\BIC{$(\mathrm{Conv}(a,\alpha) \to \varphi_{46}) \to \mathrm{Conv}(a,\alpha) \to \varphi_{20}$}
\end{ptree}
\noindent$\Pi_{23}$:
\begin{ptree}
\AXC{$\stackrel{\Pi_{14}}{\vdots}$}\noLine
\UIC{$\mathrm{Conv}(a,\alpha) \to \varphi_{46}$}
\AXC{$\stackrel{\Pi_{22}}{\vdots}$}\noLine
\UIC{$(\mathrm{Conv}(a,\alpha) \to \varphi_{46}) \to \mathrm{Conv}(a,\alpha) \to \varphi_{20}$}
\BIC{$\mathrm{Conv}(a,\alpha) \to \varphi_{20}$}
\end{ptree}
\noindent$\Pi_{24}$:
\begin{ptree}
\AXC{$\forall x\,\forall y\,(x \in \mathbb{R} \to y \in \mathbb{R} \to 0 \leq x \to 0 \leq y \to 0 \leq x + y)$\lbl{[E6]}}
\AXC{$\varphi_{91}$\lbl{A2}}
\BIC{$\forall y\,(M \in \mathbb{R} \to y \in \mathbb{R} \to 0 \leq M \to 0 \leq y \to 0 \leq M + y)$}
\AXC{$\varphi_{99}$\lbl{A2}}
\BIC{$M \in \mathbb{R} \to \lvert \beta \rvert \in \mathbb{R} \to 0 \leq M \to 0 \leq \lvert \beta \rvert \to 0 \leq M + \lvert \beta \rvert$}
\AXC{$\varphi_{237}$\lbl{A1}}
\BIC{$(\mathrm{Conv}(a,\alpha) \to \varphi_{29}) \to (\mathrm{Conv}(a,\alpha) \to \varphi_{38}) \to \varphi_{176}$}
\end{ptree}
\noindent$\Pi_{25}$:
\begin{ptree}
\AXC{$\mathrm{Conv}(a,\alpha) \to \mathrm{Conv}(b,\beta) \to \varphi_{73}$\lbl{A1}}
\AXC{$\varphi_{144}$\lbl{A1}}
\BIC{$\mathrm{Conv}(a,\alpha) \to \varphi_{29}$}
\AXC{$\stackrel{\Pi_{24}}{\vdots}$}\noLine
\UIC{$(\mathrm{Conv}(a,\alpha) \to \varphi_{29}) \to (\mathrm{Conv}(a,\alpha) \to \varphi_{38}) \to \varphi_{176}$}
\BIC{$(\mathrm{Conv}(a,\alpha) \to \varphi_{38}) \to \varphi_{176}$}
\end{ptree}
\noindent$\Pi_{26}$:
\begin{ptree}
\AXC{$\stackrel{\Pi_{23}}{\vdots}$}\noLine
\UIC{$\mathrm{Conv}(a,\alpha) \to \varphi_{20}$}
\AXC{$\stackrel{\Pi_{9}}{\vdots}$}\noLine
\UIC{$\mathrm{Conv}(a,\alpha) \to \varphi_{38}$}
\AXC{$\stackrel{\Pi_{25}}{\vdots}$}\noLine
\UIC{$(\mathrm{Conv}(a,\alpha) \to \varphi_{38}) \to \varphi_{176}$}
\BIC{$\varphi_{176}$}
\BIC{$(\mathrm{Conv}(a,\alpha) \to \varphi_{27}) \to \mathrm{Conv}(a,\alpha) \to \varphi_{36}$}
\end{ptree}
\noindent$\Pi_{27}$:
\begin{ptree}
\AXC{$\stackrel{\Pi_{2}}{\vdots}$}\noLine
\UIC{$\mathrm{Conv}(a,\alpha) \to \mathrm{Conv}(b,\beta) \to \beta \in \mathbb{R}$}
\AXC{$\stackrel{\Pi_{13}}{\vdots}$}\noLine
\UIC{$\varphi_{103}$}
\BIC{$\mathrm{Conv}(a,\alpha) \to \varphi_{27}$}
\AXC{$\stackrel{\Pi_{26}}{\vdots}$}\noLine
\UIC{$(\mathrm{Conv}(a,\alpha) \to \varphi_{27}) \to \mathrm{Conv}(a,\alpha) \to \varphi_{36}$}
\BIC{$\mathrm{Conv}(a,\alpha) \to \varphi_{36}$}
\end{ptree}
\noindent$\Pi_{28}$:
\begin{ptree}
\AXC{$\forall x\,\forall y\,\forall z\,(x \leq y \to y < z \to x < z)$\lbl{[O4]}}
\AXC{$\varphi_{14}$\lbl{A2}}
\BIC{$\forall y\,\forall z\,(0 \leq y \to y < z \to 0 < z)$}
\AXC{$\varphi_{42}$\lbl{A2}}
\BIC{$\forall z\,(0 \leq M + \lvert \beta \rvert \to M + \lvert \beta \rvert < z \to 0 < z)$}
\AXC{$\forall z\,(0 \leq M + \lvert \beta \rvert \to M + \lvert \beta \rvert < z \to 0 < z) \to \varphi_{37}$\lbl{A2}}
\BIC{$\varphi_{37}$}
\end{ptree}
\noindent$\Pi_{29}$:
\begin{ptree}
\AXC{$\stackrel{\Pi_{12}}{\vdots}$}\noLine
\UIC{$\mathrm{Conv}(a,\alpha) \to \mathrm{Conv}(b,\beta) \to \varphi_{49}$}
\AXC{$\stackrel{\Pi_{27}}{\vdots}$}\noLine
\UIC{$\mathrm{Conv}(a,\alpha) \to \varphi_{36}$}
\AXC{$\stackrel{\Pi_{28}}{\vdots}$}\noLine
\UIC{$\varphi_{37}$}
\AXC{$\varphi_{37} \to (\mathrm{Conv}(a,\alpha) \to \varphi_{36}) \to \varphi_{149}$\lbl{A1}}
\BIC{$(\mathrm{Conv}(a,\alpha) \to \varphi_{36}) \to \varphi_{149}$}
\BIC{$\varphi_{149}$}
\BIC{$\mathrm{Conv}(a,\alpha) \to \varphi_{57}$}
\end{ptree}
\noindent$\Pi_{30}$:
\begin{ptree}
\AXC{$\forall x \in \mathbb{R}\,(\lvert x \rvert \in \mathbb{R})$\lbl{[A5]}}
\AXC{$\forall x \in \mathbb{R}\,(\lvert x \rvert \in \mathbb{R}) \to \beta \in \mathbb{R} \to \lvert \beta \rvert \in \mathbb{R}$\lbl{A2}}
\BIC{$\beta \in \mathbb{R} \to \lvert \beta \rvert \in \mathbb{R}$}
\AXC{$(\beta \in \mathbb{R} \to \lvert \beta \rvert \in \mathbb{R}) \to \varphi_{115}$\lbl{A1}}
\BIC{$\varphi_{115}$}
\end{ptree}
\noindent$\Pi_{31}$:
\begin{ptree}
\AXC{$\forall x\,\forall y\,(x \in \mathbb{R} \to y \in \mathbb{R} \to x + y \in \mathbb{R})$\lbl{(R1)}}
\AXC{$\varphi_{77}$\lbl{A2}}
\BIC{$\forall y\,(M \in \mathbb{R} \to y \in \mathbb{R} \to M + y \in \mathbb{R})$}
\AXC{$\varphi_{87}$\lbl{A2}}
\BIC{$M \in \mathbb{R} \to \lvert \beta \rvert \in \mathbb{R} \to M + \lvert \beta \rvert \in \mathbb{R}$}
\AXC{$\varphi_{232}$\lbl{A1}}
\BIC{$(\mathrm{Conv}(a,\alpha) \to \varphi_{29}) \to \varphi_{171}$}
\end{ptree}
\noindent$\Pi_{32}$:
\begin{ptree}
\AXC{$\mathrm{Conv}(a,\alpha) \to \mathrm{Conv}(b,\beta) \to \varphi_{73}$\lbl{A1}}
\AXC{$\varphi_{144}$\lbl{A1}}
\BIC{$\mathrm{Conv}(a,\alpha) \to \varphi_{29}$}
\AXC{$\stackrel{\Pi_{31}}{\vdots}$}\noLine
\UIC{$(\mathrm{Conv}(a,\alpha) \to \varphi_{29}) \to \varphi_{171}$}
\BIC{$\varphi_{171}$}
\end{ptree}
\noindent$\Pi_{33}$:
\begin{ptree}
\AXC{$\forall x\,\forall y\,(x \in \mathbb{R} \to y \in \mathbb{R} \to x + y \in \mathbb{R})$\lbl{(R1)}}
\AXC{$\varphi_{100}$\lbl{A2}}
\BIC{$\forall y\,(M + \lvert \beta \rvert \in \mathbb{R} \to y \in \mathbb{R} \to M + \lvert \beta \rvert + y \in \mathbb{R})$}
\AXC{$\varphi_{109}$\lbl{A2}}
\BIC{$M + \lvert \beta \rvert \in \mathbb{R} \to 1 \in \mathbb{R} \to M + \lvert \beta \rvert + 1 \in \mathbb{R}$}
\AXC{$(M + \lvert \beta \rvert \in \mathbb{R} \to 1 \in \mathbb{R} \to M + \lvert \beta \rvert + 1 \in \mathbb{R}) \to \varphi_{181}$\lbl{A1}}
\BIC{$\varphi_{181}$}
\end{ptree}
\noindent$\Pi_{34}$:
\begin{ptree}
\AXC{$1 \in \mathbb{R}$\lbl{(R1)}}
\AXC{$\stackrel{\Pi_{2}}{\vdots}$}\noLine
\UIC{$\mathrm{Conv}(a,\alpha) \to \mathrm{Conv}(b,\beta) \to \beta \in \mathbb{R}$}
\AXC{$\stackrel{\Pi_{30}}{\vdots}$}\noLine
\UIC{$\varphi_{115}$}
\BIC{$\mathrm{Conv}(a,\alpha) \to \mathrm{Conv}(b,\beta) \to \varphi_{72}$}
\AXC{$\stackrel{\Pi_{32}}{\vdots}$}\noLine
\UIC{$\varphi_{171}$}
\BIC{$\mathrm{Conv}(a,\alpha) \to \mathrm{Conv}(b,\beta) \to \varphi_{76}$}
\AXC{$\stackrel{\Pi_{33}}{\vdots}$}\noLine
\UIC{$\varphi_{181}$}
\BIC{$1 \in \mathbb{R} \to \mathrm{Conv}(a,\alpha) \to \mathrm{Conv}(b,\beta) \to \varphi_{81}$}
\BIC{$\mathrm{Conv}(a,\alpha) \to \mathrm{Conv}(b,\beta) \to \varphi_{81}$}
\end{ptree}
\noindent$\Pi_{35}$:
\begin{ptree}
\AXC{$\forall x \in \mathbb{R}\,(0 < x \to 0 < x^{-1})$\lbl{[F34]}}
\AXC{$\forall x \in \mathbb{R}\,(0 < x \to 0 < x^{-1}) \to \varphi_{56}$\lbl{A2}}
\BIC{$\varphi_{56}$}
\end{ptree}
\noindent$\Pi_{36}$:
\begin{ptree}
\AXC{$\stackrel{\Pi_{34}}{\vdots}$}\noLine
\UIC{$\mathrm{Conv}(a,\alpha) \to \mathrm{Conv}(b,\beta) \to \varphi_{81}$}
\AXC{$\stackrel{\Pi_{35}}{\vdots}$}\noLine
\UIC{$\varphi_{56}$}
\AXC{$\varphi_{56} \to (\mathrm{Conv}(a,\alpha) \to \mathrm{Conv}(b,\beta) \to \varphi_{81}) \to \varphi_{166}$\lbl{A1}}
\BIC{$(\mathrm{Conv}(a,\alpha) \to \mathrm{Conv}(b,\beta) \to \varphi_{81}) \to \varphi_{166}$}
\BIC{$\varphi_{166}$}
\end{ptree}
\noindent$\Pi_{37}$:
\begin{ptree}
\AXC{$x = x$\lbl{A5}}
\AXC{$x = y \to x = x \to y = x$\lbl{A6}}
\AXC{$(x = y \to x = x \to y = x) \to x = x \to x = y \to y = x$\lbl{A1}}
\BIC{$x = x \to x = y \to y = x$}
\BIC{$x = y \to y = x$}
\UIC{$\forall y\,(x = y \to y = x)$}
\UIC{$\forall x\,\forall y\,(x = y \to y = x)$}
\end{ptree}
\noindent$\Pi_{38}$:
\begin{ptree}
\AXC{$\forall x\,\forall y\,(x < y \to x \neq y)$\lbl{[F26]}}
\AXC{$\forall x\,\forall y\,(x < y \to x \neq y) \to \forall y\,(0 < y \to 0 \neq y)$\lbl{A2}}
\BIC{$\forall y\,(0 < y \to 0 \neq y)$}
\AXC{$\varphi_{22}$\lbl{A2}}
\BIC{$0 < M + \lvert \beta \rvert + 1 \to 0 \neq M + \lvert \beta \rvert + 1$}
\AXC{$\varphi_{172}$\lbl{A1}}
\BIC{$(\mathrm{Conv}(a,\alpha) \to \varphi_{57}) \to \mathrm{Conv}(a,\alpha) \to \varphi_{52}$}
\end{ptree}
\noindent$\Pi_{39}$:
\begin{ptree}
\AXC{$\stackrel{\Pi_{29}}{\vdots}$}\noLine
\UIC{$\mathrm{Conv}(a,\alpha) \to \varphi_{57}$}
\AXC{$\stackrel{\Pi_{38}}{\vdots}$}\noLine
\UIC{$(\mathrm{Conv}(a,\alpha) \to \varphi_{57}) \to \mathrm{Conv}(a,\alpha) \to \varphi_{52}$}
\BIC{$\mathrm{Conv}(a,\alpha) \to \varphi_{52}$}
\AXC{$(\mathrm{Conv}(a,\alpha) \to \varphi_{52}) \to \varphi_{111}$\lbl{A1}}
\BIC{$\varphi_{111}$}
\end{ptree}
\noindent$\Pi_{40}$:
\begin{ptree}
\AXC{$\stackrel{\Pi_{37}}{\vdots}$}\noLine
\UIC{$\forall x\,\forall y\,(x = y \to y = x)$}
\AXC{$\varphi_{18}$\lbl{A2}}
\BIC{$\forall y\,(M + \lvert \beta \rvert + 1 = y \to y = M + \lvert \beta \rvert + 1)$}
\AXC{$\varphi_{70}$\lbl{A2}}
\BIC{$M + \lvert \beta \rvert + 1 = 0 \to 0 = M + \lvert \beta \rvert + 1$}
\AXC{$\stackrel{\Pi_{39}}{\vdots}$}\noLine
\UIC{$\varphi_{111}$}
\BIC{$\mathrm{Conv}(a,\alpha) \to \varphi_{53}$}
\end{ptree}
\noindent$\Pi_{41}$:
\begin{ptree}
\AXC{$\forall x \in \mathbb{R}\,(x \neq 0 \to x^{-1} \in \mathbb{R})$\lbl{(R1)}}
\AXC{$\forall x \in \mathbb{R}\,(x \neq 0 \to x^{-1} \in \mathbb{R}) \to \varphi_{62}$\lbl{A2}}
\BIC{$\varphi_{62}$}
\end{ptree}
\noindent$\Pi_{42}$:
\begin{ptree}
\AXC{$\stackrel{\Pi_{34}}{\vdots}$}\noLine
\UIC{$\mathrm{Conv}(a,\alpha) \to \mathrm{Conv}(b,\beta) \to \varphi_{81}$}
\AXC{$\stackrel{\Pi_{41}}{\vdots}$}\noLine
\UIC{$\varphi_{62}$}
\AXC{$\varphi_{62} \to (\mathrm{Conv}(a,\alpha) \to \mathrm{Conv}(b,\beta) \to \varphi_{81}) \to \varphi_{168}$\lbl{A1}}
\BIC{$(\mathrm{Conv}(a,\alpha) \to \mathrm{Conv}(b,\beta) \to \varphi_{81}) \to \varphi_{168}$}
\BIC{$\varphi_{168}$}
\end{ptree}
\noindent$\Pi_{43}$:
\begin{ptree}
\AXC{$\stackrel{\Pi_{40}}{\vdots}$}\noLine
\UIC{$\mathrm{Conv}(a,\alpha) \to \varphi_{53}$}
\AXC{$\stackrel{\Pi_{42}}{\vdots}$}\noLine
\UIC{$\varphi_{168}$}
\BIC{$\mathrm{Conv}(a,\alpha) \to \mathrm{Conv}(b,\beta) \to \varphi_{90}$}
\end{ptree}
\noindent$\Pi_{44}$:
\begin{ptree}
\AXC{$\mathrm{Conv}(a,\alpha) \to \varphi_{64}$\lbl{A1}}
\AXC{$\forall x\,\forall y\,(x \in \mathbb{R} \to y \in \mathbb{R} \to 0 < x \to 0 < y \to 0 < x \cdot y)$\lbl{(R15)}}
\AXC{$\varphi_{105}$\lbl{A2}}
\BIC{$\forall y\,(\varepsilon \in \mathbb{R} \to y \in \mathbb{R} \to 0 < \varepsilon \to 0 < y \to 0 < \varepsilon \cdot y)$}
\AXC{$\varphi_{138}$\lbl{A2}}
\BIC{$\varepsilon \in \mathbb{R} \to (M + \lvert \beta \rvert + 1)^{-1} \in \mathbb{R} \to \varphi_{39}$}
\AXC{$\varphi_{256}$\lbl{A1}}
\BIC{$(\mathrm{Conv}(a,\alpha) \to \varphi_{64}) \to \varphi_{238}$}
\BIC{$\varphi_{238}$}
\end{ptree}
\noindent$\Pi_{45}$:
\begin{ptree}
\AXC{$\mathrm{Conv}(a,\alpha) \to \mathrm{Conv}(b,\beta) \to \varphi_{55}$\lbl{A1}}
\AXC{$\stackrel{\Pi_{43}}{\vdots}$}\noLine
\UIC{$\mathrm{Conv}(a,\alpha) \to \mathrm{Conv}(b,\beta) \to \varphi_{90}$}
\AXC{$\stackrel{\Pi_{44}}{\vdots}$}\noLine
\UIC{$\varphi_{238}$}
\BIC{$(\mathrm{Conv}(a,\alpha) \to \mathrm{Conv}(b,\beta) \to \varphi_{55}) \to \varphi_{182}$}
\BIC{$\varphi_{182}$}
\end{ptree}
\noindent$\Pi_{46}$:
\begin{ptree}
\AXC{$\stackrel{\Pi_{29}}{\vdots}$}\noLine
\UIC{$\mathrm{Conv}(a,\alpha) \to \varphi_{57}$}
\AXC{$\stackrel{\Pi_{36}}{\vdots}$}\noLine
\UIC{$\varphi_{166}$}
\BIC{$\mathrm{Conv}(a,\alpha) \to \mathrm{Conv}(b,\beta) \to \varphi_{84}$}
\AXC{$\stackrel{\Pi_{45}}{\vdots}$}\noLine
\UIC{$\varphi_{182}$}
\BIC{$\mathrm{Conv}(a,\alpha) \to \mathrm{Conv}(b,\beta) \to \varphi_{93}$}
\end{ptree}
\noindent$\Pi_{47}$:
\begin{ptree}
\AXC{$\stackrel{\Pi_{6}}{\vdots}$}\noLine
\UIC{$\mathrm{Conv}(a,\alpha) \to \varphi_{191}$}
\AXC{$\forall x\,\forall y\,(x \in \mathbb{R} \to y \in \mathbb{R} \to x + y \in \mathbb{R})$\lbl{(R1)}}
\AXC{$\varphi_{100}$\lbl{A2}}
\BIC{$\forall y\,(M + \lvert \beta \rvert \in \mathbb{R} \to y \in \mathbb{R} \to M + \lvert \beta \rvert + y \in \mathbb{R})$}
\AXC{$\varphi_{109}$\lbl{A2}}
\BIC{$M + \lvert \beta \rvert \in \mathbb{R} \to 1 \in \mathbb{R} \to M + \lvert \beta \rvert + 1 \in \mathbb{R}$}
\AXC{$\varphi_{251}$\lbl{A1}}
\BIC{$(\mathrm{Conv}(a,\alpha) \to \varphi_{191}) \to 1 \in \mathbb{R} \to \mathrm{Conv}(a,\alpha) \to \varphi_{196}$}
\BIC{$1 \in \mathbb{R} \to \mathrm{Conv}(a,\alpha) \to \varphi_{196}$}
\end{ptree}
\noindent$\Pi_{48}$:
\begin{ptree}
\AXC{$1 \in \mathbb{R}$\lbl{(R1)}}
\AXC{$\stackrel{\Pi_{47}}{\vdots}$}\noLine
\UIC{$1 \in \mathbb{R} \to \mathrm{Conv}(a,\alpha) \to \varphi_{196}$}
\BIC{$\mathrm{Conv}(a,\alpha) \to \varphi_{196}$}
\end{ptree}
\noindent$\Pi_{49}$:
\begin{ptree}
\AXC{$\stackrel{\Pi_{48}}{\vdots}$}\noLine
\UIC{$\mathrm{Conv}(a,\alpha) \to \varphi_{196}$}
\AXC{$\forall x \in \mathbb{R}\,(x \neq 0 \to x^{-1} \in \mathbb{R})$\lbl{(R1)}}
\AXC{$\forall x \in \mathbb{R}\,(x \neq 0 \to x^{-1} \in \mathbb{R}) \to \varphi_{62}$\lbl{A2}}
\BIC{$\varphi_{62}$}
\AXC{$\varphi_{62} \to \varphi_{253}$\lbl{A1}}
\BIC{$\varphi_{253}$}
\BIC{$(\mathrm{Conv}(a,\alpha) \to \varphi_{53}) \to \mathrm{Conv}(a,\alpha) \to \varphi_{198}$}
\end{ptree}
\noindent$\Pi_{50}$:
\begin{ptree}
\AXC{$\forall x\,\forall y\,(x \in \mathbb{R} \to y \in \mathbb{R} \to x \cdot y \in \mathbb{R})$\lbl{(R1)}}
\AXC{$\varphi_{80}$\lbl{A2}}
\BIC{$\forall y\,(\varepsilon \in \mathbb{R} \to y \in \mathbb{R} \to \varepsilon \cdot y \in \mathbb{R})$}
\AXC{$\forall y\,(\varepsilon \in \mathbb{R} \to y \in \mathbb{R} \to \varepsilon \cdot y \in \mathbb{R}) \to \varepsilon \in \mathbb{R} \to \varphi_{33}$\lbl{A2}}
\BIC{$\varepsilon \in \mathbb{R} \to \varphi_{33}$}
\end{ptree}
\noindent$\Pi_{51}$:
\begin{ptree}
\AXC{$\mathrm{Conv}(a,\alpha) \to \varphi_{64}$\lbl{A1}}
\AXC{$\stackrel{\Pi_{50}}{\vdots}$}\noLine
\UIC{$\varepsilon \in \mathbb{R} \to \varphi_{33}$}
\AXC{$(\varepsilon \in \mathbb{R} \to \varphi_{33}) \to \varphi_{254}$\lbl{A1}}
\BIC{$\varphi_{254}$}
\BIC{$(\mathrm{Conv}(a,\alpha) \to \varphi_{198}) \to \mathrm{Conv}(a,\alpha) \to \varphi_{204}$}
\end{ptree}
\noindent$\Pi_{52}$:
\begin{ptree}
\AXC{$\stackrel{\Pi_{40}}{\vdots}$}\noLine
\UIC{$\mathrm{Conv}(a,\alpha) \to \varphi_{53}$}
\AXC{$\stackrel{\Pi_{49}}{\vdots}$}\noLine
\UIC{$(\mathrm{Conv}(a,\alpha) \to \varphi_{53}) \to \mathrm{Conv}(a,\alpha) \to \varphi_{198}$}
\BIC{$\mathrm{Conv}(a,\alpha) \to \varphi_{198}$}
\AXC{$\stackrel{\Pi_{51}}{\vdots}$}\noLine
\UIC{$(\mathrm{Conv}(a,\alpha) \to \varphi_{198}) \to \mathrm{Conv}(a,\alpha) \to \varphi_{204}$}
\BIC{$\mathrm{Conv}(a,\alpha) \to \varphi_{204}$}
\end{ptree}
\noindent$\Pi_{53}$:
\begin{ptree}
\AXC{$\stackrel{\Pi_{6}}{\vdots}$}\noLine
\UIC{$\mathrm{Conv}(a,\alpha) \to \varphi_{191}$}
\AXC{$\varphi_{58}$\lbl{[F32]}}
\AXC{$\varphi_{58} \to \varphi_{98}$\lbl{A2}}
\BIC{$\varphi_{98}$}
\AXC{$\varphi_{98} \to \varphi_{116}$\lbl{A2}}
\BIC{$\varphi_{116}$}
\AXC{$\varphi_{116} \to \varphi_{190}$\lbl{A2}}
\BIC{$\varphi_{190}$}
\AXC{$\varphi_{190} \to \varphi_{270}$\lbl{A1}}
\BIC{$\varphi_{270}$}
\BIC{$(\mathrm{Conv}(a,\alpha) \to \varphi_{196}) \to (\mathrm{Conv}(a,\alpha) \to \varphi_{204}) \to \varphi_{257}$}
\end{ptree}
\noindent$\Pi_{54}$:
\begin{ptree}
\AXC{$\stackrel{\Pi_{52}}{\vdots}$}\noLine
\UIC{$\mathrm{Conv}(a,\alpha) \to \varphi_{204}$}
\AXC{$\stackrel{\Pi_{48}}{\vdots}$}\noLine
\UIC{$\mathrm{Conv}(a,\alpha) \to \varphi_{196}$}
\AXC{$\stackrel{\Pi_{53}}{\vdots}$}\noLine
\UIC{$(\mathrm{Conv}(a,\alpha) \to \varphi_{196}) \to (\mathrm{Conv}(a,\alpha) \to \varphi_{204}) \to \varphi_{257}$}
\BIC{$(\mathrm{Conv}(a,\alpha) \to \varphi_{204}) \to \varphi_{257}$}
\BIC{$\varphi_{257}$}
\end{ptree}
\noindent$\Pi_{55}$:
\begin{ptree}
\AXC{$y = z \to x = y \to x = z$\lbl{A6}}
\AXC{$(y = z \to x = y \to x = z) \to x = y \to y = z \to x = z$\lbl{A1}}
\BIC{$x = y \to y = z \to x = z$}
\UIC{$\forall z\,(x = y \to y = z \to x = z)$}
\UIC{$\forall y\,\forall z\,(x = y \to y = z \to x = z)$}
\UIC{$\forall x\,\forall y\,\forall z\,(x = y \to y = z \to x = z)$}
\AXC{$\forall x\,\forall y\,\forall z\,(x = y \to y = z \to x = z) \to \varphi_{101}$\lbl{A2}}
\BIC{$\varphi_{101}$}
\end{ptree}
\noindent$\Pi_{56}$:
\begin{ptree}
\AXC{$\mathrm{Conv}(a,\alpha) \to \varphi_{64}$\lbl{A1}}
\AXC{$\forall x\,\forall y\,(x \in \mathbb{R} \to y \in \mathbb{R} \to y \neq 0 \to x \cdot y^{-1} \cdot y = x)$\lbl{[E23]}}
\AXC{$\varphi_{102}$\lbl{A2}}
\BIC{$\forall y\,(\varepsilon \in \mathbb{R} \to y \in \mathbb{R} \to y \neq 0 \to \varepsilon \cdot y^{-1} \cdot y = \varepsilon)$}
\AXC{$\varphi_{135}$\lbl{A2}}
\BIC{$\varepsilon \in \mathbb{R} \to M + \lvert \beta \rvert + 1 \in \mathbb{R} \to \varphi_{48}$}
\AXC{$\varphi_{263}$\lbl{A1}}
\BIC{$(\mathrm{Conv}(a,\alpha) \to \varphi_{64}) \to (\mathrm{Conv}(a,\alpha) \to \varphi_{196}) \to \varphi_{235}$}
\BIC{$(\mathrm{Conv}(a,\alpha) \to \varphi_{196}) \to \varphi_{235}$}
\end{ptree}
\noindent$\Pi_{57}$:
\begin{ptree}
\AXC{$\forall x\,\forall y\,(x \in \mathbb{R} \to y \in \mathbb{R} \to x \cdot y = y \cdot x)$\lbl{(R7)}}
\AXC{$\forall x\,\forall y\,(x \in \mathbb{R} \to y \in \mathbb{R} \to x \cdot y = y \cdot x) \to \varphi_{75}$\lbl{A2}}
\BIC{$\varphi_{75}$}
\AXC{$\varphi_{189}$\lbl{A2}}
\BIC{$M + \lvert \beta \rvert + 1 \in \mathbb{R} \to \varepsilon \cdot (M + \lvert \beta \rvert + 1)^{-1} \in \mathbb{R} \to \varphi_{88}$}
\AXC{$\varphi_{267}$\lbl{A1}}
\BIC{$\varphi_{264}$}
\end{ptree}
\noindent$\Pi_{58}$:
\begin{ptree}
\AXC{$\stackrel{\Pi_{52}}{\vdots}$}\noLine
\UIC{$\mathrm{Conv}(a,\alpha) \to \varphi_{204}$}
\AXC{$\stackrel{\Pi_{48}}{\vdots}$}\noLine
\UIC{$\mathrm{Conv}(a,\alpha) \to \varphi_{196}$}
\AXC{$\stackrel{\Pi_{57}}{\vdots}$}\noLine
\UIC{$\varphi_{264}$}
\BIC{$(\mathrm{Conv}(a,\alpha) \to \varphi_{204}) \to \mathrm{Conv}(a,\alpha) \to \varphi_{226}$}
\BIC{$\mathrm{Conv}(a,\alpha) \to \varphi_{226}$}
\AXC{$(\mathrm{Conv}(a,\alpha) \to \varphi_{226}) \to \varphi_{259}$\lbl{A1}}
\BIC{$\varphi_{259}$}
\end{ptree}
\noindent$\Pi_{59}$:
\begin{ptree}
\AXC{$\stackrel{\Pi_{40}}{\vdots}$}\noLine
\UIC{$\mathrm{Conv}(a,\alpha) \to \varphi_{53}$}
\AXC{$\stackrel{\Pi_{48}}{\vdots}$}\noLine
\UIC{$\mathrm{Conv}(a,\alpha) \to \varphi_{196}$}
\AXC{$\stackrel{\Pi_{56}}{\vdots}$}\noLine
\UIC{$(\mathrm{Conv}(a,\alpha) \to \varphi_{196}) \to \varphi_{235}$}
\BIC{$\varphi_{235}$}
\BIC{$\mathrm{Conv}(a,\alpha) \to \mathrm{Conv}(b,\beta) \to \varphi_{209}$}
\AXC{$\stackrel{\Pi_{58}}{\vdots}$}\noLine
\UIC{$\varphi_{259}$}
\BIC{$(\varphi_{88} \to \varphi_{97}) \to \varphi_{222}$}
\end{ptree}
\noindent$\Pi_{60}$:
\begin{ptree}
\AXC{$\stackrel{\Pi_{55}}{\vdots}$}\noLine
\UIC{$\varphi_{101}$}
\AXC{$\varphi_{101} \to \varphi_{159}$\lbl{A2}}
\BIC{$\varphi_{159}$}
\AXC{$\varphi_{159} \to \varphi_{88} \to \varphi_{97}$\lbl{A2}}
\BIC{$\varphi_{88} \to \varphi_{97}$}
\AXC{$\stackrel{\Pi_{59}}{\vdots}$}\noLine
\UIC{$(\varphi_{88} \to \varphi_{97}) \to \varphi_{222}$}
\BIC{$\varphi_{222}$}
\AXC{$\varphi_{222} \to \varphi_{261}$\lbl{A1}}
\BIC{$\varphi_{261}$}
\end{ptree}
\noindent$\Pi_{61}$:
\begin{ptree}
\AXC{$\stackrel{\Pi_{7}}{\vdots}$}\noLine
\UIC{$\mathrm{Conv}(a,\alpha) \to \varphi_{200}$}
\AXC{$\stackrel{\Pi_{46}}{\vdots}$}\noLine
\UIC{$\mathrm{Conv}(a,\alpha) \to \mathrm{Conv}(b,\beta) \to \varphi_{93}$}
\AXC{$\stackrel{\Pi_{54}}{\vdots}$}\noLine
\UIC{$\varphi_{257}$}
\BIC{$(\mathrm{Conv}(a,\alpha) \to \varphi_{200}) \to \mathrm{Conv}(a,\alpha) \to \varphi_{225}$}
\BIC{$\mathrm{Conv}(a,\alpha) \to \varphi_{225}$}
\AXC{$\stackrel{\Pi_{60}}{\vdots}$}\noLine
\UIC{$\varphi_{261}$}
\BIC{$\varphi_{154} \to \mathrm{Conv}(a,\alpha) \to \mathrm{Conv}(b,\beta) \to \varphi_{205}$}
\end{ptree}
\noindent$\Pi_{62}$:
\begin{ptree}
\AXC{$z = w \to \varphi_{89}$\lbl{A6}}
\UIC{$\forall w\,(z = w \to \varphi_{89})$}
\UIC{$\forall z\,\forall w\,(z = w \to \varphi_{89})$}
\AXC{$\forall z\,\forall w\,(z = w \to \varphi_{89}) \to \varphi_{153}$\lbl{A2}}
\BIC{$\varphi_{153}$}
\AXC{$\varphi_{153} \to \varphi_{154}$\lbl{A2}}
\BIC{$\varphi_{154}$}
\AXC{$\stackrel{\Pi_{61}}{\vdots}$}\noLine
\UIC{$\varphi_{154} \to \mathrm{Conv}(a,\alpha) \to \mathrm{Conv}(b,\beta) \to \varphi_{205}$}
\BIC{$\mathrm{Conv}(a,\alpha) \to \mathrm{Conv}(b,\beta) \to \varphi_{205}$}
\end{ptree}
\noindent$\Pi_{63}$:
\begin{ptree}
\AXC{$\mathrm{Conv}(a,\alpha) \to \varphi_{228}$\lbl{A1}}
\AXC{$(\mathrm{Conv}(a,\alpha) \to \varphi_{228}) \to \mathrm{Conv}(a,\alpha) \to \varphi_{230}$\lbl{A1}}
\BIC{$\mathrm{Conv}(a,\alpha) \to \varphi_{230}$}
\AXC{$(\mathrm{Conv}(a,\alpha) \to \varphi_{230}) \to \varphi_{248}$\lbl{A1}}
\BIC{$\varphi_{248}$}
\end{ptree}
\noindent$\Pi_{64}$:
\begin{ptree}
\AXC{$\mathrm{Conv}(a,\alpha) \to \varphi_{178}$\lbl{A1}}
\AXC{$\forall n \in \mathbb{N}\,(N \leq n \to \varphi_{94}) \to n \in \mathbb{N} \to N \leq n \to \varphi_{94}$\lbl{A2}}
\AXC{$\stackrel{\Pi_{63}}{\vdots}$}\noLine
\UIC{$\varphi_{248}$}
\BIC{$\mathrm{Conv}(a,\alpha) \to \varphi_{229}$}
\AXC{$(\mathrm{Conv}(a,\alpha) \to \varphi_{229}) \to \varphi_{262}$\lbl{A1}}
\BIC{$\varphi_{262}$}
\BIC{$(\mathrm{Conv}(a,\alpha) \to \varphi_{177}) \to \mathrm{Conv}(a,\alpha) \to \varphi_{227}$}
\end{ptree}
\noindent$\Pi_{65}$:
\begin{ptree}
\AXC{$\mathrm{Conv}(a,\alpha) \to \varphi_{177}$\lbl{A1}}
\AXC{$\stackrel{\Pi_{64}}{\vdots}$}\noLine
\UIC{$(\mathrm{Conv}(a,\alpha) \to \varphi_{177}) \to \mathrm{Conv}(a,\alpha) \to \varphi_{227}$}
\BIC{$\mathrm{Conv}(a,\alpha) \to \varphi_{227}$}
\end{ptree}
\noindent$\Pi_{66}$:
\begin{ptree}
\AXC{$\mathrm{Conv}(a,\alpha) \to \mathrm{Conv}(b,\beta) \to \varphi_{73}$\lbl{A1}}
\AXC{$\varphi_{164}$\lbl{A1}}
\BIC{$\mathrm{Conv}(a,\alpha) \to \mathrm{Conv}(b,\beta) \to \varphi_{44}$}
\AXC{$(\mathrm{Conv}(a,\alpha) \to \mathrm{Conv}(b,\beta) \to \varphi_{44}) \to \varphi_{117}$\lbl{A1}}
\BIC{$\varphi_{117}$}
\end{ptree}
\noindent$\Pi_{67}$:
\begin{ptree}
\AXC{$\mathrm{Conv}(a,\alpha) \to \varphi_{178}$\lbl{A1}}
\AXC{$\forall n \in \mathbb{N}\,(\lvert a_{n} \rvert \leq M) \to n \in \mathbb{N} \to \lvert a_{n} \rvert \leq M$\lbl{A2}}
\AXC{$\stackrel{\Pi_{66}}{\vdots}$}\noLine
\UIC{$\varphi_{117}$}
\BIC{$\mathrm{Conv}(a,\alpha) \to \varphi_{65}$}
\AXC{$\varphi_{246}$\lbl{A1}}
\BIC{$(\mathrm{Conv}(a,\alpha) \to \varphi_{178}) \to \mathrm{Conv}(a,\alpha) \to \varphi_{184}$}
\BIC{$\mathrm{Conv}(a,\alpha) \to \varphi_{184}$}
\end{ptree}
\noindent$\Pi_{68}$:
\begin{ptree}
\AXC{$\forall a\,\forall \alpha\,(\mathrm{Conv}(a,\alpha) \to \alpha \in \mathbb{R})$\lbl{[C2]}}
\AXC{$\varphi_{19}$\lbl{A2}}
\BIC{$\forall \alpha\,(\mathrm{Conv}(a,\alpha) \to \alpha \in \mathbb{R})$}
\AXC{$\varphi_{11}$\lbl{A2}}
\BIC{$\mathrm{Conv}(a,\alpha) \to \alpha \in \mathbb{R}$}
\AXC{$(\mathrm{Conv}(a,\alpha) \to \alpha \in \mathbb{R}) \to \varphi_{21}$\lbl{A1}}
\BIC{$\varphi_{21}$}
\end{ptree}
\noindent$\Pi_{69}$:
\begin{ptree}
\AXC{$\mathrm{Conv}(a,\alpha) \to \mathrm{Conv}(a,\alpha)$\lbl{A1}}
\AXC{$\stackrel{\Pi_{68}}{\vdots}$}\noLine
\UIC{$\varphi_{21}$}
\BIC{$\mathrm{Conv}(a,\alpha) \to \mathrm{Conv}(b,\beta) \to \alpha \in \mathbb{R}$}
\end{ptree}
\noindent$\Pi_{70}$:
\begin{ptree}
\AXC{$\forall a\,\forall \alpha\,(\mathrm{Conv}(a,\alpha) \to \mathrm{Seq}(a))$\lbl{[C1]}}
\AXC{$\varphi_{7}$\lbl{A2}}
\BIC{$\forall \alpha\,(\mathrm{Conv}(b,\alpha) \to \mathrm{Seq}(b))$}
\AXC{$\varphi_{2}$\lbl{A2}}
\BIC{$\mathrm{Conv}(b,\beta) \to \mathrm{Seq}(b)$}
\AXC{$(\mathrm{Conv}(b,\beta) \to \mathrm{Seq}(b)) \to \varphi_{31}$\lbl{A1}}
\BIC{$\varphi_{31}$}
\end{ptree}
\noindent$\Pi_{71}$:
\begin{ptree}
\AXC{$\mathrm{Conv}(a,\alpha) \to \mathrm{Conv}(b,\beta) \to \mathrm{Conv}(b,\beta)$\lbl{A1}}
\AXC{$\stackrel{\Pi_{70}}{\vdots}$}\noLine
\UIC{$\varphi_{31}$}
\BIC{$\mathrm{Conv}(a,\alpha) \to \mathrm{Conv}(b,\beta) \to \mathrm{Seq}(b)$}
\end{ptree}
\noindent$\Pi_{72}$:
\begin{ptree}
\AXC{$\forall a\,\forall n\,(\mathrm{Seq}(a) \to n \in \mathbb{N} \to a_{n} \in \mathbb{R})$\lbl{[S1]}}
\AXC{$\varphi_{71}$\lbl{A2}}
\BIC{$\forall n\,(\mathrm{Seq}(b) \to n \in \mathbb{N} \to b_{n} \in \mathbb{R})$}
\AXC{$\varphi_{59}$\lbl{A2}}
\BIC{$\mathrm{Seq}(b) \to n \in \mathbb{N} \to b_{n} \in \mathbb{R}$}
\AXC{$(\mathrm{Seq}(b) \to n \in \mathbb{N} \to b_{n} \in \mathbb{R}) \to \varphi_{241}$\lbl{A1}}
\BIC{$\varphi_{241}$}
\end{ptree}
\noindent$\Pi_{73}$:
\begin{ptree}
\AXC{$\mathrm{Conv}(a,\alpha) \to \varphi_{178}$\lbl{A1}}
\AXC{$\stackrel{\Pi_{71}}{\vdots}$}\noLine
\UIC{$\mathrm{Conv}(a,\alpha) \to \mathrm{Conv}(b,\beta) \to \mathrm{Seq}(b)$}
\AXC{$\stackrel{\Pi_{72}}{\vdots}$}\noLine
\UIC{$\varphi_{241}$}
\BIC{$(\mathrm{Conv}(a,\alpha) \to \varphi_{178}) \to \mathrm{Conv}(a,\alpha) \to \varphi_{187}$}
\BIC{$\mathrm{Conv}(a,\alpha) \to \varphi_{187}$}
\end{ptree}
\noindent$\Pi_{74}$:
\begin{ptree}
\AXC{$\forall a\,\forall n\,(\mathrm{Seq}(a) \to n \in \mathbb{N} \to a_{n} \in \mathbb{R})$\lbl{[S1]}}
\AXC{$\varphi_{68}$\lbl{A2}}
\BIC{$\forall n\,(\mathrm{Seq}(a) \to n \in \mathbb{N} \to a_{n} \in \mathbb{R})$}
\AXC{$\varphi_{60}$\lbl{A2}}
\BIC{$\mathrm{Seq}(a) \to n \in \mathbb{N} \to a_{n} \in \mathbb{R}$}
\AXC{$(\mathrm{Seq}(a) \to n \in \mathbb{N} \to a_{n} \in \mathbb{R}) \to \varphi_{242}$\lbl{A1}}
\BIC{$\varphi_{242}$}
\end{ptree}
\noindent$\Pi_{75}$:
\begin{ptree}
\AXC{$\mathrm{Conv}(a,\alpha) \to \varphi_{178}$\lbl{A1}}
\AXC{$\stackrel{\Pi_{16}}{\vdots}$}\noLine
\UIC{$\mathrm{Conv}(a,\alpha) \to \mathrm{Conv}(b,\beta) \to \mathrm{Seq}(a)$}
\AXC{$\stackrel{\Pi_{74}}{\vdots}$}\noLine
\UIC{$\varphi_{242}$}
\BIC{$(\mathrm{Conv}(a,\alpha) \to \varphi_{178}) \to \mathrm{Conv}(a,\alpha) \to \varphi_{188}$}
\BIC{$\mathrm{Conv}(a,\alpha) \to \varphi_{188}$}
\end{ptree}
\noindent$\Pi_{76}$:
\begin{ptree}
\AXC{$\varphi_{142}$\lbl{[C15]}}
\AXC{$\varphi_{142} \to \varphi_{167}$\lbl{A2}}
\BIC{$\varphi_{167}$}
\AXC{$\varphi_{167} \to \forall \alpha\,\forall \beta\,\forall M\,(\varphi_{170})$\lbl{A2}}
\BIC{$\forall \alpha\,\forall \beta\,\forall M\,(\varphi_{170})$}
\AXC{$\forall \alpha\,\forall \beta\,\forall M\,(\varphi_{170}) \to \forall \beta\,\forall M\,(\varphi_{170})$\lbl{A2}}
\BIC{$\forall \beta\,\forall M\,(\varphi_{170})$}
\end{ptree}
\noindent$\Pi_{77}$:
\begin{ptree}
\AXC{$\stackrel{\Pi_{76}}{\vdots}$}\noLine
\UIC{$\forall \beta\,\forall M\,(\varphi_{170})$}
\AXC{$\forall \beta\,\forall M\,(\varphi_{170}) \to \forall M\,(\varphi_{170})$\lbl{A2}}
\BIC{$\forall M\,(\varphi_{170})$}
\AXC{$\forall M\,(\varphi_{170}) \to \varphi_{170}$\lbl{A2}}
\BIC{$\varphi_{170}$}
\AXC{$\varphi_{170} \to a_{n} \in \mathbb{R} \to b_{n} \in \mathbb{R} \to \alpha \in \mathbb{R} \to \beta \in \mathbb{R} \to \varphi_{195}$\lbl{A2}}
\BIC{$a_{n} \in \mathbb{R} \to b_{n} \in \mathbb{R} \to \alpha \in \mathbb{R} \to \beta \in \mathbb{R} \to \varphi_{195}$}
\end{ptree}
\noindent$\Pi_{78}$:
\begin{ptree}
\AXC{$\stackrel{\Pi_{73}}{\vdots}$}\noLine
\UIC{$\mathrm{Conv}(a,\alpha) \to \varphi_{187}$}
\AXC{$\stackrel{\Pi_{75}}{\vdots}$}\noLine
\UIC{$\mathrm{Conv}(a,\alpha) \to \varphi_{188}$}
\AXC{$\stackrel{\Pi_{77}}{\vdots}$}\noLine
\UIC{$a_{n} \in \mathbb{R} \to b_{n} \in \mathbb{R} \to \alpha \in \mathbb{R} \to \beta \in \mathbb{R} \to \varphi_{195}$}
\AXC{$\varphi_{273}$\lbl{A1}}
\BIC{$(\mathrm{Conv}(a,\alpha) \to \varphi_{188}) \to \varphi_{272}$}
\BIC{$\varphi_{272}$}
\BIC{$(\mathrm{Conv}(a,\alpha) \to \mathrm{Conv}(b,\beta) \to \alpha \in \mathbb{R}) \to \varphi_{271}$}
\end{ptree}
\noindent$\Pi_{79}$:
\begin{ptree}
\AXC{$\stackrel{\Pi_{69}}{\vdots}$}\noLine
\UIC{$\mathrm{Conv}(a,\alpha) \to \mathrm{Conv}(b,\beta) \to \alpha \in \mathbb{R}$}
\AXC{$\stackrel{\Pi_{78}}{\vdots}$}\noLine
\UIC{$(\mathrm{Conv}(a,\alpha) \to \mathrm{Conv}(b,\beta) \to \alpha \in \mathbb{R}) \to \varphi_{271}$}
\BIC{$\varphi_{271}$}
\end{ptree}
\noindent$\Pi_{80}$:
\begin{ptree}
\AXC{$\mathrm{Conv}(a,\alpha) \to \mathrm{Conv}(b,\beta) \to \varphi_{73}$\lbl{A1}}
\AXC{$\varphi_{144}$\lbl{A1}}
\BIC{$\mathrm{Conv}(a,\alpha) \to \varphi_{29}$}
\AXC{$\stackrel{\Pi_{2}}{\vdots}$}\noLine
\UIC{$\mathrm{Conv}(a,\alpha) \to \mathrm{Conv}(b,\beta) \to \beta \in \mathbb{R}$}
\AXC{$\stackrel{\Pi_{79}}{\vdots}$}\noLine
\UIC{$\varphi_{271}$}
\BIC{$(\mathrm{Conv}(a,\alpha) \to \varphi_{29}) \to \varphi_{269}$}
\BIC{$\varphi_{269}$}
\end{ptree}
\noindent$\Pi_{81}$:
\begin{ptree}
\AXC{$\stackrel{\Pi_{67}}{\vdots}$}\noLine
\UIC{$\mathrm{Conv}(a,\alpha) \to \varphi_{184}$}
\AXC{$\stackrel{\Pi_{23}}{\vdots}$}\noLine
\UIC{$\mathrm{Conv}(a,\alpha) \to \varphi_{20}$}
\AXC{$\stackrel{\Pi_{52}}{\vdots}$}\noLine
\UIC{$\mathrm{Conv}(a,\alpha) \to \varphi_{204}$}
\AXC{$\stackrel{\Pi_{80}}{\vdots}$}\noLine
\UIC{$\varphi_{269}$}
\BIC{$(\mathrm{Conv}(a,\alpha) \to \varphi_{20}) \to (\mathrm{Conv}(a,\alpha) \to \varphi_{184}) \to \varphi_{265}$}
\BIC{$(\mathrm{Conv}(a,\alpha) \to \varphi_{184}) \to \varphi_{265}$}
\BIC{$\varphi_{265}$}
\end{ptree}
\noindent$\Pi_{82}$:
\begin{ptree}
\AXC{$\stackrel{\Pi_{65}}{\vdots}$}\noLine
\UIC{$\mathrm{Conv}(a,\alpha) \to \varphi_{227}$}
\AXC{$(\mathrm{Conv}(a,\alpha) \to \varphi_{227}) \to \mathrm{Conv}(a,\alpha) \to \varphi_{216}$\lbl{A1}}
\BIC{$\mathrm{Conv}(a,\alpha) \to \varphi_{216}$}
\AXC{$\stackrel{\Pi_{81}}{\vdots}$}\noLine
\UIC{$\varphi_{265}$}
\BIC{$(\mathrm{Conv}(a,\alpha) \to \varphi_{215}) \to \mathrm{Conv}(a,\alpha) \to \varphi_{224}$}
\end{ptree}
\noindent$\Pi_{83}$:
\begin{ptree}
\AXC{$\stackrel{\Pi_{65}}{\vdots}$}\noLine
\UIC{$\mathrm{Conv}(a,\alpha) \to \varphi_{227}$}
\AXC{$(\mathrm{Conv}(a,\alpha) \to \varphi_{227}) \to \mathrm{Conv}(a,\alpha) \to \varphi_{215}$\lbl{A1}}
\BIC{$\mathrm{Conv}(a,\alpha) \to \varphi_{215}$}
\AXC{$\stackrel{\Pi_{82}}{\vdots}$}\noLine
\UIC{$(\mathrm{Conv}(a,\alpha) \to \varphi_{215}) \to \mathrm{Conv}(a,\alpha) \to \varphi_{224}$}
\BIC{$\mathrm{Conv}(a,\alpha) \to \varphi_{224}$}
\end{ptree}
\noindent$\Pi_{84}$:
\begin{ptree}
\AXC{$\forall x\,\forall y\,\forall z\,(x \leq y \to y < z \to x < z)$\lbl{[O4]}}
\AXC{$\forall x\,\forall y\,\forall z\,(x \leq y \to y < z \to x < z) \to \varphi_{82}$\lbl{A2}}
\BIC{$\varphi_{82}$}
\AXC{$\varphi_{82} \to \varphi_{136}$\lbl{A2}}
\BIC{$\varphi_{136}$}
\AXC{$\varphi_{136} \to \varphi_{63} \to \varphi_{85}$\lbl{A2}}
\BIC{$\varphi_{63} \to \varphi_{85}$}
\end{ptree}
\noindent$\Pi_{85}$:
\begin{ptree}
\AXC{$\stackrel{\Pi_{83}}{\vdots}$}\noLine
\UIC{$\mathrm{Conv}(a,\alpha) \to \varphi_{224}$}
\AXC{$\stackrel{\Pi_{84}}{\vdots}$}\noLine
\UIC{$\varphi_{63} \to \varphi_{85}$}
\AXC{$(\varphi_{63} \to \varphi_{85}) \to (\mathrm{Conv}(a,\alpha) \to \varphi_{224}) \to \varphi_{244}$\lbl{A1}}
\BIC{$(\mathrm{Conv}(a,\alpha) \to \varphi_{224}) \to \varphi_{244}$}
\BIC{$\varphi_{244}$}
\end{ptree}
\noindent$\Pi_{86}$:
\begin{ptree}
\AXC{$\forall a\,\forall b\,(\mathrm{Seq}(a) \to \varphi_{61})$\lbl{[S4]}}
\AXC{$\forall a\,\forall b\,(\mathrm{Seq}(a) \to \varphi_{61}) \to \forall b\,(\mathrm{Seq}(a) \to \varphi_{61})$\lbl{A2}}
\BIC{$\forall b\,(\mathrm{Seq}(a) \to \varphi_{61})$}
\end{ptree}
\noindent$\Pi_{87}$:
\begin{ptree}
\AXC{$\stackrel{\Pi_{86}}{\vdots}$}\noLine
\UIC{$\forall b\,(\mathrm{Seq}(a) \to \varphi_{61})$}
\AXC{$\forall b\,(\mathrm{Seq}(a) \to \varphi_{61}) \to \mathrm{Seq}(a) \to \varphi_{61}$\lbl{A2}}
\BIC{$\mathrm{Seq}(a) \to \varphi_{61}$}
\AXC{$\varphi_{186}$\lbl{A1}}
\BIC{$(\mathrm{Conv}(a,\alpha) \to \mathrm{Conv}(b,\beta) \to \mathrm{Seq}(a)) \to \varphi_{110}$}
\end{ptree}
\noindent$\Pi_{88}$:
\begin{ptree}
\AXC{$\stackrel{\Pi_{16}}{\vdots}$}\noLine
\UIC{$\mathrm{Conv}(a,\alpha) \to \mathrm{Conv}(b,\beta) \to \mathrm{Seq}(a)$}
\AXC{$\stackrel{\Pi_{87}}{\vdots}$}\noLine
\UIC{$(\mathrm{Conv}(a,\alpha) \to \mathrm{Conv}(b,\beta) \to \mathrm{Seq}(a)) \to \varphi_{110}$}
\BIC{$\varphi_{110}$}
\end{ptree}
\noindent$\Pi_{89}$:
\begin{ptree}
\AXC{$\stackrel{\Pi_{71}}{\vdots}$}\noLine
\UIC{$\mathrm{Conv}(a,\alpha) \to \mathrm{Conv}(b,\beta) \to \mathrm{Seq}(b)$}
\AXC{$\stackrel{\Pi_{88}}{\vdots}$}\noLine
\UIC{$\varphi_{110}$}
\BIC{$\mathrm{Conv}(a,\alpha) \to \varphi_{66}$}
\end{ptree}
\noindent$\Pi_{90}$:
\begin{ptree}
\AXC{$\stackrel{\Pi_{89}}{\vdots}$}\noLine
\UIC{$\mathrm{Conv}(a,\alpha) \to \varphi_{66}$}
\AXC{$(\mathrm{Conv}(a,\alpha) \to \varphi_{66}) \to \mathrm{Conv}(a,\alpha) \to \varphi_{217}$\lbl{A1}}
\BIC{$\mathrm{Conv}(a,\alpha) \to \varphi_{217}$}
\end{ptree}
\noindent$\Pi_{91}$:
\begin{ptree}
\AXC{$\varphi_{95}$\lbl{A2}}
\AXC{$\stackrel{\Pi_{90}}{\vdots}$}\noLine
\UIC{$\mathrm{Conv}(a,\alpha) \to \varphi_{217}$}
\AXC{$(\mathrm{Conv}(a,\alpha) \to \varphi_{217}) \to \varphi_{95} \to \mathrm{Conv}(a,\alpha) \to \varphi_{214}$\lbl{A1}}
\BIC{$\varphi_{95} \to \mathrm{Conv}(a,\alpha) \to \varphi_{214}$}
\BIC{$\mathrm{Conv}(a,\alpha) \to \varphi_{214}$}
\AXC{$\varphi_{260}$\lbl{A1}}
\BIC{$(\mathrm{Conv}(a,\alpha) \to \varphi_{178}) \to \mathrm{Conv}(a,\alpha) \to \varphi_{206}$}
\end{ptree}
\noindent$\Pi_{92}$:
\begin{ptree}
\AXC{$\mathrm{Conv}(a,\alpha) \to \varphi_{178}$\lbl{A1}}
\AXC{$\stackrel{\Pi_{91}}{\vdots}$}\noLine
\UIC{$(\mathrm{Conv}(a,\alpha) \to \varphi_{178}) \to \mathrm{Conv}(a,\alpha) \to \varphi_{206}$}
\BIC{$\mathrm{Conv}(a,\alpha) \to \varphi_{206}$}
\AXC{$(\mathrm{Conv}(a,\alpha) \to \varphi_{206}) \to \varphi_{233}$\lbl{A1}}
\BIC{$\varphi_{233}$}
\end{ptree}
\noindent$\Pi_{93}$:
\begin{ptree}
\AXC{$\stackrel{\Pi_{37}}{\vdots}$}\noLine
\UIC{$\forall x\,\forall y\,(x = y \to y = x)$}
\AXC{$\varphi_{23}$\lbl{A2}}
\BIC{$\forall y\,(a \otimes b(n) = y \to y = a \otimes b(n))$}
\AXC{$\varphi_{107}$\lbl{A2}}
\BIC{$a \otimes b(n) = a_{n} \cdot b_{n} \to a_{n} \cdot b_{n} = a \otimes b(n)$}
\AXC{$\stackrel{\Pi_{92}}{\vdots}$}\noLine
\UIC{$\varphi_{233}$}
\BIC{$\mathrm{Conv}(a,\alpha) \to \varphi_{207}$}
\AXC{$\varphi_{266}$\lbl{A1}}
\BIC{$(\mathrm{Conv}(a,\alpha) \to \varphi_{212}) \to \varphi_{108} \to \mathrm{Conv}(a,\alpha) \to \varphi_{210}$}
\end{ptree}
\noindent$\Pi_{94}$:
\begin{ptree}
\AXC{$\stackrel{\Pi_{62}}{\vdots}$}\noLine
\UIC{$\mathrm{Conv}(a,\alpha) \to \mathrm{Conv}(b,\beta) \to \varphi_{205}$}
\AXC{$\stackrel{\Pi_{85}}{\vdots}$}\noLine
\UIC{$\varphi_{244}$}
\BIC{$\mathrm{Conv}(a,\alpha) \to \varphi_{212}$}
\AXC{$\stackrel{\Pi_{93}}{\vdots}$}\noLine
\UIC{$(\mathrm{Conv}(a,\alpha) \to \varphi_{212}) \to \varphi_{108} \to \mathrm{Conv}(a,\alpha) \to \varphi_{210}$}
\BIC{$\varphi_{108} \to \mathrm{Conv}(a,\alpha) \to \varphi_{210}$}
\end{ptree}
\noindent$\Pi_{95}$:
\begin{ptree}
\AXC{$z = w \to \lvert z - \alpha \cdot \beta \rvert < \varepsilon \to \lvert w - \alpha \cdot \beta \rvert < \varepsilon$\lbl{A6}}
\UIC{$\forall w\,(z = w \to \lvert z - \alpha \cdot \beta \rvert < \varepsilon \to \lvert w - \alpha \cdot \beta \rvert < \varepsilon)$}
\UIC{$\forall z\,\forall w\,(z = w \to \lvert z - \alpha \cdot \beta \rvert < \varepsilon \to \lvert w - \alpha \cdot \beta \rvert < \varepsilon)$}
\AXC{$\varphi_{120}$\lbl{A2}}
\BIC{$\varphi_{86}$}
\AXC{$\varphi_{86} \to \varphi_{108}$\lbl{A2}}
\BIC{$\varphi_{108}$}
\AXC{$\stackrel{\Pi_{94}}{\vdots}$}\noLine
\UIC{$\varphi_{108} \to \mathrm{Conv}(a,\alpha) \to \varphi_{210}$}
\BIC{$\mathrm{Conv}(a,\alpha) \to \varphi_{210}$}
\end{ptree}
\noindent$\Pi_{96}$:
\begin{ptree}
\AXC{$\varphi_{243}$\lbl{A3}}
\AXC{$\stackrel{\Pi_{95}}{\vdots}$}\noLine
\UIC{$\mathrm{Conv}(a,\alpha) \to \varphi_{210}$}
\AXC{$(\mathrm{Conv}(a,\alpha) \to \varphi_{210}) \to \varphi_{218}$\lbl{A1}}
\BIC{$\varphi_{218}$}
\UIC{$\forall n\,(\varphi_{218})$}
\AXC{$\forall n\,(\varphi_{218}) \to \varphi_{243} \to \varphi_{236} \to \mathrm{Conv}(a,\alpha) \to \varphi_{213}$\lbl{A1}}
\BIC{$\varphi_{243} \to \varphi_{236} \to \mathrm{Conv}(a,\alpha) \to \varphi_{213}$}
\BIC{$\varphi_{236} \to \mathrm{Conv}(a,\alpha) \to \varphi_{213}$}
\end{ptree}
\noindent$\Pi_{97}$:
\begin{ptree}
\AXC{$\mathrm{Conv}(a,\alpha) \to \varphi_{228}$\lbl{A1}}
\AXC{$(\mathrm{Conv}(a,\alpha) \to \varphi_{228}) \to \varphi_{179}$\lbl{A1}}
\BIC{$\varphi_{179}$}
\AXC{$\varphi_{179} \to (\mathrm{Conv}(a,\alpha) \to \varphi_{213}) \to \mathrm{Conv}(a,\alpha) \to \varphi_{219}$\lbl{A1}}
\BIC{$(\mathrm{Conv}(a,\alpha) \to \varphi_{213}) \to \mathrm{Conv}(a,\alpha) \to \varphi_{219}$}
\end{ptree}
\noindent$\Pi_{98}$:
\begin{ptree}
\AXC{$\varphi_{236}$\lbl{A4}}
\AXC{$\stackrel{\Pi_{96}}{\vdots}$}\noLine
\UIC{$\varphi_{236} \to \mathrm{Conv}(a,\alpha) \to \varphi_{213}$}
\BIC{$\mathrm{Conv}(a,\alpha) \to \varphi_{213}$}
\AXC{$\stackrel{\Pi_{97}}{\vdots}$}\noLine
\UIC{$(\mathrm{Conv}(a,\alpha) \to \varphi_{213}) \to \mathrm{Conv}(a,\alpha) \to \varphi_{219}$}
\BIC{$\mathrm{Conv}(a,\alpha) \to \varphi_{219}$}
\AXC{$\varphi_{258}$\lbl{A1}}
\BIC{$(\forall N\,(\varphi_{43}) \to \varphi_{43}) \to \mathrm{Conv}(a,\alpha) \to \varphi_{220}$}
\end{ptree}
\noindent$\Pi_{99}$:
\begin{ptree}
\AXC{$\forall x\,\forall y\,(x \in \mathbb{R} \to y \in \mathbb{R} \to x \cdot y \in \mathbb{R})$\lbl{(R1)}}
\AXC{$\varphi_{80}$\lbl{A2}}
\BIC{$\forall y\,(\varepsilon \in \mathbb{R} \to y \in \mathbb{R} \to \varepsilon \cdot y \in \mathbb{R})$}
\AXC{$\forall y\,(\varepsilon \in \mathbb{R} \to y \in \mathbb{R} \to \varepsilon \cdot y \in \mathbb{R}) \to \varepsilon \in \mathbb{R} \to \varphi_{33}$\lbl{A2}}
\BIC{$\varepsilon \in \mathbb{R} \to \varphi_{33}$}
\end{ptree}
\noindent$\Pi_{100}$:
\begin{ptree}
\AXC{$\mathrm{Conv}(a,\alpha) \to \varphi_{64}$\lbl{A1}}
\AXC{$\stackrel{\Pi_{99}}{\vdots}$}\noLine
\UIC{$\varepsilon \in \mathbb{R} \to \varphi_{33}$}
\AXC{$(\varepsilon \in \mathbb{R} \to \varphi_{33}) \to (\mathrm{Conv}(a,\alpha) \to \varphi_{64}) \to \varphi_{192}$\lbl{A1}}
\BIC{$(\mathrm{Conv}(a,\alpha) \to \varphi_{64}) \to \varphi_{192}$}
\BIC{$\varphi_{192}$}
\end{ptree}
\noindent$\Pi_{101}$:
\begin{ptree}
\AXC{$\forall a\,\forall \alpha\,\forall b\,\forall \beta\,(\varphi_{114})$\lbl{[C9]}}
\AXC{$\forall a\,\forall \alpha\,\forall b\,\forall \beta\,(\varphi_{114}) \to \forall \alpha\,\forall b\,\forall \beta\,(\varphi_{114})$\lbl{A2}}
\BIC{$\forall \alpha\,\forall b\,\forall \beta\,(\varphi_{114})$}
\AXC{$\forall \alpha\,\forall b\,\forall \beta\,(\varphi_{114}) \to \forall b\,\forall \beta\,(\varphi_{114})$\lbl{A2}}
\BIC{$\forall b\,\forall \beta\,(\varphi_{114})$}
\end{ptree}
\noindent$\Pi_{102}$:
\begin{ptree}
\AXC{$\stackrel{\Pi_{101}}{\vdots}$}\noLine
\UIC{$\forall b\,\forall \beta\,(\varphi_{114})$}
\AXC{$\forall b\,\forall \beta\,(\varphi_{114}) \to \forall \beta\,(\varphi_{114})$\lbl{A2}}
\BIC{$\forall \beta\,(\varphi_{114})$}
\AXC{$\forall \beta\,(\varphi_{114}) \to \varphi_{114}$\lbl{A2}}
\BIC{$\varphi_{114}$}
\AXC{$\varphi_{114} \to \varphi_{175}$\lbl{A2}}
\BIC{$\varphi_{175}$}
\end{ptree}
\noindent$\Pi_{103}$:
\begin{ptree}
\AXC{$\mathrm{Conv}(a,\alpha) \to \mathrm{Conv}(a,\alpha)$\lbl{A1}}
\AXC{$\stackrel{\Pi_{102}}{\vdots}$}\noLine
\UIC{$\varphi_{175}$}
\AXC{$\varphi_{175} \to (\mathrm{Conv}(a,\alpha) \to \mathrm{Conv}(a,\alpha)) \to \varphi_{239}$\lbl{A1}}
\BIC{$(\mathrm{Conv}(a,\alpha) \to \mathrm{Conv}(a,\alpha)) \to \varphi_{239}$}
\BIC{$\varphi_{239}$}
\end{ptree}
\noindent$\Pi_{104}$:
\begin{ptree}
\AXC{$\stackrel{\Pi_{43}}{\vdots}$}\noLine
\UIC{$\mathrm{Conv}(a,\alpha) \to \mathrm{Conv}(b,\beta) \to \varphi_{90}$}
\AXC{$\stackrel{\Pi_{100}}{\vdots}$}\noLine
\UIC{$\varphi_{192}$}
\BIC{$\mathrm{Conv}(a,\alpha) \to \mathrm{Conv}(b,\beta) \to \varphi_{96}$}
\AXC{$\mathrm{Conv}(a,\alpha) \to \mathrm{Conv}(b,\beta) \to \mathrm{Conv}(b,\beta)$\lbl{A1}}
\AXC{$\stackrel{\Pi_{103}}{\vdots}$}\noLine
\UIC{$\varphi_{239}$}
\BIC{$(\mathrm{Conv}(a,\alpha) \to \mathrm{Conv}(b,\beta) \to \varphi_{96}) \to \varphi_{231}$}
\BIC{$\varphi_{231}$}
\end{ptree}
\noindent$\Pi_{105}$:
\begin{ptree}
\AXC{$\stackrel{\Pi_{46}}{\vdots}$}\noLine
\UIC{$\mathrm{Conv}(a,\alpha) \to \mathrm{Conv}(b,\beta) \to \varphi_{93}$}
\AXC{$\stackrel{\Pi_{104}}{\vdots}$}\noLine
\UIC{$\varphi_{231}$}
\BIC{$\mathrm{Conv}(a,\alpha) \to \mathrm{Conv}(b,\beta) \to \varphi_{161}$}
\AXC{$(\mathrm{Conv}(a,\alpha) \to \mathrm{Conv}(b,\beta) \to \varphi_{161}) \to \varphi_{268}$\lbl{A1}}
\BIC{$\varphi_{268}$}
\end{ptree}
\noindent$\Pi_{106}$:
\begin{ptree}
\AXC{$\varphi_{127} \to \forall N\,(\varphi_{127})$\lbl{A4}}
\AXC{$\varphi_{252}$\lbl{A3}}
\AXC{$\forall N\,(\varphi_{43}) \to \varphi_{43}$\lbl{A2}}
\AXC{$\stackrel{\Pi_{98}}{\vdots}$}\noLine
\UIC{$(\forall N\,(\varphi_{43}) \to \varphi_{43}) \to \mathrm{Conv}(a,\alpha) \to \varphi_{220}$}
\BIC{$\mathrm{Conv}(a,\alpha) \to \varphi_{220}$}
\AXC{$\varphi_{249}$\lbl{A1}}
\BIC{$\varphi_{127} \to \neg (N \in \mathbb{N} \land \forall n \in \mathbb{N}\,(N \leq n \to \varphi_{94}))$}
\UIC{$\forall N\,(\varphi_{127} \to \neg (N \in \mathbb{N} \land \forall n \in \mathbb{N}\,(N \leq n \to \varphi_{94})))$}
\AXC{$\stackrel{\Pi_{105}}{\vdots}$}\noLine
\UIC{$\varphi_{268}$}
\BIC{$\varphi_{252} \to (\varphi_{127} \to \forall N\,(\varphi_{127})) \to \mathrm{Conv}(a,\alpha) \to \varphi_{119}$}
\BIC{$(\varphi_{127} \to \forall N\,(\varphi_{127})) \to \mathrm{Conv}(a,\alpha) \to \varphi_{119}$}
\BIC{$\mathrm{Conv}(a,\alpha) \to \varphi_{119}$}
\end{ptree}
\noindent$\Pi_{107}$:
\begin{ptree}
\AXC{$\stackrel{\Pi_{106}}{\vdots}$}\noLine
\UIC{$\mathrm{Conv}(a,\alpha) \to \varphi_{119}$}
\AXC{$(\mathrm{Conv}(a,\alpha) \to \varphi_{119}) \to \varphi_{28} \to \varepsilon \in \mathbb{R} \to \varphi_{74}$\lbl{A1}}
\BIC{$\varphi_{28} \to \varepsilon \in \mathbb{R} \to \varphi_{74}$}
\UIC{$\forall \varepsilon\,(\varphi_{28} \to \varepsilon \in \mathbb{R} \to \varphi_{74})$}
\AXC{$\forall \varepsilon\,(\varphi_{28} \to \varepsilon \in \mathbb{R} \to \varphi_{74}) \to \varphi_{255}$\lbl{A1}}
\BIC{$\varphi_{255}$}
\end{ptree}
\noindent$\Pi_{108}$:
\begin{ptree}
\AXC{$\forall x\,\forall y\,(x \in \mathbb{R} \to y \in \mathbb{R} \to x \cdot y \in \mathbb{R})$\lbl{(R1)}}
\AXC{$\varphi_{79}$\lbl{A2}}
\BIC{$\forall y\,(\alpha \in \mathbb{R} \to y \in \mathbb{R} \to \alpha \cdot y \in \mathbb{R})$}
\AXC{$\varphi_{78}$\lbl{A2}}
\BIC{$\alpha \in \mathbb{R} \to \beta \in \mathbb{R} \to \alpha \cdot \beta \in \mathbb{R}$}
\AXC{$\varphi_{162}$\lbl{A1}}
\BIC{$(\mathrm{Conv}(a,\alpha) \to \mathrm{Conv}(b,\beta) \to \alpha \in \mathbb{R}) \to \varphi_{106}$}
\end{ptree}
\noindent$\Pi_{109}$:
\begin{ptree}
\AXC{$\stackrel{\Pi_{69}}{\vdots}$}\noLine
\UIC{$\mathrm{Conv}(a,\alpha) \to \mathrm{Conv}(b,\beta) \to \alpha \in \mathbb{R}$}
\AXC{$\stackrel{\Pi_{108}}{\vdots}$}\noLine
\UIC{$(\mathrm{Conv}(a,\alpha) \to \mathrm{Conv}(b,\beta) \to \alpha \in \mathbb{R}) \to \varphi_{106}$}
\BIC{$\varphi_{106}$}
\end{ptree}
\noindent$\Pi_{110}$:
\begin{ptree}
\AXC{$\stackrel{\Pi_{89}}{\vdots}$}\noLine
\UIC{$\mathrm{Conv}(a,\alpha) \to \varphi_{66}$}
\AXC{$(\mathrm{Conv}(a,\alpha) \to \varphi_{66}) \to \mathrm{Conv}(a,\alpha) \to \varphi_{40}$\lbl{A1}}
\BIC{$\mathrm{Conv}(a,\alpha) \to \varphi_{40}$}
\end{ptree}
\noindent$\Pi_{111}$:
\begin{ptree}
\AXC{$\stackrel{\Pi_{110}}{\vdots}$}\noLine
\UIC{$\mathrm{Conv}(a,\alpha) \to \varphi_{40}$}
\AXC{$\varphi_{112}$\lbl{DefConv}}
\AXC{$\varphi_{112} \to \varphi_{118}$\lbl{A2}}
\BIC{$\varphi_{118}$}
\AXC{$\varphi_{118} \to \varphi_{126}$\lbl{A2}}
\BIC{$\varphi_{126}$}
\AXC{$\varphi_{250}$\lbl{A1}}
\BIC{$(\mathrm{Conv}(a,\alpha) \to \varphi_{40}) \to (\mathrm{Conv}(a,\alpha) \to \varphi_{45}) \to \varphi_{202}$}
\BIC{$(\mathrm{Conv}(a,\alpha) \to \varphi_{45}) \to \varphi_{202}$}
\end{ptree}
\noindent$\Pi_{112}$:
\begin{ptree}
\AXC{$\stackrel{\Pi_{2}}{\vdots}$}\noLine
\UIC{$\mathrm{Conv}(a,\alpha) \to \mathrm{Conv}(b,\beta) \to \beta \in \mathbb{R}$}
\AXC{$\stackrel{\Pi_{109}}{\vdots}$}\noLine
\UIC{$\varphi_{106}$}
\BIC{$\mathrm{Conv}(a,\alpha) \to \varphi_{45}$}
\AXC{$\stackrel{\Pi_{111}}{\vdots}$}\noLine
\UIC{$(\mathrm{Conv}(a,\alpha) \to \varphi_{45}) \to \varphi_{202}$}
\BIC{$\varphi_{202}$}
\end{ptree}
\noindent$\Pi_{113}$:
\begin{ptree}
\AXC{$\varphi_{28} \to \forall \varepsilon\,(\varphi_{28})$\lbl{A4}}
\AXC{$\forall \varepsilon\,(\varphi_{28} \to \varepsilon \in \mathbb{R} \to \varphi_{74}) \to \forall \varepsilon\,(\varphi_{28}) \to \forall \varepsilon \in \mathbb{R}\,(\varphi_{74})$\lbl{A3}}
\AXC{$\stackrel{\Pi_{107}}{\vdots}$}\noLine
\UIC{$\varphi_{255}$}
\BIC{$(\varphi_{28} \to \forall \varepsilon\,(\varphi_{28})) \to \mathrm{Conv}(a,\alpha) \to \varphi_{121}$}
\BIC{$\mathrm{Conv}(a,\alpha) \to \varphi_{121}$}
\AXC{$\stackrel{\Pi_{112}}{\vdots}$}\noLine
\UIC{$\varphi_{202}$}
\BIC{$\mathrm{Conv}(a,\alpha) \to \mathrm{Conv}(b,\beta) \to \varphi_{35}$}
\AXC{$\varphi_{148}$\lbl{A1}}
\BIC{$\varphi_{9} \to \neg (M \in \mathbb{R} \land \forall n \in \mathbb{N}\,(\lvert a_{n} \rvert \leq M))$}
\UIC{$\forall M\,(\varphi_{9} \to \neg (M \in \mathbb{R} \land \forall n \in \mathbb{N}\,(\lvert a_{n} \rvert \leq M)))$}
\end{ptree}
\noindent$\Pi_{114}$:
\begin{ptree}
\AXC{$\forall a\,\forall \alpha\,(\mathrm{Conv}(a,\alpha) \to \mathrm{SeqBdd}(a))$\lbl{定理2.11}}
\AXC{$\varphi_{13}$\lbl{A2}}
\BIC{$\forall \alpha\,(\mathrm{Conv}(a,\alpha) \to \mathrm{SeqBdd}(a))$}
\AXC{$\varphi_{5}$\lbl{A2}}
\BIC{$\mathrm{Conv}(a,\alpha) \to \mathrm{SeqBdd}(a)$}
\AXC{$(\mathrm{Conv}(a,\alpha) \to \mathrm{SeqBdd}(a)) \to \varphi_{16}$\lbl{A1}}
\BIC{$\varphi_{16}$}
\end{ptree}
\noindent$\Pi_{115}$:
\begin{ptree}
\AXC{$\forall a\,(\varphi_{25})$\lbl{DefSeqBdd}}
\AXC{$\forall a\,(\varphi_{25}) \to \varphi_{25}$\lbl{A2}}
\BIC{$\varphi_{25}$}
\AXC{$\mathrm{Conv}(a,\alpha) \to \mathrm{Conv}(a,\alpha)$\lbl{A1}}
\AXC{$\stackrel{\Pi_{114}}{\vdots}$}\noLine
\UIC{$\varphi_{16}$}
\BIC{$\mathrm{Conv}(a,\alpha) \to \mathrm{Conv}(b,\beta) \to \mathrm{SeqBdd}(a)$}
\AXC{$\varphi_{132}$\lbl{A1}}
\BIC{$\varphi_{25} \to \varphi_{32}$}
\BIC{$\varphi_{32}$}
\AXC{$\varphi_{240}$\lbl{A1}}
\BIC{$\forall M\,(\varphi_{9} \to \neg (M \in \mathbb{R} \land \forall n \in \mathbb{N}\,(\lvert a_{n} \rvert \leq M))) \to \varphi_{234}$}
\end{ptree}
\noindent$\Pi_{116}$:
\begin{ptree}
\AXC{$\stackrel{\Pi_{113}}{\vdots}$}\noLine
\UIC{$\forall M\,(\varphi_{9} \to \neg (M \in \mathbb{R} \land \forall n \in \mathbb{N}\,(\lvert a_{n} \rvert \leq M)))$}
\AXC{$\stackrel{\Pi_{115}}{\vdots}$}\noLine
\UIC{$\forall M\,(\varphi_{9} \to \neg (M \in \mathbb{R} \land \forall n \in \mathbb{N}\,(\lvert a_{n} \rvert \leq M))) \to \varphi_{234}$}
\BIC{$\varphi_{234}$}
\end{ptree}
\noindent 主木:
\begin{ptree}
\AXC{$\varphi_{9} \to \forall M\,(\varphi_{9})$\lbl{A4}}
\AXC{$\forall M\,(\varphi_{9} \to \neg (M \in \mathbb{R} \land \forall n \in \mathbb{N}\,(\lvert a_{n} \rvert \leq M))) \to \varphi_{92}$\lbl{A3}}
\AXC{$\stackrel{\Pi_{116}}{\vdots}$}\noLine
\UIC{$\varphi_{234}$}
\BIC{$(\varphi_{9} \to \forall M\,(\varphi_{9})) \to \varphi_{8}$}
\BIC{$\varphi_{8}$}
\UIC{$\forall \beta\,(\varphi_{8})$}
\UIC{$\forall b\,\forall \beta\,(\varphi_{8})$}
\UIC{$\forall \alpha\,\forall b\,\forall \beta\,(\varphi_{8})$}
\UIC{$\forall a\,\forall \alpha\,\forall b\,\forall \beta\,(\varphi_{8})$}
\end{ptree}
\begin{abbrevlist}
\item $\varphi_{1} :\equiv \forall \alpha\,(\mathrm{Conv}(a,\alpha) \to \mathrm{Seq}(a)) \to \mathrm{Conv}(a,\alpha) \to \mathrm{Seq}(a)$
\item $\varphi_{2} :\equiv \forall \alpha\,(\mathrm{Conv}(b,\alpha) \to \mathrm{Seq}(b)) \to \mathrm{Conv}(b,\beta) \to \mathrm{Seq}(b)$
\item $\varphi_{3} :\equiv \forall x\,\forall y\,\forall z\,(x \leq y \to y \leq z \to x \leq z) \to \forall y\,\forall z\,(0 \leq y \to y \leq z \to 0 \leq z)$
\item $\varphi_{4} :\equiv \forall z\,(0 \leq \lvert a_{1_{\omega}} \rvert \to \lvert a_{1_{\omega}} \rvert \leq z \to 0 \leq z)$
\item $\varphi_{5} :\equiv \forall \alpha\,(\mathrm{Conv}(a,\alpha) \to \mathrm{SeqBdd}(a)) \to \mathrm{Conv}(a,\alpha) \to \mathrm{SeqBdd}(a)$
\item $\varphi_{6} :\equiv \forall a\,\forall \alpha\,(\mathrm{Conv}(a,\alpha) \to \mathrm{Seq}(a)) \to \forall \alpha\,(\mathrm{Conv}(a,\alpha) \to \mathrm{Seq}(a))$
\item $\varphi_{7} :\equiv \forall a\,\forall \alpha\,(\mathrm{Conv}(a,\alpha) \to \mathrm{Seq}(a)) \to \forall \alpha\,(\mathrm{Conv}(b,\alpha) \to \mathrm{Seq}(b))$
\item $\varphi_{8} :\equiv \mathrm{Conv}(a,\alpha) \to \mathrm{Conv}(b,\beta) \to \mathrm{Conv}(a \otimes b,\alpha \cdot \beta)$
\item $\varphi_{9} :\equiv \mathrm{Conv}(a,\alpha) \land \mathrm{Conv}(b,\beta) \land \neg \mathrm{Conv}(a \otimes b,\alpha \cdot \beta)$
\item $\varphi_{10} :\equiv \forall \alpha\,(\mathrm{Conv}(b,\alpha) \to \alpha \in \mathbb{R}) \to \mathrm{Conv}(b,\beta) \to \beta \in \mathbb{R}$
\item $\varphi_{11} :\equiv \forall \alpha\,(\mathrm{Conv}(a,\alpha) \to \alpha \in \mathbb{R}) \to \mathrm{Conv}(a,\alpha) \to \alpha \in \mathbb{R}$
\item $\varphi_{12} :\equiv (\mathrm{Conv}(a,\alpha) \to \mathrm{Conv}(a,\alpha)) \to \mathrm{Conv}(a,\alpha) \to \mathrm{Conv}(b,\beta) \to \mathrm{Seq}(a)$
\item $\varphi_{13} :\equiv \forall a\,\forall \alpha\,(\mathrm{Conv}(a,\alpha) \to \mathrm{SeqBdd}(a)) \to \forall \alpha\,(\mathrm{Conv}(a,\alpha) \to \mathrm{SeqBdd}(a))$
\item $\varphi_{14} :\equiv \forall x\,\forall y\,\forall z\,(x \leq y \to y < z \to x < z) \to \forall y\,\forall z\,(0 \leq y \to y < z \to 0 < z)$
\item $\varphi_{15} :\equiv N \leq n \to (M + \lvert \beta \rvert + 1) \cdot (\varepsilon \cdot (M + \lvert \beta \rvert + 1)^{-1}) = \varepsilon$
\item $\varphi_{16} :\equiv (\mathrm{Conv}(a,\alpha) \to \mathrm{Conv}(a,\alpha)) \to \mathrm{Conv}(a,\alpha) \to \mathrm{Conv}(b,\beta) \to \mathrm{SeqBdd}(a)$
\item $\varphi_{17} :\equiv \forall a\,\forall \alpha\,(\mathrm{Conv}(a,\alpha) \to \alpha \in \mathbb{R}) \to \forall \alpha\,(\mathrm{Conv}(b,\alpha) \to \alpha \in \mathbb{R})$
\item $\varphi_{18} :\equiv \forall x\,\forall y\,(x = y \to y = x) \to \forall y\,(M + \lvert \beta \rvert + 1 = y \to y = M + \lvert \beta \rvert + 1)$
\item $\varphi_{19} :\equiv \forall a\,\forall \alpha\,(\mathrm{Conv}(a,\alpha) \to \alpha \in \mathbb{R}) \to \forall \alpha\,(\mathrm{Conv}(a,\alpha) \to \alpha \in \mathbb{R})$
\item $\varphi_{20} :\equiv \mathrm{Conv}(b,\beta) \to M \in \mathbb{R} \land \forall n \in \mathbb{N}\,(\lvert a_{n} \rvert \leq M) \to 0 \leq M$
\item $\varphi_{21} :\equiv (\mathrm{Conv}(a,\alpha) \to \mathrm{Conv}(a,\alpha)) \to \mathrm{Conv}(a,\alpha) \to \mathrm{Conv}(b,\beta) \to \alpha \in \mathbb{R}$
\item $\varphi_{22} :\equiv \forall y\,(0 < y \to 0 \neq y) \to 0 < M + \lvert \beta \rvert + 1 \to 0 \neq M + \lvert \beta \rvert + 1$
\item $\varphi_{23} :\equiv \forall x\,\forall y\,(x = y \to y = x) \to \forall y\,(a \otimes b(n) = y \to y = a \otimes b(n))$
\item $\varphi_{24} :\equiv \forall n \in \mathbb{N}\,(\lvert a_{n} \rvert \leq M) \to 1_{\omega} \in \mathbb{N} \to \lvert a_{1_{\omega}} \rvert \leq M$
\item $\varphi_{25} :\equiv \mathrm{SeqBdd}(a) \leftrightarrow \mathrm{Seq}(a) \land \exists M \in \mathbb{R}\,\forall n \in \mathbb{N}\,(\lvert a_{n} \rvert \leq M)$
\item $\varphi_{26} :\equiv \forall x \in \mathbb{R}\,(0 \leq \lvert x \rvert) \to a_{1_{\omega}} \in \mathbb{R} \to 0 \leq \lvert a_{1_{\omega}} \rvert$
\item $\varphi_{27} :\equiv \mathrm{Conv}(b,\beta) \to M \in \mathbb{R} \land \forall n \in \mathbb{N}\,(\lvert a_{n} \rvert \leq M) \to 0 \leq \lvert \beta \rvert$
\item $\varphi_{28} :\equiv \mathrm{Conv}(a,\alpha) \land \mathrm{Conv}(b,\beta) \land (M \in \mathbb{R} \land \forall n \in \mathbb{N}\,(\lvert a_{n} \rvert \leq M))$
\item $\varphi_{29} :\equiv \mathrm{Conv}(b,\beta) \to M \in \mathbb{R} \land \forall n \in \mathbb{N}\,(\lvert a_{n} \rvert \leq M) \to M \in \mathbb{R}$
\item $\varphi_{30} :\equiv n \in \mathbb{N} \to N \leq n \to (M + \lvert \beta \rvert) \cdot (\varepsilon \cdot (M + \lvert \beta \rvert + 1)^{-1}) < \varepsilon$
\item $\varphi_{31} :\equiv (\mathrm{Conv}(a,\alpha) \to \mathrm{Conv}(b,\beta) \to \mathrm{Conv}(b,\beta)) \to \mathrm{Conv}(a,\alpha) \to \mathrm{Conv}(b,\beta) \to \mathrm{Seq}(b)$
\item $\varphi_{32} :\equiv \mathrm{Conv}(a,\alpha) \to \mathrm{Conv}(b,\beta) \to \exists M \in \mathbb{R}\,\forall n \in \mathbb{N}\,(\lvert a_{n} \rvert \leq M)$
\item $\varphi_{33} :\equiv (M + \lvert \beta \rvert + 1)^{-1} \in \mathbb{R} \to \varepsilon \cdot (M + \lvert \beta \rvert + 1)^{-1} \in \mathbb{R}$
\item $\varphi_{34} :\equiv n \in \mathbb{N} \to N \leq n \to \varepsilon \cdot (M + \lvert \beta \rvert + 1)^{-1} \cdot (M + \lvert \beta \rvert + 1) = \varepsilon$
\item $\varphi_{35} :\equiv M \in \mathbb{R} \land \forall n \in \mathbb{N}\,(\lvert a_{n} \rvert \leq M) \to \mathrm{Conv}(a \otimes b,\alpha \cdot \beta)$
\item $\varphi_{36} :\equiv \mathrm{Conv}(b,\beta) \to M \in \mathbb{R} \land \forall n \in \mathbb{N}\,(\lvert a_{n} \rvert \leq M) \to 0 \leq M + \lvert \beta \rvert$
\item $\varphi_{37} :\equiv 0 \leq M + \lvert \beta \rvert \to M + \lvert \beta \rvert < M + \lvert \beta \rvert + 1 \to 0 < M + \lvert \beta \rvert + 1$
\item $\varphi_{38} :\equiv \mathrm{Conv}(b,\beta) \to M \in \mathbb{R} \land \forall n \in \mathbb{N}\,(\lvert a_{n} \rvert \leq M) \to \lvert \beta \rvert \in \mathbb{R}$
\item $\varphi_{39} :\equiv 0 < \varepsilon \to 0 < (M + \lvert \beta \rvert + 1)^{-1} \to 0 < \varepsilon \cdot (M + \lvert \beta \rvert + 1)^{-1}$
\item $\varphi_{40} :\equiv \mathrm{Conv}(b,\beta) \to M \in \mathbb{R} \land \forall n \in \mathbb{N}\,(\lvert a_{n} \rvert \leq M) \to \mathrm{Seq}(a \otimes b)$
\item $\varphi_{41} :\equiv (\mathrm{Conv}(a,\alpha) \to \mathrm{Conv}(b,\beta) \to \mathrm{Conv}(b,\beta)) \to \mathrm{Conv}(a,\alpha) \to \mathrm{Conv}(b,\beta) \to \beta \in \mathbb{R}$
\item $\varphi_{42} :\equiv \forall y\,\forall z\,(0 \leq y \to y < z \to 0 < z) \to \forall z\,(0 \leq M + \lvert \beta \rvert \to M + \lvert \beta \rvert < z \to 0 < z)$
\item $\varphi_{43} :\equiv \neg (N \in \mathbb{N} \land \forall n \in \mathbb{N}\,(N \leq n \to \lvert a \otimes b(n) - \alpha \cdot \beta \rvert < \varepsilon))$
\item $\varphi_{44} :\equiv M \in \mathbb{R} \land \forall n \in \mathbb{N}\,(\lvert a_{n} \rvert \leq M) \to \forall n \in \mathbb{N}\,(\lvert a_{n} \rvert \leq M)$
\item $\varphi_{45} :\equiv \mathrm{Conv}(b,\beta) \to M \in \mathbb{R} \land \forall n \in \mathbb{N}\,(\lvert a_{n} \rvert \leq M) \to \alpha \cdot \beta \in \mathbb{R}$
\item $\varphi_{46} :\equiv \mathrm{Conv}(b,\beta) \to M \in \mathbb{R} \land \forall n \in \mathbb{N}\,(\lvert a_{n} \rvert \leq M) \to \lvert a_{1_{\omega}} \rvert \leq M$
\item $\varphi_{47} :\equiv \mathrm{Conv}(b,\beta) \to M \in \mathbb{R} \land \forall n \in \mathbb{N}\,(\lvert a_{n} \rvert \leq M) \to 0 \leq \lvert a_{1_{\omega}} \rvert$
\item $\varphi_{48} :\equiv M + \lvert \beta \rvert + 1 \neq 0 \to \varepsilon \cdot (M + \lvert \beta \rvert + 1)^{-1} \cdot (M + \lvert \beta \rvert + 1) = \varepsilon$
\item $\varphi_{49} :\equiv M \in \mathbb{R} \land \forall n \in \mathbb{N}\,(\lvert a_{n} \rvert \leq M) \to M + \lvert \beta \rvert < M + \lvert \beta \rvert + 1$
\item $\varphi_{50} :\equiv M \in \mathbb{R} \land \forall n \in \mathbb{N}\,(\lvert a_{n} \rvert \leq M) \to 1_{\omega} \in \mathbb{N} \to \lvert a_{1_{\omega}} \rvert \leq M$
\item $\varphi_{51} :\equiv \mathrm{Conv}(b,\beta) \to M \in \mathbb{R} \land \forall n \in \mathbb{N}\,(\lvert a_{n} \rvert \leq M) \to a_{1_{\omega}} \in \mathbb{R}$
\item $\varphi_{52} :\equiv \mathrm{Conv}(b,\beta) \to M \in \mathbb{R} \land \forall n \in \mathbb{N}\,(\lvert a_{n} \rvert \leq M) \to 0 \neq M + \lvert \beta \rvert + 1$
\item $\varphi_{53} :\equiv \mathrm{Conv}(b,\beta) \to M \in \mathbb{R} \land \forall n \in \mathbb{N}\,(\lvert a_{n} \rvert \leq M) \to M + \lvert \beta \rvert + 1 \neq 0$
\item $\varphi_{54} :\equiv \mathrm{Conv}(b,\beta) \to M \in \mathbb{R} \land \forall n \in \mathbb{N}\,(\lvert a_{n} \rvert \leq M) \to M + \lvert \beta \rvert \in \mathbb{R}$
\item $\varphi_{55} :\equiv M \in \mathbb{R} \land \forall n \in \mathbb{N}\,(\lvert a_{n} \rvert \leq M) \to \varepsilon \in \mathbb{R} \to 0 < \varepsilon \to 0 < \varepsilon$
\item $\varphi_{56} :\equiv M + \lvert \beta \rvert + 1 \in \mathbb{R} \to 0 < M + \lvert \beta \rvert + 1 \to 0 < (M + \lvert \beta \rvert + 1)^{-1}$
\item $\varphi_{57} :\equiv \mathrm{Conv}(b,\beta) \to M \in \mathbb{R} \land \forall n \in \mathbb{N}\,(\lvert a_{n} \rvert \leq M) \to 0 < M + \lvert \beta \rvert + 1$
\item $\varphi_{58} :\equiv \forall x\,\forall y\,\forall z\,(x \in \mathbb{R} \to y \in \mathbb{R} \to z \in \mathbb{R} \to 0 < z \to x < y \to x \cdot z < y \cdot z)$
\item $\varphi_{59} :\equiv \forall n\,(\mathrm{Seq}(b) \to n \in \mathbb{N} \to b_{n} \in \mathbb{R}) \to \mathrm{Seq}(b) \to n \in \mathbb{N} \to b_{n} \in \mathbb{R}$
\item $\varphi_{60} :\equiv \forall n\,(\mathrm{Seq}(a) \to n \in \mathbb{N} \to a_{n} \in \mathbb{R}) \to \mathrm{Seq}(a) \to n \in \mathbb{N} \to a_{n} \in \mathbb{R}$
\item $\varphi_{61} :\equiv \mathrm{Seq}(b) \to \mathrm{Seq}(a \otimes b) \land \forall n \in \mathbb{N}\,(a \otimes b(n) = a_{n} \cdot b_{n})$
\item $\varphi_{62} :\equiv M + \lvert \beta \rvert + 1 \in \mathbb{R} \to M + \lvert \beta \rvert + 1 \neq 0 \to (M + \lvert \beta \rvert + 1)^{-1} \in \mathbb{R}$
\item $\varphi_{63} :\equiv \lvert a_{n} \cdot b_{n} - \alpha \cdot \beta \rvert \leq (M + \lvert \beta \rvert) \cdot (\varepsilon \cdot (M + \lvert \beta \rvert + 1)^{-1})$
\item $\varphi_{64} :\equiv \mathrm{Conv}(b,\beta) \to M \in \mathbb{R} \land \forall n \in \mathbb{N}\,(\lvert a_{n} \rvert \leq M) \to \varepsilon \in \mathbb{R} \to \varepsilon \in \mathbb{R}$
\item $\varphi_{65} :\equiv \mathrm{Conv}(b,\beta) \to M \in \mathbb{R} \land \forall n \in \mathbb{N}\,(\lvert a_{n} \rvert \leq M) \to n \in \mathbb{N} \to \lvert a_{n} \rvert \leq M$
\item $\varphi_{66} :\equiv \mathrm{Conv}(b,\beta) \to \mathrm{Seq}(a \otimes b) \land \forall n \in \mathbb{N}\,(a \otimes b(n) = a_{n} \cdot b_{n})$
\item $\varphi_{67} :\equiv \forall x \in \mathbb{R}\,(x < x + 1) \to M + \lvert \beta \rvert \in \mathbb{R} \to M + \lvert \beta \rvert < M + \lvert \beta \rvert + 1$
\item $\varphi_{68} :\equiv \forall a\,\forall n\,(\mathrm{Seq}(a) \to n \in \mathbb{N} \to a_{n} \in \mathbb{R}) \to \forall n\,(\mathrm{Seq}(a) \to n \in \mathbb{N} \to a_{n} \in \mathbb{R})$
\item $\varphi_{69} :\equiv \forall n\,(\mathrm{Seq}(a) \to n \in \mathbb{N} \to a_{n} \in \mathbb{R}) \to \mathrm{Seq}(a) \to 1_{\omega} \in \mathbb{N} \to a_{1_{\omega}} \in \mathbb{R}$
\item $\varphi_{70} :\equiv \forall y\,(M + \lvert \beta \rvert + 1 = y \to y = M + \lvert \beta \rvert + 1) \to M + \lvert \beta \rvert + 1 = 0 \to 0 = M + \lvert \beta \rvert + 1$
\item $\varphi_{71} :\equiv \forall a\,\forall n\,(\mathrm{Seq}(a) \to n \in \mathbb{N} \to a_{n} \in \mathbb{R}) \to \forall n\,(\mathrm{Seq}(b) \to n \in \mathbb{N} \to b_{n} \in \mathbb{R})$
\item $\varphi_{72} :\equiv M \in \mathbb{R} \land \forall n \in \mathbb{N}\,(\lvert a_{n} \rvert \leq M) \to \varepsilon \in \mathbb{R} \to 0 < \varepsilon \to \lvert \beta \rvert \in \mathbb{R}$
\item $\varphi_{73} :\equiv M \in \mathbb{R} \land \forall n \in \mathbb{N}\,(\lvert a_{n} \rvert \leq M) \to M \in \mathbb{R} \land \forall n \in \mathbb{N}\,(\lvert a_{n} \rvert \leq M)$
\item $\varphi_{74} :\equiv 0 < \varepsilon \to \exists N \in \mathbb{N}\,\forall n \in \mathbb{N}\,(N \leq n \to \lvert a \otimes b(n) - \alpha \cdot \beta \rvert < \varepsilon)$
\item $\varphi_{75} :\equiv \forall y\,(M + \lvert \beta \rvert + 1 \in \mathbb{R} \to y \in \mathbb{R} \to (M + \lvert \beta \rvert + 1) \cdot y = y \cdot (M + \lvert \beta \rvert + 1))$
\item $\varphi_{76} :\equiv M \in \mathbb{R} \land \forall n \in \mathbb{N}\,(\lvert a_{n} \rvert \leq M) \to \varepsilon \in \mathbb{R} \to 0 < \varepsilon \to M + \lvert \beta \rvert \in \mathbb{R}$
\item $\varphi_{77} :\equiv \forall x\,\forall y\,(x \in \mathbb{R} \to y \in \mathbb{R} \to x + y \in \mathbb{R}) \to \forall y\,(M \in \mathbb{R} \to y \in \mathbb{R} \to M + y \in \mathbb{R})$
\item $\varphi_{78} :\equiv \forall y\,(\alpha \in \mathbb{R} \to y \in \mathbb{R} \to \alpha \cdot y \in \mathbb{R}) \to \alpha \in \mathbb{R} \to \beta \in \mathbb{R} \to \alpha \cdot \beta \in \mathbb{R}$
\item $\varphi_{79} :\equiv \forall x\,\forall y\,(x \in \mathbb{R} \to y \in \mathbb{R} \to x \cdot y \in \mathbb{R}) \to \forall y\,(\alpha \in \mathbb{R} \to y \in \mathbb{R} \to \alpha \cdot y \in \mathbb{R})$
\item $\varphi_{80} :\equiv \forall x\,\forall y\,(x \in \mathbb{R} \to y \in \mathbb{R} \to x \cdot y \in \mathbb{R}) \to \forall y\,(\varepsilon \in \mathbb{R} \to y \in \mathbb{R} \to \varepsilon \cdot y \in \mathbb{R})$
\item $\varphi_{81} :\equiv M \in \mathbb{R} \land \forall n \in \mathbb{N}\,(\lvert a_{n} \rvert \leq M) \to \varepsilon \in \mathbb{R} \to 0 < \varepsilon \to M + \lvert \beta \rvert + 1 \in \mathbb{R}$
\item $\varphi_{82} :\equiv \forall y\,\forall z\,(\lvert a_{n} \cdot b_{n} - \alpha \cdot \beta \rvert \leq y \to y < z \to \lvert a_{n} \cdot b_{n} - \alpha \cdot \beta \rvert < z)$
\item $\varphi_{83} :\equiv (M + \lvert \beta \rvert) \cdot (\varepsilon \cdot (M + \lvert \beta \rvert + 1)^{-1}) < (M + \lvert \beta \rvert + 1) \cdot (\varepsilon \cdot (M + \lvert \beta \rvert + 1)^{-1})$
\item $\varphi_{84} :\equiv M \in \mathbb{R} \land \forall n \in \mathbb{N}\,(\lvert a_{n} \rvert \leq M) \to \varepsilon \in \mathbb{R} \to 0 < \varepsilon \to 0 < (M + \lvert \beta \rvert + 1)^{-1}$
\item $\varphi_{85} :\equiv (M + \lvert \beta \rvert) \cdot (\varepsilon \cdot (M + \lvert \beta \rvert + 1)^{-1}) < \varepsilon \to \lvert a_{n} \cdot b_{n} - \alpha \cdot \beta \rvert < \varepsilon$
\item $\varphi_{86} :\equiv \forall w\,(a_{n} \cdot b_{n} = w \to \lvert a_{n} \cdot b_{n} - \alpha \cdot \beta \rvert < \varepsilon \to \lvert w - \alpha \cdot \beta \rvert < \varepsilon)$
\item $\varphi_{87} :\equiv \forall y\,(M \in \mathbb{R} \to y \in \mathbb{R} \to M + y \in \mathbb{R}) \to M \in \mathbb{R} \to \lvert \beta \rvert \in \mathbb{R} \to M + \lvert \beta \rvert \in \mathbb{R}$
\item $\varphi_{88} :\equiv (M + \lvert \beta \rvert + 1) \cdot (\varepsilon \cdot (M + \lvert \beta \rvert + 1)^{-1}) = \varepsilon \cdot (M + \lvert \beta \rvert + 1)^{-1} \cdot (M + \lvert \beta \rvert + 1)$
\item $\varphi_{89} :\equiv (M + \lvert \beta \rvert) \cdot (\varepsilon \cdot (M + \lvert \beta \rvert + 1)^{-1}) < z \to (M + \lvert \beta \rvert) \cdot (\varepsilon \cdot (M + \lvert \beta \rvert + 1)^{-1}) < w$
\item $\varphi_{90} :\equiv M \in \mathbb{R} \land \forall n \in \mathbb{N}\,(\lvert a_{n} \rvert \leq M) \to \varepsilon \in \mathbb{R} \to 0 < \varepsilon \to (M + \lvert \beta \rvert + 1)^{-1} \in \mathbb{R}$
\item $\varphi_{91} :\equiv \forall x\,\forall y\,(x \in \mathbb{R} \to y \in \mathbb{R} \to 0 \leq x \to 0 \leq y \to 0 \leq x + y) \to \forall y\,(M \in \mathbb{R} \to y \in \mathbb{R} \to 0 \leq M \to 0 \leq y \to 0 \leq M + y)$
\item $\varphi_{92} :\equiv \forall M\,(\varphi_{9}) \to \forall M\,(\neg (M \in \mathbb{R} \land \forall n \in \mathbb{N}\,(\lvert a_{n} \rvert \leq M)))$
\item $\varphi_{93} :\equiv M \in \mathbb{R} \land \forall n \in \mathbb{N}\,(\lvert a_{n} \rvert \leq M) \to \varepsilon \in \mathbb{R} \to 0 < \varepsilon \to 0 < \varepsilon \cdot (M + \lvert \beta \rvert + 1)^{-1}$
\item $\varphi_{94} :\equiv \lvert a_{n} - \alpha \rvert < \varepsilon \cdot (M + \lvert \beta \rvert + 1)^{-1} \land \lvert b_{n} - \beta \rvert < \varepsilon \cdot (M + \lvert \beta \rvert + 1)^{-1}$
\item $\varphi_{95} :\equiv \forall n \in \mathbb{N}\,(a \otimes b(n) = a_{n} \cdot b_{n}) \to n \in \mathbb{N} \to a \otimes b(n) = a_{n} \cdot b_{n}$
\item $\varphi_{96} :\equiv M \in \mathbb{R} \land \forall n \in \mathbb{N}\,(\lvert a_{n} \rvert \leq M) \to \varepsilon \in \mathbb{R} \to 0 < \varepsilon \to \varepsilon \cdot (M + \lvert \beta \rvert + 1)^{-1} \in \mathbb{R}$
\item $\varphi_{97} :\equiv \varepsilon \cdot (M + \lvert \beta \rvert + 1)^{-1} \cdot (M + \lvert \beta \rvert + 1) = \varepsilon \to (M + \lvert \beta \rvert + 1) \cdot (\varepsilon \cdot (M + \lvert \beta \rvert + 1)^{-1}) = \varepsilon$
\item $\varphi_{98} :\equiv \forall y\,\forall z\,(M + \lvert \beta \rvert \in \mathbb{R} \to y \in \mathbb{R} \to z \in \mathbb{R} \to 0 < z \to M + \lvert \beta \rvert < y \to (M + \lvert \beta \rvert) \cdot z < y \cdot z)$
\item $\varphi_{99} :\equiv \forall y\,(M \in \mathbb{R} \to y \in \mathbb{R} \to 0 \leq M \to 0 \leq y \to 0 \leq M + y) \to M \in \mathbb{R} \to \lvert \beta \rvert \in \mathbb{R} \to 0 \leq M \to 0 \leq \lvert \beta \rvert \to 0 \leq M + \lvert \beta \rvert$
\item $\varphi_{100} :\equiv \forall x\,\forall y\,(x \in \mathbb{R} \to y \in \mathbb{R} \to x + y \in \mathbb{R}) \to \forall y\,(M + \lvert \beta \rvert \in \mathbb{R} \to y \in \mathbb{R} \to M + \lvert \beta \rvert + y \in \mathbb{R})$
\item $\varphi_{101} :\equiv \forall y\,\forall z\,((M + \lvert \beta \rvert + 1) \cdot (\varepsilon \cdot (M + \lvert \beta \rvert + 1)^{-1}) = y \to y = z \to (M + \lvert \beta \rvert + 1) \cdot (\varepsilon \cdot (M + \lvert \beta \rvert + 1)^{-1}) = z)$
\item $\varphi_{102} :\equiv \forall x\,\forall y\,(x \in \mathbb{R} \to y \in \mathbb{R} \to y \neq 0 \to x \cdot y^{-1} \cdot y = x) \to \forall y\,(\varepsilon \in \mathbb{R} \to y \in \mathbb{R} \to y \neq 0 \to \varepsilon \cdot y^{-1} \cdot y = \varepsilon)$
\item $\varphi_{103} :\equiv (\mathrm{Conv}(a,\alpha) \to \mathrm{Conv}(b,\beta) \to \beta \in \mathbb{R}) \to \mathrm{Conv}(a,\alpha) \to \varphi_{27}$
\item $\varphi_{104} :\equiv (\mathrm{Conv}(a,\alpha) \to \mathrm{Conv}(b,\beta) \to \beta \in \mathbb{R}) \to \mathrm{Conv}(a,\alpha) \to \varphi_{38}$
\item $\varphi_{105} :\equiv \forall x\,\forall y\,(x \in \mathbb{R} \to y \in \mathbb{R} \to 0 < x \to 0 < y \to 0 < x \cdot y) \to \forall y\,(\varepsilon \in \mathbb{R} \to y \in \mathbb{R} \to 0 < \varepsilon \to 0 < y \to 0 < \varepsilon \cdot y)$
\item $\varphi_{106} :\equiv (\mathrm{Conv}(a,\alpha) \to \mathrm{Conv}(b,\beta) \to \beta \in \mathbb{R}) \to \mathrm{Conv}(a,\alpha) \to \varphi_{45}$
\item $\varphi_{107} :\equiv \forall y\,(a \otimes b(n) = y \to y = a \otimes b(n)) \to a \otimes b(n) = a_{n} \cdot b_{n} \to a_{n} \cdot b_{n} = a \otimes b(n)$
\item $\varphi_{108} :\equiv a_{n} \cdot b_{n} = a \otimes b(n) \to \lvert a_{n} \cdot b_{n} - \alpha \cdot \beta \rvert < \varepsilon \to \lvert a \otimes b(n) - \alpha \cdot \beta \rvert < \varepsilon$
\item $\varphi_{109} :\equiv \forall y\,(M + \lvert \beta \rvert \in \mathbb{R} \to y \in \mathbb{R} \to M + \lvert \beta \rvert + y \in \mathbb{R}) \to M + \lvert \beta \rvert \in \mathbb{R} \to 1 \in \mathbb{R} \to M + \lvert \beta \rvert + 1 \in \mathbb{R}$
\item $\varphi_{110} :\equiv (\mathrm{Conv}(a,\alpha) \to \mathrm{Conv}(b,\beta) \to \mathrm{Seq}(b)) \to \mathrm{Conv}(a,\alpha) \to \varphi_{66}$
\item $\varphi_{111} :\equiv (M + \lvert \beta \rvert + 1 = 0 \to 0 = M + \lvert \beta \rvert + 1) \to \mathrm{Conv}(a,\alpha) \to \varphi_{53}$
\item $\varphi_{112} :\equiv \forall a\,\forall \alpha\,(\mathrm{Conv}(a,\alpha) \leftrightarrow \mathrm{Seq}(a) \land \alpha \in \mathbb{R} \land \forall \varepsilon \in \mathbb{R}\,(0 < \varepsilon \to \exists N \in \mathbb{N}\,\forall n \in \mathbb{N}\,(N \leq n \to \lvert a_{n} - \alpha \rvert < \varepsilon)))$
\item $\varphi_{113} :\equiv (\mathrm{Conv}(a,\alpha) \to \mathrm{Conv}(b,\beta) \to \mathrm{Seq}(a)) \to 1_{\omega} \in \mathbb{N} \to \mathrm{Conv}(a,\alpha) \to \varphi_{51}$
\item $\varphi_{114} :\equiv \forall \varepsilon\,(\mathrm{Conv}(a,\alpha) \to \mathrm{Conv}(b,\beta) \to \varepsilon \in \mathbb{R} \to 0 < \varepsilon \to \exists N \in \mathbb{N}\,\forall n \in \mathbb{N}\,(N \leq n \to \lvert a_{n} - \alpha \rvert < \varepsilon \land \lvert b_{n} - \beta \rvert < \varepsilon))$
\item $\varphi_{115} :\equiv (\mathrm{Conv}(a,\alpha) \to \mathrm{Conv}(b,\beta) \to \beta \in \mathbb{R}) \to \mathrm{Conv}(a,\alpha) \to \mathrm{Conv}(b,\beta) \to \varphi_{72}$
\item $\varphi_{116} :\equiv \forall z\,(M + \lvert \beta \rvert \in \mathbb{R} \to M + \lvert \beta \rvert + 1 \in \mathbb{R} \to z \in \mathbb{R} \to 0 < z \to M + \lvert \beta \rvert < M + \lvert \beta \rvert + 1 \to (M + \lvert \beta \rvert) \cdot z < (M + \lvert \beta \rvert + 1) \cdot z)$
\item $\varphi_{117} :\equiv (\forall n \in \mathbb{N}\,(\lvert a_{n} \rvert \leq M) \to n \in \mathbb{N} \to \lvert a_{n} \rvert \leq M) \to \mathrm{Conv}(a,\alpha) \to \varphi_{65}$
\item $\varphi_{118} :\equiv \forall \alpha\,(\mathrm{Conv}(a \otimes b,\alpha) \leftrightarrow \mathrm{Seq}(a \otimes b) \land \alpha \in \mathbb{R} \land \forall \varepsilon \in \mathbb{R}\,(0 < \varepsilon \to \exists N \in \mathbb{N}\,\forall n \in \mathbb{N}\,(N \leq n \to \lvert a \otimes b(n) - \alpha \rvert < \varepsilon)))$
\item $\varphi_{119} :\equiv \mathrm{Conv}(b,\beta) \to M \in \mathbb{R} \land \forall n \in \mathbb{N}\,(\lvert a_{n} \rvert \leq M) \to \varepsilon \in \mathbb{R} \to \varphi_{74}$
\item $\varphi_{120} :\equiv \forall z\,\forall w\,(z = w \to \lvert z - \alpha \cdot \beta \rvert < \varepsilon \to \lvert w - \alpha \cdot \beta \rvert < \varepsilon) \to \varphi_{86}$
\item $\varphi_{121} :\equiv \mathrm{Conv}(b,\beta) \to M \in \mathbb{R} \land \forall n \in \mathbb{N}\,(\lvert a_{n} \rvert \leq M) \to \forall \varepsilon \in \mathbb{R}\,(\varphi_{74})$
\item $\varphi_{122} :\equiv 0 < \varepsilon \cdot (M + \lvert \beta \rvert + 1)^{-1} \to M + \lvert \beta \rvert < M + \lvert \beta \rvert + 1 \to \varphi_{83}$
\item $\varphi_{123} :\equiv 0 < \varepsilon \to N \in \mathbb{N} \land \forall n \in \mathbb{N}\,(N \leq n \to \varphi_{94}) \to n \in \mathbb{N} \to N \leq n \to N \leq n$
\item $\varphi_{124} :\equiv 0 < \varepsilon \to N \in \mathbb{N} \land \forall n \in \mathbb{N}\,(N \leq n \to \varphi_{94}) \to n \in \mathbb{N} \to n \in \mathbb{N}$
\item $\varphi_{125} :\equiv N \in \mathbb{N} \land \forall n \in \mathbb{N}\,(N \leq n \to \varphi_{94}) \to n \in \mathbb{N} \to N \leq n \to \lvert a_{n} \rvert \leq M$
\item $\varphi_{126} :\equiv \mathrm{Conv}(a \otimes b,\alpha \cdot \beta) \leftrightarrow \mathrm{Seq}(a \otimes b) \land \alpha \cdot \beta \in \mathbb{R} \land \forall \varepsilon \in \mathbb{R}\,(\varphi_{74})$
\item $\varphi_{127} :\equiv \varphi_{28} \land \varepsilon \in \mathbb{R} \land 0 < \varepsilon \land \neg \exists N \in \mathbb{N}\,\forall n \in \mathbb{N}\,(N \leq n \to \lvert a \otimes b(n) - \alpha \cdot \beta \rvert < \varepsilon)$
\item $\varphi_{128} :\equiv N \in \mathbb{N} \land \forall n \in \mathbb{N}\,(N \leq n \to \varphi_{94}) \to n \in \mathbb{N} \to N \leq n \to \lvert \beta \rvert \in \mathbb{R}$
\item $\varphi_{129} :\equiv N \in \mathbb{N} \land \forall n \in \mathbb{N}\,(N \leq n \to \varphi_{94}) \to n \in \mathbb{N} \to N \leq n \to b_{n} \in \mathbb{R}$
\item $\varphi_{130} :\equiv N \in \mathbb{N} \land \forall n \in \mathbb{N}\,(N \leq n \to \varphi_{94}) \to n \in \mathbb{N} \to N \leq n \to a_{n} \in \mathbb{R}$
\item $\varphi_{131} :\equiv \varepsilon \in \mathbb{R} \to 0 < \varepsilon \to N \in \mathbb{N} \land \forall n \in \mathbb{N}\,(N \leq n \to \varphi_{94}) \to N \in \mathbb{N}$
\item $\varphi_{132} :\equiv (\mathrm{Conv}(a,\alpha) \to \mathrm{Conv}(b,\beta) \to \mathrm{SeqBdd}(a)) \to \varphi_{25} \to \varphi_{32}$
\item $\varphi_{133} :\equiv 0 < \varepsilon \cdot (M + \lvert \beta \rvert + 1)^{-1} \to \exists N \in \mathbb{N}\,\forall n \in \mathbb{N}\,(N \leq n \to \varphi_{94})$
\item $\varphi_{134} :\equiv N \in \mathbb{N} \land \forall n \in \mathbb{N}\,(N \leq n \to \varphi_{94}) \to n \in \mathbb{N} \to N \leq n \to M + \lvert \beta \rvert \in \mathbb{R}$
\item $\varphi_{135} :\equiv \forall y\,(\varepsilon \in \mathbb{R} \to y \in \mathbb{R} \to y \neq 0 \to \varepsilon \cdot y^{-1} \cdot y = \varepsilon) \to \varepsilon \in \mathbb{R} \to M + \lvert \beta \rvert + 1 \in \mathbb{R} \to \varphi_{48}$
\item $\varphi_{136} :\equiv \forall z\,(\varphi_{63} \to (M + \lvert \beta \rvert) \cdot (\varepsilon \cdot (M + \lvert \beta \rvert + 1)^{-1}) < z \to \lvert a_{n} \cdot b_{n} - \alpha \cdot \beta \rvert < z)$
\item $\varphi_{137} :\equiv N \in \mathbb{N} \land \forall n \in \mathbb{N}\,(N \leq n \to \varphi_{94}) \to n \in \mathbb{N} \to N \leq n \to M + \lvert \beta \rvert + 1 \in \mathbb{R}$
\item $\varphi_{138} :\equiv \forall y\,(\varepsilon \in \mathbb{R} \to y \in \mathbb{R} \to 0 < \varepsilon \to 0 < y \to 0 < \varepsilon \cdot y) \to \varepsilon \in \mathbb{R} \to (M + \lvert \beta \rvert + 1)^{-1} \in \mathbb{R} \to \varphi_{39}$
\item $\varphi_{139} :\equiv N \in \mathbb{N} \land \forall n \in \mathbb{N}\,(N \leq n \to \varphi_{94}) \to n \in \mathbb{N} \to N \leq n \to (M + \lvert \beta \rvert + 1)^{-1} \in \mathbb{R}$
\item $\varphi_{140} :\equiv \lvert a_{n} - \alpha \rvert < \varepsilon \cdot (M + \lvert \beta \rvert + 1)^{-1} \to \lvert b_{n} - \beta \rvert < \varepsilon \cdot (M + \lvert \beta \rvert + 1)^{-1} \to \varphi_{63}$
\item $\varphi_{141} :\equiv N \in \mathbb{N} \land \forall n \in \mathbb{N}\,(N \leq n \to \varphi_{94}) \to n \in \mathbb{N} \to N \leq n \to M + \lvert \beta \rvert < M + \lvert \beta \rvert + 1$
\item $\varphi_{142} :\equiv \forall x\,\forall y\,\forall \alpha\,\forall \beta\,\forall M\,\forall h\,(x \in \mathbb{R} \to y \in \mathbb{R} \to \alpha \in \mathbb{R} \to \beta \in \mathbb{R} \to M \in \mathbb{R} \to h \in \mathbb{R} \to 0 \leq M \to \lvert x \rvert \leq M \to \lvert x - \alpha \rvert < h \to \lvert y - \beta \rvert < h \to \lvert x \cdot y - \alpha \cdot \beta \rvert \leq (M + \lvert \beta \rvert) \cdot h)$
\item $\varphi_{143} :\equiv N \in \mathbb{N} \land \forall n \in \mathbb{N}\,(N \leq n \to \varphi_{94}) \to n \in \mathbb{N} \to N \leq n \to \varepsilon \cdot (M + \lvert \beta \rvert + 1)^{-1} \in \mathbb{R}$
\item $\varphi_{144} :\equiv (\mathrm{Conv}(a,\alpha) \to \mathrm{Conv}(b,\beta) \to \varphi_{73}) \to \mathrm{Conv}(a,\alpha) \to \varphi_{29}$
\item $\varphi_{145} :\equiv N \in \mathbb{N} \land \forall n \in \mathbb{N}\,(N \leq n \to \varphi_{94}) \to n \in \mathbb{N} \to N \leq n \to a \otimes b(n) = a_{n} \cdot b_{n}$
\item $\varphi_{146} :\equiv N \in \mathbb{N} \land \forall n \in \mathbb{N}\,(N \leq n \to \varphi_{94}) \to n \in \mathbb{N} \to N \leq n \to a_{n} \cdot b_{n} = a \otimes b(n)$
\item $\varphi_{147} :\equiv (\mathrm{Conv}(a,\alpha) \to \varphi_{54}) \to \mathrm{Conv}(a,\alpha) \to \mathrm{Conv}(b,\beta) \to \varphi_{49}$
\item $\varphi_{148} :\equiv (\mathrm{Conv}(a,\alpha) \to \mathrm{Conv}(b,\beta) \to \varphi_{35}) \to \varphi_{9} \to \neg (M \in \mathbb{R} \land \forall n \in \mathbb{N}\,(\lvert a_{n} \rvert \leq M))$
\item $\varphi_{149} :\equiv (\mathrm{Conv}(a,\alpha) \to \mathrm{Conv}(b,\beta) \to \varphi_{49}) \to \mathrm{Conv}(a,\alpha) \to \varphi_{57}$
\item $\varphi_{150} :\equiv N \in \mathbb{N} \land \forall n \in \mathbb{N}\,(N \leq n \to \varphi_{94}) \to n \in \mathbb{N} \to N \leq n \to \lvert a \otimes b(n) - \alpha \cdot \beta \rvert < \varepsilon$
\item $\varphi_{151} :\equiv N \in \mathbb{N} \land \forall n \in \mathbb{N}\,(N \leq n \to \varphi_{94}) \to n \in \mathbb{N} \to N \leq n \to \lvert a_{n} \cdot b_{n} - \alpha \cdot \beta \rvert < \varepsilon$
\item $\varphi_{152} :\equiv N \in \mathbb{N} \land \forall n \in \mathbb{N}\,(N \leq n \to \varphi_{94}) \to \forall n \in \mathbb{N}\,(N \leq n \to \lvert a \otimes b(n) - \alpha \cdot \beta \rvert < \varepsilon)$
\item $\varphi_{153} :\equiv \forall w\,((M + \lvert \beta \rvert + 1) \cdot (\varepsilon \cdot (M + \lvert \beta \rvert + 1)^{-1}) = w \to \varphi_{83} \to (M + \lvert \beta \rvert) \cdot (\varepsilon \cdot (M + \lvert \beta \rvert + 1)^{-1}) < w)$
\item $\varphi_{154} :\equiv (M + \lvert \beta \rvert + 1) \cdot (\varepsilon \cdot (M + \lvert \beta \rvert + 1)^{-1}) = \varepsilon \to \varphi_{83} \to (M + \lvert \beta \rvert) \cdot (\varepsilon \cdot (M + \lvert \beta \rvert + 1)^{-1}) < \varepsilon$
\item $\varphi_{155} :\equiv N \in \mathbb{N} \land \forall n \in \mathbb{N}\,(N \leq n \to \varphi_{94}) \to n \in \mathbb{N} \to N \leq n \to n \in \mathbb{N} \to a \otimes b(n) = a_{n} \cdot b_{n}$
\item $\varphi_{156} :\equiv (\mathrm{Conv}(a,\alpha) \to \mathrm{Conv}(b,\beta) \to \varphi_{50}) \to 1_{\omega} \in \mathbb{N} \to \mathrm{Conv}(a,\alpha) \to \varphi_{46}$
\item $\varphi_{157} :\equiv N \in \mathbb{N} \land \forall n \in \mathbb{N}\,(N \leq n \to \varphi_{94}) \to n \in \mathbb{N} \to N \leq n \to \lvert b_{n} - \beta \rvert < \varepsilon \cdot (M + \lvert \beta \rvert + 1)^{-1}$
\item $\varphi_{158} :\equiv N \in \mathbb{N} \land \forall n \in \mathbb{N}\,(N \leq n \to \varphi_{94}) \to n \in \mathbb{N} \to N \leq n \to \lvert a_{n} - \alpha \rvert < \varepsilon \cdot (M + \lvert \beta \rvert + 1)^{-1}$
\item $\varphi_{159} :\equiv \forall z\,(\varphi_{88} \to \varepsilon \cdot (M + \lvert \beta \rvert + 1)^{-1} \cdot (M + \lvert \beta \rvert + 1) = z \to (M + \lvert \beta \rvert + 1) \cdot (\varepsilon \cdot (M + \lvert \beta \rvert + 1)^{-1}) = z)$
\item $\varphi_{160} :\equiv N \in \mathbb{N} \land \forall n \in \mathbb{N}\,(N \leq n \to \varphi_{94}) \to n \in \mathbb{N} \to N \leq n \to \forall n \in \mathbb{N}\,(a \otimes b(n) = a_{n} \cdot b_{n})$
\item $\varphi_{161} :\equiv M \in \mathbb{R} \land \forall n \in \mathbb{N}\,(\lvert a_{n} \rvert \leq M) \to \varepsilon \in \mathbb{R} \to 0 < \varepsilon \to \exists N \in \mathbb{N}\,\forall n \in \mathbb{N}\,(N \leq n \to \varphi_{94})$
\item $\varphi_{162} :\equiv (\alpha \in \mathbb{R} \to \beta \in \mathbb{R} \to \alpha \cdot \beta \in \mathbb{R}) \to (\mathrm{Conv}(a,\alpha) \to \mathrm{Conv}(b,\beta) \to \alpha \in \mathbb{R}) \to \varphi_{106}$
\item $\varphi_{163} :\equiv N \in \mathbb{N} \land \forall n \in \mathbb{N}\,(N \leq n \to \varphi_{94}) \to N \in \mathbb{N} \land \forall n \in \mathbb{N}\,(N \leq n \to \lvert a \otimes b(n) - \alpha \cdot \beta \rvert < \varepsilon)$
\item $\varphi_{164} :\equiv (\mathrm{Conv}(a,\alpha) \to \mathrm{Conv}(b,\beta) \to \varphi_{73}) \to \mathrm{Conv}(a,\alpha) \to \mathrm{Conv}(b,\beta) \to \varphi_{44}$
\item $\varphi_{165} :\equiv N \in \mathbb{N} \land \forall n \in \mathbb{N}\,(N \leq n \to \varphi_{94}) \to \exists N \in \mathbb{N}\,\forall n \in \mathbb{N}\,(N \leq n \to \lvert a \otimes b(n) - \alpha \cdot \beta \rvert < \varepsilon)$
\item $\varphi_{166} :\equiv (\mathrm{Conv}(a,\alpha) \to \varphi_{57}) \to \mathrm{Conv}(a,\alpha) \to \mathrm{Conv}(b,\beta) \to \varphi_{84}$
\item $\varphi_{167} :\equiv \forall y\,\forall \alpha\,\forall \beta\,\forall M\,\forall h\,(a_{n} \in \mathbb{R} \to y \in \mathbb{R} \to \alpha \in \mathbb{R} \to \beta \in \mathbb{R} \to M \in \mathbb{R} \to h \in \mathbb{R} \to 0 \leq M \to \lvert a_{n} \rvert \leq M \to \lvert a_{n} - \alpha \rvert < h \to \lvert y - \beta \rvert < h \to \lvert a_{n} \cdot y - \alpha \cdot \beta \rvert \leq (M + \lvert \beta \rvert) \cdot h)$
\item $\varphi_{168} :\equiv (\mathrm{Conv}(a,\alpha) \to \varphi_{53}) \to \mathrm{Conv}(a,\alpha) \to \mathrm{Conv}(b,\beta) \to \varphi_{90}$
\item $\varphi_{169} :\equiv (a_{1_{\omega}} \in \mathbb{R} \to 0 \leq \lvert a_{1_{\omega}} \rvert) \to (\mathrm{Conv}(a,\alpha) \to \varphi_{51}) \to \mathrm{Conv}(a,\alpha) \to \varphi_{47}$
\item $\varphi_{170} :\equiv \forall h\,(a_{n} \in \mathbb{R} \to b_{n} \in \mathbb{R} \to \alpha \in \mathbb{R} \to \beta \in \mathbb{R} \to M \in \mathbb{R} \to h \in \mathbb{R} \to 0 \leq M \to \lvert a_{n} \rvert \leq M \to \lvert a_{n} - \alpha \rvert < h \to \lvert b_{n} - \beta \rvert < h \to \lvert a_{n} \cdot b_{n} - \alpha \cdot \beta \rvert \leq (M + \lvert \beta \rvert) \cdot h)$
\item $\varphi_{171} :\equiv (\mathrm{Conv}(a,\alpha) \to \mathrm{Conv}(b,\beta) \to \varphi_{72}) \to \mathrm{Conv}(a,\alpha) \to \mathrm{Conv}(b,\beta) \to \varphi_{76}$
\item $\varphi_{172} :\equiv (0 < M + \lvert \beta \rvert + 1 \to 0 \neq M + \lvert \beta \rvert + 1) \to (\mathrm{Conv}(a,\alpha) \to \varphi_{57}) \to \mathrm{Conv}(a,\alpha) \to \varphi_{52}$
\item $\varphi_{173} :\equiv \forall n\,(\varphi_{28} \land \varepsilon \in \mathbb{R} \land 0 < \varepsilon \land (N \in \mathbb{N} \land \forall n \in \mathbb{N}\,(N \leq n \to \varphi_{94})))$
\item $\varphi_{174} :\equiv \varepsilon \in \mathbb{R} \to 0 < \varepsilon \to N \in \mathbb{N} \land \forall n \in \mathbb{N}\,(N \leq n \to \varphi_{94}) \to n \in \mathbb{N} \to \varphi_{15}$
\item $\varphi_{175} :\equiv \mathrm{Conv}(a,\alpha) \to \mathrm{Conv}(b,\beta) \to \varepsilon \cdot (M + \lvert \beta \rvert + 1)^{-1} \in \mathbb{R} \to \varphi_{133}$
\item $\varphi_{176} :\equiv (\mathrm{Conv}(a,\alpha) \to \varphi_{20}) \to (\mathrm{Conv}(a,\alpha) \to \varphi_{27}) \to \mathrm{Conv}(a,\alpha) \to \varphi_{36}$
\item $\varphi_{177} :\equiv \mathrm{Conv}(b,\beta) \to M \in \mathbb{R} \land \forall n \in \mathbb{N}\,(\lvert a_{n} \rvert \leq M) \to \varepsilon \in \mathbb{R} \to \varphi_{123}$
\item $\varphi_{178} :\equiv \mathrm{Conv}(b,\beta) \to M \in \mathbb{R} \land \forall n \in \mathbb{N}\,(\lvert a_{n} \rvert \leq M) \to \varepsilon \in \mathbb{R} \to \varphi_{124}$
\item $\varphi_{179} :\equiv \mathrm{Conv}(a,\alpha) \to \mathrm{Conv}(b,\beta) \to M \in \mathbb{R} \land \forall n \in \mathbb{N}\,(\lvert a_{n} \rvert \leq M) \to \varphi_{131}$
\item $\varphi_{180} :\equiv 0 < \varepsilon \to N \in \mathbb{N} \land \forall n \in \mathbb{N}\,(N \leq n \to \varphi_{94}) \to n \in \mathbb{N} \to N \leq n \to \varphi_{63}$
\item $\varphi_{181} :\equiv (\mathrm{Conv}(a,\alpha) \to \mathrm{Conv}(b,\beta) \to \varphi_{76}) \to 1 \in \mathbb{R} \to \mathrm{Conv}(a,\alpha) \to \mathrm{Conv}(b,\beta) \to \varphi_{81}$
\item $\varphi_{182} :\equiv (\mathrm{Conv}(a,\alpha) \to \mathrm{Conv}(b,\beta) \to \varphi_{84}) \to \mathrm{Conv}(a,\alpha) \to \mathrm{Conv}(b,\beta) \to \varphi_{93}$
\item $\varphi_{183} :\equiv (\mathrm{Conv}(a,\alpha) \to \varphi_{47}) \to (\mathrm{Conv}(a,\alpha) \to \varphi_{46}) \to \mathrm{Conv}(a,\alpha) \to \varphi_{20}$
\item $\varphi_{184} :\equiv \mathrm{Conv}(b,\beta) \to M \in \mathbb{R} \land \forall n \in \mathbb{N}\,(\lvert a_{n} \rvert \leq M) \to \varepsilon \in \mathbb{R} \to 0 < \varepsilon \to \varphi_{125}$
\item $\varphi_{185} :\equiv \mathrm{Conv}(b,\beta) \to M \in \mathbb{R} \land \forall n \in \mathbb{N}\,(\lvert a_{n} \rvert \leq M) \to \varepsilon \in \mathbb{R} \to 0 < \varepsilon \to \varphi_{128}$
\item $\varphi_{186} :\equiv (\mathrm{Seq}(a) \to \varphi_{61}) \to (\mathrm{Conv}(a,\alpha) \to \mathrm{Conv}(b,\beta) \to \mathrm{Seq}(a)) \to \varphi_{110}$
\item $\varphi_{187} :\equiv \mathrm{Conv}(b,\beta) \to M \in \mathbb{R} \land \forall n \in \mathbb{N}\,(\lvert a_{n} \rvert \leq M) \to \varepsilon \in \mathbb{R} \to 0 < \varepsilon \to \varphi_{129}$
\item $\varphi_{188} :\equiv \mathrm{Conv}(b,\beta) \to M \in \mathbb{R} \land \forall n \in \mathbb{N}\,(\lvert a_{n} \rvert \leq M) \to \varepsilon \in \mathbb{R} \to 0 < \varepsilon \to \varphi_{130}$
\item $\varphi_{189} :\equiv \varphi_{75} \to M + \lvert \beta \rvert + 1 \in \mathbb{R} \to \varepsilon \cdot (M + \lvert \beta \rvert + 1)^{-1} \in \mathbb{R} \to \varphi_{88}$
\item $\varphi_{190} :\equiv M + \lvert \beta \rvert \in \mathbb{R} \to M + \lvert \beta \rvert + 1 \in \mathbb{R} \to \varepsilon \cdot (M + \lvert \beta \rvert + 1)^{-1} \in \mathbb{R} \to \varphi_{122}$
\item $\varphi_{191} :\equiv \mathrm{Conv}(b,\beta) \to M \in \mathbb{R} \land \forall n \in \mathbb{N}\,(\lvert a_{n} \rvert \leq M) \to \varepsilon \in \mathbb{R} \to 0 < \varepsilon \to \varphi_{134}$
\item $\varphi_{192} :\equiv (\mathrm{Conv}(a,\alpha) \to \mathrm{Conv}(b,\beta) \to \varphi_{90}) \to \mathrm{Conv}(a,\alpha) \to \mathrm{Conv}(b,\beta) \to \varphi_{96}$
\item $\varphi_{193} :\equiv (\mathrm{Conv}(a,\alpha) \to \varphi_{29}) \to (\mathrm{Conv}(a,\alpha) \to \varphi_{38}) \to \mathrm{Conv}(a,\alpha) \to \varphi_{54}$
\item $\varphi_{194} :\equiv 0 < \varepsilon \to N \in \mathbb{N} \land \forall n \in \mathbb{N}\,(N \leq n \to \varphi_{94}) \to n \in \mathbb{N} \to N \leq n \to \varphi_{83}$
\item $\varphi_{195} :\equiv M \in \mathbb{R} \to \varepsilon \cdot (M + \lvert \beta \rvert + 1)^{-1} \in \mathbb{R} \to 0 \leq M \to \lvert a_{n} \rvert \leq M \to \varphi_{140}$
\item $\varphi_{196} :\equiv \mathrm{Conv}(b,\beta) \to M \in \mathbb{R} \land \forall n \in \mathbb{N}\,(\lvert a_{n} \rvert \leq M) \to \varepsilon \in \mathbb{R} \to 0 < \varepsilon \to \varphi_{137}$
\item $\varphi_{197} :\equiv 0 < \varepsilon \to N \in \mathbb{N} \land \forall n \in \mathbb{N}\,(N \leq n \to \varphi_{94}) \to n \in \mathbb{N} \to N \leq n \to \varphi_{88}$
\item $\varphi_{198} :\equiv \mathrm{Conv}(b,\beta) \to M \in \mathbb{R} \land \forall n \in \mathbb{N}\,(\lvert a_{n} \rvert \leq M) \to \varepsilon \in \mathbb{R} \to 0 < \varepsilon \to \varphi_{139}$
\item $\varphi_{199} :\equiv (\mathrm{Conv}(a,\alpha) \to \mathrm{Conv}(b,\beta) \to \varphi_{44}) \to \varphi_{24} \to \mathrm{Conv}(a,\alpha) \to \mathrm{Conv}(b,\beta) \to \varphi_{50}$
\item $\varphi_{200} :\equiv \mathrm{Conv}(b,\beta) \to M \in \mathbb{R} \land \forall n \in \mathbb{N}\,(\lvert a_{n} \rvert \leq M) \to \varepsilon \in \mathbb{R} \to 0 < \varepsilon \to \varphi_{141}$
\item $\varphi_{201} :\equiv 0 < \varepsilon \to N \in \mathbb{N} \land \forall n \in \mathbb{N}\,(N \leq n \to \varphi_{94}) \to n \in \mathbb{N} \to N \leq n \to \varphi_{94}$
\item $\varphi_{202} :\equiv (\mathrm{Conv}(a,\alpha) \to \varphi_{121}) \to \mathrm{Conv}(a,\alpha) \to \mathrm{Conv}(b,\beta) \to \varphi_{35}$
\item $\varphi_{203} :\equiv N \in \mathbb{N} \land \forall n \in \mathbb{N}\,(N \leq n \to \varphi_{94}) \to N \in \mathbb{N} \land \forall n \in \mathbb{N}\,(N \leq n \to \varphi_{94})$
\item $\varphi_{204} :\equiv \mathrm{Conv}(b,\beta) \to M \in \mathbb{R} \land \forall n \in \mathbb{N}\,(\lvert a_{n} \rvert \leq M) \to \varepsilon \in \mathbb{R} \to 0 < \varepsilon \to \varphi_{143}$
\item $\varphi_{205} :\equiv M \in \mathbb{R} \land \forall n \in \mathbb{N}\,(\lvert a_{n} \rvert \leq M) \to \varepsilon \in \mathbb{R} \to 0 < \varepsilon \to N \in \mathbb{N} \land \forall n \in \mathbb{N}\,(N \leq n \to \varphi_{94}) \to \varphi_{30}$
\item $\varphi_{206} :\equiv \mathrm{Conv}(b,\beta) \to M \in \mathbb{R} \land \forall n \in \mathbb{N}\,(\lvert a_{n} \rvert \leq M) \to \varepsilon \in \mathbb{R} \to 0 < \varepsilon \to \varphi_{145}$
\item $\varphi_{207} :\equiv \mathrm{Conv}(b,\beta) \to M \in \mathbb{R} \land \forall n \in \mathbb{N}\,(\lvert a_{n} \rvert \leq M) \to \varepsilon \in \mathbb{R} \to 0 < \varepsilon \to \varphi_{146}$
\item $\varphi_{208} :\equiv N \in \mathbb{N} \land \forall n \in \mathbb{N}\,(N \leq n \to \varphi_{94}) \to n \in \mathbb{N} \to N \leq n \to n \in \mathbb{N} \to N \leq n \to \varphi_{94}$
\item $\varphi_{209} :\equiv M \in \mathbb{R} \land \forall n \in \mathbb{N}\,(\lvert a_{n} \rvert \leq M) \to \varepsilon \in \mathbb{R} \to 0 < \varepsilon \to N \in \mathbb{N} \land \forall n \in \mathbb{N}\,(N \leq n \to \varphi_{94}) \to \varphi_{34}$
\item $\varphi_{210} :\equiv \mathrm{Conv}(b,\beta) \to M \in \mathbb{R} \land \forall n \in \mathbb{N}\,(\lvert a_{n} \rvert \leq M) \to \varepsilon \in \mathbb{R} \to 0 < \varepsilon \to \varphi_{150}$
\item $\varphi_{211} :\equiv N \in \mathbb{N} \land \forall n \in \mathbb{N}\,(N \leq n \to \varphi_{94}) \to n \in \mathbb{N} \to N \leq n \to \forall n \in \mathbb{N}\,(N \leq n \to \varphi_{94})$
\item $\varphi_{212} :\equiv \mathrm{Conv}(b,\beta) \to M \in \mathbb{R} \land \forall n \in \mathbb{N}\,(\lvert a_{n} \rvert \leq M) \to \varepsilon \in \mathbb{R} \to 0 < \varepsilon \to \varphi_{151}$
\item $\varphi_{213} :\equiv \mathrm{Conv}(b,\beta) \to M \in \mathbb{R} \land \forall n \in \mathbb{N}\,(\lvert a_{n} \rvert \leq M) \to \varepsilon \in \mathbb{R} \to 0 < \varepsilon \to \varphi_{152}$
\item $\varphi_{214} :\equiv \mathrm{Conv}(b,\beta) \to M \in \mathbb{R} \land \forall n \in \mathbb{N}\,(\lvert a_{n} \rvert \leq M) \to \varepsilon \in \mathbb{R} \to 0 < \varepsilon \to \varphi_{155}$
\item $\varphi_{215} :\equiv \mathrm{Conv}(b,\beta) \to M \in \mathbb{R} \land \forall n \in \mathbb{N}\,(\lvert a_{n} \rvert \leq M) \to \varepsilon \in \mathbb{R} \to 0 < \varepsilon \to \varphi_{157}$
\item $\varphi_{216} :\equiv \mathrm{Conv}(b,\beta) \to M \in \mathbb{R} \land \forall n \in \mathbb{N}\,(\lvert a_{n} \rvert \leq M) \to \varepsilon \in \mathbb{R} \to 0 < \varepsilon \to \varphi_{158}$
\item $\varphi_{217} :\equiv \mathrm{Conv}(b,\beta) \to M \in \mathbb{R} \land \forall n \in \mathbb{N}\,(\lvert a_{n} \rvert \leq M) \to \varepsilon \in \mathbb{R} \to 0 < \varepsilon \to \varphi_{160}$
\item $\varphi_{218} :\equiv \varphi_{28} \land \varepsilon \in \mathbb{R} \land 0 < \varepsilon \land (N \in \mathbb{N} \land \forall n \in \mathbb{N}\,(N \leq n \to \varphi_{94})) \to n \in \mathbb{N} \to N \leq n \to \lvert a \otimes b(n) - \alpha \cdot \beta \rvert < \varepsilon$
\item $\varphi_{219} :\equiv \mathrm{Conv}(b,\beta) \to M \in \mathbb{R} \land \forall n \in \mathbb{N}\,(\lvert a_{n} \rvert \leq M) \to \varepsilon \in \mathbb{R} \to 0 < \varepsilon \to \varphi_{163}$
\item $\varphi_{220} :\equiv \mathrm{Conv}(b,\beta) \to M \in \mathbb{R} \land \forall n \in \mathbb{N}\,(\lvert a_{n} \rvert \leq M) \to \varepsilon \in \mathbb{R} \to 0 < \varepsilon \to \varphi_{165}$
\item $\varphi_{221} :\equiv (0 \leq \lvert a_{1_{\omega}} \rvert \to \lvert a_{1_{\omega}} \rvert \leq M \to 0 \leq M) \to \varphi_{183}$
\item $\varphi_{222} :\equiv \mathrm{Conv}(a,\alpha) \to \mathrm{Conv}(b,\beta) \to M \in \mathbb{R} \land \forall n \in \mathbb{N}\,(\lvert a_{n} \rvert \leq M) \to \varphi_{174}$
\item $\varphi_{223} :\equiv (\mathrm{Conv}(a,\alpha) \to \mathrm{Conv}(b,\beta) \to \beta \in \mathbb{R}) \to \mathrm{Conv}(a,\alpha) \to \varphi_{185}$
\item $\varphi_{224} :\equiv \mathrm{Conv}(b,\beta) \to M \in \mathbb{R} \land \forall n \in \mathbb{N}\,(\lvert a_{n} \rvert \leq M) \to \varepsilon \in \mathbb{R} \to \varphi_{180}$
\item $\varphi_{225} :\equiv \mathrm{Conv}(b,\beta) \to M \in \mathbb{R} \land \forall n \in \mathbb{N}\,(\lvert a_{n} \rvert \leq M) \to \varepsilon \in \mathbb{R} \to \varphi_{194}$
\item $\varphi_{226} :\equiv \mathrm{Conv}(b,\beta) \to M \in \mathbb{R} \land \forall n \in \mathbb{N}\,(\lvert a_{n} \rvert \leq M) \to \varepsilon \in \mathbb{R} \to \varphi_{197}$
\item $\varphi_{227} :\equiv \mathrm{Conv}(b,\beta) \to M \in \mathbb{R} \land \forall n \in \mathbb{N}\,(\lvert a_{n} \rvert \leq M) \to \varepsilon \in \mathbb{R} \to \varphi_{201}$
\item $\varphi_{228} :\equiv \mathrm{Conv}(b,\beta) \to M \in \mathbb{R} \land \forall n \in \mathbb{N}\,(\lvert a_{n} \rvert \leq M) \to \varepsilon \in \mathbb{R} \to 0 < \varepsilon \to \varphi_{203}$
\item $\varphi_{229} :\equiv \mathrm{Conv}(b,\beta) \to M \in \mathbb{R} \land \forall n \in \mathbb{N}\,(\lvert a_{n} \rvert \leq M) \to \varepsilon \in \mathbb{R} \to 0 < \varepsilon \to \varphi_{208}$
\item $\varphi_{230} :\equiv \mathrm{Conv}(b,\beta) \to M \in \mathbb{R} \land \forall n \in \mathbb{N}\,(\lvert a_{n} \rvert \leq M) \to \varepsilon \in \mathbb{R} \to 0 < \varepsilon \to \varphi_{211}$
\item $\varphi_{231} :\equiv (\mathrm{Conv}(a,\alpha) \to \mathrm{Conv}(b,\beta) \to \varphi_{93}) \to \mathrm{Conv}(a,\alpha) \to \mathrm{Conv}(b,\beta) \to \varphi_{161}$
\item $\varphi_{232} :\equiv (M \in \mathbb{R} \to \lvert \beta \rvert \in \mathbb{R} \to M + \lvert \beta \rvert \in \mathbb{R}) \to (\mathrm{Conv}(a,\alpha) \to \varphi_{29}) \to \varphi_{171}$
\item $\varphi_{233} :\equiv (a \otimes b(n) = a_{n} \cdot b_{n} \to a_{n} \cdot b_{n} = a \otimes b(n)) \to \mathrm{Conv}(a,\alpha) \to \varphi_{207}$
\item $\varphi_{234} :\equiv (\forall M\,(\varphi_{9} \to \neg (M \in \mathbb{R} \land \forall n \in \mathbb{N}\,(\lvert a_{n} \rvert \leq M))) \to \varphi_{92}) \to (\varphi_{9} \to \forall M\,(\varphi_{9})) \to \varphi_{8}$
\item $\varphi_{235} :\equiv (\mathrm{Conv}(a,\alpha) \to \varphi_{53}) \to \mathrm{Conv}(a,\alpha) \to \mathrm{Conv}(b,\beta) \to \varphi_{209}$
\item $\varphi_{236} :\equiv \varphi_{28} \land \varepsilon \in \mathbb{R} \land 0 < \varepsilon \land (N \in \mathbb{N} \land \forall n \in \mathbb{N}\,(N \leq n \to \varphi_{94})) \to \varphi_{173}$
\item $\varphi_{237} :\equiv (M \in \mathbb{R} \to \lvert \beta \rvert \in \mathbb{R} \to 0 \leq M \to 0 \leq \lvert \beta \rvert \to 0 \leq M + \lvert \beta \rvert) \to (\mathrm{Conv}(a,\alpha) \to \varphi_{29}) \to (\mathrm{Conv}(a,\alpha) \to \varphi_{38}) \to \varphi_{176}$
\item $\varphi_{238} :\equiv (\mathrm{Conv}(a,\alpha) \to \mathrm{Conv}(b,\beta) \to \varphi_{90}) \to (\mathrm{Conv}(a,\alpha) \to \mathrm{Conv}(b,\beta) \to \varphi_{55}) \to \varphi_{182}$
\item $\varphi_{239} :\equiv (\mathrm{Conv}(a,\alpha) \to \mathrm{Conv}(b,\beta) \to \mathrm{Conv}(b,\beta)) \to (\mathrm{Conv}(a,\alpha) \to \mathrm{Conv}(b,\beta) \to \varphi_{96}) \to \varphi_{231}$
\item $\varphi_{240} :\equiv \varphi_{32} \to \forall M\,(\varphi_{9} \to \neg (M \in \mathbb{R} \land \forall n \in \mathbb{N}\,(\lvert a_{n} \rvert \leq M))) \to \varphi_{234}$
\item $\varphi_{241} :\equiv (\mathrm{Conv}(a,\alpha) \to \mathrm{Conv}(b,\beta) \to \mathrm{Seq}(b)) \to (\mathrm{Conv}(a,\alpha) \to \varphi_{178}) \to \mathrm{Conv}(a,\alpha) \to \varphi_{187}$
\item $\varphi_{242} :\equiv (\mathrm{Conv}(a,\alpha) \to \mathrm{Conv}(b,\beta) \to \mathrm{Seq}(a)) \to (\mathrm{Conv}(a,\alpha) \to \varphi_{178}) \to \mathrm{Conv}(a,\alpha) \to \varphi_{188}$
\item $\varphi_{243} :\equiv \forall n\,(\varphi_{218}) \to \varphi_{173} \to \forall n \in \mathbb{N}\,(N \leq n \to \lvert a \otimes b(n) - \alpha \cdot \beta \rvert < \varepsilon)$
\item $\varphi_{244} :\equiv (\mathrm{Conv}(a,\alpha) \to \mathrm{Conv}(b,\beta) \to \varphi_{205}) \to \mathrm{Conv}(a,\alpha) \to \varphi_{212}$
\item $\varphi_{245} :\equiv (M + \lvert \beta \rvert \in \mathbb{R} \to M + \lvert \beta \rvert < M + \lvert \beta \rvert + 1) \to (\mathrm{Conv}(a,\alpha) \to \varphi_{191}) \to \mathrm{Conv}(a,\alpha) \to \varphi_{200}$
\item $\varphi_{246} :\equiv (\mathrm{Conv}(a,\alpha) \to \varphi_{65}) \to (\mathrm{Conv}(a,\alpha) \to \varphi_{178}) \to \mathrm{Conv}(a,\alpha) \to \varphi_{184}$
\item $\varphi_{247} :\equiv (\mathrm{Conv}(a,\alpha) \to \varphi_{29}) \to (\mathrm{Conv}(a,\alpha) \to \varphi_{185}) \to \mathrm{Conv}(a,\alpha) \to \varphi_{191}$
\item $\varphi_{248} :\equiv (\forall n \in \mathbb{N}\,(N \leq n \to \varphi_{94}) \to n \in \mathbb{N} \to N \leq n \to \varphi_{94}) \to \mathrm{Conv}(a,\alpha) \to \varphi_{229}$
\item $\varphi_{249} :\equiv (\mathrm{Conv}(a,\alpha) \to \varphi_{220}) \to \varphi_{127} \to \neg (N \in \mathbb{N} \land \forall n \in \mathbb{N}\,(N \leq n \to \varphi_{94}))$
\item $\varphi_{250} :\equiv \varphi_{126} \to (\mathrm{Conv}(a,\alpha) \to \varphi_{40}) \to (\mathrm{Conv}(a,\alpha) \to \varphi_{45}) \to \varphi_{202}$
\item $\varphi_{251} :\equiv (M + \lvert \beta \rvert \in \mathbb{R} \to 1 \in \mathbb{R} \to M + \lvert \beta \rvert + 1 \in \mathbb{R}) \to (\mathrm{Conv}(a,\alpha) \to \varphi_{191}) \to 1 \in \mathbb{R} \to \mathrm{Conv}(a,\alpha) \to \varphi_{196}$
\item $\varphi_{252} :\equiv \forall N\,(\varphi_{127} \to \neg (N \in \mathbb{N} \land \forall n \in \mathbb{N}\,(N \leq n \to \varphi_{94}))) \to \forall N\,(\varphi_{127}) \to \forall N\,(\neg (N \in \mathbb{N} \land \forall n \in \mathbb{N}\,(N \leq n \to \varphi_{94})))$
\item $\varphi_{253} :\equiv (\mathrm{Conv}(a,\alpha) \to \varphi_{196}) \to (\mathrm{Conv}(a,\alpha) \to \varphi_{53}) \to \mathrm{Conv}(a,\alpha) \to \varphi_{198}$
\item $\varphi_{254} :\equiv (\mathrm{Conv}(a,\alpha) \to \varphi_{64}) \to (\mathrm{Conv}(a,\alpha) \to \varphi_{198}) \to \mathrm{Conv}(a,\alpha) \to \varphi_{204}$
\item $\varphi_{255} :\equiv (\forall \varepsilon\,(\varphi_{28} \to \varepsilon \in \mathbb{R} \to \varphi_{74}) \to \forall \varepsilon\,(\varphi_{28}) \to \forall \varepsilon \in \mathbb{R}\,(\varphi_{74})) \to (\varphi_{28} \to \forall \varepsilon\,(\varphi_{28})) \to \mathrm{Conv}(a,\alpha) \to \varphi_{121}$
\item $\varphi_{256} :\equiv (\varepsilon \in \mathbb{R} \to (M + \lvert \beta \rvert + 1)^{-1} \in \mathbb{R} \to \varphi_{39}) \to (\mathrm{Conv}(a,\alpha) \to \varphi_{64}) \to \varphi_{238}$
\item $\varphi_{257} :\equiv (\mathrm{Conv}(a,\alpha) \to \mathrm{Conv}(b,\beta) \to \varphi_{93}) \to (\mathrm{Conv}(a,\alpha) \to \varphi_{200}) \to \mathrm{Conv}(a,\alpha) \to \varphi_{225}$
\item $\varphi_{258} :\equiv (\mathrm{Conv}(a,\alpha) \to \varphi_{219}) \to (\forall N\,(\varphi_{43}) \to \varphi_{43}) \to \mathrm{Conv}(a,\alpha) \to \varphi_{220}$
\item $\varphi_{259} :\equiv (\mathrm{Conv}(a,\alpha) \to \mathrm{Conv}(b,\beta) \to \varphi_{209}) \to (\varphi_{88} \to \varphi_{97}) \to \varphi_{222}$
\item $\varphi_{260} :\equiv (\mathrm{Conv}(a,\alpha) \to \varphi_{214}) \to (\mathrm{Conv}(a,\alpha) \to \varphi_{178}) \to \mathrm{Conv}(a,\alpha) \to \varphi_{206}$
\item $\varphi_{261} :\equiv (\mathrm{Conv}(a,\alpha) \to \varphi_{225}) \to \varphi_{154} \to \mathrm{Conv}(a,\alpha) \to \mathrm{Conv}(b,\beta) \to \varphi_{205}$
\item $\varphi_{262} :\equiv (\mathrm{Conv}(a,\alpha) \to \varphi_{178}) \to (\mathrm{Conv}(a,\alpha) \to \varphi_{177}) \to \mathrm{Conv}(a,\alpha) \to \varphi_{227}$
\item $\varphi_{263} :\equiv (\varepsilon \in \mathbb{R} \to M + \lvert \beta \rvert + 1 \in \mathbb{R} \to \varphi_{48}) \to (\mathrm{Conv}(a,\alpha) \to \varphi_{64}) \to (\mathrm{Conv}(a,\alpha) \to \varphi_{196}) \to \varphi_{235}$
\item $\varphi_{264} :\equiv (\mathrm{Conv}(a,\alpha) \to \varphi_{196}) \to (\mathrm{Conv}(a,\alpha) \to \varphi_{204}) \to \mathrm{Conv}(a,\alpha) \to \varphi_{226}$
\item $\varphi_{265} :\equiv (\mathrm{Conv}(a,\alpha) \to \varphi_{216}) \to (\mathrm{Conv}(a,\alpha) \to \varphi_{215}) \to \mathrm{Conv}(a,\alpha) \to \varphi_{224}$
\item $\varphi_{266} :\equiv (\mathrm{Conv}(a,\alpha) \to \varphi_{207}) \to (\mathrm{Conv}(a,\alpha) \to \varphi_{212}) \to \varphi_{108} \to \mathrm{Conv}(a,\alpha) \to \varphi_{210}$
\item $\varphi_{267} :\equiv (M + \lvert \beta \rvert + 1 \in \mathbb{R} \to \varepsilon \cdot (M + \lvert \beta \rvert + 1)^{-1} \in \mathbb{R} \to \varphi_{88}) \to \varphi_{264}$
\item $\varphi_{268} :\equiv \forall N\,(\varphi_{127} \to \neg (N \in \mathbb{N} \land \forall n \in \mathbb{N}\,(N \leq n \to \varphi_{94}))) \to \varphi_{252} \to (\varphi_{127} \to \forall N\,(\varphi_{127})) \to \mathrm{Conv}(a,\alpha) \to \varphi_{119}$
\item $\varphi_{269} :\equiv (\mathrm{Conv}(a,\alpha) \to \varphi_{204}) \to (\mathrm{Conv}(a,\alpha) \to \varphi_{20}) \to (\mathrm{Conv}(a,\alpha) \to \varphi_{184}) \to \varphi_{265}$
\item $\varphi_{270} :\equiv (\mathrm{Conv}(a,\alpha) \to \varphi_{191}) \to (\mathrm{Conv}(a,\alpha) \to \varphi_{196}) \to (\mathrm{Conv}(a,\alpha) \to \varphi_{204}) \to \varphi_{257}$
\item $\varphi_{271} :\equiv (\mathrm{Conv}(a,\alpha) \to \mathrm{Conv}(b,\beta) \to \beta \in \mathbb{R}) \to (\mathrm{Conv}(a,\alpha) \to \varphi_{29}) \to \varphi_{269}$
\item $\varphi_{272} :\equiv (\mathrm{Conv}(a,\alpha) \to \varphi_{187}) \to (\mathrm{Conv}(a,\alpha) \to \mathrm{Conv}(b,\beta) \to \alpha \in \mathbb{R}) \to \varphi_{271}$
\item $\varphi_{273} :\equiv (a_{n} \in \mathbb{R} \to b_{n} \in \mathbb{R} \to \alpha \in \mathbb{R} \to \beta \in \mathbb{R} \to \varphi_{195}) \to (\mathrm{Conv}(a,\alpha) \to \varphi_{188}) \to \varphi_{272}$
\end{abbrevlist}
\end{fullproof}

% 自動生成 (prover/build.js): LimRecip  ノード数 930
\begin{fullproof}{定理2.9(逆数)}{$\forall b\,\forall \beta\,(\mathrm{Conv}(b,\beta) \to \beta \neq 0 \to \forall n \in \mathbb{N}\,(b_{n} \neq 0) \to \mathrm{Conv}(b^{\ast},\beta^{-1}))$}{930}\phantomsection\label{tree:LimRecip}
\noindent$\Pi_{1}$:
\begin{ptree}
\AXC{$\varphi_{303}$\lbl{A1}}
\AXC{$\varphi_{304}$\lbl{A1}}
\AXC{$\varphi_{85}$\lbl{A2}}
\AXC{$\varphi_{350}$\lbl{A1}}
\AXC{$\varphi_{350} \to \varphi_{363}$\lbl{A1}}
\BIC{$\varphi_{363}$}
\AXC{$\varphi_{363} \to \varphi_{85} \to \varphi_{357}$\lbl{A1}}
\BIC{$\varphi_{85} \to \varphi_{357}$}
\BIC{$\varphi_{357}$}
\AXC{$\varphi_{357} \to \varphi_{304} \to \varphi_{303} \to \varphi_{330}$\lbl{A1}}
\BIC{$\varphi_{304} \to \varphi_{303} \to \varphi_{330}$}
\BIC{$\varphi_{303} \to \varphi_{330}$}
\BIC{$\varphi_{330}$}
\end{ptree}
\noindent$\Pi_{2}$:
\begin{ptree}
\AXC{$\forall x\,\forall y\,\forall z\,(x < y \to y \leq z \to x < z)$\lbl{[O3]}}
\AXC{$\varphi_{99}$\lbl{A2}}
\BIC{$\forall y\,\forall z\,(\lvert b_{n} - \beta \rvert < y \to y \leq z \to \lvert b_{n} - \beta \rvert < z)$}
\AXC{$\varphi_{197}$\lbl{A2}}
\BIC{$\forall z\,(\lvert b_{n} - \beta \rvert < d_{0} \to d_{0} \leq z \to \lvert b_{n} - \beta \rvert < z)$}
\AXC{$\forall z\,(\lvert b_{n} - \beta \rvert < d_{0} \to d_{0} \leq z \to \lvert b_{n} - \beta \rvert < z) \to \varphi_{82}$\lbl{A2}}
\BIC{$\varphi_{82}$}
\end{ptree}
\noindent$\Pi_{3}$:
\begin{ptree}
\AXC{$\forall a\,\forall \alpha\,(\mathrm{Conv}(a,\alpha) \to \alpha \in \mathbb{R})$\lbl{[C2]}}
\AXC{$\varphi_{9}$\lbl{A2}}
\BIC{$\forall \alpha\,(\mathrm{Conv}(b,\alpha) \to \alpha \in \mathbb{R})$}
\AXC{$\varphi_{4}$\lbl{A2}}
\BIC{$\mathrm{Conv}(b,\beta) \to \beta \in \mathbb{R}$}
\AXC{$(\mathrm{Conv}(b,\beta) \to \beta \in \mathbb{R}) \to \varphi_{71}$\lbl{A1}}
\BIC{$\varphi_{71}$}
\end{ptree}
\noindent$\Pi_{4}$:
\begin{ptree}
\AXC{$\mathrm{Conv}(b,\beta) \to \mathrm{Conv}(b,\beta)$\lbl{A1}}
\AXC{$\stackrel{\Pi_{3}}{\vdots}$}\noLine
\UIC{$\varphi_{71}$}
\BIC{$\mathrm{Conv}(b,\beta) \to \beta \neq 0 \to \forall n \in \mathbb{N}\,(b_{n} \neq 0) \to \beta \in \mathbb{R}$}
\end{ptree}
\noindent$\Pi_{5}$:
\begin{ptree}
\AXC{$\mathrm{Conv}(b,\beta) \to \mathrm{Conv}(b,\beta)$\lbl{A1}}
\AXC{$\forall a\,\forall \alpha\,(\mathrm{Conv}(a,\alpha) \to \mathrm{Seq}(a))$\lbl{[C1]}}
\AXC{$\varphi_{2}$\lbl{A2}}
\BIC{$\forall \alpha\,(\mathrm{Conv}(b,\alpha) \to \mathrm{Seq}(b))$}
\AXC{$\varphi_{1}$\lbl{A2}}
\BIC{$\mathrm{Conv}(b,\beta) \to \mathrm{Seq}(b)$}
\AXC{$\varphi_{126}$\lbl{A1}}
\BIC{$(\mathrm{Conv}(b,\beta) \to \mathrm{Conv}(b,\beta)) \to \varphi_{5}$}
\BIC{$\varphi_{5}$}
\end{ptree}
\noindent$\Pi_{6}$:
\begin{ptree}
\AXC{$\forall a\,\forall n\,(\mathrm{Seq}(a) \to n \in \mathbb{N} \to a_{n} \in \mathbb{R})$\lbl{[S1]}}
\AXC{$\varphi_{72}$\lbl{A2}}
\BIC{$\forall n\,(\mathrm{Seq}(b) \to n \in \mathbb{N} \to b_{n} \in \mathbb{R})$}
\AXC{$\varphi_{48}$\lbl{A2}}
\BIC{$\mathrm{Seq}(b) \to n \in \mathbb{N} \to b_{n} \in \mathbb{R}$}
\AXC{$(\mathrm{Seq}(b) \to n \in \mathbb{N} \to b_{n} \in \mathbb{R}) \to \varphi_{5} \to \varphi_{304} \to \varphi_{314}$\lbl{A1}}
\BIC{$\varphi_{5} \to \varphi_{304} \to \varphi_{314}$}
\end{ptree}
\noindent$\Pi_{7}$:
\begin{ptree}
\AXC{$\varphi_{304}$\lbl{A1}}
\AXC{$\stackrel{\Pi_{5}}{\vdots}$}\noLine
\UIC{$\varphi_{5}$}
\AXC{$\stackrel{\Pi_{6}}{\vdots}$}\noLine
\UIC{$\varphi_{5} \to \varphi_{304} \to \varphi_{314}$}
\BIC{$\varphi_{304} \to \varphi_{314}$}
\BIC{$\varphi_{314}$}
\end{ptree}
\noindent$\Pi_{8}$:
\begin{ptree}
\AXC{$\varphi_{105}$\lbl{[E26]}}
\AXC{$\varphi_{105} \to \forall \beta\,(b_{n} \in \mathbb{R} \to \varphi_{83})$\lbl{A2}}
\BIC{$\forall \beta\,(b_{n} \in \mathbb{R} \to \varphi_{83})$}
\AXC{$\forall \beta\,(b_{n} \in \mathbb{R} \to \varphi_{83}) \to b_{n} \in \mathbb{R} \to \varphi_{83}$\lbl{A2}}
\BIC{$b_{n} \in \mathbb{R} \to \varphi_{83}$}
\AXC{$(b_{n} \in \mathbb{R} \to \varphi_{83}) \to \varphi_{314} \to \varphi_{428}$\lbl{A1}}
\BIC{$\varphi_{314} \to \varphi_{428}$}
\end{ptree}
\noindent$\Pi_{9}$:
\begin{ptree}
\AXC{$\stackrel{\Pi_{4}}{\vdots}$}\noLine
\UIC{$\mathrm{Conv}(b,\beta) \to \beta \neq 0 \to \forall n \in \mathbb{N}\,(b_{n} \neq 0) \to \beta \in \mathbb{R}$}
\AXC{$\stackrel{\Pi_{7}}{\vdots}$}\noLine
\UIC{$\varphi_{314}$}
\AXC{$\stackrel{\Pi_{8}}{\vdots}$}\noLine
\UIC{$\varphi_{314} \to \varphi_{428}$}
\BIC{$\varphi_{428}$}
\BIC{$\varphi_{351} \to \varphi_{342}$}
\end{ptree}
\noindent$\Pi_{10}$:
\begin{ptree}
\AXC{$\forall x \in \mathbb{R}\,(x \neq 0 \to 0 < \lvert x \rvert)$\lbl{[A14]}}
\AXC{$\varphi_{20}$\lbl{A2}}
\BIC{$\beta \in \mathbb{R} \to \beta \neq 0 \to 0 < \lvert \beta \rvert$}
\AXC{$(\beta \in \mathbb{R} \to \beta \neq 0 \to 0 < \lvert \beta \rvert) \to \varphi_{233}$\lbl{A1}}
\BIC{$\varphi_{233}$}
\end{ptree}
\noindent$\Pi_{11}$:
\begin{ptree}
\AXC{$\mathrm{Conv}(b,\beta) \to \beta \neq 0 \to \beta \neq 0$\lbl{A1}}
\AXC{$\stackrel{\Pi_{4}}{\vdots}$}\noLine
\UIC{$\mathrm{Conv}(b,\beta) \to \beta \neq 0 \to \forall n \in \mathbb{N}\,(b_{n} \neq 0) \to \beta \in \mathbb{R}$}
\AXC{$\stackrel{\Pi_{10}}{\vdots}$}\noLine
\UIC{$\varphi_{233}$}
\BIC{$(\mathrm{Conv}(b,\beta) \to \beta \neq 0 \to \beta \neq 0) \to \varphi_{10}$}
\BIC{$\varphi_{10}$}
\end{ptree}
\noindent$\Pi_{12}$:
\begin{ptree}
\AXC{$\forall x \in \mathbb{R}\,(\lvert x \rvert \in \mathbb{R})$\lbl{[A5]}}
\AXC{$\forall x \in \mathbb{R}\,(\lvert x \rvert \in \mathbb{R}) \to \beta \in \mathbb{R} \to \lvert \beta \rvert \in \mathbb{R}$\lbl{A2}}
\BIC{$\beta \in \mathbb{R} \to \lvert \beta \rvert \in \mathbb{R}$}
\AXC{$(\beta \in \mathbb{R} \to \lvert \beta \rvert \in \mathbb{R}) \to \varphi_{201}$\lbl{A1}}
\BIC{$\varphi_{201}$}
\end{ptree}
\noindent$\Pi_{13}$:
\begin{ptree}
\AXC{$\stackrel{\Pi_{4}}{\vdots}$}\noLine
\UIC{$\mathrm{Conv}(b,\beta) \to \beta \neq 0 \to \forall n \in \mathbb{N}\,(b_{n} \neq 0) \to \beta \in \mathbb{R}$}
\AXC{$\stackrel{\Pi_{12}}{\vdots}$}\noLine
\UIC{$\varphi_{201}$}
\BIC{$\varphi_{15}$}
\end{ptree}
\noindent$\Pi_{14}$:
\begin{ptree}
\AXC{$\stackrel{\Pi_{13}}{\vdots}$}\noLine
\UIC{$\varphi_{15}$}
\AXC{$\forall e \in \mathbb{R}\,(0 < e \to 0 < e \cdot 2^{-1})$\lbl{[F35]}}
\AXC{$\varphi_{114}$\lbl{A2}}
\BIC{$\lvert \beta \rvert \in \mathbb{R} \to 0 < \lvert \beta \rvert \to 0 < \lvert \beta \rvert \cdot 2^{-1}$}
\AXC{$(\lvert \beta \rvert \in \mathbb{R} \to 0 < \lvert \beta \rvert \to 0 < \lvert \beta \rvert \cdot 2^{-1}) \to \varphi_{15} \to \varphi_{10} \to \varphi_{22}$\lbl{A1}}
\BIC{$\varphi_{15} \to \varphi_{10} \to \varphi_{22}$}
\BIC{$\varphi_{10} \to \varphi_{22}$}
\end{ptree}
\noindent$\Pi_{15}$:
\begin{ptree}
\AXC{$\stackrel{\Pi_{11}}{\vdots}$}\noLine
\UIC{$\varphi_{10}$}
\AXC{$\stackrel{\Pi_{14}}{\vdots}$}\noLine
\UIC{$\varphi_{10} \to \varphi_{22}$}
\BIC{$\varphi_{22}$}
\end{ptree}
\noindent$\Pi_{16}$:
\begin{ptree}
\AXC{$\forall x \in \mathbb{R}\,(\lvert x \rvert \in \mathbb{R})$\lbl{[A5]}}
\AXC{$\forall x \in \mathbb{R}\,(\lvert x \rvert \in \mathbb{R}) \to b_{n} \in \mathbb{R} \to \lvert b_{n} \rvert \in \mathbb{R}$\lbl{A2}}
\BIC{$b_{n} \in \mathbb{R} \to \lvert b_{n} \rvert \in \mathbb{R}$}
\AXC{$(b_{n} \in \mathbb{R} \to \lvert b_{n} \rvert \in \mathbb{R}) \to \varphi_{314} \to \varphi_{320}$\lbl{A1}}
\BIC{$\varphi_{314} \to \varphi_{320}$}
\end{ptree}
\noindent$\Pi_{17}$:
\begin{ptree}
\AXC{$\stackrel{\Pi_{7}}{\vdots}$}\noLine
\UIC{$\varphi_{314}$}
\AXC{$\stackrel{\Pi_{16}}{\vdots}$}\noLine
\UIC{$\varphi_{314} \to \varphi_{320}$}
\BIC{$\varphi_{320}$}
\end{ptree}
\noindent$\Pi_{18}$:
\begin{ptree}
\AXC{$\forall x \in \mathbb{R}\,(\lvert x \rvert \in \mathbb{R})$\lbl{[A5]}}
\AXC{$\forall x \in \mathbb{R}\,(\lvert x \rvert \in \mathbb{R}) \to \beta \in \mathbb{R} \to \lvert \beta \rvert \in \mathbb{R}$\lbl{A2}}
\BIC{$\beta \in \mathbb{R} \to \lvert \beta \rvert \in \mathbb{R}$}
\AXC{$(\beta \in \mathbb{R} \to \lvert \beta \rvert \in \mathbb{R}) \to \varphi_{388}$\lbl{A1}}
\BIC{$\varphi_{388}$}
\end{ptree}
\noindent$\Pi_{19}$:
\begin{ptree}
\AXC{$\stackrel{\Pi_{4}}{\vdots}$}\noLine
\UIC{$\mathrm{Conv}(b,\beta) \to \beta \neq 0 \to \forall n \in \mathbb{N}\,(b_{n} \neq 0) \to \beta \in \mathbb{R}$}
\AXC{$\stackrel{\Pi_{18}}{\vdots}$}\noLine
\UIC{$\varphi_{388}$}
\BIC{$\varphi_{310}$}
\end{ptree}
\noindent$\Pi_{20}$:
\begin{ptree}
\AXC{$2^{-1} \in \mathbb{R}$\lbl{[F31]}}
\AXC{$\stackrel{\Pi_{19}}{\vdots}$}\noLine
\UIC{$\varphi_{310}$}
\AXC{$\forall x\,\forall y\,(x \in \mathbb{R} \to y \in \mathbb{R} \to x \cdot y \in \mathbb{R})$\lbl{(R1)}}
\AXC{$\varphi_{139}$\lbl{A2}}
\BIC{$\forall y\,(\lvert \beta \rvert \in \mathbb{R} \to y \in \mathbb{R} \to \lvert \beta \rvert \cdot y \in \mathbb{R})$}
\AXC{$\varphi_{187}$\lbl{A2}}
\BIC{$\lvert \beta \rvert \in \mathbb{R} \to 2^{-1} \in \mathbb{R} \to \lvert \beta \rvert \cdot 2^{-1} \in \mathbb{R}$}
\AXC{$\varphi_{427}$\lbl{A1}}
\BIC{$\varphi_{310} \to 2^{-1} \in \mathbb{R} \to \varphi_{332}$}
\BIC{$2^{-1} \in \mathbb{R} \to \varphi_{332}$}
\BIC{$\varphi_{332}$}
\end{ptree}
\noindent$\Pi_{21}$:
\begin{ptree}
\AXC{$\stackrel{\Pi_{15}}{\vdots}$}\noLine
\UIC{$\varphi_{22}$}
\AXC{$\stackrel{\Pi_{17}}{\vdots}$}\noLine
\UIC{$\varphi_{320}$}
\AXC{$\stackrel{\Pi_{20}}{\vdots}$}\noLine
\UIC{$\varphi_{332}$}
\AXC{$\varphi_{16}$\lbl{[E19]}}
\AXC{$\varphi_{16} \to \varphi_{196}$\lbl{A2}}
\BIC{$\varphi_{196}$}
\AXC{$\varphi_{196} \to \lvert \beta \rvert \cdot 2^{-1} \in \mathbb{R} \to \lvert b_{n} \rvert \in \mathbb{R} \to \varphi_{143}$\lbl{A2}}
\BIC{$\lvert \beta \rvert \cdot 2^{-1} \in \mathbb{R} \to \lvert b_{n} \rvert \in \mathbb{R} \to \varphi_{143}$}
\AXC{$\varphi_{436}$\lbl{A1}}
\BIC{$\varphi_{332} \to \varphi_{320} \to \varphi_{22} \to \varphi_{342} \to \varphi_{360}$}
\BIC{$\varphi_{320} \to \varphi_{22} \to \varphi_{342} \to \varphi_{360}$}
\BIC{$\varphi_{22} \to \varphi_{342} \to \varphi_{360}$}
\BIC{$\varphi_{342} \to \varphi_{360}$}
\end{ptree}
\noindent$\Pi_{22}$:
\begin{ptree}
\AXC{$\varphi_{305}$\lbl{A1}}
\AXC{$\varphi_{305} \to \varphi_{318}$\lbl{A1}}
\BIC{$\varphi_{318}$}
\AXC{$\stackrel{\Pi_{1}}{\vdots}$}\noLine
\UIC{$\varphi_{330}$}
\AXC{$\stackrel{\Pi_{2}}{\vdots}$}\noLine
\UIC{$\varphi_{82}$}
\AXC{$\varphi_{82} \to \varphi_{330} \to \varphi_{318} \to \varphi_{351}$\lbl{A1}}
\BIC{$\varphi_{330} \to \varphi_{318} \to \varphi_{351}$}
\BIC{$\varphi_{318} \to \varphi_{351}$}
\BIC{$\varphi_{351}$}
\AXC{$\stackrel{\Pi_{9}}{\vdots}$}\noLine
\UIC{$\varphi_{351} \to \varphi_{342}$}
\BIC{$\varphi_{342}$}
\AXC{$\stackrel{\Pi_{21}}{\vdots}$}\noLine
\UIC{$\varphi_{342} \to \varphi_{360}$}
\BIC{$\varphi_{360}$}
\end{ptree}
\noindent$\Pi_{23}$:
\begin{ptree}
\AXC{$\stackrel{\Pi_{11}}{\vdots}$}\noLine
\UIC{$\varphi_{10}$}
\AXC{$\stackrel{\Pi_{19}}{\vdots}$}\noLine
\UIC{$\varphi_{310}$}
\AXC{$\forall x \in \mathbb{R}\,(0 < x \to 0 < x^{-1})$\lbl{[F34]}}
\AXC{$\varphi_{80}$\lbl{A2}}
\BIC{$\lvert \beta \rvert \in \mathbb{R} \to 0 < \lvert \beta \rvert \to 0 < \lvert \beta \rvert^{-1}$}
\AXC{$(\lvert \beta \rvert \in \mathbb{R} \to 0 < \lvert \beta \rvert \to 0 < \lvert \beta \rvert^{-1}) \to \varphi_{310} \to \varphi_{10} \to \varphi_{315}$\lbl{A1}}
\BIC{$\varphi_{310} \to \varphi_{10} \to \varphi_{315}$}
\BIC{$\varphi_{10} \to \varphi_{315}$}
\BIC{$\varphi_{315}$}
\end{ptree}
\noindent$\Pi_{24}$:
\begin{ptree}
\AXC{$x = x$\lbl{A5}}
\AXC{$x = y \to x = x \to y = x$\lbl{A6}}
\AXC{$(x = y \to x = x \to y = x) \to x = x \to x = y \to y = x$\lbl{A1}}
\BIC{$x = x \to x = y \to y = x$}
\BIC{$x = y \to y = x$}
\UIC{$\forall y\,(x = y \to y = x)$}
\UIC{$\forall x\,\forall y\,(x = y \to y = x)$}
\end{ptree}
\noindent$\Pi_{25}$:
\begin{ptree}
\AXC{$\stackrel{\Pi_{24}}{\vdots}$}\noLine
\UIC{$\forall x\,\forall y\,(x = y \to y = x)$}
\AXC{$\forall x\,\forall y\,(x = y \to y = x) \to \forall y\,(\lvert \beta \rvert = y \to y = \lvert \beta \rvert)$\lbl{A2}}
\BIC{$\forall y\,(\lvert \beta \rvert = y \to y = \lvert \beta \rvert)$}
\end{ptree}
\noindent$\Pi_{26}$:
\begin{ptree}
\AXC{$\stackrel{\Pi_{25}}{\vdots}$}\noLine
\UIC{$\forall y\,(\lvert \beta \rvert = y \to y = \lvert \beta \rvert)$}
\AXC{$\forall y\,(\lvert \beta \rvert = y \to y = \lvert \beta \rvert) \to \lvert \beta \rvert = 0 \to 0 = \lvert \beta \rvert$\lbl{A2}}
\BIC{$\lvert \beta \rvert = 0 \to 0 = \lvert \beta \rvert$}
\end{ptree}
\noindent$\Pi_{27}$:
\begin{ptree}
\AXC{$\forall x\,\forall y\,(x < y \to x \neq y)$\lbl{[F26]}}
\AXC{$\forall x\,\forall y\,(x < y \to x \neq y) \to \forall y\,(0 < y \to 0 \neq y)$\lbl{A2}}
\BIC{$\forall y\,(0 < y \to 0 \neq y)$}
\end{ptree}
\noindent$\Pi_{28}$:
\begin{ptree}
\AXC{$\stackrel{\Pi_{11}}{\vdots}$}\noLine
\UIC{$\varphi_{10}$}
\AXC{$\stackrel{\Pi_{27}}{\vdots}$}\noLine
\UIC{$\forall y\,(0 < y \to 0 \neq y)$}
\AXC{$\forall y\,(0 < y \to 0 \neq y) \to 0 < \lvert \beta \rvert \to 0 \neq \lvert \beta \rvert$\lbl{A2}}
\BIC{$0 < \lvert \beta \rvert \to 0 \neq \lvert \beta \rvert$}
\AXC{$(0 < \lvert \beta \rvert \to 0 \neq \lvert \beta \rvert) \to \varphi_{10} \to \varphi_{7}$\lbl{A1}}
\BIC{$\varphi_{10} \to \varphi_{7}$}
\BIC{$\varphi_{7}$}
\end{ptree}
\noindent$\Pi_{29}$:
\begin{ptree}
\AXC{$\stackrel{\Pi_{26}}{\vdots}$}\noLine
\UIC{$\lvert \beta \rvert = 0 \to 0 = \lvert \beta \rvert$}
\AXC{$\stackrel{\Pi_{28}}{\vdots}$}\noLine
\UIC{$\varphi_{7}$}
\AXC{$\varphi_{7} \to (\lvert \beta \rvert = 0 \to 0 = \lvert \beta \rvert) \to \varphi_{8}$\lbl{A1}}
\BIC{$(\lvert \beta \rvert = 0 \to 0 = \lvert \beta \rvert) \to \varphi_{8}$}
\BIC{$\varphi_{8}$}
\end{ptree}
\noindent$\Pi_{30}$:
\begin{ptree}
\AXC{$\stackrel{\Pi_{19}}{\vdots}$}\noLine
\UIC{$\varphi_{310}$}
\AXC{$\forall x \in \mathbb{R}\,(x \neq 0 \to x^{-1} \in \mathbb{R})$\lbl{(R1)}}
\AXC{$\varphi_{90}$\lbl{A2}}
\BIC{$\lvert \beta \rvert \in \mathbb{R} \to \lvert \beta \rvert \neq 0 \to \lvert \beta \rvert^{-1} \in \mathbb{R}$}
\AXC{$(\lvert \beta \rvert \in \mathbb{R} \to \lvert \beta \rvert \neq 0 \to \lvert \beta \rvert^{-1} \in \mathbb{R}) \to \varphi_{310} \to \varphi_{8} \to \varphi_{319}$\lbl{A1}}
\BIC{$\varphi_{310} \to \varphi_{8} \to \varphi_{319}$}
\BIC{$\varphi_{8} \to \varphi_{319}$}
\end{ptree}
\noindent$\Pi_{31}$:
\begin{ptree}
\AXC{$\stackrel{\Pi_{29}}{\vdots}$}\noLine
\UIC{$\varphi_{8}$}
\AXC{$\stackrel{\Pi_{30}}{\vdots}$}\noLine
\UIC{$\varphi_{8} \to \varphi_{319}$}
\BIC{$\varphi_{319}$}
\end{ptree}
\noindent$\Pi_{32}$:
\begin{ptree}
\AXC{$\forall x\,\forall y\,(x < y \to x \neq y)$\lbl{[F26]}}
\AXC{$\forall x\,\forall y\,(x < y \to x \neq y) \to \forall y\,(0 < y \to 0 \neq y)$\lbl{A2}}
\BIC{$\forall y\,(0 < y \to 0 \neq y)$}
\AXC{$\varphi_{14}$\lbl{A2}}
\BIC{$0 < \lvert \beta \rvert \cdot 2^{-1} \to 0 \neq \lvert \beta \rvert \cdot 2^{-1}$}
\end{ptree}
\noindent$\Pi_{33}$:
\begin{ptree}
\AXC{$\stackrel{\Pi_{15}}{\vdots}$}\noLine
\UIC{$\varphi_{22}$}
\AXC{$\stackrel{\Pi_{32}}{\vdots}$}\noLine
\UIC{$0 < \lvert \beta \rvert \cdot 2^{-1} \to 0 \neq \lvert \beta \rvert \cdot 2^{-1}$}
\AXC{$(0 < \lvert \beta \rvert \cdot 2^{-1} \to 0 \neq \lvert \beta \rvert \cdot 2^{-1}) \to \varphi_{22} \to \varphi_{18}$\lbl{A1}}
\BIC{$\varphi_{22} \to \varphi_{18}$}
\BIC{$\varphi_{18}$}
\end{ptree}
\noindent$\Pi_{34}$:
\begin{ptree}
\AXC{$\stackrel{\Pi_{24}}{\vdots}$}\noLine
\UIC{$\forall x\,\forall y\,(x = y \to y = x)$}
\AXC{$\varphi_{11}$\lbl{A2}}
\BIC{$\forall y\,(\lvert \beta \rvert \cdot 2^{-1} = y \to y = \lvert \beta \rvert \cdot 2^{-1})$}
\AXC{$\varphi_{73}$\lbl{A2}}
\BIC{$\lvert \beta \rvert \cdot 2^{-1} = 0 \to 0 = \lvert \beta \rvert \cdot 2^{-1}$}
\AXC{$\stackrel{\Pi_{33}}{\vdots}$}\noLine
\UIC{$\varphi_{18}$}
\AXC{$\varphi_{18} \to (\lvert \beta \rvert \cdot 2^{-1} = 0 \to 0 = \lvert \beta \rvert \cdot 2^{-1}) \to \varphi_{19}$\lbl{A1}}
\BIC{$(\lvert \beta \rvert \cdot 2^{-1} = 0 \to 0 = \lvert \beta \rvert \cdot 2^{-1}) \to \varphi_{19}$}
\BIC{$\varphi_{19}$}
\end{ptree}
\noindent$\Pi_{35}$:
\begin{ptree}
\AXC{$\stackrel{\Pi_{34}}{\vdots}$}\noLine
\UIC{$\varphi_{19}$}
\AXC{$\stackrel{\Pi_{20}}{\vdots}$}\noLine
\UIC{$\varphi_{332}$}
\AXC{$\forall x \in \mathbb{R}\,(x \neq 0 \to x^{-1} \in \mathbb{R})$\lbl{(R1)}}
\AXC{$\forall x \in \mathbb{R}\,(x \neq 0 \to x^{-1} \in \mathbb{R}) \to \varphi_{54}$\lbl{A2}}
\BIC{$\varphi_{54}$}
\AXC{$\varphi_{54} \to \varphi_{332} \to \varphi_{19} \to \varphi_{337}$\lbl{A1}}
\BIC{$\varphi_{332} \to \varphi_{19} \to \varphi_{337}$}
\BIC{$\varphi_{19} \to \varphi_{337}$}
\BIC{$\varphi_{337}$}
\end{ptree}
\noindent$\Pi_{36}$:
\begin{ptree}
\AXC{$\forall n \in \mathbb{N}\,(b_{n} \neq 0) \to n \in \mathbb{N} \to b_{n} \neq 0$\lbl{A2}}
\AXC{$\varphi_{32}$\lbl{A1}}
\AXC{$\varphi_{32} \to (\forall n \in \mathbb{N}\,(b_{n} \neq 0) \to n \in \mathbb{N} \to b_{n} \neq 0) \to \varphi_{28}$\lbl{A1}}
\BIC{$(\forall n \in \mathbb{N}\,(b_{n} \neq 0) \to n \in \mathbb{N} \to b_{n} \neq 0) \to \varphi_{28}$}
\BIC{$\varphi_{28}$}
\end{ptree}
\noindent$\Pi_{37}$:
\begin{ptree}
\AXC{$\varphi_{304}$\lbl{A1}}
\AXC{$\stackrel{\Pi_{36}}{\vdots}$}\noLine
\UIC{$\varphi_{28}$}
\AXC{$\varphi_{28} \to \varphi_{304} \to \varphi_{306}$\lbl{A1}}
\BIC{$\varphi_{304} \to \varphi_{306}$}
\BIC{$\varphi_{306}$}
\end{ptree}
\noindent$\Pi_{38}$:
\begin{ptree}
\AXC{$\stackrel{\Pi_{37}}{\vdots}$}\noLine
\UIC{$\varphi_{306}$}
\AXC{$\stackrel{\Pi_{7}}{\vdots}$}\noLine
\UIC{$\varphi_{314}$}
\AXC{$\forall x \in \mathbb{R}\,(x \neq 0 \to 0 < \lvert x \rvert)$\lbl{[A14]}}
\AXC{$\varphi_{76}$\lbl{A2}}
\BIC{$b_{n} \in \mathbb{R} \to b_{n} \neq 0 \to 0 < \lvert b_{n} \rvert$}
\AXC{$(b_{n} \in \mathbb{R} \to b_{n} \neq 0 \to 0 < \lvert b_{n} \rvert) \to \varphi_{314} \to \varphi_{306} \to \varphi_{317}$\lbl{A1}}
\BIC{$\varphi_{314} \to \varphi_{306} \to \varphi_{317}$}
\BIC{$\varphi_{306} \to \varphi_{317}$}
\BIC{$\varphi_{317}$}
\end{ptree}
\noindent$\Pi_{39}$:
\begin{ptree}
\AXC{$\forall x\,\forall y\,(x < y \to x \neq y)$\lbl{[F26]}}
\AXC{$\forall x\,\forall y\,(x < y \to x \neq y) \to \forall y\,(0 < y \to 0 \neq y)$\lbl{A2}}
\BIC{$\forall y\,(0 < y \to 0 \neq y)$}
\end{ptree}
\noindent$\Pi_{40}$:
\begin{ptree}
\AXC{$\stackrel{\Pi_{39}}{\vdots}$}\noLine
\UIC{$\forall y\,(0 < y \to 0 \neq y)$}
\AXC{$\forall y\,(0 < y \to 0 \neq y) \to 0 < \lvert b_{n} \rvert \to 0 \neq \lvert b_{n} \rvert$\lbl{A2}}
\BIC{$0 < \lvert b_{n} \rvert \to 0 \neq \lvert b_{n} \rvert$}
\AXC{$(0 < \lvert b_{n} \rvert \to 0 \neq \lvert b_{n} \rvert) \to \varphi_{317} \to \varphi_{311}$\lbl{A1}}
\BIC{$\varphi_{317} \to \varphi_{311}$}
\end{ptree}
\noindent$\Pi_{41}$:
\begin{ptree}
\AXC{$\stackrel{\Pi_{24}}{\vdots}$}\noLine
\UIC{$\forall x\,\forall y\,(x = y \to y = x)$}
\AXC{$\varphi_{3}$\lbl{A2}}
\BIC{$\forall y\,(\lvert b_{n} \rvert = y \to y = \lvert b_{n} \rvert)$}
\AXC{$\varphi_{23}$\lbl{A2}}
\BIC{$\lvert b_{n} \rvert = 0 \to 0 = \lvert b_{n} \rvert$}
\AXC{$\stackrel{\Pi_{38}}{\vdots}$}\noLine
\UIC{$\varphi_{317}$}
\AXC{$\stackrel{\Pi_{40}}{\vdots}$}\noLine
\UIC{$\varphi_{317} \to \varphi_{311}$}
\BIC{$\varphi_{311}$}
\AXC{$\varphi_{311} \to (\lvert b_{n} \rvert = 0 \to 0 = \lvert b_{n} \rvert) \to \varphi_{312}$\lbl{A1}}
\BIC{$(\lvert b_{n} \rvert = 0 \to 0 = \lvert b_{n} \rvert) \to \varphi_{312}$}
\BIC{$\varphi_{312}$}
\end{ptree}
\noindent$\Pi_{42}$:
\begin{ptree}
\AXC{$\stackrel{\Pi_{17}}{\vdots}$}\noLine
\UIC{$\varphi_{320}$}
\AXC{$\forall x \in \mathbb{R}\,(x \neq 0 \to x^{-1} \in \mathbb{R})$\lbl{(R1)}}
\AXC{$\varphi_{158}$\lbl{A2}}
\BIC{$\lvert b_{n} \rvert \in \mathbb{R} \to \lvert b_{n} \rvert \neq 0 \to \lvert b_{n} \rvert^{-1} \in \mathbb{R}$}
\AXC{$(\lvert b_{n} \rvert \in \mathbb{R} \to \lvert b_{n} \rvert \neq 0 \to \lvert b_{n} \rvert^{-1} \in \mathbb{R}) \to \varphi_{320} \to \varphi_{312} \to \varphi_{333}$\lbl{A1}}
\BIC{$\varphi_{320} \to \varphi_{312} \to \varphi_{333}$}
\BIC{$\varphi_{312} \to \varphi_{333}$}
\end{ptree}
\noindent$\Pi_{43}$:
\begin{ptree}
\AXC{$\stackrel{\Pi_{41}}{\vdots}$}\noLine
\UIC{$\varphi_{312}$}
\AXC{$\stackrel{\Pi_{42}}{\vdots}$}\noLine
\UIC{$\varphi_{312} \to \varphi_{333}$}
\BIC{$\varphi_{333}$}
\end{ptree}
\noindent$\Pi_{44}$:
\begin{ptree}
\AXC{$\stackrel{\Pi_{35}}{\vdots}$}\noLine
\UIC{$\varphi_{337}$}
\AXC{$\stackrel{\Pi_{43}}{\vdots}$}\noLine
\UIC{$\varphi_{333}$}
\AXC{$\varphi_{49}$\lbl{[F32]}}
\AXC{$\varphi_{49} \to \varphi_{202}$\lbl{A2}}
\BIC{$\varphi_{202}$}
\AXC{$\varphi_{202} \to \varphi_{272}$\lbl{A2}}
\BIC{$\varphi_{272}$}
\AXC{$\varphi_{272} \to \varphi_{292}$\lbl{A2}}
\BIC{$\varphi_{292}$}
\AXC{$\varphi_{292} \to \varphi_{333} \to \varphi_{337} \to \varphi_{319} \to \varphi_{315} \to \varphi_{360} \to \varphi_{384}$\lbl{A1}}
\BIC{$\varphi_{333} \to \varphi_{337} \to \varphi_{319} \to \varphi_{315} \to \varphi_{360} \to \varphi_{384}$}
\BIC{$\varphi_{337} \to \varphi_{319} \to \varphi_{315} \to \varphi_{360} \to \varphi_{384}$}
\BIC{$\varphi_{319} \to \varphi_{315} \to \varphi_{360} \to \varphi_{384}$}
\end{ptree}
\noindent$\Pi_{45}$:
\begin{ptree}
\AXC{$2^{-1} \in \mathbb{R}$\lbl{[F31]}}
\AXC{$\stackrel{\Pi_{13}}{\vdots}$}\noLine
\UIC{$\varphi_{15}$}
\AXC{$\forall x\,\forall y\,(x \in \mathbb{R} \to y \in \mathbb{R} \to x \cdot y \in \mathbb{R})$\lbl{(R1)}}
\AXC{$\varphi_{139}$\lbl{A2}}
\BIC{$\forall y\,(\lvert \beta \rvert \in \mathbb{R} \to y \in \mathbb{R} \to \lvert \beta \rvert \cdot y \in \mathbb{R})$}
\AXC{$\varphi_{187}$\lbl{A2}}
\BIC{$\lvert \beta \rvert \in \mathbb{R} \to 2^{-1} \in \mathbb{R} \to \lvert \beta \rvert \cdot 2^{-1} \in \mathbb{R}$}
\AXC{$\varphi_{288}$\lbl{A1}}
\BIC{$\varphi_{15} \to 2^{-1} \in \mathbb{R} \to \varphi_{33}$}
\BIC{$2^{-1} \in \mathbb{R} \to \varphi_{33}$}
\BIC{$\varphi_{33}$}
\end{ptree}
\noindent$\Pi_{46}$:
\begin{ptree}
\AXC{$\forall x\,\forall y\,(x \in \mathbb{R} \to y \in \mathbb{R} \to 0 < x \to 0 < y \to 0 < x \cdot y)$\lbl{(R15)}}
\AXC{$\varphi_{258}$\lbl{A2}}
\BIC{$\varphi_{140}$}
\AXC{$\varphi_{140} \to \lvert \beta \rvert \cdot 2^{-1} \in \mathbb{R} \to \varphi_{86}$\lbl{A2}}
\BIC{$\lvert \beta \rvert \cdot 2^{-1} \in \mathbb{R} \to \varphi_{86}$}
\end{ptree}
\noindent$\Pi_{47}$:
\begin{ptree}
\AXC{$\stackrel{\Pi_{11}}{\vdots}$}\noLine
\UIC{$\varphi_{10}$}
\AXC{$\stackrel{\Pi_{15}}{\vdots}$}\noLine
\UIC{$\varphi_{22}$}
\AXC{$\stackrel{\Pi_{13}}{\vdots}$}\noLine
\UIC{$\varphi_{15}$}
\AXC{$\stackrel{\Pi_{45}}{\vdots}$}\noLine
\UIC{$\varphi_{33}$}
\AXC{$\stackrel{\Pi_{46}}{\vdots}$}\noLine
\UIC{$\lvert \beta \rvert \cdot 2^{-1} \in \mathbb{R} \to \varphi_{86}$}
\AXC{$(\lvert \beta \rvert \cdot 2^{-1} \in \mathbb{R} \to \varphi_{86}) \to \varphi_{33} \to \varphi_{15} \to \varphi_{22} \to \varphi_{10} \to \varphi_{55}$\lbl{A1}}
\BIC{$\varphi_{33} \to \varphi_{15} \to \varphi_{22} \to \varphi_{10} \to \varphi_{55}$}
\BIC{$\varphi_{15} \to \varphi_{22} \to \varphi_{10} \to \varphi_{55}$}
\BIC{$\varphi_{22} \to \varphi_{10} \to \varphi_{55}$}
\BIC{$\varphi_{10} \to \varphi_{55}$}
\BIC{$\varphi_{55}$}
\end{ptree}
\noindent$\Pi_{48}$:
\begin{ptree}
\AXC{$\forall x \in \mathbb{R}\,(\lvert x \rvert \in \mathbb{R})$\lbl{[A5]}}
\AXC{$\forall x \in \mathbb{R}\,(\lvert x \rvert \in \mathbb{R}) \to \beta \in \mathbb{R} \to \lvert \beta \rvert \in \mathbb{R}$\lbl{A2}}
\BIC{$\beta \in \mathbb{R} \to \lvert \beta \rvert \in \mathbb{R}$}
\AXC{$(\beta \in \mathbb{R} \to \lvert \beta \rvert \in \mathbb{R}) \to \varphi_{236}$\lbl{A1}}
\BIC{$\varphi_{236}$}
\end{ptree}
\noindent$\Pi_{49}$:
\begin{ptree}
\AXC{$\stackrel{\Pi_{4}}{\vdots}$}\noLine
\UIC{$\mathrm{Conv}(b,\beta) \to \beta \neq 0 \to \forall n \in \mathbb{N}\,(b_{n} \neq 0) \to \beta \in \mathbb{R}$}
\AXC{$\stackrel{\Pi_{48}}{\vdots}$}\noLine
\UIC{$\varphi_{236}$}
\BIC{$\mathrm{Conv}(b,\beta) \to \varphi_{42}$}
\end{ptree}
\noindent$\Pi_{50}$:
\begin{ptree}
\AXC{$\forall x\,\forall y\,(x \in \mathbb{R} \to y \in \mathbb{R} \to x \cdot y \in \mathbb{R})$\lbl{(R1)}}
\AXC{$\varphi_{139}$\lbl{A2}}
\BIC{$\forall y\,(\lvert \beta \rvert \in \mathbb{R} \to y \in \mathbb{R} \to \lvert \beta \rvert \cdot y \in \mathbb{R})$}
\AXC{$\varphi_{187}$\lbl{A2}}
\BIC{$\lvert \beta \rvert \in \mathbb{R} \to 2^{-1} \in \mathbb{R} \to \lvert \beta \rvert \cdot 2^{-1} \in \mathbb{R}$}
\AXC{$(\lvert \beta \rvert \in \mathbb{R} \to 2^{-1} \in \mathbb{R} \to \lvert \beta \rvert \cdot 2^{-1} \in \mathbb{R}) \to \varphi_{283}$\lbl{A1}}
\BIC{$\varphi_{283}$}
\end{ptree}
\noindent$\Pi_{51}$:
\begin{ptree}
\AXC{$2^{-1} \in \mathbb{R}$\lbl{[F31]}}
\AXC{$\stackrel{\Pi_{49}}{\vdots}$}\noLine
\UIC{$\mathrm{Conv}(b,\beta) \to \varphi_{42}$}
\AXC{$\stackrel{\Pi_{50}}{\vdots}$}\noLine
\UIC{$\varphi_{283}$}
\BIC{$2^{-1} \in \mathbb{R} \to \mathrm{Conv}(b,\beta) \to \beta \neq 0 \to \varphi_{44}$}
\BIC{$\mathrm{Conv}(b,\beta) \to \beta \neq 0 \to \varphi_{44}$}
\end{ptree}
\noindent$\Pi_{52}$:
\begin{ptree}
\AXC{$\forall x\,\forall y\,(x \in \mathbb{R} \to y \in \mathbb{R} \to x \cdot y \in \mathbb{R})$\lbl{(R1)}}
\AXC{$\forall x\,\forall y\,(x \in \mathbb{R} \to y \in \mathbb{R} \to x \cdot y \in \mathbb{R}) \to \varphi_{34}$\lbl{A2}}
\BIC{$\varphi_{34}$}
\AXC{$\varphi_{34} \to \varphi_{61}$\lbl{A2}}
\BIC{$\varphi_{61}$}
\AXC{$\varphi_{61} \to \varphi_{381}$\lbl{A1}}
\BIC{$\varphi_{381}$}
\end{ptree}
\noindent$\Pi_{53}$:
\begin{ptree}
\AXC{$\stackrel{\Pi_{49}}{\vdots}$}\noLine
\UIC{$\mathrm{Conv}(b,\beta) \to \varphi_{42}$}
\AXC{$\stackrel{\Pi_{51}}{\vdots}$}\noLine
\UIC{$\mathrm{Conv}(b,\beta) \to \beta \neq 0 \to \varphi_{44}$}
\AXC{$\stackrel{\Pi_{52}}{\vdots}$}\noLine
\UIC{$\varphi_{381}$}
\BIC{$(\mathrm{Conv}(b,\beta) \to \varphi_{42}) \to \mathrm{Conv}(b,\beta) \to \beta \neq 0 \to \varphi_{87}$}
\BIC{$\mathrm{Conv}(b,\beta) \to \beta \neq 0 \to \varphi_{87}$}
\end{ptree}
\noindent$\Pi_{54}$:
\begin{ptree}
\AXC{$\varphi_{39}$\lbl{A1}}
\AXC{$\forall x\,\forall y\,(x \in \mathbb{R} \to y \in \mathbb{R} \to 0 < x \to 0 < y \to 0 < x \cdot y)$\lbl{(R15)}}
\AXC{$\varphi_{198}$\lbl{A2}}
\BIC{$\forall y\,(\varepsilon \in \mathbb{R} \to y \in \mathbb{R} \to 0 < \varepsilon \to 0 < y \to 0 < \varepsilon \cdot y)$}
\AXC{$\varphi_{281}$\lbl{A2}}
\BIC{$\varepsilon \in \mathbb{R} \to \lvert \beta \rvert \cdot 2^{-1} \cdot \lvert \beta \rvert \in \mathbb{R} \to 0 < \varepsilon \to \varphi_{36}$}
\AXC{$\varphi_{429}$\lbl{A1}}
\BIC{$\varphi_{39} \to \varphi_{415}$}
\BIC{$\varphi_{415}$}
\end{ptree}
\noindent$\Pi_{55}$:
\begin{ptree}
\AXC{$\stackrel{\Pi_{47}}{\vdots}$}\noLine
\UIC{$\varphi_{55}$}
\AXC{$\varphi_{56}$\lbl{A1}}
\AXC{$\stackrel{\Pi_{53}}{\vdots}$}\noLine
\UIC{$\mathrm{Conv}(b,\beta) \to \beta \neq 0 \to \varphi_{87}$}
\AXC{$\stackrel{\Pi_{54}}{\vdots}$}\noLine
\UIC{$\varphi_{415}$}
\BIC{$\varphi_{56} \to \varphi_{55} \to \mathrm{Conv}(b,\beta) \to \beta \neq 0 \to \varphi_{96}$}
\BIC{$\varphi_{55} \to \mathrm{Conv}(b,\beta) \to \beta \neq 0 \to \varphi_{96}$}
\BIC{$\mathrm{Conv}(b,\beta) \to \beta \neq 0 \to \varphi_{96}$}
\end{ptree}
\noindent$\Pi_{56}$:
\begin{ptree}
\AXC{$\forall x\,\forall y\,(x \in \mathbb{R} \to y \in \mathbb{R} \to x \cdot y \in \mathbb{R})$\lbl{(R1)}}
\AXC{$\forall x\,\forall y\,(x \in \mathbb{R} \to y \in \mathbb{R} \to x \cdot y \in \mathbb{R}) \to \varphi_{34}$\lbl{A2}}
\BIC{$\varphi_{34}$}
\AXC{$\varphi_{34} \to \varphi_{61}$\lbl{A2}}
\BIC{$\varphi_{61}$}
\end{ptree}
\noindent$\Pi_{57}$:
\begin{ptree}
\AXC{$\stackrel{\Pi_{19}}{\vdots}$}\noLine
\UIC{$\varphi_{310}$}
\AXC{$\stackrel{\Pi_{20}}{\vdots}$}\noLine
\UIC{$\varphi_{332}$}
\AXC{$\stackrel{\Pi_{56}}{\vdots}$}\noLine
\UIC{$\varphi_{61}$}
\AXC{$\varphi_{61} \to \varphi_{332} \to \varphi_{310} \to \varphi_{344}$\lbl{A1}}
\BIC{$\varphi_{332} \to \varphi_{310} \to \varphi_{344}$}
\BIC{$\varphi_{310} \to \varphi_{344}$}
\BIC{$\varphi_{344}$}
\end{ptree}
\noindent$\Pi_{58}$:
\begin{ptree}
\AXC{$\forall x\,\forall y\,(x \in \mathbb{R} \to y \in \mathbb{R} \to x \cdot y \in \mathbb{R})$\lbl{(R1)}}
\AXC{$\varphi_{108}$\lbl{A2}}
\BIC{$\forall y\,(\varepsilon \in \mathbb{R} \to y \in \mathbb{R} \to \varepsilon \cdot y \in \mathbb{R})$}
\AXC{$\forall y\,(\varepsilon \in \mathbb{R} \to y \in \mathbb{R} \to \varepsilon \cdot y \in \mathbb{R}) \to \varepsilon \in \mathbb{R} \to \varphi_{59}$\lbl{A2}}
\BIC{$\varepsilon \in \mathbb{R} \to \varphi_{59}$}
\end{ptree}
\noindent$\Pi_{59}$:
\begin{ptree}
\AXC{$\stackrel{\Pi_{57}}{\vdots}$}\noLine
\UIC{$\varphi_{344}$}
\AXC{$\varphi_{39}$\lbl{A1}}
\AXC{$\stackrel{\Pi_{58}}{\vdots}$}\noLine
\UIC{$\varepsilon \in \mathbb{R} \to \varphi_{59}$}
\AXC{$(\varepsilon \in \mathbb{R} \to \varphi_{59}) \to \varphi_{39} \to \varphi_{344} \to \varphi_{361}$\lbl{A1}}
\BIC{$\varphi_{39} \to \varphi_{344} \to \varphi_{361}$}
\BIC{$\varphi_{344} \to \varphi_{361}$}
\BIC{$\varphi_{361}$}
\end{ptree}
\noindent$\Pi_{60}$:
\begin{ptree}
\AXC{$\forall x\,\forall y\,(x \in \mathbb{R} \to y \in \mathbb{R} \to x \cdot y \in \mathbb{R})$\lbl{(R1)}}
\AXC{$\forall x\,\forall y\,(x \in \mathbb{R} \to y \in \mathbb{R} \to x \cdot y \in \mathbb{R}) \to \varphi_{66}$\lbl{A2}}
\BIC{$\varphi_{66}$}
\AXC{$\varphi_{66} \to \varphi_{132}$\lbl{A2}}
\BIC{$\varphi_{132}$}
\end{ptree}
\noindent$\Pi_{61}$:
\begin{ptree}
\AXC{$\stackrel{\Pi_{31}}{\vdots}$}\noLine
\UIC{$\varphi_{319}$}
\AXC{$\stackrel{\Pi_{43}}{\vdots}$}\noLine
\UIC{$\varphi_{333}$}
\AXC{$\forall x\,\forall y\,(x \in \mathbb{R} \to y \in \mathbb{R} \to x \cdot y \in \mathbb{R})$\lbl{(R1)}}
\AXC{$\varphi_{199}$\lbl{A2}}
\BIC{$\forall y\,(\lvert b_{n} \rvert^{-1} \in \mathbb{R} \to y \in \mathbb{R} \to \lvert b_{n} \rvert^{-1} \cdot y \in \mathbb{R})$}
\AXC{$\varphi_{263}$\lbl{A2}}
\BIC{$\lvert b_{n} \rvert^{-1} \in \mathbb{R} \to \lvert \beta \rvert^{-1} \in \mathbb{R} \to \lvert b_{n} \rvert^{-1} \cdot \lvert \beta \rvert^{-1} \in \mathbb{R}$}
\AXC{$\varphi_{434}$\lbl{A1}}
\BIC{$\varphi_{333} \to \varphi_{319} \to \varphi_{358}$}
\BIC{$\varphi_{319} \to \varphi_{358}$}
\BIC{$\varphi_{358}$}
\end{ptree}
\noindent$\Pi_{62}$:
\begin{ptree}
\AXC{$\stackrel{\Pi_{61}}{\vdots}$}\noLine
\UIC{$\varphi_{358}$}
\AXC{$\varphi_{50}$\lbl{[E8]}}
\AXC{$\varphi_{50} \to \varphi_{264}$\lbl{A2}}
\BIC{$\varphi_{264}$}
\AXC{$\varphi_{264} \to \varphi_{359}$\lbl{A2}}
\BIC{$\varphi_{359}$}
\AXC{$\varphi_{359} \to \lvert b_{n} \rvert^{-1} \cdot \lvert \beta \rvert^{-1} \in \mathbb{R} \to \varphi_{409}$\lbl{A2}}
\BIC{$\lvert b_{n} \rvert^{-1} \cdot \lvert \beta \rvert^{-1} \in \mathbb{R} \to \varphi_{409}$}
\AXC{$\varphi_{439}$\lbl{A1}}
\BIC{$\varphi_{358} \to \varphi_{364} \to \varphi_{361} \to (\mathrm{Conv}(b,\beta) \to \beta \neq 0 \to \varphi_{96}) \to \varphi_{384} \to \varphi_{414}$}
\BIC{$\varphi_{364} \to \varphi_{361} \to (\mathrm{Conv}(b,\beta) \to \beta \neq 0 \to \varphi_{96}) \to \varphi_{384} \to \varphi_{414}$}
\end{ptree}
\noindent$\Pi_{63}$:
\begin{ptree}
\AXC{$\stackrel{\Pi_{59}}{\vdots}$}\noLine
\UIC{$\varphi_{361}$}
\AXC{$\stackrel{\Pi_{31}}{\vdots}$}\noLine
\UIC{$\varphi_{319}$}
\AXC{$\stackrel{\Pi_{35}}{\vdots}$}\noLine
\UIC{$\varphi_{337}$}
\AXC{$\stackrel{\Pi_{60}}{\vdots}$}\noLine
\UIC{$\varphi_{132}$}
\AXC{$\varphi_{132} \to \varphi_{337} \to \varphi_{319} \to \varphi_{364}$\lbl{A1}}
\BIC{$\varphi_{337} \to \varphi_{319} \to \varphi_{364}$}
\BIC{$\varphi_{319} \to \varphi_{364}$}
\BIC{$\varphi_{364}$}
\AXC{$\stackrel{\Pi_{62}}{\vdots}$}\noLine
\UIC{$\varphi_{364} \to \varphi_{361} \to (\mathrm{Conv}(b,\beta) \to \beta \neq 0 \to \varphi_{96}) \to \varphi_{384} \to \varphi_{414}$}
\BIC{$\varphi_{361} \to (\mathrm{Conv}(b,\beta) \to \beta \neq 0 \to \varphi_{96}) \to \varphi_{384} \to \varphi_{414}$}
\BIC{$(\mathrm{Conv}(b,\beta) \to \beta \neq 0 \to \varphi_{96}) \to \varphi_{384} \to \varphi_{414}$}
\end{ptree}
\noindent$\Pi_{64}$:
\begin{ptree}
\AXC{$\stackrel{\Pi_{55}}{\vdots}$}\noLine
\UIC{$\mathrm{Conv}(b,\beta) \to \beta \neq 0 \to \varphi_{96}$}
\AXC{$\stackrel{\Pi_{63}}{\vdots}$}\noLine
\UIC{$(\mathrm{Conv}(b,\beta) \to \beta \neq 0 \to \varphi_{96}) \to \varphi_{384} \to \varphi_{414}$}
\BIC{$\varphi_{384} \to \varphi_{414}$}
\end{ptree}
\noindent$\Pi_{65}$:
\begin{ptree}
\AXC{$y = z \to x = y \to x = z$\lbl{A6}}
\AXC{$(y = z \to x = y \to x = z) \to x = y \to y = z \to x = z$\lbl{A1}}
\BIC{$x = y \to y = z \to x = z$}
\UIC{$\forall z\,(x = y \to y = z \to x = z)$}
\UIC{$\forall y\,\forall z\,(x = y \to y = z \to x = z)$}
\UIC{$\forall x\,\forall y\,\forall z\,(x = y \to y = z \to x = z)$}
\AXC{$\forall x\,\forall y\,\forall z\,(x = y \to y = z \to x = z) \to \varphi_{259}$\lbl{A2}}
\BIC{$\varphi_{259}$}
\end{ptree}
\noindent$\Pi_{66}$:
\begin{ptree}
\AXC{$\varphi_{39}$\lbl{A1}}
\AXC{$\forall x \in \mathbb{R}\,(x \cdot 1 = x)$\lbl{(R8)}}
\AXC{$\forall x \in \mathbb{R}\,(x \cdot 1 = x) \to \varepsilon \in \mathbb{R} \to \varepsilon \cdot 1 = \varepsilon$\lbl{A2}}
\BIC{$\varepsilon \in \mathbb{R} \to \varepsilon \cdot 1 = \varepsilon$}
\AXC{$(\varepsilon \in \mathbb{R} \to \varepsilon \cdot 1 = \varepsilon) \to \varphi_{39} \to \varphi_{308}$\lbl{A1}}
\BIC{$\varphi_{39} \to \varphi_{308}$}
\BIC{$\varphi_{308}$}
\end{ptree}
\noindent$\Pi_{67}$:
\begin{ptree}
\AXC{$y = z \to x = y \to x = z$\lbl{A6}}
\AXC{$(y = z \to x = y \to x = z) \to x = y \to y = z \to x = z$\lbl{A1}}
\BIC{$x = y \to y = z \to x = z$}
\UIC{$\forall z\,(x = y \to y = z \to x = z)$}
\UIC{$\forall y\,\forall z\,(x = y \to y = z \to x = z)$}
\UIC{$\forall x\,\forall y\,\forall z\,(x = y \to y = z \to x = z)$}
\AXC{$\forall x\,\forall y\,\forall z\,(x = y \to y = z \to x = z) \to \varphi_{259}$\lbl{A2}}
\BIC{$\varphi_{259}$}
\end{ptree}
\noindent$\Pi_{68}$:
\begin{ptree}
\AXC{$z = w \to \varphi_{47} \to \varphi_{46}$\lbl{A6}}
\UIC{$\forall w\,(z = w \to \varphi_{47} \to \varphi_{46})$}
\UIC{$\forall z\,\forall w\,(z = w \to \varphi_{47} \to \varphi_{46})$}
\AXC{$\forall z\,\forall w\,(z = w \to \varphi_{47} \to \varphi_{46}) \to \varphi_{373}$\lbl{A2}}
\BIC{$\varphi_{373}$}
\AXC{$\varphi_{373} \to \varphi_{369}$\lbl{A2}}
\BIC{$\varphi_{369}$}
\end{ptree}
\noindent$\Pi_{69}$:
\begin{ptree}
\AXC{$\stackrel{\Pi_{24}}{\vdots}$}\noLine
\UIC{$\forall x\,\forall y\,(x = y \to y = x)$}
\AXC{$\forall x\,\forall y\,(x = y \to y = x) \to \varphi_{17}$\lbl{A2}}
\BIC{$\varphi_{17}$}
\AXC{$\varphi_{17} \to \lvert \beta \rvert \cdot 2^{-1} \cdot \lvert \beta \rvert = 0 \to 0 = \lvert \beta \rvert \cdot 2^{-1} \cdot \lvert \beta \rvert$\lbl{A2}}
\BIC{$\lvert \beta \rvert \cdot 2^{-1} \cdot \lvert \beta \rvert = 0 \to 0 = \lvert \beta \rvert \cdot 2^{-1} \cdot \lvert \beta \rvert$}
\end{ptree}
\noindent$\Pi_{70}$:
\begin{ptree}
\AXC{$\stackrel{\Pi_{47}}{\vdots}$}\noLine
\UIC{$\varphi_{55}$}
\AXC{$\forall x\,\forall y\,(x < y \to x \neq y)$\lbl{[F26]}}
\AXC{$\forall x\,\forall y\,(x < y \to x \neq y) \to \forall y\,(0 < y \to 0 \neq y)$\lbl{A2}}
\BIC{$\forall y\,(0 < y \to 0 \neq y)$}
\AXC{$\varphi_{67}$\lbl{A2}}
\BIC{$0 < \lvert \beta \rvert \cdot 2^{-1} \cdot \lvert \beta \rvert \to 0 \neq \lvert \beta \rvert \cdot 2^{-1} \cdot \lvert \beta \rvert$}
\AXC{$\varphi_{416}$\lbl{A1}}
\BIC{$\varphi_{55} \to \varphi_{338}$}
\BIC{$\varphi_{338}$}
\end{ptree}
\noindent$\Pi_{71}$:
\begin{ptree}
\AXC{$\stackrel{\Pi_{57}}{\vdots}$}\noLine
\UIC{$\varphi_{344}$}
\AXC{$\forall x \in \mathbb{R}\,(x \neq 0 \to x \cdot x^{-1} = 1)$\lbl{(R9)}}
\AXC{$\varphi_{266}$\lbl{A2}}
\BIC{$\lvert \beta \rvert \cdot 2^{-1} \cdot \lvert \beta \rvert \in \mathbb{R} \to \varphi_{120}$}
\AXC{$(\lvert \beta \rvert \cdot 2^{-1} \cdot \lvert \beta \rvert \in \mathbb{R} \to \varphi_{120}) \to \varphi_{344} \to \varphi_{339} \to \varphi_{377}$\lbl{A1}}
\BIC{$\varphi_{344} \to \varphi_{339} \to \varphi_{377}$}
\BIC{$\varphi_{339} \to \varphi_{377}$}
\end{ptree}
\noindent$\Pi_{72}$:
\begin{ptree}
\AXC{$\stackrel{\Pi_{69}}{\vdots}$}\noLine
\UIC{$\lvert \beta \rvert \cdot 2^{-1} \cdot \lvert \beta \rvert = 0 \to 0 = \lvert \beta \rvert \cdot 2^{-1} \cdot \lvert \beta \rvert$}
\AXC{$\stackrel{\Pi_{70}}{\vdots}$}\noLine
\UIC{$\varphi_{338}$}
\AXC{$\varphi_{338} \to \varphi_{401}$\lbl{A1}}
\BIC{$\varphi_{401}$}
\BIC{$\varphi_{339}$}
\AXC{$\stackrel{\Pi_{71}}{\vdots}$}\noLine
\UIC{$\varphi_{339} \to \varphi_{377}$}
\BIC{$\varphi_{377}$}
\AXC{$\varphi_{377} \to \varphi_{217} \to \varphi_{369} \to \varphi_{390}$\lbl{A1}}
\BIC{$\varphi_{217} \to \varphi_{369} \to \varphi_{390}$}
\end{ptree}
\noindent$\Pi_{73}$:
\begin{ptree}
\AXC{$y = z \to x = y \to x = z$\lbl{A6}}
\AXC{$(y = z \to x = y \to x = z) \to x = y \to y = z \to x = z$\lbl{A1}}
\BIC{$x = y \to y = z \to x = z$}
\UIC{$\forall z\,(x = y \to y = z \to x = z)$}
\UIC{$\forall y\,\forall z\,(x = y \to y = z \to x = z)$}
\UIC{$\forall x\,\forall y\,\forall z\,(x = y \to y = z \to x = z)$}
\AXC{$\forall x\,\forall y\,\forall z\,(x = y \to y = z \to x = z) \to \varphi_{259}$\lbl{A2}}
\BIC{$\varphi_{259}$}
\end{ptree}
\noindent$\Pi_{74}$:
\begin{ptree}
\AXC{$\stackrel{\Pi_{24}}{\vdots}$}\noLine
\UIC{$\forall x\,\forall y\,(x = y \to y = x)$}
\AXC{$\forall x\,\forall y\,(x = y \to y = x) \to \varphi_{17}$\lbl{A2}}
\BIC{$\varphi_{17}$}
\AXC{$\varphi_{17} \to \lvert \beta \rvert \cdot 2^{-1} \cdot \lvert \beta \rvert = 0 \to 0 = \lvert \beta \rvert \cdot 2^{-1} \cdot \lvert \beta \rvert$\lbl{A2}}
\BIC{$\lvert \beta \rvert \cdot 2^{-1} \cdot \lvert \beta \rvert = 0 \to 0 = \lvert \beta \rvert \cdot 2^{-1} \cdot \lvert \beta \rvert$}
\end{ptree}
\noindent$\Pi_{75}$:
\begin{ptree}
\AXC{$\stackrel{\Pi_{47}}{\vdots}$}\noLine
\UIC{$\varphi_{55}$}
\AXC{$\forall x\,\forall y\,(x < y \to x \neq y)$\lbl{[F26]}}
\AXC{$\forall x\,\forall y\,(x < y \to x \neq y) \to \forall y\,(0 < y \to 0 \neq y)$\lbl{A2}}
\BIC{$\forall y\,(0 < y \to 0 \neq y)$}
\AXC{$\varphi_{67}$\lbl{A2}}
\BIC{$0 < \lvert \beta \rvert \cdot 2^{-1} \cdot \lvert \beta \rvert \to 0 \neq \lvert \beta \rvert \cdot 2^{-1} \cdot \lvert \beta \rvert$}
\AXC{$\varphi_{416}$\lbl{A1}}
\BIC{$\varphi_{55} \to \varphi_{338}$}
\BIC{$\varphi_{338}$}
\end{ptree}
\noindent$\Pi_{76}$:
\begin{ptree}
\AXC{$\stackrel{\Pi_{57}}{\vdots}$}\noLine
\UIC{$\varphi_{344}$}
\AXC{$\forall x \in \mathbb{R}\,(x \neq 0 \to x^{-1} \in \mathbb{R})$\lbl{(R1)}}
\AXC{$\varphi_{239}$\lbl{A2}}
\BIC{$\lvert \beta \rvert \cdot 2^{-1} \cdot \lvert \beta \rvert \in \mathbb{R} \to \varphi_{37}$}
\AXC{$(\lvert \beta \rvert \cdot 2^{-1} \cdot \lvert \beta \rvert \in \mathbb{R} \to \varphi_{37}) \to \varphi_{344} \to \varphi_{339} \to \varphi_{355}$\lbl{A1}}
\BIC{$\varphi_{344} \to \varphi_{339} \to \varphi_{355}$}
\BIC{$\varphi_{339} \to \varphi_{355}$}
\end{ptree}
\noindent$\Pi_{77}$:
\begin{ptree}
\AXC{$\varphi_{39}$\lbl{A1}}
\AXC{$\varphi_{21}$\lbl{(R6)}}
\AXC{$\varphi_{21} \to \varphi_{30}$\lbl{A2}}
\BIC{$\varphi_{30}$}
\AXC{$\varphi_{30} \to \varphi_{230}$\lbl{A2}}
\BIC{$\varphi_{230}$}
\AXC{$\varphi_{230} \to \varepsilon \in \mathbb{R} \to \varphi_{291}$\lbl{A2}}
\BIC{$\varepsilon \in \mathbb{R} \to \varphi_{291}$}
\AXC{$(\varepsilon \in \mathbb{R} \to \varphi_{291}) \to \varphi_{39} \to \varphi_{344} \to \varphi_{355} \to \varphi_{411}$\lbl{A1}}
\BIC{$\varphi_{39} \to \varphi_{344} \to \varphi_{355} \to \varphi_{411}$}
\BIC{$\varphi_{344} \to \varphi_{355} \to \varphi_{411}$}
\end{ptree}
\noindent$\Pi_{78}$:
\begin{ptree}
\AXC{$\stackrel{\Pi_{74}}{\vdots}$}\noLine
\UIC{$\lvert \beta \rvert \cdot 2^{-1} \cdot \lvert \beta \rvert = 0 \to 0 = \lvert \beta \rvert \cdot 2^{-1} \cdot \lvert \beta \rvert$}
\AXC{$\stackrel{\Pi_{75}}{\vdots}$}\noLine
\UIC{$\varphi_{338}$}
\AXC{$\varphi_{338} \to \varphi_{401}$\lbl{A1}}
\BIC{$\varphi_{401}$}
\BIC{$\varphi_{339}$}
\AXC{$\stackrel{\Pi_{76}}{\vdots}$}\noLine
\UIC{$\varphi_{339} \to \varphi_{355}$}
\BIC{$\varphi_{355}$}
\AXC{$\stackrel{\Pi_{57}}{\vdots}$}\noLine
\UIC{$\varphi_{344}$}
\AXC{$\stackrel{\Pi_{77}}{\vdots}$}\noLine
\UIC{$\varphi_{344} \to \varphi_{355} \to \varphi_{411}$}
\BIC{$\varphi_{355} \to \varphi_{411}$}
\BIC{$\varphi_{411}$}
\end{ptree}
\noindent$\Pi_{79}$:
\begin{ptree}
\AXC{$\stackrel{\Pi_{29}}{\vdots}$}\noLine
\UIC{$\varphi_{8}$}
\AXC{$\stackrel{\Pi_{34}}{\vdots}$}\noLine
\UIC{$\varphi_{19}$}
\AXC{$\stackrel{\Pi_{19}}{\vdots}$}\noLine
\UIC{$\varphi_{310}$}
\AXC{$\stackrel{\Pi_{20}}{\vdots}$}\noLine
\UIC{$\varphi_{332}$}
\AXC{$\varphi_{31}$\lbl{[E16]}}
\AXC{$\varphi_{31} \to \varphi_{208}$\lbl{A2}}
\BIC{$\varphi_{208}$}
\AXC{$\varphi_{208} \to \varphi_{237}$\lbl{A2}}
\BIC{$\varphi_{237}$}
\AXC{$\varphi_{237} \to \varphi_{332} \to \varphi_{310} \to \varphi_{19} \to \varphi_{8} \to \varphi_{382}$\lbl{A1}}
\BIC{$\varphi_{332} \to \varphi_{310} \to \varphi_{19} \to \varphi_{8} \to \varphi_{382}$}
\BIC{$\varphi_{310} \to \varphi_{19} \to \varphi_{8} \to \varphi_{382}$}
\BIC{$\varphi_{19} \to \varphi_{8} \to \varphi_{382}$}
\BIC{$\varphi_{8} \to \varphi_{382}$}
\BIC{$\varphi_{382}$}
\end{ptree}
\noindent$\Pi_{80}$:
\begin{ptree}
\AXC{$\stackrel{\Pi_{24}}{\vdots}$}\noLine
\UIC{$\forall x\,\forall y\,(x = y \to y = x)$}
\AXC{$\forall x\,\forall y\,(x = y \to y = x) \to \varphi_{35}$\lbl{A2}}
\BIC{$\varphi_{35}$}
\AXC{$\varphi_{35} \to \varphi_{27} \to \varphi_{26}$\lbl{A2}}
\BIC{$\varphi_{27} \to \varphi_{26}$}
\AXC{$\stackrel{\Pi_{79}}{\vdots}$}\noLine
\UIC{$\varphi_{382}$}
\AXC{$\varphi_{382} \to (\varphi_{27} \to \varphi_{26}) \to \varphi_{383}$\lbl{A1}}
\BIC{$(\varphi_{27} \to \varphi_{26}) \to \varphi_{383}$}
\BIC{$\varphi_{383}$}
\end{ptree}
\noindent$\Pi_{81}$:
\begin{ptree}
\AXC{$x = x$\lbl{A5}}
\UIC{$\forall x\,(x = x)$}
\AXC{$\forall x\,(x = x) \to \varphi_{234}$\lbl{A2}}
\BIC{$\varphi_{234}$}
\AXC{$\stackrel{\Pi_{80}}{\vdots}$}\noLine
\UIC{$\varphi_{383}$}
\AXC{$\varphi_{383} \to \varphi_{234} \to (\varphi_{26} \to \varphi_{234} \to \varphi_{226}) \to \varphi_{412}$\lbl{A1}}
\BIC{$\varphi_{234} \to (\varphi_{26} \to \varphi_{234} \to \varphi_{226}) \to \varphi_{412}$}
\BIC{$(\varphi_{26} \to \varphi_{234} \to \varphi_{226}) \to \varphi_{412}$}
\end{ptree}
\noindent$\Pi_{82}$:
\begin{ptree}
\AXC{$\stackrel{\Pi_{78}}{\vdots}$}\noLine
\UIC{$\varphi_{411}$}
\AXC{$z = w \to \varphi_{177} \to \varphi_{176}$\lbl{A6}}
\UIC{$\forall w\,(z = w \to \varphi_{177} \to \varphi_{176})$}
\UIC{$\forall z\,\forall w\,(z = w \to \varphi_{177} \to \varphi_{176})$}
\AXC{$\varphi_{420}$\lbl{A2}}
\BIC{$\forall w\,((\lvert \beta \rvert \cdot 2^{-1})^{-1} \cdot \lvert \beta \rvert^{-1} = w \to \varphi_{234} \to \varphi_{176})$}
\AXC{$\varphi_{432}$\lbl{A2}}
\BIC{$\varphi_{26} \to \varphi_{234} \to \varphi_{226}$}
\AXC{$\stackrel{\Pi_{81}}{\vdots}$}\noLine
\UIC{$(\varphi_{26} \to \varphi_{234} \to \varphi_{226}) \to \varphi_{412}$}
\BIC{$\varphi_{412}$}
\AXC{$\varphi_{412} \to \varphi_{411} \to (\varphi_{226} \to \varphi_{218} \to \varphi_{225}) \to \varphi_{413}$\lbl{A1}}
\BIC{$\varphi_{411} \to (\varphi_{226} \to \varphi_{218} \to \varphi_{225}) \to \varphi_{413}$}
\BIC{$(\varphi_{226} \to \varphi_{218} \to \varphi_{225}) \to \varphi_{413}$}
\end{ptree}
\noindent$\Pi_{83}$:
\begin{ptree}
\AXC{$\stackrel{\Pi_{73}}{\vdots}$}\noLine
\UIC{$\varphi_{259}$}
\AXC{$\varphi_{259} \to \varphi_{380}$\lbl{A2}}
\BIC{$\varphi_{380}$}
\AXC{$\varphi_{380} \to \varphi_{226} \to \varphi_{218} \to \varphi_{225}$\lbl{A2}}
\BIC{$\varphi_{226} \to \varphi_{218} \to \varphi_{225}$}
\AXC{$\stackrel{\Pi_{82}}{\vdots}$}\noLine
\UIC{$(\varphi_{226} \to \varphi_{218} \to \varphi_{225}) \to \varphi_{413}$}
\BIC{$\varphi_{413}$}
\AXC{$\varphi_{413} \to \varphi_{390} \to (\varphi_{225} \to \varphi_{45} \to \varphi_{69}) \to \varphi_{393}$\lbl{A1}}
\BIC{$\varphi_{390} \to (\varphi_{225} \to \varphi_{45} \to \varphi_{69}) \to \varphi_{393}$}
\end{ptree}
\noindent$\Pi_{84}$:
\begin{ptree}
\AXC{$\stackrel{\Pi_{68}}{\vdots}$}\noLine
\UIC{$\varphi_{369}$}
\AXC{$x = x$\lbl{A5}}
\UIC{$\forall x\,(x = x)$}
\AXC{$\forall x\,(x = x) \to \varphi_{217}$\lbl{A2}}
\BIC{$\varphi_{217}$}
\AXC{$\stackrel{\Pi_{72}}{\vdots}$}\noLine
\UIC{$\varphi_{217} \to \varphi_{369} \to \varphi_{390}$}
\BIC{$\varphi_{369} \to \varphi_{390}$}
\BIC{$\varphi_{390}$}
\AXC{$\stackrel{\Pi_{83}}{\vdots}$}\noLine
\UIC{$\varphi_{390} \to (\varphi_{225} \to \varphi_{45} \to \varphi_{69}) \to \varphi_{393}$}
\BIC{$(\varphi_{225} \to \varphi_{45} \to \varphi_{69}) \to \varphi_{393}$}
\end{ptree}
\noindent$\Pi_{85}$:
\begin{ptree}
\AXC{$\stackrel{\Pi_{66}}{\vdots}$}\noLine
\UIC{$\varphi_{308}$}
\AXC{$\stackrel{\Pi_{67}}{\vdots}$}\noLine
\UIC{$\varphi_{259}$}
\AXC{$\varphi_{259} \to \varphi_{379}$\lbl{A2}}
\BIC{$\varphi_{379}$}
\AXC{$\varphi_{379} \to \varphi_{225} \to \varphi_{45} \to \varphi_{69}$\lbl{A2}}
\BIC{$\varphi_{225} \to \varphi_{45} \to \varphi_{69}$}
\AXC{$\stackrel{\Pi_{84}}{\vdots}$}\noLine
\UIC{$(\varphi_{225} \to \varphi_{45} \to \varphi_{69}) \to \varphi_{393}$}
\BIC{$\varphi_{393}$}
\AXC{$\varphi_{393} \to \varphi_{308} \to (\varphi_{69} \to \varepsilon \cdot 1 = \varepsilon \to \varphi_{43}) \to \varphi_{391}$\lbl{A1}}
\BIC{$\varphi_{308} \to (\varphi_{69} \to \varepsilon \cdot 1 = \varepsilon \to \varphi_{43}) \to \varphi_{391}$}
\BIC{$(\varphi_{69} \to \varepsilon \cdot 1 = \varepsilon \to \varphi_{43}) \to \varphi_{391}$}
\end{ptree}
\noindent$\Pi_{86}$:
\begin{ptree}
\AXC{$\stackrel{\Pi_{65}}{\vdots}$}\noLine
\UIC{$\varphi_{259}$}
\AXC{$\varphi_{259} \to \varphi_{267}$\lbl{A2}}
\BIC{$\varphi_{267}$}
\AXC{$\varphi_{267} \to \varphi_{69} \to \varepsilon \cdot 1 = \varepsilon \to \varphi_{43}$\lbl{A2}}
\BIC{$\varphi_{69} \to \varepsilon \cdot 1 = \varepsilon \to \varphi_{43}$}
\AXC{$\stackrel{\Pi_{85}}{\vdots}$}\noLine
\UIC{$(\varphi_{69} \to \varepsilon \cdot 1 = \varepsilon \to \varphi_{43}) \to \varphi_{391}$}
\BIC{$\varphi_{391}$}
\AXC{$\varphi_{391} \to \varphi_{414} \to \varphi_{387} \to \varphi_{389}$\lbl{A1}}
\BIC{$\varphi_{414} \to \varphi_{387} \to \varphi_{389}$}
\end{ptree}
\noindent$\Pi_{87}$:
\begin{ptree}
\AXC{$\stackrel{\Pi_{22}}{\vdots}$}\noLine
\UIC{$\varphi_{360}$}
\AXC{$\stackrel{\Pi_{23}}{\vdots}$}\noLine
\UIC{$\varphi_{315}$}
\AXC{$\stackrel{\Pi_{31}}{\vdots}$}\noLine
\UIC{$\varphi_{319}$}
\AXC{$\stackrel{\Pi_{44}}{\vdots}$}\noLine
\UIC{$\varphi_{319} \to \varphi_{315} \to \varphi_{360} \to \varphi_{384}$}
\BIC{$\varphi_{315} \to \varphi_{360} \to \varphi_{384}$}
\BIC{$\varphi_{360} \to \varphi_{384}$}
\BIC{$\varphi_{384}$}
\AXC{$\stackrel{\Pi_{64}}{\vdots}$}\noLine
\UIC{$\varphi_{384} \to \varphi_{414}$}
\BIC{$\varphi_{414}$}
\AXC{$\stackrel{\Pi_{86}}{\vdots}$}\noLine
\UIC{$\varphi_{414} \to \varphi_{387} \to \varphi_{389}$}
\BIC{$\varphi_{387} \to \varphi_{389}$}
\end{ptree}
\noindent$\Pi_{88}$:
\begin{ptree}
\AXC{$\forall x\,\forall y\,\forall z\,(x < y \to y \leq z \to x < z)$\lbl{[O3]}}
\AXC{$\varphi_{99}$\lbl{A2}}
\BIC{$\forall y\,\forall z\,(\lvert b_{n} - \beta \rvert < y \to y \leq z \to \lvert b_{n} - \beta \rvert < z)$}
\AXC{$\varphi_{197}$\lbl{A2}}
\BIC{$\forall z\,(\lvert b_{n} - \beta \rvert < d_{0} \to d_{0} \leq z \to \lvert b_{n} - \beta \rvert < z)$}
\AXC{$\varphi_{273}$\lbl{A2}}
\BIC{$\lvert b_{n} - \beta \rvert < d_{0} \to \varphi_{112}$}
\AXC{$(\lvert b_{n} - \beta \rvert < d_{0} \to \varphi_{112}) \to \varphi_{330} \to \varphi_{346} \to \varphi_{372}$\lbl{A1}}
\BIC{$\varphi_{330} \to \varphi_{346} \to \varphi_{372}$}
\end{ptree}
\noindent$\Pi_{89}$:
\begin{ptree}
\AXC{$\stackrel{\Pi_{11}}{\vdots}$}\noLine
\UIC{$\varphi_{10}$}
\AXC{$\stackrel{\Pi_{19}}{\vdots}$}\noLine
\UIC{$\varphi_{310}$}
\AXC{$\forall x \in \mathbb{R}\,(0 < x \to 0 < x^{-1})$\lbl{[F34]}}
\AXC{$\varphi_{80}$\lbl{A2}}
\BIC{$\lvert \beta \rvert \in \mathbb{R} \to 0 < \lvert \beta \rvert \to 0 < \lvert \beta \rvert^{-1}$}
\AXC{$(\lvert \beta \rvert \in \mathbb{R} \to 0 < \lvert \beta \rvert \to 0 < \lvert \beta \rvert^{-1}) \to \varphi_{310} \to \varphi_{10} \to \varphi_{315}$\lbl{A1}}
\BIC{$\varphi_{310} \to \varphi_{10} \to \varphi_{315}$}
\BIC{$\varphi_{10} \to \varphi_{315}$}
\BIC{$\varphi_{315}$}
\end{ptree}
\noindent$\Pi_{90}$:
\begin{ptree}
\AXC{$\stackrel{\Pi_{17}}{\vdots}$}\noLine
\UIC{$\varphi_{320}$}
\AXC{$\forall x \in \mathbb{R}\,(0 < x \to 0 < x^{-1})$\lbl{[F34]}}
\AXC{$\varphi_{146}$\lbl{A2}}
\BIC{$\lvert b_{n} \rvert \in \mathbb{R} \to 0 < \lvert b_{n} \rvert \to 0 < \lvert b_{n} \rvert^{-1}$}
\AXC{$(\lvert b_{n} \rvert \in \mathbb{R} \to 0 < \lvert b_{n} \rvert \to 0 < \lvert b_{n} \rvert^{-1}) \to \varphi_{320} \to \varphi_{317} \to \varphi_{327}$\lbl{A1}}
\BIC{$\varphi_{320} \to \varphi_{317} \to \varphi_{327}$}
\BIC{$\varphi_{317} \to \varphi_{327}$}
\end{ptree}
\noindent$\Pi_{91}$:
\begin{ptree}
\AXC{$\stackrel{\Pi_{43}}{\vdots}$}\noLine
\UIC{$\varphi_{333}$}
\AXC{$\forall x\,\forall y\,(x \in \mathbb{R} \to y \in \mathbb{R} \to 0 < x \to 0 < y \to 0 < x \cdot y)$\lbl{(R15)}}
\AXC{$\varphi_{260}$\lbl{A2}}
\BIC{$\varphi_{152}$}
\AXC{$\varphi_{152} \to \varphi_{211}$\lbl{A2}}
\BIC{$\varphi_{211}$}
\AXC{$\varphi_{211} \to \varphi_{333} \to \varphi_{319} \to \varphi_{327} \to \varphi_{315} \to \varphi_{348}$\lbl{A1}}
\BIC{$\varphi_{333} \to \varphi_{319} \to \varphi_{327} \to \varphi_{315} \to \varphi_{348}$}
\BIC{$\varphi_{319} \to \varphi_{327} \to \varphi_{315} \to \varphi_{348}$}
\end{ptree}
\noindent$\Pi_{92}$:
\begin{ptree}
\AXC{$\forall x\,\forall y\,(x \in \mathbb{R} \to y \in \mathbb{R} \to x - y \in \mathbb{R})$\lbl{[F1]}}
\AXC{$\varphi_{147}$\lbl{A2}}
\BIC{$\forall y\,(b_{n} \in \mathbb{R} \to y \in \mathbb{R} \to b_{n} - y \in \mathbb{R})$}
\AXC{$\varphi_{174}$\lbl{A2}}
\BIC{$b_{n} \in \mathbb{R} \to \beta \in \mathbb{R} \to b_{n} - \beta \in \mathbb{R}$}
\AXC{$(b_{n} \in \mathbb{R} \to \beta \in \mathbb{R} \to b_{n} - \beta \in \mathbb{R}) \to \varphi_{314} \to \varphi_{394}$\lbl{A1}}
\BIC{$\varphi_{314} \to \varphi_{394}$}
\end{ptree}
\noindent$\Pi_{93}$:
\begin{ptree}
\AXC{$\forall x \in \mathbb{R}\,(\lvert x \rvert \in \mathbb{R})$\lbl{[A5]}}
\AXC{$\varphi_{74}$\lbl{A2}}
\BIC{$b_{n} - \beta \in \mathbb{R} \to \lvert b_{n} - \beta \rvert \in \mathbb{R}$}
\AXC{$(b_{n} - \beta \in \mathbb{R} \to \lvert b_{n} - \beta \rvert \in \mathbb{R}) \to \varphi_{328} \to \varphi_{334}$\lbl{A1}}
\BIC{$\varphi_{328} \to \varphi_{334}$}
\end{ptree}
\noindent$\Pi_{94}$:
\begin{ptree}
\AXC{$\stackrel{\Pi_{4}}{\vdots}$}\noLine
\UIC{$\mathrm{Conv}(b,\beta) \to \beta \neq 0 \to \forall n \in \mathbb{N}\,(b_{n} \neq 0) \to \beta \in \mathbb{R}$}
\AXC{$\stackrel{\Pi_{7}}{\vdots}$}\noLine
\UIC{$\varphi_{314}$}
\AXC{$\stackrel{\Pi_{92}}{\vdots}$}\noLine
\UIC{$\varphi_{314} \to \varphi_{394}$}
\BIC{$\varphi_{394}$}
\BIC{$\varphi_{328}$}
\AXC{$\stackrel{\Pi_{93}}{\vdots}$}\noLine
\UIC{$\varphi_{328} \to \varphi_{334}$}
\BIC{$\varphi_{334}$}
\end{ptree}
\noindent$\Pi_{95}$:
\begin{ptree}
\AXC{$\varphi_{49}$\lbl{[F32]}}
\AXC{$\varphi_{49} \to \varphi_{206}$\lbl{A2}}
\BIC{$\varphi_{206}$}
\AXC{$\varphi_{206} \to \varphi_{293}$\lbl{A2}}
\BIC{$\varphi_{293}$}
\AXC{$\varphi_{293} \to \lvert b_{n} - \beta \rvert \in \mathbb{R} \to \varphi_{400}$\lbl{A2}}
\BIC{$\lvert b_{n} - \beta \rvert \in \mathbb{R} \to \varphi_{400}$}
\end{ptree}
\noindent$\Pi_{96}$:
\begin{ptree}
\AXC{$\stackrel{\Pi_{61}}{\vdots}$}\noLine
\UIC{$\varphi_{358}$}
\AXC{$\stackrel{\Pi_{59}}{\vdots}$}\noLine
\UIC{$\varphi_{361}$}
\AXC{$\stackrel{\Pi_{94}}{\vdots}$}\noLine
\UIC{$\varphi_{334}$}
\AXC{$\stackrel{\Pi_{95}}{\vdots}$}\noLine
\UIC{$\lvert b_{n} - \beta \rvert \in \mathbb{R} \to \varphi_{400}$}
\AXC{$(\lvert b_{n} - \beta \rvert \in \mathbb{R} \to \varphi_{400}) \to \varphi_{334} \to \varphi_{361} \to \varphi_{358} \to \varphi_{348} \to \varphi_{372} \to \varphi_{410}$\lbl{A1}}
\BIC{$\varphi_{334} \to \varphi_{361} \to \varphi_{358} \to \varphi_{348} \to \varphi_{372} \to \varphi_{410}$}
\BIC{$\varphi_{361} \to \varphi_{358} \to \varphi_{348} \to \varphi_{372} \to \varphi_{410}$}
\BIC{$\varphi_{358} \to \varphi_{348} \to \varphi_{372} \to \varphi_{410}$}
\BIC{$\varphi_{348} \to \varphi_{372} \to \varphi_{410}$}
\end{ptree}
\noindent$\Pi_{97}$:
\begin{ptree}
\AXC{$\stackrel{\Pi_{89}}{\vdots}$}\noLine
\UIC{$\varphi_{315}$}
\AXC{$\stackrel{\Pi_{38}}{\vdots}$}\noLine
\UIC{$\varphi_{317}$}
\AXC{$\stackrel{\Pi_{90}}{\vdots}$}\noLine
\UIC{$\varphi_{317} \to \varphi_{327}$}
\BIC{$\varphi_{327}$}
\AXC{$\stackrel{\Pi_{31}}{\vdots}$}\noLine
\UIC{$\varphi_{319}$}
\AXC{$\stackrel{\Pi_{91}}{\vdots}$}\noLine
\UIC{$\varphi_{319} \to \varphi_{327} \to \varphi_{315} \to \varphi_{348}$}
\BIC{$\varphi_{327} \to \varphi_{315} \to \varphi_{348}$}
\BIC{$\varphi_{315} \to \varphi_{348}$}
\BIC{$\varphi_{348}$}
\AXC{$\stackrel{\Pi_{96}}{\vdots}$}\noLine
\UIC{$\varphi_{348} \to \varphi_{372} \to \varphi_{410}$}
\BIC{$\varphi_{372} \to \varphi_{410}$}
\end{ptree}
\noindent$\Pi_{98}$:
\begin{ptree}
\AXC{$\forall x\,\forall y\,\forall z\,(x < y \to y < z \to x < z)$\lbl{[O2]}}
\AXC{$\forall x\,\forall y\,\forall z\,(x < y \to y < z \to x < z) \to \varphi_{214}$\lbl{A2}}
\BIC{$\varphi_{214}$}
\AXC{$\varphi_{214} \to \varphi_{352}$\lbl{A2}}
\BIC{$\varphi_{352}$}
\AXC{$\varphi_{352} \to \varphi_{207} \to \varphi_{220}$\lbl{A2}}
\BIC{$\varphi_{207} \to \varphi_{220}$}
\end{ptree}
\noindent$\Pi_{99}$:
\begin{ptree}
\AXC{$\varphi_{305}$\lbl{A1}}
\AXC{$\varphi_{305} \to \varphi_{346}$\lbl{A1}}
\BIC{$\varphi_{346}$}
\AXC{$\stackrel{\Pi_{1}}{\vdots}$}\noLine
\UIC{$\varphi_{330}$}
\AXC{$\stackrel{\Pi_{88}}{\vdots}$}\noLine
\UIC{$\varphi_{330} \to \varphi_{346} \to \varphi_{372}$}
\BIC{$\varphi_{346} \to \varphi_{372}$}
\BIC{$\varphi_{372}$}
\AXC{$\stackrel{\Pi_{97}}{\vdots}$}\noLine
\UIC{$\varphi_{372} \to \varphi_{410}$}
\BIC{$\varphi_{410}$}
\AXC{$\stackrel{\Pi_{98}}{\vdots}$}\noLine
\UIC{$\varphi_{207} \to \varphi_{220}$}
\AXC{$(\varphi_{207} \to \varphi_{220}) \to \varphi_{410} \to \varphi_{389} \to \varphi_{375}$\lbl{A1}}
\BIC{$\varphi_{410} \to \varphi_{389} \to \varphi_{375}$}
\BIC{$\varphi_{389} \to \varphi_{375}$}
\end{ptree}
\noindent$\Pi_{100}$:
\begin{ptree}
\AXC{$\stackrel{\Pi_{24}}{\vdots}$}\noLine
\UIC{$\forall x\,\forall y\,(x = y \to y = x)$}
\AXC{$\varphi_{68}$\lbl{A2}}
\BIC{$\forall y\,(\lvert b_{n}^{-1} - \beta^{-1} \rvert = y \to y = \lvert b_{n}^{-1} - \beta^{-1} \rvert)$}
\AXC{$\forall y\,(\lvert b_{n}^{-1} - \beta^{-1} \rvert = y \to y = \lvert b_{n}^{-1} - \beta^{-1} \rvert) \to \varphi_{262}$\lbl{A2}}
\BIC{$\varphi_{262}$}
\end{ptree}
\noindent$\Pi_{101}$:
\begin{ptree}
\AXC{$y = z \to x = y \to x = z$\lbl{A6}}
\AXC{$(y = z \to x = y \to x = z) \to x = y \to y = z \to x = z$\lbl{A1}}
\BIC{$x = y \to y = z \to x = z$}
\UIC{$\forall z\,(x = y \to y = z \to x = z)$}
\UIC{$\forall y\,\forall z\,(x = y \to y = z \to x = z)$}
\UIC{$\forall x\,\forall y\,\forall z\,(x = y \to y = z \to x = z)$}
\AXC{$\forall x\,\forall y\,\forall z\,(x = y \to y = z \to x = z) \to \varphi_{40}$\lbl{A2}}
\BIC{$\varphi_{40}$}
\end{ptree}
\noindent$\Pi_{102}$:
\begin{ptree}
\AXC{$z = w \to \varphi_{279}$\lbl{A6}}
\UIC{$\forall w\,(z = w \to \varphi_{279})$}
\UIC{$\forall z\,\forall w\,(z = w \to \varphi_{279})$}
\AXC{$\forall z\,\forall w\,(z = w \to \varphi_{279}) \to \varphi_{301}$\lbl{A2}}
\BIC{$\varphi_{301}$}
\AXC{$\varphi_{301} \to \lvert \beta - b_{n} \rvert = \lvert b_{n} - \beta \rvert \to \varphi_{163} \to \varphi_{166}$\lbl{A2}}
\BIC{$\lvert \beta - b_{n} \rvert = \lvert b_{n} - \beta \rvert \to \varphi_{163} \to \varphi_{166}$}
\end{ptree}
\noindent$\Pi_{103}$:
\begin{ptree}
\AXC{$\forall x\,\forall y\,(x \in \mathbb{R} \to y \in \mathbb{R} \to \lvert x - y \rvert = \lvert y - x \rvert)$\lbl{[A13]}}
\AXC{$\varphi_{160}$\lbl{A2}}
\BIC{$\forall y\,(\beta \in \mathbb{R} \to y \in \mathbb{R} \to \lvert \beta - y \rvert = \lvert y - \beta \rvert)$}
\AXC{$\varphi_{204}$\lbl{A2}}
\BIC{$\beta \in \mathbb{R} \to b_{n} \in \mathbb{R} \to \lvert \beta - b_{n} \rvert = \lvert b_{n} - \beta \rvert$}
\end{ptree}
\noindent$\Pi_{104}$:
\begin{ptree}
\AXC{$\stackrel{\Pi_{103}}{\vdots}$}\noLine
\UIC{$\beta \in \mathbb{R} \to b_{n} \in \mathbb{R} \to \lvert \beta - b_{n} \rvert = \lvert b_{n} - \beta \rvert$}
\AXC{$(\beta \in \mathbb{R} \to b_{n} \in \mathbb{R} \to \lvert \beta - b_{n} \rvert = \lvert b_{n} - \beta \rvert) \to \varphi_{425}$\lbl{A1}}
\BIC{$\varphi_{425}$}
\end{ptree}
\noindent$\Pi_{105}$:
\begin{ptree}
\AXC{$\stackrel{\Pi_{7}}{\vdots}$}\noLine
\UIC{$\varphi_{314}$}
\AXC{$\stackrel{\Pi_{4}}{\vdots}$}\noLine
\UIC{$\mathrm{Conv}(b,\beta) \to \beta \neq 0 \to \forall n \in \mathbb{N}\,(b_{n} \neq 0) \to \beta \in \mathbb{R}$}
\AXC{$\stackrel{\Pi_{104}}{\vdots}$}\noLine
\UIC{$\varphi_{425}$}
\BIC{$\varphi_{314} \to \varphi_{349}$}
\BIC{$\varphi_{349}$}
\end{ptree}
\noindent$\Pi_{106}$:
\begin{ptree}
\AXC{$x = x$\lbl{A5}}
\UIC{$\forall x\,(x = x)$}
\AXC{$\forall x\,(x = x) \to \varphi_{163}$\lbl{A2}}
\BIC{$\varphi_{163}$}
\AXC{$\stackrel{\Pi_{105}}{\vdots}$}\noLine
\UIC{$\varphi_{349}$}
\AXC{$\varphi_{349} \to \varphi_{163} \to (\lvert \beta - b_{n} \rvert = \lvert b_{n} - \beta \rvert \to \varphi_{163} \to \varphi_{166}) \to \varphi_{405}$\lbl{A1}}
\BIC{$\varphi_{163} \to (\lvert \beta - b_{n} \rvert = \lvert b_{n} - \beta \rvert \to \varphi_{163} \to \varphi_{166}) \to \varphi_{405}$}
\BIC{$(\lvert \beta - b_{n} \rvert = \lvert b_{n} - \beta \rvert \to \varphi_{163} \to \varphi_{166}) \to \varphi_{405}$}
\end{ptree}
\noindent$\Pi_{107}$:
\begin{ptree}
\AXC{$y = z \to x = y \to x = z$\lbl{A6}}
\AXC{$(y = z \to x = y \to x = z) \to x = y \to y = z \to x = z$\lbl{A1}}
\BIC{$x = y \to y = z \to x = z$}
\UIC{$\forall z\,(x = y \to y = z \to x = z)$}
\UIC{$\forall y\,\forall z\,(x = y \to y = z \to x = z)$}
\UIC{$\forall x\,\forall y\,\forall z\,(x = y \to y = z \to x = z)$}
\AXC{$\forall x\,\forall y\,\forall z\,(x = y \to y = z \to x = z) \to \varphi_{40}$\lbl{A2}}
\BIC{$\varphi_{40}$}
\end{ptree}
\noindent$\Pi_{108}$:
\begin{ptree}
\AXC{$z = w \to \varphi_{130} \to \varphi_{129}$\lbl{A6}}
\UIC{$\forall w\,(z = w \to \varphi_{130} \to \varphi_{129})$}
\UIC{$\forall z\,\forall w\,(z = w \to \varphi_{130} \to \varphi_{129})$}
\AXC{$\forall z\,\forall w\,(z = w \to \varphi_{130} \to \varphi_{129}) \to \forall w\,(\lvert \beta^{-1} \rvert = w \to \varphi_{164} \to \varphi_{129})$\lbl{A2}}
\BIC{$\forall w\,(\lvert \beta^{-1} \rvert = w \to \varphi_{164} \to \varphi_{129})$}
\end{ptree}
\noindent$\Pi_{109}$:
\begin{ptree}
\AXC{$\forall x \in \mathbb{R}\,(x \neq 0 \to \lvert x^{-1} \rvert = \lvert x \rvert^{-1})$\lbl{[E18]}}
\AXC{$\varphi_{100}$\lbl{A2}}
\BIC{$\beta \in \mathbb{R} \to \beta \neq 0 \to \lvert \beta^{-1} \rvert = \lvert \beta \rvert^{-1}$}
\AXC{$(\beta \in \mathbb{R} \to \beta \neq 0 \to \lvert \beta^{-1} \rvert = \lvert \beta \rvert^{-1}) \to \varphi_{408}$\lbl{A1}}
\BIC{$\varphi_{408}$}
\end{ptree}
\noindent$\Pi_{110}$:
\begin{ptree}
\AXC{$\mathrm{Conv}(b,\beta) \to \beta \neq 0 \to \beta \neq 0$\lbl{A1}}
\AXC{$\stackrel{\Pi_{4}}{\vdots}$}\noLine
\UIC{$\mathrm{Conv}(b,\beta) \to \beta \neq 0 \to \forall n \in \mathbb{N}\,(b_{n} \neq 0) \to \beta \in \mathbb{R}$}
\AXC{$\stackrel{\Pi_{109}}{\vdots}$}\noLine
\UIC{$\varphi_{408}$}
\BIC{$(\mathrm{Conv}(b,\beta) \to \beta \neq 0 \to \beta \neq 0) \to \varphi_{331}$}
\BIC{$\varphi_{331}$}
\end{ptree}
\noindent$\Pi_{111}$:
\begin{ptree}
\AXC{$x = x$\lbl{A5}}
\UIC{$\forall x\,(x = x)$}
\AXC{$\forall x\,(x = x) \to \varphi_{164}$\lbl{A2}}
\BIC{$\varphi_{164}$}
\AXC{$\stackrel{\Pi_{110}}{\vdots}$}\noLine
\UIC{$\varphi_{331}$}
\AXC{$\varphi_{331} \to \varphi_{164} \to (\lvert \beta^{-1} \rvert = \lvert \beta \rvert^{-1} \to \varphi_{164} \to \varphi_{167}) \to \varphi_{406}$\lbl{A1}}
\BIC{$\varphi_{164} \to (\lvert \beta^{-1} \rvert = \lvert \beta \rvert^{-1} \to \varphi_{164} \to \varphi_{167}) \to \varphi_{406}$}
\BIC{$(\lvert \beta^{-1} \rvert = \lvert \beta \rvert^{-1} \to \varphi_{164} \to \varphi_{167}) \to \varphi_{406}$}
\end{ptree}
\noindent$\Pi_{112}$:
\begin{ptree}
\AXC{$\stackrel{\Pi_{108}}{\vdots}$}\noLine
\UIC{$\forall w\,(\lvert \beta^{-1} \rvert = w \to \varphi_{164} \to \varphi_{129})$}
\AXC{$\forall w\,(\lvert \beta^{-1} \rvert = w \to \varphi_{164} \to \varphi_{129}) \to \lvert \beta^{-1} \rvert = \lvert \beta \rvert^{-1} \to \varphi_{164} \to \varphi_{167}$\lbl{A2}}
\BIC{$\lvert \beta^{-1} \rvert = \lvert \beta \rvert^{-1} \to \varphi_{164} \to \varphi_{167}$}
\AXC{$\stackrel{\Pi_{111}}{\vdots}$}\noLine
\UIC{$(\lvert \beta^{-1} \rvert = \lvert \beta \rvert^{-1} \to \varphi_{164} \to \varphi_{167}) \to \varphi_{406}$}
\BIC{$\varphi_{406}$}
\end{ptree}
\noindent$\Pi_{113}$:
\begin{ptree}
\AXC{$y = z \to x = y \to x = z$\lbl{A6}}
\AXC{$(y = z \to x = y \to x = z) \to x = y \to y = z \to x = z$\lbl{A1}}
\BIC{$x = y \to y = z \to x = z$}
\UIC{$\forall z\,(x = y \to y = z \to x = z)$}
\UIC{$\forall y\,\forall z\,(x = y \to y = z \to x = z)$}
\UIC{$\forall x\,\forall y\,\forall z\,(x = y \to y = z \to x = z)$}
\AXC{$\forall x\,\forall y\,\forall z\,(x = y \to y = z \to x = z) \to \varphi_{40}$\lbl{A2}}
\BIC{$\varphi_{40}$}
\end{ptree}
\noindent$\Pi_{114}$:
\begin{ptree}
\AXC{$z = w \to \varphi_{110} \to \varphi_{109}$\lbl{A6}}
\UIC{$\forall w\,(z = w \to \varphi_{110} \to \varphi_{109})$}
\UIC{$\forall z\,\forall w\,(z = w \to \varphi_{110} \to \varphi_{109})$}
\AXC{$\forall z\,\forall w\,(z = w \to \varphi_{110} \to \varphi_{109}) \to \forall w\,(\lvert b_{n}^{-1} \rvert = w \to \varphi_{165} \to \varphi_{109})$\lbl{A2}}
\BIC{$\forall w\,(\lvert b_{n}^{-1} \rvert = w \to \varphi_{165} \to \varphi_{109})$}
\AXC{$\varphi_{418}$\lbl{A2}}
\BIC{$\lvert b_{n}^{-1} \rvert = \lvert b_{n} \rvert^{-1} \to \varphi_{165} \to \varphi_{168}$}
\end{ptree}
\noindent$\Pi_{115}$:
\begin{ptree}
\AXC{$\stackrel{\Pi_{7}}{\vdots}$}\noLine
\UIC{$\varphi_{314}$}
\AXC{$\forall x \in \mathbb{R}\,(x \neq 0 \to \lvert x^{-1} \rvert = \lvert x \rvert^{-1})$\lbl{[E18]}}
\AXC{$\varphi_{182}$\lbl{A2}}
\BIC{$b_{n} \in \mathbb{R} \to b_{n} \neq 0 \to \lvert b_{n}^{-1} \rvert = \lvert b_{n} \rvert^{-1}$}
\AXC{$(b_{n} \in \mathbb{R} \to b_{n} \neq 0 \to \lvert b_{n}^{-1} \rvert = \lvert b_{n} \rvert^{-1}) \to \varphi_{314} \to \varphi_{306} \to \varphi_{345}$\lbl{A1}}
\BIC{$\varphi_{314} \to \varphi_{306} \to \varphi_{345}$}
\BIC{$\varphi_{306} \to \varphi_{345}$}
\end{ptree}
\noindent$\Pi_{116}$:
\begin{ptree}
\AXC{$x = x$\lbl{A5}}
\UIC{$\forall x\,(x = x)$}
\AXC{$\forall x\,(x = x) \to \varphi_{165}$\lbl{A2}}
\BIC{$\varphi_{165}$}
\AXC{$\stackrel{\Pi_{37}}{\vdots}$}\noLine
\UIC{$\varphi_{306}$}
\AXC{$\stackrel{\Pi_{115}}{\vdots}$}\noLine
\UIC{$\varphi_{306} \to \varphi_{345}$}
\BIC{$\varphi_{345}$}
\AXC{$\varphi_{345} \to \varphi_{165} \to (\lvert b_{n}^{-1} \rvert = \lvert b_{n} \rvert^{-1} \to \varphi_{165} \to \varphi_{168}) \to \varphi_{407}$\lbl{A1}}
\BIC{$\varphi_{165} \to (\lvert b_{n}^{-1} \rvert = \lvert b_{n} \rvert^{-1} \to \varphi_{165} \to \varphi_{168}) \to \varphi_{407}$}
\BIC{$(\lvert b_{n}^{-1} \rvert = \lvert b_{n} \rvert^{-1} \to \varphi_{165} \to \varphi_{168}) \to \varphi_{407}$}
\end{ptree}
\noindent$\Pi_{117}$:
\begin{ptree}
\AXC{$y = z \to x = y \to x = z$\lbl{A6}}
\AXC{$(y = z \to x = y \to x = z) \to x = y \to y = z \to x = z$\lbl{A1}}
\BIC{$x = y \to y = z \to x = z$}
\UIC{$\forall z\,(x = y \to y = z \to x = z)$}
\UIC{$\forall y\,\forall z\,(x = y \to y = z \to x = z)$}
\UIC{$\forall x\,\forall y\,\forall z\,(x = y \to y = z \to x = z)$}
\AXC{$\forall x\,\forall y\,\forall z\,(x = y \to y = z \to x = z) \to \varphi_{40}$\lbl{A2}}
\BIC{$\varphi_{40}$}
\end{ptree}
\noindent$\Pi_{118}$:
\begin{ptree}
\AXC{$z = w \to \varphi_{240}$\lbl{A6}}
\UIC{$\forall w\,(z = w \to \varphi_{240})$}
\UIC{$\forall z\,\forall w\,(z = w \to \varphi_{240})$}
\AXC{$\forall z\,\forall w\,(z = w \to \varphi_{240}) \to \varphi_{284}$\lbl{A2}}
\BIC{$\varphi_{284}$}
\AXC{$\varphi_{284} \to \lvert b_{n}^{-1} \cdot \beta^{-1} \rvert = \lvert b_{n}^{-1} \rvert \cdot \lvert \beta^{-1} \rvert \to \varphi_{133} \to \varphi_{157}$\lbl{A2}}
\BIC{$\lvert b_{n}^{-1} \cdot \beta^{-1} \rvert = \lvert b_{n}^{-1} \rvert \cdot \lvert \beta^{-1} \rvert \to \varphi_{133} \to \varphi_{157}$}
\end{ptree}
\noindent$\Pi_{119}$:
\begin{ptree}
\AXC{$\forall x \in \mathbb{R}\,(x \neq 0 \to x^{-1} \in \mathbb{R})$\lbl{(R1)}}
\AXC{$\varphi_{41}$\lbl{A2}}
\BIC{$\beta \in \mathbb{R} \to \beta \neq 0 \to \beta^{-1} \in \mathbb{R}$}
\AXC{$(\beta \in \mathbb{R} \to \beta \neq 0 \to \beta^{-1} \in \mathbb{R}) \to \varphi_{403}$\lbl{A1}}
\BIC{$\varphi_{403}$}
\end{ptree}
\noindent$\Pi_{120}$:
\begin{ptree}
\AXC{$\mathrm{Conv}(b,\beta) \to \beta \neq 0 \to \beta \neq 0$\lbl{A1}}
\AXC{$\stackrel{\Pi_{4}}{\vdots}$}\noLine
\UIC{$\mathrm{Conv}(b,\beta) \to \beta \neq 0 \to \forall n \in \mathbb{N}\,(b_{n} \neq 0) \to \beta \in \mathbb{R}$}
\AXC{$\stackrel{\Pi_{119}}{\vdots}$}\noLine
\UIC{$\varphi_{403}$}
\BIC{$(\mathrm{Conv}(b,\beta) \to \beta \neq 0 \to \beta \neq 0) \to \varphi_{313}$}
\BIC{$\varphi_{313}$}
\end{ptree}
\noindent$\Pi_{121}$:
\begin{ptree}
\AXC{$\stackrel{\Pi_{37}}{\vdots}$}\noLine
\UIC{$\varphi_{306}$}
\AXC{$\stackrel{\Pi_{7}}{\vdots}$}\noLine
\UIC{$\varphi_{314}$}
\AXC{$\forall x \in \mathbb{R}\,(x \neq 0 \to x^{-1} \in \mathbb{R})$\lbl{(R1)}}
\AXC{$\varphi_{106}$\lbl{A2}}
\BIC{$b_{n} \in \mathbb{R} \to b_{n} \neq 0 \to b_{n}^{-1} \in \mathbb{R}$}
\AXC{$(b_{n} \in \mathbb{R} \to b_{n} \neq 0 \to b_{n}^{-1} \in \mathbb{R}) \to \varphi_{314} \to \varphi_{306} \to \varphi_{321}$\lbl{A1}}
\BIC{$\varphi_{314} \to \varphi_{306} \to \varphi_{321}$}
\BIC{$\varphi_{306} \to \varphi_{321}$}
\BIC{$\varphi_{321}$}
\end{ptree}
\noindent$\Pi_{122}$:
\begin{ptree}
\AXC{$\stackrel{\Pi_{121}}{\vdots}$}\noLine
\UIC{$\varphi_{321}$}
\AXC{$\forall x\,\forall y\,(x \in \mathbb{R} \to y \in \mathbb{R} \to \lvert x \cdot y \rvert = \lvert x \rvert \cdot \lvert y \rvert)$\lbl{[A17]}}
\AXC{$\varphi_{238}$\lbl{A2}}
\BIC{$\varphi_{88}$}
\AXC{$\varphi_{88} \to \varphi_{148}$\lbl{A2}}
\BIC{$\varphi_{148}$}
\AXC{$\varphi_{148} \to \varphi_{321} \to \varphi_{313} \to \varphi_{376}$\lbl{A1}}
\BIC{$\varphi_{321} \to \varphi_{313} \to \varphi_{376}$}
\BIC{$\varphi_{313} \to \varphi_{376}$}
\end{ptree}
\noindent$\Pi_{123}$:
\begin{ptree}
\AXC{$x = x$\lbl{A5}}
\UIC{$\forall x\,(x = x)$}
\AXC{$\forall x\,(x = x) \to \varphi_{133}$\lbl{A2}}
\BIC{$\varphi_{133}$}
\AXC{$\stackrel{\Pi_{120}}{\vdots}$}\noLine
\UIC{$\varphi_{313}$}
\AXC{$\stackrel{\Pi_{122}}{\vdots}$}\noLine
\UIC{$\varphi_{313} \to \varphi_{376}$}
\BIC{$\varphi_{376}$}
\AXC{$\varphi_{435}$\lbl{A1}}
\BIC{$\varphi_{133} \to (\lvert b_{n}^{-1} \cdot \beta^{-1} \rvert = \lvert b_{n}^{-1} \rvert \cdot \lvert \beta^{-1} \rvert \to \varphi_{133} \to \varphi_{157}) \to \varphi_{404}$}
\BIC{$(\lvert b_{n}^{-1} \cdot \beta^{-1} \rvert = \lvert b_{n}^{-1} \rvert \cdot \lvert \beta^{-1} \rvert \to \varphi_{133} \to \varphi_{157}) \to \varphi_{404}$}
\end{ptree}
\noindent$\Pi_{124}$:
\begin{ptree}
\AXC{$y = z \to x = y \to x = z$\lbl{A6}}
\AXC{$(y = z \to x = y \to x = z) \to x = y \to y = z \to x = z$\lbl{A1}}
\BIC{$x = y \to y = z \to x = z$}
\UIC{$\forall z\,(x = y \to y = z \to x = z)$}
\UIC{$\forall y\,\forall z\,(x = y \to y = z \to x = z)$}
\UIC{$\forall x\,\forall y\,\forall z\,(x = y \to y = z \to x = z)$}
\AXC{$\forall x\,\forall y\,\forall z\,(x = y \to y = z \to x = z) \to \varphi_{40}$\lbl{A2}}
\BIC{$\varphi_{40}$}
\end{ptree}
\noindent$\Pi_{125}$:
\begin{ptree}
\AXC{$\forall x\,\forall y\,(x \in \mathbb{R} \to y \in \mathbb{R} \to x \cdot y \in \mathbb{R})$\lbl{(R1)}}
\AXC{$\varphi_{175}$\lbl{A2}}
\BIC{$\forall y\,(b_{n}^{-1} \in \mathbb{R} \to y \in \mathbb{R} \to b_{n}^{-1} \cdot y \in \mathbb{R})$}
\AXC{$\varphi_{221}$\lbl{A2}}
\BIC{$b_{n}^{-1} \in \mathbb{R} \to \beta^{-1} \in \mathbb{R} \to b_{n}^{-1} \cdot \beta^{-1} \in \mathbb{R}$}
\end{ptree}
\noindent$\Pi_{126}$:
\begin{ptree}
\AXC{$\stackrel{\Pi_{120}}{\vdots}$}\noLine
\UIC{$\varphi_{313}$}
\AXC{$\stackrel{\Pi_{121}}{\vdots}$}\noLine
\UIC{$\varphi_{321}$}
\AXC{$\stackrel{\Pi_{125}}{\vdots}$}\noLine
\UIC{$b_{n}^{-1} \in \mathbb{R} \to \beta^{-1} \in \mathbb{R} \to b_{n}^{-1} \cdot \beta^{-1} \in \mathbb{R}$}
\AXC{$(b_{n}^{-1} \in \mathbb{R} \to \beta^{-1} \in \mathbb{R} \to b_{n}^{-1} \cdot \beta^{-1} \in \mathbb{R}) \to \varphi_{321} \to \varphi_{313} \to \varphi_{340}$\lbl{A1}}
\BIC{$\varphi_{321} \to \varphi_{313} \to \varphi_{340}$}
\BIC{$\varphi_{313} \to \varphi_{340}$}
\BIC{$\varphi_{340}$}
\end{ptree}
\noindent$\Pi_{127}$:
\begin{ptree}
\AXC{$\forall x\,\forall y\,(x \in \mathbb{R} \to y \in \mathbb{R} \to x - y \in \mathbb{R})$\lbl{[F1]}}
\AXC{$\varphi_{101}$\lbl{A2}}
\BIC{$\forall y\,(\beta \in \mathbb{R} \to y \in \mathbb{R} \to \beta - y \in \mathbb{R})$}
\AXC{$\varphi_{141}$\lbl{A2}}
\BIC{$\beta \in \mathbb{R} \to b_{n} \in \mathbb{R} \to \beta - b_{n} \in \mathbb{R}$}
\AXC{$(\beta \in \mathbb{R} \to b_{n} \in \mathbb{R} \to \beta - b_{n} \in \mathbb{R}) \to \varphi_{421}$\lbl{A1}}
\BIC{$\varphi_{421}$}
\end{ptree}
\noindent$\Pi_{128}$:
\begin{ptree}
\AXC{$\forall x\,\forall y\,(x \in \mathbb{R} \to y \in \mathbb{R} \to \lvert x \cdot y \rvert = \lvert x \rvert \cdot \lvert y \rvert)$\lbl{[A17]}}
\AXC{$\varphi_{243}$\lbl{A2}}
\BIC{$\varphi_{102}$}
\AXC{$\varphi_{102} \to \beta - b_{n} \in \mathbb{R} \to b_{n}^{-1} \cdot \beta^{-1} \in \mathbb{R} \to \varphi_{116}$\lbl{A2}}
\BIC{$\beta - b_{n} \in \mathbb{R} \to b_{n}^{-1} \cdot \beta^{-1} \in \mathbb{R} \to \varphi_{116}$}
\end{ptree}
\noindent$\Pi_{129}$:
\begin{ptree}
\AXC{$\stackrel{\Pi_{128}}{\vdots}$}\noLine
\UIC{$\beta - b_{n} \in \mathbb{R} \to b_{n}^{-1} \cdot \beta^{-1} \in \mathbb{R} \to \varphi_{116}$}
\AXC{$(\beta - b_{n} \in \mathbb{R} \to b_{n}^{-1} \cdot \beta^{-1} \in \mathbb{R} \to \varphi_{116}) \to \varphi_{329} \to \varphi_{340} \to \varphi_{402}$\lbl{A1}}
\BIC{$\varphi_{329} \to \varphi_{340} \to \varphi_{402}$}
\end{ptree}
\noindent$\Pi_{130}$:
\begin{ptree}
\AXC{$\stackrel{\Pi_{126}}{\vdots}$}\noLine
\UIC{$\varphi_{340}$}
\AXC{$\stackrel{\Pi_{7}}{\vdots}$}\noLine
\UIC{$\varphi_{314}$}
\AXC{$\stackrel{\Pi_{4}}{\vdots}$}\noLine
\UIC{$\mathrm{Conv}(b,\beta) \to \beta \neq 0 \to \forall n \in \mathbb{N}\,(b_{n} \neq 0) \to \beta \in \mathbb{R}$}
\AXC{$\stackrel{\Pi_{127}}{\vdots}$}\noLine
\UIC{$\varphi_{421}$}
\BIC{$\varphi_{314} \to \varphi_{329}$}
\BIC{$\varphi_{329}$}
\AXC{$\stackrel{\Pi_{129}}{\vdots}$}\noLine
\UIC{$\varphi_{329} \to \varphi_{340} \to \varphi_{402}$}
\BIC{$\varphi_{340} \to \varphi_{402}$}
\BIC{$\varphi_{402}$}
\end{ptree}
\noindent$\Pi_{131}$:
\begin{ptree}
\AXC{$z = w \to \lvert b_{n}^{-1} - \beta^{-1} \rvert = \lvert z \rvert \to \lvert b_{n}^{-1} - \beta^{-1} \rvert = \lvert w \rvert$\lbl{A6}}
\UIC{$\varphi_{57}$}
\UIC{$\forall z\,(\varphi_{57})$}
\AXC{$\forall z\,(\varphi_{57}) \to \varphi_{203}$\lbl{A2}}
\BIC{$\varphi_{203}$}
\end{ptree}
\noindent$\Pi_{132}$:
\begin{ptree}
\AXC{$\stackrel{\Pi_{7}}{\vdots}$}\noLine
\UIC{$\varphi_{314}$}
\AXC{$\varphi_{77}$\lbl{[E21]}}
\AXC{$\varphi_{77} \to \varphi_{185}$\lbl{A2}}
\BIC{$\varphi_{185}$}
\AXC{$\varphi_{185} \to b_{n} \in \mathbb{R} \to \beta \in \mathbb{R} \to \varphi_{78}$\lbl{A2}}
\BIC{$b_{n} \in \mathbb{R} \to \beta \in \mathbb{R} \to \varphi_{78}$}
\AXC{$(b_{n} \in \mathbb{R} \to \beta \in \mathbb{R} \to \varphi_{78}) \to \varphi_{314} \to \varphi_{430}$\lbl{A1}}
\BIC{$\varphi_{314} \to \varphi_{430}$}
\BIC{$\varphi_{430}$}
\end{ptree}
\noindent$\Pi_{133}$:
\begin{ptree}
\AXC{$\mathrm{Conv}(b,\beta) \to \beta \neq 0 \to \beta \neq 0$\lbl{A1}}
\AXC{$\stackrel{\Pi_{37}}{\vdots}$}\noLine
\UIC{$\varphi_{306}$}
\AXC{$\stackrel{\Pi_{4}}{\vdots}$}\noLine
\UIC{$\mathrm{Conv}(b,\beta) \to \beta \neq 0 \to \forall n \in \mathbb{N}\,(b_{n} \neq 0) \to \beta \in \mathbb{R}$}
\AXC{$\stackrel{\Pi_{132}}{\vdots}$}\noLine
\UIC{$\varphi_{430}$}
\BIC{$\varphi_{306} \to (\mathrm{Conv}(b,\beta) \to \beta \neq 0 \to \beta \neq 0) \to \varphi_{378}$}
\BIC{$(\mathrm{Conv}(b,\beta) \to \beta \neq 0 \to \beta \neq 0) \to \varphi_{378}$}
\BIC{$\varphi_{378}$}
\end{ptree}
\noindent$\Pi_{134}$:
\begin{ptree}
\AXC{$x = x$\lbl{A5}}
\UIC{$\forall x\,(x = x)$}
\AXC{$\forall x\,(x = x) \to \lvert b_{n}^{-1} - \beta^{-1} \rvert = \lvert b_{n}^{-1} - \beta^{-1} \rvert$\lbl{A2}}
\BIC{$\lvert b_{n}^{-1} - \beta^{-1} \rvert = \lvert b_{n}^{-1} - \beta^{-1} \rvert$}
\AXC{$\stackrel{\Pi_{133}}{\vdots}$}\noLine
\UIC{$\varphi_{378}$}
\AXC{$\varphi_{378} \to \varphi_{422}$\lbl{A1}}
\BIC{$\varphi_{422}$}
\BIC{$(b_{n}^{-1} - \beta^{-1} = (\beta - b_{n}) \cdot (b_{n}^{-1} \cdot \beta^{-1}) \to \varphi_{212}) \to \varphi_{386}$}
\end{ptree}
\noindent$\Pi_{135}$:
\begin{ptree}
\AXC{$\stackrel{\Pi_{131}}{\vdots}$}\noLine
\UIC{$\varphi_{203}$}
\AXC{$\varphi_{203} \to b_{n}^{-1} - \beta^{-1} = (\beta - b_{n}) \cdot (b_{n}^{-1} \cdot \beta^{-1}) \to \varphi_{212}$\lbl{A2}}
\BIC{$b_{n}^{-1} - \beta^{-1} = (\beta - b_{n}) \cdot (b_{n}^{-1} \cdot \beta^{-1}) \to \varphi_{212}$}
\AXC{$\stackrel{\Pi_{134}}{\vdots}$}\noLine
\UIC{$(b_{n}^{-1} - \beta^{-1} = (\beta - b_{n}) \cdot (b_{n}^{-1} \cdot \beta^{-1}) \to \varphi_{212}) \to \varphi_{386}$}
\BIC{$\varphi_{386}$}
\end{ptree}
\noindent$\Pi_{136}$:
\begin{ptree}
\AXC{$\stackrel{\Pi_{124}}{\vdots}$}\noLine
\UIC{$\varphi_{40}$}
\AXC{$\varphi_{40} \to \varphi_{256}$\lbl{A2}}
\BIC{$\varphi_{256}$}
\AXC{$\varphi_{256} \to \varphi_{322}$\lbl{A2}}
\BIC{$\varphi_{322}$}
\AXC{$\stackrel{\Pi_{130}}{\vdots}$}\noLine
\UIC{$\varphi_{402}$}
\AXC{$\stackrel{\Pi_{135}}{\vdots}$}\noLine
\UIC{$\varphi_{386}$}
\AXC{$\varphi_{386} \to \varphi_{402} \to \varphi_{322} \to \varphi_{392}$\lbl{A1}}
\BIC{$\varphi_{402} \to \varphi_{322} \to \varphi_{392}$}
\BIC{$\varphi_{322} \to \varphi_{392}$}
\BIC{$\varphi_{392}$}
\AXC{$\varphi_{392} \to \varphi_{404} \to \varphi_{356} \to \varphi_{395}$\lbl{A1}}
\BIC{$\varphi_{404} \to \varphi_{356} \to \varphi_{395}$}
\end{ptree}
\noindent$\Pi_{137}$:
\begin{ptree}
\AXC{$\stackrel{\Pi_{118}}{\vdots}$}\noLine
\UIC{$\lvert b_{n}^{-1} \cdot \beta^{-1} \rvert = \lvert b_{n}^{-1} \rvert \cdot \lvert \beta^{-1} \rvert \to \varphi_{133} \to \varphi_{157}$}
\AXC{$\stackrel{\Pi_{123}}{\vdots}$}\noLine
\UIC{$(\lvert b_{n}^{-1} \cdot \beta^{-1} \rvert = \lvert b_{n}^{-1} \rvert \cdot \lvert \beta^{-1} \rvert \to \varphi_{133} \to \varphi_{157}) \to \varphi_{404}$}
\BIC{$\varphi_{404}$}
\AXC{$\stackrel{\Pi_{136}}{\vdots}$}\noLine
\UIC{$\varphi_{404} \to \varphi_{356} \to \varphi_{395}$}
\BIC{$\varphi_{356} \to \varphi_{395}$}
\end{ptree}
\noindent$\Pi_{138}$:
\begin{ptree}
\AXC{$\stackrel{\Pi_{117}}{\vdots}$}\noLine
\UIC{$\varphi_{40}$}
\AXC{$\varphi_{40} \to \varphi_{261}$\lbl{A2}}
\BIC{$\varphi_{261}$}
\AXC{$\varphi_{261} \to \varphi_{356}$\lbl{A2}}
\BIC{$\varphi_{356}$}
\AXC{$\stackrel{\Pi_{137}}{\vdots}$}\noLine
\UIC{$\varphi_{356} \to \varphi_{395}$}
\BIC{$\varphi_{395}$}
\AXC{$\varphi_{395} \to \varphi_{407} \to \varphi_{367} \to \varphi_{396}$\lbl{A1}}
\BIC{$\varphi_{407} \to \varphi_{367} \to \varphi_{396}$}
\end{ptree}
\noindent$\Pi_{139}$:
\begin{ptree}
\AXC{$\stackrel{\Pi_{114}}{\vdots}$}\noLine
\UIC{$\lvert b_{n}^{-1} \rvert = \lvert b_{n} \rvert^{-1} \to \varphi_{165} \to \varphi_{168}$}
\AXC{$\stackrel{\Pi_{116}}{\vdots}$}\noLine
\UIC{$(\lvert b_{n}^{-1} \rvert = \lvert b_{n} \rvert^{-1} \to \varphi_{165} \to \varphi_{168}) \to \varphi_{407}$}
\BIC{$\varphi_{407}$}
\AXC{$\stackrel{\Pi_{138}}{\vdots}$}\noLine
\UIC{$\varphi_{407} \to \varphi_{367} \to \varphi_{396}$}
\BIC{$\varphi_{367} \to \varphi_{396}$}
\end{ptree}
\noindent$\Pi_{140}$:
\begin{ptree}
\AXC{$\stackrel{\Pi_{112}}{\vdots}$}\noLine
\UIC{$\varphi_{406}$}
\AXC{$\stackrel{\Pi_{113}}{\vdots}$}\noLine
\UIC{$\varphi_{40}$}
\AXC{$\varphi_{40} \to \varphi_{270}$\lbl{A2}}
\BIC{$\varphi_{270}$}
\AXC{$\varphi_{270} \to \varphi_{367}$\lbl{A2}}
\BIC{$\varphi_{367}$}
\AXC{$\stackrel{\Pi_{139}}{\vdots}$}\noLine
\UIC{$\varphi_{367} \to \varphi_{396}$}
\BIC{$\varphi_{396}$}
\AXC{$\varphi_{396} \to \varphi_{406} \to \varphi_{366} \to \varphi_{397}$\lbl{A1}}
\BIC{$\varphi_{406} \to \varphi_{366} \to \varphi_{397}$}
\BIC{$\varphi_{366} \to \varphi_{397}$}
\end{ptree}
\noindent$\Pi_{141}$:
\begin{ptree}
\AXC{$\stackrel{\Pi_{107}}{\vdots}$}\noLine
\UIC{$\varphi_{40}$}
\AXC{$\varphi_{40} \to \varphi_{269}$\lbl{A2}}
\BIC{$\varphi_{269}$}
\AXC{$\varphi_{269} \to \varphi_{366}$\lbl{A2}}
\BIC{$\varphi_{366}$}
\AXC{$\stackrel{\Pi_{140}}{\vdots}$}\noLine
\UIC{$\varphi_{366} \to \varphi_{397}$}
\BIC{$\varphi_{397}$}
\AXC{$\varphi_{397} \to \varphi_{405} \to \varphi_{365} \to \varphi_{398}$\lbl{A1}}
\BIC{$\varphi_{405} \to \varphi_{365} \to \varphi_{398}$}
\end{ptree}
\noindent$\Pi_{142}$:
\begin{ptree}
\AXC{$\stackrel{\Pi_{102}}{\vdots}$}\noLine
\UIC{$\lvert \beta - b_{n} \rvert = \lvert b_{n} - \beta \rvert \to \varphi_{163} \to \varphi_{166}$}
\AXC{$\stackrel{\Pi_{106}}{\vdots}$}\noLine
\UIC{$(\lvert \beta - b_{n} \rvert = \lvert b_{n} - \beta \rvert \to \varphi_{163} \to \varphi_{166}) \to \varphi_{405}$}
\BIC{$\varphi_{405}$}
\AXC{$\stackrel{\Pi_{141}}{\vdots}$}\noLine
\UIC{$\varphi_{405} \to \varphi_{365} \to \varphi_{398}$}
\BIC{$\varphi_{365} \to \varphi_{398}$}
\end{ptree}
\noindent$\Pi_{143}$:
\begin{ptree}
\AXC{$\stackrel{\Pi_{100}}{\vdots}$}\noLine
\UIC{$\varphi_{262}$}
\AXC{$\stackrel{\Pi_{101}}{\vdots}$}\noLine
\UIC{$\varphi_{40}$}
\AXC{$\varphi_{40} \to \varphi_{268}$\lbl{A2}}
\BIC{$\varphi_{268}$}
\AXC{$\varphi_{268} \to \varphi_{365}$\lbl{A2}}
\BIC{$\varphi_{365}$}
\AXC{$\stackrel{\Pi_{142}}{\vdots}$}\noLine
\UIC{$\varphi_{365} \to \varphi_{398}$}
\BIC{$\varphi_{398}$}
\AXC{$\varphi_{398} \to \varphi_{262} \to \varphi_{399}$\lbl{A1}}
\BIC{$\varphi_{262} \to \varphi_{399}$}
\BIC{$\varphi_{399}$}
\AXC{$\varphi_{399} \to \varphi_{375} \to \varphi_{274} \to \varphi_{343}$\lbl{A1}}
\BIC{$\varphi_{375} \to \varphi_{274} \to \varphi_{343}$}
\end{ptree}
\noindent$\Pi_{144}$:
\begin{ptree}
\AXC{$z = w \to \varphi_{245}$\lbl{A6}}
\UIC{$\forall w\,(z = w \to \varphi_{245})$}
\UIC{$\forall z\,\forall w\,(z = w \to \varphi_{245})$}
\AXC{$\forall z\,\forall w\,(z = w \to \varphi_{245}) \to \varphi_{385}$\lbl{A2}}
\BIC{$\varphi_{385}$}
\AXC{$\varphi_{385} \to \varphi_{387}$\lbl{A2}}
\BIC{$\varphi_{387}$}
\AXC{$\stackrel{\Pi_{87}}{\vdots}$}\noLine
\UIC{$\varphi_{387} \to \varphi_{389}$}
\BIC{$\varphi_{389}$}
\AXC{$\stackrel{\Pi_{99}}{\vdots}$}\noLine
\UIC{$\varphi_{389} \to \varphi_{375}$}
\BIC{$\varphi_{375}$}
\AXC{$\stackrel{\Pi_{143}}{\vdots}$}\noLine
\UIC{$\varphi_{375} \to \varphi_{274} \to \varphi_{343}$}
\BIC{$\varphi_{274} \to \varphi_{343}$}
\end{ptree}
\noindent$\Pi_{145}$:
\begin{ptree}
\AXC{$z = w \to z < \varepsilon \to w < \varepsilon$\lbl{A6}}
\UIC{$\forall w\,(z = w \to z < \varepsilon \to w < \varepsilon)$}
\UIC{$\forall z\,\forall w\,(z = w \to z < \varepsilon \to w < \varepsilon)$}
\AXC{$\forall z\,\forall w\,(z = w \to z < \varepsilon \to w < \varepsilon) \to \varphi_{213}$\lbl{A2}}
\BIC{$\varphi_{213}$}
\AXC{$\varphi_{213} \to \varphi_{274}$\lbl{A2}}
\BIC{$\varphi_{274}$}
\AXC{$\stackrel{\Pi_{144}}{\vdots}$}\noLine
\UIC{$\varphi_{274} \to \varphi_{343}$}
\BIC{$\varphi_{343}$}
\end{ptree}
\noindent$\Pi_{146}$:
\begin{ptree}
\AXC{$\varphi_{32}$\lbl{A1}}
\AXC{$\stackrel{\Pi_{5}}{\vdots}$}\noLine
\UIC{$\varphi_{5}$}
\AXC{$\varphi_{118}$\lbl{[S5]}}
\AXC{$\varphi_{118} \to \mathrm{Seq}(b) \to \varphi_{79}$\lbl{A2}}
\BIC{$\mathrm{Seq}(b) \to \varphi_{79}$}
\AXC{$(\mathrm{Seq}(b) \to \varphi_{79}) \to \varphi_{5} \to \varphi_{32} \to \mathrm{Conv}(b,\beta) \to \beta \neq 0 \to \varphi_{79}$\lbl{A1}}
\BIC{$\varphi_{5} \to \varphi_{32} \to \mathrm{Conv}(b,\beta) \to \beta \neq 0 \to \varphi_{79}$}
\BIC{$\varphi_{32} \to \mathrm{Conv}(b,\beta) \to \beta \neq 0 \to \varphi_{79}$}
\BIC{$\mathrm{Conv}(b,\beta) \to \beta \neq 0 \to \varphi_{79}$}
\end{ptree}
\noindent$\Pi_{147}$:
\begin{ptree}
\AXC{$\varphi_{75}$\lbl{A2}}
\AXC{$\stackrel{\Pi_{146}}{\vdots}$}\noLine
\UIC{$\mathrm{Conv}(b,\beta) \to \beta \neq 0 \to \varphi_{79}$}
\AXC{$(\mathrm{Conv}(b,\beta) \to \beta \neq 0 \to \varphi_{79}) \to \varphi_{362}$\lbl{A1}}
\BIC{$\varphi_{362}$}
\AXC{$\varphi_{362} \to \varphi_{75} \to \varphi_{353}$\lbl{A1}}
\BIC{$\varphi_{75} \to \varphi_{353}$}
\BIC{$\varphi_{353}$}
\AXC{$\varphi_{353} \to \varphi_{304} \to \varphi_{335}$\lbl{A1}}
\BIC{$\varphi_{304} \to \varphi_{335}$}
\end{ptree}
\noindent$\Pi_{148}$:
\begin{ptree}
\AXC{$\stackrel{\Pi_{24}}{\vdots}$}\noLine
\UIC{$\forall x\,\forall y\,(x = y \to y = x)$}
\AXC{$\varphi_{6}$\lbl{A2}}
\BIC{$\forall y\,(b^{\ast}(n) = y \to y = b^{\ast}(n))$}
\AXC{$\varphi_{127}$\lbl{A2}}
\BIC{$b^{\ast}(n) = b_{n}^{-1} \to b_{n}^{-1} = b^{\ast}(n)$}
\AXC{$\varphi_{304}$\lbl{A1}}
\AXC{$\stackrel{\Pi_{147}}{\vdots}$}\noLine
\UIC{$\varphi_{304} \to \varphi_{335}$}
\BIC{$\varphi_{335}$}
\AXC{$\varphi_{426}$\lbl{A1}}
\BIC{$(b^{\ast}(n) = b_{n}^{-1} \to b_{n}^{-1} = b^{\ast}(n)) \to \varphi_{336}$}
\BIC{$\varphi_{336}$}
\end{ptree}
\noindent$\Pi_{149}$:
\begin{ptree}
\AXC{$z = w \to \lvert z - \beta^{-1} \rvert < \varepsilon \to \lvert w - \beta^{-1} \rvert < \varepsilon$\lbl{A6}}
\UIC{$\forall w\,(z = w \to \lvert z - \beta^{-1} \rvert < \varepsilon \to \lvert w - \beta^{-1} \rvert < \varepsilon)$}
\UIC{$\forall z\,\forall w\,(z = w \to \lvert z - \beta^{-1} \rvert < \varepsilon \to \lvert w - \beta^{-1} \rvert < \varepsilon)$}
\AXC{$\varphi_{222}$\lbl{A2}}
\BIC{$\forall w\,(b_{n}^{-1} = w \to \lvert b_{n}^{-1} - \beta^{-1} \rvert < \varepsilon \to \lvert w - \beta^{-1} \rvert < \varepsilon)$}
\AXC{$\varphi_{278}$\lbl{A2}}
\BIC{$\varphi_{151}$}
\AXC{$\stackrel{\Pi_{145}}{\vdots}$}\noLine
\UIC{$\varphi_{343}$}
\AXC{$\stackrel{\Pi_{148}}{\vdots}$}\noLine
\UIC{$\varphi_{336}$}
\AXC{$\varphi_{336} \to \varphi_{343} \to \varphi_{151} \to \varphi_{347}$\lbl{A1}}
\BIC{$\varphi_{343} \to \varphi_{151} \to \varphi_{347}$}
\BIC{$\varphi_{151} \to \varphi_{347}$}
\BIC{$\varphi_{347}$}
\end{ptree}
\noindent$\Pi_{150}$:
\begin{ptree}
\AXC{$\varphi_{285} \to \forall n\,(\varphi_{285})$\lbl{A4}}
\AXC{$\varphi_{419}$\lbl{A3}}
\AXC{$\stackrel{\Pi_{149}}{\vdots}$}\noLine
\UIC{$\varphi_{347}$}
\AXC{$\varphi_{347} \to \varphi_{285} \to n \in \mathbb{N} \to N \leq n \to \lvert b^{\ast}(n) - \beta^{-1} \rvert < \varepsilon$\lbl{A1}}
\BIC{$\varphi_{285} \to n \in \mathbb{N} \to N \leq n \to \lvert b^{\ast}(n) - \beta^{-1} \rvert < \varepsilon$}
\UIC{$\forall n\,(\varphi_{285} \to n \in \mathbb{N} \to N \leq n \to \lvert b^{\ast}(n) - \beta^{-1} \rvert < \varepsilon)$}
\AXC{$\varphi_{437}$\lbl{A1}}
\BIC{$\varphi_{419} \to (\varphi_{285} \to \forall n\,(\varphi_{285})) \to \varphi_{354}$}
\BIC{$(\varphi_{285} \to \forall n\,(\varphi_{285})) \to \varphi_{354}$}
\BIC{$\varphi_{354}$}
\end{ptree}
\noindent$\Pi_{151}$:
\begin{ptree}
\AXC{$\forall N\,(\varphi_{25}) \to \varphi_{25}$\lbl{A2}}
\AXC{$\stackrel{\Pi_{150}}{\vdots}$}\noLine
\UIC{$\varphi_{354}$}
\AXC{$\varphi_{350}$\lbl{A1}}
\AXC{$\varphi_{350} \to \varphi_{295}$\lbl{A1}}
\BIC{$\varphi_{295}$}
\AXC{$\varphi_{295} \to \varphi_{354} \to \varphi_{368}$\lbl{A1}}
\BIC{$\varphi_{354} \to \varphi_{368}$}
\BIC{$\varphi_{368}$}
\AXC{$\varphi_{368} \to (\forall N\,(\varphi_{25}) \to \varphi_{25}) \to \varphi_{370}$\lbl{A1}}
\BIC{$(\forall N\,(\varphi_{25}) \to \varphi_{25}) \to \varphi_{370}$}
\BIC{$\varphi_{370}$}
\end{ptree}
\noindent$\Pi_{152}$:
\begin{ptree}
\AXC{$\varphi_{161}$\lbl{[C3]}}
\AXC{$\varphi_{161} \to \varphi_{153}$\lbl{A2}}
\BIC{$\varphi_{153}$}
\AXC{$\varphi_{153} \to \varphi_{142}$\lbl{A2}}
\BIC{$\varphi_{142}$}
\AXC{$\varphi_{142} \to \mathrm{Conv}(b,\beta) \to \varphi_{64}$\lbl{A2}}
\BIC{$\mathrm{Conv}(b,\beta) \to \varphi_{64}$}
\AXC{$(\mathrm{Conv}(b,\beta) \to \varphi_{64}) \to \varphi_{424}$\lbl{A1}}
\BIC{$\varphi_{424}$}
\end{ptree}
\noindent$\Pi_{153}$:
\begin{ptree}
\AXC{$\varphi_{305}$\lbl{A1}}
\AXC{$\varphi_{305} \to \mathrm{Conv}(b,\beta) \to \varphi_{252}$\lbl{A1}}
\BIC{$\mathrm{Conv}(b,\beta) \to \varphi_{252}$}
\AXC{$\mathrm{Conv}(b,\beta) \to \mathrm{Conv}(b,\beta)$\lbl{A1}}
\AXC{$\stackrel{\Pi_{152}}{\vdots}$}\noLine
\UIC{$\varphi_{424}$}
\BIC{$(\mathrm{Conv}(b,\beta) \to \varphi_{252}) \to (\mathrm{Conv}(b,\beta) \to \varphi_{242}) \to \varphi_{287}$}
\BIC{$(\mathrm{Conv}(b,\beta) \to \varphi_{242}) \to \varphi_{287}$}
\end{ptree}
\noindent$\Pi_{154}$:
\begin{ptree}
\AXC{$\varphi_{305}$\lbl{A1}}
\AXC{$\varphi_{305} \to \mathrm{Conv}(b,\beta) \to \varphi_{242}$\lbl{A1}}
\BIC{$\mathrm{Conv}(b,\beta) \to \varphi_{242}$}
\AXC{$\stackrel{\Pi_{153}}{\vdots}$}\noLine
\UIC{$(\mathrm{Conv}(b,\beta) \to \varphi_{242}) \to \varphi_{287}$}
\BIC{$\varphi_{287}$}
\end{ptree}
\noindent$\Pi_{155}$:
\begin{ptree}
\AXC{$\stackrel{\Pi_{151}}{\vdots}$}\noLine
\UIC{$\varphi_{370}$}
\AXC{$\varphi_{423}$\lbl{A1}}
\BIC{$\varphi_{300} \to \neg (N \in \mathbb{N} \land \forall n \in \mathbb{N}\,(N \leq n \to \lvert b_{n} - \beta \rvert < d_{0}))$}
\UIC{$\varphi_{374}$}
\AXC{$\stackrel{\Pi_{154}}{\vdots}$}\noLine
\UIC{$\varphi_{287}$}
\AXC{$\varphi_{287} \to \varphi_{374} \to (\varphi_{374} \to \forall N\,(\varphi_{300}) \to \varphi_{12}) \to (\varphi_{300} \to \forall N\,(\varphi_{300})) \to \varphi_{299}$\lbl{A1}}
\BIC{$\varphi_{374} \to (\varphi_{374} \to \forall N\,(\varphi_{300}) \to \varphi_{12}) \to (\varphi_{300} \to \forall N\,(\varphi_{300})) \to \varphi_{299}$}
\BIC{$(\varphi_{374} \to \forall N\,(\varphi_{300}) \to \varphi_{12}) \to (\varphi_{300} \to \forall N\,(\varphi_{300})) \to \varphi_{299}$}
\end{ptree}
\noindent$\Pi_{156}$:
\begin{ptree}
\AXC{$\varphi_{300} \to \forall N\,(\varphi_{300})$\lbl{A4}}
\AXC{$\varphi_{374} \to \forall N\,(\varphi_{300}) \to \varphi_{12}$\lbl{A3}}
\AXC{$\stackrel{\Pi_{155}}{\vdots}$}\noLine
\UIC{$(\varphi_{374} \to \forall N\,(\varphi_{300}) \to \varphi_{12}) \to (\varphi_{300} \to \forall N\,(\varphi_{300})) \to \varphi_{299}$}
\BIC{$(\varphi_{300} \to \forall N\,(\varphi_{300})) \to \varphi_{299}$}
\BIC{$\varphi_{299}$}
\AXC{$\varphi_{299} \to \varphi_{227} \to \neg \varphi_{62}$\lbl{A1}}
\BIC{$\varphi_{227} \to \neg \varphi_{62}$}
\UIC{$\forall d_{0}\,(\varphi_{227} \to \neg \varphi_{62})$}
\end{ptree}
\noindent$\Pi_{157}$:
\begin{ptree}
\AXC{$\forall x\,\forall y\,(x \in \mathbb{R} \to y \in \mathbb{R} \to x \cdot y \in \mathbb{R})$\lbl{(R1)}}
\AXC{$\varphi_{108}$\lbl{A2}}
\BIC{$\forall y\,(\varepsilon \in \mathbb{R} \to y \in \mathbb{R} \to \varepsilon \cdot y \in \mathbb{R})$}
\AXC{$\forall y\,(\varepsilon \in \mathbb{R} \to y \in \mathbb{R} \to \varepsilon \cdot y \in \mathbb{R}) \to \varepsilon \in \mathbb{R} \to \varphi_{59}$\lbl{A2}}
\BIC{$\varepsilon \in \mathbb{R} \to \varphi_{59}$}
\end{ptree}
\noindent$\Pi_{158}$:
\begin{ptree}
\AXC{$\stackrel{\Pi_{51}}{\vdots}$}\noLine
\UIC{$\mathrm{Conv}(b,\beta) \to \beta \neq 0 \to \varphi_{44}$}
\AXC{$\varphi_{58}$\lbl{[D8]}}
\AXC{$\varphi_{58} \to \varphi_{200}$\lbl{A2}}
\BIC{$\varphi_{200}$}
\AXC{$\varphi_{200} \to \lvert \beta \rvert \cdot 2^{-1} \in \mathbb{R} \to \varphi_{280}$\lbl{A2}}
\BIC{$\lvert \beta \rvert \cdot 2^{-1} \in \mathbb{R} \to \varphi_{280}$}
\AXC{$\varphi_{433}$\lbl{A1}}
\BIC{$(\mathrm{Conv}(b,\beta) \to \beta \neq 0 \to \varphi_{44}) \to \varphi_{417}$}
\BIC{$\varphi_{417}$}
\end{ptree}
\noindent$\Pi_{159}$:
\begin{ptree}
\AXC{$\stackrel{\Pi_{15}}{\vdots}$}\noLine
\UIC{$\varphi_{22}$}
\AXC{$\stackrel{\Pi_{53}}{\vdots}$}\noLine
\UIC{$\mathrm{Conv}(b,\beta) \to \beta \neq 0 \to \varphi_{87}$}
\AXC{$\varphi_{39}$\lbl{A1}}
\AXC{$\stackrel{\Pi_{157}}{\vdots}$}\noLine
\UIC{$\varepsilon \in \mathbb{R} \to \varphi_{59}$}
\AXC{$(\varepsilon \in \mathbb{R} \to \varphi_{59}) \to \varphi_{39} \to \varphi_{296}$\lbl{A1}}
\BIC{$\varphi_{39} \to \varphi_{296}$}
\BIC{$\varphi_{296}$}
\BIC{$\mathrm{Conv}(b,\beta) \to \beta \neq 0 \to \varphi_{113}$}
\AXC{$\stackrel{\Pi_{158}}{\vdots}$}\noLine
\UIC{$\varphi_{417}$}
\BIC{$\varphi_{22} \to (\mathrm{Conv}(b,\beta) \to \beta \neq 0 \to \varphi_{96}) \to \varphi_{247}$}
\BIC{$(\mathrm{Conv}(b,\beta) \to \beta \neq 0 \to \varphi_{96}) \to \varphi_{247}$}
\end{ptree}
\noindent$\Pi_{160}$:
\begin{ptree}
\AXC{$\stackrel{\Pi_{55}}{\vdots}$}\noLine
\UIC{$\mathrm{Conv}(b,\beta) \to \beta \neq 0 \to \varphi_{96}$}
\AXC{$\stackrel{\Pi_{159}}{\vdots}$}\noLine
\UIC{$(\mathrm{Conv}(b,\beta) \to \beta \neq 0 \to \varphi_{96}) \to \varphi_{247}$}
\BIC{$\varphi_{247}$}
\AXC{$\varphi_{247} \to \forall d_{0}\,(\varphi_{227} \to \neg \varphi_{62}) \to \varphi_{438}$\lbl{A1}}
\BIC{$\forall d_{0}\,(\varphi_{227} \to \neg \varphi_{62}) \to \varphi_{438}$}
\end{ptree}
\noindent$\Pi_{161}$:
\begin{ptree}
\AXC{$\varphi_{227} \to \forall d_{0}\,(\varphi_{227})$\lbl{A4}}
\AXC{$\forall d_{0}\,(\varphi_{227} \to \neg \varphi_{62}) \to \forall d_{0}\,(\varphi_{227}) \to \forall d_{0}\,(\neg \varphi_{62})$\lbl{A3}}
\AXC{$\stackrel{\Pi_{156}}{\vdots}$}\noLine
\UIC{$\forall d_{0}\,(\varphi_{227} \to \neg \varphi_{62})$}
\AXC{$\stackrel{\Pi_{160}}{\vdots}$}\noLine
\UIC{$\forall d_{0}\,(\varphi_{227} \to \neg \varphi_{62}) \to \varphi_{438}$}
\BIC{$\varphi_{438}$}
\BIC{$(\varphi_{227} \to \forall d_{0}\,(\varphi_{227})) \to \forall d\,(\forall d_{0}\,(\neg \varphi_{62}) \to \varphi_{51}) \to \varphi_{431}$}
\BIC{$\forall d\,(\forall d_{0}\,(\neg \varphi_{62}) \to \varphi_{51}) \to \varphi_{431}$}
\end{ptree}
\noindent$\Pi_{162}$:
\begin{ptree}
\AXC{$\forall d\,(\forall d_{0}\,(\neg \varphi_{62}) \to \varphi_{51}) \to \forall d\,\forall d_{0}\,(\neg \varphi_{62}) \to \forall d\,(\varphi_{51})$\lbl{A3}}
\AXC{$\forall d_{0}\,(\neg \varphi_{62}) \to \varphi_{51}$\lbl{A2}}
\UIC{$\forall d\,(\forall d_{0}\,(\neg \varphi_{62}) \to \varphi_{51})$}
\AXC{$\stackrel{\Pi_{161}}{\vdots}$}\noLine
\UIC{$\forall d\,(\forall d_{0}\,(\neg \varphi_{62}) \to \varphi_{51}) \to \varphi_{431}$}
\BIC{$\varphi_{431}$}
\BIC{$(\forall d_{0}\,(\neg \varphi_{62}) \to \forall d\,\forall d_{0}\,(\neg \varphi_{62})) \to \varphi_{223}$}
\end{ptree}
\noindent$\Pi_{163}$:
\begin{ptree}
\AXC{$\forall d_{0}\,(\neg \varphi_{62}) \to \forall d\,\forall d_{0}\,(\neg \varphi_{62})$\lbl{A4}}
\AXC{$\stackrel{\Pi_{162}}{\vdots}$}\noLine
\UIC{$(\forall d_{0}\,(\neg \varphi_{62}) \to \forall d\,\forall d_{0}\,(\neg \varphi_{62})) \to \varphi_{223}$}
\BIC{$\varphi_{223}$}
\AXC{$\varphi_{223} \to \varphi_{224}$\lbl{A1}}
\BIC{$\varphi_{224}$}
\UIC{$\forall \varepsilon\,(\varphi_{224})$}
\end{ptree}
\noindent$\Pi_{164}$:
\begin{ptree}
\AXC{$\varphi_{122}$\lbl{A4}}
\AXC{$\forall \varepsilon\,(\varphi_{224}) \to \varphi_{248}$\lbl{A3}}
\AXC{$\stackrel{\Pi_{163}}{\vdots}$}\noLine
\UIC{$\forall \varepsilon\,(\varphi_{224})$}
\AXC{$\forall \varepsilon\,(\varphi_{224}) \to (\forall \varepsilon\,(\varphi_{224}) \to \varphi_{248}) \to \varphi_{122} \to \varphi_{235}$\lbl{A1}}
\BIC{$(\forall \varepsilon\,(\varphi_{224}) \to \varphi_{248}) \to \varphi_{122} \to \varphi_{235}$}
\BIC{$\varphi_{122} \to \varphi_{235}$}
\BIC{$\varphi_{235}$}
\end{ptree}
\noindent$\Pi_{165}$:
\begin{ptree}
\AXC{$\forall x \in \mathbb{R}\,(x \neq 0 \to x^{-1} \in \mathbb{R})$\lbl{(R1)}}
\AXC{$\varphi_{41}$\lbl{A2}}
\BIC{$\beta \in \mathbb{R} \to \beta \neq 0 \to \beta^{-1} \in \mathbb{R}$}
\AXC{$(\beta \in \mathbb{R} \to \beta \neq 0 \to \beta^{-1} \in \mathbb{R}) \to \varphi_{244}$\lbl{A1}}
\BIC{$\varphi_{244}$}
\end{ptree}
\noindent$\Pi_{166}$:
\begin{ptree}
\AXC{$\mathrm{Conv}(b,\beta) \to \beta \neq 0 \to \beta \neq 0$\lbl{A1}}
\AXC{$\stackrel{\Pi_{4}}{\vdots}$}\noLine
\UIC{$\mathrm{Conv}(b,\beta) \to \beta \neq 0 \to \forall n \in \mathbb{N}\,(b_{n} \neq 0) \to \beta \in \mathbb{R}$}
\AXC{$\stackrel{\Pi_{165}}{\vdots}$}\noLine
\UIC{$\varphi_{244}$}
\BIC{$\varphi_{92}$}
\BIC{$\mathrm{Conv}(b,\beta) \to \beta \neq 0 \to \forall n \in \mathbb{N}\,(b_{n} \neq 0) \to \beta^{-1} \in \mathbb{R}$}
\end{ptree}
\noindent$\Pi_{167}$:
\begin{ptree}
\AXC{$\stackrel{\Pi_{146}}{\vdots}$}\noLine
\UIC{$\mathrm{Conv}(b,\beta) \to \beta \neq 0 \to \varphi_{79}$}
\AXC{$(\mathrm{Conv}(b,\beta) \to \beta \neq 0 \to \varphi_{79}) \to \varphi_{13}$\lbl{A1}}
\BIC{$\varphi_{13}$}
\end{ptree}
\noindent$\Pi_{168}$:
\begin{ptree}
\AXC{$\stackrel{\Pi_{167}}{\vdots}$}\noLine
\UIC{$\varphi_{13}$}
\AXC{$\varphi_{210}$\lbl{DefConv}}
\AXC{$\varphi_{210} \to \varphi_{231}$\lbl{A2}}
\BIC{$\varphi_{231}$}
\AXC{$\varphi_{231} \to \varphi_{254}$\lbl{A2}}
\BIC{$\varphi_{254}$}
\AXC{$\varphi_{254} \to \varphi_{13} \to \varphi_{371}$\lbl{A1}}
\BIC{$\varphi_{13} \to \varphi_{371}$}
\BIC{$\varphi_{371}$}
\end{ptree}
\noindent 主木:
\begin{ptree}
\AXC{$\stackrel{\Pi_{164}}{\vdots}$}\noLine
\UIC{$\varphi_{235}$}
\AXC{$\stackrel{\Pi_{166}}{\vdots}$}\noLine
\UIC{$\mathrm{Conv}(b,\beta) \to \beta \neq 0 \to \forall n \in \mathbb{N}\,(b_{n} \neq 0) \to \beta^{-1} \in \mathbb{R}$}
\AXC{$\stackrel{\Pi_{168}}{\vdots}$}\noLine
\UIC{$\varphi_{371}$}
\BIC{$\varphi_{235} \to \varphi_{24}$}
\BIC{$\varphi_{24}$}
\UIC{$\forall \beta\,(\varphi_{24})$}
\UIC{$\forall b\,\forall \beta\,(\varphi_{24})$}
\end{ptree}
\begin{abbrevlist}
\item $\varphi_{1} :\equiv \forall \alpha\,(\mathrm{Conv}(b,\alpha) \to \mathrm{Seq}(b)) \to \mathrm{Conv}(b,\beta) \to \mathrm{Seq}(b)$
\item $\varphi_{2} :\equiv \forall a\,\forall \alpha\,(\mathrm{Conv}(a,\alpha) \to \mathrm{Seq}(a)) \to \forall \alpha\,(\mathrm{Conv}(b,\alpha) \to \mathrm{Seq}(b))$
\item $\varphi_{3} :\equiv \forall x\,\forall y\,(x = y \to y = x) \to \forall y\,(\lvert b_{n} \rvert = y \to y = \lvert b_{n} \rvert)$
\item $\varphi_{4} :\equiv \forall \alpha\,(\mathrm{Conv}(b,\alpha) \to \alpha \in \mathbb{R}) \to \mathrm{Conv}(b,\beta) \to \beta \in \mathbb{R}$
\item $\varphi_{5} :\equiv \mathrm{Conv}(b,\beta) \to \beta \neq 0 \to \forall n \in \mathbb{N}\,(b_{n} \neq 0) \to \mathrm{Seq}(b)$
\item $\varphi_{6} :\equiv \forall x\,\forall y\,(x = y \to y = x) \to \forall y\,(b^{\ast}(n) = y \to y = b^{\ast}(n))$
\item $\varphi_{7} :\equiv \mathrm{Conv}(b,\beta) \to \beta \neq 0 \to \forall n \in \mathbb{N}\,(b_{n} \neq 0) \to 0 \neq \lvert \beta \rvert$
\item $\varphi_{8} :\equiv \mathrm{Conv}(b,\beta) \to \beta \neq 0 \to \forall n \in \mathbb{N}\,(b_{n} \neq 0) \to \lvert \beta \rvert \neq 0$
\item $\varphi_{9} :\equiv \forall a\,\forall \alpha\,(\mathrm{Conv}(a,\alpha) \to \alpha \in \mathbb{R}) \to \forall \alpha\,(\mathrm{Conv}(b,\alpha) \to \alpha \in \mathbb{R})$
\item $\varphi_{10} :\equiv \mathrm{Conv}(b,\beta) \to \beta \neq 0 \to \forall n \in \mathbb{N}\,(b_{n} \neq 0) \to 0 < \lvert \beta \rvert$
\item $\varphi_{11} :\equiv \forall x\,\forall y\,(x = y \to y = x) \to \forall y\,(\lvert \beta \rvert \cdot 2^{-1} = y \to y = \lvert \beta \rvert \cdot 2^{-1})$
\item $\varphi_{12} :\equiv \forall N\,(\neg (N \in \mathbb{N} \land \forall n \in \mathbb{N}\,(N \leq n \to \lvert b_{n} - \beta \rvert < d_{0})))$
\item $\varphi_{13} :\equiv \mathrm{Conv}(b,\beta) \to \beta \neq 0 \to \forall n \in \mathbb{N}\,(b_{n} \neq 0) \to \mathrm{Seq}(b^{\ast})$
\item $\varphi_{14} :\equiv \forall y\,(0 < y \to 0 \neq y) \to 0 < \lvert \beta \rvert \cdot 2^{-1} \to 0 \neq \lvert \beta \rvert \cdot 2^{-1}$
\item $\varphi_{15} :\equiv \mathrm{Conv}(b,\beta) \to \beta \neq 0 \to \forall n \in \mathbb{N}\,(b_{n} \neq 0) \to \lvert \beta \rvert \in \mathbb{R}$
\item $\varphi_{16} :\equiv \forall x\,\forall y\,(x \in \mathbb{R} \to y \in \mathbb{R} \to 0 < x \to x < y \to y^{-1} < x^{-1})$
\item $\varphi_{17} :\equiv \forall y\,(\lvert \beta \rvert \cdot 2^{-1} \cdot \lvert \beta \rvert = y \to y = \lvert \beta \rvert \cdot 2^{-1} \cdot \lvert \beta \rvert)$
\item $\varphi_{18} :\equiv \mathrm{Conv}(b,\beta) \to \beta \neq 0 \to \forall n \in \mathbb{N}\,(b_{n} \neq 0) \to 0 \neq \lvert \beta \rvert \cdot 2^{-1}$
\item $\varphi_{19} :\equiv \mathrm{Conv}(b,\beta) \to \beta \neq 0 \to \forall n \in \mathbb{N}\,(b_{n} \neq 0) \to \lvert \beta \rvert \cdot 2^{-1} \neq 0$
\item $\varphi_{20} :\equiv \forall x \in \mathbb{R}\,(x \neq 0 \to 0 < \lvert x \rvert) \to \beta \in \mathbb{R} \to \beta \neq 0 \to 0 < \lvert \beta \rvert$
\item $\varphi_{21} :\equiv \forall x\,\forall y\,\forall z\,(x \in \mathbb{R} \to y \in \mathbb{R} \to z \in \mathbb{R} \to x \cdot y \cdot z = x \cdot (y \cdot z))$
\item $\varphi_{22} :\equiv \mathrm{Conv}(b,\beta) \to \beta \neq 0 \to \forall n \in \mathbb{N}\,(b_{n} \neq 0) \to 0 < \lvert \beta \rvert \cdot 2^{-1}$
\item $\varphi_{23} :\equiv \forall y\,(\lvert b_{n} \rvert = y \to y = \lvert b_{n} \rvert) \to \lvert b_{n} \rvert = 0 \to 0 = \lvert b_{n} \rvert$
\item $\varphi_{24} :\equiv \mathrm{Conv}(b,\beta) \to \beta \neq 0 \to \forall n \in \mathbb{N}\,(b_{n} \neq 0) \to \mathrm{Conv}(b^{\ast},\beta^{-1})$
\item $\varphi_{25} :\equiv \neg (N \in \mathbb{N} \land \forall n \in \mathbb{N}\,(N \leq n \to \lvert b^{\ast}(n) - \beta^{-1} \rvert < \varepsilon))$
\item $\varphi_{26} :\equiv (\lvert \beta \rvert \cdot 2^{-1})^{-1} \cdot \lvert \beta \rvert^{-1} = (\lvert \beta \rvert \cdot 2^{-1} \cdot \lvert \beta \rvert)^{-1}$
\item $\varphi_{27} :\equiv (\lvert \beta \rvert \cdot 2^{-1} \cdot \lvert \beta \rvert)^{-1} = (\lvert \beta \rvert \cdot 2^{-1})^{-1} \cdot \lvert \beta \rvert^{-1}$
\item $\varphi_{28} :\equiv \mathrm{Conv}(b,\beta) \to \beta \neq 0 \to \forall n \in \mathbb{N}\,(b_{n} \neq 0) \to n \in \mathbb{N} \to b_{n} \neq 0$
\item $\varphi_{29} :\equiv N \leq n \to \lvert \beta \rvert \cdot 2^{-1} \cdot \lvert \beta \rvert \cdot (\lvert \beta \rvert \cdot 2^{-1} \cdot \lvert \beta \rvert)^{-1} = 1$
\item $\varphi_{30} :\equiv \forall y\,\forall z\,(\varepsilon \in \mathbb{R} \to y \in \mathbb{R} \to z \in \mathbb{R} \to \varepsilon \cdot y \cdot z = \varepsilon \cdot (y \cdot z))$
\item $\varphi_{31} :\equiv \forall x\,\forall y\,(x \in \mathbb{R} \to y \in \mathbb{R} \to x \neq 0 \to y \neq 0 \to (x \cdot y)^{-1} = x^{-1} \cdot y^{-1})$
\item $\varphi_{32} :\equiv \mathrm{Conv}(b,\beta) \to \beta \neq 0 \to \forall n \in \mathbb{N}\,(b_{n} \neq 0) \to \forall n \in \mathbb{N}\,(b_{n} \neq 0)$
\item $\varphi_{33} :\equiv \mathrm{Conv}(b,\beta) \to \beta \neq 0 \to \forall n \in \mathbb{N}\,(b_{n} \neq 0) \to \lvert \beta \rvert \cdot 2^{-1} \in \mathbb{R}$
\item $\varphi_{34} :\equiv \forall y\,(\lvert \beta \rvert \cdot 2^{-1} \in \mathbb{R} \to y \in \mathbb{R} \to \lvert \beta \rvert \cdot 2^{-1} \cdot y \in \mathbb{R})$
\item $\varphi_{35} :\equiv \forall y\,((\lvert \beta \rvert \cdot 2^{-1} \cdot \lvert \beta \rvert)^{-1} = y \to y = (\lvert \beta \rvert \cdot 2^{-1} \cdot \lvert \beta \rvert)^{-1})$
\item $\varphi_{36} :\equiv 0 < \lvert \beta \rvert \cdot 2^{-1} \cdot \lvert \beta \rvert \to 0 < \varepsilon \cdot (\lvert \beta \rvert \cdot 2^{-1} \cdot \lvert \beta \rvert)$
\item $\varphi_{37} :\equiv \lvert \beta \rvert \cdot 2^{-1} \cdot \lvert \beta \rvert \neq 0 \to (\lvert \beta \rvert \cdot 2^{-1} \cdot \lvert \beta \rvert)^{-1} \in \mathbb{R}$
\item $\varphi_{38} :\equiv n \in \mathbb{N} \to N \leq n \to \lvert b_{n} - \beta \rvert < \varepsilon \cdot (\lvert \beta \rvert \cdot 2^{-1} \cdot \lvert \beta \rvert)$
\item $\varphi_{39} :\equiv \mathrm{Conv}(b,\beta) \to \beta \neq 0 \to \forall n \in \mathbb{N}\,(b_{n} \neq 0) \to \varepsilon \in \mathbb{R} \to \varepsilon \in \mathbb{R}$
\item $\varphi_{40} :\equiv \forall y\,\forall z\,(\lvert b_{n}^{-1} - \beta^{-1} \rvert = y \to y = z \to \lvert b_{n}^{-1} - \beta^{-1} \rvert = z)$
\item $\varphi_{41} :\equiv \forall x \in \mathbb{R}\,(x \neq 0 \to x^{-1} \in \mathbb{R}) \to \beta \in \mathbb{R} \to \beta \neq 0 \to \beta^{-1} \in \mathbb{R}$
\item $\varphi_{42} :\equiv \beta \neq 0 \to \forall n \in \mathbb{N}\,(b_{n} \neq 0) \to \varepsilon \in \mathbb{R} \to 0 < \varepsilon \to \lvert \beta \rvert \in \mathbb{R}$
\item $\varphi_{43} :\equiv \varepsilon \cdot (\lvert \beta \rvert \cdot 2^{-1} \cdot \lvert \beta \rvert) \cdot ((\lvert \beta \rvert \cdot 2^{-1})^{-1} \cdot \lvert \beta \rvert^{-1}) = \varepsilon$
\item $\varphi_{44} :\equiv \forall n \in \mathbb{N}\,(b_{n} \neq 0) \to \varepsilon \in \mathbb{R} \to 0 < \varepsilon \to \lvert \beta \rvert \cdot 2^{-1} \in \mathbb{R}$
\item $\varphi_{45} :\equiv \varepsilon \cdot (\lvert \beta \rvert \cdot 2^{-1} \cdot \lvert \beta \rvert \cdot (\lvert \beta \rvert \cdot 2^{-1} \cdot \lvert \beta \rvert)^{-1}) = \varepsilon \cdot 1$
\item $\varphi_{46} :\equiv \varepsilon \cdot (\lvert \beta \rvert \cdot 2^{-1} \cdot \lvert \beta \rvert \cdot (\lvert \beta \rvert \cdot 2^{-1} \cdot \lvert \beta \rvert)^{-1}) = \varepsilon \cdot w$
\item $\varphi_{47} :\equiv \varepsilon \cdot (\lvert \beta \rvert \cdot 2^{-1} \cdot \lvert \beta \rvert \cdot (\lvert \beta \rvert \cdot 2^{-1} \cdot \lvert \beta \rvert)^{-1}) = \varepsilon \cdot z$
\item $\varphi_{48} :\equiv \forall n\,(\mathrm{Seq}(b) \to n \in \mathbb{N} \to b_{n} \in \mathbb{R}) \to \mathrm{Seq}(b) \to n \in \mathbb{N} \to b_{n} \in \mathbb{R}$
\item $\varphi_{49} :\equiv \forall x\,\forall y\,\forall z\,(x \in \mathbb{R} \to y \in \mathbb{R} \to z \in \mathbb{R} \to 0 < z \to x < y \to x \cdot z < y \cdot z)$
\item $\varphi_{50} :\equiv \forall x\,\forall y\,\forall z\,(x \in \mathbb{R} \to y \in \mathbb{R} \to z \in \mathbb{R} \to 0 < z \to x < y \to z \cdot x < z \cdot y)$
\item $\varphi_{51} :\equiv \neg (d \in \mathbb{R} \land 0 < d \land d \leq \lvert \beta \rvert \cdot 2^{-1} \land d \leq \varepsilon \cdot (\lvert \beta \rvert \cdot 2^{-1} \cdot \lvert \beta \rvert))$
\item $\varphi_{52} :\equiv N \in \mathbb{N} \land \forall n \in \mathbb{N}\,(N \leq n \to \lvert b_{n} - \beta \rvert < d_{0}) \to n \in \mathbb{N} \to N \leq n \to N \leq n$
\item $\varphi_{53} :\equiv N \in \mathbb{N} \land \forall n \in \mathbb{N}\,(N \leq n \to \lvert b_{n} - \beta \rvert < d_{0}) \to n \in \mathbb{N} \to n \in \mathbb{N}$
\item $\varphi_{54} :\equiv \lvert \beta \rvert \cdot 2^{-1} \in \mathbb{R} \to \lvert \beta \rvert \cdot 2^{-1} \neq 0 \to (\lvert \beta \rvert \cdot 2^{-1})^{-1} \in \mathbb{R}$
\item $\varphi_{55} :\equiv \mathrm{Conv}(b,\beta) \to \beta \neq 0 \to \forall n \in \mathbb{N}\,(b_{n} \neq 0) \to 0 < \lvert \beta \rvert \cdot 2^{-1} \cdot \lvert \beta \rvert$
\item $\varphi_{56} :\equiv \mathrm{Conv}(b,\beta) \to \beta \neq 0 \to \forall n \in \mathbb{N}\,(b_{n} \neq 0) \to \varepsilon \in \mathbb{R} \to 0 < \varepsilon \to 0 < \varepsilon$
\item $\varphi_{57} :\equiv \forall w\,(z = w \to \lvert b_{n}^{-1} - \beta^{-1} \rvert = \lvert z \rvert \to \lvert b_{n}^{-1} - \beta^{-1} \rvert = \lvert w \rvert)$
\item $\varphi_{58} :\equiv \forall p\,\forall q\,(p \in \mathbb{R} \to q \in \mathbb{R} \to 0 < p \to 0 < q \to \exists d\,(d \in \mathbb{R} \land 0 < d \land d \leq p \land d \leq q))$
\item $\varphi_{59} :\equiv \lvert \beta \rvert \cdot 2^{-1} \cdot \lvert \beta \rvert \in \mathbb{R} \to \varepsilon \cdot (\lvert \beta \rvert \cdot 2^{-1} \cdot \lvert \beta \rvert) \in \mathbb{R}$
\item $\varphi_{60} :\equiv \exists d\,(d \in \mathbb{R} \land 0 < d \land d \leq \lvert \beta \rvert \cdot 2^{-1} \land d \leq \varepsilon \cdot (\lvert \beta \rvert \cdot 2^{-1} \cdot \lvert \beta \rvert))$
\item $\varphi_{61} :\equiv \lvert \beta \rvert \cdot 2^{-1} \in \mathbb{R} \to \lvert \beta \rvert \in \mathbb{R} \to \lvert \beta \rvert \cdot 2^{-1} \cdot \lvert \beta \rvert \in \mathbb{R}$
\item $\varphi_{62} :\equiv d_{0} \in \mathbb{R} \land 0 < d_{0} \land d_{0} \leq \lvert \beta \rvert \cdot 2^{-1} \land d_{0} \leq \varepsilon \cdot (\lvert \beta \rvert \cdot 2^{-1} \cdot \lvert \beta \rvert)$
\item $\varphi_{63} :\equiv n \in \mathbb{N} \to N \leq n \to b_{n}^{-1} - \beta^{-1} = (\beta - b_{n}) \cdot (b_{n}^{-1} \cdot \beta^{-1})$
\item $\varphi_{64} :\equiv d_{0} \in \mathbb{R} \to 0 < d_{0} \to \exists N \in \mathbb{N}\,\forall n \in \mathbb{N}\,(N \leq n \to \lvert b_{n} - \beta \rvert < d_{0})$
\item $\varphi_{65} :\equiv 0 < \varepsilon \to \exists N \in \mathbb{N}\,\forall n \in \mathbb{N}\,(N \leq n \to \lvert b^{\ast}(n) - \beta^{-1} \rvert < \varepsilon)$
\item $\varphi_{66} :\equiv \forall y\,((\lvert \beta \rvert \cdot 2^{-1})^{-1} \in \mathbb{R} \to y \in \mathbb{R} \to (\lvert \beta \rvert \cdot 2^{-1})^{-1} \cdot y \in \mathbb{R})$
\item $\varphi_{67} :\equiv \forall y\,(0 < y \to 0 \neq y) \to 0 < \lvert \beta \rvert \cdot 2^{-1} \cdot \lvert \beta \rvert \to 0 \neq \lvert \beta \rvert \cdot 2^{-1} \cdot \lvert \beta \rvert$
\item $\varphi_{68} :\equiv \forall x\,\forall y\,(x = y \to y = x) \to \forall y\,(\lvert b_{n}^{-1} - \beta^{-1} \rvert = y \to y = \lvert b_{n}^{-1} - \beta^{-1} \rvert)$
\item $\varphi_{69} :\equiv \varepsilon \cdot (\lvert \beta \rvert \cdot 2^{-1} \cdot \lvert \beta \rvert) \cdot ((\lvert \beta \rvert \cdot 2^{-1})^{-1} \cdot \lvert \beta \rvert^{-1}) = \varepsilon \cdot 1$
\item $\varphi_{70} :\equiv N \leq n \to \varepsilon \cdot (\lvert \beta \rvert \cdot 2^{-1} \cdot \lvert \beta \rvert) \cdot (\lvert b_{n} \rvert^{-1} \cdot \lvert \beta \rvert^{-1}) < \varepsilon$
\item $\varphi_{71} :\equiv (\mathrm{Conv}(b,\beta) \to \mathrm{Conv}(b,\beta)) \to \mathrm{Conv}(b,\beta) \to \beta \neq 0 \to \forall n \in \mathbb{N}\,(b_{n} \neq 0) \to \beta \in \mathbb{R}$
\item $\varphi_{72} :\equiv \forall a\,\forall n\,(\mathrm{Seq}(a) \to n \in \mathbb{N} \to a_{n} \in \mathbb{R}) \to \forall n\,(\mathrm{Seq}(b) \to n \in \mathbb{N} \to b_{n} \in \mathbb{R})$
\item $\varphi_{73} :\equiv \forall y\,(\lvert \beta \rvert \cdot 2^{-1} = y \to y = \lvert \beta \rvert \cdot 2^{-1}) \to \lvert \beta \rvert \cdot 2^{-1} = 0 \to 0 = \lvert \beta \rvert \cdot 2^{-1}$
\item $\varphi_{74} :\equiv \forall x \in \mathbb{R}\,(\lvert x \rvert \in \mathbb{R}) \to b_{n} - \beta \in \mathbb{R} \to \lvert b_{n} - \beta \rvert \in \mathbb{R}$
\item $\varphi_{75} :\equiv \forall n \in \mathbb{N}\,(b^{\ast}(n) = b_{n}^{-1}) \to n \in \mathbb{N} \to b^{\ast}(n) = b_{n}^{-1}$
\item $\varphi_{76} :\equiv \forall x \in \mathbb{R}\,(x \neq 0 \to 0 < \lvert x \rvert) \to b_{n} \in \mathbb{R} \to b_{n} \neq 0 \to 0 < \lvert b_{n} \rvert$
\item $\varphi_{77} :\equiv \forall x\,\forall y\,(x \in \mathbb{R} \to y \in \mathbb{R} \to x \neq 0 \to y \neq 0 \to x^{-1} - y^{-1} = (y - x) \cdot (x^{-1} \cdot y^{-1}))$
\item $\varphi_{78} :\equiv b_{n} \neq 0 \to \beta \neq 0 \to b_{n}^{-1} - \beta^{-1} = (\beta - b_{n}) \cdot (b_{n}^{-1} \cdot \beta^{-1})$
\item $\varphi_{79} :\equiv \forall n \in \mathbb{N}\,(b_{n} \neq 0) \to \mathrm{Seq}(b^{\ast}) \land \forall n \in \mathbb{N}\,(b^{\ast}(n) = b_{n}^{-1})$
\item $\varphi_{80} :\equiv \forall x \in \mathbb{R}\,(0 < x \to 0 < x^{-1}) \to \lvert \beta \rvert \in \mathbb{R} \to 0 < \lvert \beta \rvert \to 0 < \lvert \beta \rvert^{-1}$
\item $\varphi_{81} :\equiv N \in \mathbb{N} \land \forall n \in \mathbb{N}\,(N \leq n \to \lvert b_{n} - \beta \rvert < d_{0}) \to n \in \mathbb{N} \to N \leq n \to b_{n} \neq 0$
\item $\varphi_{82} :\equiv \lvert b_{n} - \beta \rvert < d_{0} \to d_{0} \leq \lvert \beta \rvert \cdot 2^{-1} \to \lvert b_{n} - \beta \rvert < \lvert \beta \rvert \cdot 2^{-1}$
\item $\varphi_{83} :\equiv \beta \in \mathbb{R} \to \lvert b_{n} - \beta \rvert < \lvert \beta \rvert \cdot 2^{-1} \to \lvert \beta \rvert \cdot 2^{-1} < \lvert b_{n} \rvert$
\item $\varphi_{84} :\equiv N \in \mathbb{N} \land \forall n \in \mathbb{N}\,(N \leq n \to \lvert b_{n} - \beta \rvert < d_{0}) \to n \in \mathbb{N} \to N \leq n \to \varepsilon \cdot 1 = \varepsilon$
\item $\varphi_{85} :\equiv \forall n \in \mathbb{N}\,(N \leq n \to \lvert b_{n} - \beta \rvert < d_{0}) \to n \in \mathbb{N} \to N \leq n \to \lvert b_{n} - \beta \rvert < d_{0}$
\item $\varphi_{86} :\equiv \lvert \beta \rvert \in \mathbb{R} \to 0 < \lvert \beta \rvert \cdot 2^{-1} \to 0 < \lvert \beta \rvert \to 0 < \lvert \beta \rvert \cdot 2^{-1} \cdot \lvert \beta \rvert$
\item $\varphi_{87} :\equiv \forall n \in \mathbb{N}\,(b_{n} \neq 0) \to \varepsilon \in \mathbb{R} \to 0 < \varepsilon \to \lvert \beta \rvert \cdot 2^{-1} \cdot \lvert \beta \rvert \in \mathbb{R}$
\item $\varphi_{88} :\equiv \forall y\,(b_{n}^{-1} \in \mathbb{R} \to y \in \mathbb{R} \to \lvert b_{n}^{-1} \cdot y \rvert = \lvert b_{n}^{-1} \rvert \cdot \lvert y \rvert)$
\item $\varphi_{89} :\equiv N \leq n \to \lvert b_{n}^{-1} - \beta^{-1} \rvert = \lvert \beta - b_{n} \rvert \cdot (\lvert b_{n}^{-1} \rvert \cdot \lvert \beta^{-1} \rvert)$
\item $\varphi_{90} :\equiv \forall x \in \mathbb{R}\,(x \neq 0 \to x^{-1} \in \mathbb{R}) \to \lvert \beta \rvert \in \mathbb{R} \to \lvert \beta \rvert \neq 0 \to \lvert \beta \rvert^{-1} \in \mathbb{R}$
\item $\varphi_{91} :\equiv n \in \mathbb{N} \to N \leq n \to \lvert b_{n} \rvert^{-1} \cdot \lvert \beta \rvert^{-1} < (\lvert \beta \rvert \cdot 2^{-1})^{-1} \cdot \lvert \beta \rvert^{-1}$
\item $\varphi_{92} :\equiv (\mathrm{Conv}(b,\beta) \to \beta \neq 0 \to \beta \neq 0) \to \mathrm{Conv}(b,\beta) \to \beta \neq 0 \to \forall n \in \mathbb{N}\,(b_{n} \neq 0) \to \beta^{-1} \in \mathbb{R}$
\item $\varphi_{93} :\equiv N \in \mathbb{N} \land \forall n \in \mathbb{N}\,(N \leq n \to \lvert b_{n} - \beta \rvert < d_{0}) \to n \in \mathbb{N} \to N \leq n \to \lvert \beta \rvert \in \mathbb{R}$
\item $\varphi_{94} :\equiv N \in \mathbb{N} \land \forall n \in \mathbb{N}\,(N \leq n \to \lvert b_{n} - \beta \rvert < d_{0}) \to n \in \mathbb{N} \to N \leq n \to 0 \neq \lvert b_{n} \rvert$
\item $\varphi_{95} :\equiv N \in \mathbb{N} \land \forall n \in \mathbb{N}\,(N \leq n \to \lvert b_{n} - \beta \rvert < d_{0}) \to n \in \mathbb{N} \to N \leq n \to \lvert b_{n} \rvert \neq 0$
\item $\varphi_{96} :\equiv \forall n \in \mathbb{N}\,(b_{n} \neq 0) \to \varepsilon \in \mathbb{R} \to 0 < \varepsilon \to 0 < \varepsilon \cdot (\lvert \beta \rvert \cdot 2^{-1} \cdot \lvert \beta \rvert)$
\item $\varphi_{97} :\equiv N \in \mathbb{N} \land \forall n \in \mathbb{N}\,(N \leq n \to \lvert b_{n} - \beta \rvert < d_{0}) \to n \in \mathbb{N} \to N \leq n \to \beta^{-1} \in \mathbb{R}$
\item $\varphi_{98} :\equiv n \in \mathbb{N} \to N \leq n \to \lvert b_{n}^{-1} - \beta^{-1} \rvert = \lvert (\beta - b_{n}) \cdot (b_{n}^{-1} \cdot \beta^{-1}) \rvert$
\item $\varphi_{99} :\equiv \forall x\,\forall y\,\forall z\,(x < y \to y \leq z \to x < z) \to \forall y\,\forall z\,(\lvert b_{n} - \beta \rvert < y \to y \leq z \to \lvert b_{n} - \beta \rvert < z)$
\item $\varphi_{100} :\equiv \forall x \in \mathbb{R}\,(x \neq 0 \to \lvert x^{-1} \rvert = \lvert x \rvert^{-1}) \to \beta \in \mathbb{R} \to \beta \neq 0 \to \lvert \beta^{-1} \rvert = \lvert \beta \rvert^{-1}$
\item $\varphi_{101} :\equiv \forall x\,\forall y\,(x \in \mathbb{R} \to y \in \mathbb{R} \to x - y \in \mathbb{R}) \to \forall y\,(\beta \in \mathbb{R} \to y \in \mathbb{R} \to \beta - y \in \mathbb{R})$
\item $\varphi_{102} :\equiv \forall y\,(\beta - b_{n} \in \mathbb{R} \to y \in \mathbb{R} \to \lvert (\beta - b_{n}) \cdot y \rvert = \lvert \beta - b_{n} \rvert \cdot \lvert y \rvert)$
\item $\varphi_{103} :\equiv N \in \mathbb{N} \land \forall n \in \mathbb{N}\,(N \leq n \to \lvert b_{n} - \beta \rvert < d_{0}) \to n \in \mathbb{N} \to N \leq n \to b_{n} \in \mathbb{R}$
\item $\varphi_{104} :\equiv N \in \mathbb{N} \land \forall n \in \mathbb{N}\,(N \leq n \to \lvert b_{n} - \beta \rvert < d_{0}) \to n \in \mathbb{N} \to N \leq n \to 0 < \lvert \beta \rvert^{-1}$
\item $\varphi_{105} :\equiv \forall x\,\forall \beta\,(x \in \mathbb{R} \to \beta \in \mathbb{R} \to \lvert x - \beta \rvert < \lvert \beta \rvert \cdot 2^{-1} \to \lvert \beta \rvert \cdot 2^{-1} < \lvert x \rvert)$
\item $\varphi_{106} :\equiv \forall x \in \mathbb{R}\,(x \neq 0 \to x^{-1} \in \mathbb{R}) \to b_{n} \in \mathbb{R} \to b_{n} \neq 0 \to b_{n}^{-1} \in \mathbb{R}$
\item $\varphi_{107} :\equiv N \in \mathbb{N} \land \forall n \in \mathbb{N}\,(N \leq n \to \lvert b_{n} - \beta \rvert < d_{0}) \to n \in \mathbb{N} \to N \leq n \to 0 < \lvert b_{n} \rvert$
\item $\varphi_{108} :\equiv \forall x\,\forall y\,(x \in \mathbb{R} \to y \in \mathbb{R} \to x \cdot y \in \mathbb{R}) \to \forall y\,(\varepsilon \in \mathbb{R} \to y \in \mathbb{R} \to \varepsilon \cdot y \in \mathbb{R})$
\item $\varphi_{109} :\equiv \lvert \beta - b_{n} \rvert \cdot (\lvert b_{n}^{-1} \rvert \cdot \lvert \beta^{-1} \rvert) = \lvert \beta - b_{n} \rvert \cdot (w \cdot \lvert \beta^{-1} \rvert)$
\item $\varphi_{110} :\equiv \lvert \beta - b_{n} \rvert \cdot (\lvert b_{n}^{-1} \rvert \cdot \lvert \beta^{-1} \rvert) = \lvert \beta - b_{n} \rvert \cdot (z \cdot \lvert \beta^{-1} \rvert)$
\item $\varphi_{111} :\equiv N \in \mathbb{N} \land \forall n \in \mathbb{N}\,(N \leq n \to \lvert b_{n} - \beta \rvert < d_{0}) \to n \in \mathbb{N} \to N \leq n \to d_{0} \leq \lvert \beta \rvert \cdot 2^{-1}$
\item $\varphi_{112} :\equiv d_{0} \leq \varepsilon \cdot (\lvert \beta \rvert \cdot 2^{-1} \cdot \lvert \beta \rvert) \to \lvert b_{n} - \beta \rvert < \varepsilon \cdot (\lvert \beta \rvert \cdot 2^{-1} \cdot \lvert \beta \rvert)$
\item $\varphi_{113} :\equiv \forall n \in \mathbb{N}\,(b_{n} \neq 0) \to \varepsilon \in \mathbb{R} \to 0 < \varepsilon \to \varepsilon \cdot (\lvert \beta \rvert \cdot 2^{-1} \cdot \lvert \beta \rvert) \in \mathbb{R}$
\item $\varphi_{114} :\equiv \forall e \in \mathbb{R}\,(0 < e \to 0 < e \cdot 2^{-1}) \to \lvert \beta \rvert \in \mathbb{R} \to 0 < \lvert \beta \rvert \to 0 < \lvert \beta \rvert \cdot 2^{-1}$
\item $\varphi_{115} :\equiv N \in \mathbb{N} \land \forall n \in \mathbb{N}\,(N \leq n \to \lvert b_{n} - \beta \rvert < d_{0}) \to n \in \mathbb{N} \to N \leq n \to \lvert \beta \rvert^{-1} \in \mathbb{R}$
\item $\varphi_{116} :\equiv \lvert (\beta - b_{n}) \cdot (b_{n}^{-1} \cdot \beta^{-1}) \rvert = \lvert \beta - b_{n} \rvert \cdot \lvert b_{n}^{-1} \cdot \beta^{-1} \rvert$
\item $\varphi_{117} :\equiv n \in \mathbb{N} \to N \leq n \to \lvert b_{n}^{-1} - \beta^{-1} \rvert = \lvert \beta - b_{n} \rvert \cdot \lvert b_{n}^{-1} \cdot \beta^{-1} \rvert$
\item $\varphi_{118} :\equiv \forall a\,(\mathrm{Seq}(a) \to \forall n \in \mathbb{N}\,(a_{n} \neq 0) \to \mathrm{Seq}(a^{\ast}) \land \forall n \in \mathbb{N}\,(a^{\ast}(n) = a_{n}^{-1}))$
\item $\varphi_{119} :\equiv N \in \mathbb{N} \land \forall n \in \mathbb{N}\,(N \leq n \to \lvert b_{n} - \beta \rvert < d_{0}) \to n \in \mathbb{N} \to N \leq n \to \lvert b_{n} \rvert \in \mathbb{R}$
\item $\varphi_{120} :\equiv \lvert \beta \rvert \cdot 2^{-1} \cdot \lvert \beta \rvert \neq 0 \to \lvert \beta \rvert \cdot 2^{-1} \cdot \lvert \beta \rvert \cdot (\lvert \beta \rvert \cdot 2^{-1} \cdot \lvert \beta \rvert)^{-1} = 1$
\item $\varphi_{121} :\equiv N \in \mathbb{N} \land \forall n \in \mathbb{N}\,(N \leq n \to \lvert b_{n} - \beta \rvert < d_{0}) \to n \in \mathbb{N} \to N \leq n \to b_{n}^{-1} \in \mathbb{R}$
\item $\varphi_{122} :\equiv \mathrm{Conv}(b,\beta) \land \beta \neq 0 \land \forall n \in \mathbb{N}\,(b_{n} \neq 0) \to \forall \varepsilon\,(\mathrm{Conv}(b,\beta) \land \beta \neq 0 \land \forall n \in \mathbb{N}\,(b_{n} \neq 0))$
\item $\varphi_{123} :\equiv N \in \mathbb{N} \land \forall n \in \mathbb{N}\,(N \leq n \to \lvert b_{n} - \beta \rvert < d_{0}) \to n \in \mathbb{N} \to N \leq n \to 0 < \lvert b_{n} \rvert^{-1}$
\item $\varphi_{124} :\equiv N \in \mathbb{N} \land \forall n \in \mathbb{N}\,(N \leq n \to \lvert b_{n} - \beta \rvert < d_{0}) \to n \in \mathbb{N} \to N \leq n \to b_{n} - \beta \in \mathbb{R}$
\item $\varphi_{125} :\equiv N \in \mathbb{N} \land \forall n \in \mathbb{N}\,(N \leq n \to \lvert b_{n} - \beta \rvert < d_{0}) \to n \in \mathbb{N} \to N \leq n \to \beta - b_{n} \in \mathbb{R}$
\item $\varphi_{126} :\equiv (\mathrm{Conv}(b,\beta) \to \mathrm{Seq}(b)) \to (\mathrm{Conv}(b,\beta) \to \mathrm{Conv}(b,\beta)) \to \varphi_{5}$
\item $\varphi_{127} :\equiv \forall y\,(b^{\ast}(n) = y \to y = b^{\ast}(n)) \to b^{\ast}(n) = b_{n}^{-1} \to b_{n}^{-1} = b^{\ast}(n)$
\item $\varphi_{128} :\equiv N \in \mathbb{N} \land \forall n \in \mathbb{N}\,(N \leq n \to \lvert b_{n} - \beta \rvert < d_{0}) \to n \in \mathbb{N} \to N \leq n \to \lvert b_{n} - \beta \rvert < d_{0}$
\item $\varphi_{129} :\equiv \lvert \beta - b_{n} \rvert \cdot (\lvert b_{n} \rvert^{-1} \cdot \lvert \beta^{-1} \rvert) = \lvert \beta - b_{n} \rvert \cdot (\lvert b_{n} \rvert^{-1} \cdot w)$
\item $\varphi_{130} :\equiv \lvert \beta - b_{n} \rvert \cdot (\lvert b_{n} \rvert^{-1} \cdot \lvert \beta^{-1} \rvert) = \lvert \beta - b_{n} \rvert \cdot (\lvert b_{n} \rvert^{-1} \cdot z)$
\item $\varphi_{131} :\equiv N \in \mathbb{N} \land \forall n \in \mathbb{N}\,(N \leq n \to \lvert b_{n} - \beta \rvert < d_{0}) \to n \in \mathbb{N} \to N \leq n \to \lvert \beta^{-1} \rvert = \lvert \beta \rvert^{-1}$
\item $\varphi_{132} :\equiv (\lvert \beta \rvert \cdot 2^{-1})^{-1} \in \mathbb{R} \to \lvert \beta \rvert^{-1} \in \mathbb{R} \to (\lvert \beta \rvert \cdot 2^{-1})^{-1} \cdot \lvert \beta \rvert^{-1} \in \mathbb{R}$
\item $\varphi_{133} :\equiv \lvert \beta - b_{n} \rvert \cdot \lvert b_{n}^{-1} \cdot \beta^{-1} \rvert = \lvert \beta - b_{n} \rvert \cdot \lvert b_{n}^{-1} \cdot \beta^{-1} \rvert$
\item $\varphi_{134} :\equiv N \in \mathbb{N} \land \forall n \in \mathbb{N}\,(N \leq n \to \lvert b_{n} - \beta \rvert < d_{0}) \to n \in \mathbb{N} \to N \leq n \to \lvert \beta \rvert \cdot 2^{-1} \in \mathbb{R}$
\item $\varphi_{135} :\equiv n \in \mathbb{N} \to N \leq n \to \lvert b_{n}^{-1} - \beta^{-1} \rvert = \lvert \beta - b_{n} \rvert \cdot (\lvert b_{n} \rvert^{-1} \cdot \lvert \beta^{-1} \rvert)$
\item $\varphi_{136} :\equiv n \in \mathbb{N} \to N \leq n \to \lvert b_{n}^{-1} - \beta^{-1} \rvert = \lvert \beta - b_{n} \rvert \cdot (\lvert b_{n} \rvert^{-1} \cdot \lvert \beta \rvert^{-1})$
\item $\varphi_{137} :\equiv n \in \mathbb{N} \to N \leq n \to \lvert b_{n}^{-1} - \beta^{-1} \rvert = \lvert b_{n} - \beta \rvert \cdot (\lvert b_{n} \rvert^{-1} \cdot \lvert \beta \rvert^{-1})$
\item $\varphi_{138} :\equiv n \in \mathbb{N} \to N \leq n \to \lvert b_{n} - \beta \rvert \cdot (\lvert b_{n} \rvert^{-1} \cdot \lvert \beta \rvert^{-1}) = \lvert b_{n}^{-1} - \beta^{-1} \rvert$
\item $\varphi_{139} :\equiv \forall x\,\forall y\,(x \in \mathbb{R} \to y \in \mathbb{R} \to x \cdot y \in \mathbb{R}) \to \forall y\,(\lvert \beta \rvert \in \mathbb{R} \to y \in \mathbb{R} \to \lvert \beta \rvert \cdot y \in \mathbb{R})$
\item $\varphi_{140} :\equiv \forall y\,(\lvert \beta \rvert \cdot 2^{-1} \in \mathbb{R} \to y \in \mathbb{R} \to 0 < \lvert \beta \rvert \cdot 2^{-1} \to 0 < y \to 0 < \lvert \beta \rvert \cdot 2^{-1} \cdot y)$
\item $\varphi_{141} :\equiv \forall y\,(\beta \in \mathbb{R} \to y \in \mathbb{R} \to \beta - y \in \mathbb{R}) \to \beta \in \mathbb{R} \to b_{n} \in \mathbb{R} \to \beta - b_{n} \in \mathbb{R}$
\item $\varphi_{142} :\equiv \forall \varepsilon\,(\mathrm{Conv}(b,\beta) \to \varepsilon \in \mathbb{R} \to 0 < \varepsilon \to \exists N \in \mathbb{N}\,\forall n \in \mathbb{N}\,(N \leq n \to \lvert b_{n} - \beta \rvert < \varepsilon))$
\item $\varphi_{143} :\equiv 0 < \lvert \beta \rvert \cdot 2^{-1} \to \lvert \beta \rvert \cdot 2^{-1} < \lvert b_{n} \rvert \to \lvert b_{n} \rvert^{-1} < (\lvert \beta \rvert \cdot 2^{-1})^{-1}$
\item $\varphi_{144} :\equiv N \in \mathbb{N} \land \forall n \in \mathbb{N}\,(N \leq n \to \lvert b_{n} - \beta \rvert < d_{0}) \to n \in \mathbb{N} \to N \leq n \to \lvert b_{n} \rvert^{-1} \in \mathbb{R}$
\item $\varphi_{145} :\equiv N \in \mathbb{N} \land \forall n \in \mathbb{N}\,(N \leq n \to \lvert b_{n} - \beta \rvert < d_{0}) \to n \in \mathbb{N} \to N \leq n \to \lvert b_{n} - \beta \rvert \in \mathbb{R}$
\item $\varphi_{146} :\equiv \forall x \in \mathbb{R}\,(0 < x \to 0 < x^{-1}) \to \lvert b_{n} \rvert \in \mathbb{R} \to 0 < \lvert b_{n} \rvert \to 0 < \lvert b_{n} \rvert^{-1}$
\item $\varphi_{147} :\equiv \forall x\,\forall y\,(x \in \mathbb{R} \to y \in \mathbb{R} \to x - y \in \mathbb{R}) \to \forall y\,(b_{n} \in \mathbb{R} \to y \in \mathbb{R} \to b_{n} - y \in \mathbb{R})$
\item $\varphi_{148} :\equiv b_{n}^{-1} \in \mathbb{R} \to \beta^{-1} \in \mathbb{R} \to \lvert b_{n}^{-1} \cdot \beta^{-1} \rvert = \lvert b_{n}^{-1} \rvert \cdot \lvert \beta^{-1} \rvert$
\item $\varphi_{149} :\equiv N \in \mathbb{N} \land \forall n \in \mathbb{N}\,(N \leq n \to \lvert b_{n} - \beta \rvert < d_{0}) \to n \in \mathbb{N} \to N \leq n \to b^{\ast}(n) = b_{n}^{-1}$
\item $\varphi_{150} :\equiv N \in \mathbb{N} \land \forall n \in \mathbb{N}\,(N \leq n \to \lvert b_{n} - \beta \rvert < d_{0}) \to n \in \mathbb{N} \to N \leq n \to b_{n}^{-1} = b^{\ast}(n)$
\item $\varphi_{151} :\equiv b_{n}^{-1} = b^{\ast}(n) \to \lvert b_{n}^{-1} - \beta^{-1} \rvert < \varepsilon \to \lvert b^{\ast}(n) - \beta^{-1} \rvert < \varepsilon$
\item $\varphi_{152} :\equiv \forall y\,(\lvert b_{n} \rvert^{-1} \in \mathbb{R} \to y \in \mathbb{R} \to 0 < \lvert b_{n} \rvert^{-1} \to 0 < y \to 0 < \lvert b_{n} \rvert^{-1} \cdot y)$
\item $\varphi_{153} :\equiv \forall \alpha\,\forall \varepsilon\,(\mathrm{Conv}(b,\alpha) \to \varepsilon \in \mathbb{R} \to 0 < \varepsilon \to \exists N \in \mathbb{N}\,\forall n \in \mathbb{N}\,(N \leq n \to \lvert b_{n} - \alpha \rvert < \varepsilon))$
\item $\varphi_{154} :\equiv N \in \mathbb{N} \land \forall n \in \mathbb{N}\,(N \leq n \to \lvert b_{n} - \beta \rvert < d_{0}) \to n \in \mathbb{N} \to N \leq n \to (\lvert \beta \rvert \cdot 2^{-1})^{-1} \in \mathbb{R}$
\item $\varphi_{155} :\equiv N \in \mathbb{N} \land \forall n \in \mathbb{N}\,(N \leq n \to \lvert b_{n} - \beta \rvert < d_{0}) \to n \in \mathbb{N} \to N \leq n \to 0 \neq \lvert \beta \rvert \cdot 2^{-1} \cdot \lvert \beta \rvert$
\item $\varphi_{156} :\equiv N \in \mathbb{N} \land \forall n \in \mathbb{N}\,(N \leq n \to \lvert b_{n} - \beta \rvert < d_{0}) \to n \in \mathbb{N} \to N \leq n \to \lvert \beta \rvert \cdot 2^{-1} \cdot \lvert \beta \rvert \neq 0$
\item $\varphi_{157} :\equiv \lvert \beta - b_{n} \rvert \cdot \lvert b_{n}^{-1} \cdot \beta^{-1} \rvert = \lvert \beta - b_{n} \rvert \cdot (\lvert b_{n}^{-1} \rvert \cdot \lvert \beta^{-1} \rvert)$
\item $\varphi_{158} :\equiv \forall x \in \mathbb{R}\,(x \neq 0 \to x^{-1} \in \mathbb{R}) \to \lvert b_{n} \rvert \in \mathbb{R} \to \lvert b_{n} \rvert \neq 0 \to \lvert b_{n} \rvert^{-1} \in \mathbb{R}$
\item $\varphi_{159} :\equiv N \in \mathbb{N} \land \forall n \in \mathbb{N}\,(N \leq n \to \lvert b_{n} - \beta \rvert < d_{0}) \to n \in \mathbb{N} \to N \leq n \to b_{n}^{-1} \cdot \beta^{-1} \in \mathbb{R}$
\item $\varphi_{160} :\equiv \forall x\,\forall y\,(x \in \mathbb{R} \to y \in \mathbb{R} \to \lvert x - y \rvert = \lvert y - x \rvert) \to \forall y\,(\beta \in \mathbb{R} \to y \in \mathbb{R} \to \lvert \beta - y \rvert = \lvert y - \beta \rvert)$
\item $\varphi_{161} :\equiv \forall a\,\forall \alpha\,\forall \varepsilon\,(\mathrm{Conv}(a,\alpha) \to \varepsilon \in \mathbb{R} \to 0 < \varepsilon \to \exists N \in \mathbb{N}\,\forall n \in \mathbb{N}\,(N \leq n \to \lvert a_{n} - \alpha \rvert < \varepsilon))$
\item $\varphi_{162} :\equiv N \in \mathbb{N} \land \forall n \in \mathbb{N}\,(N \leq n \to \lvert b_{n} - \beta \rvert < d_{0}) \to n \in \mathbb{N} \to N \leq n \to \lvert \beta \rvert \cdot 2^{-1} < \lvert b_{n} \rvert$
\item $\varphi_{163} :\equiv \lvert \beta - b_{n} \rvert \cdot (\lvert b_{n} \rvert^{-1} \cdot \lvert \beta \rvert^{-1}) = \lvert \beta - b_{n} \rvert \cdot (\lvert b_{n} \rvert^{-1} \cdot \lvert \beta \rvert^{-1})$
\item $\varphi_{164} :\equiv \lvert \beta - b_{n} \rvert \cdot (\lvert b_{n} \rvert^{-1} \cdot \lvert \beta^{-1} \rvert) = \lvert \beta - b_{n} \rvert \cdot (\lvert b_{n} \rvert^{-1} \cdot \lvert \beta^{-1} \rvert)$
\item $\varphi_{165} :\equiv \lvert \beta - b_{n} \rvert \cdot (\lvert b_{n}^{-1} \rvert \cdot \lvert \beta^{-1} \rvert) = \lvert \beta - b_{n} \rvert \cdot (\lvert b_{n}^{-1} \rvert \cdot \lvert \beta^{-1} \rvert)$
\item $\varphi_{166} :\equiv \lvert \beta - b_{n} \rvert \cdot (\lvert b_{n} \rvert^{-1} \cdot \lvert \beta \rvert^{-1}) = \lvert b_{n} - \beta \rvert \cdot (\lvert b_{n} \rvert^{-1} \cdot \lvert \beta \rvert^{-1})$
\item $\varphi_{167} :\equiv \lvert \beta - b_{n} \rvert \cdot (\lvert b_{n} \rvert^{-1} \cdot \lvert \beta^{-1} \rvert) = \lvert \beta - b_{n} \rvert \cdot (\lvert b_{n} \rvert^{-1} \cdot \lvert \beta \rvert^{-1})$
\item $\varphi_{168} :\equiv \lvert \beta - b_{n} \rvert \cdot (\lvert b_{n}^{-1} \rvert \cdot \lvert \beta^{-1} \rvert) = \lvert \beta - b_{n} \rvert \cdot (\lvert b_{n} \rvert^{-1} \cdot \lvert \beta^{-1} \rvert)$
\item $\varphi_{169} :\equiv N \in \mathbb{N} \land \forall n \in \mathbb{N}\,(N \leq n \to \lvert b_{n} - \beta \rvert < d_{0}) \to n \in \mathbb{N} \to N \leq n \to \lvert b_{n}^{-1} - \beta^{-1} \rvert < \varepsilon$
\item $\varphi_{170} :\equiv N \in \mathbb{N} \land \forall n \in \mathbb{N}\,(N \leq n \to \lvert b_{n} - \beta \rvert < d_{0}) \to n \in \mathbb{N} \to N \leq n \to \lvert \beta \rvert \cdot 2^{-1} \cdot \lvert \beta \rvert \in \mathbb{R}$
\item $\varphi_{171} :\equiv N \in \mathbb{N} \land \forall n \in \mathbb{N}\,(N \leq n \to \lvert b_{n} - \beta \rvert < d_{0}) \to n \in \mathbb{N} \to N \leq n \to \lvert b_{n}^{-1} \rvert = \lvert b_{n} \rvert^{-1}$
\item $\varphi_{172} :\equiv N \in \mathbb{N} \land \forall n \in \mathbb{N}\,(N \leq n \to \lvert b_{n} - \beta \rvert < d_{0}) \to n \in \mathbb{N} \to N \leq n \to d_{0} \leq \varepsilon \cdot (\lvert \beta \rvert \cdot 2^{-1} \cdot \lvert \beta \rvert)$
\item $\varphi_{173} :\equiv N \in \mathbb{N} \land \forall n \in \mathbb{N}\,(N \leq n \to \lvert b_{n} - \beta \rvert < d_{0}) \to n \in \mathbb{N} \to N \leq n \to \lvert b^{\ast}(n) - \beta^{-1} \rvert < \varepsilon$
\item $\varphi_{174} :\equiv \forall y\,(b_{n} \in \mathbb{R} \to y \in \mathbb{R} \to b_{n} - y \in \mathbb{R}) \to b_{n} \in \mathbb{R} \to \beta \in \mathbb{R} \to b_{n} - \beta \in \mathbb{R}$
\item $\varphi_{175} :\equiv \forall x\,\forall y\,(x \in \mathbb{R} \to y \in \mathbb{R} \to x \cdot y \in \mathbb{R}) \to \forall y\,(b_{n}^{-1} \in \mathbb{R} \to y \in \mathbb{R} \to b_{n}^{-1} \cdot y \in \mathbb{R})$
\item $\varphi_{176} :\equiv \varepsilon \cdot (\lvert \beta \rvert \cdot 2^{-1} \cdot \lvert \beta \rvert) \cdot ((\lvert \beta \rvert \cdot 2^{-1})^{-1} \cdot \lvert \beta \rvert^{-1}) = \varepsilon \cdot (\lvert \beta \rvert \cdot 2^{-1} \cdot \lvert \beta \rvert) \cdot w$
\item $\varphi_{177} :\equiv \varepsilon \cdot (\lvert \beta \rvert \cdot 2^{-1} \cdot \lvert \beta \rvert) \cdot ((\lvert \beta \rvert \cdot 2^{-1})^{-1} \cdot \lvert \beta \rvert^{-1}) = \varepsilon \cdot (\lvert \beta \rvert \cdot 2^{-1} \cdot \lvert \beta \rvert) \cdot z$
\item $\varphi_{178} :\equiv N \in \mathbb{N} \land \forall n \in \mathbb{N}\,(N \leq n \to \lvert b_{n} - \beta \rvert < d_{0}) \to n \in \mathbb{N} \to N \leq n \to 0 < \lvert b_{n} \rvert^{-1} \cdot \lvert \beta \rvert^{-1}$
\item $\varphi_{179} :\equiv N \in \mathbb{N} \land \forall n \in \mathbb{N}\,(N \leq n \to \lvert b_{n} - \beta \rvert < d_{0}) \to n \in \mathbb{N} \to N \leq n \to \lvert \beta - b_{n} \rvert = \lvert b_{n} - \beta \rvert$
\item $\varphi_{180} :\equiv N \in \mathbb{N} \land \forall n \in \mathbb{N}\,(N \leq n \to \lvert b_{n} - \beta \rvert < d_{0}) \to N \in \mathbb{N} \land \forall n \in \mathbb{N}\,(N \leq n \to \lvert b_{n} - \beta \rvert < d_{0})$
\item $\varphi_{181} :\equiv N \in \mathbb{N} \land \forall n \in \mathbb{N}\,(N \leq n \to \lvert b_{n} - \beta \rvert < d_{0}) \to n \in \mathbb{N} \to N \leq n \to \lvert b_{n} - \beta \rvert < \lvert \beta \rvert \cdot 2^{-1}$
\item $\varphi_{182} :\equiv \forall x \in \mathbb{R}\,(x \neq 0 \to \lvert x^{-1} \rvert = \lvert x \rvert^{-1}) \to b_{n} \in \mathbb{R} \to b_{n} \neq 0 \to \lvert b_{n}^{-1} \rvert = \lvert b_{n} \rvert^{-1}$
\item $\varphi_{183} :\equiv N \in \mathbb{N} \land \forall n \in \mathbb{N}\,(N \leq n \to \lvert b_{n} - \beta \rvert < d_{0}) \to n \in \mathbb{N} \to N \leq n \to n \in \mathbb{N} \to b^{\ast}(n) = b_{n}^{-1}$
\item $\varphi_{184} :\equiv N \in \mathbb{N} \land \forall n \in \mathbb{N}\,(N \leq n \to \lvert b_{n} - \beta \rvert < d_{0}) \to \forall n \in \mathbb{N}\,(N \leq n \to \lvert b^{\ast}(n) - \beta^{-1} \rvert < \varepsilon)$
\item $\varphi_{185} :\equiv \forall y\,(b_{n} \in \mathbb{R} \to y \in \mathbb{R} \to b_{n} \neq 0 \to y \neq 0 \to b_{n}^{-1} - y^{-1} = (y - b_{n}) \cdot (b_{n}^{-1} \cdot y^{-1}))$
\item $\varphi_{186} :\equiv N \in \mathbb{N} \land \forall n \in \mathbb{N}\,(N \leq n \to \lvert b_{n} - \beta \rvert < d_{0}) \to n \in \mathbb{N} \to N \leq n \to (\lvert \beta \rvert \cdot 2^{-1} \cdot \lvert \beta \rvert)^{-1} \in \mathbb{R}$
\item $\varphi_{187} :\equiv \forall y\,(\lvert \beta \rvert \in \mathbb{R} \to y \in \mathbb{R} \to \lvert \beta \rvert \cdot y \in \mathbb{R}) \to \lvert \beta \rvert \in \mathbb{R} \to 2^{-1} \in \mathbb{R} \to \lvert \beta \rvert \cdot 2^{-1} \in \mathbb{R}$
\item $\varphi_{188} :\equiv N \in \mathbb{N} \land \forall n \in \mathbb{N}\,(N \leq n \to \lvert b_{n} - \beta \rvert < d_{0}) \to n \in \mathbb{N} \to N \leq n \to n \in \mathbb{N} \to N \leq n \to \lvert b_{n} - \beta \rvert < d_{0}$
\item $\varphi_{189} :\equiv N \in \mathbb{N} \land \forall n \in \mathbb{N}\,(N \leq n \to \lvert b_{n} - \beta \rvert < d_{0}) \to n \in \mathbb{N} \to N \leq n \to \lvert b_{n} \rvert^{-1} \cdot \lvert \beta \rvert^{-1} \in \mathbb{R}$
\item $\varphi_{190} :\equiv N \in \mathbb{N} \land \forall n \in \mathbb{N}\,(N \leq n \to \lvert b_{n} - \beta \rvert < d_{0}) \to n \in \mathbb{N} \to N \leq n \to \lvert b_{n} \rvert^{-1} < (\lvert \beta \rvert \cdot 2^{-1})^{-1}$
\item $\varphi_{191} :\equiv \lvert b_{n} \rvert^{-1} < (\lvert \beta \rvert \cdot 2^{-1})^{-1} \to \lvert b_{n} \rvert^{-1} \cdot \lvert \beta \rvert^{-1} < (\lvert \beta \rvert \cdot 2^{-1})^{-1} \cdot \lvert \beta \rvert^{-1}$
\item $\varphi_{192} :\equiv N \in \mathbb{N} \land \forall n \in \mathbb{N}\,(N \leq n \to \lvert b_{n} - \beta \rvert < d_{0}) \to n \in \mathbb{N} \to N \leq n \to \varepsilon \cdot (\lvert \beta \rvert \cdot 2^{-1} \cdot \lvert \beta \rvert) \in \mathbb{R}$
\item $\varphi_{193} :\equiv N \in \mathbb{N} \land \forall n \in \mathbb{N}\,(N \leq n \to \lvert b_{n} - \beta \rvert < d_{0}) \to n \in \mathbb{N} \to N \leq n \to \forall n \in \mathbb{N}\,(b^{\ast}(n) = b_{n}^{-1})$
\item $\varphi_{194} :\equiv N \in \mathbb{N} \land \forall n \in \mathbb{N}\,(N \leq n \to \lvert b_{n} - \beta \rvert < d_{0}) \to n \in \mathbb{N} \to N \leq n \to \forall n \in \mathbb{N}\,(N \leq n \to \lvert b_{n} - \beta \rvert < d_{0})$
\item $\varphi_{195} :\equiv N \in \mathbb{N} \land \forall n \in \mathbb{N}\,(N \leq n \to \lvert b_{n} - \beta \rvert < d_{0}) \to n \in \mathbb{N} \to N \leq n \to (\lvert \beta \rvert \cdot 2^{-1})^{-1} \cdot \lvert \beta \rvert^{-1} \in \mathbb{R}$
\item $\varphi_{196} :\equiv \forall y\,(\lvert \beta \rvert \cdot 2^{-1} \in \mathbb{R} \to y \in \mathbb{R} \to 0 < \lvert \beta \rvert \cdot 2^{-1} \to \lvert \beta \rvert \cdot 2^{-1} < y \to y^{-1} < (\lvert \beta \rvert \cdot 2^{-1})^{-1})$
\item $\varphi_{197} :\equiv \forall y\,\forall z\,(\lvert b_{n} - \beta \rvert < y \to y \leq z \to \lvert b_{n} - \beta \rvert < z) \to \forall z\,(\lvert b_{n} - \beta \rvert < d_{0} \to d_{0} \leq z \to \lvert b_{n} - \beta \rvert < z)$
\item $\varphi_{198} :\equiv \forall x\,\forall y\,(x \in \mathbb{R} \to y \in \mathbb{R} \to 0 < x \to 0 < y \to 0 < x \cdot y) \to \forall y\,(\varepsilon \in \mathbb{R} \to y \in \mathbb{R} \to 0 < \varepsilon \to 0 < y \to 0 < \varepsilon \cdot y)$
\item $\varphi_{199} :\equiv \forall x\,\forall y\,(x \in \mathbb{R} \to y \in \mathbb{R} \to x \cdot y \in \mathbb{R}) \to \forall y\,(\lvert b_{n} \rvert^{-1} \in \mathbb{R} \to y \in \mathbb{R} \to \lvert b_{n} \rvert^{-1} \cdot y \in \mathbb{R})$
\item $\varphi_{200} :\equiv \forall q\,(\lvert \beta \rvert \cdot 2^{-1} \in \mathbb{R} \to q \in \mathbb{R} \to 0 < \lvert \beta \rvert \cdot 2^{-1} \to 0 < q \to \exists d\,(d \in \mathbb{R} \land 0 < d \land d \leq \lvert \beta \rvert \cdot 2^{-1} \land d \leq q))$
\item $\varphi_{201} :\equiv (\mathrm{Conv}(b,\beta) \to \beta \neq 0 \to \forall n \in \mathbb{N}\,(b_{n} \neq 0) \to \beta \in \mathbb{R}) \to \varphi_{15}$
\item $\varphi_{202} :\equiv \forall y\,\forall z\,(\lvert b_{n} \rvert^{-1} \in \mathbb{R} \to y \in \mathbb{R} \to z \in \mathbb{R} \to 0 < z \to \lvert b_{n} \rvert^{-1} < y \to \lvert b_{n} \rvert^{-1} \cdot z < y \cdot z)$
\item $\varphi_{203} :\equiv \forall w\,(b_{n}^{-1} - \beta^{-1} = w \to \lvert b_{n}^{-1} - \beta^{-1} \rvert = \lvert b_{n}^{-1} - \beta^{-1} \rvert \to \lvert b_{n}^{-1} - \beta^{-1} \rvert = \lvert w \rvert)$
\item $\varphi_{204} :\equiv \forall y\,(\beta \in \mathbb{R} \to y \in \mathbb{R} \to \lvert \beta - y \rvert = \lvert y - \beta \rvert) \to \beta \in \mathbb{R} \to b_{n} \in \mathbb{R} \to \lvert \beta - b_{n} \rvert = \lvert b_{n} - \beta \rvert$
\item $\varphi_{205} :\equiv N \in \mathbb{N} \land \forall n \in \mathbb{N}\,(N \leq n \to \lvert b_{n} - \beta \rvert < d_{0}) \to N \in \mathbb{N} \land \forall n \in \mathbb{N}\,(N \leq n \to \lvert b^{\ast}(n) - \beta^{-1} \rvert < \varepsilon)$
\item $\varphi_{206} :\equiv \forall y\,\forall z\,(\lvert b_{n} - \beta \rvert \in \mathbb{R} \to y \in \mathbb{R} \to z \in \mathbb{R} \to 0 < z \to \lvert b_{n} - \beta \rvert < y \to \lvert b_{n} - \beta \rvert \cdot z < y \cdot z)$
\item $\varphi_{207} :\equiv \lvert b_{n} - \beta \rvert \cdot (\lvert b_{n} \rvert^{-1} \cdot \lvert \beta \rvert^{-1}) < \varepsilon \cdot (\lvert \beta \rvert \cdot 2^{-1} \cdot \lvert \beta \rvert) \cdot (\lvert b_{n} \rvert^{-1} \cdot \lvert \beta \rvert^{-1})$
\item $\varphi_{208} :\equiv \forall y\,(\lvert \beta \rvert \cdot 2^{-1} \in \mathbb{R} \to y \in \mathbb{R} \to \lvert \beta \rvert \cdot 2^{-1} \neq 0 \to y \neq 0 \to (\lvert \beta \rvert \cdot 2^{-1} \cdot y)^{-1} = (\lvert \beta \rvert \cdot 2^{-1})^{-1} \cdot y^{-1})$
\item $\varphi_{209} :\equiv N \in \mathbb{N} \land \forall n \in \mathbb{N}\,(N \leq n \to \lvert b_{n} - \beta \rvert < d_{0}) \to \exists N \in \mathbb{N}\,\forall n \in \mathbb{N}\,(N \leq n \to \lvert b^{\ast}(n) - \beta^{-1} \rvert < \varepsilon)$
\item $\varphi_{210} :\equiv \forall a\,\forall \alpha\,(\mathrm{Conv}(a,\alpha) \leftrightarrow \mathrm{Seq}(a) \land \alpha \in \mathbb{R} \land \forall \varepsilon \in \mathbb{R}\,(0 < \varepsilon \to \exists N \in \mathbb{N}\,\forall n \in \mathbb{N}\,(N \leq n \to \lvert a_{n} - \alpha \rvert < \varepsilon)))$
\item $\varphi_{211} :\equiv \lvert b_{n} \rvert^{-1} \in \mathbb{R} \to \lvert \beta \rvert^{-1} \in \mathbb{R} \to 0 < \lvert b_{n} \rvert^{-1} \to 0 < \lvert \beta \rvert^{-1} \to 0 < \lvert b_{n} \rvert^{-1} \cdot \lvert \beta \rvert^{-1}$
\item $\varphi_{212} :\equiv \lvert b_{n}^{-1} - \beta^{-1} \rvert = \lvert b_{n}^{-1} - \beta^{-1} \rvert \to \lvert b_{n}^{-1} - \beta^{-1} \rvert = \lvert (\beta - b_{n}) \cdot (b_{n}^{-1} \cdot \beta^{-1}) \rvert$
\item $\varphi_{213} :\equiv \forall w\,(\lvert b_{n} - \beta \rvert \cdot (\lvert b_{n} \rvert^{-1} \cdot \lvert \beta \rvert^{-1}) = w \to \lvert b_{n} - \beta \rvert \cdot (\lvert b_{n} \rvert^{-1} \cdot \lvert \beta \rvert^{-1}) < \varepsilon \to w < \varepsilon)$
\item $\varphi_{214} :\equiv \forall y\,\forall z\,(\lvert b_{n} - \beta \rvert \cdot (\lvert b_{n} \rvert^{-1} \cdot \lvert \beta \rvert^{-1}) < y \to y < z \to \lvert b_{n} - \beta \rvert \cdot (\lvert b_{n} \rvert^{-1} \cdot \lvert \beta \rvert^{-1}) < z)$
\item $\varphi_{215} :\equiv N \in \mathbb{N} \land \forall n \in \mathbb{N}\,(N \leq n \to \lvert b_{n} - \beta \rvert < d_{0}) \to n \in \mathbb{N} \to N \leq n \to \lvert b_{n} - \beta \rvert \cdot (\lvert b_{n} \rvert^{-1} \cdot \lvert \beta \rvert^{-1}) < \varepsilon$
\item $\varphi_{216} :\equiv N \in \mathbb{N} \land \forall n \in \mathbb{N}\,(N \leq n \to \lvert b_{n} - \beta \rvert < d_{0}) \to n \in \mathbb{N} \to N \leq n \to \lvert b_{n}^{-1} \cdot \beta^{-1} \rvert = \lvert b_{n}^{-1} \rvert \cdot \lvert \beta^{-1} \rvert$
\item $\varphi_{217} :\equiv \varepsilon \cdot (\lvert \beta \rvert \cdot 2^{-1} \cdot \lvert \beta \rvert \cdot (\lvert \beta \rvert \cdot 2^{-1} \cdot \lvert \beta \rvert)^{-1}) = \varepsilon \cdot (\lvert \beta \rvert \cdot 2^{-1} \cdot \lvert \beta \rvert \cdot (\lvert \beta \rvert \cdot 2^{-1} \cdot \lvert \beta \rvert)^{-1})$
\item $\varphi_{218} :\equiv \varepsilon \cdot (\lvert \beta \rvert \cdot 2^{-1} \cdot \lvert \beta \rvert) \cdot (\lvert \beta \rvert \cdot 2^{-1} \cdot \lvert \beta \rvert)^{-1} = \varepsilon \cdot (\lvert \beta \rvert \cdot 2^{-1} \cdot \lvert \beta \rvert \cdot (\lvert \beta \rvert \cdot 2^{-1} \cdot \lvert \beta \rvert)^{-1})$
\item $\varphi_{219} :\equiv N \in \mathbb{N} \land \forall n \in \mathbb{N}\,(N \leq n \to \lvert b_{n} - \beta \rvert < d_{0}) \to n \in \mathbb{N} \to \varphi_{29}$
\item $\varphi_{220} :\equiv \varepsilon \cdot (\lvert \beta \rvert \cdot 2^{-1} \cdot \lvert \beta \rvert) \cdot (\lvert b_{n} \rvert^{-1} \cdot \lvert \beta \rvert^{-1}) < \varepsilon \to \lvert b_{n} - \beta \rvert \cdot (\lvert b_{n} \rvert^{-1} \cdot \lvert \beta \rvert^{-1}) < \varepsilon$
\item $\varphi_{221} :\equiv \forall y\,(b_{n}^{-1} \in \mathbb{R} \to y \in \mathbb{R} \to b_{n}^{-1} \cdot y \in \mathbb{R}) \to b_{n}^{-1} \in \mathbb{R} \to \beta^{-1} \in \mathbb{R} \to b_{n}^{-1} \cdot \beta^{-1} \in \mathbb{R}$
\item $\varphi_{222} :\equiv \forall z\,\forall w\,(z = w \to \lvert z - \beta^{-1} \rvert < \varepsilon \to \lvert w - \beta^{-1} \rvert < \varepsilon) \to \forall w\,(b_{n}^{-1} = w \to \lvert b_{n}^{-1} - \beta^{-1} \rvert < \varepsilon \to \lvert w - \beta^{-1} \rvert < \varepsilon)$
\item $\varphi_{223} :\equiv \mathrm{Conv}(b,\beta) \to \beta \neq 0 \to \forall n \in \mathbb{N}\,(b_{n} \neq 0) \to \varepsilon \in \mathbb{R} \to \varphi_{65}$
\item $\varphi_{224} :\equiv \mathrm{Conv}(b,\beta) \land \beta \neq 0 \land \forall n \in \mathbb{N}\,(b_{n} \neq 0) \to \varepsilon \in \mathbb{R} \to \varphi_{65}$
\item $\varphi_{225} :\equiv \varepsilon \cdot (\lvert \beta \rvert \cdot 2^{-1} \cdot \lvert \beta \rvert) \cdot ((\lvert \beta \rvert \cdot 2^{-1})^{-1} \cdot \lvert \beta \rvert^{-1}) = \varepsilon \cdot (\lvert \beta \rvert \cdot 2^{-1} \cdot \lvert \beta \rvert \cdot (\lvert \beta \rvert \cdot 2^{-1} \cdot \lvert \beta \rvert)^{-1})$
\item $\varphi_{226} :\equiv \varepsilon \cdot (\lvert \beta \rvert \cdot 2^{-1} \cdot \lvert \beta \rvert) \cdot ((\lvert \beta \rvert \cdot 2^{-1})^{-1} \cdot \lvert \beta \rvert^{-1}) = \varepsilon \cdot (\lvert \beta \rvert \cdot 2^{-1} \cdot \lvert \beta \rvert) \cdot (\lvert \beta \rvert \cdot 2^{-1} \cdot \lvert \beta \rvert)^{-1}$
\item $\varphi_{227} :\equiv \mathrm{Conv}(b,\beta) \land \beta \neq 0 \land \forall n \in \mathbb{N}\,(b_{n} \neq 0) \land \varepsilon \in \mathbb{R} \land 0 < \varepsilon \land \neg \exists N \in \mathbb{N}\,\forall n \in \mathbb{N}\,(N \leq n \to \lvert b^{\ast}(n) - \beta^{-1} \rvert < \varepsilon)$
\item $\varphi_{228} :\equiv N \in \mathbb{N} \land \forall n \in \mathbb{N}\,(N \leq n \to \lvert b_{n} - \beta \rvert < d_{0}) \to n \in \mathbb{N} \to N \leq n \to \varphi_{27}$
\item $\varphi_{229} :\equiv N \in \mathbb{N} \land \forall n \in \mathbb{N}\,(N \leq n \to \lvert b_{n} - \beta \rvert < d_{0}) \to n \in \mathbb{N} \to N \leq n \to \varphi_{26}$
\item $\varphi_{230} :\equiv \forall z\,(\varepsilon \in \mathbb{R} \to \lvert \beta \rvert \cdot 2^{-1} \cdot \lvert \beta \rvert \in \mathbb{R} \to z \in \mathbb{R} \to \varepsilon \cdot (\lvert \beta \rvert \cdot 2^{-1} \cdot \lvert \beta \rvert) \cdot z = \varepsilon \cdot (\lvert \beta \rvert \cdot 2^{-1} \cdot \lvert \beta \rvert \cdot z))$
\item $\varphi_{231} :\equiv \forall \alpha\,(\mathrm{Conv}(b^{\ast},\alpha) \leftrightarrow \mathrm{Seq}(b^{\ast}) \land \alpha \in \mathbb{R} \land \forall \varepsilon \in \mathbb{R}\,(0 < \varepsilon \to \exists N \in \mathbb{N}\,\forall n \in \mathbb{N}\,(N \leq n \to \lvert b^{\ast}(n) - \alpha \rvert < \varepsilon)))$
\item $\varphi_{232} :\equiv \varepsilon \cdot (\lvert \beta \rvert \cdot 2^{-1} \cdot \lvert \beta \rvert) \cdot (\lvert b_{n} \rvert^{-1} \cdot \lvert \beta \rvert^{-1}) < \varepsilon \cdot (\lvert \beta \rvert \cdot 2^{-1} \cdot \lvert \beta \rvert) \cdot ((\lvert \beta \rvert \cdot 2^{-1})^{-1} \cdot \lvert \beta \rvert^{-1})$
\item $\varphi_{233} :\equiv (\mathrm{Conv}(b,\beta) \to \beta \neq 0 \to \forall n \in \mathbb{N}\,(b_{n} \neq 0) \to \beta \in \mathbb{R}) \to (\mathrm{Conv}(b,\beta) \to \beta \neq 0 \to \beta \neq 0) \to \varphi_{10}$
\item $\varphi_{234} :\equiv \varepsilon \cdot (\lvert \beta \rvert \cdot 2^{-1} \cdot \lvert \beta \rvert) \cdot ((\lvert \beta \rvert \cdot 2^{-1})^{-1} \cdot \lvert \beta \rvert^{-1}) = \varepsilon \cdot (\lvert \beta \rvert \cdot 2^{-1} \cdot \lvert \beta \rvert) \cdot ((\lvert \beta \rvert \cdot 2^{-1})^{-1} \cdot \lvert \beta \rvert^{-1})$
\item $\varphi_{235} :\equiv \mathrm{Conv}(b,\beta) \to \beta \neq 0 \to \forall n \in \mathbb{N}\,(b_{n} \neq 0) \to \forall \varepsilon \in \mathbb{R}\,(\varphi_{65})$
\item $\varphi_{236} :\equiv (\mathrm{Conv}(b,\beta) \to \beta \neq 0 \to \forall n \in \mathbb{N}\,(b_{n} \neq 0) \to \beta \in \mathbb{R}) \to \mathrm{Conv}(b,\beta) \to \varphi_{42}$
\item $\varphi_{237} :\equiv \lvert \beta \rvert \cdot 2^{-1} \in \mathbb{R} \to \lvert \beta \rvert \in \mathbb{R} \to \lvert \beta \rvert \cdot 2^{-1} \neq 0 \to \lvert \beta \rvert \neq 0 \to \varphi_{27}$
\item $\varphi_{238} :\equiv \forall x\,\forall y\,(x \in \mathbb{R} \to y \in \mathbb{R} \to \lvert x \cdot y \rvert = \lvert x \rvert \cdot \lvert y \rvert) \to \varphi_{88}$
\item $\varphi_{239} :\equiv \forall x \in \mathbb{R}\,(x \neq 0 \to x^{-1} \in \mathbb{R}) \to \lvert \beta \rvert \cdot 2^{-1} \cdot \lvert \beta \rvert \in \mathbb{R} \to \varphi_{37}$
\item $\varphi_{240} :\equiv \lvert \beta - b_{n} \rvert \cdot \lvert b_{n}^{-1} \cdot \beta^{-1} \rvert = \lvert \beta - b_{n} \rvert \cdot z \to \lvert \beta - b_{n} \rvert \cdot \lvert b_{n}^{-1} \cdot \beta^{-1} \rvert = \lvert \beta - b_{n} \rvert \cdot w$
\item $\varphi_{241} :\equiv \varphi_{62} \to N \in \mathbb{N} \land \forall n \in \mathbb{N}\,(N \leq n \to \lvert b_{n} - \beta \rvert < d_{0}) \to N \in \mathbb{N}$
\item $\varphi_{242} :\equiv \beta \neq 0 \to \forall n \in \mathbb{N}\,(b_{n} \neq 0) \to \varepsilon \in \mathbb{R} \to 0 < \varepsilon \to \varphi_{62} \to 0 < d_{0}$
\item $\varphi_{243} :\equiv \forall x\,\forall y\,(x \in \mathbb{R} \to y \in \mathbb{R} \to \lvert x \cdot y \rvert = \lvert x \rvert \cdot \lvert y \rvert) \to \varphi_{102}$
\item $\varphi_{244} :\equiv (\mathrm{Conv}(b,\beta) \to \beta \neq 0 \to \forall n \in \mathbb{N}\,(b_{n} \neq 0) \to \beta \in \mathbb{R}) \to \varphi_{92}$
\item $\varphi_{245} :\equiv \varepsilon \cdot (\lvert \beta \rvert \cdot 2^{-1} \cdot \lvert \beta \rvert) \cdot (\lvert b_{n} \rvert^{-1} \cdot \lvert \beta \rvert^{-1}) < z \to \varepsilon \cdot (\lvert \beta \rvert \cdot 2^{-1} \cdot \lvert \beta \rvert) \cdot (\lvert b_{n} \rvert^{-1} \cdot \lvert \beta \rvert^{-1}) < w$
\item $\varphi_{246} :\equiv N \in \mathbb{N} \land \forall n \in \mathbb{N}\,(N \leq n \to \lvert b_{n} - \beta \rvert < d_{0}) \to n \in \mathbb{N} \to \varphi_{70}$
\item $\varphi_{247} :\equiv \mathrm{Conv}(b,\beta) \to \beta \neq 0 \to \forall n \in \mathbb{N}\,(b_{n} \neq 0) \to \varepsilon \in \mathbb{R} \to 0 < \varepsilon \to \varphi_{60}$
\item $\varphi_{248} :\equiv \forall \varepsilon\,(\mathrm{Conv}(b,\beta) \land \beta \neq 0 \land \forall n \in \mathbb{N}\,(b_{n} \neq 0)) \to \forall \varepsilon \in \mathbb{R}\,(\varphi_{65})$
\item $\varphi_{249} :\equiv N \in \mathbb{N} \land \forall n \in \mathbb{N}\,(N \leq n \to \lvert b_{n} - \beta \rvert < d_{0}) \to n \in \mathbb{N} \to N \leq n \to \varphi_{45}$
\item $\varphi_{250} :\equiv N \in \mathbb{N} \land \forall n \in \mathbb{N}\,(N \leq n \to \lvert b_{n} - \beta \rvert < d_{0}) \to n \in \mathbb{N} \to N \leq n \to \varphi_{43}$
\item $\varphi_{251} :\equiv 0 < \varepsilon \to \varphi_{62} \to \exists N \in \mathbb{N}\,\forall n \in \mathbb{N}\,(N \leq n \to \lvert b_{n} - \beta \rvert < d_{0})$
\item $\varphi_{252} :\equiv \beta \neq 0 \to \forall n \in \mathbb{N}\,(b_{n} \neq 0) \to \varepsilon \in \mathbb{R} \to 0 < \varepsilon \to \varphi_{62} \to d_{0} \in \mathbb{R}$
\item $\varphi_{253} :\equiv \varphi_{62} \to \exists N \in \mathbb{N}\,\forall n \in \mathbb{N}\,(N \leq n \to \lvert b^{\ast}(n) - \beta^{-1} \rvert < \varepsilon)$
\item $\varphi_{254} :\equiv \mathrm{Conv}(b^{\ast},\beta^{-1}) \leftrightarrow \mathrm{Seq}(b^{\ast}) \land \beta^{-1} \in \mathbb{R} \land \forall \varepsilon \in \mathbb{R}\,(\varphi_{65})$
\item $\varphi_{255} :\equiv N \in \mathbb{N} \land \forall n \in \mathbb{N}\,(N \leq n \to \lvert b_{n} - \beta \rvert < d_{0}) \to n \in \mathbb{N} \to N \leq n \to \varphi_{69}$
\item $\varphi_{256} :\equiv \forall z\,(\lvert b_{n}^{-1} - \beta^{-1} \rvert = \lvert (\beta - b_{n}) \cdot (b_{n}^{-1} \cdot \beta^{-1}) \rvert \to \lvert (\beta - b_{n}) \cdot (b_{n}^{-1} \cdot \beta^{-1}) \rvert = z \to \lvert b_{n}^{-1} - \beta^{-1} \rvert = z)$
\item $\varphi_{257} :\equiv N \in \mathbb{N} \land \forall n \in \mathbb{N}\,(N \leq n \to \lvert b_{n} - \beta \rvert < d_{0}) \to n \in \mathbb{N} \to \varphi_{89}$
\item $\varphi_{258} :\equiv \forall x\,\forall y\,(x \in \mathbb{R} \to y \in \mathbb{R} \to 0 < x \to 0 < y \to 0 < x \cdot y) \to \varphi_{140}$
\item $\varphi_{259} :\equiv \forall y\,\forall z\,(\varepsilon \cdot (\lvert \beta \rvert \cdot 2^{-1} \cdot \lvert \beta \rvert) \cdot ((\lvert \beta \rvert \cdot 2^{-1})^{-1} \cdot \lvert \beta \rvert^{-1}) = y \to y = z \to \varepsilon \cdot (\lvert \beta \rvert \cdot 2^{-1} \cdot \lvert \beta \rvert) \cdot ((\lvert \beta \rvert \cdot 2^{-1})^{-1} \cdot \lvert \beta \rvert^{-1}) = z)$
\item $\varphi_{260} :\equiv \forall x\,\forall y\,(x \in \mathbb{R} \to y \in \mathbb{R} \to 0 < x \to 0 < y \to 0 < x \cdot y) \to \varphi_{152}$
\item $\varphi_{261} :\equiv \forall z\,(\lvert b_{n}^{-1} - \beta^{-1} \rvert = \lvert \beta - b_{n} \rvert \cdot \lvert b_{n}^{-1} \cdot \beta^{-1} \rvert \to \lvert \beta - b_{n} \rvert \cdot \lvert b_{n}^{-1} \cdot \beta^{-1} \rvert = z \to \lvert b_{n}^{-1} - \beta^{-1} \rvert = z)$
\item $\varphi_{262} :\equiv \lvert b_{n}^{-1} - \beta^{-1} \rvert = \lvert b_{n} - \beta \rvert \cdot (\lvert b_{n} \rvert^{-1} \cdot \lvert \beta \rvert^{-1}) \to \lvert b_{n} - \beta \rvert \cdot (\lvert b_{n} \rvert^{-1} \cdot \lvert \beta \rvert^{-1}) = \lvert b_{n}^{-1} - \beta^{-1} \rvert$
\item $\varphi_{263} :\equiv \forall y\,(\lvert b_{n} \rvert^{-1} \in \mathbb{R} \to y \in \mathbb{R} \to \lvert b_{n} \rvert^{-1} \cdot y \in \mathbb{R}) \to \lvert b_{n} \rvert^{-1} \in \mathbb{R} \to \lvert \beta \rvert^{-1} \in \mathbb{R} \to \lvert b_{n} \rvert^{-1} \cdot \lvert \beta \rvert^{-1} \in \mathbb{R}$
\item $\varphi_{264} :\equiv \forall y\,\forall z\,(\lvert b_{n} \rvert^{-1} \cdot \lvert \beta \rvert^{-1} \in \mathbb{R} \to y \in \mathbb{R} \to z \in \mathbb{R} \to 0 < z \to \lvert b_{n} \rvert^{-1} \cdot \lvert \beta \rvert^{-1} < y \to z \cdot (\lvert b_{n} \rvert^{-1} \cdot \lvert \beta \rvert^{-1}) < z \cdot y)$
\item $\varphi_{265} :\equiv N \in \mathbb{N} \land \forall n \in \mathbb{N}\,(N \leq n \to \lvert b_{n} - \beta \rvert < d_{0}) \to n \in \mathbb{N} \to N \leq n \to \varphi_{116}$
\item $\varphi_{266} :\equiv \forall x \in \mathbb{R}\,(x \neq 0 \to x \cdot x^{-1} = 1) \to \lvert \beta \rvert \cdot 2^{-1} \cdot \lvert \beta \rvert \in \mathbb{R} \to \varphi_{120}$
\item $\varphi_{267} :\equiv \forall z\,(\varphi_{69} \to \varepsilon \cdot 1 = z \to \varepsilon \cdot (\lvert \beta \rvert \cdot 2^{-1} \cdot \lvert \beta \rvert) \cdot ((\lvert \beta \rvert \cdot 2^{-1})^{-1} \cdot \lvert \beta \rvert^{-1}) = z)$
\item $\varphi_{268} :\equiv \forall z\,(\lvert b_{n}^{-1} - \beta^{-1} \rvert = \lvert \beta - b_{n} \rvert \cdot (\lvert b_{n} \rvert^{-1} \cdot \lvert \beta \rvert^{-1}) \to \lvert \beta - b_{n} \rvert \cdot (\lvert b_{n} \rvert^{-1} \cdot \lvert \beta \rvert^{-1}) = z \to \lvert b_{n}^{-1} - \beta^{-1} \rvert = z)$
\item $\varphi_{269} :\equiv \forall z\,(\lvert b_{n}^{-1} - \beta^{-1} \rvert = \lvert \beta - b_{n} \rvert \cdot (\lvert b_{n} \rvert^{-1} \cdot \lvert \beta^{-1} \rvert) \to \lvert \beta - b_{n} \rvert \cdot (\lvert b_{n} \rvert^{-1} \cdot \lvert \beta^{-1} \rvert) = z \to \lvert b_{n}^{-1} - \beta^{-1} \rvert = z)$
\item $\varphi_{270} :\equiv \forall z\,(\lvert b_{n}^{-1} - \beta^{-1} \rvert = \lvert \beta - b_{n} \rvert \cdot (\lvert b_{n}^{-1} \rvert \cdot \lvert \beta^{-1} \rvert) \to \lvert \beta - b_{n} \rvert \cdot (\lvert b_{n}^{-1} \rvert \cdot \lvert \beta^{-1} \rvert) = z \to \lvert b_{n}^{-1} - \beta^{-1} \rvert = z)$
\item $\varphi_{271} :\equiv N \in \mathbb{N} \land \forall n \in \mathbb{N}\,(N \leq n \to \lvert b_{n} - \beta \rvert < d_{0}) \to n \in \mathbb{N} \to N \leq n \to \varphi_{157}$
\item $\varphi_{272} :\equiv \forall z\,(\lvert b_{n} \rvert^{-1} \in \mathbb{R} \to (\lvert \beta \rvert \cdot 2^{-1})^{-1} \in \mathbb{R} \to z \in \mathbb{R} \to 0 < z \to \lvert b_{n} \rvert^{-1} < (\lvert \beta \rvert \cdot 2^{-1})^{-1} \to \lvert b_{n} \rvert^{-1} \cdot z < (\lvert \beta \rvert \cdot 2^{-1})^{-1} \cdot z)$
\item $\varphi_{273} :\equiv \forall z\,(\lvert b_{n} - \beta \rvert < d_{0} \to d_{0} \leq z \to \lvert b_{n} - \beta \rvert < z) \to \lvert b_{n} - \beta \rvert < d_{0} \to \varphi_{112}$
\item $\varphi_{274} :\equiv \lvert b_{n} - \beta \rvert \cdot (\lvert b_{n} \rvert^{-1} \cdot \lvert \beta \rvert^{-1}) = \lvert b_{n}^{-1} - \beta^{-1} \rvert \to \lvert b_{n} - \beta \rvert \cdot (\lvert b_{n} \rvert^{-1} \cdot \lvert \beta \rvert^{-1}) < \varepsilon \to \lvert b_{n}^{-1} - \beta^{-1} \rvert < \varepsilon$
\item $\varphi_{275} :\equiv N \in \mathbb{N} \land \forall n \in \mathbb{N}\,(N \leq n \to \lvert b_{n} - \beta \rvert < d_{0}) \to n \in \mathbb{N} \to N \leq n \to \varphi_{166}$
\item $\varphi_{276} :\equiv N \in \mathbb{N} \land \forall n \in \mathbb{N}\,(N \leq n \to \lvert b_{n} - \beta \rvert < d_{0}) \to n \in \mathbb{N} \to N \leq n \to \varphi_{167}$
\item $\varphi_{277} :\equiv N \in \mathbb{N} \land \forall n \in \mathbb{N}\,(N \leq n \to \lvert b_{n} - \beta \rvert < d_{0}) \to n \in \mathbb{N} \to N \leq n \to \varphi_{168}$
\item $\varphi_{278} :\equiv \forall w\,(b_{n}^{-1} = w \to \lvert b_{n}^{-1} - \beta^{-1} \rvert < \varepsilon \to \lvert w - \beta^{-1} \rvert < \varepsilon) \to \varphi_{151}$
\item $\varphi_{279} :\equiv \lvert \beta - b_{n} \rvert \cdot (\lvert b_{n} \rvert^{-1} \cdot \lvert \beta \rvert^{-1}) = z \cdot (\lvert b_{n} \rvert^{-1} \cdot \lvert \beta \rvert^{-1}) \to \lvert \beta - b_{n} \rvert \cdot (\lvert b_{n} \rvert^{-1} \cdot \lvert \beta \rvert^{-1}) = w \cdot (\lvert b_{n} \rvert^{-1} \cdot \lvert \beta \rvert^{-1})$
\item $\varphi_{280} :\equiv \varepsilon \cdot (\lvert \beta \rvert \cdot 2^{-1} \cdot \lvert \beta \rvert) \in \mathbb{R} \to 0 < \lvert \beta \rvert \cdot 2^{-1} \to 0 < \varepsilon \cdot (\lvert \beta \rvert \cdot 2^{-1} \cdot \lvert \beta \rvert) \to \varphi_{60}$
\item $\varphi_{281} :\equiv \forall y\,(\varepsilon \in \mathbb{R} \to y \in \mathbb{R} \to 0 < \varepsilon \to 0 < y \to 0 < \varepsilon \cdot y) \to \varepsilon \in \mathbb{R} \to \lvert \beta \rvert \cdot 2^{-1} \cdot \lvert \beta \rvert \in \mathbb{R} \to 0 < \varepsilon \to \varphi_{36}$
\item $\varphi_{282} :\equiv N \in \mathbb{N} \land \forall n \in \mathbb{N}\,(N \leq n \to \lvert b_{n} - \beta \rvert < d_{0}) \to n \in \mathbb{N} \to N \leq n \to \varphi_{207}$
\item $\varphi_{283} :\equiv (\mathrm{Conv}(b,\beta) \to \varphi_{42}) \to 2^{-1} \in \mathbb{R} \to \mathrm{Conv}(b,\beta) \to \beta \neq 0 \to \varphi_{44}$
\item $\varphi_{284} :\equiv \forall w\,(\lvert b_{n}^{-1} \cdot \beta^{-1} \rvert = w \to \varphi_{133} \to \lvert \beta - b_{n} \rvert \cdot \lvert b_{n}^{-1} \cdot \beta^{-1} \rvert = \lvert \beta - b_{n} \rvert \cdot w)$
\item $\varphi_{285} :\equiv \mathrm{Conv}(b,\beta) \land \beta \neq 0 \land \forall n \in \mathbb{N}\,(b_{n} \neq 0) \land \varepsilon \in \mathbb{R} \land 0 < \varepsilon \land \varphi_{62} \land (N \in \mathbb{N} \land \forall n \in \mathbb{N}\,(N \leq n \to \lvert b_{n} - \beta \rvert < d_{0}))$
\item $\varphi_{286} :\equiv N \in \mathbb{N} \land \forall n \in \mathbb{N}\,(N \leq n \to \lvert b_{n} - \beta \rvert < d_{0}) \to n \in \mathbb{N} \to N \leq n \to \varphi_{218}$
\item $\varphi_{287} :\equiv \mathrm{Conv}(b,\beta) \to \beta \neq 0 \to \forall n \in \mathbb{N}\,(b_{n} \neq 0) \to \varepsilon \in \mathbb{R} \to \varphi_{251}$
\item $\varphi_{288} :\equiv (\lvert \beta \rvert \in \mathbb{R} \to 2^{-1} \in \mathbb{R} \to \lvert \beta \rvert \cdot 2^{-1} \in \mathbb{R}) \to \varphi_{15} \to 2^{-1} \in \mathbb{R} \to \varphi_{33}$
\item $\varphi_{289} :\equiv N \in \mathbb{N} \land \forall n \in \mathbb{N}\,(N \leq n \to \lvert b_{n} - \beta \rvert < d_{0}) \to n \in \mathbb{N} \to N \leq n \to \varphi_{226}$
\item $\varphi_{290} :\equiv N \in \mathbb{N} \land \forall n \in \mathbb{N}\,(N \leq n \to \lvert b_{n} - \beta \rvert < d_{0}) \to n \in \mathbb{N} \to N \leq n \to \varphi_{225}$
\item $\varphi_{291} :\equiv \lvert \beta \rvert \cdot 2^{-1} \cdot \lvert \beta \rvert \in \mathbb{R} \to (\lvert \beta \rvert \cdot 2^{-1} \cdot \lvert \beta \rvert)^{-1} \in \mathbb{R} \to \varphi_{218}$
\item $\varphi_{292} :\equiv \lvert b_{n} \rvert^{-1} \in \mathbb{R} \to (\lvert \beta \rvert \cdot 2^{-1})^{-1} \in \mathbb{R} \to \lvert \beta \rvert^{-1} \in \mathbb{R} \to 0 < \lvert \beta \rvert^{-1} \to \varphi_{191}$
\item $\varphi_{293} :\equiv \forall z\,(\lvert b_{n} - \beta \rvert \in \mathbb{R} \to \varepsilon \cdot (\lvert \beta \rvert \cdot 2^{-1} \cdot \lvert \beta \rvert) \in \mathbb{R} \to z \in \mathbb{R} \to 0 < z \to \lvert b_{n} - \beta \rvert < \varepsilon \cdot (\lvert \beta \rvert \cdot 2^{-1} \cdot \lvert \beta \rvert) \to \lvert b_{n} - \beta \rvert \cdot z < \varepsilon \cdot (\lvert \beta \rvert \cdot 2^{-1} \cdot \lvert \beta \rvert) \cdot z)$
\item $\varphi_{294} :\equiv N \in \mathbb{N} \land \forall n \in \mathbb{N}\,(N \leq n \to \lvert b_{n} - \beta \rvert < d_{0}) \to n \in \mathbb{N} \to N \leq n \to \varphi_{232}$
\item $\varphi_{295} :\equiv \mathrm{Conv}(b,\beta) \to \beta \neq 0 \to \forall n \in \mathbb{N}\,(b_{n} \neq 0) \to \varepsilon \in \mathbb{R} \to 0 < \varepsilon \to \varphi_{241}$
\item $\varphi_{296} :\equiv (\mathrm{Conv}(b,\beta) \to \beta \neq 0 \to \varphi_{87}) \to \mathrm{Conv}(b,\beta) \to \beta \neq 0 \to \varphi_{113}$
\item $\varphi_{297} :\equiv 0 < \varepsilon \to \varphi_{62} \to N \in \mathbb{N} \land \forall n \in \mathbb{N}\,(N \leq n \to \lvert b_{n} - \beta \rvert < d_{0}) \to \varphi_{38}$
\item $\varphi_{298} :\equiv 0 < \lvert b_{n} \rvert^{-1} \cdot \lvert \beta \rvert^{-1} \to \lvert b_{n} - \beta \rvert < \varepsilon \cdot (\lvert \beta \rvert \cdot 2^{-1} \cdot \lvert \beta \rvert) \to \varphi_{207}$
\item $\varphi_{299} :\equiv \mathrm{Conv}(b,\beta) \to \beta \neq 0 \to \forall n \in \mathbb{N}\,(b_{n} \neq 0) \to \varepsilon \in \mathbb{R} \to 0 < \varepsilon \to \varphi_{253}$
\item $\varphi_{300} :\equiv \mathrm{Conv}(b,\beta) \land \beta \neq 0 \land \forall n \in \mathbb{N}\,(b_{n} \neq 0) \land \varepsilon \in \mathbb{R} \land 0 < \varepsilon \land \varphi_{62} \land \neg \exists N \in \mathbb{N}\,\forall n \in \mathbb{N}\,(N \leq n \to \lvert b^{\ast}(n) - \beta^{-1} \rvert < \varepsilon)$
\item $\varphi_{301} :\equiv \forall w\,(\lvert \beta - b_{n} \rvert = w \to \varphi_{163} \to \lvert \beta - b_{n} \rvert \cdot (\lvert b_{n} \rvert^{-1} \cdot \lvert \beta \rvert^{-1}) = w \cdot (\lvert b_{n} \rvert^{-1} \cdot \lvert \beta \rvert^{-1}))$
\item $\varphi_{302} :\equiv 0 < \varepsilon \to \varphi_{62} \to N \in \mathbb{N} \land \forall n \in \mathbb{N}\,(N \leq n \to \lvert b_{n} - \beta \rvert < d_{0}) \to \varphi_{63}$
\item $\varphi_{303} :\equiv \mathrm{Conv}(b,\beta) \to \beta \neq 0 \to \forall n \in \mathbb{N}\,(b_{n} \neq 0) \to \varepsilon \in \mathbb{R} \to 0 < \varepsilon \to \varphi_{62} \to \varphi_{52}$
\item $\varphi_{304} :\equiv \mathrm{Conv}(b,\beta) \to \beta \neq 0 \to \forall n \in \mathbb{N}\,(b_{n} \neq 0) \to \varepsilon \in \mathbb{R} \to 0 < \varepsilon \to \varphi_{62} \to \varphi_{53}$
\item $\varphi_{305} :\equiv \mathrm{Conv}(b,\beta) \to \beta \neq 0 \to \forall n \in \mathbb{N}\,(b_{n} \neq 0) \to \varepsilon \in \mathbb{R} \to 0 < \varepsilon \to \varphi_{62} \to \varphi_{62}$
\item $\varphi_{306} :\equiv \mathrm{Conv}(b,\beta) \to \beta \neq 0 \to \forall n \in \mathbb{N}\,(b_{n} \neq 0) \to \varepsilon \in \mathbb{R} \to 0 < \varepsilon \to \varphi_{62} \to \varphi_{81}$
\item $\varphi_{307} :\equiv 0 < \varepsilon \to \varphi_{62} \to N \in \mathbb{N} \land \forall n \in \mathbb{N}\,(N \leq n \to \lvert b_{n} - \beta \rvert < d_{0}) \to \varphi_{91}$
\item $\varphi_{308} :\equiv \mathrm{Conv}(b,\beta) \to \beta \neq 0 \to \forall n \in \mathbb{N}\,(b_{n} \neq 0) \to \varepsilon \in \mathbb{R} \to 0 < \varepsilon \to \varphi_{62} \to \varphi_{84}$
\item $\varphi_{309} :\equiv 0 < \varepsilon \to \varphi_{62} \to N \in \mathbb{N} \land \forall n \in \mathbb{N}\,(N \leq n \to \lvert b_{n} - \beta \rvert < d_{0}) \to \varphi_{98}$
\item $\varphi_{310} :\equiv \mathrm{Conv}(b,\beta) \to \beta \neq 0 \to \forall n \in \mathbb{N}\,(b_{n} \neq 0) \to \varepsilon \in \mathbb{R} \to 0 < \varepsilon \to \varphi_{62} \to \varphi_{93}$
\item $\varphi_{311} :\equiv \mathrm{Conv}(b,\beta) \to \beta \neq 0 \to \forall n \in \mathbb{N}\,(b_{n} \neq 0) \to \varepsilon \in \mathbb{R} \to 0 < \varepsilon \to \varphi_{62} \to \varphi_{94}$
\item $\varphi_{312} :\equiv \mathrm{Conv}(b,\beta) \to \beta \neq 0 \to \forall n \in \mathbb{N}\,(b_{n} \neq 0) \to \varepsilon \in \mathbb{R} \to 0 < \varepsilon \to \varphi_{62} \to \varphi_{95}$
\item $\varphi_{313} :\equiv \mathrm{Conv}(b,\beta) \to \beta \neq 0 \to \forall n \in \mathbb{N}\,(b_{n} \neq 0) \to \varepsilon \in \mathbb{R} \to 0 < \varepsilon \to \varphi_{62} \to \varphi_{97}$
\item $\varphi_{314} :\equiv \mathrm{Conv}(b,\beta) \to \beta \neq 0 \to \forall n \in \mathbb{N}\,(b_{n} \neq 0) \to \varepsilon \in \mathbb{R} \to 0 < \varepsilon \to \varphi_{62} \to \varphi_{103}$
\item $\varphi_{315} :\equiv \mathrm{Conv}(b,\beta) \to \beta \neq 0 \to \forall n \in \mathbb{N}\,(b_{n} \neq 0) \to \varepsilon \in \mathbb{R} \to 0 < \varepsilon \to \varphi_{62} \to \varphi_{104}$
\item $\varphi_{316} :\equiv 0 < \varepsilon \to \varphi_{62} \to N \in \mathbb{N} \land \forall n \in \mathbb{N}\,(N \leq n \to \lvert b_{n} - \beta \rvert < d_{0}) \to \varphi_{117}$
\item $\varphi_{317} :\equiv \mathrm{Conv}(b,\beta) \to \beta \neq 0 \to \forall n \in \mathbb{N}\,(b_{n} \neq 0) \to \varepsilon \in \mathbb{R} \to 0 < \varepsilon \to \varphi_{62} \to \varphi_{107}$
\item $\varphi_{318} :\equiv \mathrm{Conv}(b,\beta) \to \beta \neq 0 \to \forall n \in \mathbb{N}\,(b_{n} \neq 0) \to \varepsilon \in \mathbb{R} \to 0 < \varepsilon \to \varphi_{62} \to \varphi_{111}$
\item $\varphi_{319} :\equiv \mathrm{Conv}(b,\beta) \to \beta \neq 0 \to \forall n \in \mathbb{N}\,(b_{n} \neq 0) \to \varepsilon \in \mathbb{R} \to 0 < \varepsilon \to \varphi_{62} \to \varphi_{115}$
\item $\varphi_{320} :\equiv \mathrm{Conv}(b,\beta) \to \beta \neq 0 \to \forall n \in \mathbb{N}\,(b_{n} \neq 0) \to \varepsilon \in \mathbb{R} \to 0 < \varepsilon \to \varphi_{62} \to \varphi_{119}$
\item $\varphi_{321} :\equiv \mathrm{Conv}(b,\beta) \to \beta \neq 0 \to \forall n \in \mathbb{N}\,(b_{n} \neq 0) \to \varepsilon \in \mathbb{R} \to 0 < \varepsilon \to \varphi_{62} \to \varphi_{121}$
\item $\varphi_{322} :\equiv \lvert b_{n}^{-1} - \beta^{-1} \rvert = \lvert (\beta - b_{n}) \cdot (b_{n}^{-1} \cdot \beta^{-1}) \rvert \to \varphi_{116} \to \lvert b_{n}^{-1} - \beta^{-1} \rvert = \lvert \beta - b_{n} \rvert \cdot \lvert b_{n}^{-1} \cdot \beta^{-1} \rvert$
\item $\varphi_{323} :\equiv 0 < \varepsilon \to \varphi_{62} \to N \in \mathbb{N} \land \forall n \in \mathbb{N}\,(N \leq n \to \lvert b_{n} - \beta \rvert < d_{0}) \to \varphi_{135}$
\item $\varphi_{324} :\equiv 0 < \varepsilon \to \varphi_{62} \to N \in \mathbb{N} \land \forall n \in \mathbb{N}\,(N \leq n \to \lvert b_{n} - \beta \rvert < d_{0}) \to \varphi_{136}$
\item $\varphi_{325} :\equiv 0 < \varepsilon \to \varphi_{62} \to N \in \mathbb{N} \land \forall n \in \mathbb{N}\,(N \leq n \to \lvert b_{n} - \beta \rvert < d_{0}) \to \varphi_{137}$
\item $\varphi_{326} :\equiv 0 < \varepsilon \to \varphi_{62} \to N \in \mathbb{N} \land \forall n \in \mathbb{N}\,(N \leq n \to \lvert b_{n} - \beta \rvert < d_{0}) \to \varphi_{138}$
\item $\varphi_{327} :\equiv \mathrm{Conv}(b,\beta) \to \beta \neq 0 \to \forall n \in \mathbb{N}\,(b_{n} \neq 0) \to \varepsilon \in \mathbb{R} \to 0 < \varepsilon \to \varphi_{62} \to \varphi_{123}$
\item $\varphi_{328} :\equiv \mathrm{Conv}(b,\beta) \to \beta \neq 0 \to \forall n \in \mathbb{N}\,(b_{n} \neq 0) \to \varepsilon \in \mathbb{R} \to 0 < \varepsilon \to \varphi_{62} \to \varphi_{124}$
\item $\varphi_{329} :\equiv \mathrm{Conv}(b,\beta) \to \beta \neq 0 \to \forall n \in \mathbb{N}\,(b_{n} \neq 0) \to \varepsilon \in \mathbb{R} \to 0 < \varepsilon \to \varphi_{62} \to \varphi_{125}$
\item $\varphi_{330} :\equiv \mathrm{Conv}(b,\beta) \to \beta \neq 0 \to \forall n \in \mathbb{N}\,(b_{n} \neq 0) \to \varepsilon \in \mathbb{R} \to 0 < \varepsilon \to \varphi_{62} \to \varphi_{128}$
\item $\varphi_{331} :\equiv \mathrm{Conv}(b,\beta) \to \beta \neq 0 \to \forall n \in \mathbb{N}\,(b_{n} \neq 0) \to \varepsilon \in \mathbb{R} \to 0 < \varepsilon \to \varphi_{62} \to \varphi_{131}$
\item $\varphi_{332} :\equiv \mathrm{Conv}(b,\beta) \to \beta \neq 0 \to \forall n \in \mathbb{N}\,(b_{n} \neq 0) \to \varepsilon \in \mathbb{R} \to 0 < \varepsilon \to \varphi_{62} \to \varphi_{134}$
\item $\varphi_{333} :\equiv \mathrm{Conv}(b,\beta) \to \beta \neq 0 \to \forall n \in \mathbb{N}\,(b_{n} \neq 0) \to \varepsilon \in \mathbb{R} \to 0 < \varepsilon \to \varphi_{62} \to \varphi_{144}$
\item $\varphi_{334} :\equiv \mathrm{Conv}(b,\beta) \to \beta \neq 0 \to \forall n \in \mathbb{N}\,(b_{n} \neq 0) \to \varepsilon \in \mathbb{R} \to 0 < \varepsilon \to \varphi_{62} \to \varphi_{145}$
\item $\varphi_{335} :\equiv \mathrm{Conv}(b,\beta) \to \beta \neq 0 \to \forall n \in \mathbb{N}\,(b_{n} \neq 0) \to \varepsilon \in \mathbb{R} \to 0 < \varepsilon \to \varphi_{62} \to \varphi_{149}$
\item $\varphi_{336} :\equiv \mathrm{Conv}(b,\beta) \to \beta \neq 0 \to \forall n \in \mathbb{N}\,(b_{n} \neq 0) \to \varepsilon \in \mathbb{R} \to 0 < \varepsilon \to \varphi_{62} \to \varphi_{150}$
\item $\varphi_{337} :\equiv \mathrm{Conv}(b,\beta) \to \beta \neq 0 \to \forall n \in \mathbb{N}\,(b_{n} \neq 0) \to \varepsilon \in \mathbb{R} \to 0 < \varepsilon \to \varphi_{62} \to \varphi_{154}$
\item $\varphi_{338} :\equiv \mathrm{Conv}(b,\beta) \to \beta \neq 0 \to \forall n \in \mathbb{N}\,(b_{n} \neq 0) \to \varepsilon \in \mathbb{R} \to 0 < \varepsilon \to \varphi_{62} \to \varphi_{155}$
\item $\varphi_{339} :\equiv \mathrm{Conv}(b,\beta) \to \beta \neq 0 \to \forall n \in \mathbb{N}\,(b_{n} \neq 0) \to \varepsilon \in \mathbb{R} \to 0 < \varepsilon \to \varphi_{62} \to \varphi_{156}$
\item $\varphi_{340} :\equiv \mathrm{Conv}(b,\beta) \to \beta \neq 0 \to \forall n \in \mathbb{N}\,(b_{n} \neq 0) \to \varepsilon \in \mathbb{R} \to 0 < \varepsilon \to \varphi_{62} \to \varphi_{159}$
\item $\varphi_{341} :\equiv 0 < \varepsilon \cdot (\lvert \beta \rvert \cdot 2^{-1} \cdot \lvert \beta \rvert) \to \lvert b_{n} \rvert^{-1} \cdot \lvert \beta \rvert^{-1} < (\lvert \beta \rvert \cdot 2^{-1})^{-1} \cdot \lvert \beta \rvert^{-1} \to \varphi_{232}$
\item $\varphi_{342} :\equiv \mathrm{Conv}(b,\beta) \to \beta \neq 0 \to \forall n \in \mathbb{N}\,(b_{n} \neq 0) \to \varepsilon \in \mathbb{R} \to 0 < \varepsilon \to \varphi_{62} \to \varphi_{162}$
\item $\varphi_{343} :\equiv \mathrm{Conv}(b,\beta) \to \beta \neq 0 \to \forall n \in \mathbb{N}\,(b_{n} \neq 0) \to \varepsilon \in \mathbb{R} \to 0 < \varepsilon \to \varphi_{62} \to \varphi_{169}$
\item $\varphi_{344} :\equiv \mathrm{Conv}(b,\beta) \to \beta \neq 0 \to \forall n \in \mathbb{N}\,(b_{n} \neq 0) \to \varepsilon \in \mathbb{R} \to 0 < \varepsilon \to \varphi_{62} \to \varphi_{170}$
\item $\varphi_{345} :\equiv \mathrm{Conv}(b,\beta) \to \beta \neq 0 \to \forall n \in \mathbb{N}\,(b_{n} \neq 0) \to \varepsilon \in \mathbb{R} \to 0 < \varepsilon \to \varphi_{62} \to \varphi_{171}$
\item $\varphi_{346} :\equiv \mathrm{Conv}(b,\beta) \to \beta \neq 0 \to \forall n \in \mathbb{N}\,(b_{n} \neq 0) \to \varepsilon \in \mathbb{R} \to 0 < \varepsilon \to \varphi_{62} \to \varphi_{172}$
\item $\varphi_{347} :\equiv \mathrm{Conv}(b,\beta) \to \beta \neq 0 \to \forall n \in \mathbb{N}\,(b_{n} \neq 0) \to \varepsilon \in \mathbb{R} \to 0 < \varepsilon \to \varphi_{62} \to \varphi_{173}$
\item $\varphi_{348} :\equiv \mathrm{Conv}(b,\beta) \to \beta \neq 0 \to \forall n \in \mathbb{N}\,(b_{n} \neq 0) \to \varepsilon \in \mathbb{R} \to 0 < \varepsilon \to \varphi_{62} \to \varphi_{178}$
\item $\varphi_{349} :\equiv \mathrm{Conv}(b,\beta) \to \beta \neq 0 \to \forall n \in \mathbb{N}\,(b_{n} \neq 0) \to \varepsilon \in \mathbb{R} \to 0 < \varepsilon \to \varphi_{62} \to \varphi_{179}$
\item $\varphi_{350} :\equiv \mathrm{Conv}(b,\beta) \to \beta \neq 0 \to \forall n \in \mathbb{N}\,(b_{n} \neq 0) \to \varepsilon \in \mathbb{R} \to 0 < \varepsilon \to \varphi_{62} \to \varphi_{180}$
\item $\varphi_{351} :\equiv \mathrm{Conv}(b,\beta) \to \beta \neq 0 \to \forall n \in \mathbb{N}\,(b_{n} \neq 0) \to \varepsilon \in \mathbb{R} \to 0 < \varepsilon \to \varphi_{62} \to \varphi_{181}$
\item $\varphi_{352} :\equiv \forall z\,(\varphi_{207} \to \varepsilon \cdot (\lvert \beta \rvert \cdot 2^{-1} \cdot \lvert \beta \rvert) \cdot (\lvert b_{n} \rvert^{-1} \cdot \lvert \beta \rvert^{-1}) < z \to \lvert b_{n} - \beta \rvert \cdot (\lvert b_{n} \rvert^{-1} \cdot \lvert \beta \rvert^{-1}) < z)$
\item $\varphi_{353} :\equiv \mathrm{Conv}(b,\beta) \to \beta \neq 0 \to \forall n \in \mathbb{N}\,(b_{n} \neq 0) \to \varepsilon \in \mathbb{R} \to 0 < \varepsilon \to \varphi_{62} \to \varphi_{183}$
\item $\varphi_{354} :\equiv \mathrm{Conv}(b,\beta) \to \beta \neq 0 \to \forall n \in \mathbb{N}\,(b_{n} \neq 0) \to \varepsilon \in \mathbb{R} \to 0 < \varepsilon \to \varphi_{62} \to \varphi_{184}$
\item $\varphi_{355} :\equiv \mathrm{Conv}(b,\beta) \to \beta \neq 0 \to \forall n \in \mathbb{N}\,(b_{n} \neq 0) \to \varepsilon \in \mathbb{R} \to 0 < \varepsilon \to \varphi_{62} \to \varphi_{186}$
\item $\varphi_{356} :\equiv \lvert b_{n}^{-1} - \beta^{-1} \rvert = \lvert \beta - b_{n} \rvert \cdot \lvert b_{n}^{-1} \cdot \beta^{-1} \rvert \to \varphi_{157} \to \lvert b_{n}^{-1} - \beta^{-1} \rvert = \lvert \beta - b_{n} \rvert \cdot (\lvert b_{n}^{-1} \rvert \cdot \lvert \beta^{-1} \rvert)$
\item $\varphi_{357} :\equiv \mathrm{Conv}(b,\beta) \to \beta \neq 0 \to \forall n \in \mathbb{N}\,(b_{n} \neq 0) \to \varepsilon \in \mathbb{R} \to 0 < \varepsilon \to \varphi_{62} \to \varphi_{188}$
\item $\varphi_{358} :\equiv \mathrm{Conv}(b,\beta) \to \beta \neq 0 \to \forall n \in \mathbb{N}\,(b_{n} \neq 0) \to \varepsilon \in \mathbb{R} \to 0 < \varepsilon \to \varphi_{62} \to \varphi_{189}$
\item $\varphi_{359} :\equiv \forall z\,(\lvert b_{n} \rvert^{-1} \cdot \lvert \beta \rvert^{-1} \in \mathbb{R} \to (\lvert \beta \rvert \cdot 2^{-1})^{-1} \cdot \lvert \beta \rvert^{-1} \in \mathbb{R} \to z \in \mathbb{R} \to 0 < z \to \lvert b_{n} \rvert^{-1} \cdot \lvert \beta \rvert^{-1} < (\lvert \beta \rvert \cdot 2^{-1})^{-1} \cdot \lvert \beta \rvert^{-1} \to z \cdot (\lvert b_{n} \rvert^{-1} \cdot \lvert \beta \rvert^{-1}) < z \cdot ((\lvert \beta \rvert \cdot 2^{-1})^{-1} \cdot \lvert \beta \rvert^{-1}))$
\item $\varphi_{360} :\equiv \mathrm{Conv}(b,\beta) \to \beta \neq 0 \to \forall n \in \mathbb{N}\,(b_{n} \neq 0) \to \varepsilon \in \mathbb{R} \to 0 < \varepsilon \to \varphi_{62} \to \varphi_{190}$
\item $\varphi_{361} :\equiv \mathrm{Conv}(b,\beta) \to \beta \neq 0 \to \forall n \in \mathbb{N}\,(b_{n} \neq 0) \to \varepsilon \in \mathbb{R} \to 0 < \varepsilon \to \varphi_{62} \to \varphi_{192}$
\item $\varphi_{362} :\equiv \mathrm{Conv}(b,\beta) \to \beta \neq 0 \to \forall n \in \mathbb{N}\,(b_{n} \neq 0) \to \varepsilon \in \mathbb{R} \to 0 < \varepsilon \to \varphi_{62} \to \varphi_{193}$
\item $\varphi_{363} :\equiv \mathrm{Conv}(b,\beta) \to \beta \neq 0 \to \forall n \in \mathbb{N}\,(b_{n} \neq 0) \to \varepsilon \in \mathbb{R} \to 0 < \varepsilon \to \varphi_{62} \to \varphi_{194}$
\item $\varphi_{364} :\equiv \mathrm{Conv}(b,\beta) \to \beta \neq 0 \to \forall n \in \mathbb{N}\,(b_{n} \neq 0) \to \varepsilon \in \mathbb{R} \to 0 < \varepsilon \to \varphi_{62} \to \varphi_{195}$
\item $\varphi_{365} :\equiv \lvert b_{n}^{-1} - \beta^{-1} \rvert = \lvert \beta - b_{n} \rvert \cdot (\lvert b_{n} \rvert^{-1} \cdot \lvert \beta \rvert^{-1}) \to \varphi_{166} \to \lvert b_{n}^{-1} - \beta^{-1} \rvert = \lvert b_{n} - \beta \rvert \cdot (\lvert b_{n} \rvert^{-1} \cdot \lvert \beta \rvert^{-1})$
\item $\varphi_{366} :\equiv \lvert b_{n}^{-1} - \beta^{-1} \rvert = \lvert \beta - b_{n} \rvert \cdot (\lvert b_{n} \rvert^{-1} \cdot \lvert \beta^{-1} \rvert) \to \varphi_{167} \to \lvert b_{n}^{-1} - \beta^{-1} \rvert = \lvert \beta - b_{n} \rvert \cdot (\lvert b_{n} \rvert^{-1} \cdot \lvert \beta \rvert^{-1})$
\item $\varphi_{367} :\equiv \lvert b_{n}^{-1} - \beta^{-1} \rvert = \lvert \beta - b_{n} \rvert \cdot (\lvert b_{n}^{-1} \rvert \cdot \lvert \beta^{-1} \rvert) \to \varphi_{168} \to \lvert b_{n}^{-1} - \beta^{-1} \rvert = \lvert \beta - b_{n} \rvert \cdot (\lvert b_{n} \rvert^{-1} \cdot \lvert \beta^{-1} \rvert)$
\item $\varphi_{368} :\equiv \mathrm{Conv}(b,\beta) \to \beta \neq 0 \to \forall n \in \mathbb{N}\,(b_{n} \neq 0) \to \varepsilon \in \mathbb{R} \to 0 < \varepsilon \to \varphi_{62} \to \varphi_{205}$
\item $\varphi_{369} :\equiv \lvert \beta \rvert \cdot 2^{-1} \cdot \lvert \beta \rvert \cdot (\lvert \beta \rvert \cdot 2^{-1} \cdot \lvert \beta \rvert)^{-1} = 1 \to \varphi_{217} \to \varphi_{45}$
\item $\varphi_{370} :\equiv \mathrm{Conv}(b,\beta) \to \beta \neq 0 \to \forall n \in \mathbb{N}\,(b_{n} \neq 0) \to \varepsilon \in \mathbb{R} \to 0 < \varepsilon \to \varphi_{62} \to \varphi_{209}$
\item $\varphi_{371} :\equiv (\mathrm{Conv}(b,\beta) \to \beta \neq 0 \to \forall n \in \mathbb{N}\,(b_{n} \neq 0) \to \beta^{-1} \in \mathbb{R}) \to \varphi_{235} \to \varphi_{24}$
\item $\varphi_{372} :\equiv \mathrm{Conv}(b,\beta) \to \beta \neq 0 \to \forall n \in \mathbb{N}\,(b_{n} \neq 0) \to \varepsilon \in \mathbb{R} \to \varphi_{297}$
\item $\varphi_{373} :\equiv \forall w\,(\lvert \beta \rvert \cdot 2^{-1} \cdot \lvert \beta \rvert \cdot (\lvert \beta \rvert \cdot 2^{-1} \cdot \lvert \beta \rvert)^{-1} = w \to \varphi_{217} \to \varphi_{46})$
\item $\varphi_{374} :\equiv \forall N\,(\varphi_{300} \to \neg (N \in \mathbb{N} \land \forall n \in \mathbb{N}\,(N \leq n \to \lvert b_{n} - \beta \rvert < d_{0})))$
\item $\varphi_{375} :\equiv \mathrm{Conv}(b,\beta) \to \beta \neq 0 \to \forall n \in \mathbb{N}\,(b_{n} \neq 0) \to \varepsilon \in \mathbb{R} \to 0 < \varepsilon \to \varphi_{62} \to \varphi_{215}$
\item $\varphi_{376} :\equiv \mathrm{Conv}(b,\beta) \to \beta \neq 0 \to \forall n \in \mathbb{N}\,(b_{n} \neq 0) \to \varepsilon \in \mathbb{R} \to 0 < \varepsilon \to \varphi_{62} \to \varphi_{216}$
\item $\varphi_{377} :\equiv \mathrm{Conv}(b,\beta) \to \beta \neq 0 \to \forall n \in \mathbb{N}\,(b_{n} \neq 0) \to \varepsilon \in \mathbb{R} \to 0 < \varepsilon \to \varphi_{62} \to \varphi_{219}$
\item $\varphi_{378} :\equiv \mathrm{Conv}(b,\beta) \to \beta \neq 0 \to \forall n \in \mathbb{N}\,(b_{n} \neq 0) \to \varepsilon \in \mathbb{R} \to \varphi_{302}$
\item $\varphi_{379} :\equiv \forall z\,(\varphi_{225} \to \varepsilon \cdot (\lvert \beta \rvert \cdot 2^{-1} \cdot \lvert \beta \rvert \cdot (\lvert \beta \rvert \cdot 2^{-1} \cdot \lvert \beta \rvert)^{-1}) = z \to \varepsilon \cdot (\lvert \beta \rvert \cdot 2^{-1} \cdot \lvert \beta \rvert) \cdot ((\lvert \beta \rvert \cdot 2^{-1})^{-1} \cdot \lvert \beta \rvert^{-1}) = z)$
\item $\varphi_{380} :\equiv \forall z\,(\varphi_{226} \to \varepsilon \cdot (\lvert \beta \rvert \cdot 2^{-1} \cdot \lvert \beta \rvert) \cdot (\lvert \beta \rvert \cdot 2^{-1} \cdot \lvert \beta \rvert)^{-1} = z \to \varepsilon \cdot (\lvert \beta \rvert \cdot 2^{-1} \cdot \lvert \beta \rvert) \cdot ((\lvert \beta \rvert \cdot 2^{-1})^{-1} \cdot \lvert \beta \rvert^{-1}) = z)$
\item $\varphi_{381} :\equiv (\mathrm{Conv}(b,\beta) \to \beta \neq 0 \to \varphi_{44}) \to (\mathrm{Conv}(b,\beta) \to \varphi_{42}) \to \mathrm{Conv}(b,\beta) \to \beta \neq 0 \to \varphi_{87}$
\item $\varphi_{382} :\equiv \mathrm{Conv}(b,\beta) \to \beta \neq 0 \to \forall n \in \mathbb{N}\,(b_{n} \neq 0) \to \varepsilon \in \mathbb{R} \to 0 < \varepsilon \to \varphi_{62} \to \varphi_{228}$
\item $\varphi_{383} :\equiv \mathrm{Conv}(b,\beta) \to \beta \neq 0 \to \forall n \in \mathbb{N}\,(b_{n} \neq 0) \to \varepsilon \in \mathbb{R} \to 0 < \varepsilon \to \varphi_{62} \to \varphi_{229}$
\item $\varphi_{384} :\equiv \mathrm{Conv}(b,\beta) \to \beta \neq 0 \to \forall n \in \mathbb{N}\,(b_{n} \neq 0) \to \varepsilon \in \mathbb{R} \to \varphi_{307}$
\item $\varphi_{385} :\equiv \forall w\,(\varepsilon \cdot (\lvert \beta \rvert \cdot 2^{-1} \cdot \lvert \beta \rvert) \cdot ((\lvert \beta \rvert \cdot 2^{-1})^{-1} \cdot \lvert \beta \rvert^{-1}) = w \to \varphi_{232} \to \varepsilon \cdot (\lvert \beta \rvert \cdot 2^{-1} \cdot \lvert \beta \rvert) \cdot (\lvert b_{n} \rvert^{-1} \cdot \lvert \beta \rvert^{-1}) < w)$
\item $\varphi_{386} :\equiv \mathrm{Conv}(b,\beta) \to \beta \neq 0 \to \forall n \in \mathbb{N}\,(b_{n} \neq 0) \to \varepsilon \in \mathbb{R} \to \varphi_{309}$
\item $\varphi_{387} :\equiv \varphi_{43} \to \varphi_{232} \to \varepsilon \cdot (\lvert \beta \rvert \cdot 2^{-1} \cdot \lvert \beta \rvert) \cdot (\lvert b_{n} \rvert^{-1} \cdot \lvert \beta \rvert^{-1}) < \varepsilon$
\item $\varphi_{388} :\equiv (\mathrm{Conv}(b,\beta) \to \beta \neq 0 \to \forall n \in \mathbb{N}\,(b_{n} \neq 0) \to \beta \in \mathbb{R}) \to \varphi_{310}$
\item $\varphi_{389} :\equiv \mathrm{Conv}(b,\beta) \to \beta \neq 0 \to \forall n \in \mathbb{N}\,(b_{n} \neq 0) \to \varepsilon \in \mathbb{R} \to 0 < \varepsilon \to \varphi_{62} \to \varphi_{246}$
\item $\varphi_{390} :\equiv \mathrm{Conv}(b,\beta) \to \beta \neq 0 \to \forall n \in \mathbb{N}\,(b_{n} \neq 0) \to \varepsilon \in \mathbb{R} \to 0 < \varepsilon \to \varphi_{62} \to \varphi_{249}$
\item $\varphi_{391} :\equiv \mathrm{Conv}(b,\beta) \to \beta \neq 0 \to \forall n \in \mathbb{N}\,(b_{n} \neq 0) \to \varepsilon \in \mathbb{R} \to 0 < \varepsilon \to \varphi_{62} \to \varphi_{250}$
\item $\varphi_{392} :\equiv \mathrm{Conv}(b,\beta) \to \beta \neq 0 \to \forall n \in \mathbb{N}\,(b_{n} \neq 0) \to \varepsilon \in \mathbb{R} \to \varphi_{316}$
\item $\varphi_{393} :\equiv \mathrm{Conv}(b,\beta) \to \beta \neq 0 \to \forall n \in \mathbb{N}\,(b_{n} \neq 0) \to \varepsilon \in \mathbb{R} \to 0 < \varepsilon \to \varphi_{62} \to \varphi_{255}$
\item $\varphi_{394} :\equiv (\mathrm{Conv}(b,\beta) \to \beta \neq 0 \to \forall n \in \mathbb{N}\,(b_{n} \neq 0) \to \beta \in \mathbb{R}) \to \varphi_{328}$
\item $\varphi_{395} :\equiv \mathrm{Conv}(b,\beta) \to \beta \neq 0 \to \forall n \in \mathbb{N}\,(b_{n} \neq 0) \to \varepsilon \in \mathbb{R} \to 0 < \varepsilon \to \varphi_{62} \to \varphi_{257}$
\item $\varphi_{396} :\equiv \mathrm{Conv}(b,\beta) \to \beta \neq 0 \to \forall n \in \mathbb{N}\,(b_{n} \neq 0) \to \varepsilon \in \mathbb{R} \to \varphi_{323}$
\item $\varphi_{397} :\equiv \mathrm{Conv}(b,\beta) \to \beta \neq 0 \to \forall n \in \mathbb{N}\,(b_{n} \neq 0) \to \varepsilon \in \mathbb{R} \to \varphi_{324}$
\item $\varphi_{398} :\equiv \mathrm{Conv}(b,\beta) \to \beta \neq 0 \to \forall n \in \mathbb{N}\,(b_{n} \neq 0) \to \varepsilon \in \mathbb{R} \to \varphi_{325}$
\item $\varphi_{399} :\equiv \mathrm{Conv}(b,\beta) \to \beta \neq 0 \to \forall n \in \mathbb{N}\,(b_{n} \neq 0) \to \varepsilon \in \mathbb{R} \to \varphi_{326}$
\item $\varphi_{400} :\equiv \varepsilon \cdot (\lvert \beta \rvert \cdot 2^{-1} \cdot \lvert \beta \rvert) \in \mathbb{R} \to \lvert b_{n} \rvert^{-1} \cdot \lvert \beta \rvert^{-1} \in \mathbb{R} \to \varphi_{298}$
\item $\varphi_{401} :\equiv (\lvert \beta \rvert \cdot 2^{-1} \cdot \lvert \beta \rvert = 0 \to 0 = \lvert \beta \rvert \cdot 2^{-1} \cdot \lvert \beta \rvert) \to \varphi_{339}$
\item $\varphi_{402} :\equiv \mathrm{Conv}(b,\beta) \to \beta \neq 0 \to \forall n \in \mathbb{N}\,(b_{n} \neq 0) \to \varepsilon \in \mathbb{R} \to 0 < \varepsilon \to \varphi_{62} \to \varphi_{265}$
\item $\varphi_{403} :\equiv (\mathrm{Conv}(b,\beta) \to \beta \neq 0 \to \forall n \in \mathbb{N}\,(b_{n} \neq 0) \to \beta \in \mathbb{R}) \to (\mathrm{Conv}(b,\beta) \to \beta \neq 0 \to \beta \neq 0) \to \varphi_{313}$
\item $\varphi_{404} :\equiv \mathrm{Conv}(b,\beta) \to \beta \neq 0 \to \forall n \in \mathbb{N}\,(b_{n} \neq 0) \to \varepsilon \in \mathbb{R} \to 0 < \varepsilon \to \varphi_{62} \to \varphi_{271}$
\item $\varphi_{405} :\equiv \mathrm{Conv}(b,\beta) \to \beta \neq 0 \to \forall n \in \mathbb{N}\,(b_{n} \neq 0) \to \varepsilon \in \mathbb{R} \to 0 < \varepsilon \to \varphi_{62} \to \varphi_{275}$
\item $\varphi_{406} :\equiv \mathrm{Conv}(b,\beta) \to \beta \neq 0 \to \forall n \in \mathbb{N}\,(b_{n} \neq 0) \to \varepsilon \in \mathbb{R} \to 0 < \varepsilon \to \varphi_{62} \to \varphi_{276}$
\item $\varphi_{407} :\equiv \mathrm{Conv}(b,\beta) \to \beta \neq 0 \to \forall n \in \mathbb{N}\,(b_{n} \neq 0) \to \varepsilon \in \mathbb{R} \to 0 < \varepsilon \to \varphi_{62} \to \varphi_{277}$
\item $\varphi_{408} :\equiv (\mathrm{Conv}(b,\beta) \to \beta \neq 0 \to \forall n \in \mathbb{N}\,(b_{n} \neq 0) \to \beta \in \mathbb{R}) \to (\mathrm{Conv}(b,\beta) \to \beta \neq 0 \to \beta \neq 0) \to \varphi_{331}$
\item $\varphi_{409} :\equiv (\lvert \beta \rvert \cdot 2^{-1})^{-1} \cdot \lvert \beta \rvert^{-1} \in \mathbb{R} \to \varepsilon \cdot (\lvert \beta \rvert \cdot 2^{-1} \cdot \lvert \beta \rvert) \in \mathbb{R} \to \varphi_{341}$
\item $\varphi_{410} :\equiv \mathrm{Conv}(b,\beta) \to \beta \neq 0 \to \forall n \in \mathbb{N}\,(b_{n} \neq 0) \to \varepsilon \in \mathbb{R} \to 0 < \varepsilon \to \varphi_{62} \to \varphi_{282}$
\item $\varphi_{411} :\equiv \mathrm{Conv}(b,\beta) \to \beta \neq 0 \to \forall n \in \mathbb{N}\,(b_{n} \neq 0) \to \varepsilon \in \mathbb{R} \to 0 < \varepsilon \to \varphi_{62} \to \varphi_{286}$
\item $\varphi_{412} :\equiv \mathrm{Conv}(b,\beta) \to \beta \neq 0 \to \forall n \in \mathbb{N}\,(b_{n} \neq 0) \to \varepsilon \in \mathbb{R} \to 0 < \varepsilon \to \varphi_{62} \to \varphi_{289}$
\item $\varphi_{413} :\equiv \mathrm{Conv}(b,\beta) \to \beta \neq 0 \to \forall n \in \mathbb{N}\,(b_{n} \neq 0) \to \varepsilon \in \mathbb{R} \to 0 < \varepsilon \to \varphi_{62} \to \varphi_{290}$
\item $\varphi_{414} :\equiv \mathrm{Conv}(b,\beta) \to \beta \neq 0 \to \forall n \in \mathbb{N}\,(b_{n} \neq 0) \to \varepsilon \in \mathbb{R} \to 0 < \varepsilon \to \varphi_{62} \to \varphi_{294}$
\item $\varphi_{415} :\equiv (\mathrm{Conv}(b,\beta) \to \beta \neq 0 \to \varphi_{87}) \to \varphi_{56} \to \varphi_{55} \to \mathrm{Conv}(b,\beta) \to \beta \neq 0 \to \varphi_{96}$
\item $\varphi_{416} :\equiv (0 < \lvert \beta \rvert \cdot 2^{-1} \cdot \lvert \beta \rvert \to 0 \neq \lvert \beta \rvert \cdot 2^{-1} \cdot \lvert \beta \rvert) \to \varphi_{55} \to \varphi_{338}$
\item $\varphi_{417} :\equiv (\mathrm{Conv}(b,\beta) \to \beta \neq 0 \to \varphi_{113}) \to \varphi_{22} \to (\mathrm{Conv}(b,\beta) \to \beta \neq 0 \to \varphi_{96}) \to \varphi_{247}$
\item $\varphi_{418} :\equiv \forall w\,(\lvert b_{n}^{-1} \rvert = w \to \varphi_{165} \to \varphi_{109}) \to \lvert b_{n}^{-1} \rvert = \lvert b_{n} \rvert^{-1} \to \varphi_{165} \to \varphi_{168}$
\item $\varphi_{419} :\equiv \forall n\,(\varphi_{285} \to n \in \mathbb{N} \to N \leq n \to \lvert b^{\ast}(n) - \beta^{-1} \rvert < \varepsilon) \to \forall n\,(\varphi_{285}) \to \forall n \in \mathbb{N}\,(N \leq n \to \lvert b^{\ast}(n) - \beta^{-1} \rvert < \varepsilon)$
\item $\varphi_{420} :\equiv \forall z\,\forall w\,(z = w \to \varphi_{177} \to \varphi_{176}) \to \forall w\,((\lvert \beta \rvert \cdot 2^{-1})^{-1} \cdot \lvert \beta \rvert^{-1} = w \to \varphi_{234} \to \varphi_{176})$
\item $\varphi_{421} :\equiv (\mathrm{Conv}(b,\beta) \to \beta \neq 0 \to \forall n \in \mathbb{N}\,(b_{n} \neq 0) \to \beta \in \mathbb{R}) \to \varphi_{314} \to \varphi_{329}$
\item $\varphi_{422} :\equiv \lvert b_{n}^{-1} - \beta^{-1} \rvert = \lvert b_{n}^{-1} - \beta^{-1} \rvert \to (b_{n}^{-1} - \beta^{-1} = (\beta - b_{n}) \cdot (b_{n}^{-1} \cdot \beta^{-1}) \to \varphi_{212}) \to \varphi_{386}$
\item $\varphi_{423} :\equiv \varphi_{370} \to \varphi_{300} \to \neg (N \in \mathbb{N} \land \forall n \in \mathbb{N}\,(N \leq n \to \lvert b_{n} - \beta \rvert < d_{0}))$
\item $\varphi_{424} :\equiv (\mathrm{Conv}(b,\beta) \to \mathrm{Conv}(b,\beta)) \to (\mathrm{Conv}(b,\beta) \to \varphi_{252}) \to (\mathrm{Conv}(b,\beta) \to \varphi_{242}) \to \varphi_{287}$
\item $\varphi_{425} :\equiv (\mathrm{Conv}(b,\beta) \to \beta \neq 0 \to \forall n \in \mathbb{N}\,(b_{n} \neq 0) \to \beta \in \mathbb{R}) \to \varphi_{314} \to \varphi_{349}$
\item $\varphi_{426} :\equiv \varphi_{335} \to (b^{\ast}(n) = b_{n}^{-1} \to b_{n}^{-1} = b^{\ast}(n)) \to \varphi_{336}$
\item $\varphi_{427} :\equiv (\lvert \beta \rvert \in \mathbb{R} \to 2^{-1} \in \mathbb{R} \to \lvert \beta \rvert \cdot 2^{-1} \in \mathbb{R}) \to \varphi_{310} \to 2^{-1} \in \mathbb{R} \to \varphi_{332}$
\item $\varphi_{428} :\equiv (\mathrm{Conv}(b,\beta) \to \beta \neq 0 \to \forall n \in \mathbb{N}\,(b_{n} \neq 0) \to \beta \in \mathbb{R}) \to \varphi_{351} \to \varphi_{342}$
\item $\varphi_{429} :\equiv (\varepsilon \in \mathbb{R} \to \lvert \beta \rvert \cdot 2^{-1} \cdot \lvert \beta \rvert \in \mathbb{R} \to 0 < \varepsilon \to \varphi_{36}) \to \varphi_{39} \to \varphi_{415}$
\item $\varphi_{430} :\equiv (\mathrm{Conv}(b,\beta) \to \beta \neq 0 \to \forall n \in \mathbb{N}\,(b_{n} \neq 0) \to \beta \in \mathbb{R}) \to \varphi_{306} \to (\mathrm{Conv}(b,\beta) \to \beta \neq 0 \to \beta \neq 0) \to \varphi_{378}$
\item $\varphi_{431} :\equiv (\forall d\,(\forall d_{0}\,(\neg \varphi_{62}) \to \varphi_{51}) \to \forall d\,\forall d_{0}\,(\neg \varphi_{62}) \to \forall d\,(\varphi_{51})) \to (\forall d_{0}\,(\neg \varphi_{62}) \to \forall d\,\forall d_{0}\,(\neg \varphi_{62})) \to \varphi_{223}$
\item $\varphi_{432} :\equiv \forall w\,((\lvert \beta \rvert \cdot 2^{-1})^{-1} \cdot \lvert \beta \rvert^{-1} = w \to \varphi_{234} \to \varphi_{176}) \to \varphi_{26} \to \varphi_{234} \to \varphi_{226}$
\item $\varphi_{433} :\equiv (\lvert \beta \rvert \cdot 2^{-1} \in \mathbb{R} \to \varphi_{280}) \to (\mathrm{Conv}(b,\beta) \to \beta \neq 0 \to \varphi_{44}) \to \varphi_{417}$
\item $\varphi_{434} :\equiv (\lvert b_{n} \rvert^{-1} \in \mathbb{R} \to \lvert \beta \rvert^{-1} \in \mathbb{R} \to \lvert b_{n} \rvert^{-1} \cdot \lvert \beta \rvert^{-1} \in \mathbb{R}) \to \varphi_{333} \to \varphi_{319} \to \varphi_{358}$
\item $\varphi_{435} :\equiv \varphi_{376} \to \varphi_{133} \to (\lvert b_{n}^{-1} \cdot \beta^{-1} \rvert = \lvert b_{n}^{-1} \rvert \cdot \lvert \beta^{-1} \rvert \to \varphi_{133} \to \varphi_{157}) \to \varphi_{404}$
\item $\varphi_{436} :\equiv (\lvert \beta \rvert \cdot 2^{-1} \in \mathbb{R} \to \lvert b_{n} \rvert \in \mathbb{R} \to \varphi_{143}) \to \varphi_{332} \to \varphi_{320} \to \varphi_{22} \to \varphi_{342} \to \varphi_{360}$
\item $\varphi_{437} :\equiv \forall n\,(\varphi_{285} \to n \in \mathbb{N} \to N \leq n \to \lvert b^{\ast}(n) - \beta^{-1} \rvert < \varepsilon) \to \varphi_{419} \to (\varphi_{285} \to \forall n\,(\varphi_{285})) \to \varphi_{354}$
\item $\varphi_{438} :\equiv (\forall d_{0}\,(\varphi_{227} \to \neg \varphi_{62}) \to \forall d_{0}\,(\varphi_{227}) \to \forall d_{0}\,(\neg \varphi_{62})) \to (\varphi_{227} \to \forall d_{0}\,(\varphi_{227})) \to \forall d\,(\forall d_{0}\,(\neg \varphi_{62}) \to \varphi_{51}) \to \varphi_{431}$
\item $\varphi_{439} :\equiv (\lvert b_{n} \rvert^{-1} \cdot \lvert \beta \rvert^{-1} \in \mathbb{R} \to \varphi_{409}) \to \varphi_{358} \to \varphi_{364} \to \varphi_{361} \to (\mathrm{Conv}(b,\beta) \to \beta \neq 0 \to \varphi_{96}) \to \varphi_{384} \to \varphi_{414}$
\end{abbrevlist}
\end{fullproof}

% 自動生成 (prover/build.js): LimQuot  ノード数 31
\begin{fullproof}{定理2.9}{$\forall a\,\forall \alpha\,\forall b\,\forall \beta\,(\mathrm{Conv}(a,\alpha) \to \mathrm{Conv}(b,\beta) \to \beta \neq 0 \to \forall n \in \mathbb{N}\,(b_{n} \neq 0) \to \mathrm{Conv}(a \otimes b^{\ast},\alpha \cdot \beta^{-1}))$}{31}\phantomsection\label{tree:LimQuot}
\noindent$\Pi_{1}$:
\begin{ptree}
\AXC{$\forall b\,\forall \beta\,(\varphi_{3})$\lbl{定理2.9(逆数)}}
\AXC{$\forall b\,\forall \beta\,(\varphi_{3}) \to \forall \beta\,(\varphi_{3})$\lbl{A2}}
\BIC{$\forall \beta\,(\varphi_{3})$}
\AXC{$\forall \beta\,(\varphi_{3}) \to \varphi_{3}$\lbl{A2}}
\BIC{$\varphi_{3}$}
\AXC{$\varphi_{3} \to \varphi_{10}$\lbl{A1}}
\BIC{$\varphi_{10}$}
\end{ptree}
\noindent$\Pi_{2}$:
\begin{ptree}
\AXC{$\mathrm{Conv}(a,\alpha) \to \mathrm{Conv}(b,\beta) \to \beta \neq 0 \to \beta \neq 0$\lbl{A1}}
\AXC{$\mathrm{Conv}(a,\alpha) \to \mathrm{Conv}(b,\beta) \to \mathrm{Conv}(b,\beta)$\lbl{A1}}
\AXC{$\stackrel{\Pi_{1}}{\vdots}$}\noLine
\UIC{$\varphi_{10}$}
\BIC{$\varphi_{9}$}
\BIC{$(\mathrm{Conv}(a,\alpha) \to \varphi_{6}) \to \mathrm{Conv}(a,\alpha) \to \varphi_{3}$}
\end{ptree}
\noindent$\Pi_{3}$:
\begin{ptree}
\AXC{$\forall a\,\forall \alpha\,(\varphi_{1})$\lbl{定理2.7}}
\AXC{$\forall a\,\forall \alpha\,(\varphi_{1}) \to \forall \alpha\,(\varphi_{1})$\lbl{A2}}
\BIC{$\forall \alpha\,(\varphi_{1})$}
\AXC{$\forall \alpha\,(\varphi_{1}) \to \varphi_{1}$\lbl{A2}}
\BIC{$\varphi_{1}$}
\AXC{$\varphi_{1} \to \varphi_{2}$\lbl{A2}}
\BIC{$\varphi_{2}$}
\AXC{$\varphi_{2} \to \varphi_{5}$\lbl{A2}}
\BIC{$\varphi_{5}$}
\end{ptree}
\noindent$\Pi_{4}$:
\begin{ptree}
\AXC{$\mathrm{Conv}(a,\alpha) \to \mathrm{Conv}(a,\alpha)$\lbl{A1}}
\AXC{$\stackrel{\Pi_{3}}{\vdots}$}\noLine
\UIC{$\varphi_{5}$}
\AXC{$\varphi_{5} \to (\mathrm{Conv}(a,\alpha) \to \mathrm{Conv}(a,\alpha)) \to \varphi_{8}$\lbl{A1}}
\BIC{$(\mathrm{Conv}(a,\alpha) \to \mathrm{Conv}(a,\alpha)) \to \varphi_{8}$}
\BIC{$\varphi_{8}$}
\end{ptree}
\noindent 主木:
\begin{ptree}
\AXC{$\mathrm{Conv}(a,\alpha) \to \varphi_{6}$\lbl{A1}}
\AXC{$\stackrel{\Pi_{2}}{\vdots}$}\noLine
\UIC{$(\mathrm{Conv}(a,\alpha) \to \varphi_{6}) \to \mathrm{Conv}(a,\alpha) \to \varphi_{3}$}
\BIC{$\mathrm{Conv}(a,\alpha) \to \varphi_{3}$}
\AXC{$\stackrel{\Pi_{4}}{\vdots}$}\noLine
\UIC{$\varphi_{8}$}
\BIC{$\mathrm{Conv}(a,\alpha) \to \mathrm{Conv}(b,\beta) \to \varphi_{4}$}
\UIC{$\forall \beta\,(\mathrm{Conv}(a,\alpha) \to \mathrm{Conv}(b,\beta) \to \varphi_{4})$}
\UIC{$\forall b\,\forall \beta\,(\mathrm{Conv}(a,\alpha) \to \mathrm{Conv}(b,\beta) \to \varphi_{4})$}
\UIC{$\forall \alpha\,\forall b\,\forall \beta\,(\mathrm{Conv}(a,\alpha) \to \mathrm{Conv}(b,\beta) \to \varphi_{4})$}
\UIC{$\varphi_{7}$}
\end{ptree}
\begin{abbrevlist}
\item $\varphi_{1} :\equiv \forall b\,\forall \beta\,(\mathrm{Conv}(a,\alpha) \to \mathrm{Conv}(b,\beta) \to \mathrm{Conv}(a \otimes b,\alpha \cdot \beta))$
\item $\varphi_{2} :\equiv \forall \beta\,(\mathrm{Conv}(a,\alpha) \to \mathrm{Conv}(b^{\ast},\beta) \to \mathrm{Conv}(a \otimes b^{\ast},\alpha \cdot \beta))$
\item $\varphi_{3} :\equiv \mathrm{Conv}(b,\beta) \to \beta \neq 0 \to \forall n \in \mathbb{N}\,(b_{n} \neq 0) \to \mathrm{Conv}(b^{\ast},\beta^{-1})$
\item $\varphi_{4} :\equiv \beta \neq 0 \to \forall n \in \mathbb{N}\,(b_{n} \neq 0) \to \mathrm{Conv}(a \otimes b^{\ast},\alpha \cdot \beta^{-1})$
\item $\varphi_{5} :\equiv \mathrm{Conv}(a,\alpha) \to \mathrm{Conv}(b^{\ast},\beta^{-1}) \to \mathrm{Conv}(a \otimes b^{\ast},\alpha \cdot \beta^{-1})$
\item $\varphi_{6} :\equiv \mathrm{Conv}(b,\beta) \to \beta \neq 0 \to \forall n \in \mathbb{N}\,(b_{n} \neq 0) \to \forall n \in \mathbb{N}\,(b_{n} \neq 0)$
\item $\varphi_{7} :\equiv \forall a\,\forall \alpha\,\forall b\,\forall \beta\,(\mathrm{Conv}(a,\alpha) \to \mathrm{Conv}(b,\beta) \to \varphi_{4})$
\item $\varphi_{8} :\equiv (\mathrm{Conv}(a,\alpha) \to \varphi_{3}) \to \mathrm{Conv}(a,\alpha) \to \mathrm{Conv}(b,\beta) \to \varphi_{4}$
\item $\varphi_{9} :\equiv (\mathrm{Conv}(a,\alpha) \to \mathrm{Conv}(b,\beta) \to \beta \neq 0 \to \beta \neq 0) \to (\mathrm{Conv}(a,\alpha) \to \varphi_{6}) \to \mathrm{Conv}(a,\alpha) \to \varphi_{3}$
\item $\varphi_{10} :\equiv (\mathrm{Conv}(a,\alpha) \to \mathrm{Conv}(b,\beta) \to \mathrm{Conv}(b,\beta)) \to \varphi_{9}$
\end{abbrevlist}
\end{fullproof}

\section{第3章 2節(続き): 部分列, 命題2.13, 定理2.15–2.18}
% 自動生成 (prover/build.js): IdxLt  ノード数 376
\begin{fullproof}{[C16]}{$\forall s\,\forall m\,(\mathrm{SubIdx}(s) \to m \in \mathbb{N} \to \forall n \in \mathbb{N}\,(m < n \to s_{m} < s_{n}))$}{376}\phantomsection\label{tree:IdxLt}
\noindent$\Pi_{1}$:
\begin{ptree}
\AXC{$\mathrm{SubIdx}(s) \to m \in \mathbb{N} \to n \in \mathbb{N} \to n \in \mathbb{N}$\lbl{A1}}
\AXC{$(\mathrm{SubIdx}(s) \to m \in \mathbb{N} \to n \in \mathbb{N} \to n \in \mathbb{N}) \to \varphi_{63}$\lbl{A1}}
\BIC{$\varphi_{63}$}
\end{ptree}
\noindent$\Pi_{2}$:
\begin{ptree}
\AXC{$\varphi_{21}$\lbl{A4}}
\AXC{$\forall y\,(\varphi_{10}) \to \varphi_{27}$\lbl{A3}}
\AXC{$\varphi_{10}$\lbl{A2}}
\UIC{$\forall y\,(\varphi_{10})$}
\AXC{$\forall y\,(\varphi_{10}) \to (\forall y\,(\varphi_{10}) \to \varphi_{27}) \to \varphi_{21} \to \varphi_{16}$\lbl{A1}}
\BIC{$(\forall y\,(\varphi_{10}) \to \varphi_{27}) \to \varphi_{21} \to \varphi_{16}$}
\BIC{$\varphi_{21} \to \varphi_{16}$}
\BIC{$\varphi_{16}$}
\end{ptree}
\noindent$\Pi_{3}$:
\begin{ptree}
\AXC{$\varphi_{22}$\lbl{A4}}
\AXC{$\forall n\,(\varphi_{11}) \to \varphi_{28}$\lbl{A3}}
\AXC{$\varphi_{11}$\lbl{A2}}
\UIC{$\forall n\,(\varphi_{11})$}
\AXC{$\forall n\,(\varphi_{11}) \to (\forall n\,(\varphi_{11}) \to \varphi_{28}) \to \varphi_{22} \to \varphi_{17}$\lbl{A1}}
\BIC{$(\forall n\,(\varphi_{11}) \to \varphi_{28}) \to \varphi_{22} \to \varphi_{17}$}
\BIC{$\varphi_{22} \to \varphi_{17}$}
\BIC{$\varphi_{17}$}
\end{ptree}
\noindent$\Pi_{4}$:
\begin{ptree}
\AXC{$\varphi_{117}$\lbl{A2}}
\AXC{$\mathrm{SubIdx}(s) \to m \in \mathbb{N} \to \varphi_{120}$\lbl{A1}}
\AXC{$\varphi_{170}$\lbl{A1}}
\BIC{$\varphi_{117} \to \mathrm{SubIdx}(s) \to m \in \mathbb{N} \to \varphi_{117}$}
\BIC{$\mathrm{SubIdx}(s) \to m \in \mathbb{N} \to \varphi_{117}$}
\AXC{$(\mathrm{SubIdx}(s) \to m \in \mathbb{N} \to \varphi_{117}) \to \varphi_{145}$\lbl{A1}}
\BIC{$\varphi_{145}$}
\end{ptree}
\noindent$\Pi_{5}$:
\begin{ptree}
\AXC{$\mathrm{SubIdx}(s) \to m \in \mathbb{N} \to \varphi_{68}$\lbl{A1}}
\AXC{$\stackrel{\Pi_{4}}{\vdots}$}\noLine
\UIC{$\varphi_{145}$}
\BIC{$\mathrm{SubIdx}(s) \to m \in \mathbb{N} \to \varphi_{118}$}
\end{ptree}
\noindent$\Pi_{6}$:
\begin{ptree}
\AXC{$\stackrel{\Pi_{5}}{\vdots}$}\noLine
\UIC{$\mathrm{SubIdx}(s) \to m \in \mathbb{N} \to \varphi_{118}$}
\AXC{$\varphi_{146}$\lbl{A1}}
\BIC{$\mathrm{SubIdx}(s) \to m \in \mathbb{N} \to \varphi_{70}$}
\end{ptree}
\noindent$\Pi_{7}$:
\begin{ptree}
\AXC{$\stackrel{\Pi_{6}}{\vdots}$}\noLine
\UIC{$\mathrm{SubIdx}(s) \to m \in \mathbb{N} \to \varphi_{70}$}
\AXC{$\forall n\,(n \in \mathbb{N} \leftrightarrow n \in \omega \land n \neq \emptyset)$\lbl{Def$\mathbb{N}$}}
\AXC{$\varphi_{53}$\lbl{A2}}
\BIC{$n \in \mathbb{N} \leftrightarrow n \in \omega \land n \neq \emptyset$}
\AXC{$\varphi_{169}$\lbl{A1}}
\BIC{$(\mathrm{SubIdx}(s) \to m \in \mathbb{N} \to \varphi_{70}) \to \varphi_{148}$}
\BIC{$\varphi_{148}$}
\end{ptree}
\noindent$\Pi_{8}$:
\begin{ptree}
\AXC{$\mathrm{SubIdx}(s) \to m \in \mathbb{N} \to \varphi_{99}$\lbl{A1}}
\AXC{$\stackrel{\Pi_{7}}{\vdots}$}\noLine
\UIC{$\varphi_{148}$}
\BIC{$\mathrm{SubIdx}(s) \to m \in \mathbb{N} \to \varphi_{96}$}
\end{ptree}
\noindent$\Pi_{9}$:
\begin{ptree}
\AXC{$\mathrm{SubIdx}(s) \to \mathrm{SubIdx}(s)$\lbl{A1}}
\AXC{$(\mathrm{SubIdx}(s) \to \mathrm{SubIdx}(s)) \to \varphi_{98}$\lbl{A1}}
\BIC{$\varphi_{98}$}
\end{ptree}
\noindent$\Pi_{10}$:
\begin{ptree}
\AXC{$\forall k \in \mathbb{N}\,(s_{k} < s_{S(k)}) \to n \in \mathbb{N} \to s_{n} < s_{S(n)}$\lbl{A2}}
\AXC{$\forall s\,(\varphi_{23})$\lbl{DefSubIdx}}
\AXC{$\forall s\,(\varphi_{23}) \to \varphi_{23}$\lbl{A2}}
\BIC{$\varphi_{23}$}
\AXC{$\stackrel{\Pi_{9}}{\vdots}$}\noLine
\UIC{$\varphi_{98}$}
\BIC{$\mathrm{SubIdx}(s) \to m \in \mathbb{N} \to \forall k \in \mathbb{N}\,(s_{k} < s_{S(k)})$}
\AXC{$\varphi_{138}$\lbl{A1}}
\BIC{$\varphi_{115}$}
\BIC{$\mathrm{SubIdx}(s) \to m \in \mathbb{N} \to n \in \mathbb{N} \to s_{n} < s_{S(n)}$}
\end{ptree}
\noindent$\Pi_{11}$:
\begin{ptree}
\AXC{$\stackrel{\Pi_{8}}{\vdots}$}\noLine
\UIC{$\mathrm{SubIdx}(s) \to m \in \mathbb{N} \to \varphi_{96}$}
\AXC{$\stackrel{\Pi_{10}}{\vdots}$}\noLine
\UIC{$\mathrm{SubIdx}(s) \to m \in \mathbb{N} \to n \in \mathbb{N} \to s_{n} < s_{S(n)}$}
\AXC{$\varphi_{167}$\lbl{A1}}
\BIC{$\varphi_{158}$}
\BIC{$\mathrm{SubIdx}(s) \to m \in \mathbb{N} \to \varphi_{123}$}
\end{ptree}
\noindent$\Pi_{12}$:
\begin{ptree}
\AXC{$x = x$\lbl{A5}}
\AXC{$x = y \to x = x \to y = x$\lbl{A6}}
\AXC{$(x = y \to x = x \to y = x) \to x = x \to x = y \to y = x$\lbl{A1}}
\BIC{$x = x \to x = y \to y = x$}
\BIC{$x = y \to y = x$}
\UIC{$\forall y\,(x = y \to y = x)$}
\UIC{$\forall x\,\forall y\,(x = y \to y = x)$}
\end{ptree}
\noindent$\Pi_{13}$:
\begin{ptree}
\AXC{$\stackrel{\Pi_{12}}{\vdots}$}\noLine
\UIC{$\forall x\,\forall y\,(x = y \to y = x)$}
\AXC{$\forall x\,\forall y\,(x = y \to y = x) \to \forall y\,(m = y \to y = m)$\lbl{A2}}
\BIC{$\forall y\,(m = y \to y = m)$}
\AXC{$\forall y\,(m = y \to y = m) \to m = n \to n = m$\lbl{A2}}
\BIC{$m = n \to n = m$}
\end{ptree}
\noindent$\Pi_{14}$:
\begin{ptree}
\AXC{$\stackrel{\Pi_{11}}{\vdots}$}\noLine
\UIC{$\mathrm{SubIdx}(s) \to m \in \mathbb{N} \to \varphi_{123}$}
\AXC{$\stackrel{\Pi_{13}}{\vdots}$}\noLine
\UIC{$m = n \to n = m$}
\AXC{$\mathrm{SubIdx}(s) \to m \in \mathbb{N} \to \varphi_{109}$\lbl{A1}}
\AXC{$\varphi_{159}$\lbl{A1}}
\BIC{$(m = n \to n = m) \to \mathrm{SubIdx}(s) \to m \in \mathbb{N} \to \varphi_{110}$}
\BIC{$\mathrm{SubIdx}(s) \to m \in \mathbb{N} \to \varphi_{110}$}
\AXC{$\varphi_{175}$\lbl{A1}}
\BIC{$(\mathrm{SubIdx}(s) \to m \in \mathbb{N} \to \varphi_{123}) \to \varphi_{140}$}
\BIC{$\varphi_{140}$}
\end{ptree}
\noindent$\Pi_{15}$:
\begin{ptree}
\AXC{$\stackrel{\Pi_{8}}{\vdots}$}\noLine
\UIC{$\mathrm{SubIdx}(s) \to m \in \mathbb{N} \to \varphi_{96}$}
\AXC{$\stackrel{\Pi_{5}}{\vdots}$}\noLine
\UIC{$\mathrm{SubIdx}(s) \to m \in \mathbb{N} \to \varphi_{118}$}
\AXC{$\varphi_{152}$\lbl{A1}}
\BIC{$\mathrm{SubIdx}(s) \to m \in \mathbb{N} \to \varphi_{101}$}
\AXC{$(\mathrm{SubIdx}(s) \to m \in \mathbb{N} \to \varphi_{101}) \to \varphi_{153}$\lbl{A1}}
\BIC{$\varphi_{153}$}
\BIC{$\mathrm{SubIdx}(s) \to m \in \mathbb{N} \to \varphi_{121}$}
\end{ptree}
\noindent$\Pi_{16}$:
\begin{ptree}
\AXC{$\mathrm{SubIdx}(s) \to m \in \mathbb{N} \to \varphi_{119}$\lbl{A1}}
\AXC{$\stackrel{\Pi_{15}}{\vdots}$}\noLine
\UIC{$\mathrm{SubIdx}(s) \to m \in \mathbb{N} \to \varphi_{121}$}
\AXC{$(\mathrm{SubIdx}(s) \to m \in \mathbb{N} \to \varphi_{121}) \to \varphi_{161}$\lbl{A1}}
\BIC{$\varphi_{161}$}
\BIC{$\mathrm{SubIdx}(s) \to m \in \mathbb{N} \to \varphi_{127}$}
\end{ptree}
\noindent$\Pi_{17}$:
\begin{ptree}
\AXC{$\forall x\,\forall y\,\forall z\,(x < y \to y < z \to x < z)$\lbl{[O2]}}
\AXC{$\varphi_{54}$\lbl{A2}}
\BIC{$\forall y\,\forall z\,(s_{m} < y \to y < z \to s_{m} < z)$}
\AXC{$\varphi_{76}$\lbl{A2}}
\BIC{$\forall z\,(s_{m} < s_{n} \to s_{n} < z \to s_{m} < z)$}
\AXC{$\varphi_{104}$\lbl{A2}}
\BIC{$s_{m} < s_{n} \to s_{n} < s_{S(n)} \to s_{m} < s_{S(n)}$}
\AXC{$\varphi_{179}$\lbl{A1}}
\BIC{$(\mathrm{SubIdx}(s) \to m \in \mathbb{N} \to \varphi_{127}) \to \varphi_{165}$}
\end{ptree}
\noindent$\Pi_{18}$:
\begin{ptree}
\AXC{$\stackrel{\Pi_{11}}{\vdots}$}\noLine
\UIC{$\mathrm{SubIdx}(s) \to m \in \mathbb{N} \to \varphi_{123}$}
\AXC{$\stackrel{\Pi_{16}}{\vdots}$}\noLine
\UIC{$\mathrm{SubIdx}(s) \to m \in \mathbb{N} \to \varphi_{127}$}
\AXC{$\stackrel{\Pi_{17}}{\vdots}$}\noLine
\UIC{$(\mathrm{SubIdx}(s) \to m \in \mathbb{N} \to \varphi_{127}) \to \varphi_{165}$}
\BIC{$\varphi_{165}$}
\BIC{$\mathrm{SubIdx}(s) \to m \in \mathbb{N} \to \varphi_{128}$}
\end{ptree}
\noindent$\Pi_{19}$:
\begin{ptree}
\AXC{$\mathrm{SubIdx}(s) \to m \in \mathbb{N} \to \varphi_{111}$\lbl{A1}}
\AXC{$(\mathrm{SubIdx}(s) \to m \in \mathbb{N} \to \varphi_{111}) \to \varphi_{134}$\lbl{A1}}
\BIC{$\varphi_{134}$}
\end{ptree}
\noindent$\Pi_{20}$:
\begin{ptree}
\AXC{$\forall n \in \mathbb{N}\,(n \in \omega)$\lbl{[N7']}}
\AXC{$\forall n \in \mathbb{N}\,(n \in \omega) \to n \in \mathbb{N} \to n \in \omega$\lbl{A2}}
\BIC{$n \in \mathbb{N} \to n \in \omega$}
\AXC{$(n \in \mathbb{N} \to n \in \omega) \to \varphi_{150}$\lbl{A1}}
\BIC{$\varphi_{150}$}
\end{ptree}
\noindent$\Pi_{21}$:
\begin{ptree}
\AXC{$\forall n \in \mathbb{N}\,(n \in \omega)$\lbl{[N7']}}
\AXC{$\forall n \in \mathbb{N}\,(n \in \omega) \to m \in \mathbb{N} \to m \in \omega$\lbl{A2}}
\BIC{$m \in \mathbb{N} \to m \in \omega$}
\AXC{$(m \in \mathbb{N} \to m \in \omega) \to \varphi_{135}$\lbl{A1}}
\BIC{$\varphi_{135}$}
\end{ptree}
\noindent$\Pi_{22}$:
\begin{ptree}
\AXC{$\varphi_{14}$\lbl{[N9]}}
\AXC{$\varphi_{82}$\lbl{A2}}
\BIC{$\forall n\,(m \in \omega \to n \in \omega \to m \leq S(n) \to m \leq n \lor m = S(n))$}
\AXC{$\varphi_{78}$\lbl{A2}}
\BIC{$m \in \omega \to n \in \omega \to m \leq S(n) \to m \leq n \lor m = S(n)$}
\AXC{$(m \in \omega \to n \in \omega \to m \leq S(n) \to m \leq n \lor m = S(n)) \to \varphi_{180}$\lbl{A1}}
\BIC{$\varphi_{180}$}
\end{ptree}
\noindent$\Pi_{23}$:
\begin{ptree}
\AXC{$\stackrel{\Pi_{8}}{\vdots}$}\noLine
\UIC{$\mathrm{SubIdx}(s) \to m \in \mathbb{N} \to \varphi_{96}$}
\AXC{$\stackrel{\Pi_{20}}{\vdots}$}\noLine
\UIC{$\varphi_{150}$}
\BIC{$\mathrm{SubIdx}(s) \to m \in \mathbb{N} \to \varphi_{113}$}
\AXC{$\mathrm{SubIdx}(s) \to m \in \mathbb{N} \to m \in \mathbb{N}$\lbl{A1}}
\AXC{$\stackrel{\Pi_{21}}{\vdots}$}\noLine
\UIC{$\varphi_{135}$}
\BIC{$\mathrm{SubIdx}(s) \to m \in \mathbb{N} \to \varphi_{114}$}
\AXC{$\stackrel{\Pi_{22}}{\vdots}$}\noLine
\UIC{$\varphi_{180}$}
\BIC{$(\mathrm{SubIdx}(s) \to m \in \mathbb{N} \to \varphi_{113}) \to \varphi_{156}$}
\BIC{$\varphi_{156}$}
\end{ptree}
\noindent$\Pi_{24}$:
\begin{ptree}
\AXC{$\mathrm{SubIdx}(s) \to m \in \mathbb{N} \to \varphi_{111}$\lbl{A1}}
\AXC{$\varphi_{154}$\lbl{A1}}
\BIC{$\mathrm{SubIdx}(s) \to m \in \mathbb{N} \to \varphi_{106}$}
\AXC{$\stackrel{\Pi_{23}}{\vdots}$}\noLine
\UIC{$\varphi_{156}$}
\BIC{$\mathrm{SubIdx}(s) \to m \in \mathbb{N} \to \varphi_{116}$}
\AXC{$(\mathrm{SubIdx}(s) \to m \in \mathbb{N} \to \varphi_{116}) \to \varphi_{181}$\lbl{A1}}
\BIC{$\varphi_{181}$}
\end{ptree}
\noindent$\Pi_{25}$:
\begin{ptree}
\AXC{$\stackrel{\Pi_{18}}{\vdots}$}\noLine
\UIC{$\mathrm{SubIdx}(s) \to m \in \mathbb{N} \to \varphi_{128}$}
\AXC{$\forall x\,\forall y\,(x < y \leftrightarrow x \subseteq y \land x \neq y)$\lbl{Def$<$}}
\AXC{$\varphi_{24}$\lbl{A2}}
\BIC{$\forall y\,(m < y \leftrightarrow m \subseteq y \land m \neq y)$}
\AXC{$\varphi_{33}$\lbl{A2}}
\BIC{$m < S(n) \leftrightarrow m \subseteq S(n) \land m \neq S(n)$}
\AXC{$\stackrel{\Pi_{19}}{\vdots}$}\noLine
\UIC{$\varphi_{134}$}
\BIC{$\mathrm{SubIdx}(s) \to m \in \mathbb{N} \to \varphi_{108}$}
\AXC{$\stackrel{\Pi_{24}}{\vdots}$}\noLine
\UIC{$\varphi_{181}$}
\BIC{$(\mathrm{SubIdx}(s) \to m \in \mathbb{N} \to \varphi_{128}) \to \varphi_{164}$}
\BIC{$\varphi_{164}$}
\end{ptree}
\noindent$\Pi_{26}$:
\begin{ptree}
\AXC{$y = z \to x = y \to x = z$\lbl{A6}}
\AXC{$(y = z \to x = y \to x = z) \to x = y \to y = z \to x = z$\lbl{A1}}
\BIC{$x = y \to y = z \to x = z$}
\UIC{$\forall z\,(x = y \to y = z \to x = z)$}
\UIC{$\forall y\,\forall z\,(x = y \to y = z \to x = z)$}
\UIC{$\forall x\,\forall y\,\forall z\,(x = y \to y = z \to x = z)$}
\AXC{$\varphi_{12}$\lbl{A2}}
\BIC{$\forall y\,\forall z\,(S(n) = y \to y = z \to S(n) = z)$}
\AXC{$\varphi_{44}$\lbl{A2}}
\BIC{$\forall z\,(S(n) = S(\emptyset) \to S(\emptyset) = z \to S(n) = z)$}
\end{ptree}
\noindent$\Pi_{27}$:
\begin{ptree}
\AXC{$\stackrel{\Pi_{26}}{\vdots}$}\noLine
\UIC{$\forall z\,(S(n) = S(\emptyset) \to S(\emptyset) = z \to S(n) = z)$}
\AXC{$\forall z\,(S(n) = S(\emptyset) \to S(\emptyset) = z \to S(n) = z) \to \varphi_{8}$\lbl{A2}}
\BIC{$\varphi_{8}$}
\end{ptree}
\noindent$\Pi_{28}$:
\begin{ptree}
\AXC{$\stackrel{\Pi_{12}}{\vdots}$}\noLine
\UIC{$\forall x\,\forall y\,(x = y \to y = x)$}
\AXC{$\varphi_{1}$\lbl{A2}}
\BIC{$\forall y\,(1_{\omega} = y \to y = 1_{\omega})$}
\AXC{$\varphi_{15}$\lbl{A2}}
\BIC{$1_{\omega} = S(\emptyset) \to S(\emptyset) = 1_{\omega}$}
\AXC{$1_{\omega} = S(\emptyset)$\lbl{Def$1_\omega$}}
\AXC{$1_{\omega} = S(\emptyset) \to \varphi_{13}$\lbl{A1}}
\BIC{$\varphi_{13}$}
\BIC{$S(\emptyset) = 1_{\omega}$}
\end{ptree}
\noindent$\Pi_{29}$:
\begin{ptree}
\AXC{$\mathrm{SubIdx}(s) \to m \in \mathbb{N} \to \varphi_{97}$\lbl{A1}}
\AXC{$(\mathrm{SubIdx}(s) \to m \in \mathbb{N} \to \varphi_{97}) \to S(n) = S(n) \to \varphi_{136}$\lbl{A1}}
\BIC{$S(n) = S(n) \to \varphi_{136}$}
\end{ptree}
\noindent$\Pi_{30}$:
\begin{ptree}
\AXC{$x = x$\lbl{A5}}
\UIC{$\forall x\,(x = x)$}
\AXC{$\forall x\,(x = x) \to S(n) = S(n)$\lbl{A2}}
\BIC{$S(n) = S(n)$}
\AXC{$\stackrel{\Pi_{29}}{\vdots}$}\noLine
\UIC{$S(n) = S(n) \to \varphi_{136}$}
\BIC{$\varphi_{136}$}
\end{ptree}
\noindent$\Pi_{31}$:
\begin{ptree}
\AXC{$z = w \to S(n) = S(z) \to S(n) = S(w)$\lbl{A6}}
\UIC{$\forall w\,(z = w \to S(n) = S(z) \to S(n) = S(w))$}
\UIC{$\forall z\,\forall w\,(z = w \to S(n) = S(z) \to S(n) = S(w))$}
\AXC{$\varphi_{35}$\lbl{A2}}
\BIC{$\forall w\,(n = w \to S(n) = S(n) \to S(n) = S(w))$}
\AXC{$\varphi_{45}$\lbl{A2}}
\BIC{$n = \emptyset \to S(n) = S(n) \to S(n) = S(\emptyset)$}
\AXC{$\stackrel{\Pi_{30}}{\vdots}$}\noLine
\UIC{$\varphi_{136}$}
\BIC{$\mathrm{SubIdx}(s) \to m \in \mathbb{N} \to \varphi_{103}$}
\AXC{$\varphi_{166}$\lbl{A1}}
\BIC{$S(\emptyset) = 1_{\omega} \to \varphi_{8} \to \mathrm{SubIdx}(s) \to m \in \mathbb{N} \to \varphi_{94}$}
\end{ptree}
\noindent$\Pi_{32}$:
\begin{ptree}
\AXC{$\stackrel{\Pi_{27}}{\vdots}$}\noLine
\UIC{$\varphi_{8}$}
\AXC{$\stackrel{\Pi_{28}}{\vdots}$}\noLine
\UIC{$S(\emptyset) = 1_{\omega}$}
\AXC{$\stackrel{\Pi_{31}}{\vdots}$}\noLine
\UIC{$S(\emptyset) = 1_{\omega} \to \varphi_{8} \to \mathrm{SubIdx}(s) \to m \in \mathbb{N} \to \varphi_{94}$}
\BIC{$\varphi_{8} \to \mathrm{SubIdx}(s) \to m \in \mathbb{N} \to \varphi_{94}$}
\BIC{$\mathrm{SubIdx}(s) \to m \in \mathbb{N} \to \varphi_{94}$}
\end{ptree}
\noindent$\Pi_{33}$:
\begin{ptree}
\AXC{$\stackrel{\Pi_{32}}{\vdots}$}\noLine
\UIC{$\mathrm{SubIdx}(s) \to m \in \mathbb{N} \to \varphi_{94}$}
\AXC{$(\mathrm{SubIdx}(s) \to m \in \mathbb{N} \to \varphi_{94}) \to \varphi_{130}$\lbl{A1}}
\BIC{$\varphi_{130}$}
\end{ptree}
\noindent$\Pi_{34}$:
\begin{ptree}
\AXC{$\stackrel{\Pi_{12}}{\vdots}$}\noLine
\UIC{$\forall x\,\forall y\,(x = y \to y = x)$}
\AXC{$\varphi_{2}$\lbl{A2}}
\BIC{$\forall y\,(S(n) = y \to y = S(n))$}
\AXC{$\varphi_{4}$\lbl{A2}}
\BIC{$S(n) = 1_{\omega} \to 1_{\omega} = S(n)$}
\AXC{$\stackrel{\Pi_{33}}{\vdots}$}\noLine
\UIC{$\varphi_{130}$}
\BIC{$\mathrm{SubIdx}(s) \to m \in \mathbb{N} \to \varphi_{95}$}
\end{ptree}
\noindent$\Pi_{35}$:
\begin{ptree}
\AXC{$\stackrel{\Pi_{34}}{\vdots}$}\noLine
\UIC{$\mathrm{SubIdx}(s) \to m \in \mathbb{N} \to \varphi_{95}$}
\AXC{$(\mathrm{SubIdx}(s) \to m \in \mathbb{N} \to \varphi_{95}) \to \varphi_{131}$\lbl{A1}}
\BIC{$\varphi_{131}$}
\end{ptree}
\noindent$\Pi_{36}$:
\begin{ptree}
\AXC{$\forall n \in \mathbb{N}\,(1_{\omega} \leq n)$\lbl{[N12]}}
\AXC{$\forall n \in \mathbb{N}\,(1_{\omega} \leq n) \to m \in \mathbb{N} \to 1_{\omega} \leq m$\lbl{A2}}
\BIC{$m \in \mathbb{N} \to 1_{\omega} \leq m$}
\AXC{$(m \in \mathbb{N} \to 1_{\omega} \leq m) \to \varphi_{50}$\lbl{A1}}
\BIC{$\varphi_{50}$}
\end{ptree}
\noindent$\Pi_{37}$:
\begin{ptree}
\AXC{$\forall x\,\forall y\,(x \leq y \to \neg y < x)$\lbl{[O7]}}
\AXC{$\varphi_{5}$\lbl{A2}}
\BIC{$\forall y\,(1_{\omega} \leq y \to \neg y < 1_{\omega})$}
\AXC{$\varphi_{6}$\lbl{A2}}
\BIC{$1_{\omega} \leq m \to \neg m < 1_{\omega}$}
\AXC{$(1_{\omega} \leq m \to \neg m < 1_{\omega}) \to \varphi_{66}$\lbl{A1}}
\BIC{$\varphi_{66}$}
\end{ptree}
\noindent$\Pi_{38}$:
\begin{ptree}
\AXC{$\mathrm{SubIdx}(s) \to m \in \mathbb{N} \to m \in \mathbb{N}$\lbl{A1}}
\AXC{$\stackrel{\Pi_{36}}{\vdots}$}\noLine
\UIC{$\varphi_{50}$}
\BIC{$\mathrm{SubIdx}(s) \to m \in \mathbb{N} \to m < 1_{\omega} \to 1_{\omega} \leq m$}
\AXC{$\stackrel{\Pi_{37}}{\vdots}$}\noLine
\UIC{$\varphi_{66}$}
\BIC{$\varphi_{7}$}
\AXC{$\varphi_{7} \to \varphi_{80}$\lbl{A1}}
\BIC{$\varphi_{80}$}
\end{ptree}
\noindent$\Pi_{39}$:
\begin{ptree}
\AXC{$\stackrel{\Pi_{34}}{\vdots}$}\noLine
\UIC{$\mathrm{SubIdx}(s) \to m \in \mathbb{N} \to \varphi_{95}$}
\AXC{$(\mathrm{SubIdx}(s) \to m \in \mathbb{N} \to \varphi_{95}) \to \varphi_{176}$\lbl{A1}}
\BIC{$\varphi_{176}$}
\end{ptree}
\noindent$\Pi_{40}$:
\begin{ptree}
\AXC{$\mathrm{SubIdx}(s) \to m \in \mathbb{N} \to m < 1_{\omega} \to m < 1_{\omega}$\lbl{A1}}
\AXC{$\stackrel{\Pi_{38}}{\vdots}$}\noLine
\UIC{$\varphi_{80}$}
\BIC{$\varphi_{19}$}
\AXC{$\stackrel{\Pi_{39}}{\vdots}$}\noLine
\UIC{$\varphi_{176}$}
\BIC{$(1_{\omega} = S(n) \to m < 1_{\omega} \to m < S(n)) \to \varphi_{171}$}
\end{ptree}
\noindent$\Pi_{41}$:
\begin{ptree}
\AXC{$z = w \to m < z \to m < w$\lbl{A6}}
\UIC{$\forall w\,(z = w \to m < z \to m < w)$}
\UIC{$\forall z\,\forall w\,(z = w \to m < z \to m < w)$}
\AXC{$\varphi_{26}$\lbl{A2}}
\BIC{$\forall w\,(1_{\omega} = w \to m < 1_{\omega} \to m < w)$}
\AXC{$\varphi_{37}$\lbl{A2}}
\BIC{$1_{\omega} = S(n) \to m < 1_{\omega} \to m < S(n)$}
\AXC{$\stackrel{\Pi_{40}}{\vdots}$}\noLine
\UIC{$(1_{\omega} = S(n) \to m < 1_{\omega} \to m < S(n)) \to \varphi_{171}$}
\BIC{$\varphi_{171}$}
\end{ptree}
\noindent$\Pi_{42}$:
\begin{ptree}
\AXC{$z = w \to s_{m} < s_{z} \to s_{m} < s_{w}$\lbl{A6}}
\UIC{$\forall w\,(z = w \to s_{m} < s_{z} \to s_{m} < s_{w})$}
\UIC{$\forall z\,\forall w\,(z = w \to s_{m} < s_{z} \to s_{m} < s_{w})$}
\AXC{$\varphi_{88}$\lbl{A2}}
\BIC{$\forall w\,(1_{\omega} = w \to s_{m} < s_{1_{\omega}} \to s_{m} < s_{w})$}
\AXC{$\varphi_{93}$\lbl{A2}}
\BIC{$1_{\omega} = S(n) \to s_{m} < s_{1_{\omega}} \to s_{m} < s_{S(n)}$}
\AXC{$\stackrel{\Pi_{41}}{\vdots}$}\noLine
\UIC{$\varphi_{171}$}
\BIC{$(S(n) = 1_{\omega} \to m < S(n) \to m < 1_{\omega}) \to \varphi_{168}$}
\end{ptree}
\noindent$\Pi_{43}$:
\begin{ptree}
\AXC{$z = w \to m < z \to m < w$\lbl{A6}}
\UIC{$\forall w\,(z = w \to m < z \to m < w)$}
\UIC{$\forall z\,\forall w\,(z = w \to m < z \to m < w)$}
\AXC{$\varphi_{25}$\lbl{A2}}
\BIC{$\forall w\,(S(n) = w \to m < S(n) \to m < w)$}
\AXC{$\varphi_{36}$\lbl{A2}}
\BIC{$S(n) = 1_{\omega} \to m < S(n) \to m < 1_{\omega}$}
\AXC{$\stackrel{\Pi_{42}}{\vdots}$}\noLine
\UIC{$(S(n) = 1_{\omega} \to m < S(n) \to m < 1_{\omega}) \to \varphi_{168}$}
\BIC{$\varphi_{168}$}
\end{ptree}
\noindent$\Pi_{44}$:
\begin{ptree}
\AXC{$z = w \to s_{m} < s_{z} \to s_{m} < s_{w}$\lbl{A6}}
\UIC{$\forall w\,(z = w \to s_{m} < s_{z} \to s_{m} < s_{w})$}
\UIC{$\forall z\,\forall w\,(z = w \to s_{m} < s_{z} \to s_{m} < s_{w})$}
\AXC{$\varphi_{87}$\lbl{A2}}
\BIC{$\forall w\,(S(n) = w \to s_{m} < s_{S(n)} \to s_{m} < s_{w})$}
\AXC{$\varphi_{92}$\lbl{A2}}
\BIC{$S(n) = 1_{\omega} \to s_{m} < s_{S(n)} \to s_{m} < s_{1_{\omega}}$}
\AXC{$\stackrel{\Pi_{43}}{\vdots}$}\noLine
\UIC{$\varphi_{168}$}
\BIC{$\varphi_{157}$}
\end{ptree}
\noindent$\Pi_{45}$:
\begin{ptree}
\AXC{$\stackrel{\Pi_{12}}{\vdots}$}\noLine
\UIC{$\forall x\,\forall y\,(x = y \to y = x)$}
\AXC{$\varphi_{1}$\lbl{A2}}
\BIC{$\forall y\,(1_{\omega} = y \to y = 1_{\omega})$}
\AXC{$\varphi_{3}$\lbl{A2}}
\BIC{$1_{\omega} = S(n) \to S(n) = 1_{\omega}$}
\AXC{$\stackrel{\Pi_{35}}{\vdots}$}\noLine
\UIC{$\varphi_{131}$}
\BIC{$\mathrm{SubIdx}(s) \to m \in \mathbb{N} \to \varphi_{94}$}
\AXC{$\stackrel{\Pi_{44}}{\vdots}$}\noLine
\UIC{$\varphi_{157}$}
\BIC{$\mathrm{SubIdx}(s) \to m \in \mathbb{N} \to \varphi_{122}$}
\AXC{$(\mathrm{SubIdx}(s) \to m \in \mathbb{N} \to \varphi_{122}) \to \varphi_{160}$\lbl{A1}}
\BIC{$\varphi_{160}$}
\end{ptree}
\noindent$\Pi_{46}$:
\begin{ptree}
\AXC{$\forall x\,(\mathrm{Ind}(x) \leftrightarrow \emptyset \in x \land \forall y \in x\,(S(y) \in x))$\lbl{DefInd}}
\AXC{$\varphi_{90}$\lbl{A2}}
\BIC{$\mathrm{Ind}(\omega) \leftrightarrow \emptyset \in \omega \land \forall y \in \omega\,(S(y) \in \omega)$}
\AXC{$\mathrm{Ind}(\omega)$\lbl{補題(ωは帰納的)}}
\AXC{$\mathrm{Ind}(\omega) \to \varphi_{139}$\lbl{A1}}
\BIC{$\varphi_{139}$}
\BIC{$\mathrm{SubIdx}(s) \to m \in \mathbb{N} \to \varphi_{100}$}
\end{ptree}
\noindent$\Pi_{47}$:
\begin{ptree}
\AXC{$\stackrel{\Pi_{46}}{\vdots}$}\noLine
\UIC{$\mathrm{SubIdx}(s) \to m \in \mathbb{N} \to \varphi_{100}$}
\AXC{$\varphi_{147}$\lbl{A1}}
\BIC{$\mathrm{SubIdx}(s) \to m \in \mathbb{N} \to \varphi_{89}$}
\AXC{$(\mathrm{SubIdx}(s) \to m \in \mathbb{N} \to \varphi_{89}) \to \varphi_{137}$\lbl{A1}}
\BIC{$\varphi_{137}$}
\end{ptree}
\noindent$\Pi_{48}$:
\begin{ptree}
\AXC{$\forall y \in \omega\,(S(y) \in \omega) \to n \in \omega \to S(n) \in \omega$\lbl{A2}}
\AXC{$\stackrel{\Pi_{47}}{\vdots}$}\noLine
\UIC{$\varphi_{137}$}
\BIC{$\mathrm{SubIdx}(s) \to m \in \mathbb{N} \to \varphi_{84}$}
\AXC{$(\mathrm{SubIdx}(s) \to m \in \mathbb{N} \to \varphi_{84}) \to \varphi_{142}$\lbl{A1}}
\BIC{$\varphi_{142}$}
\end{ptree}
\noindent$\Pi_{49}$:
\begin{ptree}
\AXC{$\stackrel{\Pi_{6}}{\vdots}$}\noLine
\UIC{$\mathrm{SubIdx}(s) \to m \in \mathbb{N} \to \varphi_{70}$}
\AXC{$\stackrel{\Pi_{48}}{\vdots}$}\noLine
\UIC{$\varphi_{142}$}
\BIC{$\mathrm{SubIdx}(s) \to m \in \mathbb{N} \to \varphi_{73}$}
\AXC{$\varphi_{125}$\lbl{A2}}
\AXC{$\mathrm{SubIdx}(s) \to m \in \mathbb{N} \to \varphi_{120}$\lbl{A1}}
\AXC{$\varphi_{172}$\lbl{A1}}
\BIC{$\varphi_{125} \to \mathrm{SubIdx}(s) \to m \in \mathbb{N} \to \varphi_{125}$}
\BIC{$\mathrm{SubIdx}(s) \to m \in \mathbb{N} \to \varphi_{125}$}
\AXC{$\varphi_{177}$\lbl{A1}}
\BIC{$(\mathrm{SubIdx}(s) \to m \in \mathbb{N} \to \varphi_{73}) \to \varphi_{144}$}
\BIC{$\varphi_{144}$}
\end{ptree}
\noindent$\Pi_{50}$:
\begin{ptree}
\AXC{$z = w \to s_{z} < s_{S(n)} \to s_{w} < s_{S(n)}$\lbl{A6}}
\UIC{$\forall w\,(z = w \to s_{z} < s_{S(n)} \to s_{w} < s_{S(n)})$}
\UIC{$\forall z\,\forall w\,(z = w \to s_{z} < s_{S(n)} \to s_{w} < s_{S(n)})$}
\AXC{$\varphi_{91}$\lbl{A2}}
\BIC{$\forall w\,(n = w \to s_{n} < s_{S(n)} \to s_{w} < s_{S(n)})$}
\AXC{$\varphi_{86}$\lbl{A2}}
\BIC{$n = m \to s_{n} < s_{S(n)} \to s_{m} < s_{S(n)}$}
\AXC{$\stackrel{\Pi_{14}}{\vdots}$}\noLine
\UIC{$\varphi_{140}$}
\BIC{$\mathrm{SubIdx}(s) \to m \in \mathbb{N} \to \varphi_{126}$}
\AXC{$\stackrel{\Pi_{25}}{\vdots}$}\noLine
\UIC{$\varphi_{164}$}
\BIC{$\mathrm{SubIdx}(s) \to m \in \mathbb{N} \to \varphi_{124}$}
\AXC{$\stackrel{\Pi_{45}}{\vdots}$}\noLine
\UIC{$\varphi_{160}$}
\BIC{$\mathrm{SubIdx}(s) \to m \in \mathbb{N} \to \varphi_{112}$}
\AXC{$\stackrel{\Pi_{49}}{\vdots}$}\noLine
\UIC{$\varphi_{144}$}
\BIC{$\mathrm{SubIdx}(s) \to m \in \mathbb{N} \to \varphi_{71}$}
\end{ptree}
\noindent$\Pi_{51}$:
\begin{ptree}
\AXC{$\varphi_{67} \to \forall n\,(\varphi_{67})$\lbl{A4}}
\AXC{$\varphi_{143}$\lbl{A3}}
\AXC{$\stackrel{\Pi_{50}}{\vdots}$}\noLine
\UIC{$\mathrm{SubIdx}(s) \to m \in \mathbb{N} \to \varphi_{71}$}
\AXC{$\varphi_{141}$\lbl{A1}}
\BIC{$\varphi_{67} \to n \in B_{11} \to S(n) \in B_{11}$}
\UIC{$\forall n\,(\varphi_{67} \to n \in B_{11} \to S(n) \in B_{11})$}
\AXC{$\varphi_{182}$\lbl{A1}}
\BIC{$\varphi_{143} \to (\varphi_{67} \to \forall n\,(\varphi_{67})) \to \mathrm{SubIdx}(s) \to m \in \mathbb{N} \to \varphi_{74}$}
\BIC{$(\varphi_{67} \to \forall n\,(\varphi_{67})) \to \mathrm{SubIdx}(s) \to m \in \mathbb{N} \to \varphi_{74}$}
\BIC{$\mathrm{SubIdx}(s) \to m \in \mathbb{N} \to \varphi_{74}$}
\end{ptree}
\noindent$\Pi_{52}$:
\begin{ptree}
\AXC{$\forall n\,(n \in \mathbb{N} \leftrightarrow n \in \omega \land n \neq \emptyset)$\lbl{Def$\mathbb{N}$}}
\AXC{$\varphi_{75}$\lbl{A2}}
\BIC{$\emptyset \in \mathbb{N} \leftrightarrow \emptyset \in \omega \land \emptyset \neq \emptyset$}
\AXC{$(\emptyset \in \mathbb{N} \leftrightarrow \emptyset \in \omega \land \emptyset \neq \emptyset) \to \emptyset = \emptyset \to \varphi_{62}$\lbl{A1}}
\BIC{$\emptyset = \emptyset \to \varphi_{62}$}
\end{ptree}
\noindent$\Pi_{53}$:
\begin{ptree}
\AXC{$\varphi_{132}$\lbl{A2}}
\AXC{$\mathrm{SubIdx}(s) \to m \in \mathbb{N} \to \varphi_{120}$\lbl{A1}}
\AXC{$\varphi_{173}$\lbl{A1}}
\BIC{$\varphi_{132} \to \mathrm{SubIdx}(s) \to m \in \mathbb{N} \to \varphi_{132}$}
\BIC{$\mathrm{SubIdx}(s) \to m \in \mathbb{N} \to \varphi_{132}$}
\AXC{$(\mathrm{SubIdx}(s) \to m \in \mathbb{N} \to \varphi_{132}) \to \varphi_{163}$\lbl{A1}}
\BIC{$\varphi_{163}$}
\end{ptree}
\noindent$\Pi_{54}$:
\begin{ptree}
\AXC{$x = x$\lbl{A5}}
\UIC{$\forall x\,(x = x)$}
\AXC{$\forall x\,(x = x) \to \emptyset = \emptyset$\lbl{A2}}
\BIC{$\emptyset = \emptyset$}
\AXC{$\stackrel{\Pi_{52}}{\vdots}$}\noLine
\UIC{$\emptyset = \emptyset \to \varphi_{62}$}
\BIC{$\varphi_{62}$}
\AXC{$\stackrel{\Pi_{46}}{\vdots}$}\noLine
\UIC{$\mathrm{SubIdx}(s) \to m \in \mathbb{N} \to \varphi_{100}$}
\AXC{$\stackrel{\Pi_{53}}{\vdots}$}\noLine
\UIC{$\varphi_{163}$}
\BIC{$\varphi_{62} \to \mathrm{SubIdx}(s) \to m \in \mathbb{N} \to \varphi_{65}$}
\BIC{$\mathrm{SubIdx}(s) \to m \in \mathbb{N} \to \varphi_{65}$}
\end{ptree}
\noindent$\Pi_{55}$:
\begin{ptree}
\AXC{$\forall x\,(\mathrm{Ind}(x) \leftrightarrow \emptyset \in x \land \forall y \in x\,(S(y) \in x))$\lbl{DefInd}}
\AXC{$\forall x\,(\mathrm{Ind}(x) \leftrightarrow \emptyset \in x \land \forall y \in x\,(S(y) \in x)) \to \varphi_{18}$\lbl{A2}}
\BIC{$\varphi_{18}$}
\end{ptree}
\noindent$\Pi_{56}$:
\begin{ptree}
\AXC{$\stackrel{\Pi_{51}}{\vdots}$}\noLine
\UIC{$\mathrm{SubIdx}(s) \to m \in \mathbb{N} \to \varphi_{74}$}
\AXC{$\stackrel{\Pi_{54}}{\vdots}$}\noLine
\UIC{$\mathrm{SubIdx}(s) \to m \in \mathbb{N} \to \varphi_{65}$}
\AXC{$\stackrel{\Pi_{55}}{\vdots}$}\noLine
\UIC{$\varphi_{18}$}
\AXC{$\varphi_{18} \to (\mathrm{SubIdx}(s) \to m \in \mathbb{N} \to \varphi_{65}) \to \varphi_{162}$\lbl{A1}}
\BIC{$(\mathrm{SubIdx}(s) \to m \in \mathbb{N} \to \varphi_{65}) \to \varphi_{162}$}
\BIC{$\varphi_{162}$}
\BIC{$\varphi_{17} \to \varphi_{16} \to \mathrm{SubIdx}(s) \to m \in \mathbb{N} \to \varphi_{55}$}
\end{ptree}
\noindent$\Pi_{57}$:
\begin{ptree}
\AXC{$\mathrm{SubIdx}(s) \to m \in \mathbb{N} \to \varphi_{72}$\lbl{A1}}
\AXC{$(\mathrm{SubIdx}(s) \to m \in \mathbb{N} \to \varphi_{72}) \to \varphi_{129}$\lbl{A1}}
\BIC{$\varphi_{129}$}
\end{ptree}
\noindent$\Pi_{58}$:
\begin{ptree}
\AXC{$\forall n\,(n \in \omega \leftrightarrow \forall x\,(\mathrm{Ind}(x) \to n \in x))$\lbl{Def$\omega$}}
\AXC{$\varphi_{43}$\lbl{A2}}
\BIC{$n \in \omega \leftrightarrow \forall x\,(\mathrm{Ind}(x) \to n \in x)$}
\AXC{$\stackrel{\Pi_{57}}{\vdots}$}\noLine
\UIC{$\varphi_{129}$}
\BIC{$\mathrm{SubIdx}(s) \to m \in \mathbb{N} \to \varphi_{79}$}
\AXC{$(\mathrm{SubIdx}(s) \to m \in \mathbb{N} \to \varphi_{79}) \to \varphi_{133}$\lbl{A1}}
\BIC{$\varphi_{133}$}
\end{ptree}
\noindent$\Pi_{59}$:
\begin{ptree}
\AXC{$\forall x\,(\mathrm{Ind}(x) \to n \in x) \to \mathrm{Ind}(B_{11}) \to n \in B_{11}$\lbl{A2}}
\AXC{$\stackrel{\Pi_{58}}{\vdots}$}\noLine
\UIC{$\varphi_{133}$}
\BIC{$\mathrm{SubIdx}(s) \to m \in \mathbb{N} \to \varphi_{81}$}
\AXC{$\varphi_{174}$\lbl{A1}}
\BIC{$(\mathrm{SubIdx}(s) \to m \in \mathbb{N} \to \varphi_{55}) \to \varphi_{155}$}
\end{ptree}
\noindent$\Pi_{60}$:
\begin{ptree}
\AXC{$\stackrel{\Pi_{2}}{\vdots}$}\noLine
\UIC{$\varphi_{16}$}
\AXC{$\stackrel{\Pi_{3}}{\vdots}$}\noLine
\UIC{$\varphi_{17}$}
\AXC{$\stackrel{\Pi_{56}}{\vdots}$}\noLine
\UIC{$\varphi_{17} \to \varphi_{16} \to \mathrm{SubIdx}(s) \to m \in \mathbb{N} \to \varphi_{55}$}
\BIC{$\varphi_{16} \to \mathrm{SubIdx}(s) \to m \in \mathbb{N} \to \varphi_{55}$}
\BIC{$\mathrm{SubIdx}(s) \to m \in \mathbb{N} \to \varphi_{55}$}
\AXC{$\stackrel{\Pi_{59}}{\vdots}$}\noLine
\UIC{$(\mathrm{SubIdx}(s) \to m \in \mathbb{N} \to \varphi_{55}) \to \varphi_{155}$}
\BIC{$\varphi_{155}$}
\end{ptree}
\noindent$\Pi_{61}$:
\begin{ptree}
\AXC{$\varphi_{151}$\lbl{A3}}
\AXC{$\varphi_{117}$\lbl{A2}}
\AXC{$\mathrm{SubIdx}(s) \to m \in \mathbb{N} \to \varphi_{120}$\lbl{A1}}
\AXC{$\varphi_{170}$\lbl{A1}}
\BIC{$\varphi_{117} \to \mathrm{SubIdx}(s) \to m \in \mathbb{N} \to \varphi_{117}$}
\BIC{$\mathrm{SubIdx}(s) \to m \in \mathbb{N} \to \varphi_{117}$}
\AXC{$\stackrel{\Pi_{60}}{\vdots}$}\noLine
\UIC{$\varphi_{155}$}
\BIC{$\mathrm{SubIdx}(s) \to m \in \mathbb{N} \to \varphi_{102}$}
\AXC{$\varphi_{149}$\lbl{A1}}
\BIC{$\varphi_{67} \to n \in \omega \to n \in \mathbb{N} \to m < n \to s_{m} < s_{n}$}
\UIC{$\forall n\,(\varphi_{67} \to n \in \omega \to n \in \mathbb{N} \to m < n \to s_{m} < s_{n})$}
\AXC{$\varphi_{183}$\lbl{A1}}
\BIC{$\varphi_{151} \to (\varphi_{67} \to \forall n\,(\varphi_{67})) \to \mathrm{SubIdx}(s) \to m \in \mathbb{N} \to \varphi_{105}$}
\BIC{$(\varphi_{67} \to \forall n\,(\varphi_{67})) \to \mathrm{SubIdx}(s) \to m \in \mathbb{N} \to \varphi_{105}$}
\end{ptree}
\noindent$\Pi_{62}$:
\begin{ptree}
\AXC{$\varphi_{67} \to \forall n\,(\varphi_{67})$\lbl{A4}}
\AXC{$\stackrel{\Pi_{61}}{\vdots}$}\noLine
\UIC{$(\varphi_{67} \to \forall n\,(\varphi_{67})) \to \mathrm{SubIdx}(s) \to m \in \mathbb{N} \to \varphi_{105}$}
\BIC{$\mathrm{SubIdx}(s) \to m \in \mathbb{N} \to \varphi_{105}$}
\AXC{$(\mathrm{SubIdx}(s) \to m \in \mathbb{N} \to \varphi_{105}) \to \varphi_{60} \to \varphi_{42}$\lbl{A1}}
\BIC{$\varphi_{60} \to \varphi_{42}$}
\UIC{$\forall B_{11}\,(\varphi_{60} \to \varphi_{42})$}
\end{ptree}
\noindent$\Pi_{63}$:
\begin{ptree}
\AXC{$\varphi_{69}$\lbl{Sep}}
\AXC{$\varphi_{69} \to \varphi_{64}$\lbl{A2}}
\BIC{$\varphi_{64}$}
\AXC{$\varphi_{64} \to \varphi_{51}$\lbl{A2}}
\BIC{$\varphi_{51}$}
\AXC{$\varphi_{51} \to \varphi_{46}$\lbl{A2}}
\BIC{$\varphi_{46}$}
\AXC{$\varphi_{46} \to \mathrm{SubIdx}(s) \to m \in \mathbb{N} \to \varphi_{46}$\lbl{A1}}
\BIC{$\mathrm{SubIdx}(s) \to m \in \mathbb{N} \to \varphi_{46}$}
\end{ptree}
\noindent$\Pi_{64}$:
\begin{ptree}
\AXC{$\stackrel{\Pi_{62}}{\vdots}$}\noLine
\UIC{$\forall B_{11}\,(\varphi_{60} \to \varphi_{42})$}
\AXC{$\stackrel{\Pi_{63}}{\vdots}$}\noLine
\UIC{$\mathrm{SubIdx}(s) \to m \in \mathbb{N} \to \varphi_{46}$}
\AXC{$(\mathrm{SubIdx}(s) \to m \in \mathbb{N} \to \varphi_{46}) \to \forall B_{11}\,(\varphi_{60} \to \varphi_{42}) \to \varphi_{178}$\lbl{A1}}
\BIC{$\forall B_{11}\,(\varphi_{60} \to \varphi_{42}) \to \varphi_{178}$}
\BIC{$\varphi_{178}$}
\end{ptree}
\noindent$\Pi_{65}$:
\begin{ptree}
\AXC{$\varphi_{85}$\lbl{A2}}
\AXC{$\varphi_{60} \to \forall B_{11}\,(\varphi_{60})$\lbl{A4}}
\AXC{$\forall B_{11}\,(\varphi_{60} \to \varphi_{42}) \to \forall B_{11}\,(\varphi_{60}) \to \forall B_{11}\,(\varphi_{42})$\lbl{A3}}
\AXC{$\stackrel{\Pi_{64}}{\vdots}$}\noLine
\UIC{$\varphi_{178}$}
\BIC{$(\varphi_{60} \to \forall B_{11}\,(\varphi_{60})) \to \varphi_{56}$}
\BIC{$\varphi_{56}$}
\AXC{$\varphi_{56} \to \varphi_{85} \to \varphi_{52}$\lbl{A1}}
\BIC{$\varphi_{85} \to \varphi_{52}$}
\BIC{$\varphi_{52}$}
\end{ptree}
\noindent$\Pi_{66}$:
\begin{ptree}
\AXC{$\stackrel{\Pi_{65}}{\vdots}$}\noLine
\UIC{$\varphi_{52}$}
\AXC{$\varphi_{52} \to (\mathrm{SubIdx}(s) \to m \in \mathbb{N} \to n \in \mathbb{N} \to n \in \omega) \to \varphi_{83}$\lbl{A1}}
\BIC{$(\mathrm{SubIdx}(s) \to m \in \mathbb{N} \to n \in \mathbb{N} \to n \in \omega) \to \varphi_{83}$}
\end{ptree}
\noindent$\Pi_{67}$:
\begin{ptree}
\AXC{$\forall n\,(n \in \mathbb{N} \leftrightarrow n \in \omega \land n \neq \emptyset)$\lbl{Def$\mathbb{N}$}}
\AXC{$\varphi_{53}$\lbl{A2}}
\BIC{$n \in \mathbb{N} \leftrightarrow n \in \omega \land n \neq \emptyset$}
\AXC{$\stackrel{\Pi_{1}}{\vdots}$}\noLine
\UIC{$\varphi_{63}$}
\BIC{$\mathrm{SubIdx}(s) \to m \in \mathbb{N} \to n \in \mathbb{N} \to n \in \omega$}
\AXC{$\stackrel{\Pi_{66}}{\vdots}$}\noLine
\UIC{$(\mathrm{SubIdx}(s) \to m \in \mathbb{N} \to n \in \mathbb{N} \to n \in \omega) \to \varphi_{83}$}
\BIC{$\varphi_{83}$}
\end{ptree}
\noindent$\Pi_{68}$:
\begin{ptree}
\AXC{$\varphi_{38} \to \varphi_{41}$\lbl{A3}}
\AXC{$\mathrm{SubIdx}(s) \to m \in \mathbb{N} \to n \in \mathbb{N} \to n \in \mathbb{N}$\lbl{A1}}
\AXC{$\stackrel{\Pi_{67}}{\vdots}$}\noLine
\UIC{$\varphi_{83}$}
\BIC{$\mathrm{SubIdx}(s) \to m \in \mathbb{N} \to n \in \mathbb{N} \to m < n \to s_{m} < s_{n}$}
\AXC{$\varphi_{107}$\lbl{A1}}
\BIC{$\mathrm{SubIdx}(s) \land m \in \mathbb{N} \to n \in \mathbb{N} \to m < n \to s_{m} < s_{n}$}
\UIC{$\varphi_{38}$}
\AXC{$\varphi_{38} \to (\varphi_{38} \to \varphi_{41}) \to \varphi_{9} \to \varphi_{39}$\lbl{A1}}
\BIC{$(\varphi_{38} \to \varphi_{41}) \to \varphi_{9} \to \varphi_{39}$}
\BIC{$\varphi_{9} \to \varphi_{39}$}
\end{ptree}
\noindent 主木:
\begin{ptree}
\AXC{$\varphi_{9}$\lbl{A4}}
\AXC{$\stackrel{\Pi_{68}}{\vdots}$}\noLine
\UIC{$\varphi_{9} \to \varphi_{39}$}
\BIC{$\varphi_{39}$}
\UIC{$\forall m\,(\varphi_{39})$}
\UIC{$\forall s\,\forall m\,(\varphi_{39})$}
\end{ptree}
\begin{abbrevlist}
\item $\varphi_{1} :\equiv \forall x\,\forall y\,(x = y \to y = x) \to \forall y\,(1_{\omega} = y \to y = 1_{\omega})$
\item $\varphi_{2} :\equiv \forall x\,\forall y\,(x = y \to y = x) \to \forall y\,(S(n) = y \to y = S(n))$
\item $\varphi_{3} :\equiv \forall y\,(1_{\omega} = y \to y = 1_{\omega}) \to 1_{\omega} = S(n) \to S(n) = 1_{\omega}$
\item $\varphi_{4} :\equiv \forall y\,(S(n) = y \to y = S(n)) \to S(n) = 1_{\omega} \to 1_{\omega} = S(n)$
\item $\varphi_{5} :\equiv \forall x\,\forall y\,(x \leq y \to \neg y < x) \to \forall y\,(1_{\omega} \leq y \to \neg y < 1_{\omega})$
\item $\varphi_{6} :\equiv \forall y\,(1_{\omega} \leq y \to \neg y < 1_{\omega}) \to 1_{\omega} \leq m \to \neg m < 1_{\omega}$
\item $\varphi_{7} :\equiv \mathrm{SubIdx}(s) \to m \in \mathbb{N} \to m < 1_{\omega} \to \neg m < 1_{\omega}$
\item $\varphi_{8} :\equiv S(n) = S(\emptyset) \to S(\emptyset) = 1_{\omega} \to S(n) = 1_{\omega}$
\item $\varphi_{9} :\equiv \mathrm{SubIdx}(s) \land m \in \mathbb{N} \to \forall n\,(\mathrm{SubIdx}(s) \land m \in \mathbb{N})$
\item $\varphi_{10} :\equiv \forall n \in B_{11}\,(S(n) \in B_{11}) \to y \in B_{11} \to S(y) \in B_{11}$
\item $\varphi_{11} :\equiv \forall y \in B_{11}\,(S(y) \in B_{11}) \to n \in B_{11} \to S(n) \in B_{11}$
\item $\varphi_{12} :\equiv \forall x\,\forall y\,\forall z\,(x = y \to y = z \to x = z) \to \forall y\,\forall z\,(S(n) = y \to y = z \to S(n) = z)$
\item $\varphi_{13} :\equiv (1_{\omega} = S(\emptyset) \to S(\emptyset) = 1_{\omega}) \to S(\emptyset) = 1_{\omega}$
\item $\varphi_{14} :\equiv \forall m\,\forall n\,(m \in \omega \to n \in \omega \to m \leq S(n) \to m \leq n \lor m = S(n))$
\item $\varphi_{15} :\equiv \forall y\,(1_{\omega} = y \to y = 1_{\omega}) \to 1_{\omega} = S(\emptyset) \to S(\emptyset) = 1_{\omega}$
\item $\varphi_{16} :\equiv \forall n \in B_{11}\,(S(n) \in B_{11}) \to \forall y \in B_{11}\,(S(y) \in B_{11})$
\item $\varphi_{17} :\equiv \forall y \in B_{11}\,(S(y) \in B_{11}) \to \forall n \in B_{11}\,(S(n) \in B_{11})$
\item $\varphi_{18} :\equiv \mathrm{Ind}(B_{11}) \leftrightarrow \emptyset \in B_{11} \land \forall y \in B_{11}\,(S(y) \in B_{11})$
\item $\varphi_{19} :\equiv \mathrm{SubIdx}(s) \to m \in \mathbb{N} \to m < 1_{\omega} \to s_{m} < s_{1_{\omega}}$
\item $\varphi_{20} :\equiv n \in B_{11} \to S(n) \in \mathbb{N} \to n \neq \emptyset \to m < S(n) \to m \leq S(n)$
\item $\varphi_{21} :\equiv \forall n \in B_{11}\,(S(n) \in B_{11}) \to \forall y\,\forall n \in B_{11}\,(S(n) \in B_{11})$
\item $\varphi_{22} :\equiv \forall y \in B_{11}\,(S(y) \in B_{11}) \to \forall n\,\forall y \in B_{11}\,(S(y) \in B_{11})$
\item $\varphi_{23} :\equiv \mathrm{SubIdx}(s) \leftrightarrow \mathrm{NSeq}(s) \land \forall k \in \mathbb{N}\,(s_{k} < s_{S(k)})$
\item $\varphi_{24} :\equiv \forall x\,\forall y\,(x < y \leftrightarrow x \subseteq y \land x \neq y) \to \forall y\,(m < y \leftrightarrow m \subseteq y \land m \neq y)$
\item $\varphi_{25} :\equiv \forall z\,\forall w\,(z = w \to m < z \to m < w) \to \forall w\,(S(n) = w \to m < S(n) \to m < w)$
\item $\varphi_{26} :\equiv \forall z\,\forall w\,(z = w \to m < z \to m < w) \to \forall w\,(1_{\omega} = w \to m < 1_{\omega} \to m < w)$
\item $\varphi_{27} :\equiv \forall y\,\forall n \in B_{11}\,(S(n) \in B_{11}) \to \forall y \in B_{11}\,(S(y) \in B_{11})$
\item $\varphi_{28} :\equiv \forall n\,\forall y \in B_{11}\,(S(y) \in B_{11}) \to \forall n \in B_{11}\,(S(n) \in B_{11})$
\item $\varphi_{29} :\equiv n \in B_{11} \to S(n) \in \mathbb{N} \to n \neq \emptyset \to m < S(n) \to m \neq S(n)$
\item $\varphi_{30} :\equiv n \in B_{11} \to S(n) \in \mathbb{N} \to n \neq \emptyset \to m < S(n) \to m = n \to m = n$
\item $\varphi_{31} :\equiv n \in B_{11} \to S(n) \in \mathbb{N} \to n \neq \emptyset \to m < S(n) \to m = n \to n = m$
\item $\varphi_{32} :\equiv n \in B_{11} \to S(n) \in \mathbb{N} \to n \neq \emptyset \to m < S(n) \to m < S(n)$
\item $\varphi_{33} :\equiv \forall y\,(m < y \leftrightarrow m \subseteq y \land m \neq y) \to (m < S(n) \leftrightarrow m \subseteq S(n) \land m \neq S(n))$
\item $\varphi_{34} :\equiv n \in B_{11} \to S(n) \in \mathbb{N} \to n \neq \emptyset \to m < S(n) \to m \leq n \lor m = S(n)$
\item $\varphi_{35} :\equiv \forall z\,\forall w\,(z = w \to S(n) = S(z) \to S(n) = S(w)) \to \forall w\,(n = w \to S(n) = S(n) \to S(n) = S(w))$
\item $\varphi_{36} :\equiv \forall w\,(S(n) = w \to m < S(n) \to m < w) \to S(n) = 1_{\omega} \to m < S(n) \to m < 1_{\omega}$
\item $\varphi_{37} :\equiv \forall w\,(1_{\omega} = w \to m < 1_{\omega} \to m < w) \to 1_{\omega} = S(n) \to m < 1_{\omega} \to m < S(n)$
\item $\varphi_{38} :\equiv \forall n\,(\mathrm{SubIdx}(s) \land m \in \mathbb{N} \to n \in \mathbb{N} \to m < n \to s_{m} < s_{n})$
\item $\varphi_{39} :\equiv \mathrm{SubIdx}(s) \to m \in \mathbb{N} \to \forall n \in \mathbb{N}\,(m < n \to s_{m} < s_{n})$
\item $\varphi_{40} :\equiv n \in B_{11} \to S(n) \in \mathbb{N} \to n \neq \emptyset \to m < S(n) \to m < n \to m < n$
\item $\varphi_{41} :\equiv \forall n\,(\mathrm{SubIdx}(s) \land m \in \mathbb{N}) \to \forall n \in \mathbb{N}\,(m < n \to s_{m} < s_{n})$
\item $\varphi_{42} :\equiv \neg \forall n\,(n \in B_{11} \leftrightarrow n \in \omega \land (n \in \mathbb{N} \to m < n \to s_{m} < s_{n}))$
\item $\varphi_{43} :\equiv \forall n\,(n \in \omega \leftrightarrow \forall x\,(\mathrm{Ind}(x) \to n \in x)) \to (n \in \omega \leftrightarrow \forall x\,(\mathrm{Ind}(x) \to n \in x))$
\item $\varphi_{44} :\equiv \forall y\,\forall z\,(S(n) = y \to y = z \to S(n) = z) \to \forall z\,(S(n) = S(\emptyset) \to S(\emptyset) = z \to S(n) = z)$
\item $\varphi_{45} :\equiv \forall w\,(n = w \to S(n) = S(n) \to S(n) = S(w)) \to n = \emptyset \to S(n) = S(n) \to S(n) = S(\emptyset)$
\item $\varphi_{46} :\equiv \exists B_{11}\,\forall n\,(n \in B_{11} \leftrightarrow n \in \omega \land (n \in \mathbb{N} \to m < n \to s_{m} < s_{n}))$
\item $\varphi_{47} :\equiv n \in B_{11} \to S(n) \in \mathbb{N} \to n = \emptyset \to m < S(n) \to s_{m} < s_{S(n)}$
\item $\varphi_{48} :\equiv n \in B_{11} \to S(n) \in \mathbb{N} \to n \neq \emptyset \to m < S(n) \to s_{n} < s_{S(n)}$
\item $\varphi_{49} :\equiv n \in B_{11} \to S(n) \in \mathbb{N} \to n \neq \emptyset \to m < S(n) \to s_{m} < s_{S(n)}$
\item $\varphi_{50} :\equiv (\mathrm{SubIdx}(s) \to m \in \mathbb{N} \to m \in \mathbb{N}) \to \mathrm{SubIdx}(s) \to m \in \mathbb{N} \to m < 1_{\omega} \to 1_{\omega} \leq m$
\item $\varphi_{51} :\equiv \forall A_{12}\,\exists B_{11}\,\forall n\,(n \in B_{11} \leftrightarrow n \in A_{12} \land (n \in \mathbb{N} \to m < n \to s_{m} < s_{n}))$
\item $\varphi_{52} :\equiv \mathrm{SubIdx}(s) \to m \in \mathbb{N} \to n \in \omega \to n \in \mathbb{N} \to m < n \to s_{m} < s_{n}$
\item $\varphi_{53} :\equiv \forall n\,(n \in \mathbb{N} \leftrightarrow n \in \omega \land n \neq \emptyset) \to (n \in \mathbb{N} \leftrightarrow n \in \omega \land n \neq \emptyset)$
\item $\varphi_{54} :\equiv \forall x\,\forall y\,\forall z\,(x < y \to y < z \to x < z) \to \forall y\,\forall z\,(s_{m} < y \to y < z \to s_{m} < z)$
\item $\varphi_{55} :\equiv \forall n\,(n \in B_{11} \leftrightarrow n \in \omega \land (n \in \mathbb{N} \to m < n \to s_{m} < s_{n})) \to \mathrm{Ind}(B_{11})$
\item $\varphi_{56} :\equiv \mathrm{SubIdx}(s) \to m \in \mathbb{N} \to \forall n \in \omega\,(n \in \mathbb{N} \to m < n \to s_{m} < s_{n})$
\item $\varphi_{57} :\equiv n \in B_{11} \to S(n) \in \mathbb{N} \to n \neq \emptyset \to m < S(n) \to m = n \to s_{m} < s_{S(n)}$
\item $\varphi_{58} :\equiv n \in B_{11} \to S(n) \in \mathbb{N} \to n \neq \emptyset \to m < S(n) \to m < n \to s_{m} < s_{n}$
\item $\varphi_{59} :\equiv S(n) \in B_{11} \leftrightarrow S(n) \in \omega \land (S(n) \in \mathbb{N} \to m < S(n) \to s_{m} < s_{S(n)})$
\item $\varphi_{60} :\equiv \mathrm{SubIdx}(s) \land m \in \mathbb{N} \land \neg \forall n \in \omega\,(n \in \mathbb{N} \to m < n \to s_{m} < s_{n})$
\item $\varphi_{61} :\equiv n \in B_{11} \to S(n) \in \mathbb{N} \to n \neq \emptyset \to m < S(n) \to m < n \to s_{m} < s_{S(n)}$
\item $\varphi_{62} :\equiv \mathrm{SubIdx}(s) \to m \in \mathbb{N} \to \emptyset \in \mathbb{N} \to m < \emptyset \to s_{m} < s_{\emptyset}$
\item $\varphi_{63} :\equiv (n \in \mathbb{N} \leftrightarrow n \in \omega \land n \neq \emptyset) \to \mathrm{SubIdx}(s) \to m \in \mathbb{N} \to n \in \mathbb{N} \to n \in \omega$
\item $\varphi_{64} :\equiv \forall z_{14}\,\forall A_{12}\,\exists B_{11}\,\forall n\,(n \in B_{11} \leftrightarrow n \in A_{12} \land (n \in \mathbb{N} \to m < n \to z_{14}(m) < z_{14}(n)))$
\item $\varphi_{65} :\equiv \forall n\,(n \in B_{11} \leftrightarrow n \in \omega \land (n \in \mathbb{N} \to m < n \to s_{m} < s_{n})) \to \emptyset \in B_{11}$
\item $\varphi_{66} :\equiv (\mathrm{SubIdx}(s) \to m \in \mathbb{N} \to m < 1_{\omega} \to 1_{\omega} \leq m) \to \varphi_{7}$
\item $\varphi_{67} :\equiv \mathrm{SubIdx}(s) \land m \in \mathbb{N} \land \forall n\,(n \in B_{11} \leftrightarrow n \in \omega \land (n \in \mathbb{N} \to m < n \to s_{m} < s_{n}))$
\item $\varphi_{68} :\equiv \forall n\,(n \in B_{11} \leftrightarrow n \in \omega \land (n \in \mathbb{N} \to m < n \to s_{m} < s_{n})) \to n \in B_{11} \to n \in B_{11}$
\item $\varphi_{69} :\equiv \forall z_{13}\,\forall z_{14}\,\forall A_{12}\,\exists B_{11}\,\forall n\,(n \in B_{11} \leftrightarrow n \in A_{12} \land (n \in \mathbb{N} \to z_{13} < n \to z_{14}(z_{13}) < z_{14}(n)))$
\item $\varphi_{70} :\equiv \forall n\,(n \in B_{11} \leftrightarrow n \in \omega \land (n \in \mathbb{N} \to m < n \to s_{m} < s_{n})) \to n \in B_{11} \to n \in \omega$
\item $\varphi_{71} :\equiv \forall n\,(n \in B_{11} \leftrightarrow n \in \omega \land (n \in \mathbb{N} \to m < n \to s_{m} < s_{n})) \to n \in B_{11} \to S(n) \in B_{11}$
\item $\varphi_{72} :\equiv \forall n\,(n \in B_{11} \leftrightarrow n \in \omega \land (n \in \mathbb{N} \to m < n \to s_{m} < s_{n})) \to n \in \omega \to n \in \omega$
\item $\varphi_{73} :\equiv \forall n\,(n \in B_{11} \leftrightarrow n \in \omega \land (n \in \mathbb{N} \to m < n \to s_{m} < s_{n})) \to n \in B_{11} \to S(n) \in \omega$
\item $\varphi_{74} :\equiv \forall n\,(n \in B_{11} \leftrightarrow n \in \omega \land (n \in \mathbb{N} \to m < n \to s_{m} < s_{n})) \to \forall n \in B_{11}\,(S(n) \in B_{11})$
\item $\varphi_{75} :\equiv \forall n\,(n \in \mathbb{N} \leftrightarrow n \in \omega \land n \neq \emptyset) \to (\emptyset \in \mathbb{N} \leftrightarrow \emptyset \in \omega \land \emptyset \neq \emptyset)$
\item $\varphi_{76} :\equiv \forall y\,\forall z\,(s_{m} < y \to y < z \to s_{m} < z) \to \forall z\,(s_{m} < s_{n} \to s_{n} < z \to s_{m} < z)$
\item $\varphi_{77} :\equiv \emptyset \in B_{11} \leftrightarrow \emptyset \in \omega \land (\emptyset \in \mathbb{N} \to m < \emptyset \to s_{m} < s_{\emptyset})$
\item $\varphi_{78} :\equiv \forall n\,(m \in \omega \to n \in \omega \to m \leq S(n) \to m \leq n \lor m = S(n)) \to m \in \omega \to n \in \omega \to m \leq S(n) \to m \leq n \lor m = S(n)$
\item $\varphi_{79} :\equiv \forall n\,(n \in B_{11} \leftrightarrow n \in \omega \land (n \in \mathbb{N} \to m < n \to s_{m} < s_{n})) \to n \in \omega \to \forall x\,(\mathrm{Ind}(x) \to n \in x)$
\item $\varphi_{80} :\equiv (\mathrm{SubIdx}(s) \to m \in \mathbb{N} \to m < 1_{\omega} \to m < 1_{\omega}) \to \varphi_{19}$
\item $\varphi_{81} :\equiv \forall n\,(n \in B_{11} \leftrightarrow n \in \omega \land (n \in \mathbb{N} \to m < n \to s_{m} < s_{n})) \to n \in \omega \to \mathrm{Ind}(B_{11}) \to n \in B_{11}$
\item $\varphi_{82} :\equiv \varphi_{14} \to \forall n\,(m \in \omega \to n \in \omega \to m \leq S(n) \to m \leq n \lor m = S(n))$
\item $\varphi_{83} :\equiv (\mathrm{SubIdx}(s) \to m \in \mathbb{N} \to n \in \mathbb{N} \to n \in \mathbb{N}) \to \mathrm{SubIdx}(s) \to m \in \mathbb{N} \to n \in \mathbb{N} \to m < n \to s_{m} < s_{n}$
\item $\varphi_{84} :\equiv \forall n\,(n \in B_{11} \leftrightarrow n \in \omega \land (n \in \mathbb{N} \to m < n \to s_{m} < s_{n})) \to n \in B_{11} \to n \in \omega \to S(n) \in \omega$
\item $\varphi_{85} :\equiv \forall n \in \omega\,(n \in \mathbb{N} \to m < n \to s_{m} < s_{n}) \to n \in \omega \to n \in \mathbb{N} \to m < n \to s_{m} < s_{n}$
\item $\varphi_{86} :\equiv \forall w\,(n = w \to s_{n} < s_{S(n)} \to s_{w} < s_{S(n)}) \to n = m \to s_{n} < s_{S(n)} \to s_{m} < s_{S(n)}$
\item $\varphi_{87} :\equiv \forall z\,\forall w\,(z = w \to s_{m} < s_{z} \to s_{m} < s_{w}) \to \forall w\,(S(n) = w \to s_{m} < s_{S(n)} \to s_{m} < s_{w})$
\item $\varphi_{88} :\equiv \forall z\,\forall w\,(z = w \to s_{m} < s_{z} \to s_{m} < s_{w}) \to \forall w\,(1_{\omega} = w \to s_{m} < s_{1_{\omega}} \to s_{m} < s_{w})$
\item $\varphi_{89} :\equiv \forall n\,(n \in B_{11} \leftrightarrow n \in \omega \land (n \in \mathbb{N} \to m < n \to s_{m} < s_{n})) \to n \in B_{11} \to \forall y \in \omega\,(S(y) \in \omega)$
\item $\varphi_{90} :\equiv \forall x\,(\mathrm{Ind}(x) \leftrightarrow \emptyset \in x \land \forall y \in x\,(S(y) \in x)) \to (\mathrm{Ind}(\omega) \leftrightarrow \emptyset \in \omega \land \forall y \in \omega\,(S(y) \in \omega))$
\item $\varphi_{91} :\equiv \forall z\,\forall w\,(z = w \to s_{z} < s_{S(n)} \to s_{w} < s_{S(n)}) \to \forall w\,(n = w \to s_{n} < s_{S(n)} \to s_{w} < s_{S(n)})$
\item $\varphi_{92} :\equiv \forall w\,(S(n) = w \to s_{m} < s_{S(n)} \to s_{m} < s_{w}) \to S(n) = 1_{\omega} \to s_{m} < s_{S(n)} \to s_{m} < s_{1_{\omega}}$
\item $\varphi_{93} :\equiv \forall w\,(1_{\omega} = w \to s_{m} < s_{1_{\omega}} \to s_{m} < s_{w}) \to 1_{\omega} = S(n) \to s_{m} < s_{1_{\omega}} \to s_{m} < s_{S(n)}$
\item $\varphi_{94} :\equiv \forall n\,(n \in B_{11} \leftrightarrow n \in \omega \land (n \in \mathbb{N} \to m < n \to s_{m} < s_{n})) \to n \in B_{11} \to S(n) \in \mathbb{N} \to n = \emptyset \to S(n) = 1_{\omega}$
\item $\varphi_{95} :\equiv \forall n\,(n \in B_{11} \leftrightarrow n \in \omega \land (n \in \mathbb{N} \to m < n \to s_{m} < s_{n})) \to n \in B_{11} \to S(n) \in \mathbb{N} \to n = \emptyset \to 1_{\omega} = S(n)$
\item $\varphi_{96} :\equiv \forall n\,(n \in B_{11} \leftrightarrow n \in \omega \land (n \in \mathbb{N} \to m < n \to s_{m} < s_{n})) \to n \in B_{11} \to S(n) \in \mathbb{N} \to n \neq \emptyset \to n \in \mathbb{N}$
\item $\varphi_{97} :\equiv \forall n\,(n \in B_{11} \leftrightarrow n \in \omega \land (n \in \mathbb{N} \to m < n \to s_{m} < s_{n})) \to n \in B_{11} \to S(n) \in \mathbb{N} \to n = \emptyset \to n = \emptyset$
\item $\varphi_{98} :\equiv \varphi_{23} \to \mathrm{SubIdx}(s) \to m \in \mathbb{N} \to \forall k \in \mathbb{N}\,(s_{k} < s_{S(k)})$
\item $\varphi_{99} :\equiv \forall n\,(n \in B_{11} \leftrightarrow n \in \omega \land (n \in \mathbb{N} \to m < n \to s_{m} < s_{n})) \to n \in B_{11} \to S(n) \in \mathbb{N} \to n \neq \emptyset \to n \neq \emptyset$
\item $\varphi_{100} :\equiv \forall n\,(n \in B_{11} \leftrightarrow n \in \omega \land (n \in \mathbb{N} \to m < n \to s_{m} < s_{n})) \to \emptyset \in \omega \land \forall y \in \omega\,(S(y) \in \omega)$
\item $\varphi_{101} :\equiv \forall n\,(n \in B_{11} \leftrightarrow n \in \omega \land (n \in \mathbb{N} \to m < n \to s_{m} < s_{n})) \to n \in B_{11} \to n \in \mathbb{N} \to m < n \to s_{m} < s_{n}$
\item $\varphi_{102} :\equiv \forall n\,(n \in B_{11} \leftrightarrow n \in \omega \land (n \in \mathbb{N} \to m < n \to s_{m} < s_{n})) \to n \in \omega \to n \in \mathbb{N} \to m < n \to s_{m} < s_{n}$
\item $\varphi_{103} :\equiv \forall n\,(n \in B_{11} \leftrightarrow n \in \omega \land (n \in \mathbb{N} \to m < n \to s_{m} < s_{n})) \to n \in B_{11} \to S(n) \in \mathbb{N} \to n = \emptyset \to S(n) = S(\emptyset)$
\item $\varphi_{104} :\equiv \forall z\,(s_{m} < s_{n} \to s_{n} < z \to s_{m} < z) \to s_{m} < s_{n} \to s_{n} < s_{S(n)} \to s_{m} < s_{S(n)}$
\item $\varphi_{105} :\equiv \forall n\,(n \in B_{11} \leftrightarrow n \in \omega \land (n \in \mathbb{N} \to m < n \to s_{m} < s_{n})) \to \forall n \in \omega\,(n \in \mathbb{N} \to m < n \to s_{m} < s_{n})$
\item $\varphi_{106} :\equiv \forall n\,(n \in B_{11} \leftrightarrow n \in \omega \land (n \in \mathbb{N} \to m < n \to s_{m} < s_{n})) \to \varphi_{20}$
\item $\varphi_{107} :\equiv (\mathrm{SubIdx}(s) \to m \in \mathbb{N} \to n \in \mathbb{N} \to m < n \to s_{m} < s_{n}) \to \mathrm{SubIdx}(s) \land m \in \mathbb{N} \to n \in \mathbb{N} \to m < n \to s_{m} < s_{n}$
\item $\varphi_{108} :\equiv \forall n\,(n \in B_{11} \leftrightarrow n \in \omega \land (n \in \mathbb{N} \to m < n \to s_{m} < s_{n})) \to \varphi_{29}$
\item $\varphi_{109} :\equiv \forall n\,(n \in B_{11} \leftrightarrow n \in \omega \land (n \in \mathbb{N} \to m < n \to s_{m} < s_{n})) \to \varphi_{30}$
\item $\varphi_{110} :\equiv \forall n\,(n \in B_{11} \leftrightarrow n \in \omega \land (n \in \mathbb{N} \to m < n \to s_{m} < s_{n})) \to \varphi_{31}$
\item $\varphi_{111} :\equiv \forall n\,(n \in B_{11} \leftrightarrow n \in \omega \land (n \in \mathbb{N} \to m < n \to s_{m} < s_{n})) \to \varphi_{32}$
\item $\varphi_{112} :\equiv \forall n\,(n \in B_{11} \leftrightarrow n \in \omega \land (n \in \mathbb{N} \to m < n \to s_{m} < s_{n})) \to n \in B_{11} \to S(n) \in \mathbb{N} \to m < S(n) \to s_{m} < s_{S(n)}$
\item $\varphi_{113} :\equiv \forall n\,(n \in B_{11} \leftrightarrow n \in \omega \land (n \in \mathbb{N} \to m < n \to s_{m} < s_{n})) \to n \in B_{11} \to S(n) \in \mathbb{N} \to n \neq \emptyset \to m < S(n) \to n \in \omega$
\item $\varphi_{114} :\equiv \forall n\,(n \in B_{11} \leftrightarrow n \in \omega \land (n \in \mathbb{N} \to m < n \to s_{m} < s_{n})) \to n \in B_{11} \to S(n) \in \mathbb{N} \to n \neq \emptyset \to m < S(n) \to m \in \omega$
\item $\varphi_{115} :\equiv (\forall k \in \mathbb{N}\,(s_{k} < s_{S(k)}) \to n \in \mathbb{N} \to s_{n} < s_{S(n)}) \to \mathrm{SubIdx}(s) \to m \in \mathbb{N} \to n \in \mathbb{N} \to s_{n} < s_{S(n)}$
\item $\varphi_{116} :\equiv \forall n\,(n \in B_{11} \leftrightarrow n \in \omega \land (n \in \mathbb{N} \to m < n \to s_{m} < s_{n})) \to \varphi_{34}$
\item $\varphi_{117} :\equiv \forall n\,(n \in B_{11} \leftrightarrow n \in \omega \land (n \in \mathbb{N} \to m < n \to s_{m} < s_{n})) \to (n \in B_{11} \leftrightarrow n \in \omega \land (n \in \mathbb{N} \to m < n \to s_{m} < s_{n}))$
\item $\varphi_{118} :\equiv \forall n\,(n \in B_{11} \leftrightarrow n \in \omega \land (n \in \mathbb{N} \to m < n \to s_{m} < s_{n})) \to n \in B_{11} \to n \in \omega \land (n \in \mathbb{N} \to m < n \to s_{m} < s_{n})$
\item $\varphi_{119} :\equiv \forall n\,(n \in B_{11} \leftrightarrow n \in \omega \land (n \in \mathbb{N} \to m < n \to s_{m} < s_{n})) \to \varphi_{40}$
\item $\varphi_{120} :\equiv \forall n\,(n \in B_{11} \leftrightarrow n \in \omega \land (n \in \mathbb{N} \to m < n \to s_{m} < s_{n})) \to \forall n\,(n \in B_{11} \leftrightarrow n \in \omega \land (n \in \mathbb{N} \to m < n \to s_{m} < s_{n}))$
\item $\varphi_{121} :\equiv \forall n\,(n \in B_{11} \leftrightarrow n \in \omega \land (n \in \mathbb{N} \to m < n \to s_{m} < s_{n})) \to n \in B_{11} \to S(n) \in \mathbb{N} \to n \neq \emptyset \to m < n \to s_{m} < s_{n}$
\item $\varphi_{122} :\equiv \forall n\,(n \in B_{11} \leftrightarrow n \in \omega \land (n \in \mathbb{N} \to m < n \to s_{m} < s_{n})) \to \varphi_{47}$
\item $\varphi_{123} :\equiv \forall n\,(n \in B_{11} \leftrightarrow n \in \omega \land (n \in \mathbb{N} \to m < n \to s_{m} < s_{n})) \to \varphi_{48}$
\item $\varphi_{124} :\equiv \forall n\,(n \in B_{11} \leftrightarrow n \in \omega \land (n \in \mathbb{N} \to m < n \to s_{m} < s_{n})) \to \varphi_{49}$
\item $\varphi_{125} :\equiv \forall n\,(n \in B_{11} \leftrightarrow n \in \omega \land (n \in \mathbb{N} \to m < n \to s_{m} < s_{n})) \to \varphi_{59}$
\item $\varphi_{126} :\equiv \forall n\,(n \in B_{11} \leftrightarrow n \in \omega \land (n \in \mathbb{N} \to m < n \to s_{m} < s_{n})) \to \varphi_{57}$
\item $\varphi_{127} :\equiv \forall n\,(n \in B_{11} \leftrightarrow n \in \omega \land (n \in \mathbb{N} \to m < n \to s_{m} < s_{n})) \to \varphi_{58}$
\item $\varphi_{128} :\equiv \forall n\,(n \in B_{11} \leftrightarrow n \in \omega \land (n \in \mathbb{N} \to m < n \to s_{m} < s_{n})) \to \varphi_{61}$
\item $\varphi_{129} :\equiv (n \in \omega \leftrightarrow \forall x\,(\mathrm{Ind}(x) \to n \in x)) \to \mathrm{SubIdx}(s) \to m \in \mathbb{N} \to \varphi_{79}$
\item $\varphi_{130} :\equiv (S(n) = 1_{\omega} \to 1_{\omega} = S(n)) \to \mathrm{SubIdx}(s) \to m \in \mathbb{N} \to \varphi_{95}$
\item $\varphi_{131} :\equiv (1_{\omega} = S(n) \to S(n) = 1_{\omega}) \to \mathrm{SubIdx}(s) \to m \in \mathbb{N} \to \varphi_{94}$
\item $\varphi_{132} :\equiv \forall n\,(n \in B_{11} \leftrightarrow n \in \omega \land (n \in \mathbb{N} \to m < n \to s_{m} < s_{n})) \to \varphi_{77}$
\item $\varphi_{133} :\equiv (\forall x\,(\mathrm{Ind}(x) \to n \in x) \to \mathrm{Ind}(B_{11}) \to n \in B_{11}) \to \mathrm{SubIdx}(s) \to m \in \mathbb{N} \to \varphi_{81}$
\item $\varphi_{134} :\equiv (m < S(n) \leftrightarrow m \subseteq S(n) \land m \neq S(n)) \to \mathrm{SubIdx}(s) \to m \in \mathbb{N} \to \varphi_{108}$
\item $\varphi_{135} :\equiv (\mathrm{SubIdx}(s) \to m \in \mathbb{N} \to m \in \mathbb{N}) \to \mathrm{SubIdx}(s) \to m \in \mathbb{N} \to \varphi_{114}$
\item $\varphi_{136} :\equiv (n = \emptyset \to S(n) = S(n) \to S(n) = S(\emptyset)) \to \mathrm{SubIdx}(s) \to m \in \mathbb{N} \to \varphi_{103}$
\item $\varphi_{137} :\equiv (\forall y \in \omega\,(S(y) \in \omega) \to n \in \omega \to S(n) \in \omega) \to \mathrm{SubIdx}(s) \to m \in \mathbb{N} \to \varphi_{84}$
\item $\varphi_{138} :\equiv (\mathrm{SubIdx}(s) \to m \in \mathbb{N} \to \forall k \in \mathbb{N}\,(s_{k} < s_{S(k)})) \to \varphi_{115}$
\item $\varphi_{139} :\equiv (\mathrm{Ind}(\omega) \leftrightarrow \emptyset \in \omega \land \forall y \in \omega\,(S(y) \in \omega)) \to \mathrm{SubIdx}(s) \to m \in \mathbb{N} \to \varphi_{100}$
\item $\varphi_{140} :\equiv (n = m \to s_{n} < s_{S(n)} \to s_{m} < s_{S(n)}) \to \mathrm{SubIdx}(s) \to m \in \mathbb{N} \to \varphi_{126}$
\item $\varphi_{141} :\equiv (\mathrm{SubIdx}(s) \to m \in \mathbb{N} \to \varphi_{71}) \to \varphi_{67} \to n \in B_{11} \to S(n) \in B_{11}$
\item $\varphi_{142} :\equiv (\mathrm{SubIdx}(s) \to m \in \mathbb{N} \to \varphi_{70}) \to \mathrm{SubIdx}(s) \to m \in \mathbb{N} \to \varphi_{73}$
\item $\varphi_{143} :\equiv \forall n\,(\varphi_{67} \to n \in B_{11} \to S(n) \in B_{11}) \to \forall n\,(\varphi_{67}) \to \forall n \in B_{11}\,(S(n) \in B_{11})$
\item $\varphi_{144} :\equiv (\mathrm{SubIdx}(s) \to m \in \mathbb{N} \to \varphi_{112}) \to \mathrm{SubIdx}(s) \to m \in \mathbb{N} \to \varphi_{71}$
\item $\varphi_{145} :\equiv (\mathrm{SubIdx}(s) \to m \in \mathbb{N} \to \varphi_{68}) \to \mathrm{SubIdx}(s) \to m \in \mathbb{N} \to \varphi_{118}$
\item $\varphi_{146} :\equiv (\mathrm{SubIdx}(s) \to m \in \mathbb{N} \to \varphi_{118}) \to \mathrm{SubIdx}(s) \to m \in \mathbb{N} \to \varphi_{70}$
\item $\varphi_{147} :\equiv (\mathrm{SubIdx}(s) \to m \in \mathbb{N} \to \varphi_{100}) \to \mathrm{SubIdx}(s) \to m \in \mathbb{N} \to \varphi_{89}$
\item $\varphi_{148} :\equiv (\mathrm{SubIdx}(s) \to m \in \mathbb{N} \to \varphi_{99}) \to \mathrm{SubIdx}(s) \to m \in \mathbb{N} \to \varphi_{96}$
\item $\varphi_{149} :\equiv (\mathrm{SubIdx}(s) \to m \in \mathbb{N} \to \varphi_{102}) \to \varphi_{67} \to n \in \omega \to n \in \mathbb{N} \to m < n \to s_{m} < s_{n}$
\item $\varphi_{150} :\equiv (\mathrm{SubIdx}(s) \to m \in \mathbb{N} \to \varphi_{96}) \to \mathrm{SubIdx}(s) \to m \in \mathbb{N} \to \varphi_{113}$
\item $\varphi_{151} :\equiv \forall n\,(\varphi_{67} \to n \in \omega \to n \in \mathbb{N} \to m < n \to s_{m} < s_{n}) \to \forall n\,(\varphi_{67}) \to \forall n \in \omega\,(n \in \mathbb{N} \to m < n \to s_{m} < s_{n})$
\item $\varphi_{152} :\equiv (\mathrm{SubIdx}(s) \to m \in \mathbb{N} \to \varphi_{118}) \to \mathrm{SubIdx}(s) \to m \in \mathbb{N} \to \varphi_{101}$
\item $\varphi_{153} :\equiv (\mathrm{SubIdx}(s) \to m \in \mathbb{N} \to \varphi_{96}) \to \mathrm{SubIdx}(s) \to m \in \mathbb{N} \to \varphi_{121}$
\item $\varphi_{154} :\equiv (\mathrm{SubIdx}(s) \to m \in \mathbb{N} \to \varphi_{111}) \to \mathrm{SubIdx}(s) \to m \in \mathbb{N} \to \varphi_{106}$
\item $\varphi_{155} :\equiv (\mathrm{SubIdx}(s) \to m \in \mathbb{N} \to \varphi_{117}) \to \mathrm{SubIdx}(s) \to m \in \mathbb{N} \to \varphi_{102}$
\item $\varphi_{156} :\equiv (\mathrm{SubIdx}(s) \to m \in \mathbb{N} \to \varphi_{106}) \to \mathrm{SubIdx}(s) \to m \in \mathbb{N} \to \varphi_{116}$
\item $\varphi_{157} :\equiv (\mathrm{SubIdx}(s) \to m \in \mathbb{N} \to \varphi_{94}) \to \mathrm{SubIdx}(s) \to m \in \mathbb{N} \to \varphi_{122}$
\item $\varphi_{158} :\equiv (\mathrm{SubIdx}(s) \to m \in \mathbb{N} \to \varphi_{96}) \to \mathrm{SubIdx}(s) \to m \in \mathbb{N} \to \varphi_{123}$
\item $\varphi_{159} :\equiv (\mathrm{SubIdx}(s) \to m \in \mathbb{N} \to \varphi_{109}) \to (m = n \to n = m) \to \mathrm{SubIdx}(s) \to m \in \mathbb{N} \to \varphi_{110}$
\item $\varphi_{160} :\equiv (\mathrm{SubIdx}(s) \to m \in \mathbb{N} \to \varphi_{124}) \to \mathrm{SubIdx}(s) \to m \in \mathbb{N} \to \varphi_{112}$
\item $\varphi_{161} :\equiv (\mathrm{SubIdx}(s) \to m \in \mathbb{N} \to \varphi_{119}) \to \mathrm{SubIdx}(s) \to m \in \mathbb{N} \to \varphi_{127}$
\item $\varphi_{162} :\equiv (\mathrm{SubIdx}(s) \to m \in \mathbb{N} \to \varphi_{74}) \to \varphi_{17} \to \varphi_{16} \to \mathrm{SubIdx}(s) \to m \in \mathbb{N} \to \varphi_{55}$
\item $\varphi_{163} :\equiv (\mathrm{SubIdx}(s) \to m \in \mathbb{N} \to \varphi_{100}) \to \varphi_{62} \to \mathrm{SubIdx}(s) \to m \in \mathbb{N} \to \varphi_{65}$
\item $\varphi_{164} :\equiv (\mathrm{SubIdx}(s) \to m \in \mathbb{N} \to \varphi_{126}) \to \mathrm{SubIdx}(s) \to m \in \mathbb{N} \to \varphi_{124}$
\item $\varphi_{165} :\equiv (\mathrm{SubIdx}(s) \to m \in \mathbb{N} \to \varphi_{123}) \to \mathrm{SubIdx}(s) \to m \in \mathbb{N} \to \varphi_{128}$
\item $\varphi_{166} :\equiv (\mathrm{SubIdx}(s) \to m \in \mathbb{N} \to \varphi_{103}) \to S(\emptyset) = 1_{\omega} \to \varphi_{8} \to \mathrm{SubIdx}(s) \to m \in \mathbb{N} \to \varphi_{94}$
\item $\varphi_{167} :\equiv (\mathrm{SubIdx}(s) \to m \in \mathbb{N} \to n \in \mathbb{N} \to s_{n} < s_{S(n)}) \to \varphi_{158}$
\item $\varphi_{168} :\equiv (S(n) = 1_{\omega} \to s_{m} < s_{S(n)} \to s_{m} < s_{1_{\omega}}) \to \varphi_{157}$
\item $\varphi_{169} :\equiv (n \in \mathbb{N} \leftrightarrow n \in \omega \land n \neq \emptyset) \to (\mathrm{SubIdx}(s) \to m \in \mathbb{N} \to \varphi_{70}) \to \varphi_{148}$
\item $\varphi_{170} :\equiv (\mathrm{SubIdx}(s) \to m \in \mathbb{N} \to \varphi_{120}) \to \varphi_{117} \to \mathrm{SubIdx}(s) \to m \in \mathbb{N} \to \varphi_{117}$
\item $\varphi_{171} :\equiv (1_{\omega} = S(n) \to s_{m} < s_{1_{\omega}} \to s_{m} < s_{S(n)}) \to (S(n) = 1_{\omega} \to m < S(n) \to m < 1_{\omega}) \to \varphi_{168}$
\item $\varphi_{172} :\equiv (\mathrm{SubIdx}(s) \to m \in \mathbb{N} \to \varphi_{120}) \to \varphi_{125} \to \mathrm{SubIdx}(s) \to m \in \mathbb{N} \to \varphi_{125}$
\item $\varphi_{173} :\equiv (\mathrm{SubIdx}(s) \to m \in \mathbb{N} \to \varphi_{120}) \to \varphi_{132} \to \mathrm{SubIdx}(s) \to m \in \mathbb{N} \to \varphi_{132}$
\item $\varphi_{174} :\equiv (\mathrm{SubIdx}(s) \to m \in \mathbb{N} \to \varphi_{81}) \to (\mathrm{SubIdx}(s) \to m \in \mathbb{N} \to \varphi_{55}) \to \varphi_{155}$
\item $\varphi_{175} :\equiv (\mathrm{SubIdx}(s) \to m \in \mathbb{N} \to \varphi_{110}) \to (\mathrm{SubIdx}(s) \to m \in \mathbb{N} \to \varphi_{123}) \to \varphi_{140}$
\item $\varphi_{176} :\equiv \varphi_{19} \to (1_{\omega} = S(n) \to m < 1_{\omega} \to m < S(n)) \to \varphi_{171}$
\item $\varphi_{177} :\equiv (\mathrm{SubIdx}(s) \to m \in \mathbb{N} \to \varphi_{125}) \to (\mathrm{SubIdx}(s) \to m \in \mathbb{N} \to \varphi_{73}) \to \varphi_{144}$
\item $\varphi_{178} :\equiv (\forall B_{11}\,(\varphi_{60} \to \varphi_{42}) \to \forall B_{11}\,(\varphi_{60}) \to \forall B_{11}\,(\varphi_{42})) \to (\varphi_{60} \to \forall B_{11}\,(\varphi_{60})) \to \varphi_{56}$
\item $\varphi_{179} :\equiv (s_{m} < s_{n} \to s_{n} < s_{S(n)} \to s_{m} < s_{S(n)}) \to (\mathrm{SubIdx}(s) \to m \in \mathbb{N} \to \varphi_{127}) \to \varphi_{165}$
\item $\varphi_{180} :\equiv (\mathrm{SubIdx}(s) \to m \in \mathbb{N} \to \varphi_{114}) \to (\mathrm{SubIdx}(s) \to m \in \mathbb{N} \to \varphi_{113}) \to \varphi_{156}$
\item $\varphi_{181} :\equiv (\mathrm{SubIdx}(s) \to m \in \mathbb{N} \to \varphi_{108}) \to (\mathrm{SubIdx}(s) \to m \in \mathbb{N} \to \varphi_{128}) \to \varphi_{164}$
\item $\varphi_{182} :\equiv \forall n\,(\varphi_{67} \to n \in B_{11} \to S(n) \in B_{11}) \to \varphi_{143} \to (\varphi_{67} \to \forall n\,(\varphi_{67})) \to \mathrm{SubIdx}(s) \to m \in \mathbb{N} \to \varphi_{74}$
\item $\varphi_{183} :\equiv \forall n\,(\varphi_{67} \to n \in \omega \to n \in \mathbb{N} \to m < n \to s_{m} < s_{n}) \to \varphi_{151} \to (\varphi_{67} \to \forall n\,(\varphi_{67})) \to \mathrm{SubIdx}(s) \to m \in \mathbb{N} \to \varphi_{105}$
\end{abbrevlist}
\end{fullproof}

% 自動生成 (prover/build.js): IdxGe  ノード数 326
\begin{fullproof}{[C17]}{$\forall s\,(\mathrm{SubIdx}(s) \to \forall k \in \mathbb{N}\,(k \leq s_{k}))$}{326}\phantomsection\label{tree:IdxGe}
\noindent$\Pi_{1}$:
\begin{ptree}
\AXC{$\mathrm{SubIdx}(s) \to k \in \mathbb{N} \to k \in \mathbb{N}$\lbl{A1}}
\AXC{$(\mathrm{SubIdx}(s) \to k \in \mathbb{N} \to k \in \mathbb{N}) \to \varphi_{40}$\lbl{A1}}
\BIC{$\varphi_{40}$}
\end{ptree}
\noindent$\Pi_{2}$:
\begin{ptree}
\AXC{$\varphi_{66}$\lbl{A2}}
\AXC{$\mathrm{SubIdx}(s) \to \varphi_{73}$\lbl{A1}}
\AXC{$(\mathrm{SubIdx}(s) \to \varphi_{73}) \to \varphi_{66} \to \mathrm{SubIdx}(s) \to \varphi_{66}$\lbl{A1}}
\BIC{$\varphi_{66} \to \mathrm{SubIdx}(s) \to \varphi_{66}$}
\BIC{$\mathrm{SubIdx}(s) \to \varphi_{66}$}
\end{ptree}
\noindent$\Pi_{3}$:
\begin{ptree}
\AXC{$\varphi_{66}$\lbl{A2}}
\AXC{$\mathrm{SubIdx}(s) \to \varphi_{73}$\lbl{A1}}
\AXC{$(\mathrm{SubIdx}(s) \to \varphi_{73}) \to \varphi_{66} \to \mathrm{SubIdx}(s) \to \varphi_{66}$\lbl{A1}}
\BIC{$\varphi_{66} \to \mathrm{SubIdx}(s) \to \varphi_{66}$}
\BIC{$\mathrm{SubIdx}(s) \to \varphi_{66}$}
\AXC{$\varphi_{121}$\lbl{A1}}
\BIC{$(\mathrm{SubIdx}(s) \to \varphi_{41}) \to \mathrm{SubIdx}(s) \to \varphi_{67}$}
\end{ptree}
\noindent$\Pi_{4}$:
\begin{ptree}
\AXC{$\mathrm{SubIdx}(s) \to \varphi_{41}$\lbl{A1}}
\AXC{$\stackrel{\Pi_{3}}{\vdots}$}\noLine
\UIC{$(\mathrm{SubIdx}(s) \to \varphi_{41}) \to \mathrm{SubIdx}(s) \to \varphi_{67}$}
\BIC{$\mathrm{SubIdx}(s) \to \varphi_{67}$}
\end{ptree}
\noindent$\Pi_{5}$:
\begin{ptree}
\AXC{$\stackrel{\Pi_{4}}{\vdots}$}\noLine
\UIC{$\mathrm{SubIdx}(s) \to \varphi_{67}$}
\AXC{$(\mathrm{SubIdx}(s) \to \varphi_{67}) \to \mathrm{SubIdx}(s) \to \varphi_{42}$\lbl{A1}}
\BIC{$\mathrm{SubIdx}(s) \to \varphi_{42}$}
\end{ptree}
\noindent$\Pi_{6}$:
\begin{ptree}
\AXC{$\stackrel{\Pi_{5}}{\vdots}$}\noLine
\UIC{$\mathrm{SubIdx}(s) \to \varphi_{42}$}
\AXC{$\forall n\,(n \in \mathbb{N} \leftrightarrow n \in \omega \land n \neq \emptyset)$\lbl{Def$\mathbb{N}$}}
\AXC{$\varphi_{47}$\lbl{A2}}
\BIC{$k \in \mathbb{N} \leftrightarrow k \in \omega \land k \neq \emptyset$}
\AXC{$(k \in \mathbb{N} \leftrightarrow k \in \omega \land k \neq \emptyset) \to \varphi_{122}$\lbl{A1}}
\BIC{$\varphi_{122}$}
\BIC{$(\mathrm{SubIdx}(s) \to \varphi_{76}) \to \mathrm{SubIdx}(s) \to \varphi_{71}$}
\end{ptree}
\noindent$\Pi_{7}$:
\begin{ptree}
\AXC{$\mathrm{SubIdx}(s) \to \varphi_{76}$\lbl{A1}}
\AXC{$\stackrel{\Pi_{6}}{\vdots}$}\noLine
\UIC{$(\mathrm{SubIdx}(s) \to \varphi_{76}) \to \mathrm{SubIdx}(s) \to \varphi_{71}$}
\BIC{$\mathrm{SubIdx}(s) \to \varphi_{71}$}
\end{ptree}
\noindent$\Pi_{8}$:
\begin{ptree}
\AXC{$\forall s\,(\varphi_{20})$\lbl{DefSubIdx}}
\AXC{$\forall s\,(\varphi_{20}) \to \varphi_{20}$\lbl{A2}}
\BIC{$\varphi_{20}$}
\AXC{$\mathrm{SubIdx}(s) \to \mathrm{SubIdx}(s)$\lbl{A1}}
\AXC{$\varphi_{98}$\lbl{A1}}
\BIC{$\varphi_{20} \to \mathrm{SubIdx}(s) \to \forall k \in \mathbb{N}\,(s_{k} < s_{S(k)})$}
\BIC{$\mathrm{SubIdx}(s) \to \forall k \in \mathbb{N}\,(s_{k} < s_{S(k)})$}
\end{ptree}
\noindent$\Pi_{9}$:
\begin{ptree}
\AXC{$\stackrel{\Pi_{8}}{\vdots}$}\noLine
\UIC{$\mathrm{SubIdx}(s) \to \forall k \in \mathbb{N}\,(s_{k} < s_{S(k)})$}
\AXC{$(\mathrm{SubIdx}(s) \to \forall k \in \mathbb{N}\,(s_{k} < s_{S(k)})) \to \varphi_{92}$\lbl{A1}}
\BIC{$\varphi_{92}$}
\end{ptree}
\noindent$\Pi_{10}$:
\begin{ptree}
\AXC{$\stackrel{\Pi_{7}}{\vdots}$}\noLine
\UIC{$\mathrm{SubIdx}(s) \to \varphi_{71}$}
\AXC{$\forall k \in \mathbb{N}\,(s_{k} < s_{S(k)}) \to k \in \mathbb{N} \to s_{k} < s_{S(k)}$\lbl{A2}}
\AXC{$\stackrel{\Pi_{9}}{\vdots}$}\noLine
\UIC{$\varphi_{92}$}
\BIC{$\mathrm{SubIdx}(s) \to k \in \mathbb{N} \to s_{k} < s_{S(k)}$}
\AXC{$\varphi_{118}$\lbl{A1}}
\BIC{$(\mathrm{SubIdx}(s) \to \varphi_{71}) \to \mathrm{SubIdx}(s) \to \varphi_{89}$}
\BIC{$\mathrm{SubIdx}(s) \to \varphi_{89}$}
\end{ptree}
\noindent$\Pi_{11}$:
\begin{ptree}
\AXC{$\stackrel{\Pi_{7}}{\vdots}$}\noLine
\UIC{$\mathrm{SubIdx}(s) \to \varphi_{71}$}
\AXC{$\stackrel{\Pi_{4}}{\vdots}$}\noLine
\UIC{$\mathrm{SubIdx}(s) \to \varphi_{67}$}
\AXC{$(\mathrm{SubIdx}(s) \to \varphi_{67}) \to \mathrm{SubIdx}(s) \to \varphi_{54}$\lbl{A1}}
\BIC{$\mathrm{SubIdx}(s) \to \varphi_{54}$}
\AXC{$\varphi_{123}$\lbl{A1}}
\BIC{$(\mathrm{SubIdx}(s) \to \varphi_{71}) \to \mathrm{SubIdx}(s) \to \varphi_{74}$}
\BIC{$\mathrm{SubIdx}(s) \to \varphi_{74}$}
\end{ptree}
\noindent$\Pi_{12}$:
\begin{ptree}
\AXC{$\forall x\,\forall y\,\forall z\,(x \leq y \to y < z \to x < z)$\lbl{[O4]}}
\AXC{$\varphi_{21}$\lbl{A2}}
\BIC{$\forall y\,\forall z\,(k \leq y \to y < z \to k < z)$}
\AXC{$\varphi_{32}$\lbl{A2}}
\BIC{$\forall z\,(k \leq s_{k} \to s_{k} < z \to k < z)$}
\AXC{$\varphi_{61}$\lbl{A2}}
\BIC{$k \leq s_{k} \to s_{k} < s_{S(k)} \to k < s_{S(k)}$}
\end{ptree}
\noindent$\Pi_{13}$:
\begin{ptree}
\AXC{$\stackrel{\Pi_{11}}{\vdots}$}\noLine
\UIC{$\mathrm{SubIdx}(s) \to \varphi_{74}$}
\AXC{$\stackrel{\Pi_{12}}{\vdots}$}\noLine
\UIC{$k \leq s_{k} \to s_{k} < s_{S(k)} \to k < s_{S(k)}$}
\AXC{$(k \leq s_{k} \to s_{k} < s_{S(k)} \to k < s_{S(k)}) \to \varphi_{126}$\lbl{A1}}
\BIC{$\varphi_{126}$}
\BIC{$(\mathrm{SubIdx}(s) \to \varphi_{89}) \to \mathrm{SubIdx}(s) \to \varphi_{83}$}
\end{ptree}
\noindent$\Pi_{14}$:
\begin{ptree}
\AXC{$\stackrel{\Pi_{7}}{\vdots}$}\noLine
\UIC{$\mathrm{SubIdx}(s) \to \varphi_{71}$}
\AXC{$\forall n \in \mathbb{N}\,(S(n) \in \mathbb{N})$\lbl{[N2]}}
\AXC{$\forall n \in \mathbb{N}\,(S(n) \in \mathbb{N}) \to k \in \mathbb{N} \to S(k) \in \mathbb{N}$\lbl{A2}}
\BIC{$k \in \mathbb{N} \to S(k) \in \mathbb{N}$}
\AXC{$\varphi_{113}$\lbl{A1}}
\BIC{$(\mathrm{SubIdx}(s) \to \varphi_{71}) \to \mathrm{SubIdx}(s) \to \varphi_{80}$}
\BIC{$\mathrm{SubIdx}(s) \to \varphi_{80}$}
\end{ptree}
\noindent$\Pi_{15}$:
\begin{ptree}
\AXC{$\forall s\,(\varphi_{20})$\lbl{DefSubIdx}}
\AXC{$\forall s\,(\varphi_{20}) \to \varphi_{20}$\lbl{A2}}
\BIC{$\varphi_{20}$}
\AXC{$\mathrm{SubIdx}(s) \to \mathrm{SubIdx}(s)$\lbl{A1}}
\AXC{$\varphi_{75}$\lbl{A1}}
\BIC{$\varphi_{20} \to \mathrm{SubIdx}(s) \to \mathrm{NSeq}(s)$}
\BIC{$\mathrm{SubIdx}(s) \to \mathrm{NSeq}(s)$}
\end{ptree}
\noindent$\Pi_{16}$:
\begin{ptree}
\AXC{$\forall s\,\forall k\,(\mathrm{NSeq}(s) \to k \in \mathbb{N} \to s_{k} \in \mathbb{N})$\lbl{[S13]}}
\AXC{$\varphi_{51}$\lbl{A2}}
\BIC{$\forall k\,(\mathrm{NSeq}(s) \to k \in \mathbb{N} \to s_{k} \in \mathbb{N})$}
\AXC{$\varphi_{52}$\lbl{A2}}
\BIC{$\mathrm{NSeq}(s) \to S(k) \in \mathbb{N} \to s_{S(k)} \in \mathbb{N}$}
\end{ptree}
\noindent$\Pi_{17}$:
\begin{ptree}
\AXC{$\stackrel{\Pi_{16}}{\vdots}$}\noLine
\UIC{$\mathrm{NSeq}(s) \to S(k) \in \mathbb{N} \to s_{S(k)} \in \mathbb{N}$}
\AXC{$(\mathrm{NSeq}(s) \to S(k) \in \mathbb{N} \to s_{S(k)} \in \mathbb{N}) \to \varphi_{114}$\lbl{A1}}
\BIC{$\varphi_{114}$}
\end{ptree}
\noindent$\Pi_{18}$:
\begin{ptree}
\AXC{$\forall n \in \mathbb{N}\,(n \in \omega)$\lbl{[N7']}}
\AXC{$\forall n \in \mathbb{N}\,(n \in \omega) \to s_{S(k)} \in \mathbb{N} \to s_{S(k)} \in \omega$\lbl{A2}}
\BIC{$s_{S(k)} \in \mathbb{N} \to s_{S(k)} \in \omega$}
\AXC{$\varphi_{119}$\lbl{A1}}
\BIC{$(\mathrm{SubIdx}(s) \to \varphi_{85}) \to \mathrm{SubIdx}(s) \to \varphi_{87}$}
\end{ptree}
\noindent$\Pi_{19}$:
\begin{ptree}
\AXC{$\stackrel{\Pi_{14}}{\vdots}$}\noLine
\UIC{$\mathrm{SubIdx}(s) \to \varphi_{80}$}
\AXC{$\stackrel{\Pi_{15}}{\vdots}$}\noLine
\UIC{$\mathrm{SubIdx}(s) \to \mathrm{NSeq}(s)$}
\AXC{$\stackrel{\Pi_{17}}{\vdots}$}\noLine
\UIC{$\varphi_{114}$}
\BIC{$(\mathrm{SubIdx}(s) \to \varphi_{80}) \to \mathrm{SubIdx}(s) \to \varphi_{85}$}
\BIC{$\mathrm{SubIdx}(s) \to \varphi_{85}$}
\AXC{$\stackrel{\Pi_{18}}{\vdots}$}\noLine
\UIC{$(\mathrm{SubIdx}(s) \to \varphi_{85}) \to \mathrm{SubIdx}(s) \to \varphi_{87}$}
\BIC{$\mathrm{SubIdx}(s) \to \varphi_{87}$}
\end{ptree}
\noindent$\Pi_{20}$:
\begin{ptree}
\AXC{$\stackrel{\Pi_{7}}{\vdots}$}\noLine
\UIC{$\mathrm{SubIdx}(s) \to \varphi_{71}$}
\AXC{$\forall n \in \mathbb{N}\,(n \in \omega)$\lbl{[N7']}}
\AXC{$\forall n \in \mathbb{N}\,(n \in \omega) \to k \in \mathbb{N} \to k \in \omega$\lbl{A2}}
\BIC{$k \in \mathbb{N} \to k \in \omega$}
\AXC{$\varphi_{112}$\lbl{A1}}
\BIC{$(\mathrm{SubIdx}(s) \to \varphi_{71}) \to \mathrm{SubIdx}(s) \to \varphi_{77}$}
\BIC{$\mathrm{SubIdx}(s) \to \varphi_{77}$}
\end{ptree}
\noindent$\Pi_{21}$:
\begin{ptree}
\AXC{$\stackrel{\Pi_{20}}{\vdots}$}\noLine
\UIC{$\mathrm{SubIdx}(s) \to \varphi_{77}$}
\AXC{$\forall m\,\forall n\,(m \in \omega \to n \in \omega \to m < n \to S(m) \leq n)$\lbl{[N10]}}
\AXC{$\varphi_{64}$\lbl{A2}}
\BIC{$\forall n\,(k \in \omega \to n \in \omega \to k < n \to S(k) \leq n)$}
\AXC{$\varphi_{91}$\lbl{A2}}
\BIC{$k \in \omega \to s_{S(k)} \in \omega \to k < s_{S(k)} \to S(k) \leq s_{S(k)}$}
\AXC{$\varphi_{130}$\lbl{A1}}
\BIC{$(\mathrm{SubIdx}(s) \to \varphi_{77}) \to \varphi_{127}$}
\BIC{$\varphi_{127}$}
\end{ptree}
\noindent$\Pi_{22}$:
\begin{ptree}
\AXC{$\stackrel{\Pi_{10}}{\vdots}$}\noLine
\UIC{$\mathrm{SubIdx}(s) \to \varphi_{89}$}
\AXC{$\stackrel{\Pi_{13}}{\vdots}$}\noLine
\UIC{$(\mathrm{SubIdx}(s) \to \varphi_{89}) \to \mathrm{SubIdx}(s) \to \varphi_{83}$}
\BIC{$\mathrm{SubIdx}(s) \to \varphi_{83}$}
\AXC{$\stackrel{\Pi_{19}}{\vdots}$}\noLine
\UIC{$\mathrm{SubIdx}(s) \to \varphi_{87}$}
\AXC{$\stackrel{\Pi_{21}}{\vdots}$}\noLine
\UIC{$\varphi_{127}$}
\BIC{$(\mathrm{SubIdx}(s) \to \varphi_{83}) \to \mathrm{SubIdx}(s) \to \varphi_{84}$}
\BIC{$\mathrm{SubIdx}(s) \to \varphi_{84}$}
\end{ptree}
\noindent$\Pi_{23}$:
\begin{ptree}
\AXC{$\stackrel{\Pi_{15}}{\vdots}$}\noLine
\UIC{$\mathrm{SubIdx}(s) \to \mathrm{NSeq}(s)$}
\AXC{$\forall s\,\forall k\,(\mathrm{NSeq}(s) \to k \in \mathbb{N} \to s_{k} \in \mathbb{N})$\lbl{[S13]}}
\AXC{$\varphi_{51}$\lbl{A2}}
\BIC{$\forall k\,(\mathrm{NSeq}(s) \to k \in \mathbb{N} \to s_{k} \in \mathbb{N})$}
\AXC{$\varphi_{53}$\lbl{A2}}
\BIC{$\mathrm{NSeq}(s) \to 1_{\omega} \in \mathbb{N} \to s_{1_{\omega}} \in \mathbb{N}$}
\AXC{$\varphi_{86}$\lbl{A1}}
\BIC{$\varphi_{33}$}
\BIC{$1_{\omega} \in \mathbb{N} \to \mathrm{SubIdx}(s) \to s_{1_{\omega}} \in \mathbb{N}$}
\end{ptree}
\noindent$\Pi_{24}$:
\begin{ptree}
\AXC{$\forall n \in \mathbb{N}\,(1_{\omega} \leq n)$\lbl{[N12]}}
\AXC{$\varphi_{24}$\lbl{A2}}
\BIC{$s_{1_{\omega}} \in \mathbb{N} \to 1_{\omega} \leq s_{1_{\omega}}$}
\AXC{$(s_{1_{\omega}} \in \mathbb{N} \to 1_{\omega} \leq s_{1_{\omega}}) \to \varphi_{25}$\lbl{A1}}
\BIC{$\varphi_{25}$}
\end{ptree}
\noindent$\Pi_{25}$:
\begin{ptree}
\AXC{$x = x$\lbl{A5}}
\AXC{$x = y \to x = x \to y = x$\lbl{A6}}
\AXC{$(x = y \to x = x \to y = x) \to x = x \to x = y \to y = x$\lbl{A1}}
\BIC{$x = x \to x = y \to y = x$}
\BIC{$x = y \to y = x$}
\UIC{$\forall y\,(x = y \to y = x)$}
\UIC{$\forall x\,\forall y\,(x = y \to y = x)$}
\end{ptree}
\noindent$\Pi_{26}$:
\begin{ptree}
\AXC{$y = z \to x = y \to x = z$\lbl{A6}}
\AXC{$(y = z \to x = y \to x = z) \to x = y \to y = z \to x = z$\lbl{A1}}
\BIC{$x = y \to y = z \to x = z$}
\UIC{$\forall z\,(x = y \to y = z \to x = z)$}
\UIC{$\forall y\,\forall z\,(x = y \to y = z \to x = z)$}
\UIC{$\forall x\,\forall y\,\forall z\,(x = y \to y = z \to x = z)$}
\AXC{$\varphi_{9}$\lbl{A2}}
\BIC{$\forall y\,\forall z\,(S(k) = y \to y = z \to S(k) = z)$}
\AXC{$\varphi_{38}$\lbl{A2}}
\BIC{$\forall z\,(S(k) = S(\emptyset) \to S(\emptyset) = z \to S(k) = z)$}
\end{ptree}
\noindent$\Pi_{27}$:
\begin{ptree}
\AXC{$\stackrel{\Pi_{26}}{\vdots}$}\noLine
\UIC{$\forall z\,(S(k) = S(\emptyset) \to S(\emptyset) = z \to S(k) = z)$}
\AXC{$\forall z\,(S(k) = S(\emptyset) \to S(\emptyset) = z \to S(k) = z) \to \varphi_{5}$\lbl{A2}}
\BIC{$\varphi_{5}$}
\end{ptree}
\noindent$\Pi_{28}$:
\begin{ptree}
\AXC{$\stackrel{\Pi_{25}}{\vdots}$}\noLine
\UIC{$\forall x\,\forall y\,(x = y \to y = x)$}
\AXC{$\varphi_{2}$\lbl{A2}}
\BIC{$\forall y\,(1_{\omega} = y \to y = 1_{\omega})$}
\AXC{$\varphi_{11}$\lbl{A2}}
\BIC{$1_{\omega} = S(\emptyset) \to S(\emptyset) = 1_{\omega}$}
\AXC{$1_{\omega} = S(\emptyset)$\lbl{Def$1_\omega$}}
\AXC{$1_{\omega} = S(\emptyset) \to \varphi_{10}$\lbl{A1}}
\BIC{$\varphi_{10}$}
\BIC{$S(\emptyset) = 1_{\omega}$}
\end{ptree}
\noindent$\Pi_{29}$:
\begin{ptree}
\AXC{$\mathrm{SubIdx}(s) \to \varphi_{72}$\lbl{A1}}
\AXC{$(\mathrm{SubIdx}(s) \to \varphi_{72}) \to S(k) = S(k) \to \varphi_{104}$\lbl{A1}}
\BIC{$S(k) = S(k) \to \varphi_{104}$}
\end{ptree}
\noindent$\Pi_{30}$:
\begin{ptree}
\AXC{$x = x$\lbl{A5}}
\UIC{$\forall x\,(x = x)$}
\AXC{$\forall x\,(x = x) \to S(k) = S(k)$\lbl{A2}}
\BIC{$S(k) = S(k)$}
\AXC{$\stackrel{\Pi_{29}}{\vdots}$}\noLine
\UIC{$S(k) = S(k) \to \varphi_{104}$}
\BIC{$\varphi_{104}$}
\end{ptree}
\noindent$\Pi_{31}$:
\begin{ptree}
\AXC{$z = w \to S(k) = S(z) \to S(k) = S(w)$\lbl{A6}}
\UIC{$\forall w\,(z = w \to S(k) = S(z) \to S(k) = S(w))$}
\UIC{$\forall z\,\forall w\,(z = w \to S(k) = S(z) \to S(k) = S(w))$}
\AXC{$\varphi_{30}$\lbl{A2}}
\BIC{$\forall w\,(k = w \to S(k) = S(k) \to S(k) = S(w))$}
\AXC{$\varphi_{39}$\lbl{A2}}
\BIC{$k = \emptyset \to S(k) = S(k) \to S(k) = S(\emptyset)$}
\AXC{$\stackrel{\Pi_{30}}{\vdots}$}\noLine
\UIC{$\varphi_{104}$}
\BIC{$\mathrm{SubIdx}(s) \to \varphi_{81}$}
\AXC{$\varphi_{120}$\lbl{A1}}
\BIC{$S(\emptyset) = 1_{\omega} \to \varphi_{5} \to \mathrm{SubIdx}(s) \to \varphi_{69}$}
\end{ptree}
\noindent$\Pi_{32}$:
\begin{ptree}
\AXC{$\stackrel{\Pi_{27}}{\vdots}$}\noLine
\UIC{$\varphi_{5}$}
\AXC{$\stackrel{\Pi_{28}}{\vdots}$}\noLine
\UIC{$S(\emptyset) = 1_{\omega}$}
\AXC{$\stackrel{\Pi_{31}}{\vdots}$}\noLine
\UIC{$S(\emptyset) = 1_{\omega} \to \varphi_{5} \to \mathrm{SubIdx}(s) \to \varphi_{69}$}
\BIC{$\varphi_{5} \to \mathrm{SubIdx}(s) \to \varphi_{69}$}
\BIC{$\mathrm{SubIdx}(s) \to \varphi_{69}$}
\AXC{$(\mathrm{SubIdx}(s) \to \varphi_{69}) \to \varphi_{100}$\lbl{A1}}
\BIC{$\varphi_{100}$}
\end{ptree}
\noindent$\Pi_{33}$:
\begin{ptree}
\AXC{$\stackrel{\Pi_{25}}{\vdots}$}\noLine
\UIC{$\forall x\,\forall y\,(x = y \to y = x)$}
\AXC{$\varphi_{1}$\lbl{A2}}
\BIC{$\forall y\,(S(k) = y \to y = S(k))$}
\AXC{$\varphi_{3}$\lbl{A2}}
\BIC{$S(k) = 1_{\omega} \to 1_{\omega} = S(k)$}
\AXC{$\stackrel{\Pi_{32}}{\vdots}$}\noLine
\UIC{$\varphi_{100}$}
\BIC{$\mathrm{SubIdx}(s) \to \varphi_{70}$}
\AXC{$(\mathrm{SubIdx}(s) \to \varphi_{70}) \to \varphi_{111}$\lbl{A1}}
\BIC{$\varphi_{111}$}
\end{ptree}
\noindent$\Pi_{34}$:
\begin{ptree}
\AXC{$1_{\omega} \in \mathbb{N}$\lbl{[N1]}}
\AXC{$\stackrel{\Pi_{23}}{\vdots}$}\noLine
\UIC{$1_{\omega} \in \mathbb{N} \to \mathrm{SubIdx}(s) \to s_{1_{\omega}} \in \mathbb{N}$}
\BIC{$\mathrm{SubIdx}(s) \to s_{1_{\omega}} \in \mathbb{N}$}
\AXC{$\stackrel{\Pi_{24}}{\vdots}$}\noLine
\UIC{$\varphi_{25}$}
\BIC{$\mathrm{SubIdx}(s) \to 1_{\omega} \leq s_{1_{\omega}}$}
\AXC{$\stackrel{\Pi_{33}}{\vdots}$}\noLine
\UIC{$\varphi_{111}$}
\BIC{$\varphi_{6} \to \varphi_{4} \to \mathrm{SubIdx}(s) \to \varphi_{82}$}
\end{ptree}
\noindent$\Pi_{35}$:
\begin{ptree}
\AXC{$z = w \to z < s_{z} \to w < s_{w}$\lbl{A6}}
\UIC{$\forall w\,(z = w \to z < s_{z} \to w < s_{w})$}
\UIC{$\forall z\,\forall w\,(z = w \to z < s_{z} \to w < s_{w})$}
\AXC{$\varphi_{58}$\lbl{A2}}
\BIC{$\forall w\,(1_{\omega} = w \to 1_{\omega} < s_{1_{\omega}} \to w < s_{w})$}
\AXC{$\varphi_{68}$\lbl{A2}}
\BIC{$\varphi_{6}$}
\AXC{$\stackrel{\Pi_{34}}{\vdots}$}\noLine
\UIC{$\varphi_{6} \to \varphi_{4} \to \mathrm{SubIdx}(s) \to \varphi_{82}$}
\BIC{$\varphi_{4} \to \mathrm{SubIdx}(s) \to \varphi_{82}$}
\end{ptree}
\noindent$\Pi_{36}$:
\begin{ptree}
\AXC{$z = w \to z = s_{z} \to w = s_{w}$\lbl{A6}}
\UIC{$\forall w\,(z = w \to z = s_{z} \to w = s_{w})$}
\UIC{$\forall z\,\forall w\,(z = w \to z = s_{z} \to w = s_{w})$}
\AXC{$\varphi_{44}$\lbl{A2}}
\BIC{$\forall w\,(1_{\omega} = w \to 1_{\omega} = s_{1_{\omega}} \to w = s_{w})$}
\AXC{$\varphi_{57}$\lbl{A2}}
\BIC{$\varphi_{4}$}
\AXC{$\stackrel{\Pi_{35}}{\vdots}$}\noLine
\UIC{$\varphi_{4} \to \mathrm{SubIdx}(s) \to \varphi_{82}$}
\BIC{$\mathrm{SubIdx}(s) \to \varphi_{82}$}
\AXC{$\varphi_{124}$\lbl{A1}}
\BIC{$(\mathrm{SubIdx}(s) \to \varphi_{84}) \to \mathrm{SubIdx}(s) \to \varphi_{62}$}
\end{ptree}
\noindent$\Pi_{37}$:
\begin{ptree}
\AXC{$\stackrel{\Pi_{22}}{\vdots}$}\noLine
\UIC{$\mathrm{SubIdx}(s) \to \varphi_{84}$}
\AXC{$\stackrel{\Pi_{36}}{\vdots}$}\noLine
\UIC{$(\mathrm{SubIdx}(s) \to \varphi_{84}) \to \mathrm{SubIdx}(s) \to \varphi_{62}$}
\BIC{$\mathrm{SubIdx}(s) \to \varphi_{62}$}
\end{ptree}
\noindent$\Pi_{38}$:
\begin{ptree}
\AXC{$\forall x\,(\mathrm{Ind}(x) \leftrightarrow \emptyset \in x \land \forall y \in x\,(S(y) \in x))$\lbl{DefInd}}
\AXC{$\varphi_{90}$\lbl{A2}}
\BIC{$\mathrm{Ind}(\omega) \leftrightarrow \emptyset \in \omega \land \forall y \in \omega\,(S(y) \in \omega)$}
\AXC{$\mathrm{Ind}(\omega)$\lbl{補題(ωは帰納的)}}
\AXC{$\mathrm{Ind}(\omega) \to \varphi_{107}$\lbl{A1}}
\BIC{$\varphi_{107}$}
\BIC{$\mathrm{SubIdx}(s) \to \varphi_{78}$}
\end{ptree}
\noindent$\Pi_{39}$:
\begin{ptree}
\AXC{$\stackrel{\Pi_{38}}{\vdots}$}\noLine
\UIC{$\mathrm{SubIdx}(s) \to \varphi_{78}$}
\AXC{$(\mathrm{SubIdx}(s) \to \varphi_{78}) \to \mathrm{SubIdx}(s) \to \varphi_{63}$\lbl{A1}}
\BIC{$\mathrm{SubIdx}(s) \to \varphi_{63}$}
\AXC{$(\mathrm{SubIdx}(s) \to \varphi_{63}) \to \varphi_{106}$\lbl{A1}}
\BIC{$\varphi_{106}$}
\end{ptree}
\noindent$\Pi_{40}$:
\begin{ptree}
\AXC{$\stackrel{\Pi_{5}}{\vdots}$}\noLine
\UIC{$\mathrm{SubIdx}(s) \to \varphi_{42}$}
\AXC{$\forall y \in \omega\,(S(y) \in \omega) \to k \in \omega \to S(k) \in \omega$\lbl{A2}}
\AXC{$\stackrel{\Pi_{39}}{\vdots}$}\noLine
\UIC{$\varphi_{106}$}
\BIC{$\mathrm{SubIdx}(s) \to \varphi_{60}$}
\AXC{$\varphi_{115}$\lbl{A1}}
\BIC{$(\mathrm{SubIdx}(s) \to \varphi_{42}) \to \mathrm{SubIdx}(s) \to \varphi_{48}$}
\BIC{$\mathrm{SubIdx}(s) \to \varphi_{48}$}
\end{ptree}
\noindent$\Pi_{41}$:
\begin{ptree}
\AXC{$\forall k\,(k \in B_{15} \leftrightarrow k \in \omega \land (k \in \mathbb{N} \to k \leq s_{k})) \to \varphi_{28}$\lbl{A2}}
\AXC{$\mathrm{SubIdx}(s) \to \varphi_{73}$\lbl{A1}}
\AXC{$(\mathrm{SubIdx}(s) \to \varphi_{73}) \to \varphi_{110}$\lbl{A1}}
\BIC{$\varphi_{110}$}
\BIC{$\varphi_{94}$}
\end{ptree}
\noindent$\Pi_{42}$:
\begin{ptree}
\AXC{$\stackrel{\Pi_{37}}{\vdots}$}\noLine
\UIC{$\mathrm{SubIdx}(s) \to \varphi_{62}$}
\AXC{$\stackrel{\Pi_{40}}{\vdots}$}\noLine
\UIC{$\mathrm{SubIdx}(s) \to \varphi_{48}$}
\AXC{$\stackrel{\Pi_{41}}{\vdots}$}\noLine
\UIC{$\varphi_{94}$}
\AXC{$\varphi_{94} \to \varphi_{117}$\lbl{A1}}
\BIC{$\varphi_{117}$}
\BIC{$(\mathrm{SubIdx}(s) \to \varphi_{62}) \to \mathrm{SubIdx}(s) \to \varphi_{43}$}
\BIC{$\mathrm{SubIdx}(s) \to \varphi_{43}$}
\end{ptree}
\noindent$\Pi_{43}$:
\begin{ptree}
\AXC{$\varphi_{108}$\lbl{A3}}
\AXC{$\stackrel{\Pi_{42}}{\vdots}$}\noLine
\UIC{$\mathrm{SubIdx}(s) \to \varphi_{43}$}
\AXC{$(\mathrm{SubIdx}(s) \to \varphi_{43}) \to \varphi_{27} \to k \in B_{15} \to S(k) \in B_{15}$\lbl{A1}}
\BIC{$\varphi_{27} \to k \in B_{15} \to S(k) \in B_{15}$}
\UIC{$\forall k\,(\varphi_{27} \to k \in B_{15} \to S(k) \in B_{15})$}
\AXC{$\varphi_{128}$\lbl{A1}}
\BIC{$\varphi_{108} \to (\varphi_{27} \to \forall k\,(\varphi_{27})) \to \mathrm{SubIdx}(s) \to \varphi_{49}$}
\BIC{$(\varphi_{27} \to \forall k\,(\varphi_{27})) \to \mathrm{SubIdx}(s) \to \varphi_{49}$}
\end{ptree}
\noindent$\Pi_{44}$:
\begin{ptree}
\AXC{$\varphi_{99}$\lbl{A2}}
\AXC{$\mathrm{SubIdx}(s) \to \varphi_{73}$\lbl{A1}}
\AXC{$(\mathrm{SubIdx}(s) \to \varphi_{73}) \to \varphi_{99} \to \mathrm{SubIdx}(s) \to \varphi_{99}$\lbl{A1}}
\BIC{$\varphi_{99} \to \mathrm{SubIdx}(s) \to \varphi_{99}$}
\BIC{$\mathrm{SubIdx}(s) \to \varphi_{99}$}
\end{ptree}
\noindent$\Pi_{45}$:
\begin{ptree}
\AXC{$\stackrel{\Pi_{38}}{\vdots}$}\noLine
\UIC{$\mathrm{SubIdx}(s) \to \varphi_{78}$}
\AXC{$\stackrel{\Pi_{44}}{\vdots}$}\noLine
\UIC{$\mathrm{SubIdx}(s) \to \varphi_{99}$}
\AXC{$(\mathrm{SubIdx}(s) \to \varphi_{99}) \to (\mathrm{SubIdx}(s) \to \varphi_{78}) \to \varphi_{101}$\lbl{A1}}
\BIC{$(\mathrm{SubIdx}(s) \to \varphi_{78}) \to \varphi_{101}$}
\BIC{$\varphi_{101}$}
\end{ptree}
\noindent$\Pi_{46}$:
\begin{ptree}
\AXC{$x = x$\lbl{A5}}
\UIC{$\forall x\,(x = x)$}
\AXC{$\forall x\,(x = x) \to \emptyset = \emptyset$\lbl{A2}}
\BIC{$\emptyset = \emptyset$}
\AXC{$\forall n\,(n \in \mathbb{N} \leftrightarrow n \in \omega \land n \neq \emptyset)$\lbl{Def$\mathbb{N}$}}
\AXC{$\varphi_{65}$\lbl{A2}}
\BIC{$\emptyset \in \mathbb{N} \leftrightarrow \emptyset \in \omega \land \emptyset \neq \emptyset$}
\AXC{$\varphi_{97}$\lbl{A1}}
\BIC{$\emptyset = \emptyset \to \mathrm{SubIdx}(s) \to \emptyset \in \mathbb{N} \to \emptyset \leq s_{\emptyset}$}
\BIC{$\mathrm{SubIdx}(s) \to \emptyset \in \mathbb{N} \to \emptyset \leq s_{\emptyset}$}
\AXC{$\stackrel{\Pi_{45}}{\vdots}$}\noLine
\UIC{$\varphi_{101}$}
\BIC{$\mathrm{SubIdx}(s) \to \varphi_{36}$}
\end{ptree}
\noindent$\Pi_{47}$:
\begin{ptree}
\AXC{$\forall x\,(\mathrm{Ind}(x) \leftrightarrow \emptyset \in x \land \forall y \in x\,(S(y) \in x))$\lbl{DefInd}}
\AXC{$\forall x\,(\mathrm{Ind}(x) \leftrightarrow \emptyset \in x \land \forall y \in x\,(S(y) \in x)) \to \varphi_{14}$\lbl{A2}}
\BIC{$\varphi_{14}$}
\AXC{$\varphi_{14} \to \varphi_{125}$\lbl{A1}}
\BIC{$\varphi_{125}$}
\end{ptree}
\noindent$\Pi_{48}$:
\begin{ptree}
\AXC{$\varphi_{27} \to \forall k\,(\varphi_{27})$\lbl{A4}}
\AXC{$\stackrel{\Pi_{43}}{\vdots}$}\noLine
\UIC{$(\varphi_{27} \to \forall k\,(\varphi_{27})) \to \mathrm{SubIdx}(s) \to \varphi_{49}$}
\BIC{$\mathrm{SubIdx}(s) \to \varphi_{49}$}
\AXC{$\stackrel{\Pi_{46}}{\vdots}$}\noLine
\UIC{$\mathrm{SubIdx}(s) \to \varphi_{36}$}
\AXC{$\stackrel{\Pi_{47}}{\vdots}$}\noLine
\UIC{$\varphi_{125}$}
\BIC{$(\mathrm{SubIdx}(s) \to \varphi_{49}) \to \varphi_{13} \to \varphi_{12} \to \mathrm{SubIdx}(s) \to \varphi_{26}$}
\BIC{$\varphi_{13} \to \varphi_{12} \to \mathrm{SubIdx}(s) \to \varphi_{26}$}
\end{ptree}
\noindent$\Pi_{49}$:
\begin{ptree}
\AXC{$\varphi_{19}$\lbl{A4}}
\AXC{$\forall k\,(\varphi_{8}) \to \varphi_{23}$\lbl{A3}}
\AXC{$\varphi_{8}$\lbl{A2}}
\UIC{$\forall k\,(\varphi_{8})$}
\AXC{$\forall k\,(\varphi_{8}) \to (\forall k\,(\varphi_{8}) \to \varphi_{23}) \to \varphi_{19} \to \varphi_{13}$\lbl{A1}}
\BIC{$(\forall k\,(\varphi_{8}) \to \varphi_{23}) \to \varphi_{19} \to \varphi_{13}$}
\BIC{$\varphi_{19} \to \varphi_{13}$}
\BIC{$\varphi_{13}$}
\AXC{$\stackrel{\Pi_{48}}{\vdots}$}\noLine
\UIC{$\varphi_{13} \to \varphi_{12} \to \mathrm{SubIdx}(s) \to \varphi_{26}$}
\BIC{$\varphi_{12} \to \mathrm{SubIdx}(s) \to \varphi_{26}$}
\end{ptree}
\noindent$\Pi_{50}$:
\begin{ptree}
\AXC{$\varphi_{18}$\lbl{A4}}
\AXC{$\forall y\,(\varphi_{7}) \to \varphi_{22}$\lbl{A3}}
\AXC{$\varphi_{7}$\lbl{A2}}
\UIC{$\forall y\,(\varphi_{7})$}
\AXC{$\forall y\,(\varphi_{7}) \to (\forall y\,(\varphi_{7}) \to \varphi_{22}) \to \varphi_{18} \to \varphi_{12}$\lbl{A1}}
\BIC{$(\forall y\,(\varphi_{7}) \to \varphi_{22}) \to \varphi_{18} \to \varphi_{12}$}
\BIC{$\varphi_{18} \to \varphi_{12}$}
\BIC{$\varphi_{12}$}
\AXC{$\stackrel{\Pi_{49}}{\vdots}$}\noLine
\UIC{$\varphi_{12} \to \mathrm{SubIdx}(s) \to \varphi_{26}$}
\BIC{$\mathrm{SubIdx}(s) \to \varphi_{26}$}
\end{ptree}
\noindent$\Pi_{51}$:
\begin{ptree}
\AXC{$\forall n\,(n \in \omega \leftrightarrow \forall x\,(\mathrm{Ind}(x) \to n \in x))$\lbl{Def$\omega$}}
\AXC{$\varphi_{37}$\lbl{A2}}
\BIC{$k \in \omega \leftrightarrow \forall x\,(\mathrm{Ind}(x) \to k \in x)$}
\AXC{$\mathrm{SubIdx}(s) \to \varphi_{45}$\lbl{A1}}
\AXC{$(\mathrm{SubIdx}(s) \to \varphi_{45}) \to \varphi_{96}$\lbl{A1}}
\BIC{$\varphi_{96}$}
\BIC{$\mathrm{SubIdx}(s) \to \varphi_{50}$}
\end{ptree}
\noindent$\Pi_{52}$:
\begin{ptree}
\AXC{$\forall x\,(\mathrm{Ind}(x) \to k \in x) \to \mathrm{Ind}(B_{15}) \to k \in B_{15}$\lbl{A2}}
\AXC{$\stackrel{\Pi_{51}}{\vdots}$}\noLine
\UIC{$\mathrm{SubIdx}(s) \to \varphi_{50}$}
\AXC{$(\mathrm{SubIdx}(s) \to \varphi_{50}) \to \varphi_{102}$\lbl{A1}}
\BIC{$\varphi_{102}$}
\BIC{$\mathrm{SubIdx}(s) \to \varphi_{55}$}
\end{ptree}
\noindent$\Pi_{53}$:
\begin{ptree}
\AXC{$\stackrel{\Pi_{2}}{\vdots}$}\noLine
\UIC{$\mathrm{SubIdx}(s) \to \varphi_{66}$}
\AXC{$\stackrel{\Pi_{50}}{\vdots}$}\noLine
\UIC{$\mathrm{SubIdx}(s) \to \varphi_{26}$}
\AXC{$\stackrel{\Pi_{52}}{\vdots}$}\noLine
\UIC{$\mathrm{SubIdx}(s) \to \varphi_{55}$}
\AXC{$(\mathrm{SubIdx}(s) \to \varphi_{55}) \to \varphi_{116}$\lbl{A1}}
\BIC{$\varphi_{116}$}
\BIC{$(\mathrm{SubIdx}(s) \to \varphi_{66}) \to \mathrm{SubIdx}(s) \to \varphi_{56}$}
\BIC{$\mathrm{SubIdx}(s) \to \varphi_{56}$}
\end{ptree}
\noindent$\Pi_{54}$:
\begin{ptree}
\AXC{$\varphi_{109}$\lbl{A3}}
\AXC{$\stackrel{\Pi_{53}}{\vdots}$}\noLine
\UIC{$\mathrm{SubIdx}(s) \to \varphi_{56}$}
\AXC{$(\mathrm{SubIdx}(s) \to \varphi_{56}) \to \varphi_{27} \to k \in \omega \to k \in \mathbb{N} \to k \leq s_{k}$\lbl{A1}}
\BIC{$\varphi_{27} \to k \in \omega \to k \in \mathbb{N} \to k \leq s_{k}$}
\UIC{$\forall k\,(\varphi_{27} \to k \in \omega \to k \in \mathbb{N} \to k \leq s_{k})$}
\AXC{$\varphi_{129}$\lbl{A1}}
\BIC{$\varphi_{109} \to (\varphi_{27} \to \forall k\,(\varphi_{27})) \to \mathrm{SubIdx}(s) \to \varphi_{59}$}
\BIC{$(\varphi_{27} \to \forall k\,(\varphi_{27})) \to \mathrm{SubIdx}(s) \to \varphi_{59}$}
\end{ptree}
\noindent$\Pi_{55}$:
\begin{ptree}
\AXC{$\varphi_{27} \to \forall k\,(\varphi_{27})$\lbl{A4}}
\AXC{$\stackrel{\Pi_{54}}{\vdots}$}\noLine
\UIC{$(\varphi_{27} \to \forall k\,(\varphi_{27})) \to \mathrm{SubIdx}(s) \to \varphi_{59}$}
\BIC{$\mathrm{SubIdx}(s) \to \varphi_{59}$}
\AXC{$(\mathrm{SubIdx}(s) \to \varphi_{59}) \to \varphi_{79}$\lbl{A1}}
\BIC{$\varphi_{79}$}
\UIC{$\forall B_{15}\,(\varphi_{79})$}
\end{ptree}
\noindent$\Pi_{56}$:
\begin{ptree}
\AXC{$\varphi_{31}$\lbl{Sep}}
\AXC{$\varphi_{31} \to \varphi_{17}$\lbl{A2}}
\BIC{$\varphi_{17}$}
\AXC{$\varphi_{88}$\lbl{A2}}
\BIC{$\exists B_{15}\,\forall k\,(k \in B_{15} \leftrightarrow k \in \omega \land (k \in \mathbb{N} \to k \leq s_{k}))$}
\AXC{$\varphi_{93}$\lbl{A1}}
\BIC{$\varphi_{34}$}
\AXC{$\varphi_{34} \to \forall B_{15}\,(\varphi_{79}) \to (\forall B_{15}\,(\varphi_{79}) \to \varphi_{15} \to \varphi_{16}) \to \varphi_{105}$\lbl{A1}}
\BIC{$\forall B_{15}\,(\varphi_{79}) \to (\forall B_{15}\,(\varphi_{79}) \to \varphi_{15} \to \varphi_{16}) \to \varphi_{105}$}
\end{ptree}
\noindent$\Pi_{57}$:
\begin{ptree}
\AXC{$\forall B_{15}\,(\varphi_{79}) \to \varphi_{15} \to \varphi_{16}$\lbl{A3}}
\AXC{$\stackrel{\Pi_{55}}{\vdots}$}\noLine
\UIC{$\forall B_{15}\,(\varphi_{79})$}
\AXC{$\stackrel{\Pi_{56}}{\vdots}$}\noLine
\UIC{$\forall B_{15}\,(\varphi_{79}) \to (\forall B_{15}\,(\varphi_{79}) \to \varphi_{15} \to \varphi_{16}) \to \varphi_{105}$}
\BIC{$(\forall B_{15}\,(\varphi_{79}) \to \varphi_{15} \to \varphi_{16}) \to \varphi_{105}$}
\BIC{$\varphi_{105}$}
\end{ptree}
\noindent$\Pi_{58}$:
\begin{ptree}
\AXC{$\mathrm{SubIdx}(s) \land \neg \forall k \in \omega\,(k \in \mathbb{N} \to k \leq s_{k}) \to \varphi_{15}$\lbl{A4}}
\AXC{$\stackrel{\Pi_{57}}{\vdots}$}\noLine
\UIC{$\varphi_{105}$}
\BIC{$\mathrm{SubIdx}(s) \to \forall k \in \omega\,(k \in \mathbb{N} \to k \leq s_{k})$}
\AXC{$(\mathrm{SubIdx}(s) \to \forall k \in \omega\,(k \in \mathbb{N} \to k \leq s_{k})) \to \varphi_{95}$\lbl{A1}}
\BIC{$\varphi_{95}$}
\end{ptree}
\noindent$\Pi_{59}$:
\begin{ptree}
\AXC{$\forall k \in \omega\,(k \in \mathbb{N} \to k \leq s_{k}) \to k \in \omega \to k \in \mathbb{N} \to k \leq s_{k}$\lbl{A2}}
\AXC{$\stackrel{\Pi_{58}}{\vdots}$}\noLine
\UIC{$\varphi_{95}$}
\BIC{$\mathrm{SubIdx}(s) \to k \in \omega \to k \in \mathbb{N} \to k \leq s_{k}$}
\AXC{$\varphi_{103}$\lbl{A1}}
\BIC{$(\mathrm{SubIdx}(s) \to k \in \mathbb{N} \to k \in \omega) \to \varphi_{35}$}
\end{ptree}
\noindent$\Pi_{60}$:
\begin{ptree}
\AXC{$\forall n\,(n \in \mathbb{N} \leftrightarrow n \in \omega \land n \neq \emptyset)$\lbl{Def$\mathbb{N}$}}
\AXC{$\varphi_{47}$\lbl{A2}}
\BIC{$k \in \mathbb{N} \leftrightarrow k \in \omega \land k \neq \emptyset$}
\AXC{$\stackrel{\Pi_{1}}{\vdots}$}\noLine
\UIC{$\varphi_{40}$}
\BIC{$\mathrm{SubIdx}(s) \to k \in \mathbb{N} \to k \in \omega$}
\AXC{$\stackrel{\Pi_{59}}{\vdots}$}\noLine
\UIC{$(\mathrm{SubIdx}(s) \to k \in \mathbb{N} \to k \in \omega) \to \varphi_{35}$}
\BIC{$\varphi_{35}$}
\end{ptree}
\noindent$\Pi_{61}$:
\begin{ptree}
\AXC{$\varphi_{46}$\lbl{A3}}
\AXC{$\mathrm{SubIdx}(s) \to k \in \mathbb{N} \to k \in \mathbb{N}$\lbl{A1}}
\AXC{$\stackrel{\Pi_{60}}{\vdots}$}\noLine
\UIC{$\varphi_{35}$}
\BIC{$\mathrm{SubIdx}(s) \to k \in \mathbb{N} \to k \leq s_{k}$}
\UIC{$\forall k\,(\mathrm{SubIdx}(s) \to k \in \mathbb{N} \to k \leq s_{k})$}
\AXC{$\forall k\,(\mathrm{SubIdx}(s) \to k \in \mathbb{N} \to k \leq s_{k}) \to \varphi_{46} \to \varphi_{29}$\lbl{A1}}
\BIC{$\varphi_{46} \to \varphi_{29}$}
\BIC{$\varphi_{29}$}
\end{ptree}
\noindent 主木:
\begin{ptree}
\AXC{$\mathrm{SubIdx}(s) \to \forall k\,(\mathrm{SubIdx}(s))$\lbl{A4}}
\AXC{$\stackrel{\Pi_{61}}{\vdots}$}\noLine
\UIC{$\varphi_{29}$}
\BIC{$\mathrm{SubIdx}(s) \to \forall k \in \mathbb{N}\,(k \leq s_{k})$}
\UIC{$\forall s\,(\mathrm{SubIdx}(s) \to \forall k \in \mathbb{N}\,(k \leq s_{k}))$}
\end{ptree}
\begin{abbrevlist}
\item $\varphi_{1} :\equiv \forall x\,\forall y\,(x = y \to y = x) \to \forall y\,(S(k) = y \to y = S(k))$
\item $\varphi_{2} :\equiv \forall x\,\forall y\,(x = y \to y = x) \to \forall y\,(1_{\omega} = y \to y = 1_{\omega})$
\item $\varphi_{3} :\equiv \forall y\,(S(k) = y \to y = S(k)) \to S(k) = 1_{\omega} \to 1_{\omega} = S(k)$
\item $\varphi_{4} :\equiv 1_{\omega} = S(k) \to 1_{\omega} = s_{1_{\omega}} \to S(k) = s_{S(k)}$
\item $\varphi_{5} :\equiv S(k) = S(\emptyset) \to S(\emptyset) = 1_{\omega} \to S(k) = 1_{\omega}$
\item $\varphi_{6} :\equiv 1_{\omega} = S(k) \to 1_{\omega} < s_{1_{\omega}} \to S(k) < s_{S(k)}$
\item $\varphi_{7} :\equiv \forall k \in B_{15}\,(S(k) \in B_{15}) \to y \in B_{15} \to S(y) \in B_{15}$
\item $\varphi_{8} :\equiv \forall y \in B_{15}\,(S(y) \in B_{15}) \to k \in B_{15} \to S(k) \in B_{15}$
\item $\varphi_{9} :\equiv \forall x\,\forall y\,\forall z\,(x = y \to y = z \to x = z) \to \forall y\,\forall z\,(S(k) = y \to y = z \to S(k) = z)$
\item $\varphi_{10} :\equiv (1_{\omega} = S(\emptyset) \to S(\emptyset) = 1_{\omega}) \to S(\emptyset) = 1_{\omega}$
\item $\varphi_{11} :\equiv \forall y\,(1_{\omega} = y \to y = 1_{\omega}) \to 1_{\omega} = S(\emptyset) \to S(\emptyset) = 1_{\omega}$
\item $\varphi_{12} :\equiv \forall k \in B_{15}\,(S(k) \in B_{15}) \to \forall y \in B_{15}\,(S(y) \in B_{15})$
\item $\varphi_{13} :\equiv \forall y \in B_{15}\,(S(y) \in B_{15}) \to \forall k \in B_{15}\,(S(k) \in B_{15})$
\item $\varphi_{14} :\equiv \mathrm{Ind}(B_{15}) \leftrightarrow \emptyset \in B_{15} \land \forall y \in B_{15}\,(S(y) \in B_{15})$
\item $\varphi_{15} :\equiv \forall B_{15}\,(\mathrm{SubIdx}(s) \land \neg \forall k \in \omega\,(k \in \mathbb{N} \to k \leq s_{k}))$
\item $\varphi_{16} :\equiv \forall B_{15}\,(\neg \forall k\,(k \in B_{15} \leftrightarrow k \in \omega \land (k \in \mathbb{N} \to k \leq s_{k})))$
\item $\varphi_{17} :\equiv \forall A_{16}\,\exists B_{15}\,\forall k\,(k \in B_{15} \leftrightarrow k \in A_{16} \land (k \in \mathbb{N} \to k \leq s_{k}))$
\item $\varphi_{18} :\equiv \forall k \in B_{15}\,(S(k) \in B_{15}) \to \forall y\,\forall k \in B_{15}\,(S(k) \in B_{15})$
\item $\varphi_{19} :\equiv \forall y \in B_{15}\,(S(y) \in B_{15}) \to \forall k\,\forall y \in B_{15}\,(S(y) \in B_{15})$
\item $\varphi_{20} :\equiv \mathrm{SubIdx}(s) \leftrightarrow \mathrm{NSeq}(s) \land \forall k \in \mathbb{N}\,(s_{k} < s_{S(k)})$
\item $\varphi_{21} :\equiv \forall x\,\forall y\,\forall z\,(x \leq y \to y < z \to x < z) \to \forall y\,\forall z\,(k \leq y \to y < z \to k < z)$
\item $\varphi_{22} :\equiv \forall y\,\forall k \in B_{15}\,(S(k) \in B_{15}) \to \forall y \in B_{15}\,(S(y) \in B_{15})$
\item $\varphi_{23} :\equiv \forall k\,\forall y \in B_{15}\,(S(y) \in B_{15}) \to \forall k \in B_{15}\,(S(k) \in B_{15})$
\item $\varphi_{24} :\equiv \forall n \in \mathbb{N}\,(1_{\omega} \leq n) \to s_{1_{\omega}} \in \mathbb{N} \to 1_{\omega} \leq s_{1_{\omega}}$
\item $\varphi_{25} :\equiv (\mathrm{SubIdx}(s) \to s_{1_{\omega}} \in \mathbb{N}) \to \mathrm{SubIdx}(s) \to 1_{\omega} \leq s_{1_{\omega}}$
\item $\varphi_{26} :\equiv \forall k\,(k \in B_{15} \leftrightarrow k \in \omega \land (k \in \mathbb{N} \to k \leq s_{k})) \to \mathrm{Ind}(B_{15})$
\item $\varphi_{27} :\equiv \mathrm{SubIdx}(s) \land \forall k\,(k \in B_{15} \leftrightarrow k \in \omega \land (k \in \mathbb{N} \to k \leq s_{k}))$
\item $\varphi_{28} :\equiv S(k) \in B_{15} \leftrightarrow S(k) \in \omega \land (S(k) \in \mathbb{N} \to S(k) \leq s_{S(k)})$
\item $\varphi_{29} :\equiv (\mathrm{SubIdx}(s) \to \forall k\,(\mathrm{SubIdx}(s))) \to \mathrm{SubIdx}(s) \to \forall k \in \mathbb{N}\,(k \leq s_{k})$
\item $\varphi_{30} :\equiv \forall z\,\forall w\,(z = w \to S(k) = S(z) \to S(k) = S(w)) \to \forall w\,(k = w \to S(k) = S(k) \to S(k) = S(w))$
\item $\varphi_{31} :\equiv \forall z_{17}\,\forall A_{16}\,\exists B_{15}\,\forall k\,(k \in B_{15} \leftrightarrow k \in A_{16} \land (k \in \mathbb{N} \to k \leq z_{17}(k)))$
\item $\varphi_{32} :\equiv \forall y\,\forall z\,(k \leq y \to y < z \to k < z) \to \forall z\,(k \leq s_{k} \to s_{k} < z \to k < z)$
\item $\varphi_{33} :\equiv (\mathrm{SubIdx}(s) \to \mathrm{NSeq}(s)) \to 1_{\omega} \in \mathbb{N} \to \mathrm{SubIdx}(s) \to s_{1_{\omega}} \in \mathbb{N}$
\item $\varphi_{34} :\equiv \mathrm{SubIdx}(s) \to \exists B_{15}\,\forall k\,(k \in B_{15} \leftrightarrow k \in \omega \land (k \in \mathbb{N} \to k \leq s_{k}))$
\item $\varphi_{35} :\equiv (\mathrm{SubIdx}(s) \to k \in \mathbb{N} \to k \in \mathbb{N}) \to \mathrm{SubIdx}(s) \to k \in \mathbb{N} \to k \leq s_{k}$
\item $\varphi_{36} :\equiv \forall k\,(k \in B_{15} \leftrightarrow k \in \omega \land (k \in \mathbb{N} \to k \leq s_{k})) \to \emptyset \in B_{15}$
\item $\varphi_{37} :\equiv \forall n\,(n \in \omega \leftrightarrow \forall x\,(\mathrm{Ind}(x) \to n \in x)) \to (k \in \omega \leftrightarrow \forall x\,(\mathrm{Ind}(x) \to k \in x))$
\item $\varphi_{38} :\equiv \forall y\,\forall z\,(S(k) = y \to y = z \to S(k) = z) \to \forall z\,(S(k) = S(\emptyset) \to S(\emptyset) = z \to S(k) = z)$
\item $\varphi_{39} :\equiv \forall w\,(k = w \to S(k) = S(k) \to S(k) = S(w)) \to k = \emptyset \to S(k) = S(k) \to S(k) = S(\emptyset)$
\item $\varphi_{40} :\equiv (k \in \mathbb{N} \leftrightarrow k \in \omega \land k \neq \emptyset) \to \mathrm{SubIdx}(s) \to k \in \mathbb{N} \to k \in \omega$
\item $\varphi_{41} :\equiv \forall k\,(k \in B_{15} \leftrightarrow k \in \omega \land (k \in \mathbb{N} \to k \leq s_{k})) \to k \in B_{15} \to k \in B_{15}$
\item $\varphi_{42} :\equiv \forall k\,(k \in B_{15} \leftrightarrow k \in \omega \land (k \in \mathbb{N} \to k \leq s_{k})) \to k \in B_{15} \to k \in \omega$
\item $\varphi_{43} :\equiv \forall k\,(k \in B_{15} \leftrightarrow k \in \omega \land (k \in \mathbb{N} \to k \leq s_{k})) \to k \in B_{15} \to S(k) \in B_{15}$
\item $\varphi_{44} :\equiv \forall z\,\forall w\,(z = w \to z = s_{z} \to w = s_{w}) \to \forall w\,(1_{\omega} = w \to 1_{\omega} = s_{1_{\omega}} \to w = s_{w})$
\item $\varphi_{45} :\equiv \forall k\,(k \in B_{15} \leftrightarrow k \in \omega \land (k \in \mathbb{N} \to k \leq s_{k})) \to k \in \omega \to k \in \omega$
\item $\varphi_{46} :\equiv \forall k\,(\mathrm{SubIdx}(s) \to k \in \mathbb{N} \to k \leq s_{k}) \to \forall k\,(\mathrm{SubIdx}(s)) \to \forall k \in \mathbb{N}\,(k \leq s_{k})$
\item $\varphi_{47} :\equiv \forall n\,(n \in \mathbb{N} \leftrightarrow n \in \omega \land n \neq \emptyset) \to (k \in \mathbb{N} \leftrightarrow k \in \omega \land k \neq \emptyset)$
\item $\varphi_{48} :\equiv \forall k\,(k \in B_{15} \leftrightarrow k \in \omega \land (k \in \mathbb{N} \to k \leq s_{k})) \to k \in B_{15} \to S(k) \in \omega$
\item $\varphi_{49} :\equiv \forall k\,(k \in B_{15} \leftrightarrow k \in \omega \land (k \in \mathbb{N} \to k \leq s_{k})) \to \forall k \in B_{15}\,(S(k) \in B_{15})$
\item $\varphi_{50} :\equiv \forall k\,(k \in B_{15} \leftrightarrow k \in \omega \land (k \in \mathbb{N} \to k \leq s_{k})) \to k \in \omega \to \forall x\,(\mathrm{Ind}(x) \to k \in x)$
\item $\varphi_{51} :\equiv \forall s\,\forall k\,(\mathrm{NSeq}(s) \to k \in \mathbb{N} \to s_{k} \in \mathbb{N}) \to \forall k\,(\mathrm{NSeq}(s) \to k \in \mathbb{N} \to s_{k} \in \mathbb{N})$
\item $\varphi_{52} :\equiv \forall k\,(\mathrm{NSeq}(s) \to k \in \mathbb{N} \to s_{k} \in \mathbb{N}) \to \mathrm{NSeq}(s) \to S(k) \in \mathbb{N} \to s_{S(k)} \in \mathbb{N}$
\item $\varphi_{53} :\equiv \forall k\,(\mathrm{NSeq}(s) \to k \in \mathbb{N} \to s_{k} \in \mathbb{N}) \to \mathrm{NSeq}(s) \to 1_{\omega} \in \mathbb{N} \to s_{1_{\omega}} \in \mathbb{N}$
\item $\varphi_{54} :\equiv \forall k\,(k \in B_{15} \leftrightarrow k \in \omega \land (k \in \mathbb{N} \to k \leq s_{k})) \to k \in B_{15} \to k \in \mathbb{N} \to k \leq s_{k}$
\item $\varphi_{55} :\equiv \forall k\,(k \in B_{15} \leftrightarrow k \in \omega \land (k \in \mathbb{N} \to k \leq s_{k})) \to k \in \omega \to \mathrm{Ind}(B_{15}) \to k \in B_{15}$
\item $\varphi_{56} :\equiv \forall k\,(k \in B_{15} \leftrightarrow k \in \omega \land (k \in \mathbb{N} \to k \leq s_{k})) \to k \in \omega \to k \in \mathbb{N} \to k \leq s_{k}$
\item $\varphi_{57} :\equiv \forall w\,(1_{\omega} = w \to 1_{\omega} = s_{1_{\omega}} \to w = s_{w}) \to \varphi_{4}$
\item $\varphi_{58} :\equiv \forall z\,\forall w\,(z = w \to z < s_{z} \to w < s_{w}) \to \forall w\,(1_{\omega} = w \to 1_{\omega} < s_{1_{\omega}} \to w < s_{w})$
\item $\varphi_{59} :\equiv \forall k\,(k \in B_{15} \leftrightarrow k \in \omega \land (k \in \mathbb{N} \to k \leq s_{k})) \to \forall k \in \omega\,(k \in \mathbb{N} \to k \leq s_{k})$
\item $\varphi_{60} :\equiv \forall k\,(k \in B_{15} \leftrightarrow k \in \omega \land (k \in \mathbb{N} \to k \leq s_{k})) \to k \in B_{15} \to k \in \omega \to S(k) \in \omega$
\item $\varphi_{61} :\equiv \forall z\,(k \leq s_{k} \to s_{k} < z \to k < z) \to k \leq s_{k} \to s_{k} < s_{S(k)} \to k < s_{S(k)}$
\item $\varphi_{62} :\equiv \forall k\,(k \in B_{15} \leftrightarrow k \in \omega \land (k \in \mathbb{N} \to k \leq s_{k})) \to k \in B_{15} \to S(k) \in \mathbb{N} \to S(k) \leq s_{S(k)}$
\item $\varphi_{63} :\equiv \forall k\,(k \in B_{15} \leftrightarrow k \in \omega \land (k \in \mathbb{N} \to k \leq s_{k})) \to k \in B_{15} \to \forall y \in \omega\,(S(y) \in \omega)$
\item $\varphi_{64} :\equiv \forall m\,\forall n\,(m \in \omega \to n \in \omega \to m < n \to S(m) \leq n) \to \forall n\,(k \in \omega \to n \in \omega \to k < n \to S(k) \leq n)$
\item $\varphi_{65} :\equiv \forall n\,(n \in \mathbb{N} \leftrightarrow n \in \omega \land n \neq \emptyset) \to (\emptyset \in \mathbb{N} \leftrightarrow \emptyset \in \omega \land \emptyset \neq \emptyset)$
\item $\varphi_{66} :\equiv \forall k\,(k \in B_{15} \leftrightarrow k \in \omega \land (k \in \mathbb{N} \to k \leq s_{k})) \to (k \in B_{15} \leftrightarrow k \in \omega \land (k \in \mathbb{N} \to k \leq s_{k}))$
\item $\varphi_{67} :\equiv \forall k\,(k \in B_{15} \leftrightarrow k \in \omega \land (k \in \mathbb{N} \to k \leq s_{k})) \to k \in B_{15} \to k \in \omega \land (k \in \mathbb{N} \to k \leq s_{k})$
\item $\varphi_{68} :\equiv \forall w\,(1_{\omega} = w \to 1_{\omega} < s_{1_{\omega}} \to w < s_{w}) \to \varphi_{6}$
\item $\varphi_{69} :\equiv \forall k\,(k \in B_{15} \leftrightarrow k \in \omega \land (k \in \mathbb{N} \to k \leq s_{k})) \to k \in B_{15} \to S(k) \in \mathbb{N} \to k = \emptyset \to S(k) = 1_{\omega}$
\item $\varphi_{70} :\equiv \forall k\,(k \in B_{15} \leftrightarrow k \in \omega \land (k \in \mathbb{N} \to k \leq s_{k})) \to k \in B_{15} \to S(k) \in \mathbb{N} \to k = \emptyset \to 1_{\omega} = S(k)$
\item $\varphi_{71} :\equiv \forall k\,(k \in B_{15} \leftrightarrow k \in \omega \land (k \in \mathbb{N} \to k \leq s_{k})) \to k \in B_{15} \to S(k) \in \mathbb{N} \to k \neq \emptyset \to k \in \mathbb{N}$
\item $\varphi_{72} :\equiv \forall k\,(k \in B_{15} \leftrightarrow k \in \omega \land (k \in \mathbb{N} \to k \leq s_{k})) \to k \in B_{15} \to S(k) \in \mathbb{N} \to k = \emptyset \to k = \emptyset$
\item $\varphi_{73} :\equiv \forall k\,(k \in B_{15} \leftrightarrow k \in \omega \land (k \in \mathbb{N} \to k \leq s_{k})) \to \forall k\,(k \in B_{15} \leftrightarrow k \in \omega \land (k \in \mathbb{N} \to k \leq s_{k}))$
\item $\varphi_{74} :\equiv \forall k\,(k \in B_{15} \leftrightarrow k \in \omega \land (k \in \mathbb{N} \to k \leq s_{k})) \to k \in B_{15} \to S(k) \in \mathbb{N} \to k \neq \emptyset \to k \leq s_{k}$
\item $\varphi_{75} :\equiv (\mathrm{SubIdx}(s) \to \mathrm{SubIdx}(s)) \to \varphi_{20} \to \mathrm{SubIdx}(s) \to \mathrm{NSeq}(s)$
\item $\varphi_{76} :\equiv \forall k\,(k \in B_{15} \leftrightarrow k \in \omega \land (k \in \mathbb{N} \to k \leq s_{k})) \to k \in B_{15} \to S(k) \in \mathbb{N} \to k \neq \emptyset \to k \neq \emptyset$
\item $\varphi_{77} :\equiv \forall k\,(k \in B_{15} \leftrightarrow k \in \omega \land (k \in \mathbb{N} \to k \leq s_{k})) \to k \in B_{15} \to S(k) \in \mathbb{N} \to k \neq \emptyset \to k \in \omega$
\item $\varphi_{78} :\equiv \forall k\,(k \in B_{15} \leftrightarrow k \in \omega \land (k \in \mathbb{N} \to k \leq s_{k})) \to \emptyset \in \omega \land \forall y \in \omega\,(S(y) \in \omega)$
\item $\varphi_{79} :\equiv \mathrm{SubIdx}(s) \land \neg \forall k \in \omega\,(k \in \mathbb{N} \to k \leq s_{k}) \to \neg \forall k\,(k \in B_{15} \leftrightarrow k \in \omega \land (k \in \mathbb{N} \to k \leq s_{k}))$
\item $\varphi_{80} :\equiv \forall k\,(k \in B_{15} \leftrightarrow k \in \omega \land (k \in \mathbb{N} \to k \leq s_{k})) \to k \in B_{15} \to S(k) \in \mathbb{N} \to k \neq \emptyset \to S(k) \in \mathbb{N}$
\item $\varphi_{81} :\equiv \forall k\,(k \in B_{15} \leftrightarrow k \in \omega \land (k \in \mathbb{N} \to k \leq s_{k})) \to k \in B_{15} \to S(k) \in \mathbb{N} \to k = \emptyset \to S(k) = S(\emptyset)$
\item $\varphi_{82} :\equiv \forall k\,(k \in B_{15} \leftrightarrow k \in \omega \land (k \in \mathbb{N} \to k \leq s_{k})) \to k \in B_{15} \to S(k) \in \mathbb{N} \to k = \emptyset \to S(k) \leq s_{S(k)}$
\item $\varphi_{83} :\equiv \forall k\,(k \in B_{15} \leftrightarrow k \in \omega \land (k \in \mathbb{N} \to k \leq s_{k})) \to k \in B_{15} \to S(k) \in \mathbb{N} \to k \neq \emptyset \to k < s_{S(k)}$
\item $\varphi_{84} :\equiv \forall k\,(k \in B_{15} \leftrightarrow k \in \omega \land (k \in \mathbb{N} \to k \leq s_{k})) \to k \in B_{15} \to S(k) \in \mathbb{N} \to k \neq \emptyset \to S(k) \leq s_{S(k)}$
\item $\varphi_{85} :\equiv \forall k\,(k \in B_{15} \leftrightarrow k \in \omega \land (k \in \mathbb{N} \to k \leq s_{k})) \to k \in B_{15} \to S(k) \in \mathbb{N} \to k \neq \emptyset \to s_{S(k)} \in \mathbb{N}$
\item $\varphi_{86} :\equiv (\mathrm{NSeq}(s) \to 1_{\omega} \in \mathbb{N} \to s_{1_{\omega}} \in \mathbb{N}) \to \varphi_{33}$
\item $\varphi_{87} :\equiv \forall k\,(k \in B_{15} \leftrightarrow k \in \omega \land (k \in \mathbb{N} \to k \leq s_{k})) \to k \in B_{15} \to S(k) \in \mathbb{N} \to k \neq \emptyset \to s_{S(k)} \in \omega$
\item $\varphi_{88} :\equiv \varphi_{17} \to \exists B_{15}\,\forall k\,(k \in B_{15} \leftrightarrow k \in \omega \land (k \in \mathbb{N} \to k \leq s_{k}))$
\item $\varphi_{89} :\equiv \forall k\,(k \in B_{15} \leftrightarrow k \in \omega \land (k \in \mathbb{N} \to k \leq s_{k})) \to k \in B_{15} \to S(k) \in \mathbb{N} \to k \neq \emptyset \to s_{k} < s_{S(k)}$
\item $\varphi_{90} :\equiv \forall x\,(\mathrm{Ind}(x) \leftrightarrow \emptyset \in x \land \forall y \in x\,(S(y) \in x)) \to (\mathrm{Ind}(\omega) \leftrightarrow \emptyset \in \omega \land \forall y \in \omega\,(S(y) \in \omega))$
\item $\varphi_{91} :\equiv \forall n\,(k \in \omega \to n \in \omega \to k < n \to S(k) \leq n) \to k \in \omega \to s_{S(k)} \in \omega \to k < s_{S(k)} \to S(k) \leq s_{S(k)}$
\item $\varphi_{92} :\equiv (\forall k \in \mathbb{N}\,(s_{k} < s_{S(k)}) \to k \in \mathbb{N} \to s_{k} < s_{S(k)}) \to \mathrm{SubIdx}(s) \to k \in \mathbb{N} \to s_{k} < s_{S(k)}$
\item $\varphi_{93} :\equiv \exists B_{15}\,\forall k\,(k \in B_{15} \leftrightarrow k \in \omega \land (k \in \mathbb{N} \to k \leq s_{k})) \to \varphi_{34}$
\item $\varphi_{94} :\equiv \mathrm{SubIdx}(s) \to \forall k\,(k \in B_{15} \leftrightarrow k \in \omega \land (k \in \mathbb{N} \to k \leq s_{k})) \to \varphi_{28}$
\item $\varphi_{95} :\equiv (\forall k \in \omega\,(k \in \mathbb{N} \to k \leq s_{k}) \to k \in \omega \to k \in \mathbb{N} \to k \leq s_{k}) \to \mathrm{SubIdx}(s) \to k \in \omega \to k \in \mathbb{N} \to k \leq s_{k}$
\item $\varphi_{96} :\equiv (k \in \omega \leftrightarrow \forall x\,(\mathrm{Ind}(x) \to k \in x)) \to \mathrm{SubIdx}(s) \to \varphi_{50}$
\item $\varphi_{97} :\equiv (\emptyset \in \mathbb{N} \leftrightarrow \emptyset \in \omega \land \emptyset \neq \emptyset) \to \emptyset = \emptyset \to \mathrm{SubIdx}(s) \to \emptyset \in \mathbb{N} \to \emptyset \leq s_{\emptyset}$
\item $\varphi_{98} :\equiv (\mathrm{SubIdx}(s) \to \mathrm{SubIdx}(s)) \to \varphi_{20} \to \mathrm{SubIdx}(s) \to \forall k \in \mathbb{N}\,(s_{k} < s_{S(k)})$
\item $\varphi_{99} :\equiv \forall k\,(k \in B_{15} \leftrightarrow k \in \omega \land (k \in \mathbb{N} \to k \leq s_{k})) \to (\emptyset \in B_{15} \leftrightarrow \emptyset \in \omega \land (\emptyset \in \mathbb{N} \to \emptyset \leq s_{\emptyset}))$
\item $\varphi_{100} :\equiv (S(k) = 1_{\omega} \to 1_{\omega} = S(k)) \to \mathrm{SubIdx}(s) \to \varphi_{70}$
\item $\varphi_{101} :\equiv (\mathrm{SubIdx}(s) \to \emptyset \in \mathbb{N} \to \emptyset \leq s_{\emptyset}) \to \mathrm{SubIdx}(s) \to \varphi_{36}$
\item $\varphi_{102} :\equiv (\forall x\,(\mathrm{Ind}(x) \to k \in x) \to \mathrm{Ind}(B_{15}) \to k \in B_{15}) \to \mathrm{SubIdx}(s) \to \varphi_{55}$
\item $\varphi_{103} :\equiv (\mathrm{SubIdx}(s) \to k \in \omega \to k \in \mathbb{N} \to k \leq s_{k}) \to (\mathrm{SubIdx}(s) \to k \in \mathbb{N} \to k \in \omega) \to \varphi_{35}$
\item $\varphi_{104} :\equiv (k = \emptyset \to S(k) = S(k) \to S(k) = S(\emptyset)) \to \mathrm{SubIdx}(s) \to \varphi_{81}$
\item $\varphi_{105} :\equiv (\mathrm{SubIdx}(s) \land \neg \forall k \in \omega\,(k \in \mathbb{N} \to k \leq s_{k}) \to \varphi_{15}) \to \mathrm{SubIdx}(s) \to \forall k \in \omega\,(k \in \mathbb{N} \to k \leq s_{k})$
\item $\varphi_{106} :\equiv (\forall y \in \omega\,(S(y) \in \omega) \to k \in \omega \to S(k) \in \omega) \to \mathrm{SubIdx}(s) \to \varphi_{60}$
\item $\varphi_{107} :\equiv (\mathrm{Ind}(\omega) \leftrightarrow \emptyset \in \omega \land \forall y \in \omega\,(S(y) \in \omega)) \to \mathrm{SubIdx}(s) \to \varphi_{78}$
\item $\varphi_{108} :\equiv \forall k\,(\varphi_{27} \to k \in B_{15} \to S(k) \in B_{15}) \to \forall k\,(\varphi_{27}) \to \forall k \in B_{15}\,(S(k) \in B_{15})$
\item $\varphi_{109} :\equiv \forall k\,(\varphi_{27} \to k \in \omega \to k \in \mathbb{N} \to k \leq s_{k}) \to \forall k\,(\varphi_{27}) \to \forall k \in \omega\,(k \in \mathbb{N} \to k \leq s_{k})$
\item $\varphi_{110} :\equiv (\forall k\,(k \in B_{15} \leftrightarrow k \in \omega \land (k \in \mathbb{N} \to k \leq s_{k})) \to \varphi_{28}) \to \varphi_{94}$
\item $\varphi_{111} :\equiv (\mathrm{SubIdx}(s) \to 1_{\omega} \leq s_{1_{\omega}}) \to \varphi_{6} \to \varphi_{4} \to \mathrm{SubIdx}(s) \to \varphi_{82}$
\item $\varphi_{112} :\equiv (k \in \mathbb{N} \to k \in \omega) \to (\mathrm{SubIdx}(s) \to \varphi_{71}) \to \mathrm{SubIdx}(s) \to \varphi_{77}$
\item $\varphi_{113} :\equiv (k \in \mathbb{N} \to S(k) \in \mathbb{N}) \to (\mathrm{SubIdx}(s) \to \varphi_{71}) \to \mathrm{SubIdx}(s) \to \varphi_{80}$
\item $\varphi_{114} :\equiv (\mathrm{SubIdx}(s) \to \mathrm{NSeq}(s)) \to (\mathrm{SubIdx}(s) \to \varphi_{80}) \to \mathrm{SubIdx}(s) \to \varphi_{85}$
\item $\varphi_{115} :\equiv (\mathrm{SubIdx}(s) \to \varphi_{60}) \to (\mathrm{SubIdx}(s) \to \varphi_{42}) \to \mathrm{SubIdx}(s) \to \varphi_{48}$
\item $\varphi_{116} :\equiv (\mathrm{SubIdx}(s) \to \varphi_{26}) \to (\mathrm{SubIdx}(s) \to \varphi_{66}) \to \mathrm{SubIdx}(s) \to \varphi_{56}$
\item $\varphi_{117} :\equiv (\mathrm{SubIdx}(s) \to \varphi_{48}) \to (\mathrm{SubIdx}(s) \to \varphi_{62}) \to \mathrm{SubIdx}(s) \to \varphi_{43}$
\item $\varphi_{118} :\equiv (\mathrm{SubIdx}(s) \to k \in \mathbb{N} \to s_{k} < s_{S(k)}) \to (\mathrm{SubIdx}(s) \to \varphi_{71}) \to \mathrm{SubIdx}(s) \to \varphi_{89}$
\item $\varphi_{119} :\equiv (s_{S(k)} \in \mathbb{N} \to s_{S(k)} \in \omega) \to (\mathrm{SubIdx}(s) \to \varphi_{85}) \to \mathrm{SubIdx}(s) \to \varphi_{87}$
\item $\varphi_{120} :\equiv (\mathrm{SubIdx}(s) \to \varphi_{81}) \to S(\emptyset) = 1_{\omega} \to \varphi_{5} \to \mathrm{SubIdx}(s) \to \varphi_{69}$
\item $\varphi_{121} :\equiv (\mathrm{SubIdx}(s) \to \varphi_{66}) \to (\mathrm{SubIdx}(s) \to \varphi_{41}) \to \mathrm{SubIdx}(s) \to \varphi_{67}$
\item $\varphi_{122} :\equiv (\mathrm{SubIdx}(s) \to \varphi_{42}) \to (\mathrm{SubIdx}(s) \to \varphi_{76}) \to \mathrm{SubIdx}(s) \to \varphi_{71}$
\item $\varphi_{123} :\equiv (\mathrm{SubIdx}(s) \to \varphi_{54}) \to (\mathrm{SubIdx}(s) \to \varphi_{71}) \to \mathrm{SubIdx}(s) \to \varphi_{74}$
\item $\varphi_{124} :\equiv (\mathrm{SubIdx}(s) \to \varphi_{82}) \to (\mathrm{SubIdx}(s) \to \varphi_{84}) \to \mathrm{SubIdx}(s) \to \varphi_{62}$
\item $\varphi_{125} :\equiv (\mathrm{SubIdx}(s) \to \varphi_{36}) \to (\mathrm{SubIdx}(s) \to \varphi_{49}) \to \varphi_{13} \to \varphi_{12} \to \mathrm{SubIdx}(s) \to \varphi_{26}$
\item $\varphi_{126} :\equiv (\mathrm{SubIdx}(s) \to \varphi_{74}) \to (\mathrm{SubIdx}(s) \to \varphi_{89}) \to \mathrm{SubIdx}(s) \to \varphi_{83}$
\item $\varphi_{127} :\equiv (\mathrm{SubIdx}(s) \to \varphi_{87}) \to (\mathrm{SubIdx}(s) \to \varphi_{83}) \to \mathrm{SubIdx}(s) \to \varphi_{84}$
\item $\varphi_{128} :\equiv \forall k\,(\varphi_{27} \to k \in B_{15} \to S(k) \in B_{15}) \to \varphi_{108} \to (\varphi_{27} \to \forall k\,(\varphi_{27})) \to \mathrm{SubIdx}(s) \to \varphi_{49}$
\item $\varphi_{129} :\equiv \forall k\,(\varphi_{27} \to k \in \omega \to k \in \mathbb{N} \to k \leq s_{k}) \to \varphi_{109} \to (\varphi_{27} \to \forall k\,(\varphi_{27})) \to \mathrm{SubIdx}(s) \to \varphi_{59}$
\item $\varphi_{130} :\equiv (k \in \omega \to s_{S(k)} \in \omega \to k < s_{S(k)} \to S(k) \leq s_{S(k)}) \to (\mathrm{SubIdx}(s) \to \varphi_{77}) \to \varphi_{127}$
\end{abbrevlist}
\end{fullproof}

% 自動生成 (prover/build.js): SubConvIdx  ノード数 179
\begin{fullproof}{[C18]}{$\forall a\,\forall s\,\forall \alpha\,(\mathrm{Conv}(a,\alpha) \to \mathrm{SubIdx}(s) \to \mathrm{Conv}(a \circ s,\alpha))$}{179}\phantomsection\label{tree:SubConvIdx}
\noindent$\Pi_{1}$:
\begin{ptree}
\AXC{$\forall s\,(\mathrm{SubIdx}(s) \to \forall k \in \mathbb{N}\,(k \leq s_{k}))$\lbl{[C17]}}
\AXC{$\varphi_{19}$\lbl{A2}}
\BIC{$\mathrm{SubIdx}(s) \to \forall k \in \mathbb{N}\,(k \leq s_{k})$}
\AXC{$(\mathrm{SubIdx}(s) \to \forall k \in \mathbb{N}\,(k \leq s_{k})) \to \varphi_{25}$\lbl{A1}}
\BIC{$\varphi_{25}$}
\end{ptree}
\noindent$\Pi_{2}$:
\begin{ptree}
\AXC{$\mathrm{Conv}(a,\alpha) \to \mathrm{SubIdx}(s) \to \mathrm{SubIdx}(s)$\lbl{A1}}
\AXC{$\stackrel{\Pi_{1}}{\vdots}$}\noLine
\UIC{$\varphi_{25}$}
\BIC{$\varphi_{5}$}
\AXC{$\varphi_{5} \to \varphi_{32}$\lbl{A1}}
\BIC{$\varphi_{32}$}
\end{ptree}
\noindent$\Pi_{3}$:
\begin{ptree}
\AXC{$\mathrm{Conv}(a,\alpha) \to \mathrm{SubIdx}(s) \to \varepsilon \in \mathbb{R} \to 0 < \varepsilon \to \varphi_{23}$\lbl{A1}}
\AXC{$\forall k \in \mathbb{N}\,(k \leq s_{k}) \to n \in \mathbb{N} \to n \leq s_{n}$\lbl{A2}}
\AXC{$\stackrel{\Pi_{2}}{\vdots}$}\noLine
\UIC{$\varphi_{32}$}
\BIC{$\mathrm{Conv}(a,\alpha) \to \mathrm{SubIdx}(s) \to n \in \mathbb{N} \to n \leq s_{n}$}
\AXC{$\varphi_{79}$\lbl{A1}}
\BIC{$\varphi_{71}$}
\BIC{$\mathrm{Conv}(a,\alpha) \to \mathrm{SubIdx}(s) \to \varepsilon \in \mathbb{R} \to 0 < \varepsilon \to \varphi_{28}$}
\end{ptree}
\noindent$\Pi_{4}$:
\begin{ptree}
\AXC{$\forall x\,\forall y\,\forall z\,(x \leq y \to y \leq z \to x \leq z)$\lbl{[O5]}}
\AXC{$\varphi_{4}$\lbl{A2}}
\BIC{$\forall y\,\forall z\,(N \leq y \to y \leq z \to N \leq z)$}
\AXC{$\varphi_{1}$\lbl{A2}}
\BIC{$\forall z\,(N \leq n \to n \leq z \to N \leq z)$}
\AXC{$\forall z\,(N \leq n \to n \leq z \to N \leq z) \to N \leq n \to n \leq s_{n} \to N \leq s_{n}$\lbl{A2}}
\BIC{$N \leq n \to n \leq s_{n} \to N \leq s_{n}$}
\end{ptree}
\noindent$\Pi_{5}$:
\begin{ptree}
\AXC{$\mathrm{Conv}(a,\alpha) \to \mathrm{SubIdx}(s) \to \varepsilon \in \mathbb{R} \to 0 < \varepsilon \to \varphi_{24}$\lbl{A1}}
\AXC{$\stackrel{\Pi_{4}}{\vdots}$}\noLine
\UIC{$N \leq n \to n \leq s_{n} \to N \leq s_{n}$}
\AXC{$(N \leq n \to n \leq s_{n} \to N \leq s_{n}) \to \varphi_{90}$\lbl{A1}}
\BIC{$\varphi_{90}$}
\BIC{$\varphi_{73}$}
\end{ptree}
\noindent$\Pi_{6}$:
\begin{ptree}
\AXC{$\stackrel{\Pi_{3}}{\vdots}$}\noLine
\UIC{$\mathrm{Conv}(a,\alpha) \to \mathrm{SubIdx}(s) \to \varepsilon \in \mathbb{R} \to 0 < \varepsilon \to \varphi_{28}$}
\AXC{$\stackrel{\Pi_{5}}{\vdots}$}\noLine
\UIC{$\varphi_{73}$}
\BIC{$\mathrm{Conv}(a,\alpha) \to \mathrm{SubIdx}(s) \to \varepsilon \in \mathbb{R} \to 0 < \varepsilon \to \varphi_{29}$}
\end{ptree}
\noindent$\Pi_{7}$:
\begin{ptree}
\AXC{$\forall s\,(\varphi_{11})$\lbl{DefSubIdx}}
\AXC{$\forall s\,(\varphi_{11}) \to \varphi_{11}$\lbl{A2}}
\BIC{$\varphi_{11}$}
\AXC{$\mathrm{Conv}(a,\alpha) \to \mathrm{SubIdx}(s) \to \mathrm{SubIdx}(s)$\lbl{A1}}
\AXC{$\varphi_{51}$\lbl{A1}}
\BIC{$\varphi_{11} \to \mathrm{Conv}(a,\alpha) \to \mathrm{SubIdx}(s) \to \mathrm{NSeq}(s)$}
\BIC{$\mathrm{Conv}(a,\alpha) \to \mathrm{SubIdx}(s) \to \mathrm{NSeq}(s)$}
\end{ptree}
\noindent$\Pi_{8}$:
\begin{ptree}
\AXC{$\forall s\,\forall k\,(\mathrm{NSeq}(s) \to k \in \mathbb{N} \to s_{k} \in \mathbb{N})$\lbl{[S13]}}
\AXC{$\varphi_{26}$\lbl{A2}}
\BIC{$\forall k\,(\mathrm{NSeq}(s) \to k \in \mathbb{N} \to s_{k} \in \mathbb{N})$}
\AXC{$\varphi_{21}$\lbl{A2}}
\BIC{$\mathrm{NSeq}(s) \to n \in \mathbb{N} \to s_{n} \in \mathbb{N}$}
\AXC{$(\mathrm{NSeq}(s) \to n \in \mathbb{N} \to s_{n} \in \mathbb{N}) \to \varphi_{78}$\lbl{A1}}
\BIC{$\varphi_{78}$}
\end{ptree}
\noindent$\Pi_{9}$:
\begin{ptree}
\AXC{$\mathrm{Conv}(a,\alpha) \to \mathrm{SubIdx}(s) \to \varepsilon \in \mathbb{R} \to 0 < \varepsilon \to \varphi_{23}$\lbl{A1}}
\AXC{$\stackrel{\Pi_{7}}{\vdots}$}\noLine
\UIC{$\mathrm{Conv}(a,\alpha) \to \mathrm{SubIdx}(s) \to \mathrm{NSeq}(s)$}
\AXC{$\stackrel{\Pi_{8}}{\vdots}$}\noLine
\UIC{$\varphi_{78}$}
\BIC{$\varphi_{72}$}
\BIC{$\mathrm{Conv}(a,\alpha) \to \mathrm{SubIdx}(s) \to \varepsilon \in \mathbb{R} \to 0 < \varepsilon \to \varphi_{31}$}
\end{ptree}
\noindent$\Pi_{10}$:
\begin{ptree}
\AXC{$\varphi_{40}$\lbl{A2}}
\AXC{$\mathrm{Conv}(a,\alpha) \to \mathrm{SubIdx}(s) \to \varepsilon \in \mathbb{R} \to 0 < \varepsilon \to \varphi_{42}$\lbl{A1}}
\AXC{$\varphi_{84}$\lbl{A1}}
\BIC{$\mathrm{Conv}(a,\alpha) \to \mathrm{SubIdx}(s) \to \varepsilon \in \mathbb{R} \to 0 < \varepsilon \to \varphi_{46}$}
\AXC{$\varphi_{96}$\lbl{A1}}
\BIC{$\varphi_{77}$}
\BIC{$\mathrm{Conv}(a,\alpha) \to \mathrm{SubIdx}(s) \to \varepsilon \in \mathbb{R} \to 0 < \varepsilon \to \varphi_{55}$}
\AXC{$\varphi_{100}$\lbl{A1}}
\BIC{$\varphi_{94}$}
\end{ptree}
\noindent$\Pi_{11}$:
\begin{ptree}
\AXC{$\stackrel{\Pi_{6}}{\vdots}$}\noLine
\UIC{$\mathrm{Conv}(a,\alpha) \to \mathrm{SubIdx}(s) \to \varepsilon \in \mathbb{R} \to 0 < \varepsilon \to \varphi_{29}$}
\AXC{$\stackrel{\Pi_{9}}{\vdots}$}\noLine
\UIC{$\mathrm{Conv}(a,\alpha) \to \mathrm{SubIdx}(s) \to \varepsilon \in \mathbb{R} \to 0 < \varepsilon \to \varphi_{31}$}
\AXC{$\stackrel{\Pi_{10}}{\vdots}$}\noLine
\UIC{$\varphi_{94}$}
\BIC{$\varphi_{76}$}
\BIC{$\mathrm{Conv}(a,\alpha) \to \mathrm{SubIdx}(s) \to \varepsilon \in \mathbb{R} \to 0 < \varepsilon \to \varphi_{38}$}
\end{ptree}
\noindent$\Pi_{12}$:
\begin{ptree}
\AXC{$x = x$\lbl{A5}}
\AXC{$x = y \to x = x \to y = x$\lbl{A6}}
\AXC{$(x = y \to x = x \to y = x) \to x = x \to x = y \to y = x$\lbl{A1}}
\BIC{$x = x \to x = y \to y = x$}
\BIC{$x = y \to y = x$}
\UIC{$\forall y\,(x = y \to y = x)$}
\UIC{$\forall x\,\forall y\,(x = y \to y = x)$}
\AXC{$\varphi_{12}$\lbl{A2}}
\BIC{$\forall y\,(a \circ s(n) = y \to y = a \circ s(n))$}
\AXC{$\varphi_{35}$\lbl{A2}}
\BIC{$a \circ s(n) = a_{s_{n}} \to a_{s_{n}} = a \circ s(n)$}
\end{ptree}
\noindent$\Pi_{13}$:
\begin{ptree}
\AXC{$\forall a\,\forall \alpha\,(\mathrm{Conv}(a,\alpha) \to \mathrm{Seq}(a))$\lbl{[C1]}}
\AXC{$\varphi_{7}$\lbl{A2}}
\BIC{$\forall \alpha\,(\mathrm{Conv}(a,\alpha) \to \mathrm{Seq}(a))$}
\AXC{$\varphi_{3}$\lbl{A2}}
\BIC{$\mathrm{Conv}(a,\alpha) \to \mathrm{Seq}(a)$}
\AXC{$(\mathrm{Conv}(a,\alpha) \to \mathrm{Seq}(a)) \to \varphi_{9}$\lbl{A1}}
\BIC{$\varphi_{9}$}
\end{ptree}
\noindent$\Pi_{14}$:
\begin{ptree}
\AXC{$\forall a\,\forall s\,(\mathrm{Seq}(a) \to \varphi_{17})$\lbl{[S12]}}
\AXC{$\forall a\,\forall s\,(\mathrm{Seq}(a) \to \varphi_{17}) \to \forall s\,(\mathrm{Seq}(a) \to \varphi_{17})$\lbl{A2}}
\BIC{$\forall s\,(\mathrm{Seq}(a) \to \varphi_{17})$}
\end{ptree}
\noindent$\Pi_{15}$:
\begin{ptree}
\AXC{$\stackrel{\Pi_{14}}{\vdots}$}\noLine
\UIC{$\forall s\,(\mathrm{Seq}(a) \to \varphi_{17})$}
\AXC{$\forall s\,(\mathrm{Seq}(a) \to \varphi_{17}) \to \mathrm{Seq}(a) \to \varphi_{17}$\lbl{A2}}
\BIC{$\mathrm{Seq}(a) \to \varphi_{17}$}
\AXC{$\varphi_{64}$\lbl{A1}}
\BIC{$(\mathrm{Conv}(a,\alpha) \to \mathrm{SubIdx}(s) \to \mathrm{Seq}(a)) \to \varphi_{44}$}
\end{ptree}
\noindent$\Pi_{16}$:
\begin{ptree}
\AXC{$\mathrm{Conv}(a,\alpha) \to \mathrm{Conv}(a,\alpha)$\lbl{A1}}
\AXC{$\stackrel{\Pi_{13}}{\vdots}$}\noLine
\UIC{$\varphi_{9}$}
\BIC{$\mathrm{Conv}(a,\alpha) \to \mathrm{SubIdx}(s) \to \mathrm{Seq}(a)$}
\AXC{$\stackrel{\Pi_{15}}{\vdots}$}\noLine
\UIC{$(\mathrm{Conv}(a,\alpha) \to \mathrm{SubIdx}(s) \to \mathrm{Seq}(a)) \to \varphi_{44}$}
\BIC{$\varphi_{44}$}
\end{ptree}
\noindent$\Pi_{17}$:
\begin{ptree}
\AXC{$\stackrel{\Pi_{7}}{\vdots}$}\noLine
\UIC{$\mathrm{Conv}(a,\alpha) \to \mathrm{SubIdx}(s) \to \mathrm{NSeq}(s)$}
\AXC{$\stackrel{\Pi_{16}}{\vdots}$}\noLine
\UIC{$\varphi_{44}$}
\BIC{$\mathrm{Conv}(a,\alpha) \to \varphi_{18}$}
\end{ptree}
\noindent$\Pi_{18}$:
\begin{ptree}
\AXC{$\varphi_{30}$\lbl{A2}}
\AXC{$\stackrel{\Pi_{17}}{\vdots}$}\noLine
\UIC{$\mathrm{Conv}(a,\alpha) \to \varphi_{18}$}
\AXC{$\varphi_{68}$\lbl{A1}}
\BIC{$\mathrm{Conv}(a,\alpha) \to \mathrm{SubIdx}(s) \to \varepsilon \in \mathbb{R} \to 0 < \varepsilon \to \varphi_{47}$}
\AXC{$\varphi_{93}$\lbl{A1}}
\BIC{$\varphi_{69}$}
\BIC{$\mathrm{Conv}(a,\alpha) \to \mathrm{SubIdx}(s) \to \varepsilon \in \mathbb{R} \to 0 < \varepsilon \to \varphi_{45}$}
\AXC{$\varphi_{95}$\lbl{A1}}
\BIC{$\varphi_{75}$}
\end{ptree}
\noindent$\Pi_{19}$:
\begin{ptree}
\AXC{$\mathrm{Conv}(a,\alpha) \to \mathrm{SubIdx}(s) \to \varepsilon \in \mathbb{R} \to 0 < \varepsilon \to \varphi_{23}$\lbl{A1}}
\AXC{$\stackrel{\Pi_{18}}{\vdots}$}\noLine
\UIC{$\varphi_{75}$}
\BIC{$\mathrm{Conv}(a,\alpha) \to \mathrm{SubIdx}(s) \to \varepsilon \in \mathbb{R} \to 0 < \varepsilon \to \varphi_{36}$}
\AXC{$\varphi_{89}$\lbl{A1}}
\BIC{$\varphi_{63}$}
\end{ptree}
\noindent$\Pi_{20}$:
\begin{ptree}
\AXC{$\stackrel{\Pi_{11}}{\vdots}$}\noLine
\UIC{$\mathrm{Conv}(a,\alpha) \to \mathrm{SubIdx}(s) \to \varepsilon \in \mathbb{R} \to 0 < \varepsilon \to \varphi_{38}$}
\AXC{$\stackrel{\Pi_{12}}{\vdots}$}\noLine
\UIC{$a \circ s(n) = a_{s_{n}} \to a_{s_{n}} = a \circ s(n)$}
\AXC{$\stackrel{\Pi_{19}}{\vdots}$}\noLine
\UIC{$\varphi_{63}$}
\BIC{$\mathrm{Conv}(a,\alpha) \to \mathrm{SubIdx}(s) \to \varepsilon \in \mathbb{R} \to 0 < \varepsilon \to \varphi_{37}$}
\AXC{$\varphi_{99}$\lbl{A1}}
\BIC{$\varphi_{91}$}
\BIC{$\varphi_{70}$}
\end{ptree}
\noindent$\Pi_{21}$:
\begin{ptree}
\AXC{$\varphi_{82}$\lbl{A3}}
\AXC{$z = w \to \lvert z - \alpha \rvert < \varepsilon \to \lvert w - \alpha \rvert < \varepsilon$\lbl{A6}}
\UIC{$\forall w\,(z = w \to \lvert z - \alpha \rvert < \varepsilon \to \lvert w - \alpha \rvert < \varepsilon)$}
\UIC{$\forall z\,\forall w\,(z = w \to \lvert z - \alpha \rvert < \varepsilon \to \lvert w - \alpha \rvert < \varepsilon)$}
\AXC{$\varphi_{53}$\lbl{A2}}
\BIC{$\forall w\,(a_{s_{n}} = w \to \lvert a_{s_{n}} - \alpha \rvert < \varepsilon \to \lvert w - \alpha \rvert < \varepsilon)$}
\AXC{$\varphi_{58}$\lbl{A2}}
\BIC{$\varphi_{33}$}
\AXC{$\stackrel{\Pi_{20}}{\vdots}$}\noLine
\UIC{$\varphi_{70}$}
\BIC{$\mathrm{Conv}(a,\alpha) \to \mathrm{SubIdx}(s) \to \varepsilon \in \mathbb{R} \to 0 < \varepsilon \to \varphi_{39}$}
\AXC{$\varphi_{81}$\lbl{A1}}
\BIC{$\varphi_{34} \to n \in \mathbb{N} \to N \leq n \to \lvert a \circ s(n) - \alpha \rvert < \varepsilon$}
\UIC{$\forall n\,(\varphi_{34} \to n \in \mathbb{N} \to N \leq n \to \lvert a \circ s(n) - \alpha \rvert < \varepsilon)$}
\AXC{$\varphi_{101}$\lbl{A1}}
\BIC{$\varphi_{82} \to \varphi_{88}$}
\BIC{$\varphi_{88}$}
\end{ptree}
\noindent$\Pi_{22}$:
\begin{ptree}
\AXC{$\mathrm{Conv}(a,\alpha) \to \mathrm{SubIdx}(s) \to \varepsilon \in \mathbb{R} \to 0 < \varepsilon \to \varphi_{42}$\lbl{A1}}
\AXC{$\varphi_{74}$\lbl{A1}}
\BIC{$\mathrm{Conv}(a,\alpha) \to \mathrm{SubIdx}(s) \to \varepsilon \in \mathbb{R} \to \varphi_{22}$}
\end{ptree}
\noindent$\Pi_{23}$:
\begin{ptree}
\AXC{$\stackrel{\Pi_{22}}{\vdots}$}\noLine
\UIC{$\mathrm{Conv}(a,\alpha) \to \mathrm{SubIdx}(s) \to \varepsilon \in \mathbb{R} \to \varphi_{22}$}
\AXC{$(\mathrm{Conv}(a,\alpha) \to \mathrm{SubIdx}(s) \to \varepsilon \in \mathbb{R} \to \varphi_{22}) \to \varphi_{83}$\lbl{A1}}
\BIC{$\varphi_{83}$}
\end{ptree}
\noindent$\Pi_{24}$:
\begin{ptree}
\AXC{$\forall N\,(\varphi_{16}) \to \varphi_{16}$\lbl{A2}}
\AXC{$\varphi_{34} \to \forall n\,(\varphi_{34})$\lbl{A4}}
\AXC{$\stackrel{\Pi_{21}}{\vdots}$}\noLine
\UIC{$\varphi_{88}$}
\BIC{$\mathrm{Conv}(a,\alpha) \to \mathrm{SubIdx}(s) \to \varepsilon \in \mathbb{R} \to 0 < \varepsilon \to \varphi_{41}$}
\AXC{$\stackrel{\Pi_{23}}{\vdots}$}\noLine
\UIC{$\varphi_{83}$}
\BIC{$\mathrm{Conv}(a,\alpha) \to \mathrm{SubIdx}(s) \to \varepsilon \in \mathbb{R} \to 0 < \varepsilon \to \varphi_{48}$}
\AXC{$\varphi_{97}$\lbl{A1}}
\BIC{$\varphi_{80}$}
\BIC{$\mathrm{Conv}(a,\alpha) \to \mathrm{SubIdx}(s) \to \varepsilon \in \mathbb{R} \to 0 < \varepsilon \to \varphi_{52}$}
\AXC{$\varphi_{86}$\lbl{A1}}
\BIC{$\varphi_{43} \to \neg (N \in \mathbb{N} \land \forall n \in \mathbb{N}\,(N \leq n \to \lvert a_{n} - \alpha \rvert < \varepsilon))$}
\UIC{$\varphi_{59}$}
\end{ptree}
\noindent$\Pi_{25}$:
\begin{ptree}
\AXC{$\forall a\,\forall \alpha\,\forall \varepsilon\,(\mathrm{Conv}(a,\alpha) \to \varphi_{27})$\lbl{[C3]}}
\AXC{$\varphi_{61}$\lbl{A2}}
\BIC{$\forall \alpha\,\forall \varepsilon\,(\mathrm{Conv}(a,\alpha) \to \varphi_{27})$}
\AXC{$\varphi_{60}$\lbl{A2}}
\BIC{$\forall \varepsilon\,(\mathrm{Conv}(a,\alpha) \to \varphi_{27})$}
\AXC{$\forall \varepsilon\,(\mathrm{Conv}(a,\alpha) \to \varphi_{27}) \to \mathrm{Conv}(a,\alpha) \to \varphi_{27}$\lbl{A2}}
\BIC{$\mathrm{Conv}(a,\alpha) \to \varphi_{27}$}
\end{ptree}
\noindent$\Pi_{26}$:
\begin{ptree}
\AXC{$\mathrm{Conv}(a,\alpha) \to \mathrm{Conv}(a,\alpha)$\lbl{A1}}
\AXC{$\stackrel{\Pi_{25}}{\vdots}$}\noLine
\UIC{$\mathrm{Conv}(a,\alpha) \to \varphi_{27}$}
\AXC{$\varphi_{85}$\lbl{A1}}
\BIC{$(\mathrm{Conv}(a,\alpha) \to \mathrm{Conv}(a,\alpha)) \to \varphi_{62}$}
\BIC{$\varphi_{62}$}
\end{ptree}
\noindent$\Pi_{27}$:
\begin{ptree}
\AXC{$\stackrel{\Pi_{24}}{\vdots}$}\noLine
\UIC{$\varphi_{59}$}
\AXC{$\varphi_{10}$\lbl{A1}}
\AXC{$\mathrm{Conv}(a,\alpha) \to \mathrm{SubIdx}(s) \to \varepsilon \in \mathbb{R} \to \varepsilon \in \mathbb{R}$\lbl{A1}}
\AXC{$\stackrel{\Pi_{26}}{\vdots}$}\noLine
\UIC{$\varphi_{62}$}
\BIC{$\varphi_{10} \to \mathrm{Conv}(a,\alpha) \to \mathrm{SubIdx}(s) \to \varphi_{27}$}
\BIC{$\mathrm{Conv}(a,\alpha) \to \mathrm{SubIdx}(s) \to \varphi_{27}$}
\AXC{$\varphi_{102}$\lbl{A1}}
\BIC{$\varphi_{59} \to (\varphi_{59} \to \forall N\,(\varphi_{43}) \to \varphi_{15}) \to \varphi_{87}$}
\BIC{$(\varphi_{59} \to \forall N\,(\varphi_{43}) \to \varphi_{15}) \to \varphi_{87}$}
\end{ptree}
\noindent$\Pi_{28}$:
\begin{ptree}
\AXC{$\varphi_{67}$\lbl{A3}}
\AXC{$\varphi_{43} \to \forall N\,(\varphi_{43})$\lbl{A4}}
\AXC{$\varphi_{59} \to \forall N\,(\varphi_{43}) \to \varphi_{15}$\lbl{A3}}
\AXC{$\stackrel{\Pi_{27}}{\vdots}$}\noLine
\UIC{$(\varphi_{59} \to \forall N\,(\varphi_{43}) \to \varphi_{15}) \to \varphi_{87}$}
\BIC{$\varphi_{87}$}
\BIC{$\mathrm{Conv}(a,\alpha) \to \mathrm{SubIdx}(s) \to \varepsilon \in \mathbb{R} \to \varphi_{20}$}
\AXC{$\varphi_{65}$\lbl{A1}}
\BIC{$\mathrm{Conv}(a,\alpha) \land \mathrm{SubIdx}(s) \to \varepsilon \in \mathbb{R} \to \varphi_{20}$}
\UIC{$\forall \varepsilon\,(\mathrm{Conv}(a,\alpha) \land \mathrm{SubIdx}(s) \to \varepsilon \in \mathbb{R} \to \varphi_{20})$}
\AXC{$\varphi_{98}$\lbl{A1}}
\BIC{$\varphi_{92}$}
\BIC{$\varphi_{6} \to \mathrm{Conv}(a,\alpha) \to \mathrm{SubIdx}(s) \to \forall \varepsilon \in \mathbb{R}\,(\varphi_{20})$}
\end{ptree}
\noindent$\Pi_{29}$:
\begin{ptree}
\AXC{$\forall a\,\forall \alpha\,(\mathrm{Conv}(a,\alpha) \to \alpha \in \mathbb{R})$\lbl{[C2]}}
\AXC{$\varphi_{13}$\lbl{A2}}
\BIC{$\forall \alpha\,(\mathrm{Conv}(a,\alpha) \to \alpha \in \mathbb{R})$}
\AXC{$\varphi_{8}$\lbl{A2}}
\BIC{$\mathrm{Conv}(a,\alpha) \to \alpha \in \mathbb{R}$}
\AXC{$(\mathrm{Conv}(a,\alpha) \to \alpha \in \mathbb{R}) \to \varphi_{14}$\lbl{A1}}
\BIC{$\varphi_{14}$}
\end{ptree}
\noindent$\Pi_{30}$:
\begin{ptree}
\AXC{$\stackrel{\Pi_{17}}{\vdots}$}\noLine
\UIC{$\mathrm{Conv}(a,\alpha) \to \varphi_{18}$}
\AXC{$\varphi_{49}$\lbl{A1}}
\BIC{$\mathrm{Conv}(a,\alpha) \to \mathrm{SubIdx}(s) \to \mathrm{Seq}(a \circ s)$}
\AXC{$\varphi_{54}$\lbl{DefConv}}
\AXC{$\varphi_{54} \to \forall \alpha\,(\varphi_{56})$\lbl{A2}}
\BIC{$\forall \alpha\,(\varphi_{56})$}
\AXC{$\forall \alpha\,(\varphi_{56}) \to \varphi_{56}$\lbl{A2}}
\BIC{$\varphi_{56}$}
\AXC{$\varphi_{56} \to \varphi_{66}$\lbl{A1}}
\BIC{$\varphi_{66}$}
\BIC{$(\mathrm{Conv}(a,\alpha) \to \mathrm{SubIdx}(s) \to \alpha \in \mathbb{R}) \to \varphi_{57}$}
\end{ptree}
\noindent$\Pi_{31}$:
\begin{ptree}
\AXC{$\mathrm{Conv}(a,\alpha) \to \mathrm{Conv}(a,\alpha)$\lbl{A1}}
\AXC{$\stackrel{\Pi_{29}}{\vdots}$}\noLine
\UIC{$\varphi_{14}$}
\BIC{$\mathrm{Conv}(a,\alpha) \to \mathrm{SubIdx}(s) \to \alpha \in \mathbb{R}$}
\AXC{$\stackrel{\Pi_{30}}{\vdots}$}\noLine
\UIC{$(\mathrm{Conv}(a,\alpha) \to \mathrm{SubIdx}(s) \to \alpha \in \mathbb{R}) \to \varphi_{57}$}
\BIC{$\varphi_{57}$}
\end{ptree}
\noindent 主木:
\begin{ptree}
\AXC{$\varphi_{6}$\lbl{A4}}
\AXC{$\stackrel{\Pi_{28}}{\vdots}$}\noLine
\UIC{$\varphi_{6} \to \mathrm{Conv}(a,\alpha) \to \mathrm{SubIdx}(s) \to \forall \varepsilon \in \mathbb{R}\,(\varphi_{20})$}
\BIC{$\mathrm{Conv}(a,\alpha) \to \mathrm{SubIdx}(s) \to \forall \varepsilon \in \mathbb{R}\,(\varphi_{20})$}
\AXC{$\stackrel{\Pi_{31}}{\vdots}$}\noLine
\UIC{$\varphi_{57}$}
\BIC{$\varphi_{2}$}
\UIC{$\forall \alpha\,(\varphi_{2})$}
\UIC{$\forall s\,\forall \alpha\,(\varphi_{2})$}
\UIC{$\forall a\,\forall s\,\forall \alpha\,(\varphi_{2})$}
\end{ptree}
\begin{abbrevlist}
\item $\varphi_{1} :\equiv \forall y\,\forall z\,(N \leq y \to y \leq z \to N \leq z) \to \forall z\,(N \leq n \to n \leq z \to N \leq z)$
\item $\varphi_{2} :\equiv \mathrm{Conv}(a,\alpha) \to \mathrm{SubIdx}(s) \to \mathrm{Conv}(a \circ s,\alpha)$
\item $\varphi_{3} :\equiv \forall \alpha\,(\mathrm{Conv}(a,\alpha) \to \mathrm{Seq}(a)) \to \mathrm{Conv}(a,\alpha) \to \mathrm{Seq}(a)$
\item $\varphi_{4} :\equiv \forall x\,\forall y\,\forall z\,(x \leq y \to y \leq z \to x \leq z) \to \forall y\,\forall z\,(N \leq y \to y \leq z \to N \leq z)$
\item $\varphi_{5} :\equiv \mathrm{Conv}(a,\alpha) \to \mathrm{SubIdx}(s) \to \forall k \in \mathbb{N}\,(k \leq s_{k})$
\item $\varphi_{6} :\equiv \mathrm{Conv}(a,\alpha) \land \mathrm{SubIdx}(s) \to \forall \varepsilon\,(\mathrm{Conv}(a,\alpha) \land \mathrm{SubIdx}(s))$
\item $\varphi_{7} :\equiv \forall a\,\forall \alpha\,(\mathrm{Conv}(a,\alpha) \to \mathrm{Seq}(a)) \to \forall \alpha\,(\mathrm{Conv}(a,\alpha) \to \mathrm{Seq}(a))$
\item $\varphi_{8} :\equiv \forall \alpha\,(\mathrm{Conv}(a,\alpha) \to \alpha \in \mathbb{R}) \to \mathrm{Conv}(a,\alpha) \to \alpha \in \mathbb{R}$
\item $\varphi_{9} :\equiv (\mathrm{Conv}(a,\alpha) \to \mathrm{Conv}(a,\alpha)) \to \mathrm{Conv}(a,\alpha) \to \mathrm{SubIdx}(s) \to \mathrm{Seq}(a)$
\item $\varphi_{10} :\equiv \mathrm{Conv}(a,\alpha) \to \mathrm{SubIdx}(s) \to \varepsilon \in \mathbb{R} \to 0 < \varepsilon \to 0 < \varepsilon$
\item $\varphi_{11} :\equiv \mathrm{SubIdx}(s) \leftrightarrow \mathrm{NSeq}(s) \land \forall k \in \mathbb{N}\,(s_{k} < s_{S(k)})$
\item $\varphi_{12} :\equiv \forall x\,\forall y\,(x = y \to y = x) \to \forall y\,(a \circ s(n) = y \to y = a \circ s(n))$
\item $\varphi_{13} :\equiv \forall a\,\forall \alpha\,(\mathrm{Conv}(a,\alpha) \to \alpha \in \mathbb{R}) \to \forall \alpha\,(\mathrm{Conv}(a,\alpha) \to \alpha \in \mathbb{R})$
\item $\varphi_{14} :\equiv (\mathrm{Conv}(a,\alpha) \to \mathrm{Conv}(a,\alpha)) \to \mathrm{Conv}(a,\alpha) \to \mathrm{SubIdx}(s) \to \alpha \in \mathbb{R}$
\item $\varphi_{15} :\equiv \forall N\,(\neg (N \in \mathbb{N} \land \forall n \in \mathbb{N}\,(N \leq n \to \lvert a_{n} - \alpha \rvert < \varepsilon)))$
\item $\varphi_{16} :\equiv \neg (N \in \mathbb{N} \land \forall n \in \mathbb{N}\,(N \leq n \to \lvert a \circ s(n) - \alpha \rvert < \varepsilon))$
\item $\varphi_{17} :\equiv \mathrm{NSeq}(s) \to \mathrm{Seq}(a \circ s) \land \forall k \in \mathbb{N}\,(a \circ s(k) = a_{s_{k}})$
\item $\varphi_{18} :\equiv \mathrm{SubIdx}(s) \to \mathrm{Seq}(a \circ s) \land \forall k \in \mathbb{N}\,(a \circ s(k) = a_{s_{k}})$
\item $\varphi_{19} :\equiv \forall s\,(\mathrm{SubIdx}(s) \to \forall k \in \mathbb{N}\,(k \leq s_{k})) \to \mathrm{SubIdx}(s) \to \forall k \in \mathbb{N}\,(k \leq s_{k})$
\item $\varphi_{20} :\equiv 0 < \varepsilon \to \exists N \in \mathbb{N}\,\forall n \in \mathbb{N}\,(N \leq n \to \lvert a \circ s(n) - \alpha \rvert < \varepsilon)$
\item $\varphi_{21} :\equiv \forall k\,(\mathrm{NSeq}(s) \to k \in \mathbb{N} \to s_{k} \in \mathbb{N}) \to \mathrm{NSeq}(s) \to n \in \mathbb{N} \to s_{n} \in \mathbb{N}$
\item $\varphi_{22} :\equiv 0 < \varepsilon \to N \in \mathbb{N} \land \forall n \in \mathbb{N}\,(N \leq n \to \lvert a_{n} - \alpha \rvert < \varepsilon) \to N \in \mathbb{N}$
\item $\varphi_{23} :\equiv N \in \mathbb{N} \land \forall n \in \mathbb{N}\,(N \leq n \to \lvert a_{n} - \alpha \rvert < \varepsilon) \to n \in \mathbb{N} \to n \in \mathbb{N}$
\item $\varphi_{24} :\equiv N \in \mathbb{N} \land \forall n \in \mathbb{N}\,(N \leq n \to \lvert a_{n} - \alpha \rvert < \varepsilon) \to n \in \mathbb{N} \to N \leq n \to N \leq n$
\item $\varphi_{25} :\equiv (\mathrm{Conv}(a,\alpha) \to \mathrm{SubIdx}(s) \to \mathrm{SubIdx}(s)) \to \varphi_{5}$
\item $\varphi_{26} :\equiv \forall s\,\forall k\,(\mathrm{NSeq}(s) \to k \in \mathbb{N} \to s_{k} \in \mathbb{N}) \to \forall k\,(\mathrm{NSeq}(s) \to k \in \mathbb{N} \to s_{k} \in \mathbb{N})$
\item $\varphi_{27} :\equiv \varepsilon \in \mathbb{R} \to 0 < \varepsilon \to \exists N \in \mathbb{N}\,\forall n \in \mathbb{N}\,(N \leq n \to \lvert a_{n} - \alpha \rvert < \varepsilon)$
\item $\varphi_{28} :\equiv N \in \mathbb{N} \land \forall n \in \mathbb{N}\,(N \leq n \to \lvert a_{n} - \alpha \rvert < \varepsilon) \to n \in \mathbb{N} \to N \leq n \to n \leq s_{n}$
\item $\varphi_{29} :\equiv N \in \mathbb{N} \land \forall n \in \mathbb{N}\,(N \leq n \to \lvert a_{n} - \alpha \rvert < \varepsilon) \to n \in \mathbb{N} \to N \leq n \to N \leq s_{n}$
\item $\varphi_{30} :\equiv \forall k \in \mathbb{N}\,(a \circ s(k) = a_{s_{k}}) \to n \in \mathbb{N} \to a \circ s(n) = a_{s_{n}}$
\item $\varphi_{31} :\equiv N \in \mathbb{N} \land \forall n \in \mathbb{N}\,(N \leq n \to \lvert a_{n} - \alpha \rvert < \varepsilon) \to n \in \mathbb{N} \to N \leq n \to s_{n} \in \mathbb{N}$
\item $\varphi_{32} :\equiv (\forall k \in \mathbb{N}\,(k \leq s_{k}) \to n \in \mathbb{N} \to n \leq s_{n}) \to \mathrm{Conv}(a,\alpha) \to \mathrm{SubIdx}(s) \to n \in \mathbb{N} \to n \leq s_{n}$
\item $\varphi_{33} :\equiv a_{s_{n}} = a \circ s(n) \to \lvert a_{s_{n}} - \alpha \rvert < \varepsilon \to \lvert a \circ s(n) - \alpha \rvert < \varepsilon$
\item $\varphi_{34} :\equiv \mathrm{Conv}(a,\alpha) \land \mathrm{SubIdx}(s) \land \varepsilon \in \mathbb{R} \land 0 < \varepsilon \land (N \in \mathbb{N} \land \forall n \in \mathbb{N}\,(N \leq n \to \lvert a_{n} - \alpha \rvert < \varepsilon))$
\item $\varphi_{35} :\equiv \forall y\,(a \circ s(n) = y \to y = a \circ s(n)) \to a \circ s(n) = a_{s_{n}} \to a_{s_{n}} = a \circ s(n)$
\item $\varphi_{36} :\equiv N \in \mathbb{N} \land \forall n \in \mathbb{N}\,(N \leq n \to \lvert a_{n} - \alpha \rvert < \varepsilon) \to n \in \mathbb{N} \to N \leq n \to a \circ s(n) = a_{s_{n}}$
\item $\varphi_{37} :\equiv N \in \mathbb{N} \land \forall n \in \mathbb{N}\,(N \leq n \to \lvert a_{n} - \alpha \rvert < \varepsilon) \to n \in \mathbb{N} \to N \leq n \to a_{s_{n}} = a \circ s(n)$
\item $\varphi_{38} :\equiv N \in \mathbb{N} \land \forall n \in \mathbb{N}\,(N \leq n \to \lvert a_{n} - \alpha \rvert < \varepsilon) \to n \in \mathbb{N} \to N \leq n \to \lvert a_{s_{n}} - \alpha \rvert < \varepsilon$
\item $\varphi_{39} :\equiv N \in \mathbb{N} \land \forall n \in \mathbb{N}\,(N \leq n \to \lvert a_{n} - \alpha \rvert < \varepsilon) \to n \in \mathbb{N} \to N \leq n \to \lvert a \circ s(n) - \alpha \rvert < \varepsilon$
\item $\varphi_{40} :\equiv \forall n \in \mathbb{N}\,(N \leq n \to \lvert a_{n} - \alpha \rvert < \varepsilon) \to s_{n} \in \mathbb{N} \to N \leq s_{n} \to \lvert a_{s_{n}} - \alpha \rvert < \varepsilon$
\item $\varphi_{41} :\equiv N \in \mathbb{N} \land \forall n \in \mathbb{N}\,(N \leq n \to \lvert a_{n} - \alpha \rvert < \varepsilon) \to \forall n \in \mathbb{N}\,(N \leq n \to \lvert a \circ s(n) - \alpha \rvert < \varepsilon)$
\item $\varphi_{42} :\equiv N \in \mathbb{N} \land \forall n \in \mathbb{N}\,(N \leq n \to \lvert a_{n} - \alpha \rvert < \varepsilon) \to N \in \mathbb{N} \land \forall n \in \mathbb{N}\,(N \leq n \to \lvert a_{n} - \alpha \rvert < \varepsilon)$
\item $\varphi_{43} :\equiv \mathrm{Conv}(a,\alpha) \land \mathrm{SubIdx}(s) \land \varepsilon \in \mathbb{R} \land 0 < \varepsilon \land \neg \exists N \in \mathbb{N}\,\forall n \in \mathbb{N}\,(N \leq n \to \lvert a \circ s(n) - \alpha \rvert < \varepsilon)$
\item $\varphi_{44} :\equiv (\mathrm{Conv}(a,\alpha) \to \mathrm{SubIdx}(s) \to \mathrm{NSeq}(s)) \to \mathrm{Conv}(a,\alpha) \to \varphi_{18}$
\item $\varphi_{45} :\equiv N \in \mathbb{N} \land \forall n \in \mathbb{N}\,(N \leq n \to \lvert a_{n} - \alpha \rvert < \varepsilon) \to n \in \mathbb{N} \to N \leq n \to n \in \mathbb{N} \to a \circ s(n) = a_{s_{n}}$
\item $\varphi_{46} :\equiv N \in \mathbb{N} \land \forall n \in \mathbb{N}\,(N \leq n \to \lvert a_{n} - \alpha \rvert < \varepsilon) \to n \in \mathbb{N} \to N \leq n \to \forall n \in \mathbb{N}\,(N \leq n \to \lvert a_{n} - \alpha \rvert < \varepsilon)$
\item $\varphi_{47} :\equiv N \in \mathbb{N} \land \forall n \in \mathbb{N}\,(N \leq n \to \lvert a_{n} - \alpha \rvert < \varepsilon) \to n \in \mathbb{N} \to N \leq n \to \forall k \in \mathbb{N}\,(a \circ s(k) = a_{s_{k}})$
\item $\varphi_{48} :\equiv N \in \mathbb{N} \land \forall n \in \mathbb{N}\,(N \leq n \to \lvert a_{n} - \alpha \rvert < \varepsilon) \to N \in \mathbb{N} \land \forall n \in \mathbb{N}\,(N \leq n \to \lvert a \circ s(n) - \alpha \rvert < \varepsilon)$
\item $\varphi_{49} :\equiv (\mathrm{Conv}(a,\alpha) \to \varphi_{18}) \to \mathrm{Conv}(a,\alpha) \to \mathrm{SubIdx}(s) \to \mathrm{Seq}(a \circ s)$
\item $\varphi_{50} :\equiv \forall \varepsilon\,(\mathrm{Conv}(a,\alpha) \land \mathrm{SubIdx}(s)) \to \forall \varepsilon \in \mathbb{R}\,(\varphi_{20})$
\item $\varphi_{51} :\equiv (\mathrm{Conv}(a,\alpha) \to \mathrm{SubIdx}(s) \to \mathrm{SubIdx}(s)) \to \varphi_{11} \to \mathrm{Conv}(a,\alpha) \to \mathrm{SubIdx}(s) \to \mathrm{NSeq}(s)$
\item $\varphi_{52} :\equiv N \in \mathbb{N} \land \forall n \in \mathbb{N}\,(N \leq n \to \lvert a_{n} - \alpha \rvert < \varepsilon) \to \exists N \in \mathbb{N}\,\forall n \in \mathbb{N}\,(N \leq n \to \lvert a \circ s(n) - \alpha \rvert < \varepsilon)$
\item $\varphi_{53} :\equiv \forall z\,\forall w\,(z = w \to \lvert z - \alpha \rvert < \varepsilon \to \lvert w - \alpha \rvert < \varepsilon) \to \forall w\,(a_{s_{n}} = w \to \lvert a_{s_{n}} - \alpha \rvert < \varepsilon \to \lvert w - \alpha \rvert < \varepsilon)$
\item $\varphi_{54} :\equiv \forall a\,\forall \alpha\,(\mathrm{Conv}(a,\alpha) \leftrightarrow \mathrm{Seq}(a) \land \alpha \in \mathbb{R} \land \forall \varepsilon \in \mathbb{R}\,(0 < \varepsilon \to \exists N \in \mathbb{N}\,\forall n \in \mathbb{N}\,(N \leq n \to \lvert a_{n} - \alpha \rvert < \varepsilon)))$
\item $\varphi_{55} :\equiv N \in \mathbb{N} \land \forall n \in \mathbb{N}\,(N \leq n \to \lvert a_{n} - \alpha \rvert < \varepsilon) \to n \in \mathbb{N} \to N \leq n \to s_{n} \in \mathbb{N} \to N \leq s_{n} \to \lvert a_{s_{n}} - \alpha \rvert < \varepsilon$
\item $\varphi_{56} :\equiv \mathrm{Conv}(a \circ s,\alpha) \leftrightarrow \mathrm{Seq}(a \circ s) \land \alpha \in \mathbb{R} \land \forall \varepsilon \in \mathbb{R}\,(\varphi_{20})$
\item $\varphi_{57} :\equiv (\mathrm{Conv}(a,\alpha) \to \mathrm{SubIdx}(s) \to \forall \varepsilon \in \mathbb{R}\,(\varphi_{20})) \to \varphi_{2}$
\item $\varphi_{58} :\equiv \forall w\,(a_{s_{n}} = w \to \lvert a_{s_{n}} - \alpha \rvert < \varepsilon \to \lvert w - \alpha \rvert < \varepsilon) \to \varphi_{33}$
\item $\varphi_{59} :\equiv \forall N\,(\varphi_{43} \to \neg (N \in \mathbb{N} \land \forall n \in \mathbb{N}\,(N \leq n \to \lvert a_{n} - \alpha \rvert < \varepsilon)))$
\item $\varphi_{60} :\equiv \forall \alpha\,\forall \varepsilon\,(\mathrm{Conv}(a,\alpha) \to \varphi_{27}) \to \forall \varepsilon\,(\mathrm{Conv}(a,\alpha) \to \varphi_{27})$
\item $\varphi_{61} :\equiv \forall a\,\forall \alpha\,\forall \varepsilon\,(\mathrm{Conv}(a,\alpha) \to \varphi_{27}) \to \forall \alpha\,\forall \varepsilon\,(\mathrm{Conv}(a,\alpha) \to \varphi_{27})$
\item $\varphi_{62} :\equiv (\mathrm{Conv}(a,\alpha) \to \mathrm{SubIdx}(s) \to \varepsilon \in \mathbb{R} \to \varepsilon \in \mathbb{R}) \to \varphi_{10} \to \mathrm{Conv}(a,\alpha) \to \mathrm{SubIdx}(s) \to \varphi_{27}$
\item $\varphi_{63} :\equiv (a \circ s(n) = a_{s_{n}} \to a_{s_{n}} = a \circ s(n)) \to \mathrm{Conv}(a,\alpha) \to \mathrm{SubIdx}(s) \to \varepsilon \in \mathbb{R} \to 0 < \varepsilon \to \varphi_{37}$
\item $\varphi_{64} :\equiv (\mathrm{Seq}(a) \to \varphi_{17}) \to (\mathrm{Conv}(a,\alpha) \to \mathrm{SubIdx}(s) \to \mathrm{Seq}(a)) \to \varphi_{44}$
\item $\varphi_{65} :\equiv (\mathrm{Conv}(a,\alpha) \to \mathrm{SubIdx}(s) \to \varepsilon \in \mathbb{R} \to \varphi_{20}) \to \mathrm{Conv}(a,\alpha) \land \mathrm{SubIdx}(s) \to \varepsilon \in \mathbb{R} \to \varphi_{20}$
\item $\varphi_{66} :\equiv (\mathrm{Conv}(a,\alpha) \to \mathrm{SubIdx}(s) \to \mathrm{Seq}(a \circ s)) \to (\mathrm{Conv}(a,\alpha) \to \mathrm{SubIdx}(s) \to \alpha \in \mathbb{R}) \to \varphi_{57}$
\item $\varphi_{67} :\equiv \forall \varepsilon\,(\mathrm{Conv}(a,\alpha) \land \mathrm{SubIdx}(s) \to \varepsilon \in \mathbb{R} \to \varphi_{20}) \to \varphi_{50}$
\item $\varphi_{68} :\equiv (\mathrm{Conv}(a,\alpha) \to \varphi_{18}) \to \mathrm{Conv}(a,\alpha) \to \mathrm{SubIdx}(s) \to \varepsilon \in \mathbb{R} \to 0 < \varepsilon \to \varphi_{47}$
\item $\varphi_{69} :\equiv \varphi_{30} \to \mathrm{Conv}(a,\alpha) \to \mathrm{SubIdx}(s) \to \varepsilon \in \mathbb{R} \to 0 < \varepsilon \to \varphi_{45}$
\item $\varphi_{70} :\equiv \varphi_{33} \to \mathrm{Conv}(a,\alpha) \to \mathrm{SubIdx}(s) \to \varepsilon \in \mathbb{R} \to 0 < \varepsilon \to \varphi_{39}$
\item $\varphi_{71} :\equiv (\mathrm{Conv}(a,\alpha) \to \mathrm{SubIdx}(s) \to \varepsilon \in \mathbb{R} \to 0 < \varepsilon \to \varphi_{23}) \to \mathrm{Conv}(a,\alpha) \to \mathrm{SubIdx}(s) \to \varepsilon \in \mathbb{R} \to 0 < \varepsilon \to \varphi_{28}$
\item $\varphi_{72} :\equiv (\mathrm{Conv}(a,\alpha) \to \mathrm{SubIdx}(s) \to \varepsilon \in \mathbb{R} \to 0 < \varepsilon \to \varphi_{23}) \to \mathrm{Conv}(a,\alpha) \to \mathrm{SubIdx}(s) \to \varepsilon \in \mathbb{R} \to 0 < \varepsilon \to \varphi_{31}$
\item $\varphi_{73} :\equiv (\mathrm{Conv}(a,\alpha) \to \mathrm{SubIdx}(s) \to \varepsilon \in \mathbb{R} \to 0 < \varepsilon \to \varphi_{28}) \to \mathrm{Conv}(a,\alpha) \to \mathrm{SubIdx}(s) \to \varepsilon \in \mathbb{R} \to 0 < \varepsilon \to \varphi_{29}$
\item $\varphi_{74} :\equiv (\mathrm{Conv}(a,\alpha) \to \mathrm{SubIdx}(s) \to \varepsilon \in \mathbb{R} \to 0 < \varepsilon \to \varphi_{42}) \to \mathrm{Conv}(a,\alpha) \to \mathrm{SubIdx}(s) \to \varepsilon \in \mathbb{R} \to \varphi_{22}$
\item $\varphi_{75} :\equiv (\mathrm{Conv}(a,\alpha) \to \mathrm{SubIdx}(s) \to \varepsilon \in \mathbb{R} \to 0 < \varepsilon \to \varphi_{23}) \to \mathrm{Conv}(a,\alpha) \to \mathrm{SubIdx}(s) \to \varepsilon \in \mathbb{R} \to 0 < \varepsilon \to \varphi_{36}$
\item $\varphi_{76} :\equiv (\mathrm{Conv}(a,\alpha) \to \mathrm{SubIdx}(s) \to \varepsilon \in \mathbb{R} \to 0 < \varepsilon \to \varphi_{29}) \to \mathrm{Conv}(a,\alpha) \to \mathrm{SubIdx}(s) \to \varepsilon \in \mathbb{R} \to 0 < \varepsilon \to \varphi_{38}$
\item $\varphi_{77} :\equiv \varphi_{40} \to \mathrm{Conv}(a,\alpha) \to \mathrm{SubIdx}(s) \to \varepsilon \in \mathbb{R} \to 0 < \varepsilon \to \varphi_{55}$
\item $\varphi_{78} :\equiv (\mathrm{Conv}(a,\alpha) \to \mathrm{SubIdx}(s) \to \mathrm{NSeq}(s)) \to \varphi_{72}$
\item $\varphi_{79} :\equiv (\mathrm{Conv}(a,\alpha) \to \mathrm{SubIdx}(s) \to n \in \mathbb{N} \to n \leq s_{n}) \to \varphi_{71}$
\item $\varphi_{80} :\equiv (\forall N\,(\varphi_{16}) \to \varphi_{16}) \to \mathrm{Conv}(a,\alpha) \to \mathrm{SubIdx}(s) \to \varepsilon \in \mathbb{R} \to 0 < \varepsilon \to \varphi_{52}$
\item $\varphi_{81} :\equiv (\mathrm{Conv}(a,\alpha) \to \mathrm{SubIdx}(s) \to \varepsilon \in \mathbb{R} \to 0 < \varepsilon \to \varphi_{39}) \to \varphi_{34} \to n \in \mathbb{N} \to N \leq n \to \lvert a \circ s(n) - \alpha \rvert < \varepsilon$
\item $\varphi_{82} :\equiv \forall n\,(\varphi_{34} \to n \in \mathbb{N} \to N \leq n \to \lvert a \circ s(n) - \alpha \rvert < \varepsilon) \to \forall n\,(\varphi_{34}) \to \forall n \in \mathbb{N}\,(N \leq n \to \lvert a \circ s(n) - \alpha \rvert < \varepsilon)$
\item $\varphi_{83} :\equiv (\mathrm{Conv}(a,\alpha) \to \mathrm{SubIdx}(s) \to \varepsilon \in \mathbb{R} \to 0 < \varepsilon \to \varphi_{41}) \to \mathrm{Conv}(a,\alpha) \to \mathrm{SubIdx}(s) \to \varepsilon \in \mathbb{R} \to 0 < \varepsilon \to \varphi_{48}$
\item $\varphi_{84} :\equiv (\mathrm{Conv}(a,\alpha) \to \mathrm{SubIdx}(s) \to \varepsilon \in \mathbb{R} \to 0 < \varepsilon \to \varphi_{42}) \to \mathrm{Conv}(a,\alpha) \to \mathrm{SubIdx}(s) \to \varepsilon \in \mathbb{R} \to 0 < \varepsilon \to \varphi_{46}$
\item $\varphi_{85} :\equiv (\mathrm{Conv}(a,\alpha) \to \varphi_{27}) \to (\mathrm{Conv}(a,\alpha) \to \mathrm{Conv}(a,\alpha)) \to \varphi_{62}$
\item $\varphi_{86} :\equiv (\mathrm{Conv}(a,\alpha) \to \mathrm{SubIdx}(s) \to \varepsilon \in \mathbb{R} \to 0 < \varepsilon \to \varphi_{52}) \to \varphi_{43} \to \neg (N \in \mathbb{N} \land \forall n \in \mathbb{N}\,(N \leq n \to \lvert a_{n} - \alpha \rvert < \varepsilon))$
\item $\varphi_{87} :\equiv (\varphi_{43} \to \forall N\,(\varphi_{43})) \to \mathrm{Conv}(a,\alpha) \to \mathrm{SubIdx}(s) \to \varepsilon \in \mathbb{R} \to \varphi_{20}$
\item $\varphi_{88} :\equiv (\varphi_{34} \to \forall n\,(\varphi_{34})) \to \mathrm{Conv}(a,\alpha) \to \mathrm{SubIdx}(s) \to \varepsilon \in \mathbb{R} \to 0 < \varepsilon \to \varphi_{41}$
\item $\varphi_{89} :\equiv (\mathrm{Conv}(a,\alpha) \to \mathrm{SubIdx}(s) \to \varepsilon \in \mathbb{R} \to 0 < \varepsilon \to \varphi_{36}) \to \varphi_{63}$
\item $\varphi_{90} :\equiv (\mathrm{Conv}(a,\alpha) \to \mathrm{SubIdx}(s) \to \varepsilon \in \mathbb{R} \to 0 < \varepsilon \to \varphi_{24}) \to \varphi_{73}$
\item $\varphi_{91} :\equiv (\mathrm{Conv}(a,\alpha) \to \mathrm{SubIdx}(s) \to \varepsilon \in \mathbb{R} \to 0 < \varepsilon \to \varphi_{38}) \to \varphi_{70}$
\item $\varphi_{92} :\equiv \varphi_{67} \to \varphi_{6} \to \mathrm{Conv}(a,\alpha) \to \mathrm{SubIdx}(s) \to \forall \varepsilon \in \mathbb{R}\,(\varphi_{20})$
\item $\varphi_{93} :\equiv (\mathrm{Conv}(a,\alpha) \to \mathrm{SubIdx}(s) \to \varepsilon \in \mathbb{R} \to 0 < \varepsilon \to \varphi_{47}) \to \varphi_{69}$
\item $\varphi_{94} :\equiv (\mathrm{Conv}(a,\alpha) \to \mathrm{SubIdx}(s) \to \varepsilon \in \mathbb{R} \to 0 < \varepsilon \to \varphi_{31}) \to \varphi_{76}$
\item $\varphi_{95} :\equiv (\mathrm{Conv}(a,\alpha) \to \mathrm{SubIdx}(s) \to \varepsilon \in \mathbb{R} \to 0 < \varepsilon \to \varphi_{45}) \to \varphi_{75}$
\item $\varphi_{96} :\equiv (\mathrm{Conv}(a,\alpha) \to \mathrm{SubIdx}(s) \to \varepsilon \in \mathbb{R} \to 0 < \varepsilon \to \varphi_{46}) \to \varphi_{77}$
\item $\varphi_{97} :\equiv (\mathrm{Conv}(a,\alpha) \to \mathrm{SubIdx}(s) \to \varepsilon \in \mathbb{R} \to 0 < \varepsilon \to \varphi_{48}) \to \varphi_{80}$
\item $\varphi_{98} :\equiv \forall \varepsilon\,(\mathrm{Conv}(a,\alpha) \land \mathrm{SubIdx}(s) \to \varepsilon \in \mathbb{R} \to \varphi_{20}) \to \varphi_{92}$
\item $\varphi_{99} :\equiv (\mathrm{Conv}(a,\alpha) \to \mathrm{SubIdx}(s) \to \varepsilon \in \mathbb{R} \to 0 < \varepsilon \to \varphi_{37}) \to \varphi_{91}$
\item $\varphi_{100} :\equiv (\mathrm{Conv}(a,\alpha) \to \mathrm{SubIdx}(s) \to \varepsilon \in \mathbb{R} \to 0 < \varepsilon \to \varphi_{55}) \to \varphi_{94}$
\item $\varphi_{101} :\equiv \forall n\,(\varphi_{34} \to n \in \mathbb{N} \to N \leq n \to \lvert a \circ s(n) - \alpha \rvert < \varepsilon) \to \varphi_{82} \to \varphi_{88}$
\item $\varphi_{102} :\equiv (\mathrm{Conv}(a,\alpha) \to \mathrm{SubIdx}(s) \to \varphi_{27}) \to \varphi_{59} \to (\varphi_{59} \to \forall N\,(\varphi_{43}) \to \varphi_{15}) \to \varphi_{87}$
\end{abbrevlist}
\end{fullproof}

% 自動生成 (prover/build.js): SubConv  ノード数 134
\begin{fullproof}{定理6.7}{$\forall a\,\forall b\,\forall \alpha\,(\mathrm{Conv}(a,\alpha) \to \mathrm{SubSeq}(b,a) \to \mathrm{Conv}(b,\alpha))$}{134}\phantomsection\label{tree:SubConv}
\noindent$\Pi_{1}$:
\begin{ptree}
\AXC{$\forall a\,\forall s\,\forall \alpha\,(\varphi_{4})$\lbl{[C18]}}
\AXC{$\forall a\,\forall s\,\forall \alpha\,(\varphi_{4}) \to \forall s\,\forall \alpha\,(\varphi_{4})$\lbl{A2}}
\BIC{$\forall s\,\forall \alpha\,(\varphi_{4})$}
\AXC{$\forall s\,\forall \alpha\,(\varphi_{4}) \to \forall \alpha\,(\varphi_{4})$\lbl{A2}}
\BIC{$\forall \alpha\,(\varphi_{4})$}
\AXC{$\forall \alpha\,(\varphi_{4}) \to \varphi_{4}$\lbl{A2}}
\BIC{$\varphi_{4}$}
\end{ptree}
\noindent$\Pi_{2}$:
\begin{ptree}
\AXC{$\mathrm{Conv}(a,\alpha) \to \mathrm{Conv}(a,\alpha)$\lbl{A1}}
\AXC{$\stackrel{\Pi_{1}}{\vdots}$}\noLine
\UIC{$\varphi_{4}$}
\AXC{$\varphi_{4} \to (\mathrm{Conv}(a,\alpha) \to \mathrm{Conv}(a,\alpha)) \to \varphi_{46}$\lbl{A1}}
\BIC{$(\mathrm{Conv}(a,\alpha) \to \mathrm{Conv}(a,\alpha)) \to \varphi_{46}$}
\BIC{$\varphi_{46}$}
\end{ptree}
\noindent$\Pi_{3}$:
\begin{ptree}
\AXC{$y = z \to x = y \to x = z$\lbl{A6}}
\AXC{$(y = z \to x = y \to x = z) \to x = y \to y = z \to x = z$\lbl{A1}}
\BIC{$x = y \to y = z \to x = z$}
\UIC{$\forall z\,(x = y \to y = z \to x = z)$}
\UIC{$\forall y\,\forall z\,(x = y \to y = z \to x = z)$}
\UIC{$\forall x\,\forall y\,\forall z\,(x = y \to y = z \to x = z)$}
\AXC{$\varphi_{22}$\lbl{A2}}
\BIC{$\forall y\,\forall z\,(a \circ s(n) = y \to y = z \to a \circ s(n) = z)$}
\end{ptree}
\noindent$\Pi_{4}$:
\begin{ptree}
\AXC{$\stackrel{\Pi_{3}}{\vdots}$}\noLine
\UIC{$\forall y\,\forall z\,(a \circ s(n) = y \to y = z \to a \circ s(n) = z)$}
\AXC{$\forall y\,\forall z\,(a \circ s(n) = y \to y = z \to a \circ s(n) = z) \to \varphi_{20}$\lbl{A2}}
\BIC{$\varphi_{20}$}
\AXC{$\varphi_{39}$\lbl{A2}}
\BIC{$a \circ s(n) = a_{s_{n}} \to a_{s_{n}} = b_{n} \to a \circ s(n) = b_{n}$}
\end{ptree}
\noindent$\Pi_{5}$:
\begin{ptree}
\AXC{$x = x$\lbl{A5}}
\AXC{$x = y \to x = x \to y = x$\lbl{A6}}
\AXC{$(x = y \to x = x \to y = x) \to x = x \to x = y \to y = x$\lbl{A1}}
\BIC{$x = x \to x = y \to y = x$}
\BIC{$x = y \to y = x$}
\UIC{$\forall y\,(x = y \to y = x)$}
\UIC{$\forall x\,\forall y\,(x = y \to y = x)$}
\end{ptree}
\noindent$\Pi_{6}$:
\begin{ptree}
\AXC{$\stackrel{\Pi_{5}}{\vdots}$}\noLine
\UIC{$\forall x\,\forall y\,(x = y \to y = x)$}
\AXC{$\forall x\,\forall y\,(x = y \to y = x) \to \forall y\,(b_{n} = y \to y = b_{n})$\lbl{A2}}
\BIC{$\forall y\,(b_{n} = y \to y = b_{n})$}
\end{ptree}
\noindent$\Pi_{7}$:
\begin{ptree}
\AXC{$\mathrm{Conv}(a,\alpha) \to \mathrm{SubSeq}(b,a) \to \varphi_{32}$\lbl{A1}}
\AXC{$\varphi_{57}$\lbl{A1}}
\BIC{$\mathrm{Conv}(a,\alpha) \to \mathrm{SubSeq}(b,a) \to \varphi_{33}$}
\AXC{$(\mathrm{Conv}(a,\alpha) \to \mathrm{SubSeq}(b,a) \to \varphi_{33}) \to \varphi_{49}$\lbl{A1}}
\BIC{$\varphi_{49}$}
\end{ptree}
\noindent$\Pi_{8}$:
\begin{ptree}
\AXC{$\forall k \in \mathbb{N}\,(b_{k} = a_{s_{k}}) \to n \in \mathbb{N} \to b_{n} = a_{s_{n}}$\lbl{A2}}
\AXC{$\stackrel{\Pi_{7}}{\vdots}$}\noLine
\UIC{$\varphi_{49}$}
\BIC{$\mathrm{Conv}(a,\alpha) \to \mathrm{SubSeq}(b,a) \to \varphi_{31}$}
\AXC{$(\mathrm{Conv}(a,\alpha) \to \mathrm{SubSeq}(b,a) \to \varphi_{31}) \to \varphi_{50}$\lbl{A1}}
\BIC{$\varphi_{50}$}
\end{ptree}
\noindent$\Pi_{9}$:
\begin{ptree}
\AXC{$\mathrm{Conv}(a,\alpha) \to \mathrm{SubSeq}(b,a) \to \varphi_{19}$\lbl{A1}}
\AXC{$\stackrel{\Pi_{8}}{\vdots}$}\noLine
\UIC{$\varphi_{50}$}
\BIC{$\mathrm{Conv}(a,\alpha) \to \mathrm{SubSeq}(b,a) \to \varphi_{24}$}
\AXC{$(\mathrm{Conv}(a,\alpha) \to \mathrm{SubSeq}(b,a) \to \varphi_{24}) \to \varphi_{42}$\lbl{A1}}
\BIC{$\varphi_{42}$}
\end{ptree}
\noindent$\Pi_{10}$:
\begin{ptree}
\AXC{$\stackrel{\Pi_{6}}{\vdots}$}\noLine
\UIC{$\forall y\,(b_{n} = y \to y = b_{n})$}
\AXC{$\forall y\,(b_{n} = y \to y = b_{n}) \to b_{n} = a_{s_{n}} \to a_{s_{n}} = b_{n}$\lbl{A2}}
\BIC{$b_{n} = a_{s_{n}} \to a_{s_{n}} = b_{n}$}
\AXC{$\stackrel{\Pi_{9}}{\vdots}$}\noLine
\UIC{$\varphi_{42}$}
\BIC{$\mathrm{Conv}(a,\alpha) \to \mathrm{SubSeq}(b,a) \to \varphi_{25}$}
\end{ptree}
\noindent$\Pi_{11}$:
\begin{ptree}
\AXC{$\forall s\,(\varphi_{8})$\lbl{DefSubIdx}}
\AXC{$\forall s\,(\varphi_{8}) \to \varphi_{8}$\lbl{A2}}
\BIC{$\varphi_{8}$}
\AXC{$\mathrm{Conv}(a,\alpha) \to \mathrm{SubSeq}(b,a) \to \varphi_{32}$\lbl{A1}}
\AXC{$\varphi_{51}$\lbl{A1}}
\BIC{$\mathrm{Conv}(a,\alpha) \to \mathrm{SubSeq}(b,a) \to \varphi_{12}$}
\AXC{$\varphi_{55}$\lbl{A1}}
\BIC{$\varphi_{8} \to \mathrm{Conv}(a,\alpha) \to \mathrm{SubSeq}(b,a) \to \varphi_{11}$}
\BIC{$\mathrm{Conv}(a,\alpha) \to \mathrm{SubSeq}(b,a) \to \varphi_{11}$}
\end{ptree}
\noindent$\Pi_{12}$:
\begin{ptree}
\AXC{$\forall a\,\forall \alpha\,(\mathrm{Conv}(a,\alpha) \to \mathrm{Seq}(a))$\lbl{[C1]}}
\AXC{$\varphi_{6}$\lbl{A2}}
\BIC{$\forall \alpha\,(\mathrm{Conv}(a,\alpha) \to \mathrm{Seq}(a))$}
\AXC{$\varphi_{5}$\lbl{A2}}
\BIC{$\mathrm{Conv}(a,\alpha) \to \mathrm{Seq}(a)$}
\AXC{$(\mathrm{Conv}(a,\alpha) \to \mathrm{Seq}(a)) \to \varphi_{35}$\lbl{A1}}
\BIC{$\varphi_{35}$}
\end{ptree}
\noindent$\Pi_{13}$:
\begin{ptree}
\AXC{$\forall a\,\forall s\,(\mathrm{Seq}(a) \to \varphi_{17})$\lbl{[S12]}}
\AXC{$\forall a\,\forall s\,(\mathrm{Seq}(a) \to \varphi_{17}) \to \forall s\,(\mathrm{Seq}(a) \to \varphi_{17})$\lbl{A2}}
\BIC{$\forall s\,(\mathrm{Seq}(a) \to \varphi_{17})$}
\end{ptree}
\noindent$\Pi_{14}$:
\begin{ptree}
\AXC{$\stackrel{\Pi_{13}}{\vdots}$}\noLine
\UIC{$\forall s\,(\mathrm{Seq}(a) \to \varphi_{17})$}
\AXC{$\forall s\,(\mathrm{Seq}(a) \to \varphi_{17}) \to \mathrm{Seq}(a) \to \varphi_{17}$\lbl{A2}}
\BIC{$\mathrm{Seq}(a) \to \varphi_{17}$}
\AXC{$\varphi_{62}$\lbl{A1}}
\BIC{$(\mathrm{Conv}(a,\alpha) \to \mathrm{SubSeq}(b,a) \to \varphi_{9}) \to \varphi_{53}$}
\end{ptree}
\noindent$\Pi_{15}$:
\begin{ptree}
\AXC{$\mathrm{Conv}(a,\alpha) \to \mathrm{Conv}(a,\alpha)$\lbl{A1}}
\AXC{$\stackrel{\Pi_{12}}{\vdots}$}\noLine
\UIC{$\varphi_{35}$}
\BIC{$\mathrm{Conv}(a,\alpha) \to \mathrm{SubSeq}(b,a) \to \varphi_{9}$}
\AXC{$\stackrel{\Pi_{14}}{\vdots}$}\noLine
\UIC{$(\mathrm{Conv}(a,\alpha) \to \mathrm{SubSeq}(b,a) \to \varphi_{9}) \to \varphi_{53}$}
\BIC{$\varphi_{53}$}
\end{ptree}
\noindent$\Pi_{16}$:
\begin{ptree}
\AXC{$\varphi_{27}$\lbl{A2}}
\AXC{$\stackrel{\Pi_{11}}{\vdots}$}\noLine
\UIC{$\mathrm{Conv}(a,\alpha) \to \mathrm{SubSeq}(b,a) \to \varphi_{11}$}
\AXC{$\stackrel{\Pi_{15}}{\vdots}$}\noLine
\UIC{$\varphi_{53}$}
\BIC{$\mathrm{Conv}(a,\alpha) \to \mathrm{SubSeq}(b,a) \to \varphi_{37}$}
\AXC{$\varphi_{59}$\lbl{A1}}
\BIC{$\mathrm{Conv}(a,\alpha) \to \mathrm{SubSeq}(b,a) \to \varphi_{36}$}
\AXC{$\varphi_{60}$\lbl{A1}}
\BIC{$\varphi_{27} \to \mathrm{Conv}(a,\alpha) \to \mathrm{SubSeq}(b,a) \to \varphi_{34}$}
\BIC{$\mathrm{Conv}(a,\alpha) \to \mathrm{SubSeq}(b,a) \to \varphi_{34}$}
\end{ptree}
\noindent$\Pi_{17}$:
\begin{ptree}
\AXC{$\stackrel{\Pi_{16}}{\vdots}$}\noLine
\UIC{$\mathrm{Conv}(a,\alpha) \to \mathrm{SubSeq}(b,a) \to \varphi_{34}$}
\AXC{$(\mathrm{Conv}(a,\alpha) \to \mathrm{SubSeq}(b,a) \to \varphi_{34}) \to \varphi_{52}$\lbl{A1}}
\BIC{$\varphi_{52}$}
\end{ptree}
\noindent$\Pi_{18}$:
\begin{ptree}
\AXC{$\stackrel{\Pi_{10}}{\vdots}$}\noLine
\UIC{$\mathrm{Conv}(a,\alpha) \to \mathrm{SubSeq}(b,a) \to \varphi_{25}$}
\AXC{$\mathrm{Conv}(a,\alpha) \to \mathrm{SubSeq}(b,a) \to \varphi_{19}$\lbl{A1}}
\AXC{$\stackrel{\Pi_{17}}{\vdots}$}\noLine
\UIC{$\varphi_{52}$}
\BIC{$\mathrm{Conv}(a,\alpha) \to \mathrm{SubSeq}(b,a) \to \varphi_{30}$}
\AXC{$\varphi_{63}$\lbl{A1}}
\BIC{$(\mathrm{Conv}(a,\alpha) \to \mathrm{SubSeq}(b,a) \to \varphi_{25}) \to \varphi_{48}$}
\BIC{$\varphi_{48}$}
\end{ptree}
\noindent$\Pi_{19}$:
\begin{ptree}
\AXC{$\stackrel{\Pi_{4}}{\vdots}$}\noLine
\UIC{$a \circ s(n) = a_{s_{n}} \to a_{s_{n}} = b_{n} \to a \circ s(n) = b_{n}$}
\AXC{$\stackrel{\Pi_{18}}{\vdots}$}\noLine
\UIC{$\varphi_{48}$}
\BIC{$\mathrm{Conv}(a,\alpha) \to \mathrm{SubSeq}(b,a) \to \varphi_{26}$}
\AXC{$\varphi_{54}$\lbl{A1}}
\BIC{$\varphi_{23} \to n \in \mathbb{N} \to a \circ s(n) = b_{n}$}
\UIC{$\forall n\,(\varphi_{23} \to n \in \mathbb{N} \to a \circ s(n) = b_{n})$}
\end{ptree}
\noindent$\Pi_{20}$:
\begin{ptree}
\AXC{$\varphi_{23} \to \forall n\,(\varphi_{23})$\lbl{A4}}
\AXC{$\varphi_{56}$\lbl{A3}}
\AXC{$\stackrel{\Pi_{19}}{\vdots}$}\noLine
\UIC{$\forall n\,(\varphi_{23} \to n \in \mathbb{N} \to a \circ s(n) = b_{n})$}
\AXC{$\forall n\,(\varphi_{23} \to n \in \mathbb{N} \to a \circ s(n) = b_{n}) \to \varphi_{56} \to \varphi_{58}$\lbl{A1}}
\BIC{$\varphi_{56} \to \varphi_{58}$}
\BIC{$\varphi_{58}$}
\BIC{$\mathrm{Conv}(a,\alpha) \to \mathrm{SubSeq}(b,a) \to \varphi_{28}$}
\end{ptree}
\noindent$\Pi_{21}$:
\begin{ptree}
\AXC{$\forall b\,\forall a\,(\mathrm{SubSeq}(b,a) \leftrightarrow \varphi_{13})$\lbl{DefSubSeq}}
\AXC{$\varphi_{40}$\lbl{A2}}
\BIC{$\forall a\,(\mathrm{SubSeq}(b,a) \leftrightarrow \varphi_{13})$}
\AXC{$\varphi_{38}$\lbl{A2}}
\BIC{$\mathrm{SubSeq}(b,a) \leftrightarrow \varphi_{13}$}
\AXC{$\varphi_{2}$\lbl{A1}}
\AXC{$\varphi_{2} \to \varphi_{41}$\lbl{A1}}
\BIC{$\varphi_{41}$}
\BIC{$\mathrm{Conv}(a,\alpha) \to \mathrm{SubSeq}(b,a) \to \varphi_{13}$}
\end{ptree}
\noindent$\Pi_{22}$:
\begin{ptree}
\AXC{$\varphi_{21}$\lbl{[C13]}}
\AXC{$\varphi_{21} \to \forall b\,\forall \alpha\,(\mathrm{Seq}(b) \to \varphi_{18})$\lbl{A2}}
\BIC{$\forall b\,\forall \alpha\,(\mathrm{Seq}(b) \to \varphi_{18})$}
\AXC{$\forall b\,\forall \alpha\,(\mathrm{Seq}(b) \to \varphi_{18}) \to \forall \alpha\,(\mathrm{Seq}(b) \to \varphi_{18})$\lbl{A2}}
\BIC{$\forall \alpha\,(\mathrm{Seq}(b) \to \varphi_{18})$}
\end{ptree}
\noindent$\Pi_{23}$:
\begin{ptree}
\AXC{$\stackrel{\Pi_{22}}{\vdots}$}\noLine
\UIC{$\forall \alpha\,(\mathrm{Seq}(b) \to \varphi_{18})$}
\AXC{$\forall \alpha\,(\mathrm{Seq}(b) \to \varphi_{18}) \to \mathrm{Seq}(b) \to \varphi_{18}$\lbl{A2}}
\BIC{$\mathrm{Seq}(b) \to \varphi_{18}$}
\AXC{$(\mathrm{Seq}(b) \to \varphi_{18}) \to \varphi_{61}$\lbl{A1}}
\BIC{$\varphi_{61}$}
\end{ptree}
\noindent$\Pi_{24}$:
\begin{ptree}
\AXC{$\stackrel{\Pi_{20}}{\vdots}$}\noLine
\UIC{$\mathrm{Conv}(a,\alpha) \to \mathrm{SubSeq}(b,a) \to \varphi_{28}$}
\AXC{$\stackrel{\Pi_{21}}{\vdots}$}\noLine
\UIC{$\mathrm{Conv}(a,\alpha) \to \mathrm{SubSeq}(b,a) \to \varphi_{13}$}
\AXC{$\varphi_{44}$\lbl{A1}}
\BIC{$\mathrm{Conv}(a,\alpha) \to \mathrm{SubSeq}(b,a) \to \varphi_{10}$}
\AXC{$\stackrel{\Pi_{23}}{\vdots}$}\noLine
\UIC{$\varphi_{61}$}
\BIC{$(\mathrm{Conv}(a,\alpha) \to \mathrm{SubSeq}(b,a) \to \varphi_{28}) \to \varphi_{47}$}
\BIC{$\varphi_{47}$}
\end{ptree}
\noindent$\Pi_{25}$:
\begin{ptree}
\AXC{$\mathrm{Conv}(a,\alpha) \to \mathrm{SubSeq}(b,a) \to \varphi_{32}$\lbl{A1}}
\AXC{$\varphi_{51}$\lbl{A1}}
\BIC{$\mathrm{Conv}(a,\alpha) \to \mathrm{SubSeq}(b,a) \to \varphi_{12}$}
\AXC{$\stackrel{\Pi_{2}}{\vdots}$}\noLine
\UIC{$\varphi_{46}$}
\BIC{$\mathrm{Conv}(a,\alpha) \to \mathrm{SubSeq}(b,a) \to \varphi_{16}$}
\AXC{$\stackrel{\Pi_{24}}{\vdots}$}\noLine
\UIC{$\varphi_{47}$}
\BIC{$\mathrm{Conv}(a,\alpha) \to \mathrm{SubSeq}(b,a) \to \varphi_{14}$}
\AXC{$\varphi_{45}$\lbl{A1}}
\BIC{$\varphi_{3} \to \neg (\mathrm{SubIdx}(s) \land \forall k \in \mathbb{N}\,(b_{k} = a_{s_{k}}))$}
\UIC{$\varphi_{29}$}
\end{ptree}
\noindent$\Pi_{26}$:
\begin{ptree}
\AXC{$\varphi_{29} \to \forall s\,(\varphi_{3}) \to \varphi_{7}$\lbl{A3}}
\AXC{$\stackrel{\Pi_{25}}{\vdots}$}\noLine
\UIC{$\varphi_{29}$}
\AXC{$\stackrel{\Pi_{21}}{\vdots}$}\noLine
\UIC{$\mathrm{Conv}(a,\alpha) \to \mathrm{SubSeq}(b,a) \to \varphi_{13}$}
\AXC{$\varphi_{43}$\lbl{A1}}
\BIC{$\mathrm{Conv}(a,\alpha) \to \varphi_{15}$}
\AXC{$\varphi_{64}$\lbl{A1}}
\BIC{$\varphi_{29} \to (\varphi_{29} \to \forall s\,(\varphi_{3}) \to \varphi_{7}) \to (\varphi_{3} \to \forall s\,(\varphi_{3})) \to \varphi_{1}$}
\BIC{$(\varphi_{29} \to \forall s\,(\varphi_{3}) \to \varphi_{7}) \to (\varphi_{3} \to \forall s\,(\varphi_{3})) \to \varphi_{1}$}
\BIC{$(\varphi_{3} \to \forall s\,(\varphi_{3})) \to \varphi_{1}$}
\end{ptree}
\noindent 主木:
\begin{ptree}
\AXC{$\varphi_{3} \to \forall s\,(\varphi_{3})$\lbl{A4}}
\AXC{$\stackrel{\Pi_{26}}{\vdots}$}\noLine
\UIC{$(\varphi_{3} \to \forall s\,(\varphi_{3})) \to \varphi_{1}$}
\BIC{$\varphi_{1}$}
\UIC{$\forall \alpha\,(\varphi_{1})$}
\UIC{$\forall b\,\forall \alpha\,(\varphi_{1})$}
\UIC{$\forall a\,\forall b\,\forall \alpha\,(\varphi_{1})$}
\end{ptree}
\begin{abbrevlist}
\item $\varphi_{1} :\equiv \mathrm{Conv}(a,\alpha) \to \mathrm{SubSeq}(b,a) \to \mathrm{Conv}(b,\alpha)$
\item $\varphi_{2} :\equiv \mathrm{Conv}(a,\alpha) \to \mathrm{SubSeq}(b,a) \to \mathrm{SubSeq}(b,a)$
\item $\varphi_{3} :\equiv \mathrm{Conv}(a,\alpha) \land \mathrm{SubSeq}(b,a) \land \neg \mathrm{Conv}(b,\alpha)$
\item $\varphi_{4} :\equiv \mathrm{Conv}(a,\alpha) \to \mathrm{SubIdx}(s) \to \mathrm{Conv}(a \circ s,\alpha)$
\item $\varphi_{5} :\equiv \forall \alpha\,(\mathrm{Conv}(a,\alpha) \to \mathrm{Seq}(a)) \to \mathrm{Conv}(a,\alpha) \to \mathrm{Seq}(a)$
\item $\varphi_{6} :\equiv \forall a\,\forall \alpha\,(\mathrm{Conv}(a,\alpha) \to \mathrm{Seq}(a)) \to \forall \alpha\,(\mathrm{Conv}(a,\alpha) \to \mathrm{Seq}(a))$
\item $\varphi_{7} :\equiv \forall s\,(\neg (\mathrm{SubIdx}(s) \land \forall k \in \mathbb{N}\,(b_{k} = a_{s_{k}})))$
\item $\varphi_{8} :\equiv \mathrm{SubIdx}(s) \leftrightarrow \mathrm{NSeq}(s) \land \forall k \in \mathbb{N}\,(s_{k} < s_{S(k)})$
\item $\varphi_{9} :\equiv \mathrm{SubIdx}(s) \land \forall k \in \mathbb{N}\,(b_{k} = a_{s_{k}}) \to \mathrm{Seq}(a)$
\item $\varphi_{10} :\equiv \mathrm{SubIdx}(s) \land \forall k \in \mathbb{N}\,(b_{k} = a_{s_{k}}) \to \mathrm{Seq}(b)$
\item $\varphi_{11} :\equiv \mathrm{SubIdx}(s) \land \forall k \in \mathbb{N}\,(b_{k} = a_{s_{k}}) \to \mathrm{NSeq}(s)$
\item $\varphi_{12} :\equiv \mathrm{SubIdx}(s) \land \forall k \in \mathbb{N}\,(b_{k} = a_{s_{k}}) \to \mathrm{SubIdx}(s)$
\item $\varphi_{13} :\equiv \mathrm{Seq}(b) \land \exists s\,(\mathrm{SubIdx}(s) \land \forall k \in \mathbb{N}\,(b_{k} = a_{s_{k}}))$
\item $\varphi_{14} :\equiv \mathrm{SubIdx}(s) \land \forall k \in \mathbb{N}\,(b_{k} = a_{s_{k}}) \to \mathrm{Conv}(b,\alpha)$
\item $\varphi_{15} :\equiv \mathrm{SubSeq}(b,a) \to \exists s\,(\mathrm{SubIdx}(s) \land \forall k \in \mathbb{N}\,(b_{k} = a_{s_{k}}))$
\item $\varphi_{16} :\equiv \mathrm{SubIdx}(s) \land \forall k \in \mathbb{N}\,(b_{k} = a_{s_{k}}) \to \mathrm{Conv}(a \circ s,\alpha)$
\item $\varphi_{17} :\equiv \mathrm{NSeq}(s) \to \mathrm{Seq}(a \circ s) \land \forall k \in \mathbb{N}\,(a \circ s(k) = a_{s_{k}})$
\item $\varphi_{18} :\equiv \forall n \in \mathbb{N}\,(a \circ s(n) = b_{n}) \to \mathrm{Conv}(a \circ s,\alpha) \to \mathrm{Conv}(b,\alpha)$
\item $\varphi_{19} :\equiv \mathrm{SubIdx}(s) \land \forall k \in \mathbb{N}\,(b_{k} = a_{s_{k}}) \to n \in \mathbb{N} \to n \in \mathbb{N}$
\item $\varphi_{20} :\equiv \forall z\,(a \circ s(n) = a_{s_{n}} \to a_{s_{n}} = z \to a \circ s(n) = z)$
\item $\varphi_{21} :\equiv \forall a\,\forall b\,\forall \alpha\,(\mathrm{Seq}(b) \to \forall n \in \mathbb{N}\,(a_{n} = b_{n}) \to \mathrm{Conv}(a,\alpha) \to \mathrm{Conv}(b,\alpha))$
\item $\varphi_{22} :\equiv \forall x\,\forall y\,\forall z\,(x = y \to y = z \to x = z) \to \forall y\,\forall z\,(a \circ s(n) = y \to y = z \to a \circ s(n) = z)$
\item $\varphi_{23} :\equiv \mathrm{Conv}(a,\alpha) \land \mathrm{SubSeq}(b,a) \land (\mathrm{SubIdx}(s) \land \forall k \in \mathbb{N}\,(b_{k} = a_{s_{k}}))$
\item $\varphi_{24} :\equiv \mathrm{SubIdx}(s) \land \forall k \in \mathbb{N}\,(b_{k} = a_{s_{k}}) \to n \in \mathbb{N} \to b_{n} = a_{s_{n}}$
\item $\varphi_{25} :\equiv \mathrm{SubIdx}(s) \land \forall k \in \mathbb{N}\,(b_{k} = a_{s_{k}}) \to n \in \mathbb{N} \to a_{s_{n}} = b_{n}$
\item $\varphi_{26} :\equiv \mathrm{SubIdx}(s) \land \forall k \in \mathbb{N}\,(b_{k} = a_{s_{k}}) \to n \in \mathbb{N} \to a \circ s(n) = b_{n}$
\item $\varphi_{27} :\equiv \forall k \in \mathbb{N}\,(a \circ s(k) = a_{s_{k}}) \to n \in \mathbb{N} \to a \circ s(n) = a_{s_{n}}$
\item $\varphi_{28} :\equiv \mathrm{SubIdx}(s) \land \forall k \in \mathbb{N}\,(b_{k} = a_{s_{k}}) \to \forall n \in \mathbb{N}\,(a \circ s(n) = b_{n})$
\item $\varphi_{29} :\equiv \forall s\,(\varphi_{3} \to \neg (\mathrm{SubIdx}(s) \land \forall k \in \mathbb{N}\,(b_{k} = a_{s_{k}})))$
\item $\varphi_{30} :\equiv \mathrm{SubIdx}(s) \land \forall k \in \mathbb{N}\,(b_{k} = a_{s_{k}}) \to n \in \mathbb{N} \to a \circ s(n) = a_{s_{n}}$
\item $\varphi_{31} :\equiv \mathrm{SubIdx}(s) \land \forall k \in \mathbb{N}\,(b_{k} = a_{s_{k}}) \to n \in \mathbb{N} \to n \in \mathbb{N} \to b_{n} = a_{s_{n}}$
\item $\varphi_{32} :\equiv \mathrm{SubIdx}(s) \land \forall k \in \mathbb{N}\,(b_{k} = a_{s_{k}}) \to \mathrm{SubIdx}(s) \land \forall k \in \mathbb{N}\,(b_{k} = a_{s_{k}})$
\item $\varphi_{33} :\equiv \mathrm{SubIdx}(s) \land \forall k \in \mathbb{N}\,(b_{k} = a_{s_{k}}) \to n \in \mathbb{N} \to \forall k \in \mathbb{N}\,(b_{k} = a_{s_{k}})$
\item $\varphi_{34} :\equiv \mathrm{SubIdx}(s) \land \forall k \in \mathbb{N}\,(b_{k} = a_{s_{k}}) \to n \in \mathbb{N} \to n \in \mathbb{N} \to a \circ s(n) = a_{s_{n}}$
\item $\varphi_{35} :\equiv (\mathrm{Conv}(a,\alpha) \to \mathrm{Conv}(a,\alpha)) \to \mathrm{Conv}(a,\alpha) \to \mathrm{SubSeq}(b,a) \to \varphi_{9}$
\item $\varphi_{36} :\equiv \mathrm{SubIdx}(s) \land \forall k \in \mathbb{N}\,(b_{k} = a_{s_{k}}) \to n \in \mathbb{N} \to \forall k \in \mathbb{N}\,(a \circ s(k) = a_{s_{k}})$
\item $\varphi_{37} :\equiv \mathrm{SubIdx}(s) \land \forall k \in \mathbb{N}\,(b_{k} = a_{s_{k}}) \to \mathrm{Seq}(a \circ s) \land \forall k \in \mathbb{N}\,(a \circ s(k) = a_{s_{k}})$
\item $\varphi_{38} :\equiv \forall a\,(\mathrm{SubSeq}(b,a) \leftrightarrow \varphi_{13}) \to (\mathrm{SubSeq}(b,a) \leftrightarrow \varphi_{13})$
\item $\varphi_{39} :\equiv \varphi_{20} \to a \circ s(n) = a_{s_{n}} \to a_{s_{n}} = b_{n} \to a \circ s(n) = b_{n}$
\item $\varphi_{40} :\equiv \forall b\,\forall a\,(\mathrm{SubSeq}(b,a) \leftrightarrow \varphi_{13}) \to \forall a\,(\mathrm{SubSeq}(b,a) \leftrightarrow \varphi_{13})$
\item $\varphi_{41} :\equiv (\mathrm{SubSeq}(b,a) \leftrightarrow \varphi_{13}) \to \mathrm{Conv}(a,\alpha) \to \mathrm{SubSeq}(b,a) \to \varphi_{13}$
\item $\varphi_{42} :\equiv (b_{n} = a_{s_{n}} \to a_{s_{n}} = b_{n}) \to \mathrm{Conv}(a,\alpha) \to \mathrm{SubSeq}(b,a) \to \varphi_{25}$
\item $\varphi_{43} :\equiv (\mathrm{Conv}(a,\alpha) \to \mathrm{SubSeq}(b,a) \to \varphi_{13}) \to \mathrm{Conv}(a,\alpha) \to \varphi_{15}$
\item $\varphi_{44} :\equiv (\mathrm{Conv}(a,\alpha) \to \mathrm{SubSeq}(b,a) \to \varphi_{13}) \to \mathrm{Conv}(a,\alpha) \to \mathrm{SubSeq}(b,a) \to \varphi_{10}$
\item $\varphi_{45} :\equiv (\mathrm{Conv}(a,\alpha) \to \mathrm{SubSeq}(b,a) \to \varphi_{14}) \to \varphi_{3} \to \neg (\mathrm{SubIdx}(s) \land \forall k \in \mathbb{N}\,(b_{k} = a_{s_{k}}))$
\item $\varphi_{46} :\equiv (\mathrm{Conv}(a,\alpha) \to \mathrm{SubSeq}(b,a) \to \varphi_{12}) \to \mathrm{Conv}(a,\alpha) \to \mathrm{SubSeq}(b,a) \to \varphi_{16}$
\item $\varphi_{47} :\equiv (\mathrm{Conv}(a,\alpha) \to \mathrm{SubSeq}(b,a) \to \varphi_{16}) \to \mathrm{Conv}(a,\alpha) \to \mathrm{SubSeq}(b,a) \to \varphi_{14}$
\item $\varphi_{48} :\equiv (a \circ s(n) = a_{s_{n}} \to a_{s_{n}} = b_{n} \to a \circ s(n) = b_{n}) \to \mathrm{Conv}(a,\alpha) \to \mathrm{SubSeq}(b,a) \to \varphi_{26}$
\item $\varphi_{49} :\equiv (\forall k \in \mathbb{N}\,(b_{k} = a_{s_{k}}) \to n \in \mathbb{N} \to b_{n} = a_{s_{n}}) \to \mathrm{Conv}(a,\alpha) \to \mathrm{SubSeq}(b,a) \to \varphi_{31}$
\item $\varphi_{50} :\equiv (\mathrm{Conv}(a,\alpha) \to \mathrm{SubSeq}(b,a) \to \varphi_{19}) \to \mathrm{Conv}(a,\alpha) \to \mathrm{SubSeq}(b,a) \to \varphi_{24}$
\item $\varphi_{51} :\equiv (\mathrm{Conv}(a,\alpha) \to \mathrm{SubSeq}(b,a) \to \varphi_{32}) \to \mathrm{Conv}(a,\alpha) \to \mathrm{SubSeq}(b,a) \to \varphi_{12}$
\item $\varphi_{52} :\equiv (\mathrm{Conv}(a,\alpha) \to \mathrm{SubSeq}(b,a) \to \varphi_{19}) \to \mathrm{Conv}(a,\alpha) \to \mathrm{SubSeq}(b,a) \to \varphi_{30}$
\item $\varphi_{53} :\equiv (\mathrm{Conv}(a,\alpha) \to \mathrm{SubSeq}(b,a) \to \varphi_{11}) \to \mathrm{Conv}(a,\alpha) \to \mathrm{SubSeq}(b,a) \to \varphi_{37}$
\item $\varphi_{54} :\equiv (\mathrm{Conv}(a,\alpha) \to \mathrm{SubSeq}(b,a) \to \varphi_{26}) \to \varphi_{23} \to n \in \mathbb{N} \to a \circ s(n) = b_{n}$
\item $\varphi_{55} :\equiv (\mathrm{Conv}(a,\alpha) \to \mathrm{SubSeq}(b,a) \to \varphi_{12}) \to \varphi_{8} \to \mathrm{Conv}(a,\alpha) \to \mathrm{SubSeq}(b,a) \to \varphi_{11}$
\item $\varphi_{56} :\equiv \forall n\,(\varphi_{23} \to n \in \mathbb{N} \to a \circ s(n) = b_{n}) \to \forall n\,(\varphi_{23}) \to \forall n \in \mathbb{N}\,(a \circ s(n) = b_{n})$
\item $\varphi_{57} :\equiv (\mathrm{Conv}(a,\alpha) \to \mathrm{SubSeq}(b,a) \to \varphi_{32}) \to \mathrm{Conv}(a,\alpha) \to \mathrm{SubSeq}(b,a) \to \varphi_{33}$
\item $\varphi_{58} :\equiv (\varphi_{23} \to \forall n\,(\varphi_{23})) \to \mathrm{Conv}(a,\alpha) \to \mathrm{SubSeq}(b,a) \to \varphi_{28}$
\item $\varphi_{59} :\equiv (\mathrm{Conv}(a,\alpha) \to \mathrm{SubSeq}(b,a) \to \varphi_{37}) \to \mathrm{Conv}(a,\alpha) \to \mathrm{SubSeq}(b,a) \to \varphi_{36}$
\item $\varphi_{60} :\equiv (\mathrm{Conv}(a,\alpha) \to \mathrm{SubSeq}(b,a) \to \varphi_{36}) \to \varphi_{27} \to \mathrm{Conv}(a,\alpha) \to \mathrm{SubSeq}(b,a) \to \varphi_{34}$
\item $\varphi_{61} :\equiv (\mathrm{Conv}(a,\alpha) \to \mathrm{SubSeq}(b,a) \to \varphi_{10}) \to (\mathrm{Conv}(a,\alpha) \to \mathrm{SubSeq}(b,a) \to \varphi_{28}) \to \varphi_{47}$
\item $\varphi_{62} :\equiv (\mathrm{Seq}(a) \to \varphi_{17}) \to (\mathrm{Conv}(a,\alpha) \to \mathrm{SubSeq}(b,a) \to \varphi_{9}) \to \varphi_{53}$
\item $\varphi_{63} :\equiv (\mathrm{Conv}(a,\alpha) \to \mathrm{SubSeq}(b,a) \to \varphi_{30}) \to (\mathrm{Conv}(a,\alpha) \to \mathrm{SubSeq}(b,a) \to \varphi_{25}) \to \varphi_{48}$
\item $\varphi_{64} :\equiv (\mathrm{Conv}(a,\alpha) \to \varphi_{15}) \to \varphi_{29} \to (\varphi_{29} \to \forall s\,(\varphi_{3}) \to \varphi_{7}) \to (\varphi_{3} \to \forall s\,(\varphi_{3})) \to \varphi_{1}$
\end{abbrevlist}
\end{fullproof}

% 自動生成 (prover/build.js): CoverConv  ノード数 497
\begin{fullproof}{[C19]}{$\forall a\,\forall s\,\forall t\,\forall \alpha\,(\mathrm{Seq}(a) \to \mathrm{SubIdx}(s) \to \mathrm{SubIdx}(t) \to \forall n \in \mathbb{N}\,\exists k \in \mathbb{N}\,(n = s_{k} \lor n = t(k)) \to \mathrm{Conv}(a \circ s,\alpha) \to \mathrm{Conv}(a \circ t,\alpha) \to \mathrm{Conv}(a,\alpha))$}{497}\phantomsection\label{tree:CoverConv}
\noindent$\Pi_{1}$:
\begin{ptree}
\AXC{$\forall m\,\forall n\,(n \leq m \cup n)$\lbl{[C8]}}
\AXC{$\forall m\,\forall n\,(n \leq m \cup n) \to \forall n\,(n \leq s_{N_{1}} \cup n)$\lbl{A2}}
\BIC{$\forall n\,(n \leq s_{N_{1}} \cup n)$}
\end{ptree}
\noindent$\Pi_{2}$:
\begin{ptree}
\AXC{$\stackrel{\Pi_{1}}{\vdots}$}\noLine
\UIC{$\forall n\,(n \leq s_{N_{1}} \cup n)$}
\AXC{$\forall n\,(n \leq s_{N_{1}} \cup n) \to t(N_{2}) \leq s_{N_{1}} \cup t(N_{2})$\lbl{A2}}
\BIC{$t(N_{2}) \leq s_{N_{1}} \cup t(N_{2})$}
\end{ptree}
\noindent$\Pi_{3}$:
\begin{ptree}
\AXC{$\forall x\,\forall y\,\forall z\,(x \leq y \to y \leq z \to x \leq z)$\lbl{[O5]}}
\AXC{$\varphi_{14}$\lbl{A2}}
\BIC{$\forall y\,\forall z\,(t(N_{2}) \leq y \to y \leq z \to t(N_{2}) \leq z)$}
\AXC{$\forall y\,\forall z\,(t(N_{2}) \leq y \to y \leq z \to t(N_{2}) \leq z) \to \varphi_{58}$\lbl{A2}}
\BIC{$\varphi_{58}$}
\AXC{$\varphi_{58} \to \varphi_{54}$\lbl{A2}}
\BIC{$\varphi_{54}$}
\end{ptree}
\noindent$\Pi_{4}$:
\begin{ptree}
\AXC{$\stackrel{\Pi_{2}}{\vdots}$}\noLine
\UIC{$t(N_{2}) \leq s_{N_{1}} \cup t(N_{2})$}
\AXC{$\stackrel{\Pi_{3}}{\vdots}$}\noLine
\UIC{$\varphi_{54}$}
\AXC{$\varphi_{54} \to t(N_{2}) \leq s_{N_{1}} \cup t(N_{2}) \to \varphi_{440}$\lbl{A1}}
\BIC{$t(N_{2}) \leq s_{N_{1}} \cup t(N_{2}) \to \varphi_{440}$}
\BIC{$\varphi_{440}$}
\end{ptree}
\noindent$\Pi_{5}$:
\begin{ptree}
\AXC{$\varphi_{1}$\lbl{A1}}
\AXC{$\forall s\,(\varphi_{37})$\lbl{[C16]}}
\AXC{$\forall s\,(\varphi_{37}) \to \varphi_{36}$\lbl{A2}}
\BIC{$\varphi_{36}$}
\AXC{$\varphi_{36} \to \varphi_{25}$\lbl{A2}}
\BIC{$\varphi_{25}$}
\AXC{$\varphi_{25} \to \varphi_{1} \to \varphi_{461}$\lbl{A1}}
\BIC{$\varphi_{1} \to \varphi_{461}$}
\BIC{$\varphi_{461}$}
\end{ptree}
\noindent$\Pi_{6}$:
\begin{ptree}
\AXC{$\varphi_{89}$\lbl{A2}}
\AXC{$\mathrm{Seq}(a) \to \mathrm{SubIdx}(s) \to \varphi_{386}$\lbl{A1}}
\AXC{$\varphi_{446}$\lbl{A1}}
\BIC{$\mathrm{Seq}(a) \to \mathrm{SubIdx}(s) \to \varphi_{358}$}
\AXC{$\stackrel{\Pi_{5}}{\vdots}$}\noLine
\UIC{$\varphi_{461}$}
\BIC{$\mathrm{Seq}(a) \to \mathrm{SubIdx}(s) \to \varphi_{405}$}
\AXC{$\varphi_{474}$\lbl{A1}}
\BIC{$\varphi_{89} \to \mathrm{Seq}(a) \to \mathrm{SubIdx}(s) \to \varphi_{406}$}
\BIC{$\mathrm{Seq}(a) \to \mathrm{SubIdx}(s) \to \varphi_{406}$}
\AXC{$\varphi_{487}$\lbl{A1}}
\BIC{$(\mathrm{Seq}(a) \to \mathrm{SubIdx}(s) \to \varphi_{286}) \to \varphi_{459}$}
\end{ptree}
\noindent$\Pi_{7}$:
\begin{ptree}
\AXC{$\mathrm{Seq}(a) \to \mathrm{SubIdx}(s) \to \varphi_{346}$\lbl{A1}}
\AXC{$\varphi_{431}$\lbl{A1}}
\BIC{$\mathrm{Seq}(a) \to \mathrm{SubIdx}(s) \to \varphi_{286}$}
\AXC{$\stackrel{\Pi_{6}}{\vdots}$}\noLine
\UIC{$(\mathrm{Seq}(a) \to \mathrm{SubIdx}(s) \to \varphi_{286}) \to \varphi_{459}$}
\BIC{$\varphi_{459}$}
\end{ptree}
\noindent$\Pi_{8}$:
\begin{ptree}
\AXC{$\forall x\,\forall y\,\forall z\,(x < y \to y \leq z \to x < z)$\lbl{[O3]}}
\AXC{$\varphi_{42}$\lbl{A2}}
\BIC{$\forall y\,\forall z\,(t(k) < y \to y \leq z \to t(k) < z)$}
\AXC{$\varphi_{110}$\lbl{A2}}
\BIC{$\forall z\,(t(k) < t(N_{2}) \to t(N_{2}) \leq z \to t(k) < z)$}
\AXC{$\varphi_{119}$\lbl{A2}}
\BIC{$t(k) < t(N_{2}) \to t(N_{2}) \leq n \to t(k) < n$}
\AXC{$\varphi_{483}$\lbl{A1}}
\BIC{$(\mathrm{Seq}(a) \to \mathrm{SubIdx}(s) \to \varphi_{392}) \to \varphi_{451}$}
\end{ptree}
\noindent$\Pi_{9}$:
\begin{ptree}
\AXC{$\mathrm{Seq}(a) \to \mathrm{SubIdx}(s) \to \varphi_{383}$\lbl{A1}}
\AXC{$\stackrel{\Pi_{7}}{\vdots}$}\noLine
\UIC{$\varphi_{459}$}
\BIC{$\mathrm{Seq}(a) \to \mathrm{SubIdx}(s) \to \varphi_{392}$}
\AXC{$\stackrel{\Pi_{8}}{\vdots}$}\noLine
\UIC{$(\mathrm{Seq}(a) \to \mathrm{SubIdx}(s) \to \varphi_{392}) \to \varphi_{451}$}
\BIC{$\varphi_{451}$}
\end{ptree}
\noindent$\Pi_{10}$:
\begin{ptree}
\AXC{$x = x$\lbl{A5}}
\AXC{$x = y \to x = x \to y = x$\lbl{A6}}
\AXC{$(x = y \to x = x \to y = x) \to x = x \to x = y \to y = x$\lbl{A1}}
\BIC{$x = x \to x = y \to y = x$}
\BIC{$x = y \to y = x$}
\UIC{$\forall y\,(x = y \to y = x)$}
\UIC{$\forall x\,\forall y\,(x = y \to y = x)$}
\end{ptree}
\noindent$\Pi_{11}$:
\begin{ptree}
\AXC{$\stackrel{\Pi_{10}}{\vdots}$}\noLine
\UIC{$\forall x\,\forall y\,(x = y \to y = x)$}
\AXC{$\forall x\,\forall y\,(x = y \to y = x) \to \forall y\,(n = y \to y = n)$\lbl{A2}}
\BIC{$\forall y\,(n = y \to y = n)$}
\end{ptree}
\noindent$\Pi_{12}$:
\begin{ptree}
\AXC{$\mathrm{Seq}(a) \to \mathrm{SubIdx}(s) \to \varphi_{365}$\lbl{A1}}
\AXC{$(\mathrm{Seq}(a) \to \mathrm{SubIdx}(s) \to \varphi_{365}) \to \varphi_{409}$\lbl{A1}}
\BIC{$\varphi_{409}$}
\end{ptree}
\noindent$\Pi_{13}$:
\begin{ptree}
\AXC{$\stackrel{\Pi_{11}}{\vdots}$}\noLine
\UIC{$\forall y\,(n = y \to y = n)$}
\AXC{$\forall y\,(n = y \to y = n) \to n = t(k) \to t(k) = n$\lbl{A2}}
\BIC{$n = t(k) \to t(k) = n$}
\AXC{$\stackrel{\Pi_{12}}{\vdots}$}\noLine
\UIC{$\varphi_{409}$}
\BIC{$\mathrm{Seq}(a) \to \mathrm{SubIdx}(s) \to \varphi_{366}$}
\AXC{$\varphi_{481}$\lbl{A1}}
\BIC{$(\mathrm{Seq}(a) \to \mathrm{SubIdx}(s) \to \varphi_{388}) \to \varphi_{415}$}
\end{ptree}
\noindent$\Pi_{14}$:
\begin{ptree}
\AXC{$\mathrm{Seq}(a) \to \mathrm{SubIdx}(s) \to \varphi_{344}$\lbl{A1}}
\AXC{$\stackrel{\Pi_{4}}{\vdots}$}\noLine
\UIC{$\varphi_{440}$}
\BIC{$\mathrm{Seq}(a) \to \mathrm{SubIdx}(s) \to \varphi_{369}$}
\AXC{$\stackrel{\Pi_{9}}{\vdots}$}\noLine
\UIC{$\varphi_{451}$}
\BIC{$\mathrm{Seq}(a) \to \mathrm{SubIdx}(s) \to \varphi_{388}$}
\AXC{$\stackrel{\Pi_{13}}{\vdots}$}\noLine
\UIC{$(\mathrm{Seq}(a) \to \mathrm{SubIdx}(s) \to \varphi_{388}) \to \varphi_{415}$}
\BIC{$\varphi_{415}$}
\end{ptree}
\noindent$\Pi_{15}$:
\begin{ptree}
\AXC{$z = w \to z < n \to w < n$\lbl{A6}}
\UIC{$\forall w\,(z = w \to z < n \to w < n)$}
\UIC{$\forall z\,\forall w\,(z = w \to z < n \to w < n)$}
\AXC{$\varphi_{24}$\lbl{A2}}
\BIC{$\forall w\,(t(k) = w \to t(k) < n \to w < n)$}
\AXC{$\varphi_{48}$\lbl{A2}}
\BIC{$t(k) = n \to t(k) < n \to n < n$}
\AXC{$\stackrel{\Pi_{14}}{\vdots}$}\noLine
\UIC{$\varphi_{415}$}
\BIC{$\mathrm{Seq}(a) \to \mathrm{SubIdx}(s) \to \varphi_{381}$}
\AXC{$\varphi_{466}$\lbl{A1}}
\BIC{$\neg n < n \to \mathrm{Seq}(a) \to \mathrm{SubIdx}(s) \to \varphi_{397}$}
\end{ptree}
\noindent$\Pi_{16}$:
\begin{ptree}
\AXC{$\forall x\,(\neg x < x)$\lbl{[O1]}}
\AXC{$\forall x\,(\neg x < x) \to \neg n < n$\lbl{A2}}
\BIC{$\neg n < n$}
\AXC{$\stackrel{\Pi_{15}}{\vdots}$}\noLine
\UIC{$\neg n < n \to \mathrm{Seq}(a) \to \mathrm{SubIdx}(s) \to \varphi_{397}$}
\BIC{$\mathrm{Seq}(a) \to \mathrm{SubIdx}(s) \to \varphi_{397}$}
\end{ptree}
\noindent$\Pi_{17}$:
\begin{ptree}
\AXC{$z = w \to \lvert z - \alpha \rvert < \varepsilon \to \lvert w - \alpha \rvert < \varepsilon$\lbl{A6}}
\UIC{$\forall w\,(z = w \to \lvert z - \alpha \rvert < \varepsilon \to \lvert w - \alpha \rvert < \varepsilon)$}
\UIC{$\forall z\,\forall w\,(z = w \to \lvert z - \alpha \rvert < \varepsilon \to \lvert w - \alpha \rvert < \varepsilon)$}
\AXC{$\forall z\,\forall w\,(z = w \to \lvert z - \alpha \rvert < \varepsilon \to \lvert w - \alpha \rvert < \varepsilon) \to \varphi_{59}$\lbl{A2}}
\BIC{$\varphi_{59}$}
\AXC{$\varphi_{59} \to \varphi_{113}$\lbl{A2}}
\BIC{$\varphi_{113}$}
\end{ptree}
\noindent$\Pi_{18}$:
\begin{ptree}
\AXC{$\varphi_{124}$\lbl{A2}}
\AXC{$\mathrm{Seq}(a) \to \mathrm{SubIdx}(s) \to \varphi_{346}$\lbl{A1}}
\AXC{$\varphi_{457}$\lbl{A1}}
\BIC{$\mathrm{Seq}(a) \to \mathrm{SubIdx}(s) \to \varphi_{412}$}
\AXC{$\varphi_{476}$\lbl{A1}}
\BIC{$\varphi_{124} \to \mathrm{Seq}(a) \to \mathrm{SubIdx}(s) \to \varphi_{410}$}
\BIC{$\mathrm{Seq}(a) \to \mathrm{SubIdx}(s) \to \varphi_{410}$}
\AXC{$\varphi_{490}$\lbl{A1}}
\BIC{$(\mathrm{Seq}(a) \to \mathrm{SubIdx}(s) \to \varphi_{358}) \to \varphi_{463}$}
\end{ptree}
\noindent$\Pi_{19}$:
\begin{ptree}
\AXC{$\mathrm{Seq}(a) \to \mathrm{SubIdx}(s) \to \varphi_{386}$\lbl{A1}}
\AXC{$\varphi_{446}$\lbl{A1}}
\BIC{$\mathrm{Seq}(a) \to \mathrm{SubIdx}(s) \to \varphi_{358}$}
\AXC{$\stackrel{\Pi_{18}}{\vdots}$}\noLine
\UIC{$(\mathrm{Seq}(a) \to \mathrm{SubIdx}(s) \to \varphi_{358}) \to \varphi_{463}$}
\BIC{$\varphi_{463}$}
\end{ptree}
\noindent$\Pi_{20}$:
\begin{ptree}
\AXC{$\forall s\,(\varphi_{6})$\lbl{DefSubIdx}}
\AXC{$\forall s\,(\varphi_{6}) \to \varphi_{5}$\lbl{A2}}
\BIC{$\varphi_{5}$}
\AXC{$\varphi_{1}$\lbl{A1}}
\AXC{$\varphi_{1} \to \varphi_{5} \to \mathrm{Seq}(a) \to \mathrm{SubIdx}(s) \to \varphi_{126}$\lbl{A1}}
\BIC{$\varphi_{5} \to \mathrm{Seq}(a) \to \mathrm{SubIdx}(s) \to \varphi_{126}$}
\BIC{$\mathrm{Seq}(a) \to \mathrm{SubIdx}(s) \to \varphi_{126}$}
\end{ptree}
\noindent$\Pi_{21}$:
\begin{ptree}
\AXC{$\forall a\,\forall s\,(\mathrm{Seq}(a) \to \varphi_{47})$\lbl{[S12]}}
\AXC{$\forall a\,\forall s\,(\mathrm{Seq}(a) \to \varphi_{47}) \to \forall s\,(\mathrm{Seq}(a) \to \varphi_{47})$\lbl{A2}}
\BIC{$\forall s\,(\mathrm{Seq}(a) \to \varphi_{47})$}
\end{ptree}
\noindent$\Pi_{22}$:
\begin{ptree}
\AXC{$\stackrel{\Pi_{21}}{\vdots}$}\noLine
\UIC{$\forall s\,(\mathrm{Seq}(a) \to \varphi_{47})$}
\AXC{$\forall s\,(\mathrm{Seq}(a) \to \varphi_{47}) \to \mathrm{Seq}(a) \to \varphi_{45}$\lbl{A2}}
\BIC{$\mathrm{Seq}(a) \to \varphi_{45}$}
\AXC{$(\mathrm{Seq}(a) \to \varphi_{45}) \to (\mathrm{Seq}(a) \to \mathrm{Seq}(a)) \to \varphi_{331}$\lbl{A1}}
\BIC{$(\mathrm{Seq}(a) \to \mathrm{Seq}(a)) \to \varphi_{331}$}
\end{ptree}
\noindent$\Pi_{23}$:
\begin{ptree}
\AXC{$\stackrel{\Pi_{20}}{\vdots}$}\noLine
\UIC{$\mathrm{Seq}(a) \to \mathrm{SubIdx}(s) \to \varphi_{126}$}
\AXC{$\mathrm{Seq}(a) \to \mathrm{Seq}(a)$\lbl{A1}}
\AXC{$\stackrel{\Pi_{22}}{\vdots}$}\noLine
\UIC{$(\mathrm{Seq}(a) \to \mathrm{Seq}(a)) \to \varphi_{331}$}
\BIC{$\varphi_{331}$}
\BIC{$\mathrm{Seq}(a) \to \mathrm{SubIdx}(s) \to \mathrm{SubIdx}(t) \to \varphi_{162}$}
\AXC{$\varphi_{362}$\lbl{A1}}
\BIC{$\mathrm{Seq}(a) \to \mathrm{SubIdx}(s) \to \varphi_{157}$}
\end{ptree}
\noindent$\Pi_{24}$:
\begin{ptree}
\AXC{$\varphi_{105}$\lbl{A2}}
\AXC{$\stackrel{\Pi_{23}}{\vdots}$}\noLine
\UIC{$\mathrm{Seq}(a) \to \mathrm{SubIdx}(s) \to \varphi_{157}$}
\AXC{$\varphi_{418}$\lbl{A1}}
\BIC{$\varphi_{105} \to \mathrm{Seq}(a) \to \mathrm{SubIdx}(s) \to \varphi_{152}$}
\BIC{$\mathrm{Seq}(a) \to \mathrm{SubIdx}(s) \to \varphi_{152}$}
\AXC{$(\mathrm{Seq}(a) \to \mathrm{SubIdx}(s) \to \varphi_{152}) \to \varphi_{453}$\lbl{A1}}
\BIC{$\varphi_{453}$}
\end{ptree}
\noindent$\Pi_{25}$:
\begin{ptree}
\AXC{$\mathrm{Seq}(a) \to \mathrm{SubIdx}(s) \to \varphi_{386}$\lbl{A1}}
\AXC{$\varphi_{446}$\lbl{A1}}
\BIC{$\mathrm{Seq}(a) \to \mathrm{SubIdx}(s) \to \varphi_{358}$}
\AXC{$\stackrel{\Pi_{24}}{\vdots}$}\noLine
\UIC{$\varphi_{453}$}
\BIC{$\mathrm{Seq}(a) \to \mathrm{SubIdx}(s) \to \varphi_{399}$}
\AXC{$(\mathrm{Seq}(a) \to \mathrm{SubIdx}(s) \to \varphi_{399}) \to \varphi_{472}$\lbl{A1}}
\BIC{$\varphi_{472}$}
\end{ptree}
\noindent$\Pi_{26}$:
\begin{ptree}
\AXC{$\stackrel{\Pi_{10}}{\vdots}$}\noLine
\UIC{$\forall x\,\forall y\,(x = y \to y = x)$}
\AXC{$\forall x\,\forall y\,(x = y \to y = x) \to \forall y\,(n = y \to y = n)$\lbl{A2}}
\BIC{$\forall y\,(n = y \to y = n)$}
\end{ptree}
\noindent$\Pi_{27}$:
\begin{ptree}
\AXC{$\mathrm{Seq}(a) \to \mathrm{SubIdx}(s) \to \varphi_{365}$\lbl{A1}}
\AXC{$(\mathrm{Seq}(a) \to \mathrm{SubIdx}(s) \to \varphi_{365}) \to \varphi_{409}$\lbl{A1}}
\BIC{$\varphi_{409}$}
\end{ptree}
\noindent$\Pi_{28}$:
\begin{ptree}
\AXC{$\stackrel{\Pi_{26}}{\vdots}$}\noLine
\UIC{$\forall y\,(n = y \to y = n)$}
\AXC{$\forall y\,(n = y \to y = n) \to n = t(k) \to t(k) = n$\lbl{A2}}
\BIC{$n = t(k) \to t(k) = n$}
\AXC{$\stackrel{\Pi_{27}}{\vdots}$}\noLine
\UIC{$\varphi_{409}$}
\BIC{$\mathrm{Seq}(a) \to \mathrm{SubIdx}(s) \to \varphi_{366}$}
\AXC{$\varphi_{485}$\lbl{A1}}
\BIC{$(\mathrm{Seq}(a) \to \mathrm{SubIdx}(s) \to \varphi_{401}) \to \varphi_{422}$}
\end{ptree}
\noindent$\Pi_{29}$:
\begin{ptree}
\AXC{$\stackrel{\Pi_{17}}{\vdots}$}\noLine
\UIC{$\varphi_{113}$}
\AXC{$\mathrm{Seq}(a) \to \mathrm{SubIdx}(s) \to \varphi_{370}$\lbl{A1}}
\AXC{$\stackrel{\Pi_{19}}{\vdots}$}\noLine
\UIC{$\varphi_{463}$}
\BIC{$\mathrm{Seq}(a) \to \mathrm{SubIdx}(s) \to \varphi_{403}$}
\AXC{$\stackrel{\Pi_{25}}{\vdots}$}\noLine
\UIC{$\varphi_{472}$}
\BIC{$\varphi_{113} \to \mathrm{Seq}(a) \to \mathrm{SubIdx}(s) \to \varphi_{401}$}
\BIC{$\mathrm{Seq}(a) \to \mathrm{SubIdx}(s) \to \varphi_{401}$}
\AXC{$\stackrel{\Pi_{28}}{\vdots}$}\noLine
\UIC{$(\mathrm{Seq}(a) \to \mathrm{SubIdx}(s) \to \varphi_{401}) \to \varphi_{422}$}
\BIC{$\varphi_{422}$}
\end{ptree}
\noindent$\Pi_{30}$:
\begin{ptree}
\AXC{$z = w \to \lvert a_{z} - \alpha \rvert < \varepsilon \to \lvert a_{w} - \alpha \rvert < \varepsilon$\lbl{A6}}
\UIC{$\forall w\,(z = w \to \lvert a_{z} - \alpha \rvert < \varepsilon \to \lvert a_{w} - \alpha \rvert < \varepsilon)$}
\UIC{$\forall z\,\forall w\,(z = w \to \lvert a_{z} - \alpha \rvert < \varepsilon \to \lvert a_{w} - \alpha \rvert < \varepsilon)$}
\AXC{$\varphi_{148}$\lbl{A2}}
\BIC{$\forall w\,(t(k) = w \to \lvert a_{t(k)} - \alpha \rvert < \varepsilon \to \lvert a_{w} - \alpha \rvert < \varepsilon)$}
\AXC{$\varphi_{155}$\lbl{A2}}
\BIC{$t(k) = n \to \lvert a_{t(k)} - \alpha \rvert < \varepsilon \to \lvert a_{n} - \alpha \rvert < \varepsilon$}
\AXC{$\stackrel{\Pi_{29}}{\vdots}$}\noLine
\UIC{$\varphi_{422}$}
\BIC{$\mathrm{Seq}(a) \to \mathrm{SubIdx}(s) \to \varphi_{394}$}
\end{ptree}
\noindent$\Pi_{31}$:
\begin{ptree}
\AXC{$\forall n \in \mathbb{N}\,(n \in \omega)$\lbl{[N7']}}
\AXC{$\forall n \in \mathbb{N}\,(n \in \omega) \to k \in \mathbb{N} \to k \in \omega$\lbl{A2}}
\BIC{$k \in \mathbb{N} \to k \in \omega$}
\AXC{$(k \in \mathbb{N} \to k \in \omega) \to \varphi_{444}$\lbl{A1}}
\BIC{$\varphi_{444}$}
\end{ptree}
\noindent$\Pi_{32}$:
\begin{ptree}
\AXC{$\forall n \in \mathbb{N}\,(n \in \omega)$\lbl{[N7']}}
\AXC{$\forall n \in \mathbb{N}\,(n \in \omega) \to N_{2} \in \mathbb{N} \to N_{2} \in \omega$\lbl{A2}}
\BIC{$N_{2} \in \mathbb{N} \to N_{2} \in \omega$}
\AXC{$(N_{2} \in \mathbb{N} \to N_{2} \in \omega) \to \varphi_{434}$\lbl{A1}}
\BIC{$\varphi_{434}$}
\end{ptree}
\noindent$\Pi_{33}$:
\begin{ptree}
\AXC{$\forall m\,\forall n\,(m \in \omega \to n \in \omega \to m \leq n \lor n < m)$\lbl{[N11]}}
\AXC{$\varphi_{107}$\lbl{A2}}
\BIC{$\forall n\,(N_{2} \in \omega \to n \in \omega \to N_{2} \leq n \lor n < N_{2})$}
\AXC{$\varphi_{96}$\lbl{A2}}
\BIC{$N_{2} \in \omega \to k \in \omega \to N_{2} \leq k \lor k < N_{2}$}
\AXC{$\varphi_{479}$\lbl{A1}}
\BIC{$(\mathrm{Seq}(a) \to \mathrm{SubIdx}(s) \to \varphi_{375}) \to \varphi_{447}$}
\end{ptree}
\noindent$\Pi_{34}$:
\begin{ptree}
\AXC{$\mathrm{Seq}(a) \to \mathrm{SubIdx}(s) \to \varphi_{346}$\lbl{A1}}
\AXC{$\varphi_{431}$\lbl{A1}}
\BIC{$\mathrm{Seq}(a) \to \mathrm{SubIdx}(s) \to \varphi_{286}$}
\AXC{$\stackrel{\Pi_{32}}{\vdots}$}\noLine
\UIC{$\varphi_{434}$}
\BIC{$\mathrm{Seq}(a) \to \mathrm{SubIdx}(s) \to \varphi_{375}$}
\AXC{$\stackrel{\Pi_{33}}{\vdots}$}\noLine
\UIC{$(\mathrm{Seq}(a) \to \mathrm{SubIdx}(s) \to \varphi_{375}) \to \varphi_{447}$}
\BIC{$\varphi_{447}$}
\end{ptree}
\noindent$\Pi_{35}$:
\begin{ptree}
\AXC{$\mathrm{Seq}(a) \to \mathrm{SubIdx}(s) \to \varphi_{386}$\lbl{A1}}
\AXC{$\varphi_{446}$\lbl{A1}}
\BIC{$\mathrm{Seq}(a) \to \mathrm{SubIdx}(s) \to \varphi_{358}$}
\AXC{$\stackrel{\Pi_{31}}{\vdots}$}\noLine
\UIC{$\varphi_{444}$}
\BIC{$\mathrm{Seq}(a) \to \mathrm{SubIdx}(s) \to \varphi_{371}$}
\AXC{$\stackrel{\Pi_{34}}{\vdots}$}\noLine
\UIC{$\varphi_{447}$}
\BIC{$\mathrm{Seq}(a) \to \mathrm{SubIdx}(s) \to \varphi_{379}$}
\AXC{$\varphi_{492}$\lbl{A1}}
\BIC{$(\mathrm{Seq}(a) \to \mathrm{SubIdx}(s) \to \varphi_{394}) \to \varphi_{467}$}
\end{ptree}
\noindent$\Pi_{36}$:
\begin{ptree}
\AXC{$\stackrel{\Pi_{16}}{\vdots}$}\noLine
\UIC{$\mathrm{Seq}(a) \to \mathrm{SubIdx}(s) \to \varphi_{397}$}
\AXC{$\stackrel{\Pi_{30}}{\vdots}$}\noLine
\UIC{$\mathrm{Seq}(a) \to \mathrm{SubIdx}(s) \to \varphi_{394}$}
\AXC{$\stackrel{\Pi_{35}}{\vdots}$}\noLine
\UIC{$(\mathrm{Seq}(a) \to \mathrm{SubIdx}(s) \to \varphi_{394}) \to \varphi_{467}$}
\BIC{$\varphi_{467}$}
\BIC{$\mathrm{Seq}(a) \to \mathrm{SubIdx}(s) \to \varphi_{390}$}
\end{ptree}
\noindent$\Pi_{37}$:
\begin{ptree}
\AXC{$\forall m\,\forall n\,(m \leq m \cup n)$\lbl{[C7]}}
\AXC{$\forall m\,\forall n\,(m \leq m \cup n) \to \forall n\,(s_{N_{1}} \leq s_{N_{1}} \cup n)$\lbl{A2}}
\BIC{$\forall n\,(s_{N_{1}} \leq s_{N_{1}} \cup n)$}
\AXC{$\varphi_{18}$\lbl{A2}}
\BIC{$s_{N_{1}} \leq s_{N_{1}} \cup t(N_{2})$}
\end{ptree}
\noindent$\Pi_{38}$:
\begin{ptree}
\AXC{$\forall x\,\forall y\,\forall z\,(x \leq y \to y \leq z \to x \leq z)$\lbl{[O5]}}
\AXC{$\varphi_{15}$\lbl{A2}}
\BIC{$\forall y\,\forall z\,(s_{N_{1}} \leq y \to y \leq z \to s_{N_{1}} \leq z)$}
\AXC{$\forall y\,\forall z\,(s_{N_{1}} \leq y \to y \leq z \to s_{N_{1}} \leq z) \to \varphi_{60}$\lbl{A2}}
\BIC{$\varphi_{60}$}
\AXC{$\varphi_{60} \to \varphi_{55}$\lbl{A2}}
\BIC{$\varphi_{55}$}
\end{ptree}
\noindent$\Pi_{39}$:
\begin{ptree}
\AXC{$\mathrm{Seq}(a) \to \mathrm{SubIdx}(s) \to \varphi_{344}$\lbl{A1}}
\AXC{$\stackrel{\Pi_{37}}{\vdots}$}\noLine
\UIC{$s_{N_{1}} \leq s_{N_{1}} \cup t(N_{2})$}
\AXC{$\stackrel{\Pi_{38}}{\vdots}$}\noLine
\UIC{$\varphi_{55}$}
\AXC{$\varphi_{55} \to s_{N_{1}} \leq s_{N_{1}} \cup t(N_{2}) \to \varphi_{441}$\lbl{A1}}
\BIC{$s_{N_{1}} \leq s_{N_{1}} \cup t(N_{2}) \to \varphi_{441}$}
\BIC{$\varphi_{441}$}
\BIC{$\mathrm{Seq}(a) \to \mathrm{SubIdx}(s) \to \varphi_{372}$}
\end{ptree}
\noindent$\Pi_{40}$:
\begin{ptree}
\AXC{$\mathrm{Seq}(a) \to \mathrm{SubIdx}(s) \to \mathrm{SubIdx}(s)$\lbl{A1}}
\AXC{$\forall s\,(\varphi_{37})$\lbl{[C16]}}
\AXC{$\forall s\,(\varphi_{37}) \to \varphi_{37}$\lbl{A2}}
\BIC{$\varphi_{37}$}
\AXC{$\varphi_{37} \to \varphi_{27}$\lbl{A2}}
\BIC{$\varphi_{27}$}
\AXC{$\varphi_{470}$\lbl{A1}}
\BIC{$(\mathrm{Seq}(a) \to \mathrm{SubIdx}(s) \to \mathrm{SubIdx}(s)) \to \varphi_{462}$}
\BIC{$\varphi_{462}$}
\end{ptree}
\noindent$\Pi_{41}$:
\begin{ptree}
\AXC{$\varphi_{90}$\lbl{A2}}
\AXC{$\mathrm{Seq}(a) \to \mathrm{SubIdx}(s) \to \varphi_{386}$\lbl{A1}}
\AXC{$\varphi_{446}$\lbl{A1}}
\BIC{$\mathrm{Seq}(a) \to \mathrm{SubIdx}(s) \to \varphi_{358}$}
\AXC{$\stackrel{\Pi_{40}}{\vdots}$}\noLine
\UIC{$\varphi_{462}$}
\BIC{$\mathrm{Seq}(a) \to \mathrm{SubIdx}(s) \to \varphi_{407}$}
\AXC{$\varphi_{475}$\lbl{A1}}
\BIC{$\varphi_{90} \to \mathrm{Seq}(a) \to \mathrm{SubIdx}(s) \to \varphi_{408}$}
\BIC{$\mathrm{Seq}(a) \to \mathrm{SubIdx}(s) \to \varphi_{408}$}
\AXC{$\varphi_{488}$\lbl{A1}}
\BIC{$(\mathrm{Seq}(a) \to \mathrm{SubIdx}(s) \to \varphi_{289}) \to \varphi_{460}$}
\end{ptree}
\noindent$\Pi_{42}$:
\begin{ptree}
\AXC{$\mathrm{Seq}(a) \to \mathrm{SubIdx}(s) \to \varphi_{279}$\lbl{A1}}
\AXC{$\varphi_{425}$\lbl{A1}}
\BIC{$\mathrm{Seq}(a) \to \mathrm{SubIdx}(s) \to \varphi_{289}$}
\AXC{$\stackrel{\Pi_{41}}{\vdots}$}\noLine
\UIC{$(\mathrm{Seq}(a) \to \mathrm{SubIdx}(s) \to \varphi_{289}) \to \varphi_{460}$}
\BIC{$\varphi_{460}$}
\end{ptree}
\noindent$\Pi_{43}$:
\begin{ptree}
\AXC{$\forall x\,\forall y\,\forall z\,(x < y \to y \leq z \to x < z)$\lbl{[O3]}}
\AXC{$\varphi_{43}$\lbl{A2}}
\BIC{$\forall y\,\forall z\,(s_{k} < y \to y \leq z \to s_{k} < z)$}
\AXC{$\varphi_{111}$\lbl{A2}}
\BIC{$\forall z\,(s_{k} < s_{N_{1}} \to s_{N_{1}} \leq z \to s_{k} < z)$}
\AXC{$\varphi_{120}$\lbl{A2}}
\BIC{$s_{k} < s_{N_{1}} \to s_{N_{1}} \leq n \to s_{k} < n$}
\AXC{$\varphi_{484}$\lbl{A1}}
\BIC{$(\mathrm{Seq}(a) \to \mathrm{SubIdx}(s) \to \varphi_{393}) \to \varphi_{452}$}
\end{ptree}
\noindent$\Pi_{44}$:
\begin{ptree}
\AXC{$\mathrm{Seq}(a) \to \mathrm{SubIdx}(s) \to \varphi_{384}$\lbl{A1}}
\AXC{$\stackrel{\Pi_{42}}{\vdots}$}\noLine
\UIC{$\varphi_{460}$}
\BIC{$\mathrm{Seq}(a) \to \mathrm{SubIdx}(s) \to \varphi_{393}$}
\AXC{$\stackrel{\Pi_{43}}{\vdots}$}\noLine
\UIC{$(\mathrm{Seq}(a) \to \mathrm{SubIdx}(s) \to \varphi_{393}) \to \varphi_{452}$}
\BIC{$\varphi_{452}$}
\end{ptree}
\noindent$\Pi_{45}$:
\begin{ptree}
\AXC{$\stackrel{\Pi_{10}}{\vdots}$}\noLine
\UIC{$\forall x\,\forall y\,(x = y \to y = x)$}
\AXC{$\forall x\,\forall y\,(x = y \to y = x) \to \forall y\,(n = y \to y = n)$\lbl{A2}}
\BIC{$\forall y\,(n = y \to y = n)$}
\AXC{$\forall y\,(n = y \to y = n) \to n = s_{k} \to s_{k} = n$\lbl{A2}}
\BIC{$n = s_{k} \to s_{k} = n$}
\end{ptree}
\noindent$\Pi_{46}$:
\begin{ptree}
\AXC{$\stackrel{\Pi_{45}}{\vdots}$}\noLine
\UIC{$n = s_{k} \to s_{k} = n$}
\AXC{$\mathrm{Seq}(a) \to \mathrm{SubIdx}(s) \to \varphi_{367}$\lbl{A1}}
\AXC{$\varphi_{458}$\lbl{A1}}
\BIC{$(n = s_{k} \to s_{k} = n) \to \mathrm{Seq}(a) \to \mathrm{SubIdx}(s) \to \varphi_{368}$}
\BIC{$\mathrm{Seq}(a) \to \mathrm{SubIdx}(s) \to \varphi_{368}$}
\AXC{$\varphi_{482}$\lbl{A1}}
\BIC{$(\mathrm{Seq}(a) \to \mathrm{SubIdx}(s) \to \varphi_{389}) \to \varphi_{416}$}
\end{ptree}
\noindent$\Pi_{47}$:
\begin{ptree}
\AXC{$\stackrel{\Pi_{39}}{\vdots}$}\noLine
\UIC{$\mathrm{Seq}(a) \to \mathrm{SubIdx}(s) \to \varphi_{372}$}
\AXC{$\stackrel{\Pi_{44}}{\vdots}$}\noLine
\UIC{$\varphi_{452}$}
\BIC{$\mathrm{Seq}(a) \to \mathrm{SubIdx}(s) \to \varphi_{389}$}
\AXC{$\stackrel{\Pi_{46}}{\vdots}$}\noLine
\UIC{$(\mathrm{Seq}(a) \to \mathrm{SubIdx}(s) \to \varphi_{389}) \to \varphi_{416}$}
\BIC{$\varphi_{416}$}
\end{ptree}
\noindent$\Pi_{48}$:
\begin{ptree}
\AXC{$z = w \to z < n \to w < n$\lbl{A6}}
\UIC{$\forall w\,(z = w \to z < n \to w < n)$}
\UIC{$\forall z\,\forall w\,(z = w \to z < n \to w < n)$}
\AXC{$\varphi_{26}$\lbl{A2}}
\BIC{$\forall w\,(s_{k} = w \to s_{k} < n \to w < n)$}
\AXC{$\varphi_{49}$\lbl{A2}}
\BIC{$s_{k} = n \to s_{k} < n \to n < n$}
\AXC{$\stackrel{\Pi_{47}}{\vdots}$}\noLine
\UIC{$\varphi_{416}$}
\BIC{$\mathrm{Seq}(a) \to \mathrm{SubIdx}(s) \to \varphi_{382}$}
\AXC{$\varphi_{468}$\lbl{A1}}
\BIC{$\neg n < n \to \mathrm{Seq}(a) \to \mathrm{SubIdx}(s) \to \varphi_{398}$}
\end{ptree}
\noindent$\Pi_{49}$:
\begin{ptree}
\AXC{$\forall x\,(\neg x < x)$\lbl{[O1]}}
\AXC{$\forall x\,(\neg x < x) \to \neg n < n$\lbl{A2}}
\BIC{$\neg n < n$}
\AXC{$\stackrel{\Pi_{48}}{\vdots}$}\noLine
\UIC{$\neg n < n \to \mathrm{Seq}(a) \to \mathrm{SubIdx}(s) \to \varphi_{398}$}
\BIC{$\mathrm{Seq}(a) \to \mathrm{SubIdx}(s) \to \varphi_{398}$}
\end{ptree}
\noindent$\Pi_{50}$:
\begin{ptree}
\AXC{$z = w \to \lvert z - \alpha \rvert < \varepsilon \to \lvert w - \alpha \rvert < \varepsilon$\lbl{A6}}
\UIC{$\forall w\,(z = w \to \lvert z - \alpha \rvert < \varepsilon \to \lvert w - \alpha \rvert < \varepsilon)$}
\UIC{$\forall z\,\forall w\,(z = w \to \lvert z - \alpha \rvert < \varepsilon \to \lvert w - \alpha \rvert < \varepsilon)$}
\AXC{$\forall z\,\forall w\,(z = w \to \lvert z - \alpha \rvert < \varepsilon \to \lvert w - \alpha \rvert < \varepsilon) \to \varphi_{61}$\lbl{A2}}
\BIC{$\varphi_{61}$}
\AXC{$\varphi_{61} \to \varphi_{114}$\lbl{A2}}
\BIC{$\varphi_{114}$}
\end{ptree}
\noindent$\Pi_{51}$:
\begin{ptree}
\AXC{$\varphi_{125}$\lbl{A2}}
\AXC{$\mathrm{Seq}(a) \to \mathrm{SubIdx}(s) \to \varphi_{279}$\lbl{A1}}
\AXC{$\varphi_{439}$\lbl{A1}}
\BIC{$\mathrm{Seq}(a) \to \mathrm{SubIdx}(s) \to \varphi_{413}$}
\AXC{$\varphi_{477}$\lbl{A1}}
\BIC{$\varphi_{125} \to \mathrm{Seq}(a) \to \mathrm{SubIdx}(s) \to \varphi_{411}$}
\BIC{$\mathrm{Seq}(a) \to \mathrm{SubIdx}(s) \to \varphi_{411}$}
\AXC{$\varphi_{491}$\lbl{A1}}
\BIC{$(\mathrm{Seq}(a) \to \mathrm{SubIdx}(s) \to \varphi_{358}) \to \varphi_{464}$}
\end{ptree}
\noindent$\Pi_{52}$:
\begin{ptree}
\AXC{$\mathrm{Seq}(a) \to \mathrm{SubIdx}(s) \to \varphi_{386}$\lbl{A1}}
\AXC{$\varphi_{446}$\lbl{A1}}
\BIC{$\mathrm{Seq}(a) \to \mathrm{SubIdx}(s) \to \varphi_{358}$}
\AXC{$\stackrel{\Pi_{51}}{\vdots}$}\noLine
\UIC{$(\mathrm{Seq}(a) \to \mathrm{SubIdx}(s) \to \varphi_{358}) \to \varphi_{464}$}
\BIC{$\varphi_{464}$}
\end{ptree}
\noindent$\Pi_{53}$:
\begin{ptree}
\AXC{$\forall s\,(\varphi_{6})$\lbl{DefSubIdx}}
\AXC{$\forall s\,(\varphi_{6}) \to \varphi_{6}$\lbl{A2}}
\BIC{$\varphi_{6}$}
\AXC{$\mathrm{Seq}(a) \to \mathrm{SubIdx}(s) \to \mathrm{SubIdx}(s)$\lbl{A1}}
\AXC{$\varphi_{230}$\lbl{A1}}
\BIC{$\varphi_{6} \to \mathrm{Seq}(a) \to \mathrm{SubIdx}(s) \to \varphi_{127}$}
\BIC{$\mathrm{Seq}(a) \to \mathrm{SubIdx}(s) \to \varphi_{127}$}
\end{ptree}
\noindent$\Pi_{54}$:
\begin{ptree}
\AXC{$\forall a\,\forall s\,(\mathrm{Seq}(a) \to \varphi_{47})$\lbl{[S12]}}
\AXC{$\forall a\,\forall s\,(\mathrm{Seq}(a) \to \varphi_{47}) \to \forall s\,(\mathrm{Seq}(a) \to \varphi_{47})$\lbl{A2}}
\BIC{$\forall s\,(\mathrm{Seq}(a) \to \varphi_{47})$}
\end{ptree}
\noindent$\Pi_{55}$:
\begin{ptree}
\AXC{$\stackrel{\Pi_{54}}{\vdots}$}\noLine
\UIC{$\forall s\,(\mathrm{Seq}(a) \to \varphi_{47})$}
\AXC{$\forall s\,(\mathrm{Seq}(a) \to \varphi_{47}) \to \mathrm{Seq}(a) \to \varphi_{47}$\lbl{A2}}
\BIC{$\mathrm{Seq}(a) \to \varphi_{47}$}
\AXC{$(\mathrm{Seq}(a) \to \varphi_{47}) \to (\mathrm{Seq}(a) \to \mathrm{Seq}(a)) \to \varphi_{332}$\lbl{A1}}
\BIC{$(\mathrm{Seq}(a) \to \mathrm{Seq}(a)) \to \varphi_{332}$}
\end{ptree}
\noindent$\Pi_{56}$:
\begin{ptree}
\AXC{$\stackrel{\Pi_{53}}{\vdots}$}\noLine
\UIC{$\mathrm{Seq}(a) \to \mathrm{SubIdx}(s) \to \varphi_{127}$}
\AXC{$\mathrm{Seq}(a) \to \mathrm{Seq}(a)$\lbl{A1}}
\AXC{$\stackrel{\Pi_{55}}{\vdots}$}\noLine
\UIC{$(\mathrm{Seq}(a) \to \mathrm{Seq}(a)) \to \varphi_{332}$}
\BIC{$\varphi_{332}$}
\BIC{$\mathrm{Seq}(a) \to \mathrm{SubIdx}(s) \to \mathrm{SubIdx}(t) \to \varphi_{163}$}
\AXC{$\varphi_{363}$\lbl{A1}}
\BIC{$\mathrm{Seq}(a) \to \mathrm{SubIdx}(s) \to \varphi_{158}$}
\end{ptree}
\noindent$\Pi_{57}$:
\begin{ptree}
\AXC{$\varphi_{106}$\lbl{A2}}
\AXC{$\stackrel{\Pi_{56}}{\vdots}$}\noLine
\UIC{$\mathrm{Seq}(a) \to \mathrm{SubIdx}(s) \to \varphi_{158}$}
\AXC{$\varphi_{419}$\lbl{A1}}
\BIC{$\varphi_{106} \to \mathrm{Seq}(a) \to \mathrm{SubIdx}(s) \to \varphi_{154}$}
\BIC{$\mathrm{Seq}(a) \to \mathrm{SubIdx}(s) \to \varphi_{154}$}
\AXC{$(\mathrm{Seq}(a) \to \mathrm{SubIdx}(s) \to \varphi_{154}) \to \varphi_{454}$\lbl{A1}}
\BIC{$\varphi_{454}$}
\end{ptree}
\noindent$\Pi_{58}$:
\begin{ptree}
\AXC{$\mathrm{Seq}(a) \to \mathrm{SubIdx}(s) \to \varphi_{386}$\lbl{A1}}
\AXC{$\varphi_{446}$\lbl{A1}}
\BIC{$\mathrm{Seq}(a) \to \mathrm{SubIdx}(s) \to \varphi_{358}$}
\AXC{$\stackrel{\Pi_{57}}{\vdots}$}\noLine
\UIC{$\varphi_{454}$}
\BIC{$\mathrm{Seq}(a) \to \mathrm{SubIdx}(s) \to \varphi_{400}$}
\AXC{$(\mathrm{Seq}(a) \to \mathrm{SubIdx}(s) \to \varphi_{400}) \to \varphi_{473}$\lbl{A1}}
\BIC{$\varphi_{473}$}
\end{ptree}
\noindent$\Pi_{59}$:
\begin{ptree}
\AXC{$\stackrel{\Pi_{10}}{\vdots}$}\noLine
\UIC{$\forall x\,\forall y\,(x = y \to y = x)$}
\AXC{$\forall x\,\forall y\,(x = y \to y = x) \to \forall y\,(n = y \to y = n)$\lbl{A2}}
\BIC{$\forall y\,(n = y \to y = n)$}
\AXC{$\forall y\,(n = y \to y = n) \to n = s_{k} \to s_{k} = n$\lbl{A2}}
\BIC{$n = s_{k} \to s_{k} = n$}
\end{ptree}
\noindent$\Pi_{60}$:
\begin{ptree}
\AXC{$\stackrel{\Pi_{59}}{\vdots}$}\noLine
\UIC{$n = s_{k} \to s_{k} = n$}
\AXC{$\mathrm{Seq}(a) \to \mathrm{SubIdx}(s) \to \varphi_{367}$\lbl{A1}}
\AXC{$\varphi_{458}$\lbl{A1}}
\BIC{$(n = s_{k} \to s_{k} = n) \to \mathrm{Seq}(a) \to \mathrm{SubIdx}(s) \to \varphi_{368}$}
\BIC{$\mathrm{Seq}(a) \to \mathrm{SubIdx}(s) \to \varphi_{368}$}
\AXC{$\varphi_{486}$\lbl{A1}}
\BIC{$(\mathrm{Seq}(a) \to \mathrm{SubIdx}(s) \to \varphi_{402}) \to \varphi_{423}$}
\end{ptree}
\noindent$\Pi_{61}$:
\begin{ptree}
\AXC{$\stackrel{\Pi_{50}}{\vdots}$}\noLine
\UIC{$\varphi_{114}$}
\AXC{$\mathrm{Seq}(a) \to \mathrm{SubIdx}(s) \to \varphi_{373}$\lbl{A1}}
\AXC{$\stackrel{\Pi_{52}}{\vdots}$}\noLine
\UIC{$\varphi_{464}$}
\BIC{$\mathrm{Seq}(a) \to \mathrm{SubIdx}(s) \to \varphi_{404}$}
\AXC{$\stackrel{\Pi_{58}}{\vdots}$}\noLine
\UIC{$\varphi_{473}$}
\BIC{$\varphi_{114} \to \mathrm{Seq}(a) \to \mathrm{SubIdx}(s) \to \varphi_{402}$}
\BIC{$\mathrm{Seq}(a) \to \mathrm{SubIdx}(s) \to \varphi_{402}$}
\AXC{$\stackrel{\Pi_{60}}{\vdots}$}\noLine
\UIC{$(\mathrm{Seq}(a) \to \mathrm{SubIdx}(s) \to \varphi_{402}) \to \varphi_{423}$}
\BIC{$\varphi_{423}$}
\end{ptree}
\noindent$\Pi_{62}$:
\begin{ptree}
\AXC{$z = w \to \lvert a_{z} - \alpha \rvert < \varepsilon \to \lvert a_{w} - \alpha \rvert < \varepsilon$\lbl{A6}}
\UIC{$\forall w\,(z = w \to \lvert a_{z} - \alpha \rvert < \varepsilon \to \lvert a_{w} - \alpha \rvert < \varepsilon)$}
\UIC{$\forall z\,\forall w\,(z = w \to \lvert a_{z} - \alpha \rvert < \varepsilon \to \lvert a_{w} - \alpha \rvert < \varepsilon)$}
\AXC{$\varphi_{149}$\lbl{A2}}
\BIC{$\forall w\,(s_{k} = w \to \lvert a_{s_{k}} - \alpha \rvert < \varepsilon \to \lvert a_{w} - \alpha \rvert < \varepsilon)$}
\AXC{$\varphi_{156}$\lbl{A2}}
\BIC{$s_{k} = n \to \lvert a_{s_{k}} - \alpha \rvert < \varepsilon \to \lvert a_{n} - \alpha \rvert < \varepsilon$}
\AXC{$\stackrel{\Pi_{61}}{\vdots}$}\noLine
\UIC{$\varphi_{423}$}
\BIC{$\mathrm{Seq}(a) \to \mathrm{SubIdx}(s) \to \varphi_{395}$}
\end{ptree}
\noindent$\Pi_{63}$:
\begin{ptree}
\AXC{$\forall n \in \mathbb{N}\,(n \in \omega)$\lbl{[N7']}}
\AXC{$\forall n \in \mathbb{N}\,(n \in \omega) \to k \in \mathbb{N} \to k \in \omega$\lbl{A2}}
\BIC{$k \in \mathbb{N} \to k \in \omega$}
\AXC{$(k \in \mathbb{N} \to k \in \omega) \to \varphi_{445}$\lbl{A1}}
\BIC{$\varphi_{445}$}
\end{ptree}
\noindent$\Pi_{64}$:
\begin{ptree}
\AXC{$\forall n \in \mathbb{N}\,(n \in \omega)$\lbl{[N7']}}
\AXC{$\forall n \in \mathbb{N}\,(n \in \omega) \to N_{1} \in \mathbb{N} \to N_{1} \in \omega$\lbl{A2}}
\BIC{$N_{1} \in \mathbb{N} \to N_{1} \in \omega$}
\AXC{$(N_{1} \in \mathbb{N} \to N_{1} \in \omega) \to \varphi_{435}$\lbl{A1}}
\BIC{$\varphi_{435}$}
\end{ptree}
\noindent$\Pi_{65}$:
\begin{ptree}
\AXC{$\forall m\,\forall n\,(m \in \omega \to n \in \omega \to m \leq n \lor n < m)$\lbl{[N11]}}
\AXC{$\varphi_{108}$\lbl{A2}}
\BIC{$\forall n\,(N_{1} \in \omega \to n \in \omega \to N_{1} \leq n \lor n < N_{1})$}
\AXC{$\varphi_{97}$\lbl{A2}}
\BIC{$N_{1} \in \omega \to k \in \omega \to N_{1} \leq k \lor k < N_{1}$}
\AXC{$\varphi_{480}$\lbl{A1}}
\BIC{$(\mathrm{Seq}(a) \to \mathrm{SubIdx}(s) \to \varphi_{376}) \to \varphi_{448}$}
\end{ptree}
\noindent$\Pi_{66}$:
\begin{ptree}
\AXC{$\mathrm{Seq}(a) \to \mathrm{SubIdx}(s) \to \varphi_{279}$\lbl{A1}}
\AXC{$\varphi_{425}$\lbl{A1}}
\BIC{$\mathrm{Seq}(a) \to \mathrm{SubIdx}(s) \to \varphi_{289}$}
\AXC{$\stackrel{\Pi_{64}}{\vdots}$}\noLine
\UIC{$\varphi_{435}$}
\BIC{$\mathrm{Seq}(a) \to \mathrm{SubIdx}(s) \to \varphi_{376}$}
\AXC{$\stackrel{\Pi_{65}}{\vdots}$}\noLine
\UIC{$(\mathrm{Seq}(a) \to \mathrm{SubIdx}(s) \to \varphi_{376}) \to \varphi_{448}$}
\BIC{$\varphi_{448}$}
\end{ptree}
\noindent$\Pi_{67}$:
\begin{ptree}
\AXC{$\mathrm{Seq}(a) \to \mathrm{SubIdx}(s) \to \varphi_{386}$\lbl{A1}}
\AXC{$\varphi_{446}$\lbl{A1}}
\BIC{$\mathrm{Seq}(a) \to \mathrm{SubIdx}(s) \to \varphi_{358}$}
\AXC{$\stackrel{\Pi_{63}}{\vdots}$}\noLine
\UIC{$\varphi_{445}$}
\BIC{$\mathrm{Seq}(a) \to \mathrm{SubIdx}(s) \to \varphi_{374}$}
\AXC{$\stackrel{\Pi_{66}}{\vdots}$}\noLine
\UIC{$\varphi_{448}$}
\BIC{$\mathrm{Seq}(a) \to \mathrm{SubIdx}(s) \to \varphi_{380}$}
\AXC{$\varphi_{493}$\lbl{A1}}
\BIC{$(\mathrm{Seq}(a) \to \mathrm{SubIdx}(s) \to \varphi_{395}) \to \varphi_{469}$}
\end{ptree}
\noindent$\Pi_{68}$:
\begin{ptree}
\AXC{$\stackrel{\Pi_{49}}{\vdots}$}\noLine
\UIC{$\mathrm{Seq}(a) \to \mathrm{SubIdx}(s) \to \varphi_{398}$}
\AXC{$\stackrel{\Pi_{62}}{\vdots}$}\noLine
\UIC{$\mathrm{Seq}(a) \to \mathrm{SubIdx}(s) \to \varphi_{395}$}
\AXC{$\stackrel{\Pi_{67}}{\vdots}$}\noLine
\UIC{$(\mathrm{Seq}(a) \to \mathrm{SubIdx}(s) \to \varphi_{395}) \to \varphi_{469}$}
\BIC{$\varphi_{469}$}
\BIC{$\mathrm{Seq}(a) \to \mathrm{SubIdx}(s) \to \varphi_{391}$}
\end{ptree}
\noindent$\Pi_{69}$:
\begin{ptree}
\AXC{$\stackrel{\Pi_{68}}{\vdots}$}\noLine
\UIC{$\mathrm{Seq}(a) \to \mathrm{SubIdx}(s) \to \varphi_{391}$}
\AXC{$\mathrm{Seq}(a) \to \mathrm{SubIdx}(s) \to \varphi_{386}$\lbl{A1}}
\AXC{$\varphi_{489}$\lbl{A1}}
\BIC{$(\mathrm{Seq}(a) \to \mathrm{SubIdx}(s) \to \varphi_{391}) \to \varphi_{456}$}
\BIC{$\varphi_{456}$}
\end{ptree}
\noindent$\Pi_{70}$:
\begin{ptree}
\AXC{$\stackrel{\Pi_{36}}{\vdots}$}\noLine
\UIC{$\mathrm{Seq}(a) \to \mathrm{SubIdx}(s) \to \varphi_{390}$}
\AXC{$\stackrel{\Pi_{69}}{\vdots}$}\noLine
\UIC{$\varphi_{456}$}
\BIC{$\mathrm{Seq}(a) \to \mathrm{SubIdx}(s) \to \varphi_{377}$}
\AXC{$(\mathrm{Seq}(a) \to \mathrm{SubIdx}(s) \to \varphi_{377}) \to \varphi_{396}$\lbl{A1}}
\BIC{$\varphi_{396}$}
\UIC{$\forall k\,(\varphi_{396})$}
\end{ptree}
\noindent$\Pi_{71}$:
\begin{ptree}
\AXC{$\varphi_{118}$\lbl{A2}}
\AXC{$\mathrm{Seq}(a) \to \mathrm{SubIdx}(s) \to \mathrm{SubIdx}(t) \to \varphi_{122}$\lbl{A1}}
\AXC{$\varphi_{385}$\lbl{A1}}
\BIC{$\varphi_{118} \to \mathrm{Seq}(a) \to \mathrm{SubIdx}(s) \to \mathrm{SubIdx}(t) \to \varphi_{118}$}
\BIC{$\mathrm{Seq}(a) \to \mathrm{SubIdx}(s) \to \mathrm{SubIdx}(t) \to \varphi_{118}$}
\AXC{$\varphi_{465}$\lbl{A1}}
\BIC{$\varphi_{433}$}
\end{ptree}
\noindent$\Pi_{72}$:
\begin{ptree}
\AXC{$\varphi_{437}$\lbl{A4}}
\AXC{$\varphi_{455}$\lbl{A3}}
\AXC{$\stackrel{\Pi_{70}}{\vdots}$}\noLine
\UIC{$\forall k\,(\varphi_{396})$}
\AXC{$\mathrm{Seq}(a) \to \mathrm{SubIdx}(s) \to \varphi_{302}$\lbl{A1}}
\AXC{$\stackrel{\Pi_{71}}{\vdots}$}\noLine
\UIC{$\varphi_{433}$}
\BIC{$\mathrm{Seq}(a) \to \mathrm{SubIdx}(s) \to \varphi_{352}$}
\AXC{$\varphi_{498}$\lbl{A1}}
\BIC{$\forall k\,(\varphi_{396}) \to \varphi_{455} \to \varphi_{437} \to \mathrm{Seq}(a) \to \mathrm{SubIdx}(s) \to \varphi_{345}$}
\BIC{$\varphi_{455} \to \varphi_{437} \to \mathrm{Seq}(a) \to \mathrm{SubIdx}(s) \to \varphi_{345}$}
\BIC{$\varphi_{437} \to \mathrm{Seq}(a) \to \mathrm{SubIdx}(s) \to \varphi_{345}$}
\BIC{$\mathrm{Seq}(a) \to \mathrm{SubIdx}(s) \to \varphi_{345}$}
\AXC{$\varphi_{436}$\lbl{A1}}
\BIC{$\varphi_{298} \to n \in \mathbb{N} \to s_{N_{1}} \cup t(N_{2}) \leq n \to \lvert a_{n} - \alpha \rvert < \varepsilon$}
\UIC{$\varphi_{359}$}
\end{ptree}
\noindent$\Pi_{73}$:
\begin{ptree}
\AXC{$\stackrel{\Pi_{20}}{\vdots}$}\noLine
\UIC{$\mathrm{Seq}(a) \to \mathrm{SubIdx}(s) \to \varphi_{126}$}
\AXC{$\forall s\,\forall k\,(\mathrm{NSeq}(s) \to k \in \mathbb{N} \to s_{k} \in \mathbb{N})$\lbl{[S13]}}
\AXC{$\varphi_{78}$\lbl{A2}}
\BIC{$\forall k\,(\mathrm{NSeq}(t) \to k \in \mathbb{N} \to t(k) \in \mathbb{N})$}
\AXC{$\varphi_{68}$\lbl{A2}}
\BIC{$\mathrm{NSeq}(t) \to N_{2} \in \mathbb{N} \to t(N_{2}) \in \mathbb{N}$}
\AXC{$\varphi_{449}$\lbl{A1}}
\BIC{$(\mathrm{Seq}(a) \to \mathrm{SubIdx}(s) \to \varphi_{126}) \to \varphi_{426}$}
\BIC{$\varphi_{426}$}
\end{ptree}
\noindent$\Pi_{74}$:
\begin{ptree}
\AXC{$\stackrel{\Pi_{53}}{\vdots}$}\noLine
\UIC{$\mathrm{Seq}(a) \to \mathrm{SubIdx}(s) \to \varphi_{127}$}
\AXC{$\forall s\,\forall k\,(\mathrm{NSeq}(s) \to k \in \mathbb{N} \to s_{k} \in \mathbb{N})$\lbl{[S13]}}
\AXC{$\varphi_{79}$\lbl{A2}}
\BIC{$\forall k\,(\mathrm{NSeq}(s) \to k \in \mathbb{N} \to s_{k} \in \mathbb{N})$}
\AXC{$\varphi_{69}$\lbl{A2}}
\BIC{$\mathrm{NSeq}(s) \to N_{1} \in \mathbb{N} \to s_{N_{1}} \in \mathbb{N}$}
\AXC{$\varphi_{450}$\lbl{A1}}
\BIC{$(\mathrm{Seq}(a) \to \mathrm{SubIdx}(s) \to \varphi_{127}) \to \varphi_{427}$}
\BIC{$\varphi_{427}$}
\end{ptree}
\noindent$\Pi_{75}$:
\begin{ptree}
\AXC{$\forall m\,\forall n\,(m \in \mathbb{N} \to n \in \mathbb{N} \to m \cup n \in \mathbb{N})$\lbl{[N3]}}
\AXC{$\varphi_{121}$\lbl{A2}}
\BIC{$\forall n\,(s_{N_{1}} \in \mathbb{N} \to n \in \mathbb{N} \to s_{N_{1}} \cup n \in \mathbb{N})$}
\AXC{$\varphi_{131}$\lbl{A2}}
\BIC{$s_{N_{1}} \in \mathbb{N} \to t(N_{2}) \in \mathbb{N} \to s_{N_{1}} \cup t(N_{2}) \in \mathbb{N}$}
\AXC{$\varphi_{478}$\lbl{A1}}
\BIC{$(\mathrm{Seq}(a) \to \mathrm{SubIdx}(s) \to \varphi_{297}) \to \varphi_{428}$}
\end{ptree}
\noindent$\Pi_{76}$:
\begin{ptree}
\AXC{$\mathrm{Seq}(a) \to \mathrm{SubIdx}(s) \to \varphi_{279}$\lbl{A1}}
\AXC{$\varphi_{425}$\lbl{A1}}
\BIC{$\mathrm{Seq}(a) \to \mathrm{SubIdx}(s) \to \varphi_{289}$}
\AXC{$\stackrel{\Pi_{74}}{\vdots}$}\noLine
\UIC{$\varphi_{427}$}
\BIC{$\mathrm{Seq}(a) \to \mathrm{SubIdx}(s) \to \varphi_{297}$}
\AXC{$\stackrel{\Pi_{75}}{\vdots}$}\noLine
\UIC{$(\mathrm{Seq}(a) \to \mathrm{SubIdx}(s) \to \varphi_{297}) \to \varphi_{428}$}
\BIC{$\varphi_{428}$}
\end{ptree}
\noindent$\Pi_{77}$:
\begin{ptree}
\AXC{$\mathrm{Seq}(a) \to \mathrm{SubIdx}(s) \to \varphi_{346}$\lbl{A1}}
\AXC{$\varphi_{431}$\lbl{A1}}
\BIC{$\mathrm{Seq}(a) \to \mathrm{SubIdx}(s) \to \varphi_{286}$}
\AXC{$\stackrel{\Pi_{73}}{\vdots}$}\noLine
\UIC{$\varphi_{426}$}
\BIC{$\mathrm{Seq}(a) \to \mathrm{SubIdx}(s) \to \varphi_{296}$}
\AXC{$\stackrel{\Pi_{76}}{\vdots}$}\noLine
\UIC{$\varphi_{428}$}
\BIC{$\mathrm{Seq}(a) \to \mathrm{SubIdx}(s) \to \varphi_{303}$}
\AXC{$(\mathrm{Seq}(a) \to \mathrm{SubIdx}(s) \to \varphi_{303}) \to \varphi_{442}$\lbl{A1}}
\BIC{$\varphi_{442}$}
\end{ptree}
\noindent$\Pi_{78}$:
\begin{ptree}
\AXC{$\varphi_{13} \to \neg \varphi_{117}$\lbl{A2}}
\AXC{$\varphi_{298} \to \forall n\,(\varphi_{298})$\lbl{A4}}
\AXC{$\varphi_{359} \to \varphi_{361}$\lbl{A3}}
\AXC{$\stackrel{\Pi_{72}}{\vdots}$}\noLine
\UIC{$\varphi_{359}$}
\AXC{$\varphi_{359} \to \varphi_{494}$\lbl{A1}}
\BIC{$\varphi_{494}$}
\BIC{$(\varphi_{298} \to \forall n\,(\varphi_{298})) \to \mathrm{Seq}(a) \to \mathrm{SubIdx}(s) \to \varphi_{347}$}
\BIC{$\mathrm{Seq}(a) \to \mathrm{SubIdx}(s) \to \varphi_{347}$}
\AXC{$\stackrel{\Pi_{77}}{\vdots}$}\noLine
\UIC{$\varphi_{442}$}
\BIC{$\mathrm{Seq}(a) \to \mathrm{SubIdx}(s) \to \varphi_{378}$}
\AXC{$\varphi_{471}$\lbl{A1}}
\BIC{$(\varphi_{13} \to \neg \varphi_{117}) \to \mathrm{Seq}(a) \to \mathrm{SubIdx}(s) \to \varphi_{337}$}
\BIC{$\mathrm{Seq}(a) \to \mathrm{SubIdx}(s) \to \varphi_{337}$}
\end{ptree}
\noindent$\Pi_{79}$:
\begin{ptree}
\AXC{$\stackrel{\Pi_{78}}{\vdots}$}\noLine
\UIC{$\mathrm{Seq}(a) \to \mathrm{SubIdx}(s) \to \varphi_{337}$}
\AXC{$(\mathrm{Seq}(a) \to \mathrm{SubIdx}(s) \to \varphi_{337}) \to \varphi_{293} \to \varphi_{34}$\lbl{A1}}
\BIC{$\varphi_{293} \to \varphi_{34}$}
\UIC{$\forall N_{2}\,(\varphi_{293} \to \varphi_{34})$}
\end{ptree}
\noindent$\Pi_{80}$:
\begin{ptree}
\AXC{$\forall a\,\forall \alpha\,\forall \varepsilon\,(\mathrm{Conv}(a,\alpha) \to \varphi_{80})$\lbl{[C3]}}
\AXC{$\varphi_{231}$\lbl{A2}}
\BIC{$\forall \alpha\,\forall \varepsilon\,(\mathrm{Conv}(a \circ t,\alpha) \to \varepsilon \in \mathbb{R} \to \varphi_{62})$}
\AXC{$\varphi_{252}$\lbl{A2}}
\BIC{$\forall \varepsilon\,(\mathrm{Conv}(a \circ t,\alpha) \to \varepsilon \in \mathbb{R} \to \varphi_{62})$}
\AXC{$\varphi_{232}$\lbl{A2}}
\BIC{$\mathrm{Conv}(a \circ t,\alpha) \to \varepsilon \in \mathbb{R} \to \varphi_{62}$}
\AXC{$(\mathrm{Conv}(a \circ t,\alpha) \to \varepsilon \in \mathbb{R} \to \varphi_{62}) \to \varphi_{443}$\lbl{A1}}
\BIC{$\varphi_{443}$}
\end{ptree}
\noindent$\Pi_{81}$:
\begin{ptree}
\AXC{$\mathrm{Seq}(a) \to \mathrm{SubIdx}(s) \to \varphi_{136}$\lbl{A1}}
\AXC{$\mathrm{Seq}(a) \to \mathrm{SubIdx}(s) \to \varphi_{130}$\lbl{A1}}
\AXC{$\stackrel{\Pi_{80}}{\vdots}$}\noLine
\UIC{$\varphi_{443}$}
\BIC{$(\mathrm{Seq}(a) \to \mathrm{SubIdx}(s) \to \varphi_{136}) \to \varphi_{420}$}
\BIC{$\varphi_{420}$}
\end{ptree}
\noindent$\Pi_{82}$:
\begin{ptree}
\AXC{$\stackrel{\Pi_{79}}{\vdots}$}\noLine
\UIC{$\forall N_{2}\,(\varphi_{293} \to \varphi_{34})$}
\AXC{$\mathrm{Seq}(a) \to \mathrm{SubIdx}(s) \to \varphi_{147}$\lbl{A1}}
\AXC{$\stackrel{\Pi_{81}}{\vdots}$}\noLine
\UIC{$\varphi_{420}$}
\BIC{$\mathrm{Seq}(a) \to \mathrm{SubIdx}(s) \to \varphi_{281}$}
\AXC{$\varphi_{499}$\lbl{A1}}
\BIC{$\forall N_{2}\,(\varphi_{293} \to \varphi_{34}) \to \varphi_{496}$}
\BIC{$\varphi_{496}$}
\end{ptree}
\noindent$\Pi_{83}$:
\begin{ptree}
\AXC{$\forall N_{2}\,(\varphi_{34}) \to \varphi_{29}$\lbl{A2}}
\UIC{$\forall N\,(\forall N_{2}\,(\varphi_{34}) \to \varphi_{29})$}
\AXC{$\varphi_{293} \to \forall N_{2}\,(\varphi_{293})$\lbl{A4}}
\AXC{$\forall N_{2}\,(\varphi_{293} \to \varphi_{34}) \to \forall N_{2}\,(\varphi_{293}) \to \forall N_{2}\,(\varphi_{34})$\lbl{A3}}
\AXC{$\stackrel{\Pi_{82}}{\vdots}$}\noLine
\UIC{$\varphi_{496}$}
\BIC{$(\varphi_{293} \to \forall N_{2}\,(\varphi_{293})) \to \forall N\,(\forall N_{2}\,(\varphi_{34}) \to \varphi_{29}) \to \varphi_{438}$}
\BIC{$\forall N\,(\forall N_{2}\,(\varphi_{34}) \to \varphi_{29}) \to \varphi_{438}$}
\BIC{$\varphi_{438}$}
\end{ptree}
\noindent$\Pi_{84}$:
\begin{ptree}
\AXC{$\forall N_{2}\,(\varphi_{34}) \to \forall N\,\forall N_{2}\,(\varphi_{34})$\lbl{A4}}
\AXC{$\forall N\,(\forall N_{2}\,(\varphi_{34}) \to \varphi_{29}) \to \forall N\,\forall N_{2}\,(\varphi_{34}) \to \forall N\,(\varphi_{29})$\lbl{A3}}
\AXC{$\stackrel{\Pi_{83}}{\vdots}$}\noLine
\UIC{$\varphi_{438}$}
\BIC{$\varphi_{417}$}
\BIC{$\mathrm{Seq}(a) \to \mathrm{SubIdx}(s) \to \varphi_{268}$}
\end{ptree}
\noindent$\Pi_{85}$:
\begin{ptree}
\AXC{$\stackrel{\Pi_{84}}{\vdots}$}\noLine
\UIC{$\mathrm{Seq}(a) \to \mathrm{SubIdx}(s) \to \varphi_{268}$}
\AXC{$(\mathrm{Seq}(a) \to \mathrm{SubIdx}(s) \to \varphi_{268}) \to \varphi_{218} \to \varphi_{35}$\lbl{A1}}
\BIC{$\varphi_{218} \to \varphi_{35}$}
\UIC{$\forall N_{1}\,(\varphi_{218} \to \varphi_{35})$}
\end{ptree}
\noindent$\Pi_{86}$:
\begin{ptree}
\AXC{$\forall a\,\forall \alpha\,\forall \varepsilon\,(\mathrm{Conv}(a,\alpha) \to \varphi_{80})$\lbl{[C3]}}
\AXC{$\varphi_{233}$\lbl{A2}}
\BIC{$\forall \alpha\,\forall \varepsilon\,(\mathrm{Conv}(a \circ s,\alpha) \to \varepsilon \in \mathbb{R} \to \varphi_{63})$}
\AXC{$\varphi_{253}$\lbl{A2}}
\BIC{$\forall \varepsilon\,(\mathrm{Conv}(a \circ s,\alpha) \to \varepsilon \in \mathbb{R} \to \varphi_{63})$}
\AXC{$\varphi_{234}$\lbl{A2}}
\BIC{$\mathrm{Conv}(a \circ s,\alpha) \to \varepsilon \in \mathbb{R} \to \varphi_{63}$}
\AXC{$(\mathrm{Conv}(a \circ s,\alpha) \to \varepsilon \in \mathbb{R} \to \varphi_{63}) \to \varphi_{432}$\lbl{A1}}
\BIC{$\varphi_{432}$}
\end{ptree}
\noindent$\Pi_{87}$:
\begin{ptree}
\AXC{$\mathrm{Seq}(a) \to \mathrm{SubIdx}(s) \to \varphi_{136}$\lbl{A1}}
\AXC{$\mathrm{Seq}(a) \to \mathrm{SubIdx}(s) \to \mathrm{SubIdx}(t) \to \varphi_{109}$\lbl{A1}}
\AXC{$\stackrel{\Pi_{86}}{\vdots}$}\noLine
\UIC{$\varphi_{432}$}
\BIC{$(\mathrm{Seq}(a) \to \mathrm{SubIdx}(s) \to \varphi_{136}) \to \varphi_{387}$}
\BIC{$\varphi_{387}$}
\end{ptree}
\noindent$\Pi_{88}$:
\begin{ptree}
\AXC{$\stackrel{\Pi_{85}}{\vdots}$}\noLine
\UIC{$\forall N_{1}\,(\varphi_{218} \to \varphi_{35})$}
\AXC{$\mathrm{Seq}(a) \to \mathrm{SubIdx}(s) \to \varphi_{147}$\lbl{A1}}
\AXC{$\stackrel{\Pi_{87}}{\vdots}$}\noLine
\UIC{$\varphi_{387}$}
\BIC{$\mathrm{Seq}(a) \to \mathrm{SubIdx}(s) \to \mathrm{SubIdx}(t) \to \varphi_{193}$}
\AXC{$\varphi_{497}$\lbl{A1}}
\BIC{$\forall N_{1}\,(\varphi_{218} \to \varphi_{35}) \to \varphi_{495}$}
\BIC{$\varphi_{495}$}
\end{ptree}
\noindent$\Pi_{89}$:
\begin{ptree}
\AXC{$\forall N_{1}\,(\varphi_{35}) \to \varphi_{28}$\lbl{A2}}
\UIC{$\forall N\,(\forall N_{1}\,(\varphi_{35}) \to \varphi_{28})$}
\AXC{$\varphi_{218} \to \forall N_{1}\,(\varphi_{218})$\lbl{A4}}
\AXC{$\forall N_{1}\,(\varphi_{218} \to \varphi_{35}) \to \forall N_{1}\,(\varphi_{218}) \to \forall N_{1}\,(\varphi_{35})$\lbl{A3}}
\AXC{$\stackrel{\Pi_{88}}{\vdots}$}\noLine
\UIC{$\varphi_{495}$}
\BIC{$(\varphi_{218} \to \forall N_{1}\,(\varphi_{218})) \to \forall N\,(\forall N_{1}\,(\varphi_{35}) \to \varphi_{28}) \to \varphi_{430}$}
\BIC{$\forall N\,(\forall N_{1}\,(\varphi_{35}) \to \varphi_{28}) \to \varphi_{430}$}
\BIC{$\varphi_{430}$}
\end{ptree}
\noindent$\Pi_{90}$:
\begin{ptree}
\AXC{$\forall N_{1}\,(\varphi_{35}) \to \forall N\,\forall N_{1}\,(\varphi_{35})$\lbl{A4}}
\AXC{$\forall N\,(\forall N_{1}\,(\varphi_{35}) \to \varphi_{28}) \to \forall N\,\forall N_{1}\,(\varphi_{35}) \to \forall N\,(\varphi_{28})$\lbl{A3}}
\AXC{$\stackrel{\Pi_{89}}{\vdots}$}\noLine
\UIC{$\varphi_{430}$}
\BIC{$\varphi_{364}$}
\BIC{$\mathrm{Seq}(a) \to \mathrm{SubIdx}(s) \to \mathrm{SubIdx}(t) \to \varphi_{187}$}
\AXC{$\varphi_{414}$\lbl{A1}}
\BIC{$\varphi_{129} \to \varphi_{80}$}
\UIC{$\forall \varepsilon\,(\varphi_{129} \to \varphi_{80})$}
\end{ptree}
\noindent$\Pi_{91}$:
\begin{ptree}
\AXC{$\forall \varepsilon\,(\varphi_{129} \to \varphi_{80}) \to \forall \varepsilon\,(\varphi_{129}) \to \varphi_{102}$\lbl{A3}}
\AXC{$\stackrel{\Pi_{90}}{\vdots}$}\noLine
\UIC{$\forall \varepsilon\,(\varphi_{129} \to \varphi_{80})$}
\AXC{$\forall \varepsilon\,(\varphi_{129} \to \varphi_{80}) \to (\forall \varepsilon\,(\varphi_{129} \to \varphi_{80}) \to \forall \varepsilon\,(\varphi_{129}) \to \varphi_{102}) \to \varphi_{421}$\lbl{A1}}
\BIC{$(\forall \varepsilon\,(\varphi_{129} \to \varphi_{80}) \to \forall \varepsilon\,(\varphi_{129}) \to \varphi_{102}) \to \varphi_{421}$}
\BIC{$\varphi_{421}$}
\end{ptree}
\noindent$\Pi_{92}$:
\begin{ptree}
\AXC{$\forall a\,\forall \alpha\,(\mathrm{Conv}(a,\alpha) \to \alpha \in \mathbb{R})$\lbl{[C2]}}
\AXC{$\varphi_{33}$\lbl{A2}}
\BIC{$\forall \alpha\,(\mathrm{Conv}(a \circ s,\alpha) \to \alpha \in \mathbb{R})$}
\AXC{$\varphi_{38}$\lbl{A2}}
\BIC{$\mathrm{Conv}(a \circ s,\alpha) \to \alpha \in \mathbb{R}$}
\AXC{$(\mathrm{Conv}(a \circ s,\alpha) \to \alpha \in \mathbb{R}) \to \varphi_{272}$\lbl{A1}}
\BIC{$\varphi_{272}$}
\end{ptree}
\noindent$\Pi_{93}$:
\begin{ptree}
\AXC{$\forall a\,\forall \alpha\,(\mathrm{Conv}(a,\alpha) \leftrightarrow \mathrm{Seq}(a) \land \alpha \in \mathbb{R} \land \varphi_{102})$\lbl{DefConv}}
\AXC{$\varphi_{269}$\lbl{A2}}
\BIC{$\forall \alpha\,(\mathrm{Conv}(a,\alpha) \leftrightarrow \mathrm{Seq}(a) \land \alpha \in \mathbb{R} \land \varphi_{102})$}
\AXC{$\varphi_{263}$\lbl{A2}}
\BIC{$\mathrm{Conv}(a,\alpha) \leftrightarrow \mathrm{Seq}(a) \land \alpha \in \mathbb{R} \land \varphi_{102}$}
\AXC{$\varphi_{429}$\lbl{A1}}
\BIC{$(\mathrm{Seq}(a) \to \mathrm{Seq}(a)) \to \varphi_{424}$}
\end{ptree}
\noindent$\Pi_{94}$:
\begin{ptree}
\AXC{$\mathrm{Seq}(a) \to \mathrm{Seq}(a)$\lbl{A1}}
\AXC{$\stackrel{\Pi_{93}}{\vdots}$}\noLine
\UIC{$(\mathrm{Seq}(a) \to \mathrm{Seq}(a)) \to \varphi_{424}$}
\BIC{$\varphi_{424}$}
\end{ptree}
\noindent 主木:
\begin{ptree}
\AXC{$\varphi_{129} \to \forall \varepsilon\,(\varphi_{129})$\lbl{A4}}
\AXC{$\stackrel{\Pi_{91}}{\vdots}$}\noLine
\UIC{$\varphi_{421}$}
\BIC{$\mathrm{Seq}(a) \to \mathrm{SubIdx}(s) \to \mathrm{SubIdx}(t) \to \varphi_{192}$}
\AXC{$\mathrm{Seq}(a) \to \mathrm{SubIdx}(s) \to \mathrm{SubIdx}(t) \to \varphi_{109}$\lbl{A1}}
\AXC{$\stackrel{\Pi_{92}}{\vdots}$}\noLine
\UIC{$\varphi_{272}$}
\BIC{$\mathrm{Seq}(a) \to \mathrm{SubIdx}(s) \to \mathrm{SubIdx}(t) \to \varphi_{123}$}
\AXC{$\stackrel{\Pi_{94}}{\vdots}$}\noLine
\UIC{$\varphi_{424}$}
\BIC{$\varphi_{357}$}
\BIC{$\mathrm{Seq}(a) \to \mathrm{SubIdx}(s) \to \varphi_{128}$}
\UIC{$\forall \alpha\,(\mathrm{Seq}(a) \to \mathrm{SubIdx}(s) \to \varphi_{128})$}
\UIC{$\forall t\,\forall \alpha\,(\mathrm{Seq}(a) \to \mathrm{SubIdx}(s) \to \varphi_{128})$}
\UIC{$\forall s\,\forall t\,\forall \alpha\,(\mathrm{Seq}(a) \to \mathrm{SubIdx}(s) \to \varphi_{128})$}
\UIC{$\forall a\,\forall s\,\forall t\,\forall \alpha\,(\mathrm{Seq}(a) \to \mathrm{SubIdx}(s) \to \varphi_{128})$}
\end{ptree}
\begin{abbrevlist}
\item $\varphi_{1} :\equiv \mathrm{Seq}(a) \to \mathrm{SubIdx}(s) \to \mathrm{SubIdx}(t) \to \mathrm{SubIdx}(t)$
\item $\varphi_{2} :\equiv \mathrm{Conv}(a \circ s,\alpha) \to \mathrm{Conv}(a \circ t,\alpha) \to \mathrm{NSeq}(t)$
\item $\varphi_{3} :\equiv \mathrm{Conv}(a \circ s,\alpha) \to \mathrm{Conv}(a \circ t,\alpha) \to \mathrm{NSeq}(s)$
\item $\varphi_{4} :\equiv \mathrm{Conv}(a \circ s,\alpha) \to \mathrm{Conv}(a \circ t,\alpha) \to \mathrm{Conv}(a,\alpha)$
\item $\varphi_{5} :\equiv \mathrm{SubIdx}(t) \leftrightarrow \mathrm{NSeq}(t) \land \forall k \in \mathbb{N}\,(t(k) < t(S(k)))$
\item $\varphi_{6} :\equiv \mathrm{SubIdx}(s) \leftrightarrow \mathrm{NSeq}(s) \land \forall k \in \mathbb{N}\,(s_{k} < s_{S(k)})$
\item $\varphi_{7} :\equiv \mathrm{Conv}(a \circ s,\alpha) \to \mathrm{Conv}(a \circ t,\alpha) \to \mathrm{Conv}(a \circ t,\alpha)$
\item $\varphi_{8} :\equiv k \in \mathbb{N} \land (n = s_{k} \lor n = t(k)) \to n = t(k) \to n = t(k)$
\item $\varphi_{9} :\equiv k \in \mathbb{N} \land (n = s_{k} \lor n = t(k)) \to n = t(k) \to t(k) = n$
\item $\varphi_{10} :\equiv k \in \mathbb{N} \land (n = s_{k} \lor n = t(k)) \to n = t(k) \to t(N_{2}) \leq n$
\item $\varphi_{11} :\equiv k \in \mathbb{N} \land (n = s_{k} \lor n = t(k)) \to n = t(k) \to N_{2} \leq k \to N_{2} \leq k$
\item $\varphi_{12} :\equiv k \in \mathbb{N} \land (n = s_{k} \lor n = t(k)) \to n = s_{k} \to N_{1} \leq k \to N_{1} \leq k$
\item $\varphi_{13} :\equiv \forall N\,(\neg (N \in \mathbb{N} \land \forall n \in \mathbb{N}\,(N \leq n \to \lvert a_{n} - \alpha \rvert < \varepsilon)))$
\item $\varphi_{14} :\equiv \forall x\,\forall y\,\forall z\,(x \leq y \to y \leq z \to x \leq z) \to \forall y\,\forall z\,(t(N_{2}) \leq y \to y \leq z \to t(N_{2}) \leq z)$
\item $\varphi_{15} :\equiv \forall x\,\forall y\,\forall z\,(x \leq y \to y \leq z \to x \leq z) \to \forall y\,\forall z\,(s_{N_{1}} \leq y \to y \leq z \to s_{N_{1}} \leq z)$
\item $\varphi_{16} :\equiv k \in \mathbb{N} \land (n = s_{k} \lor n = t(k)) \to n = t(k) \to N_{2} \leq k \lor k < N_{2}$
\item $\varphi_{17} :\equiv k \in \mathbb{N} \land (n = s_{k} \lor n = t(k)) \to n = s_{k} \to N_{1} \leq k \lor k < N_{1}$
\item $\varphi_{18} :\equiv \forall n\,(s_{N_{1}} \leq s_{N_{1}} \cup n) \to s_{N_{1}} \leq s_{N_{1}} \cup t(N_{2})$
\item $\varphi_{19} :\equiv n = t(k) \to k < N_{2} \to \forall n \in \mathbb{N}\,(k < n \to t(k) < t(n))$
\item $\varphi_{20} :\equiv n = t(k) \to k < N_{2} \to N_{2} \in \mathbb{N} \to k < N_{2} \to t(k) < t(N_{2})$
\item $\varphi_{21} :\equiv k \in \mathbb{N} \land (n = s_{k} \lor n = t(k)) \to n = t(k) \to k < N_{2} \to n < n$
\item $\varphi_{22} :\equiv k \in \mathbb{N} \land (n = s_{k} \lor n = t(k)) \to n = s_{k} \to k < N_{1} \to n < n$
\item $\varphi_{23} :\equiv s_{N_{1}} \cup t(N_{2}) \leq n \to \exists k \in \mathbb{N}\,(n = s_{k} \lor n = t(k))$
\item $\varphi_{24} :\equiv \forall z\,\forall w\,(z = w \to z < n \to w < n) \to \forall w\,(t(k) = w \to t(k) < n \to w < n)$
\item $\varphi_{25} :\equiv \mathrm{SubIdx}(t) \to k \in \mathbb{N} \to \forall n \in \mathbb{N}\,(k < n \to t(k) < t(n))$
\item $\varphi_{26} :\equiv \forall z\,\forall w\,(z = w \to z < n \to w < n) \to \forall w\,(s_{k} = w \to s_{k} < n \to w < n)$
\item $\varphi_{27} :\equiv \mathrm{SubIdx}(s) \to k \in \mathbb{N} \to \forall n \in \mathbb{N}\,(k < n \to s_{k} < s_{n})$
\item $\varphi_{28} :\equiv \neg (N \in \mathbb{N} \land \forall n \in \mathbb{N}\,(N \leq n \to \lvert a \circ s(n) - \alpha \rvert < \varepsilon))$
\item $\varphi_{29} :\equiv \neg (N \in \mathbb{N} \land \forall n \in \mathbb{N}\,(N \leq n \to \lvert a \circ t(n) - \alpha \rvert < \varepsilon))$
\item $\varphi_{30} :\equiv k \in \mathbb{N} \land (n = s_{k} \lor n = t(k)) \to n = t(k) \to k < N_{2} \to k < N_{2}$
\item $\varphi_{31} :\equiv k \in \mathbb{N} \land (n = s_{k} \lor n = t(k)) \to n = s_{k} \to k < N_{1} \to k < N_{1}$
\item $\varphi_{32} :\equiv \mathrm{Conv}(a \circ s,\alpha) \to \mathrm{Conv}(a \circ t,\alpha) \to \varepsilon \in \mathbb{R} \to \varepsilon \in \mathbb{R}$
\item $\varphi_{33} :\equiv \forall a\,\forall \alpha\,(\mathrm{Conv}(a,\alpha) \to \alpha \in \mathbb{R}) \to \forall \alpha\,(\mathrm{Conv}(a \circ s,\alpha) \to \alpha \in \mathbb{R})$
\item $\varphi_{34} :\equiv \neg (N_{2} \in \mathbb{N} \land \forall n \in \mathbb{N}\,(N_{2} \leq n \to \lvert a \circ t(n) - \alpha \rvert < \varepsilon))$
\item $\varphi_{35} :\equiv \neg (N_{1} \in \mathbb{N} \land \forall n \in \mathbb{N}\,(N_{1} \leq n \to \lvert a \circ s(n) - \alpha \rvert < \varepsilon))$
\item $\varphi_{36} :\equiv \forall m\,(\mathrm{SubIdx}(t) \to m \in \mathbb{N} \to \forall n \in \mathbb{N}\,(m < n \to t(m) < t(n)))$
\item $\varphi_{37} :\equiv \forall m\,(\mathrm{SubIdx}(s) \to m \in \mathbb{N} \to \forall n \in \mathbb{N}\,(m < n \to s_{m} < s_{n}))$
\item $\varphi_{38} :\equiv \forall \alpha\,(\mathrm{Conv}(a \circ s,\alpha) \to \alpha \in \mathbb{R}) \to \mathrm{Conv}(a \circ s,\alpha) \to \alpha \in \mathbb{R}$
\item $\varphi_{39} :\equiv k \in \mathbb{N} \land (n = s_{k} \lor n = t(k)) \to k \in \mathbb{N} \land (n = s_{k} \lor n = t(k))$
\item $\varphi_{40} :\equiv n = t(k) \to N_{2} \leq k \to k \in \mathbb{N} \to N_{2} \leq k \to \lvert a \circ t(k) - \alpha \rvert < \varepsilon$
\item $\varphi_{41} :\equiv n = s_{k} \to N_{1} \leq k \to k \in \mathbb{N} \to N_{1} \leq k \to \lvert a \circ s(k) - \alpha \rvert < \varepsilon$
\item $\varphi_{42} :\equiv \forall x\,\forall y\,\forall z\,(x < y \to y \leq z \to x < z) \to \forall y\,\forall z\,(t(k) < y \to y \leq z \to t(k) < z)$
\item $\varphi_{43} :\equiv \forall x\,\forall y\,\forall z\,(x < y \to y \leq z \to x < z) \to \forall y\,\forall z\,(s_{k} < y \to y \leq z \to s_{k} < z)$
\item $\varphi_{44} :\equiv k \in \mathbb{N} \land (n = s_{k} \lor n = t(k)) \to n = t(k) \to k < N_{2} \to t(k) < n$
\item $\varphi_{45} :\equiv \mathrm{NSeq}(t) \to \mathrm{Seq}(a \circ t) \land \forall k \in \mathbb{N}\,(a \circ t(k) = a_{t(k)})$
\item $\varphi_{46} :\equiv k \in \mathbb{N} \land (n = s_{k} \lor n = t(k)) \to n = s_{k} \to k < N_{1} \to s_{k} < n$
\item $\varphi_{47} :\equiv \mathrm{NSeq}(s) \to \mathrm{Seq}(a \circ s) \land \forall k \in \mathbb{N}\,(a \circ s(k) = a_{s_{k}})$
\item $\varphi_{48} :\equiv \forall w\,(t(k) = w \to t(k) < n \to w < n) \to t(k) = n \to t(k) < n \to n < n$
\item $\varphi_{49} :\equiv \forall w\,(s_{k} = w \to s_{k} < n \to w < n) \to s_{k} = n \to s_{k} < n \to n < n$
\item $\varphi_{50} :\equiv s_{N_{1}} \cup t(N_{2}) \leq n \to k \in \mathbb{N} \land (n = s_{k} \lor n = t(k)) \to k \in \mathbb{N}$
\item $\varphi_{51} :\equiv n = t(k) \to N_{2} \leq k \to \forall n \in \mathbb{N}\,(N_{2} \leq n \to \lvert a \circ t(n) - \alpha \rvert < \varepsilon)$
\item $\varphi_{52} :\equiv n = s_{k} \to N_{1} \leq k \to \forall n \in \mathbb{N}\,(N_{1} \leq n \to \lvert a \circ s(n) - \alpha \rvert < \varepsilon)$
\item $\varphi_{53} :\equiv \mathrm{Conv}(a \circ s,\alpha) \to \mathrm{Conv}(a \circ t,\alpha) \to \varepsilon \in \mathbb{R} \to 0 < \varepsilon \to 0 < \varepsilon$
\item $\varphi_{54} :\equiv t(N_{2}) \leq s_{N_{1}} \cup t(N_{2}) \to s_{N_{1}} \cup t(N_{2}) \leq n \to t(N_{2}) \leq n$
\item $\varphi_{55} :\equiv s_{N_{1}} \leq s_{N_{1}} \cup t(N_{2}) \to s_{N_{1}} \cup t(N_{2}) \leq n \to s_{N_{1}} \leq n$
\item $\varphi_{56} :\equiv k \in \mathbb{N} \land (n = s_{k} \lor n = t(k)) \to n = t(k) \to \lvert a_{n} - \alpha \rvert < \varepsilon$
\item $\varphi_{57} :\equiv k \in \mathbb{N} \land (n = s_{k} \lor n = t(k)) \to n = s_{k} \to \lvert a_{n} - \alpha \rvert < \varepsilon$
\item $\varphi_{58} :\equiv \forall z\,(t(N_{2}) \leq s_{N_{1}} \cup t(N_{2}) \to s_{N_{1}} \cup t(N_{2}) \leq z \to t(N_{2}) \leq z)$
\item $\varphi_{59} :\equiv \forall w\,(a \circ t(k) = w \to \lvert a \circ t(k) - \alpha \rvert < \varepsilon \to \lvert w - \alpha \rvert < \varepsilon)$
\item $\varphi_{60} :\equiv \forall z\,(s_{N_{1}} \leq s_{N_{1}} \cup t(N_{2}) \to s_{N_{1}} \cup t(N_{2}) \leq z \to s_{N_{1}} \leq z)$
\item $\varphi_{61} :\equiv \forall w\,(a \circ s(k) = w \to \lvert a \circ s(k) - \alpha \rvert < \varepsilon \to \lvert w - \alpha \rvert < \varepsilon)$
\item $\varphi_{62} :\equiv 0 < \varepsilon \to \exists N \in \mathbb{N}\,\forall n \in \mathbb{N}\,(N \leq n \to \lvert a \circ t(n) - \alpha \rvert < \varepsilon)$
\item $\varphi_{63} :\equiv 0 < \varepsilon \to \exists N \in \mathbb{N}\,\forall n \in \mathbb{N}\,(N \leq n \to \lvert a \circ s(n) - \alpha \rvert < \varepsilon)$
\item $\varphi_{64} :\equiv N_{2} \in \mathbb{N} \land \forall n \in \mathbb{N}\,(N_{2} \leq n \to \lvert a \circ t(n) - \alpha \rvert < \varepsilon) \to N_{2} \in \mathbb{N}$
\item $\varphi_{65} :\equiv k \in \mathbb{N} \land (n = s_{k} \lor n = t(k)) \to n = t(k) \to k < N_{2} \to t(k) < t(N_{2})$
\item $\varphi_{66} :\equiv N_{2} \in \mathbb{N} \land \forall n \in \mathbb{N}\,(N_{2} \leq n \to \lvert a \circ t(n) - \alpha \rvert < \varepsilon) \to N_{1} \in \mathbb{N}$
\item $\varphi_{67} :\equiv k \in \mathbb{N} \land (n = s_{k} \lor n = t(k)) \to n = s_{k} \to k < N_{1} \to s_{k} < s_{N_{1}}$
\item $\varphi_{68} :\equiv \forall k\,(\mathrm{NSeq}(t) \to k \in \mathbb{N} \to t(k) \in \mathbb{N}) \to \mathrm{NSeq}(t) \to N_{2} \in \mathbb{N} \to t(N_{2}) \in \mathbb{N}$
\item $\varphi_{69} :\equiv \forall k\,(\mathrm{NSeq}(s) \to k \in \mathbb{N} \to s_{k} \in \mathbb{N}) \to \mathrm{NSeq}(s) \to N_{1} \in \mathbb{N} \to s_{N_{1}} \in \mathbb{N}$
\item $\varphi_{70} :\equiv k \in \mathbb{N} \land (n = s_{k} \lor n = t(k)) \to n = t(k) \to N_{2} \leq k \to \lvert a_{n} - \alpha \rvert < \varepsilon$
\item $\varphi_{71} :\equiv k \in \mathbb{N} \land (n = s_{k} \lor n = t(k)) \to n = s_{k} \to N_{1} \leq k \to \lvert a_{n} - \alpha \rvert < \varepsilon$
\item $\varphi_{72} :\equiv \mathrm{Conv}(a \circ t,\alpha) \to \mathrm{Seq}(a \circ t) \land \forall k \in \mathbb{N}\,(a \circ t(k) = a_{t(k)})$
\item $\varphi_{73} :\equiv \mathrm{Conv}(a \circ t,\alpha) \to \mathrm{Seq}(a \circ s) \land \forall k \in \mathbb{N}\,(a \circ s(k) = a_{s_{k}})$
\item $\varphi_{74} :\equiv \mathrm{Conv}(a \circ s,\alpha) \to \mathrm{Conv}(a \circ t,\alpha) \to k \in \mathbb{N} \to a \circ t(k) = a_{t(k)}$
\item $\varphi_{75} :\equiv \mathrm{Conv}(a \circ s,\alpha) \to \mathrm{Conv}(a \circ t,\alpha) \to k \in \mathbb{N} \to a \circ s(k) = a_{s_{k}}$
\item $\varphi_{76} :\equiv k \in \mathbb{N} \land (n = s_{k} \lor n = t(k)) \to n = t(k) \to k < N_{2} \to \lvert a_{n} - \alpha \rvert < \varepsilon$
\item $\varphi_{77} :\equiv k \in \mathbb{N} \land (n = s_{k} \lor n = t(k)) \to n = s_{k} \to k < N_{1} \to \lvert a_{n} - \alpha \rvert < \varepsilon$
\item $\varphi_{78} :\equiv \forall s\,\forall k\,(\mathrm{NSeq}(s) \to k \in \mathbb{N} \to s_{k} \in \mathbb{N}) \to \forall k\,(\mathrm{NSeq}(t) \to k \in \mathbb{N} \to t(k) \in \mathbb{N})$
\item $\varphi_{79} :\equiv \forall s\,\forall k\,(\mathrm{NSeq}(s) \to k \in \mathbb{N} \to s_{k} \in \mathbb{N}) \to \forall k\,(\mathrm{NSeq}(s) \to k \in \mathbb{N} \to s_{k} \in \mathbb{N})$
\item $\varphi_{80} :\equiv \varepsilon \in \mathbb{R} \to 0 < \varepsilon \to \exists N \in \mathbb{N}\,\forall n \in \mathbb{N}\,(N \leq n \to \lvert a_{n} - \alpha \rvert < \varepsilon)$
\item $\varphi_{81} :\equiv \mathrm{Conv}(a \circ s,\alpha) \to \mathrm{Conv}(a \circ t,\alpha) \to \forall k \in \mathbb{N}\,(a \circ t(k) = a_{t(k)})$
\item $\varphi_{82} :\equiv \mathrm{Conv}(a \circ s,\alpha) \to \mathrm{Conv}(a \circ t,\alpha) \to \forall k \in \mathbb{N}\,(a \circ s(k) = a_{s_{k}})$
\item $\varphi_{83} :\equiv N_{2} \in \mathbb{N} \land \forall n \in \mathbb{N}\,(N_{2} \leq n \to \lvert a \circ t(n) - \alpha \rvert < \varepsilon) \to t(N_{2}) \in \mathbb{N}$
\item $\varphi_{84} :\equiv N_{2} \in \mathbb{N} \land \forall n \in \mathbb{N}\,(N_{2} \leq n \to \lvert a \circ t(n) - \alpha \rvert < \varepsilon) \to s_{N_{1}} \in \mathbb{N}$
\item $\varphi_{85} :\equiv k \in \mathbb{N} \land (n = s_{k} \lor n = t(k)) \to n = t(k) \to N_{2} \leq k \to a \circ t(k) = a_{t(k)}$
\item $\varphi_{86} :\equiv s_{N_{1}} \cup t(N_{2}) \leq n \to k \in \mathbb{N} \land (n = s_{k} \lor n = t(k)) \to n = s_{k} \to n = s_{k}$
\item $\varphi_{87} :\equiv s_{N_{1}} \cup t(N_{2}) \leq n \to k \in \mathbb{N} \land (n = s_{k} \lor n = t(k)) \to n = s_{k} \to s_{k} = n$
\item $\varphi_{88} :\equiv k \in \mathbb{N} \land (n = s_{k} \lor n = t(k)) \to n = s_{k} \to N_{1} \leq k \to a \circ s(k) = a_{s_{k}}$
\item $\varphi_{89} :\equiv \forall n \in \mathbb{N}\,(k < n \to t(k) < t(n)) \to N_{2} \in \mathbb{N} \to k < N_{2} \to t(k) < t(N_{2})$
\item $\varphi_{90} :\equiv \forall n \in \mathbb{N}\,(k < n \to s_{k} < s_{n}) \to N_{1} \in \mathbb{N} \to k < N_{1} \to s_{k} < s_{N_{1}}$
\item $\varphi_{91} :\equiv k \in \mathbb{N} \land (n = s_{k} \lor n = t(k)) \to n = t(k) \to N_{2} \leq k \to \lvert a_{t(k)} - \alpha \rvert < \varepsilon$
\item $\varphi_{92} :\equiv s_{N_{1}} \cup t(N_{2}) \leq n \to k \in \mathbb{N} \land (n = s_{k} \lor n = t(k)) \to n = t(k) \to k \in \omega$
\item $\varphi_{93} :\equiv s_{N_{1}} \cup t(N_{2}) \leq n \to k \in \mathbb{N} \land (n = s_{k} \lor n = t(k)) \to n = s_{k} \to s_{N_{1}} \leq n$
\item $\varphi_{94} :\equiv k \in \mathbb{N} \land (n = s_{k} \lor n = t(k)) \to n = s_{k} \to N_{1} \leq k \to \lvert a_{s_{k}} - \alpha \rvert < \varepsilon$
\item $\varphi_{95} :\equiv s_{N_{1}} \cup t(N_{2}) \leq n \to k \in \mathbb{N} \land (n = s_{k} \lor n = t(k)) \to n = s_{k} \to k \in \omega$
\item $\varphi_{96} :\equiv \forall n\,(N_{2} \in \omega \to n \in \omega \to N_{2} \leq n \lor n < N_{2}) \to N_{2} \in \omega \to k \in \omega \to N_{2} \leq k \lor k < N_{2}$
\item $\varphi_{97} :\equiv \forall n\,(N_{1} \in \omega \to n \in \omega \to N_{1} \leq n \lor n < N_{1}) \to N_{1} \in \omega \to k \in \omega \to N_{1} \leq k \lor k < N_{1}$
\item $\varphi_{98} :\equiv k \in \mathbb{N} \land (n = s_{k} \lor n = t(k)) \to n = t(k) \to N_{2} \leq k \to \lvert a \circ t(k) - \alpha \rvert < \varepsilon$
\item $\varphi_{99} :\equiv s_{N_{1}} \cup t(N_{2}) \leq n \to k \in \mathbb{N} \land (n = s_{k} \lor n = t(k)) \to n = t(k) \to N_{2} \in \omega$
\item $\varphi_{100} :\equiv k \in \mathbb{N} \land (n = s_{k} \lor n = t(k)) \to n = s_{k} \to N_{1} \leq k \to \lvert a \circ s(k) - \alpha \rvert < \varepsilon$
\item $\varphi_{101} :\equiv s_{N_{1}} \cup t(N_{2}) \leq n \to k \in \mathbb{N} \land (n = s_{k} \lor n = t(k)) \to n = s_{k} \to N_{1} \in \omega$
\item $\varphi_{102} :\equiv \forall \varepsilon \in \mathbb{R}\,(0 < \varepsilon \to \exists N \in \mathbb{N}\,\forall n \in \mathbb{N}\,(N \leq n \to \lvert a_{n} - \alpha \rvert < \varepsilon))$
\item $\varphi_{103} :\equiv s_{N_{1}} \cup t(N_{2}) \leq n \to k \in \mathbb{N} \land (n = s_{k} \lor n = t(k)) \to \lvert a_{n} - \alpha \rvert < \varepsilon$
\item $\varphi_{104} :\equiv N_{2} \in \mathbb{N} \land \forall n \in \mathbb{N}\,(N_{2} \leq n \to \lvert a \circ t(n) - \alpha \rvert < \varepsilon) \to n \in \mathbb{N} \to n \in \mathbb{N}$
\item $\varphi_{105} :\equiv \forall k \in \mathbb{N}\,(a \circ t(k) = a_{t(k)}) \to k \in \mathbb{N} \to a \circ t(k) = a_{t(k)}$
\item $\varphi_{106} :\equiv \forall k \in \mathbb{N}\,(a \circ s(k) = a_{s_{k}}) \to k \in \mathbb{N} \to a \circ s(k) = a_{s_{k}}$
\item $\varphi_{107} :\equiv \forall m\,\forall n\,(m \in \omega \to n \in \omega \to m \leq n \lor n < m) \to \forall n\,(N_{2} \in \omega \to n \in \omega \to N_{2} \leq n \lor n < N_{2})$
\item $\varphi_{108} :\equiv \forall m\,\forall n\,(m \in \omega \to n \in \omega \to m \leq n \lor n < m) \to \forall n\,(N_{1} \in \omega \to n \in \omega \to N_{1} \leq n \lor n < N_{1})$
\item $\varphi_{109} :\equiv \forall n \in \mathbb{N}\,\exists k \in \mathbb{N}\,(n = s_{k} \lor n = t(k)) \to \mathrm{Conv}(a \circ s,\alpha) \to \mathrm{Conv}(a \circ s,\alpha)$
\item $\varphi_{110} :\equiv \forall y\,\forall z\,(t(k) < y \to y \leq z \to t(k) < z) \to \forall z\,(t(k) < t(N_{2}) \to t(N_{2}) \leq z \to t(k) < z)$
\item $\varphi_{111} :\equiv \forall y\,\forall z\,(s_{k} < y \to y \leq z \to s_{k} < z) \to \forall z\,(s_{k} < s_{N_{1}} \to s_{N_{1}} \leq z \to s_{k} < z)$
\item $\varphi_{112} :\equiv N_{2} \in \mathbb{N} \land \forall n \in \mathbb{N}\,(N_{2} \leq n \to \lvert a \circ t(n) - \alpha \rvert < \varepsilon) \to s_{N_{1}} \cup t(N_{2}) \in \mathbb{N}$
\item $\varphi_{113} :\equiv a \circ t(k) = a_{t(k)} \to \lvert a \circ t(k) - \alpha \rvert < \varepsilon \to \lvert a_{t(k)} - \alpha \rvert < \varepsilon$
\item $\varphi_{114} :\equiv a \circ s(k) = a_{s_{k}} \to \lvert a \circ s(k) - \alpha \rvert < \varepsilon \to \lvert a_{s_{k}} - \alpha \rvert < \varepsilon$
\item $\varphi_{115} :\equiv k \in \mathbb{N} \land (n = s_{k} \lor n = t(k)) \to n = s_{k} \to k < N_{1} \to \forall n \in \mathbb{N}\,(k < n \to s_{k} < s_{n})$
\item $\varphi_{116} :\equiv k \in \mathbb{N} \land (n = s_{k} \lor n = t(k)) \to n = s_{k} \to k < N_{1} \to N_{1} \in \mathbb{N} \to k < N_{1} \to s_{k} < s_{N_{1}}$
\item $\varphi_{117} :\equiv s_{N_{1}} \cup t(N_{2}) \in \mathbb{N} \land \forall n \in \mathbb{N}\,(s_{N_{1}} \cup t(N_{2}) \leq n \to \lvert a_{n} - \alpha \rvert < \varepsilon)$
\item $\varphi_{118} :\equiv \forall n \in \mathbb{N}\,\exists k \in \mathbb{N}\,(n = s_{k} \lor n = t(k)) \to n \in \mathbb{N} \to \exists k \in \mathbb{N}\,(n = s_{k} \lor n = t(k))$
\item $\varphi_{119} :\equiv \forall z\,(t(k) < t(N_{2}) \to t(N_{2}) \leq z \to t(k) < z) \to t(k) < t(N_{2}) \to t(N_{2}) \leq n \to t(k) < n$
\item $\varphi_{120} :\equiv \forall z\,(s_{k} < s_{N_{1}} \to s_{N_{1}} \leq z \to s_{k} < z) \to s_{k} < s_{N_{1}} \to s_{N_{1}} \leq n \to s_{k} < n$
\item $\varphi_{121} :\equiv \forall m\,\forall n\,(m \in \mathbb{N} \to n \in \mathbb{N} \to m \cup n \in \mathbb{N}) \to \forall n\,(s_{N_{1}} \in \mathbb{N} \to n \in \mathbb{N} \to s_{N_{1}} \cup n \in \mathbb{N})$
\item $\varphi_{122} :\equiv \forall n \in \mathbb{N}\,\exists k \in \mathbb{N}\,(n = s_{k} \lor n = t(k)) \to \forall n \in \mathbb{N}\,\exists k \in \mathbb{N}\,(n = s_{k} \lor n = t(k))$
\item $\varphi_{123} :\equiv \forall n \in \mathbb{N}\,\exists k \in \mathbb{N}\,(n = s_{k} \lor n = t(k)) \to \mathrm{Conv}(a \circ s,\alpha) \to \mathrm{Conv}(a \circ t,\alpha) \to \alpha \in \mathbb{R}$
\item $\varphi_{124} :\equiv \forall n \in \mathbb{N}\,(N_{2} \leq n \to \lvert a \circ t(n) - \alpha \rvert < \varepsilon) \to k \in \mathbb{N} \to N_{2} \leq k \to \lvert a \circ t(k) - \alpha \rvert < \varepsilon$
\item $\varphi_{125} :\equiv \forall n \in \mathbb{N}\,(N_{1} \leq n \to \lvert a \circ s(n) - \alpha \rvert < \varepsilon) \to k \in \mathbb{N} \to N_{1} \leq k \to \lvert a \circ s(k) - \alpha \rvert < \varepsilon$
\item $\varphi_{126} :\equiv \mathrm{SubIdx}(t) \to \forall n \in \mathbb{N}\,\exists k \in \mathbb{N}\,(n = s_{k} \lor n = t(k)) \to \varphi_{2}$
\item $\varphi_{127} :\equiv \mathrm{SubIdx}(t) \to \forall n \in \mathbb{N}\,\exists k \in \mathbb{N}\,(n = s_{k} \lor n = t(k)) \to \varphi_{3}$
\item $\varphi_{128} :\equiv \mathrm{SubIdx}(t) \to \forall n \in \mathbb{N}\,\exists k \in \mathbb{N}\,(n = s_{k} \lor n = t(k)) \to \varphi_{4}$
\item $\varphi_{129} :\equiv \mathrm{Seq}(a) \land \mathrm{SubIdx}(s) \land \mathrm{SubIdx}(t) \land \forall n \in \mathbb{N}\,\exists k \in \mathbb{N}\,(n = s_{k} \lor n = t(k)) \land \mathrm{Conv}(a \circ s,\alpha) \land \mathrm{Conv}(a \circ t,\alpha)$
\item $\varphi_{130} :\equiv \mathrm{SubIdx}(t) \to \forall n \in \mathbb{N}\,\exists k \in \mathbb{N}\,(n = s_{k} \lor n = t(k)) \to \varphi_{7}$
\item $\varphi_{131} :\equiv \forall n\,(s_{N_{1}} \in \mathbb{N} \to n \in \mathbb{N} \to s_{N_{1}} \cup n \in \mathbb{N}) \to s_{N_{1}} \in \mathbb{N} \to t(N_{2}) \in \mathbb{N} \to s_{N_{1}} \cup t(N_{2}) \in \mathbb{N}$
\item $\varphi_{132} :\equiv s_{N_{1}} \cup t(N_{2}) \leq n \to k \in \mathbb{N} \land (n = s_{k} \lor n = t(k)) \to \varphi_{19}$
\item $\varphi_{133} :\equiv s_{N_{1}} \cup t(N_{2}) \leq n \to k \in \mathbb{N} \land (n = s_{k} \lor n = t(k)) \to \varphi_{20}$
\item $\varphi_{134} :\equiv N_{2} \in \mathbb{N} \land \forall n \in \mathbb{N}\,(N_{2} \leq n \to \lvert a \circ t(n) - \alpha \rvert < \varepsilon) \to \exists N \in \mathbb{N}\,\forall n \in \mathbb{N}\,(N \leq n \to \lvert a_{n} - \alpha \rvert < \varepsilon)$
\item $\varphi_{135} :\equiv N_{1} \in \mathbb{N} \land \forall n \in \mathbb{N}\,(N_{1} \leq n \to \lvert a \circ s(n) - \alpha \rvert < \varepsilon) \to \exists N \in \mathbb{N}\,\forall n \in \mathbb{N}\,(N \leq n \to \lvert a_{n} - \alpha \rvert < \varepsilon)$
\item $\varphi_{136} :\equiv \mathrm{SubIdx}(t) \to \forall n \in \mathbb{N}\,\exists k \in \mathbb{N}\,(n = s_{k} \lor n = t(k)) \to \varphi_{32}$
\item $\varphi_{137} :\equiv N_{2} \in \mathbb{N} \land \forall n \in \mathbb{N}\,(N_{2} \leq n \to \lvert a \circ t(n) - \alpha \rvert < \varepsilon) \to n \in \mathbb{N} \to s_{N_{1}} \cup t(N_{2}) \leq n \to s_{N_{1}} \cup t(N_{2}) \leq n$
\item $\varphi_{138} :\equiv N_{2} \in \mathbb{N} \land \forall n \in \mathbb{N}\,(N_{2} \leq n \to \lvert a \circ t(n) - \alpha \rvert < \varepsilon) \to n \in \mathbb{N} \to s_{N_{1}} \cup t(N_{2}) \leq n \to \lvert a_{n} - \alpha \rvert < \varepsilon$
\item $\varphi_{139} :\equiv s_{N_{1}} \cup t(N_{2}) \leq n \to k \in \mathbb{N} \land (n = s_{k} \lor n = t(k)) \to \varphi_{40}$
\item $\varphi_{140} :\equiv s_{N_{1}} \cup t(N_{2}) \leq n \to k \in \mathbb{N} \land (n = s_{k} \lor n = t(k)) \to \varphi_{41}$
\item $\varphi_{141} :\equiv N_{2} \in \mathbb{N} \land \forall n \in \mathbb{N}\,(N_{2} \leq n \to \lvert a \circ t(n) - \alpha \rvert < \varepsilon) \to N_{2} \in \mathbb{N} \land \forall n \in \mathbb{N}\,(N_{2} \leq n \to \lvert a \circ t(n) - \alpha \rvert < \varepsilon)$
\item $\varphi_{142} :\equiv N_{1} \in \mathbb{N} \land \forall n \in \mathbb{N}\,(N_{1} \leq n \to \lvert a \circ s(n) - \alpha \rvert < \varepsilon) \to N_{1} \in \mathbb{N} \land \forall n \in \mathbb{N}\,(N_{1} \leq n \to \lvert a \circ s(n) - \alpha \rvert < \varepsilon)$
\item $\varphi_{143} :\equiv N_{2} \in \mathbb{N} \land \forall n \in \mathbb{N}\,(N_{2} \leq n \to \lvert a \circ t(n) - \alpha \rvert < \varepsilon) \to \forall n \in \mathbb{N}\,(s_{N_{1}} \cup t(N_{2}) \leq n \to \lvert a_{n} - \alpha \rvert < \varepsilon)$
\item $\varphi_{144} :\equiv N_{1} \in \mathbb{N} \land \forall n \in \mathbb{N}\,(N_{1} \leq n \to \lvert a \circ s(n) - \alpha \rvert < \varepsilon) \to \exists N \in \mathbb{N}\,\forall n \in \mathbb{N}\,(N \leq n \to \lvert a \circ t(n) - \alpha \rvert < \varepsilon)$
\item $\varphi_{145} :\equiv s_{N_{1}} \cup t(N_{2}) \leq n \to k \in \mathbb{N} \land (n = s_{k} \lor n = t(k)) \to \varphi_{51}$
\item $\varphi_{146} :\equiv s_{N_{1}} \cup t(N_{2}) \leq n \to k \in \mathbb{N} \land (n = s_{k} \lor n = t(k)) \to \varphi_{52}$
\item $\varphi_{147} :\equiv \mathrm{SubIdx}(t) \to \forall n \in \mathbb{N}\,\exists k \in \mathbb{N}\,(n = s_{k} \lor n = t(k)) \to \varphi_{53}$
\item $\varphi_{148} :\equiv \forall z\,\forall w\,(z = w \to \lvert a_{z} - \alpha \rvert < \varepsilon \to \lvert a_{w} - \alpha \rvert < \varepsilon) \to \forall w\,(t(k) = w \to \lvert a_{t(k)} - \alpha \rvert < \varepsilon \to \lvert a_{w} - \alpha \rvert < \varepsilon)$
\item $\varphi_{149} :\equiv \forall z\,\forall w\,(z = w \to \lvert a_{z} - \alpha \rvert < \varepsilon \to \lvert a_{w} - \alpha \rvert < \varepsilon) \to \forall w\,(s_{k} = w \to \lvert a_{s_{k}} - \alpha \rvert < \varepsilon \to \lvert a_{w} - \alpha \rvert < \varepsilon)$
\item $\varphi_{150} :\equiv N_{2} \in \mathbb{N} \land \forall n \in \mathbb{N}\,(N_{2} \leq n \to \lvert a \circ t(n) - \alpha \rvert < \varepsilon) \to n \in \mathbb{N} \to \varphi_{23}$
\item $\varphi_{151} :\equiv N_{1} \in \mathbb{N} \land \forall n \in \mathbb{N}\,(N_{1} \leq n \to \lvert a \circ s(n) - \alpha \rvert < \varepsilon) \to \varphi_{64}$
\item $\varphi_{152} :\equiv \mathrm{SubIdx}(t) \to \forall n \in \mathbb{N}\,\exists k \in \mathbb{N}\,(n = s_{k} \lor n = t(k)) \to \varphi_{74}$
\item $\varphi_{153} :\equiv N_{1} \in \mathbb{N} \land \forall n \in \mathbb{N}\,(N_{1} \leq n \to \lvert a \circ s(n) - \alpha \rvert < \varepsilon) \to \varphi_{66}$
\item $\varphi_{154} :\equiv \mathrm{SubIdx}(t) \to \forall n \in \mathbb{N}\,\exists k \in \mathbb{N}\,(n = s_{k} \lor n = t(k)) \to \varphi_{75}$
\item $\varphi_{155} :\equiv \forall w\,(t(k) = w \to \lvert a_{t(k)} - \alpha \rvert < \varepsilon \to \lvert a_{w} - \alpha \rvert < \varepsilon) \to t(k) = n \to \lvert a_{t(k)} - \alpha \rvert < \varepsilon \to \lvert a_{n} - \alpha \rvert < \varepsilon$
\item $\varphi_{156} :\equiv \forall w\,(s_{k} = w \to \lvert a_{s_{k}} - \alpha \rvert < \varepsilon \to \lvert a_{w} - \alpha \rvert < \varepsilon) \to s_{k} = n \to \lvert a_{s_{k}} - \alpha \rvert < \varepsilon \to \lvert a_{n} - \alpha \rvert < \varepsilon$
\item $\varphi_{157} :\equiv \mathrm{SubIdx}(t) \to \forall n \in \mathbb{N}\,\exists k \in \mathbb{N}\,(n = s_{k} \lor n = t(k)) \to \varphi_{81}$
\item $\varphi_{158} :\equiv \mathrm{SubIdx}(t) \to \forall n \in \mathbb{N}\,\exists k \in \mathbb{N}\,(n = s_{k} \lor n = t(k)) \to \varphi_{82}$
\item $\varphi_{159} :\equiv N_{2} \in \mathbb{N} \land \forall n \in \mathbb{N}\,(N_{2} \leq n \to \lvert a \circ t(n) - \alpha \rvert < \varepsilon) \to n \in \mathbb{N} \to \varphi_{50}$
\item $\varphi_{160} :\equiv N_{1} \in \mathbb{N} \land \forall n \in \mathbb{N}\,(N_{1} \leq n \to \lvert a \circ s(n) - \alpha \rvert < \varepsilon) \to \varphi_{83}$
\item $\varphi_{161} :\equiv N_{1} \in \mathbb{N} \land \forall n \in \mathbb{N}\,(N_{1} \leq n \to \lvert a \circ s(n) - \alpha \rvert < \varepsilon) \to \varphi_{84}$
\item $\varphi_{162} :\equiv \forall n \in \mathbb{N}\,\exists k \in \mathbb{N}\,(n = s_{k} \lor n = t(k)) \to \mathrm{Conv}(a \circ s,\alpha) \to \varphi_{72}$
\item $\varphi_{163} :\equiv \forall n \in \mathbb{N}\,\exists k \in \mathbb{N}\,(n = s_{k} \lor n = t(k)) \to \mathrm{Conv}(a \circ s,\alpha) \to \varphi_{73}$
\item $\varphi_{164} :\equiv N_{1} \in \mathbb{N} \land \forall n \in \mathbb{N}\,(N_{1} \leq n \to \lvert a \circ s(n) - \alpha \rvert < \varepsilon) \to \varphi_{104}$
\item $\varphi_{165} :\equiv N_{2} \in \mathbb{N} \land \forall n \in \mathbb{N}\,(N_{2} \leq n \to \lvert a \circ t(n) - \alpha \rvert < \varepsilon) \to n \in \mathbb{N} \to s_{N_{1}} \cup t(N_{2}) \leq n \to \varphi_{8}$
\item $\varphi_{166} :\equiv N_{2} \in \mathbb{N} \land \forall n \in \mathbb{N}\,(N_{2} \leq n \to \lvert a \circ t(n) - \alpha \rvert < \varepsilon) \to n \in \mathbb{N} \to s_{N_{1}} \cup t(N_{2}) \leq n \to \varphi_{9}$
\item $\varphi_{167} :\equiv N_{2} \in \mathbb{N} \land \forall n \in \mathbb{N}\,(N_{2} \leq n \to \lvert a \circ t(n) - \alpha \rvert < \varepsilon) \to n \in \mathbb{N} \to \varphi_{86}$
\item $\varphi_{168} :\equiv N_{2} \in \mathbb{N} \land \forall n \in \mathbb{N}\,(N_{2} \leq n \to \lvert a \circ t(n) - \alpha \rvert < \varepsilon) \to n \in \mathbb{N} \to \varphi_{87}$
\item $\varphi_{169} :\equiv N_{2} \in \mathbb{N} \land \forall n \in \mathbb{N}\,(N_{2} \leq n \to \lvert a \circ t(n) - \alpha \rvert < \varepsilon) \to n \in \mathbb{N} \to s_{N_{1}} \cup t(N_{2}) \leq n \to \varphi_{10}$
\item $\varphi_{170} :\equiv N_{2} \in \mathbb{N} \land \forall n \in \mathbb{N}\,(N_{2} \leq n \to \lvert a \circ t(n) - \alpha \rvert < \varepsilon) \to n \in \mathbb{N} \to s_{N_{1}} \cup t(N_{2}) \leq n \to \varphi_{11}$
\item $\varphi_{171} :\equiv N_{2} \in \mathbb{N} \land \forall n \in \mathbb{N}\,(N_{2} \leq n \to \lvert a \circ t(n) - \alpha \rvert < \varepsilon) \to n \in \mathbb{N} \to \varphi_{92}$
\item $\varphi_{172} :\equiv N_{2} \in \mathbb{N} \land \forall n \in \mathbb{N}\,(N_{2} \leq n \to \lvert a \circ t(n) - \alpha \rvert < \varepsilon) \to n \in \mathbb{N} \to \varphi_{93}$
\item $\varphi_{173} :\equiv N_{2} \in \mathbb{N} \land \forall n \in \mathbb{N}\,(N_{2} \leq n \to \lvert a \circ t(n) - \alpha \rvert < \varepsilon) \to n \in \mathbb{N} \to s_{N_{1}} \cup t(N_{2}) \leq n \to \varphi_{12}$
\item $\varphi_{174} :\equiv N_{2} \in \mathbb{N} \land \forall n \in \mathbb{N}\,(N_{2} \leq n \to \lvert a \circ t(n) - \alpha \rvert < \varepsilon) \to n \in \mathbb{N} \to \varphi_{95}$
\item $\varphi_{175} :\equiv N_{2} \in \mathbb{N} \land \forall n \in \mathbb{N}\,(N_{2} \leq n \to \lvert a \circ t(n) - \alpha \rvert < \varepsilon) \to n \in \mathbb{N} \to \varphi_{99}$
\item $\varphi_{176} :\equiv N_{2} \in \mathbb{N} \land \forall n \in \mathbb{N}\,(N_{2} \leq n \to \lvert a \circ t(n) - \alpha \rvert < \varepsilon) \to n \in \mathbb{N} \to \varphi_{101}$
\item $\varphi_{177} :\equiv N_{1} \in \mathbb{N} \land \forall n \in \mathbb{N}\,(N_{1} \leq n \to \lvert a \circ s(n) - \alpha \rvert < \varepsilon) \to \varphi_{112}$
\item $\varphi_{178} :\equiv N_{2} \in \mathbb{N} \land \forall n \in \mathbb{N}\,(N_{2} \leq n \to \lvert a \circ t(n) - \alpha \rvert < \varepsilon) \to n \in \mathbb{N} \to \varphi_{103}$
\item $\varphi_{179} :\equiv N_{2} \in \mathbb{N} \land \forall n \in \mathbb{N}\,(N_{2} \leq n \to \lvert a \circ t(n) - \alpha \rvert < \varepsilon) \to \varphi_{117}$
\item $\varphi_{180} :\equiv N_{2} \in \mathbb{N} \land \forall n \in \mathbb{N}\,(N_{2} \leq n \to \lvert a \circ t(n) - \alpha \rvert < \varepsilon) \to n \in \mathbb{N} \to s_{N_{1}} \cup t(N_{2}) \leq n \to \varphi_{16}$
\item $\varphi_{181} :\equiv N_{2} \in \mathbb{N} \land \forall n \in \mathbb{N}\,(N_{2} \leq n \to \lvert a \circ t(n) - \alpha \rvert < \varepsilon) \to n \in \mathbb{N} \to s_{N_{1}} \cup t(N_{2}) \leq n \to \varphi_{17}$
\item $\varphi_{182} :\equiv N_{2} \in \mathbb{N} \land \forall n \in \mathbb{N}\,(N_{2} \leq n \to \lvert a \circ t(n) - \alpha \rvert < \varepsilon) \to n \in \mathbb{N} \to s_{N_{1}} \cup t(N_{2}) \leq n \to \varphi_{21}$
\item $\varphi_{183} :\equiv N_{2} \in \mathbb{N} \land \forall n \in \mathbb{N}\,(N_{2} \leq n \to \lvert a \circ t(n) - \alpha \rvert < \varepsilon) \to n \in \mathbb{N} \to s_{N_{1}} \cup t(N_{2}) \leq n \to \varphi_{22}$
\item $\varphi_{184} :\equiv N_{2} \in \mathbb{N} \land \forall n \in \mathbb{N}\,(N_{2} \leq n \to \lvert a \circ t(n) - \alpha \rvert < \varepsilon) \to n \in \mathbb{N} \to s_{N_{1}} \cup t(N_{2}) \leq n \to \varphi_{30}$
\item $\varphi_{185} :\equiv N_{2} \in \mathbb{N} \land \forall n \in \mathbb{N}\,(N_{2} \leq n \to \lvert a \circ t(n) - \alpha \rvert < \varepsilon) \to n \in \mathbb{N} \to s_{N_{1}} \cup t(N_{2}) \leq n \to \varphi_{31}$
\item $\varphi_{186} :\equiv N_{2} \in \mathbb{N} \land \forall n \in \mathbb{N}\,(N_{2} \leq n \to \lvert a \circ t(n) - \alpha \rvert < \varepsilon) \to n \in \mathbb{N} \to s_{N_{1}} \cup t(N_{2}) \leq n \to \varphi_{39}$
\item $\varphi_{187} :\equiv \forall n \in \mathbb{N}\,\exists k \in \mathbb{N}\,(n = s_{k} \lor n = t(k)) \to \mathrm{Conv}(a \circ s,\alpha) \to \mathrm{Conv}(a \circ t,\alpha) \to \varphi_{80}$
\item $\varphi_{188} :\equiv N_{2} \in \mathbb{N} \land \forall n \in \mathbb{N}\,(N_{2} \leq n \to \lvert a \circ t(n) - \alpha \rvert < \varepsilon) \to n \in \mathbb{N} \to s_{N_{1}} \cup t(N_{2}) \leq n \to \varphi_{44}$
\item $\varphi_{189} :\equiv N_{2} \in \mathbb{N} \land \forall n \in \mathbb{N}\,(N_{2} \leq n \to \lvert a \circ t(n) - \alpha \rvert < \varepsilon) \to n \in \mathbb{N} \to s_{N_{1}} \cup t(N_{2}) \leq n \to \varphi_{46}$
\item $\varphi_{190} :\equiv N_{2} \in \mathbb{N} \land \forall n \in \mathbb{N}\,(N_{2} \leq n \to \lvert a \circ t(n) - \alpha \rvert < \varepsilon) \to n \in \mathbb{N} \to s_{N_{1}} \cup t(N_{2}) \leq n \to \varphi_{56}$
\item $\varphi_{191} :\equiv N_{2} \in \mathbb{N} \land \forall n \in \mathbb{N}\,(N_{2} \leq n \to \lvert a \circ t(n) - \alpha \rvert < \varepsilon) \to n \in \mathbb{N} \to s_{N_{1}} \cup t(N_{2}) \leq n \to \varphi_{57}$
\item $\varphi_{192} :\equiv \forall n \in \mathbb{N}\,\exists k \in \mathbb{N}\,(n = s_{k} \lor n = t(k)) \to \mathrm{Conv}(a \circ s,\alpha) \to \mathrm{Conv}(a \circ t,\alpha) \to \varphi_{102}$
\item $\varphi_{193} :\equiv \forall n \in \mathbb{N}\,\exists k \in \mathbb{N}\,(n = s_{k} \lor n = t(k)) \to \mathrm{Conv}(a \circ s,\alpha) \to \mathrm{Conv}(a \circ t,\alpha) \to \varepsilon \in \mathbb{R} \to \varphi_{63}$
\item $\varphi_{194} :\equiv N_{2} \in \mathbb{N} \land \forall n \in \mathbb{N}\,(N_{2} \leq n \to \lvert a \circ t(n) - \alpha \rvert < \varepsilon) \to n \in \mathbb{N} \to s_{N_{1}} \cup t(N_{2}) \leq n \to \varphi_{65}$
\item $\varphi_{195} :\equiv N_{2} \in \mathbb{N} \land \forall n \in \mathbb{N}\,(N_{2} \leq n \to \lvert a \circ t(n) - \alpha \rvert < \varepsilon) \to n \in \mathbb{N} \to s_{N_{1}} \cup t(N_{2}) \leq n \to \varphi_{67}$
\item $\varphi_{196} :\equiv N_{2} \in \mathbb{N} \land \forall n \in \mathbb{N}\,(N_{2} \leq n \to \lvert a \circ t(n) - \alpha \rvert < \varepsilon) \to n \in \mathbb{N} \to s_{N_{1}} \cup t(N_{2}) \leq n \to \varphi_{70}$
\item $\varphi_{197} :\equiv N_{2} \in \mathbb{N} \land \forall n \in \mathbb{N}\,(N_{2} \leq n \to \lvert a \circ t(n) - \alpha \rvert < \varepsilon) \to n \in \mathbb{N} \to s_{N_{1}} \cup t(N_{2}) \leq n \to \varphi_{71}$
\item $\varphi_{198} :\equiv N_{2} \in \mathbb{N} \land \forall n \in \mathbb{N}\,(N_{2} \leq n \to \lvert a \circ t(n) - \alpha \rvert < \varepsilon) \to n \in \mathbb{N} \to s_{N_{1}} \cup t(N_{2}) \leq n \to \varphi_{76}$
\item $\varphi_{199} :\equiv N_{2} \in \mathbb{N} \land \forall n \in \mathbb{N}\,(N_{2} \leq n \to \lvert a \circ t(n) - \alpha \rvert < \varepsilon) \to n \in \mathbb{N} \to s_{N_{1}} \cup t(N_{2}) \leq n \to \varphi_{77}$
\item $\varphi_{200} :\equiv \mathrm{Conv}(a \circ s,\alpha) \to \mathrm{Conv}(a \circ t,\alpha) \to \varepsilon \in \mathbb{R} \to 0 < \varepsilon \to \varphi_{135}$
\item $\varphi_{201} :\equiv N_{2} \in \mathbb{N} \land \forall n \in \mathbb{N}\,(N_{2} \leq n \to \lvert a \circ t(n) - \alpha \rvert < \varepsilon) \to n \in \mathbb{N} \to s_{N_{1}} \cup t(N_{2}) \leq n \to \varphi_{85}$
\item $\varphi_{202} :\equiv N_{2} \in \mathbb{N} \land \forall n \in \mathbb{N}\,(N_{2} \leq n \to \lvert a \circ t(n) - \alpha \rvert < \varepsilon) \to n \in \mathbb{N} \to s_{N_{1}} \cup t(N_{2}) \leq n \to \varphi_{88}$
\item $\varphi_{203} :\equiv N_{2} \in \mathbb{N} \land \forall n \in \mathbb{N}\,(N_{2} \leq n \to \lvert a \circ t(n) - \alpha \rvert < \varepsilon) \to n \in \mathbb{N} \to s_{N_{1}} \cup t(N_{2}) \leq n \to \varphi_{91}$
\item $\varphi_{204} :\equiv N_{2} \in \mathbb{N} \land \forall n \in \mathbb{N}\,(N_{2} \leq n \to \lvert a \circ t(n) - \alpha \rvert < \varepsilon) \to n \in \mathbb{N} \to s_{N_{1}} \cup t(N_{2}) \leq n \to \varphi_{94}$
\item $\varphi_{205} :\equiv N_{2} \in \mathbb{N} \land \forall n \in \mathbb{N}\,(N_{2} \leq n \to \lvert a \circ t(n) - \alpha \rvert < \varepsilon) \to n \in \mathbb{N} \to s_{N_{1}} \cup t(N_{2}) \leq n \to \varphi_{98}$
\item $\varphi_{206} :\equiv N_{2} \in \mathbb{N} \land \forall n \in \mathbb{N}\,(N_{2} \leq n \to \lvert a \circ t(n) - \alpha \rvert < \varepsilon) \to n \in \mathbb{N} \to s_{N_{1}} \cup t(N_{2}) \leq n \to \varphi_{100}$
\item $\varphi_{207} :\equiv N_{1} \in \mathbb{N} \land \forall n \in \mathbb{N}\,(N_{1} \leq n \to \lvert a \circ s(n) - \alpha \rvert < \varepsilon) \to \varphi_{134}$
\item $\varphi_{208} :\equiv \mathrm{Conv}(a \circ s,\alpha) \to \mathrm{Conv}(a \circ t,\alpha) \to \varepsilon \in \mathbb{R} \to 0 < \varepsilon \to \varphi_{142}$
\item $\varphi_{209} :\equiv \mathrm{Conv}(a \circ s,\alpha) \to \mathrm{Conv}(a \circ t,\alpha) \to \varepsilon \in \mathbb{R} \to 0 < \varepsilon \to \varphi_{144}$
\item $\varphi_{210} :\equiv N_{1} \in \mathbb{N} \land \forall n \in \mathbb{N}\,(N_{1} \leq n \to \lvert a \circ s(n) - \alpha \rvert < \varepsilon) \to \varphi_{137}$
\item $\varphi_{211} :\equiv N_{1} \in \mathbb{N} \land \forall n \in \mathbb{N}\,(N_{1} \leq n \to \lvert a \circ s(n) - \alpha \rvert < \varepsilon) \to \varphi_{138}$
\item $\varphi_{212} :\equiv N_{1} \in \mathbb{N} \land \forall n \in \mathbb{N}\,(N_{1} \leq n \to \lvert a \circ s(n) - \alpha \rvert < \varepsilon) \to \varphi_{141}$
\item $\varphi_{213} :\equiv N_{1} \in \mathbb{N} \land \forall n \in \mathbb{N}\,(N_{1} \leq n \to \lvert a \circ s(n) - \alpha \rvert < \varepsilon) \to \varphi_{143}$
\item $\varphi_{214} :\equiv N_{2} \in \mathbb{N} \land \forall n \in \mathbb{N}\,(N_{2} \leq n \to \lvert a \circ t(n) - \alpha \rvert < \varepsilon) \to n \in \mathbb{N} \to \varphi_{132}$
\item $\varphi_{215} :\equiv N_{2} \in \mathbb{N} \land \forall n \in \mathbb{N}\,(N_{2} \leq n \to \lvert a \circ t(n) - \alpha \rvert < \varepsilon) \to n \in \mathbb{N} \to \varphi_{133}$
\item $\varphi_{216} :\equiv N_{2} \in \mathbb{N} \land \forall n \in \mathbb{N}\,(N_{2} \leq n \to \lvert a \circ t(n) - \alpha \rvert < \varepsilon) \to n \in \mathbb{N} \to s_{N_{1}} \cup t(N_{2}) \leq n \to \varphi_{115}$
\item $\varphi_{217} :\equiv N_{2} \in \mathbb{N} \land \forall n \in \mathbb{N}\,(N_{2} \leq n \to \lvert a \circ t(n) - \alpha \rvert < \varepsilon) \to n \in \mathbb{N} \to s_{N_{1}} \cup t(N_{2}) \leq n \to \varphi_{116}$
\item $\varphi_{218} :\equiv \varphi_{129} \land \varepsilon \in \mathbb{R} \land 0 < \varepsilon \land \neg \exists N \in \mathbb{N}\,\forall n \in \mathbb{N}\,(N \leq n \to \lvert a_{n} - \alpha \rvert < \varepsilon)$
\item $\varphi_{219} :\equiv \mathrm{Conv}(a \circ s,\alpha) \to \mathrm{Conv}(a \circ t,\alpha) \to \varepsilon \in \mathbb{R} \to 0 < \varepsilon \to \varphi_{151}$
\item $\varphi_{220} :\equiv \mathrm{Conv}(a \circ s,\alpha) \to \mathrm{Conv}(a \circ t,\alpha) \to \varepsilon \in \mathbb{R} \to 0 < \varepsilon \to \varphi_{153}$
\item $\varphi_{221} :\equiv N_{1} \in \mathbb{N} \land \forall n \in \mathbb{N}\,(N_{1} \leq n \to \lvert a \circ s(n) - \alpha \rvert < \varepsilon) \to \varphi_{150}$
\item $\varphi_{222} :\equiv N_{2} \in \mathbb{N} \land \forall n \in \mathbb{N}\,(N_{2} \leq n \to \lvert a \circ t(n) - \alpha \rvert < \varepsilon) \to n \in \mathbb{N} \to \varphi_{139}$
\item $\varphi_{223} :\equiv N_{2} \in \mathbb{N} \land \forall n \in \mathbb{N}\,(N_{2} \leq n \to \lvert a \circ t(n) - \alpha \rvert < \varepsilon) \to n \in \mathbb{N} \to \varphi_{140}$
\item $\varphi_{224} :\equiv N_{2} \in \mathbb{N} \land \forall n \in \mathbb{N}\,(N_{2} \leq n \to \lvert a \circ t(n) - \alpha \rvert < \varepsilon) \to n \in \mathbb{N} \to \varphi_{145}$
\item $\varphi_{225} :\equiv N_{2} \in \mathbb{N} \land \forall n \in \mathbb{N}\,(N_{2} \leq n \to \lvert a \circ t(n) - \alpha \rvert < \varepsilon) \to n \in \mathbb{N} \to \varphi_{146}$
\item $\varphi_{226} :\equiv \mathrm{Conv}(a \circ s,\alpha) \to \mathrm{Conv}(a \circ t,\alpha) \to \varepsilon \in \mathbb{R} \to 0 < \varepsilon \to \varphi_{160}$
\item $\varphi_{227} :\equiv \mathrm{Conv}(a \circ s,\alpha) \to \mathrm{Conv}(a \circ t,\alpha) \to \varepsilon \in \mathbb{R} \to 0 < \varepsilon \to \varphi_{161}$
\item $\varphi_{228} :\equiv N_{1} \in \mathbb{N} \land \forall n \in \mathbb{N}\,(N_{1} \leq n \to \lvert a \circ s(n) - \alpha \rvert < \varepsilon) \to \varphi_{159}$
\item $\varphi_{229} :\equiv \mathrm{Conv}(a \circ s,\alpha) \to \mathrm{Conv}(a \circ t,\alpha) \to \varepsilon \in \mathbb{R} \to 0 < \varepsilon \to \varphi_{164}$
\item $\varphi_{230} :\equiv (\mathrm{Seq}(a) \to \mathrm{SubIdx}(s) \to \mathrm{SubIdx}(s)) \to \varphi_{6} \to \mathrm{Seq}(a) \to \mathrm{SubIdx}(s) \to \varphi_{127}$
\item $\varphi_{231} :\equiv \forall a\,\forall \alpha\,\forall \varepsilon\,(\mathrm{Conv}(a,\alpha) \to \varphi_{80}) \to \forall \alpha\,\forall \varepsilon\,(\mathrm{Conv}(a \circ t,\alpha) \to \varepsilon \in \mathbb{R} \to \varphi_{62})$
\item $\varphi_{232} :\equiv \forall \varepsilon\,(\mathrm{Conv}(a \circ t,\alpha) \to \varepsilon \in \mathbb{R} \to \varphi_{62}) \to \mathrm{Conv}(a \circ t,\alpha) \to \varepsilon \in \mathbb{R} \to \varphi_{62}$
\item $\varphi_{233} :\equiv \forall a\,\forall \alpha\,\forall \varepsilon\,(\mathrm{Conv}(a,\alpha) \to \varphi_{80}) \to \forall \alpha\,\forall \varepsilon\,(\mathrm{Conv}(a \circ s,\alpha) \to \varepsilon \in \mathbb{R} \to \varphi_{63})$
\item $\varphi_{234} :\equiv \forall \varepsilon\,(\mathrm{Conv}(a \circ s,\alpha) \to \varepsilon \in \mathbb{R} \to \varphi_{63}) \to \mathrm{Conv}(a \circ s,\alpha) \to \varepsilon \in \mathbb{R} \to \varphi_{63}$
\item $\varphi_{235} :\equiv \mathrm{Conv}(a \circ s,\alpha) \to \mathrm{Conv}(a \circ t,\alpha) \to \varepsilon \in \mathbb{R} \to 0 < \varepsilon \to \varphi_{177}$
\item $\varphi_{236} :\equiv N_{1} \in \mathbb{N} \land \forall n \in \mathbb{N}\,(N_{1} \leq n \to \lvert a \circ s(n) - \alpha \rvert < \varepsilon) \to \varphi_{165}$
\item $\varphi_{237} :\equiv N_{1} \in \mathbb{N} \land \forall n \in \mathbb{N}\,(N_{1} \leq n \to \lvert a \circ s(n) - \alpha \rvert < \varepsilon) \to \varphi_{166}$
\item $\varphi_{238} :\equiv N_{1} \in \mathbb{N} \land \forall n \in \mathbb{N}\,(N_{1} \leq n \to \lvert a \circ s(n) - \alpha \rvert < \varepsilon) \to \varphi_{167}$
\item $\varphi_{239} :\equiv N_{1} \in \mathbb{N} \land \forall n \in \mathbb{N}\,(N_{1} \leq n \to \lvert a \circ s(n) - \alpha \rvert < \varepsilon) \to \varphi_{168}$
\item $\varphi_{240} :\equiv N_{1} \in \mathbb{N} \land \forall n \in \mathbb{N}\,(N_{1} \leq n \to \lvert a \circ s(n) - \alpha \rvert < \varepsilon) \to \varphi_{169}$
\item $\varphi_{241} :\equiv N_{1} \in \mathbb{N} \land \forall n \in \mathbb{N}\,(N_{1} \leq n \to \lvert a \circ s(n) - \alpha \rvert < \varepsilon) \to \varphi_{170}$
\item $\varphi_{242} :\equiv N_{1} \in \mathbb{N} \land \forall n \in \mathbb{N}\,(N_{1} \leq n \to \lvert a \circ s(n) - \alpha \rvert < \varepsilon) \to \varphi_{171}$
\item $\varphi_{243} :\equiv N_{1} \in \mathbb{N} \land \forall n \in \mathbb{N}\,(N_{1} \leq n \to \lvert a \circ s(n) - \alpha \rvert < \varepsilon) \to \varphi_{172}$
\item $\varphi_{244} :\equiv N_{1} \in \mathbb{N} \land \forall n \in \mathbb{N}\,(N_{1} \leq n \to \lvert a \circ s(n) - \alpha \rvert < \varepsilon) \to \varphi_{173}$
\item $\varphi_{245} :\equiv N_{1} \in \mathbb{N} \land \forall n \in \mathbb{N}\,(N_{1} \leq n \to \lvert a \circ s(n) - \alpha \rvert < \varepsilon) \to \varphi_{174}$
\item $\varphi_{246} :\equiv N_{1} \in \mathbb{N} \land \forall n \in \mathbb{N}\,(N_{1} \leq n \to \lvert a \circ s(n) - \alpha \rvert < \varepsilon) \to \varphi_{175}$
\item $\varphi_{247} :\equiv N_{1} \in \mathbb{N} \land \forall n \in \mathbb{N}\,(N_{1} \leq n \to \lvert a \circ s(n) - \alpha \rvert < \varepsilon) \to \varphi_{176}$
\item $\varphi_{248} :\equiv N_{1} \in \mathbb{N} \land \forall n \in \mathbb{N}\,(N_{1} \leq n \to \lvert a \circ s(n) - \alpha \rvert < \varepsilon) \to \varphi_{178}$
\item $\varphi_{249} :\equiv N_{1} \in \mathbb{N} \land \forall n \in \mathbb{N}\,(N_{1} \leq n \to \lvert a \circ s(n) - \alpha \rvert < \varepsilon) \to \varphi_{179}$
\item $\varphi_{250} :\equiv N_{1} \in \mathbb{N} \land \forall n \in \mathbb{N}\,(N_{1} \leq n \to \lvert a \circ s(n) - \alpha \rvert < \varepsilon) \to \varphi_{180}$
\item $\varphi_{251} :\equiv N_{1} \in \mathbb{N} \land \forall n \in \mathbb{N}\,(N_{1} \leq n \to \lvert a \circ s(n) - \alpha \rvert < \varepsilon) \to \varphi_{181}$
\item $\varphi_{252} :\equiv \forall \alpha\,\forall \varepsilon\,(\mathrm{Conv}(a \circ t,\alpha) \to \varepsilon \in \mathbb{R} \to \varphi_{62}) \to \forall \varepsilon\,(\mathrm{Conv}(a \circ t,\alpha) \to \varepsilon \in \mathbb{R} \to \varphi_{62})$
\item $\varphi_{253} :\equiv \forall \alpha\,\forall \varepsilon\,(\mathrm{Conv}(a \circ s,\alpha) \to \varepsilon \in \mathbb{R} \to \varphi_{63}) \to \forall \varepsilon\,(\mathrm{Conv}(a \circ s,\alpha) \to \varepsilon \in \mathbb{R} \to \varphi_{63})$
\item $\varphi_{254} :\equiv N_{1} \in \mathbb{N} \land \forall n \in \mathbb{N}\,(N_{1} \leq n \to \lvert a \circ s(n) - \alpha \rvert < \varepsilon) \to \varphi_{182}$
\item $\varphi_{255} :\equiv N_{1} \in \mathbb{N} \land \forall n \in \mathbb{N}\,(N_{1} \leq n \to \lvert a \circ s(n) - \alpha \rvert < \varepsilon) \to \varphi_{183}$
\item $\varphi_{256} :\equiv N_{1} \in \mathbb{N} \land \forall n \in \mathbb{N}\,(N_{1} \leq n \to \lvert a \circ s(n) - \alpha \rvert < \varepsilon) \to \varphi_{184}$
\item $\varphi_{257} :\equiv N_{1} \in \mathbb{N} \land \forall n \in \mathbb{N}\,(N_{1} \leq n \to \lvert a \circ s(n) - \alpha \rvert < \varepsilon) \to \varphi_{185}$
\item $\varphi_{258} :\equiv N_{1} \in \mathbb{N} \land \forall n \in \mathbb{N}\,(N_{1} \leq n \to \lvert a \circ s(n) - \alpha \rvert < \varepsilon) \to \varphi_{186}$
\item $\varphi_{259} :\equiv N_{1} \in \mathbb{N} \land \forall n \in \mathbb{N}\,(N_{1} \leq n \to \lvert a \circ s(n) - \alpha \rvert < \varepsilon) \to \varphi_{188}$
\item $\varphi_{260} :\equiv N_{1} \in \mathbb{N} \land \forall n \in \mathbb{N}\,(N_{1} \leq n \to \lvert a \circ s(n) - \alpha \rvert < \varepsilon) \to \varphi_{189}$
\item $\varphi_{261} :\equiv N_{1} \in \mathbb{N} \land \forall n \in \mathbb{N}\,(N_{1} \leq n \to \lvert a \circ s(n) - \alpha \rvert < \varepsilon) \to \varphi_{190}$
\item $\varphi_{262} :\equiv N_{1} \in \mathbb{N} \land \forall n \in \mathbb{N}\,(N_{1} \leq n \to \lvert a \circ s(n) - \alpha \rvert < \varepsilon) \to \varphi_{191}$
\item $\varphi_{263} :\equiv \forall \alpha\,(\mathrm{Conv}(a,\alpha) \leftrightarrow \mathrm{Seq}(a) \land \alpha \in \mathbb{R} \land \varphi_{102}) \to (\mathrm{Conv}(a,\alpha) \leftrightarrow \mathrm{Seq}(a) \land \alpha \in \mathbb{R} \land \varphi_{102})$
\item $\varphi_{264} :\equiv N_{1} \in \mathbb{N} \land \forall n \in \mathbb{N}\,(N_{1} \leq n \to \lvert a \circ s(n) - \alpha \rvert < \varepsilon) \to \varphi_{194}$
\item $\varphi_{265} :\equiv N_{1} \in \mathbb{N} \land \forall n \in \mathbb{N}\,(N_{1} \leq n \to \lvert a \circ s(n) - \alpha \rvert < \varepsilon) \to \varphi_{195}$
\item $\varphi_{266} :\equiv N_{1} \in \mathbb{N} \land \forall n \in \mathbb{N}\,(N_{1} \leq n \to \lvert a \circ s(n) - \alpha \rvert < \varepsilon) \to \varphi_{196}$
\item $\varphi_{267} :\equiv N_{1} \in \mathbb{N} \land \forall n \in \mathbb{N}\,(N_{1} \leq n \to \lvert a \circ s(n) - \alpha \rvert < \varepsilon) \to \varphi_{197}$
\item $\varphi_{268} :\equiv \mathrm{SubIdx}(t) \to \forall n \in \mathbb{N}\,\exists k \in \mathbb{N}\,(n = s_{k} \lor n = t(k)) \to \varphi_{200}$
\item $\varphi_{269} :\equiv \forall a\,\forall \alpha\,(\mathrm{Conv}(a,\alpha) \leftrightarrow \mathrm{Seq}(a) \land \alpha \in \mathbb{R} \land \varphi_{102}) \to \forall \alpha\,(\mathrm{Conv}(a,\alpha) \leftrightarrow \mathrm{Seq}(a) \land \alpha \in \mathbb{R} \land \varphi_{102})$
\item $\varphi_{270} :\equiv N_{1} \in \mathbb{N} \land \forall n \in \mathbb{N}\,(N_{1} \leq n \to \lvert a \circ s(n) - \alpha \rvert < \varepsilon) \to \varphi_{198}$
\item $\varphi_{271} :\equiv N_{1} \in \mathbb{N} \land \forall n \in \mathbb{N}\,(N_{1} \leq n \to \lvert a \circ s(n) - \alpha \rvert < \varepsilon) \to \varphi_{199}$
\item $\varphi_{272} :\equiv (\mathrm{Seq}(a) \to \mathrm{SubIdx}(s) \to \mathrm{SubIdx}(t) \to \varphi_{109}) \to \mathrm{Seq}(a) \to \mathrm{SubIdx}(s) \to \mathrm{SubIdx}(t) \to \varphi_{123}$
\item $\varphi_{273} :\equiv \mathrm{Conv}(a \circ s,\alpha) \to \mathrm{Conv}(a \circ t,\alpha) \to \varepsilon \in \mathbb{R} \to 0 < \varepsilon \to \varphi_{207}$
\item $\varphi_{274} :\equiv N_{1} \in \mathbb{N} \land \forall n \in \mathbb{N}\,(N_{1} \leq n \to \lvert a \circ s(n) - \alpha \rvert < \varepsilon) \to \varphi_{201}$
\item $\varphi_{275} :\equiv N_{1} \in \mathbb{N} \land \forall n \in \mathbb{N}\,(N_{1} \leq n \to \lvert a \circ s(n) - \alpha \rvert < \varepsilon) \to \varphi_{202}$
\item $\varphi_{276} :\equiv N_{1} \in \mathbb{N} \land \forall n \in \mathbb{N}\,(N_{1} \leq n \to \lvert a \circ s(n) - \alpha \rvert < \varepsilon) \to \varphi_{203}$
\item $\varphi_{277} :\equiv N_{1} \in \mathbb{N} \land \forall n \in \mathbb{N}\,(N_{1} \leq n \to \lvert a \circ s(n) - \alpha \rvert < \varepsilon) \to \varphi_{204}$
\item $\varphi_{278} :\equiv N_{1} \in \mathbb{N} \land \forall n \in \mathbb{N}\,(N_{1} \leq n \to \lvert a \circ s(n) - \alpha \rvert < \varepsilon) \to \varphi_{205}$
\item $\varphi_{279} :\equiv \mathrm{SubIdx}(t) \to \forall n \in \mathbb{N}\,\exists k \in \mathbb{N}\,(n = s_{k} \lor n = t(k)) \to \varphi_{208}$
\item $\varphi_{280} :\equiv N_{1} \in \mathbb{N} \land \forall n \in \mathbb{N}\,(N_{1} \leq n \to \lvert a \circ s(n) - \alpha \rvert < \varepsilon) \to \varphi_{206}$
\item $\varphi_{281} :\equiv \mathrm{SubIdx}(t) \to \forall n \in \mathbb{N}\,\exists k \in \mathbb{N}\,(n = s_{k} \lor n = t(k)) \to \varphi_{209}$
\item $\varphi_{282} :\equiv \mathrm{Conv}(a \circ s,\alpha) \to \mathrm{Conv}(a \circ t,\alpha) \to \varepsilon \in \mathbb{R} \to 0 < \varepsilon \to \varphi_{210}$
\item $\varphi_{283} :\equiv \mathrm{Conv}(a \circ s,\alpha) \to \mathrm{Conv}(a \circ t,\alpha) \to \varepsilon \in \mathbb{R} \to 0 < \varepsilon \to \varphi_{211}$
\item $\varphi_{284} :\equiv \mathrm{Conv}(a \circ s,\alpha) \to \mathrm{Conv}(a \circ t,\alpha) \to \varepsilon \in \mathbb{R} \to 0 < \varepsilon \to \varphi_{212}$
\item $\varphi_{285} :\equiv \mathrm{Conv}(a \circ s,\alpha) \to \mathrm{Conv}(a \circ t,\alpha) \to \varepsilon \in \mathbb{R} \to 0 < \varepsilon \to \varphi_{213}$
\item $\varphi_{286} :\equiv \mathrm{SubIdx}(t) \to \forall n \in \mathbb{N}\,\exists k \in \mathbb{N}\,(n = s_{k} \lor n = t(k)) \to \varphi_{219}$
\item $\varphi_{287} :\equiv N_{1} \in \mathbb{N} \land \forall n \in \mathbb{N}\,(N_{1} \leq n \to \lvert a \circ s(n) - \alpha \rvert < \varepsilon) \to \varphi_{214}$
\item $\varphi_{288} :\equiv N_{1} \in \mathbb{N} \land \forall n \in \mathbb{N}\,(N_{1} \leq n \to \lvert a \circ s(n) - \alpha \rvert < \varepsilon) \to \varphi_{215}$
\item $\varphi_{289} :\equiv \mathrm{SubIdx}(t) \to \forall n \in \mathbb{N}\,\exists k \in \mathbb{N}\,(n = s_{k} \lor n = t(k)) \to \varphi_{220}$
\item $\varphi_{290} :\equiv N_{1} \in \mathbb{N} \land \forall n \in \mathbb{N}\,(N_{1} \leq n \to \lvert a \circ s(n) - \alpha \rvert < \varepsilon) \to \varphi_{216}$
\item $\varphi_{291} :\equiv N_{1} \in \mathbb{N} \land \forall n \in \mathbb{N}\,(N_{1} \leq n \to \lvert a \circ s(n) - \alpha \rvert < \varepsilon) \to \varphi_{217}$
\item $\varphi_{292} :\equiv \mathrm{Conv}(a \circ s,\alpha) \to \mathrm{Conv}(a \circ t,\alpha) \to \varepsilon \in \mathbb{R} \to 0 < \varepsilon \to \varphi_{221}$
\item $\varphi_{293} :\equiv \varphi_{129} \land \varepsilon \in \mathbb{R} \land 0 < \varepsilon \land (N_{1} \in \mathbb{N} \land \forall n \in \mathbb{N}\,(N_{1} \leq n \to \lvert a \circ s(n) - \alpha \rvert < \varepsilon)) \land \neg \exists N \in \mathbb{N}\,\forall n \in \mathbb{N}\,(N \leq n \to \lvert a_{n} - \alpha \rvert < \varepsilon)$
\item $\varphi_{294} :\equiv N_{1} \in \mathbb{N} \land \forall n \in \mathbb{N}\,(N_{1} \leq n \to \lvert a \circ s(n) - \alpha \rvert < \varepsilon) \to \varphi_{222}$
\item $\varphi_{295} :\equiv N_{1} \in \mathbb{N} \land \forall n \in \mathbb{N}\,(N_{1} \leq n \to \lvert a \circ s(n) - \alpha \rvert < \varepsilon) \to \varphi_{223}$
\item $\varphi_{296} :\equiv \mathrm{SubIdx}(t) \to \forall n \in \mathbb{N}\,\exists k \in \mathbb{N}\,(n = s_{k} \lor n = t(k)) \to \varphi_{226}$
\item $\varphi_{297} :\equiv \mathrm{SubIdx}(t) \to \forall n \in \mathbb{N}\,\exists k \in \mathbb{N}\,(n = s_{k} \lor n = t(k)) \to \varphi_{227}$
\item $\varphi_{298} :\equiv \varphi_{129} \land \varepsilon \in \mathbb{R} \land 0 < \varepsilon \land (N_{1} \in \mathbb{N} \land \forall n \in \mathbb{N}\,(N_{1} \leq n \to \lvert a \circ s(n) - \alpha \rvert < \varepsilon)) \land (N_{2} \in \mathbb{N} \land \forall n \in \mathbb{N}\,(N_{2} \leq n \to \lvert a \circ t(n) - \alpha \rvert < \varepsilon))$
\item $\varphi_{299} :\equiv N_{1} \in \mathbb{N} \land \forall n \in \mathbb{N}\,(N_{1} \leq n \to \lvert a \circ s(n) - \alpha \rvert < \varepsilon) \to \varphi_{224}$
\item $\varphi_{300} :\equiv N_{1} \in \mathbb{N} \land \forall n \in \mathbb{N}\,(N_{1} \leq n \to \lvert a \circ s(n) - \alpha \rvert < \varepsilon) \to \varphi_{225}$
\item $\varphi_{301} :\equiv \mathrm{Conv}(a \circ s,\alpha) \to \mathrm{Conv}(a \circ t,\alpha) \to \varepsilon \in \mathbb{R} \to 0 < \varepsilon \to \varphi_{228}$
\item $\varphi_{302} :\equiv \mathrm{SubIdx}(t) \to \forall n \in \mathbb{N}\,\exists k \in \mathbb{N}\,(n = s_{k} \lor n = t(k)) \to \varphi_{229}$
\item $\varphi_{303} :\equiv \mathrm{SubIdx}(t) \to \forall n \in \mathbb{N}\,\exists k \in \mathbb{N}\,(n = s_{k} \lor n = t(k)) \to \varphi_{235}$
\item $\varphi_{304} :\equiv \mathrm{Conv}(a \circ s,\alpha) \to \mathrm{Conv}(a \circ t,\alpha) \to \varepsilon \in \mathbb{R} \to 0 < \varepsilon \to \varphi_{236}$
\item $\varphi_{305} :\equiv \mathrm{Conv}(a \circ s,\alpha) \to \mathrm{Conv}(a \circ t,\alpha) \to \varepsilon \in \mathbb{R} \to 0 < \varepsilon \to \varphi_{237}$
\item $\varphi_{306} :\equiv \mathrm{Conv}(a \circ s,\alpha) \to \mathrm{Conv}(a \circ t,\alpha) \to \varepsilon \in \mathbb{R} \to 0 < \varepsilon \to \varphi_{238}$
\item $\varphi_{307} :\equiv \mathrm{Conv}(a \circ s,\alpha) \to \mathrm{Conv}(a \circ t,\alpha) \to \varepsilon \in \mathbb{R} \to 0 < \varepsilon \to \varphi_{239}$
\item $\varphi_{308} :\equiv \mathrm{Conv}(a \circ s,\alpha) \to \mathrm{Conv}(a \circ t,\alpha) \to \varepsilon \in \mathbb{R} \to 0 < \varepsilon \to \varphi_{240}$
\item $\varphi_{309} :\equiv \mathrm{Conv}(a \circ s,\alpha) \to \mathrm{Conv}(a \circ t,\alpha) \to \varepsilon \in \mathbb{R} \to 0 < \varepsilon \to \varphi_{241}$
\item $\varphi_{310} :\equiv \mathrm{Conv}(a \circ s,\alpha) \to \mathrm{Conv}(a \circ t,\alpha) \to \varepsilon \in \mathbb{R} \to 0 < \varepsilon \to \varphi_{242}$
\item $\varphi_{311} :\equiv \mathrm{Conv}(a \circ s,\alpha) \to \mathrm{Conv}(a \circ t,\alpha) \to \varepsilon \in \mathbb{R} \to 0 < \varepsilon \to \varphi_{243}$
\item $\varphi_{312} :\equiv \mathrm{Conv}(a \circ s,\alpha) \to \mathrm{Conv}(a \circ t,\alpha) \to \varepsilon \in \mathbb{R} \to 0 < \varepsilon \to \varphi_{244}$
\item $\varphi_{313} :\equiv \mathrm{Conv}(a \circ s,\alpha) \to \mathrm{Conv}(a \circ t,\alpha) \to \varepsilon \in \mathbb{R} \to 0 < \varepsilon \to \varphi_{245}$
\item $\varphi_{314} :\equiv \mathrm{Conv}(a \circ s,\alpha) \to \mathrm{Conv}(a \circ t,\alpha) \to \varepsilon \in \mathbb{R} \to 0 < \varepsilon \to \varphi_{246}$
\item $\varphi_{315} :\equiv \mathrm{Conv}(a \circ s,\alpha) \to \mathrm{Conv}(a \circ t,\alpha) \to \varepsilon \in \mathbb{R} \to 0 < \varepsilon \to \varphi_{247}$
\item $\varphi_{316} :\equiv \mathrm{Conv}(a \circ s,\alpha) \to \mathrm{Conv}(a \circ t,\alpha) \to \varepsilon \in \mathbb{R} \to 0 < \varepsilon \to \varphi_{248}$
\item $\varphi_{317} :\equiv \mathrm{Conv}(a \circ s,\alpha) \to \mathrm{Conv}(a \circ t,\alpha) \to \varepsilon \in \mathbb{R} \to 0 < \varepsilon \to \varphi_{249}$
\item $\varphi_{318} :\equiv \mathrm{Conv}(a \circ s,\alpha) \to \mathrm{Conv}(a \circ t,\alpha) \to \varepsilon \in \mathbb{R} \to 0 < \varepsilon \to \varphi_{250}$
\item $\varphi_{319} :\equiv \mathrm{Conv}(a \circ s,\alpha) \to \mathrm{Conv}(a \circ t,\alpha) \to \varepsilon \in \mathbb{R} \to 0 < \varepsilon \to \varphi_{251}$
\item $\varphi_{320} :\equiv \mathrm{Conv}(a \circ s,\alpha) \to \mathrm{Conv}(a \circ t,\alpha) \to \varepsilon \in \mathbb{R} \to 0 < \varepsilon \to \varphi_{254}$
\item $\varphi_{321} :\equiv \mathrm{Conv}(a \circ s,\alpha) \to \mathrm{Conv}(a \circ t,\alpha) \to \varepsilon \in \mathbb{R} \to 0 < \varepsilon \to \varphi_{255}$
\item $\varphi_{322} :\equiv \mathrm{Conv}(a \circ s,\alpha) \to \mathrm{Conv}(a \circ t,\alpha) \to \varepsilon \in \mathbb{R} \to 0 < \varepsilon \to \varphi_{256}$
\item $\varphi_{323} :\equiv \mathrm{Conv}(a \circ s,\alpha) \to \mathrm{Conv}(a \circ t,\alpha) \to \varepsilon \in \mathbb{R} \to 0 < \varepsilon \to \varphi_{257}$
\item $\varphi_{324} :\equiv \mathrm{Conv}(a \circ s,\alpha) \to \mathrm{Conv}(a \circ t,\alpha) \to \varepsilon \in \mathbb{R} \to 0 < \varepsilon \to \varphi_{258}$
\item $\varphi_{325} :\equiv \mathrm{Conv}(a \circ s,\alpha) \to \mathrm{Conv}(a \circ t,\alpha) \to \varepsilon \in \mathbb{R} \to 0 < \varepsilon \to \varphi_{259}$
\item $\varphi_{326} :\equiv \mathrm{Conv}(a \circ s,\alpha) \to \mathrm{Conv}(a \circ t,\alpha) \to \varepsilon \in \mathbb{R} \to 0 < \varepsilon \to \varphi_{260}$
\item $\varphi_{327} :\equiv \mathrm{Conv}(a \circ s,\alpha) \to \mathrm{Conv}(a \circ t,\alpha) \to \varepsilon \in \mathbb{R} \to 0 < \varepsilon \to \varphi_{261}$
\item $\varphi_{328} :\equiv \mathrm{Conv}(a \circ s,\alpha) \to \mathrm{Conv}(a \circ t,\alpha) \to \varepsilon \in \mathbb{R} \to 0 < \varepsilon \to \varphi_{262}$
\item $\varphi_{329} :\equiv \mathrm{Conv}(a \circ s,\alpha) \to \mathrm{Conv}(a \circ t,\alpha) \to \varepsilon \in \mathbb{R} \to 0 < \varepsilon \to \varphi_{264}$
\item $\varphi_{330} :\equiv \mathrm{Conv}(a \circ s,\alpha) \to \mathrm{Conv}(a \circ t,\alpha) \to \varepsilon \in \mathbb{R} \to 0 < \varepsilon \to \varphi_{265}$
\item $\varphi_{331} :\equiv (\mathrm{Seq}(a) \to \mathrm{SubIdx}(s) \to \varphi_{126}) \to \mathrm{Seq}(a) \to \mathrm{SubIdx}(s) \to \mathrm{SubIdx}(t) \to \varphi_{162}$
\item $\varphi_{332} :\equiv (\mathrm{Seq}(a) \to \mathrm{SubIdx}(s) \to \varphi_{127}) \to \mathrm{Seq}(a) \to \mathrm{SubIdx}(s) \to \mathrm{SubIdx}(t) \to \varphi_{163}$
\item $\varphi_{333} :\equiv \mathrm{Conv}(a \circ s,\alpha) \to \mathrm{Conv}(a \circ t,\alpha) \to \varepsilon \in \mathbb{R} \to 0 < \varepsilon \to \varphi_{266}$
\item $\varphi_{334} :\equiv \mathrm{Conv}(a \circ s,\alpha) \to \mathrm{Conv}(a \circ t,\alpha) \to \varepsilon \in \mathbb{R} \to 0 < \varepsilon \to \varphi_{267}$
\item $\varphi_{335} :\equiv \mathrm{Conv}(a \circ s,\alpha) \to \mathrm{Conv}(a \circ t,\alpha) \to \varepsilon \in \mathbb{R} \to 0 < \varepsilon \to \varphi_{270}$
\item $\varphi_{336} :\equiv \mathrm{Conv}(a \circ s,\alpha) \to \mathrm{Conv}(a \circ t,\alpha) \to \varepsilon \in \mathbb{R} \to 0 < \varepsilon \to \varphi_{271}$
\item $\varphi_{337} :\equiv \mathrm{SubIdx}(t) \to \forall n \in \mathbb{N}\,\exists k \in \mathbb{N}\,(n = s_{k} \lor n = t(k)) \to \varphi_{273}$
\item $\varphi_{338} :\equiv \mathrm{Conv}(a \circ s,\alpha) \to \mathrm{Conv}(a \circ t,\alpha) \to \varepsilon \in \mathbb{R} \to 0 < \varepsilon \to \varphi_{274}$
\item $\varphi_{339} :\equiv \mathrm{Conv}(a \circ s,\alpha) \to \mathrm{Conv}(a \circ t,\alpha) \to \varepsilon \in \mathbb{R} \to 0 < \varepsilon \to \varphi_{275}$
\item $\varphi_{340} :\equiv \mathrm{Conv}(a \circ s,\alpha) \to \mathrm{Conv}(a \circ t,\alpha) \to \varepsilon \in \mathbb{R} \to 0 < \varepsilon \to \varphi_{276}$
\item $\varphi_{341} :\equiv \mathrm{Conv}(a \circ s,\alpha) \to \mathrm{Conv}(a \circ t,\alpha) \to \varepsilon \in \mathbb{R} \to 0 < \varepsilon \to \varphi_{277}$
\item $\varphi_{342} :\equiv \mathrm{Conv}(a \circ s,\alpha) \to \mathrm{Conv}(a \circ t,\alpha) \to \varepsilon \in \mathbb{R} \to 0 < \varepsilon \to \varphi_{278}$
\item $\varphi_{343} :\equiv \mathrm{Conv}(a \circ s,\alpha) \to \mathrm{Conv}(a \circ t,\alpha) \to \varepsilon \in \mathbb{R} \to 0 < \varepsilon \to \varphi_{280}$
\item $\varphi_{344} :\equiv \mathrm{SubIdx}(t) \to \forall n \in \mathbb{N}\,\exists k \in \mathbb{N}\,(n = s_{k} \lor n = t(k)) \to \varphi_{282}$
\item $\varphi_{345} :\equiv \mathrm{SubIdx}(t) \to \forall n \in \mathbb{N}\,\exists k \in \mathbb{N}\,(n = s_{k} \lor n = t(k)) \to \varphi_{283}$
\item $\varphi_{346} :\equiv \mathrm{SubIdx}(t) \to \forall n \in \mathbb{N}\,\exists k \in \mathbb{N}\,(n = s_{k} \lor n = t(k)) \to \varphi_{284}$
\item $\varphi_{347} :\equiv \mathrm{SubIdx}(t) \to \forall n \in \mathbb{N}\,\exists k \in \mathbb{N}\,(n = s_{k} \lor n = t(k)) \to \varphi_{285}$
\item $\varphi_{348} :\equiv \mathrm{Conv}(a \circ s,\alpha) \to \mathrm{Conv}(a \circ t,\alpha) \to \varepsilon \in \mathbb{R} \to 0 < \varepsilon \to \varphi_{287}$
\item $\varphi_{349} :\equiv \mathrm{Conv}(a \circ s,\alpha) \to \mathrm{Conv}(a \circ t,\alpha) \to \varepsilon \in \mathbb{R} \to 0 < \varepsilon \to \varphi_{288}$
\item $\varphi_{350} :\equiv \mathrm{Conv}(a \circ s,\alpha) \to \mathrm{Conv}(a \circ t,\alpha) \to \varepsilon \in \mathbb{R} \to 0 < \varepsilon \to \varphi_{290}$
\item $\varphi_{351} :\equiv \mathrm{Conv}(a \circ s,\alpha) \to \mathrm{Conv}(a \circ t,\alpha) \to \varepsilon \in \mathbb{R} \to 0 < \varepsilon \to \varphi_{291}$
\item $\varphi_{352} :\equiv \mathrm{SubIdx}(t) \to \forall n \in \mathbb{N}\,\exists k \in \mathbb{N}\,(n = s_{k} \lor n = t(k)) \to \varphi_{292}$
\item $\varphi_{353} :\equiv \mathrm{Conv}(a \circ s,\alpha) \to \mathrm{Conv}(a \circ t,\alpha) \to \varepsilon \in \mathbb{R} \to 0 < \varepsilon \to \varphi_{294}$
\item $\varphi_{354} :\equiv \mathrm{Conv}(a \circ s,\alpha) \to \mathrm{Conv}(a \circ t,\alpha) \to \varepsilon \in \mathbb{R} \to 0 < \varepsilon \to \varphi_{295}$
\item $\varphi_{355} :\equiv \mathrm{Conv}(a \circ s,\alpha) \to \mathrm{Conv}(a \circ t,\alpha) \to \varepsilon \in \mathbb{R} \to 0 < \varepsilon \to \varphi_{299}$
\item $\varphi_{356} :\equiv \mathrm{Conv}(a \circ s,\alpha) \to \mathrm{Conv}(a \circ t,\alpha) \to \varepsilon \in \mathbb{R} \to 0 < \varepsilon \to \varphi_{300}$
\item $\varphi_{357} :\equiv (\mathrm{Seq}(a) \to \mathrm{SubIdx}(s) \to \mathrm{SubIdx}(t) \to \varphi_{192}) \to \mathrm{Seq}(a) \to \mathrm{SubIdx}(s) \to \varphi_{128}$
\item $\varphi_{358} :\equiv \mathrm{SubIdx}(t) \to \forall n \in \mathbb{N}\,\exists k \in \mathbb{N}\,(n = s_{k} \lor n = t(k)) \to \varphi_{301}$
\item $\varphi_{359} :\equiv \forall n\,(\varphi_{298} \to n \in \mathbb{N} \to s_{N_{1}} \cup t(N_{2}) \leq n \to \lvert a_{n} - \alpha \rvert < \varepsilon)$
\item $\varphi_{360} :\equiv \forall k\,(\varphi_{298} \land n \in \mathbb{N} \land s_{N_{1}} \cup t(N_{2}) \leq n \land \neg \lvert a_{n} - \alpha \rvert < \varepsilon)$
\item $\varphi_{361} :\equiv \forall n\,(\varphi_{298}) \to \forall n \in \mathbb{N}\,(s_{N_{1}} \cup t(N_{2}) \leq n \to \lvert a_{n} - \alpha \rvert < \varepsilon)$
\item $\varphi_{362} :\equiv (\mathrm{Seq}(a) \to \mathrm{SubIdx}(s) \to \mathrm{SubIdx}(t) \to \varphi_{162}) \to \mathrm{Seq}(a) \to \mathrm{SubIdx}(s) \to \varphi_{157}$
\item $\varphi_{363} :\equiv (\mathrm{Seq}(a) \to \mathrm{SubIdx}(s) \to \mathrm{SubIdx}(t) \to \varphi_{163}) \to \mathrm{Seq}(a) \to \mathrm{SubIdx}(s) \to \varphi_{158}$
\item $\varphi_{364} :\equiv (\forall N_{1}\,(\varphi_{35}) \to \forall N\,\forall N_{1}\,(\varphi_{35})) \to \mathrm{Seq}(a) \to \mathrm{SubIdx}(s) \to \mathrm{SubIdx}(t) \to \varphi_{187}$
\item $\varphi_{365} :\equiv \mathrm{SubIdx}(t) \to \forall n \in \mathbb{N}\,\exists k \in \mathbb{N}\,(n = s_{k} \lor n = t(k)) \to \varphi_{304}$
\item $\varphi_{366} :\equiv \mathrm{SubIdx}(t) \to \forall n \in \mathbb{N}\,\exists k \in \mathbb{N}\,(n = s_{k} \lor n = t(k)) \to \varphi_{305}$
\item $\varphi_{367} :\equiv \mathrm{SubIdx}(t) \to \forall n \in \mathbb{N}\,\exists k \in \mathbb{N}\,(n = s_{k} \lor n = t(k)) \to \varphi_{306}$
\item $\varphi_{368} :\equiv \mathrm{SubIdx}(t) \to \forall n \in \mathbb{N}\,\exists k \in \mathbb{N}\,(n = s_{k} \lor n = t(k)) \to \varphi_{307}$
\item $\varphi_{369} :\equiv \mathrm{SubIdx}(t) \to \forall n \in \mathbb{N}\,\exists k \in \mathbb{N}\,(n = s_{k} \lor n = t(k)) \to \varphi_{308}$
\item $\varphi_{370} :\equiv \mathrm{SubIdx}(t) \to \forall n \in \mathbb{N}\,\exists k \in \mathbb{N}\,(n = s_{k} \lor n = t(k)) \to \varphi_{309}$
\item $\varphi_{371} :\equiv \mathrm{SubIdx}(t) \to \forall n \in \mathbb{N}\,\exists k \in \mathbb{N}\,(n = s_{k} \lor n = t(k)) \to \varphi_{310}$
\item $\varphi_{372} :\equiv \mathrm{SubIdx}(t) \to \forall n \in \mathbb{N}\,\exists k \in \mathbb{N}\,(n = s_{k} \lor n = t(k)) \to \varphi_{311}$
\item $\varphi_{373} :\equiv \mathrm{SubIdx}(t) \to \forall n \in \mathbb{N}\,\exists k \in \mathbb{N}\,(n = s_{k} \lor n = t(k)) \to \varphi_{312}$
\item $\varphi_{374} :\equiv \mathrm{SubIdx}(t) \to \forall n \in \mathbb{N}\,\exists k \in \mathbb{N}\,(n = s_{k} \lor n = t(k)) \to \varphi_{313}$
\item $\varphi_{375} :\equiv \mathrm{SubIdx}(t) \to \forall n \in \mathbb{N}\,\exists k \in \mathbb{N}\,(n = s_{k} \lor n = t(k)) \to \varphi_{314}$
\item $\varphi_{376} :\equiv \mathrm{SubIdx}(t) \to \forall n \in \mathbb{N}\,\exists k \in \mathbb{N}\,(n = s_{k} \lor n = t(k)) \to \varphi_{315}$
\item $\varphi_{377} :\equiv \mathrm{SubIdx}(t) \to \forall n \in \mathbb{N}\,\exists k \in \mathbb{N}\,(n = s_{k} \lor n = t(k)) \to \varphi_{316}$
\item $\varphi_{378} :\equiv \mathrm{SubIdx}(t) \to \forall n \in \mathbb{N}\,\exists k \in \mathbb{N}\,(n = s_{k} \lor n = t(k)) \to \varphi_{317}$
\item $\varphi_{379} :\equiv \mathrm{SubIdx}(t) \to \forall n \in \mathbb{N}\,\exists k \in \mathbb{N}\,(n = s_{k} \lor n = t(k)) \to \varphi_{318}$
\item $\varphi_{380} :\equiv \mathrm{SubIdx}(t) \to \forall n \in \mathbb{N}\,\exists k \in \mathbb{N}\,(n = s_{k} \lor n = t(k)) \to \varphi_{319}$
\item $\varphi_{381} :\equiv \mathrm{SubIdx}(t) \to \forall n \in \mathbb{N}\,\exists k \in \mathbb{N}\,(n = s_{k} \lor n = t(k)) \to \varphi_{320}$
\item $\varphi_{382} :\equiv \mathrm{SubIdx}(t) \to \forall n \in \mathbb{N}\,\exists k \in \mathbb{N}\,(n = s_{k} \lor n = t(k)) \to \varphi_{321}$
\item $\varphi_{383} :\equiv \mathrm{SubIdx}(t) \to \forall n \in \mathbb{N}\,\exists k \in \mathbb{N}\,(n = s_{k} \lor n = t(k)) \to \varphi_{322}$
\item $\varphi_{384} :\equiv \mathrm{SubIdx}(t) \to \forall n \in \mathbb{N}\,\exists k \in \mathbb{N}\,(n = s_{k} \lor n = t(k)) \to \varphi_{323}$
\item $\varphi_{385} :\equiv (\mathrm{Seq}(a) \to \mathrm{SubIdx}(s) \to \mathrm{SubIdx}(t) \to \varphi_{122}) \to \varphi_{118} \to \mathrm{Seq}(a) \to \mathrm{SubIdx}(s) \to \mathrm{SubIdx}(t) \to \varphi_{118}$
\item $\varphi_{386} :\equiv \mathrm{SubIdx}(t) \to \forall n \in \mathbb{N}\,\exists k \in \mathbb{N}\,(n = s_{k} \lor n = t(k)) \to \varphi_{324}$
\item $\varphi_{387} :\equiv (\mathrm{Seq}(a) \to \mathrm{SubIdx}(s) \to \varphi_{147}) \to \mathrm{Seq}(a) \to \mathrm{SubIdx}(s) \to \mathrm{SubIdx}(t) \to \varphi_{193}$
\item $\varphi_{388} :\equiv \mathrm{SubIdx}(t) \to \forall n \in \mathbb{N}\,\exists k \in \mathbb{N}\,(n = s_{k} \lor n = t(k)) \to \varphi_{325}$
\item $\varphi_{389} :\equiv \mathrm{SubIdx}(t) \to \forall n \in \mathbb{N}\,\exists k \in \mathbb{N}\,(n = s_{k} \lor n = t(k)) \to \varphi_{326}$
\item $\varphi_{390} :\equiv \mathrm{SubIdx}(t) \to \forall n \in \mathbb{N}\,\exists k \in \mathbb{N}\,(n = s_{k} \lor n = t(k)) \to \varphi_{327}$
\item $\varphi_{391} :\equiv \mathrm{SubIdx}(t) \to \forall n \in \mathbb{N}\,\exists k \in \mathbb{N}\,(n = s_{k} \lor n = t(k)) \to \varphi_{328}$
\item $\varphi_{392} :\equiv \mathrm{SubIdx}(t) \to \forall n \in \mathbb{N}\,\exists k \in \mathbb{N}\,(n = s_{k} \lor n = t(k)) \to \varphi_{329}$
\item $\varphi_{393} :\equiv \mathrm{SubIdx}(t) \to \forall n \in \mathbb{N}\,\exists k \in \mathbb{N}\,(n = s_{k} \lor n = t(k)) \to \varphi_{330}$
\item $\varphi_{394} :\equiv \mathrm{SubIdx}(t) \to \forall n \in \mathbb{N}\,\exists k \in \mathbb{N}\,(n = s_{k} \lor n = t(k)) \to \varphi_{333}$
\item $\varphi_{395} :\equiv \mathrm{SubIdx}(t) \to \forall n \in \mathbb{N}\,\exists k \in \mathbb{N}\,(n = s_{k} \lor n = t(k)) \to \varphi_{334}$
\item $\varphi_{396} :\equiv \varphi_{298} \land n \in \mathbb{N} \land s_{N_{1}} \cup t(N_{2}) \leq n \land \neg \lvert a_{n} - \alpha \rvert < \varepsilon \to \neg (k \in \mathbb{N} \land (n = s_{k} \lor n = t(k)))$
\item $\varphi_{397} :\equiv \mathrm{SubIdx}(t) \to \forall n \in \mathbb{N}\,\exists k \in \mathbb{N}\,(n = s_{k} \lor n = t(k)) \to \varphi_{335}$
\item $\varphi_{398} :\equiv \mathrm{SubIdx}(t) \to \forall n \in \mathbb{N}\,\exists k \in \mathbb{N}\,(n = s_{k} \lor n = t(k)) \to \varphi_{336}$
\item $\varphi_{399} :\equiv \mathrm{SubIdx}(t) \to \forall n \in \mathbb{N}\,\exists k \in \mathbb{N}\,(n = s_{k} \lor n = t(k)) \to \varphi_{338}$
\item $\varphi_{400} :\equiv \mathrm{SubIdx}(t) \to \forall n \in \mathbb{N}\,\exists k \in \mathbb{N}\,(n = s_{k} \lor n = t(k)) \to \varphi_{339}$
\item $\varphi_{401} :\equiv \mathrm{SubIdx}(t) \to \forall n \in \mathbb{N}\,\exists k \in \mathbb{N}\,(n = s_{k} \lor n = t(k)) \to \varphi_{340}$
\item $\varphi_{402} :\equiv \mathrm{SubIdx}(t) \to \forall n \in \mathbb{N}\,\exists k \in \mathbb{N}\,(n = s_{k} \lor n = t(k)) \to \varphi_{341}$
\item $\varphi_{403} :\equiv \mathrm{SubIdx}(t) \to \forall n \in \mathbb{N}\,\exists k \in \mathbb{N}\,(n = s_{k} \lor n = t(k)) \to \varphi_{342}$
\item $\varphi_{404} :\equiv \mathrm{SubIdx}(t) \to \forall n \in \mathbb{N}\,\exists k \in \mathbb{N}\,(n = s_{k} \lor n = t(k)) \to \varphi_{343}$
\item $\varphi_{405} :\equiv \mathrm{SubIdx}(t) \to \forall n \in \mathbb{N}\,\exists k \in \mathbb{N}\,(n = s_{k} \lor n = t(k)) \to \varphi_{348}$
\item $\varphi_{406} :\equiv \mathrm{SubIdx}(t) \to \forall n \in \mathbb{N}\,\exists k \in \mathbb{N}\,(n = s_{k} \lor n = t(k)) \to \varphi_{349}$
\item $\varphi_{407} :\equiv \mathrm{SubIdx}(t) \to \forall n \in \mathbb{N}\,\exists k \in \mathbb{N}\,(n = s_{k} \lor n = t(k)) \to \varphi_{350}$
\item $\varphi_{408} :\equiv \mathrm{SubIdx}(t) \to \forall n \in \mathbb{N}\,\exists k \in \mathbb{N}\,(n = s_{k} \lor n = t(k)) \to \varphi_{351}$
\item $\varphi_{409} :\equiv (n = t(k) \to t(k) = n) \to \mathrm{Seq}(a) \to \mathrm{SubIdx}(s) \to \varphi_{366}$
\item $\varphi_{410} :\equiv \mathrm{SubIdx}(t) \to \forall n \in \mathbb{N}\,\exists k \in \mathbb{N}\,(n = s_{k} \lor n = t(k)) \to \varphi_{353}$
\item $\varphi_{411} :\equiv \mathrm{SubIdx}(t) \to \forall n \in \mathbb{N}\,\exists k \in \mathbb{N}\,(n = s_{k} \lor n = t(k)) \to \varphi_{354}$
\item $\varphi_{412} :\equiv \mathrm{SubIdx}(t) \to \forall n \in \mathbb{N}\,\exists k \in \mathbb{N}\,(n = s_{k} \lor n = t(k)) \to \varphi_{355}$
\item $\varphi_{413} :\equiv \mathrm{SubIdx}(t) \to \forall n \in \mathbb{N}\,\exists k \in \mathbb{N}\,(n = s_{k} \lor n = t(k)) \to \varphi_{356}$
\item $\varphi_{414} :\equiv (\mathrm{Seq}(a) \to \mathrm{SubIdx}(s) \to \mathrm{SubIdx}(t) \to \varphi_{187}) \to \varphi_{129} \to \varphi_{80}$
\item $\varphi_{415} :\equiv (t(k) = n \to t(k) < n \to n < n) \to \mathrm{Seq}(a) \to \mathrm{SubIdx}(s) \to \varphi_{381}$
\item $\varphi_{416} :\equiv (s_{k} = n \to s_{k} < n \to n < n) \to \mathrm{Seq}(a) \to \mathrm{SubIdx}(s) \to \varphi_{382}$
\item $\varphi_{417} :\equiv (\forall N_{2}\,(\varphi_{34}) \to \forall N\,\forall N_{2}\,(\varphi_{34})) \to \mathrm{Seq}(a) \to \mathrm{SubIdx}(s) \to \varphi_{268}$
\item $\varphi_{418} :\equiv (\mathrm{Seq}(a) \to \mathrm{SubIdx}(s) \to \varphi_{157}) \to \varphi_{105} \to \mathrm{Seq}(a) \to \mathrm{SubIdx}(s) \to \varphi_{152}$
\item $\varphi_{419} :\equiv (\mathrm{Seq}(a) \to \mathrm{SubIdx}(s) \to \varphi_{158}) \to \varphi_{106} \to \mathrm{Seq}(a) \to \mathrm{SubIdx}(s) \to \varphi_{154}$
\item $\varphi_{420} :\equiv (\mathrm{Seq}(a) \to \mathrm{SubIdx}(s) \to \varphi_{147}) \to \mathrm{Seq}(a) \to \mathrm{SubIdx}(s) \to \varphi_{281}$
\item $\varphi_{421} :\equiv (\varphi_{129} \to \forall \varepsilon\,(\varphi_{129})) \to \mathrm{Seq}(a) \to \mathrm{SubIdx}(s) \to \mathrm{SubIdx}(t) \to \varphi_{192}$
\item $\varphi_{422} :\equiv (t(k) = n \to \lvert a_{t(k)} - \alpha \rvert < \varepsilon \to \lvert a_{n} - \alpha \rvert < \varepsilon) \to \mathrm{Seq}(a) \to \mathrm{SubIdx}(s) \to \varphi_{394}$
\item $\varphi_{423} :\equiv (s_{k} = n \to \lvert a_{s_{k}} - \alpha \rvert < \varepsilon \to \lvert a_{n} - \alpha \rvert < \varepsilon) \to \mathrm{Seq}(a) \to \mathrm{SubIdx}(s) \to \varphi_{395}$
\item $\varphi_{424} :\equiv (\mathrm{Seq}(a) \to \mathrm{SubIdx}(s) \to \mathrm{SubIdx}(t) \to \varphi_{123}) \to \varphi_{357}$
\item $\varphi_{425} :\equiv (\mathrm{Seq}(a) \to \mathrm{SubIdx}(s) \to \varphi_{279}) \to \mathrm{Seq}(a) \to \mathrm{SubIdx}(s) \to \varphi_{289}$
\item $\varphi_{426} :\equiv (\mathrm{Seq}(a) \to \mathrm{SubIdx}(s) \to \varphi_{286}) \to \mathrm{Seq}(a) \to \mathrm{SubIdx}(s) \to \varphi_{296}$
\item $\varphi_{427} :\equiv (\mathrm{Seq}(a) \to \mathrm{SubIdx}(s) \to \varphi_{289}) \to \mathrm{Seq}(a) \to \mathrm{SubIdx}(s) \to \varphi_{297}$
\item $\varphi_{428} :\equiv (\mathrm{Seq}(a) \to \mathrm{SubIdx}(s) \to \varphi_{296}) \to \mathrm{Seq}(a) \to \mathrm{SubIdx}(s) \to \varphi_{303}$
\item $\varphi_{429} :\equiv (\mathrm{Conv}(a,\alpha) \leftrightarrow \mathrm{Seq}(a) \land \alpha \in \mathbb{R} \land \varphi_{102}) \to (\mathrm{Seq}(a) \to \mathrm{Seq}(a)) \to \varphi_{424}$
\item $\varphi_{430} :\equiv (\forall N\,(\forall N_{1}\,(\varphi_{35}) \to \varphi_{28}) \to \forall N\,\forall N_{1}\,(\varphi_{35}) \to \forall N\,(\varphi_{28})) \to \varphi_{364}$
\item $\varphi_{431} :\equiv (\mathrm{Seq}(a) \to \mathrm{SubIdx}(s) \to \varphi_{346}) \to \mathrm{Seq}(a) \to \mathrm{SubIdx}(s) \to \varphi_{286}$
\item $\varphi_{432} :\equiv (\mathrm{Seq}(a) \to \mathrm{SubIdx}(s) \to \mathrm{SubIdx}(t) \to \varphi_{109}) \to (\mathrm{Seq}(a) \to \mathrm{SubIdx}(s) \to \varphi_{136}) \to \varphi_{387}$
\item $\varphi_{433} :\equiv (\mathrm{Seq}(a) \to \mathrm{SubIdx}(s) \to \varphi_{302}) \to \mathrm{Seq}(a) \to \mathrm{SubIdx}(s) \to \varphi_{352}$
\item $\varphi_{434} :\equiv (\mathrm{Seq}(a) \to \mathrm{SubIdx}(s) \to \varphi_{286}) \to \mathrm{Seq}(a) \to \mathrm{SubIdx}(s) \to \varphi_{375}$
\item $\varphi_{435} :\equiv (\mathrm{Seq}(a) \to \mathrm{SubIdx}(s) \to \varphi_{289}) \to \mathrm{Seq}(a) \to \mathrm{SubIdx}(s) \to \varphi_{376}$
\item $\varphi_{436} :\equiv (\mathrm{Seq}(a) \to \mathrm{SubIdx}(s) \to \varphi_{345}) \to \varphi_{298} \to n \in \mathbb{N} \to s_{N_{1}} \cup t(N_{2}) \leq n \to \lvert a_{n} - \alpha \rvert < \varepsilon$
\item $\varphi_{437} :\equiv \varphi_{298} \land n \in \mathbb{N} \land s_{N_{1}} \cup t(N_{2}) \leq n \land \neg \lvert a_{n} - \alpha \rvert < \varepsilon \to \varphi_{360}$
\item $\varphi_{438} :\equiv (\forall N\,(\forall N_{2}\,(\varphi_{34}) \to \varphi_{29}) \to \forall N\,\forall N_{2}\,(\varphi_{34}) \to \forall N\,(\varphi_{29})) \to \varphi_{417}$
\item $\varphi_{439} :\equiv (\mathrm{Seq}(a) \to \mathrm{SubIdx}(s) \to \varphi_{279}) \to \mathrm{Seq}(a) \to \mathrm{SubIdx}(s) \to \varphi_{413}$
\item $\varphi_{440} :\equiv (\mathrm{Seq}(a) \to \mathrm{SubIdx}(s) \to \varphi_{344}) \to \mathrm{Seq}(a) \to \mathrm{SubIdx}(s) \to \varphi_{369}$
\item $\varphi_{441} :\equiv (\mathrm{Seq}(a) \to \mathrm{SubIdx}(s) \to \varphi_{344}) \to \mathrm{Seq}(a) \to \mathrm{SubIdx}(s) \to \varphi_{372}$
\item $\varphi_{442} :\equiv (\mathrm{Seq}(a) \to \mathrm{SubIdx}(s) \to \varphi_{347}) \to \mathrm{Seq}(a) \to \mathrm{SubIdx}(s) \to \varphi_{378}$
\item $\varphi_{443} :\equiv (\mathrm{Seq}(a) \to \mathrm{SubIdx}(s) \to \varphi_{130}) \to (\mathrm{Seq}(a) \to \mathrm{SubIdx}(s) \to \varphi_{136}) \to \varphi_{420}$
\item $\varphi_{444} :\equiv (\mathrm{Seq}(a) \to \mathrm{SubIdx}(s) \to \varphi_{358}) \to \mathrm{Seq}(a) \to \mathrm{SubIdx}(s) \to \varphi_{371}$
\item $\varphi_{445} :\equiv (\mathrm{Seq}(a) \to \mathrm{SubIdx}(s) \to \varphi_{358}) \to \mathrm{Seq}(a) \to \mathrm{SubIdx}(s) \to \varphi_{374}$
\item $\varphi_{446} :\equiv (\mathrm{Seq}(a) \to \mathrm{SubIdx}(s) \to \varphi_{386}) \to \mathrm{Seq}(a) \to \mathrm{SubIdx}(s) \to \varphi_{358}$
\item $\varphi_{447} :\equiv (\mathrm{Seq}(a) \to \mathrm{SubIdx}(s) \to \varphi_{371}) \to \mathrm{Seq}(a) \to \mathrm{SubIdx}(s) \to \varphi_{379}$
\item $\varphi_{448} :\equiv (\mathrm{Seq}(a) \to \mathrm{SubIdx}(s) \to \varphi_{374}) \to \mathrm{Seq}(a) \to \mathrm{SubIdx}(s) \to \varphi_{380}$
\item $\varphi_{449} :\equiv (\mathrm{NSeq}(t) \to N_{2} \in \mathbb{N} \to t(N_{2}) \in \mathbb{N}) \to (\mathrm{Seq}(a) \to \mathrm{SubIdx}(s) \to \varphi_{126}) \to \varphi_{426}$
\item $\varphi_{450} :\equiv (\mathrm{NSeq}(s) \to N_{1} \in \mathbb{N} \to s_{N_{1}} \in \mathbb{N}) \to (\mathrm{Seq}(a) \to \mathrm{SubIdx}(s) \to \varphi_{127}) \to \varphi_{427}$
\item $\varphi_{451} :\equiv (\mathrm{Seq}(a) \to \mathrm{SubIdx}(s) \to \varphi_{369}) \to \mathrm{Seq}(a) \to \mathrm{SubIdx}(s) \to \varphi_{388}$
\item $\varphi_{452} :\equiv (\mathrm{Seq}(a) \to \mathrm{SubIdx}(s) \to \varphi_{372}) \to \mathrm{Seq}(a) \to \mathrm{SubIdx}(s) \to \varphi_{389}$
\item $\varphi_{453} :\equiv (\mathrm{Seq}(a) \to \mathrm{SubIdx}(s) \to \varphi_{358}) \to \mathrm{Seq}(a) \to \mathrm{SubIdx}(s) \to \varphi_{399}$
\item $\varphi_{454} :\equiv (\mathrm{Seq}(a) \to \mathrm{SubIdx}(s) \to \varphi_{358}) \to \mathrm{Seq}(a) \to \mathrm{SubIdx}(s) \to \varphi_{400}$
\item $\varphi_{455} :\equiv \forall k\,(\varphi_{396}) \to \varphi_{360} \to \forall k\,(\neg (k \in \mathbb{N} \land (n = s_{k} \lor n = t(k))))$
\item $\varphi_{456} :\equiv (\mathrm{Seq}(a) \to \mathrm{SubIdx}(s) \to \varphi_{390}) \to \mathrm{Seq}(a) \to \mathrm{SubIdx}(s) \to \varphi_{377}$
\item $\varphi_{457} :\equiv (\mathrm{Seq}(a) \to \mathrm{SubIdx}(s) \to \varphi_{346}) \to \mathrm{Seq}(a) \to \mathrm{SubIdx}(s) \to \varphi_{412}$
\item $\varphi_{458} :\equiv (\mathrm{Seq}(a) \to \mathrm{SubIdx}(s) \to \varphi_{367}) \to (n = s_{k} \to s_{k} = n) \to \mathrm{Seq}(a) \to \mathrm{SubIdx}(s) \to \varphi_{368}$
\item $\varphi_{459} :\equiv (\mathrm{Seq}(a) \to \mathrm{SubIdx}(s) \to \varphi_{383}) \to \mathrm{Seq}(a) \to \mathrm{SubIdx}(s) \to \varphi_{392}$
\item $\varphi_{460} :\equiv (\mathrm{Seq}(a) \to \mathrm{SubIdx}(s) \to \varphi_{384}) \to \mathrm{Seq}(a) \to \mathrm{SubIdx}(s) \to \varphi_{393}$
\item $\varphi_{461} :\equiv (\mathrm{Seq}(a) \to \mathrm{SubIdx}(s) \to \varphi_{358}) \to \mathrm{Seq}(a) \to \mathrm{SubIdx}(s) \to \varphi_{405}$
\item $\varphi_{462} :\equiv (\mathrm{Seq}(a) \to \mathrm{SubIdx}(s) \to \varphi_{358}) \to \mathrm{Seq}(a) \to \mathrm{SubIdx}(s) \to \varphi_{407}$
\item $\varphi_{463} :\equiv (\mathrm{Seq}(a) \to \mathrm{SubIdx}(s) \to \varphi_{370}) \to \mathrm{Seq}(a) \to \mathrm{SubIdx}(s) \to \varphi_{403}$
\item $\varphi_{464} :\equiv (\mathrm{Seq}(a) \to \mathrm{SubIdx}(s) \to \varphi_{373}) \to \mathrm{Seq}(a) \to \mathrm{SubIdx}(s) \to \varphi_{404}$
\item $\varphi_{465} :\equiv (\mathrm{Seq}(a) \to \mathrm{SubIdx}(s) \to \mathrm{SubIdx}(t) \to \varphi_{118}) \to \varphi_{433}$
\item $\varphi_{466} :\equiv (\mathrm{Seq}(a) \to \mathrm{SubIdx}(s) \to \varphi_{381}) \to \neg n < n \to \mathrm{Seq}(a) \to \mathrm{SubIdx}(s) \to \varphi_{397}$
\item $\varphi_{467} :\equiv (\mathrm{Seq}(a) \to \mathrm{SubIdx}(s) \to \varphi_{397}) \to \mathrm{Seq}(a) \to \mathrm{SubIdx}(s) \to \varphi_{390}$
\item $\varphi_{468} :\equiv (\mathrm{Seq}(a) \to \mathrm{SubIdx}(s) \to \varphi_{382}) \to \neg n < n \to \mathrm{Seq}(a) \to \mathrm{SubIdx}(s) \to \varphi_{398}$
\item $\varphi_{469} :\equiv (\mathrm{Seq}(a) \to \mathrm{SubIdx}(s) \to \varphi_{398}) \to \mathrm{Seq}(a) \to \mathrm{SubIdx}(s) \to \varphi_{391}$
\item $\varphi_{470} :\equiv \varphi_{27} \to (\mathrm{Seq}(a) \to \mathrm{SubIdx}(s) \to \mathrm{SubIdx}(s)) \to \varphi_{462}$
\item $\varphi_{471} :\equiv (\mathrm{Seq}(a) \to \mathrm{SubIdx}(s) \to \varphi_{378}) \to (\varphi_{13} \to \neg \varphi_{117}) \to \mathrm{Seq}(a) \to \mathrm{SubIdx}(s) \to \varphi_{337}$
\item $\varphi_{472} :\equiv (\mathrm{Seq}(a) \to \mathrm{SubIdx}(s) \to \varphi_{403}) \to \varphi_{113} \to \mathrm{Seq}(a) \to \mathrm{SubIdx}(s) \to \varphi_{401}$
\item $\varphi_{473} :\equiv (\mathrm{Seq}(a) \to \mathrm{SubIdx}(s) \to \varphi_{404}) \to \varphi_{114} \to \mathrm{Seq}(a) \to \mathrm{SubIdx}(s) \to \varphi_{402}$
\item $\varphi_{474} :\equiv (\mathrm{Seq}(a) \to \mathrm{SubIdx}(s) \to \varphi_{405}) \to \varphi_{89} \to \mathrm{Seq}(a) \to \mathrm{SubIdx}(s) \to \varphi_{406}$
\item $\varphi_{475} :\equiv (\mathrm{Seq}(a) \to \mathrm{SubIdx}(s) \to \varphi_{407}) \to \varphi_{90} \to \mathrm{Seq}(a) \to \mathrm{SubIdx}(s) \to \varphi_{408}$
\item $\varphi_{476} :\equiv (\mathrm{Seq}(a) \to \mathrm{SubIdx}(s) \to \varphi_{412}) \to \varphi_{124} \to \mathrm{Seq}(a) \to \mathrm{SubIdx}(s) \to \varphi_{410}$
\item $\varphi_{477} :\equiv (\mathrm{Seq}(a) \to \mathrm{SubIdx}(s) \to \varphi_{413}) \to \varphi_{125} \to \mathrm{Seq}(a) \to \mathrm{SubIdx}(s) \to \varphi_{411}$
\item $\varphi_{478} :\equiv (s_{N_{1}} \in \mathbb{N} \to t(N_{2}) \in \mathbb{N} \to s_{N_{1}} \cup t(N_{2}) \in \mathbb{N}) \to (\mathrm{Seq}(a) \to \mathrm{SubIdx}(s) \to \varphi_{297}) \to \varphi_{428}$
\item $\varphi_{479} :\equiv (N_{2} \in \omega \to k \in \omega \to N_{2} \leq k \lor k < N_{2}) \to (\mathrm{Seq}(a) \to \mathrm{SubIdx}(s) \to \varphi_{375}) \to \varphi_{447}$
\item $\varphi_{480} :\equiv (N_{1} \in \omega \to k \in \omega \to N_{1} \leq k \lor k < N_{1}) \to (\mathrm{Seq}(a) \to \mathrm{SubIdx}(s) \to \varphi_{376}) \to \varphi_{448}$
\item $\varphi_{481} :\equiv (\mathrm{Seq}(a) \to \mathrm{SubIdx}(s) \to \varphi_{366}) \to (\mathrm{Seq}(a) \to \mathrm{SubIdx}(s) \to \varphi_{388}) \to \varphi_{415}$
\item $\varphi_{482} :\equiv (\mathrm{Seq}(a) \to \mathrm{SubIdx}(s) \to \varphi_{368}) \to (\mathrm{Seq}(a) \to \mathrm{SubIdx}(s) \to \varphi_{389}) \to \varphi_{416}$
\item $\varphi_{483} :\equiv (t(k) < t(N_{2}) \to t(N_{2}) \leq n \to t(k) < n) \to (\mathrm{Seq}(a) \to \mathrm{SubIdx}(s) \to \varphi_{392}) \to \varphi_{451}$
\item $\varphi_{484} :\equiv (s_{k} < s_{N_{1}} \to s_{N_{1}} \leq n \to s_{k} < n) \to (\mathrm{Seq}(a) \to \mathrm{SubIdx}(s) \to \varphi_{393}) \to \varphi_{452}$
\item $\varphi_{485} :\equiv (\mathrm{Seq}(a) \to \mathrm{SubIdx}(s) \to \varphi_{366}) \to (\mathrm{Seq}(a) \to \mathrm{SubIdx}(s) \to \varphi_{401}) \to \varphi_{422}$
\item $\varphi_{486} :\equiv (\mathrm{Seq}(a) \to \mathrm{SubIdx}(s) \to \varphi_{368}) \to (\mathrm{Seq}(a) \to \mathrm{SubIdx}(s) \to \varphi_{402}) \to \varphi_{423}$
\item $\varphi_{487} :\equiv (\mathrm{Seq}(a) \to \mathrm{SubIdx}(s) \to \varphi_{406}) \to (\mathrm{Seq}(a) \to \mathrm{SubIdx}(s) \to \varphi_{286}) \to \varphi_{459}$
\item $\varphi_{488} :\equiv (\mathrm{Seq}(a) \to \mathrm{SubIdx}(s) \to \varphi_{408}) \to (\mathrm{Seq}(a) \to \mathrm{SubIdx}(s) \to \varphi_{289}) \to \varphi_{460}$
\item $\varphi_{489} :\equiv (\mathrm{Seq}(a) \to \mathrm{SubIdx}(s) \to \varphi_{386}) \to (\mathrm{Seq}(a) \to \mathrm{SubIdx}(s) \to \varphi_{391}) \to \varphi_{456}$
\item $\varphi_{490} :\equiv (\mathrm{Seq}(a) \to \mathrm{SubIdx}(s) \to \varphi_{410}) \to (\mathrm{Seq}(a) \to \mathrm{SubIdx}(s) \to \varphi_{358}) \to \varphi_{463}$
\item $\varphi_{491} :\equiv (\mathrm{Seq}(a) \to \mathrm{SubIdx}(s) \to \varphi_{411}) \to (\mathrm{Seq}(a) \to \mathrm{SubIdx}(s) \to \varphi_{358}) \to \varphi_{464}$
\item $\varphi_{492} :\equiv (\mathrm{Seq}(a) \to \mathrm{SubIdx}(s) \to \varphi_{379}) \to (\mathrm{Seq}(a) \to \mathrm{SubIdx}(s) \to \varphi_{394}) \to \varphi_{467}$
\item $\varphi_{493} :\equiv (\mathrm{Seq}(a) \to \mathrm{SubIdx}(s) \to \varphi_{380}) \to (\mathrm{Seq}(a) \to \mathrm{SubIdx}(s) \to \varphi_{395}) \to \varphi_{469}$
\item $\varphi_{494} :\equiv (\varphi_{359} \to \varphi_{361}) \to (\varphi_{298} \to \forall n\,(\varphi_{298})) \to \mathrm{Seq}(a) \to \mathrm{SubIdx}(s) \to \varphi_{347}$
\item $\varphi_{495} :\equiv (\forall N_{1}\,(\varphi_{218} \to \varphi_{35}) \to \forall N_{1}\,(\varphi_{218}) \to \forall N_{1}\,(\varphi_{35})) \to (\varphi_{218} \to \forall N_{1}\,(\varphi_{218})) \to \forall N\,(\forall N_{1}\,(\varphi_{35}) \to \varphi_{28}) \to \varphi_{430}$
\item $\varphi_{496} :\equiv (\forall N_{2}\,(\varphi_{293} \to \varphi_{34}) \to \forall N_{2}\,(\varphi_{293}) \to \forall N_{2}\,(\varphi_{34})) \to (\varphi_{293} \to \forall N_{2}\,(\varphi_{293})) \to \forall N\,(\forall N_{2}\,(\varphi_{34}) \to \varphi_{29}) \to \varphi_{438}$
\item $\varphi_{497} :\equiv (\mathrm{Seq}(a) \to \mathrm{SubIdx}(s) \to \mathrm{SubIdx}(t) \to \varphi_{193}) \to \forall N_{1}\,(\varphi_{218} \to \varphi_{35}) \to \varphi_{495}$
\item $\varphi_{498} :\equiv (\mathrm{Seq}(a) \to \mathrm{SubIdx}(s) \to \varphi_{352}) \to \forall k\,(\varphi_{396}) \to \varphi_{455} \to \varphi_{437} \to \mathrm{Seq}(a) \to \mathrm{SubIdx}(s) \to \varphi_{345}$
\item $\varphi_{499} :\equiv (\mathrm{Seq}(a) \to \mathrm{SubIdx}(s) \to \varphi_{281}) \to \forall N_{2}\,(\varphi_{293} \to \varphi_{34}) \to \varphi_{496}$
\end{abbrevlist}
\end{fullproof}

% 自動生成 (prover/build.js): OddEven  ノード数 104
\begin{fullproof}{命題2.13}{$\forall a\,\forall \alpha\,(\mathrm{Conv}(a,\alpha) \leftrightarrow \mathrm{Seq}(a) \land \mathrm{Conv}(a \circ \mathsf{o},\alpha) \land \mathrm{Conv}(a \circ \mathsf{e},\alpha))$}{104}\phantomsection\label{tree:OddEven}
\noindent$\Pi_{1}$:
\begin{ptree}
\AXC{$k \in \mathbb{N} \to k \in \mathbb{N}$\lbl{A1}}
\AXC{$\varphi_{24}$\lbl{A2}}
\AXC{$\varphi_{29}$\lbl{[S15]}}
\AXC{$\varphi_{29} \to k \in \mathbb{N} \to \varphi_{16}$\lbl{A1}}
\BIC{$k \in \mathbb{N} \to \varphi_{16}$}
\AXC{$(k \in \mathbb{N} \to \varphi_{16}) \to \varphi_{24} \to \varphi_{19}$\lbl{A1}}
\BIC{$\varphi_{24} \to \varphi_{19}$}
\BIC{$\varphi_{19}$}
\AXC{$\varphi_{25}$\lbl{A1}}
\BIC{$(k \in \mathbb{N} \to k \in \mathbb{N}) \to k \in \mathbb{N} \to \mathsf{o}(k) < \mathsf{o}(S(k))$}
\BIC{$k \in \mathbb{N} \to \mathsf{o}(k) < \mathsf{o}(S(k))$}
\UIC{$\forall k \in \mathbb{N}\,(\mathsf{o}(k) < \mathsf{o}(S(k)))$}
\end{ptree}
\noindent$\Pi_{2}$:
\begin{ptree}
\AXC{$\stackrel{\Pi_{1}}{\vdots}$}\noLine
\UIC{$\forall k \in \mathbb{N}\,(\mathsf{o}(k) < \mathsf{o}(S(k)))$}
\AXC{$\varphi_{29}$\lbl{[S15]}}
\AXC{$\varphi_{10}$\lbl{DefSubIdx}}
\AXC{$\varphi_{10} \to \varphi_{14}$\lbl{A2}}
\BIC{$\varphi_{14}$}
\AXC{$\varphi_{14} \to \varphi_{30}$\lbl{A1}}
\BIC{$\varphi_{30}$}
\BIC{$\forall k \in \mathbb{N}\,(\mathsf{o}(k) < \mathsf{o}(S(k))) \to \mathrm{SubIdx}(\mathsf{o})$}
\BIC{$\mathrm{SubIdx}(\mathsf{o})$}
\end{ptree}
\noindent$\Pi_{3}$:
\begin{ptree}
\AXC{$k \in \mathbb{N} \to k \in \mathbb{N}$\lbl{A1}}
\AXC{$\varphi_{24}$\lbl{A2}}
\AXC{$\varphi_{29}$\lbl{[S15]}}
\AXC{$\varphi_{29} \to k \in \mathbb{N} \to \varphi_{16}$\lbl{A1}}
\BIC{$k \in \mathbb{N} \to \varphi_{16}$}
\AXC{$(k \in \mathbb{N} \to \varphi_{16}) \to \varphi_{24} \to \varphi_{19}$\lbl{A1}}
\BIC{$\varphi_{24} \to \varphi_{19}$}
\BIC{$\varphi_{19}$}
\AXC{$\varphi_{26}$\lbl{A1}}
\BIC{$(k \in \mathbb{N} \to k \in \mathbb{N}) \to k \in \mathbb{N} \to \mathsf{e}(k) < \mathsf{e}(S(k))$}
\BIC{$k \in \mathbb{N} \to \mathsf{e}(k) < \mathsf{e}(S(k))$}
\UIC{$\forall k \in \mathbb{N}\,(\mathsf{e}(k) < \mathsf{e}(S(k)))$}
\end{ptree}
\noindent$\Pi_{4}$:
\begin{ptree}
\AXC{$\stackrel{\Pi_{3}}{\vdots}$}\noLine
\UIC{$\forall k \in \mathbb{N}\,(\mathsf{e}(k) < \mathsf{e}(S(k)))$}
\AXC{$\varphi_{29}$\lbl{[S15]}}
\AXC{$\varphi_{10}$\lbl{DefSubIdx}}
\AXC{$\varphi_{10} \to \varphi_{15}$\lbl{A2}}
\BIC{$\varphi_{15}$}
\AXC{$\varphi_{15} \to \varphi_{31}$\lbl{A1}}
\BIC{$\varphi_{31}$}
\BIC{$\forall k \in \mathbb{N}\,(\mathsf{e}(k) < \mathsf{e}(S(k))) \to \mathrm{SubIdx}(\mathsf{e})$}
\BIC{$\mathrm{SubIdx}(\mathsf{e})$}
\end{ptree}
\noindent$\Pi_{5}$:
\begin{ptree}
\AXC{$\forall a\,(\varphi_{27})$\lbl{[C19]}}
\AXC{$\forall a\,(\varphi_{27}) \to \varphi_{27}$\lbl{A2}}
\BIC{$\varphi_{27}$}
\AXC{$\varphi_{27} \to \varphi_{28}$\lbl{A2}}
\BIC{$\varphi_{28}$}
\AXC{$\varphi_{28} \to \forall \alpha\,(\mathrm{Seq}(a) \to \mathrm{SubIdx}(\mathsf{e}) \to \varphi_{22})$\lbl{A2}}
\BIC{$\forall \alpha\,(\mathrm{Seq}(a) \to \mathrm{SubIdx}(\mathsf{e}) \to \varphi_{22})$}
\AXC{$\varphi_{32}$\lbl{A2}}
\BIC{$\mathrm{Seq}(a) \to \mathrm{SubIdx}(\mathsf{e}) \to \varphi_{22}$}
\end{ptree}
\noindent$\Pi_{6}$:
\begin{ptree}
\AXC{$\stackrel{\Pi_{4}}{\vdots}$}\noLine
\UIC{$\mathrm{SubIdx}(\mathsf{e})$}
\AXC{$\varphi_{20}$\lbl{A1}}
\AXC{$\varphi_{20} \to \varphi_{9}$\lbl{A1}}
\BIC{$\varphi_{9}$}
\AXC{$\stackrel{\Pi_{5}}{\vdots}$}\noLine
\UIC{$\mathrm{Seq}(a) \to \mathrm{SubIdx}(\mathsf{e}) \to \varphi_{22}$}
\AXC{$\varphi_{33}$\lbl{A1}}
\BIC{$\varphi_{9} \to \mathrm{SubIdx}(\mathsf{e}) \to \mathrm{SubIdx}(\mathsf{o}) \to \varphi_{21} \to \varphi_{18} \to \varphi_{17} \to \varphi_{11}$}
\BIC{$\mathrm{SubIdx}(\mathsf{e}) \to \mathrm{SubIdx}(\mathsf{o}) \to \varphi_{21} \to \varphi_{18} \to \varphi_{17} \to \varphi_{11}$}
\BIC{$\mathrm{SubIdx}(\mathsf{o}) \to \varphi_{21} \to \varphi_{18} \to \varphi_{17} \to \varphi_{11}$}
\end{ptree}
\noindent$\Pi_{7}$:
\begin{ptree}
\AXC{$\varphi_{20}$\lbl{A1}}
\AXC{$\varphi_{20} \to \varphi_{18}$\lbl{A1}}
\BIC{$\varphi_{18}$}
\AXC{$\varphi_{29}$\lbl{[S15]}}
\AXC{$\varphi_{29} \to \varphi_{21}$\lbl{A1}}
\BIC{$\varphi_{21}$}
\AXC{$\stackrel{\Pi_{2}}{\vdots}$}\noLine
\UIC{$\mathrm{SubIdx}(\mathsf{o})$}
\AXC{$\stackrel{\Pi_{6}}{\vdots}$}\noLine
\UIC{$\mathrm{SubIdx}(\mathsf{o}) \to \varphi_{21} \to \varphi_{18} \to \varphi_{17} \to \varphi_{11}$}
\BIC{$\varphi_{21} \to \varphi_{18} \to \varphi_{17} \to \varphi_{11}$}
\BIC{$\varphi_{18} \to \varphi_{17} \to \varphi_{11}$}
\BIC{$\varphi_{17} \to \varphi_{11}$}
\end{ptree}
\noindent$\Pi_{8}$:
\begin{ptree}
\AXC{$\forall a\,(\varphi_{6})$\lbl{[C18]}}
\AXC{$\forall a\,(\varphi_{6}) \to \varphi_{6}$\lbl{A2}}
\BIC{$\varphi_{6}$}
\AXC{$\varphi_{6} \to \forall \alpha\,(\varphi_{2})$\lbl{A2}}
\BIC{$\forall \alpha\,(\varphi_{2})$}
\AXC{$\forall \alpha\,(\varphi_{2}) \to \varphi_{2}$\lbl{A2}}
\BIC{$\varphi_{2}$}
\end{ptree}
\noindent$\Pi_{9}$:
\begin{ptree}
\AXC{$\stackrel{\Pi_{4}}{\vdots}$}\noLine
\UIC{$\mathrm{SubIdx}(\mathsf{e})$}
\AXC{$\mathrm{Conv}(a,\alpha) \to \mathrm{Conv}(a,\alpha)$\lbl{A1}}
\AXC{$\stackrel{\Pi_{8}}{\vdots}$}\noLine
\UIC{$\varphi_{2}$}
\AXC{$\varphi_{2} \to (\mathrm{Conv}(a,\alpha) \to \mathrm{Conv}(a,\alpha)) \to \varphi_{3}$\lbl{A1}}
\BIC{$(\mathrm{Conv}(a,\alpha) \to \mathrm{Conv}(a,\alpha)) \to \varphi_{3}$}
\BIC{$\varphi_{3}$}
\BIC{$\mathrm{Conv}(a,\alpha) \to \mathrm{Conv}(a \circ \mathsf{e},\alpha)$}
\end{ptree}
\noindent$\Pi_{10}$:
\begin{ptree}
\AXC{$\forall a\,(\varphi_{6})$\lbl{[C18]}}
\AXC{$\forall a\,(\varphi_{6}) \to \varphi_{6}$\lbl{A2}}
\BIC{$\varphi_{6}$}
\AXC{$\varphi_{6} \to \forall \alpha\,(\varphi_{4})$\lbl{A2}}
\BIC{$\forall \alpha\,(\varphi_{4})$}
\AXC{$\forall \alpha\,(\varphi_{4}) \to \varphi_{4}$\lbl{A2}}
\BIC{$\varphi_{4}$}
\end{ptree}
\noindent$\Pi_{11}$:
\begin{ptree}
\AXC{$\stackrel{\Pi_{2}}{\vdots}$}\noLine
\UIC{$\mathrm{SubIdx}(\mathsf{o})$}
\AXC{$\mathrm{Conv}(a,\alpha) \to \mathrm{Conv}(a,\alpha)$\lbl{A1}}
\AXC{$\stackrel{\Pi_{10}}{\vdots}$}\noLine
\UIC{$\varphi_{4}$}
\AXC{$\varphi_{4} \to (\mathrm{Conv}(a,\alpha) \to \mathrm{Conv}(a,\alpha)) \to \varphi_{5}$\lbl{A1}}
\BIC{$(\mathrm{Conv}(a,\alpha) \to \mathrm{Conv}(a,\alpha)) \to \varphi_{5}$}
\BIC{$\varphi_{5}$}
\BIC{$\mathrm{Conv}(a,\alpha) \to \mathrm{Conv}(a \circ \mathsf{o},\alpha)$}
\end{ptree}
\noindent$\Pi_{12}$:
\begin{ptree}
\AXC{$\forall a\,\forall \alpha\,(\mathrm{Conv}(a,\alpha) \to \mathrm{Seq}(a))$\lbl{[C1]}}
\AXC{$\varphi_{7}$\lbl{A2}}
\BIC{$\forall \alpha\,(\mathrm{Conv}(a,\alpha) \to \mathrm{Seq}(a))$}
\AXC{$\varphi_{1}$\lbl{A2}}
\BIC{$\mathrm{Conv}(a,\alpha) \to \mathrm{Seq}(a)$}
\AXC{$(\mathrm{Conv}(a,\alpha) \to \mathrm{Seq}(a)) \to \varphi_{23}$\lbl{A1}}
\BIC{$\varphi_{23}$}
\end{ptree}
\noindent$\Pi_{13}$:
\begin{ptree}
\AXC{$\stackrel{\Pi_{9}}{\vdots}$}\noLine
\UIC{$\mathrm{Conv}(a,\alpha) \to \mathrm{Conv}(a \circ \mathsf{e},\alpha)$}
\AXC{$\stackrel{\Pi_{11}}{\vdots}$}\noLine
\UIC{$\mathrm{Conv}(a,\alpha) \to \mathrm{Conv}(a \circ \mathsf{o},\alpha)$}
\AXC{$\stackrel{\Pi_{12}}{\vdots}$}\noLine
\UIC{$\varphi_{23}$}
\BIC{$(\mathrm{Conv}(a,\alpha) \to \mathrm{Conv}(a \circ \mathsf{e},\alpha)) \to \varphi_{12}$}
\BIC{$\varphi_{12}$}
\AXC{$\varphi_{12} \to \varphi_{11} \to \varphi_{13}$\lbl{A1}}
\BIC{$\varphi_{11} \to \varphi_{13}$}
\end{ptree}
\noindent 主木:
\begin{ptree}
\AXC{$\varphi_{20}$\lbl{A1}}
\AXC{$\varphi_{20} \to \varphi_{17}$\lbl{A1}}
\BIC{$\varphi_{17}$}
\AXC{$\stackrel{\Pi_{7}}{\vdots}$}\noLine
\UIC{$\varphi_{17} \to \varphi_{11}$}
\BIC{$\varphi_{11}$}
\AXC{$\stackrel{\Pi_{13}}{\vdots}$}\noLine
\UIC{$\varphi_{11} \to \varphi_{13}$}
\BIC{$\varphi_{13}$}
\UIC{$\forall \alpha\,(\varphi_{13})$}
\UIC{$\forall a\,\forall \alpha\,(\varphi_{13})$}
\end{ptree}
\begin{abbrevlist}
\item $\varphi_{1} :\equiv \forall \alpha\,(\mathrm{Conv}(a,\alpha) \to \mathrm{Seq}(a)) \to \mathrm{Conv}(a,\alpha) \to \mathrm{Seq}(a)$
\item $\varphi_{2} :\equiv \mathrm{Conv}(a,\alpha) \to \mathrm{SubIdx}(\mathsf{e}) \to \mathrm{Conv}(a \circ \mathsf{e},\alpha)$
\item $\varphi_{3} :\equiv \mathrm{SubIdx}(\mathsf{e}) \to \mathrm{Conv}(a,\alpha) \to \mathrm{Conv}(a \circ \mathsf{e},\alpha)$
\item $\varphi_{4} :\equiv \mathrm{Conv}(a,\alpha) \to \mathrm{SubIdx}(\mathsf{o}) \to \mathrm{Conv}(a \circ \mathsf{o},\alpha)$
\item $\varphi_{5} :\equiv \mathrm{SubIdx}(\mathsf{o}) \to \mathrm{Conv}(a,\alpha) \to \mathrm{Conv}(a \circ \mathsf{o},\alpha)$
\item $\varphi_{6} :\equiv \forall s\,\forall \alpha\,(\mathrm{Conv}(a,\alpha) \to \mathrm{SubIdx}(s) \to \mathrm{Conv}(a \circ s,\alpha))$
\item $\varphi_{7} :\equiv \forall a\,\forall \alpha\,(\mathrm{Conv}(a,\alpha) \to \mathrm{Seq}(a)) \to \forall \alpha\,(\mathrm{Conv}(a,\alpha) \to \mathrm{Seq}(a))$
\item $\varphi_{8} :\equiv \mathrm{Conv}(a \circ \mathsf{e},\alpha) \to \mathrm{Conv}(a \circ \mathsf{o},\alpha) \to \mathrm{Conv}(a,\alpha)$
\item $\varphi_{9} :\equiv \mathrm{Seq}(a) \land \mathrm{Conv}(a \circ \mathsf{o},\alpha) \land \mathrm{Conv}(a \circ \mathsf{e},\alpha) \to \mathrm{Seq}(a)$
\item $\varphi_{10} :\equiv \forall s\,(\mathrm{SubIdx}(s) \leftrightarrow \mathrm{NSeq}(s) \land \forall k \in \mathbb{N}\,(s_{k} < s_{S(k)}))$
\item $\varphi_{11} :\equiv \mathrm{Seq}(a) \land \mathrm{Conv}(a \circ \mathsf{o},\alpha) \land \mathrm{Conv}(a \circ \mathsf{e},\alpha) \to \mathrm{Conv}(a,\alpha)$
\item $\varphi_{12} :\equiv \mathrm{Conv}(a,\alpha) \to \mathrm{Seq}(a) \land \mathrm{Conv}(a \circ \mathsf{o},\alpha) \land \mathrm{Conv}(a \circ \mathsf{e},\alpha)$
\item $\varphi_{13} :\equiv \mathrm{Conv}(a,\alpha) \leftrightarrow \mathrm{Seq}(a) \land \mathrm{Conv}(a \circ \mathsf{o},\alpha) \land \mathrm{Conv}(a \circ \mathsf{e},\alpha)$
\item $\varphi_{14} :\equiv \mathrm{SubIdx}(\mathsf{o}) \leftrightarrow \mathrm{NSeq}(\mathsf{o}) \land \forall k \in \mathbb{N}\,(\mathsf{o}(k) < \mathsf{o}(S(k)))$
\item $\varphi_{15} :\equiv \mathrm{SubIdx}(\mathsf{e}) \leftrightarrow \mathrm{NSeq}(\mathsf{e}) \land \forall k \in \mathbb{N}\,(\mathsf{e}(k) < \mathsf{e}(S(k)))$
\item $\varphi_{16} :\equiv \forall k \in \mathbb{N}\,(\mathsf{e}(k) < \mathsf{e}(S(k)) \land \mathsf{o}(k) < \mathsf{o}(S(k)))$
\item $\varphi_{17} :\equiv \mathrm{Seq}(a) \land \mathrm{Conv}(a \circ \mathsf{o},\alpha) \land \mathrm{Conv}(a \circ \mathsf{e},\alpha) \to \mathrm{Conv}(a \circ \mathsf{o},\alpha)$
\item $\varphi_{18} :\equiv \mathrm{Seq}(a) \land \mathrm{Conv}(a \circ \mathsf{o},\alpha) \land \mathrm{Conv}(a \circ \mathsf{e},\alpha) \to \mathrm{Conv}(a \circ \mathsf{e},\alpha)$
\item $\varphi_{19} :\equiv k \in \mathbb{N} \to k \in \mathbb{N} \to \mathsf{e}(k) < \mathsf{e}(S(k)) \land \mathsf{o}(k) < \mathsf{o}(S(k))$
\item $\varphi_{20} :\equiv \mathrm{Seq}(a) \land \mathrm{Conv}(a \circ \mathsf{o},\alpha) \land \mathrm{Conv}(a \circ \mathsf{e},\alpha) \to \mathrm{Seq}(a) \land \mathrm{Conv}(a \circ \mathsf{o},\alpha) \land \mathrm{Conv}(a \circ \mathsf{e},\alpha)$
\item $\varphi_{21} :\equiv \mathrm{Seq}(a) \land \mathrm{Conv}(a \circ \mathsf{o},\alpha) \land \mathrm{Conv}(a \circ \mathsf{e},\alpha) \to \forall n \in \mathbb{N}\,\exists k \in \mathbb{N}\,(n = \mathsf{e}(k) \lor n = \mathsf{o}(k))$
\item $\varphi_{22} :\equiv \mathrm{SubIdx}(\mathsf{o}) \to \forall n \in \mathbb{N}\,\exists k \in \mathbb{N}\,(n = \mathsf{e}(k) \lor n = \mathsf{o}(k)) \to \varphi_{8}$
\item $\varphi_{23} :\equiv (\mathrm{Conv}(a,\alpha) \to \mathrm{Conv}(a \circ \mathsf{o},\alpha)) \to (\mathrm{Conv}(a,\alpha) \to \mathrm{Conv}(a \circ \mathsf{e},\alpha)) \to \varphi_{12}$
\item $\varphi_{24} :\equiv \varphi_{16} \to k \in \mathbb{N} \to \mathsf{e}(k) < \mathsf{e}(S(k)) \land \mathsf{o}(k) < \mathsf{o}(S(k))$
\item $\varphi_{25} :\equiv \varphi_{19} \to (k \in \mathbb{N} \to k \in \mathbb{N}) \to k \in \mathbb{N} \to \mathsf{o}(k) < \mathsf{o}(S(k))$
\item $\varphi_{26} :\equiv \varphi_{19} \to (k \in \mathbb{N} \to k \in \mathbb{N}) \to k \in \mathbb{N} \to \mathsf{e}(k) < \mathsf{e}(S(k))$
\item $\varphi_{27} :\equiv \forall s\,\forall t\,\forall \alpha\,(\mathrm{Seq}(a) \to \mathrm{SubIdx}(s) \to \mathrm{SubIdx}(t) \to \forall n \in \mathbb{N}\,\exists k \in \mathbb{N}\,(n = s_{k} \lor n = t(k)) \to \mathrm{Conv}(a \circ s,\alpha) \to \mathrm{Conv}(a \circ t,\alpha) \to \mathrm{Conv}(a,\alpha))$
\item $\varphi_{28} :\equiv \forall t\,\forall \alpha\,(\mathrm{Seq}(a) \to \mathrm{SubIdx}(\mathsf{e}) \to \mathrm{SubIdx}(t) \to \forall n \in \mathbb{N}\,\exists k \in \mathbb{N}\,(n = \mathsf{e}(k) \lor n = t(k)) \to \mathrm{Conv}(a \circ \mathsf{e},\alpha) \to \mathrm{Conv}(a \circ t,\alpha) \to \mathrm{Conv}(a,\alpha))$
\item $\varphi_{29} :\equiv \mathrm{NSeq}(\mathsf{e}) \land \mathrm{NSeq}(\mathsf{o}) \land \varphi_{16} \land \forall n \in \mathbb{N}\,\exists k \in \mathbb{N}\,(n = \mathsf{e}(k) \lor n = \mathsf{o}(k))$
\item $\varphi_{30} :\equiv \varphi_{29} \to \forall k \in \mathbb{N}\,(\mathsf{o}(k) < \mathsf{o}(S(k))) \to \mathrm{SubIdx}(\mathsf{o})$
\item $\varphi_{31} :\equiv \varphi_{29} \to \forall k \in \mathbb{N}\,(\mathsf{e}(k) < \mathsf{e}(S(k))) \to \mathrm{SubIdx}(\mathsf{e})$
\item $\varphi_{32} :\equiv \forall \alpha\,(\mathrm{Seq}(a) \to \mathrm{SubIdx}(\mathsf{e}) \to \varphi_{22}) \to \mathrm{Seq}(a) \to \mathrm{SubIdx}(\mathsf{e}) \to \varphi_{22}$
\item $\varphi_{33} :\equiv (\mathrm{Seq}(a) \to \mathrm{SubIdx}(\mathsf{e}) \to \varphi_{22}) \to \varphi_{9} \to \mathrm{SubIdx}(\mathsf{e}) \to \mathrm{SubIdx}(\mathsf{o}) \to \varphi_{21} \to \varphi_{18} \to \varphi_{17} \to \varphi_{11}$
\end{abbrevlist}
\end{fullproof}

% 自動生成 (prover/build.js): SubZero  ノード数 51
\begin{fullproof}{[E28]}{$\forall x \in \mathbb{R}\,(x - 0 = x)$}{51}\phantomsection\label{tree:SubZero}
\noindent$\Pi_{1}$:
\begin{ptree}
\AXC{$y = z \to x = y \to x = z$\lbl{A6}}
\AXC{$(y = z \to x = y \to x = z) \to x = y \to y = z \to x = z$\lbl{A1}}
\BIC{$x = y \to y = z \to x = z$}
\UIC{$\forall z\,(x = y \to y = z \to x = z)$}
\UIC{$\forall y\,\forall z\,(x = y \to y = z \to x = z)$}
\UIC{$\forall x\,\forall y\,\forall z\,(x = y \to y = z \to x = z)$}
\AXC{$\varphi_{1}$\lbl{A2}}
\BIC{$\forall y\,\forall z\,(x - 0 = y \to y = z \to x - 0 = z)$}
\AXC{$\varphi_{3}$\lbl{A2}}
\BIC{$\forall z\,(x - 0 = x + 0 \to x + 0 = z \to x - 0 = z)$}
\AXC{$\varphi_{4}$\lbl{A2}}
\BIC{$x - 0 = x + 0 \to x + 0 = x \to x - 0 = x$}
\end{ptree}
\noindent$\Pi_{2}$:
\begin{ptree}
\AXC{$y = z \to x = y \to x = z$\lbl{A6}}
\AXC{$(y = z \to x = y \to x = z) \to x = y \to y = z \to x = z$\lbl{A1}}
\BIC{$x = y \to y = z \to x = z$}
\UIC{$\forall z\,(x = y \to y = z \to x = z)$}
\UIC{$\forall y\,\forall z\,(x = y \to y = z \to x = z)$}
\UIC{$\forall x\,\forall y\,\forall z\,(x = y \to y = z \to x = z)$}
\AXC{$\varphi_{1}$\lbl{A2}}
\BIC{$\forall y\,\forall z\,(x - 0 = y \to y = z \to x - 0 = z)$}
\AXC{$\varphi_{5}$\lbl{A2}}
\BIC{$\forall z\,(x - 0 = x + -0 \to x + -0 = z \to x - 0 = z)$}
\AXC{$\varphi_{8}$\lbl{A2}}
\BIC{$x - 0 = x + -0 \to x + -0 = x + 0 \to x - 0 = x + 0$}
\end{ptree}
\noindent$\Pi_{3}$:
\begin{ptree}
\AXC{$x = x$\lbl{A5}}
\UIC{$\forall x\,(x = x)$}
\AXC{$\forall x\,(x = x) \to x + -0 = x + -0$\lbl{A2}}
\BIC{$x + -0 = x + -0$}
\AXC{$-0 = 0$\lbl{[F6]}}
\AXC{$-0 = 0 \to x + -0 = x + -0 \to \varphi_{7}$\lbl{A1}}
\BIC{$x + -0 = x + -0 \to \varphi_{7}$}
\BIC{$\varphi_{7}$}
\end{ptree}
\noindent$\Pi_{4}$:
\begin{ptree}
\AXC{$\forall x\,\forall y\,(x - y = x + -y)$\lbl{Def$-$}}
\AXC{$\forall x\,\forall y\,(x - y = x + -y) \to \forall y\,(x - y = x + -y)$\lbl{A2}}
\BIC{$\forall y\,(x - y = x + -y)$}
\end{ptree}
\noindent$\Pi_{5}$:
\begin{ptree}
\AXC{$\stackrel{\Pi_{4}}{\vdots}$}\noLine
\UIC{$\forall y\,(x - y = x + -y)$}
\AXC{$\forall y\,(x - y = x + -y) \to x - 0 = x + -0$\lbl{A2}}
\BIC{$x - 0 = x + -0$}
\AXC{$x - 0 = x + -0 \to (x \in \mathbb{R} \to x + -0 = x + 0) \to \varphi_{6}$\lbl{A1}}
\BIC{$(x \in \mathbb{R} \to x + -0 = x + 0) \to \varphi_{6}$}
\end{ptree}
\noindent$\Pi_{6}$:
\begin{ptree}
\AXC{$z = w \to x + -0 = x + z \to x + -0 = x + w$\lbl{A6}}
\UIC{$\forall w\,(z = w \to x + -0 = x + z \to x + -0 = x + w)$}
\UIC{$\forall z\,\forall w\,(z = w \to x + -0 = x + z \to x + -0 = x + w)$}
\AXC{$\varphi_{9}$\lbl{A2}}
\BIC{$\forall w\,(-0 = w \to x + -0 = x + -0 \to x + -0 = x + w)$}
\AXC{$\varphi_{10}$\lbl{A2}}
\BIC{$-0 = 0 \to x + -0 = x + -0 \to x + -0 = x + 0$}
\AXC{$\stackrel{\Pi_{3}}{\vdots}$}\noLine
\UIC{$\varphi_{7}$}
\BIC{$x \in \mathbb{R} \to x + -0 = x + 0$}
\AXC{$\stackrel{\Pi_{5}}{\vdots}$}\noLine
\UIC{$(x \in \mathbb{R} \to x + -0 = x + 0) \to \varphi_{6}$}
\BIC{$\varphi_{6}$}
\end{ptree}
\noindent$\Pi_{7}$:
\begin{ptree}
\AXC{$\stackrel{\Pi_{2}}{\vdots}$}\noLine
\UIC{$x - 0 = x + -0 \to x + -0 = x + 0 \to x - 0 = x + 0$}
\AXC{$\stackrel{\Pi_{6}}{\vdots}$}\noLine
\UIC{$\varphi_{6}$}
\BIC{$x \in \mathbb{R} \to x - 0 = x + 0$}
\AXC{$(x \in \mathbb{R} \to x - 0 = x + 0) \to (x \in \mathbb{R} \to x + 0 = x) \to \varphi_{2}$\lbl{A1}}
\BIC{$(x \in \mathbb{R} \to x + 0 = x) \to \varphi_{2}$}
\end{ptree}
\noindent$\Pi_{8}$:
\begin{ptree}
\AXC{$\forall x \in \mathbb{R}\,(x + 0 = x)$\lbl{(R4)}}
\AXC{$\forall x \in \mathbb{R}\,(x + 0 = x) \to x \in \mathbb{R} \to x + 0 = x$\lbl{A2}}
\BIC{$x \in \mathbb{R} \to x + 0 = x$}
\AXC{$\stackrel{\Pi_{7}}{\vdots}$}\noLine
\UIC{$(x \in \mathbb{R} \to x + 0 = x) \to \varphi_{2}$}
\BIC{$\varphi_{2}$}
\end{ptree}
\noindent 主木:
\begin{ptree}
\AXC{$\stackrel{\Pi_{1}}{\vdots}$}\noLine
\UIC{$x - 0 = x + 0 \to x + 0 = x \to x - 0 = x$}
\AXC{$\stackrel{\Pi_{8}}{\vdots}$}\noLine
\UIC{$\varphi_{2}$}
\BIC{$x \in \mathbb{R} \to x - 0 = x$}
\UIC{$\forall x \in \mathbb{R}\,(x - 0 = x)$}
\end{ptree}
\begin{abbrevlist}
\item $\varphi_{1} :\equiv \forall x\,\forall y\,\forall z\,(x = y \to y = z \to x = z) \to \forall y\,\forall z\,(x - 0 = y \to y = z \to x - 0 = z)$
\item $\varphi_{2} :\equiv (x - 0 = x + 0 \to x + 0 = x \to x - 0 = x) \to x \in \mathbb{R} \to x - 0 = x$
\item $\varphi_{3} :\equiv \forall y\,\forall z\,(x - 0 = y \to y = z \to x - 0 = z) \to \forall z\,(x - 0 = x + 0 \to x + 0 = z \to x - 0 = z)$
\item $\varphi_{4} :\equiv \forall z\,(x - 0 = x + 0 \to x + 0 = z \to x - 0 = z) \to x - 0 = x + 0 \to x + 0 = x \to x - 0 = x$
\item $\varphi_{5} :\equiv \forall y\,\forall z\,(x - 0 = y \to y = z \to x - 0 = z) \to \forall z\,(x - 0 = x + -0 \to x + -0 = z \to x - 0 = z)$
\item $\varphi_{6} :\equiv (x - 0 = x + -0 \to x + -0 = x + 0 \to x - 0 = x + 0) \to x \in \mathbb{R} \to x - 0 = x + 0$
\item $\varphi_{7} :\equiv (-0 = 0 \to x + -0 = x + -0 \to x + -0 = x + 0) \to x \in \mathbb{R} \to x + -0 = x + 0$
\item $\varphi_{8} :\equiv \forall z\,(x - 0 = x + -0 \to x + -0 = z \to x - 0 = z) \to x - 0 = x + -0 \to x + -0 = x + 0 \to x - 0 = x + 0$
\item $\varphi_{9} :\equiv \forall z\,\forall w\,(z = w \to x + -0 = x + z \to x + -0 = x + w) \to \forall w\,(-0 = w \to x + -0 = x + -0 \to x + -0 = x + w)$
\item $\varphi_{10} :\equiv \forall w\,(-0 = w \to x + -0 = x + -0 \to x + -0 = x + w) \to -0 = 0 \to x + -0 = x + -0 \to x + -0 = x + 0$
\end{abbrevlist}
\end{fullproof}

% 自動生成 (prover/build.js): AbsAbs  ノード数 10
\begin{fullproof}{[E29]}{$\forall x \in \mathbb{R}\,(\lvert \lvert x \rvert \rvert = \lvert x \rvert)$}{10}\phantomsection\label{tree:AbsAbs}
\noindent$\Pi_{1}$:
\begin{ptree}
\AXC{$\forall x \in \mathbb{R}\,(0 \leq \lvert x \rvert)$\lbl{[A6]}}
\AXC{$\forall x \in \mathbb{R}\,(0 \leq \lvert x \rvert) \to x \in \mathbb{R} \to 0 \leq \lvert x \rvert$\lbl{A2}}
\BIC{$x \in \mathbb{R} \to 0 \leq \lvert x \rvert$}
\end{ptree}
\noindent$\Pi_{2}$:
\begin{ptree}
\AXC{$\forall x\,(0 \leq x \to \lvert x \rvert = x)$\lbl{[A1]}}
\AXC{$\forall x\,(0 \leq x \to \lvert x \rvert = x) \to 0 \leq \lvert x \rvert \to \lvert \lvert x \rvert \rvert = \lvert x \rvert$\lbl{A2}}
\BIC{$0 \leq \lvert x \rvert \to \lvert \lvert x \rvert \rvert = \lvert x \rvert$}
\AXC{$\varphi_{1}$\lbl{A1}}
\BIC{$(x \in \mathbb{R} \to 0 \leq \lvert x \rvert) \to x \in \mathbb{R} \to \lvert \lvert x \rvert \rvert = \lvert x \rvert$}
\end{ptree}
\noindent 主木:
\begin{ptree}
\AXC{$\stackrel{\Pi_{1}}{\vdots}$}\noLine
\UIC{$x \in \mathbb{R} \to 0 \leq \lvert x \rvert$}
\AXC{$\stackrel{\Pi_{2}}{\vdots}$}\noLine
\UIC{$(x \in \mathbb{R} \to 0 \leq \lvert x \rvert) \to x \in \mathbb{R} \to \lvert \lvert x \rvert \rvert = \lvert x \rvert$}
\BIC{$x \in \mathbb{R} \to \lvert \lvert x \rvert \rvert = \lvert x \rvert$}
\UIC{$\forall x \in \mathbb{R}\,(\lvert \lvert x \rvert \rvert = \lvert x \rvert)$}
\end{ptree}
\begin{abbrevlist}
\item $\varphi_{1} :\equiv (0 \leq \lvert x \rvert \to \lvert \lvert x \rvert \rvert = \lvert x \rvert) \to (x \in \mathbb{R} \to 0 \leq \lvert x \rvert) \to x \in \mathbb{R} \to \lvert \lvert x \rvert \rvert = \lvert x \rvert$
\end{abbrevlist}
\end{fullproof}

% 自動生成 (prover/build.js): LimAbs  ノード数 154
\begin{fullproof}{定理2.15(1)}{$\forall a\,\forall \alpha\,(\mathrm{Conv}(a,\alpha) \to \mathrm{Conv}(\lvert a \rvert,\lvert \alpha \rvert))$}{154}\phantomsection\label{tree:LimAbs}
\noindent$\Pi_{1}$:
\begin{ptree}
\AXC{$z = w \to \lvert z - \lvert \alpha \rvert \rvert < \varepsilon \to \lvert w - \lvert \alpha \rvert \rvert < \varepsilon$\lbl{A6}}
\UIC{$\forall w\,(z = w \to \lvert z - \lvert \alpha \rvert \rvert < \varepsilon \to \lvert w - \lvert \alpha \rvert \rvert < \varepsilon)$}
\UIC{$\forall z\,\forall w\,(z = w \to \lvert z - \lvert \alpha \rvert \rvert < \varepsilon \to \lvert w - \lvert \alpha \rvert \rvert < \varepsilon)$}
\AXC{$\varphi_{47}$\lbl{A2}}
\BIC{$\varphi_{18}$}
\AXC{$\varphi_{18} \to \varphi_{28}$\lbl{A2}}
\BIC{$\varphi_{28}$}
\end{ptree}
\noindent$\Pi_{2}$:
\begin{ptree}
\AXC{$\varphi_{23}$\lbl{A2}}
\AXC{$\mathrm{Conv}(a,\alpha) \to \varepsilon \in \mathbb{R} \to 0 < \varepsilon \to \varphi_{38}$\lbl{A1}}
\AXC{$\varphi_{67}$\lbl{A1}}
\BIC{$\mathrm{Conv}(a,\alpha) \to \varepsilon \in \mathbb{R} \to 0 < \varepsilon \to \varphi_{41}$}
\AXC{$\varphi_{73}$\lbl{A1}}
\BIC{$\varphi_{23} \to \mathrm{Conv}(a,\alpha) \to \varepsilon \in \mathbb{R} \to 0 < \varepsilon \to \varphi_{40}$}
\BIC{$\mathrm{Conv}(a,\alpha) \to \varepsilon \in \mathbb{R} \to 0 < \varepsilon \to \varphi_{40}$}
\AXC{$\varphi_{75}$\lbl{A1}}
\BIC{$(\mathrm{Conv}(a,\alpha) \to \varepsilon \in \mathbb{R} \to 0 < \varepsilon \to \varphi_{17}) \to \varphi_{60}$}
\end{ptree}
\noindent$\Pi_{3}$:
\begin{ptree}
\AXC{$\mathrm{Conv}(a,\alpha) \to \varepsilon \in \mathbb{R} \to 0 < \varepsilon \to \varphi_{17}$\lbl{A1}}
\AXC{$\stackrel{\Pi_{2}}{\vdots}$}\noLine
\UIC{$(\mathrm{Conv}(a,\alpha) \to \varepsilon \in \mathbb{R} \to 0 < \varepsilon \to \varphi_{17}) \to \varphi_{60}$}
\BIC{$\varphi_{60}$}
\end{ptree}
\noindent$\Pi_{4}$:
\begin{ptree}
\AXC{$\forall a\,\forall n\,(\mathrm{Seq}(a) \to n \in \mathbb{N} \to a_{n} \in \mathbb{R})$\lbl{[S1]}}
\AXC{$\varphi_{21}$\lbl{A2}}
\BIC{$\forall n\,(\mathrm{Seq}(a) \to n \in \mathbb{N} \to a_{n} \in \mathbb{R})$}
\AXC{$\varphi_{12}$\lbl{A2}}
\BIC{$\mathrm{Seq}(a) \to n \in \mathbb{N} \to a_{n} \in \mathbb{R}$}
\AXC{$\varphi_{66}$\lbl{A1}}
\BIC{$(\mathrm{Conv}(a,\alpha) \to \mathrm{Seq}(a)) \to \varphi_{59}$}
\end{ptree}
\noindent$\Pi_{5}$:
\begin{ptree}
\AXC{$\forall a\,\forall \alpha\,(\mathrm{Conv}(a,\alpha) \to \mathrm{Seq}(a))$\lbl{[C1]}}
\AXC{$\varphi_{2}$\lbl{A2}}
\BIC{$\forall \alpha\,(\mathrm{Conv}(a,\alpha) \to \mathrm{Seq}(a))$}
\AXC{$\varphi_{1}$\lbl{A2}}
\BIC{$\mathrm{Conv}(a,\alpha) \to \mathrm{Seq}(a)$}
\AXC{$\stackrel{\Pi_{4}}{\vdots}$}\noLine
\UIC{$(\mathrm{Conv}(a,\alpha) \to \mathrm{Seq}(a)) \to \varphi_{59}$}
\BIC{$\varphi_{59}$}
\end{ptree}
\noindent$\Pi_{6}$:
\begin{ptree}
\AXC{$\varphi_{8}$\lbl{[E27]}}
\AXC{$\varphi_{8} \to \varphi_{14}$\lbl{A2}}
\BIC{$\varphi_{14}$}
\AXC{$\varphi_{14} \to a_{n} \in \mathbb{R} \to \alpha \in \mathbb{R} \to \lvert \lvert a_{n} \rvert - \lvert \alpha \rvert \rvert \leq \lvert a_{n} - \alpha \rvert$\lbl{A2}}
\BIC{$a_{n} \in \mathbb{R} \to \alpha \in \mathbb{R} \to \lvert \lvert a_{n} \rvert - \lvert \alpha \rvert \rvert \leq \lvert a_{n} - \alpha \rvert$}
\AXC{$\varphi_{71}$\lbl{A1}}
\BIC{$(\mathrm{Conv}(a,\alpha) \to \varepsilon \in \mathbb{R} \to 0 < \varepsilon \to \varphi_{24}) \to \varphi_{53}$}
\end{ptree}
\noindent$\Pi_{7}$:
\begin{ptree}
\AXC{$\mathrm{Conv}(a,\alpha) \to \varepsilon \in \mathbb{R} \to 0 < \varepsilon \to \varphi_{17}$\lbl{A1}}
\AXC{$\stackrel{\Pi_{5}}{\vdots}$}\noLine
\UIC{$\varphi_{59}$}
\BIC{$\mathrm{Conv}(a,\alpha) \to \varepsilon \in \mathbb{R} \to 0 < \varepsilon \to \varphi_{24}$}
\AXC{$\stackrel{\Pi_{6}}{\vdots}$}\noLine
\UIC{$(\mathrm{Conv}(a,\alpha) \to \varepsilon \in \mathbb{R} \to 0 < \varepsilon \to \varphi_{24}) \to \varphi_{53}$}
\BIC{$\varphi_{53}$}
\end{ptree}
\noindent$\Pi_{8}$:
\begin{ptree}
\AXC{$\forall a\,\forall \alpha\,(\mathrm{Conv}(a,\alpha) \to \alpha \in \mathbb{R})$\lbl{[C2]}}
\AXC{$\varphi_{6}$\lbl{A2}}
\BIC{$\forall \alpha\,(\mathrm{Conv}(a,\alpha) \to \alpha \in \mathbb{R})$}
\AXC{$\varphi_{3}$\lbl{A2}}
\BIC{$\mathrm{Conv}(a,\alpha) \to \alpha \in \mathbb{R}$}
\AXC{$\stackrel{\Pi_{7}}{\vdots}$}\noLine
\UIC{$\varphi_{53}$}
\BIC{$\mathrm{Conv}(a,\alpha) \to \varepsilon \in \mathbb{R} \to 0 < \varepsilon \to \varphi_{42}$}
\end{ptree}
\noindent$\Pi_{9}$:
\begin{ptree}
\AXC{$\forall x\,\forall y\,\forall z\,(x \leq y \to y < z \to x < z)$\lbl{[O4]}}
\AXC{$\forall x\,\forall y\,\forall z\,(x \leq y \to y < z \to x < z) \to \varphi_{15}$\lbl{A2}}
\BIC{$\varphi_{15}$}
\AXC{$\varphi_{15} \to \varphi_{32}$\lbl{A2}}
\BIC{$\varphi_{32}$}
\AXC{$\varphi_{32} \to \varphi_{34}$\lbl{A2}}
\BIC{$\varphi_{34}$}
\end{ptree}
\noindent$\Pi_{10}$:
\begin{ptree}
\AXC{$\stackrel{\Pi_{8}}{\vdots}$}\noLine
\UIC{$\mathrm{Conv}(a,\alpha) \to \varepsilon \in \mathbb{R} \to 0 < \varepsilon \to \varphi_{42}$}
\AXC{$\stackrel{\Pi_{9}}{\vdots}$}\noLine
\UIC{$\varphi_{34}$}
\AXC{$\varphi_{34} \to (\mathrm{Conv}(a,\alpha) \to \varepsilon \in \mathbb{R} \to 0 < \varepsilon \to \varphi_{42}) \to \varphi_{63}$\lbl{A1}}
\BIC{$(\mathrm{Conv}(a,\alpha) \to \varepsilon \in \mathbb{R} \to 0 < \varepsilon \to \varphi_{42}) \to \varphi_{63}$}
\BIC{$\varphi_{63}$}
\end{ptree}
\noindent$\Pi_{11}$:
\begin{ptree}
\AXC{$x = x$\lbl{A5}}
\AXC{$x = y \to x = x \to y = x$\lbl{A6}}
\AXC{$(x = y \to x = x \to y = x) \to x = x \to x = y \to y = x$\lbl{A1}}
\BIC{$x = x \to x = y \to y = x$}
\BIC{$x = y \to y = x$}
\UIC{$\forall y\,(x = y \to y = x)$}
\UIC{$\forall x\,\forall y\,(x = y \to y = x)$}
\AXC{$\varphi_{4}$\lbl{A2}}
\BIC{$\forall y\,(\lvert a \rvert(n) = y \to y = \lvert a \rvert(n))$}
\AXC{$\varphi_{26}$\lbl{A2}}
\BIC{$\lvert a \rvert(n) = \lvert a_{n} \rvert \to \lvert a_{n} \rvert = \lvert a \rvert(n)$}
\end{ptree}
\noindent$\Pi_{12}$:
\begin{ptree}
\AXC{$\forall a\,(\varphi_{9})$\lbl{[S11]}}
\AXC{$\forall a\,(\varphi_{9}) \to \varphi_{9}$\lbl{A2}}
\BIC{$\varphi_{9}$}
\AXC{$\varphi_{9} \to (\mathrm{Conv}(a,\alpha) \to \mathrm{Seq}(a)) \to \varphi_{10}$\lbl{A1}}
\BIC{$(\mathrm{Conv}(a,\alpha) \to \mathrm{Seq}(a)) \to \varphi_{10}$}
\end{ptree}
\noindent$\Pi_{13}$:
\begin{ptree}
\AXC{$\forall a\,\forall \alpha\,(\mathrm{Conv}(a,\alpha) \to \mathrm{Seq}(a))$\lbl{[C1]}}
\AXC{$\varphi_{2}$\lbl{A2}}
\BIC{$\forall \alpha\,(\mathrm{Conv}(a,\alpha) \to \mathrm{Seq}(a))$}
\AXC{$\varphi_{1}$\lbl{A2}}
\BIC{$\mathrm{Conv}(a,\alpha) \to \mathrm{Seq}(a)$}
\AXC{$\stackrel{\Pi_{12}}{\vdots}$}\noLine
\UIC{$(\mathrm{Conv}(a,\alpha) \to \mathrm{Seq}(a)) \to \varphi_{10}$}
\BIC{$\varphi_{10}$}
\end{ptree}
\noindent$\Pi_{14}$:
\begin{ptree}
\AXC{$\varphi_{20}$\lbl{A2}}
\AXC{$\stackrel{\Pi_{13}}{\vdots}$}\noLine
\UIC{$\varphi_{10}$}
\AXC{$\varphi_{10} \to \mathrm{Conv}(a,\alpha) \to \varepsilon \in \mathbb{R} \to 0 < \varepsilon \to \varphi_{39}$\lbl{A1}}
\BIC{$\mathrm{Conv}(a,\alpha) \to \varepsilon \in \mathbb{R} \to 0 < \varepsilon \to \varphi_{39}$}
\AXC{$\varphi_{70}$\lbl{A1}}
\BIC{$\varphi_{20} \to \mathrm{Conv}(a,\alpha) \to \varepsilon \in \mathbb{R} \to 0 < \varepsilon \to \varphi_{36}$}
\BIC{$\mathrm{Conv}(a,\alpha) \to \varepsilon \in \mathbb{R} \to 0 < \varepsilon \to \varphi_{36}$}
\end{ptree}
\noindent$\Pi_{15}$:
\begin{ptree}
\AXC{$\stackrel{\Pi_{14}}{\vdots}$}\noLine
\UIC{$\mathrm{Conv}(a,\alpha) \to \varepsilon \in \mathbb{R} \to 0 < \varepsilon \to \varphi_{36}$}
\AXC{$(\mathrm{Conv}(a,\alpha) \to \varepsilon \in \mathbb{R} \to 0 < \varepsilon \to \varphi_{36}) \to \varphi_{61}$\lbl{A1}}
\BIC{$\varphi_{61}$}
\end{ptree}
\noindent$\Pi_{16}$:
\begin{ptree}
\AXC{$\mathrm{Conv}(a,\alpha) \to \varepsilon \in \mathbb{R} \to 0 < \varepsilon \to \varphi_{17}$\lbl{A1}}
\AXC{$\stackrel{\Pi_{15}}{\vdots}$}\noLine
\UIC{$\varphi_{61}$}
\BIC{$\mathrm{Conv}(a,\alpha) \to \varepsilon \in \mathbb{R} \to 0 < \varepsilon \to \varphi_{29}$}
\AXC{$(\mathrm{Conv}(a,\alpha) \to \varepsilon \in \mathbb{R} \to 0 < \varepsilon \to \varphi_{29}) \to \varphi_{54}$\lbl{A1}}
\BIC{$\varphi_{54}$}
\end{ptree}
\noindent$\Pi_{17}$:
\begin{ptree}
\AXC{$\stackrel{\Pi_{11}}{\vdots}$}\noLine
\UIC{$\lvert a \rvert(n) = \lvert a_{n} \rvert \to \lvert a_{n} \rvert = \lvert a \rvert(n)$}
\AXC{$\stackrel{\Pi_{16}}{\vdots}$}\noLine
\UIC{$\varphi_{54}$}
\BIC{$\mathrm{Conv}(a,\alpha) \to \varepsilon \in \mathbb{R} \to 0 < \varepsilon \to \varphi_{30}$}
\AXC{$(\mathrm{Conv}(a,\alpha) \to \varepsilon \in \mathbb{R} \to 0 < \varepsilon \to \varphi_{30}) \to \varphi_{72}$\lbl{A1}}
\BIC{$\varphi_{72}$}
\end{ptree}
\noindent$\Pi_{18}$:
\begin{ptree}
\AXC{$\stackrel{\Pi_{1}}{\vdots}$}\noLine
\UIC{$\varphi_{28}$}
\AXC{$\mathrm{Conv}(a,\alpha) \to \varepsilon \in \mathbb{R} \to 0 < \varepsilon \to \varphi_{16}$\lbl{A1}}
\AXC{$\stackrel{\Pi_{3}}{\vdots}$}\noLine
\UIC{$\varphi_{60}$}
\BIC{$\mathrm{Conv}(a,\alpha) \to \varepsilon \in \mathbb{R} \to 0 < \varepsilon \to \varphi_{27}$}
\AXC{$\stackrel{\Pi_{10}}{\vdots}$}\noLine
\UIC{$\varphi_{63}$}
\BIC{$\mathrm{Conv}(a,\alpha) \to \varepsilon \in \mathbb{R} \to 0 < \varepsilon \to \varphi_{33}$}
\AXC{$\stackrel{\Pi_{17}}{\vdots}$}\noLine
\UIC{$\varphi_{72}$}
\BIC{$\varphi_{28} \to \mathrm{Conv}(a,\alpha) \to \varepsilon \in \mathbb{R} \to 0 < \varepsilon \to \varphi_{35}$}
\BIC{$\mathrm{Conv}(a,\alpha) \to \varepsilon \in \mathbb{R} \to 0 < \varepsilon \to \varphi_{35}$}
\AXC{$\varphi_{64}$\lbl{A1}}
\BIC{$\varphi_{25} \to n \in \mathbb{N} \to N \leq n \to \lvert \lvert a \rvert(n) - \lvert \alpha \rvert \rvert < \varepsilon$}
\UIC{$\forall n\,(\varphi_{25} \to n \in \mathbb{N} \to N \leq n \to \lvert \lvert a \rvert(n) - \lvert \alpha \rvert \rvert < \varepsilon)$}
\AXC{$\varphi_{78}$\lbl{A1}}
\BIC{$\varphi_{76}$}
\end{ptree}
\noindent$\Pi_{19}$:
\begin{ptree}
\AXC{$\mathrm{Conv}(a,\alpha) \to \varepsilon \in \mathbb{R} \to 0 < \varepsilon \to \varphi_{38}$\lbl{A1}}
\AXC{$\varphi_{62}$\lbl{A1}}
\BIC{$\mathrm{Conv}(a,\alpha) \to \varepsilon \in \mathbb{R} \to \varphi_{13}$}
\AXC{$(\mathrm{Conv}(a,\alpha) \to \varepsilon \in \mathbb{R} \to \varphi_{13}) \to \varphi_{68}$\lbl{A1}}
\BIC{$\varphi_{68}$}
\end{ptree}
\noindent$\Pi_{20}$:
\begin{ptree}
\AXC{$\forall N\,(\varphi_{11}) \to \varphi_{11}$\lbl{A2}}
\AXC{$\varphi_{25} \to \forall n\,(\varphi_{25})$\lbl{A4}}
\AXC{$\varphi_{65}$\lbl{A3}}
\AXC{$\stackrel{\Pi_{18}}{\vdots}$}\noLine
\UIC{$\varphi_{76}$}
\BIC{$(\varphi_{25} \to \forall n\,(\varphi_{25})) \to \mathrm{Conv}(a,\alpha) \to \varepsilon \in \mathbb{R} \to 0 < \varepsilon \to \varphi_{37}$}
\BIC{$\mathrm{Conv}(a,\alpha) \to \varepsilon \in \mathbb{R} \to 0 < \varepsilon \to \varphi_{37}$}
\AXC{$\stackrel{\Pi_{19}}{\vdots}$}\noLine
\UIC{$\varphi_{68}$}
\BIC{$\mathrm{Conv}(a,\alpha) \to \varepsilon \in \mathbb{R} \to 0 < \varepsilon \to \varphi_{43}$}
\AXC{$\varphi_{74}$\lbl{A1}}
\BIC{$(\forall N\,(\varphi_{11}) \to \varphi_{11}) \to \mathrm{Conv}(a,\alpha) \to \varepsilon \in \mathbb{R} \to 0 < \varepsilon \to \varphi_{44}$}
\BIC{$\mathrm{Conv}(a,\alpha) \to \varepsilon \in \mathbb{R} \to 0 < \varepsilon \to \varphi_{44}$}
\end{ptree}
\noindent$\Pi_{21}$:
\begin{ptree}
\AXC{$\forall a\,\forall \alpha\,\forall \varepsilon\,(\mathrm{Conv}(a,\alpha) \to \varphi_{22})$\lbl{[C3]}}
\AXC{$\varphi_{57}$\lbl{A2}}
\BIC{$\forall \alpha\,\forall \varepsilon\,(\mathrm{Conv}(a,\alpha) \to \varphi_{22})$}
\AXC{$\varphi_{56}$\lbl{A2}}
\BIC{$\forall \varepsilon\,(\mathrm{Conv}(a,\alpha) \to \varphi_{22})$}
\AXC{$\forall \varepsilon\,(\mathrm{Conv}(a,\alpha) \to \varphi_{22}) \to \mathrm{Conv}(a,\alpha) \to \varphi_{22}$\lbl{A2}}
\BIC{$\mathrm{Conv}(a,\alpha) \to \varphi_{22}$}
\end{ptree}
\noindent$\Pi_{22}$:
\begin{ptree}
\AXC{$\stackrel{\Pi_{20}}{\vdots}$}\noLine
\UIC{$\mathrm{Conv}(a,\alpha) \to \varepsilon \in \mathbb{R} \to 0 < \varepsilon \to \varphi_{44}$}
\AXC{$\varphi_{69}$\lbl{A1}}
\BIC{$\varphi_{31} \to \neg (N \in \mathbb{N} \land \forall n \in \mathbb{N}\,(N \leq n \to \lvert a_{n} - \alpha \rvert < \varepsilon))$}
\UIC{$\varphi_{52}$}
\AXC{$\stackrel{\Pi_{21}}{\vdots}$}\noLine
\UIC{$\mathrm{Conv}(a,\alpha) \to \varphi_{22}$}
\AXC{$(\mathrm{Conv}(a,\alpha) \to \varphi_{22}) \to \varphi_{52} \to \varphi_{77}$\lbl{A1}}
\BIC{$\varphi_{52} \to \varphi_{77}$}
\BIC{$\varphi_{77}$}
\end{ptree}
\noindent$\Pi_{23}$:
\begin{ptree}
\AXC{$\varphi_{31} \to \forall N\,(\varphi_{31})$\lbl{A4}}
\AXC{$\varphi_{52} \to \forall N\,(\varphi_{31}) \to \varphi_{7}$\lbl{A3}}
\AXC{$\stackrel{\Pi_{22}}{\vdots}$}\noLine
\UIC{$\varphi_{77}$}
\BIC{$(\varphi_{31} \to \forall N\,(\varphi_{31})) \to \mathrm{Conv}(a,\alpha) \to \varepsilon \in \mathbb{R} \to \varphi_{19}$}
\BIC{$\mathrm{Conv}(a,\alpha) \to \varepsilon \in \mathbb{R} \to \varphi_{19}$}
\UIC{$\forall \varepsilon\,(\mathrm{Conv}(a,\alpha) \to \varepsilon \in \mathbb{R} \to \varphi_{19})$}
\AXC{$\forall \varepsilon\,(\mathrm{Conv}(a,\alpha) \to \varepsilon \in \mathbb{R} \to \varphi_{19}) \to \varphi_{58} \to \varphi_{46}$\lbl{A1}}
\BIC{$\varphi_{58} \to \varphi_{46}$}
\end{ptree}
\noindent$\Pi_{24}$:
\begin{ptree}
\AXC{$\forall x \in \mathbb{R}\,(\lvert x \rvert \in \mathbb{R})$\lbl{[A5]}}
\AXC{$\forall x \in \mathbb{R}\,(\lvert x \rvert \in \mathbb{R}) \to \alpha \in \mathbb{R} \to \lvert \alpha \rvert \in \mathbb{R}$\lbl{A2}}
\BIC{$\alpha \in \mathbb{R} \to \lvert \alpha \rvert \in \mathbb{R}$}
\AXC{$(\alpha \in \mathbb{R} \to \lvert \alpha \rvert \in \mathbb{R}) \to \varphi_{5}$\lbl{A1}}
\BIC{$\varphi_{5}$}
\end{ptree}
\noindent$\Pi_{25}$:
\begin{ptree}
\AXC{$\stackrel{\Pi_{13}}{\vdots}$}\noLine
\UIC{$\varphi_{10}$}
\AXC{$\varphi_{10} \to \mathrm{Conv}(a,\alpha) \to \mathrm{Seq}(\lvert a \rvert)$\lbl{A1}}
\BIC{$\mathrm{Conv}(a,\alpha) \to \mathrm{Seq}(\lvert a \rvert)$}
\AXC{$\varphi_{45}$\lbl{DefConv}}
\AXC{$\varphi_{45} \to \varphi_{48}$\lbl{A2}}
\BIC{$\varphi_{48}$}
\AXC{$\varphi_{48} \to \varphi_{50}$\lbl{A2}}
\BIC{$\varphi_{50}$}
\AXC{$\varphi_{50} \to \varphi_{55}$\lbl{A1}}
\BIC{$\varphi_{55}$}
\BIC{$(\mathrm{Conv}(a,\alpha) \to \lvert \alpha \rvert \in \mathbb{R}) \to \varphi_{49}$}
\end{ptree}
\noindent$\Pi_{26}$:
\begin{ptree}
\AXC{$\forall a\,\forall \alpha\,(\mathrm{Conv}(a,\alpha) \to \alpha \in \mathbb{R})$\lbl{[C2]}}
\AXC{$\varphi_{6}$\lbl{A2}}
\BIC{$\forall \alpha\,(\mathrm{Conv}(a,\alpha) \to \alpha \in \mathbb{R})$}
\AXC{$\varphi_{3}$\lbl{A2}}
\BIC{$\mathrm{Conv}(a,\alpha) \to \alpha \in \mathbb{R}$}
\AXC{$\stackrel{\Pi_{24}}{\vdots}$}\noLine
\UIC{$\varphi_{5}$}
\BIC{$\mathrm{Conv}(a,\alpha) \to \lvert \alpha \rvert \in \mathbb{R}$}
\AXC{$\stackrel{\Pi_{25}}{\vdots}$}\noLine
\UIC{$(\mathrm{Conv}(a,\alpha) \to \lvert \alpha \rvert \in \mathbb{R}) \to \varphi_{49}$}
\BIC{$\varphi_{49}$}
\end{ptree}
\noindent 主木:
\begin{ptree}
\AXC{$\mathrm{Conv}(a,\alpha) \to \forall \varepsilon\,(\mathrm{Conv}(a,\alpha))$\lbl{A4}}
\AXC{$\varphi_{58}$\lbl{A3}}
\AXC{$\stackrel{\Pi_{23}}{\vdots}$}\noLine
\UIC{$\varphi_{58} \to \varphi_{46}$}
\BIC{$\varphi_{46}$}
\BIC{$\mathrm{Conv}(a,\alpha) \to \forall \varepsilon \in \mathbb{R}\,(\varphi_{19})$}
\AXC{$\stackrel{\Pi_{26}}{\vdots}$}\noLine
\UIC{$\varphi_{49}$}
\BIC{$\mathrm{Conv}(a,\alpha) \to \mathrm{Conv}(\lvert a \rvert,\lvert \alpha \rvert)$}
\UIC{$\forall \alpha\,(\mathrm{Conv}(a,\alpha) \to \mathrm{Conv}(\lvert a \rvert,\lvert \alpha \rvert))$}
\UIC{$\forall a\,\forall \alpha\,(\mathrm{Conv}(a,\alpha) \to \mathrm{Conv}(\lvert a \rvert,\lvert \alpha \rvert))$}
\end{ptree}
\begin{abbrevlist}
\item $\varphi_{1} :\equiv \forall \alpha\,(\mathrm{Conv}(a,\alpha) \to \mathrm{Seq}(a)) \to \mathrm{Conv}(a,\alpha) \to \mathrm{Seq}(a)$
\item $\varphi_{2} :\equiv \forall a\,\forall \alpha\,(\mathrm{Conv}(a,\alpha) \to \mathrm{Seq}(a)) \to \forall \alpha\,(\mathrm{Conv}(a,\alpha) \to \mathrm{Seq}(a))$
\item $\varphi_{3} :\equiv \forall \alpha\,(\mathrm{Conv}(a,\alpha) \to \alpha \in \mathbb{R}) \to \mathrm{Conv}(a,\alpha) \to \alpha \in \mathbb{R}$
\item $\varphi_{4} :\equiv \forall x\,\forall y\,(x = y \to y = x) \to \forall y\,(\lvert a \rvert(n) = y \to y = \lvert a \rvert(n))$
\item $\varphi_{5} :\equiv (\mathrm{Conv}(a,\alpha) \to \alpha \in \mathbb{R}) \to \mathrm{Conv}(a,\alpha) \to \lvert \alpha \rvert \in \mathbb{R}$
\item $\varphi_{6} :\equiv \forall a\,\forall \alpha\,(\mathrm{Conv}(a,\alpha) \to \alpha \in \mathbb{R}) \to \forall \alpha\,(\mathrm{Conv}(a,\alpha) \to \alpha \in \mathbb{R})$
\item $\varphi_{7} :\equiv \forall N\,(\neg (N \in \mathbb{N} \land \forall n \in \mathbb{N}\,(N \leq n \to \lvert a_{n} - \alpha \rvert < \varepsilon)))$
\item $\varphi_{8} :\equiv \forall x\,\forall y\,(x \in \mathbb{R} \to y \in \mathbb{R} \to \lvert \lvert x \rvert - \lvert y \rvert \rvert \leq \lvert x - y \rvert)$
\item $\varphi_{9} :\equiv \mathrm{Seq}(a) \to \mathrm{Seq}(\lvert a \rvert) \land \forall n \in \mathbb{N}\,(\lvert a \rvert(n) = \lvert a_{n} \rvert)$
\item $\varphi_{10} :\equiv \mathrm{Conv}(a,\alpha) \to \mathrm{Seq}(\lvert a \rvert) \land \forall n \in \mathbb{N}\,(\lvert a \rvert(n) = \lvert a_{n} \rvert)$
\item $\varphi_{11} :\equiv \neg (N \in \mathbb{N} \land \forall n \in \mathbb{N}\,(N \leq n \to \lvert \lvert a \rvert(n) - \lvert \alpha \rvert \rvert < \varepsilon))$
\item $\varphi_{12} :\equiv \forall n\,(\mathrm{Seq}(a) \to n \in \mathbb{N} \to a_{n} \in \mathbb{R}) \to \mathrm{Seq}(a) \to n \in \mathbb{N} \to a_{n} \in \mathbb{R}$
\item $\varphi_{13} :\equiv 0 < \varepsilon \to N \in \mathbb{N} \land \forall n \in \mathbb{N}\,(N \leq n \to \lvert a_{n} - \alpha \rvert < \varepsilon) \to N \in \mathbb{N}$
\item $\varphi_{14} :\equiv \forall y\,(a_{n} \in \mathbb{R} \to y \in \mathbb{R} \to \lvert \lvert a_{n} \rvert - \lvert y \rvert \rvert \leq \lvert a_{n} - y \rvert)$
\item $\varphi_{15} :\equiv \forall y\,\forall z\,(\lvert \lvert a_{n} \rvert - \lvert \alpha \rvert \rvert \leq y \to y < z \to \lvert \lvert a_{n} \rvert - \lvert \alpha \rvert \rvert < z)$
\item $\varphi_{16} :\equiv N \in \mathbb{N} \land \forall n \in \mathbb{N}\,(N \leq n \to \lvert a_{n} - \alpha \rvert < \varepsilon) \to n \in \mathbb{N} \to N \leq n \to N \leq n$
\item $\varphi_{17} :\equiv N \in \mathbb{N} \land \forall n \in \mathbb{N}\,(N \leq n \to \lvert a_{n} - \alpha \rvert < \varepsilon) \to n \in \mathbb{N} \to n \in \mathbb{N}$
\item $\varphi_{18} :\equiv \forall w\,(\lvert a_{n} \rvert = w \to \lvert \lvert a_{n} \rvert - \lvert \alpha \rvert \rvert < \varepsilon \to \lvert w - \lvert \alpha \rvert \rvert < \varepsilon)$
\item $\varphi_{19} :\equiv 0 < \varepsilon \to \exists N \in \mathbb{N}\,\forall n \in \mathbb{N}\,(N \leq n \to \lvert \lvert a \rvert(n) - \lvert \alpha \rvert \rvert < \varepsilon)$
\item $\varphi_{20} :\equiv \forall n \in \mathbb{N}\,(\lvert a \rvert(n) = \lvert a_{n} \rvert) \to n \in \mathbb{N} \to \lvert a \rvert(n) = \lvert a_{n} \rvert$
\item $\varphi_{21} :\equiv \forall a\,\forall n\,(\mathrm{Seq}(a) \to n \in \mathbb{N} \to a_{n} \in \mathbb{R}) \to \forall n\,(\mathrm{Seq}(a) \to n \in \mathbb{N} \to a_{n} \in \mathbb{R})$
\item $\varphi_{22} :\equiv \varepsilon \in \mathbb{R} \to 0 < \varepsilon \to \exists N \in \mathbb{N}\,\forall n \in \mathbb{N}\,(N \leq n \to \lvert a_{n} - \alpha \rvert < \varepsilon)$
\item $\varphi_{23} :\equiv \forall n \in \mathbb{N}\,(N \leq n \to \lvert a_{n} - \alpha \rvert < \varepsilon) \to n \in \mathbb{N} \to N \leq n \to \lvert a_{n} - \alpha \rvert < \varepsilon$
\item $\varphi_{24} :\equiv N \in \mathbb{N} \land \forall n \in \mathbb{N}\,(N \leq n \to \lvert a_{n} - \alpha \rvert < \varepsilon) \to n \in \mathbb{N} \to N \leq n \to a_{n} \in \mathbb{R}$
\item $\varphi_{25} :\equiv \mathrm{Conv}(a,\alpha) \land \varepsilon \in \mathbb{R} \land 0 < \varepsilon \land (N \in \mathbb{N} \land \forall n \in \mathbb{N}\,(N \leq n \to \lvert a_{n} - \alpha \rvert < \varepsilon))$
\item $\varphi_{26} :\equiv \forall y\,(\lvert a \rvert(n) = y \to y = \lvert a \rvert(n)) \to \lvert a \rvert(n) = \lvert a_{n} \rvert \to \lvert a_{n} \rvert = \lvert a \rvert(n)$
\item $\varphi_{27} :\equiv N \in \mathbb{N} \land \forall n \in \mathbb{N}\,(N \leq n \to \lvert a_{n} - \alpha \rvert < \varepsilon) \to n \in \mathbb{N} \to N \leq n \to \lvert a_{n} - \alpha \rvert < \varepsilon$
\item $\varphi_{28} :\equiv \lvert a_{n} \rvert = \lvert a \rvert(n) \to \lvert \lvert a_{n} \rvert - \lvert \alpha \rvert \rvert < \varepsilon \to \lvert \lvert a \rvert(n) - \lvert \alpha \rvert \rvert < \varepsilon$
\item $\varphi_{29} :\equiv N \in \mathbb{N} \land \forall n \in \mathbb{N}\,(N \leq n \to \lvert a_{n} - \alpha \rvert < \varepsilon) \to n \in \mathbb{N} \to N \leq n \to \lvert a \rvert(n) = \lvert a_{n} \rvert$
\item $\varphi_{30} :\equiv N \in \mathbb{N} \land \forall n \in \mathbb{N}\,(N \leq n \to \lvert a_{n} - \alpha \rvert < \varepsilon) \to n \in \mathbb{N} \to N \leq n \to \lvert a_{n} \rvert = \lvert a \rvert(n)$
\item $\varphi_{31} :\equiv \mathrm{Conv}(a,\alpha) \land \varepsilon \in \mathbb{R} \land 0 < \varepsilon \land \neg \exists N \in \mathbb{N}\,\forall n \in \mathbb{N}\,(N \leq n \to \lvert \lvert a \rvert(n) - \lvert \alpha \rvert \rvert < \varepsilon)$
\item $\varphi_{32} :\equiv \forall z\,(\lvert \lvert a_{n} \rvert - \lvert \alpha \rvert \rvert \leq \lvert a_{n} - \alpha \rvert \to \lvert a_{n} - \alpha \rvert < z \to \lvert \lvert a_{n} \rvert - \lvert \alpha \rvert \rvert < z)$
\item $\varphi_{33} :\equiv N \in \mathbb{N} \land \forall n \in \mathbb{N}\,(N \leq n \to \lvert a_{n} - \alpha \rvert < \varepsilon) \to n \in \mathbb{N} \to N \leq n \to \lvert \lvert a_{n} \rvert - \lvert \alpha \rvert \rvert < \varepsilon$
\item $\varphi_{34} :\equiv \lvert \lvert a_{n} \rvert - \lvert \alpha \rvert \rvert \leq \lvert a_{n} - \alpha \rvert \to \lvert a_{n} - \alpha \rvert < \varepsilon \to \lvert \lvert a_{n} \rvert - \lvert \alpha \rvert \rvert < \varepsilon$
\item $\varphi_{35} :\equiv N \in \mathbb{N} \land \forall n \in \mathbb{N}\,(N \leq n \to \lvert a_{n} - \alpha \rvert < \varepsilon) \to n \in \mathbb{N} \to N \leq n \to \lvert \lvert a \rvert(n) - \lvert \alpha \rvert \rvert < \varepsilon$
\item $\varphi_{36} :\equiv N \in \mathbb{N} \land \forall n \in \mathbb{N}\,(N \leq n \to \lvert a_{n} - \alpha \rvert < \varepsilon) \to n \in \mathbb{N} \to N \leq n \to n \in \mathbb{N} \to \lvert a \rvert(n) = \lvert a_{n} \rvert$
\item $\varphi_{37} :\equiv N \in \mathbb{N} \land \forall n \in \mathbb{N}\,(N \leq n \to \lvert a_{n} - \alpha \rvert < \varepsilon) \to \forall n \in \mathbb{N}\,(N \leq n \to \lvert \lvert a \rvert(n) - \lvert \alpha \rvert \rvert < \varepsilon)$
\item $\varphi_{38} :\equiv N \in \mathbb{N} \land \forall n \in \mathbb{N}\,(N \leq n \to \lvert a_{n} - \alpha \rvert < \varepsilon) \to N \in \mathbb{N} \land \forall n \in \mathbb{N}\,(N \leq n \to \lvert a_{n} - \alpha \rvert < \varepsilon)$
\item $\varphi_{39} :\equiv N \in \mathbb{N} \land \forall n \in \mathbb{N}\,(N \leq n \to \lvert a_{n} - \alpha \rvert < \varepsilon) \to n \in \mathbb{N} \to N \leq n \to \forall n \in \mathbb{N}\,(\lvert a \rvert(n) = \lvert a_{n} \rvert)$
\item $\varphi_{40} :\equiv N \in \mathbb{N} \land \forall n \in \mathbb{N}\,(N \leq n \to \lvert a_{n} - \alpha \rvert < \varepsilon) \to n \in \mathbb{N} \to N \leq n \to n \in \mathbb{N} \to N \leq n \to \lvert a_{n} - \alpha \rvert < \varepsilon$
\item $\varphi_{41} :\equiv N \in \mathbb{N} \land \forall n \in \mathbb{N}\,(N \leq n \to \lvert a_{n} - \alpha \rvert < \varepsilon) \to n \in \mathbb{N} \to N \leq n \to \forall n \in \mathbb{N}\,(N \leq n \to \lvert a_{n} - \alpha \rvert < \varepsilon)$
\item $\varphi_{42} :\equiv N \in \mathbb{N} \land \forall n \in \mathbb{N}\,(N \leq n \to \lvert a_{n} - \alpha \rvert < \varepsilon) \to n \in \mathbb{N} \to N \leq n \to \lvert \lvert a_{n} \rvert - \lvert \alpha \rvert \rvert \leq \lvert a_{n} - \alpha \rvert$
\item $\varphi_{43} :\equiv N \in \mathbb{N} \land \forall n \in \mathbb{N}\,(N \leq n \to \lvert a_{n} - \alpha \rvert < \varepsilon) \to N \in \mathbb{N} \land \forall n \in \mathbb{N}\,(N \leq n \to \lvert \lvert a \rvert(n) - \lvert \alpha \rvert \rvert < \varepsilon)$
\item $\varphi_{44} :\equiv N \in \mathbb{N} \land \forall n \in \mathbb{N}\,(N \leq n \to \lvert a_{n} - \alpha \rvert < \varepsilon) \to \exists N \in \mathbb{N}\,\forall n \in \mathbb{N}\,(N \leq n \to \lvert \lvert a \rvert(n) - \lvert \alpha \rvert \rvert < \varepsilon)$
\item $\varphi_{45} :\equiv \forall a\,\forall \alpha\,(\mathrm{Conv}(a,\alpha) \leftrightarrow \mathrm{Seq}(a) \land \alpha \in \mathbb{R} \land \forall \varepsilon \in \mathbb{R}\,(0 < \varepsilon \to \exists N \in \mathbb{N}\,\forall n \in \mathbb{N}\,(N \leq n \to \lvert a_{n} - \alpha \rvert < \varepsilon)))$
\item $\varphi_{46} :\equiv (\mathrm{Conv}(a,\alpha) \to \forall \varepsilon\,(\mathrm{Conv}(a,\alpha))) \to \mathrm{Conv}(a,\alpha) \to \forall \varepsilon \in \mathbb{R}\,(\varphi_{19})$
\item $\varphi_{47} :\equiv \forall z\,\forall w\,(z = w \to \lvert z - \lvert \alpha \rvert \rvert < \varepsilon \to \lvert w - \lvert \alpha \rvert \rvert < \varepsilon) \to \varphi_{18}$
\item $\varphi_{48} :\equiv \forall \alpha\,(\mathrm{Conv}(\lvert a \rvert,\alpha) \leftrightarrow \mathrm{Seq}(\lvert a \rvert) \land \alpha \in \mathbb{R} \land \forall \varepsilon \in \mathbb{R}\,(0 < \varepsilon \to \exists N \in \mathbb{N}\,\forall n \in \mathbb{N}\,(N \leq n \to \lvert \lvert a \rvert(n) - \alpha \rvert < \varepsilon)))$
\item $\varphi_{49} :\equiv (\mathrm{Conv}(a,\alpha) \to \forall \varepsilon \in \mathbb{R}\,(\varphi_{19})) \to \mathrm{Conv}(a,\alpha) \to \mathrm{Conv}(\lvert a \rvert,\lvert \alpha \rvert)$
\item $\varphi_{50} :\equiv \mathrm{Conv}(\lvert a \rvert,\lvert \alpha \rvert) \leftrightarrow \mathrm{Seq}(\lvert a \rvert) \land \lvert \alpha \rvert \in \mathbb{R} \land \forall \varepsilon \in \mathbb{R}\,(\varphi_{19})$
\item $\varphi_{51} :\equiv \forall n\,(\varphi_{25}) \to \forall n \in \mathbb{N}\,(N \leq n \to \lvert \lvert a \rvert(n) - \lvert \alpha \rvert \rvert < \varepsilon)$
\item $\varphi_{52} :\equiv \forall N\,(\varphi_{31} \to \neg (N \in \mathbb{N} \land \forall n \in \mathbb{N}\,(N \leq n \to \lvert a_{n} - \alpha \rvert < \varepsilon)))$
\item $\varphi_{53} :\equiv (\mathrm{Conv}(a,\alpha) \to \alpha \in \mathbb{R}) \to \mathrm{Conv}(a,\alpha) \to \varepsilon \in \mathbb{R} \to 0 < \varepsilon \to \varphi_{42}$
\item $\varphi_{54} :\equiv (\lvert a \rvert(n) = \lvert a_{n} \rvert \to \lvert a_{n} \rvert = \lvert a \rvert(n)) \to \mathrm{Conv}(a,\alpha) \to \varepsilon \in \mathbb{R} \to 0 < \varepsilon \to \varphi_{30}$
\item $\varphi_{55} :\equiv (\mathrm{Conv}(a,\alpha) \to \mathrm{Seq}(\lvert a \rvert)) \to (\mathrm{Conv}(a,\alpha) \to \lvert \alpha \rvert \in \mathbb{R}) \to \varphi_{49}$
\item $\varphi_{56} :\equiv \forall \alpha\,\forall \varepsilon\,(\mathrm{Conv}(a,\alpha) \to \varphi_{22}) \to \forall \varepsilon\,(\mathrm{Conv}(a,\alpha) \to \varphi_{22})$
\item $\varphi_{57} :\equiv \forall a\,\forall \alpha\,\forall \varepsilon\,(\mathrm{Conv}(a,\alpha) \to \varphi_{22}) \to \forall \alpha\,\forall \varepsilon\,(\mathrm{Conv}(a,\alpha) \to \varphi_{22})$
\item $\varphi_{58} :\equiv \forall \varepsilon\,(\mathrm{Conv}(a,\alpha) \to \varepsilon \in \mathbb{R} \to \varphi_{19}) \to \forall \varepsilon\,(\mathrm{Conv}(a,\alpha)) \to \forall \varepsilon \in \mathbb{R}\,(\varphi_{19})$
\item $\varphi_{59} :\equiv (\mathrm{Conv}(a,\alpha) \to \varepsilon \in \mathbb{R} \to 0 < \varepsilon \to \varphi_{17}) \to \mathrm{Conv}(a,\alpha) \to \varepsilon \in \mathbb{R} \to 0 < \varepsilon \to \varphi_{24}$
\item $\varphi_{60} :\equiv (\mathrm{Conv}(a,\alpha) \to \varepsilon \in \mathbb{R} \to 0 < \varepsilon \to \varphi_{16}) \to \mathrm{Conv}(a,\alpha) \to \varepsilon \in \mathbb{R} \to 0 < \varepsilon \to \varphi_{27}$
\item $\varphi_{61} :\equiv (\mathrm{Conv}(a,\alpha) \to \varepsilon \in \mathbb{R} \to 0 < \varepsilon \to \varphi_{17}) \to \mathrm{Conv}(a,\alpha) \to \varepsilon \in \mathbb{R} \to 0 < \varepsilon \to \varphi_{29}$
\item $\varphi_{62} :\equiv (\mathrm{Conv}(a,\alpha) \to \varepsilon \in \mathbb{R} \to 0 < \varepsilon \to \varphi_{38}) \to \mathrm{Conv}(a,\alpha) \to \varepsilon \in \mathbb{R} \to \varphi_{13}$
\item $\varphi_{63} :\equiv (\mathrm{Conv}(a,\alpha) \to \varepsilon \in \mathbb{R} \to 0 < \varepsilon \to \varphi_{27}) \to \mathrm{Conv}(a,\alpha) \to \varepsilon \in \mathbb{R} \to 0 < \varepsilon \to \varphi_{33}$
\item $\varphi_{64} :\equiv (\mathrm{Conv}(a,\alpha) \to \varepsilon \in \mathbb{R} \to 0 < \varepsilon \to \varphi_{35}) \to \varphi_{25} \to n \in \mathbb{N} \to N \leq n \to \lvert \lvert a \rvert(n) - \lvert \alpha \rvert \rvert < \varepsilon$
\item $\varphi_{65} :\equiv \forall n\,(\varphi_{25} \to n \in \mathbb{N} \to N \leq n \to \lvert \lvert a \rvert(n) - \lvert \alpha \rvert \rvert < \varepsilon) \to \varphi_{51}$
\item $\varphi_{66} :\equiv (\mathrm{Seq}(a) \to n \in \mathbb{N} \to a_{n} \in \mathbb{R}) \to (\mathrm{Conv}(a,\alpha) \to \mathrm{Seq}(a)) \to \varphi_{59}$
\item $\varphi_{67} :\equiv (\mathrm{Conv}(a,\alpha) \to \varepsilon \in \mathbb{R} \to 0 < \varepsilon \to \varphi_{38}) \to \mathrm{Conv}(a,\alpha) \to \varepsilon \in \mathbb{R} \to 0 < \varepsilon \to \varphi_{41}$
\item $\varphi_{68} :\equiv (\mathrm{Conv}(a,\alpha) \to \varepsilon \in \mathbb{R} \to 0 < \varepsilon \to \varphi_{37}) \to \mathrm{Conv}(a,\alpha) \to \varepsilon \in \mathbb{R} \to 0 < \varepsilon \to \varphi_{43}$
\item $\varphi_{69} :\equiv (\mathrm{Conv}(a,\alpha) \to \varepsilon \in \mathbb{R} \to 0 < \varepsilon \to \varphi_{44}) \to \varphi_{31} \to \neg (N \in \mathbb{N} \land \forall n \in \mathbb{N}\,(N \leq n \to \lvert a_{n} - \alpha \rvert < \varepsilon))$
\item $\varphi_{70} :\equiv (\mathrm{Conv}(a,\alpha) \to \varepsilon \in \mathbb{R} \to 0 < \varepsilon \to \varphi_{39}) \to \varphi_{20} \to \mathrm{Conv}(a,\alpha) \to \varepsilon \in \mathbb{R} \to 0 < \varepsilon \to \varphi_{36}$
\item $\varphi_{71} :\equiv (a_{n} \in \mathbb{R} \to \alpha \in \mathbb{R} \to \lvert \lvert a_{n} \rvert - \lvert \alpha \rvert \rvert \leq \lvert a_{n} - \alpha \rvert) \to (\mathrm{Conv}(a,\alpha) \to \varepsilon \in \mathbb{R} \to 0 < \varepsilon \to \varphi_{24}) \to \varphi_{53}$
\item $\varphi_{72} :\equiv (\mathrm{Conv}(a,\alpha) \to \varepsilon \in \mathbb{R} \to 0 < \varepsilon \to \varphi_{33}) \to \varphi_{28} \to \mathrm{Conv}(a,\alpha) \to \varepsilon \in \mathbb{R} \to 0 < \varepsilon \to \varphi_{35}$
\item $\varphi_{73} :\equiv (\mathrm{Conv}(a,\alpha) \to \varepsilon \in \mathbb{R} \to 0 < \varepsilon \to \varphi_{41}) \to \varphi_{23} \to \mathrm{Conv}(a,\alpha) \to \varepsilon \in \mathbb{R} \to 0 < \varepsilon \to \varphi_{40}$
\item $\varphi_{74} :\equiv (\mathrm{Conv}(a,\alpha) \to \varepsilon \in \mathbb{R} \to 0 < \varepsilon \to \varphi_{43}) \to (\forall N\,(\varphi_{11}) \to \varphi_{11}) \to \mathrm{Conv}(a,\alpha) \to \varepsilon \in \mathbb{R} \to 0 < \varepsilon \to \varphi_{44}$
\item $\varphi_{75} :\equiv (\mathrm{Conv}(a,\alpha) \to \varepsilon \in \mathbb{R} \to 0 < \varepsilon \to \varphi_{40}) \to (\mathrm{Conv}(a,\alpha) \to \varepsilon \in \mathbb{R} \to 0 < \varepsilon \to \varphi_{17}) \to \varphi_{60}$
\item $\varphi_{76} :\equiv \varphi_{65} \to (\varphi_{25} \to \forall n\,(\varphi_{25})) \to \mathrm{Conv}(a,\alpha) \to \varepsilon \in \mathbb{R} \to 0 < \varepsilon \to \varphi_{37}$
\item $\varphi_{77} :\equiv (\varphi_{52} \to \forall N\,(\varphi_{31}) \to \varphi_{7}) \to (\varphi_{31} \to \forall N\,(\varphi_{31})) \to \mathrm{Conv}(a,\alpha) \to \varepsilon \in \mathbb{R} \to \varphi_{19}$
\item $\varphi_{78} :\equiv \forall n\,(\varphi_{25} \to n \in \mathbb{N} \to N \leq n \to \lvert \lvert a \rvert(n) - \lvert \alpha \rvert \rvert < \varepsilon) \to \varphi_{76}$
\end{abbrevlist}
\end{fullproof}

% 自動生成 (prover/build.js): LimAbsZero  ノード数 278
\begin{fullproof}{定理2.15(2)}{$\forall a\,(\mathrm{Seq}(a) \to (\mathrm{Conv}(a,0) \leftrightarrow \mathrm{Conv}(\lvert a \rvert,0)))$}{278}\phantomsection\label{tree:LimAbsZero}
\noindent$\Pi_{1}$:
\begin{ptree}
\AXC{$z = w \to z < \varepsilon \to w < \varepsilon$\lbl{A6}}
\UIC{$\forall w\,(z = w \to z < \varepsilon \to w < \varepsilon)$}
\UIC{$\forall z\,\forall w\,(z = w \to z < \varepsilon \to w < \varepsilon)$}
\AXC{$\forall z\,\forall w\,(z = w \to z < \varepsilon \to w < \varepsilon) \to \varphi_{30}$\lbl{A2}}
\BIC{$\varphi_{30}$}
\AXC{$\varphi_{30} \to \varphi_{58}$\lbl{A2}}
\BIC{$\varphi_{58}$}
\end{ptree}
\noindent$\Pi_{2}$:
\begin{ptree}
\AXC{$\mathrm{Seq}(a) \to \mathrm{Conv}(\lvert a \rvert,0) \to \varepsilon \in \mathbb{R} \to 0 < \varepsilon \to \varphi_{82}$\lbl{A1}}
\AXC{$\varphi_{110}$\lbl{A1}}
\BIC{$\mathrm{Seq}(a) \to \mathrm{Conv}(\lvert a \rvert,0) \to \varepsilon \in \mathbb{R} \to \varphi_{86}$}
\AXC{$(\mathrm{Seq}(a) \to \mathrm{Conv}(\lvert a \rvert,0) \to \varepsilon \in \mathbb{R} \to \varphi_{86}) \to \varphi_{100}$\lbl{A1}}
\BIC{$\varphi_{100}$}
\end{ptree}
\noindent$\Pi_{3}$:
\begin{ptree}
\AXC{$\mathrm{Seq}(a) \to \mathrm{Conv}(\lvert a \rvert,0) \to \varepsilon \in \mathbb{R} \to 0 < \varepsilon \to \varphi_{49}$\lbl{A1}}
\AXC{$\varphi_{57}$\lbl{A2}}
\AXC{$\stackrel{\Pi_{2}}{\vdots}$}\noLine
\UIC{$\varphi_{100}$}
\BIC{$\mathrm{Seq}(a) \to \mathrm{Conv}(\lvert a \rvert,0) \to \varepsilon \in \mathbb{R} \to 0 < \varepsilon \to \varphi_{85}$}
\AXC{$\varphi_{118}$\lbl{A1}}
\BIC{$\varphi_{113}$}
\BIC{$\varphi_{102}$}
\end{ptree}
\noindent$\Pi_{4}$:
\begin{ptree}
\AXC{$y = z \to x = y \to x = z$\lbl{A6}}
\AXC{$(y = z \to x = y \to x = z) \to x = y \to y = z \to x = z$\lbl{A1}}
\BIC{$x = y \to y = z \to x = z$}
\UIC{$\forall z\,(x = y \to y = z \to x = z)$}
\UIC{$\forall y\,\forall z\,(x = y \to y = z \to x = z)$}
\UIC{$\forall x\,\forall y\,\forall z\,(x = y \to y = z \to x = z)$}
\AXC{$\forall x\,\forall y\,\forall z\,(x = y \to y = z \to x = z) \to \varphi_{15}$\lbl{A2}}
\BIC{$\varphi_{15}$}
\end{ptree}
\noindent$\Pi_{5}$:
\begin{ptree}
\AXC{$x = x$\lbl{A5}}
\AXC{$x = y \to x = x \to y = x$\lbl{A6}}
\AXC{$(x = y \to x = x \to y = x) \to x = x \to x = y \to y = x$\lbl{A1}}
\BIC{$x = x \to x = y \to y = x$}
\BIC{$x = y \to y = x$}
\UIC{$\forall y\,(x = y \to y = x)$}
\UIC{$\forall x\,\forall y\,(x = y \to y = x)$}
\AXC{$\varphi_{7}$\lbl{A2}}
\BIC{$\forall y\,(a_{n} - 0 = y \to y = a_{n} - 0)$}
\AXC{$\varphi_{43}$\lbl{A2}}
\BIC{$a_{n} - 0 = a_{n} \to a_{n} = a_{n} - 0$}
\end{ptree}
\noindent$\Pi_{6}$:
\begin{ptree}
\AXC{$\mathrm{Seq}(a) \to \mathrm{Seq}(a)$\lbl{A1}}
\AXC{$\forall a\,\forall n\,(\mathrm{Seq}(a) \to n \in \mathbb{N} \to a_{n} \in \mathbb{R})$\lbl{[S1]}}
\AXC{$\varphi_{47}$\lbl{A2}}
\BIC{$\forall n\,(\mathrm{Seq}(a) \to n \in \mathbb{N} \to a_{n} \in \mathbb{R})$}
\AXC{$\varphi_{39}$\lbl{A2}}
\BIC{$\mathrm{Seq}(a) \to n \in \mathbb{N} \to a_{n} \in \mathbb{R}$}
\AXC{$\varphi_{90}$\lbl{A1}}
\BIC{$(\mathrm{Seq}(a) \to \mathrm{Seq}(a)) \to \varphi_{60}$}
\BIC{$\varphi_{60}$}
\end{ptree}
\noindent$\Pi_{7}$:
\begin{ptree}
\AXC{$\mathrm{Seq}(a) \to \mathrm{Conv}(\lvert a \rvert,0) \to n \in \mathbb{N} \to n \in \mathbb{N}$\lbl{A1}}
\AXC{$\stackrel{\Pi_{6}}{\vdots}$}\noLine
\UIC{$\varphi_{60}$}
\BIC{$\mathrm{Seq}(a) \to \mathrm{Conv}(\lvert a \rvert,0) \to n \in \mathbb{N} \to a_{n} \in \mathbb{R}$}
\end{ptree}
\noindent$\Pi_{8}$:
\begin{ptree}
\AXC{$\forall x \in \mathbb{R}\,(x - 0 = x)$\lbl{[E28]}}
\AXC{$\forall x \in \mathbb{R}\,(x - 0 = x) \to a_{n} \in \mathbb{R} \to a_{n} - 0 = a_{n}$\lbl{A2}}
\BIC{$a_{n} \in \mathbb{R} \to a_{n} - 0 = a_{n}$}
\AXC{$(a_{n} \in \mathbb{R} \to a_{n} - 0 = a_{n}) \to \varphi_{74}$\lbl{A1}}
\BIC{$\varphi_{74}$}
\end{ptree}
\noindent$\Pi_{9}$:
\begin{ptree}
\AXC{$\stackrel{\Pi_{7}}{\vdots}$}\noLine
\UIC{$\mathrm{Seq}(a) \to \mathrm{Conv}(\lvert a \rvert,0) \to n \in \mathbb{N} \to a_{n} \in \mathbb{R}$}
\AXC{$\stackrel{\Pi_{8}}{\vdots}$}\noLine
\UIC{$\varphi_{74}$}
\BIC{$\varphi_{9}$}
\AXC{$\varphi_{9} \to (a_{n} - 0 = a_{n} \to a_{n} = a_{n} - 0) \to \varphi_{10}$\lbl{A1}}
\BIC{$(a_{n} - 0 = a_{n} \to a_{n} = a_{n} - 0) \to \varphi_{10}$}
\end{ptree}
\noindent$\Pi_{10}$:
\begin{ptree}
\AXC{$\stackrel{\Pi_{5}}{\vdots}$}\noLine
\UIC{$a_{n} - 0 = a_{n} \to a_{n} = a_{n} - 0$}
\AXC{$\stackrel{\Pi_{9}}{\vdots}$}\noLine
\UIC{$(a_{n} - 0 = a_{n} \to a_{n} = a_{n} - 0) \to \varphi_{10}$}
\BIC{$\varphi_{10}$}
\AXC{$\varphi_{10} \to \lvert a_{n} \rvert = \lvert a_{n} \rvert \to \varphi_{89}$\lbl{A1}}
\BIC{$\lvert a_{n} \rvert = \lvert a_{n} \rvert \to \varphi_{89}$}
\end{ptree}
\noindent$\Pi_{11}$:
\begin{ptree}
\AXC{$x = x$\lbl{A5}}
\UIC{$\forall x\,(x = x)$}
\AXC{$\forall x\,(x = x) \to \lvert a_{n} \rvert = \lvert a_{n} \rvert$\lbl{A2}}
\BIC{$\lvert a_{n} \rvert = \lvert a_{n} \rvert$}
\AXC{$\stackrel{\Pi_{10}}{\vdots}$}\noLine
\UIC{$\lvert a_{n} \rvert = \lvert a_{n} \rvert \to \varphi_{89}$}
\BIC{$\varphi_{89}$}
\end{ptree}
\noindent$\Pi_{12}$:
\begin{ptree}
\AXC{$y = z \to x = y \to x = z$\lbl{A6}}
\AXC{$(y = z \to x = y \to x = z) \to x = y \to y = z \to x = z$\lbl{A1}}
\BIC{$x = y \to y = z \to x = z$}
\UIC{$\forall z\,(x = y \to y = z \to x = z)$}
\UIC{$\forall y\,\forall z\,(x = y \to y = z \to x = z)$}
\UIC{$\forall x\,\forall y\,\forall z\,(x = y \to y = z \to x = z)$}
\AXC{$\forall x\,\forall y\,\forall z\,(x = y \to y = z \to x = z) \to \varphi_{15}$\lbl{A2}}
\BIC{$\varphi_{15}$}
\end{ptree}
\noindent$\Pi_{13}$:
\begin{ptree}
\AXC{$\forall x \in \mathbb{R}\,(\lvert \lvert x \rvert \rvert = \lvert x \rvert)$\lbl{[E29]}}
\AXC{$\varphi_{52}$\lbl{A2}}
\BIC{$a_{n} \in \mathbb{R} \to \lvert \lvert a_{n} \rvert \rvert = \lvert a_{n} \rvert$}
\AXC{$(a_{n} \in \mathbb{R} \to \lvert \lvert a_{n} \rvert \rvert = \lvert a_{n} \rvert) \to \varphi_{81}$\lbl{A1}}
\BIC{$\varphi_{81}$}
\end{ptree}
\noindent$\Pi_{14}$:
\begin{ptree}
\AXC{$y = z \to x = y \to x = z$\lbl{A6}}
\AXC{$(y = z \to x = y \to x = z) \to x = y \to y = z \to x = z$\lbl{A1}}
\BIC{$x = y \to y = z \to x = z$}
\UIC{$\forall z\,(x = y \to y = z \to x = z)$}
\UIC{$\forall y\,\forall z\,(x = y \to y = z \to x = z)$}
\UIC{$\forall x\,\forall y\,\forall z\,(x = y \to y = z \to x = z)$}
\AXC{$\forall x\,\forall y\,\forall z\,(x = y \to y = z \to x = z) \to \varphi_{15}$\lbl{A2}}
\BIC{$\varphi_{15}$}
\end{ptree}
\noindent$\Pi_{15}$:
\begin{ptree}
\AXC{$z = w \to \lvert \lvert a_{n} \rvert - 0 \rvert = \lvert z \rvert \to \lvert \lvert a_{n} \rvert - 0 \rvert = \lvert w \rvert$\lbl{A6}}
\UIC{$\varphi_{24}$}
\UIC{$\forall z\,(\varphi_{24})$}
\AXC{$\forall z\,(\varphi_{24}) \to \varphi_{59}$\lbl{A2}}
\BIC{$\varphi_{59}$}
\AXC{$\varphi_{59} \to \lvert a_{n} \rvert - 0 = \lvert a_{n} \rvert \to \varphi_{54}$\lbl{A2}}
\BIC{$\lvert a_{n} \rvert - 0 = \lvert a_{n} \rvert \to \varphi_{54}$}
\end{ptree}
\noindent$\Pi_{16}$:
\begin{ptree}
\AXC{$\forall x \in \mathbb{R}\,(\lvert x \rvert \in \mathbb{R})$\lbl{[A5]}}
\AXC{$\forall x \in \mathbb{R}\,(\lvert x \rvert \in \mathbb{R}) \to a_{n} \in \mathbb{R} \to \lvert a_{n} \rvert \in \mathbb{R}$\lbl{A2}}
\BIC{$a_{n} \in \mathbb{R} \to \lvert a_{n} \rvert \in \mathbb{R}$}
\AXC{$(a_{n} \in \mathbb{R} \to \lvert a_{n} \rvert \in \mathbb{R}) \to \varphi_{75}$\lbl{A1}}
\BIC{$\varphi_{75}$}
\end{ptree}
\noindent$\Pi_{17}$:
\begin{ptree}
\AXC{$\forall x \in \mathbb{R}\,(x - 0 = x)$\lbl{[E28]}}
\AXC{$\varphi_{42}$\lbl{A2}}
\BIC{$\lvert a_{n} \rvert \in \mathbb{R} \to \lvert a_{n} \rvert - 0 = \lvert a_{n} \rvert$}
\AXC{$(\lvert a_{n} \rvert \in \mathbb{R} \to \lvert a_{n} \rvert - 0 = \lvert a_{n} \rvert) \to \varphi_{11} \to \varphi_{23}$\lbl{A1}}
\BIC{$\varphi_{11} \to \varphi_{23}$}
\end{ptree}
\noindent$\Pi_{18}$:
\begin{ptree}
\AXC{$\stackrel{\Pi_{7}}{\vdots}$}\noLine
\UIC{$\mathrm{Seq}(a) \to \mathrm{Conv}(\lvert a \rvert,0) \to n \in \mathbb{N} \to a_{n} \in \mathbb{R}$}
\AXC{$\stackrel{\Pi_{16}}{\vdots}$}\noLine
\UIC{$\varphi_{75}$}
\BIC{$\varphi_{11}$}
\AXC{$\stackrel{\Pi_{17}}{\vdots}$}\noLine
\UIC{$\varphi_{11} \to \varphi_{23}$}
\BIC{$\varphi_{23}$}
\AXC{$\varphi_{23} \to \varphi_{96}$\lbl{A1}}
\BIC{$\varphi_{96}$}
\end{ptree}
\noindent$\Pi_{19}$:
\begin{ptree}
\AXC{$\stackrel{\Pi_{15}}{\vdots}$}\noLine
\UIC{$\lvert a_{n} \rvert - 0 = \lvert a_{n} \rvert \to \varphi_{54}$}
\AXC{$x = x$\lbl{A5}}
\UIC{$\forall x\,(x = x)$}
\AXC{$\forall x\,(x = x) \to \lvert \lvert a_{n} \rvert - 0 \rvert = \lvert \lvert a_{n} \rvert - 0 \rvert$\lbl{A2}}
\BIC{$\lvert \lvert a_{n} \rvert - 0 \rvert = \lvert \lvert a_{n} \rvert - 0 \rvert$}
\AXC{$\stackrel{\Pi_{18}}{\vdots}$}\noLine
\UIC{$\varphi_{96}$}
\BIC{$(\lvert a_{n} \rvert - 0 = \lvert a_{n} \rvert \to \varphi_{54}) \to \mathrm{Seq}(a) \to \varphi_{25}$}
\BIC{$\mathrm{Seq}(a) \to \varphi_{25}$}
\end{ptree}
\noindent$\Pi_{20}$:
\begin{ptree}
\AXC{$z = w \to \varphi_{33}$\lbl{A6}}
\UIC{$\forall w\,(z = w \to \varphi_{33})$}
\UIC{$\forall z\,\forall w\,(z = w \to \varphi_{33})$}
\AXC{$\forall z\,\forall w\,(z = w \to \varphi_{33}) \to \varphi_{62}$\lbl{A2}}
\BIC{$\varphi_{62}$}
\AXC{$\varphi_{62} \to \lvert a \rvert(n) = \lvert a_{n} \rvert \to \varphi_{56}$\lbl{A2}}
\BIC{$\lvert a \rvert(n) = \lvert a_{n} \rvert \to \varphi_{56}$}
\end{ptree}
\noindent$\Pi_{21}$:
\begin{ptree}
\AXC{$\mathrm{Seq}(a) \to \mathrm{Seq}(a)$\lbl{A1}}
\AXC{$\forall a\,(\varphi_{16})$\lbl{[S11]}}
\AXC{$\forall a\,(\varphi_{16}) \to \varphi_{16}$\lbl{A2}}
\BIC{$\varphi_{16}$}
\AXC{$\varphi_{16} \to (\mathrm{Seq}(a) \to \mathrm{Seq}(a)) \to \mathrm{Seq}(a) \to \varphi_{34}$\lbl{A1}}
\BIC{$(\mathrm{Seq}(a) \to \mathrm{Seq}(a)) \to \mathrm{Seq}(a) \to \varphi_{34}$}
\BIC{$\mathrm{Seq}(a) \to \varphi_{34}$}
\end{ptree}
\noindent$\Pi_{22}$:
\begin{ptree}
\AXC{$\varphi_{44}$\lbl{A2}}
\AXC{$\stackrel{\Pi_{21}}{\vdots}$}\noLine
\UIC{$\mathrm{Seq}(a) \to \varphi_{34}$}
\AXC{$(\mathrm{Seq}(a) \to \varphi_{34}) \to \mathrm{Seq}(a) \to \varphi_{27}$\lbl{A1}}
\BIC{$\mathrm{Seq}(a) \to \varphi_{27}$}
\AXC{$(\mathrm{Seq}(a) \to \varphi_{27}) \to \varphi_{44} \to \mathrm{Seq}(a) \to \varphi_{20}$\lbl{A1}}
\BIC{$\varphi_{44} \to \mathrm{Seq}(a) \to \varphi_{20}$}
\BIC{$\mathrm{Seq}(a) \to \varphi_{20}$}
\end{ptree}
\noindent$\Pi_{23}$:
\begin{ptree}
\AXC{$\mathrm{Seq}(a) \to \mathrm{Conv}(\lvert a \rvert,0) \to n \in \mathbb{N} \to n \in \mathbb{N}$\lbl{A1}}
\AXC{$\stackrel{\Pi_{22}}{\vdots}$}\noLine
\UIC{$\mathrm{Seq}(a) \to \varphi_{20}$}
\AXC{$(\mathrm{Seq}(a) \to \varphi_{20}) \to \varphi_{71}$\lbl{A1}}
\BIC{$\varphi_{71}$}
\BIC{$\varphi_{17}$}
\AXC{$\varphi_{17} \to \varphi_{98}$\lbl{A1}}
\BIC{$\varphi_{98}$}
\end{ptree}
\noindent$\Pi_{24}$:
\begin{ptree}
\AXC{$\stackrel{\Pi_{20}}{\vdots}$}\noLine
\UIC{$\lvert a \rvert(n) = \lvert a_{n} \rvert \to \varphi_{56}$}
\AXC{$x = x$\lbl{A5}}
\UIC{$\forall x\,(x = x)$}
\AXC{$\forall x\,(x = x) \to \lvert \lvert a \rvert(n) - 0 \rvert = \lvert \lvert a \rvert(n) - 0 \rvert$\lbl{A2}}
\BIC{$\lvert \lvert a \rvert(n) - 0 \rvert = \lvert \lvert a \rvert(n) - 0 \rvert$}
\AXC{$\stackrel{\Pi_{23}}{\vdots}$}\noLine
\UIC{$\varphi_{98}$}
\BIC{$(\lvert a \rvert(n) = \lvert a_{n} \rvert \to \varphi_{56}) \to \mathrm{Seq}(a) \to \varphi_{37}$}
\BIC{$\mathrm{Seq}(a) \to \varphi_{37}$}
\end{ptree}
\noindent$\Pi_{25}$:
\begin{ptree}
\AXC{$\stackrel{\Pi_{19}}{\vdots}$}\noLine
\UIC{$\mathrm{Seq}(a) \to \varphi_{25}$}
\AXC{$\stackrel{\Pi_{24}}{\vdots}$}\noLine
\UIC{$\mathrm{Seq}(a) \to \varphi_{37}$}
\AXC{$(\mathrm{Seq}(a) \to \varphi_{37}) \to (\mathrm{Seq}(a) \to \varphi_{25}) \to \varphi_{91}$\lbl{A1}}
\BIC{$(\mathrm{Seq}(a) \to \varphi_{25}) \to \varphi_{91}$}
\BIC{$\varphi_{91}$}
\end{ptree}
\noindent$\Pi_{26}$:
\begin{ptree}
\AXC{$\stackrel{\Pi_{14}}{\vdots}$}\noLine
\UIC{$\varphi_{15}$}
\AXC{$\varphi_{15} \to \varphi_{61}$\lbl{A2}}
\BIC{$\varphi_{61}$}
\AXC{$\varphi_{61} \to \lvert \lvert a \rvert(n) - 0 \rvert = \lvert \lvert a_{n} \rvert - 0 \rvert \to \varphi_{51}$\lbl{A2}}
\BIC{$\lvert \lvert a \rvert(n) - 0 \rvert = \lvert \lvert a_{n} \rvert - 0 \rvert \to \varphi_{51}$}
\AXC{$\stackrel{\Pi_{25}}{\vdots}$}\noLine
\UIC{$\varphi_{91}$}
\BIC{$\mathrm{Seq}(a) \to \varphi_{28}$}
\end{ptree}
\noindent$\Pi_{27}$:
\begin{ptree}
\AXC{$\stackrel{\Pi_{26}}{\vdots}$}\noLine
\UIC{$\mathrm{Seq}(a) \to \varphi_{28}$}
\AXC{$(\mathrm{Seq}(a) \to \varphi_{28}) \to \varphi_{22} \to \varphi_{73} \to \mathrm{Seq}(a) \to \varphi_{18}$\lbl{A1}}
\BIC{$\varphi_{22} \to \varphi_{73} \to \mathrm{Seq}(a) \to \varphi_{18}$}
\end{ptree}
\noindent$\Pi_{28}$:
\begin{ptree}
\AXC{$\stackrel{\Pi_{7}}{\vdots}$}\noLine
\UIC{$\mathrm{Seq}(a) \to \mathrm{Conv}(\lvert a \rvert,0) \to n \in \mathbb{N} \to a_{n} \in \mathbb{R}$}
\AXC{$\stackrel{\Pi_{13}}{\vdots}$}\noLine
\UIC{$\varphi_{81}$}
\BIC{$\varphi_{22}$}
\AXC{$\stackrel{\Pi_{27}}{\vdots}$}\noLine
\UIC{$\varphi_{22} \to \varphi_{73} \to \mathrm{Seq}(a) \to \varphi_{18}$}
\BIC{$\varphi_{73} \to \mathrm{Seq}(a) \to \varphi_{18}$}
\end{ptree}
\noindent$\Pi_{29}$:
\begin{ptree}
\AXC{$\stackrel{\Pi_{12}}{\vdots}$}\noLine
\UIC{$\varphi_{15}$}
\AXC{$\varphi_{15} \to \varphi_{55}$\lbl{A2}}
\BIC{$\varphi_{55}$}
\AXC{$\varphi_{55} \to \varphi_{73}$\lbl{A2}}
\BIC{$\varphi_{73}$}
\AXC{$\stackrel{\Pi_{28}}{\vdots}$}\noLine
\UIC{$\varphi_{73} \to \mathrm{Seq}(a) \to \varphi_{18}$}
\BIC{$\mathrm{Seq}(a) \to \varphi_{18}$}
\AXC{$(\mathrm{Seq}(a) \to \varphi_{18}) \to \varphi_{21} \to \varphi_{72} \to \mathrm{Seq}(a) \to \varphi_{29}$\lbl{A1}}
\BIC{$\varphi_{21} \to \varphi_{72} \to \mathrm{Seq}(a) \to \varphi_{29}$}
\end{ptree}
\noindent$\Pi_{30}$:
\begin{ptree}
\AXC{$z = w \to \lvert a_{n} \rvert = \lvert z \rvert \to \lvert a_{n} \rvert = \lvert w \rvert$\lbl{A6}}
\UIC{$\forall w\,(z = w \to \lvert a_{n} \rvert = \lvert z \rvert \to \lvert a_{n} \rvert = \lvert w \rvert)$}
\UIC{$\forall z\,\forall w\,(z = w \to \lvert a_{n} \rvert = \lvert z \rvert \to \lvert a_{n} \rvert = \lvert w \rvert)$}
\AXC{$\varphi_{77}$\lbl{A2}}
\BIC{$\forall w\,(a_{n} = w \to \lvert a_{n} \rvert = \lvert a_{n} \rvert \to \lvert a_{n} \rvert = \lvert w \rvert)$}
\AXC{$\varphi_{87}$\lbl{A2}}
\BIC{$a_{n} = a_{n} - 0 \to \lvert a_{n} \rvert = \lvert a_{n} \rvert \to \lvert a_{n} \rvert = \lvert a_{n} - 0 \rvert$}
\AXC{$\stackrel{\Pi_{11}}{\vdots}$}\noLine
\UIC{$\varphi_{89}$}
\BIC{$\varphi_{21}$}
\AXC{$\stackrel{\Pi_{29}}{\vdots}$}\noLine
\UIC{$\varphi_{21} \to \varphi_{72} \to \mathrm{Seq}(a) \to \varphi_{29}$}
\BIC{$\varphi_{72} \to \mathrm{Seq}(a) \to \varphi_{29}$}
\end{ptree}
\noindent$\Pi_{31}$:
\begin{ptree}
\AXC{$\stackrel{\Pi_{4}}{\vdots}$}\noLine
\UIC{$\varphi_{15}$}
\AXC{$\varphi_{15} \to \varphi_{46}$\lbl{A2}}
\BIC{$\varphi_{46}$}
\AXC{$\varphi_{46} \to \varphi_{72}$\lbl{A2}}
\BIC{$\varphi_{72}$}
\AXC{$\stackrel{\Pi_{30}}{\vdots}$}\noLine
\UIC{$\varphi_{72} \to \mathrm{Seq}(a) \to \varphi_{29}$}
\BIC{$\mathrm{Seq}(a) \to \varphi_{29}$}
\AXC{$(\mathrm{Seq}(a) \to \varphi_{29}) \to \varphi_{40}$\lbl{A1}}
\BIC{$\varphi_{40}$}
\UIC{$\forall n\,(\varphi_{40})$}
\end{ptree}
\noindent$\Pi_{32}$:
\begin{ptree}
\AXC{$\varphi_{6}$\lbl{A4}}
\AXC{$\forall n\,(\varphi_{40}) \to \varphi_{50}$\lbl{A3}}
\AXC{$\stackrel{\Pi_{31}}{\vdots}$}\noLine
\UIC{$\forall n\,(\varphi_{40})$}
\AXC{$\forall n\,(\varphi_{40}) \to (\forall n\,(\varphi_{40}) \to \varphi_{50}) \to \varphi_{6} \to \mathrm{Seq}(a) \to \varphi_{35}$\lbl{A1}}
\BIC{$(\forall n\,(\varphi_{40}) \to \varphi_{50}) \to \varphi_{6} \to \mathrm{Seq}(a) \to \varphi_{35}$}
\BIC{$\varphi_{6} \to \mathrm{Seq}(a) \to \varphi_{35}$}
\BIC{$\mathrm{Seq}(a) \to \varphi_{35}$}
\end{ptree}
\noindent$\Pi_{33}$:
\begin{ptree}
\AXC{$\varphi_{68}$\lbl{A2}}
\AXC{$\stackrel{\Pi_{32}}{\vdots}$}\noLine
\UIC{$\mathrm{Seq}(a) \to \varphi_{35}$}
\AXC{$(\mathrm{Seq}(a) \to \varphi_{35}) \to \varphi_{68} \to \mathrm{Seq}(a) \to \varphi_{29}$\lbl{A1}}
\BIC{$\varphi_{68} \to \mathrm{Seq}(a) \to \varphi_{29}$}
\BIC{$\mathrm{Seq}(a) \to \varphi_{29}$}
\AXC{$(\mathrm{Seq}(a) \to \varphi_{29}) \to \varphi_{105}$\lbl{A1}}
\BIC{$\varphi_{105}$}
\end{ptree}
\noindent$\Pi_{34}$:
\begin{ptree}
\AXC{$\mathrm{Seq}(a) \to \mathrm{Conv}(\lvert a \rvert,0) \to \varepsilon \in \mathbb{R} \to 0 < \varepsilon \to \varphi_{49}$\lbl{A1}}
\AXC{$\stackrel{\Pi_{33}}{\vdots}$}\noLine
\UIC{$\varphi_{105}$}
\BIC{$\mathrm{Seq}(a) \to \mathrm{Conv}(\lvert a \rvert,0) \to \varepsilon \in \mathbb{R} \to 0 < \varepsilon \to \varphi_{79}$}
\AXC{$\varphi_{117}$\lbl{A1}}
\BIC{$\varphi_{112}$}
\end{ptree}
\noindent$\Pi_{35}$:
\begin{ptree}
\AXC{$\stackrel{\Pi_{1}}{\vdots}$}\noLine
\UIC{$\varphi_{58}$}
\AXC{$\mathrm{Seq}(a) \to \mathrm{Conv}(\lvert a \rvert,0) \to \varepsilon \in \mathbb{R} \to 0 < \varepsilon \to \varphi_{48}$\lbl{A1}}
\AXC{$\stackrel{\Pi_{3}}{\vdots}$}\noLine
\UIC{$\varphi_{102}$}
\BIC{$\mathrm{Seq}(a) \to \mathrm{Conv}(\lvert a \rvert,0) \to \varepsilon \in \mathbb{R} \to 0 < \varepsilon \to \varphi_{69}$}
\AXC{$\stackrel{\Pi_{34}}{\vdots}$}\noLine
\UIC{$\varphi_{112}$}
\BIC{$\varphi_{97}$}
\BIC{$\mathrm{Seq}(a) \to \mathrm{Conv}(\lvert a \rvert,0) \to \varepsilon \in \mathbb{R} \to 0 < \varepsilon \to \varphi_{63}$}
\AXC{$\varphi_{106}$\lbl{A1}}
\BIC{$\varphi_{70} \to n \in \mathbb{N} \to N \leq n \to \lvert a_{n} - 0 \rvert < \varepsilon$}
\UIC{$\forall n\,(\varphi_{70} \to n \in \mathbb{N} \to N \leq n \to \lvert a_{n} - 0 \rvert < \varepsilon)$}
\end{ptree}
\noindent$\Pi_{36}$:
\begin{ptree}
\AXC{$\varphi_{107}$\lbl{A3}}
\AXC{$\stackrel{\Pi_{35}}{\vdots}$}\noLine
\UIC{$\forall n\,(\varphi_{70} \to n \in \mathbb{N} \to N \leq n \to \lvert a_{n} - 0 \rvert < \varepsilon)$}
\AXC{$\forall n\,(\varphi_{70} \to n \in \mathbb{N} \to N \leq n \to \lvert a_{n} - 0 \rvert < \varepsilon) \to \varphi_{107} \to \varphi_{111}$\lbl{A1}}
\BIC{$\varphi_{107} \to \varphi_{111}$}
\BIC{$\varphi_{111}$}
\end{ptree}
\noindent$\Pi_{37}$:
\begin{ptree}
\AXC{$\mathrm{Seq}(a) \to \mathrm{Conv}(\lvert a \rvert,0) \to \varepsilon \in \mathbb{R} \to 0 < \varepsilon \to \varphi_{82}$\lbl{A1}}
\AXC{$\varphi_{103}$\lbl{A1}}
\BIC{$\mathrm{Seq}(a) \to \mathrm{Conv}(\lvert a \rvert,0) \to \varepsilon \in \mathbb{R} \to 0 < \varepsilon \to \varphi_{36}$}
\AXC{$\varphi_{114}$\lbl{A1}}
\BIC{$\varphi_{108}$}
\end{ptree}
\noindent$\Pi_{38}$:
\begin{ptree}
\AXC{$\varphi_{70} \to \forall n\,(\varphi_{70})$\lbl{A4}}
\AXC{$\stackrel{\Pi_{36}}{\vdots}$}\noLine
\UIC{$\varphi_{111}$}
\BIC{$\mathrm{Seq}(a) \to \mathrm{Conv}(\lvert a \rvert,0) \to \varepsilon \in \mathbb{R} \to 0 < \varepsilon \to \varphi_{65}$}
\AXC{$\stackrel{\Pi_{37}}{\vdots}$}\noLine
\UIC{$\varphi_{108}$}
\BIC{$\mathrm{Seq}(a) \to \mathrm{Conv}(\lvert a \rvert,0) \to \varepsilon \in \mathbb{R} \to 0 < \varepsilon \to \varphi_{78}$}
\AXC{$\varphi_{115}$\lbl{A1}}
\BIC{$\varphi_{101}$}
\end{ptree}
\noindent$\Pi_{39}$:
\begin{ptree}
\AXC{$\varphi_{13} \to \neg (N \in \mathbb{N} \land \forall n \in \mathbb{N}\,(N \leq n \to \lvert a_{n} - 0 \rvert < \varepsilon))$\lbl{A2}}
\AXC{$\stackrel{\Pi_{38}}{\vdots}$}\noLine
\UIC{$\varphi_{101}$}
\BIC{$\mathrm{Seq}(a) \to \mathrm{Conv}(\lvert a \rvert,0) \to \varepsilon \in \mathbb{R} \to 0 < \varepsilon \to \varphi_{80}$}
\AXC{$\varphi_{109}$\lbl{A1}}
\BIC{$\varphi_{67} \to \varphi_{14}$}
\UIC{$\forall N\,(\varphi_{67} \to \varphi_{14})$}
\end{ptree}
\noindent$\Pi_{40}$:
\begin{ptree}
\AXC{$\varphi_{64}$\lbl{[C3]}}
\AXC{$\varphi_{64} \to \varphi_{76}$\lbl{A2}}
\BIC{$\varphi_{76}$}
\AXC{$\varphi_{76} \to \forall \varepsilon\,(\mathrm{Conv}(\lvert a \rvert,0) \to \varepsilon \in \mathbb{R} \to \varphi_{38})$\lbl{A2}}
\BIC{$\forall \varepsilon\,(\mathrm{Conv}(\lvert a \rvert,0) \to \varepsilon \in \mathbb{R} \to \varphi_{38})$}
\AXC{$\varphi_{92}$\lbl{A2}}
\BIC{$\mathrm{Conv}(\lvert a \rvert,0) \to \varepsilon \in \mathbb{R} \to \varphi_{38}$}
\end{ptree}
\noindent$\Pi_{41}$:
\begin{ptree}
\AXC{$\stackrel{\Pi_{40}}{\vdots}$}\noLine
\UIC{$\mathrm{Conv}(\lvert a \rvert,0) \to \varepsilon \in \mathbb{R} \to \varphi_{38}$}
\AXC{$(\mathrm{Conv}(\lvert a \rvert,0) \to \varepsilon \in \mathbb{R} \to \varphi_{38}) \to \varphi_{99}$\lbl{A1}}
\BIC{$\varphi_{99}$}
\end{ptree}
\noindent$\Pi_{42}$:
\begin{ptree}
\AXC{$\mathrm{Seq}(a) \to \mathrm{Conv}(\lvert a \rvert,0) \to \varepsilon \in \mathbb{R} \to \varepsilon \in \mathbb{R}$\lbl{A1}}
\AXC{$\mathrm{Seq}(a) \to \mathrm{Conv}(\lvert a \rvert,0) \to \mathrm{Conv}(\lvert a \rvert,0)$\lbl{A1}}
\AXC{$\stackrel{\Pi_{41}}{\vdots}$}\noLine
\UIC{$\varphi_{99}$}
\BIC{$\varphi_{94}$}
\BIC{$\varphi_{12} \to \mathrm{Seq}(a) \to \mathrm{Conv}(\lvert a \rvert,0) \to \varepsilon \in \mathbb{R} \to \varphi_{38}$}
\end{ptree}
\noindent$\Pi_{43}$:
\begin{ptree}
\AXC{$\stackrel{\Pi_{39}}{\vdots}$}\noLine
\UIC{$\forall N\,(\varphi_{67} \to \varphi_{14})$}
\AXC{$\varphi_{12}$\lbl{A1}}
\AXC{$\stackrel{\Pi_{42}}{\vdots}$}\noLine
\UIC{$\varphi_{12} \to \mathrm{Seq}(a) \to \mathrm{Conv}(\lvert a \rvert,0) \to \varepsilon \in \mathbb{R} \to \varphi_{38}$}
\BIC{$\mathrm{Seq}(a) \to \mathrm{Conv}(\lvert a \rvert,0) \to \varepsilon \in \mathbb{R} \to \varphi_{38}$}
\AXC{$\varphi_{120}$\lbl{A1}}
\BIC{$\forall N\,(\varphi_{67} \to \varphi_{14}) \to \varphi_{119}$}
\BIC{$\varphi_{119}$}
\end{ptree}
\noindent$\Pi_{44}$:
\begin{ptree}
\AXC{$\varphi_{95}$\lbl{A3}}
\AXC{$\varphi_{67} \to \forall N\,(\varphi_{67})$\lbl{A4}}
\AXC{$\forall N\,(\varphi_{67} \to \varphi_{14}) \to \forall N\,(\varphi_{67}) \to \forall N\,(\varphi_{14})$\lbl{A3}}
\AXC{$\stackrel{\Pi_{43}}{\vdots}$}\noLine
\UIC{$\varphi_{119}$}
\BIC{$(\varphi_{67} \to \forall N\,(\varphi_{67})) \to \mathrm{Seq}(a) \to \mathrm{Conv}(\lvert a \rvert,0) \to \varphi_{45}$}
\BIC{$\mathrm{Seq}(a) \to \mathrm{Conv}(\lvert a \rvert,0) \to \varphi_{45}$}
\AXC{$\varphi_{93}$\lbl{A1}}
\BIC{$\mathrm{Seq}(a) \land \mathrm{Conv}(\lvert a \rvert,0) \to \varphi_{45}$}
\UIC{$\forall \varepsilon\,(\mathrm{Seq}(a) \land \mathrm{Conv}(\lvert a \rvert,0) \to \varphi_{45})$}
\AXC{$\varphi_{116}$\lbl{A1}}
\BIC{$\varphi_{95} \to \varphi_{8} \to \mathrm{Seq}(a) \to \mathrm{Conv}(\lvert a \rvert,0) \to \varphi_{53}$}
\BIC{$\varphi_{8} \to \mathrm{Seq}(a) \to \mathrm{Conv}(\lvert a \rvert,0) \to \varphi_{53}$}
\end{ptree}
\noindent$\Pi_{45}$:
\begin{ptree}
\AXC{$\forall a\,(\varphi_{84})$\lbl{DefConv}}
\AXC{$\forall a\,(\varphi_{84}) \to \varphi_{84}$\lbl{A2}}
\BIC{$\varphi_{84}$}
\AXC{$\varphi_{84} \to (\mathrm{Conv}(a,0) \leftrightarrow \mathrm{Seq}(a) \land 0 \in \mathbb{R} \land \varphi_{53})$\lbl{A2}}
\BIC{$\mathrm{Conv}(a,0) \leftrightarrow \mathrm{Seq}(a) \land 0 \in \mathbb{R} \land \varphi_{53}$}
\AXC{$\varphi_{104}$\lbl{A1}}
\BIC{$(\mathrm{Seq}(a) \to \mathrm{Seq}(a)) \to 0 \in \mathbb{R} \to \varphi_{88}$}
\end{ptree}
\noindent$\Pi_{46}$:
\begin{ptree}
\AXC{$0 \in \mathbb{R}$\lbl{(R1)}}
\AXC{$\mathrm{Seq}(a) \to \mathrm{Seq}(a)$\lbl{A1}}
\AXC{$\stackrel{\Pi_{45}}{\vdots}$}\noLine
\UIC{$(\mathrm{Seq}(a) \to \mathrm{Seq}(a)) \to 0 \in \mathbb{R} \to \varphi_{88}$}
\BIC{$0 \in \mathbb{R} \to \varphi_{88}$}
\BIC{$\varphi_{88}$}
\end{ptree}
\noindent$\Pi_{47}$:
\begin{ptree}
\AXC{$z = w \to \mathrm{Conv}(\lvert a \rvert,z) \to \mathrm{Conv}(\lvert a \rvert,w)$\lbl{A6}}
\UIC{$\varphi_{2}$}
\UIC{$\forall z\,(\varphi_{2})$}
\AXC{$\forall z\,(\varphi_{2}) \to \varphi_{5}$\lbl{A2}}
\BIC{$\varphi_{5}$}
\AXC{$\varphi_{5} \to \varphi_{3}$\lbl{A2}}
\BIC{$\varphi_{3}$}
\end{ptree}
\noindent$\Pi_{48}$:
\begin{ptree}
\AXC{$\forall a\,\forall \alpha\,(\mathrm{Conv}(a,\alpha) \to \mathrm{Conv}(\lvert a \rvert,\lvert \alpha \rvert))$\lbl{定理2.15(1)}}
\AXC{$\varphi_{41}$\lbl{A2}}
\BIC{$\forall \alpha\,(\mathrm{Conv}(a,\alpha) \to \mathrm{Conv}(\lvert a \rvert,\lvert \alpha \rvert))$}
\AXC{$\varphi_{26}$\lbl{A2}}
\BIC{$\mathrm{Conv}(a,0) \to \mathrm{Conv}(\lvert a \rvert,\lvert 0 \rvert)$}
\end{ptree}
\noindent$\Pi_{49}$:
\begin{ptree}
\AXC{$\stackrel{\Pi_{48}}{\vdots}$}\noLine
\UIC{$\mathrm{Conv}(a,0) \to \mathrm{Conv}(\lvert a \rvert,\lvert 0 \rvert)$}
\AXC{$(\mathrm{Conv}(a,0) \to \mathrm{Conv}(\lvert a \rvert,\lvert 0 \rvert)) \to \varphi_{31}$\lbl{A1}}
\BIC{$\varphi_{31}$}
\end{ptree}
\noindent$\Pi_{50}$:
\begin{ptree}
\AXC{$\forall x\,(0 \leq x \to \lvert x \rvert = x)$\lbl{[A1]}}
\AXC{$\forall x\,(0 \leq x \to \lvert x \rvert = x) \to 0 \leq 0 \to \lvert 0 \rvert = 0$\lbl{A2}}
\BIC{$0 \leq 0 \to \lvert 0 \rvert = 0$}
\AXC{$\varphi_{4}$\lbl{A1}}
\BIC{$(\mathrm{Seq}(a) \to 0 \leq 0) \to \mathrm{Seq}(a) \to \lvert 0 \rvert = 0$}
\end{ptree}
\noindent$\Pi_{51}$:
\begin{ptree}
\AXC{$x = x$\lbl{A5}}
\UIC{$\forall x\,(x = x)$}
\AXC{$\forall x\,(x = x) \to 0 = 0$\lbl{A2}}
\BIC{$0 = 0$}
\AXC{$0 = 0 \to \mathrm{Seq}(a) \to 0 \leq 0$\lbl{A1}}
\BIC{$\mathrm{Seq}(a) \to 0 \leq 0$}
\AXC{$\stackrel{\Pi_{50}}{\vdots}$}\noLine
\UIC{$(\mathrm{Seq}(a) \to 0 \leq 0) \to \mathrm{Seq}(a) \to \lvert 0 \rvert = 0$}
\BIC{$\mathrm{Seq}(a) \to \lvert 0 \rvert = 0$}
\end{ptree}
\noindent$\Pi_{52}$:
\begin{ptree}
\AXC{$\stackrel{\Pi_{51}}{\vdots}$}\noLine
\UIC{$\mathrm{Seq}(a) \to \lvert 0 \rvert = 0$}
\AXC{$(\mathrm{Seq}(a) \to \lvert 0 \rvert = 0) \to \varphi_{83}$\lbl{A1}}
\BIC{$\varphi_{83}$}
\end{ptree}
\noindent$\Pi_{53}$:
\begin{ptree}
\AXC{$\stackrel{\Pi_{47}}{\vdots}$}\noLine
\UIC{$\varphi_{3}$}
\AXC{$\mathrm{Seq}(a) \to \mathrm{Conv}(a,0) \to \mathrm{Conv}(a,0)$\lbl{A1}}
\AXC{$\stackrel{\Pi_{49}}{\vdots}$}\noLine
\UIC{$\varphi_{31}$}
\BIC{$\mathrm{Seq}(a) \to \mathrm{Conv}(a,0) \to \mathrm{Conv}(\lvert a \rvert,\lvert 0 \rvert)$}
\AXC{$\stackrel{\Pi_{52}}{\vdots}$}\noLine
\UIC{$\varphi_{83}$}
\BIC{$\varphi_{3} \to \mathrm{Seq}(a) \to \mathrm{Conv}(a,0) \to \mathrm{Conv}(\lvert a \rvert,0)$}
\BIC{$\mathrm{Seq}(a) \to \mathrm{Conv}(a,0) \to \mathrm{Conv}(\lvert a \rvert,0)$}
\AXC{$\varphi_{66}$\lbl{A1}}
\BIC{$\varphi_{32}$}
\end{ptree}
\noindent 主木:
\begin{ptree}
\AXC{$\varphi_{8}$\lbl{A4}}
\AXC{$\stackrel{\Pi_{44}}{\vdots}$}\noLine
\UIC{$\varphi_{8} \to \mathrm{Seq}(a) \to \mathrm{Conv}(\lvert a \rvert,0) \to \varphi_{53}$}
\BIC{$\mathrm{Seq}(a) \to \mathrm{Conv}(\lvert a \rvert,0) \to \varphi_{53}$}
\AXC{$\stackrel{\Pi_{46}}{\vdots}$}\noLine
\UIC{$\varphi_{88}$}
\BIC{$\mathrm{Seq}(a) \to \mathrm{Conv}(\lvert a \rvert,0) \to \mathrm{Conv}(a,0)$}
\AXC{$\stackrel{\Pi_{53}}{\vdots}$}\noLine
\UIC{$\varphi_{32}$}
\BIC{$\mathrm{Seq}(a) \to (\mathrm{Conv}(a,0) \leftrightarrow \mathrm{Conv}(\lvert a \rvert,0))$}
\UIC{$\varphi_{1}$}
\end{ptree}
\begin{abbrevlist}
\item $\varphi_{1} :\equiv \forall a\,(\mathrm{Seq}(a) \to (\mathrm{Conv}(a,0) \leftrightarrow \mathrm{Conv}(\lvert a \rvert,0)))$
\item $\varphi_{2} :\equiv \forall w\,(z = w \to \mathrm{Conv}(\lvert a \rvert,z) \to \mathrm{Conv}(\lvert a \rvert,w))$
\item $\varphi_{3} :\equiv \lvert 0 \rvert = 0 \to \mathrm{Conv}(\lvert a \rvert,\lvert 0 \rvert) \to \mathrm{Conv}(\lvert a \rvert,0)$
\item $\varphi_{4} :\equiv (0 \leq 0 \to \lvert 0 \rvert = 0) \to (\mathrm{Seq}(a) \to 0 \leq 0) \to \mathrm{Seq}(a) \to \lvert 0 \rvert = 0$
\item $\varphi_{5} :\equiv \forall w\,(\lvert 0 \rvert = w \to \mathrm{Conv}(\lvert a \rvert,\lvert 0 \rvert) \to \mathrm{Conv}(\lvert a \rvert,w))$
\item $\varphi_{6} :\equiv \mathrm{Seq}(a) \land \mathrm{Conv}(\lvert a \rvert,0) \to \forall n\,(\mathrm{Seq}(a) \land \mathrm{Conv}(\lvert a \rvert,0))$
\item $\varphi_{7} :\equiv \forall x\,\forall y\,(x = y \to y = x) \to \forall y\,(a_{n} - 0 = y \to y = a_{n} - 0)$
\item $\varphi_{8} :\equiv \mathrm{Seq}(a) \land \mathrm{Conv}(\lvert a \rvert,0) \to \forall \varepsilon\,(\mathrm{Seq}(a) \land \mathrm{Conv}(\lvert a \rvert,0))$
\item $\varphi_{9} :\equiv \mathrm{Seq}(a) \to \mathrm{Conv}(\lvert a \rvert,0) \to n \in \mathbb{N} \to a_{n} - 0 = a_{n}$
\item $\varphi_{10} :\equiv \mathrm{Seq}(a) \to \mathrm{Conv}(\lvert a \rvert,0) \to n \in \mathbb{N} \to a_{n} = a_{n} - 0$
\item $\varphi_{11} :\equiv \mathrm{Seq}(a) \to \mathrm{Conv}(\lvert a \rvert,0) \to n \in \mathbb{N} \to \lvert a_{n} \rvert \in \mathbb{R}$
\item $\varphi_{12} :\equiv \mathrm{Seq}(a) \to \mathrm{Conv}(\lvert a \rvert,0) \to \varepsilon \in \mathbb{R} \to 0 < \varepsilon \to 0 < \varepsilon$
\item $\varphi_{13} :\equiv \forall N\,(\neg (N \in \mathbb{N} \land \forall n \in \mathbb{N}\,(N \leq n \to \lvert a_{n} - 0 \rvert < \varepsilon)))$
\item $\varphi_{14} :\equiv \neg (N \in \mathbb{N} \land \forall n \in \mathbb{N}\,(N \leq n \to \lvert \lvert a \rvert(n) - 0 \rvert < \varepsilon))$
\item $\varphi_{15} :\equiv \forall y\,\forall z\,(\lvert \lvert a \rvert(n) - 0 \rvert = y \to y = z \to \lvert \lvert a \rvert(n) - 0 \rvert = z)$
\item $\varphi_{16} :\equiv \mathrm{Seq}(a) \to \mathrm{Seq}(\lvert a \rvert) \land \forall n \in \mathbb{N}\,(\lvert a \rvert(n) = \lvert a_{n} \rvert)$
\item $\varphi_{17} :\equiv \mathrm{Seq}(a) \to \mathrm{Conv}(\lvert a \rvert,0) \to n \in \mathbb{N} \to \lvert a \rvert(n) = \lvert a_{n} \rvert$
\item $\varphi_{18} :\equiv \mathrm{Conv}(\lvert a \rvert,0) \to n \in \mathbb{N} \to \lvert \lvert a \rvert(n) - 0 \rvert = \lvert a_{n} \rvert$
\item $\varphi_{19} :\equiv n \in \mathbb{N} \to N \leq n \to \forall n \in \mathbb{N}\,(N \leq n \to \lvert \lvert a \rvert(n) - 0 \rvert < \varepsilon)$
\item $\varphi_{20} :\equiv \mathrm{Conv}(\lvert a \rvert,0) \to n \in \mathbb{N} \to n \in \mathbb{N} \to \lvert a \rvert(n) = \lvert a_{n} \rvert$
\item $\varphi_{21} :\equiv \mathrm{Seq}(a) \to \mathrm{Conv}(\lvert a \rvert,0) \to n \in \mathbb{N} \to \lvert a_{n} \rvert = \lvert a_{n} - 0 \rvert$
\item $\varphi_{22} :\equiv \mathrm{Seq}(a) \to \mathrm{Conv}(\lvert a \rvert,0) \to n \in \mathbb{N} \to \lvert \lvert a_{n} \rvert \rvert = \lvert a_{n} \rvert$
\item $\varphi_{23} :\equiv \mathrm{Seq}(a) \to \mathrm{Conv}(\lvert a \rvert,0) \to n \in \mathbb{N} \to \lvert a_{n} \rvert - 0 = \lvert a_{n} \rvert$
\item $\varphi_{24} :\equiv \forall w\,(z = w \to \lvert \lvert a_{n} \rvert - 0 \rvert = \lvert z \rvert \to \lvert \lvert a_{n} \rvert - 0 \rvert = \lvert w \rvert)$
\item $\varphi_{25} :\equiv \mathrm{Conv}(\lvert a \rvert,0) \to n \in \mathbb{N} \to \lvert \lvert a_{n} \rvert - 0 \rvert = \lvert \lvert a_{n} \rvert \rvert$
\item $\varphi_{26} :\equiv \forall \alpha\,(\mathrm{Conv}(a,\alpha) \to \mathrm{Conv}(\lvert a \rvert,\lvert \alpha \rvert)) \to \mathrm{Conv}(a,0) \to \mathrm{Conv}(\lvert a \rvert,\lvert 0 \rvert)$
\item $\varphi_{27} :\equiv \mathrm{Conv}(\lvert a \rvert,0) \to n \in \mathbb{N} \to \forall n \in \mathbb{N}\,(\lvert a \rvert(n) = \lvert a_{n} \rvert)$
\item $\varphi_{28} :\equiv \mathrm{Conv}(\lvert a \rvert,0) \to n \in \mathbb{N} \to \lvert \lvert a \rvert(n) - 0 \rvert = \lvert \lvert a_{n} \rvert \rvert$
\item $\varphi_{29} :\equiv \mathrm{Conv}(\lvert a \rvert,0) \to n \in \mathbb{N} \to \lvert \lvert a \rvert(n) - 0 \rvert = \lvert a_{n} - 0 \rvert$
\item $\varphi_{30} :\equiv \forall w\,(\lvert \lvert a \rvert(n) - 0 \rvert = w \to \lvert \lvert a \rvert(n) - 0 \rvert < \varepsilon \to w < \varepsilon)$
\item $\varphi_{31} :\equiv (\mathrm{Seq}(a) \to \mathrm{Conv}(a,0) \to \mathrm{Conv}(a,0)) \to \mathrm{Seq}(a) \to \mathrm{Conv}(a,0) \to \mathrm{Conv}(\lvert a \rvert,\lvert 0 \rvert)$
\item $\varphi_{32} :\equiv (\mathrm{Seq}(a) \to \mathrm{Conv}(\lvert a \rvert,0) \to \mathrm{Conv}(a,0)) \to \mathrm{Seq}(a) \to (\mathrm{Conv}(a,0) \leftrightarrow \mathrm{Conv}(\lvert a \rvert,0))$
\item $\varphi_{33} :\equiv \lvert \lvert a \rvert(n) - 0 \rvert = \lvert z - 0 \rvert \to \lvert \lvert a \rvert(n) - 0 \rvert = \lvert w - 0 \rvert$
\item $\varphi_{34} :\equiv \mathrm{Conv}(\lvert a \rvert,0) \to \mathrm{Seq}(\lvert a \rvert) \land \forall n \in \mathbb{N}\,(\lvert a \rvert(n) = \lvert a_{n} \rvert)$
\item $\varphi_{35} :\equiv \mathrm{Conv}(\lvert a \rvert,0) \to \forall n \in \mathbb{N}\,(\lvert \lvert a \rvert(n) - 0 \rvert = \lvert a_{n} - 0 \rvert)$
\item $\varphi_{36} :\equiv N \in \mathbb{N} \land \forall n \in \mathbb{N}\,(N \leq n \to \lvert \lvert a \rvert(n) - 0 \rvert < \varepsilon) \to N \in \mathbb{N}$
\item $\varphi_{37} :\equiv \mathrm{Conv}(\lvert a \rvert,0) \to n \in \mathbb{N} \to \lvert \lvert a \rvert(n) - 0 \rvert = \lvert \lvert a_{n} \rvert - 0 \rvert$
\item $\varphi_{38} :\equiv 0 < \varepsilon \to \exists N \in \mathbb{N}\,\forall n \in \mathbb{N}\,(N \leq n \to \lvert \lvert a \rvert(n) - 0 \rvert < \varepsilon)$
\item $\varphi_{39} :\equiv \forall n\,(\mathrm{Seq}(a) \to n \in \mathbb{N} \to a_{n} \in \mathbb{R}) \to \mathrm{Seq}(a) \to n \in \mathbb{N} \to a_{n} \in \mathbb{R}$
\item $\varphi_{40} :\equiv \mathrm{Seq}(a) \land \mathrm{Conv}(\lvert a \rvert,0) \to n \in \mathbb{N} \to \lvert \lvert a \rvert(n) - 0 \rvert = \lvert a_{n} - 0 \rvert$
\item $\varphi_{41} :\equiv \forall a\,\forall \alpha\,(\mathrm{Conv}(a,\alpha) \to \mathrm{Conv}(\lvert a \rvert,\lvert \alpha \rvert)) \to \forall \alpha\,(\mathrm{Conv}(a,\alpha) \to \mathrm{Conv}(\lvert a \rvert,\lvert \alpha \rvert))$
\item $\varphi_{42} :\equiv \forall x \in \mathbb{R}\,(x - 0 = x) \to \lvert a_{n} \rvert \in \mathbb{R} \to \lvert a_{n} \rvert - 0 = \lvert a_{n} \rvert$
\item $\varphi_{43} :\equiv \forall y\,(a_{n} - 0 = y \to y = a_{n} - 0) \to a_{n} - 0 = a_{n} \to a_{n} = a_{n} - 0$
\item $\varphi_{44} :\equiv \forall n \in \mathbb{N}\,(\lvert a \rvert(n) = \lvert a_{n} \rvert) \to n \in \mathbb{N} \to \lvert a \rvert(n) = \lvert a_{n} \rvert$
\item $\varphi_{45} :\equiv \varepsilon \in \mathbb{R} \to 0 < \varepsilon \to \exists N \in \mathbb{N}\,\forall n \in \mathbb{N}\,(N \leq n \to \lvert a_{n} - 0 \rvert < \varepsilon)$
\item $\varphi_{46} :\equiv \forall z\,(\lvert \lvert a \rvert(n) - 0 \rvert = \lvert a_{n} \rvert \to \lvert a_{n} \rvert = z \to \lvert \lvert a \rvert(n) - 0 \rvert = z)$
\item $\varphi_{47} :\equiv \forall a\,\forall n\,(\mathrm{Seq}(a) \to n \in \mathbb{N} \to a_{n} \in \mathbb{R}) \to \forall n\,(\mathrm{Seq}(a) \to n \in \mathbb{N} \to a_{n} \in \mathbb{R})$
\item $\varphi_{48} :\equiv N \in \mathbb{N} \land \forall n \in \mathbb{N}\,(N \leq n \to \lvert \lvert a \rvert(n) - 0 \rvert < \varepsilon) \to n \in \mathbb{N} \to N \leq n \to N \leq n$
\item $\varphi_{49} :\equiv N \in \mathbb{N} \land \forall n \in \mathbb{N}\,(N \leq n \to \lvert \lvert a \rvert(n) - 0 \rvert < \varepsilon) \to n \in \mathbb{N} \to n \in \mathbb{N}$
\item $\varphi_{50} :\equiv \forall n\,(\mathrm{Seq}(a) \land \mathrm{Conv}(\lvert a \rvert,0)) \to \forall n \in \mathbb{N}\,(\lvert \lvert a \rvert(n) - 0 \rvert = \lvert a_{n} - 0 \rvert)$
\item $\varphi_{51} :\equiv \lvert \lvert a_{n} \rvert - 0 \rvert = \lvert \lvert a_{n} \rvert \rvert \to \lvert \lvert a \rvert(n) - 0 \rvert = \lvert \lvert a_{n} \rvert \rvert$
\item $\varphi_{52} :\equiv \forall x \in \mathbb{R}\,(\lvert \lvert x \rvert \rvert = \lvert x \rvert) \to a_{n} \in \mathbb{R} \to \lvert \lvert a_{n} \rvert \rvert = \lvert a_{n} \rvert$
\item $\varphi_{53} :\equiv \forall \varepsilon \in \mathbb{R}\,(0 < \varepsilon \to \exists N \in \mathbb{N}\,\forall n \in \mathbb{N}\,(N \leq n \to \lvert a_{n} - 0 \rvert < \varepsilon))$
\item $\varphi_{54} :\equiv \lvert \lvert a_{n} \rvert - 0 \rvert = \lvert \lvert a_{n} \rvert - 0 \rvert \to \lvert \lvert a_{n} \rvert - 0 \rvert = \lvert \lvert a_{n} \rvert \rvert$
\item $\varphi_{55} :\equiv \forall z\,(\lvert \lvert a \rvert(n) - 0 \rvert = \lvert \lvert a_{n} \rvert \rvert \to \lvert \lvert a_{n} \rvert \rvert = z \to \lvert \lvert a \rvert(n) - 0 \rvert = z)$
\item $\varphi_{56} :\equiv \lvert \lvert a \rvert(n) - 0 \rvert = \lvert \lvert a \rvert(n) - 0 \rvert \to \lvert \lvert a \rvert(n) - 0 \rvert = \lvert \lvert a_{n} \rvert - 0 \rvert$
\item $\varphi_{57} :\equiv \forall n \in \mathbb{N}\,(N \leq n \to \lvert \lvert a \rvert(n) - 0 \rvert < \varepsilon) \to n \in \mathbb{N} \to N \leq n \to \lvert \lvert a \rvert(n) - 0 \rvert < \varepsilon$
\item $\varphi_{58} :\equiv \lvert \lvert a \rvert(n) - 0 \rvert = \lvert a_{n} - 0 \rvert \to \lvert \lvert a \rvert(n) - 0 \rvert < \varepsilon \to \lvert a_{n} - 0 \rvert < \varepsilon$
\item $\varphi_{59} :\equiv \forall w\,(\lvert a_{n} \rvert - 0 = w \to \lvert \lvert a_{n} \rvert - 0 \rvert = \lvert \lvert a_{n} \rvert - 0 \rvert \to \lvert \lvert a_{n} \rvert - 0 \rvert = \lvert w \rvert)$
\item $\varphi_{60} :\equiv (\mathrm{Seq}(a) \to \mathrm{Conv}(\lvert a \rvert,0) \to n \in \mathbb{N} \to n \in \mathbb{N}) \to \mathrm{Seq}(a) \to \mathrm{Conv}(\lvert a \rvert,0) \to n \in \mathbb{N} \to a_{n} \in \mathbb{R}$
\item $\varphi_{61} :\equiv \forall z\,(\lvert \lvert a \rvert(n) - 0 \rvert = \lvert \lvert a_{n} \rvert - 0 \rvert \to \lvert \lvert a_{n} \rvert - 0 \rvert = z \to \lvert \lvert a \rvert(n) - 0 \rvert = z)$
\item $\varphi_{62} :\equiv \forall w\,(\lvert a \rvert(n) = w \to \lvert \lvert a \rvert(n) - 0 \rvert = \lvert \lvert a \rvert(n) - 0 \rvert \to \lvert \lvert a \rvert(n) - 0 \rvert = \lvert w - 0 \rvert)$
\item $\varphi_{63} :\equiv N \in \mathbb{N} \land \forall n \in \mathbb{N}\,(N \leq n \to \lvert \lvert a \rvert(n) - 0 \rvert < \varepsilon) \to n \in \mathbb{N} \to N \leq n \to \lvert a_{n} - 0 \rvert < \varepsilon$
\item $\varphi_{64} :\equiv \forall a\,\forall \alpha\,\forall \varepsilon\,(\mathrm{Conv}(a,\alpha) \to \varepsilon \in \mathbb{R} \to 0 < \varepsilon \to \exists N \in \mathbb{N}\,\forall n \in \mathbb{N}\,(N \leq n \to \lvert a_{n} - \alpha \rvert < \varepsilon))$
\item $\varphi_{65} :\equiv N \in \mathbb{N} \land \forall n \in \mathbb{N}\,(N \leq n \to \lvert \lvert a \rvert(n) - 0 \rvert < \varepsilon) \to \forall n \in \mathbb{N}\,(N \leq n \to \lvert a_{n} - 0 \rvert < \varepsilon)$
\item $\varphi_{66} :\equiv (\mathrm{Seq}(a) \to \mathrm{Conv}(a,0) \to \mathrm{Conv}(\lvert a \rvert,0)) \to \varphi_{32}$
\item $\varphi_{67} :\equiv \mathrm{Seq}(a) \land \mathrm{Conv}(\lvert a \rvert,0) \land \varepsilon \in \mathbb{R} \land 0 < \varepsilon \land \neg \exists N \in \mathbb{N}\,\forall n \in \mathbb{N}\,(N \leq n \to \lvert a_{n} - 0 \rvert < \varepsilon)$
\item $\varphi_{68} :\equiv \forall n \in \mathbb{N}\,(\lvert \lvert a \rvert(n) - 0 \rvert = \lvert a_{n} - 0 \rvert) \to n \in \mathbb{N} \to \lvert \lvert a \rvert(n) - 0 \rvert = \lvert a_{n} - 0 \rvert$
\item $\varphi_{69} :\equiv N \in \mathbb{N} \land \forall n \in \mathbb{N}\,(N \leq n \to \lvert \lvert a \rvert(n) - 0 \rvert < \varepsilon) \to n \in \mathbb{N} \to N \leq n \to \lvert \lvert a \rvert(n) - 0 \rvert < \varepsilon$
\item $\varphi_{70} :\equiv \mathrm{Seq}(a) \land \mathrm{Conv}(\lvert a \rvert,0) \land \varepsilon \in \mathbb{R} \land 0 < \varepsilon \land (N \in \mathbb{N} \land \forall n \in \mathbb{N}\,(N \leq n \to \lvert \lvert a \rvert(n) - 0 \rvert < \varepsilon))$
\item $\varphi_{71} :\equiv (\mathrm{Seq}(a) \to \mathrm{Conv}(\lvert a \rvert,0) \to n \in \mathbb{N} \to n \in \mathbb{N}) \to \varphi_{17}$
\item $\varphi_{72} :\equiv \lvert \lvert a \rvert(n) - 0 \rvert = \lvert a_{n} \rvert \to \lvert a_{n} \rvert = \lvert a_{n} - 0 \rvert \to \lvert \lvert a \rvert(n) - 0 \rvert = \lvert a_{n} - 0 \rvert$
\item $\varphi_{73} :\equiv \lvert \lvert a \rvert(n) - 0 \rvert = \lvert \lvert a_{n} \rvert \rvert \to \lvert \lvert a_{n} \rvert \rvert = \lvert a_{n} \rvert \to \lvert \lvert a \rvert(n) - 0 \rvert = \lvert a_{n} \rvert$
\item $\varphi_{74} :\equiv (\mathrm{Seq}(a) \to \mathrm{Conv}(\lvert a \rvert,0) \to n \in \mathbb{N} \to a_{n} \in \mathbb{R}) \to \varphi_{9}$
\item $\varphi_{75} :\equiv (\mathrm{Seq}(a) \to \mathrm{Conv}(\lvert a \rvert,0) \to n \in \mathbb{N} \to a_{n} \in \mathbb{R}) \to \varphi_{11}$
\item $\varphi_{76} :\equiv \forall \alpha\,\forall \varepsilon\,(\mathrm{Conv}(\lvert a \rvert,\alpha) \to \varepsilon \in \mathbb{R} \to 0 < \varepsilon \to \exists N \in \mathbb{N}\,\forall n \in \mathbb{N}\,(N \leq n \to \lvert \lvert a \rvert(n) - \alpha \rvert < \varepsilon))$
\item $\varphi_{77} :\equiv \forall z\,\forall w\,(z = w \to \lvert a_{n} \rvert = \lvert z \rvert \to \lvert a_{n} \rvert = \lvert w \rvert) \to \forall w\,(a_{n} = w \to \lvert a_{n} \rvert = \lvert a_{n} \rvert \to \lvert a_{n} \rvert = \lvert w \rvert)$
\item $\varphi_{78} :\equiv N \in \mathbb{N} \land \forall n \in \mathbb{N}\,(N \leq n \to \lvert \lvert a \rvert(n) - 0 \rvert < \varepsilon) \to N \in \mathbb{N} \land \forall n \in \mathbb{N}\,(N \leq n \to \lvert a_{n} - 0 \rvert < \varepsilon)$
\item $\varphi_{79} :\equiv N \in \mathbb{N} \land \forall n \in \mathbb{N}\,(N \leq n \to \lvert \lvert a \rvert(n) - 0 \rvert < \varepsilon) \to n \in \mathbb{N} \to N \leq n \to \lvert \lvert a \rvert(n) - 0 \rvert = \lvert a_{n} - 0 \rvert$
\item $\varphi_{80} :\equiv N \in \mathbb{N} \land \forall n \in \mathbb{N}\,(N \leq n \to \lvert \lvert a \rvert(n) - 0 \rvert < \varepsilon) \to \exists N \in \mathbb{N}\,\forall n \in \mathbb{N}\,(N \leq n \to \lvert a_{n} - 0 \rvert < \varepsilon)$
\item $\varphi_{81} :\equiv (\mathrm{Seq}(a) \to \mathrm{Conv}(\lvert a \rvert,0) \to n \in \mathbb{N} \to a_{n} \in \mathbb{R}) \to \varphi_{22}$
\item $\varphi_{82} :\equiv N \in \mathbb{N} \land \forall n \in \mathbb{N}\,(N \leq n \to \lvert \lvert a \rvert(n) - 0 \rvert < \varepsilon) \to N \in \mathbb{N} \land \forall n \in \mathbb{N}\,(N \leq n \to \lvert \lvert a \rvert(n) - 0 \rvert < \varepsilon)$
\item $\varphi_{83} :\equiv (\mathrm{Seq}(a) \to \mathrm{Conv}(a,0) \to \mathrm{Conv}(\lvert a \rvert,\lvert 0 \rvert)) \to \varphi_{3} \to \mathrm{Seq}(a) \to \mathrm{Conv}(a,0) \to \mathrm{Conv}(\lvert a \rvert,0)$
\item $\varphi_{84} :\equiv \forall \alpha\,(\mathrm{Conv}(a,\alpha) \leftrightarrow \mathrm{Seq}(a) \land \alpha \in \mathbb{R} \land \forall \varepsilon \in \mathbb{R}\,(0 < \varepsilon \to \exists N \in \mathbb{N}\,\forall n \in \mathbb{N}\,(N \leq n \to \lvert a_{n} - \alpha \rvert < \varepsilon)))$
\item $\varphi_{85} :\equiv N \in \mathbb{N} \land \forall n \in \mathbb{N}\,(N \leq n \to \lvert \lvert a \rvert(n) - 0 \rvert < \varepsilon) \to n \in \mathbb{N} \to N \leq n \to n \in \mathbb{N} \to N \leq n \to \lvert \lvert a \rvert(n) - 0 \rvert < \varepsilon$
\item $\varphi_{86} :\equiv 0 < \varepsilon \to N \in \mathbb{N} \land \forall n \in \mathbb{N}\,(N \leq n \to \lvert \lvert a \rvert(n) - 0 \rvert < \varepsilon) \to \varphi_{19}$
\item $\varphi_{87} :\equiv \forall w\,(a_{n} = w \to \lvert a_{n} \rvert = \lvert a_{n} \rvert \to \lvert a_{n} \rvert = \lvert w \rvert) \to a_{n} = a_{n} - 0 \to \lvert a_{n} \rvert = \lvert a_{n} \rvert \to \lvert a_{n} \rvert = \lvert a_{n} - 0 \rvert$
\item $\varphi_{88} :\equiv (\mathrm{Seq}(a) \to \mathrm{Conv}(\lvert a \rvert,0) \to \varphi_{53}) \to \mathrm{Seq}(a) \to \mathrm{Conv}(\lvert a \rvert,0) \to \mathrm{Conv}(a,0)$
\item $\varphi_{89} :\equiv (a_{n} = a_{n} - 0 \to \lvert a_{n} \rvert = \lvert a_{n} \rvert \to \lvert a_{n} \rvert = \lvert a_{n} - 0 \rvert) \to \varphi_{21}$
\item $\varphi_{90} :\equiv (\mathrm{Seq}(a) \to n \in \mathbb{N} \to a_{n} \in \mathbb{R}) \to (\mathrm{Seq}(a) \to \mathrm{Seq}(a)) \to \varphi_{60}$
\item $\varphi_{91} :\equiv (\lvert \lvert a \rvert(n) - 0 \rvert = \lvert \lvert a_{n} \rvert - 0 \rvert \to \varphi_{51}) \to \mathrm{Seq}(a) \to \varphi_{28}$
\item $\varphi_{92} :\equiv \forall \varepsilon\,(\mathrm{Conv}(\lvert a \rvert,0) \to \varepsilon \in \mathbb{R} \to \varphi_{38}) \to \mathrm{Conv}(\lvert a \rvert,0) \to \varepsilon \in \mathbb{R} \to \varphi_{38}$
\item $\varphi_{93} :\equiv (\mathrm{Seq}(a) \to \mathrm{Conv}(\lvert a \rvert,0) \to \varphi_{45}) \to \mathrm{Seq}(a) \land \mathrm{Conv}(\lvert a \rvert,0) \to \varphi_{45}$
\item $\varphi_{94} :\equiv (\mathrm{Seq}(a) \to \mathrm{Conv}(\lvert a \rvert,0) \to \varepsilon \in \mathbb{R} \to \varepsilon \in \mathbb{R}) \to \varphi_{12} \to \mathrm{Seq}(a) \to \mathrm{Conv}(\lvert a \rvert,0) \to \varepsilon \in \mathbb{R} \to \varphi_{38}$
\item $\varphi_{95} :\equiv \forall \varepsilon\,(\mathrm{Seq}(a) \land \mathrm{Conv}(\lvert a \rvert,0) \to \varphi_{45}) \to \forall \varepsilon\,(\mathrm{Seq}(a) \land \mathrm{Conv}(\lvert a \rvert,0)) \to \varphi_{53}$
\item $\varphi_{96} :\equiv \lvert \lvert a_{n} \rvert - 0 \rvert = \lvert \lvert a_{n} \rvert - 0 \rvert \to (\lvert a_{n} \rvert - 0 = \lvert a_{n} \rvert \to \varphi_{54}) \to \mathrm{Seq}(a) \to \varphi_{25}$
\item $\varphi_{97} :\equiv \varphi_{58} \to \mathrm{Seq}(a) \to \mathrm{Conv}(\lvert a \rvert,0) \to \varepsilon \in \mathbb{R} \to 0 < \varepsilon \to \varphi_{63}$
\item $\varphi_{98} :\equiv \lvert \lvert a \rvert(n) - 0 \rvert = \lvert \lvert a \rvert(n) - 0 \rvert \to (\lvert a \rvert(n) = \lvert a_{n} \rvert \to \varphi_{56}) \to \mathrm{Seq}(a) \to \varphi_{37}$
\item $\varphi_{99} :\equiv (\mathrm{Seq}(a) \to \mathrm{Conv}(\lvert a \rvert,0) \to \mathrm{Conv}(\lvert a \rvert,0)) \to \varphi_{94}$
\item $\varphi_{100} :\equiv \varphi_{57} \to \mathrm{Seq}(a) \to \mathrm{Conv}(\lvert a \rvert,0) \to \varepsilon \in \mathbb{R} \to 0 < \varepsilon \to \varphi_{85}$
\item $\varphi_{101} :\equiv (\varphi_{13} \to \neg (N \in \mathbb{N} \land \forall n \in \mathbb{N}\,(N \leq n \to \lvert a_{n} - 0 \rvert < \varepsilon))) \to \mathrm{Seq}(a) \to \mathrm{Conv}(\lvert a \rvert,0) \to \varepsilon \in \mathbb{R} \to 0 < \varepsilon \to \varphi_{80}$
\item $\varphi_{102} :\equiv (\mathrm{Seq}(a) \to \mathrm{Conv}(\lvert a \rvert,0) \to \varepsilon \in \mathbb{R} \to 0 < \varepsilon \to \varphi_{48}) \to \mathrm{Seq}(a) \to \mathrm{Conv}(\lvert a \rvert,0) \to \varepsilon \in \mathbb{R} \to 0 < \varepsilon \to \varphi_{69}$
\item $\varphi_{103} :\equiv (\mathrm{Seq}(a) \to \mathrm{Conv}(\lvert a \rvert,0) \to \varepsilon \in \mathbb{R} \to 0 < \varepsilon \to \varphi_{82}) \to \mathrm{Seq}(a) \to \mathrm{Conv}(\lvert a \rvert,0) \to \varepsilon \in \mathbb{R} \to 0 < \varepsilon \to \varphi_{36}$
\item $\varphi_{104} :\equiv (\mathrm{Conv}(a,0) \leftrightarrow \mathrm{Seq}(a) \land 0 \in \mathbb{R} \land \varphi_{53}) \to (\mathrm{Seq}(a) \to \mathrm{Seq}(a)) \to 0 \in \mathbb{R} \to \varphi_{88}$
\item $\varphi_{105} :\equiv (\mathrm{Seq}(a) \to \mathrm{Conv}(\lvert a \rvert,0) \to \varepsilon \in \mathbb{R} \to 0 < \varepsilon \to \varphi_{49}) \to \mathrm{Seq}(a) \to \mathrm{Conv}(\lvert a \rvert,0) \to \varepsilon \in \mathbb{R} \to 0 < \varepsilon \to \varphi_{79}$
\item $\varphi_{106} :\equiv (\mathrm{Seq}(a) \to \mathrm{Conv}(\lvert a \rvert,0) \to \varepsilon \in \mathbb{R} \to 0 < \varepsilon \to \varphi_{63}) \to \varphi_{70} \to n \in \mathbb{N} \to N \leq n \to \lvert a_{n} - 0 \rvert < \varepsilon$
\item $\varphi_{107} :\equiv \forall n\,(\varphi_{70} \to n \in \mathbb{N} \to N \leq n \to \lvert a_{n} - 0 \rvert < \varepsilon) \to \forall n\,(\varphi_{70}) \to \forall n \in \mathbb{N}\,(N \leq n \to \lvert a_{n} - 0 \rvert < \varepsilon)$
\item $\varphi_{108} :\equiv (\mathrm{Seq}(a) \to \mathrm{Conv}(\lvert a \rvert,0) \to \varepsilon \in \mathbb{R} \to 0 < \varepsilon \to \varphi_{65}) \to \mathrm{Seq}(a) \to \mathrm{Conv}(\lvert a \rvert,0) \to \varepsilon \in \mathbb{R} \to 0 < \varepsilon \to \varphi_{78}$
\item $\varphi_{109} :\equiv (\mathrm{Seq}(a) \to \mathrm{Conv}(\lvert a \rvert,0) \to \varepsilon \in \mathbb{R} \to 0 < \varepsilon \to \varphi_{80}) \to \varphi_{67} \to \varphi_{14}$
\item $\varphi_{110} :\equiv (\mathrm{Seq}(a) \to \mathrm{Conv}(\lvert a \rvert,0) \to \varepsilon \in \mathbb{R} \to 0 < \varepsilon \to \varphi_{82}) \to \mathrm{Seq}(a) \to \mathrm{Conv}(\lvert a \rvert,0) \to \varepsilon \in \mathbb{R} \to \varphi_{86}$
\item $\varphi_{111} :\equiv (\varphi_{70} \to \forall n\,(\varphi_{70})) \to \mathrm{Seq}(a) \to \mathrm{Conv}(\lvert a \rvert,0) \to \varepsilon \in \mathbb{R} \to 0 < \varepsilon \to \varphi_{65}$
\item $\varphi_{112} :\equiv (\mathrm{Seq}(a) \to \mathrm{Conv}(\lvert a \rvert,0) \to \varepsilon \in \mathbb{R} \to 0 < \varepsilon \to \varphi_{69}) \to \varphi_{97}$
\item $\varphi_{113} :\equiv (\mathrm{Seq}(a) \to \mathrm{Conv}(\lvert a \rvert,0) \to \varepsilon \in \mathbb{R} \to 0 < \varepsilon \to \varphi_{49}) \to \varphi_{102}$
\item $\varphi_{114} :\equiv (\mathrm{Seq}(a) \to \mathrm{Conv}(\lvert a \rvert,0) \to \varepsilon \in \mathbb{R} \to 0 < \varepsilon \to \varphi_{36}) \to \varphi_{108}$
\item $\varphi_{115} :\equiv (\mathrm{Seq}(a) \to \mathrm{Conv}(\lvert a \rvert,0) \to \varepsilon \in \mathbb{R} \to 0 < \varepsilon \to \varphi_{78}) \to \varphi_{101}$
\item $\varphi_{116} :\equiv \forall \varepsilon\,(\mathrm{Seq}(a) \land \mathrm{Conv}(\lvert a \rvert,0) \to \varphi_{45}) \to \varphi_{95} \to \varphi_{8} \to \mathrm{Seq}(a) \to \mathrm{Conv}(\lvert a \rvert,0) \to \varphi_{53}$
\item $\varphi_{117} :\equiv (\mathrm{Seq}(a) \to \mathrm{Conv}(\lvert a \rvert,0) \to \varepsilon \in \mathbb{R} \to 0 < \varepsilon \to \varphi_{79}) \to \varphi_{112}$
\item $\varphi_{118} :\equiv (\mathrm{Seq}(a) \to \mathrm{Conv}(\lvert a \rvert,0) \to \varepsilon \in \mathbb{R} \to 0 < \varepsilon \to \varphi_{85}) \to \varphi_{113}$
\item $\varphi_{119} :\equiv (\forall N\,(\varphi_{67} \to \varphi_{14}) \to \forall N\,(\varphi_{67}) \to \forall N\,(\varphi_{14})) \to (\varphi_{67} \to \forall N\,(\varphi_{67})) \to \mathrm{Seq}(a) \to \mathrm{Conv}(\lvert a \rvert,0) \to \varphi_{45}$
\item $\varphi_{120} :\equiv (\mathrm{Seq}(a) \to \mathrm{Conv}(\lvert a \rvert,0) \to \varepsilon \in \mathbb{R} \to \varphi_{38}) \to \forall N\,(\varphi_{67} \to \varphi_{14}) \to \varphi_{119}$
\end{abbrevlist}
\end{fullproof}

% 自動生成 (prover/build.js): LimSubConst  ノード数 408
\begin{fullproof}{定理2.15(3)}{$\forall a\,\forall \alpha\,(\mathrm{Seq}(a) \to \alpha \in \mathbb{R} \to (\mathrm{Conv}(a,\alpha) \leftrightarrow \mathrm{Conv}(a \oplus \overline{-\alpha},0)))$}{408}\phantomsection\label{tree:LimSubConst}
\noindent$\Pi_{1}$:
\begin{ptree}
\AXC{$\varphi_{136}$\lbl{A1}}
\AXC{$\varphi_{137}$\lbl{A1}}
\AXC{$\varphi_{88}$\lbl{A2}}
\AXC{$\varphi_{162}$\lbl{A1}}
\AXC{$\varphi_{162} \to \varphi_{164}$\lbl{A1}}
\BIC{$\varphi_{164}$}
\AXC{$\varphi_{164} \to \varphi_{88} \to \varphi_{163}$\lbl{A1}}
\BIC{$\varphi_{88} \to \varphi_{163}$}
\BIC{$\varphi_{163}$}
\AXC{$\varphi_{163} \to \varphi_{137} \to \varphi_{136} \to \varphi_{159}$\lbl{A1}}
\BIC{$\varphi_{137} \to \varphi_{136} \to \varphi_{159}$}
\BIC{$\varphi_{136} \to \varphi_{159}$}
\BIC{$\varphi_{159}$}
\end{ptree}
\noindent$\Pi_{2}$:
\begin{ptree}
\AXC{$y = z \to x = y \to x = z$\lbl{A6}}
\AXC{$(y = z \to x = y \to x = z) \to x = y \to y = z \to x = z$\lbl{A1}}
\BIC{$x = y \to y = z \to x = z$}
\UIC{$\forall z\,(x = y \to y = z \to x = z)$}
\UIC{$\forall y\,\forall z\,(x = y \to y = z \to x = z)$}
\UIC{$\forall x\,\forall y\,\forall z\,(x = y \to y = z \to x = z)$}
\AXC{$\forall x\,\forall y\,\forall z\,(x = y \to y = z \to x = z) \to \varphi_{47}$\lbl{A2}}
\BIC{$\varphi_{47}$}
\end{ptree}
\noindent$\Pi_{3}$:
\begin{ptree}
\AXC{$x = x$\lbl{A5}}
\AXC{$x = y \to x = x \to y = x$\lbl{A6}}
\AXC{$(x = y \to x = x \to y = x) \to x = x \to x = y \to y = x$\lbl{A1}}
\BIC{$x = x \to x = y \to y = x$}
\BIC{$x = y \to y = x$}
\UIC{$\forall y\,(x = y \to y = x)$}
\UIC{$\forall x\,\forall y\,(x = y \to y = x)$}
\end{ptree}
\noindent$\Pi_{4}$:
\begin{ptree}
\AXC{$\forall x\,\forall y\,(x - y = x + -y)$\lbl{Def$-$}}
\AXC{$\forall x\,\forall y\,(x - y = x + -y) \to \forall y\,(a_{n} - y = a_{n} + -y)$\lbl{A2}}
\BIC{$\forall y\,(a_{n} - y = a_{n} + -y)$}
\end{ptree}
\noindent$\Pi_{5}$:
\begin{ptree}
\AXC{$\stackrel{\Pi_{4}}{\vdots}$}\noLine
\UIC{$\forall y\,(a_{n} - y = a_{n} + -y)$}
\AXC{$\forall y\,(a_{n} - y = a_{n} + -y) \to a_{n} - \alpha = a_{n} + -\alpha$\lbl{A2}}
\BIC{$a_{n} - \alpha = a_{n} + -\alpha$}
\AXC{$a_{n} - \alpha = a_{n} + -\alpha \to \varphi_{64}$\lbl{A1}}
\BIC{$\varphi_{64}$}
\end{ptree}
\noindent$\Pi_{6}$:
\begin{ptree}
\AXC{$\stackrel{\Pi_{3}}{\vdots}$}\noLine
\UIC{$\forall x\,\forall y\,(x = y \to y = x)$}
\AXC{$\varphi_{2}$\lbl{A2}}
\BIC{$\forall y\,(a_{n} - \alpha = y \to y = a_{n} - \alpha)$}
\AXC{$\varphi_{68}$\lbl{A2}}
\BIC{$a_{n} - \alpha = a_{n} + -\alpha \to a_{n} + -\alpha = a_{n} - \alpha$}
\AXC{$\stackrel{\Pi_{5}}{\vdots}$}\noLine
\UIC{$\varphi_{64}$}
\BIC{$a_{n} + -\alpha = a_{n} - \alpha$}
\AXC{$a_{n} + -\alpha = a_{n} - \alpha \to \lvert a_{n} + -\alpha \rvert = \lvert a_{n} + -\alpha \rvert \to \varphi_{87} \to \varphi_{16}$\lbl{A1}}
\BIC{$\lvert a_{n} + -\alpha \rvert = \lvert a_{n} + -\alpha \rvert \to \varphi_{87} \to \varphi_{16}$}
\end{ptree}
\noindent$\Pi_{7}$:
\begin{ptree}
\AXC{$x = x$\lbl{A5}}
\UIC{$\forall x\,(x = x)$}
\AXC{$\forall x\,(x = x) \to \lvert a_{n} + -\alpha \rvert = \lvert a_{n} + -\alpha \rvert$\lbl{A2}}
\BIC{$\lvert a_{n} + -\alpha \rvert = \lvert a_{n} + -\alpha \rvert$}
\AXC{$\stackrel{\Pi_{6}}{\vdots}$}\noLine
\UIC{$\lvert a_{n} + -\alpha \rvert = \lvert a_{n} + -\alpha \rvert \to \varphi_{87} \to \varphi_{16}$}
\BIC{$\varphi_{87} \to \varphi_{16}$}
\end{ptree}
\noindent$\Pi_{8}$:
\begin{ptree}
\AXC{$z = w \to \lvert a_{n} + -\alpha \rvert = \lvert z \rvert \to \lvert a_{n} + -\alpha \rvert = \lvert w \rvert$\lbl{A6}}
\UIC{$\varphi_{10}$}
\UIC{$\forall z\,(\varphi_{10})$}
\AXC{$\forall z\,(\varphi_{10}) \to \varphi_{66}$\lbl{A2}}
\BIC{$\varphi_{66}$}
\AXC{$\varphi_{66} \to \varphi_{87}$\lbl{A2}}
\BIC{$\varphi_{87}$}
\AXC{$\stackrel{\Pi_{7}}{\vdots}$}\noLine
\UIC{$\varphi_{87} \to \varphi_{16}$}
\BIC{$\varphi_{16}$}
\end{ptree}
\noindent$\Pi_{9}$:
\begin{ptree}
\AXC{$y = z \to x = y \to x = z$\lbl{A6}}
\AXC{$(y = z \to x = y \to x = z) \to x = y \to y = z \to x = z$\lbl{A1}}
\BIC{$x = y \to y = z \to x = z$}
\UIC{$\forall z\,(x = y \to y = z \to x = z)$}
\UIC{$\forall y\,\forall z\,(x = y \to y = z \to x = z)$}
\UIC{$\forall x\,\forall y\,\forall z\,(x = y \to y = z \to x = z)$}
\AXC{$\forall x\,\forall y\,\forall z\,(x = y \to y = z \to x = z) \to \varphi_{47}$\lbl{A2}}
\BIC{$\varphi_{47}$}
\end{ptree}
\noindent$\Pi_{10}$:
\begin{ptree}
\AXC{$\stackrel{\Pi_{9}}{\vdots}$}\noLine
\UIC{$\varphi_{47}$}
\AXC{$\varphi_{47} \to \varphi_{113}$\lbl{A2}}
\BIC{$\varphi_{113}$}
\AXC{$\varphi_{113} \to \lvert a \oplus \overline{-\alpha}(n) - 0 \rvert = \lvert a_{n} + -\alpha - 0 \rvert \to \varphi_{76}$\lbl{A2}}
\BIC{$\lvert a \oplus \overline{-\alpha}(n) - 0 \rvert = \lvert a_{n} + -\alpha - 0 \rvert \to \varphi_{76}$}
\end{ptree}
\noindent$\Pi_{11}$:
\begin{ptree}
\AXC{$\forall x \in \mathbb{R}\,(-x \in \mathbb{R})$\lbl{(R1)}}
\AXC{$\forall x \in \mathbb{R}\,(-x \in \mathbb{R}) \to \alpha \in \mathbb{R} \to -\alpha \in \mathbb{R}$\lbl{A2}}
\BIC{$\alpha \in \mathbb{R} \to -\alpha \in \mathbb{R}$}
\AXC{$(\alpha \in \mathbb{R} \to -\alpha \in \mathbb{R}) \to \varphi_{30}$\lbl{A1}}
\BIC{$\varphi_{30}$}
\end{ptree}
\noindent$\Pi_{12}$:
\begin{ptree}
\AXC{$\mathrm{Seq}(a) \to \mathrm{Seq}(a)$\lbl{A1}}
\AXC{$\forall a\,\forall n\,(\mathrm{Seq}(a) \to n \in \mathbb{N} \to a_{n} \in \mathbb{R})$\lbl{[S1]}}
\AXC{$\varphi_{31}$\lbl{A2}}
\BIC{$\forall n\,(\mathrm{Seq}(a) \to n \in \mathbb{N} \to a_{n} \in \mathbb{R})$}
\AXC{$\varphi_{20}$\lbl{A2}}
\BIC{$\mathrm{Seq}(a) \to n \in \mathbb{N} \to a_{n} \in \mathbb{R}$}
\AXC{$\varphi_{118}$\lbl{A1}}
\BIC{$(\mathrm{Seq}(a) \to \mathrm{Seq}(a)) \to \varphi_{51}$}
\BIC{$\varphi_{51}$}
\end{ptree}
\noindent$\Pi_{13}$:
\begin{ptree}
\AXC{$\forall x\,\forall y\,(x \in \mathbb{R} \to y \in \mathbb{R} \to x + y \in \mathbb{R})$\lbl{(R1)}}
\AXC{$\varphi_{65}$\lbl{A2}}
\BIC{$\forall y\,(a_{n} \in \mathbb{R} \to y \in \mathbb{R} \to a_{n} + y \in \mathbb{R})$}
\AXC{$\varphi_{83}$\lbl{A2}}
\BIC{$a_{n} \in \mathbb{R} \to -\alpha \in \mathbb{R} \to a_{n} + -\alpha \in \mathbb{R}$}
\AXC{$\varphi_{161}$\lbl{A1}}
\BIC{$(\mathrm{Seq}(a) \to \alpha \in \mathbb{R} \to n \in \mathbb{N} \to a_{n} \in \mathbb{R}) \to \varphi_{73}$}
\end{ptree}
\noindent$\Pi_{14}$:
\begin{ptree}
\AXC{$\mathrm{Seq}(a) \to \alpha \in \mathbb{R} \to n \in \mathbb{N} \to n \in \mathbb{N}$\lbl{A1}}
\AXC{$\stackrel{\Pi_{12}}{\vdots}$}\noLine
\UIC{$\varphi_{51}$}
\BIC{$\mathrm{Seq}(a) \to \alpha \in \mathbb{R} \to n \in \mathbb{N} \to a_{n} \in \mathbb{R}$}
\AXC{$\stackrel{\Pi_{13}}{\vdots}$}\noLine
\UIC{$(\mathrm{Seq}(a) \to \alpha \in \mathbb{R} \to n \in \mathbb{N} \to a_{n} \in \mathbb{R}) \to \varphi_{73}$}
\BIC{$\varphi_{73}$}
\end{ptree}
\noindent$\Pi_{15}$:
\begin{ptree}
\AXC{$\forall x \in \mathbb{R}\,(x - 0 = x)$\lbl{[E28]}}
\AXC{$\varphi_{55}$\lbl{A2}}
\BIC{$a_{n} + -\alpha \in \mathbb{R} \to a_{n} + -\alpha - 0 = a_{n} + -\alpha$}
\AXC{$(a_{n} + -\alpha \in \mathbb{R} \to a_{n} + -\alpha - 0 = a_{n} + -\alpha) \to \varphi_{102}$\lbl{A1}}
\BIC{$\varphi_{102}$}
\end{ptree}
\noindent$\Pi_{16}$:
\begin{ptree}
\AXC{$\mathrm{Seq}(a) \to \alpha \in \mathbb{R} \to \alpha \in \mathbb{R}$\lbl{A1}}
\AXC{$\stackrel{\Pi_{11}}{\vdots}$}\noLine
\UIC{$\varphi_{30}$}
\BIC{$\mathrm{Seq}(a) \to \alpha \in \mathbb{R} \to n \in \mathbb{N} \to -\alpha \in \mathbb{R}$}
\AXC{$\stackrel{\Pi_{14}}{\vdots}$}\noLine
\UIC{$\varphi_{73}$}
\BIC{$\mathrm{Seq}(a) \to \alpha \in \mathbb{R} \to n \in \mathbb{N} \to a_{n} + -\alpha \in \mathbb{R}$}
\AXC{$\stackrel{\Pi_{15}}{\vdots}$}\noLine
\UIC{$\varphi_{102}$}
\BIC{$\mathrm{Seq}(a) \to \alpha \in \mathbb{R} \to n \in \mathbb{N} \to a_{n} + -\alpha - 0 = a_{n} + -\alpha$}
\AXC{$\varphi_{180}$\lbl{A1}}
\BIC{$\varphi_{173}$}
\end{ptree}
\noindent$\Pi_{17}$:
\begin{ptree}
\AXC{$x = x$\lbl{A5}}
\UIC{$\forall x\,(x = x)$}
\AXC{$\forall x\,(x = x) \to \lvert a_{n} + -\alpha - 0 \rvert = \lvert a_{n} + -\alpha - 0 \rvert$\lbl{A2}}
\BIC{$\lvert a_{n} + -\alpha - 0 \rvert = \lvert a_{n} + -\alpha - 0 \rvert$}
\AXC{$\stackrel{\Pi_{16}}{\vdots}$}\noLine
\UIC{$\varphi_{173}$}
\BIC{$(a_{n} + -\alpha - 0 = a_{n} + -\alpha \to \varphi_{72}) \to \varphi_{36}$}
\end{ptree}
\noindent$\Pi_{18}$:
\begin{ptree}
\AXC{$\varphi_{18}$\lbl{A6}}
\UIC{$\forall w\,(\varphi_{18})$}
\UIC{$\forall z\,\forall w\,(\varphi_{18})$}
\AXC{$\forall z\,\forall w\,(\varphi_{18}) \to \varphi_{86}$\lbl{A2}}
\BIC{$\varphi_{86}$}
\AXC{$\varphi_{86} \to a_{n} + -\alpha - 0 = a_{n} + -\alpha \to \varphi_{72}$\lbl{A2}}
\BIC{$a_{n} + -\alpha - 0 = a_{n} + -\alpha \to \varphi_{72}$}
\AXC{$\stackrel{\Pi_{17}}{\vdots}$}\noLine
\UIC{$(a_{n} + -\alpha - 0 = a_{n} + -\alpha \to \varphi_{72}) \to \varphi_{36}$}
\BIC{$\varphi_{36}$}
\end{ptree}
\noindent$\Pi_{19}$:
\begin{ptree}
\AXC{$y = z \to x = y \to x = z$\lbl{A6}}
\AXC{$(y = z \to x = y \to x = z) \to x = y \to y = z \to x = z$\lbl{A1}}
\BIC{$x = y \to y = z \to x = z$}
\UIC{$\forall z\,(x = y \to y = z \to x = z)$}
\UIC{$\forall y\,\forall z\,(x = y \to y = z \to x = z)$}
\UIC{$\forall x\,\forall y\,\forall z\,(x = y \to y = z \to x = z)$}
\AXC{$\forall x\,\forall y\,\forall z\,(x = y \to y = z \to x = z) \to \varphi_{47}$\lbl{A2}}
\BIC{$\varphi_{47}$}
\end{ptree}
\noindent$\Pi_{20}$:
\begin{ptree}
\AXC{$z = w \to \varphi_{89}$\lbl{A6}}
\UIC{$\forall w\,(z = w \to \varphi_{89})$}
\UIC{$\forall z\,\forall w\,(z = w \to \varphi_{89})$}
\AXC{$\forall z\,\forall w\,(z = w \to \varphi_{89}) \to \varphi_{133}$\lbl{A2}}
\BIC{$\varphi_{133}$}
\AXC{$\varphi_{133} \to \overline{-\alpha}(n) = -\alpha \to \varphi_{120}$\lbl{A2}}
\BIC{$\overline{-\alpha}(n) = -\alpha \to \varphi_{120}$}
\end{ptree}
\noindent$\Pi_{21}$:
\begin{ptree}
\AXC{$\forall x \in \mathbb{R}\,(-x \in \mathbb{R})$\lbl{(R1)}}
\AXC{$\forall x \in \mathbb{R}\,(-x \in \mathbb{R}) \to \alpha \in \mathbb{R} \to -\alpha \in \mathbb{R}$\lbl{A2}}
\BIC{$\alpha \in \mathbb{R} \to -\alpha \in \mathbb{R}$}
\AXC{$(\alpha \in \mathbb{R} \to -\alpha \in \mathbb{R}) \to \varphi_{13}$\lbl{A1}}
\BIC{$\varphi_{13}$}
\end{ptree}
\noindent$\Pi_{22}$:
\begin{ptree}
\AXC{$\forall c \in \mathbb{R}\,(\mathrm{Seq}(\overline{c}) \land \forall n \in \mathbb{N}\,(\overline{c}(n) = c))$\lbl{[S10]}}
\AXC{$\varphi_{110}$\lbl{A2}}
\BIC{$\varphi_{23}$}
\AXC{$\varphi_{23} \to (\mathrm{Seq}(a) \to \alpha \in \mathbb{R} \to -\alpha \in \mathbb{R}) \to \varphi_{29}$\lbl{A1}}
\BIC{$(\mathrm{Seq}(a) \to \alpha \in \mathbb{R} \to -\alpha \in \mathbb{R}) \to \varphi_{29}$}
\end{ptree}
\noindent$\Pi_{23}$:
\begin{ptree}
\AXC{$\mathrm{Seq}(a) \to \alpha \in \mathbb{R} \to \alpha \in \mathbb{R}$\lbl{A1}}
\AXC{$\stackrel{\Pi_{21}}{\vdots}$}\noLine
\UIC{$\varphi_{13}$}
\BIC{$\mathrm{Seq}(a) \to \alpha \in \mathbb{R} \to -\alpha \in \mathbb{R}$}
\AXC{$\stackrel{\Pi_{22}}{\vdots}$}\noLine
\UIC{$(\mathrm{Seq}(a) \to \alpha \in \mathbb{R} \to -\alpha \in \mathbb{R}) \to \varphi_{29}$}
\BIC{$\varphi_{29}$}
\end{ptree}
\noindent$\Pi_{24}$:
\begin{ptree}
\AXC{$\mathrm{Seq}(a) \to \alpha \in \mathbb{R} \to n \in \mathbb{N} \to n \in \mathbb{N}$\lbl{A1}}
\AXC{$\varphi_{28}$\lbl{A2}}
\AXC{$\stackrel{\Pi_{23}}{\vdots}$}\noLine
\UIC{$\varphi_{29}$}
\AXC{$\varphi_{29} \to \varphi_{15}$\lbl{A1}}
\BIC{$\varphi_{15}$}
\AXC{$\varphi_{15} \to \varphi_{28} \to \varphi_{12}$\lbl{A1}}
\BIC{$\varphi_{28} \to \varphi_{12}$}
\BIC{$\varphi_{12}$}
\AXC{$\varphi_{12} \to \varphi_{67}$\lbl{A1}}
\BIC{$\varphi_{67}$}
\BIC{$\mathrm{Seq}(a) \to \alpha \in \mathbb{R} \to n \in \mathbb{N} \to \overline{-\alpha}(n) = -\alpha$}
\end{ptree}
\noindent$\Pi_{25}$:
\begin{ptree}
\AXC{$\stackrel{\Pi_{24}}{\vdots}$}\noLine
\UIC{$\mathrm{Seq}(a) \to \alpha \in \mathbb{R} \to n \in \mathbb{N} \to \overline{-\alpha}(n) = -\alpha$}
\AXC{$(\mathrm{Seq}(a) \to \alpha \in \mathbb{R} \to n \in \mathbb{N} \to \overline{-\alpha}(n) = -\alpha) \to \varphi_{179}$\lbl{A1}}
\BIC{$\varphi_{179}$}
\end{ptree}
\noindent$\Pi_{26}$:
\begin{ptree}
\AXC{$\stackrel{\Pi_{20}}{\vdots}$}\noLine
\UIC{$\overline{-\alpha}(n) = -\alpha \to \varphi_{120}$}
\AXC{$x = x$\lbl{A5}}
\UIC{$\forall x\,(x = x)$}
\AXC{$\varphi_{41}$\lbl{A2}}
\BIC{$\lvert a_{n} + \overline{-\alpha}(n) - 0 \rvert = \lvert a_{n} + \overline{-\alpha}(n) - 0 \rvert$}
\AXC{$\stackrel{\Pi_{25}}{\vdots}$}\noLine
\UIC{$\varphi_{179}$}
\BIC{$(\overline{-\alpha}(n) = -\alpha \to \varphi_{120}) \to \mathrm{Seq}(a) \to \varphi_{50}$}
\BIC{$\mathrm{Seq}(a) \to \varphi_{50}$}
\end{ptree}
\noindent$\Pi_{27}$:
\begin{ptree}
\AXC{$\forall a\,(\varphi_{38})$\lbl{[S2]}}
\AXC{$\forall a\,(\varphi_{38}) \to \varphi_{38}$\lbl{A2}}
\BIC{$\varphi_{38}$}
\AXC{$\varphi_{38} \to \mathrm{Seq}(a) \to \mathrm{Seq}(\overline{-\alpha}) \to \varphi_{63}$\lbl{A2}}
\BIC{$\mathrm{Seq}(a) \to \mathrm{Seq}(\overline{-\alpha}) \to \varphi_{63}$}
\AXC{$\varphi_{174}$\lbl{A1}}
\BIC{$(\mathrm{Seq}(a) \to \mathrm{Seq}(a)) \to \varphi_{131}$}
\end{ptree}
\noindent$\Pi_{28}$:
\begin{ptree}
\AXC{$\stackrel{\Pi_{23}}{\vdots}$}\noLine
\UIC{$\varphi_{29}$}
\AXC{$\varphi_{29} \to \mathrm{Seq}(a) \to \alpha \in \mathbb{R} \to \mathrm{Seq}(\overline{-\alpha})$\lbl{A1}}
\BIC{$\mathrm{Seq}(a) \to \alpha \in \mathbb{R} \to \mathrm{Seq}(\overline{-\alpha})$}
\AXC{$\mathrm{Seq}(a) \to \mathrm{Seq}(a)$\lbl{A1}}
\AXC{$\stackrel{\Pi_{27}}{\vdots}$}\noLine
\UIC{$(\mathrm{Seq}(a) \to \mathrm{Seq}(a)) \to \varphi_{131}$}
\BIC{$\varphi_{131}$}
\BIC{$\mathrm{Seq}(a) \to \alpha \in \mathbb{R} \to \varphi_{63}$}
\end{ptree}
\noindent$\Pi_{29}$:
\begin{ptree}
\AXC{$\stackrel{\Pi_{28}}{\vdots}$}\noLine
\UIC{$\mathrm{Seq}(a) \to \alpha \in \mathbb{R} \to \varphi_{63}$}
\AXC{$(\mathrm{Seq}(a) \to \alpha \in \mathbb{R} \to \varphi_{63}) \to \mathrm{Seq}(a) \to \varphi_{58}$\lbl{A1}}
\BIC{$\mathrm{Seq}(a) \to \varphi_{58}$}
\AXC{$(\mathrm{Seq}(a) \to \varphi_{58}) \to \varphi_{115} \to \mathrm{Seq}(a) \to \varphi_{54}$\lbl{A1}}
\BIC{$\varphi_{115} \to \mathrm{Seq}(a) \to \varphi_{54}$}
\end{ptree}
\noindent$\Pi_{30}$:
\begin{ptree}
\AXC{$\mathrm{Seq}(a) \to \alpha \in \mathbb{R} \to n \in \mathbb{N} \to n \in \mathbb{N}$\lbl{A1}}
\AXC{$\varphi_{115}$\lbl{A2}}
\AXC{$\stackrel{\Pi_{29}}{\vdots}$}\noLine
\UIC{$\varphi_{115} \to \mathrm{Seq}(a) \to \varphi_{54}$}
\BIC{$\mathrm{Seq}(a) \to \varphi_{54}$}
\AXC{$\varphi_{165}$\lbl{A1}}
\BIC{$(\mathrm{Seq}(a) \to \alpha \in \mathbb{R} \to n \in \mathbb{N} \to n \in \mathbb{N}) \to \varphi_{48}$}
\BIC{$\varphi_{48}$}
\end{ptree}
\noindent$\Pi_{31}$:
\begin{ptree}
\AXC{$x = x$\lbl{A5}}
\UIC{$\forall x\,(x = x)$}
\AXC{$\varphi_{32}$\lbl{A2}}
\BIC{$\lvert a \oplus \overline{-\alpha}(n) - 0 \rvert = \lvert a \oplus \overline{-\alpha}(n) - 0 \rvert$}
\AXC{$\stackrel{\Pi_{30}}{\vdots}$}\noLine
\UIC{$\varphi_{48}$}
\AXC{$\varphi_{48} \to \lvert a \oplus \overline{-\alpha}(n) - 0 \rvert = \lvert a \oplus \overline{-\alpha}(n) - 0 \rvert \to \varphi_{176}$\lbl{A1}}
\BIC{$\lvert a \oplus \overline{-\alpha}(n) - 0 \rvert = \lvert a \oplus \overline{-\alpha}(n) - 0 \rvert \to \varphi_{176}$}
\BIC{$\varphi_{176}$}
\end{ptree}
\noindent$\Pi_{32}$:
\begin{ptree}
\AXC{$z = w \to \varphi_{60}$\lbl{A6}}
\UIC{$\forall w\,(z = w \to \varphi_{60})$}
\UIC{$\forall z\,\forall w\,(z = w \to \varphi_{60})$}
\AXC{$\forall z\,\forall w\,(z = w \to \varphi_{60}) \to \varphi_{129}$\lbl{A2}}
\BIC{$\varphi_{129}$}
\AXC{$\varphi_{129} \to a \oplus \overline{-\alpha}(n) = a_{n} + \overline{-\alpha}(n) \to \varphi_{126}$\lbl{A2}}
\BIC{$a \oplus \overline{-\alpha}(n) = a_{n} + \overline{-\alpha}(n) \to \varphi_{126}$}
\AXC{$\stackrel{\Pi_{31}}{\vdots}$}\noLine
\UIC{$\varphi_{176}$}
\BIC{$\mathrm{Seq}(a) \to \alpha \in \mathbb{R} \to \varphi_{45}$}
\end{ptree}
\noindent$\Pi_{33}$:
\begin{ptree}
\AXC{$\stackrel{\Pi_{26}}{\vdots}$}\noLine
\UIC{$\mathrm{Seq}(a) \to \varphi_{50}$}
\AXC{$\stackrel{\Pi_{32}}{\vdots}$}\noLine
\UIC{$\mathrm{Seq}(a) \to \alpha \in \mathbb{R} \to \varphi_{45}$}
\AXC{$(\mathrm{Seq}(a) \to \alpha \in \mathbb{R} \to \varphi_{45}) \to (\mathrm{Seq}(a) \to \varphi_{50}) \to \varphi_{175}$\lbl{A1}}
\BIC{$(\mathrm{Seq}(a) \to \varphi_{50}) \to \varphi_{175}$}
\BIC{$\varphi_{175}$}
\end{ptree}
\noindent$\Pi_{34}$:
\begin{ptree}
\AXC{$\stackrel{\Pi_{18}}{\vdots}$}\noLine
\UIC{$\varphi_{36}$}
\AXC{$\stackrel{\Pi_{19}}{\vdots}$}\noLine
\UIC{$\varphi_{47}$}
\AXC{$\varphi_{47} \to \varphi_{132}$\lbl{A2}}
\BIC{$\varphi_{132}$}
\AXC{$\varphi_{132} \to \lvert a \oplus \overline{-\alpha}(n) - 0 \rvert = \lvert a_{n} + \overline{-\alpha}(n) - 0 \rvert \to \varphi_{99}$\lbl{A2}}
\BIC{$\lvert a \oplus \overline{-\alpha}(n) - 0 \rvert = \lvert a_{n} + \overline{-\alpha}(n) - 0 \rvert \to \varphi_{99}$}
\AXC{$\stackrel{\Pi_{33}}{\vdots}$}\noLine
\UIC{$\varphi_{175}$}
\BIC{$\mathrm{Seq}(a) \to \varphi_{46}$}
\AXC{$\varphi_{186}$\lbl{A1}}
\BIC{$\varphi_{36} \to (\lvert a \oplus \overline{-\alpha}(n) - 0 \rvert = \lvert a_{n} + -\alpha - 0 \rvert \to \varphi_{76}) \to \varphi_{52}$}
\BIC{$(\lvert a \oplus \overline{-\alpha}(n) - 0 \rvert = \lvert a_{n} + -\alpha - 0 \rvert \to \varphi_{76}) \to \varphi_{52}$}
\end{ptree}
\noindent$\Pi_{35}$:
\begin{ptree}
\AXC{$\stackrel{\Pi_{10}}{\vdots}$}\noLine
\UIC{$\lvert a \oplus \overline{-\alpha}(n) - 0 \rvert = \lvert a_{n} + -\alpha - 0 \rvert \to \varphi_{76}$}
\AXC{$\stackrel{\Pi_{34}}{\vdots}$}\noLine
\UIC{$(\lvert a \oplus \overline{-\alpha}(n) - 0 \rvert = \lvert a_{n} + -\alpha - 0 \rvert \to \varphi_{76}) \to \varphi_{52}$}
\BIC{$\varphi_{52}$}
\end{ptree}
\noindent$\Pi_{36}$:
\begin{ptree}
\AXC{$\stackrel{\Pi_{8}}{\vdots}$}\noLine
\UIC{$\varphi_{16}$}
\AXC{$\stackrel{\Pi_{35}}{\vdots}$}\noLine
\UIC{$\varphi_{52}$}
\AXC{$\varphi_{52} \to \varphi_{16} \to (\lvert a \oplus \overline{-\alpha}(n) - 0 \rvert = \lvert a_{n} + -\alpha \rvert \to \varphi_{61}) \to \varphi_{42}$\lbl{A1}}
\BIC{$\varphi_{16} \to (\lvert a \oplus \overline{-\alpha}(n) - 0 \rvert = \lvert a_{n} + -\alpha \rvert \to \varphi_{61}) \to \varphi_{42}$}
\BIC{$(\lvert a \oplus \overline{-\alpha}(n) - 0 \rvert = \lvert a_{n} + -\alpha \rvert \to \varphi_{61}) \to \varphi_{42}$}
\end{ptree}
\noindent$\Pi_{37}$:
\begin{ptree}
\AXC{$\stackrel{\Pi_{2}}{\vdots}$}\noLine
\UIC{$\varphi_{47}$}
\AXC{$\varphi_{47} \to \varphi_{94}$\lbl{A2}}
\BIC{$\varphi_{94}$}
\AXC{$\varphi_{94} \to \lvert a \oplus \overline{-\alpha}(n) - 0 \rvert = \lvert a_{n} + -\alpha \rvert \to \varphi_{61}$\lbl{A2}}
\BIC{$\lvert a \oplus \overline{-\alpha}(n) - 0 \rvert = \lvert a_{n} + -\alpha \rvert \to \varphi_{61}$}
\AXC{$\stackrel{\Pi_{36}}{\vdots}$}\noLine
\UIC{$(\lvert a \oplus \overline{-\alpha}(n) - 0 \rvert = \lvert a_{n} + -\alpha \rvert \to \varphi_{61}) \to \varphi_{42}$}
\BIC{$\varphi_{42}$}
\end{ptree}
\noindent$\Pi_{38}$:
\begin{ptree}
\AXC{$\forall n\,(\varphi_{43}) \to \varphi_{53}$\lbl{A3}}
\AXC{$\stackrel{\Pi_{37}}{\vdots}$}\noLine
\UIC{$\varphi_{42}$}
\AXC{$\varphi_{42} \to \varphi_{43}$\lbl{A1}}
\BIC{$\varphi_{43}$}
\UIC{$\forall n\,(\varphi_{43})$}
\AXC{$\forall n\,(\varphi_{43}) \to (\forall n\,(\varphi_{43}) \to \varphi_{53}) \to \varphi_{107}$\lbl{A1}}
\BIC{$(\forall n\,(\varphi_{43}) \to \varphi_{53}) \to \varphi_{107}$}
\BIC{$\varphi_{107}$}
\end{ptree}
\noindent$\Pi_{39}$:
\begin{ptree}
\AXC{$\mathrm{Seq}(a) \land \alpha \in \mathbb{R} \to \forall n\,(\mathrm{Seq}(a) \land \alpha \in \mathbb{R})$\lbl{A4}}
\AXC{$\stackrel{\Pi_{38}}{\vdots}$}\noLine
\UIC{$\varphi_{107}$}
\BIC{$\mathrm{Seq}(a) \to \varphi_{34}$}
\end{ptree}
\noindent$\Pi_{40}$:
\begin{ptree}
\AXC{$\varphi_{137}$\lbl{A1}}
\AXC{$\varphi_{109}$\lbl{A2}}
\AXC{$\stackrel{\Pi_{39}}{\vdots}$}\noLine
\UIC{$\mathrm{Seq}(a) \to \varphi_{34}$}
\AXC{$(\mathrm{Seq}(a) \to \varphi_{34}) \to \varphi_{109} \to \varphi_{42}$\lbl{A1}}
\BIC{$\varphi_{109} \to \varphi_{42}$}
\BIC{$\varphi_{42}$}
\AXC{$\varphi_{42} \to \varphi_{137} \to \varphi_{160}$\lbl{A1}}
\BIC{$\varphi_{137} \to \varphi_{160}$}
\BIC{$\varphi_{160}$}
\AXC{$\varphi_{160} \to \varphi_{159} \to \varphi_{90} \to \varphi_{149}$\lbl{A1}}
\BIC{$\varphi_{159} \to \varphi_{90} \to \varphi_{149}$}
\end{ptree}
\noindent$\Pi_{41}$:
\begin{ptree}
\AXC{$z = w \to z < \varepsilon \to w < \varepsilon$\lbl{A6}}
\UIC{$\forall w\,(z = w \to z < \varepsilon \to w < \varepsilon)$}
\UIC{$\forall z\,\forall w\,(z = w \to z < \varepsilon \to w < \varepsilon)$}
\AXC{$\forall z\,\forall w\,(z = w \to z < \varepsilon \to w < \varepsilon) \to \varphi_{59}$\lbl{A2}}
\BIC{$\varphi_{59}$}
\AXC{$\varphi_{59} \to \varphi_{90}$\lbl{A2}}
\BIC{$\varphi_{90}$}
\AXC{$\stackrel{\Pi_{1}}{\vdots}$}\noLine
\UIC{$\varphi_{159}$}
\AXC{$\stackrel{\Pi_{40}}{\vdots}$}\noLine
\UIC{$\varphi_{159} \to \varphi_{90} \to \varphi_{149}$}
\BIC{$\varphi_{90} \to \varphi_{149}$}
\BIC{$\varphi_{149}$}
\end{ptree}
\noindent$\Pi_{42}$:
\begin{ptree}
\AXC{$\varphi_{121} \to \forall n\,(\varphi_{121})$\lbl{A4}}
\AXC{$\varphi_{183}$\lbl{A3}}
\AXC{$\stackrel{\Pi_{41}}{\vdots}$}\noLine
\UIC{$\varphi_{149}$}
\AXC{$\varphi_{149} \to \varphi_{121} \to n \in \mathbb{N} \to N \leq n \to \lvert a_{n} - \alpha \rvert < \varepsilon$\lbl{A1}}
\BIC{$\varphi_{121} \to n \in \mathbb{N} \to N \leq n \to \lvert a_{n} - \alpha \rvert < \varepsilon$}
\UIC{$\forall n\,(\varphi_{121} \to n \in \mathbb{N} \to N \leq n \to \lvert a_{n} - \alpha \rvert < \varepsilon)$}
\AXC{$\varphi_{192}$\lbl{A1}}
\BIC{$\varphi_{183} \to (\varphi_{121} \to \forall n\,(\varphi_{121})) \to \varphi_{150}$}
\BIC{$(\varphi_{121} \to \forall n\,(\varphi_{121})) \to \varphi_{150}$}
\BIC{$\varphi_{150}$}
\end{ptree}
\noindent$\Pi_{43}$:
\begin{ptree}
\AXC{$\stackrel{\Pi_{42}}{\vdots}$}\noLine
\UIC{$\varphi_{150}$}
\AXC{$\varphi_{162}$\lbl{A1}}
\AXC{$\varphi_{162} \to \varphi_{130}$\lbl{A1}}
\BIC{$\varphi_{130}$}
\AXC{$\varphi_{130} \to \varphi_{150} \to \varphi_{155}$\lbl{A1}}
\BIC{$\varphi_{150} \to \varphi_{155}$}
\BIC{$\varphi_{155}$}
\AXC{$\varphi_{155} \to \varphi_{177}$\lbl{A1}}
\BIC{$\varphi_{177}$}
\end{ptree}
\noindent$\Pi_{44}$:
\begin{ptree}
\AXC{$\varphi_{4} \to \neg (N \in \mathbb{N} \land \forall n \in \mathbb{N}\,(N \leq n \to \lvert a_{n} - \alpha \rvert < \varepsilon))$\lbl{A2}}
\AXC{$\stackrel{\Pi_{43}}{\vdots}$}\noLine
\UIC{$\varphi_{177}$}
\BIC{$\varphi_{156}$}
\AXC{$\varphi_{156} \to \varphi_{105} \to \varphi_{17}$\lbl{A1}}
\BIC{$\varphi_{105} \to \varphi_{17}$}
\UIC{$\forall N\,(\varphi_{105} \to \varphi_{17})$}
\end{ptree}
\noindent$\Pi_{45}$:
\begin{ptree}
\AXC{$\forall a\,\forall \alpha\,\forall \varepsilon\,(\mathrm{Conv}(a,\alpha) \to \varphi_{35})$\lbl{[C3]}}
\AXC{$\forall a\,\forall \alpha\,\forall \varepsilon\,(\mathrm{Conv}(a,\alpha) \to \varphi_{35}) \to \varphi_{70}$\lbl{A2}}
\BIC{$\varphi_{70}$}
\AXC{$\varphi_{70} \to \varphi_{71}$\lbl{A2}}
\BIC{$\varphi_{71}$}
\end{ptree}
\noindent$\Pi_{46}$:
\begin{ptree}
\AXC{$\stackrel{\Pi_{45}}{\vdots}$}\noLine
\UIC{$\varphi_{71}$}
\AXC{$\varphi_{71} \to \mathrm{Conv}(v_{700},v_{701}) \to v_{702} \in \mathbb{R} \to 0 < v_{702} \to \varphi_{8}$\lbl{A2}}
\BIC{$\mathrm{Conv}(v_{700},v_{701}) \to v_{702} \in \mathbb{R} \to 0 < v_{702} \to \varphi_{8}$}
\UIC{$\varphi_{74}$}
\UIC{$\forall v_{701}\,(\varphi_{74})$}
\UIC{$\forall v_{700}\,\forall v_{701}\,(\varphi_{74})$}
\AXC{$\forall v_{700}\,\forall v_{701}\,(\varphi_{74}) \to \varphi_{123}$\lbl{A2}}
\BIC{$\varphi_{123}$}
\AXC{$\varphi_{123} \to \varphi_{114}$\lbl{A2}}
\BIC{$\varphi_{114}$}
\end{ptree}
\noindent$\Pi_{47}$:
\begin{ptree}
\AXC{$\varphi_{24}$\lbl{A1}}
\AXC{$\mathrm{Seq}(a) \to \varphi_{14}$\lbl{A1}}
\AXC{$\stackrel{\Pi_{46}}{\vdots}$}\noLine
\UIC{$\varphi_{114}$}
\AXC{$\varphi_{114} \to \mathrm{Conv}(a \oplus \overline{-\alpha},0) \to \varepsilon \in \mathbb{R} \to \varphi_{40}$\lbl{A2}}
\BIC{$\mathrm{Conv}(a \oplus \overline{-\alpha},0) \to \varepsilon \in \mathbb{R} \to \varphi_{40}$}
\AXC{$\varphi_{189}$\lbl{A1}}
\BIC{$(\mathrm{Seq}(a) \to \varphi_{14}) \to \varphi_{24} \to (\mathrm{Seq}(a) \to \varphi_{19}) \to \varphi_{124}$}
\BIC{$\varphi_{24} \to (\mathrm{Seq}(a) \to \varphi_{19}) \to \varphi_{124}$}
\BIC{$(\mathrm{Seq}(a) \to \varphi_{19}) \to \varphi_{124}$}
\end{ptree}
\noindent$\Pi_{48}$:
\begin{ptree}
\AXC{$\stackrel{\Pi_{44}}{\vdots}$}\noLine
\UIC{$\forall N\,(\varphi_{105} \to \varphi_{17})$}
\AXC{$\mathrm{Seq}(a) \to \varphi_{19}$\lbl{A1}}
\AXC{$\stackrel{\Pi_{47}}{\vdots}$}\noLine
\UIC{$(\mathrm{Seq}(a) \to \varphi_{19}) \to \varphi_{124}$}
\BIC{$\varphi_{124}$}
\AXC{$\varphi_{124} \to \varphi_{191}$\lbl{A1}}
\BIC{$\varphi_{191}$}
\BIC{$(\forall N\,(\varphi_{105} \to \varphi_{17}) \to \forall N\,(\varphi_{105}) \to \forall N\,(\varphi_{17})) \to \varphi_{184}$}
\end{ptree}
\noindent$\Pi_{49}$:
\begin{ptree}
\AXC{$\varphi_{105} \to \forall N\,(\varphi_{105})$\lbl{A4}}
\AXC{$\forall N\,(\varphi_{105} \to \varphi_{17}) \to \forall N\,(\varphi_{105}) \to \forall N\,(\varphi_{17})$\lbl{A3}}
\AXC{$\stackrel{\Pi_{48}}{\vdots}$}\noLine
\UIC{$(\forall N\,(\varphi_{105} \to \varphi_{17}) \to \forall N\,(\varphi_{105}) \to \forall N\,(\varphi_{17})) \to \varphi_{184}$}
\BIC{$\varphi_{184}$}
\BIC{$\mathrm{Seq}(a) \to \alpha \in \mathbb{R} \to \mathrm{Conv}(a \oplus \overline{-\alpha},0) \to \varphi_{35}$}
\AXC{$\varphi_{171}$\lbl{A1}}
\BIC{$\mathrm{Seq}(a) \land \alpha \in \mathbb{R} \land \mathrm{Conv}(a \oplus \overline{-\alpha},0) \to \varphi_{35}$}
\UIC{$\varphi_{111}$}
\end{ptree}
\noindent$\Pi_{50}$:
\begin{ptree}
\AXC{$\forall a\,\forall \alpha\,(\mathrm{Conv}(a,\alpha) \leftrightarrow \mathrm{Seq}(a) \land \alpha \in \mathbb{R} \land \varphi_{39})$\lbl{DefConv}}
\AXC{$\varphi_{167}$\lbl{A2}}
\BIC{$\forall \alpha\,(\mathrm{Conv}(a,\alpha) \leftrightarrow \mathrm{Seq}(a) \land \alpha \in \mathbb{R} \land \varphi_{39})$}
\AXC{$\varphi_{166}$\lbl{A2}}
\BIC{$\mathrm{Conv}(a,\alpha) \leftrightarrow \mathrm{Seq}(a) \land \alpha \in \mathbb{R} \land \varphi_{39}$}
\end{ptree}
\noindent$\Pi_{51}$:
\begin{ptree}
\AXC{$\stackrel{\Pi_{50}}{\vdots}$}\noLine
\UIC{$\mathrm{Conv}(a,\alpha) \leftrightarrow \mathrm{Seq}(a) \land \alpha \in \mathbb{R} \land \varphi_{39}$}
\AXC{$(\mathrm{Conv}(a,\alpha) \leftrightarrow \mathrm{Seq}(a) \land \alpha \in \mathbb{R} \land \varphi_{39}) \to \varphi_{168}$\lbl{A1}}
\BIC{$\varphi_{168}$}
\end{ptree}
\noindent$\Pi_{52}$:
\begin{ptree}
\AXC{$\mathrm{Seq}(a) \to \alpha \in \mathbb{R} \to \alpha \in \mathbb{R}$\lbl{A1}}
\AXC{$\mathrm{Seq}(a) \to \mathrm{Seq}(a)$\lbl{A1}}
\AXC{$\stackrel{\Pi_{51}}{\vdots}$}\noLine
\UIC{$\varphi_{168}$}
\BIC{$(\mathrm{Seq}(a) \to \alpha \in \mathbb{R} \to \alpha \in \mathbb{R}) \to \varphi_{152}$}
\BIC{$\varphi_{152}$}
\end{ptree}
\noindent$\Pi_{53}$:
\begin{ptree}
\AXC{$z = w \to z < \varepsilon \to w < \varepsilon$\lbl{A6}}
\UIC{$\forall w\,(z = w \to z < \varepsilon \to w < \varepsilon)$}
\UIC{$\forall z\,\forall w\,(z = w \to z < \varepsilon \to w < \varepsilon)$}
\AXC{$\varphi_{49}$\lbl{A2}}
\BIC{$\forall w\,(\lvert a_{n} - \alpha \rvert = w \to \lvert a_{n} - \alpha \rvert < \varepsilon \to w < \varepsilon)$}
\AXC{$\forall w\,(\lvert a_{n} - \alpha \rvert = w \to \lvert a_{n} - \alpha \rvert < \varepsilon \to w < \varepsilon) \to \varphi_{91}$\lbl{A2}}
\BIC{$\varphi_{91}$}
\end{ptree}
\noindent$\Pi_{54}$:
\begin{ptree}
\AXC{$\varphi_{103}$\lbl{A1}}
\AXC{$\varphi_{104}$\lbl{A1}}
\AXC{$\varphi_{44}$\lbl{A2}}
\AXC{$\varphi_{134}$\lbl{A1}}
\AXC{$\varphi_{134} \to \varphi_{140}$\lbl{A1}}
\BIC{$\varphi_{140}$}
\AXC{$\varphi_{140} \to \varphi_{44} \to \varphi_{138}$\lbl{A1}}
\BIC{$\varphi_{44} \to \varphi_{138}$}
\BIC{$\varphi_{138}$}
\AXC{$\varphi_{138} \to \varphi_{104} \to \varphi_{103} \to \varphi_{128}$\lbl{A1}}
\BIC{$\varphi_{104} \to \varphi_{103} \to \varphi_{128}$}
\BIC{$\varphi_{103} \to \varphi_{128}$}
\BIC{$\varphi_{128}$}
\end{ptree}
\noindent$\Pi_{55}$:
\begin{ptree}
\AXC{$\varphi_{104}$\lbl{A1}}
\AXC{$\varphi_{109}$\lbl{A2}}
\AXC{$\stackrel{\Pi_{39}}{\vdots}$}\noLine
\UIC{$\mathrm{Seq}(a) \to \varphi_{34}$}
\AXC{$(\mathrm{Seq}(a) \to \varphi_{34}) \to \varphi_{109} \to \varphi_{42}$\lbl{A1}}
\BIC{$\varphi_{109} \to \varphi_{42}$}
\BIC{$\varphi_{42}$}
\AXC{$\varphi_{42} \to \varphi_{104} \to \varphi_{144}$\lbl{A1}}
\BIC{$\varphi_{104} \to \varphi_{144}$}
\BIC{$\varphi_{144}$}
\AXC{$\varphi_{144} \to \varphi_{75} \to \varphi_{145}$\lbl{A1}}
\BIC{$\varphi_{75} \to \varphi_{145}$}
\end{ptree}
\noindent$\Pi_{56}$:
\begin{ptree}
\AXC{$\stackrel{\Pi_{3}}{\vdots}$}\noLine
\UIC{$\forall x\,\forall y\,(x = y \to y = x)$}
\AXC{$\forall x\,\forall y\,(x = y \to y = x) \to \varphi_{33}$\lbl{A2}}
\BIC{$\varphi_{33}$}
\AXC{$\varphi_{33} \to \varphi_{75}$\lbl{A2}}
\BIC{$\varphi_{75}$}
\AXC{$\stackrel{\Pi_{55}}{\vdots}$}\noLine
\UIC{$\varphi_{75} \to \varphi_{145}$}
\BIC{$\varphi_{145}$}
\AXC{$\varphi_{145} \to \varphi_{128} \to \varphi_{91} \to \varphi_{139}$\lbl{A1}}
\BIC{$\varphi_{128} \to \varphi_{91} \to \varphi_{139}$}
\end{ptree}
\noindent$\Pi_{57}$:
\begin{ptree}
\AXC{$\varphi_{178}$\lbl{A3}}
\AXC{$\stackrel{\Pi_{53}}{\vdots}$}\noLine
\UIC{$\varphi_{91}$}
\AXC{$\stackrel{\Pi_{54}}{\vdots}$}\noLine
\UIC{$\varphi_{128}$}
\AXC{$\stackrel{\Pi_{56}}{\vdots}$}\noLine
\UIC{$\varphi_{128} \to \varphi_{91} \to \varphi_{139}$}
\BIC{$\varphi_{91} \to \varphi_{139}$}
\BIC{$\varphi_{139}$}
\AXC{$\varphi_{139} \to \varphi_{77} \to n \in \mathbb{N} \to N \leq n \to \lvert a \oplus \overline{-\alpha}(n) - 0 \rvert < \varepsilon$\lbl{A1}}
\BIC{$\varphi_{77} \to n \in \mathbb{N} \to N \leq n \to \lvert a \oplus \overline{-\alpha}(n) - 0 \rvert < \varepsilon$}
\UIC{$\forall n\,(\varphi_{77} \to n \in \mathbb{N} \to N \leq n \to \lvert a \oplus \overline{-\alpha}(n) - 0 \rvert < \varepsilon)$}
\AXC{$\varphi_{190}$\lbl{A1}}
\BIC{$\varphi_{178} \to (\varphi_{77} \to \forall n\,(\varphi_{77})) \to \varphi_{142}$}
\BIC{$(\varphi_{77} \to \forall n\,(\varphi_{77})) \to \varphi_{142}$}
\end{ptree}
\noindent$\Pi_{58}$:
\begin{ptree}
\AXC{$\varphi_{134}$\lbl{A1}}
\AXC{$\varphi_{134} \to \mathrm{Seq}(a) \to \alpha \in \mathbb{R} \to \mathrm{Conv}(a,\alpha) \to \varepsilon \in \mathbb{R} \to \varphi_{22}$\lbl{A1}}
\BIC{$\mathrm{Seq}(a) \to \alpha \in \mathbb{R} \to \mathrm{Conv}(a,\alpha) \to \varepsilon \in \mathbb{R} \to \varphi_{22}$}
\AXC{$\varphi_{187}$\lbl{A1}}
\BIC{$\varphi_{142} \to \varphi_{146}$}
\end{ptree}
\noindent$\Pi_{59}$:
\begin{ptree}
\AXC{$\forall N\,(\varphi_{17}) \to \varphi_{17}$\lbl{A2}}
\AXC{$\varphi_{77} \to \forall n\,(\varphi_{77})$\lbl{A4}}
\AXC{$\stackrel{\Pi_{57}}{\vdots}$}\noLine
\UIC{$(\varphi_{77} \to \forall n\,(\varphi_{77})) \to \varphi_{142}$}
\BIC{$\varphi_{142}$}
\AXC{$\stackrel{\Pi_{58}}{\vdots}$}\noLine
\UIC{$\varphi_{142} \to \varphi_{146}$}
\BIC{$\varphi_{146}$}
\AXC{$\varphi_{146} \to (\forall N\,(\varphi_{17}) \to \varphi_{17}) \to \varphi_{148}$\lbl{A1}}
\BIC{$(\forall N\,(\varphi_{17}) \to \varphi_{17}) \to \varphi_{148}$}
\BIC{$\varphi_{148}$}
\end{ptree}
\noindent$\Pi_{60}$:
\begin{ptree}
\AXC{$\forall a\,\forall \alpha\,\forall \varepsilon\,(\mathrm{Conv}(a,\alpha) \to \varphi_{35})$\lbl{[C3]}}
\AXC{$\varphi_{157}$\lbl{A2}}
\BIC{$\forall \alpha\,\forall \varepsilon\,(\mathrm{Conv}(a,\alpha) \to \varphi_{35})$}
\AXC{$\varphi_{154}$\lbl{A2}}
\BIC{$\forall \varepsilon\,(\mathrm{Conv}(a,\alpha) \to \varphi_{35})$}
\AXC{$\forall \varepsilon\,(\mathrm{Conv}(a,\alpha) \to \varphi_{35}) \to \mathrm{Conv}(a,\alpha) \to \varphi_{35}$\lbl{A2}}
\BIC{$\mathrm{Conv}(a,\alpha) \to \varphi_{35}$}
\end{ptree}
\noindent$\Pi_{61}$:
\begin{ptree}
\AXC{$\varphi_{9}$\lbl{A1}}
\AXC{$\varphi_{3}$\lbl{A1}}
\AXC{$\varphi_{1}$\lbl{A1}}
\AXC{$\stackrel{\Pi_{60}}{\vdots}$}\noLine
\UIC{$\mathrm{Conv}(a,\alpha) \to \varphi_{35}$}
\AXC{$\varphi_{181}$\lbl{A1}}
\BIC{$\varphi_{1} \to \varphi_{3} \to \varphi_{9} \to \mathrm{Seq}(a) \to \alpha \in \mathbb{R} \to \mathrm{Conv}(a,\alpha) \to \varphi_{35}$}
\BIC{$\varphi_{3} \to \varphi_{9} \to \mathrm{Seq}(a) \to \alpha \in \mathbb{R} \to \mathrm{Conv}(a,\alpha) \to \varphi_{35}$}
\BIC{$\varphi_{9} \to \mathrm{Seq}(a) \to \alpha \in \mathbb{R} \to \mathrm{Conv}(a,\alpha) \to \varphi_{35}$}
\BIC{$\mathrm{Seq}(a) \to \alpha \in \mathbb{R} \to \mathrm{Conv}(a,\alpha) \to \varphi_{35}$}
\AXC{$\varphi_{193}$\lbl{A1}}
\BIC{$\varphi_{151} \to (\varphi_{151} \to \forall N\,(\varphi_{106}) \to \varphi_{4}) \to \varphi_{185}$}
\end{ptree}
\noindent$\Pi_{62}$:
\begin{ptree}
\AXC{$\varphi_{151} \to \forall N\,(\varphi_{106}) \to \varphi_{4}$\lbl{A3}}
\AXC{$\stackrel{\Pi_{59}}{\vdots}$}\noLine
\UIC{$\varphi_{148}$}
\AXC{$\varphi_{182}$\lbl{A1}}
\BIC{$\varphi_{106} \to \neg (N \in \mathbb{N} \land \forall n \in \mathbb{N}\,(N \leq n \to \lvert a_{n} - \alpha \rvert < \varepsilon))$}
\UIC{$\varphi_{151}$}
\AXC{$\stackrel{\Pi_{61}}{\vdots}$}\noLine
\UIC{$\varphi_{151} \to (\varphi_{151} \to \forall N\,(\varphi_{106}) \to \varphi_{4}) \to \varphi_{185}$}
\BIC{$(\varphi_{151} \to \forall N\,(\varphi_{106}) \to \varphi_{4}) \to \varphi_{185}$}
\BIC{$\varphi_{185}$}
\end{ptree}
\noindent$\Pi_{63}$:
\begin{ptree}
\AXC{$\varphi_{11}$\lbl{A4}}
\AXC{$\varphi_{112} \to \varphi_{117}$\lbl{A3}}
\AXC{$\varphi_{106} \to \forall N\,(\varphi_{106})$\lbl{A4}}
\AXC{$\stackrel{\Pi_{62}}{\vdots}$}\noLine
\UIC{$\varphi_{185}$}
\BIC{$\mathrm{Seq}(a) \to \alpha \in \mathbb{R} \to \mathrm{Conv}(a,\alpha) \to \varepsilon \in \mathbb{R} \to \varphi_{40}$}
\AXC{$\varphi_{172}$\lbl{A1}}
\BIC{$\mathrm{Seq}(a) \land \alpha \in \mathbb{R} \land \mathrm{Conv}(a,\alpha) \to \varepsilon \in \mathbb{R} \to \varphi_{40}$}
\UIC{$\varphi_{112}$}
\AXC{$\varphi_{112} \to (\varphi_{112} \to \varphi_{117}) \to \varphi_{158}$\lbl{A1}}
\BIC{$(\varphi_{112} \to \varphi_{117}) \to \varphi_{158}$}
\BIC{$\varphi_{158}$}
\BIC{$\mathrm{Seq}(a) \to \alpha \in \mathbb{R} \to \mathrm{Conv}(a,\alpha) \to \forall \varepsilon \in \mathbb{R}\,(\varphi_{40})$}
\end{ptree}
\noindent$\Pi_{64}$:
\begin{ptree}
\AXC{$\forall a\,\forall \alpha\,(\mathrm{Conv}(a,\alpha) \leftrightarrow \mathrm{Seq}(a) \land \alpha \in \mathbb{R} \land \varphi_{39})$\lbl{DefConv}}
\AXC{$\varphi_{169}$\lbl{A2}}
\BIC{$\varphi_{97}$}
\AXC{$\varphi_{97} \to \varphi_{98}$\lbl{A2}}
\BIC{$\varphi_{98}$}
\UIC{$\forall v_{699}\,(\varphi_{98})$}
\UIC{$\forall v_{698}\,\forall v_{699}\,(\varphi_{98})$}
\AXC{$\forall v_{698}\,\forall v_{699}\,(\varphi_{98}) \to \varphi_{147}$\lbl{A2}}
\BIC{$\varphi_{147}$}
\end{ptree}
\noindent$\Pi_{65}$:
\begin{ptree}
\AXC{$0 \in \mathbb{R}$\lbl{(R1)}}
\AXC{$\stackrel{\Pi_{28}}{\vdots}$}\noLine
\UIC{$\mathrm{Seq}(a) \to \alpha \in \mathbb{R} \to \varphi_{63}$}
\AXC{$\varphi_{135}$\lbl{A1}}
\BIC{$\mathrm{Seq}(a) \to \alpha \in \mathbb{R} \to \mathrm{Seq}(a \oplus \overline{-\alpha})$}
\AXC{$\stackrel{\Pi_{64}}{\vdots}$}\noLine
\UIC{$\varphi_{147}$}
\AXC{$\varphi_{147} \to \varphi_{141}$\lbl{A2}}
\BIC{$\varphi_{141}$}
\AXC{$\varphi_{141} \to \varphi_{170}$\lbl{A1}}
\BIC{$\varphi_{170}$}
\BIC{$0 \in \mathbb{R} \to \varphi_{153}$}
\BIC{$\varphi_{153}$}
\end{ptree}
\noindent$\Pi_{66}$:
\begin{ptree}
\AXC{$\stackrel{\Pi_{63}}{\vdots}$}\noLine
\UIC{$\mathrm{Seq}(a) \to \alpha \in \mathbb{R} \to \mathrm{Conv}(a,\alpha) \to \forall \varepsilon \in \mathbb{R}\,(\varphi_{40})$}
\AXC{$\stackrel{\Pi_{65}}{\vdots}$}\noLine
\UIC{$\varphi_{153}$}
\BIC{$\varphi_{6}$}
\AXC{$\varphi_{6} \to \varphi_{5} \to \varphi_{7}$\lbl{A1}}
\BIC{$\varphi_{5} \to \varphi_{7}$}
\end{ptree}
\noindent 主木:
\begin{ptree}
\AXC{$\varphi_{69}$\lbl{A4}}
\AXC{$\varphi_{111} \to \varphi_{116}$\lbl{A3}}
\AXC{$\stackrel{\Pi_{49}}{\vdots}$}\noLine
\UIC{$\varphi_{111}$}
\AXC{$\varphi_{111} \to \varphi_{188}$\lbl{A1}}
\BIC{$\varphi_{188}$}
\BIC{$\varphi_{69} \to \mathrm{Seq}(a) \to \alpha \in \mathbb{R} \to \mathrm{Conv}(a \oplus \overline{-\alpha},0) \to \varphi_{39}$}
\BIC{$\mathrm{Seq}(a) \to \alpha \in \mathbb{R} \to \mathrm{Conv}(a \oplus \overline{-\alpha},0) \to \varphi_{39}$}
\AXC{$\stackrel{\Pi_{52}}{\vdots}$}\noLine
\UIC{$\varphi_{152}$}
\BIC{$\varphi_{5}$}
\AXC{$\stackrel{\Pi_{66}}{\vdots}$}\noLine
\UIC{$\varphi_{5} \to \varphi_{7}$}
\BIC{$\varphi_{7}$}
\UIC{$\forall \alpha\,(\varphi_{7})$}
\UIC{$\forall a\,\forall \alpha\,(\varphi_{7})$}
\end{ptree}
\begin{abbrevlist}
\item $\varphi_{1} :\equiv \mathrm{Seq}(a) \to \alpha \in \mathbb{R} \to \mathrm{Conv}(a,\alpha) \to \mathrm{Conv}(a,\alpha)$
\item $\varphi_{2} :\equiv \forall x\,\forall y\,(x = y \to y = x) \to \forall y\,(a_{n} - \alpha = y \to y = a_{n} - \alpha)$
\item $\varphi_{3} :\equiv \mathrm{Seq}(a) \to \alpha \in \mathbb{R} \to \mathrm{Conv}(a,\alpha) \to \varepsilon \in \mathbb{R} \to \varepsilon \in \mathbb{R}$
\item $\varphi_{4} :\equiv \forall N\,(\neg (N \in \mathbb{N} \land \forall n \in \mathbb{N}\,(N \leq n \to \lvert a_{n} - \alpha \rvert < \varepsilon)))$
\item $\varphi_{5} :\equiv \mathrm{Seq}(a) \to \alpha \in \mathbb{R} \to \mathrm{Conv}(a \oplus \overline{-\alpha},0) \to \mathrm{Conv}(a,\alpha)$
\item $\varphi_{6} :\equiv \mathrm{Seq}(a) \to \alpha \in \mathbb{R} \to \mathrm{Conv}(a,\alpha) \to \mathrm{Conv}(a \oplus \overline{-\alpha},0)$
\item $\varphi_{7} :\equiv \mathrm{Seq}(a) \to \alpha \in \mathbb{R} \to (\mathrm{Conv}(a,\alpha) \leftrightarrow \mathrm{Conv}(a \oplus \overline{-\alpha},0))$
\item $\varphi_{8} :\equiv \exists N \in \mathbb{N}\,\forall n \in \mathbb{N}\,(N \leq n \to \lvert v_{700}(n) - v_{701} \rvert < v_{702})$
\item $\varphi_{9} :\equiv \mathrm{Seq}(a) \to \alpha \in \mathbb{R} \to \mathrm{Conv}(a,\alpha) \to \varepsilon \in \mathbb{R} \to 0 < \varepsilon \to 0 < \varepsilon$
\item $\varphi_{10} :\equiv \forall w\,(z = w \to \lvert a_{n} + -\alpha \rvert = \lvert z \rvert \to \lvert a_{n} + -\alpha \rvert = \lvert w \rvert)$
\item $\varphi_{11} :\equiv \mathrm{Seq}(a) \land \alpha \in \mathbb{R} \land \mathrm{Conv}(a,\alpha) \to \forall \varepsilon\,(\mathrm{Seq}(a) \land \alpha \in \mathbb{R} \land \mathrm{Conv}(a,\alpha))$
\item $\varphi_{12} :\equiv \mathrm{Seq}(a) \to \alpha \in \mathbb{R} \to n \in \mathbb{N} \to n \in \mathbb{N} \to \overline{-\alpha}(n) = -\alpha$
\item $\varphi_{13} :\equiv (\mathrm{Seq}(a) \to \alpha \in \mathbb{R} \to \alpha \in \mathbb{R}) \to \mathrm{Seq}(a) \to \alpha \in \mathbb{R} \to -\alpha \in \mathbb{R}$
\item $\varphi_{14} :\equiv \alpha \in \mathbb{R} \to \mathrm{Conv}(a \oplus \overline{-\alpha},0) \to \mathrm{Conv}(a \oplus \overline{-\alpha},0)$
\item $\varphi_{15} :\equiv \mathrm{Seq}(a) \to \alpha \in \mathbb{R} \to n \in \mathbb{N} \to \forall n \in \mathbb{N}\,(\overline{-\alpha}(n) = -\alpha)$
\item $\varphi_{16} :\equiv \mathrm{Seq}(a) \to \alpha \in \mathbb{R} \to n \in \mathbb{N} \to \lvert a_{n} + -\alpha \rvert = \lvert a_{n} - \alpha \rvert$
\item $\varphi_{17} :\equiv \neg (N \in \mathbb{N} \land \forall n \in \mathbb{N}\,(N \leq n \to \lvert a \oplus \overline{-\alpha}(n) - 0 \rvert < \varepsilon))$
\item $\varphi_{18} :\equiv z = w \to \lvert a_{n} + -\alpha - 0 \rvert = \lvert z \rvert \to \lvert a_{n} + -\alpha - 0 \rvert = \lvert w \rvert$
\item $\varphi_{19} :\equiv \alpha \in \mathbb{R} \to \mathrm{Conv}(a \oplus \overline{-\alpha},0) \to \varepsilon \in \mathbb{R} \to 0 < \varepsilon \to 0 < \varepsilon$
\item $\varphi_{20} :\equiv \forall n\,(\mathrm{Seq}(a) \to n \in \mathbb{N} \to a_{n} \in \mathbb{R}) \to \mathrm{Seq}(a) \to n \in \mathbb{N} \to a_{n} \in \mathbb{R}$
\item $\varphi_{21} :\equiv n \in \mathbb{N} \to N \leq n \to n \in \mathbb{N} \to N \leq n \to \lvert a \oplus \overline{-\alpha}(n) - 0 \rvert < \varepsilon$
\item $\varphi_{22} :\equiv 0 < \varepsilon \to N \in \mathbb{N} \land \forall n \in \mathbb{N}\,(N \leq n \to \lvert a_{n} - \alpha \rvert < \varepsilon) \to N \in \mathbb{N}$
\item $\varphi_{23} :\equiv -\alpha \in \mathbb{R} \to \mathrm{Seq}(\overline{-\alpha}) \land \forall n \in \mathbb{N}\,(\overline{-\alpha}(n) = -\alpha)$
\item $\varphi_{24} :\equiv \mathrm{Seq}(a) \to \alpha \in \mathbb{R} \to \mathrm{Conv}(a \oplus \overline{-\alpha},0) \to \varepsilon \in \mathbb{R} \to \varepsilon \in \mathbb{R}$
\item $\varphi_{25} :\equiv N \in \mathbb{N} \land \forall n \in \mathbb{N}\,(N \leq n \to \lvert a_{n} - \alpha \rvert < \varepsilon) \to n \in \mathbb{N} \to N \leq n \to N \leq n$
\item $\varphi_{26} :\equiv N \in \mathbb{N} \land \forall n \in \mathbb{N}\,(N \leq n \to \lvert a_{n} - \alpha \rvert < \varepsilon) \to n \in \mathbb{N} \to n \in \mathbb{N}$
\item $\varphi_{27} :\equiv n \in \mathbb{N} \to N \leq n \to \forall n \in \mathbb{N}\,(N \leq n \to \lvert a \oplus \overline{-\alpha}(n) - 0 \rvert < \varepsilon)$
\item $\varphi_{28} :\equiv \forall n \in \mathbb{N}\,(\overline{-\alpha}(n) = -\alpha) \to n \in \mathbb{N} \to \overline{-\alpha}(n) = -\alpha$
\item $\varphi_{29} :\equiv \mathrm{Seq}(a) \to \alpha \in \mathbb{R} \to \mathrm{Seq}(\overline{-\alpha}) \land \forall n \in \mathbb{N}\,(\overline{-\alpha}(n) = -\alpha)$
\item $\varphi_{30} :\equiv (\mathrm{Seq}(a) \to \alpha \in \mathbb{R} \to \alpha \in \mathbb{R}) \to \mathrm{Seq}(a) \to \alpha \in \mathbb{R} \to n \in \mathbb{N} \to -\alpha \in \mathbb{R}$
\item $\varphi_{31} :\equiv \forall a\,\forall n\,(\mathrm{Seq}(a) \to n \in \mathbb{N} \to a_{n} \in \mathbb{R}) \to \forall n\,(\mathrm{Seq}(a) \to n \in \mathbb{N} \to a_{n} \in \mathbb{R})$
\item $\varphi_{32} :\equiv \forall x\,(x = x) \to \lvert a \oplus \overline{-\alpha}(n) - 0 \rvert = \lvert a \oplus \overline{-\alpha}(n) - 0 \rvert$
\item $\varphi_{33} :\equiv \forall y\,(\lvert a \oplus \overline{-\alpha}(n) - 0 \rvert = y \to y = \lvert a \oplus \overline{-\alpha}(n) - 0 \rvert)$
\item $\varphi_{34} :\equiv \alpha \in \mathbb{R} \to \forall n \in \mathbb{N}\,(\lvert a \oplus \overline{-\alpha}(n) - 0 \rvert = \lvert a_{n} - \alpha \rvert)$
\item $\varphi_{35} :\equiv \varepsilon \in \mathbb{R} \to 0 < \varepsilon \to \exists N \in \mathbb{N}\,\forall n \in \mathbb{N}\,(N \leq n \to \lvert a_{n} - \alpha \rvert < \varepsilon)$
\item $\varphi_{36} :\equiv \mathrm{Seq}(a) \to \alpha \in \mathbb{R} \to n \in \mathbb{N} \to \lvert a_{n} + -\alpha - 0 \rvert = \lvert a_{n} + -\alpha \rvert$
\item $\varphi_{37} :\equiv N \in \mathbb{N} \land \forall n \in \mathbb{N}\,(N \leq n \to \lvert a \oplus \overline{-\alpha}(n) - 0 \rvert < \varepsilon) \to N \in \mathbb{N}$
\item $\varphi_{38} :\equiv \forall b\,(\mathrm{Seq}(a) \to \mathrm{Seq}(b) \to \mathrm{Seq}(a \oplus b) \land \forall n \in \mathbb{N}\,(a \oplus b(n) = a_{n} + b_{n}))$
\item $\varphi_{39} :\equiv \forall \varepsilon \in \mathbb{R}\,(0 < \varepsilon \to \exists N \in \mathbb{N}\,\forall n \in \mathbb{N}\,(N \leq n \to \lvert a_{n} - \alpha \rvert < \varepsilon))$
\item $\varphi_{40} :\equiv 0 < \varepsilon \to \exists N \in \mathbb{N}\,\forall n \in \mathbb{N}\,(N \leq n \to \lvert a \oplus \overline{-\alpha}(n) - 0 \rvert < \varepsilon)$
\item $\varphi_{41} :\equiv \forall x\,(x = x) \to \lvert a_{n} + \overline{-\alpha}(n) - 0 \rvert = \lvert a_{n} + \overline{-\alpha}(n) - 0 \rvert$
\item $\varphi_{42} :\equiv \mathrm{Seq}(a) \to \alpha \in \mathbb{R} \to n \in \mathbb{N} \to \lvert a \oplus \overline{-\alpha}(n) - 0 \rvert = \lvert a_{n} - \alpha \rvert$
\item $\varphi_{43} :\equiv \mathrm{Seq}(a) \land \alpha \in \mathbb{R} \to n \in \mathbb{N} \to \lvert a \oplus \overline{-\alpha}(n) - 0 \rvert = \lvert a_{n} - \alpha \rvert$
\item $\varphi_{44} :\equiv \forall n \in \mathbb{N}\,(N \leq n \to \lvert a_{n} - \alpha \rvert < \varepsilon) \to n \in \mathbb{N} \to N \leq n \to \lvert a_{n} - \alpha \rvert < \varepsilon$
\item $\varphi_{45} :\equiv n \in \mathbb{N} \to \lvert a \oplus \overline{-\alpha}(n) - 0 \rvert = \lvert a_{n} + \overline{-\alpha}(n) - 0 \rvert$
\item $\varphi_{46} :\equiv \alpha \in \mathbb{R} \to n \in \mathbb{N} \to \lvert a \oplus \overline{-\alpha}(n) - 0 \rvert = \lvert a_{n} + -\alpha - 0 \rvert$
\item $\varphi_{47} :\equiv \forall y\,\forall z\,(\lvert a \oplus \overline{-\alpha}(n) - 0 \rvert = y \to y = z \to \lvert a \oplus \overline{-\alpha}(n) - 0 \rvert = z)$
\item $\varphi_{48} :\equiv \mathrm{Seq}(a) \to \alpha \in \mathbb{R} \to n \in \mathbb{N} \to a \oplus \overline{-\alpha}(n) = a_{n} + \overline{-\alpha}(n)$
\item $\varphi_{49} :\equiv \forall z\,\forall w\,(z = w \to z < \varepsilon \to w < \varepsilon) \to \forall w\,(\lvert a_{n} - \alpha \rvert = w \to \lvert a_{n} - \alpha \rvert < \varepsilon \to w < \varepsilon)$
\item $\varphi_{50} :\equiv \alpha \in \mathbb{R} \to n \in \mathbb{N} \to \lvert a_{n} + \overline{-\alpha}(n) - 0 \rvert = \lvert a_{n} + -\alpha - 0 \rvert$
\item $\varphi_{51} :\equiv (\mathrm{Seq}(a) \to \alpha \in \mathbb{R} \to n \in \mathbb{N} \to n \in \mathbb{N}) \to \mathrm{Seq}(a) \to \alpha \in \mathbb{R} \to n \in \mathbb{N} \to a_{n} \in \mathbb{R}$
\item $\varphi_{52} :\equiv \mathrm{Seq}(a) \to \alpha \in \mathbb{R} \to n \in \mathbb{N} \to \lvert a \oplus \overline{-\alpha}(n) - 0 \rvert = \lvert a_{n} + -\alpha \rvert$
\item $\varphi_{53} :\equiv \forall n\,(\mathrm{Seq}(a) \land \alpha \in \mathbb{R}) \to \forall n \in \mathbb{N}\,(\lvert a \oplus \overline{-\alpha}(n) - 0 \rvert = \lvert a_{n} - \alpha \rvert)$
\item $\varphi_{54} :\equiv \alpha \in \mathbb{R} \to n \in \mathbb{N} \to n \in \mathbb{N} \to a \oplus \overline{-\alpha}(n) = a_{n} + \overline{-\alpha}(n)$
\item $\varphi_{55} :\equiv \forall x \in \mathbb{R}\,(x - 0 = x) \to a_{n} + -\alpha \in \mathbb{R} \to a_{n} + -\alpha - 0 = a_{n} + -\alpha$
\item $\varphi_{56} :\equiv N \in \mathbb{N} \land \forall n \in \mathbb{N}\,(N \leq n \to \lvert a \oplus \overline{-\alpha}(n) - 0 \rvert < \varepsilon) \to n \in \mathbb{N} \to N \leq n \to N \leq n$
\item $\varphi_{57} :\equiv N \in \mathbb{N} \land \forall n \in \mathbb{N}\,(N \leq n \to \lvert a \oplus \overline{-\alpha}(n) - 0 \rvert < \varepsilon) \to n \in \mathbb{N} \to n \in \mathbb{N}$
\item $\varphi_{58} :\equiv \alpha \in \mathbb{R} \to n \in \mathbb{N} \to \forall n \in \mathbb{N}\,(a \oplus \overline{-\alpha}(n) = a_{n} + \overline{-\alpha}(n))$
\item $\varphi_{59} :\equiv \forall w\,(\lvert a \oplus \overline{-\alpha}(n) - 0 \rvert = w \to \lvert a \oplus \overline{-\alpha}(n) - 0 \rvert < \varepsilon \to w < \varepsilon)$
\item $\varphi_{60} :\equiv \lvert a \oplus \overline{-\alpha}(n) - 0 \rvert = \lvert z - 0 \rvert \to \lvert a \oplus \overline{-\alpha}(n) - 0 \rvert = \lvert w - 0 \rvert$
\item $\varphi_{61} :\equiv \lvert a_{n} + -\alpha \rvert = \lvert a_{n} - \alpha \rvert \to \lvert a \oplus \overline{-\alpha}(n) - 0 \rvert = \lvert a_{n} - \alpha \rvert$
\item $\varphi_{62} :\equiv N \in \mathbb{N} \land \forall n \in \mathbb{N}\,(N \leq n \to \lvert a_{n} - \alpha \rvert < \varepsilon) \to n \in \mathbb{N} \to N \leq n \to \lvert a_{n} - \alpha \rvert < \varepsilon$
\item $\varphi_{63} :\equiv \mathrm{Seq}(a \oplus \overline{-\alpha}) \land \forall n \in \mathbb{N}\,(a \oplus \overline{-\alpha}(n) = a_{n} + \overline{-\alpha}(n))$
\item $\varphi_{64} :\equiv (a_{n} - \alpha = a_{n} + -\alpha \to a_{n} + -\alpha = a_{n} - \alpha) \to a_{n} + -\alpha = a_{n} - \alpha$
\item $\varphi_{65} :\equiv \forall x\,\forall y\,(x \in \mathbb{R} \to y \in \mathbb{R} \to x + y \in \mathbb{R}) \to \forall y\,(a_{n} \in \mathbb{R} \to y \in \mathbb{R} \to a_{n} + y \in \mathbb{R})$
\item $\varphi_{66} :\equiv \forall w\,(a_{n} + -\alpha = w \to \lvert a_{n} + -\alpha \rvert = \lvert a_{n} + -\alpha \rvert \to \lvert a_{n} + -\alpha \rvert = \lvert w \rvert)$
\item $\varphi_{67} :\equiv (\mathrm{Seq}(a) \to \alpha \in \mathbb{R} \to n \in \mathbb{N} \to n \in \mathbb{N}) \to \mathrm{Seq}(a) \to \alpha \in \mathbb{R} \to n \in \mathbb{N} \to \overline{-\alpha}(n) = -\alpha$
\item $\varphi_{68} :\equiv \forall y\,(a_{n} - \alpha = y \to y = a_{n} - \alpha) \to a_{n} - \alpha = a_{n} + -\alpha \to a_{n} + -\alpha = a_{n} - \alpha$
\item $\varphi_{69} :\equiv \mathrm{Seq}(a) \land \alpha \in \mathbb{R} \land \mathrm{Conv}(a \oplus \overline{-\alpha},0) \to \forall \varepsilon\,(\mathrm{Seq}(a) \land \alpha \in \mathbb{R} \land \mathrm{Conv}(a \oplus \overline{-\alpha},0))$
\item $\varphi_{70} :\equiv \forall \alpha\,\forall \varepsilon\,(\mathrm{Conv}(v_{700},\alpha) \to \varepsilon \in \mathbb{R} \to 0 < \varepsilon \to \exists N \in \mathbb{N}\,\forall n \in \mathbb{N}\,(N \leq n \to \lvert v_{700}(n) - \alpha \rvert < \varepsilon))$
\item $\varphi_{71} :\equiv \forall \varepsilon\,(\mathrm{Conv}(v_{700},v_{701}) \to \varepsilon \in \mathbb{R} \to 0 < \varepsilon \to \exists N \in \mathbb{N}\,\forall n \in \mathbb{N}\,(N \leq n \to \lvert v_{700}(n) - v_{701} \rvert < \varepsilon))$
\item $\varphi_{72} :\equiv \lvert a_{n} + -\alpha - 0 \rvert = \lvert a_{n} + -\alpha - 0 \rvert \to \lvert a_{n} + -\alpha - 0 \rvert = \lvert a_{n} + -\alpha \rvert$
\item $\varphi_{73} :\equiv (\mathrm{Seq}(a) \to \alpha \in \mathbb{R} \to n \in \mathbb{N} \to -\alpha \in \mathbb{R}) \to \mathrm{Seq}(a) \to \alpha \in \mathbb{R} \to n \in \mathbb{N} \to a_{n} + -\alpha \in \mathbb{R}$
\item $\varphi_{74} :\equiv \forall v_{702}\,(\mathrm{Conv}(v_{700},v_{701}) \to v_{702} \in \mathbb{R} \to 0 < v_{702} \to \varphi_{8})$
\item $\varphi_{75} :\equiv \lvert a \oplus \overline{-\alpha}(n) - 0 \rvert = \lvert a_{n} - \alpha \rvert \to \lvert a_{n} - \alpha \rvert = \lvert a \oplus \overline{-\alpha}(n) - 0 \rvert$
\item $\varphi_{76} :\equiv \lvert a_{n} + -\alpha - 0 \rvert = \lvert a_{n} + -\alpha \rvert \to \lvert a \oplus \overline{-\alpha}(n) - 0 \rvert = \lvert a_{n} + -\alpha \rvert$
\item $\varphi_{77} :\equiv \mathrm{Seq}(a) \land \alpha \in \mathbb{R} \land \mathrm{Conv}(a,\alpha) \land \varepsilon \in \mathbb{R} \land 0 < \varepsilon \land (N \in \mathbb{N} \land \forall n \in \mathbb{N}\,(N \leq n \to \lvert a_{n} - \alpha \rvert < \varepsilon))$
\item $\varphi_{78} :\equiv N \in \mathbb{N} \land \forall n \in \mathbb{N}\,(N \leq n \to \lvert a_{n} - \alpha \rvert < \varepsilon) \to N \in \mathbb{N} \land \forall n \in \mathbb{N}\,(N \leq n \to \lvert a_{n} - \alpha \rvert < \varepsilon)$
\item $\varphi_{79} :\equiv N \in \mathbb{N} \land \forall n \in \mathbb{N}\,(N \leq n \to \lvert a_{n} - \alpha \rvert < \varepsilon) \to n \in \mathbb{N} \to N \leq n \to n \in \mathbb{N} \to N \leq n \to \lvert a_{n} - \alpha \rvert < \varepsilon$
\item $\varphi_{80} :\equiv N \in \mathbb{N} \land \forall n \in \mathbb{N}\,(N \leq n \to \lvert a \oplus \overline{-\alpha}(n) - 0 \rvert < \varepsilon) \to n \in \mathbb{N} \to N \leq n \to \lvert a_{n} - \alpha \rvert < \varepsilon$
\item $\varphi_{81} :\equiv N \in \mathbb{N} \land \forall n \in \mathbb{N}\,(N \leq n \to \lvert a_{n} - \alpha \rvert < \varepsilon) \to n \in \mathbb{N} \to N \leq n \to \lvert a \oplus \overline{-\alpha}(n) - 0 \rvert < \varepsilon$
\item $\varphi_{82} :\equiv N \in \mathbb{N} \land \forall n \in \mathbb{N}\,(N \leq n \to \lvert a_{n} - \alpha \rvert < \varepsilon) \to n \in \mathbb{N} \to N \leq n \to \forall n \in \mathbb{N}\,(N \leq n \to \lvert a_{n} - \alpha \rvert < \varepsilon)$
\item $\varphi_{83} :\equiv \forall y\,(a_{n} \in \mathbb{R} \to y \in \mathbb{R} \to a_{n} + y \in \mathbb{R}) \to a_{n} \in \mathbb{R} \to -\alpha \in \mathbb{R} \to a_{n} + -\alpha \in \mathbb{R}$
\item $\varphi_{84} :\equiv N \in \mathbb{N} \land \forall n \in \mathbb{N}\,(N \leq n \to \lvert a \oplus \overline{-\alpha}(n) - 0 \rvert < \varepsilon) \to \forall n \in \mathbb{N}\,(N \leq n \to \lvert a_{n} - \alpha \rvert < \varepsilon)$
\item $\varphi_{85} :\equiv N \in \mathbb{N} \land \forall n \in \mathbb{N}\,(N \leq n \to \lvert a_{n} - \alpha \rvert < \varepsilon) \to \forall n \in \mathbb{N}\,(N \leq n \to \lvert a \oplus \overline{-\alpha}(n) - 0 \rvert < \varepsilon)$
\item $\varphi_{86} :\equiv \forall w\,(a_{n} + -\alpha - 0 = w \to \lvert a_{n} + -\alpha - 0 \rvert = \lvert a_{n} + -\alpha - 0 \rvert \to \lvert a_{n} + -\alpha - 0 \rvert = \lvert w \rvert)$
\item $\varphi_{87} :\equiv a_{n} + -\alpha = a_{n} - \alpha \to \lvert a_{n} + -\alpha \rvert = \lvert a_{n} + -\alpha \rvert \to \lvert a_{n} + -\alpha \rvert = \lvert a_{n} - \alpha \rvert$
\item $\varphi_{88} :\equiv \forall n \in \mathbb{N}\,(N \leq n \to \lvert a \oplus \overline{-\alpha}(n) - 0 \rvert < \varepsilon) \to n \in \mathbb{N} \to N \leq n \to \lvert a \oplus \overline{-\alpha}(n) - 0 \rvert < \varepsilon$
\item $\varphi_{89} :\equiv \lvert a_{n} + \overline{-\alpha}(n) - 0 \rvert = \lvert a_{n} + z - 0 \rvert \to \lvert a_{n} + \overline{-\alpha}(n) - 0 \rvert = \lvert a_{n} + w - 0 \rvert$
\item $\varphi_{90} :\equiv \lvert a \oplus \overline{-\alpha}(n) - 0 \rvert = \lvert a_{n} - \alpha \rvert \to \lvert a \oplus \overline{-\alpha}(n) - 0 \rvert < \varepsilon \to \lvert a_{n} - \alpha \rvert < \varepsilon$
\item $\varphi_{91} :\equiv \lvert a_{n} - \alpha \rvert = \lvert a \oplus \overline{-\alpha}(n) - 0 \rvert \to \lvert a_{n} - \alpha \rvert < \varepsilon \to \lvert a \oplus \overline{-\alpha}(n) - 0 \rvert < \varepsilon$
\item $\varphi_{92} :\equiv N \in \mathbb{N} \land \forall n \in \mathbb{N}\,(N \leq n \to \lvert a_{n} - \alpha \rvert < \varepsilon) \to n \in \mathbb{N} \to N \leq n \to \lvert a \oplus \overline{-\alpha}(n) - 0 \rvert = \lvert a_{n} - \alpha \rvert$
\item $\varphi_{93} :\equiv N \in \mathbb{N} \land \forall n \in \mathbb{N}\,(N \leq n \to \lvert a_{n} - \alpha \rvert < \varepsilon) \to n \in \mathbb{N} \to N \leq n \to \lvert a_{n} - \alpha \rvert = \lvert a \oplus \overline{-\alpha}(n) - 0 \rvert$
\item $\varphi_{94} :\equiv \forall z\,(\lvert a \oplus \overline{-\alpha}(n) - 0 \rvert = \lvert a_{n} + -\alpha \rvert \to \lvert a_{n} + -\alpha \rvert = z \to \lvert a \oplus \overline{-\alpha}(n) - 0 \rvert = z)$
\item $\varphi_{95} :\equiv N \in \mathbb{N} \land \forall n \in \mathbb{N}\,(N \leq n \to \lvert a \oplus \overline{-\alpha}(n) - 0 \rvert < \varepsilon) \to N \in \mathbb{N} \land \forall n \in \mathbb{N}\,(N \leq n \to \lvert a_{n} - \alpha \rvert < \varepsilon)$
\item $\varphi_{96} :\equiv N \in \mathbb{N} \land \forall n \in \mathbb{N}\,(N \leq n \to \lvert a_{n} - \alpha \rvert < \varepsilon) \to N \in \mathbb{N} \land \forall n \in \mathbb{N}\,(N \leq n \to \lvert a \oplus \overline{-\alpha}(n) - 0 \rvert < \varepsilon)$
\item $\varphi_{97} :\equiv \forall \alpha\,(\mathrm{Conv}(v_{698},\alpha) \leftrightarrow \mathrm{Seq}(v_{698}) \land \alpha \in \mathbb{R} \land \forall \varepsilon \in \mathbb{R}\,(0 < \varepsilon \to \exists N \in \mathbb{N}\,\forall n \in \mathbb{N}\,(N \leq n \to \lvert v_{698}(n) - \alpha \rvert < \varepsilon)))$
\item $\varphi_{98} :\equiv \mathrm{Conv}(v_{698},v_{699}) \leftrightarrow \mathrm{Seq}(v_{698}) \land v_{699} \in \mathbb{R} \land \forall \varepsilon \in \mathbb{R}\,(0 < \varepsilon \to \exists N \in \mathbb{N}\,\forall n \in \mathbb{N}\,(N \leq n \to \lvert v_{698}(n) - v_{699} \rvert < \varepsilon))$
\item $\varphi_{99} :\equiv \lvert a_{n} + \overline{-\alpha}(n) - 0 \rvert = \lvert a_{n} + -\alpha - 0 \rvert \to \lvert a \oplus \overline{-\alpha}(n) - 0 \rvert = \lvert a_{n} + -\alpha - 0 \rvert$
\item $\varphi_{100} :\equiv N \in \mathbb{N} \land \forall n \in \mathbb{N}\,(N \leq n \to \lvert a \oplus \overline{-\alpha}(n) - 0 \rvert < \varepsilon) \to \exists N \in \mathbb{N}\,\forall n \in \mathbb{N}\,(N \leq n \to \lvert a_{n} - \alpha \rvert < \varepsilon)$
\item $\varphi_{101} :\equiv N \in \mathbb{N} \land \forall n \in \mathbb{N}\,(N \leq n \to \lvert a_{n} - \alpha \rvert < \varepsilon) \to \exists N \in \mathbb{N}\,\forall n \in \mathbb{N}\,(N \leq n \to \lvert a \oplus \overline{-\alpha}(n) - 0 \rvert < \varepsilon)$
\item $\varphi_{102} :\equiv (\mathrm{Seq}(a) \to \alpha \in \mathbb{R} \to n \in \mathbb{N} \to a_{n} + -\alpha \in \mathbb{R}) \to \mathrm{Seq}(a) \to \alpha \in \mathbb{R} \to n \in \mathbb{N} \to a_{n} + -\alpha - 0 = a_{n} + -\alpha$
\item $\varphi_{103} :\equiv \mathrm{Seq}(a) \to \alpha \in \mathbb{R} \to \mathrm{Conv}(a,\alpha) \to \varepsilon \in \mathbb{R} \to 0 < \varepsilon \to \varphi_{25}$
\item $\varphi_{104} :\equiv \mathrm{Seq}(a) \to \alpha \in \mathbb{R} \to \mathrm{Conv}(a,\alpha) \to \varepsilon \in \mathbb{R} \to 0 < \varepsilon \to \varphi_{26}$
\item $\varphi_{105} :\equiv \mathrm{Seq}(a) \land \alpha \in \mathbb{R} \land \mathrm{Conv}(a \oplus \overline{-\alpha},0) \land \varepsilon \in \mathbb{R} \land 0 < \varepsilon \land \neg \exists N \in \mathbb{N}\,\forall n \in \mathbb{N}\,(N \leq n \to \lvert a_{n} - \alpha \rvert < \varepsilon)$
\item $\varphi_{106} :\equiv \mathrm{Seq}(a) \land \alpha \in \mathbb{R} \land \mathrm{Conv}(a,\alpha) \land \varepsilon \in \mathbb{R} \land 0 < \varepsilon \land \neg \exists N \in \mathbb{N}\,\forall n \in \mathbb{N}\,(N \leq n \to \lvert a \oplus \overline{-\alpha}(n) - 0 \rvert < \varepsilon)$
\item $\varphi_{107} :\equiv (\mathrm{Seq}(a) \land \alpha \in \mathbb{R} \to \forall n\,(\mathrm{Seq}(a) \land \alpha \in \mathbb{R})) \to \mathrm{Seq}(a) \to \varphi_{34}$
\item $\varphi_{108} :\equiv N \in \mathbb{N} \land \forall n \in \mathbb{N}\,(N \leq n \to \lvert a \oplus \overline{-\alpha}(n) - 0 \rvert < \varepsilon) \to n \in \mathbb{N} \to N \leq n \to \lvert a \oplus \overline{-\alpha}(n) - 0 \rvert < \varepsilon$
\item $\varphi_{109} :\equiv \forall n \in \mathbb{N}\,(\lvert a \oplus \overline{-\alpha}(n) - 0 \rvert = \lvert a_{n} - \alpha \rvert) \to n \in \mathbb{N} \to \lvert a \oplus \overline{-\alpha}(n) - 0 \rvert = \lvert a_{n} - \alpha \rvert$
\item $\varphi_{110} :\equiv \forall c \in \mathbb{R}\,(\mathrm{Seq}(\overline{c}) \land \forall n \in \mathbb{N}\,(\overline{c}(n) = c)) \to \varphi_{23}$
\item $\varphi_{111} :\equiv \forall \varepsilon\,(\mathrm{Seq}(a) \land \alpha \in \mathbb{R} \land \mathrm{Conv}(a \oplus \overline{-\alpha},0) \to \varphi_{35})$
\item $\varphi_{112} :\equiv \forall \varepsilon\,(\mathrm{Seq}(a) \land \alpha \in \mathbb{R} \land \mathrm{Conv}(a,\alpha) \to \varepsilon \in \mathbb{R} \to \varphi_{40})$
\item $\varphi_{113} :\equiv \forall z\,(\lvert a \oplus \overline{-\alpha}(n) - 0 \rvert = \lvert a_{n} + -\alpha - 0 \rvert \to \lvert a_{n} + -\alpha - 0 \rvert = z \to \lvert a \oplus \overline{-\alpha}(n) - 0 \rvert = z)$
\item $\varphi_{114} :\equiv \forall v_{702}\,(\mathrm{Conv}(a \oplus \overline{-\alpha},0) \to v_{702} \in \mathbb{R} \to 0 < v_{702} \to \exists N \in \mathbb{N}\,\forall n \in \mathbb{N}\,(N \leq n \to \lvert a \oplus \overline{-\alpha}(n) - 0 \rvert < v_{702}))$
\item $\varphi_{115} :\equiv \forall n \in \mathbb{N}\,(a \oplus \overline{-\alpha}(n) = a_{n} + \overline{-\alpha}(n)) \to n \in \mathbb{N} \to a \oplus \overline{-\alpha}(n) = a_{n} + \overline{-\alpha}(n)$
\item $\varphi_{116} :\equiv \forall \varepsilon\,(\mathrm{Seq}(a) \land \alpha \in \mathbb{R} \land \mathrm{Conv}(a \oplus \overline{-\alpha},0)) \to \varphi_{39}$
\item $\varphi_{117} :\equiv \forall \varepsilon\,(\mathrm{Seq}(a) \land \alpha \in \mathbb{R} \land \mathrm{Conv}(a,\alpha)) \to \forall \varepsilon \in \mathbb{R}\,(\varphi_{40})$
\item $\varphi_{118} :\equiv (\mathrm{Seq}(a) \to n \in \mathbb{N} \to a_{n} \in \mathbb{R}) \to (\mathrm{Seq}(a) \to \mathrm{Seq}(a)) \to \varphi_{51}$
\item $\varphi_{119} :\equiv N \in \mathbb{N} \land \forall n \in \mathbb{N}\,(N \leq n \to \lvert a \oplus \overline{-\alpha}(n) - 0 \rvert < \varepsilon) \to n \in \mathbb{N} \to N \leq n \to \lvert a \oplus \overline{-\alpha}(n) - 0 \rvert = \lvert a_{n} - \alpha \rvert$
\item $\varphi_{120} :\equiv \lvert a_{n} + \overline{-\alpha}(n) - 0 \rvert = \lvert a_{n} + \overline{-\alpha}(n) - 0 \rvert \to \lvert a_{n} + \overline{-\alpha}(n) - 0 \rvert = \lvert a_{n} + -\alpha - 0 \rvert$
\item $\varphi_{121} :\equiv \mathrm{Seq}(a) \land \alpha \in \mathbb{R} \land \mathrm{Conv}(a \oplus \overline{-\alpha},0) \land \varepsilon \in \mathbb{R} \land 0 < \varepsilon \land (N \in \mathbb{N} \land \forall n \in \mathbb{N}\,(N \leq n \to \lvert a \oplus \overline{-\alpha}(n) - 0 \rvert < \varepsilon))$
\item $\varphi_{122} :\equiv N \in \mathbb{N} \land \forall n \in \mathbb{N}\,(N \leq n \to \lvert a \oplus \overline{-\alpha}(n) - 0 \rvert < \varepsilon) \to N \in \mathbb{N} \land \forall n \in \mathbb{N}\,(N \leq n \to \lvert a \oplus \overline{-\alpha}(n) - 0 \rvert < \varepsilon)$
\item $\varphi_{123} :\equiv \forall v_{701}\,\forall v_{702}\,(\mathrm{Conv}(a \oplus \overline{-\alpha},v_{701}) \to v_{702} \in \mathbb{R} \to 0 < v_{702} \to \exists N \in \mathbb{N}\,\forall n \in \mathbb{N}\,(N \leq n \to \lvert a \oplus \overline{-\alpha}(n) - v_{701} \rvert < v_{702}))$
\item $\varphi_{124} :\equiv \mathrm{Seq}(a) \to \alpha \in \mathbb{R} \to \mathrm{Conv}(a \oplus \overline{-\alpha},0) \to \varepsilon \in \mathbb{R} \to \varphi_{40}$
\item $\varphi_{125} :\equiv N \in \mathbb{N} \land \forall n \in \mathbb{N}\,(N \leq n \to \lvert a \oplus \overline{-\alpha}(n) - 0 \rvert < \varepsilon) \to \varphi_{21}$
\item $\varphi_{126} :\equiv \lvert a \oplus \overline{-\alpha}(n) - 0 \rvert = \lvert a \oplus \overline{-\alpha}(n) - 0 \rvert \to \lvert a \oplus \overline{-\alpha}(n) - 0 \rvert = \lvert a_{n} + \overline{-\alpha}(n) - 0 \rvert$
\item $\varphi_{127} :\equiv N \in \mathbb{N} \land \forall n \in \mathbb{N}\,(N \leq n \to \lvert a \oplus \overline{-\alpha}(n) - 0 \rvert < \varepsilon) \to \varphi_{27}$
\item $\varphi_{128} :\equiv \mathrm{Seq}(a) \to \alpha \in \mathbb{R} \to \mathrm{Conv}(a,\alpha) \to \varepsilon \in \mathbb{R} \to 0 < \varepsilon \to \varphi_{62}$
\item $\varphi_{129} :\equiv \forall w\,(a \oplus \overline{-\alpha}(n) = w \to \lvert a \oplus \overline{-\alpha}(n) - 0 \rvert = \lvert a \oplus \overline{-\alpha}(n) - 0 \rvert \to \lvert a \oplus \overline{-\alpha}(n) - 0 \rvert = \lvert w - 0 \rvert)$
\item $\varphi_{130} :\equiv \mathrm{Seq}(a) \to \alpha \in \mathbb{R} \to \mathrm{Conv}(a \oplus \overline{-\alpha},0) \to \varepsilon \in \mathbb{R} \to 0 < \varepsilon \to \varphi_{37}$
\item $\varphi_{131} :\equiv (\mathrm{Seq}(a) \to \alpha \in \mathbb{R} \to \mathrm{Seq}(\overline{-\alpha})) \to \mathrm{Seq}(a) \to \alpha \in \mathbb{R} \to \varphi_{63}$
\item $\varphi_{132} :\equiv \forall z\,(\lvert a \oplus \overline{-\alpha}(n) - 0 \rvert = \lvert a_{n} + \overline{-\alpha}(n) - 0 \rvert \to \lvert a_{n} + \overline{-\alpha}(n) - 0 \rvert = z \to \lvert a \oplus \overline{-\alpha}(n) - 0 \rvert = z)$
\item $\varphi_{133} :\equiv \forall w\,(\overline{-\alpha}(n) = w \to \lvert a_{n} + \overline{-\alpha}(n) - 0 \rvert = \lvert a_{n} + \overline{-\alpha}(n) - 0 \rvert \to \lvert a_{n} + \overline{-\alpha}(n) - 0 \rvert = \lvert a_{n} + w - 0 \rvert)$
\item $\varphi_{134} :\equiv \mathrm{Seq}(a) \to \alpha \in \mathbb{R} \to \mathrm{Conv}(a,\alpha) \to \varepsilon \in \mathbb{R} \to 0 < \varepsilon \to \varphi_{78}$
\item $\varphi_{135} :\equiv (\mathrm{Seq}(a) \to \alpha \in \mathbb{R} \to \varphi_{63}) \to \mathrm{Seq}(a) \to \alpha \in \mathbb{R} \to \mathrm{Seq}(a \oplus \overline{-\alpha})$
\item $\varphi_{136} :\equiv \mathrm{Seq}(a) \to \alpha \in \mathbb{R} \to \mathrm{Conv}(a \oplus \overline{-\alpha},0) \to \varepsilon \in \mathbb{R} \to 0 < \varepsilon \to \varphi_{56}$
\item $\varphi_{137} :\equiv \mathrm{Seq}(a) \to \alpha \in \mathbb{R} \to \mathrm{Conv}(a \oplus \overline{-\alpha},0) \to \varepsilon \in \mathbb{R} \to 0 < \varepsilon \to \varphi_{57}$
\item $\varphi_{138} :\equiv \mathrm{Seq}(a) \to \alpha \in \mathbb{R} \to \mathrm{Conv}(a,\alpha) \to \varepsilon \in \mathbb{R} \to 0 < \varepsilon \to \varphi_{79}$
\item $\varphi_{139} :\equiv \mathrm{Seq}(a) \to \alpha \in \mathbb{R} \to \mathrm{Conv}(a,\alpha) \to \varepsilon \in \mathbb{R} \to 0 < \varepsilon \to \varphi_{81}$
\item $\varphi_{140} :\equiv \mathrm{Seq}(a) \to \alpha \in \mathbb{R} \to \mathrm{Conv}(a,\alpha) \to \varepsilon \in \mathbb{R} \to 0 < \varepsilon \to \varphi_{82}$
\item $\varphi_{141} :\equiv \mathrm{Conv}(a \oplus \overline{-\alpha},0) \leftrightarrow \mathrm{Seq}(a \oplus \overline{-\alpha}) \land 0 \in \mathbb{R} \land \forall \varepsilon \in \mathbb{R}\,(\varphi_{40})$
\item $\varphi_{142} :\equiv \mathrm{Seq}(a) \to \alpha \in \mathbb{R} \to \mathrm{Conv}(a,\alpha) \to \varepsilon \in \mathbb{R} \to 0 < \varepsilon \to \varphi_{85}$
\item $\varphi_{143} :\equiv \forall n\,(\varphi_{77}) \to \forall n \in \mathbb{N}\,(N \leq n \to \lvert a \oplus \overline{-\alpha}(n) - 0 \rvert < \varepsilon)$
\item $\varphi_{144} :\equiv \mathrm{Seq}(a) \to \alpha \in \mathbb{R} \to \mathrm{Conv}(a,\alpha) \to \varepsilon \in \mathbb{R} \to 0 < \varepsilon \to \varphi_{92}$
\item $\varphi_{145} :\equiv \mathrm{Seq}(a) \to \alpha \in \mathbb{R} \to \mathrm{Conv}(a,\alpha) \to \varepsilon \in \mathbb{R} \to 0 < \varepsilon \to \varphi_{93}$
\item $\varphi_{146} :\equiv \mathrm{Seq}(a) \to \alpha \in \mathbb{R} \to \mathrm{Conv}(a,\alpha) \to \varepsilon \in \mathbb{R} \to 0 < \varepsilon \to \varphi_{96}$
\item $\varphi_{147} :\equiv \forall v_{699}\,(\mathrm{Conv}(a \oplus \overline{-\alpha},v_{699}) \leftrightarrow \mathrm{Seq}(a \oplus \overline{-\alpha}) \land v_{699} \in \mathbb{R} \land \forall \varepsilon \in \mathbb{R}\,(0 < \varepsilon \to \exists N \in \mathbb{N}\,\forall n \in \mathbb{N}\,(N \leq n \to \lvert a \oplus \overline{-\alpha}(n) - v_{699} \rvert < \varepsilon)))$
\item $\varphi_{148} :\equiv \mathrm{Seq}(a) \to \alpha \in \mathbb{R} \to \mathrm{Conv}(a,\alpha) \to \varepsilon \in \mathbb{R} \to 0 < \varepsilon \to \varphi_{101}$
\item $\varphi_{149} :\equiv \mathrm{Seq}(a) \to \alpha \in \mathbb{R} \to \mathrm{Conv}(a \oplus \overline{-\alpha},0) \to \varepsilon \in \mathbb{R} \to 0 < \varepsilon \to \varphi_{80}$
\item $\varphi_{150} :\equiv \mathrm{Seq}(a) \to \alpha \in \mathbb{R} \to \mathrm{Conv}(a \oplus \overline{-\alpha},0) \to \varepsilon \in \mathbb{R} \to 0 < \varepsilon \to \varphi_{84}$
\item $\varphi_{151} :\equiv \forall N\,(\varphi_{106} \to \neg (N \in \mathbb{N} \land \forall n \in \mathbb{N}\,(N \leq n \to \lvert a_{n} - \alpha \rvert < \varepsilon)))$
\item $\varphi_{152} :\equiv (\mathrm{Seq}(a) \to \alpha \in \mathbb{R} \to \mathrm{Conv}(a \oplus \overline{-\alpha},0) \to \varphi_{39}) \to \varphi_{5}$
\item $\varphi_{153} :\equiv (\mathrm{Seq}(a) \to \alpha \in \mathbb{R} \to \mathrm{Conv}(a,\alpha) \to \forall \varepsilon \in \mathbb{R}\,(\varphi_{40})) \to \varphi_{6}$
\item $\varphi_{154} :\equiv \forall \alpha\,\forall \varepsilon\,(\mathrm{Conv}(a,\alpha) \to \varphi_{35}) \to \forall \varepsilon\,(\mathrm{Conv}(a,\alpha) \to \varphi_{35})$
\item $\varphi_{155} :\equiv \mathrm{Seq}(a) \to \alpha \in \mathbb{R} \to \mathrm{Conv}(a \oplus \overline{-\alpha},0) \to \varepsilon \in \mathbb{R} \to 0 < \varepsilon \to \varphi_{95}$
\item $\varphi_{156} :\equiv \mathrm{Seq}(a) \to \alpha \in \mathbb{R} \to \mathrm{Conv}(a \oplus \overline{-\alpha},0) \to \varepsilon \in \mathbb{R} \to 0 < \varepsilon \to \varphi_{100}$
\item $\varphi_{157} :\equiv \forall a\,\forall \alpha\,\forall \varepsilon\,(\mathrm{Conv}(a,\alpha) \to \varphi_{35}) \to \forall \alpha\,\forall \varepsilon\,(\mathrm{Conv}(a,\alpha) \to \varphi_{35})$
\item $\varphi_{158} :\equiv \varphi_{11} \to \mathrm{Seq}(a) \to \alpha \in \mathbb{R} \to \mathrm{Conv}(a,\alpha) \to \forall \varepsilon \in \mathbb{R}\,(\varphi_{40})$
\item $\varphi_{159} :\equiv \mathrm{Seq}(a) \to \alpha \in \mathbb{R} \to \mathrm{Conv}(a \oplus \overline{-\alpha},0) \to \varepsilon \in \mathbb{R} \to 0 < \varepsilon \to \varphi_{108}$
\item $\varphi_{160} :\equiv \mathrm{Seq}(a) \to \alpha \in \mathbb{R} \to \mathrm{Conv}(a \oplus \overline{-\alpha},0) \to \varepsilon \in \mathbb{R} \to 0 < \varepsilon \to \varphi_{119}$
\item $\varphi_{161} :\equiv (a_{n} \in \mathbb{R} \to -\alpha \in \mathbb{R} \to a_{n} + -\alpha \in \mathbb{R}) \to (\mathrm{Seq}(a) \to \alpha \in \mathbb{R} \to n \in \mathbb{N} \to a_{n} \in \mathbb{R}) \to \varphi_{73}$
\item $\varphi_{162} :\equiv \mathrm{Seq}(a) \to \alpha \in \mathbb{R} \to \mathrm{Conv}(a \oplus \overline{-\alpha},0) \to \varepsilon \in \mathbb{R} \to 0 < \varepsilon \to \varphi_{122}$
\item $\varphi_{163} :\equiv \mathrm{Seq}(a) \to \alpha \in \mathbb{R} \to \mathrm{Conv}(a \oplus \overline{-\alpha},0) \to \varepsilon \in \mathbb{R} \to 0 < \varepsilon \to \varphi_{125}$
\item $\varphi_{164} :\equiv \mathrm{Seq}(a) \to \alpha \in \mathbb{R} \to \mathrm{Conv}(a \oplus \overline{-\alpha},0) \to \varepsilon \in \mathbb{R} \to 0 < \varepsilon \to \varphi_{127}$
\item $\varphi_{165} :\equiv (\mathrm{Seq}(a) \to \varphi_{54}) \to (\mathrm{Seq}(a) \to \alpha \in \mathbb{R} \to n \in \mathbb{N} \to n \in \mathbb{N}) \to \varphi_{48}$
\item $\varphi_{166} :\equiv \forall \alpha\,(\mathrm{Conv}(a,\alpha) \leftrightarrow \mathrm{Seq}(a) \land \alpha \in \mathbb{R} \land \varphi_{39}) \to (\mathrm{Conv}(a,\alpha) \leftrightarrow \mathrm{Seq}(a) \land \alpha \in \mathbb{R} \land \varphi_{39})$
\item $\varphi_{167} :\equiv \forall a\,\forall \alpha\,(\mathrm{Conv}(a,\alpha) \leftrightarrow \mathrm{Seq}(a) \land \alpha \in \mathbb{R} \land \varphi_{39}) \to \forall \alpha\,(\mathrm{Conv}(a,\alpha) \leftrightarrow \mathrm{Seq}(a) \land \alpha \in \mathbb{R} \land \varphi_{39})$
\item $\varphi_{168} :\equiv (\mathrm{Seq}(a) \to \mathrm{Seq}(a)) \to (\mathrm{Seq}(a) \to \alpha \in \mathbb{R} \to \alpha \in \mathbb{R}) \to \varphi_{152}$
\item $\varphi_{169} :\equiv \forall a\,\forall \alpha\,(\mathrm{Conv}(a,\alpha) \leftrightarrow \mathrm{Seq}(a) \land \alpha \in \mathbb{R} \land \varphi_{39}) \to \varphi_{97}$
\item $\varphi_{170} :\equiv (\mathrm{Seq}(a) \to \alpha \in \mathbb{R} \to \mathrm{Seq}(a \oplus \overline{-\alpha})) \to 0 \in \mathbb{R} \to \varphi_{153}$
\item $\varphi_{171} :\equiv (\mathrm{Seq}(a) \to \alpha \in \mathbb{R} \to \mathrm{Conv}(a \oplus \overline{-\alpha},0) \to \varphi_{35}) \to \mathrm{Seq}(a) \land \alpha \in \mathbb{R} \land \mathrm{Conv}(a \oplus \overline{-\alpha},0) \to \varphi_{35}$
\item $\varphi_{172} :\equiv (\mathrm{Seq}(a) \to \alpha \in \mathbb{R} \to \mathrm{Conv}(a,\alpha) \to \varepsilon \in \mathbb{R} \to \varphi_{40}) \to \mathrm{Seq}(a) \land \alpha \in \mathbb{R} \land \mathrm{Conv}(a,\alpha) \to \varepsilon \in \mathbb{R} \to \varphi_{40}$
\item $\varphi_{173} :\equiv \lvert a_{n} + -\alpha - 0 \rvert = \lvert a_{n} + -\alpha - 0 \rvert \to (a_{n} + -\alpha - 0 = a_{n} + -\alpha \to \varphi_{72}) \to \varphi_{36}$
\item $\varphi_{174} :\equiv (\mathrm{Seq}(a) \to \mathrm{Seq}(\overline{-\alpha}) \to \varphi_{63}) \to (\mathrm{Seq}(a) \to \mathrm{Seq}(a)) \to \varphi_{131}$
\item $\varphi_{175} :\equiv (\lvert a \oplus \overline{-\alpha}(n) - 0 \rvert = \lvert a_{n} + \overline{-\alpha}(n) - 0 \rvert \to \varphi_{99}) \to \mathrm{Seq}(a) \to \varphi_{46}$
\item $\varphi_{176} :\equiv (a \oplus \overline{-\alpha}(n) = a_{n} + \overline{-\alpha}(n) \to \varphi_{126}) \to \mathrm{Seq}(a) \to \alpha \in \mathbb{R} \to \varphi_{45}$
\item $\varphi_{177} :\equiv (\varphi_{4} \to \neg (N \in \mathbb{N} \land \forall n \in \mathbb{N}\,(N \leq n \to \lvert a_{n} - \alpha \rvert < \varepsilon))) \to \varphi_{156}$
\item $\varphi_{178} :\equiv \forall n\,(\varphi_{77} \to n \in \mathbb{N} \to N \leq n \to \lvert a \oplus \overline{-\alpha}(n) - 0 \rvert < \varepsilon) \to \varphi_{143}$
\item $\varphi_{179} :\equiv \lvert a_{n} + \overline{-\alpha}(n) - 0 \rvert = \lvert a_{n} + \overline{-\alpha}(n) - 0 \rvert \to (\overline{-\alpha}(n) = -\alpha \to \varphi_{120}) \to \mathrm{Seq}(a) \to \varphi_{50}$
\item $\varphi_{180} :\equiv (\mathrm{Seq}(a) \to \alpha \in \mathbb{R} \to n \in \mathbb{N} \to a_{n} + -\alpha - 0 = a_{n} + -\alpha) \to \varphi_{173}$
\item $\varphi_{181} :\equiv (\mathrm{Conv}(a,\alpha) \to \varphi_{35}) \to \varphi_{1} \to \varphi_{3} \to \varphi_{9} \to \mathrm{Seq}(a) \to \alpha \in \mathbb{R} \to \mathrm{Conv}(a,\alpha) \to \varphi_{35}$
\item $\varphi_{182} :\equiv \varphi_{148} \to \varphi_{106} \to \neg (N \in \mathbb{N} \land \forall n \in \mathbb{N}\,(N \leq n \to \lvert a_{n} - \alpha \rvert < \varepsilon))$
\item $\varphi_{183} :\equiv \forall n\,(\varphi_{121} \to n \in \mathbb{N} \to N \leq n \to \lvert a_{n} - \alpha \rvert < \varepsilon) \to \forall n\,(\varphi_{121}) \to \forall n \in \mathbb{N}\,(N \leq n \to \lvert a_{n} - \alpha \rvert < \varepsilon)$
\item $\varphi_{184} :\equiv (\varphi_{105} \to \forall N\,(\varphi_{105})) \to \mathrm{Seq}(a) \to \alpha \in \mathbb{R} \to \mathrm{Conv}(a \oplus \overline{-\alpha},0) \to \varphi_{35}$
\item $\varphi_{185} :\equiv (\varphi_{106} \to \forall N\,(\varphi_{106})) \to \mathrm{Seq}(a) \to \alpha \in \mathbb{R} \to \mathrm{Conv}(a,\alpha) \to \varepsilon \in \mathbb{R} \to \varphi_{40}$
\item $\varphi_{186} :\equiv (\mathrm{Seq}(a) \to \varphi_{46}) \to \varphi_{36} \to (\lvert a \oplus \overline{-\alpha}(n) - 0 \rvert = \lvert a_{n} + -\alpha - 0 \rvert \to \varphi_{76}) \to \varphi_{52}$
\item $\varphi_{187} :\equiv (\mathrm{Seq}(a) \to \alpha \in \mathbb{R} \to \mathrm{Conv}(a,\alpha) \to \varepsilon \in \mathbb{R} \to \varphi_{22}) \to \varphi_{142} \to \varphi_{146}$
\item $\varphi_{188} :\equiv (\varphi_{111} \to \varphi_{116}) \to \varphi_{69} \to \mathrm{Seq}(a) \to \alpha \in \mathbb{R} \to \mathrm{Conv}(a \oplus \overline{-\alpha},0) \to \varphi_{39}$
\item $\varphi_{189} :\equiv (\mathrm{Conv}(a \oplus \overline{-\alpha},0) \to \varepsilon \in \mathbb{R} \to \varphi_{40}) \to (\mathrm{Seq}(a) \to \varphi_{14}) \to \varphi_{24} \to (\mathrm{Seq}(a) \to \varphi_{19}) \to \varphi_{124}$
\item $\varphi_{190} :\equiv \forall n\,(\varphi_{77} \to n \in \mathbb{N} \to N \leq n \to \lvert a \oplus \overline{-\alpha}(n) - 0 \rvert < \varepsilon) \to \varphi_{178} \to (\varphi_{77} \to \forall n\,(\varphi_{77})) \to \varphi_{142}$
\item $\varphi_{191} :\equiv \forall N\,(\varphi_{105} \to \varphi_{17}) \to (\forall N\,(\varphi_{105} \to \varphi_{17}) \to \forall N\,(\varphi_{105}) \to \forall N\,(\varphi_{17})) \to \varphi_{184}$
\item $\varphi_{192} :\equiv \forall n\,(\varphi_{121} \to n \in \mathbb{N} \to N \leq n \to \lvert a_{n} - \alpha \rvert < \varepsilon) \to \varphi_{183} \to (\varphi_{121} \to \forall n\,(\varphi_{121})) \to \varphi_{150}$
\item $\varphi_{193} :\equiv (\mathrm{Seq}(a) \to \alpha \in \mathbb{R} \to \mathrm{Conv}(a,\alpha) \to \varphi_{35}) \to \varphi_{151} \to (\varphi_{151} \to \forall N\,(\varphi_{106}) \to \varphi_{4}) \to \varphi_{185}$
\end{abbrevlist}
\end{fullproof}

% 自動生成 (prover/build.js): InvInv  ノード数 36
\begin{fullproof}{[E30]}{$\forall x \in \mathbb{R}\,(x \neq 0 \to (x^{-1})^{-1} = x)$}{36}\phantomsection\label{tree:InvInv}
\noindent$\Pi_{1}$:
\begin{ptree}
\AXC{$x = x$\lbl{A5}}
\AXC{$x = y \to x = x \to y = x$\lbl{A6}}
\AXC{$(x = y \to x = x \to y = x) \to x = x \to x = y \to y = x$\lbl{A1}}
\BIC{$x = x \to x = y \to y = x$}
\BIC{$x = y \to y = x$}
\UIC{$\forall y\,(x = y \to y = x)$}
\UIC{$\forall x\,\forall y\,(x = y \to y = x)$}
\end{ptree}
\noindent$\Pi_{2}$:
\begin{ptree}
\AXC{$\varphi_{2}$\lbl{[E13]}}
\AXC{$\varphi_{2} \to \varphi_{8}$\lbl{A2}}
\BIC{$\varphi_{8}$}
\AXC{$\varphi_{8} \to \varphi_{7}$\lbl{A2}}
\BIC{$\varphi_{7}$}
\AXC{$\varphi_{7} \to (x \in \mathbb{R} \to x \neq 0 \to x^{-1} \in \mathbb{R}) \to \varphi_{10}$\lbl{A1}}
\BIC{$(x \in \mathbb{R} \to x \neq 0 \to x^{-1} \in \mathbb{R}) \to \varphi_{10}$}
\end{ptree}
\noindent$\Pi_{3}$:
\begin{ptree}
\AXC{$x \in \mathbb{R} \to x \neq 0 \to x \in \mathbb{R}$\lbl{A1}}
\AXC{$\forall x \in \mathbb{R}\,(x \neq 0 \to x^{-1} \in \mathbb{R})$\lbl{(R1)}}
\AXC{$\varphi_{9}$\lbl{A2}}
\BIC{$x \in \mathbb{R} \to x \neq 0 \to x^{-1} \in \mathbb{R}$}
\AXC{$\stackrel{\Pi_{2}}{\vdots}$}\noLine
\UIC{$(x \in \mathbb{R} \to x \neq 0 \to x^{-1} \in \mathbb{R}) \to \varphi_{10}$}
\BIC{$\varphi_{10}$}
\BIC{$(x \in \mathbb{R} \to x \neq 0 \to x^{-1} \neq 0) \to \varphi_{5}$}
\end{ptree}
\noindent$\Pi_{4}$:
\begin{ptree}
\AXC{$\forall x \in \mathbb{R}\,(x \neq 0 \to x^{-1} \neq 0)$\lbl{[E14]}}
\AXC{$\varphi_{3}$\lbl{A2}}
\BIC{$x \in \mathbb{R} \to x \neq 0 \to x^{-1} \neq 0$}
\AXC{$\stackrel{\Pi_{3}}{\vdots}$}\noLine
\UIC{$(x \in \mathbb{R} \to x \neq 0 \to x^{-1} \neq 0) \to \varphi_{5}$}
\BIC{$\varphi_{5}$}
\end{ptree}
\noindent$\Pi_{5}$:
\begin{ptree}
\AXC{$\forall x \in \mathbb{R}\,(x \neq 0 \to x^{-1} \cdot x = 1)$\lbl{[E22]}}
\AXC{$\varphi_{6}$\lbl{A2}}
\BIC{$x \in \mathbb{R} \to x \neq 0 \to x^{-1} \cdot x = 1$}
\AXC{$\stackrel{\Pi_{4}}{\vdots}$}\noLine
\UIC{$\varphi_{5}$}
\BIC{$x \in \mathbb{R} \to x \neq 0 \to x = (x^{-1})^{-1}$}
\AXC{$(x \in \mathbb{R} \to x \neq 0 \to x = (x^{-1})^{-1}) \to \varphi_{4}$\lbl{A1}}
\BIC{$\varphi_{4}$}
\end{ptree}
\noindent 主木:
\begin{ptree}
\AXC{$\stackrel{\Pi_{1}}{\vdots}$}\noLine
\UIC{$\forall x\,\forall y\,(x = y \to y = x)$}
\AXC{$\forall x\,\forall y\,(x = y \to y = x) \to \forall y\,(x = y \to y = x)$\lbl{A2}}
\BIC{$\forall y\,(x = y \to y = x)$}
\AXC{$\varphi_{1}$\lbl{A2}}
\BIC{$x = (x^{-1})^{-1} \to (x^{-1})^{-1} = x$}
\AXC{$\stackrel{\Pi_{5}}{\vdots}$}\noLine
\UIC{$\varphi_{4}$}
\BIC{$x \in \mathbb{R} \to x \neq 0 \to (x^{-1})^{-1} = x$}
\UIC{$\forall x \in \mathbb{R}\,(x \neq 0 \to (x^{-1})^{-1} = x)$}
\end{ptree}
\begin{abbrevlist}
\item $\varphi_{1} :\equiv \forall y\,(x = y \to y = x) \to x = (x^{-1})^{-1} \to (x^{-1})^{-1} = x$
\item $\varphi_{2} :\equiv \forall x\,\forall y\,(x \in \mathbb{R} \to y \in \mathbb{R} \to x \neq 0 \to x \cdot y = 1 \to y = x^{-1})$
\item $\varphi_{3} :\equiv \forall x \in \mathbb{R}\,(x \neq 0 \to x^{-1} \neq 0) \to x \in \mathbb{R} \to x \neq 0 \to x^{-1} \neq 0$
\item $\varphi_{4} :\equiv (x = (x^{-1})^{-1} \to (x^{-1})^{-1} = x) \to x \in \mathbb{R} \to x \neq 0 \to (x^{-1})^{-1} = x$
\item $\varphi_{5} :\equiv (x \in \mathbb{R} \to x \neq 0 \to x^{-1} \cdot x = 1) \to x \in \mathbb{R} \to x \neq 0 \to x = (x^{-1})^{-1}$
\item $\varphi_{6} :\equiv \forall x \in \mathbb{R}\,(x \neq 0 \to x^{-1} \cdot x = 1) \to x \in \mathbb{R} \to x \neq 0 \to x^{-1} \cdot x = 1$
\item $\varphi_{7} :\equiv x^{-1} \in \mathbb{R} \to x \in \mathbb{R} \to x^{-1} \neq 0 \to x^{-1} \cdot x = 1 \to x = (x^{-1})^{-1}$
\item $\varphi_{8} :\equiv \forall y\,(x^{-1} \in \mathbb{R} \to y \in \mathbb{R} \to x^{-1} \neq 0 \to x^{-1} \cdot y = 1 \to y = (x^{-1})^{-1})$
\item $\varphi_{9} :\equiv \forall x \in \mathbb{R}\,(x \neq 0 \to x^{-1} \in \mathbb{R}) \to x \in \mathbb{R} \to x \neq 0 \to x^{-1} \in \mathbb{R}$
\item $\varphi_{10} :\equiv (x \in \mathbb{R} \to x \neq 0 \to x \in \mathbb{R}) \to (x \in \mathbb{R} \to x \neq 0 \to x^{-1} \neq 0) \to \varphi_{5}$
\end{abbrevlist}
\end{fullproof}

% 自動生成 (prover/build.js): DvgRecip  ノード数 298
\begin{fullproof}{定理2.16}{$\forall a\,(\mathrm{DvgP}(a) \to \forall n \in \mathbb{N}\,(a_{n} \neq 0) \to \mathrm{Conv}(a^{\ast},0))$}{298}\phantomsection\label{tree:DvgRecip}
\noindent$\Pi_{1}$:
\begin{ptree}
\AXC{$z = w \to \lvert z - 0 \rvert < \varepsilon \to \lvert w - 0 \rvert < \varepsilon$\lbl{A6}}
\UIC{$\forall w\,(z = w \to \lvert z - 0 \rvert < \varepsilon \to \lvert w - 0 \rvert < \varepsilon)$}
\UIC{$\forall z\,\forall w\,(z = w \to \lvert z - 0 \rvert < \varepsilon \to \lvert w - 0 \rvert < \varepsilon)$}
\AXC{$\varphi_{73}$\lbl{A2}}
\BIC{$\forall w\,(a_{n}^{-1} = w \to \lvert a_{n}^{-1} - 0 \rvert < \varepsilon \to \lvert w - 0 \rvert < \varepsilon)$}
\AXC{$\varphi_{85}$\lbl{A2}}
\BIC{$\varphi_{46}$}
\end{ptree}
\noindent$\Pi_{2}$:
\begin{ptree}
\AXC{$\mathrm{DvgP}(a) \to \forall n \in \mathbb{N}\,(a_{n} \neq 0) \to \varepsilon \in \mathbb{R} \to 0 < \varepsilon \to \varphi_{58}$\lbl{A1}}
\AXC{$\varphi_{126}$\lbl{A1}}
\BIC{$\mathrm{DvgP}(a) \to \forall n \in \mathbb{N}\,(a_{n} \neq 0) \to \varepsilon \in \mathbb{R} \to 0 < \varepsilon \to \varphi_{64}$}
\AXC{$\varphi_{139}$\lbl{A1}}
\BIC{$\varphi_{99}$}
\end{ptree}
\noindent$\Pi_{3}$:
\begin{ptree}
\AXC{$\forall n \in \mathbb{N}\,(N \leq n \to \varepsilon^{-1} < a_{n}) \to n \in \mathbb{N} \to N \leq n \to \varepsilon^{-1} < a_{n}$\lbl{A2}}
\AXC{$\stackrel{\Pi_{2}}{\vdots}$}\noLine
\UIC{$\varphi_{99}$}
\BIC{$\mathrm{DvgP}(a) \to \forall n \in \mathbb{N}\,(a_{n} \neq 0) \to \varepsilon \in \mathbb{R} \to 0 < \varepsilon \to \varphi_{63}$}
\AXC{$\varphi_{154}$\lbl{A1}}
\BIC{$(\mathrm{DvgP}(a) \to \forall n \in \mathbb{N}\,(a_{n} \neq 0) \to \varepsilon \in \mathbb{R} \to \varphi_{30}) \to \varphi_{112}$}
\end{ptree}
\noindent$\Pi_{4}$:
\begin{ptree}
\AXC{$\mathrm{DvgP}(a) \to \forall n \in \mathbb{N}\,(a_{n} \neq 0) \to \varepsilon \in \mathbb{R} \to \varphi_{30}$\lbl{A1}}
\AXC{$\stackrel{\Pi_{3}}{\vdots}$}\noLine
\UIC{$(\mathrm{DvgP}(a) \to \forall n \in \mathbb{N}\,(a_{n} \neq 0) \to \varepsilon \in \mathbb{R} \to \varphi_{30}) \to \varphi_{112}$}
\BIC{$\varphi_{112}$}
\end{ptree}
\noindent$\Pi_{5}$:
\begin{ptree}
\AXC{$\mathrm{DvgP}(a) \to \forall n \in \mathbb{N}\,(a_{n} \neq 0) \to \varepsilon \in \mathbb{R} \to 0 < \varepsilon \to \varphi_{16}$\lbl{A1}}
\AXC{$\stackrel{\Pi_{4}}{\vdots}$}\noLine
\UIC{$\varphi_{112}$}
\BIC{$\mathrm{DvgP}(a) \to \forall n \in \mathbb{N}\,(a_{n} \neq 0) \to \varepsilon \in \mathbb{R} \to 0 < \varepsilon \to \varphi_{42}$}
\end{ptree}
\noindent$\Pi_{6}$:
\begin{ptree}
\AXC{$\forall x \in \mathbb{R}\,(0 < x \to 0 < x^{-1})$\lbl{[F34]}}
\AXC{$\forall x \in \mathbb{R}\,(0 < x \to 0 < x^{-1}) \to \varepsilon \in \mathbb{R} \to 0 < \varepsilon \to 0 < \varepsilon^{-1}$\lbl{A2}}
\BIC{$\varepsilon \in \mathbb{R} \to 0 < \varepsilon \to 0 < \varepsilon^{-1}$}
\end{ptree}
\noindent$\Pi_{7}$:
\begin{ptree}
\AXC{$\mathrm{DvgP}(a) \to \forall n \in \mathbb{N}\,(a_{n} \neq 0) \to \varepsilon \in \mathbb{R} \to \varepsilon \in \mathbb{R}$\lbl{A1}}
\AXC{$\stackrel{\Pi_{6}}{\vdots}$}\noLine
\UIC{$\varepsilon \in \mathbb{R} \to 0 < \varepsilon \to 0 < \varepsilon^{-1}$}
\AXC{$(\varepsilon \in \mathbb{R} \to 0 < \varepsilon \to 0 < \varepsilon^{-1}) \to \varphi_{91}$\lbl{A1}}
\BIC{$\varphi_{91}$}
\BIC{$\varphi_{12} \to \varphi_{17}$}
\end{ptree}
\noindent$\Pi_{8}$:
\begin{ptree}
\AXC{$\varphi_{12}$\lbl{A1}}
\AXC{$\stackrel{\Pi_{7}}{\vdots}$}\noLine
\UIC{$\varphi_{12} \to \varphi_{17}$}
\BIC{$\varphi_{17}$}
\end{ptree}
\noindent$\Pi_{9}$:
\begin{ptree}
\AXC{$\mathrm{DvgP}(a) \to \mathrm{DvgP}(a)$\lbl{A1}}
\AXC{$(\mathrm{DvgP}(a) \to \mathrm{DvgP}(a)) \to \varphi_{78}$\lbl{A1}}
\BIC{$\varphi_{78}$}
\end{ptree}
\noindent$\Pi_{10}$:
\begin{ptree}
\AXC{$\forall a\,(\mathrm{DvgP}(a) \leftrightarrow \mathrm{Seq}(a) \land \varphi_{15})$\lbl{Def$\mathrm{DvgP}$}}
\AXC{$\varphi_{86}$\lbl{A2}}
\BIC{$\mathrm{DvgP}(a) \leftrightarrow \mathrm{Seq}(a) \land \varphi_{15}$}
\AXC{$\stackrel{\Pi_{9}}{\vdots}$}\noLine
\UIC{$\varphi_{78}$}
\BIC{$\mathrm{DvgP}(a) \to \forall n \in \mathbb{N}\,(a_{n} \neq 0) \to \mathrm{Seq}(a)$}
\end{ptree}
\noindent$\Pi_{11}$:
\begin{ptree}
\AXC{$\forall a\,\forall n\,(\mathrm{Seq}(a) \to n \in \mathbb{N} \to a_{n} \in \mathbb{R})$\lbl{[S1]}}
\AXC{$\varphi_{28}$\lbl{A2}}
\BIC{$\forall n\,(\mathrm{Seq}(a) \to n \in \mathbb{N} \to a_{n} \in \mathbb{R})$}
\AXC{$\varphi_{21}$\lbl{A2}}
\BIC{$\mathrm{Seq}(a) \to n \in \mathbb{N} \to a_{n} \in \mathbb{R}$}
\AXC{$\varphi_{130}$\lbl{A1}}
\BIC{$(\mathrm{DvgP}(a) \to \forall n \in \mathbb{N}\,(a_{n} \neq 0) \to \mathrm{Seq}(a)) \to \varphi_{111}$}
\end{ptree}
\noindent$\Pi_{12}$:
\begin{ptree}
\AXC{$\stackrel{\Pi_{10}}{\vdots}$}\noLine
\UIC{$\mathrm{DvgP}(a) \to \forall n \in \mathbb{N}\,(a_{n} \neq 0) \to \mathrm{Seq}(a)$}
\AXC{$\stackrel{\Pi_{11}}{\vdots}$}\noLine
\UIC{$(\mathrm{DvgP}(a) \to \forall n \in \mathbb{N}\,(a_{n} \neq 0) \to \mathrm{Seq}(a)) \to \varphi_{111}$}
\BIC{$\varphi_{111}$}
\end{ptree}
\noindent$\Pi_{13}$:
\begin{ptree}
\AXC{$\mathrm{DvgP}(a) \to \forall n \in \mathbb{N}\,(a_{n} \neq 0) \to \varepsilon \in \mathbb{R} \to \varphi_{30}$\lbl{A1}}
\AXC{$\stackrel{\Pi_{12}}{\vdots}$}\noLine
\UIC{$\varphi_{111}$}
\BIC{$\mathrm{DvgP}(a) \to \forall n \in \mathbb{N}\,(a_{n} \neq 0) \to \varepsilon \in \mathbb{R} \to 0 < \varepsilon \to \varphi_{34}$}
\end{ptree}
\noindent$\Pi_{14}$:
\begin{ptree}
\AXC{$x = x$\lbl{A5}}
\AXC{$x = y \to x = x \to y = x$\lbl{A6}}
\AXC{$(x = y \to x = x \to y = x) \to x = x \to x = y \to y = x$\lbl{A1}}
\BIC{$x = x \to x = y \to y = x$}
\BIC{$x = y \to y = x$}
\UIC{$\forall y\,(x = y \to y = x)$}
\UIC{$\forall x\,\forall y\,(x = y \to y = x)$}
\end{ptree}
\noindent$\Pi_{15}$:
\begin{ptree}
\AXC{$\stackrel{\Pi_{14}}{\vdots}$}\noLine
\UIC{$\forall x\,\forall y\,(x = y \to y = x)$}
\AXC{$\forall x\,\forall y\,(x = y \to y = x) \to \forall y\,(\varepsilon = y \to y = \varepsilon)$\lbl{A2}}
\BIC{$\forall y\,(\varepsilon = y \to y = \varepsilon)$}
\AXC{$\forall y\,(\varepsilon = y \to y = \varepsilon) \to \varepsilon = 0 \to 0 = \varepsilon$\lbl{A2}}
\BIC{$\varepsilon = 0 \to 0 = \varepsilon$}
\end{ptree}
\noindent$\Pi_{16}$:
\begin{ptree}
\AXC{$\forall x\,\forall y\,(x < y \to x \neq y)$\lbl{[F26]}}
\AXC{$\forall x\,\forall y\,(x < y \to x \neq y) \to \forall y\,(0 < y \to 0 \neq y)$\lbl{A2}}
\BIC{$\forall y\,(0 < y \to 0 \neq y)$}
\AXC{$\forall y\,(0 < y \to 0 \neq y) \to 0 < \varepsilon \to 0 \neq \varepsilon$\lbl{A2}}
\BIC{$0 < \varepsilon \to 0 \neq \varepsilon$}
\end{ptree}
\noindent$\Pi_{17}$:
\begin{ptree}
\AXC{$\stackrel{\Pi_{15}}{\vdots}$}\noLine
\UIC{$\varepsilon = 0 \to 0 = \varepsilon$}
\AXC{$\varphi_{12}$\lbl{A1}}
\AXC{$\stackrel{\Pi_{16}}{\vdots}$}\noLine
\UIC{$0 < \varepsilon \to 0 \neq \varepsilon$}
\AXC{$(0 < \varepsilon \to 0 \neq \varepsilon) \to \varphi_{12} \to \varphi_{9}$\lbl{A1}}
\BIC{$\varphi_{12} \to \varphi_{9}$}
\BIC{$\varphi_{9}$}
\AXC{$\varphi_{9} \to (\varepsilon = 0 \to 0 = \varepsilon) \to \varphi_{10}$\lbl{A1}}
\BIC{$(\varepsilon = 0 \to 0 = \varepsilon) \to \varphi_{10}$}
\BIC{$\varphi_{10}$}
\end{ptree}
\noindent$\Pi_{18}$:
\begin{ptree}
\AXC{$\forall x \in \mathbb{R}\,(x \neq 0 \to x^{-1} \in \mathbb{R})$\lbl{(R1)}}
\AXC{$\varphi_{22}$\lbl{A2}}
\BIC{$\varepsilon \in \mathbb{R} \to \varepsilon \neq 0 \to \varepsilon^{-1} \in \mathbb{R}$}
\AXC{$(\varepsilon \in \mathbb{R} \to \varepsilon \neq 0 \to \varepsilon^{-1} \in \mathbb{R}) \to \varphi_{109}$\lbl{A1}}
\BIC{$\varphi_{109}$}
\end{ptree}
\noindent$\Pi_{19}$:
\begin{ptree}
\AXC{$\stackrel{\Pi_{17}}{\vdots}$}\noLine
\UIC{$\varphi_{10}$}
\AXC{$\mathrm{DvgP}(a) \to \forall n \in \mathbb{N}\,(a_{n} \neq 0) \to \varepsilon \in \mathbb{R} \to \varepsilon \in \mathbb{R}$\lbl{A1}}
\AXC{$\stackrel{\Pi_{18}}{\vdots}$}\noLine
\UIC{$\varphi_{109}$}
\BIC{$\varphi_{93}$}
\BIC{$\mathrm{DvgP}(a) \to \forall n \in \mathbb{N}\,(a_{n} \neq 0) \to \varepsilon \in \mathbb{R} \to 0 < \varepsilon \to \varphi_{35}$}
\end{ptree}
\noindent$\Pi_{20}$:
\begin{ptree}
\AXC{$\stackrel{\Pi_{19}}{\vdots}$}\noLine
\UIC{$\mathrm{DvgP}(a) \to \forall n \in \mathbb{N}\,(a_{n} \neq 0) \to \varepsilon \in \mathbb{R} \to 0 < \varepsilon \to \varphi_{35}$}
\AXC{$\varphi_{5}$\lbl{[E19]}}
\AXC{$\varphi_{5} \to \varphi_{33}$\lbl{A2}}
\BIC{$\varphi_{33}$}
\AXC{$\varphi_{89}$\lbl{A2}}
\BIC{$\varepsilon^{-1} \in \mathbb{R} \to a_{n} \in \mathbb{R} \to 0 < \varepsilon^{-1} \to \varepsilon^{-1} < a_{n} \to a_{n}^{-1} < (\varepsilon^{-1})^{-1}$}
\AXC{$\varphi_{158}$\lbl{A1}}
\BIC{$\varphi_{157}$}
\BIC{$\varphi_{151}$}
\end{ptree}
\noindent$\Pi_{21}$:
\begin{ptree}
\AXC{$\stackrel{\Pi_{8}}{\vdots}$}\noLine
\UIC{$\varphi_{17}$}
\AXC{$\stackrel{\Pi_{13}}{\vdots}$}\noLine
\UIC{$\mathrm{DvgP}(a) \to \forall n \in \mathbb{N}\,(a_{n} \neq 0) \to \varepsilon \in \mathbb{R} \to 0 < \varepsilon \to \varphi_{34}$}
\AXC{$\stackrel{\Pi_{20}}{\vdots}$}\noLine
\UIC{$\varphi_{151}$}
\BIC{$\varphi_{17} \to \varphi_{122}$}
\BIC{$\varphi_{122}$}
\end{ptree}
\noindent$\Pi_{22}$:
\begin{ptree}
\AXC{$\forall x \in \mathbb{R}\,(x \neq 0 \to (x^{-1})^{-1} = x)$\lbl{[E30]}}
\AXC{$\varphi_{20}$\lbl{A2}}
\BIC{$\varepsilon \in \mathbb{R} \to \varepsilon \neq 0 \to (\varepsilon^{-1})^{-1} = \varepsilon$}
\AXC{$(\varepsilon \in \mathbb{R} \to \varepsilon \neq 0 \to (\varepsilon^{-1})^{-1} = \varepsilon) \to \varphi_{110}$\lbl{A1}}
\BIC{$\varphi_{110}$}
\end{ptree}
\noindent$\Pi_{23}$:
\begin{ptree}
\AXC{$\stackrel{\Pi_{17}}{\vdots}$}\noLine
\UIC{$\varphi_{10}$}
\AXC{$\mathrm{DvgP}(a) \to \forall n \in \mathbb{N}\,(a_{n} \neq 0) \to \varepsilon \in \mathbb{R} \to \varepsilon \in \mathbb{R}$\lbl{A1}}
\AXC{$\stackrel{\Pi_{22}}{\vdots}$}\noLine
\UIC{$\varphi_{110}$}
\BIC{$\varphi_{94}$}
\BIC{$\mathrm{DvgP}(a) \to \forall n \in \mathbb{N}\,(a_{n} \neq 0) \to \varepsilon \in \mathbb{R} \to 0 < \varepsilon \to \varphi_{36}$}
\AXC{$\varphi_{150}$\lbl{A1}}
\BIC{$\varphi_{133}$}
\end{ptree}
\noindent$\Pi_{24}$:
\begin{ptree}
\AXC{$\stackrel{\Pi_{5}}{\vdots}$}\noLine
\UIC{$\mathrm{DvgP}(a) \to \forall n \in \mathbb{N}\,(a_{n} \neq 0) \to \varepsilon \in \mathbb{R} \to 0 < \varepsilon \to \varphi_{42}$}
\AXC{$\stackrel{\Pi_{21}}{\vdots}$}\noLine
\UIC{$\varphi_{122}$}
\BIC{$\mathrm{DvgP}(a) \to \forall n \in \mathbb{N}\,(a_{n} \neq 0) \to \varepsilon \in \mathbb{R} \to 0 < \varepsilon \to \varphi_{52}$}
\AXC{$\stackrel{\Pi_{23}}{\vdots}$}\noLine
\UIC{$\varphi_{133}$}
\BIC{$\varphi_{96}$}
\end{ptree}
\noindent$\Pi_{25}$:
\begin{ptree}
\AXC{$y = z \to x = y \to x = z$\lbl{A6}}
\AXC{$(y = z \to x = y \to x = z) \to x = y \to y = z \to x = z$\lbl{A1}}
\BIC{$x = y \to y = z \to x = z$}
\UIC{$\forall z\,(x = y \to y = z \to x = z)$}
\UIC{$\forall y\,\forall z\,(x = y \to y = z \to x = z)$}
\UIC{$\forall x\,\forall y\,\forall z\,(x = y \to y = z \to x = z)$}
\AXC{$\forall x\,\forall y\,\forall z\,(x = y \to y = z \to x = z) \to \varphi_{3}$\lbl{A2}}
\BIC{$\varphi_{3}$}
\end{ptree}
\noindent$\Pi_{26}$:
\begin{ptree}
\AXC{$\forall x\,\forall y\,\forall z\,(x < y \to y < z \to x < z)$\lbl{[O2]}}
\AXC{$\varphi_{6}$\lbl{A2}}
\BIC{$\forall y\,\forall z\,(0 < y \to y < z \to 0 < z)$}
\AXC{$\varphi_{14}$\lbl{A2}}
\BIC{$\forall z\,(0 < \varepsilon^{-1} \to \varepsilon^{-1} < z \to 0 < z)$}
\AXC{$\forall z\,(0 < \varepsilon^{-1} \to \varepsilon^{-1} < z \to 0 < z) \to 0 < \varepsilon^{-1} \to \varepsilon^{-1} < a_{n} \to 0 < a_{n}$\lbl{A2}}
\BIC{$0 < \varepsilon^{-1} \to \varepsilon^{-1} < a_{n} \to 0 < a_{n}$}
\end{ptree}
\noindent$\Pi_{27}$:
\begin{ptree}
\AXC{$\stackrel{\Pi_{8}}{\vdots}$}\noLine
\UIC{$\varphi_{17}$}
\AXC{$\stackrel{\Pi_{26}}{\vdots}$}\noLine
\UIC{$0 < \varepsilon^{-1} \to \varepsilon^{-1} < a_{n} \to 0 < a_{n}$}
\AXC{$(0 < \varepsilon^{-1} \to \varepsilon^{-1} < a_{n} \to 0 < a_{n}) \to \varphi_{17} \to \varphi_{117}$\lbl{A1}}
\BIC{$\varphi_{17} \to \varphi_{117}$}
\BIC{$\varphi_{117}$}
\end{ptree}
\noindent$\Pi_{28}$:
\begin{ptree}
\AXC{$\stackrel{\Pi_{5}}{\vdots}$}\noLine
\UIC{$\mathrm{DvgP}(a) \to \forall n \in \mathbb{N}\,(a_{n} \neq 0) \to \varepsilon \in \mathbb{R} \to 0 < \varepsilon \to \varphi_{42}$}
\AXC{$\stackrel{\Pi_{27}}{\vdots}$}\noLine
\UIC{$\varphi_{117}$}
\BIC{$\mathrm{DvgP}(a) \to \forall n \in \mathbb{N}\,(a_{n} \neq 0) \to \varepsilon \in \mathbb{R} \to 0 < \varepsilon \to \varphi_{27}$}
\end{ptree}
\noindent$\Pi_{29}$:
\begin{ptree}
\AXC{$\forall x \in \mathbb{R}\,(0 < x \to 0 < x^{-1})$\lbl{[F34]}}
\AXC{$\varphi_{41}$\lbl{A2}}
\BIC{$a_{n} \in \mathbb{R} \to 0 < a_{n} \to 0 < a_{n}^{-1}$}
\AXC{$(a_{n} \in \mathbb{R} \to 0 < a_{n} \to 0 < a_{n}^{-1}) \to \varphi_{143}$\lbl{A1}}
\BIC{$\varphi_{143}$}
\end{ptree}
\noindent$\Pi_{30}$:
\begin{ptree}
\AXC{$\stackrel{\Pi_{28}}{\vdots}$}\noLine
\UIC{$\mathrm{DvgP}(a) \to \forall n \in \mathbb{N}\,(a_{n} \neq 0) \to \varepsilon \in \mathbb{R} \to 0 < \varepsilon \to \varphi_{27}$}
\AXC{$\stackrel{\Pi_{13}}{\vdots}$}\noLine
\UIC{$\mathrm{DvgP}(a) \to \forall n \in \mathbb{N}\,(a_{n} \neq 0) \to \varepsilon \in \mathbb{R} \to 0 < \varepsilon \to \varphi_{34}$}
\AXC{$\stackrel{\Pi_{29}}{\vdots}$}\noLine
\UIC{$\varphi_{143}$}
\BIC{$\varphi_{115}$}
\BIC{$\mathrm{DvgP}(a) \to \forall n \in \mathbb{N}\,(a_{n} \neq 0) \to \varepsilon \in \mathbb{R} \to 0 < \varepsilon \to \varphi_{38}$}
\end{ptree}
\noindent$\Pi_{31}$:
\begin{ptree}
\AXC{$\forall x\,(0 \leq x \to \lvert x \rvert = x)$\lbl{[A1]}}
\AXC{$\varphi_{8}$\lbl{A2}}
\BIC{$0 \leq a_{n}^{-1} \to \lvert a_{n}^{-1} \rvert = a_{n}^{-1}$}
\AXC{$(0 \leq a_{n}^{-1} \to \lvert a_{n}^{-1} \rvert = a_{n}^{-1}) \to \varphi_{120}$\lbl{A1}}
\BIC{$\varphi_{120}$}
\end{ptree}
\noindent$\Pi_{32}$:
\begin{ptree}
\AXC{$\stackrel{\Pi_{30}}{\vdots}$}\noLine
\UIC{$\mathrm{DvgP}(a) \to \forall n \in \mathbb{N}\,(a_{n} \neq 0) \to \varepsilon \in \mathbb{R} \to 0 < \varepsilon \to \varphi_{38}$}
\AXC{$\varphi_{118}$\lbl{A1}}
\BIC{$\mathrm{DvgP}(a) \to \forall n \in \mathbb{N}\,(a_{n} \neq 0) \to \varepsilon \in \mathbb{R} \to 0 < \varepsilon \to \varphi_{31}$}
\AXC{$\stackrel{\Pi_{31}}{\vdots}$}\noLine
\UIC{$\varphi_{120}$}
\BIC{$\mathrm{DvgP}(a) \to \forall n \in \mathbb{N}\,(a_{n} \neq 0) \to \varepsilon \in \mathbb{R} \to 0 < \varepsilon \to \varphi_{55}$}
\end{ptree}
\noindent$\Pi_{33}$:
\begin{ptree}
\AXC{$z = w \to \lvert a_{n}^{-1} - 0 \rvert = \lvert z \rvert \to \lvert a_{n}^{-1} - 0 \rvert = \lvert w \rvert$\lbl{A6}}
\UIC{$\varphi_{11}$}
\UIC{$\forall z\,(\varphi_{11})$}
\AXC{$\forall z\,(\varphi_{11}) \to \varphi_{54}$\lbl{A2}}
\BIC{$\varphi_{54}$}
\AXC{$\varphi_{54} \to a_{n}^{-1} - 0 = a_{n}^{-1} \to \varphi_{39}$\lbl{A2}}
\BIC{$a_{n}^{-1} - 0 = a_{n}^{-1} \to \varphi_{39}$}
\end{ptree}
\noindent$\Pi_{34}$:
\begin{ptree}
\AXC{$\forall n \in \mathbb{N}\,(a_{n} \neq 0) \to n \in \mathbb{N} \to a_{n} \neq 0$\lbl{A2}}
\AXC{$\varphi_{4}$\lbl{A1}}
\AXC{$\varphi_{4} \to \varphi_{68}$\lbl{A1}}
\BIC{$\varphi_{68}$}
\BIC{$\mathrm{DvgP}(a) \to \forall n \in \mathbb{N}\,(a_{n} \neq 0) \to n \in \mathbb{N} \to a_{n} \neq 0$}
\AXC{$\varphi_{127}$\lbl{A1}}
\BIC{$\varphi_{107}$}
\end{ptree}
\noindent$\Pi_{35}$:
\begin{ptree}
\AXC{$\forall x \in \mathbb{R}\,(x \neq 0 \to x^{-1} \in \mathbb{R})$\lbl{(R1)}}
\AXC{$\varphi_{47}$\lbl{A2}}
\BIC{$a_{n} \in \mathbb{R} \to a_{n} \neq 0 \to a_{n}^{-1} \in \mathbb{R}$}
\AXC{$(a_{n} \in \mathbb{R} \to a_{n} \neq 0 \to a_{n}^{-1} \in \mathbb{R}) \to \varphi_{145}$\lbl{A1}}
\BIC{$\varphi_{145}$}
\end{ptree}
\noindent$\Pi_{36}$:
\begin{ptree}
\AXC{$\stackrel{\Pi_{13}}{\vdots}$}\noLine
\UIC{$\mathrm{DvgP}(a) \to \forall n \in \mathbb{N}\,(a_{n} \neq 0) \to \varepsilon \in \mathbb{R} \to 0 < \varepsilon \to \varphi_{34}$}
\AXC{$\stackrel{\Pi_{35}}{\vdots}$}\noLine
\UIC{$\varphi_{145}$}
\BIC{$\varphi_{119}$}
\end{ptree}
\noindent$\Pi_{37}$:
\begin{ptree}
\AXC{$\forall x \in \mathbb{R}\,(x - 0 = x)$\lbl{[E28]}}
\AXC{$\varphi_{24}$\lbl{A2}}
\BIC{$a_{n}^{-1} \in \mathbb{R} \to a_{n}^{-1} - 0 = a_{n}^{-1}$}
\AXC{$(a_{n}^{-1} \in \mathbb{R} \to a_{n}^{-1} - 0 = a_{n}^{-1}) \to \varphi_{123}$\lbl{A1}}
\BIC{$\varphi_{123}$}
\end{ptree}
\noindent$\Pi_{38}$:
\begin{ptree}
\AXC{$\mathrm{DvgP}(a) \to \forall n \in \mathbb{N}\,(a_{n} \neq 0) \to \varepsilon \in \mathbb{R} \to \varphi_{30}$\lbl{A1}}
\AXC{$\stackrel{\Pi_{34}}{\vdots}$}\noLine
\UIC{$\varphi_{107}$}
\BIC{$\mathrm{DvgP}(a) \to \forall n \in \mathbb{N}\,(a_{n} \neq 0) \to \varepsilon \in \mathbb{R} \to 0 < \varepsilon \to \varphi_{25}$}
\AXC{$\stackrel{\Pi_{36}}{\vdots}$}\noLine
\UIC{$\varphi_{119}$}
\BIC{$\mathrm{DvgP}(a) \to \forall n \in \mathbb{N}\,(a_{n} \neq 0) \to \varepsilon \in \mathbb{R} \to 0 < \varepsilon \to \varphi_{45}$}
\AXC{$\stackrel{\Pi_{37}}{\vdots}$}\noLine
\UIC{$\varphi_{123}$}
\BIC{$\mathrm{DvgP}(a) \to \forall n \in \mathbb{N}\,(a_{n} \neq 0) \to \varepsilon \in \mathbb{R} \to 0 < \varepsilon \to \varphi_{56}$}
\AXC{$\varphi_{149}$\lbl{A1}}
\BIC{$\lvert a_{n}^{-1} - 0 \rvert = \lvert a_{n}^{-1} - 0 \rvert \to \varphi_{114}$}
\end{ptree}
\noindent$\Pi_{39}$:
\begin{ptree}
\AXC{$x = x$\lbl{A5}}
\UIC{$\forall x\,(x = x)$}
\AXC{$\forall x\,(x = x) \to \lvert a_{n}^{-1} - 0 \rvert = \lvert a_{n}^{-1} - 0 \rvert$\lbl{A2}}
\BIC{$\lvert a_{n}^{-1} - 0 \rvert = \lvert a_{n}^{-1} - 0 \rvert$}
\AXC{$\stackrel{\Pi_{38}}{\vdots}$}\noLine
\UIC{$\lvert a_{n}^{-1} - 0 \rvert = \lvert a_{n}^{-1} - 0 \rvert \to \varphi_{114}$}
\BIC{$\varphi_{114}$}
\end{ptree}
\noindent$\Pi_{40}$:
\begin{ptree}
\AXC{$\stackrel{\Pi_{32}}{\vdots}$}\noLine
\UIC{$\mathrm{DvgP}(a) \to \forall n \in \mathbb{N}\,(a_{n} \neq 0) \to \varepsilon \in \mathbb{R} \to 0 < \varepsilon \to \varphi_{55}$}
\AXC{$\stackrel{\Pi_{33}}{\vdots}$}\noLine
\UIC{$a_{n}^{-1} - 0 = a_{n}^{-1} \to \varphi_{39}$}
\AXC{$\stackrel{\Pi_{39}}{\vdots}$}\noLine
\UIC{$\varphi_{114}$}
\BIC{$\mathrm{DvgP}(a) \to \forall n \in \mathbb{N}\,(a_{n} \neq 0) \to \varepsilon \in \mathbb{R} \to 0 < \varepsilon \to \varphi_{66}$}
\AXC{$\varphi_{156}$\lbl{A1}}
\BIC{$\varphi_{142}$}
\BIC{$\varphi_{108}$}
\end{ptree}
\noindent$\Pi_{41}$:
\begin{ptree}
\AXC{$\stackrel{\Pi_{25}}{\vdots}$}\noLine
\UIC{$\varphi_{3}$}
\AXC{$\varphi_{3} \to \varphi_{44}$\lbl{A2}}
\BIC{$\varphi_{44}$}
\AXC{$\varphi_{44} \to \varphi_{67}$\lbl{A2}}
\BIC{$\varphi_{67}$}
\AXC{$\stackrel{\Pi_{40}}{\vdots}$}\noLine
\UIC{$\varphi_{108}$}
\BIC{$\mathrm{DvgP}(a) \to \forall n \in \mathbb{N}\,(a_{n} \neq 0) \to \varepsilon \in \mathbb{R} \to 0 < \varepsilon \to \varphi_{60}$}
\AXC{$\varphi_{137}$\lbl{A1}}
\BIC{$\varphi_{98}$}
\end{ptree}
\noindent$\Pi_{42}$:
\begin{ptree}
\AXC{$\stackrel{\Pi_{14}}{\vdots}$}\noLine
\UIC{$\forall x\,\forall y\,(x = y \to y = x)$}
\AXC{$\varphi_{13}$\lbl{A2}}
\BIC{$\forall y\,(\lvert a_{n}^{-1} - 0 \rvert = y \to y = \lvert a_{n}^{-1} - 0 \rvert)$}
\AXC{$\varphi_{76}$\lbl{A2}}
\BIC{$\lvert a_{n}^{-1} - 0 \rvert = a_{n}^{-1} \to a_{n}^{-1} = \lvert a_{n}^{-1} - 0 \rvert$}
\AXC{$\stackrel{\Pi_{41}}{\vdots}$}\noLine
\UIC{$\varphi_{98}$}
\BIC{$\mathrm{DvgP}(a) \to \forall n \in \mathbb{N}\,(a_{n} \neq 0) \to \varepsilon \in \mathbb{R} \to 0 < \varepsilon \to \varphi_{61}$}
\AXC{$\varphi_{152}$\lbl{A1}}
\BIC{$\varphi_{136}$}
\end{ptree}
\noindent$\Pi_{43}$:
\begin{ptree}
\AXC{$z = w \to a_{n}^{-1} < z \to a_{n}^{-1} < w$\lbl{A6}}
\UIC{$\forall w\,(z = w \to a_{n}^{-1} < z \to a_{n}^{-1} < w)$}
\UIC{$\forall z\,\forall w\,(z = w \to a_{n}^{-1} < z \to a_{n}^{-1} < w)$}
\AXC{$\varphi_{74}$\lbl{A2}}
\BIC{$\forall w\,((\varepsilon^{-1})^{-1} = w \to a_{n}^{-1} < (\varepsilon^{-1})^{-1} \to a_{n}^{-1} < w)$}
\AXC{$\varphi_{81}$\lbl{A2}}
\BIC{$(\varepsilon^{-1})^{-1} = \varepsilon \to a_{n}^{-1} < (\varepsilon^{-1})^{-1} \to a_{n}^{-1} < \varepsilon$}
\AXC{$\stackrel{\Pi_{24}}{\vdots}$}\noLine
\UIC{$\varphi_{96}$}
\BIC{$\mathrm{DvgP}(a) \to \forall n \in \mathbb{N}\,(a_{n} \neq 0) \to \varepsilon \in \mathbb{R} \to 0 < \varepsilon \to \varphi_{43}$}
\AXC{$\stackrel{\Pi_{42}}{\vdots}$}\noLine
\UIC{$\varphi_{136}$}
\BIC{$\varphi_{101}$}
\end{ptree}
\noindent$\Pi_{44}$:
\begin{ptree}
\AXC{$z = w \to z < \varepsilon \to w < \varepsilon$\lbl{A6}}
\UIC{$\forall w\,(z = w \to z < \varepsilon \to w < \varepsilon)$}
\UIC{$\forall z\,\forall w\,(z = w \to z < \varepsilon \to w < \varepsilon)$}
\AXC{$\varphi_{26}$\lbl{A2}}
\BIC{$\forall w\,(a_{n}^{-1} = w \to a_{n}^{-1} < \varepsilon \to w < \varepsilon)$}
\AXC{$\forall w\,(a_{n}^{-1} = w \to a_{n}^{-1} < \varepsilon \to w < \varepsilon) \to \varphi_{40}$\lbl{A2}}
\BIC{$\varphi_{40}$}
\AXC{$\stackrel{\Pi_{43}}{\vdots}$}\noLine
\UIC{$\varphi_{101}$}
\BIC{$\mathrm{DvgP}(a) \to \forall n \in \mathbb{N}\,(a_{n} \neq 0) \to \varepsilon \in \mathbb{R} \to 0 < \varepsilon \to \varphi_{53}$}
\end{ptree}
\noindent$\Pi_{45}$:
\begin{ptree}
\AXC{$\forall a\,(\mathrm{Seq}(a) \to \varphi_{32})$\lbl{[S5]}}
\AXC{$\forall a\,(\mathrm{Seq}(a) \to \varphi_{32}) \to \mathrm{Seq}(a) \to \varphi_{32}$\lbl{A2}}
\BIC{$\mathrm{Seq}(a) \to \varphi_{32}$}
\AXC{$(\mathrm{Seq}(a) \to \varphi_{32}) \to \varphi_{90}$\lbl{A1}}
\BIC{$\varphi_{90}$}
\end{ptree}
\noindent$\Pi_{46}$:
\begin{ptree}
\AXC{$\varphi_{4}$\lbl{A1}}
\AXC{$\stackrel{\Pi_{10}}{\vdots}$}\noLine
\UIC{$\mathrm{DvgP}(a) \to \forall n \in \mathbb{N}\,(a_{n} \neq 0) \to \mathrm{Seq}(a)$}
\AXC{$\stackrel{\Pi_{45}}{\vdots}$}\noLine
\UIC{$\varphi_{90}$}
\BIC{$\varphi_{4} \to \mathrm{DvgP}(a) \to \varphi_{32}$}
\BIC{$\mathrm{DvgP}(a) \to \varphi_{32}$}
\end{ptree}
\noindent$\Pi_{47}$:
\begin{ptree}
\AXC{$\varphi_{29}$\lbl{A2}}
\AXC{$\stackrel{\Pi_{46}}{\vdots}$}\noLine
\UIC{$\mathrm{DvgP}(a) \to \varphi_{32}$}
\AXC{$\varphi_{106}$\lbl{A1}}
\BIC{$\mathrm{DvgP}(a) \to \forall n \in \mathbb{N}\,(a_{n} \neq 0) \to \varepsilon \in \mathbb{R} \to 0 < \varepsilon \to \varphi_{70}$}
\AXC{$\varphi_{141}$\lbl{A1}}
\BIC{$\varphi_{103}$}
\BIC{$\mathrm{DvgP}(a) \to \forall n \in \mathbb{N}\,(a_{n} \neq 0) \to \varepsilon \in \mathbb{R} \to 0 < \varepsilon \to \varphi_{65}$}
\AXC{$\varphi_{146}$\lbl{A1}}
\BIC{$\varphi_{116}$}
\end{ptree}
\noindent$\Pi_{48}$:
\begin{ptree}
\AXC{$\mathrm{DvgP}(a) \to \forall n \in \mathbb{N}\,(a_{n} \neq 0) \to \varepsilon \in \mathbb{R} \to \varphi_{30}$\lbl{A1}}
\AXC{$\stackrel{\Pi_{47}}{\vdots}$}\noLine
\UIC{$\varphi_{116}$}
\BIC{$\mathrm{DvgP}(a) \to \forall n \in \mathbb{N}\,(a_{n} \neq 0) \to \varepsilon \in \mathbb{R} \to 0 < \varepsilon \to \varphi_{48}$}
\AXC{$\varphi_{132}$\lbl{A1}}
\BIC{$\varphi_{95}$}
\end{ptree}
\noindent$\Pi_{49}$:
\begin{ptree}
\AXC{$\stackrel{\Pi_{14}}{\vdots}$}\noLine
\UIC{$\forall x\,\forall y\,(x = y \to y = x)$}
\AXC{$\varphi_{2}$\lbl{A2}}
\BIC{$\forall y\,(a^{\ast}(n) = y \to y = a^{\ast}(n))$}
\AXC{$\varphi_{51}$\lbl{A2}}
\BIC{$a^{\ast}(n) = a_{n}^{-1} \to a_{n}^{-1} = a^{\ast}(n)$}
\AXC{$\stackrel{\Pi_{48}}{\vdots}$}\noLine
\UIC{$\varphi_{95}$}
\BIC{$\mathrm{DvgP}(a) \to \forall n \in \mathbb{N}\,(a_{n} \neq 0) \to \varepsilon \in \mathbb{R} \to 0 < \varepsilon \to \varphi_{49}$}
\AXC{$\varphi_{153}$\lbl{A1}}
\BIC{$\varphi_{138}$}
\end{ptree}
\noindent$\Pi_{50}$:
\begin{ptree}
\AXC{$\varphi_{125}$\lbl{A3}}
\AXC{$\stackrel{\Pi_{1}}{\vdots}$}\noLine
\UIC{$\varphi_{46}$}
\AXC{$\stackrel{\Pi_{44}}{\vdots}$}\noLine
\UIC{$\mathrm{DvgP}(a) \to \forall n \in \mathbb{N}\,(a_{n} \neq 0) \to \varepsilon \in \mathbb{R} \to 0 < \varepsilon \to \varphi_{53}$}
\AXC{$\stackrel{\Pi_{49}}{\vdots}$}\noLine
\UIC{$\varphi_{138}$}
\BIC{$\varphi_{104}$}
\BIC{$\mathrm{DvgP}(a) \to \forall n \in \mathbb{N}\,(a_{n} \neq 0) \to \varepsilon \in \mathbb{R} \to 0 < \varepsilon \to \varphi_{57}$}
\AXC{$\varphi_{124}$\lbl{A1}}
\BIC{$\varphi_{69} \to n \in \mathbb{N} \to N \leq n \to \lvert a^{\ast}(n) - 0 \rvert < \varepsilon$}
\UIC{$\forall n\,(\varphi_{69} \to n \in \mathbb{N} \to N \leq n \to \lvert a^{\ast}(n) - 0 \rvert < \varepsilon)$}
\AXC{$\varphi_{159}$\lbl{A1}}
\BIC{$\varphi_{125} \to \varphi_{135}$}
\BIC{$\varphi_{135}$}
\end{ptree}
\noindent$\Pi_{51}$:
\begin{ptree}
\AXC{$\mathrm{DvgP}(a) \to \forall n \in \mathbb{N}\,(a_{n} \neq 0) \to \varepsilon \in \mathbb{R} \to 0 < \varepsilon \to \varphi_{58}$\lbl{A1}}
\AXC{$\varphi_{113}$\lbl{A1}}
\BIC{$\mathrm{DvgP}(a) \to \forall n \in \mathbb{N}\,(a_{n} \neq 0) \to \varphi_{37}$}
\AXC{$(\mathrm{DvgP}(a) \to \forall n \in \mathbb{N}\,(a_{n} \neq 0) \to \varphi_{37}) \to \varphi_{128}$\lbl{A1}}
\BIC{$\varphi_{128}$}
\end{ptree}
\noindent$\Pi_{52}$:
\begin{ptree}
\AXC{$\forall N\,(\varphi_{7}) \to \varphi_{7}$\lbl{A2}}
\AXC{$\varphi_{69} \to \forall n\,(\varphi_{69})$\lbl{A4}}
\AXC{$\stackrel{\Pi_{50}}{\vdots}$}\noLine
\UIC{$\varphi_{135}$}
\BIC{$\mathrm{DvgP}(a) \to \forall n \in \mathbb{N}\,(a_{n} \neq 0) \to \varepsilon \in \mathbb{R} \to 0 < \varepsilon \to \varphi_{59}$}
\AXC{$\stackrel{\Pi_{51}}{\vdots}$}\noLine
\UIC{$\varphi_{128}$}
\BIC{$\mathrm{DvgP}(a) \to \forall n \in \mathbb{N}\,(a_{n} \neq 0) \to \varepsilon \in \mathbb{R} \to 0 < \varepsilon \to \varphi_{71}$}
\AXC{$\varphi_{148}$\lbl{A1}}
\BIC{$\varphi_{121}$}
\BIC{$\mathrm{DvgP}(a) \to \forall n \in \mathbb{N}\,(a_{n} \neq 0) \to \varepsilon \in \mathbb{R} \to 0 < \varepsilon \to \varphi_{72}$}
\AXC{$\varphi_{129}$\lbl{A1}}
\BIC{$\varphi_{75} \to \neg (N \in \mathbb{N} \land \forall n \in \mathbb{N}\,(N \leq n \to \varepsilon^{-1} < a_{n}))$}
\UIC{$\forall N\,(\varphi_{75} \to \neg (N \in \mathbb{N} \land \forall n \in \mathbb{N}\,(N \leq n \to \varepsilon^{-1} < a_{n})))$}
\end{ptree}
\noindent$\Pi_{53}$:
\begin{ptree}
\AXC{$\forall x \in \mathbb{R}\,(x \neq 0 \to x^{-1} \in \mathbb{R})$\lbl{(R1)}}
\AXC{$\varphi_{22}$\lbl{A2}}
\BIC{$\varepsilon \in \mathbb{R} \to \varepsilon \neq 0 \to \varepsilon^{-1} \in \mathbb{R}$}
\AXC{$(\varepsilon \in \mathbb{R} \to \varepsilon \neq 0 \to \varepsilon^{-1} \in \mathbb{R}) \to \varphi_{92}$\lbl{A1}}
\BIC{$\varphi_{92}$}
\end{ptree}
\noindent$\Pi_{54}$:
\begin{ptree}
\AXC{$\mathrm{DvgP}(a) \to \mathrm{DvgP}(a)$\lbl{A1}}
\AXC{$(\mathrm{DvgP}(a) \to \mathrm{DvgP}(a)) \to \varphi_{97}$\lbl{A1}}
\BIC{$\varphi_{97}$}
\end{ptree}
\noindent$\Pi_{55}$:
\begin{ptree}
\AXC{$\varphi_{84}$\lbl{A2}}
\AXC{$\forall a\,(\mathrm{DvgP}(a) \leftrightarrow \mathrm{Seq}(a) \land \varphi_{15})$\lbl{Def$\mathrm{DvgP}$}}
\AXC{$\varphi_{86}$\lbl{A2}}
\BIC{$\mathrm{DvgP}(a) \leftrightarrow \mathrm{Seq}(a) \land \varphi_{15}$}
\AXC{$\stackrel{\Pi_{54}}{\vdots}$}\noLine
\UIC{$\varphi_{97}$}
\BIC{$\mathrm{DvgP}(a) \to \forall n \in \mathbb{N}\,(a_{n} \neq 0) \to \varepsilon \in \mathbb{R} \to 0 < \varepsilon \to \varphi_{15}$}
\AXC{$\varphi_{144}$\lbl{A1}}
\BIC{$\varphi_{84} \to \mathrm{DvgP}(a) \to \forall n \in \mathbb{N}\,(a_{n} \neq 0) \to \varepsilon \in \mathbb{R} \to \varphi_{50}$}
\BIC{$\mathrm{DvgP}(a) \to \forall n \in \mathbb{N}\,(a_{n} \neq 0) \to \varepsilon \in \mathbb{R} \to \varphi_{50}$}
\AXC{$\varphi_{140}$\lbl{A1}}
\BIC{$\varphi_{23} \to \varphi_{17} \to \mathrm{DvgP}(a) \to \varphi_{62}$}
\end{ptree}
\noindent$\Pi_{56}$:
\begin{ptree}
\AXC{$\stackrel{\Pi_{8}}{\vdots}$}\noLine
\UIC{$\varphi_{17}$}
\AXC{$\stackrel{\Pi_{17}}{\vdots}$}\noLine
\UIC{$\varphi_{10}$}
\AXC{$\mathrm{DvgP}(a) \to \forall n \in \mathbb{N}\,(a_{n} \neq 0) \to \varepsilon \in \mathbb{R} \to \varepsilon \in \mathbb{R}$\lbl{A1}}
\AXC{$\stackrel{\Pi_{53}}{\vdots}$}\noLine
\UIC{$\varphi_{92}$}
\BIC{$\varphi_{10} \to \varphi_{23}$}
\BIC{$\varphi_{23}$}
\AXC{$\stackrel{\Pi_{55}}{\vdots}$}\noLine
\UIC{$\varphi_{23} \to \varphi_{17} \to \mathrm{DvgP}(a) \to \varphi_{62}$}
\BIC{$\varphi_{17} \to \mathrm{DvgP}(a) \to \varphi_{62}$}
\BIC{$\mathrm{DvgP}(a) \to \varphi_{62}$}
\end{ptree}
\noindent$\Pi_{57}$:
\begin{ptree}
\AXC{$\varphi_{131}$\lbl{A3}}
\AXC{$\stackrel{\Pi_{52}}{\vdots}$}\noLine
\UIC{$\forall N\,(\varphi_{75} \to \neg (N \in \mathbb{N} \land \forall n \in \mathbb{N}\,(N \leq n \to \varepsilon^{-1} < a_{n})))$}
\AXC{$\stackrel{\Pi_{56}}{\vdots}$}\noLine
\UIC{$\mathrm{DvgP}(a) \to \varphi_{62}$}
\AXC{$(\mathrm{DvgP}(a) \to \varphi_{62}) \to \varphi_{160}$\lbl{A1}}
\BIC{$\varphi_{160}$}
\BIC{$\varphi_{131} \to \varphi_{134}$}
\BIC{$\varphi_{134}$}
\end{ptree}
\noindent$\Pi_{58}$:
\begin{ptree}
\AXC{$\varphi_{19}$\lbl{A4}}
\AXC{$\varphi_{105}$\lbl{A3}}
\AXC{$\varphi_{75} \to \forall N\,(\varphi_{75})$\lbl{A4}}
\AXC{$\stackrel{\Pi_{57}}{\vdots}$}\noLine
\UIC{$\varphi_{134}$}
\BIC{$\mathrm{DvgP}(a) \to \forall n \in \mathbb{N}\,(a_{n} \neq 0) \to \varepsilon \in \mathbb{R} \to \varphi_{18}$}
\AXC{$\varphi_{102}$\lbl{A1}}
\BIC{$\mathrm{DvgP}(a) \land \forall n \in \mathbb{N}\,(a_{n} \neq 0) \to \varepsilon \in \mathbb{R} \to \varphi_{18}$}
\UIC{$\forall \varepsilon\,(\mathrm{DvgP}(a) \land \forall n \in \mathbb{N}\,(a_{n} \neq 0) \to \varepsilon \in \mathbb{R} \to \varphi_{18})$}
\AXC{$\varphi_{155}$\lbl{A1}}
\BIC{$\varphi_{147}$}
\BIC{$\varphi_{19} \to \mathrm{DvgP}(a) \to \forall n \in \mathbb{N}\,(a_{n} \neq 0) \to \forall \varepsilon \in \mathbb{R}\,(\varphi_{18})$}
\BIC{$\mathrm{DvgP}(a) \to \forall n \in \mathbb{N}\,(a_{n} \neq 0) \to \forall \varepsilon \in \mathbb{R}\,(\varphi_{18})$}
\end{ptree}
\noindent$\Pi_{59}$:
\begin{ptree}
\AXC{$\stackrel{\Pi_{46}}{\vdots}$}\noLine
\UIC{$\mathrm{DvgP}(a) \to \varphi_{32}$}
\AXC{$\varphi_{82}$\lbl{A1}}
\BIC{$\mathrm{DvgP}(a) \to \forall n \in \mathbb{N}\,(a_{n} \neq 0) \to \mathrm{Seq}(a^{\ast})$}
\AXC{$\varphi_{77}$\lbl{DefConv}}
\AXC{$\varphi_{77} \to \varphi_{83}$\lbl{A2}}
\BIC{$\varphi_{83}$}
\AXC{$\varphi_{83} \to \varphi_{80}$\lbl{A2}}
\BIC{$\varphi_{80}$}
\AXC{$\varphi_{80} \to \varphi_{100}$\lbl{A1}}
\BIC{$\varphi_{100}$}
\BIC{$0 \in \mathbb{R} \to \varphi_{87}$}
\end{ptree}
\noindent 主木:
\begin{ptree}
\AXC{$\stackrel{\Pi_{58}}{\vdots}$}\noLine
\UIC{$\mathrm{DvgP}(a) \to \forall n \in \mathbb{N}\,(a_{n} \neq 0) \to \forall \varepsilon \in \mathbb{R}\,(\varphi_{18})$}
\AXC{$0 \in \mathbb{R}$\lbl{(R1)}}
\AXC{$\stackrel{\Pi_{59}}{\vdots}$}\noLine
\UIC{$0 \in \mathbb{R} \to \varphi_{87}$}
\BIC{$\varphi_{87}$}
\BIC{$\varphi_{1}$}
\UIC{$\forall a\,(\varphi_{1})$}
\end{ptree}
\begin{abbrevlist}
\item $\varphi_{1} :\equiv \mathrm{DvgP}(a) \to \forall n \in \mathbb{N}\,(a_{n} \neq 0) \to \mathrm{Conv}(a^{\ast},0)$
\item $\varphi_{2} :\equiv \forall x\,\forall y\,(x = y \to y = x) \to \forall y\,(a^{\ast}(n) = y \to y = a^{\ast}(n))$
\item $\varphi_{3} :\equiv \forall y\,\forall z\,(\lvert a_{n}^{-1} - 0 \rvert = y \to y = z \to \lvert a_{n}^{-1} - 0 \rvert = z)$
\item $\varphi_{4} :\equiv \mathrm{DvgP}(a) \to \forall n \in \mathbb{N}\,(a_{n} \neq 0) \to \forall n \in \mathbb{N}\,(a_{n} \neq 0)$
\item $\varphi_{5} :\equiv \forall x\,\forall y\,(x \in \mathbb{R} \to y \in \mathbb{R} \to 0 < x \to x < y \to y^{-1} < x^{-1})$
\item $\varphi_{6} :\equiv \forall x\,\forall y\,\forall z\,(x < y \to y < z \to x < z) \to \forall y\,\forall z\,(0 < y \to y < z \to 0 < z)$
\item $\varphi_{7} :\equiv \neg (N \in \mathbb{N} \land \forall n \in \mathbb{N}\,(N \leq n \to \lvert a^{\ast}(n) - 0 \rvert < \varepsilon))$
\item $\varphi_{8} :\equiv \forall x\,(0 \leq x \to \lvert x \rvert = x) \to 0 \leq a_{n}^{-1} \to \lvert a_{n}^{-1} \rvert = a_{n}^{-1}$
\item $\varphi_{9} :\equiv \mathrm{DvgP}(a) \to \forall n \in \mathbb{N}\,(a_{n} \neq 0) \to \varepsilon \in \mathbb{R} \to 0 < \varepsilon \to 0 \neq \varepsilon$
\item $\varphi_{10} :\equiv \mathrm{DvgP}(a) \to \forall n \in \mathbb{N}\,(a_{n} \neq 0) \to \varepsilon \in \mathbb{R} \to 0 < \varepsilon \to \varepsilon \neq 0$
\item $\varphi_{11} :\equiv \forall w\,(z = w \to \lvert a_{n}^{-1} - 0 \rvert = \lvert z \rvert \to \lvert a_{n}^{-1} - 0 \rvert = \lvert w \rvert)$
\item $\varphi_{12} :\equiv \mathrm{DvgP}(a) \to \forall n \in \mathbb{N}\,(a_{n} \neq 0) \to \varepsilon \in \mathbb{R} \to 0 < \varepsilon \to 0 < \varepsilon$
\item $\varphi_{13} :\equiv \forall x\,\forall y\,(x = y \to y = x) \to \forall y\,(\lvert a_{n}^{-1} - 0 \rvert = y \to y = \lvert a_{n}^{-1} - 0 \rvert)$
\item $\varphi_{14} :\equiv \forall y\,\forall z\,(0 < y \to y < z \to 0 < z) \to \forall z\,(0 < \varepsilon^{-1} \to \varepsilon^{-1} < z \to 0 < z)$
\item $\varphi_{15} :\equiv \forall K \in \mathbb{R}\,(0 < K \to \exists N \in \mathbb{N}\,\forall n \in \mathbb{N}\,(N \leq n \to K < a_{n}))$
\item $\varphi_{16} :\equiv N \in \mathbb{N} \land \forall n \in \mathbb{N}\,(N \leq n \to \varepsilon^{-1} < a_{n}) \to n \in \mathbb{N} \to N \leq n \to N \leq n$
\item $\varphi_{17} :\equiv \mathrm{DvgP}(a) \to \forall n \in \mathbb{N}\,(a_{n} \neq 0) \to \varepsilon \in \mathbb{R} \to 0 < \varepsilon \to 0 < \varepsilon^{-1}$
\item $\varphi_{18} :\equiv 0 < \varepsilon \to \exists N \in \mathbb{N}\,\forall n \in \mathbb{N}\,(N \leq n \to \lvert a^{\ast}(n) - 0 \rvert < \varepsilon)$
\item $\varphi_{19} :\equiv \mathrm{DvgP}(a) \land \forall n \in \mathbb{N}\,(a_{n} \neq 0) \to \forall \varepsilon\,(\mathrm{DvgP}(a) \land \forall n \in \mathbb{N}\,(a_{n} \neq 0))$
\item $\varphi_{20} :\equiv \forall x \in \mathbb{R}\,(x \neq 0 \to (x^{-1})^{-1} = x) \to \varepsilon \in \mathbb{R} \to \varepsilon \neq 0 \to (\varepsilon^{-1})^{-1} = \varepsilon$
\item $\varphi_{21} :\equiv \forall n\,(\mathrm{Seq}(a) \to n \in \mathbb{N} \to a_{n} \in \mathbb{R}) \to \mathrm{Seq}(a) \to n \in \mathbb{N} \to a_{n} \in \mathbb{R}$
\item $\varphi_{22} :\equiv \forall x \in \mathbb{R}\,(x \neq 0 \to x^{-1} \in \mathbb{R}) \to \varepsilon \in \mathbb{R} \to \varepsilon \neq 0 \to \varepsilon^{-1} \in \mathbb{R}$
\item $\varphi_{23} :\equiv \mathrm{DvgP}(a) \to \forall n \in \mathbb{N}\,(a_{n} \neq 0) \to \varepsilon \in \mathbb{R} \to 0 < \varepsilon \to \varepsilon^{-1} \in \mathbb{R}$
\item $\varphi_{24} :\equiv \forall x \in \mathbb{R}\,(x - 0 = x) \to a_{n}^{-1} \in \mathbb{R} \to a_{n}^{-1} - 0 = a_{n}^{-1}$
\item $\varphi_{25} :\equiv N \in \mathbb{N} \land \forall n \in \mathbb{N}\,(N \leq n \to \varepsilon^{-1} < a_{n}) \to n \in \mathbb{N} \to N \leq n \to a_{n} \neq 0$
\item $\varphi_{26} :\equiv \forall z\,\forall w\,(z = w \to z < \varepsilon \to w < \varepsilon) \to \forall w\,(a_{n}^{-1} = w \to a_{n}^{-1} < \varepsilon \to w < \varepsilon)$
\item $\varphi_{27} :\equiv N \in \mathbb{N} \land \forall n \in \mathbb{N}\,(N \leq n \to \varepsilon^{-1} < a_{n}) \to n \in \mathbb{N} \to N \leq n \to 0 < a_{n}$
\item $\varphi_{28} :\equiv \forall a\,\forall n\,(\mathrm{Seq}(a) \to n \in \mathbb{N} \to a_{n} \in \mathbb{R}) \to \forall n\,(\mathrm{Seq}(a) \to n \in \mathbb{N} \to a_{n} \in \mathbb{R})$
\item $\varphi_{29} :\equiv \forall n \in \mathbb{N}\,(a^{\ast}(n) = a_{n}^{-1}) \to n \in \mathbb{N} \to a^{\ast}(n) = a_{n}^{-1}$
\item $\varphi_{30} :\equiv 0 < \varepsilon \to N \in \mathbb{N} \land \forall n \in \mathbb{N}\,(N \leq n \to \varepsilon^{-1} < a_{n}) \to n \in \mathbb{N} \to n \in \mathbb{N}$
\item $\varphi_{31} :\equiv N \in \mathbb{N} \land \forall n \in \mathbb{N}\,(N \leq n \to \varepsilon^{-1} < a_{n}) \to n \in \mathbb{N} \to N \leq n \to 0 \leq a_{n}^{-1}$
\item $\varphi_{32} :\equiv \forall n \in \mathbb{N}\,(a_{n} \neq 0) \to \mathrm{Seq}(a^{\ast}) \land \forall n \in \mathbb{N}\,(a^{\ast}(n) = a_{n}^{-1})$
\item $\varphi_{33} :\equiv \forall y\,(\varepsilon^{-1} \in \mathbb{R} \to y \in \mathbb{R} \to 0 < \varepsilon^{-1} \to \varepsilon^{-1} < y \to y^{-1} < (\varepsilon^{-1})^{-1})$
\item $\varphi_{34} :\equiv N \in \mathbb{N} \land \forall n \in \mathbb{N}\,(N \leq n \to \varepsilon^{-1} < a_{n}) \to n \in \mathbb{N} \to N \leq n \to a_{n} \in \mathbb{R}$
\item $\varphi_{35} :\equiv N \in \mathbb{N} \land \forall n \in \mathbb{N}\,(N \leq n \to \varepsilon^{-1} < a_{n}) \to n \in \mathbb{N} \to N \leq n \to \varepsilon^{-1} \in \mathbb{R}$
\item $\varphi_{36} :\equiv N \in \mathbb{N} \land \forall n \in \mathbb{N}\,(N \leq n \to \varepsilon^{-1} < a_{n}) \to n \in \mathbb{N} \to N \leq n \to (\varepsilon^{-1})^{-1} = \varepsilon$
\item $\varphi_{37} :\equiv \varepsilon \in \mathbb{R} \to 0 < \varepsilon \to N \in \mathbb{N} \land \forall n \in \mathbb{N}\,(N \leq n \to \varepsilon^{-1} < a_{n}) \to N \in \mathbb{N}$
\item $\varphi_{38} :\equiv N \in \mathbb{N} \land \forall n \in \mathbb{N}\,(N \leq n \to \varepsilon^{-1} < a_{n}) \to n \in \mathbb{N} \to N \leq n \to 0 < a_{n}^{-1}$
\item $\varphi_{39} :\equiv \lvert a_{n}^{-1} - 0 \rvert = \lvert a_{n}^{-1} - 0 \rvert \to \lvert a_{n}^{-1} - 0 \rvert = \lvert a_{n}^{-1} \rvert$
\item $\varphi_{40} :\equiv a_{n}^{-1} = \lvert a_{n}^{-1} - 0 \rvert \to a_{n}^{-1} < \varepsilon \to \lvert a_{n}^{-1} - 0 \rvert < \varepsilon$
\item $\varphi_{41} :\equiv \forall x \in \mathbb{R}\,(0 < x \to 0 < x^{-1}) \to a_{n} \in \mathbb{R} \to 0 < a_{n} \to 0 < a_{n}^{-1}$
\item $\varphi_{42} :\equiv N \in \mathbb{N} \land \forall n \in \mathbb{N}\,(N \leq n \to \varepsilon^{-1} < a_{n}) \to n \in \mathbb{N} \to N \leq n \to \varepsilon^{-1} < a_{n}$
\item $\varphi_{43} :\equiv N \in \mathbb{N} \land \forall n \in \mathbb{N}\,(N \leq n \to \varepsilon^{-1} < a_{n}) \to n \in \mathbb{N} \to N \leq n \to a_{n}^{-1} < \varepsilon$
\item $\varphi_{44} :\equiv \forall z\,(\lvert a_{n}^{-1} - 0 \rvert = \lvert a_{n}^{-1} \rvert \to \lvert a_{n}^{-1} \rvert = z \to \lvert a_{n}^{-1} - 0 \rvert = z)$
\item $\varphi_{45} :\equiv N \in \mathbb{N} \land \forall n \in \mathbb{N}\,(N \leq n \to \varepsilon^{-1} < a_{n}) \to n \in \mathbb{N} \to N \leq n \to a_{n}^{-1} \in \mathbb{R}$
\item $\varphi_{46} :\equiv a_{n}^{-1} = a^{\ast}(n) \to \lvert a_{n}^{-1} - 0 \rvert < \varepsilon \to \lvert a^{\ast}(n) - 0 \rvert < \varepsilon$
\item $\varphi_{47} :\equiv \forall x \in \mathbb{R}\,(x \neq 0 \to x^{-1} \in \mathbb{R}) \to a_{n} \in \mathbb{R} \to a_{n} \neq 0 \to a_{n}^{-1} \in \mathbb{R}$
\item $\varphi_{48} :\equiv N \in \mathbb{N} \land \forall n \in \mathbb{N}\,(N \leq n \to \varepsilon^{-1} < a_{n}) \to n \in \mathbb{N} \to N \leq n \to a^{\ast}(n) = a_{n}^{-1}$
\item $\varphi_{49} :\equiv N \in \mathbb{N} \land \forall n \in \mathbb{N}\,(N \leq n \to \varepsilon^{-1} < a_{n}) \to n \in \mathbb{N} \to N \leq n \to a_{n}^{-1} = a^{\ast}(n)$
\item $\varphi_{50} :\equiv 0 < \varepsilon \to \varepsilon^{-1} \in \mathbb{R} \to 0 < \varepsilon^{-1} \to \exists N \in \mathbb{N}\,\forall n \in \mathbb{N}\,(N \leq n \to \varepsilon^{-1} < a_{n})$
\item $\varphi_{51} :\equiv \forall y\,(a^{\ast}(n) = y \to y = a^{\ast}(n)) \to a^{\ast}(n) = a_{n}^{-1} \to a_{n}^{-1} = a^{\ast}(n)$
\item $\varphi_{52} :\equiv N \in \mathbb{N} \land \forall n \in \mathbb{N}\,(N \leq n \to \varepsilon^{-1} < a_{n}) \to n \in \mathbb{N} \to N \leq n \to a_{n}^{-1} < (\varepsilon^{-1})^{-1}$
\item $\varphi_{53} :\equiv N \in \mathbb{N} \land \forall n \in \mathbb{N}\,(N \leq n \to \varepsilon^{-1} < a_{n}) \to n \in \mathbb{N} \to N \leq n \to \lvert a_{n}^{-1} - 0 \rvert < \varepsilon$
\item $\varphi_{54} :\equiv \forall w\,(a_{n}^{-1} - 0 = w \to \lvert a_{n}^{-1} - 0 \rvert = \lvert a_{n}^{-1} - 0 \rvert \to \lvert a_{n}^{-1} - 0 \rvert = \lvert w \rvert)$
\item $\varphi_{55} :\equiv N \in \mathbb{N} \land \forall n \in \mathbb{N}\,(N \leq n \to \varepsilon^{-1} < a_{n}) \to n \in \mathbb{N} \to N \leq n \to \lvert a_{n}^{-1} \rvert = a_{n}^{-1}$
\item $\varphi_{56} :\equiv N \in \mathbb{N} \land \forall n \in \mathbb{N}\,(N \leq n \to \varepsilon^{-1} < a_{n}) \to n \in \mathbb{N} \to N \leq n \to a_{n}^{-1} - 0 = a_{n}^{-1}$
\item $\varphi_{57} :\equiv N \in \mathbb{N} \land \forall n \in \mathbb{N}\,(N \leq n \to \varepsilon^{-1} < a_{n}) \to n \in \mathbb{N} \to N \leq n \to \lvert a^{\ast}(n) - 0 \rvert < \varepsilon$
\item $\varphi_{58} :\equiv N \in \mathbb{N} \land \forall n \in \mathbb{N}\,(N \leq n \to \varepsilon^{-1} < a_{n}) \to N \in \mathbb{N} \land \forall n \in \mathbb{N}\,(N \leq n \to \varepsilon^{-1} < a_{n})$
\item $\varphi_{59} :\equiv N \in \mathbb{N} \land \forall n \in \mathbb{N}\,(N \leq n \to \varepsilon^{-1} < a_{n}) \to \forall n \in \mathbb{N}\,(N \leq n \to \lvert a^{\ast}(n) - 0 \rvert < \varepsilon)$
\item $\varphi_{60} :\equiv N \in \mathbb{N} \land \forall n \in \mathbb{N}\,(N \leq n \to \varepsilon^{-1} < a_{n}) \to n \in \mathbb{N} \to N \leq n \to \lvert a_{n}^{-1} - 0 \rvert = a_{n}^{-1}$
\item $\varphi_{61} :\equiv N \in \mathbb{N} \land \forall n \in \mathbb{N}\,(N \leq n \to \varepsilon^{-1} < a_{n}) \to n \in \mathbb{N} \to N \leq n \to a_{n}^{-1} = \lvert a_{n}^{-1} - 0 \rvert$
\item $\varphi_{62} :\equiv \forall n \in \mathbb{N}\,(a_{n} \neq 0) \to \varepsilon \in \mathbb{R} \to 0 < \varepsilon \to \exists N \in \mathbb{N}\,\forall n \in \mathbb{N}\,(N \leq n \to \varepsilon^{-1} < a_{n})$
\item $\varphi_{63} :\equiv N \in \mathbb{N} \land \forall n \in \mathbb{N}\,(N \leq n \to \varepsilon^{-1} < a_{n}) \to n \in \mathbb{N} \to N \leq n \to n \in \mathbb{N} \to N \leq n \to \varepsilon^{-1} < a_{n}$
\item $\varphi_{64} :\equiv N \in \mathbb{N} \land \forall n \in \mathbb{N}\,(N \leq n \to \varepsilon^{-1} < a_{n}) \to n \in \mathbb{N} \to N \leq n \to \forall n \in \mathbb{N}\,(N \leq n \to \varepsilon^{-1} < a_{n})$
\item $\varphi_{65} :\equiv N \in \mathbb{N} \land \forall n \in \mathbb{N}\,(N \leq n \to \varepsilon^{-1} < a_{n}) \to n \in \mathbb{N} \to N \leq n \to n \in \mathbb{N} \to a^{\ast}(n) = a_{n}^{-1}$
\item $\varphi_{66} :\equiv N \in \mathbb{N} \land \forall n \in \mathbb{N}\,(N \leq n \to \varepsilon^{-1} < a_{n}) \to n \in \mathbb{N} \to N \leq n \to \lvert a_{n}^{-1} - 0 \rvert = \lvert a_{n}^{-1} \rvert$
\item $\varphi_{67} :\equiv \lvert a_{n}^{-1} - 0 \rvert = \lvert a_{n}^{-1} \rvert \to \lvert a_{n}^{-1} \rvert = a_{n}^{-1} \to \lvert a_{n}^{-1} - 0 \rvert = a_{n}^{-1}$
\item $\varphi_{68} :\equiv (\forall n \in \mathbb{N}\,(a_{n} \neq 0) \to n \in \mathbb{N} \to a_{n} \neq 0) \to \mathrm{DvgP}(a) \to \forall n \in \mathbb{N}\,(a_{n} \neq 0) \to n \in \mathbb{N} \to a_{n} \neq 0$
\item $\varphi_{69} :\equiv \mathrm{DvgP}(a) \land \forall n \in \mathbb{N}\,(a_{n} \neq 0) \land \varepsilon \in \mathbb{R} \land 0 < \varepsilon \land (N \in \mathbb{N} \land \forall n \in \mathbb{N}\,(N \leq n \to \varepsilon^{-1} < a_{n}))$
\item $\varphi_{70} :\equiv N \in \mathbb{N} \land \forall n \in \mathbb{N}\,(N \leq n \to \varepsilon^{-1} < a_{n}) \to n \in \mathbb{N} \to N \leq n \to \forall n \in \mathbb{N}\,(a^{\ast}(n) = a_{n}^{-1})$
\item $\varphi_{71} :\equiv N \in \mathbb{N} \land \forall n \in \mathbb{N}\,(N \leq n \to \varepsilon^{-1} < a_{n}) \to N \in \mathbb{N} \land \forall n \in \mathbb{N}\,(N \leq n \to \lvert a^{\ast}(n) - 0 \rvert < \varepsilon)$
\item $\varphi_{72} :\equiv N \in \mathbb{N} \land \forall n \in \mathbb{N}\,(N \leq n \to \varepsilon^{-1} < a_{n}) \to \exists N \in \mathbb{N}\,\forall n \in \mathbb{N}\,(N \leq n \to \lvert a^{\ast}(n) - 0 \rvert < \varepsilon)$
\item $\varphi_{73} :\equiv \forall z\,\forall w\,(z = w \to \lvert z - 0 \rvert < \varepsilon \to \lvert w - 0 \rvert < \varepsilon) \to \forall w\,(a_{n}^{-1} = w \to \lvert a_{n}^{-1} - 0 \rvert < \varepsilon \to \lvert w - 0 \rvert < \varepsilon)$
\item $\varphi_{74} :\equiv \forall z\,\forall w\,(z = w \to a_{n}^{-1} < z \to a_{n}^{-1} < w) \to \forall w\,((\varepsilon^{-1})^{-1} = w \to a_{n}^{-1} < (\varepsilon^{-1})^{-1} \to a_{n}^{-1} < w)$
\item $\varphi_{75} :\equiv \mathrm{DvgP}(a) \land \forall n \in \mathbb{N}\,(a_{n} \neq 0) \land \varepsilon \in \mathbb{R} \land 0 < \varepsilon \land \neg \exists N \in \mathbb{N}\,\forall n \in \mathbb{N}\,(N \leq n \to \lvert a^{\ast}(n) - 0 \rvert < \varepsilon)$
\item $\varphi_{76} :\equiv \forall y\,(\lvert a_{n}^{-1} - 0 \rvert = y \to y = \lvert a_{n}^{-1} - 0 \rvert) \to \lvert a_{n}^{-1} - 0 \rvert = a_{n}^{-1} \to a_{n}^{-1} = \lvert a_{n}^{-1} - 0 \rvert$
\item $\varphi_{77} :\equiv \forall a\,\forall \alpha\,(\mathrm{Conv}(a,\alpha) \leftrightarrow \mathrm{Seq}(a) \land \alpha \in \mathbb{R} \land \forall \varepsilon \in \mathbb{R}\,(0 < \varepsilon \to \exists N \in \mathbb{N}\,\forall n \in \mathbb{N}\,(N \leq n \to \lvert a_{n} - \alpha \rvert < \varepsilon)))$
\item $\varphi_{78} :\equiv (\mathrm{DvgP}(a) \leftrightarrow \mathrm{Seq}(a) \land \varphi_{15}) \to \mathrm{DvgP}(a) \to \forall n \in \mathbb{N}\,(a_{n} \neq 0) \to \mathrm{Seq}(a)$
\item $\varphi_{79} :\equiv \forall \varepsilon\,(\mathrm{DvgP}(a) \land \forall n \in \mathbb{N}\,(a_{n} \neq 0)) \to \forall \varepsilon \in \mathbb{R}\,(\varphi_{18})$
\item $\varphi_{80} :\equiv \mathrm{Conv}(a^{\ast},0) \leftrightarrow \mathrm{Seq}(a^{\ast}) \land 0 \in \mathbb{R} \land \forall \varepsilon \in \mathbb{R}\,(\varphi_{18})$
\item $\varphi_{81} :\equiv \forall w\,((\varepsilon^{-1})^{-1} = w \to a_{n}^{-1} < (\varepsilon^{-1})^{-1} \to a_{n}^{-1} < w) \to (\varepsilon^{-1})^{-1} = \varepsilon \to a_{n}^{-1} < (\varepsilon^{-1})^{-1} \to a_{n}^{-1} < \varepsilon$
\item $\varphi_{82} :\equiv (\mathrm{DvgP}(a) \to \varphi_{32}) \to \mathrm{DvgP}(a) \to \forall n \in \mathbb{N}\,(a_{n} \neq 0) \to \mathrm{Seq}(a^{\ast})$
\item $\varphi_{83} :\equiv \forall \alpha\,(\mathrm{Conv}(a^{\ast},\alpha) \leftrightarrow \mathrm{Seq}(a^{\ast}) \land \alpha \in \mathbb{R} \land \forall \varepsilon \in \mathbb{R}\,(0 < \varepsilon \to \exists N \in \mathbb{N}\,\forall n \in \mathbb{N}\,(N \leq n \to \lvert a^{\ast}(n) - \alpha \rvert < \varepsilon)))$
\item $\varphi_{84} :\equiv \varphi_{15} \to \varepsilon^{-1} \in \mathbb{R} \to 0 < \varepsilon^{-1} \to \exists N \in \mathbb{N}\,\forall n \in \mathbb{N}\,(N \leq n \to \varepsilon^{-1} < a_{n})$
\item $\varphi_{85} :\equiv \forall w\,(a_{n}^{-1} = w \to \lvert a_{n}^{-1} - 0 \rvert < \varepsilon \to \lvert w - 0 \rvert < \varepsilon) \to \varphi_{46}$
\item $\varphi_{86} :\equiv \forall a\,(\mathrm{DvgP}(a) \leftrightarrow \mathrm{Seq}(a) \land \varphi_{15}) \to (\mathrm{DvgP}(a) \leftrightarrow \mathrm{Seq}(a) \land \varphi_{15})$
\item $\varphi_{87} :\equiv (\mathrm{DvgP}(a) \to \forall n \in \mathbb{N}\,(a_{n} \neq 0) \to \forall \varepsilon \in \mathbb{R}\,(\varphi_{18})) \to \varphi_{1}$
\item $\varphi_{88} :\equiv \forall N\,(\varphi_{75}) \to \forall N\,(\neg (N \in \mathbb{N} \land \forall n \in \mathbb{N}\,(N \leq n \to \varepsilon^{-1} < a_{n})))$
\item $\varphi_{89} :\equiv \varphi_{33} \to \varepsilon^{-1} \in \mathbb{R} \to a_{n} \in \mathbb{R} \to 0 < \varepsilon^{-1} \to \varepsilon^{-1} < a_{n} \to a_{n}^{-1} < (\varepsilon^{-1})^{-1}$
\item $\varphi_{90} :\equiv (\mathrm{DvgP}(a) \to \forall n \in \mathbb{N}\,(a_{n} \neq 0) \to \mathrm{Seq}(a)) \to \varphi_{4} \to \mathrm{DvgP}(a) \to \varphi_{32}$
\item $\varphi_{91} :\equiv (\mathrm{DvgP}(a) \to \forall n \in \mathbb{N}\,(a_{n} \neq 0) \to \varepsilon \in \mathbb{R} \to \varepsilon \in \mathbb{R}) \to \varphi_{12} \to \varphi_{17}$
\item $\varphi_{92} :\equiv (\mathrm{DvgP}(a) \to \forall n \in \mathbb{N}\,(a_{n} \neq 0) \to \varepsilon \in \mathbb{R} \to \varepsilon \in \mathbb{R}) \to \varphi_{10} \to \varphi_{23}$
\item $\varphi_{93} :\equiv \varphi_{10} \to \mathrm{DvgP}(a) \to \forall n \in \mathbb{N}\,(a_{n} \neq 0) \to \varepsilon \in \mathbb{R} \to 0 < \varepsilon \to \varphi_{35}$
\item $\varphi_{94} :\equiv \varphi_{10} \to \mathrm{DvgP}(a) \to \forall n \in \mathbb{N}\,(a_{n} \neq 0) \to \varepsilon \in \mathbb{R} \to 0 < \varepsilon \to \varphi_{36}$
\item $\varphi_{95} :\equiv (a^{\ast}(n) = a_{n}^{-1} \to a_{n}^{-1} = a^{\ast}(n)) \to \mathrm{DvgP}(a) \to \forall n \in \mathbb{N}\,(a_{n} \neq 0) \to \varepsilon \in \mathbb{R} \to 0 < \varepsilon \to \varphi_{49}$
\item $\varphi_{96} :\equiv ((\varepsilon^{-1})^{-1} = \varepsilon \to a_{n}^{-1} < (\varepsilon^{-1})^{-1} \to a_{n}^{-1} < \varepsilon) \to \mathrm{DvgP}(a) \to \forall n \in \mathbb{N}\,(a_{n} \neq 0) \to \varepsilon \in \mathbb{R} \to 0 < \varepsilon \to \varphi_{43}$
\item $\varphi_{97} :\equiv (\mathrm{DvgP}(a) \leftrightarrow \mathrm{Seq}(a) \land \varphi_{15}) \to \mathrm{DvgP}(a) \to \forall n \in \mathbb{N}\,(a_{n} \neq 0) \to \varepsilon \in \mathbb{R} \to 0 < \varepsilon \to \varphi_{15}$
\item $\varphi_{98} :\equiv (\lvert a_{n}^{-1} - 0 \rvert = a_{n}^{-1} \to a_{n}^{-1} = \lvert a_{n}^{-1} - 0 \rvert) \to \mathrm{DvgP}(a) \to \forall n \in \mathbb{N}\,(a_{n} \neq 0) \to \varepsilon \in \mathbb{R} \to 0 < \varepsilon \to \varphi_{61}$
\item $\varphi_{99} :\equiv (\forall n \in \mathbb{N}\,(N \leq n \to \varepsilon^{-1} < a_{n}) \to n \in \mathbb{N} \to N \leq n \to \varepsilon^{-1} < a_{n}) \to \mathrm{DvgP}(a) \to \forall n \in \mathbb{N}\,(a_{n} \neq 0) \to \varepsilon \in \mathbb{R} \to 0 < \varepsilon \to \varphi_{63}$
\item $\varphi_{100} :\equiv (\mathrm{DvgP}(a) \to \forall n \in \mathbb{N}\,(a_{n} \neq 0) \to \mathrm{Seq}(a^{\ast})) \to 0 \in \mathbb{R} \to \varphi_{87}$
\item $\varphi_{101} :\equiv \varphi_{40} \to \mathrm{DvgP}(a) \to \forall n \in \mathbb{N}\,(a_{n} \neq 0) \to \varepsilon \in \mathbb{R} \to 0 < \varepsilon \to \varphi_{53}$
\item $\varphi_{102} :\equiv (\mathrm{DvgP}(a) \to \forall n \in \mathbb{N}\,(a_{n} \neq 0) \to \varepsilon \in \mathbb{R} \to \varphi_{18}) \to \mathrm{DvgP}(a) \land \forall n \in \mathbb{N}\,(a_{n} \neq 0) \to \varepsilon \in \mathbb{R} \to \varphi_{18}$
\item $\varphi_{103} :\equiv \varphi_{29} \to \mathrm{DvgP}(a) \to \forall n \in \mathbb{N}\,(a_{n} \neq 0) \to \varepsilon \in \mathbb{R} \to 0 < \varepsilon \to \varphi_{65}$
\item $\varphi_{104} :\equiv \varphi_{46} \to \mathrm{DvgP}(a) \to \forall n \in \mathbb{N}\,(a_{n} \neq 0) \to \varepsilon \in \mathbb{R} \to 0 < \varepsilon \to \varphi_{57}$
\item $\varphi_{105} :\equiv \forall \varepsilon\,(\mathrm{DvgP}(a) \land \forall n \in \mathbb{N}\,(a_{n} \neq 0) \to \varepsilon \in \mathbb{R} \to \varphi_{18}) \to \varphi_{79}$
\item $\varphi_{106} :\equiv (\mathrm{DvgP}(a) \to \varphi_{32}) \to \mathrm{DvgP}(a) \to \forall n \in \mathbb{N}\,(a_{n} \neq 0) \to \varepsilon \in \mathbb{R} \to 0 < \varepsilon \to \varphi_{70}$
\item $\varphi_{107} :\equiv (\mathrm{DvgP}(a) \to \forall n \in \mathbb{N}\,(a_{n} \neq 0) \to \varepsilon \in \mathbb{R} \to \varphi_{30}) \to \mathrm{DvgP}(a) \to \forall n \in \mathbb{N}\,(a_{n} \neq 0) \to \varepsilon \in \mathbb{R} \to 0 < \varepsilon \to \varphi_{25}$
\item $\varphi_{108} :\equiv \varphi_{67} \to \mathrm{DvgP}(a) \to \forall n \in \mathbb{N}\,(a_{n} \neq 0) \to \varepsilon \in \mathbb{R} \to 0 < \varepsilon \to \varphi_{60}$
\item $\varphi_{109} :\equiv (\mathrm{DvgP}(a) \to \forall n \in \mathbb{N}\,(a_{n} \neq 0) \to \varepsilon \in \mathbb{R} \to \varepsilon \in \mathbb{R}) \to \varphi_{93}$
\item $\varphi_{110} :\equiv (\mathrm{DvgP}(a) \to \forall n \in \mathbb{N}\,(a_{n} \neq 0) \to \varepsilon \in \mathbb{R} \to \varepsilon \in \mathbb{R}) \to \varphi_{94}$
\item $\varphi_{111} :\equiv (\mathrm{DvgP}(a) \to \forall n \in \mathbb{N}\,(a_{n} \neq 0) \to \varepsilon \in \mathbb{R} \to \varphi_{30}) \to \mathrm{DvgP}(a) \to \forall n \in \mathbb{N}\,(a_{n} \neq 0) \to \varepsilon \in \mathbb{R} \to 0 < \varepsilon \to \varphi_{34}$
\item $\varphi_{112} :\equiv (\mathrm{DvgP}(a) \to \forall n \in \mathbb{N}\,(a_{n} \neq 0) \to \varepsilon \in \mathbb{R} \to 0 < \varepsilon \to \varphi_{16}) \to \mathrm{DvgP}(a) \to \forall n \in \mathbb{N}\,(a_{n} \neq 0) \to \varepsilon \in \mathbb{R} \to 0 < \varepsilon \to \varphi_{42}$
\item $\varphi_{113} :\equiv (\mathrm{DvgP}(a) \to \forall n \in \mathbb{N}\,(a_{n} \neq 0) \to \varepsilon \in \mathbb{R} \to 0 < \varepsilon \to \varphi_{58}) \to \mathrm{DvgP}(a) \to \forall n \in \mathbb{N}\,(a_{n} \neq 0) \to \varphi_{37}$
\item $\varphi_{114} :\equiv (a_{n}^{-1} - 0 = a_{n}^{-1} \to \varphi_{39}) \to \mathrm{DvgP}(a) \to \forall n \in \mathbb{N}\,(a_{n} \neq 0) \to \varepsilon \in \mathbb{R} \to 0 < \varepsilon \to \varphi_{66}$
\item $\varphi_{115} :\equiv (\mathrm{DvgP}(a) \to \forall n \in \mathbb{N}\,(a_{n} \neq 0) \to \varepsilon \in \mathbb{R} \to 0 < \varepsilon \to \varphi_{27}) \to \mathrm{DvgP}(a) \to \forall n \in \mathbb{N}\,(a_{n} \neq 0) \to \varepsilon \in \mathbb{R} \to 0 < \varepsilon \to \varphi_{38}$
\item $\varphi_{116} :\equiv (\mathrm{DvgP}(a) \to \forall n \in \mathbb{N}\,(a_{n} \neq 0) \to \varepsilon \in \mathbb{R} \to \varphi_{30}) \to \mathrm{DvgP}(a) \to \forall n \in \mathbb{N}\,(a_{n} \neq 0) \to \varepsilon \in \mathbb{R} \to 0 < \varepsilon \to \varphi_{48}$
\item $\varphi_{117} :\equiv (\mathrm{DvgP}(a) \to \forall n \in \mathbb{N}\,(a_{n} \neq 0) \to \varepsilon \in \mathbb{R} \to 0 < \varepsilon \to \varphi_{42}) \to \mathrm{DvgP}(a) \to \forall n \in \mathbb{N}\,(a_{n} \neq 0) \to \varepsilon \in \mathbb{R} \to 0 < \varepsilon \to \varphi_{27}$
\item $\varphi_{118} :\equiv (\mathrm{DvgP}(a) \to \forall n \in \mathbb{N}\,(a_{n} \neq 0) \to \varepsilon \in \mathbb{R} \to 0 < \varepsilon \to \varphi_{38}) \to \mathrm{DvgP}(a) \to \forall n \in \mathbb{N}\,(a_{n} \neq 0) \to \varepsilon \in \mathbb{R} \to 0 < \varepsilon \to \varphi_{31}$
\item $\varphi_{119} :\equiv (\mathrm{DvgP}(a) \to \forall n \in \mathbb{N}\,(a_{n} \neq 0) \to \varepsilon \in \mathbb{R} \to 0 < \varepsilon \to \varphi_{25}) \to \mathrm{DvgP}(a) \to \forall n \in \mathbb{N}\,(a_{n} \neq 0) \to \varepsilon \in \mathbb{R} \to 0 < \varepsilon \to \varphi_{45}$
\item $\varphi_{120} :\equiv (\mathrm{DvgP}(a) \to \forall n \in \mathbb{N}\,(a_{n} \neq 0) \to \varepsilon \in \mathbb{R} \to 0 < \varepsilon \to \varphi_{31}) \to \mathrm{DvgP}(a) \to \forall n \in \mathbb{N}\,(a_{n} \neq 0) \to \varepsilon \in \mathbb{R} \to 0 < \varepsilon \to \varphi_{55}$
\item $\varphi_{121} :\equiv (\forall N\,(\varphi_{7}) \to \varphi_{7}) \to \mathrm{DvgP}(a) \to \forall n \in \mathbb{N}\,(a_{n} \neq 0) \to \varepsilon \in \mathbb{R} \to 0 < \varepsilon \to \varphi_{72}$
\item $\varphi_{122} :\equiv (\mathrm{DvgP}(a) \to \forall n \in \mathbb{N}\,(a_{n} \neq 0) \to \varepsilon \in \mathbb{R} \to 0 < \varepsilon \to \varphi_{42}) \to \mathrm{DvgP}(a) \to \forall n \in \mathbb{N}\,(a_{n} \neq 0) \to \varepsilon \in \mathbb{R} \to 0 < \varepsilon \to \varphi_{52}$
\item $\varphi_{123} :\equiv (\mathrm{DvgP}(a) \to \forall n \in \mathbb{N}\,(a_{n} \neq 0) \to \varepsilon \in \mathbb{R} \to 0 < \varepsilon \to \varphi_{45}) \to \mathrm{DvgP}(a) \to \forall n \in \mathbb{N}\,(a_{n} \neq 0) \to \varepsilon \in \mathbb{R} \to 0 < \varepsilon \to \varphi_{56}$
\item $\varphi_{124} :\equiv (\mathrm{DvgP}(a) \to \forall n \in \mathbb{N}\,(a_{n} \neq 0) \to \varepsilon \in \mathbb{R} \to 0 < \varepsilon \to \varphi_{57}) \to \varphi_{69} \to n \in \mathbb{N} \to N \leq n \to \lvert a^{\ast}(n) - 0 \rvert < \varepsilon$
\item $\varphi_{125} :\equiv \forall n\,(\varphi_{69} \to n \in \mathbb{N} \to N \leq n \to \lvert a^{\ast}(n) - 0 \rvert < \varepsilon) \to \forall n\,(\varphi_{69}) \to \forall n \in \mathbb{N}\,(N \leq n \to \lvert a^{\ast}(n) - 0 \rvert < \varepsilon)$
\item $\varphi_{126} :\equiv (\mathrm{DvgP}(a) \to \forall n \in \mathbb{N}\,(a_{n} \neq 0) \to \varepsilon \in \mathbb{R} \to 0 < \varepsilon \to \varphi_{58}) \to \mathrm{DvgP}(a) \to \forall n \in \mathbb{N}\,(a_{n} \neq 0) \to \varepsilon \in \mathbb{R} \to 0 < \varepsilon \to \varphi_{64}$
\item $\varphi_{127} :\equiv (\mathrm{DvgP}(a) \to \forall n \in \mathbb{N}\,(a_{n} \neq 0) \to n \in \mathbb{N} \to a_{n} \neq 0) \to \varphi_{107}$
\item $\varphi_{128} :\equiv (\mathrm{DvgP}(a) \to \forall n \in \mathbb{N}\,(a_{n} \neq 0) \to \varepsilon \in \mathbb{R} \to 0 < \varepsilon \to \varphi_{59}) \to \mathrm{DvgP}(a) \to \forall n \in \mathbb{N}\,(a_{n} \neq 0) \to \varepsilon \in \mathbb{R} \to 0 < \varepsilon \to \varphi_{71}$
\item $\varphi_{129} :\equiv (\mathrm{DvgP}(a) \to \forall n \in \mathbb{N}\,(a_{n} \neq 0) \to \varepsilon \in \mathbb{R} \to 0 < \varepsilon \to \varphi_{72}) \to \varphi_{75} \to \neg (N \in \mathbb{N} \land \forall n \in \mathbb{N}\,(N \leq n \to \varepsilon^{-1} < a_{n}))$
\item $\varphi_{130} :\equiv (\mathrm{Seq}(a) \to n \in \mathbb{N} \to a_{n} \in \mathbb{R}) \to (\mathrm{DvgP}(a) \to \forall n \in \mathbb{N}\,(a_{n} \neq 0) \to \mathrm{Seq}(a)) \to \varphi_{111}$
\item $\varphi_{131} :\equiv \forall N\,(\varphi_{75} \to \neg (N \in \mathbb{N} \land \forall n \in \mathbb{N}\,(N \leq n \to \varepsilon^{-1} < a_{n}))) \to \varphi_{88}$
\item $\varphi_{132} :\equiv (\mathrm{DvgP}(a) \to \forall n \in \mathbb{N}\,(a_{n} \neq 0) \to \varepsilon \in \mathbb{R} \to 0 < \varepsilon \to \varphi_{48}) \to \varphi_{95}$
\item $\varphi_{133} :\equiv (\mathrm{DvgP}(a) \to \forall n \in \mathbb{N}\,(a_{n} \neq 0) \to \varepsilon \in \mathbb{R} \to 0 < \varepsilon \to \varphi_{52}) \to \varphi_{96}$
\item $\varphi_{134} :\equiv (\varphi_{75} \to \forall N\,(\varphi_{75})) \to \mathrm{DvgP}(a) \to \forall n \in \mathbb{N}\,(a_{n} \neq 0) \to \varepsilon \in \mathbb{R} \to \varphi_{18}$
\item $\varphi_{135} :\equiv (\varphi_{69} \to \forall n\,(\varphi_{69})) \to \mathrm{DvgP}(a) \to \forall n \in \mathbb{N}\,(a_{n} \neq 0) \to \varepsilon \in \mathbb{R} \to 0 < \varepsilon \to \varphi_{59}$
\item $\varphi_{136} :\equiv (\mathrm{DvgP}(a) \to \forall n \in \mathbb{N}\,(a_{n} \neq 0) \to \varepsilon \in \mathbb{R} \to 0 < \varepsilon \to \varphi_{43}) \to \varphi_{101}$
\item $\varphi_{137} :\equiv (\mathrm{DvgP}(a) \to \forall n \in \mathbb{N}\,(a_{n} \neq 0) \to \varepsilon \in \mathbb{R} \to 0 < \varepsilon \to \varphi_{60}) \to \varphi_{98}$
\item $\varphi_{138} :\equiv (\mathrm{DvgP}(a) \to \forall n \in \mathbb{N}\,(a_{n} \neq 0) \to \varepsilon \in \mathbb{R} \to 0 < \varepsilon \to \varphi_{53}) \to \varphi_{104}$
\item $\varphi_{139} :\equiv (\mathrm{DvgP}(a) \to \forall n \in \mathbb{N}\,(a_{n} \neq 0) \to \varepsilon \in \mathbb{R} \to 0 < \varepsilon \to \varphi_{64}) \to \varphi_{99}$
\item $\varphi_{140} :\equiv (\mathrm{DvgP}(a) \to \forall n \in \mathbb{N}\,(a_{n} \neq 0) \to \varepsilon \in \mathbb{R} \to \varphi_{50}) \to \varphi_{23} \to \varphi_{17} \to \mathrm{DvgP}(a) \to \varphi_{62}$
\item $\varphi_{141} :\equiv (\mathrm{DvgP}(a) \to \forall n \in \mathbb{N}\,(a_{n} \neq 0) \to \varepsilon \in \mathbb{R} \to 0 < \varepsilon \to \varphi_{70}) \to \varphi_{103}$
\item $\varphi_{142} :\equiv (\mathrm{DvgP}(a) \to \forall n \in \mathbb{N}\,(a_{n} \neq 0) \to \varepsilon \in \mathbb{R} \to 0 < \varepsilon \to \varphi_{55}) \to \varphi_{108}$
\item $\varphi_{143} :\equiv (\mathrm{DvgP}(a) \to \forall n \in \mathbb{N}\,(a_{n} \neq 0) \to \varepsilon \in \mathbb{R} \to 0 < \varepsilon \to \varphi_{34}) \to \varphi_{115}$
\item $\varphi_{144} :\equiv (\mathrm{DvgP}(a) \to \forall n \in \mathbb{N}\,(a_{n} \neq 0) \to \varepsilon \in \mathbb{R} \to 0 < \varepsilon \to \varphi_{15}) \to \varphi_{84} \to \mathrm{DvgP}(a) \to \forall n \in \mathbb{N}\,(a_{n} \neq 0) \to \varepsilon \in \mathbb{R} \to \varphi_{50}$
\item $\varphi_{145} :\equiv (\mathrm{DvgP}(a) \to \forall n \in \mathbb{N}\,(a_{n} \neq 0) \to \varepsilon \in \mathbb{R} \to 0 < \varepsilon \to \varphi_{34}) \to \varphi_{119}$
\item $\varphi_{146} :\equiv (\mathrm{DvgP}(a) \to \forall n \in \mathbb{N}\,(a_{n} \neq 0) \to \varepsilon \in \mathbb{R} \to 0 < \varepsilon \to \varphi_{65}) \to \varphi_{116}$
\item $\varphi_{147} :\equiv \varphi_{105} \to \varphi_{19} \to \mathrm{DvgP}(a) \to \forall n \in \mathbb{N}\,(a_{n} \neq 0) \to \forall \varepsilon \in \mathbb{R}\,(\varphi_{18})$
\item $\varphi_{148} :\equiv (\mathrm{DvgP}(a) \to \forall n \in \mathbb{N}\,(a_{n} \neq 0) \to \varepsilon \in \mathbb{R} \to 0 < \varepsilon \to \varphi_{71}) \to \varphi_{121}$
\item $\varphi_{149} :\equiv (\mathrm{DvgP}(a) \to \forall n \in \mathbb{N}\,(a_{n} \neq 0) \to \varepsilon \in \mathbb{R} \to 0 < \varepsilon \to \varphi_{56}) \to \lvert a_{n}^{-1} - 0 \rvert = \lvert a_{n}^{-1} - 0 \rvert \to \varphi_{114}$
\item $\varphi_{150} :\equiv (\mathrm{DvgP}(a) \to \forall n \in \mathbb{N}\,(a_{n} \neq 0) \to \varepsilon \in \mathbb{R} \to 0 < \varepsilon \to \varphi_{36}) \to \varphi_{133}$
\item $\varphi_{151} :\equiv (\mathrm{DvgP}(a) \to \forall n \in \mathbb{N}\,(a_{n} \neq 0) \to \varepsilon \in \mathbb{R} \to 0 < \varepsilon \to \varphi_{34}) \to \varphi_{17} \to \varphi_{122}$
\item $\varphi_{152} :\equiv (\mathrm{DvgP}(a) \to \forall n \in \mathbb{N}\,(a_{n} \neq 0) \to \varepsilon \in \mathbb{R} \to 0 < \varepsilon \to \varphi_{61}) \to \varphi_{136}$
\item $\varphi_{153} :\equiv (\mathrm{DvgP}(a) \to \forall n \in \mathbb{N}\,(a_{n} \neq 0) \to \varepsilon \in \mathbb{R} \to 0 < \varepsilon \to \varphi_{49}) \to \varphi_{138}$
\item $\varphi_{154} :\equiv (\mathrm{DvgP}(a) \to \forall n \in \mathbb{N}\,(a_{n} \neq 0) \to \varepsilon \in \mathbb{R} \to 0 < \varepsilon \to \varphi_{63}) \to (\mathrm{DvgP}(a) \to \forall n \in \mathbb{N}\,(a_{n} \neq 0) \to \varepsilon \in \mathbb{R} \to \varphi_{30}) \to \varphi_{112}$
\item $\varphi_{155} :\equiv \forall \varepsilon\,(\mathrm{DvgP}(a) \land \forall n \in \mathbb{N}\,(a_{n} \neq 0) \to \varepsilon \in \mathbb{R} \to \varphi_{18}) \to \varphi_{147}$
\item $\varphi_{156} :\equiv (\mathrm{DvgP}(a) \to \forall n \in \mathbb{N}\,(a_{n} \neq 0) \to \varepsilon \in \mathbb{R} \to 0 < \varepsilon \to \varphi_{66}) \to \varphi_{142}$
\item $\varphi_{157} :\equiv (\mathrm{DvgP}(a) \to \forall n \in \mathbb{N}\,(a_{n} \neq 0) \to \varepsilon \in \mathbb{R} \to 0 < \varepsilon \to \varphi_{35}) \to \varphi_{151}$
\item $\varphi_{158} :\equiv (\varepsilon^{-1} \in \mathbb{R} \to a_{n} \in \mathbb{R} \to 0 < \varepsilon^{-1} \to \varepsilon^{-1} < a_{n} \to a_{n}^{-1} < (\varepsilon^{-1})^{-1}) \to \varphi_{157}$
\item $\varphi_{159} :\equiv \forall n\,(\varphi_{69} \to n \in \mathbb{N} \to N \leq n \to \lvert a^{\ast}(n) - 0 \rvert < \varepsilon) \to \varphi_{125} \to \varphi_{135}$
\item $\varphi_{160} :\equiv \forall N\,(\varphi_{75} \to \neg (N \in \mathbb{N} \land \forall n \in \mathbb{N}\,(N \leq n \to \varepsilon^{-1} < a_{n}))) \to \varphi_{131} \to \varphi_{134}$
\end{abbrevlist}
\end{fullproof}

% 自動生成 (prover/build.js): RecipDvg  ノード数 287
\begin{fullproof}{定理2.17}{$\forall a\,(\forall n \in \mathbb{N}\,(0 < a_{n}) \to \mathrm{Conv}(a,0) \to \mathrm{DvgP}(a^{\ast}))$}{287}\phantomsection\label{tree:RecipDvg}
\noindent$\Pi_{1}$:
\begin{ptree}
\AXC{$z = w \to z < a_{n}^{-1} \to w < a_{n}^{-1}$\lbl{A6}}
\UIC{$\forall w\,(z = w \to z < a_{n}^{-1} \to w < a_{n}^{-1})$}
\UIC{$\forall z\,\forall w\,(z = w \to z < a_{n}^{-1} \to w < a_{n}^{-1})$}
\AXC{$\forall z\,\forall w\,(z = w \to z < a_{n}^{-1} \to w < a_{n}^{-1}) \to \varphi_{16}$\lbl{A2}}
\BIC{$\varphi_{16}$}
\AXC{$\varphi_{83}$\lbl{A2}}
\BIC{$(K^{-1})^{-1} = K \to (K^{-1})^{-1} < a_{n}^{-1} \to K < a_{n}^{-1}$}
\end{ptree}
\noindent$\Pi_{2}$:
\begin{ptree}
\AXC{$\varphi_{86}$\lbl{A1}}
\AXC{$\varphi_{87}$\lbl{A1}}
\AXC{$\varphi_{54}$\lbl{A2}}
\AXC{$\varphi_{110}$\lbl{A1}}
\AXC{$\varphi_{110} \to \varphi_{113}$\lbl{A1}}
\BIC{$\varphi_{113}$}
\AXC{$\varphi_{113} \to \varphi_{54} \to \varphi_{112}$\lbl{A1}}
\BIC{$\varphi_{54} \to \varphi_{112}$}
\BIC{$\varphi_{112}$}
\AXC{$\varphi_{112} \to \varphi_{87} \to \varphi_{86} \to \varphi_{101}$\lbl{A1}}
\BIC{$\varphi_{87} \to \varphi_{86} \to \varphi_{101}$}
\BIC{$\varphi_{86} \to \varphi_{101}$}
\BIC{$\varphi_{101}$}
\end{ptree}
\noindent$\Pi_{3}$:
\begin{ptree}
\AXC{$y = z \to x = y \to x = z$\lbl{A6}}
\AXC{$(y = z \to x = y \to x = z) \to x = y \to y = z \to x = z$\lbl{A1}}
\BIC{$x = y \to y = z \to x = z$}
\UIC{$\forall z\,(x = y \to y = z \to x = z)$}
\UIC{$\forall y\,\forall z\,(x = y \to y = z \to x = z)$}
\UIC{$\forall x\,\forall y\,\forall z\,(x = y \to y = z \to x = z)$}
\AXC{$\varphi_{35}$\lbl{A2}}
\BIC{$\forall y\,\forall z\,(\lvert a_{n} - 0 \rvert = y \to y = z \to \lvert a_{n} - 0 \rvert = z)$}
\end{ptree}
\noindent$\Pi_{4}$:
\begin{ptree}
\AXC{$\stackrel{\Pi_{3}}{\vdots}$}\noLine
\UIC{$\forall y\,\forall z\,(\lvert a_{n} - 0 \rvert = y \to y = z \to \lvert a_{n} - 0 \rvert = z)$}
\AXC{$\forall y\,\forall z\,(\lvert a_{n} - 0 \rvert = y \to y = z \to \lvert a_{n} - 0 \rvert = z) \to \varphi_{27}$\lbl{A2}}
\BIC{$\varphi_{27}$}
\AXC{$\varphi_{95}$\lbl{A2}}
\BIC{$\lvert a_{n} - 0 \rvert = \lvert a_{n} \rvert \to \lvert a_{n} \rvert = a_{n} \to \lvert a_{n} - 0 \rvert = a_{n}$}
\end{ptree}
\noindent$\Pi_{5}$:
\begin{ptree}
\AXC{$\varphi_{87}$\lbl{A1}}
\AXC{$\forall n \in \mathbb{N}\,(0 < a_{n}) \to n \in \mathbb{N} \to 0 < a_{n}$\lbl{A2}}
\AXC{$(\forall n \in \mathbb{N}\,(0 < a_{n}) \to n \in \mathbb{N} \to 0 < a_{n}) \to \varphi_{87} \to \varphi_{89}$\lbl{A1}}
\BIC{$\varphi_{87} \to \varphi_{89}$}
\BIC{$\varphi_{89}$}
\end{ptree}
\noindent$\Pi_{6}$:
\begin{ptree}
\AXC{$\stackrel{\Pi_{5}}{\vdots}$}\noLine
\UIC{$\varphi_{89}$}
\AXC{$\varphi_{89} \to \varphi_{88}$\lbl{A1}}
\BIC{$\varphi_{88}$}
\end{ptree}
\noindent$\Pi_{7}$:
\begin{ptree}
\AXC{$\forall x\,(0 \leq x \to \lvert x \rvert = x)$\lbl{[A1]}}
\AXC{$\forall x\,(0 \leq x \to \lvert x \rvert = x) \to 0 \leq a_{n} \to \lvert a_{n} \rvert = a_{n}$\lbl{A2}}
\BIC{$0 \leq a_{n} \to \lvert a_{n} \rvert = a_{n}$}
\AXC{$(0 \leq a_{n} \to \lvert a_{n} \rvert = a_{n}) \to \varphi_{88} \to \varphi_{97}$\lbl{A1}}
\BIC{$\varphi_{88} \to \varphi_{97}$}
\end{ptree}
\noindent$\Pi_{8}$:
\begin{ptree}
\AXC{$\varphi_{2}$\lbl{A1}}
\AXC{$\forall a\,\forall \alpha\,(\mathrm{Conv}(a,\alpha) \to \mathrm{Seq}(a))$\lbl{[C1]}}
\AXC{$\varphi_{3}$\lbl{A2}}
\BIC{$\forall \alpha\,(\mathrm{Conv}(a,\alpha) \to \mathrm{Seq}(a))$}
\AXC{$\varphi_{1}$\lbl{A2}}
\BIC{$\mathrm{Conv}(a,0) \to \mathrm{Seq}(a)$}
\AXC{$\varphi_{77}$\lbl{A1}}
\BIC{$\varphi_{2} \to \forall n \in \mathbb{N}\,(0 < a_{n}) \to \mathrm{Conv}(a,0) \to \mathrm{Seq}(a)$}
\BIC{$\forall n \in \mathbb{N}\,(0 < a_{n}) \to \mathrm{Conv}(a,0) \to \mathrm{Seq}(a)$}
\end{ptree}
\noindent$\Pi_{9}$:
\begin{ptree}
\AXC{$\forall a\,\forall n\,(\mathrm{Seq}(a) \to n \in \mathbb{N} \to a_{n} \in \mathbb{R})$\lbl{[S1]}}
\AXC{$\varphi_{44}$\lbl{A2}}
\BIC{$\forall n\,(\mathrm{Seq}(a) \to n \in \mathbb{N} \to a_{n} \in \mathbb{R})$}
\AXC{$\varphi_{36}$\lbl{A2}}
\BIC{$\mathrm{Seq}(a) \to n \in \mathbb{N} \to a_{n} \in \mathbb{R}$}
\AXC{$(\mathrm{Seq}(a) \to n \in \mathbb{N} \to a_{n} \in \mathbb{R}) \to \varphi_{124}$\lbl{A1}}
\BIC{$\varphi_{124}$}
\end{ptree}
\noindent$\Pi_{10}$:
\begin{ptree}
\AXC{$\varphi_{87}$\lbl{A1}}
\AXC{$\stackrel{\Pi_{8}}{\vdots}$}\noLine
\UIC{$\forall n \in \mathbb{N}\,(0 < a_{n}) \to \mathrm{Conv}(a,0) \to \mathrm{Seq}(a)$}
\AXC{$\stackrel{\Pi_{9}}{\vdots}$}\noLine
\UIC{$\varphi_{124}$}
\BIC{$\varphi_{87} \to \varphi_{92}$}
\BIC{$\varphi_{92}$}
\end{ptree}
\noindent$\Pi_{11}$:
\begin{ptree}
\AXC{$\forall x \in \mathbb{R}\,(x - 0 = x)$\lbl{[E28]}}
\AXC{$\forall x \in \mathbb{R}\,(x - 0 = x) \to a_{n} \in \mathbb{R} \to a_{n} - 0 = a_{n}$\lbl{A2}}
\BIC{$a_{n} \in \mathbb{R} \to a_{n} - 0 = a_{n}$}
\AXC{$(a_{n} \in \mathbb{R} \to a_{n} - 0 = a_{n}) \to \varphi_{92} \to \varphi_{98}$\lbl{A1}}
\BIC{$\varphi_{92} \to \varphi_{98}$}
\end{ptree}
\noindent$\Pi_{12}$:
\begin{ptree}
\AXC{$\stackrel{\Pi_{10}}{\vdots}$}\noLine
\UIC{$\varphi_{92}$}
\AXC{$\stackrel{\Pi_{11}}{\vdots}$}\noLine
\UIC{$\varphi_{92} \to \varphi_{98}$}
\BIC{$\varphi_{98}$}
\AXC{$\varphi_{98} \to \lvert a_{n} - 0 \rvert = \lvert a_{n} - 0 \rvert \to \varphi_{52} \to \varphi_{105}$\lbl{A1}}
\BIC{$\lvert a_{n} - 0 \rvert = \lvert a_{n} - 0 \rvert \to \varphi_{52} \to \varphi_{105}$}
\end{ptree}
\noindent$\Pi_{13}$:
\begin{ptree}
\AXC{$x = x$\lbl{A5}}
\UIC{$\forall x\,(x = x)$}
\AXC{$\forall x\,(x = x) \to \lvert a_{n} - 0 \rvert = \lvert a_{n} - 0 \rvert$\lbl{A2}}
\BIC{$\lvert a_{n} - 0 \rvert = \lvert a_{n} - 0 \rvert$}
\AXC{$\stackrel{\Pi_{12}}{\vdots}$}\noLine
\UIC{$\lvert a_{n} - 0 \rvert = \lvert a_{n} - 0 \rvert \to \varphi_{52} \to \varphi_{105}$}
\BIC{$\varphi_{52} \to \varphi_{105}$}
\end{ptree}
\noindent$\Pi_{14}$:
\begin{ptree}
\AXC{$z = w \to \lvert a_{n} - 0 \rvert = \lvert z \rvert \to \lvert a_{n} - 0 \rvert = \lvert w \rvert$\lbl{A6}}
\UIC{$\forall w\,(z = w \to \lvert a_{n} - 0 \rvert = \lvert z \rvert \to \lvert a_{n} - 0 \rvert = \lvert w \rvert)$}
\UIC{$\varphi_{9}$}
\AXC{$\varphi_{9} \to \varphi_{41}$\lbl{A2}}
\BIC{$\varphi_{41}$}
\AXC{$\varphi_{41} \to \varphi_{52}$\lbl{A2}}
\BIC{$\varphi_{52}$}
\AXC{$\stackrel{\Pi_{13}}{\vdots}$}\noLine
\UIC{$\varphi_{52} \to \varphi_{105}$}
\BIC{$\varphi_{105}$}
\AXC{$\varphi_{105} \to \varphi_{97} \to \varphi_{120}$\lbl{A1}}
\BIC{$\varphi_{97} \to \varphi_{120}$}
\end{ptree}
\noindent$\Pi_{15}$:
\begin{ptree}
\AXC{$\stackrel{\Pi_{4}}{\vdots}$}\noLine
\UIC{$\lvert a_{n} - 0 \rvert = \lvert a_{n} \rvert \to \lvert a_{n} \rvert = a_{n} \to \lvert a_{n} - 0 \rvert = a_{n}$}
\AXC{$\stackrel{\Pi_{6}}{\vdots}$}\noLine
\UIC{$\varphi_{88}$}
\AXC{$\stackrel{\Pi_{7}}{\vdots}$}\noLine
\UIC{$\varphi_{88} \to \varphi_{97}$}
\BIC{$\varphi_{97}$}
\AXC{$\stackrel{\Pi_{14}}{\vdots}$}\noLine
\UIC{$\varphi_{97} \to \varphi_{120}$}
\BIC{$\varphi_{120}$}
\BIC{$\varphi_{100}$}
\AXC{$\varphi_{100} \to \varphi_{101} \to \varphi_{118}$\lbl{A1}}
\BIC{$\varphi_{101} \to \varphi_{118}$}
\end{ptree}
\noindent$\Pi_{16}$:
\begin{ptree}
\AXC{$x = x$\lbl{A5}}
\AXC{$x = y \to x = x \to y = x$\lbl{A6}}
\AXC{$(x = y \to x = x \to y = x) \to x = x \to x = y \to y = x$\lbl{A1}}
\BIC{$x = x \to x = y \to y = x$}
\BIC{$x = y \to y = x$}
\UIC{$\forall y\,(x = y \to y = x)$}
\UIC{$\forall x\,\forall y\,(x = y \to y = x)$}
\end{ptree}
\noindent$\Pi_{17}$:
\begin{ptree}
\AXC{$\stackrel{\Pi_{16}}{\vdots}$}\noLine
\UIC{$\forall x\,\forall y\,(x = y \to y = x)$}
\AXC{$\forall x\,\forall y\,(x = y \to y = x) \to \forall y\,(K = y \to y = K)$\lbl{A2}}
\BIC{$\forall y\,(K = y \to y = K)$}
\AXC{$\forall y\,(K = y \to y = K) \to K = 0 \to 0 = K$\lbl{A2}}
\BIC{$K = 0 \to 0 = K$}
\end{ptree}
\noindent$\Pi_{18}$:
\begin{ptree}
\AXC{$\forall x\,\forall y\,(x < y \to x \neq y)$\lbl{[F26]}}
\AXC{$\forall x\,\forall y\,(x < y \to x \neq y) \to \forall y\,(0 < y \to 0 \neq y)$\lbl{A2}}
\BIC{$\forall y\,(0 < y \to 0 \neq y)$}
\AXC{$\forall y\,(0 < y \to 0 \neq y) \to 0 < K \to 0 \neq K$\lbl{A2}}
\BIC{$0 < K \to 0 \neq K$}
\end{ptree}
\noindent$\Pi_{19}$:
\begin{ptree}
\AXC{$\stackrel{\Pi_{17}}{\vdots}$}\noLine
\UIC{$K = 0 \to 0 = K$}
\AXC{$\varphi_{21}$\lbl{A1}}
\AXC{$\stackrel{\Pi_{18}}{\vdots}$}\noLine
\UIC{$0 < K \to 0 \neq K$}
\AXC{$(0 < K \to 0 \neq K) \to \varphi_{21} \to \varphi_{18}$\lbl{A1}}
\BIC{$\varphi_{21} \to \varphi_{18}$}
\BIC{$\varphi_{18}$}
\AXC{$\varphi_{18} \to (K = 0 \to 0 = K) \to \varphi_{19}$\lbl{A1}}
\BIC{$(K = 0 \to 0 = K) \to \varphi_{19}$}
\BIC{$\varphi_{19}$}
\end{ptree}
\noindent$\Pi_{20}$:
\begin{ptree}
\AXC{$\stackrel{\Pi_{19}}{\vdots}$}\noLine
\UIC{$\varphi_{19}$}
\AXC{$\varphi_{11}$\lbl{A1}}
\AXC{$\forall x \in \mathbb{R}\,(x \neq 0 \to x^{-1} \in \mathbb{R})$\lbl{(R1)}}
\AXC{$\varphi_{28}$\lbl{A2}}
\BIC{$K \in \mathbb{R} \to K \neq 0 \to K^{-1} \in \mathbb{R}$}
\AXC{$(K \in \mathbb{R} \to K \neq 0 \to K^{-1} \in \mathbb{R}) \to \varphi_{11} \to \varphi_{19} \to \varphi_{91}$\lbl{A1}}
\BIC{$\varphi_{11} \to \varphi_{19} \to \varphi_{91}$}
\BIC{$\varphi_{19} \to \varphi_{91}$}
\BIC{$\varphi_{91}$}
\end{ptree}
\noindent$\Pi_{21}$:
\begin{ptree}
\AXC{$\stackrel{\Pi_{20}}{\vdots}$}\noLine
\UIC{$\varphi_{91}$}
\AXC{$\stackrel{\Pi_{10}}{\vdots}$}\noLine
\UIC{$\varphi_{92}$}
\AXC{$\varphi_{12}$\lbl{[E19]}}
\AXC{$\varphi_{12} \to \varphi_{47}$\lbl{A2}}
\BIC{$\varphi_{47}$}
\AXC{$\varphi_{47} \to a_{n} \in \mathbb{R} \to \varphi_{37}$\lbl{A2}}
\BIC{$a_{n} \in \mathbb{R} \to \varphi_{37}$}
\AXC{$(a_{n} \in \mathbb{R} \to \varphi_{37}) \to \varphi_{92} \to \varphi_{91} \to \varphi_{89} \to \varphi_{93} \to \varphi_{102}$\lbl{A1}}
\BIC{$\varphi_{92} \to \varphi_{91} \to \varphi_{89} \to \varphi_{93} \to \varphi_{102}$}
\BIC{$\varphi_{91} \to \varphi_{89} \to \varphi_{93} \to \varphi_{102}$}
\BIC{$\varphi_{89} \to \varphi_{93} \to \varphi_{102}$}
\end{ptree}
\noindent$\Pi_{22}$:
\begin{ptree}
\AXC{$z = w \to z < K^{-1} \to w < K^{-1}$\lbl{A6}}
\UIC{$\forall w\,(z = w \to z < K^{-1} \to w < K^{-1})$}
\UIC{$\forall z\,\forall w\,(z = w \to z < K^{-1} \to w < K^{-1})$}
\AXC{$\varphi_{64}$\lbl{A2}}
\BIC{$\forall w\,(\lvert a_{n} - 0 \rvert = w \to \lvert a_{n} - 0 \rvert < K^{-1} \to w < K^{-1})$}
\AXC{$\varphi_{85}$\lbl{A2}}
\BIC{$\lvert a_{n} - 0 \rvert = a_{n} \to \lvert a_{n} - 0 \rvert < K^{-1} \to a_{n} < K^{-1}$}
\AXC{$\stackrel{\Pi_{2}}{\vdots}$}\noLine
\UIC{$\varphi_{101}$}
\AXC{$\stackrel{\Pi_{15}}{\vdots}$}\noLine
\UIC{$\varphi_{101} \to \varphi_{118}$}
\BIC{$\varphi_{118}$}
\BIC{$\varphi_{93}$}
\AXC{$\stackrel{\Pi_{5}}{\vdots}$}\noLine
\UIC{$\varphi_{89}$}
\AXC{$\stackrel{\Pi_{21}}{\vdots}$}\noLine
\UIC{$\varphi_{89} \to \varphi_{93} \to \varphi_{102}$}
\BIC{$\varphi_{93} \to \varphi_{102}$}
\BIC{$\varphi_{102}$}
\end{ptree}
\noindent$\Pi_{23}$:
\begin{ptree}
\AXC{$\varphi_{11}$\lbl{A1}}
\AXC{$\forall x \in \mathbb{R}\,(x \neq 0 \to (x^{-1})^{-1} = x)$\lbl{[E30]}}
\AXC{$\varphi_{22}$\lbl{A2}}
\BIC{$K \in \mathbb{R} \to K \neq 0 \to (K^{-1})^{-1} = K$}
\AXC{$(K \in \mathbb{R} \to K \neq 0 \to (K^{-1})^{-1} = K) \to \varphi_{11} \to \varphi_{19} \to \varphi_{90}$\lbl{A1}}
\BIC{$\varphi_{11} \to \varphi_{19} \to \varphi_{90}$}
\BIC{$\varphi_{19} \to \varphi_{90}$}
\end{ptree}
\noindent$\Pi_{24}$:
\begin{ptree}
\AXC{$\stackrel{\Pi_{1}}{\vdots}$}\noLine
\UIC{$(K^{-1})^{-1} = K \to (K^{-1})^{-1} < a_{n}^{-1} \to K < a_{n}^{-1}$}
\AXC{$\stackrel{\Pi_{22}}{\vdots}$}\noLine
\UIC{$\varphi_{102}$}
\AXC{$\stackrel{\Pi_{19}}{\vdots}$}\noLine
\UIC{$\varphi_{19}$}
\AXC{$\stackrel{\Pi_{23}}{\vdots}$}\noLine
\UIC{$\varphi_{19} \to \varphi_{90}$}
\BIC{$\varphi_{90}$}
\AXC{$\varphi_{90} \to \varphi_{102} \to \varphi_{115}$\lbl{A1}}
\BIC{$\varphi_{102} \to \varphi_{115}$}
\BIC{$\varphi_{115}$}
\BIC{$\varphi_{94}$}
\end{ptree}
\noindent$\Pi_{25}$:
\begin{ptree}
\AXC{$\stackrel{\Pi_{16}}{\vdots}$}\noLine
\UIC{$\forall x\,\forall y\,(x = y \to y = x)$}
\AXC{$\forall x\,\forall y\,(x = y \to y = x) \to \forall y\,(a_{n} = y \to y = a_{n})$\lbl{A2}}
\BIC{$\forall y\,(a_{n} = y \to y = a_{n})$}
\end{ptree}
\noindent$\Pi_{26}$:
\begin{ptree}
\AXC{$\varphi_{8}$\lbl{A1}}
\AXC{$\forall n \in \mathbb{N}\,(0 < a_{n}) \to n \in \mathbb{N} \to 0 < a_{n}$\lbl{A2}}
\AXC{$(\forall n \in \mathbb{N}\,(0 < a_{n}) \to n \in \mathbb{N} \to 0 < a_{n}) \to \varphi_{8} \to \varphi_{17}$\lbl{A1}}
\BIC{$\varphi_{8} \to \varphi_{17}$}
\BIC{$\varphi_{17}$}
\end{ptree}
\noindent$\Pi_{27}$:
\begin{ptree}
\AXC{$\forall x\,\forall y\,(x < y \to x \neq y)$\lbl{[F26]}}
\AXC{$\forall x\,\forall y\,(x < y \to x \neq y) \to \forall y\,(0 < y \to 0 \neq y)$\lbl{A2}}
\BIC{$\forall y\,(0 < y \to 0 \neq y)$}
\AXC{$\forall y\,(0 < y \to 0 \neq y) \to 0 < a_{n} \to 0 \neq a_{n}$\lbl{A2}}
\BIC{$0 < a_{n} \to 0 \neq a_{n}$}
\end{ptree}
\noindent$\Pi_{28}$:
\begin{ptree}
\AXC{$\stackrel{\Pi_{26}}{\vdots}$}\noLine
\UIC{$\varphi_{17}$}
\AXC{$\stackrel{\Pi_{27}}{\vdots}$}\noLine
\UIC{$0 < a_{n} \to 0 \neq a_{n}$}
\AXC{$(0 < a_{n} \to 0 \neq a_{n}) \to \varphi_{17} \to \varphi_{13}$\lbl{A1}}
\BIC{$\varphi_{17} \to \varphi_{13}$}
\BIC{$\varphi_{13}$}
\AXC{$\varphi_{13} \to (a_{n} = 0 \to 0 = a_{n}) \to \varphi_{14}$\lbl{A1}}
\BIC{$(a_{n} = 0 \to 0 = a_{n}) \to \varphi_{14}$}
\end{ptree}
\noindent$\Pi_{29}$:
\begin{ptree}
\AXC{$\stackrel{\Pi_{25}}{\vdots}$}\noLine
\UIC{$\forall y\,(a_{n} = y \to y = a_{n})$}
\AXC{$\forall y\,(a_{n} = y \to y = a_{n}) \to a_{n} = 0 \to 0 = a_{n}$\lbl{A2}}
\BIC{$a_{n} = 0 \to 0 = a_{n}$}
\AXC{$\stackrel{\Pi_{28}}{\vdots}$}\noLine
\UIC{$(a_{n} = 0 \to 0 = a_{n}) \to \varphi_{14}$}
\BIC{$\varphi_{14}$}
\AXC{$\varphi_{14} \to \varphi_{15}$\lbl{A1}}
\BIC{$\varphi_{15}$}
\UIC{$\forall n\,(\varphi_{15})$}
\end{ptree}
\noindent$\Pi_{30}$:
\begin{ptree}
\AXC{$\varphi_{42}$\lbl{A4}}
\AXC{$\forall n\,(\varphi_{15}) \to \varphi_{24}$\lbl{A3}}
\AXC{$\stackrel{\Pi_{29}}{\vdots}$}\noLine
\UIC{$\forall n\,(\varphi_{15})$}
\AXC{$\forall n\,(\varphi_{15}) \to (\forall n\,(\varphi_{15}) \to \varphi_{24}) \to \varphi_{42} \to \varphi_{20}$\lbl{A1}}
\BIC{$(\forall n\,(\varphi_{15}) \to \varphi_{24}) \to \varphi_{42} \to \varphi_{20}$}
\BIC{$\varphi_{42} \to \varphi_{20}$}
\BIC{$\varphi_{20}$}
\end{ptree}
\noindent$\Pi_{31}$:
\begin{ptree}
\AXC{$\forall a\,(\mathrm{Seq}(a) \to \varphi_{46})$\lbl{[S5]}}
\AXC{$\forall a\,(\mathrm{Seq}(a) \to \varphi_{46}) \to \mathrm{Seq}(a) \to \varphi_{46}$\lbl{A2}}
\BIC{$\mathrm{Seq}(a) \to \varphi_{46}$}
\AXC{$(\mathrm{Seq}(a) \to \varphi_{46}) \to \varphi_{114}$\lbl{A1}}
\BIC{$\varphi_{114}$}
\end{ptree}
\noindent$\Pi_{32}$:
\begin{ptree}
\AXC{$\stackrel{\Pi_{30}}{\vdots}$}\noLine
\UIC{$\varphi_{20}$}
\AXC{$\stackrel{\Pi_{8}}{\vdots}$}\noLine
\UIC{$\forall n \in \mathbb{N}\,(0 < a_{n}) \to \mathrm{Conv}(a,0) \to \mathrm{Seq}(a)$}
\AXC{$\stackrel{\Pi_{31}}{\vdots}$}\noLine
\UIC{$\varphi_{114}$}
\BIC{$\varphi_{20} \to \forall n \in \mathbb{N}\,(0 < a_{n}) \to \varphi_{25}$}
\BIC{$\forall n \in \mathbb{N}\,(0 < a_{n}) \to \varphi_{25}$}
\end{ptree}
\noindent$\Pi_{33}$:
\begin{ptree}
\AXC{$\varphi_{45}$\lbl{A2}}
\AXC{$\stackrel{\Pi_{32}}{\vdots}$}\noLine
\UIC{$\forall n \in \mathbb{N}\,(0 < a_{n}) \to \varphi_{25}$}
\AXC{$(\forall n \in \mathbb{N}\,(0 < a_{n}) \to \varphi_{25}) \to \varphi_{111}$\lbl{A1}}
\BIC{$\varphi_{111}$}
\AXC{$\varphi_{111} \to \varphi_{45} \to \varphi_{109}$\lbl{A1}}
\BIC{$\varphi_{45} \to \varphi_{109}$}
\BIC{$\varphi_{109}$}
\AXC{$\varphi_{109} \to \varphi_{87} \to \varphi_{103}$\lbl{A1}}
\BIC{$\varphi_{87} \to \varphi_{103}$}
\end{ptree}
\noindent$\Pi_{34}$:
\begin{ptree}
\AXC{$\stackrel{\Pi_{16}}{\vdots}$}\noLine
\UIC{$\forall x\,\forall y\,(x = y \to y = x)$}
\AXC{$\varphi_{7}$\lbl{A2}}
\BIC{$\forall y\,(a^{\ast}(n) = y \to y = a^{\ast}(n))$}
\AXC{$\varphi_{59}$\lbl{A2}}
\BIC{$a^{\ast}(n) = a_{n}^{-1} \to a_{n}^{-1} = a^{\ast}(n)$}
\AXC{$\varphi_{87}$\lbl{A1}}
\AXC{$\stackrel{\Pi_{33}}{\vdots}$}\noLine
\UIC{$\varphi_{87} \to \varphi_{103}$}
\BIC{$\varphi_{103}$}
\AXC{$\varphi_{126}$\lbl{A1}}
\BIC{$(a^{\ast}(n) = a_{n}^{-1} \to a_{n}^{-1} = a^{\ast}(n)) \to \varphi_{104}$}
\BIC{$\varphi_{104}$}
\AXC{$\varphi_{104} \to \varphi_{94} \to \varphi_{116}$\lbl{A1}}
\BIC{$\varphi_{94} \to \varphi_{116}$}
\end{ptree}
\noindent$\Pi_{35}$:
\begin{ptree}
\AXC{$z = w \to K < z \to K < w$\lbl{A6}}
\UIC{$\forall w\,(z = w \to K < z \to K < w)$}
\UIC{$\forall z\,\forall w\,(z = w \to K < z \to K < w)$}
\AXC{$\varphi_{30}$\lbl{A2}}
\BIC{$\forall w\,(a_{n}^{-1} = w \to K < a_{n}^{-1} \to K < w)$}
\AXC{$\varphi_{76}$\lbl{A2}}
\BIC{$a_{n}^{-1} = a^{\ast}(n) \to K < a_{n}^{-1} \to K < a^{\ast}(n)$}
\AXC{$\stackrel{\Pi_{24}}{\vdots}$}\noLine
\UIC{$\varphi_{94}$}
\AXC{$\stackrel{\Pi_{34}}{\vdots}$}\noLine
\UIC{$\varphi_{94} \to \varphi_{116}$}
\BIC{$\varphi_{116}$}
\BIC{$\varphi_{96}$}
\end{ptree}
\noindent$\Pi_{36}$:
\begin{ptree}
\AXC{$\varphi_{79} \to \forall n\,(\varphi_{79})$\lbl{A4}}
\AXC{$\varphi_{122}$\lbl{A3}}
\AXC{$\stackrel{\Pi_{35}}{\vdots}$}\noLine
\UIC{$\varphi_{96}$}
\AXC{$\varphi_{96} \to \varphi_{79} \to n \in \mathbb{N} \to N \leq n \to K < a^{\ast}(n)$\lbl{A1}}
\BIC{$\varphi_{79} \to n \in \mathbb{N} \to N \leq n \to K < a^{\ast}(n)$}
\UIC{$\forall n\,(\varphi_{79} \to n \in \mathbb{N} \to N \leq n \to K < a^{\ast}(n))$}
\AXC{$\varphi_{129}$\lbl{A1}}
\BIC{$\varphi_{122} \to (\varphi_{79} \to \forall n\,(\varphi_{79})) \to \varphi_{99}$}
\BIC{$(\varphi_{79} \to \forall n\,(\varphi_{79})) \to \varphi_{99}$}
\BIC{$\varphi_{99}$}
\end{ptree}
\noindent$\Pi_{37}$:
\begin{ptree}
\AXC{$\stackrel{\Pi_{36}}{\vdots}$}\noLine
\UIC{$\varphi_{99}$}
\AXC{$\varphi_{110}$\lbl{A1}}
\AXC{$\varphi_{110} \to \varphi_{84}$\lbl{A1}}
\BIC{$\varphi_{84}$}
\AXC{$\varphi_{84} \to \varphi_{99} \to \varphi_{106}$\lbl{A1}}
\BIC{$\varphi_{99} \to \varphi_{106}$}
\BIC{$\varphi_{106}$}
\AXC{$\varphi_{106} \to \varphi_{121}$\lbl{A1}}
\BIC{$\varphi_{121}$}
\end{ptree}
\noindent$\Pi_{38}$:
\begin{ptree}
\AXC{$\varphi_{6} \to \neg (N \in \mathbb{N} \land \forall n \in \mathbb{N}\,(N \leq n \to K < a^{\ast}(n)))$\lbl{A2}}
\AXC{$\stackrel{\Pi_{37}}{\vdots}$}\noLine
\UIC{$\varphi_{121}$}
\BIC{$\varphi_{107}$}
\AXC{$\varphi_{107} \to \varphi_{108}$\lbl{A1}}
\BIC{$\varphi_{108}$}
\UIC{$\forall N\,(\varphi_{108})$}
\end{ptree}
\noindent$\Pi_{39}$:
\begin{ptree}
\AXC{$\varphi_{21}$\lbl{A1}}
\AXC{$\varphi_{11}$\lbl{A1}}
\AXC{$\forall x \in \mathbb{R}\,(0 < x \to 0 < x^{-1})$\lbl{[F34]}}
\AXC{$\varphi_{23}$\lbl{A2}}
\BIC{$K \in \mathbb{R} \to 0 < K \to 0 < K^{-1}$}
\AXC{$(K \in \mathbb{R} \to 0 < K \to 0 < K^{-1}) \to \varphi_{11} \to \varphi_{21} \to \varphi_{29}$\lbl{A1}}
\BIC{$\varphi_{11} \to \varphi_{21} \to \varphi_{29}$}
\BIC{$\varphi_{21} \to \varphi_{29}$}
\BIC{$\varphi_{29}$}
\end{ptree}
\noindent$\Pi_{40}$:
\begin{ptree}
\AXC{$\stackrel{\Pi_{19}}{\vdots}$}\noLine
\UIC{$\varphi_{19}$}
\AXC{$\varphi_{11}$\lbl{A1}}
\AXC{$\forall x \in \mathbb{R}\,(x \neq 0 \to x^{-1} \in \mathbb{R})$\lbl{(R1)}}
\AXC{$\varphi_{28}$\lbl{A2}}
\BIC{$K \in \mathbb{R} \to K \neq 0 \to K^{-1} \in \mathbb{R}$}
\AXC{$(K \in \mathbb{R} \to K \neq 0 \to K^{-1} \in \mathbb{R}) \to \varphi_{11} \to \varphi_{19} \to \varphi_{33}$\lbl{A1}}
\BIC{$\varphi_{11} \to \varphi_{19} \to \varphi_{33}$}
\BIC{$\varphi_{19} \to \varphi_{33}$}
\BIC{$\varphi_{33}$}
\end{ptree}
\noindent$\Pi_{41}$:
\begin{ptree}
\AXC{$\forall a\,(\varphi_{66})$\lbl{[C3]}}
\AXC{$\forall a\,(\varphi_{66}) \to \varphi_{66}$\lbl{A2}}
\BIC{$\varphi_{66}$}
\AXC{$\varphi_{66} \to \varphi_{60}$\lbl{A2}}
\BIC{$\varphi_{60}$}
\AXC{$\varphi_{60} \to \mathrm{Conv}(a,0) \to K^{-1} \in \mathbb{R} \to \varphi_{31}$\lbl{A2}}
\BIC{$\mathrm{Conv}(a,0) \to K^{-1} \in \mathbb{R} \to \varphi_{31}$}
\AXC{$\varphi_{127}$\lbl{A1}}
\BIC{$\varphi_{2} \to \varphi_{33} \to \varphi_{29} \to \forall n \in \mathbb{N}\,(0 < a_{n}) \to \mathrm{Conv}(a,0) \to \varphi_{43}$}
\end{ptree}
\noindent$\Pi_{42}$:
\begin{ptree}
\AXC{$\stackrel{\Pi_{39}}{\vdots}$}\noLine
\UIC{$\varphi_{29}$}
\AXC{$\stackrel{\Pi_{40}}{\vdots}$}\noLine
\UIC{$\varphi_{33}$}
\AXC{$\varphi_{2}$\lbl{A1}}
\AXC{$\stackrel{\Pi_{41}}{\vdots}$}\noLine
\UIC{$\varphi_{2} \to \varphi_{33} \to \varphi_{29} \to \forall n \in \mathbb{N}\,(0 < a_{n}) \to \mathrm{Conv}(a,0) \to \varphi_{43}$}
\BIC{$\varphi_{33} \to \varphi_{29} \to \forall n \in \mathbb{N}\,(0 < a_{n}) \to \mathrm{Conv}(a,0) \to \varphi_{43}$}
\BIC{$\varphi_{29} \to \forall n \in \mathbb{N}\,(0 < a_{n}) \to \mathrm{Conv}(a,0) \to \varphi_{43}$}
\BIC{$\forall n \in \mathbb{N}\,(0 < a_{n}) \to \mathrm{Conv}(a,0) \to \varphi_{43}$}
\AXC{$\varphi_{130}$\lbl{A1}}
\BIC{$\forall N\,(\varphi_{108}) \to (\forall N\,(\varphi_{108}) \to \forall N\,(\varphi_{74}) \to \varphi_{10}) \to \varphi_{125}$}
\end{ptree}
\noindent$\Pi_{43}$:
\begin{ptree}
\AXC{$\varphi_{74} \to \forall N\,(\varphi_{74})$\lbl{A4}}
\AXC{$\forall N\,(\varphi_{108}) \to \forall N\,(\varphi_{74}) \to \varphi_{10}$\lbl{A3}}
\AXC{$\stackrel{\Pi_{38}}{\vdots}$}\noLine
\UIC{$\forall N\,(\varphi_{108})$}
\AXC{$\stackrel{\Pi_{42}}{\vdots}$}\noLine
\UIC{$\forall N\,(\varphi_{108}) \to (\forall N\,(\varphi_{108}) \to \forall N\,(\varphi_{74}) \to \varphi_{10}) \to \varphi_{125}$}
\BIC{$(\forall N\,(\varphi_{108}) \to \forall N\,(\varphi_{74}) \to \varphi_{10}) \to \varphi_{125}$}
\BIC{$\varphi_{125}$}
\BIC{$\forall n \in \mathbb{N}\,(0 < a_{n}) \to \mathrm{Conv}(a,0) \to \varphi_{32}$}
\end{ptree}
\noindent$\Pi_{44}$:
\begin{ptree}
\AXC{$\varphi_{40}$\lbl{A4}}
\AXC{$\varphi_{119}$\lbl{A3}}
\AXC{$\stackrel{\Pi_{43}}{\vdots}$}\noLine
\UIC{$\forall n \in \mathbb{N}\,(0 < a_{n}) \to \mathrm{Conv}(a,0) \to \varphi_{32}$}
\AXC{$\varphi_{117}$\lbl{A1}}
\BIC{$\forall n \in \mathbb{N}\,(0 < a_{n}) \land \mathrm{Conv}(a,0) \to \varphi_{32}$}
\UIC{$\forall K\,(\forall n \in \mathbb{N}\,(0 < a_{n}) \land \mathrm{Conv}(a,0) \to \varphi_{32})$}
\AXC{$\varphi_{128}$\lbl{A1}}
\BIC{$\varphi_{119} \to \varphi_{40} \to \forall n \in \mathbb{N}\,(0 < a_{n}) \to \mathrm{Conv}(a,0) \to \varphi_{34}$}
\BIC{$\varphi_{40} \to \forall n \in \mathbb{N}\,(0 < a_{n}) \to \mathrm{Conv}(a,0) \to \varphi_{34}$}
\BIC{$\forall n \in \mathbb{N}\,(0 < a_{n}) \to \mathrm{Conv}(a,0) \to \varphi_{34}$}
\end{ptree}
\noindent$\Pi_{45}$:
\begin{ptree}
\AXC{$\varphi_{51}$\lbl{Def$\mathrm{DvgP}$}}
\AXC{$\varphi_{51} \to (\mathrm{DvgP}(a^{\ast}) \leftrightarrow \mathrm{Seq}(a^{\ast}) \land \varphi_{34})$\lbl{A2}}
\BIC{$\mathrm{DvgP}(a^{\ast}) \leftrightarrow \mathrm{Seq}(a^{\ast}) \land \varphi_{34}$}
\AXC{$\varphi_{123}$\lbl{A1}}
\BIC{$\varphi_{4} \to (\forall n \in \mathbb{N}\,(0 < a_{n}) \to \mathrm{Conv}(a,0) \to \varphi_{34}) \to \varphi_{5}$}
\end{ptree}
\noindent$\Pi_{46}$:
\begin{ptree}
\AXC{$\stackrel{\Pi_{32}}{\vdots}$}\noLine
\UIC{$\forall n \in \mathbb{N}\,(0 < a_{n}) \to \varphi_{25}$}
\AXC{$(\forall n \in \mathbb{N}\,(0 < a_{n}) \to \varphi_{25}) \to \varphi_{4}$\lbl{A1}}
\BIC{$\varphi_{4}$}
\AXC{$\stackrel{\Pi_{45}}{\vdots}$}\noLine
\UIC{$\varphi_{4} \to (\forall n \in \mathbb{N}\,(0 < a_{n}) \to \mathrm{Conv}(a,0) \to \varphi_{34}) \to \varphi_{5}$}
\BIC{$(\forall n \in \mathbb{N}\,(0 < a_{n}) \to \mathrm{Conv}(a,0) \to \varphi_{34}) \to \varphi_{5}$}
\end{ptree}
\noindent 主木:
\begin{ptree}
\AXC{$\stackrel{\Pi_{44}}{\vdots}$}\noLine
\UIC{$\forall n \in \mathbb{N}\,(0 < a_{n}) \to \mathrm{Conv}(a,0) \to \varphi_{34}$}
\AXC{$\stackrel{\Pi_{46}}{\vdots}$}\noLine
\UIC{$(\forall n \in \mathbb{N}\,(0 < a_{n}) \to \mathrm{Conv}(a,0) \to \varphi_{34}) \to \varphi_{5}$}
\BIC{$\varphi_{5}$}
\UIC{$\forall a\,(\varphi_{5})$}
\end{ptree}
\begin{abbrevlist}
\item $\varphi_{1} :\equiv \forall \alpha\,(\mathrm{Conv}(a,\alpha) \to \mathrm{Seq}(a)) \to \mathrm{Conv}(a,0) \to \mathrm{Seq}(a)$
\item $\varphi_{2} :\equiv \forall n \in \mathbb{N}\,(0 < a_{n}) \to \mathrm{Conv}(a,0) \to \mathrm{Conv}(a,0)$
\item $\varphi_{3} :\equiv \forall a\,\forall \alpha\,(\mathrm{Conv}(a,\alpha) \to \mathrm{Seq}(a)) \to \forall \alpha\,(\mathrm{Conv}(a,\alpha) \to \mathrm{Seq}(a))$
\item $\varphi_{4} :\equiv \forall n \in \mathbb{N}\,(0 < a_{n}) \to \mathrm{Conv}(a,0) \to \mathrm{Seq}(a^{\ast})$
\item $\varphi_{5} :\equiv \forall n \in \mathbb{N}\,(0 < a_{n}) \to \mathrm{Conv}(a,0) \to \mathrm{DvgP}(a^{\ast})$
\item $\varphi_{6} :\equiv \forall N\,(\neg (N \in \mathbb{N} \land \forall n \in \mathbb{N}\,(N \leq n \to K < a^{\ast}(n))))$
\item $\varphi_{7} :\equiv \forall x\,\forall y\,(x = y \to y = x) \to \forall y\,(a^{\ast}(n) = y \to y = a^{\ast}(n))$
\item $\varphi_{8} :\equiv \forall n \in \mathbb{N}\,(0 < a_{n}) \to \mathrm{Conv}(a,0) \to n \in \mathbb{N} \to n \in \mathbb{N}$
\item $\varphi_{9} :\equiv \forall z\,\forall w\,(z = w \to \lvert a_{n} - 0 \rvert = \lvert z \rvert \to \lvert a_{n} - 0 \rvert = \lvert w \rvert)$
\item $\varphi_{10} :\equiv \forall N\,(\neg (N \in \mathbb{N} \land \forall n \in \mathbb{N}\,(N \leq n \to \lvert a_{n} - 0 \rvert < K^{-1})))$
\item $\varphi_{11} :\equiv \forall n \in \mathbb{N}\,(0 < a_{n}) \to \mathrm{Conv}(a,0) \to K \in \mathbb{R} \to K \in \mathbb{R}$
\item $\varphi_{12} :\equiv \forall x\,\forall y\,(x \in \mathbb{R} \to y \in \mathbb{R} \to 0 < x \to x < y \to y^{-1} < x^{-1})$
\item $\varphi_{13} :\equiv \forall n \in \mathbb{N}\,(0 < a_{n}) \to \mathrm{Conv}(a,0) \to n \in \mathbb{N} \to 0 \neq a_{n}$
\item $\varphi_{14} :\equiv \forall n \in \mathbb{N}\,(0 < a_{n}) \to \mathrm{Conv}(a,0) \to n \in \mathbb{N} \to a_{n} \neq 0$
\item $\varphi_{15} :\equiv \forall n \in \mathbb{N}\,(0 < a_{n}) \land \mathrm{Conv}(a,0) \to n \in \mathbb{N} \to a_{n} \neq 0$
\item $\varphi_{16} :\equiv \forall w\,((K^{-1})^{-1} = w \to (K^{-1})^{-1} < a_{n}^{-1} \to w < a_{n}^{-1})$
\item $\varphi_{17} :\equiv \forall n \in \mathbb{N}\,(0 < a_{n}) \to \mathrm{Conv}(a,0) \to n \in \mathbb{N} \to 0 < a_{n}$
\item $\varphi_{18} :\equiv \forall n \in \mathbb{N}\,(0 < a_{n}) \to \mathrm{Conv}(a,0) \to K \in \mathbb{R} \to 0 < K \to 0 \neq K$
\item $\varphi_{19} :\equiv \forall n \in \mathbb{N}\,(0 < a_{n}) \to \mathrm{Conv}(a,0) \to K \in \mathbb{R} \to 0 < K \to K \neq 0$
\item $\varphi_{20} :\equiv \forall n \in \mathbb{N}\,(0 < a_{n}) \to \mathrm{Conv}(a,0) \to \forall n \in \mathbb{N}\,(a_{n} \neq 0)$
\item $\varphi_{21} :\equiv \forall n \in \mathbb{N}\,(0 < a_{n}) \to \mathrm{Conv}(a,0) \to K \in \mathbb{R} \to 0 < K \to 0 < K$
\item $\varphi_{22} :\equiv \forall x \in \mathbb{R}\,(x \neq 0 \to (x^{-1})^{-1} = x) \to K \in \mathbb{R} \to K \neq 0 \to (K^{-1})^{-1} = K$
\item $\varphi_{23} :\equiv \forall x \in \mathbb{R}\,(0 < x \to 0 < x^{-1}) \to K \in \mathbb{R} \to 0 < K \to 0 < K^{-1}$
\item $\varphi_{24} :\equiv \forall n\,(\forall n \in \mathbb{N}\,(0 < a_{n}) \land \mathrm{Conv}(a,0)) \to \forall n \in \mathbb{N}\,(a_{n} \neq 0)$
\item $\varphi_{25} :\equiv \mathrm{Conv}(a,0) \to \mathrm{Seq}(a^{\ast}) \land \forall n \in \mathbb{N}\,(a^{\ast}(n) = a_{n}^{-1})$
\item $\varphi_{26} :\equiv N \in \mathbb{N} \land \forall n \in \mathbb{N}\,(N \leq n \to \lvert a_{n} - 0 \rvert < K^{-1}) \to N \in \mathbb{N}$
\item $\varphi_{27} :\equiv \forall z\,(\lvert a_{n} - 0 \rvert = \lvert a_{n} \rvert \to \lvert a_{n} \rvert = z \to \lvert a_{n} - 0 \rvert = z)$
\item $\varphi_{28} :\equiv \forall x \in \mathbb{R}\,(x \neq 0 \to x^{-1} \in \mathbb{R}) \to K \in \mathbb{R} \to K \neq 0 \to K^{-1} \in \mathbb{R}$
\item $\varphi_{29} :\equiv \forall n \in \mathbb{N}\,(0 < a_{n}) \to \mathrm{Conv}(a,0) \to K \in \mathbb{R} \to 0 < K \to 0 < K^{-1}$
\item $\varphi_{30} :\equiv \forall z\,\forall w\,(z = w \to K < z \to K < w) \to \forall w\,(a_{n}^{-1} = w \to K < a_{n}^{-1} \to K < w)$
\item $\varphi_{31} :\equiv 0 < K^{-1} \to \exists N \in \mathbb{N}\,\forall n \in \mathbb{N}\,(N \leq n \to \lvert a_{n} - 0 \rvert < K^{-1})$
\item $\varphi_{32} :\equiv K \in \mathbb{R} \to 0 < K \to \exists N \in \mathbb{N}\,\forall n \in \mathbb{N}\,(N \leq n \to K < a^{\ast}(n))$
\item $\varphi_{33} :\equiv \forall n \in \mathbb{N}\,(0 < a_{n}) \to \mathrm{Conv}(a,0) \to K \in \mathbb{R} \to 0 < K \to K^{-1} \in \mathbb{R}$
\item $\varphi_{34} :\equiv \forall K \in \mathbb{R}\,(0 < K \to \exists N \in \mathbb{N}\,\forall n \in \mathbb{N}\,(N \leq n \to K < a^{\ast}(n)))$
\item $\varphi_{35} :\equiv \forall x\,\forall y\,\forall z\,(x = y \to y = z \to x = z) \to \forall y\,\forall z\,(\lvert a_{n} - 0 \rvert = y \to y = z \to \lvert a_{n} - 0 \rvert = z)$
\item $\varphi_{36} :\equiv \forall n\,(\mathrm{Seq}(a) \to n \in \mathbb{N} \to a_{n} \in \mathbb{R}) \to \mathrm{Seq}(a) \to n \in \mathbb{N} \to a_{n} \in \mathbb{R}$
\item $\varphi_{37} :\equiv K^{-1} \in \mathbb{R} \to 0 < a_{n} \to a_{n} < K^{-1} \to (K^{-1})^{-1} < a_{n}^{-1}$
\item $\varphi_{38} :\equiv N \in \mathbb{N} \land \forall n \in \mathbb{N}\,(N \leq n \to \lvert a_{n} - 0 \rvert < K^{-1}) \to n \in \mathbb{N} \to N \leq n \to N \leq n$
\item $\varphi_{39} :\equiv N \in \mathbb{N} \land \forall n \in \mathbb{N}\,(N \leq n \to \lvert a_{n} - 0 \rvert < K^{-1}) \to n \in \mathbb{N} \to n \in \mathbb{N}$
\item $\varphi_{40} :\equiv \forall n \in \mathbb{N}\,(0 < a_{n}) \land \mathrm{Conv}(a,0) \to \forall K\,(\forall n \in \mathbb{N}\,(0 < a_{n}) \land \mathrm{Conv}(a,0))$
\item $\varphi_{41} :\equiv \forall w\,(a_{n} - 0 = w \to \lvert a_{n} - 0 \rvert = \lvert a_{n} - 0 \rvert \to \lvert a_{n} - 0 \rvert = \lvert w \rvert)$
\item $\varphi_{42} :\equiv \forall n \in \mathbb{N}\,(0 < a_{n}) \land \mathrm{Conv}(a,0) \to \forall n\,(\forall n \in \mathbb{N}\,(0 < a_{n}) \land \mathrm{Conv}(a,0))$
\item $\varphi_{43} :\equiv K \in \mathbb{R} \to 0 < K \to \exists N \in \mathbb{N}\,\forall n \in \mathbb{N}\,(N \leq n \to \lvert a_{n} - 0 \rvert < K^{-1})$
\item $\varphi_{44} :\equiv \forall a\,\forall n\,(\mathrm{Seq}(a) \to n \in \mathbb{N} \to a_{n} \in \mathbb{R}) \to \forall n\,(\mathrm{Seq}(a) \to n \in \mathbb{N} \to a_{n} \in \mathbb{R})$
\item $\varphi_{45} :\equiv \forall n \in \mathbb{N}\,(a^{\ast}(n) = a_{n}^{-1}) \to n \in \mathbb{N} \to a^{\ast}(n) = a_{n}^{-1}$
\item $\varphi_{46} :\equiv \forall n \in \mathbb{N}\,(a_{n} \neq 0) \to \mathrm{Seq}(a^{\ast}) \land \forall n \in \mathbb{N}\,(a^{\ast}(n) = a_{n}^{-1})$
\item $\varphi_{47} :\equiv \forall y\,(a_{n} \in \mathbb{R} \to y \in \mathbb{R} \to 0 < a_{n} \to a_{n} < y \to y^{-1} < a_{n}^{-1})$
\item $\varphi_{48} :\equiv N \in \mathbb{N} \land \forall n \in \mathbb{N}\,(N \leq n \to \lvert a_{n} - 0 \rvert < K^{-1}) \to n \in \mathbb{N} \to N \leq n \to 0 \leq a_{n}$
\item $\varphi_{49} :\equiv N \in \mathbb{N} \land \forall n \in \mathbb{N}\,(N \leq n \to \lvert a_{n} - 0 \rvert < K^{-1}) \to n \in \mathbb{N} \to N \leq n \to 0 < a_{n}$
\item $\varphi_{50} :\equiv N \in \mathbb{N} \land \forall n \in \mathbb{N}\,(N \leq n \to \lvert a_{n} - 0 \rvert < K^{-1}) \to n \in \mathbb{N} \to N \leq n \to (K^{-1})^{-1} = K$
\item $\varphi_{51} :\equiv \forall a\,(\mathrm{DvgP}(a) \leftrightarrow \mathrm{Seq}(a) \land \forall K \in \mathbb{R}\,(0 < K \to \exists N \in \mathbb{N}\,\forall n \in \mathbb{N}\,(N \leq n \to K < a_{n})))$
\item $\varphi_{52} :\equiv a_{n} - 0 = a_{n} \to \lvert a_{n} - 0 \rvert = \lvert a_{n} - 0 \rvert \to \lvert a_{n} - 0 \rvert = \lvert a_{n} \rvert$
\item $\varphi_{53} :\equiv N \in \mathbb{N} \land \forall n \in \mathbb{N}\,(N \leq n \to \lvert a_{n} - 0 \rvert < K^{-1}) \to n \in \mathbb{N} \to N \leq n \to K^{-1} \in \mathbb{R}$
\item $\varphi_{54} :\equiv \forall n \in \mathbb{N}\,(N \leq n \to \lvert a_{n} - 0 \rvert < K^{-1}) \to n \in \mathbb{N} \to N \leq n \to \lvert a_{n} - 0 \rvert < K^{-1}$
\item $\varphi_{55} :\equiv N \in \mathbb{N} \land \forall n \in \mathbb{N}\,(N \leq n \to \lvert a_{n} - 0 \rvert < K^{-1}) \to n \in \mathbb{N} \to N \leq n \to a_{n} \in \mathbb{R}$
\item $\varphi_{56} :\equiv N \in \mathbb{N} \land \forall n \in \mathbb{N}\,(N \leq n \to \lvert a_{n} - 0 \rvert < K^{-1}) \to n \in \mathbb{N} \to N \leq n \to a_{n} < K^{-1}$
\item $\varphi_{57} :\equiv N \in \mathbb{N} \land \forall n \in \mathbb{N}\,(N \leq n \to \lvert a_{n} - 0 \rvert < K^{-1}) \to n \in \mathbb{N} \to N \leq n \to K < a_{n}^{-1}$
\item $\varphi_{58} :\equiv N \in \mathbb{N} \land \forall n \in \mathbb{N}\,(N \leq n \to \lvert a_{n} - 0 \rvert < K^{-1}) \to n \in \mathbb{N} \to N \leq n \to K < a^{\ast}(n)$
\item $\varphi_{59} :\equiv \forall y\,(a^{\ast}(n) = y \to y = a^{\ast}(n)) \to a^{\ast}(n) = a_{n}^{-1} \to a_{n}^{-1} = a^{\ast}(n)$
\item $\varphi_{60} :\equiv \forall \varepsilon\,(\mathrm{Conv}(a,0) \to \varepsilon \in \mathbb{R} \to 0 < \varepsilon \to \exists N \in \mathbb{N}\,\forall n \in \mathbb{N}\,(N \leq n \to \lvert a_{n} - 0 \rvert < \varepsilon))$
\item $\varphi_{61} :\equiv N \in \mathbb{N} \land \forall n \in \mathbb{N}\,(N \leq n \to \lvert a_{n} - 0 \rvert < K^{-1}) \to n \in \mathbb{N} \to N \leq n \to \lvert a_{n} \rvert = a_{n}$
\item $\varphi_{62} :\equiv N \in \mathbb{N} \land \forall n \in \mathbb{N}\,(N \leq n \to \lvert a_{n} - 0 \rvert < K^{-1}) \to n \in \mathbb{N} \to N \leq n \to a_{n} - 0 = a_{n}$
\item $\varphi_{63} :\equiv N \in \mathbb{N} \land \forall n \in \mathbb{N}\,(N \leq n \to \lvert a_{n} - 0 \rvert < K^{-1}) \to \forall n \in \mathbb{N}\,(N \leq n \to K < a^{\ast}(n))$
\item $\varphi_{64} :\equiv \forall z\,\forall w\,(z = w \to z < K^{-1} \to w < K^{-1}) \to \forall w\,(\lvert a_{n} - 0 \rvert = w \to \lvert a_{n} - 0 \rvert < K^{-1} \to w < K^{-1})$
\item $\varphi_{65} :\equiv N \in \mathbb{N} \land \forall n \in \mathbb{N}\,(N \leq n \to \lvert a_{n} - 0 \rvert < K^{-1}) \to n \in \mathbb{N} \to N \leq n \to \lvert a_{n} - 0 \rvert = a_{n}$
\item $\varphi_{66} :\equiv \forall \alpha\,\forall \varepsilon\,(\mathrm{Conv}(a,\alpha) \to \varepsilon \in \mathbb{R} \to 0 < \varepsilon \to \exists N \in \mathbb{N}\,\forall n \in \mathbb{N}\,(N \leq n \to \lvert a_{n} - \alpha \rvert < \varepsilon))$
\item $\varphi_{67} :\equiv N \in \mathbb{N} \land \forall n \in \mathbb{N}\,(N \leq n \to \lvert a_{n} - 0 \rvert < K^{-1}) \to n \in \mathbb{N} \to N \leq n \to \lvert a_{n} - 0 \rvert < K^{-1}$
\item $\varphi_{68} :\equiv N \in \mathbb{N} \land \forall n \in \mathbb{N}\,(N \leq n \to \lvert a_{n} - 0 \rvert < K^{-1}) \to n \in \mathbb{N} \to N \leq n \to (K^{-1})^{-1} < a_{n}^{-1}$
\item $\varphi_{69} :\equiv N \in \mathbb{N} \land \forall n \in \mathbb{N}\,(N \leq n \to \lvert a_{n} - 0 \rvert < K^{-1}) \to n \in \mathbb{N} \to N \leq n \to a^{\ast}(n) = a_{n}^{-1}$
\item $\varphi_{70} :\equiv N \in \mathbb{N} \land \forall n \in \mathbb{N}\,(N \leq n \to \lvert a_{n} - 0 \rvert < K^{-1}) \to n \in \mathbb{N} \to N \leq n \to a_{n}^{-1} = a^{\ast}(n)$
\item $\varphi_{71} :\equiv N \in \mathbb{N} \land \forall n \in \mathbb{N}\,(N \leq n \to \lvert a_{n} - 0 \rvert < K^{-1}) \to n \in \mathbb{N} \to N \leq n \to \lvert a_{n} - 0 \rvert = \lvert a_{n} \rvert$
\item $\varphi_{72} :\equiv N \in \mathbb{N} \land \forall n \in \mathbb{N}\,(N \leq n \to \lvert a_{n} - 0 \rvert < K^{-1}) \to N \in \mathbb{N} \land \forall n \in \mathbb{N}\,(N \leq n \to K < a^{\ast}(n))$
\item $\varphi_{73} :\equiv N \in \mathbb{N} \land \forall n \in \mathbb{N}\,(N \leq n \to \lvert a_{n} - 0 \rvert < K^{-1}) \to \exists N \in \mathbb{N}\,\forall n \in \mathbb{N}\,(N \leq n \to K < a^{\ast}(n))$
\item $\varphi_{74} :\equiv \forall n \in \mathbb{N}\,(0 < a_{n}) \land \mathrm{Conv}(a,0) \land K \in \mathbb{R} \land 0 < K \land \neg \exists N \in \mathbb{N}\,\forall n \in \mathbb{N}\,(N \leq n \to K < a^{\ast}(n))$
\item $\varphi_{75} :\equiv N \in \mathbb{N} \land \forall n \in \mathbb{N}\,(N \leq n \to \lvert a_{n} - 0 \rvert < K^{-1}) \to n \in \mathbb{N} \to N \leq n \to n \in \mathbb{N} \to a^{\ast}(n) = a_{n}^{-1}$
\item $\varphi_{76} :\equiv \forall w\,(a_{n}^{-1} = w \to K < a_{n}^{-1} \to K < w) \to a_{n}^{-1} = a^{\ast}(n) \to K < a_{n}^{-1} \to K < a^{\ast}(n)$
\item $\varphi_{77} :\equiv (\mathrm{Conv}(a,0) \to \mathrm{Seq}(a)) \to \varphi_{2} \to \forall n \in \mathbb{N}\,(0 < a_{n}) \to \mathrm{Conv}(a,0) \to \mathrm{Seq}(a)$
\item $\varphi_{78} :\equiv N \in \mathbb{N} \land \forall n \in \mathbb{N}\,(N \leq n \to \lvert a_{n} - 0 \rvert < K^{-1}) \to N \in \mathbb{N} \land \forall n \in \mathbb{N}\,(N \leq n \to \lvert a_{n} - 0 \rvert < K^{-1})$
\item $\varphi_{79} :\equiv \forall n \in \mathbb{N}\,(0 < a_{n}) \land \mathrm{Conv}(a,0) \land K \in \mathbb{R} \land 0 < K \land (N \in \mathbb{N} \land \forall n \in \mathbb{N}\,(N \leq n \to \lvert a_{n} - 0 \rvert < K^{-1}))$
\item $\varphi_{80} :\equiv N \in \mathbb{N} \land \forall n \in \mathbb{N}\,(N \leq n \to \lvert a_{n} - 0 \rvert < K^{-1}) \to n \in \mathbb{N} \to N \leq n \to \forall n \in \mathbb{N}\,(a^{\ast}(n) = a_{n}^{-1})$
\item $\varphi_{81} :\equiv N \in \mathbb{N} \land \forall n \in \mathbb{N}\,(N \leq n \to \lvert a_{n} - 0 \rvert < K^{-1}) \to n \in \mathbb{N} \to N \leq n \to n \in \mathbb{N} \to N \leq n \to \lvert a_{n} - 0 \rvert < K^{-1}$
\item $\varphi_{82} :\equiv N \in \mathbb{N} \land \forall n \in \mathbb{N}\,(N \leq n \to \lvert a_{n} - 0 \rvert < K^{-1}) \to n \in \mathbb{N} \to N \leq n \to \forall n \in \mathbb{N}\,(N \leq n \to \lvert a_{n} - 0 \rvert < K^{-1})$
\item $\varphi_{83} :\equiv \varphi_{16} \to (K^{-1})^{-1} = K \to (K^{-1})^{-1} < a_{n}^{-1} \to K < a_{n}^{-1}$
\item $\varphi_{84} :\equiv \forall n \in \mathbb{N}\,(0 < a_{n}) \to \mathrm{Conv}(a,0) \to K \in \mathbb{R} \to 0 < K \to \varphi_{26}$
\item $\varphi_{85} :\equiv \forall w\,(\lvert a_{n} - 0 \rvert = w \to \lvert a_{n} - 0 \rvert < K^{-1} \to w < K^{-1}) \to \lvert a_{n} - 0 \rvert = a_{n} \to \lvert a_{n} - 0 \rvert < K^{-1} \to a_{n} < K^{-1}$
\item $\varphi_{86} :\equiv \forall n \in \mathbb{N}\,(0 < a_{n}) \to \mathrm{Conv}(a,0) \to K \in \mathbb{R} \to 0 < K \to \varphi_{38}$
\item $\varphi_{87} :\equiv \forall n \in \mathbb{N}\,(0 < a_{n}) \to \mathrm{Conv}(a,0) \to K \in \mathbb{R} \to 0 < K \to \varphi_{39}$
\item $\varphi_{88} :\equiv \forall n \in \mathbb{N}\,(0 < a_{n}) \to \mathrm{Conv}(a,0) \to K \in \mathbb{R} \to 0 < K \to \varphi_{48}$
\item $\varphi_{89} :\equiv \forall n \in \mathbb{N}\,(0 < a_{n}) \to \mathrm{Conv}(a,0) \to K \in \mathbb{R} \to 0 < K \to \varphi_{49}$
\item $\varphi_{90} :\equiv \forall n \in \mathbb{N}\,(0 < a_{n}) \to \mathrm{Conv}(a,0) \to K \in \mathbb{R} \to 0 < K \to \varphi_{50}$
\item $\varphi_{91} :\equiv \forall n \in \mathbb{N}\,(0 < a_{n}) \to \mathrm{Conv}(a,0) \to K \in \mathbb{R} \to 0 < K \to \varphi_{53}$
\item $\varphi_{92} :\equiv \forall n \in \mathbb{N}\,(0 < a_{n}) \to \mathrm{Conv}(a,0) \to K \in \mathbb{R} \to 0 < K \to \varphi_{55}$
\item $\varphi_{93} :\equiv \forall n \in \mathbb{N}\,(0 < a_{n}) \to \mathrm{Conv}(a,0) \to K \in \mathbb{R} \to 0 < K \to \varphi_{56}$
\item $\varphi_{94} :\equiv \forall n \in \mathbb{N}\,(0 < a_{n}) \to \mathrm{Conv}(a,0) \to K \in \mathbb{R} \to 0 < K \to \varphi_{57}$
\item $\varphi_{95} :\equiv \varphi_{27} \to \lvert a_{n} - 0 \rvert = \lvert a_{n} \rvert \to \lvert a_{n} \rvert = a_{n} \to \lvert a_{n} - 0 \rvert = a_{n}$
\item $\varphi_{96} :\equiv \forall n \in \mathbb{N}\,(0 < a_{n}) \to \mathrm{Conv}(a,0) \to K \in \mathbb{R} \to 0 < K \to \varphi_{58}$
\item $\varphi_{97} :\equiv \forall n \in \mathbb{N}\,(0 < a_{n}) \to \mathrm{Conv}(a,0) \to K \in \mathbb{R} \to 0 < K \to \varphi_{61}$
\item $\varphi_{98} :\equiv \forall n \in \mathbb{N}\,(0 < a_{n}) \to \mathrm{Conv}(a,0) \to K \in \mathbb{R} \to 0 < K \to \varphi_{62}$
\item $\varphi_{99} :\equiv \forall n \in \mathbb{N}\,(0 < a_{n}) \to \mathrm{Conv}(a,0) \to K \in \mathbb{R} \to 0 < K \to \varphi_{63}$
\item $\varphi_{100} :\equiv \forall n \in \mathbb{N}\,(0 < a_{n}) \to \mathrm{Conv}(a,0) \to K \in \mathbb{R} \to 0 < K \to \varphi_{65}$
\item $\varphi_{101} :\equiv \forall n \in \mathbb{N}\,(0 < a_{n}) \to \mathrm{Conv}(a,0) \to K \in \mathbb{R} \to 0 < K \to \varphi_{67}$
\item $\varphi_{102} :\equiv \forall n \in \mathbb{N}\,(0 < a_{n}) \to \mathrm{Conv}(a,0) \to K \in \mathbb{R} \to 0 < K \to \varphi_{68}$
\item $\varphi_{103} :\equiv \forall n \in \mathbb{N}\,(0 < a_{n}) \to \mathrm{Conv}(a,0) \to K \in \mathbb{R} \to 0 < K \to \varphi_{69}$
\item $\varphi_{104} :\equiv \forall n \in \mathbb{N}\,(0 < a_{n}) \to \mathrm{Conv}(a,0) \to K \in \mathbb{R} \to 0 < K \to \varphi_{70}$
\item $\varphi_{105} :\equiv \forall n \in \mathbb{N}\,(0 < a_{n}) \to \mathrm{Conv}(a,0) \to K \in \mathbb{R} \to 0 < K \to \varphi_{71}$
\item $\varphi_{106} :\equiv \forall n \in \mathbb{N}\,(0 < a_{n}) \to \mathrm{Conv}(a,0) \to K \in \mathbb{R} \to 0 < K \to \varphi_{72}$
\item $\varphi_{107} :\equiv \forall n \in \mathbb{N}\,(0 < a_{n}) \to \mathrm{Conv}(a,0) \to K \in \mathbb{R} \to 0 < K \to \varphi_{73}$
\item $\varphi_{108} :\equiv \varphi_{74} \to \neg (N \in \mathbb{N} \land \forall n \in \mathbb{N}\,(N \leq n \to \lvert a_{n} - 0 \rvert < K^{-1}))$
\item $\varphi_{109} :\equiv \forall n \in \mathbb{N}\,(0 < a_{n}) \to \mathrm{Conv}(a,0) \to K \in \mathbb{R} \to 0 < K \to \varphi_{75}$
\item $\varphi_{110} :\equiv \forall n \in \mathbb{N}\,(0 < a_{n}) \to \mathrm{Conv}(a,0) \to K \in \mathbb{R} \to 0 < K \to \varphi_{78}$
\item $\varphi_{111} :\equiv \forall n \in \mathbb{N}\,(0 < a_{n}) \to \mathrm{Conv}(a,0) \to K \in \mathbb{R} \to 0 < K \to \varphi_{80}$
\item $\varphi_{112} :\equiv \forall n \in \mathbb{N}\,(0 < a_{n}) \to \mathrm{Conv}(a,0) \to K \in \mathbb{R} \to 0 < K \to \varphi_{81}$
\item $\varphi_{113} :\equiv \forall n \in \mathbb{N}\,(0 < a_{n}) \to \mathrm{Conv}(a,0) \to K \in \mathbb{R} \to 0 < K \to \varphi_{82}$
\item $\varphi_{114} :\equiv (\forall n \in \mathbb{N}\,(0 < a_{n}) \to \mathrm{Conv}(a,0) \to \mathrm{Seq}(a)) \to \varphi_{20} \to \forall n \in \mathbb{N}\,(0 < a_{n}) \to \varphi_{25}$
\item $\varphi_{115} :\equiv ((K^{-1})^{-1} = K \to (K^{-1})^{-1} < a_{n}^{-1} \to K < a_{n}^{-1}) \to \varphi_{94}$
\item $\varphi_{116} :\equiv (a_{n}^{-1} = a^{\ast}(n) \to K < a_{n}^{-1} \to K < a^{\ast}(n)) \to \varphi_{96}$
\item $\varphi_{117} :\equiv (\forall n \in \mathbb{N}\,(0 < a_{n}) \to \mathrm{Conv}(a,0) \to \varphi_{32}) \to \forall n \in \mathbb{N}\,(0 < a_{n}) \land \mathrm{Conv}(a,0) \to \varphi_{32}$
\item $\varphi_{118} :\equiv (\lvert a_{n} - 0 \rvert = a_{n} \to \lvert a_{n} - 0 \rvert < K^{-1} \to a_{n} < K^{-1}) \to \varphi_{93}$
\item $\varphi_{119} :\equiv \forall K\,(\forall n \in \mathbb{N}\,(0 < a_{n}) \land \mathrm{Conv}(a,0) \to \varphi_{32}) \to \forall K\,(\forall n \in \mathbb{N}\,(0 < a_{n}) \land \mathrm{Conv}(a,0)) \to \varphi_{34}$
\item $\varphi_{120} :\equiv (\lvert a_{n} - 0 \rvert = \lvert a_{n} \rvert \to \lvert a_{n} \rvert = a_{n} \to \lvert a_{n} - 0 \rvert = a_{n}) \to \varphi_{100}$
\item $\varphi_{121} :\equiv (\varphi_{6} \to \neg (N \in \mathbb{N} \land \forall n \in \mathbb{N}\,(N \leq n \to K < a^{\ast}(n)))) \to \varphi_{107}$
\item $\varphi_{122} :\equiv \forall n\,(\varphi_{79} \to n \in \mathbb{N} \to N \leq n \to K < a^{\ast}(n)) \to \forall n\,(\varphi_{79}) \to \forall n \in \mathbb{N}\,(N \leq n \to K < a^{\ast}(n))$
\item $\varphi_{123} :\equiv (\mathrm{DvgP}(a^{\ast}) \leftrightarrow \mathrm{Seq}(a^{\ast}) \land \varphi_{34}) \to \varphi_{4} \to (\forall n \in \mathbb{N}\,(0 < a_{n}) \to \mathrm{Conv}(a,0) \to \varphi_{34}) \to \varphi_{5}$
\item $\varphi_{124} :\equiv (\forall n \in \mathbb{N}\,(0 < a_{n}) \to \mathrm{Conv}(a,0) \to \mathrm{Seq}(a)) \to \varphi_{87} \to \varphi_{92}$
\item $\varphi_{125} :\equiv (\varphi_{74} \to \forall N\,(\varphi_{74})) \to \forall n \in \mathbb{N}\,(0 < a_{n}) \to \mathrm{Conv}(a,0) \to \varphi_{32}$
\item $\varphi_{126} :\equiv \varphi_{103} \to (a^{\ast}(n) = a_{n}^{-1} \to a_{n}^{-1} = a^{\ast}(n)) \to \varphi_{104}$
\item $\varphi_{127} :\equiv (\mathrm{Conv}(a,0) \to K^{-1} \in \mathbb{R} \to \varphi_{31}) \to \varphi_{2} \to \varphi_{33} \to \varphi_{29} \to \forall n \in \mathbb{N}\,(0 < a_{n}) \to \mathrm{Conv}(a,0) \to \varphi_{43}$
\item $\varphi_{128} :\equiv \forall K\,(\forall n \in \mathbb{N}\,(0 < a_{n}) \land \mathrm{Conv}(a,0) \to \varphi_{32}) \to \varphi_{119} \to \varphi_{40} \to \forall n \in \mathbb{N}\,(0 < a_{n}) \to \mathrm{Conv}(a,0) \to \varphi_{34}$
\item $\varphi_{129} :\equiv \forall n\,(\varphi_{79} \to n \in \mathbb{N} \to N \leq n \to K < a^{\ast}(n)) \to \varphi_{122} \to (\varphi_{79} \to \forall n\,(\varphi_{79})) \to \varphi_{99}$
\item $\varphi_{130} :\equiv (\forall n \in \mathbb{N}\,(0 < a_{n}) \to \mathrm{Conv}(a,0) \to \varphi_{43}) \to \forall N\,(\varphi_{108}) \to (\forall N\,(\varphi_{108}) \to \forall N\,(\varphi_{74}) \to \varphi_{10}) \to \varphi_{125}$
\end{abbrevlist}
\end{fullproof}

% 自動生成 (prover/build.js): DvgLe  ノード数 96
\begin{fullproof}{定理2.18(1)}{$\forall a\,\forall b\,(\mathrm{DvgP}(a) \to \mathrm{Seq}(b) \to \forall n \in \mathbb{N}\,(a_{n} \leq b_{n}) \to \mathrm{DvgP}(b))$}{96}\phantomsection\label{tree:DvgLe}
\noindent$\Pi_{1}$:
\begin{ptree}
\AXC{$\varphi_{31}$\lbl{A1}}
\AXC{$\forall n \in \mathbb{N}\,(a_{n} \leq b_{n}) \to n \in \mathbb{N} \to a_{n} \leq b_{n}$\lbl{A2}}
\AXC{$\mathrm{DvgP}(a) \to \varphi_{9}$\lbl{A1}}
\AXC{$\varphi_{49}$\lbl{A1}}
\BIC{$(\forall n \in \mathbb{N}\,(a_{n} \leq b_{n}) \to n \in \mathbb{N} \to a_{n} \leq b_{n}) \to \varphi_{16}$}
\BIC{$\varphi_{16}$}
\AXC{$\varphi_{16} \to \varphi_{31} \to \varphi_{39}$\lbl{A1}}
\BIC{$\varphi_{31} \to \varphi_{39}$}
\BIC{$\varphi_{39}$}
\end{ptree}
\noindent$\Pi_{2}$:
\begin{ptree}
\AXC{$\forall n \in \mathbb{N}\,(N \leq n \to K < a_{n}) \to n \in \mathbb{N} \to N \leq n \to K < a_{n}$\lbl{A2}}
\AXC{$\varphi_{40}$\lbl{A1}}
\AXC{$\varphi_{40} \to \varphi_{44}$\lbl{A1}}
\BIC{$\varphi_{44}$}
\AXC{$\varphi_{44} \to \varphi_{51}$\lbl{A1}}
\BIC{$\varphi_{51}$}
\BIC{$\varphi_{42}$}
\AXC{$\varphi_{42} \to \varphi_{31} \to \varphi_{32} \to \varphi_{36}$\lbl{A1}}
\BIC{$\varphi_{31} \to \varphi_{32} \to \varphi_{36}$}
\end{ptree}
\noindent$\Pi_{3}$:
\begin{ptree}
\AXC{$\forall x\,\forall y\,\forall z\,(x < y \to y \leq z \to x < z)$\lbl{[O3]}}
\AXC{$\varphi_{2}$\lbl{A2}}
\BIC{$\forall y\,\forall z\,(K < y \to y \leq z \to K < z)$}
\AXC{$\varphi_{4}$\lbl{A2}}
\BIC{$\forall z\,(K < a_{n} \to a_{n} \leq z \to K < z)$}
\AXC{$\varphi_{19}$\lbl{A2}}
\BIC{$K < a_{n} \to a_{n} \leq b_{n} \to K < b_{n}$}
\end{ptree}
\noindent$\Pi_{4}$:
\begin{ptree}
\AXC{$\varphi_{32}$\lbl{A1}}
\AXC{$\varphi_{31}$\lbl{A1}}
\AXC{$\stackrel{\Pi_{2}}{\vdots}$}\noLine
\UIC{$\varphi_{31} \to \varphi_{32} \to \varphi_{36}$}
\BIC{$\varphi_{32} \to \varphi_{36}$}
\BIC{$\varphi_{36}$}
\AXC{$\stackrel{\Pi_{3}}{\vdots}$}\noLine
\UIC{$K < a_{n} \to a_{n} \leq b_{n} \to K < b_{n}$}
\AXC{$(K < a_{n} \to a_{n} \leq b_{n} \to K < b_{n}) \to \varphi_{36} \to \varphi_{39} \to \varphi_{37}$\lbl{A1}}
\BIC{$\varphi_{36} \to \varphi_{39} \to \varphi_{37}$}
\BIC{$\varphi_{39} \to \varphi_{37}$}
\end{ptree}
\noindent$\Pi_{5}$:
\begin{ptree}
\AXC{$\varphi_{28} \to \forall n\,(\varphi_{28})$\lbl{A4}}
\AXC{$\varphi_{55}$\lbl{A3}}
\AXC{$\stackrel{\Pi_{1}}{\vdots}$}\noLine
\UIC{$\varphi_{39}$}
\AXC{$\stackrel{\Pi_{4}}{\vdots}$}\noLine
\UIC{$\varphi_{39} \to \varphi_{37}$}
\BIC{$\varphi_{37}$}
\AXC{$\varphi_{37} \to \varphi_{28} \to n \in \mathbb{N} \to N \leq n \to K < b_{n}$\lbl{A1}}
\BIC{$\varphi_{28} \to n \in \mathbb{N} \to N \leq n \to K < b_{n}$}
\UIC{$\forall n\,(\varphi_{28} \to n \in \mathbb{N} \to N \leq n \to K < b_{n})$}
\AXC{$\varphi_{59}$\lbl{A1}}
\BIC{$\varphi_{55} \to (\varphi_{28} \to \forall n\,(\varphi_{28})) \to \varphi_{38}$}
\BIC{$(\varphi_{28} \to \forall n\,(\varphi_{28})) \to \varphi_{38}$}
\BIC{$\varphi_{38}$}
\end{ptree}
\noindent$\Pi_{6}$:
\begin{ptree}
\AXC{$\varphi_{24}$\lbl{A2}}
\AXC{$\stackrel{\Pi_{5}}{\vdots}$}\noLine
\UIC{$\varphi_{38}$}
\AXC{$\varphi_{40}$\lbl{A1}}
\AXC{$\varphi_{53}$\lbl{A1}}
\BIC{$\mathrm{DvgP}(a) \to \mathrm{Seq}(b) \to \forall n \in \mathbb{N}\,(a_{n} \leq b_{n}) \to \varphi_{14}$}
\AXC{$\varphi_{58}$\lbl{A1}}
\BIC{$\varphi_{38} \to \varphi_{41}$}
\BIC{$\varphi_{41}$}
\AXC{$\varphi_{41} \to \varphi_{24} \to \varphi_{43}$\lbl{A1}}
\BIC{$\varphi_{24} \to \varphi_{43}$}
\BIC{$\varphi_{43}$}
\end{ptree}
\noindent$\Pi_{7}$:
\begin{ptree}
\AXC{$\forall a\,(\mathrm{DvgP}(a) \leftrightarrow \mathrm{Seq}(a) \land \varphi_{8})$\lbl{Def$\mathrm{DvgP}$}}
\AXC{$\varphi_{46}$\lbl{A2}}
\BIC{$\mathrm{DvgP}(a) \leftrightarrow \mathrm{Seq}(a) \land \varphi_{8}$}
\AXC{$\mathrm{DvgP}(a) \to \mathrm{DvgP}(a)$\lbl{A1}}
\AXC{$\varphi_{52}$\lbl{A1}}
\BIC{$(\mathrm{DvgP}(a) \leftrightarrow \mathrm{Seq}(a) \land \varphi_{8}) \to \varphi_{35}$}
\BIC{$\varphi_{35}$}
\AXC{$\varphi_{35} \to \varphi_{34} \to \varphi_{33}$\lbl{A1}}
\BIC{$\varphi_{34} \to \varphi_{33}$}
\end{ptree}
\noindent$\Pi_{8}$:
\begin{ptree}
\AXC{$\mathrm{DvgP}(a) \to \varphi_{6}$\lbl{A1}}
\AXC{$\varphi_{10}$\lbl{A1}}
\AXC{$\varphi_{34}$\lbl{A2}}
\AXC{$\stackrel{\Pi_{7}}{\vdots}$}\noLine
\UIC{$\varphi_{34} \to \varphi_{33}$}
\BIC{$\varphi_{33}$}
\AXC{$\varphi_{57}$\lbl{A1}}
\BIC{$\varphi_{10} \to (\mathrm{DvgP}(a) \to \varphi_{6}) \to \mathrm{DvgP}(a) \to \mathrm{Seq}(b) \to \varphi_{26}$}
\BIC{$(\mathrm{DvgP}(a) \to \varphi_{6}) \to \mathrm{DvgP}(a) \to \mathrm{Seq}(b) \to \varphi_{26}$}
\BIC{$\mathrm{DvgP}(a) \to \mathrm{Seq}(b) \to \varphi_{26}$}
\end{ptree}
\noindent$\Pi_{9}$:
\begin{ptree}
\AXC{$\stackrel{\Pi_{8}}{\vdots}$}\noLine
\UIC{$\mathrm{DvgP}(a) \to \mathrm{Seq}(b) \to \varphi_{26}$}
\AXC{$(\mathrm{DvgP}(a) \to \mathrm{Seq}(b) \to \varphi_{26}) \to \varphi_{60}$\lbl{A1}}
\BIC{$\varphi_{60}$}
\end{ptree}
\noindent$\Pi_{10}$:
\begin{ptree}
\AXC{$\varphi_{30} \to \forall N\,(\varphi_{30})$\lbl{A4}}
\AXC{$\varphi_{56}$\lbl{A3}}
\AXC{$\stackrel{\Pi_{6}}{\vdots}$}\noLine
\UIC{$\varphi_{43}$}
\AXC{$\varphi_{43} \to \varphi_{30} \to \neg (N \in \mathbb{N} \land \forall n \in \mathbb{N}\,(N \leq n \to K < a_{n}))$\lbl{A1}}
\BIC{$\varphi_{30} \to \neg (N \in \mathbb{N} \land \forall n \in \mathbb{N}\,(N \leq n \to K < a_{n}))$}
\UIC{$\forall N\,(\varphi_{30} \to \neg (N \in \mathbb{N} \land \forall n \in \mathbb{N}\,(N \leq n \to K < a_{n})))$}
\AXC{$\stackrel{\Pi_{9}}{\vdots}$}\noLine
\UIC{$\varphi_{60}$}
\BIC{$\varphi_{56} \to (\varphi_{30} \to \forall N\,(\varphi_{30})) \to \mathrm{DvgP}(a) \to \mathrm{Seq}(b) \to \varphi_{27}$}
\BIC{$(\varphi_{30} \to \forall N\,(\varphi_{30})) \to \mathrm{DvgP}(a) \to \mathrm{Seq}(b) \to \varphi_{27}$}
\BIC{$\mathrm{DvgP}(a) \to \mathrm{Seq}(b) \to \varphi_{27}$}
\end{ptree}
\noindent$\Pi_{11}$:
\begin{ptree}
\AXC{$\stackrel{\Pi_{10}}{\vdots}$}\noLine
\UIC{$\mathrm{DvgP}(a) \to \mathrm{Seq}(b) \to \varphi_{27}$}
\AXC{$(\mathrm{DvgP}(a) \to \mathrm{Seq}(b) \to \varphi_{27}) \to \varphi_{29}$\lbl{A1}}
\BIC{$\varphi_{29}$}
\UIC{$\forall K\,(\varphi_{29})$}
\AXC{$\forall K\,(\varphi_{29}) \to (\forall K\,(\varphi_{29}) \to \varphi_{1} \to \varphi_{7}) \to \varphi_{50}$\lbl{A1}}
\BIC{$(\forall K\,(\varphi_{29}) \to \varphi_{1} \to \varphi_{7}) \to \varphi_{50}$}
\end{ptree}
\noindent$\Pi_{12}$:
\begin{ptree}
\AXC{$\mathrm{DvgP}(a) \to \mathrm{Seq}(b) \to \mathrm{Seq}(b)$\lbl{A1}}
\AXC{$\forall a\,(\mathrm{DvgP}(a) \leftrightarrow \mathrm{Seq}(a) \land \varphi_{8})$\lbl{Def$\mathrm{DvgP}$}}
\AXC{$\varphi_{47}$\lbl{A2}}
\BIC{$\mathrm{DvgP}(b) \leftrightarrow \mathrm{Seq}(b) \land \varphi_{7}$}
\AXC{$\varphi_{54}$\lbl{A1}}
\BIC{$(\mathrm{DvgP}(a) \to \mathrm{Seq}(b) \to \mathrm{Seq}(b)) \to \varphi_{48}$}
\BIC{$\varphi_{48}$}
\end{ptree}
\noindent 主木:
\begin{ptree}
\AXC{$\mathrm{DvgP}(a) \land \mathrm{Seq}(b) \land \forall n \in \mathbb{N}\,(a_{n} \leq b_{n}) \to \varphi_{1}$\lbl{A4}}
\AXC{$\forall K\,(\varphi_{29}) \to \varphi_{1} \to \varphi_{7}$\lbl{A3}}
\AXC{$\stackrel{\Pi_{11}}{\vdots}$}\noLine
\UIC{$(\forall K\,(\varphi_{29}) \to \varphi_{1} \to \varphi_{7}) \to \varphi_{50}$}
\BIC{$\varphi_{50}$}
\BIC{$\mathrm{DvgP}(a) \to \mathrm{Seq}(b) \to \forall n \in \mathbb{N}\,(a_{n} \leq b_{n}) \to \varphi_{7}$}
\AXC{$\stackrel{\Pi_{12}}{\vdots}$}\noLine
\UIC{$\varphi_{48}$}
\BIC{$\varphi_{3}$}
\UIC{$\forall b\,(\varphi_{3})$}
\UIC{$\forall a\,\forall b\,(\varphi_{3})$}
\end{ptree}
\begin{abbrevlist}
\item $\varphi_{1} :\equiv \forall K\,(\mathrm{DvgP}(a) \land \mathrm{Seq}(b) \land \forall n \in \mathbb{N}\,(a_{n} \leq b_{n}))$
\item $\varphi_{2} :\equiv \forall x\,\forall y\,\forall z\,(x < y \to y \leq z \to x < z) \to \forall y\,\forall z\,(K < y \to y \leq z \to K < z)$
\item $\varphi_{3} :\equiv \mathrm{DvgP}(a) \to \mathrm{Seq}(b) \to \forall n \in \mathbb{N}\,(a_{n} \leq b_{n}) \to \mathrm{DvgP}(b)$
\item $\varphi_{4} :\equiv \forall y\,\forall z\,(K < y \to y \leq z \to K < z) \to \forall z\,(K < a_{n} \to a_{n} \leq z \to K < z)$
\item $\varphi_{5} :\equiv N \in \mathbb{N} \land \forall n \in \mathbb{N}\,(N \leq n \to K < a_{n}) \to n \in \mathbb{N} \to N \leq n \to N \leq n$
\item $\varphi_{6} :\equiv \mathrm{Seq}(b) \to \forall n \in \mathbb{N}\,(a_{n} \leq b_{n}) \to K \in \mathbb{R} \to 0 < K \to 0 < K$
\item $\varphi_{7} :\equiv \forall K \in \mathbb{R}\,(0 < K \to \exists N \in \mathbb{N}\,\forall n \in \mathbb{N}\,(N \leq n \to K < b_{n}))$
\item $\varphi_{8} :\equiv \forall K \in \mathbb{R}\,(0 < K \to \exists N \in \mathbb{N}\,\forall n \in \mathbb{N}\,(N \leq n \to K < a_{n}))$
\item $\varphi_{9} :\equiv \mathrm{Seq}(b) \to \forall n \in \mathbb{N}\,(a_{n} \leq b_{n}) \to \forall n \in \mathbb{N}\,(a_{n} \leq b_{n})$
\item $\varphi_{10} :\equiv \mathrm{DvgP}(a) \to \mathrm{Seq}(b) \to \forall n \in \mathbb{N}\,(a_{n} \leq b_{n}) \to K \in \mathbb{R} \to K \in \mathbb{R}$
\item $\varphi_{11} :\equiv 0 < K \to N \in \mathbb{N} \land \forall n \in \mathbb{N}\,(N \leq n \to K < a_{n}) \to n \in \mathbb{N} \to n \in \mathbb{N}$
\item $\varphi_{12} :\equiv N \in \mathbb{N} \land \forall n \in \mathbb{N}\,(N \leq n \to K < a_{n}) \to n \in \mathbb{N} \to N \leq n \to K < a_{n}$
\item $\varphi_{13} :\equiv N \in \mathbb{N} \land \forall n \in \mathbb{N}\,(N \leq n \to K < a_{n}) \to n \in \mathbb{N} \to N \leq n \to K < b_{n}$
\item $\varphi_{14} :\equiv K \in \mathbb{R} \to 0 < K \to N \in \mathbb{N} \land \forall n \in \mathbb{N}\,(N \leq n \to K < a_{n}) \to N \in \mathbb{N}$
\item $\varphi_{15} :\equiv 0 < K \to K \in \mathbb{R} \to 0 < K \to \exists N \in \mathbb{N}\,\forall n \in \mathbb{N}\,(N \leq n \to K < a_{n})$
\item $\varphi_{16} :\equiv \mathrm{DvgP}(a) \to \mathrm{Seq}(b) \to \forall n \in \mathbb{N}\,(a_{n} \leq b_{n}) \to n \in \mathbb{N} \to a_{n} \leq b_{n}$
\item $\varphi_{17} :\equiv N \in \mathbb{N} \land \forall n \in \mathbb{N}\,(N \leq n \to K < a_{n}) \to \forall n \in \mathbb{N}\,(N \leq n \to K < b_{n})$
\item $\varphi_{18} :\equiv N \in \mathbb{N} \land \forall n \in \mathbb{N}\,(N \leq n \to K < a_{n}) \to n \in \mathbb{N} \to N \leq n \to a_{n} \leq b_{n}$
\item $\varphi_{19} :\equiv \forall z\,(K < a_{n} \to a_{n} \leq z \to K < z) \to K < a_{n} \to a_{n} \leq b_{n} \to K < b_{n}$
\item $\varphi_{20} :\equiv N \in \mathbb{N} \land \forall n \in \mathbb{N}\,(N \leq n \to K < a_{n}) \to N \in \mathbb{N} \land \forall n \in \mathbb{N}\,(N \leq n \to K < a_{n})$
\item $\varphi_{21} :\equiv N \in \mathbb{N} \land \forall n \in \mathbb{N}\,(N \leq n \to K < a_{n}) \to N \in \mathbb{N} \land \forall n \in \mathbb{N}\,(N \leq n \to K < b_{n})$
\item $\varphi_{22} :\equiv N \in \mathbb{N} \land \forall n \in \mathbb{N}\,(N \leq n \to K < a_{n}) \to n \in \mathbb{N} \to N \leq n \to n \in \mathbb{N} \to N \leq n \to K < a_{n}$
\item $\varphi_{23} :\equiv N \in \mathbb{N} \land \forall n \in \mathbb{N}\,(N \leq n \to K < a_{n}) \to \exists N \in \mathbb{N}\,\forall n \in \mathbb{N}\,(N \leq n \to K < b_{n})$
\item $\varphi_{24} :\equiv \forall N\,(\neg (N \in \mathbb{N} \land \forall n \in \mathbb{N}\,(N \leq n \to K < b_{n}))) \to \neg (N \in \mathbb{N} \land \forall n \in \mathbb{N}\,(N \leq n \to K < b_{n}))$
\item $\varphi_{25} :\equiv N \in \mathbb{N} \land \forall n \in \mathbb{N}\,(N \leq n \to K < a_{n}) \to n \in \mathbb{N} \to N \leq n \to \forall n \in \mathbb{N}\,(N \leq n \to K < a_{n})$
\item $\varphi_{26} :\equiv \forall n \in \mathbb{N}\,(a_{n} \leq b_{n}) \to K \in \mathbb{R} \to 0 < K \to \exists N \in \mathbb{N}\,\forall n \in \mathbb{N}\,(N \leq n \to K < a_{n})$
\item $\varphi_{27} :\equiv \forall n \in \mathbb{N}\,(a_{n} \leq b_{n}) \to K \in \mathbb{R} \to 0 < K \to \exists N \in \mathbb{N}\,\forall n \in \mathbb{N}\,(N \leq n \to K < b_{n})$
\item $\varphi_{28} :\equiv \mathrm{DvgP}(a) \land \mathrm{Seq}(b) \land \forall n \in \mathbb{N}\,(a_{n} \leq b_{n}) \land K \in \mathbb{R} \land 0 < K \land (N \in \mathbb{N} \land \forall n \in \mathbb{N}\,(N \leq n \to K < a_{n}))$
\item $\varphi_{29} :\equiv \mathrm{DvgP}(a) \land \mathrm{Seq}(b) \land \forall n \in \mathbb{N}\,(a_{n} \leq b_{n}) \to K \in \mathbb{R} \to 0 < K \to \exists N \in \mathbb{N}\,\forall n \in \mathbb{N}\,(N \leq n \to K < b_{n})$
\item $\varphi_{30} :\equiv \mathrm{DvgP}(a) \land \mathrm{Seq}(b) \land \forall n \in \mathbb{N}\,(a_{n} \leq b_{n}) \land K \in \mathbb{R} \land 0 < K \land \neg \exists N \in \mathbb{N}\,\forall n \in \mathbb{N}\,(N \leq n \to K < b_{n})$
\item $\varphi_{31} :\equiv \mathrm{DvgP}(a) \to \mathrm{Seq}(b) \to \forall n \in \mathbb{N}\,(a_{n} \leq b_{n}) \to K \in \mathbb{R} \to \varphi_{11}$
\item $\varphi_{32} :\equiv \mathrm{DvgP}(a) \to \mathrm{Seq}(b) \to \forall n \in \mathbb{N}\,(a_{n} \leq b_{n}) \to K \in \mathbb{R} \to 0 < K \to \varphi_{5}$
\item $\varphi_{33} :\equiv \mathrm{DvgP}(a) \to \mathrm{Seq}(b) \to \forall n \in \mathbb{N}\,(a_{n} \leq b_{n}) \to K \in \mathbb{R} \to \varphi_{15}$
\item $\varphi_{34} :\equiv \varphi_{8} \to K \in \mathbb{R} \to 0 < K \to \exists N \in \mathbb{N}\,\forall n \in \mathbb{N}\,(N \leq n \to K < a_{n})$
\item $\varphi_{35} :\equiv \mathrm{DvgP}(a) \to \mathrm{Seq}(b) \to \forall n \in \mathbb{N}\,(a_{n} \leq b_{n}) \to K \in \mathbb{R} \to 0 < K \to \varphi_{8}$
\item $\varphi_{36} :\equiv \mathrm{DvgP}(a) \to \mathrm{Seq}(b) \to \forall n \in \mathbb{N}\,(a_{n} \leq b_{n}) \to K \in \mathbb{R} \to 0 < K \to \varphi_{12}$
\item $\varphi_{37} :\equiv \mathrm{DvgP}(a) \to \mathrm{Seq}(b) \to \forall n \in \mathbb{N}\,(a_{n} \leq b_{n}) \to K \in \mathbb{R} \to 0 < K \to \varphi_{13}$
\item $\varphi_{38} :\equiv \mathrm{DvgP}(a) \to \mathrm{Seq}(b) \to \forall n \in \mathbb{N}\,(a_{n} \leq b_{n}) \to K \in \mathbb{R} \to 0 < K \to \varphi_{17}$
\item $\varphi_{39} :\equiv \mathrm{DvgP}(a) \to \mathrm{Seq}(b) \to \forall n \in \mathbb{N}\,(a_{n} \leq b_{n}) \to K \in \mathbb{R} \to 0 < K \to \varphi_{18}$
\item $\varphi_{40} :\equiv \mathrm{DvgP}(a) \to \mathrm{Seq}(b) \to \forall n \in \mathbb{N}\,(a_{n} \leq b_{n}) \to K \in \mathbb{R} \to 0 < K \to \varphi_{20}$
\item $\varphi_{41} :\equiv \mathrm{DvgP}(a) \to \mathrm{Seq}(b) \to \forall n \in \mathbb{N}\,(a_{n} \leq b_{n}) \to K \in \mathbb{R} \to 0 < K \to \varphi_{21}$
\item $\varphi_{42} :\equiv \mathrm{DvgP}(a) \to \mathrm{Seq}(b) \to \forall n \in \mathbb{N}\,(a_{n} \leq b_{n}) \to K \in \mathbb{R} \to 0 < K \to \varphi_{22}$
\item $\varphi_{43} :\equiv \mathrm{DvgP}(a) \to \mathrm{Seq}(b) \to \forall n \in \mathbb{N}\,(a_{n} \leq b_{n}) \to K \in \mathbb{R} \to 0 < K \to \varphi_{23}$
\item $\varphi_{44} :\equiv \mathrm{DvgP}(a) \to \mathrm{Seq}(b) \to \forall n \in \mathbb{N}\,(a_{n} \leq b_{n}) \to K \in \mathbb{R} \to 0 < K \to \varphi_{25}$
\item $\varphi_{45} :\equiv \forall N\,(\varphi_{30}) \to \forall N\,(\neg (N \in \mathbb{N} \land \forall n \in \mathbb{N}\,(N \leq n \to K < a_{n})))$
\item $\varphi_{46} :\equiv \forall a\,(\mathrm{DvgP}(a) \leftrightarrow \mathrm{Seq}(a) \land \varphi_{8}) \to (\mathrm{DvgP}(a) \leftrightarrow \mathrm{Seq}(a) \land \varphi_{8})$
\item $\varphi_{47} :\equiv \forall a\,(\mathrm{DvgP}(a) \leftrightarrow \mathrm{Seq}(a) \land \varphi_{8}) \to (\mathrm{DvgP}(b) \leftrightarrow \mathrm{Seq}(b) \land \varphi_{7})$
\item $\varphi_{48} :\equiv (\mathrm{DvgP}(a) \to \mathrm{Seq}(b) \to \forall n \in \mathbb{N}\,(a_{n} \leq b_{n}) \to \varphi_{7}) \to \varphi_{3}$
\item $\varphi_{49} :\equiv (\mathrm{DvgP}(a) \to \varphi_{9}) \to (\forall n \in \mathbb{N}\,(a_{n} \leq b_{n}) \to n \in \mathbb{N} \to a_{n} \leq b_{n}) \to \varphi_{16}$
\item $\varphi_{50} :\equiv (\mathrm{DvgP}(a) \land \mathrm{Seq}(b) \land \forall n \in \mathbb{N}\,(a_{n} \leq b_{n}) \to \varphi_{1}) \to \mathrm{DvgP}(a) \to \mathrm{Seq}(b) \to \forall n \in \mathbb{N}\,(a_{n} \leq b_{n}) \to \varphi_{7}$
\item $\varphi_{51} :\equiv (\forall n \in \mathbb{N}\,(N \leq n \to K < a_{n}) \to n \in \mathbb{N} \to N \leq n \to K < a_{n}) \to \varphi_{42}$
\item $\varphi_{52} :\equiv (\mathrm{DvgP}(a) \to \mathrm{DvgP}(a)) \to (\mathrm{DvgP}(a) \leftrightarrow \mathrm{Seq}(a) \land \varphi_{8}) \to \varphi_{35}$
\item $\varphi_{53} :\equiv \varphi_{40} \to \mathrm{DvgP}(a) \to \mathrm{Seq}(b) \to \forall n \in \mathbb{N}\,(a_{n} \leq b_{n}) \to \varphi_{14}$
\item $\varphi_{54} :\equiv (\mathrm{DvgP}(b) \leftrightarrow \mathrm{Seq}(b) \land \varphi_{7}) \to (\mathrm{DvgP}(a) \to \mathrm{Seq}(b) \to \mathrm{Seq}(b)) \to \varphi_{48}$
\item $\varphi_{55} :\equiv \forall n\,(\varphi_{28} \to n \in \mathbb{N} \to N \leq n \to K < b_{n}) \to \forall n\,(\varphi_{28}) \to \forall n \in \mathbb{N}\,(N \leq n \to K < b_{n})$
\item $\varphi_{56} :\equiv \forall N\,(\varphi_{30} \to \neg (N \in \mathbb{N} \land \forall n \in \mathbb{N}\,(N \leq n \to K < a_{n}))) \to \varphi_{45}$
\item $\varphi_{57} :\equiv \varphi_{33} \to \varphi_{10} \to (\mathrm{DvgP}(a) \to \varphi_{6}) \to \mathrm{DvgP}(a) \to \mathrm{Seq}(b) \to \varphi_{26}$
\item $\varphi_{58} :\equiv (\mathrm{DvgP}(a) \to \mathrm{Seq}(b) \to \forall n \in \mathbb{N}\,(a_{n} \leq b_{n}) \to \varphi_{14}) \to \varphi_{38} \to \varphi_{41}$
\item $\varphi_{59} :\equiv \forall n\,(\varphi_{28} \to n \in \mathbb{N} \to N \leq n \to K < b_{n}) \to \varphi_{55} \to (\varphi_{28} \to \forall n\,(\varphi_{28})) \to \varphi_{38}$
\item $\varphi_{60} :\equiv \forall N\,(\varphi_{30} \to \neg (N \in \mathbb{N} \land \forall n \in \mathbb{N}\,(N \leq n \to K < a_{n}))) \to \varphi_{56} \to (\varphi_{30} \to \forall N\,(\varphi_{30})) \to \mathrm{DvgP}(a) \to \mathrm{Seq}(b) \to \varphi_{27}$
\end{abbrevlist}
\end{fullproof}

% 自動生成 (prover/build.js): DvgAdd  ノード数 291
\begin{fullproof}{定理2.18(2)}{$\forall a\,\forall b\,\forall C\,(\mathrm{DvgP}(a) \to \mathrm{Seq}(b) \to C \in \mathbb{R} \to \forall n \in \mathbb{N}\,(C \leq b_{n}) \to \mathrm{DvgP}(a \oplus b))$}{291}\phantomsection\label{tree:DvgAdd}
\noindent$\Pi_{1}$:
\begin{ptree}
\AXC{$z = w \to z < a_{n} + b_{n} \to w < a_{n} + b_{n}$\lbl{A6}}
\UIC{$\forall w\,(z = w \to z < a_{n} + b_{n} \to w < a_{n} + b_{n})$}
\UIC{$\forall z\,\forall w\,(z = w \to z < a_{n} + b_{n} \to w < a_{n} + b_{n})$}
\AXC{$\forall z\,\forall w\,(z = w \to z < a_{n} + b_{n} \to w < a_{n} + b_{n}) \to \varphi_{24}$\lbl{A2}}
\BIC{$\varphi_{24}$}
\end{ptree}
\noindent$\Pi_{2}$:
\begin{ptree}
\AXC{$\forall n \in \mathbb{N}\,(C \leq b_{n}) \to n \in \mathbb{N} \to C \leq b_{n}$\lbl{A2}}
\AXC{$\mathrm{DvgP}(a) \to \varphi_{16}$\lbl{A1}}
\AXC{$(\mathrm{DvgP}(a) \to \varphi_{16}) \to \varphi_{82}$\lbl{A1}}
\BIC{$\varphi_{82}$}
\BIC{$\mathrm{DvgP}(a) \to \varphi_{11}$}
\AXC{$\varphi_{150}$\lbl{A1}}
\BIC{$(\mathrm{DvgP}(a) \to \varphi_{93}) \to \mathrm{DvgP}(a) \to \varphi_{96}$}
\end{ptree}
\noindent$\Pi_{3}$:
\begin{ptree}
\AXC{$\mathrm{DvgP}(a) \to \varphi_{119}$\lbl{A1}}
\AXC{$(\mathrm{DvgP}(a) \to \varphi_{119}) \to \mathrm{DvgP}(a) \to \varphi_{123}$\lbl{A1}}
\BIC{$\mathrm{DvgP}(a) \to \varphi_{123}$}
\AXC{$(\mathrm{DvgP}(a) \to \varphi_{123}) \to \varphi_{49} \to \mathrm{DvgP}(a) \to \varphi_{121}$\lbl{A1}}
\BIC{$\varphi_{49} \to \mathrm{DvgP}(a) \to \varphi_{121}$}
\end{ptree}
\noindent$\Pi_{4}$:
\begin{ptree}
\AXC{$\mathrm{DvgP}(a) \to \varphi_{94}$\lbl{A1}}
\AXC{$\mathrm{DvgP}(a) \to \varphi_{93}$\lbl{A1}}
\AXC{$\varphi_{49}$\lbl{A2}}
\AXC{$\stackrel{\Pi_{3}}{\vdots}$}\noLine
\UIC{$\varphi_{49} \to \mathrm{DvgP}(a) \to \varphi_{121}$}
\BIC{$\mathrm{DvgP}(a) \to \varphi_{121}$}
\AXC{$(\mathrm{DvgP}(a) \to \varphi_{121}) \to \varphi_{154}$\lbl{A1}}
\BIC{$\varphi_{154}$}
\BIC{$(\mathrm{DvgP}(a) \to \varphi_{94}) \to \mathrm{DvgP}(a) \to \varphi_{110}$}
\BIC{$\mathrm{DvgP}(a) \to \varphi_{110}$}
\end{ptree}
\noindent$\Pi_{5}$:
\begin{ptree}
\AXC{$\mathrm{DvgP}(a) \to \varphi_{12}$\lbl{A1}}
\AXC{$\forall x\,\forall y\,(x \in \mathbb{R} \to y \in \mathbb{R} \to x - y \in \mathbb{R})$\lbl{[F1]}}
\AXC{$\varphi_{44}$\lbl{A2}}
\BIC{$\forall y\,(K \in \mathbb{R} \to y \in \mathbb{R} \to K - y \in \mathbb{R})$}
\AXC{$\varphi_{36}$\lbl{A2}}
\BIC{$K \in \mathbb{R} \to C \in \mathbb{R} \to K - C \in \mathbb{R}$}
\AXC{$\varphi_{142}$\lbl{A1}}
\BIC{$(\mathrm{DvgP}(a) \to \varphi_{12}) \to \varphi_{132}$}
\BIC{$\varphi_{132}$}
\end{ptree}
\noindent$\Pi_{6}$:
\begin{ptree}
\AXC{$\mathrm{DvgP}(a) \to \mathrm{Seq}(b) \to C \in \mathbb{R} \to C \in \mathbb{R}$\lbl{A1}}
\AXC{$\stackrel{\Pi_{5}}{\vdots}$}\noLine
\UIC{$\varphi_{132}$}
\BIC{$\mathrm{DvgP}(a) \to \varphi_{99}$}
\end{ptree}
\noindent$\Pi_{7}$:
\begin{ptree}
\AXC{$\stackrel{\Pi_{6}}{\vdots}$}\noLine
\UIC{$\mathrm{DvgP}(a) \to \varphi_{99}$}
\AXC{$\forall x \in \mathbb{R}\,(\lvert x \rvert \in \mathbb{R})$\lbl{[A5]}}
\AXC{$\forall x \in \mathbb{R}\,(\lvert x \rvert \in \mathbb{R}) \to K - C \in \mathbb{R} \to \lvert K - C \rvert \in \mathbb{R}$\lbl{A2}}
\BIC{$K - C \in \mathbb{R} \to \lvert K - C \rvert \in \mathbb{R}$}
\AXC{$\varphi_{148}$\lbl{A1}}
\BIC{$(\mathrm{DvgP}(a) \to \varphi_{99}) \to \mathrm{DvgP}(a) \to \varphi_{104}$}
\BIC{$\mathrm{DvgP}(a) \to \varphi_{104}$}
\end{ptree}
\noindent$\Pi_{8}$:
\begin{ptree}
\AXC{$\forall x\,\forall y\,\forall z\,(x < y \to y < z \to x < z)$\lbl{[O2]}}
\AXC{$\varphi_{34}$\lbl{A2}}
\BIC{$\forall y\,\forall z\,(\lvert K - C \rvert < y \to y < z \to \lvert K - C \rvert < z)$}
\AXC{$\forall y\,\forall z\,(\lvert K - C \rvert < y \to y < z \to \lvert K - C \rvert < z) \to \varphi_{19}$\lbl{A2}}
\BIC{$\varphi_{19}$}
\AXC{$\varphi_{19} \to \varphi_{37}$\lbl{A2}}
\BIC{$\varphi_{37}$}
\end{ptree}
\noindent$\Pi_{9}$:
\begin{ptree}
\AXC{$\stackrel{\Pi_{7}}{\vdots}$}\noLine
\UIC{$\mathrm{DvgP}(a) \to \varphi_{104}$}
\AXC{$\forall x \in \mathbb{R}\,(x < x + 1)$\lbl{[N13]}}
\AXC{$\varphi_{29}$\lbl{A2}}
\BIC{$\lvert K - C \rvert \in \mathbb{R} \to \lvert K - C \rvert < \lvert K - C \rvert + 1$}
\AXC{$\varphi_{151}$\lbl{A1}}
\BIC{$(\mathrm{DvgP}(a) \to \varphi_{104}) \to \mathrm{DvgP}(a) \to \varphi_{111}$}
\BIC{$\mathrm{DvgP}(a) \to \varphi_{111}$}
\AXC{$\stackrel{\Pi_{8}}{\vdots}$}\noLine
\UIC{$\varphi_{37}$}
\AXC{$\varphi_{37} \to \varphi_{158}$\lbl{A1}}
\BIC{$\varphi_{158}$}
\BIC{$(\mathrm{DvgP}(a) \to \varphi_{110}) \to \mathrm{DvgP}(a) \to \varphi_{106}$}
\end{ptree}
\noindent$\Pi_{10}$:
\begin{ptree}
\AXC{$\stackrel{\Pi_{6}}{\vdots}$}\noLine
\UIC{$\mathrm{DvgP}(a) \to \varphi_{99}$}
\AXC{$\forall x \in \mathbb{R}\,(-\lvert x \rvert \leq x \land x \leq \lvert x \rvert)$\lbl{[A8]}}
\AXC{$\varphi_{58}$\lbl{A2}}
\BIC{$K - C \in \mathbb{R} \to -\lvert K - C \rvert \leq K - C \land K - C \leq \lvert K - C \rvert$}
\AXC{$\varphi_{152}$\lbl{A1}}
\BIC{$(\mathrm{DvgP}(a) \to \varphi_{99}) \to \mathrm{DvgP}(a) \to \varphi_{120}$}
\BIC{$\mathrm{DvgP}(a) \to \varphi_{120}$}
\end{ptree}
\noindent$\Pi_{11}$:
\begin{ptree}
\AXC{$\forall x\,\forall y\,\forall z\,(x \leq y \to y < z \to x < z)$\lbl{[O4]}}
\AXC{$\varphi_{9}$\lbl{A2}}
\BIC{$\forall y\,\forall z\,(K - C \leq y \to y < z \to K - C < z)$}
\AXC{$\varphi_{45}$\lbl{A2}}
\BIC{$\forall z\,(K - C \leq \lvert K - C \rvert \to \lvert K - C \rvert < z \to K - C < z)$}
\AXC{$\varphi_{71}$\lbl{A2}}
\BIC{$K - C \leq \lvert K - C \rvert \to \lvert K - C \rvert < a_{n} \to K - C < a_{n}$}
\end{ptree}
\noindent$\Pi_{12}$:
\begin{ptree}
\AXC{$\stackrel{\Pi_{11}}{\vdots}$}\noLine
\UIC{$K - C \leq \lvert K - C \rvert \to \lvert K - C \rvert < a_{n} \to K - C < a_{n}$}
\AXC{$(K - C \leq \lvert K - C \rvert \to \lvert K - C \rvert < a_{n} \to K - C < a_{n}) \to \varphi_{156}$\lbl{A1}}
\BIC{$\varphi_{156}$}
\end{ptree}
\noindent$\Pi_{13}$:
\begin{ptree}
\AXC{$\stackrel{\Pi_{10}}{\vdots}$}\noLine
\UIC{$\mathrm{DvgP}(a) \to \varphi_{120}$}
\AXC{$(\mathrm{DvgP}(a) \to \varphi_{120}) \to \mathrm{DvgP}(a) \to \varphi_{102}$\lbl{A1}}
\BIC{$\mathrm{DvgP}(a) \to \varphi_{102}$}
\AXC{$\stackrel{\Pi_{12}}{\vdots}$}\noLine
\UIC{$\varphi_{156}$}
\BIC{$(\mathrm{DvgP}(a) \to \varphi_{106}) \to \mathrm{DvgP}(a) \to \varphi_{103}$}
\end{ptree}
\noindent$\Pi_{14}$:
\begin{ptree}
\AXC{$\stackrel{\Pi_{4}}{\vdots}$}\noLine
\UIC{$\mathrm{DvgP}(a) \to \varphi_{110}$}
\AXC{$\stackrel{\Pi_{9}}{\vdots}$}\noLine
\UIC{$(\mathrm{DvgP}(a) \to \varphi_{110}) \to \mathrm{DvgP}(a) \to \varphi_{106}$}
\BIC{$\mathrm{DvgP}(a) \to \varphi_{106}$}
\AXC{$\stackrel{\Pi_{13}}{\vdots}$}\noLine
\UIC{$(\mathrm{DvgP}(a) \to \varphi_{106}) \to \mathrm{DvgP}(a) \to \varphi_{103}$}
\BIC{$\mathrm{DvgP}(a) \to \varphi_{103}$}
\end{ptree}
\noindent$\Pi_{15}$:
\begin{ptree}
\AXC{$\forall a\,\forall n\,(\mathrm{Seq}(a) \to n \in \mathbb{N} \to a_{n} \in \mathbb{R})$\lbl{[S1]}}
\AXC{$\varphi_{38}$\lbl{A2}}
\BIC{$\forall n\,(\mathrm{Seq}(b) \to n \in \mathbb{N} \to b_{n} \in \mathbb{R})$}
\AXC{$\varphi_{25}$\lbl{A2}}
\BIC{$\mathrm{Seq}(b) \to n \in \mathbb{N} \to b_{n} \in \mathbb{R}$}
\AXC{$(\mathrm{Seq}(b) \to n \in \mathbb{N} \to b_{n} \in \mathbb{R}) \to \varphi_{145}$\lbl{A1}}
\BIC{$\varphi_{145}$}
\end{ptree}
\noindent$\Pi_{16}$:
\begin{ptree}
\AXC{$\mathrm{DvgP}(a) \to \varphi_{93}$\lbl{A1}}
\AXC{$\mathrm{DvgP}(a) \to \mathrm{Seq}(b) \to \mathrm{Seq}(b)$\lbl{A1}}
\AXC{$\stackrel{\Pi_{15}}{\vdots}$}\noLine
\UIC{$\varphi_{145}$}
\BIC{$(\mathrm{DvgP}(a) \to \varphi_{93}) \to \mathrm{DvgP}(a) \to \varphi_{100}$}
\BIC{$\mathrm{DvgP}(a) \to \varphi_{100}$}
\end{ptree}
\noindent$\Pi_{17}$:
\begin{ptree}
\AXC{$\forall a\,(\mathrm{DvgP}(a) \leftrightarrow \mathrm{Seq}(a) \land \varphi_{15})$\lbl{Def$\mathrm{DvgP}$}}
\AXC{$\varphi_{114}$\lbl{A2}}
\BIC{$\mathrm{DvgP}(a) \leftrightarrow \mathrm{Seq}(a) \land \varphi_{15}$}
\AXC{$\mathrm{DvgP}(a) \to \mathrm{DvgP}(a)$\lbl{A1}}
\AXC{$\varphi_{109}$\lbl{A1}}
\BIC{$(\mathrm{DvgP}(a) \leftrightarrow \mathrm{Seq}(a) \land \varphi_{15}) \to \varphi_{7}$}
\BIC{$\varphi_{7}$}
\end{ptree}
\noindent$\Pi_{18}$:
\begin{ptree}
\AXC{$\mathrm{DvgP}(a) \to \varphi_{93}$\lbl{A1}}
\AXC{$\stackrel{\Pi_{17}}{\vdots}$}\noLine
\UIC{$\varphi_{7}$}
\AXC{$\forall a\,\forall n\,(\mathrm{Seq}(a) \to n \in \mathbb{N} \to a_{n} \in \mathbb{R})$\lbl{[S1]}}
\AXC{$\varphi_{39}$\lbl{A2}}
\BIC{$\forall n\,(\mathrm{Seq}(a) \to n \in \mathbb{N} \to a_{n} \in \mathbb{R})$}
\AXC{$\varphi_{26}$\lbl{A2}}
\BIC{$\mathrm{Seq}(a) \to n \in \mathbb{N} \to a_{n} \in \mathbb{R}$}
\AXC{$\varphi_{153}$\lbl{A1}}
\BIC{$\varphi_{7} \to (\mathrm{DvgP}(a) \to \varphi_{93}) \to \mathrm{DvgP}(a) \to \varphi_{101}$}
\BIC{$(\mathrm{DvgP}(a) \to \varphi_{93}) \to \mathrm{DvgP}(a) \to \varphi_{101}$}
\BIC{$\mathrm{DvgP}(a) \to \varphi_{101}$}
\end{ptree}
\noindent$\Pi_{19}$:
\begin{ptree}
\AXC{$\varphi_{46}$\lbl{[F20]}}
\AXC{$\varphi_{46} \to \varphi_{61}$\lbl{A2}}
\BIC{$\varphi_{61}$}
\AXC{$\varphi_{61} \to \varphi_{74}$\lbl{A2}}
\BIC{$\varphi_{74}$}
\AXC{$\varphi_{74} \to \varphi_{72}$\lbl{A2}}
\BIC{$\varphi_{72}$}
\AXC{$\varphi_{72} \to K - C \in \mathbb{R} \to a_{n} \in \mathbb{R} \to \varphi_{56}$\lbl{A2}}
\BIC{$K - C \in \mathbb{R} \to a_{n} \in \mathbb{R} \to \varphi_{56}$}
\AXC{$\varphi_{164}$\lbl{A1}}
\BIC{$(\mathrm{DvgP}(a) \to \varphi_{99}) \to (\mathrm{DvgP}(a) \to \varphi_{101}) \to \varphi_{161}$}
\end{ptree}
\noindent$\Pi_{20}$:
\begin{ptree}
\AXC{$\stackrel{\Pi_{18}}{\vdots}$}\noLine
\UIC{$\mathrm{DvgP}(a) \to \varphi_{101}$}
\AXC{$\stackrel{\Pi_{6}}{\vdots}$}\noLine
\UIC{$\mathrm{DvgP}(a) \to \varphi_{99}$}
\AXC{$\stackrel{\Pi_{19}}{\vdots}$}\noLine
\UIC{$(\mathrm{DvgP}(a) \to \varphi_{99}) \to (\mathrm{DvgP}(a) \to \varphi_{101}) \to \varphi_{161}$}
\BIC{$(\mathrm{DvgP}(a) \to \varphi_{101}) \to \varphi_{161}$}
\BIC{$\varphi_{161}$}
\end{ptree}
\noindent$\Pi_{21}$:
\begin{ptree}
\AXC{$\stackrel{\Pi_{14}}{\vdots}$}\noLine
\UIC{$\mathrm{DvgP}(a) \to \varphi_{103}$}
\AXC{$\stackrel{\Pi_{16}}{\vdots}$}\noLine
\UIC{$\mathrm{DvgP}(a) \to \varphi_{100}$}
\AXC{$\mathrm{DvgP}(a) \to \mathrm{Seq}(b) \to C \in \mathbb{R} \to C \in \mathbb{R}$\lbl{A1}}
\AXC{$\stackrel{\Pi_{20}}{\vdots}$}\noLine
\UIC{$\varphi_{161}$}
\BIC{$(\mathrm{DvgP}(a) \to \varphi_{100}) \to \varphi_{157}$}
\BIC{$\varphi_{157}$}
\BIC{$(\mathrm{DvgP}(a) \to \varphi_{96}) \to \mathrm{DvgP}(a) \to \varphi_{112}$}
\end{ptree}
\noindent$\Pi_{22}$:
\begin{ptree}
\AXC{$\mathrm{DvgP}(a) \to \varphi_{93}$\lbl{A1}}
\AXC{$\stackrel{\Pi_{2}}{\vdots}$}\noLine
\UIC{$(\mathrm{DvgP}(a) \to \varphi_{93}) \to \mathrm{DvgP}(a) \to \varphi_{96}$}
\BIC{$\mathrm{DvgP}(a) \to \varphi_{96}$}
\AXC{$\stackrel{\Pi_{21}}{\vdots}$}\noLine
\UIC{$(\mathrm{DvgP}(a) \to \varphi_{96}) \to \mathrm{DvgP}(a) \to \varphi_{112}$}
\BIC{$\mathrm{DvgP}(a) \to \varphi_{112}$}
\end{ptree}
\noindent$\Pi_{23}$:
\begin{ptree}
\AXC{$\mathrm{DvgP}(a) \to \varphi_{12}$\lbl{A1}}
\AXC{$\forall x\,\forall y\,(x \in \mathbb{R} \to y \in \mathbb{R} \to x - y + y = x)$\lbl{[F12]}}
\AXC{$\varphi_{41}$\lbl{A2}}
\BIC{$\forall y\,(K \in \mathbb{R} \to y \in \mathbb{R} \to K - y + y = K)$}
\AXC{$\varphi_{30}$\lbl{A2}}
\BIC{$K \in \mathbb{R} \to C \in \mathbb{R} \to K - C + C = K$}
\AXC{$\varphi_{140}$\lbl{A1}}
\BIC{$(\mathrm{DvgP}(a) \to \varphi_{12}) \to \varphi_{131}$}
\BIC{$\varphi_{131}$}
\end{ptree}
\noindent$\Pi_{24}$:
\begin{ptree}
\AXC{$\mathrm{DvgP}(a) \to \mathrm{Seq}(b) \to C \in \mathbb{R} \to C \in \mathbb{R}$\lbl{A1}}
\AXC{$\stackrel{\Pi_{23}}{\vdots}$}\noLine
\UIC{$\varphi_{131}$}
\BIC{$\mathrm{DvgP}(a) \to \varphi_{98}$}
\AXC{$(\mathrm{DvgP}(a) \to \varphi_{98}) \to (\mathrm{DvgP}(a) \to \varphi_{112}) \to \varphi_{134}$\lbl{A1}}
\BIC{$(\mathrm{DvgP}(a) \to \varphi_{112}) \to \varphi_{134}$}
\end{ptree}
\noindent$\Pi_{25}$:
\begin{ptree}
\AXC{$\stackrel{\Pi_{1}}{\vdots}$}\noLine
\UIC{$\varphi_{24}$}
\AXC{$\varphi_{24} \to K - C + C = K \to K - C + C < a_{n} + b_{n} \to K < a_{n} + b_{n}$\lbl{A2}}
\BIC{$K - C + C = K \to K - C + C < a_{n} + b_{n} \to K < a_{n} + b_{n}$}
\AXC{$\stackrel{\Pi_{22}}{\vdots}$}\noLine
\UIC{$\mathrm{DvgP}(a) \to \varphi_{112}$}
\AXC{$\stackrel{\Pi_{24}}{\vdots}$}\noLine
\UIC{$(\mathrm{DvgP}(a) \to \varphi_{112}) \to \varphi_{134}$}
\BIC{$\varphi_{134}$}
\BIC{$\mathrm{DvgP}(a) \to \varphi_{107}$}
\end{ptree}
\noindent$\Pi_{26}$:
\begin{ptree}
\AXC{$x = x$\lbl{A5}}
\AXC{$x = y \to x = x \to y = x$\lbl{A6}}
\AXC{$(x = y \to x = x \to y = x) \to x = x \to x = y \to y = x$\lbl{A1}}
\BIC{$x = x \to x = y \to y = x$}
\BIC{$x = y \to y = x$}
\UIC{$\forall y\,(x = y \to y = x)$}
\UIC{$\forall x\,\forall y\,(x = y \to y = x)$}
\AXC{$\varphi_{4}$\lbl{A2}}
\BIC{$\forall y\,(a \oplus b(n) = y \to y = a \oplus b(n))$}
\AXC{$\varphi_{76}$\lbl{A2}}
\BIC{$a \oplus b(n) = a_{n} + b_{n} \to a_{n} + b_{n} = a \oplus b(n)$}
\end{ptree}
\noindent$\Pi_{27}$:
\begin{ptree}
\AXC{$\forall a\,\forall b\,(\mathrm{Seq}(a) \to \varphi_{21})$\lbl{[S2]}}
\AXC{$\forall a\,\forall b\,(\mathrm{Seq}(a) \to \varphi_{21}) \to \forall b\,(\mathrm{Seq}(a) \to \varphi_{21})$\lbl{A2}}
\BIC{$\forall b\,(\mathrm{Seq}(a) \to \varphi_{21})$}
\end{ptree}
\noindent$\Pi_{28}$:
\begin{ptree}
\AXC{$\stackrel{\Pi_{17}}{\vdots}$}\noLine
\UIC{$\varphi_{7}$}
\AXC{$\stackrel{\Pi_{27}}{\vdots}$}\noLine
\UIC{$\forall b\,(\mathrm{Seq}(a) \to \varphi_{21})$}
\AXC{$\forall b\,(\mathrm{Seq}(a) \to \varphi_{21}) \to \mathrm{Seq}(a) \to \varphi_{21}$\lbl{A2}}
\BIC{$\mathrm{Seq}(a) \to \varphi_{21}$}
\AXC{$(\mathrm{Seq}(a) \to \varphi_{21}) \to \varphi_{7} \to \varphi_{95}$\lbl{A1}}
\BIC{$\varphi_{7} \to \varphi_{95}$}
\BIC{$\varphi_{95}$}
\end{ptree}
\noindent$\Pi_{29}$:
\begin{ptree}
\AXC{$\mathrm{DvgP}(a) \to \mathrm{Seq}(b) \to \mathrm{Seq}(b)$\lbl{A1}}
\AXC{$\stackrel{\Pi_{28}}{\vdots}$}\noLine
\UIC{$\varphi_{95}$}
\BIC{$\mathrm{DvgP}(a) \to \mathrm{Seq}(b) \to C \in \mathbb{R} \to \varphi_{52}$}
\end{ptree}
\noindent$\Pi_{30}$:
\begin{ptree}
\AXC{$\varphi_{60}$\lbl{A2}}
\AXC{$\stackrel{\Pi_{29}}{\vdots}$}\noLine
\UIC{$\mathrm{DvgP}(a) \to \mathrm{Seq}(b) \to C \in \mathbb{R} \to \varphi_{52}$}
\AXC{$\varphi_{141}$\lbl{A1}}
\BIC{$\mathrm{DvgP}(a) \to \varphi_{126}$}
\AXC{$(\mathrm{DvgP}(a) \to \varphi_{126}) \to \varphi_{60} \to \mathrm{DvgP}(a) \to \varphi_{124}$\lbl{A1}}
\BIC{$\varphi_{60} \to \mathrm{DvgP}(a) \to \varphi_{124}$}
\BIC{$\mathrm{DvgP}(a) \to \varphi_{124}$}
\end{ptree}
\noindent$\Pi_{31}$:
\begin{ptree}
\AXC{$\mathrm{DvgP}(a) \to \varphi_{93}$\lbl{A1}}
\AXC{$\stackrel{\Pi_{30}}{\vdots}$}\noLine
\UIC{$\mathrm{DvgP}(a) \to \varphi_{124}$}
\AXC{$\varphi_{159}$\lbl{A1}}
\BIC{$(\mathrm{DvgP}(a) \to \varphi_{93}) \to \mathrm{DvgP}(a) \to \varphi_{115}$}
\BIC{$\mathrm{DvgP}(a) \to \varphi_{115}$}
\AXC{$(\mathrm{DvgP}(a) \to \varphi_{115}) \to \varphi_{138}$\lbl{A1}}
\BIC{$\varphi_{138}$}
\end{ptree}
\noindent$\Pi_{32}$:
\begin{ptree}
\AXC{$\stackrel{\Pi_{26}}{\vdots}$}\noLine
\UIC{$a \oplus b(n) = a_{n} + b_{n} \to a_{n} + b_{n} = a \oplus b(n)$}
\AXC{$\stackrel{\Pi_{31}}{\vdots}$}\noLine
\UIC{$\varphi_{138}$}
\BIC{$\mathrm{DvgP}(a) \to \varphi_{116}$}
\AXC{$(\mathrm{DvgP}(a) \to \varphi_{116}) \to (\mathrm{DvgP}(a) \to \varphi_{107}) \to \varphi_{137}$\lbl{A1}}
\BIC{$(\mathrm{DvgP}(a) \to \varphi_{107}) \to \varphi_{137}$}
\end{ptree}
\noindent$\Pi_{33}$:
\begin{ptree}
\AXC{$z = w \to K < z \to K < w$\lbl{A6}}
\UIC{$\forall w\,(z = w \to K < z \to K < w)$}
\UIC{$\forall z\,\forall w\,(z = w \to K < z \to K < w)$}
\AXC{$\varphi_{42}$\lbl{A2}}
\BIC{$\forall w\,(a_{n} + b_{n} = w \to K < a_{n} + b_{n} \to K < w)$}
\AXC{$\varphi_{92}$\lbl{A2}}
\BIC{$a_{n} + b_{n} = a \oplus b(n) \to K < a_{n} + b_{n} \to K < a \oplus b(n)$}
\AXC{$\stackrel{\Pi_{25}}{\vdots}$}\noLine
\UIC{$\mathrm{DvgP}(a) \to \varphi_{107}$}
\AXC{$\stackrel{\Pi_{32}}{\vdots}$}\noLine
\UIC{$(\mathrm{DvgP}(a) \to \varphi_{107}) \to \varphi_{137}$}
\BIC{$\varphi_{137}$}
\BIC{$\mathrm{DvgP}(a) \to \varphi_{105}$}
\AXC{$\varphi_{144}$\lbl{A1}}
\BIC{$\varphi_{88} \to n \in \mathbb{N} \to N \leq n \to K < a \oplus b(n)$}
\UIC{$\forall n\,(\varphi_{88} \to n \in \mathbb{N} \to N \leq n \to K < a \oplus b(n))$}
\end{ptree}
\noindent$\Pi_{34}$:
\begin{ptree}
\AXC{$\varphi_{88} \to \forall n\,(\varphi_{88})$\lbl{A4}}
\AXC{$\varphi_{146}$\lbl{A3}}
\AXC{$\stackrel{\Pi_{33}}{\vdots}$}\noLine
\UIC{$\forall n\,(\varphi_{88} \to n \in \mathbb{N} \to N \leq n \to K < a \oplus b(n))$}
\AXC{$\varphi_{162}$\lbl{A1}}
\BIC{$\varphi_{146} \to (\varphi_{88} \to \forall n\,(\varphi_{88})) \to \mathrm{DvgP}(a) \to \varphi_{108}$}
\BIC{$(\varphi_{88} \to \forall n\,(\varphi_{88})) \to \mathrm{DvgP}(a) \to \varphi_{108}$}
\BIC{$\mathrm{DvgP}(a) \to \varphi_{108}$}
\end{ptree}
\noindent$\Pi_{35}$:
\begin{ptree}
\AXC{$\stackrel{\Pi_{34}}{\vdots}$}\noLine
\UIC{$\mathrm{DvgP}(a) \to \varphi_{108}$}
\AXC{$\mathrm{DvgP}(a) \to \varphi_{119}$\lbl{A1}}
\AXC{$(\mathrm{DvgP}(a) \to \varphi_{119}) \to \mathrm{DvgP}(a) \to \varphi_{89}$\lbl{A1}}
\BIC{$\mathrm{DvgP}(a) \to \varphi_{89}$}
\AXC{$\varphi_{155}$\lbl{A1}}
\BIC{$(\mathrm{DvgP}(a) \to \varphi_{108}) \to \mathrm{DvgP}(a) \to \varphi_{113}$}
\BIC{$\mathrm{DvgP}(a) \to \varphi_{113}$}
\end{ptree}
\noindent$\Pi_{36}$:
\begin{ptree}
\AXC{$\varphi_{2} \to \neg (N \in \mathbb{N} \land \forall n \in \mathbb{N}\,(N \leq n \to K < a \oplus b(n)))$\lbl{A2}}
\AXC{$\stackrel{\Pi_{35}}{\vdots}$}\noLine
\UIC{$\mathrm{DvgP}(a) \to \varphi_{113}$}
\AXC{$(\mathrm{DvgP}(a) \to \varphi_{113}) \to \varphi_{139}$\lbl{A1}}
\BIC{$\varphi_{139}$}
\BIC{$\mathrm{DvgP}(a) \to \varphi_{117}$}
\end{ptree}
\noindent$\Pi_{37}$:
\begin{ptree}
\AXC{$\mathrm{DvgP}(a) \to \varphi_{12}$\lbl{A1}}
\AXC{$\forall x\,\forall y\,(x \in \mathbb{R} \to y \in \mathbb{R} \to x - y \in \mathbb{R})$\lbl{[F1]}}
\AXC{$\varphi_{44}$\lbl{A2}}
\BIC{$\forall y\,(K \in \mathbb{R} \to y \in \mathbb{R} \to K - y \in \mathbb{R})$}
\AXC{$\varphi_{36}$\lbl{A2}}
\BIC{$K \in \mathbb{R} \to C \in \mathbb{R} \to K - C \in \mathbb{R}$}
\AXC{$\varphi_{135}$\lbl{A1}}
\BIC{$(\mathrm{DvgP}(a) \to \varphi_{12}) \to \varphi_{86}$}
\BIC{$\varphi_{86}$}
\end{ptree}
\noindent$\Pi_{38}$:
\begin{ptree}
\AXC{$\mathrm{DvgP}(a) \to \mathrm{Seq}(b) \to C \in \mathbb{R} \to C \in \mathbb{R}$\lbl{A1}}
\AXC{$\stackrel{\Pi_{37}}{\vdots}$}\noLine
\UIC{$\varphi_{86}$}
\BIC{$\mathrm{DvgP}(a) \to \mathrm{Seq}(b) \to \varphi_{20}$}
\end{ptree}
\noindent$\Pi_{39}$:
\begin{ptree}
\AXC{$\forall x \in \mathbb{R}\,(\lvert x \rvert \in \mathbb{R})$\lbl{[A5]}}
\AXC{$\forall x \in \mathbb{R}\,(\lvert x \rvert \in \mathbb{R}) \to K - C \in \mathbb{R} \to \lvert K - C \rvert \in \mathbb{R}$\lbl{A2}}
\BIC{$K - C \in \mathbb{R} \to \lvert K - C \rvert \in \mathbb{R}$}
\end{ptree}
\noindent$\Pi_{40}$:
\begin{ptree}
\AXC{$\stackrel{\Pi_{38}}{\vdots}$}\noLine
\UIC{$\mathrm{DvgP}(a) \to \mathrm{Seq}(b) \to \varphi_{20}$}
\AXC{$\stackrel{\Pi_{39}}{\vdots}$}\noLine
\UIC{$K - C \in \mathbb{R} \to \lvert K - C \rvert \in \mathbb{R}$}
\AXC{$(K - C \in \mathbb{R} \to \lvert K - C \rvert \in \mathbb{R}) \to \varphi_{122}$\lbl{A1}}
\BIC{$\varphi_{122}$}
\BIC{$\mathrm{DvgP}(a) \to \mathrm{Seq}(b) \to \varphi_{28}$}
\end{ptree}
\noindent$\Pi_{41}$:
\begin{ptree}
\AXC{$\forall x \in \mathbb{R}\,(x < x + 1)$\lbl{[N13]}}
\AXC{$\varphi_{29}$\lbl{A2}}
\BIC{$\lvert K - C \rvert \in \mathbb{R} \to \lvert K - C \rvert < \lvert K - C \rvert + 1$}
\AXC{$(\lvert K - C \rvert \in \mathbb{R} \to \lvert K - C \rvert < \lvert K - C \rvert + 1) \to \varphi_{128}$\lbl{A1}}
\BIC{$\varphi_{128}$}
\end{ptree}
\noindent$\Pi_{42}$:
\begin{ptree}
\AXC{$\forall x \in \mathbb{R}\,(0 \leq \lvert x \rvert)$\lbl{[A6]}}
\AXC{$\forall x \in \mathbb{R}\,(0 \leq \lvert x \rvert) \to K - C \in \mathbb{R} \to 0 \leq \lvert K - C \rvert$\lbl{A2}}
\BIC{$K - C \in \mathbb{R} \to 0 \leq \lvert K - C \rvert$}
\AXC{$(K - C \in \mathbb{R} \to 0 \leq \lvert K - C \rvert) \to \varphi_{118}$\lbl{A1}}
\BIC{$\varphi_{118}$}
\end{ptree}
\noindent$\Pi_{43}$:
\begin{ptree}
\AXC{$\forall x\,\forall y\,\forall z\,(x \leq y \to y < z \to x < z)$\lbl{[O4]}}
\AXC{$\varphi_{3}$\lbl{A2}}
\BIC{$\forall y\,\forall z\,(0 \leq y \to y < z \to 0 < z)$}
\AXC{$\varphi_{14}$\lbl{A2}}
\BIC{$\forall z\,(0 \leq \lvert K - C \rvert \to \lvert K - C \rvert < z \to 0 < z)$}
\AXC{$\forall z\,(0 \leq \lvert K - C \rvert \to \lvert K - C \rvert < z \to 0 < z) \to \varphi_{10}$\lbl{A2}}
\BIC{$\varphi_{10}$}
\end{ptree}
\noindent$\Pi_{44}$:
\begin{ptree}
\AXC{$\stackrel{\Pi_{38}}{\vdots}$}\noLine
\UIC{$\mathrm{DvgP}(a) \to \mathrm{Seq}(b) \to \varphi_{20}$}
\AXC{$\stackrel{\Pi_{42}}{\vdots}$}\noLine
\UIC{$\varphi_{118}$}
\BIC{$\mathrm{DvgP}(a) \to \mathrm{Seq}(b) \to \varphi_{18}$}
\AXC{$\stackrel{\Pi_{43}}{\vdots}$}\noLine
\UIC{$\varphi_{10}$}
\AXC{$\varphi_{10} \to (\mathrm{DvgP}(a) \to \mathrm{Seq}(b) \to \varphi_{18}) \to \varphi_{129}$\lbl{A1}}
\BIC{$(\mathrm{DvgP}(a) \to \mathrm{Seq}(b) \to \varphi_{18}) \to \varphi_{129}$}
\BIC{$\varphi_{129}$}
\end{ptree}
\noindent$\Pi_{45}$:
\begin{ptree}
\AXC{$\forall x\,\forall y\,(x \in \mathbb{R} \to y \in \mathbb{R} \to x + y \in \mathbb{R})$\lbl{(R1)}}
\AXC{$\varphi_{66}$\lbl{A2}}
\BIC{$\forall y\,(\lvert K - C \rvert \in \mathbb{R} \to y \in \mathbb{R} \to \lvert K - C \rvert + y \in \mathbb{R})$}
\AXC{$\varphi_{78}$\lbl{A2}}
\BIC{$\lvert K - C \rvert \in \mathbb{R} \to 1 \in \mathbb{R} \to \lvert K - C \rvert + 1 \in \mathbb{R}$}
\AXC{$(\lvert K - C \rvert \in \mathbb{R} \to 1 \in \mathbb{R} \to \lvert K - C \rvert + 1 \in \mathbb{R}) \to \varphi_{133}$\lbl{A1}}
\BIC{$\varphi_{133}$}
\end{ptree}
\noindent$\Pi_{46}$:
\begin{ptree}
\AXC{$\forall a\,(\mathrm{DvgP}(a) \leftrightarrow \mathrm{Seq}(a) \land \varphi_{15})$\lbl{Def$\mathrm{DvgP}$}}
\AXC{$\varphi_{114}$\lbl{A2}}
\BIC{$\mathrm{DvgP}(a) \leftrightarrow \mathrm{Seq}(a) \land \varphi_{15}$}
\AXC{$\mathrm{DvgP}(a) \to \mathrm{DvgP}(a)$\lbl{A1}}
\AXC{$\varphi_{136}$\lbl{A1}}
\BIC{$(\mathrm{DvgP}(a) \leftrightarrow \mathrm{Seq}(a) \land \varphi_{15}) \to \mathrm{DvgP}(a) \to \varphi_{91}$}
\BIC{$\mathrm{DvgP}(a) \to \varphi_{91}$}
\end{ptree}
\noindent$\Pi_{47}$:
\begin{ptree}
\AXC{$\varphi_{15} \to \lvert K - C \rvert + 1 \in \mathbb{R} \to \varphi_{40}$\lbl{A2}}
\AXC{$\stackrel{\Pi_{46}}{\vdots}$}\noLine
\UIC{$\mathrm{DvgP}(a) \to \varphi_{91}$}
\AXC{$(\mathrm{DvgP}(a) \to \varphi_{91}) \to \varphi_{147}$\lbl{A1}}
\BIC{$\varphi_{147}$}
\BIC{$\mathrm{DvgP}(a) \to \mathrm{Seq}(b) \to C \in \mathbb{R} \to \varphi_{97}$}
\end{ptree}
\noindent$\Pi_{48}$:
\begin{ptree}
\AXC{$\stackrel{\Pi_{47}}{\vdots}$}\noLine
\UIC{$\mathrm{DvgP}(a) \to \mathrm{Seq}(b) \to C \in \mathbb{R} \to \varphi_{97}$}
\AXC{$(\mathrm{DvgP}(a) \to \mathrm{Seq}(b) \to C \in \mathbb{R} \to \varphi_{97}) \to \varphi_{143}$\lbl{A1}}
\BIC{$\varphi_{143}$}
\end{ptree}
\noindent$\Pi_{49}$:
\begin{ptree}
\AXC{$1 \in \mathbb{R}$\lbl{(R1)}}
\AXC{$\stackrel{\Pi_{40}}{\vdots}$}\noLine
\UIC{$\mathrm{DvgP}(a) \to \mathrm{Seq}(b) \to \varphi_{28}$}
\AXC{$\stackrel{\Pi_{45}}{\vdots}$}\noLine
\UIC{$\varphi_{133}$}
\BIC{$1 \in \mathbb{R} \to \mathrm{DvgP}(a) \to \mathrm{Seq}(b) \to C \in \mathbb{R} \to \varphi_{17}$}
\BIC{$\mathrm{DvgP}(a) \to \mathrm{Seq}(b) \to C \in \mathbb{R} \to \varphi_{17}$}
\AXC{$\stackrel{\Pi_{48}}{\vdots}$}\noLine
\UIC{$\varphi_{143}$}
\BIC{$(\mathrm{DvgP}(a) \to \mathrm{Seq}(b) \to \varphi_{31}) \to \varphi_{90}$}
\end{ptree}
\noindent$\Pi_{50}$:
\begin{ptree}
\AXC{$\stackrel{\Pi_{40}}{\vdots}$}\noLine
\UIC{$\mathrm{DvgP}(a) \to \mathrm{Seq}(b) \to \varphi_{28}$}
\AXC{$\stackrel{\Pi_{41}}{\vdots}$}\noLine
\UIC{$\varphi_{128}$}
\BIC{$\mathrm{DvgP}(a) \to \mathrm{Seq}(b) \to C \in \mathbb{R} \to \varphi_{23}$}
\AXC{$\stackrel{\Pi_{44}}{\vdots}$}\noLine
\UIC{$\varphi_{129}$}
\BIC{$\mathrm{DvgP}(a) \to \mathrm{Seq}(b) \to \varphi_{31}$}
\AXC{$\stackrel{\Pi_{49}}{\vdots}$}\noLine
\UIC{$(\mathrm{DvgP}(a) \to \mathrm{Seq}(b) \to \varphi_{31}) \to \varphi_{90}$}
\BIC{$\varphi_{90}$}
\AXC{$\varphi_{90} \to \varphi_{163}$\lbl{A1}}
\BIC{$\varphi_{163}$}
\end{ptree}
\noindent$\Pi_{51}$:
\begin{ptree}
\AXC{$\varphi_{85} \to \forall N\,(\varphi_{85})$\lbl{A4}}
\AXC{$\forall N\,(\varphi_{125}) \to \forall N\,(\varphi_{85}) \to \varphi_{5}$\lbl{A3}}
\AXC{$\stackrel{\Pi_{36}}{\vdots}$}\noLine
\UIC{$\mathrm{DvgP}(a) \to \varphi_{117}$}
\AXC{$(\mathrm{DvgP}(a) \to \varphi_{117}) \to \varphi_{125}$\lbl{A1}}
\BIC{$\varphi_{125}$}
\UIC{$\forall N\,(\varphi_{125})$}
\AXC{$\stackrel{\Pi_{50}}{\vdots}$}\noLine
\UIC{$\varphi_{163}$}
\BIC{$(\forall N\,(\varphi_{125}) \to \forall N\,(\varphi_{85}) \to \varphi_{5}) \to (\varphi_{85} \to \forall N\,(\varphi_{85})) \to \varphi_{84}$}
\BIC{$(\varphi_{85} \to \forall N\,(\varphi_{85})) \to \varphi_{84}$}
\BIC{$\varphi_{84}$}
\end{ptree}
\noindent$\Pi_{52}$:
\begin{ptree}
\AXC{$\varphi_{1} \to \forall K\,(\varphi_{1})$\lbl{A4}}
\AXC{$\forall K\,(\varphi_{1} \to \varphi_{22}) \to \forall K\,(\varphi_{1}) \to \varphi_{27}$\lbl{A3}}
\AXC{$\stackrel{\Pi_{51}}{\vdots}$}\noLine
\UIC{$\varphi_{84}$}
\AXC{$\varphi_{84} \to \varphi_{1} \to \varphi_{22}$\lbl{A1}}
\BIC{$\varphi_{1} \to \varphi_{22}$}
\UIC{$\forall K\,(\varphi_{1} \to \varphi_{22})$}
\AXC{$\varphi_{160}$\lbl{A1}}
\BIC{$(\forall K\,(\varphi_{1} \to \varphi_{22}) \to \forall K\,(\varphi_{1}) \to \varphi_{27}) \to (\varphi_{1} \to \forall K\,(\varphi_{1})) \to \varphi_{87}$}
\BIC{$(\varphi_{1} \to \forall K\,(\varphi_{1})) \to \varphi_{87}$}
\BIC{$\varphi_{87}$}
\end{ptree}
\noindent$\Pi_{53}$:
\begin{ptree}
\AXC{$\stackrel{\Pi_{29}}{\vdots}$}\noLine
\UIC{$\mathrm{DvgP}(a) \to \mathrm{Seq}(b) \to C \in \mathbb{R} \to \varphi_{52}$}
\AXC{$\varphi_{127}$\lbl{A1}}
\BIC{$\mathrm{DvgP}(a) \to \varphi_{6}$}
\AXC{$\forall a\,(\mathrm{DvgP}(a) \leftrightarrow \mathrm{Seq}(a) \land \varphi_{15})$\lbl{Def$\mathrm{DvgP}$}}
\AXC{$\varphi_{130}$\lbl{A2}}
\BIC{$\mathrm{DvgP}(a \oplus b) \leftrightarrow \mathrm{Seq}(a \oplus b) \land \varphi_{27}$}
\AXC{$\varphi_{149}$\lbl{A1}}
\BIC{$(\mathrm{DvgP}(a) \to \varphi_{6}) \to \varphi_{87} \to \mathrm{DvgP}(a) \to \varphi_{8}$}
\BIC{$\varphi_{87} \to \mathrm{DvgP}(a) \to \varphi_{8}$}
\end{ptree}
\noindent 主木:
\begin{ptree}
\AXC{$\stackrel{\Pi_{52}}{\vdots}$}\noLine
\UIC{$\varphi_{87}$}
\AXC{$\stackrel{\Pi_{53}}{\vdots}$}\noLine
\UIC{$\varphi_{87} \to \mathrm{DvgP}(a) \to \varphi_{8}$}
\BIC{$\mathrm{DvgP}(a) \to \varphi_{8}$}
\UIC{$\forall C\,(\mathrm{DvgP}(a) \to \varphi_{8})$}
\UIC{$\forall b\,\forall C\,(\mathrm{DvgP}(a) \to \varphi_{8})$}
\UIC{$\forall a\,\forall b\,\forall C\,(\mathrm{DvgP}(a) \to \varphi_{8})$}
\end{ptree}
\begin{abbrevlist}
\item $\varphi_{1} :\equiv \mathrm{DvgP}(a) \land \mathrm{Seq}(b) \land C \in \mathbb{R} \land \forall n \in \mathbb{N}\,(C \leq b_{n})$
\item $\varphi_{2} :\equiv \forall N\,(\neg (N \in \mathbb{N} \land \forall n \in \mathbb{N}\,(N \leq n \to K < a \oplus b(n))))$
\item $\varphi_{3} :\equiv \forall x\,\forall y\,\forall z\,(x \leq y \to y < z \to x < z) \to \forall y\,\forall z\,(0 \leq y \to y < z \to 0 < z)$
\item $\varphi_{4} :\equiv \forall x\,\forall y\,(x = y \to y = x) \to \forall y\,(a \oplus b(n) = y \to y = a \oplus b(n))$
\item $\varphi_{5} :\equiv \forall N\,(\neg (N \in \mathbb{N} \land \forall n \in \mathbb{N}\,(N \leq n \to \lvert K - C \rvert + 1 < a_{n})))$
\item $\varphi_{6} :\equiv \mathrm{Seq}(b) \to C \in \mathbb{R} \to \forall n \in \mathbb{N}\,(C \leq b_{n}) \to \mathrm{Seq}(a \oplus b)$
\item $\varphi_{7} :\equiv \mathrm{DvgP}(a) \to \mathrm{Seq}(b) \to C \in \mathbb{R} \to \forall n \in \mathbb{N}\,(C \leq b_{n}) \to \mathrm{Seq}(a)$
\item $\varphi_{8} :\equiv \mathrm{Seq}(b) \to C \in \mathbb{R} \to \forall n \in \mathbb{N}\,(C \leq b_{n}) \to \mathrm{DvgP}(a \oplus b)$
\item $\varphi_{9} :\equiv \forall x\,\forall y\,\forall z\,(x \leq y \to y < z \to x < z) \to \forall y\,\forall z\,(K - C \leq y \to y < z \to K - C < z)$
\item $\varphi_{10} :\equiv 0 \leq \lvert K - C \rvert \to \lvert K - C \rvert < \lvert K - C \rvert + 1 \to 0 < \lvert K - C \rvert + 1$
\item $\varphi_{11} :\equiv \mathrm{Seq}(b) \to C \in \mathbb{R} \to \forall n \in \mathbb{N}\,(C \leq b_{n}) \to n \in \mathbb{N} \to C \leq b_{n}$
\item $\varphi_{12} :\equiv \mathrm{Seq}(b) \to C \in \mathbb{R} \to \forall n \in \mathbb{N}\,(C \leq b_{n}) \to K \in \mathbb{R} \to K \in \mathbb{R}$
\item $\varphi_{13} :\equiv N \in \mathbb{N} \land \forall n \in \mathbb{N}\,(N \leq n \to \lvert K - C \rvert + 1 < a_{n}) \to N \in \mathbb{N}$
\item $\varphi_{14} :\equiv \forall y\,\forall z\,(0 \leq y \to y < z \to 0 < z) \to \forall z\,(0 \leq \lvert K - C \rvert \to \lvert K - C \rvert < z \to 0 < z)$
\item $\varphi_{15} :\equiv \forall K \in \mathbb{R}\,(0 < K \to \exists N \in \mathbb{N}\,\forall n \in \mathbb{N}\,(N \leq n \to K < a_{n}))$
\item $\varphi_{16} :\equiv \mathrm{Seq}(b) \to C \in \mathbb{R} \to \forall n \in \mathbb{N}\,(C \leq b_{n}) \to \forall n \in \mathbb{N}\,(C \leq b_{n})$
\item $\varphi_{17} :\equiv \forall n \in \mathbb{N}\,(C \leq b_{n}) \to K \in \mathbb{R} \to 0 < K \to \lvert K - C \rvert + 1 \in \mathbb{R}$
\item $\varphi_{18} :\equiv C \in \mathbb{R} \to \forall n \in \mathbb{N}\,(C \leq b_{n}) \to K \in \mathbb{R} \to 0 < K \to 0 \leq \lvert K - C \rvert$
\item $\varphi_{19} :\equiv \forall z\,(\lvert K - C \rvert < \lvert K - C \rvert + 1 \to \lvert K - C \rvert + 1 < z \to \lvert K - C \rvert < z)$
\item $\varphi_{20} :\equiv C \in \mathbb{R} \to \forall n \in \mathbb{N}\,(C \leq b_{n}) \to K \in \mathbb{R} \to 0 < K \to K - C \in \mathbb{R}$
\item $\varphi_{21} :\equiv \mathrm{Seq}(b) \to \mathrm{Seq}(a \oplus b) \land \forall n \in \mathbb{N}\,(a \oplus b(n) = a_{n} + b_{n})$
\item $\varphi_{22} :\equiv K \in \mathbb{R} \to 0 < K \to \exists N \in \mathbb{N}\,\forall n \in \mathbb{N}\,(N \leq n \to K < a \oplus b(n))$
\item $\varphi_{23} :\equiv \forall n \in \mathbb{N}\,(C \leq b_{n}) \to K \in \mathbb{R} \to 0 < K \to \lvert K - C \rvert < \lvert K - C \rvert + 1$
\item $\varphi_{24} :\equiv \forall w\,(K - C + C = w \to K - C + C < a_{n} + b_{n} \to w < a_{n} + b_{n})$
\item $\varphi_{25} :\equiv \forall n\,(\mathrm{Seq}(b) \to n \in \mathbb{N} \to b_{n} \in \mathbb{R}) \to \mathrm{Seq}(b) \to n \in \mathbb{N} \to b_{n} \in \mathbb{R}$
\item $\varphi_{26} :\equiv \forall n\,(\mathrm{Seq}(a) \to n \in \mathbb{N} \to a_{n} \in \mathbb{R}) \to \mathrm{Seq}(a) \to n \in \mathbb{N} \to a_{n} \in \mathbb{R}$
\item $\varphi_{27} :\equiv \forall K \in \mathbb{R}\,(0 < K \to \exists N \in \mathbb{N}\,\forall n \in \mathbb{N}\,(N \leq n \to K < a \oplus b(n)))$
\item $\varphi_{28} :\equiv C \in \mathbb{R} \to \forall n \in \mathbb{N}\,(C \leq b_{n}) \to K \in \mathbb{R} \to 0 < K \to \lvert K - C \rvert \in \mathbb{R}$
\item $\varphi_{29} :\equiv \forall x \in \mathbb{R}\,(x < x + 1) \to \lvert K - C \rvert \in \mathbb{R} \to \lvert K - C \rvert < \lvert K - C \rvert + 1$
\item $\varphi_{30} :\equiv \forall y\,(K \in \mathbb{R} \to y \in \mathbb{R} \to K - y + y = K) \to K \in \mathbb{R} \to C \in \mathbb{R} \to K - C + C = K$
\item $\varphi_{31} :\equiv C \in \mathbb{R} \to \forall n \in \mathbb{N}\,(C \leq b_{n}) \to K \in \mathbb{R} \to 0 < K \to 0 < \lvert K - C \rvert + 1$
\item $\varphi_{32} :\equiv N \in \mathbb{N} \land \forall n \in \mathbb{N}\,(N \leq n \to \lvert K - C \rvert + 1 < a_{n}) \to n \in \mathbb{N} \to n \in \mathbb{N}$
\item $\varphi_{33} :\equiv N \in \mathbb{N} \land \forall n \in \mathbb{N}\,(N \leq n \to \lvert K - C \rvert + 1 < a_{n}) \to n \in \mathbb{N} \to N \leq n \to N \leq n$
\item $\varphi_{34} :\equiv \forall x\,\forall y\,\forall z\,(x < y \to y < z \to x < z) \to \forall y\,\forall z\,(\lvert K - C \rvert < y \to y < z \to \lvert K - C \rvert < z)$
\item $\varphi_{35} :\equiv K \in \mathbb{R} \to 0 < K \to \exists N \in \mathbb{N}\,\forall n \in \mathbb{N}\,(N \leq n \to \lvert K - C \rvert + 1 < a_{n})$
\item $\varphi_{36} :\equiv \forall y\,(K \in \mathbb{R} \to y \in \mathbb{R} \to K - y \in \mathbb{R}) \to K \in \mathbb{R} \to C \in \mathbb{R} \to K - C \in \mathbb{R}$
\item $\varphi_{37} :\equiv \lvert K - C \rvert < \lvert K - C \rvert + 1 \to \lvert K - C \rvert + 1 < a_{n} \to \lvert K - C \rvert < a_{n}$
\item $\varphi_{38} :\equiv \forall a\,\forall n\,(\mathrm{Seq}(a) \to n \in \mathbb{N} \to a_{n} \in \mathbb{R}) \to \forall n\,(\mathrm{Seq}(b) \to n \in \mathbb{N} \to b_{n} \in \mathbb{R})$
\item $\varphi_{39} :\equiv \forall a\,\forall n\,(\mathrm{Seq}(a) \to n \in \mathbb{N} \to a_{n} \in \mathbb{R}) \to \forall n\,(\mathrm{Seq}(a) \to n \in \mathbb{N} \to a_{n} \in \mathbb{R})$
\item $\varphi_{40} :\equiv 0 < \lvert K - C \rvert + 1 \to \exists N \in \mathbb{N}\,\forall n \in \mathbb{N}\,(N \leq n \to \lvert K - C \rvert + 1 < a_{n})$
\item $\varphi_{41} :\equiv \forall x\,\forall y\,(x \in \mathbb{R} \to y \in \mathbb{R} \to x - y + y = x) \to \forall y\,(K \in \mathbb{R} \to y \in \mathbb{R} \to K - y + y = K)$
\item $\varphi_{42} :\equiv \forall z\,\forall w\,(z = w \to K < z \to K < w) \to \forall w\,(a_{n} + b_{n} = w \to K < a_{n} + b_{n} \to K < w)$
\item $\varphi_{43} :\equiv N \in \mathbb{N} \land \forall n \in \mathbb{N}\,(N \leq n \to \lvert K - C \rvert + 1 < a_{n}) \to n \in \mathbb{N} \to N \leq n \to C \leq b_{n}$
\item $\varphi_{44} :\equiv \forall x\,\forall y\,(x \in \mathbb{R} \to y \in \mathbb{R} \to x - y \in \mathbb{R}) \to \forall y\,(K \in \mathbb{R} \to y \in \mathbb{R} \to K - y \in \mathbb{R})$
\item $\varphi_{45} :\equiv \forall y\,\forall z\,(K - C \leq y \to y < z \to K - C < z) \to \forall z\,(K - C \leq \lvert K - C \rvert \to \lvert K - C \rvert < z \to K - C < z)$
\item $\varphi_{46} :\equiv \forall x\,\forall y\,\forall u\,\forall v\,(x \in \mathbb{R} \to y \in \mathbb{R} \to u \in \mathbb{R} \to v \in \mathbb{R} \to x < y \to u \leq v \to x + u < y + v)$
\item $\varphi_{47} :\equiv N \in \mathbb{N} \land \forall n \in \mathbb{N}\,(N \leq n \to \lvert K - C \rvert + 1 < a_{n}) \to n \in \mathbb{N} \to N \leq n \to K - C + C = K$
\item $\varphi_{48} :\equiv N \in \mathbb{N} \land \forall n \in \mathbb{N}\,(N \leq n \to \lvert K - C \rvert + 1 < a_{n}) \to n \in \mathbb{N} \to N \leq n \to K - C \in \mathbb{R}$
\item $\varphi_{49} :\equiv \forall n \in \mathbb{N}\,(N \leq n \to \lvert K - C \rvert + 1 < a_{n}) \to n \in \mathbb{N} \to N \leq n \to \lvert K - C \rvert + 1 < a_{n}$
\item $\varphi_{50} :\equiv N \in \mathbb{N} \land \forall n \in \mathbb{N}\,(N \leq n \to \lvert K - C \rvert + 1 < a_{n}) \to n \in \mathbb{N} \to N \leq n \to b_{n} \in \mathbb{R}$
\item $\varphi_{51} :\equiv N \in \mathbb{N} \land \forall n \in \mathbb{N}\,(N \leq n \to \lvert K - C \rvert + 1 < a_{n}) \to n \in \mathbb{N} \to N \leq n \to a_{n} \in \mathbb{R}$
\item $\varphi_{52} :\equiv \forall n \in \mathbb{N}\,(C \leq b_{n}) \to \mathrm{Seq}(a \oplus b) \land \forall n \in \mathbb{N}\,(a \oplus b(n) = a_{n} + b_{n})$
\item $\varphi_{53} :\equiv N \in \mathbb{N} \land \forall n \in \mathbb{N}\,(N \leq n \to \lvert K - C \rvert + 1 < a_{n}) \to n \in \mathbb{N} \to N \leq n \to K - C \leq \lvert K - C \rvert$
\item $\varphi_{54} :\equiv N \in \mathbb{N} \land \forall n \in \mathbb{N}\,(N \leq n \to \lvert K - C \rvert + 1 < a_{n}) \to n \in \mathbb{N} \to N \leq n \to K - C < a_{n}$
\item $\varphi_{55} :\equiv N \in \mathbb{N} \land \forall n \in \mathbb{N}\,(N \leq n \to \lvert K - C \rvert + 1 < a_{n}) \to n \in \mathbb{N} \to N \leq n \to \lvert K - C \rvert \in \mathbb{R}$
\item $\varphi_{56} :\equiv C \in \mathbb{R} \to b_{n} \in \mathbb{R} \to K - C < a_{n} \to C \leq b_{n} \to K - C + C < a_{n} + b_{n}$
\item $\varphi_{57} :\equiv N \in \mathbb{N} \land \forall n \in \mathbb{N}\,(N \leq n \to \lvert K - C \rvert + 1 < a_{n}) \to n \in \mathbb{N} \to N \leq n \to K < a \oplus b(n)$
\item $\varphi_{58} :\equiv \forall x \in \mathbb{R}\,(-\lvert x \rvert \leq x \land x \leq \lvert x \rvert) \to K - C \in \mathbb{R} \to -\lvert K - C \rvert \leq K - C \land K - C \leq \lvert K - C \rvert$
\item $\varphi_{59} :\equiv N \in \mathbb{N} \land \forall n \in \mathbb{N}\,(N \leq n \to \lvert K - C \rvert + 1 < a_{n}) \to n \in \mathbb{N} \to N \leq n \to \lvert K - C \rvert < a_{n}$
\item $\varphi_{60} :\equiv \forall n \in \mathbb{N}\,(a \oplus b(n) = a_{n} + b_{n}) \to n \in \mathbb{N} \to a \oplus b(n) = a_{n} + b_{n}$
\item $\varphi_{61} :\equiv \forall y\,\forall u\,\forall v\,(K - C \in \mathbb{R} \to y \in \mathbb{R} \to u \in \mathbb{R} \to v \in \mathbb{R} \to K - C < y \to u \leq v \to K - C + u < y + v)$
\item $\varphi_{62} :\equiv N \in \mathbb{N} \land \forall n \in \mathbb{N}\,(N \leq n \to \lvert K - C \rvert + 1 < a_{n}) \to n \in \mathbb{N} \to N \leq n \to K < a_{n} + b_{n}$
\item $\varphi_{63} :\equiv N \in \mathbb{N} \land \forall n \in \mathbb{N}\,(N \leq n \to \lvert K - C \rvert + 1 < a_{n}) \to \forall n \in \mathbb{N}\,(N \leq n \to K < a \oplus b(n))$
\item $\varphi_{64} :\equiv N \in \mathbb{N} \land \forall n \in \mathbb{N}\,(N \leq n \to \lvert K - C \rvert + 1 < a_{n}) \to n \in \mathbb{N} \to N \leq n \to \lvert K - C \rvert + 1 < a_{n}$
\item $\varphi_{65} :\equiv N \in \mathbb{N} \land \forall n \in \mathbb{N}\,(N \leq n \to \lvert K - C \rvert + 1 < a_{n}) \to n \in \mathbb{N} \to N \leq n \to \lvert K - C \rvert < \lvert K - C \rvert + 1$
\item $\varphi_{66} :\equiv \forall x\,\forall y\,(x \in \mathbb{R} \to y \in \mathbb{R} \to x + y \in \mathbb{R}) \to \forall y\,(\lvert K - C \rvert \in \mathbb{R} \to y \in \mathbb{R} \to \lvert K - C \rvert + y \in \mathbb{R})$
\item $\varphi_{67} :\equiv N \in \mathbb{N} \land \forall n \in \mathbb{N}\,(N \leq n \to \lvert K - C \rvert + 1 < a_{n}) \to n \in \mathbb{N} \to N \leq n \to K - C + C < a_{n} + b_{n}$
\item $\varphi_{68} :\equiv N \in \mathbb{N} \land \forall n \in \mathbb{N}\,(N \leq n \to \lvert K - C \rvert + 1 < a_{n}) \to N \in \mathbb{N} \land \forall n \in \mathbb{N}\,(N \leq n \to K < a \oplus b(n))$
\item $\varphi_{69} :\equiv N \in \mathbb{N} \land \forall n \in \mathbb{N}\,(N \leq n \to \lvert K - C \rvert + 1 < a_{n}) \to n \in \mathbb{N} \to N \leq n \to a \oplus b(n) = a_{n} + b_{n}$
\item $\varphi_{70} :\equiv N \in \mathbb{N} \land \forall n \in \mathbb{N}\,(N \leq n \to \lvert K - C \rvert + 1 < a_{n}) \to n \in \mathbb{N} \to N \leq n \to a_{n} + b_{n} = a \oplus b(n)$
\item $\varphi_{71} :\equiv \forall z\,(K - C \leq \lvert K - C \rvert \to \lvert K - C \rvert < z \to K - C < z) \to K - C \leq \lvert K - C \rvert \to \lvert K - C \rvert < a_{n} \to K - C < a_{n}$
\item $\varphi_{72} :\equiv \forall v\,(K - C \in \mathbb{R} \to a_{n} \in \mathbb{R} \to C \in \mathbb{R} \to v \in \mathbb{R} \to K - C < a_{n} \to C \leq v \to K - C + C < a_{n} + v)$
\item $\varphi_{73} :\equiv N \in \mathbb{N} \land \forall n \in \mathbb{N}\,(N \leq n \to \lvert K - C \rvert + 1 < a_{n}) \to \exists N \in \mathbb{N}\,\forall n \in \mathbb{N}\,(N \leq n \to K < a \oplus b(n))$
\item $\varphi_{74} :\equiv \forall u\,\forall v\,(K - C \in \mathbb{R} \to a_{n} \in \mathbb{R} \to u \in \mathbb{R} \to v \in \mathbb{R} \to K - C < a_{n} \to u \leq v \to K - C + u < a_{n} + v)$
\item $\varphi_{75} :\equiv N \in \mathbb{N} \land \forall n \in \mathbb{N}\,(N \leq n \to \lvert K - C \rvert + 1 < a_{n}) \to N \in \mathbb{N} \land \forall n \in \mathbb{N}\,(N \leq n \to \lvert K - C \rvert + 1 < a_{n})$
\item $\varphi_{76} :\equiv \forall y\,(a \oplus b(n) = y \to y = a \oplus b(n)) \to a \oplus b(n) = a_{n} + b_{n} \to a_{n} + b_{n} = a \oplus b(n)$
\item $\varphi_{77} :\equiv N \in \mathbb{N} \land \forall n \in \mathbb{N}\,(N \leq n \to \lvert K - C \rvert + 1 < a_{n}) \to n \in \mathbb{N} \to N \leq n \to -\lvert K - C \rvert \leq K - C \land K - C \leq \lvert K - C \rvert$
\item $\varphi_{78} :\equiv \forall y\,(\lvert K - C \rvert \in \mathbb{R} \to y \in \mathbb{R} \to \lvert K - C \rvert + y \in \mathbb{R}) \to \lvert K - C \rvert \in \mathbb{R} \to 1 \in \mathbb{R} \to \lvert K - C \rvert + 1 \in \mathbb{R}$
\item $\varphi_{79} :\equiv N \in \mathbb{N} \land \forall n \in \mathbb{N}\,(N \leq n \to \lvert K - C \rvert + 1 < a_{n}) \to n \in \mathbb{N} \to N \leq n \to n \in \mathbb{N} \to N \leq n \to \lvert K - C \rvert + 1 < a_{n}$
\item $\varphi_{80} :\equiv N \in \mathbb{N} \land \forall n \in \mathbb{N}\,(N \leq n \to \lvert K - C \rvert + 1 < a_{n}) \to n \in \mathbb{N} \to N \leq n \to \forall n \in \mathbb{N}\,(N \leq n \to \lvert K - C \rvert + 1 < a_{n})$
\item $\varphi_{81} :\equiv N \in \mathbb{N} \land \forall n \in \mathbb{N}\,(N \leq n \to \lvert K - C \rvert + 1 < a_{n}) \to n \in \mathbb{N} \to N \leq n \to n \in \mathbb{N} \to a \oplus b(n) = a_{n} + b_{n}$
\item $\varphi_{82} :\equiv (\forall n \in \mathbb{N}\,(C \leq b_{n}) \to n \in \mathbb{N} \to C \leq b_{n}) \to \mathrm{DvgP}(a) \to \varphi_{11}$
\item $\varphi_{83} :\equiv N \in \mathbb{N} \land \forall n \in \mathbb{N}\,(N \leq n \to \lvert K - C \rvert + 1 < a_{n}) \to n \in \mathbb{N} \to N \leq n \to \forall n \in \mathbb{N}\,(a \oplus b(n) = a_{n} + b_{n})$
\item $\varphi_{84} :\equiv \mathrm{DvgP}(a) \to \mathrm{Seq}(b) \to C \in \mathbb{R} \to \forall n \in \mathbb{N}\,(C \leq b_{n}) \to \varphi_{22}$
\item $\varphi_{85} :\equiv \varphi_{1} \land K \in \mathbb{R} \land 0 < K \land \neg \exists N \in \mathbb{N}\,\forall n \in \mathbb{N}\,(N \leq n \to K < a \oplus b(n))$
\item $\varphi_{86} :\equiv (\mathrm{DvgP}(a) \to \mathrm{Seq}(b) \to C \in \mathbb{R} \to C \in \mathbb{R}) \to \mathrm{DvgP}(a) \to \mathrm{Seq}(b) \to \varphi_{20}$
\item $\varphi_{87} :\equiv \mathrm{DvgP}(a) \to \mathrm{Seq}(b) \to C \in \mathbb{R} \to \forall n \in \mathbb{N}\,(C \leq b_{n}) \to \varphi_{27}$
\item $\varphi_{88} :\equiv \varphi_{1} \land K \in \mathbb{R} \land 0 < K \land (N \in \mathbb{N} \land \forall n \in \mathbb{N}\,(N \leq n \to \lvert K - C \rvert + 1 < a_{n}))$
\item $\varphi_{89} :\equiv \mathrm{Seq}(b) \to C \in \mathbb{R} \to \forall n \in \mathbb{N}\,(C \leq b_{n}) \to K \in \mathbb{R} \to 0 < K \to \varphi_{13}$
\item $\varphi_{90} :\equiv \mathrm{DvgP}(a) \to \mathrm{Seq}(b) \to C \in \mathbb{R} \to \forall n \in \mathbb{N}\,(C \leq b_{n}) \to \varphi_{35}$
\item $\varphi_{91} :\equiv \mathrm{Seq}(b) \to C \in \mathbb{R} \to \forall n \in \mathbb{N}\,(C \leq b_{n}) \to K \in \mathbb{R} \to 0 < K \to \varphi_{15}$
\item $\varphi_{92} :\equiv \forall w\,(a_{n} + b_{n} = w \to K < a_{n} + b_{n} \to K < w) \to a_{n} + b_{n} = a \oplus b(n) \to K < a_{n} + b_{n} \to K < a \oplus b(n)$
\item $\varphi_{93} :\equiv \mathrm{Seq}(b) \to C \in \mathbb{R} \to \forall n \in \mathbb{N}\,(C \leq b_{n}) \to K \in \mathbb{R} \to 0 < K \to \varphi_{32}$
\item $\varphi_{94} :\equiv \mathrm{Seq}(b) \to C \in \mathbb{R} \to \forall n \in \mathbb{N}\,(C \leq b_{n}) \to K \in \mathbb{R} \to 0 < K \to \varphi_{33}$
\item $\varphi_{95} :\equiv (\mathrm{DvgP}(a) \to \mathrm{Seq}(b) \to \mathrm{Seq}(b)) \to \mathrm{DvgP}(a) \to \mathrm{Seq}(b) \to C \in \mathbb{R} \to \varphi_{52}$
\item $\varphi_{96} :\equiv \mathrm{Seq}(b) \to C \in \mathbb{R} \to \forall n \in \mathbb{N}\,(C \leq b_{n}) \to K \in \mathbb{R} \to 0 < K \to \varphi_{43}$
\item $\varphi_{97} :\equiv \forall n \in \mathbb{N}\,(C \leq b_{n}) \to K \in \mathbb{R} \to 0 < K \to \lvert K - C \rvert + 1 \in \mathbb{R} \to \varphi_{40}$
\item $\varphi_{98} :\equiv \mathrm{Seq}(b) \to C \in \mathbb{R} \to \forall n \in \mathbb{N}\,(C \leq b_{n}) \to K \in \mathbb{R} \to 0 < K \to \varphi_{47}$
\item $\varphi_{99} :\equiv \mathrm{Seq}(b) \to C \in \mathbb{R} \to \forall n \in \mathbb{N}\,(C \leq b_{n}) \to K \in \mathbb{R} \to 0 < K \to \varphi_{48}$
\item $\varphi_{100} :\equiv \mathrm{Seq}(b) \to C \in \mathbb{R} \to \forall n \in \mathbb{N}\,(C \leq b_{n}) \to K \in \mathbb{R} \to 0 < K \to \varphi_{50}$
\item $\varphi_{101} :\equiv \mathrm{Seq}(b) \to C \in \mathbb{R} \to \forall n \in \mathbb{N}\,(C \leq b_{n}) \to K \in \mathbb{R} \to 0 < K \to \varphi_{51}$
\item $\varphi_{102} :\equiv \mathrm{Seq}(b) \to C \in \mathbb{R} \to \forall n \in \mathbb{N}\,(C \leq b_{n}) \to K \in \mathbb{R} \to 0 < K \to \varphi_{53}$
\item $\varphi_{103} :\equiv \mathrm{Seq}(b) \to C \in \mathbb{R} \to \forall n \in \mathbb{N}\,(C \leq b_{n}) \to K \in \mathbb{R} \to 0 < K \to \varphi_{54}$
\item $\varphi_{104} :\equiv \mathrm{Seq}(b) \to C \in \mathbb{R} \to \forall n \in \mathbb{N}\,(C \leq b_{n}) \to K \in \mathbb{R} \to 0 < K \to \varphi_{55}$
\item $\varphi_{105} :\equiv \mathrm{Seq}(b) \to C \in \mathbb{R} \to \forall n \in \mathbb{N}\,(C \leq b_{n}) \to K \in \mathbb{R} \to 0 < K \to \varphi_{57}$
\item $\varphi_{106} :\equiv \mathrm{Seq}(b) \to C \in \mathbb{R} \to \forall n \in \mathbb{N}\,(C \leq b_{n}) \to K \in \mathbb{R} \to 0 < K \to \varphi_{59}$
\item $\varphi_{107} :\equiv \mathrm{Seq}(b) \to C \in \mathbb{R} \to \forall n \in \mathbb{N}\,(C \leq b_{n}) \to K \in \mathbb{R} \to 0 < K \to \varphi_{62}$
\item $\varphi_{108} :\equiv \mathrm{Seq}(b) \to C \in \mathbb{R} \to \forall n \in \mathbb{N}\,(C \leq b_{n}) \to K \in \mathbb{R} \to 0 < K \to \varphi_{63}$
\item $\varphi_{109} :\equiv (\mathrm{DvgP}(a) \to \mathrm{DvgP}(a)) \to (\mathrm{DvgP}(a) \leftrightarrow \mathrm{Seq}(a) \land \varphi_{15}) \to \varphi_{7}$
\item $\varphi_{110} :\equiv \mathrm{Seq}(b) \to C \in \mathbb{R} \to \forall n \in \mathbb{N}\,(C \leq b_{n}) \to K \in \mathbb{R} \to 0 < K \to \varphi_{64}$
\item $\varphi_{111} :\equiv \mathrm{Seq}(b) \to C \in \mathbb{R} \to \forall n \in \mathbb{N}\,(C \leq b_{n}) \to K \in \mathbb{R} \to 0 < K \to \varphi_{65}$
\item $\varphi_{112} :\equiv \mathrm{Seq}(b) \to C \in \mathbb{R} \to \forall n \in \mathbb{N}\,(C \leq b_{n}) \to K \in \mathbb{R} \to 0 < K \to \varphi_{67}$
\item $\varphi_{113} :\equiv \mathrm{Seq}(b) \to C \in \mathbb{R} \to \forall n \in \mathbb{N}\,(C \leq b_{n}) \to K \in \mathbb{R} \to 0 < K \to \varphi_{68}$
\item $\varphi_{114} :\equiv \forall a\,(\mathrm{DvgP}(a) \leftrightarrow \mathrm{Seq}(a) \land \varphi_{15}) \to (\mathrm{DvgP}(a) \leftrightarrow \mathrm{Seq}(a) \land \varphi_{15})$
\item $\varphi_{115} :\equiv \mathrm{Seq}(b) \to C \in \mathbb{R} \to \forall n \in \mathbb{N}\,(C \leq b_{n}) \to K \in \mathbb{R} \to 0 < K \to \varphi_{69}$
\item $\varphi_{116} :\equiv \mathrm{Seq}(b) \to C \in \mathbb{R} \to \forall n \in \mathbb{N}\,(C \leq b_{n}) \to K \in \mathbb{R} \to 0 < K \to \varphi_{70}$
\item $\varphi_{117} :\equiv \mathrm{Seq}(b) \to C \in \mathbb{R} \to \forall n \in \mathbb{N}\,(C \leq b_{n}) \to K \in \mathbb{R} \to 0 < K \to \varphi_{73}$
\item $\varphi_{118} :\equiv (\mathrm{DvgP}(a) \to \mathrm{Seq}(b) \to \varphi_{20}) \to \mathrm{DvgP}(a) \to \mathrm{Seq}(b) \to \varphi_{18}$
\item $\varphi_{119} :\equiv \mathrm{Seq}(b) \to C \in \mathbb{R} \to \forall n \in \mathbb{N}\,(C \leq b_{n}) \to K \in \mathbb{R} \to 0 < K \to \varphi_{75}$
\item $\varphi_{120} :\equiv \mathrm{Seq}(b) \to C \in \mathbb{R} \to \forall n \in \mathbb{N}\,(C \leq b_{n}) \to K \in \mathbb{R} \to 0 < K \to \varphi_{77}$
\item $\varphi_{121} :\equiv \mathrm{Seq}(b) \to C \in \mathbb{R} \to \forall n \in \mathbb{N}\,(C \leq b_{n}) \to K \in \mathbb{R} \to 0 < K \to \varphi_{79}$
\item $\varphi_{122} :\equiv (\mathrm{DvgP}(a) \to \mathrm{Seq}(b) \to \varphi_{20}) \to \mathrm{DvgP}(a) \to \mathrm{Seq}(b) \to \varphi_{28}$
\item $\varphi_{123} :\equiv \mathrm{Seq}(b) \to C \in \mathbb{R} \to \forall n \in \mathbb{N}\,(C \leq b_{n}) \to K \in \mathbb{R} \to 0 < K \to \varphi_{80}$
\item $\varphi_{124} :\equiv \mathrm{Seq}(b) \to C \in \mathbb{R} \to \forall n \in \mathbb{N}\,(C \leq b_{n}) \to K \in \mathbb{R} \to 0 < K \to \varphi_{81}$
\item $\varphi_{125} :\equiv \varphi_{85} \to \neg (N \in \mathbb{N} \land \forall n \in \mathbb{N}\,(N \leq n \to \lvert K - C \rvert + 1 < a_{n}))$
\item $\varphi_{126} :\equiv \mathrm{Seq}(b) \to C \in \mathbb{R} \to \forall n \in \mathbb{N}\,(C \leq b_{n}) \to K \in \mathbb{R} \to 0 < K \to \varphi_{83}$
\item $\varphi_{127} :\equiv (\mathrm{DvgP}(a) \to \mathrm{Seq}(b) \to C \in \mathbb{R} \to \varphi_{52}) \to \mathrm{DvgP}(a) \to \varphi_{6}$
\item $\varphi_{128} :\equiv (\mathrm{DvgP}(a) \to \mathrm{Seq}(b) \to \varphi_{28}) \to \mathrm{DvgP}(a) \to \mathrm{Seq}(b) \to C \in \mathbb{R} \to \varphi_{23}$
\item $\varphi_{129} :\equiv (\mathrm{DvgP}(a) \to \mathrm{Seq}(b) \to C \in \mathbb{R} \to \varphi_{23}) \to \mathrm{DvgP}(a) \to \mathrm{Seq}(b) \to \varphi_{31}$
\item $\varphi_{130} :\equiv \forall a\,(\mathrm{DvgP}(a) \leftrightarrow \mathrm{Seq}(a) \land \varphi_{15}) \to (\mathrm{DvgP}(a \oplus b) \leftrightarrow \mathrm{Seq}(a \oplus b) \land \varphi_{27})$
\item $\varphi_{131} :\equiv (\mathrm{DvgP}(a) \to \mathrm{Seq}(b) \to C \in \mathbb{R} \to C \in \mathbb{R}) \to \mathrm{DvgP}(a) \to \varphi_{98}$
\item $\varphi_{132} :\equiv (\mathrm{DvgP}(a) \to \mathrm{Seq}(b) \to C \in \mathbb{R} \to C \in \mathbb{R}) \to \mathrm{DvgP}(a) \to \varphi_{99}$
\item $\varphi_{133} :\equiv (\mathrm{DvgP}(a) \to \mathrm{Seq}(b) \to \varphi_{28}) \to 1 \in \mathbb{R} \to \mathrm{DvgP}(a) \to \mathrm{Seq}(b) \to C \in \mathbb{R} \to \varphi_{17}$
\item $\varphi_{134} :\equiv (K - C + C = K \to K - C + C < a_{n} + b_{n} \to K < a_{n} + b_{n}) \to \mathrm{DvgP}(a) \to \varphi_{107}$
\item $\varphi_{135} :\equiv (K \in \mathbb{R} \to C \in \mathbb{R} \to K - C \in \mathbb{R}) \to (\mathrm{DvgP}(a) \to \varphi_{12}) \to \varphi_{86}$
\item $\varphi_{136} :\equiv (\mathrm{DvgP}(a) \to \mathrm{DvgP}(a)) \to (\mathrm{DvgP}(a) \leftrightarrow \mathrm{Seq}(a) \land \varphi_{15}) \to \mathrm{DvgP}(a) \to \varphi_{91}$
\item $\varphi_{137} :\equiv (a_{n} + b_{n} = a \oplus b(n) \to K < a_{n} + b_{n} \to K < a \oplus b(n)) \to \mathrm{DvgP}(a) \to \varphi_{105}$
\item $\varphi_{138} :\equiv (a \oplus b(n) = a_{n} + b_{n} \to a_{n} + b_{n} = a \oplus b(n)) \to \mathrm{DvgP}(a) \to \varphi_{116}$
\item $\varphi_{139} :\equiv (\varphi_{2} \to \neg (N \in \mathbb{N} \land \forall n \in \mathbb{N}\,(N \leq n \to K < a \oplus b(n)))) \to \mathrm{DvgP}(a) \to \varphi_{117}$
\item $\varphi_{140} :\equiv (K \in \mathbb{R} \to C \in \mathbb{R} \to K - C + C = K) \to (\mathrm{DvgP}(a) \to \varphi_{12}) \to \varphi_{131}$
\item $\varphi_{141} :\equiv (\mathrm{DvgP}(a) \to \mathrm{Seq}(b) \to C \in \mathbb{R} \to \varphi_{52}) \to \mathrm{DvgP}(a) \to \varphi_{126}$
\item $\varphi_{142} :\equiv (K \in \mathbb{R} \to C \in \mathbb{R} \to K - C \in \mathbb{R}) \to (\mathrm{DvgP}(a) \to \varphi_{12}) \to \varphi_{132}$
\item $\varphi_{143} :\equiv (\mathrm{DvgP}(a) \to \mathrm{Seq}(b) \to C \in \mathbb{R} \to \varphi_{17}) \to (\mathrm{DvgP}(a) \to \mathrm{Seq}(b) \to \varphi_{31}) \to \varphi_{90}$
\item $\varphi_{144} :\equiv (\mathrm{DvgP}(a) \to \varphi_{105}) \to \varphi_{88} \to n \in \mathbb{N} \to N \leq n \to K < a \oplus b(n)$
\item $\varphi_{145} :\equiv (\mathrm{DvgP}(a) \to \mathrm{Seq}(b) \to \mathrm{Seq}(b)) \to (\mathrm{DvgP}(a) \to \varphi_{93}) \to \mathrm{DvgP}(a) \to \varphi_{100}$
\item $\varphi_{146} :\equiv \forall n\,(\varphi_{88} \to n \in \mathbb{N} \to N \leq n \to K < a \oplus b(n)) \to \forall n\,(\varphi_{88}) \to \forall n \in \mathbb{N}\,(N \leq n \to K < a \oplus b(n))$
\item $\varphi_{147} :\equiv (\varphi_{15} \to \lvert K - C \rvert + 1 \in \mathbb{R} \to \varphi_{40}) \to \mathrm{DvgP}(a) \to \mathrm{Seq}(b) \to C \in \mathbb{R} \to \varphi_{97}$
\item $\varphi_{148} :\equiv (K - C \in \mathbb{R} \to \lvert K - C \rvert \in \mathbb{R}) \to (\mathrm{DvgP}(a) \to \varphi_{99}) \to \mathrm{DvgP}(a) \to \varphi_{104}$
\item $\varphi_{149} :\equiv (\mathrm{DvgP}(a \oplus b) \leftrightarrow \mathrm{Seq}(a \oplus b) \land \varphi_{27}) \to (\mathrm{DvgP}(a) \to \varphi_{6}) \to \varphi_{87} \to \mathrm{DvgP}(a) \to \varphi_{8}$
\item $\varphi_{150} :\equiv (\mathrm{DvgP}(a) \to \varphi_{11}) \to (\mathrm{DvgP}(a) \to \varphi_{93}) \to \mathrm{DvgP}(a) \to \varphi_{96}$
\item $\varphi_{151} :\equiv (\lvert K - C \rvert \in \mathbb{R} \to \lvert K - C \rvert < \lvert K - C \rvert + 1) \to (\mathrm{DvgP}(a) \to \varphi_{104}) \to \mathrm{DvgP}(a) \to \varphi_{111}$
\item $\varphi_{152} :\equiv (K - C \in \mathbb{R} \to -\lvert K - C \rvert \leq K - C \land K - C \leq \lvert K - C \rvert) \to (\mathrm{DvgP}(a) \to \varphi_{99}) \to \mathrm{DvgP}(a) \to \varphi_{120}$
\item $\varphi_{153} :\equiv (\mathrm{Seq}(a) \to n \in \mathbb{N} \to a_{n} \in \mathbb{R}) \to \varphi_{7} \to (\mathrm{DvgP}(a) \to \varphi_{93}) \to \mathrm{DvgP}(a) \to \varphi_{101}$
\item $\varphi_{154} :\equiv (\mathrm{DvgP}(a) \to \varphi_{93}) \to (\mathrm{DvgP}(a) \to \varphi_{94}) \to \mathrm{DvgP}(a) \to \varphi_{110}$
\item $\varphi_{155} :\equiv (\mathrm{DvgP}(a) \to \varphi_{89}) \to (\mathrm{DvgP}(a) \to \varphi_{108}) \to \mathrm{DvgP}(a) \to \varphi_{113}$
\item $\varphi_{156} :\equiv (\mathrm{DvgP}(a) \to \varphi_{102}) \to (\mathrm{DvgP}(a) \to \varphi_{106}) \to \mathrm{DvgP}(a) \to \varphi_{103}$
\item $\varphi_{157} :\equiv (\mathrm{DvgP}(a) \to \varphi_{103}) \to (\mathrm{DvgP}(a) \to \varphi_{96}) \to \mathrm{DvgP}(a) \to \varphi_{112}$
\item $\varphi_{158} :\equiv (\mathrm{DvgP}(a) \to \varphi_{111}) \to (\mathrm{DvgP}(a) \to \varphi_{110}) \to \mathrm{DvgP}(a) \to \varphi_{106}$
\item $\varphi_{159} :\equiv (\mathrm{DvgP}(a) \to \varphi_{124}) \to (\mathrm{DvgP}(a) \to \varphi_{93}) \to \mathrm{DvgP}(a) \to \varphi_{115}$
\item $\varphi_{160} :\equiv \forall K\,(\varphi_{1} \to \varphi_{22}) \to (\forall K\,(\varphi_{1} \to \varphi_{22}) \to \forall K\,(\varphi_{1}) \to \varphi_{27}) \to (\varphi_{1} \to \forall K\,(\varphi_{1})) \to \varphi_{87}$
\item $\varphi_{161} :\equiv (\mathrm{DvgP}(a) \to \mathrm{Seq}(b) \to C \in \mathbb{R} \to C \in \mathbb{R}) \to (\mathrm{DvgP}(a) \to \varphi_{100}) \to \varphi_{157}$
\item $\varphi_{162} :\equiv \forall n\,(\varphi_{88} \to n \in \mathbb{N} \to N \leq n \to K < a \oplus b(n)) \to \varphi_{146} \to (\varphi_{88} \to \forall n\,(\varphi_{88})) \to \mathrm{DvgP}(a) \to \varphi_{108}$
\item $\varphi_{163} :\equiv \forall N\,(\varphi_{125}) \to (\forall N\,(\varphi_{125}) \to \forall N\,(\varphi_{85}) \to \varphi_{5}) \to (\varphi_{85} \to \forall N\,(\varphi_{85})) \to \varphi_{84}$
\item $\varphi_{164} :\equiv (K - C \in \mathbb{R} \to a_{n} \in \mathbb{R} \to \varphi_{56}) \to (\mathrm{DvgP}(a) \to \varphi_{99}) \to (\mathrm{DvgP}(a) \to \varphi_{101}) \to \varphi_{161}$
\end{abbrevlist}
\end{fullproof}

% 自動生成 (prover/build.js): DvgMul  ノード数 295
\begin{fullproof}{定理2.18(3)}{$\forall a\,\forall b\,\forall C\,(\mathrm{DvgP}(a) \to \mathrm{Seq}(b) \to C \in \mathbb{R} \to 0 < C \to \forall n \in \mathbb{N}\,(C \leq b_{n}) \to \mathrm{DvgP}(a \otimes b))$}{295}\phantomsection\label{tree:DvgMul}
\noindent$\Pi_{1}$:
\begin{ptree}
\AXC{$\forall n \in \mathbb{N}\,(C \leq b_{n}) \to n \in \mathbb{N} \to C \leq b_{n}$\lbl{A2}}
\AXC{$\mathrm{DvgP}(a) \to \mathrm{Seq}(b) \to \varphi_{15}$\lbl{A1}}
\AXC{$(\mathrm{DvgP}(a) \to \mathrm{Seq}(b) \to \varphi_{15}) \to \varphi_{84}$\lbl{A1}}
\BIC{$\varphi_{84}$}
\BIC{$\mathrm{DvgP}(a) \to \varphi_{29}$}
\end{ptree}
\noindent$\Pi_{2}$:
\begin{ptree}
\AXC{$\mathrm{DvgP}(a) \to \mathrm{Seq}(b) \to \varphi_{93}$\lbl{A1}}
\AXC{$\varphi_{41}$\lbl{A2}}
\AXC{$\mathrm{DvgP}(a) \to \mathrm{Seq}(b) \to \varphi_{116}$\lbl{A1}}
\AXC{$\varphi_{159}$\lbl{A1}}
\BIC{$\mathrm{DvgP}(a) \to \mathrm{Seq}(b) \to \varphi_{124}$}
\AXC{$\varphi_{167}$\lbl{A1}}
\BIC{$\varphi_{41} \to \mathrm{DvgP}(a) \to \mathrm{Seq}(b) \to \varphi_{120}$}
\BIC{$\mathrm{DvgP}(a) \to \mathrm{Seq}(b) \to \varphi_{120}$}
\AXC{$\varphi_{176}$\lbl{A1}}
\BIC{$(\mathrm{DvgP}(a) \to \mathrm{Seq}(b) \to \varphi_{93}) \to \varphi_{145}$}
\BIC{$\varphi_{145}$}
\end{ptree}
\noindent$\Pi_{3}$:
\begin{ptree}
\AXC{$\mathrm{DvgP}(a) \to \mathrm{Seq}(b) \to \varphi_{94}$\lbl{A1}}
\AXC{$\stackrel{\Pi_{2}}{\vdots}$}\noLine
\UIC{$\varphi_{145}$}
\BIC{$\mathrm{DvgP}(a) \to \mathrm{Seq}(b) \to \varphi_{109}$}
\end{ptree}
\noindent$\Pi_{4}$:
\begin{ptree}
\AXC{$\mathrm{DvgP}(a) \to \mathrm{Seq}(b) \to C \in \mathbb{R} \to C \in \mathbb{R}$\lbl{A1}}
\AXC{$\forall x \in \mathbb{R}\,(0 < x \to 0 < x^{-1})$\lbl{[F34]}}
\AXC{$\varphi_{11}$\lbl{A2}}
\BIC{$C \in \mathbb{R} \to 0 < C \to 0 < C^{-1}$}
\AXC{$\varphi_{132}$\lbl{A1}}
\BIC{$(\mathrm{DvgP}(a) \to \mathrm{Seq}(b) \to C \in \mathbb{R} \to C \in \mathbb{R}) \to \varphi_{92}$}
\BIC{$\varphi_{92}$}
\end{ptree}
\noindent$\Pi_{5}$:
\begin{ptree}
\AXC{$x = x$\lbl{A5}}
\AXC{$x = y \to x = x \to y = x$\lbl{A6}}
\AXC{$(x = y \to x = x \to y = x) \to x = x \to x = y \to y = x$\lbl{A1}}
\BIC{$x = x \to x = y \to y = x$}
\BIC{$x = y \to y = x$}
\UIC{$\forall y\,(x = y \to y = x)$}
\UIC{$\forall x\,\forall y\,(x = y \to y = x)$}
\end{ptree}
\noindent$\Pi_{6}$:
\begin{ptree}
\AXC{$\stackrel{\Pi_{5}}{\vdots}$}\noLine
\UIC{$\forall x\,\forall y\,(x = y \to y = x)$}
\AXC{$\forall x\,\forall y\,(x = y \to y = x) \to \forall y\,(C = y \to y = C)$\lbl{A2}}
\BIC{$\forall y\,(C = y \to y = C)$}
\AXC{$\forall y\,(C = y \to y = C) \to C = 0 \to 0 = C$\lbl{A2}}
\BIC{$C = 0 \to 0 = C$}
\end{ptree}
\noindent$\Pi_{7}$:
\begin{ptree}
\AXC{$\forall x\,\forall y\,(x < y \to x \neq y)$\lbl{[F26]}}
\AXC{$\forall x\,\forall y\,(x < y \to x \neq y) \to \forall y\,(0 < y \to 0 \neq y)$\lbl{A2}}
\BIC{$\forall y\,(0 < y \to 0 \neq y)$}
\AXC{$\forall y\,(0 < y \to 0 \neq y) \to 0 < C \to 0 \neq C$\lbl{A2}}
\BIC{$0 < C \to 0 \neq C$}
\end{ptree}
\noindent$\Pi_{8}$:
\begin{ptree}
\AXC{$\mathrm{DvgP}(a) \to \mathrm{Seq}(b) \to C \in \mathbb{R} \to 0 < C \to 0 < C$\lbl{A1}}
\AXC{$\stackrel{\Pi_{7}}{\vdots}$}\noLine
\UIC{$0 < C \to 0 \neq C$}
\AXC{$(0 < C \to 0 \neq C) \to \varphi_{77}$\lbl{A1}}
\BIC{$\varphi_{77}$}
\BIC{$\mathrm{DvgP}(a) \to \varphi_{4}$}
\AXC{$\varphi_{100}$\lbl{A1}}
\BIC{$(C = 0 \to 0 = C) \to \mathrm{DvgP}(a) \to \varphi_{5}$}
\end{ptree}
\noindent$\Pi_{9}$:
\begin{ptree}
\AXC{$\stackrel{\Pi_{6}}{\vdots}$}\noLine
\UIC{$C = 0 \to 0 = C$}
\AXC{$\stackrel{\Pi_{8}}{\vdots}$}\noLine
\UIC{$(C = 0 \to 0 = C) \to \mathrm{DvgP}(a) \to \varphi_{5}$}
\BIC{$\mathrm{DvgP}(a) \to \varphi_{5}$}
\end{ptree}
\noindent$\Pi_{10}$:
\begin{ptree}
\AXC{$\mathrm{DvgP}(a) \to \mathrm{Seq}(b) \to C \in \mathbb{R} \to C \in \mathbb{R}$\lbl{A1}}
\AXC{$\forall x \in \mathbb{R}\,(x \neq 0 \to x^{-1} \in \mathbb{R})$\lbl{(R1)}}
\AXC{$\varphi_{18}$\lbl{A2}}
\BIC{$C \in \mathbb{R} \to C \neq 0 \to C^{-1} \in \mathbb{R}$}
\AXC{$\varphi_{135}$\lbl{A1}}
\BIC{$(\mathrm{DvgP}(a) \to \mathrm{Seq}(b) \to C \in \mathbb{R} \to C \in \mathbb{R}) \to \varphi_{113}$}
\BIC{$\varphi_{113}$}
\end{ptree}
\noindent$\Pi_{11}$:
\begin{ptree}
\AXC{$\stackrel{\Pi_{9}}{\vdots}$}\noLine
\UIC{$\mathrm{DvgP}(a) \to \varphi_{5}$}
\AXC{$\stackrel{\Pi_{10}}{\vdots}$}\noLine
\UIC{$\varphi_{113}$}
\BIC{$\mathrm{DvgP}(a) \to \mathrm{Seq}(b) \to C \in \mathbb{R} \to \varphi_{13}$}
\end{ptree}
\noindent$\Pi_{12}$:
\begin{ptree}
\AXC{$\mathrm{DvgP}(a) \to \varphi_{35}$\lbl{A1}}
\AXC{$\forall x\,\forall y\,(x \in \mathbb{R} \to y \in \mathbb{R} \to 0 < x \to 0 < y \to 0 < x \cdot y)$\lbl{(R15)}}
\AXC{$\varphi_{76}$\lbl{A2}}
\BIC{$\forall y\,(K \in \mathbb{R} \to y \in \mathbb{R} \to 0 < K \to 0 < y \to 0 < K \cdot y)$}
\AXC{$\varphi_{81}$\lbl{A2}}
\BIC{$K \in \mathbb{R} \to C^{-1} \in \mathbb{R} \to 0 < K \to 0 < C^{-1} \to 0 < K \cdot C^{-1}$}
\AXC{$\varphi_{171}$\lbl{A1}}
\BIC{$(\mathrm{DvgP}(a) \to \varphi_{35}) \to \varphi_{160}$}
\BIC{$\varphi_{160}$}
\end{ptree}
\noindent$\Pi_{13}$:
\begin{ptree}
\AXC{$\mathrm{DvgP}(a) \to \mathrm{Seq}(b) \to \varphi_{23}$\lbl{A1}}
\AXC{$\stackrel{\Pi_{11}}{\vdots}$}\noLine
\UIC{$\mathrm{DvgP}(a) \to \mathrm{Seq}(b) \to C \in \mathbb{R} \to \varphi_{13}$}
\AXC{$\stackrel{\Pi_{12}}{\vdots}$}\noLine
\UIC{$\varphi_{160}$}
\BIC{$(\mathrm{DvgP}(a) \to \mathrm{Seq}(b) \to \varphi_{23}) \to \varphi_{127}$}
\BIC{$\varphi_{127}$}
\end{ptree}
\noindent$\Pi_{14}$:
\begin{ptree}
\AXC{$\mathrm{DvgP}(a) \to \mathrm{Seq}(b) \to C \in \mathbb{R} \to 0 < C \to 0 < C$\lbl{A1}}
\AXC{$\stackrel{\Pi_{4}}{\vdots}$}\noLine
\UIC{$\varphi_{92}$}
\BIC{$\mathrm{DvgP}(a) \to \mathrm{Seq}(b) \to \varphi_{37}$}
\AXC{$\stackrel{\Pi_{13}}{\vdots}$}\noLine
\UIC{$\varphi_{127}$}
\BIC{$\mathrm{DvgP}(a) \to \mathrm{Seq}(b) \to C \in \mathbb{R} \to \varphi_{16}$}
\end{ptree}
\noindent$\Pi_{15}$:
\begin{ptree}
\AXC{$\forall x\,\forall y\,\forall z\,(x < y \to y < z \to x < z)$\lbl{[O2]}}
\AXC{$\varphi_{8}$\lbl{A2}}
\BIC{$\forall y\,\forall z\,(0 < y \to y < z \to 0 < z)$}
\AXC{$\varphi_{26}$\lbl{A2}}
\BIC{$\forall z\,(0 < K \cdot C^{-1} \to K \cdot C^{-1} < z \to 0 < z)$}
\AXC{$\varphi_{62}$\lbl{A2}}
\BIC{$0 < K \cdot C^{-1} \to K \cdot C^{-1} < a_{n} \to 0 < a_{n}$}
\AXC{$\varphi_{169}$\lbl{A1}}
\BIC{$(\mathrm{DvgP}(a) \to \mathrm{Seq}(b) \to C \in \mathbb{R} \to \varphi_{16}) \to \varphi_{151}$}
\end{ptree}
\noindent$\Pi_{16}$:
\begin{ptree}
\AXC{$\stackrel{\Pi_{3}}{\vdots}$}\noLine
\UIC{$\mathrm{DvgP}(a) \to \mathrm{Seq}(b) \to \varphi_{109}$}
\AXC{$\stackrel{\Pi_{14}}{\vdots}$}\noLine
\UIC{$\mathrm{DvgP}(a) \to \mathrm{Seq}(b) \to C \in \mathbb{R} \to \varphi_{16}$}
\AXC{$\stackrel{\Pi_{15}}{\vdots}$}\noLine
\UIC{$(\mathrm{DvgP}(a) \to \mathrm{Seq}(b) \to C \in \mathbb{R} \to \varphi_{16}) \to \varphi_{151}$}
\BIC{$\varphi_{151}$}
\BIC{$\mathrm{DvgP}(a) \to \mathrm{Seq}(b) \to \varphi_{101}$}
\end{ptree}
\noindent$\Pi_{17}$:
\begin{ptree}
\AXC{$\forall a\,(\mathrm{DvgP}(a) \leftrightarrow \mathrm{Seq}(a) \land \varphi_{14})$\lbl{Def$\mathrm{DvgP}$}}
\AXC{$\varphi_{123}$\lbl{A2}}
\BIC{$\mathrm{DvgP}(a) \leftrightarrow \mathrm{Seq}(a) \land \varphi_{14}$}
\AXC{$\mathrm{DvgP}(a) \to \mathrm{DvgP}(a)$\lbl{A1}}
\AXC{$\varphi_{122}$\lbl{A1}}
\BIC{$(\mathrm{DvgP}(a) \leftrightarrow \mathrm{Seq}(a) \land \varphi_{14}) \to \mathrm{DvgP}(a) \to \varphi_{7}$}
\BIC{$\mathrm{DvgP}(a) \to \varphi_{7}$}
\end{ptree}
\noindent$\Pi_{18}$:
\begin{ptree}
\AXC{$\stackrel{\Pi_{17}}{\vdots}$}\noLine
\UIC{$\mathrm{DvgP}(a) \to \varphi_{7}$}
\AXC{$\forall a\,\forall n\,(\mathrm{Seq}(a) \to n \in \mathbb{N} \to a_{n} \in \mathbb{R})$\lbl{[S1]}}
\AXC{$\varphi_{42}$\lbl{A2}}
\BIC{$\forall n\,(\mathrm{Seq}(a) \to n \in \mathbb{N} \to a_{n} \in \mathbb{R})$}
\AXC{$\varphi_{32}$\lbl{A2}}
\BIC{$\mathrm{Seq}(a) \to n \in \mathbb{N} \to a_{n} \in \mathbb{R}$}
\AXC{$\varphi_{166}$\lbl{A1}}
\BIC{$(\mathrm{DvgP}(a) \to \varphi_{7}) \to \varphi_{142}$}
\BIC{$\varphi_{142}$}
\end{ptree}
\noindent$\Pi_{19}$:
\begin{ptree}
\AXC{$\mathrm{DvgP}(a) \to \mathrm{Seq}(b) \to \varphi_{93}$\lbl{A1}}
\AXC{$\stackrel{\Pi_{18}}{\vdots}$}\noLine
\UIC{$\varphi_{142}$}
\BIC{$\mathrm{DvgP}(a) \to \mathrm{Seq}(b) \to \varphi_{103}$}
\end{ptree}
\noindent$\Pi_{20}$:
\begin{ptree}
\AXC{$\forall a\,\forall n\,(\mathrm{Seq}(a) \to n \in \mathbb{N} \to a_{n} \in \mathbb{R})$\lbl{[S1]}}
\AXC{$\varphi_{43}$\lbl{A2}}
\BIC{$\forall n\,(\mathrm{Seq}(b) \to n \in \mathbb{N} \to b_{n} \in \mathbb{R})$}
\AXC{$\varphi_{33}$\lbl{A2}}
\BIC{$\mathrm{Seq}(b) \to n \in \mathbb{N} \to b_{n} \in \mathbb{R}$}
\AXC{$\varphi_{162}$\lbl{A1}}
\BIC{$(\mathrm{DvgP}(a) \to \mathrm{Seq}(b) \to \mathrm{Seq}(b)) \to \varphi_{143}$}
\end{ptree}
\noindent$\Pi_{21}$:
\begin{ptree}
\AXC{$\mathrm{DvgP}(a) \to \mathrm{Seq}(b) \to \mathrm{Seq}(b)$\lbl{A1}}
\AXC{$\stackrel{\Pi_{20}}{\vdots}$}\noLine
\UIC{$(\mathrm{DvgP}(a) \to \mathrm{Seq}(b) \to \mathrm{Seq}(b)) \to \varphi_{143}$}
\BIC{$\varphi_{143}$}
\end{ptree}
\noindent$\Pi_{22}$:
\begin{ptree}
\AXC{$\varphi_{12}$\lbl{[E9]}}
\AXC{$\varphi_{12} \to \varphi_{10}$\lbl{A2}}
\BIC{$\varphi_{10}$}
\AXC{$\varphi_{10} \to \varphi_{44}$\lbl{A2}}
\BIC{$\varphi_{44}$}
\AXC{$\varphi_{44} \to C \in \mathbb{R} \to \varphi_{56}$\lbl{A2}}
\BIC{$C \in \mathbb{R} \to \varphi_{56}$}
\AXC{$(C \in \mathbb{R} \to \varphi_{56}) \to \varphi_{179}$\lbl{A1}}
\BIC{$\varphi_{179}$}
\end{ptree}
\noindent$\Pi_{23}$:
\begin{ptree}
\AXC{$\mathrm{DvgP}(a) \to \mathrm{Seq}(b) \to \varphi_{93}$\lbl{A1}}
\AXC{$\stackrel{\Pi_{21}}{\vdots}$}\noLine
\UIC{$\varphi_{143}$}
\BIC{$\mathrm{DvgP}(a) \to \mathrm{Seq}(b) \to \varphi_{104}$}
\AXC{$\mathrm{DvgP}(a) \to \mathrm{Seq}(b) \to C \in \mathbb{R} \to C \in \mathbb{R}$\lbl{A1}}
\AXC{$\stackrel{\Pi_{22}}{\vdots}$}\noLine
\UIC{$\varphi_{179}$}
\BIC{$(\mathrm{DvgP}(a) \to \mathrm{Seq}(b) \to \varphi_{104}) \to \varphi_{177}$}
\BIC{$\varphi_{177}$}
\end{ptree}
\noindent$\Pi_{24}$:
\begin{ptree}
\AXC{$\stackrel{\Pi_{16}}{\vdots}$}\noLine
\UIC{$\mathrm{DvgP}(a) \to \mathrm{Seq}(b) \to \varphi_{101}$}
\AXC{$\varphi_{144}$\lbl{A1}}
\BIC{$\mathrm{DvgP}(a) \to \mathrm{Seq}(b) \to \varphi_{99}$}
\AXC{$\stackrel{\Pi_{19}}{\vdots}$}\noLine
\UIC{$\mathrm{DvgP}(a) \to \mathrm{Seq}(b) \to \varphi_{103}$}
\AXC{$\stackrel{\Pi_{23}}{\vdots}$}\noLine
\UIC{$\varphi_{177}$}
\BIC{$(\mathrm{DvgP}(a) \to \mathrm{Seq}(b) \to \varphi_{99}) \to \varphi_{153}$}
\BIC{$\varphi_{153}$}
\end{ptree}
\noindent$\Pi_{25}$:
\begin{ptree}
\AXC{$z = w \to z < a_{n} \cdot C \to w < a_{n} \cdot C$\lbl{A6}}
\UIC{$\forall w\,(z = w \to z < a_{n} \cdot C \to w < a_{n} \cdot C)$}
\UIC{$\forall z\,\forall w\,(z = w \to z < a_{n} \cdot C \to w < a_{n} \cdot C)$}
\AXC{$\forall z\,\forall w\,(z = w \to z < a_{n} \cdot C \to w < a_{n} \cdot C) \to \varphi_{21}$\lbl{A2}}
\BIC{$\varphi_{21}$}
\AXC{$\varphi_{96}$\lbl{A2}}
\BIC{$K \cdot C^{-1} \cdot C = K \to K \cdot C^{-1} \cdot C < a_{n} \cdot C \to K < a_{n} \cdot C$}
\end{ptree}
\noindent$\Pi_{26}$:
\begin{ptree}
\AXC{$\forall x \in \mathbb{R}\,(x \neq 0 \to x^{-1} \in \mathbb{R})$\lbl{(R1)}}
\AXC{$\varphi_{18}$\lbl{A2}}
\BIC{$C \in \mathbb{R} \to C \neq 0 \to C^{-1} \in \mathbb{R}$}
\AXC{$(C \in \mathbb{R} \to C \neq 0 \to C^{-1} \in \mathbb{R}) \to \varphi_{138}$\lbl{A1}}
\BIC{$\varphi_{138}$}
\end{ptree}
\noindent$\Pi_{27}$:
\begin{ptree}
\AXC{$\mathrm{DvgP}(a) \to \varphi_{35}$\lbl{A1}}
\AXC{$\forall x\,\forall y\,(x \in \mathbb{R} \to y \in \mathbb{R} \to x \cdot y \in \mathbb{R})$\lbl{(R1)}}
\AXC{$\varphi_{48}$\lbl{A2}}
\BIC{$\forall y\,(K \in \mathbb{R} \to y \in \mathbb{R} \to K \cdot y \in \mathbb{R})$}
\AXC{$\varphi_{54}$\lbl{A2}}
\BIC{$K \in \mathbb{R} \to C^{-1} \in \mathbb{R} \to K \cdot C^{-1} \in \mathbb{R}$}
\AXC{$\varphi_{168}$\lbl{A1}}
\BIC{$(\mathrm{DvgP}(a) \to \varphi_{35}) \to \varphi_{149}$}
\BIC{$\varphi_{149}$}
\end{ptree}
\noindent$\Pi_{28}$:
\begin{ptree}
\AXC{$\stackrel{\Pi_{9}}{\vdots}$}\noLine
\UIC{$\mathrm{DvgP}(a) \to \varphi_{5}$}
\AXC{$\mathrm{DvgP}(a) \to \mathrm{Seq}(b) \to C \in \mathbb{R} \to C \in \mathbb{R}$\lbl{A1}}
\AXC{$\stackrel{\Pi_{26}}{\vdots}$}\noLine
\UIC{$\varphi_{138}$}
\BIC{$(\mathrm{DvgP}(a) \to \varphi_{5}) \to \mathrm{DvgP}(a) \to \mathrm{Seq}(b) \to \varphi_{102}$}
\BIC{$\mathrm{DvgP}(a) \to \mathrm{Seq}(b) \to \varphi_{102}$}
\AXC{$\stackrel{\Pi_{27}}{\vdots}$}\noLine
\UIC{$\varphi_{149}$}
\BIC{$\mathrm{DvgP}(a) \to \mathrm{Seq}(b) \to \varphi_{107}$}
\end{ptree}
\noindent$\Pi_{29}$:
\begin{ptree}
\AXC{$\stackrel{\Pi_{28}}{\vdots}$}\noLine
\UIC{$\mathrm{DvgP}(a) \to \mathrm{Seq}(b) \to \varphi_{107}$}
\AXC{$\varphi_{34}$\lbl{[F32]}}
\AXC{$\varphi_{34} \to \varphi_{65}$\lbl{A2}}
\BIC{$\varphi_{65}$}
\AXC{$\varphi_{65} \to \varphi_{78}$\lbl{A2}}
\BIC{$\varphi_{78}$}
\AXC{$\varphi_{78} \to K \cdot C^{-1} \in \mathbb{R} \to \varphi_{58}$\lbl{A2}}
\BIC{$K \cdot C^{-1} \in \mathbb{R} \to \varphi_{58}$}
\AXC{$\varphi_{180}$\lbl{A1}}
\BIC{$(\mathrm{DvgP}(a) \to \mathrm{Seq}(b) \to \varphi_{107}) \to \varphi_{174}$}
\BIC{$\varphi_{174}$}
\end{ptree}
\noindent$\Pi_{30}$:
\begin{ptree}
\AXC{$\mathrm{DvgP}(a) \to \mathrm{Seq}(b) \to C \in \mathbb{R} \to C \in \mathbb{R}$\lbl{A1}}
\AXC{$\stackrel{\Pi_{19}}{\vdots}$}\noLine
\UIC{$\mathrm{DvgP}(a) \to \mathrm{Seq}(b) \to \varphi_{103}$}
\AXC{$\stackrel{\Pi_{29}}{\vdots}$}\noLine
\UIC{$\varphi_{174}$}
\BIC{$(\mathrm{DvgP}(a) \to \mathrm{Seq}(b) \to C \in \mathbb{R} \to C \in \mathbb{R}) \to \varphi_{163}$}
\BIC{$\varphi_{163}$}
\end{ptree}
\noindent$\Pi_{31}$:
\begin{ptree}
\AXC{$\stackrel{\Pi_{3}}{\vdots}$}\noLine
\UIC{$\mathrm{DvgP}(a) \to \mathrm{Seq}(b) \to \varphi_{109}$}
\AXC{$\mathrm{DvgP}(a) \to \mathrm{Seq}(b) \to C \in \mathbb{R} \to 0 < C \to 0 < C$\lbl{A1}}
\AXC{$\stackrel{\Pi_{30}}{\vdots}$}\noLine
\UIC{$\varphi_{163}$}
\BIC{$\varphi_{156}$}
\BIC{$\mathrm{DvgP}(a) \to \mathrm{Seq}(b) \to \varphi_{114}$}
\end{ptree}
\noindent$\Pi_{32}$:
\begin{ptree}
\AXC{$\mathrm{DvgP}(a) \to \varphi_{35}$\lbl{A1}}
\AXC{$\forall x\,\forall y\,(x \in \mathbb{R} \to y \in \mathbb{R} \to y \neq 0 \to x \cdot y^{-1} \cdot y = x)$\lbl{[E23]}}
\AXC{$\varphi_{68}$\lbl{A2}}
\BIC{$\forall y\,(K \in \mathbb{R} \to y \in \mathbb{R} \to y \neq 0 \to K \cdot y^{-1} \cdot y = K)$}
\AXC{$\varphi_{63}$\lbl{A2}}
\BIC{$K \in \mathbb{R} \to C \in \mathbb{R} \to C \neq 0 \to K \cdot C^{-1} \cdot C = K$}
\AXC{$\varphi_{164}$\lbl{A1}}
\BIC{$(\mathrm{DvgP}(a) \to \varphi_{35}) \to \varphi_{139}$}
\BIC{$\varphi_{139}$}
\end{ptree}
\noindent$\Pi_{33}$:
\begin{ptree}
\AXC{$\stackrel{\Pi_{9}}{\vdots}$}\noLine
\UIC{$\mathrm{DvgP}(a) \to \varphi_{5}$}
\AXC{$\mathrm{DvgP}(a) \to \mathrm{Seq}(b) \to C \in \mathbb{R} \to C \in \mathbb{R}$\lbl{A1}}
\AXC{$\stackrel{\Pi_{32}}{\vdots}$}\noLine
\UIC{$\varphi_{139}$}
\BIC{$(\mathrm{DvgP}(a) \to \varphi_{5}) \to \mathrm{DvgP}(a) \to \mathrm{Seq}(b) \to \varphi_{105}$}
\BIC{$\mathrm{DvgP}(a) \to \mathrm{Seq}(b) \to \varphi_{105}$}
\AXC{$\varphi_{172}$\lbl{A1}}
\BIC{$(\mathrm{DvgP}(a) \to \mathrm{Seq}(b) \to \varphi_{114}) \to \varphi_{134}$}
\end{ptree}
\noindent$\Pi_{34}$:
\begin{ptree}
\AXC{$\stackrel{\Pi_{31}}{\vdots}$}\noLine
\UIC{$\mathrm{DvgP}(a) \to \mathrm{Seq}(b) \to \varphi_{114}$}
\AXC{$\stackrel{\Pi_{33}}{\vdots}$}\noLine
\UIC{$(\mathrm{DvgP}(a) \to \mathrm{Seq}(b) \to \varphi_{114}) \to \varphi_{134}$}
\BIC{$\varphi_{134}$}
\end{ptree}
\noindent$\Pi_{35}$:
\begin{ptree}
\AXC{$\forall x\,\forall y\,\forall z\,(x < y \to y \leq z \to x < z)$\lbl{[O3]}}
\AXC{$\varphi_{1}$\lbl{A2}}
\BIC{$\forall y\,\forall z\,(K < y \to y \leq z \to K < z)$}
\AXC{$\varphi_{22}$\lbl{A2}}
\BIC{$\forall z\,(K < a_{n} \cdot C \to a_{n} \cdot C \leq z \to K < z)$}
\AXC{$\varphi_{83}$\lbl{A2}}
\BIC{$K < a_{n} \cdot C \to a_{n} \cdot C \leq a_{n} \cdot b_{n} \to K < a_{n} \cdot b_{n}$}
\AXC{$\varphi_{173}$\lbl{A1}}
\BIC{$(\mathrm{DvgP}(a) \to \mathrm{Seq}(b) \to \varphi_{106}) \to \varphi_{157}$}
\end{ptree}
\noindent$\Pi_{36}$:
\begin{ptree}
\AXC{$\stackrel{\Pi_{25}}{\vdots}$}\noLine
\UIC{$K \cdot C^{-1} \cdot C = K \to K \cdot C^{-1} \cdot C < a_{n} \cdot C \to K < a_{n} \cdot C$}
\AXC{$\stackrel{\Pi_{34}}{\vdots}$}\noLine
\UIC{$\varphi_{134}$}
\BIC{$\mathrm{DvgP}(a) \to \mathrm{Seq}(b) \to \varphi_{106}$}
\AXC{$\stackrel{\Pi_{35}}{\vdots}$}\noLine
\UIC{$(\mathrm{DvgP}(a) \to \mathrm{Seq}(b) \to \varphi_{106}) \to \varphi_{157}$}
\BIC{$\varphi_{157}$}
\end{ptree}
\noindent$\Pi_{37}$:
\begin{ptree}
\AXC{$\mathrm{DvgP}(a) \to \mathrm{Seq}(b) \to \varphi_{93}$\lbl{A1}}
\AXC{$\stackrel{\Pi_{1}}{\vdots}$}\noLine
\UIC{$\mathrm{DvgP}(a) \to \varphi_{29}$}
\AXC{$(\mathrm{DvgP}(a) \to \varphi_{29}) \to \varphi_{141}$\lbl{A1}}
\BIC{$\varphi_{141}$}
\BIC{$\mathrm{DvgP}(a) \to \mathrm{Seq}(b) \to \varphi_{98}$}
\AXC{$\stackrel{\Pi_{24}}{\vdots}$}\noLine
\UIC{$\varphi_{153}$}
\BIC{$\mathrm{DvgP}(a) \to \mathrm{Seq}(b) \to \varphi_{115}$}
\AXC{$\stackrel{\Pi_{36}}{\vdots}$}\noLine
\UIC{$\varphi_{157}$}
\BIC{$\mathrm{DvgP}(a) \to \mathrm{Seq}(b) \to \varphi_{111}$}
\end{ptree}
\noindent$\Pi_{38}$:
\begin{ptree}
\AXC{$\forall a\,\forall b\,(\mathrm{Seq}(a) \to \varphi_{36})$\lbl{[S4]}}
\AXC{$\forall a\,\forall b\,(\mathrm{Seq}(a) \to \varphi_{36}) \to \forall b\,(\mathrm{Seq}(a) \to \varphi_{36})$\lbl{A2}}
\BIC{$\forall b\,(\mathrm{Seq}(a) \to \varphi_{36})$}
\end{ptree}
\noindent$\Pi_{39}$:
\begin{ptree}
\AXC{$\stackrel{\Pi_{38}}{\vdots}$}\noLine
\UIC{$\forall b\,(\mathrm{Seq}(a) \to \varphi_{36})$}
\AXC{$\forall b\,(\mathrm{Seq}(a) \to \varphi_{36}) \to \mathrm{Seq}(a) \to \varphi_{36}$\lbl{A2}}
\BIC{$\mathrm{Seq}(a) \to \varphi_{36}$}
\AXC{$(\mathrm{Seq}(a) \to \varphi_{36}) \to (\mathrm{DvgP}(a) \to \varphi_{7}) \to \varphi_{108}$\lbl{A1}}
\BIC{$(\mathrm{DvgP}(a) \to \varphi_{7}) \to \varphi_{108}$}
\end{ptree}
\noindent$\Pi_{40}$:
\begin{ptree}
\AXC{$\mathrm{DvgP}(a) \to \mathrm{Seq}(b) \to \mathrm{Seq}(b)$\lbl{A1}}
\AXC{$\stackrel{\Pi_{17}}{\vdots}$}\noLine
\UIC{$\mathrm{DvgP}(a) \to \varphi_{7}$}
\AXC{$\stackrel{\Pi_{39}}{\vdots}$}\noLine
\UIC{$(\mathrm{DvgP}(a) \to \varphi_{7}) \to \varphi_{108}$}
\BIC{$\varphi_{108}$}
\BIC{$\mathrm{DvgP}(a) \to \mathrm{Seq}(b) \to C \in \mathbb{R} \to 0 < C \to \varphi_{59}$}
\end{ptree}
\noindent$\Pi_{41}$:
\begin{ptree}
\AXC{$\varphi_{64}$\lbl{A2}}
\AXC{$\stackrel{\Pi_{40}}{\vdots}$}\noLine
\UIC{$\mathrm{DvgP}(a) \to \mathrm{Seq}(b) \to C \in \mathbb{R} \to 0 < C \to \varphi_{59}$}
\AXC{$\varphi_{146}$\lbl{A1}}
\BIC{$\mathrm{DvgP}(a) \to \mathrm{Seq}(b) \to \varphi_{126}$}
\AXC{$\varphi_{170}$\lbl{A1}}
\BIC{$\varphi_{64} \to \mathrm{DvgP}(a) \to \mathrm{Seq}(b) \to \varphi_{125}$}
\BIC{$\mathrm{DvgP}(a) \to \mathrm{Seq}(b) \to \varphi_{125}$}
\end{ptree}
\noindent$\Pi_{42}$:
\begin{ptree}
\AXC{$\mathrm{DvgP}(a) \to \mathrm{Seq}(b) \to \varphi_{93}$\lbl{A1}}
\AXC{$\stackrel{\Pi_{41}}{\vdots}$}\noLine
\UIC{$\mathrm{DvgP}(a) \to \mathrm{Seq}(b) \to \varphi_{125}$}
\AXC{$(\mathrm{DvgP}(a) \to \mathrm{Seq}(b) \to \varphi_{125}) \to \varphi_{152}$\lbl{A1}}
\BIC{$\varphi_{152}$}
\BIC{$\mathrm{DvgP}(a) \to \mathrm{Seq}(b) \to \varphi_{118}$}
\end{ptree}
\noindent$\Pi_{43}$:
\begin{ptree}
\AXC{$\stackrel{\Pi_{42}}{\vdots}$}\noLine
\UIC{$\mathrm{DvgP}(a) \to \mathrm{Seq}(b) \to \varphi_{118}$}
\AXC{$(\mathrm{DvgP}(a) \to \mathrm{Seq}(b) \to \varphi_{118}) \to \varphi_{137}$\lbl{A1}}
\BIC{$\varphi_{137}$}
\end{ptree}
\noindent$\Pi_{44}$:
\begin{ptree}
\AXC{$\stackrel{\Pi_{37}}{\vdots}$}\noLine
\UIC{$\mathrm{DvgP}(a) \to \mathrm{Seq}(b) \to \varphi_{111}$}
\AXC{$\stackrel{\Pi_{5}}{\vdots}$}\noLine
\UIC{$\forall x\,\forall y\,(x = y \to y = x)$}
\AXC{$\varphi_{6}$\lbl{A2}}
\BIC{$\forall y\,(a \otimes b(n) = y \to y = a \otimes b(n))$}
\AXC{$\varphi_{80}$\lbl{A2}}
\BIC{$a \otimes b(n) = a_{n} \cdot b_{n} \to a_{n} \cdot b_{n} = a \otimes b(n)$}
\AXC{$\stackrel{\Pi_{43}}{\vdots}$}\noLine
\UIC{$\varphi_{137}$}
\BIC{$\mathrm{DvgP}(a) \to \mathrm{Seq}(b) \to \varphi_{119}$}
\AXC{$\varphi_{175}$\lbl{A1}}
\BIC{$(\mathrm{DvgP}(a) \to \mathrm{Seq}(b) \to \varphi_{111}) \to \varphi_{136}$}
\BIC{$\varphi_{136}$}
\end{ptree}
\noindent$\Pi_{45}$:
\begin{ptree}
\AXC{$\varphi_{155}$\lbl{A3}}
\AXC{$z = w \to K < z \to K < w$\lbl{A6}}
\UIC{$\forall w\,(z = w \to K < z \to K < w)$}
\UIC{$\forall z\,\forall w\,(z = w \to K < z \to K < w)$}
\AXC{$\varphi_{46}$\lbl{A2}}
\BIC{$\forall w\,(a_{n} \cdot b_{n} = w \to K < a_{n} \cdot b_{n} \to K < w)$}
\AXC{$\varphi_{97}$\lbl{A2}}
\BIC{$a_{n} \cdot b_{n} = a \otimes b(n) \to K < a_{n} \cdot b_{n} \to K < a \otimes b(n)$}
\AXC{$\stackrel{\Pi_{44}}{\vdots}$}\noLine
\UIC{$\varphi_{136}$}
\BIC{$\mathrm{DvgP}(a) \to \mathrm{Seq}(b) \to \varphi_{110}$}
\AXC{$\varphi_{154}$\lbl{A1}}
\BIC{$\varphi_{86} \to n \in \mathbb{N} \to N \leq n \to K < a \otimes b(n)$}
\UIC{$\forall n\,(\varphi_{86} \to n \in \mathbb{N} \to N \leq n \to K < a \otimes b(n))$}
\AXC{$\varphi_{181}$\lbl{A1}}
\BIC{$\varphi_{155} \to (\varphi_{86} \to \forall n\,(\varphi_{86})) \to \mathrm{DvgP}(a) \to \mathrm{Seq}(b) \to \varphi_{112}$}
\BIC{$(\varphi_{86} \to \forall n\,(\varphi_{86})) \to \mathrm{DvgP}(a) \to \mathrm{Seq}(b) \to \varphi_{112}$}
\end{ptree}
\noindent$\Pi_{46}$:
\begin{ptree}
\AXC{$\mathrm{DvgP}(a) \to \mathrm{Seq}(b) \to \varphi_{116}$\lbl{A1}}
\AXC{$(\mathrm{DvgP}(a) \to \mathrm{Seq}(b) \to \varphi_{116}) \to \mathrm{DvgP}(a) \to \varphi_{87}$\lbl{A1}}
\BIC{$\mathrm{DvgP}(a) \to \varphi_{87}$}
\AXC{$(\mathrm{DvgP}(a) \to \varphi_{87}) \to \varphi_{158}$\lbl{A1}}
\BIC{$\varphi_{158}$}
\end{ptree}
\noindent$\Pi_{47}$:
\begin{ptree}
\AXC{$\varphi_{86} \to \forall n\,(\varphi_{86})$\lbl{A4}}
\AXC{$\stackrel{\Pi_{45}}{\vdots}$}\noLine
\UIC{$(\varphi_{86} \to \forall n\,(\varphi_{86})) \to \mathrm{DvgP}(a) \to \mathrm{Seq}(b) \to \varphi_{112}$}
\BIC{$\mathrm{DvgP}(a) \to \mathrm{Seq}(b) \to \varphi_{112}$}
\AXC{$\stackrel{\Pi_{46}}{\vdots}$}\noLine
\UIC{$\varphi_{158}$}
\BIC{$\mathrm{DvgP}(a) \to \mathrm{Seq}(b) \to \varphi_{117}$}
\AXC{$(\mathrm{DvgP}(a) \to \mathrm{Seq}(b) \to \varphi_{117}) \to \varphi_{140}$\lbl{A1}}
\BIC{$\varphi_{140}$}
\end{ptree}
\noindent$\Pi_{48}$:
\begin{ptree}
\AXC{$\varphi_{3} \to \neg (N \in \mathbb{N} \land \forall n \in \mathbb{N}\,(N \leq n \to K < a \otimes b(n)))$\lbl{A2}}
\AXC{$\stackrel{\Pi_{47}}{\vdots}$}\noLine
\UIC{$\varphi_{140}$}
\BIC{$\mathrm{DvgP}(a) \to \mathrm{Seq}(b) \to \varphi_{121}$}
\AXC{$\varphi_{161}$\lbl{A1}}
\BIC{$\varphi_{90} \to \neg (N \in \mathbb{N} \land \forall n \in \mathbb{N}\,(N \leq n \to K \cdot C^{-1} < a_{n}))$}
\UIC{$\varphi_{128}$}
\end{ptree}
\noindent$\Pi_{49}$:
\begin{ptree}
\AXC{$\mathrm{DvgP}(a) \to \varphi_{35}$\lbl{A1}}
\AXC{$\forall x\,\forall y\,(x \in \mathbb{R} \to y \in \mathbb{R} \to x \cdot y \in \mathbb{R})$\lbl{(R1)}}
\AXC{$\varphi_{48}$\lbl{A2}}
\BIC{$\forall y\,(K \in \mathbb{R} \to y \in \mathbb{R} \to K \cdot y \in \mathbb{R})$}
\AXC{$\varphi_{54}$\lbl{A2}}
\BIC{$K \in \mathbb{R} \to C^{-1} \in \mathbb{R} \to K \cdot C^{-1} \in \mathbb{R}$}
\AXC{$\varphi_{147}$\lbl{A1}}
\BIC{$(\mathrm{DvgP}(a) \to \varphi_{35}) \to \varphi_{129}$}
\BIC{$\varphi_{129}$}
\end{ptree}
\noindent$\Pi_{50}$:
\begin{ptree}
\AXC{$\mathrm{DvgP}(a) \to \mathrm{DvgP}(a)$\lbl{A1}}
\AXC{$(\mathrm{DvgP}(a) \to \mathrm{DvgP}(a)) \to \varphi_{133}$\lbl{A1}}
\BIC{$\varphi_{133}$}
\end{ptree}
\noindent$\Pi_{51}$:
\begin{ptree}
\AXC{$\forall a\,(\mathrm{DvgP}(a) \leftrightarrow \mathrm{Seq}(a) \land \varphi_{14})$\lbl{Def$\mathrm{DvgP}$}}
\AXC{$\varphi_{123}$\lbl{A2}}
\BIC{$\mathrm{DvgP}(a) \leftrightarrow \mathrm{Seq}(a) \land \varphi_{14}$}
\AXC{$\stackrel{\Pi_{50}}{\vdots}$}\noLine
\UIC{$\varphi_{133}$}
\BIC{$\mathrm{DvgP}(a) \to \mathrm{Seq}(b) \to \varphi_{85}$}
\AXC{$(\mathrm{DvgP}(a) \to \mathrm{Seq}(b) \to \varphi_{85}) \to \varphi_{148}$\lbl{A1}}
\BIC{$\varphi_{148}$}
\end{ptree}
\noindent$\Pi_{52}$:
\begin{ptree}
\AXC{$\varphi_{14} \to K \cdot C^{-1} \in \mathbb{R} \to \varphi_{24}$\lbl{A2}}
\AXC{$\stackrel{\Pi_{51}}{\vdots}$}\noLine
\UIC{$\varphi_{148}$}
\BIC{$\mathrm{DvgP}(a) \to \mathrm{Seq}(b) \to C \in \mathbb{R} \to 0 < C \to \varphi_{91}$}
\AXC{$(\mathrm{DvgP}(a) \to \mathrm{Seq}(b) \to C \in \mathbb{R} \to 0 < C \to \varphi_{91}) \to \varphi_{150}$\lbl{A1}}
\BIC{$\varphi_{150}$}
\end{ptree}
\noindent$\Pi_{53}$:
\begin{ptree}
\AXC{$\stackrel{\Pi_{14}}{\vdots}$}\noLine
\UIC{$\mathrm{DvgP}(a) \to \mathrm{Seq}(b) \to C \in \mathbb{R} \to \varphi_{16}$}
\AXC{$\stackrel{\Pi_{11}}{\vdots}$}\noLine
\UIC{$\mathrm{DvgP}(a) \to \mathrm{Seq}(b) \to C \in \mathbb{R} \to \varphi_{13}$}
\AXC{$\stackrel{\Pi_{49}}{\vdots}$}\noLine
\UIC{$\varphi_{129}$}
\BIC{$\mathrm{DvgP}(a) \to \mathrm{Seq}(b) \to C \in \mathbb{R} \to \varphi_{25}$}
\AXC{$\stackrel{\Pi_{52}}{\vdots}$}\noLine
\UIC{$\varphi_{150}$}
\BIC{$(\mathrm{DvgP}(a) \to \mathrm{Seq}(b) \to C \in \mathbb{R} \to \varphi_{16}) \to \varphi_{88}$}
\BIC{$\varphi_{88}$}
\end{ptree}
\noindent$\Pi_{54}$:
\begin{ptree}
\AXC{$\varphi_{90} \to \forall N\,(\varphi_{90})$\lbl{A4}}
\AXC{$\varphi_{128} \to \forall N\,(\varphi_{90}) \to \varphi_{2}$\lbl{A3}}
\AXC{$\stackrel{\Pi_{48}}{\vdots}$}\noLine
\UIC{$\varphi_{128}$}
\AXC{$\stackrel{\Pi_{53}}{\vdots}$}\noLine
\UIC{$\varphi_{88}$}
\AXC{$\varphi_{88} \to \varphi_{128} \to (\varphi_{128} \to \forall N\,(\varphi_{90}) \to \varphi_{2}) \to (\varphi_{90} \to \forall N\,(\varphi_{90})) \to \varphi_{89}$\lbl{A1}}
\BIC{$\varphi_{128} \to (\varphi_{128} \to \forall N\,(\varphi_{90}) \to \varphi_{2}) \to (\varphi_{90} \to \forall N\,(\varphi_{90})) \to \varphi_{89}$}
\BIC{$(\varphi_{128} \to \forall N\,(\varphi_{90}) \to \varphi_{2}) \to (\varphi_{90} \to \forall N\,(\varphi_{90})) \to \varphi_{89}$}
\BIC{$(\varphi_{90} \to \forall N\,(\varphi_{90})) \to \varphi_{89}$}
\BIC{$\varphi_{89}$}
\end{ptree}
\noindent$\Pi_{55}$:
\begin{ptree}
\AXC{$\varphi_{9} \to \forall K\,(\varphi_{9})$\lbl{A4}}
\AXC{$\forall K\,(\varphi_{9} \to \varphi_{31}) \to \forall K\,(\varphi_{9}) \to \varphi_{38}$\lbl{A3}}
\AXC{$\stackrel{\Pi_{54}}{\vdots}$}\noLine
\UIC{$\varphi_{89}$}
\AXC{$\varphi_{89} \to \varphi_{9} \to \varphi_{31}$\lbl{A1}}
\BIC{$\varphi_{9} \to \varphi_{31}$}
\UIC{$\forall K\,(\varphi_{9} \to \varphi_{31})$}
\AXC{$\varphi_{178}$\lbl{A1}}
\BIC{$(\forall K\,(\varphi_{9} \to \varphi_{31}) \to \forall K\,(\varphi_{9}) \to \varphi_{38}) \to (\varphi_{9} \to \forall K\,(\varphi_{9})) \to \varphi_{95}$}
\BIC{$(\varphi_{9} \to \forall K\,(\varphi_{9})) \to \varphi_{95}$}
\BIC{$\varphi_{95}$}
\end{ptree}
\noindent$\Pi_{56}$:
\begin{ptree}
\AXC{$\forall a\,(\mathrm{DvgP}(a) \leftrightarrow \mathrm{Seq}(a) \land \varphi_{14})$\lbl{Def$\mathrm{DvgP}$}}
\AXC{$\varphi_{130}$\lbl{A2}}
\BIC{$\mathrm{DvgP}(a \otimes b) \leftrightarrow \mathrm{Seq}(a \otimes b) \land \varphi_{38}$}
\AXC{$\varphi_{165}$\lbl{A1}}
\BIC{$(\mathrm{DvgP}(a) \to \varphi_{17}) \to \varphi_{95} \to \mathrm{DvgP}(a) \to \varphi_{19}$}
\end{ptree}
\noindent 主木:
\begin{ptree}
\AXC{$\stackrel{\Pi_{55}}{\vdots}$}\noLine
\UIC{$\varphi_{95}$}
\AXC{$\stackrel{\Pi_{40}}{\vdots}$}\noLine
\UIC{$\mathrm{DvgP}(a) \to \mathrm{Seq}(b) \to C \in \mathbb{R} \to 0 < C \to \varphi_{59}$}
\AXC{$\varphi_{131}$\lbl{A1}}
\BIC{$\mathrm{DvgP}(a) \to \varphi_{17}$}
\AXC{$\stackrel{\Pi_{56}}{\vdots}$}\noLine
\UIC{$(\mathrm{DvgP}(a) \to \varphi_{17}) \to \varphi_{95} \to \mathrm{DvgP}(a) \to \varphi_{19}$}
\BIC{$\varphi_{95} \to \mathrm{DvgP}(a) \to \varphi_{19}$}
\BIC{$\mathrm{DvgP}(a) \to \varphi_{19}$}
\UIC{$\forall C\,(\mathrm{DvgP}(a) \to \varphi_{19})$}
\UIC{$\forall b\,\forall C\,(\mathrm{DvgP}(a) \to \varphi_{19})$}
\UIC{$\forall a\,\forall b\,\forall C\,(\mathrm{DvgP}(a) \to \varphi_{19})$}
\end{ptree}
\begin{abbrevlist}
\item $\varphi_{1} :\equiv \forall x\,\forall y\,\forall z\,(x < y \to y \leq z \to x < z) \to \forall y\,\forall z\,(K < y \to y \leq z \to K < z)$
\item $\varphi_{2} :\equiv \forall N\,(\neg (N \in \mathbb{N} \land \forall n \in \mathbb{N}\,(N \leq n \to K \cdot C^{-1} < a_{n})))$
\item $\varphi_{3} :\equiv \forall N\,(\neg (N \in \mathbb{N} \land \forall n \in \mathbb{N}\,(N \leq n \to K < a \otimes b(n))))$
\item $\varphi_{4} :\equiv \mathrm{Seq}(b) \to C \in \mathbb{R} \to 0 < C \to \forall n \in \mathbb{N}\,(C \leq b_{n}) \to 0 \neq C$
\item $\varphi_{5} :\equiv \mathrm{Seq}(b) \to C \in \mathbb{R} \to 0 < C \to \forall n \in \mathbb{N}\,(C \leq b_{n}) \to C \neq 0$
\item $\varphi_{6} :\equiv \forall x\,\forall y\,(x = y \to y = x) \to \forall y\,(a \otimes b(n) = y \to y = a \otimes b(n))$
\item $\varphi_{7} :\equiv \mathrm{Seq}(b) \to C \in \mathbb{R} \to 0 < C \to \forall n \in \mathbb{N}\,(C \leq b_{n}) \to \mathrm{Seq}(a)$
\item $\varphi_{8} :\equiv \forall x\,\forall y\,\forall z\,(x < y \to y < z \to x < z) \to \forall y\,\forall z\,(0 < y \to y < z \to 0 < z)$
\item $\varphi_{9} :\equiv \mathrm{DvgP}(a) \land \mathrm{Seq}(b) \land C \in \mathbb{R} \land 0 < C \land \forall n \in \mathbb{N}\,(C \leq b_{n})$
\item $\varphi_{10} :\equiv \forall y\,\forall z\,(C \in \mathbb{R} \to y \in \mathbb{R} \to z \in \mathbb{R} \to 0 \leq z \to C \leq y \to z \cdot C \leq z \cdot y)$
\item $\varphi_{11} :\equiv \forall x \in \mathbb{R}\,(0 < x \to 0 < x^{-1}) \to C \in \mathbb{R} \to 0 < C \to 0 < C^{-1}$
\item $\varphi_{12} :\equiv \forall x\,\forall y\,\forall z\,(x \in \mathbb{R} \to y \in \mathbb{R} \to z \in \mathbb{R} \to 0 \leq z \to x \leq y \to z \cdot x \leq z \cdot y)$
\item $\varphi_{13} :\equiv 0 < C \to \forall n \in \mathbb{N}\,(C \leq b_{n}) \to K \in \mathbb{R} \to 0 < K \to C^{-1} \in \mathbb{R}$
\item $\varphi_{14} :\equiv \forall K \in \mathbb{R}\,(0 < K \to \exists N \in \mathbb{N}\,\forall n \in \mathbb{N}\,(N \leq n \to K < a_{n}))$
\item $\varphi_{15} :\equiv C \in \mathbb{R} \to 0 < C \to \forall n \in \mathbb{N}\,(C \leq b_{n}) \to \forall n \in \mathbb{N}\,(C \leq b_{n})$
\item $\varphi_{16} :\equiv 0 < C \to \forall n \in \mathbb{N}\,(C \leq b_{n}) \to K \in \mathbb{R} \to 0 < K \to 0 < K \cdot C^{-1}$
\item $\varphi_{17} :\equiv \mathrm{Seq}(b) \to C \in \mathbb{R} \to 0 < C \to \forall n \in \mathbb{N}\,(C \leq b_{n}) \to \mathrm{Seq}(a \otimes b)$
\item $\varphi_{18} :\equiv \forall x \in \mathbb{R}\,(x \neq 0 \to x^{-1} \in \mathbb{R}) \to C \in \mathbb{R} \to C \neq 0 \to C^{-1} \in \mathbb{R}$
\item $\varphi_{19} :\equiv \mathrm{Seq}(b) \to C \in \mathbb{R} \to 0 < C \to \forall n \in \mathbb{N}\,(C \leq b_{n}) \to \mathrm{DvgP}(a \otimes b)$
\item $\varphi_{20} :\equiv 0 < K \to N \in \mathbb{N} \land \forall n \in \mathbb{N}\,(N \leq n \to K \cdot C^{-1} < a_{n}) \to N \in \mathbb{N}$
\item $\varphi_{21} :\equiv \forall w\,(K \cdot C^{-1} \cdot C = w \to K \cdot C^{-1} \cdot C < a_{n} \cdot C \to w < a_{n} \cdot C)$
\item $\varphi_{22} :\equiv \forall y\,\forall z\,(K < y \to y \leq z \to K < z) \to \forall z\,(K < a_{n} \cdot C \to a_{n} \cdot C \leq z \to K < z)$
\item $\varphi_{23} :\equiv C \in \mathbb{R} \to 0 < C \to \forall n \in \mathbb{N}\,(C \leq b_{n}) \to K \in \mathbb{R} \to 0 < K \to 0 < K$
\item $\varphi_{24} :\equiv 0 < K \cdot C^{-1} \to \exists N \in \mathbb{N}\,\forall n \in \mathbb{N}\,(N \leq n \to K \cdot C^{-1} < a_{n})$
\item $\varphi_{25} :\equiv 0 < C \to \forall n \in \mathbb{N}\,(C \leq b_{n}) \to K \in \mathbb{R} \to 0 < K \to K \cdot C^{-1} \in \mathbb{R}$
\item $\varphi_{26} :\equiv \forall y\,\forall z\,(0 < y \to y < z \to 0 < z) \to \forall z\,(0 < K \cdot C^{-1} \to K \cdot C^{-1} < z \to 0 < z)$
\item $\varphi_{27} :\equiv N \in \mathbb{N} \land \forall n \in \mathbb{N}\,(N \leq n \to K \cdot C^{-1} < a_{n}) \to n \in \mathbb{N} \to n \in \mathbb{N}$
\item $\varphi_{28} :\equiv N \in \mathbb{N} \land \forall n \in \mathbb{N}\,(N \leq n \to K \cdot C^{-1} < a_{n}) \to n \in \mathbb{N} \to N \leq n \to N \leq n$
\item $\varphi_{29} :\equiv \mathrm{Seq}(b) \to C \in \mathbb{R} \to 0 < C \to \forall n \in \mathbb{N}\,(C \leq b_{n}) \to n \in \mathbb{N} \to C \leq b_{n}$
\item $\varphi_{30} :\equiv K \in \mathbb{R} \to 0 < K \to \exists N \in \mathbb{N}\,\forall n \in \mathbb{N}\,(N \leq n \to K \cdot C^{-1} < a_{n})$
\item $\varphi_{31} :\equiv K \in \mathbb{R} \to 0 < K \to \exists N \in \mathbb{N}\,\forall n \in \mathbb{N}\,(N \leq n \to K < a \otimes b(n))$
\item $\varphi_{32} :\equiv \forall n\,(\mathrm{Seq}(a) \to n \in \mathbb{N} \to a_{n} \in \mathbb{R}) \to \mathrm{Seq}(a) \to n \in \mathbb{N} \to a_{n} \in \mathbb{R}$
\item $\varphi_{33} :\equiv \forall n\,(\mathrm{Seq}(b) \to n \in \mathbb{N} \to b_{n} \in \mathbb{R}) \to \mathrm{Seq}(b) \to n \in \mathbb{N} \to b_{n} \in \mathbb{R}$
\item $\varphi_{34} :\equiv \forall x\,\forall y\,\forall z\,(x \in \mathbb{R} \to y \in \mathbb{R} \to z \in \mathbb{R} \to 0 < z \to x < y \to x \cdot z < y \cdot z)$
\item $\varphi_{35} :\equiv \mathrm{Seq}(b) \to C \in \mathbb{R} \to 0 < C \to \forall n \in \mathbb{N}\,(C \leq b_{n}) \to K \in \mathbb{R} \to K \in \mathbb{R}$
\item $\varphi_{36} :\equiv \mathrm{Seq}(b) \to \mathrm{Seq}(a \otimes b) \land \forall n \in \mathbb{N}\,(a \otimes b(n) = a_{n} \cdot b_{n})$
\item $\varphi_{37} :\equiv C \in \mathbb{R} \to 0 < C \to \forall n \in \mathbb{N}\,(C \leq b_{n}) \to K \in \mathbb{R} \to 0 < K \to 0 < C^{-1}$
\item $\varphi_{38} :\equiv \forall K \in \mathbb{R}\,(0 < K \to \exists N \in \mathbb{N}\,\forall n \in \mathbb{N}\,(N \leq n \to K < a \otimes b(n)))$
\item $\varphi_{39} :\equiv N \in \mathbb{N} \land \forall n \in \mathbb{N}\,(N \leq n \to K \cdot C^{-1} < a_{n}) \to n \in \mathbb{N} \to N \leq n \to C \leq b_{n}$
\item $\varphi_{40} :\equiv N \in \mathbb{N} \land \forall n \in \mathbb{N}\,(N \leq n \to K \cdot C^{-1} < a_{n}) \to n \in \mathbb{N} \to N \leq n \to 0 \leq a_{n}$
\item $\varphi_{41} :\equiv \forall n \in \mathbb{N}\,(N \leq n \to K \cdot C^{-1} < a_{n}) \to n \in \mathbb{N} \to N \leq n \to K \cdot C^{-1} < a_{n}$
\item $\varphi_{42} :\equiv \forall a\,\forall n\,(\mathrm{Seq}(a) \to n \in \mathbb{N} \to a_{n} \in \mathbb{R}) \to \forall n\,(\mathrm{Seq}(a) \to n \in \mathbb{N} \to a_{n} \in \mathbb{R})$
\item $\varphi_{43} :\equiv \forall a\,\forall n\,(\mathrm{Seq}(a) \to n \in \mathbb{N} \to a_{n} \in \mathbb{R}) \to \forall n\,(\mathrm{Seq}(b) \to n \in \mathbb{N} \to b_{n} \in \mathbb{R})$
\item $\varphi_{44} :\equiv \forall z\,(C \in \mathbb{R} \to b_{n} \in \mathbb{R} \to z \in \mathbb{R} \to 0 \leq z \to C \leq b_{n} \to z \cdot C \leq z \cdot b_{n})$
\item $\varphi_{45} :\equiv N \in \mathbb{N} \land \forall n \in \mathbb{N}\,(N \leq n \to K \cdot C^{-1} < a_{n}) \to n \in \mathbb{N} \to N \leq n \to 0 < a_{n}$
\item $\varphi_{46} :\equiv \forall z\,\forall w\,(z = w \to K < z \to K < w) \to \forall w\,(a_{n} \cdot b_{n} = w \to K < a_{n} \cdot b_{n} \to K < w)$
\item $\varphi_{47} :\equiv N \in \mathbb{N} \land \forall n \in \mathbb{N}\,(N \leq n \to K \cdot C^{-1} < a_{n}) \to n \in \mathbb{N} \to N \leq n \to C^{-1} \in \mathbb{R}$
\item $\varphi_{48} :\equiv \forall x\,\forall y\,(x \in \mathbb{R} \to y \in \mathbb{R} \to x \cdot y \in \mathbb{R}) \to \forall y\,(K \in \mathbb{R} \to y \in \mathbb{R} \to K \cdot y \in \mathbb{R})$
\item $\varphi_{49} :\equiv N \in \mathbb{N} \land \forall n \in \mathbb{N}\,(N \leq n \to K \cdot C^{-1} < a_{n}) \to n \in \mathbb{N} \to N \leq n \to a_{n} \in \mathbb{R}$
\item $\varphi_{50} :\equiv N \in \mathbb{N} \land \forall n \in \mathbb{N}\,(N \leq n \to K \cdot C^{-1} < a_{n}) \to n \in \mathbb{N} \to N \leq n \to b_{n} \in \mathbb{R}$
\item $\varphi_{51} :\equiv N \in \mathbb{N} \land \forall n \in \mathbb{N}\,(N \leq n \to K \cdot C^{-1} < a_{n}) \to n \in \mathbb{N} \to N \leq n \to K \cdot C^{-1} \cdot C = K$
\item $\varphi_{52} :\equiv N \in \mathbb{N} \land \forall n \in \mathbb{N}\,(N \leq n \to K \cdot C^{-1} < a_{n}) \to n \in \mathbb{N} \to N \leq n \to K < a_{n} \cdot C$
\item $\varphi_{53} :\equiv N \in \mathbb{N} \land \forall n \in \mathbb{N}\,(N \leq n \to K \cdot C^{-1} < a_{n}) \to n \in \mathbb{N} \to N \leq n \to K \cdot C^{-1} \in \mathbb{R}$
\item $\varphi_{54} :\equiv \forall y\,(K \in \mathbb{R} \to y \in \mathbb{R} \to K \cdot y \in \mathbb{R}) \to K \in \mathbb{R} \to C^{-1} \in \mathbb{R} \to K \cdot C^{-1} \in \mathbb{R}$
\item $\varphi_{55} :\equiv N \in \mathbb{N} \land \forall n \in \mathbb{N}\,(N \leq n \to K \cdot C^{-1} < a_{n}) \to n \in \mathbb{N} \to N \leq n \to K \cdot C^{-1} < a_{n}$
\item $\varphi_{56} :\equiv b_{n} \in \mathbb{R} \to a_{n} \in \mathbb{R} \to 0 \leq a_{n} \to C \leq b_{n} \to a_{n} \cdot C \leq a_{n} \cdot b_{n}$
\item $\varphi_{57} :\equiv N \in \mathbb{N} \land \forall n \in \mathbb{N}\,(N \leq n \to K \cdot C^{-1} < a_{n}) \to n \in \mathbb{N} \to N \leq n \to K < a \otimes b(n)$
\item $\varphi_{58} :\equiv a_{n} \in \mathbb{R} \to C \in \mathbb{R} \to 0 < C \to K \cdot C^{-1} < a_{n} \to K \cdot C^{-1} \cdot C < a_{n} \cdot C$
\item $\varphi_{59} :\equiv \forall n \in \mathbb{N}\,(C \leq b_{n}) \to \mathrm{Seq}(a \otimes b) \land \forall n \in \mathbb{N}\,(a \otimes b(n) = a_{n} \cdot b_{n})$
\item $\varphi_{60} :\equiv N \in \mathbb{N} \land \forall n \in \mathbb{N}\,(N \leq n \to K \cdot C^{-1} < a_{n}) \to n \in \mathbb{N} \to N \leq n \to K < a_{n} \cdot b_{n}$
\item $\varphi_{61} :\equiv N \in \mathbb{N} \land \forall n \in \mathbb{N}\,(N \leq n \to K \cdot C^{-1} < a_{n}) \to \forall n \in \mathbb{N}\,(N \leq n \to K < a \otimes b(n))$
\item $\varphi_{62} :\equiv \forall z\,(0 < K \cdot C^{-1} \to K \cdot C^{-1} < z \to 0 < z) \to 0 < K \cdot C^{-1} \to K \cdot C^{-1} < a_{n} \to 0 < a_{n}$
\item $\varphi_{63} :\equiv \forall y\,(K \in \mathbb{R} \to y \in \mathbb{R} \to y \neq 0 \to K \cdot y^{-1} \cdot y = K) \to K \in \mathbb{R} \to C \in \mathbb{R} \to C \neq 0 \to K \cdot C^{-1} \cdot C = K$
\item $\varphi_{64} :\equiv \forall n \in \mathbb{N}\,(a \otimes b(n) = a_{n} \cdot b_{n}) \to n \in \mathbb{N} \to a \otimes b(n) = a_{n} \cdot b_{n}$
\item $\varphi_{65} :\equiv \forall y\,\forall z\,(K \cdot C^{-1} \in \mathbb{R} \to y \in \mathbb{R} \to z \in \mathbb{R} \to 0 < z \to K \cdot C^{-1} < y \to K \cdot C^{-1} \cdot z < y \cdot z)$
\item $\varphi_{66} :\equiv N \in \mathbb{N} \land \forall n \in \mathbb{N}\,(N \leq n \to K \cdot C^{-1} < a_{n}) \to n \in \mathbb{N} \to N \leq n \to K \cdot C^{-1} \cdot C < a_{n} \cdot C$
\item $\varphi_{67} :\equiv N \in \mathbb{N} \land \forall n \in \mathbb{N}\,(N \leq n \to K \cdot C^{-1} < a_{n}) \to n \in \mathbb{N} \to N \leq n \to a_{n} \cdot C \leq a_{n} \cdot b_{n}$
\item $\varphi_{68} :\equiv \forall x\,\forall y\,(x \in \mathbb{R} \to y \in \mathbb{R} \to y \neq 0 \to x \cdot y^{-1} \cdot y = x) \to \forall y\,(K \in \mathbb{R} \to y \in \mathbb{R} \to y \neq 0 \to K \cdot y^{-1} \cdot y = K)$
\item $\varphi_{69} :\equiv N \in \mathbb{N} \land \forall n \in \mathbb{N}\,(N \leq n \to K \cdot C^{-1} < a_{n}) \to N \in \mathbb{N} \land \forall n \in \mathbb{N}\,(N \leq n \to K \cdot C^{-1} < a_{n})$
\item $\varphi_{70} :\equiv N \in \mathbb{N} \land \forall n \in \mathbb{N}\,(N \leq n \to K \cdot C^{-1} < a_{n}) \to N \in \mathbb{N} \land \forall n \in \mathbb{N}\,(N \leq n \to K < a \otimes b(n))$
\item $\varphi_{71} :\equiv N \in \mathbb{N} \land \forall n \in \mathbb{N}\,(N \leq n \to K \cdot C^{-1} < a_{n}) \to n \in \mathbb{N} \to N \leq n \to a \otimes b(n) = a_{n} \cdot b_{n}$
\item $\varphi_{72} :\equiv N \in \mathbb{N} \land \forall n \in \mathbb{N}\,(N \leq n \to K \cdot C^{-1} < a_{n}) \to n \in \mathbb{N} \to N \leq n \to a_{n} \cdot b_{n} = a \otimes b(n)$
\item $\varphi_{73} :\equiv N \in \mathbb{N} \land \forall n \in \mathbb{N}\,(N \leq n \to K \cdot C^{-1} < a_{n}) \to n \in \mathbb{N} \to N \leq n \to n \in \mathbb{N} \to N \leq n \to K \cdot C^{-1} < a_{n}$
\item $\varphi_{74} :\equiv N \in \mathbb{N} \land \forall n \in \mathbb{N}\,(N \leq n \to K \cdot C^{-1} < a_{n}) \to \exists N \in \mathbb{N}\,\forall n \in \mathbb{N}\,(N \leq n \to K < a \otimes b(n))$
\item $\varphi_{75} :\equiv N \in \mathbb{N} \land \forall n \in \mathbb{N}\,(N \leq n \to K \cdot C^{-1} < a_{n}) \to n \in \mathbb{N} \to N \leq n \to \forall n \in \mathbb{N}\,(N \leq n \to K \cdot C^{-1} < a_{n})$
\item $\varphi_{76} :\equiv \forall x\,\forall y\,(x \in \mathbb{R} \to y \in \mathbb{R} \to 0 < x \to 0 < y \to 0 < x \cdot y) \to \forall y\,(K \in \mathbb{R} \to y \in \mathbb{R} \to 0 < K \to 0 < y \to 0 < K \cdot y)$
\item $\varphi_{77} :\equiv (\mathrm{DvgP}(a) \to \mathrm{Seq}(b) \to C \in \mathbb{R} \to 0 < C \to 0 < C) \to \mathrm{DvgP}(a) \to \varphi_{4}$
\item $\varphi_{78} :\equiv \forall z\,(K \cdot C^{-1} \in \mathbb{R} \to a_{n} \in \mathbb{R} \to z \in \mathbb{R} \to 0 < z \to K \cdot C^{-1} < a_{n} \to K \cdot C^{-1} \cdot z < a_{n} \cdot z)$
\item $\varphi_{79} :\equiv N \in \mathbb{N} \land \forall n \in \mathbb{N}\,(N \leq n \to K \cdot C^{-1} < a_{n}) \to n \in \mathbb{N} \to N \leq n \to n \in \mathbb{N} \to a \otimes b(n) = a_{n} \cdot b_{n}$
\item $\varphi_{80} :\equiv \forall y\,(a \otimes b(n) = y \to y = a \otimes b(n)) \to a \otimes b(n) = a_{n} \cdot b_{n} \to a_{n} \cdot b_{n} = a \otimes b(n)$
\item $\varphi_{81} :\equiv \forall y\,(K \in \mathbb{R} \to y \in \mathbb{R} \to 0 < K \to 0 < y \to 0 < K \cdot y) \to K \in \mathbb{R} \to C^{-1} \in \mathbb{R} \to 0 < K \to 0 < C^{-1} \to 0 < K \cdot C^{-1}$
\item $\varphi_{82} :\equiv N \in \mathbb{N} \land \forall n \in \mathbb{N}\,(N \leq n \to K \cdot C^{-1} < a_{n}) \to n \in \mathbb{N} \to N \leq n \to \forall n \in \mathbb{N}\,(a \otimes b(n) = a_{n} \cdot b_{n})$
\item $\varphi_{83} :\equiv \forall z\,(K < a_{n} \cdot C \to a_{n} \cdot C \leq z \to K < z) \to K < a_{n} \cdot C \to a_{n} \cdot C \leq a_{n} \cdot b_{n} \to K < a_{n} \cdot b_{n}$
\item $\varphi_{84} :\equiv (\forall n \in \mathbb{N}\,(C \leq b_{n}) \to n \in \mathbb{N} \to C \leq b_{n}) \to \mathrm{DvgP}(a) \to \varphi_{29}$
\item $\varphi_{85} :\equiv C \in \mathbb{R} \to 0 < C \to \forall n \in \mathbb{N}\,(C \leq b_{n}) \to K \in \mathbb{R} \to 0 < K \to \varphi_{14}$
\item $\varphi_{86} :\equiv \varphi_{9} \land K \in \mathbb{R} \land 0 < K \land (N \in \mathbb{N} \land \forall n \in \mathbb{N}\,(N \leq n \to K \cdot C^{-1} < a_{n}))$
\item $\varphi_{87} :\equiv \mathrm{Seq}(b) \to C \in \mathbb{R} \to 0 < C \to \forall n \in \mathbb{N}\,(C \leq b_{n}) \to K \in \mathbb{R} \to \varphi_{20}$
\item $\varphi_{88} :\equiv \mathrm{DvgP}(a) \to \mathrm{Seq}(b) \to C \in \mathbb{R} \to 0 < C \to \forall n \in \mathbb{N}\,(C \leq b_{n}) \to \varphi_{30}$
\item $\varphi_{89} :\equiv \mathrm{DvgP}(a) \to \mathrm{Seq}(b) \to C \in \mathbb{R} \to 0 < C \to \forall n \in \mathbb{N}\,(C \leq b_{n}) \to \varphi_{31}$
\item $\varphi_{90} :\equiv \varphi_{9} \land K \in \mathbb{R} \land 0 < K \land \neg \exists N \in \mathbb{N}\,\forall n \in \mathbb{N}\,(N \leq n \to K < a \otimes b(n))$
\item $\varphi_{91} :\equiv \forall n \in \mathbb{N}\,(C \leq b_{n}) \to K \in \mathbb{R} \to 0 < K \to K \cdot C^{-1} \in \mathbb{R} \to \varphi_{24}$
\item $\varphi_{92} :\equiv (\mathrm{DvgP}(a) \to \mathrm{Seq}(b) \to C \in \mathbb{R} \to 0 < C \to 0 < C) \to \mathrm{DvgP}(a) \to \mathrm{Seq}(b) \to \varphi_{37}$
\item $\varphi_{93} :\equiv C \in \mathbb{R} \to 0 < C \to \forall n \in \mathbb{N}\,(C \leq b_{n}) \to K \in \mathbb{R} \to 0 < K \to \varphi_{27}$
\item $\varphi_{94} :\equiv C \in \mathbb{R} \to 0 < C \to \forall n \in \mathbb{N}\,(C \leq b_{n}) \to K \in \mathbb{R} \to 0 < K \to \varphi_{28}$
\item $\varphi_{95} :\equiv \mathrm{DvgP}(a) \to \mathrm{Seq}(b) \to C \in \mathbb{R} \to 0 < C \to \forall n \in \mathbb{N}\,(C \leq b_{n}) \to \varphi_{38}$
\item $\varphi_{96} :\equiv \varphi_{21} \to K \cdot C^{-1} \cdot C = K \to K \cdot C^{-1} \cdot C < a_{n} \cdot C \to K < a_{n} \cdot C$
\item $\varphi_{97} :\equiv \forall w\,(a_{n} \cdot b_{n} = w \to K < a_{n} \cdot b_{n} \to K < w) \to a_{n} \cdot b_{n} = a \otimes b(n) \to K < a_{n} \cdot b_{n} \to K < a \otimes b(n)$
\item $\varphi_{98} :\equiv C \in \mathbb{R} \to 0 < C \to \forall n \in \mathbb{N}\,(C \leq b_{n}) \to K \in \mathbb{R} \to 0 < K \to \varphi_{39}$
\item $\varphi_{99} :\equiv C \in \mathbb{R} \to 0 < C \to \forall n \in \mathbb{N}\,(C \leq b_{n}) \to K \in \mathbb{R} \to 0 < K \to \varphi_{40}$
\item $\varphi_{100} :\equiv (\mathrm{DvgP}(a) \to \varphi_{4}) \to (C = 0 \to 0 = C) \to \mathrm{DvgP}(a) \to \varphi_{5}$
\item $\varphi_{101} :\equiv C \in \mathbb{R} \to 0 < C \to \forall n \in \mathbb{N}\,(C \leq b_{n}) \to K \in \mathbb{R} \to 0 < K \to \varphi_{45}$
\item $\varphi_{102} :\equiv C \in \mathbb{R} \to 0 < C \to \forall n \in \mathbb{N}\,(C \leq b_{n}) \to K \in \mathbb{R} \to 0 < K \to \varphi_{47}$
\item $\varphi_{103} :\equiv C \in \mathbb{R} \to 0 < C \to \forall n \in \mathbb{N}\,(C \leq b_{n}) \to K \in \mathbb{R} \to 0 < K \to \varphi_{49}$
\item $\varphi_{104} :\equiv C \in \mathbb{R} \to 0 < C \to \forall n \in \mathbb{N}\,(C \leq b_{n}) \to K \in \mathbb{R} \to 0 < K \to \varphi_{50}$
\item $\varphi_{105} :\equiv C \in \mathbb{R} \to 0 < C \to \forall n \in \mathbb{N}\,(C \leq b_{n}) \to K \in \mathbb{R} \to 0 < K \to \varphi_{51}$
\item $\varphi_{106} :\equiv C \in \mathbb{R} \to 0 < C \to \forall n \in \mathbb{N}\,(C \leq b_{n}) \to K \in \mathbb{R} \to 0 < K \to \varphi_{52}$
\item $\varphi_{107} :\equiv C \in \mathbb{R} \to 0 < C \to \forall n \in \mathbb{N}\,(C \leq b_{n}) \to K \in \mathbb{R} \to 0 < K \to \varphi_{53}$
\item $\varphi_{108} :\equiv (\mathrm{DvgP}(a) \to \mathrm{Seq}(b) \to \mathrm{Seq}(b)) \to \mathrm{DvgP}(a) \to \mathrm{Seq}(b) \to C \in \mathbb{R} \to 0 < C \to \varphi_{59}$
\item $\varphi_{109} :\equiv C \in \mathbb{R} \to 0 < C \to \forall n \in \mathbb{N}\,(C \leq b_{n}) \to K \in \mathbb{R} \to 0 < K \to \varphi_{55}$
\item $\varphi_{110} :\equiv C \in \mathbb{R} \to 0 < C \to \forall n \in \mathbb{N}\,(C \leq b_{n}) \to K \in \mathbb{R} \to 0 < K \to \varphi_{57}$
\item $\varphi_{111} :\equiv C \in \mathbb{R} \to 0 < C \to \forall n \in \mathbb{N}\,(C \leq b_{n}) \to K \in \mathbb{R} \to 0 < K \to \varphi_{60}$
\item $\varphi_{112} :\equiv C \in \mathbb{R} \to 0 < C \to \forall n \in \mathbb{N}\,(C \leq b_{n}) \to K \in \mathbb{R} \to 0 < K \to \varphi_{61}$
\item $\varphi_{113} :\equiv (\mathrm{DvgP}(a) \to \varphi_{5}) \to \mathrm{DvgP}(a) \to \mathrm{Seq}(b) \to C \in \mathbb{R} \to \varphi_{13}$
\item $\varphi_{114} :\equiv C \in \mathbb{R} \to 0 < C \to \forall n \in \mathbb{N}\,(C \leq b_{n}) \to K \in \mathbb{R} \to 0 < K \to \varphi_{66}$
\item $\varphi_{115} :\equiv C \in \mathbb{R} \to 0 < C \to \forall n \in \mathbb{N}\,(C \leq b_{n}) \to K \in \mathbb{R} \to 0 < K \to \varphi_{67}$
\item $\varphi_{116} :\equiv C \in \mathbb{R} \to 0 < C \to \forall n \in \mathbb{N}\,(C \leq b_{n}) \to K \in \mathbb{R} \to 0 < K \to \varphi_{69}$
\item $\varphi_{117} :\equiv C \in \mathbb{R} \to 0 < C \to \forall n \in \mathbb{N}\,(C \leq b_{n}) \to K \in \mathbb{R} \to 0 < K \to \varphi_{70}$
\item $\varphi_{118} :\equiv C \in \mathbb{R} \to 0 < C \to \forall n \in \mathbb{N}\,(C \leq b_{n}) \to K \in \mathbb{R} \to 0 < K \to \varphi_{71}$
\item $\varphi_{119} :\equiv C \in \mathbb{R} \to 0 < C \to \forall n \in \mathbb{N}\,(C \leq b_{n}) \to K \in \mathbb{R} \to 0 < K \to \varphi_{72}$
\item $\varphi_{120} :\equiv C \in \mathbb{R} \to 0 < C \to \forall n \in \mathbb{N}\,(C \leq b_{n}) \to K \in \mathbb{R} \to 0 < K \to \varphi_{73}$
\item $\varphi_{121} :\equiv C \in \mathbb{R} \to 0 < C \to \forall n \in \mathbb{N}\,(C \leq b_{n}) \to K \in \mathbb{R} \to 0 < K \to \varphi_{74}$
\item $\varphi_{122} :\equiv (\mathrm{DvgP}(a) \to \mathrm{DvgP}(a)) \to (\mathrm{DvgP}(a) \leftrightarrow \mathrm{Seq}(a) \land \varphi_{14}) \to \mathrm{DvgP}(a) \to \varphi_{7}$
\item $\varphi_{123} :\equiv \forall a\,(\mathrm{DvgP}(a) \leftrightarrow \mathrm{Seq}(a) \land \varphi_{14}) \to (\mathrm{DvgP}(a) \leftrightarrow \mathrm{Seq}(a) \land \varphi_{14})$
\item $\varphi_{124} :\equiv C \in \mathbb{R} \to 0 < C \to \forall n \in \mathbb{N}\,(C \leq b_{n}) \to K \in \mathbb{R} \to 0 < K \to \varphi_{75}$
\item $\varphi_{125} :\equiv C \in \mathbb{R} \to 0 < C \to \forall n \in \mathbb{N}\,(C \leq b_{n}) \to K \in \mathbb{R} \to 0 < K \to \varphi_{79}$
\item $\varphi_{126} :\equiv C \in \mathbb{R} \to 0 < C \to \forall n \in \mathbb{N}\,(C \leq b_{n}) \to K \in \mathbb{R} \to 0 < K \to \varphi_{82}$
\item $\varphi_{127} :\equiv (\mathrm{DvgP}(a) \to \mathrm{Seq}(b) \to \varphi_{37}) \to \mathrm{DvgP}(a) \to \mathrm{Seq}(b) \to C \in \mathbb{R} \to \varphi_{16}$
\item $\varphi_{128} :\equiv \forall N\,(\varphi_{90} \to \neg (N \in \mathbb{N} \land \forall n \in \mathbb{N}\,(N \leq n \to K \cdot C^{-1} < a_{n})))$
\item $\varphi_{129} :\equiv (\mathrm{DvgP}(a) \to \mathrm{Seq}(b) \to C \in \mathbb{R} \to \varphi_{13}) \to \mathrm{DvgP}(a) \to \mathrm{Seq}(b) \to C \in \mathbb{R} \to \varphi_{25}$
\item $\varphi_{130} :\equiv \forall a\,(\mathrm{DvgP}(a) \leftrightarrow \mathrm{Seq}(a) \land \varphi_{14}) \to (\mathrm{DvgP}(a \otimes b) \leftrightarrow \mathrm{Seq}(a \otimes b) \land \varphi_{38})$
\item $\varphi_{131} :\equiv (\mathrm{DvgP}(a) \to \mathrm{Seq}(b) \to C \in \mathbb{R} \to 0 < C \to \varphi_{59}) \to \mathrm{DvgP}(a) \to \varphi_{17}$
\item $\varphi_{132} :\equiv (C \in \mathbb{R} \to 0 < C \to 0 < C^{-1}) \to (\mathrm{DvgP}(a) \to \mathrm{Seq}(b) \to C \in \mathbb{R} \to C \in \mathbb{R}) \to \varphi_{92}$
\item $\varphi_{133} :\equiv (\mathrm{DvgP}(a) \leftrightarrow \mathrm{Seq}(a) \land \varphi_{14}) \to \mathrm{DvgP}(a) \to \mathrm{Seq}(b) \to \varphi_{85}$
\item $\varphi_{134} :\equiv (K \cdot C^{-1} \cdot C = K \to K \cdot C^{-1} \cdot C < a_{n} \cdot C \to K < a_{n} \cdot C) \to \mathrm{DvgP}(a) \to \mathrm{Seq}(b) \to \varphi_{106}$
\item $\varphi_{135} :\equiv (C \in \mathbb{R} \to C \neq 0 \to C^{-1} \in \mathbb{R}) \to (\mathrm{DvgP}(a) \to \mathrm{Seq}(b) \to C \in \mathbb{R} \to C \in \mathbb{R}) \to \varphi_{113}$
\item $\varphi_{136} :\equiv (a_{n} \cdot b_{n} = a \otimes b(n) \to K < a_{n} \cdot b_{n} \to K < a \otimes b(n)) \to \mathrm{DvgP}(a) \to \mathrm{Seq}(b) \to \varphi_{110}$
\item $\varphi_{137} :\equiv (a \otimes b(n) = a_{n} \cdot b_{n} \to a_{n} \cdot b_{n} = a \otimes b(n)) \to \mathrm{DvgP}(a) \to \mathrm{Seq}(b) \to \varphi_{119}$
\item $\varphi_{138} :\equiv (\mathrm{DvgP}(a) \to \mathrm{Seq}(b) \to C \in \mathbb{R} \to C \in \mathbb{R}) \to (\mathrm{DvgP}(a) \to \varphi_{5}) \to \mathrm{DvgP}(a) \to \mathrm{Seq}(b) \to \varphi_{102}$
\item $\varphi_{139} :\equiv (\mathrm{DvgP}(a) \to \mathrm{Seq}(b) \to C \in \mathbb{R} \to C \in \mathbb{R}) \to (\mathrm{DvgP}(a) \to \varphi_{5}) \to \mathrm{DvgP}(a) \to \mathrm{Seq}(b) \to \varphi_{105}$
\item $\varphi_{140} :\equiv (\varphi_{3} \to \neg (N \in \mathbb{N} \land \forall n \in \mathbb{N}\,(N \leq n \to K < a \otimes b(n)))) \to \mathrm{DvgP}(a) \to \mathrm{Seq}(b) \to \varphi_{121}$
\item $\varphi_{141} :\equiv (\mathrm{DvgP}(a) \to \mathrm{Seq}(b) \to \varphi_{93}) \to \mathrm{DvgP}(a) \to \mathrm{Seq}(b) \to \varphi_{98}$
\item $\varphi_{142} :\equiv (\mathrm{DvgP}(a) \to \mathrm{Seq}(b) \to \varphi_{93}) \to \mathrm{DvgP}(a) \to \mathrm{Seq}(b) \to \varphi_{103}$
\item $\varphi_{143} :\equiv (\mathrm{DvgP}(a) \to \mathrm{Seq}(b) \to \varphi_{93}) \to \mathrm{DvgP}(a) \to \mathrm{Seq}(b) \to \varphi_{104}$
\item $\varphi_{144} :\equiv (\mathrm{DvgP}(a) \to \mathrm{Seq}(b) \to \varphi_{101}) \to \mathrm{DvgP}(a) \to \mathrm{Seq}(b) \to \varphi_{99}$
\item $\varphi_{145} :\equiv (\mathrm{DvgP}(a) \to \mathrm{Seq}(b) \to \varphi_{94}) \to \mathrm{DvgP}(a) \to \mathrm{Seq}(b) \to \varphi_{109}$
\item $\varphi_{146} :\equiv (\mathrm{DvgP}(a) \to \mathrm{Seq}(b) \to C \in \mathbb{R} \to 0 < C \to \varphi_{59}) \to \mathrm{DvgP}(a) \to \mathrm{Seq}(b) \to \varphi_{126}$
\item $\varphi_{147} :\equiv (K \in \mathbb{R} \to C^{-1} \in \mathbb{R} \to K \cdot C^{-1} \in \mathbb{R}) \to (\mathrm{DvgP}(a) \to \varphi_{35}) \to \varphi_{129}$
\item $\varphi_{148} :\equiv (\varphi_{14} \to K \cdot C^{-1} \in \mathbb{R} \to \varphi_{24}) \to \mathrm{DvgP}(a) \to \mathrm{Seq}(b) \to C \in \mathbb{R} \to 0 < C \to \varphi_{91}$
\item $\varphi_{149} :\equiv (\mathrm{DvgP}(a) \to \mathrm{Seq}(b) \to \varphi_{102}) \to \mathrm{DvgP}(a) \to \mathrm{Seq}(b) \to \varphi_{107}$
\item $\varphi_{150} :\equiv (\mathrm{DvgP}(a) \to \mathrm{Seq}(b) \to C \in \mathbb{R} \to \varphi_{25}) \to (\mathrm{DvgP}(a) \to \mathrm{Seq}(b) \to C \in \mathbb{R} \to \varphi_{16}) \to \varphi_{88}$
\item $\varphi_{151} :\equiv (\mathrm{DvgP}(a) \to \mathrm{Seq}(b) \to \varphi_{109}) \to \mathrm{DvgP}(a) \to \mathrm{Seq}(b) \to \varphi_{101}$
\item $\varphi_{152} :\equiv (\mathrm{DvgP}(a) \to \mathrm{Seq}(b) \to \varphi_{93}) \to \mathrm{DvgP}(a) \to \mathrm{Seq}(b) \to \varphi_{118}$
\item $\varphi_{153} :\equiv (\mathrm{DvgP}(a) \to \mathrm{Seq}(b) \to \varphi_{98}) \to \mathrm{DvgP}(a) \to \mathrm{Seq}(b) \to \varphi_{115}$
\item $\varphi_{154} :\equiv (\mathrm{DvgP}(a) \to \mathrm{Seq}(b) \to \varphi_{110}) \to \varphi_{86} \to n \in \mathbb{N} \to N \leq n \to K < a \otimes b(n)$
\item $\varphi_{155} :\equiv \forall n\,(\varphi_{86} \to n \in \mathbb{N} \to N \leq n \to K < a \otimes b(n)) \to \forall n\,(\varphi_{86}) \to \forall n \in \mathbb{N}\,(N \leq n \to K < a \otimes b(n))$
\item $\varphi_{156} :\equiv (\mathrm{DvgP}(a) \to \mathrm{Seq}(b) \to \varphi_{109}) \to \mathrm{DvgP}(a) \to \mathrm{Seq}(b) \to \varphi_{114}$
\item $\varphi_{157} :\equiv (\mathrm{DvgP}(a) \to \mathrm{Seq}(b) \to \varphi_{115}) \to \mathrm{DvgP}(a) \to \mathrm{Seq}(b) \to \varphi_{111}$
\item $\varphi_{158} :\equiv (\mathrm{DvgP}(a) \to \mathrm{Seq}(b) \to \varphi_{112}) \to \mathrm{DvgP}(a) \to \mathrm{Seq}(b) \to \varphi_{117}$
\item $\varphi_{159} :\equiv (\mathrm{DvgP}(a) \to \mathrm{Seq}(b) \to \varphi_{116}) \to \mathrm{DvgP}(a) \to \mathrm{Seq}(b) \to \varphi_{124}$
\item $\varphi_{160} :\equiv (\mathrm{DvgP}(a) \to \mathrm{Seq}(b) \to C \in \mathbb{R} \to \varphi_{13}) \to (\mathrm{DvgP}(a) \to \mathrm{Seq}(b) \to \varphi_{23}) \to \varphi_{127}$
\item $\varphi_{161} :\equiv (\mathrm{DvgP}(a) \to \mathrm{Seq}(b) \to \varphi_{121}) \to \varphi_{90} \to \neg (N \in \mathbb{N} \land \forall n \in \mathbb{N}\,(N \leq n \to K \cdot C^{-1} < a_{n}))$
\item $\varphi_{162} :\equiv (\mathrm{Seq}(b) \to n \in \mathbb{N} \to b_{n} \in \mathbb{R}) \to (\mathrm{DvgP}(a) \to \mathrm{Seq}(b) \to \mathrm{Seq}(b)) \to \varphi_{143}$
\item $\varphi_{163} :\equiv (\mathrm{DvgP}(a) \to \mathrm{Seq}(b) \to C \in \mathbb{R} \to 0 < C \to 0 < C) \to \varphi_{156}$
\item $\varphi_{164} :\equiv (K \in \mathbb{R} \to C \in \mathbb{R} \to C \neq 0 \to K \cdot C^{-1} \cdot C = K) \to (\mathrm{DvgP}(a) \to \varphi_{35}) \to \varphi_{139}$
\item $\varphi_{165} :\equiv (\mathrm{DvgP}(a \otimes b) \leftrightarrow \mathrm{Seq}(a \otimes b) \land \varphi_{38}) \to (\mathrm{DvgP}(a) \to \varphi_{17}) \to \varphi_{95} \to \mathrm{DvgP}(a) \to \varphi_{19}$
\item $\varphi_{166} :\equiv (\mathrm{Seq}(a) \to n \in \mathbb{N} \to a_{n} \in \mathbb{R}) \to (\mathrm{DvgP}(a) \to \varphi_{7}) \to \varphi_{142}$
\item $\varphi_{167} :\equiv (\mathrm{DvgP}(a) \to \mathrm{Seq}(b) \to \varphi_{124}) \to \varphi_{41} \to \mathrm{DvgP}(a) \to \mathrm{Seq}(b) \to \varphi_{120}$
\item $\varphi_{168} :\equiv (K \in \mathbb{R} \to C^{-1} \in \mathbb{R} \to K \cdot C^{-1} \in \mathbb{R}) \to (\mathrm{DvgP}(a) \to \varphi_{35}) \to \varphi_{149}$
\item $\varphi_{169} :\equiv (0 < K \cdot C^{-1} \to K \cdot C^{-1} < a_{n} \to 0 < a_{n}) \to (\mathrm{DvgP}(a) \to \mathrm{Seq}(b) \to C \in \mathbb{R} \to \varphi_{16}) \to \varphi_{151}$
\item $\varphi_{170} :\equiv (\mathrm{DvgP}(a) \to \mathrm{Seq}(b) \to \varphi_{126}) \to \varphi_{64} \to \mathrm{DvgP}(a) \to \mathrm{Seq}(b) \to \varphi_{125}$
\item $\varphi_{171} :\equiv (K \in \mathbb{R} \to C^{-1} \in \mathbb{R} \to 0 < K \to 0 < C^{-1} \to 0 < K \cdot C^{-1}) \to (\mathrm{DvgP}(a) \to \varphi_{35}) \to \varphi_{160}$
\item $\varphi_{172} :\equiv (\mathrm{DvgP}(a) \to \mathrm{Seq}(b) \to \varphi_{105}) \to (\mathrm{DvgP}(a) \to \mathrm{Seq}(b) \to \varphi_{114}) \to \varphi_{134}$
\item $\varphi_{173} :\equiv (K < a_{n} \cdot C \to a_{n} \cdot C \leq a_{n} \cdot b_{n} \to K < a_{n} \cdot b_{n}) \to (\mathrm{DvgP}(a) \to \mathrm{Seq}(b) \to \varphi_{106}) \to \varphi_{157}$
\item $\varphi_{174} :\equiv (\mathrm{DvgP}(a) \to \mathrm{Seq}(b) \to \varphi_{103}) \to (\mathrm{DvgP}(a) \to \mathrm{Seq}(b) \to C \in \mathbb{R} \to C \in \mathbb{R}) \to \varphi_{163}$
\item $\varphi_{175} :\equiv (\mathrm{DvgP}(a) \to \mathrm{Seq}(b) \to \varphi_{119}) \to (\mathrm{DvgP}(a) \to \mathrm{Seq}(b) \to \varphi_{111}) \to \varphi_{136}$
\item $\varphi_{176} :\equiv (\mathrm{DvgP}(a) \to \mathrm{Seq}(b) \to \varphi_{120}) \to (\mathrm{DvgP}(a) \to \mathrm{Seq}(b) \to \varphi_{93}) \to \varphi_{145}$
\item $\varphi_{177} :\equiv (\mathrm{DvgP}(a) \to \mathrm{Seq}(b) \to \varphi_{103}) \to (\mathrm{DvgP}(a) \to \mathrm{Seq}(b) \to \varphi_{99}) \to \varphi_{153}$
\item $\varphi_{178} :\equiv \forall K\,(\varphi_{9} \to \varphi_{31}) \to (\forall K\,(\varphi_{9} \to \varphi_{31}) \to \forall K\,(\varphi_{9}) \to \varphi_{38}) \to (\varphi_{9} \to \forall K\,(\varphi_{9})) \to \varphi_{95}$
\item $\varphi_{179} :\equiv (\mathrm{DvgP}(a) \to \mathrm{Seq}(b) \to C \in \mathbb{R} \to C \in \mathbb{R}) \to (\mathrm{DvgP}(a) \to \mathrm{Seq}(b) \to \varphi_{104}) \to \varphi_{177}$
\item $\varphi_{180} :\equiv (K \cdot C^{-1} \in \mathbb{R} \to \varphi_{58}) \to (\mathrm{DvgP}(a) \to \mathrm{Seq}(b) \to \varphi_{107}) \to \varphi_{174}$
\item $\varphi_{181} :\equiv \forall n\,(\varphi_{86} \to n \in \mathbb{N} \to N \leq n \to K < a \otimes b(n)) \to \varphi_{155} \to (\varphi_{86} \to \forall n\,(\varphi_{86})) \to \mathrm{DvgP}(a) \to \mathrm{Seq}(b) \to \varphi_{112}$
\end{abbrevlist}
\end{fullproof}

\section{第3章 3節: 累乗と等比数列の極限 (命題3.1)}
% 自動生成 (prover/build.js): MulZeroCancel  ノード数 181
\begin{fullproof}{[E48]}{$\forall x\,\forall y\,(x \in \mathbb{R} \to y \in \mathbb{R} \to y \neq 0 \to x \cdot y = 0 \to x = 0)$}{181}\phantomsection\label{tree:MulZeroCancel}
\noindent$\Pi_{1}$:
\begin{ptree}
\AXC{$y = z \to x = y \to x = z$\lbl{A6}}
\AXC{$(y = z \to x = y \to x = z) \to x = y \to y = z \to x = z$\lbl{A1}}
\BIC{$x = y \to y = z \to x = z$}
\UIC{$\forall z\,(x = y \to y = z \to x = z)$}
\UIC{$\forall y\,\forall z\,(x = y \to y = z \to x = z)$}
\UIC{$\forall x\,\forall y\,\forall z\,(x = y \to y = z \to x = z)$}
\AXC{$\varphi_{3}$\lbl{A2}}
\BIC{$\forall y\,\forall z\,(x = y \to y = z \to x = z)$}
\AXC{$\varphi_{7}$\lbl{A2}}
\BIC{$\forall z\,(x = y^{-1} \cdot 0 \to y^{-1} \cdot 0 = z \to x = z)$}
\AXC{$\varphi_{25}$\lbl{A2}}
\BIC{$x = y^{-1} \cdot 0 \to y^{-1} \cdot 0 = 0 \to x = 0$}
\end{ptree}
\noindent$\Pi_{2}$:
\begin{ptree}
\AXC{$x \in \mathbb{R} \to y \in \mathbb{R} \to y \neq 0 \to x \cdot y = 0 \to y \in \mathbb{R}$\lbl{A1}}
\AXC{$\forall x \in \mathbb{R}\,(x \neq 0 \to x^{-1} \in \mathbb{R})$\lbl{(R1)}}
\AXC{$\varphi_{26}$\lbl{A2}}
\BIC{$y \in \mathbb{R} \to y \neq 0 \to y^{-1} \in \mathbb{R}$}
\AXC{$\varphi_{53}$\lbl{A1}}
\BIC{$(x \in \mathbb{R} \to y \in \mathbb{R} \to y \neq 0 \to x \cdot y = 0 \to y \in \mathbb{R}) \to \varphi_{40}$}
\BIC{$\varphi_{40}$}
\end{ptree}
\noindent$\Pi_{3}$:
\begin{ptree}
\AXC{$x \in \mathbb{R} \to y \in \mathbb{R} \to y \neq 0 \to x \cdot y = 0 \to y \neq 0$\lbl{A1}}
\AXC{$\stackrel{\Pi_{2}}{\vdots}$}\noLine
\UIC{$\varphi_{40}$}
\BIC{$x \in \mathbb{R} \to y \in \mathbb{R} \to y \neq 0 \to x \cdot y = 0 \to y^{-1} \in \mathbb{R}$}
\end{ptree}
\noindent$\Pi_{4}$:
\begin{ptree}
\AXC{$\forall x \in \mathbb{R}\,(x \cdot 0 = 0)$\lbl{[F8]}}
\AXC{$\forall x \in \mathbb{R}\,(x \cdot 0 = 0) \to y^{-1} \in \mathbb{R} \to y^{-1} \cdot 0 = 0$\lbl{A2}}
\BIC{$y^{-1} \in \mathbb{R} \to y^{-1} \cdot 0 = 0$}
\AXC{$(y^{-1} \in \mathbb{R} \to y^{-1} \cdot 0 = 0) \to \varphi_{42}$\lbl{A1}}
\BIC{$\varphi_{42}$}
\end{ptree}
\noindent$\Pi_{5}$:
\begin{ptree}
\AXC{$y = z \to x = y \to x = z$\lbl{A6}}
\AXC{$(y = z \to x = y \to x = z) \to x = y \to y = z \to x = z$\lbl{A1}}
\BIC{$x = y \to y = z \to x = z$}
\UIC{$\forall z\,(x = y \to y = z \to x = z)$}
\UIC{$\forall y\,\forall z\,(x = y \to y = z \to x = z)$}
\UIC{$\forall x\,\forall y\,\forall z\,(x = y \to y = z \to x = z)$}
\AXC{$\varphi_{3}$\lbl{A2}}
\BIC{$\forall y\,\forall z\,(x = y \to y = z \to x = z)$}
\AXC{$\varphi_{8}$\lbl{A2}}
\BIC{$\forall z\,(x = 0 \cdot y^{-1} \to 0 \cdot y^{-1} = z \to x = z)$}
\AXC{$\varphi_{35}$\lbl{A2}}
\BIC{$x = 0 \cdot y^{-1} \to 0 \cdot y^{-1} = y^{-1} \cdot 0 \to x = y^{-1} \cdot 0$}
\end{ptree}
\noindent$\Pi_{6}$:
\begin{ptree}
\AXC{$\forall x\,\forall y\,(x \in \mathbb{R} \to y \in \mathbb{R} \to x \cdot y = y \cdot x)$\lbl{(R7)}}
\AXC{$\varphi_{34}$\lbl{A2}}
\BIC{$\forall y\,(0 \in \mathbb{R} \to y \in \mathbb{R} \to 0 \cdot y = y \cdot 0)$}
\AXC{$\varphi_{37}$\lbl{A2}}
\BIC{$0 \in \mathbb{R} \to y^{-1} \in \mathbb{R} \to 0 \cdot y^{-1} = y^{-1} \cdot 0$}
\AXC{$(0 \in \mathbb{R} \to y^{-1} \in \mathbb{R} \to 0 \cdot y^{-1} = y^{-1} \cdot 0) \to \varphi_{47}$\lbl{A1}}
\BIC{$\varphi_{47}$}
\end{ptree}
\noindent$\Pi_{7}$:
\begin{ptree}
\AXC{$y = z \to x = y \to x = z$\lbl{A6}}
\AXC{$(y = z \to x = y \to x = z) \to x = y \to y = z \to x = z$\lbl{A1}}
\BIC{$x = y \to y = z \to x = z$}
\UIC{$\forall z\,(x = y \to y = z \to x = z)$}
\UIC{$\forall y\,\forall z\,(x = y \to y = z \to x = z)$}
\UIC{$\forall x\,\forall y\,\forall z\,(x = y \to y = z \to x = z)$}
\AXC{$\varphi_{3}$\lbl{A2}}
\BIC{$\forall y\,\forall z\,(x = y \to y = z \to x = z)$}
\AXC{$\varphi_{18}$\lbl{A2}}
\BIC{$\forall z\,(x = x \cdot y \cdot y^{-1} \to x \cdot y \cdot y^{-1} = z \to x = z)$}
\end{ptree}
\noindent$\Pi_{8}$:
\begin{ptree}
\AXC{$\stackrel{\Pi_{7}}{\vdots}$}\noLine
\UIC{$\forall z\,(x = x \cdot y \cdot y^{-1} \to x \cdot y \cdot y^{-1} = z \to x = z)$}
\AXC{$\forall z\,(x = x \cdot y \cdot y^{-1} \to x \cdot y \cdot y^{-1} = z \to x = z) \to \varphi_{12}$\lbl{A2}}
\BIC{$\varphi_{12}$}
\end{ptree}
\noindent$\Pi_{9}$:
\begin{ptree}
\AXC{$x = x$\lbl{A5}}
\UIC{$\forall x\,(x = x)$}
\AXC{$\forall x\,(x = x) \to x \cdot y \cdot y^{-1} = x \cdot y \cdot y^{-1}$\lbl{A2}}
\BIC{$x \cdot y \cdot y^{-1} = x \cdot y \cdot y^{-1}$}
\AXC{$x \in \mathbb{R} \to y \in \mathbb{R} \to y \neq 0 \to x \cdot y = 0 \to x \cdot y = 0$\lbl{A1}}
\AXC{$\varphi_{55}$\lbl{A1}}
\BIC{$x \cdot y \cdot y^{-1} = x \cdot y \cdot y^{-1} \to \varphi_{23} \to \varphi_{19}$}
\BIC{$\varphi_{23} \to \varphi_{19}$}
\end{ptree}
\noindent$\Pi_{10}$:
\begin{ptree}
\AXC{$y = z \to x = y \to x = z$\lbl{A6}}
\AXC{$(y = z \to x = y \to x = z) \to x = y \to y = z \to x = z$\lbl{A1}}
\BIC{$x = y \to y = z \to x = z$}
\UIC{$\forall z\,(x = y \to y = z \to x = z)$}
\UIC{$\forall y\,\forall z\,(x = y \to y = z \to x = z)$}
\UIC{$\forall x\,\forall y\,\forall z\,(x = y \to y = z \to x = z)$}
\AXC{$\varphi_{3}$\lbl{A2}}
\BIC{$\forall y\,\forall z\,(x = y \to y = z \to x = z)$}
\end{ptree}
\noindent$\Pi_{11}$:
\begin{ptree}
\AXC{$\stackrel{\Pi_{10}}{\vdots}$}\noLine
\UIC{$\forall y\,\forall z\,(x = y \to y = z \to x = z)$}
\AXC{$\forall y\,\forall z\,(x = y \to y = z \to x = z) \to \varphi_{4}$\lbl{A2}}
\BIC{$\varphi_{4}$}
\AXC{$\varphi_{4} \to \varphi_{24}$\lbl{A2}}
\BIC{$\varphi_{24}$}
\end{ptree}
\noindent$\Pi_{12}$:
\begin{ptree}
\AXC{$x = x$\lbl{A5}}
\AXC{$x = y \to x = x \to y = x$\lbl{A6}}
\AXC{$(x = y \to x = x \to y = x) \to x = x \to x = y \to y = x$\lbl{A1}}
\BIC{$x = x \to x = y \to y = x$}
\BIC{$x = y \to y = x$}
\UIC{$\forall y\,(x = y \to y = x)$}
\UIC{$\forall x\,\forall y\,(x = y \to y = x)$}
\end{ptree}
\noindent$\Pi_{13}$:
\begin{ptree}
\AXC{$\stackrel{\Pi_{12}}{\vdots}$}\noLine
\UIC{$\forall x\,\forall y\,(x = y \to y = x)$}
\AXC{$\forall x\,\forall y\,(x = y \to y = x) \to \forall y\,(v_{851} = y \to y = v_{851})$\lbl{A2}}
\BIC{$\forall y\,(v_{851} = y \to y = v_{851})$}
\AXC{$\varphi_{1}$\lbl{A2}}
\BIC{$v_{851} = v_{852} \to v_{852} = v_{851}$}
\UIC{$\forall v_{852}\,(v_{851} = v_{852} \to v_{852} = v_{851})$}
\UIC{$\forall v_{851}\,\forall v_{852}\,(v_{851} = v_{852} \to v_{852} = v_{851})$}
\AXC{$\varphi_{32}$\lbl{A2}}
\BIC{$\forall v_{852}\,(x \cdot y \cdot y^{-1} = v_{852} \to v_{852} = x \cdot y \cdot y^{-1})$}
\AXC{$\varphi_{43}$\lbl{A2}}
\BIC{$\varphi_{17}$}
\end{ptree}
\noindent$\Pi_{14}$:
\begin{ptree}
\AXC{$\forall x\,\forall y\,(\varphi_{14})$\lbl{(R6)}}
\AXC{$\forall x\,\forall y\,(\varphi_{14}) \to \forall y\,(\varphi_{14})$\lbl{A2}}
\BIC{$\forall y\,(\varphi_{14})$}
\AXC{$\forall y\,(\varphi_{14}) \to \varphi_{14}$\lbl{A2}}
\BIC{$\varphi_{14}$}
\AXC{$\varphi_{14} \to \varphi_{27}$\lbl{A2}}
\BIC{$\varphi_{27}$}
\AXC{$\varphi_{27} \to \varphi_{58}$\lbl{A1}}
\BIC{$\varphi_{58}$}
\end{ptree}
\noindent$\Pi_{15}$:
\begin{ptree}
\AXC{$x \in \mathbb{R} \to y \in \mathbb{R} \to y \neq 0 \to x \cdot y = 0 \to y \in \mathbb{R}$\lbl{A1}}
\AXC{$x \in \mathbb{R} \to y \in \mathbb{R} \to y \neq 0 \to x \cdot y = 0 \to x \in \mathbb{R}$\lbl{A1}}
\AXC{$\stackrel{\Pi_{14}}{\vdots}$}\noLine
\UIC{$\varphi_{58}$}
\BIC{$(x \in \mathbb{R} \to y \in \mathbb{R} \to y \neq 0 \to x \cdot y = 0 \to y \in \mathbb{R}) \to \varphi_{45}$}
\BIC{$\varphi_{45}$}
\end{ptree}
\noindent$\Pi_{16}$:
\begin{ptree}
\AXC{$\stackrel{\Pi_{13}}{\vdots}$}\noLine
\UIC{$\varphi_{17}$}
\AXC{$\stackrel{\Pi_{3}}{\vdots}$}\noLine
\UIC{$x \in \mathbb{R} \to y \in \mathbb{R} \to y \neq 0 \to x \cdot y = 0 \to y^{-1} \in \mathbb{R}$}
\AXC{$\stackrel{\Pi_{15}}{\vdots}$}\noLine
\UIC{$\varphi_{45}$}
\BIC{$x \in \mathbb{R} \to \varphi_{10}$}
\AXC{$(x \in \mathbb{R} \to \varphi_{10}) \to \varphi_{17} \to x \in \mathbb{R} \to \varphi_{11}$\lbl{A1}}
\BIC{$\varphi_{17} \to x \in \mathbb{R} \to \varphi_{11}$}
\BIC{$x \in \mathbb{R} \to \varphi_{11}$}
\end{ptree}
\noindent$\Pi_{17}$:
\begin{ptree}
\AXC{$y = z \to x = y \to x = z$\lbl{A6}}
\AXC{$(y = z \to x = y \to x = z) \to x = y \to y = z \to x = z$\lbl{A1}}
\BIC{$x = y \to y = z \to x = z$}
\UIC{$\forall z\,(x = y \to y = z \to x = z)$}
\UIC{$\forall y\,\forall z\,(x = y \to y = z \to x = z)$}
\UIC{$\forall x\,\forall y\,\forall z\,(x = y \to y = z \to x = z)$}
\AXC{$\varphi_{3}$\lbl{A2}}
\BIC{$\forall y\,\forall z\,(x = y \to y = z \to x = z)$}
\AXC{$\varphi_{5}$\lbl{A2}}
\BIC{$\forall z\,(x = x \cdot 1 \to x \cdot 1 = z \to x = z)$}
\end{ptree}
\noindent$\Pi_{18}$:
\begin{ptree}
\AXC{$\stackrel{\Pi_{12}}{\vdots}$}\noLine
\UIC{$\forall x\,\forall y\,(x = y \to y = x)$}
\AXC{$\forall x\,\forall y\,(x = y \to y = x) \to \forall y\,(v_{849} = y \to y = v_{849})$\lbl{A2}}
\BIC{$\forall y\,(v_{849} = y \to y = v_{849})$}
\AXC{$\varphi_{2}$\lbl{A2}}
\BIC{$v_{849} = v_{850} \to v_{850} = v_{849}$}
\UIC{$\forall v_{850}\,(v_{849} = v_{850} \to v_{850} = v_{849})$}
\UIC{$\forall v_{849}\,\forall v_{850}\,(v_{849} = v_{850} \to v_{850} = v_{849})$}
\AXC{$\varphi_{28}$\lbl{A2}}
\BIC{$\forall v_{850}\,(y \cdot y^{-1} = v_{850} \to v_{850} = y \cdot y^{-1})$}
\AXC{$\varphi_{21}$\lbl{A2}}
\BIC{$y \cdot y^{-1} = 1 \to 1 = y \cdot y^{-1}$}
\end{ptree}
\noindent$\Pi_{19}$:
\begin{ptree}
\AXC{$x \in \mathbb{R} \to y \in \mathbb{R} \to y \neq 0 \to x \cdot y = 0 \to y \in \mathbb{R}$\lbl{A1}}
\AXC{$\forall x \in \mathbb{R}\,(x \neq 0 \to x \cdot x^{-1} = 1)$\lbl{(R9)}}
\AXC{$\varphi_{20}$\lbl{A2}}
\BIC{$y \in \mathbb{R} \to y \neq 0 \to y \cdot y^{-1} = 1$}
\AXC{$\varphi_{51}$\lbl{A1}}
\BIC{$(x \in \mathbb{R} \to y \in \mathbb{R} \to y \neq 0 \to x \cdot y = 0 \to y \in \mathbb{R}) \to \varphi_{38}$}
\BIC{$\varphi_{38}$}
\end{ptree}
\noindent$\Pi_{20}$:
\begin{ptree}
\AXC{$\stackrel{\Pi_{18}}{\vdots}$}\noLine
\UIC{$y \cdot y^{-1} = 1 \to 1 = y \cdot y^{-1}$}
\AXC{$x \in \mathbb{R} \to y \in \mathbb{R} \to y \neq 0 \to x \cdot y = 0 \to y \neq 0$\lbl{A1}}
\AXC{$\stackrel{\Pi_{19}}{\vdots}$}\noLine
\UIC{$\varphi_{38}$}
\BIC{$x \in \mathbb{R} \to y \in \mathbb{R} \to y \neq 0 \to x \cdot y = 0 \to y \cdot y^{-1} = 1$}
\AXC{$\varphi_{49}$\lbl{A1}}
\BIC{$\varphi_{33}$}
\BIC{$x \in \mathbb{R} \to y \in \mathbb{R} \to y \neq 0 \to x \cdot y = 0 \to 1 = y \cdot y^{-1}$}
\AXC{$\varphi_{52}$\lbl{A1}}
\BIC{$x \cdot 1 = x \cdot 1 \to \varphi_{44}$}
\end{ptree}
\noindent$\Pi_{21}$:
\begin{ptree}
\AXC{$x = x$\lbl{A5}}
\UIC{$\forall x\,(x = x)$}
\AXC{$\forall x\,(x = x) \to x \cdot 1 = x \cdot 1$\lbl{A2}}
\BIC{$x \cdot 1 = x \cdot 1$}
\AXC{$\stackrel{\Pi_{20}}{\vdots}$}\noLine
\UIC{$x \cdot 1 = x \cdot 1 \to \varphi_{44}$}
\BIC{$\varphi_{44}$}
\end{ptree}
\noindent$\Pi_{22}$:
\begin{ptree}
\AXC{$\stackrel{\Pi_{12}}{\vdots}$}\noLine
\UIC{$\forall x\,\forall y\,(x = y \to y = x)$}
\AXC{$\forall x\,\forall y\,(x = y \to y = x) \to \forall y\,(x \cdot 1 = y \to y = x \cdot 1)$\lbl{A2}}
\BIC{$\forall y\,(x \cdot 1 = y \to y = x \cdot 1)$}
\end{ptree}
\noindent$\Pi_{23}$:
\begin{ptree}
\AXC{$\forall x \in \mathbb{R}\,(x \cdot 1 = x)$\lbl{(R8)}}
\AXC{$\forall x \in \mathbb{R}\,(x \cdot 1 = x) \to x \in \mathbb{R} \to x \cdot 1 = x$\lbl{A2}}
\BIC{$x \in \mathbb{R} \to x \cdot 1 = x$}
\AXC{$(x \in \mathbb{R} \to x \cdot 1 = x) \to \varphi_{39}$\lbl{A1}}
\BIC{$\varphi_{39}$}
\end{ptree}
\noindent$\Pi_{24}$:
\begin{ptree}
\AXC{$x \in \mathbb{R} \to y \in \mathbb{R} \to y \neq 0 \to x \cdot y = 0 \to x \in \mathbb{R}$\lbl{A1}}
\AXC{$\stackrel{\Pi_{23}}{\vdots}$}\noLine
\UIC{$\varphi_{39}$}
\BIC{$x \in \mathbb{R} \to y \in \mathbb{R} \to y \neq 0 \to x \cdot y = 0 \to x \cdot 1 = x$}
\AXC{$(x \in \mathbb{R} \to y \in \mathbb{R} \to y \neq 0 \to x \cdot y = 0 \to x \cdot 1 = x) \to \varphi_{22}$\lbl{A1}}
\BIC{$\varphi_{22}$}
\end{ptree}
\noindent$\Pi_{25}$:
\begin{ptree}
\AXC{$\stackrel{\Pi_{22}}{\vdots}$}\noLine
\UIC{$\forall y\,(x \cdot 1 = y \to y = x \cdot 1)$}
\AXC{$\forall y\,(x \cdot 1 = y \to y = x \cdot 1) \to x \cdot 1 = x \to x = x \cdot 1$\lbl{A2}}
\BIC{$x \cdot 1 = x \to x = x \cdot 1$}
\AXC{$\stackrel{\Pi_{24}}{\vdots}$}\noLine
\UIC{$\varphi_{22}$}
\BIC{$x \in \mathbb{R} \to y \in \mathbb{R} \to y \neq 0 \to x \cdot y = 0 \to x = x \cdot 1$}
\AXC{$\varphi_{56}$\lbl{A1}}
\BIC{$\varphi_{16} \to \varphi_{6} \to \varphi_{9}$}
\end{ptree}
\noindent$\Pi_{26}$:
\begin{ptree}
\AXC{$z = w \to x \cdot 1 = x \cdot z \to x \cdot 1 = x \cdot w$\lbl{A6}}
\UIC{$\forall w\,(z = w \to x \cdot 1 = x \cdot z \to x \cdot 1 = x \cdot w)$}
\UIC{$\forall z\,\forall w\,(z = w \to x \cdot 1 = x \cdot z \to x \cdot 1 = x \cdot w)$}
\AXC{$\varphi_{30}$\lbl{A2}}
\BIC{$\forall w\,(1 = w \to x \cdot 1 = x \cdot 1 \to x \cdot 1 = x \cdot w)$}
\AXC{$\varphi_{36}$\lbl{A2}}
\BIC{$1 = y \cdot y^{-1} \to x \cdot 1 = x \cdot 1 \to x \cdot 1 = x \cdot (y \cdot y^{-1})$}
\AXC{$\stackrel{\Pi_{21}}{\vdots}$}\noLine
\UIC{$\varphi_{44}$}
\BIC{$\varphi_{16}$}
\AXC{$\stackrel{\Pi_{25}}{\vdots}$}\noLine
\UIC{$\varphi_{16} \to \varphi_{6} \to \varphi_{9}$}
\BIC{$\varphi_{6} \to \varphi_{9}$}
\end{ptree}
\noindent$\Pi_{27}$:
\begin{ptree}
\AXC{$\stackrel{\Pi_{17}}{\vdots}$}\noLine
\UIC{$\forall z\,(x = x \cdot 1 \to x \cdot 1 = z \to x = z)$}
\AXC{$\forall z\,(x = x \cdot 1 \to x \cdot 1 = z \to x = z) \to \varphi_{6}$\lbl{A2}}
\BIC{$\varphi_{6}$}
\AXC{$\stackrel{\Pi_{26}}{\vdots}$}\noLine
\UIC{$\varphi_{6} \to \varphi_{9}$}
\BIC{$\varphi_{9}$}
\end{ptree}
\noindent$\Pi_{28}$:
\begin{ptree}
\AXC{$\stackrel{\Pi_{11}}{\vdots}$}\noLine
\UIC{$\varphi_{24}$}
\AXC{$\stackrel{\Pi_{16}}{\vdots}$}\noLine
\UIC{$x \in \mathbb{R} \to \varphi_{11}$}
\AXC{$\stackrel{\Pi_{27}}{\vdots}$}\noLine
\UIC{$\varphi_{9}$}
\AXC{$\varphi_{9} \to (x \in \mathbb{R} \to \varphi_{11}) \to \varphi_{46}$\lbl{A1}}
\BIC{$(x \in \mathbb{R} \to \varphi_{11}) \to \varphi_{46}$}
\BIC{$\varphi_{46}$}
\BIC{$x \in \mathbb{R} \to y \in \mathbb{R} \to y \neq 0 \to x \cdot y = 0 \to x = x \cdot y \cdot y^{-1}$}
\AXC{$\varphi_{59}$\lbl{A1}}
\BIC{$\varphi_{50}$}
\end{ptree}
\noindent$\Pi_{29}$:
\begin{ptree}
\AXC{$\stackrel{\Pi_{8}}{\vdots}$}\noLine
\UIC{$\varphi_{12}$}
\AXC{$\varphi_{13}$\lbl{A6}}
\UIC{$\forall w\,(\varphi_{13})$}
\UIC{$\forall z\,\forall w\,(\varphi_{13})$}
\AXC{$\forall z\,\forall w\,(\varphi_{13}) \to \varphi_{29}$\lbl{A2}}
\BIC{$\varphi_{29}$}
\AXC{$\varphi_{29} \to \varphi_{23}$\lbl{A2}}
\BIC{$\varphi_{23}$}
\AXC{$\stackrel{\Pi_{9}}{\vdots}$}\noLine
\UIC{$\varphi_{23} \to \varphi_{19}$}
\BIC{$\varphi_{19}$}
\AXC{$\stackrel{\Pi_{28}}{\vdots}$}\noLine
\UIC{$\varphi_{50}$}
\BIC{$\varphi_{12} \to x \in \mathbb{R} \to y \in \mathbb{R} \to y \neq 0 \to x \cdot y = 0 \to x = 0 \cdot y^{-1}$}
\BIC{$x \in \mathbb{R} \to y \in \mathbb{R} \to y \neq 0 \to x \cdot y = 0 \to x = 0 \cdot y^{-1}$}
\AXC{$\varphi_{57}$\lbl{A1}}
\BIC{$\varphi_{15} \to \varphi_{41}$}
\end{ptree}
\noindent$\Pi_{30}$:
\begin{ptree}
\AXC{$\stackrel{\Pi_{3}}{\vdots}$}\noLine
\UIC{$x \in \mathbb{R} \to y \in \mathbb{R} \to y \neq 0 \to x \cdot y = 0 \to y^{-1} \in \mathbb{R}$}
\AXC{$0 \in \mathbb{R}$\lbl{(R1)}}
\AXC{$\stackrel{\Pi_{6}}{\vdots}$}\noLine
\UIC{$\varphi_{47}$}
\BIC{$(x \in \mathbb{R} \to y \in \mathbb{R} \to y \neq 0 \to x \cdot y = 0 \to y^{-1} \in \mathbb{R}) \to \varphi_{15}$}
\BIC{$\varphi_{15}$}
\AXC{$\stackrel{\Pi_{29}}{\vdots}$}\noLine
\UIC{$\varphi_{15} \to \varphi_{41}$}
\BIC{$\varphi_{41}$}
\end{ptree}
\noindent$\Pi_{31}$:
\begin{ptree}
\AXC{$\stackrel{\Pi_{5}}{\vdots}$}\noLine
\UIC{$x = 0 \cdot y^{-1} \to 0 \cdot y^{-1} = y^{-1} \cdot 0 \to x = y^{-1} \cdot 0$}
\AXC{$\stackrel{\Pi_{30}}{\vdots}$}\noLine
\UIC{$\varphi_{41}$}
\BIC{$x \in \mathbb{R} \to y \in \mathbb{R} \to y \neq 0 \to x \cdot y = 0 \to x = y^{-1} \cdot 0$}
\AXC{$\varphi_{54}$\lbl{A1}}
\BIC{$\varphi_{48}$}
\end{ptree}
\noindent 主木:
\begin{ptree}
\AXC{$\stackrel{\Pi_{1}}{\vdots}$}\noLine
\UIC{$x = y^{-1} \cdot 0 \to y^{-1} \cdot 0 = 0 \to x = 0$}
\AXC{$\stackrel{\Pi_{3}}{\vdots}$}\noLine
\UIC{$x \in \mathbb{R} \to y \in \mathbb{R} \to y \neq 0 \to x \cdot y = 0 \to y^{-1} \in \mathbb{R}$}
\AXC{$\stackrel{\Pi_{4}}{\vdots}$}\noLine
\UIC{$\varphi_{42}$}
\BIC{$x \in \mathbb{R} \to y \in \mathbb{R} \to y \neq 0 \to x \cdot y = 0 \to y^{-1} \cdot 0 = 0$}
\AXC{$\stackrel{\Pi_{31}}{\vdots}$}\noLine
\UIC{$\varphi_{48}$}
\BIC{$\varphi_{31}$}
\BIC{$x \in \mathbb{R} \to y \in \mathbb{R} \to y \neq 0 \to x \cdot y = 0 \to x = 0$}
\UIC{$\forall y\,(x \in \mathbb{R} \to y \in \mathbb{R} \to y \neq 0 \to x \cdot y = 0 \to x = 0)$}
\UIC{$\forall x\,\forall y\,(x \in \mathbb{R} \to y \in \mathbb{R} \to y \neq 0 \to x \cdot y = 0 \to x = 0)$}
\end{ptree}
\begin{abbrevlist}
\item $\varphi_{1} :\equiv \forall y\,(v_{851} = y \to y = v_{851}) \to v_{851} = v_{852} \to v_{852} = v_{851}$
\item $\varphi_{2} :\equiv \forall y\,(v_{849} = y \to y = v_{849}) \to v_{849} = v_{850} \to v_{850} = v_{849}$
\item $\varphi_{3} :\equiv \forall x\,\forall y\,\forall z\,(x = y \to y = z \to x = z) \to \forall y\,\forall z\,(x = y \to y = z \to x = z)$
\item $\varphi_{4} :\equiv \forall z\,(x = x \cdot (y \cdot y^{-1}) \to x \cdot (y \cdot y^{-1}) = z \to x = z)$
\item $\varphi_{5} :\equiv \forall y\,\forall z\,(x = y \to y = z \to x = z) \to \forall z\,(x = x \cdot 1 \to x \cdot 1 = z \to x = z)$
\item $\varphi_{6} :\equiv x = x \cdot 1 \to x \cdot 1 = x \cdot (y \cdot y^{-1}) \to x = x \cdot (y \cdot y^{-1})$
\item $\varphi_{7} :\equiv \forall y\,\forall z\,(x = y \to y = z \to x = z) \to \forall z\,(x = y^{-1} \cdot 0 \to y^{-1} \cdot 0 = z \to x = z)$
\item $\varphi_{8} :\equiv \forall y\,\forall z\,(x = y \to y = z \to x = z) \to \forall z\,(x = 0 \cdot y^{-1} \to 0 \cdot y^{-1} = z \to x = z)$
\item $\varphi_{9} :\equiv x \in \mathbb{R} \to y \in \mathbb{R} \to y \neq 0 \to x \cdot y = 0 \to x = x \cdot (y \cdot y^{-1})$
\item $\varphi_{10} :\equiv y \in \mathbb{R} \to y \neq 0 \to x \cdot y = 0 \to x \cdot y \cdot y^{-1} = x \cdot (y \cdot y^{-1})$
\item $\varphi_{11} :\equiv y \in \mathbb{R} \to y \neq 0 \to x \cdot y = 0 \to x \cdot (y \cdot y^{-1}) = x \cdot y \cdot y^{-1}$
\item $\varphi_{12} :\equiv x = x \cdot y \cdot y^{-1} \to x \cdot y \cdot y^{-1} = 0 \cdot y^{-1} \to x = 0 \cdot y^{-1}$
\item $\varphi_{13} :\equiv z = w \to x \cdot y \cdot y^{-1} = z \cdot y^{-1} \to x \cdot y \cdot y^{-1} = w \cdot y^{-1}$
\item $\varphi_{14} :\equiv \forall z\,(x \in \mathbb{R} \to y \in \mathbb{R} \to z \in \mathbb{R} \to x \cdot y \cdot z = x \cdot (y \cdot z))$
\item $\varphi_{15} :\equiv x \in \mathbb{R} \to y \in \mathbb{R} \to y \neq 0 \to x \cdot y = 0 \to 0 \cdot y^{-1} = y^{-1} \cdot 0$
\item $\varphi_{16} :\equiv x \in \mathbb{R} \to y \in \mathbb{R} \to y \neq 0 \to x \cdot y = 0 \to x \cdot 1 = x \cdot (y \cdot y^{-1})$
\item $\varphi_{17} :\equiv x \cdot y \cdot y^{-1} = x \cdot (y \cdot y^{-1}) \to x \cdot (y \cdot y^{-1}) = x \cdot y \cdot y^{-1}$
\item $\varphi_{18} :\equiv \forall y\,\forall z\,(x = y \to y = z \to x = z) \to \forall z\,(x = x \cdot y \cdot y^{-1} \to x \cdot y \cdot y^{-1} = z \to x = z)$
\item $\varphi_{19} :\equiv x \in \mathbb{R} \to y \in \mathbb{R} \to y \neq 0 \to x \cdot y = 0 \to x \cdot y \cdot y^{-1} = 0 \cdot y^{-1}$
\item $\varphi_{20} :\equiv \forall x \in \mathbb{R}\,(x \neq 0 \to x \cdot x^{-1} = 1) \to y \in \mathbb{R} \to y \neq 0 \to y \cdot y^{-1} = 1$
\item $\varphi_{21} :\equiv \forall v_{850}\,(y \cdot y^{-1} = v_{850} \to v_{850} = y \cdot y^{-1}) \to y \cdot y^{-1} = 1 \to 1 = y \cdot y^{-1}$
\item $\varphi_{22} :\equiv (x \cdot 1 = x \to x = x \cdot 1) \to x \in \mathbb{R} \to y \in \mathbb{R} \to y \neq 0 \to x \cdot y = 0 \to x = x \cdot 1$
\item $\varphi_{23} :\equiv x \cdot y = 0 \to x \cdot y \cdot y^{-1} = x \cdot y \cdot y^{-1} \to x \cdot y \cdot y^{-1} = 0 \cdot y^{-1}$
\item $\varphi_{24} :\equiv x = x \cdot (y \cdot y^{-1}) \to x \cdot (y \cdot y^{-1}) = x \cdot y \cdot y^{-1} \to x = x \cdot y \cdot y^{-1}$
\item $\varphi_{25} :\equiv \forall z\,(x = y^{-1} \cdot 0 \to y^{-1} \cdot 0 = z \to x = z) \to x = y^{-1} \cdot 0 \to y^{-1} \cdot 0 = 0 \to x = 0$
\item $\varphi_{26} :\equiv \forall x \in \mathbb{R}\,(x \neq 0 \to x^{-1} \in \mathbb{R}) \to y \in \mathbb{R} \to y \neq 0 \to y^{-1} \in \mathbb{R}$
\item $\varphi_{27} :\equiv x \in \mathbb{R} \to y \in \mathbb{R} \to y^{-1} \in \mathbb{R} \to x \cdot y \cdot y^{-1} = x \cdot (y \cdot y^{-1})$
\item $\varphi_{28} :\equiv \forall v_{849}\,\forall v_{850}\,(v_{849} = v_{850} \to v_{850} = v_{849}) \to \forall v_{850}\,(y \cdot y^{-1} = v_{850} \to v_{850} = y \cdot y^{-1})$
\item $\varphi_{29} :\equiv \forall w\,(x \cdot y = w \to x \cdot y \cdot y^{-1} = x \cdot y \cdot y^{-1} \to x \cdot y \cdot y^{-1} = w \cdot y^{-1})$
\item $\varphi_{30} :\equiv \forall z\,\forall w\,(z = w \to x \cdot 1 = x \cdot z \to x \cdot 1 = x \cdot w) \to \forall w\,(1 = w \to x \cdot 1 = x \cdot 1 \to x \cdot 1 = x \cdot w)$
\item $\varphi_{31} :\equiv (x = y^{-1} \cdot 0 \to y^{-1} \cdot 0 = 0 \to x = 0) \to x \in \mathbb{R} \to y \in \mathbb{R} \to y \neq 0 \to x \cdot y = 0 \to x = 0$
\item $\varphi_{32} :\equiv \forall v_{851}\,\forall v_{852}\,(v_{851} = v_{852} \to v_{852} = v_{851}) \to \forall v_{852}\,(x \cdot y \cdot y^{-1} = v_{852} \to v_{852} = x \cdot y \cdot y^{-1})$
\item $\varphi_{33} :\equiv (y \cdot y^{-1} = 1 \to 1 = y \cdot y^{-1}) \to x \in \mathbb{R} \to y \in \mathbb{R} \to y \neq 0 \to x \cdot y = 0 \to 1 = y \cdot y^{-1}$
\item $\varphi_{34} :\equiv \forall x\,\forall y\,(x \in \mathbb{R} \to y \in \mathbb{R} \to x \cdot y = y \cdot x) \to \forall y\,(0 \in \mathbb{R} \to y \in \mathbb{R} \to 0 \cdot y = y \cdot 0)$
\item $\varphi_{35} :\equiv \forall z\,(x = 0 \cdot y^{-1} \to 0 \cdot y^{-1} = z \to x = z) \to x = 0 \cdot y^{-1} \to 0 \cdot y^{-1} = y^{-1} \cdot 0 \to x = y^{-1} \cdot 0$
\item $\varphi_{36} :\equiv \forall w\,(1 = w \to x \cdot 1 = x \cdot 1 \to x \cdot 1 = x \cdot w) \to 1 = y \cdot y^{-1} \to x \cdot 1 = x \cdot 1 \to x \cdot 1 = x \cdot (y \cdot y^{-1})$
\item $\varphi_{37} :\equiv \forall y\,(0 \in \mathbb{R} \to y \in \mathbb{R} \to 0 \cdot y = y \cdot 0) \to 0 \in \mathbb{R} \to y^{-1} \in \mathbb{R} \to 0 \cdot y^{-1} = y^{-1} \cdot 0$
\item $\varphi_{38} :\equiv (x \in \mathbb{R} \to y \in \mathbb{R} \to y \neq 0 \to x \cdot y = 0 \to y \neq 0) \to x \in \mathbb{R} \to y \in \mathbb{R} \to y \neq 0 \to x \cdot y = 0 \to y \cdot y^{-1} = 1$
\item $\varphi_{39} :\equiv (x \in \mathbb{R} \to y \in \mathbb{R} \to y \neq 0 \to x \cdot y = 0 \to x \in \mathbb{R}) \to x \in \mathbb{R} \to y \in \mathbb{R} \to y \neq 0 \to x \cdot y = 0 \to x \cdot 1 = x$
\item $\varphi_{40} :\equiv (x \in \mathbb{R} \to y \in \mathbb{R} \to y \neq 0 \to x \cdot y = 0 \to y \neq 0) \to x \in \mathbb{R} \to y \in \mathbb{R} \to y \neq 0 \to x \cdot y = 0 \to y^{-1} \in \mathbb{R}$
\item $\varphi_{41} :\equiv (x = 0 \cdot y^{-1} \to 0 \cdot y^{-1} = y^{-1} \cdot 0 \to x = y^{-1} \cdot 0) \to x \in \mathbb{R} \to y \in \mathbb{R} \to y \neq 0 \to x \cdot y = 0 \to x = y^{-1} \cdot 0$
\item $\varphi_{42} :\equiv (x \in \mathbb{R} \to y \in \mathbb{R} \to y \neq 0 \to x \cdot y = 0 \to y^{-1} \in \mathbb{R}) \to x \in \mathbb{R} \to y \in \mathbb{R} \to y \neq 0 \to x \cdot y = 0 \to y^{-1} \cdot 0 = 0$
\item $\varphi_{43} :\equiv \forall v_{852}\,(x \cdot y \cdot y^{-1} = v_{852} \to v_{852} = x \cdot y \cdot y^{-1}) \to \varphi_{17}$
\item $\varphi_{44} :\equiv (1 = y \cdot y^{-1} \to x \cdot 1 = x \cdot 1 \to x \cdot 1 = x \cdot (y \cdot y^{-1})) \to \varphi_{16}$
\item $\varphi_{45} :\equiv (x \in \mathbb{R} \to y \in \mathbb{R} \to y \neq 0 \to x \cdot y = 0 \to y^{-1} \in \mathbb{R}) \to x \in \mathbb{R} \to \varphi_{10}$
\item $\varphi_{46} :\equiv \varphi_{24} \to x \in \mathbb{R} \to y \in \mathbb{R} \to y \neq 0 \to x \cdot y = 0 \to x = x \cdot y \cdot y^{-1}$
\item $\varphi_{47} :\equiv 0 \in \mathbb{R} \to (x \in \mathbb{R} \to y \in \mathbb{R} \to y \neq 0 \to x \cdot y = 0 \to y^{-1} \in \mathbb{R}) \to \varphi_{15}$
\item $\varphi_{48} :\equiv (x \in \mathbb{R} \to y \in \mathbb{R} \to y \neq 0 \to x \cdot y = 0 \to y^{-1} \cdot 0 = 0) \to \varphi_{31}$
\item $\varphi_{49} :\equiv (x \in \mathbb{R} \to y \in \mathbb{R} \to y \neq 0 \to x \cdot y = 0 \to y \cdot y^{-1} = 1) \to \varphi_{33}$
\item $\varphi_{50} :\equiv \varphi_{19} \to \varphi_{12} \to x \in \mathbb{R} \to y \in \mathbb{R} \to y \neq 0 \to x \cdot y = 0 \to x = 0 \cdot y^{-1}$
\item $\varphi_{51} :\equiv (y \in \mathbb{R} \to y \neq 0 \to y \cdot y^{-1} = 1) \to (x \in \mathbb{R} \to y \in \mathbb{R} \to y \neq 0 \to x \cdot y = 0 \to y \in \mathbb{R}) \to \varphi_{38}$
\item $\varphi_{52} :\equiv (x \in \mathbb{R} \to y \in \mathbb{R} \to y \neq 0 \to x \cdot y = 0 \to 1 = y \cdot y^{-1}) \to x \cdot 1 = x \cdot 1 \to \varphi_{44}$
\item $\varphi_{53} :\equiv (y \in \mathbb{R} \to y \neq 0 \to y^{-1} \in \mathbb{R}) \to (x \in \mathbb{R} \to y \in \mathbb{R} \to y \neq 0 \to x \cdot y = 0 \to y \in \mathbb{R}) \to \varphi_{40}$
\item $\varphi_{54} :\equiv (x \in \mathbb{R} \to y \in \mathbb{R} \to y \neq 0 \to x \cdot y = 0 \to x = y^{-1} \cdot 0) \to \varphi_{48}$
\item $\varphi_{55} :\equiv (x \in \mathbb{R} \to y \in \mathbb{R} \to y \neq 0 \to x \cdot y = 0 \to x \cdot y = 0) \to x \cdot y \cdot y^{-1} = x \cdot y \cdot y^{-1} \to \varphi_{23} \to \varphi_{19}$
\item $\varphi_{56} :\equiv (x \in \mathbb{R} \to y \in \mathbb{R} \to y \neq 0 \to x \cdot y = 0 \to x = x \cdot 1) \to \varphi_{16} \to \varphi_{6} \to \varphi_{9}$
\item $\varphi_{57} :\equiv (x \in \mathbb{R} \to y \in \mathbb{R} \to y \neq 0 \to x \cdot y = 0 \to x = 0 \cdot y^{-1}) \to \varphi_{15} \to \varphi_{41}$
\item $\varphi_{58} :\equiv (x \in \mathbb{R} \to y \in \mathbb{R} \to y \neq 0 \to x \cdot y = 0 \to x \in \mathbb{R}) \to (x \in \mathbb{R} \to y \in \mathbb{R} \to y \neq 0 \to x \cdot y = 0 \to y \in \mathbb{R}) \to \varphi_{45}$
\item $\varphi_{59} :\equiv (x \in \mathbb{R} \to y \in \mathbb{R} \to y \neq 0 \to x \cdot y = 0 \to x = x \cdot y \cdot y^{-1}) \to \varphi_{50}$
\end{abbrevlist}
\end{fullproof}

% 自動生成 (prover/build.js): PowAbs  ノード数 409
\begin{fullproof}{[P1]}{$\forall r \in \mathbb{R}\,\forall n \in \mathbb{N}\,(\lvert {{r}^{n}} \rvert = {{\lvert r \rvert}^{n}})$}{409}\phantomsection\label{tree:PowAbs}
\noindent$\Pi_{1}$:
\begin{ptree}
\AXC{$\forall n\,(n \in \mathbb{N} \leftrightarrow n \in \omega \land n \neq \emptyset)$\lbl{Def$\mathbb{N}$}}
\AXC{$\varphi_{27}$\lbl{A2}}
\BIC{$n \in \mathbb{N} \leftrightarrow n \in \omega \land n \neq \emptyset$}
\AXC{$r \in \mathbb{R} \to n \in \mathbb{N} \to n \in \mathbb{N}$\lbl{A1}}
\AXC{$\varphi_{83}$\lbl{A1}}
\BIC{$(n \in \mathbb{N} \leftrightarrow n \in \omega \land n \neq \emptyset) \to r \in \mathbb{R} \to n \in \mathbb{N} \to n \in \omega$}
\BIC{$r \in \mathbb{R} \to n \in \mathbb{N} \to n \in \omega$}
\end{ptree}
\noindent$\Pi_{2}$:
\begin{ptree}
\AXC{$\varphi_{111}$\lbl{A2}}
\AXC{$r \in \mathbb{R} \to \varphi_{115}$\lbl{A1}}
\AXC{$(r \in \mathbb{R} \to \varphi_{115}) \to \varphi_{111} \to r \in \mathbb{R} \to \varphi_{111}$\lbl{A1}}
\BIC{$\varphi_{111} \to r \in \mathbb{R} \to \varphi_{111}$}
\BIC{$r \in \mathbb{R} \to \varphi_{111}$}
\end{ptree}
\noindent$\Pi_{3}$:
\begin{ptree}
\AXC{$y = z \to x = y \to x = z$\lbl{A6}}
\AXC{$(y = z \to x = y \to x = z) \to x = y \to y = z \to x = z$\lbl{A1}}
\BIC{$x = y \to y = z \to x = z$}
\UIC{$\forall z\,(x = y \to y = z \to x = z)$}
\UIC{$\forall y\,\forall z\,(x = y \to y = z \to x = z)$}
\UIC{$\forall x\,\forall y\,\forall z\,(x = y \to y = z \to x = z)$}
\AXC{$\varphi_{35}$\lbl{A2}}
\BIC{$\forall y\,\forall z\,(\lvert {{r}^{S(n)}} \rvert = y \to y = z \to \lvert {{r}^{S(n)}} \rvert = z)$}
\end{ptree}
\noindent$\Pi_{4}$:
\begin{ptree}
\AXC{$\stackrel{\Pi_{3}}{\vdots}$}\noLine
\UIC{$\forall y\,\forall z\,(\lvert {{r}^{S(n)}} \rvert = y \to y = z \to \lvert {{r}^{S(n)}} \rvert = z)$}
\AXC{$\forall y\,\forall z\,(\lvert {{r}^{S(n)}} \rvert = y \to y = z \to \lvert {{r}^{S(n)}} \rvert = z) \to \varphi_{77}$\lbl{A2}}
\BIC{$\varphi_{77}$}
\AXC{$\varphi_{77} \to \varphi_{108}$\lbl{A2}}
\BIC{$\varphi_{108}$}
\end{ptree}
\noindent$\Pi_{5}$:
\begin{ptree}
\AXC{$x = x$\lbl{A5}}
\AXC{$x = y \to x = x \to y = x$\lbl{A6}}
\AXC{$(x = y \to x = x \to y = x) \to x = x \to x = y \to y = x$\lbl{A1}}
\BIC{$x = x \to x = y \to y = x$}
\BIC{$x = y \to y = x$}
\UIC{$\forall y\,(x = y \to y = x)$}
\UIC{$\forall x\,\forall y\,(x = y \to y = x)$}
\end{ptree}
\noindent$\Pi_{6}$:
\begin{ptree}
\AXC{$\stackrel{\Pi_{5}}{\vdots}$}\noLine
\UIC{$\forall x\,\forall y\,(x = y \to y = x)$}
\AXC{$\varphi_{19}$\lbl{A2}}
\BIC{$\forall y\,({{\lvert r \rvert}^{S(n)}} = y \to y = {{\lvert r \rvert}^{S(n)}})$}
\AXC{$\forall y\,({{\lvert r \rvert}^{S(n)}} = y \to y = {{\lvert r \rvert}^{S(n)}}) \to \varphi_{64}$\lbl{A2}}
\BIC{$\varphi_{64}$}
\end{ptree}
\noindent$\Pi_{7}$:
\begin{ptree}
\AXC{$\varphi_{111}$\lbl{A2}}
\AXC{$r \in \mathbb{R} \to \varphi_{115}$\lbl{A1}}
\AXC{$(r \in \mathbb{R} \to \varphi_{115}) \to \varphi_{111} \to r \in \mathbb{R} \to \varphi_{111}$\lbl{A1}}
\BIC{$\varphi_{111} \to r \in \mathbb{R} \to \varphi_{111}$}
\BIC{$r \in \mathbb{R} \to \varphi_{111}$}
\AXC{$(r \in \mathbb{R} \to \varphi_{111}) \to (r \in \mathbb{R} \to \varphi_{66}) \to r \in \mathbb{R} \to \varphi_{112}$\lbl{A1}}
\BIC{$(r \in \mathbb{R} \to \varphi_{66}) \to r \in \mathbb{R} \to \varphi_{112}$}
\end{ptree}
\noindent$\Pi_{8}$:
\begin{ptree}
\AXC{$r \in \mathbb{R} \to \varphi_{66}$\lbl{A1}}
\AXC{$\stackrel{\Pi_{7}}{\vdots}$}\noLine
\UIC{$(r \in \mathbb{R} \to \varphi_{66}) \to r \in \mathbb{R} \to \varphi_{112}$}
\BIC{$r \in \mathbb{R} \to \varphi_{112}$}
\end{ptree}
\noindent$\Pi_{9}$:
\begin{ptree}
\AXC{$\stackrel{\Pi_{8}}{\vdots}$}\noLine
\UIC{$r \in \mathbb{R} \to \varphi_{112}$}
\AXC{$(r \in \mathbb{R} \to \varphi_{112}) \to r \in \mathbb{R} \to \varphi_{71}$\lbl{A1}}
\BIC{$r \in \mathbb{R} \to \varphi_{71}$}
\end{ptree}
\noindent$\Pi_{10}$:
\begin{ptree}
\AXC{$r \in \mathbb{R} \to \varphi_{100}$\lbl{A1}}
\AXC{$\stackrel{\Pi_{9}}{\vdots}$}\noLine
\UIC{$r \in \mathbb{R} \to \varphi_{71}$}
\AXC{$\forall n\,(n \in \mathbb{N} \leftrightarrow n \in \omega \land n \neq \emptyset)$\lbl{Def$\mathbb{N}$}}
\AXC{$\varphi_{27}$\lbl{A2}}
\BIC{$n \in \mathbb{N} \leftrightarrow n \in \omega \land n \neq \emptyset$}
\AXC{$\varphi_{155}$\lbl{A1}}
\BIC{$(r \in \mathbb{R} \to \varphi_{71}) \to (r \in \mathbb{R} \to \varphi_{100}) \to r \in \mathbb{R} \to \varphi_{96}$}
\BIC{$(r \in \mathbb{R} \to \varphi_{100}) \to r \in \mathbb{R} \to \varphi_{96}$}
\BIC{$r \in \mathbb{R} \to \varphi_{96}$}
\end{ptree}
\noindent$\Pi_{11}$:
\begin{ptree}
\AXC{$\forall x \in \mathbb{R}\,(\lvert x \rvert \in \mathbb{R})$\lbl{[A5]}}
\AXC{$\forall x \in \mathbb{R}\,(\lvert x \rvert \in \mathbb{R}) \to r \in \mathbb{R} \to \lvert r \rvert \in \mathbb{R}$\lbl{A2}}
\BIC{$r \in \mathbb{R} \to \lvert r \rvert \in \mathbb{R}$}
\end{ptree}
\noindent$\Pi_{12}$:
\begin{ptree}
\AXC{$r \in \mathbb{R} \to \varphi_{98}$\lbl{A1}}
\AXC{$\stackrel{\Pi_{11}}{\vdots}$}\noLine
\UIC{$r \in \mathbb{R} \to \lvert r \rvert \in \mathbb{R}$}
\AXC{$(r \in \mathbb{R} \to \lvert r \rvert \in \mathbb{R}) \to (r \in \mathbb{R} \to \varphi_{98}) \to r \in \mathbb{R} \to \varphi_{103}$\lbl{A1}}
\BIC{$(r \in \mathbb{R} \to \varphi_{98}) \to r \in \mathbb{R} \to \varphi_{103}$}
\BIC{$r \in \mathbb{R} \to \varphi_{103}$}
\end{ptree}
\noindent$\Pi_{13}$:
\begin{ptree}
\AXC{$\forall r \in \mathbb{R}\,(\varphi_{29})$\lbl{[S24]}}
\AXC{$\forall r \in \mathbb{R}\,(\varphi_{29}) \to \lvert r \rvert \in \mathbb{R} \to \varphi_{81}$\lbl{A2}}
\BIC{$\lvert r \rvert \in \mathbb{R} \to \varphi_{81}$}
\AXC{$(\lvert r \rvert \in \mathbb{R} \to \varphi_{81}) \to (r \in \mathbb{R} \to \varphi_{103}) \to r \in \mathbb{R} \to \varphi_{145}$\lbl{A1}}
\BIC{$(r \in \mathbb{R} \to \varphi_{103}) \to r \in \mathbb{R} \to \varphi_{145}$}
\end{ptree}
\noindent$\Pi_{14}$:
\begin{ptree}
\AXC{$\stackrel{\Pi_{12}}{\vdots}$}\noLine
\UIC{$r \in \mathbb{R} \to \varphi_{103}$}
\AXC{$\stackrel{\Pi_{13}}{\vdots}$}\noLine
\UIC{$(r \in \mathbb{R} \to \varphi_{103}) \to r \in \mathbb{R} \to \varphi_{145}$}
\BIC{$r \in \mathbb{R} \to \varphi_{145}$}
\AXC{$(r \in \mathbb{R} \to \varphi_{145}) \to r \in \mathbb{R} \to \varphi_{136}$\lbl{A1}}
\BIC{$r \in \mathbb{R} \to \varphi_{136}$}
\AXC{$(r \in \mathbb{R} \to \varphi_{136}) \to \varphi_{91} \to r \in \mathbb{R} \to \varphi_{135}$\lbl{A1}}
\BIC{$\varphi_{91} \to r \in \mathbb{R} \to \varphi_{135}$}
\end{ptree}
\noindent$\Pi_{15}$:
\begin{ptree}
\AXC{$\stackrel{\Pi_{10}}{\vdots}$}\noLine
\UIC{$r \in \mathbb{R} \to \varphi_{96}$}
\AXC{$\varphi_{91}$\lbl{A2}}
\AXC{$\stackrel{\Pi_{14}}{\vdots}$}\noLine
\UIC{$\varphi_{91} \to r \in \mathbb{R} \to \varphi_{135}$}
\BIC{$r \in \mathbb{R} \to \varphi_{135}$}
\AXC{$(r \in \mathbb{R} \to \varphi_{135}) \to (r \in \mathbb{R} \to \varphi_{96}) \to r \in \mathbb{R} \to \varphi_{123}$\lbl{A1}}
\BIC{$(r \in \mathbb{R} \to \varphi_{96}) \to r \in \mathbb{R} \to \varphi_{123}$}
\BIC{$r \in \mathbb{R} \to \varphi_{123}$}
\end{ptree}
\noindent$\Pi_{16}$:
\begin{ptree}
\AXC{$y = z \to x = y \to x = z$\lbl{A6}}
\AXC{$(y = z \to x = y \to x = z) \to x = y \to y = z \to x = z$\lbl{A1}}
\BIC{$x = y \to y = z \to x = z$}
\UIC{$\forall z\,(x = y \to y = z \to x = z)$}
\UIC{$\forall y\,\forall z\,(x = y \to y = z \to x = z)$}
\UIC{$\forall x\,\forall y\,\forall z\,(x = y \to y = z \to x = z)$}
\AXC{$\varphi_{35}$\lbl{A2}}
\BIC{$\forall y\,\forall z\,(\lvert {{r}^{S(n)}} \rvert = y \to y = z \to \lvert {{r}^{S(n)}} \rvert = z)$}
\end{ptree}
\noindent$\Pi_{17}$:
\begin{ptree}
\AXC{$\stackrel{\Pi_{16}}{\vdots}$}\noLine
\UIC{$\forall y\,\forall z\,(\lvert {{r}^{S(n)}} \rvert = y \to y = z \to \lvert {{r}^{S(n)}} \rvert = z)$}
\AXC{$\forall y\,\forall z\,(\lvert {{r}^{S(n)}} \rvert = y \to y = z \to \lvert {{r}^{S(n)}} \rvert = z) \to \varphi_{78}$\lbl{A2}}
\BIC{$\varphi_{78}$}
\end{ptree}
\noindent$\Pi_{18}$:
\begin{ptree}
\AXC{$\stackrel{\Pi_{10}}{\vdots}$}\noLine
\UIC{$r \in \mathbb{R} \to \varphi_{96}$}
\AXC{$\stackrel{\Pi_{8}}{\vdots}$}\noLine
\UIC{$r \in \mathbb{R} \to \varphi_{112}$}
\AXC{$(r \in \mathbb{R} \to \varphi_{112}) \to r \in \mathbb{R} \to \varphi_{105}$\lbl{A1}}
\BIC{$r \in \mathbb{R} \to \varphi_{105}$}
\AXC{$(r \in \mathbb{R} \to \varphi_{105}) \to (r \in \mathbb{R} \to \varphi_{96}) \to r \in \mathbb{R} \to \varphi_{116}$\lbl{A1}}
\BIC{$(r \in \mathbb{R} \to \varphi_{96}) \to r \in \mathbb{R} \to \varphi_{116}$}
\BIC{$r \in \mathbb{R} \to \varphi_{116}$}
\end{ptree}
\noindent$\Pi_{19}$:
\begin{ptree}
\AXC{$x = x$\lbl{A5}}
\UIC{$\forall x\,(x = x)$}
\AXC{$\forall x\,(x = x) \to \lvert {{r}^{n}} \rvert \cdot \lvert r \rvert = \lvert {{r}^{n}} \rvert \cdot \lvert r \rvert$\lbl{A2}}
\BIC{$\lvert {{r}^{n}} \rvert \cdot \lvert r \rvert = \lvert {{r}^{n}} \rvert \cdot \lvert r \rvert$}
\AXC{$\stackrel{\Pi_{18}}{\vdots}$}\noLine
\UIC{$r \in \mathbb{R} \to \varphi_{116}$}
\AXC{$(r \in \mathbb{R} \to \varphi_{116}) \to \varphi_{154}$\lbl{A1}}
\BIC{$\varphi_{154}$}
\BIC{$(\lvert {{r}^{n}} \rvert = {{\lvert r \rvert}^{n}} \to \varphi_{84}) \to r \in \mathbb{R} \to \varphi_{131}$}
\end{ptree}
\noindent$\Pi_{20}$:
\begin{ptree}
\AXC{$y = z \to x = y \to x = z$\lbl{A6}}
\AXC{$(y = z \to x = y \to x = z) \to x = y \to y = z \to x = z$\lbl{A1}}
\BIC{$x = y \to y = z \to x = z$}
\UIC{$\forall z\,(x = y \to y = z \to x = z)$}
\UIC{$\forall y\,\forall z\,(x = y \to y = z \to x = z)$}
\UIC{$\forall x\,\forall y\,\forall z\,(x = y \to y = z \to x = z)$}
\AXC{$\varphi_{35}$\lbl{A2}}
\BIC{$\forall y\,\forall z\,(\lvert {{r}^{S(n)}} \rvert = y \to y = z \to \lvert {{r}^{S(n)}} \rvert = z)$}
\end{ptree}
\noindent$\Pi_{21}$:
\begin{ptree}
\AXC{$\stackrel{\Pi_{20}}{\vdots}$}\noLine
\UIC{$\forall y\,\forall z\,(\lvert {{r}^{S(n)}} \rvert = y \to y = z \to \lvert {{r}^{S(n)}} \rvert = z)$}
\AXC{$\forall y\,\forall z\,(\lvert {{r}^{S(n)}} \rvert = y \to y = z \to \lvert {{r}^{S(n)}} \rvert = z) \to \varphi_{62}$\lbl{A2}}
\BIC{$\varphi_{62}$}
\end{ptree}
\noindent$\Pi_{22}$:
\begin{ptree}
\AXC{$\forall r \in \mathbb{R}\,(\varphi_{29})$\lbl{[S24]}}
\AXC{$\forall r \in \mathbb{R}\,(\varphi_{29}) \to r \in \mathbb{R} \to \varphi_{29}$\lbl{A2}}
\BIC{$r \in \mathbb{R} \to \varphi_{29}$}
\AXC{$(r \in \mathbb{R} \to \varphi_{29}) \to (r \in \mathbb{R} \to \varphi_{98}) \to r \in \mathbb{R} \to \varphi_{142}$\lbl{A1}}
\BIC{$(r \in \mathbb{R} \to \varphi_{98}) \to r \in \mathbb{R} \to \varphi_{142}$}
\end{ptree}
\noindent$\Pi_{23}$:
\begin{ptree}
\AXC{$r \in \mathbb{R} \to \varphi_{98}$\lbl{A1}}
\AXC{$\stackrel{\Pi_{22}}{\vdots}$}\noLine
\UIC{$(r \in \mathbb{R} \to \varphi_{98}) \to r \in \mathbb{R} \to \varphi_{142}$}
\BIC{$r \in \mathbb{R} \to \varphi_{142}$}
\AXC{$(r \in \mathbb{R} \to \varphi_{142}) \to r \in \mathbb{R} \to \varphi_{99}$\lbl{A1}}
\BIC{$r \in \mathbb{R} \to \varphi_{99}$}
\end{ptree}
\noindent$\Pi_{24}$:
\begin{ptree}
\AXC{$\stackrel{\Pi_{10}}{\vdots}$}\noLine
\UIC{$r \in \mathbb{R} \to \varphi_{96}$}
\AXC{$\stackrel{\Pi_{23}}{\vdots}$}\noLine
\UIC{$r \in \mathbb{R} \to \varphi_{99}$}
\AXC{$\forall a\,\forall n\,(\mathrm{Seq}(a) \to n \in \mathbb{N} \to a_{n} \in \mathbb{R})$\lbl{[S1]}}
\AXC{$\varphi_{57}$\lbl{A2}}
\BIC{$\forall n\,(\mathrm{Seq}({(r)^{\bullet}}) \to n \in \mathbb{N} \to {{r}^{n}} \in \mathbb{R})$}
\AXC{$\varphi_{63}$\lbl{A2}}
\BIC{$\mathrm{Seq}({(r)^{\bullet}}) \to n \in \mathbb{N} \to {{r}^{n}} \in \mathbb{R}$}
\AXC{$\varphi_{157}$\lbl{A1}}
\BIC{$(r \in \mathbb{R} \to \varphi_{99}) \to (r \in \mathbb{R} \to \varphi_{96}) \to r \in \mathbb{R} \to \varphi_{107}$}
\BIC{$(r \in \mathbb{R} \to \varphi_{96}) \to r \in \mathbb{R} \to \varphi_{107}$}
\BIC{$r \in \mathbb{R} \to \varphi_{107}$}
\end{ptree}
\noindent$\Pi_{25}$:
\begin{ptree}
\AXC{$\stackrel{\Pi_{24}}{\vdots}$}\noLine
\UIC{$r \in \mathbb{R} \to \varphi_{107}$}
\AXC{$\forall x\,\forall y\,(x \in \mathbb{R} \to y \in \mathbb{R} \to \lvert x \cdot y \rvert = \lvert x \rvert \cdot \lvert y \rvert)$\lbl{[A17]}}
\AXC{$\varphi_{113}$\lbl{A2}}
\BIC{$\forall y\,({{r}^{n}} \in \mathbb{R} \to y \in \mathbb{R} \to \lvert {{r}^{n}} \cdot y \rvert = \lvert {{r}^{n}} \rvert \cdot \lvert y \rvert)$}
\AXC{$\varphi_{134}$\lbl{A2}}
\BIC{${{r}^{n}} \in \mathbb{R} \to r \in \mathbb{R} \to \lvert {{r}^{n}} \cdot r \rvert = \lvert {{r}^{n}} \rvert \cdot \lvert r \rvert$}
\AXC{$\varphi_{162}$\lbl{A1}}
\BIC{$(r \in \mathbb{R} \to \varphi_{107}) \to (r \in \mathbb{R} \to \varphi_{98}) \to r \in \mathbb{R} \to \varphi_{129}$}
\BIC{$(r \in \mathbb{R} \to \varphi_{98}) \to r \in \mathbb{R} \to \varphi_{129}$}
\end{ptree}
\noindent$\Pi_{26}$:
\begin{ptree}
\AXC{$\forall r \in \mathbb{R}\,(\varphi_{29})$\lbl{[S24]}}
\AXC{$\forall r \in \mathbb{R}\,(\varphi_{29}) \to r \in \mathbb{R} \to \varphi_{29}$\lbl{A2}}
\BIC{$r \in \mathbb{R} \to \varphi_{29}$}
\AXC{$(r \in \mathbb{R} \to \varphi_{29}) \to (r \in \mathbb{R} \to \varphi_{98}) \to r \in \mathbb{R} \to \varphi_{142}$\lbl{A1}}
\BIC{$(r \in \mathbb{R} \to \varphi_{98}) \to r \in \mathbb{R} \to \varphi_{142}$}
\end{ptree}
\noindent$\Pi_{27}$:
\begin{ptree}
\AXC{$r \in \mathbb{R} \to \varphi_{98}$\lbl{A1}}
\AXC{$\stackrel{\Pi_{26}}{\vdots}$}\noLine
\UIC{$(r \in \mathbb{R} \to \varphi_{98}) \to r \in \mathbb{R} \to \varphi_{142}$}
\BIC{$r \in \mathbb{R} \to \varphi_{142}$}
\AXC{$(r \in \mathbb{R} \to \varphi_{142}) \to r \in \mathbb{R} \to \varphi_{125}$\lbl{A1}}
\BIC{$r \in \mathbb{R} \to \varphi_{125}$}
\AXC{$(r \in \mathbb{R} \to \varphi_{125}) \to \varphi_{58} \to r \in \mathbb{R} \to \varphi_{121}$\lbl{A1}}
\BIC{$\varphi_{58} \to r \in \mathbb{R} \to \varphi_{121}$}
\end{ptree}
\noindent$\Pi_{28}$:
\begin{ptree}
\AXC{$\stackrel{\Pi_{10}}{\vdots}$}\noLine
\UIC{$r \in \mathbb{R} \to \varphi_{96}$}
\AXC{$\varphi_{58}$\lbl{A2}}
\AXC{$\stackrel{\Pi_{27}}{\vdots}$}\noLine
\UIC{$\varphi_{58} \to r \in \mathbb{R} \to \varphi_{121}$}
\BIC{$r \in \mathbb{R} \to \varphi_{121}$}
\AXC{$(r \in \mathbb{R} \to \varphi_{121}) \to (r \in \mathbb{R} \to \varphi_{96}) \to r \in \mathbb{R} \to \varphi_{114}$\lbl{A1}}
\BIC{$(r \in \mathbb{R} \to \varphi_{96}) \to r \in \mathbb{R} \to \varphi_{114}$}
\BIC{$r \in \mathbb{R} \to \varphi_{114}$}
\AXC{$\varphi_{158}$\lbl{A1}}
\BIC{$\lvert {{r}^{S(n)}} \rvert = \lvert {{r}^{S(n)}} \rvert \to \varphi_{90} \to r \in \mathbb{R} \to \varphi_{120}$}
\end{ptree}
\noindent$\Pi_{29}$:
\begin{ptree}
\AXC{$x = x$\lbl{A5}}
\UIC{$\forall x\,(x = x)$}
\AXC{$\forall x\,(x = x) \to \lvert {{r}^{S(n)}} \rvert = \lvert {{r}^{S(n)}} \rvert$\lbl{A2}}
\BIC{$\lvert {{r}^{S(n)}} \rvert = \lvert {{r}^{S(n)}} \rvert$}
\AXC{$\stackrel{\Pi_{28}}{\vdots}$}\noLine
\UIC{$\lvert {{r}^{S(n)}} \rvert = \lvert {{r}^{S(n)}} \rvert \to \varphi_{90} \to r \in \mathbb{R} \to \varphi_{120}$}
\BIC{$\varphi_{90} \to r \in \mathbb{R} \to \varphi_{120}$}
\end{ptree}
\noindent$\Pi_{30}$:
\begin{ptree}
\AXC{$z = w \to \lvert {{r}^{S(n)}} \rvert = \lvert z \rvert \to \lvert {{r}^{S(n)}} \rvert = \lvert w \rvert$\lbl{A6}}
\UIC{$\forall w\,(z = w \to \lvert {{r}^{S(n)}} \rvert = \lvert z \rvert \to \lvert {{r}^{S(n)}} \rvert = \lvert w \rvert)$}
\UIC{$\varphi_{21}$}
\AXC{$\varphi_{21} \to \varphi_{54}$\lbl{A2}}
\BIC{$\varphi_{54}$}
\AXC{$\varphi_{54} \to \varphi_{90}$\lbl{A2}}
\BIC{$\varphi_{90}$}
\AXC{$\stackrel{\Pi_{29}}{\vdots}$}\noLine
\UIC{$\varphi_{90} \to r \in \mathbb{R} \to \varphi_{120}$}
\BIC{$r \in \mathbb{R} \to \varphi_{120}$}
\end{ptree}
\noindent$\Pi_{31}$:
\begin{ptree}
\AXC{$r \in \mathbb{R} \to \varphi_{98}$\lbl{A1}}
\AXC{$\stackrel{\Pi_{25}}{\vdots}$}\noLine
\UIC{$(r \in \mathbb{R} \to \varphi_{98}) \to r \in \mathbb{R} \to \varphi_{129}$}
\BIC{$r \in \mathbb{R} \to \varphi_{129}$}
\AXC{$\stackrel{\Pi_{30}}{\vdots}$}\noLine
\UIC{$r \in \mathbb{R} \to \varphi_{120}$}
\AXC{$(r \in \mathbb{R} \to \varphi_{120}) \to \varphi_{159}$\lbl{A1}}
\BIC{$\varphi_{159}$}
\BIC{$(\lvert {{r}^{S(n)}} \rvert = \lvert {{r}^{n}} \cdot r \rvert \to \varphi_{70}) \to r \in \mathbb{R} \to \varphi_{126}$}
\end{ptree}
\noindent$\Pi_{32}$:
\begin{ptree}
\AXC{$\stackrel{\Pi_{21}}{\vdots}$}\noLine
\UIC{$\varphi_{62}$}
\AXC{$\varphi_{62} \to \lvert {{r}^{S(n)}} \rvert = \lvert {{r}^{n}} \cdot r \rvert \to \varphi_{70}$\lbl{A2}}
\BIC{$\lvert {{r}^{S(n)}} \rvert = \lvert {{r}^{n}} \cdot r \rvert \to \varphi_{70}$}
\AXC{$\stackrel{\Pi_{31}}{\vdots}$}\noLine
\UIC{$(\lvert {{r}^{S(n)}} \rvert = \lvert {{r}^{n}} \cdot r \rvert \to \varphi_{70}) \to r \in \mathbb{R} \to \varphi_{126}$}
\BIC{$r \in \mathbb{R} \to \varphi_{126}$}
\AXC{$(r \in \mathbb{R} \to \varphi_{126}) \to \varphi_{160}$\lbl{A1}}
\BIC{$\varphi_{160}$}
\end{ptree}
\noindent$\Pi_{33}$:
\begin{ptree}
\AXC{$\varphi_{38}$\lbl{A6}}
\UIC{$\forall w\,(\varphi_{38})$}
\UIC{$\forall z\,\forall w\,(\varphi_{38})$}
\AXC{$\forall z\,\forall w\,(\varphi_{38}) \to \varphi_{88}$\lbl{A2}}
\BIC{$\varphi_{88}$}
\AXC{$\varphi_{88} \to \lvert {{r}^{n}} \rvert = {{\lvert r \rvert}^{n}} \to \varphi_{84}$\lbl{A2}}
\BIC{$\lvert {{r}^{n}} \rvert = {{\lvert r \rvert}^{n}} \to \varphi_{84}$}
\AXC{$\stackrel{\Pi_{19}}{\vdots}$}\noLine
\UIC{$(\lvert {{r}^{n}} \rvert = {{\lvert r \rvert}^{n}} \to \varphi_{84}) \to r \in \mathbb{R} \to \varphi_{131}$}
\BIC{$r \in \mathbb{R} \to \varphi_{131}$}
\AXC{$\stackrel{\Pi_{32}}{\vdots}$}\noLine
\UIC{$\varphi_{160}$}
\BIC{$(\lvert {{r}^{S(n)}} \rvert = \lvert {{r}^{n}} \rvert \cdot \lvert r \rvert \to \varphi_{75}) \to r \in \mathbb{R} \to \varphi_{127}$}
\end{ptree}
\noindent$\Pi_{34}$:
\begin{ptree}
\AXC{$\stackrel{\Pi_{17}}{\vdots}$}\noLine
\UIC{$\varphi_{78}$}
\AXC{$\varphi_{78} \to \lvert {{r}^{S(n)}} \rvert = \lvert {{r}^{n}} \rvert \cdot \lvert r \rvert \to \varphi_{75}$\lbl{A2}}
\BIC{$\lvert {{r}^{S(n)}} \rvert = \lvert {{r}^{n}} \rvert \cdot \lvert r \rvert \to \varphi_{75}$}
\AXC{$\stackrel{\Pi_{33}}{\vdots}$}\noLine
\UIC{$(\lvert {{r}^{S(n)}} \rvert = \lvert {{r}^{n}} \rvert \cdot \lvert r \rvert \to \varphi_{75}) \to r \in \mathbb{R} \to \varphi_{127}$}
\BIC{$r \in \mathbb{R} \to \varphi_{127}$}
\end{ptree}
\noindent$\Pi_{35}$:
\begin{ptree}
\AXC{$\stackrel{\Pi_{34}}{\vdots}$}\noLine
\UIC{$r \in \mathbb{R} \to \varphi_{127}$}
\AXC{$(r \in \mathbb{R} \to \varphi_{127}) \to (r \in \mathbb{R} \to \varphi_{124}) \to \varphi_{108} \to r \in \mathbb{R} \to \varphi_{118}$\lbl{A1}}
\BIC{$(r \in \mathbb{R} \to \varphi_{124}) \to \varphi_{108} \to r \in \mathbb{R} \to \varphi_{118}$}
\end{ptree}
\noindent$\Pi_{36}$:
\begin{ptree}
\AXC{$\stackrel{\Pi_{4}}{\vdots}$}\noLine
\UIC{$\varphi_{108}$}
\AXC{$\stackrel{\Pi_{6}}{\vdots}$}\noLine
\UIC{$\varphi_{64}$}
\AXC{$\stackrel{\Pi_{15}}{\vdots}$}\noLine
\UIC{$r \in \mathbb{R} \to \varphi_{123}$}
\AXC{$(r \in \mathbb{R} \to \varphi_{123}) \to \varphi_{64} \to r \in \mathbb{R} \to \varphi_{124}$\lbl{A1}}
\BIC{$\varphi_{64} \to r \in \mathbb{R} \to \varphi_{124}$}
\BIC{$r \in \mathbb{R} \to \varphi_{124}$}
\AXC{$\stackrel{\Pi_{35}}{\vdots}$}\noLine
\UIC{$(r \in \mathbb{R} \to \varphi_{124}) \to \varphi_{108} \to r \in \mathbb{R} \to \varphi_{118}$}
\BIC{$\varphi_{108} \to r \in \mathbb{R} \to \varphi_{118}$}
\BIC{$r \in \mathbb{R} \to \varphi_{118}$}
\end{ptree}
\noindent$\Pi_{37}$:
\begin{ptree}
\AXC{$y = z \to x = y \to x = z$\lbl{A6}}
\AXC{$(y = z \to x = y \to x = z) \to x = y \to y = z \to x = z$\lbl{A1}}
\BIC{$x = y \to y = z \to x = z$}
\UIC{$\forall z\,(x = y \to y = z \to x = z)$}
\UIC{$\forall y\,\forall z\,(x = y \to y = z \to x = z)$}
\UIC{$\forall x\,\forall y\,\forall z\,(x = y \to y = z \to x = z)$}
\AXC{$\varphi_{36}$\lbl{A2}}
\BIC{$\forall y\,\forall z\,(\lvert {{r}^{1_{\omega}}} \rvert = y \to y = z \to \lvert {{r}^{1_{\omega}}} \rvert = z)$}
\AXC{$\varphi_{93}$\lbl{A2}}
\BIC{$\forall z\,(\lvert {{r}^{1_{\omega}}} \rvert = \lvert r \rvert \to \lvert r \rvert = z \to \lvert {{r}^{1_{\omega}}} \rvert = z)$}
\AXC{$\varphi_{128}$\lbl{A2}}
\BIC{$\varphi_{67}$}
\end{ptree}
\noindent$\Pi_{38}$:
\begin{ptree}
\AXC{$\forall x \in \mathbb{R}\,(\lvert x \rvert \in \mathbb{R})$\lbl{[A5]}}
\AXC{$\forall x \in \mathbb{R}\,(\lvert x \rvert \in \mathbb{R}) \to r \in \mathbb{R} \to \lvert r \rvert \in \mathbb{R}$\lbl{A2}}
\BIC{$r \in \mathbb{R} \to \lvert r \rvert \in \mathbb{R}$}
\end{ptree}
\noindent$\Pi_{39}$:
\begin{ptree}
\AXC{$\stackrel{\Pi_{38}}{\vdots}$}\noLine
\UIC{$r \in \mathbb{R} \to \lvert r \rvert \in \mathbb{R}$}
\AXC{$\forall r \in \mathbb{R}\,(\varphi_{29})$\lbl{[S24]}}
\AXC{$\forall r \in \mathbb{R}\,(\varphi_{29}) \to \lvert r \rvert \in \mathbb{R} \to \varphi_{81}$\lbl{A2}}
\BIC{$\lvert r \rvert \in \mathbb{R} \to \varphi_{81}$}
\AXC{$\varphi_{150}$\lbl{A1}}
\BIC{$(r \in \mathbb{R} \to \lvert r \rvert \in \mathbb{R}) \to r \in \mathbb{R} \to \varphi_{81}$}
\BIC{$r \in \mathbb{R} \to \varphi_{81}$}
\end{ptree}
\noindent$\Pi_{40}$:
\begin{ptree}
\AXC{$\stackrel{\Pi_{39}}{\vdots}$}\noLine
\UIC{$r \in \mathbb{R} \to \varphi_{81}$}
\AXC{$(r \in \mathbb{R} \to \varphi_{81}) \to r \in \mathbb{R} \to {{\lvert r \rvert}^{1_{\omega}}} = \lvert r \rvert$\lbl{A1}}
\BIC{$r \in \mathbb{R} \to {{\lvert r \rvert}^{1_{\omega}}} = \lvert r \rvert$}
\AXC{$(r \in \mathbb{R} \to {{\lvert r \rvert}^{1_{\omega}}} = \lvert r \rvert) \to \varphi_{65}$\lbl{A1}}
\BIC{$\varphi_{65}$}
\end{ptree}
\noindent$\Pi_{41}$:
\begin{ptree}
\AXC{$\forall r \in \mathbb{R}\,(\varphi_{29})$\lbl{[S24]}}
\AXC{$\forall r \in \mathbb{R}\,(\varphi_{29}) \to r \in \mathbb{R} \to \varphi_{29}$\lbl{A2}}
\BIC{$r \in \mathbb{R} \to \varphi_{29}$}
\AXC{$(r \in \mathbb{R} \to \varphi_{29}) \to r \in \mathbb{R} \to {{r}^{1_{\omega}}} = r$\lbl{A1}}
\BIC{$r \in \mathbb{R} \to {{r}^{1_{\omega}}} = r$}
\end{ptree}
\noindent$\Pi_{42}$:
\begin{ptree}
\AXC{$x = x$\lbl{A5}}
\UIC{$\forall x\,(x = x)$}
\AXC{$\forall x\,(x = x) \to \lvert {{r}^{1_{\omega}}} \rvert = \lvert {{r}^{1_{\omega}}} \rvert$\lbl{A2}}
\BIC{$\lvert {{r}^{1_{\omega}}} \rvert = \lvert {{r}^{1_{\omega}}} \rvert$}
\AXC{$\stackrel{\Pi_{41}}{\vdots}$}\noLine
\UIC{$r \in \mathbb{R} \to {{r}^{1_{\omega}}} = r$}
\AXC{$(r \in \mathbb{R} \to {{r}^{1_{\omega}}} = r) \to \varphi_{132}$\lbl{A1}}
\BIC{$\varphi_{132}$}
\BIC{$\varphi_{49} \to r \in \mathbb{R} \to \lvert {{r}^{1_{\omega}}} \rvert = \lvert r \rvert$}
\end{ptree}
\noindent$\Pi_{43}$:
\begin{ptree}
\AXC{$z = w \to \lvert {{r}^{1_{\omega}}} \rvert = \lvert z \rvert \to \lvert {{r}^{1_{\omega}}} \rvert = \lvert w \rvert$\lbl{A6}}
\UIC{$\forall w\,(z = w \to \lvert {{r}^{1_{\omega}}} \rvert = \lvert z \rvert \to \lvert {{r}^{1_{\omega}}} \rvert = \lvert w \rvert)$}
\UIC{$\varphi_{22}$}
\AXC{$\varphi_{22} \to \varphi_{55}$\lbl{A2}}
\BIC{$\varphi_{55}$}
\AXC{$\varphi_{55} \to \varphi_{49}$\lbl{A2}}
\BIC{$\varphi_{49}$}
\AXC{$\stackrel{\Pi_{42}}{\vdots}$}\noLine
\UIC{$\varphi_{49} \to r \in \mathbb{R} \to \lvert {{r}^{1_{\omega}}} \rvert = \lvert r \rvert$}
\BIC{$r \in \mathbb{R} \to \lvert {{r}^{1_{\omega}}} \rvert = \lvert r \rvert$}
\end{ptree}
\noindent$\Pi_{44}$:
\begin{ptree}
\AXC{$\stackrel{\Pi_{43}}{\vdots}$}\noLine
\UIC{$r \in \mathbb{R} \to \lvert {{r}^{1_{\omega}}} \rvert = \lvert r \rvert$}
\AXC{$(r \in \mathbb{R} \to \lvert {{r}^{1_{\omega}}} \rvert = \lvert r \rvert) \to \varphi_{141}$\lbl{A1}}
\BIC{$\varphi_{141}$}
\end{ptree}
\noindent$\Pi_{45}$:
\begin{ptree}
\AXC{$\stackrel{\Pi_{37}}{\vdots}$}\noLine
\UIC{$\varphi_{67}$}
\AXC{$\stackrel{\Pi_{5}}{\vdots}$}\noLine
\UIC{$\forall x\,\forall y\,(x = y \to y = x)$}
\AXC{$\varphi_{20}$\lbl{A2}}
\BIC{$\forall y\,({{\lvert r \rvert}^{1_{\omega}}} = y \to y = {{\lvert r \rvert}^{1_{\omega}}})$}
\AXC{$\varphi_{80}$\lbl{A2}}
\BIC{${{\lvert r \rvert}^{1_{\omega}}} = \lvert r \rvert \to \lvert r \rvert = {{\lvert r \rvert}^{1_{\omega}}}$}
\AXC{$\stackrel{\Pi_{40}}{\vdots}$}\noLine
\UIC{$\varphi_{65}$}
\BIC{$r \in \mathbb{R} \to \lvert r \rvert = {{\lvert r \rvert}^{1_{\omega}}}$}
\AXC{$\stackrel{\Pi_{44}}{\vdots}$}\noLine
\UIC{$\varphi_{141}$}
\BIC{$\varphi_{67} \to r \in \mathbb{R} \to \lvert {{r}^{1_{\omega}}} \rvert = {{\lvert r \rvert}^{1_{\omega}}}$}
\BIC{$r \in \mathbb{R} \to \lvert {{r}^{1_{\omega}}} \rvert = {{\lvert r \rvert}^{1_{\omega}}}$}
\end{ptree}
\noindent$\Pi_{46}$:
\begin{ptree}
\AXC{$y = z \to x = y \to x = z$\lbl{A6}}
\AXC{$(y = z \to x = y \to x = z) \to x = y \to y = z \to x = z$\lbl{A1}}
\BIC{$x = y \to y = z \to x = z$}
\UIC{$\forall z\,(x = y \to y = z \to x = z)$}
\UIC{$\forall y\,\forall z\,(x = y \to y = z \to x = z)$}
\UIC{$\forall x\,\forall y\,\forall z\,(x = y \to y = z \to x = z)$}
\AXC{$\varphi_{8}$\lbl{A2}}
\BIC{$\forall y\,\forall z\,(S(n) = y \to y = z \to S(n) = z)$}
\AXC{$\varphi_{24}$\lbl{A2}}
\BIC{$\forall z\,(S(n) = S(\emptyset) \to S(\emptyset) = z \to S(n) = z)$}
\end{ptree}
\noindent$\Pi_{47}$:
\begin{ptree}
\AXC{$\stackrel{\Pi_{46}}{\vdots}$}\noLine
\UIC{$\forall z\,(S(n) = S(\emptyset) \to S(\emptyset) = z \to S(n) = z)$}
\AXC{$\forall z\,(S(n) = S(\emptyset) \to S(\emptyset) = z \to S(n) = z) \to \varphi_{4}$\lbl{A2}}
\BIC{$\varphi_{4}$}
\end{ptree}
\noindent$\Pi_{48}$:
\begin{ptree}
\AXC{$\stackrel{\Pi_{5}}{\vdots}$}\noLine
\UIC{$\forall x\,\forall y\,(x = y \to y = x)$}
\AXC{$\varphi_{2}$\lbl{A2}}
\BIC{$\forall y\,(1_{\omega} = y \to y = 1_{\omega})$}
\AXC{$\varphi_{10}$\lbl{A2}}
\BIC{$1_{\omega} = S(\emptyset) \to S(\emptyset) = 1_{\omega}$}
\AXC{$1_{\omega} = S(\emptyset)$\lbl{Def$1_\omega$}}
\AXC{$1_{\omega} = S(\emptyset) \to \varphi_{9}$\lbl{A1}}
\BIC{$\varphi_{9}$}
\BIC{$S(\emptyset) = 1_{\omega}$}
\end{ptree}
\noindent$\Pi_{49}$:
\begin{ptree}
\AXC{$x = x$\lbl{A5}}
\UIC{$\forall x\,(x = x)$}
\AXC{$\forall x\,(x = x) \to S(n) = S(n)$\lbl{A2}}
\BIC{$S(n) = S(n)$}
\AXC{$r \in \mathbb{R} \to \varphi_{97}$\lbl{A1}}
\AXC{$(r \in \mathbb{R} \to \varphi_{97}) \to S(n) = S(n) \to \varphi_{138}$\lbl{A1}}
\BIC{$S(n) = S(n) \to \varphi_{138}$}
\BIC{$\varphi_{138}$}
\end{ptree}
\noindent$\Pi_{50}$:
\begin{ptree}
\AXC{$z = w \to S(n) = S(z) \to S(n) = S(w)$\lbl{A6}}
\UIC{$\forall w\,(z = w \to S(n) = S(z) \to S(n) = S(w))$}
\UIC{$\forall z\,\forall w\,(z = w \to S(n) = S(z) \to S(n) = S(w))$}
\AXC{$\varphi_{18}$\lbl{A2}}
\BIC{$\forall w\,(n = w \to S(n) = S(n) \to S(n) = S(w))$}
\AXC{$\varphi_{25}$\lbl{A2}}
\BIC{$n = \emptyset \to S(n) = S(n) \to S(n) = S(\emptyset)$}
\AXC{$\stackrel{\Pi_{49}}{\vdots}$}\noLine
\UIC{$\varphi_{138}$}
\BIC{$r \in \mathbb{R} \to \varphi_{104}$}
\end{ptree}
\noindent$\Pi_{51}$:
\begin{ptree}
\AXC{$\stackrel{\Pi_{47}}{\vdots}$}\noLine
\UIC{$\varphi_{4}$}
\AXC{$\stackrel{\Pi_{48}}{\vdots}$}\noLine
\UIC{$S(\emptyset) = 1_{\omega}$}
\AXC{$\stackrel{\Pi_{50}}{\vdots}$}\noLine
\UIC{$r \in \mathbb{R} \to \varphi_{104}$}
\AXC{$(r \in \mathbb{R} \to \varphi_{104}) \to S(\emptyset) = 1_{\omega} \to \varphi_{4} \to r \in \mathbb{R} \to \varphi_{94}$\lbl{A1}}
\BIC{$S(\emptyset) = 1_{\omega} \to \varphi_{4} \to r \in \mathbb{R} \to \varphi_{94}$}
\BIC{$\varphi_{4} \to r \in \mathbb{R} \to \varphi_{94}$}
\BIC{$r \in \mathbb{R} \to \varphi_{94}$}
\AXC{$\varphi_{151}$\lbl{A1}}
\BIC{$(S(n) = 1_{\omega} \to 1_{\omega} = S(n)) \to r \in \mathbb{R} \to \varphi_{95}$}
\end{ptree}
\noindent$\Pi_{52}$:
\begin{ptree}
\AXC{$\stackrel{\Pi_{5}}{\vdots}$}\noLine
\UIC{$\forall x\,\forall y\,(x = y \to y = x)$}
\AXC{$\varphi_{1}$\lbl{A2}}
\BIC{$\forall y\,(S(n) = y \to y = S(n))$}
\AXC{$\varphi_{3}$\lbl{A2}}
\BIC{$S(n) = 1_{\omega} \to 1_{\omega} = S(n)$}
\AXC{$\stackrel{\Pi_{51}}{\vdots}$}\noLine
\UIC{$(S(n) = 1_{\omega} \to 1_{\omega} = S(n)) \to r \in \mathbb{R} \to \varphi_{95}$}
\BIC{$r \in \mathbb{R} \to \varphi_{95}$}
\AXC{$\varphi_{156}$\lbl{A1}}
\BIC{$(r \in \mathbb{R} \to \lvert {{r}^{1_{\omega}}} \rvert = {{\lvert r \rvert}^{1_{\omega}}}) \to \varphi_{60} \to r \in \mathbb{R} \to \varphi_{117}$}
\end{ptree}
\noindent$\Pi_{53}$:
\begin{ptree}
\AXC{$\stackrel{\Pi_{45}}{\vdots}$}\noLine
\UIC{$r \in \mathbb{R} \to \lvert {{r}^{1_{\omega}}} \rvert = {{\lvert r \rvert}^{1_{\omega}}}$}
\AXC{$\stackrel{\Pi_{52}}{\vdots}$}\noLine
\UIC{$(r \in \mathbb{R} \to \lvert {{r}^{1_{\omega}}} \rvert = {{\lvert r \rvert}^{1_{\omega}}}) \to \varphi_{60} \to r \in \mathbb{R} \to \varphi_{117}$}
\BIC{$\varphi_{60} \to r \in \mathbb{R} \to \varphi_{117}$}
\end{ptree}
\noindent$\Pi_{54}$:
\begin{ptree}
\AXC{$z = w \to \lvert {{r}^{z}} \rvert = {{\lvert r \rvert}^{z}} \to \lvert {{r}^{w}} \rvert = {{\lvert r \rvert}^{w}}$\lbl{A6}}
\UIC{$\forall w\,(z = w \to \lvert {{r}^{z}} \rvert = {{\lvert r \rvert}^{z}} \to \lvert {{r}^{w}} \rvert = {{\lvert r \rvert}^{w}})$}
\UIC{$\forall z\,\forall w\,(z = w \to \lvert {{r}^{z}} \rvert = {{\lvert r \rvert}^{z}} \to \lvert {{r}^{w}} \rvert = {{\lvert r \rvert}^{w}})$}
\AXC{$\varphi_{133}$\lbl{A2}}
\BIC{$\varphi_{48}$}
\AXC{$\varphi_{48} \to \varphi_{60}$\lbl{A2}}
\BIC{$\varphi_{60}$}
\AXC{$\stackrel{\Pi_{53}}{\vdots}$}\noLine
\UIC{$\varphi_{60} \to r \in \mathbb{R} \to \varphi_{117}$}
\BIC{$r \in \mathbb{R} \to \varphi_{117}$}
\end{ptree}
\noindent$\Pi_{55}$:
\begin{ptree}
\AXC{$\forall x\,(\mathrm{Ind}(x) \leftrightarrow \emptyset \in x \land \forall y \in x\,(S(y) \in x))$\lbl{DefInd}}
\AXC{$\varphi_{87}$\lbl{A2}}
\BIC{$\mathrm{Ind}(\omega) \leftrightarrow \emptyset \in \omega \land \forall y \in \omega\,(S(y) \in \omega)$}
\AXC{$\mathrm{Ind}(\omega)$\lbl{補題(ωは帰納的)}}
\AXC{$\mathrm{Ind}(\omega) \to \varphi_{143}$\lbl{A1}}
\BIC{$\varphi_{143}$}
\BIC{$r \in \mathbb{R} \to \varphi_{101}$}
\end{ptree}
\noindent$\Pi_{56}$:
\begin{ptree}
\AXC{$\stackrel{\Pi_{55}}{\vdots}$}\noLine
\UIC{$r \in \mathbb{R} \to \varphi_{101}$}
\AXC{$(r \in \mathbb{R} \to \varphi_{101}) \to r \in \mathbb{R} \to \varphi_{89}$\lbl{A1}}
\BIC{$r \in \mathbb{R} \to \varphi_{89}$}
\AXC{$(r \in \mathbb{R} \to \varphi_{89}) \to \varphi_{140}$\lbl{A1}}
\BIC{$\varphi_{140}$}
\end{ptree}
\noindent$\Pi_{57}$:
\begin{ptree}
\AXC{$\stackrel{\Pi_{9}}{\vdots}$}\noLine
\UIC{$r \in \mathbb{R} \to \varphi_{71}$}
\AXC{$\forall y \in \omega\,(S(y) \in \omega) \to n \in \omega \to S(n) \in \omega$\lbl{A2}}
\AXC{$\stackrel{\Pi_{56}}{\vdots}$}\noLine
\UIC{$\varphi_{140}$}
\BIC{$r \in \mathbb{R} \to \varphi_{86}$}
\AXC{$(r \in \mathbb{R} \to \varphi_{86}) \to (r \in \mathbb{R} \to \varphi_{71}) \to r \in \mathbb{R} \to \varphi_{74}$\lbl{A1}}
\BIC{$(r \in \mathbb{R} \to \varphi_{71}) \to r \in \mathbb{R} \to \varphi_{74}$}
\BIC{$r \in \mathbb{R} \to \varphi_{74}$}
\end{ptree}
\noindent$\Pi_{58}$:
\begin{ptree}
\AXC{$\stackrel{\Pi_{57}}{\vdots}$}\noLine
\UIC{$r \in \mathbb{R} \to \varphi_{74}$}
\AXC{$\varphi_{122}$\lbl{A2}}
\AXC{$r \in \mathbb{R} \to \varphi_{115}$\lbl{A1}}
\AXC{$(r \in \mathbb{R} \to \varphi_{115}) \to \varphi_{122} \to r \in \mathbb{R} \to \varphi_{122}$\lbl{A1}}
\BIC{$\varphi_{122} \to r \in \mathbb{R} \to \varphi_{122}$}
\BIC{$r \in \mathbb{R} \to \varphi_{122}$}
\AXC{$\varphi_{163}$\lbl{A1}}
\BIC{$(r \in \mathbb{R} \to \varphi_{74}) \to (r \in \mathbb{R} \to \varphi_{110}) \to r \in \mathbb{R} \to \varphi_{72}$}
\BIC{$(r \in \mathbb{R} \to \varphi_{110}) \to r \in \mathbb{R} \to \varphi_{72}$}
\end{ptree}
\noindent$\Pi_{59}$:
\begin{ptree}
\AXC{$\stackrel{\Pi_{36}}{\vdots}$}\noLine
\UIC{$r \in \mathbb{R} \to \varphi_{118}$}
\AXC{$\stackrel{\Pi_{54}}{\vdots}$}\noLine
\UIC{$r \in \mathbb{R} \to \varphi_{117}$}
\AXC{$(r \in \mathbb{R} \to \varphi_{117}) \to (r \in \mathbb{R} \to \varphi_{118}) \to r \in \mathbb{R} \to \varphi_{110}$\lbl{A1}}
\BIC{$(r \in \mathbb{R} \to \varphi_{118}) \to r \in \mathbb{R} \to \varphi_{110}$}
\BIC{$r \in \mathbb{R} \to \varphi_{110}$}
\AXC{$\stackrel{\Pi_{58}}{\vdots}$}\noLine
\UIC{$(r \in \mathbb{R} \to \varphi_{110}) \to r \in \mathbb{R} \to \varphi_{72}$}
\BIC{$r \in \mathbb{R} \to \varphi_{72}$}
\end{ptree}
\noindent$\Pi_{60}$:
\begin{ptree}
\AXC{$\varphi_{41} \to \forall n\,(\varphi_{41})$\lbl{A4}}
\AXC{$\varphi_{146}$\lbl{A3}}
\AXC{$\stackrel{\Pi_{59}}{\vdots}$}\noLine
\UIC{$r \in \mathbb{R} \to \varphi_{72}$}
\AXC{$(r \in \mathbb{R} \to \varphi_{72}) \to \varphi_{41} \to n \in B_{26} \to S(n) \in B_{26}$\lbl{A1}}
\BIC{$\varphi_{41} \to n \in B_{26} \to S(n) \in B_{26}$}
\UIC{$\forall n\,(\varphi_{41} \to n \in B_{26} \to S(n) \in B_{26})$}
\AXC{$\varphi_{164}$\lbl{A1}}
\BIC{$\varphi_{146} \to (\varphi_{41} \to \forall n\,(\varphi_{41})) \to r \in \mathbb{R} \to \varphi_{76}$}
\BIC{$(\varphi_{41} \to \forall n\,(\varphi_{41})) \to r \in \mathbb{R} \to \varphi_{76}$}
\BIC{$r \in \mathbb{R} \to \varphi_{76}$}
\end{ptree}
\noindent$\Pi_{61}$:
\begin{ptree}
\AXC{$\forall n\,(n \in \mathbb{N} \leftrightarrow n \in \omega \land n \neq \emptyset)$\lbl{Def$\mathbb{N}$}}
\AXC{$\varphi_{68}$\lbl{A2}}
\BIC{$\emptyset \in \mathbb{N} \leftrightarrow \emptyset \in \omega \land \emptyset \neq \emptyset$}
\AXC{$(\emptyset \in \mathbb{N} \leftrightarrow \emptyset \in \omega \land \emptyset \neq \emptyset) \to \varphi_{69}$\lbl{A1}}
\BIC{$\varphi_{69}$}
\end{ptree}
\noindent$\Pi_{62}$:
\begin{ptree}
\AXC{$\varphi_{137}$\lbl{A2}}
\AXC{$r \in \mathbb{R} \to \varphi_{115}$\lbl{A1}}
\AXC{$(r \in \mathbb{R} \to \varphi_{115}) \to \varphi_{137} \to r \in \mathbb{R} \to \varphi_{137}$\lbl{A1}}
\BIC{$\varphi_{137} \to r \in \mathbb{R} \to \varphi_{137}$}
\BIC{$r \in \mathbb{R} \to \varphi_{137}$}
\AXC{$(r \in \mathbb{R} \to \varphi_{137}) \to (r \in \mathbb{R} \to \varphi_{101}) \to \varphi_{139}$\lbl{A1}}
\BIC{$(r \in \mathbb{R} \to \varphi_{101}) \to \varphi_{139}$}
\end{ptree}
\noindent$\Pi_{63}$:
\begin{ptree}
\AXC{$x = x$\lbl{A5}}
\UIC{$\forall x\,(x = x)$}
\AXC{$\forall x\,(x = x) \to \emptyset = \emptyset$\lbl{A2}}
\BIC{$\emptyset = \emptyset$}
\AXC{$\stackrel{\Pi_{61}}{\vdots}$}\noLine
\UIC{$\varphi_{69}$}
\BIC{$r \in \mathbb{R} \to \emptyset \in \mathbb{N} \to \lvert {{r}^{\emptyset}} \rvert = {{\lvert r \rvert}^{\emptyset}}$}
\AXC{$\stackrel{\Pi_{55}}{\vdots}$}\noLine
\UIC{$r \in \mathbb{R} \to \varphi_{101}$}
\AXC{$\stackrel{\Pi_{62}}{\vdots}$}\noLine
\UIC{$(r \in \mathbb{R} \to \varphi_{101}) \to \varphi_{139}$}
\BIC{$\varphi_{139}$}
\BIC{$r \in \mathbb{R} \to \varphi_{56}$}
\end{ptree}
\noindent$\Pi_{64}$:
\begin{ptree}
\AXC{$\forall x\,(\mathrm{Ind}(x) \leftrightarrow \emptyset \in x \land \forall y \in x\,(S(y) \in x))$\lbl{DefInd}}
\AXC{$\forall x\,(\mathrm{Ind}(x) \leftrightarrow \emptyset \in x \land \forall y \in x\,(S(y) \in x)) \to \varphi_{13}$\lbl{A2}}
\BIC{$\varphi_{13}$}
\end{ptree}
\noindent$\Pi_{65}$:
\begin{ptree}
\AXC{$\stackrel{\Pi_{60}}{\vdots}$}\noLine
\UIC{$r \in \mathbb{R} \to \varphi_{76}$}
\AXC{$\stackrel{\Pi_{63}}{\vdots}$}\noLine
\UIC{$r \in \mathbb{R} \to \varphi_{56}$}
\AXC{$\stackrel{\Pi_{64}}{\vdots}$}\noLine
\UIC{$\varphi_{13}$}
\AXC{$\varphi_{13} \to (r \in \mathbb{R} \to \varphi_{56}) \to (r \in \mathbb{R} \to \varphi_{76}) \to \varphi_{12} \to \varphi_{11} \to r \in \mathbb{R} \to \varphi_{39}$\lbl{A1}}
\BIC{$(r \in \mathbb{R} \to \varphi_{56}) \to (r \in \mathbb{R} \to \varphi_{76}) \to \varphi_{12} \to \varphi_{11} \to r \in \mathbb{R} \to \varphi_{39}$}
\BIC{$(r \in \mathbb{R} \to \varphi_{76}) \to \varphi_{12} \to \varphi_{11} \to r \in \mathbb{R} \to \varphi_{39}$}
\BIC{$\varphi_{12} \to \varphi_{11} \to r \in \mathbb{R} \to \varphi_{39}$}
\end{ptree}
\noindent$\Pi_{66}$:
\begin{ptree}
\AXC{$\varphi_{15}$\lbl{A4}}
\AXC{$\forall n\,(\varphi_{7}) \to \varphi_{17}$\lbl{A3}}
\AXC{$\varphi_{7}$\lbl{A2}}
\UIC{$\forall n\,(\varphi_{7})$}
\AXC{$\forall n\,(\varphi_{7}) \to (\forall n\,(\varphi_{7}) \to \varphi_{17}) \to \varphi_{15} \to \varphi_{12}$\lbl{A1}}
\BIC{$(\forall n\,(\varphi_{7}) \to \varphi_{17}) \to \varphi_{15} \to \varphi_{12}$}
\BIC{$\varphi_{15} \to \varphi_{12}$}
\BIC{$\varphi_{12}$}
\AXC{$\stackrel{\Pi_{65}}{\vdots}$}\noLine
\UIC{$\varphi_{12} \to \varphi_{11} \to r \in \mathbb{R} \to \varphi_{39}$}
\BIC{$\varphi_{11} \to r \in \mathbb{R} \to \varphi_{39}$}
\end{ptree}
\noindent$\Pi_{67}$:
\begin{ptree}
\AXC{$\varphi_{14}$\lbl{A4}}
\AXC{$\forall y\,(\varphi_{6}) \to \varphi_{16}$\lbl{A3}}
\AXC{$\varphi_{6}$\lbl{A2}}
\UIC{$\forall y\,(\varphi_{6})$}
\AXC{$\forall y\,(\varphi_{6}) \to (\forall y\,(\varphi_{6}) \to \varphi_{16}) \to \varphi_{14} \to \varphi_{11}$\lbl{A1}}
\BIC{$(\forall y\,(\varphi_{6}) \to \varphi_{16}) \to \varphi_{14} \to \varphi_{11}$}
\BIC{$\varphi_{14} \to \varphi_{11}$}
\BIC{$\varphi_{11}$}
\AXC{$\stackrel{\Pi_{66}}{\vdots}$}\noLine
\UIC{$\varphi_{11} \to r \in \mathbb{R} \to \varphi_{39}$}
\BIC{$r \in \mathbb{R} \to \varphi_{39}$}
\end{ptree}
\noindent$\Pi_{68}$:
\begin{ptree}
\AXC{$\forall n\,(n \in \omega \leftrightarrow \forall x\,(\mathrm{Ind}(x) \to n \in x))$\lbl{Def$\omega$}}
\AXC{$\varphi_{23}$\lbl{A2}}
\BIC{$n \in \omega \leftrightarrow \forall x\,(\mathrm{Ind}(x) \to n \in x)$}
\AXC{$r \in \mathbb{R} \to \varphi_{73}$\lbl{A1}}
\AXC{$\varphi_{148}$\lbl{A1}}
\BIC{$(n \in \omega \leftrightarrow \forall x\,(\mathrm{Ind}(x) \to n \in x)) \to r \in \mathbb{R} \to \varphi_{82}$}
\BIC{$r \in \mathbb{R} \to \varphi_{82}$}
\end{ptree}
\noindent$\Pi_{69}$:
\begin{ptree}
\AXC{$\forall x\,(\mathrm{Ind}(x) \to n \in x) \to \mathrm{Ind}(B_{26}) \to n \in B_{26}$\lbl{A2}}
\AXC{$\stackrel{\Pi_{68}}{\vdots}$}\noLine
\UIC{$r \in \mathbb{R} \to \varphi_{82}$}
\AXC{$(r \in \mathbb{R} \to \varphi_{82}) \to \varphi_{130}$\lbl{A1}}
\BIC{$\varphi_{130}$}
\BIC{$r \in \mathbb{R} \to \varphi_{85}$}
\AXC{$\varphi_{161}$\lbl{A1}}
\BIC{$(r \in \mathbb{R} \to \varphi_{39}) \to (r \in \mathbb{R} \to \varphi_{111}) \to r \in \mathbb{R} \to \varphi_{106}$}
\end{ptree}
\noindent$\Pi_{70}$:
\begin{ptree}
\AXC{$\stackrel{\Pi_{2}}{\vdots}$}\noLine
\UIC{$r \in \mathbb{R} \to \varphi_{111}$}
\AXC{$\stackrel{\Pi_{67}}{\vdots}$}\noLine
\UIC{$r \in \mathbb{R} \to \varphi_{39}$}
\AXC{$\stackrel{\Pi_{69}}{\vdots}$}\noLine
\UIC{$(r \in \mathbb{R} \to \varphi_{39}) \to (r \in \mathbb{R} \to \varphi_{111}) \to r \in \mathbb{R} \to \varphi_{106}$}
\BIC{$(r \in \mathbb{R} \to \varphi_{111}) \to r \in \mathbb{R} \to \varphi_{106}$}
\BIC{$r \in \mathbb{R} \to \varphi_{106}$}
\end{ptree}
\noindent$\Pi_{71}$:
\begin{ptree}
\AXC{$\varphi_{41} \to \forall n\,(\varphi_{41})$\lbl{A4}}
\AXC{$\varphi_{153}$\lbl{A3}}
\AXC{$\stackrel{\Pi_{70}}{\vdots}$}\noLine
\UIC{$r \in \mathbb{R} \to \varphi_{106}$}
\AXC{$(r \in \mathbb{R} \to \varphi_{106}) \to \varphi_{41} \to n \in \omega \to n \in \mathbb{N} \to \lvert {{r}^{n}} \rvert = {{\lvert r \rvert}^{n}}$\lbl{A1}}
\BIC{$\varphi_{41} \to n \in \omega \to n \in \mathbb{N} \to \lvert {{r}^{n}} \rvert = {{\lvert r \rvert}^{n}}$}
\UIC{$\forall n\,(\varphi_{41} \to n \in \omega \to n \in \mathbb{N} \to \lvert {{r}^{n}} \rvert = {{\lvert r \rvert}^{n}})$}
\AXC{$\varphi_{165}$\lbl{A1}}
\BIC{$\varphi_{153} \to (\varphi_{41} \to \forall n\,(\varphi_{41})) \to r \in \mathbb{R} \to \varphi_{109}$}
\BIC{$(\varphi_{41} \to \forall n\,(\varphi_{41})) \to r \in \mathbb{R} \to \varphi_{109}$}
\BIC{$r \in \mathbb{R} \to \varphi_{109}$}
\end{ptree}
\noindent$\Pi_{72}$:
\begin{ptree}
\AXC{$\varphi_{53}$\lbl{Sep}}
\AXC{$\varphi_{53} \to \varphi_{33}$\lbl{A2}}
\BIC{$\varphi_{33}$}
\AXC{$\varphi_{33} \to \varphi_{28}$\lbl{A2}}
\BIC{$\varphi_{28}$}
\AXC{$\varphi_{28} \to r \in \mathbb{R} \to \varphi_{28}$\lbl{A1}}
\BIC{$r \in \mathbb{R} \to \varphi_{28}$}
\AXC{$(r \in \mathbb{R} \to \varphi_{28}) \to \varphi_{119} \to (\varphi_{119} \to \varphi_{32} \to \forall B_{26}\,(\varphi_{26})) \to \varphi_{147}$\lbl{A1}}
\BIC{$\varphi_{119} \to (\varphi_{119} \to \varphi_{32} \to \forall B_{26}\,(\varphi_{26})) \to \varphi_{147}$}
\end{ptree}
\noindent$\Pi_{73}$:
\begin{ptree}
\AXC{$\varphi_{119} \to \varphi_{32} \to \forall B_{26}\,(\varphi_{26})$\lbl{A3}}
\AXC{$\stackrel{\Pi_{71}}{\vdots}$}\noLine
\UIC{$r \in \mathbb{R} \to \varphi_{109}$}
\AXC{$\varphi_{152}$\lbl{A1}}
\BIC{$r \in \mathbb{R} \land \neg \forall n \in \omega\,(n \in \mathbb{N} \to \lvert {{r}^{n}} \rvert = {{\lvert r \rvert}^{n}}) \to \varphi_{26}$}
\UIC{$\varphi_{119}$}
\AXC{$\stackrel{\Pi_{72}}{\vdots}$}\noLine
\UIC{$\varphi_{119} \to (\varphi_{119} \to \varphi_{32} \to \forall B_{26}\,(\varphi_{26})) \to \varphi_{147}$}
\BIC{$(\varphi_{119} \to \varphi_{32} \to \forall B_{26}\,(\varphi_{26})) \to \varphi_{147}$}
\BIC{$\varphi_{147}$}
\end{ptree}
\noindent$\Pi_{74}$:
\begin{ptree}
\AXC{$\varphi_{92}$\lbl{A2}}
\AXC{$r \in \mathbb{R} \land \neg \forall n \in \omega\,(n \in \mathbb{N} \to \lvert {{r}^{n}} \rvert = {{\lvert r \rvert}^{n}}) \to \varphi_{32}$\lbl{A4}}
\AXC{$\stackrel{\Pi_{73}}{\vdots}$}\noLine
\UIC{$\varphi_{147}$}
\BIC{$r \in \mathbb{R} \to \forall n \in \omega\,(n \in \mathbb{N} \to \lvert {{r}^{n}} \rvert = {{\lvert r \rvert}^{n}})$}
\AXC{$\varphi_{149}$\lbl{A1}}
\BIC{$\varphi_{92} \to r \in \mathbb{R} \to n \in \omega \to n \in \mathbb{N} \to \lvert {{r}^{n}} \rvert = {{\lvert r \rvert}^{n}}$}
\BIC{$r \in \mathbb{R} \to n \in \omega \to n \in \mathbb{N} \to \lvert {{r}^{n}} \rvert = {{\lvert r \rvert}^{n}}$}
\AXC{$\varphi_{144}$\lbl{A1}}
\BIC{$(r \in \mathbb{R} \to n \in \mathbb{N} \to n \in \omega) \to \varphi_{61}$}
\end{ptree}
\noindent$\Pi_{75}$:
\begin{ptree}
\AXC{$r \in \mathbb{R} \to n \in \mathbb{N} \to n \in \mathbb{N}$\lbl{A1}}
\AXC{$\stackrel{\Pi_{1}}{\vdots}$}\noLine
\UIC{$r \in \mathbb{R} \to n \in \mathbb{N} \to n \in \omega$}
\AXC{$\stackrel{\Pi_{74}}{\vdots}$}\noLine
\UIC{$(r \in \mathbb{R} \to n \in \mathbb{N} \to n \in \omega) \to \varphi_{61}$}
\BIC{$\varphi_{61}$}
\BIC{$r \in \mathbb{R} \to n \in \mathbb{N} \to \lvert {{r}^{n}} \rvert = {{\lvert r \rvert}^{n}}$}
\UIC{$\forall n\,(r \in \mathbb{R} \to n \in \mathbb{N} \to \lvert {{r}^{n}} \rvert = {{\lvert r \rvert}^{n}})$}
\end{ptree}
\noindent$\Pi_{76}$:
\begin{ptree}
\AXC{$\varphi_{102}$\lbl{A3}}
\AXC{$\stackrel{\Pi_{75}}{\vdots}$}\noLine
\UIC{$\forall n\,(r \in \mathbb{R} \to n \in \mathbb{N} \to \lvert {{r}^{n}} \rvert = {{\lvert r \rvert}^{n}})$}
\AXC{$\forall n\,(r \in \mathbb{R} \to n \in \mathbb{N} \to \lvert {{r}^{n}} \rvert = {{\lvert r \rvert}^{n}}) \to \varphi_{102} \to \varphi_{51}$\lbl{A1}}
\BIC{$\varphi_{102} \to \varphi_{51}$}
\BIC{$\varphi_{51}$}
\end{ptree}
\noindent 主木:
\begin{ptree}
\AXC{$r \in \mathbb{R} \to \forall n\,(r \in \mathbb{R})$\lbl{A4}}
\AXC{$\stackrel{\Pi_{76}}{\vdots}$}\noLine
\UIC{$\varphi_{51}$}
\BIC{$r \in \mathbb{R} \to \forall n \in \mathbb{N}\,(\lvert {{r}^{n}} \rvert = {{\lvert r \rvert}^{n}})$}
\UIC{$\forall r \in \mathbb{R}\,\forall n \in \mathbb{N}\,(\lvert {{r}^{n}} \rvert = {{\lvert r \rvert}^{n}})$}
\end{ptree}
\begin{abbrevlist}
\item $\varphi_{1} :\equiv \forall x\,\forall y\,(x = y \to y = x) \to \forall y\,(S(n) = y \to y = S(n))$
\item $\varphi_{2} :\equiv \forall x\,\forall y\,(x = y \to y = x) \to \forall y\,(1_{\omega} = y \to y = 1_{\omega})$
\item $\varphi_{3} :\equiv \forall y\,(S(n) = y \to y = S(n)) \to S(n) = 1_{\omega} \to 1_{\omega} = S(n)$
\item $\varphi_{4} :\equiv S(n) = S(\emptyset) \to S(\emptyset) = 1_{\omega} \to S(n) = 1_{\omega}$
\item $\varphi_{5} :\equiv n \in B_{26} \to S(n) \in \mathbb{N} \to n \neq \emptyset \to \mathrm{Seq}({(r)^{\bullet}})$
\item $\varphi_{6} :\equiv \forall n \in B_{26}\,(S(n) \in B_{26}) \to y \in B_{26} \to S(y) \in B_{26}$
\item $\varphi_{7} :\equiv \forall y \in B_{26}\,(S(y) \in B_{26}) \to n \in B_{26} \to S(n) \in B_{26}$
\item $\varphi_{8} :\equiv \forall x\,\forall y\,\forall z\,(x = y \to y = z \to x = z) \to \forall y\,\forall z\,(S(n) = y \to y = z \to S(n) = z)$
\item $\varphi_{9} :\equiv (1_{\omega} = S(\emptyset) \to S(\emptyset) = 1_{\omega}) \to S(\emptyset) = 1_{\omega}$
\item $\varphi_{10} :\equiv \forall y\,(1_{\omega} = y \to y = 1_{\omega}) \to 1_{\omega} = S(\emptyset) \to S(\emptyset) = 1_{\omega}$
\item $\varphi_{11} :\equiv \forall n \in B_{26}\,(S(n) \in B_{26}) \to \forall y \in B_{26}\,(S(y) \in B_{26})$
\item $\varphi_{12} :\equiv \forall y \in B_{26}\,(S(y) \in B_{26}) \to \forall n \in B_{26}\,(S(n) \in B_{26})$
\item $\varphi_{13} :\equiv \mathrm{Ind}(B_{26}) \leftrightarrow \emptyset \in B_{26} \land \forall y \in B_{26}\,(S(y) \in B_{26})$
\item $\varphi_{14} :\equiv \forall n \in B_{26}\,(S(n) \in B_{26}) \to \forall y\,\forall n \in B_{26}\,(S(n) \in B_{26})$
\item $\varphi_{15} :\equiv \forall y \in B_{26}\,(S(y) \in B_{26}) \to \forall n\,\forall y \in B_{26}\,(S(y) \in B_{26})$
\item $\varphi_{16} :\equiv \forall y\,\forall n \in B_{26}\,(S(n) \in B_{26}) \to \forall y \in B_{26}\,(S(y) \in B_{26})$
\item $\varphi_{17} :\equiv \forall n\,\forall y \in B_{26}\,(S(y) \in B_{26}) \to \forall n \in B_{26}\,(S(n) \in B_{26})$
\item $\varphi_{18} :\equiv \forall z\,\forall w\,(z = w \to S(n) = S(z) \to S(n) = S(w)) \to \forall w\,(n = w \to S(n) = S(n) \to S(n) = S(w))$
\item $\varphi_{19} :\equiv \forall x\,\forall y\,(x = y \to y = x) \to \forall y\,({{\lvert r \rvert}^{S(n)}} = y \to y = {{\lvert r \rvert}^{S(n)}})$
\item $\varphi_{20} :\equiv \forall x\,\forall y\,(x = y \to y = x) \to \forall y\,({{\lvert r \rvert}^{1_{\omega}}} = y \to y = {{\lvert r \rvert}^{1_{\omega}}})$
\item $\varphi_{21} :\equiv \forall z\,\forall w\,(z = w \to \lvert {{r}^{S(n)}} \rvert = \lvert z \rvert \to \lvert {{r}^{S(n)}} \rvert = \lvert w \rvert)$
\item $\varphi_{22} :\equiv \forall z\,\forall w\,(z = w \to \lvert {{r}^{1_{\omega}}} \rvert = \lvert z \rvert \to \lvert {{r}^{1_{\omega}}} \rvert = \lvert w \rvert)$
\item $\varphi_{23} :\equiv \forall n\,(n \in \omega \leftrightarrow \forall x\,(\mathrm{Ind}(x) \to n \in x)) \to (n \in \omega \leftrightarrow \forall x\,(\mathrm{Ind}(x) \to n \in x))$
\item $\varphi_{24} :\equiv \forall y\,\forall z\,(S(n) = y \to y = z \to S(n) = z) \to \forall z\,(S(n) = S(\emptyset) \to S(\emptyset) = z \to S(n) = z)$
\item $\varphi_{25} :\equiv \forall w\,(n = w \to S(n) = S(n) \to S(n) = S(w)) \to n = \emptyset \to S(n) = S(n) \to S(n) = S(\emptyset)$
\item $\varphi_{26} :\equiv \neg \forall n\,(n \in B_{26} \leftrightarrow n \in \omega \land (n \in \mathbb{N} \to \lvert {{r}^{n}} \rvert = {{\lvert r \rvert}^{n}}))$
\item $\varphi_{27} :\equiv \forall n\,(n \in \mathbb{N} \leftrightarrow n \in \omega \land n \neq \emptyset) \to (n \in \mathbb{N} \leftrightarrow n \in \omega \land n \neq \emptyset)$
\item $\varphi_{28} :\equiv \exists B_{26}\,\forall n\,(n \in B_{26} \leftrightarrow n \in \omega \land (n \in \mathbb{N} \to \lvert {{r}^{n}} \rvert = {{\lvert r \rvert}^{n}}))$
\item $\varphi_{29} :\equiv \mathrm{Seq}({(r)^{\bullet}}) \land {{r}^{1_{\omega}}} = r \land \forall n \in \mathbb{N}\,({{r}^{S(n)}} = {{r}^{n}} \cdot r)$
\item $\varphi_{30} :\equiv n \in B_{26} \to S(n) \in \mathbb{N} \to n = \emptyset \to \lvert {{r}^{S(n)}} \rvert = {{\lvert r \rvert}^{S(n)}}$
\item $\varphi_{31} :\equiv n \in B_{26} \to S(n) \in \mathbb{N} \to n \neq \emptyset \to \lvert {{r}^{S(n)}} \rvert = {{\lvert r \rvert}^{S(n)}}$
\item $\varphi_{32} :\equiv \forall B_{26}\,(r \in \mathbb{R} \land \neg \forall n \in \omega\,(n \in \mathbb{N} \to \lvert {{r}^{n}} \rvert = {{\lvert r \rvert}^{n}}))$
\item $\varphi_{33} :\equiv \forall A_{27}\,\exists B_{26}\,\forall n\,(n \in B_{26} \leftrightarrow n \in A_{27} \land (n \in \mathbb{N} \to \lvert {{r}^{n}} \rvert = {{\lvert r \rvert}^{n}}))$
\item $\varphi_{34} :\equiv n \in B_{26} \to S(n) \in \mathbb{N} \to n \neq \emptyset \to \lvert {{r}^{S(n)}} \rvert = \lvert {{r}^{n}} \cdot r \rvert$
\item $\varphi_{35} :\equiv \forall x\,\forall y\,\forall z\,(x = y \to y = z \to x = z) \to \forall y\,\forall z\,(\lvert {{r}^{S(n)}} \rvert = y \to y = z \to \lvert {{r}^{S(n)}} \rvert = z)$
\item $\varphi_{36} :\equiv \forall x\,\forall y\,\forall z\,(x = y \to y = z \to x = z) \to \forall y\,\forall z\,(\lvert {{r}^{1_{\omega}}} \rvert = y \to y = z \to \lvert {{r}^{1_{\omega}}} \rvert = z)$
\item $\varphi_{37} :\equiv n \in B_{26} \to S(n) \in \mathbb{N} \to n \neq \emptyset \to n \in \mathbb{N} \to {{r}^{S(n)}} = {{r}^{n}} \cdot r$
\item $\varphi_{38} :\equiv z = w \to \lvert {{r}^{n}} \rvert \cdot \lvert r \rvert = z \cdot \lvert r \rvert \to \lvert {{r}^{n}} \rvert \cdot \lvert r \rvert = w \cdot \lvert r \rvert$
\item $\varphi_{39} :\equiv \forall n\,(n \in B_{26} \leftrightarrow n \in \omega \land (n \in \mathbb{N} \to \lvert {{r}^{n}} \rvert = {{\lvert r \rvert}^{n}})) \to \mathrm{Ind}(B_{26})$
\item $\varphi_{40} :\equiv S(n) \in B_{26} \leftrightarrow S(n) \in \omega \land (S(n) \in \mathbb{N} \to \lvert {{r}^{S(n)}} \rvert = {{\lvert r \rvert}^{S(n)}})$
\item $\varphi_{41} :\equiv r \in \mathbb{R} \land \forall n\,(n \in B_{26} \leftrightarrow n \in \omega \land (n \in \mathbb{N} \to \lvert {{r}^{n}} \rvert = {{\lvert r \rvert}^{n}}))$
\item $\varphi_{42} :\equiv S(n) \in \mathbb{N} \to n \neq \emptyset \to n \in \mathbb{N} \to {{\lvert r \rvert}^{S(n)}} = {{\lvert r \rvert}^{n}} \cdot \lvert r \rvert$
\item $\varphi_{43} :\equiv n \in B_{26} \to S(n) \in \mathbb{N} \to n \neq \emptyset \to {{\lvert r \rvert}^{S(n)}} = {{\lvert r \rvert}^{n}} \cdot \lvert r \rvert$
\item $\varphi_{44} :\equiv n \in B_{26} \to S(n) \in \mathbb{N} \to n \neq \emptyset \to {{\lvert r \rvert}^{n}} \cdot \lvert r \rvert = {{\lvert r \rvert}^{S(n)}}$
\item $\varphi_{45} :\equiv n \in B_{26} \to S(n) \in \mathbb{N} \to n \neq \emptyset \to \forall n \in \mathbb{N}\,({{r}^{S(n)}} = {{r}^{n}} \cdot r)$
\item $\varphi_{46} :\equiv n \in B_{26} \to S(n) \in \mathbb{N} \to n \neq \emptyset \to \lvert {{r}^{S(n)}} \rvert = \lvert {{r}^{n}} \rvert \cdot \lvert r \rvert$
\item $\varphi_{47} :\equiv n \in B_{26} \to S(n) \in \mathbb{N} \to n \neq \emptyset \to \lvert {{r}^{S(n)}} \rvert = {{\lvert r \rvert}^{n}} \cdot \lvert r \rvert$
\item $\varphi_{48} :\equiv \forall w\,(1_{\omega} = w \to \lvert {{r}^{1_{\omega}}} \rvert = {{\lvert r \rvert}^{1_{\omega}}} \to \lvert {{r}^{w}} \rvert = {{\lvert r \rvert}^{w}})$
\item $\varphi_{49} :\equiv {{r}^{1_{\omega}}} = r \to \lvert {{r}^{1_{\omega}}} \rvert = \lvert {{r}^{1_{\omega}}} \rvert \to \lvert {{r}^{1_{\omega}}} \rvert = \lvert r \rvert$
\item $\varphi_{50} :\equiv n \in B_{26} \to S(n) \in \mathbb{N} \to n \neq \emptyset \to \lvert {{r}^{n}} \cdot r \rvert = \lvert {{r}^{n}} \rvert \cdot \lvert r \rvert$
\item $\varphi_{51} :\equiv (r \in \mathbb{R} \to \forall n\,(r \in \mathbb{R})) \to r \in \mathbb{R} \to \forall n \in \mathbb{N}\,(\lvert {{r}^{n}} \rvert = {{\lvert r \rvert}^{n}})$
\item $\varphi_{52} :\equiv S(n) \in \mathbb{N} \to n \neq \emptyset \to \forall n \in \mathbb{N}\,({{\lvert r \rvert}^{S(n)}} = {{\lvert r \rvert}^{n}} \cdot \lvert r \rvert)$
\item $\varphi_{53} :\equiv \forall z_{28}\,\forall A_{27}\,\exists B_{26}\,\forall n\,(n \in B_{26} \leftrightarrow n \in A_{27} \land (n \in \mathbb{N} \to \lvert {{z_{28}}^{n}} \rvert = {{\lvert z_{28} \rvert}^{n}}))$
\item $\varphi_{54} :\equiv \forall w\,({{r}^{S(n)}} = w \to \lvert {{r}^{S(n)}} \rvert = \lvert {{r}^{S(n)}} \rvert \to \lvert {{r}^{S(n)}} \rvert = \lvert w \rvert)$
\item $\varphi_{55} :\equiv \forall w\,({{r}^{1_{\omega}}} = w \to \lvert {{r}^{1_{\omega}}} \rvert = \lvert {{r}^{1_{\omega}}} \rvert \to \lvert {{r}^{1_{\omega}}} \rvert = \lvert w \rvert)$
\item $\varphi_{56} :\equiv \forall n\,(n \in B_{26} \leftrightarrow n \in \omega \land (n \in \mathbb{N} \to \lvert {{r}^{n}} \rvert = {{\lvert r \rvert}^{n}})) \to \emptyset \in B_{26}$
\item $\varphi_{57} :\equiv \forall a\,\forall n\,(\mathrm{Seq}(a) \to n \in \mathbb{N} \to a_{n} \in \mathbb{R}) \to \forall n\,(\mathrm{Seq}({(r)^{\bullet}}) \to n \in \mathbb{N} \to {{r}^{n}} \in \mathbb{R})$
\item $\varphi_{58} :\equiv \forall n \in \mathbb{N}\,({{r}^{S(n)}} = {{r}^{n}} \cdot r) \to n \in \mathbb{N} \to {{r}^{S(n)}} = {{r}^{n}} \cdot r$
\item $\varphi_{59} :\equiv n \in B_{26} \to S(n) \in \mathbb{N} \to n \neq \emptyset \to \lvert {{r}^{n}} \rvert \cdot \lvert r \rvert = {{\lvert r \rvert}^{n}} \cdot \lvert r \rvert$
\item $\varphi_{60} :\equiv 1_{\omega} = S(n) \to \lvert {{r}^{1_{\omega}}} \rvert = {{\lvert r \rvert}^{1_{\omega}}} \to \lvert {{r}^{S(n)}} \rvert = {{\lvert r \rvert}^{S(n)}}$
\item $\varphi_{61} :\equiv (r \in \mathbb{R} \to n \in \mathbb{N} \to n \in \mathbb{N}) \to r \in \mathbb{R} \to n \in \mathbb{N} \to \lvert {{r}^{n}} \rvert = {{\lvert r \rvert}^{n}}$
\item $\varphi_{62} :\equiv \forall z\,(\lvert {{r}^{S(n)}} \rvert = \lvert {{r}^{n}} \cdot r \rvert \to \lvert {{r}^{n}} \cdot r \rvert = z \to \lvert {{r}^{S(n)}} \rvert = z)$
\item $\varphi_{63} :\equiv \forall n\,(\mathrm{Seq}({(r)^{\bullet}}) \to n \in \mathbb{N} \to {{r}^{n}} \in \mathbb{R}) \to \mathrm{Seq}({(r)^{\bullet}}) \to n \in \mathbb{N} \to {{r}^{n}} \in \mathbb{R}$
\item $\varphi_{64} :\equiv {{\lvert r \rvert}^{S(n)}} = {{\lvert r \rvert}^{n}} \cdot \lvert r \rvert \to {{\lvert r \rvert}^{n}} \cdot \lvert r \rvert = {{\lvert r \rvert}^{S(n)}}$
\item $\varphi_{65} :\equiv ({{\lvert r \rvert}^{1_{\omega}}} = \lvert r \rvert \to \lvert r \rvert = {{\lvert r \rvert}^{1_{\omega}}}) \to r \in \mathbb{R} \to \lvert r \rvert = {{\lvert r \rvert}^{1_{\omega}}}$
\item $\varphi_{66} :\equiv \forall n\,(n \in B_{26} \leftrightarrow n \in \omega \land (n \in \mathbb{N} \to \lvert {{r}^{n}} \rvert = {{\lvert r \rvert}^{n}})) \to n \in B_{26} \to n \in B_{26}$
\item $\varphi_{67} :\equiv \lvert {{r}^{1_{\omega}}} \rvert = \lvert r \rvert \to \lvert r \rvert = {{\lvert r \rvert}^{1_{\omega}}} \to \lvert {{r}^{1_{\omega}}} \rvert = {{\lvert r \rvert}^{1_{\omega}}}$
\item $\varphi_{68} :\equiv \forall n\,(n \in \mathbb{N} \leftrightarrow n \in \omega \land n \neq \emptyset) \to (\emptyset \in \mathbb{N} \leftrightarrow \emptyset \in \omega \land \emptyset \neq \emptyset)$
\item $\varphi_{69} :\equiv \emptyset = \emptyset \to r \in \mathbb{R} \to \emptyset \in \mathbb{N} \to \lvert {{r}^{\emptyset}} \rvert = {{\lvert r \rvert}^{\emptyset}}$
\item $\varphi_{70} :\equiv \lvert {{r}^{n}} \cdot r \rvert = \lvert {{r}^{n}} \rvert \cdot \lvert r \rvert \to \lvert {{r}^{S(n)}} \rvert = \lvert {{r}^{n}} \rvert \cdot \lvert r \rvert$
\item $\varphi_{71} :\equiv \forall n\,(n \in B_{26} \leftrightarrow n \in \omega \land (n \in \mathbb{N} \to \lvert {{r}^{n}} \rvert = {{\lvert r \rvert}^{n}})) \to n \in B_{26} \to n \in \omega$
\item $\varphi_{72} :\equiv \forall n\,(n \in B_{26} \leftrightarrow n \in \omega \land (n \in \mathbb{N} \to \lvert {{r}^{n}} \rvert = {{\lvert r \rvert}^{n}})) \to n \in B_{26} \to S(n) \in B_{26}$
\item $\varphi_{73} :\equiv \forall n\,(n \in B_{26} \leftrightarrow n \in \omega \land (n \in \mathbb{N} \to \lvert {{r}^{n}} \rvert = {{\lvert r \rvert}^{n}})) \to n \in \omega \to n \in \omega$
\item $\varphi_{74} :\equiv \forall n\,(n \in B_{26} \leftrightarrow n \in \omega \land (n \in \mathbb{N} \to \lvert {{r}^{n}} \rvert = {{\lvert r \rvert}^{n}})) \to n \in B_{26} \to S(n) \in \omega$
\item $\varphi_{75} :\equiv \lvert {{r}^{n}} \rvert \cdot \lvert r \rvert = {{\lvert r \rvert}^{n}} \cdot \lvert r \rvert \to \lvert {{r}^{S(n)}} \rvert = {{\lvert r \rvert}^{n}} \cdot \lvert r \rvert$
\item $\varphi_{76} :\equiv \forall n\,(n \in B_{26} \leftrightarrow n \in \omega \land (n \in \mathbb{N} \to \lvert {{r}^{n}} \rvert = {{\lvert r \rvert}^{n}})) \to \forall n \in B_{26}\,(S(n) \in B_{26})$
\item $\varphi_{77} :\equiv \forall z\,(\lvert {{r}^{S(n)}} \rvert = {{\lvert r \rvert}^{n}} \cdot \lvert r \rvert \to {{\lvert r \rvert}^{n}} \cdot \lvert r \rvert = z \to \lvert {{r}^{S(n)}} \rvert = z)$
\item $\varphi_{78} :\equiv \forall z\,(\lvert {{r}^{S(n)}} \rvert = \lvert {{r}^{n}} \rvert \cdot \lvert r \rvert \to \lvert {{r}^{n}} \rvert \cdot \lvert r \rvert = z \to \lvert {{r}^{S(n)}} \rvert = z)$
\item $\varphi_{79} :\equiv \emptyset \in B_{26} \leftrightarrow \emptyset \in \omega \land (\emptyset \in \mathbb{N} \to \lvert {{r}^{\emptyset}} \rvert = {{\lvert r \rvert}^{\emptyset}})$
\item $\varphi_{80} :\equiv \forall y\,({{\lvert r \rvert}^{1_{\omega}}} = y \to y = {{\lvert r \rvert}^{1_{\omega}}}) \to {{\lvert r \rvert}^{1_{\omega}}} = \lvert r \rvert \to \lvert r \rvert = {{\lvert r \rvert}^{1_{\omega}}}$
\item $\varphi_{81} :\equiv \mathrm{Seq}({(\lvert r \rvert)^{\bullet}}) \land {{\lvert r \rvert}^{1_{\omega}}} = \lvert r \rvert \land \forall n \in \mathbb{N}\,({{\lvert r \rvert}^{S(n)}} = {{\lvert r \rvert}^{n}} \cdot \lvert r \rvert)$
\item $\varphi_{82} :\equiv \forall n\,(n \in B_{26} \leftrightarrow n \in \omega \land (n \in \mathbb{N} \to \lvert {{r}^{n}} \rvert = {{\lvert r \rvert}^{n}})) \to n \in \omega \to \forall x\,(\mathrm{Ind}(x) \to n \in x)$
\item $\varphi_{83} :\equiv (r \in \mathbb{R} \to n \in \mathbb{N} \to n \in \mathbb{N}) \to (n \in \mathbb{N} \leftrightarrow n \in \omega \land n \neq \emptyset) \to r \in \mathbb{R} \to n \in \mathbb{N} \to n \in \omega$
\item $\varphi_{84} :\equiv \lvert {{r}^{n}} \rvert \cdot \lvert r \rvert = \lvert {{r}^{n}} \rvert \cdot \lvert r \rvert \to \lvert {{r}^{n}} \rvert \cdot \lvert r \rvert = {{\lvert r \rvert}^{n}} \cdot \lvert r \rvert$
\item $\varphi_{85} :\equiv \forall n\,(n \in B_{26} \leftrightarrow n \in \omega \land (n \in \mathbb{N} \to \lvert {{r}^{n}} \rvert = {{\lvert r \rvert}^{n}})) \to n \in \omega \to \mathrm{Ind}(B_{26}) \to n \in B_{26}$
\item $\varphi_{86} :\equiv \forall n\,(n \in B_{26} \leftrightarrow n \in \omega \land (n \in \mathbb{N} \to \lvert {{r}^{n}} \rvert = {{\lvert r \rvert}^{n}})) \to n \in B_{26} \to n \in \omega \to S(n) \in \omega$
\item $\varphi_{87} :\equiv \forall x\,(\mathrm{Ind}(x) \leftrightarrow \emptyset \in x \land \forall y \in x\,(S(y) \in x)) \to (\mathrm{Ind}(\omega) \leftrightarrow \emptyset \in \omega \land \forall y \in \omega\,(S(y) \in \omega))$
\item $\varphi_{88} :\equiv \forall w\,(\lvert {{r}^{n}} \rvert = w \to \lvert {{r}^{n}} \rvert \cdot \lvert r \rvert = \lvert {{r}^{n}} \rvert \cdot \lvert r \rvert \to \lvert {{r}^{n}} \rvert \cdot \lvert r \rvert = w \cdot \lvert r \rvert)$
\item $\varphi_{89} :\equiv \forall n\,(n \in B_{26} \leftrightarrow n \in \omega \land (n \in \mathbb{N} \to \lvert {{r}^{n}} \rvert = {{\lvert r \rvert}^{n}})) \to n \in B_{26} \to \forall y \in \omega\,(S(y) \in \omega)$
\item $\varphi_{90} :\equiv {{r}^{S(n)}} = {{r}^{n}} \cdot r \to \lvert {{r}^{S(n)}} \rvert = \lvert {{r}^{S(n)}} \rvert \to \lvert {{r}^{S(n)}} \rvert = \lvert {{r}^{n}} \cdot r \rvert$
\item $\varphi_{91} :\equiv \forall n \in \mathbb{N}\,({{\lvert r \rvert}^{S(n)}} = {{\lvert r \rvert}^{n}} \cdot \lvert r \rvert) \to n \in \mathbb{N} \to {{\lvert r \rvert}^{S(n)}} = {{\lvert r \rvert}^{n}} \cdot \lvert r \rvert$
\item $\varphi_{92} :\equiv \forall n \in \omega\,(n \in \mathbb{N} \to \lvert {{r}^{n}} \rvert = {{\lvert r \rvert}^{n}}) \to n \in \omega \to n \in \mathbb{N} \to \lvert {{r}^{n}} \rvert = {{\lvert r \rvert}^{n}}$
\item $\varphi_{93} :\equiv \forall y\,\forall z\,(\lvert {{r}^{1_{\omega}}} \rvert = y \to y = z \to \lvert {{r}^{1_{\omega}}} \rvert = z) \to \forall z\,(\lvert {{r}^{1_{\omega}}} \rvert = \lvert r \rvert \to \lvert r \rvert = z \to \lvert {{r}^{1_{\omega}}} \rvert = z)$
\item $\varphi_{94} :\equiv \forall n\,(n \in B_{26} \leftrightarrow n \in \omega \land (n \in \mathbb{N} \to \lvert {{r}^{n}} \rvert = {{\lvert r \rvert}^{n}})) \to n \in B_{26} \to S(n) \in \mathbb{N} \to n = \emptyset \to S(n) = 1_{\omega}$
\item $\varphi_{95} :\equiv \forall n\,(n \in B_{26} \leftrightarrow n \in \omega \land (n \in \mathbb{N} \to \lvert {{r}^{n}} \rvert = {{\lvert r \rvert}^{n}})) \to n \in B_{26} \to S(n) \in \mathbb{N} \to n = \emptyset \to 1_{\omega} = S(n)$
\item $\varphi_{96} :\equiv \forall n\,(n \in B_{26} \leftrightarrow n \in \omega \land (n \in \mathbb{N} \to \lvert {{r}^{n}} \rvert = {{\lvert r \rvert}^{n}})) \to n \in B_{26} \to S(n) \in \mathbb{N} \to n \neq \emptyset \to n \in \mathbb{N}$
\item $\varphi_{97} :\equiv \forall n\,(n \in B_{26} \leftrightarrow n \in \omega \land (n \in \mathbb{N} \to \lvert {{r}^{n}} \rvert = {{\lvert r \rvert}^{n}})) \to n \in B_{26} \to S(n) \in \mathbb{N} \to n = \emptyset \to n = \emptyset$
\item $\varphi_{98} :\equiv \forall n\,(n \in B_{26} \leftrightarrow n \in \omega \land (n \in \mathbb{N} \to \lvert {{r}^{n}} \rvert = {{\lvert r \rvert}^{n}})) \to n \in B_{26} \to S(n) \in \mathbb{N} \to n \neq \emptyset \to r \in \mathbb{R}$
\item $\varphi_{99} :\equiv \forall n\,(n \in B_{26} \leftrightarrow n \in \omega \land (n \in \mathbb{N} \to \lvert {{r}^{n}} \rvert = {{\lvert r \rvert}^{n}})) \to \varphi_{5}$
\item $\varphi_{100} :\equiv \forall n\,(n \in B_{26} \leftrightarrow n \in \omega \land (n \in \mathbb{N} \to \lvert {{r}^{n}} \rvert = {{\lvert r \rvert}^{n}})) \to n \in B_{26} \to S(n) \in \mathbb{N} \to n \neq \emptyset \to n \neq \emptyset$
\item $\varphi_{101} :\equiv \forall n\,(n \in B_{26} \leftrightarrow n \in \omega \land (n \in \mathbb{N} \to \lvert {{r}^{n}} \rvert = {{\lvert r \rvert}^{n}})) \to \emptyset \in \omega \land \forall y \in \omega\,(S(y) \in \omega)$
\item $\varphi_{102} :\equiv \forall n\,(r \in \mathbb{R} \to n \in \mathbb{N} \to \lvert {{r}^{n}} \rvert = {{\lvert r \rvert}^{n}}) \to \forall n\,(r \in \mathbb{R}) \to \forall n \in \mathbb{N}\,(\lvert {{r}^{n}} \rvert = {{\lvert r \rvert}^{n}})$
\item $\varphi_{103} :\equiv \forall n\,(n \in B_{26} \leftrightarrow n \in \omega \land (n \in \mathbb{N} \to \lvert {{r}^{n}} \rvert = {{\lvert r \rvert}^{n}})) \to n \in B_{26} \to S(n) \in \mathbb{N} \to n \neq \emptyset \to \lvert r \rvert \in \mathbb{R}$
\item $\varphi_{104} :\equiv \forall n\,(n \in B_{26} \leftrightarrow n \in \omega \land (n \in \mathbb{N} \to \lvert {{r}^{n}} \rvert = {{\lvert r \rvert}^{n}})) \to n \in B_{26} \to S(n) \in \mathbb{N} \to n = \emptyset \to S(n) = S(\emptyset)$
\item $\varphi_{105} :\equiv \forall n\,(n \in B_{26} \leftrightarrow n \in \omega \land (n \in \mathbb{N} \to \lvert {{r}^{n}} \rvert = {{\lvert r \rvert}^{n}})) \to n \in B_{26} \to n \in \mathbb{N} \to \lvert {{r}^{n}} \rvert = {{\lvert r \rvert}^{n}}$
\item $\varphi_{106} :\equiv \forall n\,(n \in B_{26} \leftrightarrow n \in \omega \land (n \in \mathbb{N} \to \lvert {{r}^{n}} \rvert = {{\lvert r \rvert}^{n}})) \to n \in \omega \to n \in \mathbb{N} \to \lvert {{r}^{n}} \rvert = {{\lvert r \rvert}^{n}}$
\item $\varphi_{107} :\equiv \forall n\,(n \in B_{26} \leftrightarrow n \in \omega \land (n \in \mathbb{N} \to \lvert {{r}^{n}} \rvert = {{\lvert r \rvert}^{n}})) \to n \in B_{26} \to S(n) \in \mathbb{N} \to n \neq \emptyset \to {{r}^{n}} \in \mathbb{R}$
\item $\varphi_{108} :\equiv \lvert {{r}^{S(n)}} \rvert = {{\lvert r \rvert}^{n}} \cdot \lvert r \rvert \to {{\lvert r \rvert}^{n}} \cdot \lvert r \rvert = {{\lvert r \rvert}^{S(n)}} \to \lvert {{r}^{S(n)}} \rvert = {{\lvert r \rvert}^{S(n)}}$
\item $\varphi_{109} :\equiv \forall n\,(n \in B_{26} \leftrightarrow n \in \omega \land (n \in \mathbb{N} \to \lvert {{r}^{n}} \rvert = {{\lvert r \rvert}^{n}})) \to \forall n \in \omega\,(n \in \mathbb{N} \to \lvert {{r}^{n}} \rvert = {{\lvert r \rvert}^{n}})$
\item $\varphi_{110} :\equiv \forall n\,(n \in B_{26} \leftrightarrow n \in \omega \land (n \in \mathbb{N} \to \lvert {{r}^{n}} \rvert = {{\lvert r \rvert}^{n}})) \to n \in B_{26} \to S(n) \in \mathbb{N} \to \lvert {{r}^{S(n)}} \rvert = {{\lvert r \rvert}^{S(n)}}$
\item $\varphi_{111} :\equiv \forall n\,(n \in B_{26} \leftrightarrow n \in \omega \land (n \in \mathbb{N} \to \lvert {{r}^{n}} \rvert = {{\lvert r \rvert}^{n}})) \to (n \in B_{26} \leftrightarrow n \in \omega \land (n \in \mathbb{N} \to \lvert {{r}^{n}} \rvert = {{\lvert r \rvert}^{n}}))$
\item $\varphi_{112} :\equiv \forall n\,(n \in B_{26} \leftrightarrow n \in \omega \land (n \in \mathbb{N} \to \lvert {{r}^{n}} \rvert = {{\lvert r \rvert}^{n}})) \to n \in B_{26} \to n \in \omega \land (n \in \mathbb{N} \to \lvert {{r}^{n}} \rvert = {{\lvert r \rvert}^{n}})$
\item $\varphi_{113} :\equiv \forall x\,\forall y\,(x \in \mathbb{R} \to y \in \mathbb{R} \to \lvert x \cdot y \rvert = \lvert x \rvert \cdot \lvert y \rvert) \to \forall y\,({{r}^{n}} \in \mathbb{R} \to y \in \mathbb{R} \to \lvert {{r}^{n}} \cdot y \rvert = \lvert {{r}^{n}} \rvert \cdot \lvert y \rvert)$
\item $\varphi_{114} :\equiv \forall n\,(n \in B_{26} \leftrightarrow n \in \omega \land (n \in \mathbb{N} \to \lvert {{r}^{n}} \rvert = {{\lvert r \rvert}^{n}})) \to n \in B_{26} \to S(n) \in \mathbb{N} \to n \neq \emptyset \to {{r}^{S(n)}} = {{r}^{n}} \cdot r$
\item $\varphi_{115} :\equiv \forall n\,(n \in B_{26} \leftrightarrow n \in \omega \land (n \in \mathbb{N} \to \lvert {{r}^{n}} \rvert = {{\lvert r \rvert}^{n}})) \to \forall n\,(n \in B_{26} \leftrightarrow n \in \omega \land (n \in \mathbb{N} \to \lvert {{r}^{n}} \rvert = {{\lvert r \rvert}^{n}}))$
\item $\varphi_{116} :\equiv \forall n\,(n \in B_{26} \leftrightarrow n \in \omega \land (n \in \mathbb{N} \to \lvert {{r}^{n}} \rvert = {{\lvert r \rvert}^{n}})) \to n \in B_{26} \to S(n) \in \mathbb{N} \to n \neq \emptyset \to \lvert {{r}^{n}} \rvert = {{\lvert r \rvert}^{n}}$
\item $\varphi_{117} :\equiv \forall n\,(n \in B_{26} \leftrightarrow n \in \omega \land (n \in \mathbb{N} \to \lvert {{r}^{n}} \rvert = {{\lvert r \rvert}^{n}})) \to \varphi_{30}$
\item $\varphi_{118} :\equiv \forall n\,(n \in B_{26} \leftrightarrow n \in \omega \land (n \in \mathbb{N} \to \lvert {{r}^{n}} \rvert = {{\lvert r \rvert}^{n}})) \to \varphi_{31}$
\item $\varphi_{119} :\equiv \forall B_{26}\,(r \in \mathbb{R} \land \neg \forall n \in \omega\,(n \in \mathbb{N} \to \lvert {{r}^{n}} \rvert = {{\lvert r \rvert}^{n}}) \to \varphi_{26})$
\item $\varphi_{120} :\equiv \forall n\,(n \in B_{26} \leftrightarrow n \in \omega \land (n \in \mathbb{N} \to \lvert {{r}^{n}} \rvert = {{\lvert r \rvert}^{n}})) \to \varphi_{34}$
\item $\varphi_{121} :\equiv \forall n\,(n \in B_{26} \leftrightarrow n \in \omega \land (n \in \mathbb{N} \to \lvert {{r}^{n}} \rvert = {{\lvert r \rvert}^{n}})) \to \varphi_{37}$
\item $\varphi_{122} :\equiv \forall n\,(n \in B_{26} \leftrightarrow n \in \omega \land (n \in \mathbb{N} \to \lvert {{r}^{n}} \rvert = {{\lvert r \rvert}^{n}})) \to \varphi_{40}$
\item $\varphi_{123} :\equiv \forall n\,(n \in B_{26} \leftrightarrow n \in \omega \land (n \in \mathbb{N} \to \lvert {{r}^{n}} \rvert = {{\lvert r \rvert}^{n}})) \to \varphi_{43}$
\item $\varphi_{124} :\equiv \forall n\,(n \in B_{26} \leftrightarrow n \in \omega \land (n \in \mathbb{N} \to \lvert {{r}^{n}} \rvert = {{\lvert r \rvert}^{n}})) \to \varphi_{44}$
\item $\varphi_{125} :\equiv \forall n\,(n \in B_{26} \leftrightarrow n \in \omega \land (n \in \mathbb{N} \to \lvert {{r}^{n}} \rvert = {{\lvert r \rvert}^{n}})) \to \varphi_{45}$
\item $\varphi_{126} :\equiv \forall n\,(n \in B_{26} \leftrightarrow n \in \omega \land (n \in \mathbb{N} \to \lvert {{r}^{n}} \rvert = {{\lvert r \rvert}^{n}})) \to \varphi_{46}$
\item $\varphi_{127} :\equiv \forall n\,(n \in B_{26} \leftrightarrow n \in \omega \land (n \in \mathbb{N} \to \lvert {{r}^{n}} \rvert = {{\lvert r \rvert}^{n}})) \to \varphi_{47}$
\item $\varphi_{128} :\equiv \forall z\,(\lvert {{r}^{1_{\omega}}} \rvert = \lvert r \rvert \to \lvert r \rvert = z \to \lvert {{r}^{1_{\omega}}} \rvert = z) \to \varphi_{67}$
\item $\varphi_{129} :\equiv \forall n\,(n \in B_{26} \leftrightarrow n \in \omega \land (n \in \mathbb{N} \to \lvert {{r}^{n}} \rvert = {{\lvert r \rvert}^{n}})) \to \varphi_{50}$
\item $\varphi_{130} :\equiv (\forall x\,(\mathrm{Ind}(x) \to n \in x) \to \mathrm{Ind}(B_{26}) \to n \in B_{26}) \to r \in \mathbb{R} \to \varphi_{85}$
\item $\varphi_{131} :\equiv \forall n\,(n \in B_{26} \leftrightarrow n \in \omega \land (n \in \mathbb{N} \to \lvert {{r}^{n}} \rvert = {{\lvert r \rvert}^{n}})) \to \varphi_{59}$
\item $\varphi_{132} :\equiv \lvert {{r}^{1_{\omega}}} \rvert = \lvert {{r}^{1_{\omega}}} \rvert \to \varphi_{49} \to r \in \mathbb{R} \to \lvert {{r}^{1_{\omega}}} \rvert = \lvert r \rvert$
\item $\varphi_{133} :\equiv \forall z\,\forall w\,(z = w \to \lvert {{r}^{z}} \rvert = {{\lvert r \rvert}^{z}} \to \lvert {{r}^{w}} \rvert = {{\lvert r \rvert}^{w}}) \to \varphi_{48}$
\item $\varphi_{134} :\equiv \forall y\,({{r}^{n}} \in \mathbb{R} \to y \in \mathbb{R} \to \lvert {{r}^{n}} \cdot y \rvert = \lvert {{r}^{n}} \rvert \cdot \lvert y \rvert) \to {{r}^{n}} \in \mathbb{R} \to r \in \mathbb{R} \to \lvert {{r}^{n}} \cdot r \rvert = \lvert {{r}^{n}} \rvert \cdot \lvert r \rvert$
\item $\varphi_{135} :\equiv \forall n\,(n \in B_{26} \leftrightarrow n \in \omega \land (n \in \mathbb{N} \to \lvert {{r}^{n}} \rvert = {{\lvert r \rvert}^{n}})) \to n \in B_{26} \to \varphi_{42}$
\item $\varphi_{136} :\equiv \forall n\,(n \in B_{26} \leftrightarrow n \in \omega \land (n \in \mathbb{N} \to \lvert {{r}^{n}} \rvert = {{\lvert r \rvert}^{n}})) \to n \in B_{26} \to \varphi_{52}$
\item $\varphi_{137} :\equiv \forall n\,(n \in B_{26} \leftrightarrow n \in \omega \land (n \in \mathbb{N} \to \lvert {{r}^{n}} \rvert = {{\lvert r \rvert}^{n}})) \to \varphi_{79}$
\item $\varphi_{138} :\equiv (n = \emptyset \to S(n) = S(n) \to S(n) = S(\emptyset)) \to r \in \mathbb{R} \to \varphi_{104}$
\item $\varphi_{139} :\equiv (r \in \mathbb{R} \to \emptyset \in \mathbb{N} \to \lvert {{r}^{\emptyset}} \rvert = {{\lvert r \rvert}^{\emptyset}}) \to r \in \mathbb{R} \to \varphi_{56}$
\item $\varphi_{140} :\equiv (\forall y \in \omega\,(S(y) \in \omega) \to n \in \omega \to S(n) \in \omega) \to r \in \mathbb{R} \to \varphi_{86}$
\item $\varphi_{141} :\equiv (r \in \mathbb{R} \to \lvert r \rvert = {{\lvert r \rvert}^{1_{\omega}}}) \to \varphi_{67} \to r \in \mathbb{R} \to \lvert {{r}^{1_{\omega}}} \rvert = {{\lvert r \rvert}^{1_{\omega}}}$
\item $\varphi_{142} :\equiv \forall n\,(n \in B_{26} \leftrightarrow n \in \omega \land (n \in \mathbb{N} \to \lvert {{r}^{n}} \rvert = {{\lvert r \rvert}^{n}})) \to n \in B_{26} \to S(n) \in \mathbb{N} \to n \neq \emptyset \to \varphi_{29}$
\item $\varphi_{143} :\equiv (\mathrm{Ind}(\omega) \leftrightarrow \emptyset \in \omega \land \forall y \in \omega\,(S(y) \in \omega)) \to r \in \mathbb{R} \to \varphi_{101}$
\item $\varphi_{144} :\equiv (r \in \mathbb{R} \to n \in \omega \to n \in \mathbb{N} \to \lvert {{r}^{n}} \rvert = {{\lvert r \rvert}^{n}}) \to (r \in \mathbb{R} \to n \in \mathbb{N} \to n \in \omega) \to \varphi_{61}$
\item $\varphi_{145} :\equiv \forall n\,(n \in B_{26} \leftrightarrow n \in \omega \land (n \in \mathbb{N} \to \lvert {{r}^{n}} \rvert = {{\lvert r \rvert}^{n}})) \to n \in B_{26} \to S(n) \in \mathbb{N} \to n \neq \emptyset \to \varphi_{81}$
\item $\varphi_{146} :\equiv \forall n\,(\varphi_{41} \to n \in B_{26} \to S(n) \in B_{26}) \to \forall n\,(\varphi_{41}) \to \forall n \in B_{26}\,(S(n) \in B_{26})$
\item $\varphi_{147} :\equiv (r \in \mathbb{R} \land \neg \forall n \in \omega\,(n \in \mathbb{N} \to \lvert {{r}^{n}} \rvert = {{\lvert r \rvert}^{n}}) \to \varphi_{32}) \to r \in \mathbb{R} \to \forall n \in \omega\,(n \in \mathbb{N} \to \lvert {{r}^{n}} \rvert = {{\lvert r \rvert}^{n}})$
\item $\varphi_{148} :\equiv (r \in \mathbb{R} \to \varphi_{73}) \to (n \in \omega \leftrightarrow \forall x\,(\mathrm{Ind}(x) \to n \in x)) \to r \in \mathbb{R} \to \varphi_{82}$
\item $\varphi_{149} :\equiv (r \in \mathbb{R} \to \forall n \in \omega\,(n \in \mathbb{N} \to \lvert {{r}^{n}} \rvert = {{\lvert r \rvert}^{n}})) \to \varphi_{92} \to r \in \mathbb{R} \to n \in \omega \to n \in \mathbb{N} \to \lvert {{r}^{n}} \rvert = {{\lvert r \rvert}^{n}}$
\item $\varphi_{150} :\equiv (\lvert r \rvert \in \mathbb{R} \to \varphi_{81}) \to (r \in \mathbb{R} \to \lvert r \rvert \in \mathbb{R}) \to r \in \mathbb{R} \to \varphi_{81}$
\item $\varphi_{151} :\equiv (r \in \mathbb{R} \to \varphi_{94}) \to (S(n) = 1_{\omega} \to 1_{\omega} = S(n)) \to r \in \mathbb{R} \to \varphi_{95}$
\item $\varphi_{152} :\equiv (r \in \mathbb{R} \to \varphi_{109}) \to r \in \mathbb{R} \land \neg \forall n \in \omega\,(n \in \mathbb{N} \to \lvert {{r}^{n}} \rvert = {{\lvert r \rvert}^{n}}) \to \varphi_{26}$
\item $\varphi_{153} :\equiv \forall n\,(\varphi_{41} \to n \in \omega \to n \in \mathbb{N} \to \lvert {{r}^{n}} \rvert = {{\lvert r \rvert}^{n}}) \to \forall n\,(\varphi_{41}) \to \forall n \in \omega\,(n \in \mathbb{N} \to \lvert {{r}^{n}} \rvert = {{\lvert r \rvert}^{n}})$
\item $\varphi_{154} :\equiv \lvert {{r}^{n}} \rvert \cdot \lvert r \rvert = \lvert {{r}^{n}} \rvert \cdot \lvert r \rvert \to (\lvert {{r}^{n}} \rvert = {{\lvert r \rvert}^{n}} \to \varphi_{84}) \to r \in \mathbb{R} \to \varphi_{131}$
\item $\varphi_{155} :\equiv (n \in \mathbb{N} \leftrightarrow n \in \omega \land n \neq \emptyset) \to (r \in \mathbb{R} \to \varphi_{71}) \to (r \in \mathbb{R} \to \varphi_{100}) \to r \in \mathbb{R} \to \varphi_{96}$
\item $\varphi_{156} :\equiv (r \in \mathbb{R} \to \varphi_{95}) \to (r \in \mathbb{R} \to \lvert {{r}^{1_{\omega}}} \rvert = {{\lvert r \rvert}^{1_{\omega}}}) \to \varphi_{60} \to r \in \mathbb{R} \to \varphi_{117}$
\item $\varphi_{157} :\equiv (\mathrm{Seq}({(r)^{\bullet}}) \to n \in \mathbb{N} \to {{r}^{n}} \in \mathbb{R}) \to (r \in \mathbb{R} \to \varphi_{99}) \to (r \in \mathbb{R} \to \varphi_{96}) \to r \in \mathbb{R} \to \varphi_{107}$
\item $\varphi_{158} :\equiv (r \in \mathbb{R} \to \varphi_{114}) \to \lvert {{r}^{S(n)}} \rvert = \lvert {{r}^{S(n)}} \rvert \to \varphi_{90} \to r \in \mathbb{R} \to \varphi_{120}$
\item $\varphi_{159} :\equiv (r \in \mathbb{R} \to \varphi_{129}) \to (\lvert {{r}^{S(n)}} \rvert = \lvert {{r}^{n}} \cdot r \rvert \to \varphi_{70}) \to r \in \mathbb{R} \to \varphi_{126}$
\item $\varphi_{160} :\equiv (r \in \mathbb{R} \to \varphi_{131}) \to (\lvert {{r}^{S(n)}} \rvert = \lvert {{r}^{n}} \rvert \cdot \lvert r \rvert \to \varphi_{75}) \to r \in \mathbb{R} \to \varphi_{127}$
\item $\varphi_{161} :\equiv (r \in \mathbb{R} \to \varphi_{85}) \to (r \in \mathbb{R} \to \varphi_{39}) \to (r \in \mathbb{R} \to \varphi_{111}) \to r \in \mathbb{R} \to \varphi_{106}$
\item $\varphi_{162} :\equiv ({{r}^{n}} \in \mathbb{R} \to r \in \mathbb{R} \to \lvert {{r}^{n}} \cdot r \rvert = \lvert {{r}^{n}} \rvert \cdot \lvert r \rvert) \to (r \in \mathbb{R} \to \varphi_{107}) \to (r \in \mathbb{R} \to \varphi_{98}) \to r \in \mathbb{R} \to \varphi_{129}$
\item $\varphi_{163} :\equiv (r \in \mathbb{R} \to \varphi_{122}) \to (r \in \mathbb{R} \to \varphi_{74}) \to (r \in \mathbb{R} \to \varphi_{110}) \to r \in \mathbb{R} \to \varphi_{72}$
\item $\varphi_{164} :\equiv \forall n\,(\varphi_{41} \to n \in B_{26} \to S(n) \in B_{26}) \to \varphi_{146} \to (\varphi_{41} \to \forall n\,(\varphi_{41})) \to r \in \mathbb{R} \to \varphi_{76}$
\item $\varphi_{165} :\equiv \forall n\,(\varphi_{41} \to n \in \omega \to n \in \mathbb{N} \to \lvert {{r}^{n}} \rvert = {{\lvert r \rvert}^{n}}) \to \varphi_{153} \to (\varphi_{41} \to \forall n\,(\varphi_{41})) \to r \in \mathbb{R} \to \varphi_{109}$
\end{abbrevlist}
\end{fullproof}

% 自動生成 (prover/build.js): PowNonneg  ノード数 352
\begin{fullproof}{[P2]}{$\forall r \in \mathbb{R}\,(0 \leq r \to \forall n \in \mathbb{N}\,(0 \leq {{r}^{n}}))$}{352}\phantomsection\label{tree:PowNonneg}
\noindent$\Pi_{1}$:
\begin{ptree}
\AXC{$r \in \mathbb{R} \to 0 \leq r \to n \in \mathbb{N} \to n \in \mathbb{N}$\lbl{A1}}
\AXC{$(r \in \mathbb{R} \to 0 \leq r \to n \in \mathbb{N} \to n \in \mathbb{N}) \to \varphi_{44}$\lbl{A1}}
\BIC{$\varphi_{44}$}
\end{ptree}
\noindent$\Pi_{2}$:
\begin{ptree}
\AXC{$\varphi_{77}$\lbl{A2}}
\AXC{$r \in \mathbb{R} \to 0 \leq r \to \varphi_{83}$\lbl{A1}}
\AXC{$(r \in \mathbb{R} \to 0 \leq r \to \varphi_{83}) \to \varphi_{77} \to r \in \mathbb{R} \to 0 \leq r \to \varphi_{77}$\lbl{A1}}
\BIC{$\varphi_{77} \to r \in \mathbb{R} \to 0 \leq r \to \varphi_{77}$}
\BIC{$r \in \mathbb{R} \to 0 \leq r \to \varphi_{77}$}
\end{ptree}
\noindent$\Pi_{3}$:
\begin{ptree}
\AXC{$\varphi_{77}$\lbl{A2}}
\AXC{$r \in \mathbb{R} \to 0 \leq r \to \varphi_{83}$\lbl{A1}}
\AXC{$(r \in \mathbb{R} \to 0 \leq r \to \varphi_{83}) \to \varphi_{77} \to r \in \mathbb{R} \to 0 \leq r \to \varphi_{77}$\lbl{A1}}
\BIC{$\varphi_{77} \to r \in \mathbb{R} \to 0 \leq r \to \varphi_{77}$}
\BIC{$r \in \mathbb{R} \to 0 \leq r \to \varphi_{77}$}
\AXC{$\varphi_{131}$\lbl{A1}}
\BIC{$(r \in \mathbb{R} \to 0 \leq r \to \varphi_{39}) \to r \in \mathbb{R} \to 0 \leq r \to \varphi_{80}$}
\end{ptree}
\noindent$\Pi_{4}$:
\begin{ptree}
\AXC{$r \in \mathbb{R} \to 0 \leq r \to \varphi_{39}$\lbl{A1}}
\AXC{$\stackrel{\Pi_{3}}{\vdots}$}\noLine
\UIC{$(r \in \mathbb{R} \to 0 \leq r \to \varphi_{39}) \to r \in \mathbb{R} \to 0 \leq r \to \varphi_{80}$}
\BIC{$r \in \mathbb{R} \to 0 \leq r \to \varphi_{80}$}
\end{ptree}
\noindent$\Pi_{5}$:
\begin{ptree}
\AXC{$\stackrel{\Pi_{4}}{\vdots}$}\noLine
\UIC{$r \in \mathbb{R} \to 0 \leq r \to \varphi_{80}$}
\AXC{$(r \in \mathbb{R} \to 0 \leq r \to \varphi_{80}) \to r \in \mathbb{R} \to 0 \leq r \to \varphi_{42}$\lbl{A1}}
\BIC{$r \in \mathbb{R} \to 0 \leq r \to \varphi_{42}$}
\end{ptree}
\noindent$\Pi_{6}$:
\begin{ptree}
\AXC{$\stackrel{\Pi_{5}}{\vdots}$}\noLine
\UIC{$r \in \mathbb{R} \to 0 \leq r \to \varphi_{42}$}
\AXC{$\forall n\,(n \in \mathbb{N} \leftrightarrow n \in \omega \land n \neq \emptyset)$\lbl{Def$\mathbb{N}$}}
\AXC{$\varphi_{40}$\lbl{A2}}
\BIC{$n \in \mathbb{N} \leftrightarrow n \in \omega \land n \neq \emptyset$}
\AXC{$(n \in \mathbb{N} \leftrightarrow n \in \omega \land n \neq \emptyset) \to \varphi_{132}$\lbl{A1}}
\BIC{$\varphi_{132}$}
\BIC{$(r \in \mathbb{R} \to 0 \leq r \to \varphi_{79}) \to r \in \mathbb{R} \to 0 \leq r \to \varphi_{74}$}
\end{ptree}
\noindent$\Pi_{7}$:
\begin{ptree}
\AXC{$r \in \mathbb{R} \to 0 \leq r \to \varphi_{79}$\lbl{A1}}
\AXC{$\stackrel{\Pi_{6}}{\vdots}$}\noLine
\UIC{$(r \in \mathbb{R} \to 0 \leq r \to \varphi_{79}) \to r \in \mathbb{R} \to 0 \leq r \to \varphi_{74}$}
\BIC{$r \in \mathbb{R} \to 0 \leq r \to \varphi_{74}$}
\end{ptree}
\noindent$\Pi_{8}$:
\begin{ptree}
\AXC{$\stackrel{\Pi_{7}}{\vdots}$}\noLine
\UIC{$r \in \mathbb{R} \to 0 \leq r \to \varphi_{74}$}
\AXC{$\stackrel{\Pi_{4}}{\vdots}$}\noLine
\UIC{$r \in \mathbb{R} \to 0 \leq r \to \varphi_{80}$}
\AXC{$(r \in \mathbb{R} \to 0 \leq r \to \varphi_{80}) \to r \in \mathbb{R} \to 0 \leq r \to \varphi_{61}$\lbl{A1}}
\BIC{$r \in \mathbb{R} \to 0 \leq r \to \varphi_{61}$}
\AXC{$\varphi_{133}$\lbl{A1}}
\BIC{$(r \in \mathbb{R} \to 0 \leq r \to \varphi_{74}) \to r \in \mathbb{R} \to 0 \leq r \to \varphi_{84}$}
\BIC{$r \in \mathbb{R} \to 0 \leq r \to \varphi_{84}$}
\end{ptree}
\noindent$\Pi_{9}$:
\begin{ptree}
\AXC{$r \in \mathbb{R} \to 0 \leq r \to \varphi_{76}$\lbl{A1}}
\AXC{$\forall r \in \mathbb{R}\,(\varphi_{43})$\lbl{[S24]}}
\AXC{$\forall r \in \mathbb{R}\,(\varphi_{43}) \to r \in \mathbb{R} \to \varphi_{43}$\lbl{A2}}
\BIC{$r \in \mathbb{R} \to \varphi_{43}$}
\AXC{$\varphi_{137}$\lbl{A1}}
\BIC{$(r \in \mathbb{R} \to 0 \leq r \to \varphi_{76}) \to r \in \mathbb{R} \to 0 \leq r \to \varphi_{113}$}
\BIC{$r \in \mathbb{R} \to 0 \leq r \to \varphi_{113}$}
\end{ptree}
\noindent$\Pi_{10}$:
\begin{ptree}
\AXC{$\forall a\,\forall n\,(\mathrm{Seq}(a) \to n \in \mathbb{N} \to a_{n} \in \mathbb{R})$\lbl{[S1]}}
\AXC{$\varphi_{59}$\lbl{A2}}
\BIC{$\forall n\,(\mathrm{Seq}({(r)^{\bullet}}) \to n \in \mathbb{N} \to {{r}^{n}} \in \mathbb{R})$}
\AXC{$\varphi_{62}$\lbl{A2}}
\BIC{$\mathrm{Seq}({(r)^{\bullet}}) \to n \in \mathbb{N} \to {{r}^{n}} \in \mathbb{R}$}
\end{ptree}
\noindent$\Pi_{11}$:
\begin{ptree}
\AXC{$\stackrel{\Pi_{10}}{\vdots}$}\noLine
\UIC{$\mathrm{Seq}({(r)^{\bullet}}) \to n \in \mathbb{N} \to {{r}^{n}} \in \mathbb{R}$}
\AXC{$(\mathrm{Seq}({(r)^{\bullet}}) \to n \in \mathbb{N} \to {{r}^{n}} \in \mathbb{R}) \to \varphi_{135}$\lbl{A1}}
\BIC{$\varphi_{135}$}
\end{ptree}
\noindent$\Pi_{12}$:
\begin{ptree}
\AXC{$\stackrel{\Pi_{7}}{\vdots}$}\noLine
\UIC{$r \in \mathbb{R} \to 0 \leq r \to \varphi_{74}$}
\AXC{$\stackrel{\Pi_{9}}{\vdots}$}\noLine
\UIC{$r \in \mathbb{R} \to 0 \leq r \to \varphi_{113}$}
\AXC{$(r \in \mathbb{R} \to 0 \leq r \to \varphi_{113}) \to r \in \mathbb{R} \to \varphi_{86}$\lbl{A1}}
\BIC{$r \in \mathbb{R} \to \varphi_{86}$}
\AXC{$\stackrel{\Pi_{11}}{\vdots}$}\noLine
\UIC{$\varphi_{135}$}
\BIC{$(r \in \mathbb{R} \to 0 \leq r \to \varphi_{74}) \to r \in \mathbb{R} \to 0 \leq r \to \varphi_{91}$}
\BIC{$r \in \mathbb{R} \to 0 \leq r \to \varphi_{91}$}
\end{ptree}
\noindent$\Pi_{13}$:
\begin{ptree}
\AXC{$\stackrel{\Pi_{12}}{\vdots}$}\noLine
\UIC{$r \in \mathbb{R} \to 0 \leq r \to \varphi_{91}$}
\AXC{$\forall x\,\forall y\,(x \in \mathbb{R} \to y \in \mathbb{R} \to 0 \leq x \to 0 \leq y \to 0 \leq x \cdot y)$\lbl{[A16]}}
\AXC{$\varphi_{98}$\lbl{A2}}
\BIC{$\forall y\,({{r}^{n}} \in \mathbb{R} \to y \in \mathbb{R} \to 0 \leq {{r}^{n}} \to 0 \leq y \to 0 \leq {{r}^{n}} \cdot y)$}
\AXC{$\varphi_{111}$\lbl{A2}}
\BIC{${{r}^{n}} \in \mathbb{R} \to r \in \mathbb{R} \to 0 \leq {{r}^{n}} \to 0 \leq r \to 0 \leq {{r}^{n}} \cdot r$}
\AXC{$\varphi_{142}$\lbl{A1}}
\BIC{$(r \in \mathbb{R} \to 0 \leq r \to \varphi_{91}) \to (r \in \mathbb{R} \to 0 \leq r \to \varphi_{76}) \to \varphi_{126}$}
\BIC{$(r \in \mathbb{R} \to 0 \leq r \to \varphi_{76}) \to \varphi_{126}$}
\end{ptree}
\noindent$\Pi_{14}$:
\begin{ptree}
\AXC{$r \in \mathbb{R} \to 0 \leq r \to 0 \leq r$\lbl{A1}}
\AXC{$\stackrel{\Pi_{8}}{\vdots}$}\noLine
\UIC{$r \in \mathbb{R} \to 0 \leq r \to \varphi_{84}$}
\AXC{$r \in \mathbb{R} \to 0 \leq r \to \varphi_{76}$\lbl{A1}}
\AXC{$\stackrel{\Pi_{13}}{\vdots}$}\noLine
\UIC{$(r \in \mathbb{R} \to 0 \leq r \to \varphi_{76}) \to \varphi_{126}$}
\BIC{$\varphi_{126}$}
\BIC{$(r \in \mathbb{R} \to 0 \leq r \to 0 \leq r) \to r \in \mathbb{R} \to 0 \leq r \to \varphi_{89}$}
\BIC{$r \in \mathbb{R} \to 0 \leq r \to \varphi_{89}$}
\end{ptree}
\noindent$\Pi_{15}$:
\begin{ptree}
\AXC{$x = x$\lbl{A5}}
\AXC{$x = y \to x = x \to y = x$\lbl{A6}}
\AXC{$(x = y \to x = x \to y = x) \to x = x \to x = y \to y = x$\lbl{A1}}
\BIC{$x = x \to x = y \to y = x$}
\BIC{$x = y \to y = x$}
\UIC{$\forall y\,(x = y \to y = x)$}
\UIC{$\forall x\,\forall y\,(x = y \to y = x)$}
\end{ptree}
\noindent$\Pi_{16}$:
\begin{ptree}
\AXC{$r \in \mathbb{R} \to 0 \leq r \to \varphi_{76}$\lbl{A1}}
\AXC{$\forall r \in \mathbb{R}\,(\varphi_{43})$\lbl{[S24]}}
\AXC{$\forall r \in \mathbb{R}\,(\varphi_{43}) \to r \in \mathbb{R} \to \varphi_{43}$\lbl{A2}}
\BIC{$r \in \mathbb{R} \to \varphi_{43}$}
\AXC{$\varphi_{137}$\lbl{A1}}
\BIC{$(r \in \mathbb{R} \to 0 \leq r \to \varphi_{76}) \to r \in \mathbb{R} \to 0 \leq r \to \varphi_{113}$}
\BIC{$r \in \mathbb{R} \to 0 \leq r \to \varphi_{113}$}
\end{ptree}
\noindent$\Pi_{17}$:
\begin{ptree}
\AXC{$\varphi_{60}$\lbl{A2}}
\AXC{$\stackrel{\Pi_{16}}{\vdots}$}\noLine
\UIC{$r \in \mathbb{R} \to 0 \leq r \to \varphi_{113}$}
\AXC{$(r \in \mathbb{R} \to 0 \leq r \to \varphi_{113}) \to r \in \mathbb{R} \to \varphi_{107}$\lbl{A1}}
\BIC{$r \in \mathbb{R} \to \varphi_{107}$}
\AXC{$(r \in \mathbb{R} \to \varphi_{107}) \to \varphi_{60} \to r \in \mathbb{R} \to \varphi_{106}$\lbl{A1}}
\BIC{$\varphi_{60} \to r \in \mathbb{R} \to \varphi_{106}$}
\BIC{$r \in \mathbb{R} \to \varphi_{106}$}
\end{ptree}
\noindent$\Pi_{18}$:
\begin{ptree}
\AXC{$\stackrel{\Pi_{7}}{\vdots}$}\noLine
\UIC{$r \in \mathbb{R} \to 0 \leq r \to \varphi_{74}$}
\AXC{$\stackrel{\Pi_{17}}{\vdots}$}\noLine
\UIC{$r \in \mathbb{R} \to \varphi_{106}$}
\AXC{$\varphi_{139}$\lbl{A1}}
\BIC{$(r \in \mathbb{R} \to 0 \leq r \to \varphi_{74}) \to r \in \mathbb{R} \to 0 \leq r \to \varphi_{96}$}
\BIC{$r \in \mathbb{R} \to 0 \leq r \to \varphi_{96}$}
\AXC{$(r \in \mathbb{R} \to 0 \leq r \to \varphi_{96}) \to \varphi_{121}$\lbl{A1}}
\BIC{$\varphi_{121}$}
\end{ptree}
\noindent$\Pi_{19}$:
\begin{ptree}
\AXC{$\stackrel{\Pi_{15}}{\vdots}$}\noLine
\UIC{$\forall x\,\forall y\,(x = y \to y = x)$}
\AXC{$\varphi_{16}$\lbl{A2}}
\BIC{$\forall y\,({{r}^{S(n)}} = y \to y = {{r}^{S(n)}})$}
\AXC{$\varphi_{82}$\lbl{A2}}
\BIC{${{r}^{S(n)}} = {{r}^{n}} \cdot r \to {{r}^{n}} \cdot r = {{r}^{S(n)}}$}
\AXC{$\stackrel{\Pi_{18}}{\vdots}$}\noLine
\UIC{$\varphi_{121}$}
\BIC{$r \in \mathbb{R} \to 0 \leq r \to \varphi_{97}$}
\AXC{$(r \in \mathbb{R} \to 0 \leq r \to \varphi_{97}) \to \varphi_{138}$\lbl{A1}}
\BIC{$\varphi_{138}$}
\end{ptree}
\noindent$\Pi_{20}$:
\begin{ptree}
\AXC{$z = w \to 0 < z \to 0 < w$\lbl{A6}}
\UIC{$\forall w\,(z = w \to 0 < z \to 0 < w)$}
\UIC{$\forall z\,\forall w\,(z = w \to 0 < z \to 0 < w)$}
\AXC{$\varphi_{50}$\lbl{A2}}
\BIC{$\forall w\,({{r}^{n}} \cdot r = w \to 0 < {{r}^{n}} \cdot r \to 0 < w)$}
\AXC{$\varphi_{100}$\lbl{A2}}
\BIC{${{r}^{n}} \cdot r = {{r}^{S(n)}} \to 0 < {{r}^{n}} \cdot r \to 0 < {{r}^{S(n)}}$}
\AXC{$\stackrel{\Pi_{14}}{\vdots}$}\noLine
\UIC{$r \in \mathbb{R} \to 0 \leq r \to \varphi_{89}$}
\AXC{$\stackrel{\Pi_{19}}{\vdots}$}\noLine
\UIC{$\varphi_{138}$}
\BIC{$({{r}^{n}} \cdot r = {{r}^{S(n)}} \to 0 < {{r}^{n}} \cdot r \to 0 < {{r}^{S(n)}}) \to \varphi_{119}$}
\BIC{$\varphi_{119}$}
\end{ptree}
\noindent$\Pi_{21}$:
\begin{ptree}
\AXC{$r \in \mathbb{R} \to 0 \leq r \to r \in \mathbb{R}$\lbl{A1}}
\AXC{$\forall r \in \mathbb{R}\,(\varphi_{43})$\lbl{[S24]}}
\AXC{$\forall r \in \mathbb{R}\,(\varphi_{43}) \to r \in \mathbb{R} \to \varphi_{43}$\lbl{A2}}
\BIC{$r \in \mathbb{R} \to \varphi_{43}$}
\AXC{$\varphi_{122}$\lbl{A1}}
\BIC{$(r \in \mathbb{R} \to 0 \leq r \to r \in \mathbb{R}) \to r \in \mathbb{R} \to 0 \leq r \to \varphi_{43}$}
\BIC{$r \in \mathbb{R} \to 0 \leq r \to \varphi_{43}$}
\end{ptree}
\noindent$\Pi_{22}$:
\begin{ptree}
\AXC{$\stackrel{\Pi_{21}}{\vdots}$}\noLine
\UIC{$r \in \mathbb{R} \to 0 \leq r \to \varphi_{43}$}
\AXC{$(r \in \mathbb{R} \to 0 \leq r \to \varphi_{43}) \to r \in \mathbb{R} \to 0 \leq r \to {{r}^{1_{\omega}}} = r$\lbl{A1}}
\BIC{$r \in \mathbb{R} \to 0 \leq r \to {{r}^{1_{\omega}}} = r$}
\AXC{$(r \in \mathbb{R} \to 0 \leq r \to {{r}^{1_{\omega}}} = r) \to \varphi_{38}$\lbl{A1}}
\BIC{$\varphi_{38}$}
\end{ptree}
\noindent$\Pi_{23}$:
\begin{ptree}
\AXC{$\stackrel{\Pi_{15}}{\vdots}$}\noLine
\UIC{$\forall x\,\forall y\,(x = y \to y = x)$}
\AXC{$\varphi_{17}$\lbl{A2}}
\BIC{$\forall y\,({{r}^{1_{\omega}}} = y \to y = {{r}^{1_{\omega}}})$}
\AXC{$\varphi_{41}$\lbl{A2}}
\BIC{${{r}^{1_{\omega}}} = r \to r = {{r}^{1_{\omega}}}$}
\AXC{$\stackrel{\Pi_{22}}{\vdots}$}\noLine
\UIC{$\varphi_{38}$}
\BIC{$r \in \mathbb{R} \to 0 \leq r \to r = {{r}^{1_{\omega}}}$}
\AXC{$(r \in \mathbb{R} \to 0 \leq r \to r = {{r}^{1_{\omega}}}) \to \varphi_{110}$\lbl{A1}}
\BIC{$\varphi_{110}$}
\end{ptree}
\noindent$\Pi_{24}$:
\begin{ptree}
\AXC{$z = w \to 0 < z \to 0 < w$\lbl{A6}}
\UIC{$\forall w\,(z = w \to 0 < z \to 0 < w)$}
\UIC{$\forall z\,\forall w\,(z = w \to 0 < z \to 0 < w)$}
\AXC{$\varphi_{9}$\lbl{A2}}
\BIC{$\forall w\,(r = w \to 0 < r \to 0 < w)$}
\AXC{$\varphi_{37}$\lbl{A2}}
\BIC{$r = {{r}^{1_{\omega}}} \to 0 < r \to 0 < {{r}^{1_{\omega}}}$}
\AXC{$r \in \mathbb{R} \to 0 \leq r \to 0 \leq r$\lbl{A1}}
\AXC{$\stackrel{\Pi_{23}}{\vdots}$}\noLine
\UIC{$\varphi_{110}$}
\BIC{$(r = {{r}^{1_{\omega}}} \to 0 < r \to 0 < {{r}^{1_{\omega}}}) \to \varphi_{45}$}
\BIC{$\varphi_{45}$}
\end{ptree}
\noindent$\Pi_{25}$:
\begin{ptree}
\AXC{$y = z \to x = y \to x = z$\lbl{A6}}
\AXC{$(y = z \to x = y \to x = z) \to x = y \to y = z \to x = z$\lbl{A1}}
\BIC{$x = y \to y = z \to x = z$}
\UIC{$\forall z\,(x = y \to y = z \to x = z)$}
\UIC{$\forall y\,\forall z\,(x = y \to y = z \to x = z)$}
\UIC{$\forall x\,\forall y\,\forall z\,(x = y \to y = z \to x = z)$}
\AXC{$\varphi_{10}$\lbl{A2}}
\BIC{$\forall y\,\forall z\,(S(n) = y \to y = z \to S(n) = z)$}
\AXC{$\varphi_{33}$\lbl{A2}}
\BIC{$\forall z\,(S(n) = S(\emptyset) \to S(\emptyset) = z \to S(n) = z)$}
\end{ptree}
\noindent$\Pi_{26}$:
\begin{ptree}
\AXC{$\stackrel{\Pi_{25}}{\vdots}$}\noLine
\UIC{$\forall z\,(S(n) = S(\emptyset) \to S(\emptyset) = z \to S(n) = z)$}
\AXC{$\forall z\,(S(n) = S(\emptyset) \to S(\emptyset) = z \to S(n) = z) \to \varphi_{5}$\lbl{A2}}
\BIC{$\varphi_{5}$}
\end{ptree}
\noindent$\Pi_{27}$:
\begin{ptree}
\AXC{$\stackrel{\Pi_{15}}{\vdots}$}\noLine
\UIC{$\forall x\,\forall y\,(x = y \to y = x)$}
\AXC{$\varphi_{2}$\lbl{A2}}
\BIC{$\forall y\,(1_{\omega} = y \to y = 1_{\omega})$}
\AXC{$\varphi_{12}$\lbl{A2}}
\BIC{$1_{\omega} = S(\emptyset) \to S(\emptyset) = 1_{\omega}$}
\AXC{$1_{\omega} = S(\emptyset)$\lbl{Def$1_\omega$}}
\AXC{$1_{\omega} = S(\emptyset) \to \varphi_{11}$\lbl{A1}}
\BIC{$\varphi_{11}$}
\BIC{$S(\emptyset) = 1_{\omega}$}
\end{ptree}
\noindent$\Pi_{28}$:
\begin{ptree}
\AXC{$r \in \mathbb{R} \to 0 \leq r \to \varphi_{75}$\lbl{A1}}
\AXC{$(r \in \mathbb{R} \to 0 \leq r \to \varphi_{75}) \to S(n) = S(n) \to \varphi_{112}$\lbl{A1}}
\BIC{$S(n) = S(n) \to \varphi_{112}$}
\end{ptree}
\noindent$\Pi_{29}$:
\begin{ptree}
\AXC{$x = x$\lbl{A5}}
\UIC{$\forall x\,(x = x)$}
\AXC{$\forall x\,(x = x) \to S(n) = S(n)$\lbl{A2}}
\BIC{$S(n) = S(n)$}
\AXC{$\stackrel{\Pi_{28}}{\vdots}$}\noLine
\UIC{$S(n) = S(n) \to \varphi_{112}$}
\BIC{$\varphi_{112}$}
\end{ptree}
\noindent$\Pi_{30}$:
\begin{ptree}
\AXC{$z = w \to S(n) = S(z) \to S(n) = S(w)$\lbl{A6}}
\UIC{$\forall w\,(z = w \to S(n) = S(z) \to S(n) = S(w))$}
\UIC{$\forall z\,\forall w\,(z = w \to S(n) = S(z) \to S(n) = S(w))$}
\AXC{$\varphi_{26}$\lbl{A2}}
\BIC{$\forall w\,(n = w \to S(n) = S(n) \to S(n) = S(w))$}
\AXC{$\varphi_{34}$\lbl{A2}}
\BIC{$n = \emptyset \to S(n) = S(n) \to S(n) = S(\emptyset)$}
\AXC{$\stackrel{\Pi_{29}}{\vdots}$}\noLine
\UIC{$\varphi_{112}$}
\BIC{$r \in \mathbb{R} \to 0 \leq r \to \varphi_{85}$}
\AXC{$\varphi_{128}$\lbl{A1}}
\BIC{$S(\emptyset) = 1_{\omega} \to \varphi_{5} \to r \in \mathbb{R} \to 0 \leq r \to \varphi_{72}$}
\end{ptree}
\noindent$\Pi_{31}$:
\begin{ptree}
\AXC{$\stackrel{\Pi_{26}}{\vdots}$}\noLine
\UIC{$\varphi_{5}$}
\AXC{$\stackrel{\Pi_{27}}{\vdots}$}\noLine
\UIC{$S(\emptyset) = 1_{\omega}$}
\AXC{$\stackrel{\Pi_{30}}{\vdots}$}\noLine
\UIC{$S(\emptyset) = 1_{\omega} \to \varphi_{5} \to r \in \mathbb{R} \to 0 \leq r \to \varphi_{72}$}
\BIC{$\varphi_{5} \to r \in \mathbb{R} \to 0 \leq r \to \varphi_{72}$}
\BIC{$r \in \mathbb{R} \to 0 \leq r \to \varphi_{72}$}
\AXC{$(r \in \mathbb{R} \to 0 \leq r \to \varphi_{72}) \to \varphi_{105}$\lbl{A1}}
\BIC{$\varphi_{105}$}
\end{ptree}
\noindent$\Pi_{32}$:
\begin{ptree}
\AXC{$\stackrel{\Pi_{15}}{\vdots}$}\noLine
\UIC{$\forall x\,\forall y\,(x = y \to y = x)$}
\AXC{$\varphi_{1}$\lbl{A2}}
\BIC{$\forall y\,(S(n) = y \to y = S(n))$}
\AXC{$\varphi_{4}$\lbl{A2}}
\BIC{$S(n) = 1_{\omega} \to 1_{\omega} = S(n)$}
\AXC{$\stackrel{\Pi_{31}}{\vdots}$}\noLine
\UIC{$\varphi_{105}$}
\BIC{$r \in \mathbb{R} \to 0 \leq r \to \varphi_{73}$}
\AXC{$(r \in \mathbb{R} \to 0 \leq r \to \varphi_{73}) \to \varphi_{125}$\lbl{A1}}
\BIC{$\varphi_{125}$}
\end{ptree}
\noindent$\Pi_{33}$:
\begin{ptree}
\AXC{$z = w \to 0 = z \to 0 = w$\lbl{A6}}
\UIC{$\forall w\,(z = w \to 0 = z \to 0 = w)$}
\UIC{$\forall z\,\forall w\,(z = w \to 0 = z \to 0 = w)$}
\AXC{$\varphi_{3}$\lbl{A2}}
\BIC{$\forall w\,(r = w \to 0 = r \to 0 = w)$}
\AXC{$\varphi_{24}$\lbl{A2}}
\BIC{$r = {{r}^{1_{\omega}}} \to 0 = r \to 0 = {{r}^{1_{\omega}}}$}
\AXC{$\stackrel{\Pi_{24}}{\vdots}$}\noLine
\UIC{$\varphi_{45}$}
\BIC{$r \in \mathbb{R} \to 0 \leq r \to 0 \leq {{r}^{1_{\omega}}}$}
\AXC{$\stackrel{\Pi_{32}}{\vdots}$}\noLine
\UIC{$\varphi_{125}$}
\BIC{$(1_{\omega} = S(n) \to 0 < {{r}^{1_{\omega}}} \to 0 < {{r}^{S(n)}}) \to \varphi_{115}$}
\end{ptree}
\noindent$\Pi_{34}$:
\begin{ptree}
\AXC{$z = w \to 0 < {{r}^{z}} \to 0 < {{r}^{w}}$\lbl{A6}}
\UIC{$\forall w\,(z = w \to 0 < {{r}^{z}} \to 0 < {{r}^{w}})$}
\UIC{$\forall z\,\forall w\,(z = w \to 0 < {{r}^{z}} \to 0 < {{r}^{w}})$}
\AXC{$\varphi_{71}$\lbl{A2}}
\BIC{$\forall w\,(1_{\omega} = w \to 0 < {{r}^{1_{\omega}}} \to 0 < {{r}^{w}})$}
\AXC{$\varphi_{78}$\lbl{A2}}
\BIC{$1_{\omega} = S(n) \to 0 < {{r}^{1_{\omega}}} \to 0 < {{r}^{S(n)}}$}
\AXC{$\stackrel{\Pi_{33}}{\vdots}$}\noLine
\UIC{$(1_{\omega} = S(n) \to 0 < {{r}^{1_{\omega}}} \to 0 < {{r}^{S(n)}}) \to \varphi_{115}$}
\BIC{$\varphi_{115}$}
\end{ptree}
\noindent$\Pi_{35}$:
\begin{ptree}
\AXC{$z = w \to 0 = {{r}^{z}} \to 0 = {{r}^{w}}$\lbl{A6}}
\UIC{$\forall w\,(z = w \to 0 = {{r}^{z}} \to 0 = {{r}^{w}})$}
\UIC{$\forall z\,\forall w\,(z = w \to 0 = {{r}^{z}} \to 0 = {{r}^{w}})$}
\AXC{$\varphi_{56}$\lbl{A2}}
\BIC{$\forall w\,(1_{\omega} = w \to 0 = {{r}^{1_{\omega}}} \to 0 = {{r}^{w}})$}
\AXC{$\varphi_{63}$\lbl{A2}}
\BIC{$1_{\omega} = S(n) \to 0 = {{r}^{1_{\omega}}} \to 0 = {{r}^{S(n)}}$}
\AXC{$\stackrel{\Pi_{34}}{\vdots}$}\noLine
\UIC{$\varphi_{115}$}
\BIC{$r \in \mathbb{R} \to 0 \leq r \to \varphi_{87}$}
\AXC{$\varphi_{134}$\lbl{A1}}
\BIC{$(r \in \mathbb{R} \to 0 \leq r \to \varphi_{88}) \to r \in \mathbb{R} \to 0 \leq r \to \varphi_{70}$}
\end{ptree}
\noindent$\Pi_{36}$:
\begin{ptree}
\AXC{$z = w \to 0 = z \to 0 = w$\lbl{A6}}
\UIC{$\forall w\,(z = w \to 0 = z \to 0 = w)$}
\UIC{$\forall z\,\forall w\,(z = w \to 0 = z \to 0 = w)$}
\AXC{$\varphi_{32}$\lbl{A2}}
\BIC{$\forall w\,({{r}^{n}} \cdot r = w \to 0 = {{r}^{n}} \cdot r \to 0 = w)$}
\AXC{$\varphi_{94}$\lbl{A2}}
\BIC{${{r}^{n}} \cdot r = {{r}^{S(n)}} \to 0 = {{r}^{n}} \cdot r \to 0 = {{r}^{S(n)}}$}
\AXC{$\stackrel{\Pi_{20}}{\vdots}$}\noLine
\UIC{$\varphi_{119}$}
\BIC{$r \in \mathbb{R} \to 0 \leq r \to \varphi_{88}$}
\AXC{$\stackrel{\Pi_{35}}{\vdots}$}\noLine
\UIC{$(r \in \mathbb{R} \to 0 \leq r \to \varphi_{88}) \to r \in \mathbb{R} \to 0 \leq r \to \varphi_{70}$}
\BIC{$r \in \mathbb{R} \to 0 \leq r \to \varphi_{70}$}
\end{ptree}
\noindent$\Pi_{37}$:
\begin{ptree}
\AXC{$\forall x\,(\mathrm{Ind}(x) \leftrightarrow \emptyset \in x \land \forall y \in x\,(S(y) \in x))$\lbl{DefInd}}
\AXC{$\varphi_{90}$\lbl{A2}}
\BIC{$\mathrm{Ind}(\omega) \leftrightarrow \emptyset \in \omega \land \forall y \in \omega\,(S(y) \in \omega)$}
\AXC{$\mathrm{Ind}(\omega)$\lbl{補題(ωは帰納的)}}
\AXC{$\mathrm{Ind}(\omega) \to \varphi_{117}$\lbl{A1}}
\BIC{$\varphi_{117}$}
\BIC{$r \in \mathbb{R} \to 0 \leq r \to \varphi_{81}$}
\end{ptree}
\noindent$\Pi_{38}$:
\begin{ptree}
\AXC{$\stackrel{\Pi_{37}}{\vdots}$}\noLine
\UIC{$r \in \mathbb{R} \to 0 \leq r \to \varphi_{81}$}
\AXC{$(r \in \mathbb{R} \to 0 \leq r \to \varphi_{81}) \to r \in \mathbb{R} \to 0 \leq r \to \varphi_{67}$\lbl{A1}}
\BIC{$r \in \mathbb{R} \to 0 \leq r \to \varphi_{67}$}
\AXC{$(r \in \mathbb{R} \to 0 \leq r \to \varphi_{67}) \to \varphi_{114}$\lbl{A1}}
\BIC{$\varphi_{114}$}
\end{ptree}
\noindent$\Pi_{39}$:
\begin{ptree}
\AXC{$\stackrel{\Pi_{5}}{\vdots}$}\noLine
\UIC{$r \in \mathbb{R} \to 0 \leq r \to \varphi_{42}$}
\AXC{$\forall y \in \omega\,(S(y) \in \omega) \to n \in \omega \to S(n) \in \omega$\lbl{A2}}
\AXC{$\stackrel{\Pi_{38}}{\vdots}$}\noLine
\UIC{$\varphi_{114}$}
\BIC{$r \in \mathbb{R} \to 0 \leq r \to \varphi_{64}$}
\AXC{$\varphi_{127}$\lbl{A1}}
\BIC{$(r \in \mathbb{R} \to 0 \leq r \to \varphi_{42}) \to r \in \mathbb{R} \to 0 \leq r \to \varphi_{48}$}
\BIC{$r \in \mathbb{R} \to 0 \leq r \to \varphi_{48}$}
\end{ptree}
\noindent$\Pi_{40}$:
\begin{ptree}
\AXC{$\forall n\,(n \in B_{29} \leftrightarrow n \in \omega \land (n \in \mathbb{N} \to 0 \leq {{r}^{n}})) \to \varphi_{25}$\lbl{A2}}
\AXC{$r \in \mathbb{R} \to 0 \leq r \to \varphi_{83}$\lbl{A1}}
\AXC{$(r \in \mathbb{R} \to 0 \leq r \to \varphi_{83}) \to \varphi_{124}$\lbl{A1}}
\BIC{$\varphi_{124}$}
\BIC{$r \in \mathbb{R} \to \varphi_{95}$}
\end{ptree}
\noindent$\Pi_{41}$:
\begin{ptree}
\AXC{$\stackrel{\Pi_{36}}{\vdots}$}\noLine
\UIC{$r \in \mathbb{R} \to 0 \leq r \to \varphi_{70}$}
\AXC{$\stackrel{\Pi_{39}}{\vdots}$}\noLine
\UIC{$r \in \mathbb{R} \to 0 \leq r \to \varphi_{48}$}
\AXC{$\stackrel{\Pi_{40}}{\vdots}$}\noLine
\UIC{$r \in \mathbb{R} \to \varphi_{95}$}
\AXC{$(r \in \mathbb{R} \to \varphi_{95}) \to \varphi_{129}$\lbl{A1}}
\BIC{$\varphi_{129}$}
\BIC{$(r \in \mathbb{R} \to 0 \leq r \to \varphi_{70}) \to r \in \mathbb{R} \to 0 \leq r \to \varphi_{46}$}
\BIC{$r \in \mathbb{R} \to 0 \leq r \to \varphi_{46}$}
\end{ptree}
\noindent$\Pi_{42}$:
\begin{ptree}
\AXC{$\varphi_{120}$\lbl{A3}}
\AXC{$\stackrel{\Pi_{41}}{\vdots}$}\noLine
\UIC{$r \in \mathbb{R} \to 0 \leq r \to \varphi_{46}$}
\AXC{$(r \in \mathbb{R} \to 0 \leq r \to \varphi_{46}) \to \varphi_{35} \to n \in B_{29} \to S(n) \in B_{29}$\lbl{A1}}
\BIC{$\varphi_{35} \to n \in B_{29} \to S(n) \in B_{29}$}
\UIC{$\forall n\,(\varphi_{35} \to n \in B_{29} \to S(n) \in B_{29})$}
\AXC{$\varphi_{140}$\lbl{A1}}
\BIC{$\varphi_{120} \to (\varphi_{35} \to \forall n\,(\varphi_{35})) \to r \in \mathbb{R} \to 0 \leq r \to \varphi_{49}$}
\BIC{$(\varphi_{35} \to \forall n\,(\varphi_{35})) \to r \in \mathbb{R} \to 0 \leq r \to \varphi_{49}$}
\end{ptree}
\noindent$\Pi_{43}$:
\begin{ptree}
\AXC{$x = x$\lbl{A5}}
\UIC{$\forall x\,(x = x)$}
\AXC{$\forall x\,(x = x) \to \emptyset = \emptyset$\lbl{A2}}
\BIC{$\emptyset = \emptyset$}
\AXC{$\forall n\,(n \in \mathbb{N} \leftrightarrow n \in \omega \land n \neq \emptyset)$\lbl{Def$\mathbb{N}$}}
\AXC{$\varphi_{66}$\lbl{A2}}
\BIC{$\emptyset \in \mathbb{N} \leftrightarrow \emptyset \in \omega \land \emptyset \neq \emptyset$}
\AXC{$\varphi_{101}$\lbl{A1}}
\BIC{$\emptyset = \emptyset \to r \in \mathbb{R} \to 0 \leq r \to \emptyset \in \mathbb{N} \to 0 \leq {{r}^{\emptyset}}$}
\BIC{$r \in \mathbb{R} \to 0 \leq r \to \emptyset \in \mathbb{N} \to 0 \leq {{r}^{\emptyset}}$}
\end{ptree}
\noindent$\Pi_{44}$:
\begin{ptree}
\AXC{$\varphi_{99}$\lbl{A2}}
\AXC{$r \in \mathbb{R} \to 0 \leq r \to \varphi_{83}$\lbl{A1}}
\AXC{$(r \in \mathbb{R} \to 0 \leq r \to \varphi_{83}) \to \varphi_{99} \to r \in \mathbb{R} \to 0 \leq r \to \varphi_{99}$\lbl{A1}}
\BIC{$\varphi_{99} \to r \in \mathbb{R} \to 0 \leq r \to \varphi_{99}$}
\BIC{$r \in \mathbb{R} \to 0 \leq r \to \varphi_{99}$}
\end{ptree}
\noindent$\Pi_{45}$:
\begin{ptree}
\AXC{$\stackrel{\Pi_{37}}{\vdots}$}\noLine
\UIC{$r \in \mathbb{R} \to 0 \leq r \to \varphi_{81}$}
\AXC{$\stackrel{\Pi_{44}}{\vdots}$}\noLine
\UIC{$r \in \mathbb{R} \to 0 \leq r \to \varphi_{99}$}
\AXC{$(r \in \mathbb{R} \to 0 \leq r \to \varphi_{99}) \to (r \in \mathbb{R} \to 0 \leq r \to \varphi_{81}) \to \varphi_{108}$\lbl{A1}}
\BIC{$(r \in \mathbb{R} \to 0 \leq r \to \varphi_{81}) \to \varphi_{108}$}
\BIC{$\varphi_{108}$}
\end{ptree}
\noindent$\Pi_{46}$:
\begin{ptree}
\AXC{$\forall x\,(\mathrm{Ind}(x) \leftrightarrow \emptyset \in x \land \forall y \in x\,(S(y) \in x))$\lbl{DefInd}}
\AXC{$\forall x\,(\mathrm{Ind}(x) \leftrightarrow \emptyset \in x \land \forall y \in x\,(S(y) \in x)) \to \varphi_{15}$\lbl{A2}}
\BIC{$\varphi_{15}$}
\AXC{$\varphi_{15} \to \varphi_{136}$\lbl{A1}}
\BIC{$\varphi_{136}$}
\end{ptree}
\noindent$\Pi_{47}$:
\begin{ptree}
\AXC{$\stackrel{\Pi_{43}}{\vdots}$}\noLine
\UIC{$r \in \mathbb{R} \to 0 \leq r \to \emptyset \in \mathbb{N} \to 0 \leq {{r}^{\emptyset}}$}
\AXC{$\stackrel{\Pi_{45}}{\vdots}$}\noLine
\UIC{$\varphi_{108}$}
\BIC{$r \in \mathbb{R} \to 0 \leq r \to \varphi_{36}$}
\AXC{$\stackrel{\Pi_{46}}{\vdots}$}\noLine
\UIC{$\varphi_{136}$}
\BIC{$(r \in \mathbb{R} \to 0 \leq r \to \varphi_{49}) \to \varphi_{14} \to \varphi_{13} \to r \in \mathbb{R} \to 0 \leq r \to \varphi_{27}$}
\end{ptree}
\noindent$\Pi_{48}$:
\begin{ptree}
\AXC{$\varphi_{35} \to \forall n\,(\varphi_{35})$\lbl{A4}}
\AXC{$\stackrel{\Pi_{42}}{\vdots}$}\noLine
\UIC{$(\varphi_{35} \to \forall n\,(\varphi_{35})) \to r \in \mathbb{R} \to 0 \leq r \to \varphi_{49}$}
\BIC{$r \in \mathbb{R} \to 0 \leq r \to \varphi_{49}$}
\AXC{$\stackrel{\Pi_{47}}{\vdots}$}\noLine
\UIC{$(r \in \mathbb{R} \to 0 \leq r \to \varphi_{49}) \to \varphi_{14} \to \varphi_{13} \to r \in \mathbb{R} \to 0 \leq r \to \varphi_{27}$}
\BIC{$\varphi_{14} \to \varphi_{13} \to r \in \mathbb{R} \to 0 \leq r \to \varphi_{27}$}
\end{ptree}
\noindent$\Pi_{49}$:
\begin{ptree}
\AXC{$\varphi_{19}$\lbl{A4}}
\AXC{$\forall n\,(\varphi_{8}) \to \varphi_{21}$\lbl{A3}}
\AXC{$\varphi_{8}$\lbl{A2}}
\UIC{$\forall n\,(\varphi_{8})$}
\AXC{$\forall n\,(\varphi_{8}) \to (\forall n\,(\varphi_{8}) \to \varphi_{21}) \to \varphi_{19} \to \varphi_{14}$\lbl{A1}}
\BIC{$(\forall n\,(\varphi_{8}) \to \varphi_{21}) \to \varphi_{19} \to \varphi_{14}$}
\BIC{$\varphi_{19} \to \varphi_{14}$}
\BIC{$\varphi_{14}$}
\AXC{$\stackrel{\Pi_{48}}{\vdots}$}\noLine
\UIC{$\varphi_{14} \to \varphi_{13} \to r \in \mathbb{R} \to 0 \leq r \to \varphi_{27}$}
\BIC{$\varphi_{13} \to r \in \mathbb{R} \to 0 \leq r \to \varphi_{27}$}
\end{ptree}
\noindent$\Pi_{50}$:
\begin{ptree}
\AXC{$\varphi_{18}$\lbl{A4}}
\AXC{$\forall y\,(\varphi_{7}) \to \varphi_{20}$\lbl{A3}}
\AXC{$\varphi_{7}$\lbl{A2}}
\UIC{$\forall y\,(\varphi_{7})$}
\AXC{$\forall y\,(\varphi_{7}) \to (\forall y\,(\varphi_{7}) \to \varphi_{20}) \to \varphi_{18} \to \varphi_{13}$\lbl{A1}}
\BIC{$(\forall y\,(\varphi_{7}) \to \varphi_{20}) \to \varphi_{18} \to \varphi_{13}$}
\BIC{$\varphi_{18} \to \varphi_{13}$}
\BIC{$\varphi_{13}$}
\AXC{$\stackrel{\Pi_{49}}{\vdots}$}\noLine
\UIC{$\varphi_{13} \to r \in \mathbb{R} \to 0 \leq r \to \varphi_{27}$}
\BIC{$r \in \mathbb{R} \to 0 \leq r \to \varphi_{27}$}
\end{ptree}
\noindent$\Pi_{51}$:
\begin{ptree}
\AXC{$\forall n\,(n \in \omega \leftrightarrow \forall x\,(\mathrm{Ind}(x) \to n \in x))$\lbl{Def$\omega$}}
\AXC{$\varphi_{31}$\lbl{A2}}
\BIC{$n \in \omega \leftrightarrow \forall x\,(\mathrm{Ind}(x) \to n \in x)$}
\AXC{$r \in \mathbb{R} \to 0 \leq r \to \varphi_{47}$\lbl{A1}}
\AXC{$(r \in \mathbb{R} \to 0 \leq r \to \varphi_{47}) \to \varphi_{103}$\lbl{A1}}
\BIC{$\varphi_{103}$}
\BIC{$r \in \mathbb{R} \to 0 \leq r \to \varphi_{55}$}
\end{ptree}
\noindent$\Pi_{52}$:
\begin{ptree}
\AXC{$\forall x\,(\mathrm{Ind}(x) \to n \in x) \to \mathrm{Ind}(B_{29}) \to n \in B_{29}$\lbl{A2}}
\AXC{$\stackrel{\Pi_{51}}{\vdots}$}\noLine
\UIC{$r \in \mathbb{R} \to 0 \leq r \to \varphi_{55}$}
\AXC{$(r \in \mathbb{R} \to 0 \leq r \to \varphi_{55}) \to \varphi_{109}$\lbl{A1}}
\BIC{$\varphi_{109}$}
\BIC{$r \in \mathbb{R} \to 0 \leq r \to \varphi_{57}$}
\end{ptree}
\noindent$\Pi_{53}$:
\begin{ptree}
\AXC{$\stackrel{\Pi_{2}}{\vdots}$}\noLine
\UIC{$r \in \mathbb{R} \to 0 \leq r \to \varphi_{77}$}
\AXC{$\stackrel{\Pi_{50}}{\vdots}$}\noLine
\UIC{$r \in \mathbb{R} \to 0 \leq r \to \varphi_{27}$}
\AXC{$\stackrel{\Pi_{52}}{\vdots}$}\noLine
\UIC{$r \in \mathbb{R} \to 0 \leq r \to \varphi_{57}$}
\AXC{$(r \in \mathbb{R} \to 0 \leq r \to \varphi_{57}) \to \varphi_{130}$\lbl{A1}}
\BIC{$\varphi_{130}$}
\BIC{$(r \in \mathbb{R} \to 0 \leq r \to \varphi_{77}) \to r \in \mathbb{R} \to 0 \leq r \to \varphi_{65}$}
\BIC{$r \in \mathbb{R} \to 0 \leq r \to \varphi_{65}$}
\end{ptree}
\noindent$\Pi_{54}$:
\begin{ptree}
\AXC{$\varphi_{123}$\lbl{A3}}
\AXC{$\stackrel{\Pi_{53}}{\vdots}$}\noLine
\UIC{$r \in \mathbb{R} \to 0 \leq r \to \varphi_{65}$}
\AXC{$(r \in \mathbb{R} \to 0 \leq r \to \varphi_{65}) \to \varphi_{35} \to n \in \omega \to n \in \mathbb{N} \to 0 \leq {{r}^{n}}$\lbl{A1}}
\BIC{$\varphi_{35} \to n \in \omega \to n \in \mathbb{N} \to 0 \leq {{r}^{n}}$}
\UIC{$\forall n\,(\varphi_{35} \to n \in \omega \to n \in \mathbb{N} \to 0 \leq {{r}^{n}})$}
\AXC{$\varphi_{141}$\lbl{A1}}
\BIC{$\varphi_{123} \to (\varphi_{35} \to \forall n\,(\varphi_{35})) \to r \in \mathbb{R} \to 0 \leq r \to \varphi_{68}$}
\BIC{$(\varphi_{35} \to \forall n\,(\varphi_{35})) \to r \in \mathbb{R} \to 0 \leq r \to \varphi_{68}$}
\end{ptree}
\noindent$\Pi_{55}$:
\begin{ptree}
\AXC{$\varphi_{35} \to \forall n\,(\varphi_{35})$\lbl{A4}}
\AXC{$\stackrel{\Pi_{54}}{\vdots}$}\noLine
\UIC{$(\varphi_{35} \to \forall n\,(\varphi_{35})) \to r \in \mathbb{R} \to 0 \leq r \to \varphi_{68}$}
\BIC{$r \in \mathbb{R} \to 0 \leq r \to \varphi_{68}$}
\AXC{$(r \in \mathbb{R} \to 0 \leq r \to \varphi_{68}) \to \varphi_{92}$\lbl{A1}}
\BIC{$\varphi_{92}$}
\UIC{$\forall B_{29}\,(\varphi_{92})$}
\end{ptree}
\noindent$\Pi_{56}$:
\begin{ptree}
\AXC{$\forall B_{29}\,(\varphi_{92}) \to \varphi_{29} \to \varphi_{22}$\lbl{A3}}
\AXC{$\stackrel{\Pi_{55}}{\vdots}$}\noLine
\UIC{$\forall B_{29}\,(\varphi_{92})$}
\AXC{$\varphi_{30}$\lbl{Sep}}
\AXC{$\varphi_{30} \to \varphi_{23}$\lbl{A2}}
\BIC{$\varphi_{23}$}
\AXC{$\varphi_{93}$\lbl{A2}}
\BIC{$\exists B_{29}\,\forall n\,(n \in B_{29} \leftrightarrow n \in \omega \land (n \in \mathbb{N} \to 0 \leq {{r}^{n}}))$}
\AXC{$\varphi_{102}$\lbl{A1}}
\BIC{$r \in \mathbb{R} \to \varphi_{28}$}
\AXC{$\varphi_{143}$\lbl{A1}}
\BIC{$\forall B_{29}\,(\varphi_{92}) \to (\forall B_{29}\,(\varphi_{92}) \to \varphi_{29} \to \varphi_{22}) \to \varphi_{118}$}
\BIC{$(\forall B_{29}\,(\varphi_{92}) \to \varphi_{29} \to \varphi_{22}) \to \varphi_{118}$}
\BIC{$\varphi_{118}$}
\end{ptree}
\noindent$\Pi_{57}$:
\begin{ptree}
\AXC{$r \in \mathbb{R} \land 0 \leq r \land \neg \forall n \in \omega\,(n \in \mathbb{N} \to 0 \leq {{r}^{n}}) \to \varphi_{29}$\lbl{A4}}
\AXC{$\stackrel{\Pi_{56}}{\vdots}$}\noLine
\UIC{$\varphi_{118}$}
\BIC{$r \in \mathbb{R} \to 0 \leq r \to \forall n \in \omega\,(n \in \mathbb{N} \to 0 \leq {{r}^{n}})$}
\AXC{$(r \in \mathbb{R} \to 0 \leq r \to \forall n \in \omega\,(n \in \mathbb{N} \to 0 \leq {{r}^{n}})) \to \varphi_{104}$\lbl{A1}}
\BIC{$\varphi_{104}$}
\end{ptree}
\noindent$\Pi_{58}$:
\begin{ptree}
\AXC{$\forall n \in \omega\,(n \in \mathbb{N} \to 0 \leq {{r}^{n}}) \to n \in \omega \to n \in \mathbb{N} \to 0 \leq {{r}^{n}}$\lbl{A2}}
\AXC{$\stackrel{\Pi_{57}}{\vdots}$}\noLine
\UIC{$\varphi_{104}$}
\BIC{$r \in \mathbb{R} \to 0 \leq r \to n \in \omega \to n \in \mathbb{N} \to 0 \leq {{r}^{n}}$}
\AXC{$\varphi_{116}$\lbl{A1}}
\BIC{$(r \in \mathbb{R} \to 0 \leq r \to n \in \mathbb{N} \to n \in \omega) \to \varphi_{52}$}
\end{ptree}
\noindent$\Pi_{59}$:
\begin{ptree}
\AXC{$\forall n\,(n \in \mathbb{N} \leftrightarrow n \in \omega \land n \neq \emptyset)$\lbl{Def$\mathbb{N}$}}
\AXC{$\varphi_{40}$\lbl{A2}}
\BIC{$n \in \mathbb{N} \leftrightarrow n \in \omega \land n \neq \emptyset$}
\AXC{$\stackrel{\Pi_{1}}{\vdots}$}\noLine
\UIC{$\varphi_{44}$}
\BIC{$r \in \mathbb{R} \to 0 \leq r \to n \in \mathbb{N} \to n \in \omega$}
\AXC{$\stackrel{\Pi_{58}}{\vdots}$}\noLine
\UIC{$(r \in \mathbb{R} \to 0 \leq r \to n \in \mathbb{N} \to n \in \omega) \to \varphi_{52}$}
\BIC{$\varphi_{52}$}
\end{ptree}
\noindent$\Pi_{60}$:
\begin{ptree}
\AXC{$r \in \mathbb{R} \to 0 \leq r \to n \in \mathbb{N} \to n \in \mathbb{N}$\lbl{A1}}
\AXC{$\stackrel{\Pi_{59}}{\vdots}$}\noLine
\UIC{$\varphi_{52}$}
\BIC{$r \in \mathbb{R} \to 0 \leq r \to n \in \mathbb{N} \to 0 \leq {{r}^{n}}$}
\AXC{$\varphi_{58}$\lbl{A1}}
\BIC{$r \in \mathbb{R} \land 0 \leq r \to n \in \mathbb{N} \to 0 \leq {{r}^{n}}$}
\UIC{$\forall n\,(r \in \mathbb{R} \land 0 \leq r \to n \in \mathbb{N} \to 0 \leq {{r}^{n}})$}
\AXC{$\forall n\,(r \in \mathbb{R} \land 0 \leq r \to n \in \mathbb{N} \to 0 \leq {{r}^{n}}) \to \varphi_{69} \to \varphi_{53}$\lbl{A1}}
\BIC{$\varphi_{69} \to \varphi_{53}$}
\end{ptree}
\noindent 主木:
\begin{ptree}
\AXC{$r \in \mathbb{R} \land 0 \leq r \to \forall n\,(r \in \mathbb{R} \land 0 \leq r)$\lbl{A4}}
\AXC{$\varphi_{69}$\lbl{A3}}
\AXC{$\stackrel{\Pi_{60}}{\vdots}$}\noLine
\UIC{$\varphi_{69} \to \varphi_{53}$}
\BIC{$\varphi_{53}$}
\BIC{$r \in \mathbb{R} \to 0 \leq r \to \forall n \in \mathbb{N}\,(0 \leq {{r}^{n}})$}
\UIC{$\forall r \in \mathbb{R}\,(0 \leq r \to \forall n \in \mathbb{N}\,(0 \leq {{r}^{n}}))$}
\end{ptree}
\begin{abbrevlist}
\item $\varphi_{1} :\equiv \forall x\,\forall y\,(x = y \to y = x) \to \forall y\,(S(n) = y \to y = S(n))$
\item $\varphi_{2} :\equiv \forall x\,\forall y\,(x = y \to y = x) \to \forall y\,(1_{\omega} = y \to y = 1_{\omega})$
\item $\varphi_{3} :\equiv \forall z\,\forall w\,(z = w \to 0 = z \to 0 = w) \to \forall w\,(r = w \to 0 = r \to 0 = w)$
\item $\varphi_{4} :\equiv \forall y\,(S(n) = y \to y = S(n)) \to S(n) = 1_{\omega} \to 1_{\omega} = S(n)$
\item $\varphi_{5} :\equiv S(n) = S(\emptyset) \to S(\emptyset) = 1_{\omega} \to S(n) = 1_{\omega}$
\item $\varphi_{6} :\equiv n \in B_{29} \to S(n) \in \mathbb{N} \to n \neq \emptyset \to \mathrm{Seq}({(r)^{\bullet}})$
\item $\varphi_{7} :\equiv \forall n \in B_{29}\,(S(n) \in B_{29}) \to y \in B_{29} \to S(y) \in B_{29}$
\item $\varphi_{8} :\equiv \forall y \in B_{29}\,(S(y) \in B_{29}) \to n \in B_{29} \to S(n) \in B_{29}$
\item $\varphi_{9} :\equiv \forall z\,\forall w\,(z = w \to 0 < z \to 0 < w) \to \forall w\,(r = w \to 0 < r \to 0 < w)$
\item $\varphi_{10} :\equiv \forall x\,\forall y\,\forall z\,(x = y \to y = z \to x = z) \to \forall y\,\forall z\,(S(n) = y \to y = z \to S(n) = z)$
\item $\varphi_{11} :\equiv (1_{\omega} = S(\emptyset) \to S(\emptyset) = 1_{\omega}) \to S(\emptyset) = 1_{\omega}$
\item $\varphi_{12} :\equiv \forall y\,(1_{\omega} = y \to y = 1_{\omega}) \to 1_{\omega} = S(\emptyset) \to S(\emptyset) = 1_{\omega}$
\item $\varphi_{13} :\equiv \forall n \in B_{29}\,(S(n) \in B_{29}) \to \forall y \in B_{29}\,(S(y) \in B_{29})$
\item $\varphi_{14} :\equiv \forall y \in B_{29}\,(S(y) \in B_{29}) \to \forall n \in B_{29}\,(S(n) \in B_{29})$
\item $\varphi_{15} :\equiv \mathrm{Ind}(B_{29}) \leftrightarrow \emptyset \in B_{29} \land \forall y \in B_{29}\,(S(y) \in B_{29})$
\item $\varphi_{16} :\equiv \forall x\,\forall y\,(x = y \to y = x) \to \forall y\,({{r}^{S(n)}} = y \to y = {{r}^{S(n)}})$
\item $\varphi_{17} :\equiv \forall x\,\forall y\,(x = y \to y = x) \to \forall y\,({{r}^{1_{\omega}}} = y \to y = {{r}^{1_{\omega}}})$
\item $\varphi_{18} :\equiv \forall n \in B_{29}\,(S(n) \in B_{29}) \to \forall y\,\forall n \in B_{29}\,(S(n) \in B_{29})$
\item $\varphi_{19} :\equiv \forall y \in B_{29}\,(S(y) \in B_{29}) \to \forall n\,\forall y \in B_{29}\,(S(y) \in B_{29})$
\item $\varphi_{20} :\equiv \forall y\,\forall n \in B_{29}\,(S(n) \in B_{29}) \to \forall y \in B_{29}\,(S(y) \in B_{29})$
\item $\varphi_{21} :\equiv \forall n\,\forall y \in B_{29}\,(S(y) \in B_{29}) \to \forall n \in B_{29}\,(S(n) \in B_{29})$
\item $\varphi_{22} :\equiv \forall B_{29}\,(\neg \forall n\,(n \in B_{29} \leftrightarrow n \in \omega \land (n \in \mathbb{N} \to 0 \leq {{r}^{n}})))$
\item $\varphi_{23} :\equiv \forall A_{30}\,\exists B_{29}\,\forall n\,(n \in B_{29} \leftrightarrow n \in A_{30} \land (n \in \mathbb{N} \to 0 \leq {{r}^{n}}))$
\item $\varphi_{24} :\equiv \forall w\,(r = w \to 0 = r \to 0 = w) \to r = {{r}^{1_{\omega}}} \to 0 = r \to 0 = {{r}^{1_{\omega}}}$
\item $\varphi_{25} :\equiv S(n) \in B_{29} \leftrightarrow S(n) \in \omega \land (S(n) \in \mathbb{N} \to 0 \leq {{r}^{S(n)}})$
\item $\varphi_{26} :\equiv \forall z\,\forall w\,(z = w \to S(n) = S(z) \to S(n) = S(w)) \to \forall w\,(n = w \to S(n) = S(n) \to S(n) = S(w))$
\item $\varphi_{27} :\equiv \forall n\,(n \in B_{29} \leftrightarrow n \in \omega \land (n \in \mathbb{N} \to 0 \leq {{r}^{n}})) \to \mathrm{Ind}(B_{29})$
\item $\varphi_{28} :\equiv 0 \leq r \to \exists B_{29}\,\forall n\,(n \in B_{29} \leftrightarrow n \in \omega \land (n \in \mathbb{N} \to 0 \leq {{r}^{n}}))$
\item $\varphi_{29} :\equiv \forall B_{29}\,(r \in \mathbb{R} \land 0 \leq r \land \neg \forall n \in \omega\,(n \in \mathbb{N} \to 0 \leq {{r}^{n}}))$
\item $\varphi_{30} :\equiv \forall z_{31}\,\forall A_{30}\,\exists B_{29}\,\forall n\,(n \in B_{29} \leftrightarrow n \in A_{30} \land (n \in \mathbb{N} \to 0 \leq {{z_{31}}^{n}}))$
\item $\varphi_{31} :\equiv \forall n\,(n \in \omega \leftrightarrow \forall x\,(\mathrm{Ind}(x) \to n \in x)) \to (n \in \omega \leftrightarrow \forall x\,(\mathrm{Ind}(x) \to n \in x))$
\item $\varphi_{32} :\equiv \forall z\,\forall w\,(z = w \to 0 = z \to 0 = w) \to \forall w\,({{r}^{n}} \cdot r = w \to 0 = {{r}^{n}} \cdot r \to 0 = w)$
\item $\varphi_{33} :\equiv \forall y\,\forall z\,(S(n) = y \to y = z \to S(n) = z) \to \forall z\,(S(n) = S(\emptyset) \to S(\emptyset) = z \to S(n) = z)$
\item $\varphi_{34} :\equiv \forall w\,(n = w \to S(n) = S(n) \to S(n) = S(w)) \to n = \emptyset \to S(n) = S(n) \to S(n) = S(\emptyset)$
\item $\varphi_{35} :\equiv r \in \mathbb{R} \land 0 \leq r \land \forall n\,(n \in B_{29} \leftrightarrow n \in \omega \land (n \in \mathbb{N} \to 0 \leq {{r}^{n}}))$
\item $\varphi_{36} :\equiv \forall n\,(n \in B_{29} \leftrightarrow n \in \omega \land (n \in \mathbb{N} \to 0 \leq {{r}^{n}})) \to \emptyset \in B_{29}$
\item $\varphi_{37} :\equiv \forall w\,(r = w \to 0 < r \to 0 < w) \to r = {{r}^{1_{\omega}}} \to 0 < r \to 0 < {{r}^{1_{\omega}}}$
\item $\varphi_{38} :\equiv ({{r}^{1_{\omega}}} = r \to r = {{r}^{1_{\omega}}}) \to r \in \mathbb{R} \to 0 \leq r \to r = {{r}^{1_{\omega}}}$
\item $\varphi_{39} :\equiv \forall n\,(n \in B_{29} \leftrightarrow n \in \omega \land (n \in \mathbb{N} \to 0 \leq {{r}^{n}})) \to n \in B_{29} \to n \in B_{29}$
\item $\varphi_{40} :\equiv \forall n\,(n \in \mathbb{N} \leftrightarrow n \in \omega \land n \neq \emptyset) \to (n \in \mathbb{N} \leftrightarrow n \in \omega \land n \neq \emptyset)$
\item $\varphi_{41} :\equiv \forall y\,({{r}^{1_{\omega}}} = y \to y = {{r}^{1_{\omega}}}) \to {{r}^{1_{\omega}}} = r \to r = {{r}^{1_{\omega}}}$
\item $\varphi_{42} :\equiv \forall n\,(n \in B_{29} \leftrightarrow n \in \omega \land (n \in \mathbb{N} \to 0 \leq {{r}^{n}})) \to n \in B_{29} \to n \in \omega$
\item $\varphi_{43} :\equiv \mathrm{Seq}({(r)^{\bullet}}) \land {{r}^{1_{\omega}}} = r \land \forall n \in \mathbb{N}\,({{r}^{S(n)}} = {{r}^{n}} \cdot r)$
\item $\varphi_{44} :\equiv (n \in \mathbb{N} \leftrightarrow n \in \omega \land n \neq \emptyset) \to r \in \mathbb{R} \to 0 \leq r \to n \in \mathbb{N} \to n \in \omega$
\item $\varphi_{45} :\equiv (r = {{r}^{1_{\omega}}} \to 0 = r \to 0 = {{r}^{1_{\omega}}}) \to r \in \mathbb{R} \to 0 \leq r \to 0 \leq {{r}^{1_{\omega}}}$
\item $\varphi_{46} :\equiv \forall n\,(n \in B_{29} \leftrightarrow n \in \omega \land (n \in \mathbb{N} \to 0 \leq {{r}^{n}})) \to n \in B_{29} \to S(n) \in B_{29}$
\item $\varphi_{47} :\equiv \forall n\,(n \in B_{29} \leftrightarrow n \in \omega \land (n \in \mathbb{N} \to 0 \leq {{r}^{n}})) \to n \in \omega \to n \in \omega$
\item $\varphi_{48} :\equiv \forall n\,(n \in B_{29} \leftrightarrow n \in \omega \land (n \in \mathbb{N} \to 0 \leq {{r}^{n}})) \to n \in B_{29} \to S(n) \in \omega$
\item $\varphi_{49} :\equiv \forall n\,(n \in B_{29} \leftrightarrow n \in \omega \land (n \in \mathbb{N} \to 0 \leq {{r}^{n}})) \to \forall n \in B_{29}\,(S(n) \in B_{29})$
\item $\varphi_{50} :\equiv \forall z\,\forall w\,(z = w \to 0 < z \to 0 < w) \to \forall w\,({{r}^{n}} \cdot r = w \to 0 < {{r}^{n}} \cdot r \to 0 < w)$
\item $\varphi_{51} :\equiv n \in B_{29} \to S(n) \in \mathbb{N} \to n \neq \emptyset \to n \in \mathbb{N} \to {{r}^{S(n)}} = {{r}^{n}} \cdot r$
\item $\varphi_{52} :\equiv (r \in \mathbb{R} \to 0 \leq r \to n \in \mathbb{N} \to n \in \mathbb{N}) \to r \in \mathbb{R} \to 0 \leq r \to n \in \mathbb{N} \to 0 \leq {{r}^{n}}$
\item $\varphi_{53} :\equiv (r \in \mathbb{R} \land 0 \leq r \to \forall n\,(r \in \mathbb{R} \land 0 \leq r)) \to r \in \mathbb{R} \to 0 \leq r \to \forall n \in \mathbb{N}\,(0 \leq {{r}^{n}})$
\item $\varphi_{54} :\equiv n \in B_{29} \to S(n) \in \mathbb{N} \to n \neq \emptyset \to \forall n \in \mathbb{N}\,({{r}^{S(n)}} = {{r}^{n}} \cdot r)$
\item $\varphi_{55} :\equiv \forall n\,(n \in B_{29} \leftrightarrow n \in \omega \land (n \in \mathbb{N} \to 0 \leq {{r}^{n}})) \to n \in \omega \to \forall x\,(\mathrm{Ind}(x) \to n \in x)$
\item $\varphi_{56} :\equiv \forall z\,\forall w\,(z = w \to 0 = {{r}^{z}} \to 0 = {{r}^{w}}) \to \forall w\,(1_{\omega} = w \to 0 = {{r}^{1_{\omega}}} \to 0 = {{r}^{w}})$
\item $\varphi_{57} :\equiv \forall n\,(n \in B_{29} \leftrightarrow n \in \omega \land (n \in \mathbb{N} \to 0 \leq {{r}^{n}})) \to n \in \omega \to \mathrm{Ind}(B_{29}) \to n \in B_{29}$
\item $\varphi_{58} :\equiv (r \in \mathbb{R} \to 0 \leq r \to n \in \mathbb{N} \to 0 \leq {{r}^{n}}) \to r \in \mathbb{R} \land 0 \leq r \to n \in \mathbb{N} \to 0 \leq {{r}^{n}}$
\item $\varphi_{59} :\equiv \forall a\,\forall n\,(\mathrm{Seq}(a) \to n \in \mathbb{N} \to a_{n} \in \mathbb{R}) \to \forall n\,(\mathrm{Seq}({(r)^{\bullet}}) \to n \in \mathbb{N} \to {{r}^{n}} \in \mathbb{R})$
\item $\varphi_{60} :\equiv \forall n \in \mathbb{N}\,({{r}^{S(n)}} = {{r}^{n}} \cdot r) \to n \in \mathbb{N} \to {{r}^{S(n)}} = {{r}^{n}} \cdot r$
\item $\varphi_{61} :\equiv \forall n\,(n \in B_{29} \leftrightarrow n \in \omega \land (n \in \mathbb{N} \to 0 \leq {{r}^{n}})) \to n \in B_{29} \to n \in \mathbb{N} \to 0 \leq {{r}^{n}}$
\item $\varphi_{62} :\equiv \forall n\,(\mathrm{Seq}({(r)^{\bullet}}) \to n \in \mathbb{N} \to {{r}^{n}} \in \mathbb{R}) \to \mathrm{Seq}({(r)^{\bullet}}) \to n \in \mathbb{N} \to {{r}^{n}} \in \mathbb{R}$
\item $\varphi_{63} :\equiv \forall w\,(1_{\omega} = w \to 0 = {{r}^{1_{\omega}}} \to 0 = {{r}^{w}}) \to 1_{\omega} = S(n) \to 0 = {{r}^{1_{\omega}}} \to 0 = {{r}^{S(n)}}$
\item $\varphi_{64} :\equiv \forall n\,(n \in B_{29} \leftrightarrow n \in \omega \land (n \in \mathbb{N} \to 0 \leq {{r}^{n}})) \to n \in B_{29} \to n \in \omega \to S(n) \in \omega$
\item $\varphi_{65} :\equiv \forall n\,(n \in B_{29} \leftrightarrow n \in \omega \land (n \in \mathbb{N} \to 0 \leq {{r}^{n}})) \to n \in \omega \to n \in \mathbb{N} \to 0 \leq {{r}^{n}}$
\item $\varphi_{66} :\equiv \forall n\,(n \in \mathbb{N} \leftrightarrow n \in \omega \land n \neq \emptyset) \to (\emptyset \in \mathbb{N} \leftrightarrow \emptyset \in \omega \land \emptyset \neq \emptyset)$
\item $\varphi_{67} :\equiv \forall n\,(n \in B_{29} \leftrightarrow n \in \omega \land (n \in \mathbb{N} \to 0 \leq {{r}^{n}})) \to n \in B_{29} \to \forall y \in \omega\,(S(y) \in \omega)$
\item $\varphi_{68} :\equiv \forall n\,(n \in B_{29} \leftrightarrow n \in \omega \land (n \in \mathbb{N} \to 0 \leq {{r}^{n}})) \to \forall n \in \omega\,(n \in \mathbb{N} \to 0 \leq {{r}^{n}})$
\item $\varphi_{69} :\equiv \forall n\,(r \in \mathbb{R} \land 0 \leq r \to n \in \mathbb{N} \to 0 \leq {{r}^{n}}) \to \forall n\,(r \in \mathbb{R} \land 0 \leq r) \to \forall n \in \mathbb{N}\,(0 \leq {{r}^{n}})$
\item $\varphi_{70} :\equiv \forall n\,(n \in B_{29} \leftrightarrow n \in \omega \land (n \in \mathbb{N} \to 0 \leq {{r}^{n}})) \to n \in B_{29} \to S(n) \in \mathbb{N} \to 0 \leq {{r}^{S(n)}}$
\item $\varphi_{71} :\equiv \forall z\,\forall w\,(z = w \to 0 < {{r}^{z}} \to 0 < {{r}^{w}}) \to \forall w\,(1_{\omega} = w \to 0 < {{r}^{1_{\omega}}} \to 0 < {{r}^{w}})$
\item $\varphi_{72} :\equiv \forall n\,(n \in B_{29} \leftrightarrow n \in \omega \land (n \in \mathbb{N} \to 0 \leq {{r}^{n}})) \to n \in B_{29} \to S(n) \in \mathbb{N} \to n = \emptyset \to S(n) = 1_{\omega}$
\item $\varphi_{73} :\equiv \forall n\,(n \in B_{29} \leftrightarrow n \in \omega \land (n \in \mathbb{N} \to 0 \leq {{r}^{n}})) \to n \in B_{29} \to S(n) \in \mathbb{N} \to n = \emptyset \to 1_{\omega} = S(n)$
\item $\varphi_{74} :\equiv \forall n\,(n \in B_{29} \leftrightarrow n \in \omega \land (n \in \mathbb{N} \to 0 \leq {{r}^{n}})) \to n \in B_{29} \to S(n) \in \mathbb{N} \to n \neq \emptyset \to n \in \mathbb{N}$
\item $\varphi_{75} :\equiv \forall n\,(n \in B_{29} \leftrightarrow n \in \omega \land (n \in \mathbb{N} \to 0 \leq {{r}^{n}})) \to n \in B_{29} \to S(n) \in \mathbb{N} \to n = \emptyset \to n = \emptyset$
\item $\varphi_{76} :\equiv \forall n\,(n \in B_{29} \leftrightarrow n \in \omega \land (n \in \mathbb{N} \to 0 \leq {{r}^{n}})) \to n \in B_{29} \to S(n) \in \mathbb{N} \to n \neq \emptyset \to r \in \mathbb{R}$
\item $\varphi_{77} :\equiv \forall n\,(n \in B_{29} \leftrightarrow n \in \omega \land (n \in \mathbb{N} \to 0 \leq {{r}^{n}})) \to (n \in B_{29} \leftrightarrow n \in \omega \land (n \in \mathbb{N} \to 0 \leq {{r}^{n}}))$
\item $\varphi_{78} :\equiv \forall w\,(1_{\omega} = w \to 0 < {{r}^{1_{\omega}}} \to 0 < {{r}^{w}}) \to 1_{\omega} = S(n) \to 0 < {{r}^{1_{\omega}}} \to 0 < {{r}^{S(n)}}$
\item $\varphi_{79} :\equiv \forall n\,(n \in B_{29} \leftrightarrow n \in \omega \land (n \in \mathbb{N} \to 0 \leq {{r}^{n}})) \to n \in B_{29} \to S(n) \in \mathbb{N} \to n \neq \emptyset \to n \neq \emptyset$
\item $\varphi_{80} :\equiv \forall n\,(n \in B_{29} \leftrightarrow n \in \omega \land (n \in \mathbb{N} \to 0 \leq {{r}^{n}})) \to n \in B_{29} \to n \in \omega \land (n \in \mathbb{N} \to 0 \leq {{r}^{n}})$
\item $\varphi_{81} :\equiv \forall n\,(n \in B_{29} \leftrightarrow n \in \omega \land (n \in \mathbb{N} \to 0 \leq {{r}^{n}})) \to \emptyset \in \omega \land \forall y \in \omega\,(S(y) \in \omega)$
\item $\varphi_{82} :\equiv \forall y\,({{r}^{S(n)}} = y \to y = {{r}^{S(n)}}) \to {{r}^{S(n)}} = {{r}^{n}} \cdot r \to {{r}^{n}} \cdot r = {{r}^{S(n)}}$
\item $\varphi_{83} :\equiv \forall n\,(n \in B_{29} \leftrightarrow n \in \omega \land (n \in \mathbb{N} \to 0 \leq {{r}^{n}})) \to \forall n\,(n \in B_{29} \leftrightarrow n \in \omega \land (n \in \mathbb{N} \to 0 \leq {{r}^{n}}))$
\item $\varphi_{84} :\equiv \forall n\,(n \in B_{29} \leftrightarrow n \in \omega \land (n \in \mathbb{N} \to 0 \leq {{r}^{n}})) \to n \in B_{29} \to S(n) \in \mathbb{N} \to n \neq \emptyset \to 0 \leq {{r}^{n}}$
\item $\varphi_{85} :\equiv \forall n\,(n \in B_{29} \leftrightarrow n \in \omega \land (n \in \mathbb{N} \to 0 \leq {{r}^{n}})) \to n \in B_{29} \to S(n) \in \mathbb{N} \to n = \emptyset \to S(n) = S(\emptyset)$
\item $\varphi_{86} :\equiv 0 \leq r \to \forall n\,(n \in B_{29} \leftrightarrow n \in \omega \land (n \in \mathbb{N} \to 0 \leq {{r}^{n}})) \to \varphi_{6}$
\item $\varphi_{87} :\equiv \forall n\,(n \in B_{29} \leftrightarrow n \in \omega \land (n \in \mathbb{N} \to 0 \leq {{r}^{n}})) \to n \in B_{29} \to S(n) \in \mathbb{N} \to n = \emptyset \to 0 \leq {{r}^{S(n)}}$
\item $\varphi_{88} :\equiv \forall n\,(n \in B_{29} \leftrightarrow n \in \omega \land (n \in \mathbb{N} \to 0 \leq {{r}^{n}})) \to n \in B_{29} \to S(n) \in \mathbb{N} \to n \neq \emptyset \to 0 \leq {{r}^{S(n)}}$
\item $\varphi_{89} :\equiv \forall n\,(n \in B_{29} \leftrightarrow n \in \omega \land (n \in \mathbb{N} \to 0 \leq {{r}^{n}})) \to n \in B_{29} \to S(n) \in \mathbb{N} \to n \neq \emptyset \to 0 \leq {{r}^{n}} \cdot r$
\item $\varphi_{90} :\equiv \forall x\,(\mathrm{Ind}(x) \leftrightarrow \emptyset \in x \land \forall y \in x\,(S(y) \in x)) \to (\mathrm{Ind}(\omega) \leftrightarrow \emptyset \in \omega \land \forall y \in \omega\,(S(y) \in \omega))$
\item $\varphi_{91} :\equiv \forall n\,(n \in B_{29} \leftrightarrow n \in \omega \land (n \in \mathbb{N} \to 0 \leq {{r}^{n}})) \to n \in B_{29} \to S(n) \in \mathbb{N} \to n \neq \emptyset \to {{r}^{n}} \in \mathbb{R}$
\item $\varphi_{92} :\equiv r \in \mathbb{R} \land 0 \leq r \land \neg \forall n \in \omega\,(n \in \mathbb{N} \to 0 \leq {{r}^{n}}) \to \neg \forall n\,(n \in B_{29} \leftrightarrow n \in \omega \land (n \in \mathbb{N} \to 0 \leq {{r}^{n}}))$
\item $\varphi_{93} :\equiv \varphi_{23} \to \exists B_{29}\,\forall n\,(n \in B_{29} \leftrightarrow n \in \omega \land (n \in \mathbb{N} \to 0 \leq {{r}^{n}}))$
\item $\varphi_{94} :\equiv \forall w\,({{r}^{n}} \cdot r = w \to 0 = {{r}^{n}} \cdot r \to 0 = w) \to {{r}^{n}} \cdot r = {{r}^{S(n)}} \to 0 = {{r}^{n}} \cdot r \to 0 = {{r}^{S(n)}}$
\item $\varphi_{95} :\equiv 0 \leq r \to \forall n\,(n \in B_{29} \leftrightarrow n \in \omega \land (n \in \mathbb{N} \to 0 \leq {{r}^{n}})) \to \varphi_{25}$
\item $\varphi_{96} :\equiv \forall n\,(n \in B_{29} \leftrightarrow n \in \omega \land (n \in \mathbb{N} \to 0 \leq {{r}^{n}})) \to n \in B_{29} \to S(n) \in \mathbb{N} \to n \neq \emptyset \to {{r}^{S(n)}} = {{r}^{n}} \cdot r$
\item $\varphi_{97} :\equiv \forall n\,(n \in B_{29} \leftrightarrow n \in \omega \land (n \in \mathbb{N} \to 0 \leq {{r}^{n}})) \to n \in B_{29} \to S(n) \in \mathbb{N} \to n \neq \emptyset \to {{r}^{n}} \cdot r = {{r}^{S(n)}}$
\item $\varphi_{98} :\equiv \forall x\,\forall y\,(x \in \mathbb{R} \to y \in \mathbb{R} \to 0 \leq x \to 0 \leq y \to 0 \leq x \cdot y) \to \forall y\,({{r}^{n}} \in \mathbb{R} \to y \in \mathbb{R} \to 0 \leq {{r}^{n}} \to 0 \leq y \to 0 \leq {{r}^{n}} \cdot y)$
\item $\varphi_{99} :\equiv \forall n\,(n \in B_{29} \leftrightarrow n \in \omega \land (n \in \mathbb{N} \to 0 \leq {{r}^{n}})) \to (\emptyset \in B_{29} \leftrightarrow \emptyset \in \omega \land (\emptyset \in \mathbb{N} \to 0 \leq {{r}^{\emptyset}}))$
\item $\varphi_{100} :\equiv \forall w\,({{r}^{n}} \cdot r = w \to 0 < {{r}^{n}} \cdot r \to 0 < w) \to {{r}^{n}} \cdot r = {{r}^{S(n)}} \to 0 < {{r}^{n}} \cdot r \to 0 < {{r}^{S(n)}}$
\item $\varphi_{101} :\equiv (\emptyset \in \mathbb{N} \leftrightarrow \emptyset \in \omega \land \emptyset \neq \emptyset) \to \emptyset = \emptyset \to r \in \mathbb{R} \to 0 \leq r \to \emptyset \in \mathbb{N} \to 0 \leq {{r}^{\emptyset}}$
\item $\varphi_{102} :\equiv \exists B_{29}\,\forall n\,(n \in B_{29} \leftrightarrow n \in \omega \land (n \in \mathbb{N} \to 0 \leq {{r}^{n}})) \to r \in \mathbb{R} \to \varphi_{28}$
\item $\varphi_{103} :\equiv (n \in \omega \leftrightarrow \forall x\,(\mathrm{Ind}(x) \to n \in x)) \to r \in \mathbb{R} \to 0 \leq r \to \varphi_{55}$
\item $\varphi_{104} :\equiv (\forall n \in \omega\,(n \in \mathbb{N} \to 0 \leq {{r}^{n}}) \to n \in \omega \to n \in \mathbb{N} \to 0 \leq {{r}^{n}}) \to r \in \mathbb{R} \to 0 \leq r \to n \in \omega \to n \in \mathbb{N} \to 0 \leq {{r}^{n}}$
\item $\varphi_{105} :\equiv (S(n) = 1_{\omega} \to 1_{\omega} = S(n)) \to r \in \mathbb{R} \to 0 \leq r \to \varphi_{73}$
\item $\varphi_{106} :\equiv 0 \leq r \to \forall n\,(n \in B_{29} \leftrightarrow n \in \omega \land (n \in \mathbb{N} \to 0 \leq {{r}^{n}})) \to \varphi_{51}$
\item $\varphi_{107} :\equiv 0 \leq r \to \forall n\,(n \in B_{29} \leftrightarrow n \in \omega \land (n \in \mathbb{N} \to 0 \leq {{r}^{n}})) \to \varphi_{54}$
\item $\varphi_{108} :\equiv (r \in \mathbb{R} \to 0 \leq r \to \emptyset \in \mathbb{N} \to 0 \leq {{r}^{\emptyset}}) \to r \in \mathbb{R} \to 0 \leq r \to \varphi_{36}$
\item $\varphi_{109} :\equiv (\forall x\,(\mathrm{Ind}(x) \to n \in x) \to \mathrm{Ind}(B_{29}) \to n \in B_{29}) \to r \in \mathbb{R} \to 0 \leq r \to \varphi_{57}$
\item $\varphi_{110} :\equiv (r \in \mathbb{R} \to 0 \leq r \to 0 \leq r) \to (r = {{r}^{1_{\omega}}} \to 0 < r \to 0 < {{r}^{1_{\omega}}}) \to \varphi_{45}$
\item $\varphi_{111} :\equiv \forall y\,({{r}^{n}} \in \mathbb{R} \to y \in \mathbb{R} \to 0 \leq {{r}^{n}} \to 0 \leq y \to 0 \leq {{r}^{n}} \cdot y) \to {{r}^{n}} \in \mathbb{R} \to r \in \mathbb{R} \to 0 \leq {{r}^{n}} \to 0 \leq r \to 0 \leq {{r}^{n}} \cdot r$
\item $\varphi_{112} :\equiv (n = \emptyset \to S(n) = S(n) \to S(n) = S(\emptyset)) \to r \in \mathbb{R} \to 0 \leq r \to \varphi_{85}$
\item $\varphi_{113} :\equiv \forall n\,(n \in B_{29} \leftrightarrow n \in \omega \land (n \in \mathbb{N} \to 0 \leq {{r}^{n}})) \to n \in B_{29} \to S(n) \in \mathbb{N} \to n \neq \emptyset \to \varphi_{43}$
\item $\varphi_{114} :\equiv (\forall y \in \omega\,(S(y) \in \omega) \to n \in \omega \to S(n) \in \omega) \to r \in \mathbb{R} \to 0 \leq r \to \varphi_{64}$
\item $\varphi_{115} :\equiv (1_{\omega} = S(n) \to 0 = {{r}^{1_{\omega}}} \to 0 = {{r}^{S(n)}}) \to r \in \mathbb{R} \to 0 \leq r \to \varphi_{87}$
\item $\varphi_{116} :\equiv (r \in \mathbb{R} \to 0 \leq r \to n \in \omega \to n \in \mathbb{N} \to 0 \leq {{r}^{n}}) \to (r \in \mathbb{R} \to 0 \leq r \to n \in \mathbb{N} \to n \in \omega) \to \varphi_{52}$
\item $\varphi_{117} :\equiv (\mathrm{Ind}(\omega) \leftrightarrow \emptyset \in \omega \land \forall y \in \omega\,(S(y) \in \omega)) \to r \in \mathbb{R} \to 0 \leq r \to \varphi_{81}$
\item $\varphi_{118} :\equiv (r \in \mathbb{R} \land 0 \leq r \land \neg \forall n \in \omega\,(n \in \mathbb{N} \to 0 \leq {{r}^{n}}) \to \varphi_{29}) \to r \in \mathbb{R} \to 0 \leq r \to \forall n \in \omega\,(n \in \mathbb{N} \to 0 \leq {{r}^{n}})$
\item $\varphi_{119} :\equiv ({{r}^{n}} \cdot r = {{r}^{S(n)}} \to 0 = {{r}^{n}} \cdot r \to 0 = {{r}^{S(n)}}) \to r \in \mathbb{R} \to 0 \leq r \to \varphi_{88}$
\item $\varphi_{120} :\equiv \forall n\,(\varphi_{35} \to n \in B_{29} \to S(n) \in B_{29}) \to \forall n\,(\varphi_{35}) \to \forall n \in B_{29}\,(S(n) \in B_{29})$
\item $\varphi_{121} :\equiv ({{r}^{S(n)}} = {{r}^{n}} \cdot r \to {{r}^{n}} \cdot r = {{r}^{S(n)}}) \to r \in \mathbb{R} \to 0 \leq r \to \varphi_{97}$
\item $\varphi_{122} :\equiv (r \in \mathbb{R} \to \varphi_{43}) \to (r \in \mathbb{R} \to 0 \leq r \to r \in \mathbb{R}) \to r \in \mathbb{R} \to 0 \leq r \to \varphi_{43}$
\item $\varphi_{123} :\equiv \forall n\,(\varphi_{35} \to n \in \omega \to n \in \mathbb{N} \to 0 \leq {{r}^{n}}) \to \forall n\,(\varphi_{35}) \to \forall n \in \omega\,(n \in \mathbb{N} \to 0 \leq {{r}^{n}})$
\item $\varphi_{124} :\equiv (\forall n\,(n \in B_{29} \leftrightarrow n \in \omega \land (n \in \mathbb{N} \to 0 \leq {{r}^{n}})) \to \varphi_{25}) \to r \in \mathbb{R} \to \varphi_{95}$
\item $\varphi_{125} :\equiv (r \in \mathbb{R} \to 0 \leq r \to 0 \leq {{r}^{1_{\omega}}}) \to (1_{\omega} = S(n) \to 0 < {{r}^{1_{\omega}}} \to 0 < {{r}^{S(n)}}) \to \varphi_{115}$
\item $\varphi_{126} :\equiv (r \in \mathbb{R} \to 0 \leq r \to \varphi_{84}) \to (r \in \mathbb{R} \to 0 \leq r \to 0 \leq r) \to r \in \mathbb{R} \to 0 \leq r \to \varphi_{89}$
\item $\varphi_{127} :\equiv (r \in \mathbb{R} \to 0 \leq r \to \varphi_{64}) \to (r \in \mathbb{R} \to 0 \leq r \to \varphi_{42}) \to r \in \mathbb{R} \to 0 \leq r \to \varphi_{48}$
\item $\varphi_{128} :\equiv (r \in \mathbb{R} \to 0 \leq r \to \varphi_{85}) \to S(\emptyset) = 1_{\omega} \to \varphi_{5} \to r \in \mathbb{R} \to 0 \leq r \to \varphi_{72}$
\item $\varphi_{129} :\equiv (r \in \mathbb{R} \to 0 \leq r \to \varphi_{48}) \to (r \in \mathbb{R} \to 0 \leq r \to \varphi_{70}) \to r \in \mathbb{R} \to 0 \leq r \to \varphi_{46}$
\item $\varphi_{130} :\equiv (r \in \mathbb{R} \to 0 \leq r \to \varphi_{27}) \to (r \in \mathbb{R} \to 0 \leq r \to \varphi_{77}) \to r \in \mathbb{R} \to 0 \leq r \to \varphi_{65}$
\item $\varphi_{131} :\equiv (r \in \mathbb{R} \to 0 \leq r \to \varphi_{77}) \to (r \in \mathbb{R} \to 0 \leq r \to \varphi_{39}) \to r \in \mathbb{R} \to 0 \leq r \to \varphi_{80}$
\item $\varphi_{132} :\equiv (r \in \mathbb{R} \to 0 \leq r \to \varphi_{42}) \to (r \in \mathbb{R} \to 0 \leq r \to \varphi_{79}) \to r \in \mathbb{R} \to 0 \leq r \to \varphi_{74}$
\item $\varphi_{133} :\equiv (r \in \mathbb{R} \to 0 \leq r \to \varphi_{61}) \to (r \in \mathbb{R} \to 0 \leq r \to \varphi_{74}) \to r \in \mathbb{R} \to 0 \leq r \to \varphi_{84}$
\item $\varphi_{134} :\equiv (r \in \mathbb{R} \to 0 \leq r \to \varphi_{87}) \to (r \in \mathbb{R} \to 0 \leq r \to \varphi_{88}) \to r \in \mathbb{R} \to 0 \leq r \to \varphi_{70}$
\item $\varphi_{135} :\equiv (r \in \mathbb{R} \to \varphi_{86}) \to (r \in \mathbb{R} \to 0 \leq r \to \varphi_{74}) \to r \in \mathbb{R} \to 0 \leq r \to \varphi_{91}$
\item $\varphi_{136} :\equiv (r \in \mathbb{R} \to 0 \leq r \to \varphi_{36}) \to (r \in \mathbb{R} \to 0 \leq r \to \varphi_{49}) \to \varphi_{14} \to \varphi_{13} \to r \in \mathbb{R} \to 0 \leq r \to \varphi_{27}$
\item $\varphi_{137} :\equiv (r \in \mathbb{R} \to \varphi_{43}) \to (r \in \mathbb{R} \to 0 \leq r \to \varphi_{76}) \to r \in \mathbb{R} \to 0 \leq r \to \varphi_{113}$
\item $\varphi_{138} :\equiv (r \in \mathbb{R} \to 0 \leq r \to \varphi_{89}) \to ({{r}^{n}} \cdot r = {{r}^{S(n)}} \to 0 < {{r}^{n}} \cdot r \to 0 < {{r}^{S(n)}}) \to \varphi_{119}$
\item $\varphi_{139} :\equiv (r \in \mathbb{R} \to \varphi_{106}) \to (r \in \mathbb{R} \to 0 \leq r \to \varphi_{74}) \to r \in \mathbb{R} \to 0 \leq r \to \varphi_{96}$
\item $\varphi_{140} :\equiv \forall n\,(\varphi_{35} \to n \in B_{29} \to S(n) \in B_{29}) \to \varphi_{120} \to (\varphi_{35} \to \forall n\,(\varphi_{35})) \to r \in \mathbb{R} \to 0 \leq r \to \varphi_{49}$
\item $\varphi_{141} :\equiv \forall n\,(\varphi_{35} \to n \in \omega \to n \in \mathbb{N} \to 0 \leq {{r}^{n}}) \to \varphi_{123} \to (\varphi_{35} \to \forall n\,(\varphi_{35})) \to r \in \mathbb{R} \to 0 \leq r \to \varphi_{68}$
\item $\varphi_{142} :\equiv ({{r}^{n}} \in \mathbb{R} \to r \in \mathbb{R} \to 0 \leq {{r}^{n}} \to 0 \leq r \to 0 \leq {{r}^{n}} \cdot r) \to (r \in \mathbb{R} \to 0 \leq r \to \varphi_{91}) \to (r \in \mathbb{R} \to 0 \leq r \to \varphi_{76}) \to \varphi_{126}$
\item $\varphi_{143} :\equiv (r \in \mathbb{R} \to \varphi_{28}) \to \forall B_{29}\,(\varphi_{92}) \to (\forall B_{29}\,(\varphi_{92}) \to \varphi_{29} \to \varphi_{22}) \to \varphi_{118}$
\end{abbrevlist}
\end{fullproof}

% 自動生成 (prover/build.js): PowPos  ノード数 328
\begin{fullproof}{[P3]}{$\forall r \in \mathbb{R}\,(0 < r \to \forall n \in \mathbb{N}\,(0 < {{r}^{n}}))$}{328}\phantomsection\label{tree:PowPos}
\noindent$\Pi_{1}$:
\begin{ptree}
\AXC{$r \in \mathbb{R} \to 0 < r \to n \in \mathbb{N} \to n \in \mathbb{N}$\lbl{A1}}
\AXC{$(r \in \mathbb{R} \to 0 < r \to n \in \mathbb{N} \to n \in \mathbb{N}) \to \varphi_{41}$\lbl{A1}}
\BIC{$\varphi_{41}$}
\end{ptree}
\noindent$\Pi_{2}$:
\begin{ptree}
\AXC{$\varphi_{76}$\lbl{A2}}
\AXC{$r \in \mathbb{R} \to 0 < r \to \varphi_{79}$\lbl{A1}}
\AXC{$(r \in \mathbb{R} \to 0 < r \to \varphi_{79}) \to \varphi_{76} \to r \in \mathbb{R} \to 0 < r \to \varphi_{76}$\lbl{A1}}
\BIC{$\varphi_{76} \to r \in \mathbb{R} \to 0 < r \to \varphi_{76}$}
\BIC{$r \in \mathbb{R} \to 0 < r \to \varphi_{76}$}
\end{ptree}
\noindent$\Pi_{3}$:
\begin{ptree}
\AXC{$\varphi_{76}$\lbl{A2}}
\AXC{$r \in \mathbb{R} \to 0 < r \to \varphi_{79}$\lbl{A1}}
\AXC{$(r \in \mathbb{R} \to 0 < r \to \varphi_{79}) \to \varphi_{76} \to r \in \mathbb{R} \to 0 < r \to \varphi_{76}$\lbl{A1}}
\BIC{$\varphi_{76} \to r \in \mathbb{R} \to 0 < r \to \varphi_{76}$}
\BIC{$r \in \mathbb{R} \to 0 < r \to \varphi_{76}$}
\AXC{$\varphi_{126}$\lbl{A1}}
\BIC{$(r \in \mathbb{R} \to 0 < r \to \varphi_{39}) \to r \in \mathbb{R} \to 0 < r \to \varphi_{77}$}
\end{ptree}
\noindent$\Pi_{4}$:
\begin{ptree}
\AXC{$r \in \mathbb{R} \to 0 < r \to \varphi_{39}$\lbl{A1}}
\AXC{$\stackrel{\Pi_{3}}{\vdots}$}\noLine
\UIC{$(r \in \mathbb{R} \to 0 < r \to \varphi_{39}) \to r \in \mathbb{R} \to 0 < r \to \varphi_{77}$}
\BIC{$r \in \mathbb{R} \to 0 < r \to \varphi_{77}$}
\end{ptree}
\noindent$\Pi_{5}$:
\begin{ptree}
\AXC{$\stackrel{\Pi_{4}}{\vdots}$}\noLine
\UIC{$r \in \mathbb{R} \to 0 < r \to \varphi_{77}$}
\AXC{$(r \in \mathbb{R} \to 0 < r \to \varphi_{77}) \to r \in \mathbb{R} \to 0 < r \to \varphi_{40}$\lbl{A1}}
\BIC{$r \in \mathbb{R} \to 0 < r \to \varphi_{40}$}
\end{ptree}
\noindent$\Pi_{6}$:
\begin{ptree}
\AXC{$\stackrel{\Pi_{5}}{\vdots}$}\noLine
\UIC{$r \in \mathbb{R} \to 0 < r \to \varphi_{40}$}
\AXC{$\forall n\,(n \in \mathbb{N} \leftrightarrow n \in \omega \land n \neq \emptyset)$\lbl{Def$\mathbb{N}$}}
\AXC{$\varphi_{35}$\lbl{A2}}
\BIC{$n \in \mathbb{N} \leftrightarrow n \in \omega \land n \neq \emptyset$}
\AXC{$(n \in \mathbb{N} \leftrightarrow n \in \omega \land n \neq \emptyset) \to \varphi_{124}$\lbl{A1}}
\BIC{$\varphi_{124}$}
\BIC{$(r \in \mathbb{R} \to 0 < r \to \varphi_{73}) \to r \in \mathbb{R} \to 0 < r \to \varphi_{69}$}
\end{ptree}
\noindent$\Pi_{7}$:
\begin{ptree}
\AXC{$r \in \mathbb{R} \to 0 < r \to \varphi_{73}$\lbl{A1}}
\AXC{$\stackrel{\Pi_{6}}{\vdots}$}\noLine
\UIC{$(r \in \mathbb{R} \to 0 < r \to \varphi_{73}) \to r \in \mathbb{R} \to 0 < r \to \varphi_{69}$}
\BIC{$r \in \mathbb{R} \to 0 < r \to \varphi_{69}$}
\end{ptree}
\noindent$\Pi_{8}$:
\begin{ptree}
\AXC{$\stackrel{\Pi_{7}}{\vdots}$}\noLine
\UIC{$r \in \mathbb{R} \to 0 < r \to \varphi_{69}$}
\AXC{$\stackrel{\Pi_{4}}{\vdots}$}\noLine
\UIC{$r \in \mathbb{R} \to 0 < r \to \varphi_{77}$}
\AXC{$(r \in \mathbb{R} \to 0 < r \to \varphi_{77}) \to r \in \mathbb{R} \to 0 < r \to \varphi_{59}$\lbl{A1}}
\BIC{$r \in \mathbb{R} \to 0 < r \to \varphi_{59}$}
\AXC{$\varphi_{127}$\lbl{A1}}
\BIC{$(r \in \mathbb{R} \to 0 < r \to \varphi_{69}) \to r \in \mathbb{R} \to 0 < r \to \varphi_{80}$}
\BIC{$r \in \mathbb{R} \to 0 < r \to \varphi_{80}$}
\end{ptree}
\noindent$\Pi_{9}$:
\begin{ptree}
\AXC{$r \in \mathbb{R} \to 0 < r \to \varphi_{72}$\lbl{A1}}
\AXC{$\forall r \in \mathbb{R}\,(\varphi_{38})$\lbl{[S24]}}
\AXC{$\forall r \in \mathbb{R}\,(\varphi_{38}) \to r \in \mathbb{R} \to \varphi_{38}$\lbl{A2}}
\BIC{$r \in \mathbb{R} \to \varphi_{38}$}
\AXC{$\varphi_{131}$\lbl{A1}}
\BIC{$(r \in \mathbb{R} \to 0 < r \to \varphi_{72}) \to r \in \mathbb{R} \to 0 < r \to \varphi_{107}$}
\BIC{$r \in \mathbb{R} \to 0 < r \to \varphi_{107}$}
\end{ptree}
\noindent$\Pi_{10}$:
\begin{ptree}
\AXC{$\forall a\,\forall n\,(\mathrm{Seq}(a) \to n \in \mathbb{N} \to a_{n} \in \mathbb{R})$\lbl{[S1]}}
\AXC{$\varphi_{52}$\lbl{A2}}
\BIC{$\forall n\,(\mathrm{Seq}({(r)^{\bullet}}) \to n \in \mathbb{N} \to {{r}^{n}} \in \mathbb{R})$}
\AXC{$\varphi_{55}$\lbl{A2}}
\BIC{$\mathrm{Seq}({(r)^{\bullet}}) \to n \in \mathbb{N} \to {{r}^{n}} \in \mathbb{R}$}
\end{ptree}
\noindent$\Pi_{11}$:
\begin{ptree}
\AXC{$\stackrel{\Pi_{10}}{\vdots}$}\noLine
\UIC{$\mathrm{Seq}({(r)^{\bullet}}) \to n \in \mathbb{N} \to {{r}^{n}} \in \mathbb{R}$}
\AXC{$(\mathrm{Seq}({(r)^{\bullet}}) \to n \in \mathbb{N} \to {{r}^{n}} \in \mathbb{R}) \to \varphi_{128}$\lbl{A1}}
\BIC{$\varphi_{128}$}
\end{ptree}
\noindent$\Pi_{12}$:
\begin{ptree}
\AXC{$\stackrel{\Pi_{7}}{\vdots}$}\noLine
\UIC{$r \in \mathbb{R} \to 0 < r \to \varphi_{69}$}
\AXC{$\stackrel{\Pi_{9}}{\vdots}$}\noLine
\UIC{$r \in \mathbb{R} \to 0 < r \to \varphi_{107}$}
\AXC{$(r \in \mathbb{R} \to 0 < r \to \varphi_{107}) \to r \in \mathbb{R} \to \varphi_{82}$\lbl{A1}}
\BIC{$r \in \mathbb{R} \to \varphi_{82}$}
\AXC{$\stackrel{\Pi_{11}}{\vdots}$}\noLine
\UIC{$\varphi_{128}$}
\BIC{$(r \in \mathbb{R} \to 0 < r \to \varphi_{69}) \to r \in \mathbb{R} \to 0 < r \to \varphi_{85}$}
\BIC{$r \in \mathbb{R} \to 0 < r \to \varphi_{85}$}
\end{ptree}
\noindent$\Pi_{13}$:
\begin{ptree}
\AXC{$\stackrel{\Pi_{12}}{\vdots}$}\noLine
\UIC{$r \in \mathbb{R} \to 0 < r \to \varphi_{85}$}
\AXC{$\forall x\,\forall y\,(x \in \mathbb{R} \to y \in \mathbb{R} \to 0 < x \to 0 < y \to 0 < x \cdot y)$\lbl{[E11]}}
\AXC{$\varphi_{97}$\lbl{A2}}
\BIC{$\forall y\,({{r}^{n}} \in \mathbb{R} \to y \in \mathbb{R} \to 0 < {{r}^{n}} \to 0 < y \to 0 < {{r}^{n}} \cdot y)$}
\AXC{$\varphi_{105}$\lbl{A2}}
\BIC{${{r}^{n}} \in \mathbb{R} \to r \in \mathbb{R} \to 0 < {{r}^{n}} \to 0 < r \to 0 < {{r}^{n}} \cdot r$}
\AXC{$\varphi_{134}$\lbl{A1}}
\BIC{$(r \in \mathbb{R} \to 0 < r \to \varphi_{85}) \to (r \in \mathbb{R} \to 0 < r \to \varphi_{72}) \to \varphi_{119}$}
\BIC{$(r \in \mathbb{R} \to 0 < r \to \varphi_{72}) \to \varphi_{119}$}
\end{ptree}
\noindent$\Pi_{14}$:
\begin{ptree}
\AXC{$r \in \mathbb{R} \to 0 < r \to 0 < r$\lbl{A1}}
\AXC{$\stackrel{\Pi_{8}}{\vdots}$}\noLine
\UIC{$r \in \mathbb{R} \to 0 < r \to \varphi_{80}$}
\AXC{$r \in \mathbb{R} \to 0 < r \to \varphi_{72}$\lbl{A1}}
\AXC{$\stackrel{\Pi_{13}}{\vdots}$}\noLine
\UIC{$(r \in \mathbb{R} \to 0 < r \to \varphi_{72}) \to \varphi_{119}$}
\BIC{$\varphi_{119}$}
\BIC{$(r \in \mathbb{R} \to 0 < r \to 0 < r) \to r \in \mathbb{R} \to 0 < r \to \varphi_{86}$}
\BIC{$r \in \mathbb{R} \to 0 < r \to \varphi_{86}$}
\end{ptree}
\noindent$\Pi_{15}$:
\begin{ptree}
\AXC{$x = x$\lbl{A5}}
\AXC{$x = y \to x = x \to y = x$\lbl{A6}}
\AXC{$(x = y \to x = x \to y = x) \to x = x \to x = y \to y = x$\lbl{A1}}
\BIC{$x = x \to x = y \to y = x$}
\BIC{$x = y \to y = x$}
\UIC{$\forall y\,(x = y \to y = x)$}
\UIC{$\forall x\,\forall y\,(x = y \to y = x)$}
\end{ptree}
\noindent$\Pi_{16}$:
\begin{ptree}
\AXC{$r \in \mathbb{R} \to 0 < r \to \varphi_{72}$\lbl{A1}}
\AXC{$\forall r \in \mathbb{R}\,(\varphi_{38})$\lbl{[S24]}}
\AXC{$\forall r \in \mathbb{R}\,(\varphi_{38}) \to r \in \mathbb{R} \to \varphi_{38}$\lbl{A2}}
\BIC{$r \in \mathbb{R} \to \varphi_{38}$}
\AXC{$\varphi_{131}$\lbl{A1}}
\BIC{$(r \in \mathbb{R} \to 0 < r \to \varphi_{72}) \to r \in \mathbb{R} \to 0 < r \to \varphi_{107}$}
\BIC{$r \in \mathbb{R} \to 0 < r \to \varphi_{107}$}
\end{ptree}
\noindent$\Pi_{17}$:
\begin{ptree}
\AXC{$\varphi_{53}$\lbl{A2}}
\AXC{$\stackrel{\Pi_{16}}{\vdots}$}\noLine
\UIC{$r \in \mathbb{R} \to 0 < r \to \varphi_{107}$}
\AXC{$(r \in \mathbb{R} \to 0 < r \to \varphi_{107}) \to r \in \mathbb{R} \to \varphi_{100}$\lbl{A1}}
\BIC{$r \in \mathbb{R} \to \varphi_{100}$}
\AXC{$(r \in \mathbb{R} \to \varphi_{100}) \to \varphi_{53} \to r \in \mathbb{R} \to \varphi_{99}$\lbl{A1}}
\BIC{$\varphi_{53} \to r \in \mathbb{R} \to \varphi_{99}$}
\BIC{$r \in \mathbb{R} \to \varphi_{99}$}
\end{ptree}
\noindent$\Pi_{18}$:
\begin{ptree}
\AXC{$\stackrel{\Pi_{7}}{\vdots}$}\noLine
\UIC{$r \in \mathbb{R} \to 0 < r \to \varphi_{69}$}
\AXC{$\stackrel{\Pi_{17}}{\vdots}$}\noLine
\UIC{$r \in \mathbb{R} \to \varphi_{99}$}
\AXC{$\varphi_{132}$\lbl{A1}}
\BIC{$(r \in \mathbb{R} \to 0 < r \to \varphi_{69}) \to r \in \mathbb{R} \to 0 < r \to \varphi_{89}$}
\BIC{$r \in \mathbb{R} \to 0 < r \to \varphi_{89}$}
\AXC{$(r \in \mathbb{R} \to 0 < r \to \varphi_{89}) \to \varphi_{115}$\lbl{A1}}
\BIC{$\varphi_{115}$}
\end{ptree}
\noindent$\Pi_{19}$:
\begin{ptree}
\AXC{$\stackrel{\Pi_{15}}{\vdots}$}\noLine
\UIC{$\forall x\,\forall y\,(x = y \to y = x)$}
\AXC{$\varphi_{15}$\lbl{A2}}
\BIC{$\forall y\,({{r}^{S(n)}} = y \to y = {{r}^{S(n)}})$}
\AXC{$\varphi_{71}$\lbl{A2}}
\BIC{${{r}^{S(n)}} = {{r}^{n}} \cdot r \to {{r}^{n}} \cdot r = {{r}^{S(n)}}$}
\AXC{$\stackrel{\Pi_{18}}{\vdots}$}\noLine
\UIC{$\varphi_{115}$}
\BIC{$r \in \mathbb{R} \to 0 < r \to \varphi_{90}$}
\AXC{$(r \in \mathbb{R} \to 0 < r \to \varphi_{90}) \to (r \in \mathbb{R} \to 0 < r \to \varphi_{86}) \to \varphi_{114}$\lbl{A1}}
\BIC{$(r \in \mathbb{R} \to 0 < r \to \varphi_{86}) \to \varphi_{114}$}
\end{ptree}
\noindent$\Pi_{20}$:
\begin{ptree}
\AXC{$z = w \to 0 < z \to 0 < w$\lbl{A6}}
\UIC{$\forall w\,(z = w \to 0 < z \to 0 < w)$}
\UIC{$\forall z\,\forall w\,(z = w \to 0 < z \to 0 < w)$}
\AXC{$\varphi_{43}$\lbl{A2}}
\BIC{$\forall w\,({{r}^{n}} \cdot r = w \to 0 < {{r}^{n}} \cdot r \to 0 < w)$}
\AXC{$\varphi_{92}$\lbl{A2}}
\BIC{${{r}^{n}} \cdot r = {{r}^{S(n)}} \to 0 < {{r}^{n}} \cdot r \to 0 < {{r}^{S(n)}}$}
\AXC{$\stackrel{\Pi_{14}}{\vdots}$}\noLine
\UIC{$r \in \mathbb{R} \to 0 < r \to \varphi_{86}$}
\AXC{$\stackrel{\Pi_{19}}{\vdots}$}\noLine
\UIC{$(r \in \mathbb{R} \to 0 < r \to \varphi_{86}) \to \varphi_{114}$}
\BIC{$\varphi_{114}$}
\BIC{$r \in \mathbb{R} \to 0 < r \to \varphi_{84}$}
\end{ptree}
\noindent$\Pi_{21}$:
\begin{ptree}
\AXC{$r \in \mathbb{R} \to 0 < r \to r \in \mathbb{R}$\lbl{A1}}
\AXC{$\forall r \in \mathbb{R}\,(\varphi_{38})$\lbl{[S24]}}
\AXC{$\forall r \in \mathbb{R}\,(\varphi_{38}) \to r \in \mathbb{R} \to \varphi_{38}$\lbl{A2}}
\BIC{$r \in \mathbb{R} \to \varphi_{38}$}
\AXC{$\varphi_{116}$\lbl{A1}}
\BIC{$(r \in \mathbb{R} \to 0 < r \to r \in \mathbb{R}) \to r \in \mathbb{R} \to 0 < r \to \varphi_{38}$}
\BIC{$r \in \mathbb{R} \to 0 < r \to \varphi_{38}$}
\end{ptree}
\noindent$\Pi_{22}$:
\begin{ptree}
\AXC{$\stackrel{\Pi_{21}}{\vdots}$}\noLine
\UIC{$r \in \mathbb{R} \to 0 < r \to \varphi_{38}$}
\AXC{$(r \in \mathbb{R} \to 0 < r \to \varphi_{38}) \to r \in \mathbb{R} \to 0 < r \to {{r}^{1_{\omega}}} = r$\lbl{A1}}
\BIC{$r \in \mathbb{R} \to 0 < r \to {{r}^{1_{\omega}}} = r$}
\AXC{$(r \in \mathbb{R} \to 0 < r \to {{r}^{1_{\omega}}} = r) \to \varphi_{34}$\lbl{A1}}
\BIC{$\varphi_{34}$}
\end{ptree}
\noindent$\Pi_{23}$:
\begin{ptree}
\AXC{$r \in \mathbb{R} \to 0 < r \to 0 < r$\lbl{A1}}
\AXC{$\stackrel{\Pi_{15}}{\vdots}$}\noLine
\UIC{$\forall x\,\forall y\,(x = y \to y = x)$}
\AXC{$\varphi_{16}$\lbl{A2}}
\BIC{$\forall y\,({{r}^{1_{\omega}}} = y \to y = {{r}^{1_{\omega}}})$}
\AXC{$\varphi_{36}$\lbl{A2}}
\BIC{${{r}^{1_{\omega}}} = r \to r = {{r}^{1_{\omega}}}$}
\AXC{$\stackrel{\Pi_{22}}{\vdots}$}\noLine
\UIC{$\varphi_{34}$}
\BIC{$r \in \mathbb{R} \to 0 < r \to r = {{r}^{1_{\omega}}}$}
\AXC{$\varphi_{103}$\lbl{A1}}
\BIC{$(r \in \mathbb{R} \to 0 < r \to 0 < r) \to \varphi_{49}$}
\BIC{$\varphi_{49}$}
\end{ptree}
\noindent$\Pi_{24}$:
\begin{ptree}
\AXC{$y = z \to x = y \to x = z$\lbl{A6}}
\AXC{$(y = z \to x = y \to x = z) \to x = y \to y = z \to x = z$\lbl{A1}}
\BIC{$x = y \to y = z \to x = z$}
\UIC{$\forall z\,(x = y \to y = z \to x = z)$}
\UIC{$\forall y\,\forall z\,(x = y \to y = z \to x = z)$}
\UIC{$\forall x\,\forall y\,\forall z\,(x = y \to y = z \to x = z)$}
\AXC{$\varphi_{9}$\lbl{A2}}
\BIC{$\forall y\,\forall z\,(S(n) = y \to y = z \to S(n) = z)$}
\AXC{$\varphi_{29}$\lbl{A2}}
\BIC{$\forall z\,(S(n) = S(\emptyset) \to S(\emptyset) = z \to S(n) = z)$}
\end{ptree}
\noindent$\Pi_{25}$:
\begin{ptree}
\AXC{$\stackrel{\Pi_{24}}{\vdots}$}\noLine
\UIC{$\forall z\,(S(n) = S(\emptyset) \to S(\emptyset) = z \to S(n) = z)$}
\AXC{$\forall z\,(S(n) = S(\emptyset) \to S(\emptyset) = z \to S(n) = z) \to \varphi_{4}$\lbl{A2}}
\BIC{$\varphi_{4}$}
\end{ptree}
\noindent$\Pi_{26}$:
\begin{ptree}
\AXC{$\stackrel{\Pi_{15}}{\vdots}$}\noLine
\UIC{$\forall x\,\forall y\,(x = y \to y = x)$}
\AXC{$\varphi_{2}$\lbl{A2}}
\BIC{$\forall y\,(1_{\omega} = y \to y = 1_{\omega})$}
\AXC{$\varphi_{11}$\lbl{A2}}
\BIC{$1_{\omega} = S(\emptyset) \to S(\emptyset) = 1_{\omega}$}
\AXC{$1_{\omega} = S(\emptyset)$\lbl{Def$1_\omega$}}
\AXC{$1_{\omega} = S(\emptyset) \to \varphi_{10}$\lbl{A1}}
\BIC{$\varphi_{10}$}
\BIC{$S(\emptyset) = 1_{\omega}$}
\end{ptree}
\noindent$\Pi_{27}$:
\begin{ptree}
\AXC{$r \in \mathbb{R} \to 0 < r \to \varphi_{70}$\lbl{A1}}
\AXC{$(r \in \mathbb{R} \to 0 < r \to \varphi_{70}) \to S(n) = S(n) \to \varphi_{106}$\lbl{A1}}
\BIC{$S(n) = S(n) \to \varphi_{106}$}
\end{ptree}
\noindent$\Pi_{28}$:
\begin{ptree}
\AXC{$x = x$\lbl{A5}}
\UIC{$\forall x\,(x = x)$}
\AXC{$\forall x\,(x = x) \to S(n) = S(n)$\lbl{A2}}
\BIC{$S(n) = S(n)$}
\AXC{$\stackrel{\Pi_{27}}{\vdots}$}\noLine
\UIC{$S(n) = S(n) \to \varphi_{106}$}
\BIC{$\varphi_{106}$}
\end{ptree}
\noindent$\Pi_{29}$:
\begin{ptree}
\AXC{$z = w \to S(n) = S(z) \to S(n) = S(w)$\lbl{A6}}
\UIC{$\forall w\,(z = w \to S(n) = S(z) \to S(n) = S(w))$}
\UIC{$\forall z\,\forall w\,(z = w \to S(n) = S(z) \to S(n) = S(w))$}
\AXC{$\varphi_{22}$\lbl{A2}}
\BIC{$\forall w\,(n = w \to S(n) = S(n) \to S(n) = S(w))$}
\AXC{$\varphi_{30}$\lbl{A2}}
\BIC{$n = \emptyset \to S(n) = S(n) \to S(n) = S(\emptyset)$}
\AXC{$\stackrel{\Pi_{28}}{\vdots}$}\noLine
\UIC{$\varphi_{106}$}
\BIC{$r \in \mathbb{R} \to 0 < r \to \varphi_{78}$}
\AXC{$\varphi_{120}$\lbl{A1}}
\BIC{$S(\emptyset) = 1_{\omega} \to \varphi_{4} \to r \in \mathbb{R} \to 0 < r \to \varphi_{67}$}
\end{ptree}
\noindent$\Pi_{30}$:
\begin{ptree}
\AXC{$\stackrel{\Pi_{25}}{\vdots}$}\noLine
\UIC{$\varphi_{4}$}
\AXC{$\stackrel{\Pi_{26}}{\vdots}$}\noLine
\UIC{$S(\emptyset) = 1_{\omega}$}
\AXC{$\stackrel{\Pi_{29}}{\vdots}$}\noLine
\UIC{$S(\emptyset) = 1_{\omega} \to \varphi_{4} \to r \in \mathbb{R} \to 0 < r \to \varphi_{67}$}
\BIC{$\varphi_{4} \to r \in \mathbb{R} \to 0 < r \to \varphi_{67}$}
\BIC{$r \in \mathbb{R} \to 0 < r \to \varphi_{67}$}
\AXC{$(r \in \mathbb{R} \to 0 < r \to \varphi_{67}) \to \varphi_{98}$\lbl{A1}}
\BIC{$\varphi_{98}$}
\end{ptree}
\noindent$\Pi_{31}$:
\begin{ptree}
\AXC{$\stackrel{\Pi_{15}}{\vdots}$}\noLine
\UIC{$\forall x\,\forall y\,(x = y \to y = x)$}
\AXC{$\varphi_{1}$\lbl{A2}}
\BIC{$\forall y\,(S(n) = y \to y = S(n))$}
\AXC{$\varphi_{3}$\lbl{A2}}
\BIC{$S(n) = 1_{\omega} \to 1_{\omega} = S(n)$}
\AXC{$\stackrel{\Pi_{30}}{\vdots}$}\noLine
\UIC{$\varphi_{98}$}
\BIC{$r \in \mathbb{R} \to 0 < r \to \varphi_{68}$}
\AXC{$\varphi_{125}$\lbl{A1}}
\BIC{$(r \in \mathbb{R} \to 0 < r \to 0 < {{r}^{1_{\omega}}}) \to \varphi_{109}$}
\end{ptree}
\noindent$\Pi_{32}$:
\begin{ptree}
\AXC{$z = w \to 0 < z \to 0 < w$\lbl{A6}}
\UIC{$\forall w\,(z = w \to 0 < z \to 0 < w)$}
\UIC{$\forall z\,\forall w\,(z = w \to 0 < z \to 0 < w)$}
\AXC{$\varphi_{8}$\lbl{A2}}
\BIC{$\forall w\,(r = w \to 0 < r \to 0 < w)$}
\AXC{$\varphi_{32}$\lbl{A2}}
\BIC{$r = {{r}^{1_{\omega}}} \to 0 < r \to 0 < {{r}^{1_{\omega}}}$}
\AXC{$\stackrel{\Pi_{23}}{\vdots}$}\noLine
\UIC{$\varphi_{49}$}
\BIC{$r \in \mathbb{R} \to 0 < r \to 0 < {{r}^{1_{\omega}}}$}
\AXC{$\stackrel{\Pi_{31}}{\vdots}$}\noLine
\UIC{$(r \in \mathbb{R} \to 0 < r \to 0 < {{r}^{1_{\omega}}}) \to \varphi_{109}$}
\BIC{$\varphi_{109}$}
\end{ptree}
\noindent$\Pi_{33}$:
\begin{ptree}
\AXC{$\stackrel{\Pi_{20}}{\vdots}$}\noLine
\UIC{$r \in \mathbb{R} \to 0 < r \to \varphi_{84}$}
\AXC{$z = w \to 0 < {{r}^{z}} \to 0 < {{r}^{w}}$\lbl{A6}}
\UIC{$\forall w\,(z = w \to 0 < {{r}^{z}} \to 0 < {{r}^{w}})$}
\UIC{$\forall z\,\forall w\,(z = w \to 0 < {{r}^{z}} \to 0 < {{r}^{w}})$}
\AXC{$\varphi_{60}$\lbl{A2}}
\BIC{$\forall w\,(1_{\omega} = w \to 0 < {{r}^{1_{\omega}}} \to 0 < {{r}^{w}})$}
\AXC{$\varphi_{66}$\lbl{A2}}
\BIC{$1_{\omega} = S(n) \to 0 < {{r}^{1_{\omega}}} \to 0 < {{r}^{S(n)}}$}
\AXC{$\stackrel{\Pi_{32}}{\vdots}$}\noLine
\UIC{$\varphi_{109}$}
\BIC{$r \in \mathbb{R} \to 0 < r \to \varphi_{83}$}
\AXC{$\varphi_{129}$\lbl{A1}}
\BIC{$(r \in \mathbb{R} \to 0 < r \to \varphi_{84}) \to r \in \mathbb{R} \to 0 < r \to \varphi_{65}$}
\BIC{$r \in \mathbb{R} \to 0 < r \to \varphi_{65}$}
\end{ptree}
\noindent$\Pi_{34}$:
\begin{ptree}
\AXC{$\forall x\,(\mathrm{Ind}(x) \leftrightarrow \emptyset \in x \land \forall y \in x\,(S(y) \in x))$\lbl{DefInd}}
\AXC{$\varphi_{81}$\lbl{A2}}
\BIC{$\mathrm{Ind}(\omega) \leftrightarrow \emptyset \in \omega \land \forall y \in \omega\,(S(y) \in \omega)$}
\AXC{$\mathrm{Ind}(\omega)$\lbl{補題(ωは帰納的)}}
\AXC{$\mathrm{Ind}(\omega) \to \varphi_{110}$\lbl{A1}}
\BIC{$\varphi_{110}$}
\BIC{$r \in \mathbb{R} \to 0 < r \to \varphi_{74}$}
\end{ptree}
\noindent$\Pi_{35}$:
\begin{ptree}
\AXC{$\stackrel{\Pi_{34}}{\vdots}$}\noLine
\UIC{$r \in \mathbb{R} \to 0 < r \to \varphi_{74}$}
\AXC{$(r \in \mathbb{R} \to 0 < r \to \varphi_{74}) \to r \in \mathbb{R} \to 0 < r \to \varphi_{61}$\lbl{A1}}
\BIC{$r \in \mathbb{R} \to 0 < r \to \varphi_{61}$}
\AXC{$(r \in \mathbb{R} \to 0 < r \to \varphi_{61}) \to \varphi_{108}$\lbl{A1}}
\BIC{$\varphi_{108}$}
\end{ptree}
\noindent$\Pi_{36}$:
\begin{ptree}
\AXC{$\stackrel{\Pi_{5}}{\vdots}$}\noLine
\UIC{$r \in \mathbb{R} \to 0 < r \to \varphi_{40}$}
\AXC{$\forall y \in \omega\,(S(y) \in \omega) \to n \in \omega \to S(n) \in \omega$\lbl{A2}}
\AXC{$\stackrel{\Pi_{35}}{\vdots}$}\noLine
\UIC{$\varphi_{108}$}
\BIC{$r \in \mathbb{R} \to 0 < r \to \varphi_{58}$}
\AXC{$\varphi_{121}$\lbl{A1}}
\BIC{$(r \in \mathbb{R} \to 0 < r \to \varphi_{40}) \to r \in \mathbb{R} \to 0 < r \to \varphi_{46}$}
\BIC{$r \in \mathbb{R} \to 0 < r \to \varphi_{46}$}
\end{ptree}
\noindent$\Pi_{37}$:
\begin{ptree}
\AXC{$\forall n\,(n \in B_{32} \leftrightarrow n \in \omega \land (n \in \mathbb{N} \to 0 < {{r}^{n}})) \to \varphi_{24}$\lbl{A2}}
\AXC{$r \in \mathbb{R} \to 0 < r \to \varphi_{79}$\lbl{A1}}
\AXC{$(r \in \mathbb{R} \to 0 < r \to \varphi_{79}) \to \varphi_{118}$\lbl{A1}}
\BIC{$\varphi_{118}$}
\BIC{$r \in \mathbb{R} \to \varphi_{91}$}
\end{ptree}
\noindent$\Pi_{38}$:
\begin{ptree}
\AXC{$\stackrel{\Pi_{33}}{\vdots}$}\noLine
\UIC{$r \in \mathbb{R} \to 0 < r \to \varphi_{65}$}
\AXC{$\stackrel{\Pi_{36}}{\vdots}$}\noLine
\UIC{$r \in \mathbb{R} \to 0 < r \to \varphi_{46}$}
\AXC{$\stackrel{\Pi_{37}}{\vdots}$}\noLine
\UIC{$r \in \mathbb{R} \to \varphi_{91}$}
\AXC{$(r \in \mathbb{R} \to \varphi_{91}) \to \varphi_{122}$\lbl{A1}}
\BIC{$\varphi_{122}$}
\BIC{$(r \in \mathbb{R} \to 0 < r \to \varphi_{65}) \to r \in \mathbb{R} \to 0 < r \to \varphi_{42}$}
\BIC{$r \in \mathbb{R} \to 0 < r \to \varphi_{42}$}
\end{ptree}
\noindent$\Pi_{39}$:
\begin{ptree}
\AXC{$\varphi_{113}$\lbl{A3}}
\AXC{$\stackrel{\Pi_{38}}{\vdots}$}\noLine
\UIC{$r \in \mathbb{R} \to 0 < r \to \varphi_{42}$}
\AXC{$(r \in \mathbb{R} \to 0 < r \to \varphi_{42}) \to \varphi_{37} \to n \in B_{32} \to S(n) \in B_{32}$\lbl{A1}}
\BIC{$\varphi_{37} \to n \in B_{32} \to S(n) \in B_{32}$}
\UIC{$\forall n\,(\varphi_{37} \to n \in B_{32} \to S(n) \in B_{32})$}
\AXC{$\varphi_{133}$\lbl{A1}}
\BIC{$\varphi_{113} \to (\varphi_{37} \to \forall n\,(\varphi_{37})) \to r \in \mathbb{R} \to 0 < r \to \varphi_{47}$}
\BIC{$(\varphi_{37} \to \forall n\,(\varphi_{37})) \to r \in \mathbb{R} \to 0 < r \to \varphi_{47}$}
\end{ptree}
\noindent$\Pi_{40}$:
\begin{ptree}
\AXC{$x = x$\lbl{A5}}
\UIC{$\forall x\,(x = x)$}
\AXC{$\forall x\,(x = x) \to \emptyset = \emptyset$\lbl{A2}}
\BIC{$\emptyset = \emptyset$}
\AXC{$\forall n\,(n \in \mathbb{N} \leftrightarrow n \in \omega \land n \neq \emptyset)$\lbl{Def$\mathbb{N}$}}
\AXC{$\varphi_{56}$\lbl{A2}}
\BIC{$\emptyset \in \mathbb{N} \leftrightarrow \emptyset \in \omega \land \emptyset \neq \emptyset$}
\AXC{$\varphi_{94}$\lbl{A1}}
\BIC{$\emptyset = \emptyset \to r \in \mathbb{R} \to 0 < r \to \emptyset \in \mathbb{N} \to 0 < {{r}^{\emptyset}}$}
\BIC{$r \in \mathbb{R} \to 0 < r \to \emptyset \in \mathbb{N} \to 0 < {{r}^{\emptyset}}$}
\end{ptree}
\noindent$\Pi_{41}$:
\begin{ptree}
\AXC{$\varphi_{93}$\lbl{A2}}
\AXC{$r \in \mathbb{R} \to 0 < r \to \varphi_{79}$\lbl{A1}}
\AXC{$(r \in \mathbb{R} \to 0 < r \to \varphi_{79}) \to \varphi_{93} \to r \in \mathbb{R} \to 0 < r \to \varphi_{93}$\lbl{A1}}
\BIC{$\varphi_{93} \to r \in \mathbb{R} \to 0 < r \to \varphi_{93}$}
\BIC{$r \in \mathbb{R} \to 0 < r \to \varphi_{93}$}
\end{ptree}
\noindent$\Pi_{42}$:
\begin{ptree}
\AXC{$\stackrel{\Pi_{34}}{\vdots}$}\noLine
\UIC{$r \in \mathbb{R} \to 0 < r \to \varphi_{74}$}
\AXC{$\stackrel{\Pi_{41}}{\vdots}$}\noLine
\UIC{$r \in \mathbb{R} \to 0 < r \to \varphi_{93}$}
\AXC{$(r \in \mathbb{R} \to 0 < r \to \varphi_{93}) \to (r \in \mathbb{R} \to 0 < r \to \varphi_{74}) \to \varphi_{104}$\lbl{A1}}
\BIC{$(r \in \mathbb{R} \to 0 < r \to \varphi_{74}) \to \varphi_{104}$}
\BIC{$\varphi_{104}$}
\end{ptree}
\noindent$\Pi_{43}$:
\begin{ptree}
\AXC{$\forall x\,(\mathrm{Ind}(x) \leftrightarrow \emptyset \in x \land \forall y \in x\,(S(y) \in x))$\lbl{DefInd}}
\AXC{$\forall x\,(\mathrm{Ind}(x) \leftrightarrow \emptyset \in x \land \forall y \in x\,(S(y) \in x)) \to \varphi_{14}$\lbl{A2}}
\BIC{$\varphi_{14}$}
\AXC{$\varphi_{14} \to \varphi_{130}$\lbl{A1}}
\BIC{$\varphi_{130}$}
\end{ptree}
\noindent$\Pi_{44}$:
\begin{ptree}
\AXC{$\stackrel{\Pi_{40}}{\vdots}$}\noLine
\UIC{$r \in \mathbb{R} \to 0 < r \to \emptyset \in \mathbb{N} \to 0 < {{r}^{\emptyset}}$}
\AXC{$\stackrel{\Pi_{42}}{\vdots}$}\noLine
\UIC{$\varphi_{104}$}
\BIC{$r \in \mathbb{R} \to 0 < r \to \varphi_{33}$}
\AXC{$\stackrel{\Pi_{43}}{\vdots}$}\noLine
\UIC{$\varphi_{130}$}
\BIC{$(r \in \mathbb{R} \to 0 < r \to \varphi_{47}) \to \varphi_{13} \to \varphi_{12} \to r \in \mathbb{R} \to 0 < r \to \varphi_{25}$}
\end{ptree}
\noindent$\Pi_{45}$:
\begin{ptree}
\AXC{$\varphi_{37} \to \forall n\,(\varphi_{37})$\lbl{A4}}
\AXC{$\stackrel{\Pi_{39}}{\vdots}$}\noLine
\UIC{$(\varphi_{37} \to \forall n\,(\varphi_{37})) \to r \in \mathbb{R} \to 0 < r \to \varphi_{47}$}
\BIC{$r \in \mathbb{R} \to 0 < r \to \varphi_{47}$}
\AXC{$\stackrel{\Pi_{44}}{\vdots}$}\noLine
\UIC{$(r \in \mathbb{R} \to 0 < r \to \varphi_{47}) \to \varphi_{13} \to \varphi_{12} \to r \in \mathbb{R} \to 0 < r \to \varphi_{25}$}
\BIC{$\varphi_{13} \to \varphi_{12} \to r \in \mathbb{R} \to 0 < r \to \varphi_{25}$}
\end{ptree}
\noindent$\Pi_{46}$:
\begin{ptree}
\AXC{$\varphi_{18}$\lbl{A4}}
\AXC{$\forall n\,(\varphi_{7}) \to \varphi_{20}$\lbl{A3}}
\AXC{$\varphi_{7}$\lbl{A2}}
\UIC{$\forall n\,(\varphi_{7})$}
\AXC{$\forall n\,(\varphi_{7}) \to (\forall n\,(\varphi_{7}) \to \varphi_{20}) \to \varphi_{18} \to \varphi_{13}$\lbl{A1}}
\BIC{$(\forall n\,(\varphi_{7}) \to \varphi_{20}) \to \varphi_{18} \to \varphi_{13}$}
\BIC{$\varphi_{18} \to \varphi_{13}$}
\BIC{$\varphi_{13}$}
\AXC{$\stackrel{\Pi_{45}}{\vdots}$}\noLine
\UIC{$\varphi_{13} \to \varphi_{12} \to r \in \mathbb{R} \to 0 < r \to \varphi_{25}$}
\BIC{$\varphi_{12} \to r \in \mathbb{R} \to 0 < r \to \varphi_{25}$}
\end{ptree}
\noindent$\Pi_{47}$:
\begin{ptree}
\AXC{$\varphi_{17}$\lbl{A4}}
\AXC{$\forall y\,(\varphi_{6}) \to \varphi_{19}$\lbl{A3}}
\AXC{$\varphi_{6}$\lbl{A2}}
\UIC{$\forall y\,(\varphi_{6})$}
\AXC{$\forall y\,(\varphi_{6}) \to (\forall y\,(\varphi_{6}) \to \varphi_{19}) \to \varphi_{17} \to \varphi_{12}$\lbl{A1}}
\BIC{$(\forall y\,(\varphi_{6}) \to \varphi_{19}) \to \varphi_{17} \to \varphi_{12}$}
\BIC{$\varphi_{17} \to \varphi_{12}$}
\BIC{$\varphi_{12}$}
\AXC{$\stackrel{\Pi_{46}}{\vdots}$}\noLine
\UIC{$\varphi_{12} \to r \in \mathbb{R} \to 0 < r \to \varphi_{25}$}
\BIC{$r \in \mathbb{R} \to 0 < r \to \varphi_{25}$}
\end{ptree}
\noindent$\Pi_{48}$:
\begin{ptree}
\AXC{$\forall n\,(n \in \omega \leftrightarrow \forall x\,(\mathrm{Ind}(x) \to n \in x))$\lbl{Def$\omega$}}
\AXC{$\varphi_{26}$\lbl{A2}}
\BIC{$n \in \omega \leftrightarrow \forall x\,(\mathrm{Ind}(x) \to n \in x)$}
\AXC{$r \in \mathbb{R} \to 0 < r \to \varphi_{45}$\lbl{A1}}
\AXC{$(r \in \mathbb{R} \to 0 < r \to \varphi_{45}) \to \varphi_{96}$\lbl{A1}}
\BIC{$\varphi_{96}$}
\BIC{$r \in \mathbb{R} \to 0 < r \to \varphi_{50}$}
\end{ptree}
\noindent$\Pi_{49}$:
\begin{ptree}
\AXC{$\forall x\,(\mathrm{Ind}(x) \to n \in x) \to \mathrm{Ind}(B_{32}) \to n \in B_{32}$\lbl{A2}}
\AXC{$\stackrel{\Pi_{48}}{\vdots}$}\noLine
\UIC{$r \in \mathbb{R} \to 0 < r \to \varphi_{50}$}
\AXC{$(r \in \mathbb{R} \to 0 < r \to \varphi_{50}) \to \varphi_{102}$\lbl{A1}}
\BIC{$\varphi_{102}$}
\BIC{$r \in \mathbb{R} \to 0 < r \to \varphi_{51}$}
\end{ptree}
\noindent$\Pi_{50}$:
\begin{ptree}
\AXC{$\stackrel{\Pi_{2}}{\vdots}$}\noLine
\UIC{$r \in \mathbb{R} \to 0 < r \to \varphi_{76}$}
\AXC{$\stackrel{\Pi_{47}}{\vdots}$}\noLine
\UIC{$r \in \mathbb{R} \to 0 < r \to \varphi_{25}$}
\AXC{$\stackrel{\Pi_{49}}{\vdots}$}\noLine
\UIC{$r \in \mathbb{R} \to 0 < r \to \varphi_{51}$}
\AXC{$(r \in \mathbb{R} \to 0 < r \to \varphi_{51}) \to \varphi_{123}$\lbl{A1}}
\BIC{$\varphi_{123}$}
\BIC{$(r \in \mathbb{R} \to 0 < r \to \varphi_{76}) \to r \in \mathbb{R} \to 0 < r \to \varphi_{62}$}
\BIC{$r \in \mathbb{R} \to 0 < r \to \varphi_{62}$}
\end{ptree}
\noindent$\Pi_{51}$:
\begin{ptree}
\AXC{$\varphi_{117}$\lbl{A3}}
\AXC{$\stackrel{\Pi_{50}}{\vdots}$}\noLine
\UIC{$r \in \mathbb{R} \to 0 < r \to \varphi_{62}$}
\AXC{$(r \in \mathbb{R} \to 0 < r \to \varphi_{62}) \to \varphi_{37} \to n \in \omega \to n \in \mathbb{N} \to 0 < {{r}^{n}}$\lbl{A1}}
\BIC{$\varphi_{37} \to n \in \omega \to n \in \mathbb{N} \to 0 < {{r}^{n}}$}
\UIC{$\forall n\,(\varphi_{37} \to n \in \omega \to n \in \mathbb{N} \to 0 < {{r}^{n}})$}
\AXC{$\varphi_{135}$\lbl{A1}}
\BIC{$\varphi_{117} \to (\varphi_{37} \to \forall n\,(\varphi_{37})) \to r \in \mathbb{R} \to 0 < r \to \varphi_{64}$}
\BIC{$(\varphi_{37} \to \forall n\,(\varphi_{37})) \to r \in \mathbb{R} \to 0 < r \to \varphi_{64}$}
\end{ptree}
\noindent$\Pi_{52}$:
\begin{ptree}
\AXC{$\varphi_{37} \to \forall n\,(\varphi_{37})$\lbl{A4}}
\AXC{$\stackrel{\Pi_{51}}{\vdots}$}\noLine
\UIC{$(\varphi_{37} \to \forall n\,(\varphi_{37})) \to r \in \mathbb{R} \to 0 < r \to \varphi_{64}$}
\BIC{$r \in \mathbb{R} \to 0 < r \to \varphi_{64}$}
\AXC{$(r \in \mathbb{R} \to 0 < r \to \varphi_{64}) \to \varphi_{87}$\lbl{A1}}
\BIC{$\varphi_{87}$}
\UIC{$\forall B_{32}\,(\varphi_{87})$}
\end{ptree}
\noindent$\Pi_{53}$:
\begin{ptree}
\AXC{$\forall B_{32}\,(\varphi_{87}) \to \varphi_{31} \to \varphi_{21}$\lbl{A3}}
\AXC{$\stackrel{\Pi_{52}}{\vdots}$}\noLine
\UIC{$\forall B_{32}\,(\varphi_{87})$}
\AXC{$\varphi_{28}$\lbl{Sep}}
\AXC{$\varphi_{28} \to \varphi_{23}$\lbl{A2}}
\BIC{$\varphi_{23}$}
\AXC{$\varphi_{88}$\lbl{A2}}
\BIC{$\exists B_{32}\,\forall n\,(n \in B_{32} \leftrightarrow n \in \omega \land (n \in \mathbb{N} \to 0 < {{r}^{n}}))$}
\AXC{$\varphi_{95}$\lbl{A1}}
\BIC{$r \in \mathbb{R} \to \varphi_{27}$}
\AXC{$\varphi_{136}$\lbl{A1}}
\BIC{$\forall B_{32}\,(\varphi_{87}) \to (\forall B_{32}\,(\varphi_{87}) \to \varphi_{31} \to \varphi_{21}) \to \varphi_{112}$}
\BIC{$(\forall B_{32}\,(\varphi_{87}) \to \varphi_{31} \to \varphi_{21}) \to \varphi_{112}$}
\BIC{$\varphi_{112}$}
\end{ptree}
\noindent$\Pi_{54}$:
\begin{ptree}
\AXC{$r \in \mathbb{R} \land 0 < r \land \neg \forall n \in \omega\,(n \in \mathbb{N} \to 0 < {{r}^{n}}) \to \varphi_{31}$\lbl{A4}}
\AXC{$\stackrel{\Pi_{53}}{\vdots}$}\noLine
\UIC{$\varphi_{112}$}
\BIC{$r \in \mathbb{R} \to 0 < r \to \forall n \in \omega\,(n \in \mathbb{N} \to 0 < {{r}^{n}})$}
\AXC{$(r \in \mathbb{R} \to 0 < r \to \forall n \in \omega\,(n \in \mathbb{N} \to 0 < {{r}^{n}})) \to \varphi_{101}$\lbl{A1}}
\BIC{$\varphi_{101}$}
\end{ptree}
\noindent$\Pi_{55}$:
\begin{ptree}
\AXC{$\forall n \in \omega\,(n \in \mathbb{N} \to 0 < {{r}^{n}}) \to n \in \omega \to n \in \mathbb{N} \to 0 < {{r}^{n}}$\lbl{A2}}
\AXC{$\stackrel{\Pi_{54}}{\vdots}$}\noLine
\UIC{$\varphi_{101}$}
\BIC{$r \in \mathbb{R} \to 0 < r \to n \in \omega \to n \in \mathbb{N} \to 0 < {{r}^{n}}$}
\AXC{$\varphi_{111}$\lbl{A1}}
\BIC{$(r \in \mathbb{R} \to 0 < r \to n \in \mathbb{N} \to n \in \omega) \to \varphi_{54}$}
\end{ptree}
\noindent$\Pi_{56}$:
\begin{ptree}
\AXC{$\forall n\,(n \in \mathbb{N} \leftrightarrow n \in \omega \land n \neq \emptyset)$\lbl{Def$\mathbb{N}$}}
\AXC{$\varphi_{35}$\lbl{A2}}
\BIC{$n \in \mathbb{N} \leftrightarrow n \in \omega \land n \neq \emptyset$}
\AXC{$\stackrel{\Pi_{1}}{\vdots}$}\noLine
\UIC{$\varphi_{41}$}
\BIC{$r \in \mathbb{R} \to 0 < r \to n \in \mathbb{N} \to n \in \omega$}
\AXC{$\stackrel{\Pi_{55}}{\vdots}$}\noLine
\UIC{$(r \in \mathbb{R} \to 0 < r \to n \in \mathbb{N} \to n \in \omega) \to \varphi_{54}$}
\BIC{$\varphi_{54}$}
\end{ptree}
\noindent$\Pi_{57}$:
\begin{ptree}
\AXC{$r \in \mathbb{R} \to 0 < r \to n \in \mathbb{N} \to n \in \mathbb{N}$\lbl{A1}}
\AXC{$\stackrel{\Pi_{56}}{\vdots}$}\noLine
\UIC{$\varphi_{54}$}
\BIC{$r \in \mathbb{R} \to 0 < r \to n \in \mathbb{N} \to 0 < {{r}^{n}}$}
\AXC{$\varphi_{63}$\lbl{A1}}
\BIC{$r \in \mathbb{R} \land 0 < r \to n \in \mathbb{N} \to 0 < {{r}^{n}}$}
\UIC{$\forall n\,(r \in \mathbb{R} \land 0 < r \to n \in \mathbb{N} \to 0 < {{r}^{n}})$}
\AXC{$\forall n\,(r \in \mathbb{R} \land 0 < r \to n \in \mathbb{N} \to 0 < {{r}^{n}}) \to \varphi_{75} \to \varphi_{57}$\lbl{A1}}
\BIC{$\varphi_{75} \to \varphi_{57}$}
\end{ptree}
\noindent 主木:
\begin{ptree}
\AXC{$r \in \mathbb{R} \land 0 < r \to \forall n\,(r \in \mathbb{R} \land 0 < r)$\lbl{A4}}
\AXC{$\varphi_{75}$\lbl{A3}}
\AXC{$\stackrel{\Pi_{57}}{\vdots}$}\noLine
\UIC{$\varphi_{75} \to \varphi_{57}$}
\BIC{$\varphi_{57}$}
\BIC{$r \in \mathbb{R} \to 0 < r \to \forall n \in \mathbb{N}\,(0 < {{r}^{n}})$}
\UIC{$\forall r \in \mathbb{R}\,(0 < r \to \forall n \in \mathbb{N}\,(0 < {{r}^{n}}))$}
\end{ptree}
\begin{abbrevlist}
\item $\varphi_{1} :\equiv \forall x\,\forall y\,(x = y \to y = x) \to \forall y\,(S(n) = y \to y = S(n))$
\item $\varphi_{2} :\equiv \forall x\,\forall y\,(x = y \to y = x) \to \forall y\,(1_{\omega} = y \to y = 1_{\omega})$
\item $\varphi_{3} :\equiv \forall y\,(S(n) = y \to y = S(n)) \to S(n) = 1_{\omega} \to 1_{\omega} = S(n)$
\item $\varphi_{4} :\equiv S(n) = S(\emptyset) \to S(\emptyset) = 1_{\omega} \to S(n) = 1_{\omega}$
\item $\varphi_{5} :\equiv n \in B_{32} \to S(n) \in \mathbb{N} \to n \neq \emptyset \to \mathrm{Seq}({(r)^{\bullet}})$
\item $\varphi_{6} :\equiv \forall n \in B_{32}\,(S(n) \in B_{32}) \to y \in B_{32} \to S(y) \in B_{32}$
\item $\varphi_{7} :\equiv \forall y \in B_{32}\,(S(y) \in B_{32}) \to n \in B_{32} \to S(n) \in B_{32}$
\item $\varphi_{8} :\equiv \forall z\,\forall w\,(z = w \to 0 < z \to 0 < w) \to \forall w\,(r = w \to 0 < r \to 0 < w)$
\item $\varphi_{9} :\equiv \forall x\,\forall y\,\forall z\,(x = y \to y = z \to x = z) \to \forall y\,\forall z\,(S(n) = y \to y = z \to S(n) = z)$
\item $\varphi_{10} :\equiv (1_{\omega} = S(\emptyset) \to S(\emptyset) = 1_{\omega}) \to S(\emptyset) = 1_{\omega}$
\item $\varphi_{11} :\equiv \forall y\,(1_{\omega} = y \to y = 1_{\omega}) \to 1_{\omega} = S(\emptyset) \to S(\emptyset) = 1_{\omega}$
\item $\varphi_{12} :\equiv \forall n \in B_{32}\,(S(n) \in B_{32}) \to \forall y \in B_{32}\,(S(y) \in B_{32})$
\item $\varphi_{13} :\equiv \forall y \in B_{32}\,(S(y) \in B_{32}) \to \forall n \in B_{32}\,(S(n) \in B_{32})$
\item $\varphi_{14} :\equiv \mathrm{Ind}(B_{32}) \leftrightarrow \emptyset \in B_{32} \land \forall y \in B_{32}\,(S(y) \in B_{32})$
\item $\varphi_{15} :\equiv \forall x\,\forall y\,(x = y \to y = x) \to \forall y\,({{r}^{S(n)}} = y \to y = {{r}^{S(n)}})$
\item $\varphi_{16} :\equiv \forall x\,\forall y\,(x = y \to y = x) \to \forall y\,({{r}^{1_{\omega}}} = y \to y = {{r}^{1_{\omega}}})$
\item $\varphi_{17} :\equiv \forall n \in B_{32}\,(S(n) \in B_{32}) \to \forall y\,\forall n \in B_{32}\,(S(n) \in B_{32})$
\item $\varphi_{18} :\equiv \forall y \in B_{32}\,(S(y) \in B_{32}) \to \forall n\,\forall y \in B_{32}\,(S(y) \in B_{32})$
\item $\varphi_{19} :\equiv \forall y\,\forall n \in B_{32}\,(S(n) \in B_{32}) \to \forall y \in B_{32}\,(S(y) \in B_{32})$
\item $\varphi_{20} :\equiv \forall n\,\forall y \in B_{32}\,(S(y) \in B_{32}) \to \forall n \in B_{32}\,(S(n) \in B_{32})$
\item $\varphi_{21} :\equiv \forall B_{32}\,(\neg \forall n\,(n \in B_{32} \leftrightarrow n \in \omega \land (n \in \mathbb{N} \to 0 < {{r}^{n}})))$
\item $\varphi_{22} :\equiv \forall z\,\forall w\,(z = w \to S(n) = S(z) \to S(n) = S(w)) \to \forall w\,(n = w \to S(n) = S(n) \to S(n) = S(w))$
\item $\varphi_{23} :\equiv \forall A_{33}\,\exists B_{32}\,\forall n\,(n \in B_{32} \leftrightarrow n \in A_{33} \land (n \in \mathbb{N} \to 0 < {{r}^{n}}))$
\item $\varphi_{24} :\equiv S(n) \in B_{32} \leftrightarrow S(n) \in \omega \land (S(n) \in \mathbb{N} \to 0 < {{r}^{S(n)}})$
\item $\varphi_{25} :\equiv \forall n\,(n \in B_{32} \leftrightarrow n \in \omega \land (n \in \mathbb{N} \to 0 < {{r}^{n}})) \to \mathrm{Ind}(B_{32})$
\item $\varphi_{26} :\equiv \forall n\,(n \in \omega \leftrightarrow \forall x\,(\mathrm{Ind}(x) \to n \in x)) \to (n \in \omega \leftrightarrow \forall x\,(\mathrm{Ind}(x) \to n \in x))$
\item $\varphi_{27} :\equiv 0 < r \to \exists B_{32}\,\forall n\,(n \in B_{32} \leftrightarrow n \in \omega \land (n \in \mathbb{N} \to 0 < {{r}^{n}}))$
\item $\varphi_{28} :\equiv \forall z_{34}\,\forall A_{33}\,\exists B_{32}\,\forall n\,(n \in B_{32} \leftrightarrow n \in A_{33} \land (n \in \mathbb{N} \to 0 < {{z_{34}}^{n}}))$
\item $\varphi_{29} :\equiv \forall y\,\forall z\,(S(n) = y \to y = z \to S(n) = z) \to \forall z\,(S(n) = S(\emptyset) \to S(\emptyset) = z \to S(n) = z)$
\item $\varphi_{30} :\equiv \forall w\,(n = w \to S(n) = S(n) \to S(n) = S(w)) \to n = \emptyset \to S(n) = S(n) \to S(n) = S(\emptyset)$
\item $\varphi_{31} :\equiv \forall B_{32}\,(r \in \mathbb{R} \land 0 < r \land \neg \forall n \in \omega\,(n \in \mathbb{N} \to 0 < {{r}^{n}}))$
\item $\varphi_{32} :\equiv \forall w\,(r = w \to 0 < r \to 0 < w) \to r = {{r}^{1_{\omega}}} \to 0 < r \to 0 < {{r}^{1_{\omega}}}$
\item $\varphi_{33} :\equiv \forall n\,(n \in B_{32} \leftrightarrow n \in \omega \land (n \in \mathbb{N} \to 0 < {{r}^{n}})) \to \emptyset \in B_{32}$
\item $\varphi_{34} :\equiv ({{r}^{1_{\omega}}} = r \to r = {{r}^{1_{\omega}}}) \to r \in \mathbb{R} \to 0 < r \to r = {{r}^{1_{\omega}}}$
\item $\varphi_{35} :\equiv \forall n\,(n \in \mathbb{N} \leftrightarrow n \in \omega \land n \neq \emptyset) \to (n \in \mathbb{N} \leftrightarrow n \in \omega \land n \neq \emptyset)$
\item $\varphi_{36} :\equiv \forall y\,({{r}^{1_{\omega}}} = y \to y = {{r}^{1_{\omega}}}) \to {{r}^{1_{\omega}}} = r \to r = {{r}^{1_{\omega}}}$
\item $\varphi_{37} :\equiv r \in \mathbb{R} \land 0 < r \land \forall n\,(n \in B_{32} \leftrightarrow n \in \omega \land (n \in \mathbb{N} \to 0 < {{r}^{n}}))$
\item $\varphi_{38} :\equiv \mathrm{Seq}({(r)^{\bullet}}) \land {{r}^{1_{\omega}}} = r \land \forall n \in \mathbb{N}\,({{r}^{S(n)}} = {{r}^{n}} \cdot r)$
\item $\varphi_{39} :\equiv \forall n\,(n \in B_{32} \leftrightarrow n \in \omega \land (n \in \mathbb{N} \to 0 < {{r}^{n}})) \to n \in B_{32} \to n \in B_{32}$
\item $\varphi_{40} :\equiv \forall n\,(n \in B_{32} \leftrightarrow n \in \omega \land (n \in \mathbb{N} \to 0 < {{r}^{n}})) \to n \in B_{32} \to n \in \omega$
\item $\varphi_{41} :\equiv (n \in \mathbb{N} \leftrightarrow n \in \omega \land n \neq \emptyset) \to r \in \mathbb{R} \to 0 < r \to n \in \mathbb{N} \to n \in \omega$
\item $\varphi_{42} :\equiv \forall n\,(n \in B_{32} \leftrightarrow n \in \omega \land (n \in \mathbb{N} \to 0 < {{r}^{n}})) \to n \in B_{32} \to S(n) \in B_{32}$
\item $\varphi_{43} :\equiv \forall z\,\forall w\,(z = w \to 0 < z \to 0 < w) \to \forall w\,({{r}^{n}} \cdot r = w \to 0 < {{r}^{n}} \cdot r \to 0 < w)$
\item $\varphi_{44} :\equiv n \in B_{32} \to S(n) \in \mathbb{N} \to n \neq \emptyset \to n \in \mathbb{N} \to {{r}^{S(n)}} = {{r}^{n}} \cdot r$
\item $\varphi_{45} :\equiv \forall n\,(n \in B_{32} \leftrightarrow n \in \omega \land (n \in \mathbb{N} \to 0 < {{r}^{n}})) \to n \in \omega \to n \in \omega$
\item $\varphi_{46} :\equiv \forall n\,(n \in B_{32} \leftrightarrow n \in \omega \land (n \in \mathbb{N} \to 0 < {{r}^{n}})) \to n \in B_{32} \to S(n) \in \omega$
\item $\varphi_{47} :\equiv \forall n\,(n \in B_{32} \leftrightarrow n \in \omega \land (n \in \mathbb{N} \to 0 < {{r}^{n}})) \to \forall n \in B_{32}\,(S(n) \in B_{32})$
\item $\varphi_{48} :\equiv n \in B_{32} \to S(n) \in \mathbb{N} \to n \neq \emptyset \to \forall n \in \mathbb{N}\,({{r}^{S(n)}} = {{r}^{n}} \cdot r)$
\item $\varphi_{49} :\equiv (r = {{r}^{1_{\omega}}} \to 0 < r \to 0 < {{r}^{1_{\omega}}}) \to r \in \mathbb{R} \to 0 < r \to 0 < {{r}^{1_{\omega}}}$
\item $\varphi_{50} :\equiv \forall n\,(n \in B_{32} \leftrightarrow n \in \omega \land (n \in \mathbb{N} \to 0 < {{r}^{n}})) \to n \in \omega \to \forall x\,(\mathrm{Ind}(x) \to n \in x)$
\item $\varphi_{51} :\equiv \forall n\,(n \in B_{32} \leftrightarrow n \in \omega \land (n \in \mathbb{N} \to 0 < {{r}^{n}})) \to n \in \omega \to \mathrm{Ind}(B_{32}) \to n \in B_{32}$
\item $\varphi_{52} :\equiv \forall a\,\forall n\,(\mathrm{Seq}(a) \to n \in \mathbb{N} \to a_{n} \in \mathbb{R}) \to \forall n\,(\mathrm{Seq}({(r)^{\bullet}}) \to n \in \mathbb{N} \to {{r}^{n}} \in \mathbb{R})$
\item $\varphi_{53} :\equiv \forall n \in \mathbb{N}\,({{r}^{S(n)}} = {{r}^{n}} \cdot r) \to n \in \mathbb{N} \to {{r}^{S(n)}} = {{r}^{n}} \cdot r$
\item $\varphi_{54} :\equiv (r \in \mathbb{R} \to 0 < r \to n \in \mathbb{N} \to n \in \mathbb{N}) \to r \in \mathbb{R} \to 0 < r \to n \in \mathbb{N} \to 0 < {{r}^{n}}$
\item $\varphi_{55} :\equiv \forall n\,(\mathrm{Seq}({(r)^{\bullet}}) \to n \in \mathbb{N} \to {{r}^{n}} \in \mathbb{R}) \to \mathrm{Seq}({(r)^{\bullet}}) \to n \in \mathbb{N} \to {{r}^{n}} \in \mathbb{R}$
\item $\varphi_{56} :\equiv \forall n\,(n \in \mathbb{N} \leftrightarrow n \in \omega \land n \neq \emptyset) \to (\emptyset \in \mathbb{N} \leftrightarrow \emptyset \in \omega \land \emptyset \neq \emptyset)$
\item $\varphi_{57} :\equiv (r \in \mathbb{R} \land 0 < r \to \forall n\,(r \in \mathbb{R} \land 0 < r)) \to r \in \mathbb{R} \to 0 < r \to \forall n \in \mathbb{N}\,(0 < {{r}^{n}})$
\item $\varphi_{58} :\equiv \forall n\,(n \in B_{32} \leftrightarrow n \in \omega \land (n \in \mathbb{N} \to 0 < {{r}^{n}})) \to n \in B_{32} \to n \in \omega \to S(n) \in \omega$
\item $\varphi_{59} :\equiv \forall n\,(n \in B_{32} \leftrightarrow n \in \omega \land (n \in \mathbb{N} \to 0 < {{r}^{n}})) \to n \in B_{32} \to n \in \mathbb{N} \to 0 < {{r}^{n}}$
\item $\varphi_{60} :\equiv \forall z\,\forall w\,(z = w \to 0 < {{r}^{z}} \to 0 < {{r}^{w}}) \to \forall w\,(1_{\omega} = w \to 0 < {{r}^{1_{\omega}}} \to 0 < {{r}^{w}})$
\item $\varphi_{61} :\equiv \forall n\,(n \in B_{32} \leftrightarrow n \in \omega \land (n \in \mathbb{N} \to 0 < {{r}^{n}})) \to n \in B_{32} \to \forall y \in \omega\,(S(y) \in \omega)$
\item $\varphi_{62} :\equiv \forall n\,(n \in B_{32} \leftrightarrow n \in \omega \land (n \in \mathbb{N} \to 0 < {{r}^{n}})) \to n \in \omega \to n \in \mathbb{N} \to 0 < {{r}^{n}}$
\item $\varphi_{63} :\equiv (r \in \mathbb{R} \to 0 < r \to n \in \mathbb{N} \to 0 < {{r}^{n}}) \to r \in \mathbb{R} \land 0 < r \to n \in \mathbb{N} \to 0 < {{r}^{n}}$
\item $\varphi_{64} :\equiv \forall n\,(n \in B_{32} \leftrightarrow n \in \omega \land (n \in \mathbb{N} \to 0 < {{r}^{n}})) \to \forall n \in \omega\,(n \in \mathbb{N} \to 0 < {{r}^{n}})$
\item $\varphi_{65} :\equiv \forall n\,(n \in B_{32} \leftrightarrow n \in \omega \land (n \in \mathbb{N} \to 0 < {{r}^{n}})) \to n \in B_{32} \to S(n) \in \mathbb{N} \to 0 < {{r}^{S(n)}}$
\item $\varphi_{66} :\equiv \forall w\,(1_{\omega} = w \to 0 < {{r}^{1_{\omega}}} \to 0 < {{r}^{w}}) \to 1_{\omega} = S(n) \to 0 < {{r}^{1_{\omega}}} \to 0 < {{r}^{S(n)}}$
\item $\varphi_{67} :\equiv \forall n\,(n \in B_{32} \leftrightarrow n \in \omega \land (n \in \mathbb{N} \to 0 < {{r}^{n}})) \to n \in B_{32} \to S(n) \in \mathbb{N} \to n = \emptyset \to S(n) = 1_{\omega}$
\item $\varphi_{68} :\equiv \forall n\,(n \in B_{32} \leftrightarrow n \in \omega \land (n \in \mathbb{N} \to 0 < {{r}^{n}})) \to n \in B_{32} \to S(n) \in \mathbb{N} \to n = \emptyset \to 1_{\omega} = S(n)$
\item $\varphi_{69} :\equiv \forall n\,(n \in B_{32} \leftrightarrow n \in \omega \land (n \in \mathbb{N} \to 0 < {{r}^{n}})) \to n \in B_{32} \to S(n) \in \mathbb{N} \to n \neq \emptyset \to n \in \mathbb{N}$
\item $\varphi_{70} :\equiv \forall n\,(n \in B_{32} \leftrightarrow n \in \omega \land (n \in \mathbb{N} \to 0 < {{r}^{n}})) \to n \in B_{32} \to S(n) \in \mathbb{N} \to n = \emptyset \to n = \emptyset$
\item $\varphi_{71} :\equiv \forall y\,({{r}^{S(n)}} = y \to y = {{r}^{S(n)}}) \to {{r}^{S(n)}} = {{r}^{n}} \cdot r \to {{r}^{n}} \cdot r = {{r}^{S(n)}}$
\item $\varphi_{72} :\equiv \forall n\,(n \in B_{32} \leftrightarrow n \in \omega \land (n \in \mathbb{N} \to 0 < {{r}^{n}})) \to n \in B_{32} \to S(n) \in \mathbb{N} \to n \neq \emptyset \to r \in \mathbb{R}$
\item $\varphi_{73} :\equiv \forall n\,(n \in B_{32} \leftrightarrow n \in \omega \land (n \in \mathbb{N} \to 0 < {{r}^{n}})) \to n \in B_{32} \to S(n) \in \mathbb{N} \to n \neq \emptyset \to n \neq \emptyset$
\item $\varphi_{74} :\equiv \forall n\,(n \in B_{32} \leftrightarrow n \in \omega \land (n \in \mathbb{N} \to 0 < {{r}^{n}})) \to \emptyset \in \omega \land \forall y \in \omega\,(S(y) \in \omega)$
\item $\varphi_{75} :\equiv \forall n\,(r \in \mathbb{R} \land 0 < r \to n \in \mathbb{N} \to 0 < {{r}^{n}}) \to \forall n\,(r \in \mathbb{R} \land 0 < r) \to \forall n \in \mathbb{N}\,(0 < {{r}^{n}})$
\item $\varphi_{76} :\equiv \forall n\,(n \in B_{32} \leftrightarrow n \in \omega \land (n \in \mathbb{N} \to 0 < {{r}^{n}})) \to (n \in B_{32} \leftrightarrow n \in \omega \land (n \in \mathbb{N} \to 0 < {{r}^{n}}))$
\item $\varphi_{77} :\equiv \forall n\,(n \in B_{32} \leftrightarrow n \in \omega \land (n \in \mathbb{N} \to 0 < {{r}^{n}})) \to n \in B_{32} \to n \in \omega \land (n \in \mathbb{N} \to 0 < {{r}^{n}})$
\item $\varphi_{78} :\equiv \forall n\,(n \in B_{32} \leftrightarrow n \in \omega \land (n \in \mathbb{N} \to 0 < {{r}^{n}})) \to n \in B_{32} \to S(n) \in \mathbb{N} \to n = \emptyset \to S(n) = S(\emptyset)$
\item $\varphi_{79} :\equiv \forall n\,(n \in B_{32} \leftrightarrow n \in \omega \land (n \in \mathbb{N} \to 0 < {{r}^{n}})) \to \forall n\,(n \in B_{32} \leftrightarrow n \in \omega \land (n \in \mathbb{N} \to 0 < {{r}^{n}}))$
\item $\varphi_{80} :\equiv \forall n\,(n \in B_{32} \leftrightarrow n \in \omega \land (n \in \mathbb{N} \to 0 < {{r}^{n}})) \to n \in B_{32} \to S(n) \in \mathbb{N} \to n \neq \emptyset \to 0 < {{r}^{n}}$
\item $\varphi_{81} :\equiv \forall x\,(\mathrm{Ind}(x) \leftrightarrow \emptyset \in x \land \forall y \in x\,(S(y) \in x)) \to (\mathrm{Ind}(\omega) \leftrightarrow \emptyset \in \omega \land \forall y \in \omega\,(S(y) \in \omega))$
\item $\varphi_{82} :\equiv 0 < r \to \forall n\,(n \in B_{32} \leftrightarrow n \in \omega \land (n \in \mathbb{N} \to 0 < {{r}^{n}})) \to \varphi_{5}$
\item $\varphi_{83} :\equiv \forall n\,(n \in B_{32} \leftrightarrow n \in \omega \land (n \in \mathbb{N} \to 0 < {{r}^{n}})) \to n \in B_{32} \to S(n) \in \mathbb{N} \to n = \emptyset \to 0 < {{r}^{S(n)}}$
\item $\varphi_{84} :\equiv \forall n\,(n \in B_{32} \leftrightarrow n \in \omega \land (n \in \mathbb{N} \to 0 < {{r}^{n}})) \to n \in B_{32} \to S(n) \in \mathbb{N} \to n \neq \emptyset \to 0 < {{r}^{S(n)}}$
\item $\varphi_{85} :\equiv \forall n\,(n \in B_{32} \leftrightarrow n \in \omega \land (n \in \mathbb{N} \to 0 < {{r}^{n}})) \to n \in B_{32} \to S(n) \in \mathbb{N} \to n \neq \emptyset \to {{r}^{n}} \in \mathbb{R}$
\item $\varphi_{86} :\equiv \forall n\,(n \in B_{32} \leftrightarrow n \in \omega \land (n \in \mathbb{N} \to 0 < {{r}^{n}})) \to n \in B_{32} \to S(n) \in \mathbb{N} \to n \neq \emptyset \to 0 < {{r}^{n}} \cdot r$
\item $\varphi_{87} :\equiv r \in \mathbb{R} \land 0 < r \land \neg \forall n \in \omega\,(n \in \mathbb{N} \to 0 < {{r}^{n}}) \to \neg \forall n\,(n \in B_{32} \leftrightarrow n \in \omega \land (n \in \mathbb{N} \to 0 < {{r}^{n}}))$
\item $\varphi_{88} :\equiv \varphi_{23} \to \exists B_{32}\,\forall n\,(n \in B_{32} \leftrightarrow n \in \omega \land (n \in \mathbb{N} \to 0 < {{r}^{n}}))$
\item $\varphi_{89} :\equiv \forall n\,(n \in B_{32} \leftrightarrow n \in \omega \land (n \in \mathbb{N} \to 0 < {{r}^{n}})) \to n \in B_{32} \to S(n) \in \mathbb{N} \to n \neq \emptyset \to {{r}^{S(n)}} = {{r}^{n}} \cdot r$
\item $\varphi_{90} :\equiv \forall n\,(n \in B_{32} \leftrightarrow n \in \omega \land (n \in \mathbb{N} \to 0 < {{r}^{n}})) \to n \in B_{32} \to S(n) \in \mathbb{N} \to n \neq \emptyset \to {{r}^{n}} \cdot r = {{r}^{S(n)}}$
\item $\varphi_{91} :\equiv 0 < r \to \forall n\,(n \in B_{32} \leftrightarrow n \in \omega \land (n \in \mathbb{N} \to 0 < {{r}^{n}})) \to \varphi_{24}$
\item $\varphi_{92} :\equiv \forall w\,({{r}^{n}} \cdot r = w \to 0 < {{r}^{n}} \cdot r \to 0 < w) \to {{r}^{n}} \cdot r = {{r}^{S(n)}} \to 0 < {{r}^{n}} \cdot r \to 0 < {{r}^{S(n)}}$
\item $\varphi_{93} :\equiv \forall n\,(n \in B_{32} \leftrightarrow n \in \omega \land (n \in \mathbb{N} \to 0 < {{r}^{n}})) \to (\emptyset \in B_{32} \leftrightarrow \emptyset \in \omega \land (\emptyset \in \mathbb{N} \to 0 < {{r}^{\emptyset}}))$
\item $\varphi_{94} :\equiv (\emptyset \in \mathbb{N} \leftrightarrow \emptyset \in \omega \land \emptyset \neq \emptyset) \to \emptyset = \emptyset \to r \in \mathbb{R} \to 0 < r \to \emptyset \in \mathbb{N} \to 0 < {{r}^{\emptyset}}$
\item $\varphi_{95} :\equiv \exists B_{32}\,\forall n\,(n \in B_{32} \leftrightarrow n \in \omega \land (n \in \mathbb{N} \to 0 < {{r}^{n}})) \to r \in \mathbb{R} \to \varphi_{27}$
\item $\varphi_{96} :\equiv (n \in \omega \leftrightarrow \forall x\,(\mathrm{Ind}(x) \to n \in x)) \to r \in \mathbb{R} \to 0 < r \to \varphi_{50}$
\item $\varphi_{97} :\equiv \forall x\,\forall y\,(x \in \mathbb{R} \to y \in \mathbb{R} \to 0 < x \to 0 < y \to 0 < x \cdot y) \to \forall y\,({{r}^{n}} \in \mathbb{R} \to y \in \mathbb{R} \to 0 < {{r}^{n}} \to 0 < y \to 0 < {{r}^{n}} \cdot y)$
\item $\varphi_{98} :\equiv (S(n) = 1_{\omega} \to 1_{\omega} = S(n)) \to r \in \mathbb{R} \to 0 < r \to \varphi_{68}$
\item $\varphi_{99} :\equiv 0 < r \to \forall n\,(n \in B_{32} \leftrightarrow n \in \omega \land (n \in \mathbb{N} \to 0 < {{r}^{n}})) \to \varphi_{44}$
\item $\varphi_{100} :\equiv 0 < r \to \forall n\,(n \in B_{32} \leftrightarrow n \in \omega \land (n \in \mathbb{N} \to 0 < {{r}^{n}})) \to \varphi_{48}$
\item $\varphi_{101} :\equiv (\forall n \in \omega\,(n \in \mathbb{N} \to 0 < {{r}^{n}}) \to n \in \omega \to n \in \mathbb{N} \to 0 < {{r}^{n}}) \to r \in \mathbb{R} \to 0 < r \to n \in \omega \to n \in \mathbb{N} \to 0 < {{r}^{n}}$
\item $\varphi_{102} :\equiv (\forall x\,(\mathrm{Ind}(x) \to n \in x) \to \mathrm{Ind}(B_{32}) \to n \in B_{32}) \to r \in \mathbb{R} \to 0 < r \to \varphi_{51}$
\item $\varphi_{103} :\equiv (r \in \mathbb{R} \to 0 < r \to r = {{r}^{1_{\omega}}}) \to (r \in \mathbb{R} \to 0 < r \to 0 < r) \to \varphi_{49}$
\item $\varphi_{104} :\equiv (r \in \mathbb{R} \to 0 < r \to \emptyset \in \mathbb{N} \to 0 < {{r}^{\emptyset}}) \to r \in \mathbb{R} \to 0 < r \to \varphi_{33}$
\item $\varphi_{105} :\equiv \forall y\,({{r}^{n}} \in \mathbb{R} \to y \in \mathbb{R} \to 0 < {{r}^{n}} \to 0 < y \to 0 < {{r}^{n}} \cdot y) \to {{r}^{n}} \in \mathbb{R} \to r \in \mathbb{R} \to 0 < {{r}^{n}} \to 0 < r \to 0 < {{r}^{n}} \cdot r$
\item $\varphi_{106} :\equiv (n = \emptyset \to S(n) = S(n) \to S(n) = S(\emptyset)) \to r \in \mathbb{R} \to 0 < r \to \varphi_{78}$
\item $\varphi_{107} :\equiv \forall n\,(n \in B_{32} \leftrightarrow n \in \omega \land (n \in \mathbb{N} \to 0 < {{r}^{n}})) \to n \in B_{32} \to S(n) \in \mathbb{N} \to n \neq \emptyset \to \varphi_{38}$
\item $\varphi_{108} :\equiv (\forall y \in \omega\,(S(y) \in \omega) \to n \in \omega \to S(n) \in \omega) \to r \in \mathbb{R} \to 0 < r \to \varphi_{58}$
\item $\varphi_{109} :\equiv (1_{\omega} = S(n) \to 0 < {{r}^{1_{\omega}}} \to 0 < {{r}^{S(n)}}) \to r \in \mathbb{R} \to 0 < r \to \varphi_{83}$
\item $\varphi_{110} :\equiv (\mathrm{Ind}(\omega) \leftrightarrow \emptyset \in \omega \land \forall y \in \omega\,(S(y) \in \omega)) \to r \in \mathbb{R} \to 0 < r \to \varphi_{74}$
\item $\varphi_{111} :\equiv (r \in \mathbb{R} \to 0 < r \to n \in \omega \to n \in \mathbb{N} \to 0 < {{r}^{n}}) \to (r \in \mathbb{R} \to 0 < r \to n \in \mathbb{N} \to n \in \omega) \to \varphi_{54}$
\item $\varphi_{112} :\equiv (r \in \mathbb{R} \land 0 < r \land \neg \forall n \in \omega\,(n \in \mathbb{N} \to 0 < {{r}^{n}}) \to \varphi_{31}) \to r \in \mathbb{R} \to 0 < r \to \forall n \in \omega\,(n \in \mathbb{N} \to 0 < {{r}^{n}})$
\item $\varphi_{113} :\equiv \forall n\,(\varphi_{37} \to n \in B_{32} \to S(n) \in B_{32}) \to \forall n\,(\varphi_{37}) \to \forall n \in B_{32}\,(S(n) \in B_{32})$
\item $\varphi_{114} :\equiv ({{r}^{n}} \cdot r = {{r}^{S(n)}} \to 0 < {{r}^{n}} \cdot r \to 0 < {{r}^{S(n)}}) \to r \in \mathbb{R} \to 0 < r \to \varphi_{84}$
\item $\varphi_{115} :\equiv ({{r}^{S(n)}} = {{r}^{n}} \cdot r \to {{r}^{n}} \cdot r = {{r}^{S(n)}}) \to r \in \mathbb{R} \to 0 < r \to \varphi_{90}$
\item $\varphi_{116} :\equiv (r \in \mathbb{R} \to \varphi_{38}) \to (r \in \mathbb{R} \to 0 < r \to r \in \mathbb{R}) \to r \in \mathbb{R} \to 0 < r \to \varphi_{38}$
\item $\varphi_{117} :\equiv \forall n\,(\varphi_{37} \to n \in \omega \to n \in \mathbb{N} \to 0 < {{r}^{n}}) \to \forall n\,(\varphi_{37}) \to \forall n \in \omega\,(n \in \mathbb{N} \to 0 < {{r}^{n}})$
\item $\varphi_{118} :\equiv (\forall n\,(n \in B_{32} \leftrightarrow n \in \omega \land (n \in \mathbb{N} \to 0 < {{r}^{n}})) \to \varphi_{24}) \to r \in \mathbb{R} \to \varphi_{91}$
\item $\varphi_{119} :\equiv (r \in \mathbb{R} \to 0 < r \to \varphi_{80}) \to (r \in \mathbb{R} \to 0 < r \to 0 < r) \to r \in \mathbb{R} \to 0 < r \to \varphi_{86}$
\item $\varphi_{120} :\equiv (r \in \mathbb{R} \to 0 < r \to \varphi_{78}) \to S(\emptyset) = 1_{\omega} \to \varphi_{4} \to r \in \mathbb{R} \to 0 < r \to \varphi_{67}$
\item $\varphi_{121} :\equiv (r \in \mathbb{R} \to 0 < r \to \varphi_{58}) \to (r \in \mathbb{R} \to 0 < r \to \varphi_{40}) \to r \in \mathbb{R} \to 0 < r \to \varphi_{46}$
\item $\varphi_{122} :\equiv (r \in \mathbb{R} \to 0 < r \to \varphi_{46}) \to (r \in \mathbb{R} \to 0 < r \to \varphi_{65}) \to r \in \mathbb{R} \to 0 < r \to \varphi_{42}$
\item $\varphi_{123} :\equiv (r \in \mathbb{R} \to 0 < r \to \varphi_{25}) \to (r \in \mathbb{R} \to 0 < r \to \varphi_{76}) \to r \in \mathbb{R} \to 0 < r \to \varphi_{62}$
\item $\varphi_{124} :\equiv (r \in \mathbb{R} \to 0 < r \to \varphi_{40}) \to (r \in \mathbb{R} \to 0 < r \to \varphi_{73}) \to r \in \mathbb{R} \to 0 < r \to \varphi_{69}$
\item $\varphi_{125} :\equiv (r \in \mathbb{R} \to 0 < r \to \varphi_{68}) \to (r \in \mathbb{R} \to 0 < r \to 0 < {{r}^{1_{\omega}}}) \to \varphi_{109}$
\item $\varphi_{126} :\equiv (r \in \mathbb{R} \to 0 < r \to \varphi_{76}) \to (r \in \mathbb{R} \to 0 < r \to \varphi_{39}) \to r \in \mathbb{R} \to 0 < r \to \varphi_{77}$
\item $\varphi_{127} :\equiv (r \in \mathbb{R} \to 0 < r \to \varphi_{59}) \to (r \in \mathbb{R} \to 0 < r \to \varphi_{69}) \to r \in \mathbb{R} \to 0 < r \to \varphi_{80}$
\item $\varphi_{128} :\equiv (r \in \mathbb{R} \to \varphi_{82}) \to (r \in \mathbb{R} \to 0 < r \to \varphi_{69}) \to r \in \mathbb{R} \to 0 < r \to \varphi_{85}$
\item $\varphi_{129} :\equiv (r \in \mathbb{R} \to 0 < r \to \varphi_{83}) \to (r \in \mathbb{R} \to 0 < r \to \varphi_{84}) \to r \in \mathbb{R} \to 0 < r \to \varphi_{65}$
\item $\varphi_{130} :\equiv (r \in \mathbb{R} \to 0 < r \to \varphi_{33}) \to (r \in \mathbb{R} \to 0 < r \to \varphi_{47}) \to \varphi_{13} \to \varphi_{12} \to r \in \mathbb{R} \to 0 < r \to \varphi_{25}$
\item $\varphi_{131} :\equiv (r \in \mathbb{R} \to \varphi_{38}) \to (r \in \mathbb{R} \to 0 < r \to \varphi_{72}) \to r \in \mathbb{R} \to 0 < r \to \varphi_{107}$
\item $\varphi_{132} :\equiv (r \in \mathbb{R} \to \varphi_{99}) \to (r \in \mathbb{R} \to 0 < r \to \varphi_{69}) \to r \in \mathbb{R} \to 0 < r \to \varphi_{89}$
\item $\varphi_{133} :\equiv \forall n\,(\varphi_{37} \to n \in B_{32} \to S(n) \in B_{32}) \to \varphi_{113} \to (\varphi_{37} \to \forall n\,(\varphi_{37})) \to r \in \mathbb{R} \to 0 < r \to \varphi_{47}$
\item $\varphi_{134} :\equiv ({{r}^{n}} \in \mathbb{R} \to r \in \mathbb{R} \to 0 < {{r}^{n}} \to 0 < r \to 0 < {{r}^{n}} \cdot r) \to (r \in \mathbb{R} \to 0 < r \to \varphi_{85}) \to (r \in \mathbb{R} \to 0 < r \to \varphi_{72}) \to \varphi_{119}$
\item $\varphi_{135} :\equiv \forall n\,(\varphi_{37} \to n \in \omega \to n \in \mathbb{N} \to 0 < {{r}^{n}}) \to \varphi_{117} \to (\varphi_{37} \to \forall n\,(\varphi_{37})) \to r \in \mathbb{R} \to 0 < r \to \varphi_{64}$
\item $\varphi_{136} :\equiv (r \in \mathbb{R} \to \varphi_{27}) \to \forall B_{32}\,(\varphi_{87}) \to (\forall B_{32}\,(\varphi_{87}) \to \varphi_{31} \to \varphi_{21}) \to \varphi_{112}$
\end{abbrevlist}
\end{fullproof}

% 自動生成 (prover/build.js): PowNe0  ノード数 376
\begin{fullproof}{[P4]}{$\forall r \in \mathbb{R}\,(r \neq 0 \to \forall n \in \mathbb{N}\,({{r}^{n}} \neq 0))$}{376}\phantomsection\label{tree:PowNe0}
\noindent$\Pi_{1}$:
\begin{ptree}
\AXC{$r \in \mathbb{R} \to r \neq 0 \to n \in \mathbb{N} \to n \in \mathbb{N}$\lbl{A1}}
\AXC{$(r \in \mathbb{R} \to r \neq 0 \to n \in \mathbb{N} \to n \in \mathbb{N}) \to \varphi_{48}$\lbl{A1}}
\BIC{$\varphi_{48}$}
\end{ptree}
\noindent$\Pi_{2}$:
\begin{ptree}
\AXC{$\varphi_{83}$\lbl{A2}}
\AXC{$r \in \mathbb{R} \to r \neq 0 \to \varphi_{86}$\lbl{A1}}
\AXC{$(r \in \mathbb{R} \to r \neq 0 \to \varphi_{86}) \to \varphi_{83} \to r \in \mathbb{R} \to r \neq 0 \to \varphi_{83}$\lbl{A1}}
\BIC{$\varphi_{83} \to r \in \mathbb{R} \to r \neq 0 \to \varphi_{83}$}
\BIC{$r \in \mathbb{R} \to r \neq 0 \to \varphi_{83}$}
\end{ptree}
\noindent$\Pi_{3}$:
\begin{ptree}
\AXC{$x = x$\lbl{A5}}
\AXC{$x = y \to x = x \to y = x$\lbl{A6}}
\AXC{$(x = y \to x = x \to y = x) \to x = x \to x = y \to y = x$\lbl{A1}}
\BIC{$x = x \to x = y \to y = x$}
\BIC{$x = y \to y = x$}
\UIC{$\forall y\,(x = y \to y = x)$}
\UIC{$\forall x\,\forall y\,(x = y \to y = x)$}
\end{ptree}
\noindent$\Pi_{4}$:
\begin{ptree}
\AXC{$\varphi_{83}$\lbl{A2}}
\AXC{$r \in \mathbb{R} \to r \neq 0 \to \varphi_{86}$\lbl{A1}}
\AXC{$(r \in \mathbb{R} \to r \neq 0 \to \varphi_{86}) \to \varphi_{83} \to r \in \mathbb{R} \to r \neq 0 \to \varphi_{83}$\lbl{A1}}
\BIC{$\varphi_{83} \to r \in \mathbb{R} \to r \neq 0 \to \varphi_{83}$}
\BIC{$r \in \mathbb{R} \to r \neq 0 \to \varphi_{83}$}
\AXC{$\varphi_{137}$\lbl{A1}}
\BIC{$(r \in \mathbb{R} \to r \neq 0 \to \varphi_{45}) \to r \in \mathbb{R} \to r \neq 0 \to \varphi_{85}$}
\end{ptree}
\noindent$\Pi_{5}$:
\begin{ptree}
\AXC{$r \in \mathbb{R} \to r \neq 0 \to \varphi_{45}$\lbl{A1}}
\AXC{$\stackrel{\Pi_{4}}{\vdots}$}\noLine
\UIC{$(r \in \mathbb{R} \to r \neq 0 \to \varphi_{45}) \to r \in \mathbb{R} \to r \neq 0 \to \varphi_{85}$}
\BIC{$r \in \mathbb{R} \to r \neq 0 \to \varphi_{85}$}
\end{ptree}
\noindent$\Pi_{6}$:
\begin{ptree}
\AXC{$\stackrel{\Pi_{5}}{\vdots}$}\noLine
\UIC{$r \in \mathbb{R} \to r \neq 0 \to \varphi_{85}$}
\AXC{$(r \in \mathbb{R} \to r \neq 0 \to \varphi_{85}) \to r \in \mathbb{R} \to r \neq 0 \to \varphi_{47}$\lbl{A1}}
\BIC{$r \in \mathbb{R} \to r \neq 0 \to \varphi_{47}$}
\end{ptree}
\noindent$\Pi_{7}$:
\begin{ptree}
\AXC{$\stackrel{\Pi_{6}}{\vdots}$}\noLine
\UIC{$r \in \mathbb{R} \to r \neq 0 \to \varphi_{47}$}
\AXC{$\forall n\,(n \in \mathbb{N} \leftrightarrow n \in \omega \land n \neq \emptyset)$\lbl{Def$\mathbb{N}$}}
\AXC{$\varphi_{43}$\lbl{A2}}
\BIC{$n \in \mathbb{N} \leftrightarrow n \in \omega \land n \neq \emptyset$}
\AXC{$(n \in \mathbb{N} \leftrightarrow n \in \omega \land n \neq \emptyset) \to \varphi_{136}$\lbl{A1}}
\BIC{$\varphi_{136}$}
\BIC{$(r \in \mathbb{R} \to r \neq 0 \to \varphi_{81}) \to r \in \mathbb{R} \to r \neq 0 \to \varphi_{78}$}
\end{ptree}
\noindent$\Pi_{8}$:
\begin{ptree}
\AXC{$r \in \mathbb{R} \to r \neq 0 \to \varphi_{81}$\lbl{A1}}
\AXC{$\stackrel{\Pi_{7}}{\vdots}$}\noLine
\UIC{$(r \in \mathbb{R} \to r \neq 0 \to \varphi_{81}) \to r \in \mathbb{R} \to r \neq 0 \to \varphi_{78}$}
\BIC{$r \in \mathbb{R} \to r \neq 0 \to \varphi_{78}$}
\end{ptree}
\noindent$\Pi_{9}$:
\begin{ptree}
\AXC{$r \in \mathbb{R} \to r \neq 0 \to \varphi_{80}$\lbl{A1}}
\AXC{$\forall r \in \mathbb{R}\,(\varphi_{46})$\lbl{[S24]}}
\AXC{$\forall r \in \mathbb{R}\,(\varphi_{46}) \to r \in \mathbb{R} \to \varphi_{46}$\lbl{A2}}
\BIC{$r \in \mathbb{R} \to \varphi_{46}$}
\AXC{$\varphi_{144}$\lbl{A1}}
\BIC{$(r \in \mathbb{R} \to r \neq 0 \to \varphi_{80}) \to r \in \mathbb{R} \to r \neq 0 \to \varphi_{118}$}
\BIC{$r \in \mathbb{R} \to r \neq 0 \to \varphi_{118}$}
\end{ptree}
\noindent$\Pi_{10}$:
\begin{ptree}
\AXC{$\varphi_{62}$\lbl{A2}}
\AXC{$\stackrel{\Pi_{9}}{\vdots}$}\noLine
\UIC{$r \in \mathbb{R} \to r \neq 0 \to \varphi_{118}$}
\AXC{$(r \in \mathbb{R} \to r \neq 0 \to \varphi_{118}) \to r \in \mathbb{R} \to \varphi_{112}$\lbl{A1}}
\BIC{$r \in \mathbb{R} \to \varphi_{112}$}
\AXC{$(r \in \mathbb{R} \to \varphi_{112}) \to \varphi_{62} \to r \in \mathbb{R} \to \varphi_{111}$\lbl{A1}}
\BIC{$\varphi_{62} \to r \in \mathbb{R} \to \varphi_{111}$}
\BIC{$r \in \mathbb{R} \to \varphi_{111}$}
\end{ptree}
\noindent$\Pi_{11}$:
\begin{ptree}
\AXC{$\stackrel{\Pi_{8}}{\vdots}$}\noLine
\UIC{$r \in \mathbb{R} \to r \neq 0 \to \varphi_{78}$}
\AXC{$\stackrel{\Pi_{10}}{\vdots}$}\noLine
\UIC{$r \in \mathbb{R} \to \varphi_{111}$}
\AXC{$\varphi_{145}$\lbl{A1}}
\BIC{$(r \in \mathbb{R} \to r \neq 0 \to \varphi_{78}) \to r \in \mathbb{R} \to r \neq 0 \to \varphi_{101}$}
\BIC{$r \in \mathbb{R} \to r \neq 0 \to \varphi_{101}$}
\AXC{$(r \in \mathbb{R} \to r \neq 0 \to \varphi_{101}) \to \varphi_{124}$\lbl{A1}}
\BIC{$\varphi_{124}$}
\end{ptree}
\noindent$\Pi_{12}$:
\begin{ptree}
\AXC{$\stackrel{\Pi_{3}}{\vdots}$}\noLine
\UIC{$\forall x\,\forall y\,(x = y \to y = x)$}
\AXC{$\varphi_{16}$\lbl{A2}}
\BIC{$\forall y\,({{r}^{S(n)}} = y \to y = {{r}^{S(n)}})$}
\AXC{$\varphi_{84}$\lbl{A2}}
\BIC{${{r}^{S(n)}} = {{r}^{n}} \cdot r \to {{r}^{n}} \cdot r = {{r}^{S(n)}}$}
\AXC{$\stackrel{\Pi_{11}}{\vdots}$}\noLine
\UIC{$\varphi_{124}$}
\BIC{$r \in \mathbb{R} \to r \neq 0 \to \varphi_{102}$}
\end{ptree}
\noindent$\Pi_{13}$:
\begin{ptree}
\AXC{$\stackrel{\Pi_{3}}{\vdots}$}\noLine
\UIC{$\forall x\,\forall y\,(x = y \to y = x)$}
\AXC{$\varphi_{22}$\lbl{A2}}
\BIC{$\forall y\,({{r}^{n}} \cdot r = y \to y = {{r}^{n}} \cdot r)$}
\AXC{$\varphi_{89}$\lbl{A2}}
\BIC{${{r}^{n}} \cdot r = {{r}^{S(n)}} \to {{r}^{S(n)}} = {{r}^{n}} \cdot r$}
\AXC{$\stackrel{\Pi_{12}}{\vdots}$}\noLine
\UIC{$r \in \mathbb{R} \to r \neq 0 \to \varphi_{102}$}
\AXC{$(r \in \mathbb{R} \to r \neq 0 \to \varphi_{102}) \to \varphi_{125}$\lbl{A1}}
\BIC{$\varphi_{125}$}
\BIC{$r \in \mathbb{R} \to r \neq 0 \to \varphi_{101}$}
\end{ptree}
\noindent$\Pi_{14}$:
\begin{ptree}
\AXC{$\stackrel{\Pi_{8}}{\vdots}$}\noLine
\UIC{$r \in \mathbb{R} \to r \neq 0 \to \varphi_{78}$}
\AXC{$\stackrel{\Pi_{5}}{\vdots}$}\noLine
\UIC{$r \in \mathbb{R} \to r \neq 0 \to \varphi_{85}$}
\AXC{$(r \in \mathbb{R} \to r \neq 0 \to \varphi_{85}) \to r \in \mathbb{R} \to r \neq 0 \to \varphi_{68}$\lbl{A1}}
\BIC{$r \in \mathbb{R} \to r \neq 0 \to \varphi_{68}$}
\AXC{$\varphi_{139}$\lbl{A1}}
\BIC{$(r \in \mathbb{R} \to r \neq 0 \to \varphi_{78}) \to r \in \mathbb{R} \to r \neq 0 \to \varphi_{87}$}
\BIC{$r \in \mathbb{R} \to r \neq 0 \to \varphi_{87}$}
\end{ptree}
\noindent$\Pi_{15}$:
\begin{ptree}
\AXC{$r \in \mathbb{R} \to r \neq 0 \to \varphi_{80}$\lbl{A1}}
\AXC{$\forall r \in \mathbb{R}\,(\varphi_{46})$\lbl{[S24]}}
\AXC{$\forall r \in \mathbb{R}\,(\varphi_{46}) \to r \in \mathbb{R} \to \varphi_{46}$\lbl{A2}}
\BIC{$r \in \mathbb{R} \to \varphi_{46}$}
\AXC{$\varphi_{144}$\lbl{A1}}
\BIC{$(r \in \mathbb{R} \to r \neq 0 \to \varphi_{80}) \to r \in \mathbb{R} \to r \neq 0 \to \varphi_{118}$}
\BIC{$r \in \mathbb{R} \to r \neq 0 \to \varphi_{118}$}
\end{ptree}
\noindent$\Pi_{16}$:
\begin{ptree}
\AXC{$\forall a\,\forall n\,(\mathrm{Seq}(a) \to n \in \mathbb{N} \to a_{n} \in \mathbb{R})$\lbl{[S1]}}
\AXC{$\varphi_{63}$\lbl{A2}}
\BIC{$\forall n\,(\mathrm{Seq}({(r)^{\bullet}}) \to n \in \mathbb{N} \to {{r}^{n}} \in \mathbb{R})$}
\AXC{$\varphi_{65}$\lbl{A2}}
\BIC{$\mathrm{Seq}({(r)^{\bullet}}) \to n \in \mathbb{N} \to {{r}^{n}} \in \mathbb{R}$}
\end{ptree}
\noindent$\Pi_{17}$:
\begin{ptree}
\AXC{$\stackrel{\Pi_{16}}{\vdots}$}\noLine
\UIC{$\mathrm{Seq}({(r)^{\bullet}}) \to n \in \mathbb{N} \to {{r}^{n}} \in \mathbb{R}$}
\AXC{$(\mathrm{Seq}({(r)^{\bullet}}) \to n \in \mathbb{N} \to {{r}^{n}} \in \mathbb{R}) \to \varphi_{141}$\lbl{A1}}
\BIC{$\varphi_{141}$}
\end{ptree}
\noindent$\Pi_{18}$:
\begin{ptree}
\AXC{$\stackrel{\Pi_{8}}{\vdots}$}\noLine
\UIC{$r \in \mathbb{R} \to r \neq 0 \to \varphi_{78}$}
\AXC{$\stackrel{\Pi_{15}}{\vdots}$}\noLine
\UIC{$r \in \mathbb{R} \to r \neq 0 \to \varphi_{118}$}
\AXC{$(r \in \mathbb{R} \to r \neq 0 \to \varphi_{118}) \to r \in \mathbb{R} \to \varphi_{90}$\lbl{A1}}
\BIC{$r \in \mathbb{R} \to \varphi_{90}$}
\AXC{$\stackrel{\Pi_{17}}{\vdots}$}\noLine
\UIC{$\varphi_{141}$}
\BIC{$(r \in \mathbb{R} \to r \neq 0 \to \varphi_{78}) \to r \in \mathbb{R} \to r \neq 0 \to \varphi_{95}$}
\BIC{$r \in \mathbb{R} \to r \neq 0 \to \varphi_{95}$}
\end{ptree}
\noindent$\Pi_{19}$:
\begin{ptree}
\AXC{$\stackrel{\Pi_{18}}{\vdots}$}\noLine
\UIC{$r \in \mathbb{R} \to r \neq 0 \to \varphi_{95}$}
\AXC{$\forall x\,\forall y\,(x \in \mathbb{R} \to y \in \mathbb{R} \to x \neq 0 \to y \neq 0 \to x \cdot y \neq 0)$\lbl{[E15]}}
\AXC{$\varphi_{105}$\lbl{A2}}
\BIC{$\forall y\,({{r}^{n}} \in \mathbb{R} \to y \in \mathbb{R} \to {{r}^{n}} \neq 0 \to y \neq 0 \to {{r}^{n}} \cdot y \neq 0)$}
\AXC{$\varphi_{115}$\lbl{A2}}
\BIC{${{r}^{n}} \in \mathbb{R} \to r \in \mathbb{R} \to {{r}^{n}} \neq 0 \to r \neq 0 \to {{r}^{n}} \cdot r \neq 0$}
\AXC{$\varphi_{149}$\lbl{A1}}
\BIC{$(r \in \mathbb{R} \to r \neq 0 \to \varphi_{95}) \to (r \in \mathbb{R} \to r \neq 0 \to \varphi_{80}) \to \varphi_{130}$}
\BIC{$(r \in \mathbb{R} \to r \neq 0 \to \varphi_{80}) \to \varphi_{130}$}
\end{ptree}
\noindent$\Pi_{20}$:
\begin{ptree}
\AXC{$r \in \mathbb{R} \to r \neq 0 \to r \neq 0$\lbl{A1}}
\AXC{$\stackrel{\Pi_{14}}{\vdots}$}\noLine
\UIC{$r \in \mathbb{R} \to r \neq 0 \to \varphi_{87}$}
\AXC{$r \in \mathbb{R} \to r \neq 0 \to \varphi_{80}$\lbl{A1}}
\AXC{$\stackrel{\Pi_{19}}{\vdots}$}\noLine
\UIC{$(r \in \mathbb{R} \to r \neq 0 \to \varphi_{80}) \to \varphi_{130}$}
\BIC{$\varphi_{130}$}
\BIC{$(r \in \mathbb{R} \to r \neq 0 \to r \neq 0) \to r \in \mathbb{R} \to r \neq 0 \to \varphi_{94}$}
\BIC{$r \in \mathbb{R} \to r \neq 0 \to \varphi_{94}$}
\end{ptree}
\noindent$\Pi_{21}$:
\begin{ptree}
\AXC{$\stackrel{\Pi_{20}}{\vdots}$}\noLine
\UIC{$r \in \mathbb{R} \to r \neq 0 \to \varphi_{94}$}
\AXC{$\stackrel{\Pi_{12}}{\vdots}$}\noLine
\UIC{$r \in \mathbb{R} \to r \neq 0 \to \varphi_{102}$}
\AXC{$(r \in \mathbb{R} \to r \neq 0 \to \varphi_{102}) \to \varphi_{146}$\lbl{A1}}
\BIC{$\varphi_{146}$}
\BIC{$({{r}^{n}} \cdot r = {{r}^{S(n)}} \to {{r}^{n}} \cdot r = 0 \to {{r}^{S(n)}} = 0) \to \varphi_{138}$}
\end{ptree}
\noindent$\Pi_{22}$:
\begin{ptree}
\AXC{$z = w \to z = 0 \to w = 0$\lbl{A6}}
\UIC{$\forall w\,(z = w \to z = 0 \to w = 0)$}
\UIC{$\forall z\,\forall w\,(z = w \to z = 0 \to w = 0)$}
\AXC{$\varphi_{35}$\lbl{A2}}
\BIC{$\forall w\,({{r}^{n}} \cdot r = w \to {{r}^{n}} \cdot r = 0 \to w = 0)$}
\AXC{$\varphi_{98}$\lbl{A2}}
\BIC{${{r}^{n}} \cdot r = {{r}^{S(n)}} \to {{r}^{n}} \cdot r = 0 \to {{r}^{S(n)}} = 0$}
\AXC{$\stackrel{\Pi_{21}}{\vdots}$}\noLine
\UIC{$({{r}^{n}} \cdot r = {{r}^{S(n)}} \to {{r}^{n}} \cdot r = 0 \to {{r}^{S(n)}} = 0) \to \varphi_{138}$}
\BIC{$\varphi_{138}$}
\end{ptree}
\noindent$\Pi_{23}$:
\begin{ptree}
\AXC{$\stackrel{\Pi_{13}}{\vdots}$}\noLine
\UIC{$r \in \mathbb{R} \to r \neq 0 \to \varphi_{101}$}
\AXC{$z = w \to z = 0 \to w = 0$\lbl{A6}}
\UIC{$\forall w\,(z = w \to z = 0 \to w = 0)$}
\UIC{$\forall z\,\forall w\,(z = w \to z = 0 \to w = 0)$}
\AXC{$\varphi_{30}$\lbl{A2}}
\BIC{$\forall w\,({{r}^{S(n)}} = w \to {{r}^{S(n)}} = 0 \to w = 0)$}
\AXC{$\varphi_{96}$\lbl{A2}}
\BIC{${{r}^{S(n)}} = {{r}^{n}} \cdot r \to {{r}^{S(n)}} = 0 \to {{r}^{n}} \cdot r = 0$}
\AXC{$\stackrel{\Pi_{22}}{\vdots}$}\noLine
\UIC{$\varphi_{138}$}
\BIC{$(r \in \mathbb{R} \to r \neq 0 \to \varphi_{101}) \to r \in \mathbb{R} \to r \neq 0 \to \varphi_{92}$}
\BIC{$r \in \mathbb{R} \to r \neq 0 \to \varphi_{92}$}
\end{ptree}
\noindent$\Pi_{24}$:
\begin{ptree}
\AXC{$y = z \to x = y \to x = z$\lbl{A6}}
\AXC{$(y = z \to x = y \to x = z) \to x = y \to y = z \to x = z$\lbl{A1}}
\BIC{$x = y \to y = z \to x = z$}
\UIC{$\forall z\,(x = y \to y = z \to x = z)$}
\UIC{$\forall y\,\forall z\,(x = y \to y = z \to x = z)$}
\UIC{$\forall x\,\forall y\,\forall z\,(x = y \to y = z \to x = z)$}
\AXC{$\varphi_{10}$\lbl{A2}}
\BIC{$\forall y\,\forall z\,(S(n) = y \to y = z \to S(n) = z)$}
\AXC{$\varphi_{36}$\lbl{A2}}
\BIC{$\forall z\,(S(n) = S(\emptyset) \to S(\emptyset) = z \to S(n) = z)$}
\end{ptree}
\noindent$\Pi_{25}$:
\begin{ptree}
\AXC{$\stackrel{\Pi_{24}}{\vdots}$}\noLine
\UIC{$\forall z\,(S(n) = S(\emptyset) \to S(\emptyset) = z \to S(n) = z)$}
\AXC{$\forall z\,(S(n) = S(\emptyset) \to S(\emptyset) = z \to S(n) = z) \to \varphi_{6}$\lbl{A2}}
\BIC{$\varphi_{6}$}
\end{ptree}
\noindent$\Pi_{26}$:
\begin{ptree}
\AXC{$\stackrel{\Pi_{3}}{\vdots}$}\noLine
\UIC{$\forall x\,\forall y\,(x = y \to y = x)$}
\AXC{$\varphi_{1}$\lbl{A2}}
\BIC{$\forall y\,(1_{\omega} = y \to y = 1_{\omega})$}
\AXC{$\varphi_{12}$\lbl{A2}}
\BIC{$1_{\omega} = S(\emptyset) \to S(\emptyset) = 1_{\omega}$}
\AXC{$1_{\omega} = S(\emptyset)$\lbl{Def$1_\omega$}}
\AXC{$1_{\omega} = S(\emptyset) \to \varphi_{11}$\lbl{A1}}
\BIC{$\varphi_{11}$}
\BIC{$S(\emptyset) = 1_{\omega}$}
\end{ptree}
\noindent$\Pi_{27}$:
\begin{ptree}
\AXC{$r \in \mathbb{R} \to r \neq 0 \to \varphi_{79}$\lbl{A1}}
\AXC{$(r \in \mathbb{R} \to r \neq 0 \to \varphi_{79}) \to S(n) = S(n) \to \varphi_{117}$\lbl{A1}}
\BIC{$S(n) = S(n) \to \varphi_{117}$}
\end{ptree}
\noindent$\Pi_{28}$:
\begin{ptree}
\AXC{$x = x$\lbl{A5}}
\UIC{$\forall x\,(x = x)$}
\AXC{$\forall x\,(x = x) \to S(n) = S(n)$\lbl{A2}}
\BIC{$S(n) = S(n)$}
\AXC{$\stackrel{\Pi_{27}}{\vdots}$}\noLine
\UIC{$S(n) = S(n) \to \varphi_{117}$}
\BIC{$\varphi_{117}$}
\end{ptree}
\noindent$\Pi_{29}$:
\begin{ptree}
\AXC{$z = w \to S(n) = S(z) \to S(n) = S(w)$\lbl{A6}}
\UIC{$\forall w\,(z = w \to S(n) = S(z) \to S(n) = S(w))$}
\UIC{$\forall z\,\forall w\,(z = w \to S(n) = S(z) \to S(n) = S(w))$}
\AXC{$\varphi_{27}$\lbl{A2}}
\BIC{$\forall w\,(n = w \to S(n) = S(n) \to S(n) = S(w))$}
\AXC{$\varphi_{37}$\lbl{A2}}
\BIC{$n = \emptyset \to S(n) = S(n) \to S(n) = S(\emptyset)$}
\AXC{$\stackrel{\Pi_{28}}{\vdots}$}\noLine
\UIC{$\varphi_{117}$}
\BIC{$r \in \mathbb{R} \to r \neq 0 \to \varphi_{88}$}
\AXC{$\varphi_{133}$\lbl{A1}}
\BIC{$S(\emptyset) = 1_{\omega} \to \varphi_{6} \to r \in \mathbb{R} \to r \neq 0 \to \varphi_{76}$}
\end{ptree}
\noindent$\Pi_{30}$:
\begin{ptree}
\AXC{$\stackrel{\Pi_{25}}{\vdots}$}\noLine
\UIC{$\varphi_{6}$}
\AXC{$\stackrel{\Pi_{26}}{\vdots}$}\noLine
\UIC{$S(\emptyset) = 1_{\omega}$}
\AXC{$\stackrel{\Pi_{29}}{\vdots}$}\noLine
\UIC{$S(\emptyset) = 1_{\omega} \to \varphi_{6} \to r \in \mathbb{R} \to r \neq 0 \to \varphi_{76}$}
\BIC{$\varphi_{6} \to r \in \mathbb{R} \to r \neq 0 \to \varphi_{76}$}
\BIC{$r \in \mathbb{R} \to r \neq 0 \to \varphi_{76}$}
\AXC{$(r \in \mathbb{R} \to r \neq 0 \to \varphi_{76}) \to \varphi_{108}$\lbl{A1}}
\BIC{$\varphi_{108}$}
\end{ptree}
\noindent$\Pi_{31}$:
\begin{ptree}
\AXC{$\stackrel{\Pi_{3}}{\vdots}$}\noLine
\UIC{$\forall x\,\forall y\,(x = y \to y = x)$}
\AXC{$\varphi_{2}$\lbl{A2}}
\BIC{$\forall y\,(S(n) = y \to y = S(n))$}
\AXC{$\varphi_{4}$\lbl{A2}}
\BIC{$S(n) = 1_{\omega} \to 1_{\omega} = S(n)$}
\AXC{$\stackrel{\Pi_{30}}{\vdots}$}\noLine
\UIC{$\varphi_{108}$}
\BIC{$r \in \mathbb{R} \to r \neq 0 \to \varphi_{77}$}
\end{ptree}
\noindent$\Pi_{32}$:
\begin{ptree}
\AXC{$\stackrel{\Pi_{3}}{\vdots}$}\noLine
\UIC{$\forall x\,\forall y\,(x = y \to y = x)$}
\AXC{$\varphi_{1}$\lbl{A2}}
\BIC{$\forall y\,(1_{\omega} = y \to y = 1_{\omega})$}
\AXC{$\varphi_{3}$\lbl{A2}}
\BIC{$1_{\omega} = S(n) \to S(n) = 1_{\omega}$}
\AXC{$\stackrel{\Pi_{31}}{\vdots}$}\noLine
\UIC{$r \in \mathbb{R} \to r \neq 0 \to \varphi_{77}$}
\AXC{$(r \in \mathbb{R} \to r \neq 0 \to \varphi_{77}) \to \varphi_{109}$\lbl{A1}}
\BIC{$\varphi_{109}$}
\BIC{$r \in \mathbb{R} \to r \neq 0 \to \varphi_{76}$}
\end{ptree}
\noindent$\Pi_{33}$:
\begin{ptree}
\AXC{$\stackrel{\Pi_{3}}{\vdots}$}\noLine
\UIC{$\forall x\,\forall y\,(x = y \to y = x)$}
\AXC{$\forall x\,\forall y\,(x = y \to y = x) \to \forall y\,(r = y \to y = r)$\lbl{A2}}
\BIC{$\forall y\,(r = y \to y = r)$}
\end{ptree}
\noindent$\Pi_{34}$:
\begin{ptree}
\AXC{$\stackrel{\Pi_{33}}{\vdots}$}\noLine
\UIC{$\forall y\,(r = y \to y = r)$}
\AXC{$\forall y\,(r = y \to y = r) \to r = {{r}^{1_{\omega}}} \to {{r}^{1_{\omega}}} = r$\lbl{A2}}
\BIC{$r = {{r}^{1_{\omega}}} \to {{r}^{1_{\omega}}} = r$}
\end{ptree}
\noindent$\Pi_{35}$:
\begin{ptree}
\AXC{$r \in \mathbb{R} \to r \neq 0 \to r \in \mathbb{R}$\lbl{A1}}
\AXC{$\forall r \in \mathbb{R}\,(\varphi_{46})$\lbl{[S24]}}
\AXC{$\forall r \in \mathbb{R}\,(\varphi_{46}) \to r \in \mathbb{R} \to \varphi_{46}$\lbl{A2}}
\BIC{$r \in \mathbb{R} \to \varphi_{46}$}
\AXC{$\varphi_{126}$\lbl{A1}}
\BIC{$(r \in \mathbb{R} \to r \neq 0 \to r \in \mathbb{R}) \to r \in \mathbb{R} \to r \neq 0 \to \varphi_{46}$}
\BIC{$r \in \mathbb{R} \to r \neq 0 \to \varphi_{46}$}
\end{ptree}
\noindent$\Pi_{36}$:
\begin{ptree}
\AXC{$\stackrel{\Pi_{35}}{\vdots}$}\noLine
\UIC{$r \in \mathbb{R} \to r \neq 0 \to \varphi_{46}$}
\AXC{$(r \in \mathbb{R} \to r \neq 0 \to \varphi_{46}) \to r \in \mathbb{R} \to r \neq 0 \to {{r}^{1_{\omega}}} = r$\lbl{A1}}
\BIC{$r \in \mathbb{R} \to r \neq 0 \to {{r}^{1_{\omega}}} = r$}
\AXC{$(r \in \mathbb{R} \to r \neq 0 \to {{r}^{1_{\omega}}} = r) \to \varphi_{40}$\lbl{A1}}
\BIC{$\varphi_{40}$}
\end{ptree}
\noindent$\Pi_{37}$:
\begin{ptree}
\AXC{$\stackrel{\Pi_{3}}{\vdots}$}\noLine
\UIC{$\forall x\,\forall y\,(x = y \to y = x)$}
\AXC{$\varphi_{17}$\lbl{A2}}
\BIC{$\forall y\,({{r}^{1_{\omega}}} = y \to y = {{r}^{1_{\omega}}})$}
\AXC{$\varphi_{44}$\lbl{A2}}
\BIC{${{r}^{1_{\omega}}} = r \to r = {{r}^{1_{\omega}}}$}
\AXC{$\stackrel{\Pi_{36}}{\vdots}$}\noLine
\UIC{$\varphi_{40}$}
\BIC{$r \in \mathbb{R} \to r \neq 0 \to r = {{r}^{1_{\omega}}}$}
\end{ptree}
\noindent$\Pi_{38}$:
\begin{ptree}
\AXC{$\stackrel{\Pi_{34}}{\vdots}$}\noLine
\UIC{$r = {{r}^{1_{\omega}}} \to {{r}^{1_{\omega}}} = r$}
\AXC{$\stackrel{\Pi_{37}}{\vdots}$}\noLine
\UIC{$r \in \mathbb{R} \to r \neq 0 \to r = {{r}^{1_{\omega}}}$}
\AXC{$(r \in \mathbb{R} \to r \neq 0 \to r = {{r}^{1_{\omega}}}) \to \varphi_{41}$\lbl{A1}}
\BIC{$\varphi_{41}$}
\BIC{$r \in \mathbb{R} \to r \neq 0 \to {{r}^{1_{\omega}}} = r$}
\end{ptree}
\noindent$\Pi_{39}$:
\begin{ptree}
\AXC{$r \in \mathbb{R} \to r \neq 0 \to r \neq 0$\lbl{A1}}
\AXC{$\stackrel{\Pi_{37}}{\vdots}$}\noLine
\UIC{$r \in \mathbb{R} \to r \neq 0 \to r = {{r}^{1_{\omega}}}$}
\AXC{$\varphi_{127}$\lbl{A1}}
\BIC{$(r \in \mathbb{R} \to r \neq 0 \to r \neq 0) \to \varphi_{116}$}
\BIC{$\varphi_{116}$}
\end{ptree}
\noindent$\Pi_{40}$:
\begin{ptree}
\AXC{$z = w \to z = 0 \to w = 0$\lbl{A6}}
\UIC{$\forall w\,(z = w \to z = 0 \to w = 0)$}
\UIC{$\forall z\,\forall w\,(z = w \to z = 0 \to w = 0)$}
\AXC{$\varphi_{5}$\lbl{A2}}
\BIC{$\forall w\,(r = w \to r = 0 \to w = 0)$}
\AXC{$\varphi_{24}$\lbl{A2}}
\BIC{$r = {{r}^{1_{\omega}}} \to r = 0 \to {{r}^{1_{\omega}}} = 0$}
\AXC{$\stackrel{\Pi_{39}}{\vdots}$}\noLine
\UIC{$\varphi_{116}$}
\BIC{$({{r}^{1_{\omega}}} = r \to {{r}^{1_{\omega}}} = 0 \to r = 0) \to \varphi_{42}$}
\end{ptree}
\noindent$\Pi_{41}$:
\begin{ptree}
\AXC{$z = w \to z = 0 \to w = 0$\lbl{A6}}
\UIC{$\forall w\,(z = w \to z = 0 \to w = 0)$}
\UIC{$\forall z\,\forall w\,(z = w \to z = 0 \to w = 0)$}
\AXC{$\varphi_{31}$\lbl{A2}}
\BIC{$\forall w\,({{r}^{1_{\omega}}} = w \to {{r}^{1_{\omega}}} = 0 \to w = 0)$}
\AXC{$\varphi_{58}$\lbl{A2}}
\BIC{${{r}^{1_{\omega}}} = r \to {{r}^{1_{\omega}}} = 0 \to r = 0$}
\AXC{$\stackrel{\Pi_{40}}{\vdots}$}\noLine
\UIC{$({{r}^{1_{\omega}}} = r \to {{r}^{1_{\omega}}} = 0 \to r = 0) \to \varphi_{42}$}
\BIC{$\varphi_{42}$}
\end{ptree}
\noindent$\Pi_{42}$:
\begin{ptree}
\AXC{$\stackrel{\Pi_{38}}{\vdots}$}\noLine
\UIC{$r \in \mathbb{R} \to r \neq 0 \to {{r}^{1_{\omega}}} = r$}
\AXC{$\stackrel{\Pi_{41}}{\vdots}$}\noLine
\UIC{$\varphi_{42}$}
\BIC{$r \in \mathbb{R} \to r \neq 0 \to {{r}^{1_{\omega}}} \neq 0$}
\AXC{$\stackrel{\Pi_{31}}{\vdots}$}\noLine
\UIC{$r \in \mathbb{R} \to r \neq 0 \to \varphi_{77}$}
\AXC{$(r \in \mathbb{R} \to r \neq 0 \to \varphi_{77}) \to \varphi_{142}$\lbl{A1}}
\BIC{$\varphi_{142}$}
\BIC{$(1_{\omega} = S(n) \to {{r}^{1_{\omega}}} = 0 \to {{r}^{S(n)}} = 0) \to \varphi_{131}$}
\end{ptree}
\noindent$\Pi_{43}$:
\begin{ptree}
\AXC{$z = w \to {{r}^{z}} = 0 \to {{r}^{w}} = 0$\lbl{A6}}
\UIC{$\forall w\,(z = w \to {{r}^{z}} = 0 \to {{r}^{w}} = 0)$}
\UIC{$\forall z\,\forall w\,(z = w \to {{r}^{z}} = 0 \to {{r}^{w}} = 0)$}
\AXC{$\varphi_{57}$\lbl{A2}}
\BIC{$\forall w\,(1_{\omega} = w \to {{r}^{1_{\omega}}} = 0 \to {{r}^{w}} = 0)$}
\AXC{$\varphi_{67}$\lbl{A2}}
\BIC{$1_{\omega} = S(n) \to {{r}^{1_{\omega}}} = 0 \to {{r}^{S(n)}} = 0$}
\AXC{$\stackrel{\Pi_{42}}{\vdots}$}\noLine
\UIC{$(1_{\omega} = S(n) \to {{r}^{1_{\omega}}} = 0 \to {{r}^{S(n)}} = 0) \to \varphi_{131}$}
\BIC{$\varphi_{131}$}
\end{ptree}
\noindent$\Pi_{44}$:
\begin{ptree}
\AXC{$\stackrel{\Pi_{32}}{\vdots}$}\noLine
\UIC{$r \in \mathbb{R} \to r \neq 0 \to \varphi_{76}$}
\AXC{$z = w \to {{r}^{z}} = 0 \to {{r}^{w}} = 0$\lbl{A6}}
\UIC{$\forall w\,(z = w \to {{r}^{z}} = 0 \to {{r}^{w}} = 0)$}
\UIC{$\forall z\,\forall w\,(z = w \to {{r}^{z}} = 0 \to {{r}^{w}} = 0)$}
\AXC{$\varphi_{56}$\lbl{A2}}
\BIC{$\forall w\,(S(n) = w \to {{r}^{S(n)}} = 0 \to {{r}^{w}} = 0)$}
\AXC{$\varphi_{66}$\lbl{A2}}
\BIC{$S(n) = 1_{\omega} \to {{r}^{S(n)}} = 0 \to {{r}^{1_{\omega}}} = 0$}
\AXC{$\stackrel{\Pi_{43}}{\vdots}$}\noLine
\UIC{$\varphi_{131}$}
\BIC{$(r \in \mathbb{R} \to r \neq 0 \to \varphi_{76}) \to r \in \mathbb{R} \to r \neq 0 \to \varphi_{91}$}
\BIC{$r \in \mathbb{R} \to r \neq 0 \to \varphi_{91}$}
\AXC{$\varphi_{140}$\lbl{A1}}
\BIC{$(r \in \mathbb{R} \to r \neq 0 \to \varphi_{92}) \to r \in \mathbb{R} \to r \neq 0 \to \varphi_{74}$}
\end{ptree}
\noindent$\Pi_{45}$:
\begin{ptree}
\AXC{$\forall x\,(\mathrm{Ind}(x) \leftrightarrow \emptyset \in x \land \forall y \in x\,(S(y) \in x))$\lbl{DefInd}}
\AXC{$\varphi_{93}$\lbl{A2}}
\BIC{$\mathrm{Ind}(\omega) \leftrightarrow \emptyset \in \omega \land \forall y \in \omega\,(S(y) \in \omega)$}
\AXC{$\mathrm{Ind}(\omega)$\lbl{補題(ωは帰納的)}}
\AXC{$\mathrm{Ind}(\omega) \to \varphi_{120}$\lbl{A1}}
\BIC{$\varphi_{120}$}
\BIC{$r \in \mathbb{R} \to r \neq 0 \to \varphi_{82}$}
\end{ptree}
\noindent$\Pi_{46}$:
\begin{ptree}
\AXC{$\stackrel{\Pi_{45}}{\vdots}$}\noLine
\UIC{$r \in \mathbb{R} \to r \neq 0 \to \varphi_{82}$}
\AXC{$(r \in \mathbb{R} \to r \neq 0 \to \varphi_{82}) \to r \in \mathbb{R} \to r \neq 0 \to \varphi_{72}$\lbl{A1}}
\BIC{$r \in \mathbb{R} \to r \neq 0 \to \varphi_{72}$}
\AXC{$(r \in \mathbb{R} \to r \neq 0 \to \varphi_{72}) \to \varphi_{119}$\lbl{A1}}
\BIC{$\varphi_{119}$}
\end{ptree}
\noindent$\Pi_{47}$:
\begin{ptree}
\AXC{$\stackrel{\Pi_{6}}{\vdots}$}\noLine
\UIC{$r \in \mathbb{R} \to r \neq 0 \to \varphi_{47}$}
\AXC{$\forall y \in \omega\,(S(y) \in \omega) \to n \in \omega \to S(n) \in \omega$\lbl{A2}}
\AXC{$\stackrel{\Pi_{46}}{\vdots}$}\noLine
\UIC{$\varphi_{119}$}
\BIC{$r \in \mathbb{R} \to r \neq 0 \to \varphi_{69}$}
\AXC{$\varphi_{132}$\lbl{A1}}
\BIC{$(r \in \mathbb{R} \to r \neq 0 \to \varphi_{47}) \to r \in \mathbb{R} \to r \neq 0 \to \varphi_{51}$}
\BIC{$r \in \mathbb{R} \to r \neq 0 \to \varphi_{51}$}
\end{ptree}
\noindent$\Pi_{48}$:
\begin{ptree}
\AXC{$\forall n\,(n \in B_{35} \leftrightarrow n \in \omega \land (n \in \mathbb{N} \to {{r}^{n}} \neq 0)) \to \varphi_{26}$\lbl{A2}}
\AXC{$r \in \mathbb{R} \to r \neq 0 \to \varphi_{86}$\lbl{A1}}
\AXC{$(r \in \mathbb{R} \to r \neq 0 \to \varphi_{86}) \to \varphi_{129}$\lbl{A1}}
\BIC{$\varphi_{129}$}
\BIC{$r \in \mathbb{R} \to \varphi_{100}$}
\end{ptree}
\noindent$\Pi_{49}$:
\begin{ptree}
\AXC{$\stackrel{\Pi_{47}}{\vdots}$}\noLine
\UIC{$r \in \mathbb{R} \to r \neq 0 \to \varphi_{51}$}
\AXC{$\stackrel{\Pi_{48}}{\vdots}$}\noLine
\UIC{$r \in \mathbb{R} \to \varphi_{100}$}
\AXC{$(r \in \mathbb{R} \to \varphi_{100}) \to \varphi_{134}$\lbl{A1}}
\BIC{$\varphi_{134}$}
\BIC{$(r \in \mathbb{R} \to r \neq 0 \to \varphi_{74}) \to r \in \mathbb{R} \to r \neq 0 \to \varphi_{49}$}
\end{ptree}
\noindent$\Pi_{50}$:
\begin{ptree}
\AXC{$\stackrel{\Pi_{23}}{\vdots}$}\noLine
\UIC{$r \in \mathbb{R} \to r \neq 0 \to \varphi_{92}$}
\AXC{$\stackrel{\Pi_{44}}{\vdots}$}\noLine
\UIC{$(r \in \mathbb{R} \to r \neq 0 \to \varphi_{92}) \to r \in \mathbb{R} \to r \neq 0 \to \varphi_{74}$}
\BIC{$r \in \mathbb{R} \to r \neq 0 \to \varphi_{74}$}
\AXC{$\stackrel{\Pi_{49}}{\vdots}$}\noLine
\UIC{$(r \in \mathbb{R} \to r \neq 0 \to \varphi_{74}) \to r \in \mathbb{R} \to r \neq 0 \to \varphi_{49}$}
\BIC{$r \in \mathbb{R} \to r \neq 0 \to \varphi_{49}$}
\end{ptree}
\noindent$\Pi_{51}$:
\begin{ptree}
\AXC{$\varphi_{123}$\lbl{A3}}
\AXC{$\stackrel{\Pi_{50}}{\vdots}$}\noLine
\UIC{$r \in \mathbb{R} \to r \neq 0 \to \varphi_{49}$}
\AXC{$(r \in \mathbb{R} \to r \neq 0 \to \varphi_{49}) \to \varphi_{39} \to n \in B_{35} \to S(n) \in B_{35}$\lbl{A1}}
\BIC{$\varphi_{39} \to n \in B_{35} \to S(n) \in B_{35}$}
\UIC{$\forall n\,(\varphi_{39} \to n \in B_{35} \to S(n) \in B_{35})$}
\AXC{$\varphi_{147}$\lbl{A1}}
\BIC{$\varphi_{123} \to (\varphi_{39} \to \forall n\,(\varphi_{39})) \to r \in \mathbb{R} \to r \neq 0 \to \varphi_{53}$}
\BIC{$(\varphi_{39} \to \forall n\,(\varphi_{39})) \to r \in \mathbb{R} \to r \neq 0 \to \varphi_{53}$}
\end{ptree}
\noindent$\Pi_{52}$:
\begin{ptree}
\AXC{$x = x$\lbl{A5}}
\UIC{$\forall x\,(x = x)$}
\AXC{$\forall x\,(x = x) \to \emptyset = \emptyset$\lbl{A2}}
\BIC{$\emptyset = \emptyset$}
\AXC{$\forall n\,(n \in \mathbb{N} \leftrightarrow n \in \omega \land n \neq \emptyset)$\lbl{Def$\mathbb{N}$}}
\AXC{$\varphi_{71}$\lbl{A2}}
\BIC{$\emptyset \in \mathbb{N} \leftrightarrow \emptyset \in \omega \land \emptyset \neq \emptyset$}
\AXC{$\varphi_{104}$\lbl{A1}}
\BIC{$\emptyset = \emptyset \to r \in \mathbb{R} \to r \neq 0 \to \emptyset \in \mathbb{N} \to {{r}^{\emptyset}} \neq 0$}
\BIC{$r \in \mathbb{R} \to r \neq 0 \to \emptyset \in \mathbb{N} \to {{r}^{\emptyset}} \neq 0$}
\end{ptree}
\noindent$\Pi_{53}$:
\begin{ptree}
\AXC{$\varphi_{103}$\lbl{A2}}
\AXC{$r \in \mathbb{R} \to r \neq 0 \to \varphi_{86}$\lbl{A1}}
\AXC{$(r \in \mathbb{R} \to r \neq 0 \to \varphi_{86}) \to \varphi_{103} \to r \in \mathbb{R} \to r \neq 0 \to \varphi_{103}$\lbl{A1}}
\BIC{$\varphi_{103} \to r \in \mathbb{R} \to r \neq 0 \to \varphi_{103}$}
\BIC{$r \in \mathbb{R} \to r \neq 0 \to \varphi_{103}$}
\end{ptree}
\noindent$\Pi_{54}$:
\begin{ptree}
\AXC{$\stackrel{\Pi_{45}}{\vdots}$}\noLine
\UIC{$r \in \mathbb{R} \to r \neq 0 \to \varphi_{82}$}
\AXC{$\stackrel{\Pi_{53}}{\vdots}$}\noLine
\UIC{$r \in \mathbb{R} \to r \neq 0 \to \varphi_{103}$}
\AXC{$(r \in \mathbb{R} \to r \neq 0 \to \varphi_{103}) \to (r \in \mathbb{R} \to r \neq 0 \to \varphi_{82}) \to \varphi_{114}$\lbl{A1}}
\BIC{$(r \in \mathbb{R} \to r \neq 0 \to \varphi_{82}) \to \varphi_{114}$}
\BIC{$\varphi_{114}$}
\end{ptree}
\noindent$\Pi_{55}$:
\begin{ptree}
\AXC{$\forall x\,(\mathrm{Ind}(x) \leftrightarrow \emptyset \in x \land \forall y \in x\,(S(y) \in x))$\lbl{DefInd}}
\AXC{$\forall x\,(\mathrm{Ind}(x) \leftrightarrow \emptyset \in x \land \forall y \in x\,(S(y) \in x)) \to \varphi_{15}$\lbl{A2}}
\BIC{$\varphi_{15}$}
\AXC{$\varphi_{15} \to \varphi_{143}$\lbl{A1}}
\BIC{$\varphi_{143}$}
\end{ptree}
\noindent$\Pi_{56}$:
\begin{ptree}
\AXC{$\stackrel{\Pi_{52}}{\vdots}$}\noLine
\UIC{$r \in \mathbb{R} \to r \neq 0 \to \emptyset \in \mathbb{N} \to {{r}^{\emptyset}} \neq 0$}
\AXC{$\stackrel{\Pi_{54}}{\vdots}$}\noLine
\UIC{$\varphi_{114}$}
\BIC{$r \in \mathbb{R} \to r \neq 0 \to \varphi_{38}$}
\AXC{$\stackrel{\Pi_{55}}{\vdots}$}\noLine
\UIC{$\varphi_{143}$}
\BIC{$(r \in \mathbb{R} \to r \neq 0 \to \varphi_{53}) \to \varphi_{14} \to \varphi_{13} \to r \in \mathbb{R} \to r \neq 0 \to \varphi_{28}$}
\end{ptree}
\noindent$\Pi_{57}$:
\begin{ptree}
\AXC{$\varphi_{39} \to \forall n\,(\varphi_{39})$\lbl{A4}}
\AXC{$\stackrel{\Pi_{51}}{\vdots}$}\noLine
\UIC{$(\varphi_{39} \to \forall n\,(\varphi_{39})) \to r \in \mathbb{R} \to r \neq 0 \to \varphi_{53}$}
\BIC{$r \in \mathbb{R} \to r \neq 0 \to \varphi_{53}$}
\AXC{$\stackrel{\Pi_{56}}{\vdots}$}\noLine
\UIC{$(r \in \mathbb{R} \to r \neq 0 \to \varphi_{53}) \to \varphi_{14} \to \varphi_{13} \to r \in \mathbb{R} \to r \neq 0 \to \varphi_{28}$}
\BIC{$\varphi_{14} \to \varphi_{13} \to r \in \mathbb{R} \to r \neq 0 \to \varphi_{28}$}
\end{ptree}
\noindent$\Pi_{58}$:
\begin{ptree}
\AXC{$\varphi_{19}$\lbl{A4}}
\AXC{$\forall n\,(\varphi_{9}) \to \varphi_{21}$\lbl{A3}}
\AXC{$\varphi_{9}$\lbl{A2}}
\UIC{$\forall n\,(\varphi_{9})$}
\AXC{$\forall n\,(\varphi_{9}) \to (\forall n\,(\varphi_{9}) \to \varphi_{21}) \to \varphi_{19} \to \varphi_{14}$\lbl{A1}}
\BIC{$(\forall n\,(\varphi_{9}) \to \varphi_{21}) \to \varphi_{19} \to \varphi_{14}$}
\BIC{$\varphi_{19} \to \varphi_{14}$}
\BIC{$\varphi_{14}$}
\AXC{$\stackrel{\Pi_{57}}{\vdots}$}\noLine
\UIC{$\varphi_{14} \to \varphi_{13} \to r \in \mathbb{R} \to r \neq 0 \to \varphi_{28}$}
\BIC{$\varphi_{13} \to r \in \mathbb{R} \to r \neq 0 \to \varphi_{28}$}
\end{ptree}
\noindent$\Pi_{59}$:
\begin{ptree}
\AXC{$\varphi_{18}$\lbl{A4}}
\AXC{$\forall y\,(\varphi_{8}) \to \varphi_{20}$\lbl{A3}}
\AXC{$\varphi_{8}$\lbl{A2}}
\UIC{$\forall y\,(\varphi_{8})$}
\AXC{$\forall y\,(\varphi_{8}) \to (\forall y\,(\varphi_{8}) \to \varphi_{20}) \to \varphi_{18} \to \varphi_{13}$\lbl{A1}}
\BIC{$(\forall y\,(\varphi_{8}) \to \varphi_{20}) \to \varphi_{18} \to \varphi_{13}$}
\BIC{$\varphi_{18} \to \varphi_{13}$}
\BIC{$\varphi_{13}$}
\AXC{$\stackrel{\Pi_{58}}{\vdots}$}\noLine
\UIC{$\varphi_{13} \to r \in \mathbb{R} \to r \neq 0 \to \varphi_{28}$}
\BIC{$r \in \mathbb{R} \to r \neq 0 \to \varphi_{28}$}
\end{ptree}
\noindent$\Pi_{60}$:
\begin{ptree}
\AXC{$\forall n\,(n \in \omega \leftrightarrow \forall x\,(\mathrm{Ind}(x) \to n \in x))$\lbl{Def$\omega$}}
\AXC{$\varphi_{34}$\lbl{A2}}
\BIC{$n \in \omega \leftrightarrow \forall x\,(\mathrm{Ind}(x) \to n \in x)$}
\AXC{$r \in \mathbb{R} \to r \neq 0 \to \varphi_{50}$\lbl{A1}}
\AXC{$(r \in \mathbb{R} \to r \neq 0 \to \varphi_{50}) \to \varphi_{107}$\lbl{A1}}
\BIC{$\varphi_{107}$}
\BIC{$r \in \mathbb{R} \to r \neq 0 \to \varphi_{59}$}
\end{ptree}
\noindent$\Pi_{61}$:
\begin{ptree}
\AXC{$\forall x\,(\mathrm{Ind}(x) \to n \in x) \to \mathrm{Ind}(B_{35}) \to n \in B_{35}$\lbl{A2}}
\AXC{$\stackrel{\Pi_{60}}{\vdots}$}\noLine
\UIC{$r \in \mathbb{R} \to r \neq 0 \to \varphi_{59}$}
\AXC{$(r \in \mathbb{R} \to r \neq 0 \to \varphi_{59}) \to \varphi_{113}$\lbl{A1}}
\BIC{$\varphi_{113}$}
\BIC{$r \in \mathbb{R} \to r \neq 0 \to \varphi_{61}$}
\end{ptree}
\noindent$\Pi_{62}$:
\begin{ptree}
\AXC{$\stackrel{\Pi_{2}}{\vdots}$}\noLine
\UIC{$r \in \mathbb{R} \to r \neq 0 \to \varphi_{83}$}
\AXC{$\stackrel{\Pi_{59}}{\vdots}$}\noLine
\UIC{$r \in \mathbb{R} \to r \neq 0 \to \varphi_{28}$}
\AXC{$\stackrel{\Pi_{61}}{\vdots}$}\noLine
\UIC{$r \in \mathbb{R} \to r \neq 0 \to \varphi_{61}$}
\AXC{$(r \in \mathbb{R} \to r \neq 0 \to \varphi_{61}) \to \varphi_{135}$\lbl{A1}}
\BIC{$\varphi_{135}$}
\BIC{$(r \in \mathbb{R} \to r \neq 0 \to \varphi_{83}) \to r \in \mathbb{R} \to r \neq 0 \to \varphi_{70}$}
\BIC{$r \in \mathbb{R} \to r \neq 0 \to \varphi_{70}$}
\end{ptree}
\noindent$\Pi_{63}$:
\begin{ptree}
\AXC{$\varphi_{128}$\lbl{A3}}
\AXC{$\stackrel{\Pi_{62}}{\vdots}$}\noLine
\UIC{$r \in \mathbb{R} \to r \neq 0 \to \varphi_{70}$}
\AXC{$(r \in \mathbb{R} \to r \neq 0 \to \varphi_{70}) \to \varphi_{39} \to n \in \omega \to n \in \mathbb{N} \to {{r}^{n}} \neq 0$\lbl{A1}}
\BIC{$\varphi_{39} \to n \in \omega \to n \in \mathbb{N} \to {{r}^{n}} \neq 0$}
\UIC{$\forall n\,(\varphi_{39} \to n \in \omega \to n \in \mathbb{N} \to {{r}^{n}} \neq 0)$}
\AXC{$\varphi_{148}$\lbl{A1}}
\BIC{$\varphi_{128} \to (\varphi_{39} \to \forall n\,(\varphi_{39})) \to r \in \mathbb{R} \to r \neq 0 \to \varphi_{73}$}
\BIC{$(\varphi_{39} \to \forall n\,(\varphi_{39})) \to r \in \mathbb{R} \to r \neq 0 \to \varphi_{73}$}
\end{ptree}
\noindent$\Pi_{64}$:
\begin{ptree}
\AXC{$\varphi_{39} \to \forall n\,(\varphi_{39})$\lbl{A4}}
\AXC{$\stackrel{\Pi_{63}}{\vdots}$}\noLine
\UIC{$(\varphi_{39} \to \forall n\,(\varphi_{39})) \to r \in \mathbb{R} \to r \neq 0 \to \varphi_{73}$}
\BIC{$r \in \mathbb{R} \to r \neq 0 \to \varphi_{73}$}
\AXC{$(r \in \mathbb{R} \to r \neq 0 \to \varphi_{73}) \to \varphi_{97}$\lbl{A1}}
\BIC{$\varphi_{97}$}
\UIC{$\forall B_{35}\,(\varphi_{97})$}
\end{ptree}
\noindent$\Pi_{65}$:
\begin{ptree}
\AXC{$\forall B_{35}\,(\varphi_{97}) \to \varphi_{32} \to \varphi_{23}$\lbl{A3}}
\AXC{$\stackrel{\Pi_{64}}{\vdots}$}\noLine
\UIC{$\forall B_{35}\,(\varphi_{97})$}
\AXC{$\varphi_{33}$\lbl{Sep}}
\AXC{$\varphi_{33} \to \varphi_{25}$\lbl{A2}}
\BIC{$\varphi_{25}$}
\AXC{$\varphi_{99}$\lbl{A2}}
\BIC{$\exists B_{35}\,\forall n\,(n \in B_{35} \leftrightarrow n \in \omega \land (n \in \mathbb{N} \to {{r}^{n}} \neq 0))$}
\AXC{$\varphi_{106}$\lbl{A1}}
\BIC{$r \in \mathbb{R} \to \varphi_{29}$}
\AXC{$\varphi_{150}$\lbl{A1}}
\BIC{$\forall B_{35}\,(\varphi_{97}) \to (\forall B_{35}\,(\varphi_{97}) \to \varphi_{32} \to \varphi_{23}) \to \varphi_{122}$}
\BIC{$(\forall B_{35}\,(\varphi_{97}) \to \varphi_{32} \to \varphi_{23}) \to \varphi_{122}$}
\BIC{$\varphi_{122}$}
\end{ptree}
\noindent$\Pi_{66}$:
\begin{ptree}
\AXC{$r \in \mathbb{R} \land r \neq 0 \land \neg \forall n \in \omega\,(n \in \mathbb{N} \to {{r}^{n}} \neq 0) \to \varphi_{32}$\lbl{A4}}
\AXC{$\stackrel{\Pi_{65}}{\vdots}$}\noLine
\UIC{$\varphi_{122}$}
\BIC{$r \in \mathbb{R} \to r \neq 0 \to \forall n \in \omega\,(n \in \mathbb{N} \to {{r}^{n}} \neq 0)$}
\AXC{$(r \in \mathbb{R} \to r \neq 0 \to \forall n \in \omega\,(n \in \mathbb{N} \to {{r}^{n}} \neq 0)) \to \varphi_{110}$\lbl{A1}}
\BIC{$\varphi_{110}$}
\end{ptree}
\noindent$\Pi_{67}$:
\begin{ptree}
\AXC{$\forall n \in \omega\,(n \in \mathbb{N} \to {{r}^{n}} \neq 0) \to n \in \omega \to n \in \mathbb{N} \to {{r}^{n}} \neq 0$\lbl{A2}}
\AXC{$\stackrel{\Pi_{66}}{\vdots}$}\noLine
\UIC{$\varphi_{110}$}
\BIC{$r \in \mathbb{R} \to r \neq 0 \to n \in \omega \to n \in \mathbb{N} \to {{r}^{n}} \neq 0$}
\AXC{$\varphi_{121}$\lbl{A1}}
\BIC{$(r \in \mathbb{R} \to r \neq 0 \to n \in \mathbb{N} \to n \in \omega) \to \varphi_{55}$}
\end{ptree}
\noindent$\Pi_{68}$:
\begin{ptree}
\AXC{$\forall n\,(n \in \mathbb{N} \leftrightarrow n \in \omega \land n \neq \emptyset)$\lbl{Def$\mathbb{N}$}}
\AXC{$\varphi_{43}$\lbl{A2}}
\BIC{$n \in \mathbb{N} \leftrightarrow n \in \omega \land n \neq \emptyset$}
\AXC{$\stackrel{\Pi_{1}}{\vdots}$}\noLine
\UIC{$\varphi_{48}$}
\BIC{$r \in \mathbb{R} \to r \neq 0 \to n \in \mathbb{N} \to n \in \omega$}
\AXC{$\stackrel{\Pi_{67}}{\vdots}$}\noLine
\UIC{$(r \in \mathbb{R} \to r \neq 0 \to n \in \mathbb{N} \to n \in \omega) \to \varphi_{55}$}
\BIC{$\varphi_{55}$}
\end{ptree}
\noindent$\Pi_{69}$:
\begin{ptree}
\AXC{$r \in \mathbb{R} \to r \neq 0 \to n \in \mathbb{N} \to n \in \mathbb{N}$\lbl{A1}}
\AXC{$\stackrel{\Pi_{68}}{\vdots}$}\noLine
\UIC{$\varphi_{55}$}
\BIC{$r \in \mathbb{R} \to r \neq 0 \to n \in \mathbb{N} \to {{r}^{n}} \neq 0$}
\AXC{$\varphi_{64}$\lbl{A1}}
\BIC{$r \in \mathbb{R} \land r \neq 0 \to n \in \mathbb{N} \to {{r}^{n}} \neq 0$}
\UIC{$\forall n\,(r \in \mathbb{R} \land r \neq 0 \to n \in \mathbb{N} \to {{r}^{n}} \neq 0)$}
\AXC{$\forall n\,(r \in \mathbb{R} \land r \neq 0 \to n \in \mathbb{N} \to {{r}^{n}} \neq 0) \to \varphi_{75} \to \varphi_{60}$\lbl{A1}}
\BIC{$\varphi_{75} \to \varphi_{60}$}
\end{ptree}
\noindent 主木:
\begin{ptree}
\AXC{$r \in \mathbb{R} \land r \neq 0 \to \forall n\,(r \in \mathbb{R} \land r \neq 0)$\lbl{A4}}
\AXC{$\varphi_{75}$\lbl{A3}}
\AXC{$\stackrel{\Pi_{69}}{\vdots}$}\noLine
\UIC{$\varphi_{75} \to \varphi_{60}$}
\BIC{$\varphi_{60}$}
\BIC{$r \in \mathbb{R} \to r \neq 0 \to \forall n \in \mathbb{N}\,({{r}^{n}} \neq 0)$}
\UIC{$\forall r \in \mathbb{R}\,(r \neq 0 \to \forall n \in \mathbb{N}\,({{r}^{n}} \neq 0))$}
\end{ptree}
\begin{abbrevlist}
\item $\varphi_{1} :\equiv \forall x\,\forall y\,(x = y \to y = x) \to \forall y\,(1_{\omega} = y \to y = 1_{\omega})$
\item $\varphi_{2} :\equiv \forall x\,\forall y\,(x = y \to y = x) \to \forall y\,(S(n) = y \to y = S(n))$
\item $\varphi_{3} :\equiv \forall y\,(1_{\omega} = y \to y = 1_{\omega}) \to 1_{\omega} = S(n) \to S(n) = 1_{\omega}$
\item $\varphi_{4} :\equiv \forall y\,(S(n) = y \to y = S(n)) \to S(n) = 1_{\omega} \to 1_{\omega} = S(n)$
\item $\varphi_{5} :\equiv \forall z\,\forall w\,(z = w \to z = 0 \to w = 0) \to \forall w\,(r = w \to r = 0 \to w = 0)$
\item $\varphi_{6} :\equiv S(n) = S(\emptyset) \to S(\emptyset) = 1_{\omega} \to S(n) = 1_{\omega}$
\item $\varphi_{7} :\equiv n \in B_{35} \to S(n) \in \mathbb{N} \to n \neq \emptyset \to \mathrm{Seq}({(r)^{\bullet}})$
\item $\varphi_{8} :\equiv \forall n \in B_{35}\,(S(n) \in B_{35}) \to y \in B_{35} \to S(y) \in B_{35}$
\item $\varphi_{9} :\equiv \forall y \in B_{35}\,(S(y) \in B_{35}) \to n \in B_{35} \to S(n) \in B_{35}$
\item $\varphi_{10} :\equiv \forall x\,\forall y\,\forall z\,(x = y \to y = z \to x = z) \to \forall y\,\forall z\,(S(n) = y \to y = z \to S(n) = z)$
\item $\varphi_{11} :\equiv (1_{\omega} = S(\emptyset) \to S(\emptyset) = 1_{\omega}) \to S(\emptyset) = 1_{\omega}$
\item $\varphi_{12} :\equiv \forall y\,(1_{\omega} = y \to y = 1_{\omega}) \to 1_{\omega} = S(\emptyset) \to S(\emptyset) = 1_{\omega}$
\item $\varphi_{13} :\equiv \forall n \in B_{35}\,(S(n) \in B_{35}) \to \forall y \in B_{35}\,(S(y) \in B_{35})$
\item $\varphi_{14} :\equiv \forall y \in B_{35}\,(S(y) \in B_{35}) \to \forall n \in B_{35}\,(S(n) \in B_{35})$
\item $\varphi_{15} :\equiv \mathrm{Ind}(B_{35}) \leftrightarrow \emptyset \in B_{35} \land \forall y \in B_{35}\,(S(y) \in B_{35})$
\item $\varphi_{16} :\equiv \forall x\,\forall y\,(x = y \to y = x) \to \forall y\,({{r}^{S(n)}} = y \to y = {{r}^{S(n)}})$
\item $\varphi_{17} :\equiv \forall x\,\forall y\,(x = y \to y = x) \to \forall y\,({{r}^{1_{\omega}}} = y \to y = {{r}^{1_{\omega}}})$
\item $\varphi_{18} :\equiv \forall n \in B_{35}\,(S(n) \in B_{35}) \to \forall y\,\forall n \in B_{35}\,(S(n) \in B_{35})$
\item $\varphi_{19} :\equiv \forall y \in B_{35}\,(S(y) \in B_{35}) \to \forall n\,\forall y \in B_{35}\,(S(y) \in B_{35})$
\item $\varphi_{20} :\equiv \forall y\,\forall n \in B_{35}\,(S(n) \in B_{35}) \to \forall y \in B_{35}\,(S(y) \in B_{35})$
\item $\varphi_{21} :\equiv \forall n\,\forall y \in B_{35}\,(S(y) \in B_{35}) \to \forall n \in B_{35}\,(S(n) \in B_{35})$
\item $\varphi_{22} :\equiv \forall x\,\forall y\,(x = y \to y = x) \to \forall y\,({{r}^{n}} \cdot r = y \to y = {{r}^{n}} \cdot r)$
\item $\varphi_{23} :\equiv \forall B_{35}\,(\neg \forall n\,(n \in B_{35} \leftrightarrow n \in \omega \land (n \in \mathbb{N} \to {{r}^{n}} \neq 0)))$
\item $\varphi_{24} :\equiv \forall w\,(r = w \to r = 0 \to w = 0) \to r = {{r}^{1_{\omega}}} \to r = 0 \to {{r}^{1_{\omega}}} = 0$
\item $\varphi_{25} :\equiv \forall A_{36}\,\exists B_{35}\,\forall n\,(n \in B_{35} \leftrightarrow n \in A_{36} \land (n \in \mathbb{N} \to {{r}^{n}} \neq 0))$
\item $\varphi_{26} :\equiv S(n) \in B_{35} \leftrightarrow S(n) \in \omega \land (S(n) \in \mathbb{N} \to {{r}^{S(n)}} \neq 0)$
\item $\varphi_{27} :\equiv \forall z\,\forall w\,(z = w \to S(n) = S(z) \to S(n) = S(w)) \to \forall w\,(n = w \to S(n) = S(n) \to S(n) = S(w))$
\item $\varphi_{28} :\equiv \forall n\,(n \in B_{35} \leftrightarrow n \in \omega \land (n \in \mathbb{N} \to {{r}^{n}} \neq 0)) \to \mathrm{Ind}(B_{35})$
\item $\varphi_{29} :\equiv r \neq 0 \to \exists B_{35}\,\forall n\,(n \in B_{35} \leftrightarrow n \in \omega \land (n \in \mathbb{N} \to {{r}^{n}} \neq 0))$
\item $\varphi_{30} :\equiv \forall z\,\forall w\,(z = w \to z = 0 \to w = 0) \to \forall w\,({{r}^{S(n)}} = w \to {{r}^{S(n)}} = 0 \to w = 0)$
\item $\varphi_{31} :\equiv \forall z\,\forall w\,(z = w \to z = 0 \to w = 0) \to \forall w\,({{r}^{1_{\omega}}} = w \to {{r}^{1_{\omega}}} = 0 \to w = 0)$
\item $\varphi_{32} :\equiv \forall B_{35}\,(r \in \mathbb{R} \land r \neq 0 \land \neg \forall n \in \omega\,(n \in \mathbb{N} \to {{r}^{n}} \neq 0))$
\item $\varphi_{33} :\equiv \forall z_{37}\,\forall A_{36}\,\exists B_{35}\,\forall n\,(n \in B_{35} \leftrightarrow n \in A_{36} \land (n \in \mathbb{N} \to {{z_{37}}^{n}} \neq 0))$
\item $\varphi_{34} :\equiv \forall n\,(n \in \omega \leftrightarrow \forall x\,(\mathrm{Ind}(x) \to n \in x)) \to (n \in \omega \leftrightarrow \forall x\,(\mathrm{Ind}(x) \to n \in x))$
\item $\varphi_{35} :\equiv \forall z\,\forall w\,(z = w \to z = 0 \to w = 0) \to \forall w\,({{r}^{n}} \cdot r = w \to {{r}^{n}} \cdot r = 0 \to w = 0)$
\item $\varphi_{36} :\equiv \forall y\,\forall z\,(S(n) = y \to y = z \to S(n) = z) \to \forall z\,(S(n) = S(\emptyset) \to S(\emptyset) = z \to S(n) = z)$
\item $\varphi_{37} :\equiv \forall w\,(n = w \to S(n) = S(n) \to S(n) = S(w)) \to n = \emptyset \to S(n) = S(n) \to S(n) = S(\emptyset)$
\item $\varphi_{38} :\equiv \forall n\,(n \in B_{35} \leftrightarrow n \in \omega \land (n \in \mathbb{N} \to {{r}^{n}} \neq 0)) \to \emptyset \in B_{35}$
\item $\varphi_{39} :\equiv r \in \mathbb{R} \land r \neq 0 \land \forall n\,(n \in B_{35} \leftrightarrow n \in \omega \land (n \in \mathbb{N} \to {{r}^{n}} \neq 0))$
\item $\varphi_{40} :\equiv ({{r}^{1_{\omega}}} = r \to r = {{r}^{1_{\omega}}}) \to r \in \mathbb{R} \to r \neq 0 \to r = {{r}^{1_{\omega}}}$
\item $\varphi_{41} :\equiv (r = {{r}^{1_{\omega}}} \to {{r}^{1_{\omega}}} = r) \to r \in \mathbb{R} \to r \neq 0 \to {{r}^{1_{\omega}}} = r$
\item $\varphi_{42} :\equiv (r \in \mathbb{R} \to r \neq 0 \to {{r}^{1_{\omega}}} = r) \to r \in \mathbb{R} \to r \neq 0 \to {{r}^{1_{\omega}}} \neq 0$
\item $\varphi_{43} :\equiv \forall n\,(n \in \mathbb{N} \leftrightarrow n \in \omega \land n \neq \emptyset) \to (n \in \mathbb{N} \leftrightarrow n \in \omega \land n \neq \emptyset)$
\item $\varphi_{44} :\equiv \forall y\,({{r}^{1_{\omega}}} = y \to y = {{r}^{1_{\omega}}}) \to {{r}^{1_{\omega}}} = r \to r = {{r}^{1_{\omega}}}$
\item $\varphi_{45} :\equiv \forall n\,(n \in B_{35} \leftrightarrow n \in \omega \land (n \in \mathbb{N} \to {{r}^{n}} \neq 0)) \to n \in B_{35} \to n \in B_{35}$
\item $\varphi_{46} :\equiv \mathrm{Seq}({(r)^{\bullet}}) \land {{r}^{1_{\omega}}} = r \land \forall n \in \mathbb{N}\,({{r}^{S(n)}} = {{r}^{n}} \cdot r)$
\item $\varphi_{47} :\equiv \forall n\,(n \in B_{35} \leftrightarrow n \in \omega \land (n \in \mathbb{N} \to {{r}^{n}} \neq 0)) \to n \in B_{35} \to n \in \omega$
\item $\varphi_{48} :\equiv (n \in \mathbb{N} \leftrightarrow n \in \omega \land n \neq \emptyset) \to r \in \mathbb{R} \to r \neq 0 \to n \in \mathbb{N} \to n \in \omega$
\item $\varphi_{49} :\equiv \forall n\,(n \in B_{35} \leftrightarrow n \in \omega \land (n \in \mathbb{N} \to {{r}^{n}} \neq 0)) \to n \in B_{35} \to S(n) \in B_{35}$
\item $\varphi_{50} :\equiv \forall n\,(n \in B_{35} \leftrightarrow n \in \omega \land (n \in \mathbb{N} \to {{r}^{n}} \neq 0)) \to n \in \omega \to n \in \omega$
\item $\varphi_{51} :\equiv \forall n\,(n \in B_{35} \leftrightarrow n \in \omega \land (n \in \mathbb{N} \to {{r}^{n}} \neq 0)) \to n \in B_{35} \to S(n) \in \omega$
\item $\varphi_{52} :\equiv n \in B_{35} \to S(n) \in \mathbb{N} \to n \neq \emptyset \to n \in \mathbb{N} \to {{r}^{S(n)}} = {{r}^{n}} \cdot r$
\item $\varphi_{53} :\equiv \forall n\,(n \in B_{35} \leftrightarrow n \in \omega \land (n \in \mathbb{N} \to {{r}^{n}} \neq 0)) \to \forall n \in B_{35}\,(S(n) \in B_{35})$
\item $\varphi_{54} :\equiv n \in B_{35} \to S(n) \in \mathbb{N} \to n \neq \emptyset \to \forall n \in \mathbb{N}\,({{r}^{S(n)}} = {{r}^{n}} \cdot r)$
\item $\varphi_{55} :\equiv (r \in \mathbb{R} \to r \neq 0 \to n \in \mathbb{N} \to n \in \mathbb{N}) \to r \in \mathbb{R} \to r \neq 0 \to n \in \mathbb{N} \to {{r}^{n}} \neq 0$
\item $\varphi_{56} :\equiv \forall z\,\forall w\,(z = w \to {{r}^{z}} = 0 \to {{r}^{w}} = 0) \to \forall w\,(S(n) = w \to {{r}^{S(n)}} = 0 \to {{r}^{w}} = 0)$
\item $\varphi_{57} :\equiv \forall z\,\forall w\,(z = w \to {{r}^{z}} = 0 \to {{r}^{w}} = 0) \to \forall w\,(1_{\omega} = w \to {{r}^{1_{\omega}}} = 0 \to {{r}^{w}} = 0)$
\item $\varphi_{58} :\equiv \forall w\,({{r}^{1_{\omega}}} = w \to {{r}^{1_{\omega}}} = 0 \to w = 0) \to {{r}^{1_{\omega}}} = r \to {{r}^{1_{\omega}}} = 0 \to r = 0$
\item $\varphi_{59} :\equiv \forall n\,(n \in B_{35} \leftrightarrow n \in \omega \land (n \in \mathbb{N} \to {{r}^{n}} \neq 0)) \to n \in \omega \to \forall x\,(\mathrm{Ind}(x) \to n \in x)$
\item $\varphi_{60} :\equiv (r \in \mathbb{R} \land r \neq 0 \to \forall n\,(r \in \mathbb{R} \land r \neq 0)) \to r \in \mathbb{R} \to r \neq 0 \to \forall n \in \mathbb{N}\,({{r}^{n}} \neq 0)$
\item $\varphi_{61} :\equiv \forall n\,(n \in B_{35} \leftrightarrow n \in \omega \land (n \in \mathbb{N} \to {{r}^{n}} \neq 0)) \to n \in \omega \to \mathrm{Ind}(B_{35}) \to n \in B_{35}$
\item $\varphi_{62} :\equiv \forall n \in \mathbb{N}\,({{r}^{S(n)}} = {{r}^{n}} \cdot r) \to n \in \mathbb{N} \to {{r}^{S(n)}} = {{r}^{n}} \cdot r$
\item $\varphi_{63} :\equiv \forall a\,\forall n\,(\mathrm{Seq}(a) \to n \in \mathbb{N} \to a_{n} \in \mathbb{R}) \to \forall n\,(\mathrm{Seq}({(r)^{\bullet}}) \to n \in \mathbb{N} \to {{r}^{n}} \in \mathbb{R})$
\item $\varphi_{64} :\equiv (r \in \mathbb{R} \to r \neq 0 \to n \in \mathbb{N} \to {{r}^{n}} \neq 0) \to r \in \mathbb{R} \land r \neq 0 \to n \in \mathbb{N} \to {{r}^{n}} \neq 0$
\item $\varphi_{65} :\equiv \forall n\,(\mathrm{Seq}({(r)^{\bullet}}) \to n \in \mathbb{N} \to {{r}^{n}} \in \mathbb{R}) \to \mathrm{Seq}({(r)^{\bullet}}) \to n \in \mathbb{N} \to {{r}^{n}} \in \mathbb{R}$
\item $\varphi_{66} :\equiv \forall w\,(S(n) = w \to {{r}^{S(n)}} = 0 \to {{r}^{w}} = 0) \to S(n) = 1_{\omega} \to {{r}^{S(n)}} = 0 \to {{r}^{1_{\omega}}} = 0$
\item $\varphi_{67} :\equiv \forall w\,(1_{\omega} = w \to {{r}^{1_{\omega}}} = 0 \to {{r}^{w}} = 0) \to 1_{\omega} = S(n) \to {{r}^{1_{\omega}}} = 0 \to {{r}^{S(n)}} = 0$
\item $\varphi_{68} :\equiv \forall n\,(n \in B_{35} \leftrightarrow n \in \omega \land (n \in \mathbb{N} \to {{r}^{n}} \neq 0)) \to n \in B_{35} \to n \in \mathbb{N} \to {{r}^{n}} \neq 0$
\item $\varphi_{69} :\equiv \forall n\,(n \in B_{35} \leftrightarrow n \in \omega \land (n \in \mathbb{N} \to {{r}^{n}} \neq 0)) \to n \in B_{35} \to n \in \omega \to S(n) \in \omega$
\item $\varphi_{70} :\equiv \forall n\,(n \in B_{35} \leftrightarrow n \in \omega \land (n \in \mathbb{N} \to {{r}^{n}} \neq 0)) \to n \in \omega \to n \in \mathbb{N} \to {{r}^{n}} \neq 0$
\item $\varphi_{71} :\equiv \forall n\,(n \in \mathbb{N} \leftrightarrow n \in \omega \land n \neq \emptyset) \to (\emptyset \in \mathbb{N} \leftrightarrow \emptyset \in \omega \land \emptyset \neq \emptyset)$
\item $\varphi_{72} :\equiv \forall n\,(n \in B_{35} \leftrightarrow n \in \omega \land (n \in \mathbb{N} \to {{r}^{n}} \neq 0)) \to n \in B_{35} \to \forall y \in \omega\,(S(y) \in \omega)$
\item $\varphi_{73} :\equiv \forall n\,(n \in B_{35} \leftrightarrow n \in \omega \land (n \in \mathbb{N} \to {{r}^{n}} \neq 0)) \to \forall n \in \omega\,(n \in \mathbb{N} \to {{r}^{n}} \neq 0)$
\item $\varphi_{74} :\equiv \forall n\,(n \in B_{35} \leftrightarrow n \in \omega \land (n \in \mathbb{N} \to {{r}^{n}} \neq 0)) \to n \in B_{35} \to S(n) \in \mathbb{N} \to {{r}^{S(n)}} \neq 0$
\item $\varphi_{75} :\equiv \forall n\,(r \in \mathbb{R} \land r \neq 0 \to n \in \mathbb{N} \to {{r}^{n}} \neq 0) \to \forall n\,(r \in \mathbb{R} \land r \neq 0) \to \forall n \in \mathbb{N}\,({{r}^{n}} \neq 0)$
\item $\varphi_{76} :\equiv \forall n\,(n \in B_{35} \leftrightarrow n \in \omega \land (n \in \mathbb{N} \to {{r}^{n}} \neq 0)) \to n \in B_{35} \to S(n) \in \mathbb{N} \to n = \emptyset \to S(n) = 1_{\omega}$
\item $\varphi_{77} :\equiv \forall n\,(n \in B_{35} \leftrightarrow n \in \omega \land (n \in \mathbb{N} \to {{r}^{n}} \neq 0)) \to n \in B_{35} \to S(n) \in \mathbb{N} \to n = \emptyset \to 1_{\omega} = S(n)$
\item $\varphi_{78} :\equiv \forall n\,(n \in B_{35} \leftrightarrow n \in \omega \land (n \in \mathbb{N} \to {{r}^{n}} \neq 0)) \to n \in B_{35} \to S(n) \in \mathbb{N} \to n \neq \emptyset \to n \in \mathbb{N}$
\item $\varphi_{79} :\equiv \forall n\,(n \in B_{35} \leftrightarrow n \in \omega \land (n \in \mathbb{N} \to {{r}^{n}} \neq 0)) \to n \in B_{35} \to S(n) \in \mathbb{N} \to n = \emptyset \to n = \emptyset$
\item $\varphi_{80} :\equiv \forall n\,(n \in B_{35} \leftrightarrow n \in \omega \land (n \in \mathbb{N} \to {{r}^{n}} \neq 0)) \to n \in B_{35} \to S(n) \in \mathbb{N} \to n \neq \emptyset \to r \in \mathbb{R}$
\item $\varphi_{81} :\equiv \forall n\,(n \in B_{35} \leftrightarrow n \in \omega \land (n \in \mathbb{N} \to {{r}^{n}} \neq 0)) \to n \in B_{35} \to S(n) \in \mathbb{N} \to n \neq \emptyset \to n \neq \emptyset$
\item $\varphi_{82} :\equiv \forall n\,(n \in B_{35} \leftrightarrow n \in \omega \land (n \in \mathbb{N} \to {{r}^{n}} \neq 0)) \to \emptyset \in \omega \land \forall y \in \omega\,(S(y) \in \omega)$
\item $\varphi_{83} :\equiv \forall n\,(n \in B_{35} \leftrightarrow n \in \omega \land (n \in \mathbb{N} \to {{r}^{n}} \neq 0)) \to (n \in B_{35} \leftrightarrow n \in \omega \land (n \in \mathbb{N} \to {{r}^{n}} \neq 0))$
\item $\varphi_{84} :\equiv \forall y\,({{r}^{S(n)}} = y \to y = {{r}^{S(n)}}) \to {{r}^{S(n)}} = {{r}^{n}} \cdot r \to {{r}^{n}} \cdot r = {{r}^{S(n)}}$
\item $\varphi_{85} :\equiv \forall n\,(n \in B_{35} \leftrightarrow n \in \omega \land (n \in \mathbb{N} \to {{r}^{n}} \neq 0)) \to n \in B_{35} \to n \in \omega \land (n \in \mathbb{N} \to {{r}^{n}} \neq 0)$
\item $\varphi_{86} :\equiv \forall n\,(n \in B_{35} \leftrightarrow n \in \omega \land (n \in \mathbb{N} \to {{r}^{n}} \neq 0)) \to \forall n\,(n \in B_{35} \leftrightarrow n \in \omega \land (n \in \mathbb{N} \to {{r}^{n}} \neq 0))$
\item $\varphi_{87} :\equiv \forall n\,(n \in B_{35} \leftrightarrow n \in \omega \land (n \in \mathbb{N} \to {{r}^{n}} \neq 0)) \to n \in B_{35} \to S(n) \in \mathbb{N} \to n \neq \emptyset \to {{r}^{n}} \neq 0$
\item $\varphi_{88} :\equiv \forall n\,(n \in B_{35} \leftrightarrow n \in \omega \land (n \in \mathbb{N} \to {{r}^{n}} \neq 0)) \to n \in B_{35} \to S(n) \in \mathbb{N} \to n = \emptyset \to S(n) = S(\emptyset)$
\item $\varphi_{89} :\equiv \forall y\,({{r}^{n}} \cdot r = y \to y = {{r}^{n}} \cdot r) \to {{r}^{n}} \cdot r = {{r}^{S(n)}} \to {{r}^{S(n)}} = {{r}^{n}} \cdot r$
\item $\varphi_{90} :\equiv r \neq 0 \to \forall n\,(n \in B_{35} \leftrightarrow n \in \omega \land (n \in \mathbb{N} \to {{r}^{n}} \neq 0)) \to \varphi_{7}$
\item $\varphi_{91} :\equiv \forall n\,(n \in B_{35} \leftrightarrow n \in \omega \land (n \in \mathbb{N} \to {{r}^{n}} \neq 0)) \to n \in B_{35} \to S(n) \in \mathbb{N} \to n = \emptyset \to {{r}^{S(n)}} \neq 0$
\item $\varphi_{92} :\equiv \forall n\,(n \in B_{35} \leftrightarrow n \in \omega \land (n \in \mathbb{N} \to {{r}^{n}} \neq 0)) \to n \in B_{35} \to S(n) \in \mathbb{N} \to n \neq \emptyset \to {{r}^{S(n)}} \neq 0$
\item $\varphi_{93} :\equiv \forall x\,(\mathrm{Ind}(x) \leftrightarrow \emptyset \in x \land \forall y \in x\,(S(y) \in x)) \to (\mathrm{Ind}(\omega) \leftrightarrow \emptyset \in \omega \land \forall y \in \omega\,(S(y) \in \omega))$
\item $\varphi_{94} :\equiv \forall n\,(n \in B_{35} \leftrightarrow n \in \omega \land (n \in \mathbb{N} \to {{r}^{n}} \neq 0)) \to n \in B_{35} \to S(n) \in \mathbb{N} \to n \neq \emptyset \to {{r}^{n}} \cdot r \neq 0$
\item $\varphi_{95} :\equiv \forall n\,(n \in B_{35} \leftrightarrow n \in \omega \land (n \in \mathbb{N} \to {{r}^{n}} \neq 0)) \to n \in B_{35} \to S(n) \in \mathbb{N} \to n \neq \emptyset \to {{r}^{n}} \in \mathbb{R}$
\item $\varphi_{96} :\equiv \forall w\,({{r}^{S(n)}} = w \to {{r}^{S(n)}} = 0 \to w = 0) \to {{r}^{S(n)}} = {{r}^{n}} \cdot r \to {{r}^{S(n)}} = 0 \to {{r}^{n}} \cdot r = 0$
\item $\varphi_{97} :\equiv r \in \mathbb{R} \land r \neq 0 \land \neg \forall n \in \omega\,(n \in \mathbb{N} \to {{r}^{n}} \neq 0) \to \neg \forall n\,(n \in B_{35} \leftrightarrow n \in \omega \land (n \in \mathbb{N} \to {{r}^{n}} \neq 0))$
\item $\varphi_{98} :\equiv \forall w\,({{r}^{n}} \cdot r = w \to {{r}^{n}} \cdot r = 0 \to w = 0) \to {{r}^{n}} \cdot r = {{r}^{S(n)}} \to {{r}^{n}} \cdot r = 0 \to {{r}^{S(n)}} = 0$
\item $\varphi_{99} :\equiv \varphi_{25} \to \exists B_{35}\,\forall n\,(n \in B_{35} \leftrightarrow n \in \omega \land (n \in \mathbb{N} \to {{r}^{n}} \neq 0))$
\item $\varphi_{100} :\equiv r \neq 0 \to \forall n\,(n \in B_{35} \leftrightarrow n \in \omega \land (n \in \mathbb{N} \to {{r}^{n}} \neq 0)) \to \varphi_{26}$
\item $\varphi_{101} :\equiv \forall n\,(n \in B_{35} \leftrightarrow n \in \omega \land (n \in \mathbb{N} \to {{r}^{n}} \neq 0)) \to n \in B_{35} \to S(n) \in \mathbb{N} \to n \neq \emptyset \to {{r}^{S(n)}} = {{r}^{n}} \cdot r$
\item $\varphi_{102} :\equiv \forall n\,(n \in B_{35} \leftrightarrow n \in \omega \land (n \in \mathbb{N} \to {{r}^{n}} \neq 0)) \to n \in B_{35} \to S(n) \in \mathbb{N} \to n \neq \emptyset \to {{r}^{n}} \cdot r = {{r}^{S(n)}}$
\item $\varphi_{103} :\equiv \forall n\,(n \in B_{35} \leftrightarrow n \in \omega \land (n \in \mathbb{N} \to {{r}^{n}} \neq 0)) \to (\emptyset \in B_{35} \leftrightarrow \emptyset \in \omega \land (\emptyset \in \mathbb{N} \to {{r}^{\emptyset}} \neq 0))$
\item $\varphi_{104} :\equiv (\emptyset \in \mathbb{N} \leftrightarrow \emptyset \in \omega \land \emptyset \neq \emptyset) \to \emptyset = \emptyset \to r \in \mathbb{R} \to r \neq 0 \to \emptyset \in \mathbb{N} \to {{r}^{\emptyset}} \neq 0$
\item $\varphi_{105} :\equiv \forall x\,\forall y\,(x \in \mathbb{R} \to y \in \mathbb{R} \to x \neq 0 \to y \neq 0 \to x \cdot y \neq 0) \to \forall y\,({{r}^{n}} \in \mathbb{R} \to y \in \mathbb{R} \to {{r}^{n}} \neq 0 \to y \neq 0 \to {{r}^{n}} \cdot y \neq 0)$
\item $\varphi_{106} :\equiv \exists B_{35}\,\forall n\,(n \in B_{35} \leftrightarrow n \in \omega \land (n \in \mathbb{N} \to {{r}^{n}} \neq 0)) \to r \in \mathbb{R} \to \varphi_{29}$
\item $\varphi_{107} :\equiv (n \in \omega \leftrightarrow \forall x\,(\mathrm{Ind}(x) \to n \in x)) \to r \in \mathbb{R} \to r \neq 0 \to \varphi_{59}$
\item $\varphi_{108} :\equiv (S(n) = 1_{\omega} \to 1_{\omega} = S(n)) \to r \in \mathbb{R} \to r \neq 0 \to \varphi_{77}$
\item $\varphi_{109} :\equiv (1_{\omega} = S(n) \to S(n) = 1_{\omega}) \to r \in \mathbb{R} \to r \neq 0 \to \varphi_{76}$
\item $\varphi_{110} :\equiv (\forall n \in \omega\,(n \in \mathbb{N} \to {{r}^{n}} \neq 0) \to n \in \omega \to n \in \mathbb{N} \to {{r}^{n}} \neq 0) \to r \in \mathbb{R} \to r \neq 0 \to n \in \omega \to n \in \mathbb{N} \to {{r}^{n}} \neq 0$
\item $\varphi_{111} :\equiv r \neq 0 \to \forall n\,(n \in B_{35} \leftrightarrow n \in \omega \land (n \in \mathbb{N} \to {{r}^{n}} \neq 0)) \to \varphi_{52}$
\item $\varphi_{112} :\equiv r \neq 0 \to \forall n\,(n \in B_{35} \leftrightarrow n \in \omega \land (n \in \mathbb{N} \to {{r}^{n}} \neq 0)) \to \varphi_{54}$
\item $\varphi_{113} :\equiv (\forall x\,(\mathrm{Ind}(x) \to n \in x) \to \mathrm{Ind}(B_{35}) \to n \in B_{35}) \to r \in \mathbb{R} \to r \neq 0 \to \varphi_{61}$
\item $\varphi_{114} :\equiv (r \in \mathbb{R} \to r \neq 0 \to \emptyset \in \mathbb{N} \to {{r}^{\emptyset}} \neq 0) \to r \in \mathbb{R} \to r \neq 0 \to \varphi_{38}$
\item $\varphi_{115} :\equiv \forall y\,({{r}^{n}} \in \mathbb{R} \to y \in \mathbb{R} \to {{r}^{n}} \neq 0 \to y \neq 0 \to {{r}^{n}} \cdot y \neq 0) \to {{r}^{n}} \in \mathbb{R} \to r \in \mathbb{R} \to {{r}^{n}} \neq 0 \to r \neq 0 \to {{r}^{n}} \cdot r \neq 0$
\item $\varphi_{116} :\equiv (r = {{r}^{1_{\omega}}} \to r = 0 \to {{r}^{1_{\omega}}} = 0) \to ({{r}^{1_{\omega}}} = r \to {{r}^{1_{\omega}}} = 0 \to r = 0) \to \varphi_{42}$
\item $\varphi_{117} :\equiv (n = \emptyset \to S(n) = S(n) \to S(n) = S(\emptyset)) \to r \in \mathbb{R} \to r \neq 0 \to \varphi_{88}$
\item $\varphi_{118} :\equiv \forall n\,(n \in B_{35} \leftrightarrow n \in \omega \land (n \in \mathbb{N} \to {{r}^{n}} \neq 0)) \to n \in B_{35} \to S(n) \in \mathbb{N} \to n \neq \emptyset \to \varphi_{46}$
\item $\varphi_{119} :\equiv (\forall y \in \omega\,(S(y) \in \omega) \to n \in \omega \to S(n) \in \omega) \to r \in \mathbb{R} \to r \neq 0 \to \varphi_{69}$
\item $\varphi_{120} :\equiv (\mathrm{Ind}(\omega) \leftrightarrow \emptyset \in \omega \land \forall y \in \omega\,(S(y) \in \omega)) \to r \in \mathbb{R} \to r \neq 0 \to \varphi_{82}$
\item $\varphi_{121} :\equiv (r \in \mathbb{R} \to r \neq 0 \to n \in \omega \to n \in \mathbb{N} \to {{r}^{n}} \neq 0) \to (r \in \mathbb{R} \to r \neq 0 \to n \in \mathbb{N} \to n \in \omega) \to \varphi_{55}$
\item $\varphi_{122} :\equiv (r \in \mathbb{R} \land r \neq 0 \land \neg \forall n \in \omega\,(n \in \mathbb{N} \to {{r}^{n}} \neq 0) \to \varphi_{32}) \to r \in \mathbb{R} \to r \neq 0 \to \forall n \in \omega\,(n \in \mathbb{N} \to {{r}^{n}} \neq 0)$
\item $\varphi_{123} :\equiv \forall n\,(\varphi_{39} \to n \in B_{35} \to S(n) \in B_{35}) \to \forall n\,(\varphi_{39}) \to \forall n \in B_{35}\,(S(n) \in B_{35})$
\item $\varphi_{124} :\equiv ({{r}^{S(n)}} = {{r}^{n}} \cdot r \to {{r}^{n}} \cdot r = {{r}^{S(n)}}) \to r \in \mathbb{R} \to r \neq 0 \to \varphi_{102}$
\item $\varphi_{125} :\equiv ({{r}^{n}} \cdot r = {{r}^{S(n)}} \to {{r}^{S(n)}} = {{r}^{n}} \cdot r) \to r \in \mathbb{R} \to r \neq 0 \to \varphi_{101}$
\item $\varphi_{126} :\equiv (r \in \mathbb{R} \to \varphi_{46}) \to (r \in \mathbb{R} \to r \neq 0 \to r \in \mathbb{R}) \to r \in \mathbb{R} \to r \neq 0 \to \varphi_{46}$
\item $\varphi_{127} :\equiv (r \in \mathbb{R} \to r \neq 0 \to r = {{r}^{1_{\omega}}}) \to (r \in \mathbb{R} \to r \neq 0 \to r \neq 0) \to \varphi_{116}$
\item $\varphi_{128} :\equiv \forall n\,(\varphi_{39} \to n \in \omega \to n \in \mathbb{N} \to {{r}^{n}} \neq 0) \to \forall n\,(\varphi_{39}) \to \forall n \in \omega\,(n \in \mathbb{N} \to {{r}^{n}} \neq 0)$
\item $\varphi_{129} :\equiv (\forall n\,(n \in B_{35} \leftrightarrow n \in \omega \land (n \in \mathbb{N} \to {{r}^{n}} \neq 0)) \to \varphi_{26}) \to r \in \mathbb{R} \to \varphi_{100}$
\item $\varphi_{130} :\equiv (r \in \mathbb{R} \to r \neq 0 \to \varphi_{87}) \to (r \in \mathbb{R} \to r \neq 0 \to r \neq 0) \to r \in \mathbb{R} \to r \neq 0 \to \varphi_{94}$
\item $\varphi_{131} :\equiv (S(n) = 1_{\omega} \to {{r}^{S(n)}} = 0 \to {{r}^{1_{\omega}}} = 0) \to (r \in \mathbb{R} \to r \neq 0 \to \varphi_{76}) \to r \in \mathbb{R} \to r \neq 0 \to \varphi_{91}$
\item $\varphi_{132} :\equiv (r \in \mathbb{R} \to r \neq 0 \to \varphi_{69}) \to (r \in \mathbb{R} \to r \neq 0 \to \varphi_{47}) \to r \in \mathbb{R} \to r \neq 0 \to \varphi_{51}$
\item $\varphi_{133} :\equiv (r \in \mathbb{R} \to r \neq 0 \to \varphi_{88}) \to S(\emptyset) = 1_{\omega} \to \varphi_{6} \to r \in \mathbb{R} \to r \neq 0 \to \varphi_{76}$
\item $\varphi_{134} :\equiv (r \in \mathbb{R} \to r \neq 0 \to \varphi_{51}) \to (r \in \mathbb{R} \to r \neq 0 \to \varphi_{74}) \to r \in \mathbb{R} \to r \neq 0 \to \varphi_{49}$
\item $\varphi_{135} :\equiv (r \in \mathbb{R} \to r \neq 0 \to \varphi_{28}) \to (r \in \mathbb{R} \to r \neq 0 \to \varphi_{83}) \to r \in \mathbb{R} \to r \neq 0 \to \varphi_{70}$
\item $\varphi_{136} :\equiv (r \in \mathbb{R} \to r \neq 0 \to \varphi_{47}) \to (r \in \mathbb{R} \to r \neq 0 \to \varphi_{81}) \to r \in \mathbb{R} \to r \neq 0 \to \varphi_{78}$
\item $\varphi_{137} :\equiv (r \in \mathbb{R} \to r \neq 0 \to \varphi_{83}) \to (r \in \mathbb{R} \to r \neq 0 \to \varphi_{45}) \to r \in \mathbb{R} \to r \neq 0 \to \varphi_{85}$
\item $\varphi_{138} :\equiv ({{r}^{S(n)}} = {{r}^{n}} \cdot r \to {{r}^{S(n)}} = 0 \to {{r}^{n}} \cdot r = 0) \to (r \in \mathbb{R} \to r \neq 0 \to \varphi_{101}) \to r \in \mathbb{R} \to r \neq 0 \to \varphi_{92}$
\item $\varphi_{139} :\equiv (r \in \mathbb{R} \to r \neq 0 \to \varphi_{68}) \to (r \in \mathbb{R} \to r \neq 0 \to \varphi_{78}) \to r \in \mathbb{R} \to r \neq 0 \to \varphi_{87}$
\item $\varphi_{140} :\equiv (r \in \mathbb{R} \to r \neq 0 \to \varphi_{91}) \to (r \in \mathbb{R} \to r \neq 0 \to \varphi_{92}) \to r \in \mathbb{R} \to r \neq 0 \to \varphi_{74}$
\item $\varphi_{141} :\equiv (r \in \mathbb{R} \to \varphi_{90}) \to (r \in \mathbb{R} \to r \neq 0 \to \varphi_{78}) \to r \in \mathbb{R} \to r \neq 0 \to \varphi_{95}$
\item $\varphi_{142} :\equiv (r \in \mathbb{R} \to r \neq 0 \to {{r}^{1_{\omega}}} \neq 0) \to (1_{\omega} = S(n) \to {{r}^{1_{\omega}}} = 0 \to {{r}^{S(n)}} = 0) \to \varphi_{131}$
\item $\varphi_{143} :\equiv (r \in \mathbb{R} \to r \neq 0 \to \varphi_{38}) \to (r \in \mathbb{R} \to r \neq 0 \to \varphi_{53}) \to \varphi_{14} \to \varphi_{13} \to r \in \mathbb{R} \to r \neq 0 \to \varphi_{28}$
\item $\varphi_{144} :\equiv (r \in \mathbb{R} \to \varphi_{46}) \to (r \in \mathbb{R} \to r \neq 0 \to \varphi_{80}) \to r \in \mathbb{R} \to r \neq 0 \to \varphi_{118}$
\item $\varphi_{145} :\equiv (r \in \mathbb{R} \to \varphi_{111}) \to (r \in \mathbb{R} \to r \neq 0 \to \varphi_{78}) \to r \in \mathbb{R} \to r \neq 0 \to \varphi_{101}$
\item $\varphi_{146} :\equiv (r \in \mathbb{R} \to r \neq 0 \to \varphi_{94}) \to ({{r}^{n}} \cdot r = {{r}^{S(n)}} \to {{r}^{n}} \cdot r = 0 \to {{r}^{S(n)}} = 0) \to \varphi_{138}$
\item $\varphi_{147} :\equiv \forall n\,(\varphi_{39} \to n \in B_{35} \to S(n) \in B_{35}) \to \varphi_{123} \to (\varphi_{39} \to \forall n\,(\varphi_{39})) \to r \in \mathbb{R} \to r \neq 0 \to \varphi_{53}$
\item $\varphi_{148} :\equiv \forall n\,(\varphi_{39} \to n \in \omega \to n \in \mathbb{N} \to {{r}^{n}} \neq 0) \to \varphi_{128} \to (\varphi_{39} \to \forall n\,(\varphi_{39})) \to r \in \mathbb{R} \to r \neq 0 \to \varphi_{73}$
\item $\varphi_{149} :\equiv ({{r}^{n}} \in \mathbb{R} \to r \in \mathbb{R} \to {{r}^{n}} \neq 0 \to r \neq 0 \to {{r}^{n}} \cdot r \neq 0) \to (r \in \mathbb{R} \to r \neq 0 \to \varphi_{95}) \to (r \in \mathbb{R} \to r \neq 0 \to \varphi_{80}) \to \varphi_{130}$
\item $\varphi_{150} :\equiv (r \in \mathbb{R} \to \varphi_{29}) \to \forall B_{35}\,(\varphi_{97}) \to (\forall B_{35}\,(\varphi_{97}) \to \varphi_{32} \to \varphi_{23}) \to \varphi_{122}$
\end{abbrevlist}
\end{fullproof}

% 自動生成 (prover/build.js): PowOneOne  ノード数 272
\begin{fullproof}{[P5]}{$\forall n \in \mathbb{N}\,({{1}^{n}} = 1)$}{272}\phantomsection\label{tree:PowOneOne}
\noindent$\Pi_{1}$:
\begin{ptree}
\AXC{$\varphi_{15}$\lbl{A4}}
\AXC{$\forall n\,(\varphi_{7}) \to \varphi_{17}$\lbl{A3}}
\AXC{$\varphi_{7}$\lbl{A2}}
\UIC{$\forall n\,(\varphi_{7})$}
\AXC{$\forall n\,(\varphi_{7}) \to (\forall n\,(\varphi_{7}) \to \varphi_{17}) \to \varphi_{15} \to \varphi_{12}$\lbl{A1}}
\BIC{$(\forall n\,(\varphi_{7}) \to \varphi_{17}) \to \varphi_{15} \to \varphi_{12}$}
\BIC{$\varphi_{15} \to \varphi_{12}$}
\BIC{$\varphi_{12}$}
\end{ptree}
\noindent$\Pi_{2}$:
\begin{ptree}
\AXC{$y = z \to x = y \to x = z$\lbl{A6}}
\AXC{$(y = z \to x = y \to x = z) \to x = y \to y = z \to x = z$\lbl{A1}}
\BIC{$x = y \to y = z \to x = z$}
\UIC{$\forall z\,(x = y \to y = z \to x = z)$}
\UIC{$\forall y\,\forall z\,(x = y \to y = z \to x = z)$}
\UIC{$\forall x\,\forall y\,\forall z\,(x = y \to y = z \to x = z)$}
\AXC{$\varphi_{27}$\lbl{A2}}
\BIC{$\forall y\,\forall z\,({{1}^{S(n)}} = y \to y = z \to {{1}^{S(n)}} = z)$}
\AXC{$\varphi_{72}$\lbl{A2}}
\BIC{$\forall z\,({{1}^{S(n)}} = {{1}^{n}} \to {{1}^{n}} = z \to {{1}^{S(n)}} = z)$}
\AXC{$\varphi_{81}$\lbl{A2}}
\BIC{${{1}^{S(n)}} = {{1}^{n}} \to {{1}^{n}} = 1 \to {{1}^{S(n)}} = 1$}
\end{ptree}
\noindent$\Pi_{3}$:
\begin{ptree}
\AXC{$\varphi_{29}$\lbl{A1}}
\AXC{$\varphi_{63}$\lbl{A2}}
\AXC{$\varphi_{63} \to \varphi_{29} \to \varphi_{65}$\lbl{A1}}
\BIC{$\varphi_{29} \to \varphi_{65}$}
\BIC{$\varphi_{65}$}
\AXC{$\varphi_{65} \to \varphi_{31}$\lbl{A1}}
\BIC{$\varphi_{31}$}
\end{ptree}
\noindent$\Pi_{4}$:
\begin{ptree}
\AXC{$\varphi_{64}$\lbl{A1}}
\AXC{$\stackrel{\Pi_{3}}{\vdots}$}\noLine
\UIC{$\varphi_{31}$}
\AXC{$\forall n\,(n \in \mathbb{N} \leftrightarrow n \in \omega \land n \neq \emptyset)$\lbl{Def$\mathbb{N}$}}
\AXC{$\varphi_{30}$\lbl{A2}}
\BIC{$n \in \mathbb{N} \leftrightarrow n \in \omega \land n \neq \emptyset$}
\AXC{$(n \in \mathbb{N} \leftrightarrow n \in \omega \land n \neq \emptyset) \to \varphi_{31} \to \varphi_{64} \to \varphi_{61}$\lbl{A1}}
\BIC{$\varphi_{31} \to \varphi_{64} \to \varphi_{61}$}
\BIC{$\varphi_{64} \to \varphi_{61}$}
\BIC{$\varphi_{61}$}
\end{ptree}
\noindent$\Pi_{5}$:
\begin{ptree}
\AXC{$\stackrel{\Pi_{4}}{\vdots}$}\noLine
\UIC{$\varphi_{61}$}
\AXC{$\varphi_{29}$\lbl{A1}}
\AXC{$\varphi_{63}$\lbl{A2}}
\AXC{$\varphi_{63} \to \varphi_{29} \to \varphi_{65}$\lbl{A1}}
\BIC{$\varphi_{29} \to \varphi_{65}$}
\BIC{$\varphi_{65}$}
\AXC{$\varphi_{65} \to \varphi_{45}$\lbl{A1}}
\BIC{$\varphi_{45}$}
\AXC{$\varphi_{45} \to \varphi_{61} \to \varphi_{67}$\lbl{A1}}
\BIC{$\varphi_{61} \to \varphi_{67}$}
\BIC{$\varphi_{67}$}
\end{ptree}
\noindent$\Pi_{6}$:
\begin{ptree}
\AXC{$y = z \to x = y \to x = z$\lbl{A6}}
\AXC{$(y = z \to x = y \to x = z) \to x = y \to y = z \to x = z$\lbl{A1}}
\BIC{$x = y \to y = z \to x = z$}
\UIC{$\forall z\,(x = y \to y = z \to x = z)$}
\UIC{$\forall y\,\forall z\,(x = y \to y = z \to x = z)$}
\UIC{$\forall x\,\forall y\,\forall z\,(x = y \to y = z \to x = z)$}
\AXC{$\varphi_{27}$\lbl{A2}}
\BIC{$\forall y\,\forall z\,({{1}^{S(n)}} = y \to y = z \to {{1}^{S(n)}} = z)$}
\AXC{$\varphi_{79}$\lbl{A2}}
\BIC{$\forall z\,({{1}^{S(n)}} = {{1}^{n}} \cdot 1 \to {{1}^{n}} \cdot 1 = z \to {{1}^{S(n)}} = z)$}
\AXC{$\varphi_{85}$\lbl{A2}}
\BIC{${{1}^{S(n)}} = {{1}^{n}} \cdot 1 \to {{1}^{n}} \cdot 1 = {{1}^{n}} \to {{1}^{S(n)}} = {{1}^{n}}$}
\end{ptree}
\noindent$\Pi_{7}$:
\begin{ptree}
\AXC{$1 \in \mathbb{R}$\lbl{(R1)}}
\AXC{$\varphi_{48}$\lbl{[S24]}}
\AXC{$\varphi_{48} \to 1 \in \mathbb{R} \to \varphi_{32}$\lbl{A2}}
\BIC{$1 \in \mathbb{R} \to \varphi_{32}$}
\AXC{$(1 \in \mathbb{R} \to \varphi_{32}) \to 1 \in \mathbb{R} \to \varphi_{86}$\lbl{A1}}
\BIC{$1 \in \mathbb{R} \to \varphi_{86}$}
\BIC{$\varphi_{86}$}
\end{ptree}
\noindent$\Pi_{8}$:
\begin{ptree}
\AXC{$\forall a\,\forall n\,(\mathrm{Seq}(a) \to n \in \mathbb{N} \to a_{n} \in \mathbb{R})$\lbl{[S1]}}
\AXC{$\varphi_{46}$\lbl{A2}}
\BIC{$\forall n\,(\mathrm{Seq}({(1)^{\bullet}}) \to n \in \mathbb{N} \to {{1}^{n}} \in \mathbb{R})$}
\AXC{$\varphi_{49}$\lbl{A2}}
\BIC{$\mathrm{Seq}({(1)^{\bullet}}) \to n \in \mathbb{N} \to {{1}^{n}} \in \mathbb{R}$}
\end{ptree}
\noindent$\Pi_{9}$:
\begin{ptree}
\AXC{$\stackrel{\Pi_{8}}{\vdots}$}\noLine
\UIC{$\mathrm{Seq}({(1)^{\bullet}}) \to n \in \mathbb{N} \to {{1}^{n}} \in \mathbb{R}$}
\AXC{$(\mathrm{Seq}({(1)^{\bullet}}) \to n \in \mathbb{N} \to {{1}^{n}} \in \mathbb{R}) \to \varphi_{90}$\lbl{A1}}
\BIC{$\varphi_{90}$}
\end{ptree}
\noindent$\Pi_{10}$:
\begin{ptree}
\AXC{$\forall x \in \mathbb{R}\,(x \cdot 1 = x)$\lbl{(R8)}}
\AXC{$\forall x \in \mathbb{R}\,(x \cdot 1 = x) \to {{1}^{n}} \in \mathbb{R} \to {{1}^{n}} \cdot 1 = {{1}^{n}}$\lbl{A2}}
\BIC{${{1}^{n}} \in \mathbb{R} \to {{1}^{n}} \cdot 1 = {{1}^{n}}$}
\AXC{$({{1}^{n}} \in \mathbb{R} \to {{1}^{n}} \cdot 1 = {{1}^{n}}) \to \varphi_{73} \to \varphi_{78}$\lbl{A1}}
\BIC{$\varphi_{73} \to \varphi_{78}$}
\end{ptree}
\noindent$\Pi_{11}$:
\begin{ptree}
\AXC{$\stackrel{\Pi_{4}}{\vdots}$}\noLine
\UIC{$\varphi_{61}$}
\AXC{$\stackrel{\Pi_{7}}{\vdots}$}\noLine
\UIC{$\varphi_{86}$}
\AXC{$\varphi_{86} \to \forall n\,(n \in B_{38} \leftrightarrow n \in \omega \land (n \in \mathbb{N} \to {{1}^{n}} = 1)) \to \varphi_{5}$\lbl{A1}}
\BIC{$\forall n\,(n \in B_{38} \leftrightarrow n \in \omega \land (n \in \mathbb{N} \to {{1}^{n}} = 1)) \to \varphi_{5}$}
\AXC{$\stackrel{\Pi_{9}}{\vdots}$}\noLine
\UIC{$\varphi_{90}$}
\BIC{$\varphi_{61} \to \varphi_{73}$}
\BIC{$\varphi_{73}$}
\AXC{$\stackrel{\Pi_{10}}{\vdots}$}\noLine
\UIC{$\varphi_{73} \to \varphi_{78}$}
\BIC{$\varphi_{78}$}
\end{ptree}
\noindent$\Pi_{12}$:
\begin{ptree}
\AXC{$1 \in \mathbb{R}$\lbl{(R1)}}
\AXC{$\varphi_{48}$\lbl{[S24]}}
\AXC{$\varphi_{48} \to 1 \in \mathbb{R} \to \varphi_{32}$\lbl{A2}}
\BIC{$1 \in \mathbb{R} \to \varphi_{32}$}
\AXC{$(1 \in \mathbb{R} \to \varphi_{32}) \to 1 \in \mathbb{R} \to \varphi_{86}$\lbl{A1}}
\BIC{$1 \in \mathbb{R} \to \varphi_{86}$}
\BIC{$\varphi_{86}$}
\end{ptree}
\noindent$\Pi_{13}$:
\begin{ptree}
\AXC{$\stackrel{\Pi_{4}}{\vdots}$}\noLine
\UIC{$\varphi_{61}$}
\AXC{$\varphi_{47}$\lbl{A2}}
\AXC{$\stackrel{\Pi_{12}}{\vdots}$}\noLine
\UIC{$\varphi_{86}$}
\AXC{$\varphi_{86} \to \forall n\,(n \in B_{38} \leftrightarrow n \in \omega \land (n \in \mathbb{N} \to {{1}^{n}} = 1)) \to \varphi_{40}$\lbl{A1}}
\BIC{$\forall n\,(n \in B_{38} \leftrightarrow n \in \omega \land (n \in \mathbb{N} \to {{1}^{n}} = 1)) \to \varphi_{40}$}
\AXC{$\varphi_{93}$\lbl{A1}}
\BIC{$\varphi_{47} \to \forall n\,(n \in B_{38} \leftrightarrow n \in \omega \land (n \in \mathbb{N} \to {{1}^{n}} = 1)) \to \varphi_{37}$}
\BIC{$\forall n\,(n \in B_{38} \leftrightarrow n \in \omega \land (n \in \mathbb{N} \to {{1}^{n}} = 1)) \to \varphi_{37}$}
\AXC{$\varphi_{92}$\lbl{A1}}
\BIC{$\varphi_{61} \to \varphi_{80}$}
\BIC{$\varphi_{80}$}
\end{ptree}
\noindent$\Pi_{14}$:
\begin{ptree}
\AXC{$\stackrel{\Pi_{6}}{\vdots}$}\noLine
\UIC{${{1}^{S(n)}} = {{1}^{n}} \cdot 1 \to {{1}^{n}} \cdot 1 = {{1}^{n}} \to {{1}^{S(n)}} = {{1}^{n}}$}
\AXC{$\stackrel{\Pi_{11}}{\vdots}$}\noLine
\UIC{$\varphi_{78}$}
\AXC{$\stackrel{\Pi_{13}}{\vdots}$}\noLine
\UIC{$\varphi_{80}$}
\AXC{$\varphi_{80} \to \varphi_{91}$\lbl{A1}}
\BIC{$\varphi_{91}$}
\BIC{$({{1}^{S(n)}} = {{1}^{n}} \cdot 1 \to {{1}^{n}} \cdot 1 = {{1}^{n}} \to {{1}^{S(n)}} = {{1}^{n}}) \to \varphi_{77}$}
\BIC{$\varphi_{77}$}
\end{ptree}
\noindent$\Pi_{15}$:
\begin{ptree}
\AXC{$\stackrel{\Pi_{2}}{\vdots}$}\noLine
\UIC{${{1}^{S(n)}} = {{1}^{n}} \to {{1}^{n}} = 1 \to {{1}^{S(n)}} = 1$}
\AXC{$\stackrel{\Pi_{5}}{\vdots}$}\noLine
\UIC{$\varphi_{67}$}
\AXC{$\stackrel{\Pi_{14}}{\vdots}$}\noLine
\UIC{$\varphi_{77}$}
\AXC{$\varphi_{77} \to \varphi_{67} \to ({{1}^{S(n)}} = {{1}^{n}} \to {{1}^{n}} = 1 \to {{1}^{S(n)}} = 1) \to \varphi_{71}$\lbl{A1}}
\BIC{$\varphi_{67} \to ({{1}^{S(n)}} = {{1}^{n}} \to {{1}^{n}} = 1 \to {{1}^{S(n)}} = 1) \to \varphi_{71}$}
\BIC{$({{1}^{S(n)}} = {{1}^{n}} \to {{1}^{n}} = 1 \to {{1}^{S(n)}} = 1) \to \varphi_{71}$}
\BIC{$\varphi_{71}$}
\end{ptree}
\noindent$\Pi_{16}$:
\begin{ptree}
\AXC{$1 \in \mathbb{R}$\lbl{(R1)}}
\AXC{$\varphi_{48}$\lbl{[S24]}}
\AXC{$\varphi_{48} \to 1 \in \mathbb{R} \to \varphi_{32}$\lbl{A2}}
\BIC{$1 \in \mathbb{R} \to \varphi_{32}$}
\AXC{$(1 \in \mathbb{R} \to \varphi_{32}) \to 1 \in \mathbb{R} \to \varphi_{32}$\lbl{A1}}
\BIC{$1 \in \mathbb{R} \to \varphi_{32}$}
\BIC{$\varphi_{32}$}
\AXC{$\varphi_{32} \to {{1}^{1_{\omega}}} = 1$\lbl{A1}}
\BIC{${{1}^{1_{\omega}}} = 1$}
\end{ptree}
\noindent$\Pi_{17}$:
\begin{ptree}
\AXC{$x = x$\lbl{A5}}
\AXC{$x = y \to x = x \to y = x$\lbl{A6}}
\AXC{$(x = y \to x = x \to y = x) \to x = x \to x = y \to y = x$\lbl{A1}}
\BIC{$x = x \to x = y \to y = x$}
\BIC{$x = y \to y = x$}
\UIC{$\forall y\,(x = y \to y = x)$}
\UIC{$\forall x\,\forall y\,(x = y \to y = x)$}
\end{ptree}
\noindent$\Pi_{18}$:
\begin{ptree}
\AXC{$y = z \to x = y \to x = z$\lbl{A6}}
\AXC{$(y = z \to x = y \to x = z) \to x = y \to y = z \to x = z$\lbl{A1}}
\BIC{$x = y \to y = z \to x = z$}
\UIC{$\forall z\,(x = y \to y = z \to x = z)$}
\UIC{$\forall y\,\forall z\,(x = y \to y = z \to x = z)$}
\UIC{$\forall x\,\forall y\,\forall z\,(x = y \to y = z \to x = z)$}
\AXC{$\varphi_{8}$\lbl{A2}}
\BIC{$\forall y\,\forall z\,(S(n) = y \to y = z \to S(n) = z)$}
\AXC{$\varphi_{24}$\lbl{A2}}
\BIC{$\forall z\,(S(n) = S(\emptyset) \to S(\emptyset) = z \to S(n) = z)$}
\end{ptree}
\noindent$\Pi_{19}$:
\begin{ptree}
\AXC{$\stackrel{\Pi_{18}}{\vdots}$}\noLine
\UIC{$\forall z\,(S(n) = S(\emptyset) \to S(\emptyset) = z \to S(n) = z)$}
\AXC{$\forall z\,(S(n) = S(\emptyset) \to S(\emptyset) = z \to S(n) = z) \to \varphi_{4}$\lbl{A2}}
\BIC{$\varphi_{4}$}
\end{ptree}
\noindent$\Pi_{20}$:
\begin{ptree}
\AXC{$\stackrel{\Pi_{17}}{\vdots}$}\noLine
\UIC{$\forall x\,\forall y\,(x = y \to y = x)$}
\AXC{$\varphi_{2}$\lbl{A2}}
\BIC{$\forall y\,(1_{\omega} = y \to y = 1_{\omega})$}
\AXC{$\varphi_{10}$\lbl{A2}}
\BIC{$1_{\omega} = S(\emptyset) \to S(\emptyset) = 1_{\omega}$}
\AXC{$1_{\omega} = S(\emptyset)$\lbl{Def$1_\omega$}}
\AXC{$1_{\omega} = S(\emptyset) \to \varphi_{9}$\lbl{A1}}
\BIC{$\varphi_{9}$}
\BIC{$S(\emptyset) = 1_{\omega}$}
\end{ptree}
\noindent$\Pi_{21}$:
\begin{ptree}
\AXC{$z = w \to S(n) = S(z) \to S(n) = S(w)$\lbl{A6}}
\UIC{$\forall w\,(z = w \to S(n) = S(z) \to S(n) = S(w))$}
\UIC{$\forall z\,\forall w\,(z = w \to S(n) = S(z) \to S(n) = S(w))$}
\AXC{$\varphi_{21}$\lbl{A2}}
\BIC{$\forall w\,(n = w \to S(n) = S(n) \to S(n) = S(w))$}
\AXC{$\varphi_{25}$\lbl{A2}}
\BIC{$n = \emptyset \to S(n) = S(n) \to S(n) = S(\emptyset)$}
\end{ptree}
\noindent$\Pi_{22}$:
\begin{ptree}
\AXC{$x = x$\lbl{A5}}
\UIC{$\forall x\,(x = x)$}
\AXC{$\forall x\,(x = x) \to S(n) = S(n)$\lbl{A2}}
\BIC{$S(n) = S(n)$}
\AXC{$\varphi_{62}$\lbl{A1}}
\AXC{$\varphi_{62} \to \varphi_{84}$\lbl{A1}}
\BIC{$\varphi_{84}$}
\BIC{$(n = \emptyset \to S(n) = S(n) \to S(n) = S(\emptyset)) \to \varphi_{68}$}
\end{ptree}
\noindent$\Pi_{23}$:
\begin{ptree}
\AXC{$\stackrel{\Pi_{21}}{\vdots}$}\noLine
\UIC{$n = \emptyset \to S(n) = S(n) \to S(n) = S(\emptyset)$}
\AXC{$\stackrel{\Pi_{22}}{\vdots}$}\noLine
\UIC{$(n = \emptyset \to S(n) = S(n) \to S(n) = S(\emptyset)) \to \varphi_{68}$}
\BIC{$\varphi_{68}$}
\end{ptree}
\noindent$\Pi_{24}$:
\begin{ptree}
\AXC{$\stackrel{\Pi_{19}}{\vdots}$}\noLine
\UIC{$\varphi_{4}$}
\AXC{$\stackrel{\Pi_{20}}{\vdots}$}\noLine
\UIC{$S(\emptyset) = 1_{\omega}$}
\AXC{$\stackrel{\Pi_{23}}{\vdots}$}\noLine
\UIC{$\varphi_{68}$}
\AXC{$\varphi_{68} \to S(\emptyset) = 1_{\omega} \to \varphi_{4} \to \varphi_{59}$\lbl{A1}}
\BIC{$S(\emptyset) = 1_{\omega} \to \varphi_{4} \to \varphi_{59}$}
\BIC{$\varphi_{4} \to \varphi_{59}$}
\BIC{$\varphi_{59}$}
\end{ptree}
\noindent$\Pi_{25}$:
\begin{ptree}
\AXC{$\stackrel{\Pi_{17}}{\vdots}$}\noLine
\UIC{$\forall x\,\forall y\,(x = y \to y = x)$}
\AXC{$\varphi_{1}$\lbl{A2}}
\BIC{$\forall y\,(S(n) = y \to y = S(n))$}
\AXC{$\varphi_{3}$\lbl{A2}}
\BIC{$S(n) = 1_{\omega} \to 1_{\omega} = S(n)$}
\AXC{$\stackrel{\Pi_{24}}{\vdots}$}\noLine
\UIC{$\varphi_{59}$}
\AXC{$\varphi_{59} \to (S(n) = 1_{\omega} \to 1_{\omega} = S(n)) \to \varphi_{60}$\lbl{A1}}
\BIC{$(S(n) = 1_{\omega} \to 1_{\omega} = S(n)) \to \varphi_{60}$}
\BIC{$\varphi_{60}$}
\end{ptree}
\noindent$\Pi_{26}$:
\begin{ptree}
\AXC{$z = w \to {{1}^{z}} = 1 \to {{1}^{w}} = 1$\lbl{A6}}
\UIC{$\forall w\,(z = w \to {{1}^{z}} = 1 \to {{1}^{w}} = 1)$}
\UIC{$\forall z\,\forall w\,(z = w \to {{1}^{z}} = 1 \to {{1}^{w}} = 1)$}
\AXC{$\varphi_{42}$\lbl{A2}}
\BIC{$\forall w\,(1_{\omega} = w \to {{1}^{1_{\omega}}} = 1 \to {{1}^{w}} = 1)$}
\AXC{$\varphi_{50}$\lbl{A2}}
\BIC{$1_{\omega} = S(n) \to {{1}^{1_{\omega}}} = 1 \to {{1}^{S(n)}} = 1$}
\AXC{$\stackrel{\Pi_{16}}{\vdots}$}\noLine
\UIC{${{1}^{1_{\omega}}} = 1$}
\AXC{$\stackrel{\Pi_{25}}{\vdots}$}\noLine
\UIC{$\varphi_{60}$}
\AXC{$\varphi_{60} \to \varphi_{87}$\lbl{A1}}
\BIC{$\varphi_{87}$}
\BIC{$(1_{\omega} = S(n) \to {{1}^{1_{\omega}}} = 1 \to {{1}^{S(n)}} = 1) \to \varphi_{69}$}
\BIC{$\varphi_{69}$}
\end{ptree}
\noindent$\Pi_{27}$:
\begin{ptree}
\AXC{$\forall x\,(\mathrm{Ind}(x) \leftrightarrow \emptyset \in x \land \forall y \in x\,(S(y) \in x))$\lbl{DefInd}}
\AXC{$\varphi_{74}$\lbl{A2}}
\BIC{$\mathrm{Ind}(\omega) \leftrightarrow \emptyset \in \omega \land \forall y \in \omega\,(S(y) \in \omega)$}
\AXC{$\mathrm{Ind}(\omega)$\lbl{補題(ωは帰納的)}}
\AXC{$\varphi_{88}$\lbl{A1}}
\BIC{$(\mathrm{Ind}(\omega) \leftrightarrow \emptyset \in \omega \land \forall y \in \omega\,(S(y) \in \omega)) \to \varphi_{66}$}
\BIC{$\varphi_{66}$}
\end{ptree}
\noindent$\Pi_{28}$:
\begin{ptree}
\AXC{$\stackrel{\Pi_{27}}{\vdots}$}\noLine
\UIC{$\varphi_{66}$}
\AXC{$\varphi_{66} \to \varphi_{53}$\lbl{A1}}
\BIC{$\varphi_{53}$}
\AXC{$\varphi_{53} \to (\forall y \in \omega\,(S(y) \in \omega) \to n \in \omega \to S(n) \in \omega) \to \varphi_{51}$\lbl{A1}}
\BIC{$(\forall y \in \omega\,(S(y) \in \omega) \to n \in \omega \to S(n) \in \omega) \to \varphi_{51}$}
\end{ptree}
\noindent$\Pi_{29}$:
\begin{ptree}
\AXC{$\forall y \in \omega\,(S(y) \in \omega) \to n \in \omega \to S(n) \in \omega$\lbl{A2}}
\AXC{$\stackrel{\Pi_{28}}{\vdots}$}\noLine
\UIC{$(\forall y \in \omega\,(S(y) \in \omega) \to n \in \omega \to S(n) \in \omega) \to \varphi_{51}$}
\BIC{$\varphi_{51}$}
\AXC{$\varphi_{51} \to \varphi_{31} \to \varphi_{35}$\lbl{A1}}
\BIC{$\varphi_{31} \to \varphi_{35}$}
\end{ptree}
\noindent$\Pi_{30}$:
\begin{ptree}
\AXC{$\stackrel{\Pi_{3}}{\vdots}$}\noLine
\UIC{$\varphi_{31}$}
\AXC{$\stackrel{\Pi_{29}}{\vdots}$}\noLine
\UIC{$\varphi_{31} \to \varphi_{35}$}
\BIC{$\varphi_{35}$}
\AXC{$\forall n\,(n \in B_{38} \leftrightarrow n \in \omega \land (n \in \mathbb{N} \to {{1}^{n}} = 1)) \to \varphi_{20}$\lbl{A2}}
\AXC{$\varphi_{94}$\lbl{A1}}
\BIC{$\varphi_{35} \to \varphi_{56} \to \varphi_{33}$}
\BIC{$\varphi_{56} \to \varphi_{33}$}
\end{ptree}
\noindent$\Pi_{31}$:
\begin{ptree}
\AXC{$\forall n\,(\varphi_{33}) \to \varphi_{39}$\lbl{A3}}
\AXC{$\stackrel{\Pi_{15}}{\vdots}$}\noLine
\UIC{$\varphi_{71}$}
\AXC{$\stackrel{\Pi_{26}}{\vdots}$}\noLine
\UIC{$\varphi_{69}$}
\AXC{$\varphi_{69} \to \varphi_{71} \to \varphi_{56}$\lbl{A1}}
\BIC{$\varphi_{71} \to \varphi_{56}$}
\BIC{$\varphi_{56}$}
\AXC{$\stackrel{\Pi_{30}}{\vdots}$}\noLine
\UIC{$\varphi_{56} \to \varphi_{33}$}
\BIC{$\varphi_{33}$}
\UIC{$\forall n\,(\varphi_{33})$}
\AXC{$\forall n\,(\varphi_{33}) \to (\forall n\,(\varphi_{33}) \to \varphi_{39}) \to \varphi_{70} \to \varphi_{36}$\lbl{A1}}
\BIC{$(\forall n\,(\varphi_{33}) \to \varphi_{39}) \to \varphi_{70} \to \varphi_{36}$}
\BIC{$\varphi_{70} \to \varphi_{36}$}
\end{ptree}
\noindent$\Pi_{32}$:
\begin{ptree}
\AXC{$x = x$\lbl{A5}}
\UIC{$\forall x\,(x = x)$}
\AXC{$\forall x\,(x = x) \to \emptyset = \emptyset$\lbl{A2}}
\BIC{$\emptyset = \emptyset$}
\AXC{$\forall n\,(n \in \mathbb{N} \leftrightarrow n \in \omega \land n \neq \emptyset)$\lbl{Def$\mathbb{N}$}}
\AXC{$\varphi_{54}$\lbl{A2}}
\BIC{$\emptyset \in \mathbb{N} \leftrightarrow \emptyset \in \omega \land \emptyset \neq \emptyset$}
\AXC{$\varphi_{75}$\lbl{A1}}
\BIC{$\emptyset = \emptyset \to \emptyset \in \mathbb{N} \to {{1}^{\emptyset}} = 1$}
\BIC{$\emptyset \in \mathbb{N} \to {{1}^{\emptyset}} = 1$}
\end{ptree}
\noindent$\Pi_{33}$:
\begin{ptree}
\AXC{$\stackrel{\Pi_{32}}{\vdots}$}\noLine
\UIC{$\emptyset \in \mathbb{N} \to {{1}^{\emptyset}} = 1$}
\AXC{$\stackrel{\Pi_{27}}{\vdots}$}\noLine
\UIC{$\varphi_{66}$}
\AXC{$\varphi_{83}$\lbl{A2}}
\AXC{$\varphi_{83} \to \varphi_{66} \to (\emptyset \in \mathbb{N} \to {{1}^{\emptyset}} = 1) \to \varphi_{26}$\lbl{A1}}
\BIC{$\varphi_{66} \to (\emptyset \in \mathbb{N} \to {{1}^{\emptyset}} = 1) \to \varphi_{26}$}
\BIC{$(\emptyset \in \mathbb{N} \to {{1}^{\emptyset}} = 1) \to \varphi_{26}$}
\BIC{$\varphi_{26}$}
\end{ptree}
\noindent$\Pi_{34}$:
\begin{ptree}
\AXC{$\forall x\,(\mathrm{Ind}(x) \leftrightarrow \emptyset \in x \land \forall y \in x\,(S(y) \in x))$\lbl{DefInd}}
\AXC{$\forall x\,(\mathrm{Ind}(x) \leftrightarrow \emptyset \in x \land \forall y \in x\,(S(y) \in x)) \to \varphi_{13}$\lbl{A2}}
\BIC{$\varphi_{13}$}
\end{ptree}
\noindent$\Pi_{35}$:
\begin{ptree}
\AXC{$\stackrel{\Pi_{1}}{\vdots}$}\noLine
\UIC{$\varphi_{12}$}
\AXC{$\varphi_{70}$\lbl{A4}}
\AXC{$\stackrel{\Pi_{31}}{\vdots}$}\noLine
\UIC{$\varphi_{70} \to \varphi_{36}$}
\BIC{$\varphi_{36}$}
\AXC{$\stackrel{\Pi_{33}}{\vdots}$}\noLine
\UIC{$\varphi_{26}$}
\AXC{$\stackrel{\Pi_{34}}{\vdots}$}\noLine
\UIC{$\varphi_{13}$}
\AXC{$\varphi_{13} \to \varphi_{26} \to \varphi_{36} \to \varphi_{12} \to \varphi_{11} \to \varphi_{22}$\lbl{A1}}
\BIC{$\varphi_{26} \to \varphi_{36} \to \varphi_{12} \to \varphi_{11} \to \varphi_{22}$}
\BIC{$\varphi_{36} \to \varphi_{12} \to \varphi_{11} \to \varphi_{22}$}
\BIC{$\varphi_{12} \to \varphi_{11} \to \varphi_{22}$}
\BIC{$\varphi_{11} \to \varphi_{22}$}
\end{ptree}
\noindent$\Pi_{36}$:
\begin{ptree}
\AXC{$\forall n\,(n \in \omega \leftrightarrow \forall x\,(\mathrm{Ind}(x) \to n \in x))$\lbl{Def$\omega$}}
\AXC{$\varphi_{23}$\lbl{A2}}
\BIC{$n \in \omega \leftrightarrow \forall x\,(\mathrm{Ind}(x) \to n \in x)$}
\AXC{$\varphi_{34}$\lbl{A1}}
\AXC{$\varphi_{34} \to (n \in \omega \leftrightarrow \forall x\,(\mathrm{Ind}(x) \to n \in x)) \to \varphi_{41}$\lbl{A1}}
\BIC{$(n \in \omega \leftrightarrow \forall x\,(\mathrm{Ind}(x) \to n \in x)) \to \varphi_{41}$}
\BIC{$\varphi_{41}$}
\end{ptree}
\noindent$\Pi_{37}$:
\begin{ptree}
\AXC{$\forall x\,(\mathrm{Ind}(x) \to n \in x) \to \mathrm{Ind}(B_{38}) \to n \in B_{38}$\lbl{A2}}
\AXC{$\stackrel{\Pi_{36}}{\vdots}$}\noLine
\UIC{$\varphi_{41}$}
\AXC{$\varphi_{41} \to \varphi_{82}$\lbl{A1}}
\BIC{$\varphi_{82}$}
\BIC{$\varphi_{43}$}
\AXC{$\varphi_{43} \to \varphi_{22} \to \varphi_{63} \to \varphi_{52}$\lbl{A1}}
\BIC{$\varphi_{22} \to \varphi_{63} \to \varphi_{52}$}
\end{ptree}
\noindent$\Pi_{38}$:
\begin{ptree}
\AXC{$\varphi_{14}$\lbl{A4}}
\AXC{$\forall y\,(\varphi_{6}) \to \varphi_{16}$\lbl{A3}}
\AXC{$\varphi_{6}$\lbl{A2}}
\UIC{$\forall y\,(\varphi_{6})$}
\AXC{$\forall y\,(\varphi_{6}) \to (\forall y\,(\varphi_{6}) \to \varphi_{16}) \to \varphi_{14} \to \varphi_{11}$\lbl{A1}}
\BIC{$(\forall y\,(\varphi_{6}) \to \varphi_{16}) \to \varphi_{14} \to \varphi_{11}$}
\BIC{$\varphi_{14} \to \varphi_{11}$}
\BIC{$\varphi_{11}$}
\AXC{$\stackrel{\Pi_{35}}{\vdots}$}\noLine
\UIC{$\varphi_{11} \to \varphi_{22}$}
\BIC{$\varphi_{22}$}
\AXC{$\stackrel{\Pi_{37}}{\vdots}$}\noLine
\UIC{$\varphi_{22} \to \varphi_{63} \to \varphi_{52}$}
\BIC{$\varphi_{63} \to \varphi_{52}$}
\end{ptree}
\noindent$\Pi_{39}$:
\begin{ptree}
\AXC{$\varphi_{70}$\lbl{A4}}
\AXC{$\forall n\,(\varphi_{52}) \to \varphi_{58}$\lbl{A3}}
\AXC{$\varphi_{63}$\lbl{A2}}
\AXC{$\stackrel{\Pi_{38}}{\vdots}$}\noLine
\UIC{$\varphi_{63} \to \varphi_{52}$}
\BIC{$\varphi_{52}$}
\UIC{$\forall n\,(\varphi_{52})$}
\AXC{$\forall n\,(\varphi_{52}) \to (\forall n\,(\varphi_{52}) \to \varphi_{58}) \to \varphi_{70} \to \varphi_{55}$\lbl{A1}}
\BIC{$(\forall n\,(\varphi_{52}) \to \varphi_{58}) \to \varphi_{70} \to \varphi_{55}$}
\BIC{$\varphi_{70} \to \varphi_{55}$}
\BIC{$\varphi_{55}$}
\AXC{$\varphi_{55} \to \varphi_{57}$\lbl{A1}}
\BIC{$\varphi_{57}$}
\UIC{$\forall B_{38}\,(\varphi_{57})$}
\end{ptree}
\noindent$\Pi_{40}$:
\begin{ptree}
\AXC{$\varphi_{44}$\lbl{A4}}
\AXC{$\varphi_{89}$\lbl{A3}}
\AXC{$\stackrel{\Pi_{39}}{\vdots}$}\noLine
\UIC{$\forall B_{38}\,(\varphi_{57})$}
\AXC{$\varphi_{19}$\lbl{Sep}}
\AXC{$\varphi_{76}$\lbl{A2}}
\BIC{$\exists B_{38}\,\forall n\,(n \in B_{38} \leftrightarrow n \in \omega \land (n \in \mathbb{N} \to {{1}^{n}} = 1))$}
\AXC{$\varphi_{95}$\lbl{A1}}
\BIC{$\forall B_{38}\,(\varphi_{57}) \to \varphi_{89} \to \varphi_{44} \to \forall n \in \omega\,(n \in \mathbb{N} \to {{1}^{n}} = 1)$}
\BIC{$\varphi_{89} \to \varphi_{44} \to \forall n \in \omega\,(n \in \mathbb{N} \to {{1}^{n}} = 1)$}
\BIC{$\varphi_{44} \to \forall n \in \omega\,(n \in \mathbb{N} \to {{1}^{n}} = 1)$}
\BIC{$\forall n \in \omega\,(n \in \mathbb{N} \to {{1}^{n}} = 1)$}
\end{ptree}
\noindent$\Pi_{41}$:
\begin{ptree}
\AXC{$\stackrel{\Pi_{40}}{\vdots}$}\noLine
\UIC{$\forall n \in \omega\,(n \in \mathbb{N} \to {{1}^{n}} = 1)$}
\AXC{$\forall n \in \omega\,(n \in \mathbb{N} \to {{1}^{n}} = 1) \to n \in \omega \to n \in \mathbb{N} \to {{1}^{n}} = 1$\lbl{A2}}
\BIC{$n \in \omega \to n \in \mathbb{N} \to {{1}^{n}} = 1$}
\AXC{$(n \in \omega \to n \in \mathbb{N} \to {{1}^{n}} = 1) \to \varphi_{28}$\lbl{A1}}
\BIC{$\varphi_{28}$}
\end{ptree}
\noindent$\Pi_{42}$:
\begin{ptree}
\AXC{$\forall n\,(n \in \mathbb{N} \leftrightarrow n \in \omega \land n \neq \emptyset)$\lbl{Def$\mathbb{N}$}}
\AXC{$\varphi_{30}$\lbl{A2}}
\BIC{$n \in \mathbb{N} \leftrightarrow n \in \omega \land n \neq \emptyset$}
\AXC{$n \in \mathbb{N} \to n \in \mathbb{N}$\lbl{A1}}
\AXC{$\varphi_{38}$\lbl{A1}}
\BIC{$(n \in \mathbb{N} \leftrightarrow n \in \omega \land n \neq \emptyset) \to n \in \mathbb{N} \to n \in \omega$}
\BIC{$n \in \mathbb{N} \to n \in \omega$}
\AXC{$\stackrel{\Pi_{41}}{\vdots}$}\noLine
\UIC{$\varphi_{28}$}
\BIC{$(n \in \mathbb{N} \to n \in \mathbb{N}) \to n \in \mathbb{N} \to {{1}^{n}} = 1$}
\end{ptree}
\noindent 主木:
\begin{ptree}
\AXC{$n \in \mathbb{N} \to n \in \mathbb{N}$\lbl{A1}}
\AXC{$\stackrel{\Pi_{42}}{\vdots}$}\noLine
\UIC{$(n \in \mathbb{N} \to n \in \mathbb{N}) \to n \in \mathbb{N} \to {{1}^{n}} = 1$}
\BIC{$n \in \mathbb{N} \to {{1}^{n}} = 1$}
\UIC{$\forall n \in \mathbb{N}\,({{1}^{n}} = 1)$}
\end{ptree}
\begin{abbrevlist}
\item $\varphi_{1} :\equiv \forall x\,\forall y\,(x = y \to y = x) \to \forall y\,(S(n) = y \to y = S(n))$
\item $\varphi_{2} :\equiv \forall x\,\forall y\,(x = y \to y = x) \to \forall y\,(1_{\omega} = y \to y = 1_{\omega})$
\item $\varphi_{3} :\equiv \forall y\,(S(n) = y \to y = S(n)) \to S(n) = 1_{\omega} \to 1_{\omega} = S(n)$
\item $\varphi_{4} :\equiv S(n) = S(\emptyset) \to S(\emptyset) = 1_{\omega} \to S(n) = 1_{\omega}$
\item $\varphi_{5} :\equiv n \in B_{38} \to S(n) \in \mathbb{N} \to n \neq \emptyset \to \mathrm{Seq}({(1)^{\bullet}})$
\item $\varphi_{6} :\equiv \forall n \in B_{38}\,(S(n) \in B_{38}) \to y \in B_{38} \to S(y) \in B_{38}$
\item $\varphi_{7} :\equiv \forall y \in B_{38}\,(S(y) \in B_{38}) \to n \in B_{38} \to S(n) \in B_{38}$
\item $\varphi_{8} :\equiv \forall x\,\forall y\,\forall z\,(x = y \to y = z \to x = z) \to \forall y\,\forall z\,(S(n) = y \to y = z \to S(n) = z)$
\item $\varphi_{9} :\equiv (1_{\omega} = S(\emptyset) \to S(\emptyset) = 1_{\omega}) \to S(\emptyset) = 1_{\omega}$
\item $\varphi_{10} :\equiv \forall y\,(1_{\omega} = y \to y = 1_{\omega}) \to 1_{\omega} = S(\emptyset) \to S(\emptyset) = 1_{\omega}$
\item $\varphi_{11} :\equiv \forall n \in B_{38}\,(S(n) \in B_{38}) \to \forall y \in B_{38}\,(S(y) \in B_{38})$
\item $\varphi_{12} :\equiv \forall y \in B_{38}\,(S(y) \in B_{38}) \to \forall n \in B_{38}\,(S(n) \in B_{38})$
\item $\varphi_{13} :\equiv \mathrm{Ind}(B_{38}) \leftrightarrow \emptyset \in B_{38} \land \forall y \in B_{38}\,(S(y) \in B_{38})$
\item $\varphi_{14} :\equiv \forall n \in B_{38}\,(S(n) \in B_{38}) \to \forall y\,\forall n \in B_{38}\,(S(n) \in B_{38})$
\item $\varphi_{15} :\equiv \forall y \in B_{38}\,(S(y) \in B_{38}) \to \forall n\,\forall y \in B_{38}\,(S(y) \in B_{38})$
\item $\varphi_{16} :\equiv \forall y\,\forall n \in B_{38}\,(S(n) \in B_{38}) \to \forall y \in B_{38}\,(S(y) \in B_{38})$
\item $\varphi_{17} :\equiv \forall n\,\forall y \in B_{38}\,(S(y) \in B_{38}) \to \forall n \in B_{38}\,(S(n) \in B_{38})$
\item $\varphi_{18} :\equiv \forall B_{38}\,(\neg \forall n\,(n \in B_{38} \leftrightarrow n \in \omega \land (n \in \mathbb{N} \to {{1}^{n}} = 1)))$
\item $\varphi_{19} :\equiv \forall A_{39}\,\exists B_{38}\,\forall n\,(n \in B_{38} \leftrightarrow n \in A_{39} \land (n \in \mathbb{N} \to {{1}^{n}} = 1))$
\item $\varphi_{20} :\equiv S(n) \in B_{38} \leftrightarrow S(n) \in \omega \land (S(n) \in \mathbb{N} \to {{1}^{S(n)}} = 1)$
\item $\varphi_{21} :\equiv \forall z\,\forall w\,(z = w \to S(n) = S(z) \to S(n) = S(w)) \to \forall w\,(n = w \to S(n) = S(n) \to S(n) = S(w))$
\item $\varphi_{22} :\equiv \forall n\,(n \in B_{38} \leftrightarrow n \in \omega \land (n \in \mathbb{N} \to {{1}^{n}} = 1)) \to \mathrm{Ind}(B_{38})$
\item $\varphi_{23} :\equiv \forall n\,(n \in \omega \leftrightarrow \forall x\,(\mathrm{Ind}(x) \to n \in x)) \to (n \in \omega \leftrightarrow \forall x\,(\mathrm{Ind}(x) \to n \in x))$
\item $\varphi_{24} :\equiv \forall y\,\forall z\,(S(n) = y \to y = z \to S(n) = z) \to \forall z\,(S(n) = S(\emptyset) \to S(\emptyset) = z \to S(n) = z)$
\item $\varphi_{25} :\equiv \forall w\,(n = w \to S(n) = S(n) \to S(n) = S(w)) \to n = \emptyset \to S(n) = S(n) \to S(n) = S(\emptyset)$
\item $\varphi_{26} :\equiv \forall n\,(n \in B_{38} \leftrightarrow n \in \omega \land (n \in \mathbb{N} \to {{1}^{n}} = 1)) \to \emptyset \in B_{38}$
\item $\varphi_{27} :\equiv \forall x\,\forall y\,\forall z\,(x = y \to y = z \to x = z) \to \forall y\,\forall z\,({{1}^{S(n)}} = y \to y = z \to {{1}^{S(n)}} = z)$
\item $\varphi_{28} :\equiv (n \in \mathbb{N} \to n \in \omega) \to (n \in \mathbb{N} \to n \in \mathbb{N}) \to n \in \mathbb{N} \to {{1}^{n}} = 1$
\item $\varphi_{29} :\equiv \forall n\,(n \in B_{38} \leftrightarrow n \in \omega \land (n \in \mathbb{N} \to {{1}^{n}} = 1)) \to n \in B_{38} \to n \in B_{38}$
\item $\varphi_{30} :\equiv \forall n\,(n \in \mathbb{N} \leftrightarrow n \in \omega \land n \neq \emptyset) \to (n \in \mathbb{N} \leftrightarrow n \in \omega \land n \neq \emptyset)$
\item $\varphi_{31} :\equiv \forall n\,(n \in B_{38} \leftrightarrow n \in \omega \land (n \in \mathbb{N} \to {{1}^{n}} = 1)) \to n \in B_{38} \to n \in \omega$
\item $\varphi_{32} :\equiv \mathrm{Seq}({(1)^{\bullet}}) \land {{1}^{1_{\omega}}} = 1 \land \forall n \in \mathbb{N}\,({{1}^{S(n)}} = {{1}^{n}} \cdot 1)$
\item $\varphi_{33} :\equiv \forall n\,(n \in B_{38} \leftrightarrow n \in \omega \land (n \in \mathbb{N} \to {{1}^{n}} = 1)) \to n \in B_{38} \to S(n) \in B_{38}$
\item $\varphi_{34} :\equiv \forall n\,(n \in B_{38} \leftrightarrow n \in \omega \land (n \in \mathbb{N} \to {{1}^{n}} = 1)) \to n \in \omega \to n \in \omega$
\item $\varphi_{35} :\equiv \forall n\,(n \in B_{38} \leftrightarrow n \in \omega \land (n \in \mathbb{N} \to {{1}^{n}} = 1)) \to n \in B_{38} \to S(n) \in \omega$
\item $\varphi_{36} :\equiv \forall n\,(n \in B_{38} \leftrightarrow n \in \omega \land (n \in \mathbb{N} \to {{1}^{n}} = 1)) \to \forall n \in B_{38}\,(S(n) \in B_{38})$
\item $\varphi_{37} :\equiv n \in B_{38} \to S(n) \in \mathbb{N} \to n \neq \emptyset \to n \in \mathbb{N} \to {{1}^{S(n)}} = {{1}^{n}} \cdot 1$
\item $\varphi_{38} :\equiv (n \in \mathbb{N} \to n \in \mathbb{N}) \to (n \in \mathbb{N} \leftrightarrow n \in \omega \land n \neq \emptyset) \to n \in \mathbb{N} \to n \in \omega$
\item $\varphi_{39} :\equiv \forall n\,\forall n\,(n \in B_{38} \leftrightarrow n \in \omega \land (n \in \mathbb{N} \to {{1}^{n}} = 1)) \to \forall n \in B_{38}\,(S(n) \in B_{38})$
\item $\varphi_{40} :\equiv n \in B_{38} \to S(n) \in \mathbb{N} \to n \neq \emptyset \to \forall n \in \mathbb{N}\,({{1}^{S(n)}} = {{1}^{n}} \cdot 1)$
\item $\varphi_{41} :\equiv \forall n\,(n \in B_{38} \leftrightarrow n \in \omega \land (n \in \mathbb{N} \to {{1}^{n}} = 1)) \to n \in \omega \to \forall x\,(\mathrm{Ind}(x) \to n \in x)$
\item $\varphi_{42} :\equiv \forall z\,\forall w\,(z = w \to {{1}^{z}} = 1 \to {{1}^{w}} = 1) \to \forall w\,(1_{\omega} = w \to {{1}^{1_{\omega}}} = 1 \to {{1}^{w}} = 1)$
\item $\varphi_{43} :\equiv \forall n\,(n \in B_{38} \leftrightarrow n \in \omega \land (n \in \mathbb{N} \to {{1}^{n}} = 1)) \to n \in \omega \to \mathrm{Ind}(B_{38}) \to n \in B_{38}$
\item $\varphi_{44} :\equiv \neg \forall n \in \omega\,(n \in \mathbb{N} \to {{1}^{n}} = 1) \to \forall B_{38}\,(\neg \forall n \in \omega\,(n \in \mathbb{N} \to {{1}^{n}} = 1))$
\item $\varphi_{45} :\equiv \forall n\,(n \in B_{38} \leftrightarrow n \in \omega \land (n \in \mathbb{N} \to {{1}^{n}} = 1)) \to n \in B_{38} \to n \in \mathbb{N} \to {{1}^{n}} = 1$
\item $\varphi_{46} :\equiv \forall a\,\forall n\,(\mathrm{Seq}(a) \to n \in \mathbb{N} \to a_{n} \in \mathbb{R}) \to \forall n\,(\mathrm{Seq}({(1)^{\bullet}}) \to n \in \mathbb{N} \to {{1}^{n}} \in \mathbb{R})$
\item $\varphi_{47} :\equiv \forall n \in \mathbb{N}\,({{1}^{S(n)}} = {{1}^{n}} \cdot 1) \to n \in \mathbb{N} \to {{1}^{S(n)}} = {{1}^{n}} \cdot 1$
\item $\varphi_{48} :\equiv \forall r \in \mathbb{R}\,(\mathrm{Seq}({(r)^{\bullet}}) \land {{r}^{1_{\omega}}} = r \land \forall n \in \mathbb{N}\,({{r}^{S(n)}} = {{r}^{n}} \cdot r))$
\item $\varphi_{49} :\equiv \forall n\,(\mathrm{Seq}({(1)^{\bullet}}) \to n \in \mathbb{N} \to {{1}^{n}} \in \mathbb{R}) \to \mathrm{Seq}({(1)^{\bullet}}) \to n \in \mathbb{N} \to {{1}^{n}} \in \mathbb{R}$
\item $\varphi_{50} :\equiv \forall w\,(1_{\omega} = w \to {{1}^{1_{\omega}}} = 1 \to {{1}^{w}} = 1) \to 1_{\omega} = S(n) \to {{1}^{1_{\omega}}} = 1 \to {{1}^{S(n)}} = 1$
\item $\varphi_{51} :\equiv \forall n\,(n \in B_{38} \leftrightarrow n \in \omega \land (n \in \mathbb{N} \to {{1}^{n}} = 1)) \to n \in B_{38} \to n \in \omega \to S(n) \in \omega$
\item $\varphi_{52} :\equiv \forall n\,(n \in B_{38} \leftrightarrow n \in \omega \land (n \in \mathbb{N} \to {{1}^{n}} = 1)) \to n \in \omega \to n \in \mathbb{N} \to {{1}^{n}} = 1$
\item $\varphi_{53} :\equiv \forall n\,(n \in B_{38} \leftrightarrow n \in \omega \land (n \in \mathbb{N} \to {{1}^{n}} = 1)) \to n \in B_{38} \to \forall y \in \omega\,(S(y) \in \omega)$
\item $\varphi_{54} :\equiv \forall n\,(n \in \mathbb{N} \leftrightarrow n \in \omega \land n \neq \emptyset) \to (\emptyset \in \mathbb{N} \leftrightarrow \emptyset \in \omega \land \emptyset \neq \emptyset)$
\item $\varphi_{55} :\equiv \forall n\,(n \in B_{38} \leftrightarrow n \in \omega \land (n \in \mathbb{N} \to {{1}^{n}} = 1)) \to \forall n \in \omega\,(n \in \mathbb{N} \to {{1}^{n}} = 1)$
\item $\varphi_{56} :\equiv \forall n\,(n \in B_{38} \leftrightarrow n \in \omega \land (n \in \mathbb{N} \to {{1}^{n}} = 1)) \to n \in B_{38} \to S(n) \in \mathbb{N} \to {{1}^{S(n)}} = 1$
\item $\varphi_{57} :\equiv \neg \forall n \in \omega\,(n \in \mathbb{N} \to {{1}^{n}} = 1) \to \neg \forall n\,(n \in B_{38} \leftrightarrow n \in \omega \land (n \in \mathbb{N} \to {{1}^{n}} = 1))$
\item $\varphi_{58} :\equiv \forall n\,\forall n\,(n \in B_{38} \leftrightarrow n \in \omega \land (n \in \mathbb{N} \to {{1}^{n}} = 1)) \to \forall n \in \omega\,(n \in \mathbb{N} \to {{1}^{n}} = 1)$
\item $\varphi_{59} :\equiv \forall n\,(n \in B_{38} \leftrightarrow n \in \omega \land (n \in \mathbb{N} \to {{1}^{n}} = 1)) \to n \in B_{38} \to S(n) \in \mathbb{N} \to n = \emptyset \to S(n) = 1_{\omega}$
\item $\varphi_{60} :\equiv \forall n\,(n \in B_{38} \leftrightarrow n \in \omega \land (n \in \mathbb{N} \to {{1}^{n}} = 1)) \to n \in B_{38} \to S(n) \in \mathbb{N} \to n = \emptyset \to 1_{\omega} = S(n)$
\item $\varphi_{61} :\equiv \forall n\,(n \in B_{38} \leftrightarrow n \in \omega \land (n \in \mathbb{N} \to {{1}^{n}} = 1)) \to n \in B_{38} \to S(n) \in \mathbb{N} \to n \neq \emptyset \to n \in \mathbb{N}$
\item $\varphi_{62} :\equiv \forall n\,(n \in B_{38} \leftrightarrow n \in \omega \land (n \in \mathbb{N} \to {{1}^{n}} = 1)) \to n \in B_{38} \to S(n) \in \mathbb{N} \to n = \emptyset \to n = \emptyset$
\item $\varphi_{63} :\equiv \forall n\,(n \in B_{38} \leftrightarrow n \in \omega \land (n \in \mathbb{N} \to {{1}^{n}} = 1)) \to (n \in B_{38} \leftrightarrow n \in \omega \land (n \in \mathbb{N} \to {{1}^{n}} = 1))$
\item $\varphi_{64} :\equiv \forall n\,(n \in B_{38} \leftrightarrow n \in \omega \land (n \in \mathbb{N} \to {{1}^{n}} = 1)) \to n \in B_{38} \to S(n) \in \mathbb{N} \to n \neq \emptyset \to n \neq \emptyset$
\item $\varphi_{65} :\equiv \forall n\,(n \in B_{38} \leftrightarrow n \in \omega \land (n \in \mathbb{N} \to {{1}^{n}} = 1)) \to n \in B_{38} \to n \in \omega \land (n \in \mathbb{N} \to {{1}^{n}} = 1)$
\item $\varphi_{66} :\equiv \forall n\,(n \in B_{38} \leftrightarrow n \in \omega \land (n \in \mathbb{N} \to {{1}^{n}} = 1)) \to \emptyset \in \omega \land \forall y \in \omega\,(S(y) \in \omega)$
\item $\varphi_{67} :\equiv \forall n\,(n \in B_{38} \leftrightarrow n \in \omega \land (n \in \mathbb{N} \to {{1}^{n}} = 1)) \to n \in B_{38} \to S(n) \in \mathbb{N} \to n \neq \emptyset \to {{1}^{n}} = 1$
\item $\varphi_{68} :\equiv \forall n\,(n \in B_{38} \leftrightarrow n \in \omega \land (n \in \mathbb{N} \to {{1}^{n}} = 1)) \to n \in B_{38} \to S(n) \in \mathbb{N} \to n = \emptyset \to S(n) = S(\emptyset)$
\item $\varphi_{69} :\equiv \forall n\,(n \in B_{38} \leftrightarrow n \in \omega \land (n \in \mathbb{N} \to {{1}^{n}} = 1)) \to n \in B_{38} \to S(n) \in \mathbb{N} \to n = \emptyset \to {{1}^{S(n)}} = 1$
\item $\varphi_{70} :\equiv \forall n\,(n \in B_{38} \leftrightarrow n \in \omega \land (n \in \mathbb{N} \to {{1}^{n}} = 1)) \to \forall n\,\forall n\,(n \in B_{38} \leftrightarrow n \in \omega \land (n \in \mathbb{N} \to {{1}^{n}} = 1))$
\item $\varphi_{71} :\equiv \forall n\,(n \in B_{38} \leftrightarrow n \in \omega \land (n \in \mathbb{N} \to {{1}^{n}} = 1)) \to n \in B_{38} \to S(n) \in \mathbb{N} \to n \neq \emptyset \to {{1}^{S(n)}} = 1$
\item $\varphi_{72} :\equiv \forall y\,\forall z\,({{1}^{S(n)}} = y \to y = z \to {{1}^{S(n)}} = z) \to \forall z\,({{1}^{S(n)}} = {{1}^{n}} \to {{1}^{n}} = z \to {{1}^{S(n)}} = z)$
\item $\varphi_{73} :\equiv \forall n\,(n \in B_{38} \leftrightarrow n \in \omega \land (n \in \mathbb{N} \to {{1}^{n}} = 1)) \to n \in B_{38} \to S(n) \in \mathbb{N} \to n \neq \emptyset \to {{1}^{n}} \in \mathbb{R}$
\item $\varphi_{74} :\equiv \forall x\,(\mathrm{Ind}(x) \leftrightarrow \emptyset \in x \land \forall y \in x\,(S(y) \in x)) \to (\mathrm{Ind}(\omega) \leftrightarrow \emptyset \in \omega \land \forall y \in \omega\,(S(y) \in \omega))$
\item $\varphi_{75} :\equiv (\emptyset \in \mathbb{N} \leftrightarrow \emptyset \in \omega \land \emptyset \neq \emptyset) \to \emptyset = \emptyset \to \emptyset \in \mathbb{N} \to {{1}^{\emptyset}} = 1$
\item $\varphi_{76} :\equiv \varphi_{19} \to \exists B_{38}\,\forall n\,(n \in B_{38} \leftrightarrow n \in \omega \land (n \in \mathbb{N} \to {{1}^{n}} = 1))$
\item $\varphi_{77} :\equiv \forall n\,(n \in B_{38} \leftrightarrow n \in \omega \land (n \in \mathbb{N} \to {{1}^{n}} = 1)) \to n \in B_{38} \to S(n) \in \mathbb{N} \to n \neq \emptyset \to {{1}^{S(n)}} = {{1}^{n}}$
\item $\varphi_{78} :\equiv \forall n\,(n \in B_{38} \leftrightarrow n \in \omega \land (n \in \mathbb{N} \to {{1}^{n}} = 1)) \to n \in B_{38} \to S(n) \in \mathbb{N} \to n \neq \emptyset \to {{1}^{n}} \cdot 1 = {{1}^{n}}$
\item $\varphi_{79} :\equiv \forall y\,\forall z\,({{1}^{S(n)}} = y \to y = z \to {{1}^{S(n)}} = z) \to \forall z\,({{1}^{S(n)}} = {{1}^{n}} \cdot 1 \to {{1}^{n}} \cdot 1 = z \to {{1}^{S(n)}} = z)$
\item $\varphi_{80} :\equiv \forall n\,(n \in B_{38} \leftrightarrow n \in \omega \land (n \in \mathbb{N} \to {{1}^{n}} = 1)) \to n \in B_{38} \to S(n) \in \mathbb{N} \to n \neq \emptyset \to {{1}^{S(n)}} = {{1}^{n}} \cdot 1$
\item $\varphi_{81} :\equiv \forall z\,({{1}^{S(n)}} = {{1}^{n}} \to {{1}^{n}} = z \to {{1}^{S(n)}} = z) \to {{1}^{S(n)}} = {{1}^{n}} \to {{1}^{n}} = 1 \to {{1}^{S(n)}} = 1$
\item $\varphi_{82} :\equiv (\forall x\,(\mathrm{Ind}(x) \to n \in x) \to \mathrm{Ind}(B_{38}) \to n \in B_{38}) \to \varphi_{43}$
\item $\varphi_{83} :\equiv \forall n\,(n \in B_{38} \leftrightarrow n \in \omega \land (n \in \mathbb{N} \to {{1}^{n}} = 1)) \to (\emptyset \in B_{38} \leftrightarrow \emptyset \in \omega \land (\emptyset \in \mathbb{N} \to {{1}^{\emptyset}} = 1))$
\item $\varphi_{84} :\equiv S(n) = S(n) \to (n = \emptyset \to S(n) = S(n) \to S(n) = S(\emptyset)) \to \varphi_{68}$
\item $\varphi_{85} :\equiv \forall z\,({{1}^{S(n)}} = {{1}^{n}} \cdot 1 \to {{1}^{n}} \cdot 1 = z \to {{1}^{S(n)}} = z) \to {{1}^{S(n)}} = {{1}^{n}} \cdot 1 \to {{1}^{n}} \cdot 1 = {{1}^{n}} \to {{1}^{S(n)}} = {{1}^{n}}$
\item $\varphi_{86} :\equiv \forall n\,(n \in B_{38} \leftrightarrow n \in \omega \land (n \in \mathbb{N} \to {{1}^{n}} = 1)) \to n \in B_{38} \to S(n) \in \mathbb{N} \to n \neq \emptyset \to \varphi_{32}$
\item $\varphi_{87} :\equiv {{1}^{1_{\omega}}} = 1 \to (1_{\omega} = S(n) \to {{1}^{1_{\omega}}} = 1 \to {{1}^{S(n)}} = 1) \to \varphi_{69}$
\item $\varphi_{88} :\equiv \mathrm{Ind}(\omega) \to (\mathrm{Ind}(\omega) \leftrightarrow \emptyset \in \omega \land \forall y \in \omega\,(S(y) \in \omega)) \to \varphi_{66}$
\item $\varphi_{89} :\equiv \forall B_{38}\,(\varphi_{57}) \to \forall B_{38}\,(\neg \forall n \in \omega\,(n \in \mathbb{N} \to {{1}^{n}} = 1)) \to \varphi_{18}$
\item $\varphi_{90} :\equiv (\forall n\,(n \in B_{38} \leftrightarrow n \in \omega \land (n \in \mathbb{N} \to {{1}^{n}} = 1)) \to \varphi_{5}) \to \varphi_{61} \to \varphi_{73}$
\item $\varphi_{91} :\equiv \varphi_{78} \to ({{1}^{S(n)}} = {{1}^{n}} \cdot 1 \to {{1}^{n}} \cdot 1 = {{1}^{n}} \to {{1}^{S(n)}} = {{1}^{n}}) \to \varphi_{77}$
\item $\varphi_{92} :\equiv (\forall n\,(n \in B_{38} \leftrightarrow n \in \omega \land (n \in \mathbb{N} \to {{1}^{n}} = 1)) \to \varphi_{37}) \to \varphi_{61} \to \varphi_{80}$
\item $\varphi_{93} :\equiv (\forall n\,(n \in B_{38} \leftrightarrow n \in \omega \land (n \in \mathbb{N} \to {{1}^{n}} = 1)) \to \varphi_{40}) \to \varphi_{47} \to \forall n\,(n \in B_{38} \leftrightarrow n \in \omega \land (n \in \mathbb{N} \to {{1}^{n}} = 1)) \to \varphi_{37}$
\item $\varphi_{94} :\equiv (\forall n\,(n \in B_{38} \leftrightarrow n \in \omega \land (n \in \mathbb{N} \to {{1}^{n}} = 1)) \to \varphi_{20}) \to \varphi_{35} \to \varphi_{56} \to \varphi_{33}$
\item $\varphi_{95} :\equiv \exists B_{38}\,\forall n\,(n \in B_{38} \leftrightarrow n \in \omega \land (n \in \mathbb{N} \to {{1}^{n}} = 1)) \to \forall B_{38}\,(\varphi_{57}) \to \varphi_{89} \to \varphi_{44} \to \forall n \in \omega\,(n \in \mathbb{N} \to {{1}^{n}} = 1)$
\end{abbrevlist}
\end{fullproof}

% 自動生成 (prover/build.js): PowInv  ノード数 441
\begin{fullproof}{[P6]}{$\forall r \in \mathbb{R}\,(r \neq 0 \to \forall n \in \mathbb{N}\,({{(r^{-1})}^{n}} = {{r}^{n}}^{-1}))$}{441}\phantomsection\label{tree:PowInv}
\noindent$\Pi_{1}$:
\begin{ptree}
\AXC{$r \in \mathbb{R} \to r \neq 0 \to n \in \mathbb{N} \to n \in \mathbb{N}$\lbl{A1}}
\AXC{$(r \in \mathbb{R} \to r \neq 0 \to n \in \mathbb{N} \to n \in \mathbb{N}) \to \varphi_{43}$\lbl{A1}}
\BIC{$\varphi_{43}$}
\end{ptree}
\noindent$\Pi_{2}$:
\begin{ptree}
\AXC{$\varphi_{103}$\lbl{A2}}
\AXC{$r \in \mathbb{R} \to r \neq 0 \to \varphi_{32} \to \varphi_{32}$\lbl{A1}}
\AXC{$(r \in \mathbb{R} \to r \neq 0 \to \varphi_{32} \to \varphi_{32}) \to \varphi_{103} \to r \in \mathbb{R} \to r \neq 0 \to \varphi_{103}$\lbl{A1}}
\BIC{$\varphi_{103} \to r \in \mathbb{R} \to r \neq 0 \to \varphi_{103}$}
\BIC{$r \in \mathbb{R} \to r \neq 0 \to \varphi_{103}$}
\end{ptree}
\noindent$\Pi_{3}$:
\begin{ptree}
\AXC{$\varphi_{18}$\lbl{A4}}
\AXC{$\forall y\,(\varphi_{6}) \to \varphi_{20}$\lbl{A3}}
\AXC{$\varphi_{6}$\lbl{A2}}
\UIC{$\forall y\,(\varphi_{6})$}
\AXC{$\forall y\,(\varphi_{6}) \to (\forall y\,(\varphi_{6}) \to \varphi_{20}) \to \varphi_{18} \to \varphi_{11}$\lbl{A1}}
\BIC{$(\forall y\,(\varphi_{6}) \to \varphi_{20}) \to \varphi_{18} \to \varphi_{11}$}
\BIC{$\varphi_{18} \to \varphi_{11}$}
\BIC{$\varphi_{11}$}
\end{ptree}
\noindent$\Pi_{4}$:
\begin{ptree}
\AXC{$\varphi_{19}$\lbl{A4}}
\AXC{$\forall n\,(\varphi_{7}) \to \varphi_{21}$\lbl{A3}}
\AXC{$\varphi_{7}$\lbl{A2}}
\UIC{$\forall n\,(\varphi_{7})$}
\AXC{$\forall n\,(\varphi_{7}) \to (\forall n\,(\varphi_{7}) \to \varphi_{21}) \to \varphi_{19} \to \varphi_{12}$\lbl{A1}}
\BIC{$(\forall n\,(\varphi_{7}) \to \varphi_{21}) \to \varphi_{19} \to \varphi_{12}$}
\BIC{$\varphi_{19} \to \varphi_{12}$}
\BIC{$\varphi_{12}$}
\end{ptree}
\noindent$\Pi_{5}$:
\begin{ptree}
\AXC{$y = z \to x = y \to x = z$\lbl{A6}}
\AXC{$(y = z \to x = y \to x = z) \to x = y \to y = z \to x = z$\lbl{A1}}
\BIC{$x = y \to y = z \to x = z$}
\UIC{$\forall z\,(x = y \to y = z \to x = z)$}
\UIC{$\forall y\,\forall z\,(x = y \to y = z \to x = z)$}
\UIC{$\forall x\,\forall y\,\forall z\,(x = y \to y = z \to x = z)$}
\AXC{$\forall x\,\forall y\,\forall z\,(x = y \to y = z \to x = z) \to \varphi_{16}$\lbl{A2}}
\BIC{$\varphi_{16}$}
\end{ptree}
\noindent$\Pi_{6}$:
\begin{ptree}
\AXC{$z = w \to ({{r}^{n}} \cdot r)^{-1} = z^{-1} \to ({{r}^{n}} \cdot r)^{-1} = w^{-1}$\lbl{A6}}
\UIC{$\varphi_{24}$}
\UIC{$\forall z\,(\varphi_{24})$}
\AXC{$\forall z\,(\varphi_{24}) \to \varphi_{72}$\lbl{A2}}
\BIC{$\varphi_{72}$}
\AXC{$\varphi_{72} \to {{r}^{n}} \cdot r = {{r}^{S(n)}} \to \varphi_{60}$\lbl{A2}}
\BIC{${{r}^{n}} \cdot r = {{r}^{S(n)}} \to \varphi_{60}$}
\end{ptree}
\noindent$\Pi_{7}$:
\begin{ptree}
\AXC{$x = x$\lbl{A5}}
\AXC{$x = y \to x = x \to y = x$\lbl{A6}}
\AXC{$(x = y \to x = x \to y = x) \to x = x \to x = y \to y = x$\lbl{A1}}
\BIC{$x = x \to x = y \to y = x$}
\BIC{$x = y \to y = x$}
\UIC{$\forall y\,(x = y \to y = x)$}
\UIC{$\forall x\,\forall y\,(x = y \to y = x)$}
\end{ptree}
\noindent$\Pi_{8}$:
\begin{ptree}
\AXC{$\varphi_{103}$\lbl{A2}}
\AXC{$r \in \mathbb{R} \to r \neq 0 \to \varphi_{32} \to \varphi_{32}$\lbl{A1}}
\AXC{$(r \in \mathbb{R} \to r \neq 0 \to \varphi_{32} \to \varphi_{32}) \to \varphi_{103} \to r \in \mathbb{R} \to r \neq 0 \to \varphi_{103}$\lbl{A1}}
\BIC{$\varphi_{103} \to r \in \mathbb{R} \to r \neq 0 \to \varphi_{103}$}
\BIC{$r \in \mathbb{R} \to r \neq 0 \to \varphi_{103}$}
\end{ptree}
\noindent$\Pi_{9}$:
\begin{ptree}
\AXC{$r \in \mathbb{R} \to r \neq 0 \to \varphi_{32} \to n \in B_{40} \to n \in B_{40}$\lbl{A1}}
\AXC{$\stackrel{\Pi_{8}}{\vdots}$}\noLine
\UIC{$r \in \mathbb{R} \to r \neq 0 \to \varphi_{103}$}
\AXC{$(r \in \mathbb{R} \to r \neq 0 \to \varphi_{103}) \to \varphi_{131}$\lbl{A1}}
\BIC{$\varphi_{131}$}
\BIC{$r \in \mathbb{R} \to \varphi_{109}$}
\end{ptree}
\noindent$\Pi_{10}$:
\begin{ptree}
\AXC{$\stackrel{\Pi_{9}}{\vdots}$}\noLine
\UIC{$r \in \mathbb{R} \to \varphi_{109}$}
\AXC{$(r \in \mathbb{R} \to \varphi_{109}) \to r \in \mathbb{R} \to r \neq 0 \to \varphi_{32} \to n \in B_{40} \to n \in \omega$\lbl{A1}}
\BIC{$r \in \mathbb{R} \to r \neq 0 \to \varphi_{32} \to n \in B_{40} \to n \in \omega$}
\end{ptree}
\noindent$\Pi_{11}$:
\begin{ptree}
\AXC{$\varphi_{102}$\lbl{A1}}
\AXC{$\stackrel{\Pi_{10}}{\vdots}$}\noLine
\UIC{$r \in \mathbb{R} \to r \neq 0 \to \varphi_{32} \to n \in B_{40} \to n \in \omega$}
\AXC{$\forall n\,(n \in \mathbb{N} \leftrightarrow n \in \omega \land n \neq \emptyset)$\lbl{Def$\mathbb{N}$}}
\AXC{$\varphi_{35}$\lbl{A2}}
\BIC{$n \in \mathbb{N} \leftrightarrow n \in \omega \land n \neq \emptyset$}
\AXC{$\varphi_{158}$\lbl{A1}}
\BIC{$(r \in \mathbb{R} \to r \neq 0 \to \varphi_{32} \to n \in B_{40} \to n \in \omega) \to \varphi_{102} \to \varphi_{98}$}
\BIC{$\varphi_{102} \to \varphi_{98}$}
\BIC{$\varphi_{98}$}
\end{ptree}
\noindent$\Pi_{12}$:
\begin{ptree}
\AXC{$\varphi_{100}$\lbl{A1}}
\AXC{$\forall r \in \mathbb{R}\,(\varphi_{37})$\lbl{[S24]}}
\AXC{$\forall r \in \mathbb{R}\,(\varphi_{37}) \to r \in \mathbb{R} \to \varphi_{37}$\lbl{A2}}
\BIC{$r \in \mathbb{R} \to \varphi_{37}$}
\AXC{$(r \in \mathbb{R} \to \varphi_{37}) \to \varphi_{148}$\lbl{A1}}
\BIC{$\varphi_{148}$}
\BIC{$r \in \mathbb{R} \to r \neq 0 \to \varphi_{32} \to n \in B_{40} \to S(n) \in \mathbb{N} \to n \neq \emptyset \to \varphi_{37}$}
\AXC{$\varphi_{154}$\lbl{A1}}
\BIC{$r \in \mathbb{R} \to r \neq 0 \to \varphi_{32} \to \varphi_{56}$}
\end{ptree}
\noindent$\Pi_{13}$:
\begin{ptree}
\AXC{$\varphi_{65}$\lbl{A2}}
\AXC{$\stackrel{\Pi_{12}}{\vdots}$}\noLine
\UIC{$r \in \mathbb{R} \to r \neq 0 \to \varphi_{32} \to \varphi_{56}$}
\AXC{$(r \in \mathbb{R} \to r \neq 0 \to \varphi_{32} \to \varphi_{56}) \to \varphi_{65} \to r \in \mathbb{R} \to r \neq 0 \to \varphi_{32} \to \varphi_{51}$\lbl{A1}}
\BIC{$\varphi_{65} \to r \in \mathbb{R} \to r \neq 0 \to \varphi_{32} \to \varphi_{51}$}
\BIC{$r \in \mathbb{R} \to r \neq 0 \to \varphi_{32} \to \varphi_{51}$}
\end{ptree}
\noindent$\Pi_{14}$:
\begin{ptree}
\AXC{$\stackrel{\Pi_{11}}{\vdots}$}\noLine
\UIC{$\varphi_{98}$}
\AXC{$\stackrel{\Pi_{13}}{\vdots}$}\noLine
\UIC{$r \in \mathbb{R} \to r \neq 0 \to \varphi_{32} \to \varphi_{51}$}
\AXC{$(r \in \mathbb{R} \to r \neq 0 \to \varphi_{32} \to \varphi_{51}) \to \varphi_{98} \to r \in \mathbb{R} \to \varphi_{112}$\lbl{A1}}
\BIC{$\varphi_{98} \to r \in \mathbb{R} \to \varphi_{112}$}
\BIC{$r \in \mathbb{R} \to \varphi_{112}$}
\AXC{$\varphi_{155}$\lbl{A1}}
\BIC{$({{r}^{S(n)}} = {{r}^{n}} \cdot r \to {{r}^{n}} \cdot r = {{r}^{S(n)}}) \to r \in \mathbb{R} \to \varphi_{113}$}
\end{ptree}
\noindent$\Pi_{15}$:
\begin{ptree}
\AXC{$\stackrel{\Pi_{7}}{\vdots}$}\noLine
\UIC{$\forall x\,\forall y\,(x = y \to y = x)$}
\AXC{$\varphi_{14}$\lbl{A2}}
\BIC{$\forall y\,({{r}^{S(n)}} = y \to y = {{r}^{S(n)}})$}
\AXC{$\varphi_{82}$\lbl{A2}}
\BIC{${{r}^{S(n)}} = {{r}^{n}} \cdot r \to {{r}^{n}} \cdot r = {{r}^{S(n)}}$}
\AXC{$\stackrel{\Pi_{14}}{\vdots}$}\noLine
\UIC{$({{r}^{S(n)}} = {{r}^{n}} \cdot r \to {{r}^{n}} \cdot r = {{r}^{S(n)}}) \to r \in \mathbb{R} \to \varphi_{113}$}
\BIC{$r \in \mathbb{R} \to \varphi_{113}$}
\end{ptree}
\noindent$\Pi_{16}$:
\begin{ptree}
\AXC{$x = x$\lbl{A5}}
\UIC{$\forall x\,(x = x)$}
\AXC{$\forall x\,(x = x) \to ({{r}^{n}} \cdot r)^{-1} = ({{r}^{n}} \cdot r)^{-1}$\lbl{A2}}
\BIC{$({{r}^{n}} \cdot r)^{-1} = ({{r}^{n}} \cdot r)^{-1}$}
\AXC{$\stackrel{\Pi_{15}}{\vdots}$}\noLine
\UIC{$r \in \mathbb{R} \to \varphi_{113}$}
\AXC{$(r \in \mathbb{R} \to \varphi_{113}) \to \varphi_{146}$\lbl{A1}}
\BIC{$\varphi_{146}$}
\BIC{$({{r}^{n}} \cdot r = {{r}^{S(n)}} \to \varphi_{60}) \to r \in \mathbb{R} \to r \neq 0 \to \varphi_{32} \to \varphi_{46}$}
\end{ptree}
\noindent$\Pi_{17}$:
\begin{ptree}
\AXC{$y = z \to x = y \to x = z$\lbl{A6}}
\AXC{$(y = z \to x = y \to x = z) \to x = y \to y = z \to x = z$\lbl{A1}}
\BIC{$x = y \to y = z \to x = z$}
\UIC{$\forall z\,(x = y \to y = z \to x = z)$}
\UIC{$\forall y\,\forall z\,(x = y \to y = z \to x = z)$}
\UIC{$\forall x\,\forall y\,\forall z\,(x = y \to y = z \to x = z)$}
\AXC{$\forall x\,\forall y\,\forall z\,(x = y \to y = z \to x = z) \to \varphi_{16}$\lbl{A2}}
\BIC{$\varphi_{16}$}
\end{ptree}
\noindent$\Pi_{18}$:
\begin{ptree}
\AXC{$\stackrel{\Pi_{17}}{\vdots}$}\noLine
\UIC{$\varphi_{16}$}
\AXC{$\varphi_{16} \to \varphi_{77}$\lbl{A2}}
\BIC{$\varphi_{77}$}
\AXC{$\varphi_{77} \to {{(r^{-1})}^{S(n)}} = {{r}^{n}}^{-1} \cdot r^{-1} \to \varphi_{67}$\lbl{A2}}
\BIC{${{(r^{-1})}^{S(n)}} = {{r}^{n}}^{-1} \cdot r^{-1} \to \varphi_{67}$}
\end{ptree}
\noindent$\Pi_{19}$:
\begin{ptree}
\AXC{$\stackrel{\Pi_{7}}{\vdots}$}\noLine
\UIC{$\forall x\,\forall y\,(x = y \to y = x)$}
\AXC{$\varphi_{25}$\lbl{A2}}
\BIC{$\forall y\,(({{r}^{n}} \cdot r)^{-1} = y \to y = ({{r}^{n}} \cdot r)^{-1})$}
\AXC{$\forall y\,(({{r}^{n}} \cdot r)^{-1} = y \to y = ({{r}^{n}} \cdot r)^{-1}) \to \varphi_{73}$\lbl{A2}}
\BIC{$\varphi_{73}$}
\end{ptree}
\noindent$\Pi_{20}$:
\begin{ptree}
\AXC{$r \in \mathbb{R} \to r \neq 0 \to r \neq 0$\lbl{A1}}
\AXC{$\varphi_{100}$\lbl{A1}}
\AXC{$\forall r \in \mathbb{R}\,(r \neq 0 \to \forall n \in \mathbb{N}\,({{r}^{n}} \neq 0))$\lbl{[P4]}}
\AXC{$\varphi_{75}$\lbl{A2}}
\BIC{$r \in \mathbb{R} \to r \neq 0 \to \forall n \in \mathbb{N}\,({{r}^{n}} \neq 0)$}
\AXC{$\varphi_{151}$\lbl{A1}}
\BIC{$\varphi_{100} \to (r \in \mathbb{R} \to r \neq 0 \to r \neq 0) \to r \in \mathbb{R} \to r \neq 0 \to \varphi_{101}$}
\BIC{$(r \in \mathbb{R} \to r \neq 0 \to r \neq 0) \to r \in \mathbb{R} \to r \neq 0 \to \varphi_{101}$}
\BIC{$r \in \mathbb{R} \to r \neq 0 \to \varphi_{101}$}
\end{ptree}
\noindent$\Pi_{21}$:
\begin{ptree}
\AXC{$\forall n \in \mathbb{N}\,({{r}^{n}} \neq 0) \to n \in \mathbb{N} \to {{r}^{n}} \neq 0$\lbl{A2}}
\AXC{$\stackrel{\Pi_{20}}{\vdots}$}\noLine
\UIC{$r \in \mathbb{R} \to r \neq 0 \to \varphi_{101}$}
\AXC{$\varphi_{152}$\lbl{A1}}
\BIC{$(\forall n \in \mathbb{N}\,({{r}^{n}} \neq 0) \to n \in \mathbb{N} \to {{r}^{n}} \neq 0) \to r \in \mathbb{R} \to \varphi_{104}$}
\BIC{$r \in \mathbb{R} \to \varphi_{104}$}
\AXC{$(r \in \mathbb{R} \to \varphi_{104}) \to \varphi_{98} \to \varphi_{105}$\lbl{A1}}
\BIC{$\varphi_{98} \to \varphi_{105}$}
\end{ptree}
\noindent$\Pi_{22}$:
\begin{ptree}
\AXC{$\varphi_{100}$\lbl{A1}}
\AXC{$\forall r \in \mathbb{R}\,(\varphi_{37})$\lbl{[S24]}}
\AXC{$\forall r \in \mathbb{R}\,(\varphi_{37}) \to r \in \mathbb{R} \to \varphi_{37}$\lbl{A2}}
\BIC{$r \in \mathbb{R} \to \varphi_{37}$}
\AXC{$(r \in \mathbb{R} \to \varphi_{37}) \to \varphi_{148}$\lbl{A1}}
\BIC{$\varphi_{148}$}
\BIC{$r \in \mathbb{R} \to r \neq 0 \to \varphi_{32} \to n \in B_{40} \to S(n) \in \mathbb{N} \to n \neq \emptyset \to \varphi_{37}$}
\AXC{$\varphi_{149}$\lbl{A1}}
\BIC{$r \in \mathbb{R} \to r \neq 0 \to \varphi_{32} \to \varphi_{5}$}
\end{ptree}
\noindent$\Pi_{23}$:
\begin{ptree}
\AXC{$\stackrel{\Pi_{22}}{\vdots}$}\noLine
\UIC{$r \in \mathbb{R} \to r \neq 0 \to \varphi_{32} \to \varphi_{5}$}
\AXC{$\forall a\,\forall n\,(\mathrm{Seq}(a) \to n \in \mathbb{N} \to a_{n} \in \mathbb{R})$\lbl{[S1]}}
\AXC{$\varphi_{66}$\lbl{A2}}
\BIC{$\forall n\,(\mathrm{Seq}({(r)^{\bullet}}) \to n \in \mathbb{N} \to {{r}^{n}} \in \mathbb{R})$}
\AXC{$\varphi_{69}$\lbl{A2}}
\BIC{$\mathrm{Seq}({(r)^{\bullet}}) \to n \in \mathbb{N} \to {{r}^{n}} \in \mathbb{R}$}
\AXC{$\varphi_{159}$\lbl{A1}}
\BIC{$(r \in \mathbb{R} \to r \neq 0 \to \varphi_{32} \to \varphi_{5}) \to \varphi_{98} \to \varphi_{114}$}
\BIC{$\varphi_{98} \to \varphi_{114}$}
\end{ptree}
\noindent$\Pi_{24}$:
\begin{ptree}
\AXC{$\stackrel{\Pi_{11}}{\vdots}$}\noLine
\UIC{$\varphi_{98}$}
\AXC{$\stackrel{\Pi_{23}}{\vdots}$}\noLine
\UIC{$\varphi_{98} \to \varphi_{114}$}
\BIC{$\varphi_{114}$}
\AXC{$\varphi_{28}$\lbl{[E16]}}
\AXC{$\varphi_{28} \to \varphi_{86}$\lbl{A2}}
\BIC{$\varphi_{86}$}
\AXC{$\varphi_{86} \to {{r}^{n}} \in \mathbb{R} \to \varphi_{54}$\lbl{A2}}
\BIC{${{r}^{n}} \in \mathbb{R} \to \varphi_{54}$}
\AXC{$({{r}^{n}} \in \mathbb{R} \to \varphi_{54}) \to \varphi_{114} \to \varphi_{160}$\lbl{A1}}
\BIC{$\varphi_{114} \to \varphi_{160}$}
\BIC{$\varphi_{160}$}
\end{ptree}
\noindent$\Pi_{25}$:
\begin{ptree}
\AXC{$r \in \mathbb{R} \to r \neq 0 \to r \neq 0$\lbl{A1}}
\AXC{$\stackrel{\Pi_{11}}{\vdots}$}\noLine
\UIC{$\varphi_{98}$}
\AXC{$\stackrel{\Pi_{21}}{\vdots}$}\noLine
\UIC{$\varphi_{98} \to \varphi_{105}$}
\BIC{$\varphi_{105}$}
\AXC{$\varphi_{100}$\lbl{A1}}
\AXC{$\stackrel{\Pi_{24}}{\vdots}$}\noLine
\UIC{$\varphi_{160}$}
\BIC{$\varphi_{105} \to (r \in \mathbb{R} \to r \neq 0 \to r \neq 0) \to r \in \mathbb{R} \to r \neq 0 \to \varphi_{32} \to \varphi_{61}$}
\BIC{$(r \in \mathbb{R} \to r \neq 0 \to r \neq 0) \to r \in \mathbb{R} \to r \neq 0 \to \varphi_{32} \to \varphi_{61}$}
\BIC{$r \in \mathbb{R} \to r \neq 0 \to \varphi_{32} \to \varphi_{61}$}
\end{ptree}
\noindent$\Pi_{26}$:
\begin{ptree}
\AXC{$\stackrel{\Pi_{19}}{\vdots}$}\noLine
\UIC{$\varphi_{73}$}
\AXC{$\stackrel{\Pi_{25}}{\vdots}$}\noLine
\UIC{$r \in \mathbb{R} \to r \neq 0 \to \varphi_{32} \to \varphi_{61}$}
\AXC{$(r \in \mathbb{R} \to r \neq 0 \to \varphi_{32} \to \varphi_{61}) \to \varphi_{73} \to r \in \mathbb{R} \to r \neq 0 \to \varphi_{32} \to \varphi_{62}$\lbl{A1}}
\BIC{$\varphi_{73} \to r \in \mathbb{R} \to r \neq 0 \to \varphi_{32} \to \varphi_{62}$}
\BIC{$r \in \mathbb{R} \to r \neq 0 \to \varphi_{32} \to \varphi_{62}$}
\end{ptree}
\noindent$\Pi_{27}$:
\begin{ptree}
\AXC{$y = z \to x = y \to x = z$\lbl{A6}}
\AXC{$(y = z \to x = y \to x = z) \to x = y \to y = z \to x = z$\lbl{A1}}
\BIC{$x = y \to y = z \to x = z$}
\UIC{$\forall z\,(x = y \to y = z \to x = z)$}
\UIC{$\forall y\,\forall z\,(x = y \to y = z \to x = z)$}
\UIC{$\forall x\,\forall y\,\forall z\,(x = y \to y = z \to x = z)$}
\AXC{$\forall x\,\forall y\,\forall z\,(x = y \to y = z \to x = z) \to \varphi_{16}$\lbl{A2}}
\BIC{$\varphi_{16}$}
\end{ptree}
\noindent$\Pi_{28}$:
\begin{ptree}
\AXC{$\stackrel{\Pi_{11}}{\vdots}$}\noLine
\UIC{$\varphi_{98}$}
\AXC{$\stackrel{\Pi_{9}}{\vdots}$}\noLine
\UIC{$r \in \mathbb{R} \to \varphi_{109}$}
\AXC{$(r \in \mathbb{R} \to \varphi_{109}) \to \varphi_{108}$\lbl{A1}}
\BIC{$\varphi_{108}$}
\AXC{$\varphi_{108} \to \varphi_{98} \to r \in \mathbb{R} \to r \neq 0 \to \varphi_{107}$\lbl{A1}}
\BIC{$\varphi_{98} \to r \in \mathbb{R} \to r \neq 0 \to \varphi_{107}$}
\BIC{$r \in \mathbb{R} \to r \neq 0 \to \varphi_{107}$}
\AXC{$\varphi_{167}$\lbl{A1}}
\BIC{${{(r^{-1})}^{n}} \cdot r^{-1} = {{(r^{-1})}^{n}} \cdot r^{-1} \to \varphi_{142}$}
\end{ptree}
\noindent$\Pi_{29}$:
\begin{ptree}
\AXC{$x = x$\lbl{A5}}
\UIC{$\forall x\,(x = x)$}
\AXC{$\forall x\,(x = x) \to {{(r^{-1})}^{n}} \cdot r^{-1} = {{(r^{-1})}^{n}} \cdot r^{-1}$\lbl{A2}}
\BIC{${{(r^{-1})}^{n}} \cdot r^{-1} = {{(r^{-1})}^{n}} \cdot r^{-1}$}
\AXC{$\stackrel{\Pi_{28}}{\vdots}$}\noLine
\UIC{${{(r^{-1})}^{n}} \cdot r^{-1} = {{(r^{-1})}^{n}} \cdot r^{-1} \to \varphi_{142}$}
\BIC{$\varphi_{142}$}
\end{ptree}
\noindent$\Pi_{30}$:
\begin{ptree}
\AXC{$z = w \to \varphi_{38}$\lbl{A6}}
\UIC{$\forall w\,(z = w \to \varphi_{38})$}
\UIC{$\forall z\,\forall w\,(z = w \to \varphi_{38})$}
\AXC{$\forall z\,\forall w\,(z = w \to \varphi_{38}) \to \varphi_{88}$\lbl{A2}}
\BIC{$\varphi_{88}$}
\AXC{$\varphi_{88} \to {{(r^{-1})}^{n}} = {{r}^{n}}^{-1} \to \varphi_{85}$\lbl{A2}}
\BIC{${{(r^{-1})}^{n}} = {{r}^{n}}^{-1} \to \varphi_{85}$}
\AXC{$\stackrel{\Pi_{29}}{\vdots}$}\noLine
\UIC{$\varphi_{142}$}
\BIC{$r \in \mathbb{R} \to r \neq 0 \to \varphi_{32} \to n \in B_{40} \to \varphi_{47}$}
\end{ptree}
\noindent$\Pi_{31}$:
\begin{ptree}
\AXC{$r \in \mathbb{R} \to r \neq 0 \to r \neq 0$\lbl{A1}}
\AXC{$\varphi_{100}$\lbl{A1}}
\AXC{$\forall x \in \mathbb{R}\,(x \neq 0 \to x^{-1} \in \mathbb{R})$\lbl{(R1)}}
\AXC{$\varphi_{31}$\lbl{A2}}
\BIC{$r \in \mathbb{R} \to r \neq 0 \to r^{-1} \in \mathbb{R}$}
\AXC{$\varphi_{147}$\lbl{A1}}
\BIC{$\varphi_{100} \to (r \in \mathbb{R} \to r \neq 0 \to r \neq 0) \to \varphi_{106}$}
\BIC{$(r \in \mathbb{R} \to r \neq 0 \to r \neq 0) \to \varphi_{106}$}
\BIC{$\varphi_{106}$}
\end{ptree}
\noindent$\Pi_{32}$:
\begin{ptree}
\AXC{$\stackrel{\Pi_{31}}{\vdots}$}\noLine
\UIC{$\varphi_{106}$}
\AXC{$\forall r \in \mathbb{R}\,(\varphi_{37})$\lbl{[S24]}}
\AXC{$\forall r \in \mathbb{R}\,(\varphi_{37}) \to r^{-1} \in \mathbb{R} \to \varphi_{81}$\lbl{A2}}
\BIC{$r^{-1} \in \mathbb{R} \to \varphi_{81}$}
\AXC{$(r^{-1} \in \mathbb{R} \to \varphi_{81}) \to \varphi_{153}$\lbl{A1}}
\BIC{$\varphi_{153}$}
\BIC{$r \in \mathbb{R} \to r \neq 0 \to \varphi_{32} \to n \in B_{40} \to S(n) \in \mathbb{N} \to n \neq \emptyset \to \varphi_{81}$}
\AXC{$\varphi_{157}$\lbl{A1}}
\BIC{$r \in \mathbb{R} \to r \neq 0 \to \varphi_{32} \to n \in B_{40} \to \varphi_{63}$}
\end{ptree}
\noindent$\Pi_{33}$:
\begin{ptree}
\AXC{$\stackrel{\Pi_{11}}{\vdots}$}\noLine
\UIC{$\varphi_{98}$}
\AXC{$\varphi_{89}$\lbl{A2}}
\AXC{$\stackrel{\Pi_{32}}{\vdots}$}\noLine
\UIC{$r \in \mathbb{R} \to r \neq 0 \to \varphi_{32} \to n \in B_{40} \to \varphi_{63}$}
\AXC{$\varphi_{162}$\lbl{A1}}
\BIC{$\varphi_{89} \to r \in \mathbb{R} \to r \neq 0 \to \varphi_{32} \to n \in B_{40} \to \varphi_{57}$}
\BIC{$r \in \mathbb{R} \to r \neq 0 \to \varphi_{32} \to n \in B_{40} \to \varphi_{57}$}
\AXC{$\varphi_{164}$\lbl{A1}}
\BIC{$\varphi_{98} \to r \in \mathbb{R} \to r \neq 0 \to \varphi_{32} \to n \in B_{40} \to \varphi_{34}$}
\BIC{$r \in \mathbb{R} \to r \neq 0 \to \varphi_{32} \to n \in B_{40} \to \varphi_{34}$}
\AXC{$\varphi_{174}$\lbl{A1}}
\BIC{$(r \in \mathbb{R} \to r \neq 0 \to \varphi_{32} \to n \in B_{40} \to \varphi_{47}) \to \varphi_{141}$}
\end{ptree}
\noindent$\Pi_{34}$:
\begin{ptree}
\AXC{$\stackrel{\Pi_{30}}{\vdots}$}\noLine
\UIC{$r \in \mathbb{R} \to r \neq 0 \to \varphi_{32} \to n \in B_{40} \to \varphi_{47}$}
\AXC{$\stackrel{\Pi_{33}}{\vdots}$}\noLine
\UIC{$(r \in \mathbb{R} \to r \neq 0 \to \varphi_{32} \to n \in B_{40} \to \varphi_{47}) \to \varphi_{141}$}
\BIC{$\varphi_{141}$}
\end{ptree}
\noindent$\Pi_{35}$:
\begin{ptree}
\AXC{$\stackrel{\Pi_{27}}{\vdots}$}\noLine
\UIC{$\varphi_{16}$}
\AXC{$\varphi_{16} \to \varphi_{78}$\lbl{A2}}
\BIC{$\varphi_{78}$}
\AXC{$\varphi_{78} \to {{(r^{-1})}^{S(n)}} = {{(r^{-1})}^{n}} \cdot r^{-1} \to \varphi_{76}$\lbl{A2}}
\BIC{${{(r^{-1})}^{S(n)}} = {{(r^{-1})}^{n}} \cdot r^{-1} \to \varphi_{76}$}
\AXC{$\stackrel{\Pi_{34}}{\vdots}$}\noLine
\UIC{$\varphi_{141}$}
\BIC{$r \in \mathbb{R} \to r \neq 0 \to \varphi_{32} \to \varphi_{58}$}
\AXC{$\varphi_{172}$\lbl{A1}}
\BIC{$(r \in \mathbb{R} \to r \neq 0 \to \varphi_{32} \to \varphi_{62}) \to \varphi_{138}$}
\end{ptree}
\noindent$\Pi_{36}$:
\begin{ptree}
\AXC{$\stackrel{\Pi_{18}}{\vdots}$}\noLine
\UIC{${{(r^{-1})}^{S(n)}} = {{r}^{n}}^{-1} \cdot r^{-1} \to \varphi_{67}$}
\AXC{$\stackrel{\Pi_{26}}{\vdots}$}\noLine
\UIC{$r \in \mathbb{R} \to r \neq 0 \to \varphi_{32} \to \varphi_{62}$}
\AXC{$\stackrel{\Pi_{35}}{\vdots}$}\noLine
\UIC{$(r \in \mathbb{R} \to r \neq 0 \to \varphi_{32} \to \varphi_{62}) \to \varphi_{138}$}
\BIC{$\varphi_{138}$}
\BIC{$r \in \mathbb{R} \to r \neq 0 \to \varphi_{32} \to \varphi_{48}$}
\AXC{$\varphi_{170}$\lbl{A1}}
\BIC{$(r \in \mathbb{R} \to r \neq 0 \to \varphi_{32} \to \varphi_{46}) \to \varphi_{135}$}
\end{ptree}
\noindent$\Pi_{37}$:
\begin{ptree}
\AXC{$\stackrel{\Pi_{6}}{\vdots}$}\noLine
\UIC{${{r}^{n}} \cdot r = {{r}^{S(n)}} \to \varphi_{60}$}
\AXC{$\stackrel{\Pi_{16}}{\vdots}$}\noLine
\UIC{$({{r}^{n}} \cdot r = {{r}^{S(n)}} \to \varphi_{60}) \to r \in \mathbb{R} \to r \neq 0 \to \varphi_{32} \to \varphi_{46}$}
\BIC{$r \in \mathbb{R} \to r \neq 0 \to \varphi_{32} \to \varphi_{46}$}
\AXC{$\stackrel{\Pi_{36}}{\vdots}$}\noLine
\UIC{$(r \in \mathbb{R} \to r \neq 0 \to \varphi_{32} \to \varphi_{46}) \to \varphi_{135}$}
\BIC{$\varphi_{135}$}
\end{ptree}
\noindent$\Pi_{38}$:
\begin{ptree}
\AXC{$\stackrel{\Pi_{5}}{\vdots}$}\noLine
\UIC{$\varphi_{16}$}
\AXC{$\varphi_{16} \to \varphi_{68}$\lbl{A2}}
\BIC{$\varphi_{68}$}
\AXC{$\varphi_{68} \to {{(r^{-1})}^{S(n)}} = ({{r}^{n}} \cdot r)^{-1} \to \varphi_{52}$\lbl{A2}}
\BIC{${{(r^{-1})}^{S(n)}} = ({{r}^{n}} \cdot r)^{-1} \to \varphi_{52}$}
\AXC{$\stackrel{\Pi_{37}}{\vdots}$}\noLine
\UIC{$\varphi_{135}$}
\BIC{$r \in \mathbb{R} \to r \neq 0 \to \varphi_{32} \to \varphi_{42}$}
\end{ptree}
\noindent$\Pi_{39}$:
\begin{ptree}
\AXC{$z = w \to {{(r^{-1})}^{z}} = {{r}^{z}}^{-1} \to {{(r^{-1})}^{w}} = {{r}^{w}}^{-1}$\lbl{A6}}
\UIC{$\varphi_{40}$}
\UIC{$\forall z\,(\varphi_{40})$}
\AXC{$\forall z\,(\varphi_{40}) \to \varphi_{59}$\lbl{A2}}
\BIC{$\varphi_{59}$}
\AXC{$\varphi_{59} \to 1_{\omega} = S(n) \to \varphi_{50}$\lbl{A2}}
\BIC{$1_{\omega} = S(n) \to \varphi_{50}$}
\end{ptree}
\noindent$\Pi_{40}$:
\begin{ptree}
\AXC{$y = z \to x = y \to x = z$\lbl{A6}}
\AXC{$(y = z \to x = y \to x = z) \to x = y \to y = z \to x = z$\lbl{A1}}
\BIC{$x = y \to y = z \to x = z$}
\UIC{$\forall z\,(x = y \to y = z \to x = z)$}
\UIC{$\forall y\,\forall z\,(x = y \to y = z \to x = z)$}
\UIC{$\forall x\,\forall y\,\forall z\,(x = y \to y = z \to x = z)$}
\AXC{$\forall x\,\forall y\,\forall z\,(x = y \to y = z \to x = z) \to \varphi_{17}$\lbl{A2}}
\BIC{$\varphi_{17}$}
\end{ptree}
\noindent$\Pi_{41}$:
\begin{ptree}
\AXC{$r \in \mathbb{R} \to r \neq 0 \to r \in \mathbb{R}$\lbl{A1}}
\AXC{$\forall r \in \mathbb{R}\,(\varphi_{37})$\lbl{[S24]}}
\AXC{$\forall r \in \mathbb{R}\,(\varphi_{37}) \to r \in \mathbb{R} \to \varphi_{37}$\lbl{A2}}
\BIC{$r \in \mathbb{R} \to \varphi_{37}$}
\AXC{$\varphi_{125}$\lbl{A1}}
\BIC{$(r \in \mathbb{R} \to r \neq 0 \to r \in \mathbb{R}) \to r \in \mathbb{R} \to r \neq 0 \to \varphi_{37}$}
\BIC{$r \in \mathbb{R} \to r \neq 0 \to \varphi_{37}$}
\end{ptree}
\noindent$\Pi_{42}$:
\begin{ptree}
\AXC{$\stackrel{\Pi_{41}}{\vdots}$}\noLine
\UIC{$r \in \mathbb{R} \to r \neq 0 \to \varphi_{37}$}
\AXC{$(r \in \mathbb{R} \to r \neq 0 \to \varphi_{37}) \to r \in \mathbb{R} \to r \neq 0 \to {{r}^{1_{\omega}}} = r$\lbl{A1}}
\BIC{$r \in \mathbb{R} \to r \neq 0 \to {{r}^{1_{\omega}}} = r$}
\AXC{$(r \in \mathbb{R} \to r \neq 0 \to {{r}^{1_{\omega}}} = r) \to \varphi_{33}$\lbl{A1}}
\BIC{$\varphi_{33}$}
\end{ptree}
\noindent$\Pi_{43}$:
\begin{ptree}
\AXC{$\stackrel{\Pi_{7}}{\vdots}$}\noLine
\UIC{$\forall x\,\forall y\,(x = y \to y = x)$}
\AXC{$\varphi_{15}$\lbl{A2}}
\BIC{$\forall y\,({{r}^{1_{\omega}}} = y \to y = {{r}^{1_{\omega}}})$}
\AXC{$\varphi_{36}$\lbl{A2}}
\BIC{${{r}^{1_{\omega}}} = r \to r = {{r}^{1_{\omega}}}$}
\AXC{$\stackrel{\Pi_{42}}{\vdots}$}\noLine
\UIC{$\varphi_{33}$}
\BIC{$r \in \mathbb{R} \to r \neq 0 \to r = {{r}^{1_{\omega}}}$}
\AXC{$(r \in \mathbb{R} \to r \neq 0 \to r = {{r}^{1_{\omega}}}) \to r^{-1} = r^{-1} \to \varphi_{83}$\lbl{A1}}
\BIC{$r^{-1} = r^{-1} \to \varphi_{83}$}
\end{ptree}
\noindent$\Pi_{44}$:
\begin{ptree}
\AXC{$x = x$\lbl{A5}}
\UIC{$\forall x\,(x = x)$}
\AXC{$\forall x\,(x = x) \to r^{-1} = r^{-1}$\lbl{A2}}
\BIC{$r^{-1} = r^{-1}$}
\AXC{$\stackrel{\Pi_{43}}{\vdots}$}\noLine
\UIC{$r^{-1} = r^{-1} \to \varphi_{83}$}
\BIC{$\varphi_{83}$}
\end{ptree}
\noindent$\Pi_{45}$:
\begin{ptree}
\AXC{$\forall x \in \mathbb{R}\,(x \neq 0 \to x^{-1} \in \mathbb{R})$\lbl{(R1)}}
\AXC{$\varphi_{31}$\lbl{A2}}
\BIC{$r \in \mathbb{R} \to r \neq 0 \to r^{-1} \in \mathbb{R}$}
\AXC{$\forall r \in \mathbb{R}\,(\varphi_{37})$\lbl{[S24]}}
\AXC{$\forall r \in \mathbb{R}\,(\varphi_{37}) \to r^{-1} \in \mathbb{R} \to \varphi_{81}$\lbl{A2}}
\BIC{$r^{-1} \in \mathbb{R} \to \varphi_{81}$}
\AXC{$\varphi_{133}$\lbl{A1}}
\BIC{$(r \in \mathbb{R} \to r \neq 0 \to r^{-1} \in \mathbb{R}) \to r \in \mathbb{R} \to r \neq 0 \to \varphi_{81}$}
\BIC{$r \in \mathbb{R} \to r \neq 0 \to \varphi_{81}$}
\end{ptree}
\noindent$\Pi_{46}$:
\begin{ptree}
\AXC{$\stackrel{\Pi_{45}}{\vdots}$}\noLine
\UIC{$r \in \mathbb{R} \to r \neq 0 \to \varphi_{81}$}
\AXC{$\varphi_{119}$\lbl{A1}}
\BIC{$r \in \mathbb{R} \to r \neq 0 \to {{(r^{-1})}^{1_{\omega}}} = r^{-1}$}
\AXC{$(r \in \mathbb{R} \to r \neq 0 \to {{(r^{-1})}^{1_{\omega}}} = r^{-1}) \to \varphi_{123}$\lbl{A1}}
\BIC{$\varphi_{123}$}
\end{ptree}
\noindent$\Pi_{47}$:
\begin{ptree}
\AXC{$z = w \to r^{-1} = z^{-1} \to r^{-1} = w^{-1}$\lbl{A6}}
\UIC{$\forall w\,(z = w \to r^{-1} = z^{-1} \to r^{-1} = w^{-1})$}
\UIC{$\forall z\,\forall w\,(z = w \to r^{-1} = z^{-1} \to r^{-1} = w^{-1})$}
\AXC{$\varphi_{41}$\lbl{A2}}
\BIC{$\forall w\,(r = w \to r^{-1} = r^{-1} \to r^{-1} = w^{-1})$}
\AXC{$\varphi_{74}$\lbl{A2}}
\BIC{$r = {{r}^{1_{\omega}}} \to r^{-1} = r^{-1} \to r^{-1} = {{r}^{1_{\omega}}}^{-1}$}
\AXC{$\stackrel{\Pi_{44}}{\vdots}$}\noLine
\UIC{$\varphi_{83}$}
\BIC{$r \in \mathbb{R} \to r \neq 0 \to r^{-1} = {{r}^{1_{\omega}}}^{-1}$}
\AXC{$\stackrel{\Pi_{46}}{\vdots}$}\noLine
\UIC{$\varphi_{123}$}
\BIC{$\varphi_{70} \to r \in \mathbb{R} \to r \neq 0 \to {{(r^{-1})}^{1_{\omega}}} = {{r}^{1_{\omega}}}^{-1}$}
\end{ptree}
\noindent$\Pi_{48}$:
\begin{ptree}
\AXC{$\stackrel{\Pi_{40}}{\vdots}$}\noLine
\UIC{$\varphi_{17}$}
\AXC{$\varphi_{17} \to \varphi_{22}$\lbl{A2}}
\BIC{$\varphi_{22}$}
\AXC{$\varphi_{22} \to \varphi_{70}$\lbl{A2}}
\BIC{$\varphi_{70}$}
\AXC{$\stackrel{\Pi_{47}}{\vdots}$}\noLine
\UIC{$\varphi_{70} \to r \in \mathbb{R} \to r \neq 0 \to {{(r^{-1})}^{1_{\omega}}} = {{r}^{1_{\omega}}}^{-1}$}
\BIC{$r \in \mathbb{R} \to r \neq 0 \to {{(r^{-1})}^{1_{\omega}}} = {{r}^{1_{\omega}}}^{-1}$}
\end{ptree}
\noindent$\Pi_{49}$:
\begin{ptree}
\AXC{$y = z \to x = y \to x = z$\lbl{A6}}
\AXC{$(y = z \to x = y \to x = z) \to x = y \to y = z \to x = z$\lbl{A1}}
\BIC{$x = y \to y = z \to x = z$}
\UIC{$\forall z\,(x = y \to y = z \to x = z)$}
\UIC{$\forall y\,\forall z\,(x = y \to y = z \to x = z)$}
\UIC{$\forall x\,\forall y\,\forall z\,(x = y \to y = z \to x = z)$}
\AXC{$\varphi_{8}$\lbl{A2}}
\BIC{$\forall y\,\forall z\,(S(n) = y \to y = z \to S(n) = z)$}
\AXC{$\varphi_{29}$\lbl{A2}}
\BIC{$\forall z\,(S(n) = S(\emptyset) \to S(\emptyset) = z \to S(n) = z)$}
\end{ptree}
\noindent$\Pi_{50}$:
\begin{ptree}
\AXC{$\stackrel{\Pi_{49}}{\vdots}$}\noLine
\UIC{$\forall z\,(S(n) = S(\emptyset) \to S(\emptyset) = z \to S(n) = z)$}
\AXC{$\forall z\,(S(n) = S(\emptyset) \to S(\emptyset) = z \to S(n) = z) \to \varphi_{4}$\lbl{A2}}
\BIC{$\varphi_{4}$}
\end{ptree}
\noindent$\Pi_{51}$:
\begin{ptree}
\AXC{$\stackrel{\Pi_{7}}{\vdots}$}\noLine
\UIC{$\forall x\,\forall y\,(x = y \to y = x)$}
\AXC{$\varphi_{2}$\lbl{A2}}
\BIC{$\forall y\,(1_{\omega} = y \to y = 1_{\omega})$}
\AXC{$\varphi_{10}$\lbl{A2}}
\BIC{$1_{\omega} = S(\emptyset) \to S(\emptyset) = 1_{\omega}$}
\AXC{$1_{\omega} = S(\emptyset)$\lbl{Def$1_\omega$}}
\AXC{$1_{\omega} = S(\emptyset) \to \varphi_{9}$\lbl{A1}}
\BIC{$\varphi_{9}$}
\BIC{$S(\emptyset) = 1_{\omega}$}
\end{ptree}
\noindent$\Pi_{52}$:
\begin{ptree}
\AXC{$x = x$\lbl{A5}}
\UIC{$\forall x\,(x = x)$}
\AXC{$\forall x\,(x = x) \to S(n) = S(n)$\lbl{A2}}
\BIC{$S(n) = S(n)$}
\AXC{$\varphi_{99}$\lbl{A1}}
\AXC{$\varphi_{99} \to S(n) = S(n) \to \varphi_{120}$\lbl{A1}}
\BIC{$S(n) = S(n) \to \varphi_{120}$}
\BIC{$\varphi_{120}$}
\end{ptree}
\noindent$\Pi_{53}$:
\begin{ptree}
\AXC{$z = w \to S(n) = S(z) \to S(n) = S(w)$\lbl{A6}}
\UIC{$\forall w\,(z = w \to S(n) = S(z) \to S(n) = S(w))$}
\UIC{$\forall z\,\forall w\,(z = w \to S(n) = S(z) \to S(n) = S(w))$}
\AXC{$\varphi_{23}$\lbl{A2}}
\BIC{$\forall w\,(n = w \to S(n) = S(n) \to S(n) = S(w))$}
\AXC{$\varphi_{30}$\lbl{A2}}
\BIC{$n = \emptyset \to S(n) = S(n) \to S(n) = S(\emptyset)$}
\AXC{$\stackrel{\Pi_{52}}{\vdots}$}\noLine
\UIC{$\varphi_{120}$}
\BIC{$r \in \mathbb{R} \to \varphi_{95}$}
\end{ptree}
\noindent$\Pi_{54}$:
\begin{ptree}
\AXC{$\stackrel{\Pi_{50}}{\vdots}$}\noLine
\UIC{$\varphi_{4}$}
\AXC{$\stackrel{\Pi_{51}}{\vdots}$}\noLine
\UIC{$S(\emptyset) = 1_{\omega}$}
\AXC{$\stackrel{\Pi_{53}}{\vdots}$}\noLine
\UIC{$r \in \mathbb{R} \to \varphi_{95}$}
\AXC{$(r \in \mathbb{R} \to \varphi_{95}) \to S(\emptyset) = 1_{\omega} \to \varphi_{4} \to r \in \mathbb{R} \to \varphi_{92}$\lbl{A1}}
\BIC{$S(\emptyset) = 1_{\omega} \to \varphi_{4} \to r \in \mathbb{R} \to \varphi_{92}$}
\BIC{$\varphi_{4} \to r \in \mathbb{R} \to \varphi_{92}$}
\BIC{$r \in \mathbb{R} \to \varphi_{92}$}
\AXC{$\varphi_{136}$\lbl{A1}}
\BIC{$(S(n) = 1_{\omega} \to 1_{\omega} = S(n)) \to r \in \mathbb{R} \to \varphi_{93}$}
\end{ptree}
\noindent$\Pi_{55}$:
\begin{ptree}
\AXC{$\stackrel{\Pi_{7}}{\vdots}$}\noLine
\UIC{$\forall x\,\forall y\,(x = y \to y = x)$}
\AXC{$\varphi_{1}$\lbl{A2}}
\BIC{$\forall y\,(S(n) = y \to y = S(n))$}
\AXC{$\varphi_{3}$\lbl{A2}}
\BIC{$S(n) = 1_{\omega} \to 1_{\omega} = S(n)$}
\AXC{$\stackrel{\Pi_{54}}{\vdots}$}\noLine
\UIC{$(S(n) = 1_{\omega} \to 1_{\omega} = S(n)) \to r \in \mathbb{R} \to \varphi_{93}$}
\BIC{$r \in \mathbb{R} \to \varphi_{93}$}
\end{ptree}
\noindent$\Pi_{56}$:
\begin{ptree}
\AXC{$\stackrel{\Pi_{39}}{\vdots}$}\noLine
\UIC{$1_{\omega} = S(n) \to \varphi_{50}$}
\AXC{$\stackrel{\Pi_{48}}{\vdots}$}\noLine
\UIC{$r \in \mathbb{R} \to r \neq 0 \to {{(r^{-1})}^{1_{\omega}}} = {{r}^{1_{\omega}}}^{-1}$}
\AXC{$\stackrel{\Pi_{55}}{\vdots}$}\noLine
\UIC{$r \in \mathbb{R} \to \varphi_{93}$}
\AXC{$(r \in \mathbb{R} \to \varphi_{93}) \to \varphi_{140}$\lbl{A1}}
\BIC{$\varphi_{140}$}
\BIC{$(1_{\omega} = S(n) \to \varphi_{50}) \to r \in \mathbb{R} \to r \neq 0 \to \varphi_{32} \to \varphi_{39}$}
\BIC{$r \in \mathbb{R} \to r \neq 0 \to \varphi_{32} \to \varphi_{39}$}
\end{ptree}
\noindent$\Pi_{57}$:
\begin{ptree}
\AXC{$\forall x\,(\mathrm{Ind}(x) \leftrightarrow \emptyset \in x \land \forall y \in x\,(S(y) \in x))$\lbl{DefInd}}
\AXC{$\varphi_{87}$\lbl{A2}}
\BIC{$\mathrm{Ind}(\omega) \leftrightarrow \emptyset \in \omega \land \forall y \in \omega\,(S(y) \in \omega)$}
\AXC{$\mathrm{Ind}(\omega)$\lbl{補題(ωは帰納的)}}
\AXC{$\mathrm{Ind}(\omega) \to \varphi_{124}$\lbl{A1}}
\BIC{$\varphi_{124}$}
\BIC{$r \in \mathbb{R} \to r \neq 0 \to \varphi_{32} \to \emptyset \in \omega \land \forall y \in \omega\,(S(y) \in \omega)$}
\end{ptree}
\noindent$\Pi_{58}$:
\begin{ptree}
\AXC{$\forall y \in \omega\,(S(y) \in \omega) \to n \in \omega \to S(n) \in \omega$\lbl{A2}}
\AXC{$\stackrel{\Pi_{57}}{\vdots}$}\noLine
\UIC{$r \in \mathbb{R} \to r \neq 0 \to \varphi_{32} \to \emptyset \in \omega \land \forall y \in \omega\,(S(y) \in \omega)$}
\AXC{$\varphi_{132}$\lbl{A1}}
\BIC{$r \in \mathbb{R} \to r \neq 0 \to \varphi_{32} \to n \in B_{40} \to \forall y \in \omega\,(S(y) \in \omega)$}
\AXC{$\varphi_{143}$\lbl{A1}}
\BIC{$\varphi_{121}$}
\BIC{$r \in \mathbb{R} \to r \neq 0 \to \varphi_{32} \to n \in B_{40} \to n \in \omega \to S(n) \in \omega$}
\end{ptree}
\noindent$\Pi_{59}$:
\begin{ptree}
\AXC{$\stackrel{\Pi_{10}}{\vdots}$}\noLine
\UIC{$r \in \mathbb{R} \to r \neq 0 \to \varphi_{32} \to n \in B_{40} \to n \in \omega$}
\AXC{$\stackrel{\Pi_{58}}{\vdots}$}\noLine
\UIC{$r \in \mathbb{R} \to r \neq 0 \to \varphi_{32} \to n \in B_{40} \to n \in \omega \to S(n) \in \omega$}
\AXC{$\varphi_{150}$\lbl{A1}}
\BIC{$\varphi_{128}$}
\BIC{$r \in \mathbb{R} \to r \neq 0 \to \varphi_{32} \to n \in B_{40} \to S(n) \in \omega$}
\end{ptree}
\noindent$\Pi_{60}$:
\begin{ptree}
\AXC{$\stackrel{\Pi_{59}}{\vdots}$}\noLine
\UIC{$r \in \mathbb{R} \to r \neq 0 \to \varphi_{32} \to n \in B_{40} \to S(n) \in \omega$}
\AXC{$\varphi_{32} \to \varphi_{55}$\lbl{A2}}
\AXC{$r \in \mathbb{R} \to r \neq 0 \to \varphi_{32} \to \varphi_{32}$\lbl{A1}}
\AXC{$\varphi_{163}$\lbl{A1}}
\BIC{$(\varphi_{32} \to \varphi_{55}) \to r \in \mathbb{R} \to r \neq 0 \to \varphi_{32} \to \varphi_{55}$}
\BIC{$r \in \mathbb{R} \to r \neq 0 \to \varphi_{32} \to \varphi_{55}$}
\AXC{$\varphi_{168}$\lbl{A1}}
\BIC{$(r \in \mathbb{R} \to r \neq 0 \to \varphi_{32} \to n \in B_{40} \to S(n) \in \omega) \to \varphi_{130}$}
\BIC{$\varphi_{130}$}
\end{ptree}
\noindent$\Pi_{61}$:
\begin{ptree}
\AXC{$\varphi_{129}$\lbl{A3}}
\AXC{$\stackrel{\Pi_{38}}{\vdots}$}\noLine
\UIC{$r \in \mathbb{R} \to r \neq 0 \to \varphi_{32} \to \varphi_{42}$}
\AXC{$\stackrel{\Pi_{56}}{\vdots}$}\noLine
\UIC{$r \in \mathbb{R} \to r \neq 0 \to \varphi_{32} \to \varphi_{39}$}
\AXC{$\varphi_{161}$\lbl{A1}}
\BIC{$(r \in \mathbb{R} \to r \neq 0 \to \varphi_{32} \to \varphi_{42}) \to r \in \mathbb{R} \to r \neq 0 \to \varphi_{97}$}
\BIC{$r \in \mathbb{R} \to r \neq 0 \to \varphi_{97}$}
\AXC{$\stackrel{\Pi_{60}}{\vdots}$}\noLine
\UIC{$\varphi_{130}$}
\BIC{$r \in \mathbb{R} \to r \neq 0 \to \varphi_{32} \to n \in B_{40} \to S(n) \in B_{40}$}
\AXC{$\varphi_{127}$\lbl{A1}}
\BIC{$r \in \mathbb{R} \land r \neq 0 \land \varphi_{32} \to n \in B_{40} \to S(n) \in B_{40}$}
\UIC{$\forall n\,(r \in \mathbb{R} \land r \neq 0 \land \varphi_{32} \to n \in B_{40} \to S(n) \in B_{40})$}
\AXC{$\varphi_{173}$\lbl{A1}}
\BIC{$\varphi_{129} \to \varphi_{137}$}
\BIC{$\varphi_{137}$}
\end{ptree}
\noindent$\Pi_{62}$:
\begin{ptree}
\AXC{$\forall n\,(n \in \mathbb{N} \leftrightarrow n \in \omega \land n \neq \emptyset)$\lbl{Def$\mathbb{N}$}}
\AXC{$\varphi_{71}$\lbl{A2}}
\BIC{$\emptyset \in \mathbb{N} \leftrightarrow \emptyset \in \omega \land \emptyset \neq \emptyset$}
\AXC{$(\emptyset \in \mathbb{N} \leftrightarrow \emptyset \in \omega \land \emptyset \neq \emptyset) \to \emptyset = \emptyset \to \varphi_{53}$\lbl{A1}}
\BIC{$\emptyset = \emptyset \to \varphi_{53}$}
\end{ptree}
\noindent$\Pi_{63}$:
\begin{ptree}
\AXC{$\varphi_{32} \to \varphi_{79}$\lbl{A2}}
\AXC{$r \in \mathbb{R} \to r \neq 0 \to \varphi_{32} \to \varphi_{32}$\lbl{A1}}
\AXC{$\varphi_{165}$\lbl{A1}}
\BIC{$(\varphi_{32} \to \varphi_{79}) \to r \in \mathbb{R} \to r \neq 0 \to \varphi_{32} \to \varphi_{79}$}
\BIC{$r \in \mathbb{R} \to r \neq 0 \to \varphi_{32} \to \varphi_{79}$}
\AXC{$(r \in \mathbb{R} \to r \neq 0 \to \varphi_{32} \to \varphi_{79}) \to \varphi_{145}$\lbl{A1}}
\BIC{$\varphi_{145}$}
\end{ptree}
\noindent$\Pi_{64}$:
\begin{ptree}
\AXC{$x = x$\lbl{A5}}
\UIC{$\forall x\,(x = x)$}
\AXC{$\forall x\,(x = x) \to \emptyset = \emptyset$\lbl{A2}}
\BIC{$\emptyset = \emptyset$}
\AXC{$\stackrel{\Pi_{62}}{\vdots}$}\noLine
\UIC{$\emptyset = \emptyset \to \varphi_{53}$}
\BIC{$\varphi_{53}$}
\AXC{$\stackrel{\Pi_{57}}{\vdots}$}\noLine
\UIC{$r \in \mathbb{R} \to r \neq 0 \to \varphi_{32} \to \emptyset \in \omega \land \forall y \in \omega\,(S(y) \in \omega)$}
\AXC{$\stackrel{\Pi_{63}}{\vdots}$}\noLine
\UIC{$\varphi_{145}$}
\BIC{$\varphi_{53} \to r \in \mathbb{R} \to r \neq 0 \to \varphi_{32} \to \emptyset \in B_{40}$}
\BIC{$r \in \mathbb{R} \to r \neq 0 \to \varphi_{32} \to \emptyset \in B_{40}$}
\end{ptree}
\noindent$\Pi_{65}$:
\begin{ptree}
\AXC{$\forall x\,(\mathrm{Ind}(x) \leftrightarrow \emptyset \in x \land \forall y \in x\,(S(y) \in x))$\lbl{DefInd}}
\AXC{$\forall x\,(\mathrm{Ind}(x) \leftrightarrow \emptyset \in x \land \forall y \in x\,(S(y) \in x)) \to \varphi_{13}$\lbl{A2}}
\BIC{$\varphi_{13}$}
\end{ptree}
\noindent$\Pi_{66}$:
\begin{ptree}
\AXC{$\stackrel{\Pi_{64}}{\vdots}$}\noLine
\UIC{$r \in \mathbb{R} \to r \neq 0 \to \varphi_{32} \to \emptyset \in B_{40}$}
\AXC{$\stackrel{\Pi_{65}}{\vdots}$}\noLine
\UIC{$\varphi_{13}$}
\AXC{$\varphi_{13} \to (r \in \mathbb{R} \to r \neq 0 \to \varphi_{32} \to \emptyset \in B_{40}) \to \varphi_{144}$\lbl{A1}}
\BIC{$(r \in \mathbb{R} \to r \neq 0 \to \varphi_{32} \to \emptyset \in B_{40}) \to \varphi_{144}$}
\BIC{$\varphi_{144}$}
\end{ptree}
\noindent$\Pi_{67}$:
\begin{ptree}
\AXC{$\stackrel{\Pi_{3}}{\vdots}$}\noLine
\UIC{$\varphi_{11}$}
\AXC{$\stackrel{\Pi_{4}}{\vdots}$}\noLine
\UIC{$\varphi_{12}$}
\AXC{$r \in \mathbb{R} \land r \neq 0 \land \varphi_{32} \to \forall n\,(r \in \mathbb{R} \land r \neq 0 \land \varphi_{32})$\lbl{A4}}
\AXC{$\stackrel{\Pi_{61}}{\vdots}$}\noLine
\UIC{$\varphi_{137}$}
\BIC{$r \in \mathbb{R} \to r \neq 0 \to \varphi_{32} \to \forall n \in B_{40}\,(S(n) \in B_{40})$}
\AXC{$\stackrel{\Pi_{66}}{\vdots}$}\noLine
\UIC{$\varphi_{144}$}
\BIC{$\varphi_{12} \to \varphi_{11} \to r \in \mathbb{R} \to r \neq 0 \to \varphi_{32} \to \mathrm{Ind}(B_{40})$}
\BIC{$\varphi_{11} \to r \in \mathbb{R} \to r \neq 0 \to \varphi_{32} \to \mathrm{Ind}(B_{40})$}
\BIC{$r \in \mathbb{R} \to r \neq 0 \to \varphi_{32} \to \mathrm{Ind}(B_{40})$}
\end{ptree}
\noindent$\Pi_{68}$:
\begin{ptree}
\AXC{$r \in \mathbb{R} \to r \neq 0 \to \varphi_{32} \to n \in \omega \to n \in \omega$\lbl{A1}}
\AXC{$(r \in \mathbb{R} \to r \neq 0 \to \varphi_{32} \to n \in \omega \to n \in \omega) \to \varphi_{117}$\lbl{A1}}
\BIC{$\varphi_{117}$}
\end{ptree}
\noindent$\Pi_{69}$:
\begin{ptree}
\AXC{$\forall n\,(n \in \omega \leftrightarrow \forall x\,(\mathrm{Ind}(x) \to n \in x))$\lbl{Def$\omega$}}
\AXC{$\varphi_{26}$\lbl{A2}}
\BIC{$n \in \omega \leftrightarrow \forall x\,(\mathrm{Ind}(x) \to n \in x)$}
\AXC{$\stackrel{\Pi_{68}}{\vdots}$}\noLine
\UIC{$\varphi_{117}$}
\BIC{$r \in \mathbb{R} \to r \neq 0 \to \varphi_{32} \to n \in \omega \to \forall x\,(\mathrm{Ind}(x) \to n \in x)$}
\AXC{$\varphi_{134}$\lbl{A1}}
\BIC{$\varphi_{118}$}
\end{ptree}
\noindent$\Pi_{70}$:
\begin{ptree}
\AXC{$\stackrel{\Pi_{67}}{\vdots}$}\noLine
\UIC{$r \in \mathbb{R} \to r \neq 0 \to \varphi_{32} \to \mathrm{Ind}(B_{40})$}
\AXC{$\forall x\,(\mathrm{Ind}(x) \to n \in x) \to \mathrm{Ind}(B_{40}) \to n \in B_{40}$\lbl{A2}}
\AXC{$\stackrel{\Pi_{69}}{\vdots}$}\noLine
\UIC{$\varphi_{118}$}
\BIC{$r \in \mathbb{R} \to r \neq 0 \to \varphi_{32} \to n \in \omega \to \mathrm{Ind}(B_{40}) \to n \in B_{40}$}
\AXC{$\varphi_{166}$\lbl{A1}}
\BIC{$\varphi_{156}$}
\BIC{$(r \in \mathbb{R} \to r \neq 0 \to \varphi_{103}) \to \varphi_{110}$}
\end{ptree}
\noindent$\Pi_{71}$:
\begin{ptree}
\AXC{$\forall n\,(\varphi_{111}) \to \varphi_{116}$\lbl{A3}}
\AXC{$\stackrel{\Pi_{2}}{\vdots}$}\noLine
\UIC{$r \in \mathbb{R} \to r \neq 0 \to \varphi_{103}$}
\AXC{$\stackrel{\Pi_{70}}{\vdots}$}\noLine
\UIC{$(r \in \mathbb{R} \to r \neq 0 \to \varphi_{103}) \to \varphi_{110}$}
\BIC{$\varphi_{110}$}
\AXC{$\varphi_{110} \to \varphi_{111}$\lbl{A1}}
\BIC{$\varphi_{111}$}
\UIC{$\forall n\,(\varphi_{111})$}
\AXC{$\forall n\,(\varphi_{111}) \to \varphi_{171}$\lbl{A1}}
\BIC{$\varphi_{171}$}
\BIC{$(r \in \mathbb{R} \land r \neq 0 \land \varphi_{32} \to \forall n\,(r \in \mathbb{R} \land r \neq 0 \land \varphi_{32})) \to \varphi_{115}$}
\end{ptree}
\noindent$\Pi_{72}$:
\begin{ptree}
\AXC{$r \in \mathbb{R} \land r \neq 0 \land \varphi_{32} \to \forall n\,(r \in \mathbb{R} \land r \neq 0 \land \varphi_{32})$\lbl{A4}}
\AXC{$\stackrel{\Pi_{71}}{\vdots}$}\noLine
\UIC{$(r \in \mathbb{R} \land r \neq 0 \land \varphi_{32} \to \forall n\,(r \in \mathbb{R} \land r \neq 0 \land \varphi_{32})) \to \varphi_{115}$}
\BIC{$\varphi_{115}$}
\AXC{$\varphi_{115} \to \varphi_{49} \to \neg \varphi_{32}$\lbl{A1}}
\BIC{$\varphi_{49} \to \neg \varphi_{32}$}
\UIC{$\forall B_{40}\,(\varphi_{49} \to \neg \varphi_{32})$}
\end{ptree}
\noindent$\Pi_{73}$:
\begin{ptree}
\AXC{$\varphi_{64}$\lbl{Sep}}
\AXC{$\varphi_{64} \to \varphi_{44}$\lbl{A2}}
\BIC{$\varphi_{44}$}
\AXC{$\varphi_{44} \to \exists B_{40}\,(\varphi_{32})$\lbl{A2}}
\BIC{$\exists B_{40}\,(\varphi_{32})$}
\AXC{$\exists B_{40}\,(\varphi_{32}) \to r \in \mathbb{R} \to r \neq 0 \to \exists B_{40}\,(\varphi_{32})$\lbl{A1}}
\BIC{$r \in \mathbb{R} \to r \neq 0 \to \exists B_{40}\,(\varphi_{32})$}
\AXC{$\varphi_{175}$\lbl{A1}}
\BIC{$\forall B_{40}\,(\varphi_{49} \to \neg \varphi_{32}) \to \varphi_{169}$}
\end{ptree}
\noindent$\Pi_{74}$:
\begin{ptree}
\AXC{$\forall B_{40}\,(\varphi_{49} \to \neg \varphi_{32}) \to \forall B_{40}\,(\varphi_{49}) \to \forall B_{40}\,(\neg \varphi_{32})$\lbl{A3}}
\AXC{$\stackrel{\Pi_{72}}{\vdots}$}\noLine
\UIC{$\forall B_{40}\,(\varphi_{49} \to \neg \varphi_{32})$}
\AXC{$\stackrel{\Pi_{73}}{\vdots}$}\noLine
\UIC{$\forall B_{40}\,(\varphi_{49} \to \neg \varphi_{32}) \to \varphi_{169}$}
\BIC{$\varphi_{169}$}
\BIC{$(\varphi_{49} \to \forall B_{40}\,(\varphi_{49})) \to \varphi_{45}$}
\end{ptree}
\noindent$\Pi_{75}$:
\begin{ptree}
\AXC{$\varphi_{90}$\lbl{A2}}
\AXC{$\varphi_{49} \to \forall B_{40}\,(\varphi_{49})$\lbl{A4}}
\AXC{$\stackrel{\Pi_{74}}{\vdots}$}\noLine
\UIC{$(\varphi_{49} \to \forall B_{40}\,(\varphi_{49})) \to \varphi_{45}$}
\BIC{$\varphi_{45}$}
\AXC{$\varphi_{45} \to \varphi_{122}$\lbl{A1}}
\BIC{$\varphi_{122}$}
\BIC{$r \in \mathbb{R} \to r \neq 0 \to n \in \omega \to n \in \mathbb{N} \to {{(r^{-1})}^{n}} = {{r}^{n}}^{-1}$}
\AXC{$\varphi_{126}$\lbl{A1}}
\BIC{$(r \in \mathbb{R} \to r \neq 0 \to n \in \mathbb{N} \to n \in \omega) \to \varphi_{80}$}
\end{ptree}
\noindent$\Pi_{76}$:
\begin{ptree}
\AXC{$\forall n\,(n \in \mathbb{N} \leftrightarrow n \in \omega \land n \neq \emptyset)$\lbl{Def$\mathbb{N}$}}
\AXC{$\varphi_{35}$\lbl{A2}}
\BIC{$n \in \mathbb{N} \leftrightarrow n \in \omega \land n \neq \emptyset$}
\AXC{$\stackrel{\Pi_{1}}{\vdots}$}\noLine
\UIC{$\varphi_{43}$}
\BIC{$r \in \mathbb{R} \to r \neq 0 \to n \in \mathbb{N} \to n \in \omega$}
\AXC{$\stackrel{\Pi_{75}}{\vdots}$}\noLine
\UIC{$(r \in \mathbb{R} \to r \neq 0 \to n \in \mathbb{N} \to n \in \omega) \to \varphi_{80}$}
\BIC{$\varphi_{80}$}
\end{ptree}
\noindent$\Pi_{77}$:
\begin{ptree}
\AXC{$\varphi_{96}$\lbl{A3}}
\AXC{$r \in \mathbb{R} \to r \neq 0 \to n \in \mathbb{N} \to n \in \mathbb{N}$\lbl{A1}}
\AXC{$\stackrel{\Pi_{76}}{\vdots}$}\noLine
\UIC{$\varphi_{80}$}
\BIC{$r \in \mathbb{R} \to r \neq 0 \to n \in \mathbb{N} \to {{(r^{-1})}^{n}} = {{r}^{n}}^{-1}$}
\AXC{$\varphi_{94}$\lbl{A1}}
\BIC{$r \in \mathbb{R} \land r \neq 0 \to n \in \mathbb{N} \to {{(r^{-1})}^{n}} = {{r}^{n}}^{-1}$}
\UIC{$\forall n\,(r \in \mathbb{R} \land r \neq 0 \to n \in \mathbb{N} \to {{(r^{-1})}^{n}} = {{r}^{n}}^{-1})$}
\AXC{$\varphi_{139}$\lbl{A1}}
\BIC{$\varphi_{96} \to \varphi_{84}$}
\BIC{$\varphi_{84}$}
\end{ptree}
\noindent 主木:
\begin{ptree}
\AXC{$r \in \mathbb{R} \land r \neq 0 \to \forall n\,(r \in \mathbb{R} \land r \neq 0)$\lbl{A4}}
\AXC{$\stackrel{\Pi_{77}}{\vdots}$}\noLine
\UIC{$\varphi_{84}$}
\BIC{$r \in \mathbb{R} \to r \neq 0 \to \forall n \in \mathbb{N}\,({{(r^{-1})}^{n}} = {{r}^{n}}^{-1})$}
\UIC{$\forall r \in \mathbb{R}\,(r \neq 0 \to \forall n \in \mathbb{N}\,({{(r^{-1})}^{n}} = {{r}^{n}}^{-1}))$}
\end{ptree}
\begin{abbrevlist}
\item $\varphi_{1} :\equiv \forall x\,\forall y\,(x = y \to y = x) \to \forall y\,(S(n) = y \to y = S(n))$
\item $\varphi_{2} :\equiv \forall x\,\forall y\,(x = y \to y = x) \to \forall y\,(1_{\omega} = y \to y = 1_{\omega})$
\item $\varphi_{3} :\equiv \forall y\,(S(n) = y \to y = S(n)) \to S(n) = 1_{\omega} \to 1_{\omega} = S(n)$
\item $\varphi_{4} :\equiv S(n) = S(\emptyset) \to S(\emptyset) = 1_{\omega} \to S(n) = 1_{\omega}$
\item $\varphi_{5} :\equiv n \in B_{40} \to S(n) \in \mathbb{N} \to n \neq \emptyset \to \mathrm{Seq}({(r)^{\bullet}})$
\item $\varphi_{6} :\equiv \forall n \in B_{40}\,(S(n) \in B_{40}) \to y \in B_{40} \to S(y) \in B_{40}$
\item $\varphi_{7} :\equiv \forall y \in B_{40}\,(S(y) \in B_{40}) \to n \in B_{40} \to S(n) \in B_{40}$
\item $\varphi_{8} :\equiv \forall x\,\forall y\,\forall z\,(x = y \to y = z \to x = z) \to \forall y\,\forall z\,(S(n) = y \to y = z \to S(n) = z)$
\item $\varphi_{9} :\equiv (1_{\omega} = S(\emptyset) \to S(\emptyset) = 1_{\omega}) \to S(\emptyset) = 1_{\omega}$
\item $\varphi_{10} :\equiv \forall y\,(1_{\omega} = y \to y = 1_{\omega}) \to 1_{\omega} = S(\emptyset) \to S(\emptyset) = 1_{\omega}$
\item $\varphi_{11} :\equiv \forall n \in B_{40}\,(S(n) \in B_{40}) \to \forall y \in B_{40}\,(S(y) \in B_{40})$
\item $\varphi_{12} :\equiv \forall y \in B_{40}\,(S(y) \in B_{40}) \to \forall n \in B_{40}\,(S(n) \in B_{40})$
\item $\varphi_{13} :\equiv \mathrm{Ind}(B_{40}) \leftrightarrow \emptyset \in B_{40} \land \forall y \in B_{40}\,(S(y) \in B_{40})$
\item $\varphi_{14} :\equiv \forall x\,\forall y\,(x = y \to y = x) \to \forall y\,({{r}^{S(n)}} = y \to y = {{r}^{S(n)}})$
\item $\varphi_{15} :\equiv \forall x\,\forall y\,(x = y \to y = x) \to \forall y\,({{r}^{1_{\omega}}} = y \to y = {{r}^{1_{\omega}}})$
\item $\varphi_{16} :\equiv \forall y\,\forall z\,({{(r^{-1})}^{S(n)}} = y \to y = z \to {{(r^{-1})}^{S(n)}} = z)$
\item $\varphi_{17} :\equiv \forall y\,\forall z\,({{(r^{-1})}^{1_{\omega}}} = y \to y = z \to {{(r^{-1})}^{1_{\omega}}} = z)$
\item $\varphi_{18} :\equiv \forall n \in B_{40}\,(S(n) \in B_{40}) \to \forall y\,\forall n \in B_{40}\,(S(n) \in B_{40})$
\item $\varphi_{19} :\equiv \forall y \in B_{40}\,(S(y) \in B_{40}) \to \forall n\,\forall y \in B_{40}\,(S(y) \in B_{40})$
\item $\varphi_{20} :\equiv \forall y\,\forall n \in B_{40}\,(S(n) \in B_{40}) \to \forall y \in B_{40}\,(S(y) \in B_{40})$
\item $\varphi_{21} :\equiv \forall n\,\forall y \in B_{40}\,(S(y) \in B_{40}) \to \forall n \in B_{40}\,(S(n) \in B_{40})$
\item $\varphi_{22} :\equiv \forall z\,({{(r^{-1})}^{1_{\omega}}} = r^{-1} \to r^{-1} = z \to {{(r^{-1})}^{1_{\omega}}} = z)$
\item $\varphi_{23} :\equiv \forall z\,\forall w\,(z = w \to S(n) = S(z) \to S(n) = S(w)) \to \forall w\,(n = w \to S(n) = S(n) \to S(n) = S(w))$
\item $\varphi_{24} :\equiv \forall w\,(z = w \to ({{r}^{n}} \cdot r)^{-1} = z^{-1} \to ({{r}^{n}} \cdot r)^{-1} = w^{-1})$
\item $\varphi_{25} :\equiv \forall x\,\forall y\,(x = y \to y = x) \to \forall y\,(({{r}^{n}} \cdot r)^{-1} = y \to y = ({{r}^{n}} \cdot r)^{-1})$
\item $\varphi_{26} :\equiv \forall n\,(n \in \omega \leftrightarrow \forall x\,(\mathrm{Ind}(x) \to n \in x)) \to (n \in \omega \leftrightarrow \forall x\,(\mathrm{Ind}(x) \to n \in x))$
\item $\varphi_{27} :\equiv \forall n\,(r \in \mathbb{R} \land r \neq 0) \to \forall n \in \mathbb{N}\,({{(r^{-1})}^{n}} = {{r}^{n}}^{-1})$
\item $\varphi_{28} :\equiv \forall x\,\forall y\,(x \in \mathbb{R} \to y \in \mathbb{R} \to x \neq 0 \to y \neq 0 \to (x \cdot y)^{-1} = x^{-1} \cdot y^{-1})$
\item $\varphi_{29} :\equiv \forall y\,\forall z\,(S(n) = y \to y = z \to S(n) = z) \to \forall z\,(S(n) = S(\emptyset) \to S(\emptyset) = z \to S(n) = z)$
\item $\varphi_{30} :\equiv \forall w\,(n = w \to S(n) = S(n) \to S(n) = S(w)) \to n = \emptyset \to S(n) = S(n) \to S(n) = S(\emptyset)$
\item $\varphi_{31} :\equiv \forall x \in \mathbb{R}\,(x \neq 0 \to x^{-1} \in \mathbb{R}) \to r \in \mathbb{R} \to r \neq 0 \to r^{-1} \in \mathbb{R}$
\item $\varphi_{32} :\equiv \forall n\,(n \in B_{40} \leftrightarrow n \in \omega \land (n \in \mathbb{N} \to {{(r^{-1})}^{n}} = {{r}^{n}}^{-1}))$
\item $\varphi_{33} :\equiv ({{r}^{1_{\omega}}} = r \to r = {{r}^{1_{\omega}}}) \to r \in \mathbb{R} \to r \neq 0 \to r = {{r}^{1_{\omega}}}$
\item $\varphi_{34} :\equiv S(n) \in \mathbb{N} \to n \neq \emptyset \to {{(r^{-1})}^{S(n)}} = {{(r^{-1})}^{n}} \cdot r^{-1}$
\item $\varphi_{35} :\equiv \forall n\,(n \in \mathbb{N} \leftrightarrow n \in \omega \land n \neq \emptyset) \to (n \in \mathbb{N} \leftrightarrow n \in \omega \land n \neq \emptyset)$
\item $\varphi_{36} :\equiv \forall y\,({{r}^{1_{\omega}}} = y \to y = {{r}^{1_{\omega}}}) \to {{r}^{1_{\omega}}} = r \to r = {{r}^{1_{\omega}}}$
\item $\varphi_{37} :\equiv \mathrm{Seq}({(r)^{\bullet}}) \land {{r}^{1_{\omega}}} = r \land \forall n \in \mathbb{N}\,({{r}^{S(n)}} = {{r}^{n}} \cdot r)$
\item $\varphi_{38} :\equiv {{(r^{-1})}^{n}} \cdot r^{-1} = z \cdot r^{-1} \to {{(r^{-1})}^{n}} \cdot r^{-1} = w \cdot r^{-1}$
\item $\varphi_{39} :\equiv n \in B_{40} \to S(n) \in \mathbb{N} \to n = \emptyset \to {{(r^{-1})}^{S(n)}} = {{r}^{S(n)}}^{-1}$
\item $\varphi_{40} :\equiv \forall w\,(z = w \to {{(r^{-1})}^{z}} = {{r}^{z}}^{-1} \to {{(r^{-1})}^{w}} = {{r}^{w}}^{-1})$
\item $\varphi_{41} :\equiv \forall z\,\forall w\,(z = w \to r^{-1} = z^{-1} \to r^{-1} = w^{-1}) \to \forall w\,(r = w \to r^{-1} = r^{-1} \to r^{-1} = w^{-1})$
\item $\varphi_{42} :\equiv n \in B_{40} \to S(n) \in \mathbb{N} \to n \neq \emptyset \to {{(r^{-1})}^{S(n)}} = {{r}^{S(n)}}^{-1}$
\item $\varphi_{43} :\equiv (n \in \mathbb{N} \leftrightarrow n \in \omega \land n \neq \emptyset) \to r \in \mathbb{R} \to r \neq 0 \to n \in \mathbb{N} \to n \in \omega$
\item $\varphi_{44} :\equiv \forall A_{41}\,\exists B_{40}\,\forall n\,(n \in B_{40} \leftrightarrow n \in A_{41} \land (n \in \mathbb{N} \to {{(r^{-1})}^{n}} = {{r}^{n}}^{-1}))$
\item $\varphi_{45} :\equiv r \in \mathbb{R} \to r \neq 0 \to \forall n \in \omega\,(n \in \mathbb{N} \to {{(r^{-1})}^{n}} = {{r}^{n}}^{-1})$
\item $\varphi_{46} :\equiv n \in B_{40} \to S(n) \in \mathbb{N} \to n \neq \emptyset \to ({{r}^{n}} \cdot r)^{-1} = {{r}^{S(n)}}^{-1}$
\item $\varphi_{47} :\equiv S(n) \in \mathbb{N} \to n \neq \emptyset \to {{(r^{-1})}^{n}} \cdot r^{-1} = {{r}^{n}}^{-1} \cdot r^{-1}$
\item $\varphi_{48} :\equiv n \in B_{40} \to S(n) \in \mathbb{N} \to n \neq \emptyset \to {{(r^{-1})}^{S(n)}} = ({{r}^{n}} \cdot r)^{-1}$
\item $\varphi_{49} :\equiv r \in \mathbb{R} \land r \neq 0 \land \neg \forall n \in \omega\,(n \in \mathbb{N} \to {{(r^{-1})}^{n}} = {{r}^{n}}^{-1})$
\item $\varphi_{50} :\equiv {{(r^{-1})}^{1_{\omega}}} = {{r}^{1_{\omega}}}^{-1} \to {{(r^{-1})}^{S(n)}} = {{r}^{S(n)}}^{-1}$
\item $\varphi_{51} :\equiv n \in B_{40} \to S(n) \in \mathbb{N} \to n \neq \emptyset \to n \in \mathbb{N} \to {{r}^{S(n)}} = {{r}^{n}} \cdot r$
\item $\varphi_{52} :\equiv ({{r}^{n}} \cdot r)^{-1} = {{r}^{S(n)}}^{-1} \to {{(r^{-1})}^{S(n)}} = {{r}^{S(n)}}^{-1}$
\item $\varphi_{53} :\equiv r \in \mathbb{R} \to r \neq 0 \to \emptyset \in \mathbb{N} \to {{(r^{-1})}^{\emptyset}} = {{r}^{\emptyset}}^{-1}$
\item $\varphi_{54} :\equiv r \in \mathbb{R} \to {{r}^{n}} \neq 0 \to r \neq 0 \to ({{r}^{n}} \cdot r)^{-1} = {{r}^{n}}^{-1} \cdot r^{-1}$
\item $\varphi_{55} :\equiv S(n) \in B_{40} \leftrightarrow S(n) \in \omega \land (S(n) \in \mathbb{N} \to {{(r^{-1})}^{S(n)}} = {{r}^{S(n)}}^{-1})$
\item $\varphi_{56} :\equiv n \in B_{40} \to S(n) \in \mathbb{N} \to n \neq \emptyset \to \forall n \in \mathbb{N}\,({{r}^{S(n)}} = {{r}^{n}} \cdot r)$
\item $\varphi_{57} :\equiv S(n) \in \mathbb{N} \to n \neq \emptyset \to n \in \mathbb{N} \to {{(r^{-1})}^{S(n)}} = {{(r^{-1})}^{n}} \cdot r^{-1}$
\item $\varphi_{58} :\equiv n \in B_{40} \to S(n) \in \mathbb{N} \to n \neq \emptyset \to {{(r^{-1})}^{S(n)}} = {{r}^{n}}^{-1} \cdot r^{-1}$
\item $\varphi_{59} :\equiv \forall w\,(1_{\omega} = w \to {{(r^{-1})}^{1_{\omega}}} = {{r}^{1_{\omega}}}^{-1} \to {{(r^{-1})}^{w}} = {{r}^{w}}^{-1})$
\item $\varphi_{60} :\equiv ({{r}^{n}} \cdot r)^{-1} = ({{r}^{n}} \cdot r)^{-1} \to ({{r}^{n}} \cdot r)^{-1} = {{r}^{S(n)}}^{-1}$
\item $\varphi_{61} :\equiv n \in B_{40} \to S(n) \in \mathbb{N} \to n \neq \emptyset \to ({{r}^{n}} \cdot r)^{-1} = {{r}^{n}}^{-1} \cdot r^{-1}$
\item $\varphi_{62} :\equiv n \in B_{40} \to S(n) \in \mathbb{N} \to n \neq \emptyset \to {{r}^{n}}^{-1} \cdot r^{-1} = ({{r}^{n}} \cdot r)^{-1}$
\item $\varphi_{63} :\equiv S(n) \in \mathbb{N} \to n \neq \emptyset \to \forall n \in \mathbb{N}\,({{(r^{-1})}^{S(n)}} = {{(r^{-1})}^{n}} \cdot r^{-1})$
\item $\varphi_{64} :\equiv \forall z_{42}\,\forall A_{41}\,\exists B_{40}\,\forall n\,(n \in B_{40} \leftrightarrow n \in A_{41} \land (n \in \mathbb{N} \to {{(z_{42}^{-1})}^{n}} = {{z_{42}}^{n}}^{-1}))$
\item $\varphi_{65} :\equiv \forall n \in \mathbb{N}\,({{r}^{S(n)}} = {{r}^{n}} \cdot r) \to n \in \mathbb{N} \to {{r}^{S(n)}} = {{r}^{n}} \cdot r$
\item $\varphi_{66} :\equiv \forall a\,\forall n\,(\mathrm{Seq}(a) \to n \in \mathbb{N} \to a_{n} \in \mathbb{R}) \to \forall n\,(\mathrm{Seq}({(r)^{\bullet}}) \to n \in \mathbb{N} \to {{r}^{n}} \in \mathbb{R})$
\item $\varphi_{67} :\equiv {{r}^{n}}^{-1} \cdot r^{-1} = ({{r}^{n}} \cdot r)^{-1} \to {{(r^{-1})}^{S(n)}} = ({{r}^{n}} \cdot r)^{-1}$
\item $\varphi_{68} :\equiv \forall z\,({{(r^{-1})}^{S(n)}} = ({{r}^{n}} \cdot r)^{-1} \to ({{r}^{n}} \cdot r)^{-1} = z \to {{(r^{-1})}^{S(n)}} = z)$
\item $\varphi_{69} :\equiv \forall n\,(\mathrm{Seq}({(r)^{\bullet}}) \to n \in \mathbb{N} \to {{r}^{n}} \in \mathbb{R}) \to \mathrm{Seq}({(r)^{\bullet}}) \to n \in \mathbb{N} \to {{r}^{n}} \in \mathbb{R}$
\item $\varphi_{70} :\equiv {{(r^{-1})}^{1_{\omega}}} = r^{-1} \to r^{-1} = {{r}^{1_{\omega}}}^{-1} \to {{(r^{-1})}^{1_{\omega}}} = {{r}^{1_{\omega}}}^{-1}$
\item $\varphi_{71} :\equiv \forall n\,(n \in \mathbb{N} \leftrightarrow n \in \omega \land n \neq \emptyset) \to (\emptyset \in \mathbb{N} \leftrightarrow \emptyset \in \omega \land \emptyset \neq \emptyset)$
\item $\varphi_{72} :\equiv \forall w\,({{r}^{n}} \cdot r = w \to ({{r}^{n}} \cdot r)^{-1} = ({{r}^{n}} \cdot r)^{-1} \to ({{r}^{n}} \cdot r)^{-1} = w^{-1})$
\item $\varphi_{73} :\equiv ({{r}^{n}} \cdot r)^{-1} = {{r}^{n}}^{-1} \cdot r^{-1} \to {{r}^{n}}^{-1} \cdot r^{-1} = ({{r}^{n}} \cdot r)^{-1}$
\item $\varphi_{74} :\equiv \forall w\,(r = w \to r^{-1} = r^{-1} \to r^{-1} = w^{-1}) \to r = {{r}^{1_{\omega}}} \to r^{-1} = r^{-1} \to r^{-1} = {{r}^{1_{\omega}}}^{-1}$
\item $\varphi_{75} :\equiv \forall r \in \mathbb{R}\,(r \neq 0 \to \forall n \in \mathbb{N}\,({{r}^{n}} \neq 0)) \to r \in \mathbb{R} \to r \neq 0 \to \forall n \in \mathbb{N}\,({{r}^{n}} \neq 0)$
\item $\varphi_{76} :\equiv {{(r^{-1})}^{n}} \cdot r^{-1} = {{r}^{n}}^{-1} \cdot r^{-1} \to {{(r^{-1})}^{S(n)}} = {{r}^{n}}^{-1} \cdot r^{-1}$
\item $\varphi_{77} :\equiv \forall z\,({{(r^{-1})}^{S(n)}} = {{r}^{n}}^{-1} \cdot r^{-1} \to {{r}^{n}}^{-1} \cdot r^{-1} = z \to {{(r^{-1})}^{S(n)}} = z)$
\item $\varphi_{78} :\equiv \forall z\,({{(r^{-1})}^{S(n)}} = {{(r^{-1})}^{n}} \cdot r^{-1} \to {{(r^{-1})}^{n}} \cdot r^{-1} = z \to {{(r^{-1})}^{S(n)}} = z)$
\item $\varphi_{79} :\equiv \emptyset \in B_{40} \leftrightarrow \emptyset \in \omega \land (\emptyset \in \mathbb{N} \to {{(r^{-1})}^{\emptyset}} = {{r}^{\emptyset}}^{-1})$
\item $\varphi_{80} :\equiv (r \in \mathbb{R} \to r \neq 0 \to n \in \mathbb{N} \to n \in \mathbb{N}) \to r \in \mathbb{R} \to r \neq 0 \to n \in \mathbb{N} \to {{(r^{-1})}^{n}} = {{r}^{n}}^{-1}$
\item $\varphi_{81} :\equiv \mathrm{Seq}({(r^{-1})^{\bullet}}) \land {{(r^{-1})}^{1_{\omega}}} = r^{-1} \land \forall n \in \mathbb{N}\,({{(r^{-1})}^{S(n)}} = {{(r^{-1})}^{n}} \cdot r^{-1})$
\item $\varphi_{82} :\equiv \forall y\,({{r}^{S(n)}} = y \to y = {{r}^{S(n)}}) \to {{r}^{S(n)}} = {{r}^{n}} \cdot r \to {{r}^{n}} \cdot r = {{r}^{S(n)}}$
\item $\varphi_{83} :\equiv (r = {{r}^{1_{\omega}}} \to r^{-1} = r^{-1} \to r^{-1} = {{r}^{1_{\omega}}}^{-1}) \to r \in \mathbb{R} \to r \neq 0 \to r^{-1} = {{r}^{1_{\omega}}}^{-1}$
\item $\varphi_{84} :\equiv (r \in \mathbb{R} \land r \neq 0 \to \forall n\,(r \in \mathbb{R} \land r \neq 0)) \to r \in \mathbb{R} \to r \neq 0 \to \forall n \in \mathbb{N}\,({{(r^{-1})}^{n}} = {{r}^{n}}^{-1})$
\item $\varphi_{85} :\equiv {{(r^{-1})}^{n}} \cdot r^{-1} = {{(r^{-1})}^{n}} \cdot r^{-1} \to {{(r^{-1})}^{n}} \cdot r^{-1} = {{r}^{n}}^{-1} \cdot r^{-1}$
\item $\varphi_{86} :\equiv \forall y\,({{r}^{n}} \in \mathbb{R} \to y \in \mathbb{R} \to {{r}^{n}} \neq 0 \to y \neq 0 \to ({{r}^{n}} \cdot y)^{-1} = {{r}^{n}}^{-1} \cdot y^{-1})$
\item $\varphi_{87} :\equiv \forall x\,(\mathrm{Ind}(x) \leftrightarrow \emptyset \in x \land \forall y \in x\,(S(y) \in x)) \to (\mathrm{Ind}(\omega) \leftrightarrow \emptyset \in \omega \land \forall y \in \omega\,(S(y) \in \omega))$
\item $\varphi_{88} :\equiv \forall w\,({{(r^{-1})}^{n}} = w \to {{(r^{-1})}^{n}} \cdot r^{-1} = {{(r^{-1})}^{n}} \cdot r^{-1} \to {{(r^{-1})}^{n}} \cdot r^{-1} = w \cdot r^{-1})$
\item $\varphi_{89} :\equiv \forall n \in \mathbb{N}\,({{(r^{-1})}^{S(n)}} = {{(r^{-1})}^{n}} \cdot r^{-1}) \to n \in \mathbb{N} \to {{(r^{-1})}^{S(n)}} = {{(r^{-1})}^{n}} \cdot r^{-1}$
\item $\varphi_{90} :\equiv \forall n \in \omega\,(n \in \mathbb{N} \to {{(r^{-1})}^{n}} = {{r}^{n}}^{-1}) \to n \in \omega \to n \in \mathbb{N} \to {{(r^{-1})}^{n}} = {{r}^{n}}^{-1}$
\item $\varphi_{91} :\equiv \forall n\,(r \in \mathbb{R} \land r \neq 0 \land \varphi_{32}) \to \forall n \in B_{40}\,(S(n) \in B_{40})$
\item $\varphi_{92} :\equiv r \neq 0 \to \varphi_{32} \to n \in B_{40} \to S(n) \in \mathbb{N} \to n = \emptyset \to S(n) = 1_{\omega}$
\item $\varphi_{93} :\equiv r \neq 0 \to \varphi_{32} \to n \in B_{40} \to S(n) \in \mathbb{N} \to n = \emptyset \to 1_{\omega} = S(n)$
\item $\varphi_{94} :\equiv (r \in \mathbb{R} \to r \neq 0 \to n \in \mathbb{N} \to {{(r^{-1})}^{n}} = {{r}^{n}}^{-1}) \to r \in \mathbb{R} \land r \neq 0 \to n \in \mathbb{N} \to {{(r^{-1})}^{n}} = {{r}^{n}}^{-1}$
\item $\varphi_{95} :\equiv r \neq 0 \to \varphi_{32} \to n \in B_{40} \to S(n) \in \mathbb{N} \to n = \emptyset \to S(n) = S(\emptyset)$
\item $\varphi_{96} :\equiv \forall n\,(r \in \mathbb{R} \land r \neq 0 \to n \in \mathbb{N} \to {{(r^{-1})}^{n}} = {{r}^{n}}^{-1}) \to \varphi_{27}$
\item $\varphi_{97} :\equiv \varphi_{32} \to n \in B_{40} \to S(n) \in \mathbb{N} \to {{(r^{-1})}^{S(n)}} = {{r}^{S(n)}}^{-1}$
\item $\varphi_{98} :\equiv r \in \mathbb{R} \to r \neq 0 \to \varphi_{32} \to n \in B_{40} \to S(n) \in \mathbb{N} \to n \neq \emptyset \to n \in \mathbb{N}$
\item $\varphi_{99} :\equiv r \in \mathbb{R} \to r \neq 0 \to \varphi_{32} \to n \in B_{40} \to S(n) \in \mathbb{N} \to n = \emptyset \to n = \emptyset$
\item $\varphi_{100} :\equiv r \in \mathbb{R} \to r \neq 0 \to \varphi_{32} \to n \in B_{40} \to S(n) \in \mathbb{N} \to n \neq \emptyset \to r \in \mathbb{R}$
\item $\varphi_{101} :\equiv \varphi_{32} \to n \in B_{40} \to S(n) \in \mathbb{N} \to n \neq \emptyset \to \forall n \in \mathbb{N}\,({{r}^{n}} \neq 0)$
\item $\varphi_{102} :\equiv r \in \mathbb{R} \to r \neq 0 \to \varphi_{32} \to n \in B_{40} \to S(n) \in \mathbb{N} \to n \neq \emptyset \to n \neq \emptyset$
\item $\varphi_{103} :\equiv \varphi_{32} \to (n \in B_{40} \leftrightarrow n \in \omega \land (n \in \mathbb{N} \to {{(r^{-1})}^{n}} = {{r}^{n}}^{-1}))$
\item $\varphi_{104} :\equiv r \neq 0 \to \varphi_{32} \to n \in B_{40} \to S(n) \in \mathbb{N} \to n \neq \emptyset \to n \in \mathbb{N} \to {{r}^{n}} \neq 0$
\item $\varphi_{105} :\equiv r \in \mathbb{R} \to r \neq 0 \to \varphi_{32} \to n \in B_{40} \to S(n) \in \mathbb{N} \to n \neq \emptyset \to {{r}^{n}} \neq 0$
\item $\varphi_{106} :\equiv r \in \mathbb{R} \to r \neq 0 \to \varphi_{32} \to n \in B_{40} \to S(n) \in \mathbb{N} \to n \neq \emptyset \to r^{-1} \in \mathbb{R}$
\item $\varphi_{107} :\equiv \varphi_{32} \to n \in B_{40} \to S(n) \in \mathbb{N} \to n \neq \emptyset \to {{(r^{-1})}^{n}} = {{r}^{n}}^{-1}$
\item $\varphi_{108} :\equiv r \in \mathbb{R} \to r \neq 0 \to \varphi_{32} \to n \in B_{40} \to n \in \mathbb{N} \to {{(r^{-1})}^{n}} = {{r}^{n}}^{-1}$
\item $\varphi_{109} :\equiv r \neq 0 \to \varphi_{32} \to n \in B_{40} \to n \in \omega \land (n \in \mathbb{N} \to {{(r^{-1})}^{n}} = {{r}^{n}}^{-1})$
\item $\varphi_{110} :\equiv r \in \mathbb{R} \to r \neq 0 \to \varphi_{32} \to n \in \omega \to n \in \mathbb{N} \to {{(r^{-1})}^{n}} = {{r}^{n}}^{-1}$
\item $\varphi_{111} :\equiv r \in \mathbb{R} \land r \neq 0 \land \varphi_{32} \to n \in \omega \to n \in \mathbb{N} \to {{(r^{-1})}^{n}} = {{r}^{n}}^{-1}$
\item $\varphi_{112} :\equiv r \neq 0 \to \varphi_{32} \to n \in B_{40} \to S(n) \in \mathbb{N} \to n \neq \emptyset \to {{r}^{S(n)}} = {{r}^{n}} \cdot r$
\item $\varphi_{113} :\equiv r \neq 0 \to \varphi_{32} \to n \in B_{40} \to S(n) \in \mathbb{N} \to n \neq \emptyset \to {{r}^{n}} \cdot r = {{r}^{S(n)}}$
\item $\varphi_{114} :\equiv r \in \mathbb{R} \to r \neq 0 \to \varphi_{32} \to n \in B_{40} \to S(n) \in \mathbb{N} \to n \neq \emptyset \to {{r}^{n}} \in \mathbb{R}$
\item $\varphi_{115} :\equiv r \in \mathbb{R} \to r \neq 0 \to \varphi_{32} \to \forall n \in \omega\,(n \in \mathbb{N} \to {{(r^{-1})}^{n}} = {{r}^{n}}^{-1})$
\item $\varphi_{116} :\equiv \forall n\,(r \in \mathbb{R} \land r \neq 0 \land \varphi_{32}) \to \forall n \in \omega\,(n \in \mathbb{N} \to {{(r^{-1})}^{n}} = {{r}^{n}}^{-1})$
\item $\varphi_{117} :\equiv (n \in \omega \leftrightarrow \forall x\,(\mathrm{Ind}(x) \to n \in x)) \to r \in \mathbb{R} \to r \neq 0 \to \varphi_{32} \to n \in \omega \to \forall x\,(\mathrm{Ind}(x) \to n \in x)$
\item $\varphi_{118} :\equiv (\forall x\,(\mathrm{Ind}(x) \to n \in x) \to \mathrm{Ind}(B_{40}) \to n \in B_{40}) \to r \in \mathbb{R} \to r \neq 0 \to \varphi_{32} \to n \in \omega \to \mathrm{Ind}(B_{40}) \to n \in B_{40}$
\item $\varphi_{119} :\equiv (r \in \mathbb{R} \to r \neq 0 \to \varphi_{81}) \to r \in \mathbb{R} \to r \neq 0 \to {{(r^{-1})}^{1_{\omega}}} = r^{-1}$
\item $\varphi_{120} :\equiv (n = \emptyset \to S(n) = S(n) \to S(n) = S(\emptyset)) \to r \in \mathbb{R} \to \varphi_{95}$
\item $\varphi_{121} :\equiv (\forall y \in \omega\,(S(y) \in \omega) \to n \in \omega \to S(n) \in \omega) \to r \in \mathbb{R} \to r \neq 0 \to \varphi_{32} \to n \in B_{40} \to n \in \omega \to S(n) \in \omega$
\item $\varphi_{122} :\equiv \varphi_{90} \to r \in \mathbb{R} \to r \neq 0 \to n \in \omega \to n \in \mathbb{N} \to {{(r^{-1})}^{n}} = {{r}^{n}}^{-1}$
\item $\varphi_{123} :\equiv (r \in \mathbb{R} \to r \neq 0 \to r^{-1} = {{r}^{1_{\omega}}}^{-1}) \to \varphi_{70} \to r \in \mathbb{R} \to r \neq 0 \to {{(r^{-1})}^{1_{\omega}}} = {{r}^{1_{\omega}}}^{-1}$
\item $\varphi_{124} :\equiv (\mathrm{Ind}(\omega) \leftrightarrow \emptyset \in \omega \land \forall y \in \omega\,(S(y) \in \omega)) \to r \in \mathbb{R} \to r \neq 0 \to \varphi_{32} \to \emptyset \in \omega \land \forall y \in \omega\,(S(y) \in \omega)$
\item $\varphi_{125} :\equiv (r \in \mathbb{R} \to \varphi_{37}) \to (r \in \mathbb{R} \to r \neq 0 \to r \in \mathbb{R}) \to r \in \mathbb{R} \to r \neq 0 \to \varphi_{37}$
\item $\varphi_{126} :\equiv (r \in \mathbb{R} \to r \neq 0 \to n \in \omega \to n \in \mathbb{N} \to {{(r^{-1})}^{n}} = {{r}^{n}}^{-1}) \to (r \in \mathbb{R} \to r \neq 0 \to n \in \mathbb{N} \to n \in \omega) \to \varphi_{80}$
\item $\varphi_{127} :\equiv (r \in \mathbb{R} \to r \neq 0 \to \varphi_{32} \to n \in B_{40} \to S(n) \in B_{40}) \to r \in \mathbb{R} \land r \neq 0 \land \varphi_{32} \to n \in B_{40} \to S(n) \in B_{40}$
\item $\varphi_{128} :\equiv (r \in \mathbb{R} \to r \neq 0 \to \varphi_{32} \to n \in B_{40} \to n \in \omega) \to r \in \mathbb{R} \to r \neq 0 \to \varphi_{32} \to n \in B_{40} \to S(n) \in \omega$
\item $\varphi_{129} :\equiv \forall n\,(r \in \mathbb{R} \land r \neq 0 \land \varphi_{32} \to n \in B_{40} \to S(n) \in B_{40}) \to \varphi_{91}$
\item $\varphi_{130} :\equiv (r \in \mathbb{R} \to r \neq 0 \to \varphi_{97}) \to r \in \mathbb{R} \to r \neq 0 \to \varphi_{32} \to n \in B_{40} \to S(n) \in B_{40}$
\item $\varphi_{131} :\equiv (r \in \mathbb{R} \to r \neq 0 \to \varphi_{32} \to n \in B_{40} \to n \in B_{40}) \to r \in \mathbb{R} \to \varphi_{109}$
\item $\varphi_{132} :\equiv (r \in \mathbb{R} \to r \neq 0 \to \varphi_{32} \to \emptyset \in \omega \land \forall y \in \omega\,(S(y) \in \omega)) \to r \in \mathbb{R} \to r \neq 0 \to \varphi_{32} \to n \in B_{40} \to \forall y \in \omega\,(S(y) \in \omega)$
\item $\varphi_{133} :\equiv (r^{-1} \in \mathbb{R} \to \varphi_{81}) \to (r \in \mathbb{R} \to r \neq 0 \to r^{-1} \in \mathbb{R}) \to r \in \mathbb{R} \to r \neq 0 \to \varphi_{81}$
\item $\varphi_{134} :\equiv (r \in \mathbb{R} \to r \neq 0 \to \varphi_{32} \to n \in \omega \to \forall x\,(\mathrm{Ind}(x) \to n \in x)) \to \varphi_{118}$
\item $\varphi_{135} :\equiv ({{(r^{-1})}^{S(n)}} = ({{r}^{n}} \cdot r)^{-1} \to \varphi_{52}) \to r \in \mathbb{R} \to r \neq 0 \to \varphi_{32} \to \varphi_{42}$
\item $\varphi_{136} :\equiv (r \in \mathbb{R} \to \varphi_{92}) \to (S(n) = 1_{\omega} \to 1_{\omega} = S(n)) \to r \in \mathbb{R} \to \varphi_{93}$
\item $\varphi_{137} :\equiv (r \in \mathbb{R} \land r \neq 0 \land \varphi_{32} \to \forall n\,(r \in \mathbb{R} \land r \neq 0 \land \varphi_{32})) \to r \in \mathbb{R} \to r \neq 0 \to \varphi_{32} \to \forall n \in B_{40}\,(S(n) \in B_{40})$
\item $\varphi_{138} :\equiv ({{(r^{-1})}^{S(n)}} = {{r}^{n}}^{-1} \cdot r^{-1} \to \varphi_{67}) \to r \in \mathbb{R} \to r \neq 0 \to \varphi_{32} \to \varphi_{48}$
\item $\varphi_{139} :\equiv \forall n\,(r \in \mathbb{R} \land r \neq 0 \to n \in \mathbb{N} \to {{(r^{-1})}^{n}} = {{r}^{n}}^{-1}) \to \varphi_{96} \to \varphi_{84}$
\item $\varphi_{140} :\equiv (r \in \mathbb{R} \to r \neq 0 \to {{(r^{-1})}^{1_{\omega}}} = {{r}^{1_{\omega}}}^{-1}) \to (1_{\omega} = S(n) \to \varphi_{50}) \to r \in \mathbb{R} \to r \neq 0 \to \varphi_{32} \to \varphi_{39}$
\item $\varphi_{141} :\equiv ({{(r^{-1})}^{S(n)}} = {{(r^{-1})}^{n}} \cdot r^{-1} \to \varphi_{76}) \to r \in \mathbb{R} \to r \neq 0 \to \varphi_{32} \to \varphi_{58}$
\item $\varphi_{142} :\equiv ({{(r^{-1})}^{n}} = {{r}^{n}}^{-1} \to \varphi_{85}) \to r \in \mathbb{R} \to r \neq 0 \to \varphi_{32} \to n \in B_{40} \to \varphi_{47}$
\item $\varphi_{143} :\equiv (r \in \mathbb{R} \to r \neq 0 \to \varphi_{32} \to n \in B_{40} \to \forall y \in \omega\,(S(y) \in \omega)) \to \varphi_{121}$
\item $\varphi_{144} :\equiv (r \in \mathbb{R} \to r \neq 0 \to \varphi_{32} \to \forall n \in B_{40}\,(S(n) \in B_{40})) \to \varphi_{12} \to \varphi_{11} \to r \in \mathbb{R} \to r \neq 0 \to \varphi_{32} \to \mathrm{Ind}(B_{40})$
\item $\varphi_{145} :\equiv (r \in \mathbb{R} \to r \neq 0 \to \varphi_{32} \to \emptyset \in \omega \land \forall y \in \omega\,(S(y) \in \omega)) \to \varphi_{53} \to r \in \mathbb{R} \to r \neq 0 \to \varphi_{32} \to \emptyset \in B_{40}$
\item $\varphi_{146} :\equiv ({{r}^{n}} \cdot r)^{-1} = ({{r}^{n}} \cdot r)^{-1} \to ({{r}^{n}} \cdot r = {{r}^{S(n)}} \to \varphi_{60}) \to r \in \mathbb{R} \to r \neq 0 \to \varphi_{32} \to \varphi_{46}$
\item $\varphi_{147} :\equiv (r \in \mathbb{R} \to r \neq 0 \to r^{-1} \in \mathbb{R}) \to \varphi_{100} \to (r \in \mathbb{R} \to r \neq 0 \to r \neq 0) \to \varphi_{106}$
\item $\varphi_{148} :\equiv \varphi_{100} \to r \in \mathbb{R} \to r \neq 0 \to \varphi_{32} \to n \in B_{40} \to S(n) \in \mathbb{N} \to n \neq \emptyset \to \varphi_{37}$
\item $\varphi_{149} :\equiv (r \in \mathbb{R} \to r \neq 0 \to \varphi_{32} \to n \in B_{40} \to S(n) \in \mathbb{N} \to n \neq \emptyset \to \varphi_{37}) \to r \in \mathbb{R} \to r \neq 0 \to \varphi_{32} \to \varphi_{5}$
\item $\varphi_{150} :\equiv (r \in \mathbb{R} \to r \neq 0 \to \varphi_{32} \to n \in B_{40} \to n \in \omega \to S(n) \in \omega) \to \varphi_{128}$
\item $\varphi_{151} :\equiv (r \in \mathbb{R} \to r \neq 0 \to \forall n \in \mathbb{N}\,({{r}^{n}} \neq 0)) \to \varphi_{100} \to (r \in \mathbb{R} \to r \neq 0 \to r \neq 0) \to r \in \mathbb{R} \to r \neq 0 \to \varphi_{101}$
\item $\varphi_{152} :\equiv (r \in \mathbb{R} \to r \neq 0 \to \varphi_{101}) \to (\forall n \in \mathbb{N}\,({{r}^{n}} \neq 0) \to n \in \mathbb{N} \to {{r}^{n}} \neq 0) \to r \in \mathbb{R} \to \varphi_{104}$
\item $\varphi_{153} :\equiv \varphi_{106} \to r \in \mathbb{R} \to r \neq 0 \to \varphi_{32} \to n \in B_{40} \to S(n) \in \mathbb{N} \to n \neq \emptyset \to \varphi_{81}$
\item $\varphi_{154} :\equiv (r \in \mathbb{R} \to r \neq 0 \to \varphi_{32} \to n \in B_{40} \to S(n) \in \mathbb{N} \to n \neq \emptyset \to \varphi_{37}) \to r \in \mathbb{R} \to r \neq 0 \to \varphi_{32} \to \varphi_{56}$
\item $\varphi_{155} :\equiv (r \in \mathbb{R} \to \varphi_{112}) \to ({{r}^{S(n)}} = {{r}^{n}} \cdot r \to {{r}^{n}} \cdot r = {{r}^{S(n)}}) \to r \in \mathbb{R} \to \varphi_{113}$
\item $\varphi_{156} :\equiv (r \in \mathbb{R} \to r \neq 0 \to \varphi_{32} \to \mathrm{Ind}(B_{40})) \to (r \in \mathbb{R} \to r \neq 0 \to \varphi_{103}) \to \varphi_{110}$
\item $\varphi_{157} :\equiv (r \in \mathbb{R} \to r \neq 0 \to \varphi_{32} \to n \in B_{40} \to S(n) \in \mathbb{N} \to n \neq \emptyset \to \varphi_{81}) \to r \in \mathbb{R} \to r \neq 0 \to \varphi_{32} \to n \in B_{40} \to \varphi_{63}$
\item $\varphi_{158} :\equiv (n \in \mathbb{N} \leftrightarrow n \in \omega \land n \neq \emptyset) \to (r \in \mathbb{R} \to r \neq 0 \to \varphi_{32} \to n \in B_{40} \to n \in \omega) \to \varphi_{102} \to \varphi_{98}$
\item $\varphi_{159} :\equiv (\mathrm{Seq}({(r)^{\bullet}}) \to n \in \mathbb{N} \to {{r}^{n}} \in \mathbb{R}) \to (r \in \mathbb{R} \to r \neq 0 \to \varphi_{32} \to \varphi_{5}) \to \varphi_{98} \to \varphi_{114}$
\item $\varphi_{160} :\equiv \varphi_{100} \to \varphi_{105} \to (r \in \mathbb{R} \to r \neq 0 \to r \neq 0) \to r \in \mathbb{R} \to r \neq 0 \to \varphi_{32} \to \varphi_{61}$
\item $\varphi_{161} :\equiv (r \in \mathbb{R} \to r \neq 0 \to \varphi_{32} \to \varphi_{39}) \to (r \in \mathbb{R} \to r \neq 0 \to \varphi_{32} \to \varphi_{42}) \to r \in \mathbb{R} \to r \neq 0 \to \varphi_{97}$
\item $\varphi_{162} :\equiv (r \in \mathbb{R} \to r \neq 0 \to \varphi_{32} \to n \in B_{40} \to \varphi_{63}) \to \varphi_{89} \to r \in \mathbb{R} \to r \neq 0 \to \varphi_{32} \to n \in B_{40} \to \varphi_{57}$
\item $\varphi_{163} :\equiv (r \in \mathbb{R} \to r \neq 0 \to \varphi_{32} \to \varphi_{32}) \to (\varphi_{32} \to \varphi_{55}) \to r \in \mathbb{R} \to r \neq 0 \to \varphi_{32} \to \varphi_{55}$
\item $\varphi_{164} :\equiv (r \in \mathbb{R} \to r \neq 0 \to \varphi_{32} \to n \in B_{40} \to \varphi_{57}) \to \varphi_{98} \to r \in \mathbb{R} \to r \neq 0 \to \varphi_{32} \to n \in B_{40} \to \varphi_{34}$
\item $\varphi_{165} :\equiv (r \in \mathbb{R} \to r \neq 0 \to \varphi_{32} \to \varphi_{32}) \to (\varphi_{32} \to \varphi_{79}) \to r \in \mathbb{R} \to r \neq 0 \to \varphi_{32} \to \varphi_{79}$
\item $\varphi_{166} :\equiv (r \in \mathbb{R} \to r \neq 0 \to \varphi_{32} \to n \in \omega \to \mathrm{Ind}(B_{40}) \to n \in B_{40}) \to \varphi_{156}$
\item $\varphi_{167} :\equiv (r \in \mathbb{R} \to r \neq 0 \to \varphi_{107}) \to {{(r^{-1})}^{n}} \cdot r^{-1} = {{(r^{-1})}^{n}} \cdot r^{-1} \to \varphi_{142}$
\item $\varphi_{168} :\equiv (r \in \mathbb{R} \to r \neq 0 \to \varphi_{32} \to \varphi_{55}) \to (r \in \mathbb{R} \to r \neq 0 \to \varphi_{32} \to n \in B_{40} \to S(n) \in \omega) \to \varphi_{130}$
\item $\varphi_{169} :\equiv (\forall B_{40}\,(\varphi_{49} \to \neg \varphi_{32}) \to \forall B_{40}\,(\varphi_{49}) \to \forall B_{40}\,(\neg \varphi_{32})) \to (\varphi_{49} \to \forall B_{40}\,(\varphi_{49})) \to \varphi_{45}$
\item $\varphi_{170} :\equiv (r \in \mathbb{R} \to r \neq 0 \to \varphi_{32} \to \varphi_{48}) \to (r \in \mathbb{R} \to r \neq 0 \to \varphi_{32} \to \varphi_{46}) \to \varphi_{135}$
\item $\varphi_{171} :\equiv (\forall n\,(\varphi_{111}) \to \varphi_{116}) \to (r \in \mathbb{R} \land r \neq 0 \land \varphi_{32} \to \forall n\,(r \in \mathbb{R} \land r \neq 0 \land \varphi_{32})) \to \varphi_{115}$
\item $\varphi_{172} :\equiv (r \in \mathbb{R} \to r \neq 0 \to \varphi_{32} \to \varphi_{58}) \to (r \in \mathbb{R} \to r \neq 0 \to \varphi_{32} \to \varphi_{62}) \to \varphi_{138}$
\item $\varphi_{173} :\equiv \forall n\,(r \in \mathbb{R} \land r \neq 0 \land \varphi_{32} \to n \in B_{40} \to S(n) \in B_{40}) \to \varphi_{129} \to \varphi_{137}$
\item $\varphi_{174} :\equiv (r \in \mathbb{R} \to r \neq 0 \to \varphi_{32} \to n \in B_{40} \to \varphi_{34}) \to (r \in \mathbb{R} \to r \neq 0 \to \varphi_{32} \to n \in B_{40} \to \varphi_{47}) \to \varphi_{141}$
\item $\varphi_{175} :\equiv (r \in \mathbb{R} \to r \neq 0 \to \exists B_{40}\,(\varphi_{32})) \to \forall B_{40}\,(\varphi_{49} \to \neg \varphi_{32}) \to \varphi_{169}$
\end{abbrevlist}
\end{fullproof}

% 自動生成 (prover/build.js): PowDec  ノード数 152
\begin{fullproof}{[P7]}{$\forall r \in \mathbb{R}\,(0 \leq r \to r \leq 1 \to \mathrm{MonoDec}({(r)^{\bullet}}))$}{152}\phantomsection\label{tree:PowDec}
\noindent$\Pi_{1}$:
\begin{ptree}
\AXC{$r \in \mathbb{R} \to 0 \leq r \to r \leq 1 \to n \in \mathbb{N} \to r \in \mathbb{R}$\lbl{A1}}
\AXC{$\forall r \in \mathbb{R}\,(0 \leq r \to \forall n \in \mathbb{N}\,(0 \leq {{r}^{n}}))$\lbl{[P2]}}
\AXC{$\varphi_{18}$\lbl{A2}}
\BIC{$r \in \mathbb{R} \to 0 \leq r \to \forall n \in \mathbb{N}\,(0 \leq {{r}^{n}})$}
\AXC{$\varphi_{45}$\lbl{A1}}
\BIC{$(r \in \mathbb{R} \to 0 \leq r \to r \leq 1 \to n \in \mathbb{N} \to r \in \mathbb{R}) \to \varphi_{15}$}
\BIC{$\varphi_{15}$}
\end{ptree}
\noindent$\Pi_{2}$:
\begin{ptree}
\AXC{$\forall n \in \mathbb{N}\,(0 \leq {{r}^{n}}) \to n \in \mathbb{N} \to 0 \leq {{r}^{n}}$\lbl{A2}}
\AXC{$r \in \mathbb{R} \to 0 \leq r \to r \leq 1 \to 0 \leq r$\lbl{A1}}
\AXC{$\stackrel{\Pi_{1}}{\vdots}$}\noLine
\UIC{$\varphi_{15}$}
\BIC{$r \in \mathbb{R} \to 0 \leq r \to r \leq 1 \to n \in \mathbb{N} \to \forall n \in \mathbb{N}\,(0 \leq {{r}^{n}})$}
\AXC{$\varphi_{44}$\lbl{A1}}
\BIC{$\varphi_{26}$}
\BIC{$r \in \mathbb{R} \to 0 \leq r \to r \leq 1 \to n \in \mathbb{N} \to n \in \mathbb{N} \to 0 \leq {{r}^{n}}$}
\AXC{$\varphi_{38}$\lbl{A1}}
\BIC{$\varphi_{17}$}
\end{ptree}
\noindent$\Pi_{3}$:
\begin{ptree}
\AXC{$\forall r \in \mathbb{R}\,(\varphi_{10})$\lbl{[S24]}}
\AXC{$\forall r \in \mathbb{R}\,(\varphi_{10}) \to r \in \mathbb{R} \to \varphi_{10}$\lbl{A2}}
\BIC{$r \in \mathbb{R} \to \varphi_{10}$}
\AXC{$(r \in \mathbb{R} \to \varphi_{10}) \to \varphi_{39}$\lbl{A1}}
\BIC{$\varphi_{39}$}
\end{ptree}
\noindent$\Pi_{4}$:
\begin{ptree}
\AXC{$r \in \mathbb{R} \to 0 \leq r \to r \leq 1 \to n \in \mathbb{N} \to r \in \mathbb{R}$\lbl{A1}}
\AXC{$\stackrel{\Pi_{3}}{\vdots}$}\noLine
\UIC{$\varphi_{39}$}
\BIC{$r \in \mathbb{R} \to 0 \leq r \to r \leq 1 \to n \in \mathbb{N} \to \varphi_{10}$}
\AXC{$(r \in \mathbb{R} \to 0 \leq r \to r \leq 1 \to n \in \mathbb{N} \to \varphi_{10}) \to \varphi_{2}$\lbl{A1}}
\BIC{$\varphi_{2}$}
\end{ptree}
\noindent$\Pi_{5}$:
\begin{ptree}
\AXC{$\forall a\,\forall n\,(\mathrm{Seq}(a) \to n \in \mathbb{N} \to a_{n} \in \mathbb{R})$\lbl{[S1]}}
\AXC{$\varphi_{13}$\lbl{A2}}
\BIC{$\forall n\,(\mathrm{Seq}({(r)^{\bullet}}) \to n \in \mathbb{N} \to {{r}^{n}} \in \mathbb{R})$}
\AXC{$\varphi_{16}$\lbl{A2}}
\BIC{$\mathrm{Seq}({(r)^{\bullet}}) \to n \in \mathbb{N} \to {{r}^{n}} \in \mathbb{R}$}
\end{ptree}
\noindent$\Pi_{6}$:
\begin{ptree}
\AXC{$\stackrel{\Pi_{4}}{\vdots}$}\noLine
\UIC{$\varphi_{2}$}
\AXC{$\stackrel{\Pi_{5}}{\vdots}$}\noLine
\UIC{$\mathrm{Seq}({(r)^{\bullet}}) \to n \in \mathbb{N} \to {{r}^{n}} \in \mathbb{R}$}
\AXC{$(\mathrm{Seq}({(r)^{\bullet}}) \to n \in \mathbb{N} \to {{r}^{n}} \in \mathbb{R}) \to \varphi_{2} \to \varphi_{20}$\lbl{A1}}
\BIC{$\varphi_{2} \to \varphi_{20}$}
\BIC{$\varphi_{20}$}
\end{ptree}
\noindent$\Pi_{7}$:
\begin{ptree}
\AXC{$r \in \mathbb{R} \to 0 \leq r \to r \leq 1 \to n \in \mathbb{N} \to n \in \mathbb{N}$\lbl{A1}}
\AXC{$\stackrel{\Pi_{6}}{\vdots}$}\noLine
\UIC{$\varphi_{20}$}
\BIC{$r \in \mathbb{R} \to 0 \leq r \to r \leq 1 \to n \in \mathbb{N} \to {{r}^{n}} \in \mathbb{R}$}
\end{ptree}
\noindent$\Pi_{8}$:
\begin{ptree}
\AXC{$r \in \mathbb{R} \to 0 \leq r \to r \leq 1 \to n \in \mathbb{N} \to r \in \mathbb{R}$\lbl{A1}}
\AXC{$\varphi_{8}$\lbl{[E9]}}
\AXC{$\varphi_{8} \to \varphi_{6}$\lbl{A2}}
\BIC{$\varphi_{6}$}
\AXC{$\varphi_{6} \to \varphi_{5}$\lbl{A2}}
\BIC{$\varphi_{5}$}
\AXC{$\varphi_{5} \to \varphi_{21}$\lbl{A2}}
\BIC{$\varphi_{21}$}
\AXC{$\varphi_{21} \to \varphi_{50}$\lbl{A1}}
\BIC{$\varphi_{50}$}
\BIC{$1 \in \mathbb{R} \to \varphi_{48}$}
\end{ptree}
\noindent$\Pi_{9}$:
\begin{ptree}
\AXC{$\stackrel{\Pi_{7}}{\vdots}$}\noLine
\UIC{$r \in \mathbb{R} \to 0 \leq r \to r \leq 1 \to n \in \mathbb{N} \to {{r}^{n}} \in \mathbb{R}$}
\AXC{$1 \in \mathbb{R}$\lbl{(R1)}}
\AXC{$\stackrel{\Pi_{8}}{\vdots}$}\noLine
\UIC{$1 \in \mathbb{R} \to \varphi_{48}$}
\BIC{$\varphi_{48}$}
\BIC{$(r \in \mathbb{R} \to 0 \leq r \to r \leq 1 \to n \in \mathbb{N} \to 0 \leq {{r}^{n}}) \to \varphi_{19}$}
\end{ptree}
\noindent$\Pi_{10}$:
\begin{ptree}
\AXC{$r \in \mathbb{R} \to 0 \leq r \to r \leq 1 \to n \in \mathbb{N} \to n \in \mathbb{N}$\lbl{A1}}
\AXC{$\stackrel{\Pi_{2}}{\vdots}$}\noLine
\UIC{$\varphi_{17}$}
\BIC{$r \in \mathbb{R} \to 0 \leq r \to r \leq 1 \to n \in \mathbb{N} \to 0 \leq {{r}^{n}}$}
\AXC{$\stackrel{\Pi_{9}}{\vdots}$}\noLine
\UIC{$(r \in \mathbb{R} \to 0 \leq r \to r \leq 1 \to n \in \mathbb{N} \to 0 \leq {{r}^{n}}) \to \varphi_{19}$}
\BIC{$\varphi_{19}$}
\end{ptree}
\noindent$\Pi_{11}$:
\begin{ptree}
\AXC{$x = x$\lbl{A5}}
\AXC{$x = y \to x = x \to y = x$\lbl{A6}}
\AXC{$(x = y \to x = x \to y = x) \to x = x \to x = y \to y = x$\lbl{A1}}
\BIC{$x = x \to x = y \to y = x$}
\BIC{$x = y \to y = x$}
\UIC{$\forall y\,(x = y \to y = x)$}
\UIC{$\forall x\,\forall y\,(x = y \to y = x)$}
\AXC{$\varphi_{3}$\lbl{A2}}
\BIC{$\forall y\,({{r}^{S(n)}} = y \to y = {{r}^{S(n)}})$}
\AXC{$\varphi_{23}$\lbl{A2}}
\BIC{${{r}^{S(n)}} = {{r}^{n}} \cdot r \to {{r}^{n}} \cdot r = {{r}^{S(n)}}$}
\end{ptree}
\noindent$\Pi_{12}$:
\begin{ptree}
\AXC{$\forall r \in \mathbb{R}\,(\varphi_{10})$\lbl{[S24]}}
\AXC{$\forall r \in \mathbb{R}\,(\varphi_{10}) \to r \in \mathbb{R} \to \varphi_{10}$\lbl{A2}}
\BIC{$r \in \mathbb{R} \to \varphi_{10}$}
\AXC{$(r \in \mathbb{R} \to \varphi_{10}) \to \varphi_{39}$\lbl{A1}}
\BIC{$\varphi_{39}$}
\end{ptree}
\noindent$\Pi_{13}$:
\begin{ptree}
\AXC{$r \in \mathbb{R} \to 0 \leq r \to r \leq 1 \to n \in \mathbb{N} \to r \in \mathbb{R}$\lbl{A1}}
\AXC{$\stackrel{\Pi_{12}}{\vdots}$}\noLine
\UIC{$\varphi_{39}$}
\BIC{$r \in \mathbb{R} \to 0 \leq r \to r \leq 1 \to n \in \mathbb{N} \to \varphi_{10}$}
\AXC{$(r \in \mathbb{R} \to 0 \leq r \to r \leq 1 \to n \in \mathbb{N} \to \varphi_{10}) \to \varphi_{12}$\lbl{A1}}
\BIC{$\varphi_{12}$}
\AXC{$\varphi_{12} \to \varphi_{14} \to \varphi_{11}$\lbl{A1}}
\BIC{$\varphi_{14} \to \varphi_{11}$}
\end{ptree}
\noindent$\Pi_{14}$:
\begin{ptree}
\AXC{$r \in \mathbb{R} \to 0 \leq r \to r \leq 1 \to n \in \mathbb{N} \to n \in \mathbb{N}$\lbl{A1}}
\AXC{$\varphi_{14}$\lbl{A2}}
\AXC{$\stackrel{\Pi_{13}}{\vdots}$}\noLine
\UIC{$\varphi_{14} \to \varphi_{11}$}
\BIC{$\varphi_{11}$}
\AXC{$\varphi_{11} \to \varphi_{24}$\lbl{A1}}
\BIC{$\varphi_{24}$}
\BIC{$r \in \mathbb{R} \to 0 \leq r \to r \leq 1 \to n \in \mathbb{N} \to {{r}^{S(n)}} = {{r}^{n}} \cdot r$}
\AXC{$\varphi_{47}$\lbl{A1}}
\BIC{$\varphi_{33}$}
\end{ptree}
\noindent$\Pi_{15}$:
\begin{ptree}
\AXC{$r \in \mathbb{R} \to 0 \leq r \to r \leq 1 \to r \leq 1$\lbl{A1}}
\AXC{$\stackrel{\Pi_{10}}{\vdots}$}\noLine
\UIC{$\varphi_{19}$}
\BIC{$r \in \mathbb{R} \to 0 \leq r \to r \leq 1 \to n \in \mathbb{N} \to {{r}^{n}} \cdot r \leq {{r}^{n}} \cdot 1$}
\AXC{$\stackrel{\Pi_{11}}{\vdots}$}\noLine
\UIC{${{r}^{S(n)}} = {{r}^{n}} \cdot r \to {{r}^{n}} \cdot r = {{r}^{S(n)}}$}
\AXC{$\stackrel{\Pi_{14}}{\vdots}$}\noLine
\UIC{$\varphi_{33}$}
\BIC{$r \in \mathbb{R} \to 0 \leq r \to r \leq 1 \to n \in \mathbb{N} \to {{r}^{n}} \cdot r = {{r}^{S(n)}}$}
\AXC{$\varphi_{55}$\lbl{A1}}
\BIC{$\varphi_{54}$}
\BIC{$\varphi_{51}$}
\end{ptree}
\noindent$\Pi_{16}$:
\begin{ptree}
\AXC{$z = w \to z < {{r}^{n}} \cdot 1 \to w < {{r}^{n}} \cdot 1$\lbl{A6}}
\UIC{$\forall w\,(z = w \to z < {{r}^{n}} \cdot 1 \to w < {{r}^{n}} \cdot 1)$}
\UIC{$\forall z\,\forall w\,(z = w \to z < {{r}^{n}} \cdot 1 \to w < {{r}^{n}} \cdot 1)$}
\AXC{$\varphi_{36}$\lbl{A2}}
\BIC{$\forall w\,({{r}^{n}} \cdot r = w \to {{r}^{n}} \cdot r < {{r}^{n}} \cdot 1 \to w < {{r}^{n}} \cdot 1)$}
\AXC{$\varphi_{46}$\lbl{A2}}
\BIC{${{r}^{n}} \cdot r = {{r}^{S(n)}} \to {{r}^{n}} \cdot r < {{r}^{n}} \cdot 1 \to {{r}^{S(n)}} < {{r}^{n}} \cdot 1$}
\AXC{$\stackrel{\Pi_{15}}{\vdots}$}\noLine
\UIC{$\varphi_{51}$}
\BIC{$\varphi_{41}$}
\end{ptree}
\noindent$\Pi_{17}$:
\begin{ptree}
\AXC{$\forall x \in \mathbb{R}\,(x \cdot 1 = x)$\lbl{(R8)}}
\AXC{$\forall x \in \mathbb{R}\,(x \cdot 1 = x) \to {{r}^{n}} \in \mathbb{R} \to {{r}^{n}} \cdot 1 = {{r}^{n}}$\lbl{A2}}
\BIC{${{r}^{n}} \in \mathbb{R} \to {{r}^{n}} \cdot 1 = {{r}^{n}}$}
\AXC{$({{r}^{n}} \in \mathbb{R} \to {{r}^{n}} \cdot 1 = {{r}^{n}}) \to \varphi_{27}$\lbl{A1}}
\BIC{$\varphi_{27}$}
\end{ptree}
\noindent$\Pi_{18}$:
\begin{ptree}
\AXC{$\stackrel{\Pi_{7}}{\vdots}$}\noLine
\UIC{$r \in \mathbb{R} \to 0 \leq r \to r \leq 1 \to n \in \mathbb{N} \to {{r}^{n}} \in \mathbb{R}$}
\AXC{$\stackrel{\Pi_{17}}{\vdots}$}\noLine
\UIC{$\varphi_{27}$}
\BIC{$r \in \mathbb{R} \to 0 \leq r \to r \leq 1 \to n \in \mathbb{N} \to {{r}^{n}} \cdot 1 = {{r}^{n}}$}
\AXC{$(r \in \mathbb{R} \to 0 \leq r \to r \leq 1 \to n \in \mathbb{N} \to {{r}^{n}} \cdot 1 = {{r}^{n}}) \to \varphi_{53}$\lbl{A1}}
\BIC{$\varphi_{53}$}
\end{ptree}
\noindent$\Pi_{19}$:
\begin{ptree}
\AXC{$z = w \to z = {{r}^{n}} \cdot 1 \to w = {{r}^{n}} \cdot 1$\lbl{A6}}
\UIC{$\forall w\,(z = w \to z = {{r}^{n}} \cdot 1 \to w = {{r}^{n}} \cdot 1)$}
\UIC{$\forall z\,\forall w\,(z = w \to z = {{r}^{n}} \cdot 1 \to w = {{r}^{n}} \cdot 1)$}
\AXC{$\varphi_{30}$\lbl{A2}}
\BIC{$\forall w\,({{r}^{n}} \cdot r = w \to {{r}^{n}} \cdot r = {{r}^{n}} \cdot 1 \to w = {{r}^{n}} \cdot 1)$}
\AXC{$\varphi_{43}$\lbl{A2}}
\BIC{${{r}^{n}} \cdot r = {{r}^{S(n)}} \to {{r}^{n}} \cdot r = {{r}^{n}} \cdot 1 \to {{r}^{S(n)}} = {{r}^{n}} \cdot 1$}
\AXC{$\stackrel{\Pi_{16}}{\vdots}$}\noLine
\UIC{$\varphi_{41}$}
\BIC{$r \in \mathbb{R} \to 0 \leq r \to r \leq 1 \to n \in \mathbb{N} \to {{r}^{S(n)}} \leq {{r}^{n}} \cdot 1$}
\AXC{$\stackrel{\Pi_{18}}{\vdots}$}\noLine
\UIC{$\varphi_{53}$}
\BIC{$({{r}^{n}} \cdot 1 = {{r}^{n}} \to {{r}^{S(n)}} < {{r}^{n}} \cdot 1 \to {{r}^{S(n)}} < {{r}^{n}}) \to \varphi_{37}$}
\end{ptree}
\noindent$\Pi_{20}$:
\begin{ptree}
\AXC{$z = w \to {{r}^{S(n)}} < z \to {{r}^{S(n)}} < w$\lbl{A6}}
\UIC{$\forall w\,(z = w \to {{r}^{S(n)}} < z \to {{r}^{S(n)}} < w)$}
\UIC{$\forall z\,\forall w\,(z = w \to {{r}^{S(n)}} < z \to {{r}^{S(n)}} < w)$}
\AXC{$\varphi_{32}$\lbl{A2}}
\BIC{$\forall w\,({{r}^{n}} \cdot 1 = w \to {{r}^{S(n)}} < {{r}^{n}} \cdot 1 \to {{r}^{S(n)}} < w)$}
\AXC{$\varphi_{42}$\lbl{A2}}
\BIC{${{r}^{n}} \cdot 1 = {{r}^{n}} \to {{r}^{S(n)}} < {{r}^{n}} \cdot 1 \to {{r}^{S(n)}} < {{r}^{n}}$}
\AXC{$\stackrel{\Pi_{19}}{\vdots}$}\noLine
\UIC{$({{r}^{n}} \cdot 1 = {{r}^{n}} \to {{r}^{S(n)}} < {{r}^{n}} \cdot 1 \to {{r}^{S(n)}} < {{r}^{n}}) \to \varphi_{37}$}
\BIC{$\varphi_{37}$}
\end{ptree}
\noindent$\Pi_{21}$:
\begin{ptree}
\AXC{$\varphi_{31}$\lbl{A3}}
\AXC{$z = w \to {{r}^{S(n)}} = z \to {{r}^{S(n)}} = w$\lbl{A6}}
\UIC{$\forall w\,(z = w \to {{r}^{S(n)}} = z \to {{r}^{S(n)}} = w)$}
\UIC{$\forall z\,\forall w\,(z = w \to {{r}^{S(n)}} = z \to {{r}^{S(n)}} = w)$}
\AXC{$\varphi_{28}$\lbl{A2}}
\BIC{$\forall w\,({{r}^{n}} \cdot 1 = w \to {{r}^{S(n)}} = {{r}^{n}} \cdot 1 \to {{r}^{S(n)}} = w)$}
\AXC{$\varphi_{40}$\lbl{A2}}
\BIC{${{r}^{n}} \cdot 1 = {{r}^{n}} \to {{r}^{S(n)}} = {{r}^{n}} \cdot 1 \to {{r}^{S(n)}} = {{r}^{n}}$}
\AXC{$\stackrel{\Pi_{20}}{\vdots}$}\noLine
\UIC{$\varphi_{37}$}
\BIC{$r \in \mathbb{R} \to 0 \leq r \to r \leq 1 \to n \in \mathbb{N} \to {{r}^{S(n)}} \leq {{r}^{n}}$}
\AXC{$\varphi_{29}$\lbl{A1}}
\BIC{$r \in \mathbb{R} \land 0 \leq r \land r \leq 1 \to n \in \mathbb{N} \to {{r}^{S(n)}} \leq {{r}^{n}}$}
\UIC{$\forall n\,(r \in \mathbb{R} \land 0 \leq r \land r \leq 1 \to n \in \mathbb{N} \to {{r}^{S(n)}} \leq {{r}^{n}})$}
\AXC{$\varphi_{52}$\lbl{A1}}
\BIC{$\varphi_{31} \to \varphi_{25}$}
\BIC{$\varphi_{25}$}
\end{ptree}
\noindent$\Pi_{22}$:
\begin{ptree}
\AXC{$r \in \mathbb{R} \to 0 \leq r \to r \leq 1 \to r \in \mathbb{R}$\lbl{A1}}
\AXC{$\forall r \in \mathbb{R}\,(\varphi_{10})$\lbl{[S24]}}
\AXC{$\forall r \in \mathbb{R}\,(\varphi_{10}) \to r \in \mathbb{R} \to \varphi_{10}$\lbl{A2}}
\BIC{$r \in \mathbb{R} \to \varphi_{10}$}
\AXC{$(r \in \mathbb{R} \to \varphi_{10}) \to \varphi_{34}$\lbl{A1}}
\BIC{$\varphi_{34}$}
\BIC{$r \in \mathbb{R} \to 0 \leq r \to r \leq 1 \to \varphi_{10}$}
\end{ptree}
\noindent$\Pi_{23}$:
\begin{ptree}
\AXC{$\stackrel{\Pi_{22}}{\vdots}$}\noLine
\UIC{$r \in \mathbb{R} \to 0 \leq r \to r \leq 1 \to \varphi_{10}$}
\AXC{$\varphi_{35}$\lbl{A1}}
\BIC{$r \in \mathbb{R} \to 0 \leq r \to r \leq 1 \to \mathrm{Seq}({(r)^{\bullet}})$}
\AXC{$\varphi_{4}$\lbl{DefMonoDec}}
\AXC{$\varphi_{4} \to \varphi_{9}$\lbl{A2}}
\BIC{$\varphi_{9}$}
\AXC{$\varphi_{49}$\lbl{A1}}
\BIC{$(r \in \mathbb{R} \to 0 \leq r \to r \leq 1 \to \mathrm{Seq}({(r)^{\bullet}})) \to \varphi_{22}$}
\BIC{$\varphi_{22}$}
\end{ptree}
\noindent 主木:
\begin{ptree}
\AXC{$r \in \mathbb{R} \land 0 \leq r \land r \leq 1 \to \forall n\,(r \in \mathbb{R} \land 0 \leq r \land r \leq 1)$\lbl{A4}}
\AXC{$\stackrel{\Pi_{21}}{\vdots}$}\noLine
\UIC{$\varphi_{25}$}
\BIC{$r \in \mathbb{R} \to 0 \leq r \to r \leq 1 \to \forall n \in \mathbb{N}\,({{r}^{S(n)}} \leq {{r}^{n}})$}
\AXC{$\stackrel{\Pi_{23}}{\vdots}$}\noLine
\UIC{$\varphi_{22}$}
\BIC{$r \in \mathbb{R} \to 0 \leq r \to r \leq 1 \to \mathrm{MonoDec}({(r)^{\bullet}})$}
\UIC{$\varphi_{1}$}
\end{ptree}
\begin{abbrevlist}
\item $\varphi_{1} :\equiv \forall r \in \mathbb{R}\,(0 \leq r \to r \leq 1 \to \mathrm{MonoDec}({(r)^{\bullet}}))$
\item $\varphi_{2} :\equiv r \in \mathbb{R} \to 0 \leq r \to r \leq 1 \to n \in \mathbb{N} \to \mathrm{Seq}({(r)^{\bullet}})$
\item $\varphi_{3} :\equiv \forall x\,\forall y\,(x = y \to y = x) \to \forall y\,({{r}^{S(n)}} = y \to y = {{r}^{S(n)}})$
\item $\varphi_{4} :\equiv \forall a\,(\mathrm{MonoDec}(a) \leftrightarrow \mathrm{Seq}(a) \land \forall n \in \mathbb{N}\,(a_{S(n)} \leq a_{n}))$
\item $\varphi_{5} :\equiv \forall z\,(r \in \mathbb{R} \to 1 \in \mathbb{R} \to z \in \mathbb{R} \to 0 \leq z \to r \leq 1 \to z \cdot r \leq z \cdot 1)$
\item $\varphi_{6} :\equiv \forall y\,\forall z\,(r \in \mathbb{R} \to y \in \mathbb{R} \to z \in \mathbb{R} \to 0 \leq z \to r \leq y \to z \cdot r \leq z \cdot y)$
\item $\varphi_{7} :\equiv \forall n\,(r \in \mathbb{R} \land 0 \leq r \land r \leq 1) \to \forall n \in \mathbb{N}\,({{r}^{S(n)}} \leq {{r}^{n}})$
\item $\varphi_{8} :\equiv \forall x\,\forall y\,\forall z\,(x \in \mathbb{R} \to y \in \mathbb{R} \to z \in \mathbb{R} \to 0 \leq z \to x \leq y \to z \cdot x \leq z \cdot y)$
\item $\varphi_{9} :\equiv \mathrm{MonoDec}({(r)^{\bullet}}) \leftrightarrow \mathrm{Seq}({(r)^{\bullet}}) \land \forall n \in \mathbb{N}\,({{r}^{S(n)}} \leq {{r}^{n}})$
\item $\varphi_{10} :\equiv \mathrm{Seq}({(r)^{\bullet}}) \land {{r}^{1_{\omega}}} = r \land \forall n \in \mathbb{N}\,({{r}^{S(n)}} = {{r}^{n}} \cdot r)$
\item $\varphi_{11} :\equiv r \in \mathbb{R} \to 0 \leq r \to r \leq 1 \to n \in \mathbb{N} \to n \in \mathbb{N} \to {{r}^{S(n)}} = {{r}^{n}} \cdot r$
\item $\varphi_{12} :\equiv r \in \mathbb{R} \to 0 \leq r \to r \leq 1 \to n \in \mathbb{N} \to \forall n \in \mathbb{N}\,({{r}^{S(n)}} = {{r}^{n}} \cdot r)$
\item $\varphi_{13} :\equiv \forall a\,\forall n\,(\mathrm{Seq}(a) \to n \in \mathbb{N} \to a_{n} \in \mathbb{R}) \to \forall n\,(\mathrm{Seq}({(r)^{\bullet}}) \to n \in \mathbb{N} \to {{r}^{n}} \in \mathbb{R})$
\item $\varphi_{14} :\equiv \forall n \in \mathbb{N}\,({{r}^{S(n)}} = {{r}^{n}} \cdot r) \to n \in \mathbb{N} \to {{r}^{S(n)}} = {{r}^{n}} \cdot r$
\item $\varphi_{15} :\equiv (r \in \mathbb{R} \to 0 \leq r \to r \leq 1 \to 0 \leq r) \to r \in \mathbb{R} \to 0 \leq r \to r \leq 1 \to n \in \mathbb{N} \to \forall n \in \mathbb{N}\,(0 \leq {{r}^{n}})$
\item $\varphi_{16} :\equiv \forall n\,(\mathrm{Seq}({(r)^{\bullet}}) \to n \in \mathbb{N} \to {{r}^{n}} \in \mathbb{R}) \to \mathrm{Seq}({(r)^{\bullet}}) \to n \in \mathbb{N} \to {{r}^{n}} \in \mathbb{R}$
\item $\varphi_{17} :\equiv (r \in \mathbb{R} \to 0 \leq r \to r \leq 1 \to n \in \mathbb{N} \to n \in \mathbb{N}) \to r \in \mathbb{R} \to 0 \leq r \to r \leq 1 \to n \in \mathbb{N} \to 0 \leq {{r}^{n}}$
\item $\varphi_{18} :\equiv \forall r \in \mathbb{R}\,(0 \leq r \to \forall n \in \mathbb{N}\,(0 \leq {{r}^{n}})) \to r \in \mathbb{R} \to 0 \leq r \to \forall n \in \mathbb{N}\,(0 \leq {{r}^{n}})$
\item $\varphi_{19} :\equiv (r \in \mathbb{R} \to 0 \leq r \to r \leq 1 \to r \leq 1) \to r \in \mathbb{R} \to 0 \leq r \to r \leq 1 \to n \in \mathbb{N} \to {{r}^{n}} \cdot r \leq {{r}^{n}} \cdot 1$
\item $\varphi_{20} :\equiv (r \in \mathbb{R} \to 0 \leq r \to r \leq 1 \to n \in \mathbb{N} \to n \in \mathbb{N}) \to r \in \mathbb{R} \to 0 \leq r \to r \leq 1 \to n \in \mathbb{N} \to {{r}^{n}} \in \mathbb{R}$
\item $\varphi_{21} :\equiv r \in \mathbb{R} \to 1 \in \mathbb{R} \to {{r}^{n}} \in \mathbb{R} \to 0 \leq {{r}^{n}} \to r \leq 1 \to {{r}^{n}} \cdot r \leq {{r}^{n}} \cdot 1$
\item $\varphi_{22} :\equiv (r \in \mathbb{R} \to 0 \leq r \to r \leq 1 \to \forall n \in \mathbb{N}\,({{r}^{S(n)}} \leq {{r}^{n}})) \to r \in \mathbb{R} \to 0 \leq r \to r \leq 1 \to \mathrm{MonoDec}({(r)^{\bullet}})$
\item $\varphi_{23} :\equiv \forall y\,({{r}^{S(n)}} = y \to y = {{r}^{S(n)}}) \to {{r}^{S(n)}} = {{r}^{n}} \cdot r \to {{r}^{n}} \cdot r = {{r}^{S(n)}}$
\item $\varphi_{24} :\equiv (r \in \mathbb{R} \to 0 \leq r \to r \leq 1 \to n \in \mathbb{N} \to n \in \mathbb{N}) \to r \in \mathbb{R} \to 0 \leq r \to r \leq 1 \to n \in \mathbb{N} \to {{r}^{S(n)}} = {{r}^{n}} \cdot r$
\item $\varphi_{25} :\equiv (r \in \mathbb{R} \land 0 \leq r \land r \leq 1 \to \forall n\,(r \in \mathbb{R} \land 0 \leq r \land r \leq 1)) \to r \in \mathbb{R} \to 0 \leq r \to r \leq 1 \to \forall n \in \mathbb{N}\,({{r}^{S(n)}} \leq {{r}^{n}})$
\item $\varphi_{26} :\equiv (\forall n \in \mathbb{N}\,(0 \leq {{r}^{n}}) \to n \in \mathbb{N} \to 0 \leq {{r}^{n}}) \to r \in \mathbb{R} \to 0 \leq r \to r \leq 1 \to n \in \mathbb{N} \to n \in \mathbb{N} \to 0 \leq {{r}^{n}}$
\item $\varphi_{27} :\equiv (r \in \mathbb{R} \to 0 \leq r \to r \leq 1 \to n \in \mathbb{N} \to {{r}^{n}} \in \mathbb{R}) \to r \in \mathbb{R} \to 0 \leq r \to r \leq 1 \to n \in \mathbb{N} \to {{r}^{n}} \cdot 1 = {{r}^{n}}$
\item $\varphi_{28} :\equiv \forall z\,\forall w\,(z = w \to {{r}^{S(n)}} = z \to {{r}^{S(n)}} = w) \to \forall w\,({{r}^{n}} \cdot 1 = w \to {{r}^{S(n)}} = {{r}^{n}} \cdot 1 \to {{r}^{S(n)}} = w)$
\item $\varphi_{29} :\equiv (r \in \mathbb{R} \to 0 \leq r \to r \leq 1 \to n \in \mathbb{N} \to {{r}^{S(n)}} \leq {{r}^{n}}) \to r \in \mathbb{R} \land 0 \leq r \land r \leq 1 \to n \in \mathbb{N} \to {{r}^{S(n)}} \leq {{r}^{n}}$
\item $\varphi_{30} :\equiv \forall z\,\forall w\,(z = w \to z = {{r}^{n}} \cdot 1 \to w = {{r}^{n}} \cdot 1) \to \forall w\,({{r}^{n}} \cdot r = w \to {{r}^{n}} \cdot r = {{r}^{n}} \cdot 1 \to w = {{r}^{n}} \cdot 1)$
\item $\varphi_{31} :\equiv \forall n\,(r \in \mathbb{R} \land 0 \leq r \land r \leq 1 \to n \in \mathbb{N} \to {{r}^{S(n)}} \leq {{r}^{n}}) \to \varphi_{7}$
\item $\varphi_{32} :\equiv \forall z\,\forall w\,(z = w \to {{r}^{S(n)}} < z \to {{r}^{S(n)}} < w) \to \forall w\,({{r}^{n}} \cdot 1 = w \to {{r}^{S(n)}} < {{r}^{n}} \cdot 1 \to {{r}^{S(n)}} < w)$
\item $\varphi_{33} :\equiv ({{r}^{S(n)}} = {{r}^{n}} \cdot r \to {{r}^{n}} \cdot r = {{r}^{S(n)}}) \to r \in \mathbb{R} \to 0 \leq r \to r \leq 1 \to n \in \mathbb{N} \to {{r}^{n}} \cdot r = {{r}^{S(n)}}$
\item $\varphi_{34} :\equiv (r \in \mathbb{R} \to 0 \leq r \to r \leq 1 \to r \in \mathbb{R}) \to r \in \mathbb{R} \to 0 \leq r \to r \leq 1 \to \varphi_{10}$
\item $\varphi_{35} :\equiv (r \in \mathbb{R} \to 0 \leq r \to r \leq 1 \to \varphi_{10}) \to r \in \mathbb{R} \to 0 \leq r \to r \leq 1 \to \mathrm{Seq}({(r)^{\bullet}})$
\item $\varphi_{36} :\equiv \forall z\,\forall w\,(z = w \to z < {{r}^{n}} \cdot 1 \to w < {{r}^{n}} \cdot 1) \to \forall w\,({{r}^{n}} \cdot r = w \to {{r}^{n}} \cdot r < {{r}^{n}} \cdot 1 \to w < {{r}^{n}} \cdot 1)$
\item $\varphi_{37} :\equiv ({{r}^{n}} \cdot 1 = {{r}^{n}} \to {{r}^{S(n)}} = {{r}^{n}} \cdot 1 \to {{r}^{S(n)}} = {{r}^{n}}) \to r \in \mathbb{R} \to 0 \leq r \to r \leq 1 \to n \in \mathbb{N} \to {{r}^{S(n)}} \leq {{r}^{n}}$
\item $\varphi_{38} :\equiv (r \in \mathbb{R} \to 0 \leq r \to r \leq 1 \to n \in \mathbb{N} \to n \in \mathbb{N} \to 0 \leq {{r}^{n}}) \to \varphi_{17}$
\item $\varphi_{39} :\equiv (r \in \mathbb{R} \to 0 \leq r \to r \leq 1 \to n \in \mathbb{N} \to r \in \mathbb{R}) \to r \in \mathbb{R} \to 0 \leq r \to r \leq 1 \to n \in \mathbb{N} \to \varphi_{10}$
\item $\varphi_{40} :\equiv \forall w\,({{r}^{n}} \cdot 1 = w \to {{r}^{S(n)}} = {{r}^{n}} \cdot 1 \to {{r}^{S(n)}} = w) \to {{r}^{n}} \cdot 1 = {{r}^{n}} \to {{r}^{S(n)}} = {{r}^{n}} \cdot 1 \to {{r}^{S(n)}} = {{r}^{n}}$
\item $\varphi_{41} :\equiv ({{r}^{n}} \cdot r = {{r}^{S(n)}} \to {{r}^{n}} \cdot r = {{r}^{n}} \cdot 1 \to {{r}^{S(n)}} = {{r}^{n}} \cdot 1) \to r \in \mathbb{R} \to 0 \leq r \to r \leq 1 \to n \in \mathbb{N} \to {{r}^{S(n)}} \leq {{r}^{n}} \cdot 1$
\item $\varphi_{42} :\equiv \forall w\,({{r}^{n}} \cdot 1 = w \to {{r}^{S(n)}} < {{r}^{n}} \cdot 1 \to {{r}^{S(n)}} < w) \to {{r}^{n}} \cdot 1 = {{r}^{n}} \to {{r}^{S(n)}} < {{r}^{n}} \cdot 1 \to {{r}^{S(n)}} < {{r}^{n}}$
\item $\varphi_{43} :\equiv \forall w\,({{r}^{n}} \cdot r = w \to {{r}^{n}} \cdot r = {{r}^{n}} \cdot 1 \to w = {{r}^{n}} \cdot 1) \to {{r}^{n}} \cdot r = {{r}^{S(n)}} \to {{r}^{n}} \cdot r = {{r}^{n}} \cdot 1 \to {{r}^{S(n)}} = {{r}^{n}} \cdot 1$
\item $\varphi_{44} :\equiv (r \in \mathbb{R} \to 0 \leq r \to r \leq 1 \to n \in \mathbb{N} \to \forall n \in \mathbb{N}\,(0 \leq {{r}^{n}})) \to \varphi_{26}$
\item $\varphi_{45} :\equiv (r \in \mathbb{R} \to 0 \leq r \to \forall n \in \mathbb{N}\,(0 \leq {{r}^{n}})) \to (r \in \mathbb{R} \to 0 \leq r \to r \leq 1 \to n \in \mathbb{N} \to r \in \mathbb{R}) \to \varphi_{15}$
\item $\varphi_{46} :\equiv \forall w\,({{r}^{n}} \cdot r = w \to {{r}^{n}} \cdot r < {{r}^{n}} \cdot 1 \to w < {{r}^{n}} \cdot 1) \to {{r}^{n}} \cdot r = {{r}^{S(n)}} \to {{r}^{n}} \cdot r < {{r}^{n}} \cdot 1 \to {{r}^{S(n)}} < {{r}^{n}} \cdot 1$
\item $\varphi_{47} :\equiv (r \in \mathbb{R} \to 0 \leq r \to r \leq 1 \to n \in \mathbb{N} \to {{r}^{S(n)}} = {{r}^{n}} \cdot r) \to \varphi_{33}$
\item $\varphi_{48} :\equiv (r \in \mathbb{R} \to 0 \leq r \to r \leq 1 \to n \in \mathbb{N} \to {{r}^{n}} \in \mathbb{R}) \to (r \in \mathbb{R} \to 0 \leq r \to r \leq 1 \to n \in \mathbb{N} \to 0 \leq {{r}^{n}}) \to \varphi_{19}$
\item $\varphi_{49} :\equiv \varphi_{9} \to (r \in \mathbb{R} \to 0 \leq r \to r \leq 1 \to \mathrm{Seq}({(r)^{\bullet}})) \to \varphi_{22}$
\item $\varphi_{50} :\equiv (r \in \mathbb{R} \to 0 \leq r \to r \leq 1 \to n \in \mathbb{N} \to r \in \mathbb{R}) \to 1 \in \mathbb{R} \to \varphi_{48}$
\item $\varphi_{51} :\equiv ({{r}^{n}} \cdot r = {{r}^{S(n)}} \to {{r}^{n}} \cdot r < {{r}^{n}} \cdot 1 \to {{r}^{S(n)}} < {{r}^{n}} \cdot 1) \to \varphi_{41}$
\item $\varphi_{52} :\equiv \forall n\,(r \in \mathbb{R} \land 0 \leq r \land r \leq 1 \to n \in \mathbb{N} \to {{r}^{S(n)}} \leq {{r}^{n}}) \to \varphi_{31} \to \varphi_{25}$
\item $\varphi_{53} :\equiv (r \in \mathbb{R} \to 0 \leq r \to r \leq 1 \to n \in \mathbb{N} \to {{r}^{S(n)}} \leq {{r}^{n}} \cdot 1) \to ({{r}^{n}} \cdot 1 = {{r}^{n}} \to {{r}^{S(n)}} < {{r}^{n}} \cdot 1 \to {{r}^{S(n)}} < {{r}^{n}}) \to \varphi_{37}$
\item $\varphi_{54} :\equiv (r \in \mathbb{R} \to 0 \leq r \to r \leq 1 \to n \in \mathbb{N} \to {{r}^{n}} \cdot r \leq {{r}^{n}} \cdot 1) \to \varphi_{51}$
\item $\varphi_{55} :\equiv (r \in \mathbb{R} \to 0 \leq r \to r \leq 1 \to n \in \mathbb{N} \to {{r}^{n}} \cdot r = {{r}^{S(n)}}) \to \varphi_{54}$
\end{abbrevlist}
\end{fullproof}

% 自動生成 (prover/build.js): SucIdx  ノード数 88
\begin{fullproof}{[C28]}{$\mathrm{SubIdx}(\mathsf{s})$}{88}\phantomsection\label{tree:SucIdx}
\noindent$\Pi_{1}$:
\begin{ptree}
\AXC{$z = w \to \mathsf{s}(k) < z \to \mathsf{s}(k) < w$\lbl{A6}}
\UIC{$\forall w\,(z = w \to \mathsf{s}(k) < z \to \mathsf{s}(k) < w)$}
\UIC{$\forall z\,\forall w\,(z = w \to \mathsf{s}(k) < z \to \mathsf{s}(k) < w)$}
\AXC{$\forall z\,\forall w\,(z = w \to \mathsf{s}(k) < z \to \mathsf{s}(k) < w) \to \varphi_{4}$\lbl{A2}}
\BIC{$\varphi_{4}$}
\AXC{$\varphi_{4} \to \varphi_{16}$\lbl{A2}}
\BIC{$\varphi_{16}$}
\end{ptree}
\noindent$\Pi_{2}$:
\begin{ptree}
\AXC{$\forall n \in \mathbb{N}\,(S(n) \in \mathbb{N})$\lbl{[N2]}}
\AXC{$\forall n \in \mathbb{N}\,(S(n) \in \mathbb{N}) \to k \in \mathbb{N} \to S(k) \in \mathbb{N}$\lbl{A2}}
\BIC{$k \in \mathbb{N} \to S(k) \in \mathbb{N}$}
\end{ptree}
\noindent$\Pi_{3}$:
\begin{ptree}
\AXC{$\forall n \in \mathbb{N}\,(n \in \omega)$\lbl{[N7']}}
\AXC{$\forall n \in \mathbb{N}\,(n \in \omega) \to S(k) \in \mathbb{N} \to S(k) \in \omega$\lbl{A2}}
\BIC{$S(k) \in \mathbb{N} \to S(k) \in \omega$}
\AXC{$\varphi_{18}$\lbl{A1}}
\BIC{$(k \in \mathbb{N} \to S(k) \in \mathbb{N}) \to k \in \mathbb{N} \to S(k) \in \omega$}
\end{ptree}
\noindent$\Pi_{4}$:
\begin{ptree}
\AXC{$\forall n \in \omega\,(n < S(n))$\lbl{[N17]}}
\AXC{$\varphi_{8}$\lbl{A2}}
\BIC{$S(k) \in \omega \to S(k) < S(S(k))$}
\AXC{$(S(k) \in \omega \to S(k) < S(S(k))) \to \varphi_{5}$\lbl{A1}}
\BIC{$\varphi_{5}$}
\end{ptree}
\noindent$\Pi_{5}$:
\begin{ptree}
\AXC{$x = x$\lbl{A5}}
\AXC{$x = y \to x = x \to y = x$\lbl{A6}}
\AXC{$(x = y \to x = x \to y = x) \to x = x \to x = y \to y = x$\lbl{A1}}
\BIC{$x = x \to x = y \to y = x$}
\BIC{$x = y \to y = x$}
\UIC{$\forall y\,(x = y \to y = x)$}
\UIC{$\forall x\,\forall y\,(x = y \to y = x)$}
\end{ptree}
\noindent$\Pi_{6}$:
\begin{ptree}
\AXC{$\varphi_{10}$\lbl{A2}}
\AXC{$\mathrm{NSeq}(\mathsf{s}) \land \forall k \in \mathbb{N}\,(\mathsf{s}(k) = S(k))$\lbl{[S14]}}
\AXC{$\varphi_{19}$\lbl{A1}}
\BIC{$k \in \mathbb{N} \to \forall k \in \mathbb{N}\,(\mathsf{s}(k) = S(k))$}
\AXC{$\varphi_{25}$\lbl{A1}}
\BIC{$\varphi_{10} \to k \in \mathbb{N} \to k \in \mathbb{N} \to \mathsf{s}(k) = S(k)$}
\BIC{$k \in \mathbb{N} \to k \in \mathbb{N} \to \mathsf{s}(k) = S(k)$}
\AXC{$\varphi_{20}$\lbl{A1}}
\BIC{$(k \in \mathbb{N} \to k \in \mathbb{N}) \to k \in \mathbb{N} \to \mathsf{s}(k) = S(k)$}
\end{ptree}
\noindent$\Pi_{7}$:
\begin{ptree}
\AXC{$k \in \mathbb{N} \to k \in \mathbb{N}$\lbl{A1}}
\AXC{$\stackrel{\Pi_{6}}{\vdots}$}\noLine
\UIC{$(k \in \mathbb{N} \to k \in \mathbb{N}) \to k \in \mathbb{N} \to \mathsf{s}(k) = S(k)$}
\BIC{$k \in \mathbb{N} \to \mathsf{s}(k) = S(k)$}
\AXC{$(k \in \mathbb{N} \to \mathsf{s}(k) = S(k)) \to \varphi_{12}$\lbl{A1}}
\BIC{$\varphi_{12}$}
\end{ptree}
\noindent$\Pi_{8}$:
\begin{ptree}
\AXC{$\stackrel{\Pi_{5}}{\vdots}$}\noLine
\UIC{$\forall x\,\forall y\,(x = y \to y = x)$}
\AXC{$\varphi_{3}$\lbl{A2}}
\BIC{$\forall y\,(\mathsf{s}(k) = y \to y = \mathsf{s}(k))$}
\AXC{$\varphi_{13}$\lbl{A2}}
\BIC{$\mathsf{s}(k) = S(k) \to S(k) = \mathsf{s}(k)$}
\AXC{$\stackrel{\Pi_{7}}{\vdots}$}\noLine
\UIC{$\varphi_{12}$}
\BIC{$k \in \mathbb{N} \to S(k) = \mathsf{s}(k)$}
\AXC{$(k \in \mathbb{N} \to S(k) = \mathsf{s}(k)) \to \varphi_{24}$\lbl{A1}}
\BIC{$\varphi_{24}$}
\end{ptree}
\noindent$\Pi_{9}$:
\begin{ptree}
\AXC{$\stackrel{\Pi_{2}}{\vdots}$}\noLine
\UIC{$k \in \mathbb{N} \to S(k) \in \mathbb{N}$}
\AXC{$\stackrel{\Pi_{3}}{\vdots}$}\noLine
\UIC{$(k \in \mathbb{N} \to S(k) \in \mathbb{N}) \to k \in \mathbb{N} \to S(k) \in \omega$}
\BIC{$k \in \mathbb{N} \to S(k) \in \omega$}
\AXC{$\stackrel{\Pi_{4}}{\vdots}$}\noLine
\UIC{$\varphi_{5}$}
\BIC{$k \in \mathbb{N} \to S(k) < S(S(k))$}
\AXC{$\stackrel{\Pi_{8}}{\vdots}$}\noLine
\UIC{$\varphi_{24}$}
\BIC{$\varphi_{7} \to k \in \mathbb{N} \to \mathsf{s}(k) < S(S(k))$}
\end{ptree}
\noindent$\Pi_{10}$:
\begin{ptree}
\AXC{$z = w \to z < S(S(k)) \to w < S(S(k))$\lbl{A6}}
\UIC{$\forall w\,(z = w \to z < S(S(k)) \to w < S(S(k)))$}
\UIC{$\forall z\,\forall w\,(z = w \to z < S(S(k)) \to w < S(S(k)))$}
\AXC{$\varphi_{17}$\lbl{A2}}
\BIC{$\varphi_{1}$}
\AXC{$\varphi_{1} \to \varphi_{7}$\lbl{A2}}
\BIC{$\varphi_{7}$}
\AXC{$\stackrel{\Pi_{9}}{\vdots}$}\noLine
\UIC{$\varphi_{7} \to k \in \mathbb{N} \to \mathsf{s}(k) < S(S(k))$}
\BIC{$k \in \mathbb{N} \to \mathsf{s}(k) < S(S(k))$}
\end{ptree}
\noindent$\Pi_{11}$:
\begin{ptree}
\AXC{$\stackrel{\Pi_{5}}{\vdots}$}\noLine
\UIC{$\forall x\,\forall y\,(x = y \to y = x)$}
\AXC{$\varphi_{6}$\lbl{A2}}
\BIC{$\forall y\,(\mathsf{s}(S(k)) = y \to y = \mathsf{s}(S(k)))$}
\AXC{$\forall y\,(\mathsf{s}(S(k)) = y \to y = \mathsf{s}(S(k))) \to \varphi_{2}$\lbl{A2}}
\BIC{$\varphi_{2}$}
\end{ptree}
\noindent$\Pi_{12}$:
\begin{ptree}
\AXC{$\varphi_{14}$\lbl{A2}}
\AXC{$\mathrm{NSeq}(\mathsf{s}) \land \forall k \in \mathbb{N}\,(\mathsf{s}(k) = S(k))$\lbl{[S14]}}
\AXC{$\varphi_{19}$\lbl{A1}}
\BIC{$k \in \mathbb{N} \to \forall k \in \mathbb{N}\,(\mathsf{s}(k) = S(k))$}
\AXC{$\varphi_{27}$\lbl{A1}}
\BIC{$\varphi_{14} \to k \in \mathbb{N} \to S(k) \in \mathbb{N} \to \mathsf{s}(S(k)) = S(S(k))$}
\BIC{$k \in \mathbb{N} \to S(k) \in \mathbb{N} \to \mathsf{s}(S(k)) = S(S(k))$}
\AXC{$\varphi_{22}$\lbl{A1}}
\BIC{$\varphi_{9}$}
\end{ptree}
\noindent$\Pi_{13}$:
\begin{ptree}
\AXC{$\stackrel{\Pi_{11}}{\vdots}$}\noLine
\UIC{$\varphi_{2}$}
\AXC{$\stackrel{\Pi_{2}}{\vdots}$}\noLine
\UIC{$k \in \mathbb{N} \to S(k) \in \mathbb{N}$}
\AXC{$\stackrel{\Pi_{12}}{\vdots}$}\noLine
\UIC{$\varphi_{9}$}
\BIC{$k \in \mathbb{N} \to \mathsf{s}(S(k)) = S(S(k))$}
\AXC{$\varphi_{23}$\lbl{A1}}
\BIC{$\varphi_{2} \to k \in \mathbb{N} \to S(S(k)) = \mathsf{s}(S(k))$}
\BIC{$k \in \mathbb{N} \to S(S(k)) = \mathsf{s}(S(k))$}
\AXC{$(k \in \mathbb{N} \to S(S(k)) = \mathsf{s}(S(k))) \to \varphi_{26}$\lbl{A1}}
\BIC{$\varphi_{26}$}
\end{ptree}
\noindent$\Pi_{14}$:
\begin{ptree}
\AXC{$\mathrm{NSeq}(\mathsf{s}) \land \forall k \in \mathbb{N}\,(\mathsf{s}(k) = S(k))$\lbl{[S14]}}
\AXC{$\varphi_{11}$\lbl{DefSubIdx}}
\AXC{$\varphi_{11} \to \varphi_{15}$\lbl{A2}}
\BIC{$\varphi_{15}$}
\AXC{$\varphi_{15} \to \varphi_{21}$\lbl{A1}}
\BIC{$\varphi_{21}$}
\BIC{$\forall k \in \mathbb{N}\,(\mathsf{s}(k) < \mathsf{s}(S(k))) \to \mathrm{SubIdx}(\mathsf{s})$}
\end{ptree}
\noindent 主木:
\begin{ptree}
\AXC{$\stackrel{\Pi_{1}}{\vdots}$}\noLine
\UIC{$\varphi_{16}$}
\AXC{$\stackrel{\Pi_{10}}{\vdots}$}\noLine
\UIC{$k \in \mathbb{N} \to \mathsf{s}(k) < S(S(k))$}
\AXC{$\stackrel{\Pi_{13}}{\vdots}$}\noLine
\UIC{$\varphi_{26}$}
\BIC{$\varphi_{16} \to k \in \mathbb{N} \to \mathsf{s}(k) < \mathsf{s}(S(k))$}
\BIC{$k \in \mathbb{N} \to \mathsf{s}(k) < \mathsf{s}(S(k))$}
\UIC{$\forall k \in \mathbb{N}\,(\mathsf{s}(k) < \mathsf{s}(S(k)))$}
\AXC{$\stackrel{\Pi_{14}}{\vdots}$}\noLine
\UIC{$\forall k \in \mathbb{N}\,(\mathsf{s}(k) < \mathsf{s}(S(k))) \to \mathrm{SubIdx}(\mathsf{s})$}
\BIC{$\mathrm{SubIdx}(\mathsf{s})$}
\end{ptree}
\begin{abbrevlist}
\item $\varphi_{1} :\equiv \forall w\,(S(k) = w \to S(k) < S(S(k)) \to w < S(S(k)))$
\item $\varphi_{2} :\equiv \mathsf{s}(S(k)) = S(S(k)) \to S(S(k)) = \mathsf{s}(S(k))$
\item $\varphi_{3} :\equiv \forall x\,\forall y\,(x = y \to y = x) \to \forall y\,(\mathsf{s}(k) = y \to y = \mathsf{s}(k))$
\item $\varphi_{4} :\equiv \forall w\,(S(S(k)) = w \to \mathsf{s}(k) < S(S(k)) \to \mathsf{s}(k) < w)$
\item $\varphi_{5} :\equiv (k \in \mathbb{N} \to S(k) \in \omega) \to k \in \mathbb{N} \to S(k) < S(S(k))$
\item $\varphi_{6} :\equiv \forall x\,\forall y\,(x = y \to y = x) \to \forall y\,(\mathsf{s}(S(k)) = y \to y = \mathsf{s}(S(k)))$
\item $\varphi_{7} :\equiv S(k) = \mathsf{s}(k) \to S(k) < S(S(k)) \to \mathsf{s}(k) < S(S(k))$
\item $\varphi_{8} :\equiv \forall n \in \omega\,(n < S(n)) \to S(k) \in \omega \to S(k) < S(S(k))$
\item $\varphi_{9} :\equiv (k \in \mathbb{N} \to S(k) \in \mathbb{N}) \to k \in \mathbb{N} \to \mathsf{s}(S(k)) = S(S(k))$
\item $\varphi_{10} :\equiv \forall k \in \mathbb{N}\,(\mathsf{s}(k) = S(k)) \to k \in \mathbb{N} \to \mathsf{s}(k) = S(k)$
\item $\varphi_{11} :\equiv \forall s\,(\mathrm{SubIdx}(s) \leftrightarrow \mathrm{NSeq}(s) \land \forall k \in \mathbb{N}\,(s_{k} < s_{S(k)}))$
\item $\varphi_{12} :\equiv (\mathsf{s}(k) = S(k) \to S(k) = \mathsf{s}(k)) \to k \in \mathbb{N} \to S(k) = \mathsf{s}(k)$
\item $\varphi_{13} :\equiv \forall y\,(\mathsf{s}(k) = y \to y = \mathsf{s}(k)) \to \mathsf{s}(k) = S(k) \to S(k) = \mathsf{s}(k)$
\item $\varphi_{14} :\equiv \forall k \in \mathbb{N}\,(\mathsf{s}(k) = S(k)) \to S(k) \in \mathbb{N} \to \mathsf{s}(S(k)) = S(S(k))$
\item $\varphi_{15} :\equiv \mathrm{SubIdx}(\mathsf{s}) \leftrightarrow \mathrm{NSeq}(\mathsf{s}) \land \forall k \in \mathbb{N}\,(\mathsf{s}(k) < \mathsf{s}(S(k)))$
\item $\varphi_{16} :\equiv S(S(k)) = \mathsf{s}(S(k)) \to \mathsf{s}(k) < S(S(k)) \to \mathsf{s}(k) < \mathsf{s}(S(k))$
\item $\varphi_{17} :\equiv \forall z\,\forall w\,(z = w \to z < S(S(k)) \to w < S(S(k))) \to \varphi_{1}$
\item $\varphi_{18} :\equiv (S(k) \in \mathbb{N} \to S(k) \in \omega) \to (k \in \mathbb{N} \to S(k) \in \mathbb{N}) \to k \in \mathbb{N} \to S(k) \in \omega$
\item $\varphi_{19} :\equiv \mathrm{NSeq}(\mathsf{s}) \land \forall k \in \mathbb{N}\,(\mathsf{s}(k) = S(k)) \to k \in \mathbb{N} \to \forall k \in \mathbb{N}\,(\mathsf{s}(k) = S(k))$
\item $\varphi_{20} :\equiv (k \in \mathbb{N} \to k \in \mathbb{N} \to \mathsf{s}(k) = S(k)) \to (k \in \mathbb{N} \to k \in \mathbb{N}) \to k \in \mathbb{N} \to \mathsf{s}(k) = S(k)$
\item $\varphi_{21} :\equiv \mathrm{NSeq}(\mathsf{s}) \land \forall k \in \mathbb{N}\,(\mathsf{s}(k) = S(k)) \to \forall k \in \mathbb{N}\,(\mathsf{s}(k) < \mathsf{s}(S(k))) \to \mathrm{SubIdx}(\mathsf{s})$
\item $\varphi_{22} :\equiv (k \in \mathbb{N} \to S(k) \in \mathbb{N} \to \mathsf{s}(S(k)) = S(S(k))) \to \varphi_{9}$
\item $\varphi_{23} :\equiv (k \in \mathbb{N} \to \mathsf{s}(S(k)) = S(S(k))) \to \varphi_{2} \to k \in \mathbb{N} \to S(S(k)) = \mathsf{s}(S(k))$
\item $\varphi_{24} :\equiv (k \in \mathbb{N} \to S(k) < S(S(k))) \to \varphi_{7} \to k \in \mathbb{N} \to \mathsf{s}(k) < S(S(k))$
\item $\varphi_{25} :\equiv (k \in \mathbb{N} \to \forall k \in \mathbb{N}\,(\mathsf{s}(k) = S(k))) \to \varphi_{10} \to k \in \mathbb{N} \to k \in \mathbb{N} \to \mathsf{s}(k) = S(k)$
\item $\varphi_{26} :\equiv (k \in \mathbb{N} \to \mathsf{s}(k) < S(S(k))) \to \varphi_{16} \to k \in \mathbb{N} \to \mathsf{s}(k) < \mathsf{s}(S(k))$
\item $\varphi_{27} :\equiv (k \in \mathbb{N} \to \forall k \in \mathbb{N}\,(\mathsf{s}(k) = S(k))) \to \varphi_{14} \to k \in \mathbb{N} \to S(k) \in \mathbb{N} \to \mathsf{s}(S(k)) = S(S(k))$
\end{abbrevlist}
\end{fullproof}

% 自動生成 (prover/build.js): GeoLim  ノード数 607
\begin{fullproof}{命題3.1(3)}{$\forall r \in \mathbb{R}\,(\lvert r \rvert < 1 \to \mathrm{Conv}({(r)^{\bullet}},0))$}{607}\phantomsection\label{tree:GeoLim}
\noindent$\Pi_{1}$:
\begin{ptree}
\AXC{$y = z \to x = y \to x = z$\lbl{A6}}
\AXC{$(y = z \to x = y \to x = z) \to x = y \to y = z \to x = z$\lbl{A1}}
\BIC{$x = y \to y = z \to x = z$}
\UIC{$\forall z\,(x = y \to y = z \to x = z)$}
\UIC{$\forall y\,\forall z\,(x = y \to y = z \to x = z)$}
\UIC{$\forall x\,\forall y\,\forall z\,(x = y \to y = z \to x = z)$}
\AXC{$\forall x\,\forall y\,\forall z\,(x = y \to y = z \to x = z) \to \varphi_{36}$\lbl{A2}}
\BIC{$\varphi_{36}$}
\end{ptree}
\noindent$\Pi_{2}$:
\begin{ptree}
\AXC{$r \in \mathbb{R} \to \lvert r \rvert < 1 \to \varphi_{9}$\lbl{A1}}
\AXC{$\forall a\,\forall \alpha\,(\mathrm{Conv}(a,\alpha) \to \alpha \in \mathbb{R})$\lbl{[C2]}}
\AXC{$\varphi_{96}$\lbl{A2}}
\BIC{$\forall \alpha\,(\mathrm{Conv}({(\lvert r \rvert)^{\bullet}},\alpha) \to \alpha \in \mathbb{R})$}
\AXC{$\varphi_{97}$\lbl{A2}}
\BIC{$\mathrm{Conv}({(\lvert r \rvert)^{\bullet}},L) \to L \in \mathbb{R}$}
\AXC{$\varphi_{215}$\lbl{A1}}
\BIC{$(r \in \mathbb{R} \to \lvert r \rvert < 1 \to \varphi_{9}) \to \varphi_{53}$}
\BIC{$\varphi_{53}$}
\end{ptree}
\noindent$\Pi_{3}$:
\begin{ptree}
\AXC{$\stackrel{\Pi_{2}}{\vdots}$}\noLine
\UIC{$\varphi_{53}$}
\AXC{$\forall x \in \mathbb{R}\,(x - x = 0)$\lbl{[F13]}}
\AXC{$\forall x \in \mathbb{R}\,(x - x = 0) \to L \in \mathbb{R} \to L - L = 0$\lbl{A2}}
\BIC{$L \in \mathbb{R} \to L - L = 0$}
\AXC{$(L \in \mathbb{R} \to L - L = 0) \to \varphi_{53} \to \varphi_{48}$\lbl{A1}}
\BIC{$\varphi_{53} \to \varphi_{48}$}
\BIC{$\varphi_{48}$}
\end{ptree}
\noindent$\Pi_{4}$:
\begin{ptree}
\AXC{$y = z \to x = y \to x = z$\lbl{A6}}
\AXC{$(y = z \to x = y \to x = z) \to x = y \to y = z \to x = z$\lbl{A1}}
\BIC{$x = y \to y = z \to x = z$}
\UIC{$\forall z\,(x = y \to y = z \to x = z)$}
\UIC{$\forall y\,\forall z\,(x = y \to y = z \to x = z)$}
\UIC{$\forall x\,\forall y\,\forall z\,(x = y \to y = z \to x = z)$}
\AXC{$\forall x\,\forall y\,\forall z\,(x = y \to y = z \to x = z) \to \varphi_{36}$\lbl{A2}}
\BIC{$\varphi_{36}$}
\end{ptree}
\noindent$\Pi_{5}$:
\begin{ptree}
\AXC{$x = x$\lbl{A5}}
\AXC{$x = y \to x = x \to y = x$\lbl{A6}}
\AXC{$(x = y \to x = x \to y = x) \to x = x \to x = y \to y = x$\lbl{A1}}
\BIC{$x = x \to x = y \to y = x$}
\BIC{$x = y \to y = x$}
\UIC{$\forall y\,(x = y \to y = x)$}
\UIC{$\forall x\,\forall y\,(x = y \to y = x)$}
\end{ptree}
\noindent$\Pi_{6}$:
\begin{ptree}
\AXC{$\stackrel{\Pi_{5}}{\vdots}$}\noLine
\UIC{$\forall x\,\forall y\,(x = y \to y = x)$}
\AXC{$\forall x\,\forall y\,(x = y \to y = x) \to \forall y\,(L = y \to y = L)$\lbl{A2}}
\BIC{$\forall y\,(L = y \to y = L)$}
\end{ptree}
\noindent$\Pi_{7}$:
\begin{ptree}
\AXC{$\forall x \in \mathbb{R}\,(\lvert x \rvert \in \mathbb{R})$\lbl{[A5]}}
\AXC{$\forall x \in \mathbb{R}\,(\lvert x \rvert \in \mathbb{R}) \to r \in \mathbb{R} \to \lvert r \rvert \in \mathbb{R}$\lbl{A2}}
\BIC{$r \in \mathbb{R} \to \lvert r \rvert \in \mathbb{R}$}
\AXC{$(r \in \mathbb{R} \to \lvert r \rvert \in \mathbb{R}) \to \varphi_{55} \to r \in \mathbb{R} \to \varphi_{33}$\lbl{A1}}
\BIC{$\varphi_{55} \to r \in \mathbb{R} \to \varphi_{33}$}
\end{ptree}
\noindent$\Pi_{8}$:
\begin{ptree}
\AXC{$\varphi_{55}$\lbl{A1}}
\AXC{$\stackrel{\Pi_{7}}{\vdots}$}\noLine
\UIC{$\varphi_{55} \to r \in \mathbb{R} \to \varphi_{33}$}
\BIC{$r \in \mathbb{R} \to \varphi_{33}$}
\end{ptree}
\noindent$\Pi_{9}$:
\begin{ptree}
\AXC{$\stackrel{\Pi_{8}}{\vdots}$}\noLine
\UIC{$r \in \mathbb{R} \to \varphi_{33}$}
\AXC{$\varphi_{59}$\lbl{定理2.1(2)}}
\AXC{$\varphi_{59} \to \varphi_{98}$\lbl{A2}}
\BIC{$\varphi_{98}$}
\AXC{$\varphi_{98} \to \varphi_{139}$\lbl{A2}}
\BIC{$\varphi_{139}$}
\AXC{$\varphi_{139} \to \lvert r \rvert \in \mathbb{R} \to \varphi_{71}$\lbl{A2}}
\BIC{$\lvert r \rvert \in \mathbb{R} \to \varphi_{71}$}
\AXC{$(\lvert r \rvert \in \mathbb{R} \to \varphi_{71}) \to \varphi_{234}$\lbl{A1}}
\BIC{$\varphi_{234}$}
\BIC{$(r \in \mathbb{R} \to \lvert r \rvert < 1 \to \varphi_{9}) \to r \in \mathbb{R} \to \lvert r \rvert < 1 \to \varphi_{71}$}
\end{ptree}
\noindent$\Pi_{10}$:
\begin{ptree}
\AXC{$r \in \mathbb{R} \to \lvert r \rvert < 1 \to \varphi_{9}$\lbl{A1}}
\AXC{$\stackrel{\Pi_{9}}{\vdots}$}\noLine
\UIC{$(r \in \mathbb{R} \to \lvert r \rvert < 1 \to \varphi_{9}) \to r \in \mathbb{R} \to \lvert r \rvert < 1 \to \varphi_{71}$}
\BIC{$r \in \mathbb{R} \to \lvert r \rvert < 1 \to \varphi_{71}$}
\end{ptree}
\noindent$\Pi_{11}$:
\begin{ptree}
\AXC{$y = z \to x = y \to x = z$\lbl{A6}}
\AXC{$(y = z \to x = y \to x = z) \to x = y \to y = z \to x = z$\lbl{A1}}
\BIC{$x = y \to y = z \to x = z$}
\UIC{$\forall z\,(x = y \to y = z \to x = z)$}
\UIC{$\forall y\,\forall z\,(x = y \to y = z \to x = z)$}
\UIC{$\forall x\,\forall y\,\forall z\,(x = y \to y = z \to x = z)$}
\AXC{$\forall x\,\forall y\,\forall z\,(x = y \to y = z \to x = z) \to \varphi_{117}$\lbl{A2}}
\BIC{$\varphi_{117}$}
\end{ptree}
\noindent$\Pi_{12}$:
\begin{ptree}
\AXC{$\stackrel{\Pi_{8}}{\vdots}$}\noLine
\UIC{$r \in \mathbb{R} \to \varphi_{33}$}
\AXC{$\forall r \in \mathbb{R}\,(\varphi_{127})$\lbl{[S24]}}
\AXC{$\forall r \in \mathbb{R}\,(\varphi_{127}) \to \lvert r \rvert \in \mathbb{R} \to \varphi_{175}$\lbl{A2}}
\BIC{$\lvert r \rvert \in \mathbb{R} \to \varphi_{175}$}
\AXC{$(\lvert r \rvert \in \mathbb{R} \to \varphi_{175}) \to (r \in \mathbb{R} \to \varphi_{33}) \to \varphi_{213}$\lbl{A1}}
\BIC{$(r \in \mathbb{R} \to \varphi_{33}) \to \varphi_{213}$}
\BIC{$\varphi_{213}$}
\end{ptree}
\noindent$\Pi_{13}$:
\begin{ptree}
\AXC{$\stackrel{\Pi_{12}}{\vdots}$}\noLine
\UIC{$\varphi_{213}$}
\AXC{$\varphi_{213} \to r \in \mathbb{R} \to \lvert r \rvert < 1 \to \varphi_{4}$\lbl{A1}}
\BIC{$r \in \mathbb{R} \to \lvert r \rvert < 1 \to \varphi_{4}$}
\end{ptree}
\noindent$\Pi_{14}$:
\begin{ptree}
\AXC{$\stackrel{\Pi_{13}}{\vdots}$}\noLine
\UIC{$r \in \mathbb{R} \to \lvert r \rvert < 1 \to \varphi_{4}$}
\AXC{$\varphi_{143}$\lbl{[S12]}}
\AXC{$\varphi_{143} \to \varphi_{192}$\lbl{A2}}
\BIC{$\varphi_{192}$}
\AXC{$\varphi_{192} \to \mathrm{Seq}({(\lvert r \rvert)^{\bullet}}) \to \mathrm{NSeq}(\mathsf{s}) \to \varphi_{163}$\lbl{A2}}
\BIC{$\mathrm{Seq}({(\lvert r \rvert)^{\bullet}}) \to \mathrm{NSeq}(\mathsf{s}) \to \varphi_{163}$}
\AXC{$\varphi_{262}$\lbl{A1}}
\BIC{$(r \in \mathbb{R} \to \lvert r \rvert < 1 \to \varphi_{4}) \to (r \in \mathbb{R} \to \varphi_{18}) \to \varphi_{210}$}
\BIC{$(r \in \mathbb{R} \to \varphi_{18}) \to \varphi_{210}$}
\end{ptree}
\noindent$\Pi_{15}$:
\begin{ptree}
\AXC{$\mathrm{NSeq}(\mathsf{s}) \land \forall k \in \mathbb{N}\,(\mathsf{s}(k) = S(k))$\lbl{[S14]}}
\AXC{$\varphi_{167}$\lbl{A1}}
\BIC{$r \in \mathbb{R} \to \varphi_{18}$}
\AXC{$\stackrel{\Pi_{14}}{\vdots}$}\noLine
\UIC{$(r \in \mathbb{R} \to \varphi_{18}) \to \varphi_{210}$}
\BIC{$\varphi_{210}$}
\end{ptree}
\noindent$\Pi_{16}$:
\begin{ptree}
\AXC{$\varphi_{209}$\lbl{A2}}
\AXC{$\stackrel{\Pi_{15}}{\vdots}$}\noLine
\UIC{$\varphi_{210}$}
\AXC{$\varphi_{210} \to \varphi_{203}$\lbl{A1}}
\BIC{$\varphi_{203}$}
\AXC{$\varphi_{203} \to \varphi_{209} \to \varphi_{197}$\lbl{A1}}
\BIC{$\varphi_{209} \to \varphi_{197}$}
\BIC{$\varphi_{197}$}
\AXC{$\varphi_{197} \to (r \in \mathbb{R} \to \varphi_{50}) \to r \in \mathbb{R} \to \lvert r \rvert < 1 \to \varphi_{144}$\lbl{A1}}
\BIC{$(r \in \mathbb{R} \to \varphi_{50}) \to r \in \mathbb{R} \to \lvert r \rvert < 1 \to \varphi_{144}$}
\end{ptree}
\noindent$\Pi_{17}$:
\begin{ptree}
\AXC{$r \in \mathbb{R} \to \varphi_{50}$\lbl{A1}}
\AXC{$\stackrel{\Pi_{16}}{\vdots}$}\noLine
\UIC{$(r \in \mathbb{R} \to \varphi_{50}) \to r \in \mathbb{R} \to \lvert r \rvert < 1 \to \varphi_{144}$}
\BIC{$r \in \mathbb{R} \to \lvert r \rvert < 1 \to \varphi_{144}$}
\AXC{$(r \in \mathbb{R} \to \lvert r \rvert < 1 \to \varphi_{144}) \to \varphi_{179} \to r \in \mathbb{R} \to \lvert r \rvert < 1 \to \varphi_{145}$\lbl{A1}}
\BIC{$\varphi_{179} \to r \in \mathbb{R} \to \lvert r \rvert < 1 \to \varphi_{145}$}
\end{ptree}
\noindent$\Pi_{18}$:
\begin{ptree}
\AXC{$\stackrel{\Pi_{5}}{\vdots}$}\noLine
\UIC{$\forall x\,\forall y\,(x = y \to y = x)$}
\AXC{$\forall x\,\forall y\,(x = y \to y = x) \to \varphi_{78}$\lbl{A2}}
\BIC{$\varphi_{78}$}
\AXC{$\varphi_{78} \to \varphi_{179}$\lbl{A2}}
\BIC{$\varphi_{179}$}
\AXC{$\stackrel{\Pi_{17}}{\vdots}$}\noLine
\UIC{$\varphi_{179} \to r \in \mathbb{R} \to \lvert r \rvert < 1 \to \varphi_{145}$}
\BIC{$r \in \mathbb{R} \to \lvert r \rvert < 1 \to \varphi_{145}$}
\end{ptree}
\noindent$\Pi_{19}$:
\begin{ptree}
\AXC{$y = z \to x = y \to x = z$\lbl{A6}}
\AXC{$(y = z \to x = y \to x = z) \to x = y \to y = z \to x = z$\lbl{A1}}
\BIC{$x = y \to y = z \to x = z$}
\UIC{$\forall z\,(x = y \to y = z \to x = z)$}
\UIC{$\forall y\,\forall z\,(x = y \to y = z \to x = z)$}
\UIC{$\forall x\,\forall y\,\forall z\,(x = y \to y = z \to x = z)$}
\AXC{$\forall x\,\forall y\,\forall z\,(x = y \to y = z \to x = z) \to \varphi_{117}$\lbl{A2}}
\BIC{$\varphi_{117}$}
\end{ptree}
\noindent$\Pi_{20}$:
\begin{ptree}
\AXC{$\mathrm{NSeq}(\mathsf{s}) \land \forall k \in \mathbb{N}\,(\mathsf{s}(k) = S(k))$\lbl{[S14]}}
\AXC{$\varphi_{207}$\lbl{A1}}
\BIC{$r \in \mathbb{R} \to \lvert r \rvert < 1 \to \varphi_{75}$}
\AXC{$(r \in \mathbb{R} \to \lvert r \rvert < 1 \to \varphi_{75}) \to \varphi_{69} \to r \in \mathbb{R} \to \lvert r \rvert < 1 \to \varphi_{68}$\lbl{A1}}
\BIC{$\varphi_{69} \to r \in \mathbb{R} \to \lvert r \rvert < 1 \to \varphi_{68}$}
\end{ptree}
\noindent$\Pi_{21}$:
\begin{ptree}
\AXC{$r \in \mathbb{R} \to \varphi_{50}$\lbl{A1}}
\AXC{$\varphi_{69}$\lbl{A2}}
\AXC{$\stackrel{\Pi_{20}}{\vdots}$}\noLine
\UIC{$\varphi_{69} \to r \in \mathbb{R} \to \lvert r \rvert < 1 \to \varphi_{68}$}
\BIC{$r \in \mathbb{R} \to \lvert r \rvert < 1 \to \varphi_{68}$}
\AXC{$\varphi_{239}$\lbl{A1}}
\BIC{$(r \in \mathbb{R} \to \varphi_{50}) \to r \in \mathbb{R} \to \lvert r \rvert < 1 \to \varphi_{26}$}
\BIC{$r \in \mathbb{R} \to \lvert r \rvert < 1 \to \varphi_{26}$}
\AXC{$(r \in \mathbb{R} \to \lvert r \rvert < 1 \to \varphi_{26}) \to \varphi_{178}$\lbl{A1}}
\BIC{$\varphi_{178}$}
\end{ptree}
\noindent$\Pi_{22}$:
\begin{ptree}
\AXC{$\stackrel{\Pi_{5}}{\vdots}$}\noLine
\UIC{$\forall x\,\forall y\,(x = y \to y = x)$}
\AXC{$\varphi_{30}$\lbl{A2}}
\BIC{$\forall y\,(\mathsf{s}(n) = y \to y = \mathsf{s}(n))$}
\AXC{$\varphi_{89}$\lbl{A2}}
\BIC{$\mathsf{s}(n) = S(n) \to S(n) = \mathsf{s}(n)$}
\AXC{$\stackrel{\Pi_{21}}{\vdots}$}\noLine
\UIC{$\varphi_{178}$}
\BIC{$r \in \mathbb{R} \to \lvert r \rvert < 1 \to \varphi_{27}$}
\AXC{$\varphi_{257}$\lbl{A1}}
\BIC{${{\lvert r \rvert}^{S(n)}} = {{\lvert r \rvert}^{S(n)}} \to \varphi_{174} \to r \in \mathbb{R} \to \lvert r \rvert < 1 \to \varphi_{124}$}
\end{ptree}
\noindent$\Pi_{23}$:
\begin{ptree}
\AXC{$x = x$\lbl{A5}}
\UIC{$\forall x\,(x = x)$}
\AXC{$\forall x\,(x = x) \to {{\lvert r \rvert}^{S(n)}} = {{\lvert r \rvert}^{S(n)}}$\lbl{A2}}
\BIC{${{\lvert r \rvert}^{S(n)}} = {{\lvert r \rvert}^{S(n)}}$}
\AXC{$\stackrel{\Pi_{22}}{\vdots}$}\noLine
\UIC{${{\lvert r \rvert}^{S(n)}} = {{\lvert r \rvert}^{S(n)}} \to \varphi_{174} \to r \in \mathbb{R} \to \lvert r \rvert < 1 \to \varphi_{124}$}
\BIC{$\varphi_{174} \to r \in \mathbb{R} \to \lvert r \rvert < 1 \to \varphi_{124}$}
\end{ptree}
\noindent$\Pi_{24}$:
\begin{ptree}
\AXC{$z = w \to {{\lvert r \rvert}^{S(n)}} = {{\lvert r \rvert}^{z}} \to {{\lvert r \rvert}^{S(n)}} = {{\lvert r \rvert}^{w}}$\lbl{A6}}
\UIC{$\varphi_{140}$}
\UIC{$\forall z\,(\varphi_{140})$}
\AXC{$\forall z\,(\varphi_{140}) \to \varphi_{151}$\lbl{A2}}
\BIC{$\varphi_{151}$}
\AXC{$\varphi_{151} \to \varphi_{174}$\lbl{A2}}
\BIC{$\varphi_{174}$}
\AXC{$\stackrel{\Pi_{23}}{\vdots}$}\noLine
\UIC{$\varphi_{174} \to r \in \mathbb{R} \to \lvert r \rvert < 1 \to \varphi_{124}$}
\BIC{$r \in \mathbb{R} \to \lvert r \rvert < 1 \to \varphi_{124}$}
\end{ptree}
\noindent$\Pi_{25}$:
\begin{ptree}
\AXC{$y = z \to x = y \to x = z$\lbl{A6}}
\AXC{$(y = z \to x = y \to x = z) \to x = y \to y = z \to x = z$\lbl{A1}}
\BIC{$x = y \to y = z \to x = z$}
\UIC{$\forall z\,(x = y \to y = z \to x = z)$}
\UIC{$\forall y\,\forall z\,(x = y \to y = z \to x = z)$}
\UIC{$\forall x\,\forall y\,\forall z\,(x = y \to y = z \to x = z)$}
\AXC{$\forall x\,\forall y\,\forall z\,(x = y \to y = z \to x = z) \to \varphi_{117}$\lbl{A2}}
\BIC{$\varphi_{117}$}
\end{ptree}
\noindent$\Pi_{26}$:
\begin{ptree}
\AXC{$\stackrel{\Pi_{25}}{\vdots}$}\noLine
\UIC{$\varphi_{117}$}
\AXC{$\varphi_{117} \to \varphi_{198}$\lbl{A2}}
\BIC{$\varphi_{198}$}
\AXC{$\varphi_{198} \to \lvert r \rvert \odot {(\lvert r \rvert)^{\bullet}}(n) = {{\lvert r \rvert}^{n}} \cdot \lvert r \rvert \to \varphi_{166}$\lbl{A2}}
\BIC{$\lvert r \rvert \odot {(\lvert r \rvert)^{\bullet}}(n) = {{\lvert r \rvert}^{n}} \cdot \lvert r \rvert \to \varphi_{166}$}
\end{ptree}
\noindent$\Pi_{27}$:
\begin{ptree}
\AXC{$\forall x \in \mathbb{R}\,(\lvert x \rvert \in \mathbb{R})$\lbl{[A5]}}
\AXC{$\forall x \in \mathbb{R}\,(\lvert x \rvert \in \mathbb{R}) \to r \in \mathbb{R} \to \lvert r \rvert \in \mathbb{R}$\lbl{A2}}
\BIC{$r \in \mathbb{R} \to \lvert r \rvert \in \mathbb{R}$}
\end{ptree}
\noindent$\Pi_{28}$:
\begin{ptree}
\AXC{$r \in \mathbb{R} \to \varphi_{51}$\lbl{A1}}
\AXC{$\stackrel{\Pi_{27}}{\vdots}$}\noLine
\UIC{$r \in \mathbb{R} \to \lvert r \rvert \in \mathbb{R}$}
\AXC{$(r \in \mathbb{R} \to \lvert r \rvert \in \mathbb{R}) \to (r \in \mathbb{R} \to \varphi_{51}) \to r \in \mathbb{R} \to \varphi_{72}$\lbl{A1}}
\BIC{$(r \in \mathbb{R} \to \varphi_{51}) \to r \in \mathbb{R} \to \varphi_{72}$}
\BIC{$r \in \mathbb{R} \to \varphi_{72}$}
\end{ptree}
\noindent$\Pi_{29}$:
\begin{ptree}
\AXC{$\forall r \in \mathbb{R}\,(\varphi_{127})$\lbl{[S24]}}
\AXC{$\forall r \in \mathbb{R}\,(\varphi_{127}) \to \lvert r \rvert \in \mathbb{R} \to \varphi_{175}$\lbl{A2}}
\BIC{$\lvert r \rvert \in \mathbb{R} \to \varphi_{175}$}
\AXC{$(\lvert r \rvert \in \mathbb{R} \to \varphi_{175}) \to (r \in \mathbb{R} \to \varphi_{72}) \to r \in \mathbb{R} \to \varphi_{212}$\lbl{A1}}
\BIC{$(r \in \mathbb{R} \to \varphi_{72}) \to r \in \mathbb{R} \to \varphi_{212}$}
\end{ptree}
\noindent$\Pi_{30}$:
\begin{ptree}
\AXC{$\stackrel{\Pi_{28}}{\vdots}$}\noLine
\UIC{$r \in \mathbb{R} \to \varphi_{72}$}
\AXC{$\stackrel{\Pi_{29}}{\vdots}$}\noLine
\UIC{$(r \in \mathbb{R} \to \varphi_{72}) \to r \in \mathbb{R} \to \varphi_{212}$}
\BIC{$r \in \mathbb{R} \to \varphi_{212}$}
\AXC{$(r \in \mathbb{R} \to \varphi_{212}) \to r \in \mathbb{R} \to \lvert r \rvert < 1 \to \varphi_{155}$\lbl{A1}}
\BIC{$r \in \mathbb{R} \to \lvert r \rvert < 1 \to \varphi_{155}$}
\end{ptree}
\noindent$\Pi_{31}$:
\begin{ptree}
\AXC{$\varphi_{194}$\lbl{A2}}
\AXC{$\stackrel{\Pi_{30}}{\vdots}$}\noLine
\UIC{$r \in \mathbb{R} \to \lvert r \rvert < 1 \to \varphi_{155}$}
\AXC{$(r \in \mathbb{R} \to \lvert r \rvert < 1 \to \varphi_{155}) \to \varphi_{194} \to r \in \mathbb{R} \to \lvert r \rvert < 1 \to \varphi_{150}$\lbl{A1}}
\BIC{$\varphi_{194} \to r \in \mathbb{R} \to \lvert r \rvert < 1 \to \varphi_{150}$}
\BIC{$r \in \mathbb{R} \to \lvert r \rvert < 1 \to \varphi_{150}$}
\AXC{$\varphi_{251}$\lbl{A1}}
\BIC{$(r \in \mathbb{R} \to \varphi_{50}) \to r \in \mathbb{R} \to \lvert r \rvert < 1 \to \varphi_{125}$}
\end{ptree}
\noindent$\Pi_{32}$:
\begin{ptree}
\AXC{$r \in \mathbb{R} \to \varphi_{50}$\lbl{A1}}
\AXC{$\stackrel{\Pi_{31}}{\vdots}$}\noLine
\UIC{$(r \in \mathbb{R} \to \varphi_{50}) \to r \in \mathbb{R} \to \lvert r \rvert < 1 \to \varphi_{125}$}
\BIC{$r \in \mathbb{R} \to \lvert r \rvert < 1 \to \varphi_{125}$}
\AXC{$(r \in \mathbb{R} \to \lvert r \rvert < 1 \to \varphi_{125}) \to \varphi_{159} \to r \in \mathbb{R} \to \lvert r \rvert < 1 \to \varphi_{126}$\lbl{A1}}
\BIC{$\varphi_{159} \to r \in \mathbb{R} \to \lvert r \rvert < 1 \to \varphi_{126}$}
\end{ptree}
\noindent$\Pi_{33}$:
\begin{ptree}
\AXC{$\stackrel{\Pi_{5}}{\vdots}$}\noLine
\UIC{$\forall x\,\forall y\,(x = y \to y = x)$}
\AXC{$\varphi_{79}$\lbl{A2}}
\BIC{$\forall y\,({{\lvert r \rvert}^{S(n)}} = y \to y = {{\lvert r \rvert}^{S(n)}})$}
\AXC{$\forall y\,({{\lvert r \rvert}^{S(n)}} = y \to y = {{\lvert r \rvert}^{S(n)}}) \to \varphi_{159}$\lbl{A2}}
\BIC{$\varphi_{159}$}
\AXC{$\stackrel{\Pi_{32}}{\vdots}$}\noLine
\UIC{$\varphi_{159} \to r \in \mathbb{R} \to \lvert r \rvert < 1 \to \varphi_{126}$}
\BIC{$r \in \mathbb{R} \to \lvert r \rvert < 1 \to \varphi_{126}$}
\end{ptree}
\noindent$\Pi_{34}$:
\begin{ptree}
\AXC{$y = z \to x = y \to x = z$\lbl{A6}}
\AXC{$(y = z \to x = y \to x = z) \to x = y \to y = z \to x = z$\lbl{A1}}
\BIC{$x = y \to y = z \to x = z$}
\UIC{$\forall z\,(x = y \to y = z \to x = z)$}
\UIC{$\forall y\,\forall z\,(x = y \to y = z \to x = z)$}
\UIC{$\forall x\,\forall y\,\forall z\,(x = y \to y = z \to x = z)$}
\AXC{$\forall x\,\forall y\,\forall z\,(x = y \to y = z \to x = z) \to \varphi_{117}$\lbl{A2}}
\BIC{$\varphi_{117}$}
\end{ptree}
\noindent$\Pi_{35}$:
\begin{ptree}
\AXC{$\forall r \in \mathbb{R}\,(\varphi_{127})$\lbl{[S24]}}
\AXC{$\forall r \in \mathbb{R}\,(\varphi_{127}) \to \lvert r \rvert \in \mathbb{R} \to \varphi_{175}$\lbl{A2}}
\BIC{$\lvert r \rvert \in \mathbb{R} \to \varphi_{175}$}
\AXC{$(\lvert r \rvert \in \mathbb{R} \to \varphi_{175}) \to (r \in \mathbb{R} \to \varphi_{72}) \to r \in \mathbb{R} \to \varphi_{212}$\lbl{A1}}
\BIC{$(r \in \mathbb{R} \to \varphi_{72}) \to r \in \mathbb{R} \to \varphi_{212}$}
\end{ptree}
\noindent$\Pi_{36}$:
\begin{ptree}
\AXC{$\stackrel{\Pi_{28}}{\vdots}$}\noLine
\UIC{$r \in \mathbb{R} \to \varphi_{72}$}
\AXC{$\stackrel{\Pi_{35}}{\vdots}$}\noLine
\UIC{$(r \in \mathbb{R} \to \varphi_{72}) \to r \in \mathbb{R} \to \varphi_{212}$}
\BIC{$r \in \mathbb{R} \to \varphi_{212}$}
\AXC{$(r \in \mathbb{R} \to \varphi_{212}) \to r \in \mathbb{R} \to \lvert r \rvert < 1 \to \varphi_{23}$\lbl{A1}}
\BIC{$r \in \mathbb{R} \to \lvert r \rvert < 1 \to \varphi_{23}$}
\end{ptree}
\noindent$\Pi_{37}$:
\begin{ptree}
\AXC{$\forall a\,\forall n\,(\mathrm{Seq}(a) \to n \in \mathbb{N} \to a_{n} \in \mathbb{R})$\lbl{[S1]}}
\AXC{$\forall a\,\forall n\,(\mathrm{Seq}(a) \to n \in \mathbb{N} \to a_{n} \in \mathbb{R}) \to \varphi_{60}$\lbl{A2}}
\BIC{$\varphi_{60}$}
\end{ptree}
\noindent$\Pi_{38}$:
\begin{ptree}
\AXC{$\stackrel{\Pi_{36}}{\vdots}$}\noLine
\UIC{$r \in \mathbb{R} \to \lvert r \rvert < 1 \to \varphi_{23}$}
\AXC{$\stackrel{\Pi_{37}}{\vdots}$}\noLine
\UIC{$\varphi_{60}$}
\AXC{$\varphi_{60} \to \mathrm{Seq}({(\lvert r \rvert)^{\bullet}}) \to n \in \mathbb{N} \to {{\lvert r \rvert}^{n}} \in \mathbb{R}$\lbl{A2}}
\BIC{$\mathrm{Seq}({(\lvert r \rvert)^{\bullet}}) \to n \in \mathbb{N} \to {{\lvert r \rvert}^{n}} \in \mathbb{R}$}
\AXC{$\varphi_{249}$\lbl{A1}}
\BIC{$\varphi_{237}$}
\BIC{$(r \in \mathbb{R} \to \varphi_{50}) \to r \in \mathbb{R} \to \lvert r \rvert < 1 \to \varphi_{62}$}
\end{ptree}
\noindent$\Pi_{39}$:
\begin{ptree}
\AXC{$\forall x\,\forall y\,(x \in \mathbb{R} \to y \in \mathbb{R} \to x \cdot y = y \cdot x)$\lbl{(R7)}}
\AXC{$\varphi_{168}$\lbl{A2}}
\BIC{$\forall y\,(\lvert r \rvert \in \mathbb{R} \to y \in \mathbb{R} \to \lvert r \rvert \cdot y = y \cdot \lvert r \rvert)$}
\AXC{$\forall y\,(\lvert r \rvert \in \mathbb{R} \to y \in \mathbb{R} \to \lvert r \rvert \cdot y = y \cdot \lvert r \rvert) \to \varphi_{164}$\lbl{A2}}
\BIC{$\varphi_{164}$}
\end{ptree}
\noindent$\Pi_{40}$:
\begin{ptree}
\AXC{$r \in \mathbb{R} \to \varphi_{50}$\lbl{A1}}
\AXC{$\stackrel{\Pi_{38}}{\vdots}$}\noLine
\UIC{$(r \in \mathbb{R} \to \varphi_{50}) \to r \in \mathbb{R} \to \lvert r \rvert < 1 \to \varphi_{62}$}
\BIC{$r \in \mathbb{R} \to \lvert r \rvert < 1 \to \varphi_{62}$}
\AXC{$\stackrel{\Pi_{28}}{\vdots}$}\noLine
\UIC{$r \in \mathbb{R} \to \varphi_{72}$}
\AXC{$\stackrel{\Pi_{39}}{\vdots}$}\noLine
\UIC{$\varphi_{164}$}
\AXC{$\varphi_{164} \to \varphi_{247}$\lbl{A1}}
\BIC{$\varphi_{247}$}
\BIC{$(r \in \mathbb{R} \to \lvert r \rvert < 1 \to \varphi_{62}) \to r \in \mathbb{R} \to \lvert r \rvert < 1 \to \varphi_{142}$}
\BIC{$r \in \mathbb{R} \to \lvert r \rvert < 1 \to \varphi_{142}$}
\end{ptree}
\noindent$\Pi_{41}$:
\begin{ptree}
\AXC{$\stackrel{\Pi_{28}}{\vdots}$}\noLine
\UIC{$r \in \mathbb{R} \to \varphi_{72}$}
\AXC{$\varphi_{146}$\lbl{[S3]}}
\AXC{$\varphi_{146} \to \varphi_{171}$\lbl{A2}}
\BIC{$\varphi_{171}$}
\AXC{$\varphi_{171} \to \lvert r \rvert \in \mathbb{R} \to \varphi_{193}$\lbl{A2}}
\BIC{$\lvert r \rvert \in \mathbb{R} \to \varphi_{193}$}
\AXC{$(\lvert r \rvert \in \mathbb{R} \to \varphi_{193}) \to \varphi_{253}$\lbl{A1}}
\BIC{$\varphi_{253}$}
\BIC{$(r \in \mathbb{R} \to \lvert r \rvert < 1 \to \varphi_{4}) \to r \in \mathbb{R} \to \lvert r \rvert < 1 \to \varphi_{205}$}
\end{ptree}
\noindent$\Pi_{42}$:
\begin{ptree}
\AXC{$\stackrel{\Pi_{13}}{\vdots}$}\noLine
\UIC{$r \in \mathbb{R} \to \lvert r \rvert < 1 \to \varphi_{4}$}
\AXC{$\stackrel{\Pi_{41}}{\vdots}$}\noLine
\UIC{$(r \in \mathbb{R} \to \lvert r \rvert < 1 \to \varphi_{4}) \to r \in \mathbb{R} \to \lvert r \rvert < 1 \to \varphi_{205}$}
\BIC{$r \in \mathbb{R} \to \lvert r \rvert < 1 \to \varphi_{205}$}
\AXC{$(r \in \mathbb{R} \to \lvert r \rvert < 1 \to \varphi_{205}) \to r \in \mathbb{R} \to \varphi_{191}$\lbl{A1}}
\BIC{$r \in \mathbb{R} \to \varphi_{191}$}
\end{ptree}
\noindent$\Pi_{43}$:
\begin{ptree}
\AXC{$\varphi_{103} \to \varphi_{93}$\lbl{A2}}
\AXC{$\stackrel{\Pi_{42}}{\vdots}$}\noLine
\UIC{$r \in \mathbb{R} \to \varphi_{191}$}
\AXC{$(r \in \mathbb{R} \to \varphi_{191}) \to (\varphi_{103} \to \varphi_{93}) \to r \in \mathbb{R} \to \varphi_{182}$\lbl{A1}}
\BIC{$(\varphi_{103} \to \varphi_{93}) \to r \in \mathbb{R} \to \varphi_{182}$}
\BIC{$r \in \mathbb{R} \to \varphi_{182}$}
\AXC{$(r \in \mathbb{R} \to \varphi_{182}) \to (r \in \mathbb{R} \to \varphi_{50}) \to \varphi_{183}$\lbl{A1}}
\BIC{$(r \in \mathbb{R} \to \varphi_{50}) \to \varphi_{183}$}
\end{ptree}
\noindent$\Pi_{44}$:
\begin{ptree}
\AXC{$\stackrel{\Pi_{40}}{\vdots}$}\noLine
\UIC{$r \in \mathbb{R} \to \lvert r \rvert < 1 \to \varphi_{142}$}
\AXC{$r \in \mathbb{R} \to \varphi_{50}$\lbl{A1}}
\AXC{$\stackrel{\Pi_{43}}{\vdots}$}\noLine
\UIC{$(r \in \mathbb{R} \to \varphi_{50}) \to \varphi_{183}$}
\BIC{$\varphi_{183}$}
\AXC{$\varphi_{183} \to (r \in \mathbb{R} \to \lvert r \rvert < 1 \to \varphi_{142}) \to \varphi_{250}$\lbl{A1}}
\BIC{$(r \in \mathbb{R} \to \lvert r \rvert < 1 \to \varphi_{142}) \to \varphi_{250}$}
\BIC{$\varphi_{250}$}
\end{ptree}
\noindent$\Pi_{45}$:
\begin{ptree}
\AXC{$\stackrel{\Pi_{34}}{\vdots}$}\noLine
\UIC{$\varphi_{117}$}
\AXC{$\varphi_{117} \to \varphi_{199}$\lbl{A2}}
\BIC{$\varphi_{199}$}
\AXC{$\varphi_{199} \to \lvert r \rvert \odot {(\lvert r \rvert)^{\bullet}}(n) = \lvert r \rvert \cdot {{\lvert r \rvert}^{n}} \to \varphi_{180}$\lbl{A2}}
\BIC{$\lvert r \rvert \odot {(\lvert r \rvert)^{\bullet}}(n) = \lvert r \rvert \cdot {{\lvert r \rvert}^{n}} \to \varphi_{180}$}
\AXC{$\stackrel{\Pi_{44}}{\vdots}$}\noLine
\UIC{$\varphi_{250}$}
\BIC{$\varphi_{184}$}
\end{ptree}
\noindent$\Pi_{46}$:
\begin{ptree}
\AXC{$\stackrel{\Pi_{33}}{\vdots}$}\noLine
\UIC{$r \in \mathbb{R} \to \lvert r \rvert < 1 \to \varphi_{126}$}
\AXC{$\stackrel{\Pi_{45}}{\vdots}$}\noLine
\UIC{$\varphi_{184}$}
\AXC{$\varphi_{184} \to (r \in \mathbb{R} \to \lvert r \rvert < 1 \to \varphi_{126}) \to \varphi_{244}$\lbl{A1}}
\BIC{$(r \in \mathbb{R} \to \lvert r \rvert < 1 \to \varphi_{126}) \to \varphi_{244}$}
\BIC{$\varphi_{244}$}
\end{ptree}
\noindent$\Pi_{47}$:
\begin{ptree}
\AXC{$\stackrel{\Pi_{24}}{\vdots}$}\noLine
\UIC{$r \in \mathbb{R} \to \lvert r \rvert < 1 \to \varphi_{124}$}
\AXC{$\stackrel{\Pi_{26}}{\vdots}$}\noLine
\UIC{$\lvert r \rvert \odot {(\lvert r \rvert)^{\bullet}}(n) = {{\lvert r \rvert}^{n}} \cdot \lvert r \rvert \to \varphi_{166}$}
\AXC{$\stackrel{\Pi_{46}}{\vdots}$}\noLine
\UIC{$\varphi_{244}$}
\BIC{$r \in \mathbb{R} \to \lvert r \rvert < 1 \to \varphi_{135}$}
\AXC{$\varphi_{263}$\lbl{A1}}
\BIC{$(r \in \mathbb{R} \to \lvert r \rvert < 1 \to \varphi_{124}) \to \varphi_{248}$}
\BIC{$\varphi_{248}$}
\end{ptree}
\noindent$\Pi_{48}$:
\begin{ptree}
\AXC{$\stackrel{\Pi_{19}}{\vdots}$}\noLine
\UIC{$\varphi_{117}$}
\AXC{$\varphi_{117} \to \varphi_{177}$\lbl{A2}}
\BIC{$\varphi_{177}$}
\AXC{$\varphi_{177} \to \lvert r \rvert \odot {(\lvert r \rvert)^{\bullet}}(n) = {{\lvert r \rvert}^{S(n)}} \to \varphi_{172}$\lbl{A2}}
\BIC{$\lvert r \rvert \odot {(\lvert r \rvert)^{\bullet}}(n) = {{\lvert r \rvert}^{S(n)}} \to \varphi_{172}$}
\AXC{$\stackrel{\Pi_{47}}{\vdots}$}\noLine
\UIC{$\varphi_{248}$}
\BIC{$r \in \mathbb{R} \to \lvert r \rvert < 1 \to \varphi_{148}$}
\AXC{$\varphi_{264}$\lbl{A1}}
\BIC{$(r \in \mathbb{R} \to \lvert r \rvert < 1 \to \varphi_{145}) \to \varphi_{222} \to r \in \mathbb{R} \to \varphi_{170}$}
\end{ptree}
\noindent$\Pi_{49}$:
\begin{ptree}
\AXC{$\stackrel{\Pi_{11}}{\vdots}$}\noLine
\UIC{$\varphi_{117}$}
\AXC{$\varphi_{117} \to \varphi_{196}$\lbl{A2}}
\BIC{$\varphi_{196}$}
\AXC{$\varphi_{196} \to \varphi_{222}$\lbl{A2}}
\BIC{$\varphi_{222}$}
\AXC{$\stackrel{\Pi_{18}}{\vdots}$}\noLine
\UIC{$r \in \mathbb{R} \to \lvert r \rvert < 1 \to \varphi_{145}$}
\AXC{$\stackrel{\Pi_{48}}{\vdots}$}\noLine
\UIC{$(r \in \mathbb{R} \to \lvert r \rvert < 1 \to \varphi_{145}) \to \varphi_{222} \to r \in \mathbb{R} \to \varphi_{170}$}
\BIC{$\varphi_{222} \to r \in \mathbb{R} \to \varphi_{170}$}
\BIC{$r \in \mathbb{R} \to \varphi_{170}$}
\end{ptree}
\noindent$\Pi_{50}$:
\begin{ptree}
\AXC{$\varphi_{153}$\lbl{A4}}
\AXC{$\forall n\,(\varphi_{187}) \to \varphi_{22} \to \forall n \in \mathbb{N}\,(\varphi_{61})$\lbl{A3}}
\AXC{$\stackrel{\Pi_{49}}{\vdots}$}\noLine
\UIC{$r \in \mathbb{R} \to \varphi_{170}$}
\AXC{$(r \in \mathbb{R} \to \varphi_{170}) \to \varphi_{187}$\lbl{A1}}
\BIC{$\varphi_{187}$}
\UIC{$\forall n\,(\varphi_{187})$}
\AXC{$\varphi_{266}$\lbl{A1}}
\BIC{$(\forall n\,(\varphi_{187}) \to \varphi_{22} \to \forall n \in \mathbb{N}\,(\varphi_{61})) \to \varphi_{153} \to r \in \mathbb{R} \to \varphi_{173}$}
\BIC{$\varphi_{153} \to r \in \mathbb{R} \to \varphi_{173}$}
\BIC{$r \in \mathbb{R} \to \varphi_{173}$}
\end{ptree}
\noindent$\Pi_{51}$:
\begin{ptree}
\AXC{$\varphi_{112}$\lbl{[C13]}}
\AXC{$\varphi_{112} \to \varphi_{185}$\lbl{A2}}
\BIC{$\varphi_{185}$}
\AXC{$\varphi_{185} \to \varphi_{220}$\lbl{A2}}
\BIC{$\varphi_{220}$}
\AXC{$\varphi_{220} \to \mathrm{Seq}({(\lvert r \rvert)^{\bullet}} \circ \mathsf{s}) \to \forall n \in \mathbb{N}\,(\varphi_{61}) \to \varphi_{132}$\lbl{A2}}
\BIC{$\mathrm{Seq}({(\lvert r \rvert)^{\bullet}} \circ \mathsf{s}) \to \forall n \in \mathbb{N}\,(\varphi_{61}) \to \varphi_{132}$}
\AXC{$\varphi_{265}$\lbl{A1}}
\BIC{$(r \in \mathbb{R} \to \lvert r \rvert < 1 \to \varphi_{20}) \to \varphi_{252}$}
\end{ptree}
\noindent$\Pi_{52}$:
\begin{ptree}
\AXC{$\stackrel{\Pi_{50}}{\vdots}$}\noLine
\UIC{$r \in \mathbb{R} \to \varphi_{173}$}
\AXC{$\stackrel{\Pi_{15}}{\vdots}$}\noLine
\UIC{$\varphi_{210}$}
\AXC{$\varphi_{210} \to r \in \mathbb{R} \to \lvert r \rvert < 1 \to \varphi_{20}$\lbl{A1}}
\BIC{$r \in \mathbb{R} \to \lvert r \rvert < 1 \to \varphi_{20}$}
\AXC{$\stackrel{\Pi_{51}}{\vdots}$}\noLine
\UIC{$(r \in \mathbb{R} \to \lvert r \rvert < 1 \to \varphi_{20}) \to \varphi_{252}$}
\BIC{$\varphi_{252}$}
\BIC{$(r \in \mathbb{R} \to \lvert r \rvert < 1 \to \varphi_{71}) \to r \in \mathbb{R} \to \lvert r \rvert < 1 \to \varphi_{66}$}
\end{ptree}
\noindent$\Pi_{53}$:
\begin{ptree}
\AXC{$r \in \mathbb{R} \to \lvert r \rvert < 1 \to \varphi_{9}$\lbl{A1}}
\AXC{$\varphi_{39}$\lbl{[C18]}}
\AXC{$\varphi_{39} \to \varphi_{80}$\lbl{A2}}
\BIC{$\varphi_{80}$}
\AXC{$\varphi_{80} \to \varphi_{99}$\lbl{A2}}
\BIC{$\varphi_{99}$}
\AXC{$\varphi_{99} \to \varphi_{77}$\lbl{A2}}
\BIC{$\varphi_{77}$}
\AXC{$\varphi_{77} \to \varphi_{216}$\lbl{A1}}
\BIC{$\varphi_{216}$}
\BIC{$\mathrm{SubIdx}(\mathsf{s}) \to r \in \mathbb{R} \to \lvert r \rvert < 1 \to \varphi_{34}$}
\end{ptree}
\noindent$\Pi_{54}$:
\begin{ptree}
\AXC{$\varphi_{13}$\lbl{命題1.3}}
\AXC{$\varphi_{13} \to \varphi_{118}$\lbl{A2}}
\BIC{$\varphi_{118}$}
\AXC{$\varphi_{118} \to \varphi_{101}$\lbl{A2}}
\BIC{$\varphi_{101}$}
\AXC{$\varphi_{101} \to \mathrm{Conv}({(\lvert r \rvert)^{\bullet}} \circ \mathsf{s},L) \to \varphi_{46}$\lbl{A2}}
\BIC{$\mathrm{Conv}({(\lvert r \rvert)^{\bullet}} \circ \mathsf{s},L) \to \varphi_{46}$}
\end{ptree}
\noindent$\Pi_{55}$:
\begin{ptree}
\AXC{$\stackrel{\Pi_{54}}{\vdots}$}\noLine
\UIC{$\mathrm{Conv}({(\lvert r \rvert)^{\bullet}} \circ \mathsf{s},L) \to \varphi_{46}$}
\AXC{$(\mathrm{Conv}({(\lvert r \rvert)^{\bullet}} \circ \mathsf{s},L) \to \varphi_{46}) \to \varphi_{236}$\lbl{A1}}
\BIC{$\varphi_{236}$}
\end{ptree}
\noindent$\Pi_{56}$:
\begin{ptree}
\AXC{$\mathrm{SubIdx}(\mathsf{s})$\lbl{[C28]}}
\AXC{$\stackrel{\Pi_{53}}{\vdots}$}\noLine
\UIC{$\mathrm{SubIdx}(\mathsf{s}) \to r \in \mathbb{R} \to \lvert r \rvert < 1 \to \varphi_{34}$}
\BIC{$r \in \mathbb{R} \to \lvert r \rvert < 1 \to \varphi_{34}$}
\AXC{$\stackrel{\Pi_{55}}{\vdots}$}\noLine
\UIC{$\varphi_{236}$}
\BIC{$(r \in \mathbb{R} \to \lvert r \rvert < 1 \to \varphi_{66}) \to r \in \mathbb{R} \to \varphi_{24}$}
\end{ptree}
\noindent$\Pi_{57}$:
\begin{ptree}
\AXC{$\stackrel{\Pi_{10}}{\vdots}$}\noLine
\UIC{$r \in \mathbb{R} \to \lvert r \rvert < 1 \to \varphi_{71}$}
\AXC{$\stackrel{\Pi_{52}}{\vdots}$}\noLine
\UIC{$(r \in \mathbb{R} \to \lvert r \rvert < 1 \to \varphi_{71}) \to r \in \mathbb{R} \to \lvert r \rvert < 1 \to \varphi_{66}$}
\BIC{$r \in \mathbb{R} \to \lvert r \rvert < 1 \to \varphi_{66}$}
\AXC{$\stackrel{\Pi_{56}}{\vdots}$}\noLine
\UIC{$(r \in \mathbb{R} \to \lvert r \rvert < 1 \to \varphi_{66}) \to r \in \mathbb{R} \to \varphi_{24}$}
\BIC{$r \in \mathbb{R} \to \varphi_{24}$}
\end{ptree}
\noindent$\Pi_{58}$:
\begin{ptree}
\AXC{$\stackrel{\Pi_{6}}{\vdots}$}\noLine
\UIC{$\forall y\,(L = y \to y = L)$}
\AXC{$\forall y\,(L = y \to y = L) \to L = \lvert r \rvert \cdot L \to \lvert r \rvert \cdot L = L$\lbl{A2}}
\BIC{$L = \lvert r \rvert \cdot L \to \lvert r \rvert \cdot L = L$}
\AXC{$\stackrel{\Pi_{57}}{\vdots}$}\noLine
\UIC{$r \in \mathbb{R} \to \varphi_{24}$}
\AXC{$\varphi_{211}$\lbl{A1}}
\BIC{$(L = \lvert r \rvert \cdot L \to \lvert r \rvert \cdot L = L) \to r \in \mathbb{R} \to \varphi_{25}$}
\BIC{$r \in \mathbb{R} \to \varphi_{25}$}
\end{ptree}
\noindent$\Pi_{59}$:
\begin{ptree}
\AXC{$x = x$\lbl{A5}}
\UIC{$\forall x\,(x = x)$}
\AXC{$\forall x\,(x = x) \to L - \lvert r \rvert \cdot L = L - \lvert r \rvert \cdot L$\lbl{A2}}
\BIC{$L - \lvert r \rvert \cdot L = L - \lvert r \rvert \cdot L$}
\AXC{$\stackrel{\Pi_{58}}{\vdots}$}\noLine
\UIC{$r \in \mathbb{R} \to \varphi_{25}$}
\AXC{$\varphi_{240}$\lbl{A1}}
\BIC{$L - \lvert r \rvert \cdot L = L - \lvert r \rvert \cdot L \to \varphi_{106} \to r \in \mathbb{R} \to \varphi_{57}$}
\BIC{$\varphi_{106} \to r \in \mathbb{R} \to \varphi_{57}$}
\end{ptree}
\noindent$\Pi_{60}$:
\begin{ptree}
\AXC{$z = w \to L - \lvert r \rvert \cdot L = L - z \to L - \lvert r \rvert \cdot L = L - w$\lbl{A6}}
\UIC{$\varphi_{54}$}
\UIC{$\forall z\,(\varphi_{54})$}
\AXC{$\forall z\,(\varphi_{54}) \to \varphi_{114}$\lbl{A2}}
\BIC{$\varphi_{114}$}
\AXC{$\varphi_{114} \to \varphi_{106}$\lbl{A2}}
\BIC{$\varphi_{106}$}
\AXC{$\stackrel{\Pi_{59}}{\vdots}$}\noLine
\UIC{$\varphi_{106} \to r \in \mathbb{R} \to \varphi_{57}$}
\BIC{$r \in \mathbb{R} \to \varphi_{57}$}
\end{ptree}
\noindent$\Pi_{61}$:
\begin{ptree}
\AXC{$y = z \to x = y \to x = z$\lbl{A6}}
\AXC{$(y = z \to x = y \to x = z) \to x = y \to y = z \to x = z$\lbl{A1}}
\BIC{$x = y \to y = z \to x = z$}
\UIC{$\forall z\,(x = y \to y = z \to x = z)$}
\UIC{$\forall y\,\forall z\,(x = y \to y = z \to x = z)$}
\UIC{$\forall x\,\forall y\,\forall z\,(x = y \to y = z \to x = z)$}
\AXC{$\forall x\,\forall y\,\forall z\,(x = y \to y = z \to x = z) \to \varphi_{36}$\lbl{A2}}
\BIC{$\varphi_{36}$}
\end{ptree}
\noindent$\Pi_{62}$:
\begin{ptree}
\AXC{$z = w \to L - L \cdot \lvert r \rvert = L - z \to L - L \cdot \lvert r \rvert = L - w$\lbl{A6}}
\UIC{$\varphi_{56}$}
\UIC{$\forall z\,(\varphi_{56})$}
\AXC{$\forall z\,(\varphi_{56}) \to \varphi_{116}$\lbl{A2}}
\BIC{$\varphi_{116}$}
\AXC{$\varphi_{116} \to L \cdot \lvert r \rvert = \lvert r \rvert \cdot L \to \varphi_{84}$\lbl{A2}}
\BIC{$L \cdot \lvert r \rvert = \lvert r \rvert \cdot L \to \varphi_{84}$}
\end{ptree}
\noindent$\Pi_{63}$:
\begin{ptree}
\AXC{$\stackrel{\Pi_{8}}{\vdots}$}\noLine
\UIC{$r \in \mathbb{R} \to \varphi_{33}$}
\AXC{$\stackrel{\Pi_{2}}{\vdots}$}\noLine
\UIC{$\varphi_{53}$}
\AXC{$\forall x\,\forall y\,(x \in \mathbb{R} \to y \in \mathbb{R} \to x \cdot y = y \cdot x)$\lbl{(R7)}}
\AXC{$\varphi_{147}$\lbl{A2}}
\BIC{$\forall y\,(L \in \mathbb{R} \to y \in \mathbb{R} \to L \cdot y = y \cdot L)$}
\AXC{$\varphi_{160}$\lbl{A2}}
\BIC{$L \in \mathbb{R} \to \lvert r \rvert \in \mathbb{R} \to L \cdot \lvert r \rvert = \lvert r \rvert \cdot L$}
\AXC{$\varphi_{242}$\lbl{A1}}
\BIC{$\varphi_{53} \to (r \in \mathbb{R} \to \varphi_{33}) \to r \in \mathbb{R} \to \varphi_{58}$}
\BIC{$(r \in \mathbb{R} \to \varphi_{33}) \to r \in \mathbb{R} \to \varphi_{58}$}
\BIC{$r \in \mathbb{R} \to \varphi_{58}$}
\end{ptree}
\noindent$\Pi_{64}$:
\begin{ptree}
\AXC{$x = x$\lbl{A5}}
\UIC{$\forall x\,(x = x)$}
\AXC{$\forall x\,(x = x) \to L - L \cdot \lvert r \rvert = L - L \cdot \lvert r \rvert$\lbl{A2}}
\BIC{$L - L \cdot \lvert r \rvert = L - L \cdot \lvert r \rvert$}
\AXC{$\stackrel{\Pi_{63}}{\vdots}$}\noLine
\UIC{$r \in \mathbb{R} \to \varphi_{58}$}
\AXC{$(r \in \mathbb{R} \to \varphi_{58}) \to \varphi_{233}$\lbl{A1}}
\BIC{$\varphi_{233}$}
\BIC{$(L \cdot \lvert r \rvert = \lvert r \rvert \cdot L \to \varphi_{84}) \to r \in \mathbb{R} \to \lvert r \rvert < 1 \to \varphi_{42}$}
\end{ptree}
\noindent$\Pi_{65}$:
\begin{ptree}
\AXC{$y = z \to x = y \to x = z$\lbl{A6}}
\AXC{$(y = z \to x = y \to x = z) \to x = y \to y = z \to x = z$\lbl{A1}}
\BIC{$x = y \to y = z \to x = z$}
\UIC{$\forall z\,(x = y \to y = z \to x = z)$}
\UIC{$\forall y\,\forall z\,(x = y \to y = z \to x = z)$}
\UIC{$\forall x\,\forall y\,\forall z\,(x = y \to y = z \to x = z)$}
\AXC{$\forall x\,\forall y\,\forall z\,(x = y \to y = z \to x = z) \to \varphi_{36}$\lbl{A2}}
\BIC{$\varphi_{36}$}
\end{ptree}
\noindent$\Pi_{66}$:
\begin{ptree}
\AXC{$\stackrel{\Pi_{2}}{\vdots}$}\noLine
\UIC{$\varphi_{53}$}
\AXC{$\forall x \in \mathbb{R}\,(x \cdot 1 = x)$\lbl{(R8)}}
\AXC{$\forall x \in \mathbb{R}\,(x \cdot 1 = x) \to L \in \mathbb{R} \to L \cdot 1 = L$\lbl{A2}}
\BIC{$L \in \mathbb{R} \to L \cdot 1 = L$}
\AXC{$(L \in \mathbb{R} \to L \cdot 1 = L) \to \varphi_{53} \to \varphi_{49}$\lbl{A1}}
\BIC{$\varphi_{53} \to \varphi_{49}$}
\BIC{$\varphi_{49}$}
\end{ptree}
\noindent$\Pi_{67}$:
\begin{ptree}
\AXC{$x = x$\lbl{A5}}
\UIC{$\forall x\,(x = x)$}
\AXC{$\forall x\,(x = x) \to L \cdot 1 - L \cdot \lvert r \rvert = L \cdot 1 - L \cdot \lvert r \rvert$\lbl{A2}}
\BIC{$L \cdot 1 - L \cdot \lvert r \rvert = L \cdot 1 - L \cdot \lvert r \rvert$}
\AXC{$\stackrel{\Pi_{66}}{\vdots}$}\noLine
\UIC{$\varphi_{49}$}
\AXC{$\varphi_{49} \to \varphi_{235}$\lbl{A1}}
\BIC{$\varphi_{235}$}
\BIC{$(L \cdot 1 = L \to \varphi_{136}) \to r \in \mathbb{R} \to \lvert r \rvert < 1 \to \varphi_{63}$}
\end{ptree}
\noindent$\Pi_{68}$:
\begin{ptree}
\AXC{$z = w \to \varphi_{120}$\lbl{A6}}
\UIC{$\forall w\,(z = w \to \varphi_{120})$}
\UIC{$\forall z\,\forall w\,(z = w \to \varphi_{120})$}
\AXC{$\forall z\,\forall w\,(z = w \to \varphi_{120}) \to \varphi_{156}$\lbl{A2}}
\BIC{$\varphi_{156}$}
\AXC{$\varphi_{156} \to L \cdot 1 = L \to \varphi_{136}$\lbl{A2}}
\BIC{$L \cdot 1 = L \to \varphi_{136}$}
\AXC{$\stackrel{\Pi_{67}}{\vdots}$}\noLine
\UIC{$(L \cdot 1 = L \to \varphi_{136}) \to r \in \mathbb{R} \to \lvert r \rvert < 1 \to \varphi_{63}$}
\BIC{$r \in \mathbb{R} \to \lvert r \rvert < 1 \to \varphi_{63}$}
\end{ptree}
\noindent$\Pi_{69}$:
\begin{ptree}
\AXC{$\stackrel{\Pi_{5}}{\vdots}$}\noLine
\UIC{$\forall x\,\forall y\,(x = y \to y = x)$}
\AXC{$\varphi_{90}$\lbl{A2}}
\BIC{$\forall y\,(L \cdot 1 - L \cdot \lvert r \rvert = y \to y = L \cdot 1 - L \cdot \lvert r \rvert)$}
\AXC{$\forall y\,(L \cdot 1 - L \cdot \lvert r \rvert = y \to y = L \cdot 1 - L \cdot \lvert r \rvert) \to \varphi_{119}$\lbl{A2}}
\BIC{$\varphi_{119}$}
\end{ptree}
\noindent$\Pi_{70}$:
\begin{ptree}
\AXC{$1 \in \mathbb{R}$\lbl{(R1)}}
\AXC{$\stackrel{\Pi_{2}}{\vdots}$}\noLine
\UIC{$\varphi_{53}$}
\AXC{$\varphi_{110}$\lbl{[F17]}}
\AXC{$\varphi_{110} \to \varphi_{102}$\lbl{A2}}
\BIC{$\varphi_{102}$}
\AXC{$\varphi_{102} \to \varphi_{91}$\lbl{A2}}
\BIC{$\varphi_{91}$}
\AXC{$\varphi_{91} \to L \in \mathbb{R} \to \varphi_{85}$\lbl{A2}}
\BIC{$L \in \mathbb{R} \to \varphi_{85}$}
\AXC{$\varphi_{254}$\lbl{A1}}
\BIC{$\varphi_{53} \to 1 \in \mathbb{R} \to (r \in \mathbb{R} \to \varphi_{33}) \to r \in \mathbb{R} \to \lvert r \rvert < 1 \to \varphi_{64}$}
\BIC{$1 \in \mathbb{R} \to (r \in \mathbb{R} \to \varphi_{33}) \to r \in \mathbb{R} \to \lvert r \rvert < 1 \to \varphi_{64}$}
\BIC{$(r \in \mathbb{R} \to \varphi_{33}) \to r \in \mathbb{R} \to \lvert r \rvert < 1 \to \varphi_{64}$}
\end{ptree}
\noindent$\Pi_{71}$:
\begin{ptree}
\AXC{$\stackrel{\Pi_{8}}{\vdots}$}\noLine
\UIC{$r \in \mathbb{R} \to \varphi_{33}$}
\AXC{$\stackrel{\Pi_{70}}{\vdots}$}\noLine
\UIC{$(r \in \mathbb{R} \to \varphi_{33}) \to r \in \mathbb{R} \to \lvert r \rvert < 1 \to \varphi_{64}$}
\BIC{$r \in \mathbb{R} \to \lvert r \rvert < 1 \to \varphi_{64}$}
\AXC{$(r \in \mathbb{R} \to \lvert r \rvert < 1 \to \varphi_{64}) \to \varphi_{119} \to r \in \mathbb{R} \to \lvert r \rvert < 1 \to \varphi_{65}$\lbl{A1}}
\BIC{$\varphi_{119} \to r \in \mathbb{R} \to \lvert r \rvert < 1 \to \varphi_{65}$}
\end{ptree}
\noindent$\Pi_{72}$:
\begin{ptree}
\AXC{$\stackrel{\Pi_{68}}{\vdots}$}\noLine
\UIC{$r \in \mathbb{R} \to \lvert r \rvert < 1 \to \varphi_{63}$}
\AXC{$\stackrel{\Pi_{69}}{\vdots}$}\noLine
\UIC{$\varphi_{119}$}
\AXC{$\stackrel{\Pi_{71}}{\vdots}$}\noLine
\UIC{$\varphi_{119} \to r \in \mathbb{R} \to \lvert r \rvert < 1 \to \varphi_{65}$}
\BIC{$r \in \mathbb{R} \to \lvert r \rvert < 1 \to \varphi_{65}$}
\AXC{$\varphi_{260}$\lbl{A1}}
\BIC{$(r \in \mathbb{R} \to \lvert r \rvert < 1 \to \varphi_{63}) \to \varphi_{228}$}
\BIC{$\varphi_{228}$}
\end{ptree}
\noindent$\Pi_{73}$:
\begin{ptree}
\AXC{$\stackrel{\Pi_{65}}{\vdots}$}\noLine
\UIC{$\varphi_{36}$}
\AXC{$\varphi_{36} \to \varphi_{141}$\lbl{A2}}
\BIC{$\varphi_{141}$}
\AXC{$\varphi_{141} \to L \cdot (1 - \lvert r \rvert) = L \cdot 1 - L \cdot \lvert r \rvert \to \varphi_{105}$\lbl{A2}}
\BIC{$L \cdot (1 - \lvert r \rvert) = L \cdot 1 - L \cdot \lvert r \rvert \to \varphi_{105}$}
\AXC{$\stackrel{\Pi_{72}}{\vdots}$}\noLine
\UIC{$\varphi_{228}$}
\BIC{$r \in \mathbb{R} \to \lvert r \rvert < 1 \to \varphi_{43}$}
\AXC{$\varphi_{259}$\lbl{A1}}
\BIC{$(r \in \mathbb{R} \to \lvert r \rvert < 1 \to \varphi_{42}) \to \varphi_{225}$}
\end{ptree}
\noindent$\Pi_{74}$:
\begin{ptree}
\AXC{$\stackrel{\Pi_{62}}{\vdots}$}\noLine
\UIC{$L \cdot \lvert r \rvert = \lvert r \rvert \cdot L \to \varphi_{84}$}
\AXC{$\stackrel{\Pi_{64}}{\vdots}$}\noLine
\UIC{$(L \cdot \lvert r \rvert = \lvert r \rvert \cdot L \to \varphi_{84}) \to r \in \mathbb{R} \to \lvert r \rvert < 1 \to \varphi_{42}$}
\BIC{$r \in \mathbb{R} \to \lvert r \rvert < 1 \to \varphi_{42}$}
\AXC{$\stackrel{\Pi_{73}}{\vdots}$}\noLine
\UIC{$(r \in \mathbb{R} \to \lvert r \rvert < 1 \to \varphi_{42}) \to \varphi_{225}$}
\BIC{$\varphi_{225}$}
\end{ptree}
\noindent$\Pi_{75}$:
\begin{ptree}
\AXC{$\stackrel{\Pi_{60}}{\vdots}$}\noLine
\UIC{$r \in \mathbb{R} \to \varphi_{57}$}
\AXC{$\stackrel{\Pi_{61}}{\vdots}$}\noLine
\UIC{$\varphi_{36}$}
\AXC{$\varphi_{36} \to \varphi_{115}$\lbl{A2}}
\BIC{$\varphi_{115}$}
\AXC{$\varphi_{115} \to L \cdot (1 - \lvert r \rvert) = L - L \cdot \lvert r \rvert \to \varphi_{83}$\lbl{A2}}
\BIC{$L \cdot (1 - \lvert r \rvert) = L - L \cdot \lvert r \rvert \to \varphi_{83}$}
\AXC{$\stackrel{\Pi_{74}}{\vdots}$}\noLine
\UIC{$\varphi_{225}$}
\BIC{$r \in \mathbb{R} \to \lvert r \rvert < 1 \to \varphi_{44}$}
\AXC{$\varphi_{256}$\lbl{A1}}
\BIC{$(r \in \mathbb{R} \to \varphi_{57}) \to \varphi_{138} \to r \in \mathbb{R} \to \lvert r \rvert < 1 \to \varphi_{16}$}
\BIC{$\varphi_{138} \to r \in \mathbb{R} \to \lvert r \rvert < 1 \to \varphi_{16}$}
\end{ptree}
\noindent$\Pi_{76}$:
\begin{ptree}
\AXC{$\stackrel{\Pi_{4}}{\vdots}$}\noLine
\UIC{$\varphi_{36}$}
\AXC{$\varphi_{36} \to \varphi_{113}$\lbl{A2}}
\BIC{$\varphi_{113}$}
\AXC{$\varphi_{113} \to \varphi_{138}$\lbl{A2}}
\BIC{$\varphi_{138}$}
\AXC{$\stackrel{\Pi_{75}}{\vdots}$}\noLine
\UIC{$\varphi_{138} \to r \in \mathbb{R} \to \lvert r \rvert < 1 \to \varphi_{16}$}
\BIC{$r \in \mathbb{R} \to \lvert r \rvert < 1 \to \varphi_{16}$}
\AXC{$(r \in \mathbb{R} \to \lvert r \rvert < 1 \to \varphi_{16}) \to \varphi_{48} \to \varphi_{45} \to r \in \mathbb{R} \to \varphi_{41}$\lbl{A1}}
\BIC{$\varphi_{48} \to \varphi_{45} \to r \in \mathbb{R} \to \varphi_{41}$}
\end{ptree}
\noindent$\Pi_{77}$:
\begin{ptree}
\AXC{$\forall x\,\forall y\,(x \in \mathbb{R} \to y \in \mathbb{R} \to x < y \to 0 < y - x)$\lbl{[F24]}}
\AXC{$\varphi_{188}$\lbl{A2}}
\BIC{$\forall y\,(\lvert r \rvert \in \mathbb{R} \to y \in \mathbb{R} \to \lvert r \rvert < y \to 0 < y - \lvert r \rvert)$}
\AXC{$\varphi_{200}$\lbl{A2}}
\BIC{$\lvert r \rvert \in \mathbb{R} \to 1 \in \mathbb{R} \to \lvert r \rvert < 1 \to 0 < 1 - \lvert r \rvert$}
\end{ptree}
\noindent$\Pi_{78}$:
\begin{ptree}
\AXC{$\stackrel{\Pi_{8}}{\vdots}$}\noLine
\UIC{$r \in \mathbb{R} \to \varphi_{33}$}
\AXC{$\stackrel{\Pi_{77}}{\vdots}$}\noLine
\UIC{$\lvert r \rvert \in \mathbb{R} \to 1 \in \mathbb{R} \to \lvert r \rvert < 1 \to 0 < 1 - \lvert r \rvert$}
\AXC{$(\lvert r \rvert \in \mathbb{R} \to 1 \in \mathbb{R} \to \lvert r \rvert < 1 \to 0 < 1 - \lvert r \rvert) \to \varphi_{224}$\lbl{A1}}
\BIC{$\varphi_{224}$}
\BIC{$1 \in \mathbb{R} \to (r \in \mathbb{R} \to \lvert r \rvert < 1 \to \lvert r \rvert < 1) \to r \in \mathbb{R} \to \varphi_{38}$}
\end{ptree}
\noindent$\Pi_{79}$:
\begin{ptree}
\AXC{$r \in \mathbb{R} \to \lvert r \rvert < 1 \to \lvert r \rvert < 1$\lbl{A1}}
\AXC{$1 \in \mathbb{R}$\lbl{(R1)}}
\AXC{$\stackrel{\Pi_{78}}{\vdots}$}\noLine
\UIC{$1 \in \mathbb{R} \to (r \in \mathbb{R} \to \lvert r \rvert < 1 \to \lvert r \rvert < 1) \to r \in \mathbb{R} \to \varphi_{38}$}
\BIC{$(r \in \mathbb{R} \to \lvert r \rvert < 1 \to \lvert r \rvert < 1) \to r \in \mathbb{R} \to \varphi_{38}$}
\BIC{$r \in \mathbb{R} \to \varphi_{38}$}
\end{ptree}
\noindent$\Pi_{80}$:
\begin{ptree}
\AXC{$\forall x\,\forall y\,(x < y \to x \neq y)$\lbl{[F26]}}
\AXC{$\forall x\,\forall y\,(x < y \to x \neq y) \to \forall y\,(0 < y \to 0 \neq y)$\lbl{A2}}
\BIC{$\forall y\,(0 < y \to 0 \neq y)$}
\end{ptree}
\noindent$\Pi_{81}$:
\begin{ptree}
\AXC{$\stackrel{\Pi_{79}}{\vdots}$}\noLine
\UIC{$r \in \mathbb{R} \to \varphi_{38}$}
\AXC{$\stackrel{\Pi_{80}}{\vdots}$}\noLine
\UIC{$\forall y\,(0 < y \to 0 \neq y)$}
\AXC{$\forall y\,(0 < y \to 0 \neq y) \to 0 < 1 - \lvert r \rvert \to 0 \neq 1 - \lvert r \rvert$\lbl{A2}}
\BIC{$0 < 1 - \lvert r \rvert \to 0 \neq 1 - \lvert r \rvert$}
\AXC{$\varphi_{217}$\lbl{A1}}
\BIC{$(r \in \mathbb{R} \to \varphi_{38}) \to r \in \mathbb{R} \to \varphi_{28}$}
\BIC{$r \in \mathbb{R} \to \varphi_{28}$}
\AXC{$\varphi_{214}$\lbl{A1}}
\BIC{$(1 - \lvert r \rvert = 0 \to 0 = 1 - \lvert r \rvert) \to r \in \mathbb{R} \to \varphi_{29}$}
\end{ptree}
\noindent$\Pi_{82}$:
\begin{ptree}
\AXC{$\stackrel{\Pi_{5}}{\vdots}$}\noLine
\UIC{$\forall x\,\forall y\,(x = y \to y = x)$}
\AXC{$\varphi_{31}$\lbl{A2}}
\BIC{$\forall y\,(1 - \lvert r \rvert = y \to y = 1 - \lvert r \rvert)$}
\AXC{$\varphi_{74}$\lbl{A2}}
\BIC{$1 - \lvert r \rvert = 0 \to 0 = 1 - \lvert r \rvert$}
\AXC{$\stackrel{\Pi_{81}}{\vdots}$}\noLine
\UIC{$(1 - \lvert r \rvert = 0 \to 0 = 1 - \lvert r \rvert) \to r \in \mathbb{R} \to \varphi_{29}$}
\BIC{$r \in \mathbb{R} \to \varphi_{29}$}
\end{ptree}
\noindent$\Pi_{83}$:
\begin{ptree}
\AXC{$\stackrel{\Pi_{8}}{\vdots}$}\noLine
\UIC{$r \in \mathbb{R} \to \varphi_{33}$}
\AXC{$1 \in \mathbb{R}$\lbl{(R1)}}
\AXC{$\forall x\,\forall y\,(x \in \mathbb{R} \to y \in \mathbb{R} \to x - y \in \mathbb{R})$\lbl{[F1]}}
\AXC{$\varphi_{152}$\lbl{A2}}
\BIC{$\forall y\,(1 \in \mathbb{R} \to y \in \mathbb{R} \to 1 - y \in \mathbb{R})$}
\AXC{$\varphi_{157}$\lbl{A2}}
\BIC{$1 \in \mathbb{R} \to \lvert r \rvert \in \mathbb{R} \to 1 - \lvert r \rvert \in \mathbb{R}$}
\AXC{$\varphi_{229}$\lbl{A1}}
\BIC{$1 \in \mathbb{R} \to (r \in \mathbb{R} \to \varphi_{33}) \to r \in \mathbb{R} \to \varphi_{47}$}
\BIC{$(r \in \mathbb{R} \to \varphi_{33}) \to r \in \mathbb{R} \to \varphi_{47}$}
\BIC{$r \in \mathbb{R} \to \varphi_{47}$}
\end{ptree}
\noindent$\Pi_{84}$:
\begin{ptree}
\AXC{$\stackrel{\Pi_{2}}{\vdots}$}\noLine
\UIC{$\varphi_{53}$}
\AXC{$\forall x\,\forall y\,(x \in \mathbb{R} \to y \in \mathbb{R} \to y \neq 0 \to x \cdot y = 0 \to x = 0)$\lbl{[E48]}}
\AXC{$\varphi_{169}$\lbl{A2}}
\BIC{$\forall y\,(L \in \mathbb{R} \to y \in \mathbb{R} \to y \neq 0 \to L \cdot y = 0 \to L = 0)$}
\AXC{$\varphi_{202}$\lbl{A2}}
\BIC{$L \in \mathbb{R} \to \varphi_{76}$}
\AXC{$\varphi_{261}$\lbl{A1}}
\BIC{$\varphi_{53} \to (r \in \mathbb{R} \to \varphi_{47}) \to (r \in \mathbb{R} \to \varphi_{29}) \to (r \in \mathbb{R} \to \varphi_{41}) \to \varphi_{37}$}
\BIC{$(r \in \mathbb{R} \to \varphi_{47}) \to (r \in \mathbb{R} \to \varphi_{29}) \to (r \in \mathbb{R} \to \varphi_{41}) \to \varphi_{37}$}
\end{ptree}
\noindent$\Pi_{85}$:
\begin{ptree}
\AXC{$\stackrel{\Pi_{82}}{\vdots}$}\noLine
\UIC{$r \in \mathbb{R} \to \varphi_{29}$}
\AXC{$\stackrel{\Pi_{83}}{\vdots}$}\noLine
\UIC{$r \in \mathbb{R} \to \varphi_{47}$}
\AXC{$\stackrel{\Pi_{84}}{\vdots}$}\noLine
\UIC{$(r \in \mathbb{R} \to \varphi_{47}) \to (r \in \mathbb{R} \to \varphi_{29}) \to (r \in \mathbb{R} \to \varphi_{41}) \to \varphi_{37}$}
\BIC{$(r \in \mathbb{R} \to \varphi_{29}) \to (r \in \mathbb{R} \to \varphi_{41}) \to \varphi_{37}$}
\BIC{$(r \in \mathbb{R} \to \varphi_{41}) \to \varphi_{37}$}
\end{ptree}
\noindent$\Pi_{86}$:
\begin{ptree}
\AXC{$\stackrel{\Pi_{1}}{\vdots}$}\noLine
\UIC{$\varphi_{36}$}
\AXC{$\varphi_{36} \to \varphi_{52}$\lbl{A2}}
\BIC{$\varphi_{52}$}
\AXC{$\varphi_{52} \to \varphi_{45}$\lbl{A2}}
\BIC{$\varphi_{45}$}
\AXC{$\stackrel{\Pi_{3}}{\vdots}$}\noLine
\UIC{$\varphi_{48}$}
\AXC{$\stackrel{\Pi_{76}}{\vdots}$}\noLine
\UIC{$\varphi_{48} \to \varphi_{45} \to r \in \mathbb{R} \to \varphi_{41}$}
\BIC{$\varphi_{45} \to r \in \mathbb{R} \to \varphi_{41}$}
\BIC{$r \in \mathbb{R} \to \varphi_{41}$}
\AXC{$\stackrel{\Pi_{85}}{\vdots}$}\noLine
\UIC{$(r \in \mathbb{R} \to \varphi_{41}) \to \varphi_{37}$}
\BIC{$\varphi_{37}$}
\AXC{$\varphi_{37} \to \varphi_{223}$\lbl{A1}}
\BIC{$\varphi_{223}$}
\end{ptree}
\noindent$\Pi_{87}$:
\begin{ptree}
\AXC{$z = w \to \varphi_{7}$\lbl{A6}}
\UIC{$\forall w\,(z = w \to \varphi_{7})$}
\UIC{$\forall z\,\forall w\,(z = w \to \varphi_{7})$}
\AXC{$\forall z\,\forall w\,(z = w \to \varphi_{7}) \to \varphi_{21}$\lbl{A2}}
\BIC{$\varphi_{21}$}
\AXC{$\varphi_{21} \to L = 0 \to \varphi_{8}$\lbl{A2}}
\BIC{$L = 0 \to \varphi_{8}$}
\AXC{$r \in \mathbb{R} \to \lvert r \rvert < 1 \to \varphi_{9}$\lbl{A1}}
\AXC{$\stackrel{\Pi_{86}}{\vdots}$}\noLine
\UIC{$\varphi_{223}$}
\BIC{$(L = 0 \to \varphi_{8}) \to r \in \mathbb{R} \to \lvert r \rvert < 1 \to \varphi_{8}$}
\BIC{$r \in \mathbb{R} \to \lvert r \rvert < 1 \to \varphi_{8}$}
\end{ptree}
\noindent$\Pi_{88}$:
\begin{ptree}
\AXC{$y = z \to x = y \to x = z$\lbl{A6}}
\AXC{$(y = z \to x = y \to x = z) \to x = y \to y = z \to x = z$\lbl{A1}}
\BIC{$x = y \to y = z \to x = z$}
\UIC{$\forall z\,(x = y \to y = z \to x = z)$}
\UIC{$\forall y\,\forall z\,(x = y \to y = z \to x = z)$}
\UIC{$\forall x\,\forall y\,\forall z\,(x = y \to y = z \to x = z)$}
\AXC{$\varphi_{121}$\lbl{A2}}
\BIC{$\forall y\,\forall z\,({{\lvert r \rvert}^{n}} = y \to y = z \to {{\lvert r \rvert}^{n}} = z)$}
\AXC{$\varphi_{208}$\lbl{A2}}
\BIC{$\forall z\,({{\lvert r \rvert}^{n}} = \lvert {{r}^{n}} \rvert \to \lvert {{r}^{n}} \rvert = z \to {{\lvert r \rvert}^{n}} = z)$}
\AXC{$\varphi_{227}$\lbl{A2}}
\BIC{${{\lvert r \rvert}^{n}} = \lvert {{r}^{n}} \rvert \to \varphi_{109}$}
\end{ptree}
\noindent$\Pi_{89}$:
\begin{ptree}
\AXC{$\varphi_{55}$\lbl{A1}}
\AXC{$\forall r \in \mathbb{R}\,(\varphi_{127})$\lbl{[S24]}}
\AXC{$\forall r \in \mathbb{R}\,(\varphi_{127}) \to r \in \mathbb{R} \to \varphi_{127}$\lbl{A2}}
\BIC{$r \in \mathbb{R} \to \varphi_{127}$}
\AXC{$(r \in \mathbb{R} \to \varphi_{127}) \to \varphi_{55} \to \varphi_{201}$\lbl{A1}}
\BIC{$\varphi_{55} \to \varphi_{201}$}
\BIC{$\varphi_{201}$}
\end{ptree}
\noindent$\Pi_{90}$:
\begin{ptree}
\AXC{$\stackrel{\Pi_{89}}{\vdots}$}\noLine
\UIC{$\varphi_{201}$}
\AXC{$\varphi_{201} \to r \in \mathbb{R} \to \lvert r \rvert < 1 \to \varphi_{1}$\lbl{A1}}
\BIC{$r \in \mathbb{R} \to \lvert r \rvert < 1 \to \varphi_{1}$}
\end{ptree}
\noindent$\Pi_{91}$:
\begin{ptree}
\AXC{$\stackrel{\Pi_{90}}{\vdots}$}\noLine
\UIC{$r \in \mathbb{R} \to \lvert r \rvert < 1 \to \varphi_{1}$}
\AXC{$\varphi_{92}$\lbl{[S11]}}
\AXC{$\varphi_{92} \to \mathrm{Seq}({(r)^{\bullet}}) \to \varphi_{95}$\lbl{A2}}
\BIC{$\mathrm{Seq}({(r)^{\bullet}}) \to \varphi_{95}$}
\AXC{$\varphi_{243}$\lbl{A1}}
\BIC{$(r \in \mathbb{R} \to \lvert r \rvert < 1 \to \varphi_{1}) \to \varphi_{186}$}
\BIC{$\varphi_{186}$}
\end{ptree}
\noindent$\Pi_{92}$:
\begin{ptree}
\AXC{$\varphi_{165}$\lbl{A2}}
\AXC{$\stackrel{\Pi_{91}}{\vdots}$}\noLine
\UIC{$\varphi_{186}$}
\AXC{$\varphi_{186} \to r \in \mathbb{R} \to \lvert r \rvert < 1 \to \varphi_{137}$\lbl{A1}}
\BIC{$r \in \mathbb{R} \to \lvert r \rvert < 1 \to \varphi_{137}$}
\AXC{$(r \in \mathbb{R} \to \lvert r \rvert < 1 \to \varphi_{137}) \to \varphi_{165} \to r \in \mathbb{R} \to \lvert r \rvert < 1 \to \varphi_{128}$\lbl{A1}}
\BIC{$\varphi_{165} \to r \in \mathbb{R} \to \lvert r \rvert < 1 \to \varphi_{128}$}
\BIC{$r \in \mathbb{R} \to \lvert r \rvert < 1 \to \varphi_{128}$}
\end{ptree}
\noindent$\Pi_{93}$:
\begin{ptree}
\AXC{$r \in \mathbb{R} \to \varphi_{50}$\lbl{A1}}
\AXC{$\stackrel{\Pi_{92}}{\vdots}$}\noLine
\UIC{$r \in \mathbb{R} \to \lvert r \rvert < 1 \to \varphi_{128}$}
\AXC{$\varphi_{246}$\lbl{A1}}
\BIC{$(r \in \mathbb{R} \to \varphi_{50}) \to r \in \mathbb{R} \to \lvert r \rvert < 1 \to \varphi_{86}$}
\BIC{$r \in \mathbb{R} \to \lvert r \rvert < 1 \to \varphi_{86}$}
\AXC{$(r \in \mathbb{R} \to \lvert r \rvert < 1 \to \varphi_{86}) \to \varphi_{107} \to r \in \mathbb{R} \to \lvert r \rvert < 1 \to \varphi_{87}$\lbl{A1}}
\BIC{$\varphi_{107} \to r \in \mathbb{R} \to \lvert r \rvert < 1 \to \varphi_{87}$}
\end{ptree}
\noindent$\Pi_{94}$:
\begin{ptree}
\AXC{$\stackrel{\Pi_{5}}{\vdots}$}\noLine
\UIC{$\forall x\,\forall y\,(x = y \to y = x)$}
\AXC{$\forall x\,\forall y\,(x = y \to y = x) \to \varphi_{19}$\lbl{A2}}
\BIC{$\varphi_{19}$}
\AXC{$\varphi_{19} \to \varphi_{107}$\lbl{A2}}
\BIC{$\varphi_{107}$}
\AXC{$\stackrel{\Pi_{93}}{\vdots}$}\noLine
\UIC{$\varphi_{107} \to r \in \mathbb{R} \to \lvert r \rvert < 1 \to \varphi_{87}$}
\BIC{$r \in \mathbb{R} \to \lvert r \rvert < 1 \to \varphi_{87}$}
\end{ptree}
\noindent$\Pi_{95}$:
\begin{ptree}
\AXC{$r \in \mathbb{R} \to \varphi_{51}$\lbl{A1}}
\AXC{$\forall r \in \mathbb{R}\,\forall n \in \mathbb{N}\,(\lvert {{r}^{n}} \rvert = {{\lvert r \rvert}^{n}})$\lbl{[P1]}}
\AXC{$\varphi_{204}$\lbl{A2}}
\BIC{$r \in \mathbb{R} \to \forall n \in \mathbb{N}\,(\lvert {{r}^{n}} \rvert = {{\lvert r \rvert}^{n}})$}
\AXC{$\varphi_{241}$\lbl{A1}}
\BIC{$(r \in \mathbb{R} \to \varphi_{51}) \to r \in \mathbb{R} \to \lvert r \rvert < 1 \to \varphi_{133}$}
\BIC{$r \in \mathbb{R} \to \lvert r \rvert < 1 \to \varphi_{133}$}
\end{ptree}
\noindent$\Pi_{96}$:
\begin{ptree}
\AXC{$\varphi_{162}$\lbl{A2}}
\AXC{$\stackrel{\Pi_{95}}{\vdots}$}\noLine
\UIC{$r \in \mathbb{R} \to \lvert r \rvert < 1 \to \varphi_{133}$}
\AXC{$(r \in \mathbb{R} \to \lvert r \rvert < 1 \to \varphi_{133}) \to \varphi_{162} \to r \in \mathbb{R} \to \lvert r \rvert < 1 \to \varphi_{123}$\lbl{A1}}
\BIC{$\varphi_{162} \to r \in \mathbb{R} \to \lvert r \rvert < 1 \to \varphi_{123}$}
\BIC{$r \in \mathbb{R} \to \lvert r \rvert < 1 \to \varphi_{123}$}
\AXC{$\varphi_{245}$\lbl{A1}}
\BIC{$(r \in \mathbb{R} \to \varphi_{50}) \to r \in \mathbb{R} \to \lvert r \rvert < 1 \to \varphi_{81}$}
\end{ptree}
\noindent$\Pi_{97}$:
\begin{ptree}
\AXC{$r \in \mathbb{R} \to \varphi_{50}$\lbl{A1}}
\AXC{$\stackrel{\Pi_{96}}{\vdots}$}\noLine
\UIC{$(r \in \mathbb{R} \to \varphi_{50}) \to r \in \mathbb{R} \to \lvert r \rvert < 1 \to \varphi_{81}$}
\BIC{$r \in \mathbb{R} \to \lvert r \rvert < 1 \to \varphi_{81}$}
\AXC{$(r \in \mathbb{R} \to \lvert r \rvert < 1 \to \varphi_{81}) \to \varphi_{219}$\lbl{A1}}
\BIC{$\varphi_{219}$}
\end{ptree}
\noindent$\Pi_{98}$:
\begin{ptree}
\AXC{$\stackrel{\Pi_{5}}{\vdots}$}\noLine
\UIC{$\forall x\,\forall y\,(x = y \to y = x)$}
\AXC{$\varphi_{70}$\lbl{A2}}
\BIC{$\forall y\,(\lvert {{r}^{n}} \rvert = y \to y = \lvert {{r}^{n}} \rvert)$}
\AXC{$\varphi_{190}$\lbl{A2}}
\BIC{$\lvert {{r}^{n}} \rvert = {{\lvert r \rvert}^{n}} \to {{\lvert r \rvert}^{n}} = \lvert {{r}^{n}} \rvert$}
\AXC{$\stackrel{\Pi_{97}}{\vdots}$}\noLine
\UIC{$\varphi_{219}$}
\BIC{$r \in \mathbb{R} \to \lvert r \rvert < 1 \to \varphi_{82}$}
\AXC{$(r \in \mathbb{R} \to \lvert r \rvert < 1 \to \varphi_{82}) \to \varphi_{255}$\lbl{A1}}
\BIC{$\varphi_{255}$}
\end{ptree}
\noindent$\Pi_{99}$:
\begin{ptree}
\AXC{$\stackrel{\Pi_{88}}{\vdots}$}\noLine
\UIC{${{\lvert r \rvert}^{n}} = \lvert {{r}^{n}} \rvert \to \varphi_{109}$}
\AXC{$\stackrel{\Pi_{94}}{\vdots}$}\noLine
\UIC{$r \in \mathbb{R} \to \lvert r \rvert < 1 \to \varphi_{87}$}
\AXC{$\stackrel{\Pi_{98}}{\vdots}$}\noLine
\UIC{$\varphi_{255}$}
\BIC{$({{\lvert r \rvert}^{n}} = \lvert {{r}^{n}} \rvert \to \varphi_{109}) \to r \in \mathbb{R} \to \lvert r \rvert < 1 \to \varphi_{88}$}
\BIC{$r \in \mathbb{R} \to \lvert r \rvert < 1 \to \varphi_{88}$}
\AXC{$(r \in \mathbb{R} \to \lvert r \rvert < 1 \to \varphi_{88}) \to \varphi_{154}$\lbl{A1}}
\BIC{$\varphi_{154}$}
\UIC{$\forall n\,(\varphi_{154})$}
\end{ptree}
\noindent$\Pi_{100}$:
\begin{ptree}
\AXC{$\varphi_{153}$\lbl{A4}}
\AXC{$\varphi_{231}$\lbl{A3}}
\AXC{$\stackrel{\Pi_{99}}{\vdots}$}\noLine
\UIC{$\forall n\,(\varphi_{154})$}
\AXC{$\forall n\,(\varphi_{154}) \to \varphi_{231} \to \varphi_{153} \to r \in \mathbb{R} \to \lvert r \rvert < 1 \to \varphi_{100}$\lbl{A1}}
\BIC{$\varphi_{231} \to \varphi_{153} \to r \in \mathbb{R} \to \lvert r \rvert < 1 \to \varphi_{100}$}
\BIC{$\varphi_{153} \to r \in \mathbb{R} \to \lvert r \rvert < 1 \to \varphi_{100}$}
\BIC{$r \in \mathbb{R} \to \lvert r \rvert < 1 \to \varphi_{100}$}
\end{ptree}
\noindent$\Pi_{101}$:
\begin{ptree}
\AXC{$\varphi_{112}$\lbl{[C13]}}
\AXC{$\varphi_{112} \to \varphi_{149}$\lbl{A2}}
\BIC{$\varphi_{149}$}
\AXC{$\varphi_{149} \to \varphi_{189}$\lbl{A2}}
\BIC{$\varphi_{189}$}
\AXC{$\varphi_{189} \to \varphi_{176}$\lbl{A2}}
\BIC{$\varphi_{176}$}
\AXC{$\varphi_{176} \to (r \in \mathbb{R} \to \lvert r \rvert < 1 \to \varphi_{5}) \to \varphi_{238}$\lbl{A1}}
\BIC{$(r \in \mathbb{R} \to \lvert r \rvert < 1 \to \varphi_{5}) \to \varphi_{238}$}
\end{ptree}
\noindent$\Pi_{102}$:
\begin{ptree}
\AXC{$\stackrel{\Pi_{100}}{\vdots}$}\noLine
\UIC{$r \in \mathbb{R} \to \lvert r \rvert < 1 \to \varphi_{100}$}
\AXC{$\stackrel{\Pi_{91}}{\vdots}$}\noLine
\UIC{$\varphi_{186}$}
\AXC{$\varphi_{186} \to r \in \mathbb{R} \to \lvert r \rvert < 1 \to \varphi_{5}$\lbl{A1}}
\BIC{$r \in \mathbb{R} \to \lvert r \rvert < 1 \to \varphi_{5}$}
\AXC{$\stackrel{\Pi_{101}}{\vdots}$}\noLine
\UIC{$(r \in \mathbb{R} \to \lvert r \rvert < 1 \to \varphi_{5}) \to \varphi_{238}$}
\BIC{$\varphi_{238}$}
\BIC{$(r \in \mathbb{R} \to \lvert r \rvert < 1 \to \varphi_{8}) \to r \in \mathbb{R} \to \lvert r \rvert < 1 \to \varphi_{11}$}
\end{ptree}
\noindent$\Pi_{103}$:
\begin{ptree}
\AXC{$\stackrel{\Pi_{87}}{\vdots}$}\noLine
\UIC{$r \in \mathbb{R} \to \lvert r \rvert < 1 \to \varphi_{8}$}
\AXC{$\stackrel{\Pi_{102}}{\vdots}$}\noLine
\UIC{$(r \in \mathbb{R} \to \lvert r \rvert < 1 \to \varphi_{8}) \to r \in \mathbb{R} \to \lvert r \rvert < 1 \to \varphi_{11}$}
\BIC{$r \in \mathbb{R} \to \lvert r \rvert < 1 \to \varphi_{11}$}
\end{ptree}
\noindent$\Pi_{104}$:
\begin{ptree}
\AXC{$\stackrel{\Pi_{90}}{\vdots}$}\noLine
\UIC{$r \in \mathbb{R} \to \lvert r \rvert < 1 \to \varphi_{1}$}
\AXC{$\varphi_{6}$\lbl{定理2.15(2)}}
\AXC{$\varphi_{6} \to \mathrm{Seq}({(r)^{\bullet}}) \to \varphi_{3}$\lbl{A2}}
\BIC{$\mathrm{Seq}({(r)^{\bullet}}) \to \varphi_{3}$}
\AXC{$\varphi_{226}$\lbl{A1}}
\BIC{$(r \in \mathbb{R} \to \lvert r \rvert < 1 \to \varphi_{1}) \to \varphi_{130}$}
\BIC{$\varphi_{130}$}
\end{ptree}
\noindent$\Pi_{105}$:
\begin{ptree}
\AXC{$\stackrel{\Pi_{103}}{\vdots}$}\noLine
\UIC{$r \in \mathbb{R} \to \lvert r \rvert < 1 \to \varphi_{11}$}
\AXC{$\stackrel{\Pi_{104}}{\vdots}$}\noLine
\UIC{$\varphi_{130}$}
\AXC{$\varphi_{130} \to (r \in \mathbb{R} \to \lvert r \rvert < 1 \to \varphi_{11}) \to r \in \mathbb{R} \to \lvert r \rvert < 1 \to \varphi_{2}$\lbl{A1}}
\BIC{$(r \in \mathbb{R} \to \lvert r \rvert < 1 \to \varphi_{11}) \to r \in \mathbb{R} \to \lvert r \rvert < 1 \to \varphi_{2}$}
\BIC{$r \in \mathbb{R} \to \lvert r \rvert < 1 \to \varphi_{2}$}
\end{ptree}
\noindent$\Pi_{106}$:
\begin{ptree}
\AXC{$\forall x \in \mathbb{R}\,(0 \leq \lvert x \rvert)$\lbl{[A6]}}
\AXC{$\forall x \in \mathbb{R}\,(0 \leq \lvert x \rvert) \to r \in \mathbb{R} \to 0 \leq \lvert r \rvert$\lbl{A2}}
\BIC{$r \in \mathbb{R} \to 0 \leq \lvert r \rvert$}
\AXC{$(r \in \mathbb{R} \to 0 \leq \lvert r \rvert) \to \varphi_{108}$\lbl{A1}}
\BIC{$\varphi_{108}$}
\end{ptree}
\noindent$\Pi_{107}$:
\begin{ptree}
\AXC{$r \in \mathbb{R} \to \lvert r \rvert < 1 \to r \in \mathbb{R}$\lbl{A1}}
\AXC{$\stackrel{\Pi_{106}}{\vdots}$}\noLine
\UIC{$\varphi_{108}$}
\BIC{$r \in \mathbb{R} \to \lvert r \rvert < 1 \to 0 \leq \lvert r \rvert$}
\end{ptree}
\noindent$\Pi_{108}$:
\begin{ptree}
\AXC{$\forall x \in \mathbb{R}\,(\lvert x \rvert \in \mathbb{R})$\lbl{[A5]}}
\AXC{$\forall x \in \mathbb{R}\,(\lvert x \rvert \in \mathbb{R}) \to r \in \mathbb{R} \to \lvert r \rvert \in \mathbb{R}$\lbl{A2}}
\BIC{$r \in \mathbb{R} \to \lvert r \rvert \in \mathbb{R}$}
\AXC{$(r \in \mathbb{R} \to \lvert r \rvert \in \mathbb{R}) \to \varphi_{129}$\lbl{A1}}
\BIC{$\varphi_{129}$}
\end{ptree}
\noindent$\Pi_{109}$:
\begin{ptree}
\AXC{$r \in \mathbb{R} \to \lvert r \rvert < 1 \to r \in \mathbb{R}$\lbl{A1}}
\AXC{$\stackrel{\Pi_{108}}{\vdots}$}\noLine
\UIC{$\varphi_{129}$}
\BIC{$r \in \mathbb{R} \to \lvert r \rvert < 1 \to \lvert r \rvert \in \mathbb{R}$}
\end{ptree}
\noindent$\Pi_{110}$:
\begin{ptree}
\AXC{$\stackrel{\Pi_{109}}{\vdots}$}\noLine
\UIC{$r \in \mathbb{R} \to \lvert r \rvert < 1 \to \lvert r \rvert \in \mathbb{R}$}
\AXC{$\forall r \in \mathbb{R}\,(0 \leq r \to \forall n \in \mathbb{N}\,(0 \leq {{r}^{n}}))$\lbl{[P2]}}
\AXC{$\varphi_{181}$\lbl{A2}}
\BIC{$\lvert r \rvert \in \mathbb{R} \to 0 \leq \lvert r \rvert \to \forall n \in \mathbb{N}\,(0 \leq {{\lvert r \rvert}^{n}})$}
\AXC{$\varphi_{232}$\lbl{A1}}
\BIC{$(r \in \mathbb{R} \to \lvert r \rvert < 1 \to \lvert r \rvert \in \mathbb{R}) \to \varphi_{161}$}
\BIC{$\varphi_{161}$}
\end{ptree}
\noindent$\Pi_{111}$:
\begin{ptree}
\AXC{$\stackrel{\Pi_{109}}{\vdots}$}\noLine
\UIC{$r \in \mathbb{R} \to \lvert r \rvert < 1 \to \lvert r \rvert \in \mathbb{R}$}
\AXC{$\varphi_{14}$\lbl{[P7]}}
\AXC{$\varphi_{14} \to \lvert r \rvert \in \mathbb{R} \to \varphi_{12}$\lbl{A2}}
\BIC{$\lvert r \rvert \in \mathbb{R} \to \varphi_{12}$}
\AXC{$(\lvert r \rvert \in \mathbb{R} \to \varphi_{12}) \to \varphi_{218}$\lbl{A1}}
\BIC{$\varphi_{218}$}
\BIC{$(r \in \mathbb{R} \to \lvert r \rvert < 1 \to 0 \leq \lvert r \rvert) \to \varphi_{134}$}
\end{ptree}
\noindent$\Pi_{112}$:
\begin{ptree}
\AXC{$r \in \mathbb{R} \to \lvert r \rvert < 1 \to \lvert r \rvert < 1$\lbl{A1}}
\AXC{$\varphi_{111}$\lbl{A1}}
\BIC{$r \in \mathbb{R} \to \lvert r \rvert < 1 \to \lvert r \rvert \leq 1$}
\AXC{$\stackrel{\Pi_{107}}{\vdots}$}\noLine
\UIC{$r \in \mathbb{R} \to \lvert r \rvert < 1 \to 0 \leq \lvert r \rvert$}
\AXC{$\stackrel{\Pi_{111}}{\vdots}$}\noLine
\UIC{$(r \in \mathbb{R} \to \lvert r \rvert < 1 \to 0 \leq \lvert r \rvert) \to \varphi_{134}$}
\BIC{$\varphi_{134}$}
\BIC{$r \in \mathbb{R} \to \lvert r \rvert < 1 \to \mathrm{MonoDec}({(\lvert r \rvert)^{\bullet}})$}
\end{ptree}
\noindent$\Pi_{113}$:
\begin{ptree}
\AXC{$\varphi_{104}$\lbl{定理5.3(2)}}
\AXC{$\varphi_{104} \to \varphi_{158}$\lbl{A2}}
\BIC{$\varphi_{158}$}
\AXC{$\varphi_{158} \to \mathrm{MonoDec}({(\lvert r \rvert)^{\bullet}}) \to 0 \in \mathbb{R} \to \varphi_{67}$\lbl{A2}}
\BIC{$\mathrm{MonoDec}({(\lvert r \rvert)^{\bullet}}) \to 0 \in \mathbb{R} \to \varphi_{67}$}
\end{ptree}
\noindent$\Pi_{114}$:
\begin{ptree}
\AXC{$\stackrel{\Pi_{113}}{\vdots}$}\noLine
\UIC{$\mathrm{MonoDec}({(\lvert r \rvert)^{\bullet}}) \to 0 \in \mathbb{R} \to \varphi_{67}$}
\AXC{$(\mathrm{MonoDec}({(\lvert r \rvert)^{\bullet}}) \to 0 \in \mathbb{R} \to \varphi_{67}) \to \varphi_{221}$\lbl{A1}}
\BIC{$\varphi_{221}$}
\end{ptree}
\noindent$\Pi_{115}$:
\begin{ptree}
\AXC{$0 \in \mathbb{R}$\lbl{(R1)}}
\AXC{$\stackrel{\Pi_{112}}{\vdots}$}\noLine
\UIC{$r \in \mathbb{R} \to \lvert r \rvert < 1 \to \mathrm{MonoDec}({(\lvert r \rvert)^{\bullet}})$}
\AXC{$\stackrel{\Pi_{114}}{\vdots}$}\noLine
\UIC{$\varphi_{221}$}
\BIC{$\varphi_{195}$}
\BIC{$(r \in \mathbb{R} \to \lvert r \rvert < 1 \to \forall n \in \mathbb{N}\,(0 \leq {{\lvert r \rvert}^{n}})) \to \varphi_{32}$}
\end{ptree}
\noindent$\Pi_{116}$:
\begin{ptree}
\AXC{$\stackrel{\Pi_{107}}{\vdots}$}\noLine
\UIC{$r \in \mathbb{R} \to \lvert r \rvert < 1 \to 0 \leq \lvert r \rvert$}
\AXC{$\stackrel{\Pi_{110}}{\vdots}$}\noLine
\UIC{$\varphi_{161}$}
\BIC{$r \in \mathbb{R} \to \lvert r \rvert < 1 \to \forall n \in \mathbb{N}\,(0 \leq {{\lvert r \rvert}^{n}})$}
\AXC{$\stackrel{\Pi_{115}}{\vdots}$}\noLine
\UIC{$(r \in \mathbb{R} \to \lvert r \rvert < 1 \to \forall n \in \mathbb{N}\,(0 \leq {{\lvert r \rvert}^{n}})) \to \varphi_{32}$}
\BIC{$\varphi_{32}$}
\end{ptree}
\noindent$\Pi_{117}$:
\begin{ptree}
\AXC{$\varphi_{206}$\lbl{A3}}
\AXC{$\stackrel{\Pi_{105}}{\vdots}$}\noLine
\UIC{$r \in \mathbb{R} \to \lvert r \rvert < 1 \to \varphi_{2}$}
\AXC{$(r \in \mathbb{R} \to \lvert r \rvert < 1 \to \varphi_{2}) \to \varphi_{73}$\lbl{A1}}
\BIC{$\varphi_{73}$}
\UIC{$\forall L\,(\varphi_{73})$}
\AXC{$\stackrel{\Pi_{116}}{\vdots}$}\noLine
\UIC{$\varphi_{32}$}
\AXC{$\varphi_{32} \to \forall L\,(\varphi_{73}) \to \varphi_{206} \to \varphi_{258}$\lbl{A1}}
\BIC{$\forall L\,(\varphi_{73}) \to \varphi_{206} \to \varphi_{258}$}
\BIC{$\varphi_{206} \to \varphi_{258}$}
\BIC{$\varphi_{258}$}
\end{ptree}
\noindent 主木:
\begin{ptree}
\AXC{$\varphi_{35}$\lbl{A4}}
\AXC{$\forall \alpha\,(\varphi_{17}) \to \varphi_{40}$\lbl{A3}}
\AXC{$\varphi_{17}$\lbl{A2}}
\UIC{$\forall \alpha\,(\varphi_{17})$}
\AXC{$r \in \mathbb{R} \land \lvert r \rvert < 1 \land \neg \mathrm{Conv}({(r)^{\bullet}},0) \to \varphi_{15}$\lbl{A4}}
\AXC{$\stackrel{\Pi_{117}}{\vdots}$}\noLine
\UIC{$\varphi_{258}$}
\BIC{$\forall \alpha\,(\varphi_{17}) \to \varphi_{230}$}
\BIC{$\varphi_{230}$}
\BIC{$\varphi_{35} \to r \in \mathbb{R} \to \lvert r \rvert < 1 \to \mathrm{Conv}({(r)^{\bullet}},0)$}
\BIC{$r \in \mathbb{R} \to \lvert r \rvert < 1 \to \mathrm{Conv}({(r)^{\bullet}},0)$}
\UIC{$\forall r \in \mathbb{R}\,(\lvert r \rvert < 1 \to \mathrm{Conv}({(r)^{\bullet}},0))$}
\end{ptree}
\begin{abbrevlist}
\item $\varphi_{1} :\equiv \mathrm{Conv}({(\lvert r \rvert)^{\bullet}},L) \to \mathrm{Seq}({(r)^{\bullet}})$
\item $\varphi_{2} :\equiv \mathrm{Conv}({(\lvert r \rvert)^{\bullet}},L) \to \mathrm{Conv}({(r)^{\bullet}},0)$
\item $\varphi_{3} :\equiv \mathrm{Conv}({(r)^{\bullet}},0) \leftrightarrow \mathrm{Conv}(\lvert {(r)^{\bullet}} \rvert,0)$
\item $\varphi_{4} :\equiv \mathrm{Conv}({(\lvert r \rvert)^{\bullet}},L) \to \mathrm{Seq}({(\lvert r \rvert)^{\bullet}})$
\item $\varphi_{5} :\equiv \mathrm{Conv}({(\lvert r \rvert)^{\bullet}},L) \to \mathrm{Seq}(\lvert {(r)^{\bullet}} \rvert)$
\item $\varphi_{6} :\equiv \forall a\,(\mathrm{Seq}(a) \to (\mathrm{Conv}(a,0) \leftrightarrow \mathrm{Conv}(\lvert a \rvert,0)))$
\item $\varphi_{7} :\equiv \mathrm{Conv}({(\lvert r \rvert)^{\bullet}},z) \to \mathrm{Conv}({(\lvert r \rvert)^{\bullet}},w)$
\item $\varphi_{8} :\equiv \mathrm{Conv}({(\lvert r \rvert)^{\bullet}},L) \to \mathrm{Conv}({(\lvert r \rvert)^{\bullet}},0)$
\item $\varphi_{9} :\equiv \mathrm{Conv}({(\lvert r \rvert)^{\bullet}},L) \to \mathrm{Conv}({(\lvert r \rvert)^{\bullet}},L)$
\item $\varphi_{10} :\equiv \mathrm{Conv}({(\lvert r \rvert)^{\bullet}},0) \to \mathrm{Conv}(\lvert {(r)^{\bullet}} \rvert,0)$
\item $\varphi_{11} :\equiv \mathrm{Conv}({(\lvert r \rvert)^{\bullet}},L) \to \mathrm{Conv}(\lvert {(r)^{\bullet}} \rvert,0)$
\item $\varphi_{12} :\equiv 0 \leq \lvert r \rvert \to \lvert r \rvert \leq 1 \to \mathrm{MonoDec}({(\lvert r \rvert)^{\bullet}})$
\item $\varphi_{13} :\equiv \forall a\,\forall \alpha\,\forall \beta\,(\mathrm{Conv}(a,\alpha) \to \mathrm{Conv}(a,\beta) \to \alpha = \beta)$
\item $\varphi_{14} :\equiv \forall r \in \mathbb{R}\,(0 \leq r \to r \leq 1 \to \mathrm{MonoDec}({(r)^{\bullet}}))$
\item $\varphi_{15} :\equiv \forall L\,(r \in \mathbb{R} \land \lvert r \rvert < 1 \land \neg \mathrm{Conv}({(r)^{\bullet}},0))$
\item $\varphi_{16} :\equiv \mathrm{Conv}({(\lvert r \rvert)^{\bullet}},L) \to L \cdot (1 - \lvert r \rvert) = L - L$
\item $\varphi_{17} :\equiv \forall L\,(\neg \mathrm{Conv}({(\lvert r \rvert)^{\bullet}},L)) \to \neg \mathrm{Conv}({(\lvert r \rvert)^{\bullet}},\alpha)$
\item $\varphi_{18} :\equiv \lvert r \rvert < 1 \to \mathrm{Conv}({(\lvert r \rvert)^{\bullet}},L) \to \mathrm{NSeq}(\mathsf{s})$
\item $\varphi_{19} :\equiv \forall y\,(\lvert {(r)^{\bullet}} \rvert(n) = y \to y = \lvert {(r)^{\bullet}} \rvert(n))$
\item $\varphi_{20} :\equiv \mathrm{Conv}({(\lvert r \rvert)^{\bullet}},L) \to \mathrm{Seq}({(\lvert r \rvert)^{\bullet}} \circ \mathsf{s})$
\item $\varphi_{21} :\equiv \forall w\,(L = w \to \mathrm{Conv}({(\lvert r \rvert)^{\bullet}},L) \to \mathrm{Conv}({(\lvert r \rvert)^{\bullet}},w))$
\item $\varphi_{22} :\equiv \forall n\,(r \in \mathbb{R} \land \lvert r \rvert < 1 \land \mathrm{Conv}({(\lvert r \rvert)^{\bullet}},L))$
\item $\varphi_{23} :\equiv \mathrm{Conv}({(\lvert r \rvert)^{\bullet}},L) \to n \in \mathbb{N} \to \mathrm{Seq}({(\lvert r \rvert)^{\bullet}})$
\item $\varphi_{24} :\equiv \lvert r \rvert < 1 \to \mathrm{Conv}({(\lvert r \rvert)^{\bullet}},L) \to L = \lvert r \rvert \cdot L$
\item $\varphi_{25} :\equiv \lvert r \rvert < 1 \to \mathrm{Conv}({(\lvert r \rvert)^{\bullet}},L) \to \lvert r \rvert \cdot L = L$
\item $\varphi_{26} :\equiv \mathrm{Conv}({(\lvert r \rvert)^{\bullet}},L) \to n \in \mathbb{N} \to \mathsf{s}(n) = S(n)$
\item $\varphi_{27} :\equiv \mathrm{Conv}({(\lvert r \rvert)^{\bullet}},L) \to n \in \mathbb{N} \to S(n) = \mathsf{s}(n)$
\item $\varphi_{28} :\equiv \lvert r \rvert < 1 \to \mathrm{Conv}({(\lvert r \rvert)^{\bullet}},L) \to 0 \neq 1 - \lvert r \rvert$
\item $\varphi_{29} :\equiv \lvert r \rvert < 1 \to \mathrm{Conv}({(\lvert r \rvert)^{\bullet}},L) \to 1 - \lvert r \rvert \neq 0$
\item $\varphi_{30} :\equiv \forall x\,\forall y\,(x = y \to y = x) \to \forall y\,(\mathsf{s}(n) = y \to y = \mathsf{s}(n))$
\item $\varphi_{31} :\equiv \forall x\,\forall y\,(x = y \to y = x) \to \forall y\,(1 - \lvert r \rvert = y \to y = 1 - \lvert r \rvert)$
\item $\varphi_{32} :\equiv r \in \mathbb{R} \to \lvert r \rvert < 1 \to \exists \alpha\,(\mathrm{Conv}({(\lvert r \rvert)^{\bullet}},\alpha))$
\item $\varphi_{33} :\equiv \lvert r \rvert < 1 \to \mathrm{Conv}({(\lvert r \rvert)^{\bullet}},L) \to \lvert r \rvert \in \mathbb{R}$
\item $\varphi_{34} :\equiv \mathrm{Conv}({(\lvert r \rvert)^{\bullet}},L) \to \mathrm{Conv}({(\lvert r \rvert)^{\bullet}} \circ \mathsf{s},L)$
\item $\varphi_{35} :\equiv \forall L\,(\neg \mathrm{Conv}({(\lvert r \rvert)^{\bullet}},L)) \to \forall \alpha\,\forall L\,(\neg \mathrm{Conv}({(\lvert r \rvert)^{\bullet}},L))$
\item $\varphi_{36} :\equiv \forall y\,\forall z\,(L \cdot (1 - \lvert r \rvert) = y \to y = z \to L \cdot (1 - \lvert r \rvert) = z)$
\item $\varphi_{37} :\equiv r \in \mathbb{R} \to \lvert r \rvert < 1 \to \mathrm{Conv}({(\lvert r \rvert)^{\bullet}},L) \to L = 0$
\item $\varphi_{38} :\equiv \lvert r \rvert < 1 \to \mathrm{Conv}({(\lvert r \rvert)^{\bullet}},L) \to 0 < 1 - \lvert r \rvert$
\item $\varphi_{39} :\equiv \forall a\,\forall s\,\forall \alpha\,(\mathrm{Conv}(a,\alpha) \to \mathrm{SubIdx}(s) \to \mathrm{Conv}(a \circ s,\alpha))$
\item $\varphi_{40} :\equiv \forall \alpha\,\forall L\,(\neg \mathrm{Conv}({(\lvert r \rvert)^{\bullet}},L)) \to \forall \alpha\,(\neg \mathrm{Conv}({(\lvert r \rvert)^{\bullet}},\alpha))$
\item $\varphi_{41} :\equiv \lvert r \rvert < 1 \to \mathrm{Conv}({(\lvert r \rvert)^{\bullet}},L) \to L \cdot (1 - \lvert r \rvert) = 0$
\item $\varphi_{42} :\equiv \mathrm{Conv}({(\lvert r \rvert)^{\bullet}},L) \to L - L \cdot \lvert r \rvert = L - \lvert r \rvert \cdot L$
\item $\varphi_{43} :\equiv \mathrm{Conv}({(\lvert r \rvert)^{\bullet}},L) \to L \cdot (1 - \lvert r \rvert) = L - L \cdot \lvert r \rvert$
\item $\varphi_{44} :\equiv \mathrm{Conv}({(\lvert r \rvert)^{\bullet}},L) \to L \cdot (1 - \lvert r \rvert) = L - \lvert r \rvert \cdot L$
\item $\varphi_{45} :\equiv L \cdot (1 - \lvert r \rvert) = L - L \to L - L = 0 \to L \cdot (1 - \lvert r \rvert) = 0$
\item $\varphi_{46} :\equiv \mathrm{Conv}({(\lvert r \rvert)^{\bullet}} \circ \mathsf{s},\lvert r \rvert \cdot L) \to L = \lvert r \rvert \cdot L$
\item $\varphi_{47} :\equiv \lvert r \rvert < 1 \to \mathrm{Conv}({(\lvert r \rvert)^{\bullet}},L) \to 1 - \lvert r \rvert \in \mathbb{R}$
\item $\varphi_{48} :\equiv r \in \mathbb{R} \to \lvert r \rvert < 1 \to \mathrm{Conv}({(\lvert r \rvert)^{\bullet}},L) \to L - L = 0$
\item $\varphi_{49} :\equiv r \in \mathbb{R} \to \lvert r \rvert < 1 \to \mathrm{Conv}({(\lvert r \rvert)^{\bullet}},L) \to L \cdot 1 = L$
\item $\varphi_{50} :\equiv \lvert r \rvert < 1 \to \mathrm{Conv}({(\lvert r \rvert)^{\bullet}},L) \to n \in \mathbb{N} \to n \in \mathbb{N}$
\item $\varphi_{51} :\equiv \lvert r \rvert < 1 \to \mathrm{Conv}({(\lvert r \rvert)^{\bullet}},L) \to n \in \mathbb{N} \to r \in \mathbb{R}$
\item $\varphi_{52} :\equiv \forall z\,(L \cdot (1 - \lvert r \rvert) = L - L \to L - L = z \to L \cdot (1 - \lvert r \rvert) = z)$
\item $\varphi_{53} :\equiv r \in \mathbb{R} \to \lvert r \rvert < 1 \to \mathrm{Conv}({(\lvert r \rvert)^{\bullet}},L) \to L \in \mathbb{R}$
\item $\varphi_{54} :\equiv \forall w\,(z = w \to L - \lvert r \rvert \cdot L = L - z \to L - \lvert r \rvert \cdot L = L - w)$
\item $\varphi_{55} :\equiv r \in \mathbb{R} \to \lvert r \rvert < 1 \to \mathrm{Conv}({(\lvert r \rvert)^{\bullet}},L) \to r \in \mathbb{R}$
\item $\varphi_{56} :\equiv \forall w\,(z = w \to L - L \cdot \lvert r \rvert = L - z \to L - L \cdot \lvert r \rvert = L - w)$
\item $\varphi_{57} :\equiv \lvert r \rvert < 1 \to \mathrm{Conv}({(\lvert r \rvert)^{\bullet}},L) \to L - \lvert r \rvert \cdot L = L - L$
\item $\varphi_{58} :\equiv \lvert r \rvert < 1 \to \mathrm{Conv}({(\lvert r \rvert)^{\bullet}},L) \to L \cdot \lvert r \rvert = \lvert r \rvert \cdot L$
\item $\varphi_{59} :\equiv \forall c\,\forall a\,\forall \alpha\,(c \in \mathbb{R} \to \mathrm{Conv}(a,\alpha) \to \mathrm{Conv}(c \odot a,c \cdot \alpha))$
\item $\varphi_{60} :\equiv \forall n\,(\mathrm{Seq}({(\lvert r \rvert)^{\bullet}}) \to n \in \mathbb{N} \to {{\lvert r \rvert}^{n}} \in \mathbb{R})$
\item $\varphi_{61} :\equiv \lvert r \rvert \odot {(\lvert r \rvert)^{\bullet}}(n) = {(\lvert r \rvert)^{\bullet}} \circ \mathsf{s}(n)$
\item $\varphi_{62} :\equiv \mathrm{Conv}({(\lvert r \rvert)^{\bullet}},L) \to n \in \mathbb{N} \to {{\lvert r \rvert}^{n}} \in \mathbb{R}$
\item $\varphi_{63} :\equiv \mathrm{Conv}({(\lvert r \rvert)^{\bullet}},L) \to L \cdot 1 - L \cdot \lvert r \rvert = L - L \cdot \lvert r \rvert$
\item $\varphi_{64} :\equiv \mathrm{Conv}({(\lvert r \rvert)^{\bullet}},L) \to L \cdot 1 - L \cdot \lvert r \rvert = L \cdot (1 - \lvert r \rvert)$
\item $\varphi_{65} :\equiv \mathrm{Conv}({(\lvert r \rvert)^{\bullet}},L) \to L \cdot (1 - \lvert r \rvert) = L \cdot 1 - L \cdot \lvert r \rvert$
\item $\varphi_{66} :\equiv \mathrm{Conv}({(\lvert r \rvert)^{\bullet}},L) \to \mathrm{Conv}({(\lvert r \rvert)^{\bullet}} \circ \mathsf{s},\lvert r \rvert \cdot L)$
\item $\varphi_{67} :\equiv \forall n \in \mathbb{N}\,(0 \leq {{\lvert r \rvert}^{n}}) \to \exists \alpha\,(\mathrm{Conv}({(\lvert r \rvert)^{\bullet}},\alpha))$
\item $\varphi_{68} :\equiv \mathrm{Conv}({(\lvert r \rvert)^{\bullet}},L) \to n \in \mathbb{N} \to n \in \mathbb{N} \to \mathsf{s}(n) = S(n)$
\item $\varphi_{69} :\equiv \forall k \in \mathbb{N}\,(\mathsf{s}(k) = S(k)) \to n \in \mathbb{N} \to \mathsf{s}(n) = S(n)$
\item $\varphi_{70} :\equiv \forall x\,\forall y\,(x = y \to y = x) \to \forall y\,(\lvert {{r}^{n}} \rvert = y \to y = \lvert {{r}^{n}} \rvert)$
\item $\varphi_{71} :\equiv \mathrm{Conv}({(\lvert r \rvert)^{\bullet}},L) \to \mathrm{Conv}(\lvert r \rvert \odot {(\lvert r \rvert)^{\bullet}},\lvert r \rvert \cdot L)$
\item $\varphi_{72} :\equiv \lvert r \rvert < 1 \to \mathrm{Conv}({(\lvert r \rvert)^{\bullet}},L) \to n \in \mathbb{N} \to \lvert r \rvert \in \mathbb{R}$
\item $\varphi_{73} :\equiv r \in \mathbb{R} \land \lvert r \rvert < 1 \land \neg \mathrm{Conv}({(r)^{\bullet}},0) \to \neg \mathrm{Conv}({(\lvert r \rvert)^{\bullet}},L)$
\item $\varphi_{74} :\equiv \forall y\,(1 - \lvert r \rvert = y \to y = 1 - \lvert r \rvert) \to 1 - \lvert r \rvert = 0 \to 0 = 1 - \lvert r \rvert$
\item $\varphi_{75} :\equiv \mathrm{Conv}({(\lvert r \rvert)^{\bullet}},L) \to n \in \mathbb{N} \to \forall k \in \mathbb{N}\,(\mathsf{s}(k) = S(k))$
\item $\varphi_{76} :\equiv 1 - \lvert r \rvert \in \mathbb{R} \to 1 - \lvert r \rvert \neq 0 \to L \cdot (1 - \lvert r \rvert) = 0 \to L = 0$
\item $\varphi_{77} :\equiv \mathrm{Conv}({(\lvert r \rvert)^{\bullet}},L) \to \mathrm{SubIdx}(\mathsf{s}) \to \mathrm{Conv}({(\lvert r \rvert)^{\bullet}} \circ \mathsf{s},L)$
\item $\varphi_{78} :\equiv \forall y\,({(\lvert r \rvert)^{\bullet}} \circ \mathsf{s}(n) = y \to y = {(\lvert r \rvert)^{\bullet}} \circ \mathsf{s}(n))$
\item $\varphi_{79} :\equiv \forall x\,\forall y\,(x = y \to y = x) \to \forall y\,({{\lvert r \rvert}^{S(n)}} = y \to y = {{\lvert r \rvert}^{S(n)}})$
\item $\varphi_{80} :\equiv \forall s\,\forall \alpha\,(\mathrm{Conv}({(\lvert r \rvert)^{\bullet}},\alpha) \to \mathrm{SubIdx}(s) \to \mathrm{Conv}({(\lvert r \rvert)^{\bullet}} \circ s,\alpha))$
\item $\varphi_{81} :\equiv \mathrm{Conv}({(\lvert r \rvert)^{\bullet}},L) \to n \in \mathbb{N} \to \lvert {{r}^{n}} \rvert = {{\lvert r \rvert}^{n}}$
\item $\varphi_{82} :\equiv \mathrm{Conv}({(\lvert r \rvert)^{\bullet}},L) \to n \in \mathbb{N} \to {{\lvert r \rvert}^{n}} = \lvert {{r}^{n}} \rvert$
\item $\varphi_{83} :\equiv L - L \cdot \lvert r \rvert = L - \lvert r \rvert \cdot L \to L \cdot (1 - \lvert r \rvert) = L - \lvert r \rvert \cdot L$
\item $\varphi_{84} :\equiv L - L \cdot \lvert r \rvert = L - L \cdot \lvert r \rvert \to L - L \cdot \lvert r \rvert = L - \lvert r \rvert \cdot L$
\item $\varphi_{85} :\equiv 1 \in \mathbb{R} \to \lvert r \rvert \in \mathbb{R} \to L \cdot 1 - L \cdot \lvert r \rvert = L \cdot (1 - \lvert r \rvert)$
\item $\varphi_{86} :\equiv \mathrm{Conv}({(\lvert r \rvert)^{\bullet}},L) \to n \in \mathbb{N} \to \lvert {(r)^{\bullet}} \rvert(n) = \lvert {{r}^{n}} \rvert$
\item $\varphi_{87} :\equiv \mathrm{Conv}({(\lvert r \rvert)^{\bullet}},L) \to n \in \mathbb{N} \to \lvert {{r}^{n}} \rvert = \lvert {(r)^{\bullet}} \rvert(n)$
\item $\varphi_{88} :\equiv \mathrm{Conv}({(\lvert r \rvert)^{\bullet}},L) \to n \in \mathbb{N} \to {{\lvert r \rvert}^{n}} = \lvert {(r)^{\bullet}} \rvert(n)$
\item $\varphi_{89} :\equiv \forall y\,(\mathsf{s}(n) = y \to y = \mathsf{s}(n)) \to \mathsf{s}(n) = S(n) \to S(n) = \mathsf{s}(n)$
\item $\varphi_{90} :\equiv \forall x\,\forall y\,(x = y \to y = x) \to \forall y\,(L \cdot 1 - L \cdot \lvert r \rvert = y \to y = L \cdot 1 - L \cdot \lvert r \rvert)$
\item $\varphi_{91} :\equiv \forall a\,(L \in \mathbb{R} \to 1 \in \mathbb{R} \to a \in \mathbb{R} \to L \cdot 1 - L \cdot a = L \cdot (1 - a))$
\item $\varphi_{92} :\equiv \forall a\,(\mathrm{Seq}(a) \to \mathrm{Seq}(\lvert a \rvert) \land \forall n \in \mathbb{N}\,(\lvert a \rvert(n) = \lvert a_{n} \rvert))$
\item $\varphi_{93} :\equiv n \in \mathbb{N} \to \lvert r \rvert \odot {(\lvert r \rvert)^{\bullet}}(n) = \lvert r \rvert \cdot {{\lvert r \rvert}^{n}}$
\item $\varphi_{94} :\equiv n \in \mathbb{N} \to \lvert r \rvert \odot {(\lvert r \rvert)^{\bullet}}(n) = {{\lvert r \rvert}^{n}} \cdot \lvert r \rvert$
\item $\varphi_{95} :\equiv \mathrm{Seq}(\lvert {(r)^{\bullet}} \rvert) \land \forall n \in \mathbb{N}\,(\lvert {(r)^{\bullet}} \rvert(n) = \lvert {{r}^{n}} \rvert)$
\item $\varphi_{96} :\equiv \forall a\,\forall \alpha\,(\mathrm{Conv}(a,\alpha) \to \alpha \in \mathbb{R}) \to \forall \alpha\,(\mathrm{Conv}({(\lvert r \rvert)^{\bullet}},\alpha) \to \alpha \in \mathbb{R})$
\item $\varphi_{97} :\equiv \forall \alpha\,(\mathrm{Conv}({(\lvert r \rvert)^{\bullet}},\alpha) \to \alpha \in \mathbb{R}) \to \mathrm{Conv}({(\lvert r \rvert)^{\bullet}},L) \to L \in \mathbb{R}$
\item $\varphi_{98} :\equiv \forall a\,\forall \alpha\,(\lvert r \rvert \in \mathbb{R} \to \mathrm{Conv}(a,\alpha) \to \mathrm{Conv}(\lvert r \rvert \odot a,\lvert r \rvert \cdot \alpha))$
\item $\varphi_{99} :\equiv \forall \alpha\,(\mathrm{Conv}({(\lvert r \rvert)^{\bullet}},\alpha) \to \mathrm{SubIdx}(\mathsf{s}) \to \mathrm{Conv}({(\lvert r \rvert)^{\bullet}} \circ \mathsf{s},\alpha))$
\item $\varphi_{100} :\equiv \mathrm{Conv}({(\lvert r \rvert)^{\bullet}},L) \to \forall n \in \mathbb{N}\,({{\lvert r \rvert}^{n}} = \lvert {(r)^{\bullet}} \rvert(n))$
\item $\varphi_{101} :\equiv \forall \beta\,(\mathrm{Conv}({(\lvert r \rvert)^{\bullet}} \circ \mathsf{s},L) \to \mathrm{Conv}({(\lvert r \rvert)^{\bullet}} \circ \mathsf{s},\beta) \to L = \beta)$
\item $\varphi_{102} :\equiv \forall x\,\forall a\,(L \in \mathbb{R} \to x \in \mathbb{R} \to a \in \mathbb{R} \to L \cdot x - L \cdot a = L \cdot (x - a))$
\item $\varphi_{103} :\equiv \forall n \in \mathbb{N}\,(\lvert r \rvert \odot {(\lvert r \rvert)^{\bullet}}(n) = \lvert r \rvert \cdot {{\lvert r \rvert}^{n}})$
\item $\varphi_{104} :\equiv \forall a\,\forall M\,(\mathrm{MonoDec}(a) \to M \in \mathbb{R} \to \forall n \in \mathbb{N}\,(M \leq a_{n}) \to \exists \alpha\,(\mathrm{Conv}(a,\alpha)))$
\item $\varphi_{105} :\equiv L \cdot 1 - L \cdot \lvert r \rvert = L - L \cdot \lvert r \rvert \to L \cdot (1 - \lvert r \rvert) = L - L \cdot \lvert r \rvert$
\item $\varphi_{106} :\equiv \lvert r \rvert \cdot L = L \to L - \lvert r \rvert \cdot L = L - \lvert r \rvert \cdot L \to L - \lvert r \rvert \cdot L = L - L$
\item $\varphi_{107} :\equiv \lvert {(r)^{\bullet}} \rvert(n) = \lvert {{r}^{n}} \rvert \to \lvert {{r}^{n}} \rvert = \lvert {(r)^{\bullet}} \rvert(n)$
\item $\varphi_{108} :\equiv (r \in \mathbb{R} \to \lvert r \rvert < 1 \to r \in \mathbb{R}) \to r \in \mathbb{R} \to \lvert r \rvert < 1 \to 0 \leq \lvert r \rvert$
\item $\varphi_{109} :\equiv \lvert {{r}^{n}} \rvert = \lvert {(r)^{\bullet}} \rvert(n) \to {{\lvert r \rvert}^{n}} = \lvert {(r)^{\bullet}} \rvert(n)$
\item $\varphi_{110} :\equiv \forall c\,\forall x\,\forall a\,(c \in \mathbb{R} \to x \in \mathbb{R} \to a \in \mathbb{R} \to c \cdot x - c \cdot a = c \cdot (x - a))$
\item $\varphi_{111} :\equiv (r \in \mathbb{R} \to \lvert r \rvert < 1 \to \lvert r \rvert < 1) \to r \in \mathbb{R} \to \lvert r \rvert < 1 \to \lvert r \rvert \leq 1$
\item $\varphi_{112} :\equiv \forall a\,\forall b\,\forall \alpha\,(\mathrm{Seq}(b) \to \forall n \in \mathbb{N}\,(a_{n} = b_{n}) \to \mathrm{Conv}(a,\alpha) \to \mathrm{Conv}(b,\alpha))$
\item $\varphi_{113} :\equiv \forall z\,(L \cdot (1 - \lvert r \rvert) = L - \lvert r \rvert \cdot L \to L - \lvert r \rvert \cdot L = z \to L \cdot (1 - \lvert r \rvert) = z)$
\item $\varphi_{114} :\equiv \forall w\,(\lvert r \rvert \cdot L = w \to L - \lvert r \rvert \cdot L = L - \lvert r \rvert \cdot L \to L - \lvert r \rvert \cdot L = L - w)$
\item $\varphi_{115} :\equiv \forall z\,(L \cdot (1 - \lvert r \rvert) = L - L \cdot \lvert r \rvert \to L - L \cdot \lvert r \rvert = z \to L \cdot (1 - \lvert r \rvert) = z)$
\item $\varphi_{116} :\equiv \forall w\,(L \cdot \lvert r \rvert = w \to L - L \cdot \lvert r \rvert = L - L \cdot \lvert r \rvert \to L - L \cdot \lvert r \rvert = L - w)$
\item $\varphi_{117} :\equiv \forall y\,\forall z\,(\lvert r \rvert \odot {(\lvert r \rvert)^{\bullet}}(n) = y \to y = z \to \lvert r \rvert \odot {(\lvert r \rvert)^{\bullet}}(n) = z)$
\item $\varphi_{118} :\equiv \forall \alpha\,\forall \beta\,(\mathrm{Conv}({(\lvert r \rvert)^{\bullet}} \circ \mathsf{s},\alpha) \to \mathrm{Conv}({(\lvert r \rvert)^{\bullet}} \circ \mathsf{s},\beta) \to \alpha = \beta)$
\item $\varphi_{119} :\equiv L \cdot 1 - L \cdot \lvert r \rvert = L \cdot (1 - \lvert r \rvert) \to L \cdot (1 - \lvert r \rvert) = L \cdot 1 - L \cdot \lvert r \rvert$
\item $\varphi_{120} :\equiv L \cdot 1 - L \cdot \lvert r \rvert = z - L \cdot \lvert r \rvert \to L \cdot 1 - L \cdot \lvert r \rvert = w - L \cdot \lvert r \rvert$
\item $\varphi_{121} :\equiv \forall x\,\forall y\,\forall z\,(x = y \to y = z \to x = z) \to \forall y\,\forall z\,({{\lvert r \rvert}^{n}} = y \to y = z \to {{\lvert r \rvert}^{n}} = z)$
\item $\varphi_{122} :\equiv n \in \mathbb{N} \to n \in \mathbb{N} \to {(\lvert r \rvert)^{\bullet}} \circ \mathsf{s}(n) = {{\lvert r \rvert}^{\mathsf{s}(n)}}$
\item $\varphi_{123} :\equiv \mathrm{Conv}({(\lvert r \rvert)^{\bullet}},L) \to n \in \mathbb{N} \to n \in \mathbb{N} \to \lvert {{r}^{n}} \rvert = {{\lvert r \rvert}^{n}}$
\item $\varphi_{124} :\equiv \mathrm{Conv}({(\lvert r \rvert)^{\bullet}},L) \to n \in \mathbb{N} \to {{\lvert r \rvert}^{S(n)}} = {{\lvert r \rvert}^{\mathsf{s}(n)}}$
\item $\varphi_{125} :\equiv \mathrm{Conv}({(\lvert r \rvert)^{\bullet}},L) \to n \in \mathbb{N} \to {{\lvert r \rvert}^{S(n)}} = {{\lvert r \rvert}^{n}} \cdot \lvert r \rvert$
\item $\varphi_{126} :\equiv \mathrm{Conv}({(\lvert r \rvert)^{\bullet}},L) \to n \in \mathbb{N} \to {{\lvert r \rvert}^{n}} \cdot \lvert r \rvert = {{\lvert r \rvert}^{S(n)}}$
\item $\varphi_{127} :\equiv \mathrm{Seq}({(r)^{\bullet}}) \land {{r}^{1_{\omega}}} = r \land \forall n \in \mathbb{N}\,({{r}^{S(n)}} = {{r}^{n}} \cdot r)$
\item $\varphi_{128} :\equiv \mathrm{Conv}({(\lvert r \rvert)^{\bullet}},L) \to n \in \mathbb{N} \to n \in \mathbb{N} \to \lvert {(r)^{\bullet}} \rvert(n) = \lvert {{r}^{n}} \rvert$
\item $\varphi_{129} :\equiv (r \in \mathbb{R} \to \lvert r \rvert < 1 \to r \in \mathbb{R}) \to r \in \mathbb{R} \to \lvert r \rvert < 1 \to \lvert r \rvert \in \mathbb{R}$
\item $\varphi_{130} :\equiv r \in \mathbb{R} \to \lvert r \rvert < 1 \to \mathrm{Conv}({(\lvert r \rvert)^{\bullet}},L) \to \varphi_{3}$
\item $\varphi_{131} :\equiv n \in \mathbb{N} \to \forall k \in \mathbb{N}\,({(\lvert r \rvert)^{\bullet}} \circ \mathsf{s}(k) = {{\lvert r \rvert}^{\mathsf{s}(k)}})$
\item $\varphi_{132} :\equiv \mathrm{Conv}(\lvert r \rvert \odot {(\lvert r \rvert)^{\bullet}},\lvert r \rvert \cdot L) \to \mathrm{Conv}({(\lvert r \rvert)^{\bullet}} \circ \mathsf{s},\lvert r \rvert \cdot L)$
\item $\varphi_{133} :\equiv \mathrm{Conv}({(\lvert r \rvert)^{\bullet}},L) \to n \in \mathbb{N} \to \forall n \in \mathbb{N}\,(\lvert {{r}^{n}} \rvert = {{\lvert r \rvert}^{n}})$
\item $\varphi_{134} :\equiv (r \in \mathbb{R} \to \lvert r \rvert < 1 \to \lvert r \rvert \leq 1) \to r \in \mathbb{R} \to \lvert r \rvert < 1 \to \mathrm{MonoDec}({(\lvert r \rvert)^{\bullet}})$
\item $\varphi_{135} :\equiv \mathrm{Conv}({(\lvert r \rvert)^{\bullet}},L) \to n \in \mathbb{N} \to \lvert r \rvert \odot {(\lvert r \rvert)^{\bullet}}(n) = {{\lvert r \rvert}^{S(n)}}$
\item $\varphi_{136} :\equiv L \cdot 1 - L \cdot \lvert r \rvert = L \cdot 1 - L \cdot \lvert r \rvert \to L \cdot 1 - L \cdot \lvert r \rvert = L - L \cdot \lvert r \rvert$
\item $\varphi_{137} :\equiv \mathrm{Conv}({(\lvert r \rvert)^{\bullet}},L) \to n \in \mathbb{N} \to \forall n \in \mathbb{N}\,(\lvert {(r)^{\bullet}} \rvert(n) = \lvert {{r}^{n}} \rvert)$
\item $\varphi_{138} :\equiv L \cdot (1 - \lvert r \rvert) = L - \lvert r \rvert \cdot L \to L - \lvert r \rvert \cdot L = L - L \to L \cdot (1 - \lvert r \rvert) = L - L$
\item $\varphi_{139} :\equiv \forall \alpha\,(\lvert r \rvert \in \mathbb{R} \to \mathrm{Conv}({(\lvert r \rvert)^{\bullet}},\alpha) \to \mathrm{Conv}(\lvert r \rvert \odot {(\lvert r \rvert)^{\bullet}},\lvert r \rvert \cdot \alpha))$
\item $\varphi_{140} :\equiv \forall w\,(z = w \to {{\lvert r \rvert}^{S(n)}} = {{\lvert r \rvert}^{z}} \to {{\lvert r \rvert}^{S(n)}} = {{\lvert r \rvert}^{w}})$
\item $\varphi_{141} :\equiv \forall z\,(L \cdot (1 - \lvert r \rvert) = L \cdot 1 - L \cdot \lvert r \rvert \to L \cdot 1 - L \cdot \lvert r \rvert = z \to L \cdot (1 - \lvert r \rvert) = z)$
\item $\varphi_{142} :\equiv \mathrm{Conv}({(\lvert r \rvert)^{\bullet}},L) \to n \in \mathbb{N} \to \lvert r \rvert \cdot {{\lvert r \rvert}^{n}} = {{\lvert r \rvert}^{n}} \cdot \lvert r \rvert$
\item $\varphi_{143} :\equiv \forall a\,\forall s\,(\mathrm{Seq}(a) \to \mathrm{NSeq}(s) \to \mathrm{Seq}(a \circ s) \land \forall k \in \mathbb{N}\,(a \circ s(k) = a_{s_{k}}))$
\item $\varphi_{144} :\equiv \mathrm{Conv}({(\lvert r \rvert)^{\bullet}},L) \to n \in \mathbb{N} \to {(\lvert r \rvert)^{\bullet}} \circ \mathsf{s}(n) = {{\lvert r \rvert}^{\mathsf{s}(n)}}$
\item $\varphi_{145} :\equiv \mathrm{Conv}({(\lvert r \rvert)^{\bullet}},L) \to n \in \mathbb{N} \to {{\lvert r \rvert}^{\mathsf{s}(n)}} = {(\lvert r \rvert)^{\bullet}} \circ \mathsf{s}(n)$
\item $\varphi_{146} :\equiv \forall c\,\forall a\,(c \in \mathbb{R} \to \mathrm{Seq}(a) \to \mathrm{Seq}(c \odot a) \land \forall n \in \mathbb{N}\,(c \odot a(n) = c \cdot a_{n}))$
\item $\varphi_{147} :\equiv \forall x\,\forall y\,(x \in \mathbb{R} \to y \in \mathbb{R} \to x \cdot y = y \cdot x) \to \forall y\,(L \in \mathbb{R} \to y \in \mathbb{R} \to L \cdot y = y \cdot L)$
\item $\varphi_{148} :\equiv \mathrm{Conv}({(\lvert r \rvert)^{\bullet}},L) \to n \in \mathbb{N} \to \lvert r \rvert \odot {(\lvert r \rvert)^{\bullet}}(n) = {{\lvert r \rvert}^{\mathsf{s}(n)}}$
\item $\varphi_{149} :\equiv \forall b\,\forall \alpha\,(\mathrm{Seq}(b) \to \forall n \in \mathbb{N}\,({{\lvert r \rvert}^{n}} = b_{n}) \to \mathrm{Conv}({(\lvert r \rvert)^{\bullet}},\alpha) \to \mathrm{Conv}(b,\alpha))$
\item $\varphi_{150} :\equiv \mathrm{Conv}({(\lvert r \rvert)^{\bullet}},L) \to n \in \mathbb{N} \to n \in \mathbb{N} \to {{\lvert r \rvert}^{S(n)}} = {{\lvert r \rvert}^{n}} \cdot \lvert r \rvert$
\item $\varphi_{151} :\equiv \forall w\,(S(n) = w \to {{\lvert r \rvert}^{S(n)}} = {{\lvert r \rvert}^{S(n)}} \to {{\lvert r \rvert}^{S(n)}} = {{\lvert r \rvert}^{w}})$
\item $\varphi_{152} :\equiv \forall x\,\forall y\,(x \in \mathbb{R} \to y \in \mathbb{R} \to x - y \in \mathbb{R}) \to \forall y\,(1 \in \mathbb{R} \to y \in \mathbb{R} \to 1 - y \in \mathbb{R})$
\item $\varphi_{153} :\equiv r \in \mathbb{R} \land \lvert r \rvert < 1 \land \mathrm{Conv}({(\lvert r \rvert)^{\bullet}},L) \to \varphi_{22}$
\item $\varphi_{154} :\equiv r \in \mathbb{R} \land \lvert r \rvert < 1 \land \mathrm{Conv}({(\lvert r \rvert)^{\bullet}},L) \to n \in \mathbb{N} \to {{\lvert r \rvert}^{n}} = \lvert {(r)^{\bullet}} \rvert(n)$
\item $\varphi_{155} :\equiv \mathrm{Conv}({(\lvert r \rvert)^{\bullet}},L) \to n \in \mathbb{N} \to \forall n \in \mathbb{N}\,({{\lvert r \rvert}^{S(n)}} = {{\lvert r \rvert}^{n}} \cdot \lvert r \rvert)$
\item $\varphi_{156} :\equiv \forall w\,(L \cdot 1 = w \to L \cdot 1 - L \cdot \lvert r \rvert = L \cdot 1 - L \cdot \lvert r \rvert \to L \cdot 1 - L \cdot \lvert r \rvert = w - L \cdot \lvert r \rvert)$
\item $\varphi_{157} :\equiv \forall y\,(1 \in \mathbb{R} \to y \in \mathbb{R} \to 1 - y \in \mathbb{R}) \to 1 \in \mathbb{R} \to \lvert r \rvert \in \mathbb{R} \to 1 - \lvert r \rvert \in \mathbb{R}$
\item $\varphi_{158} :\equiv \forall M\,(\mathrm{MonoDec}({(\lvert r \rvert)^{\bullet}}) \to M \in \mathbb{R} \to \forall n \in \mathbb{N}\,(M \leq {{\lvert r \rvert}^{n}}) \to \exists \alpha\,(\mathrm{Conv}({(\lvert r \rvert)^{\bullet}},\alpha)))$
\item $\varphi_{159} :\equiv {{\lvert r \rvert}^{S(n)}} = {{\lvert r \rvert}^{n}} \cdot \lvert r \rvert \to {{\lvert r \rvert}^{n}} \cdot \lvert r \rvert = {{\lvert r \rvert}^{S(n)}}$
\item $\varphi_{160} :\equiv \forall y\,(L \in \mathbb{R} \to y \in \mathbb{R} \to L \cdot y = y \cdot L) \to L \in \mathbb{R} \to \lvert r \rvert \in \mathbb{R} \to L \cdot \lvert r \rvert = \lvert r \rvert \cdot L$
\item $\varphi_{161} :\equiv (r \in \mathbb{R} \to \lvert r \rvert < 1 \to 0 \leq \lvert r \rvert) \to r \in \mathbb{R} \to \lvert r \rvert < 1 \to \forall n \in \mathbb{N}\,(0 \leq {{\lvert r \rvert}^{n}})$
\item $\varphi_{162} :\equiv \forall n \in \mathbb{N}\,(\lvert {{r}^{n}} \rvert = {{\lvert r \rvert}^{n}}) \to n \in \mathbb{N} \to \lvert {{r}^{n}} \rvert = {{\lvert r \rvert}^{n}}$
\item $\varphi_{163} :\equiv \mathrm{Seq}({(\lvert r \rvert)^{\bullet}} \circ \mathsf{s}) \land \forall k \in \mathbb{N}\,({(\lvert r \rvert)^{\bullet}} \circ \mathsf{s}(k) = {{\lvert r \rvert}^{\mathsf{s}(k)}})$
\item $\varphi_{164} :\equiv \lvert r \rvert \in \mathbb{R} \to {{\lvert r \rvert}^{n}} \in \mathbb{R} \to \lvert r \rvert \cdot {{\lvert r \rvert}^{n}} = {{\lvert r \rvert}^{n}} \cdot \lvert r \rvert$
\item $\varphi_{165} :\equiv \forall n \in \mathbb{N}\,(\lvert {(r)^{\bullet}} \rvert(n) = \lvert {{r}^{n}} \rvert) \to n \in \mathbb{N} \to \lvert {(r)^{\bullet}} \rvert(n) = \lvert {{r}^{n}} \rvert$
\item $\varphi_{166} :\equiv {{\lvert r \rvert}^{n}} \cdot \lvert r \rvert = {{\lvert r \rvert}^{S(n)}} \to \lvert r \rvert \odot {(\lvert r \rvert)^{\bullet}}(n) = {{\lvert r \rvert}^{S(n)}}$
\item $\varphi_{167} :\equiv \mathrm{NSeq}(\mathsf{s}) \land \forall k \in \mathbb{N}\,(\mathsf{s}(k) = S(k)) \to r \in \mathbb{R} \to \varphi_{18}$
\item $\varphi_{168} :\equiv \forall x\,\forall y\,(x \in \mathbb{R} \to y \in \mathbb{R} \to x \cdot y = y \cdot x) \to \forall y\,(\lvert r \rvert \in \mathbb{R} \to y \in \mathbb{R} \to \lvert r \rvert \cdot y = y \cdot \lvert r \rvert)$
\item $\varphi_{169} :\equiv \forall x\,\forall y\,(x \in \mathbb{R} \to y \in \mathbb{R} \to y \neq 0 \to x \cdot y = 0 \to x = 0) \to \forall y\,(L \in \mathbb{R} \to y \in \mathbb{R} \to y \neq 0 \to L \cdot y = 0 \to L = 0)$
\item $\varphi_{170} :\equiv \lvert r \rvert < 1 \to \mathrm{Conv}({(\lvert r \rvert)^{\bullet}},L) \to n \in \mathbb{N} \to \varphi_{61}$
\item $\varphi_{171} :\equiv \forall a\,(\lvert r \rvert \in \mathbb{R} \to \mathrm{Seq}(a) \to \mathrm{Seq}(\lvert r \rvert \odot a) \land \forall n \in \mathbb{N}\,(\lvert r \rvert \odot a(n) = \lvert r \rvert \cdot a_{n}))$
\item $\varphi_{172} :\equiv {{\lvert r \rvert}^{S(n)}} = {{\lvert r \rvert}^{\mathsf{s}(n)}} \to \lvert r \rvert \odot {(\lvert r \rvert)^{\bullet}}(n) = {{\lvert r \rvert}^{\mathsf{s}(n)}}$
\item $\varphi_{173} :\equiv \lvert r \rvert < 1 \to \mathrm{Conv}({(\lvert r \rvert)^{\bullet}},L) \to \forall n \in \mathbb{N}\,(\varphi_{61})$
\item $\varphi_{174} :\equiv S(n) = \mathsf{s}(n) \to {{\lvert r \rvert}^{S(n)}} = {{\lvert r \rvert}^{S(n)}} \to {{\lvert r \rvert}^{S(n)}} = {{\lvert r \rvert}^{\mathsf{s}(n)}}$
\item $\varphi_{175} :\equiv \mathrm{Seq}({(\lvert r \rvert)^{\bullet}}) \land {{\lvert r \rvert}^{1_{\omega}}} = \lvert r \rvert \land \forall n \in \mathbb{N}\,({{\lvert r \rvert}^{S(n)}} = {{\lvert r \rvert}^{n}} \cdot \lvert r \rvert)$
\item $\varphi_{176} :\equiv \mathrm{Seq}(\lvert {(r)^{\bullet}} \rvert) \to \forall n \in \mathbb{N}\,({{\lvert r \rvert}^{n}} = \lvert {(r)^{\bullet}} \rvert(n)) \to \varphi_{10}$
\item $\varphi_{177} :\equiv \forall z\,(\lvert r \rvert \odot {(\lvert r \rvert)^{\bullet}}(n) = {{\lvert r \rvert}^{S(n)}} \to {{\lvert r \rvert}^{S(n)}} = z \to \lvert r \rvert \odot {(\lvert r \rvert)^{\bullet}}(n) = z)$
\item $\varphi_{178} :\equiv (\mathsf{s}(n) = S(n) \to S(n) = \mathsf{s}(n)) \to r \in \mathbb{R} \to \lvert r \rvert < 1 \to \varphi_{27}$
\item $\varphi_{179} :\equiv {(\lvert r \rvert)^{\bullet}} \circ \mathsf{s}(n) = {{\lvert r \rvert}^{\mathsf{s}(n)}} \to {{\lvert r \rvert}^{\mathsf{s}(n)}} = {(\lvert r \rvert)^{\bullet}} \circ \mathsf{s}(n)$
\item $\varphi_{180} :\equiv \lvert r \rvert \cdot {{\lvert r \rvert}^{n}} = {{\lvert r \rvert}^{n}} \cdot \lvert r \rvert \to \lvert r \rvert \odot {(\lvert r \rvert)^{\bullet}}(n) = {{\lvert r \rvert}^{n}} \cdot \lvert r \rvert$
\item $\varphi_{181} :\equiv \forall r \in \mathbb{R}\,(0 \leq r \to \forall n \in \mathbb{N}\,(0 \leq {{r}^{n}})) \to \lvert r \rvert \in \mathbb{R} \to 0 \leq \lvert r \rvert \to \forall n \in \mathbb{N}\,(0 \leq {{\lvert r \rvert}^{n}})$
\item $\varphi_{182} :\equiv \lvert r \rvert < 1 \to \mathrm{Conv}({(\lvert r \rvert)^{\bullet}},L) \to n \in \mathbb{N} \to \varphi_{93}$
\item $\varphi_{183} :\equiv r \in \mathbb{R} \to \lvert r \rvert < 1 \to \mathrm{Conv}({(\lvert r \rvert)^{\bullet}},L) \to \varphi_{93}$
\item $\varphi_{184} :\equiv r \in \mathbb{R} \to \lvert r \rvert < 1 \to \mathrm{Conv}({(\lvert r \rvert)^{\bullet}},L) \to \varphi_{94}$
\item $\varphi_{185} :\equiv \forall b\,\forall \alpha\,(\mathrm{Seq}(b) \to \forall n \in \mathbb{N}\,(\lvert r \rvert \odot {(\lvert r \rvert)^{\bullet}}(n) = b_{n}) \to \mathrm{Conv}(\lvert r \rvert \odot {(\lvert r \rvert)^{\bullet}},\alpha) \to \mathrm{Conv}(b,\alpha))$
\item $\varphi_{186} :\equiv r \in \mathbb{R} \to \lvert r \rvert < 1 \to \mathrm{Conv}({(\lvert r \rvert)^{\bullet}},L) \to \varphi_{95}$
\item $\varphi_{187} :\equiv r \in \mathbb{R} \land \lvert r \rvert < 1 \land \mathrm{Conv}({(\lvert r \rvert)^{\bullet}},L) \to n \in \mathbb{N} \to \varphi_{61}$
\item $\varphi_{188} :\equiv \forall x\,\forall y\,(x \in \mathbb{R} \to y \in \mathbb{R} \to x < y \to 0 < y - x) \to \forall y\,(\lvert r \rvert \in \mathbb{R} \to y \in \mathbb{R} \to \lvert r \rvert < y \to 0 < y - \lvert r \rvert)$
\item $\varphi_{189} :\equiv \forall \alpha\,(\mathrm{Seq}(\lvert {(r)^{\bullet}} \rvert) \to \forall n \in \mathbb{N}\,({{\lvert r \rvert}^{n}} = \lvert {(r)^{\bullet}} \rvert(n)) \to \mathrm{Conv}({(\lvert r \rvert)^{\bullet}},\alpha) \to \mathrm{Conv}(\lvert {(r)^{\bullet}} \rvert,\alpha))$
\item $\varphi_{190} :\equiv \forall y\,(\lvert {{r}^{n}} \rvert = y \to y = \lvert {{r}^{n}} \rvert) \to \lvert {{r}^{n}} \rvert = {{\lvert r \rvert}^{n}} \to {{\lvert r \rvert}^{n}} = \lvert {{r}^{n}} \rvert$
\item $\varphi_{191} :\equiv \lvert r \rvert < 1 \to \mathrm{Conv}({(\lvert r \rvert)^{\bullet}},L) \to n \in \mathbb{N} \to \varphi_{103}$
\item $\varphi_{192} :\equiv \forall s\,(\mathrm{Seq}({(\lvert r \rvert)^{\bullet}}) \to \mathrm{NSeq}(s) \to \mathrm{Seq}({(\lvert r \rvert)^{\bullet}} \circ s) \land \forall k \in \mathbb{N}\,({(\lvert r \rvert)^{\bullet}} \circ s(k) = {{\lvert r \rvert}^{s_{k}}}))$
\item $\varphi_{193} :\equiv \mathrm{Seq}({(\lvert r \rvert)^{\bullet}}) \to \mathrm{Seq}(\lvert r \rvert \odot {(\lvert r \rvert)^{\bullet}}) \land \varphi_{103}$
\item $\varphi_{194} :\equiv \forall n \in \mathbb{N}\,({{\lvert r \rvert}^{S(n)}} = {{\lvert r \rvert}^{n}} \cdot \lvert r \rvert) \to n \in \mathbb{N} \to {{\lvert r \rvert}^{S(n)}} = {{\lvert r \rvert}^{n}} \cdot \lvert r \rvert$
\item $\varphi_{195} :\equiv 0 \in \mathbb{R} \to (r \in \mathbb{R} \to \lvert r \rvert < 1 \to \forall n \in \mathbb{N}\,(0 \leq {{\lvert r \rvert}^{n}})) \to \varphi_{32}$
\item $\varphi_{196} :\equiv \forall z\,(\lvert r \rvert \odot {(\lvert r \rvert)^{\bullet}}(n) = {{\lvert r \rvert}^{\mathsf{s}(n)}} \to {{\lvert r \rvert}^{\mathsf{s}(n)}} = z \to \lvert r \rvert \odot {(\lvert r \rvert)^{\bullet}}(n) = z)$
\item $\varphi_{197} :\equiv r \in \mathbb{R} \to \lvert r \rvert < 1 \to \mathrm{Conv}({(\lvert r \rvert)^{\bullet}},L) \to \varphi_{122}$
\item $\varphi_{198} :\equiv \forall z\,(\lvert r \rvert \odot {(\lvert r \rvert)^{\bullet}}(n) = {{\lvert r \rvert}^{n}} \cdot \lvert r \rvert \to {{\lvert r \rvert}^{n}} \cdot \lvert r \rvert = z \to \lvert r \rvert \odot {(\lvert r \rvert)^{\bullet}}(n) = z)$
\item $\varphi_{199} :\equiv \forall z\,(\lvert r \rvert \odot {(\lvert r \rvert)^{\bullet}}(n) = \lvert r \rvert \cdot {{\lvert r \rvert}^{n}} \to \lvert r \rvert \cdot {{\lvert r \rvert}^{n}} = z \to \lvert r \rvert \odot {(\lvert r \rvert)^{\bullet}}(n) = z)$
\item $\varphi_{200} :\equiv \forall y\,(\lvert r \rvert \in \mathbb{R} \to y \in \mathbb{R} \to \lvert r \rvert < y \to 0 < y - \lvert r \rvert) \to \lvert r \rvert \in \mathbb{R} \to 1 \in \mathbb{R} \to \lvert r \rvert < 1 \to 0 < 1 - \lvert r \rvert$
\item $\varphi_{201} :\equiv r \in \mathbb{R} \to \lvert r \rvert < 1 \to \mathrm{Conv}({(\lvert r \rvert)^{\bullet}},L) \to \varphi_{127}$
\item $\varphi_{202} :\equiv \forall y\,(L \in \mathbb{R} \to y \in \mathbb{R} \to y \neq 0 \to L \cdot y = 0 \to L = 0) \to L \in \mathbb{R} \to \varphi_{76}$
\item $\varphi_{203} :\equiv r \in \mathbb{R} \to \lvert r \rvert < 1 \to \mathrm{Conv}({(\lvert r \rvert)^{\bullet}},L) \to \varphi_{131}$
\item $\varphi_{204} :\equiv \forall r \in \mathbb{R}\,\forall n \in \mathbb{N}\,(\lvert {{r}^{n}} \rvert = {{\lvert r \rvert}^{n}}) \to r \in \mathbb{R} \to \forall n \in \mathbb{N}\,(\lvert {{r}^{n}} \rvert = {{\lvert r \rvert}^{n}})$
\item $\varphi_{205} :\equiv \mathrm{Conv}({(\lvert r \rvert)^{\bullet}},L) \to n \in \mathbb{N} \to \mathrm{Seq}(\lvert r \rvert \odot {(\lvert r \rvert)^{\bullet}}) \land \varphi_{103}$
\item $\varphi_{206} :\equiv \forall L\,(\varphi_{73}) \to \varphi_{15} \to \forall L\,(\neg \mathrm{Conv}({(\lvert r \rvert)^{\bullet}},L))$
\item $\varphi_{207} :\equiv \mathrm{NSeq}(\mathsf{s}) \land \forall k \in \mathbb{N}\,(\mathsf{s}(k) = S(k)) \to r \in \mathbb{R} \to \lvert r \rvert < 1 \to \varphi_{75}$
\item $\varphi_{208} :\equiv \forall y\,\forall z\,({{\lvert r \rvert}^{n}} = y \to y = z \to {{\lvert r \rvert}^{n}} = z) \to \forall z\,({{\lvert r \rvert}^{n}} = \lvert {{r}^{n}} \rvert \to \lvert {{r}^{n}} \rvert = z \to {{\lvert r \rvert}^{n}} = z)$
\item $\varphi_{209} :\equiv \forall k \in \mathbb{N}\,({(\lvert r \rvert)^{\bullet}} \circ \mathsf{s}(k) = {{\lvert r \rvert}^{\mathsf{s}(k)}}) \to n \in \mathbb{N} \to {(\lvert r \rvert)^{\bullet}} \circ \mathsf{s}(n) = {{\lvert r \rvert}^{\mathsf{s}(n)}}$
\item $\varphi_{210} :\equiv r \in \mathbb{R} \to \lvert r \rvert < 1 \to \mathrm{Conv}({(\lvert r \rvert)^{\bullet}},L) \to \varphi_{163}$
\item $\varphi_{211} :\equiv (r \in \mathbb{R} \to \varphi_{24}) \to (L = \lvert r \rvert \cdot L \to \lvert r \rvert \cdot L = L) \to r \in \mathbb{R} \to \varphi_{25}$
\item $\varphi_{212} :\equiv \lvert r \rvert < 1 \to \mathrm{Conv}({(\lvert r \rvert)^{\bullet}},L) \to n \in \mathbb{N} \to \varphi_{175}$
\item $\varphi_{213} :\equiv r \in \mathbb{R} \to \lvert r \rvert < 1 \to \mathrm{Conv}({(\lvert r \rvert)^{\bullet}},L) \to \varphi_{175}$
\item $\varphi_{214} :\equiv (r \in \mathbb{R} \to \varphi_{28}) \to (1 - \lvert r \rvert = 0 \to 0 = 1 - \lvert r \rvert) \to r \in \mathbb{R} \to \varphi_{29}$
\item $\varphi_{215} :\equiv (\mathrm{Conv}({(\lvert r \rvert)^{\bullet}},L) \to L \in \mathbb{R}) \to (r \in \mathbb{R} \to \lvert r \rvert < 1 \to \varphi_{9}) \to \varphi_{53}$
\item $\varphi_{216} :\equiv (r \in \mathbb{R} \to \lvert r \rvert < 1 \to \varphi_{9}) \to \mathrm{SubIdx}(\mathsf{s}) \to r \in \mathbb{R} \to \lvert r \rvert < 1 \to \varphi_{34}$
\item $\varphi_{217} :\equiv (0 < 1 - \lvert r \rvert \to 0 \neq 1 - \lvert r \rvert) \to (r \in \mathbb{R} \to \varphi_{38}) \to r \in \mathbb{R} \to \varphi_{28}$
\item $\varphi_{218} :\equiv (r \in \mathbb{R} \to \lvert r \rvert < 1 \to \lvert r \rvert \in \mathbb{R}) \to (r \in \mathbb{R} \to \lvert r \rvert < 1 \to 0 \leq \lvert r \rvert) \to \varphi_{134}$
\item $\varphi_{219} :\equiv (\lvert {{r}^{n}} \rvert = {{\lvert r \rvert}^{n}} \to {{\lvert r \rvert}^{n}} = \lvert {{r}^{n}} \rvert) \to r \in \mathbb{R} \to \lvert r \rvert < 1 \to \varphi_{82}$
\item $\varphi_{220} :\equiv \forall \alpha\,(\mathrm{Seq}({(\lvert r \rvert)^{\bullet}} \circ \mathsf{s}) \to \forall n \in \mathbb{N}\,(\varphi_{61}) \to \mathrm{Conv}(\lvert r \rvert \odot {(\lvert r \rvert)^{\bullet}},\alpha) \to \mathrm{Conv}({(\lvert r \rvert)^{\bullet}} \circ \mathsf{s},\alpha))$
\item $\varphi_{221} :\equiv (r \in \mathbb{R} \to \lvert r \rvert < 1 \to \mathrm{MonoDec}({(\lvert r \rvert)^{\bullet}})) \to \varphi_{195}$
\item $\varphi_{222} :\equiv \lvert r \rvert \odot {(\lvert r \rvert)^{\bullet}}(n) = {{\lvert r \rvert}^{\mathsf{s}(n)}} \to {{\lvert r \rvert}^{\mathsf{s}(n)}} = {(\lvert r \rvert)^{\bullet}} \circ \mathsf{s}(n) \to \varphi_{61}$
\item $\varphi_{223} :\equiv (r \in \mathbb{R} \to \lvert r \rvert < 1 \to \varphi_{9}) \to (L = 0 \to \varphi_{8}) \to r \in \mathbb{R} \to \lvert r \rvert < 1 \to \varphi_{8}$
\item $\varphi_{224} :\equiv (r \in \mathbb{R} \to \varphi_{33}) \to 1 \in \mathbb{R} \to (r \in \mathbb{R} \to \lvert r \rvert < 1 \to \lvert r \rvert < 1) \to r \in \mathbb{R} \to \varphi_{38}$
\item $\varphi_{225} :\equiv (L \cdot (1 - \lvert r \rvert) = L - L \cdot \lvert r \rvert \to \varphi_{83}) \to r \in \mathbb{R} \to \lvert r \rvert < 1 \to \varphi_{44}$
\item $\varphi_{226} :\equiv (\mathrm{Seq}({(r)^{\bullet}}) \to \varphi_{3}) \to (r \in \mathbb{R} \to \lvert r \rvert < 1 \to \varphi_{1}) \to \varphi_{130}$
\item $\varphi_{227} :\equiv \forall z\,({{\lvert r \rvert}^{n}} = \lvert {{r}^{n}} \rvert \to \lvert {{r}^{n}} \rvert = z \to {{\lvert r \rvert}^{n}} = z) \to {{\lvert r \rvert}^{n}} = \lvert {{r}^{n}} \rvert \to \varphi_{109}$
\item $\varphi_{228} :\equiv (L \cdot (1 - \lvert r \rvert) = L \cdot 1 - L \cdot \lvert r \rvert \to \varphi_{105}) \to r \in \mathbb{R} \to \lvert r \rvert < 1 \to \varphi_{43}$
\item $\varphi_{229} :\equiv (1 \in \mathbb{R} \to \lvert r \rvert \in \mathbb{R} \to 1 - \lvert r \rvert \in \mathbb{R}) \to 1 \in \mathbb{R} \to (r \in \mathbb{R} \to \varphi_{33}) \to r \in \mathbb{R} \to \varphi_{47}$
\item $\varphi_{230} :\equiv (\forall \alpha\,(\varphi_{17}) \to \varphi_{40}) \to \varphi_{35} \to r \in \mathbb{R} \to \lvert r \rvert < 1 \to \mathrm{Conv}({(r)^{\bullet}},0)$
\item $\varphi_{231} :\equiv \forall n\,(\varphi_{154}) \to \varphi_{22} \to \forall n \in \mathbb{N}\,({{\lvert r \rvert}^{n}} = \lvert {(r)^{\bullet}} \rvert(n))$
\item $\varphi_{232} :\equiv (\lvert r \rvert \in \mathbb{R} \to 0 \leq \lvert r \rvert \to \forall n \in \mathbb{N}\,(0 \leq {{\lvert r \rvert}^{n}})) \to (r \in \mathbb{R} \to \lvert r \rvert < 1 \to \lvert r \rvert \in \mathbb{R}) \to \varphi_{161}$
\item $\varphi_{233} :\equiv L - L \cdot \lvert r \rvert = L - L \cdot \lvert r \rvert \to (L \cdot \lvert r \rvert = \lvert r \rvert \cdot L \to \varphi_{84}) \to r \in \mathbb{R} \to \lvert r \rvert < 1 \to \varphi_{42}$
\item $\varphi_{234} :\equiv (r \in \mathbb{R} \to \varphi_{33}) \to (r \in \mathbb{R} \to \lvert r \rvert < 1 \to \varphi_{9}) \to r \in \mathbb{R} \to \lvert r \rvert < 1 \to \varphi_{71}$
\item $\varphi_{235} :\equiv L \cdot 1 - L \cdot \lvert r \rvert = L \cdot 1 - L \cdot \lvert r \rvert \to (L \cdot 1 = L \to \varphi_{136}) \to r \in \mathbb{R} \to \lvert r \rvert < 1 \to \varphi_{63}$
\item $\varphi_{236} :\equiv (r \in \mathbb{R} \to \lvert r \rvert < 1 \to \varphi_{34}) \to (r \in \mathbb{R} \to \lvert r \rvert < 1 \to \varphi_{66}) \to r \in \mathbb{R} \to \varphi_{24}$
\item $\varphi_{237} :\equiv (r \in \mathbb{R} \to \lvert r \rvert < 1 \to \varphi_{23}) \to (r \in \mathbb{R} \to \varphi_{50}) \to r \in \mathbb{R} \to \lvert r \rvert < 1 \to \varphi_{62}$
\item $\varphi_{238} :\equiv (r \in \mathbb{R} \to \lvert r \rvert < 1 \to \varphi_{100}) \to (r \in \mathbb{R} \to \lvert r \rvert < 1 \to \varphi_{8}) \to r \in \mathbb{R} \to \lvert r \rvert < 1 \to \varphi_{11}$
\item $\varphi_{239} :\equiv (r \in \mathbb{R} \to \lvert r \rvert < 1 \to \varphi_{68}) \to (r \in \mathbb{R} \to \varphi_{50}) \to r \in \mathbb{R} \to \lvert r \rvert < 1 \to \varphi_{26}$
\item $\varphi_{240} :\equiv (r \in \mathbb{R} \to \varphi_{25}) \to L - \lvert r \rvert \cdot L = L - \lvert r \rvert \cdot L \to \varphi_{106} \to r \in \mathbb{R} \to \varphi_{57}$
\item $\varphi_{241} :\equiv (r \in \mathbb{R} \to \forall n \in \mathbb{N}\,(\lvert {{r}^{n}} \rvert = {{\lvert r \rvert}^{n}})) \to (r \in \mathbb{R} \to \varphi_{51}) \to r \in \mathbb{R} \to \lvert r \rvert < 1 \to \varphi_{133}$
\item $\varphi_{242} :\equiv (L \in \mathbb{R} \to \lvert r \rvert \in \mathbb{R} \to L \cdot \lvert r \rvert = \lvert r \rvert \cdot L) \to \varphi_{53} \to (r \in \mathbb{R} \to \varphi_{33}) \to r \in \mathbb{R} \to \varphi_{58}$
\item $\varphi_{243} :\equiv (\mathrm{Seq}({(r)^{\bullet}}) \to \varphi_{95}) \to (r \in \mathbb{R} \to \lvert r \rvert < 1 \to \varphi_{1}) \to \varphi_{186}$
\item $\varphi_{244} :\equiv (\lvert r \rvert \odot {(\lvert r \rvert)^{\bullet}}(n) = {{\lvert r \rvert}^{n}} \cdot \lvert r \rvert \to \varphi_{166}) \to r \in \mathbb{R} \to \lvert r \rvert < 1 \to \varphi_{135}$
\item $\varphi_{245} :\equiv (r \in \mathbb{R} \to \lvert r \rvert < 1 \to \varphi_{123}) \to (r \in \mathbb{R} \to \varphi_{50}) \to r \in \mathbb{R} \to \lvert r \rvert < 1 \to \varphi_{81}$
\item $\varphi_{246} :\equiv (r \in \mathbb{R} \to \lvert r \rvert < 1 \to \varphi_{128}) \to (r \in \mathbb{R} \to \varphi_{50}) \to r \in \mathbb{R} \to \lvert r \rvert < 1 \to \varphi_{86}$
\item $\varphi_{247} :\equiv (r \in \mathbb{R} \to \varphi_{72}) \to (r \in \mathbb{R} \to \lvert r \rvert < 1 \to \varphi_{62}) \to r \in \mathbb{R} \to \lvert r \rvert < 1 \to \varphi_{142}$
\item $\varphi_{248} :\equiv (\lvert r \rvert \odot {(\lvert r \rvert)^{\bullet}}(n) = {{\lvert r \rvert}^{S(n)}} \to \varphi_{172}) \to r \in \mathbb{R} \to \lvert r \rvert < 1 \to \varphi_{148}$
\item $\varphi_{249} :\equiv (\mathrm{Seq}({(\lvert r \rvert)^{\bullet}}) \to n \in \mathbb{N} \to {{\lvert r \rvert}^{n}} \in \mathbb{R}) \to \varphi_{237}$
\item $\varphi_{250} :\equiv (\lvert r \rvert \odot {(\lvert r \rvert)^{\bullet}}(n) = \lvert r \rvert \cdot {{\lvert r \rvert}^{n}} \to \varphi_{180}) \to \varphi_{184}$
\item $\varphi_{251} :\equiv (r \in \mathbb{R} \to \lvert r \rvert < 1 \to \varphi_{150}) \to (r \in \mathbb{R} \to \varphi_{50}) \to r \in \mathbb{R} \to \lvert r \rvert < 1 \to \varphi_{125}$
\item $\varphi_{252} :\equiv (r \in \mathbb{R} \to \varphi_{173}) \to (r \in \mathbb{R} \to \lvert r \rvert < 1 \to \varphi_{71}) \to r \in \mathbb{R} \to \lvert r \rvert < 1 \to \varphi_{66}$
\item $\varphi_{253} :\equiv (r \in \mathbb{R} \to \varphi_{72}) \to (r \in \mathbb{R} \to \lvert r \rvert < 1 \to \varphi_{4}) \to r \in \mathbb{R} \to \lvert r \rvert < 1 \to \varphi_{205}$
\item $\varphi_{254} :\equiv (L \in \mathbb{R} \to \varphi_{85}) \to \varphi_{53} \to 1 \in \mathbb{R} \to (r \in \mathbb{R} \to \varphi_{33}) \to r \in \mathbb{R} \to \lvert r \rvert < 1 \to \varphi_{64}$
\item $\varphi_{255} :\equiv (r \in \mathbb{R} \to \lvert r \rvert < 1 \to \varphi_{87}) \to ({{\lvert r \rvert}^{n}} = \lvert {{r}^{n}} \rvert \to \varphi_{109}) \to r \in \mathbb{R} \to \lvert r \rvert < 1 \to \varphi_{88}$
\item $\varphi_{256} :\equiv (r \in \mathbb{R} \to \lvert r \rvert < 1 \to \varphi_{44}) \to (r \in \mathbb{R} \to \varphi_{57}) \to \varphi_{138} \to r \in \mathbb{R} \to \lvert r \rvert < 1 \to \varphi_{16}$
\item $\varphi_{257} :\equiv (r \in \mathbb{R} \to \lvert r \rvert < 1 \to \varphi_{27}) \to {{\lvert r \rvert}^{S(n)}} = {{\lvert r \rvert}^{S(n)}} \to \varphi_{174} \to r \in \mathbb{R} \to \lvert r \rvert < 1 \to \varphi_{124}$
\item $\varphi_{258} :\equiv (r \in \mathbb{R} \land \lvert r \rvert < 1 \land \neg \mathrm{Conv}({(r)^{\bullet}},0) \to \varphi_{15}) \to \forall \alpha\,(\varphi_{17}) \to \varphi_{230}$
\item $\varphi_{259} :\equiv (r \in \mathbb{R} \to \lvert r \rvert < 1 \to \varphi_{43}) \to (r \in \mathbb{R} \to \lvert r \rvert < 1 \to \varphi_{42}) \to \varphi_{225}$
\item $\varphi_{260} :\equiv (r \in \mathbb{R} \to \lvert r \rvert < 1 \to \varphi_{65}) \to (r \in \mathbb{R} \to \lvert r \rvert < 1 \to \varphi_{63}) \to \varphi_{228}$
\item $\varphi_{261} :\equiv (L \in \mathbb{R} \to \varphi_{76}) \to \varphi_{53} \to (r \in \mathbb{R} \to \varphi_{47}) \to (r \in \mathbb{R} \to \varphi_{29}) \to (r \in \mathbb{R} \to \varphi_{41}) \to \varphi_{37}$
\item $\varphi_{262} :\equiv (\mathrm{Seq}({(\lvert r \rvert)^{\bullet}}) \to \mathrm{NSeq}(\mathsf{s}) \to \varphi_{163}) \to (r \in \mathbb{R} \to \lvert r \rvert < 1 \to \varphi_{4}) \to (r \in \mathbb{R} \to \varphi_{18}) \to \varphi_{210}$
\item $\varphi_{263} :\equiv (r \in \mathbb{R} \to \lvert r \rvert < 1 \to \varphi_{135}) \to (r \in \mathbb{R} \to \lvert r \rvert < 1 \to \varphi_{124}) \to \varphi_{248}$
\item $\varphi_{264} :\equiv (r \in \mathbb{R} \to \lvert r \rvert < 1 \to \varphi_{148}) \to (r \in \mathbb{R} \to \lvert r \rvert < 1 \to \varphi_{145}) \to \varphi_{222} \to r \in \mathbb{R} \to \varphi_{170}$
\item $\varphi_{265} :\equiv (\mathrm{Seq}({(\lvert r \rvert)^{\bullet}} \circ \mathsf{s}) \to \forall n \in \mathbb{N}\,(\varphi_{61}) \to \varphi_{132}) \to (r \in \mathbb{R} \to \lvert r \rvert < 1 \to \varphi_{20}) \to \varphi_{252}$
\item $\varphi_{266} :\equiv \forall n\,(\varphi_{187}) \to (\forall n\,(\varphi_{187}) \to \varphi_{22} \to \forall n \in \mathbb{N}\,(\varphi_{61})) \to \varphi_{153} \to r \in \mathbb{R} \to \varphi_{173}$
\end{abbrevlist}
\end{fullproof}

% 自動生成 (prover/build.js): GeoOne  ノード数 77
\begin{fullproof}{命題3.1(2)}{$\mathrm{Conv}({(1)^{\bullet}},1)$}{77}\phantomsection\label{tree:GeoOne}
\noindent$\Pi_{1}$:
\begin{ptree}
\AXC{$1 \in \mathbb{R}$\lbl{(R1)}}
\AXC{$\forall c \in \mathbb{R}\,(\mathrm{Conv}(\overline{c},c))$\lbl{[G4]}}
\AXC{$\varphi_{2}$\lbl{A2}}
\BIC{$1 \in \mathbb{R} \to \mathrm{Conv}(\overline{1},1)$}
\AXC{$\varphi_{1}$\lbl{A1}}
\BIC{$1 \in \mathbb{R} \to \mathrm{Conv}(\overline{1},1)$}
\BIC{$\mathrm{Conv}(\overline{1},1)$}
\end{ptree}
\noindent$\Pi_{2}$:
\begin{ptree}
\AXC{$y = z \to x = y \to x = z$\lbl{A6}}
\AXC{$(y = z \to x = y \to x = z) \to x = y \to y = z \to x = z$\lbl{A1}}
\BIC{$x = y \to y = z \to x = z$}
\UIC{$\forall z\,(x = y \to y = z \to x = z)$}
\UIC{$\forall y\,\forall z\,(x = y \to y = z \to x = z)$}
\UIC{$\forall x\,\forall y\,\forall z\,(x = y \to y = z \to x = z)$}
\AXC{$\varphi_{3}$\lbl{A2}}
\BIC{$\forall y\,\forall z\,(\overline{1}(n) = y \to y = z \to \overline{1}(n) = z)$}
\AXC{$\varphi_{10}$\lbl{A2}}
\BIC{$\forall z\,(\overline{1}(n) = 1 \to 1 = z \to \overline{1}(n) = z)$}
\AXC{$\varphi_{16}$\lbl{A2}}
\BIC{$\overline{1}(n) = 1 \to 1 = {{1}^{n}} \to \overline{1}(n) = {{1}^{n}}$}
\end{ptree}
\noindent$\Pi_{3}$:
\begin{ptree}
\AXC{$x = x$\lbl{A5}}
\AXC{$x = y \to x = x \to y = x$\lbl{A6}}
\AXC{$(x = y \to x = x \to y = x) \to x = x \to x = y \to y = x$\lbl{A1}}
\BIC{$x = x \to x = y \to y = x$}
\BIC{$x = y \to y = x$}
\UIC{$\forall y\,(x = y \to y = x)$}
\UIC{$\forall x\,\forall y\,(x = y \to y = x)$}
\end{ptree}
\noindent$\Pi_{4}$:
\begin{ptree}
\AXC{$\stackrel{\Pi_{3}}{\vdots}$}\noLine
\UIC{$\forall x\,\forall y\,(x = y \to y = x)$}
\AXC{$\forall x\,\forall y\,(x = y \to y = x) \to \forall y\,({{1}^{n}} = y \to y = {{1}^{n}})$\lbl{A2}}
\BIC{$\forall y\,({{1}^{n}} = y \to y = {{1}^{n}})$}
\end{ptree}
\noindent$\Pi_{5}$:
\begin{ptree}
\AXC{$\forall n \in \mathbb{N}\,({{1}^{n}} = 1)$\lbl{[P5]}}
\AXC{$\forall n \in \mathbb{N}\,({{1}^{n}} = 1) \to n \in \mathbb{N} \to {{1}^{n}} = 1$\lbl{A2}}
\BIC{$n \in \mathbb{N} \to {{1}^{n}} = 1$}
\AXC{$\varphi_{8}$\lbl{A1}}
\BIC{$({{1}^{n}} = 1 \to 1 = {{1}^{n}}) \to n \in \mathbb{N} \to 1 = {{1}^{n}}$}
\end{ptree}
\noindent$\Pi_{6}$:
\begin{ptree}
\AXC{$\stackrel{\Pi_{4}}{\vdots}$}\noLine
\UIC{$\forall y\,({{1}^{n}} = y \to y = {{1}^{n}})$}
\AXC{$\forall y\,({{1}^{n}} = y \to y = {{1}^{n}}) \to {{1}^{n}} = 1 \to 1 = {{1}^{n}}$\lbl{A2}}
\BIC{${{1}^{n}} = 1 \to 1 = {{1}^{n}}$}
\AXC{$\stackrel{\Pi_{5}}{\vdots}$}\noLine
\UIC{$({{1}^{n}} = 1 \to 1 = {{1}^{n}}) \to n \in \mathbb{N} \to 1 = {{1}^{n}}$}
\BIC{$n \in \mathbb{N} \to 1 = {{1}^{n}}$}
\end{ptree}
\noindent$\Pi_{7}$:
\begin{ptree}
\AXC{$1 \in \mathbb{R}$\lbl{(R1)}}
\AXC{$\forall c \in \mathbb{R}\,(\mathrm{Seq}(\overline{c}) \land \forall n \in \mathbb{N}\,(\overline{c}(n) = c))$\lbl{[S10]}}
\AXC{$\varphi_{18}$\lbl{A2}}
\BIC{$1 \in \mathbb{R} \to \mathrm{Seq}(\overline{1}) \land \forall n \in \mathbb{N}\,(\overline{1}(n) = 1)$}
\AXC{$\varphi_{17}$\lbl{A1}}
\BIC{$1 \in \mathbb{R} \to \mathrm{Seq}(\overline{1}) \land \forall n \in \mathbb{N}\,(\overline{1}(n) = 1)$}
\BIC{$\mathrm{Seq}(\overline{1}) \land \forall n \in \mathbb{N}\,(\overline{1}(n) = 1)$}
\AXC{$\varphi_{9}$\lbl{A1}}
\BIC{$n \in \mathbb{N} \to \forall n \in \mathbb{N}\,(\overline{1}(n) = 1)$}
\end{ptree}
\noindent$\Pi_{8}$:
\begin{ptree}
\AXC{$\stackrel{\Pi_{7}}{\vdots}$}\noLine
\UIC{$n \in \mathbb{N} \to \forall n \in \mathbb{N}\,(\overline{1}(n) = 1)$}
\AXC{$(n \in \mathbb{N} \to \forall n \in \mathbb{N}\,(\overline{1}(n) = 1)) \to \varphi_{14}$\lbl{A1}}
\BIC{$\varphi_{14}$}
\end{ptree}
\noindent$\Pi_{9}$:
\begin{ptree}
\AXC{$n \in \mathbb{N} \to n \in \mathbb{N}$\lbl{A1}}
\AXC{$\forall n \in \mathbb{N}\,(\overline{1}(n) = 1) \to n \in \mathbb{N} \to \overline{1}(n) = 1$\lbl{A2}}
\AXC{$\stackrel{\Pi_{8}}{\vdots}$}\noLine
\UIC{$\varphi_{14}$}
\BIC{$n \in \mathbb{N} \to n \in \mathbb{N} \to \overline{1}(n) = 1$}
\AXC{$\varphi_{11}$\lbl{A1}}
\BIC{$(n \in \mathbb{N} \to n \in \mathbb{N}) \to n \in \mathbb{N} \to \overline{1}(n) = 1$}
\BIC{$n \in \mathbb{N} \to \overline{1}(n) = 1$}
\end{ptree}
\noindent$\Pi_{10}$:
\begin{ptree}
\AXC{$\stackrel{\Pi_{6}}{\vdots}$}\noLine
\UIC{$n \in \mathbb{N} \to 1 = {{1}^{n}}$}
\AXC{$\stackrel{\Pi_{9}}{\vdots}$}\noLine
\UIC{$n \in \mathbb{N} \to \overline{1}(n) = 1$}
\AXC{$(n \in \mathbb{N} \to \overline{1}(n) = 1) \to (n \in \mathbb{N} \to 1 = {{1}^{n}}) \to \varphi_{15}$\lbl{A1}}
\BIC{$(n \in \mathbb{N} \to 1 = {{1}^{n}}) \to \varphi_{15}$}
\BIC{$\varphi_{15}$}
\end{ptree}
\noindent$\Pi_{11}$:
\begin{ptree}
\AXC{$1 \in \mathbb{R}$\lbl{(R1)}}
\AXC{$\varphi_{13}$\lbl{[S24]}}
\AXC{$\varphi_{13} \to 1 \in \mathbb{R} \to \varphi_{6}$\lbl{A2}}
\BIC{$1 \in \mathbb{R} \to \varphi_{6}$}
\AXC{$(1 \in \mathbb{R} \to \varphi_{6}) \to 1 \in \mathbb{R} \to \varphi_{6}$\lbl{A1}}
\BIC{$1 \in \mathbb{R} \to \varphi_{6}$}
\BIC{$\varphi_{6}$}
\AXC{$\varphi_{6} \to \mathrm{Seq}({(1)^{\bullet}})$\lbl{A1}}
\BIC{$\mathrm{Seq}({(1)^{\bullet}})$}
\end{ptree}
\noindent$\Pi_{12}$:
\begin{ptree}
\AXC{$\varphi_{5}$\lbl{[C13]}}
\AXC{$\varphi_{5} \to \varphi_{7}$\lbl{A2}}
\BIC{$\varphi_{7}$}
\AXC{$\varphi_{7} \to \varphi_{12}$\lbl{A2}}
\BIC{$\varphi_{12}$}
\AXC{$\varphi_{12} \to \mathrm{Seq}({(1)^{\bullet}}) \to \varphi_{4}$\lbl{A2}}
\BIC{$\mathrm{Seq}({(1)^{\bullet}}) \to \varphi_{4}$}
\AXC{$\varphi_{19}$\lbl{A1}}
\BIC{$\mathrm{Seq}({(1)^{\bullet}}) \to \varphi_{4}$}
\end{ptree}
\noindent$\Pi_{13}$:
\begin{ptree}
\AXC{$\stackrel{\Pi_{2}}{\vdots}$}\noLine
\UIC{$\overline{1}(n) = 1 \to 1 = {{1}^{n}} \to \overline{1}(n) = {{1}^{n}}$}
\AXC{$\stackrel{\Pi_{10}}{\vdots}$}\noLine
\UIC{$\varphi_{15}$}
\BIC{$n \in \mathbb{N} \to \overline{1}(n) = {{1}^{n}}$}
\UIC{$\forall n \in \mathbb{N}\,(\overline{1}(n) = {{1}^{n}})$}
\AXC{$\stackrel{\Pi_{11}}{\vdots}$}\noLine
\UIC{$\mathrm{Seq}({(1)^{\bullet}})$}
\AXC{$\stackrel{\Pi_{12}}{\vdots}$}\noLine
\UIC{$\mathrm{Seq}({(1)^{\bullet}}) \to \varphi_{4}$}
\BIC{$\varphi_{4}$}
\BIC{$\mathrm{Conv}(\overline{1},1) \to \mathrm{Conv}({(1)^{\bullet}},1)$}
\end{ptree}
\noindent 主木:
\begin{ptree}
\AXC{$\stackrel{\Pi_{1}}{\vdots}$}\noLine
\UIC{$\mathrm{Conv}(\overline{1},1)$}
\AXC{$\stackrel{\Pi_{13}}{\vdots}$}\noLine
\UIC{$\mathrm{Conv}(\overline{1},1) \to \mathrm{Conv}({(1)^{\bullet}},1)$}
\BIC{$\mathrm{Conv}({(1)^{\bullet}},1)$}
\end{ptree}
\begin{abbrevlist}
\item $\varphi_{1} :\equiv (1 \in \mathbb{R} \to \mathrm{Conv}(\overline{1},1)) \to 1 \in \mathbb{R} \to \mathrm{Conv}(\overline{1},1)$
\item $\varphi_{2} :\equiv \forall c \in \mathbb{R}\,(\mathrm{Conv}(\overline{c},c)) \to 1 \in \mathbb{R} \to \mathrm{Conv}(\overline{1},1)$
\item $\varphi_{3} :\equiv \forall x\,\forall y\,\forall z\,(x = y \to y = z \to x = z) \to \forall y\,\forall z\,(\overline{1}(n) = y \to y = z \to \overline{1}(n) = z)$
\item $\varphi_{4} :\equiv \forall n \in \mathbb{N}\,(\overline{1}(n) = {{1}^{n}}) \to \mathrm{Conv}(\overline{1},1) \to \mathrm{Conv}({(1)^{\bullet}},1)$
\item $\varphi_{5} :\equiv \forall a\,\forall b\,\forall \alpha\,(\mathrm{Seq}(b) \to \forall n \in \mathbb{N}\,(a_{n} = b_{n}) \to \mathrm{Conv}(a,\alpha) \to \mathrm{Conv}(b,\alpha))$
\item $\varphi_{6} :\equiv \mathrm{Seq}({(1)^{\bullet}}) \land {{1}^{1_{\omega}}} = 1 \land \forall n \in \mathbb{N}\,({{1}^{S(n)}} = {{1}^{n}} \cdot 1)$
\item $\varphi_{7} :\equiv \forall b\,\forall \alpha\,(\mathrm{Seq}(b) \to \forall n \in \mathbb{N}\,(\overline{1}(n) = b_{n}) \to \mathrm{Conv}(\overline{1},\alpha) \to \mathrm{Conv}(b,\alpha))$
\item $\varphi_{8} :\equiv (n \in \mathbb{N} \to {{1}^{n}} = 1) \to ({{1}^{n}} = 1 \to 1 = {{1}^{n}}) \to n \in \mathbb{N} \to 1 = {{1}^{n}}$
\item $\varphi_{9} :\equiv \mathrm{Seq}(\overline{1}) \land \forall n \in \mathbb{N}\,(\overline{1}(n) = 1) \to n \in \mathbb{N} \to \forall n \in \mathbb{N}\,(\overline{1}(n) = 1)$
\item $\varphi_{10} :\equiv \forall y\,\forall z\,(\overline{1}(n) = y \to y = z \to \overline{1}(n) = z) \to \forall z\,(\overline{1}(n) = 1 \to 1 = z \to \overline{1}(n) = z)$
\item $\varphi_{11} :\equiv (n \in \mathbb{N} \to n \in \mathbb{N} \to \overline{1}(n) = 1) \to (n \in \mathbb{N} \to n \in \mathbb{N}) \to n \in \mathbb{N} \to \overline{1}(n) = 1$
\item $\varphi_{12} :\equiv \forall \alpha\,(\mathrm{Seq}({(1)^{\bullet}}) \to \forall n \in \mathbb{N}\,(\overline{1}(n) = {{1}^{n}}) \to \mathrm{Conv}(\overline{1},\alpha) \to \mathrm{Conv}({(1)^{\bullet}},\alpha))$
\item $\varphi_{13} :\equiv \forall r \in \mathbb{R}\,(\mathrm{Seq}({(r)^{\bullet}}) \land {{r}^{1_{\omega}}} = r \land \forall n \in \mathbb{N}\,({{r}^{S(n)}} = {{r}^{n}} \cdot r))$
\item $\varphi_{14} :\equiv (\forall n \in \mathbb{N}\,(\overline{1}(n) = 1) \to n \in \mathbb{N} \to \overline{1}(n) = 1) \to n \in \mathbb{N} \to n \in \mathbb{N} \to \overline{1}(n) = 1$
\item $\varphi_{15} :\equiv (\overline{1}(n) = 1 \to 1 = {{1}^{n}} \to \overline{1}(n) = {{1}^{n}}) \to n \in \mathbb{N} \to \overline{1}(n) = {{1}^{n}}$
\item $\varphi_{16} :\equiv \forall z\,(\overline{1}(n) = 1 \to 1 = z \to \overline{1}(n) = z) \to \overline{1}(n) = 1 \to 1 = {{1}^{n}} \to \overline{1}(n) = {{1}^{n}}$
\item $\varphi_{17} :\equiv (1 \in \mathbb{R} \to \mathrm{Seq}(\overline{1}) \land \forall n \in \mathbb{N}\,(\overline{1}(n) = 1)) \to 1 \in \mathbb{R} \to \mathrm{Seq}(\overline{1}) \land \forall n \in \mathbb{N}\,(\overline{1}(n) = 1)$
\item $\varphi_{18} :\equiv \forall c \in \mathbb{R}\,(\mathrm{Seq}(\overline{c}) \land \forall n \in \mathbb{N}\,(\overline{c}(n) = c)) \to 1 \in \mathbb{R} \to \mathrm{Seq}(\overline{1}) \land \forall n \in \mathbb{N}\,(\overline{1}(n) = 1)$
\item $\varphi_{19} :\equiv (\mathrm{Seq}({(1)^{\bullet}}) \to \varphi_{4}) \to \mathrm{Seq}({(1)^{\bullet}}) \to \varphi_{4}$
\end{abbrevlist}
\end{fullproof}

% 自動生成 (prover/build.js): GeoDvg  ノード数 308
\begin{fullproof}{命題3.1(1)}{$\forall r \in \mathbb{R}\,(1 < r \to \mathrm{DvgP}({(r)^{\bullet}}))$}{308}\phantomsection\label{tree:GeoDvg}
\noindent$\Pi_{1}$:
\begin{ptree}
\AXC{$y = z \to x = y \to x = z$\lbl{A6}}
\AXC{$(y = z \to x = y \to x = z) \to x = y \to y = z \to x = z$\lbl{A1}}
\BIC{$x = y \to y = z \to x = z$}
\UIC{$\forall z\,(x = y \to y = z \to x = z)$}
\UIC{$\forall y\,\forall z\,(x = y \to y = z \to x = z)$}
\UIC{$\forall x\,\forall y\,\forall z\,(x = y \to y = z \to x = z)$}
\AXC{$\forall x\,\forall y\,\forall z\,(x = y \to y = z \to x = z) \to \varphi_{13}$\lbl{A2}}
\BIC{$\varphi_{13}$}
\end{ptree}
\noindent$\Pi_{2}$:
\begin{ptree}
\AXC{$x = x$\lbl{A5}}
\AXC{$x = y \to x = x \to y = x$\lbl{A6}}
\AXC{$(x = y \to x = x \to y = x) \to x = x \to x = y \to y = x$\lbl{A1}}
\BIC{$x = x \to x = y \to y = x$}
\BIC{$x = y \to y = x$}
\UIC{$\forall y\,(x = y \to y = x)$}
\UIC{$\forall x\,\forall y\,(x = y \to y = x)$}
\end{ptree}
\noindent$\Pi_{3}$:
\begin{ptree}
\AXC{$\stackrel{\Pi_{2}}{\vdots}$}\noLine
\UIC{$\forall x\,\forall y\,(x = y \to y = x)$}
\AXC{$\forall x\,\forall y\,(x = y \to y = x) \to \forall y\,(r = y \to y = r)$\lbl{A2}}
\BIC{$\forall y\,(r = y \to y = r)$}
\AXC{$\forall y\,(r = y \to y = r) \to r = 0 \to 0 = r$\lbl{A2}}
\BIC{$r = 0 \to 0 = r$}
\end{ptree}
\noindent$\Pi_{4}$:
\begin{ptree}
\AXC{$\forall x\,\forall y\,\forall z\,(x < y \to y < z \to x < z)$\lbl{[O2]}}
\AXC{$\varphi_{9}$\lbl{A2}}
\BIC{$\forall y\,\forall z\,(0 < y \to y < z \to 0 < z)$}
\AXC{$\varphi_{4}$\lbl{A2}}
\BIC{$\forall z\,(0 < 1 \to 1 < z \to 0 < z)$}
\AXC{$\forall z\,(0 < 1 \to 1 < z \to 0 < z) \to 0 < 1 \to 1 < r \to 0 < r$\lbl{A2}}
\BIC{$0 < 1 \to 1 < r \to 0 < r$}
\end{ptree}
\noindent$\Pi_{5}$:
\begin{ptree}
\AXC{$r \in \mathbb{R} \to 1 < r \to 1 < r$\lbl{A1}}
\AXC{$0 < 1$\lbl{[F28]}}
\AXC{$\stackrel{\Pi_{4}}{\vdots}$}\noLine
\UIC{$0 < 1 \to 1 < r \to 0 < r$}
\AXC{$(0 < 1 \to 1 < r \to 0 < r) \to \varphi_{16}$\lbl{A1}}
\BIC{$\varphi_{16}$}
\BIC{$(r \in \mathbb{R} \to 1 < r \to 1 < r) \to r \in \mathbb{R} \to 1 < r \to 0 < r$}
\BIC{$r \in \mathbb{R} \to 1 < r \to 0 < r$}
\end{ptree}
\noindent$\Pi_{6}$:
\begin{ptree}
\AXC{$\forall x\,\forall y\,(x < y \to x \neq y)$\lbl{[F26]}}
\AXC{$\forall x\,\forall y\,(x < y \to x \neq y) \to \forall y\,(0 < y \to 0 \neq y)$\lbl{A2}}
\BIC{$\forall y\,(0 < y \to 0 \neq y)$}
\AXC{$\forall y\,(0 < y \to 0 \neq y) \to 0 < r \to 0 \neq r$\lbl{A2}}
\BIC{$0 < r \to 0 \neq r$}
\end{ptree}
\noindent$\Pi_{7}$:
\begin{ptree}
\AXC{$\stackrel{\Pi_{3}}{\vdots}$}\noLine
\UIC{$r = 0 \to 0 = r$}
\AXC{$\stackrel{\Pi_{5}}{\vdots}$}\noLine
\UIC{$r \in \mathbb{R} \to 1 < r \to 0 < r$}
\AXC{$\stackrel{\Pi_{6}}{\vdots}$}\noLine
\UIC{$0 < r \to 0 \neq r$}
\AXC{$\varphi_{26}$\lbl{A1}}
\BIC{$(r \in \mathbb{R} \to 1 < r \to 0 < r) \to r \in \mathbb{R} \to 1 < r \to 0 \neq r$}
\BIC{$r \in \mathbb{R} \to 1 < r \to 0 \neq r$}
\AXC{$\varphi_{14}$\lbl{A1}}
\BIC{$(r = 0 \to 0 = r) \to r \in \mathbb{R} \to 1 < r \to r \neq 0$}
\BIC{$r \in \mathbb{R} \to 1 < r \to r \neq 0$}
\end{ptree}
\noindent$\Pi_{8}$:
\begin{ptree}
\AXC{$r \in \mathbb{R} \to 1 < r \to n \in \mathbb{N} \to r \in \mathbb{R}$\lbl{A1}}
\AXC{$\forall r \in \mathbb{R}\,(r \neq 0 \to \forall n \in \mathbb{N}\,({{r}^{n}} \neq 0))$\lbl{[P4]}}
\AXC{$\varphi_{64}$\lbl{A2}}
\BIC{$r \in \mathbb{R} \to r \neq 0 \to \forall n \in \mathbb{N}\,({{r}^{n}} \neq 0)$}
\AXC{$\varphi_{92}$\lbl{A1}}
\BIC{$(r \in \mathbb{R} \to 1 < r \to n \in \mathbb{N} \to r \in \mathbb{R}) \to \varphi_{45}$}
\BIC{$\varphi_{45}$}
\end{ptree}
\noindent$\Pi_{9}$:
\begin{ptree}
\AXC{$\stackrel{\Pi_{7}}{\vdots}$}\noLine
\UIC{$r \in \mathbb{R} \to 1 < r \to r \neq 0$}
\AXC{$\stackrel{\Pi_{8}}{\vdots}$}\noLine
\UIC{$\varphi_{45}$}
\BIC{$r \in \mathbb{R} \to 1 < r \to n \in \mathbb{N} \to \forall n \in \mathbb{N}\,({{r}^{n}} \neq 0)$}
\AXC{$(r \in \mathbb{R} \to 1 < r \to n \in \mathbb{N} \to \forall n \in \mathbb{N}\,({{r}^{n}} \neq 0)) \to \varphi_{73}$\lbl{A1}}
\BIC{$\varphi_{73}$}
\end{ptree}
\noindent$\Pi_{10}$:
\begin{ptree}
\AXC{$\forall n \in \mathbb{N}\,({{r}^{n}} \neq 0) \to n \in \mathbb{N} \to {{r}^{n}} \neq 0$\lbl{A2}}
\AXC{$\stackrel{\Pi_{9}}{\vdots}$}\noLine
\UIC{$\varphi_{73}$}
\BIC{$r \in \mathbb{R} \to 1 < r \to n \in \mathbb{N} \to n \in \mathbb{N} \to {{r}^{n}} \neq 0$}
\AXC{$(r \in \mathbb{R} \to 1 < r \to n \in \mathbb{N} \to n \in \mathbb{N} \to {{r}^{n}} \neq 0) \to \varphi_{47}$\lbl{A1}}
\BIC{$\varphi_{47}$}
\end{ptree}
\noindent$\Pi_{11}$:
\begin{ptree}
\AXC{$r \in \mathbb{R} \to 1 < r \to n \in \mathbb{N} \to r \in \mathbb{R}$\lbl{A1}}
\AXC{$\forall r \in \mathbb{R}\,(\varphi_{36})$\lbl{[S24]}}
\AXC{$\forall r \in \mathbb{R}\,(\varphi_{36}) \to r \in \mathbb{R} \to \varphi_{36}$\lbl{A2}}
\BIC{$r \in \mathbb{R} \to \varphi_{36}$}
\AXC{$(r \in \mathbb{R} \to \varphi_{36}) \to \varphi_{85}$\lbl{A1}}
\BIC{$\varphi_{85}$}
\BIC{$r \in \mathbb{R} \to 1 < r \to n \in \mathbb{N} \to \varphi_{36}$}
\end{ptree}
\noindent$\Pi_{12}$:
\begin{ptree}
\AXC{$\forall a\,\forall n\,(\mathrm{Seq}(a) \to n \in \mathbb{N} \to a_{n} \in \mathbb{R})$\lbl{[S1]}}
\AXC{$\varphi_{50}$\lbl{A2}}
\BIC{$\forall n\,(\mathrm{Seq}({(r)^{\bullet}}) \to n \in \mathbb{N} \to {{r}^{n}} \in \mathbb{R})$}
\AXC{$\varphi_{52}$\lbl{A2}}
\BIC{$\mathrm{Seq}({(r)^{\bullet}}) \to n \in \mathbb{N} \to {{r}^{n}} \in \mathbb{R}$}
\end{ptree}
\noindent$\Pi_{13}$:
\begin{ptree}
\AXC{$\stackrel{\Pi_{12}}{\vdots}$}\noLine
\UIC{$\mathrm{Seq}({(r)^{\bullet}}) \to n \in \mathbb{N} \to {{r}^{n}} \in \mathbb{R}$}
\AXC{$(\mathrm{Seq}({(r)^{\bullet}}) \to n \in \mathbb{N} \to {{r}^{n}} \in \mathbb{R}) \to \varphi_{83}$\lbl{A1}}
\BIC{$\varphi_{83}$}
\end{ptree}
\noindent$\Pi_{14}$:
\begin{ptree}
\AXC{$r \in \mathbb{R} \to 1 < r \to n \in \mathbb{N} \to n \in \mathbb{N}$\lbl{A1}}
\AXC{$\stackrel{\Pi_{11}}{\vdots}$}\noLine
\UIC{$r \in \mathbb{R} \to 1 < r \to n \in \mathbb{N} \to \varphi_{36}$}
\AXC{$\varphi_{86}$\lbl{A1}}
\BIC{$r \in \mathbb{R} \to 1 < r \to n \in \mathbb{N} \to \mathrm{Seq}({(r)^{\bullet}})$}
\AXC{$\stackrel{\Pi_{13}}{\vdots}$}\noLine
\UIC{$\varphi_{83}$}
\BIC{$\varphi_{56}$}
\BIC{$r \in \mathbb{R} \to 1 < r \to n \in \mathbb{N} \to {{r}^{n}} \in \mathbb{R}$}
\end{ptree}
\noindent$\Pi_{15}$:
\begin{ptree}
\AXC{$\stackrel{\Pi_{14}}{\vdots}$}\noLine
\UIC{$r \in \mathbb{R} \to 1 < r \to n \in \mathbb{N} \to {{r}^{n}} \in \mathbb{R}$}
\AXC{$\forall x \in \mathbb{R}\,(x \neq 0 \to (x^{-1})^{-1} = x)$\lbl{[E30]}}
\AXC{$\varphi_{70}$\lbl{A2}}
\BIC{${{r}^{n}} \in \mathbb{R} \to {{r}^{n}} \neq 0 \to ({{r}^{n}}^{-1})^{-1} = {{r}^{n}}$}
\AXC{$\varphi_{103}$\lbl{A1}}
\BIC{$(r \in \mathbb{R} \to 1 < r \to n \in \mathbb{N} \to {{r}^{n}} \in \mathbb{R}) \to \varphi_{72}$}
\BIC{$\varphi_{72}$}
\end{ptree}
\noindent$\Pi_{16}$:
\begin{ptree}
\AXC{$y = z \to x = y \to x = z$\lbl{A6}}
\AXC{$(y = z \to x = y \to x = z) \to x = y \to y = z \to x = z$\lbl{A1}}
\BIC{$x = y \to y = z \to x = z$}
\UIC{$\forall z\,(x = y \to y = z \to x = z)$}
\UIC{$\forall y\,\forall z\,(x = y \to y = z \to x = z)$}
\UIC{$\forall x\,\forall y\,\forall z\,(x = y \to y = z \to x = z)$}
\AXC{$\forall x\,\forall y\,\forall z\,(x = y \to y = z \to x = z) \to \varphi_{13}$\lbl{A2}}
\BIC{$\varphi_{13}$}
\end{ptree}
\noindent$\Pi_{17}$:
\begin{ptree}
\AXC{$z = w \to {{(r^{-1})}^{n}}^{-1} = z^{-1} \to {{(r^{-1})}^{n}}^{-1} = w^{-1}$\lbl{A6}}
\UIC{$\varphi_{20}$}
\UIC{$\forall z\,(\varphi_{20})$}
\AXC{$\forall z\,(\varphi_{20}) \to \varphi_{60}$\lbl{A2}}
\BIC{$\varphi_{60}$}
\AXC{$\varphi_{60} \to {{(r^{-1})}^{n}} = {{r}^{n}}^{-1} \to \varphi_{46}$\lbl{A2}}
\BIC{${{(r^{-1})}^{n}} = {{r}^{n}}^{-1} \to \varphi_{46}$}
\end{ptree}
\noindent$\Pi_{18}$:
\begin{ptree}
\AXC{$r \in \mathbb{R} \to 1 < r \to n \in \mathbb{N} \to r \in \mathbb{R}$\lbl{A1}}
\AXC{$\forall r \in \mathbb{R}\,(r \neq 0 \to \forall n \in \mathbb{N}\,({{(r^{-1})}^{n}} = {{r}^{n}}^{-1}))$\lbl{[P6]}}
\AXC{$\varphi_{82}$\lbl{A2}}
\BIC{$r \in \mathbb{R} \to r \neq 0 \to \forall n \in \mathbb{N}\,({{(r^{-1})}^{n}} = {{r}^{n}}^{-1})$}
\AXC{$\varphi_{100}$\lbl{A1}}
\BIC{$\varphi_{84}$}
\BIC{$(r \in \mathbb{R} \to 1 < r \to r \neq 0) \to \varphi_{38}$}
\end{ptree}
\noindent$\Pi_{19}$:
\begin{ptree}
\AXC{$\varphi_{55}$\lbl{A2}}
\AXC{$\stackrel{\Pi_{7}}{\vdots}$}\noLine
\UIC{$r \in \mathbb{R} \to 1 < r \to r \neq 0$}
\AXC{$\stackrel{\Pi_{18}}{\vdots}$}\noLine
\UIC{$(r \in \mathbb{R} \to 1 < r \to r \neq 0) \to \varphi_{38}$}
\BIC{$\varphi_{38}$}
\AXC{$\varphi_{38} \to \varphi_{90}$\lbl{A1}}
\BIC{$\varphi_{90}$}
\BIC{$r \in \mathbb{R} \to 1 < r \to n \in \mathbb{N} \to n \in \mathbb{N} \to {{(r^{-1})}^{n}} = {{r}^{n}}^{-1}$}
\AXC{$\varphi_{94}$\lbl{A1}}
\BIC{$\varphi_{69}$}
\end{ptree}
\noindent$\Pi_{20}$:
\begin{ptree}
\AXC{$r \in \mathbb{R} \to 1 < r \to n \in \mathbb{N} \to n \in \mathbb{N}$\lbl{A1}}
\AXC{$\stackrel{\Pi_{19}}{\vdots}$}\noLine
\UIC{$\varphi_{69}$}
\BIC{$r \in \mathbb{R} \to 1 < r \to n \in \mathbb{N} \to {{(r^{-1})}^{n}} = {{r}^{n}}^{-1}$}
\AXC{$(r \in \mathbb{R} \to 1 < r \to n \in \mathbb{N} \to {{(r^{-1})}^{n}} = {{r}^{n}}^{-1}) \to \varphi_{104}$\lbl{A1}}
\BIC{$\varphi_{104}$}
\end{ptree}
\noindent$\Pi_{21}$:
\begin{ptree}
\AXC{$r \in \mathbb{R} \to 1 < r \to r \in \mathbb{R}$\lbl{A1}}
\AXC{$\forall x \in \mathbb{R}\,(0 < x \to 0 < x^{-1})$\lbl{[F34]}}
\AXC{$\varphi_{25}$\lbl{A2}}
\BIC{$r \in \mathbb{R} \to 0 < r \to 0 < r^{-1}$}
\AXC{$(r \in \mathbb{R} \to 0 < r \to 0 < r^{-1}) \to \varphi_{54}$\lbl{A1}}
\BIC{$\varphi_{54}$}
\BIC{$(r \in \mathbb{R} \to 1 < r \to 0 < r) \to r \in \mathbb{R} \to 1 < r \to 0 < r^{-1}$}
\end{ptree}
\noindent$\Pi_{22}$:
\begin{ptree}
\AXC{$\stackrel{\Pi_{5}}{\vdots}$}\noLine
\UIC{$r \in \mathbb{R} \to 1 < r \to 0 < r$}
\AXC{$\stackrel{\Pi_{21}}{\vdots}$}\noLine
\UIC{$(r \in \mathbb{R} \to 1 < r \to 0 < r) \to r \in \mathbb{R} \to 1 < r \to 0 < r^{-1}$}
\BIC{$r \in \mathbb{R} \to 1 < r \to 0 < r^{-1}$}
\end{ptree}
\noindent$\Pi_{23}$:
\begin{ptree}
\AXC{$r \in \mathbb{R} \to 1 < r \to r \in \mathbb{R}$\lbl{A1}}
\AXC{$\forall x \in \mathbb{R}\,(x \neq 0 \to x^{-1} \in \mathbb{R})$\lbl{(R1)}}
\AXC{$\varphi_{31}$\lbl{A2}}
\BIC{$r \in \mathbb{R} \to r \neq 0 \to r^{-1} \in \mathbb{R}$}
\AXC{$(r \in \mathbb{R} \to r \neq 0 \to r^{-1} \in \mathbb{R}) \to \varphi_{57}$\lbl{A1}}
\BIC{$\varphi_{57}$}
\BIC{$(r \in \mathbb{R} \to 1 < r \to r \neq 0) \to r \in \mathbb{R} \to 1 < r \to r^{-1} \in \mathbb{R}$}
\end{ptree}
\noindent$\Pi_{24}$:
\begin{ptree}
\AXC{$\stackrel{\Pi_{7}}{\vdots}$}\noLine
\UIC{$r \in \mathbb{R} \to 1 < r \to r \neq 0$}
\AXC{$\stackrel{\Pi_{23}}{\vdots}$}\noLine
\UIC{$(r \in \mathbb{R} \to 1 < r \to r \neq 0) \to r \in \mathbb{R} \to 1 < r \to r^{-1} \in \mathbb{R}$}
\BIC{$r \in \mathbb{R} \to 1 < r \to r^{-1} \in \mathbb{R}$}
\end{ptree}
\noindent$\Pi_{25}$:
\begin{ptree}
\AXC{$\stackrel{\Pi_{24}}{\vdots}$}\noLine
\UIC{$r \in \mathbb{R} \to 1 < r \to r^{-1} \in \mathbb{R}$}
\AXC{$\forall r \in \mathbb{R}\,(0 < r \to \forall n \in \mathbb{N}\,(0 < {{r}^{n}}))$\lbl{[P3]}}
\AXC{$\varphi_{77}$\lbl{A2}}
\BIC{$r^{-1} \in \mathbb{R} \to 0 < r^{-1} \to \forall n \in \mathbb{N}\,(0 < {{(r^{-1})}^{n}})$}
\AXC{$\varphi_{96}$\lbl{A1}}
\BIC{$(r \in \mathbb{R} \to 1 < r \to r^{-1} \in \mathbb{R}) \to \varphi_{48}$}
\BIC{$\varphi_{48}$}
\end{ptree}
\noindent$\Pi_{26}$:
\begin{ptree}
\AXC{$\stackrel{\Pi_{22}}{\vdots}$}\noLine
\UIC{$r \in \mathbb{R} \to 1 < r \to 0 < r^{-1}$}
\AXC{$\stackrel{\Pi_{25}}{\vdots}$}\noLine
\UIC{$\varphi_{48}$}
\BIC{$r \in \mathbb{R} \to 1 < r \to \forall n \in \mathbb{N}\,(0 < {{(r^{-1})}^{n}})$}
\end{ptree}
\noindent$\Pi_{27}$:
\begin{ptree}
\AXC{$\stackrel{\Pi_{26}}{\vdots}$}\noLine
\UIC{$r \in \mathbb{R} \to 1 < r \to \forall n \in \mathbb{N}\,(0 < {{(r^{-1})}^{n}})$}
\AXC{$(r \in \mathbb{R} \to 1 < r \to \forall n \in \mathbb{N}\,(0 < {{(r^{-1})}^{n}})) \to \varphi_{79}$\lbl{A1}}
\BIC{$\varphi_{79}$}
\end{ptree}
\noindent$\Pi_{28}$:
\begin{ptree}
\AXC{$\forall x\,\forall y\,(x < y \to x \neq y)$\lbl{[F26]}}
\AXC{$\forall x\,\forall y\,(x < y \to x \neq y) \to \forall y\,(0 < y \to 0 \neq y)$\lbl{A2}}
\BIC{$\forall y\,(0 < y \to 0 \neq y)$}
\end{ptree}
\noindent$\Pi_{29}$:
\begin{ptree}
\AXC{$\stackrel{\Pi_{28}}{\vdots}$}\noLine
\UIC{$\forall y\,(0 < y \to 0 \neq y)$}
\AXC{$\forall y\,(0 < y \to 0 \neq y) \to 0 < {{(r^{-1})}^{n}} \to 0 \neq {{(r^{-1})}^{n}}$\lbl{A2}}
\BIC{$0 < {{(r^{-1})}^{n}} \to 0 \neq {{(r^{-1})}^{n}}$}
\AXC{$(0 < {{(r^{-1})}^{n}} \to 0 \neq {{(r^{-1})}^{n}}) \to \varphi_{68}$\lbl{A1}}
\BIC{$\varphi_{68}$}
\end{ptree}
\noindent$\Pi_{30}$:
\begin{ptree}
\AXC{$\forall n \in \mathbb{N}\,(0 < {{(r^{-1})}^{n}}) \to n \in \mathbb{N} \to 0 < {{(r^{-1})}^{n}}$\lbl{A2}}
\AXC{$\stackrel{\Pi_{27}}{\vdots}$}\noLine
\UIC{$\varphi_{79}$}
\BIC{$r \in \mathbb{R} \to 1 < r \to n \in \mathbb{N} \to 0 < {{(r^{-1})}^{n}}$}
\AXC{$\stackrel{\Pi_{29}}{\vdots}$}\noLine
\UIC{$\varphi_{68}$}
\BIC{$r \in \mathbb{R} \to 1 < r \to n \in \mathbb{N} \to 0 \neq {{(r^{-1})}^{n}}$}
\end{ptree}
\noindent$\Pi_{31}$:
\begin{ptree}
\AXC{$\stackrel{\Pi_{30}}{\vdots}$}\noLine
\UIC{$r \in \mathbb{R} \to 1 < r \to n \in \mathbb{N} \to 0 \neq {{(r^{-1})}^{n}}$}
\AXC{$(r \in \mathbb{R} \to 1 < r \to n \in \mathbb{N} \to 0 \neq {{(r^{-1})}^{n}}) \to \varphi_{53}$\lbl{A1}}
\BIC{$\varphi_{53}$}
\end{ptree}
\noindent$\Pi_{32}$:
\begin{ptree}
\AXC{$\stackrel{\Pi_{2}}{\vdots}$}\noLine
\UIC{$\forall x\,\forall y\,(x = y \to y = x)$}
\AXC{$\varphi_{6}$\lbl{A2}}
\BIC{$\forall y\,({{(r^{-1})}^{n}} = y \to y = {{(r^{-1})}^{n}})$}
\AXC{$\varphi_{40}$\lbl{A2}}
\BIC{${{(r^{-1})}^{n}} = 0 \to 0 = {{(r^{-1})}^{n}}$}
\AXC{$\stackrel{\Pi_{31}}{\vdots}$}\noLine
\UIC{$\varphi_{53}$}
\BIC{$r \in \mathbb{R} \to 1 < r \to n \in \mathbb{N} \to {{(r^{-1})}^{n}} \neq 0$}
\AXC{$\varphi_{66}$\lbl{A1}}
\BIC{$r \in \mathbb{R} \land 1 < r \to n \in \mathbb{N} \to {{(r^{-1})}^{n}} \neq 0$}
\UIC{$\forall n\,(r \in \mathbb{R} \land 1 < r \to n \in \mathbb{N} \to {{(r^{-1})}^{n}} \neq 0)$}
\end{ptree}
\noindent$\Pi_{33}$:
\begin{ptree}
\AXC{$\varphi_{71}$\lbl{A3}}
\AXC{$\stackrel{\Pi_{32}}{\vdots}$}\noLine
\UIC{$\forall n\,(r \in \mathbb{R} \land 1 < r \to n \in \mathbb{N} \to {{(r^{-1})}^{n}} \neq 0)$}
\AXC{$\forall n\,(r \in \mathbb{R} \land 1 < r \to n \in \mathbb{N} \to {{(r^{-1})}^{n}} \neq 0) \to \varphi_{71} \to \varphi_{61}$\lbl{A1}}
\BIC{$\varphi_{71} \to \varphi_{61}$}
\BIC{$\varphi_{61}$}
\end{ptree}
\noindent$\Pi_{34}$:
\begin{ptree}
\AXC{$\stackrel{\Pi_{24}}{\vdots}$}\noLine
\UIC{$r \in \mathbb{R} \to 1 < r \to r^{-1} \in \mathbb{R}$}
\AXC{$\forall r \in \mathbb{R}\,(\varphi_{36})$\lbl{[S24]}}
\AXC{$\forall r \in \mathbb{R}\,(\varphi_{36}) \to r^{-1} \in \mathbb{R} \to \varphi_{67}$\lbl{A2}}
\BIC{$r^{-1} \in \mathbb{R} \to \varphi_{67}$}
\AXC{$\varphi_{106}$\lbl{A1}}
\BIC{$(r \in \mathbb{R} \to 1 < r \to r^{-1} \in \mathbb{R}) \to r \in \mathbb{R} \to 1 < r \to \varphi_{67}$}
\BIC{$r \in \mathbb{R} \to 1 < r \to \varphi_{67}$}
\AXC{$\varphi_{89}$\lbl{A1}}
\BIC{$r \in \mathbb{R} \to 1 < r \to \mathrm{Seq}({(r^{-1})^{\bullet}})$}
\end{ptree}
\noindent$\Pi_{35}$:
\begin{ptree}
\AXC{$\stackrel{\Pi_{34}}{\vdots}$}\noLine
\UIC{$r \in \mathbb{R} \to 1 < r \to \mathrm{Seq}({(r^{-1})^{\bullet}})$}
\AXC{$\varphi_{58}$\lbl{[S5]}}
\AXC{$\varphi_{58} \to \mathrm{Seq}({(r^{-1})^{\bullet}}) \to \varphi_{75}$\lbl{A2}}
\BIC{$\mathrm{Seq}({(r^{-1})^{\bullet}}) \to \varphi_{75}$}
\AXC{$\varphi_{107}$\lbl{A1}}
\BIC{$(r \in \mathbb{R} \to 1 < r \to \mathrm{Seq}({(r^{-1})^{\bullet}})) \to \varphi_{87}$}
\BIC{$\varphi_{87}$}
\end{ptree}
\noindent$\Pi_{36}$:
\begin{ptree}
\AXC{$r \in \mathbb{R} \land 1 < r \to \forall n\,(r \in \mathbb{R} \land 1 < r)$\lbl{A4}}
\AXC{$\stackrel{\Pi_{33}}{\vdots}$}\noLine
\UIC{$\varphi_{61}$}
\BIC{$r \in \mathbb{R} \to 1 < r \to \forall n \in \mathbb{N}\,({{(r^{-1})}^{n}} \neq 0)$}
\AXC{$\stackrel{\Pi_{35}}{\vdots}$}\noLine
\UIC{$\varphi_{87}$}
\BIC{$r \in \mathbb{R} \to 1 < r \to \mathrm{Seq}({(r^{-1})^{\bullet}}^{\ast}) \land \varphi_{7}$}
\AXC{$\varphi_{97}$\lbl{A1}}
\BIC{$r \in \mathbb{R} \to 1 < r \to n \in \mathbb{N} \to \varphi_{7}$}
\end{ptree}
\noindent$\Pi_{37}$:
\begin{ptree}
\AXC{$\varphi_{74}$\lbl{A2}}
\AXC{$\stackrel{\Pi_{36}}{\vdots}$}\noLine
\UIC{$r \in \mathbb{R} \to 1 < r \to n \in \mathbb{N} \to \varphi_{7}$}
\AXC{$(r \in \mathbb{R} \to 1 < r \to n \in \mathbb{N} \to \varphi_{7}) \to \varphi_{74} \to r \in \mathbb{R} \to 1 < r \to \varphi_{24}$\lbl{A1}}
\BIC{$\varphi_{74} \to r \in \mathbb{R} \to 1 < r \to \varphi_{24}$}
\BIC{$r \in \mathbb{R} \to 1 < r \to \varphi_{24}$}
\AXC{$\varphi_{98}$\lbl{A1}}
\BIC{$(r \in \mathbb{R} \to 1 < r \to n \in \mathbb{N} \to n \in \mathbb{N}) \to r \in \mathbb{R} \to \varphi_{17}$}
\end{ptree}
\noindent$\Pi_{38}$:
\begin{ptree}
\AXC{$r \in \mathbb{R} \to 1 < r \to n \in \mathbb{N} \to n \in \mathbb{N}$\lbl{A1}}
\AXC{$\stackrel{\Pi_{37}}{\vdots}$}\noLine
\UIC{$(r \in \mathbb{R} \to 1 < r \to n \in \mathbb{N} \to n \in \mathbb{N}) \to r \in \mathbb{R} \to \varphi_{17}$}
\BIC{$r \in \mathbb{R} \to \varphi_{17}$}
\AXC{$(r \in \mathbb{R} \to \varphi_{17}) \to \varphi_{34} \to \varphi_{99}$\lbl{A1}}
\BIC{$\varphi_{34} \to \varphi_{99}$}
\end{ptree}
\noindent$\Pi_{39}$:
\begin{ptree}
\AXC{$\stackrel{\Pi_{17}}{\vdots}$}\noLine
\UIC{${{(r^{-1})}^{n}} = {{r}^{n}}^{-1} \to \varphi_{46}$}
\AXC{$x = x$\lbl{A5}}
\UIC{$\forall x\,(x = x)$}
\AXC{$\forall x\,(x = x) \to {{(r^{-1})}^{n}}^{-1} = {{(r^{-1})}^{n}}^{-1}$\lbl{A2}}
\BIC{${{(r^{-1})}^{n}}^{-1} = {{(r^{-1})}^{n}}^{-1}$}
\AXC{$\stackrel{\Pi_{20}}{\vdots}$}\noLine
\UIC{$\varphi_{104}$}
\BIC{$({{(r^{-1})}^{n}} = {{r}^{n}}^{-1} \to \varphi_{46}) \to \varphi_{34}$}
\BIC{$\varphi_{34}$}
\AXC{$\stackrel{\Pi_{38}}{\vdots}$}\noLine
\UIC{$\varphi_{34} \to \varphi_{99}$}
\BIC{$\varphi_{99}$}
\end{ptree}
\noindent$\Pi_{40}$:
\begin{ptree}
\AXC{$\stackrel{\Pi_{16}}{\vdots}$}\noLine
\UIC{$\varphi_{13}$}
\AXC{$\varphi_{13} \to \varphi_{65}$\lbl{A2}}
\BIC{$\varphi_{65}$}
\AXC{$\varphi_{65} \to {(r^{-1})^{\bullet}}^{\ast}(n) = {{(r^{-1})}^{n}}^{-1} \to \varphi_{49}$\lbl{A2}}
\BIC{${(r^{-1})^{\bullet}}^{\ast}(n) = {{(r^{-1})}^{n}}^{-1} \to \varphi_{49}$}
\AXC{$\stackrel{\Pi_{39}}{\vdots}$}\noLine
\UIC{$\varphi_{99}$}
\BIC{$r \in \mathbb{R} \to \varphi_{18}$}
\AXC{$\varphi_{108}$\lbl{A1}}
\BIC{$(r \in \mathbb{R} \to 1 < r \to n \in \mathbb{N} \to ({{r}^{n}}^{-1})^{-1} = {{r}^{n}}) \to \varphi_{95}$}
\end{ptree}
\noindent$\Pi_{41}$:
\begin{ptree}
\AXC{$r \in \mathbb{R} \to 1 < r \to n \in \mathbb{N} \to n \in \mathbb{N}$\lbl{A1}}
\AXC{$\stackrel{\Pi_{10}}{\vdots}$}\noLine
\UIC{$\varphi_{47}$}
\BIC{$r \in \mathbb{R} \to 1 < r \to n \in \mathbb{N} \to {{r}^{n}} \neq 0$}
\AXC{$\stackrel{\Pi_{15}}{\vdots}$}\noLine
\UIC{$\varphi_{72}$}
\BIC{$r \in \mathbb{R} \to 1 < r \to n \in \mathbb{N} \to ({{r}^{n}}^{-1})^{-1} = {{r}^{n}}$}
\AXC{$\stackrel{\Pi_{40}}{\vdots}$}\noLine
\UIC{$(r \in \mathbb{R} \to 1 < r \to n \in \mathbb{N} \to ({{r}^{n}}^{-1})^{-1} = {{r}^{n}}) \to \varphi_{95}$}
\BIC{$\varphi_{95}$}
\end{ptree}
\noindent$\Pi_{42}$:
\begin{ptree}
\AXC{$\stackrel{\Pi_{1}}{\vdots}$}\noLine
\UIC{$\varphi_{13}$}
\AXC{$\varphi_{13} \to \varphi_{63}$\lbl{A2}}
\BIC{$\varphi_{63}$}
\AXC{$\varphi_{63} \to {(r^{-1})^{\bullet}}^{\ast}(n) = ({{r}^{n}}^{-1})^{-1} \to \varphi_{30}$\lbl{A2}}
\BIC{${(r^{-1})^{\bullet}}^{\ast}(n) = ({{r}^{n}}^{-1})^{-1} \to \varphi_{30}$}
\AXC{$\stackrel{\Pi_{41}}{\vdots}$}\noLine
\UIC{$\varphi_{95}$}
\BIC{$\varphi_{21}$}
\AXC{$\varphi_{21} \to \varphi_{22}$\lbl{A1}}
\BIC{$\varphi_{22}$}
\end{ptree}
\noindent$\Pi_{43}$:
\begin{ptree}
\AXC{$r \in \mathbb{R} \land 1 < r \to \forall n\,(r \in \mathbb{R} \land 1 < r)$\lbl{A4}}
\AXC{$\forall n\,(\varphi_{23}) \to \varphi_{33}$\lbl{A3}}
\AXC{$\stackrel{\Pi_{42}}{\vdots}$}\noLine
\UIC{$\varphi_{22}$}
\AXC{$\varphi_{22} \to \varphi_{23}$\lbl{A1}}
\BIC{$\varphi_{23}$}
\UIC{$\forall n\,(\varphi_{23})$}
\AXC{$\forall n\,(\varphi_{23}) \to \varphi_{105}$\lbl{A1}}
\BIC{$\varphi_{105}$}
\BIC{$(r \in \mathbb{R} \land 1 < r \to \forall n\,(r \in \mathbb{R} \land 1 < r)) \to \varphi_{28}$}
\BIC{$\varphi_{28}$}
\end{ptree}
\noindent$\Pi_{44}$:
\begin{ptree}
\AXC{$r \in \mathbb{R} \to 1 < r \to r \in \mathbb{R}$\lbl{A1}}
\AXC{$\forall r \in \mathbb{R}\,(\varphi_{36})$\lbl{[S24]}}
\AXC{$\forall r \in \mathbb{R}\,(\varphi_{36}) \to r \in \mathbb{R} \to \varphi_{36}$\lbl{A2}}
\BIC{$r \in \mathbb{R} \to \varphi_{36}$}
\AXC{$\varphi_{101}$\lbl{A1}}
\BIC{$(r \in \mathbb{R} \to 1 < r \to r \in \mathbb{R}) \to r \in \mathbb{R} \to 1 < r \to \varphi_{36}$}
\BIC{$r \in \mathbb{R} \to 1 < r \to \varphi_{36}$}
\AXC{$\varphi_{80}$\lbl{A1}}
\BIC{$r \in \mathbb{R} \to 1 < r \to \mathrm{Seq}({(r)^{\bullet}})$}
\end{ptree}
\noindent$\Pi_{45}$:
\begin{ptree}
\AXC{$r \in \mathbb{R} \to 1 < r \to r \in \mathbb{R}$\lbl{A1}}
\AXC{$1 \in \mathbb{R}$\lbl{(R1)}}
\AXC{$\varphi_{10}$\lbl{[E19]}}
\AXC{$\varphi_{78}$\lbl{A2}}
\BIC{$\forall y\,(1 \in \mathbb{R} \to y \in \mathbb{R} \to 0 < 1 \to 1 < y \to y^{-1} < 1^{-1})$}
\AXC{$\varphi_{76}$\lbl{A2}}
\BIC{$1 \in \mathbb{R} \to r \in \mathbb{R} \to 0 < 1 \to 1 < r \to r^{-1} < 1^{-1}$}
\AXC{$\varphi_{91}$\lbl{A1}}
\BIC{$1 \in \mathbb{R} \to (r \in \mathbb{R} \to 1 < r \to r \in \mathbb{R}) \to 0 < 1 \to \varphi_{19}$}
\BIC{$(r \in \mathbb{R} \to 1 < r \to r \in \mathbb{R}) \to 0 < 1 \to \varphi_{19}$}
\BIC{$0 < 1 \to \varphi_{19}$}
\end{ptree}
\noindent$\Pi_{46}$:
\begin{ptree}
\AXC{$1^{-1} = 1$\lbl{[E35]}}
\AXC{$1^{-1} = 1 \to (r \in \mathbb{R} \to 1 < r \to r^{-1} < 1^{-1}) \to \varphi_{32}$\lbl{A1}}
\BIC{$(r \in \mathbb{R} \to 1 < r \to r^{-1} < 1^{-1}) \to \varphi_{32}$}
\end{ptree}
\noindent$\Pi_{47}$:
\begin{ptree}
\AXC{$r \in \mathbb{R} \to 1 < r \to 1 < r$\lbl{A1}}
\AXC{$0 < 1$\lbl{[F28]}}
\AXC{$\stackrel{\Pi_{45}}{\vdots}$}\noLine
\UIC{$0 < 1 \to \varphi_{19}$}
\BIC{$\varphi_{19}$}
\BIC{$r \in \mathbb{R} \to 1 < r \to r^{-1} < 1^{-1}$}
\AXC{$\stackrel{\Pi_{46}}{\vdots}$}\noLine
\UIC{$(r \in \mathbb{R} \to 1 < r \to r^{-1} < 1^{-1}) \to \varphi_{32}$}
\BIC{$\varphi_{32}$}
\end{ptree}
\noindent$\Pi_{48}$:
\begin{ptree}
\AXC{$\forall x\,(0 < x \to 0 \leq x)$\lbl{[E10]}}
\AXC{$\forall x\,(0 < x \to 0 \leq x) \to 0 < r^{-1} \to 0 \leq r^{-1}$\lbl{A2}}
\BIC{$0 < r^{-1} \to 0 \leq r^{-1}$}
\AXC{$(0 < r^{-1} \to 0 \leq r^{-1}) \to \varphi_{11}$\lbl{A1}}
\BIC{$\varphi_{11}$}
\end{ptree}
\noindent$\Pi_{49}$:
\begin{ptree}
\AXC{$\forall x\,(0 \leq x \to \lvert x \rvert = x)$\lbl{[A1]}}
\AXC{$\forall x\,(0 \leq x \to \lvert x \rvert = x) \to 0 \leq r^{-1} \to \lvert r^{-1} \rvert = r^{-1}$\lbl{A2}}
\BIC{$0 \leq r^{-1} \to \lvert r^{-1} \rvert = r^{-1}$}
\end{ptree}
\noindent$\Pi_{50}$:
\begin{ptree}
\AXC{$\stackrel{\Pi_{22}}{\vdots}$}\noLine
\UIC{$r \in \mathbb{R} \to 1 < r \to 0 < r^{-1}$}
\AXC{$\stackrel{\Pi_{48}}{\vdots}$}\noLine
\UIC{$\varphi_{11}$}
\BIC{$r \in \mathbb{R} \to 1 < r \to 0 \leq r^{-1}$}
\AXC{$\stackrel{\Pi_{49}}{\vdots}$}\noLine
\UIC{$0 \leq r^{-1} \to \lvert r^{-1} \rvert = r^{-1}$}
\AXC{$(0 \leq r^{-1} \to \lvert r^{-1} \rvert = r^{-1}) \to \varphi_{27}$\lbl{A1}}
\BIC{$\varphi_{27}$}
\BIC{$r \in \mathbb{R} \to 1 < r \to \lvert r^{-1} \rvert = r^{-1}$}
\end{ptree}
\noindent$\Pi_{51}$:
\begin{ptree}
\AXC{$\stackrel{\Pi_{50}}{\vdots}$}\noLine
\UIC{$r \in \mathbb{R} \to 1 < r \to \lvert r^{-1} \rvert = r^{-1}$}
\AXC{$(r \in \mathbb{R} \to 1 < r \to \lvert r^{-1} \rvert = r^{-1}) \to \varphi_{37}$\lbl{A1}}
\BIC{$\varphi_{37}$}
\end{ptree}
\noindent$\Pi_{52}$:
\begin{ptree}
\AXC{$\stackrel{\Pi_{2}}{\vdots}$}\noLine
\UIC{$\forall x\,\forall y\,(x = y \to y = x)$}
\AXC{$\varphi_{2}$\lbl{A2}}
\BIC{$\forall y\,(\lvert r^{-1} \rvert = y \to y = \lvert r^{-1} \rvert)$}
\AXC{$\varphi_{29}$\lbl{A2}}
\BIC{$\lvert r^{-1} \rvert = r^{-1} \to r^{-1} = \lvert r^{-1} \rvert$}
\AXC{$\stackrel{\Pi_{51}}{\vdots}$}\noLine
\UIC{$\varphi_{37}$}
\BIC{$r \in \mathbb{R} \to 1 < r \to r^{-1} = \lvert r^{-1} \rvert$}
\AXC{$\varphi_{88}$\lbl{A1}}
\BIC{$(r \in \mathbb{R} \to 1 < r \to r^{-1} < 1) \to \varphi_{44}$}
\end{ptree}
\noindent$\Pi_{53}$:
\begin{ptree}
\AXC{$z = w \to r^{-1} < z \to r^{-1} < w$\lbl{A6}}
\UIC{$\forall w\,(z = w \to r^{-1} < z \to r^{-1} < w)$}
\UIC{$\forall z\,\forall w\,(z = w \to r^{-1} < z \to r^{-1} < w)$}
\AXC{$\varphi_{39}$\lbl{A2}}
\BIC{$\forall w\,(1^{-1} = w \to r^{-1} < 1^{-1} \to r^{-1} < w)$}
\AXC{$\varphi_{42}$\lbl{A2}}
\BIC{$1^{-1} = 1 \to r^{-1} < 1^{-1} \to r^{-1} < 1$}
\AXC{$\stackrel{\Pi_{47}}{\vdots}$}\noLine
\UIC{$\varphi_{32}$}
\BIC{$r \in \mathbb{R} \to 1 < r \to r^{-1} < 1$}
\AXC{$\stackrel{\Pi_{52}}{\vdots}$}\noLine
\UIC{$(r \in \mathbb{R} \to 1 < r \to r^{-1} < 1) \to \varphi_{44}$}
\BIC{$\varphi_{44}$}
\end{ptree}
\noindent$\Pi_{54}$:
\begin{ptree}
\AXC{$z = w \to z < 1 \to w < 1$\lbl{A6}}
\UIC{$\forall w\,(z = w \to z < 1 \to w < 1)$}
\UIC{$\forall z\,\forall w\,(z = w \to z < 1 \to w < 1)$}
\AXC{$\varphi_{8}$\lbl{A2}}
\BIC{$\forall w\,(r^{-1} = w \to r^{-1} < 1 \to w < 1)$}
\AXC{$\varphi_{41}$\lbl{A2}}
\BIC{$r^{-1} = \lvert r^{-1} \rvert \to r^{-1} < 1 \to \lvert r^{-1} \rvert < 1$}
\AXC{$\stackrel{\Pi_{53}}{\vdots}$}\noLine
\UIC{$\varphi_{44}$}
\BIC{$r \in \mathbb{R} \to 1 < r \to \lvert r^{-1} \rvert < 1$}
\end{ptree}
\noindent$\Pi_{55}$:
\begin{ptree}
\AXC{$\forall r \in \mathbb{R}\,(\lvert r \rvert < 1 \to \mathrm{Conv}({(r)^{\bullet}},0))$\lbl{命題3.1(3)}}
\AXC{$\varphi_{59}$\lbl{A2}}
\BIC{$\varphi_{3}$}
\AXC{$\varphi_{3} \to (r \in \mathbb{R} \to 1 < r \to r^{-1} \in \mathbb{R}) \to \varphi_{35}$\lbl{A1}}
\BIC{$(r \in \mathbb{R} \to 1 < r \to r^{-1} \in \mathbb{R}) \to \varphi_{35}$}
\end{ptree}
\noindent$\Pi_{56}$:
\begin{ptree}
\AXC{$\stackrel{\Pi_{26}}{\vdots}$}\noLine
\UIC{$r \in \mathbb{R} \to 1 < r \to \forall n \in \mathbb{N}\,(0 < {{(r^{-1})}^{n}})$}
\AXC{$\varphi_{5}$\lbl{定理2.17}}
\AXC{$\varphi_{5} \to \forall n \in \mathbb{N}\,(0 < {{(r^{-1})}^{n}}) \to \varphi_{1}$\lbl{A2}}
\BIC{$\forall n \in \mathbb{N}\,(0 < {{(r^{-1})}^{n}}) \to \varphi_{1}$}
\AXC{$\varphi_{102}$\lbl{A1}}
\BIC{$(r \in \mathbb{R} \to 1 < r \to \forall n \in \mathbb{N}\,(0 < {{(r^{-1})}^{n}})) \to \varphi_{43}$}
\BIC{$\varphi_{43}$}
\end{ptree}
\noindent$\Pi_{57}$:
\begin{ptree}
\AXC{$\stackrel{\Pi_{54}}{\vdots}$}\noLine
\UIC{$r \in \mathbb{R} \to 1 < r \to \lvert r^{-1} \rvert < 1$}
\AXC{$\stackrel{\Pi_{24}}{\vdots}$}\noLine
\UIC{$r \in \mathbb{R} \to 1 < r \to r^{-1} \in \mathbb{R}$}
\AXC{$\stackrel{\Pi_{55}}{\vdots}$}\noLine
\UIC{$(r \in \mathbb{R} \to 1 < r \to r^{-1} \in \mathbb{R}) \to \varphi_{35}$}
\BIC{$\varphi_{35}$}
\BIC{$r \in \mathbb{R} \to 1 < r \to \mathrm{Conv}({(r^{-1})^{\bullet}},0)$}
\AXC{$\stackrel{\Pi_{56}}{\vdots}$}\noLine
\UIC{$\varphi_{43}$}
\BIC{$r \in \mathbb{R} \to 1 < r \to \mathrm{DvgP}({(r^{-1})^{\bullet}}^{\ast})$}
\end{ptree}
\noindent$\Pi_{58}$:
\begin{ptree}
\AXC{$\stackrel{\Pi_{57}}{\vdots}$}\noLine
\UIC{$r \in \mathbb{R} \to 1 < r \to \mathrm{DvgP}({(r^{-1})^{\bullet}}^{\ast})$}
\AXC{$\varphi_{15}$\lbl{定理2.18(1)}}
\AXC{$\varphi_{15} \to \varphi_{51}$\lbl{A2}}
\BIC{$\varphi_{51}$}
\AXC{$\varphi_{51} \to \varphi_{62}$\lbl{A2}}
\BIC{$\varphi_{62}$}
\AXC{$\varphi_{62} \to \varphi_{93}$\lbl{A1}}
\BIC{$\varphi_{93}$}
\BIC{$\varphi_{81}$}
\end{ptree}
\noindent 主木:
\begin{ptree}
\AXC{$\stackrel{\Pi_{43}}{\vdots}$}\noLine
\UIC{$\varphi_{28}$}
\AXC{$\stackrel{\Pi_{44}}{\vdots}$}\noLine
\UIC{$r \in \mathbb{R} \to 1 < r \to \mathrm{Seq}({(r)^{\bullet}})$}
\AXC{$\stackrel{\Pi_{58}}{\vdots}$}\noLine
\UIC{$\varphi_{81}$}
\BIC{$\varphi_{28} \to r \in \mathbb{R} \to 1 < r \to \mathrm{DvgP}({(r)^{\bullet}})$}
\BIC{$r \in \mathbb{R} \to 1 < r \to \mathrm{DvgP}({(r)^{\bullet}})$}
\UIC{$\forall r \in \mathbb{R}\,(1 < r \to \mathrm{DvgP}({(r)^{\bullet}}))$}
\end{ptree}
\begin{abbrevlist}
\item $\varphi_{1} :\equiv \mathrm{Conv}({(r^{-1})^{\bullet}},0) \to \mathrm{DvgP}({(r^{-1})^{\bullet}}^{\ast})$
\item $\varphi_{2} :\equiv \forall x\,\forall y\,(x = y \to y = x) \to \forall y\,(\lvert r^{-1} \rvert = y \to y = \lvert r^{-1} \rvert)$
\item $\varphi_{3} :\equiv r^{-1} \in \mathbb{R} \to \lvert r^{-1} \rvert < 1 \to \mathrm{Conv}({(r^{-1})^{\bullet}},0)$
\item $\varphi_{4} :\equiv \forall y\,\forall z\,(0 < y \to y < z \to 0 < z) \to \forall z\,(0 < 1 \to 1 < z \to 0 < z)$
\item $\varphi_{5} :\equiv \forall a\,(\forall n \in \mathbb{N}\,(0 < a_{n}) \to \mathrm{Conv}(a,0) \to \mathrm{DvgP}(a^{\ast}))$
\item $\varphi_{6} :\equiv \forall x\,\forall y\,(x = y \to y = x) \to \forall y\,({{(r^{-1})}^{n}} = y \to y = {{(r^{-1})}^{n}})$
\item $\varphi_{7} :\equiv \forall n \in \mathbb{N}\,({(r^{-1})^{\bullet}}^{\ast}(n) = {{(r^{-1})}^{n}}^{-1})$
\item $\varphi_{8} :\equiv \forall z\,\forall w\,(z = w \to z < 1 \to w < 1) \to \forall w\,(r^{-1} = w \to r^{-1} < 1 \to w < 1)$
\item $\varphi_{9} :\equiv \forall x\,\forall y\,\forall z\,(x < y \to y < z \to x < z) \to \forall y\,\forall z\,(0 < y \to y < z \to 0 < z)$
\item $\varphi_{10} :\equiv \forall x\,\forall y\,(x \in \mathbb{R} \to y \in \mathbb{R} \to 0 < x \to x < y \to y^{-1} < x^{-1})$
\item $\varphi_{11} :\equiv (r \in \mathbb{R} \to 1 < r \to 0 < r^{-1}) \to r \in \mathbb{R} \to 1 < r \to 0 \leq r^{-1}$
\item $\varphi_{12} :\equiv \forall n \in \mathbb{N}\,({(r^{-1})^{\bullet}}^{\ast}(n) \leq {{r}^{n}}) \to \mathrm{DvgP}({(r)^{\bullet}})$
\item $\varphi_{13} :\equiv \forall y\,\forall z\,({(r^{-1})^{\bullet}}^{\ast}(n) = y \to y = z \to {(r^{-1})^{\bullet}}^{\ast}(n) = z)$
\item $\varphi_{14} :\equiv (r \in \mathbb{R} \to 1 < r \to 0 \neq r) \to (r = 0 \to 0 = r) \to r \in \mathbb{R} \to 1 < r \to r \neq 0$
\item $\varphi_{15} :\equiv \forall a\,\forall b\,(\mathrm{DvgP}(a) \to \mathrm{Seq}(b) \to \forall n \in \mathbb{N}\,(a_{n} \leq b_{n}) \to \mathrm{DvgP}(b))$
\item $\varphi_{16} :\equiv 0 < 1 \to (r \in \mathbb{R} \to 1 < r \to 1 < r) \to r \in \mathbb{R} \to 1 < r \to 0 < r$
\item $\varphi_{17} :\equiv 1 < r \to n \in \mathbb{N} \to {(r^{-1})^{\bullet}}^{\ast}(n) = {{(r^{-1})}^{n}}^{-1}$
\item $\varphi_{18} :\equiv 1 < r \to n \in \mathbb{N} \to {(r^{-1})^{\bullet}}^{\ast}(n) = ({{r}^{n}}^{-1})^{-1}$
\item $\varphi_{19} :\equiv (r \in \mathbb{R} \to 1 < r \to 1 < r) \to r \in \mathbb{R} \to 1 < r \to r^{-1} < 1^{-1}$
\item $\varphi_{20} :\equiv \forall w\,(z = w \to {{(r^{-1})}^{n}}^{-1} = z^{-1} \to {{(r^{-1})}^{n}}^{-1} = w^{-1})$
\item $\varphi_{21} :\equiv r \in \mathbb{R} \to 1 < r \to n \in \mathbb{N} \to {(r^{-1})^{\bullet}}^{\ast}(n) = {{r}^{n}}$
\item $\varphi_{22} :\equiv r \in \mathbb{R} \to 1 < r \to n \in \mathbb{N} \to {(r^{-1})^{\bullet}}^{\ast}(n) \leq {{r}^{n}}$
\item $\varphi_{23} :\equiv r \in \mathbb{R} \land 1 < r \to n \in \mathbb{N} \to {(r^{-1})^{\bullet}}^{\ast}(n) \leq {{r}^{n}}$
\item $\varphi_{24} :\equiv n \in \mathbb{N} \to n \in \mathbb{N} \to {(r^{-1})^{\bullet}}^{\ast}(n) = {{(r^{-1})}^{n}}^{-1}$
\item $\varphi_{25} :\equiv \forall x \in \mathbb{R}\,(0 < x \to 0 < x^{-1}) \to r \in \mathbb{R} \to 0 < r \to 0 < r^{-1}$
\item $\varphi_{26} :\equiv (0 < r \to 0 \neq r) \to (r \in \mathbb{R} \to 1 < r \to 0 < r) \to r \in \mathbb{R} \to 1 < r \to 0 \neq r$
\item $\varphi_{27} :\equiv (r \in \mathbb{R} \to 1 < r \to 0 \leq r^{-1}) \to r \in \mathbb{R} \to 1 < r \to \lvert r^{-1} \rvert = r^{-1}$
\item $\varphi_{28} :\equiv r \in \mathbb{R} \to 1 < r \to \forall n \in \mathbb{N}\,({(r^{-1})^{\bullet}}^{\ast}(n) \leq {{r}^{n}})$
\item $\varphi_{29} :\equiv \forall y\,(\lvert r^{-1} \rvert = y \to y = \lvert r^{-1} \rvert) \to \lvert r^{-1} \rvert = r^{-1} \to r^{-1} = \lvert r^{-1} \rvert$
\item $\varphi_{30} :\equiv ({{r}^{n}}^{-1})^{-1} = {{r}^{n}} \to {(r^{-1})^{\bullet}}^{\ast}(n) = {{r}^{n}}$
\item $\varphi_{31} :\equiv \forall x \in \mathbb{R}\,(x \neq 0 \to x^{-1} \in \mathbb{R}) \to r \in \mathbb{R} \to r \neq 0 \to r^{-1} \in \mathbb{R}$
\item $\varphi_{32} :\equiv (1^{-1} = 1 \to r^{-1} < 1^{-1} \to r^{-1} < 1) \to r \in \mathbb{R} \to 1 < r \to r^{-1} < 1$
\item $\varphi_{33} :\equiv \forall n\,(r \in \mathbb{R} \land 1 < r) \to \forall n \in \mathbb{N}\,({(r^{-1})^{\bullet}}^{\ast}(n) \leq {{r}^{n}})$
\item $\varphi_{34} :\equiv r \in \mathbb{R} \to 1 < r \to n \in \mathbb{N} \to {{(r^{-1})}^{n}}^{-1} = ({{r}^{n}}^{-1})^{-1}$
\item $\varphi_{35} :\equiv (r \in \mathbb{R} \to 1 < r \to \lvert r^{-1} \rvert < 1) \to r \in \mathbb{R} \to 1 < r \to \mathrm{Conv}({(r^{-1})^{\bullet}},0)$
\item $\varphi_{36} :\equiv \mathrm{Seq}({(r)^{\bullet}}) \land {{r}^{1_{\omega}}} = r \land \forall n \in \mathbb{N}\,({{r}^{S(n)}} = {{r}^{n}} \cdot r)$
\item $\varphi_{37} :\equiv (\lvert r^{-1} \rvert = r^{-1} \to r^{-1} = \lvert r^{-1} \rvert) \to r \in \mathbb{R} \to 1 < r \to r^{-1} = \lvert r^{-1} \rvert$
\item $\varphi_{38} :\equiv r \in \mathbb{R} \to 1 < r \to n \in \mathbb{N} \to \forall n \in \mathbb{N}\,({{(r^{-1})}^{n}} = {{r}^{n}}^{-1})$
\item $\varphi_{39} :\equiv \forall z\,\forall w\,(z = w \to r^{-1} < z \to r^{-1} < w) \to \forall w\,(1^{-1} = w \to r^{-1} < 1^{-1} \to r^{-1} < w)$
\item $\varphi_{40} :\equiv \forall y\,({{(r^{-1})}^{n}} = y \to y = {{(r^{-1})}^{n}}) \to {{(r^{-1})}^{n}} = 0 \to 0 = {{(r^{-1})}^{n}}$
\item $\varphi_{41} :\equiv \forall w\,(r^{-1} = w \to r^{-1} < 1 \to w < 1) \to r^{-1} = \lvert r^{-1} \rvert \to r^{-1} < 1 \to \lvert r^{-1} \rvert < 1$
\item $\varphi_{42} :\equiv \forall w\,(1^{-1} = w \to r^{-1} < 1^{-1} \to r^{-1} < w) \to 1^{-1} = 1 \to r^{-1} < 1^{-1} \to r^{-1} < 1$
\item $\varphi_{43} :\equiv (r \in \mathbb{R} \to 1 < r \to \mathrm{Conv}({(r^{-1})^{\bullet}},0)) \to r \in \mathbb{R} \to 1 < r \to \mathrm{DvgP}({(r^{-1})^{\bullet}}^{\ast})$
\item $\varphi_{44} :\equiv (r^{-1} = \lvert r^{-1} \rvert \to r^{-1} < 1 \to \lvert r^{-1} \rvert < 1) \to r \in \mathbb{R} \to 1 < r \to \lvert r^{-1} \rvert < 1$
\item $\varphi_{45} :\equiv (r \in \mathbb{R} \to 1 < r \to r \neq 0) \to r \in \mathbb{R} \to 1 < r \to n \in \mathbb{N} \to \forall n \in \mathbb{N}\,({{r}^{n}} \neq 0)$
\item $\varphi_{46} :\equiv {{(r^{-1})}^{n}}^{-1} = {{(r^{-1})}^{n}}^{-1} \to {{(r^{-1})}^{n}}^{-1} = ({{r}^{n}}^{-1})^{-1}$
\item $\varphi_{47} :\equiv (r \in \mathbb{R} \to 1 < r \to n \in \mathbb{N} \to n \in \mathbb{N}) \to r \in \mathbb{R} \to 1 < r \to n \in \mathbb{N} \to {{r}^{n}} \neq 0$
\item $\varphi_{48} :\equiv (r \in \mathbb{R} \to 1 < r \to 0 < r^{-1}) \to r \in \mathbb{R} \to 1 < r \to \forall n \in \mathbb{N}\,(0 < {{(r^{-1})}^{n}})$
\item $\varphi_{49} :\equiv {{(r^{-1})}^{n}}^{-1} = ({{r}^{n}}^{-1})^{-1} \to {(r^{-1})^{\bullet}}^{\ast}(n) = ({{r}^{n}}^{-1})^{-1}$
\item $\varphi_{50} :\equiv \forall a\,\forall n\,(\mathrm{Seq}(a) \to n \in \mathbb{N} \to a_{n} \in \mathbb{R}) \to \forall n\,(\mathrm{Seq}({(r)^{\bullet}}) \to n \in \mathbb{N} \to {{r}^{n}} \in \mathbb{R})$
\item $\varphi_{51} :\equiv \forall b\,(\mathrm{DvgP}({(r^{-1})^{\bullet}}^{\ast}) \to \mathrm{Seq}(b) \to \forall n \in \mathbb{N}\,({(r^{-1})^{\bullet}}^{\ast}(n) \leq b_{n}) \to \mathrm{DvgP}(b))$
\item $\varphi_{52} :\equiv \forall n\,(\mathrm{Seq}({(r)^{\bullet}}) \to n \in \mathbb{N} \to {{r}^{n}} \in \mathbb{R}) \to \mathrm{Seq}({(r)^{\bullet}}) \to n \in \mathbb{N} \to {{r}^{n}} \in \mathbb{R}$
\item $\varphi_{53} :\equiv ({{(r^{-1})}^{n}} = 0 \to 0 = {{(r^{-1})}^{n}}) \to r \in \mathbb{R} \to 1 < r \to n \in \mathbb{N} \to {{(r^{-1})}^{n}} \neq 0$
\item $\varphi_{54} :\equiv (r \in \mathbb{R} \to 1 < r \to r \in \mathbb{R}) \to (r \in \mathbb{R} \to 1 < r \to 0 < r) \to r \in \mathbb{R} \to 1 < r \to 0 < r^{-1}$
\item $\varphi_{55} :\equiv \forall n \in \mathbb{N}\,({{(r^{-1})}^{n}} = {{r}^{n}}^{-1}) \to n \in \mathbb{N} \to {{(r^{-1})}^{n}} = {{r}^{n}}^{-1}$
\item $\varphi_{56} :\equiv (r \in \mathbb{R} \to 1 < r \to n \in \mathbb{N} \to n \in \mathbb{N}) \to r \in \mathbb{R} \to 1 < r \to n \in \mathbb{N} \to {{r}^{n}} \in \mathbb{R}$
\item $\varphi_{57} :\equiv (r \in \mathbb{R} \to 1 < r \to r \in \mathbb{R}) \to (r \in \mathbb{R} \to 1 < r \to r \neq 0) \to r \in \mathbb{R} \to 1 < r \to r^{-1} \in \mathbb{R}$
\item $\varphi_{58} :\equiv \forall a\,(\mathrm{Seq}(a) \to \forall n \in \mathbb{N}\,(a_{n} \neq 0) \to \mathrm{Seq}(a^{\ast}) \land \forall n \in \mathbb{N}\,(a^{\ast}(n) = a_{n}^{-1}))$
\item $\varphi_{59} :\equiv \forall r \in \mathbb{R}\,(\lvert r \rvert < 1 \to \mathrm{Conv}({(r)^{\bullet}},0)) \to \varphi_{3}$
\item $\varphi_{60} :\equiv \forall w\,({{(r^{-1})}^{n}} = w \to {{(r^{-1})}^{n}}^{-1} = {{(r^{-1})}^{n}}^{-1} \to {{(r^{-1})}^{n}}^{-1} = w^{-1})$
\item $\varphi_{61} :\equiv (r \in \mathbb{R} \land 1 < r \to \forall n\,(r \in \mathbb{R} \land 1 < r)) \to r \in \mathbb{R} \to 1 < r \to \forall n \in \mathbb{N}\,({{(r^{-1})}^{n}} \neq 0)$
\item $\varphi_{62} :\equiv \mathrm{DvgP}({(r^{-1})^{\bullet}}^{\ast}) \to \mathrm{Seq}({(r)^{\bullet}}) \to \varphi_{12}$
\item $\varphi_{63} :\equiv \forall z\,({(r^{-1})^{\bullet}}^{\ast}(n) = ({{r}^{n}}^{-1})^{-1} \to ({{r}^{n}}^{-1})^{-1} = z \to {(r^{-1})^{\bullet}}^{\ast}(n) = z)$
\item $\varphi_{64} :\equiv \forall r \in \mathbb{R}\,(r \neq 0 \to \forall n \in \mathbb{N}\,({{r}^{n}} \neq 0)) \to r \in \mathbb{R} \to r \neq 0 \to \forall n \in \mathbb{N}\,({{r}^{n}} \neq 0)$
\item $\varphi_{65} :\equiv \forall z\,({(r^{-1})^{\bullet}}^{\ast}(n) = {{(r^{-1})}^{n}}^{-1} \to {{(r^{-1})}^{n}}^{-1} = z \to {(r^{-1})^{\bullet}}^{\ast}(n) = z)$
\item $\varphi_{66} :\equiv (r \in \mathbb{R} \to 1 < r \to n \in \mathbb{N} \to {{(r^{-1})}^{n}} \neq 0) \to r \in \mathbb{R} \land 1 < r \to n \in \mathbb{N} \to {{(r^{-1})}^{n}} \neq 0$
\item $\varphi_{67} :\equiv \mathrm{Seq}({(r^{-1})^{\bullet}}) \land {{(r^{-1})}^{1_{\omega}}} = r^{-1} \land \forall n \in \mathbb{N}\,({{(r^{-1})}^{S(n)}} = {{(r^{-1})}^{n}} \cdot r^{-1})$
\item $\varphi_{68} :\equiv (r \in \mathbb{R} \to 1 < r \to n \in \mathbb{N} \to 0 < {{(r^{-1})}^{n}}) \to r \in \mathbb{R} \to 1 < r \to n \in \mathbb{N} \to 0 \neq {{(r^{-1})}^{n}}$
\item $\varphi_{69} :\equiv (r \in \mathbb{R} \to 1 < r \to n \in \mathbb{N} \to n \in \mathbb{N}) \to r \in \mathbb{R} \to 1 < r \to n \in \mathbb{N} \to {{(r^{-1})}^{n}} = {{r}^{n}}^{-1}$
\item $\varphi_{70} :\equiv \forall x \in \mathbb{R}\,(x \neq 0 \to (x^{-1})^{-1} = x) \to {{r}^{n}} \in \mathbb{R} \to {{r}^{n}} \neq 0 \to ({{r}^{n}}^{-1})^{-1} = {{r}^{n}}$
\item $\varphi_{71} :\equiv \forall n\,(r \in \mathbb{R} \land 1 < r \to n \in \mathbb{N} \to {{(r^{-1})}^{n}} \neq 0) \to \forall n\,(r \in \mathbb{R} \land 1 < r) \to \forall n \in \mathbb{N}\,({{(r^{-1})}^{n}} \neq 0)$
\item $\varphi_{72} :\equiv (r \in \mathbb{R} \to 1 < r \to n \in \mathbb{N} \to {{r}^{n}} \neq 0) \to r \in \mathbb{R} \to 1 < r \to n \in \mathbb{N} \to ({{r}^{n}}^{-1})^{-1} = {{r}^{n}}$
\item $\varphi_{73} :\equiv (\forall n \in \mathbb{N}\,({{r}^{n}} \neq 0) \to n \in \mathbb{N} \to {{r}^{n}} \neq 0) \to r \in \mathbb{R} \to 1 < r \to n \in \mathbb{N} \to n \in \mathbb{N} \to {{r}^{n}} \neq 0$
\item $\varphi_{74} :\equiv \varphi_{7} \to n \in \mathbb{N} \to {(r^{-1})^{\bullet}}^{\ast}(n) = {{(r^{-1})}^{n}}^{-1}$
\item $\varphi_{75} :\equiv \forall n \in \mathbb{N}\,({{(r^{-1})}^{n}} \neq 0) \to \mathrm{Seq}({(r^{-1})^{\bullet}}^{\ast}) \land \varphi_{7}$
\item $\varphi_{76} :\equiv \forall y\,(1 \in \mathbb{R} \to y \in \mathbb{R} \to 0 < 1 \to 1 < y \to y^{-1} < 1^{-1}) \to 1 \in \mathbb{R} \to r \in \mathbb{R} \to 0 < 1 \to 1 < r \to r^{-1} < 1^{-1}$
\item $\varphi_{77} :\equiv \forall r \in \mathbb{R}\,(0 < r \to \forall n \in \mathbb{N}\,(0 < {{r}^{n}})) \to r^{-1} \in \mathbb{R} \to 0 < r^{-1} \to \forall n \in \mathbb{N}\,(0 < {{(r^{-1})}^{n}})$
\item $\varphi_{78} :\equiv \varphi_{10} \to \forall y\,(1 \in \mathbb{R} \to y \in \mathbb{R} \to 0 < 1 \to 1 < y \to y^{-1} < 1^{-1})$
\item $\varphi_{79} :\equiv (\forall n \in \mathbb{N}\,(0 < {{(r^{-1})}^{n}}) \to n \in \mathbb{N} \to 0 < {{(r^{-1})}^{n}}) \to r \in \mathbb{R} \to 1 < r \to n \in \mathbb{N} \to 0 < {{(r^{-1})}^{n}}$
\item $\varphi_{80} :\equiv (r \in \mathbb{R} \to 1 < r \to \varphi_{36}) \to r \in \mathbb{R} \to 1 < r \to \mathrm{Seq}({(r)^{\bullet}})$
\item $\varphi_{81} :\equiv (r \in \mathbb{R} \to 1 < r \to \mathrm{Seq}({(r)^{\bullet}})) \to \varphi_{28} \to r \in \mathbb{R} \to 1 < r \to \mathrm{DvgP}({(r)^{\bullet}})$
\item $\varphi_{82} :\equiv \forall r \in \mathbb{R}\,(r \neq 0 \to \forall n \in \mathbb{N}\,({{(r^{-1})}^{n}} = {{r}^{n}}^{-1})) \to r \in \mathbb{R} \to r \neq 0 \to \forall n \in \mathbb{N}\,({{(r^{-1})}^{n}} = {{r}^{n}}^{-1})$
\item $\varphi_{83} :\equiv (r \in \mathbb{R} \to 1 < r \to n \in \mathbb{N} \to \mathrm{Seq}({(r)^{\bullet}})) \to \varphi_{56}$
\item $\varphi_{84} :\equiv (r \in \mathbb{R} \to 1 < r \to n \in \mathbb{N} \to r \in \mathbb{R}) \to (r \in \mathbb{R} \to 1 < r \to r \neq 0) \to \varphi_{38}$
\item $\varphi_{85} :\equiv (r \in \mathbb{R} \to 1 < r \to n \in \mathbb{N} \to r \in \mathbb{R}) \to r \in \mathbb{R} \to 1 < r \to n \in \mathbb{N} \to \varphi_{36}$
\item $\varphi_{86} :\equiv (r \in \mathbb{R} \to 1 < r \to n \in \mathbb{N} \to \varphi_{36}) \to r \in \mathbb{R} \to 1 < r \to n \in \mathbb{N} \to \mathrm{Seq}({(r)^{\bullet}})$
\item $\varphi_{87} :\equiv (r \in \mathbb{R} \to 1 < r \to \forall n \in \mathbb{N}\,({{(r^{-1})}^{n}} \neq 0)) \to r \in \mathbb{R} \to 1 < r \to \mathrm{Seq}({(r^{-1})^{\bullet}}^{\ast}) \land \varphi_{7}$
\item $\varphi_{88} :\equiv (r \in \mathbb{R} \to 1 < r \to r^{-1} = \lvert r^{-1} \rvert) \to (r \in \mathbb{R} \to 1 < r \to r^{-1} < 1) \to \varphi_{44}$
\item $\varphi_{89} :\equiv (r \in \mathbb{R} \to 1 < r \to \varphi_{67}) \to r \in \mathbb{R} \to 1 < r \to \mathrm{Seq}({(r^{-1})^{\bullet}})$
\item $\varphi_{90} :\equiv \varphi_{55} \to r \in \mathbb{R} \to 1 < r \to n \in \mathbb{N} \to n \in \mathbb{N} \to {{(r^{-1})}^{n}} = {{r}^{n}}^{-1}$
\item $\varphi_{91} :\equiv (1 \in \mathbb{R} \to r \in \mathbb{R} \to 0 < 1 \to 1 < r \to r^{-1} < 1^{-1}) \to 1 \in \mathbb{R} \to (r \in \mathbb{R} \to 1 < r \to r \in \mathbb{R}) \to 0 < 1 \to \varphi_{19}$
\item $\varphi_{92} :\equiv (r \in \mathbb{R} \to r \neq 0 \to \forall n \in \mathbb{N}\,({{r}^{n}} \neq 0)) \to (r \in \mathbb{R} \to 1 < r \to n \in \mathbb{N} \to r \in \mathbb{R}) \to \varphi_{45}$
\item $\varphi_{93} :\equiv (r \in \mathbb{R} \to 1 < r \to \mathrm{DvgP}({(r^{-1})^{\bullet}}^{\ast})) \to \varphi_{81}$
\item $\varphi_{94} :\equiv (r \in \mathbb{R} \to 1 < r \to n \in \mathbb{N} \to n \in \mathbb{N} \to {{(r^{-1})}^{n}} = {{r}^{n}}^{-1}) \to \varphi_{69}$
\item $\varphi_{95} :\equiv ({(r^{-1})^{\bullet}}^{\ast}(n) = ({{r}^{n}}^{-1})^{-1} \to \varphi_{30}) \to \varphi_{21}$
\item $\varphi_{96} :\equiv (r^{-1} \in \mathbb{R} \to 0 < r^{-1} \to \forall n \in \mathbb{N}\,(0 < {{(r^{-1})}^{n}})) \to (r \in \mathbb{R} \to 1 < r \to r^{-1} \in \mathbb{R}) \to \varphi_{48}$
\item $\varphi_{97} :\equiv (r \in \mathbb{R} \to 1 < r \to \mathrm{Seq}({(r^{-1})^{\bullet}}^{\ast}) \land \varphi_{7}) \to r \in \mathbb{R} \to 1 < r \to n \in \mathbb{N} \to \varphi_{7}$
\item $\varphi_{98} :\equiv (r \in \mathbb{R} \to 1 < r \to \varphi_{24}) \to (r \in \mathbb{R} \to 1 < r \to n \in \mathbb{N} \to n \in \mathbb{N}) \to r \in \mathbb{R} \to \varphi_{17}$
\item $\varphi_{99} :\equiv ({(r^{-1})^{\bullet}}^{\ast}(n) = {{(r^{-1})}^{n}}^{-1} \to \varphi_{49}) \to r \in \mathbb{R} \to \varphi_{18}$
\item $\varphi_{100} :\equiv (r \in \mathbb{R} \to r \neq 0 \to \forall n \in \mathbb{N}\,({{(r^{-1})}^{n}} = {{r}^{n}}^{-1})) \to \varphi_{84}$
\item $\varphi_{101} :\equiv (r \in \mathbb{R} \to \varphi_{36}) \to (r \in \mathbb{R} \to 1 < r \to r \in \mathbb{R}) \to r \in \mathbb{R} \to 1 < r \to \varphi_{36}$
\item $\varphi_{102} :\equiv (\forall n \in \mathbb{N}\,(0 < {{(r^{-1})}^{n}}) \to \varphi_{1}) \to (r \in \mathbb{R} \to 1 < r \to \forall n \in \mathbb{N}\,(0 < {{(r^{-1})}^{n}})) \to \varphi_{43}$
\item $\varphi_{103} :\equiv ({{r}^{n}} \in \mathbb{R} \to {{r}^{n}} \neq 0 \to ({{r}^{n}}^{-1})^{-1} = {{r}^{n}}) \to (r \in \mathbb{R} \to 1 < r \to n \in \mathbb{N} \to {{r}^{n}} \in \mathbb{R}) \to \varphi_{72}$
\item $\varphi_{104} :\equiv {{(r^{-1})}^{n}}^{-1} = {{(r^{-1})}^{n}}^{-1} \to ({{(r^{-1})}^{n}} = {{r}^{n}}^{-1} \to \varphi_{46}) \to \varphi_{34}$
\item $\varphi_{105} :\equiv (\forall n\,(\varphi_{23}) \to \varphi_{33}) \to (r \in \mathbb{R} \land 1 < r \to \forall n\,(r \in \mathbb{R} \land 1 < r)) \to \varphi_{28}$
\item $\varphi_{106} :\equiv (r^{-1} \in \mathbb{R} \to \varphi_{67}) \to (r \in \mathbb{R} \to 1 < r \to r^{-1} \in \mathbb{R}) \to r \in \mathbb{R} \to 1 < r \to \varphi_{67}$
\item $\varphi_{107} :\equiv (\mathrm{Seq}({(r^{-1})^{\bullet}}) \to \varphi_{75}) \to (r \in \mathbb{R} \to 1 < r \to \mathrm{Seq}({(r^{-1})^{\bullet}})) \to \varphi_{87}$
\item $\varphi_{108} :\equiv (r \in \mathbb{R} \to \varphi_{18}) \to (r \in \mathbb{R} \to 1 < r \to n \in \mathbb{N} \to ({{r}^{n}}^{-1})^{-1} = {{r}^{n}}) \to \varphi_{95}$
\end{abbrevlist}
\end{fullproof}

\section{第3章 4節: 級数 (定義4.1, 命題4.3, 定理4.4, 定理4.7)}
% 自動生成 (prover/build.js): AddSwap4  ノード数 272
\begin{fullproof}{[E49]}{$\forall x\,\forall y\,\forall u\,\forall v\,(x \in \mathbb{R} \to y \in \mathbb{R} \to u \in \mathbb{R} \to v \in \mathbb{R} \to x + y + (u + v) = x + u + (y + v))$}{272}\phantomsection\label{tree:AddSwap4}
\noindent$\Pi_{1}$:
\begin{ptree}
\AXC{$y = z \to x = y \to x = z$\lbl{A6}}
\AXC{$(y = z \to x = y \to x = z) \to x = y \to y = z \to x = z$\lbl{A1}}
\BIC{$x = y \to y = z \to x = z$}
\UIC{$\forall z\,(x = y \to y = z \to x = z)$}
\UIC{$\forall y\,\forall z\,(x = y \to y = z \to x = z)$}
\UIC{$\forall x\,\forall y\,\forall z\,(x = y \to y = z \to x = z)$}
\AXC{$\varphi_{10}$\lbl{A2}}
\BIC{$\forall y\,\forall z\,(v_{873} = y \to y = z \to v_{873} = z)$}
\AXC{$\varphi_{23}$\lbl{A2}}
\BIC{$\forall z\,(v_{873} = v_{874} \to v_{874} = z \to v_{873} = z)$}
\AXC{$\varphi_{34}$\lbl{A2}}
\BIC{$v_{873} = v_{874} \to v_{874} = v_{875} \to v_{873} = v_{875}$}
\UIC{$\forall v_{875}\,(v_{873} = v_{874} \to v_{874} = v_{875} \to v_{873} = v_{875})$}
\UIC{$\varphi_{4}$}
\UIC{$\forall v_{873}\,(\varphi_{4})$}
\end{ptree}
\noindent$\Pi_{2}$:
\begin{ptree}
\AXC{$\stackrel{\Pi_{1}}{\vdots}$}\noLine
\UIC{$\forall v_{873}\,(\varphi_{4})$}
\AXC{$\forall v_{873}\,(\varphi_{4}) \to \varphi_{38}$\lbl{A2}}
\BIC{$\varphi_{38}$}
\AXC{$\varphi_{38} \to \varphi_{64}$\lbl{A2}}
\BIC{$\varphi_{64}$}
\AXC{$\varphi_{64} \to x + y + (u + v) = x + (u + (y + v)) \to \varphi_{43}$\lbl{A2}}
\BIC{$x + y + (u + v) = x + (u + (y + v)) \to \varphi_{43}$}
\end{ptree}
\noindent$\Pi_{3}$:
\begin{ptree}
\AXC{$x = x$\lbl{A5}}
\AXC{$x = y \to x = x \to y = x$\lbl{A6}}
\AXC{$(x = y \to x = x \to y = x) \to x = x \to x = y \to y = x$\lbl{A1}}
\BIC{$x = x \to x = y \to y = x$}
\BIC{$x = y \to y = x$}
\UIC{$\forall y\,(x = y \to y = x)$}
\UIC{$\forall x\,\forall y\,(x = y \to y = x)$}
\end{ptree}
\noindent$\Pi_{4}$:
\begin{ptree}
\AXC{$\stackrel{\Pi_{3}}{\vdots}$}\noLine
\UIC{$\forall x\,\forall y\,(x = y \to y = x)$}
\AXC{$\forall x\,\forall y\,(x = y \to y = x) \to \forall y\,(v_{862} = y \to y = v_{862})$\lbl{A2}}
\BIC{$\forall y\,(v_{862} = y \to y = v_{862})$}
\AXC{$\varphi_{1}$\lbl{A2}}
\BIC{$v_{862} = v_{863} \to v_{863} = v_{862}$}
\UIC{$\forall v_{863}\,(v_{862} = v_{863} \to v_{863} = v_{862})$}
\UIC{$\forall v_{862}\,\forall v_{863}\,(v_{862} = v_{863} \to v_{863} = v_{862})$}
\end{ptree}
\noindent$\Pi_{5}$:
\begin{ptree}
\AXC{$\stackrel{\Pi_{4}}{\vdots}$}\noLine
\UIC{$\forall v_{862}\,\forall v_{863}\,(v_{862} = v_{863} \to v_{863} = v_{862})$}
\AXC{$\forall v_{862}\,\forall v_{863}\,(v_{862} = v_{863} \to v_{863} = v_{862}) \to \varphi_{9}$\lbl{A2}}
\BIC{$\varphi_{9}$}
\AXC{$\varphi_{9} \to \varphi_{39}$\lbl{A2}}
\BIC{$\varphi_{39}$}
\end{ptree}
\noindent$\Pi_{6}$:
\begin{ptree}
\AXC{$\forall x\,\forall y\,(x \in \mathbb{R} \to y \in \mathbb{R} \to x + y \in \mathbb{R})$\lbl{(R1)}}
\AXC{$\varphi_{75}$\lbl{A2}}
\BIC{$\forall y\,(v_{858} \in \mathbb{R} \to y \in \mathbb{R} \to v_{858} + y \in \mathbb{R})$}
\AXC{$\varphi_{76}$\lbl{A2}}
\BIC{$v_{858} \in \mathbb{R} \to v_{859} \in \mathbb{R} \to v_{858} + v_{859} \in \mathbb{R}$}
\UIC{$\forall v_{859}\,(v_{858} \in \mathbb{R} \to v_{859} \in \mathbb{R} \to v_{858} + v_{859} \in \mathbb{R})$}
\UIC{$\varphi_{30}$}
\AXC{$\varphi_{30} \to \forall v_{859}\,(y \in \mathbb{R} \to v_{859} \in \mathbb{R} \to y + v_{859} \in \mathbb{R})$\lbl{A2}}
\BIC{$\forall v_{859}\,(y \in \mathbb{R} \to v_{859} \in \mathbb{R} \to y + v_{859} \in \mathbb{R})$}
\end{ptree}
\noindent$\Pi_{7}$:
\begin{ptree}
\AXC{$x \in \mathbb{R} \to y \in \mathbb{R} \to u \in \mathbb{R} \to v \in \mathbb{R} \to y \in \mathbb{R}$\lbl{A1}}
\AXC{$\stackrel{\Pi_{6}}{\vdots}$}\noLine
\UIC{$\forall v_{859}\,(y \in \mathbb{R} \to v_{859} \in \mathbb{R} \to y + v_{859} \in \mathbb{R})$}
\AXC{$\varphi_{73}$\lbl{A2}}
\BIC{$y \in \mathbb{R} \to v \in \mathbb{R} \to y + v \in \mathbb{R}$}
\AXC{$\varphi_{95}$\lbl{A1}}
\BIC{$(x \in \mathbb{R} \to y \in \mathbb{R} \to u \in \mathbb{R} \to v \in \mathbb{R} \to y \in \mathbb{R}) \to \varphi_{81}$}
\BIC{$\varphi_{81}$}
\end{ptree}
\noindent$\Pi_{8}$:
\begin{ptree}
\AXC{$\forall x\,\forall y\,(\varphi_{33})$\lbl{(R2)}}
\AXC{$\forall x\,\forall y\,(\varphi_{33}) \to \forall y\,(\varphi_{33})$\lbl{A2}}
\BIC{$\forall y\,(\varphi_{33})$}
\AXC{$\forall y\,(\varphi_{33}) \to \varphi_{31}$\lbl{A2}}
\BIC{$\varphi_{31}$}
\AXC{$\varphi_{31} \to x \in \mathbb{R} \to \varphi_{28}$\lbl{A2}}
\BIC{$x \in \mathbb{R} \to \varphi_{28}$}
\end{ptree}
\noindent$\Pi_{9}$:
\begin{ptree}
\AXC{$x \in \mathbb{R} \to y \in \mathbb{R} \to u \in \mathbb{R} \to v \in \mathbb{R} \to x \in \mathbb{R}$\lbl{A1}}
\AXC{$\stackrel{\Pi_{8}}{\vdots}$}\noLine
\UIC{$x \in \mathbb{R} \to \varphi_{28}$}
\AXC{$(x \in \mathbb{R} \to \varphi_{28}) \to \varphi_{99}$\lbl{A1}}
\BIC{$\varphi_{99}$}
\BIC{$(x \in \mathbb{R} \to y \in \mathbb{R} \to u \in \mathbb{R} \to v \in \mathbb{R} \to u \in \mathbb{R}) \to \varphi_{83}$}
\end{ptree}
\noindent$\Pi_{10}$:
\begin{ptree}
\AXC{$x \in \mathbb{R} \to y \in \mathbb{R} \to u \in \mathbb{R} \to v \in \mathbb{R} \to u \in \mathbb{R}$\lbl{A1}}
\AXC{$\stackrel{\Pi_{9}}{\vdots}$}\noLine
\UIC{$(x \in \mathbb{R} \to y \in \mathbb{R} \to u \in \mathbb{R} \to v \in \mathbb{R} \to u \in \mathbb{R}) \to \varphi_{83}$}
\BIC{$\varphi_{83}$}
\end{ptree}
\noindent$\Pi_{11}$:
\begin{ptree}
\AXC{$x \in \mathbb{R} \to y \in \mathbb{R} \to u \in \mathbb{R} \to v \in \mathbb{R} \to v \in \mathbb{R}$\lbl{A1}}
\AXC{$\stackrel{\Pi_{7}}{\vdots}$}\noLine
\UIC{$\varphi_{81}$}
\BIC{$x \in \mathbb{R} \to y \in \mathbb{R} \to u \in \mathbb{R} \to v \in \mathbb{R} \to y + v \in \mathbb{R}$}
\AXC{$\stackrel{\Pi_{10}}{\vdots}$}\noLine
\UIC{$\varphi_{83}$}
\BIC{$x \in \mathbb{R} \to y \in \mathbb{R} \to \varphi_{16}$}
\AXC{$(x \in \mathbb{R} \to y \in \mathbb{R} \to \varphi_{16}) \to \varphi_{39} \to x \in \mathbb{R} \to y \in \mathbb{R} \to \varphi_{17}$\lbl{A1}}
\BIC{$\varphi_{39} \to x \in \mathbb{R} \to y \in \mathbb{R} \to \varphi_{17}$}
\end{ptree}
\noindent$\Pi_{12}$:
\begin{ptree}
\AXC{$y = z \to x = y \to x = z$\lbl{A6}}
\AXC{$(y = z \to x = y \to x = z) \to x = y \to y = z \to x = z$\lbl{A1}}
\BIC{$x = y \to y = z \to x = z$}
\UIC{$\forall z\,(x = y \to y = z \to x = z)$}
\UIC{$\forall y\,\forall z\,(x = y \to y = z \to x = z)$}
\UIC{$\forall x\,\forall y\,\forall z\,(x = y \to y = z \to x = z)$}
\AXC{$\varphi_{11}$\lbl{A2}}
\BIC{$\forall y\,\forall z\,(v_{870} = y \to y = z \to v_{870} = z)$}
\AXC{$\varphi_{24}$\lbl{A2}}
\BIC{$\forall z\,(v_{870} = v_{871} \to v_{871} = z \to v_{870} = z)$}
\AXC{$\varphi_{35}$\lbl{A2}}
\BIC{$v_{870} = v_{871} \to v_{871} = v_{872} \to v_{870} = v_{872}$}
\UIC{$\forall v_{872}\,(v_{870} = v_{871} \to v_{871} = v_{872} \to v_{870} = v_{872})$}
\UIC{$\varphi_{5}$}
\UIC{$\forall v_{870}\,(\varphi_{5})$}
\end{ptree}
\noindent$\Pi_{13}$:
\begin{ptree}
\AXC{$\stackrel{\Pi_{12}}{\vdots}$}\noLine
\UIC{$\forall v_{870}\,(\varphi_{5})$}
\AXC{$\forall v_{870}\,(\varphi_{5}) \to \varphi_{40}$\lbl{A2}}
\BIC{$\varphi_{40}$}
\AXC{$\varphi_{40} \to \varphi_{65}$\lbl{A2}}
\BIC{$\varphi_{65}$}
\AXC{$\varphi_{65} \to x + y + (u + v) = x + (u + y + v) \to \varphi_{44}$\lbl{A2}}
\BIC{$x + y + (u + v) = x + (u + y + v) \to \varphi_{44}$}
\end{ptree}
\noindent$\Pi_{14}$:
\begin{ptree}
\AXC{$\varphi_{14}$\lbl{A6}}
\UIC{$\forall w\,(\varphi_{14})$}
\UIC{$\forall z\,\forall w\,(\varphi_{14})$}
\AXC{$\forall z\,\forall w\,(\varphi_{14}) \to \varphi_{60}$\lbl{A2}}
\BIC{$\varphi_{60}$}
\AXC{$\varphi_{60} \to u + y + v = u + (y + v) \to \varphi_{45}$\lbl{A2}}
\BIC{$u + y + v = u + (y + v) \to \varphi_{45}$}
\end{ptree}
\noindent$\Pi_{15}$:
\begin{ptree}
\AXC{$\forall x\,\forall y\,(\varphi_{33})$\lbl{(R2)}}
\AXC{$\forall x\,\forall y\,(\varphi_{33}) \to \forall y\,(\varphi_{32})$\lbl{A2}}
\BIC{$\forall y\,(\varphi_{32})$}
\AXC{$\forall y\,(\varphi_{32}) \to \varphi_{32}$\lbl{A2}}
\BIC{$\varphi_{32}$}
\end{ptree}
\noindent$\Pi_{16}$:
\begin{ptree}
\AXC{$\stackrel{\Pi_{15}}{\vdots}$}\noLine
\UIC{$\varphi_{32}$}
\AXC{$\varphi_{32} \to u \in \mathbb{R} \to y \in \mathbb{R} \to v \in \mathbb{R} \to u + y + v = u + (y + v)$\lbl{A2}}
\BIC{$u \in \mathbb{R} \to y \in \mathbb{R} \to v \in \mathbb{R} \to u + y + v = u + (y + v)$}
\AXC{$\varphi_{101}$\lbl{A1}}
\BIC{$(x \in \mathbb{R} \to y \in \mathbb{R} \to u \in \mathbb{R} \to v \in \mathbb{R} \to u \in \mathbb{R}) \to \varphi_{89}$}
\end{ptree}
\noindent$\Pi_{17}$:
\begin{ptree}
\AXC{$x \in \mathbb{R} \to y \in \mathbb{R} \to u \in \mathbb{R} \to v \in \mathbb{R} \to u \in \mathbb{R}$\lbl{A1}}
\AXC{$\stackrel{\Pi_{16}}{\vdots}$}\noLine
\UIC{$(x \in \mathbb{R} \to y \in \mathbb{R} \to u \in \mathbb{R} \to v \in \mathbb{R} \to u \in \mathbb{R}) \to \varphi_{89}$}
\BIC{$\varphi_{89}$}
\end{ptree}
\noindent$\Pi_{18}$:
\begin{ptree}
\AXC{$x \in \mathbb{R} \to y \in \mathbb{R} \to u \in \mathbb{R} \to v \in \mathbb{R} \to v \in \mathbb{R}$\lbl{A1}}
\AXC{$x \in \mathbb{R} \to y \in \mathbb{R} \to u \in \mathbb{R} \to v \in \mathbb{R} \to y \in \mathbb{R}$\lbl{A1}}
\AXC{$\stackrel{\Pi_{17}}{\vdots}$}\noLine
\UIC{$\varphi_{89}$}
\BIC{$(x \in \mathbb{R} \to y \in \mathbb{R} \to u \in \mathbb{R} \to v \in \mathbb{R} \to v \in \mathbb{R}) \to \varphi_{54}$}
\BIC{$\varphi_{54}$}
\AXC{$\varphi_{54} \to \varphi_{92}$\lbl{A1}}
\BIC{$\varphi_{92}$}
\end{ptree}
\noindent$\Pi_{19}$:
\begin{ptree}
\AXC{$x = x$\lbl{A5}}
\UIC{$\forall x\,(x = x)$}
\AXC{$\forall x\,(x = x) \to x + (u + y + v) = x + (u + y + v)$\lbl{A2}}
\BIC{$x + (u + y + v) = x + (u + y + v)$}
\AXC{$\stackrel{\Pi_{18}}{\vdots}$}\noLine
\UIC{$\varphi_{92}$}
\BIC{$(u + y + v = u + (y + v) \to \varphi_{45}) \to x \in \mathbb{R} \to y \in \mathbb{R} \to \varphi_{18}$}
\end{ptree}
\noindent$\Pi_{20}$:
\begin{ptree}
\AXC{$\stackrel{\Pi_{14}}{\vdots}$}\noLine
\UIC{$u + y + v = u + (y + v) \to \varphi_{45}$}
\AXC{$\stackrel{\Pi_{19}}{\vdots}$}\noLine
\UIC{$(u + y + v = u + (y + v) \to \varphi_{45}) \to x \in \mathbb{R} \to y \in \mathbb{R} \to \varphi_{18}$}
\BIC{$x \in \mathbb{R} \to y \in \mathbb{R} \to \varphi_{18}$}
\end{ptree}
\noindent$\Pi_{21}$:
\begin{ptree}
\AXC{$y = z \to x = y \to x = z$\lbl{A6}}
\AXC{$(y = z \to x = y \to x = z) \to x = y \to y = z \to x = z$\lbl{A1}}
\BIC{$x = y \to y = z \to x = z$}
\UIC{$\forall z\,(x = y \to y = z \to x = z)$}
\UIC{$\forall y\,\forall z\,(x = y \to y = z \to x = z)$}
\UIC{$\forall x\,\forall y\,\forall z\,(x = y \to y = z \to x = z)$}
\AXC{$\varphi_{12}$\lbl{A2}}
\BIC{$\forall y\,\forall z\,(v_{867} = y \to y = z \to v_{867} = z)$}
\AXC{$\varphi_{25}$\lbl{A2}}
\BIC{$\forall z\,(v_{867} = v_{868} \to v_{868} = z \to v_{867} = z)$}
\AXC{$\varphi_{36}$\lbl{A2}}
\BIC{$v_{867} = v_{868} \to v_{868} = v_{869} \to v_{867} = v_{869}$}
\UIC{$\forall v_{869}\,(v_{867} = v_{868} \to v_{868} = v_{869} \to v_{867} = v_{869})$}
\UIC{$\varphi_{6}$}
\UIC{$\forall v_{867}\,(\varphi_{6})$}
\end{ptree}
\noindent$\Pi_{22}$:
\begin{ptree}
\AXC{$\stackrel{\Pi_{21}}{\vdots}$}\noLine
\UIC{$\forall v_{867}\,(\varphi_{6})$}
\AXC{$\forall v_{867}\,(\varphi_{6}) \to \varphi_{41}$\lbl{A2}}
\BIC{$\varphi_{41}$}
\AXC{$\varphi_{41} \to \varphi_{66}$\lbl{A2}}
\BIC{$\varphi_{66}$}
\AXC{$\varphi_{66} \to x + y + (u + v) = x + (y + u + v) \to \varphi_{46}$\lbl{A2}}
\BIC{$x + y + (u + v) = x + (y + u + v) \to \varphi_{46}$}
\end{ptree}
\noindent$\Pi_{23}$:
\begin{ptree}
\AXC{$z = w \to \varphi_{19}$\lbl{A6}}
\UIC{$\forall w\,(z = w \to \varphi_{19})$}
\UIC{$\forall z\,\forall w\,(z = w \to \varphi_{19})$}
\AXC{$\forall z\,\forall w\,(z = w \to \varphi_{19}) \to \varphi_{61}$\lbl{A2}}
\BIC{$\varphi_{61}$}
\AXC{$\varphi_{61} \to y + u = u + y \to \varphi_{47}$\lbl{A2}}
\BIC{$y + u = u + y \to \varphi_{47}$}
\end{ptree}
\noindent$\Pi_{24}$:
\begin{ptree}
\AXC{$\forall x\,\forall y\,(x \in \mathbb{R} \to y \in \mathbb{R} \to x + y = y + x)$\lbl{(R3)}}
\AXC{$\varphi_{74}$\lbl{A2}}
\BIC{$\forall y\,(v_{856} \in \mathbb{R} \to y \in \mathbb{R} \to v_{856} + y = y + v_{856})$}
\AXC{$\varphi_{77}$\lbl{A2}}
\BIC{$v_{856} \in \mathbb{R} \to v_{857} \in \mathbb{R} \to v_{856} + v_{857} = v_{857} + v_{856}$}
\UIC{$\varphi_{27}$}
\UIC{$\forall v_{856}\,(\varphi_{27})$}
\AXC{$\varphi_{78}$\lbl{A2}}
\BIC{$\forall v_{857}\,(y \in \mathbb{R} \to v_{857} \in \mathbb{R} \to y + v_{857} = v_{857} + y)$}
\AXC{$\varphi_{72}$\lbl{A2}}
\BIC{$y \in \mathbb{R} \to u \in \mathbb{R} \to y + u = u + y$}
\AXC{$\varphi_{94}$\lbl{A1}}
\BIC{$(x \in \mathbb{R} \to y \in \mathbb{R} \to u \in \mathbb{R} \to v \in \mathbb{R} \to y \in \mathbb{R}) \to \varphi_{80}$}
\end{ptree}
\noindent$\Pi_{25}$:
\begin{ptree}
\AXC{$x \in \mathbb{R} \to y \in \mathbb{R} \to u \in \mathbb{R} \to v \in \mathbb{R} \to y \in \mathbb{R}$\lbl{A1}}
\AXC{$\stackrel{\Pi_{24}}{\vdots}$}\noLine
\UIC{$(x \in \mathbb{R} \to y \in \mathbb{R} \to u \in \mathbb{R} \to v \in \mathbb{R} \to y \in \mathbb{R}) \to \varphi_{80}$}
\BIC{$\varphi_{80}$}
\end{ptree}
\noindent$\Pi_{26}$:
\begin{ptree}
\AXC{$x \in \mathbb{R} \to y \in \mathbb{R} \to u \in \mathbb{R} \to v \in \mathbb{R} \to u \in \mathbb{R}$\lbl{A1}}
\AXC{$\stackrel{\Pi_{25}}{\vdots}$}\noLine
\UIC{$\varphi_{80}$}
\BIC{$x \in \mathbb{R} \to y \in \mathbb{R} \to u \in \mathbb{R} \to v \in \mathbb{R} \to y + u = u + y$}
\AXC{$\varphi_{97}$\lbl{A1}}
\BIC{$\varphi_{91}$}
\end{ptree}
\noindent$\Pi_{27}$:
\begin{ptree}
\AXC{$\stackrel{\Pi_{23}}{\vdots}$}\noLine
\UIC{$y + u = u + y \to \varphi_{47}$}
\AXC{$x = x$\lbl{A5}}
\UIC{$\forall x\,(x = x)$}
\AXC{$\forall x\,(x = x) \to x + (y + u + v) = x + (y + u + v)$\lbl{A2}}
\BIC{$x + (y + u + v) = x + (y + u + v)$}
\AXC{$\stackrel{\Pi_{26}}{\vdots}$}\noLine
\UIC{$\varphi_{91}$}
\BIC{$(y + u = u + y \to \varphi_{47}) \to x \in \mathbb{R} \to \varphi_{50}$}
\BIC{$x \in \mathbb{R} \to \varphi_{50}$}
\end{ptree}
\noindent$\Pi_{28}$:
\begin{ptree}
\AXC{$y = z \to x = y \to x = z$\lbl{A6}}
\AXC{$(y = z \to x = y \to x = z) \to x = y \to y = z \to x = z$\lbl{A1}}
\BIC{$x = y \to y = z \to x = z$}
\UIC{$\forall z\,(x = y \to y = z \to x = z)$}
\UIC{$\forall y\,\forall z\,(x = y \to y = z \to x = z)$}
\UIC{$\forall x\,\forall y\,\forall z\,(x = y \to y = z \to x = z)$}
\AXC{$\varphi_{13}$\lbl{A2}}
\BIC{$\forall y\,\forall z\,(v_{864} = y \to y = z \to v_{864} = z)$}
\AXC{$\varphi_{26}$\lbl{A2}}
\BIC{$\forall z\,(v_{864} = v_{865} \to v_{865} = z \to v_{864} = z)$}
\AXC{$\varphi_{37}$\lbl{A2}}
\BIC{$v_{864} = v_{865} \to v_{865} = v_{866} \to v_{864} = v_{866}$}
\UIC{$\forall v_{866}\,(v_{864} = v_{865} \to v_{865} = v_{866} \to v_{864} = v_{866})$}
\UIC{$\varphi_{7}$}
\UIC{$\forall v_{864}\,(\varphi_{7})$}
\end{ptree}
\noindent$\Pi_{29}$:
\begin{ptree}
\AXC{$\stackrel{\Pi_{28}}{\vdots}$}\noLine
\UIC{$\forall v_{864}\,(\varphi_{7})$}
\AXC{$\forall v_{864}\,(\varphi_{7}) \to \varphi_{42}$\lbl{A2}}
\BIC{$\varphi_{42}$}
\AXC{$\varphi_{42} \to \varphi_{67}$\lbl{A2}}
\BIC{$\varphi_{67}$}
\AXC{$\varphi_{67} \to x + y + (u + v) = x + (y + (u + v)) \to \varphi_{48}$\lbl{A2}}
\BIC{$x + y + (u + v) = x + (y + (u + v)) \to \varphi_{48}$}
\end{ptree}
\noindent$\Pi_{30}$:
\begin{ptree}
\AXC{$z = w \to \varphi_{8}$\lbl{A6}}
\UIC{$\forall w\,(z = w \to \varphi_{8})$}
\UIC{$\forall z\,\forall w\,(z = w \to \varphi_{8})$}
\AXC{$\forall z\,\forall w\,(z = w \to \varphi_{8}) \to \varphi_{62}$\lbl{A2}}
\BIC{$\varphi_{62}$}
\AXC{$\varphi_{62} \to y + (u + v) = y + u + v \to \varphi_{49}$\lbl{A2}}
\BIC{$y + (u + v) = y + u + v \to \varphi_{49}$}
\end{ptree}
\noindent$\Pi_{31}$:
\begin{ptree}
\AXC{$\stackrel{\Pi_{3}}{\vdots}$}\noLine
\UIC{$\forall x\,\forall y\,(x = y \to y = x)$}
\AXC{$\forall x\,\forall y\,(x = y \to y = x) \to \forall y\,(v_{860} = y \to y = v_{860})$\lbl{A2}}
\BIC{$\forall y\,(v_{860} = y \to y = v_{860})$}
\AXC{$\varphi_{3}$\lbl{A2}}
\BIC{$v_{860} = v_{861} \to v_{861} = v_{860}$}
\UIC{$\forall v_{861}\,(v_{860} = v_{861} \to v_{861} = v_{860})$}
\UIC{$\forall v_{860}\,\forall v_{861}\,(v_{860} = v_{861} \to v_{861} = v_{860})$}
\AXC{$\varphi_{57}$\lbl{A2}}
\BIC{$\forall v_{861}\,(y + u + v = v_{861} \to v_{861} = y + u + v)$}
\AXC{$\varphi_{70}$\lbl{A2}}
\BIC{$y + u + v = y + (u + v) \to y + (u + v) = y + u + v$}
\end{ptree}
\noindent$\Pi_{32}$:
\begin{ptree}
\AXC{$\forall x\,\forall y\,(\varphi_{33})$\lbl{(R2)}}
\AXC{$\forall x\,\forall y\,(\varphi_{33}) \to \varphi_{58}$\lbl{A2}}
\BIC{$\varphi_{58}$}
\AXC{$\varphi_{58} \to \varphi_{63}$\lbl{A2}}
\BIC{$\varphi_{63}$}
\AXC{$\varphi_{63} \to v_{853} \in \mathbb{R} \to v_{854} \in \mathbb{R} \to \varphi_{15}$\lbl{A2}}
\BIC{$v_{853} \in \mathbb{R} \to v_{854} \in \mathbb{R} \to \varphi_{15}$}
\UIC{$\forall v_{855}\,(v_{853} \in \mathbb{R} \to v_{854} \in \mathbb{R} \to \varphi_{15})$}
\UIC{$\forall v_{854}\,\forall v_{855}\,(v_{853} \in \mathbb{R} \to v_{854} \in \mathbb{R} \to \varphi_{15})$}
\UIC{$\forall v_{853}\,\forall v_{854}\,\forall v_{855}\,(v_{853} \in \mathbb{R} \to v_{854} \in \mathbb{R} \to \varphi_{15})$}
\AXC{$\varphi_{88}$\lbl{A2}}
\BIC{$\varphi_{69}$}
\end{ptree}
\noindent$\Pi_{33}$:
\begin{ptree}
\AXC{$\stackrel{\Pi_{32}}{\vdots}$}\noLine
\UIC{$\varphi_{69}$}
\AXC{$\varphi_{69} \to \varphi_{59}$\lbl{A2}}
\BIC{$\varphi_{59}$}
\AXC{$\varphi_{59} \to y \in \mathbb{R} \to u \in \mathbb{R} \to v \in \mathbb{R} \to y + u + v = y + (u + v)$\lbl{A2}}
\BIC{$y \in \mathbb{R} \to u \in \mathbb{R} \to v \in \mathbb{R} \to y + u + v = y + (u + v)$}
\AXC{$\varphi_{102}$\lbl{A1}}
\BIC{$(x \in \mathbb{R} \to y \in \mathbb{R} \to u \in \mathbb{R} \to v \in \mathbb{R} \to y \in \mathbb{R}) \to \varphi_{90}$}
\end{ptree}
\noindent$\Pi_{34}$:
\begin{ptree}
\AXC{$x \in \mathbb{R} \to y \in \mathbb{R} \to u \in \mathbb{R} \to v \in \mathbb{R} \to y \in \mathbb{R}$\lbl{A1}}
\AXC{$\stackrel{\Pi_{33}}{\vdots}$}\noLine
\UIC{$(x \in \mathbb{R} \to y \in \mathbb{R} \to u \in \mathbb{R} \to v \in \mathbb{R} \to y \in \mathbb{R}) \to \varphi_{90}$}
\BIC{$\varphi_{90}$}
\end{ptree}
\noindent$\Pi_{35}$:
\begin{ptree}
\AXC{$x \in \mathbb{R} \to y \in \mathbb{R} \to u \in \mathbb{R} \to v \in \mathbb{R} \to v \in \mathbb{R}$\lbl{A1}}
\AXC{$x \in \mathbb{R} \to y \in \mathbb{R} \to u \in \mathbb{R} \to v \in \mathbb{R} \to u \in \mathbb{R}$\lbl{A1}}
\AXC{$\stackrel{\Pi_{34}}{\vdots}$}\noLine
\UIC{$\varphi_{90}$}
\BIC{$(x \in \mathbb{R} \to y \in \mathbb{R} \to u \in \mathbb{R} \to v \in \mathbb{R} \to v \in \mathbb{R}) \to \varphi_{55}$}
\BIC{$\varphi_{55}$}
\AXC{$\varphi_{55} \to \varphi_{79}$\lbl{A1}}
\BIC{$\varphi_{79}$}
\end{ptree}
\noindent$\Pi_{36}$:
\begin{ptree}
\AXC{$x = x$\lbl{A5}}
\UIC{$\forall x\,(x = x)$}
\AXC{$\varphi_{2}$\lbl{A2}}
\BIC{$x + (y + (u + v)) = x + (y + (u + v))$}
\AXC{$\stackrel{\Pi_{31}}{\vdots}$}\noLine
\UIC{$y + u + v = y + (u + v) \to y + (u + v) = y + u + v$}
\AXC{$\stackrel{\Pi_{35}}{\vdots}$}\noLine
\UIC{$\varphi_{79}$}
\BIC{$\varphi_{56}$}
\AXC{$\varphi_{56} \to \varphi_{93}$\lbl{A1}}
\BIC{$\varphi_{93}$}
\BIC{$(y + (u + v) = y + u + v \to \varphi_{49}) \to x \in \mathbb{R} \to y \in \mathbb{R} \to \varphi_{20}$}
\end{ptree}
\noindent$\Pi_{37}$:
\begin{ptree}
\AXC{$\stackrel{\Pi_{30}}{\vdots}$}\noLine
\UIC{$y + (u + v) = y + u + v \to \varphi_{49}$}
\AXC{$\stackrel{\Pi_{36}}{\vdots}$}\noLine
\UIC{$(y + (u + v) = y + u + v \to \varphi_{49}) \to x \in \mathbb{R} \to y \in \mathbb{R} \to \varphi_{20}$}
\BIC{$x \in \mathbb{R} \to y \in \mathbb{R} \to \varphi_{20}$}
\end{ptree}
\noindent$\Pi_{38}$:
\begin{ptree}
\AXC{$\forall x\,\forall y\,(x \in \mathbb{R} \to y \in \mathbb{R} \to x + y \in \mathbb{R})$\lbl{(R1)}}
\AXC{$\varphi_{71}$\lbl{A2}}
\BIC{$\forall y\,(u \in \mathbb{R} \to y \in \mathbb{R} \to u + y \in \mathbb{R})$}
\AXC{$\varphi_{68}$\lbl{A2}}
\BIC{$u \in \mathbb{R} \to v \in \mathbb{R} \to u + v \in \mathbb{R}$}
\AXC{$\varphi_{96}$\lbl{A1}}
\BIC{$(x \in \mathbb{R} \to y \in \mathbb{R} \to u \in \mathbb{R} \to v \in \mathbb{R} \to u \in \mathbb{R}) \to \varphi_{82}$}
\end{ptree}
\noindent$\Pi_{39}$:
\begin{ptree}
\AXC{$x \in \mathbb{R} \to y \in \mathbb{R} \to u \in \mathbb{R} \to v \in \mathbb{R} \to u \in \mathbb{R}$\lbl{A1}}
\AXC{$\stackrel{\Pi_{38}}{\vdots}$}\noLine
\UIC{$(x \in \mathbb{R} \to y \in \mathbb{R} \to u \in \mathbb{R} \to v \in \mathbb{R} \to u \in \mathbb{R}) \to \varphi_{82}$}
\BIC{$\varphi_{82}$}
\end{ptree}
\noindent$\Pi_{40}$:
\begin{ptree}
\AXC{$\forall x\,\forall y\,(\varphi_{33})$\lbl{(R2)}}
\AXC{$\forall x\,\forall y\,(\varphi_{33}) \to \forall y\,(\varphi_{33})$\lbl{A2}}
\BIC{$\forall y\,(\varphi_{33})$}
\AXC{$\forall y\,(\varphi_{33}) \to \varphi_{33}$\lbl{A2}}
\BIC{$\varphi_{33}$}
\AXC{$\varphi_{33} \to x \in \mathbb{R} \to \varphi_{29}$\lbl{A2}}
\BIC{$x \in \mathbb{R} \to \varphi_{29}$}
\end{ptree}
\noindent$\Pi_{41}$:
\begin{ptree}
\AXC{$x \in \mathbb{R} \to y \in \mathbb{R} \to u \in \mathbb{R} \to v \in \mathbb{R} \to x \in \mathbb{R}$\lbl{A1}}
\AXC{$\stackrel{\Pi_{40}}{\vdots}$}\noLine
\UIC{$x \in \mathbb{R} \to \varphi_{29}$}
\AXC{$(x \in \mathbb{R} \to \varphi_{29}) \to \varphi_{100}$\lbl{A1}}
\BIC{$\varphi_{100}$}
\BIC{$(x \in \mathbb{R} \to y \in \mathbb{R} \to u \in \mathbb{R} \to v \in \mathbb{R} \to y \in \mathbb{R}) \to \varphi_{84}$}
\end{ptree}
\noindent$\Pi_{42}$:
\begin{ptree}
\AXC{$x \in \mathbb{R} \to y \in \mathbb{R} \to u \in \mathbb{R} \to v \in \mathbb{R} \to y \in \mathbb{R}$\lbl{A1}}
\AXC{$\stackrel{\Pi_{41}}{\vdots}$}\noLine
\UIC{$(x \in \mathbb{R} \to y \in \mathbb{R} \to u \in \mathbb{R} \to v \in \mathbb{R} \to y \in \mathbb{R}) \to \varphi_{84}$}
\BIC{$\varphi_{84}$}
\end{ptree}
\noindent$\Pi_{43}$:
\begin{ptree}
\AXC{$x \in \mathbb{R} \to y \in \mathbb{R} \to u \in \mathbb{R} \to v \in \mathbb{R} \to v \in \mathbb{R}$\lbl{A1}}
\AXC{$\stackrel{\Pi_{39}}{\vdots}$}\noLine
\UIC{$\varphi_{82}$}
\BIC{$x \in \mathbb{R} \to y \in \mathbb{R} \to u \in \mathbb{R} \to v \in \mathbb{R} \to u + v \in \mathbb{R}$}
\AXC{$\stackrel{\Pi_{42}}{\vdots}$}\noLine
\UIC{$\varphi_{84}$}
\BIC{$x \in \mathbb{R} \to y \in \mathbb{R} \to \varphi_{21}$}
\end{ptree}
\noindent$\Pi_{44}$:
\begin{ptree}
\AXC{$\stackrel{\Pi_{37}}{\vdots}$}\noLine
\UIC{$x \in \mathbb{R} \to y \in \mathbb{R} \to \varphi_{20}$}
\AXC{$\stackrel{\Pi_{43}}{\vdots}$}\noLine
\UIC{$x \in \mathbb{R} \to y \in \mathbb{R} \to \varphi_{21}$}
\AXC{$(x \in \mathbb{R} \to y \in \mathbb{R} \to \varphi_{21}) \to (x \in \mathbb{R} \to y \in \mathbb{R} \to \varphi_{20}) \to \varphi_{85}$\lbl{A1}}
\BIC{$(x \in \mathbb{R} \to y \in \mathbb{R} \to \varphi_{20}) \to \varphi_{85}$}
\BIC{$\varphi_{85}$}
\end{ptree}
\noindent$\Pi_{45}$:
\begin{ptree}
\AXC{$\stackrel{\Pi_{27}}{\vdots}$}\noLine
\UIC{$x \in \mathbb{R} \to \varphi_{50}$}
\AXC{$\stackrel{\Pi_{29}}{\vdots}$}\noLine
\UIC{$x + y + (u + v) = x + (y + (u + v)) \to \varphi_{48}$}
\AXC{$\stackrel{\Pi_{44}}{\vdots}$}\noLine
\UIC{$\varphi_{85}$}
\BIC{$x \in \mathbb{R} \to \varphi_{51}$}
\AXC{$(x \in \mathbb{R} \to \varphi_{51}) \to \varphi_{98}$\lbl{A1}}
\BIC{$\varphi_{98}$}
\BIC{$(x + y + (u + v) = x + (y + u + v) \to \varphi_{46}) \to x \in \mathbb{R} \to \varphi_{52}$}
\end{ptree}
\noindent$\Pi_{46}$:
\begin{ptree}
\AXC{$\stackrel{\Pi_{22}}{\vdots}$}\noLine
\UIC{$x + y + (u + v) = x + (y + u + v) \to \varphi_{46}$}
\AXC{$\stackrel{\Pi_{45}}{\vdots}$}\noLine
\UIC{$(x + y + (u + v) = x + (y + u + v) \to \varphi_{46}) \to x \in \mathbb{R} \to \varphi_{52}$}
\BIC{$x \in \mathbb{R} \to \varphi_{52}$}
\end{ptree}
\noindent$\Pi_{47}$:
\begin{ptree}
\AXC{$\stackrel{\Pi_{20}}{\vdots}$}\noLine
\UIC{$x \in \mathbb{R} \to y \in \mathbb{R} \to \varphi_{18}$}
\AXC{$\stackrel{\Pi_{46}}{\vdots}$}\noLine
\UIC{$x \in \mathbb{R} \to \varphi_{52}$}
\AXC{$(x \in \mathbb{R} \to \varphi_{52}) \to (x \in \mathbb{R} \to y \in \mathbb{R} \to \varphi_{18}) \to \varphi_{86}$\lbl{A1}}
\BIC{$(x \in \mathbb{R} \to y \in \mathbb{R} \to \varphi_{18}) \to \varphi_{86}$}
\BIC{$\varphi_{86}$}
\end{ptree}
\noindent$\Pi_{48}$:
\begin{ptree}
\AXC{$\stackrel{\Pi_{13}}{\vdots}$}\noLine
\UIC{$x + y + (u + v) = x + (u + y + v) \to \varphi_{44}$}
\AXC{$\stackrel{\Pi_{47}}{\vdots}$}\noLine
\UIC{$\varphi_{86}$}
\BIC{$x \in \mathbb{R} \to y \in \mathbb{R} \to \varphi_{22}$}
\AXC{$(x \in \mathbb{R} \to y \in \mathbb{R} \to \varphi_{22}) \to (x \in \mathbb{R} \to y \in \mathbb{R} \to \varphi_{17}) \to \varphi_{87}$\lbl{A1}}
\BIC{$(x \in \mathbb{R} \to y \in \mathbb{R} \to \varphi_{17}) \to \varphi_{87}$}
\end{ptree}
\noindent 主木:
\begin{ptree}
\AXC{$\stackrel{\Pi_{2}}{\vdots}$}\noLine
\UIC{$x + y + (u + v) = x + (u + (y + v)) \to \varphi_{43}$}
\AXC{$\stackrel{\Pi_{5}}{\vdots}$}\noLine
\UIC{$\varphi_{39}$}
\AXC{$\stackrel{\Pi_{11}}{\vdots}$}\noLine
\UIC{$\varphi_{39} \to x \in \mathbb{R} \to y \in \mathbb{R} \to \varphi_{17}$}
\BIC{$x \in \mathbb{R} \to y \in \mathbb{R} \to \varphi_{17}$}
\AXC{$\stackrel{\Pi_{48}}{\vdots}$}\noLine
\UIC{$(x \in \mathbb{R} \to y \in \mathbb{R} \to \varphi_{17}) \to \varphi_{87}$}
\BIC{$\varphi_{87}$}
\BIC{$x \in \mathbb{R} \to \varphi_{53}$}
\UIC{$\forall v\,(x \in \mathbb{R} \to \varphi_{53})$}
\UIC{$\forall u\,\forall v\,(x \in \mathbb{R} \to \varphi_{53})$}
\UIC{$\forall y\,\forall u\,\forall v\,(x \in \mathbb{R} \to \varphi_{53})$}
\UIC{$\forall x\,\forall y\,\forall u\,\forall v\,(x \in \mathbb{R} \to \varphi_{53})$}
\end{ptree}
\begin{abbrevlist}
\item $\varphi_{1} :\equiv \forall y\,(v_{862} = y \to y = v_{862}) \to v_{862} = v_{863} \to v_{863} = v_{862}$
\item $\varphi_{2} :\equiv \forall x\,(x = x) \to x + (y + (u + v)) = x + (y + (u + v))$
\item $\varphi_{3} :\equiv \forall y\,(v_{860} = y \to y = v_{860}) \to v_{860} = v_{861} \to v_{861} = v_{860}$
\item $\varphi_{4} :\equiv \forall v_{874}\,\forall v_{875}\,(v_{873} = v_{874} \to v_{874} = v_{875} \to v_{873} = v_{875})$
\item $\varphi_{5} :\equiv \forall v_{871}\,\forall v_{872}\,(v_{870} = v_{871} \to v_{871} = v_{872} \to v_{870} = v_{872})$
\item $\varphi_{6} :\equiv \forall v_{868}\,\forall v_{869}\,(v_{867} = v_{868} \to v_{868} = v_{869} \to v_{867} = v_{869})$
\item $\varphi_{7} :\equiv \forall v_{865}\,\forall v_{866}\,(v_{864} = v_{865} \to v_{865} = v_{866} \to v_{864} = v_{866})$
\item $\varphi_{8} :\equiv x + (y + (u + v)) = x + z \to x + (y + (u + v)) = x + w$
\item $\varphi_{9} :\equiv \forall v_{863}\,(x + u + (y + v) = v_{863} \to v_{863} = x + u + (y + v))$
\item $\varphi_{10} :\equiv \forall x\,\forall y\,\forall z\,(x = y \to y = z \to x = z) \to \forall y\,\forall z\,(v_{873} = y \to y = z \to v_{873} = z)$
\item $\varphi_{11} :\equiv \forall x\,\forall y\,\forall z\,(x = y \to y = z \to x = z) \to \forall y\,\forall z\,(v_{870} = y \to y = z \to v_{870} = z)$
\item $\varphi_{12} :\equiv \forall x\,\forall y\,\forall z\,(x = y \to y = z \to x = z) \to \forall y\,\forall z\,(v_{867} = y \to y = z \to v_{867} = z)$
\item $\varphi_{13} :\equiv \forall x\,\forall y\,\forall z\,(x = y \to y = z \to x = z) \to \forall y\,\forall z\,(v_{864} = y \to y = z \to v_{864} = z)$
\item $\varphi_{14} :\equiv z = w \to x + (u + y + v) = x + z \to x + (u + y + v) = x + w$
\item $\varphi_{15} :\equiv v_{855} \in \mathbb{R} \to v_{853} + v_{854} + v_{855} = v_{853} + (v_{854} + v_{855})$
\item $\varphi_{16} :\equiv u \in \mathbb{R} \to v \in \mathbb{R} \to x + u + (y + v) = x + (u + (y + v))$
\item $\varphi_{17} :\equiv u \in \mathbb{R} \to v \in \mathbb{R} \to x + (u + (y + v)) = x + u + (y + v)$
\item $\varphi_{18} :\equiv u \in \mathbb{R} \to v \in \mathbb{R} \to x + (u + y + v) = x + (u + (y + v))$
\item $\varphi_{19} :\equiv x + (y + u + v) = x + (z + v) \to x + (y + u + v) = x + (w + v)$
\item $\varphi_{20} :\equiv u \in \mathbb{R} \to v \in \mathbb{R} \to x + (y + (u + v)) = x + (y + u + v)$
\item $\varphi_{21} :\equiv u \in \mathbb{R} \to v \in \mathbb{R} \to x + y + (u + v) = x + (y + (u + v))$
\item $\varphi_{22} :\equiv u \in \mathbb{R} \to v \in \mathbb{R} \to x + y + (u + v) = x + (u + (y + v))$
\item $\varphi_{23} :\equiv \forall y\,\forall z\,(v_{873} = y \to y = z \to v_{873} = z) \to \forall z\,(v_{873} = v_{874} \to v_{874} = z \to v_{873} = z)$
\item $\varphi_{24} :\equiv \forall y\,\forall z\,(v_{870} = y \to y = z \to v_{870} = z) \to \forall z\,(v_{870} = v_{871} \to v_{871} = z \to v_{870} = z)$
\item $\varphi_{25} :\equiv \forall y\,\forall z\,(v_{867} = y \to y = z \to v_{867} = z) \to \forall z\,(v_{867} = v_{868} \to v_{868} = z \to v_{867} = z)$
\item $\varphi_{26} :\equiv \forall y\,\forall z\,(v_{864} = y \to y = z \to v_{864} = z) \to \forall z\,(v_{864} = v_{865} \to v_{865} = z \to v_{864} = z)$
\item $\varphi_{27} :\equiv \forall v_{857}\,(v_{856} \in \mathbb{R} \to v_{857} \in \mathbb{R} \to v_{856} + v_{857} = v_{857} + v_{856})$
\item $\varphi_{28} :\equiv u \in \mathbb{R} \to y + v \in \mathbb{R} \to x + u + (y + v) = x + (u + (y + v))$
\item $\varphi_{29} :\equiv y \in \mathbb{R} \to u + v \in \mathbb{R} \to x + y + (u + v) = x + (y + (u + v))$
\item $\varphi_{30} :\equiv \forall v_{858}\,\forall v_{859}\,(v_{858} \in \mathbb{R} \to v_{859} \in \mathbb{R} \to v_{858} + v_{859} \in \mathbb{R})$
\item $\varphi_{31} :\equiv \forall z\,(x \in \mathbb{R} \to u \in \mathbb{R} \to z \in \mathbb{R} \to x + u + z = x + (u + z))$
\item $\varphi_{32} :\equiv \forall z\,(u \in \mathbb{R} \to y \in \mathbb{R} \to z \in \mathbb{R} \to u + y + z = u + (y + z))$
\item $\varphi_{33} :\equiv \forall z\,(x \in \mathbb{R} \to y \in \mathbb{R} \to z \in \mathbb{R} \to x + y + z = x + (y + z))$
\item $\varphi_{34} :\equiv \forall z\,(v_{873} = v_{874} \to v_{874} = z \to v_{873} = z) \to v_{873} = v_{874} \to v_{874} = v_{875} \to v_{873} = v_{875}$
\item $\varphi_{35} :\equiv \forall z\,(v_{870} = v_{871} \to v_{871} = z \to v_{870} = z) \to v_{870} = v_{871} \to v_{871} = v_{872} \to v_{870} = v_{872}$
\item $\varphi_{36} :\equiv \forall z\,(v_{867} = v_{868} \to v_{868} = z \to v_{867} = z) \to v_{867} = v_{868} \to v_{868} = v_{869} \to v_{867} = v_{869}$
\item $\varphi_{37} :\equiv \forall z\,(v_{864} = v_{865} \to v_{865} = z \to v_{864} = z) \to v_{864} = v_{865} \to v_{865} = v_{866} \to v_{864} = v_{866}$
\item $\varphi_{38} :\equiv \forall v_{874}\,\forall v_{875}\,(x + y + (u + v) = v_{874} \to v_{874} = v_{875} \to x + y + (u + v) = v_{875})$
\item $\varphi_{39} :\equiv x + u + (y + v) = x + (u + (y + v)) \to x + (u + (y + v)) = x + u + (y + v)$
\item $\varphi_{40} :\equiv \forall v_{871}\,\forall v_{872}\,(x + y + (u + v) = v_{871} \to v_{871} = v_{872} \to x + y + (u + v) = v_{872})$
\item $\varphi_{41} :\equiv \forall v_{868}\,\forall v_{869}\,(x + y + (u + v) = v_{868} \to v_{868} = v_{869} \to x + y + (u + v) = v_{869})$
\item $\varphi_{42} :\equiv \forall v_{865}\,\forall v_{866}\,(x + y + (u + v) = v_{865} \to v_{865} = v_{866} \to x + y + (u + v) = v_{866})$
\item $\varphi_{43} :\equiv x + (u + (y + v)) = x + u + (y + v) \to x + y + (u + v) = x + u + (y + v)$
\item $\varphi_{44} :\equiv x + (u + y + v) = x + (u + (y + v)) \to x + y + (u + v) = x + (u + (y + v))$
\item $\varphi_{45} :\equiv x + (u + y + v) = x + (u + y + v) \to x + (u + y + v) = x + (u + (y + v))$
\item $\varphi_{46} :\equiv x + (y + u + v) = x + (u + y + v) \to x + y + (u + v) = x + (u + y + v)$
\item $\varphi_{47} :\equiv x + (y + u + v) = x + (y + u + v) \to x + (y + u + v) = x + (u + y + v)$
\item $\varphi_{48} :\equiv x + (y + (u + v)) = x + (y + u + v) \to x + y + (u + v) = x + (y + u + v)$
\item $\varphi_{49} :\equiv x + (y + (u + v)) = x + (y + (u + v)) \to x + (y + (u + v)) = x + (y + u + v)$
\item $\varphi_{50} :\equiv y \in \mathbb{R} \to u \in \mathbb{R} \to v \in \mathbb{R} \to x + (y + u + v) = x + (u + y + v)$
\item $\varphi_{51} :\equiv y \in \mathbb{R} \to u \in \mathbb{R} \to v \in \mathbb{R} \to x + y + (u + v) = x + (y + u + v)$
\item $\varphi_{52} :\equiv y \in \mathbb{R} \to u \in \mathbb{R} \to v \in \mathbb{R} \to x + y + (u + v) = x + (u + y + v)$
\item $\varphi_{53} :\equiv y \in \mathbb{R} \to u \in \mathbb{R} \to v \in \mathbb{R} \to x + y + (u + v) = x + u + (y + v)$
\item $\varphi_{54} :\equiv x \in \mathbb{R} \to y \in \mathbb{R} \to u \in \mathbb{R} \to v \in \mathbb{R} \to u + y + v = u + (y + v)$
\item $\varphi_{55} :\equiv x \in \mathbb{R} \to y \in \mathbb{R} \to u \in \mathbb{R} \to v \in \mathbb{R} \to y + u + v = y + (u + v)$
\item $\varphi_{56} :\equiv x \in \mathbb{R} \to y \in \mathbb{R} \to u \in \mathbb{R} \to v \in \mathbb{R} \to y + (u + v) = y + u + v$
\item $\varphi_{57} :\equiv \forall v_{860}\,\forall v_{861}\,(v_{860} = v_{861} \to v_{861} = v_{860}) \to \forall v_{861}\,(y + u + v = v_{861} \to v_{861} = y + u + v)$
\item $\varphi_{58} :\equiv \forall y\,\forall z\,(v_{853} \in \mathbb{R} \to y \in \mathbb{R} \to z \in \mathbb{R} \to v_{853} + y + z = v_{853} + (y + z))$
\item $\varphi_{59} :\equiv \forall v_{855}\,(y \in \mathbb{R} \to u \in \mathbb{R} \to v_{855} \in \mathbb{R} \to y + u + v_{855} = y + (u + v_{855}))$
\item $\varphi_{60} :\equiv \forall w\,(u + y + v = w \to x + (u + y + v) = x + (u + y + v) \to x + (u + y + v) = x + w)$
\item $\varphi_{61} :\equiv \forall w\,(y + u = w \to x + (y + u + v) = x + (y + u + v) \to x + (y + u + v) = x + (w + v))$
\item $\varphi_{62} :\equiv \forall w\,(y + (u + v) = w \to x + (y + (u + v)) = x + (y + (u + v)) \to x + (y + (u + v)) = x + w)$
\item $\varphi_{63} :\equiv \forall z\,(v_{853} \in \mathbb{R} \to v_{854} \in \mathbb{R} \to z \in \mathbb{R} \to v_{853} + v_{854} + z = v_{853} + (v_{854} + z))$
\item $\varphi_{64} :\equiv \forall v_{875}\,(x + y + (u + v) = x + (u + (y + v)) \to x + (u + (y + v)) = v_{875} \to x + y + (u + v) = v_{875})$
\item $\varphi_{65} :\equiv \forall v_{872}\,(x + y + (u + v) = x + (u + y + v) \to x + (u + y + v) = v_{872} \to x + y + (u + v) = v_{872})$
\item $\varphi_{66} :\equiv \forall v_{869}\,(x + y + (u + v) = x + (y + u + v) \to x + (y + u + v) = v_{869} \to x + y + (u + v) = v_{869})$
\item $\varphi_{67} :\equiv \forall v_{866}\,(x + y + (u + v) = x + (y + (u + v)) \to x + (y + (u + v)) = v_{866} \to x + y + (u + v) = v_{866})$
\item $\varphi_{68} :\equiv \forall y\,(u \in \mathbb{R} \to y \in \mathbb{R} \to u + y \in \mathbb{R}) \to u \in \mathbb{R} \to v \in \mathbb{R} \to u + v \in \mathbb{R}$
\item $\varphi_{69} :\equiv \forall v_{854}\,\forall v_{855}\,(y \in \mathbb{R} \to v_{854} \in \mathbb{R} \to v_{855} \in \mathbb{R} \to y + v_{854} + v_{855} = y + (v_{854} + v_{855}))$
\item $\varphi_{70} :\equiv \forall v_{861}\,(y + u + v = v_{861} \to v_{861} = y + u + v) \to y + u + v = y + (u + v) \to y + (u + v) = y + u + v$
\item $\varphi_{71} :\equiv \forall x\,\forall y\,(x \in \mathbb{R} \to y \in \mathbb{R} \to x + y \in \mathbb{R}) \to \forall y\,(u \in \mathbb{R} \to y \in \mathbb{R} \to u + y \in \mathbb{R})$
\item $\varphi_{72} :\equiv \forall v_{857}\,(y \in \mathbb{R} \to v_{857} \in \mathbb{R} \to y + v_{857} = v_{857} + y) \to y \in \mathbb{R} \to u \in \mathbb{R} \to y + u = u + y$
\item $\varphi_{73} :\equiv \forall v_{859}\,(y \in \mathbb{R} \to v_{859} \in \mathbb{R} \to y + v_{859} \in \mathbb{R}) \to y \in \mathbb{R} \to v \in \mathbb{R} \to y + v \in \mathbb{R}$
\item $\varphi_{74} :\equiv \forall x\,\forall y\,(x \in \mathbb{R} \to y \in \mathbb{R} \to x + y = y + x) \to \forall y\,(v_{856} \in \mathbb{R} \to y \in \mathbb{R} \to v_{856} + y = y + v_{856})$
\item $\varphi_{75} :\equiv \forall x\,\forall y\,(x \in \mathbb{R} \to y \in \mathbb{R} \to x + y \in \mathbb{R}) \to \forall y\,(v_{858} \in \mathbb{R} \to y \in \mathbb{R} \to v_{858} + y \in \mathbb{R})$
\item $\varphi_{76} :\equiv \forall y\,(v_{858} \in \mathbb{R} \to y \in \mathbb{R} \to v_{858} + y \in \mathbb{R}) \to v_{858} \in \mathbb{R} \to v_{859} \in \mathbb{R} \to v_{858} + v_{859} \in \mathbb{R}$
\item $\varphi_{77} :\equiv \forall y\,(v_{856} \in \mathbb{R} \to y \in \mathbb{R} \to v_{856} + y = y + v_{856}) \to v_{856} \in \mathbb{R} \to v_{857} \in \mathbb{R} \to v_{856} + v_{857} = v_{857} + v_{856}$
\item $\varphi_{78} :\equiv \forall v_{856}\,(\varphi_{27}) \to \forall v_{857}\,(y \in \mathbb{R} \to v_{857} \in \mathbb{R} \to y + v_{857} = v_{857} + y)$
\item $\varphi_{79} :\equiv (y + u + v = y + (u + v) \to y + (u + v) = y + u + v) \to \varphi_{56}$
\item $\varphi_{80} :\equiv (x \in \mathbb{R} \to y \in \mathbb{R} \to u \in \mathbb{R} \to v \in \mathbb{R} \to u \in \mathbb{R}) \to x \in \mathbb{R} \to y \in \mathbb{R} \to u \in \mathbb{R} \to v \in \mathbb{R} \to y + u = u + y$
\item $\varphi_{81} :\equiv (x \in \mathbb{R} \to y \in \mathbb{R} \to u \in \mathbb{R} \to v \in \mathbb{R} \to v \in \mathbb{R}) \to x \in \mathbb{R} \to y \in \mathbb{R} \to u \in \mathbb{R} \to v \in \mathbb{R} \to y + v \in \mathbb{R}$
\item $\varphi_{82} :\equiv (x \in \mathbb{R} \to y \in \mathbb{R} \to u \in \mathbb{R} \to v \in \mathbb{R} \to v \in \mathbb{R}) \to x \in \mathbb{R} \to y \in \mathbb{R} \to u \in \mathbb{R} \to v \in \mathbb{R} \to u + v \in \mathbb{R}$
\item $\varphi_{83} :\equiv (x \in \mathbb{R} \to y \in \mathbb{R} \to u \in \mathbb{R} \to v \in \mathbb{R} \to y + v \in \mathbb{R}) \to x \in \mathbb{R} \to y \in \mathbb{R} \to \varphi_{16}$
\item $\varphi_{84} :\equiv (x \in \mathbb{R} \to y \in \mathbb{R} \to u \in \mathbb{R} \to v \in \mathbb{R} \to u + v \in \mathbb{R}) \to x \in \mathbb{R} \to y \in \mathbb{R} \to \varphi_{21}$
\item $\varphi_{85} :\equiv (x + y + (u + v) = x + (y + (u + v)) \to \varphi_{48}) \to x \in \mathbb{R} \to \varphi_{51}$
\item $\varphi_{86} :\equiv (x + y + (u + v) = x + (u + y + v) \to \varphi_{44}) \to x \in \mathbb{R} \to y \in \mathbb{R} \to \varphi_{22}$
\item $\varphi_{87} :\equiv (x + y + (u + v) = x + (u + (y + v)) \to \varphi_{43}) \to x \in \mathbb{R} \to \varphi_{53}$
\item $\varphi_{88} :\equiv \forall v_{853}\,\forall v_{854}\,\forall v_{855}\,(v_{853} \in \mathbb{R} \to v_{854} \in \mathbb{R} \to \varphi_{15}) \to \varphi_{69}$
\item $\varphi_{89} :\equiv (x \in \mathbb{R} \to y \in \mathbb{R} \to u \in \mathbb{R} \to v \in \mathbb{R} \to y \in \mathbb{R}) \to (x \in \mathbb{R} \to y \in \mathbb{R} \to u \in \mathbb{R} \to v \in \mathbb{R} \to v \in \mathbb{R}) \to \varphi_{54}$
\item $\varphi_{90} :\equiv (x \in \mathbb{R} \to y \in \mathbb{R} \to u \in \mathbb{R} \to v \in \mathbb{R} \to u \in \mathbb{R}) \to (x \in \mathbb{R} \to y \in \mathbb{R} \to u \in \mathbb{R} \to v \in \mathbb{R} \to v \in \mathbb{R}) \to \varphi_{55}$
\item $\varphi_{91} :\equiv x + (y + u + v) = x + (y + u + v) \to (y + u = u + y \to \varphi_{47}) \to x \in \mathbb{R} \to \varphi_{50}$
\item $\varphi_{92} :\equiv x + (u + y + v) = x + (u + y + v) \to (u + y + v = u + (y + v) \to \varphi_{45}) \to x \in \mathbb{R} \to y \in \mathbb{R} \to \varphi_{18}$
\item $\varphi_{93} :\equiv x + (y + (u + v)) = x + (y + (u + v)) \to (y + (u + v) = y + u + v \to \varphi_{49}) \to x \in \mathbb{R} \to y \in \mathbb{R} \to \varphi_{20}$
\item $\varphi_{94} :\equiv (y \in \mathbb{R} \to u \in \mathbb{R} \to y + u = u + y) \to (x \in \mathbb{R} \to y \in \mathbb{R} \to u \in \mathbb{R} \to v \in \mathbb{R} \to y \in \mathbb{R}) \to \varphi_{80}$
\item $\varphi_{95} :\equiv (y \in \mathbb{R} \to v \in \mathbb{R} \to y + v \in \mathbb{R}) \to (x \in \mathbb{R} \to y \in \mathbb{R} \to u \in \mathbb{R} \to v \in \mathbb{R} \to y \in \mathbb{R}) \to \varphi_{81}$
\item $\varphi_{96} :\equiv (u \in \mathbb{R} \to v \in \mathbb{R} \to u + v \in \mathbb{R}) \to (x \in \mathbb{R} \to y \in \mathbb{R} \to u \in \mathbb{R} \to v \in \mathbb{R} \to u \in \mathbb{R}) \to \varphi_{82}$
\item $\varphi_{97} :\equiv (x \in \mathbb{R} \to y \in \mathbb{R} \to u \in \mathbb{R} \to v \in \mathbb{R} \to y + u = u + y) \to \varphi_{91}$
\item $\varphi_{98} :\equiv (x \in \mathbb{R} \to \varphi_{50}) \to (x + y + (u + v) = x + (y + u + v) \to \varphi_{46}) \to x \in \mathbb{R} \to \varphi_{52}$
\item $\varphi_{99} :\equiv (x \in \mathbb{R} \to y \in \mathbb{R} \to u \in \mathbb{R} \to v \in \mathbb{R} \to x \in \mathbb{R}) \to (x \in \mathbb{R} \to y \in \mathbb{R} \to u \in \mathbb{R} \to v \in \mathbb{R} \to u \in \mathbb{R}) \to \varphi_{83}$
\item $\varphi_{100} :\equiv (x \in \mathbb{R} \to y \in \mathbb{R} \to u \in \mathbb{R} \to v \in \mathbb{R} \to x \in \mathbb{R}) \to (x \in \mathbb{R} \to y \in \mathbb{R} \to u \in \mathbb{R} \to v \in \mathbb{R} \to y \in \mathbb{R}) \to \varphi_{84}$
\item $\varphi_{101} :\equiv (u \in \mathbb{R} \to y \in \mathbb{R} \to v \in \mathbb{R} \to u + y + v = u + (y + v)) \to (x \in \mathbb{R} \to y \in \mathbb{R} \to u \in \mathbb{R} \to v \in \mathbb{R} \to u \in \mathbb{R}) \to \varphi_{89}$
\item $\varphi_{102} :\equiv (y \in \mathbb{R} \to u \in \mathbb{R} \to v \in \mathbb{R} \to y + u + v = y + (u + v)) \to (x \in \mathbb{R} \to y \in \mathbb{R} \to u \in \mathbb{R} \to v \in \mathbb{R} \to y \in \mathbb{R}) \to \varphi_{90}$
\end{abbrevlist}
\end{fullproof}

% 自動生成 (prover/build.js): NegOneMul  ノード数 79
\begin{fullproof}{[E50]}{$\forall x \in \mathbb{R}\,(-1 \cdot x = -x)$}{79}\phantomsection\label{tree:NegOneMul}
\noindent$\Pi_{1}$:
\begin{ptree}
\AXC{$y = z \to x = y \to x = z$\lbl{A6}}
\AXC{$(y = z \to x = y \to x = z) \to x = y \to y = z \to x = z$\lbl{A1}}
\BIC{$x = y \to y = z \to x = z$}
\UIC{$\forall z\,(x = y \to y = z \to x = z)$}
\UIC{$\forall y\,\forall z\,(x = y \to y = z \to x = z)$}
\UIC{$\forall x\,\forall y\,\forall z\,(x = y \to y = z \to x = z)$}
\AXC{$\varphi_{6}$\lbl{A2}}
\BIC{$\forall y\,\forall z\,(-1 \cdot x = y \to y = z \to -1 \cdot x = z)$}
\end{ptree}
\noindent$\Pi_{2}$:
\begin{ptree}
\AXC{$\forall x \in \mathbb{R}\,(x \cdot 1 = x)$\lbl{(R8)}}
\AXC{$\forall x \in \mathbb{R}\,(x \cdot 1 = x) \to x \in \mathbb{R} \to x \cdot 1 = x$\lbl{A2}}
\BIC{$x \in \mathbb{R} \to x \cdot 1 = x$}
\end{ptree}
\noindent$\Pi_{3}$:
\begin{ptree}
\AXC{$\stackrel{\Pi_{2}}{\vdots}$}\noLine
\UIC{$x \in \mathbb{R} \to x \cdot 1 = x$}
\AXC{$(x \in \mathbb{R} \to x \cdot 1 = x) \to -(x \cdot 1) = -(x \cdot 1) \to \varphi_{10}$\lbl{A1}}
\BIC{$-(x \cdot 1) = -(x \cdot 1) \to \varphi_{10}$}
\end{ptree}
\noindent$\Pi_{4}$:
\begin{ptree}
\AXC{$x = x$\lbl{A5}}
\UIC{$\forall x\,(x = x)$}
\AXC{$\forall x\,(x = x) \to -(x \cdot 1) = -(x \cdot 1)$\lbl{A2}}
\BIC{$-(x \cdot 1) = -(x \cdot 1)$}
\AXC{$\stackrel{\Pi_{3}}{\vdots}$}\noLine
\UIC{$-(x \cdot 1) = -(x \cdot 1) \to \varphi_{10}$}
\BIC{$\varphi_{10}$}
\end{ptree}
\noindent$\Pi_{5}$:
\begin{ptree}
\AXC{$y = z \to x = y \to x = z$\lbl{A6}}
\AXC{$(y = z \to x = y \to x = z) \to x = y \to y = z \to x = z$\lbl{A1}}
\BIC{$x = y \to y = z \to x = z$}
\UIC{$\forall z\,(x = y \to y = z \to x = z)$}
\UIC{$\forall y\,\forall z\,(x = y \to y = z \to x = z)$}
\UIC{$\forall x\,\forall y\,\forall z\,(x = y \to y = z \to x = z)$}
\AXC{$\varphi_{6}$\lbl{A2}}
\BIC{$\forall y\,\forall z\,(-1 \cdot x = y \to y = z \to -1 \cdot x = z)$}
\end{ptree}
\noindent$\Pi_{6}$:
\begin{ptree}
\AXC{$\stackrel{\Pi_{5}}{\vdots}$}\noLine
\UIC{$\forall y\,\forall z\,(-1 \cdot x = y \to y = z \to -1 \cdot x = z)$}
\AXC{$\forall y\,\forall z\,(-1 \cdot x = y \to y = z \to -1 \cdot x = z) \to \varphi_{3}$\lbl{A2}}
\BIC{$\varphi_{3}$}
\AXC{$\varphi_{3} \to \varphi_{7}$\lbl{A2}}
\BIC{$\varphi_{7}$}
\end{ptree}
\noindent$\Pi_{7}$:
\begin{ptree}
\AXC{$z = w \to -(1 \cdot x) = -z \to -(1 \cdot x) = -w$\lbl{A6}}
\UIC{$\forall w\,(z = w \to -(1 \cdot x) = -z \to -(1 \cdot x) = -w)$}
\UIC{$\forall z\,\forall w\,(z = w \to -(1 \cdot x) = -z \to -(1 \cdot x) = -w)$}
\AXC{$\varphi_{17}$\lbl{A2}}
\BIC{$\varphi_{4}$}
\AXC{$\varphi_{4} \to \varphi_{5}$\lbl{A2}}
\BIC{$\varphi_{5}$}
\end{ptree}
\noindent$\Pi_{8}$:
\begin{ptree}
\AXC{$1 \in \mathbb{R}$\lbl{(R1)}}
\AXC{$\forall x\,\forall y\,(x \in \mathbb{R} \to y \in \mathbb{R} \to x \cdot y = y \cdot x)$\lbl{(R7)}}
\AXC{$\varphi_{11}$\lbl{A2}}
\BIC{$\forall y\,(1 \in \mathbb{R} \to y \in \mathbb{R} \to 1 \cdot y = y \cdot 1)$}
\AXC{$\varphi_{9}$\lbl{A2}}
\BIC{$1 \in \mathbb{R} \to x \in \mathbb{R} \to 1 \cdot x = x \cdot 1$}
\AXC{$\varphi_{14}$\lbl{A1}}
\BIC{$1 \in \mathbb{R} \to (x \in \mathbb{R} \to x \in \mathbb{R}) \to x \in \mathbb{R} \to 1 \cdot x = x \cdot 1$}
\BIC{$(x \in \mathbb{R} \to x \in \mathbb{R}) \to x \in \mathbb{R} \to 1 \cdot x = x \cdot 1$}
\end{ptree}
\noindent$\Pi_{9}$:
\begin{ptree}
\AXC{$x \in \mathbb{R} \to x \in \mathbb{R}$\lbl{A1}}
\AXC{$\stackrel{\Pi_{8}}{\vdots}$}\noLine
\UIC{$(x \in \mathbb{R} \to x \in \mathbb{R}) \to x \in \mathbb{R} \to 1 \cdot x = x \cdot 1$}
\BIC{$x \in \mathbb{R} \to 1 \cdot x = x \cdot 1$}
\AXC{$(x \in \mathbb{R} \to 1 \cdot x = x \cdot 1) \to \varphi_{20}$\lbl{A1}}
\BIC{$\varphi_{20}$}
\end{ptree}
\noindent$\Pi_{10}$:
\begin{ptree}
\AXC{$\forall x\,\forall y\,(x \in \mathbb{R} \to y \in \mathbb{R} \to -x \cdot y = -(x \cdot y))$\lbl{[F9]}}
\AXC{$\varphi_{15}$\lbl{A2}}
\BIC{$\forall y\,(1 \in \mathbb{R} \to y \in \mathbb{R} \to -1 \cdot y = -(1 \cdot y))$}
\AXC{$\varphi_{13}$\lbl{A2}}
\BIC{$1 \in \mathbb{R} \to x \in \mathbb{R} \to -1 \cdot x = -(1 \cdot x)$}
\end{ptree}
\noindent$\Pi_{11}$:
\begin{ptree}
\AXC{$1 \in \mathbb{R}$\lbl{(R1)}}
\AXC{$\stackrel{\Pi_{10}}{\vdots}$}\noLine
\UIC{$1 \in \mathbb{R} \to x \in \mathbb{R} \to -1 \cdot x = -(1 \cdot x)$}
\AXC{$(1 \in \mathbb{R} \to x \in \mathbb{R} \to -1 \cdot x = -(1 \cdot x)) \to \varphi_{8}$\lbl{A1}}
\BIC{$\varphi_{8}$}
\BIC{$(x \in \mathbb{R} \to x \in \mathbb{R}) \to x \in \mathbb{R} \to -1 \cdot x = -(1 \cdot x)$}
\end{ptree}
\noindent$\Pi_{12}$:
\begin{ptree}
\AXC{$x \in \mathbb{R} \to x \in \mathbb{R}$\lbl{A1}}
\AXC{$\stackrel{\Pi_{11}}{\vdots}$}\noLine
\UIC{$(x \in \mathbb{R} \to x \in \mathbb{R}) \to x \in \mathbb{R} \to -1 \cdot x = -(1 \cdot x)$}
\BIC{$x \in \mathbb{R} \to -1 \cdot x = -(1 \cdot x)$}
\AXC{$(x \in \mathbb{R} \to -1 \cdot x = -(1 \cdot x)) \to \varphi_{21}$\lbl{A1}}
\BIC{$\varphi_{21}$}
\end{ptree}
\noindent$\Pi_{13}$:
\begin{ptree}
\AXC{$\stackrel{\Pi_{6}}{\vdots}$}\noLine
\UIC{$\varphi_{7}$}
\AXC{$\stackrel{\Pi_{7}}{\vdots}$}\noLine
\UIC{$\varphi_{5}$}
\AXC{$x = x$\lbl{A5}}
\UIC{$\forall x\,(x = x)$}
\AXC{$\forall x\,(x = x) \to -(1 \cdot x) = -(1 \cdot x)$\lbl{A2}}
\BIC{$-(1 \cdot x) = -(1 \cdot x)$}
\AXC{$\stackrel{\Pi_{9}}{\vdots}$}\noLine
\UIC{$\varphi_{20}$}
\BIC{$\varphi_{5} \to x \in \mathbb{R} \to -(1 \cdot x) = -(x \cdot 1)$}
\BIC{$x \in \mathbb{R} \to -(1 \cdot x) = -(x \cdot 1)$}
\AXC{$\stackrel{\Pi_{12}}{\vdots}$}\noLine
\UIC{$\varphi_{21}$}
\BIC{$\varphi_{7} \to x \in \mathbb{R} \to -1 \cdot x = -(x \cdot 1)$}
\BIC{$x \in \mathbb{R} \to -1 \cdot x = -(x \cdot 1)$}
\AXC{$\varphi_{22}$\lbl{A1}}
\BIC{$(x \in \mathbb{R} \to -(x \cdot 1) = -x) \to \varphi_{12}$}
\end{ptree}
\noindent$\Pi_{14}$:
\begin{ptree}
\AXC{$z = w \to -(x \cdot 1) = -z \to -(x \cdot 1) = -w$\lbl{A6}}
\UIC{$\forall w\,(z = w \to -(x \cdot 1) = -z \to -(x \cdot 1) = -w)$}
\UIC{$\forall z\,\forall w\,(z = w \to -(x \cdot 1) = -z \to -(x \cdot 1) = -w)$}
\AXC{$\varphi_{16}$\lbl{A2}}
\BIC{$\varphi_{2}$}
\AXC{$\varphi_{18}$\lbl{A2}}
\BIC{$x \cdot 1 = x \to -(x \cdot 1) = -(x \cdot 1) \to -(x \cdot 1) = -x$}
\AXC{$\stackrel{\Pi_{4}}{\vdots}$}\noLine
\UIC{$\varphi_{10}$}
\BIC{$x \in \mathbb{R} \to -(x \cdot 1) = -x$}
\AXC{$\stackrel{\Pi_{13}}{\vdots}$}\noLine
\UIC{$(x \in \mathbb{R} \to -(x \cdot 1) = -x) \to \varphi_{12}$}
\BIC{$\varphi_{12}$}
\end{ptree}
\noindent 主木:
\begin{ptree}
\AXC{$\stackrel{\Pi_{1}}{\vdots}$}\noLine
\UIC{$\forall y\,\forall z\,(-1 \cdot x = y \to y = z \to -1 \cdot x = z)$}
\AXC{$\forall y\,\forall z\,(-1 \cdot x = y \to y = z \to -1 \cdot x = z) \to \varphi_{1}$\lbl{A2}}
\BIC{$\varphi_{1}$}
\AXC{$\varphi_{19}$\lbl{A2}}
\BIC{$-1 \cdot x = -(x \cdot 1) \to -(x \cdot 1) = -x \to -1 \cdot x = -x$}
\AXC{$\stackrel{\Pi_{14}}{\vdots}$}\noLine
\UIC{$\varphi_{12}$}
\BIC{$x \in \mathbb{R} \to -1 \cdot x = -x$}
\UIC{$\forall x \in \mathbb{R}\,(-1 \cdot x = -x)$}
\end{ptree}
\begin{abbrevlist}
\item $\varphi_{1} :\equiv \forall z\,(-1 \cdot x = -(x \cdot 1) \to -(x \cdot 1) = z \to -1 \cdot x = z)$
\item $\varphi_{2} :\equiv \forall w\,(x \cdot 1 = w \to -(x \cdot 1) = -(x \cdot 1) \to -(x \cdot 1) = -w)$
\item $\varphi_{3} :\equiv \forall z\,(-1 \cdot x = -(1 \cdot x) \to -(1 \cdot x) = z \to -1 \cdot x = z)$
\item $\varphi_{4} :\equiv \forall w\,(1 \cdot x = w \to -(1 \cdot x) = -(1 \cdot x) \to -(1 \cdot x) = -w)$
\item $\varphi_{5} :\equiv 1 \cdot x = x \cdot 1 \to -(1 \cdot x) = -(1 \cdot x) \to -(1 \cdot x) = -(x \cdot 1)$
\item $\varphi_{6} :\equiv \forall x\,\forall y\,\forall z\,(x = y \to y = z \to x = z) \to \forall y\,\forall z\,(-1 \cdot x = y \to y = z \to -1 \cdot x = z)$
\item $\varphi_{7} :\equiv -1 \cdot x = -(1 \cdot x) \to -(1 \cdot x) = -(x \cdot 1) \to -1 \cdot x = -(x \cdot 1)$
\item $\varphi_{8} :\equiv 1 \in \mathbb{R} \to (x \in \mathbb{R} \to x \in \mathbb{R}) \to x \in \mathbb{R} \to -1 \cdot x = -(1 \cdot x)$
\item $\varphi_{9} :\equiv \forall y\,(1 \in \mathbb{R} \to y \in \mathbb{R} \to 1 \cdot y = y \cdot 1) \to 1 \in \mathbb{R} \to x \in \mathbb{R} \to 1 \cdot x = x \cdot 1$
\item $\varphi_{10} :\equiv (x \cdot 1 = x \to -(x \cdot 1) = -(x \cdot 1) \to -(x \cdot 1) = -x) \to x \in \mathbb{R} \to -(x \cdot 1) = -x$
\item $\varphi_{11} :\equiv \forall x\,\forall y\,(x \in \mathbb{R} \to y \in \mathbb{R} \to x \cdot y = y \cdot x) \to \forall y\,(1 \in \mathbb{R} \to y \in \mathbb{R} \to 1 \cdot y = y \cdot 1)$
\item $\varphi_{12} :\equiv (-1 \cdot x = -(x \cdot 1) \to -(x \cdot 1) = -x \to -1 \cdot x = -x) \to x \in \mathbb{R} \to -1 \cdot x = -x$
\item $\varphi_{13} :\equiv \forall y\,(1 \in \mathbb{R} \to y \in \mathbb{R} \to -1 \cdot y = -(1 \cdot y)) \to 1 \in \mathbb{R} \to x \in \mathbb{R} \to -1 \cdot x = -(1 \cdot x)$
\item $\varphi_{14} :\equiv (1 \in \mathbb{R} \to x \in \mathbb{R} \to 1 \cdot x = x \cdot 1) \to 1 \in \mathbb{R} \to (x \in \mathbb{R} \to x \in \mathbb{R}) \to x \in \mathbb{R} \to 1 \cdot x = x \cdot 1$
\item $\varphi_{15} :\equiv \forall x\,\forall y\,(x \in \mathbb{R} \to y \in \mathbb{R} \to -x \cdot y = -(x \cdot y)) \to \forall y\,(1 \in \mathbb{R} \to y \in \mathbb{R} \to -1 \cdot y = -(1 \cdot y))$
\item $\varphi_{16} :\equiv \forall z\,\forall w\,(z = w \to -(x \cdot 1) = -z \to -(x \cdot 1) = -w) \to \varphi_{2}$
\item $\varphi_{17} :\equiv \forall z\,\forall w\,(z = w \to -(1 \cdot x) = -z \to -(1 \cdot x) = -w) \to \varphi_{4}$
\item $\varphi_{18} :\equiv \varphi_{2} \to x \cdot 1 = x \to -(x \cdot 1) = -(x \cdot 1) \to -(x \cdot 1) = -x$
\item $\varphi_{19} :\equiv \varphi_{1} \to -1 \cdot x = -(x \cdot 1) \to -(x \cdot 1) = -x \to -1 \cdot x = -x$
\item $\varphi_{20} :\equiv -(1 \cdot x) = -(1 \cdot x) \to \varphi_{5} \to x \in \mathbb{R} \to -(1 \cdot x) = -(x \cdot 1)$
\item $\varphi_{21} :\equiv (x \in \mathbb{R} \to -(1 \cdot x) = -(x \cdot 1)) \to \varphi_{7} \to x \in \mathbb{R} \to -1 \cdot x = -(x \cdot 1)$
\item $\varphi_{22} :\equiv (x \in \mathbb{R} \to -1 \cdot x = -(x \cdot 1)) \to (x \in \mathbb{R} \to -(x \cdot 1) = -x) \to \varphi_{12}$
\end{abbrevlist}
\end{fullproof}

% 自動生成 (prover/build.js): ShiftConv  ノード数 370
\begin{fullproof}{[C29]}{$\forall a\,\forall \alpha\,(\mathrm{Seq}(a) \to \mathrm{Conv}(a \circ \mathsf{s},\alpha) \to \mathrm{Conv}(a,\alpha))$}{370}\phantomsection\label{tree:ShiftConv}
\noindent$\Pi_{1}$:
\begin{ptree}
\AXC{$z = w \to \lvert z - \alpha \rvert < \varepsilon \to \lvert w - \alpha \rvert < \varepsilon$\lbl{A6}}
\UIC{$\forall w\,(z = w \to \lvert z - \alpha \rvert < \varepsilon \to \lvert w - \alpha \rvert < \varepsilon)$}
\UIC{$\forall z\,\forall w\,(z = w \to \lvert z - \alpha \rvert < \varepsilon \to \lvert w - \alpha \rvert < \varepsilon)$}
\AXC{$\forall z\,\forall w\,(z = w \to \lvert z - \alpha \rvert < \varepsilon \to \lvert w - \alpha \rvert < \varepsilon) \to \varphi_{62}$\lbl{A2}}
\BIC{$\varphi_{62}$}
\AXC{$\varphi_{62} \to \varphi_{71}$\lbl{A2}}
\BIC{$\varphi_{71}$}
\end{ptree}
\noindent$\Pi_{2}$:
\begin{ptree}
\AXC{$\forall n \in \mathbb{N}\,(n \in \omega)$\lbl{[N7']}}
\AXC{$\forall n \in \mathbb{N}\,(n \in \omega) \to N \in \mathbb{N} \to N \in \omega$\lbl{A2}}
\BIC{$N \in \mathbb{N} \to N \in \omega$}
\AXC{$(N \in \mathbb{N} \to N \in \omega) \to \varphi_{149}$\lbl{A1}}
\BIC{$\varphi_{149}$}
\end{ptree}
\noindent$\Pi_{3}$:
\begin{ptree}
\AXC{$\mathrm{Seq}(a) \to \mathrm{Conv}(a \circ \mathsf{s},\alpha) \to \varepsilon \in \mathbb{R} \to 0 < \varepsilon \to \varphi_{109}$\lbl{A1}}
\AXC{$\varphi_{164}$\lbl{A1}}
\BIC{$\mathrm{Seq}(a) \to \mathrm{Conv}(a \circ \mathsf{s},\alpha) \to \varepsilon \in \mathbb{R} \to 0 < \varepsilon \to \varphi_{55}$}
\AXC{$\stackrel{\Pi_{2}}{\vdots}$}\noLine
\UIC{$\varphi_{149}$}
\BIC{$\mathrm{Seq}(a) \to \mathrm{Conv}(a \circ \mathsf{s},\alpha) \to \varepsilon \in \mathbb{R} \to 0 < \varepsilon \to \varphi_{59}$}
\end{ptree}
\noindent$\Pi_{4}$:
\begin{ptree}
\AXC{$\forall n \in \omega\,(n < S(n))$\lbl{[N17]}}
\AXC{$\forall n \in \omega\,(n < S(n)) \to N \in \omega \to N < S(N)$\lbl{A2}}
\BIC{$N \in \omega \to N < S(N)$}
\AXC{$(N \in \omega \to N < S(N)) \to \varphi_{150}$\lbl{A1}}
\BIC{$\varphi_{150}$}
\end{ptree}
\noindent$\Pi_{5}$:
\begin{ptree}
\AXC{$\stackrel{\Pi_{3}}{\vdots}$}\noLine
\UIC{$\mathrm{Seq}(a) \to \mathrm{Conv}(a \circ \mathsf{s},\alpha) \to \varepsilon \in \mathbb{R} \to 0 < \varepsilon \to \varphi_{59}$}
\AXC{$\stackrel{\Pi_{4}}{\vdots}$}\noLine
\UIC{$\varphi_{150}$}
\BIC{$\mathrm{Seq}(a) \to \mathrm{Conv}(a \circ \mathsf{s},\alpha) \to \varepsilon \in \mathbb{R} \to 0 < \varepsilon \to \varphi_{56}$}
\end{ptree}
\noindent$\Pi_{6}$:
\begin{ptree}
\AXC{$z = w \to S(N) < z \to S(N) < w$\lbl{A6}}
\UIC{$\forall w\,(z = w \to S(N) < z \to S(N) < w)$}
\UIC{$\forall z\,\forall w\,(z = w \to S(N) < z \to S(N) < w)$}
\AXC{$\varphi_{42}$\lbl{A2}}
\BIC{$\forall w\,(S(m) = w \to S(N) < S(m) \to S(N) < w)$}
\AXC{$\varphi_{43}$\lbl{A2}}
\BIC{$S(m) = N \to S(N) < S(m) \to S(N) < N$}
\end{ptree}
\noindent$\Pi_{7}$:
\begin{ptree}
\AXC{$z = w \to S(N) < z \to S(N) < w$\lbl{A6}}
\UIC{$\forall w\,(z = w \to S(N) < z \to S(N) < w)$}
\UIC{$\forall z\,\forall w\,(z = w \to S(N) < z \to S(N) < w)$}
\AXC{$\varphi_{29}$\lbl{A2}}
\BIC{$\forall w\,(n = w \to S(N) < n \to S(N) < w)$}
\AXC{$\varphi_{30}$\lbl{A2}}
\BIC{$n = S(m) \to S(N) < n \to S(N) < S(m)$}
\end{ptree}
\noindent$\Pi_{8}$:
\begin{ptree}
\AXC{$\mathrm{Seq}(a) \to \mathrm{Conv}(a \circ \mathsf{s},\alpha) \to \varepsilon \in \mathbb{R} \to \varphi_{111}$\lbl{A1}}
\AXC{$\varphi_{173}$\lbl{A1}}
\BIC{$\mathrm{Seq}(a) \to \mathrm{Conv}(a \circ \mathsf{s},\alpha) \to \varepsilon \in \mathbb{R} \to 0 < \varepsilon \to \varphi_{89}$}
\AXC{$\varphi_{208}$\lbl{A1}}
\BIC{$\varphi_{184}$}
\end{ptree}
\noindent$\Pi_{9}$:
\begin{ptree}
\AXC{$\stackrel{\Pi_{7}}{\vdots}$}\noLine
\UIC{$n = S(m) \to S(N) < n \to S(N) < S(m)$}
\AXC{$\mathrm{Seq}(a) \to \mathrm{Conv}(a \circ \mathsf{s},\alpha) \to \varepsilon \in \mathbb{R} \to 0 < \varepsilon \to \varphi_{72}$\lbl{A1}}
\AXC{$\stackrel{\Pi_{8}}{\vdots}$}\noLine
\UIC{$\varphi_{184}$}
\BIC{$(n = S(m) \to S(N) < n \to S(N) < S(m)) \to \varphi_{134}$}
\BIC{$\varphi_{134}$}
\end{ptree}
\noindent$\Pi_{10}$:
\begin{ptree}
\AXC{$z = w \to S(N) = z \to S(N) = w$\lbl{A6}}
\UIC{$\forall w\,(z = w \to S(N) = z \to S(N) = w)$}
\UIC{$\forall z\,\forall w\,(z = w \to S(N) = z \to S(N) = w)$}
\AXC{$\varphi_{14}$\lbl{A2}}
\BIC{$\forall w\,(n = w \to S(N) = n \to S(N) = w)$}
\AXC{$\varphi_{15}$\lbl{A2}}
\BIC{$n = S(m) \to S(N) = n \to S(N) = S(m)$}
\AXC{$\stackrel{\Pi_{9}}{\vdots}$}\noLine
\UIC{$\varphi_{134}$}
\BIC{$\mathrm{Seq}(a) \to \mathrm{Conv}(a \circ \mathsf{s},\alpha) \to \varepsilon \in \mathbb{R} \to \varphi_{102}$}
\end{ptree}
\noindent$\Pi_{11}$:
\begin{ptree}
\AXC{$x = x$\lbl{A5}}
\AXC{$x = y \to x = x \to y = x$\lbl{A6}}
\AXC{$(x = y \to x = x \to y = x) \to x = x \to x = y \to y = x$\lbl{A1}}
\BIC{$x = x \to x = y \to y = x$}
\BIC{$x = y \to y = x$}
\UIC{$\forall y\,(x = y \to y = x)$}
\UIC{$\forall x\,\forall y\,(x = y \to y = x)$}
\end{ptree}
\noindent$\Pi_{12}$:
\begin{ptree}
\AXC{$\stackrel{\Pi_{11}}{\vdots}$}\noLine
\UIC{$\forall x\,\forall y\,(x = y \to y = x)$}
\AXC{$\forall x\,\forall y\,(x = y \to y = x) \to \forall y\,(N = y \to y = N)$\lbl{A2}}
\BIC{$\forall y\,(N = y \to y = N)$}
\end{ptree}
\noindent$\Pi_{13}$:
\begin{ptree}
\AXC{$\stackrel{\Pi_{12}}{\vdots}$}\noLine
\UIC{$\forall y\,(N = y \to y = N)$}
\AXC{$\forall y\,(N = y \to y = N) \to N = S(m) \to S(m) = N$\lbl{A2}}
\BIC{$N = S(m) \to S(m) = N$}
\end{ptree}
\noindent$\Pi_{14}$:
\begin{ptree}
\AXC{$\stackrel{\Pi_{13}}{\vdots}$}\noLine
\UIC{$N = S(m) \to S(m) = N$}
\AXC{$\mathrm{Seq}(a) \to \mathrm{Conv}(a \circ \mathsf{s},\alpha) \to \varepsilon \in \mathbb{R} \to 0 < \varepsilon \to \varphi_{98}$\lbl{A1}}
\AXC{$\varphi_{182}$\lbl{A1}}
\BIC{$\varphi_{132}$}
\BIC{$\mathrm{Seq}(a) \to \mathrm{Conv}(a \circ \mathsf{s},\alpha) \to \varepsilon \in \mathbb{R} \to 0 < \varepsilon \to \varphi_{99}$}
\AXC{$\varphi_{211}$\lbl{A1}}
\BIC{$\varphi_{187}$}
\end{ptree}
\noindent$\Pi_{15}$:
\begin{ptree}
\AXC{$\stackrel{\Pi_{6}}{\vdots}$}\noLine
\UIC{$S(m) = N \to S(N) < S(m) \to S(N) < N$}
\AXC{$\stackrel{\Pi_{10}}{\vdots}$}\noLine
\UIC{$\mathrm{Seq}(a) \to \mathrm{Conv}(a \circ \mathsf{s},\alpha) \to \varepsilon \in \mathbb{R} \to \varphi_{102}$}
\AXC{$\stackrel{\Pi_{14}}{\vdots}$}\noLine
\UIC{$\varphi_{187}$}
\BIC{$(S(m) = N \to S(N) < S(m) \to S(N) < N) \to \varphi_{135}$}
\BIC{$\varphi_{135}$}
\end{ptree}
\noindent$\Pi_{16}$:
\begin{ptree}
\AXC{$\forall x\,\forall y\,(x \leq y \to \neg y < x)$\lbl{[O7]}}
\AXC{$\varphi_{1}$\lbl{A2}}
\BIC{$\forall y\,(S(N) \leq y \to \neg y < S(N))$}
\AXC{$\varphi_{2}$\lbl{A2}}
\BIC{$S(N) \leq N \to \neg N < S(N)$}
\AXC{$(S(N) \leq N \to \neg N < S(N)) \to \varphi_{178}$\lbl{A1}}
\BIC{$\varphi_{178}$}
\end{ptree}
\noindent$\Pi_{17}$:
\begin{ptree}
\AXC{$z = w \to S(N) = z \to S(N) = w$\lbl{A6}}
\UIC{$\forall w\,(z = w \to S(N) = z \to S(N) = w)$}
\UIC{$\forall z\,\forall w\,(z = w \to S(N) = z \to S(N) = w)$}
\AXC{$\varphi_{22}$\lbl{A2}}
\BIC{$\forall w\,(S(m) = w \to S(N) = S(m) \to S(N) = w)$}
\AXC{$\varphi_{23}$\lbl{A2}}
\BIC{$S(m) = N \to S(N) = S(m) \to S(N) = N$}
\AXC{$\stackrel{\Pi_{15}}{\vdots}$}\noLine
\UIC{$\varphi_{135}$}
\BIC{$\mathrm{Seq}(a) \to \mathrm{Conv}(a \circ \mathsf{s},\alpha) \to \varepsilon \in \mathbb{R} \to 0 < \varepsilon \to \varphi_{100}$}
\AXC{$\stackrel{\Pi_{16}}{\vdots}$}\noLine
\UIC{$\varphi_{178}$}
\BIC{$\mathrm{Seq}(a) \to \mathrm{Conv}(a \circ \mathsf{s},\alpha) \to \varepsilon \in \mathbb{R} \to 0 < \varepsilon \to \varphi_{101}$}
\AXC{$\varphi_{197}$\lbl{A1}}
\BIC{$\varphi_{160}$}
\end{ptree}
\noindent$\Pi_{18}$:
\begin{ptree}
\AXC{$\forall x\,\forall y\,\forall z\,(x \leq y \to y \leq z \to x \leq z)$\lbl{[O5]}}
\AXC{$\varphi_{9}$\lbl{A2}}
\BIC{$\forall y\,\forall z\,(N \leq y \to y \leq z \to N \leq z)$}
\AXC{$\varphi_{10}$\lbl{A2}}
\BIC{$\forall z\,(N \leq S(N) \to S(N) \leq z \to N \leq z)$}
\AXC{$\varphi_{16}$\lbl{A2}}
\BIC{$N \leq S(N) \to S(N) \leq S(m) \to N \leq S(m)$}
\AXC{$\varphi_{183}$\lbl{A1}}
\BIC{$N \leq S(N) \to \varphi_{168}$}
\end{ptree}
\noindent$\Pi_{19}$:
\begin{ptree}
\AXC{$\forall n\,(n \leq S(n))$\lbl{[N8]}}
\AXC{$\forall n\,(n \leq S(n)) \to N \leq S(N)$\lbl{A2}}
\BIC{$N \leq S(N)$}
\AXC{$\stackrel{\Pi_{18}}{\vdots}$}\noLine
\UIC{$N \leq S(N) \to \varphi_{168}$}
\BIC{$\varphi_{168}$}
\end{ptree}
\noindent$\Pi_{20}$:
\begin{ptree}
\AXC{$\forall n \in \mathbb{N}\,(n \in \omega)$\lbl{[N7']}}
\AXC{$\forall n \in \mathbb{N}\,(n \in \omega) \to m \in \mathbb{N} \to m \in \omega$\lbl{A2}}
\BIC{$m \in \mathbb{N} \to m \in \omega$}
\AXC{$(m \in \mathbb{N} \to m \in \omega) \to \varphi_{172}$\lbl{A1}}
\BIC{$\varphi_{172}$}
\end{ptree}
\noindent$\Pi_{21}$:
\begin{ptree}
\AXC{$\mathrm{Seq}(a) \to \mathrm{Conv}(a \circ \mathsf{s},\alpha) \to \varepsilon \in \mathbb{R} \to \varphi_{111}$\lbl{A1}}
\AXC{$\varphi_{175}$\lbl{A1}}
\BIC{$\mathrm{Seq}(a) \to \mathrm{Conv}(a \circ \mathsf{s},\alpha) \to \varepsilon \in \mathbb{R} \to 0 < \varepsilon \to \varphi_{95}$}
\AXC{$\stackrel{\Pi_{20}}{\vdots}$}\noLine
\UIC{$\varphi_{172}$}
\BIC{$\mathrm{Seq}(a) \to \mathrm{Conv}(a \circ \mathsf{s},\alpha) \to \varepsilon \in \mathbb{R} \to 0 < \varepsilon \to \varphi_{96}$}
\end{ptree}
\noindent$\Pi_{22}$:
\begin{ptree}
\AXC{$\varphi_{19}$\lbl{[N9]}}
\AXC{$\varphi_{82}$\lbl{A2}}
\BIC{$\forall n\,(N \in \omega \to n \in \omega \to N \leq S(n) \to N \leq n \lor N = S(n))$}
\AXC{$\varphi_{78}$\lbl{A2}}
\BIC{$N \in \omega \to m \in \omega \to N \leq S(m) \to N \leq m \lor N = S(m)$}
\AXC{$(N \in \omega \to m \in \omega \to N \leq S(m) \to N \leq m \lor N = S(m)) \to \varphi_{213}$\lbl{A1}}
\BIC{$\varphi_{213}$}
\end{ptree}
\noindent$\Pi_{23}$:
\begin{ptree}
\AXC{$\stackrel{\Pi_{21}}{\vdots}$}\noLine
\UIC{$\mathrm{Seq}(a) \to \mathrm{Conv}(a \circ \mathsf{s},\alpha) \to \varepsilon \in \mathbb{R} \to 0 < \varepsilon \to \varphi_{96}$}
\AXC{$\stackrel{\Pi_{3}}{\vdots}$}\noLine
\UIC{$\mathrm{Seq}(a) \to \mathrm{Conv}(a \circ \mathsf{s},\alpha) \to \varepsilon \in \mathbb{R} \to 0 < \varepsilon \to \varphi_{59}$}
\AXC{$\stackrel{\Pi_{22}}{\vdots}$}\noLine
\UIC{$\varphi_{213}$}
\BIC{$\varphi_{202}$}
\BIC{$\varphi_{170}$}
\end{ptree}
\noindent$\Pi_{24}$:
\begin{ptree}
\AXC{$\stackrel{\Pi_{10}}{\vdots}$}\noLine
\UIC{$\mathrm{Seq}(a) \to \mathrm{Conv}(a \circ \mathsf{s},\alpha) \to \varepsilon \in \mathbb{R} \to \varphi_{102}$}
\AXC{$\stackrel{\Pi_{19}}{\vdots}$}\noLine
\UIC{$\varphi_{168}$}
\BIC{$\mathrm{Seq}(a) \to \mathrm{Conv}(a \circ \mathsf{s},\alpha) \to \varepsilon \in \mathbb{R} \to 0 < \varepsilon \to \varphi_{90}$}
\AXC{$\stackrel{\Pi_{23}}{\vdots}$}\noLine
\UIC{$\varphi_{170}$}
\BIC{$\mathrm{Seq}(a) \to \mathrm{Conv}(a \circ \mathsf{s},\alpha) \to \varepsilon \in \mathbb{R} \to \varphi_{106}$}
\end{ptree}
\noindent$\Pi_{25}$:
\begin{ptree}
\AXC{$\stackrel{\Pi_{24}}{\vdots}$}\noLine
\UIC{$\mathrm{Seq}(a) \to \mathrm{Conv}(a \circ \mathsf{s},\alpha) \to \varepsilon \in \mathbb{R} \to \varphi_{106}$}
\AXC{$(\mathrm{Seq}(a) \to \mathrm{Conv}(a \circ \mathsf{s},\alpha) \to \varepsilon \in \mathbb{R} \to \varphi_{106}) \to \varphi_{169}$\lbl{A1}}
\BIC{$\varphi_{169}$}
\end{ptree}
\noindent$\Pi_{26}$:
\begin{ptree}
\AXC{$\stackrel{\Pi_{5}}{\vdots}$}\noLine
\UIC{$\mathrm{Seq}(a) \to \mathrm{Conv}(a \circ \mathsf{s},\alpha) \to \varepsilon \in \mathbb{R} \to 0 < \varepsilon \to \varphi_{56}$}
\AXC{$\stackrel{\Pi_{17}}{\vdots}$}\noLine
\UIC{$\varphi_{160}$}
\BIC{$\mathrm{Seq}(a) \to \mathrm{Conv}(a \circ \mathsf{s},\alpha) \to \varepsilon \in \mathbb{R} \to \varphi_{105}$}
\AXC{$\stackrel{\Pi_{25}}{\vdots}$}\noLine
\UIC{$\varphi_{169}$}
\BIC{$\mathrm{Seq}(a) \to \mathrm{Conv}(a \circ \mathsf{s},\alpha) \to \varepsilon \in \mathbb{R} \to 0 < \varepsilon \to \varphi_{87}$}
\end{ptree}
\noindent$\Pi_{27}$:
\begin{ptree}
\AXC{$\varphi_{84}$\lbl{A2}}
\AXC{$\mathrm{Seq}(a) \to \mathrm{Conv}(a \circ \mathsf{s},\alpha) \to \varepsilon \in \mathbb{R} \to 0 < \varepsilon \to \varphi_{109}$\lbl{A1}}
\AXC{$\varphi_{186}$\lbl{A1}}
\BIC{$\mathrm{Seq}(a) \to \mathrm{Conv}(a \circ \mathsf{s},\alpha) \to \varepsilon \in \mathbb{R} \to 0 < \varepsilon \to \varphi_{126}$}
\AXC{$\varphi_{206}$\lbl{A1}}
\BIC{$\varphi_{165}$}
\BIC{$\mathrm{Seq}(a) \to \mathrm{Conv}(a \circ \mathsf{s},\alpha) \to \varepsilon \in \mathbb{R} \to 0 < \varepsilon \to \varphi_{124}$}
\AXC{$\varphi_{215}$\lbl{A1}}
\BIC{$\varphi_{204}$}
\end{ptree}
\noindent$\Pi_{28}$:
\begin{ptree}
\AXC{$\mathrm{Seq}(a) \to \mathrm{Conv}(a \circ \mathsf{s},\alpha) \to \varepsilon \in \mathbb{R} \to \varphi_{111}$\lbl{A1}}
\AXC{$\varphi_{175}$\lbl{A1}}
\BIC{$\mathrm{Seq}(a) \to \mathrm{Conv}(a \circ \mathsf{s},\alpha) \to \varepsilon \in \mathbb{R} \to 0 < \varepsilon \to \varphi_{95}$}
\AXC{$\stackrel{\Pi_{27}}{\vdots}$}\noLine
\UIC{$\varphi_{204}$}
\BIC{$\varphi_{179}$}
\end{ptree}
\noindent$\Pi_{29}$:
\begin{ptree}
\AXC{$y = z \to x = y \to x = z$\lbl{A6}}
\AXC{$(y = z \to x = y \to x = z) \to x = y \to y = z \to x = z$\lbl{A1}}
\BIC{$x = y \to y = z \to x = z$}
\UIC{$\forall z\,(x = y \to y = z \to x = z)$}
\UIC{$\forall y\,\forall z\,(x = y \to y = z \to x = z)$}
\UIC{$\forall x\,\forall y\,\forall z\,(x = y \to y = z \to x = z)$}
\AXC{$\varphi_{58}$\lbl{A2}}
\BIC{$\forall y\,\forall z\,(a \circ \mathsf{s}(m) = y \to y = z \to a \circ \mathsf{s}(m) = z)$}
\end{ptree}
\noindent$\Pi_{30}$:
\begin{ptree}
\AXC{$\stackrel{\Pi_{29}}{\vdots}$}\noLine
\UIC{$\forall y\,\forall z\,(a \circ \mathsf{s}(m) = y \to y = z \to a \circ \mathsf{s}(m) = z)$}
\AXC{$\forall y\,\forall z\,(a \circ \mathsf{s}(m) = y \to y = z \to a \circ \mathsf{s}(m) = z) \to \varphi_{49}$\lbl{A2}}
\BIC{$\varphi_{49}$}
\AXC{$\varphi_{120}$\lbl{A2}}
\BIC{$a \circ \mathsf{s}(m) = a_{S(m)} \to a_{S(m)} = a_{n} \to a \circ \mathsf{s}(m) = a_{n}$}
\end{ptree}
\noindent$\Pi_{31}$:
\begin{ptree}
\AXC{$\stackrel{\Pi_{11}}{\vdots}$}\noLine
\UIC{$\forall x\,\forall y\,(x = y \to y = x)$}
\AXC{$\forall x\,\forall y\,(x = y \to y = x) \to \forall y\,(n = y \to y = n)$\lbl{A2}}
\BIC{$\forall y\,(n = y \to y = n)$}
\end{ptree}
\noindent$\Pi_{32}$:
\begin{ptree}
\AXC{$\stackrel{\Pi_{31}}{\vdots}$}\noLine
\UIC{$\forall y\,(n = y \to y = n)$}
\AXC{$\forall y\,(n = y \to y = n) \to n = S(m) \to S(m) = n$\lbl{A2}}
\BIC{$n = S(m) \to S(m) = n$}
\end{ptree}
\noindent$\Pi_{33}$:
\begin{ptree}
\AXC{$\stackrel{\Pi_{32}}{\vdots}$}\noLine
\UIC{$n = S(m) \to S(m) = n$}
\AXC{$\mathrm{Seq}(a) \to \mathrm{Conv}(a \circ \mathsf{s},\alpha) \to \varepsilon \in \mathbb{R} \to \varphi_{111}$\lbl{A1}}
\AXC{$\varphi_{173}$\lbl{A1}}
\BIC{$\mathrm{Seq}(a) \to \mathrm{Conv}(a \circ \mathsf{s},\alpha) \to \varepsilon \in \mathbb{R} \to 0 < \varepsilon \to \varphi_{89}$}
\AXC{$\varphi_{176}$\lbl{A1}}
\BIC{$\varphi_{130}$}
\BIC{$\mathrm{Seq}(a) \to \mathrm{Conv}(a \circ \mathsf{s},\alpha) \to \varepsilon \in \mathbb{R} \to 0 < \varepsilon \to \varphi_{91}$}
\AXC{$\varphi_{189}$\lbl{A1}}
\BIC{$a_{S(m)} = a_{S(m)} \to \varphi_{140}$}
\end{ptree}
\noindent$\Pi_{34}$:
\begin{ptree}
\AXC{$x = x$\lbl{A5}}
\UIC{$\forall x\,(x = x)$}
\AXC{$\forall x\,(x = x) \to a_{S(m)} = a_{S(m)}$\lbl{A2}}
\BIC{$a_{S(m)} = a_{S(m)}$}
\AXC{$\stackrel{\Pi_{33}}{\vdots}$}\noLine
\UIC{$a_{S(m)} = a_{S(m)} \to \varphi_{140}$}
\BIC{$\varphi_{140}$}
\end{ptree}
\noindent$\Pi_{35}$:
\begin{ptree}
\AXC{$z = w \to a_{S(m)} = a_{z} \to a_{S(m)} = a_{w}$\lbl{A6}}
\UIC{$\forall w\,(z = w \to a_{S(m)} = a_{z} \to a_{S(m)} = a_{w})$}
\UIC{$\forall z\,\forall w\,(z = w \to a_{S(m)} = a_{z} \to a_{S(m)} = a_{w})$}
\AXC{$\varphi_{85}$\lbl{A2}}
\BIC{$\forall w\,(S(m) = w \to a_{S(m)} = a_{S(m)} \to a_{S(m)} = a_{w})$}
\AXC{$\varphi_{86}$\lbl{A2}}
\BIC{$S(m) = n \to a_{S(m)} = a_{S(m)} \to a_{S(m)} = a_{n}$}
\AXC{$\stackrel{\Pi_{34}}{\vdots}$}\noLine
\UIC{$\varphi_{140}$}
\BIC{$\mathrm{Seq}(a) \to \mathrm{Conv}(a \circ \mathsf{s},\alpha) \to \varepsilon \in \mathbb{R} \to \varphi_{112}$}
\end{ptree}
\noindent$\Pi_{36}$:
\begin{ptree}
\AXC{$y = z \to x = y \to x = z$\lbl{A6}}
\AXC{$(y = z \to x = y \to x = z) \to x = y \to y = z \to x = z$\lbl{A1}}
\BIC{$x = y \to y = z \to x = z$}
\UIC{$\forall z\,(x = y \to y = z \to x = z)$}
\UIC{$\forall y\,\forall z\,(x = y \to y = z \to x = z)$}
\UIC{$\forall x\,\forall y\,\forall z\,(x = y \to y = z \to x = z)$}
\AXC{$\varphi_{58}$\lbl{A2}}
\BIC{$\forall y\,\forall z\,(a \circ \mathsf{s}(m) = y \to y = z \to a \circ \mathsf{s}(m) = z)$}
\end{ptree}
\noindent$\Pi_{37}$:
\begin{ptree}
\AXC{$\stackrel{\Pi_{36}}{\vdots}$}\noLine
\UIC{$\forall y\,\forall z\,(a \circ \mathsf{s}(m) = y \to y = z \to a \circ \mathsf{s}(m) = z)$}
\AXC{$\forall y\,\forall z\,(a \circ \mathsf{s}(m) = y \to y = z \to a \circ \mathsf{s}(m) = z) \to \varphi_{65}$\lbl{A2}}
\BIC{$\varphi_{65}$}
\AXC{$\varphi_{65} \to \varphi_{79}$\lbl{A2}}
\BIC{$\varphi_{79}$}
\end{ptree}
\noindent$\Pi_{38}$:
\begin{ptree}
\AXC{$\varphi_{26}$\lbl{A2}}
\AXC{$\mathrm{NSeq}(\mathsf{s}) \land \forall k \in \mathbb{N}\,(\mathsf{s}(k) = S(k))$\lbl{[S14]}}
\AXC{$\varphi_{141}$\lbl{A1}}
\BIC{$\mathrm{Seq}(a) \to \mathrm{Conv}(a \circ \mathsf{s},\alpha) \to \varepsilon \in \mathbb{R} \to 0 < \varepsilon \to \varphi_{113}$}
\AXC{$\varphi_{191}$\lbl{A1}}
\BIC{$\varphi_{147}$}
\BIC{$\mathrm{Seq}(a) \to \mathrm{Conv}(a \circ \mathsf{s},\alpha) \to \varepsilon \in \mathbb{R} \to 0 < \varepsilon \to \varphi_{110}$}
\AXC{$\varphi_{205}$\lbl{A1}}
\BIC{$\varphi_{174}$}
\end{ptree}
\noindent$\Pi_{39}$:
\begin{ptree}
\AXC{$\mathrm{Seq}(a) \to \mathrm{Conv}(a \circ \mathsf{s},\alpha) \to \varepsilon \in \mathbb{R} \to \varphi_{111}$\lbl{A1}}
\AXC{$\varphi_{175}$\lbl{A1}}
\BIC{$\mathrm{Seq}(a) \to \mathrm{Conv}(a \circ \mathsf{s},\alpha) \to \varepsilon \in \mathbb{R} \to 0 < \varepsilon \to \varphi_{95}$}
\AXC{$\stackrel{\Pi_{38}}{\vdots}$}\noLine
\UIC{$\varphi_{174}$}
\BIC{$\mathrm{Seq}(a) \to \mathrm{Conv}(a \circ \mathsf{s},\alpha) \to \varepsilon \in \mathbb{R} \to \varphi_{108}$}
\AXC{$\varphi_{198}$\lbl{A1}}
\BIC{$a_{\mathsf{s}(m)} = a_{\mathsf{s}(m)} \to \varphi_{153}$}
\end{ptree}
\noindent$\Pi_{40}$:
\begin{ptree}
\AXC{$x = x$\lbl{A5}}
\UIC{$\forall x\,(x = x)$}
\AXC{$\forall x\,(x = x) \to a_{\mathsf{s}(m)} = a_{\mathsf{s}(m)}$\lbl{A2}}
\BIC{$a_{\mathsf{s}(m)} = a_{\mathsf{s}(m)}$}
\AXC{$\stackrel{\Pi_{39}}{\vdots}$}\noLine
\UIC{$a_{\mathsf{s}(m)} = a_{\mathsf{s}(m)} \to \varphi_{153}$}
\BIC{$\varphi_{153}$}
\end{ptree}
\noindent$\Pi_{41}$:
\begin{ptree}
\AXC{$\forall a\,(\varphi_{63})$\lbl{[S12]}}
\AXC{$\forall a\,(\varphi_{63}) \to \varphi_{63}$\lbl{A2}}
\BIC{$\varphi_{63}$}
\AXC{$\varphi_{63} \to \mathrm{Seq}(a) \to \varphi_{64}$\lbl{A2}}
\BIC{$\mathrm{Seq}(a) \to \varphi_{64}$}
\AXC{$(\mathrm{Seq}(a) \to \varphi_{64}) \to \varphi_{209}$\lbl{A1}}
\BIC{$\varphi_{209}$}
\end{ptree}
\noindent$\Pi_{42}$:
\begin{ptree}
\AXC{$\mathrm{Seq}(a) \to \mathrm{Conv}(a \circ \mathsf{s},\alpha) \to \varepsilon \in \mathbb{R} \to 0 < \varepsilon \to \varphi_{92}$\lbl{A1}}
\AXC{$\stackrel{\Pi_{41}}{\vdots}$}\noLine
\UIC{$\varphi_{209}$}
\BIC{$\varphi_{185}$}
\end{ptree}
\noindent$\Pi_{43}$:
\begin{ptree}
\AXC{$\varphi_{80}$\lbl{A2}}
\AXC{$\mathrm{NSeq}(\mathsf{s}) \land \forall k \in \mathbb{N}\,(\mathsf{s}(k) = S(k))$\lbl{[S14]}}
\AXC{$\varphi_{136}$\lbl{A1}}
\BIC{$\mathrm{Seq}(a) \to \mathrm{Conv}(a \circ \mathsf{s},\alpha) \to \varepsilon \in \mathbb{R} \to \varphi_{104}$}
\AXC{$\stackrel{\Pi_{42}}{\vdots}$}\noLine
\UIC{$\varphi_{185}$}
\BIC{$\mathrm{Seq}(a) \to \mathrm{Conv}(a \circ \mathsf{s},\alpha) \to \varepsilon \in \mathbb{R} \to 0 < \varepsilon \to \varphi_{129}$}
\AXC{$\varphi_{192}$\lbl{A1}}
\BIC{$\mathrm{Seq}(a) \to \mathrm{Conv}(a \circ \mathsf{s},\alpha) \to \varepsilon \in \mathbb{R} \to 0 < \varepsilon \to \varphi_{123}$}
\AXC{$\varphi_{203}$\lbl{A1}}
\BIC{$\varphi_{161}$}
\BIC{$\mathrm{Seq}(a) \to \mathrm{Conv}(a \circ \mathsf{s},\alpha) \to \varepsilon \in \mathbb{R} \to 0 < \varepsilon \to \varphi_{122}$}
\end{ptree}
\noindent$\Pi_{44}$:
\begin{ptree}
\AXC{$\stackrel{\Pi_{43}}{\vdots}$}\noLine
\UIC{$\mathrm{Seq}(a) \to \mathrm{Conv}(a \circ \mathsf{s},\alpha) \to \varepsilon \in \mathbb{R} \to 0 < \varepsilon \to \varphi_{122}$}
\AXC{$\varphi_{210}$\lbl{A1}}
\BIC{$\varphi_{181}$}
\end{ptree}
\noindent$\Pi_{45}$:
\begin{ptree}
\AXC{$\mathrm{Seq}(a) \to \mathrm{Conv}(a \circ \mathsf{s},\alpha) \to \varepsilon \in \mathbb{R} \to \varphi_{111}$\lbl{A1}}
\AXC{$\varphi_{175}$\lbl{A1}}
\BIC{$\mathrm{Seq}(a) \to \mathrm{Conv}(a \circ \mathsf{s},\alpha) \to \varepsilon \in \mathbb{R} \to 0 < \varepsilon \to \varphi_{95}$}
\AXC{$\stackrel{\Pi_{44}}{\vdots}$}\noLine
\UIC{$\varphi_{181}$}
\BIC{$\mathrm{Seq}(a) \to \mathrm{Conv}(a \circ \mathsf{s},\alpha) \to \varepsilon \in \mathbb{R} \to \varphi_{121}$}
\end{ptree}
\noindent$\Pi_{46}$:
\begin{ptree}
\AXC{$\stackrel{\Pi_{45}}{\vdots}$}\noLine
\UIC{$\mathrm{Seq}(a) \to \mathrm{Conv}(a \circ \mathsf{s},\alpha) \to \varepsilon \in \mathbb{R} \to \varphi_{121}$}
\AXC{$(\mathrm{Seq}(a) \to \mathrm{Conv}(a \circ \mathsf{s},\alpha) \to \varepsilon \in \mathbb{R} \to \varphi_{121}) \to \varphi_{196}$\lbl{A1}}
\BIC{$\varphi_{196}$}
\end{ptree}
\noindent$\Pi_{47}$:
\begin{ptree}
\AXC{$\stackrel{\Pi_{37}}{\vdots}$}\noLine
\UIC{$\varphi_{79}$}
\AXC{$z = w \to a_{\mathsf{s}(m)} = a_{z} \to a_{\mathsf{s}(m)} = a_{w}$\lbl{A6}}
\UIC{$\forall w\,(z = w \to a_{\mathsf{s}(m)} = a_{z} \to a_{\mathsf{s}(m)} = a_{w})$}
\UIC{$\forall z\,\forall w\,(z = w \to a_{\mathsf{s}(m)} = a_{z} \to a_{\mathsf{s}(m)} = a_{w})$}
\AXC{$\varphi_{119}$\lbl{A2}}
\BIC{$\forall w\,(\mathsf{s}(m) = w \to a_{\mathsf{s}(m)} = a_{\mathsf{s}(m)} \to a_{\mathsf{s}(m)} = a_{w})$}
\AXC{$\varphi_{127}$\lbl{A2}}
\BIC{$\mathsf{s}(m) = S(m) \to a_{\mathsf{s}(m)} = a_{\mathsf{s}(m)} \to a_{\mathsf{s}(m)} = a_{S(m)}$}
\AXC{$\stackrel{\Pi_{40}}{\vdots}$}\noLine
\UIC{$\varphi_{153}$}
\BIC{$\mathrm{Seq}(a) \to \mathrm{Conv}(a \circ \mathsf{s},\alpha) \to \varepsilon \in \mathbb{R} \to \varphi_{117}$}
\AXC{$\stackrel{\Pi_{46}}{\vdots}$}\noLine
\UIC{$\varphi_{196}$}
\BIC{$\varphi_{79} \to \mathrm{Seq}(a) \to \mathrm{Conv}(a \circ \mathsf{s},\alpha) \to \varepsilon \in \mathbb{R} \to \varphi_{118}$}
\BIC{$\mathrm{Seq}(a) \to \mathrm{Conv}(a \circ \mathsf{s},\alpha) \to \varepsilon \in \mathbb{R} \to \varphi_{118}$}
\AXC{$\varphi_{212}$\lbl{A1}}
\BIC{$(\mathrm{Seq}(a) \to \mathrm{Conv}(a \circ \mathsf{s},\alpha) \to \varepsilon \in \mathbb{R} \to \varphi_{112}) \to \varphi_{152}$}
\end{ptree}
\noindent$\Pi_{48}$:
\begin{ptree}
\AXC{$\stackrel{\Pi_{35}}{\vdots}$}\noLine
\UIC{$\mathrm{Seq}(a) \to \mathrm{Conv}(a \circ \mathsf{s},\alpha) \to \varepsilon \in \mathbb{R} \to \varphi_{112}$}
\AXC{$\stackrel{\Pi_{47}}{\vdots}$}\noLine
\UIC{$(\mathrm{Seq}(a) \to \mathrm{Conv}(a \circ \mathsf{s},\alpha) \to \varepsilon \in \mathbb{R} \to \varphi_{112}) \to \varphi_{152}$}
\BIC{$\varphi_{152}$}
\end{ptree}
\noindent$\Pi_{49}$:
\begin{ptree}
\AXC{$\stackrel{\Pi_{30}}{\vdots}$}\noLine
\UIC{$a \circ \mathsf{s}(m) = a_{S(m)} \to a_{S(m)} = a_{n} \to a \circ \mathsf{s}(m) = a_{n}$}
\AXC{$\stackrel{\Pi_{48}}{\vdots}$}\noLine
\UIC{$\varphi_{152}$}
\BIC{$\mathrm{Seq}(a) \to \mathrm{Conv}(a \circ \mathsf{s},\alpha) \to \varepsilon \in \mathbb{R} \to \varphi_{116}$}
\AXC{$(\mathrm{Seq}(a) \to \mathrm{Conv}(a \circ \mathsf{s},\alpha) \to \varepsilon \in \mathbb{R} \to \varphi_{116}) \to \varphi_{195}$\lbl{A1}}
\BIC{$\varphi_{195}$}
\end{ptree}
\noindent$\Pi_{50}$:
\begin{ptree}
\AXC{$\stackrel{\Pi_{1}}{\vdots}$}\noLine
\UIC{$\varphi_{71}$}
\AXC{$\stackrel{\Pi_{26}}{\vdots}$}\noLine
\UIC{$\mathrm{Seq}(a) \to \mathrm{Conv}(a \circ \mathsf{s},\alpha) \to \varepsilon \in \mathbb{R} \to 0 < \varepsilon \to \varphi_{87}$}
\AXC{$\stackrel{\Pi_{28}}{\vdots}$}\noLine
\UIC{$\varphi_{179}$}
\BIC{$\mathrm{Seq}(a) \to \mathrm{Conv}(a \circ \mathsf{s},\alpha) \to \varepsilon \in \mathbb{R} \to 0 < \varepsilon \to \varphi_{114}$}
\AXC{$\stackrel{\Pi_{49}}{\vdots}$}\noLine
\UIC{$\varphi_{195}$}
\BIC{$\varphi_{71} \to \mathrm{Seq}(a) \to \mathrm{Conv}(a \circ \mathsf{s},\alpha) \to \varepsilon \in \mathbb{R} \to \varphi_{115}$}
\BIC{$\mathrm{Seq}(a) \to \mathrm{Conv}(a \circ \mathsf{s},\alpha) \to \varepsilon \in \mathbb{R} \to \varphi_{115}$}
\end{ptree}
\noindent$\Pi_{51}$:
\begin{ptree}
\AXC{$\stackrel{\Pi_{50}}{\vdots}$}\noLine
\UIC{$\mathrm{Seq}(a) \to \mathrm{Conv}(a \circ \mathsf{s},\alpha) \to \varepsilon \in \mathbb{R} \to \varphi_{115}$}
\AXC{$(\mathrm{Seq}(a) \to \mathrm{Conv}(a \circ \mathsf{s},\alpha) \to \varepsilon \in \mathbb{R} \to \varphi_{115}) \to \varphi_{131}$\lbl{A1}}
\BIC{$\varphi_{131}$}
\UIC{$\forall m\,(\varphi_{131})$}
\end{ptree}
\noindent$\Pi_{52}$:
\begin{ptree}
\AXC{$\stackrel{\Pi_{11}}{\vdots}$}\noLine
\UIC{$\forall x\,\forall y\,(x = y \to y = x)$}
\AXC{$\forall x\,\forall y\,(x = y \to y = x) \to \forall y\,(n = y \to y = n)$\lbl{A2}}
\BIC{$\forall y\,(n = y \to y = n)$}
\end{ptree}
\noindent$\Pi_{53}$:
\begin{ptree}
\AXC{$\forall n \in \mathbb{N}\,(1_{\omega} \leq n)$\lbl{[N12]}}
\AXC{$\forall n \in \mathbb{N}\,(1_{\omega} \leq n) \to N \in \mathbb{N} \to 1_{\omega} \leq N$\lbl{A2}}
\BIC{$N \in \mathbb{N} \to 1_{\omega} \leq N$}
\AXC{$(N \in \mathbb{N} \to 1_{\omega} \leq N) \to \varphi_{154}$\lbl{A1}}
\BIC{$\varphi_{154}$}
\end{ptree}
\noindent$\Pi_{54}$:
\begin{ptree}
\AXC{$\forall x\,\forall y\,\forall z\,(x \leq y \to y < z \to x < z)$\lbl{[O4]}}
\AXC{$\varphi_{31}$\lbl{A2}}
\BIC{$\forall y\,\forall z\,(1_{\omega} \leq y \to y < z \to 1_{\omega} < z)$}
\AXC{$\varphi_{32}$\lbl{A2}}
\BIC{$\forall z\,(1_{\omega} \leq N \to N < z \to 1_{\omega} < z)$}
\AXC{$\varphi_{33}$\lbl{A2}}
\BIC{$1_{\omega} \leq N \to N < S(N) \to 1_{\omega} < S(N)$}
\end{ptree}
\noindent$\Pi_{55}$:
\begin{ptree}
\AXC{$\stackrel{\Pi_{54}}{\vdots}$}\noLine
\UIC{$1_{\omega} \leq N \to N < S(N) \to 1_{\omega} < S(N)$}
\AXC{$(1_{\omega} \leq N \to N < S(N) \to 1_{\omega} < S(N)) \to \varphi_{190}$\lbl{A1}}
\BIC{$\varphi_{190}$}
\end{ptree}
\noindent$\Pi_{56}$:
\begin{ptree}
\AXC{$\mathrm{Seq}(a) \to \mathrm{Conv}(a \circ \mathsf{s},\alpha) \to \varepsilon \in \mathbb{R} \to 0 < \varepsilon \to \varphi_{109}$\lbl{A1}}
\AXC{$\varphi_{164}$\lbl{A1}}
\BIC{$\mathrm{Seq}(a) \to \mathrm{Conv}(a \circ \mathsf{s},\alpha) \to \varepsilon \in \mathbb{R} \to 0 < \varepsilon \to \varphi_{55}$}
\AXC{$\stackrel{\Pi_{53}}{\vdots}$}\noLine
\UIC{$\varphi_{154}$}
\BIC{$\mathrm{Seq}(a) \to \mathrm{Conv}(a \circ \mathsf{s},\alpha) \to \varepsilon \in \mathbb{R} \to 0 < \varepsilon \to \varphi_{73}$}
\AXC{$\stackrel{\Pi_{55}}{\vdots}$}\noLine
\UIC{$\varphi_{190}$}
\BIC{$\varphi_{155}$}
\end{ptree}
\noindent$\Pi_{57}$:
\begin{ptree}
\AXC{$\forall x\,\forall y\,\forall z\,(x < y \to y \leq z \to x < z)$\lbl{[O3]}}
\AXC{$\varphi_{34}$\lbl{A2}}
\BIC{$\forall y\,\forall z\,(1_{\omega} < y \to y \leq z \to 1_{\omega} < z)$}
\AXC{$\varphi_{44}$\lbl{A2}}
\BIC{$\forall z\,(1_{\omega} < S(N) \to S(N) \leq z \to 1_{\omega} < z)$}
\AXC{$\varphi_{45}$\lbl{A2}}
\BIC{$1_{\omega} < S(N) \to S(N) \leq n \to 1_{\omega} < n$}
\end{ptree}
\noindent$\Pi_{58}$:
\begin{ptree}
\AXC{$\stackrel{\Pi_{57}}{\vdots}$}\noLine
\UIC{$1_{\omega} < S(N) \to S(N) \leq n \to 1_{\omega} < n$}
\AXC{$(1_{\omega} < S(N) \to S(N) \leq n \to 1_{\omega} < n) \to \varphi_{193}$\lbl{A1}}
\BIC{$\varphi_{193}$}
\end{ptree}
\noindent$\Pi_{59}$:
\begin{ptree}
\AXC{$\stackrel{\Pi_{5}}{\vdots}$}\noLine
\UIC{$\mathrm{Seq}(a) \to \mathrm{Conv}(a \circ \mathsf{s},\alpha) \to \varepsilon \in \mathbb{R} \to 0 < \varepsilon \to \varphi_{56}$}
\AXC{$\stackrel{\Pi_{56}}{\vdots}$}\noLine
\UIC{$\varphi_{155}$}
\BIC{$\mathrm{Seq}(a) \to \mathrm{Conv}(a \circ \mathsf{s},\alpha) \to \varepsilon \in \mathbb{R} \to 0 < \varepsilon \to \varphi_{81}$}
\AXC{$\stackrel{\Pi_{58}}{\vdots}$}\noLine
\UIC{$\varphi_{193}$}
\BIC{$\varphi_{157}$}
\end{ptree}
\noindent$\Pi_{60}$:
\begin{ptree}
\AXC{$\forall x\,\forall y\,(x < y \to x \neq y)$\lbl{[F26]}}
\AXC{$\varphi_{3}$\lbl{A2}}
\BIC{$\forall y\,(1_{\omega} < y \to 1_{\omega} \neq y)$}
\AXC{$\varphi_{4}$\lbl{A2}}
\BIC{$1_{\omega} < n \to 1_{\omega} \neq n$}
\AXC{$(1_{\omega} < n \to 1_{\omega} \neq n) \to \varphi_{158}$\lbl{A1}}
\BIC{$\varphi_{158}$}
\end{ptree}
\noindent$\Pi_{61}$:
\begin{ptree}
\AXC{$\mathrm{Seq}(a) \to \mathrm{Conv}(a \circ \mathsf{s},\alpha) \to \varepsilon \in \mathbb{R} \to 0 < \varepsilon \to \varphi_{72}$\lbl{A1}}
\AXC{$\stackrel{\Pi_{59}}{\vdots}$}\noLine
\UIC{$\varphi_{157}$}
\BIC{$\mathrm{Seq}(a) \to \mathrm{Conv}(a \circ \mathsf{s},\alpha) \to \varepsilon \in \mathbb{R} \to 0 < \varepsilon \to \varphi_{77}$}
\AXC{$\stackrel{\Pi_{60}}{\vdots}$}\noLine
\UIC{$\varphi_{158}$}
\BIC{$\mathrm{Seq}(a) \to \mathrm{Conv}(a \circ \mathsf{s},\alpha) \to \varepsilon \in \mathbb{R} \to 0 < \varepsilon \to \varphi_{75}$}
\AXC{$\varphi_{163}$\lbl{A1}}
\BIC{$\varphi_{128}$}
\end{ptree}
\noindent$\Pi_{62}$:
\begin{ptree}
\AXC{$\stackrel{\Pi_{52}}{\vdots}$}\noLine
\UIC{$\forall y\,(n = y \to y = n)$}
\AXC{$\forall y\,(n = y \to y = n) \to n = 1_{\omega} \to 1_{\omega} = n$\lbl{A2}}
\BIC{$n = 1_{\omega} \to 1_{\omega} = n$}
\AXC{$\stackrel{\Pi_{61}}{\vdots}$}\noLine
\UIC{$\varphi_{128}$}
\BIC{$\mathrm{Seq}(a) \to \mathrm{Conv}(a \circ \mathsf{s},\alpha) \to \varepsilon \in \mathbb{R} \to 0 < \varepsilon \to \varphi_{76}$}
\end{ptree}
\noindent$\Pi_{63}$:
\begin{ptree}
\AXC{$\forall n \in \mathbb{N}\,(n \neq 1_{\omega} \to \exists m \in \mathbb{N}\,(n = S(m)))$\lbl{[N18]}}
\AXC{$\varphi_{67}$\lbl{A2}}
\BIC{$n \in \mathbb{N} \to n \neq 1_{\omega} \to \exists m \in \mathbb{N}\,(n = S(m))$}
\AXC{$(n \in \mathbb{N} \to n \neq 1_{\omega} \to \exists m \in \mathbb{N}\,(n = S(m))) \to \varphi_{194}$\lbl{A1}}
\BIC{$\varphi_{194}$}
\end{ptree}
\noindent$\Pi_{64}$:
\begin{ptree}
\AXC{$\stackrel{\Pi_{62}}{\vdots}$}\noLine
\UIC{$\mathrm{Seq}(a) \to \mathrm{Conv}(a \circ \mathsf{s},\alpha) \to \varepsilon \in \mathbb{R} \to 0 < \varepsilon \to \varphi_{76}$}
\AXC{$\mathrm{Seq}(a) \to \mathrm{Conv}(a \circ \mathsf{s},\alpha) \to \varepsilon \in \mathbb{R} \to 0 < \varepsilon \to \varphi_{70}$\lbl{A1}}
\AXC{$\stackrel{\Pi_{63}}{\vdots}$}\noLine
\UIC{$\varphi_{194}$}
\BIC{$\varphi_{162}$}
\BIC{$\mathrm{Seq}(a) \to \mathrm{Conv}(a \circ \mathsf{s},\alpha) \to \varepsilon \in \mathbb{R} \to 0 < \varepsilon \to \varphi_{83}$}
\end{ptree}
\noindent$\Pi_{65}$:
\begin{ptree}
\AXC{$\varphi_{167}$\lbl{A4}}
\AXC{$\forall m\,(\varphi_{131}) \to \varphi_{133}$\lbl{A3}}
\AXC{$\stackrel{\Pi_{51}}{\vdots}$}\noLine
\UIC{$\forall m\,(\varphi_{131})$}
\AXC{$\stackrel{\Pi_{64}}{\vdots}$}\noLine
\UIC{$\mathrm{Seq}(a) \to \mathrm{Conv}(a \circ \mathsf{s},\alpha) \to \varepsilon \in \mathbb{R} \to 0 < \varepsilon \to \varphi_{83}$}
\AXC{$\varphi_{217}$\lbl{A1}}
\BIC{$\forall m\,(\varphi_{131}) \to (\forall m\,(\varphi_{131}) \to \varphi_{133}) \to \varphi_{201}$}
\BIC{$(\forall m\,(\varphi_{131}) \to \varphi_{133}) \to \varphi_{201}$}
\BIC{$\varphi_{201}$}
\BIC{$\mathrm{Seq}(a) \to \mathrm{Conv}(a \circ \mathsf{s},\alpha) \to \varepsilon \in \mathbb{R} \to 0 < \varepsilon \to \varphi_{88}$}
\end{ptree}
\noindent$\Pi_{66}$:
\begin{ptree}
\AXC{$\stackrel{\Pi_{65}}{\vdots}$}\noLine
\UIC{$\mathrm{Seq}(a) \to \mathrm{Conv}(a \circ \mathsf{s},\alpha) \to \varepsilon \in \mathbb{R} \to 0 < \varepsilon \to \varphi_{88}$}
\AXC{$\varphi_{166}$\lbl{A1}}
\BIC{$\varphi_{97} \to n \in \mathbb{N} \to S(N) \leq n \to \lvert a_{n} - \alpha \rvert < \varepsilon$}
\UIC{$\forall n\,(\varphi_{97} \to n \in \mathbb{N} \to S(N) \leq n \to \lvert a_{n} - \alpha \rvert < \varepsilon)$}
\AXC{$\forall n\,(\varphi_{97} \to n \in \mathbb{N} \to S(N) \leq n \to \lvert a_{n} - \alpha \rvert < \varepsilon) \to \varphi_{171} \to \varphi_{188}$\lbl{A1}}
\BIC{$\varphi_{171} \to \varphi_{188}$}
\end{ptree}
\noindent$\Pi_{67}$:
\begin{ptree}
\AXC{$\forall n \in \mathbb{N}\,(S(n) \in \mathbb{N})$\lbl{[N2]}}
\AXC{$\forall n \in \mathbb{N}\,(S(n) \in \mathbb{N}) \to N \in \mathbb{N} \to S(N) \in \mathbb{N}$\lbl{A2}}
\BIC{$N \in \mathbb{N} \to S(N) \in \mathbb{N}$}
\AXC{$(N \in \mathbb{N} \to S(N) \in \mathbb{N}) \to \varphi_{151}$\lbl{A1}}
\BIC{$\varphi_{151}$}
\end{ptree}
\noindent$\Pi_{68}$:
\begin{ptree}
\AXC{$\mathrm{Seq}(a) \to \mathrm{Conv}(a \circ \mathsf{s},\alpha) \to \varepsilon \in \mathbb{R} \to 0 < \varepsilon \to \varphi_{109}$\lbl{A1}}
\AXC{$\varphi_{164}$\lbl{A1}}
\BIC{$\mathrm{Seq}(a) \to \mathrm{Conv}(a \circ \mathsf{s},\alpha) \to \varepsilon \in \mathbb{R} \to 0 < \varepsilon \to \varphi_{55}$}
\AXC{$\stackrel{\Pi_{67}}{\vdots}$}\noLine
\UIC{$\varphi_{151}$}
\BIC{$\mathrm{Seq}(a) \to \mathrm{Conv}(a \circ \mathsf{s},\alpha) \to \varepsilon \in \mathbb{R} \to 0 < \varepsilon \to \varphi_{61}$}
\AXC{$\varphi_{199}$\lbl{A1}}
\BIC{$\varphi_{177}$}
\end{ptree}
\noindent$\Pi_{69}$:
\begin{ptree}
\AXC{$\varphi_{28} \to \varphi_{35}$\lbl{A2}}
\AXC{$\varphi_{97} \to \forall n\,(\varphi_{97})$\lbl{A4}}
\AXC{$\varphi_{171}$\lbl{A3}}
\AXC{$\stackrel{\Pi_{66}}{\vdots}$}\noLine
\UIC{$\varphi_{171} \to \varphi_{188}$}
\BIC{$\varphi_{188}$}
\BIC{$\mathrm{Seq}(a) \to \mathrm{Conv}(a \circ \mathsf{s},\alpha) \to \varepsilon \in \mathbb{R} \to 0 < \varepsilon \to \varphi_{94}$}
\AXC{$\stackrel{\Pi_{68}}{\vdots}$}\noLine
\UIC{$\varphi_{177}$}
\BIC{$\mathrm{Seq}(a) \to \mathrm{Conv}(a \circ \mathsf{s},\alpha) \to \varepsilon \in \mathbb{R} \to 0 < \varepsilon \to \varphi_{107}$}
\AXC{$\varphi_{200}$\lbl{A1}}
\BIC{$\varphi_{159}$}
\BIC{$\mathrm{Seq}(a) \to \mathrm{Conv}(a \circ \mathsf{s},\alpha) \to \varepsilon \in \mathbb{R} \to 0 < \varepsilon \to \varphi_{103}$}
\AXC{$\varphi_{180}$\lbl{A1}}
\BIC{$\varphi_{93} \to \varphi_{41}$}
\UIC{$\forall N\,(\varphi_{93} \to \varphi_{41})$}
\end{ptree}
\noindent$\Pi_{70}$:
\begin{ptree}
\AXC{$\forall a\,\forall \alpha\,\forall \varepsilon\,(\mathrm{Conv}(a,\alpha) \to \varphi_{66})$\lbl{[C3]}}
\AXC{$\varphi_{139}$\lbl{A2}}
\BIC{$\forall \alpha\,\forall \varepsilon\,(\mathrm{Conv}(a \circ \mathsf{s},\alpha) \to \varepsilon \in \mathbb{R} \to \varphi_{60})$}
\AXC{$\varphi_{144}$\lbl{A2}}
\BIC{$\forall \varepsilon\,(\mathrm{Conv}(a \circ \mathsf{s},\alpha) \to \varepsilon \in \mathbb{R} \to \varphi_{60})$}
\AXC{$\varphi_{142}$\lbl{A2}}
\BIC{$\mathrm{Conv}(a \circ \mathsf{s},\alpha) \to \varepsilon \in \mathbb{R} \to \varphi_{60}$}
\AXC{$(\mathrm{Conv}(a \circ \mathsf{s},\alpha) \to \varepsilon \in \mathbb{R} \to \varphi_{60}) \to \varphi_{156}$\lbl{A1}}
\BIC{$\varphi_{156}$}
\end{ptree}
\noindent$\Pi_{71}$:
\begin{ptree}
\AXC{$\mathrm{Seq}(a) \to \mathrm{Conv}(a \circ \mathsf{s},\alpha) \to \varepsilon \in \mathbb{R} \to \varepsilon \in \mathbb{R}$\lbl{A1}}
\AXC{$\mathrm{Seq}(a) \to \mathrm{Conv}(a \circ \mathsf{s},\alpha) \to \mathrm{Conv}(a \circ \mathsf{s},\alpha)$\lbl{A1}}
\AXC{$\stackrel{\Pi_{70}}{\vdots}$}\noLine
\UIC{$\varphi_{156}$}
\BIC{$\varphi_{145}$}
\BIC{$\varphi_{36} \to \mathrm{Seq}(a) \to \mathrm{Conv}(a \circ \mathsf{s},\alpha) \to \varepsilon \in \mathbb{R} \to \varphi_{60}$}
\end{ptree}
\noindent$\Pi_{72}$:
\begin{ptree}
\AXC{$\stackrel{\Pi_{69}}{\vdots}$}\noLine
\UIC{$\forall N\,(\varphi_{93} \to \varphi_{41})$}
\AXC{$\varphi_{36}$\lbl{A1}}
\AXC{$\stackrel{\Pi_{71}}{\vdots}$}\noLine
\UIC{$\varphi_{36} \to \mathrm{Seq}(a) \to \mathrm{Conv}(a \circ \mathsf{s},\alpha) \to \varepsilon \in \mathbb{R} \to \varphi_{60}$}
\BIC{$\mathrm{Seq}(a) \to \mathrm{Conv}(a \circ \mathsf{s},\alpha) \to \varepsilon \in \mathbb{R} \to \varphi_{60}$}
\AXC{$\varphi_{216}$\lbl{A1}}
\BIC{$\forall N\,(\varphi_{93} \to \varphi_{41}) \to \varphi_{214}$}
\BIC{$\varphi_{214}$}
\end{ptree}
\noindent$\Pi_{73}$:
\begin{ptree}
\AXC{$\varphi_{143}$\lbl{A3}}
\AXC{$\varphi_{93} \to \forall N\,(\varphi_{93})$\lbl{A4}}
\AXC{$\forall N\,(\varphi_{93} \to \varphi_{41}) \to \forall N\,(\varphi_{93}) \to \forall N\,(\varphi_{41})$\lbl{A3}}
\AXC{$\stackrel{\Pi_{72}}{\vdots}$}\noLine
\UIC{$\varphi_{214}$}
\BIC{$(\varphi_{93} \to \forall N\,(\varphi_{93})) \to \mathrm{Seq}(a) \to \mathrm{Conv}(a \circ \mathsf{s},\alpha) \to \varphi_{66}$}
\BIC{$\mathrm{Seq}(a) \to \mathrm{Conv}(a \circ \mathsf{s},\alpha) \to \varphi_{66}$}
\AXC{$\varphi_{138}$\lbl{A1}}
\BIC{$\mathrm{Seq}(a) \land \mathrm{Conv}(a \circ \mathsf{s},\alpha) \to \varphi_{66}$}
\UIC{$\forall \varepsilon\,(\mathrm{Seq}(a) \land \mathrm{Conv}(a \circ \mathsf{s},\alpha) \to \varphi_{66})$}
\AXC{$\varphi_{207}$\lbl{A1}}
\BIC{$\varphi_{143} \to \varphi_{38} \to \mathrm{Seq}(a) \to \mathrm{Conv}(a \circ \mathsf{s},\alpha) \to \varphi_{68}$}
\BIC{$\varphi_{38} \to \mathrm{Seq}(a) \to \mathrm{Conv}(a \circ \mathsf{s},\alpha) \to \varphi_{68}$}
\end{ptree}
\noindent$\Pi_{74}$:
\begin{ptree}
\AXC{$\forall a\,\forall \alpha\,(\mathrm{Conv}(a,\alpha) \to \alpha \in \mathbb{R})$\lbl{[C2]}}
\AXC{$\varphi_{47}$\lbl{A2}}
\BIC{$\forall \alpha\,(\mathrm{Conv}(a \circ \mathsf{s},\alpha) \to \alpha \in \mathbb{R})$}
\AXC{$\varphi_{51}$\lbl{A2}}
\BIC{$\mathrm{Conv}(a \circ \mathsf{s},\alpha) \to \alpha \in \mathbb{R}$}
\AXC{$(\mathrm{Conv}(a \circ \mathsf{s},\alpha) \to \alpha \in \mathbb{R}) \to \varphi_{74}$\lbl{A1}}
\BIC{$\varphi_{74}$}
\end{ptree}
\noindent$\Pi_{75}$:
\begin{ptree}
\AXC{$\forall a\,\forall \alpha\,(\mathrm{Conv}(a,\alpha) \leftrightarrow \mathrm{Seq}(a) \land \alpha \in \mathbb{R} \land \varphi_{68})$\lbl{DefConv}}
\AXC{$\varphi_{148}$\lbl{A2}}
\BIC{$\forall \alpha\,(\mathrm{Conv}(a,\alpha) \leftrightarrow \mathrm{Seq}(a) \land \alpha \in \mathbb{R} \land \varphi_{68})$}
\AXC{$\varphi_{146}$\lbl{A2}}
\BIC{$\mathrm{Conv}(a,\alpha) \leftrightarrow \mathrm{Seq}(a) \land \alpha \in \mathbb{R} \land \varphi_{68}$}
\end{ptree}
\noindent$\Pi_{76}$:
\begin{ptree}
\AXC{$\stackrel{\Pi_{75}}{\vdots}$}\noLine
\UIC{$\mathrm{Conv}(a,\alpha) \leftrightarrow \mathrm{Seq}(a) \land \alpha \in \mathbb{R} \land \varphi_{68}$}
\AXC{$(\mathrm{Conv}(a,\alpha) \leftrightarrow \mathrm{Seq}(a) \land \alpha \in \mathbb{R} \land \varphi_{68}) \to \varphi_{137}$\lbl{A1}}
\BIC{$\varphi_{137}$}
\end{ptree}
\noindent$\Pi_{77}$:
\begin{ptree}
\AXC{$\mathrm{Seq}(a) \to \mathrm{Conv}(a \circ \mathsf{s},\alpha) \to \mathrm{Conv}(a \circ \mathsf{s},\alpha)$\lbl{A1}}
\AXC{$\stackrel{\Pi_{74}}{\vdots}$}\noLine
\UIC{$\varphi_{74}$}
\BIC{$\mathrm{Seq}(a) \to \mathrm{Conv}(a \circ \mathsf{s},\alpha) \to \alpha \in \mathbb{R}$}
\AXC{$\mathrm{Seq}(a) \to \mathrm{Seq}(a)$\lbl{A1}}
\AXC{$\stackrel{\Pi_{76}}{\vdots}$}\noLine
\UIC{$\varphi_{137}$}
\BIC{$(\mathrm{Seq}(a) \to \mathrm{Conv}(a \circ \mathsf{s},\alpha) \to \alpha \in \mathbb{R}) \to \varphi_{125}$}
\BIC{$\varphi_{125}$}
\end{ptree}
\noindent 主木:
\begin{ptree}
\AXC{$\varphi_{38}$\lbl{A4}}
\AXC{$\stackrel{\Pi_{73}}{\vdots}$}\noLine
\UIC{$\varphi_{38} \to \mathrm{Seq}(a) \to \mathrm{Conv}(a \circ \mathsf{s},\alpha) \to \varphi_{68}$}
\BIC{$\mathrm{Seq}(a) \to \mathrm{Conv}(a \circ \mathsf{s},\alpha) \to \varphi_{68}$}
\AXC{$\stackrel{\Pi_{77}}{\vdots}$}\noLine
\UIC{$\varphi_{125}$}
\BIC{$\mathrm{Seq}(a) \to \mathrm{Conv}(a \circ \mathsf{s},\alpha) \to \mathrm{Conv}(a,\alpha)$}
\UIC{$\varphi_{8}$}
\UIC{$\forall a\,(\varphi_{8})$}
\end{ptree}
\begin{abbrevlist}
\item $\varphi_{1} :\equiv \forall x\,\forall y\,(x \leq y \to \neg y < x) \to \forall y\,(S(N) \leq y \to \neg y < S(N))$
\item $\varphi_{2} :\equiv \forall y\,(S(N) \leq y \to \neg y < S(N)) \to S(N) \leq N \to \neg N < S(N)$
\item $\varphi_{3} :\equiv \forall x\,\forall y\,(x < y \to x \neq y) \to \forall y\,(1_{\omega} < y \to 1_{\omega} \neq y)$
\item $\varphi_{4} :\equiv \forall y\,(1_{\omega} < y \to 1_{\omega} \neq y) \to 1_{\omega} < n \to 1_{\omega} \neq n$
\item $\varphi_{5} :\equiv S(N) \leq n \to m \in \mathbb{N} \land n = S(m) \to N = S(m) \to N = S(m)$
\item $\varphi_{6} :\equiv S(N) \leq n \to m \in \mathbb{N} \land n = S(m) \to N = S(m) \to S(m) = N$
\item $\varphi_{7} :\equiv S(N) \leq n \to m \in \mathbb{N} \land n = S(m) \to N = S(m) \to S(N) \leq N$
\item $\varphi_{8} :\equiv \forall \alpha\,(\mathrm{Seq}(a) \to \mathrm{Conv}(a \circ \mathsf{s},\alpha) \to \mathrm{Conv}(a,\alpha))$
\item $\varphi_{9} :\equiv \forall x\,\forall y\,\forall z\,(x \leq y \to y \leq z \to x \leq z) \to \forall y\,\forall z\,(N \leq y \to y \leq z \to N \leq z)$
\item $\varphi_{10} :\equiv \forall y\,\forall z\,(N \leq y \to y \leq z \to N \leq z) \to \forall z\,(N \leq S(N) \to S(N) \leq z \to N \leq z)$
\item $\varphi_{11} :\equiv S(N) \leq n \to m \in \mathbb{N} \land n = S(m) \to N = S(m) \to \neg N < S(N)$
\item $\varphi_{12} :\equiv n \in \mathbb{N} \to S(N) \leq n \to m \in \mathbb{N} \land n = S(m) \to S(N) \leq S(m)$
\item $\varphi_{13} :\equiv n \in \mathbb{N} \to S(N) \leq n \to m \in \mathbb{N} \land n = S(m) \to \mathrm{NSeq}(\mathsf{s})$
\item $\varphi_{14} :\equiv \forall z\,\forall w\,(z = w \to S(N) = z \to S(N) = w) \to \forall w\,(n = w \to S(N) = n \to S(N) = w)$
\item $\varphi_{15} :\equiv \forall w\,(n = w \to S(N) = n \to S(N) = w) \to n = S(m) \to S(N) = n \to S(N) = S(m)$
\item $\varphi_{16} :\equiv \forall z\,(N \leq S(N) \to S(N) \leq z \to N \leq z) \to N \leq S(N) \to S(N) \leq S(m) \to N \leq S(m)$
\item $\varphi_{17} :\equiv n \in \mathbb{N} \to S(N) \leq n \to m \in \mathbb{N} \land n = S(m) \to N = S(m) \to N \leq m$
\item $\varphi_{18} :\equiv n \in \mathbb{N} \to S(N) \leq n \to m \in \mathbb{N} \land n = S(m) \to N \leq m \lor N = S(m)$
\item $\varphi_{19} :\equiv \forall m\,\forall n\,(m \in \omega \to n \in \omega \to m \leq S(n) \to m \leq n \lor m = S(n))$
\item $\varphi_{20} :\equiv S(N) \leq n \to m \in \mathbb{N} \land n = S(m) \to m \in \mathbb{N} \to \mathsf{s}(m) = S(m)$
\item $\varphi_{21} :\equiv n \in \mathbb{N} \to S(N) \leq n \to m \in \mathbb{N} \land n = S(m) \to \mathsf{s}(m) = S(m)$
\item $\varphi_{22} :\equiv \forall z\,\forall w\,(z = w \to S(N) = z \to S(N) = w) \to \forall w\,(S(m) = w \to S(N) = S(m) \to S(N) = w)$
\item $\varphi_{23} :\equiv \forall w\,(S(m) = w \to S(N) = S(m) \to S(N) = w) \to S(m) = N \to S(N) = S(m) \to S(N) = N$
\item $\varphi_{24} :\equiv n \in \mathbb{N} \to S(N) \leq n \to m \in \mathbb{N} \land n = S(m) \to m \in \mathbb{N} \land n = S(m)$
\item $\varphi_{25} :\equiv S(N) \leq n \to m \in \mathbb{N} \land n = S(m) \to \forall k \in \mathbb{N}\,(\mathsf{s}(k) = S(k))$
\item $\varphi_{26} :\equiv \forall k \in \mathbb{N}\,(\mathsf{s}(k) = S(k)) \to m \in \mathbb{N} \to \mathsf{s}(m) = S(m)$
\item $\varphi_{27} :\equiv n \in \mathbb{N} \to S(N) \leq n \to m \in \mathbb{N} \land n = S(m) \to a_{S(m)} = a_{n}$
\item $\varphi_{28} :\equiv \forall N\,(\neg (N \in \mathbb{N} \land \forall n \in \mathbb{N}\,(N \leq n \to \lvert a_{n} - \alpha \rvert < \varepsilon)))$
\item $\varphi_{29} :\equiv \forall z\,\forall w\,(z = w \to S(N) < z \to S(N) < w) \to \forall w\,(n = w \to S(N) < n \to S(N) < w)$
\item $\varphi_{30} :\equiv \forall w\,(n = w \to S(N) < n \to S(N) < w) \to n = S(m) \to S(N) < n \to S(N) < S(m)$
\item $\varphi_{31} :\equiv \forall x\,\forall y\,\forall z\,(x \leq y \to y < z \to x < z) \to \forall y\,\forall z\,(1_{\omega} \leq y \to y < z \to 1_{\omega} < z)$
\item $\varphi_{32} :\equiv \forall y\,\forall z\,(1_{\omega} \leq y \to y < z \to 1_{\omega} < z) \to \forall z\,(1_{\omega} \leq N \to N < z \to 1_{\omega} < z)$
\item $\varphi_{33} :\equiv \forall z\,(1_{\omega} \leq N \to N < z \to 1_{\omega} < z) \to 1_{\omega} \leq N \to N < S(N) \to 1_{\omega} < S(N)$
\item $\varphi_{34} :\equiv \forall x\,\forall y\,\forall z\,(x < y \to y \leq z \to x < z) \to \forall y\,\forall z\,(1_{\omega} < y \to y \leq z \to 1_{\omega} < z)$
\item $\varphi_{35} :\equiv \neg (S(N) \in \mathbb{N} \land \forall n \in \mathbb{N}\,(S(N) \leq n \to \lvert a_{n} - \alpha \rvert < \varepsilon))$
\item $\varphi_{36} :\equiv \mathrm{Seq}(a) \to \mathrm{Conv}(a \circ \mathsf{s},\alpha) \to \varepsilon \in \mathbb{R} \to 0 < \varepsilon \to 0 < \varepsilon$
\item $\varphi_{37} :\equiv S(N) \leq n \to m \in \mathbb{N} \land n = S(m) \to \lvert a \circ \mathsf{s}(m) - \alpha \rvert < \varepsilon$
\item $\varphi_{38} :\equiv \mathrm{Seq}(a) \land \mathrm{Conv}(a \circ \mathsf{s},\alpha) \to \forall \varepsilon\,(\mathrm{Seq}(a) \land \mathrm{Conv}(a \circ \mathsf{s},\alpha))$
\item $\varphi_{39} :\equiv n \in \mathbb{N} \to S(N) \leq n \to m \in \mathbb{N} \land n = S(m) \to \lvert a_{n} - \alpha \rvert < \varepsilon$
\item $\varphi_{40} :\equiv n \in \mathbb{N} \to S(N) \leq n \to m \in \mathbb{N} \land n = S(m) \to a \circ \mathsf{s}(m) = a_{n}$
\item $\varphi_{41} :\equiv \neg (N \in \mathbb{N} \land \forall n \in \mathbb{N}\,(N \leq n \to \lvert a \circ \mathsf{s}(n) - \alpha \rvert < \varepsilon))$
\item $\varphi_{42} :\equiv \forall z\,\forall w\,(z = w \to S(N) < z \to S(N) < w) \to \forall w\,(S(m) = w \to S(N) < S(m) \to S(N) < w)$
\item $\varphi_{43} :\equiv \forall w\,(S(m) = w \to S(N) < S(m) \to S(N) < w) \to S(m) = N \to S(N) < S(m) \to S(N) < N$
\item $\varphi_{44} :\equiv \forall y\,\forall z\,(1_{\omega} < y \to y \leq z \to 1_{\omega} < z) \to \forall z\,(1_{\omega} < S(N) \to S(N) \leq z \to 1_{\omega} < z)$
\item $\varphi_{45} :\equiv \forall z\,(1_{\omega} < S(N) \to S(N) \leq z \to 1_{\omega} < z) \to 1_{\omega} < S(N) \to S(N) \leq n \to 1_{\omega} < n$
\item $\varphi_{46} :\equiv n \in \mathbb{N} \to S(N) \leq n \to m \in \mathbb{N} \land n = S(m) \to a_{\mathsf{s}(m)} = a_{S(m)}$
\item $\varphi_{47} :\equiv \forall a\,\forall \alpha\,(\mathrm{Conv}(a,\alpha) \to \alpha \in \mathbb{R}) \to \forall \alpha\,(\mathrm{Conv}(a \circ \mathsf{s},\alpha) \to \alpha \in \mathbb{R})$
\item $\varphi_{48} :\equiv n \in \mathbb{N} \to S(N) \leq n \to m \in \mathbb{N} \land n = S(m) \to a \circ \mathsf{s}(m) = a_{S(m)}$
\item $\varphi_{49} :\equiv \forall z\,(a \circ \mathsf{s}(m) = a_{S(m)} \to a_{S(m)} = z \to a \circ \mathsf{s}(m) = z)$
\item $\varphi_{50} :\equiv m \in \mathbb{N} \land n = S(m) \to m \in \mathbb{N} \to N \leq m \to \lvert a \circ \mathsf{s}(m) - \alpha \rvert < \varepsilon$
\item $\varphi_{51} :\equiv \forall \alpha\,(\mathrm{Conv}(a \circ \mathsf{s},\alpha) \to \alpha \in \mathbb{R}) \to \mathrm{Conv}(a \circ \mathsf{s},\alpha) \to \alpha \in \mathbb{R}$
\item $\varphi_{52} :\equiv m \in \mathbb{N} \land n = S(m) \to \forall n \in \mathbb{N}\,(N \leq n \to \lvert a \circ \mathsf{s}(n) - \alpha \rvert < \varepsilon)$
\item $\varphi_{53} :\equiv S(N) \leq n \to m \in \mathbb{N} \land n = S(m) \to m \in \mathbb{N} \to a \circ \mathsf{s}(m) = a_{\mathsf{s}(m)}$
\item $\varphi_{54} :\equiv n \in \mathbb{N} \to S(N) \leq n \to m \in \mathbb{N} \land n = S(m) \to a \circ \mathsf{s}(m) = a_{\mathsf{s}(m)}$
\item $\varphi_{55} :\equiv N \in \mathbb{N} \land \forall n \in \mathbb{N}\,(N \leq n \to \lvert a \circ \mathsf{s}(n) - \alpha \rvert < \varepsilon) \to N \in \mathbb{N}$
\item $\varphi_{56} :\equiv N \in \mathbb{N} \land \forall n \in \mathbb{N}\,(N \leq n \to \lvert a \circ \mathsf{s}(n) - \alpha \rvert < \varepsilon) \to N < S(N)$
\item $\varphi_{57} :\equiv S(N) \leq n \to m \in \mathbb{N} \land n = S(m) \to \forall k \in \mathbb{N}\,(a \circ \mathsf{s}(k) = a_{\mathsf{s}(k)})$
\item $\varphi_{58} :\equiv \forall x\,\forall y\,\forall z\,(x = y \to y = z \to x = z) \to \forall y\,\forall z\,(a \circ \mathsf{s}(m) = y \to y = z \to a \circ \mathsf{s}(m) = z)$
\item $\varphi_{59} :\equiv N \in \mathbb{N} \land \forall n \in \mathbb{N}\,(N \leq n \to \lvert a \circ \mathsf{s}(n) - \alpha \rvert < \varepsilon) \to N \in \omega$
\item $\varphi_{60} :\equiv 0 < \varepsilon \to \exists N \in \mathbb{N}\,\forall n \in \mathbb{N}\,(N \leq n \to \lvert a \circ \mathsf{s}(n) - \alpha \rvert < \varepsilon)$
\item $\varphi_{61} :\equiv N \in \mathbb{N} \land \forall n \in \mathbb{N}\,(N \leq n \to \lvert a \circ \mathsf{s}(n) - \alpha \rvert < \varepsilon) \to S(N) \in \mathbb{N}$
\item $\varphi_{62} :\equiv \forall w\,(a \circ \mathsf{s}(m) = w \to \lvert a \circ \mathsf{s}(m) - \alpha \rvert < \varepsilon \to \lvert w - \alpha \rvert < \varepsilon)$
\item $\varphi_{63} :\equiv \forall s\,(\mathrm{Seq}(a) \to \mathrm{NSeq}(s) \to \mathrm{Seq}(a \circ s) \land \forall k \in \mathbb{N}\,(a \circ s(k) = a_{s_{k}}))$
\item $\varphi_{64} :\equiv \mathrm{NSeq}(\mathsf{s}) \to \mathrm{Seq}(a \circ \mathsf{s}) \land \forall k \in \mathbb{N}\,(a \circ \mathsf{s}(k) = a_{\mathsf{s}(k)})$
\item $\varphi_{65} :\equiv \forall z\,(a \circ \mathsf{s}(m) = a_{\mathsf{s}(m)} \to a_{\mathsf{s}(m)} = z \to a \circ \mathsf{s}(m) = z)$
\item $\varphi_{66} :\equiv \varepsilon \in \mathbb{R} \to 0 < \varepsilon \to \exists N \in \mathbb{N}\,\forall n \in \mathbb{N}\,(N \leq n \to \lvert a_{n} - \alpha \rvert < \varepsilon)$
\item $\varphi_{67} :\equiv \forall n \in \mathbb{N}\,(n \neq 1_{\omega} \to \exists m \in \mathbb{N}\,(n = S(m))) \to n \in \mathbb{N} \to n \neq 1_{\omega} \to \exists m \in \mathbb{N}\,(n = S(m))$
\item $\varphi_{68} :\equiv \forall \varepsilon \in \mathbb{R}\,(0 < \varepsilon \to \exists N \in \mathbb{N}\,\forall n \in \mathbb{N}\,(N \leq n \to \lvert a_{n} - \alpha \rvert < \varepsilon))$
\item $\varphi_{69} :\equiv m \in \mathbb{N} \land n = S(m) \to \mathrm{Seq}(a \circ \mathsf{s}) \land \forall k \in \mathbb{N}\,(a \circ \mathsf{s}(k) = a_{\mathsf{s}(k)})$
\item $\varphi_{70} :\equiv N \in \mathbb{N} \land \forall n \in \mathbb{N}\,(N \leq n \to \lvert a \circ \mathsf{s}(n) - \alpha \rvert < \varepsilon) \to n \in \mathbb{N} \to n \in \mathbb{N}$
\item $\varphi_{71} :\equiv a \circ \mathsf{s}(m) = a_{n} \to \lvert a \circ \mathsf{s}(m) - \alpha \rvert < \varepsilon \to \lvert a_{n} - \alpha \rvert < \varepsilon$
\item $\varphi_{72} :\equiv N \in \mathbb{N} \land \forall n \in \mathbb{N}\,(N \leq n \to \lvert a \circ \mathsf{s}(n) - \alpha \rvert < \varepsilon) \to n \in \mathbb{N} \to S(N) \leq n \to S(N) \leq n$
\item $\varphi_{73} :\equiv N \in \mathbb{N} \land \forall n \in \mathbb{N}\,(N \leq n \to \lvert a \circ \mathsf{s}(n) - \alpha \rvert < \varepsilon) \to n \in \mathbb{N} \to S(N) \leq n \to 1_{\omega} \leq N$
\item $\varphi_{74} :\equiv (\mathrm{Seq}(a) \to \mathrm{Conv}(a \circ \mathsf{s},\alpha) \to \mathrm{Conv}(a \circ \mathsf{s},\alpha)) \to \mathrm{Seq}(a) \to \mathrm{Conv}(a \circ \mathsf{s},\alpha) \to \alpha \in \mathbb{R}$
\item $\varphi_{75} :\equiv N \in \mathbb{N} \land \forall n \in \mathbb{N}\,(N \leq n \to \lvert a \circ \mathsf{s}(n) - \alpha \rvert < \varepsilon) \to n \in \mathbb{N} \to S(N) \leq n \to 1_{\omega} \neq n$
\item $\varphi_{76} :\equiv N \in \mathbb{N} \land \forall n \in \mathbb{N}\,(N \leq n \to \lvert a \circ \mathsf{s}(n) - \alpha \rvert < \varepsilon) \to n \in \mathbb{N} \to S(N) \leq n \to n \neq 1_{\omega}$
\item $\varphi_{77} :\equiv N \in \mathbb{N} \land \forall n \in \mathbb{N}\,(N \leq n \to \lvert a \circ \mathsf{s}(n) - \alpha \rvert < \varepsilon) \to n \in \mathbb{N} \to S(N) \leq n \to 1_{\omega} < n$
\item $\varphi_{78} :\equiv \forall n\,(N \in \omega \to n \in \omega \to N \leq S(n) \to N \leq n \lor N = S(n)) \to N \in \omega \to m \in \omega \to N \leq S(m) \to N \leq m \lor N = S(m)$
\item $\varphi_{79} :\equiv a \circ \mathsf{s}(m) = a_{\mathsf{s}(m)} \to a_{\mathsf{s}(m)} = a_{S(m)} \to a \circ \mathsf{s}(m) = a_{S(m)}$
\item $\varphi_{80} :\equiv \forall k \in \mathbb{N}\,(a \circ \mathsf{s}(k) = a_{\mathsf{s}(k)}) \to m \in \mathbb{N} \to a \circ \mathsf{s}(m) = a_{\mathsf{s}(m)}$
\item $\varphi_{81} :\equiv N \in \mathbb{N} \land \forall n \in \mathbb{N}\,(N \leq n \to \lvert a \circ \mathsf{s}(n) - \alpha \rvert < \varepsilon) \to n \in \mathbb{N} \to S(N) \leq n \to 1_{\omega} < S(N)$
\item $\varphi_{82} :\equiv \varphi_{19} \to \forall n\,(N \in \omega \to n \in \omega \to N \leq S(n) \to N \leq n \lor N = S(n))$
\item $\varphi_{83} :\equiv N \in \mathbb{N} \land \forall n \in \mathbb{N}\,(N \leq n \to \lvert a \circ \mathsf{s}(n) - \alpha \rvert < \varepsilon) \to n \in \mathbb{N} \to S(N) \leq n \to \exists m \in \mathbb{N}\,(n = S(m))$
\item $\varphi_{84} :\equiv \forall n \in \mathbb{N}\,(N \leq n \to \lvert a \circ \mathsf{s}(n) - \alpha \rvert < \varepsilon) \to m \in \mathbb{N} \to N \leq m \to \lvert a \circ \mathsf{s}(m) - \alpha \rvert < \varepsilon$
\item $\varphi_{85} :\equiv \forall z\,\forall w\,(z = w \to a_{S(m)} = a_{z} \to a_{S(m)} = a_{w}) \to \forall w\,(S(m) = w \to a_{S(m)} = a_{S(m)} \to a_{S(m)} = a_{w})$
\item $\varphi_{86} :\equiv \forall w\,(S(m) = w \to a_{S(m)} = a_{S(m)} \to a_{S(m)} = a_{w}) \to S(m) = n \to a_{S(m)} = a_{S(m)} \to a_{S(m)} = a_{n}$
\item $\varphi_{87} :\equiv N \in \mathbb{N} \land \forall n \in \mathbb{N}\,(N \leq n \to \lvert a \circ \mathsf{s}(n) - \alpha \rvert < \varepsilon) \to n \in \mathbb{N} \to S(N) \leq n \to m \in \mathbb{N} \land n = S(m) \to N \leq m$
\item $\varphi_{88} :\equiv N \in \mathbb{N} \land \forall n \in \mathbb{N}\,(N \leq n \to \lvert a \circ \mathsf{s}(n) - \alpha \rvert < \varepsilon) \to n \in \mathbb{N} \to S(N) \leq n \to \lvert a_{n} - \alpha \rvert < \varepsilon$
\item $\varphi_{89} :\equiv N \in \mathbb{N} \land \forall n \in \mathbb{N}\,(N \leq n \to \lvert a \circ \mathsf{s}(n) - \alpha \rvert < \varepsilon) \to n \in \mathbb{N} \to S(N) \leq n \to m \in \mathbb{N} \land n = S(m) \to n = S(m)$
\item $\varphi_{90} :\equiv N \in \mathbb{N} \land \forall n \in \mathbb{N}\,(N \leq n \to \lvert a \circ \mathsf{s}(n) - \alpha \rvert < \varepsilon) \to n \in \mathbb{N} \to S(N) \leq n \to m \in \mathbb{N} \land n = S(m) \to N \leq S(m)$
\item $\varphi_{91} :\equiv N \in \mathbb{N} \land \forall n \in \mathbb{N}\,(N \leq n \to \lvert a \circ \mathsf{s}(n) - \alpha \rvert < \varepsilon) \to n \in \mathbb{N} \to S(N) \leq n \to m \in \mathbb{N} \land n = S(m) \to S(m) = n$
\item $\varphi_{92} :\equiv N \in \mathbb{N} \land \forall n \in \mathbb{N}\,(N \leq n \to \lvert a \circ \mathsf{s}(n) - \alpha \rvert < \varepsilon) \to n \in \mathbb{N} \to S(N) \leq n \to m \in \mathbb{N} \land n = S(m) \to \mathrm{Seq}(a)$
\item $\varphi_{93} :\equiv \mathrm{Seq}(a) \land \mathrm{Conv}(a \circ \mathsf{s},\alpha) \land \varepsilon \in \mathbb{R} \land 0 < \varepsilon \land \neg \exists N \in \mathbb{N}\,\forall n \in \mathbb{N}\,(N \leq n \to \lvert a_{n} - \alpha \rvert < \varepsilon)$
\item $\varphi_{94} :\equiv N \in \mathbb{N} \land \forall n \in \mathbb{N}\,(N \leq n \to \lvert a \circ \mathsf{s}(n) - \alpha \rvert < \varepsilon) \to \forall n \in \mathbb{N}\,(S(N) \leq n \to \lvert a_{n} - \alpha \rvert < \varepsilon)$
\item $\varphi_{95} :\equiv N \in \mathbb{N} \land \forall n \in \mathbb{N}\,(N \leq n \to \lvert a \circ \mathsf{s}(n) - \alpha \rvert < \varepsilon) \to n \in \mathbb{N} \to S(N) \leq n \to m \in \mathbb{N} \land n = S(m) \to m \in \mathbb{N}$
\item $\varphi_{96} :\equiv N \in \mathbb{N} \land \forall n \in \mathbb{N}\,(N \leq n \to \lvert a \circ \mathsf{s}(n) - \alpha \rvert < \varepsilon) \to n \in \mathbb{N} \to S(N) \leq n \to m \in \mathbb{N} \land n = S(m) \to m \in \omega$
\item $\varphi_{97} :\equiv \mathrm{Seq}(a) \land \mathrm{Conv}(a \circ \mathsf{s},\alpha) \land \varepsilon \in \mathbb{R} \land 0 < \varepsilon \land (N \in \mathbb{N} \land \forall n \in \mathbb{N}\,(N \leq n \to \lvert a \circ \mathsf{s}(n) - \alpha \rvert < \varepsilon))$
\item $\varphi_{98} :\equiv N \in \mathbb{N} \land \forall n \in \mathbb{N}\,(N \leq n \to \lvert a \circ \mathsf{s}(n) - \alpha \rvert < \varepsilon) \to n \in \mathbb{N} \to \varphi_{5}$
\item $\varphi_{99} :\equiv N \in \mathbb{N} \land \forall n \in \mathbb{N}\,(N \leq n \to \lvert a \circ \mathsf{s}(n) - \alpha \rvert < \varepsilon) \to n \in \mathbb{N} \to \varphi_{6}$
\item $\varphi_{100} :\equiv N \in \mathbb{N} \land \forall n \in \mathbb{N}\,(N \leq n \to \lvert a \circ \mathsf{s}(n) - \alpha \rvert < \varepsilon) \to n \in \mathbb{N} \to \varphi_{7}$
\item $\varphi_{101} :\equiv N \in \mathbb{N} \land \forall n \in \mathbb{N}\,(N \leq n \to \lvert a \circ \mathsf{s}(n) - \alpha \rvert < \varepsilon) \to n \in \mathbb{N} \to \varphi_{11}$
\item $\varphi_{102} :\equiv 0 < \varepsilon \to N \in \mathbb{N} \land \forall n \in \mathbb{N}\,(N \leq n \to \lvert a \circ \mathsf{s}(n) - \alpha \rvert < \varepsilon) \to \varphi_{12}$
\item $\varphi_{103} :\equiv N \in \mathbb{N} \land \forall n \in \mathbb{N}\,(N \leq n \to \lvert a \circ \mathsf{s}(n) - \alpha \rvert < \varepsilon) \to \exists N \in \mathbb{N}\,\forall n \in \mathbb{N}\,(N \leq n \to \lvert a_{n} - \alpha \rvert < \varepsilon)$
\item $\varphi_{104} :\equiv 0 < \varepsilon \to N \in \mathbb{N} \land \forall n \in \mathbb{N}\,(N \leq n \to \lvert a \circ \mathsf{s}(n) - \alpha \rvert < \varepsilon) \to \varphi_{13}$
\item $\varphi_{105} :\equiv 0 < \varepsilon \to N \in \mathbb{N} \land \forall n \in \mathbb{N}\,(N \leq n \to \lvert a \circ \mathsf{s}(n) - \alpha \rvert < \varepsilon) \to \varphi_{17}$
\item $\varphi_{106} :\equiv 0 < \varepsilon \to N \in \mathbb{N} \land \forall n \in \mathbb{N}\,(N \leq n \to \lvert a \circ \mathsf{s}(n) - \alpha \rvert < \varepsilon) \to \varphi_{18}$
\item $\varphi_{107} :\equiv N \in \mathbb{N} \land \forall n \in \mathbb{N}\,(N \leq n \to \lvert a \circ \mathsf{s}(n) - \alpha \rvert < \varepsilon) \to S(N) \in \mathbb{N} \land \forall n \in \mathbb{N}\,(S(N) \leq n \to \lvert a_{n} - \alpha \rvert < \varepsilon)$
\item $\varphi_{108} :\equiv 0 < \varepsilon \to N \in \mathbb{N} \land \forall n \in \mathbb{N}\,(N \leq n \to \lvert a \circ \mathsf{s}(n) - \alpha \rvert < \varepsilon) \to \varphi_{21}$
\item $\varphi_{109} :\equiv N \in \mathbb{N} \land \forall n \in \mathbb{N}\,(N \leq n \to \lvert a \circ \mathsf{s}(n) - \alpha \rvert < \varepsilon) \to N \in \mathbb{N} \land \forall n \in \mathbb{N}\,(N \leq n \to \lvert a \circ \mathsf{s}(n) - \alpha \rvert < \varepsilon)$
\item $\varphi_{110} :\equiv N \in \mathbb{N} \land \forall n \in \mathbb{N}\,(N \leq n \to \lvert a \circ \mathsf{s}(n) - \alpha \rvert < \varepsilon) \to n \in \mathbb{N} \to \varphi_{20}$
\item $\varphi_{111} :\equiv 0 < \varepsilon \to N \in \mathbb{N} \land \forall n \in \mathbb{N}\,(N \leq n \to \lvert a \circ \mathsf{s}(n) - \alpha \rvert < \varepsilon) \to \varphi_{24}$
\item $\varphi_{112} :\equiv 0 < \varepsilon \to N \in \mathbb{N} \land \forall n \in \mathbb{N}\,(N \leq n \to \lvert a \circ \mathsf{s}(n) - \alpha \rvert < \varepsilon) \to \varphi_{27}$
\item $\varphi_{113} :\equiv N \in \mathbb{N} \land \forall n \in \mathbb{N}\,(N \leq n \to \lvert a \circ \mathsf{s}(n) - \alpha \rvert < \varepsilon) \to n \in \mathbb{N} \to \varphi_{25}$
\item $\varphi_{114} :\equiv N \in \mathbb{N} \land \forall n \in \mathbb{N}\,(N \leq n \to \lvert a \circ \mathsf{s}(n) - \alpha \rvert < \varepsilon) \to n \in \mathbb{N} \to \varphi_{37}$
\item $\varphi_{115} :\equiv 0 < \varepsilon \to N \in \mathbb{N} \land \forall n \in \mathbb{N}\,(N \leq n \to \lvert a \circ \mathsf{s}(n) - \alpha \rvert < \varepsilon) \to \varphi_{39}$
\item $\varphi_{116} :\equiv 0 < \varepsilon \to N \in \mathbb{N} \land \forall n \in \mathbb{N}\,(N \leq n \to \lvert a \circ \mathsf{s}(n) - \alpha \rvert < \varepsilon) \to \varphi_{40}$
\item $\varphi_{117} :\equiv 0 < \varepsilon \to N \in \mathbb{N} \land \forall n \in \mathbb{N}\,(N \leq n \to \lvert a \circ \mathsf{s}(n) - \alpha \rvert < \varepsilon) \to \varphi_{46}$
\item $\varphi_{118} :\equiv 0 < \varepsilon \to N \in \mathbb{N} \land \forall n \in \mathbb{N}\,(N \leq n \to \lvert a \circ \mathsf{s}(n) - \alpha \rvert < \varepsilon) \to \varphi_{48}$
\item $\varphi_{119} :\equiv \forall z\,\forall w\,(z = w \to a_{\mathsf{s}(m)} = a_{z} \to a_{\mathsf{s}(m)} = a_{w}) \to \forall w\,(\mathsf{s}(m) = w \to a_{\mathsf{s}(m)} = a_{\mathsf{s}(m)} \to a_{\mathsf{s}(m)} = a_{w})$
\item $\varphi_{120} :\equiv \varphi_{49} \to a \circ \mathsf{s}(m) = a_{S(m)} \to a_{S(m)} = a_{n} \to a \circ \mathsf{s}(m) = a_{n}$
\item $\varphi_{121} :\equiv 0 < \varepsilon \to N \in \mathbb{N} \land \forall n \in \mathbb{N}\,(N \leq n \to \lvert a \circ \mathsf{s}(n) - \alpha \rvert < \varepsilon) \to \varphi_{54}$
\item $\varphi_{122} :\equiv N \in \mathbb{N} \land \forall n \in \mathbb{N}\,(N \leq n \to \lvert a \circ \mathsf{s}(n) - \alpha \rvert < \varepsilon) \to n \in \mathbb{N} \to \varphi_{53}$
\item $\varphi_{123} :\equiv N \in \mathbb{N} \land \forall n \in \mathbb{N}\,(N \leq n \to \lvert a \circ \mathsf{s}(n) - \alpha \rvert < \varepsilon) \to n \in \mathbb{N} \to \varphi_{57}$
\item $\varphi_{124} :\equiv N \in \mathbb{N} \land \forall n \in \mathbb{N}\,(N \leq n \to \lvert a \circ \mathsf{s}(n) - \alpha \rvert < \varepsilon) \to n \in \mathbb{N} \to S(N) \leq n \to \varphi_{50}$
\item $\varphi_{125} :\equiv (\mathrm{Seq}(a) \to \mathrm{Conv}(a \circ \mathsf{s},\alpha) \to \varphi_{68}) \to \mathrm{Seq}(a) \to \mathrm{Conv}(a \circ \mathsf{s},\alpha) \to \mathrm{Conv}(a,\alpha)$
\item $\varphi_{126} :\equiv N \in \mathbb{N} \land \forall n \in \mathbb{N}\,(N \leq n \to \lvert a \circ \mathsf{s}(n) - \alpha \rvert < \varepsilon) \to n \in \mathbb{N} \to S(N) \leq n \to \varphi_{52}$
\item $\varphi_{127} :\equiv \forall w\,(\mathsf{s}(m) = w \to a_{\mathsf{s}(m)} = a_{\mathsf{s}(m)} \to a_{\mathsf{s}(m)} = a_{w}) \to \mathsf{s}(m) = S(m) \to a_{\mathsf{s}(m)} = a_{\mathsf{s}(m)} \to a_{\mathsf{s}(m)} = a_{S(m)}$
\item $\varphi_{128} :\equiv (n = 1_{\omega} \to 1_{\omega} = n) \to \mathrm{Seq}(a) \to \mathrm{Conv}(a \circ \mathsf{s},\alpha) \to \varepsilon \in \mathbb{R} \to 0 < \varepsilon \to \varphi_{76}$
\item $\varphi_{129} :\equiv N \in \mathbb{N} \land \forall n \in \mathbb{N}\,(N \leq n \to \lvert a \circ \mathsf{s}(n) - \alpha \rvert < \varepsilon) \to n \in \mathbb{N} \to S(N) \leq n \to \varphi_{69}$
\item $\varphi_{130} :\equiv (n = S(m) \to S(m) = n) \to \mathrm{Seq}(a) \to \mathrm{Conv}(a \circ \mathsf{s},\alpha) \to \varepsilon \in \mathbb{R} \to 0 < \varepsilon \to \varphi_{91}$
\item $\varphi_{131} :\equiv \varphi_{97} \land n \in \mathbb{N} \land S(N) \leq n \land \neg \lvert a_{n} - \alpha \rvert < \varepsilon \to \neg (m \in \mathbb{N} \land n = S(m))$
\item $\varphi_{132} :\equiv (N = S(m) \to S(m) = N) \to \mathrm{Seq}(a) \to \mathrm{Conv}(a \circ \mathsf{s},\alpha) \to \varepsilon \in \mathbb{R} \to 0 < \varepsilon \to \varphi_{99}$
\item $\varphi_{133} :\equiv \forall m\,(\varphi_{97} \land n \in \mathbb{N} \land S(N) \leq n \land \neg \lvert a_{n} - \alpha \rvert < \varepsilon) \to \forall m\,(\neg (m \in \mathbb{N} \land n = S(m)))$
\item $\varphi_{134} :\equiv (n = S(m) \to S(N) = n \to S(N) = S(m)) \to \mathrm{Seq}(a) \to \mathrm{Conv}(a \circ \mathsf{s},\alpha) \to \varepsilon \in \mathbb{R} \to \varphi_{102}$
\item $\varphi_{135} :\equiv (S(m) = N \to S(N) = S(m) \to S(N) = N) \to \mathrm{Seq}(a) \to \mathrm{Conv}(a \circ \mathsf{s},\alpha) \to \varepsilon \in \mathbb{R} \to 0 < \varepsilon \to \varphi_{100}$
\item $\varphi_{136} :\equiv \mathrm{NSeq}(\mathsf{s}) \land \forall k \in \mathbb{N}\,(\mathsf{s}(k) = S(k)) \to \mathrm{Seq}(a) \to \mathrm{Conv}(a \circ \mathsf{s},\alpha) \to \varepsilon \in \mathbb{R} \to \varphi_{104}$
\item $\varphi_{137} :\equiv (\mathrm{Seq}(a) \to \mathrm{Seq}(a)) \to (\mathrm{Seq}(a) \to \mathrm{Conv}(a \circ \mathsf{s},\alpha) \to \alpha \in \mathbb{R}) \to \varphi_{125}$
\item $\varphi_{138} :\equiv (\mathrm{Seq}(a) \to \mathrm{Conv}(a \circ \mathsf{s},\alpha) \to \varphi_{66}) \to \mathrm{Seq}(a) \land \mathrm{Conv}(a \circ \mathsf{s},\alpha) \to \varphi_{66}$
\item $\varphi_{139} :\equiv \forall a\,\forall \alpha\,\forall \varepsilon\,(\mathrm{Conv}(a,\alpha) \to \varphi_{66}) \to \forall \alpha\,\forall \varepsilon\,(\mathrm{Conv}(a \circ \mathsf{s},\alpha) \to \varepsilon \in \mathbb{R} \to \varphi_{60})$
\item $\varphi_{140} :\equiv (S(m) = n \to a_{S(m)} = a_{S(m)} \to a_{S(m)} = a_{n}) \to \mathrm{Seq}(a) \to \mathrm{Conv}(a \circ \mathsf{s},\alpha) \to \varepsilon \in \mathbb{R} \to \varphi_{112}$
\item $\varphi_{141} :\equiv \mathrm{NSeq}(\mathsf{s}) \land \forall k \in \mathbb{N}\,(\mathsf{s}(k) = S(k)) \to \mathrm{Seq}(a) \to \mathrm{Conv}(a \circ \mathsf{s},\alpha) \to \varepsilon \in \mathbb{R} \to 0 < \varepsilon \to \varphi_{113}$
\item $\varphi_{142} :\equiv \forall \varepsilon\,(\mathrm{Conv}(a \circ \mathsf{s},\alpha) \to \varepsilon \in \mathbb{R} \to \varphi_{60}) \to \mathrm{Conv}(a \circ \mathsf{s},\alpha) \to \varepsilon \in \mathbb{R} \to \varphi_{60}$
\item $\varphi_{143} :\equiv \forall \varepsilon\,(\mathrm{Seq}(a) \land \mathrm{Conv}(a \circ \mathsf{s},\alpha) \to \varphi_{66}) \to \forall \varepsilon\,(\mathrm{Seq}(a) \land \mathrm{Conv}(a \circ \mathsf{s},\alpha)) \to \varphi_{68}$
\item $\varphi_{144} :\equiv \forall \alpha\,\forall \varepsilon\,(\mathrm{Conv}(a \circ \mathsf{s},\alpha) \to \varepsilon \in \mathbb{R} \to \varphi_{60}) \to \forall \varepsilon\,(\mathrm{Conv}(a \circ \mathsf{s},\alpha) \to \varepsilon \in \mathbb{R} \to \varphi_{60})$
\item $\varphi_{145} :\equiv (\mathrm{Seq}(a) \to \mathrm{Conv}(a \circ \mathsf{s},\alpha) \to \varepsilon \in \mathbb{R} \to \varepsilon \in \mathbb{R}) \to \varphi_{36} \to \mathrm{Seq}(a) \to \mathrm{Conv}(a \circ \mathsf{s},\alpha) \to \varepsilon \in \mathbb{R} \to \varphi_{60}$
\item $\varphi_{146} :\equiv \forall \alpha\,(\mathrm{Conv}(a,\alpha) \leftrightarrow \mathrm{Seq}(a) \land \alpha \in \mathbb{R} \land \varphi_{68}) \to (\mathrm{Conv}(a,\alpha) \leftrightarrow \mathrm{Seq}(a) \land \alpha \in \mathbb{R} \land \varphi_{68})$
\item $\varphi_{147} :\equiv \varphi_{26} \to \mathrm{Seq}(a) \to \mathrm{Conv}(a \circ \mathsf{s},\alpha) \to \varepsilon \in \mathbb{R} \to 0 < \varepsilon \to \varphi_{110}$
\item $\varphi_{148} :\equiv \forall a\,\forall \alpha\,(\mathrm{Conv}(a,\alpha) \leftrightarrow \mathrm{Seq}(a) \land \alpha \in \mathbb{R} \land \varphi_{68}) \to \forall \alpha\,(\mathrm{Conv}(a,\alpha) \leftrightarrow \mathrm{Seq}(a) \land \alpha \in \mathbb{R} \land \varphi_{68})$
\item $\varphi_{149} :\equiv (\mathrm{Seq}(a) \to \mathrm{Conv}(a \circ \mathsf{s},\alpha) \to \varepsilon \in \mathbb{R} \to 0 < \varepsilon \to \varphi_{55}) \to \mathrm{Seq}(a) \to \mathrm{Conv}(a \circ \mathsf{s},\alpha) \to \varepsilon \in \mathbb{R} \to 0 < \varepsilon \to \varphi_{59}$
\item $\varphi_{150} :\equiv (\mathrm{Seq}(a) \to \mathrm{Conv}(a \circ \mathsf{s},\alpha) \to \varepsilon \in \mathbb{R} \to 0 < \varepsilon \to \varphi_{59}) \to \mathrm{Seq}(a) \to \mathrm{Conv}(a \circ \mathsf{s},\alpha) \to \varepsilon \in \mathbb{R} \to 0 < \varepsilon \to \varphi_{56}$
\item $\varphi_{151} :\equiv (\mathrm{Seq}(a) \to \mathrm{Conv}(a \circ \mathsf{s},\alpha) \to \varepsilon \in \mathbb{R} \to 0 < \varepsilon \to \varphi_{55}) \to \mathrm{Seq}(a) \to \mathrm{Conv}(a \circ \mathsf{s},\alpha) \to \varepsilon \in \mathbb{R} \to 0 < \varepsilon \to \varphi_{61}$
\item $\varphi_{152} :\equiv (a \circ \mathsf{s}(m) = a_{S(m)} \to a_{S(m)} = a_{n} \to a \circ \mathsf{s}(m) = a_{n}) \to \mathrm{Seq}(a) \to \mathrm{Conv}(a \circ \mathsf{s},\alpha) \to \varepsilon \in \mathbb{R} \to \varphi_{116}$
\item $\varphi_{153} :\equiv (\mathsf{s}(m) = S(m) \to a_{\mathsf{s}(m)} = a_{\mathsf{s}(m)} \to a_{\mathsf{s}(m)} = a_{S(m)}) \to \mathrm{Seq}(a) \to \mathrm{Conv}(a \circ \mathsf{s},\alpha) \to \varepsilon \in \mathbb{R} \to \varphi_{117}$
\item $\varphi_{154} :\equiv (\mathrm{Seq}(a) \to \mathrm{Conv}(a \circ \mathsf{s},\alpha) \to \varepsilon \in \mathbb{R} \to 0 < \varepsilon \to \varphi_{55}) \to \mathrm{Seq}(a) \to \mathrm{Conv}(a \circ \mathsf{s},\alpha) \to \varepsilon \in \mathbb{R} \to 0 < \varepsilon \to \varphi_{73}$
\item $\varphi_{155} :\equiv (\mathrm{Seq}(a) \to \mathrm{Conv}(a \circ \mathsf{s},\alpha) \to \varepsilon \in \mathbb{R} \to 0 < \varepsilon \to \varphi_{56}) \to \mathrm{Seq}(a) \to \mathrm{Conv}(a \circ \mathsf{s},\alpha) \to \varepsilon \in \mathbb{R} \to 0 < \varepsilon \to \varphi_{81}$
\item $\varphi_{156} :\equiv (\mathrm{Seq}(a) \to \mathrm{Conv}(a \circ \mathsf{s},\alpha) \to \mathrm{Conv}(a \circ \mathsf{s},\alpha)) \to \varphi_{145}$
\item $\varphi_{157} :\equiv (\mathrm{Seq}(a) \to \mathrm{Conv}(a \circ \mathsf{s},\alpha) \to \varepsilon \in \mathbb{R} \to 0 < \varepsilon \to \varphi_{72}) \to \mathrm{Seq}(a) \to \mathrm{Conv}(a \circ \mathsf{s},\alpha) \to \varepsilon \in \mathbb{R} \to 0 < \varepsilon \to \varphi_{77}$
\item $\varphi_{158} :\equiv (\mathrm{Seq}(a) \to \mathrm{Conv}(a \circ \mathsf{s},\alpha) \to \varepsilon \in \mathbb{R} \to 0 < \varepsilon \to \varphi_{77}) \to \mathrm{Seq}(a) \to \mathrm{Conv}(a \circ \mathsf{s},\alpha) \to \varepsilon \in \mathbb{R} \to 0 < \varepsilon \to \varphi_{75}$
\item $\varphi_{159} :\equiv (\varphi_{28} \to \varphi_{35}) \to \mathrm{Seq}(a) \to \mathrm{Conv}(a \circ \mathsf{s},\alpha) \to \varepsilon \in \mathbb{R} \to 0 < \varepsilon \to \varphi_{103}$
\item $\varphi_{160} :\equiv (\mathrm{Seq}(a) \to \mathrm{Conv}(a \circ \mathsf{s},\alpha) \to \varepsilon \in \mathbb{R} \to 0 < \varepsilon \to \varphi_{56}) \to \mathrm{Seq}(a) \to \mathrm{Conv}(a \circ \mathsf{s},\alpha) \to \varepsilon \in \mathbb{R} \to \varphi_{105}$
\item $\varphi_{161} :\equiv \varphi_{80} \to \mathrm{Seq}(a) \to \mathrm{Conv}(a \circ \mathsf{s},\alpha) \to \varepsilon \in \mathbb{R} \to 0 < \varepsilon \to \varphi_{122}$
\item $\varphi_{162} :\equiv (\mathrm{Seq}(a) \to \mathrm{Conv}(a \circ \mathsf{s},\alpha) \to \varepsilon \in \mathbb{R} \to 0 < \varepsilon \to \varphi_{76}) \to \mathrm{Seq}(a) \to \mathrm{Conv}(a \circ \mathsf{s},\alpha) \to \varepsilon \in \mathbb{R} \to 0 < \varepsilon \to \varphi_{83}$
\item $\varphi_{163} :\equiv (\mathrm{Seq}(a) \to \mathrm{Conv}(a \circ \mathsf{s},\alpha) \to \varepsilon \in \mathbb{R} \to 0 < \varepsilon \to \varphi_{75}) \to \varphi_{128}$
\item $\varphi_{164} :\equiv (\mathrm{Seq}(a) \to \mathrm{Conv}(a \circ \mathsf{s},\alpha) \to \varepsilon \in \mathbb{R} \to 0 < \varepsilon \to \varphi_{109}) \to \mathrm{Seq}(a) \to \mathrm{Conv}(a \circ \mathsf{s},\alpha) \to \varepsilon \in \mathbb{R} \to 0 < \varepsilon \to \varphi_{55}$
\item $\varphi_{165} :\equiv \varphi_{84} \to \mathrm{Seq}(a) \to \mathrm{Conv}(a \circ \mathsf{s},\alpha) \to \varepsilon \in \mathbb{R} \to 0 < \varepsilon \to \varphi_{124}$
\item $\varphi_{166} :\equiv (\mathrm{Seq}(a) \to \mathrm{Conv}(a \circ \mathsf{s},\alpha) \to \varepsilon \in \mathbb{R} \to 0 < \varepsilon \to \varphi_{88}) \to \varphi_{97} \to n \in \mathbb{N} \to S(N) \leq n \to \lvert a_{n} - \alpha \rvert < \varepsilon$
\item $\varphi_{167} :\equiv \varphi_{97} \land n \in \mathbb{N} \land S(N) \leq n \land \neg \lvert a_{n} - \alpha \rvert < \varepsilon \to \forall m\,(\varphi_{97} \land n \in \mathbb{N} \land S(N) \leq n \land \neg \lvert a_{n} - \alpha \rvert < \varepsilon)$
\item $\varphi_{168} :\equiv (\mathrm{Seq}(a) \to \mathrm{Conv}(a \circ \mathsf{s},\alpha) \to \varepsilon \in \mathbb{R} \to \varphi_{102}) \to \mathrm{Seq}(a) \to \mathrm{Conv}(a \circ \mathsf{s},\alpha) \to \varepsilon \in \mathbb{R} \to 0 < \varepsilon \to \varphi_{90}$
\item $\varphi_{169} :\equiv (\mathrm{Seq}(a) \to \mathrm{Conv}(a \circ \mathsf{s},\alpha) \to \varepsilon \in \mathbb{R} \to \varphi_{105}) \to \mathrm{Seq}(a) \to \mathrm{Conv}(a \circ \mathsf{s},\alpha) \to \varepsilon \in \mathbb{R} \to 0 < \varepsilon \to \varphi_{87}$
\item $\varphi_{170} :\equiv (\mathrm{Seq}(a) \to \mathrm{Conv}(a \circ \mathsf{s},\alpha) \to \varepsilon \in \mathbb{R} \to 0 < \varepsilon \to \varphi_{90}) \to \mathrm{Seq}(a) \to \mathrm{Conv}(a \circ \mathsf{s},\alpha) \to \varepsilon \in \mathbb{R} \to \varphi_{106}$
\item $\varphi_{171} :\equiv \forall n\,(\varphi_{97} \to n \in \mathbb{N} \to S(N) \leq n \to \lvert a_{n} - \alpha \rvert < \varepsilon) \to \forall n\,(\varphi_{97}) \to \forall n \in \mathbb{N}\,(S(N) \leq n \to \lvert a_{n} - \alpha \rvert < \varepsilon)$
\item $\varphi_{172} :\equiv (\mathrm{Seq}(a) \to \mathrm{Conv}(a \circ \mathsf{s},\alpha) \to \varepsilon \in \mathbb{R} \to 0 < \varepsilon \to \varphi_{95}) \to \mathrm{Seq}(a) \to \mathrm{Conv}(a \circ \mathsf{s},\alpha) \to \varepsilon \in \mathbb{R} \to 0 < \varepsilon \to \varphi_{96}$
\item $\varphi_{173} :\equiv (\mathrm{Seq}(a) \to \mathrm{Conv}(a \circ \mathsf{s},\alpha) \to \varepsilon \in \mathbb{R} \to \varphi_{111}) \to \mathrm{Seq}(a) \to \mathrm{Conv}(a \circ \mathsf{s},\alpha) \to \varepsilon \in \mathbb{R} \to 0 < \varepsilon \to \varphi_{89}$
\item $\varphi_{174} :\equiv (\mathrm{Seq}(a) \to \mathrm{Conv}(a \circ \mathsf{s},\alpha) \to \varepsilon \in \mathbb{R} \to 0 < \varepsilon \to \varphi_{95}) \to \mathrm{Seq}(a) \to \mathrm{Conv}(a \circ \mathsf{s},\alpha) \to \varepsilon \in \mathbb{R} \to \varphi_{108}$
\item $\varphi_{175} :\equiv (\mathrm{Seq}(a) \to \mathrm{Conv}(a \circ \mathsf{s},\alpha) \to \varepsilon \in \mathbb{R} \to \varphi_{111}) \to \mathrm{Seq}(a) \to \mathrm{Conv}(a \circ \mathsf{s},\alpha) \to \varepsilon \in \mathbb{R} \to 0 < \varepsilon \to \varphi_{95}$
\item $\varphi_{176} :\equiv (\mathrm{Seq}(a) \to \mathrm{Conv}(a \circ \mathsf{s},\alpha) \to \varepsilon \in \mathbb{R} \to 0 < \varepsilon \to \varphi_{89}) \to \varphi_{130}$
\item $\varphi_{177} :\equiv (\mathrm{Seq}(a) \to \mathrm{Conv}(a \circ \mathsf{s},\alpha) \to \varepsilon \in \mathbb{R} \to 0 < \varepsilon \to \varphi_{94}) \to \mathrm{Seq}(a) \to \mathrm{Conv}(a \circ \mathsf{s},\alpha) \to \varepsilon \in \mathbb{R} \to 0 < \varepsilon \to \varphi_{107}$
\item $\varphi_{178} :\equiv (\mathrm{Seq}(a) \to \mathrm{Conv}(a \circ \mathsf{s},\alpha) \to \varepsilon \in \mathbb{R} \to 0 < \varepsilon \to \varphi_{100}) \to \mathrm{Seq}(a) \to \mathrm{Conv}(a \circ \mathsf{s},\alpha) \to \varepsilon \in \mathbb{R} \to 0 < \varepsilon \to \varphi_{101}$
\item $\varphi_{179} :\equiv (\mathrm{Seq}(a) \to \mathrm{Conv}(a \circ \mathsf{s},\alpha) \to \varepsilon \in \mathbb{R} \to 0 < \varepsilon \to \varphi_{87}) \to \mathrm{Seq}(a) \to \mathrm{Conv}(a \circ \mathsf{s},\alpha) \to \varepsilon \in \mathbb{R} \to 0 < \varepsilon \to \varphi_{114}$
\item $\varphi_{180} :\equiv (\mathrm{Seq}(a) \to \mathrm{Conv}(a \circ \mathsf{s},\alpha) \to \varepsilon \in \mathbb{R} \to 0 < \varepsilon \to \varphi_{103}) \to \varphi_{93} \to \varphi_{41}$
\item $\varphi_{181} :\equiv (\mathrm{Seq}(a) \to \mathrm{Conv}(a \circ \mathsf{s},\alpha) \to \varepsilon \in \mathbb{R} \to 0 < \varepsilon \to \varphi_{95}) \to \mathrm{Seq}(a) \to \mathrm{Conv}(a \circ \mathsf{s},\alpha) \to \varepsilon \in \mathbb{R} \to \varphi_{121}$
\item $\varphi_{182} :\equiv (\mathrm{Seq}(a) \to \mathrm{Conv}(a \circ \mathsf{s},\alpha) \to \varepsilon \in \mathbb{R} \to 0 < \varepsilon \to \varphi_{98}) \to \varphi_{132}$
\item $\varphi_{183} :\equiv (N \leq S(N) \to S(N) \leq S(m) \to N \leq S(m)) \to N \leq S(N) \to \varphi_{168}$
\item $\varphi_{184} :\equiv (\mathrm{Seq}(a) \to \mathrm{Conv}(a \circ \mathsf{s},\alpha) \to \varepsilon \in \mathbb{R} \to 0 < \varepsilon \to \varphi_{72}) \to (n = S(m) \to S(N) < n \to S(N) < S(m)) \to \varphi_{134}$
\item $\varphi_{185} :\equiv (\mathrm{Seq}(a) \to \mathrm{Conv}(a \circ \mathsf{s},\alpha) \to \varepsilon \in \mathbb{R} \to \varphi_{104}) \to \mathrm{Seq}(a) \to \mathrm{Conv}(a \circ \mathsf{s},\alpha) \to \varepsilon \in \mathbb{R} \to 0 < \varepsilon \to \varphi_{129}$
\item $\varphi_{186} :\equiv (\mathrm{Seq}(a) \to \mathrm{Conv}(a \circ \mathsf{s},\alpha) \to \varepsilon \in \mathbb{R} \to 0 < \varepsilon \to \varphi_{109}) \to \mathrm{Seq}(a) \to \mathrm{Conv}(a \circ \mathsf{s},\alpha) \to \varepsilon \in \mathbb{R} \to 0 < \varepsilon \to \varphi_{126}$
\item $\varphi_{187} :\equiv (\mathrm{Seq}(a) \to \mathrm{Conv}(a \circ \mathsf{s},\alpha) \to \varepsilon \in \mathbb{R} \to \varphi_{102}) \to (S(m) = N \to S(N) < S(m) \to S(N) < N) \to \varphi_{135}$
\item $\varphi_{188} :\equiv (\varphi_{97} \to \forall n\,(\varphi_{97})) \to \mathrm{Seq}(a) \to \mathrm{Conv}(a \circ \mathsf{s},\alpha) \to \varepsilon \in \mathbb{R} \to 0 < \varepsilon \to \varphi_{94}$
\item $\varphi_{189} :\equiv (\mathrm{Seq}(a) \to \mathrm{Conv}(a \circ \mathsf{s},\alpha) \to \varepsilon \in \mathbb{R} \to 0 < \varepsilon \to \varphi_{91}) \to a_{S(m)} = a_{S(m)} \to \varphi_{140}$
\item $\varphi_{190} :\equiv (\mathrm{Seq}(a) \to \mathrm{Conv}(a \circ \mathsf{s},\alpha) \to \varepsilon \in \mathbb{R} \to 0 < \varepsilon \to \varphi_{73}) \to \varphi_{155}$
\item $\varphi_{191} :\equiv (\mathrm{Seq}(a) \to \mathrm{Conv}(a \circ \mathsf{s},\alpha) \to \varepsilon \in \mathbb{R} \to 0 < \varepsilon \to \varphi_{113}) \to \varphi_{147}$
\item $\varphi_{192} :\equiv (\mathrm{Seq}(a) \to \mathrm{Conv}(a \circ \mathsf{s},\alpha) \to \varepsilon \in \mathbb{R} \to 0 < \varepsilon \to \varphi_{129}) \to \mathrm{Seq}(a) \to \mathrm{Conv}(a \circ \mathsf{s},\alpha) \to \varepsilon \in \mathbb{R} \to 0 < \varepsilon \to \varphi_{123}$
\item $\varphi_{193} :\equiv (\mathrm{Seq}(a) \to \mathrm{Conv}(a \circ \mathsf{s},\alpha) \to \varepsilon \in \mathbb{R} \to 0 < \varepsilon \to \varphi_{81}) \to \varphi_{157}$
\item $\varphi_{194} :\equiv (\mathrm{Seq}(a) \to \mathrm{Conv}(a \circ \mathsf{s},\alpha) \to \varepsilon \in \mathbb{R} \to 0 < \varepsilon \to \varphi_{70}) \to \varphi_{162}$
\item $\varphi_{195} :\equiv (\mathrm{Seq}(a) \to \mathrm{Conv}(a \circ \mathsf{s},\alpha) \to \varepsilon \in \mathbb{R} \to 0 < \varepsilon \to \varphi_{114}) \to \varphi_{71} \to \mathrm{Seq}(a) \to \mathrm{Conv}(a \circ \mathsf{s},\alpha) \to \varepsilon \in \mathbb{R} \to \varphi_{115}$
\item $\varphi_{196} :\equiv (\mathrm{Seq}(a) \to \mathrm{Conv}(a \circ \mathsf{s},\alpha) \to \varepsilon \in \mathbb{R} \to \varphi_{117}) \to \varphi_{79} \to \mathrm{Seq}(a) \to \mathrm{Conv}(a \circ \mathsf{s},\alpha) \to \varepsilon \in \mathbb{R} \to \varphi_{118}$
\item $\varphi_{197} :\equiv (\mathrm{Seq}(a) \to \mathrm{Conv}(a \circ \mathsf{s},\alpha) \to \varepsilon \in \mathbb{R} \to 0 < \varepsilon \to \varphi_{101}) \to \varphi_{160}$
\item $\varphi_{198} :\equiv (\mathrm{Seq}(a) \to \mathrm{Conv}(a \circ \mathsf{s},\alpha) \to \varepsilon \in \mathbb{R} \to \varphi_{108}) \to a_{\mathsf{s}(m)} = a_{\mathsf{s}(m)} \to \varphi_{153}$
\item $\varphi_{199} :\equiv (\mathrm{Seq}(a) \to \mathrm{Conv}(a \circ \mathsf{s},\alpha) \to \varepsilon \in \mathbb{R} \to 0 < \varepsilon \to \varphi_{61}) \to \varphi_{177}$
\item $\varphi_{200} :\equiv (\mathrm{Seq}(a) \to \mathrm{Conv}(a \circ \mathsf{s},\alpha) \to \varepsilon \in \mathbb{R} \to 0 < \varepsilon \to \varphi_{107}) \to \varphi_{159}$
\item $\varphi_{201} :\equiv \varphi_{167} \to \mathrm{Seq}(a) \to \mathrm{Conv}(a \circ \mathsf{s},\alpha) \to \varepsilon \in \mathbb{R} \to 0 < \varepsilon \to \varphi_{88}$
\item $\varphi_{202} :\equiv (\mathrm{Seq}(a) \to \mathrm{Conv}(a \circ \mathsf{s},\alpha) \to \varepsilon \in \mathbb{R} \to 0 < \varepsilon \to \varphi_{96}) \to \varphi_{170}$
\item $\varphi_{203} :\equiv (\mathrm{Seq}(a) \to \mathrm{Conv}(a \circ \mathsf{s},\alpha) \to \varepsilon \in \mathbb{R} \to 0 < \varepsilon \to \varphi_{123}) \to \varphi_{161}$
\item $\varphi_{204} :\equiv (\mathrm{Seq}(a) \to \mathrm{Conv}(a \circ \mathsf{s},\alpha) \to \varepsilon \in \mathbb{R} \to 0 < \varepsilon \to \varphi_{95}) \to \varphi_{179}$
\item $\varphi_{205} :\equiv (\mathrm{Seq}(a) \to \mathrm{Conv}(a \circ \mathsf{s},\alpha) \to \varepsilon \in \mathbb{R} \to 0 < \varepsilon \to \varphi_{110}) \to \varphi_{174}$
\item $\varphi_{206} :\equiv (\mathrm{Seq}(a) \to \mathrm{Conv}(a \circ \mathsf{s},\alpha) \to \varepsilon \in \mathbb{R} \to 0 < \varepsilon \to \varphi_{126}) \to \varphi_{165}$
\item $\varphi_{207} :\equiv \forall \varepsilon\,(\mathrm{Seq}(a) \land \mathrm{Conv}(a \circ \mathsf{s},\alpha) \to \varphi_{66}) \to \varphi_{143} \to \varphi_{38} \to \mathrm{Seq}(a) \to \mathrm{Conv}(a \circ \mathsf{s},\alpha) \to \varphi_{68}$
\item $\varphi_{208} :\equiv (\mathrm{Seq}(a) \to \mathrm{Conv}(a \circ \mathsf{s},\alpha) \to \varepsilon \in \mathbb{R} \to 0 < \varepsilon \to \varphi_{89}) \to \varphi_{184}$
\item $\varphi_{209} :\equiv (\mathrm{Seq}(a) \to \mathrm{Conv}(a \circ \mathsf{s},\alpha) \to \varepsilon \in \mathbb{R} \to 0 < \varepsilon \to \varphi_{92}) \to \varphi_{185}$
\item $\varphi_{210} :\equiv (\mathrm{Seq}(a) \to \mathrm{Conv}(a \circ \mathsf{s},\alpha) \to \varepsilon \in \mathbb{R} \to 0 < \varepsilon \to \varphi_{122}) \to \varphi_{181}$
\item $\varphi_{211} :\equiv (\mathrm{Seq}(a) \to \mathrm{Conv}(a \circ \mathsf{s},\alpha) \to \varepsilon \in \mathbb{R} \to 0 < \varepsilon \to \varphi_{99}) \to \varphi_{187}$
\item $\varphi_{212} :\equiv (\mathrm{Seq}(a) \to \mathrm{Conv}(a \circ \mathsf{s},\alpha) \to \varepsilon \in \mathbb{R} \to \varphi_{118}) \to (\mathrm{Seq}(a) \to \mathrm{Conv}(a \circ \mathsf{s},\alpha) \to \varepsilon \in \mathbb{R} \to \varphi_{112}) \to \varphi_{152}$
\item $\varphi_{213} :\equiv (\mathrm{Seq}(a) \to \mathrm{Conv}(a \circ \mathsf{s},\alpha) \to \varepsilon \in \mathbb{R} \to 0 < \varepsilon \to \varphi_{59}) \to \varphi_{202}$
\item $\varphi_{214} :\equiv (\forall N\,(\varphi_{93} \to \varphi_{41}) \to \forall N\,(\varphi_{93}) \to \forall N\,(\varphi_{41})) \to (\varphi_{93} \to \forall N\,(\varphi_{93})) \to \mathrm{Seq}(a) \to \mathrm{Conv}(a \circ \mathsf{s},\alpha) \to \varphi_{66}$
\item $\varphi_{215} :\equiv (\mathrm{Seq}(a) \to \mathrm{Conv}(a \circ \mathsf{s},\alpha) \to \varepsilon \in \mathbb{R} \to 0 < \varepsilon \to \varphi_{124}) \to \varphi_{204}$
\item $\varphi_{216} :\equiv (\mathrm{Seq}(a) \to \mathrm{Conv}(a \circ \mathsf{s},\alpha) \to \varepsilon \in \mathbb{R} \to \varphi_{60}) \to \forall N\,(\varphi_{93} \to \varphi_{41}) \to \varphi_{214}$
\item $\varphi_{217} :\equiv (\mathrm{Seq}(a) \to \mathrm{Conv}(a \circ \mathsf{s},\alpha) \to \varepsilon \in \mathbb{R} \to 0 < \varepsilon \to \varphi_{83}) \to \forall m\,(\varphi_{131}) \to (\forall m\,(\varphi_{131}) \to \varphi_{133}) \to \varphi_{201}$
\end{abbrevlist}
\end{fullproof}

% 自動生成 (prover/build.js): PSumAdd  ノード数 630
\begin{fullproof}{[Q1]}{$\forall a\,\forall b\,(\mathrm{Seq}(a) \to \mathrm{Seq}(b) \to \forall n \in \mathbb{N}\,({\textstyle\sum_{k=1}^{n} a \oplus b(k)} = {\textstyle\sum_{k=1}^{n} a_{k}} + {\textstyle\sum_{k=1}^{n} b_{k}}))$}{630}\phantomsection\label{tree:PSumAdd}
\noindent$\Pi_{1}$:
\begin{ptree}
\AXC{$\mathrm{Seq}(a) \to \mathrm{Seq}(b) \to n \in \mathbb{N} \to n \in \mathbb{N}$\lbl{A1}}
\AXC{$(\mathrm{Seq}(a) \to \mathrm{Seq}(b) \to n \in \mathbb{N} \to n \in \mathbb{N}) \to \varphi_{42}$\lbl{A1}}
\BIC{$\varphi_{42}$}
\end{ptree}
\noindent$\Pi_{2}$:
\begin{ptree}
\AXC{$\mathrm{Seq}(a) \to \mathrm{Seq}(b) \to \forall n\,(\varphi_{68}) \to \forall n\,(\varphi_{68})$\lbl{A1}}
\AXC{$(\mathrm{Seq}(a) \to \mathrm{Seq}(b) \to \forall n\,(\varphi_{68}) \to \forall n\,(\varphi_{68})) \to \varphi_{233}$\lbl{A1}}
\BIC{$\varphi_{233}$}
\end{ptree}
\noindent$\Pi_{3}$:
\begin{ptree}
\AXC{$\varphi_{16}$\lbl{A4}}
\AXC{$\forall y\,(\varphi_{7}) \to \varphi_{18}$\lbl{A3}}
\AXC{$\varphi_{7}$\lbl{A2}}
\UIC{$\forall y\,(\varphi_{7})$}
\AXC{$\forall y\,(\varphi_{7}) \to (\forall y\,(\varphi_{7}) \to \varphi_{18}) \to \varphi_{16} \to \varphi_{13}$\lbl{A1}}
\BIC{$(\forall y\,(\varphi_{7}) \to \varphi_{18}) \to \varphi_{16} \to \varphi_{13}$}
\BIC{$\varphi_{16} \to \varphi_{13}$}
\BIC{$\varphi_{13}$}
\end{ptree}
\noindent$\Pi_{4}$:
\begin{ptree}
\AXC{$\varphi_{17}$\lbl{A4}}
\AXC{$\forall n\,(\varphi_{8}) \to \varphi_{19}$\lbl{A3}}
\AXC{$\varphi_{8}$\lbl{A2}}
\UIC{$\forall n\,(\varphi_{8})$}
\AXC{$\forall n\,(\varphi_{8}) \to (\forall n\,(\varphi_{8}) \to \varphi_{19}) \to \varphi_{17} \to \varphi_{14}$\lbl{A1}}
\BIC{$(\forall n\,(\varphi_{8}) \to \varphi_{19}) \to \varphi_{17} \to \varphi_{14}$}
\BIC{$\varphi_{17} \to \varphi_{14}$}
\BIC{$\varphi_{14}$}
\end{ptree}
\noindent$\Pi_{5}$:
\begin{ptree}
\AXC{$y = z \to x = y \to x = z$\lbl{A6}}
\AXC{$(y = z \to x = y \to x = z) \to x = y \to y = z \to x = z$\lbl{A1}}
\BIC{$x = y \to y = z \to x = z$}
\UIC{$\forall z\,(x = y \to y = z \to x = z)$}
\UIC{$\forall y\,\forall z\,(x = y \to y = z \to x = z)$}
\UIC{$\forall x\,\forall y\,\forall z\,(x = y \to y = z \to x = z)$}
\AXC{$\forall x\,\forall y\,\forall z\,(x = y \to y = z \to x = z) \to \varphi_{29}$\lbl{A2}}
\BIC{$\varphi_{29}$}
\end{ptree}
\noindent$\Pi_{6}$:
\begin{ptree}
\AXC{$\stackrel{\Pi_{5}}{\vdots}$}\noLine
\UIC{$\varphi_{29}$}
\AXC{$\varphi_{29} \to \varphi_{166}$\lbl{A2}}
\BIC{$\varphi_{166}$}
\AXC{$\varphi_{166} \to {\textstyle\sum_{k=1}^{S(n)} a \oplus b(k)} = {\textstyle\sum_{k=1}^{S(n)} a_{k}} + ({\textstyle\sum_{k=1}^{n} b_{k}} + b_{S(n)}) \to \varphi_{154}$\lbl{A2}}
\BIC{${\textstyle\sum_{k=1}^{S(n)} a \oplus b(k)} = {\textstyle\sum_{k=1}^{S(n)} a_{k}} + ({\textstyle\sum_{k=1}^{n} b_{k}} + b_{S(n)}) \to \varphi_{154}$}
\end{ptree}
\noindent$\Pi_{7}$:
\begin{ptree}
\AXC{$x = x$\lbl{A5}}
\AXC{$x = y \to x = x \to y = x$\lbl{A6}}
\AXC{$(x = y \to x = x \to y = x) \to x = x \to x = y \to y = x$\lbl{A1}}
\BIC{$x = x \to x = y \to y = x$}
\BIC{$x = y \to y = x$}
\UIC{$\forall y\,(x = y \to y = x)$}
\UIC{$\forall x\,\forall y\,(x = y \to y = x)$}
\end{ptree}
\noindent$\Pi_{8}$:
\begin{ptree}
\AXC{$\mathrm{Seq}(a) \to \mathrm{Seq}(b) \to \forall n\,(\varphi_{68}) \to \forall n\,(\varphi_{68})$\lbl{A1}}
\AXC{$(\mathrm{Seq}(a) \to \mathrm{Seq}(b) \to \forall n\,(\varphi_{68}) \to \forall n\,(\varphi_{68})) \to \varphi_{233}$\lbl{A1}}
\BIC{$\varphi_{233}$}
\end{ptree}
\noindent$\Pi_{9}$:
\begin{ptree}
\AXC{$\varphi_{126}$\lbl{A1}}
\AXC{$\forall n\,(\varphi_{68}) \to \varphi_{68}$\lbl{A2}}
\AXC{$\stackrel{\Pi_{8}}{\vdots}$}\noLine
\UIC{$\varphi_{233}$}
\BIC{$\mathrm{Seq}(a) \to \mathrm{Seq}(b) \to \forall n\,(\varphi_{68}) \to \varphi_{68}$}
\AXC{$\varphi_{261}$\lbl{A1}}
\BIC{$\varphi_{126} \to \mathrm{Seq}(a) \to \mathrm{Seq}(b) \to \forall n\,(\varphi_{68}) \to \varphi_{69}$}
\BIC{$\mathrm{Seq}(a) \to \mathrm{Seq}(b) \to \forall n\,(\varphi_{68}) \to \varphi_{69}$}
\end{ptree}
\noindent$\Pi_{10}$:
\begin{ptree}
\AXC{$\stackrel{\Pi_{9}}{\vdots}$}\noLine
\UIC{$\mathrm{Seq}(a) \to \mathrm{Seq}(b) \to \forall n\,(\varphi_{68}) \to \varphi_{69}$}
\AXC{$\varphi_{225}$\lbl{A1}}
\BIC{$\mathrm{Seq}(a) \to \mathrm{Seq}(b) \to \forall n\,(\varphi_{68}) \to n \in B_{43} \to n \in \omega$}
\end{ptree}
\noindent$\Pi_{11}$:
\begin{ptree}
\AXC{$\forall n\,(n \in \mathbb{N} \leftrightarrow n \in \omega \land n \neq \emptyset)$\lbl{Def$\mathbb{N}$}}
\AXC{$\varphi_{39}$\lbl{A2}}
\BIC{$n \in \mathbb{N} \leftrightarrow n \in \omega \land n \neq \emptyset$}
\AXC{$(n \in \mathbb{N} \leftrightarrow n \in \omega \land n \neq \emptyset) \to \varphi_{245}$\lbl{A1}}
\BIC{$\varphi_{245}$}
\end{ptree}
\noindent$\Pi_{12}$:
\begin{ptree}
\AXC{$\mathrm{Seq}(a) \to \varphi_{159}$\lbl{A1}}
\AXC{$\stackrel{\Pi_{10}}{\vdots}$}\noLine
\UIC{$\mathrm{Seq}(a) \to \mathrm{Seq}(b) \to \forall n\,(\varphi_{68}) \to n \in B_{43} \to n \in \omega$}
\AXC{$\stackrel{\Pi_{11}}{\vdots}$}\noLine
\UIC{$\varphi_{245}$}
\BIC{$(\mathrm{Seq}(a) \to \varphi_{159}) \to \mathrm{Seq}(a) \to \varphi_{156}$}
\BIC{$\mathrm{Seq}(a) \to \varphi_{156}$}
\end{ptree}
\noindent$\Pi_{13}$:
\begin{ptree}
\AXC{$\forall a\,(\mathrm{Seq}(a) \to \varphi_{98})$\lbl{[S25]}}
\AXC{$\forall a\,(\mathrm{Seq}(a) \to \varphi_{98}) \to \mathrm{Seq}(b) \to \varphi_{97}$\lbl{A2}}
\BIC{$\mathrm{Seq}(b) \to \varphi_{97}$}
\end{ptree}
\noindent$\Pi_{14}$:
\begin{ptree}
\AXC{$\mathrm{Seq}(a) \to \varphi_{150}$\lbl{A1}}
\AXC{$\stackrel{\Pi_{13}}{\vdots}$}\noLine
\UIC{$\mathrm{Seq}(b) \to \varphi_{97}$}
\AXC{$(\mathrm{Seq}(b) \to \varphi_{97}) \to (\mathrm{Seq}(a) \to \varphi_{150}) \to \mathrm{Seq}(a) \to \varphi_{207}$\lbl{A1}}
\BIC{$(\mathrm{Seq}(a) \to \varphi_{150}) \to \mathrm{Seq}(a) \to \varphi_{207}$}
\BIC{$\mathrm{Seq}(a) \to \varphi_{207}$}
\AXC{$\varphi_{246}$\lbl{A1}}
\BIC{$\mathrm{Seq}(a) \to \mathrm{Seq}(b) \to \forall n\,(\varphi_{68}) \to n \in B_{43} \to \varphi_{63}$}
\end{ptree}
\noindent$\Pi_{15}$:
\begin{ptree}
\AXC{$\stackrel{\Pi_{12}}{\vdots}$}\noLine
\UIC{$\mathrm{Seq}(a) \to \varphi_{156}$}
\AXC{$\varphi_{116}$\lbl{A2}}
\AXC{$\stackrel{\Pi_{14}}{\vdots}$}\noLine
\UIC{$\mathrm{Seq}(a) \to \mathrm{Seq}(b) \to \forall n\,(\varphi_{68}) \to n \in B_{43} \to \varphi_{63}$}
\AXC{$\varphi_{262}$\lbl{A1}}
\BIC{$\varphi_{116} \to \mathrm{Seq}(a) \to \mathrm{Seq}(b) \to \forall n\,(\varphi_{68}) \to n \in B_{43} \to \varphi_{53}$}
\BIC{$\mathrm{Seq}(a) \to \mathrm{Seq}(b) \to \forall n\,(\varphi_{68}) \to n \in B_{43} \to \varphi_{53}$}
\AXC{$\varphi_{264}$\lbl{A1}}
\BIC{$\varphi_{228}$}
\BIC{$\mathrm{Seq}(a) \to \mathrm{Seq}(b) \to \forall n\,(\varphi_{68}) \to \varphi_{54}$}
\end{ptree}
\noindent$\Pi_{16}$:
\begin{ptree}
\AXC{$\stackrel{\Pi_{15}}{\vdots}$}\noLine
\UIC{$\mathrm{Seq}(a) \to \mathrm{Seq}(b) \to \forall n\,(\varphi_{68}) \to \varphi_{54}$}
\AXC{$(\mathrm{Seq}(a) \to \mathrm{Seq}(b) \to \forall n\,(\varphi_{68}) \to \varphi_{54}) \to \varphi_{215}$\lbl{A1}}
\BIC{$\varphi_{215}$}
\end{ptree}
\noindent$\Pi_{17}$:
\begin{ptree}
\AXC{$\stackrel{\Pi_{7}}{\vdots}$}\noLine
\UIC{$\forall x\,\forall y\,(x = y \to y = x)$}
\AXC{$\varphi_{20}$\lbl{A2}}
\BIC{$\forall y\,({\textstyle\sum_{k=1}^{S(n)} b_{k}} = y \to y = {\textstyle\sum_{k=1}^{S(n)} b_{k}})$}
\AXC{$\varphi_{134}$\lbl{A2}}
\BIC{${\textstyle\sum_{k=1}^{S(n)} b_{k}} = {\textstyle\sum_{k=1}^{n} b_{k}} + b_{S(n)} \to {\textstyle\sum_{k=1}^{n} b_{k}} + b_{S(n)} = {\textstyle\sum_{k=1}^{S(n)} b_{k}}$}
\AXC{$\stackrel{\Pi_{16}}{\vdots}$}\noLine
\UIC{$\varphi_{215}$}
\BIC{$\mathrm{Seq}(a) \to \mathrm{Seq}(b) \to \forall n\,(\varphi_{68}) \to \varphi_{55}$}
\AXC{$(\mathrm{Seq}(a) \to \mathrm{Seq}(b) \to \forall n\,(\varphi_{68}) \to \varphi_{55}) \to \varphi_{265}$\lbl{A1}}
\BIC{$\varphi_{265}$}
\end{ptree}
\noindent$\Pi_{18}$:
\begin{ptree}
\AXC{$z = w \to \varphi_{140}$\lbl{A6}}
\UIC{$\forall w\,(z = w \to \varphi_{140})$}
\UIC{$\forall z\,\forall w\,(z = w \to \varphi_{140})$}
\AXC{$\forall z\,\forall w\,(z = w \to \varphi_{140}) \to \varphi_{187}$\lbl{A2}}
\BIC{$\varphi_{187}$}
\AXC{$\varphi_{187} \to \varphi_{201}$\lbl{A2}}
\BIC{$\varphi_{201}$}
\AXC{$x = x$\lbl{A5}}
\UIC{$\forall x\,(x = x)$}
\AXC{$\forall x\,(x = x) \to \varphi_{81}$\lbl{A2}}
\BIC{$\varphi_{81}$}
\AXC{$\stackrel{\Pi_{17}}{\vdots}$}\noLine
\UIC{$\varphi_{265}$}
\BIC{$\varphi_{201} \to \mathrm{Seq}(a) \to \mathrm{Seq}(b) \to \forall n\,(\varphi_{68}) \to n \in B_{43} \to \varphi_{111}$}
\BIC{$\mathrm{Seq}(a) \to \mathrm{Seq}(b) \to \forall n\,(\varphi_{68}) \to n \in B_{43} \to \varphi_{111}$}
\end{ptree}
\noindent$\Pi_{19}$:
\begin{ptree}
\AXC{$y = z \to x = y \to x = z$\lbl{A6}}
\AXC{$(y = z \to x = y \to x = z) \to x = y \to y = z \to x = z$\lbl{A1}}
\BIC{$x = y \to y = z \to x = z$}
\UIC{$\forall z\,(x = y \to y = z \to x = z)$}
\UIC{$\forall y\,\forall z\,(x = y \to y = z \to x = z)$}
\UIC{$\forall x\,\forall y\,\forall z\,(x = y \to y = z \to x = z)$}
\AXC{$\forall x\,\forall y\,\forall z\,(x = y \to y = z \to x = z) \to \varphi_{29}$\lbl{A2}}
\BIC{$\varphi_{29}$}
\end{ptree}
\noindent$\Pi_{20}$:
\begin{ptree}
\AXC{$\forall a\,(\mathrm{Seq}(a) \to \varphi_{98})$\lbl{[S25]}}
\AXC{$\forall a\,(\mathrm{Seq}(a) \to \varphi_{98}) \to \mathrm{Seq}(a) \to \varphi_{98}$\lbl{A2}}
\BIC{$\mathrm{Seq}(a) \to \varphi_{98}$}
\end{ptree}
\noindent$\Pi_{21}$:
\begin{ptree}
\AXC{$\mathrm{Seq}(a) \to \varphi_{151}$\lbl{A1}}
\AXC{$\stackrel{\Pi_{20}}{\vdots}$}\noLine
\UIC{$\mathrm{Seq}(a) \to \varphi_{98}$}
\AXC{$(\mathrm{Seq}(a) \to \varphi_{98}) \to (\mathrm{Seq}(a) \to \varphi_{151}) \to \mathrm{Seq}(a) \to \varphi_{208}$\lbl{A1}}
\BIC{$(\mathrm{Seq}(a) \to \varphi_{151}) \to \mathrm{Seq}(a) \to \varphi_{208}$}
\BIC{$\mathrm{Seq}(a) \to \varphi_{208}$}
\AXC{$\varphi_{247}$\lbl{A1}}
\BIC{$\mathrm{Seq}(a) \to \mathrm{Seq}(b) \to \forall n\,(\varphi_{68}) \to n \in B_{43} \to \varphi_{64}$}
\end{ptree}
\noindent$\Pi_{22}$:
\begin{ptree}
\AXC{$\stackrel{\Pi_{12}}{\vdots}$}\noLine
\UIC{$\mathrm{Seq}(a) \to \varphi_{156}$}
\AXC{$\varphi_{117}$\lbl{A2}}
\AXC{$\stackrel{\Pi_{21}}{\vdots}$}\noLine
\UIC{$\mathrm{Seq}(a) \to \mathrm{Seq}(b) \to \forall n\,(\varphi_{68}) \to n \in B_{43} \to \varphi_{64}$}
\AXC{$\varphi_{263}$\lbl{A1}}
\BIC{$\varphi_{117} \to \mathrm{Seq}(a) \to \mathrm{Seq}(b) \to \forall n\,(\varphi_{68}) \to n \in B_{43} \to \varphi_{56}$}
\BIC{$\mathrm{Seq}(a) \to \mathrm{Seq}(b) \to \forall n\,(\varphi_{68}) \to n \in B_{43} \to \varphi_{56}$}
\AXC{$\varphi_{266}$\lbl{A1}}
\BIC{$\varphi_{229}$}
\BIC{$\mathrm{Seq}(a) \to \mathrm{Seq}(b) \to \forall n\,(\varphi_{68}) \to \varphi_{57}$}
\end{ptree}
\noindent$\Pi_{23}$:
\begin{ptree}
\AXC{$\stackrel{\Pi_{22}}{\vdots}$}\noLine
\UIC{$\mathrm{Seq}(a) \to \mathrm{Seq}(b) \to \forall n\,(\varphi_{68}) \to \varphi_{57}$}
\AXC{$(\mathrm{Seq}(a) \to \mathrm{Seq}(b) \to \forall n\,(\varphi_{68}) \to \varphi_{57}) \to \varphi_{216}$\lbl{A1}}
\BIC{$\varphi_{216}$}
\end{ptree}
\noindent$\Pi_{24}$:
\begin{ptree}
\AXC{$\stackrel{\Pi_{7}}{\vdots}$}\noLine
\UIC{$\forall x\,\forall y\,(x = y \to y = x)$}
\AXC{$\varphi_{21}$\lbl{A2}}
\BIC{$\forall y\,({\textstyle\sum_{k=1}^{S(n)} a_{k}} = y \to y = {\textstyle\sum_{k=1}^{S(n)} a_{k}})$}
\AXC{$\varphi_{135}$\lbl{A2}}
\BIC{${\textstyle\sum_{k=1}^{S(n)} a_{k}} = {\textstyle\sum_{k=1}^{n} a_{k}} + a_{S(n)} \to {\textstyle\sum_{k=1}^{n} a_{k}} + a_{S(n)} = {\textstyle\sum_{k=1}^{S(n)} a_{k}}$}
\AXC{$\stackrel{\Pi_{23}}{\vdots}$}\noLine
\UIC{$\varphi_{216}$}
\BIC{$\mathrm{Seq}(a) \to \mathrm{Seq}(b) \to \forall n\,(\varphi_{68}) \to \varphi_{58}$}
\AXC{$(\mathrm{Seq}(a) \to \mathrm{Seq}(b) \to \forall n\,(\varphi_{68}) \to \varphi_{58}) \to \varphi_{271}$\lbl{A1}}
\BIC{$\varphi_{271}$}
\end{ptree}
\noindent$\Pi_{25}$:
\begin{ptree}
\AXC{$z = w \to \varphi_{76} \to \varphi_{75}$\lbl{A6}}
\UIC{$\forall w\,(z = w \to \varphi_{76} \to \varphi_{75})$}
\UIC{$\forall z\,\forall w\,(z = w \to \varphi_{76} \to \varphi_{75})$}
\AXC{$\varphi_{243}$\lbl{A2}}
\BIC{$\forall w\,({\textstyle\sum_{k=1}^{n} a_{k}} + a_{S(n)} = w \to \varphi_{118} \to \varphi_{75})$}
\AXC{$\varphi_{256}$\lbl{A2}}
\BIC{${\textstyle\sum_{k=1}^{n} a_{k}} + a_{S(n)} = {\textstyle\sum_{k=1}^{S(n)} a_{k}} \to \varphi_{118} \to \varphi_{103}$}
\AXC{$x = x$\lbl{A5}}
\UIC{$\forall x\,(x = x)$}
\AXC{$\forall x\,(x = x) \to \varphi_{118}$\lbl{A2}}
\BIC{$\varphi_{118}$}
\AXC{$\stackrel{\Pi_{24}}{\vdots}$}\noLine
\UIC{$\varphi_{271}$}
\BIC{$({\textstyle\sum_{k=1}^{n} a_{k}} + a_{S(n)} = {\textstyle\sum_{k=1}^{S(n)} a_{k}} \to \varphi_{118} \to \varphi_{103}) \to \mathrm{Seq}(a) \to \varphi_{209}$}
\BIC{$\mathrm{Seq}(a) \to \varphi_{209}$}
\end{ptree}
\noindent$\Pi_{26}$:
\begin{ptree}
\AXC{$y = z \to x = y \to x = z$\lbl{A6}}
\AXC{$(y = z \to x = y \to x = z) \to x = y \to y = z \to x = z$\lbl{A1}}
\BIC{$x = y \to y = z \to x = z$}
\UIC{$\forall z\,(x = y \to y = z \to x = z)$}
\UIC{$\forall y\,\forall z\,(x = y \to y = z \to x = z)$}
\UIC{$\forall x\,\forall y\,\forall z\,(x = y \to y = z \to x = z)$}
\AXC{$\forall x\,\forall y\,\forall z\,(x = y \to y = z \to x = z) \to \varphi_{29}$\lbl{A2}}
\BIC{$\varphi_{29}$}
\end{ptree}
\noindent$\Pi_{27}$:
\begin{ptree}
\AXC{$\stackrel{\Pi_{12}}{\vdots}$}\noLine
\UIC{$\mathrm{Seq}(a) \to \varphi_{156}$}
\AXC{$\forall n \in \mathbb{N}\,(S(n) \in \mathbb{N})$\lbl{[N2]}}
\AXC{$\forall n \in \mathbb{N}\,(S(n) \in \mathbb{N}) \to n \in \mathbb{N} \to S(n) \in \mathbb{N}$\lbl{A2}}
\BIC{$n \in \mathbb{N} \to S(n) \in \mathbb{N}$}
\AXC{$\varphi_{226}$\lbl{A1}}
\BIC{$(\mathrm{Seq}(a) \to \varphi_{156}) \to \mathrm{Seq}(a) \to \mathrm{Seq}(b) \to \varphi_{145}$}
\BIC{$\mathrm{Seq}(a) \to \mathrm{Seq}(b) \to \varphi_{145}$}
\end{ptree}
\noindent$\Pi_{28}$:
\begin{ptree}
\AXC{$\mathrm{Seq}(a) \to \varphi_{150}$\lbl{A1}}
\AXC{$\forall a\,\forall n\,(\mathrm{Seq}(a) \to n \in \mathbb{N} \to a_{n} \in \mathbb{R})$\lbl{[S1]}}
\AXC{$\varphi_{49}$\lbl{A2}}
\BIC{$\forall n\,(\mathrm{Seq}(b) \to n \in \mathbb{N} \to b_{n} \in \mathbb{R})$}
\AXC{$\varphi_{50}$\lbl{A2}}
\BIC{$\mathrm{Seq}(b) \to S(n) \in \mathbb{N} \to b_{S(n)} \in \mathbb{R}$}
\AXC{$\varphi_{259}$\lbl{A1}}
\BIC{$(\mathrm{Seq}(a) \to \varphi_{150}) \to \varphi_{223}$}
\BIC{$\varphi_{223}$}
\end{ptree}
\noindent$\Pi_{29}$:
\begin{ptree}
\AXC{$\mathrm{Seq}(a) \to \varphi_{151}$\lbl{A1}}
\AXC{$\forall a\,\forall n\,(\mathrm{Seq}(a) \to n \in \mathbb{N} \to a_{n} \in \mathbb{R})$\lbl{[S1]}}
\AXC{$\varphi_{51}$\lbl{A2}}
\BIC{$\forall n\,(\mathrm{Seq}(a) \to n \in \mathbb{N} \to a_{n} \in \mathbb{R})$}
\AXC{$\varphi_{52}$\lbl{A2}}
\BIC{$\mathrm{Seq}(a) \to S(n) \in \mathbb{N} \to a_{S(n)} \in \mathbb{R}$}
\AXC{$\varphi_{260}$\lbl{A1}}
\BIC{$(\mathrm{Seq}(a) \to \varphi_{151}) \to \varphi_{224}$}
\BIC{$\varphi_{224}$}
\end{ptree}
\noindent$\Pi_{30}$:
\begin{ptree}
\AXC{$\forall a\,(\mathrm{Seq}(a) \to \varphi_{98})$\lbl{[S25]}}
\AXC{$\forall a\,(\mathrm{Seq}(a) \to \varphi_{98}) \to \mathrm{Seq}(b) \to \varphi_{97}$\lbl{A2}}
\BIC{$\mathrm{Seq}(b) \to \varphi_{97}$}
\end{ptree}
\noindent$\Pi_{31}$:
\begin{ptree}
\AXC{$\mathrm{Seq}(a) \to \varphi_{150}$\lbl{A1}}
\AXC{$\stackrel{\Pi_{30}}{\vdots}$}\noLine
\UIC{$\mathrm{Seq}(b) \to \varphi_{97}$}
\AXC{$(\mathrm{Seq}(b) \to \varphi_{97}) \to (\mathrm{Seq}(a) \to \varphi_{150}) \to \mathrm{Seq}(a) \to \varphi_{207}$\lbl{A1}}
\BIC{$(\mathrm{Seq}(a) \to \varphi_{150}) \to \mathrm{Seq}(a) \to \varphi_{207}$}
\BIC{$\mathrm{Seq}(a) \to \varphi_{207}$}
\end{ptree}
\noindent$\Pi_{32}$:
\begin{ptree}
\AXC{$\forall a\,\forall n\,(\mathrm{Seq}(a) \to n \in \mathbb{N} \to a_{n} \in \mathbb{R})$\lbl{[S1]}}
\AXC{$\varphi_{79}$\lbl{A2}}
\BIC{$\forall n\,(\mathrm{Seq}(\Sigma b) \to n \in \mathbb{N} \to {\textstyle\sum_{k=1}^{n} b_{k}} \in \mathbb{R})$}
\AXC{$\varphi_{90}$\lbl{A2}}
\BIC{$\mathrm{Seq}(\Sigma b) \to n \in \mathbb{N} \to {\textstyle\sum_{k=1}^{n} b_{k}} \in \mathbb{R}$}
\AXC{$(\mathrm{Seq}(\Sigma b) \to n \in \mathbb{N} \to {\textstyle\sum_{k=1}^{n} b_{k}} \in \mathbb{R}) \to \varphi_{250}$\lbl{A1}}
\BIC{$\varphi_{250}$}
\end{ptree}
\noindent$\Pi_{33}$:
\begin{ptree}
\AXC{$\stackrel{\Pi_{12}}{\vdots}$}\noLine
\UIC{$\mathrm{Seq}(a) \to \varphi_{156}$}
\AXC{$\stackrel{\Pi_{31}}{\vdots}$}\noLine
\UIC{$\mathrm{Seq}(a) \to \varphi_{207}$}
\AXC{$(\mathrm{Seq}(a) \to \varphi_{207}) \to \mathrm{Seq}(a) \to \mathrm{Seq}(b) \to \varphi_{146}$\lbl{A1}}
\BIC{$\mathrm{Seq}(a) \to \mathrm{Seq}(b) \to \varphi_{146}$}
\AXC{$\stackrel{\Pi_{32}}{\vdots}$}\noLine
\UIC{$\varphi_{250}$}
\BIC{$(\mathrm{Seq}(a) \to \varphi_{156}) \to \mathrm{Seq}(a) \to \mathrm{Seq}(b) \to \varphi_{164}$}
\BIC{$\mathrm{Seq}(a) \to \mathrm{Seq}(b) \to \varphi_{164}$}
\end{ptree}
\noindent$\Pi_{34}$:
\begin{ptree}
\AXC{$\forall a\,(\mathrm{Seq}(a) \to \varphi_{98})$\lbl{[S25]}}
\AXC{$\forall a\,(\mathrm{Seq}(a) \to \varphi_{98}) \to \mathrm{Seq}(a) \to \varphi_{98}$\lbl{A2}}
\BIC{$\mathrm{Seq}(a) \to \varphi_{98}$}
\end{ptree}
\noindent$\Pi_{35}$:
\begin{ptree}
\AXC{$\mathrm{Seq}(a) \to \varphi_{151}$\lbl{A1}}
\AXC{$\stackrel{\Pi_{34}}{\vdots}$}\noLine
\UIC{$\mathrm{Seq}(a) \to \varphi_{98}$}
\AXC{$(\mathrm{Seq}(a) \to \varphi_{98}) \to (\mathrm{Seq}(a) \to \varphi_{151}) \to \mathrm{Seq}(a) \to \varphi_{208}$\lbl{A1}}
\BIC{$(\mathrm{Seq}(a) \to \varphi_{151}) \to \mathrm{Seq}(a) \to \varphi_{208}$}
\BIC{$\mathrm{Seq}(a) \to \varphi_{208}$}
\end{ptree}
\noindent$\Pi_{36}$:
\begin{ptree}
\AXC{$\forall a\,\forall n\,(\mathrm{Seq}(a) \to n \in \mathbb{N} \to a_{n} \in \mathbb{R})$\lbl{[S1]}}
\AXC{$\varphi_{80}$\lbl{A2}}
\BIC{$\forall n\,(\mathrm{Seq}(\Sigma a) \to n \in \mathbb{N} \to {\textstyle\sum_{k=1}^{n} a_{k}} \in \mathbb{R})$}
\AXC{$\varphi_{91}$\lbl{A2}}
\BIC{$\mathrm{Seq}(\Sigma a) \to n \in \mathbb{N} \to {\textstyle\sum_{k=1}^{n} a_{k}} \in \mathbb{R}$}
\AXC{$(\mathrm{Seq}(\Sigma a) \to n \in \mathbb{N} \to {\textstyle\sum_{k=1}^{n} a_{k}} \in \mathbb{R}) \to \varphi_{251}$\lbl{A1}}
\BIC{$\varphi_{251}$}
\end{ptree}
\noindent$\Pi_{37}$:
\begin{ptree}
\AXC{$\stackrel{\Pi_{12}}{\vdots}$}\noLine
\UIC{$\mathrm{Seq}(a) \to \varphi_{156}$}
\AXC{$\stackrel{\Pi_{35}}{\vdots}$}\noLine
\UIC{$\mathrm{Seq}(a) \to \varphi_{208}$}
\AXC{$(\mathrm{Seq}(a) \to \varphi_{208}) \to \mathrm{Seq}(a) \to \mathrm{Seq}(b) \to \varphi_{147}$\lbl{A1}}
\BIC{$\mathrm{Seq}(a) \to \mathrm{Seq}(b) \to \varphi_{147}$}
\AXC{$\stackrel{\Pi_{36}}{\vdots}$}\noLine
\UIC{$\varphi_{251}$}
\BIC{$(\mathrm{Seq}(a) \to \varphi_{156}) \to \mathrm{Seq}(a) \to \mathrm{Seq}(b) \to \varphi_{165}$}
\BIC{$\mathrm{Seq}(a) \to \mathrm{Seq}(b) \to \varphi_{165}$}
\end{ptree}
\noindent$\Pi_{38}$:
\begin{ptree}
\AXC{$\varphi_{74}$\lbl{[E49]}}
\AXC{$\varphi_{74} \to \varphi_{122}$\lbl{A2}}
\BIC{$\varphi_{122}$}
\AXC{$\varphi_{122} \to \varphi_{167}$\lbl{A2}}
\BIC{$\varphi_{167}$}
\AXC{$\varphi_{167} \to \varphi_{180}$\lbl{A2}}
\BIC{$\varphi_{180}$}
\AXC{$\varphi_{180} \to \varphi_{192}$\lbl{A2}}
\BIC{$\varphi_{192}$}
\AXC{$\varphi_{192} \to \varphi_{278}$\lbl{A1}}
\BIC{$\varphi_{278}$}
\end{ptree}
\noindent$\Pi_{39}$:
\begin{ptree}
\AXC{$\stackrel{\Pi_{27}}{\vdots}$}\noLine
\UIC{$\mathrm{Seq}(a) \to \mathrm{Seq}(b) \to \varphi_{145}$}
\AXC{$\stackrel{\Pi_{29}}{\vdots}$}\noLine
\UIC{$\varphi_{224}$}
\BIC{$\mathrm{Seq}(a) \to \mathrm{Seq}(b) \to \varphi_{161}$}
\AXC{$\stackrel{\Pi_{33}}{\vdots}$}\noLine
\UIC{$\mathrm{Seq}(a) \to \mathrm{Seq}(b) \to \varphi_{164}$}
\AXC{$\stackrel{\Pi_{37}}{\vdots}$}\noLine
\UIC{$\mathrm{Seq}(a) \to \mathrm{Seq}(b) \to \varphi_{165}$}
\AXC{$\stackrel{\Pi_{38}}{\vdots}$}\noLine
\UIC{$\varphi_{278}$}
\BIC{$(\mathrm{Seq}(a) \to \mathrm{Seq}(b) \to \varphi_{164}) \to \varphi_{267}$}
\BIC{$\varphi_{267}$}
\BIC{$(\mathrm{Seq}(a) \to \mathrm{Seq}(b) \to \varphi_{160}) \to \mathrm{Seq}(a) \to \varphi_{212}$}
\end{ptree}
\noindent$\Pi_{40}$:
\begin{ptree}
\AXC{$\stackrel{\Pi_{27}}{\vdots}$}\noLine
\UIC{$\mathrm{Seq}(a) \to \mathrm{Seq}(b) \to \varphi_{145}$}
\AXC{$\stackrel{\Pi_{28}}{\vdots}$}\noLine
\UIC{$\varphi_{223}$}
\BIC{$\mathrm{Seq}(a) \to \mathrm{Seq}(b) \to \varphi_{160}$}
\AXC{$\stackrel{\Pi_{39}}{\vdots}$}\noLine
\UIC{$(\mathrm{Seq}(a) \to \mathrm{Seq}(b) \to \varphi_{160}) \to \mathrm{Seq}(a) \to \varphi_{212}$}
\BIC{$\mathrm{Seq}(a) \to \varphi_{212}$}
\end{ptree}
\noindent$\Pi_{41}$:
\begin{ptree}
\AXC{$y = z \to x = y \to x = z$\lbl{A6}}
\AXC{$(y = z \to x = y \to x = z) \to x = y \to y = z \to x = z$\lbl{A1}}
\BIC{$x = y \to y = z \to x = z$}
\UIC{$\forall z\,(x = y \to y = z \to x = z)$}
\UIC{$\forall y\,\forall z\,(x = y \to y = z \to x = z)$}
\UIC{$\forall x\,\forall y\,\forall z\,(x = y \to y = z \to x = z)$}
\AXC{$\forall x\,\forall y\,\forall z\,(x = y \to y = z \to x = z) \to \varphi_{29}$\lbl{A2}}
\BIC{$\varphi_{29}$}
\end{ptree}
\noindent$\Pi_{42}$:
\begin{ptree}
\AXC{$\forall a\,\forall b\,(\mathrm{Seq}(a) \to \varphi_{40})$\lbl{[S2]}}
\AXC{$\forall a\,\forall b\,(\mathrm{Seq}(a) \to \varphi_{40}) \to \forall b\,(\mathrm{Seq}(a) \to \varphi_{40})$\lbl{A2}}
\BIC{$\forall b\,(\mathrm{Seq}(a) \to \varphi_{40})$}
\end{ptree}
\noindent$\Pi_{43}$:
\begin{ptree}
\AXC{$\mathrm{Seq}(a) \to \varphi_{151}$\lbl{A1}}
\AXC{$\stackrel{\Pi_{42}}{\vdots}$}\noLine
\UIC{$\forall b\,(\mathrm{Seq}(a) \to \varphi_{40})$}
\AXC{$\forall b\,(\mathrm{Seq}(a) \to \varphi_{40}) \to \mathrm{Seq}(a) \to \varphi_{40}$\lbl{A2}}
\BIC{$\mathrm{Seq}(a) \to \varphi_{40}$}
\AXC{$\varphi_{268}$\lbl{A1}}
\BIC{$(\mathrm{Seq}(a) \to \varphi_{151}) \to (\mathrm{Seq}(a) \to \varphi_{150}) \to \varphi_{195}$}
\BIC{$(\mathrm{Seq}(a) \to \varphi_{150}) \to \varphi_{195}$}
\end{ptree}
\noindent$\Pi_{44}$:
\begin{ptree}
\AXC{$\mathrm{Seq}(a) \to \varphi_{150}$\lbl{A1}}
\AXC{$\stackrel{\Pi_{43}}{\vdots}$}\noLine
\UIC{$(\mathrm{Seq}(a) \to \varphi_{150}) \to \varphi_{195}$}
\BIC{$\varphi_{195}$}
\AXC{$\varphi_{195} \to \mathrm{Seq}(a) \to \mathrm{Seq}(b) \to \forall n\,(\varphi_{68}) \to n \in B_{43} \to \varphi_{43}$\lbl{A1}}
\BIC{$\mathrm{Seq}(a) \to \mathrm{Seq}(b) \to \forall n\,(\varphi_{68}) \to n \in B_{43} \to \varphi_{43}$}
\AXC{$\varphi_{257}$\lbl{A1}}
\BIC{$\varphi_{108} \to \varphi_{191}$}
\end{ptree}
\noindent$\Pi_{45}$:
\begin{ptree}
\AXC{$\stackrel{\Pi_{27}}{\vdots}$}\noLine
\UIC{$\mathrm{Seq}(a) \to \mathrm{Seq}(b) \to \varphi_{145}$}
\AXC{$\varphi_{108}$\lbl{A2}}
\AXC{$\stackrel{\Pi_{44}}{\vdots}$}\noLine
\UIC{$\varphi_{108} \to \varphi_{191}$}
\BIC{$\varphi_{191}$}
\AXC{$\varphi_{191} \to \varphi_{230}$\lbl{A1}}
\BIC{$\varphi_{230}$}
\BIC{$\mathrm{Seq}(a) \to \mathrm{Seq}(b) \to \forall n\,(\varphi_{68}) \to n \in B_{43} \to \varphi_{37}$}
\AXC{$\varphi_{275}$\lbl{A1}}
\BIC{$\varphi_{102} \to \varphi_{258}$}
\end{ptree}
\noindent$\Pi_{46}$:
\begin{ptree}
\AXC{$z = w \to \varphi_{71} \to \varphi_{70}$\lbl{A6}}
\UIC{$\forall w\,(z = w \to \varphi_{71} \to \varphi_{70})$}
\UIC{$\forall z\,\forall w\,(z = w \to \varphi_{71} \to \varphi_{70})$}
\AXC{$\varphi_{239}$\lbl{A2}}
\BIC{$\forall w\,(a \oplus b(S(n)) = w \to \varphi_{102} \to \varphi_{70})$}
\AXC{$\varphi_{249}$\lbl{A2}}
\BIC{$a \oplus b(S(n)) = a_{S(n)} + b_{S(n)} \to \varphi_{102} \to \varphi_{112}$}
\AXC{$x = x$\lbl{A5}}
\UIC{$\forall x\,(x = x)$}
\AXC{$\forall x\,(x = x) \to \varphi_{102}$\lbl{A2}}
\BIC{$\varphi_{102}$}
\AXC{$\stackrel{\Pi_{45}}{\vdots}$}\noLine
\UIC{$\varphi_{102} \to \varphi_{258}$}
\BIC{$\varphi_{258}$}
\BIC{$\mathrm{Seq}(a) \to \varphi_{211}$}
\end{ptree}
\noindent$\Pi_{47}$:
\begin{ptree}
\AXC{$y = z \to x = y \to x = z$\lbl{A6}}
\AXC{$(y = z \to x = y \to x = z) \to x = y \to y = z \to x = z$\lbl{A1}}
\BIC{$x = y \to y = z \to x = z$}
\UIC{$\forall z\,(x = y \to y = z \to x = z)$}
\UIC{$\forall y\,\forall z\,(x = y \to y = z \to x = z)$}
\UIC{$\forall x\,\forall y\,\forall z\,(x = y \to y = z \to x = z)$}
\AXC{$\forall x\,\forall y\,\forall z\,(x = y \to y = z \to x = z) \to \varphi_{29}$\lbl{A2}}
\BIC{$\varphi_{29}$}
\end{ptree}
\noindent$\Pi_{48}$:
\begin{ptree}
\AXC{$z = w \to \varphi_{141}$\lbl{A6}}
\UIC{$\forall w\,(z = w \to \varphi_{141})$}
\UIC{$\forall z\,\forall w\,(z = w \to \varphi_{141})$}
\AXC{$\forall z\,\forall w\,(z = w \to \varphi_{141}) \to \varphi_{184}$\lbl{A2}}
\BIC{$\varphi_{184}$}
\AXC{$\varphi_{184} \to {\textstyle\sum_{k=1}^{n} a \oplus b(k)} = {\textstyle\sum_{k=1}^{n} a_{k}} + {\textstyle\sum_{k=1}^{n} b_{k}} \to \varphi_{73} \to \varphi_{87}$\lbl{A2}}
\BIC{${\textstyle\sum_{k=1}^{n} a \oplus b(k)} = {\textstyle\sum_{k=1}^{n} a_{k}} + {\textstyle\sum_{k=1}^{n} b_{k}} \to \varphi_{73} \to \varphi_{87}$}
\end{ptree}
\noindent$\Pi_{49}$:
\begin{ptree}
\AXC{$\stackrel{\Pi_{12}}{\vdots}$}\noLine
\UIC{$\mathrm{Seq}(a) \to \varphi_{156}$}
\AXC{$\stackrel{\Pi_{9}}{\vdots}$}\noLine
\UIC{$\mathrm{Seq}(a) \to \mathrm{Seq}(b) \to \forall n\,(\varphi_{68}) \to \varphi_{69}$}
\AXC{$\varphi_{234}$\lbl{A1}}
\BIC{$\mathrm{Seq}(a) \to \mathrm{Seq}(b) \to \varphi_{174}$}
\AXC{$(\mathrm{Seq}(a) \to \mathrm{Seq}(b) \to \varphi_{174}) \to \varphi_{231}$\lbl{A1}}
\BIC{$\varphi_{231}$}
\BIC{$\mathrm{Seq}(a) \to \mathrm{Seq}(b) \to \forall n\,(\varphi_{68}) \to \varphi_{77}$}
\end{ptree}
\noindent$\Pi_{50}$:
\begin{ptree}
\AXC{$\stackrel{\Pi_{49}}{\vdots}$}\noLine
\UIC{$\mathrm{Seq}(a) \to \mathrm{Seq}(b) \to \forall n\,(\varphi_{68}) \to \varphi_{77}$}
\AXC{$(\mathrm{Seq}(a) \to \mathrm{Seq}(b) \to \forall n\,(\varphi_{68}) \to \varphi_{77}) \to \varphi_{73} \to \varphi_{252}$\lbl{A1}}
\BIC{$\varphi_{73} \to \varphi_{252}$}
\end{ptree}
\noindent$\Pi_{51}$:
\begin{ptree}
\AXC{$\stackrel{\Pi_{48}}{\vdots}$}\noLine
\UIC{${\textstyle\sum_{k=1}^{n} a \oplus b(k)} = {\textstyle\sum_{k=1}^{n} a_{k}} + {\textstyle\sum_{k=1}^{n} b_{k}} \to \varphi_{73} \to \varphi_{87}$}
\AXC{$x = x$\lbl{A5}}
\UIC{$\forall x\,(x = x)$}
\AXC{$\forall x\,(x = x) \to \varphi_{73}$\lbl{A2}}
\BIC{$\varphi_{73}$}
\AXC{$\stackrel{\Pi_{50}}{\vdots}$}\noLine
\UIC{$\varphi_{73} \to \varphi_{252}$}
\BIC{$\varphi_{252}$}
\BIC{$\mathrm{Seq}(a) \to \varphi_{206}$}
\end{ptree}
\noindent$\Pi_{52}$:
\begin{ptree}
\AXC{$\forall a\,\forall b\,(\mathrm{Seq}(a) \to \varphi_{40})$\lbl{[S2]}}
\AXC{$\forall a\,\forall b\,(\mathrm{Seq}(a) \to \varphi_{40}) \to \forall b\,(\mathrm{Seq}(a) \to \varphi_{40})$\lbl{A2}}
\BIC{$\forall b\,(\mathrm{Seq}(a) \to \varphi_{40})$}
\end{ptree}
\noindent$\Pi_{53}$:
\begin{ptree}
\AXC{$\mathrm{Seq}(a) \to \varphi_{151}$\lbl{A1}}
\AXC{$\stackrel{\Pi_{52}}{\vdots}$}\noLine
\UIC{$\forall b\,(\mathrm{Seq}(a) \to \varphi_{40})$}
\AXC{$\forall b\,(\mathrm{Seq}(a) \to \varphi_{40}) \to \mathrm{Seq}(a) \to \varphi_{40}$\lbl{A2}}
\BIC{$\mathrm{Seq}(a) \to \varphi_{40}$}
\AXC{$\varphi_{268}$\lbl{A1}}
\BIC{$(\mathrm{Seq}(a) \to \varphi_{151}) \to (\mathrm{Seq}(a) \to \varphi_{150}) \to \varphi_{195}$}
\BIC{$(\mathrm{Seq}(a) \to \varphi_{150}) \to \varphi_{195}$}
\end{ptree}
\noindent$\Pi_{54}$:
\begin{ptree}
\AXC{$\mathrm{Seq}(a) \to \varphi_{150}$\lbl{A1}}
\AXC{$\stackrel{\Pi_{53}}{\vdots}$}\noLine
\UIC{$(\mathrm{Seq}(a) \to \varphi_{150}) \to \varphi_{195}$}
\BIC{$\varphi_{195}$}
\AXC{$\varphi_{195} \to \mathrm{Seq}(a) \to \mathrm{Seq}(b) \to \varphi_{152}$\lbl{A1}}
\BIC{$\mathrm{Seq}(a) \to \mathrm{Seq}(b) \to \varphi_{152}$}
\end{ptree}
\noindent$\Pi_{55}$:
\begin{ptree}
\AXC{$\stackrel{\Pi_{54}}{\vdots}$}\noLine
\UIC{$\mathrm{Seq}(a) \to \mathrm{Seq}(b) \to \varphi_{152}$}
\AXC{$\forall a\,(\mathrm{Seq}(a) \to \varphi_{98})$\lbl{[S25]}}
\AXC{$\forall a\,(\mathrm{Seq}(a) \to \varphi_{98}) \to \mathrm{Seq}(a \oplus b) \to \varphi_{163}$\lbl{A2}}
\BIC{$\mathrm{Seq}(a \oplus b) \to \varphi_{163}$}
\AXC{$\varphi_{270}$\lbl{A1}}
\BIC{$(\mathrm{Seq}(a) \to \mathrm{Seq}(b) \to \varphi_{152}) \to \mathrm{Seq}(a) \to \varphi_{218}$}
\BIC{$\mathrm{Seq}(a) \to \varphi_{218}$}
\end{ptree}
\noindent$\Pi_{56}$:
\begin{ptree}
\AXC{$\stackrel{\Pi_{55}}{\vdots}$}\noLine
\UIC{$\mathrm{Seq}(a) \to \varphi_{218}$}
\AXC{$(\mathrm{Seq}(a) \to \varphi_{218}) \to \mathrm{Seq}(a) \to \varphi_{199}$\lbl{A1}}
\BIC{$\mathrm{Seq}(a) \to \varphi_{199}$}
\AXC{$(\mathrm{Seq}(a) \to \varphi_{199}) \to (\varphi_{61} \to \varphi_{48}) \to \mathrm{Seq}(a) \to \varphi_{198}$\lbl{A1}}
\BIC{$(\varphi_{61} \to \varphi_{48}) \to \mathrm{Seq}(a) \to \varphi_{198}$}
\end{ptree}
\noindent$\Pi_{57}$:
\begin{ptree}
\AXC{$\stackrel{\Pi_{12}}{\vdots}$}\noLine
\UIC{$\mathrm{Seq}(a) \to \varphi_{156}$}
\AXC{$\varphi_{61} \to \varphi_{48}$\lbl{A2}}
\AXC{$\stackrel{\Pi_{56}}{\vdots}$}\noLine
\UIC{$(\varphi_{61} \to \varphi_{48}) \to \mathrm{Seq}(a) \to \varphi_{198}$}
\BIC{$\mathrm{Seq}(a) \to \varphi_{198}$}
\AXC{$(\mathrm{Seq}(a) \to \varphi_{198}) \to (\mathrm{Seq}(a) \to \varphi_{156}) \to \varphi_{196}$\lbl{A1}}
\BIC{$(\mathrm{Seq}(a) \to \varphi_{156}) \to \varphi_{196}$}
\BIC{$\varphi_{196}$}
\end{ptree}
\noindent$\Pi_{58}$:
\begin{ptree}
\AXC{$\stackrel{\Pi_{47}}{\vdots}$}\noLine
\UIC{$\varphi_{29}$}
\AXC{$\varphi_{29} \to \varphi_{158}$\lbl{A2}}
\BIC{$\varphi_{158}$}
\AXC{$\varphi_{158} \to \varphi_{210}$\lbl{A2}}
\BIC{$\varphi_{210}$}
\AXC{$\stackrel{\Pi_{51}}{\vdots}$}\noLine
\UIC{$\mathrm{Seq}(a) \to \varphi_{206}$}
\AXC{$\stackrel{\Pi_{57}}{\vdots}$}\noLine
\UIC{$\varphi_{196}$}
\AXC{$\varphi_{196} \to (\mathrm{Seq}(a) \to \varphi_{206}) \to \varphi_{210} \to \varphi_{200}$\lbl{A1}}
\BIC{$(\mathrm{Seq}(a) \to \varphi_{206}) \to \varphi_{210} \to \varphi_{200}$}
\BIC{$\varphi_{210} \to \varphi_{200}$}
\BIC{$\varphi_{200}$}
\end{ptree}
\noindent$\Pi_{59}$:
\begin{ptree}
\AXC{$\stackrel{\Pi_{41}}{\vdots}$}\noLine
\UIC{$\varphi_{29}$}
\AXC{$\varphi_{29} \to \varphi_{173}$\lbl{A2}}
\BIC{$\varphi_{173}$}
\AXC{$\varphi_{244}$\lbl{A2}}
\BIC{${\textstyle\sum_{k=1}^{S(n)} a \oplus b(k)} = {\textstyle\sum_{k=1}^{n} a_{k}} + {\textstyle\sum_{k=1}^{n} b_{k}} + a \oplus b(S(n)) \to \varphi_{186}$}
\AXC{$\stackrel{\Pi_{46}}{\vdots}$}\noLine
\UIC{$\mathrm{Seq}(a) \to \varphi_{211}$}
\AXC{$\stackrel{\Pi_{58}}{\vdots}$}\noLine
\UIC{$\varphi_{200}$}
\AXC{$\varphi_{200} \to (\mathrm{Seq}(a) \to \varphi_{211}) \to \varphi_{255}$\lbl{A1}}
\BIC{$(\mathrm{Seq}(a) \to \varphi_{211}) \to \varphi_{255}$}
\BIC{$\varphi_{255}$}
\BIC{$\varphi_{203}$}
\end{ptree}
\noindent$\Pi_{60}$:
\begin{ptree}
\AXC{$\stackrel{\Pi_{26}}{\vdots}$}\noLine
\UIC{$\varphi_{29}$}
\AXC{$\varphi_{29} \to \varphi_{178}$\lbl{A2}}
\BIC{$\varphi_{178}$}
\AXC{$\varphi_{178} \to \varphi_{217}$\lbl{A2}}
\BIC{$\varphi_{217}$}
\AXC{$\stackrel{\Pi_{40}}{\vdots}$}\noLine
\UIC{$\mathrm{Seq}(a) \to \varphi_{212}$}
\AXC{$\stackrel{\Pi_{59}}{\vdots}$}\noLine
\UIC{$\varphi_{203}$}
\AXC{$\varphi_{203} \to (\mathrm{Seq}(a) \to \varphi_{212}) \to \varphi_{217} \to \varphi_{204}$\lbl{A1}}
\BIC{$(\mathrm{Seq}(a) \to \varphi_{212}) \to \varphi_{217} \to \varphi_{204}$}
\BIC{$\varphi_{217} \to \varphi_{204}$}
\BIC{$\varphi_{204}$}
\end{ptree}
\noindent$\Pi_{61}$:
\begin{ptree}
\AXC{$\stackrel{\Pi_{19}}{\vdots}$}\noLine
\UIC{$\varphi_{29}$}
\AXC{$\varphi_{29} \to \varphi_{177}$\lbl{A2}}
\BIC{$\varphi_{177}$}
\AXC{$\varphi_{177} \to \varphi_{213}$\lbl{A2}}
\BIC{$\varphi_{213}$}
\AXC{$\stackrel{\Pi_{25}}{\vdots}$}\noLine
\UIC{$\mathrm{Seq}(a) \to \varphi_{209}$}
\AXC{$\stackrel{\Pi_{60}}{\vdots}$}\noLine
\UIC{$\varphi_{204}$}
\AXC{$\varphi_{204} \to \varphi_{274}$\lbl{A1}}
\BIC{$\varphi_{274}$}
\BIC{$\varphi_{213} \to \mathrm{Seq}(a) \to \mathrm{Seq}(b) \to \forall n\,(\varphi_{68}) \to n \in B_{43} \to \varphi_{93}$}
\BIC{$\mathrm{Seq}(a) \to \mathrm{Seq}(b) \to \forall n\,(\varphi_{68}) \to n \in B_{43} \to \varphi_{93}$}
\AXC{$\varphi_{276}$\lbl{A1}}
\BIC{$\varphi_{272}$}
\end{ptree}
\noindent$\Pi_{62}$:
\begin{ptree}
\AXC{$\stackrel{\Pi_{6}}{\vdots}$}\noLine
\UIC{${\textstyle\sum_{k=1}^{S(n)} a \oplus b(k)} = {\textstyle\sum_{k=1}^{S(n)} a_{k}} + ({\textstyle\sum_{k=1}^{n} b_{k}} + b_{S(n)}) \to \varphi_{154}$}
\AXC{$\stackrel{\Pi_{18}}{\vdots}$}\noLine
\UIC{$\mathrm{Seq}(a) \to \mathrm{Seq}(b) \to \forall n\,(\varphi_{68}) \to n \in B_{43} \to \varphi_{111}$}
\AXC{$\stackrel{\Pi_{61}}{\vdots}$}\noLine
\UIC{$\varphi_{272}$}
\BIC{$\varphi_{242}$}
\BIC{$\mathrm{Seq}(a) \to \mathrm{Seq}(b) \to \forall n\,(\varphi_{68}) \to \varphi_{94}$}
\end{ptree}
\noindent$\Pi_{63}$:
\begin{ptree}
\AXC{$y = z \to x = y \to x = z$\lbl{A6}}
\AXC{$(y = z \to x = y \to x = z) \to x = y \to y = z \to x = z$\lbl{A1}}
\BIC{$x = y \to y = z \to x = z$}
\UIC{$\forall z\,(x = y \to y = z \to x = z)$}
\UIC{$\forall y\,\forall z\,(x = y \to y = z \to x = z)$}
\UIC{$\forall x\,\forall y\,\forall z\,(x = y \to y = z \to x = z)$}
\AXC{$\forall x\,\forall y\,\forall z\,(x = y \to y = z \to x = z) \to \varphi_{30}$\lbl{A2}}
\BIC{$\varphi_{30}$}
\end{ptree}
\noindent$\Pi_{64}$:
\begin{ptree}
\AXC{$\stackrel{\Pi_{63}}{\vdots}$}\noLine
\UIC{$\varphi_{30}$}
\AXC{$\varphi_{30} \to \varphi_{121}$\lbl{A2}}
\BIC{$\varphi_{121}$}
\AXC{$\varphi_{121} \to {\textstyle\sum_{k=1}^{1_{\omega}} a \oplus b(k)} = {\textstyle\sum_{k=1}^{1_{\omega}} a_{k}} + b_{1_{\omega}} \to \varphi_{127}$\lbl{A2}}
\BIC{${\textstyle\sum_{k=1}^{1_{\omega}} a \oplus b(k)} = {\textstyle\sum_{k=1}^{1_{\omega}} a_{k}} + b_{1_{\omega}} \to \varphi_{127}$}
\end{ptree}
\noindent$\Pi_{65}$:
\begin{ptree}
\AXC{$z = w \to \varphi_{104}$\lbl{A6}}
\UIC{$\forall w\,(z = w \to \varphi_{104})$}
\UIC{$\forall z\,\forall w\,(z = w \to \varphi_{104})$}
\AXC{$\forall z\,\forall w\,(z = w \to \varphi_{104}) \to \varphi_{129}$\lbl{A2}}
\BIC{$\varphi_{129}$}
\AXC{$\varphi_{129} \to b_{1_{\omega}} = {\textstyle\sum_{k=1}^{1_{\omega}} b_{k}} \to \varphi_{128}$\lbl{A2}}
\BIC{$b_{1_{\omega}} = {\textstyle\sum_{k=1}^{1_{\omega}} b_{k}} \to \varphi_{128}$}
\end{ptree}
\noindent$\Pi_{66}$:
\begin{ptree}
\AXC{$\forall a\,(\mathrm{Seq}(a) \to \varphi_{98})$\lbl{[S25]}}
\AXC{$\forall a\,(\mathrm{Seq}(a) \to \varphi_{98}) \to \mathrm{Seq}(b) \to \varphi_{97}$\lbl{A2}}
\BIC{$\mathrm{Seq}(b) \to \varphi_{97}$}
\AXC{$(\mathrm{Seq}(b) \to \varphi_{97}) \to \varphi_{148}$\lbl{A1}}
\BIC{$\varphi_{148}$}
\end{ptree}
\noindent$\Pi_{67}$:
\begin{ptree}
\AXC{$\mathrm{Seq}(a) \to \mathrm{Seq}(b) \to \mathrm{Seq}(b)$\lbl{A1}}
\AXC{$\stackrel{\Pi_{66}}{\vdots}$}\noLine
\UIC{$\varphi_{148}$}
\BIC{$\mathrm{Seq}(a) \to \mathrm{Seq}(b) \to \varphi_{97}$}
\AXC{$\varphi_{171}$\lbl{A1}}
\BIC{$\mathrm{Seq}(a) \to \mathrm{Seq}(b) \to {\textstyle\sum_{k=1}^{1_{\omega}} b_{k}} = b_{1_{\omega}}$}
\AXC{$(\mathrm{Seq}(a) \to \mathrm{Seq}(b) \to {\textstyle\sum_{k=1}^{1_{\omega}} b_{k}} = b_{1_{\omega}}) \to \varphi_{99}$\lbl{A1}}
\BIC{$\varphi_{99}$}
\end{ptree}
\noindent$\Pi_{68}$:
\begin{ptree}
\AXC{$\stackrel{\Pi_{7}}{\vdots}$}\noLine
\UIC{$\forall x\,\forall y\,(x = y \to y = x)$}
\AXC{$\varphi_{22}$\lbl{A2}}
\BIC{$\forall y\,({\textstyle\sum_{k=1}^{1_{\omega}} b_{k}} = y \to y = {\textstyle\sum_{k=1}^{1_{\omega}} b_{k}})$}
\AXC{$\varphi_{95}$\lbl{A2}}
\BIC{${\textstyle\sum_{k=1}^{1_{\omega}} b_{k}} = b_{1_{\omega}} \to b_{1_{\omega}} = {\textstyle\sum_{k=1}^{1_{\omega}} b_{k}}$}
\AXC{$\stackrel{\Pi_{67}}{\vdots}$}\noLine
\UIC{$\varphi_{99}$}
\BIC{$\mathrm{Seq}(a) \to \mathrm{Seq}(b) \to b_{1_{\omega}} = {\textstyle\sum_{k=1}^{1_{\omega}} b_{k}}$}
\AXC{$(\mathrm{Seq}(a) \to \mathrm{Seq}(b) \to b_{1_{\omega}} = {\textstyle\sum_{k=1}^{1_{\omega}} b_{k}}) \to \varphi_{219}$\lbl{A1}}
\BIC{$\varphi_{219}$}
\end{ptree}
\noindent$\Pi_{69}$:
\begin{ptree}
\AXC{$\stackrel{\Pi_{65}}{\vdots}$}\noLine
\UIC{$b_{1_{\omega}} = {\textstyle\sum_{k=1}^{1_{\omega}} b_{k}} \to \varphi_{128}$}
\AXC{$x = x$\lbl{A5}}
\UIC{$\forall x\,(x = x)$}
\AXC{$\varphi_{33}$\lbl{A2}}
\BIC{${\textstyle\sum_{k=1}^{1_{\omega}} a_{k}} + b_{1_{\omega}} = {\textstyle\sum_{k=1}^{1_{\omega}} a_{k}} + b_{1_{\omega}}$}
\AXC{$\stackrel{\Pi_{68}}{\vdots}$}\noLine
\UIC{$\varphi_{219}$}
\BIC{$(b_{1_{\omega}} = {\textstyle\sum_{k=1}^{1_{\omega}} b_{k}} \to \varphi_{128}) \to \varphi_{60}$}
\BIC{$\varphi_{60}$}
\end{ptree}
\noindent$\Pi_{70}$:
\begin{ptree}
\AXC{$y = z \to x = y \to x = z$\lbl{A6}}
\AXC{$(y = z \to x = y \to x = z) \to x = y \to y = z \to x = z$\lbl{A1}}
\BIC{$x = y \to y = z \to x = z$}
\UIC{$\forall z\,(x = y \to y = z \to x = z)$}
\UIC{$\forall y\,\forall z\,(x = y \to y = z \to x = z)$}
\UIC{$\forall x\,\forall y\,\forall z\,(x = y \to y = z \to x = z)$}
\AXC{$\forall x\,\forall y\,\forall z\,(x = y \to y = z \to x = z) \to \varphi_{30}$\lbl{A2}}
\BIC{$\varphi_{30}$}
\end{ptree}
\noindent$\Pi_{71}$:
\begin{ptree}
\AXC{$\stackrel{\Pi_{70}}{\vdots}$}\noLine
\UIC{$\varphi_{30}$}
\AXC{$\varphi_{30} \to \varphi_{109}$\lbl{A2}}
\BIC{$\varphi_{109}$}
\AXC{$\varphi_{109} \to {\textstyle\sum_{k=1}^{1_{\omega}} a \oplus b(k)} = a_{1_{\omega}} + b_{1_{\omega}} \to \varphi_{114}$\lbl{A2}}
\BIC{${\textstyle\sum_{k=1}^{1_{\omega}} a \oplus b(k)} = a_{1_{\omega}} + b_{1_{\omega}} \to \varphi_{114}$}
\end{ptree}
\noindent$\Pi_{72}$:
\begin{ptree}
\AXC{$\forall a\,(\mathrm{Seq}(a) \to \varphi_{98})$\lbl{[S25]}}
\AXC{$\forall a\,(\mathrm{Seq}(a) \to \varphi_{98}) \to \mathrm{Seq}(a) \to \varphi_{98}$\lbl{A2}}
\BIC{$\mathrm{Seq}(a) \to \varphi_{98}$}
\AXC{$(\mathrm{Seq}(a) \to \varphi_{98}) \to \varphi_{149}$\lbl{A1}}
\BIC{$\varphi_{149}$}
\end{ptree}
\noindent$\Pi_{73}$:
\begin{ptree}
\AXC{$\mathrm{Seq}(a) \to \mathrm{Seq}(b) \to \mathrm{Seq}(a)$\lbl{A1}}
\AXC{$\stackrel{\Pi_{72}}{\vdots}$}\noLine
\UIC{$\varphi_{149}$}
\BIC{$\mathrm{Seq}(a) \to \mathrm{Seq}(b) \to \varphi_{98}$}
\AXC{$\varphi_{172}$\lbl{A1}}
\BIC{$\mathrm{Seq}(a) \to \mathrm{Seq}(b) \to {\textstyle\sum_{k=1}^{1_{\omega}} a_{k}} = a_{1_{\omega}}$}
\AXC{$(\mathrm{Seq}(a) \to \mathrm{Seq}(b) \to {\textstyle\sum_{k=1}^{1_{\omega}} a_{k}} = a_{1_{\omega}}) \to \varphi_{100}$\lbl{A1}}
\BIC{$\varphi_{100}$}
\end{ptree}
\noindent$\Pi_{74}$:
\begin{ptree}
\AXC{$x = x$\lbl{A5}}
\UIC{$\forall x\,(x = x)$}
\AXC{$\forall x\,(x = x) \to \varphi_{6}$\lbl{A2}}
\BIC{$\varphi_{6}$}
\AXC{$\stackrel{\Pi_{7}}{\vdots}$}\noLine
\UIC{$\forall x\,\forall y\,(x = y \to y = x)$}
\AXC{$\varphi_{23}$\lbl{A2}}
\BIC{$\forall y\,({\textstyle\sum_{k=1}^{1_{\omega}} a_{k}} = y \to y = {\textstyle\sum_{k=1}^{1_{\omega}} a_{k}})$}
\AXC{$\varphi_{96}$\lbl{A2}}
\BIC{${\textstyle\sum_{k=1}^{1_{\omega}} a_{k}} = a_{1_{\omega}} \to a_{1_{\omega}} = {\textstyle\sum_{k=1}^{1_{\omega}} a_{k}}$}
\AXC{$\stackrel{\Pi_{73}}{\vdots}$}\noLine
\UIC{$\varphi_{100}$}
\BIC{$\mathrm{Seq}(a) \to \mathrm{Seq}(b) \to a_{1_{\omega}} = {\textstyle\sum_{k=1}^{1_{\omega}} a_{k}}$}
\AXC{$\varphi_{220}$\lbl{A1}}
\BIC{$\varphi_{6} \to \varphi_{138} \to \mathrm{Seq}(a) \to \varphi_{28}$}
\BIC{$\varphi_{138} \to \mathrm{Seq}(a) \to \varphi_{28}$}
\end{ptree}
\noindent$\Pi_{75}$:
\begin{ptree}
\AXC{$y = z \to x = y \to x = z$\lbl{A6}}
\AXC{$(y = z \to x = y \to x = z) \to x = y \to y = z \to x = z$\lbl{A1}}
\BIC{$x = y \to y = z \to x = z$}
\UIC{$\forall z\,(x = y \to y = z \to x = z)$}
\UIC{$\forall y\,\forall z\,(x = y \to y = z \to x = z)$}
\UIC{$\forall x\,\forall y\,\forall z\,(x = y \to y = z \to x = z)$}
\AXC{$\forall x\,\forall y\,\forall z\,(x = y \to y = z \to x = z) \to \varphi_{30}$\lbl{A2}}
\BIC{$\varphi_{30}$}
\end{ptree}
\noindent$\Pi_{76}$:
\begin{ptree}
\AXC{$\forall a\,\forall b\,(\mathrm{Seq}(a) \to \varphi_{40})$\lbl{[S2]}}
\AXC{$\forall a\,\forall b\,(\mathrm{Seq}(a) \to \varphi_{40}) \to \forall b\,(\mathrm{Seq}(a) \to \varphi_{40})$\lbl{A2}}
\BIC{$\forall b\,(\mathrm{Seq}(a) \to \varphi_{40})$}
\end{ptree}
\noindent$\Pi_{77}$:
\begin{ptree}
\AXC{$\stackrel{\Pi_{76}}{\vdots}$}\noLine
\UIC{$\forall b\,(\mathrm{Seq}(a) \to \varphi_{40})$}
\AXC{$\forall b\,(\mathrm{Seq}(a) \to \varphi_{40}) \to \mathrm{Seq}(a) \to \varphi_{40}$\lbl{A2}}
\BIC{$\mathrm{Seq}(a) \to \varphi_{40}$}
\AXC{$(\mathrm{Seq}(a) \to \varphi_{40}) \to \varphi_{31}$\lbl{A1}}
\BIC{$\varphi_{31}$}
\AXC{$\varphi_{31} \to \varphi_{183}$\lbl{A1}}
\BIC{$\varphi_{183}$}
\end{ptree}
\noindent$\Pi_{78}$:
\begin{ptree}
\AXC{$\forall a\,\forall b\,(\mathrm{Seq}(a) \to \varphi_{40})$\lbl{[S2]}}
\AXC{$\forall a\,\forall b\,(\mathrm{Seq}(a) \to \varphi_{40}) \to \forall b\,(\mathrm{Seq}(a) \to \varphi_{40})$\lbl{A2}}
\BIC{$\forall b\,(\mathrm{Seq}(a) \to \varphi_{40})$}
\end{ptree}
\noindent$\Pi_{79}$:
\begin{ptree}
\AXC{$\forall a\,(\mathrm{Seq}(a) \to \varphi_{98})$\lbl{[S25]}}
\AXC{$\forall a\,(\mathrm{Seq}(a) \to \varphi_{98}) \to \mathrm{Seq}(a \oplus b) \to \varphi_{163}$\lbl{A2}}
\BIC{$\mathrm{Seq}(a \oplus b) \to \varphi_{163}$}
\AXC{$(\mathrm{Seq}(a \oplus b) \to \varphi_{163}) \to \varphi_{188}$\lbl{A1}}
\BIC{$\varphi_{188}$}
\end{ptree}
\noindent$\Pi_{80}$:
\begin{ptree}
\AXC{$\stackrel{\Pi_{78}}{\vdots}$}\noLine
\UIC{$\forall b\,(\mathrm{Seq}(a) \to \varphi_{40})$}
\AXC{$\forall b\,(\mathrm{Seq}(a) \to \varphi_{40}) \to \mathrm{Seq}(a) \to \varphi_{40}$\lbl{A2}}
\BIC{$\mathrm{Seq}(a) \to \varphi_{40}$}
\AXC{$\varphi_{119}$\lbl{A1}}
\BIC{$\mathrm{Seq}(a) \to \mathrm{Seq}(b) \to \mathrm{Seq}(a \oplus b)$}
\AXC{$\stackrel{\Pi_{79}}{\vdots}$}\noLine
\UIC{$\varphi_{188}$}
\BIC{$\mathrm{Seq}(a) \to \mathrm{Seq}(b) \to \varphi_{163}$}
\end{ptree}
\noindent$\Pi_{81}$:
\begin{ptree}
\AXC{$\stackrel{\Pi_{80}}{\vdots}$}\noLine
\UIC{$\mathrm{Seq}(a) \to \mathrm{Seq}(b) \to \varphi_{163}$}
\AXC{$(\mathrm{Seq}(a) \to \mathrm{Seq}(b) \to \varphi_{163}) \to \varphi_{26}$\lbl{A1}}
\BIC{$\varphi_{26}$}
\AXC{$\varphi_{26} \to (\mathrm{Seq}(a) \to \varphi_{12}) \to \varphi_{194}$\lbl{A1}}
\BIC{$(\mathrm{Seq}(a) \to \varphi_{12}) \to \varphi_{194}$}
\end{ptree}
\noindent$\Pi_{82}$:
\begin{ptree}
\AXC{$1_{\omega} \in \mathbb{N}$\lbl{[N1]}}
\AXC{$\forall n \in \mathbb{N}\,(a \oplus b(n) = a_{n} + b_{n}) \to \varphi_{24}$\lbl{A2}}
\AXC{$\stackrel{\Pi_{77}}{\vdots}$}\noLine
\UIC{$\varphi_{183}$}
\BIC{$\mathrm{Seq}(a) \to \mathrm{Seq}(b) \to \varphi_{24}$}
\AXC{$\varphi_{162}$\lbl{A1}}
\BIC{$1_{\omega} \in \mathbb{N} \to \mathrm{Seq}(a) \to \varphi_{12}$}
\BIC{$\mathrm{Seq}(a) \to \varphi_{12}$}
\AXC{$\stackrel{\Pi_{81}}{\vdots}$}\noLine
\UIC{$(\mathrm{Seq}(a) \to \varphi_{12}) \to \varphi_{194}$}
\BIC{$\varphi_{194}$}
\end{ptree}
\noindent$\Pi_{83}$:
\begin{ptree}
\AXC{$\stackrel{\Pi_{75}}{\vdots}$}\noLine
\UIC{$\varphi_{30}$}
\AXC{$\varphi_{30} \to \varphi_{92}$\lbl{A2}}
\BIC{$\varphi_{92}$}
\AXC{$\varphi_{92} \to {\textstyle\sum_{k=1}^{1_{\omega}} a \oplus b(k)} = a \oplus b(1_{\omega}) \to \varphi_{85}$\lbl{A2}}
\BIC{${\textstyle\sum_{k=1}^{1_{\omega}} a \oplus b(k)} = a \oplus b(1_{\omega}) \to \varphi_{85}$}
\AXC{$\stackrel{\Pi_{82}}{\vdots}$}\noLine
\UIC{$\varphi_{194}$}
\BIC{$\mathrm{Seq}(a) \to \varphi_{25}$}
\AXC{$(\mathrm{Seq}(a) \to \varphi_{25}) \to \varphi_{222}$\lbl{A1}}
\BIC{$\varphi_{222}$}
\end{ptree}
\noindent$\Pi_{84}$:
\begin{ptree}
\AXC{$z = w \to \varphi_{62}$\lbl{A6}}
\UIC{$\forall w\,(z = w \to \varphi_{62})$}
\UIC{$\forall z\,\forall w\,(z = w \to \varphi_{62})$}
\AXC{$\forall z\,\forall w\,(z = w \to \varphi_{62}) \to \varphi_{110}$\lbl{A2}}
\BIC{$\varphi_{110}$}
\AXC{$\varphi_{110} \to \varphi_{138}$\lbl{A2}}
\BIC{$\varphi_{138}$}
\AXC{$\stackrel{\Pi_{74}}{\vdots}$}\noLine
\UIC{$\varphi_{138} \to \mathrm{Seq}(a) \to \varphi_{28}$}
\BIC{$\mathrm{Seq}(a) \to \varphi_{28}$}
\AXC{$\stackrel{\Pi_{83}}{\vdots}$}\noLine
\UIC{$\varphi_{222}$}
\BIC{$({\textstyle\sum_{k=1}^{1_{\omega}} a \oplus b(k)} = a_{1_{\omega}} + b_{1_{\omega}} \to \varphi_{114}) \to \varphi_{38}$}
\end{ptree}
\noindent$\Pi_{85}$:
\begin{ptree}
\AXC{$\stackrel{\Pi_{71}}{\vdots}$}\noLine
\UIC{${\textstyle\sum_{k=1}^{1_{\omega}} a \oplus b(k)} = a_{1_{\omega}} + b_{1_{\omega}} \to \varphi_{114}$}
\AXC{$\stackrel{\Pi_{84}}{\vdots}$}\noLine
\UIC{$({\textstyle\sum_{k=1}^{1_{\omega}} a \oplus b(k)} = a_{1_{\omega}} + b_{1_{\omega}} \to \varphi_{114}) \to \varphi_{38}$}
\BIC{$\varphi_{38}$}
\AXC{$\varphi_{241}$\lbl{A1}}
\BIC{$\varphi_{60} \to ({\textstyle\sum_{k=1}^{1_{\omega}} a \oplus b(k)} = {\textstyle\sum_{k=1}^{1_{\omega}} a_{k}} + b_{1_{\omega}} \to \varphi_{127}) \to \varphi_{44}$}
\end{ptree}
\noindent$\Pi_{86}$:
\begin{ptree}
\AXC{$y = z \to x = y \to x = z$\lbl{A6}}
\AXC{$(y = z \to x = y \to x = z) \to x = y \to y = z \to x = z$\lbl{A1}}
\BIC{$x = y \to y = z \to x = z$}
\UIC{$\forall z\,(x = y \to y = z \to x = z)$}
\UIC{$\forall y\,\forall z\,(x = y \to y = z \to x = z)$}
\UIC{$\forall x\,\forall y\,\forall z\,(x = y \to y = z \to x = z)$}
\AXC{$\varphi_{9}$\lbl{A2}}
\BIC{$\forall y\,\forall z\,(S(n) = y \to y = z \to S(n) = z)$}
\AXC{$\varphi_{34}$\lbl{A2}}
\BIC{$\forall z\,(S(n) = S(\emptyset) \to S(\emptyset) = z \to S(n) = z)$}
\end{ptree}
\noindent$\Pi_{87}$:
\begin{ptree}
\AXC{$\stackrel{\Pi_{86}}{\vdots}$}\noLine
\UIC{$\forall z\,(S(n) = S(\emptyset) \to S(\emptyset) = z \to S(n) = z)$}
\AXC{$\forall z\,(S(n) = S(\emptyset) \to S(\emptyset) = z \to S(n) = z) \to \varphi_{5}$\lbl{A2}}
\BIC{$\varphi_{5}$}
\end{ptree}
\noindent$\Pi_{88}$:
\begin{ptree}
\AXC{$\stackrel{\Pi_{7}}{\vdots}$}\noLine
\UIC{$\forall x\,\forall y\,(x = y \to y = x)$}
\AXC{$\varphi_{3}$\lbl{A2}}
\BIC{$\forall y\,(1_{\omega} = y \to y = 1_{\omega})$}
\AXC{$\varphi_{11}$\lbl{A2}}
\BIC{$1_{\omega} = S(\emptyset) \to S(\emptyset) = 1_{\omega}$}
\AXC{$1_{\omega} = S(\emptyset)$\lbl{Def$1_\omega$}}
\AXC{$1_{\omega} = S(\emptyset) \to \varphi_{10}$\lbl{A1}}
\BIC{$\varphi_{10}$}
\BIC{$S(\emptyset) = 1_{\omega}$}
\end{ptree}
\noindent$\Pi_{89}$:
\begin{ptree}
\AXC{$x = x$\lbl{A5}}
\UIC{$\forall x\,(x = x)$}
\AXC{$\forall x\,(x = x) \to S(n) = S(n)$\lbl{A2}}
\BIC{$S(n) = S(n)$}
\AXC{$\mathrm{Seq}(a) \to \varphi_{157}$\lbl{A1}}
\AXC{$(\mathrm{Seq}(a) \to \varphi_{157}) \to S(n) = S(n) \to \varphi_{190}$\lbl{A1}}
\BIC{$S(n) = S(n) \to \varphi_{190}$}
\BIC{$\varphi_{190}$}
\end{ptree}
\noindent$\Pi_{90}$:
\begin{ptree}
\AXC{$z = w \to S(n) = S(z) \to S(n) = S(w)$\lbl{A6}}
\UIC{$\forall w\,(z = w \to S(n) = S(z) \to S(n) = S(w))$}
\UIC{$\forall z\,\forall w\,(z = w \to S(n) = S(z) \to S(n) = S(w))$}
\AXC{$\varphi_{27}$\lbl{A2}}
\BIC{$\forall w\,(n = w \to S(n) = S(n) \to S(n) = S(w))$}
\AXC{$\varphi_{35}$\lbl{A2}}
\BIC{$n = \emptyset \to S(n) = S(n) \to S(n) = S(\emptyset)$}
\AXC{$\stackrel{\Pi_{89}}{\vdots}$}\noLine
\UIC{$\varphi_{190}$}
\BIC{$\mathrm{Seq}(a) \to \mathrm{Seq}(b) \to \varphi_{153}$}
\AXC{$\varphi_{235}$\lbl{A1}}
\BIC{$S(\emptyset) = 1_{\omega} \to \varphi_{5} \to \mathrm{Seq}(a) \to \mathrm{Seq}(b) \to \varphi_{142}$}
\end{ptree}
\noindent$\Pi_{91}$:
\begin{ptree}
\AXC{$\stackrel{\Pi_{87}}{\vdots}$}\noLine
\UIC{$\varphi_{5}$}
\AXC{$\stackrel{\Pi_{88}}{\vdots}$}\noLine
\UIC{$S(\emptyset) = 1_{\omega}$}
\AXC{$\stackrel{\Pi_{90}}{\vdots}$}\noLine
\UIC{$S(\emptyset) = 1_{\omega} \to \varphi_{5} \to \mathrm{Seq}(a) \to \mathrm{Seq}(b) \to \varphi_{142}$}
\BIC{$\varphi_{5} \to \mathrm{Seq}(a) \to \mathrm{Seq}(b) \to \varphi_{142}$}
\BIC{$\mathrm{Seq}(a) \to \mathrm{Seq}(b) \to \varphi_{142}$}
\end{ptree}
\noindent$\Pi_{92}$:
\begin{ptree}
\AXC{$\stackrel{\Pi_{91}}{\vdots}$}\noLine
\UIC{$\mathrm{Seq}(a) \to \mathrm{Seq}(b) \to \varphi_{142}$}
\AXC{$(\mathrm{Seq}(a) \to \mathrm{Seq}(b) \to \varphi_{142}) \to \varphi_{179}$\lbl{A1}}
\BIC{$\varphi_{179}$}
\end{ptree}
\noindent$\Pi_{93}$:
\begin{ptree}
\AXC{$\stackrel{\Pi_{7}}{\vdots}$}\noLine
\UIC{$\forall x\,\forall y\,(x = y \to y = x)$}
\AXC{$\varphi_{2}$\lbl{A2}}
\BIC{$\forall y\,(S(n) = y \to y = S(n))$}
\AXC{$\varphi_{4}$\lbl{A2}}
\BIC{$S(n) = 1_{\omega} \to 1_{\omega} = S(n)$}
\AXC{$\stackrel{\Pi_{92}}{\vdots}$}\noLine
\UIC{$\varphi_{179}$}
\BIC{$\mathrm{Seq}(a) \to \mathrm{Seq}(b) \to \varphi_{143}$}
\AXC{$(\mathrm{Seq}(a) \to \mathrm{Seq}(b) \to \varphi_{143}) \to \varphi_{44} \to \varphi_{227}$\lbl{A1}}
\BIC{$\varphi_{44} \to \varphi_{227}$}
\end{ptree}
\noindent$\Pi_{94}$:
\begin{ptree}
\AXC{$\stackrel{\Pi_{64}}{\vdots}$}\noLine
\UIC{${\textstyle\sum_{k=1}^{1_{\omega}} a \oplus b(k)} = {\textstyle\sum_{k=1}^{1_{\omega}} a_{k}} + b_{1_{\omega}} \to \varphi_{127}$}
\AXC{$\stackrel{\Pi_{69}}{\vdots}$}\noLine
\UIC{$\varphi_{60}$}
\AXC{$\stackrel{\Pi_{85}}{\vdots}$}\noLine
\UIC{$\varphi_{60} \to ({\textstyle\sum_{k=1}^{1_{\omega}} a \oplus b(k)} = {\textstyle\sum_{k=1}^{1_{\omega}} a_{k}} + b_{1_{\omega}} \to \varphi_{127}) \to \varphi_{44}$}
\BIC{$({\textstyle\sum_{k=1}^{1_{\omega}} a \oplus b(k)} = {\textstyle\sum_{k=1}^{1_{\omega}} a_{k}} + b_{1_{\omega}} \to \varphi_{127}) \to \varphi_{44}$}
\BIC{$\varphi_{44}$}
\AXC{$\stackrel{\Pi_{93}}{\vdots}$}\noLine
\UIC{$\varphi_{44} \to \varphi_{227}$}
\BIC{$\varphi_{227}$}
\end{ptree}
\noindent$\Pi_{95}$:
\begin{ptree}
\AXC{$z = w \to \varphi_{107}$\lbl{A6}}
\UIC{$\forall w\,(z = w \to \varphi_{107})$}
\UIC{$\forall z\,\forall w\,(z = w \to \varphi_{107})$}
\AXC{$\forall z\,\forall w\,(z = w \to \varphi_{107}) \to \varphi_{124}$\lbl{A2}}
\BIC{$\varphi_{124}$}
\AXC{$\varphi_{124} \to 1_{\omega} = S(n) \to \varphi_{123}$\lbl{A2}}
\BIC{$1_{\omega} = S(n) \to \varphi_{123}$}
\AXC{$\stackrel{\Pi_{94}}{\vdots}$}\noLine
\UIC{$\varphi_{227}$}
\BIC{$\mathrm{Seq}(a) \to \mathrm{Seq}(b) \to \forall n\,(\varphi_{68}) \to \varphi_{89}$}
\end{ptree}
\noindent$\Pi_{96}$:
\begin{ptree}
\AXC{$\stackrel{\Pi_{95}}{\vdots}$}\noLine
\UIC{$\mathrm{Seq}(a) \to \mathrm{Seq}(b) \to \forall n\,(\varphi_{68}) \to \varphi_{89}$}
\AXC{$(\mathrm{Seq}(a) \to \mathrm{Seq}(b) \to \forall n\,(\varphi_{68}) \to \varphi_{89}) \to \varphi_{238}$\lbl{A1}}
\BIC{$\varphi_{238}$}
\end{ptree}
\noindent$\Pi_{97}$:
\begin{ptree}
\AXC{$\forall x\,(\mathrm{Ind}(x) \leftrightarrow \emptyset \in x \land \forall y \in x\,(S(y) \in x))$\lbl{DefInd}}
\AXC{$\varphi_{113}$\lbl{A2}}
\BIC{$\mathrm{Ind}(\omega) \leftrightarrow \emptyset \in \omega \land \forall y \in \omega\,(S(y) \in \omega)$}
\AXC{$\mathrm{Ind}(\omega)$\lbl{補題(ωは帰納的)}}
\AXC{$\varphi_{205}$\lbl{A1}}
\BIC{$(\mathrm{Ind}(\omega) \leftrightarrow \emptyset \in \omega \land \forall y \in \omega\,(S(y) \in \omega)) \to \varphi_{169}$}
\BIC{$\varphi_{169}$}
\end{ptree}
\noindent$\Pi_{98}$:
\begin{ptree}
\AXC{$\forall y \in \omega\,(S(y) \in \omega) \to n \in \omega \to S(n) \in \omega$\lbl{A2}}
\AXC{$\stackrel{\Pi_{97}}{\vdots}$}\noLine
\UIC{$\varphi_{169}$}
\AXC{$\varphi_{169} \to \mathrm{Seq}(a) \to \varphi_{144}$\lbl{A1}}
\BIC{$\mathrm{Seq}(a) \to \varphi_{144}$}
\AXC{$\varphi_{232}$\lbl{A1}}
\BIC{$(\forall y \in \omega\,(S(y) \in \omega) \to n \in \omega \to S(n) \in \omega) \to \varphi_{155}$}
\BIC{$\varphi_{155}$}
\AXC{$\varphi_{155} \to \varphi_{214}$\lbl{A1}}
\BIC{$\varphi_{214}$}
\end{ptree}
\noindent$\Pi_{99}$:
\begin{ptree}
\AXC{$\mathrm{Seq}(a) \to \mathrm{Seq}(b) \to \forall n\,(\varphi_{68}) \to \forall n\,(\varphi_{68})$\lbl{A1}}
\AXC{$(\mathrm{Seq}(a) \to \mathrm{Seq}(b) \to \forall n\,(\varphi_{68}) \to \forall n\,(\varphi_{68})) \to \varphi_{240}$\lbl{A1}}
\BIC{$\varphi_{240}$}
\end{ptree}
\noindent$\Pi_{100}$:
\begin{ptree}
\AXC{$\stackrel{\Pi_{10}}{\vdots}$}\noLine
\UIC{$\mathrm{Seq}(a) \to \mathrm{Seq}(b) \to \forall n\,(\varphi_{68}) \to n \in B_{43} \to n \in \omega$}
\AXC{$\stackrel{\Pi_{98}}{\vdots}$}\noLine
\UIC{$\varphi_{214}$}
\BIC{$\varphi_{136}$}
\AXC{$\forall n\,(\varphi_{68}) \to \varphi_{101}$\lbl{A2}}
\AXC{$\stackrel{\Pi_{99}}{\vdots}$}\noLine
\UIC{$\varphi_{240}$}
\BIC{$\mathrm{Seq}(a) \to \mathrm{Seq}(b) \to \forall n\,(\varphi_{68}) \to \varphi_{101}$}
\AXC{$\varphi_{273}$\lbl{A1}}
\BIC{$\varphi_{136} \to (\mathrm{Seq}(a) \to \mathrm{Seq}(b) \to \forall n\,(\varphi_{68}) \to \varphi_{67}) \to \varphi_{130}$}
\BIC{$(\mathrm{Seq}(a) \to \mathrm{Seq}(b) \to \forall n\,(\varphi_{68}) \to \varphi_{67}) \to \varphi_{130}$}
\end{ptree}
\noindent$\Pi_{101}$:
\begin{ptree}
\AXC{$\stackrel{\Pi_{62}}{\vdots}$}\noLine
\UIC{$\mathrm{Seq}(a) \to \mathrm{Seq}(b) \to \forall n\,(\varphi_{68}) \to \varphi_{94}$}
\AXC{$\stackrel{\Pi_{96}}{\vdots}$}\noLine
\UIC{$\varphi_{238}$}
\BIC{$\mathrm{Seq}(a) \to \mathrm{Seq}(b) \to \forall n\,(\varphi_{68}) \to \varphi_{67}$}
\AXC{$\stackrel{\Pi_{100}}{\vdots}$}\noLine
\UIC{$(\mathrm{Seq}(a) \to \mathrm{Seq}(b) \to \forall n\,(\varphi_{68}) \to \varphi_{67}) \to \varphi_{130}$}
\BIC{$\varphi_{130}$}
\AXC{$\varphi_{130} \to \varphi_{131}$\lbl{A1}}
\BIC{$\varphi_{131}$}
\UIC{$\forall n\,(\varphi_{131})$}
\end{ptree}
\noindent$\Pi_{102}$:
\begin{ptree}
\AXC{$\varphi_{197}$\lbl{A4}}
\AXC{$\forall n\,(\varphi_{131}) \to \varphi_{139}$\lbl{A3}}
\AXC{$\stackrel{\Pi_{101}}{\vdots}$}\noLine
\UIC{$\forall n\,(\varphi_{131})$}
\AXC{$\forall n\,(\varphi_{131}) \to (\forall n\,(\varphi_{131}) \to \varphi_{139}) \to \varphi_{197} \to \varphi_{137}$\lbl{A1}}
\BIC{$(\forall n\,(\varphi_{131}) \to \varphi_{139}) \to \varphi_{197} \to \varphi_{137}$}
\BIC{$\varphi_{197} \to \varphi_{137}$}
\BIC{$\varphi_{137}$}
\end{ptree}
\noindent$\Pi_{103}$:
\begin{ptree}
\AXC{$\forall n\,(n \in \mathbb{N} \leftrightarrow n \in \omega \land n \neq \emptyset)$\lbl{Def$\mathbb{N}$}}
\AXC{$\varphi_{82}$\lbl{A2}}
\BIC{$\emptyset \in \mathbb{N} \leftrightarrow \emptyset \in \omega \land \emptyset \neq \emptyset$}
\AXC{$(\emptyset \in \mathbb{N} \leftrightarrow \emptyset \in \omega \land \emptyset \neq \emptyset) \to \emptyset = \emptyset \to \varphi_{105}$\lbl{A1}}
\BIC{$\emptyset = \emptyset \to \varphi_{105}$}
\end{ptree}
\noindent$\Pi_{104}$:
\begin{ptree}
\AXC{$\mathrm{Seq}(a) \to \mathrm{Seq}(b) \to \forall n\,(\varphi_{68}) \to \forall n\,(\varphi_{68})$\lbl{A1}}
\AXC{$(\mathrm{Seq}(a) \to \mathrm{Seq}(b) \to \forall n\,(\varphi_{68}) \to \forall n\,(\varphi_{68})) \to \varphi_{248}$\lbl{A1}}
\BIC{$\varphi_{248}$}
\end{ptree}
\noindent$\Pi_{105}$:
\begin{ptree}
\AXC{$x = x$\lbl{A5}}
\UIC{$\forall x\,(x = x)$}
\AXC{$\forall x\,(x = x) \to \emptyset = \emptyset$\lbl{A2}}
\BIC{$\emptyset = \emptyset$}
\AXC{$\stackrel{\Pi_{103}}{\vdots}$}\noLine
\UIC{$\emptyset = \emptyset \to \varphi_{105}$}
\BIC{$\varphi_{105}$}
\AXC{$\stackrel{\Pi_{97}}{\vdots}$}\noLine
\UIC{$\varphi_{169}$}
\AXC{$\forall n\,(\varphi_{68}) \to \varphi_{125}$\lbl{A2}}
\AXC{$\stackrel{\Pi_{104}}{\vdots}$}\noLine
\UIC{$\varphi_{248}$}
\BIC{$\mathrm{Seq}(a) \to \mathrm{Seq}(b) \to \forall n\,(\varphi_{68}) \to \varphi_{125}$}
\AXC{$\varphi_{269}$\lbl{A1}}
\BIC{$\varphi_{169} \to \varphi_{105} \to \mathrm{Seq}(a) \to \mathrm{Seq}(b) \to \forall n\,(\varphi_{68}) \to \emptyset \in B_{43}$}
\BIC{$\varphi_{105} \to \mathrm{Seq}(a) \to \mathrm{Seq}(b) \to \forall n\,(\varphi_{68}) \to \emptyset \in B_{43}$}
\BIC{$\mathrm{Seq}(a) \to \mathrm{Seq}(b) \to \forall n\,(\varphi_{68}) \to \emptyset \in B_{43}$}
\end{ptree}
\noindent$\Pi_{106}$:
\begin{ptree}
\AXC{$\forall x\,(\mathrm{Ind}(x) \leftrightarrow \emptyset \in x \land \forall y \in x\,(S(y) \in x))$\lbl{DefInd}}
\AXC{$\forall x\,(\mathrm{Ind}(x) \leftrightarrow \emptyset \in x \land \forall y \in x\,(S(y) \in x)) \to \varphi_{15}$\lbl{A2}}
\BIC{$\varphi_{15}$}
\AXC{$\varphi_{15} \to \varphi_{253}$\lbl{A1}}
\BIC{$\varphi_{253}$}
\end{ptree}
\noindent$\Pi_{107}$:
\begin{ptree}
\AXC{$\stackrel{\Pi_{3}}{\vdots}$}\noLine
\UIC{$\varphi_{13}$}
\AXC{$\stackrel{\Pi_{4}}{\vdots}$}\noLine
\UIC{$\varphi_{14}$}
\AXC{$\stackrel{\Pi_{102}}{\vdots}$}\noLine
\UIC{$\varphi_{137}$}
\AXC{$\stackrel{\Pi_{105}}{\vdots}$}\noLine
\UIC{$\mathrm{Seq}(a) \to \mathrm{Seq}(b) \to \forall n\,(\varphi_{68}) \to \emptyset \in B_{43}$}
\AXC{$\stackrel{\Pi_{106}}{\vdots}$}\noLine
\UIC{$\varphi_{253}$}
\BIC{$\varphi_{137} \to \varphi_{202}$}
\BIC{$\varphi_{202}$}
\BIC{$\varphi_{13} \to \mathrm{Seq}(a) \to \mathrm{Seq}(b) \to \forall n\,(\varphi_{68}) \to \mathrm{Ind}(B_{43})$}
\BIC{$\mathrm{Seq}(a) \to \mathrm{Seq}(b) \to \forall n\,(\varphi_{68}) \to \mathrm{Ind}(B_{43})$}
\end{ptree}
\noindent$\Pi_{108}$:
\begin{ptree}
\AXC{$\forall n\,(n \in \omega \leftrightarrow \forall x\,(\mathrm{Ind}(x) \to n \in x))$\lbl{Def$\omega$}}
\AXC{$\varphi_{32}$\lbl{A2}}
\BIC{$n \in \omega \leftrightarrow \forall x\,(\mathrm{Ind}(x) \to n \in x)$}
\AXC{$\mathrm{Seq}(a) \to \mathrm{Seq}(b) \to \forall n\,(\varphi_{68}) \to n \in \omega \to n \in \omega$\lbl{A1}}
\AXC{$\varphi_{221}$\lbl{A1}}
\BIC{$(n \in \omega \leftrightarrow \forall x\,(\mathrm{Ind}(x) \to n \in x)) \to \mathrm{Seq}(a) \to \varphi_{132}$}
\BIC{$\mathrm{Seq}(a) \to \varphi_{132}$}
\end{ptree}
\noindent$\Pi_{109}$:
\begin{ptree}
\AXC{$\forall x\,(\mathrm{Ind}(x) \to n \in x) \to \mathrm{Ind}(B_{43}) \to n \in B_{43}$\lbl{A2}}
\AXC{$\stackrel{\Pi_{108}}{\vdots}$}\noLine
\UIC{$\mathrm{Seq}(a) \to \varphi_{132}$}
\AXC{$(\mathrm{Seq}(a) \to \varphi_{132}) \to \varphi_{181}$\lbl{A1}}
\BIC{$\varphi_{181}$}
\BIC{$\mathrm{Seq}(a) \to \varphi_{133}$}
\end{ptree}
\noindent$\Pi_{110}$:
\begin{ptree}
\AXC{$\stackrel{\Pi_{107}}{\vdots}$}\noLine
\UIC{$\mathrm{Seq}(a) \to \mathrm{Seq}(b) \to \forall n\,(\varphi_{68}) \to \mathrm{Ind}(B_{43})$}
\AXC{$\stackrel{\Pi_{109}}{\vdots}$}\noLine
\UIC{$\mathrm{Seq}(a) \to \varphi_{133}$}
\AXC{$(\mathrm{Seq}(a) \to \varphi_{133}) \to \varphi_{254}$\lbl{A1}}
\BIC{$\varphi_{254}$}
\BIC{$\varphi_{236}$}
\end{ptree}
\noindent$\Pi_{111}$:
\begin{ptree}
\AXC{$\forall n\,(\varphi_{182}) \to \varphi_{185}$\lbl{A3}}
\AXC{$\forall n\,(\varphi_{68}) \to \varphi_{68}$\lbl{A2}}
\AXC{$\stackrel{\Pi_{2}}{\vdots}$}\noLine
\UIC{$\varphi_{233}$}
\BIC{$\mathrm{Seq}(a) \to \mathrm{Seq}(b) \to \forall n\,(\varphi_{68}) \to \varphi_{68}$}
\AXC{$\stackrel{\Pi_{110}}{\vdots}$}\noLine
\UIC{$\varphi_{236}$}
\BIC{$\mathrm{Seq}(a) \to \mathrm{Seq}(b) \to \varphi_{175}$}
\AXC{$(\mathrm{Seq}(a) \to \mathrm{Seq}(b) \to \varphi_{175}) \to \varphi_{182}$\lbl{A1}}
\BIC{$\varphi_{182}$}
\UIC{$\forall n\,(\varphi_{182})$}
\AXC{$\varphi_{279}$\lbl{A1}}
\BIC{$(\forall n\,(\varphi_{182}) \to \varphi_{185}) \to \varphi_{197} \to \mathrm{Seq}(a) \to \mathrm{Seq}(b) \to \varphi_{176}$}
\BIC{$\varphi_{197} \to \mathrm{Seq}(a) \to \mathrm{Seq}(b) \to \varphi_{176}$}
\end{ptree}
\noindent$\Pi_{112}$:
\begin{ptree}
\AXC{$\varphi_{197}$\lbl{A4}}
\AXC{$\stackrel{\Pi_{111}}{\vdots}$}\noLine
\UIC{$\varphi_{197} \to \mathrm{Seq}(a) \to \mathrm{Seq}(b) \to \varphi_{176}$}
\BIC{$\mathrm{Seq}(a) \to \mathrm{Seq}(b) \to \varphi_{176}$}
\AXC{$(\mathrm{Seq}(a) \to \mathrm{Seq}(b) \to \varphi_{176}) \to \varphi_{86} \to \neg \forall n\,(\varphi_{68})$\lbl{A1}}
\BIC{$\varphi_{86} \to \neg \forall n\,(\varphi_{68})$}
\UIC{$\forall B_{43}\,(\varphi_{86} \to \neg \forall n\,(\varphi_{68}))$}
\end{ptree}
\noindent$\Pi_{113}$:
\begin{ptree}
\AXC{$\varphi_{115}$\lbl{Sep}}
\AXC{$\varphi_{115} \to \varphi_{106}$\lbl{A2}}
\BIC{$\varphi_{106}$}
\AXC{$\varphi_{106} \to \varphi_{88}$\lbl{A2}}
\BIC{$\varphi_{88}$}
\AXC{$\varphi_{88} \to \exists B_{43}\,\forall n\,(\varphi_{68})$\lbl{A2}}
\BIC{$\exists B_{43}\,\forall n\,(\varphi_{68})$}
\AXC{$\varphi_{193}$\lbl{A1}}
\BIC{$\mathrm{Seq}(a) \to \mathrm{Seq}(b) \to \exists B_{43}\,\forall n\,(\varphi_{68})$}
\end{ptree}
\noindent$\Pi_{114}$:
\begin{ptree}
\AXC{$\stackrel{\Pi_{113}}{\vdots}$}\noLine
\UIC{$\mathrm{Seq}(a) \to \mathrm{Seq}(b) \to \exists B_{43}\,\forall n\,(\varphi_{68})$}
\AXC{$(\mathrm{Seq}(a) \to \mathrm{Seq}(b) \to \exists B_{43}\,\forall n\,(\varphi_{68})) \to \varphi_{277}$\lbl{A1}}
\BIC{$\varphi_{277}$}
\end{ptree}
\noindent$\Pi_{115}$:
\begin{ptree}
\AXC{$\varphi_{86} \to \forall B_{43}\,(\varphi_{86})$\lbl{A4}}
\AXC{$\varphi_{237}$\lbl{A3}}
\AXC{$\stackrel{\Pi_{112}}{\vdots}$}\noLine
\UIC{$\forall B_{43}\,(\varphi_{86} \to \neg \forall n\,(\varphi_{68}))$}
\AXC{$\stackrel{\Pi_{114}}{\vdots}$}\noLine
\UIC{$\varphi_{277}$}
\BIC{$\varphi_{237} \to (\varphi_{86} \to \forall B_{43}\,(\varphi_{86})) \to \mathrm{Seq}(a) \to \varphi_{72}$}
\BIC{$(\varphi_{86} \to \forall B_{43}\,(\varphi_{86})) \to \mathrm{Seq}(a) \to \varphi_{72}$}
\BIC{$\mathrm{Seq}(a) \to \varphi_{72}$}
\AXC{$(\mathrm{Seq}(a) \to \varphi_{72}) \to \varphi_{168} \to \mathrm{Seq}(a) \to \varphi_{65}$\lbl{A1}}
\BIC{$\varphi_{168} \to \mathrm{Seq}(a) \to \varphi_{65}$}
\end{ptree}
\noindent$\Pi_{116}$:
\begin{ptree}
\AXC{$\varphi_{168}$\lbl{A2}}
\AXC{$\stackrel{\Pi_{115}}{\vdots}$}\noLine
\UIC{$\varphi_{168} \to \mathrm{Seq}(a) \to \varphi_{65}$}
\BIC{$\mathrm{Seq}(a) \to \varphi_{65}$}
\AXC{$(\mathrm{Seq}(a) \to \varphi_{65}) \to \varphi_{170}$\lbl{A1}}
\BIC{$\varphi_{170}$}
\end{ptree}
\noindent$\Pi_{117}$:
\begin{ptree}
\AXC{$\mathrm{Seq}(a) \to \mathrm{Seq}(b) \to n \in \mathbb{N} \to n \in \mathbb{N}$\lbl{A1}}
\AXC{$\forall n\,(n \in \mathbb{N} \leftrightarrow n \in \omega \land n \neq \emptyset)$\lbl{Def$\mathbb{N}$}}
\AXC{$\varphi_{39}$\lbl{A2}}
\BIC{$n \in \mathbb{N} \leftrightarrow n \in \omega \land n \neq \emptyset$}
\AXC{$\stackrel{\Pi_{1}}{\vdots}$}\noLine
\UIC{$\varphi_{42}$}
\BIC{$\mathrm{Seq}(a) \to \mathrm{Seq}(b) \to n \in \mathbb{N} \to n \in \omega$}
\AXC{$\stackrel{\Pi_{116}}{\vdots}$}\noLine
\UIC{$\varphi_{170}$}
\BIC{$(\mathrm{Seq}(a) \to \mathrm{Seq}(b) \to n \in \mathbb{N} \to n \in \mathbb{N}) \to \varphi_{46}$}
\BIC{$\varphi_{46}$}
\end{ptree}
\noindent 主木:
\begin{ptree}
\AXC{$\varphi_{1}$\lbl{A4}}
\AXC{$\forall n\,(\varphi_{47}) \to \varphi_{66}$\lbl{A3}}
\AXC{$\stackrel{\Pi_{117}}{\vdots}$}\noLine
\UIC{$\varphi_{46}$}
\AXC{$\varphi_{46} \to \varphi_{47}$\lbl{A1}}
\BIC{$\varphi_{47}$}
\UIC{$\forall n\,(\varphi_{47})$}
\AXC{$\forall n\,(\varphi_{47}) \to (\forall n\,(\varphi_{47}) \to \varphi_{66}) \to \varphi_{1} \to \mathrm{Seq}(a) \to \varphi_{41}$\lbl{A1}}
\BIC{$(\forall n\,(\varphi_{47}) \to \varphi_{66}) \to \varphi_{1} \to \mathrm{Seq}(a) \to \varphi_{41}$}
\BIC{$\varphi_{1} \to \mathrm{Seq}(a) \to \varphi_{41}$}
\BIC{$\mathrm{Seq}(a) \to \varphi_{41}$}
\UIC{$\forall b\,(\mathrm{Seq}(a) \to \varphi_{41})$}
\UIC{$\forall a\,\forall b\,(\mathrm{Seq}(a) \to \varphi_{41})$}
\end{ptree}
\begin{abbrevlist}
\item $\varphi_{1} :\equiv \mathrm{Seq}(a) \land \mathrm{Seq}(b) \to \forall n\,(\mathrm{Seq}(a) \land \mathrm{Seq}(b))$
\item $\varphi_{2} :\equiv \forall x\,\forall y\,(x = y \to y = x) \to \forall y\,(S(n) = y \to y = S(n))$
\item $\varphi_{3} :\equiv \forall x\,\forall y\,(x = y \to y = x) \to \forall y\,(1_{\omega} = y \to y = 1_{\omega})$
\item $\varphi_{4} :\equiv \forall y\,(S(n) = y \to y = S(n)) \to S(n) = 1_{\omega} \to 1_{\omega} = S(n)$
\item $\varphi_{5} :\equiv S(n) = S(\emptyset) \to S(\emptyset) = 1_{\omega} \to S(n) = 1_{\omega}$
\item $\varphi_{6} :\equiv a_{1_{\omega}} + b_{1_{\omega}} = a_{1_{\omega}} + b_{1_{\omega}}$
\item $\varphi_{7} :\equiv \forall n \in B_{43}\,(S(n) \in B_{43}) \to y \in B_{43} \to S(y) \in B_{43}$
\item $\varphi_{8} :\equiv \forall y \in B_{43}\,(S(y) \in B_{43}) \to n \in B_{43} \to S(n) \in B_{43}$
\item $\varphi_{9} :\equiv \forall x\,\forall y\,\forall z\,(x = y \to y = z \to x = z) \to \forall y\,\forall z\,(S(n) = y \to y = z \to S(n) = z)$
\item $\varphi_{10} :\equiv (1_{\omega} = S(\emptyset) \to S(\emptyset) = 1_{\omega}) \to S(\emptyset) = 1_{\omega}$
\item $\varphi_{11} :\equiv \forall y\,(1_{\omega} = y \to y = 1_{\omega}) \to 1_{\omega} = S(\emptyset) \to S(\emptyset) = 1_{\omega}$
\item $\varphi_{12} :\equiv \mathrm{Seq}(b) \to a \oplus b(1_{\omega}) = a_{1_{\omega}} + b_{1_{\omega}}$
\item $\varphi_{13} :\equiv \forall n \in B_{43}\,(S(n) \in B_{43}) \to \forall y \in B_{43}\,(S(y) \in B_{43})$
\item $\varphi_{14} :\equiv \forall y \in B_{43}\,(S(y) \in B_{43}) \to \forall n \in B_{43}\,(S(n) \in B_{43})$
\item $\varphi_{15} :\equiv \mathrm{Ind}(B_{43}) \leftrightarrow \emptyset \in B_{43} \land \forall y \in B_{43}\,(S(y) \in B_{43})$
\item $\varphi_{16} :\equiv \forall n \in B_{43}\,(S(n) \in B_{43}) \to \forall y\,\forall n \in B_{43}\,(S(n) \in B_{43})$
\item $\varphi_{17} :\equiv \forall y \in B_{43}\,(S(y) \in B_{43}) \to \forall n\,\forall y \in B_{43}\,(S(y) \in B_{43})$
\item $\varphi_{18} :\equiv \forall y\,\forall n \in B_{43}\,(S(n) \in B_{43}) \to \forall y \in B_{43}\,(S(y) \in B_{43})$
\item $\varphi_{19} :\equiv \forall n\,\forall y \in B_{43}\,(S(y) \in B_{43}) \to \forall n \in B_{43}\,(S(n) \in B_{43})$
\item $\varphi_{20} :\equiv \forall x\,\forall y\,(x = y \to y = x) \to \forall y\,({\textstyle\sum_{k=1}^{S(n)} b_{k}} = y \to y = {\textstyle\sum_{k=1}^{S(n)} b_{k}})$
\item $\varphi_{21} :\equiv \forall x\,\forall y\,(x = y \to y = x) \to \forall y\,({\textstyle\sum_{k=1}^{S(n)} a_{k}} = y \to y = {\textstyle\sum_{k=1}^{S(n)} a_{k}})$
\item $\varphi_{22} :\equiv \forall x\,\forall y\,(x = y \to y = x) \to \forall y\,({\textstyle\sum_{k=1}^{1_{\omega}} b_{k}} = y \to y = {\textstyle\sum_{k=1}^{1_{\omega}} b_{k}})$
\item $\varphi_{23} :\equiv \forall x\,\forall y\,(x = y \to y = x) \to \forall y\,({\textstyle\sum_{k=1}^{1_{\omega}} a_{k}} = y \to y = {\textstyle\sum_{k=1}^{1_{\omega}} a_{k}})$
\item $\varphi_{24} :\equiv 1_{\omega} \in \mathbb{N} \to a \oplus b(1_{\omega}) = a_{1_{\omega}} + b_{1_{\omega}}$
\item $\varphi_{25} :\equiv \mathrm{Seq}(b) \to {\textstyle\sum_{k=1}^{1_{\omega}} a \oplus b(k)} = a_{1_{\omega}} + b_{1_{\omega}}$
\item $\varphi_{26} :\equiv \mathrm{Seq}(a) \to \mathrm{Seq}(b) \to {\textstyle\sum_{k=1}^{1_{\omega}} a \oplus b(k)} = a \oplus b(1_{\omega})$
\item $\varphi_{27} :\equiv \forall z\,\forall w\,(z = w \to S(n) = S(z) \to S(n) = S(w)) \to \forall w\,(n = w \to S(n) = S(n) \to S(n) = S(w))$
\item $\varphi_{28} :\equiv \mathrm{Seq}(b) \to a_{1_{\omega}} + b_{1_{\omega}} = {\textstyle\sum_{k=1}^{1_{\omega}} a_{k}} + b_{1_{\omega}}$
\item $\varphi_{29} :\equiv \forall y\,\forall z\,({\textstyle\sum_{k=1}^{S(n)} a \oplus b(k)} = y \to y = z \to {\textstyle\sum_{k=1}^{S(n)} a \oplus b(k)} = z)$
\item $\varphi_{30} :\equiv \forall y\,\forall z\,({\textstyle\sum_{k=1}^{1_{\omega}} a \oplus b(k)} = y \to y = z \to {\textstyle\sum_{k=1}^{1_{\omega}} a \oplus b(k)} = z)$
\item $\varphi_{31} :\equiv \mathrm{Seq}(a) \to \mathrm{Seq}(b) \to \forall n \in \mathbb{N}\,(a \oplus b(n) = a_{n} + b_{n})$
\item $\varphi_{32} :\equiv \forall n\,(n \in \omega \leftrightarrow \forall x\,(\mathrm{Ind}(x) \to n \in x)) \to (n \in \omega \leftrightarrow \forall x\,(\mathrm{Ind}(x) \to n \in x))$
\item $\varphi_{33} :\equiv \forall x\,(x = x) \to {\textstyle\sum_{k=1}^{1_{\omega}} a_{k}} + b_{1_{\omega}} = {\textstyle\sum_{k=1}^{1_{\omega}} a_{k}} + b_{1_{\omega}}$
\item $\varphi_{34} :\equiv \forall y\,\forall z\,(S(n) = y \to y = z \to S(n) = z) \to \forall z\,(S(n) = S(\emptyset) \to S(\emptyset) = z \to S(n) = z)$
\item $\varphi_{35} :\equiv \forall w\,(n = w \to S(n) = S(n) \to S(n) = S(w)) \to n = \emptyset \to S(n) = S(n) \to S(n) = S(\emptyset)$
\item $\varphi_{36} :\equiv n \neq \emptyset \to S(n) \in \mathbb{N} \to a \oplus b(S(n)) = a_{S(n)} + b_{S(n)}$
\item $\varphi_{37} :\equiv S(n) \in \mathbb{N} \to n \neq \emptyset \to a \oplus b(S(n)) = a_{S(n)} + b_{S(n)}$
\item $\varphi_{38} :\equiv \mathrm{Seq}(a) \to \mathrm{Seq}(b) \to {\textstyle\sum_{k=1}^{1_{\omega}} a \oplus b(k)} = {\textstyle\sum_{k=1}^{1_{\omega}} a_{k}} + b_{1_{\omega}}$
\item $\varphi_{39} :\equiv \forall n\,(n \in \mathbb{N} \leftrightarrow n \in \omega \land n \neq \emptyset) \to (n \in \mathbb{N} \leftrightarrow n \in \omega \land n \neq \emptyset)$
\item $\varphi_{40} :\equiv \mathrm{Seq}(b) \to \mathrm{Seq}(a \oplus b) \land \forall n \in \mathbb{N}\,(a \oplus b(n) = a_{n} + b_{n})$
\item $\varphi_{41} :\equiv \mathrm{Seq}(b) \to \forall n \in \mathbb{N}\,({\textstyle\sum_{k=1}^{n} a \oplus b(k)} = {\textstyle\sum_{k=1}^{n} a_{k}} + {\textstyle\sum_{k=1}^{n} b_{k}})$
\item $\varphi_{42} :\equiv (n \in \mathbb{N} \leftrightarrow n \in \omega \land n \neq \emptyset) \to \mathrm{Seq}(a) \to \mathrm{Seq}(b) \to n \in \mathbb{N} \to n \in \omega$
\item $\varphi_{43} :\equiv S(n) \in \mathbb{N} \to n \neq \emptyset \to \forall n \in \mathbb{N}\,(a \oplus b(n) = a_{n} + b_{n})$
\item $\varphi_{44} :\equiv \mathrm{Seq}(a) \to \mathrm{Seq}(b) \to {\textstyle\sum_{k=1}^{1_{\omega}} a \oplus b(k)} = {\textstyle\sum_{k=1}^{1_{\omega}} a_{k}} + {\textstyle\sum_{k=1}^{1_{\omega}} b_{k}}$
\item $\varphi_{45} :\equiv n \neq \emptyset \to \mathrm{Seq}(a \oplus b) \land \forall n \in \mathbb{N}\,(a \oplus b(n) = a_{n} + b_{n})$
\item $\varphi_{46} :\equiv \mathrm{Seq}(a) \to \mathrm{Seq}(b) \to n \in \mathbb{N} \to {\textstyle\sum_{k=1}^{n} a \oplus b(k)} = {\textstyle\sum_{k=1}^{n} a_{k}} + {\textstyle\sum_{k=1}^{n} b_{k}}$
\item $\varphi_{47} :\equiv \mathrm{Seq}(a) \land \mathrm{Seq}(b) \to n \in \mathbb{N} \to {\textstyle\sum_{k=1}^{n} a \oplus b(k)} = {\textstyle\sum_{k=1}^{n} a_{k}} + {\textstyle\sum_{k=1}^{n} b_{k}}$
\item $\varphi_{48} :\equiv n \in \mathbb{N} \to {\textstyle\sum_{k=1}^{S(n)} a \oplus b(k)} = {\textstyle\sum_{k=1}^{n} a \oplus b(k)} + a \oplus b(S(n))$
\item $\varphi_{49} :\equiv \forall a\,\forall n\,(\mathrm{Seq}(a) \to n \in \mathbb{N} \to a_{n} \in \mathbb{R}) \to \forall n\,(\mathrm{Seq}(b) \to n \in \mathbb{N} \to b_{n} \in \mathbb{R})$
\item $\varphi_{50} :\equiv \forall n\,(\mathrm{Seq}(b) \to n \in \mathbb{N} \to b_{n} \in \mathbb{R}) \to \mathrm{Seq}(b) \to S(n) \in \mathbb{N} \to b_{S(n)} \in \mathbb{R}$
\item $\varphi_{51} :\equiv \forall a\,\forall n\,(\mathrm{Seq}(a) \to n \in \mathbb{N} \to a_{n} \in \mathbb{R}) \to \forall n\,(\mathrm{Seq}(a) \to n \in \mathbb{N} \to a_{n} \in \mathbb{R})$
\item $\varphi_{52} :\equiv \forall n\,(\mathrm{Seq}(a) \to n \in \mathbb{N} \to a_{n} \in \mathbb{R}) \to \mathrm{Seq}(a) \to S(n) \in \mathbb{N} \to a_{S(n)} \in \mathbb{R}$
\item $\varphi_{53} :\equiv S(n) \in \mathbb{N} \to n \neq \emptyset \to n \in \mathbb{N} \to {\textstyle\sum_{k=1}^{S(n)} b_{k}} = {\textstyle\sum_{k=1}^{n} b_{k}} + b_{S(n)}$
\item $\varphi_{54} :\equiv n \in B_{43} \to S(n) \in \mathbb{N} \to n \neq \emptyset \to {\textstyle\sum_{k=1}^{S(n)} b_{k}} = {\textstyle\sum_{k=1}^{n} b_{k}} + b_{S(n)}$
\item $\varphi_{55} :\equiv n \in B_{43} \to S(n) \in \mathbb{N} \to n \neq \emptyset \to {\textstyle\sum_{k=1}^{n} b_{k}} + b_{S(n)} = {\textstyle\sum_{k=1}^{S(n)} b_{k}}$
\item $\varphi_{56} :\equiv S(n) \in \mathbb{N} \to n \neq \emptyset \to n \in \mathbb{N} \to {\textstyle\sum_{k=1}^{S(n)} a_{k}} = {\textstyle\sum_{k=1}^{n} a_{k}} + a_{S(n)}$
\item $\varphi_{57} :\equiv n \in B_{43} \to S(n) \in \mathbb{N} \to n \neq \emptyset \to {\textstyle\sum_{k=1}^{S(n)} a_{k}} = {\textstyle\sum_{k=1}^{n} a_{k}} + a_{S(n)}$
\item $\varphi_{58} :\equiv n \in B_{43} \to S(n) \in \mathbb{N} \to n \neq \emptyset \to {\textstyle\sum_{k=1}^{n} a_{k}} + a_{S(n)} = {\textstyle\sum_{k=1}^{S(n)} a_{k}}$
\item $\varphi_{59} :\equiv n \neq \emptyset \to {\textstyle\sum_{k=1}^{S(n)} a \oplus b(k)} = {\textstyle\sum_{k=1}^{n} a \oplus b(k)} + a \oplus b(S(n))$
\item $\varphi_{60} :\equiv \mathrm{Seq}(a) \to \mathrm{Seq}(b) \to {\textstyle\sum_{k=1}^{1_{\omega}} a_{k}} + b_{1_{\omega}} = {\textstyle\sum_{k=1}^{1_{\omega}} a_{k}} + {\textstyle\sum_{k=1}^{1_{\omega}} b_{k}}$
\item $\varphi_{61} :\equiv \forall n \in \mathbb{N}\,({\textstyle\sum_{k=1}^{S(n)} a \oplus b(k)} = {\textstyle\sum_{k=1}^{n} a \oplus b(k)} + a \oplus b(S(n)))$
\item $\varphi_{62} :\equiv a_{1_{\omega}} + b_{1_{\omega}} = z + b_{1_{\omega}} \to a_{1_{\omega}} + b_{1_{\omega}} = w + b_{1_{\omega}}$
\item $\varphi_{63} :\equiv S(n) \in \mathbb{N} \to n \neq \emptyset \to \forall n \in \mathbb{N}\,({\textstyle\sum_{k=1}^{S(n)} b_{k}} = {\textstyle\sum_{k=1}^{n} b_{k}} + b_{S(n)})$
\item $\varphi_{64} :\equiv S(n) \in \mathbb{N} \to n \neq \emptyset \to \forall n \in \mathbb{N}\,({\textstyle\sum_{k=1}^{S(n)} a_{k}} = {\textstyle\sum_{k=1}^{n} a_{k}} + a_{S(n)})$
\item $\varphi_{65} :\equiv \mathrm{Seq}(b) \to n \in \omega \to n \in \mathbb{N} \to {\textstyle\sum_{k=1}^{n} a \oplus b(k)} = {\textstyle\sum_{k=1}^{n} a_{k}} + {\textstyle\sum_{k=1}^{n} b_{k}}$
\item $\varphi_{66} :\equiv \forall n\,(\mathrm{Seq}(a) \land \mathrm{Seq}(b)) \to \forall n \in \mathbb{N}\,({\textstyle\sum_{k=1}^{n} a \oplus b(k)} = {\textstyle\sum_{k=1}^{n} a_{k}} + {\textstyle\sum_{k=1}^{n} b_{k}})$
\item $\varphi_{67} :\equiv n \in B_{43} \to S(n) \in \mathbb{N} \to {\textstyle\sum_{k=1}^{S(n)} a \oplus b(k)} = {\textstyle\sum_{k=1}^{S(n)} a_{k}} + {\textstyle\sum_{k=1}^{S(n)} b_{k}}$
\item $\varphi_{68} :\equiv n \in B_{43} \leftrightarrow n \in \omega \land (n \in \mathbb{N} \to {\textstyle\sum_{k=1}^{n} a \oplus b(k)} = {\textstyle\sum_{k=1}^{n} a_{k}} + {\textstyle\sum_{k=1}^{n} b_{k}})$
\item $\varphi_{69} :\equiv n \in B_{43} \to n \in \omega \land (n \in \mathbb{N} \to {\textstyle\sum_{k=1}^{n} a \oplus b(k)} = {\textstyle\sum_{k=1}^{n} a_{k}} + {\textstyle\sum_{k=1}^{n} b_{k}})$
\item $\varphi_{70} :\equiv {\textstyle\sum_{k=1}^{n} a_{k}} + {\textstyle\sum_{k=1}^{n} b_{k}} + a \oplus b(S(n)) = {\textstyle\sum_{k=1}^{n} a_{k}} + {\textstyle\sum_{k=1}^{n} b_{k}} + w$
\item $\varphi_{71} :\equiv {\textstyle\sum_{k=1}^{n} a_{k}} + {\textstyle\sum_{k=1}^{n} b_{k}} + a \oplus b(S(n)) = {\textstyle\sum_{k=1}^{n} a_{k}} + {\textstyle\sum_{k=1}^{n} b_{k}} + z$
\item $\varphi_{72} :\equiv \mathrm{Seq}(b) \to \forall n \in \omega\,(n \in \mathbb{N} \to {\textstyle\sum_{k=1}^{n} a \oplus b(k)} = {\textstyle\sum_{k=1}^{n} a_{k}} + {\textstyle\sum_{k=1}^{n} b_{k}})$
\item $\varphi_{73} :\equiv {\textstyle\sum_{k=1}^{n} a \oplus b(k)} + a \oplus b(S(n)) = {\textstyle\sum_{k=1}^{n} a \oplus b(k)} + a \oplus b(S(n))$
\item $\varphi_{74} :\equiv \forall x\,\forall y\,\forall u\,\forall v\,(x \in \mathbb{R} \to y \in \mathbb{R} \to u \in \mathbb{R} \to v \in \mathbb{R} \to x + y + (u + v) = x + u + (y + v))$
\item $\varphi_{75} :\equiv {\textstyle\sum_{k=1}^{n} a_{k}} + a_{S(n)} + ({\textstyle\sum_{k=1}^{n} b_{k}} + b_{S(n)}) = w + ({\textstyle\sum_{k=1}^{n} b_{k}} + b_{S(n)})$
\item $\varphi_{76} :\equiv {\textstyle\sum_{k=1}^{n} a_{k}} + a_{S(n)} + ({\textstyle\sum_{k=1}^{n} b_{k}} + b_{S(n)}) = z + ({\textstyle\sum_{k=1}^{n} b_{k}} + b_{S(n)})$
\item $\varphi_{77} :\equiv n \in B_{43} \to S(n) \in \mathbb{N} \to n \neq \emptyset \to {\textstyle\sum_{k=1}^{n} a \oplus b(k)} = {\textstyle\sum_{k=1}^{n} a_{k}} + {\textstyle\sum_{k=1}^{n} b_{k}}$
\item $\varphi_{78} :\equiv n \neq \emptyset \to {\textstyle\sum_{k=1}^{S(n)} a \oplus b(k)} = {\textstyle\sum_{k=1}^{n} a_{k}} + {\textstyle\sum_{k=1}^{n} b_{k}} + a \oplus b(S(n))$
\item $\varphi_{79} :\equiv \forall a\,\forall n\,(\mathrm{Seq}(a) \to n \in \mathbb{N} \to a_{n} \in \mathbb{R}) \to \forall n\,(\mathrm{Seq}(\Sigma b) \to n \in \mathbb{N} \to {\textstyle\sum_{k=1}^{n} b_{k}} \in \mathbb{R})$
\item $\varphi_{80} :\equiv \forall a\,\forall n\,(\mathrm{Seq}(a) \to n \in \mathbb{N} \to a_{n} \in \mathbb{R}) \to \forall n\,(\mathrm{Seq}(\Sigma a) \to n \in \mathbb{N} \to {\textstyle\sum_{k=1}^{n} a_{k}} \in \mathbb{R})$
\item $\varphi_{81} :\equiv {\textstyle\sum_{k=1}^{S(n)} a_{k}} + ({\textstyle\sum_{k=1}^{n} b_{k}} + b_{S(n)}) = {\textstyle\sum_{k=1}^{S(n)} a_{k}} + ({\textstyle\sum_{k=1}^{n} b_{k}} + b_{S(n)})$
\item $\varphi_{82} :\equiv \forall n\,(n \in \mathbb{N} \leftrightarrow n \in \omega \land n \neq \emptyset) \to (\emptyset \in \mathbb{N} \leftrightarrow \emptyset \in \omega \land \emptyset \neq \emptyset)$
\item $\varphi_{83} :\equiv n \neq \emptyset \to {\textstyle\sum_{k=1}^{S(n)} a \oplus b(k)} = {\textstyle\sum_{k=1}^{n} a_{k}} + {\textstyle\sum_{k=1}^{n} b_{k}} + (a_{S(n)} + b_{S(n)})$
\item $\varphi_{84} :\equiv n \neq \emptyset \to {\textstyle\sum_{k=1}^{S(n)} a \oplus b(k)} = {\textstyle\sum_{k=1}^{n} a_{k}} + a_{S(n)} + ({\textstyle\sum_{k=1}^{n} b_{k}} + b_{S(n)})$
\item $\varphi_{85} :\equiv a \oplus b(1_{\omega}) = a_{1_{\omega}} + b_{1_{\omega}} \to {\textstyle\sum_{k=1}^{1_{\omega}} a \oplus b(k)} = a_{1_{\omega}} + b_{1_{\omega}}$
\item $\varphi_{86} :\equiv \mathrm{Seq}(a) \land \mathrm{Seq}(b) \land \neg \forall n \in \omega\,(n \in \mathbb{N} \to {\textstyle\sum_{k=1}^{n} a \oplus b(k)} = {\textstyle\sum_{k=1}^{n} a_{k}} + {\textstyle\sum_{k=1}^{n} b_{k}})$
\item $\varphi_{87} :\equiv {\textstyle\sum_{k=1}^{n} a \oplus b(k)} + a \oplus b(S(n)) = {\textstyle\sum_{k=1}^{n} a_{k}} + {\textstyle\sum_{k=1}^{n} b_{k}} + a \oplus b(S(n))$
\item $\varphi_{88} :\equiv \forall A_{44}\,\exists B_{43}\,\forall n\,(n \in B_{43} \leftrightarrow n \in A_{44} \land (n \in \mathbb{N} \to {\textstyle\sum_{k=1}^{n} a \oplus b(k)} = {\textstyle\sum_{k=1}^{n} a_{k}} + {\textstyle\sum_{k=1}^{n} b_{k}}))$
\item $\varphi_{89} :\equiv n \in B_{43} \to S(n) \in \mathbb{N} \to n = \emptyset \to {\textstyle\sum_{k=1}^{S(n)} a \oplus b(k)} = {\textstyle\sum_{k=1}^{S(n)} a_{k}} + {\textstyle\sum_{k=1}^{S(n)} b_{k}}$
\item $\varphi_{90} :\equiv \forall n\,(\mathrm{Seq}(\Sigma b) \to n \in \mathbb{N} \to {\textstyle\sum_{k=1}^{n} b_{k}} \in \mathbb{R}) \to \mathrm{Seq}(\Sigma b) \to n \in \mathbb{N} \to {\textstyle\sum_{k=1}^{n} b_{k}} \in \mathbb{R}$
\item $\varphi_{91} :\equiv \forall n\,(\mathrm{Seq}(\Sigma a) \to n \in \mathbb{N} \to {\textstyle\sum_{k=1}^{n} a_{k}} \in \mathbb{R}) \to \mathrm{Seq}(\Sigma a) \to n \in \mathbb{N} \to {\textstyle\sum_{k=1}^{n} a_{k}} \in \mathbb{R}$
\item $\varphi_{92} :\equiv \forall z\,({\textstyle\sum_{k=1}^{1_{\omega}} a \oplus b(k)} = a \oplus b(1_{\omega}) \to a \oplus b(1_{\omega}) = z \to {\textstyle\sum_{k=1}^{1_{\omega}} a \oplus b(k)} = z)$
\item $\varphi_{93} :\equiv S(n) \in \mathbb{N} \to n \neq \emptyset \to {\textstyle\sum_{k=1}^{S(n)} a \oplus b(k)} = {\textstyle\sum_{k=1}^{S(n)} a_{k}} + ({\textstyle\sum_{k=1}^{n} b_{k}} + b_{S(n)})$
\item $\varphi_{94} :\equiv n \in B_{43} \to S(n) \in \mathbb{N} \to n \neq \emptyset \to {\textstyle\sum_{k=1}^{S(n)} a \oplus b(k)} = {\textstyle\sum_{k=1}^{S(n)} a_{k}} + {\textstyle\sum_{k=1}^{S(n)} b_{k}}$
\item $\varphi_{95} :\equiv \forall y\,({\textstyle\sum_{k=1}^{1_{\omega}} b_{k}} = y \to y = {\textstyle\sum_{k=1}^{1_{\omega}} b_{k}}) \to {\textstyle\sum_{k=1}^{1_{\omega}} b_{k}} = b_{1_{\omega}} \to b_{1_{\omega}} = {\textstyle\sum_{k=1}^{1_{\omega}} b_{k}}$
\item $\varphi_{96} :\equiv \forall y\,({\textstyle\sum_{k=1}^{1_{\omega}} a_{k}} = y \to y = {\textstyle\sum_{k=1}^{1_{\omega}} a_{k}}) \to {\textstyle\sum_{k=1}^{1_{\omega}} a_{k}} = a_{1_{\omega}} \to a_{1_{\omega}} = {\textstyle\sum_{k=1}^{1_{\omega}} a_{k}}$
\item $\varphi_{97} :\equiv \mathrm{Seq}(\Sigma b) \land {\textstyle\sum_{k=1}^{1_{\omega}} b_{k}} = b_{1_{\omega}} \land \forall n \in \mathbb{N}\,({\textstyle\sum_{k=1}^{S(n)} b_{k}} = {\textstyle\sum_{k=1}^{n} b_{k}} + b_{S(n)})$
\item $\varphi_{98} :\equiv \mathrm{Seq}(\Sigma a) \land {\textstyle\sum_{k=1}^{1_{\omega}} a_{k}} = a_{1_{\omega}} \land \forall n \in \mathbb{N}\,({\textstyle\sum_{k=1}^{S(n)} a_{k}} = {\textstyle\sum_{k=1}^{n} a_{k}} + a_{S(n)})$
\item $\varphi_{99} :\equiv ({\textstyle\sum_{k=1}^{1_{\omega}} b_{k}} = b_{1_{\omega}} \to b_{1_{\omega}} = {\textstyle\sum_{k=1}^{1_{\omega}} b_{k}}) \to \mathrm{Seq}(a) \to \mathrm{Seq}(b) \to b_{1_{\omega}} = {\textstyle\sum_{k=1}^{1_{\omega}} b_{k}}$
\item $\varphi_{100} :\equiv ({\textstyle\sum_{k=1}^{1_{\omega}} a_{k}} = a_{1_{\omega}} \to a_{1_{\omega}} = {\textstyle\sum_{k=1}^{1_{\omega}} a_{k}}) \to \mathrm{Seq}(a) \to \mathrm{Seq}(b) \to a_{1_{\omega}} = {\textstyle\sum_{k=1}^{1_{\omega}} a_{k}}$
\item $\varphi_{101} :\equiv S(n) \in B_{43} \leftrightarrow S(n) \in \omega \land (S(n) \in \mathbb{N} \to {\textstyle\sum_{k=1}^{S(n)} a \oplus b(k)} = {\textstyle\sum_{k=1}^{S(n)} a_{k}} + {\textstyle\sum_{k=1}^{S(n)} b_{k}})$
\item $\varphi_{102} :\equiv {\textstyle\sum_{k=1}^{n} a_{k}} + {\textstyle\sum_{k=1}^{n} b_{k}} + a \oplus b(S(n)) = {\textstyle\sum_{k=1}^{n} a_{k}} + {\textstyle\sum_{k=1}^{n} b_{k}} + a \oplus b(S(n))$
\item $\varphi_{103} :\equiv {\textstyle\sum_{k=1}^{n} a_{k}} + a_{S(n)} + ({\textstyle\sum_{k=1}^{n} b_{k}} + b_{S(n)}) = {\textstyle\sum_{k=1}^{S(n)} a_{k}} + ({\textstyle\sum_{k=1}^{n} b_{k}} + b_{S(n)})$
\item $\varphi_{104} :\equiv {\textstyle\sum_{k=1}^{1_{\omega}} a_{k}} + b_{1_{\omega}} = {\textstyle\sum_{k=1}^{1_{\omega}} a_{k}} + z \to {\textstyle\sum_{k=1}^{1_{\omega}} a_{k}} + b_{1_{\omega}} = {\textstyle\sum_{k=1}^{1_{\omega}} a_{k}} + w$
\item $\varphi_{105} :\equiv \mathrm{Seq}(a) \to \mathrm{Seq}(b) \to \emptyset \in \mathbb{N} \to {\textstyle\sum_{k=1}^{\emptyset} a \oplus b(k)} = {\textstyle\sum_{k=1}^{\emptyset} a_{k}} + {\textstyle\sum_{k=1}^{\emptyset} b_{k}}$
\item $\varphi_{106} :\equiv \forall z_{46}\,\forall A_{44}\,\exists B_{43}\,\forall n\,(n \in B_{43} \leftrightarrow n \in A_{44} \land (n \in \mathbb{N} \to {\textstyle\sum_{k=1}^{n} a \oplus z_{46}(k)} = {\textstyle\sum_{k=1}^{n} a_{k}} + {\textstyle\sum_{k=1}^{n} z_{46}(k)}))$
\item $\varphi_{107} :\equiv {\textstyle\sum_{k=1}^{z} a \oplus b(k)} = {\textstyle\sum_{k=1}^{z} a_{k}} + {\textstyle\sum_{k=1}^{z} b_{k}} \to {\textstyle\sum_{k=1}^{w} a \oplus b(k)} = {\textstyle\sum_{k=1}^{w} a_{k}} + {\textstyle\sum_{k=1}^{w} b_{k}}$
\item $\varphi_{108} :\equiv \forall n \in \mathbb{N}\,(a \oplus b(n) = a_{n} + b_{n}) \to S(n) \in \mathbb{N} \to a \oplus b(S(n)) = a_{S(n)} + b_{S(n)}$
\item $\varphi_{109} :\equiv \forall z\,({\textstyle\sum_{k=1}^{1_{\omega}} a \oplus b(k)} = a_{1_{\omega}} + b_{1_{\omega}} \to a_{1_{\omega}} + b_{1_{\omega}} = z \to {\textstyle\sum_{k=1}^{1_{\omega}} a \oplus b(k)} = z)$
\item $\varphi_{110} :\equiv \forall w\,(a_{1_{\omega}} = w \to \varphi_{6} \to a_{1_{\omega}} + b_{1_{\omega}} = w + b_{1_{\omega}})$
\item $\varphi_{111} :\equiv S(n) \in \mathbb{N} \to n \neq \emptyset \to {\textstyle\sum_{k=1}^{S(n)} a_{k}} + ({\textstyle\sum_{k=1}^{n} b_{k}} + b_{S(n)}) = {\textstyle\sum_{k=1}^{S(n)} a_{k}} + {\textstyle\sum_{k=1}^{S(n)} b_{k}}$
\item $\varphi_{112} :\equiv {\textstyle\sum_{k=1}^{n} a_{k}} + {\textstyle\sum_{k=1}^{n} b_{k}} + a \oplus b(S(n)) = {\textstyle\sum_{k=1}^{n} a_{k}} + {\textstyle\sum_{k=1}^{n} b_{k}} + (a_{S(n)} + b_{S(n)})$
\item $\varphi_{113} :\equiv \forall x\,(\mathrm{Ind}(x) \leftrightarrow \emptyset \in x \land \forall y \in x\,(S(y) \in x)) \to (\mathrm{Ind}(\omega) \leftrightarrow \emptyset \in \omega \land \forall y \in \omega\,(S(y) \in \omega))$
\item $\varphi_{114} :\equiv a_{1_{\omega}} + b_{1_{\omega}} = {\textstyle\sum_{k=1}^{1_{\omega}} a_{k}} + b_{1_{\omega}} \to {\textstyle\sum_{k=1}^{1_{\omega}} a \oplus b(k)} = {\textstyle\sum_{k=1}^{1_{\omega}} a_{k}} + b_{1_{\omega}}$
\item $\varphi_{115} :\equiv \forall z_{45}\,\forall z_{46}\,\forall A_{44}\,\exists B_{43}\,\forall n\,(n \in B_{43} \leftrightarrow n \in A_{44} \land (n \in \mathbb{N} \to {\textstyle\sum_{k=1}^{n} z_{45} \oplus z_{46}(k)} = {\textstyle\sum_{k=1}^{n} z_{45}(k)} + {\textstyle\sum_{k=1}^{n} z_{46}(k)}))$
\item $\varphi_{116} :\equiv \forall n \in \mathbb{N}\,({\textstyle\sum_{k=1}^{S(n)} b_{k}} = {\textstyle\sum_{k=1}^{n} b_{k}} + b_{S(n)}) \to n \in \mathbb{N} \to {\textstyle\sum_{k=1}^{S(n)} b_{k}} = {\textstyle\sum_{k=1}^{n} b_{k}} + b_{S(n)}$
\item $\varphi_{117} :\equiv \forall n \in \mathbb{N}\,({\textstyle\sum_{k=1}^{S(n)} a_{k}} = {\textstyle\sum_{k=1}^{n} a_{k}} + a_{S(n)}) \to n \in \mathbb{N} \to {\textstyle\sum_{k=1}^{S(n)} a_{k}} = {\textstyle\sum_{k=1}^{n} a_{k}} + a_{S(n)}$
\item $\varphi_{118} :\equiv {\textstyle\sum_{k=1}^{n} a_{k}} + a_{S(n)} + ({\textstyle\sum_{k=1}^{n} b_{k}} + b_{S(n)}) = {\textstyle\sum_{k=1}^{n} a_{k}} + a_{S(n)} + ({\textstyle\sum_{k=1}^{n} b_{k}} + b_{S(n)})$
\item $\varphi_{119} :\equiv (\mathrm{Seq}(a) \to \varphi_{40}) \to \mathrm{Seq}(a) \to \mathrm{Seq}(b) \to \mathrm{Seq}(a \oplus b)$
\item $\varphi_{120} :\equiv {\textstyle\sum_{k=1}^{n} a_{k}} + {\textstyle\sum_{k=1}^{n} b_{k}} + (a_{S(n)} + b_{S(n)}) = {\textstyle\sum_{k=1}^{n} a_{k}} + a_{S(n)} + ({\textstyle\sum_{k=1}^{n} b_{k}} + b_{S(n)})$
\item $\varphi_{121} :\equiv \forall z\,({\textstyle\sum_{k=1}^{1_{\omega}} a \oplus b(k)} = {\textstyle\sum_{k=1}^{1_{\omega}} a_{k}} + b_{1_{\omega}} \to {\textstyle\sum_{k=1}^{1_{\omega}} a_{k}} + b_{1_{\omega}} = z \to {\textstyle\sum_{k=1}^{1_{\omega}} a \oplus b(k)} = z)$
\item $\varphi_{122} :\equiv \forall y\,\forall u\,\forall v\,({\textstyle\sum_{k=1}^{n} a_{k}} \in \mathbb{R} \to y \in \mathbb{R} \to u \in \mathbb{R} \to v \in \mathbb{R} \to {\textstyle\sum_{k=1}^{n} a_{k}} + y + (u + v) = {\textstyle\sum_{k=1}^{n} a_{k}} + u + (y + v))$
\item $\varphi_{123} :\equiv {\textstyle\sum_{k=1}^{1_{\omega}} a \oplus b(k)} = {\textstyle\sum_{k=1}^{1_{\omega}} a_{k}} + {\textstyle\sum_{k=1}^{1_{\omega}} b_{k}} \to {\textstyle\sum_{k=1}^{S(n)} a \oplus b(k)} = {\textstyle\sum_{k=1}^{S(n)} a_{k}} + {\textstyle\sum_{k=1}^{S(n)} b_{k}}$
\item $\varphi_{124} :\equiv \forall w\,(1_{\omega} = w \to {\textstyle\sum_{k=1}^{1_{\omega}} a \oplus b(k)} = {\textstyle\sum_{k=1}^{1_{\omega}} a_{k}} + {\textstyle\sum_{k=1}^{1_{\omega}} b_{k}} \to {\textstyle\sum_{k=1}^{w} a \oplus b(k)} = {\textstyle\sum_{k=1}^{w} a_{k}} + {\textstyle\sum_{k=1}^{w} b_{k}})$
\item $\varphi_{125} :\equiv \emptyset \in B_{43} \leftrightarrow \emptyset \in \omega \land (\emptyset \in \mathbb{N} \to {\textstyle\sum_{k=1}^{\emptyset} a \oplus b(k)} = {\textstyle\sum_{k=1}^{\emptyset} a_{k}} + {\textstyle\sum_{k=1}^{\emptyset} b_{k}})$
\item $\varphi_{126} :\equiv \mathrm{Seq}(a) \to \mathrm{Seq}(b) \to \forall n\,(\varphi_{68}) \to n \in B_{43} \to n \in B_{43}$
\item $\varphi_{127} :\equiv {\textstyle\sum_{k=1}^{1_{\omega}} a_{k}} + b_{1_{\omega}} = {\textstyle\sum_{k=1}^{1_{\omega}} a_{k}} + {\textstyle\sum_{k=1}^{1_{\omega}} b_{k}} \to {\textstyle\sum_{k=1}^{1_{\omega}} a \oplus b(k)} = {\textstyle\sum_{k=1}^{1_{\omega}} a_{k}} + {\textstyle\sum_{k=1}^{1_{\omega}} b_{k}}$
\item $\varphi_{128} :\equiv {\textstyle\sum_{k=1}^{1_{\omega}} a_{k}} + b_{1_{\omega}} = {\textstyle\sum_{k=1}^{1_{\omega}} a_{k}} + b_{1_{\omega}} \to {\textstyle\sum_{k=1}^{1_{\omega}} a_{k}} + b_{1_{\omega}} = {\textstyle\sum_{k=1}^{1_{\omega}} a_{k}} + {\textstyle\sum_{k=1}^{1_{\omega}} b_{k}}$
\item $\varphi_{129} :\equiv \forall w\,(b_{1_{\omega}} = w \to {\textstyle\sum_{k=1}^{1_{\omega}} a_{k}} + b_{1_{\omega}} = {\textstyle\sum_{k=1}^{1_{\omega}} a_{k}} + b_{1_{\omega}} \to {\textstyle\sum_{k=1}^{1_{\omega}} a_{k}} + b_{1_{\omega}} = {\textstyle\sum_{k=1}^{1_{\omega}} a_{k}} + w)$
\item $\varphi_{130} :\equiv \mathrm{Seq}(a) \to \mathrm{Seq}(b) \to \forall n\,(\varphi_{68}) \to n \in B_{43} \to S(n) \in B_{43}$
\item $\varphi_{131} :\equiv \mathrm{Seq}(a) \land \mathrm{Seq}(b) \land \forall n\,(\varphi_{68}) \to n \in B_{43} \to S(n) \in B_{43}$
\item $\varphi_{132} :\equiv \mathrm{Seq}(b) \to \forall n\,(\varphi_{68}) \to n \in \omega \to \forall x\,(\mathrm{Ind}(x) \to n \in x)$
\item $\varphi_{133} :\equiv \mathrm{Seq}(b) \to \forall n\,(\varphi_{68}) \to n \in \omega \to \mathrm{Ind}(B_{43}) \to n \in B_{43}$
\item $\varphi_{134} :\equiv \forall y\,({\textstyle\sum_{k=1}^{S(n)} b_{k}} = y \to y = {\textstyle\sum_{k=1}^{S(n)} b_{k}}) \to {\textstyle\sum_{k=1}^{S(n)} b_{k}} = {\textstyle\sum_{k=1}^{n} b_{k}} + b_{S(n)} \to {\textstyle\sum_{k=1}^{n} b_{k}} + b_{S(n)} = {\textstyle\sum_{k=1}^{S(n)} b_{k}}$
\item $\varphi_{135} :\equiv \forall y\,({\textstyle\sum_{k=1}^{S(n)} a_{k}} = y \to y = {\textstyle\sum_{k=1}^{S(n)} a_{k}}) \to {\textstyle\sum_{k=1}^{S(n)} a_{k}} = {\textstyle\sum_{k=1}^{n} a_{k}} + a_{S(n)} \to {\textstyle\sum_{k=1}^{n} a_{k}} + a_{S(n)} = {\textstyle\sum_{k=1}^{S(n)} a_{k}}$
\item $\varphi_{136} :\equiv \mathrm{Seq}(a) \to \mathrm{Seq}(b) \to \forall n\,(\varphi_{68}) \to n \in B_{43} \to S(n) \in \omega$
\item $\varphi_{137} :\equiv \mathrm{Seq}(a) \to \mathrm{Seq}(b) \to \forall n\,(\varphi_{68}) \to \forall n \in B_{43}\,(S(n) \in B_{43})$
\item $\varphi_{138} :\equiv a_{1_{\omega}} = {\textstyle\sum_{k=1}^{1_{\omega}} a_{k}} \to \varphi_{6} \to a_{1_{\omega}} + b_{1_{\omega}} = {\textstyle\sum_{k=1}^{1_{\omega}} a_{k}} + b_{1_{\omega}}$
\item $\varphi_{139} :\equiv \forall n\,(\mathrm{Seq}(a) \land \mathrm{Seq}(b) \land \forall n\,(\varphi_{68})) \to \forall n \in B_{43}\,(S(n) \in B_{43})$
\item $\varphi_{140} :\equiv {\textstyle\sum_{k=1}^{S(n)} a_{k}} + ({\textstyle\sum_{k=1}^{n} b_{k}} + b_{S(n)}) = {\textstyle\sum_{k=1}^{S(n)} a_{k}} + z \to {\textstyle\sum_{k=1}^{S(n)} a_{k}} + ({\textstyle\sum_{k=1}^{n} b_{k}} + b_{S(n)}) = {\textstyle\sum_{k=1}^{S(n)} a_{k}} + w$
\item $\varphi_{141} :\equiv {\textstyle\sum_{k=1}^{n} a \oplus b(k)} + a \oplus b(S(n)) = z + a \oplus b(S(n)) \to {\textstyle\sum_{k=1}^{n} a \oplus b(k)} + a \oplus b(S(n)) = w + a \oplus b(S(n))$
\item $\varphi_{142} :\equiv \forall n\,(\varphi_{68}) \to n \in B_{43} \to S(n) \in \mathbb{N} \to n = \emptyset \to S(n) = 1_{\omega}$
\item $\varphi_{143} :\equiv \forall n\,(\varphi_{68}) \to n \in B_{43} \to S(n) \in \mathbb{N} \to n = \emptyset \to 1_{\omega} = S(n)$
\item $\varphi_{144} :\equiv \mathrm{Seq}(b) \to \forall n\,(\varphi_{68}) \to n \in B_{43} \to \forall y \in \omega\,(S(y) \in \omega)$
\item $\varphi_{145} :\equiv \forall n\,(\varphi_{68}) \to n \in B_{43} \to S(n) \in \mathbb{N} \to n \neq \emptyset \to S(n) \in \mathbb{N}$
\item $\varphi_{146} :\equiv \forall n\,(\varphi_{68}) \to n \in B_{43} \to S(n) \in \mathbb{N} \to n \neq \emptyset \to \mathrm{Seq}(\Sigma b)$
\item $\varphi_{147} :\equiv \forall n\,(\varphi_{68}) \to n \in B_{43} \to S(n) \in \mathbb{N} \to n \neq \emptyset \to \mathrm{Seq}(\Sigma a)$
\item $\varphi_{148} :\equiv (\mathrm{Seq}(a) \to \mathrm{Seq}(b) \to \mathrm{Seq}(b)) \to \mathrm{Seq}(a) \to \mathrm{Seq}(b) \to \varphi_{97}$
\item $\varphi_{149} :\equiv (\mathrm{Seq}(a) \to \mathrm{Seq}(b) \to \mathrm{Seq}(a)) \to \mathrm{Seq}(a) \to \mathrm{Seq}(b) \to \varphi_{98}$
\item $\varphi_{150} :\equiv \mathrm{Seq}(b) \to \forall n\,(\varphi_{68}) \to n \in B_{43} \to S(n) \in \mathbb{N} \to n \neq \emptyset \to \mathrm{Seq}(b)$
\item $\varphi_{151} :\equiv \mathrm{Seq}(b) \to \forall n\,(\varphi_{68}) \to n \in B_{43} \to S(n) \in \mathbb{N} \to n \neq \emptyset \to \mathrm{Seq}(a)$
\item $\varphi_{152} :\equiv \forall n\,(\varphi_{68}) \to n \in B_{43} \to S(n) \in \mathbb{N} \to n \neq \emptyset \to \mathrm{Seq}(a \oplus b)$
\item $\varphi_{153} :\equiv \forall n\,(\varphi_{68}) \to n \in B_{43} \to S(n) \in \mathbb{N} \to n = \emptyset \to S(n) = S(\emptyset)$
\item $\varphi_{154} :\equiv {\textstyle\sum_{k=1}^{S(n)} a_{k}} + ({\textstyle\sum_{k=1}^{n} b_{k}} + b_{S(n)}) = {\textstyle\sum_{k=1}^{S(n)} a_{k}} + {\textstyle\sum_{k=1}^{S(n)} b_{k}} \to {\textstyle\sum_{k=1}^{S(n)} a \oplus b(k)} = {\textstyle\sum_{k=1}^{S(n)} a_{k}} + {\textstyle\sum_{k=1}^{S(n)} b_{k}}$
\item $\varphi_{155} :\equiv \mathrm{Seq}(a) \to \mathrm{Seq}(b) \to \forall n\,(\varphi_{68}) \to n \in B_{43} \to n \in \omega \to S(n) \in \omega$
\item $\varphi_{156} :\equiv \mathrm{Seq}(b) \to \forall n\,(\varphi_{68}) \to n \in B_{43} \to S(n) \in \mathbb{N} \to n \neq \emptyset \to n \in \mathbb{N}$
\item $\varphi_{157} :\equiv \mathrm{Seq}(b) \to \forall n\,(\varphi_{68}) \to n \in B_{43} \to S(n) \in \mathbb{N} \to n = \emptyset \to n = \emptyset$
\item $\varphi_{158} :\equiv \forall z\,({\textstyle\sum_{k=1}^{S(n)} a \oplus b(k)} = {\textstyle\sum_{k=1}^{n} a \oplus b(k)} + a \oplus b(S(n)) \to {\textstyle\sum_{k=1}^{n} a \oplus b(k)} + a \oplus b(S(n)) = z \to {\textstyle\sum_{k=1}^{S(n)} a \oplus b(k)} = z)$
\item $\varphi_{159} :\equiv \mathrm{Seq}(b) \to \forall n\,(\varphi_{68}) \to n \in B_{43} \to S(n) \in \mathbb{N} \to n \neq \emptyset \to n \neq \emptyset$
\item $\varphi_{160} :\equiv \forall n\,(\varphi_{68}) \to n \in B_{43} \to S(n) \in \mathbb{N} \to n \neq \emptyset \to b_{S(n)} \in \mathbb{R}$
\item $\varphi_{161} :\equiv \forall n\,(\varphi_{68}) \to n \in B_{43} \to S(n) \in \mathbb{N} \to n \neq \emptyset \to a_{S(n)} \in \mathbb{R}$
\item $\varphi_{162} :\equiv (\mathrm{Seq}(a) \to \mathrm{Seq}(b) \to \varphi_{24}) \to 1_{\omega} \in \mathbb{N} \to \mathrm{Seq}(a) \to \varphi_{12}$
\item $\varphi_{163} :\equiv \mathrm{Seq}(\Sigma a \oplus b) \land {\textstyle\sum_{k=1}^{1_{\omega}} a \oplus b(k)} = a \oplus b(1_{\omega}) \land \varphi_{61}$
\item $\varphi_{164} :\equiv \forall n\,(\varphi_{68}) \to n \in B_{43} \to S(n) \in \mathbb{N} \to n \neq \emptyset \to {\textstyle\sum_{k=1}^{n} b_{k}} \in \mathbb{R}$
\item $\varphi_{165} :\equiv \forall n\,(\varphi_{68}) \to n \in B_{43} \to S(n) \in \mathbb{N} \to n \neq \emptyset \to {\textstyle\sum_{k=1}^{n} a_{k}} \in \mathbb{R}$
\item $\varphi_{166} :\equiv \forall z\,({\textstyle\sum_{k=1}^{S(n)} a \oplus b(k)} = {\textstyle\sum_{k=1}^{S(n)} a_{k}} + ({\textstyle\sum_{k=1}^{n} b_{k}} + b_{S(n)}) \to {\textstyle\sum_{k=1}^{S(n)} a_{k}} + ({\textstyle\sum_{k=1}^{n} b_{k}} + b_{S(n)}) = z \to {\textstyle\sum_{k=1}^{S(n)} a \oplus b(k)} = z)$
\item $\varphi_{167} :\equiv \forall u\,\forall v\,({\textstyle\sum_{k=1}^{n} a_{k}} \in \mathbb{R} \to {\textstyle\sum_{k=1}^{n} b_{k}} \in \mathbb{R} \to u \in \mathbb{R} \to v \in \mathbb{R} \to {\textstyle\sum_{k=1}^{n} a_{k}} + {\textstyle\sum_{k=1}^{n} b_{k}} + (u + v) = {\textstyle\sum_{k=1}^{n} a_{k}} + u + ({\textstyle\sum_{k=1}^{n} b_{k}} + v))$
\item $\varphi_{168} :\equiv \forall n \in \omega\,(n \in \mathbb{N} \to {\textstyle\sum_{k=1}^{n} a \oplus b(k)} = {\textstyle\sum_{k=1}^{n} a_{k}} + {\textstyle\sum_{k=1}^{n} b_{k}}) \to n \in \omega \to n \in \mathbb{N} \to {\textstyle\sum_{k=1}^{n} a \oplus b(k)} = {\textstyle\sum_{k=1}^{n} a_{k}} + {\textstyle\sum_{k=1}^{n} b_{k}}$
\item $\varphi_{169} :\equiv \mathrm{Seq}(a) \to \mathrm{Seq}(b) \to \forall n\,(\varphi_{68}) \to \emptyset \in \omega \land \forall y \in \omega\,(S(y) \in \omega)$
\item $\varphi_{170} :\equiv (\mathrm{Seq}(a) \to \mathrm{Seq}(b) \to n \in \mathbb{N} \to n \in \omega) \to (\mathrm{Seq}(a) \to \mathrm{Seq}(b) \to n \in \mathbb{N} \to n \in \mathbb{N}) \to \varphi_{46}$
\item $\varphi_{171} :\equiv (\mathrm{Seq}(a) \to \mathrm{Seq}(b) \to \varphi_{97}) \to \mathrm{Seq}(a) \to \mathrm{Seq}(b) \to {\textstyle\sum_{k=1}^{1_{\omega}} b_{k}} = b_{1_{\omega}}$
\item $\varphi_{172} :\equiv (\mathrm{Seq}(a) \to \mathrm{Seq}(b) \to \varphi_{98}) \to \mathrm{Seq}(a) \to \mathrm{Seq}(b) \to {\textstyle\sum_{k=1}^{1_{\omega}} a_{k}} = a_{1_{\omega}}$
\item $\varphi_{173} :\equiv \forall z\,({\textstyle\sum_{k=1}^{S(n)} a \oplus b(k)} = {\textstyle\sum_{k=1}^{n} a_{k}} + {\textstyle\sum_{k=1}^{n} b_{k}} + a \oplus b(S(n)) \to {\textstyle\sum_{k=1}^{n} a_{k}} + {\textstyle\sum_{k=1}^{n} b_{k}} + a \oplus b(S(n)) = z \to {\textstyle\sum_{k=1}^{S(n)} a \oplus b(k)} = z)$
\item $\varphi_{174} :\equiv \forall n\,(\varphi_{68}) \to n \in B_{43} \to n \in \mathbb{N} \to {\textstyle\sum_{k=1}^{n} a \oplus b(k)} = {\textstyle\sum_{k=1}^{n} a_{k}} + {\textstyle\sum_{k=1}^{n} b_{k}}$
\item $\varphi_{175} :\equiv \forall n\,(\varphi_{68}) \to n \in \omega \to n \in \mathbb{N} \to {\textstyle\sum_{k=1}^{n} a \oplus b(k)} = {\textstyle\sum_{k=1}^{n} a_{k}} + {\textstyle\sum_{k=1}^{n} b_{k}}$
\item $\varphi_{176} :\equiv \forall n\,(\varphi_{68}) \to \forall n \in \omega\,(n \in \mathbb{N} \to {\textstyle\sum_{k=1}^{n} a \oplus b(k)} = {\textstyle\sum_{k=1}^{n} a_{k}} + {\textstyle\sum_{k=1}^{n} b_{k}})$
\item $\varphi_{177} :\equiv \forall z\,({\textstyle\sum_{k=1}^{S(n)} a \oplus b(k)} = {\textstyle\sum_{k=1}^{n} a_{k}} + a_{S(n)} + ({\textstyle\sum_{k=1}^{n} b_{k}} + b_{S(n)}) \to {\textstyle\sum_{k=1}^{n} a_{k}} + a_{S(n)} + ({\textstyle\sum_{k=1}^{n} b_{k}} + b_{S(n)}) = z \to {\textstyle\sum_{k=1}^{S(n)} a \oplus b(k)} = z)$
\item $\varphi_{178} :\equiv \forall z\,({\textstyle\sum_{k=1}^{S(n)} a \oplus b(k)} = {\textstyle\sum_{k=1}^{n} a_{k}} + {\textstyle\sum_{k=1}^{n} b_{k}} + (a_{S(n)} + b_{S(n)}) \to {\textstyle\sum_{k=1}^{n} a_{k}} + {\textstyle\sum_{k=1}^{n} b_{k}} + (a_{S(n)} + b_{S(n)}) = z \to {\textstyle\sum_{k=1}^{S(n)} a \oplus b(k)} = z)$
\item $\varphi_{179} :\equiv (S(n) = 1_{\omega} \to 1_{\omega} = S(n)) \to \mathrm{Seq}(a) \to \mathrm{Seq}(b) \to \varphi_{143}$
\item $\varphi_{180} :\equiv \forall v\,({\textstyle\sum_{k=1}^{n} a_{k}} \in \mathbb{R} \to {\textstyle\sum_{k=1}^{n} b_{k}} \in \mathbb{R} \to a_{S(n)} \in \mathbb{R} \to v \in \mathbb{R} \to {\textstyle\sum_{k=1}^{n} a_{k}} + {\textstyle\sum_{k=1}^{n} b_{k}} + (a_{S(n)} + v) = {\textstyle\sum_{k=1}^{n} a_{k}} + a_{S(n)} + ({\textstyle\sum_{k=1}^{n} b_{k}} + v))$
\item $\varphi_{181} :\equiv (\forall x\,(\mathrm{Ind}(x) \to n \in x) \to \mathrm{Ind}(B_{43}) \to n \in B_{43}) \to \mathrm{Seq}(a) \to \varphi_{133}$
\item $\varphi_{182} :\equiv \mathrm{Seq}(a) \land \mathrm{Seq}(b) \land \forall n\,(\varphi_{68}) \to n \in \omega \to n \in \mathbb{N} \to {\textstyle\sum_{k=1}^{n} a \oplus b(k)} = {\textstyle\sum_{k=1}^{n} a_{k}} + {\textstyle\sum_{k=1}^{n} b_{k}}$
\item $\varphi_{183} :\equiv (\forall n \in \mathbb{N}\,(a \oplus b(n) = a_{n} + b_{n}) \to \varphi_{24}) \to \mathrm{Seq}(a) \to \mathrm{Seq}(b) \to \varphi_{24}$
\item $\varphi_{184} :\equiv \forall w\,({\textstyle\sum_{k=1}^{n} a \oplus b(k)} = w \to \varphi_{73} \to {\textstyle\sum_{k=1}^{n} a \oplus b(k)} + a \oplus b(S(n)) = w + a \oplus b(S(n)))$
\item $\varphi_{185} :\equiv \forall n\,(\mathrm{Seq}(a) \land \mathrm{Seq}(b) \land \forall n\,(\varphi_{68})) \to \forall n \in \omega\,(n \in \mathbb{N} \to {\textstyle\sum_{k=1}^{n} a \oplus b(k)} = {\textstyle\sum_{k=1}^{n} a_{k}} + {\textstyle\sum_{k=1}^{n} b_{k}})$
\item $\varphi_{186} :\equiv \varphi_{112} \to {\textstyle\sum_{k=1}^{S(n)} a \oplus b(k)} = {\textstyle\sum_{k=1}^{n} a_{k}} + {\textstyle\sum_{k=1}^{n} b_{k}} + (a_{S(n)} + b_{S(n)})$
\item $\varphi_{187} :\equiv \forall w\,({\textstyle\sum_{k=1}^{n} b_{k}} + b_{S(n)} = w \to \varphi_{81} \to {\textstyle\sum_{k=1}^{S(n)} a_{k}} + ({\textstyle\sum_{k=1}^{n} b_{k}} + b_{S(n)}) = {\textstyle\sum_{k=1}^{S(n)} a_{k}} + w)$
\item $\varphi_{188} :\equiv (\mathrm{Seq}(a) \to \mathrm{Seq}(b) \to \mathrm{Seq}(a \oplus b)) \to \mathrm{Seq}(a) \to \mathrm{Seq}(b) \to \varphi_{163}$
\item $\varphi_{189} :\equiv \varphi_{120} \to {\textstyle\sum_{k=1}^{S(n)} a \oplus b(k)} = {\textstyle\sum_{k=1}^{n} a_{k}} + a_{S(n)} + ({\textstyle\sum_{k=1}^{n} b_{k}} + b_{S(n)})$
\item $\varphi_{190} :\equiv (n = \emptyset \to S(n) = S(n) \to S(n) = S(\emptyset)) \to \mathrm{Seq}(a) \to \mathrm{Seq}(b) \to \varphi_{153}$
\item $\varphi_{191} :\equiv \mathrm{Seq}(a) \to \mathrm{Seq}(b) \to \forall n\,(\varphi_{68}) \to n \in B_{43} \to S(n) \in \mathbb{N} \to \varphi_{36}$
\item $\varphi_{192} :\equiv {\textstyle\sum_{k=1}^{n} a_{k}} \in \mathbb{R} \to {\textstyle\sum_{k=1}^{n} b_{k}} \in \mathbb{R} \to a_{S(n)} \in \mathbb{R} \to b_{S(n)} \in \mathbb{R} \to \varphi_{120}$
\item $\varphi_{193} :\equiv \exists B_{43}\,\forall n\,(\varphi_{68}) \to \mathrm{Seq}(a) \to \mathrm{Seq}(b) \to \exists B_{43}\,\forall n\,(\varphi_{68})$
\item $\varphi_{194} :\equiv ({\textstyle\sum_{k=1}^{1_{\omega}} a \oplus b(k)} = a \oplus b(1_{\omega}) \to \varphi_{85}) \to \mathrm{Seq}(a) \to \varphi_{25}$
\item $\varphi_{195} :\equiv \mathrm{Seq}(a) \to \mathrm{Seq}(b) \to \forall n\,(\varphi_{68}) \to n \in B_{43} \to S(n) \in \mathbb{N} \to \varphi_{45}$
\item $\varphi_{196} :\equiv \mathrm{Seq}(a) \to \mathrm{Seq}(b) \to \forall n\,(\varphi_{68}) \to n \in B_{43} \to S(n) \in \mathbb{N} \to \varphi_{59}$
\item $\varphi_{197} :\equiv \mathrm{Seq}(a) \land \mathrm{Seq}(b) \land \forall n\,(\varphi_{68}) \to \forall n\,(\mathrm{Seq}(a) \land \mathrm{Seq}(b) \land \forall n\,(\varphi_{68}))$
\item $\varphi_{198} :\equiv \mathrm{Seq}(b) \to \forall n\,(\varphi_{68}) \to n \in B_{43} \to S(n) \in \mathbb{N} \to n \neq \emptyset \to \varphi_{48}$
\item $\varphi_{199} :\equiv \mathrm{Seq}(b) \to \forall n\,(\varphi_{68}) \to n \in B_{43} \to S(n) \in \mathbb{N} \to n \neq \emptyset \to \varphi_{61}$
\item $\varphi_{200} :\equiv \mathrm{Seq}(a) \to \mathrm{Seq}(b) \to \forall n\,(\varphi_{68}) \to n \in B_{43} \to S(n) \in \mathbb{N} \to \varphi_{78}$
\item $\varphi_{201} :\equiv {\textstyle\sum_{k=1}^{n} b_{k}} + b_{S(n)} = {\textstyle\sum_{k=1}^{S(n)} b_{k}} \to \varphi_{81} \to {\textstyle\sum_{k=1}^{S(n)} a_{k}} + ({\textstyle\sum_{k=1}^{n} b_{k}} + b_{S(n)}) = {\textstyle\sum_{k=1}^{S(n)} a_{k}} + {\textstyle\sum_{k=1}^{S(n)} b_{k}}$
\item $\varphi_{202} :\equiv \varphi_{14} \to \varphi_{13} \to \mathrm{Seq}(a) \to \mathrm{Seq}(b) \to \forall n\,(\varphi_{68}) \to \mathrm{Ind}(B_{43})$
\item $\varphi_{203} :\equiv \mathrm{Seq}(a) \to \mathrm{Seq}(b) \to \forall n\,(\varphi_{68}) \to n \in B_{43} \to S(n) \in \mathbb{N} \to \varphi_{83}$
\item $\varphi_{204} :\equiv \mathrm{Seq}(a) \to \mathrm{Seq}(b) \to \forall n\,(\varphi_{68}) \to n \in B_{43} \to S(n) \in \mathbb{N} \to \varphi_{84}$
\item $\varphi_{205} :\equiv \mathrm{Ind}(\omega) \to (\mathrm{Ind}(\omega) \leftrightarrow \emptyset \in \omega \land \forall y \in \omega\,(S(y) \in \omega)) \to \varphi_{169}$
\item $\varphi_{206} :\equiv \mathrm{Seq}(b) \to \forall n\,(\varphi_{68}) \to n \in B_{43} \to S(n) \in \mathbb{N} \to n \neq \emptyset \to \varphi_{87}$
\item $\varphi_{207} :\equiv \mathrm{Seq}(b) \to \forall n\,(\varphi_{68}) \to n \in B_{43} \to S(n) \in \mathbb{N} \to n \neq \emptyset \to \varphi_{97}$
\item $\varphi_{208} :\equiv \mathrm{Seq}(b) \to \forall n\,(\varphi_{68}) \to n \in B_{43} \to S(n) \in \mathbb{N} \to n \neq \emptyset \to \varphi_{98}$
\item $\varphi_{209} :\equiv \mathrm{Seq}(b) \to \forall n\,(\varphi_{68}) \to n \in B_{43} \to S(n) \in \mathbb{N} \to n \neq \emptyset \to \varphi_{103}$
\item $\varphi_{210} :\equiv {\textstyle\sum_{k=1}^{S(n)} a \oplus b(k)} = {\textstyle\sum_{k=1}^{n} a \oplus b(k)} + a \oplus b(S(n)) \to \varphi_{87} \to {\textstyle\sum_{k=1}^{S(n)} a \oplus b(k)} = {\textstyle\sum_{k=1}^{n} a_{k}} + {\textstyle\sum_{k=1}^{n} b_{k}} + a \oplus b(S(n))$
\item $\varphi_{211} :\equiv \mathrm{Seq}(b) \to \forall n\,(\varphi_{68}) \to n \in B_{43} \to S(n) \in \mathbb{N} \to n \neq \emptyset \to \varphi_{112}$
\item $\varphi_{212} :\equiv \mathrm{Seq}(b) \to \forall n\,(\varphi_{68}) \to n \in B_{43} \to S(n) \in \mathbb{N} \to n \neq \emptyset \to \varphi_{120}$
\item $\varphi_{213} :\equiv {\textstyle\sum_{k=1}^{S(n)} a \oplus b(k)} = {\textstyle\sum_{k=1}^{n} a_{k}} + a_{S(n)} + ({\textstyle\sum_{k=1}^{n} b_{k}} + b_{S(n)}) \to \varphi_{103} \to {\textstyle\sum_{k=1}^{S(n)} a \oplus b(k)} = {\textstyle\sum_{k=1}^{S(n)} a_{k}} + ({\textstyle\sum_{k=1}^{n} b_{k}} + b_{S(n)})$
\item $\varphi_{214} :\equiv (\mathrm{Seq}(a) \to \mathrm{Seq}(b) \to \forall n\,(\varphi_{68}) \to n \in B_{43} \to n \in \omega) \to \varphi_{136}$
\item $\varphi_{215} :\equiv ({\textstyle\sum_{k=1}^{S(n)} b_{k}} = {\textstyle\sum_{k=1}^{n} b_{k}} + b_{S(n)} \to {\textstyle\sum_{k=1}^{n} b_{k}} + b_{S(n)} = {\textstyle\sum_{k=1}^{S(n)} b_{k}}) \to \mathrm{Seq}(a) \to \mathrm{Seq}(b) \to \forall n\,(\varphi_{68}) \to \varphi_{55}$
\item $\varphi_{216} :\equiv ({\textstyle\sum_{k=1}^{S(n)} a_{k}} = {\textstyle\sum_{k=1}^{n} a_{k}} + a_{S(n)} \to {\textstyle\sum_{k=1}^{n} a_{k}} + a_{S(n)} = {\textstyle\sum_{k=1}^{S(n)} a_{k}}) \to \mathrm{Seq}(a) \to \mathrm{Seq}(b) \to \forall n\,(\varphi_{68}) \to \varphi_{58}$
\item $\varphi_{217} :\equiv {\textstyle\sum_{k=1}^{S(n)} a \oplus b(k)} = {\textstyle\sum_{k=1}^{n} a_{k}} + {\textstyle\sum_{k=1}^{n} b_{k}} + (a_{S(n)} + b_{S(n)}) \to \varphi_{189}$
\item $\varphi_{218} :\equiv \mathrm{Seq}(b) \to \forall n\,(\varphi_{68}) \to n \in B_{43} \to S(n) \in \mathbb{N} \to n \neq \emptyset \to \varphi_{163}$
\item $\varphi_{219} :\equiv {\textstyle\sum_{k=1}^{1_{\omega}} a_{k}} + b_{1_{\omega}} = {\textstyle\sum_{k=1}^{1_{\omega}} a_{k}} + b_{1_{\omega}} \to (b_{1_{\omega}} = {\textstyle\sum_{k=1}^{1_{\omega}} b_{k}} \to \varphi_{128}) \to \varphi_{60}$
\item $\varphi_{220} :\equiv (\mathrm{Seq}(a) \to \mathrm{Seq}(b) \to a_{1_{\omega}} = {\textstyle\sum_{k=1}^{1_{\omega}} a_{k}}) \to \varphi_{6} \to \varphi_{138} \to \mathrm{Seq}(a) \to \varphi_{28}$
\item $\varphi_{221} :\equiv (\mathrm{Seq}(a) \to \mathrm{Seq}(b) \to \forall n\,(\varphi_{68}) \to n \in \omega \to n \in \omega) \to (n \in \omega \leftrightarrow \forall x\,(\mathrm{Ind}(x) \to n \in x)) \to \mathrm{Seq}(a) \to \varphi_{132}$
\item $\varphi_{222} :\equiv (\mathrm{Seq}(a) \to \varphi_{28}) \to ({\textstyle\sum_{k=1}^{1_{\omega}} a \oplus b(k)} = a_{1_{\omega}} + b_{1_{\omega}} \to \varphi_{114}) \to \varphi_{38}$
\item $\varphi_{223} :\equiv (\mathrm{Seq}(a) \to \mathrm{Seq}(b) \to \varphi_{145}) \to \mathrm{Seq}(a) \to \mathrm{Seq}(b) \to \varphi_{160}$
\item $\varphi_{224} :\equiv (\mathrm{Seq}(a) \to \mathrm{Seq}(b) \to \varphi_{145}) \to \mathrm{Seq}(a) \to \mathrm{Seq}(b) \to \varphi_{161}$
\item $\varphi_{225} :\equiv (\mathrm{Seq}(a) \to \mathrm{Seq}(b) \to \forall n\,(\varphi_{68}) \to \varphi_{69}) \to \mathrm{Seq}(a) \to \mathrm{Seq}(b) \to \forall n\,(\varphi_{68}) \to n \in B_{43} \to n \in \omega$
\item $\varphi_{226} :\equiv (n \in \mathbb{N} \to S(n) \in \mathbb{N}) \to (\mathrm{Seq}(a) \to \varphi_{156}) \to \mathrm{Seq}(a) \to \mathrm{Seq}(b) \to \varphi_{145}$
\item $\varphi_{227} :\equiv (1_{\omega} = S(n) \to \varphi_{123}) \to \mathrm{Seq}(a) \to \mathrm{Seq}(b) \to \forall n\,(\varphi_{68}) \to \varphi_{89}$
\item $\varphi_{228} :\equiv (\mathrm{Seq}(a) \to \varphi_{156}) \to \mathrm{Seq}(a) \to \mathrm{Seq}(b) \to \forall n\,(\varphi_{68}) \to \varphi_{54}$
\item $\varphi_{229} :\equiv (\mathrm{Seq}(a) \to \varphi_{156}) \to \mathrm{Seq}(a) \to \mathrm{Seq}(b) \to \forall n\,(\varphi_{68}) \to \varphi_{57}$
\item $\varphi_{230} :\equiv (\mathrm{Seq}(a) \to \mathrm{Seq}(b) \to \varphi_{145}) \to \mathrm{Seq}(a) \to \mathrm{Seq}(b) \to \forall n\,(\varphi_{68}) \to n \in B_{43} \to \varphi_{37}$
\item $\varphi_{231} :\equiv (\mathrm{Seq}(a) \to \varphi_{156}) \to \mathrm{Seq}(a) \to \mathrm{Seq}(b) \to \forall n\,(\varphi_{68}) \to \varphi_{77}$
\item $\varphi_{232} :\equiv (\mathrm{Seq}(a) \to \varphi_{144}) \to (\forall y \in \omega\,(S(y) \in \omega) \to n \in \omega \to S(n) \in \omega) \to \varphi_{155}$
\item $\varphi_{233} :\equiv (\forall n\,(\varphi_{68}) \to \varphi_{68}) \to \mathrm{Seq}(a) \to \mathrm{Seq}(b) \to \forall n\,(\varphi_{68}) \to \varphi_{68}$
\item $\varphi_{234} :\equiv (\mathrm{Seq}(a) \to \mathrm{Seq}(b) \to \forall n\,(\varphi_{68}) \to \varphi_{69}) \to \mathrm{Seq}(a) \to \mathrm{Seq}(b) \to \varphi_{174}$
\item $\varphi_{235} :\equiv (\mathrm{Seq}(a) \to \mathrm{Seq}(b) \to \varphi_{153}) \to S(\emptyset) = 1_{\omega} \to \varphi_{5} \to \mathrm{Seq}(a) \to \mathrm{Seq}(b) \to \varphi_{142}$
\item $\varphi_{236} :\equiv (\mathrm{Seq}(a) \to \mathrm{Seq}(b) \to \forall n\,(\varphi_{68}) \to \varphi_{68}) \to \mathrm{Seq}(a) \to \mathrm{Seq}(b) \to \varphi_{175}$
\item $\varphi_{237} :\equiv \forall B_{43}\,(\varphi_{86} \to \neg \forall n\,(\varphi_{68})) \to \forall B_{43}\,(\varphi_{86}) \to \forall B_{43}\,(\neg \forall n\,(\varphi_{68}))$
\item $\varphi_{238} :\equiv (\mathrm{Seq}(a) \to \mathrm{Seq}(b) \to \forall n\,(\varphi_{68}) \to \varphi_{94}) \to \mathrm{Seq}(a) \to \mathrm{Seq}(b) \to \forall n\,(\varphi_{68}) \to \varphi_{67}$
\item $\varphi_{239} :\equiv \forall z\,\forall w\,(z = w \to \varphi_{71} \to \varphi_{70}) \to \forall w\,(a \oplus b(S(n)) = w \to \varphi_{102} \to \varphi_{70})$
\item $\varphi_{240} :\equiv (\forall n\,(\varphi_{68}) \to \varphi_{101}) \to \mathrm{Seq}(a) \to \mathrm{Seq}(b) \to \forall n\,(\varphi_{68}) \to \varphi_{101}$
\item $\varphi_{241} :\equiv \varphi_{38} \to \varphi_{60} \to ({\textstyle\sum_{k=1}^{1_{\omega}} a \oplus b(k)} = {\textstyle\sum_{k=1}^{1_{\omega}} a_{k}} + b_{1_{\omega}} \to \varphi_{127}) \to \varphi_{44}$
\item $\varphi_{242} :\equiv ({\textstyle\sum_{k=1}^{S(n)} a \oplus b(k)} = {\textstyle\sum_{k=1}^{S(n)} a_{k}} + ({\textstyle\sum_{k=1}^{n} b_{k}} + b_{S(n)}) \to \varphi_{154}) \to \mathrm{Seq}(a) \to \mathrm{Seq}(b) \to \forall n\,(\varphi_{68}) \to \varphi_{94}$
\item $\varphi_{243} :\equiv \forall z\,\forall w\,(z = w \to \varphi_{76} \to \varphi_{75}) \to \forall w\,({\textstyle\sum_{k=1}^{n} a_{k}} + a_{S(n)} = w \to \varphi_{118} \to \varphi_{75})$
\item $\varphi_{244} :\equiv \varphi_{173} \to {\textstyle\sum_{k=1}^{S(n)} a \oplus b(k)} = {\textstyle\sum_{k=1}^{n} a_{k}} + {\textstyle\sum_{k=1}^{n} b_{k}} + a \oplus b(S(n)) \to \varphi_{186}$
\item $\varphi_{245} :\equiv (\mathrm{Seq}(a) \to \mathrm{Seq}(b) \to \forall n\,(\varphi_{68}) \to n \in B_{43} \to n \in \omega) \to (\mathrm{Seq}(a) \to \varphi_{159}) \to \mathrm{Seq}(a) \to \varphi_{156}$
\item $\varphi_{246} :\equiv (\mathrm{Seq}(a) \to \varphi_{207}) \to \mathrm{Seq}(a) \to \mathrm{Seq}(b) \to \forall n\,(\varphi_{68}) \to n \in B_{43} \to \varphi_{63}$
\item $\varphi_{247} :\equiv (\mathrm{Seq}(a) \to \varphi_{208}) \to \mathrm{Seq}(a) \to \mathrm{Seq}(b) \to \forall n\,(\varphi_{68}) \to n \in B_{43} \to \varphi_{64}$
\item $\varphi_{248} :\equiv (\forall n\,(\varphi_{68}) \to \varphi_{125}) \to \mathrm{Seq}(a) \to \mathrm{Seq}(b) \to \forall n\,(\varphi_{68}) \to \varphi_{125}$
\item $\varphi_{249} :\equiv \forall w\,(a \oplus b(S(n)) = w \to \varphi_{102} \to \varphi_{70}) \to a \oplus b(S(n)) = a_{S(n)} + b_{S(n)} \to \varphi_{102} \to \varphi_{112}$
\item $\varphi_{250} :\equiv (\mathrm{Seq}(a) \to \mathrm{Seq}(b) \to \varphi_{146}) \to (\mathrm{Seq}(a) \to \varphi_{156}) \to \mathrm{Seq}(a) \to \mathrm{Seq}(b) \to \varphi_{164}$
\item $\varphi_{251} :\equiv (\mathrm{Seq}(a) \to \mathrm{Seq}(b) \to \varphi_{147}) \to (\mathrm{Seq}(a) \to \varphi_{156}) \to \mathrm{Seq}(a) \to \mathrm{Seq}(b) \to \varphi_{165}$
\item $\varphi_{252} :\equiv ({\textstyle\sum_{k=1}^{n} a \oplus b(k)} = {\textstyle\sum_{k=1}^{n} a_{k}} + {\textstyle\sum_{k=1}^{n} b_{k}} \to \varphi_{73} \to \varphi_{87}) \to \mathrm{Seq}(a) \to \varphi_{206}$
\item $\varphi_{253} :\equiv (\mathrm{Seq}(a) \to \mathrm{Seq}(b) \to \forall n\,(\varphi_{68}) \to \emptyset \in B_{43}) \to \varphi_{137} \to \varphi_{202}$
\item $\varphi_{254} :\equiv (\mathrm{Seq}(a) \to \mathrm{Seq}(b) \to \forall n\,(\varphi_{68}) \to \mathrm{Ind}(B_{43})) \to \varphi_{236}$
\item $\varphi_{255} :\equiv ({\textstyle\sum_{k=1}^{S(n)} a \oplus b(k)} = {\textstyle\sum_{k=1}^{n} a_{k}} + {\textstyle\sum_{k=1}^{n} b_{k}} + a \oplus b(S(n)) \to \varphi_{186}) \to \varphi_{203}$
\item $\varphi_{256} :\equiv \forall w\,({\textstyle\sum_{k=1}^{n} a_{k}} + a_{S(n)} = w \to \varphi_{118} \to \varphi_{75}) \to {\textstyle\sum_{k=1}^{n} a_{k}} + a_{S(n)} = {\textstyle\sum_{k=1}^{S(n)} a_{k}} \to \varphi_{118} \to \varphi_{103}$
\item $\varphi_{257} :\equiv (\mathrm{Seq}(a) \to \mathrm{Seq}(b) \to \forall n\,(\varphi_{68}) \to n \in B_{43} \to \varphi_{43}) \to \varphi_{108} \to \varphi_{191}$
\item $\varphi_{258} :\equiv (a \oplus b(S(n)) = a_{S(n)} + b_{S(n)} \to \varphi_{102} \to \varphi_{112}) \to \mathrm{Seq}(a) \to \varphi_{211}$
\item $\varphi_{259} :\equiv (\mathrm{Seq}(b) \to S(n) \in \mathbb{N} \to b_{S(n)} \in \mathbb{R}) \to (\mathrm{Seq}(a) \to \varphi_{150}) \to \varphi_{223}$
\item $\varphi_{260} :\equiv (\mathrm{Seq}(a) \to S(n) \in \mathbb{N} \to a_{S(n)} \in \mathbb{R}) \to (\mathrm{Seq}(a) \to \varphi_{151}) \to \varphi_{224}$
\item $\varphi_{261} :\equiv (\mathrm{Seq}(a) \to \mathrm{Seq}(b) \to \forall n\,(\varphi_{68}) \to \varphi_{68}) \to \varphi_{126} \to \mathrm{Seq}(a) \to \mathrm{Seq}(b) \to \forall n\,(\varphi_{68}) \to \varphi_{69}$
\item $\varphi_{262} :\equiv (\mathrm{Seq}(a) \to \mathrm{Seq}(b) \to \forall n\,(\varphi_{68}) \to n \in B_{43} \to \varphi_{63}) \to \varphi_{116} \to \mathrm{Seq}(a) \to \mathrm{Seq}(b) \to \forall n\,(\varphi_{68}) \to n \in B_{43} \to \varphi_{53}$
\item $\varphi_{263} :\equiv (\mathrm{Seq}(a) \to \mathrm{Seq}(b) \to \forall n\,(\varphi_{68}) \to n \in B_{43} \to \varphi_{64}) \to \varphi_{117} \to \mathrm{Seq}(a) \to \mathrm{Seq}(b) \to \forall n\,(\varphi_{68}) \to n \in B_{43} \to \varphi_{56}$
\item $\varphi_{264} :\equiv (\mathrm{Seq}(a) \to \mathrm{Seq}(b) \to \forall n\,(\varphi_{68}) \to n \in B_{43} \to \varphi_{53}) \to \varphi_{228}$
\item $\varphi_{265} :\equiv \varphi_{81} \to \varphi_{201} \to \mathrm{Seq}(a) \to \mathrm{Seq}(b) \to \forall n\,(\varphi_{68}) \to n \in B_{43} \to \varphi_{111}$
\item $\varphi_{266} :\equiv (\mathrm{Seq}(a) \to \mathrm{Seq}(b) \to \forall n\,(\varphi_{68}) \to n \in B_{43} \to \varphi_{56}) \to \varphi_{229}$
\item $\varphi_{267} :\equiv (\mathrm{Seq}(a) \to \mathrm{Seq}(b) \to \varphi_{161}) \to (\mathrm{Seq}(a) \to \mathrm{Seq}(b) \to \varphi_{160}) \to \mathrm{Seq}(a) \to \varphi_{212}$
\item $\varphi_{268} :\equiv (\mathrm{Seq}(a) \to \varphi_{40}) \to (\mathrm{Seq}(a) \to \varphi_{151}) \to (\mathrm{Seq}(a) \to \varphi_{150}) \to \varphi_{195}$
\item $\varphi_{269} :\equiv (\mathrm{Seq}(a) \to \mathrm{Seq}(b) \to \forall n\,(\varphi_{68}) \to \varphi_{125}) \to \varphi_{169} \to \varphi_{105} \to \mathrm{Seq}(a) \to \mathrm{Seq}(b) \to \forall n\,(\varphi_{68}) \to \emptyset \in B_{43}$
\item $\varphi_{270} :\equiv (\mathrm{Seq}(a \oplus b) \to \varphi_{163}) \to (\mathrm{Seq}(a) \to \mathrm{Seq}(b) \to \varphi_{152}) \to \mathrm{Seq}(a) \to \varphi_{218}$
\item $\varphi_{271} :\equiv \varphi_{118} \to ({\textstyle\sum_{k=1}^{n} a_{k}} + a_{S(n)} = {\textstyle\sum_{k=1}^{S(n)} a_{k}} \to \varphi_{118} \to \varphi_{103}) \to \mathrm{Seq}(a) \to \varphi_{209}$
\item $\varphi_{272} :\equiv (\mathrm{Seq}(a) \to \mathrm{Seq}(b) \to \forall n\,(\varphi_{68}) \to n \in B_{43} \to \varphi_{111}) \to \varphi_{242}$
\item $\varphi_{273} :\equiv (\mathrm{Seq}(a) \to \mathrm{Seq}(b) \to \forall n\,(\varphi_{68}) \to \varphi_{101}) \to \varphi_{136} \to (\mathrm{Seq}(a) \to \mathrm{Seq}(b) \to \forall n\,(\varphi_{68}) \to \varphi_{67}) \to \varphi_{130}$
\item $\varphi_{274} :\equiv (\mathrm{Seq}(a) \to \varphi_{209}) \to \varphi_{213} \to \mathrm{Seq}(a) \to \mathrm{Seq}(b) \to \forall n\,(\varphi_{68}) \to n \in B_{43} \to \varphi_{93}$
\item $\varphi_{275} :\equiv (\mathrm{Seq}(a) \to \mathrm{Seq}(b) \to \forall n\,(\varphi_{68}) \to n \in B_{43} \to \varphi_{37}) \to \varphi_{102} \to \varphi_{258}$
\item $\varphi_{276} :\equiv (\mathrm{Seq}(a) \to \mathrm{Seq}(b) \to \forall n\,(\varphi_{68}) \to n \in B_{43} \to \varphi_{93}) \to \varphi_{272}$
\item $\varphi_{277} :\equiv \forall B_{43}\,(\varphi_{86} \to \neg \forall n\,(\varphi_{68})) \to \varphi_{237} \to (\varphi_{86} \to \forall B_{43}\,(\varphi_{86})) \to \mathrm{Seq}(a) \to \varphi_{72}$
\item $\varphi_{278} :\equiv (\mathrm{Seq}(a) \to \mathrm{Seq}(b) \to \varphi_{165}) \to (\mathrm{Seq}(a) \to \mathrm{Seq}(b) \to \varphi_{164}) \to \varphi_{267}$
\item $\varphi_{279} :\equiv \forall n\,(\varphi_{182}) \to (\forall n\,(\varphi_{182}) \to \varphi_{185}) \to \varphi_{197} \to \mathrm{Seq}(a) \to \mathrm{Seq}(b) \to \varphi_{176}$
\end{abbrevlist}
\end{fullproof}

% 自動生成 (prover/build.js): PSumScal  ノード数 522
\begin{fullproof}{[Q2]}{$\forall l\,\forall a\,(l \in \mathbb{R} \to \mathrm{Seq}(a) \to \forall n \in \mathbb{N}\,({\textstyle\sum_{k=1}^{n} l \odot a(k)} = l \cdot {\textstyle\sum_{k=1}^{n} a_{k}}))$}{522}\phantomsection\label{tree:PSumScal}
\noindent$\Pi_{1}$:
\begin{ptree}
\AXC{$l \in \mathbb{R} \to \mathrm{Seq}(a) \to n \in \mathbb{N} \to n \in \mathbb{N}$\lbl{A1}}
\AXC{$(l \in \mathbb{R} \to \mathrm{Seq}(a) \to n \in \mathbb{N} \to n \in \mathbb{N}) \to \varphi_{51}$\lbl{A1}}
\BIC{$\varphi_{51}$}
\end{ptree}
\noindent$\Pi_{2}$:
\begin{ptree}
\AXC{$\forall n\,(\varphi_{43}) \to \varphi_{43}$\lbl{A2}}
\AXC{$l \in \mathbb{R} \to \mathrm{Seq}(a) \to \forall n\,(\varphi_{43}) \to \forall n\,(\varphi_{43})$\lbl{A1}}
\AXC{$\varphi_{204}$\lbl{A1}}
\BIC{$(\forall n\,(\varphi_{43}) \to \varphi_{43}) \to l \in \mathbb{R} \to \mathrm{Seq}(a) \to \forall n\,(\varphi_{43}) \to \varphi_{43}$}
\BIC{$l \in \mathbb{R} \to \mathrm{Seq}(a) \to \forall n\,(\varphi_{43}) \to \varphi_{43}$}
\end{ptree}
\noindent$\Pi_{3}$:
\begin{ptree}
\AXC{$\varphi_{15}$\lbl{A4}}
\AXC{$\forall y\,(\varphi_{5}) \to \varphi_{17}$\lbl{A3}}
\AXC{$\varphi_{5}$\lbl{A2}}
\UIC{$\forall y\,(\varphi_{5})$}
\AXC{$\forall y\,(\varphi_{5}) \to (\forall y\,(\varphi_{5}) \to \varphi_{17}) \to \varphi_{15} \to \varphi_{11}$\lbl{A1}}
\BIC{$(\forall y\,(\varphi_{5}) \to \varphi_{17}) \to \varphi_{15} \to \varphi_{11}$}
\BIC{$\varphi_{15} \to \varphi_{11}$}
\BIC{$\varphi_{11}$}
\end{ptree}
\noindent$\Pi_{4}$:
\begin{ptree}
\AXC{$\varphi_{16}$\lbl{A4}}
\AXC{$\forall n\,(\varphi_{6}) \to \varphi_{18}$\lbl{A3}}
\AXC{$\varphi_{6}$\lbl{A2}}
\UIC{$\forall n\,(\varphi_{6})$}
\AXC{$\forall n\,(\varphi_{6}) \to (\forall n\,(\varphi_{6}) \to \varphi_{18}) \to \varphi_{16} \to \varphi_{12}$\lbl{A1}}
\BIC{$(\forall n\,(\varphi_{6}) \to \varphi_{18}) \to \varphi_{16} \to \varphi_{12}$}
\BIC{$\varphi_{16} \to \varphi_{12}$}
\BIC{$\varphi_{12}$}
\end{ptree}
\noindent$\Pi_{5}$:
\begin{ptree}
\AXC{$y = z \to x = y \to x = z$\lbl{A6}}
\AXC{$(y = z \to x = y \to x = z) \to x = y \to y = z \to x = z$\lbl{A1}}
\BIC{$x = y \to y = z \to x = z$}
\UIC{$\forall z\,(x = y \to y = z \to x = z)$}
\UIC{$\forall y\,\forall z\,(x = y \to y = z \to x = z)$}
\UIC{$\forall x\,\forall y\,\forall z\,(x = y \to y = z \to x = z)$}
\AXC{$\forall x\,\forall y\,\forall z\,(x = y \to y = z \to x = z) \to \varphi_{24}$\lbl{A2}}
\BIC{$\varphi_{24}$}
\end{ptree}
\noindent$\Pi_{6}$:
\begin{ptree}
\AXC{$x = x$\lbl{A5}}
\AXC{$x = y \to x = x \to y = x$\lbl{A6}}
\AXC{$(x = y \to x = x \to y = x) \to x = x \to x = y \to y = x$\lbl{A1}}
\BIC{$x = x \to x = y \to y = x$}
\BIC{$x = y \to y = x$}
\UIC{$\forall y\,(x = y \to y = x)$}
\UIC{$\forall x\,\forall y\,(x = y \to y = x)$}
\end{ptree}
\noindent$\Pi_{7}$:
\begin{ptree}
\AXC{$\forall n\,(\varphi_{43}) \to \varphi_{43}$\lbl{A2}}
\AXC{$l \in \mathbb{R} \to \mathrm{Seq}(a) \to \forall n\,(\varphi_{43}) \to \forall n\,(\varphi_{43})$\lbl{A1}}
\AXC{$\varphi_{204}$\lbl{A1}}
\BIC{$(\forall n\,(\varphi_{43}) \to \varphi_{43}) \to l \in \mathbb{R} \to \mathrm{Seq}(a) \to \forall n\,(\varphi_{43}) \to \varphi_{43}$}
\BIC{$l \in \mathbb{R} \to \mathrm{Seq}(a) \to \forall n\,(\varphi_{43}) \to \varphi_{43}$}
\end{ptree}
\noindent$\Pi_{8}$:
\begin{ptree}
\AXC{$\stackrel{\Pi_{7}}{\vdots}$}\noLine
\UIC{$l \in \mathbb{R} \to \mathrm{Seq}(a) \to \forall n\,(\varphi_{43}) \to \varphi_{43}$}
\AXC{$(l \in \mathbb{R} \to \mathrm{Seq}(a) \to \forall n\,(\varphi_{43}) \to \varphi_{43}) \to \varphi_{179}$\lbl{A1}}
\BIC{$\varphi_{179}$}
\end{ptree}
\noindent$\Pi_{9}$:
\begin{ptree}
\AXC{$l \in \mathbb{R} \to \mathrm{Seq}(a) \to \forall n\,(\varphi_{43}) \to n \in B_{47} \to n \in B_{47}$\lbl{A1}}
\AXC{$\stackrel{\Pi_{8}}{\vdots}$}\noLine
\UIC{$\varphi_{179}$}
\BIC{$l \in \mathbb{R} \to \mathrm{Seq}(a) \to \forall n\,(\varphi_{43}) \to \varphi_{44}$}
\end{ptree}
\noindent$\Pi_{10}$:
\begin{ptree}
\AXC{$\stackrel{\Pi_{9}}{\vdots}$}\noLine
\UIC{$l \in \mathbb{R} \to \mathrm{Seq}(a) \to \forall n\,(\varphi_{43}) \to \varphi_{44}$}
\AXC{$\varphi_{180}$\lbl{A1}}
\BIC{$l \in \mathbb{R} \to \mathrm{Seq}(a) \to \forall n\,(\varphi_{43}) \to n \in B_{47} \to n \in \omega$}
\end{ptree}
\noindent$\Pi_{11}$:
\begin{ptree}
\AXC{$\forall n\,(n \in \mathbb{N} \leftrightarrow n \in \omega \land n \neq \emptyset)$\lbl{Def$\mathbb{N}$}}
\AXC{$\varphi_{42}$\lbl{A2}}
\BIC{$n \in \mathbb{N} \leftrightarrow n \in \omega \land n \neq \emptyset$}
\AXC{$(n \in \mathbb{N} \leftrightarrow n \in \omega \land n \neq \emptyset) \to \varphi_{195}$\lbl{A1}}
\BIC{$\varphi_{195}$}
\end{ptree}
\noindent$\Pi_{12}$:
\begin{ptree}
\AXC{$l \in \mathbb{R} \to \varphi_{127}$\lbl{A1}}
\AXC{$\stackrel{\Pi_{10}}{\vdots}$}\noLine
\UIC{$l \in \mathbb{R} \to \mathrm{Seq}(a) \to \forall n\,(\varphi_{43}) \to n \in B_{47} \to n \in \omega$}
\AXC{$\stackrel{\Pi_{11}}{\vdots}$}\noLine
\UIC{$\varphi_{195}$}
\BIC{$(l \in \mathbb{R} \to \varphi_{127}) \to l \in \mathbb{R} \to \varphi_{124}$}
\BIC{$l \in \mathbb{R} \to \varphi_{124}$}
\end{ptree}
\noindent$\Pi_{13}$:
\begin{ptree}
\AXC{$\forall a\,(\mathrm{Seq}(a) \to \varphi_{96})$\lbl{[S25]}}
\AXC{$\forall a\,(\mathrm{Seq}(a) \to \varphi_{96}) \to \mathrm{Seq}(a) \to \varphi_{96}$\lbl{A2}}
\BIC{$\mathrm{Seq}(a) \to \varphi_{96}$}
\AXC{$(\mathrm{Seq}(a) \to \varphi_{96}) \to (l \in \mathbb{R} \to \varphi_{120}) \to l \in \mathbb{R} \to \varphi_{171}$\lbl{A1}}
\BIC{$(l \in \mathbb{R} \to \varphi_{120}) \to l \in \mathbb{R} \to \varphi_{171}$}
\end{ptree}
\noindent$\Pi_{14}$:
\begin{ptree}
\AXC{$\varphi_{100}$\lbl{A2}}
\AXC{$l \in \mathbb{R} \to \varphi_{120}$\lbl{A1}}
\AXC{$\stackrel{\Pi_{13}}{\vdots}$}\noLine
\UIC{$(l \in \mathbb{R} \to \varphi_{120}) \to l \in \mathbb{R} \to \varphi_{171}$}
\BIC{$l \in \mathbb{R} \to \varphi_{171}$}
\AXC{$\varphi_{197}$\lbl{A1}}
\BIC{$l \in \mathbb{R} \to \mathrm{Seq}(a) \to \forall n\,(\varphi_{43}) \to n \in B_{47} \to \varphi_{74}$}
\AXC{$\varphi_{203}$\lbl{A1}}
\BIC{$\varphi_{100} \to l \in \mathbb{R} \to \mathrm{Seq}(a) \to \forall n\,(\varphi_{43}) \to n \in B_{47} \to \varphi_{61}$}
\BIC{$l \in \mathbb{R} \to \mathrm{Seq}(a) \to \forall n\,(\varphi_{43}) \to n \in B_{47} \to \varphi_{61}$}
\AXC{$\varphi_{205}$\lbl{A1}}
\BIC{$(l \in \mathbb{R} \to \varphi_{124}) \to l \in \mathbb{R} \to \mathrm{Seq}(a) \to \forall n\,(\varphi_{43}) \to \varphi_{62}$}
\end{ptree}
\noindent$\Pi_{15}$:
\begin{ptree}
\AXC{$\stackrel{\Pi_{12}}{\vdots}$}\noLine
\UIC{$l \in \mathbb{R} \to \varphi_{124}$}
\AXC{$\stackrel{\Pi_{14}}{\vdots}$}\noLine
\UIC{$(l \in \mathbb{R} \to \varphi_{124}) \to l \in \mathbb{R} \to \mathrm{Seq}(a) \to \forall n\,(\varphi_{43}) \to \varphi_{62}$}
\BIC{$l \in \mathbb{R} \to \mathrm{Seq}(a) \to \forall n\,(\varphi_{43}) \to \varphi_{62}$}
\AXC{$(l \in \mathbb{R} \to \mathrm{Seq}(a) \to \forall n\,(\varphi_{43}) \to \varphi_{62}) \to \varphi_{176}$\lbl{A1}}
\BIC{$\varphi_{176}$}
\end{ptree}
\noindent$\Pi_{16}$:
\begin{ptree}
\AXC{$\stackrel{\Pi_{6}}{\vdots}$}\noLine
\UIC{$\forall x\,\forall y\,(x = y \to y = x)$}
\AXC{$\varphi_{20}$\lbl{A2}}
\BIC{$\forall y\,({\textstyle\sum_{k=1}^{S(n)} a_{k}} = y \to y = {\textstyle\sum_{k=1}^{S(n)} a_{k}})$}
\AXC{$\varphi_{123}$\lbl{A2}}
\BIC{${\textstyle\sum_{k=1}^{S(n)} a_{k}} = {\textstyle\sum_{k=1}^{n} a_{k}} + a_{S(n)} \to {\textstyle\sum_{k=1}^{n} a_{k}} + a_{S(n)} = {\textstyle\sum_{k=1}^{S(n)} a_{k}}$}
\AXC{$\stackrel{\Pi_{15}}{\vdots}$}\noLine
\UIC{$\varphi_{176}$}
\BIC{$l \in \mathbb{R} \to \mathrm{Seq}(a) \to \forall n\,(\varphi_{43}) \to \varphi_{63}$}
\AXC{$\varphi_{211}$\lbl{A1}}
\BIC{$l \cdot ({\textstyle\sum_{k=1}^{n} a_{k}} + a_{S(n)}) = l \cdot ({\textstyle\sum_{k=1}^{n} a_{k}} + a_{S(n)}) \to \varphi_{191}$}
\end{ptree}
\noindent$\Pi_{17}$:
\begin{ptree}
\AXC{$x = x$\lbl{A5}}
\UIC{$\forall x\,(x = x)$}
\AXC{$\varphi_{38}$\lbl{A2}}
\BIC{$l \cdot ({\textstyle\sum_{k=1}^{n} a_{k}} + a_{S(n)}) = l \cdot ({\textstyle\sum_{k=1}^{n} a_{k}} + a_{S(n)})$}
\AXC{$\stackrel{\Pi_{16}}{\vdots}$}\noLine
\UIC{$l \cdot ({\textstyle\sum_{k=1}^{n} a_{k}} + a_{S(n)}) = l \cdot ({\textstyle\sum_{k=1}^{n} a_{k}} + a_{S(n)}) \to \varphi_{191}$}
\BIC{$\varphi_{191}$}
\end{ptree}
\noindent$\Pi_{18}$:
\begin{ptree}
\AXC{$z = w \to \varphi_{53}$\lbl{A6}}
\UIC{$\forall w\,(z = w \to \varphi_{53})$}
\UIC{$\forall z\,\forall w\,(z = w \to \varphi_{53})$}
\AXC{$\forall z\,\forall w\,(z = w \to \varphi_{53}) \to \varphi_{131}$\lbl{A2}}
\BIC{$\varphi_{131}$}
\AXC{$\varphi_{131} \to {\textstyle\sum_{k=1}^{n} a_{k}} + a_{S(n)} = {\textstyle\sum_{k=1}^{S(n)} a_{k}} \to \varphi_{103}$\lbl{A2}}
\BIC{${\textstyle\sum_{k=1}^{n} a_{k}} + a_{S(n)} = {\textstyle\sum_{k=1}^{S(n)} a_{k}} \to \varphi_{103}$}
\AXC{$\stackrel{\Pi_{17}}{\vdots}$}\noLine
\UIC{$\varphi_{191}$}
\BIC{$l \in \mathbb{R} \to \mathrm{Seq}(a) \to \forall n\,(\varphi_{43}) \to n \in B_{47} \to \varphi_{55}$}
\end{ptree}
\noindent$\Pi_{19}$:
\begin{ptree}
\AXC{$y = z \to x = y \to x = z$\lbl{A6}}
\AXC{$(y = z \to x = y \to x = z) \to x = y \to y = z \to x = z$\lbl{A1}}
\BIC{$x = y \to y = z \to x = z$}
\UIC{$\forall z\,(x = y \to y = z \to x = z)$}
\UIC{$\forall y\,\forall z\,(x = y \to y = z \to x = z)$}
\UIC{$\forall x\,\forall y\,\forall z\,(x = y \to y = z \to x = z)$}
\AXC{$\forall x\,\forall y\,\forall z\,(x = y \to y = z \to x = z) \to \varphi_{24}$\lbl{A2}}
\BIC{$\varphi_{24}$}
\end{ptree}
\noindent$\Pi_{20}$:
\begin{ptree}
\AXC{$\stackrel{\Pi_{19}}{\vdots}$}\noLine
\UIC{$\varphi_{24}$}
\AXC{$\varphi_{24} \to \varphi_{119}$\lbl{A2}}
\BIC{$\varphi_{119}$}
\AXC{$\varphi_{119} \to {\textstyle\sum_{k=1}^{S(n)} l \odot a(k)} = l \cdot {\textstyle\sum_{k=1}^{n} a_{k}} + l \cdot a_{S(n)} \to \varphi_{111}$\lbl{A2}}
\BIC{${\textstyle\sum_{k=1}^{S(n)} l \odot a(k)} = l \cdot {\textstyle\sum_{k=1}^{n} a_{k}} + l \cdot a_{S(n)} \to \varphi_{111}$}
\end{ptree}
\noindent$\Pi_{21}$:
\begin{ptree}
\AXC{$\stackrel{\Pi_{12}}{\vdots}$}\noLine
\UIC{$l \in \mathbb{R} \to \varphi_{124}$}
\AXC{$\forall n \in \mathbb{N}\,(S(n) \in \mathbb{N})$\lbl{[N2]}}
\AXC{$\forall n \in \mathbb{N}\,(S(n) \in \mathbb{N}) \to n \in \mathbb{N} \to S(n) \in \mathbb{N}$\lbl{A2}}
\BIC{$n \in \mathbb{N} \to S(n) \in \mathbb{N}$}
\AXC{$\varphi_{182}$\lbl{A1}}
\BIC{$(l \in \mathbb{R} \to \varphi_{124}) \to l \in \mathbb{R} \to \mathrm{Seq}(a) \to \varphi_{115}$}
\BIC{$l \in \mathbb{R} \to \mathrm{Seq}(a) \to \varphi_{115}$}
\end{ptree}
\noindent$\Pi_{22}$:
\begin{ptree}
\AXC{$\forall a\,\forall n\,(\mathrm{Seq}(a) \to n \in \mathbb{N} \to a_{n} \in \mathbb{R})$\lbl{[S1]}}
\AXC{$\varphi_{59}$\lbl{A2}}
\BIC{$\forall n\,(\mathrm{Seq}(a) \to n \in \mathbb{N} \to a_{n} \in \mathbb{R})$}
\AXC{$\varphi_{60}$\lbl{A2}}
\BIC{$\mathrm{Seq}(a) \to S(n) \in \mathbb{N} \to a_{S(n)} \in \mathbb{R}$}
\AXC{$(\mathrm{Seq}(a) \to S(n) \in \mathbb{N} \to a_{S(n)} \in \mathbb{R}) \to \varphi_{199}$\lbl{A1}}
\BIC{$\varphi_{199}$}
\end{ptree}
\noindent$\Pi_{23}$:
\begin{ptree}
\AXC{$\stackrel{\Pi_{21}}{\vdots}$}\noLine
\UIC{$l \in \mathbb{R} \to \mathrm{Seq}(a) \to \varphi_{115}$}
\AXC{$l \in \mathbb{R} \to \varphi_{120}$\lbl{A1}}
\AXC{$\stackrel{\Pi_{22}}{\vdots}$}\noLine
\UIC{$\varphi_{199}$}
\BIC{$(l \in \mathbb{R} \to \mathrm{Seq}(a) \to \varphi_{115}) \to l \in \mathbb{R} \to \mathrm{Seq}(a) \to \varphi_{128}$}
\BIC{$l \in \mathbb{R} \to \mathrm{Seq}(a) \to \varphi_{128}$}
\end{ptree}
\noindent$\Pi_{24}$:
\begin{ptree}
\AXC{$\forall a\,(\mathrm{Seq}(a) \to \varphi_{96})$\lbl{[S25]}}
\AXC{$\forall a\,(\mathrm{Seq}(a) \to \varphi_{96}) \to \mathrm{Seq}(a) \to \varphi_{96}$\lbl{A2}}
\BIC{$\mathrm{Seq}(a) \to \varphi_{96}$}
\AXC{$(\mathrm{Seq}(a) \to \varphi_{96}) \to (l \in \mathbb{R} \to \varphi_{120}) \to l \in \mathbb{R} \to \varphi_{171}$\lbl{A1}}
\BIC{$(l \in \mathbb{R} \to \varphi_{120}) \to l \in \mathbb{R} \to \varphi_{171}$}
\end{ptree}
\noindent$\Pi_{25}$:
\begin{ptree}
\AXC{$\forall a\,\forall n\,(\mathrm{Seq}(a) \to n \in \mathbb{N} \to a_{n} \in \mathbb{R})$\lbl{[S1]}}
\AXC{$\varphi_{86}$\lbl{A2}}
\BIC{$\forall n\,(\mathrm{Seq}(\Sigma a) \to n \in \mathbb{N} \to {\textstyle\sum_{k=1}^{n} a_{k}} \in \mathbb{R})$}
\AXC{$\varphi_{91}$\lbl{A2}}
\BIC{$\mathrm{Seq}(\Sigma a) \to n \in \mathbb{N} \to {\textstyle\sum_{k=1}^{n} a_{k}} \in \mathbb{R}$}
\AXC{$(\mathrm{Seq}(\Sigma a) \to n \in \mathbb{N} \to {\textstyle\sum_{k=1}^{n} a_{k}} \in \mathbb{R}) \to \varphi_{200}$\lbl{A1}}
\BIC{$\varphi_{200}$}
\end{ptree}
\noindent$\Pi_{26}$:
\begin{ptree}
\AXC{$l \in \mathbb{R} \to \varphi_{120}$\lbl{A1}}
\AXC{$\stackrel{\Pi_{24}}{\vdots}$}\noLine
\UIC{$(l \in \mathbb{R} \to \varphi_{120}) \to l \in \mathbb{R} \to \varphi_{171}$}
\BIC{$l \in \mathbb{R} \to \varphi_{171}$}
\AXC{$(l \in \mathbb{R} \to \varphi_{171}) \to l \in \mathbb{R} \to \mathrm{Seq}(a) \to \varphi_{116}$\lbl{A1}}
\BIC{$l \in \mathbb{R} \to \mathrm{Seq}(a) \to \varphi_{116}$}
\AXC{$\stackrel{\Pi_{25}}{\vdots}$}\noLine
\UIC{$\varphi_{200}$}
\BIC{$(l \in \mathbb{R} \to \varphi_{124}) \to l \in \mathbb{R} \to \mathrm{Seq}(a) \to \varphi_{132}$}
\end{ptree}
\noindent$\Pi_{27}$:
\begin{ptree}
\AXC{$\varphi_{32}$\lbl{(R10)}}
\AXC{$\varphi_{32} \to \varphi_{28}$\lbl{A2}}
\BIC{$\varphi_{28}$}
\AXC{$\varphi_{28} \to \varphi_{94}$\lbl{A2}}
\BIC{$\varphi_{94}$}
\AXC{$\varphi_{94} \to l \in \mathbb{R} \to {\textstyle\sum_{k=1}^{n} a_{k}} \in \mathbb{R} \to \varphi_{81}$\lbl{A2}}
\BIC{$l \in \mathbb{R} \to {\textstyle\sum_{k=1}^{n} a_{k}} \in \mathbb{R} \to \varphi_{81}$}
\end{ptree}
\noindent$\Pi_{28}$:
\begin{ptree}
\AXC{$l \in \mathbb{R} \to \varphi_{126}$\lbl{A1}}
\AXC{$\stackrel{\Pi_{27}}{\vdots}$}\noLine
\UIC{$l \in \mathbb{R} \to {\textstyle\sum_{k=1}^{n} a_{k}} \in \mathbb{R} \to \varphi_{81}$}
\AXC{$(l \in \mathbb{R} \to {\textstyle\sum_{k=1}^{n} a_{k}} \in \mathbb{R} \to \varphi_{81}) \to (l \in \mathbb{R} \to \varphi_{126}) \to \varphi_{206}$\lbl{A1}}
\BIC{$(l \in \mathbb{R} \to \varphi_{126}) \to \varphi_{206}$}
\BIC{$\varphi_{206}$}
\end{ptree}
\noindent$\Pi_{29}$:
\begin{ptree}
\AXC{$\stackrel{\Pi_{23}}{\vdots}$}\noLine
\UIC{$l \in \mathbb{R} \to \mathrm{Seq}(a) \to \varphi_{128}$}
\AXC{$\stackrel{\Pi_{12}}{\vdots}$}\noLine
\UIC{$l \in \mathbb{R} \to \varphi_{124}$}
\AXC{$\stackrel{\Pi_{26}}{\vdots}$}\noLine
\UIC{$(l \in \mathbb{R} \to \varphi_{124}) \to l \in \mathbb{R} \to \mathrm{Seq}(a) \to \varphi_{132}$}
\BIC{$l \in \mathbb{R} \to \mathrm{Seq}(a) \to \varphi_{132}$}
\AXC{$\stackrel{\Pi_{28}}{\vdots}$}\noLine
\UIC{$\varphi_{206}$}
\BIC{$(l \in \mathbb{R} \to \mathrm{Seq}(a) \to \varphi_{128}) \to \varphi_{164}$}
\BIC{$\varphi_{164}$}
\AXC{$\varphi_{164} \to \varphi_{134} \to \varphi_{165}$\lbl{A1}}
\BIC{$\varphi_{134} \to \varphi_{165}$}
\end{ptree}
\noindent$\Pi_{30}$:
\begin{ptree}
\AXC{$\stackrel{\Pi_{6}}{\vdots}$}\noLine
\UIC{$\forall x\,\forall y\,(x = y \to y = x)$}
\AXC{$\forall x\,\forall y\,(x = y \to y = x) \to \varphi_{39}$\lbl{A2}}
\BIC{$\varphi_{39}$}
\AXC{$\varphi_{39} \to \varphi_{134}$\lbl{A2}}
\BIC{$\varphi_{134}$}
\AXC{$\stackrel{\Pi_{29}}{\vdots}$}\noLine
\UIC{$\varphi_{134} \to \varphi_{165}$}
\BIC{$\varphi_{165}$}
\end{ptree}
\noindent$\Pi_{31}$:
\begin{ptree}
\AXC{$y = z \to x = y \to x = z$\lbl{A6}}
\AXC{$(y = z \to x = y \to x = z) \to x = y \to y = z \to x = z$\lbl{A1}}
\BIC{$x = y \to y = z \to x = z$}
\UIC{$\forall z\,(x = y \to y = z \to x = z)$}
\UIC{$\forall y\,\forall z\,(x = y \to y = z \to x = z)$}
\UIC{$\forall x\,\forall y\,\forall z\,(x = y \to y = z \to x = z)$}
\AXC{$\forall x\,\forall y\,\forall z\,(x = y \to y = z \to x = z) \to \varphi_{24}$\lbl{A2}}
\BIC{$\varphi_{24}$}
\end{ptree}
\noindent$\Pi_{32}$:
\begin{ptree}
\AXC{$\stackrel{\Pi_{31}}{\vdots}$}\noLine
\UIC{$\varphi_{24}$}
\AXC{$\varphi_{24} \to \varphi_{129}$\lbl{A2}}
\BIC{$\varphi_{129}$}
\AXC{$\varphi_{129} \to {\textstyle\sum_{k=1}^{S(n)} l \odot a(k)} = l \cdot {\textstyle\sum_{k=1}^{n} a_{k}} + l \odot a(S(n)) \to \varphi_{133}$\lbl{A2}}
\BIC{${\textstyle\sum_{k=1}^{S(n)} l \odot a(k)} = l \cdot {\textstyle\sum_{k=1}^{n} a_{k}} + l \odot a(S(n)) \to \varphi_{133}$}
\end{ptree}
\noindent$\Pi_{33}$:
\begin{ptree}
\AXC{$z = w \to \varphi_{112}$\lbl{A6}}
\UIC{$\forall w\,(z = w \to \varphi_{112})$}
\UIC{$\forall z\,\forall w\,(z = w \to \varphi_{112})$}
\AXC{$\forall z\,\forall w\,(z = w \to \varphi_{112}) \to \varphi_{149}$\lbl{A2}}
\BIC{$\varphi_{149}$}
\AXC{$\varphi_{149} \to l \odot a(S(n)) = l \cdot a_{S(n)} \to \varphi_{145}$\lbl{A2}}
\BIC{$l \odot a(S(n)) = l \cdot a_{S(n)} \to \varphi_{145}$}
\end{ptree}
\noindent$\Pi_{34}$:
\begin{ptree}
\AXC{$\varphi_{70}$\lbl{[S3]}}
\AXC{$\varphi_{70} \to \forall a\,(l \in \mathbb{R} \to \varphi_{26})$\lbl{A2}}
\BIC{$\forall a\,(l \in \mathbb{R} \to \varphi_{26})$}
\AXC{$\forall a\,(l \in \mathbb{R} \to \varphi_{26}) \to l \in \mathbb{R} \to \varphi_{26}$\lbl{A2}}
\BIC{$l \in \mathbb{R} \to \varphi_{26}$}
\end{ptree}
\noindent$\Pi_{35}$:
\begin{ptree}
\AXC{$l \in \mathbb{R} \to \varphi_{120}$\lbl{A1}}
\AXC{$l \in \mathbb{R} \to \varphi_{126}$\lbl{A1}}
\AXC{$\stackrel{\Pi_{34}}{\vdots}$}\noLine
\UIC{$l \in \mathbb{R} \to \varphi_{26}$}
\AXC{$(l \in \mathbb{R} \to \varphi_{26}) \to (l \in \mathbb{R} \to \varphi_{126}) \to (l \in \mathbb{R} \to \varphi_{120}) \to \varphi_{158}$\lbl{A1}}
\BIC{$(l \in \mathbb{R} \to \varphi_{126}) \to (l \in \mathbb{R} \to \varphi_{120}) \to \varphi_{158}$}
\BIC{$(l \in \mathbb{R} \to \varphi_{120}) \to \varphi_{158}$}
\BIC{$\varphi_{158}$}
\end{ptree}
\noindent$\Pi_{36}$:
\begin{ptree}
\AXC{$\varphi_{83}$\lbl{A2}}
\AXC{$\stackrel{\Pi_{35}}{\vdots}$}\noLine
\UIC{$\varphi_{158}$}
\AXC{$\varphi_{158} \to l \in \mathbb{R} \to \mathrm{Seq}(a) \to \forall n\,(\varphi_{43}) \to \varphi_{56}$\lbl{A1}}
\BIC{$l \in \mathbb{R} \to \mathrm{Seq}(a) \to \forall n\,(\varphi_{43}) \to \varphi_{56}$}
\AXC{$\varphi_{201}$\lbl{A1}}
\BIC{$\varphi_{83} \to l \in \mathbb{R} \to \mathrm{Seq}(a) \to \forall n\,(\varphi_{43}) \to n \in B_{47} \to \varphi_{45}$}
\BIC{$l \in \mathbb{R} \to \mathrm{Seq}(a) \to \forall n\,(\varphi_{43}) \to n \in B_{47} \to \varphi_{45}$}
\end{ptree}
\noindent$\Pi_{37}$:
\begin{ptree}
\AXC{$\stackrel{\Pi_{36}}{\vdots}$}\noLine
\UIC{$l \in \mathbb{R} \to \mathrm{Seq}(a) \to \forall n\,(\varphi_{43}) \to n \in B_{47} \to \varphi_{45}$}
\AXC{$(l \in \mathbb{R} \to \mathrm{Seq}(a) \to \forall n\,(\varphi_{43}) \to n \in B_{47} \to \varphi_{45}) \to \varphi_{183}$\lbl{A1}}
\BIC{$\varphi_{183}$}
\end{ptree}
\noindent$\Pi_{38}$:
\begin{ptree}
\AXC{$x = x$\lbl{A5}}
\UIC{$\forall x\,(x = x)$}
\AXC{$\forall x\,(x = x) \to \varphi_{52}$\lbl{A2}}
\BIC{$\varphi_{52}$}
\AXC{$\stackrel{\Pi_{21}}{\vdots}$}\noLine
\UIC{$l \in \mathbb{R} \to \mathrm{Seq}(a) \to \varphi_{115}$}
\AXC{$\stackrel{\Pi_{37}}{\vdots}$}\noLine
\UIC{$\varphi_{183}$}
\BIC{$l \in \mathbb{R} \to \mathrm{Seq}(a) \to \forall n\,(\varphi_{43}) \to \varphi_{36}$}
\AXC{$\varphi_{213}$\lbl{A1}}
\BIC{$\varphi_{52} \to (l \odot a(S(n)) = l \cdot a_{S(n)} \to \varphi_{145}) \to \varphi_{167}$}
\BIC{$(l \odot a(S(n)) = l \cdot a_{S(n)} \to \varphi_{145}) \to \varphi_{167}$}
\end{ptree}
\noindent$\Pi_{39}$:
\begin{ptree}
\AXC{$\stackrel{\Pi_{33}}{\vdots}$}\noLine
\UIC{$l \odot a(S(n)) = l \cdot a_{S(n)} \to \varphi_{145}$}
\AXC{$\stackrel{\Pi_{38}}{\vdots}$}\noLine
\UIC{$(l \odot a(S(n)) = l \cdot a_{S(n)} \to \varphi_{145}) \to \varphi_{167}$}
\BIC{$\varphi_{167}$}
\end{ptree}
\noindent$\Pi_{40}$:
\begin{ptree}
\AXC{$y = z \to x = y \to x = z$\lbl{A6}}
\AXC{$(y = z \to x = y \to x = z) \to x = y \to y = z \to x = z$\lbl{A1}}
\BIC{$x = y \to y = z \to x = z$}
\UIC{$\forall z\,(x = y \to y = z \to x = z)$}
\UIC{$\forall y\,\forall z\,(x = y \to y = z \to x = z)$}
\UIC{$\forall x\,\forall y\,\forall z\,(x = y \to y = z \to x = z)$}
\AXC{$\forall x\,\forall y\,\forall z\,(x = y \to y = z \to x = z) \to \varphi_{24}$\lbl{A2}}
\BIC{$\varphi_{24}$}
\end{ptree}
\noindent$\Pi_{41}$:
\begin{ptree}
\AXC{$z = w \to \varphi_{122}$\lbl{A6}}
\UIC{$\forall w\,(z = w \to \varphi_{122})$}
\UIC{$\forall z\,\forall w\,(z = w \to \varphi_{122})$}
\AXC{$\forall z\,\forall w\,(z = w \to \varphi_{122}) \to \varphi_{153}$\lbl{A2}}
\BIC{$\varphi_{153}$}
\AXC{$\varphi_{153} \to {\textstyle\sum_{k=1}^{n} l \odot a(k)} = l \cdot {\textstyle\sum_{k=1}^{n} a_{k}} \to \varphi_{72} \to \varphi_{66}$\lbl{A2}}
\BIC{${\textstyle\sum_{k=1}^{n} l \odot a(k)} = l \cdot {\textstyle\sum_{k=1}^{n} a_{k}} \to \varphi_{72} \to \varphi_{66}$}
\end{ptree}
\noindent$\Pi_{42}$:
\begin{ptree}
\AXC{$\stackrel{\Pi_{12}}{\vdots}$}\noLine
\UIC{$l \in \mathbb{R} \to \varphi_{124}$}
\AXC{$\stackrel{\Pi_{9}}{\vdots}$}\noLine
\UIC{$l \in \mathbb{R} \to \mathrm{Seq}(a) \to \forall n\,(\varphi_{43}) \to \varphi_{44}$}
\AXC{$\varphi_{184}$\lbl{A1}}
\BIC{$l \in \mathbb{R} \to \mathrm{Seq}(a) \to \varphi_{135}$}
\AXC{$\varphi_{202}$\lbl{A1}}
\BIC{$(l \in \mathbb{R} \to \varphi_{124}) \to l \in \mathbb{R} \to \mathrm{Seq}(a) \to \forall n\,(\varphi_{43}) \to \varphi_{49}$}
\BIC{$l \in \mathbb{R} \to \mathrm{Seq}(a) \to \forall n\,(\varphi_{43}) \to \varphi_{49}$}
\AXC{$\varphi_{214}$\lbl{A1}}
\BIC{$\varphi_{72} \to ({\textstyle\sum_{k=1}^{n} l \odot a(k)} = l \cdot {\textstyle\sum_{k=1}^{n} a_{k}} \to \varphi_{72} \to \varphi_{66}) \to l \in \mathbb{R} \to \varphi_{166}$}
\end{ptree}
\noindent$\Pi_{43}$:
\begin{ptree}
\AXC{$x = x$\lbl{A5}}
\UIC{$\forall x\,(x = x)$}
\AXC{$\forall x\,(x = x) \to \varphi_{72}$\lbl{A2}}
\BIC{$\varphi_{72}$}
\AXC{$\stackrel{\Pi_{42}}{\vdots}$}\noLine
\UIC{$\varphi_{72} \to ({\textstyle\sum_{k=1}^{n} l \odot a(k)} = l \cdot {\textstyle\sum_{k=1}^{n} a_{k}} \to \varphi_{72} \to \varphi_{66}) \to l \in \mathbb{R} \to \varphi_{166}$}
\BIC{$({\textstyle\sum_{k=1}^{n} l \odot a(k)} = l \cdot {\textstyle\sum_{k=1}^{n} a_{k}} \to \varphi_{72} \to \varphi_{66}) \to l \in \mathbb{R} \to \varphi_{166}$}
\end{ptree}
\noindent$\Pi_{44}$:
\begin{ptree}
\AXC{$\stackrel{\Pi_{41}}{\vdots}$}\noLine
\UIC{${\textstyle\sum_{k=1}^{n} l \odot a(k)} = l \cdot {\textstyle\sum_{k=1}^{n} a_{k}} \to \varphi_{72} \to \varphi_{66}$}
\AXC{$\stackrel{\Pi_{43}}{\vdots}$}\noLine
\UIC{$({\textstyle\sum_{k=1}^{n} l \odot a(k)} = l \cdot {\textstyle\sum_{k=1}^{n} a_{k}} \to \varphi_{72} \to \varphi_{66}) \to l \in \mathbb{R} \to \varphi_{166}$}
\BIC{$l \in \mathbb{R} \to \varphi_{166}$}
\end{ptree}
\noindent$\Pi_{45}$:
\begin{ptree}
\AXC{$\varphi_{70}$\lbl{[S3]}}
\AXC{$\varphi_{70} \to \forall a\,(l \in \mathbb{R} \to \varphi_{26})$\lbl{A2}}
\BIC{$\forall a\,(l \in \mathbb{R} \to \varphi_{26})$}
\AXC{$\forall a\,(l \in \mathbb{R} \to \varphi_{26}) \to l \in \mathbb{R} \to \varphi_{26}$\lbl{A2}}
\BIC{$l \in \mathbb{R} \to \varphi_{26}$}
\end{ptree}
\noindent$\Pi_{46}$:
\begin{ptree}
\AXC{$l \in \mathbb{R} \to \varphi_{120}$\lbl{A1}}
\AXC{$l \in \mathbb{R} \to \varphi_{126}$\lbl{A1}}
\AXC{$\stackrel{\Pi_{45}}{\vdots}$}\noLine
\UIC{$l \in \mathbb{R} \to \varphi_{26}$}
\AXC{$(l \in \mathbb{R} \to \varphi_{26}) \to (l \in \mathbb{R} \to \varphi_{126}) \to (l \in \mathbb{R} \to \varphi_{120}) \to \varphi_{158}$\lbl{A1}}
\BIC{$(l \in \mathbb{R} \to \varphi_{126}) \to (l \in \mathbb{R} \to \varphi_{120}) \to \varphi_{158}$}
\BIC{$(l \in \mathbb{R} \to \varphi_{120}) \to \varphi_{158}$}
\BIC{$\varphi_{158}$}
\end{ptree}
\noindent$\Pi_{47}$:
\begin{ptree}
\AXC{$\forall a\,(\mathrm{Seq}(a) \to \varphi_{96})$\lbl{[S25]}}
\AXC{$\forall a\,(\mathrm{Seq}(a) \to \varphi_{96}) \to \mathrm{Seq}(l \odot a) \to \varphi_{138}$\lbl{A2}}
\BIC{$\mathrm{Seq}(l \odot a) \to \varphi_{138}$}
\AXC{$\varphi_{210}$\lbl{A1}}
\BIC{$(l \in \mathbb{R} \to \mathrm{Seq}(a) \to \varphi_{117}) \to l \in \mathbb{R} \to \varphi_{175}$}
\end{ptree}
\noindent$\Pi_{48}$:
\begin{ptree}
\AXC{$\stackrel{\Pi_{46}}{\vdots}$}\noLine
\UIC{$\varphi_{158}$}
\AXC{$\varphi_{158} \to l \in \mathbb{R} \to \mathrm{Seq}(a) \to \varphi_{117}$\lbl{A1}}
\BIC{$l \in \mathbb{R} \to \mathrm{Seq}(a) \to \varphi_{117}$}
\AXC{$\stackrel{\Pi_{47}}{\vdots}$}\noLine
\UIC{$(l \in \mathbb{R} \to \mathrm{Seq}(a) \to \varphi_{117}) \to l \in \mathbb{R} \to \varphi_{175}$}
\BIC{$l \in \mathbb{R} \to \varphi_{175}$}
\AXC{$(l \in \mathbb{R} \to \varphi_{175}) \to l \in \mathbb{R} \to \varphi_{163}$\lbl{A1}}
\BIC{$l \in \mathbb{R} \to \varphi_{163}$}
\end{ptree}
\noindent$\Pi_{49}$:
\begin{ptree}
\AXC{$\varphi_{57} \to \varphi_{50}$\lbl{A2}}
\AXC{$\stackrel{\Pi_{48}}{\vdots}$}\noLine
\UIC{$l \in \mathbb{R} \to \varphi_{163}$}
\AXC{$(l \in \mathbb{R} \to \varphi_{163}) \to (\varphi_{57} \to \varphi_{50}) \to l \in \mathbb{R} \to \varphi_{161}$\lbl{A1}}
\BIC{$(\varphi_{57} \to \varphi_{50}) \to l \in \mathbb{R} \to \varphi_{161}$}
\BIC{$l \in \mathbb{R} \to \varphi_{161}$}
\AXC{$(l \in \mathbb{R} \to \varphi_{161}) \to (l \in \mathbb{R} \to \varphi_{124}) \to \varphi_{162}$\lbl{A1}}
\BIC{$(l \in \mathbb{R} \to \varphi_{124}) \to \varphi_{162}$}
\end{ptree}
\noindent$\Pi_{50}$:
\begin{ptree}
\AXC{$\stackrel{\Pi_{40}}{\vdots}$}\noLine
\UIC{$\varphi_{24}$}
\AXC{$\varphi_{24} \to \varphi_{136}$\lbl{A2}}
\BIC{$\varphi_{136}$}
\AXC{$\varphi_{136} \to \varphi_{168}$\lbl{A2}}
\BIC{$\varphi_{168}$}
\AXC{$\stackrel{\Pi_{44}}{\vdots}$}\noLine
\UIC{$l \in \mathbb{R} \to \varphi_{166}$}
\AXC{$\stackrel{\Pi_{12}}{\vdots}$}\noLine
\UIC{$l \in \mathbb{R} \to \varphi_{124}$}
\AXC{$\stackrel{\Pi_{49}}{\vdots}$}\noLine
\UIC{$(l \in \mathbb{R} \to \varphi_{124}) \to \varphi_{162}$}
\BIC{$\varphi_{162}$}
\AXC{$\varphi_{162} \to \varphi_{212}$\lbl{A1}}
\BIC{$\varphi_{212}$}
\BIC{$\varphi_{168} \to l \in \mathbb{R} \to \mathrm{Seq}(a) \to \forall n\,(\varphi_{43}) \to n \in B_{47} \to \varphi_{82}$}
\BIC{$l \in \mathbb{R} \to \mathrm{Seq}(a) \to \forall n\,(\varphi_{43}) \to n \in B_{47} \to \varphi_{82}$}
\end{ptree}
\noindent$\Pi_{51}$:
\begin{ptree}
\AXC{$\stackrel{\Pi_{32}}{\vdots}$}\noLine
\UIC{${\textstyle\sum_{k=1}^{S(n)} l \odot a(k)} = l \cdot {\textstyle\sum_{k=1}^{n} a_{k}} + l \odot a(S(n)) \to \varphi_{133}$}
\AXC{$\stackrel{\Pi_{39}}{\vdots}$}\noLine
\UIC{$\varphi_{167}$}
\AXC{$\stackrel{\Pi_{50}}{\vdots}$}\noLine
\UIC{$l \in \mathbb{R} \to \mathrm{Seq}(a) \to \forall n\,(\varphi_{43}) \to n \in B_{47} \to \varphi_{82}$}
\AXC{$\varphi_{218}$\lbl{A1}}
\BIC{$\varphi_{167} \to \varphi_{194}$}
\BIC{$\varphi_{194}$}
\BIC{$l \in \mathbb{R} \to \mathrm{Seq}(a) \to \forall n\,(\varphi_{43}) \to n \in B_{47} \to \varphi_{76}$}
\end{ptree}
\noindent$\Pi_{52}$:
\begin{ptree}
\AXC{$\stackrel{\Pi_{20}}{\vdots}$}\noLine
\UIC{${\textstyle\sum_{k=1}^{S(n)} l \odot a(k)} = l \cdot {\textstyle\sum_{k=1}^{n} a_{k}} + l \cdot a_{S(n)} \to \varphi_{111}$}
\AXC{$\stackrel{\Pi_{30}}{\vdots}$}\noLine
\UIC{$\varphi_{165}$}
\AXC{$\stackrel{\Pi_{51}}{\vdots}$}\noLine
\UIC{$l \in \mathbb{R} \to \mathrm{Seq}(a) \to \forall n\,(\varphi_{43}) \to n \in B_{47} \to \varphi_{76}$}
\AXC{$\varphi_{216}$\lbl{A1}}
\BIC{$\varphi_{165} \to \varphi_{193}$}
\BIC{$\varphi_{193}$}
\BIC{$l \in \mathbb{R} \to \mathrm{Seq}(a) \to \forall n\,(\varphi_{43}) \to n \in B_{47} \to \varphi_{68}$}
\end{ptree}
\noindent$\Pi_{53}$:
\begin{ptree}
\AXC{$\stackrel{\Pi_{18}}{\vdots}$}\noLine
\UIC{$l \in \mathbb{R} \to \mathrm{Seq}(a) \to \forall n\,(\varphi_{43}) \to n \in B_{47} \to \varphi_{55}$}
\AXC{$\stackrel{\Pi_{52}}{\vdots}$}\noLine
\UIC{$l \in \mathbb{R} \to \mathrm{Seq}(a) \to \forall n\,(\varphi_{43}) \to n \in B_{47} \to \varphi_{68}$}
\AXC{$\varphi_{215}$\lbl{A1}}
\BIC{$(l \in \mathbb{R} \to \mathrm{Seq}(a) \to \forall n\,(\varphi_{43}) \to n \in B_{47} \to \varphi_{55}) \to \varphi_{186}$}
\BIC{$\varphi_{186}$}
\end{ptree}
\noindent$\Pi_{54}$:
\begin{ptree}
\AXC{$\stackrel{\Pi_{5}}{\vdots}$}\noLine
\UIC{$\varphi_{24}$}
\AXC{$\varphi_{24} \to \varphi_{105}$\lbl{A2}}
\BIC{$\varphi_{105}$}
\AXC{$\varphi_{105} \to {\textstyle\sum_{k=1}^{S(n)} l \odot a(k)} = l \cdot ({\textstyle\sum_{k=1}^{n} a_{k}} + a_{S(n)}) \to \varphi_{93}$\lbl{A2}}
\BIC{${\textstyle\sum_{k=1}^{S(n)} l \odot a(k)} = l \cdot ({\textstyle\sum_{k=1}^{n} a_{k}} + a_{S(n)}) \to \varphi_{93}$}
\AXC{$\stackrel{\Pi_{53}}{\vdots}$}\noLine
\UIC{$\varphi_{186}$}
\BIC{$l \in \mathbb{R} \to \mathrm{Seq}(a) \to \forall n\,(\varphi_{43}) \to \varphi_{69}$}
\end{ptree}
\noindent$\Pi_{55}$:
\begin{ptree}
\AXC{$y = z \to x = y \to x = z$\lbl{A6}}
\AXC{$(y = z \to x = y \to x = z) \to x = y \to y = z \to x = z$\lbl{A1}}
\BIC{$x = y \to y = z \to x = z$}
\UIC{$\forall z\,(x = y \to y = z \to x = z)$}
\UIC{$\forall y\,\forall z\,(x = y \to y = z \to x = z)$}
\UIC{$\forall x\,\forall y\,\forall z\,(x = y \to y = z \to x = z)$}
\AXC{$\forall x\,\forall y\,\forall z\,(x = y \to y = z \to x = z) \to \varphi_{25}$\lbl{A2}}
\BIC{$\varphi_{25}$}
\end{ptree}
\noindent$\Pi_{56}$:
\begin{ptree}
\AXC{$\forall a\,(\mathrm{Seq}(a) \to \varphi_{96})$\lbl{[S25]}}
\AXC{$\forall a\,(\mathrm{Seq}(a) \to \varphi_{96}) \to \mathrm{Seq}(a) \to \varphi_{96}$\lbl{A2}}
\BIC{$\mathrm{Seq}(a) \to \varphi_{96}$}
\AXC{$(\mathrm{Seq}(a) \to \varphi_{96}) \to \varphi_{142}$\lbl{A1}}
\BIC{$\varphi_{142}$}
\end{ptree}
\noindent$\Pi_{57}$:
\begin{ptree}
\AXC{$l \in \mathbb{R} \to \mathrm{Seq}(a) \to \mathrm{Seq}(a)$\lbl{A1}}
\AXC{$\stackrel{\Pi_{56}}{\vdots}$}\noLine
\UIC{$\varphi_{142}$}
\BIC{$l \in \mathbb{R} \to \mathrm{Seq}(a) \to \varphi_{96}$}
\AXC{$\varphi_{151}$\lbl{A1}}
\BIC{$l \in \mathbb{R} \to \mathrm{Seq}(a) \to {\textstyle\sum_{k=1}^{1_{\omega}} a_{k}} = a_{1_{\omega}}$}
\AXC{$(l \in \mathbb{R} \to \mathrm{Seq}(a) \to {\textstyle\sum_{k=1}^{1_{\omega}} a_{k}} = a_{1_{\omega}}) \to \varphi_{97}$\lbl{A1}}
\BIC{$\varphi_{97}$}
\end{ptree}
\noindent$\Pi_{58}$:
\begin{ptree}
\AXC{$\stackrel{\Pi_{6}}{\vdots}$}\noLine
\UIC{$\forall x\,\forall y\,(x = y \to y = x)$}
\AXC{$\varphi_{21}$\lbl{A2}}
\BIC{$\forall y\,({\textstyle\sum_{k=1}^{1_{\omega}} a_{k}} = y \to y = {\textstyle\sum_{k=1}^{1_{\omega}} a_{k}})$}
\AXC{$\varphi_{92}$\lbl{A2}}
\BIC{${\textstyle\sum_{k=1}^{1_{\omega}} a_{k}} = a_{1_{\omega}} \to a_{1_{\omega}} = {\textstyle\sum_{k=1}^{1_{\omega}} a_{k}}$}
\AXC{$\stackrel{\Pi_{57}}{\vdots}$}\noLine
\UIC{$\varphi_{97}$}
\BIC{$l \in \mathbb{R} \to \mathrm{Seq}(a) \to a_{1_{\omega}} = {\textstyle\sum_{k=1}^{1_{\omega}} a_{k}}$}
\AXC{$\varphi_{172}$\lbl{A1}}
\BIC{$l \cdot a_{1_{\omega}} = l \cdot a_{1_{\omega}} \to \varphi_{146}$}
\end{ptree}
\noindent$\Pi_{59}$:
\begin{ptree}
\AXC{$x = x$\lbl{A5}}
\UIC{$\forall x\,(x = x)$}
\AXC{$\forall x\,(x = x) \to l \cdot a_{1_{\omega}} = l \cdot a_{1_{\omega}}$\lbl{A2}}
\BIC{$l \cdot a_{1_{\omega}} = l \cdot a_{1_{\omega}}$}
\AXC{$\stackrel{\Pi_{58}}{\vdots}$}\noLine
\UIC{$l \cdot a_{1_{\omega}} = l \cdot a_{1_{\omega}} \to \varphi_{146}$}
\BIC{$\varphi_{146}$}
\end{ptree}
\noindent$\Pi_{60}$:
\begin{ptree}
\AXC{$y = z \to x = y \to x = z$\lbl{A6}}
\AXC{$(y = z \to x = y \to x = z) \to x = y \to y = z \to x = z$\lbl{A1}}
\BIC{$x = y \to y = z \to x = z$}
\UIC{$\forall z\,(x = y \to y = z \to x = z)$}
\UIC{$\forall y\,\forall z\,(x = y \to y = z \to x = z)$}
\UIC{$\forall x\,\forall y\,\forall z\,(x = y \to y = z \to x = z)$}
\AXC{$\forall x\,\forall y\,\forall z\,(x = y \to y = z \to x = z) \to \varphi_{25}$\lbl{A2}}
\BIC{$\varphi_{25}$}
\end{ptree}
\noindent$\Pi_{61}$:
\begin{ptree}
\AXC{$\varphi_{84}$\lbl{A2}}
\AXC{$\varphi_{70}$\lbl{[S3]}}
\AXC{$\varphi_{70} \to \forall a\,(l \in \mathbb{R} \to \varphi_{26})$\lbl{A2}}
\BIC{$\forall a\,(l \in \mathbb{R} \to \varphi_{26})$}
\AXC{$\forall a\,(l \in \mathbb{R} \to \varphi_{26}) \to l \in \mathbb{R} \to \varphi_{26}$\lbl{A2}}
\BIC{$l \in \mathbb{R} \to \varphi_{26}$}
\AXC{$\varphi_{137}$\lbl{A1}}
\BIC{$l \in \mathbb{R} \to \mathrm{Seq}(a) \to \forall n \in \mathbb{N}\,(l \odot a(n) = l \cdot a_{n})$}
\AXC{$\varphi_{170}$\lbl{A1}}
\BIC{$\varphi_{84} \to l \in \mathbb{R} \to \varphi_{19}$}
\BIC{$l \in \mathbb{R} \to \varphi_{19}$}
\end{ptree}
\noindent$\Pi_{62}$:
\begin{ptree}
\AXC{$\varphi_{70}$\lbl{[S3]}}
\AXC{$\varphi_{70} \to \forall a\,(l \in \mathbb{R} \to \varphi_{26})$\lbl{A2}}
\BIC{$\forall a\,(l \in \mathbb{R} \to \varphi_{26})$}
\AXC{$\forall a\,(l \in \mathbb{R} \to \varphi_{26}) \to l \in \mathbb{R} \to \varphi_{26}$\lbl{A2}}
\BIC{$l \in \mathbb{R} \to \varphi_{26}$}
\end{ptree}
\noindent$\Pi_{63}$:
\begin{ptree}
\AXC{$\forall a\,(\mathrm{Seq}(a) \to \varphi_{96})$\lbl{[S25]}}
\AXC{$\forall a\,(\mathrm{Seq}(a) \to \varphi_{96}) \to \mathrm{Seq}(l \odot a) \to \varphi_{138}$\lbl{A2}}
\BIC{$\mathrm{Seq}(l \odot a) \to \varphi_{138}$}
\AXC{$(\mathrm{Seq}(l \odot a) \to \varphi_{138}) \to \varphi_{160}$\lbl{A1}}
\BIC{$\varphi_{160}$}
\end{ptree}
\noindent$\Pi_{64}$:
\begin{ptree}
\AXC{$\stackrel{\Pi_{62}}{\vdots}$}\noLine
\UIC{$l \in \mathbb{R} \to \varphi_{26}$}
\AXC{$(l \in \mathbb{R} \to \varphi_{26}) \to l \in \mathbb{R} \to \mathrm{Seq}(a) \to \mathrm{Seq}(l \odot a)$\lbl{A1}}
\BIC{$l \in \mathbb{R} \to \mathrm{Seq}(a) \to \mathrm{Seq}(l \odot a)$}
\AXC{$\stackrel{\Pi_{63}}{\vdots}$}\noLine
\UIC{$\varphi_{160}$}
\BIC{$l \in \mathbb{R} \to \mathrm{Seq}(a) \to \varphi_{138}$}
\AXC{$(l \in \mathbb{R} \to \mathrm{Seq}(a) \to \varphi_{138}) \to \varphi_{23}$\lbl{A1}}
\BIC{$\varphi_{23}$}
\end{ptree}
\noindent$\Pi_{65}$:
\begin{ptree}
\AXC{$1_{\omega} \in \mathbb{N}$\lbl{[N1]}}
\AXC{$\stackrel{\Pi_{61}}{\vdots}$}\noLine
\UIC{$l \in \mathbb{R} \to \varphi_{19}$}
\AXC{$(l \in \mathbb{R} \to \varphi_{19}) \to 1_{\omega} \in \mathbb{N} \to \varphi_{14}$\lbl{A1}}
\BIC{$1_{\omega} \in \mathbb{N} \to \varphi_{14}$}
\BIC{$\varphi_{14}$}
\AXC{$\stackrel{\Pi_{64}}{\vdots}$}\noLine
\UIC{$\varphi_{23}$}
\AXC{$\varphi_{23} \to \varphi_{14} \to \varphi_{154}$\lbl{A1}}
\BIC{$\varphi_{14} \to \varphi_{154}$}
\BIC{$\varphi_{154}$}
\end{ptree}
\noindent$\Pi_{66}$:
\begin{ptree}
\AXC{$\stackrel{\Pi_{60}}{\vdots}$}\noLine
\UIC{$\varphi_{25}$}
\AXC{$\varphi_{25} \to \varphi_{89}$\lbl{A2}}
\BIC{$\varphi_{89}$}
\AXC{$\varphi_{89} \to {\textstyle\sum_{k=1}^{1_{\omega}} l \odot a(k)} = l \odot a(1_{\omega}) \to \varphi_{47}$\lbl{A2}}
\BIC{${\textstyle\sum_{k=1}^{1_{\omega}} l \odot a(k)} = l \odot a(1_{\omega}) \to \varphi_{47}$}
\AXC{$\stackrel{\Pi_{65}}{\vdots}$}\noLine
\UIC{$\varphi_{154}$}
\BIC{$l \in \mathbb{R} \to \mathrm{Seq}(a) \to {\textstyle\sum_{k=1}^{1_{\omega}} l \odot a(k)} = l \cdot a_{1_{\omega}}$}
\AXC{$\varphi_{181}$\lbl{A1}}
\BIC{$\varphi_{174}$}
\end{ptree}
\noindent$\Pi_{67}$:
\begin{ptree}
\AXC{$\varphi_{9}$\lbl{A6}}
\UIC{$\forall w\,(\varphi_{9})$}
\UIC{$\forall z\,\forall w\,(\varphi_{9})$}
\AXC{$\forall z\,\forall w\,(\varphi_{9}) \to \varphi_{40}$\lbl{A2}}
\BIC{$\varphi_{40}$}
\AXC{$\varphi_{40} \to a_{1_{\omega}} = {\textstyle\sum_{k=1}^{1_{\omega}} a_{k}} \to \varphi_{31}$\lbl{A2}}
\BIC{$a_{1_{\omega}} = {\textstyle\sum_{k=1}^{1_{\omega}} a_{k}} \to \varphi_{31}$}
\AXC{$\stackrel{\Pi_{59}}{\vdots}$}\noLine
\UIC{$\varphi_{146}$}
\BIC{$l \in \mathbb{R} \to \mathrm{Seq}(a) \to l \cdot a_{1_{\omega}} = l \cdot {\textstyle\sum_{k=1}^{1_{\omega}} a_{k}}$}
\AXC{$\stackrel{\Pi_{66}}{\vdots}$}\noLine
\UIC{$\varphi_{174}$}
\BIC{$\varphi_{156}$}
\end{ptree}
\noindent$\Pi_{68}$:
\begin{ptree}
\AXC{$y = z \to x = y \to x = z$\lbl{A6}}
\AXC{$(y = z \to x = y \to x = z) \to x = y \to y = z \to x = z$\lbl{A1}}
\BIC{$x = y \to y = z \to x = z$}
\UIC{$\forall z\,(x = y \to y = z \to x = z)$}
\UIC{$\forall y\,\forall z\,(x = y \to y = z \to x = z)$}
\UIC{$\forall x\,\forall y\,\forall z\,(x = y \to y = z \to x = z)$}
\AXC{$\varphi_{7}$\lbl{A2}}
\BIC{$\forall y\,\forall z\,(S(n) = y \to y = z \to S(n) = z)$}
\AXC{$\varphi_{29}$\lbl{A2}}
\BIC{$\forall z\,(S(n) = S(\emptyset) \to S(\emptyset) = z \to S(n) = z)$}
\end{ptree}
\noindent$\Pi_{69}$:
\begin{ptree}
\AXC{$\stackrel{\Pi_{68}}{\vdots}$}\noLine
\UIC{$\forall z\,(S(n) = S(\emptyset) \to S(\emptyset) = z \to S(n) = z)$}
\AXC{$\forall z\,(S(n) = S(\emptyset) \to S(\emptyset) = z \to S(n) = z) \to \varphi_{4}$\lbl{A2}}
\BIC{$\varphi_{4}$}
\end{ptree}
\noindent$\Pi_{70}$:
\begin{ptree}
\AXC{$\stackrel{\Pi_{6}}{\vdots}$}\noLine
\UIC{$\forall x\,\forall y\,(x = y \to y = x)$}
\AXC{$\varphi_{2}$\lbl{A2}}
\BIC{$\forall y\,(1_{\omega} = y \to y = 1_{\omega})$}
\AXC{$\varphi_{10}$\lbl{A2}}
\BIC{$1_{\omega} = S(\emptyset) \to S(\emptyset) = 1_{\omega}$}
\AXC{$1_{\omega} = S(\emptyset)$\lbl{Def$1_\omega$}}
\AXC{$1_{\omega} = S(\emptyset) \to \varphi_{8}$\lbl{A1}}
\BIC{$\varphi_{8}$}
\BIC{$S(\emptyset) = 1_{\omega}$}
\end{ptree}
\noindent$\Pi_{71}$:
\begin{ptree}
\AXC{$x = x$\lbl{A5}}
\UIC{$\forall x\,(x = x)$}
\AXC{$\forall x\,(x = x) \to S(n) = S(n)$\lbl{A2}}
\BIC{$S(n) = S(n)$}
\AXC{$l \in \mathbb{R} \to \varphi_{125}$\lbl{A1}}
\AXC{$(l \in \mathbb{R} \to \varphi_{125}) \to S(n) = S(n) \to \varphi_{157}$\lbl{A1}}
\BIC{$S(n) = S(n) \to \varphi_{157}$}
\BIC{$\varphi_{157}$}
\end{ptree}
\noindent$\Pi_{72}$:
\begin{ptree}
\AXC{$z = w \to S(n) = S(z) \to S(n) = S(w)$\lbl{A6}}
\UIC{$\forall w\,(z = w \to S(n) = S(z) \to S(n) = S(w))$}
\UIC{$\forall z\,\forall w\,(z = w \to S(n) = S(z) \to S(n) = S(w))$}
\AXC{$\varphi_{22}$\lbl{A2}}
\BIC{$\forall w\,(n = w \to S(n) = S(n) \to S(n) = S(w))$}
\AXC{$\varphi_{30}$\lbl{A2}}
\BIC{$n = \emptyset \to S(n) = S(n) \to S(n) = S(\emptyset)$}
\AXC{$\stackrel{\Pi_{71}}{\vdots}$}\noLine
\UIC{$\varphi_{157}$}
\BIC{$l \in \mathbb{R} \to \mathrm{Seq}(a) \to \varphi_{121}$}
\AXC{$\varphi_{192}$\lbl{A1}}
\BIC{$S(\emptyset) = 1_{\omega} \to \varphi_{4} \to l \in \mathbb{R} \to \mathrm{Seq}(a) \to \varphi_{109}$}
\end{ptree}
\noindent$\Pi_{73}$:
\begin{ptree}
\AXC{$\stackrel{\Pi_{69}}{\vdots}$}\noLine
\UIC{$\varphi_{4}$}
\AXC{$\stackrel{\Pi_{70}}{\vdots}$}\noLine
\UIC{$S(\emptyset) = 1_{\omega}$}
\AXC{$\stackrel{\Pi_{72}}{\vdots}$}\noLine
\UIC{$S(\emptyset) = 1_{\omega} \to \varphi_{4} \to l \in \mathbb{R} \to \mathrm{Seq}(a) \to \varphi_{109}$}
\BIC{$\varphi_{4} \to l \in \mathbb{R} \to \mathrm{Seq}(a) \to \varphi_{109}$}
\BIC{$l \in \mathbb{R} \to \mathrm{Seq}(a) \to \varphi_{109}$}
\AXC{$(l \in \mathbb{R} \to \mathrm{Seq}(a) \to \varphi_{109}) \to \varphi_{148}$\lbl{A1}}
\BIC{$\varphi_{148}$}
\end{ptree}
\noindent$\Pi_{74}$:
\begin{ptree}
\AXC{$\stackrel{\Pi_{6}}{\vdots}$}\noLine
\UIC{$\forall x\,\forall y\,(x = y \to y = x)$}
\AXC{$\varphi_{1}$\lbl{A2}}
\BIC{$\forall y\,(S(n) = y \to y = S(n))$}
\AXC{$\varphi_{3}$\lbl{A2}}
\BIC{$S(n) = 1_{\omega} \to 1_{\omega} = S(n)$}
\AXC{$\stackrel{\Pi_{73}}{\vdots}$}\noLine
\UIC{$\varphi_{148}$}
\BIC{$l \in \mathbb{R} \to \mathrm{Seq}(a) \to \varphi_{110}$}
\AXC{$(l \in \mathbb{R} \to \mathrm{Seq}(a) \to \varphi_{110}) \to \varphi_{189}$\lbl{A1}}
\BIC{$\varphi_{189}$}
\end{ptree}
\noindent$\Pi_{75}$:
\begin{ptree}
\AXC{$\stackrel{\Pi_{55}}{\vdots}$}\noLine
\UIC{$\varphi_{25}$}
\AXC{$\varphi_{25} \to \varphi_{78}$\lbl{A2}}
\BIC{$\varphi_{78}$}
\AXC{$\varphi_{78} \to {\textstyle\sum_{k=1}^{1_{\omega}} l \odot a(k)} = l \cdot a_{1_{\omega}} \to \varphi_{71}$\lbl{A2}}
\BIC{${\textstyle\sum_{k=1}^{1_{\omega}} l \odot a(k)} = l \cdot a_{1_{\omega}} \to \varphi_{71}$}
\AXC{$\stackrel{\Pi_{67}}{\vdots}$}\noLine
\UIC{$\varphi_{156}$}
\BIC{$l \in \mathbb{R} \to \mathrm{Seq}(a) \to {\textstyle\sum_{k=1}^{1_{\omega}} l \odot a(k)} = l \cdot {\textstyle\sum_{k=1}^{1_{\omega}} a_{k}}$}
\AXC{$\stackrel{\Pi_{74}}{\vdots}$}\noLine
\UIC{$\varphi_{189}$}
\BIC{$\varphi_{177}$}
\end{ptree}
\noindent$\Pi_{76}$:
\begin{ptree}
\AXC{$z = w \to \varphi_{58}$\lbl{A6}}
\UIC{$\forall w\,(z = w \to \varphi_{58})$}
\UIC{$\forall z\,\forall w\,(z = w \to \varphi_{58})$}
\AXC{$\forall z\,\forall w\,(z = w \to \varphi_{58}) \to \varphi_{90}$\lbl{A2}}
\BIC{$\varphi_{90}$}
\AXC{$\varphi_{90} \to 1_{\omega} = S(n) \to \varphi_{85}$\lbl{A2}}
\BIC{$1_{\omega} = S(n) \to \varphi_{85}$}
\AXC{$\stackrel{\Pi_{75}}{\vdots}$}\noLine
\UIC{$\varphi_{177}$}
\BIC{$l \in \mathbb{R} \to \mathrm{Seq}(a) \to \forall n\,(\varphi_{43}) \to \varphi_{67}$}
\end{ptree}
\noindent$\Pi_{77}$:
\begin{ptree}
\AXC{$\stackrel{\Pi_{54}}{\vdots}$}\noLine
\UIC{$l \in \mathbb{R} \to \mathrm{Seq}(a) \to \forall n\,(\varphi_{43}) \to \varphi_{69}$}
\AXC{$\stackrel{\Pi_{76}}{\vdots}$}\noLine
\UIC{$l \in \mathbb{R} \to \mathrm{Seq}(a) \to \forall n\,(\varphi_{43}) \to \varphi_{67}$}
\AXC{$(l \in \mathbb{R} \to \mathrm{Seq}(a) \to \forall n\,(\varphi_{43}) \to \varphi_{67}) \to \varphi_{190}$\lbl{A1}}
\BIC{$\varphi_{190}$}
\BIC{$l \in \mathbb{R} \to \mathrm{Seq}(a) \to \varphi_{141}$}
\end{ptree}
\noindent$\Pi_{78}$:
\begin{ptree}
\AXC{$\forall x\,(\mathrm{Ind}(x) \leftrightarrow \emptyset \in x \land \forall y \in x\,(S(y) \in x))$\lbl{DefInd}}
\AXC{$\varphi_{99}$\lbl{A2}}
\BIC{$\mathrm{Ind}(\omega) \leftrightarrow \emptyset \in \omega \land \forall y \in \omega\,(S(y) \in \omega)$}
\AXC{$\mathrm{Ind}(\omega)$\lbl{補題(ωは帰納的)}}
\AXC{$\varphi_{169}$\lbl{A1}}
\BIC{$(\mathrm{Ind}(\omega) \leftrightarrow \emptyset \in \omega \land \forall y \in \omega\,(S(y) \in \omega)) \to \varphi_{140}$}
\BIC{$\varphi_{140}$}
\end{ptree}
\noindent$\Pi_{79}$:
\begin{ptree}
\AXC{$\forall y \in \omega\,(S(y) \in \omega) \to n \in \omega \to S(n) \in \omega$\lbl{A2}}
\AXC{$\stackrel{\Pi_{78}}{\vdots}$}\noLine
\UIC{$\varphi_{140}$}
\AXC{$\varphi_{140} \to l \in \mathbb{R} \to \varphi_{114}$\lbl{A1}}
\BIC{$l \in \mathbb{R} \to \varphi_{114}$}
\AXC{$\varphi_{185}$\lbl{A1}}
\BIC{$(\forall y \in \omega\,(S(y) \in \omega) \to n \in \omega \to S(n) \in \omega) \to \varphi_{130}$}
\BIC{$\varphi_{130}$}
\AXC{$\varphi_{130} \to \varphi_{173}$\lbl{A1}}
\BIC{$\varphi_{173}$}
\end{ptree}
\noindent$\Pi_{80}$:
\begin{ptree}
\AXC{$\forall n\,(\varphi_{43}) \to \varphi_{77}$\lbl{A2}}
\AXC{$l \in \mathbb{R} \to \mathrm{Seq}(a) \to \forall n\,(\varphi_{43}) \to \forall n\,(\varphi_{43})$\lbl{A1}}
\AXC{$\varphi_{207}$\lbl{A1}}
\BIC{$(\forall n\,(\varphi_{43}) \to \varphi_{77}) \to l \in \mathbb{R} \to \mathrm{Seq}(a) \to \forall n\,(\varphi_{43}) \to \varphi_{77}$}
\BIC{$l \in \mathbb{R} \to \mathrm{Seq}(a) \to \forall n\,(\varphi_{43}) \to \varphi_{77}$}
\end{ptree}
\noindent$\Pi_{81}$:
\begin{ptree}
\AXC{$\stackrel{\Pi_{80}}{\vdots}$}\noLine
\UIC{$l \in \mathbb{R} \to \mathrm{Seq}(a) \to \forall n\,(\varphi_{43}) \to \varphi_{77}$}
\AXC{$(l \in \mathbb{R} \to \mathrm{Seq}(a) \to \forall n\,(\varphi_{43}) \to \varphi_{77}) \to \varphi_{196}$\lbl{A1}}
\BIC{$\varphi_{196}$}
\end{ptree}
\noindent$\Pi_{82}$:
\begin{ptree}
\AXC{$\stackrel{\Pi_{77}}{\vdots}$}\noLine
\UIC{$l \in \mathbb{R} \to \mathrm{Seq}(a) \to \varphi_{141}$}
\AXC{$\stackrel{\Pi_{10}}{\vdots}$}\noLine
\UIC{$l \in \mathbb{R} \to \mathrm{Seq}(a) \to \forall n\,(\varphi_{43}) \to n \in B_{47} \to n \in \omega$}
\AXC{$\stackrel{\Pi_{79}}{\vdots}$}\noLine
\UIC{$\varphi_{173}$}
\BIC{$l \in \mathbb{R} \to \mathrm{Seq}(a) \to \forall n\,(\varphi_{43}) \to n \in B_{47} \to S(n) \in \omega$}
\AXC{$\stackrel{\Pi_{81}}{\vdots}$}\noLine
\UIC{$\varphi_{196}$}
\BIC{$(l \in \mathbb{R} \to \mathrm{Seq}(a) \to \varphi_{141}) \to \varphi_{106}$}
\BIC{$\varphi_{106}$}
\AXC{$\varphi_{106} \to \varphi_{107}$\lbl{A1}}
\BIC{$\varphi_{107}$}
\UIC{$\forall n\,(\varphi_{107})$}
\end{ptree}
\noindent$\Pi_{83}$:
\begin{ptree}
\AXC{$\forall n\,(n \in \mathbb{N} \leftrightarrow n \in \omega \land n \neq \emptyset)$\lbl{Def$\mathbb{N}$}}
\AXC{$\varphi_{88}$\lbl{A2}}
\BIC{$\emptyset \in \mathbb{N} \leftrightarrow \emptyset \in \omega \land \emptyset \neq \emptyset$}
\AXC{$(\emptyset \in \mathbb{N} \leftrightarrow \emptyset \in \omega \land \emptyset \neq \emptyset) \to \emptyset = \emptyset \to \varphi_{80}$\lbl{A1}}
\BIC{$\emptyset = \emptyset \to \varphi_{80}$}
\end{ptree}
\noindent$\Pi_{84}$:
\begin{ptree}
\AXC{$\stackrel{\Pi_{78}}{\vdots}$}\noLine
\UIC{$\varphi_{140}$}
\AXC{$\forall n\,(\varphi_{43}) \to \varphi_{98}$\lbl{A2}}
\AXC{$l \in \mathbb{R} \to \mathrm{Seq}(a) \to \forall n\,(\varphi_{43}) \to \forall n\,(\varphi_{43})$\lbl{A1}}
\AXC{$\varphi_{209}$\lbl{A1}}
\BIC{$(\forall n\,(\varphi_{43}) \to \varphi_{98}) \to l \in \mathbb{R} \to \mathrm{Seq}(a) \to \forall n\,(\varphi_{43}) \to \varphi_{98}$}
\BIC{$l \in \mathbb{R} \to \mathrm{Seq}(a) \to \forall n\,(\varphi_{43}) \to \varphi_{98}$}
\AXC{$\varphi_{208}$\lbl{A1}}
\BIC{$\varphi_{140} \to \varphi_{80} \to l \in \mathbb{R} \to \mathrm{Seq}(a) \to \forall n\,(\varphi_{43}) \to \emptyset \in B_{47}$}
\BIC{$\varphi_{80} \to l \in \mathbb{R} \to \mathrm{Seq}(a) \to \forall n\,(\varphi_{43}) \to \emptyset \in B_{47}$}
\end{ptree}
\noindent$\Pi_{85}$:
\begin{ptree}
\AXC{$x = x$\lbl{A5}}
\UIC{$\forall x\,(x = x)$}
\AXC{$\forall x\,(x = x) \to \emptyset = \emptyset$\lbl{A2}}
\BIC{$\emptyset = \emptyset$}
\AXC{$\stackrel{\Pi_{83}}{\vdots}$}\noLine
\UIC{$\emptyset = \emptyset \to \varphi_{80}$}
\BIC{$\varphi_{80}$}
\AXC{$\stackrel{\Pi_{84}}{\vdots}$}\noLine
\UIC{$\varphi_{80} \to l \in \mathbb{R} \to \mathrm{Seq}(a) \to \forall n\,(\varphi_{43}) \to \emptyset \in B_{47}$}
\BIC{$l \in \mathbb{R} \to \mathrm{Seq}(a) \to \forall n\,(\varphi_{43}) \to \emptyset \in B_{47}$}
\end{ptree}
\noindent$\Pi_{86}$:
\begin{ptree}
\AXC{$\forall x\,(\mathrm{Ind}(x) \leftrightarrow \emptyset \in x \land \forall y \in x\,(S(y) \in x))$\lbl{DefInd}}
\AXC{$\forall x\,(\mathrm{Ind}(x) \leftrightarrow \emptyset \in x \land \forall y \in x\,(S(y) \in x)) \to \varphi_{13}$\lbl{A2}}
\BIC{$\varphi_{13}$}
\end{ptree}
\noindent$\Pi_{87}$:
\begin{ptree}
\AXC{$\stackrel{\Pi_{85}}{\vdots}$}\noLine
\UIC{$l \in \mathbb{R} \to \mathrm{Seq}(a) \to \forall n\,(\varphi_{43}) \to \emptyset \in B_{47}$}
\AXC{$\stackrel{\Pi_{86}}{\vdots}$}\noLine
\UIC{$\varphi_{13}$}
\AXC{$\varphi_{13} \to (l \in \mathbb{R} \to \mathrm{Seq}(a) \to \forall n\,(\varphi_{43}) \to \emptyset \in B_{47}) \to \varphi_{187}$\lbl{A1}}
\BIC{$(l \in \mathbb{R} \to \mathrm{Seq}(a) \to \forall n\,(\varphi_{43}) \to \emptyset \in B_{47}) \to \varphi_{187}$}
\BIC{$\varphi_{187}$}
\end{ptree}
\noindent$\Pi_{88}$:
\begin{ptree}
\AXC{$\varphi_{159}$\lbl{A4}}
\AXC{$\forall n\,(\varphi_{107}) \to \varphi_{113}$\lbl{A3}}
\AXC{$\stackrel{\Pi_{82}}{\vdots}$}\noLine
\UIC{$\forall n\,(\varphi_{107})$}
\AXC{$\forall n\,(\varphi_{107}) \to (\forall n\,(\varphi_{107}) \to \varphi_{113}) \to \varphi_{159} \to \varphi_{108}$\lbl{A1}}
\BIC{$(\forall n\,(\varphi_{107}) \to \varphi_{113}) \to \varphi_{159} \to \varphi_{108}$}
\BIC{$\varphi_{159} \to \varphi_{108}$}
\BIC{$\varphi_{108}$}
\AXC{$\stackrel{\Pi_{87}}{\vdots}$}\noLine
\UIC{$\varphi_{187}$}
\BIC{$\varphi_{12} \to \varphi_{11} \to l \in \mathbb{R} \to \mathrm{Seq}(a) \to \forall n\,(\varphi_{43}) \to \mathrm{Ind}(B_{47})$}
\end{ptree}
\noindent$\Pi_{89}$:
\begin{ptree}
\AXC{$\forall n\,(n \in \omega \leftrightarrow \forall x\,(\mathrm{Ind}(x) \to n \in x))$\lbl{Def$\omega$}}
\AXC{$\varphi_{27}$\lbl{A2}}
\BIC{$n \in \omega \leftrightarrow \forall x\,(\mathrm{Ind}(x) \to n \in x)$}
\AXC{$l \in \mathbb{R} \to \mathrm{Seq}(a) \to \forall n\,(\varphi_{43}) \to n \in \omega \to n \in \omega$\lbl{A1}}
\AXC{$\varphi_{178}$\lbl{A1}}
\BIC{$(n \in \omega \leftrightarrow \forall x\,(\mathrm{Ind}(x) \to n \in x)) \to l \in \mathbb{R} \to \varphi_{102}$}
\BIC{$l \in \mathbb{R} \to \varphi_{102}$}
\end{ptree}
\noindent$\Pi_{90}$:
\begin{ptree}
\AXC{$\forall x\,(\mathrm{Ind}(x) \to n \in x) \to \mathrm{Ind}(B_{47}) \to n \in B_{47}$\lbl{A2}}
\AXC{$\stackrel{\Pi_{89}}{\vdots}$}\noLine
\UIC{$l \in \mathbb{R} \to \varphi_{102}$}
\AXC{$(l \in \mathbb{R} \to \varphi_{102}) \to \varphi_{152}$\lbl{A1}}
\BIC{$\varphi_{152}$}
\BIC{$l \in \mathbb{R} \to \varphi_{104}$}
\end{ptree}
\noindent$\Pi_{91}$:
\begin{ptree}
\AXC{$\stackrel{\Pi_{3}}{\vdots}$}\noLine
\UIC{$\varphi_{11}$}
\AXC{$\stackrel{\Pi_{4}}{\vdots}$}\noLine
\UIC{$\varphi_{12}$}
\AXC{$\stackrel{\Pi_{88}}{\vdots}$}\noLine
\UIC{$\varphi_{12} \to \varphi_{11} \to l \in \mathbb{R} \to \mathrm{Seq}(a) \to \forall n\,(\varphi_{43}) \to \mathrm{Ind}(B_{47})$}
\BIC{$\varphi_{11} \to l \in \mathbb{R} \to \mathrm{Seq}(a) \to \forall n\,(\varphi_{43}) \to \mathrm{Ind}(B_{47})$}
\BIC{$l \in \mathbb{R} \to \mathrm{Seq}(a) \to \forall n\,(\varphi_{43}) \to \mathrm{Ind}(B_{47})$}
\AXC{$\stackrel{\Pi_{90}}{\vdots}$}\noLine
\UIC{$l \in \mathbb{R} \to \varphi_{104}$}
\AXC{$(l \in \mathbb{R} \to \varphi_{104}) \to \varphi_{198}$\lbl{A1}}
\BIC{$\varphi_{198}$}
\BIC{$(l \in \mathbb{R} \to \mathrm{Seq}(a) \to \forall n\,(\varphi_{43}) \to \varphi_{43}) \to l \in \mathbb{R} \to \varphi_{143}$}
\end{ptree}
\noindent$\Pi_{92}$:
\begin{ptree}
\AXC{$\stackrel{\Pi_{2}}{\vdots}$}\noLine
\UIC{$l \in \mathbb{R} \to \mathrm{Seq}(a) \to \forall n\,(\varphi_{43}) \to \varphi_{43}$}
\AXC{$\stackrel{\Pi_{91}}{\vdots}$}\noLine
\UIC{$(l \in \mathbb{R} \to \mathrm{Seq}(a) \to \forall n\,(\varphi_{43}) \to \varphi_{43}) \to l \in \mathbb{R} \to \varphi_{143}$}
\BIC{$l \in \mathbb{R} \to \varphi_{143}$}
\AXC{$(l \in \mathbb{R} \to \varphi_{143}) \to \varphi_{147}$\lbl{A1}}
\BIC{$\varphi_{147}$}
\UIC{$\forall n\,(\varphi_{147})$}
\end{ptree}
\noindent$\Pi_{93}$:
\begin{ptree}
\AXC{$\varphi_{159}$\lbl{A4}}
\AXC{$\forall n\,(\varphi_{147}) \to \varphi_{150}$\lbl{A3}}
\AXC{$\stackrel{\Pi_{92}}{\vdots}$}\noLine
\UIC{$\forall n\,(\varphi_{147})$}
\AXC{$\varphi_{219}$\lbl{A1}}
\BIC{$(\forall n\,(\varphi_{147}) \to \varphi_{150}) \to \varphi_{159} \to l \in \mathbb{R} \to \mathrm{Seq}(a) \to \varphi_{139}$}
\BIC{$\varphi_{159} \to l \in \mathbb{R} \to \mathrm{Seq}(a) \to \varphi_{139}$}
\BIC{$l \in \mathbb{R} \to \mathrm{Seq}(a) \to \varphi_{139}$}
\end{ptree}
\noindent$\Pi_{94}$:
\begin{ptree}
\AXC{$\varphi_{95}$\lbl{Sep}}
\AXC{$\varphi_{95} \to \varphi_{87}$\lbl{A2}}
\BIC{$\varphi_{87}$}
\AXC{$\varphi_{87} \to \varphi_{73}$\lbl{A2}}
\BIC{$\varphi_{73}$}
\AXC{$\varphi_{73} \to \exists B_{47}\,\forall n\,(\varphi_{43})$\lbl{A2}}
\BIC{$\exists B_{47}\,\forall n\,(\varphi_{43})$}
\AXC{$\varphi_{155}$\lbl{A1}}
\BIC{$l \in \mathbb{R} \to \mathrm{Seq}(a) \to \exists B_{47}\,\forall n\,(\varphi_{43})$}
\end{ptree}
\noindent$\Pi_{95}$:
\begin{ptree}
\AXC{$\stackrel{\Pi_{94}}{\vdots}$}\noLine
\UIC{$l \in \mathbb{R} \to \mathrm{Seq}(a) \to \exists B_{47}\,\forall n\,(\varphi_{43})$}
\AXC{$(l \in \mathbb{R} \to \mathrm{Seq}(a) \to \exists B_{47}\,\forall n\,(\varphi_{43})) \to \varphi_{217}$\lbl{A1}}
\BIC{$\varphi_{217}$}
\end{ptree}
\noindent$\Pi_{96}$:
\begin{ptree}
\AXC{$\varphi_{188}$\lbl{A3}}
\AXC{$\stackrel{\Pi_{93}}{\vdots}$}\noLine
\UIC{$l \in \mathbb{R} \to \mathrm{Seq}(a) \to \varphi_{139}$}
\AXC{$(l \in \mathbb{R} \to \mathrm{Seq}(a) \to \varphi_{139}) \to \varphi_{75} \to \neg \forall n\,(\varphi_{43})$\lbl{A1}}
\BIC{$\varphi_{75} \to \neg \forall n\,(\varphi_{43})$}
\UIC{$\forall B_{47}\,(\varphi_{75} \to \neg \forall n\,(\varphi_{43}))$}
\AXC{$\stackrel{\Pi_{95}}{\vdots}$}\noLine
\UIC{$\varphi_{217}$}
\BIC{$\varphi_{188} \to (\varphi_{75} \to \forall B_{47}\,(\varphi_{75})) \to l \in \mathbb{R} \to \varphi_{46}$}
\BIC{$(\varphi_{75} \to \forall B_{47}\,(\varphi_{75})) \to l \in \mathbb{R} \to \varphi_{46}$}
\end{ptree}
\noindent$\Pi_{97}$:
\begin{ptree}
\AXC{$\varphi_{118}$\lbl{A2}}
\AXC{$\varphi_{75} \to \forall B_{47}\,(\varphi_{75})$\lbl{A4}}
\AXC{$\stackrel{\Pi_{96}}{\vdots}$}\noLine
\UIC{$(\varphi_{75} \to \forall B_{47}\,(\varphi_{75})) \to l \in \mathbb{R} \to \varphi_{46}$}
\BIC{$l \in \mathbb{R} \to \varphi_{46}$}
\AXC{$(l \in \mathbb{R} \to \varphi_{46}) \to \varphi_{118} \to l \in \mathbb{R} \to \varphi_{37}$\lbl{A1}}
\BIC{$\varphi_{118} \to l \in \mathbb{R} \to \varphi_{37}$}
\BIC{$l \in \mathbb{R} \to \varphi_{37}$}
\end{ptree}
\noindent$\Pi_{98}$:
\begin{ptree}
\AXC{$\forall n\,(n \in \mathbb{N} \leftrightarrow n \in \omega \land n \neq \emptyset)$\lbl{Def$\mathbb{N}$}}
\AXC{$\varphi_{42}$\lbl{A2}}
\BIC{$n \in \mathbb{N} \leftrightarrow n \in \omega \land n \neq \emptyset$}
\AXC{$\stackrel{\Pi_{1}}{\vdots}$}\noLine
\UIC{$\varphi_{51}$}
\BIC{$l \in \mathbb{R} \to \mathrm{Seq}(a) \to n \in \mathbb{N} \to n \in \omega$}
\AXC{$\stackrel{\Pi_{97}}{\vdots}$}\noLine
\UIC{$l \in \mathbb{R} \to \varphi_{37}$}
\AXC{$(l \in \mathbb{R} \to \varphi_{37}) \to \varphi_{144}$\lbl{A1}}
\BIC{$\varphi_{144}$}
\BIC{$(l \in \mathbb{R} \to \mathrm{Seq}(a) \to n \in \mathbb{N} \to n \in \mathbb{N}) \to \varphi_{33}$}
\end{ptree}
\noindent$\Pi_{99}$:
\begin{ptree}
\AXC{$l \in \mathbb{R} \to \mathrm{Seq}(a) \to n \in \mathbb{N} \to n \in \mathbb{N}$\lbl{A1}}
\AXC{$\stackrel{\Pi_{98}}{\vdots}$}\noLine
\UIC{$(l \in \mathbb{R} \to \mathrm{Seq}(a) \to n \in \mathbb{N} \to n \in \mathbb{N}) \to \varphi_{33}$}
\BIC{$\varphi_{33}$}
\AXC{$\varphi_{33} \to \varphi_{34}$\lbl{A1}}
\BIC{$\varphi_{34}$}
\UIC{$\forall n\,(\varphi_{34})$}
\end{ptree}
\noindent$\Pi_{100}$:
\begin{ptree}
\AXC{$\forall n\,(\varphi_{34}) \to \varphi_{48}$\lbl{A3}}
\AXC{$\stackrel{\Pi_{99}}{\vdots}$}\noLine
\UIC{$\forall n\,(\varphi_{34})$}
\AXC{$\forall n\,(\varphi_{34}) \to (\forall n\,(\varphi_{34}) \to \varphi_{48}) \to \varphi_{101}$\lbl{A1}}
\BIC{$(\forall n\,(\varphi_{34}) \to \varphi_{48}) \to \varphi_{101}$}
\BIC{$\varphi_{101}$}
\end{ptree}
\noindent 主木:
\begin{ptree}
\AXC{$l \in \mathbb{R} \land \mathrm{Seq}(a) \to \forall n\,(l \in \mathbb{R} \land \mathrm{Seq}(a))$\lbl{A4}}
\AXC{$\stackrel{\Pi_{100}}{\vdots}$}\noLine
\UIC{$\varphi_{101}$}
\BIC{$\varphi_{41}$}
\UIC{$\forall a\,(\varphi_{41})$}
\UIC{$\forall l\,\forall a\,(\varphi_{41})$}
\end{ptree}
\begin{abbrevlist}
\item $\varphi_{1} :\equiv \forall x\,\forall y\,(x = y \to y = x) \to \forall y\,(S(n) = y \to y = S(n))$
\item $\varphi_{2} :\equiv \forall x\,\forall y\,(x = y \to y = x) \to \forall y\,(1_{\omega} = y \to y = 1_{\omega})$
\item $\varphi_{3} :\equiv \forall y\,(S(n) = y \to y = S(n)) \to S(n) = 1_{\omega} \to 1_{\omega} = S(n)$
\item $\varphi_{4} :\equiv S(n) = S(\emptyset) \to S(\emptyset) = 1_{\omega} \to S(n) = 1_{\omega}$
\item $\varphi_{5} :\equiv \forall n \in B_{47}\,(S(n) \in B_{47}) \to y \in B_{47} \to S(y) \in B_{47}$
\item $\varphi_{6} :\equiv \forall y \in B_{47}\,(S(y) \in B_{47}) \to n \in B_{47} \to S(n) \in B_{47}$
\item $\varphi_{7} :\equiv \forall x\,\forall y\,\forall z\,(x = y \to y = z \to x = z) \to \forall y\,\forall z\,(S(n) = y \to y = z \to S(n) = z)$
\item $\varphi_{8} :\equiv (1_{\omega} = S(\emptyset) \to S(\emptyset) = 1_{\omega}) \to S(\emptyset) = 1_{\omega}$
\item $\varphi_{9} :\equiv z = w \to l \cdot a_{1_{\omega}} = l \cdot z \to l \cdot a_{1_{\omega}} = l \cdot w$
\item $\varphi_{10} :\equiv \forall y\,(1_{\omega} = y \to y = 1_{\omega}) \to 1_{\omega} = S(\emptyset) \to S(\emptyset) = 1_{\omega}$
\item $\varphi_{11} :\equiv \forall n \in B_{47}\,(S(n) \in B_{47}) \to \forall y \in B_{47}\,(S(y) \in B_{47})$
\item $\varphi_{12} :\equiv \forall y \in B_{47}\,(S(y) \in B_{47}) \to \forall n \in B_{47}\,(S(n) \in B_{47})$
\item $\varphi_{13} :\equiv \mathrm{Ind}(B_{47}) \leftrightarrow \emptyset \in B_{47} \land \forall y \in B_{47}\,(S(y) \in B_{47})$
\item $\varphi_{14} :\equiv l \in \mathbb{R} \to \mathrm{Seq}(a) \to l \odot a(1_{\omega}) = l \cdot a_{1_{\omega}}$
\item $\varphi_{15} :\equiv \forall n \in B_{47}\,(S(n) \in B_{47}) \to \forall y\,\forall n \in B_{47}\,(S(n) \in B_{47})$
\item $\varphi_{16} :\equiv \forall y \in B_{47}\,(S(y) \in B_{47}) \to \forall n\,\forall y \in B_{47}\,(S(y) \in B_{47})$
\item $\varphi_{17} :\equiv \forall y\,\forall n \in B_{47}\,(S(n) \in B_{47}) \to \forall y \in B_{47}\,(S(y) \in B_{47})$
\item $\varphi_{18} :\equiv \forall n\,\forall y \in B_{47}\,(S(y) \in B_{47}) \to \forall n \in B_{47}\,(S(n) \in B_{47})$
\item $\varphi_{19} :\equiv \mathrm{Seq}(a) \to 1_{\omega} \in \mathbb{N} \to l \odot a(1_{\omega}) = l \cdot a_{1_{\omega}}$
\item $\varphi_{20} :\equiv \forall x\,\forall y\,(x = y \to y = x) \to \forall y\,({\textstyle\sum_{k=1}^{S(n)} a_{k}} = y \to y = {\textstyle\sum_{k=1}^{S(n)} a_{k}})$
\item $\varphi_{21} :\equiv \forall x\,\forall y\,(x = y \to y = x) \to \forall y\,({\textstyle\sum_{k=1}^{1_{\omega}} a_{k}} = y \to y = {\textstyle\sum_{k=1}^{1_{\omega}} a_{k}})$
\item $\varphi_{22} :\equiv \forall z\,\forall w\,(z = w \to S(n) = S(z) \to S(n) = S(w)) \to \forall w\,(n = w \to S(n) = S(n) \to S(n) = S(w))$
\item $\varphi_{23} :\equiv l \in \mathbb{R} \to \mathrm{Seq}(a) \to {\textstyle\sum_{k=1}^{1_{\omega}} l \odot a(k)} = l \odot a(1_{\omega})$
\item $\varphi_{24} :\equiv \forall y\,\forall z\,({\textstyle\sum_{k=1}^{S(n)} l \odot a(k)} = y \to y = z \to {\textstyle\sum_{k=1}^{S(n)} l \odot a(k)} = z)$
\item $\varphi_{25} :\equiv \forall y\,\forall z\,({\textstyle\sum_{k=1}^{1_{\omega}} l \odot a(k)} = y \to y = z \to {\textstyle\sum_{k=1}^{1_{\omega}} l \odot a(k)} = z)$
\item $\varphi_{26} :\equiv \mathrm{Seq}(a) \to \mathrm{Seq}(l \odot a) \land \forall n \in \mathbb{N}\,(l \odot a(n) = l \cdot a_{n})$
\item $\varphi_{27} :\equiv \forall n\,(n \in \omega \leftrightarrow \forall x\,(\mathrm{Ind}(x) \to n \in x)) \to (n \in \omega \leftrightarrow \forall x\,(\mathrm{Ind}(x) \to n \in x))$
\item $\varphi_{28} :\equiv \forall y\,\forall z\,(l \in \mathbb{R} \to y \in \mathbb{R} \to z \in \mathbb{R} \to l \cdot (y + z) = l \cdot y + l \cdot z)$
\item $\varphi_{29} :\equiv \forall y\,\forall z\,(S(n) = y \to y = z \to S(n) = z) \to \forall z\,(S(n) = S(\emptyset) \to S(\emptyset) = z \to S(n) = z)$
\item $\varphi_{30} :\equiv \forall w\,(n = w \to S(n) = S(n) \to S(n) = S(w)) \to n = \emptyset \to S(n) = S(n) \to S(n) = S(\emptyset)$
\item $\varphi_{31} :\equiv l \cdot a_{1_{\omega}} = l \cdot a_{1_{\omega}} \to l \cdot a_{1_{\omega}} = l \cdot {\textstyle\sum_{k=1}^{1_{\omega}} a_{k}}$
\item $\varphi_{32} :\equiv \forall x\,\forall y\,\forall z\,(x \in \mathbb{R} \to y \in \mathbb{R} \to z \in \mathbb{R} \to x \cdot (y + z) = x \cdot y + x \cdot z)$
\item $\varphi_{33} :\equiv l \in \mathbb{R} \to \mathrm{Seq}(a) \to n \in \mathbb{N} \to {\textstyle\sum_{k=1}^{n} l \odot a(k)} = l \cdot {\textstyle\sum_{k=1}^{n} a_{k}}$
\item $\varphi_{34} :\equiv l \in \mathbb{R} \land \mathrm{Seq}(a) \to n \in \mathbb{N} \to {\textstyle\sum_{k=1}^{n} l \odot a(k)} = l \cdot {\textstyle\sum_{k=1}^{n} a_{k}}$
\item $\varphi_{35} :\equiv n \neq \emptyset \to \mathrm{Seq}(l \odot a) \land \forall n \in \mathbb{N}\,(l \odot a(n) = l \cdot a_{n})$
\item $\varphi_{36} :\equiv n \in B_{47} \to S(n) \in \mathbb{N} \to n \neq \emptyset \to l \odot a(S(n)) = l \cdot a_{S(n)}$
\item $\varphi_{37} :\equiv \mathrm{Seq}(a) \to n \in \omega \to n \in \mathbb{N} \to {\textstyle\sum_{k=1}^{n} l \odot a(k)} = l \cdot {\textstyle\sum_{k=1}^{n} a_{k}}$
\item $\varphi_{38} :\equiv \forall x\,(x = x) \to l \cdot ({\textstyle\sum_{k=1}^{n} a_{k}} + a_{S(n)}) = l \cdot ({\textstyle\sum_{k=1}^{n} a_{k}} + a_{S(n)})$
\item $\varphi_{39} :\equiv \forall y\,(l \cdot ({\textstyle\sum_{k=1}^{n} a_{k}} + a_{S(n)}) = y \to y = l \cdot ({\textstyle\sum_{k=1}^{n} a_{k}} + a_{S(n)}))$
\item $\varphi_{40} :\equiv \forall w\,(a_{1_{\omega}} = w \to l \cdot a_{1_{\omega}} = l \cdot a_{1_{\omega}} \to l \cdot a_{1_{\omega}} = l \cdot w)$
\item $\varphi_{41} :\equiv l \in \mathbb{R} \to \mathrm{Seq}(a) \to \forall n \in \mathbb{N}\,({\textstyle\sum_{k=1}^{n} l \odot a(k)} = l \cdot {\textstyle\sum_{k=1}^{n} a_{k}})$
\item $\varphi_{42} :\equiv \forall n\,(n \in \mathbb{N} \leftrightarrow n \in \omega \land n \neq \emptyset) \to (n \in \mathbb{N} \leftrightarrow n \in \omega \land n \neq \emptyset)$
\item $\varphi_{43} :\equiv n \in B_{47} \leftrightarrow n \in \omega \land (n \in \mathbb{N} \to {\textstyle\sum_{k=1}^{n} l \odot a(k)} = l \cdot {\textstyle\sum_{k=1}^{n} a_{k}})$
\item $\varphi_{44} :\equiv n \in B_{47} \to n \in \omega \land (n \in \mathbb{N} \to {\textstyle\sum_{k=1}^{n} l \odot a(k)} = l \cdot {\textstyle\sum_{k=1}^{n} a_{k}})$
\item $\varphi_{45} :\equiv S(n) \in \mathbb{N} \to n \neq \emptyset \to S(n) \in \mathbb{N} \to l \odot a(S(n)) = l \cdot a_{S(n)}$
\item $\varphi_{46} :\equiv \mathrm{Seq}(a) \to \forall n \in \omega\,(n \in \mathbb{N} \to {\textstyle\sum_{k=1}^{n} l \odot a(k)} = l \cdot {\textstyle\sum_{k=1}^{n} a_{k}})$
\item $\varphi_{47} :\equiv l \odot a(1_{\omega}) = l \cdot a_{1_{\omega}} \to {\textstyle\sum_{k=1}^{1_{\omega}} l \odot a(k)} = l \cdot a_{1_{\omega}}$
\item $\varphi_{48} :\equiv \forall n\,(l \in \mathbb{R} \land \mathrm{Seq}(a)) \to \forall n \in \mathbb{N}\,({\textstyle\sum_{k=1}^{n} l \odot a(k)} = l \cdot {\textstyle\sum_{k=1}^{n} a_{k}})$
\item $\varphi_{49} :\equiv n \in B_{47} \to S(n) \in \mathbb{N} \to n \neq \emptyset \to {\textstyle\sum_{k=1}^{n} l \odot a(k)} = l \cdot {\textstyle\sum_{k=1}^{n} a_{k}}$
\item $\varphi_{50} :\equiv n \in \mathbb{N} \to {\textstyle\sum_{k=1}^{S(n)} l \odot a(k)} = {\textstyle\sum_{k=1}^{n} l \odot a(k)} + l \odot a(S(n))$
\item $\varphi_{51} :\equiv (n \in \mathbb{N} \leftrightarrow n \in \omega \land n \neq \emptyset) \to l \in \mathbb{R} \to \mathrm{Seq}(a) \to n \in \mathbb{N} \to n \in \omega$
\item $\varphi_{52} :\equiv l \cdot {\textstyle\sum_{k=1}^{n} a_{k}} + l \odot a(S(n)) = l \cdot {\textstyle\sum_{k=1}^{n} a_{k}} + l \odot a(S(n))$
\item $\varphi_{53} :\equiv l \cdot ({\textstyle\sum_{k=1}^{n} a_{k}} + a_{S(n)}) = l \cdot z \to l \cdot ({\textstyle\sum_{k=1}^{n} a_{k}} + a_{S(n)}) = l \cdot w$
\item $\varphi_{54} :\equiv n \neq \emptyset \to {\textstyle\sum_{k=1}^{S(n)} l \odot a(k)} = {\textstyle\sum_{k=1}^{n} l \odot a(k)} + l \odot a(S(n))$
\item $\varphi_{55} :\equiv S(n) \in \mathbb{N} \to n \neq \emptyset \to l \cdot ({\textstyle\sum_{k=1}^{n} a_{k}} + a_{S(n)}) = l \cdot {\textstyle\sum_{k=1}^{S(n)} a_{k}}$
\item $\varphi_{56} :\equiv n \in B_{47} \to S(n) \in \mathbb{N} \to n \neq \emptyset \to \forall n \in \mathbb{N}\,(l \odot a(n) = l \cdot a_{n})$
\item $\varphi_{57} :\equiv \forall n \in \mathbb{N}\,({\textstyle\sum_{k=1}^{S(n)} l \odot a(k)} = {\textstyle\sum_{k=1}^{n} l \odot a(k)} + l \odot a(S(n)))$
\item $\varphi_{58} :\equiv {\textstyle\sum_{k=1}^{z} l \odot a(k)} = l \cdot {\textstyle\sum_{k=1}^{z} a_{k}} \to {\textstyle\sum_{k=1}^{w} l \odot a(k)} = l \cdot {\textstyle\sum_{k=1}^{w} a_{k}}$
\item $\varphi_{59} :\equiv \forall a\,\forall n\,(\mathrm{Seq}(a) \to n \in \mathbb{N} \to a_{n} \in \mathbb{R}) \to \forall n\,(\mathrm{Seq}(a) \to n \in \mathbb{N} \to a_{n} \in \mathbb{R})$
\item $\varphi_{60} :\equiv \forall n\,(\mathrm{Seq}(a) \to n \in \mathbb{N} \to a_{n} \in \mathbb{R}) \to \mathrm{Seq}(a) \to S(n) \in \mathbb{N} \to a_{S(n)} \in \mathbb{R}$
\item $\varphi_{61} :\equiv S(n) \in \mathbb{N} \to n \neq \emptyset \to n \in \mathbb{N} \to {\textstyle\sum_{k=1}^{S(n)} a_{k}} = {\textstyle\sum_{k=1}^{n} a_{k}} + a_{S(n)}$
\item $\varphi_{62} :\equiv n \in B_{47} \to S(n) \in \mathbb{N} \to n \neq \emptyset \to {\textstyle\sum_{k=1}^{S(n)} a_{k}} = {\textstyle\sum_{k=1}^{n} a_{k}} + a_{S(n)}$
\item $\varphi_{63} :\equiv n \in B_{47} \to S(n) \in \mathbb{N} \to n \neq \emptyset \to {\textstyle\sum_{k=1}^{n} a_{k}} + a_{S(n)} = {\textstyle\sum_{k=1}^{S(n)} a_{k}}$
\item $\varphi_{64} :\equiv n \neq \emptyset \to l \cdot ({\textstyle\sum_{k=1}^{n} a_{k}} + a_{S(n)}) = l \cdot {\textstyle\sum_{k=1}^{n} a_{k}} + l \cdot a_{S(n)}$
\item $\varphi_{65} :\equiv n \neq \emptyset \to l \cdot {\textstyle\sum_{k=1}^{n} a_{k}} + l \cdot a_{S(n)} = l \cdot ({\textstyle\sum_{k=1}^{n} a_{k}} + a_{S(n)})$
\item $\varphi_{66} :\equiv {\textstyle\sum_{k=1}^{n} l \odot a(k)} + l \odot a(S(n)) = l \cdot {\textstyle\sum_{k=1}^{n} a_{k}} + l \odot a(S(n))$
\item $\varphi_{67} :\equiv n \in B_{47} \to S(n) \in \mathbb{N} \to n = \emptyset \to {\textstyle\sum_{k=1}^{S(n)} l \odot a(k)} = l \cdot {\textstyle\sum_{k=1}^{S(n)} a_{k}}$
\item $\varphi_{68} :\equiv S(n) \in \mathbb{N} \to n \neq \emptyset \to {\textstyle\sum_{k=1}^{S(n)} l \odot a(k)} = l \cdot ({\textstyle\sum_{k=1}^{n} a_{k}} + a_{S(n)})$
\item $\varphi_{69} :\equiv n \in B_{47} \to S(n) \in \mathbb{N} \to n \neq \emptyset \to {\textstyle\sum_{k=1}^{S(n)} l \odot a(k)} = l \cdot {\textstyle\sum_{k=1}^{S(n)} a_{k}}$
\item $\varphi_{70} :\equiv \forall c\,\forall a\,(c \in \mathbb{R} \to \mathrm{Seq}(a) \to \mathrm{Seq}(c \odot a) \land \forall n \in \mathbb{N}\,(c \odot a(n) = c \cdot a_{n}))$
\item $\varphi_{71} :\equiv l \cdot a_{1_{\omega}} = l \cdot {\textstyle\sum_{k=1}^{1_{\omega}} a_{k}} \to {\textstyle\sum_{k=1}^{1_{\omega}} l \odot a(k)} = l \cdot {\textstyle\sum_{k=1}^{1_{\omega}} a_{k}}$
\item $\varphi_{72} :\equiv {\textstyle\sum_{k=1}^{n} l \odot a(k)} + l \odot a(S(n)) = {\textstyle\sum_{k=1}^{n} l \odot a(k)} + l \odot a(S(n))$
\item $\varphi_{73} :\equiv \forall A_{48}\,\exists B_{47}\,\forall n\,(n \in B_{47} \leftrightarrow n \in A_{48} \land (n \in \mathbb{N} \to {\textstyle\sum_{k=1}^{n} l \odot a(k)} = l \cdot {\textstyle\sum_{k=1}^{n} a_{k}}))$
\item $\varphi_{74} :\equiv S(n) \in \mathbb{N} \to n \neq \emptyset \to \forall n \in \mathbb{N}\,({\textstyle\sum_{k=1}^{S(n)} a_{k}} = {\textstyle\sum_{k=1}^{n} a_{k}} + a_{S(n)})$
\item $\varphi_{75} :\equiv l \in \mathbb{R} \land \mathrm{Seq}(a) \land \neg \forall n \in \omega\,(n \in \mathbb{N} \to {\textstyle\sum_{k=1}^{n} l \odot a(k)} = l \cdot {\textstyle\sum_{k=1}^{n} a_{k}})$
\item $\varphi_{76} :\equiv S(n) \in \mathbb{N} \to n \neq \emptyset \to {\textstyle\sum_{k=1}^{S(n)} l \odot a(k)} = l \cdot {\textstyle\sum_{k=1}^{n} a_{k}} + l \cdot a_{S(n)}$
\item $\varphi_{77} :\equiv S(n) \in B_{47} \leftrightarrow S(n) \in \omega \land (S(n) \in \mathbb{N} \to {\textstyle\sum_{k=1}^{S(n)} l \odot a(k)} = l \cdot {\textstyle\sum_{k=1}^{S(n)} a_{k}})$
\item $\varphi_{78} :\equiv \forall z\,({\textstyle\sum_{k=1}^{1_{\omega}} l \odot a(k)} = l \cdot a_{1_{\omega}} \to l \cdot a_{1_{\omega}} = z \to {\textstyle\sum_{k=1}^{1_{\omega}} l \odot a(k)} = z)$
\item $\varphi_{79} :\equiv n \neq \emptyset \to l \cdot {\textstyle\sum_{k=1}^{n} a_{k}} + l \odot a(S(n)) = l \cdot {\textstyle\sum_{k=1}^{n} a_{k}} + l \cdot a_{S(n)}$
\item $\varphi_{80} :\equiv l \in \mathbb{R} \to \mathrm{Seq}(a) \to \emptyset \in \mathbb{N} \to {\textstyle\sum_{k=1}^{\emptyset} l \odot a(k)} = l \cdot {\textstyle\sum_{k=1}^{\emptyset} a_{k}}$
\item $\varphi_{81} :\equiv a_{S(n)} \in \mathbb{R} \to l \cdot ({\textstyle\sum_{k=1}^{n} a_{k}} + a_{S(n)}) = l \cdot {\textstyle\sum_{k=1}^{n} a_{k}} + l \cdot a_{S(n)}$
\item $\varphi_{82} :\equiv S(n) \in \mathbb{N} \to n \neq \emptyset \to {\textstyle\sum_{k=1}^{S(n)} l \odot a(k)} = l \cdot {\textstyle\sum_{k=1}^{n} a_{k}} + l \odot a(S(n))$
\item $\varphi_{83} :\equiv \forall n \in \mathbb{N}\,(l \odot a(n) = l \cdot a_{n}) \to S(n) \in \mathbb{N} \to l \odot a(S(n)) = l \cdot a_{S(n)}$
\item $\varphi_{84} :\equiv \forall n \in \mathbb{N}\,(l \odot a(n) = l \cdot a_{n}) \to 1_{\omega} \in \mathbb{N} \to l \odot a(1_{\omega}) = l \cdot a_{1_{\omega}}$
\item $\varphi_{85} :\equiv {\textstyle\sum_{k=1}^{1_{\omega}} l \odot a(k)} = l \cdot {\textstyle\sum_{k=1}^{1_{\omega}} a_{k}} \to {\textstyle\sum_{k=1}^{S(n)} l \odot a(k)} = l \cdot {\textstyle\sum_{k=1}^{S(n)} a_{k}}$
\item $\varphi_{86} :\equiv \forall a\,\forall n\,(\mathrm{Seq}(a) \to n \in \mathbb{N} \to a_{n} \in \mathbb{R}) \to \forall n\,(\mathrm{Seq}(\Sigma a) \to n \in \mathbb{N} \to {\textstyle\sum_{k=1}^{n} a_{k}} \in \mathbb{R})$
\item $\varphi_{87} :\equiv \forall z_{50}\,\forall A_{48}\,\exists B_{47}\,\forall n\,(n \in B_{47} \leftrightarrow n \in A_{48} \land (n \in \mathbb{N} \to {\textstyle\sum_{k=1}^{n} l \odot z_{50}(k)} = l \cdot {\textstyle\sum_{k=1}^{n} z_{50}(k)}))$
\item $\varphi_{88} :\equiv \forall n\,(n \in \mathbb{N} \leftrightarrow n \in \omega \land n \neq \emptyset) \to (\emptyset \in \mathbb{N} \leftrightarrow \emptyset \in \omega \land \emptyset \neq \emptyset)$
\item $\varphi_{89} :\equiv \forall z\,({\textstyle\sum_{k=1}^{1_{\omega}} l \odot a(k)} = l \odot a(1_{\omega}) \to l \odot a(1_{\omega}) = z \to {\textstyle\sum_{k=1}^{1_{\omega}} l \odot a(k)} = z)$
\item $\varphi_{90} :\equiv \forall w\,(1_{\omega} = w \to {\textstyle\sum_{k=1}^{1_{\omega}} l \odot a(k)} = l \cdot {\textstyle\sum_{k=1}^{1_{\omega}} a_{k}} \to {\textstyle\sum_{k=1}^{w} l \odot a(k)} = l \cdot {\textstyle\sum_{k=1}^{w} a_{k}})$
\item $\varphi_{91} :\equiv \forall n\,(\mathrm{Seq}(\Sigma a) \to n \in \mathbb{N} \to {\textstyle\sum_{k=1}^{n} a_{k}} \in \mathbb{R}) \to \mathrm{Seq}(\Sigma a) \to n \in \mathbb{N} \to {\textstyle\sum_{k=1}^{n} a_{k}} \in \mathbb{R}$
\item $\varphi_{92} :\equiv \forall y\,({\textstyle\sum_{k=1}^{1_{\omega}} a_{k}} = y \to y = {\textstyle\sum_{k=1}^{1_{\omega}} a_{k}}) \to {\textstyle\sum_{k=1}^{1_{\omega}} a_{k}} = a_{1_{\omega}} \to a_{1_{\omega}} = {\textstyle\sum_{k=1}^{1_{\omega}} a_{k}}$
\item $\varphi_{93} :\equiv l \cdot ({\textstyle\sum_{k=1}^{n} a_{k}} + a_{S(n)}) = l \cdot {\textstyle\sum_{k=1}^{S(n)} a_{k}} \to {\textstyle\sum_{k=1}^{S(n)} l \odot a(k)} = l \cdot {\textstyle\sum_{k=1}^{S(n)} a_{k}}$
\item $\varphi_{94} :\equiv \forall z\,(l \in \mathbb{R} \to {\textstyle\sum_{k=1}^{n} a_{k}} \in \mathbb{R} \to z \in \mathbb{R} \to l \cdot ({\textstyle\sum_{k=1}^{n} a_{k}} + z) = l \cdot {\textstyle\sum_{k=1}^{n} a_{k}} + l \cdot z)$
\item $\varphi_{95} :\equiv \forall z_{49}\,\forall z_{50}\,\forall A_{48}\,\exists B_{47}\,\forall n\,(n \in B_{47} \leftrightarrow n \in A_{48} \land (n \in \mathbb{N} \to {\textstyle\sum_{k=1}^{n} z_{49} \odot z_{50}(k)} = z_{49} \cdot {\textstyle\sum_{k=1}^{n} z_{50}(k)}))$
\item $\varphi_{96} :\equiv \mathrm{Seq}(\Sigma a) \land {\textstyle\sum_{k=1}^{1_{\omega}} a_{k}} = a_{1_{\omega}} \land \forall n \in \mathbb{N}\,({\textstyle\sum_{k=1}^{S(n)} a_{k}} = {\textstyle\sum_{k=1}^{n} a_{k}} + a_{S(n)})$
\item $\varphi_{97} :\equiv ({\textstyle\sum_{k=1}^{1_{\omega}} a_{k}} = a_{1_{\omega}} \to a_{1_{\omega}} = {\textstyle\sum_{k=1}^{1_{\omega}} a_{k}}) \to l \in \mathbb{R} \to \mathrm{Seq}(a) \to a_{1_{\omega}} = {\textstyle\sum_{k=1}^{1_{\omega}} a_{k}}$
\item $\varphi_{98} :\equiv \emptyset \in B_{47} \leftrightarrow \emptyset \in \omega \land (\emptyset \in \mathbb{N} \to {\textstyle\sum_{k=1}^{\emptyset} l \odot a(k)} = l \cdot {\textstyle\sum_{k=1}^{\emptyset} a_{k}})$
\item $\varphi_{99} :\equiv \forall x\,(\mathrm{Ind}(x) \leftrightarrow \emptyset \in x \land \forall y \in x\,(S(y) \in x)) \to (\mathrm{Ind}(\omega) \leftrightarrow \emptyset \in \omega \land \forall y \in \omega\,(S(y) \in \omega))$
\item $\varphi_{100} :\equiv \forall n \in \mathbb{N}\,({\textstyle\sum_{k=1}^{S(n)} a_{k}} = {\textstyle\sum_{k=1}^{n} a_{k}} + a_{S(n)}) \to n \in \mathbb{N} \to {\textstyle\sum_{k=1}^{S(n)} a_{k}} = {\textstyle\sum_{k=1}^{n} a_{k}} + a_{S(n)}$
\item $\varphi_{101} :\equiv (l \in \mathbb{R} \land \mathrm{Seq}(a) \to \forall n\,(l \in \mathbb{R} \land \mathrm{Seq}(a))) \to \varphi_{41}$
\item $\varphi_{102} :\equiv \mathrm{Seq}(a) \to \forall n\,(\varphi_{43}) \to n \in \omega \to \forall x\,(\mathrm{Ind}(x) \to n \in x)$
\item $\varphi_{103} :\equiv l \cdot ({\textstyle\sum_{k=1}^{n} a_{k}} + a_{S(n)}) = l \cdot ({\textstyle\sum_{k=1}^{n} a_{k}} + a_{S(n)}) \to l \cdot ({\textstyle\sum_{k=1}^{n} a_{k}} + a_{S(n)}) = l \cdot {\textstyle\sum_{k=1}^{S(n)} a_{k}}$
\item $\varphi_{104} :\equiv \mathrm{Seq}(a) \to \forall n\,(\varphi_{43}) \to n \in \omega \to \mathrm{Ind}(B_{47}) \to n \in B_{47}$
\item $\varphi_{105} :\equiv \forall z\,({\textstyle\sum_{k=1}^{S(n)} l \odot a(k)} = l \cdot ({\textstyle\sum_{k=1}^{n} a_{k}} + a_{S(n)}) \to l \cdot ({\textstyle\sum_{k=1}^{n} a_{k}} + a_{S(n)}) = z \to {\textstyle\sum_{k=1}^{S(n)} l \odot a(k)} = z)$
\item $\varphi_{106} :\equiv l \in \mathbb{R} \to \mathrm{Seq}(a) \to \forall n\,(\varphi_{43}) \to n \in B_{47} \to S(n) \in B_{47}$
\item $\varphi_{107} :\equiv l \in \mathbb{R} \land \mathrm{Seq}(a) \land \forall n\,(\varphi_{43}) \to n \in B_{47} \to S(n) \in B_{47}$
\item $\varphi_{108} :\equiv l \in \mathbb{R} \to \mathrm{Seq}(a) \to \forall n\,(\varphi_{43}) \to \forall n \in B_{47}\,(S(n) \in B_{47})$
\item $\varphi_{109} :\equiv \forall n\,(\varphi_{43}) \to n \in B_{47} \to S(n) \in \mathbb{N} \to n = \emptyset \to S(n) = 1_{\omega}$
\item $\varphi_{110} :\equiv \forall n\,(\varphi_{43}) \to n \in B_{47} \to S(n) \in \mathbb{N} \to n = \emptyset \to 1_{\omega} = S(n)$
\item $\varphi_{111} :\equiv l \cdot {\textstyle\sum_{k=1}^{n} a_{k}} + l \cdot a_{S(n)} = l \cdot ({\textstyle\sum_{k=1}^{n} a_{k}} + a_{S(n)}) \to {\textstyle\sum_{k=1}^{S(n)} l \odot a(k)} = l \cdot ({\textstyle\sum_{k=1}^{n} a_{k}} + a_{S(n)})$
\item $\varphi_{112} :\equiv l \cdot {\textstyle\sum_{k=1}^{n} a_{k}} + l \odot a(S(n)) = l \cdot {\textstyle\sum_{k=1}^{n} a_{k}} + z \to l \cdot {\textstyle\sum_{k=1}^{n} a_{k}} + l \odot a(S(n)) = l \cdot {\textstyle\sum_{k=1}^{n} a_{k}} + w$
\item $\varphi_{113} :\equiv \forall n\,(l \in \mathbb{R} \land \mathrm{Seq}(a) \land \forall n\,(\varphi_{43})) \to \forall n \in B_{47}\,(S(n) \in B_{47})$
\item $\varphi_{114} :\equiv \mathrm{Seq}(a) \to \forall n\,(\varphi_{43}) \to n \in B_{47} \to \forall y \in \omega\,(S(y) \in \omega)$
\item $\varphi_{115} :\equiv \forall n\,(\varphi_{43}) \to n \in B_{47} \to S(n) \in \mathbb{N} \to n \neq \emptyset \to S(n) \in \mathbb{N}$
\item $\varphi_{116} :\equiv \forall n\,(\varphi_{43}) \to n \in B_{47} \to S(n) \in \mathbb{N} \to n \neq \emptyset \to \mathrm{Seq}(\Sigma a)$
\item $\varphi_{117} :\equiv \forall n\,(\varphi_{43}) \to n \in B_{47} \to S(n) \in \mathbb{N} \to n \neq \emptyset \to \mathrm{Seq}(l \odot a)$
\item $\varphi_{118} :\equiv \forall n \in \omega\,(n \in \mathbb{N} \to {\textstyle\sum_{k=1}^{n} l \odot a(k)} = l \cdot {\textstyle\sum_{k=1}^{n} a_{k}}) \to n \in \omega \to n \in \mathbb{N} \to {\textstyle\sum_{k=1}^{n} l \odot a(k)} = l \cdot {\textstyle\sum_{k=1}^{n} a_{k}}$
\item $\varphi_{119} :\equiv \forall z\,({\textstyle\sum_{k=1}^{S(n)} l \odot a(k)} = l \cdot {\textstyle\sum_{k=1}^{n} a_{k}} + l \cdot a_{S(n)} \to l \cdot {\textstyle\sum_{k=1}^{n} a_{k}} + l \cdot a_{S(n)} = z \to {\textstyle\sum_{k=1}^{S(n)} l \odot a(k)} = z)$
\item $\varphi_{120} :\equiv \mathrm{Seq}(a) \to \forall n\,(\varphi_{43}) \to n \in B_{47} \to S(n) \in \mathbb{N} \to n \neq \emptyset \to \mathrm{Seq}(a)$
\item $\varphi_{121} :\equiv \forall n\,(\varphi_{43}) \to n \in B_{47} \to S(n) \in \mathbb{N} \to n = \emptyset \to S(n) = S(\emptyset)$
\item $\varphi_{122} :\equiv {\textstyle\sum_{k=1}^{n} l \odot a(k)} + l \odot a(S(n)) = z + l \odot a(S(n)) \to {\textstyle\sum_{k=1}^{n} l \odot a(k)} + l \odot a(S(n)) = w + l \odot a(S(n))$
\item $\varphi_{123} :\equiv \forall y\,({\textstyle\sum_{k=1}^{S(n)} a_{k}} = y \to y = {\textstyle\sum_{k=1}^{S(n)} a_{k}}) \to {\textstyle\sum_{k=1}^{S(n)} a_{k}} = {\textstyle\sum_{k=1}^{n} a_{k}} + a_{S(n)} \to {\textstyle\sum_{k=1}^{n} a_{k}} + a_{S(n)} = {\textstyle\sum_{k=1}^{S(n)} a_{k}}$
\item $\varphi_{124} :\equiv \mathrm{Seq}(a) \to \forall n\,(\varphi_{43}) \to n \in B_{47} \to S(n) \in \mathbb{N} \to n \neq \emptyset \to n \in \mathbb{N}$
\item $\varphi_{125} :\equiv \mathrm{Seq}(a) \to \forall n\,(\varphi_{43}) \to n \in B_{47} \to S(n) \in \mathbb{N} \to n = \emptyset \to n = \emptyset$
\item $\varphi_{126} :\equiv \mathrm{Seq}(a) \to \forall n\,(\varphi_{43}) \to n \in B_{47} \to S(n) \in \mathbb{N} \to n \neq \emptyset \to l \in \mathbb{R}$
\item $\varphi_{127} :\equiv \mathrm{Seq}(a) \to \forall n\,(\varphi_{43}) \to n \in B_{47} \to S(n) \in \mathbb{N} \to n \neq \emptyset \to n \neq \emptyset$
\item $\varphi_{128} :\equiv \forall n\,(\varphi_{43}) \to n \in B_{47} \to S(n) \in \mathbb{N} \to n \neq \emptyset \to a_{S(n)} \in \mathbb{R}$
\item $\varphi_{129} :\equiv \forall z\,({\textstyle\sum_{k=1}^{S(n)} l \odot a(k)} = l \cdot {\textstyle\sum_{k=1}^{n} a_{k}} + l \odot a(S(n)) \to l \cdot {\textstyle\sum_{k=1}^{n} a_{k}} + l \odot a(S(n)) = z \to {\textstyle\sum_{k=1}^{S(n)} l \odot a(k)} = z)$
\item $\varphi_{130} :\equiv l \in \mathbb{R} \to \mathrm{Seq}(a) \to \forall n\,(\varphi_{43}) \to n \in B_{47} \to n \in \omega \to S(n) \in \omega$
\item $\varphi_{131} :\equiv \forall w\,({\textstyle\sum_{k=1}^{n} a_{k}} + a_{S(n)} = w \to l \cdot ({\textstyle\sum_{k=1}^{n} a_{k}} + a_{S(n)}) = l \cdot ({\textstyle\sum_{k=1}^{n} a_{k}} + a_{S(n)}) \to l \cdot ({\textstyle\sum_{k=1}^{n} a_{k}} + a_{S(n)}) = l \cdot w)$
\item $\varphi_{132} :\equiv \forall n\,(\varphi_{43}) \to n \in B_{47} \to S(n) \in \mathbb{N} \to n \neq \emptyset \to {\textstyle\sum_{k=1}^{n} a_{k}} \in \mathbb{R}$
\item $\varphi_{133} :\equiv l \cdot {\textstyle\sum_{k=1}^{n} a_{k}} + l \odot a(S(n)) = l \cdot {\textstyle\sum_{k=1}^{n} a_{k}} + l \cdot a_{S(n)} \to {\textstyle\sum_{k=1}^{S(n)} l \odot a(k)} = l \cdot {\textstyle\sum_{k=1}^{n} a_{k}} + l \cdot a_{S(n)}$
\item $\varphi_{134} :\equiv l \cdot ({\textstyle\sum_{k=1}^{n} a_{k}} + a_{S(n)}) = l \cdot {\textstyle\sum_{k=1}^{n} a_{k}} + l \cdot a_{S(n)} \to l \cdot {\textstyle\sum_{k=1}^{n} a_{k}} + l \cdot a_{S(n)} = l \cdot ({\textstyle\sum_{k=1}^{n} a_{k}} + a_{S(n)})$
\item $\varphi_{135} :\equiv \forall n\,(\varphi_{43}) \to n \in B_{47} \to n \in \mathbb{N} \to {\textstyle\sum_{k=1}^{n} l \odot a(k)} = l \cdot {\textstyle\sum_{k=1}^{n} a_{k}}$
\item $\varphi_{136} :\equiv \forall z\,({\textstyle\sum_{k=1}^{S(n)} l \odot a(k)} = {\textstyle\sum_{k=1}^{n} l \odot a(k)} + l \odot a(S(n)) \to {\textstyle\sum_{k=1}^{n} l \odot a(k)} + l \odot a(S(n)) = z \to {\textstyle\sum_{k=1}^{S(n)} l \odot a(k)} = z)$
\item $\varphi_{137} :\equiv (l \in \mathbb{R} \to \varphi_{26}) \to l \in \mathbb{R} \to \mathrm{Seq}(a) \to \forall n \in \mathbb{N}\,(l \odot a(n) = l \cdot a_{n})$
\item $\varphi_{138} :\equiv \mathrm{Seq}(\Sigma l \odot a) \land {\textstyle\sum_{k=1}^{1_{\omega}} l \odot a(k)} = l \odot a(1_{\omega}) \land \varphi_{57}$
\item $\varphi_{139} :\equiv \forall n\,(\varphi_{43}) \to \forall n \in \omega\,(n \in \mathbb{N} \to {\textstyle\sum_{k=1}^{n} l \odot a(k)} = l \cdot {\textstyle\sum_{k=1}^{n} a_{k}})$
\item $\varphi_{140} :\equiv l \in \mathbb{R} \to \mathrm{Seq}(a) \to \forall n\,(\varphi_{43}) \to \emptyset \in \omega \land \forall y \in \omega\,(S(y) \in \omega)$
\item $\varphi_{141} :\equiv \forall n\,(\varphi_{43}) \to n \in B_{47} \to S(n) \in \mathbb{N} \to {\textstyle\sum_{k=1}^{S(n)} l \odot a(k)} = l \cdot {\textstyle\sum_{k=1}^{S(n)} a_{k}}$
\item $\varphi_{142} :\equiv (l \in \mathbb{R} \to \mathrm{Seq}(a) \to \mathrm{Seq}(a)) \to l \in \mathbb{R} \to \mathrm{Seq}(a) \to \varphi_{96}$
\item $\varphi_{143} :\equiv \mathrm{Seq}(a) \to \forall n\,(\varphi_{43}) \to n \in \omega \to n \in \mathbb{N} \to {\textstyle\sum_{k=1}^{n} l \odot a(k)} = l \cdot {\textstyle\sum_{k=1}^{n} a_{k}}$
\item $\varphi_{144} :\equiv (l \in \mathbb{R} \to \mathrm{Seq}(a) \to n \in \mathbb{N} \to n \in \omega) \to (l \in \mathbb{R} \to \mathrm{Seq}(a) \to n \in \mathbb{N} \to n \in \mathbb{N}) \to \varphi_{33}$
\item $\varphi_{145} :\equiv \varphi_{52} \to l \cdot {\textstyle\sum_{k=1}^{n} a_{k}} + l \odot a(S(n)) = l \cdot {\textstyle\sum_{k=1}^{n} a_{k}} + l \cdot a_{S(n)}$
\item $\varphi_{146} :\equiv (a_{1_{\omega}} = {\textstyle\sum_{k=1}^{1_{\omega}} a_{k}} \to \varphi_{31}) \to l \in \mathbb{R} \to \mathrm{Seq}(a) \to l \cdot a_{1_{\omega}} = l \cdot {\textstyle\sum_{k=1}^{1_{\omega}} a_{k}}$
\item $\varphi_{147} :\equiv l \in \mathbb{R} \land \mathrm{Seq}(a) \land \forall n\,(\varphi_{43}) \to n \in \omega \to n \in \mathbb{N} \to {\textstyle\sum_{k=1}^{n} l \odot a(k)} = l \cdot {\textstyle\sum_{k=1}^{n} a_{k}}$
\item $\varphi_{148} :\equiv (S(n) = 1_{\omega} \to 1_{\omega} = S(n)) \to l \in \mathbb{R} \to \mathrm{Seq}(a) \to \varphi_{110}$
\item $\varphi_{149} :\equiv \forall w\,(l \odot a(S(n)) = w \to \varphi_{52} \to l \cdot {\textstyle\sum_{k=1}^{n} a_{k}} + l \odot a(S(n)) = l \cdot {\textstyle\sum_{k=1}^{n} a_{k}} + w)$
\item $\varphi_{150} :\equiv \forall n\,(l \in \mathbb{R} \land \mathrm{Seq}(a) \land \forall n\,(\varphi_{43})) \to \forall n \in \omega\,(n \in \mathbb{N} \to {\textstyle\sum_{k=1}^{n} l \odot a(k)} = l \cdot {\textstyle\sum_{k=1}^{n} a_{k}})$
\item $\varphi_{151} :\equiv (l \in \mathbb{R} \to \mathrm{Seq}(a) \to \varphi_{96}) \to l \in \mathbb{R} \to \mathrm{Seq}(a) \to {\textstyle\sum_{k=1}^{1_{\omega}} a_{k}} = a_{1_{\omega}}$
\item $\varphi_{152} :\equiv (\forall x\,(\mathrm{Ind}(x) \to n \in x) \to \mathrm{Ind}(B_{47}) \to n \in B_{47}) \to l \in \mathbb{R} \to \varphi_{104}$
\item $\varphi_{153} :\equiv \forall w\,({\textstyle\sum_{k=1}^{n} l \odot a(k)} = w \to \varphi_{72} \to {\textstyle\sum_{k=1}^{n} l \odot a(k)} + l \odot a(S(n)) = w + l \odot a(S(n)))$
\item $\varphi_{154} :\equiv ({\textstyle\sum_{k=1}^{1_{\omega}} l \odot a(k)} = l \odot a(1_{\omega}) \to \varphi_{47}) \to l \in \mathbb{R} \to \mathrm{Seq}(a) \to {\textstyle\sum_{k=1}^{1_{\omega}} l \odot a(k)} = l \cdot a_{1_{\omega}}$
\item $\varphi_{155} :\equiv \exists B_{47}\,\forall n\,(\varphi_{43}) \to l \in \mathbb{R} \to \mathrm{Seq}(a) \to \exists B_{47}\,\forall n\,(\varphi_{43})$
\item $\varphi_{156} :\equiv ({\textstyle\sum_{k=1}^{1_{\omega}} l \odot a(k)} = l \cdot a_{1_{\omega}} \to \varphi_{71}) \to l \in \mathbb{R} \to \mathrm{Seq}(a) \to {\textstyle\sum_{k=1}^{1_{\omega}} l \odot a(k)} = l \cdot {\textstyle\sum_{k=1}^{1_{\omega}} a_{k}}$
\item $\varphi_{157} :\equiv (n = \emptyset \to S(n) = S(n) \to S(n) = S(\emptyset)) \to l \in \mathbb{R} \to \mathrm{Seq}(a) \to \varphi_{121}$
\item $\varphi_{158} :\equiv l \in \mathbb{R} \to \mathrm{Seq}(a) \to \forall n\,(\varphi_{43}) \to n \in B_{47} \to S(n) \in \mathbb{N} \to \varphi_{35}$
\item $\varphi_{159} :\equiv l \in \mathbb{R} \land \mathrm{Seq}(a) \land \forall n\,(\varphi_{43}) \to \forall n\,(l \in \mathbb{R} \land \mathrm{Seq}(a) \land \forall n\,(\varphi_{43}))$
\item $\varphi_{160} :\equiv (l \in \mathbb{R} \to \mathrm{Seq}(a) \to \mathrm{Seq}(l \odot a)) \to l \in \mathbb{R} \to \mathrm{Seq}(a) \to \varphi_{138}$
\item $\varphi_{161} :\equiv \mathrm{Seq}(a) \to \forall n\,(\varphi_{43}) \to n \in B_{47} \to S(n) \in \mathbb{N} \to n \neq \emptyset \to \varphi_{50}$
\item $\varphi_{162} :\equiv l \in \mathbb{R} \to \mathrm{Seq}(a) \to \forall n\,(\varphi_{43}) \to n \in B_{47} \to S(n) \in \mathbb{N} \to \varphi_{54}$
\item $\varphi_{163} :\equiv \mathrm{Seq}(a) \to \forall n\,(\varphi_{43}) \to n \in B_{47} \to S(n) \in \mathbb{N} \to n \neq \emptyset \to \varphi_{57}$
\item $\varphi_{164} :\equiv l \in \mathbb{R} \to \mathrm{Seq}(a) \to \forall n\,(\varphi_{43}) \to n \in B_{47} \to S(n) \in \mathbb{N} \to \varphi_{64}$
\item $\varphi_{165} :\equiv l \in \mathbb{R} \to \mathrm{Seq}(a) \to \forall n\,(\varphi_{43}) \to n \in B_{47} \to S(n) \in \mathbb{N} \to \varphi_{65}$
\item $\varphi_{166} :\equiv \mathrm{Seq}(a) \to \forall n\,(\varphi_{43}) \to n \in B_{47} \to S(n) \in \mathbb{N} \to n \neq \emptyset \to \varphi_{66}$
\item $\varphi_{167} :\equiv l \in \mathbb{R} \to \mathrm{Seq}(a) \to \forall n\,(\varphi_{43}) \to n \in B_{47} \to S(n) \in \mathbb{N} \to \varphi_{79}$
\item $\varphi_{168} :\equiv {\textstyle\sum_{k=1}^{S(n)} l \odot a(k)} = {\textstyle\sum_{k=1}^{n} l \odot a(k)} + l \odot a(S(n)) \to \varphi_{66} \to {\textstyle\sum_{k=1}^{S(n)} l \odot a(k)} = l \cdot {\textstyle\sum_{k=1}^{n} a_{k}} + l \odot a(S(n))$
\item $\varphi_{169} :\equiv \mathrm{Ind}(\omega) \to (\mathrm{Ind}(\omega) \leftrightarrow \emptyset \in \omega \land \forall y \in \omega\,(S(y) \in \omega)) \to \varphi_{140}$
\item $\varphi_{170} :\equiv (l \in \mathbb{R} \to \mathrm{Seq}(a) \to \forall n \in \mathbb{N}\,(l \odot a(n) = l \cdot a_{n})) \to \varphi_{84} \to l \in \mathbb{R} \to \varphi_{19}$
\item $\varphi_{171} :\equiv \mathrm{Seq}(a) \to \forall n\,(\varphi_{43}) \to n \in B_{47} \to S(n) \in \mathbb{N} \to n \neq \emptyset \to \varphi_{96}$
\item $\varphi_{172} :\equiv (l \in \mathbb{R} \to \mathrm{Seq}(a) \to a_{1_{\omega}} = {\textstyle\sum_{k=1}^{1_{\omega}} a_{k}}) \to l \cdot a_{1_{\omega}} = l \cdot a_{1_{\omega}} \to \varphi_{146}$
\item $\varphi_{173} :\equiv (l \in \mathbb{R} \to \mathrm{Seq}(a) \to \forall n\,(\varphi_{43}) \to n \in B_{47} \to n \in \omega) \to l \in \mathbb{R} \to \mathrm{Seq}(a) \to \forall n\,(\varphi_{43}) \to n \in B_{47} \to S(n) \in \omega$
\item $\varphi_{174} :\equiv (l \in \mathbb{R} \to \mathrm{Seq}(a) \to l \cdot a_{1_{\omega}} = l \cdot {\textstyle\sum_{k=1}^{1_{\omega}} a_{k}}) \to \varphi_{156}$
\item $\varphi_{175} :\equiv \mathrm{Seq}(a) \to \forall n\,(\varphi_{43}) \to n \in B_{47} \to S(n) \in \mathbb{N} \to n \neq \emptyset \to \varphi_{138}$
\item $\varphi_{176} :\equiv ({\textstyle\sum_{k=1}^{S(n)} a_{k}} = {\textstyle\sum_{k=1}^{n} a_{k}} + a_{S(n)} \to {\textstyle\sum_{k=1}^{n} a_{k}} + a_{S(n)} = {\textstyle\sum_{k=1}^{S(n)} a_{k}}) \to l \in \mathbb{R} \to \mathrm{Seq}(a) \to \forall n\,(\varphi_{43}) \to \varphi_{63}$
\item $\varphi_{177} :\equiv (1_{\omega} = S(n) \to \varphi_{85}) \to l \in \mathbb{R} \to \mathrm{Seq}(a) \to \forall n\,(\varphi_{43}) \to \varphi_{67}$
\item $\varphi_{178} :\equiv (l \in \mathbb{R} \to \mathrm{Seq}(a) \to \forall n\,(\varphi_{43}) \to n \in \omega \to n \in \omega) \to (n \in \omega \leftrightarrow \forall x\,(\mathrm{Ind}(x) \to n \in x)) \to l \in \mathbb{R} \to \varphi_{102}$
\item $\varphi_{179} :\equiv (l \in \mathbb{R} \to \mathrm{Seq}(a) \to \forall n\,(\varphi_{43}) \to n \in B_{47} \to n \in B_{47}) \to l \in \mathbb{R} \to \mathrm{Seq}(a) \to \forall n\,(\varphi_{43}) \to \varphi_{44}$
\item $\varphi_{180} :\equiv (l \in \mathbb{R} \to \mathrm{Seq}(a) \to \forall n\,(\varphi_{43}) \to \varphi_{44}) \to l \in \mathbb{R} \to \mathrm{Seq}(a) \to \forall n\,(\varphi_{43}) \to n \in B_{47} \to n \in \omega$
\item $\varphi_{181} :\equiv (l \in \mathbb{R} \to \mathrm{Seq}(a) \to {\textstyle\sum_{k=1}^{1_{\omega}} l \odot a(k)} = l \cdot a_{1_{\omega}}) \to \varphi_{174}$
\item $\varphi_{182} :\equiv (n \in \mathbb{N} \to S(n) \in \mathbb{N}) \to (l \in \mathbb{R} \to \varphi_{124}) \to l \in \mathbb{R} \to \mathrm{Seq}(a) \to \varphi_{115}$
\item $\varphi_{183} :\equiv (l \in \mathbb{R} \to \mathrm{Seq}(a) \to \varphi_{115}) \to l \in \mathbb{R} \to \mathrm{Seq}(a) \to \forall n\,(\varphi_{43}) \to \varphi_{36}$
\item $\varphi_{184} :\equiv (l \in \mathbb{R} \to \mathrm{Seq}(a) \to \forall n\,(\varphi_{43}) \to \varphi_{44}) \to l \in \mathbb{R} \to \mathrm{Seq}(a) \to \varphi_{135}$
\item $\varphi_{185} :\equiv (l \in \mathbb{R} \to \varphi_{114}) \to (\forall y \in \omega\,(S(y) \in \omega) \to n \in \omega \to S(n) \in \omega) \to \varphi_{130}$
\item $\varphi_{186} :\equiv ({\textstyle\sum_{k=1}^{S(n)} l \odot a(k)} = l \cdot ({\textstyle\sum_{k=1}^{n} a_{k}} + a_{S(n)}) \to \varphi_{93}) \to l \in \mathbb{R} \to \mathrm{Seq}(a) \to \forall n\,(\varphi_{43}) \to \varphi_{69}$
\item $\varphi_{187} :\equiv \varphi_{108} \to \varphi_{12} \to \varphi_{11} \to l \in \mathbb{R} \to \mathrm{Seq}(a) \to \forall n\,(\varphi_{43}) \to \mathrm{Ind}(B_{47})$
\item $\varphi_{188} :\equiv \forall B_{47}\,(\varphi_{75} \to \neg \forall n\,(\varphi_{43})) \to \forall B_{47}\,(\varphi_{75}) \to \forall B_{47}\,(\neg \forall n\,(\varphi_{43}))$
\item $\varphi_{189} :\equiv (l \in \mathbb{R} \to \mathrm{Seq}(a) \to {\textstyle\sum_{k=1}^{1_{\omega}} l \odot a(k)} = l \cdot {\textstyle\sum_{k=1}^{1_{\omega}} a_{k}}) \to \varphi_{177}$
\item $\varphi_{190} :\equiv (l \in \mathbb{R} \to \mathrm{Seq}(a) \to \forall n\,(\varphi_{43}) \to \varphi_{69}) \to l \in \mathbb{R} \to \mathrm{Seq}(a) \to \varphi_{141}$
\item $\varphi_{191} :\equiv ({\textstyle\sum_{k=1}^{n} a_{k}} + a_{S(n)} = {\textstyle\sum_{k=1}^{S(n)} a_{k}} \to \varphi_{103}) \to l \in \mathbb{R} \to \mathrm{Seq}(a) \to \forall n\,(\varphi_{43}) \to n \in B_{47} \to \varphi_{55}$
\item $\varphi_{192} :\equiv (l \in \mathbb{R} \to \mathrm{Seq}(a) \to \varphi_{121}) \to S(\emptyset) = 1_{\omega} \to \varphi_{4} \to l \in \mathbb{R} \to \mathrm{Seq}(a) \to \varphi_{109}$
\item $\varphi_{193} :\equiv ({\textstyle\sum_{k=1}^{S(n)} l \odot a(k)} = l \cdot {\textstyle\sum_{k=1}^{n} a_{k}} + l \cdot a_{S(n)} \to \varphi_{111}) \to l \in \mathbb{R} \to \mathrm{Seq}(a) \to \forall n\,(\varphi_{43}) \to n \in B_{47} \to \varphi_{68}$
\item $\varphi_{194} :\equiv ({\textstyle\sum_{k=1}^{S(n)} l \odot a(k)} = l \cdot {\textstyle\sum_{k=1}^{n} a_{k}} + l \odot a(S(n)) \to \varphi_{133}) \to l \in \mathbb{R} \to \mathrm{Seq}(a) \to \forall n\,(\varphi_{43}) \to n \in B_{47} \to \varphi_{76}$
\item $\varphi_{195} :\equiv (l \in \mathbb{R} \to \mathrm{Seq}(a) \to \forall n\,(\varphi_{43}) \to n \in B_{47} \to n \in \omega) \to (l \in \mathbb{R} \to \varphi_{127}) \to l \in \mathbb{R} \to \varphi_{124}$
\item $\varphi_{196} :\equiv (l \in \mathbb{R} \to \mathrm{Seq}(a) \to \forall n\,(\varphi_{43}) \to n \in B_{47} \to S(n) \in \omega) \to (l \in \mathbb{R} \to \mathrm{Seq}(a) \to \varphi_{141}) \to \varphi_{106}$
\item $\varphi_{197} :\equiv (l \in \mathbb{R} \to \varphi_{171}) \to l \in \mathbb{R} \to \mathrm{Seq}(a) \to \forall n\,(\varphi_{43}) \to n \in B_{47} \to \varphi_{74}$
\item $\varphi_{198} :\equiv (l \in \mathbb{R} \to \mathrm{Seq}(a) \to \forall n\,(\varphi_{43}) \to \mathrm{Ind}(B_{47})) \to (l \in \mathbb{R} \to \mathrm{Seq}(a) \to \forall n\,(\varphi_{43}) \to \varphi_{43}) \to l \in \mathbb{R} \to \varphi_{143}$
\item $\varphi_{199} :\equiv (l \in \mathbb{R} \to \varphi_{120}) \to (l \in \mathbb{R} \to \mathrm{Seq}(a) \to \varphi_{115}) \to l \in \mathbb{R} \to \mathrm{Seq}(a) \to \varphi_{128}$
\item $\varphi_{200} :\equiv (l \in \mathbb{R} \to \mathrm{Seq}(a) \to \varphi_{116}) \to (l \in \mathbb{R} \to \varphi_{124}) \to l \in \mathbb{R} \to \mathrm{Seq}(a) \to \varphi_{132}$
\item $\varphi_{201} :\equiv (l \in \mathbb{R} \to \mathrm{Seq}(a) \to \forall n\,(\varphi_{43}) \to \varphi_{56}) \to \varphi_{83} \to l \in \mathbb{R} \to \mathrm{Seq}(a) \to \forall n\,(\varphi_{43}) \to n \in B_{47} \to \varphi_{45}$
\item $\varphi_{202} :\equiv (l \in \mathbb{R} \to \mathrm{Seq}(a) \to \varphi_{135}) \to (l \in \mathbb{R} \to \varphi_{124}) \to l \in \mathbb{R} \to \mathrm{Seq}(a) \to \forall n\,(\varphi_{43}) \to \varphi_{49}$
\item $\varphi_{203} :\equiv (l \in \mathbb{R} \to \mathrm{Seq}(a) \to \forall n\,(\varphi_{43}) \to n \in B_{47} \to \varphi_{74}) \to \varphi_{100} \to l \in \mathbb{R} \to \mathrm{Seq}(a) \to \forall n\,(\varphi_{43}) \to n \in B_{47} \to \varphi_{61}$
\item $\varphi_{204} :\equiv (l \in \mathbb{R} \to \mathrm{Seq}(a) \to \forall n\,(\varphi_{43}) \to \forall n\,(\varphi_{43})) \to (\forall n\,(\varphi_{43}) \to \varphi_{43}) \to l \in \mathbb{R} \to \mathrm{Seq}(a) \to \forall n\,(\varphi_{43}) \to \varphi_{43}$
\item $\varphi_{205} :\equiv (l \in \mathbb{R} \to \mathrm{Seq}(a) \to \forall n\,(\varphi_{43}) \to n \in B_{47} \to \varphi_{61}) \to (l \in \mathbb{R} \to \varphi_{124}) \to l \in \mathbb{R} \to \mathrm{Seq}(a) \to \forall n\,(\varphi_{43}) \to \varphi_{62}$
\item $\varphi_{206} :\equiv (l \in \mathbb{R} \to \mathrm{Seq}(a) \to \varphi_{132}) \to (l \in \mathbb{R} \to \mathrm{Seq}(a) \to \varphi_{128}) \to \varphi_{164}$
\item $\varphi_{207} :\equiv (l \in \mathbb{R} \to \mathrm{Seq}(a) \to \forall n\,(\varphi_{43}) \to \forall n\,(\varphi_{43})) \to (\forall n\,(\varphi_{43}) \to \varphi_{77}) \to l \in \mathbb{R} \to \mathrm{Seq}(a) \to \forall n\,(\varphi_{43}) \to \varphi_{77}$
\item $\varphi_{208} :\equiv (l \in \mathbb{R} \to \mathrm{Seq}(a) \to \forall n\,(\varphi_{43}) \to \varphi_{98}) \to \varphi_{140} \to \varphi_{80} \to l \in \mathbb{R} \to \mathrm{Seq}(a) \to \forall n\,(\varphi_{43}) \to \emptyset \in B_{47}$
\item $\varphi_{209} :\equiv (l \in \mathbb{R} \to \mathrm{Seq}(a) \to \forall n\,(\varphi_{43}) \to \forall n\,(\varphi_{43})) \to (\forall n\,(\varphi_{43}) \to \varphi_{98}) \to l \in \mathbb{R} \to \mathrm{Seq}(a) \to \forall n\,(\varphi_{43}) \to \varphi_{98}$
\item $\varphi_{210} :\equiv (\mathrm{Seq}(l \odot a) \to \varphi_{138}) \to (l \in \mathbb{R} \to \mathrm{Seq}(a) \to \varphi_{117}) \to l \in \mathbb{R} \to \varphi_{175}$
\item $\varphi_{211} :\equiv (l \in \mathbb{R} \to \mathrm{Seq}(a) \to \forall n\,(\varphi_{43}) \to \varphi_{63}) \to l \cdot ({\textstyle\sum_{k=1}^{n} a_{k}} + a_{S(n)}) = l \cdot ({\textstyle\sum_{k=1}^{n} a_{k}} + a_{S(n)}) \to \varphi_{191}$
\item $\varphi_{212} :\equiv (l \in \mathbb{R} \to \varphi_{166}) \to \varphi_{168} \to l \in \mathbb{R} \to \mathrm{Seq}(a) \to \forall n\,(\varphi_{43}) \to n \in B_{47} \to \varphi_{82}$
\item $\varphi_{213} :\equiv (l \in \mathbb{R} \to \mathrm{Seq}(a) \to \forall n\,(\varphi_{43}) \to \varphi_{36}) \to \varphi_{52} \to (l \odot a(S(n)) = l \cdot a_{S(n)} \to \varphi_{145}) \to \varphi_{167}$
\item $\varphi_{214} :\equiv (l \in \mathbb{R} \to \mathrm{Seq}(a) \to \forall n\,(\varphi_{43}) \to \varphi_{49}) \to \varphi_{72} \to ({\textstyle\sum_{k=1}^{n} l \odot a(k)} = l \cdot {\textstyle\sum_{k=1}^{n} a_{k}} \to \varphi_{72} \to \varphi_{66}) \to l \in \mathbb{R} \to \varphi_{166}$
\item $\varphi_{215} :\equiv (l \in \mathbb{R} \to \mathrm{Seq}(a) \to \forall n\,(\varphi_{43}) \to n \in B_{47} \to \varphi_{68}) \to (l \in \mathbb{R} \to \mathrm{Seq}(a) \to \forall n\,(\varphi_{43}) \to n \in B_{47} \to \varphi_{55}) \to \varphi_{186}$
\item $\varphi_{216} :\equiv (l \in \mathbb{R} \to \mathrm{Seq}(a) \to \forall n\,(\varphi_{43}) \to n \in B_{47} \to \varphi_{76}) \to \varphi_{165} \to \varphi_{193}$
\item $\varphi_{217} :\equiv \forall B_{47}\,(\varphi_{75} \to \neg \forall n\,(\varphi_{43})) \to \varphi_{188} \to (\varphi_{75} \to \forall B_{47}\,(\varphi_{75})) \to l \in \mathbb{R} \to \varphi_{46}$
\item $\varphi_{218} :\equiv (l \in \mathbb{R} \to \mathrm{Seq}(a) \to \forall n\,(\varphi_{43}) \to n \in B_{47} \to \varphi_{82}) \to \varphi_{167} \to \varphi_{194}$
\item $\varphi_{219} :\equiv \forall n\,(\varphi_{147}) \to (\forall n\,(\varphi_{147}) \to \varphi_{150}) \to \varphi_{159} \to l \in \mathbb{R} \to \mathrm{Seq}(a) \to \varphi_{139}$
\end{abbrevlist}
\end{fullproof}

% 自動生成 (prover/build.js): SerAdd  ノード数 155
\begin{fullproof}{定理4.4(和)}{$\forall a\,\forall b\,\forall s\,\forall t\,(\mathrm{Seq}(a) \to \mathrm{Seq}(b) \to \mathrm{Conv}(\Sigma a,s) \to \mathrm{Conv}(\Sigma b,t) \to \mathrm{Conv}(\Sigma a \oplus b,s + t))$}{155}\phantomsection\label{tree:SerAdd}
\noindent$\Pi_{1}$:
\begin{ptree}
\AXC{$\varphi_{1}$\lbl{A1}}
\AXC{$\varphi_{13}$\lbl{定理2.1(1)}}
\AXC{$\varphi_{13} \to \varphi_{17}$\lbl{A2}}
\BIC{$\varphi_{17}$}
\AXC{$\varphi_{17} \to \varphi_{15}$\lbl{A2}}
\BIC{$\varphi_{15}$}
\AXC{$\varphi_{15} \to \varphi_{20}$\lbl{A2}}
\BIC{$\varphi_{20}$}
\AXC{$\varphi_{20} \to \varphi_{16}$\lbl{A2}}
\BIC{$\varphi_{16}$}
\AXC{$\varphi_{16} \to \varphi_{1} \to \varphi_{48}$\lbl{A1}}
\BIC{$\varphi_{1} \to \varphi_{48}$}
\BIC{$\varphi_{48}$}
\end{ptree}
\noindent$\Pi_{2}$:
\begin{ptree}
\AXC{$y = z \to x = y \to x = z$\lbl{A6}}
\AXC{$(y = z \to x = y \to x = z) \to x = y \to y = z \to x = z$\lbl{A1}}
\BIC{$x = y \to y = z \to x = z$}
\UIC{$\forall z\,(x = y \to y = z \to x = z)$}
\UIC{$\forall y\,\forall z\,(x = y \to y = z \to x = z)$}
\UIC{$\forall x\,\forall y\,\forall z\,(x = y \to y = z \to x = z)$}
\AXC{$\forall x\,\forall y\,\forall z\,(x = y \to y = z \to x = z) \to \varphi_{21}$\lbl{A2}}
\BIC{$\varphi_{21}$}
\end{ptree}
\noindent$\Pi_{3}$:
\begin{ptree}
\AXC{$x = x$\lbl{A5}}
\AXC{$x = y \to x = x \to y = x$\lbl{A6}}
\AXC{$(x = y \to x = x \to y = x) \to x = x \to x = y \to y = x$\lbl{A1}}
\BIC{$x = x \to x = y \to y = x$}
\BIC{$x = y \to y = x$}
\UIC{$\forall y\,(x = y \to y = x)$}
\UIC{$\forall x\,\forall y\,(x = y \to y = x)$}
\AXC{$\varphi_{18}$\lbl{A2}}
\BIC{$\forall y\,({\textstyle\sum_{k=1}^{n} a \oplus b(k)} = y \to y = {\textstyle\sum_{k=1}^{n} a \oplus b(k)})$}
\end{ptree}
\noindent$\Pi_{4}$:
\begin{ptree}
\AXC{$\stackrel{\Pi_{3}}{\vdots}$}\noLine
\UIC{$\forall y\,({\textstyle\sum_{k=1}^{n} a \oplus b(k)} = y \to y = {\textstyle\sum_{k=1}^{n} a \oplus b(k)})$}
\AXC{$\forall y\,({\textstyle\sum_{k=1}^{n} a \oplus b(k)} = y \to y = {\textstyle\sum_{k=1}^{n} a \oplus b(k)}) \to \varphi_{41}$\lbl{A2}}
\BIC{$\varphi_{41}$}
\end{ptree}
\noindent$\Pi_{5}$:
\begin{ptree}
\AXC{$\forall a\,\forall b\,(\mathrm{Seq}(a) \to \varphi_{24})$\lbl{[Q1]}}
\AXC{$\forall a\,\forall b\,(\mathrm{Seq}(a) \to \varphi_{24}) \to \forall b\,(\mathrm{Seq}(a) \to \varphi_{24})$\lbl{A2}}
\BIC{$\forall b\,(\mathrm{Seq}(a) \to \varphi_{24})$}
\end{ptree}
\noindent$\Pi_{6}$:
\begin{ptree}
\AXC{$\stackrel{\Pi_{5}}{\vdots}$}\noLine
\UIC{$\forall b\,(\mathrm{Seq}(a) \to \varphi_{24})$}
\AXC{$\forall b\,(\mathrm{Seq}(a) \to \varphi_{24}) \to \mathrm{Seq}(a) \to \varphi_{24}$\lbl{A2}}
\BIC{$\mathrm{Seq}(a) \to \varphi_{24}$}
\end{ptree}
\noindent$\Pi_{7}$:
\begin{ptree}
\AXC{$\mathrm{Seq}(a) \to \mathrm{Seq}(b) \to \varphi_{7}$\lbl{A1}}
\AXC{$\stackrel{\Pi_{6}}{\vdots}$}\noLine
\UIC{$\mathrm{Seq}(a) \to \varphi_{24}$}
\AXC{$(\mathrm{Seq}(a) \to \varphi_{24}) \to (\mathrm{Seq}(a) \to \mathrm{Seq}(b) \to \varphi_{7}) \to \varphi_{59}$\lbl{A1}}
\BIC{$(\mathrm{Seq}(a) \to \mathrm{Seq}(b) \to \varphi_{7}) \to \varphi_{59}$}
\BIC{$\varphi_{59}$}
\end{ptree}
\noindent$\Pi_{8}$:
\begin{ptree}
\AXC{$\varphi_{46}$\lbl{A2}}
\AXC{$\mathrm{Seq}(a) \to \mathrm{Seq}(b) \to \varphi_{6}$\lbl{A1}}
\AXC{$\stackrel{\Pi_{7}}{\vdots}$}\noLine
\UIC{$\varphi_{59}$}
\BIC{$\mathrm{Seq}(a) \to \mathrm{Seq}(b) \to \mathrm{Conv}(\Sigma a,s) \to \varphi_{34}$}
\AXC{$\varphi_{71}$\lbl{A1}}
\BIC{$\varphi_{46} \to \mathrm{Seq}(a) \to \mathrm{Seq}(b) \to \mathrm{Conv}(\Sigma a,s) \to \varphi_{33}$}
\BIC{$\mathrm{Seq}(a) \to \mathrm{Seq}(b) \to \mathrm{Conv}(\Sigma a,s) \to \varphi_{33}$}
\end{ptree}
\noindent$\Pi_{9}$:
\begin{ptree}
\AXC{$\stackrel{\Pi_{8}}{\vdots}$}\noLine
\UIC{$\mathrm{Seq}(a) \to \mathrm{Seq}(b) \to \mathrm{Conv}(\Sigma a,s) \to \varphi_{33}$}
\AXC{$(\mathrm{Seq}(a) \to \mathrm{Seq}(b) \to \mathrm{Conv}(\Sigma a,s) \to \varphi_{33}) \to \varphi_{57}$\lbl{A1}}
\BIC{$\varphi_{57}$}
\end{ptree}
\noindent$\Pi_{10}$:
\begin{ptree}
\AXC{$\stackrel{\Pi_{4}}{\vdots}$}\noLine
\UIC{$\varphi_{41}$}
\AXC{$\mathrm{Seq}(a) \to \mathrm{Seq}(b) \to \varphi_{8}$\lbl{A1}}
\AXC{$\stackrel{\Pi_{9}}{\vdots}$}\noLine
\UIC{$\varphi_{57}$}
\BIC{$\mathrm{Seq}(a) \to \mathrm{Seq}(b) \to \mathrm{Conv}(\Sigma a,s) \to \varphi_{28}$}
\AXC{$\varphi_{67}$\lbl{A1}}
\BIC{$\varphi_{41} \to \mathrm{Seq}(a) \to \mathrm{Seq}(b) \to \mathrm{Conv}(\Sigma a,s) \to \varphi_{29}$}
\BIC{$\mathrm{Seq}(a) \to \mathrm{Seq}(b) \to \mathrm{Conv}(\Sigma a,s) \to \varphi_{29}$}
\end{ptree}
\noindent$\Pi_{11}$:
\begin{ptree}
\AXC{$\mathrm{Seq}(a) \to \mathrm{Seq}(b) \to \varphi_{6}$\lbl{A1}}
\AXC{$\forall a\,(\mathrm{Seq}(a) \to \varphi_{39})$\lbl{[S25]}}
\AXC{$\forall a\,(\mathrm{Seq}(a) \to \varphi_{39}) \to \mathrm{Seq}(b) \to \varphi_{38}$\lbl{A2}}
\BIC{$\mathrm{Seq}(b) \to \varphi_{38}$}
\AXC{$\varphi_{68}$\lbl{A1}}
\BIC{$(\mathrm{Seq}(a) \to \mathrm{Seq}(b) \to \varphi_{6}) \to \mathrm{Seq}(a) \to \varphi_{52}$}
\BIC{$\mathrm{Seq}(a) \to \varphi_{52}$}
\end{ptree}
\noindent$\Pi_{12}$:
\begin{ptree}
\AXC{$\mathrm{Seq}(a) \to \mathrm{Seq}(b) \to \varphi_{7}$\lbl{A1}}
\AXC{$\forall a\,(\mathrm{Seq}(a) \to \varphi_{39})$\lbl{[S25]}}
\AXC{$\forall a\,(\mathrm{Seq}(a) \to \varphi_{39}) \to \mathrm{Seq}(a) \to \varphi_{39}$\lbl{A2}}
\BIC{$\mathrm{Seq}(a) \to \varphi_{39}$}
\AXC{$\varphi_{69}$\lbl{A1}}
\BIC{$(\mathrm{Seq}(a) \to \mathrm{Seq}(b) \to \varphi_{7}) \to \mathrm{Seq}(a) \to \varphi_{53}$}
\BIC{$\mathrm{Seq}(a) \to \varphi_{53}$}
\end{ptree}
\noindent$\Pi_{13}$:
\begin{ptree}
\AXC{$\forall a\,\forall b\,(\mathrm{Seq}(a) \to \varphi_{25})$\lbl{[S2]}}
\AXC{$\forall a\,\forall b\,(\mathrm{Seq}(a) \to \varphi_{25}) \to \varphi_{42}$\lbl{A2}}
\BIC{$\varphi_{42}$}
\AXC{$\varphi_{42} \to \mathrm{Seq}(\Sigma a) \to \mathrm{Seq}(\Sigma b) \to \varphi_{36}$\lbl{A2}}
\BIC{$\mathrm{Seq}(\Sigma a) \to \mathrm{Seq}(\Sigma b) \to \varphi_{36}$}
\end{ptree}
\noindent$\Pi_{14}$:
\begin{ptree}
\AXC{$\stackrel{\Pi_{13}}{\vdots}$}\noLine
\UIC{$\mathrm{Seq}(\Sigma a) \to \mathrm{Seq}(\Sigma b) \to \varphi_{36}$}
\AXC{$(\mathrm{Seq}(\Sigma a) \to \mathrm{Seq}(\Sigma b) \to \varphi_{36}) \to \varphi_{66}$\lbl{A1}}
\BIC{$\varphi_{66}$}
\end{ptree}
\noindent$\Pi_{15}$:
\begin{ptree}
\AXC{$\stackrel{\Pi_{12}}{\vdots}$}\noLine
\UIC{$\mathrm{Seq}(a) \to \varphi_{53}$}
\AXC{$(\mathrm{Seq}(a) \to \varphi_{53}) \to \mathrm{Seq}(a) \to \mathrm{Seq}(b) \to \varphi_{11}$\lbl{A1}}
\BIC{$\mathrm{Seq}(a) \to \mathrm{Seq}(b) \to \varphi_{11}$}
\AXC{$\stackrel{\Pi_{14}}{\vdots}$}\noLine
\UIC{$\varphi_{66}$}
\BIC{$(\mathrm{Seq}(a) \to \mathrm{Seq}(b) \to \varphi_{10}) \to \mathrm{Seq}(a) \to \varphi_{51}$}
\end{ptree}
\noindent$\Pi_{16}$:
\begin{ptree}
\AXC{$\stackrel{\Pi_{11}}{\vdots}$}\noLine
\UIC{$\mathrm{Seq}(a) \to \varphi_{52}$}
\AXC{$(\mathrm{Seq}(a) \to \varphi_{52}) \to \mathrm{Seq}(a) \to \mathrm{Seq}(b) \to \varphi_{10}$\lbl{A1}}
\BIC{$\mathrm{Seq}(a) \to \mathrm{Seq}(b) \to \varphi_{10}$}
\AXC{$\stackrel{\Pi_{15}}{\vdots}$}\noLine
\UIC{$(\mathrm{Seq}(a) \to \mathrm{Seq}(b) \to \varphi_{10}) \to \mathrm{Seq}(a) \to \varphi_{51}$}
\BIC{$\mathrm{Seq}(a) \to \varphi_{51}$}
\end{ptree}
\noindent$\Pi_{17}$:
\begin{ptree}
\AXC{$\varphi_{47}$\lbl{A2}}
\AXC{$\stackrel{\Pi_{16}}{\vdots}$}\noLine
\UIC{$\mathrm{Seq}(a) \to \varphi_{51}$}
\AXC{$\varphi_{65}$\lbl{A1}}
\BIC{$\mathrm{Seq}(a) \to \mathrm{Seq}(b) \to \mathrm{Conv}(\Sigma a,s) \to \varphi_{37}$}
\AXC{$\varphi_{72}$\lbl{A1}}
\BIC{$\varphi_{47} \to \mathrm{Seq}(a) \to \mathrm{Seq}(b) \to \mathrm{Conv}(\Sigma a,s) \to \varphi_{35}$}
\BIC{$\mathrm{Seq}(a) \to \mathrm{Seq}(b) \to \mathrm{Conv}(\Sigma a,s) \to \varphi_{35}$}
\end{ptree}
\noindent$\Pi_{18}$:
\begin{ptree}
\AXC{$\stackrel{\Pi_{17}}{\vdots}$}\noLine
\UIC{$\mathrm{Seq}(a) \to \mathrm{Seq}(b) \to \mathrm{Conv}(\Sigma a,s) \to \varphi_{35}$}
\AXC{$(\mathrm{Seq}(a) \to \mathrm{Seq}(b) \to \mathrm{Conv}(\Sigma a,s) \to \varphi_{35}) \to \varphi_{58}$\lbl{A1}}
\BIC{$\varphi_{58}$}
\end{ptree}
\noindent$\Pi_{19}$:
\begin{ptree}
\AXC{$\stackrel{\Pi_{10}}{\vdots}$}\noLine
\UIC{$\mathrm{Seq}(a) \to \mathrm{Seq}(b) \to \mathrm{Conv}(\Sigma a,s) \to \varphi_{29}$}
\AXC{$\mathrm{Seq}(a) \to \mathrm{Seq}(b) \to \varphi_{8}$\lbl{A1}}
\AXC{$\stackrel{\Pi_{18}}{\vdots}$}\noLine
\UIC{$\varphi_{58}$}
\BIC{$\mathrm{Seq}(a) \to \mathrm{Seq}(b) \to \mathrm{Conv}(\Sigma a,s) \to \varphi_{32}$}
\AXC{$\varphi_{73}$\lbl{A1}}
\BIC{$(\mathrm{Seq}(a) \to \mathrm{Seq}(b) \to \mathrm{Conv}(\Sigma a,s) \to \varphi_{29}) \to \varphi_{64}$}
\BIC{$\varphi_{64}$}
\end{ptree}
\noindent$\Pi_{20}$:
\begin{ptree}
\AXC{$\varphi_{63}$\lbl{A3}}
\AXC{$\stackrel{\Pi_{2}}{\vdots}$}\noLine
\UIC{$\varphi_{21}$}
\AXC{$\varphi_{21} \to \varphi_{44}$\lbl{A2}}
\BIC{$\varphi_{44}$}
\AXC{$\varphi_{44} \to \Sigma a \oplus \Sigma b(n) = {\textstyle\sum_{k=1}^{n} a_{k}} + {\textstyle\sum_{k=1}^{n} b_{k}} \to \varphi_{40}$\lbl{A2}}
\BIC{$\Sigma a \oplus \Sigma b(n) = {\textstyle\sum_{k=1}^{n} a_{k}} + {\textstyle\sum_{k=1}^{n} b_{k}} \to \varphi_{40}$}
\AXC{$\stackrel{\Pi_{19}}{\vdots}$}\noLine
\UIC{$\varphi_{64}$}
\BIC{$\mathrm{Seq}(a) \to \mathrm{Seq}(b) \to \mathrm{Conv}(\Sigma a,s) \to \varphi_{26}$}
\AXC{$\varphi_{62}$\lbl{A1}}
\BIC{$\varphi_{3} \to n \in \mathbb{N} \to \Sigma a \oplus \Sigma b(n) = {\textstyle\sum_{k=1}^{n} a \oplus b(k)}$}
\UIC{$\forall n\,(\varphi_{3} \to n \in \mathbb{N} \to \Sigma a \oplus \Sigma b(n) = {\textstyle\sum_{k=1}^{n} a \oplus b(k)})$}
\AXC{$\varphi_{74}$\lbl{A1}}
\BIC{$\varphi_{63} \to \varphi_{61}$}
\BIC{$\varphi_{61}$}
\end{ptree}
\noindent$\Pi_{21}$:
\begin{ptree}
\AXC{$\forall a\,\forall b\,(\mathrm{Seq}(a) \to \varphi_{25})$\lbl{[S2]}}
\AXC{$\forall a\,\forall b\,(\mathrm{Seq}(a) \to \varphi_{25}) \to \forall b\,(\mathrm{Seq}(a) \to \varphi_{25})$\lbl{A2}}
\BIC{$\forall b\,(\mathrm{Seq}(a) \to \varphi_{25})$}
\end{ptree}
\noindent$\Pi_{22}$:
\begin{ptree}
\AXC{$\stackrel{\Pi_{21}}{\vdots}$}\noLine
\UIC{$\forall b\,(\mathrm{Seq}(a) \to \varphi_{25})$}
\AXC{$\forall b\,(\mathrm{Seq}(a) \to \varphi_{25}) \to \mathrm{Seq}(a) \to \varphi_{25}$\lbl{A2}}
\BIC{$\mathrm{Seq}(a) \to \varphi_{25}$}
\AXC{$(\mathrm{Seq}(a) \to \varphi_{25}) \to (\mathrm{Seq}(a) \to \varphi_{5}) \to \varphi_{55}$\lbl{A1}}
\BIC{$(\mathrm{Seq}(a) \to \varphi_{5}) \to \varphi_{55}$}
\end{ptree}
\noindent$\Pi_{23}$:
\begin{ptree}
\AXC{$\forall a\,(\mathrm{Seq}(a) \to \varphi_{39})$\lbl{[S25]}}
\AXC{$\forall a\,(\mathrm{Seq}(a) \to \varphi_{39}) \to \mathrm{Seq}(a \oplus b) \to \varphi_{49}$\lbl{A2}}
\BIC{$\mathrm{Seq}(a \oplus b) \to \varphi_{49}$}
\AXC{$(\mathrm{Seq}(a \oplus b) \to \varphi_{49}) \to (\mathrm{Seq}(a) \to \varphi_{12}) \to \varphi_{60}$\lbl{A1}}
\BIC{$(\mathrm{Seq}(a) \to \varphi_{12}) \to \varphi_{60}$}
\end{ptree}
\noindent$\Pi_{24}$:
\begin{ptree}
\AXC{$\mathrm{Seq}(a) \to \varphi_{4}$\lbl{A1}}
\AXC{$\mathrm{Seq}(a) \to \varphi_{5}$\lbl{A1}}
\AXC{$\stackrel{\Pi_{22}}{\vdots}$}\noLine
\UIC{$(\mathrm{Seq}(a) \to \varphi_{5}) \to \varphi_{55}$}
\BIC{$\varphi_{55}$}
\BIC{$\mathrm{Seq}(a) \to \mathrm{Seq}(b) \to \mathrm{Conv}(\Sigma a,s) \to \varphi_{31}$}
\AXC{$\varphi_{56}$\lbl{A1}}
\BIC{$\mathrm{Seq}(a) \to \varphi_{12}$}
\AXC{$\stackrel{\Pi_{23}}{\vdots}$}\noLine
\UIC{$(\mathrm{Seq}(a) \to \varphi_{12}) \to \varphi_{60}$}
\BIC{$\varphi_{60}$}
\end{ptree}
\noindent$\Pi_{25}$:
\begin{ptree}
\AXC{$\varphi_{22}$\lbl{[C13]}}
\AXC{$\varphi_{22} \to \varphi_{23}$\lbl{A2}}
\BIC{$\varphi_{23}$}
\AXC{$\varphi_{23} \to \varphi_{27}$\lbl{A2}}
\BIC{$\varphi_{27}$}
\AXC{$\varphi_{27} \to \mathrm{Seq}(v_{877}) \to \varphi_{19}$\lbl{A2}}
\BIC{$\mathrm{Seq}(v_{877}) \to \varphi_{19}$}
\UIC{$\forall v_{878}\,(\mathrm{Seq}(v_{877}) \to \varphi_{19})$}
\UIC{$\forall v_{877}\,\forall v_{878}\,(\mathrm{Seq}(v_{877}) \to \varphi_{19})$}
\UIC{$\forall v_{876}\,\forall v_{877}\,\forall v_{878}\,(\mathrm{Seq}(v_{877}) \to \varphi_{19})$}
\end{ptree}
\noindent$\Pi_{26}$:
\begin{ptree}
\AXC{$\stackrel{\Pi_{25}}{\vdots}$}\noLine
\UIC{$\forall v_{876}\,\forall v_{877}\,\forall v_{878}\,(\mathrm{Seq}(v_{877}) \to \varphi_{19})$}
\AXC{$\forall v_{876}\,\forall v_{877}\,\forall v_{878}\,(\mathrm{Seq}(v_{877}) \to \varphi_{19}) \to \varphi_{43}$\lbl{A2}}
\BIC{$\varphi_{43}$}
\AXC{$\varphi_{43} \to \varphi_{50}$\lbl{A2}}
\BIC{$\varphi_{50}$}
\end{ptree}
\noindent$\Pi_{27}$:
\begin{ptree}
\AXC{$\stackrel{\Pi_{26}}{\vdots}$}\noLine
\UIC{$\varphi_{50}$}
\AXC{$\varphi_{50} \to \mathrm{Seq}(\Sigma a \oplus b) \to \varphi_{45}$\lbl{A2}}
\BIC{$\mathrm{Seq}(\Sigma a \oplus b) \to \varphi_{45}$}
\AXC{$(\mathrm{Seq}(\Sigma a \oplus b) \to \varphi_{45}) \to \varphi_{70}$\lbl{A1}}
\BIC{$\varphi_{70}$}
\end{ptree}
\noindent$\Pi_{28}$:
\begin{ptree}
\AXC{$\varphi_{3} \to \forall n\,(\varphi_{3})$\lbl{A4}}
\AXC{$\stackrel{\Pi_{20}}{\vdots}$}\noLine
\UIC{$\varphi_{61}$}
\BIC{$\mathrm{Seq}(a) \to \mathrm{Seq}(b) \to \mathrm{Conv}(\Sigma a,s) \to \varphi_{30}$}
\AXC{$\stackrel{\Pi_{24}}{\vdots}$}\noLine
\UIC{$\varphi_{60}$}
\AXC{$\varphi_{60} \to \mathrm{Seq}(a) \to \mathrm{Seq}(b) \to \varphi_{9}$\lbl{A1}}
\BIC{$\mathrm{Seq}(a) \to \mathrm{Seq}(b) \to \varphi_{9}$}
\AXC{$\stackrel{\Pi_{27}}{\vdots}$}\noLine
\UIC{$\varphi_{70}$}
\BIC{$(\mathrm{Seq}(a) \to \mathrm{Seq}(b) \to \mathrm{Conv}(\Sigma a,s) \to \varphi_{30}) \to \varphi_{54}$}
\BIC{$\varphi_{54}$}
\end{ptree}
\noindent 主木:
\begin{ptree}
\AXC{$\mathrm{Seq}(a) \to \mathrm{Seq}(b) \to \varphi_{2}$\lbl{A1}}
\AXC{$\stackrel{\Pi_{1}}{\vdots}$}\noLine
\UIC{$\varphi_{48}$}
\BIC{$\mathrm{Seq}(a) \to \mathrm{Seq}(b) \to \varphi_{16}$}
\AXC{$\stackrel{\Pi_{28}}{\vdots}$}\noLine
\UIC{$\varphi_{54}$}
\BIC{$\mathrm{Seq}(a) \to \mathrm{Seq}(b) \to \varphi_{14}$}
\UIC{$\forall t\,(\mathrm{Seq}(a) \to \mathrm{Seq}(b) \to \varphi_{14})$}
\UIC{$\forall s\,\forall t\,(\mathrm{Seq}(a) \to \mathrm{Seq}(b) \to \varphi_{14})$}
\UIC{$\forall b\,\forall s\,\forall t\,(\mathrm{Seq}(a) \to \mathrm{Seq}(b) \to \varphi_{14})$}
\UIC{$\forall a\,\forall b\,\forall s\,\forall t\,(\mathrm{Seq}(a) \to \mathrm{Seq}(b) \to \varphi_{14})$}
\end{ptree}
\begin{abbrevlist}
\item $\varphi_{1} :\equiv \mathrm{Seq}(a) \to \mathrm{Seq}(b) \to \mathrm{Conv}(\Sigma a,s) \to \mathrm{Conv}(\Sigma a,s)$
\item $\varphi_{2} :\equiv \mathrm{Conv}(\Sigma a,s) \to \mathrm{Conv}(\Sigma b,t) \to \mathrm{Conv}(\Sigma b,t)$
\item $\varphi_{3} :\equiv \mathrm{Seq}(a) \land \mathrm{Seq}(b) \land \mathrm{Conv}(\Sigma a,s) \land \mathrm{Conv}(\Sigma b,t)$
\item $\varphi_{4} :\equiv \mathrm{Seq}(b) \to \mathrm{Conv}(\Sigma a,s) \to \mathrm{Conv}(\Sigma b,t) \to \mathrm{Seq}(b)$
\item $\varphi_{5} :\equiv \mathrm{Seq}(b) \to \mathrm{Conv}(\Sigma a,s) \to \mathrm{Conv}(\Sigma b,t) \to \mathrm{Seq}(a)$
\item $\varphi_{6} :\equiv \mathrm{Conv}(\Sigma a,s) \to \mathrm{Conv}(\Sigma b,t) \to n \in \mathbb{N} \to \mathrm{Seq}(b)$
\item $\varphi_{7} :\equiv \mathrm{Conv}(\Sigma a,s) \to \mathrm{Conv}(\Sigma b,t) \to n \in \mathbb{N} \to \mathrm{Seq}(a)$
\item $\varphi_{8} :\equiv \mathrm{Conv}(\Sigma a,s) \to \mathrm{Conv}(\Sigma b,t) \to n \in \mathbb{N} \to n \in \mathbb{N}$
\item $\varphi_{9} :\equiv \mathrm{Conv}(\Sigma a,s) \to \mathrm{Conv}(\Sigma b,t) \to \mathrm{Seq}(\Sigma a \oplus b)$
\item $\varphi_{10} :\equiv \mathrm{Conv}(\Sigma a,s) \to \mathrm{Conv}(\Sigma b,t) \to n \in \mathbb{N} \to \mathrm{Seq}(\Sigma b)$
\item $\varphi_{11} :\equiv \mathrm{Conv}(\Sigma a,s) \to \mathrm{Conv}(\Sigma b,t) \to n \in \mathbb{N} \to \mathrm{Seq}(\Sigma a)$
\item $\varphi_{12} :\equiv \mathrm{Seq}(b) \to \mathrm{Conv}(\Sigma a,s) \to \mathrm{Conv}(\Sigma b,t) \to \mathrm{Seq}(a \oplus b)$
\item $\varphi_{13} :\equiv \forall a\,\forall \alpha\,\forall b\,\forall \beta\,(\mathrm{Conv}(a,\alpha) \to \mathrm{Conv}(b,\beta) \to \mathrm{Conv}(a \oplus b,\alpha + \beta))$
\item $\varphi_{14} :\equiv \mathrm{Conv}(\Sigma a,s) \to \mathrm{Conv}(\Sigma b,t) \to \mathrm{Conv}(\Sigma a \oplus b,s + t)$
\item $\varphi_{15} :\equiv \forall b\,\forall \beta\,(\mathrm{Conv}(\Sigma a,s) \to \mathrm{Conv}(b,\beta) \to \mathrm{Conv}(\Sigma a \oplus b,s + \beta))$
\item $\varphi_{16} :\equiv \mathrm{Conv}(\Sigma a,s) \to \mathrm{Conv}(\Sigma b,t) \to \mathrm{Conv}(\Sigma a \oplus \Sigma b,s + t)$
\item $\varphi_{17} :\equiv \forall \alpha\,\forall b\,\forall \beta\,(\mathrm{Conv}(\Sigma a,\alpha) \to \mathrm{Conv}(b,\beta) \to \mathrm{Conv}(\Sigma a \oplus b,\alpha + \beta))$
\item $\varphi_{18} :\equiv \forall x\,\forall y\,(x = y \to y = x) \to \forall y\,({\textstyle\sum_{k=1}^{n} a \oplus b(k)} = y \to y = {\textstyle\sum_{k=1}^{n} a \oplus b(k)})$
\item $\varphi_{19} :\equiv \forall n \in \mathbb{N}\,(v_{876}(n) = v_{877}(n)) \to \mathrm{Conv}(v_{876},v_{878}) \to \mathrm{Conv}(v_{877},v_{878})$
\item $\varphi_{20} :\equiv \forall \beta\,(\mathrm{Conv}(\Sigma a,s) \to \mathrm{Conv}(\Sigma b,\beta) \to \mathrm{Conv}(\Sigma a \oplus \Sigma b,s + \beta))$
\item $\varphi_{21} :\equiv \forall y\,\forall z\,(\Sigma a \oplus \Sigma b(n) = y \to y = z \to \Sigma a \oplus \Sigma b(n) = z)$
\item $\varphi_{22} :\equiv \forall a\,\forall b\,\forall \alpha\,(\mathrm{Seq}(b) \to \forall n \in \mathbb{N}\,(a_{n} = b_{n}) \to \mathrm{Conv}(a,\alpha) \to \mathrm{Conv}(b,\alpha))$
\item $\varphi_{23} :\equiv \forall b\,\forall \alpha\,(\mathrm{Seq}(b) \to \forall n \in \mathbb{N}\,(v_{876}(n) = b_{n}) \to \mathrm{Conv}(v_{876},\alpha) \to \mathrm{Conv}(b,\alpha))$
\item $\varphi_{24} :\equiv \mathrm{Seq}(b) \to \forall n \in \mathbb{N}\,({\textstyle\sum_{k=1}^{n} a \oplus b(k)} = {\textstyle\sum_{k=1}^{n} a_{k}} + {\textstyle\sum_{k=1}^{n} b_{k}})$
\item $\varphi_{25} :\equiv \mathrm{Seq}(b) \to \mathrm{Seq}(a \oplus b) \land \forall n \in \mathbb{N}\,(a \oplus b(n) = a_{n} + b_{n})$
\item $\varphi_{26} :\equiv \mathrm{Conv}(\Sigma b,t) \to n \in \mathbb{N} \to \Sigma a \oplus \Sigma b(n) = {\textstyle\sum_{k=1}^{n} a \oplus b(k)}$
\item $\varphi_{27} :\equiv \forall \alpha\,(\mathrm{Seq}(v_{877}) \to \forall n \in \mathbb{N}\,(v_{876}(n) = v_{877}(n)) \to \mathrm{Conv}(v_{876},\alpha) \to \mathrm{Conv}(v_{877},\alpha))$
\item $\varphi_{28} :\equiv \mathrm{Conv}(\Sigma b,t) \to n \in \mathbb{N} \to {\textstyle\sum_{k=1}^{n} a \oplus b(k)} = {\textstyle\sum_{k=1}^{n} a_{k}} + {\textstyle\sum_{k=1}^{n} b_{k}}$
\item $\varphi_{29} :\equiv \mathrm{Conv}(\Sigma b,t) \to n \in \mathbb{N} \to {\textstyle\sum_{k=1}^{n} a_{k}} + {\textstyle\sum_{k=1}^{n} b_{k}} = {\textstyle\sum_{k=1}^{n} a \oplus b(k)}$
\item $\varphi_{30} :\equiv \mathrm{Conv}(\Sigma b,t) \to \forall n \in \mathbb{N}\,(\Sigma a \oplus \Sigma b(n) = {\textstyle\sum_{k=1}^{n} a \oplus b(k)})$
\item $\varphi_{31} :\equiv \mathrm{Conv}(\Sigma b,t) \to \mathrm{Seq}(a \oplus b) \land \forall n \in \mathbb{N}\,(a \oplus b(n) = a_{n} + b_{n})$
\item $\varphi_{32} :\equiv \mathrm{Conv}(\Sigma b,t) \to n \in \mathbb{N} \to \Sigma a \oplus \Sigma b(n) = {\textstyle\sum_{k=1}^{n} a_{k}} + {\textstyle\sum_{k=1}^{n} b_{k}}$
\item $\varphi_{33} :\equiv \mathrm{Conv}(\Sigma b,t) \to n \in \mathbb{N} \to n \in \mathbb{N} \to {\textstyle\sum_{k=1}^{n} a \oplus b(k)} = {\textstyle\sum_{k=1}^{n} a_{k}} + {\textstyle\sum_{k=1}^{n} b_{k}}$
\item $\varphi_{34} :\equiv \mathrm{Conv}(\Sigma b,t) \to n \in \mathbb{N} \to \forall n \in \mathbb{N}\,({\textstyle\sum_{k=1}^{n} a \oplus b(k)} = {\textstyle\sum_{k=1}^{n} a_{k}} + {\textstyle\sum_{k=1}^{n} b_{k}})$
\item $\varphi_{35} :\equiv \mathrm{Conv}(\Sigma b,t) \to n \in \mathbb{N} \to n \in \mathbb{N} \to \Sigma a \oplus \Sigma b(n) = {\textstyle\sum_{k=1}^{n} a_{k}} + {\textstyle\sum_{k=1}^{n} b_{k}}$
\item $\varphi_{36} :\equiv \mathrm{Seq}(\Sigma a \oplus \Sigma b) \land \forall n \in \mathbb{N}\,(\Sigma a \oplus \Sigma b(n) = {\textstyle\sum_{k=1}^{n} a_{k}} + {\textstyle\sum_{k=1}^{n} b_{k}})$
\item $\varphi_{37} :\equiv \mathrm{Conv}(\Sigma b,t) \to n \in \mathbb{N} \to \forall n \in \mathbb{N}\,(\Sigma a \oplus \Sigma b(n) = {\textstyle\sum_{k=1}^{n} a_{k}} + {\textstyle\sum_{k=1}^{n} b_{k}})$
\item $\varphi_{38} :\equiv \mathrm{Seq}(\Sigma b) \land {\textstyle\sum_{k=1}^{1_{\omega}} b_{k}} = b_{1_{\omega}} \land \forall n \in \mathbb{N}\,({\textstyle\sum_{k=1}^{S(n)} b_{k}} = {\textstyle\sum_{k=1}^{n} b_{k}} + b_{S(n)})$
\item $\varphi_{39} :\equiv \mathrm{Seq}(\Sigma a) \land {\textstyle\sum_{k=1}^{1_{\omega}} a_{k}} = a_{1_{\omega}} \land \forall n \in \mathbb{N}\,({\textstyle\sum_{k=1}^{S(n)} a_{k}} = {\textstyle\sum_{k=1}^{n} a_{k}} + a_{S(n)})$
\item $\varphi_{40} :\equiv {\textstyle\sum_{k=1}^{n} a_{k}} + {\textstyle\sum_{k=1}^{n} b_{k}} = {\textstyle\sum_{k=1}^{n} a \oplus b(k)} \to \Sigma a \oplus \Sigma b(n) = {\textstyle\sum_{k=1}^{n} a \oplus b(k)}$
\item $\varphi_{41} :\equiv {\textstyle\sum_{k=1}^{n} a \oplus b(k)} = {\textstyle\sum_{k=1}^{n} a_{k}} + {\textstyle\sum_{k=1}^{n} b_{k}} \to {\textstyle\sum_{k=1}^{n} a_{k}} + {\textstyle\sum_{k=1}^{n} b_{k}} = {\textstyle\sum_{k=1}^{n} a \oplus b(k)}$
\item $\varphi_{42} :\equiv \forall b\,(\mathrm{Seq}(\Sigma a) \to \mathrm{Seq}(b) \to \mathrm{Seq}(\Sigma a \oplus b) \land \forall n \in \mathbb{N}\,(\Sigma a \oplus b(n) = {\textstyle\sum_{k=1}^{n} a_{k}} + b_{n}))$
\item $\varphi_{43} :\equiv \forall v_{877}\,\forall v_{878}\,(\mathrm{Seq}(v_{877}) \to \forall n \in \mathbb{N}\,(\Sigma a \oplus \Sigma b(n) = v_{877}(n)) \to \mathrm{Conv}(\Sigma a \oplus \Sigma b,v_{878}) \to \mathrm{Conv}(v_{877},v_{878}))$
\item $\varphi_{44} :\equiv \forall z\,(\Sigma a \oplus \Sigma b(n) = {\textstyle\sum_{k=1}^{n} a_{k}} + {\textstyle\sum_{k=1}^{n} b_{k}} \to {\textstyle\sum_{k=1}^{n} a_{k}} + {\textstyle\sum_{k=1}^{n} b_{k}} = z \to \Sigma a \oplus \Sigma b(n) = z)$
\item $\varphi_{45} :\equiv \forall n \in \mathbb{N}\,(\Sigma a \oplus \Sigma b(n) = {\textstyle\sum_{k=1}^{n} a \oplus b(k)}) \to \mathrm{Conv}(\Sigma a \oplus \Sigma b,s + t) \to \mathrm{Conv}(\Sigma a \oplus b,s + t)$
\item $\varphi_{46} :\equiv \forall n \in \mathbb{N}\,({\textstyle\sum_{k=1}^{n} a \oplus b(k)} = {\textstyle\sum_{k=1}^{n} a_{k}} + {\textstyle\sum_{k=1}^{n} b_{k}}) \to n \in \mathbb{N} \to {\textstyle\sum_{k=1}^{n} a \oplus b(k)} = {\textstyle\sum_{k=1}^{n} a_{k}} + {\textstyle\sum_{k=1}^{n} b_{k}}$
\item $\varphi_{47} :\equiv \forall n \in \mathbb{N}\,(\Sigma a \oplus \Sigma b(n) = {\textstyle\sum_{k=1}^{n} a_{k}} + {\textstyle\sum_{k=1}^{n} b_{k}}) \to n \in \mathbb{N} \to \Sigma a \oplus \Sigma b(n) = {\textstyle\sum_{k=1}^{n} a_{k}} + {\textstyle\sum_{k=1}^{n} b_{k}}$
\item $\varphi_{48} :\equiv (\mathrm{Seq}(a) \to \mathrm{Seq}(b) \to \varphi_{2}) \to \mathrm{Seq}(a) \to \mathrm{Seq}(b) \to \varphi_{16}$
\item $\varphi_{49} :\equiv \mathrm{Seq}(\Sigma a \oplus b) \land {\textstyle\sum_{k=1}^{1_{\omega}} a \oplus b(k)} = a \oplus b(1_{\omega}) \land \forall n \in \mathbb{N}\,({\textstyle\sum_{k=1}^{S(n)} a \oplus b(k)} = {\textstyle\sum_{k=1}^{n} a \oplus b(k)} + a \oplus b(S(n)))$
\item $\varphi_{50} :\equiv \forall v_{878}\,(\mathrm{Seq}(\Sigma a \oplus b) \to \forall n \in \mathbb{N}\,(\Sigma a \oplus \Sigma b(n) = {\textstyle\sum_{k=1}^{n} a \oplus b(k)}) \to \mathrm{Conv}(\Sigma a \oplus \Sigma b,v_{878}) \to \mathrm{Conv}(\Sigma a \oplus b,v_{878}))$
\item $\varphi_{51} :\equiv \mathrm{Seq}(b) \to \mathrm{Conv}(\Sigma a,s) \to \mathrm{Conv}(\Sigma b,t) \to n \in \mathbb{N} \to \varphi_{36}$
\item $\varphi_{52} :\equiv \mathrm{Seq}(b) \to \mathrm{Conv}(\Sigma a,s) \to \mathrm{Conv}(\Sigma b,t) \to n \in \mathbb{N} \to \varphi_{38}$
\item $\varphi_{53} :\equiv \mathrm{Seq}(b) \to \mathrm{Conv}(\Sigma a,s) \to \mathrm{Conv}(\Sigma b,t) \to n \in \mathbb{N} \to \varphi_{39}$
\item $\varphi_{54} :\equiv (\mathrm{Seq}(a) \to \mathrm{Seq}(b) \to \varphi_{16}) \to \mathrm{Seq}(a) \to \mathrm{Seq}(b) \to \varphi_{14}$
\item $\varphi_{55} :\equiv (\mathrm{Seq}(a) \to \varphi_{4}) \to \mathrm{Seq}(a) \to \mathrm{Seq}(b) \to \mathrm{Conv}(\Sigma a,s) \to \varphi_{31}$
\item $\varphi_{56} :\equiv (\mathrm{Seq}(a) \to \mathrm{Seq}(b) \to \mathrm{Conv}(\Sigma a,s) \to \varphi_{31}) \to \mathrm{Seq}(a) \to \varphi_{12}$
\item $\varphi_{57} :\equiv (\mathrm{Seq}(a) \to \mathrm{Seq}(b) \to \varphi_{8}) \to \mathrm{Seq}(a) \to \mathrm{Seq}(b) \to \mathrm{Conv}(\Sigma a,s) \to \varphi_{28}$
\item $\varphi_{58} :\equiv (\mathrm{Seq}(a) \to \mathrm{Seq}(b) \to \varphi_{8}) \to \mathrm{Seq}(a) \to \mathrm{Seq}(b) \to \mathrm{Conv}(\Sigma a,s) \to \varphi_{32}$
\item $\varphi_{59} :\equiv (\mathrm{Seq}(a) \to \mathrm{Seq}(b) \to \varphi_{6}) \to \mathrm{Seq}(a) \to \mathrm{Seq}(b) \to \mathrm{Conv}(\Sigma a,s) \to \varphi_{34}$
\item $\varphi_{60} :\equiv \mathrm{Seq}(a) \to \mathrm{Seq}(b) \to \mathrm{Conv}(\Sigma a,s) \to \mathrm{Conv}(\Sigma b,t) \to \varphi_{49}$
\item $\varphi_{61} :\equiv (\varphi_{3} \to \forall n\,(\varphi_{3})) \to \mathrm{Seq}(a) \to \mathrm{Seq}(b) \to \mathrm{Conv}(\Sigma a,s) \to \varphi_{30}$
\item $\varphi_{62} :\equiv (\mathrm{Seq}(a) \to \mathrm{Seq}(b) \to \mathrm{Conv}(\Sigma a,s) \to \varphi_{26}) \to \varphi_{3} \to n \in \mathbb{N} \to \Sigma a \oplus \Sigma b(n) = {\textstyle\sum_{k=1}^{n} a \oplus b(k)}$
\item $\varphi_{63} :\equiv \forall n\,(\varphi_{3} \to n \in \mathbb{N} \to \Sigma a \oplus \Sigma b(n) = {\textstyle\sum_{k=1}^{n} a \oplus b(k)}) \to \forall n\,(\varphi_{3}) \to \forall n \in \mathbb{N}\,(\Sigma a \oplus \Sigma b(n) = {\textstyle\sum_{k=1}^{n} a \oplus b(k)})$
\item $\varphi_{64} :\equiv (\Sigma a \oplus \Sigma b(n) = {\textstyle\sum_{k=1}^{n} a_{k}} + {\textstyle\sum_{k=1}^{n} b_{k}} \to \varphi_{40}) \to \mathrm{Seq}(a) \to \mathrm{Seq}(b) \to \mathrm{Conv}(\Sigma a,s) \to \varphi_{26}$
\item $\varphi_{65} :\equiv (\mathrm{Seq}(a) \to \varphi_{51}) \to \mathrm{Seq}(a) \to \mathrm{Seq}(b) \to \mathrm{Conv}(\Sigma a,s) \to \varphi_{37}$
\item $\varphi_{66} :\equiv (\mathrm{Seq}(a) \to \mathrm{Seq}(b) \to \varphi_{11}) \to (\mathrm{Seq}(a) \to \mathrm{Seq}(b) \to \varphi_{10}) \to \mathrm{Seq}(a) \to \varphi_{51}$
\item $\varphi_{67} :\equiv (\mathrm{Seq}(a) \to \mathrm{Seq}(b) \to \mathrm{Conv}(\Sigma a,s) \to \varphi_{28}) \to \varphi_{41} \to \mathrm{Seq}(a) \to \mathrm{Seq}(b) \to \mathrm{Conv}(\Sigma a,s) \to \varphi_{29}$
\item $\varphi_{68} :\equiv (\mathrm{Seq}(b) \to \varphi_{38}) \to (\mathrm{Seq}(a) \to \mathrm{Seq}(b) \to \varphi_{6}) \to \mathrm{Seq}(a) \to \varphi_{52}$
\item $\varphi_{69} :\equiv (\mathrm{Seq}(a) \to \varphi_{39}) \to (\mathrm{Seq}(a) \to \mathrm{Seq}(b) \to \varphi_{7}) \to \mathrm{Seq}(a) \to \varphi_{53}$
\item $\varphi_{70} :\equiv (\mathrm{Seq}(a) \to \mathrm{Seq}(b) \to \varphi_{9}) \to (\mathrm{Seq}(a) \to \mathrm{Seq}(b) \to \mathrm{Conv}(\Sigma a,s) \to \varphi_{30}) \to \varphi_{54}$
\item $\varphi_{71} :\equiv (\mathrm{Seq}(a) \to \mathrm{Seq}(b) \to \mathrm{Conv}(\Sigma a,s) \to \varphi_{34}) \to \varphi_{46} \to \mathrm{Seq}(a) \to \mathrm{Seq}(b) \to \mathrm{Conv}(\Sigma a,s) \to \varphi_{33}$
\item $\varphi_{72} :\equiv (\mathrm{Seq}(a) \to \mathrm{Seq}(b) \to \mathrm{Conv}(\Sigma a,s) \to \varphi_{37}) \to \varphi_{47} \to \mathrm{Seq}(a) \to \mathrm{Seq}(b) \to \mathrm{Conv}(\Sigma a,s) \to \varphi_{35}$
\item $\varphi_{73} :\equiv (\mathrm{Seq}(a) \to \mathrm{Seq}(b) \to \mathrm{Conv}(\Sigma a,s) \to \varphi_{32}) \to (\mathrm{Seq}(a) \to \mathrm{Seq}(b) \to \mathrm{Conv}(\Sigma a,s) \to \varphi_{29}) \to \varphi_{64}$
\item $\varphi_{74} :\equiv \forall n\,(\varphi_{3} \to n \in \mathbb{N} \to \Sigma a \oplus \Sigma b(n) = {\textstyle\sum_{k=1}^{n} a \oplus b(k)}) \to \varphi_{63} \to \varphi_{61}$
\end{abbrevlist}
\end{fullproof}

% 自動生成 (prover/build.js): SerScal  ノード数 135
\begin{fullproof}{定理4.4(定数倍)}{$\forall l\,\forall a\,\forall s\,(l \in \mathbb{R} \to \mathrm{Seq}(a) \to \mathrm{Conv}(\Sigma a,s) \to \mathrm{Conv}(\Sigma l \odot a,l \cdot s))$}{135}\phantomsection\label{tree:SerScal}
\noindent$\Pi_{1}$:
\begin{ptree}
\AXC{$l \in \mathbb{R} \to \mathrm{Seq}(a) \to \mathrm{Conv}(\Sigma a,s) \to l \in \mathbb{R}$\lbl{A1}}
\AXC{$\varphi_{7}$\lbl{定理2.1(2)}}
\AXC{$\varphi_{7} \to \varphi_{5}$\lbl{A2}}
\BIC{$\varphi_{5}$}
\AXC{$\varphi_{5} \to \varphi_{12}$\lbl{A2}}
\BIC{$\varphi_{12}$}
\AXC{$\varphi_{41}$\lbl{A2}}
\BIC{$l \in \mathbb{R} \to \mathrm{Conv}(\Sigma a,s) \to \mathrm{Conv}(l \odot \Sigma a,l \cdot s)$}
\AXC{$\varphi_{52}$\lbl{A1}}
\BIC{$\varphi_{46}$}
\BIC{$\varphi_{2} \to l \in \mathbb{R} \to \varphi_{3}$}
\end{ptree}
\noindent$\Pi_{2}$:
\begin{ptree}
\AXC{$y = z \to x = y \to x = z$\lbl{A6}}
\AXC{$(y = z \to x = y \to x = z) \to x = y \to y = z \to x = z$\lbl{A1}}
\BIC{$x = y \to y = z \to x = z$}
\UIC{$\forall z\,(x = y \to y = z \to x = z)$}
\UIC{$\forall y\,\forall z\,(x = y \to y = z \to x = z)$}
\UIC{$\forall x\,\forall y\,\forall z\,(x = y \to y = z \to x = z)$}
\AXC{$\varphi_{30}$\lbl{A2}}
\BIC{$\forall y\,\forall z\,(l \odot \Sigma a(n) = y \to y = z \to l \odot \Sigma a(n) = z)$}
\AXC{$\varphi_{45}$\lbl{A2}}
\BIC{$\varphi_{33}$}
\end{ptree}
\noindent$\Pi_{3}$:
\begin{ptree}
\AXC{$x = x$\lbl{A5}}
\AXC{$x = y \to x = x \to y = x$\lbl{A6}}
\AXC{$(x = y \to x = x \to y = x) \to x = x \to x = y \to y = x$\lbl{A1}}
\BIC{$x = x \to x = y \to y = x$}
\BIC{$x = y \to y = x$}
\UIC{$\forall y\,(x = y \to y = x)$}
\UIC{$\forall x\,\forall y\,(x = y \to y = x)$}
\AXC{$\varphi_{14}$\lbl{A2}}
\BIC{$\forall y\,({\textstyle\sum_{k=1}^{n} l \odot a(k)} = y \to y = {\textstyle\sum_{k=1}^{n} l \odot a(k)})$}
\end{ptree}
\noindent$\Pi_{4}$:
\begin{ptree}
\AXC{$\stackrel{\Pi_{3}}{\vdots}$}\noLine
\UIC{$\forall y\,({\textstyle\sum_{k=1}^{n} l \odot a(k)} = y \to y = {\textstyle\sum_{k=1}^{n} l \odot a(k)})$}
\AXC{$\forall y\,({\textstyle\sum_{k=1}^{n} l \odot a(k)} = y \to y = {\textstyle\sum_{k=1}^{n} l \odot a(k)}) \to \varphi_{27}$\lbl{A2}}
\BIC{$\varphi_{27}$}
\end{ptree}
\noindent$\Pi_{5}$:
\begin{ptree}
\AXC{$\forall l\,\forall a\,(\varphi_{20})$\lbl{[Q2]}}
\AXC{$\forall l\,\forall a\,(\varphi_{20}) \to \forall a\,(\varphi_{20})$\lbl{A2}}
\BIC{$\forall a\,(\varphi_{20})$}
\AXC{$\forall a\,(\varphi_{20}) \to \varphi_{20}$\lbl{A2}}
\BIC{$\varphi_{20}$}
\AXC{$\varphi_{20} \to \varphi_{10} \to \varphi_{6} \to l \in \mathbb{R} \to \mathrm{Seq}(a) \to \varphi_{28}$\lbl{A1}}
\BIC{$\varphi_{10} \to \varphi_{6} \to l \in \mathbb{R} \to \mathrm{Seq}(a) \to \varphi_{28}$}
\end{ptree}
\noindent$\Pi_{6}$:
\begin{ptree}
\AXC{$\varphi_{9}$\lbl{A1}}
\AXC{$\varphi_{38}$\lbl{A2}}
\AXC{$\varphi_{6}$\lbl{A1}}
\AXC{$\varphi_{10}$\lbl{A1}}
\AXC{$\stackrel{\Pi_{5}}{\vdots}$}\noLine
\UIC{$\varphi_{10} \to \varphi_{6} \to l \in \mathbb{R} \to \mathrm{Seq}(a) \to \varphi_{28}$}
\BIC{$\varphi_{6} \to l \in \mathbb{R} \to \mathrm{Seq}(a) \to \varphi_{28}$}
\BIC{$l \in \mathbb{R} \to \mathrm{Seq}(a) \to \varphi_{28}$}
\AXC{$\varphi_{57}$\lbl{A1}}
\BIC{$\varphi_{38} \to l \in \mathbb{R} \to \mathrm{Seq}(a) \to \varphi_{25}$}
\BIC{$l \in \mathbb{R} \to \mathrm{Seq}(a) \to \varphi_{25}$}
\AXC{$\varphi_{55}$\lbl{A1}}
\BIC{$\varphi_{9} \to l \in \mathbb{R} \to \mathrm{Seq}(a) \to \varphi_{15}$}
\BIC{$l \in \mathbb{R} \to \mathrm{Seq}(a) \to \varphi_{15}$}
\end{ptree}
\noindent$\Pi_{7}$:
\begin{ptree}
\AXC{$\varphi_{6}$\lbl{A1}}
\AXC{$\forall a\,(\mathrm{Seq}(a) \to \varphi_{37})$\lbl{[S25]}}
\AXC{$\forall a\,(\mathrm{Seq}(a) \to \varphi_{37}) \to \mathrm{Seq}(a) \to \varphi_{37}$\lbl{A2}}
\BIC{$\mathrm{Seq}(a) \to \varphi_{37}$}
\AXC{$(\mathrm{Seq}(a) \to \varphi_{37}) \to \varphi_{49}$\lbl{A1}}
\BIC{$\varphi_{49}$}
\BIC{$l \in \mathbb{R} \to \mathrm{Seq}(a) \to \mathrm{Conv}(\Sigma a,s) \to n \in \mathbb{N} \to \varphi_{37}$}
\end{ptree}
\noindent$\Pi_{8}$:
\begin{ptree}
\AXC{$\varphi_{31}$\lbl{[S3]}}
\AXC{$\varphi_{31} \to \forall a\,(l \in \mathbb{R} \to \varphi_{13})$\lbl{A2}}
\BIC{$\forall a\,(l \in \mathbb{R} \to \varphi_{13})$}
\AXC{$\forall a\,(l \in \mathbb{R} \to \varphi_{13}) \to l \in \mathbb{R} \to \mathrm{Seq}(\Sigma a) \to \varphi_{22}$\lbl{A2}}
\BIC{$l \in \mathbb{R} \to \mathrm{Seq}(\Sigma a) \to \varphi_{22}$}
\end{ptree}
\noindent$\Pi_{9}$:
\begin{ptree}
\AXC{$\varphi_{10}$\lbl{A1}}
\AXC{$\stackrel{\Pi_{8}}{\vdots}$}\noLine
\UIC{$l \in \mathbb{R} \to \mathrm{Seq}(\Sigma a) \to \varphi_{22}$}
\AXC{$(l \in \mathbb{R} \to \mathrm{Seq}(\Sigma a) \to \varphi_{22}) \to \varphi_{10} \to \varphi_{47}$\lbl{A1}}
\BIC{$\varphi_{10} \to \varphi_{47}$}
\BIC{$\varphi_{47}$}
\end{ptree}
\noindent$\Pi_{10}$:
\begin{ptree}
\AXC{$\varphi_{39}$\lbl{A2}}
\AXC{$\stackrel{\Pi_{7}}{\vdots}$}\noLine
\UIC{$l \in \mathbb{R} \to \mathrm{Seq}(a) \to \mathrm{Conv}(\Sigma a,s) \to n \in \mathbb{N} \to \varphi_{37}$}
\AXC{$\varphi_{50}$\lbl{A1}}
\BIC{$\varphi_{11}$}
\AXC{$\stackrel{\Pi_{9}}{\vdots}$}\noLine
\UIC{$\varphi_{47}$}
\BIC{$l \in \mathbb{R} \to \mathrm{Seq}(a) \to \mathrm{Conv}(\Sigma a,s) \to n \in \mathbb{N} \to \varphi_{22}$}
\AXC{$\varphi_{54}$\lbl{A1}}
\BIC{$l \in \mathbb{R} \to \mathrm{Seq}(a) \to \varphi_{29}$}
\AXC{$\varphi_{58}$\lbl{A1}}
\BIC{$\varphi_{39} \to l \in \mathbb{R} \to \mathrm{Seq}(a) \to \varphi_{26}$}
\BIC{$l \in \mathbb{R} \to \mathrm{Seq}(a) \to \varphi_{26}$}
\end{ptree}
\noindent$\Pi_{11}$:
\begin{ptree}
\AXC{$\varphi_{9}$\lbl{A1}}
\AXC{$\stackrel{\Pi_{10}}{\vdots}$}\noLine
\UIC{$l \in \mathbb{R} \to \mathrm{Seq}(a) \to \varphi_{26}$}
\AXC{$(l \in \mathbb{R} \to \mathrm{Seq}(a) \to \varphi_{26}) \to \varphi_{9} \to l \in \mathbb{R} \to \varphi_{24}$\lbl{A1}}
\BIC{$\varphi_{9} \to l \in \mathbb{R} \to \varphi_{24}$}
\BIC{$l \in \mathbb{R} \to \varphi_{24}$}
\end{ptree}
\noindent$\Pi_{12}$:
\begin{ptree}
\AXC{$\stackrel{\Pi_{4}}{\vdots}$}\noLine
\UIC{$\varphi_{27}$}
\AXC{$\stackrel{\Pi_{6}}{\vdots}$}\noLine
\UIC{$l \in \mathbb{R} \to \mathrm{Seq}(a) \to \varphi_{15}$}
\AXC{$\varphi_{56}$\lbl{A1}}
\BIC{$\varphi_{27} \to l \in \mathbb{R} \to \mathrm{Seq}(a) \to \varphi_{16}$}
\BIC{$l \in \mathbb{R} \to \mathrm{Seq}(a) \to \varphi_{16}$}
\AXC{$\stackrel{\Pi_{11}}{\vdots}$}\noLine
\UIC{$l \in \mathbb{R} \to \varphi_{24}$}
\AXC{$(l \in \mathbb{R} \to \varphi_{24}) \to (l \in \mathbb{R} \to \mathrm{Seq}(a) \to \varphi_{16}) \to \varphi_{51}$\lbl{A1}}
\BIC{$(l \in \mathbb{R} \to \mathrm{Seq}(a) \to \varphi_{16}) \to \varphi_{51}$}
\BIC{$\varphi_{51}$}
\end{ptree}
\noindent$\Pi_{13}$:
\begin{ptree}
\AXC{$\stackrel{\Pi_{2}}{\vdots}$}\noLine
\UIC{$\varphi_{33}$}
\AXC{$\varphi_{33} \to l \odot \Sigma a(n) = l \cdot {\textstyle\sum_{k=1}^{n} a_{k}} \to \varphi_{32}$\lbl{A2}}
\BIC{$l \odot \Sigma a(n) = l \cdot {\textstyle\sum_{k=1}^{n} a_{k}} \to \varphi_{32}$}
\AXC{$\stackrel{\Pi_{12}}{\vdots}$}\noLine
\UIC{$\varphi_{51}$}
\BIC{$l \in \mathbb{R} \to \mathrm{Seq}(a) \to \varphi_{17}$}
\AXC{$(l \in \mathbb{R} \to \mathrm{Seq}(a) \to \varphi_{17}) \to \varphi_{34}$\lbl{A1}}
\BIC{$\varphi_{34}$}
\UIC{$\forall n\,(\varphi_{34})$}
\AXC{$\varphi_{60}$\lbl{A1}}
\BIC{$(\forall n\,(\varphi_{34}) \to \varphi_{36}) \to \varphi_{21} \to l \in \mathbb{R} \to \mathrm{Seq}(a) \to \varphi_{19}$}
\end{ptree}
\noindent$\Pi_{14}$:
\begin{ptree}
\AXC{$\varphi_{31}$\lbl{[S3]}}
\AXC{$\varphi_{31} \to \forall a\,(l \in \mathbb{R} \to \varphi_{13})$\lbl{A2}}
\BIC{$\forall a\,(l \in \mathbb{R} \to \varphi_{13})$}
\AXC{$\forall a\,(l \in \mathbb{R} \to \varphi_{13}) \to l \in \mathbb{R} \to \varphi_{13}$\lbl{A2}}
\BIC{$l \in \mathbb{R} \to \varphi_{13}$}
\AXC{$(l \in \mathbb{R} \to \varphi_{13}) \to \varphi_{48}$\lbl{A1}}
\BIC{$\varphi_{48}$}
\end{ptree}
\noindent$\Pi_{15}$:
\begin{ptree}
\AXC{$l \in \mathbb{R} \to \mathrm{Seq}(a) \to \mathrm{Conv}(\Sigma a,s) \to \mathrm{Seq}(a)$\lbl{A1}}
\AXC{$l \in \mathbb{R} \to \mathrm{Seq}(a) \to \mathrm{Conv}(\Sigma a,s) \to l \in \mathbb{R}$\lbl{A1}}
\AXC{$\stackrel{\Pi_{14}}{\vdots}$}\noLine
\UIC{$\varphi_{48}$}
\BIC{$\varphi_{43}$}
\BIC{$l \in \mathbb{R} \to \mathrm{Seq}(a) \to \varphi_{23}$}
\end{ptree}
\noindent$\Pi_{16}$:
\begin{ptree}
\AXC{$\forall a\,(\mathrm{Seq}(a) \to \varphi_{37})$\lbl{[S25]}}
\AXC{$\forall a\,(\mathrm{Seq}(a) \to \varphi_{37}) \to \mathrm{Seq}(l \odot a) \to \varphi_{44}$\lbl{A2}}
\BIC{$\mathrm{Seq}(l \odot a) \to \varphi_{44}$}
\AXC{$\varphi_{59}$\lbl{A1}}
\BIC{$\varphi_{1} \to l \in \mathbb{R} \to \mathrm{Seq}(a) \to \mathrm{Conv}(\Sigma a,s) \to \varphi_{44}$}
\end{ptree}
\noindent$\Pi_{17}$:
\begin{ptree}
\AXC{$\stackrel{\Pi_{15}}{\vdots}$}\noLine
\UIC{$l \in \mathbb{R} \to \mathrm{Seq}(a) \to \varphi_{23}$}
\AXC{$(l \in \mathbb{R} \to \mathrm{Seq}(a) \to \varphi_{23}) \to \varphi_{1}$\lbl{A1}}
\BIC{$\varphi_{1}$}
\AXC{$\stackrel{\Pi_{16}}{\vdots}$}\noLine
\UIC{$\varphi_{1} \to l \in \mathbb{R} \to \mathrm{Seq}(a) \to \mathrm{Conv}(\Sigma a,s) \to \varphi_{44}$}
\BIC{$l \in \mathbb{R} \to \mathrm{Seq}(a) \to \mathrm{Conv}(\Sigma a,s) \to \varphi_{44}$}
\end{ptree}
\noindent$\Pi_{18}$:
\begin{ptree}
\AXC{$\varphi_{18}$\lbl{[C13]}}
\AXC{$\varphi_{18} \to \varphi_{35}$\lbl{A2}}
\BIC{$\varphi_{35}$}
\AXC{$\varphi_{35} \to \varphi_{42}$\lbl{A2}}
\BIC{$\varphi_{42}$}
\AXC{$\varphi_{42} \to \mathrm{Seq}(\Sigma l \odot a) \to \varphi_{40}$\lbl{A2}}
\BIC{$\mathrm{Seq}(\Sigma l \odot a) \to \varphi_{40}$}
\AXC{$(\mathrm{Seq}(\Sigma l \odot a) \to \varphi_{40}) \to \varphi_{8} \to \varphi_{53}$\lbl{A1}}
\BIC{$\varphi_{8} \to \varphi_{53}$}
\end{ptree}
\noindent$\Pi_{19}$:
\begin{ptree}
\AXC{$\stackrel{\Pi_{17}}{\vdots}$}\noLine
\UIC{$l \in \mathbb{R} \to \mathrm{Seq}(a) \to \mathrm{Conv}(\Sigma a,s) \to \varphi_{44}$}
\AXC{$(l \in \mathbb{R} \to \mathrm{Seq}(a) \to \mathrm{Conv}(\Sigma a,s) \to \varphi_{44}) \to \varphi_{8}$\lbl{A1}}
\BIC{$\varphi_{8}$}
\AXC{$\stackrel{\Pi_{18}}{\vdots}$}\noLine
\UIC{$\varphi_{8} \to \varphi_{53}$}
\BIC{$\varphi_{53}$}
\end{ptree}
\noindent$\Pi_{20}$:
\begin{ptree}
\AXC{$\varphi_{21}$\lbl{A4}}
\AXC{$\forall n\,(\varphi_{34}) \to \varphi_{36}$\lbl{A3}}
\AXC{$\stackrel{\Pi_{13}}{\vdots}$}\noLine
\UIC{$(\forall n\,(\varphi_{34}) \to \varphi_{36}) \to \varphi_{21} \to l \in \mathbb{R} \to \mathrm{Seq}(a) \to \varphi_{19}$}
\BIC{$\varphi_{21} \to l \in \mathbb{R} \to \mathrm{Seq}(a) \to \varphi_{19}$}
\BIC{$l \in \mathbb{R} \to \mathrm{Seq}(a) \to \varphi_{19}$}
\AXC{$\stackrel{\Pi_{19}}{\vdots}$}\noLine
\UIC{$\varphi_{53}$}
\BIC{$(l \in \mathbb{R} \to \varphi_{3}) \to l \in \mathbb{R} \to \varphi_{4}$}
\end{ptree}
\noindent 主木:
\begin{ptree}
\AXC{$\varphi_{2}$\lbl{A1}}
\AXC{$\stackrel{\Pi_{1}}{\vdots}$}\noLine
\UIC{$\varphi_{2} \to l \in \mathbb{R} \to \varphi_{3}$}
\BIC{$l \in \mathbb{R} \to \varphi_{3}$}
\AXC{$\stackrel{\Pi_{20}}{\vdots}$}\noLine
\UIC{$(l \in \mathbb{R} \to \varphi_{3}) \to l \in \mathbb{R} \to \varphi_{4}$}
\BIC{$l \in \mathbb{R} \to \varphi_{4}$}
\UIC{$\forall s\,(l \in \mathbb{R} \to \varphi_{4})$}
\UIC{$\forall a\,\forall s\,(l \in \mathbb{R} \to \varphi_{4})$}
\UIC{$\forall l\,\forall a\,\forall s\,(l \in \mathbb{R} \to \varphi_{4})$}
\end{ptree}
\begin{abbrevlist}
\item $\varphi_{1} :\equiv l \in \mathbb{R} \to \mathrm{Seq}(a) \to \mathrm{Conv}(\Sigma a,s) \to \mathrm{Seq}(l \odot a)$
\item $\varphi_{2} :\equiv l \in \mathbb{R} \to \mathrm{Seq}(a) \to \mathrm{Conv}(\Sigma a,s) \to \mathrm{Conv}(\Sigma a,s)$
\item $\varphi_{3} :\equiv \mathrm{Seq}(a) \to \mathrm{Conv}(\Sigma a,s) \to \mathrm{Conv}(l \odot \Sigma a,l \cdot s)$
\item $\varphi_{4} :\equiv \mathrm{Seq}(a) \to \mathrm{Conv}(\Sigma a,s) \to \mathrm{Conv}(\Sigma l \odot a,l \cdot s)$
\item $\varphi_{5} :\equiv \forall a\,\forall \alpha\,(l \in \mathbb{R} \to \mathrm{Conv}(a,\alpha) \to \mathrm{Conv}(l \odot a,l \cdot \alpha))$
\item $\varphi_{6} :\equiv l \in \mathbb{R} \to \mathrm{Seq}(a) \to \mathrm{Conv}(\Sigma a,s) \to n \in \mathbb{N} \to \mathrm{Seq}(a)$
\item $\varphi_{7} :\equiv \forall c\,\forall a\,\forall \alpha\,(c \in \mathbb{R} \to \mathrm{Conv}(a,\alpha) \to \mathrm{Conv}(c \odot a,c \cdot \alpha))$
\item $\varphi_{8} :\equiv l \in \mathbb{R} \to \mathrm{Seq}(a) \to \mathrm{Conv}(\Sigma a,s) \to \mathrm{Seq}(\Sigma l \odot a)$
\item $\varphi_{9} :\equiv l \in \mathbb{R} \to \mathrm{Seq}(a) \to \mathrm{Conv}(\Sigma a,s) \to n \in \mathbb{N} \to n \in \mathbb{N}$
\item $\varphi_{10} :\equiv l \in \mathbb{R} \to \mathrm{Seq}(a) \to \mathrm{Conv}(\Sigma a,s) \to n \in \mathbb{N} \to l \in \mathbb{R}$
\item $\varphi_{11} :\equiv l \in \mathbb{R} \to \mathrm{Seq}(a) \to \mathrm{Conv}(\Sigma a,s) \to n \in \mathbb{N} \to \mathrm{Seq}(\Sigma a)$
\item $\varphi_{12} :\equiv \forall \alpha\,(l \in \mathbb{R} \to \mathrm{Conv}(\Sigma a,\alpha) \to \mathrm{Conv}(l \odot \Sigma a,l \cdot \alpha))$
\item $\varphi_{13} :\equiv \mathrm{Seq}(a) \to \mathrm{Seq}(l \odot a) \land \forall n \in \mathbb{N}\,(l \odot a(n) = l \cdot a_{n})$
\item $\varphi_{14} :\equiv \forall x\,\forall y\,(x = y \to y = x) \to \forall y\,({\textstyle\sum_{k=1}^{n} l \odot a(k)} = y \to y = {\textstyle\sum_{k=1}^{n} l \odot a(k)})$
\item $\varphi_{15} :\equiv \mathrm{Conv}(\Sigma a,s) \to n \in \mathbb{N} \to {\textstyle\sum_{k=1}^{n} l \odot a(k)} = l \cdot {\textstyle\sum_{k=1}^{n} a_{k}}$
\item $\varphi_{16} :\equiv \mathrm{Conv}(\Sigma a,s) \to n \in \mathbb{N} \to l \cdot {\textstyle\sum_{k=1}^{n} a_{k}} = {\textstyle\sum_{k=1}^{n} l \odot a(k)}$
\item $\varphi_{17} :\equiv \mathrm{Conv}(\Sigma a,s) \to n \in \mathbb{N} \to l \odot \Sigma a(n) = {\textstyle\sum_{k=1}^{n} l \odot a(k)}$
\item $\varphi_{18} :\equiv \forall a\,\forall b\,\forall \alpha\,(\mathrm{Seq}(b) \to \forall n \in \mathbb{N}\,(a_{n} = b_{n}) \to \mathrm{Conv}(a,\alpha) \to \mathrm{Conv}(b,\alpha))$
\item $\varphi_{19} :\equiv \mathrm{Conv}(\Sigma a,s) \to \forall n \in \mathbb{N}\,(l \odot \Sigma a(n) = {\textstyle\sum_{k=1}^{n} l \odot a(k)})$
\item $\varphi_{20} :\equiv l \in \mathbb{R} \to \mathrm{Seq}(a) \to \forall n \in \mathbb{N}\,({\textstyle\sum_{k=1}^{n} l \odot a(k)} = l \cdot {\textstyle\sum_{k=1}^{n} a_{k}})$
\item $\varphi_{21} :\equiv l \in \mathbb{R} \land \mathrm{Seq}(a) \land \mathrm{Conv}(\Sigma a,s) \to \forall n\,(l \in \mathbb{R} \land \mathrm{Seq}(a) \land \mathrm{Conv}(\Sigma a,s))$
\item $\varphi_{22} :\equiv \mathrm{Seq}(l \odot \Sigma a) \land \forall n \in \mathbb{N}\,(l \odot \Sigma a(n) = l \cdot {\textstyle\sum_{k=1}^{n} a_{k}})$
\item $\varphi_{23} :\equiv \mathrm{Conv}(\Sigma a,s) \to \mathrm{Seq}(l \odot a) \land \forall n \in \mathbb{N}\,(l \odot a(n) = l \cdot a_{n})$
\item $\varphi_{24} :\equiv \mathrm{Seq}(a) \to \mathrm{Conv}(\Sigma a,s) \to n \in \mathbb{N} \to l \odot \Sigma a(n) = l \cdot {\textstyle\sum_{k=1}^{n} a_{k}}$
\item $\varphi_{25} :\equiv \mathrm{Conv}(\Sigma a,s) \to n \in \mathbb{N} \to n \in \mathbb{N} \to {\textstyle\sum_{k=1}^{n} l \odot a(k)} = l \cdot {\textstyle\sum_{k=1}^{n} a_{k}}$
\item $\varphi_{26} :\equiv \mathrm{Conv}(\Sigma a,s) \to n \in \mathbb{N} \to n \in \mathbb{N} \to l \odot \Sigma a(n) = l \cdot {\textstyle\sum_{k=1}^{n} a_{k}}$
\item $\varphi_{27} :\equiv {\textstyle\sum_{k=1}^{n} l \odot a(k)} = l \cdot {\textstyle\sum_{k=1}^{n} a_{k}} \to l \cdot {\textstyle\sum_{k=1}^{n} a_{k}} = {\textstyle\sum_{k=1}^{n} l \odot a(k)}$
\item $\varphi_{28} :\equiv \mathrm{Conv}(\Sigma a,s) \to n \in \mathbb{N} \to \forall n \in \mathbb{N}\,({\textstyle\sum_{k=1}^{n} l \odot a(k)} = l \cdot {\textstyle\sum_{k=1}^{n} a_{k}})$
\item $\varphi_{29} :\equiv \mathrm{Conv}(\Sigma a,s) \to n \in \mathbb{N} \to \forall n \in \mathbb{N}\,(l \odot \Sigma a(n) = l \cdot {\textstyle\sum_{k=1}^{n} a_{k}})$
\item $\varphi_{30} :\equiv \forall x\,\forall y\,\forall z\,(x = y \to y = z \to x = z) \to \forall y\,\forall z\,(l \odot \Sigma a(n) = y \to y = z \to l \odot \Sigma a(n) = z)$
\item $\varphi_{31} :\equiv \forall c\,\forall a\,(c \in \mathbb{R} \to \mathrm{Seq}(a) \to \mathrm{Seq}(c \odot a) \land \forall n \in \mathbb{N}\,(c \odot a(n) = c \cdot a_{n}))$
\item $\varphi_{32} :\equiv l \cdot {\textstyle\sum_{k=1}^{n} a_{k}} = {\textstyle\sum_{k=1}^{n} l \odot a(k)} \to l \odot \Sigma a(n) = {\textstyle\sum_{k=1}^{n} l \odot a(k)}$
\item $\varphi_{33} :\equiv \forall z\,(l \odot \Sigma a(n) = l \cdot {\textstyle\sum_{k=1}^{n} a_{k}} \to l \cdot {\textstyle\sum_{k=1}^{n} a_{k}} = z \to l \odot \Sigma a(n) = z)$
\item $\varphi_{34} :\equiv l \in \mathbb{R} \land \mathrm{Seq}(a) \land \mathrm{Conv}(\Sigma a,s) \to n \in \mathbb{N} \to l \odot \Sigma a(n) = {\textstyle\sum_{k=1}^{n} l \odot a(k)}$
\item $\varphi_{35} :\equiv \forall b\,\forall \alpha\,(\mathrm{Seq}(b) \to \forall n \in \mathbb{N}\,(l \odot \Sigma a(n) = b_{n}) \to \mathrm{Conv}(l \odot \Sigma a,\alpha) \to \mathrm{Conv}(b,\alpha))$
\item $\varphi_{36} :\equiv \forall n\,(l \in \mathbb{R} \land \mathrm{Seq}(a) \land \mathrm{Conv}(\Sigma a,s)) \to \forall n \in \mathbb{N}\,(l \odot \Sigma a(n) = {\textstyle\sum_{k=1}^{n} l \odot a(k)})$
\item $\varphi_{37} :\equiv \mathrm{Seq}(\Sigma a) \land {\textstyle\sum_{k=1}^{1_{\omega}} a_{k}} = a_{1_{\omega}} \land \forall n \in \mathbb{N}\,({\textstyle\sum_{k=1}^{S(n)} a_{k}} = {\textstyle\sum_{k=1}^{n} a_{k}} + a_{S(n)})$
\item $\varphi_{38} :\equiv \forall n \in \mathbb{N}\,({\textstyle\sum_{k=1}^{n} l \odot a(k)} = l \cdot {\textstyle\sum_{k=1}^{n} a_{k}}) \to n \in \mathbb{N} \to {\textstyle\sum_{k=1}^{n} l \odot a(k)} = l \cdot {\textstyle\sum_{k=1}^{n} a_{k}}$
\item $\varphi_{39} :\equiv \forall n \in \mathbb{N}\,(l \odot \Sigma a(n) = l \cdot {\textstyle\sum_{k=1}^{n} a_{k}}) \to n \in \mathbb{N} \to l \odot \Sigma a(n) = l \cdot {\textstyle\sum_{k=1}^{n} a_{k}}$
\item $\varphi_{40} :\equiv \forall n \in \mathbb{N}\,(l \odot \Sigma a(n) = {\textstyle\sum_{k=1}^{n} l \odot a(k)}) \to \mathrm{Conv}(l \odot \Sigma a,l \cdot s) \to \mathrm{Conv}(\Sigma l \odot a,l \cdot s)$
\item $\varphi_{41} :\equiv \varphi_{12} \to l \in \mathbb{R} \to \mathrm{Conv}(\Sigma a,s) \to \mathrm{Conv}(l \odot \Sigma a,l \cdot s)$
\item $\varphi_{42} :\equiv \forall \alpha\,(\mathrm{Seq}(\Sigma l \odot a) \to \forall n \in \mathbb{N}\,(l \odot \Sigma a(n) = {\textstyle\sum_{k=1}^{n} l \odot a(k)}) \to \mathrm{Conv}(l \odot \Sigma a,\alpha) \to \mathrm{Conv}(\Sigma l \odot a,\alpha))$
\item $\varphi_{43} :\equiv (l \in \mathbb{R} \to \mathrm{Seq}(a) \to \mathrm{Conv}(\Sigma a,s) \to \mathrm{Seq}(a)) \to l \in \mathbb{R} \to \mathrm{Seq}(a) \to \varphi_{23}$
\item $\varphi_{44} :\equiv \mathrm{Seq}(\Sigma l \odot a) \land {\textstyle\sum_{k=1}^{1_{\omega}} l \odot a(k)} = l \odot a(1_{\omega}) \land \forall n \in \mathbb{N}\,({\textstyle\sum_{k=1}^{S(n)} l \odot a(k)} = {\textstyle\sum_{k=1}^{n} l \odot a(k)} + l \odot a(S(n)))$
\item $\varphi_{45} :\equiv \forall y\,\forall z\,(l \odot \Sigma a(n) = y \to y = z \to l \odot \Sigma a(n) = z) \to \varphi_{33}$
\item $\varphi_{46} :\equiv (l \in \mathbb{R} \to \mathrm{Seq}(a) \to \mathrm{Conv}(\Sigma a,s) \to l \in \mathbb{R}) \to \varphi_{2} \to l \in \mathbb{R} \to \varphi_{3}$
\item $\varphi_{47} :\equiv \varphi_{11} \to l \in \mathbb{R} \to \mathrm{Seq}(a) \to \mathrm{Conv}(\Sigma a,s) \to n \in \mathbb{N} \to \varphi_{22}$
\item $\varphi_{48} :\equiv (l \in \mathbb{R} \to \mathrm{Seq}(a) \to \mathrm{Conv}(\Sigma a,s) \to l \in \mathbb{R}) \to \varphi_{43}$
\item $\varphi_{49} :\equiv \varphi_{6} \to l \in \mathbb{R} \to \mathrm{Seq}(a) \to \mathrm{Conv}(\Sigma a,s) \to n \in \mathbb{N} \to \varphi_{37}$
\item $\varphi_{50} :\equiv (l \in \mathbb{R} \to \mathrm{Seq}(a) \to \mathrm{Conv}(\Sigma a,s) \to n \in \mathbb{N} \to \varphi_{37}) \to \varphi_{11}$
\item $\varphi_{51} :\equiv (l \odot \Sigma a(n) = l \cdot {\textstyle\sum_{k=1}^{n} a_{k}} \to \varphi_{32}) \to l \in \mathbb{R} \to \mathrm{Seq}(a) \to \varphi_{17}$
\item $\varphi_{52} :\equiv (l \in \mathbb{R} \to \mathrm{Conv}(\Sigma a,s) \to \mathrm{Conv}(l \odot \Sigma a,l \cdot s)) \to \varphi_{46}$
\item $\varphi_{53} :\equiv (l \in \mathbb{R} \to \mathrm{Seq}(a) \to \varphi_{19}) \to (l \in \mathbb{R} \to \varphi_{3}) \to l \in \mathbb{R} \to \varphi_{4}$
\item $\varphi_{54} :\equiv (l \in \mathbb{R} \to \mathrm{Seq}(a) \to \mathrm{Conv}(\Sigma a,s) \to n \in \mathbb{N} \to \varphi_{22}) \to l \in \mathbb{R} \to \mathrm{Seq}(a) \to \varphi_{29}$
\item $\varphi_{55} :\equiv (l \in \mathbb{R} \to \mathrm{Seq}(a) \to \varphi_{25}) \to \varphi_{9} \to l \in \mathbb{R} \to \mathrm{Seq}(a) \to \varphi_{15}$
\item $\varphi_{56} :\equiv (l \in \mathbb{R} \to \mathrm{Seq}(a) \to \varphi_{15}) \to \varphi_{27} \to l \in \mathbb{R} \to \mathrm{Seq}(a) \to \varphi_{16}$
\item $\varphi_{57} :\equiv (l \in \mathbb{R} \to \mathrm{Seq}(a) \to \varphi_{28}) \to \varphi_{38} \to l \in \mathbb{R} \to \mathrm{Seq}(a) \to \varphi_{25}$
\item $\varphi_{58} :\equiv (l \in \mathbb{R} \to \mathrm{Seq}(a) \to \varphi_{29}) \to \varphi_{39} \to l \in \mathbb{R} \to \mathrm{Seq}(a) \to \varphi_{26}$
\item $\varphi_{59} :\equiv (\mathrm{Seq}(l \odot a) \to \varphi_{44}) \to \varphi_{1} \to l \in \mathbb{R} \to \mathrm{Seq}(a) \to \mathrm{Conv}(\Sigma a,s) \to \varphi_{44}$
\item $\varphi_{60} :\equiv \forall n\,(\varphi_{34}) \to (\forall n\,(\varphi_{34}) \to \varphi_{36}) \to \varphi_{21} \to l \in \mathbb{R} \to \mathrm{Seq}(a) \to \varphi_{19}$
\end{abbrevlist}
\end{fullproof}

% 自動生成 (prover/build.js): SerTerm  ノード数 613
\begin{fullproof}{定理4.7(1)}{$\forall a\,\forall s\,(\mathrm{Seq}(a) \to \mathrm{Conv}(\Sigma a,s) \to \mathrm{Conv}(a,0))$}{613}\phantomsection\label{tree:SerTerm}
\noindent$\Pi_{1}$:
\begin{ptree}
\AXC{$z = w \to \varphi_{102}$\lbl{A6}}
\UIC{$\forall w\,(z = w \to \varphi_{102})$}
\UIC{$\forall z\,\forall w\,(z = w \to \varphi_{102})$}
\AXC{$\forall z\,\forall w\,(z = w \to \varphi_{102}) \to \varphi_{147}$\lbl{A2}}
\BIC{$\varphi_{147}$}
\AXC{$\varphi_{147} \to s + -1 \cdot s = 0 \to \varphi_{127}$\lbl{A2}}
\BIC{$s + -1 \cdot s = 0 \to \varphi_{127}$}
\end{ptree}
\noindent$\Pi_{2}$:
\begin{ptree}
\AXC{$1 \in \mathbb{R}$\lbl{(R1)}}
\AXC{$\forall x \in \mathbb{R}\,(-x \in \mathbb{R})$\lbl{(R1)}}
\AXC{$\forall x \in \mathbb{R}\,(-x \in \mathbb{R}) \to 1 \in \mathbb{R} \to -1 \in \mathbb{R}$\lbl{A2}}
\BIC{$1 \in \mathbb{R} \to -1 \in \mathbb{R}$}
\AXC{$\varphi_{48}$\lbl{A1}}
\BIC{$1 \in \mathbb{R} \to \mathrm{Seq}(a) \to \mathrm{Conv}(\Sigma a,s) \to -1 \in \mathbb{R}$}
\BIC{$\mathrm{Seq}(a) \to \mathrm{Conv}(\Sigma a,s) \to -1 \in \mathbb{R}$}
\end{ptree}
\noindent$\Pi_{3}$:
\begin{ptree}
\AXC{$\varphi_{6}$\lbl{定理2.1(2)}}
\AXC{$\varphi_{6} \to \varphi_{27}$\lbl{A2}}
\BIC{$\varphi_{27}$}
\AXC{$\varphi_{27} \to \varphi_{33}$\lbl{A2}}
\BIC{$\varphi_{33}$}
\AXC{$\varphi_{33} \to \varphi_{28}$\lbl{A2}}
\BIC{$\varphi_{28}$}
\AXC{$\varphi_{28} \to (\mathrm{Seq}(a) \to \mathrm{Conv}(\Sigma a,s) \to -1 \in \mathbb{R}) \to \varphi_{99}$\lbl{A1}}
\BIC{$(\mathrm{Seq}(a) \to \mathrm{Conv}(\Sigma a,s) \to -1 \in \mathbb{R}) \to \varphi_{99}$}
\end{ptree}
\noindent$\Pi_{4}$:
\begin{ptree}
\AXC{$\stackrel{\Pi_{2}}{\vdots}$}\noLine
\UIC{$\mathrm{Seq}(a) \to \mathrm{Conv}(\Sigma a,s) \to -1 \in \mathbb{R}$}
\AXC{$\stackrel{\Pi_{3}}{\vdots}$}\noLine
\UIC{$(\mathrm{Seq}(a) \to \mathrm{Conv}(\Sigma a,s) \to -1 \in \mathbb{R}) \to \varphi_{99}$}
\BIC{$\varphi_{99}$}
\end{ptree}
\noindent$\Pi_{5}$:
\begin{ptree}
\AXC{$\varphi_{3}$\lbl{[C18]}}
\AXC{$\varphi_{3} \to \varphi_{10}$\lbl{A2}}
\BIC{$\varphi_{10}$}
\AXC{$\varphi_{10} \to \varphi_{15}$\lbl{A2}}
\BIC{$\varphi_{15}$}
\AXC{$\varphi_{15} \to \varphi_{7}$\lbl{A2}}
\BIC{$\varphi_{7}$}
\AXC{$\varphi_{7} \to \varphi_{106}$\lbl{A1}}
\BIC{$\varphi_{106}$}
\end{ptree}
\noindent$\Pi_{6}$:
\begin{ptree}
\AXC{$\mathrm{SubIdx}(\mathsf{s})$\lbl{[C28]}}
\AXC{$\mathrm{Seq}(a) \to \mathrm{Conv}(\Sigma a,s) \to \mathrm{Conv}(\Sigma a,s)$\lbl{A1}}
\AXC{$\stackrel{\Pi_{5}}{\vdots}$}\noLine
\UIC{$\varphi_{106}$}
\BIC{$\varphi_{21}$}
\BIC{$\mathrm{Seq}(a) \to \mathrm{Conv}(\Sigma a,s) \to \mathrm{Conv}(\Sigma a \circ \mathsf{s},s)$}
\end{ptree}
\noindent$\Pi_{7}$:
\begin{ptree}
\AXC{$\varphi_{16}$\lbl{定理2.1(1)}}
\AXC{$\varphi_{16} \to \varphi_{71}$\lbl{A2}}
\BIC{$\varphi_{71}$}
\AXC{$\varphi_{71} \to \varphi_{61}$\lbl{A2}}
\BIC{$\varphi_{61}$}
\AXC{$\varphi_{61} \to \varphi_{124}$\lbl{A2}}
\BIC{$\varphi_{124}$}
\AXC{$\varphi_{124} \to \mathrm{Conv}(\Sigma a \circ \mathsf{s},s) \to \varphi_{101}$\lbl{A2}}
\BIC{$\mathrm{Conv}(\Sigma a \circ \mathsf{s},s) \to \varphi_{101}$}
\end{ptree}
\noindent$\Pi_{8}$:
\begin{ptree}
\AXC{$\stackrel{\Pi_{6}}{\vdots}$}\noLine
\UIC{$\mathrm{Seq}(a) \to \mathrm{Conv}(\Sigma a,s) \to \mathrm{Conv}(\Sigma a \circ \mathsf{s},s)$}
\AXC{$\stackrel{\Pi_{7}}{\vdots}$}\noLine
\UIC{$\mathrm{Conv}(\Sigma a \circ \mathsf{s},s) \to \varphi_{101}$}
\AXC{$(\mathrm{Conv}(\Sigma a \circ \mathsf{s},s) \to \varphi_{101}) \to \varphi_{193}$\lbl{A1}}
\BIC{$\varphi_{193}$}
\BIC{$\varphi_{14} \to \mathrm{Seq}(a) \to \varphi_{53}$}
\end{ptree}
\noindent$\Pi_{9}$:
\begin{ptree}
\AXC{$\mathrm{Seq}(a) \to \mathrm{Conv}(\Sigma a,s) \to \mathrm{Conv}(\Sigma a,s)$\lbl{A1}}
\AXC{$\stackrel{\Pi_{4}}{\vdots}$}\noLine
\UIC{$\varphi_{99}$}
\BIC{$\varphi_{14}$}
\AXC{$\stackrel{\Pi_{8}}{\vdots}$}\noLine
\UIC{$\varphi_{14} \to \mathrm{Seq}(a) \to \varphi_{53}$}
\BIC{$\mathrm{Seq}(a) \to \varphi_{53}$}
\end{ptree}
\noindent$\Pi_{10}$:
\begin{ptree}
\AXC{$y = z \to x = y \to x = z$\lbl{A6}}
\AXC{$(y = z \to x = y \to x = z) \to x = y \to y = z \to x = z$\lbl{A1}}
\BIC{$x = y \to y = z \to x = z$}
\UIC{$\forall z\,(x = y \to y = z \to x = z)$}
\UIC{$\forall y\,\forall z\,(x = y \to y = z \to x = z)$}
\UIC{$\forall x\,\forall y\,\forall z\,(x = y \to y = z \to x = z)$}
\AXC{$\varphi_{31}$\lbl{A2}}
\BIC{$\forall y\,\forall z\,(s + -1 \cdot s = y \to y = z \to s + -1 \cdot s = z)$}
\end{ptree}
\noindent$\Pi_{11}$:
\begin{ptree}
\AXC{$\stackrel{\Pi_{10}}{\vdots}$}\noLine
\UIC{$\forall y\,\forall z\,(s + -1 \cdot s = y \to y = z \to s + -1 \cdot s = z)$}
\AXC{$\forall y\,\forall z\,(s + -1 \cdot s = y \to y = z \to s + -1 \cdot s = z) \to \varphi_{22}$\lbl{A2}}
\BIC{$\varphi_{22}$}
\end{ptree}
\noindent$\Pi_{12}$:
\begin{ptree}
\AXC{$\forall a\,\forall \alpha\,(\mathrm{Conv}(a,\alpha) \to \alpha \in \mathbb{R})$\lbl{[C2]}}
\AXC{$\varphi_{23}$\lbl{A2}}
\BIC{$\forall \alpha\,(\mathrm{Conv}(\Sigma a,\alpha) \to \alpha \in \mathbb{R})$}
\AXC{$\varphi_{17}$\lbl{A2}}
\BIC{$\mathrm{Conv}(\Sigma a,s) \to s \in \mathbb{R}$}
\AXC{$(\mathrm{Conv}(\Sigma a,s) \to s \in \mathbb{R}) \to \varphi_{38}$\lbl{A1}}
\BIC{$\varphi_{38}$}
\end{ptree}
\noindent$\Pi_{13}$:
\begin{ptree}
\AXC{$\mathrm{Seq}(a) \to \mathrm{Conv}(\Sigma a,s) \to \mathrm{Conv}(\Sigma a,s)$\lbl{A1}}
\AXC{$\stackrel{\Pi_{12}}{\vdots}$}\noLine
\UIC{$\varphi_{38}$}
\BIC{$\mathrm{Seq}(a) \to \mathrm{Conv}(\Sigma a,s) \to s \in \mathbb{R}$}
\end{ptree}
\noindent$\Pi_{14}$:
\begin{ptree}
\AXC{$\forall x \in \mathbb{R}\,(x + -x = 0)$\lbl{(R5)}}
\AXC{$\forall x \in \mathbb{R}\,(x + -x = 0) \to s \in \mathbb{R} \to s + -s = 0$\lbl{A2}}
\BIC{$s \in \mathbb{R} \to s + -s = 0$}
\AXC{$(s \in \mathbb{R} \to s + -s = 0) \to \varphi_{34}$\lbl{A1}}
\BIC{$\varphi_{34}$}
\end{ptree}
\noindent$\Pi_{15}$:
\begin{ptree}
\AXC{$\forall x \in \mathbb{R}\,(-1 \cdot x = -x)$\lbl{[E50]}}
\AXC{$\forall x \in \mathbb{R}\,(-1 \cdot x = -x) \to s \in \mathbb{R} \to -1 \cdot s = -s$\lbl{A2}}
\BIC{$s \in \mathbb{R} \to -1 \cdot s = -s$}
\AXC{$(s \in \mathbb{R} \to -1 \cdot s = -s) \to \varphi_{46}$\lbl{A1}}
\BIC{$\varphi_{46}$}
\end{ptree}
\noindent$\Pi_{16}$:
\begin{ptree}
\AXC{$\stackrel{\Pi_{13}}{\vdots}$}\noLine
\UIC{$\mathrm{Seq}(a) \to \mathrm{Conv}(\Sigma a,s) \to s \in \mathbb{R}$}
\AXC{$\stackrel{\Pi_{15}}{\vdots}$}\noLine
\UIC{$\varphi_{46}$}
\BIC{$\mathrm{Seq}(a) \to \mathrm{Conv}(\Sigma a,s) \to -1 \cdot s = -s$}
\AXC{$\varphi_{201}$\lbl{A1}}
\BIC{$s + -1 \cdot s = s + -1 \cdot s \to \varphi_{142}$}
\end{ptree}
\noindent$\Pi_{17}$:
\begin{ptree}
\AXC{$x = x$\lbl{A5}}
\UIC{$\forall x\,(x = x)$}
\AXC{$\forall x\,(x = x) \to s + -1 \cdot s = s + -1 \cdot s$\lbl{A2}}
\BIC{$s + -1 \cdot s = s + -1 \cdot s$}
\AXC{$\stackrel{\Pi_{16}}{\vdots}$}\noLine
\UIC{$s + -1 \cdot s = s + -1 \cdot s \to \varphi_{142}$}
\BIC{$\varphi_{142}$}
\end{ptree}
\noindent$\Pi_{18}$:
\begin{ptree}
\AXC{$z = w \to s + -1 \cdot s = s + z \to s + -1 \cdot s = s + w$\lbl{A6}}
\UIC{$\forall w\,(z = w \to s + -1 \cdot s = s + z \to s + -1 \cdot s = s + w)$}
\UIC{$\varphi_{8}$}
\AXC{$\varphi_{8} \to \varphi_{32}$\lbl{A2}}
\BIC{$\varphi_{32}$}
\AXC{$\varphi_{32} \to \varphi_{45}$\lbl{A2}}
\BIC{$\varphi_{45}$}
\AXC{$\stackrel{\Pi_{17}}{\vdots}$}\noLine
\UIC{$\varphi_{142}$}
\BIC{$\mathrm{Seq}(a) \to \mathrm{Conv}(\Sigma a,s) \to s + -1 \cdot s = s + -s$}
\AXC{$\varphi_{189}$\lbl{A1}}
\BIC{$(\mathrm{Seq}(a) \to \mathrm{Conv}(\Sigma a,s) \to s + -s = 0) \to \varphi_{107}$}
\end{ptree}
\noindent$\Pi_{19}$:
\begin{ptree}
\AXC{$\stackrel{\Pi_{13}}{\vdots}$}\noLine
\UIC{$\mathrm{Seq}(a) \to \mathrm{Conv}(\Sigma a,s) \to s \in \mathbb{R}$}
\AXC{$\stackrel{\Pi_{14}}{\vdots}$}\noLine
\UIC{$\varphi_{34}$}
\BIC{$\mathrm{Seq}(a) \to \mathrm{Conv}(\Sigma a,s) \to s + -s = 0$}
\AXC{$\stackrel{\Pi_{18}}{\vdots}$}\noLine
\UIC{$(\mathrm{Seq}(a) \to \mathrm{Conv}(\Sigma a,s) \to s + -s = 0) \to \varphi_{107}$}
\BIC{$\varphi_{107}$}
\end{ptree}
\noindent$\Pi_{20}$:
\begin{ptree}
\AXC{$\stackrel{\Pi_{11}}{\vdots}$}\noLine
\UIC{$\varphi_{22}$}
\AXC{$\varphi_{22} \to s + -1 \cdot s = s + -s \to s + -s = 0 \to s + -1 \cdot s = 0$\lbl{A2}}
\BIC{$s + -1 \cdot s = s + -s \to s + -s = 0 \to s + -1 \cdot s = 0$}
\AXC{$\stackrel{\Pi_{19}}{\vdots}$}\noLine
\UIC{$\varphi_{107}$}
\BIC{$\mathrm{Seq}(a) \to \mathrm{Conv}(\Sigma a,s) \to s + -1 \cdot s = 0$}
\AXC{$\varphi_{218}$\lbl{A1}}
\BIC{$(\mathrm{Seq}(a) \to \varphi_{53}) \to (s + -1 \cdot s = 0 \to \varphi_{127}) \to \varphi_{36}$}
\end{ptree}
\noindent$\Pi_{21}$:
\begin{ptree}
\AXC{$\stackrel{\Pi_{1}}{\vdots}$}\noLine
\UIC{$s + -1 \cdot s = 0 \to \varphi_{127}$}
\AXC{$\stackrel{\Pi_{9}}{\vdots}$}\noLine
\UIC{$\mathrm{Seq}(a) \to \varphi_{53}$}
\AXC{$\stackrel{\Pi_{20}}{\vdots}$}\noLine
\UIC{$(\mathrm{Seq}(a) \to \varphi_{53}) \to (s + -1 \cdot s = 0 \to \varphi_{127}) \to \varphi_{36}$}
\BIC{$(s + -1 \cdot s = 0 \to \varphi_{127}) \to \varphi_{36}$}
\BIC{$\varphi_{36}$}
\end{ptree}
\noindent$\Pi_{22}$:
\begin{ptree}
\AXC{$y = z \to x = y \to x = z$\lbl{A6}}
\AXC{$(y = z \to x = y \to x = z) \to x = y \to y = z \to x = z$\lbl{A1}}
\BIC{$x = y \to y = z \to x = z$}
\UIC{$\forall z\,(x = y \to y = z \to x = z)$}
\UIC{$\forall y\,\forall z\,(x = y \to y = z \to x = z)$}
\UIC{$\forall x\,\forall y\,\forall z\,(x = y \to y = z \to x = z)$}
\AXC{$\forall x\,\forall y\,\forall z\,(x = y \to y = z \to x = z) \to \varphi_{125}$\lbl{A2}}
\BIC{$\varphi_{125}$}
\end{ptree}
\noindent$\Pi_{23}$:
\begin{ptree}
\AXC{$\stackrel{\Pi_{22}}{\vdots}$}\noLine
\UIC{$\varphi_{125}$}
\AXC{$\varphi_{125} \to \varphi_{140}$\lbl{A2}}
\BIC{$\varphi_{140}$}
\AXC{$\varphi_{140} \to \Sigma a \circ \mathsf{s} \oplus -1 \odot \Sigma a(n) = a_{S(n)} \to \varphi_{95}$\lbl{A2}}
\BIC{$\Sigma a \circ \mathsf{s} \oplus -1 \odot \Sigma a(n) = a_{S(n)} \to \varphi_{95}$}
\end{ptree}
\noindent$\Pi_{24}$:
\begin{ptree}
\AXC{$x = x$\lbl{A5}}
\AXC{$x = y \to x = x \to y = x$\lbl{A6}}
\AXC{$(x = y \to x = x \to y = x) \to x = x \to x = y \to y = x$\lbl{A1}}
\BIC{$x = x \to x = y \to y = x$}
\BIC{$x = y \to y = x$}
\UIC{$\forall y\,(x = y \to y = x)$}
\UIC{$\forall x\,\forall y\,(x = y \to y = x)$}
\end{ptree}
\noindent$\Pi_{25}$:
\begin{ptree}
\AXC{$y = z \to x = y \to x = z$\lbl{A6}}
\AXC{$(y = z \to x = y \to x = z) \to x = y \to y = z \to x = z$\lbl{A1}}
\BIC{$x = y \to y = z \to x = z$}
\UIC{$\forall z\,(x = y \to y = z \to x = z)$}
\UIC{$\forall y\,\forall z\,(x = y \to y = z \to x = z)$}
\UIC{$\forall x\,\forall y\,\forall z\,(x = y \to y = z \to x = z)$}
\AXC{$\varphi_{43}$\lbl{A2}}
\BIC{$\forall y\,\forall z\,(a \circ \mathsf{s}(n) = y \to y = z \to a \circ \mathsf{s}(n) = z)$}
\end{ptree}
\noindent$\Pi_{26}$:
\begin{ptree}
\AXC{$\stackrel{\Pi_{25}}{\vdots}$}\noLine
\UIC{$\forall y\,\forall z\,(a \circ \mathsf{s}(n) = y \to y = z \to a \circ \mathsf{s}(n) = z)$}
\AXC{$\forall y\,\forall z\,(a \circ \mathsf{s}(n) = y \to y = z \to a \circ \mathsf{s}(n) = z) \to \varphi_{54}$\lbl{A2}}
\BIC{$\varphi_{54}$}
\AXC{$\varphi_{54} \to \varphi_{83}$\lbl{A2}}
\BIC{$\varphi_{83}$}
\end{ptree}
\noindent$\Pi_{27}$:
\begin{ptree}
\AXC{$\varphi_{11}$\lbl{A2}}
\AXC{$\mathrm{NSeq}(\mathsf{s}) \land \forall k \in \mathbb{N}\,(\mathsf{s}(k) = S(k))$\lbl{[S14]}}
\AXC{$\varphi_{115}$\lbl{A1}}
\BIC{$\mathrm{Seq}(a) \to \varphi_{9}$}
\AXC{$(\mathrm{Seq}(a) \to \varphi_{9}) \to \varphi_{11} \to \varphi_{19}$\lbl{A1}}
\BIC{$\varphi_{11} \to \varphi_{19}$}
\BIC{$\varphi_{19}$}
\AXC{$\varphi_{19} \to \varphi_{87}$\lbl{A1}}
\BIC{$\varphi_{87}$}
\end{ptree}
\noindent$\Pi_{28}$:
\begin{ptree}
\AXC{$\mathrm{Seq}(a) \to \mathrm{Conv}(\Sigma a,s) \to n \in \mathbb{N} \to n \in \mathbb{N}$\lbl{A1}}
\AXC{$\stackrel{\Pi_{27}}{\vdots}$}\noLine
\UIC{$\varphi_{87}$}
\BIC{$\varphi_{5}$}
\AXC{$\varphi_{5} \to a_{\mathsf{s}(n)} = a_{\mathsf{s}(n)} \to \varphi_{152}$\lbl{A1}}
\BIC{$a_{\mathsf{s}(n)} = a_{\mathsf{s}(n)} \to \varphi_{152}$}
\end{ptree}
\noindent$\Pi_{29}$:
\begin{ptree}
\AXC{$x = x$\lbl{A5}}
\UIC{$\forall x\,(x = x)$}
\AXC{$\forall x\,(x = x) \to a_{\mathsf{s}(n)} = a_{\mathsf{s}(n)}$\lbl{A2}}
\BIC{$a_{\mathsf{s}(n)} = a_{\mathsf{s}(n)}$}
\AXC{$\stackrel{\Pi_{28}}{\vdots}$}\noLine
\UIC{$a_{\mathsf{s}(n)} = a_{\mathsf{s}(n)} \to \varphi_{152}$}
\BIC{$\varphi_{152}$}
\end{ptree}
\noindent$\Pi_{30}$:
\begin{ptree}
\AXC{$\forall a\,(\varphi_{49})$\lbl{[S12]}}
\AXC{$\forall a\,(\varphi_{49}) \to \varphi_{49}$\lbl{A2}}
\BIC{$\varphi_{49}$}
\AXC{$\varphi_{49} \to \mathrm{Seq}(a) \to \varphi_{52}$\lbl{A2}}
\BIC{$\mathrm{Seq}(a) \to \varphi_{52}$}
\AXC{$(\mathrm{Seq}(a) \to \varphi_{52}) \to \varphi_{188}$\lbl{A1}}
\BIC{$\varphi_{188}$}
\end{ptree}
\noindent$\Pi_{31}$:
\begin{ptree}
\AXC{$\mathrm{Seq}(a) \to \mathrm{Conv}(\Sigma a,s) \to n \in \mathbb{N} \to \mathrm{Seq}(a)$\lbl{A1}}
\AXC{$\stackrel{\Pi_{30}}{\vdots}$}\noLine
\UIC{$\varphi_{188}$}
\BIC{$\varphi_{155}$}
\end{ptree}
\noindent$\Pi_{32}$:
\begin{ptree}
\AXC{$\mathrm{NSeq}(\mathsf{s}) \land \forall k \in \mathbb{N}\,(\mathsf{s}(k) = S(k))$\lbl{[S14]}}
\AXC{$\varphi_{72}$\lbl{A1}}
\BIC{$\mathrm{Seq}(a) \to \mathrm{Conv}(\Sigma a,s) \to n \in \mathbb{N} \to \mathrm{NSeq}(\mathsf{s})$}
\AXC{$\stackrel{\Pi_{31}}{\vdots}$}\noLine
\UIC{$\varphi_{155}$}
\BIC{$\mathrm{Seq}(a) \to \mathrm{Conv}(\Sigma a,s) \to \varphi_{47}$}
\end{ptree}
\noindent$\Pi_{33}$:
\begin{ptree}
\AXC{$\stackrel{\Pi_{32}}{\vdots}$}\noLine
\UIC{$\mathrm{Seq}(a) \to \mathrm{Conv}(\Sigma a,s) \to \varphi_{47}$}
\AXC{$(\mathrm{Seq}(a) \to \mathrm{Conv}(\Sigma a,s) \to \varphi_{47}) \to \mathrm{Seq}(a) \to \varphi_{42}$\lbl{A1}}
\BIC{$\mathrm{Seq}(a) \to \varphi_{42}$}
\end{ptree}
\noindent$\Pi_{34}$:
\begin{ptree}
\AXC{$\varphi_{88}$\lbl{A2}}
\AXC{$\stackrel{\Pi_{33}}{\vdots}$}\noLine
\UIC{$\mathrm{Seq}(a) \to \varphi_{42}$}
\AXC{$(\mathrm{Seq}(a) \to \varphi_{42}) \to \varphi_{88} \to \mathrm{Seq}(a) \to \varphi_{35}$\lbl{A1}}
\BIC{$\varphi_{88} \to \mathrm{Seq}(a) \to \varphi_{35}$}
\BIC{$\mathrm{Seq}(a) \to \varphi_{35}$}
\AXC{$(\mathrm{Seq}(a) \to \varphi_{35}) \to \varphi_{129}$\lbl{A1}}
\BIC{$\varphi_{129}$}
\end{ptree}
\noindent$\Pi_{35}$:
\begin{ptree}
\AXC{$\mathrm{Seq}(a) \to \mathrm{Conv}(\Sigma a,s) \to n \in \mathbb{N} \to n \in \mathbb{N}$\lbl{A1}}
\AXC{$\stackrel{\Pi_{34}}{\vdots}$}\noLine
\UIC{$\varphi_{129}$}
\BIC{$\varphi_{29}$}
\AXC{$\varphi_{29} \to \varphi_{24} \to \varphi_{83} \to \varphi_{25}$\lbl{A1}}
\BIC{$\varphi_{24} \to \varphi_{83} \to \varphi_{25}$}
\end{ptree}
\noindent$\Pi_{36}$:
\begin{ptree}
\AXC{$\stackrel{\Pi_{26}}{\vdots}$}\noLine
\UIC{$\varphi_{83}$}
\AXC{$z = w \to a_{\mathsf{s}(n)} = a_{z} \to a_{\mathsf{s}(n)} = a_{w}$\lbl{A6}}
\UIC{$\forall w\,(z = w \to a_{\mathsf{s}(n)} = a_{z} \to a_{\mathsf{s}(n)} = a_{w})$}
\UIC{$\forall z\,\forall w\,(z = w \to a_{\mathsf{s}(n)} = a_{z} \to a_{\mathsf{s}(n)} = a_{w})$}
\AXC{$\varphi_{149}$\lbl{A2}}
\BIC{$\forall w\,(\mathsf{s}(n) = w \to a_{\mathsf{s}(n)} = a_{\mathsf{s}(n)} \to a_{\mathsf{s}(n)} = a_{w})$}
\AXC{$\varphi_{161}$\lbl{A2}}
\BIC{$\mathsf{s}(n) = S(n) \to a_{\mathsf{s}(n)} = a_{\mathsf{s}(n)} \to a_{\mathsf{s}(n)} = a_{S(n)}$}
\AXC{$\stackrel{\Pi_{29}}{\vdots}$}\noLine
\UIC{$\varphi_{152}$}
\BIC{$\varphi_{24}$}
\AXC{$\stackrel{\Pi_{35}}{\vdots}$}\noLine
\UIC{$\varphi_{24} \to \varphi_{83} \to \varphi_{25}$}
\BIC{$\varphi_{83} \to \varphi_{25}$}
\BIC{$\varphi_{25}$}
\end{ptree}
\noindent$\Pi_{37}$:
\begin{ptree}
\AXC{$\stackrel{\Pi_{24}}{\vdots}$}\noLine
\UIC{$\forall x\,\forall y\,(x = y \to y = x)$}
\AXC{$\varphi_{18}$\lbl{A2}}
\BIC{$\forall y\,(a \circ \mathsf{s}(n) = y \to y = a \circ \mathsf{s}(n))$}
\AXC{$\varphi_{113}$\lbl{A2}}
\BIC{$a \circ \mathsf{s}(n) = a_{S(n)} \to a_{S(n)} = a \circ \mathsf{s}(n)$}
\AXC{$\stackrel{\Pi_{36}}{\vdots}$}\noLine
\UIC{$\varphi_{25}$}
\AXC{$\varphi_{25} \to (a \circ \mathsf{s}(n) = a_{S(n)} \to a_{S(n)} = a \circ \mathsf{s}(n)) \to \varphi_{26}$\lbl{A1}}
\BIC{$(a \circ \mathsf{s}(n) = a_{S(n)} \to a_{S(n)} = a \circ \mathsf{s}(n)) \to \varphi_{26}$}
\BIC{$\varphi_{26}$}
\end{ptree}
\noindent$\Pi_{38}$:
\begin{ptree}
\AXC{$y = z \to x = y \to x = z$\lbl{A6}}
\AXC{$(y = z \to x = y \to x = z) \to x = y \to y = z \to x = z$\lbl{A1}}
\BIC{$x = y \to y = z \to x = z$}
\UIC{$\forall z\,(x = y \to y = z \to x = z)$}
\UIC{$\forall y\,\forall z\,(x = y \to y = z \to x = z)$}
\UIC{$\forall x\,\forall y\,\forall z\,(x = y \to y = z \to x = z)$}
\AXC{$\forall x\,\forall y\,\forall z\,(x = y \to y = z \to x = z) \to \varphi_{125}$\lbl{A2}}
\BIC{$\varphi_{125}$}
\end{ptree}
\noindent$\Pi_{39}$:
\begin{ptree}
\AXC{$\forall a\,(\mathrm{Seq}(a) \to \varphi_{98})$\lbl{[S25]}}
\AXC{$\forall a\,(\mathrm{Seq}(a) \to \varphi_{98}) \to \mathrm{Seq}(a) \to \varphi_{98}$\lbl{A2}}
\BIC{$\mathrm{Seq}(a) \to \varphi_{98}$}
\AXC{$(\mathrm{Seq}(a) \to \varphi_{98}) \to \varphi_{174}$\lbl{A1}}
\BIC{$\varphi_{174}$}
\end{ptree}
\noindent$\Pi_{40}$:
\begin{ptree}
\AXC{$\forall a\,\forall n\,(\mathrm{Seq}(a) \to n \in \mathbb{N} \to a_{n} \in \mathbb{R})$\lbl{[S1]}}
\AXC{$\varphi_{75}$\lbl{A2}}
\BIC{$\forall n\,(\mathrm{Seq}(\Sigma a) \to n \in \mathbb{N} \to {\textstyle\sum_{k=1}^{n} a_{k}} \in \mathbb{R})$}
\AXC{$\varphi_{89}$\lbl{A2}}
\BIC{$\mathrm{Seq}(\Sigma a) \to n \in \mathbb{N} \to {\textstyle\sum_{k=1}^{n} a_{k}} \in \mathbb{R}$}
\AXC{$(\mathrm{Seq}(\Sigma a) \to n \in \mathbb{N} \to {\textstyle\sum_{k=1}^{n} a_{k}} \in \mathbb{R}) \to \varphi_{157}$\lbl{A1}}
\BIC{$\varphi_{157}$}
\end{ptree}
\noindent$\Pi_{41}$:
\begin{ptree}
\AXC{$\mathrm{Seq}(a) \to \mathrm{Conv}(\Sigma a,s) \to n \in \mathbb{N} \to \mathrm{Seq}(a)$\lbl{A1}}
\AXC{$\stackrel{\Pi_{39}}{\vdots}$}\noLine
\UIC{$\varphi_{174}$}
\BIC{$\mathrm{Seq}(a) \to \mathrm{Conv}(\Sigma a,s) \to n \in \mathbb{N} \to \varphi_{98}$}
\AXC{$\varphi_{179}$\lbl{A1}}
\BIC{$\mathrm{Seq}(a) \to \mathrm{Conv}(\Sigma a,s) \to n \in \mathbb{N} \to \mathrm{Seq}(\Sigma a)$}
\AXC{$\stackrel{\Pi_{40}}{\vdots}$}\noLine
\UIC{$\varphi_{157}$}
\BIC{$\varphi_{108}$}
\end{ptree}
\noindent$\Pi_{42}$:
\begin{ptree}
\AXC{$\mathrm{Seq}(a) \to \mathrm{Conv}(\Sigma a,s) \to n \in \mathbb{N} \to n \in \mathbb{N}$\lbl{A1}}
\AXC{$\stackrel{\Pi_{41}}{\vdots}$}\noLine
\UIC{$\varphi_{108}$}
\BIC{$\mathrm{Seq}(a) \to \mathrm{Conv}(\Sigma a,s) \to n \in \mathbb{N} \to {\textstyle\sum_{k=1}^{n} a_{k}} \in \mathbb{R}$}
\end{ptree}
\noindent$\Pi_{43}$:
\begin{ptree}
\AXC{$\forall n \in \mathbb{N}\,(S(n) \in \mathbb{N})$\lbl{[N2]}}
\AXC{$\forall n \in \mathbb{N}\,(S(n) \in \mathbb{N}) \to n \in \mathbb{N} \to S(n) \in \mathbb{N}$\lbl{A2}}
\BIC{$n \in \mathbb{N} \to S(n) \in \mathbb{N}$}
\AXC{$(n \in \mathbb{N} \to S(n) \in \mathbb{N}) \to \varphi_{77}$\lbl{A1}}
\BIC{$\varphi_{77}$}
\end{ptree}
\noindent$\Pi_{44}$:
\begin{ptree}
\AXC{$\forall a\,\forall n\,(\mathrm{Seq}(a) \to n \in \mathbb{N} \to a_{n} \in \mathbb{R})$\lbl{[S1]}}
\AXC{$\varphi_{55}$\lbl{A2}}
\BIC{$\forall n\,(\mathrm{Seq}(a) \to n \in \mathbb{N} \to a_{n} \in \mathbb{R})$}
\AXC{$\varphi_{56}$\lbl{A2}}
\BIC{$\mathrm{Seq}(a) \to S(n) \in \mathbb{N} \to a_{S(n)} \in \mathbb{R}$}
\AXC{$(\mathrm{Seq}(a) \to S(n) \in \mathbb{N} \to a_{S(n)} \in \mathbb{R}) \to \varphi_{153}$\lbl{A1}}
\BIC{$\varphi_{153}$}
\end{ptree}
\noindent$\Pi_{45}$:
\begin{ptree}
\AXC{$\mathrm{Seq}(a) \to \mathrm{Conv}(\Sigma a,s) \to n \in \mathbb{N} \to \mathrm{Seq}(a)$\lbl{A1}}
\AXC{$\stackrel{\Pi_{44}}{\vdots}$}\noLine
\UIC{$\varphi_{153}$}
\BIC{$\varphi_{109}$}
\end{ptree}
\noindent$\Pi_{46}$:
\begin{ptree}
\AXC{$\mathrm{Seq}(a) \to \mathrm{Conv}(\Sigma a,s) \to n \in \mathbb{N} \to n \in \mathbb{N}$\lbl{A1}}
\AXC{$\stackrel{\Pi_{43}}{\vdots}$}\noLine
\UIC{$\varphi_{77}$}
\BIC{$\mathrm{Seq}(a) \to \mathrm{Conv}(\Sigma a,s) \to n \in \mathbb{N} \to S(n) \in \mathbb{N}$}
\AXC{$\stackrel{\Pi_{45}}{\vdots}$}\noLine
\UIC{$\varphi_{109}$}
\BIC{$\mathrm{Seq}(a) \to \mathrm{Conv}(\Sigma a,s) \to n \in \mathbb{N} \to a_{S(n)} \in \mathbb{R}$}
\end{ptree}
\noindent$\Pi_{47}$:
\begin{ptree}
\AXC{$\forall x\,\forall y\,(x \in \mathbb{R} \to y \in \mathbb{R} \to x + y - y = x)$\lbl{[E38]}}
\AXC{$\varphi_{120}$\lbl{A2}}
\BIC{$\forall y\,(a_{S(n)} \in \mathbb{R} \to y \in \mathbb{R} \to a_{S(n)} + y - y = a_{S(n)})$}
\AXC{$\varphi_{166}$\lbl{A2}}
\BIC{$\varphi_{96}$}
\AXC{$\varphi_{96} \to \varphi_{198}$\lbl{A1}}
\BIC{$\varphi_{198}$}
\end{ptree}
\noindent$\Pi_{48}$:
\begin{ptree}
\AXC{$\stackrel{\Pi_{42}}{\vdots}$}\noLine
\UIC{$\mathrm{Seq}(a) \to \mathrm{Conv}(\Sigma a,s) \to n \in \mathbb{N} \to {\textstyle\sum_{k=1}^{n} a_{k}} \in \mathbb{R}$}
\AXC{$\stackrel{\Pi_{46}}{\vdots}$}\noLine
\UIC{$\mathrm{Seq}(a) \to \mathrm{Conv}(\Sigma a,s) \to n \in \mathbb{N} \to a_{S(n)} \in \mathbb{R}$}
\AXC{$\stackrel{\Pi_{47}}{\vdots}$}\noLine
\UIC{$\varphi_{198}$}
\BIC{$\varphi_{154}$}
\BIC{$\mathrm{Seq}(a) \to \varphi_{58}$}
\end{ptree}
\noindent$\Pi_{49}$:
\begin{ptree}
\AXC{$y = z \to x = y \to x = z$\lbl{A6}}
\AXC{$(y = z \to x = y \to x = z) \to x = y \to y = z \to x = z$\lbl{A1}}
\BIC{$x = y \to y = z \to x = z$}
\UIC{$\forall z\,(x = y \to y = z \to x = z)$}
\UIC{$\forall y\,\forall z\,(x = y \to y = z \to x = z)$}
\UIC{$\forall x\,\forall y\,\forall z\,(x = y \to y = z \to x = z)$}
\AXC{$\forall x\,\forall y\,\forall z\,(x = y \to y = z \to x = z) \to \varphi_{125}$\lbl{A2}}
\BIC{$\varphi_{125}$}
\end{ptree}
\noindent$\Pi_{50}$:
\begin{ptree}
\AXC{$\stackrel{\Pi_{24}}{\vdots}$}\noLine
\UIC{$\forall x\,\forall y\,(x = y \to y = x)$}
\AXC{$\forall x\,\forall y\,(x = y \to y = x) \to \varphi_{80}$\lbl{A2}}
\BIC{$\varphi_{80}$}
\AXC{$\varphi_{80} \to \varphi_{181}$\lbl{A2}}
\BIC{$\varphi_{181}$}
\end{ptree}
\noindent$\Pi_{51}$:
\begin{ptree}
\AXC{$\stackrel{\Pi_{50}}{\vdots}$}\noLine
\UIC{$\varphi_{181}$}
\AXC{$\forall x\,\forall y\,(x - y = x + -y)$\lbl{Def$-$}}
\AXC{$\varphi_{81}$\lbl{A2}}
\BIC{$\forall y\,(a_{S(n)} + {\textstyle\sum_{k=1}^{n} a_{k}} - y = a_{S(n)} + {\textstyle\sum_{k=1}^{n} a_{k}} + -y)$}
\AXC{$\varphi_{162}$\lbl{A2}}
\BIC{$a_{S(n)} + {\textstyle\sum_{k=1}^{n} a_{k}} - {\textstyle\sum_{k=1}^{n} a_{k}} = a_{S(n)} + {\textstyle\sum_{k=1}^{n} a_{k}} + -{\textstyle\sum_{k=1}^{n} a_{k}}$}
\AXC{$\varphi_{221}$\lbl{A1}}
\BIC{$\varphi_{212}$}
\BIC{$a_{S(n)} + {\textstyle\sum_{k=1}^{n} a_{k}} + -{\textstyle\sum_{k=1}^{n} a_{k}} = a_{S(n)} + {\textstyle\sum_{k=1}^{n} a_{k}} - {\textstyle\sum_{k=1}^{n} a_{k}}$}
\end{ptree}
\noindent$\Pi_{52}$:
\begin{ptree}
\AXC{$y = z \to x = y \to x = z$\lbl{A6}}
\AXC{$(y = z \to x = y \to x = z) \to x = y \to y = z \to x = z$\lbl{A1}}
\BIC{$x = y \to y = z \to x = z$}
\UIC{$\forall z\,(x = y \to y = z \to x = z)$}
\UIC{$\forall y\,\forall z\,(x = y \to y = z \to x = z)$}
\UIC{$\forall x\,\forall y\,\forall z\,(x = y \to y = z \to x = z)$}
\AXC{$\forall x\,\forall y\,\forall z\,(x = y \to y = z \to x = z) \to \varphi_{125}$\lbl{A2}}
\BIC{$\varphi_{125}$}
\end{ptree}
\noindent$\Pi_{53}$:
\begin{ptree}
\AXC{$z = w \to \varphi_{148}$\lbl{A6}}
\UIC{$\forall w\,(z = w \to \varphi_{148})$}
\UIC{$\forall z\,\forall w\,(z = w \to \varphi_{148})$}
\AXC{$\forall z\,\forall w\,(z = w \to \varphi_{148}) \to \varphi_{195}$\lbl{A2}}
\BIC{$\varphi_{195}$}
\AXC{$\varphi_{195} \to {\textstyle\sum_{k=1}^{n} a_{k}} + a_{S(n)} = a_{S(n)} + {\textstyle\sum_{k=1}^{n} a_{k}} \to \varphi_{84} \to \varphi_{85}$\lbl{A2}}
\BIC{${\textstyle\sum_{k=1}^{n} a_{k}} + a_{S(n)} = a_{S(n)} + {\textstyle\sum_{k=1}^{n} a_{k}} \to \varphi_{84} \to \varphi_{85}$}
\end{ptree}
\noindent$\Pi_{54}$:
\begin{ptree}
\AXC{$\forall x\,\forall y\,(x \in \mathbb{R} \to y \in \mathbb{R} \to x + y = y + x)$\lbl{(R3)}}
\AXC{$\varphi_{130}$\lbl{A2}}
\BIC{$\forall y\,({\textstyle\sum_{k=1}^{n} a_{k}} \in \mathbb{R} \to y \in \mathbb{R} \to {\textstyle\sum_{k=1}^{n} a_{k}} + y = y + {\textstyle\sum_{k=1}^{n} a_{k}})$}
\AXC{$\varphi_{170}$\lbl{A2}}
\BIC{$\varphi_{97}$}
\AXC{$\varphi_{97} \to \varphi_{199}$\lbl{A1}}
\BIC{$\varphi_{199}$}
\end{ptree}
\noindent$\Pi_{55}$:
\begin{ptree}
\AXC{$\stackrel{\Pi_{46}}{\vdots}$}\noLine
\UIC{$\mathrm{Seq}(a) \to \mathrm{Conv}(\Sigma a,s) \to n \in \mathbb{N} \to a_{S(n)} \in \mathbb{R}$}
\AXC{$\stackrel{\Pi_{42}}{\vdots}$}\noLine
\UIC{$\mathrm{Seq}(a) \to \mathrm{Conv}(\Sigma a,s) \to n \in \mathbb{N} \to {\textstyle\sum_{k=1}^{n} a_{k}} \in \mathbb{R}$}
\AXC{$\stackrel{\Pi_{54}}{\vdots}$}\noLine
\UIC{$\varphi_{199}$}
\BIC{$\varphi_{150}$}
\BIC{$\mathrm{Seq}(a) \to \varphi_{59}$}
\end{ptree}
\noindent$\Pi_{56}$:
\begin{ptree}
\AXC{$x = x$\lbl{A5}}
\UIC{$\forall x\,(x = x)$}
\AXC{$\forall x\,(x = x) \to \varphi_{84}$\lbl{A2}}
\BIC{$\varphi_{84}$}
\AXC{$\stackrel{\Pi_{55}}{\vdots}$}\noLine
\UIC{$\mathrm{Seq}(a) \to \varphi_{59}$}
\AXC{$(\mathrm{Seq}(a) \to \varphi_{59}) \to \varphi_{84} \to \varphi_{224}$\lbl{A1}}
\BIC{$\varphi_{84} \to \varphi_{224}$}
\BIC{$\varphi_{224}$}
\end{ptree}
\noindent$\Pi_{57}$:
\begin{ptree}
\AXC{$y = z \to x = y \to x = z$\lbl{A6}}
\AXC{$(y = z \to x = y \to x = z) \to x = y \to y = z \to x = z$\lbl{A1}}
\BIC{$x = y \to y = z \to x = z$}
\UIC{$\forall z\,(x = y \to y = z \to x = z)$}
\UIC{$\forall y\,\forall z\,(x = y \to y = z \to x = z)$}
\UIC{$\forall x\,\forall y\,\forall z\,(x = y \to y = z \to x = z)$}
\AXC{$\forall x\,\forall y\,\forall z\,(x = y \to y = z \to x = z) \to \varphi_{125}$\lbl{A2}}
\BIC{$\varphi_{125}$}
\end{ptree}
\noindent$\Pi_{58}$:
\begin{ptree}
\AXC{$x = x$\lbl{A5}}
\UIC{$\forall x\,(x = x)$}
\AXC{$\forall x\,(x = x) \to {\textstyle\sum_{k=1}^{S(n)} a_{k}} + -{\textstyle\sum_{k=1}^{n} a_{k}} = {\textstyle\sum_{k=1}^{S(n)} a_{k}} + -{\textstyle\sum_{k=1}^{n} a_{k}}$\lbl{A2}}
\BIC{${\textstyle\sum_{k=1}^{S(n)} a_{k}} + -{\textstyle\sum_{k=1}^{n} a_{k}} = {\textstyle\sum_{k=1}^{S(n)} a_{k}} + -{\textstyle\sum_{k=1}^{n} a_{k}}$}
\end{ptree}
\noindent$\Pi_{59}$:
\begin{ptree}
\AXC{$\forall a\,(\mathrm{Seq}(a) \to \varphi_{98})$\lbl{[S25]}}
\AXC{$\forall a\,(\mathrm{Seq}(a) \to \varphi_{98}) \to \mathrm{Seq}(a) \to \varphi_{98}$\lbl{A2}}
\BIC{$\mathrm{Seq}(a) \to \varphi_{98}$}
\AXC{$(\mathrm{Seq}(a) \to \varphi_{98}) \to \varphi_{174}$\lbl{A1}}
\BIC{$\varphi_{174}$}
\end{ptree}
\noindent$\Pi_{60}$:
\begin{ptree}
\AXC{$\mathrm{Seq}(a) \to \mathrm{Conv}(\Sigma a,s) \to n \in \mathbb{N} \to \mathrm{Seq}(a)$\lbl{A1}}
\AXC{$\stackrel{\Pi_{59}}{\vdots}$}\noLine
\UIC{$\varphi_{174}$}
\BIC{$\mathrm{Seq}(a) \to \mathrm{Conv}(\Sigma a,s) \to n \in \mathbb{N} \to \varphi_{98}$}
\AXC{$\varphi_{208}$\lbl{A1}}
\BIC{$\mathrm{Seq}(a) \to \varphi_{64}$}
\AXC{$(\mathrm{Seq}(a) \to \varphi_{64}) \to \varphi_{122} \to \mathrm{Seq}(a) \to \varphi_{60}$\lbl{A1}}
\BIC{$\varphi_{122} \to \mathrm{Seq}(a) \to \varphi_{60}$}
\end{ptree}
\noindent$\Pi_{61}$:
\begin{ptree}
\AXC{$\mathrm{Seq}(a) \to \mathrm{Conv}(\Sigma a,s) \to n \in \mathbb{N} \to n \in \mathbb{N}$\lbl{A1}}
\AXC{$\varphi_{122}$\lbl{A2}}
\AXC{$\stackrel{\Pi_{60}}{\vdots}$}\noLine
\UIC{$\varphi_{122} \to \mathrm{Seq}(a) \to \varphi_{60}$}
\BIC{$\mathrm{Seq}(a) \to \varphi_{60}$}
\AXC{$(\mathrm{Seq}(a) \to \varphi_{60}) \to \varphi_{135}$\lbl{A1}}
\BIC{$\varphi_{135}$}
\BIC{$\varphi_{50}$}
\end{ptree}
\noindent$\Pi_{62}$:
\begin{ptree}
\AXC{$\stackrel{\Pi_{58}}{\vdots}$}\noLine
\UIC{${\textstyle\sum_{k=1}^{S(n)} a_{k}} + -{\textstyle\sum_{k=1}^{n} a_{k}} = {\textstyle\sum_{k=1}^{S(n)} a_{k}} + -{\textstyle\sum_{k=1}^{n} a_{k}}$}
\AXC{$\stackrel{\Pi_{61}}{\vdots}$}\noLine
\UIC{$\varphi_{50}$}
\AXC{$\varphi_{50} \to {\textstyle\sum_{k=1}^{S(n)} a_{k}} + -{\textstyle\sum_{k=1}^{n} a_{k}} = {\textstyle\sum_{k=1}^{S(n)} a_{k}} + -{\textstyle\sum_{k=1}^{n} a_{k}} \to \varphi_{217}$\lbl{A1}}
\BIC{${\textstyle\sum_{k=1}^{S(n)} a_{k}} + -{\textstyle\sum_{k=1}^{n} a_{k}} = {\textstyle\sum_{k=1}^{S(n)} a_{k}} + -{\textstyle\sum_{k=1}^{n} a_{k}} \to \varphi_{217}$}
\BIC{$\varphi_{217}$}
\end{ptree}
\noindent$\Pi_{63}$:
\begin{ptree}
\AXC{$z = w \to \varphi_{133}$\lbl{A6}}
\UIC{$\forall w\,(z = w \to \varphi_{133})$}
\UIC{$\forall z\,\forall w\,(z = w \to \varphi_{133})$}
\AXC{$\forall z\,\forall w\,(z = w \to \varphi_{133}) \to \varphi_{159}$\lbl{A2}}
\BIC{$\varphi_{159}$}
\AXC{$\varphi_{159} \to {\textstyle\sum_{k=1}^{S(n)} a_{k}} = {\textstyle\sum_{k=1}^{n} a_{k}} + a_{S(n)} \to \varphi_{160}$\lbl{A2}}
\BIC{${\textstyle\sum_{k=1}^{S(n)} a_{k}} = {\textstyle\sum_{k=1}^{n} a_{k}} + a_{S(n)} \to \varphi_{160}$}
\AXC{$\stackrel{\Pi_{62}}{\vdots}$}\noLine
\UIC{$\varphi_{217}$}
\BIC{$\mathrm{Seq}(a) \to \mathrm{Conv}(\Sigma a,s) \to \varphi_{86}$}
\end{ptree}
\noindent$\Pi_{64}$:
\begin{ptree}
\AXC{$y = z \to x = y \to x = z$\lbl{A6}}
\AXC{$(y = z \to x = y \to x = z) \to x = y \to y = z \to x = z$\lbl{A1}}
\BIC{$x = y \to y = z \to x = z$}
\UIC{$\forall z\,(x = y \to y = z \to x = z)$}
\UIC{$\forall y\,\forall z\,(x = y \to y = z \to x = z)$}
\UIC{$\forall x\,\forall y\,\forall z\,(x = y \to y = z \to x = z)$}
\AXC{$\forall x\,\forall y\,\forall z\,(x = y \to y = z \to x = z) \to \varphi_{125}$\lbl{A2}}
\BIC{$\varphi_{125}$}
\end{ptree}
\noindent$\Pi_{65}$:
\begin{ptree}
\AXC{$z = w \to \varphi_{136}$\lbl{A6}}
\UIC{$\forall w\,(z = w \to \varphi_{136})$}
\UIC{$\forall z\,\forall w\,(z = w \to \varphi_{136})$}
\AXC{$\forall z\,\forall w\,(z = w \to \varphi_{136}) \to \varphi_{171}$\lbl{A2}}
\BIC{$\varphi_{171}$}
\AXC{$\varphi_{171} \to -1 \cdot {\textstyle\sum_{k=1}^{n} a_{k}} = -{\textstyle\sum_{k=1}^{n} a_{k}} \to \varphi_{163}$\lbl{A2}}
\BIC{$-1 \cdot {\textstyle\sum_{k=1}^{n} a_{k}} = -{\textstyle\sum_{k=1}^{n} a_{k}} \to \varphi_{163}$}
\end{ptree}
\noindent$\Pi_{66}$:
\begin{ptree}
\AXC{$\forall x \in \mathbb{R}\,(-1 \cdot x = -x)$\lbl{[E50]}}
\AXC{$\varphi_{93}$\lbl{A2}}
\BIC{${\textstyle\sum_{k=1}^{n} a_{k}} \in \mathbb{R} \to -1 \cdot {\textstyle\sum_{k=1}^{n} a_{k}} = -{\textstyle\sum_{k=1}^{n} a_{k}}$}
\AXC{$({\textstyle\sum_{k=1}^{n} a_{k}} \in \mathbb{R} \to -1 \cdot {\textstyle\sum_{k=1}^{n} a_{k}} = -{\textstyle\sum_{k=1}^{n} a_{k}}) \to \varphi_{145}$\lbl{A1}}
\BIC{$\varphi_{145}$}
\end{ptree}
\noindent$\Pi_{67}$:
\begin{ptree}
\AXC{$\stackrel{\Pi_{42}}{\vdots}$}\noLine
\UIC{$\mathrm{Seq}(a) \to \mathrm{Conv}(\Sigma a,s) \to n \in \mathbb{N} \to {\textstyle\sum_{k=1}^{n} a_{k}} \in \mathbb{R}$}
\AXC{$\stackrel{\Pi_{66}}{\vdots}$}\noLine
\UIC{$\varphi_{145}$}
\BIC{$\varphi_{41}$}
\AXC{$\varphi_{41} \to \varphi_{228}$\lbl{A1}}
\BIC{$\varphi_{228}$}
\end{ptree}
\noindent$\Pi_{68}$:
\begin{ptree}
\AXC{$\stackrel{\Pi_{65}}{\vdots}$}\noLine
\UIC{$-1 \cdot {\textstyle\sum_{k=1}^{n} a_{k}} = -{\textstyle\sum_{k=1}^{n} a_{k}} \to \varphi_{163}$}
\AXC{$x = x$\lbl{A5}}
\UIC{$\forall x\,(x = x)$}
\AXC{$\varphi_{76}$\lbl{A2}}
\BIC{${\textstyle\sum_{k=1}^{S(n)} a_{k}} + -1 \cdot {\textstyle\sum_{k=1}^{n} a_{k}} = {\textstyle\sum_{k=1}^{S(n)} a_{k}} + -1 \cdot {\textstyle\sum_{k=1}^{n} a_{k}}$}
\AXC{$\stackrel{\Pi_{67}}{\vdots}$}\noLine
\UIC{$\varphi_{228}$}
\BIC{$(-1 \cdot {\textstyle\sum_{k=1}^{n} a_{k}} = -{\textstyle\sum_{k=1}^{n} a_{k}} \to \varphi_{163}) \to \mathrm{Seq}(a) \to \varphi_{114}$}
\BIC{$\mathrm{Seq}(a) \to \varphi_{114}$}
\end{ptree}
\noindent$\Pi_{69}$:
\begin{ptree}
\AXC{$y = z \to x = y \to x = z$\lbl{A6}}
\AXC{$(y = z \to x = y \to x = z) \to x = y \to y = z \to x = z$\lbl{A1}}
\BIC{$x = y \to y = z \to x = z$}
\UIC{$\forall z\,(x = y \to y = z \to x = z)$}
\UIC{$\forall y\,\forall z\,(x = y \to y = z \to x = z)$}
\UIC{$\forall x\,\forall y\,\forall z\,(x = y \to y = z \to x = z)$}
\AXC{$\forall x\,\forall y\,\forall z\,(x = y \to y = z \to x = z) \to \varphi_{125}$\lbl{A2}}
\BIC{$\varphi_{125}$}
\end{ptree}
\noindent$\Pi_{70}$:
\begin{ptree}
\AXC{$\forall a\,(\mathrm{Seq}(a) \to \varphi_{98})$\lbl{[S25]}}
\AXC{$\forall a\,(\mathrm{Seq}(a) \to \varphi_{98}) \to \mathrm{Seq}(a) \to \varphi_{98}$\lbl{A2}}
\BIC{$\mathrm{Seq}(a) \to \varphi_{98}$}
\AXC{$(\mathrm{Seq}(a) \to \varphi_{98}) \to \varphi_{174}$\lbl{A1}}
\BIC{$\varphi_{174}$}
\end{ptree}
\noindent$\Pi_{71}$:
\begin{ptree}
\AXC{$\mathrm{Seq}(a) \to \mathrm{Conv}(\Sigma a,s) \to n \in \mathbb{N} \to \mathrm{Seq}(a)$\lbl{A1}}
\AXC{$\stackrel{\Pi_{70}}{\vdots}$}\noLine
\UIC{$\varphi_{174}$}
\BIC{$\mathrm{Seq}(a) \to \mathrm{Conv}(\Sigma a,s) \to n \in \mathbb{N} \to \varphi_{98}$}
\AXC{$\varphi_{179}$\lbl{A1}}
\BIC{$\mathrm{Seq}(a) \to \mathrm{Conv}(\Sigma a,s) \to n \in \mathbb{N} \to \mathrm{Seq}(\Sigma a)$}
\end{ptree}
\noindent$\Pi_{72}$:
\begin{ptree}
\AXC{$\forall x \in \mathbb{R}\,(-x \in \mathbb{R})$\lbl{(R1)}}
\AXC{$\forall x \in \mathbb{R}\,(-x \in \mathbb{R}) \to 1 \in \mathbb{R} \to -1 \in \mathbb{R}$\lbl{A2}}
\BIC{$1 \in \mathbb{R} \to -1 \in \mathbb{R}$}
\AXC{$(1 \in \mathbb{R} \to -1 \in \mathbb{R}) \to \varphi_{20}$\lbl{A1}}
\BIC{$\varphi_{20}$}
\end{ptree}
\noindent$\Pi_{73}$:
\begin{ptree}
\AXC{$1 \in \mathbb{R}$\lbl{(R1)}}
\AXC{$\stackrel{\Pi_{72}}{\vdots}$}\noLine
\UIC{$\varphi_{20}$}
\BIC{$\mathrm{Seq}(a) \to \mathrm{Conv}(\Sigma a,s) \to n \in \mathbb{N} \to -1 \in \mathbb{R}$}
\end{ptree}
\noindent$\Pi_{74}$:
\begin{ptree}
\AXC{$\varphi_{62}$\lbl{[S3]}}
\AXC{$\varphi_{62} \to \varphi_{94}$\lbl{A2}}
\BIC{$\varphi_{94}$}
\AXC{$\varphi_{94} \to -1 \in \mathbb{R} \to \mathrm{Seq}(\Sigma a) \to \varphi_{68}$\lbl{A2}}
\BIC{$-1 \in \mathbb{R} \to \mathrm{Seq}(\Sigma a) \to \varphi_{68}$}
\AXC{$(-1 \in \mathbb{R} \to \mathrm{Seq}(\Sigma a) \to \varphi_{68}) \to \varphi_{207}$\lbl{A1}}
\BIC{$\varphi_{207}$}
\end{ptree}
\noindent$\Pi_{75}$:
\begin{ptree}
\AXC{$\stackrel{\Pi_{71}}{\vdots}$}\noLine
\UIC{$\mathrm{Seq}(a) \to \mathrm{Conv}(\Sigma a,s) \to n \in \mathbb{N} \to \mathrm{Seq}(\Sigma a)$}
\AXC{$\stackrel{\Pi_{73}}{\vdots}$}\noLine
\UIC{$\mathrm{Seq}(a) \to \mathrm{Conv}(\Sigma a,s) \to n \in \mathbb{N} \to -1 \in \mathbb{R}$}
\AXC{$\stackrel{\Pi_{74}}{\vdots}$}\noLine
\UIC{$\varphi_{207}$}
\BIC{$\varphi_{169}$}
\BIC{$\mathrm{Seq}(a) \to \mathrm{Conv}(\Sigma a,s) \to n \in \mathbb{N} \to \varphi_{68}$}
\end{ptree}
\noindent$\Pi_{76}$:
\begin{ptree}
\AXC{$\varphi_{134}$\lbl{A2}}
\AXC{$\stackrel{\Pi_{75}}{\vdots}$}\noLine
\UIC{$\mathrm{Seq}(a) \to \mathrm{Conv}(\Sigma a,s) \to n \in \mathbb{N} \to \varphi_{68}$}
\AXC{$\varphi_{206}$\lbl{A1}}
\BIC{$\mathrm{Seq}(a) \to \varphi_{74}$}
\AXC{$(\mathrm{Seq}(a) \to \varphi_{74}) \to \varphi_{134} \to \mathrm{Seq}(a) \to \varphi_{66}$\lbl{A1}}
\BIC{$\varphi_{134} \to \mathrm{Seq}(a) \to \varphi_{66}$}
\BIC{$\mathrm{Seq}(a) \to \varphi_{66}$}
\end{ptree}
\noindent$\Pi_{77}$:
\begin{ptree}
\AXC{$\mathrm{Seq}(a) \to \mathrm{Conv}(\Sigma a,s) \to n \in \mathbb{N} \to n \in \mathbb{N}$\lbl{A1}}
\AXC{$\stackrel{\Pi_{76}}{\vdots}$}\noLine
\UIC{$\mathrm{Seq}(a) \to \varphi_{66}$}
\AXC{$(\mathrm{Seq}(a) \to \varphi_{66}) \to \varphi_{141}$\lbl{A1}}
\BIC{$\varphi_{141}$}
\BIC{$\mathrm{Seq}(a) \to \varphi_{40}$}
\AXC{$\varphi_{231}$\lbl{A1}}
\BIC{${\textstyle\sum_{k=1}^{S(n)} a_{k}} + -1 \odot \Sigma a(n) = {\textstyle\sum_{k=1}^{S(n)} a_{k}} + -1 \odot \Sigma a(n) \to \varphi_{219}$}
\end{ptree}
\noindent$\Pi_{78}$:
\begin{ptree}
\AXC{$x = x$\lbl{A5}}
\UIC{$\forall x\,(x = x)$}
\AXC{$\varphi_{90}$\lbl{A2}}
\BIC{${\textstyle\sum_{k=1}^{S(n)} a_{k}} + -1 \odot \Sigma a(n) = {\textstyle\sum_{k=1}^{S(n)} a_{k}} + -1 \odot \Sigma a(n)$}
\AXC{$\stackrel{\Pi_{77}}{\vdots}$}\noLine
\UIC{${\textstyle\sum_{k=1}^{S(n)} a_{k}} + -1 \odot \Sigma a(n) = {\textstyle\sum_{k=1}^{S(n)} a_{k}} + -1 \odot \Sigma a(n) \to \varphi_{219}$}
\BIC{$\varphi_{219}$}
\end{ptree}
\noindent$\Pi_{79}$:
\begin{ptree}
\AXC{$z = w \to \varphi_{143}$\lbl{A6}}
\UIC{$\forall w\,(z = w \to \varphi_{143})$}
\UIC{$\forall z\,\forall w\,(z = w \to \varphi_{143})$}
\AXC{$\forall z\,\forall w\,(z = w \to \varphi_{143}) \to \varphi_{183}$\lbl{A2}}
\BIC{$\varphi_{183}$}
\AXC{$\varphi_{183} \to -1 \odot \Sigma a(n) = -1 \cdot {\textstyle\sum_{k=1}^{n} a_{k}} \to \varphi_{173}$\lbl{A2}}
\BIC{$-1 \odot \Sigma a(n) = -1 \cdot {\textstyle\sum_{k=1}^{n} a_{k}} \to \varphi_{173}$}
\AXC{$\stackrel{\Pi_{78}}{\vdots}$}\noLine
\UIC{$\varphi_{219}$}
\BIC{$\mathrm{Seq}(a) \to \mathrm{Conv}(\Sigma a,s) \to \varphi_{91}$}
\end{ptree}
\noindent$\Pi_{80}$:
\begin{ptree}
\AXC{$y = z \to x = y \to x = z$\lbl{A6}}
\AXC{$(y = z \to x = y \to x = z) \to x = y \to y = z \to x = z$\lbl{A1}}
\BIC{$x = y \to y = z \to x = z$}
\UIC{$\forall z\,(x = y \to y = z \to x = z)$}
\UIC{$\forall y\,\forall z\,(x = y \to y = z \to x = z)$}
\UIC{$\forall x\,\forall y\,\forall z\,(x = y \to y = z \to x = z)$}
\AXC{$\forall x\,\forall y\,\forall z\,(x = y \to y = z \to x = z) \to \varphi_{125}$\lbl{A2}}
\BIC{$\varphi_{125}$}
\end{ptree}
\noindent$\Pi_{81}$:
\begin{ptree}
\AXC{$z = w \to \varphi_{164}$\lbl{A6}}
\UIC{$\forall w\,(z = w \to \varphi_{164})$}
\UIC{$\forall z\,\forall w\,(z = w \to \varphi_{164})$}
\AXC{$\forall z\,\forall w\,(z = w \to \varphi_{164}) \to \varphi_{205}$\lbl{A2}}
\BIC{$\varphi_{205}$}
\AXC{$\varphi_{205} \to \Sigma a \circ \mathsf{s}(n) = {\textstyle\sum_{k=1}^{S(n)} a_{k}} \to \varphi_{194}$\lbl{A2}}
\BIC{$\Sigma a \circ \mathsf{s}(n) = {\textstyle\sum_{k=1}^{S(n)} a_{k}} \to \varphi_{194}$}
\end{ptree}
\noindent$\Pi_{82}$:
\begin{ptree}
\AXC{$y = z \to x = y \to x = z$\lbl{A6}}
\AXC{$(y = z \to x = y \to x = z) \to x = y \to y = z \to x = z$\lbl{A1}}
\BIC{$x = y \to y = z \to x = z$}
\UIC{$\forall z\,(x = y \to y = z \to x = z)$}
\UIC{$\forall y\,\forall z\,(x = y \to y = z \to x = z)$}
\UIC{$\forall x\,\forall y\,\forall z\,(x = y \to y = z \to x = z)$}
\AXC{$\varphi_{73}$\lbl{A2}}
\BIC{$\forall y\,\forall z\,(\Sigma a \circ \mathsf{s}(n) = y \to y = z \to \Sigma a \circ \mathsf{s}(n) = z)$}
\AXC{$\varphi_{168}$\lbl{A2}}
\BIC{$\varphi_{104}$}
\AXC{$\varphi_{104} \to \varphi_{139}$\lbl{A2}}
\BIC{$\varphi_{139}$}
\end{ptree}
\noindent$\Pi_{83}$:
\begin{ptree}
\AXC{$\varphi_{11}$\lbl{A2}}
\AXC{$\mathrm{NSeq}(\mathsf{s}) \land \forall k \in \mathbb{N}\,(\mathsf{s}(k) = S(k))$\lbl{[S14]}}
\AXC{$\varphi_{115}$\lbl{A1}}
\BIC{$\mathrm{Seq}(a) \to \varphi_{9}$}
\AXC{$(\mathrm{Seq}(a) \to \varphi_{9}) \to \varphi_{11} \to \varphi_{19}$\lbl{A1}}
\BIC{$\varphi_{11} \to \varphi_{19}$}
\BIC{$\varphi_{19}$}
\AXC{$\varphi_{19} \to \varphi_{87}$\lbl{A1}}
\BIC{$\varphi_{87}$}
\end{ptree}
\noindent$\Pi_{84}$:
\begin{ptree}
\AXC{$\mathrm{Seq}(a) \to \mathrm{Conv}(\Sigma a,s) \to n \in \mathbb{N} \to n \in \mathbb{N}$\lbl{A1}}
\AXC{$\stackrel{\Pi_{83}}{\vdots}$}\noLine
\UIC{$\varphi_{87}$}
\BIC{$\varphi_{5}$}
\AXC{$\varphi_{5} \to {\textstyle\sum_{k=1}^{\mathsf{s}(n)} a_{k}} = {\textstyle\sum_{k=1}^{\mathsf{s}(n)} a_{k}} \to \varphi_{176}$\lbl{A1}}
\BIC{${\textstyle\sum_{k=1}^{\mathsf{s}(n)} a_{k}} = {\textstyle\sum_{k=1}^{\mathsf{s}(n)} a_{k}} \to \varphi_{176}$}
\end{ptree}
\noindent$\Pi_{85}$:
\begin{ptree}
\AXC{$x = x$\lbl{A5}}
\UIC{$\forall x\,(x = x)$}
\AXC{$\forall x\,(x = x) \to {\textstyle\sum_{k=1}^{\mathsf{s}(n)} a_{k}} = {\textstyle\sum_{k=1}^{\mathsf{s}(n)} a_{k}}$\lbl{A2}}
\BIC{${\textstyle\sum_{k=1}^{\mathsf{s}(n)} a_{k}} = {\textstyle\sum_{k=1}^{\mathsf{s}(n)} a_{k}}$}
\AXC{$\stackrel{\Pi_{84}}{\vdots}$}\noLine
\UIC{${\textstyle\sum_{k=1}^{\mathsf{s}(n)} a_{k}} = {\textstyle\sum_{k=1}^{\mathsf{s}(n)} a_{k}} \to \varphi_{176}$}
\BIC{$\varphi_{176}$}
\end{ptree}
\noindent$\Pi_{86}$:
\begin{ptree}
\AXC{$z = w \to {\textstyle\sum_{k=1}^{\mathsf{s}(n)} a_{k}} = {\textstyle\sum_{k=1}^{z} a_{k}} \to {\textstyle\sum_{k=1}^{\mathsf{s}(n)} a_{k}} = {\textstyle\sum_{k=1}^{w} a_{k}}$\lbl{A6}}
\UIC{$\forall w\,(z = w \to {\textstyle\sum_{k=1}^{\mathsf{s}(n)} a_{k}} = {\textstyle\sum_{k=1}^{z} a_{k}} \to {\textstyle\sum_{k=1}^{\mathsf{s}(n)} a_{k}} = {\textstyle\sum_{k=1}^{w} a_{k}})$}
\UIC{$\forall z\,\forall w\,(z = w \to {\textstyle\sum_{k=1}^{\mathsf{s}(n)} a_{k}} = {\textstyle\sum_{k=1}^{z} a_{k}} \to {\textstyle\sum_{k=1}^{\mathsf{s}(n)} a_{k}} = {\textstyle\sum_{k=1}^{w} a_{k}})$}
\AXC{$\varphi_{184}$\lbl{A2}}
\BIC{$\forall w\,(\mathsf{s}(n) = w \to {\textstyle\sum_{k=1}^{\mathsf{s}(n)} a_{k}} = {\textstyle\sum_{k=1}^{\mathsf{s}(n)} a_{k}} \to {\textstyle\sum_{k=1}^{\mathsf{s}(n)} a_{k}} = {\textstyle\sum_{k=1}^{w} a_{k}})$}
\AXC{$\varphi_{200}$\lbl{A2}}
\BIC{$\mathsf{s}(n) = S(n) \to {\textstyle\sum_{k=1}^{\mathsf{s}(n)} a_{k}} = {\textstyle\sum_{k=1}^{\mathsf{s}(n)} a_{k}} \to {\textstyle\sum_{k=1}^{\mathsf{s}(n)} a_{k}} = {\textstyle\sum_{k=1}^{S(n)} a_{k}}$}
\AXC{$\stackrel{\Pi_{85}}{\vdots}$}\noLine
\UIC{$\varphi_{176}$}
\BIC{$\varphi_{37}$}
\end{ptree}
\noindent$\Pi_{87}$:
\begin{ptree}
\AXC{$\forall a\,(\mathrm{Seq}(a) \to \varphi_{98})$\lbl{[S25]}}
\AXC{$\forall a\,(\mathrm{Seq}(a) \to \varphi_{98}) \to \mathrm{Seq}(a) \to \varphi_{98}$\lbl{A2}}
\BIC{$\mathrm{Seq}(a) \to \varphi_{98}$}
\AXC{$(\mathrm{Seq}(a) \to \varphi_{98}) \to \varphi_{174}$\lbl{A1}}
\BIC{$\varphi_{174}$}
\end{ptree}
\noindent$\Pi_{88}$:
\begin{ptree}
\AXC{$\forall a\,(\varphi_{49})$\lbl{[S12]}}
\AXC{$\forall a\,(\varphi_{49}) \to \varphi_{100}$\lbl{A2}}
\BIC{$\varphi_{100}$}
\AXC{$\varphi_{100} \to \mathrm{Seq}(\Sigma a) \to \varphi_{82}$\lbl{A2}}
\BIC{$\mathrm{Seq}(\Sigma a) \to \varphi_{82}$}
\AXC{$(\mathrm{Seq}(\Sigma a) \to \varphi_{82}) \to \varphi_{204}$\lbl{A1}}
\BIC{$\varphi_{204}$}
\end{ptree}
\noindent$\Pi_{89}$:
\begin{ptree}
\AXC{$\mathrm{Seq}(a) \to \mathrm{Conv}(\Sigma a,s) \to n \in \mathbb{N} \to \mathrm{Seq}(a)$\lbl{A1}}
\AXC{$\stackrel{\Pi_{87}}{\vdots}$}\noLine
\UIC{$\varphi_{174}$}
\BIC{$\mathrm{Seq}(a) \to \mathrm{Conv}(\Sigma a,s) \to n \in \mathbb{N} \to \varphi_{98}$}
\AXC{$\varphi_{179}$\lbl{A1}}
\BIC{$\mathrm{Seq}(a) \to \mathrm{Conv}(\Sigma a,s) \to n \in \mathbb{N} \to \mathrm{Seq}(\Sigma a)$}
\AXC{$\stackrel{\Pi_{88}}{\vdots}$}\noLine
\UIC{$\varphi_{204}$}
\BIC{$\varphi_{165}$}
\end{ptree}
\noindent$\Pi_{90}$:
\begin{ptree}
\AXC{$\mathrm{NSeq}(\mathsf{s}) \land \forall k \in \mathbb{N}\,(\mathsf{s}(k) = S(k))$\lbl{[S14]}}
\AXC{$\varphi_{72}$\lbl{A1}}
\BIC{$\mathrm{Seq}(a) \to \mathrm{Conv}(\Sigma a,s) \to n \in \mathbb{N} \to \mathrm{NSeq}(\mathsf{s})$}
\AXC{$\stackrel{\Pi_{89}}{\vdots}$}\noLine
\UIC{$\varphi_{165}$}
\BIC{$\mathrm{Seq}(a) \to \mathrm{Conv}(\Sigma a,s) \to \varphi_{78}$}
\end{ptree}
\noindent$\Pi_{91}$:
\begin{ptree}
\AXC{$\stackrel{\Pi_{90}}{\vdots}$}\noLine
\UIC{$\mathrm{Seq}(a) \to \mathrm{Conv}(\Sigma a,s) \to \varphi_{78}$}
\AXC{$(\mathrm{Seq}(a) \to \mathrm{Conv}(\Sigma a,s) \to \varphi_{78}) \to \mathrm{Seq}(a) \to \varphi_{69}$\lbl{A1}}
\BIC{$\mathrm{Seq}(a) \to \varphi_{69}$}
\end{ptree}
\noindent$\Pi_{92}$:
\begin{ptree}
\AXC{$\varphi_{131}$\lbl{A2}}
\AXC{$\stackrel{\Pi_{91}}{\vdots}$}\noLine
\UIC{$\mathrm{Seq}(a) \to \varphi_{69}$}
\AXC{$(\mathrm{Seq}(a) \to \varphi_{69}) \to \varphi_{131} \to \mathrm{Seq}(a) \to \varphi_{65}$\lbl{A1}}
\BIC{$\varphi_{131} \to \mathrm{Seq}(a) \to \varphi_{65}$}
\BIC{$\mathrm{Seq}(a) \to \varphi_{65}$}
\AXC{$(\mathrm{Seq}(a) \to \varphi_{65}) \to \varphi_{137}$\lbl{A1}}
\BIC{$\varphi_{137}$}
\end{ptree}
\noindent$\Pi_{93}$:
\begin{ptree}
\AXC{$\stackrel{\Pi_{82}}{\vdots}$}\noLine
\UIC{$\varphi_{139}$}
\AXC{$\stackrel{\Pi_{86}}{\vdots}$}\noLine
\UIC{$\varphi_{37}$}
\AXC{$\mathrm{Seq}(a) \to \mathrm{Conv}(\Sigma a,s) \to n \in \mathbb{N} \to n \in \mathbb{N}$\lbl{A1}}
\AXC{$\stackrel{\Pi_{92}}{\vdots}$}\noLine
\UIC{$\varphi_{137}$}
\BIC{$\varphi_{57}$}
\AXC{$\varphi_{57} \to \varphi_{37} \to \varphi_{139} \to \varphi_{39}$\lbl{A1}}
\BIC{$\varphi_{37} \to \varphi_{139} \to \varphi_{39}$}
\BIC{$\varphi_{139} \to \varphi_{39}$}
\BIC{$\varphi_{39}$}
\end{ptree}
\noindent$\Pi_{94}$:
\begin{ptree}
\AXC{$\forall a\,(\mathrm{Seq}(a) \to \varphi_{98})$\lbl{[S25]}}
\AXC{$\forall a\,(\mathrm{Seq}(a) \to \varphi_{98}) \to \mathrm{Seq}(a) \to \varphi_{98}$\lbl{A2}}
\BIC{$\mathrm{Seq}(a) \to \varphi_{98}$}
\AXC{$(\mathrm{Seq}(a) \to \varphi_{98}) \to \varphi_{174}$\lbl{A1}}
\BIC{$\varphi_{174}$}
\end{ptree}
\noindent$\Pi_{95}$:
\begin{ptree}
\AXC{$\mathrm{Seq}(a) \to \mathrm{Conv}(\Sigma a,s) \to n \in \mathbb{N} \to \mathrm{Seq}(a)$\lbl{A1}}
\AXC{$\stackrel{\Pi_{94}}{\vdots}$}\noLine
\UIC{$\varphi_{174}$}
\BIC{$\mathrm{Seq}(a) \to \mathrm{Conv}(\Sigma a,s) \to n \in \mathbb{N} \to \varphi_{98}$}
\AXC{$\varphi_{179}$\lbl{A1}}
\BIC{$\mathrm{Seq}(a) \to \mathrm{Conv}(\Sigma a,s) \to n \in \mathbb{N} \to \mathrm{Seq}(\Sigma a)$}
\end{ptree}
\noindent$\Pi_{96}$:
\begin{ptree}
\AXC{$\varphi_{62}$\lbl{[S3]}}
\AXC{$\varphi_{62} \to \varphi_{94}$\lbl{A2}}
\BIC{$\varphi_{94}$}
\AXC{$\varphi_{94} \to -1 \in \mathbb{R} \to \mathrm{Seq}(\Sigma a) \to \varphi_{68}$\lbl{A2}}
\BIC{$-1 \in \mathbb{R} \to \mathrm{Seq}(\Sigma a) \to \varphi_{68}$}
\AXC{$(-1 \in \mathbb{R} \to \mathrm{Seq}(\Sigma a) \to \varphi_{68}) \to \varphi_{207}$\lbl{A1}}
\BIC{$\varphi_{207}$}
\end{ptree}
\noindent$\Pi_{97}$:
\begin{ptree}
\AXC{$\stackrel{\Pi_{95}}{\vdots}$}\noLine
\UIC{$\mathrm{Seq}(a) \to \mathrm{Conv}(\Sigma a,s) \to n \in \mathbb{N} \to \mathrm{Seq}(\Sigma a)$}
\AXC{$\stackrel{\Pi_{73}}{\vdots}$}\noLine
\UIC{$\mathrm{Seq}(a) \to \mathrm{Conv}(\Sigma a,s) \to n \in \mathbb{N} \to -1 \in \mathbb{R}$}
\AXC{$\stackrel{\Pi_{96}}{\vdots}$}\noLine
\UIC{$\varphi_{207}$}
\BIC{$\varphi_{169}$}
\BIC{$\mathrm{Seq}(a) \to \mathrm{Conv}(\Sigma a,s) \to n \in \mathbb{N} \to \varphi_{68}$}
\end{ptree}
\noindent$\Pi_{98}$:
\begin{ptree}
\AXC{$\stackrel{\Pi_{97}}{\vdots}$}\noLine
\UIC{$\mathrm{Seq}(a) \to \mathrm{Conv}(\Sigma a,s) \to n \in \mathbb{N} \to \varphi_{68}$}
\AXC{$(\mathrm{Seq}(a) \to \mathrm{Conv}(\Sigma a,s) \to n \in \mathbb{N} \to \varphi_{68}) \to \varphi_{13}$\lbl{A1}}
\BIC{$\varphi_{13}$}
\end{ptree}
\noindent$\Pi_{99}$:
\begin{ptree}
\AXC{$\forall a\,(\mathrm{Seq}(a) \to \varphi_{98})$\lbl{[S25]}}
\AXC{$\forall a\,(\mathrm{Seq}(a) \to \varphi_{98}) \to \mathrm{Seq}(a) \to \varphi_{98}$\lbl{A2}}
\BIC{$\mathrm{Seq}(a) \to \varphi_{98}$}
\AXC{$(\mathrm{Seq}(a) \to \varphi_{98}) \to \varphi_{174}$\lbl{A1}}
\BIC{$\varphi_{174}$}
\end{ptree}
\noindent$\Pi_{100}$:
\begin{ptree}
\AXC{$\forall a\,(\varphi_{49})$\lbl{[S12]}}
\AXC{$\forall a\,(\varphi_{49}) \to \varphi_{100}$\lbl{A2}}
\BIC{$\varphi_{100}$}
\AXC{$\varphi_{100} \to \mathrm{Seq}(\Sigma a) \to \varphi_{82}$\lbl{A2}}
\BIC{$\mathrm{Seq}(\Sigma a) \to \varphi_{82}$}
\AXC{$(\mathrm{Seq}(\Sigma a) \to \varphi_{82}) \to \varphi_{204}$\lbl{A1}}
\BIC{$\varphi_{204}$}
\end{ptree}
\noindent$\Pi_{101}$:
\begin{ptree}
\AXC{$\mathrm{Seq}(a) \to \mathrm{Conv}(\Sigma a,s) \to n \in \mathbb{N} \to \mathrm{Seq}(a)$\lbl{A1}}
\AXC{$\stackrel{\Pi_{99}}{\vdots}$}\noLine
\UIC{$\varphi_{174}$}
\BIC{$\mathrm{Seq}(a) \to \mathrm{Conv}(\Sigma a,s) \to n \in \mathbb{N} \to \varphi_{98}$}
\AXC{$\varphi_{179}$\lbl{A1}}
\BIC{$\mathrm{Seq}(a) \to \mathrm{Conv}(\Sigma a,s) \to n \in \mathbb{N} \to \mathrm{Seq}(\Sigma a)$}
\AXC{$\stackrel{\Pi_{100}}{\vdots}$}\noLine
\UIC{$\varphi_{204}$}
\BIC{$\varphi_{165}$}
\end{ptree}
\noindent$\Pi_{102}$:
\begin{ptree}
\AXC{$\mathrm{NSeq}(\mathsf{s}) \land \forall k \in \mathbb{N}\,(\mathsf{s}(k) = S(k))$\lbl{[S14]}}
\AXC{$\varphi_{72}$\lbl{A1}}
\BIC{$\mathrm{Seq}(a) \to \mathrm{Conv}(\Sigma a,s) \to n \in \mathbb{N} \to \mathrm{NSeq}(\mathsf{s})$}
\AXC{$\stackrel{\Pi_{101}}{\vdots}$}\noLine
\UIC{$\varphi_{165}$}
\BIC{$\mathrm{Seq}(a) \to \mathrm{Conv}(\Sigma a,s) \to \varphi_{78}$}
\AXC{$(\mathrm{Seq}(a) \to \mathrm{Conv}(\Sigma a,s) \to \varphi_{78}) \to \varphi_{12}$\lbl{A1}}
\BIC{$\varphi_{12}$}
\end{ptree}
\noindent$\Pi_{103}$:
\begin{ptree}
\AXC{$\stackrel{\Pi_{98}}{\vdots}$}\noLine
\UIC{$\varphi_{13}$}
\AXC{$\stackrel{\Pi_{102}}{\vdots}$}\noLine
\UIC{$\varphi_{12}$}
\AXC{$\varphi_{67}$\lbl{[S2]}}
\AXC{$\varphi_{67} \to \varphi_{151}$\lbl{A2}}
\BIC{$\varphi_{151}$}
\AXC{$\varphi_{151} \to \varphi_{203}$\lbl{A2}}
\BIC{$\varphi_{203}$}
\AXC{$\varphi_{203} \to \varphi_{12} \to \varphi_{13} \to \mathrm{Seq}(a) \to \varphi_{185}$\lbl{A1}}
\BIC{$\varphi_{12} \to \varphi_{13} \to \mathrm{Seq}(a) \to \varphi_{185}$}
\BIC{$\varphi_{13} \to \mathrm{Seq}(a) \to \varphi_{185}$}
\BIC{$\mathrm{Seq}(a) \to \varphi_{185}$}
\AXC{$\varphi_{220}$\lbl{A1}}
\BIC{$\mathrm{Seq}(a) \to \mathrm{Conv}(\Sigma a,s) \to n \in \mathbb{N} \to \varphi_{126}$}
\end{ptree}
\noindent$\Pi_{104}$:
\begin{ptree}
\AXC{$\varphi_{126} \to \varphi_{121}$\lbl{A2}}
\AXC{$\stackrel{\Pi_{103}}{\vdots}$}\noLine
\UIC{$\mathrm{Seq}(a) \to \mathrm{Conv}(\Sigma a,s) \to n \in \mathbb{N} \to \varphi_{126}$}
\AXC{$\varphi_{232}$\lbl{A1}}
\BIC{$(\varphi_{126} \to \varphi_{121}) \to \mathrm{Seq}(a) \to \mathrm{Conv}(\Sigma a,s) \to n \in \mathbb{N} \to \varphi_{121}$}
\BIC{$\mathrm{Seq}(a) \to \mathrm{Conv}(\Sigma a,s) \to n \in \mathbb{N} \to \varphi_{121}$}
\end{ptree}
\noindent$\Pi_{105}$:
\begin{ptree}
\AXC{$\stackrel{\Pi_{104}}{\vdots}$}\noLine
\UIC{$\mathrm{Seq}(a) \to \mathrm{Conv}(\Sigma a,s) \to n \in \mathbb{N} \to \varphi_{121}$}
\AXC{$(\mathrm{Seq}(a) \to \mathrm{Conv}(\Sigma a,s) \to n \in \mathbb{N} \to \varphi_{121}) \to \varphi_{178}$\lbl{A1}}
\BIC{$\varphi_{178}$}
\end{ptree}
\noindent$\Pi_{106}$:
\begin{ptree}
\AXC{$\mathrm{Seq}(a) \to \mathrm{Conv}(\Sigma a,s) \to n \in \mathbb{N} \to n \in \mathbb{N}$\lbl{A1}}
\AXC{$\stackrel{\Pi_{105}}{\vdots}$}\noLine
\UIC{$\varphi_{178}$}
\BIC{$\mathrm{Seq}(a) \to \mathrm{Conv}(\Sigma a,s) \to \varphi_{121}$}
\AXC{$(\mathrm{Seq}(a) \to \mathrm{Conv}(\Sigma a,s) \to \varphi_{121}) \to \varphi_{234}$\lbl{A1}}
\BIC{$\varphi_{234}$}
\end{ptree}
\noindent$\Pi_{107}$:
\begin{ptree}
\AXC{$\stackrel{\Pi_{81}}{\vdots}$}\noLine
\UIC{$\Sigma a \circ \mathsf{s}(n) = {\textstyle\sum_{k=1}^{S(n)} a_{k}} \to \varphi_{194}$}
\AXC{$x = x$\lbl{A5}}
\UIC{$\forall x\,(x = x)$}
\AXC{$\forall x\,(x = x) \to \varphi_{110}$\lbl{A2}}
\BIC{$\varphi_{110}$}
\AXC{$\stackrel{\Pi_{93}}{\vdots}$}\noLine
\UIC{$\varphi_{39}$}
\AXC{$\varphi_{39} \to \varphi_{110} \to \varphi_{222}$\lbl{A1}}
\BIC{$\varphi_{110} \to \varphi_{222}$}
\BIC{$\varphi_{222}$}
\BIC{$\mathrm{Seq}(a) \to \mathrm{Conv}(\Sigma a,s) \to \varphi_{116}$}
\AXC{$\stackrel{\Pi_{106}}{\vdots}$}\noLine
\UIC{$\varphi_{234}$}
\BIC{$\varphi_{215} \to \mathrm{Seq}(a) \to \mathrm{Conv}(\Sigma a,s) \to \varphi_{112}$}
\end{ptree}
\noindent$\Pi_{108}$:
\begin{ptree}
\AXC{$\stackrel{\Pi_{80}}{\vdots}$}\noLine
\UIC{$\varphi_{125}$}
\AXC{$\varphi_{125} \to \varphi_{202}$\lbl{A2}}
\BIC{$\varphi_{202}$}
\AXC{$\varphi_{202} \to \varphi_{215}$\lbl{A2}}
\BIC{$\varphi_{215}$}
\AXC{$\stackrel{\Pi_{107}}{\vdots}$}\noLine
\UIC{$\varphi_{215} \to \mathrm{Seq}(a) \to \mathrm{Conv}(\Sigma a,s) \to \varphi_{112}$}
\BIC{$\mathrm{Seq}(a) \to \mathrm{Conv}(\Sigma a,s) \to \varphi_{112}$}
\AXC{$\varphi_{238}$\lbl{A1}}
\BIC{$(\mathrm{Seq}(a) \to \mathrm{Conv}(\Sigma a,s) \to \varphi_{91}) \to \varphi_{225}$}
\end{ptree}
\noindent$\Pi_{109}$:
\begin{ptree}
\AXC{$\stackrel{\Pi_{79}}{\vdots}$}\noLine
\UIC{$\mathrm{Seq}(a) \to \mathrm{Conv}(\Sigma a,s) \to \varphi_{91}$}
\AXC{$\stackrel{\Pi_{108}}{\vdots}$}\noLine
\UIC{$(\mathrm{Seq}(a) \to \mathrm{Conv}(\Sigma a,s) \to \varphi_{91}) \to \varphi_{225}$}
\BIC{$\varphi_{225}$}
\end{ptree}
\noindent$\Pi_{110}$:
\begin{ptree}
\AXC{$\stackrel{\Pi_{68}}{\vdots}$}\noLine
\UIC{$\mathrm{Seq}(a) \to \varphi_{114}$}
\AXC{$\stackrel{\Pi_{69}}{\vdots}$}\noLine
\UIC{$\varphi_{125}$}
\AXC{$\varphi_{125} \to \varphi_{192}$\lbl{A2}}
\BIC{$\varphi_{192}$}
\AXC{$\varphi_{229}$\lbl{A2}}
\BIC{$\Sigma a \circ \mathsf{s} \oplus -1 \odot \Sigma a(n) = {\textstyle\sum_{k=1}^{S(n)} a_{k}} + -1 \odot \Sigma a(n) \to \varphi_{175}$}
\AXC{$\stackrel{\Pi_{109}}{\vdots}$}\noLine
\UIC{$\varphi_{225}$}
\BIC{$\mathrm{Seq}(a) \to \mathrm{Conv}(\Sigma a,s) \to \varphi_{103}$}
\AXC{$\varphi_{236}$\lbl{A1}}
\BIC{$(\mathrm{Seq}(a) \to \varphi_{114}) \to \varphi_{223}$}
\BIC{$\varphi_{223}$}
\end{ptree}
\noindent$\Pi_{111}$:
\begin{ptree}
\AXC{$\stackrel{\Pi_{64}}{\vdots}$}\noLine
\UIC{$\varphi_{125}$}
\AXC{$\varphi_{125} \to \varphi_{186}$\lbl{A2}}
\BIC{$\varphi_{186}$}
\AXC{$\varphi_{186} \to \Sigma a \circ \mathsf{s} \oplus -1 \odot \Sigma a(n) = {\textstyle\sum_{k=1}^{S(n)} a_{k}} + -1 \cdot {\textstyle\sum_{k=1}^{n} a_{k}} \to \varphi_{167}$\lbl{A2}}
\BIC{$\Sigma a \circ \mathsf{s} \oplus -1 \odot \Sigma a(n) = {\textstyle\sum_{k=1}^{S(n)} a_{k}} + -1 \cdot {\textstyle\sum_{k=1}^{n} a_{k}} \to \varphi_{167}$}
\AXC{$\stackrel{\Pi_{110}}{\vdots}$}\noLine
\UIC{$\varphi_{223}$}
\BIC{$\mathrm{Seq}(a) \to \varphi_{123}$}
\AXC{$\varphi_{237}$\lbl{A1}}
\BIC{$(\mathrm{Seq}(a) \to \mathrm{Conv}(\Sigma a,s) \to \varphi_{86}) \to \varphi_{226}$}
\end{ptree}
\noindent$\Pi_{112}$:
\begin{ptree}
\AXC{$\stackrel{\Pi_{63}}{\vdots}$}\noLine
\UIC{$\mathrm{Seq}(a) \to \mathrm{Conv}(\Sigma a,s) \to \varphi_{86}$}
\AXC{$\stackrel{\Pi_{111}}{\vdots}$}\noLine
\UIC{$(\mathrm{Seq}(a) \to \mathrm{Conv}(\Sigma a,s) \to \varphi_{86}) \to \varphi_{226}$}
\BIC{$\varphi_{226}$}
\end{ptree}
\noindent$\Pi_{113}$:
\begin{ptree}
\AXC{$\stackrel{\Pi_{57}}{\vdots}$}\noLine
\UIC{$\varphi_{125}$}
\AXC{$\varphi_{125} \to \varphi_{177}$\lbl{A2}}
\BIC{$\varphi_{177}$}
\AXC{$\varphi_{177} \to \Sigma a \circ \mathsf{s} \oplus -1 \odot \Sigma a(n) = {\textstyle\sum_{k=1}^{S(n)} a_{k}} + -{\textstyle\sum_{k=1}^{n} a_{k}} \to \varphi_{180}$\lbl{A2}}
\BIC{$\Sigma a \circ \mathsf{s} \oplus -1 \odot \Sigma a(n) = {\textstyle\sum_{k=1}^{S(n)} a_{k}} + -{\textstyle\sum_{k=1}^{n} a_{k}} \to \varphi_{180}$}
\AXC{$\stackrel{\Pi_{112}}{\vdots}$}\noLine
\UIC{$\varphi_{226}$}
\BIC{$\mathrm{Seq}(a) \to \mathrm{Conv}(\Sigma a,s) \to \varphi_{118}$}
\end{ptree}
\noindent$\Pi_{114}$:
\begin{ptree}
\AXC{$\stackrel{\Pi_{113}}{\vdots}$}\noLine
\UIC{$\mathrm{Seq}(a) \to \mathrm{Conv}(\Sigma a,s) \to \varphi_{118}$}
\AXC{$(\mathrm{Seq}(a) \to \mathrm{Conv}(\Sigma a,s) \to \varphi_{118}) \to \varphi_{235}$\lbl{A1}}
\BIC{$\varphi_{235}$}
\end{ptree}
\noindent$\Pi_{115}$:
\begin{ptree}
\AXC{$\stackrel{\Pi_{52}}{\vdots}$}\noLine
\UIC{$\varphi_{125}$}
\AXC{$\varphi_{125} \to \varphi_{197}$\lbl{A2}}
\BIC{$\varphi_{197}$}
\AXC{$\varphi_{197} \to \varphi_{214}$\lbl{A2}}
\BIC{$\varphi_{214}$}
\AXC{$\stackrel{\Pi_{53}}{\vdots}$}\noLine
\UIC{${\textstyle\sum_{k=1}^{n} a_{k}} + a_{S(n)} = a_{S(n)} + {\textstyle\sum_{k=1}^{n} a_{k}} \to \varphi_{84} \to \varphi_{85}$}
\AXC{$\stackrel{\Pi_{56}}{\vdots}$}\noLine
\UIC{$\varphi_{224}$}
\BIC{$\mathrm{Seq}(a) \to \mathrm{Conv}(\Sigma a,s) \to n \in \mathbb{N} \to \varphi_{85}$}
\AXC{$\stackrel{\Pi_{114}}{\vdots}$}\noLine
\UIC{$\varphi_{235}$}
\BIC{$\varphi_{214} \to \mathrm{Seq}(a) \to \mathrm{Conv}(\Sigma a,s) \to \varphi_{119}$}
\BIC{$\mathrm{Seq}(a) \to \mathrm{Conv}(\Sigma a,s) \to \varphi_{119}$}
\end{ptree}
\noindent$\Pi_{116}$:
\begin{ptree}
\AXC{$\stackrel{\Pi_{115}}{\vdots}$}\noLine
\UIC{$\mathrm{Seq}(a) \to \mathrm{Conv}(\Sigma a,s) \to \varphi_{119}$}
\AXC{$(\mathrm{Seq}(a) \to \mathrm{Conv}(\Sigma a,s) \to \varphi_{119}) \to \varphi_{230}$\lbl{A1}}
\BIC{$\varphi_{230}$}
\end{ptree}
\noindent$\Pi_{117}$:
\begin{ptree}
\AXC{$\stackrel{\Pi_{49}}{\vdots}$}\noLine
\UIC{$\varphi_{125}$}
\AXC{$\varphi_{125} \to \varphi_{196}$\lbl{A2}}
\BIC{$\varphi_{196}$}
\AXC{$\varphi_{196} \to \varphi_{213}$\lbl{A2}}
\BIC{$\varphi_{213}$}
\AXC{$\stackrel{\Pi_{51}}{\vdots}$}\noLine
\UIC{$a_{S(n)} + {\textstyle\sum_{k=1}^{n} a_{k}} + -{\textstyle\sum_{k=1}^{n} a_{k}} = a_{S(n)} + {\textstyle\sum_{k=1}^{n} a_{k}} - {\textstyle\sum_{k=1}^{n} a_{k}}$}
\AXC{$\stackrel{\Pi_{116}}{\vdots}$}\noLine
\UIC{$\varphi_{230}$}
\BIC{$\varphi_{213} \to \mathrm{Seq}(a) \to \mathrm{Conv}(\Sigma a,s) \to \varphi_{111}$}
\BIC{$\mathrm{Seq}(a) \to \mathrm{Conv}(\Sigma a,s) \to \varphi_{111}$}
\AXC{$\varphi_{233}$\lbl{A1}}
\BIC{$(\mathrm{Seq}(a) \to \varphi_{58}) \to \varphi_{216}$}
\end{ptree}
\noindent$\Pi_{118}$:
\begin{ptree}
\AXC{$\stackrel{\Pi_{38}}{\vdots}$}\noLine
\UIC{$\varphi_{125}$}
\AXC{$\varphi_{125} \to \varphi_{190}$\lbl{A2}}
\BIC{$\varphi_{190}$}
\AXC{$\varphi_{227}$\lbl{A2}}
\BIC{$\Sigma a \circ \mathsf{s} \oplus -1 \odot \Sigma a(n) = a_{S(n)} + {\textstyle\sum_{k=1}^{n} a_{k}} - {\textstyle\sum_{k=1}^{n} a_{k}} \to \varphi_{128}$}
\AXC{$\stackrel{\Pi_{48}}{\vdots}$}\noLine
\UIC{$\mathrm{Seq}(a) \to \varphi_{58}$}
\AXC{$\stackrel{\Pi_{117}}{\vdots}$}\noLine
\UIC{$(\mathrm{Seq}(a) \to \varphi_{58}) \to \varphi_{216}$}
\BIC{$\varphi_{216}$}
\BIC{$\mathrm{Seq}(a) \to \varphi_{70}$}
\end{ptree}
\noindent$\Pi_{119}$:
\begin{ptree}
\AXC{$\stackrel{\Pi_{23}}{\vdots}$}\noLine
\UIC{$\Sigma a \circ \mathsf{s} \oplus -1 \odot \Sigma a(n) = a_{S(n)} \to \varphi_{95}$}
\AXC{$\stackrel{\Pi_{37}}{\vdots}$}\noLine
\UIC{$\varphi_{26}$}
\AXC{$\stackrel{\Pi_{118}}{\vdots}$}\noLine
\UIC{$\mathrm{Seq}(a) \to \varphi_{70}$}
\AXC{$(\mathrm{Seq}(a) \to \varphi_{70}) \to \varphi_{26} \to \varphi_{210}$\lbl{A1}}
\BIC{$\varphi_{26} \to \varphi_{210}$}
\BIC{$\varphi_{210}$}
\BIC{$\mathrm{Seq}(a) \to \varphi_{79}$}
\end{ptree}
\noindent$\Pi_{120}$:
\begin{ptree}
\AXC{$\forall n\,(\varphi_{105}) \to \varphi_{117}$\lbl{A3}}
\AXC{$\stackrel{\Pi_{119}}{\vdots}$}\noLine
\UIC{$\mathrm{Seq}(a) \to \varphi_{79}$}
\AXC{$(\mathrm{Seq}(a) \to \varphi_{79}) \to \varphi_{105}$\lbl{A1}}
\BIC{$\varphi_{105}$}
\UIC{$\forall n\,(\varphi_{105})$}
\AXC{$\forall n\,(\varphi_{105}) \to (\forall n\,(\varphi_{105}) \to \varphi_{117}) \to \varphi_{4} \to \mathrm{Seq}(a) \to \varphi_{92}$\lbl{A1}}
\BIC{$(\forall n\,(\varphi_{105}) \to \varphi_{117}) \to \varphi_{4} \to \mathrm{Seq}(a) \to \varphi_{92}$}
\BIC{$\varphi_{4} \to \mathrm{Seq}(a) \to \varphi_{92}$}
\end{ptree}
\noindent$\Pi_{121}$:
\begin{ptree}
\AXC{$\forall a\,(\varphi_{49})$\lbl{[S12]}}
\AXC{$\forall a\,(\varphi_{49}) \to \varphi_{49}$\lbl{A2}}
\BIC{$\varphi_{49}$}
\AXC{$\varphi_{49} \to \mathrm{Seq}(a) \to \varphi_{52}$\lbl{A2}}
\BIC{$\mathrm{Seq}(a) \to \varphi_{52}$}
\AXC{$\varphi_{209}$\lbl{A1}}
\BIC{$(\mathrm{Seq}(a) \to \mathrm{Conv}(\Sigma a,s) \to \mathrm{Seq}(a)) \to \varphi_{138}$}
\end{ptree}
\noindent$\Pi_{122}$:
\begin{ptree}
\AXC{$\mathrm{Seq}(a) \to \mathrm{Conv}(\Sigma a,s) \to \mathrm{Seq}(a)$\lbl{A1}}
\AXC{$\stackrel{\Pi_{121}}{\vdots}$}\noLine
\UIC{$(\mathrm{Seq}(a) \to \mathrm{Conv}(\Sigma a,s) \to \mathrm{Seq}(a)) \to \varphi_{138}$}
\BIC{$\varphi_{138}$}
\end{ptree}
\noindent$\Pi_{123}$:
\begin{ptree}
\AXC{$\varphi_{30}$\lbl{[C13]}}
\AXC{$\varphi_{30} \to \varphi_{156}$\lbl{A2}}
\BIC{$\varphi_{156}$}
\AXC{$\varphi_{156} \to \varphi_{172}$\lbl{A2}}
\BIC{$\varphi_{172}$}
\AXC{$\varphi_{172} \to \mathrm{Seq}(a \circ \mathsf{s}) \to \varphi_{158}$\lbl{A2}}
\BIC{$\mathrm{Seq}(a \circ \mathsf{s}) \to \varphi_{158}$}
\AXC{$(\mathrm{Seq}(a \circ \mathsf{s}) \to \varphi_{158}) \to \varphi_{211}$\lbl{A1}}
\BIC{$\varphi_{211}$}
\end{ptree}
\noindent$\Pi_{124}$:
\begin{ptree}
\AXC{$\mathrm{NSeq}(\mathsf{s}) \land \forall k \in \mathbb{N}\,(\mathsf{s}(k) = S(k))$\lbl{[S14]}}
\AXC{$\varphi_{44}$\lbl{A1}}
\BIC{$\mathrm{Seq}(a) \to \mathrm{Conv}(\Sigma a,s) \to \mathrm{NSeq}(\mathsf{s})$}
\AXC{$\stackrel{\Pi_{122}}{\vdots}$}\noLine
\UIC{$\varphi_{138}$}
\BIC{$\mathrm{Seq}(a) \to \varphi_{63}$}
\AXC{$\varphi_{144}$\lbl{A1}}
\BIC{$\mathrm{Seq}(a) \to \mathrm{Conv}(\Sigma a,s) \to \mathrm{Seq}(a \circ \mathsf{s})$}
\AXC{$\stackrel{\Pi_{123}}{\vdots}$}\noLine
\UIC{$\varphi_{211}$}
\BIC{$(\mathrm{Seq}(a) \to \varphi_{92}) \to \varphi_{132}$}
\end{ptree}
\noindent$\Pi_{125}$:
\begin{ptree}
\AXC{$\stackrel{\Pi_{21}}{\vdots}$}\noLine
\UIC{$\varphi_{36}$}
\AXC{$\varphi_{4}$\lbl{A4}}
\AXC{$\stackrel{\Pi_{120}}{\vdots}$}\noLine
\UIC{$\varphi_{4} \to \mathrm{Seq}(a) \to \varphi_{92}$}
\BIC{$\mathrm{Seq}(a) \to \varphi_{92}$}
\AXC{$\stackrel{\Pi_{124}}{\vdots}$}\noLine
\UIC{$(\mathrm{Seq}(a) \to \varphi_{92}) \to \varphi_{132}$}
\BIC{$\varphi_{132}$}
\BIC{$\mathrm{Seq}(a) \to \mathrm{Conv}(\Sigma a,s) \to \mathrm{Conv}(a \circ \mathsf{s},0)$}
\end{ptree}
\noindent$\Pi_{126}$:
\begin{ptree}
\AXC{$\forall a\,(\varphi_{2})$\lbl{[C29]}}
\AXC{$\forall a\,(\varphi_{2}) \to \varphi_{2}$\lbl{A2}}
\BIC{$\varphi_{2}$}
\AXC{$\varphi_{2} \to \mathrm{Seq}(a) \to \mathrm{Conv}(a \circ \mathsf{s},0) \to \mathrm{Conv}(a,0)$\lbl{A2}}
\BIC{$\mathrm{Seq}(a) \to \mathrm{Conv}(a \circ \mathsf{s},0) \to \mathrm{Conv}(a,0)$}
\AXC{$\varphi_{146}$\lbl{A1}}
\BIC{$(\mathrm{Seq}(a) \to \mathrm{Seq}(a)) \to \varphi_{51}$}
\end{ptree}
\noindent 主木:
\begin{ptree}
\AXC{$\stackrel{\Pi_{125}}{\vdots}$}\noLine
\UIC{$\mathrm{Seq}(a) \to \mathrm{Conv}(\Sigma a,s) \to \mathrm{Conv}(a \circ \mathsf{s},0)$}
\AXC{$\mathrm{Seq}(a) \to \mathrm{Seq}(a)$\lbl{A1}}
\AXC{$\stackrel{\Pi_{126}}{\vdots}$}\noLine
\UIC{$(\mathrm{Seq}(a) \to \mathrm{Seq}(a)) \to \varphi_{51}$}
\BIC{$\varphi_{51}$}
\BIC{$\mathrm{Seq}(a) \to \mathrm{Conv}(\Sigma a,s) \to \mathrm{Conv}(a,0)$}
\UIC{$\varphi_{1}$}
\UIC{$\forall a\,(\varphi_{1})$}
\end{ptree}
\begin{abbrevlist}
\item $\varphi_{1} :\equiv \forall s\,(\mathrm{Seq}(a) \to \mathrm{Conv}(\Sigma a,s) \to \mathrm{Conv}(a,0))$
\item $\varphi_{2} :\equiv \forall \alpha\,(\mathrm{Seq}(a) \to \mathrm{Conv}(a \circ \mathsf{s},\alpha) \to \mathrm{Conv}(a,\alpha))$
\item $\varphi_{3} :\equiv \forall a\,\forall s\,\forall \alpha\,(\mathrm{Conv}(a,\alpha) \to \mathrm{SubIdx}(s) \to \mathrm{Conv}(a \circ s,\alpha))$
\item $\varphi_{4} :\equiv \mathrm{Seq}(a) \land \mathrm{Conv}(\Sigma a,s) \to \forall n\,(\mathrm{Seq}(a) \land \mathrm{Conv}(\Sigma a,s))$
\item $\varphi_{5} :\equiv \mathrm{Seq}(a) \to \mathrm{Conv}(\Sigma a,s) \to n \in \mathbb{N} \to \mathsf{s}(n) = S(n)$
\item $\varphi_{6} :\equiv \forall c\,\forall a\,\forall \alpha\,(c \in \mathbb{R} \to \mathrm{Conv}(a,\alpha) \to \mathrm{Conv}(c \odot a,c \cdot \alpha))$
\item $\varphi_{7} :\equiv \mathrm{Conv}(\Sigma a,s) \to \mathrm{SubIdx}(\mathsf{s}) \to \mathrm{Conv}(\Sigma a \circ \mathsf{s},s)$
\item $\varphi_{8} :\equiv \forall z\,\forall w\,(z = w \to s + -1 \cdot s = s + z \to s + -1 \cdot s = s + w)$
\item $\varphi_{9} :\equiv \mathrm{Conv}(\Sigma a,s) \to n \in \mathbb{N} \to \forall k \in \mathbb{N}\,(\mathsf{s}(k) = S(k))$
\item $\varphi_{10} :\equiv \forall s\,\forall \alpha\,(\mathrm{Conv}(\Sigma a,\alpha) \to \mathrm{SubIdx}(s) \to \mathrm{Conv}(\Sigma a \circ s,\alpha))$
\item $\varphi_{11} :\equiv \forall k \in \mathbb{N}\,(\mathsf{s}(k) = S(k)) \to n \in \mathbb{N} \to \mathsf{s}(n) = S(n)$
\item $\varphi_{12} :\equiv \mathrm{Seq}(a) \to \mathrm{Conv}(\Sigma a,s) \to n \in \mathbb{N} \to \mathrm{Seq}(\Sigma a \circ \mathsf{s})$
\item $\varphi_{13} :\equiv \mathrm{Seq}(a) \to \mathrm{Conv}(\Sigma a,s) \to n \in \mathbb{N} \to \mathrm{Seq}(-1 \odot \Sigma a)$
\item $\varphi_{14} :\equiv \mathrm{Seq}(a) \to \mathrm{Conv}(\Sigma a,s) \to \mathrm{Conv}(-1 \odot \Sigma a,-1 \cdot s)$
\item $\varphi_{15} :\equiv \forall \alpha\,(\mathrm{Conv}(\Sigma a,\alpha) \to \mathrm{SubIdx}(\mathsf{s}) \to \mathrm{Conv}(\Sigma a \circ \mathsf{s},\alpha))$
\item $\varphi_{16} :\equiv \forall a\,\forall \alpha\,\forall b\,\forall \beta\,(\mathrm{Conv}(a,\alpha) \to \mathrm{Conv}(b,\beta) \to \mathrm{Conv}(a \oplus b,\alpha + \beta))$
\item $\varphi_{17} :\equiv \forall \alpha\,(\mathrm{Conv}(\Sigma a,\alpha) \to \alpha \in \mathbb{R}) \to \mathrm{Conv}(\Sigma a,s) \to s \in \mathbb{R}$
\item $\varphi_{18} :\equiv \forall x\,\forall y\,(x = y \to y = x) \to \forall y\,(a \circ \mathsf{s}(n) = y \to y = a \circ \mathsf{s}(n))$
\item $\varphi_{19} :\equiv \mathrm{Seq}(a) \to \mathrm{Conv}(\Sigma a,s) \to n \in \mathbb{N} \to n \in \mathbb{N} \to \mathsf{s}(n) = S(n)$
\item $\varphi_{20} :\equiv 1 \in \mathbb{R} \to \mathrm{Seq}(a) \to \mathrm{Conv}(\Sigma a,s) \to n \in \mathbb{N} \to -1 \in \mathbb{R}$
\item $\varphi_{21} :\equiv \mathrm{SubIdx}(\mathsf{s}) \to \mathrm{Seq}(a) \to \mathrm{Conv}(\Sigma a,s) \to \mathrm{Conv}(\Sigma a \circ \mathsf{s},s)$
\item $\varphi_{22} :\equiv \forall z\,(s + -1 \cdot s = s + -s \to s + -s = z \to s + -1 \cdot s = z)$
\item $\varphi_{23} :\equiv \forall a\,\forall \alpha\,(\mathrm{Conv}(a,\alpha) \to \alpha \in \mathbb{R}) \to \forall \alpha\,(\mathrm{Conv}(\Sigma a,\alpha) \to \alpha \in \mathbb{R})$
\item $\varphi_{24} :\equiv \mathrm{Seq}(a) \to \mathrm{Conv}(\Sigma a,s) \to n \in \mathbb{N} \to a_{\mathsf{s}(n)} = a_{S(n)}$
\item $\varphi_{25} :\equiv \mathrm{Seq}(a) \to \mathrm{Conv}(\Sigma a,s) \to n \in \mathbb{N} \to a \circ \mathsf{s}(n) = a_{S(n)}$
\item $\varphi_{26} :\equiv \mathrm{Seq}(a) \to \mathrm{Conv}(\Sigma a,s) \to n \in \mathbb{N} \to a_{S(n)} = a \circ \mathsf{s}(n)$
\item $\varphi_{27} :\equiv \forall a\,\forall \alpha\,(-1 \in \mathbb{R} \to \mathrm{Conv}(a,\alpha) \to \mathrm{Conv}(-1 \odot a,-1 \cdot \alpha))$
\item $\varphi_{28} :\equiv -1 \in \mathbb{R} \to \mathrm{Conv}(\Sigma a,s) \to \mathrm{Conv}(-1 \odot \Sigma a,-1 \cdot s)$
\item $\varphi_{29} :\equiv \mathrm{Seq}(a) \to \mathrm{Conv}(\Sigma a,s) \to n \in \mathbb{N} \to a \circ \mathsf{s}(n) = a_{\mathsf{s}(n)}$
\item $\varphi_{30} :\equiv \forall a\,\forall b\,\forall \alpha\,(\mathrm{Seq}(b) \to \forall n \in \mathbb{N}\,(a_{n} = b_{n}) \to \mathrm{Conv}(a,\alpha) \to \mathrm{Conv}(b,\alpha))$
\item $\varphi_{31} :\equiv \forall x\,\forall y\,\forall z\,(x = y \to y = z \to x = z) \to \forall y\,\forall z\,(s + -1 \cdot s = y \to y = z \to s + -1 \cdot s = z)$
\item $\varphi_{32} :\equiv \forall w\,(-1 \cdot s = w \to s + -1 \cdot s = s + -1 \cdot s \to s + -1 \cdot s = s + w)$
\item $\varphi_{33} :\equiv \forall \alpha\,(-1 \in \mathbb{R} \to \mathrm{Conv}(\Sigma a,\alpha) \to \mathrm{Conv}(-1 \odot \Sigma a,-1 \cdot \alpha))$
\item $\varphi_{34} :\equiv (\mathrm{Seq}(a) \to \mathrm{Conv}(\Sigma a,s) \to s \in \mathbb{R}) \to \mathrm{Seq}(a) \to \mathrm{Conv}(\Sigma a,s) \to s + -s = 0$
\item $\varphi_{35} :\equiv \mathrm{Conv}(\Sigma a,s) \to n \in \mathbb{N} \to n \in \mathbb{N} \to a \circ \mathsf{s}(n) = a_{\mathsf{s}(n)}$
\item $\varphi_{36} :\equiv \mathrm{Seq}(a) \to \mathrm{Conv}(\Sigma a,s) \to \mathrm{Conv}(\Sigma a \circ \mathsf{s} \oplus -1 \odot \Sigma a,0)$
\item $\varphi_{37} :\equiv \mathrm{Seq}(a) \to \mathrm{Conv}(\Sigma a,s) \to n \in \mathbb{N} \to {\textstyle\sum_{k=1}^{\mathsf{s}(n)} a_{k}} = {\textstyle\sum_{k=1}^{S(n)} a_{k}}$
\item $\varphi_{38} :\equiv (\mathrm{Seq}(a) \to \mathrm{Conv}(\Sigma a,s) \to \mathrm{Conv}(\Sigma a,s)) \to \mathrm{Seq}(a) \to \mathrm{Conv}(\Sigma a,s) \to s \in \mathbb{R}$
\item $\varphi_{39} :\equiv \mathrm{Seq}(a) \to \mathrm{Conv}(\Sigma a,s) \to n \in \mathbb{N} \to \Sigma a \circ \mathsf{s}(n) = {\textstyle\sum_{k=1}^{S(n)} a_{k}}$
\item $\varphi_{40} :\equiv \mathrm{Conv}(\Sigma a,s) \to n \in \mathbb{N} \to -1 \odot \Sigma a(n) = -1 \cdot {\textstyle\sum_{k=1}^{n} a_{k}}$
\item $\varphi_{41} :\equiv \mathrm{Seq}(a) \to \mathrm{Conv}(\Sigma a,s) \to n \in \mathbb{N} \to -1 \cdot {\textstyle\sum_{k=1}^{n} a_{k}} = -{\textstyle\sum_{k=1}^{n} a_{k}}$
\item $\varphi_{42} :\equiv \mathrm{Conv}(\Sigma a,s) \to n \in \mathbb{N} \to \forall k \in \mathbb{N}\,(a \circ \mathsf{s}(k) = a_{\mathsf{s}(k)})$
\item $\varphi_{43} :\equiv \forall x\,\forall y\,\forall z\,(x = y \to y = z \to x = z) \to \forall y\,\forall z\,(a \circ \mathsf{s}(n) = y \to y = z \to a \circ \mathsf{s}(n) = z)$
\item $\varphi_{44} :\equiv \mathrm{NSeq}(\mathsf{s}) \land \forall k \in \mathbb{N}\,(\mathsf{s}(k) = S(k)) \to \mathrm{Seq}(a) \to \mathrm{Conv}(\Sigma a,s) \to \mathrm{NSeq}(\mathsf{s})$
\item $\varphi_{45} :\equiv -1 \cdot s = -s \to s + -1 \cdot s = s + -1 \cdot s \to s + -1 \cdot s = s + -s$
\item $\varphi_{46} :\equiv (\mathrm{Seq}(a) \to \mathrm{Conv}(\Sigma a,s) \to s \in \mathbb{R}) \to \mathrm{Seq}(a) \to \mathrm{Conv}(\Sigma a,s) \to -1 \cdot s = -s$
\item $\varphi_{47} :\equiv n \in \mathbb{N} \to \mathrm{Seq}(a \circ \mathsf{s}) \land \forall k \in \mathbb{N}\,(a \circ \mathsf{s}(k) = a_{\mathsf{s}(k)})$
\item $\varphi_{48} :\equiv (1 \in \mathbb{R} \to -1 \in \mathbb{R}) \to 1 \in \mathbb{R} \to \mathrm{Seq}(a) \to \mathrm{Conv}(\Sigma a,s) \to -1 \in \mathbb{R}$
\item $\varphi_{49} :\equiv \forall s\,(\mathrm{Seq}(a) \to \mathrm{NSeq}(s) \to \mathrm{Seq}(a \circ s) \land \forall k \in \mathbb{N}\,(a \circ s(k) = a_{s_{k}}))$
\item $\varphi_{50} :\equiv \mathrm{Seq}(a) \to \mathrm{Conv}(\Sigma a,s) \to n \in \mathbb{N} \to {\textstyle\sum_{k=1}^{S(n)} a_{k}} = {\textstyle\sum_{k=1}^{n} a_{k}} + a_{S(n)}$
\item $\varphi_{51} :\equiv (\mathrm{Seq}(a) \to \mathrm{Conv}(\Sigma a,s) \to \mathrm{Conv}(a \circ \mathsf{s},0)) \to \mathrm{Seq}(a) \to \mathrm{Conv}(\Sigma a,s) \to \mathrm{Conv}(a,0)$
\item $\varphi_{52} :\equiv \mathrm{NSeq}(\mathsf{s}) \to \mathrm{Seq}(a \circ \mathsf{s}) \land \forall k \in \mathbb{N}\,(a \circ \mathsf{s}(k) = a_{\mathsf{s}(k)})$
\item $\varphi_{53} :\equiv \mathrm{Conv}(\Sigma a,s) \to \mathrm{Conv}(\Sigma a \circ \mathsf{s} \oplus -1 \odot \Sigma a,s + -1 \cdot s)$
\item $\varphi_{54} :\equiv \forall z\,(a \circ \mathsf{s}(n) = a_{\mathsf{s}(n)} \to a_{\mathsf{s}(n)} = z \to a \circ \mathsf{s}(n) = z)$
\item $\varphi_{55} :\equiv \forall a\,\forall n\,(\mathrm{Seq}(a) \to n \in \mathbb{N} \to a_{n} \in \mathbb{R}) \to \forall n\,(\mathrm{Seq}(a) \to n \in \mathbb{N} \to a_{n} \in \mathbb{R})$
\item $\varphi_{56} :\equiv \forall n\,(\mathrm{Seq}(a) \to n \in \mathbb{N} \to a_{n} \in \mathbb{R}) \to \mathrm{Seq}(a) \to S(n) \in \mathbb{N} \to a_{S(n)} \in \mathbb{R}$
\item $\varphi_{57} :\equiv \mathrm{Seq}(a) \to \mathrm{Conv}(\Sigma a,s) \to n \in \mathbb{N} \to \Sigma a \circ \mathsf{s}(n) = {\textstyle\sum_{k=1}^{\mathsf{s}(n)} a_{k}}$
\item $\varphi_{58} :\equiv \mathrm{Conv}(\Sigma a,s) \to n \in \mathbb{N} \to a_{S(n)} + {\textstyle\sum_{k=1}^{n} a_{k}} - {\textstyle\sum_{k=1}^{n} a_{k}} = a_{S(n)}$
\item $\varphi_{59} :\equiv \mathrm{Conv}(\Sigma a,s) \to n \in \mathbb{N} \to {\textstyle\sum_{k=1}^{n} a_{k}} + a_{S(n)} = a_{S(n)} + {\textstyle\sum_{k=1}^{n} a_{k}}$
\item $\varphi_{60} :\equiv \mathrm{Conv}(\Sigma a,s) \to n \in \mathbb{N} \to n \in \mathbb{N} \to {\textstyle\sum_{k=1}^{S(n)} a_{k}} = {\textstyle\sum_{k=1}^{n} a_{k}} + a_{S(n)}$
\item $\varphi_{61} :\equiv \forall b\,\forall \beta\,(\mathrm{Conv}(\Sigma a \circ \mathsf{s},s) \to \mathrm{Conv}(b,\beta) \to \mathrm{Conv}(\Sigma a \circ \mathsf{s} \oplus b,s + \beta))$
\item $\varphi_{62} :\equiv \forall c\,\forall a\,(c \in \mathbb{R} \to \mathrm{Seq}(a) \to \mathrm{Seq}(c \odot a) \land \forall n \in \mathbb{N}\,(c \odot a(n) = c \cdot a_{n}))$
\item $\varphi_{63} :\equiv \mathrm{Conv}(\Sigma a,s) \to \mathrm{Seq}(a \circ \mathsf{s}) \land \forall k \in \mathbb{N}\,(a \circ \mathsf{s}(k) = a_{\mathsf{s}(k)})$
\item $\varphi_{64} :\equiv \mathrm{Conv}(\Sigma a,s) \to n \in \mathbb{N} \to \forall n \in \mathbb{N}\,({\textstyle\sum_{k=1}^{S(n)} a_{k}} = {\textstyle\sum_{k=1}^{n} a_{k}} + a_{S(n)})$
\item $\varphi_{65} :\equiv \mathrm{Conv}(\Sigma a,s) \to n \in \mathbb{N} \to n \in \mathbb{N} \to \Sigma a \circ \mathsf{s}(n) = {\textstyle\sum_{k=1}^{\mathsf{s}(n)} a_{k}}$
\item $\varphi_{66} :\equiv \mathrm{Conv}(\Sigma a,s) \to n \in \mathbb{N} \to n \in \mathbb{N} \to -1 \odot \Sigma a(n) = -1 \cdot {\textstyle\sum_{k=1}^{n} a_{k}}$
\item $\varphi_{67} :\equiv \forall a\,\forall b\,(\mathrm{Seq}(a) \to \mathrm{Seq}(b) \to \mathrm{Seq}(a \oplus b) \land \forall n \in \mathbb{N}\,(a \oplus b(n) = a_{n} + b_{n}))$
\item $\varphi_{68} :\equiv \mathrm{Seq}(-1 \odot \Sigma a) \land \forall n \in \mathbb{N}\,(-1 \odot \Sigma a(n) = -1 \cdot {\textstyle\sum_{k=1}^{n} a_{k}})$
\item $\varphi_{69} :\equiv \mathrm{Conv}(\Sigma a,s) \to n \in \mathbb{N} \to \forall k \in \mathbb{N}\,(\Sigma a \circ \mathsf{s}(k) = {\textstyle\sum_{k=1}^{\mathsf{s}(k)} a_{k}})$
\item $\varphi_{70} :\equiv \mathrm{Conv}(\Sigma a,s) \to n \in \mathbb{N} \to \Sigma a \circ \mathsf{s} \oplus -1 \odot \Sigma a(n) = a_{S(n)}$
\item $\varphi_{71} :\equiv \forall \alpha\,\forall b\,\forall \beta\,(\mathrm{Conv}(\Sigma a \circ \mathsf{s},\alpha) \to \mathrm{Conv}(b,\beta) \to \mathrm{Conv}(\Sigma a \circ \mathsf{s} \oplus b,\alpha + \beta))$
\item $\varphi_{72} :\equiv \mathrm{NSeq}(\mathsf{s}) \land \forall k \in \mathbb{N}\,(\mathsf{s}(k) = S(k)) \to \mathrm{Seq}(a) \to \mathrm{Conv}(\Sigma a,s) \to n \in \mathbb{N} \to \mathrm{NSeq}(\mathsf{s})$
\item $\varphi_{73} :\equiv \forall x\,\forall y\,\forall z\,(x = y \to y = z \to x = z) \to \forall y\,\forall z\,(\Sigma a \circ \mathsf{s}(n) = y \to y = z \to \Sigma a \circ \mathsf{s}(n) = z)$
\item $\varphi_{74} :\equiv \mathrm{Conv}(\Sigma a,s) \to n \in \mathbb{N} \to \forall n \in \mathbb{N}\,(-1 \odot \Sigma a(n) = -1 \cdot {\textstyle\sum_{k=1}^{n} a_{k}})$
\item $\varphi_{75} :\equiv \forall a\,\forall n\,(\mathrm{Seq}(a) \to n \in \mathbb{N} \to a_{n} \in \mathbb{R}) \to \forall n\,(\mathrm{Seq}(\Sigma a) \to n \in \mathbb{N} \to {\textstyle\sum_{k=1}^{n} a_{k}} \in \mathbb{R})$
\item $\varphi_{76} :\equiv \forall x\,(x = x) \to {\textstyle\sum_{k=1}^{S(n)} a_{k}} + -1 \cdot {\textstyle\sum_{k=1}^{n} a_{k}} = {\textstyle\sum_{k=1}^{S(n)} a_{k}} + -1 \cdot {\textstyle\sum_{k=1}^{n} a_{k}}$
\item $\varphi_{77} :\equiv (\mathrm{Seq}(a) \to \mathrm{Conv}(\Sigma a,s) \to n \in \mathbb{N} \to n \in \mathbb{N}) \to \mathrm{Seq}(a) \to \mathrm{Conv}(\Sigma a,s) \to n \in \mathbb{N} \to S(n) \in \mathbb{N}$
\item $\varphi_{78} :\equiv n \in \mathbb{N} \to \mathrm{Seq}(\Sigma a \circ \mathsf{s}) \land \forall k \in \mathbb{N}\,(\Sigma a \circ \mathsf{s}(k) = {\textstyle\sum_{k=1}^{\mathsf{s}(k)} a_{k}})$
\item $\varphi_{79} :\equiv \mathrm{Conv}(\Sigma a,s) \to n \in \mathbb{N} \to \Sigma a \circ \mathsf{s} \oplus -1 \odot \Sigma a(n) = a \circ \mathsf{s}(n)$
\item $\varphi_{80} :\equiv \forall y\,(a_{S(n)} + {\textstyle\sum_{k=1}^{n} a_{k}} - {\textstyle\sum_{k=1}^{n} a_{k}} = y \to y = a_{S(n)} + {\textstyle\sum_{k=1}^{n} a_{k}} - {\textstyle\sum_{k=1}^{n} a_{k}})$
\item $\varphi_{81} :\equiv \forall x\,\forall y\,(x - y = x + -y) \to \forall y\,(a_{S(n)} + {\textstyle\sum_{k=1}^{n} a_{k}} - y = a_{S(n)} + {\textstyle\sum_{k=1}^{n} a_{k}} + -y)$
\item $\varphi_{82} :\equiv \mathrm{NSeq}(\mathsf{s}) \to \mathrm{Seq}(\Sigma a \circ \mathsf{s}) \land \forall k \in \mathbb{N}\,(\Sigma a \circ \mathsf{s}(k) = {\textstyle\sum_{k=1}^{\mathsf{s}(k)} a_{k}})$
\item $\varphi_{83} :\equiv a \circ \mathsf{s}(n) = a_{\mathsf{s}(n)} \to a_{\mathsf{s}(n)} = a_{S(n)} \to a \circ \mathsf{s}(n) = a_{S(n)}$
\item $\varphi_{84} :\equiv {\textstyle\sum_{k=1}^{n} a_{k}} + a_{S(n)} + -{\textstyle\sum_{k=1}^{n} a_{k}} = {\textstyle\sum_{k=1}^{n} a_{k}} + a_{S(n)} + -{\textstyle\sum_{k=1}^{n} a_{k}}$
\item $\varphi_{85} :\equiv {\textstyle\sum_{k=1}^{n} a_{k}} + a_{S(n)} + -{\textstyle\sum_{k=1}^{n} a_{k}} = a_{S(n)} + {\textstyle\sum_{k=1}^{n} a_{k}} + -{\textstyle\sum_{k=1}^{n} a_{k}}$
\item $\varphi_{86} :\equiv n \in \mathbb{N} \to {\textstyle\sum_{k=1}^{S(n)} a_{k}} + -{\textstyle\sum_{k=1}^{n} a_{k}} = {\textstyle\sum_{k=1}^{n} a_{k}} + a_{S(n)} + -{\textstyle\sum_{k=1}^{n} a_{k}}$
\item $\varphi_{87} :\equiv (\mathrm{Seq}(a) \to \mathrm{Conv}(\Sigma a,s) \to n \in \mathbb{N} \to n \in \mathbb{N}) \to \varphi_{5}$
\item $\varphi_{88} :\equiv \forall k \in \mathbb{N}\,(a \circ \mathsf{s}(k) = a_{\mathsf{s}(k)}) \to n \in \mathbb{N} \to a \circ \mathsf{s}(n) = a_{\mathsf{s}(n)}$
\item $\varphi_{89} :\equiv \forall n\,(\mathrm{Seq}(\Sigma a) \to n \in \mathbb{N} \to {\textstyle\sum_{k=1}^{n} a_{k}} \in \mathbb{R}) \to \mathrm{Seq}(\Sigma a) \to n \in \mathbb{N} \to {\textstyle\sum_{k=1}^{n} a_{k}} \in \mathbb{R}$
\item $\varphi_{90} :\equiv \forall x\,(x = x) \to {\textstyle\sum_{k=1}^{S(n)} a_{k}} + -1 \odot \Sigma a(n) = {\textstyle\sum_{k=1}^{S(n)} a_{k}} + -1 \odot \Sigma a(n)$
\item $\varphi_{91} :\equiv n \in \mathbb{N} \to {\textstyle\sum_{k=1}^{S(n)} a_{k}} + -1 \odot \Sigma a(n) = {\textstyle\sum_{k=1}^{S(n)} a_{k}} + -1 \cdot {\textstyle\sum_{k=1}^{n} a_{k}}$
\item $\varphi_{92} :\equiv \mathrm{Conv}(\Sigma a,s) \to \forall n \in \mathbb{N}\,(\Sigma a \circ \mathsf{s} \oplus -1 \odot \Sigma a(n) = a \circ \mathsf{s}(n))$
\item $\varphi_{93} :\equiv \forall x \in \mathbb{R}\,(-1 \cdot x = -x) \to {\textstyle\sum_{k=1}^{n} a_{k}} \in \mathbb{R} \to -1 \cdot {\textstyle\sum_{k=1}^{n} a_{k}} = -{\textstyle\sum_{k=1}^{n} a_{k}}$
\item $\varphi_{94} :\equiv \forall a\,(-1 \in \mathbb{R} \to \mathrm{Seq}(a) \to \mathrm{Seq}(-1 \odot a) \land \forall n \in \mathbb{N}\,(-1 \odot a(n) = -1 \cdot a_{n}))$
\item $\varphi_{95} :\equiv a_{S(n)} = a \circ \mathsf{s}(n) \to \Sigma a \circ \mathsf{s} \oplus -1 \odot \Sigma a(n) = a \circ \mathsf{s}(n)$
\item $\varphi_{96} :\equiv a_{S(n)} \in \mathbb{R} \to {\textstyle\sum_{k=1}^{n} a_{k}} \in \mathbb{R} \to a_{S(n)} + {\textstyle\sum_{k=1}^{n} a_{k}} - {\textstyle\sum_{k=1}^{n} a_{k}} = a_{S(n)}$
\item $\varphi_{97} :\equiv {\textstyle\sum_{k=1}^{n} a_{k}} \in \mathbb{R} \to a_{S(n)} \in \mathbb{R} \to {\textstyle\sum_{k=1}^{n} a_{k}} + a_{S(n)} = a_{S(n)} + {\textstyle\sum_{k=1}^{n} a_{k}}$
\item $\varphi_{98} :\equiv \mathrm{Seq}(\Sigma a) \land {\textstyle\sum_{k=1}^{1_{\omega}} a_{k}} = a_{1_{\omega}} \land \forall n \in \mathbb{N}\,({\textstyle\sum_{k=1}^{S(n)} a_{k}} = {\textstyle\sum_{k=1}^{n} a_{k}} + a_{S(n)})$
\item $\varphi_{99} :\equiv (\mathrm{Seq}(a) \to \mathrm{Conv}(\Sigma a,s) \to \mathrm{Conv}(\Sigma a,s)) \to \varphi_{14}$
\item $\varphi_{100} :\equiv \forall s\,(\mathrm{Seq}(\Sigma a) \to \mathrm{NSeq}(s) \to \mathrm{Seq}(\Sigma a \circ s) \land \forall k \in \mathbb{N}\,(\Sigma a \circ s(k) = {\textstyle\sum_{k=1}^{s_{k}} a_{k}}))$
\item $\varphi_{101} :\equiv \mathrm{Conv}(-1 \odot \Sigma a,-1 \cdot s) \to \mathrm{Conv}(\Sigma a \circ \mathsf{s} \oplus -1 \odot \Sigma a,s + -1 \cdot s)$
\item $\varphi_{102} :\equiv \mathrm{Conv}(\Sigma a \circ \mathsf{s} \oplus -1 \odot \Sigma a,z) \to \mathrm{Conv}(\Sigma a \circ \mathsf{s} \oplus -1 \odot \Sigma a,w)$
\item $\varphi_{103} :\equiv n \in \mathbb{N} \to \Sigma a \circ \mathsf{s} \oplus -1 \odot \Sigma a(n) = {\textstyle\sum_{k=1}^{S(n)} a_{k}} + -1 \cdot {\textstyle\sum_{k=1}^{n} a_{k}}$
\item $\varphi_{104} :\equiv \forall z\,(\Sigma a \circ \mathsf{s}(n) = {\textstyle\sum_{k=1}^{\mathsf{s}(n)} a_{k}} \to {\textstyle\sum_{k=1}^{\mathsf{s}(n)} a_{k}} = z \to \Sigma a \circ \mathsf{s}(n) = z)$
\item $\varphi_{105} :\equiv \mathrm{Seq}(a) \land \mathrm{Conv}(\Sigma a,s) \to n \in \mathbb{N} \to \Sigma a \circ \mathsf{s} \oplus -1 \odot \Sigma a(n) = a \circ \mathsf{s}(n)$
\item $\varphi_{106} :\equiv (\mathrm{Seq}(a) \to \mathrm{Conv}(\Sigma a,s) \to \mathrm{Conv}(\Sigma a,s)) \to \varphi_{21}$
\item $\varphi_{107} :\equiv (s + -1 \cdot s = s + -s \to s + -s = 0 \to s + -1 \cdot s = 0) \to \mathrm{Seq}(a) \to \mathrm{Conv}(\Sigma a,s) \to s + -1 \cdot s = 0$
\item $\varphi_{108} :\equiv (\mathrm{Seq}(a) \to \mathrm{Conv}(\Sigma a,s) \to n \in \mathbb{N} \to n \in \mathbb{N}) \to \mathrm{Seq}(a) \to \mathrm{Conv}(\Sigma a,s) \to n \in \mathbb{N} \to {\textstyle\sum_{k=1}^{n} a_{k}} \in \mathbb{R}$
\item $\varphi_{109} :\equiv (\mathrm{Seq}(a) \to \mathrm{Conv}(\Sigma a,s) \to n \in \mathbb{N} \to S(n) \in \mathbb{N}) \to \mathrm{Seq}(a) \to \mathrm{Conv}(\Sigma a,s) \to n \in \mathbb{N} \to a_{S(n)} \in \mathbb{R}$
\item $\varphi_{110} :\equiv \Sigma a \circ \mathsf{s}(n) + -1 \odot \Sigma a(n) = \Sigma a \circ \mathsf{s}(n) + -1 \odot \Sigma a(n)$
\item $\varphi_{111} :\equiv n \in \mathbb{N} \to \Sigma a \circ \mathsf{s} \oplus -1 \odot \Sigma a(n) = a_{S(n)} + {\textstyle\sum_{k=1}^{n} a_{k}} - {\textstyle\sum_{k=1}^{n} a_{k}}$
\item $\varphi_{112} :\equiv n \in \mathbb{N} \to \Sigma a \circ \mathsf{s} \oplus -1 \odot \Sigma a(n) = {\textstyle\sum_{k=1}^{S(n)} a_{k}} + -1 \odot \Sigma a(n)$
\item $\varphi_{113} :\equiv \forall y\,(a \circ \mathsf{s}(n) = y \to y = a \circ \mathsf{s}(n)) \to a \circ \mathsf{s}(n) = a_{S(n)} \to a_{S(n)} = a \circ \mathsf{s}(n)$
\item $\varphi_{114} :\equiv \mathrm{Conv}(\Sigma a,s) \to n \in \mathbb{N} \to {\textstyle\sum_{k=1}^{S(n)} a_{k}} + -1 \cdot {\textstyle\sum_{k=1}^{n} a_{k}} = {\textstyle\sum_{k=1}^{S(n)} a_{k}} + -{\textstyle\sum_{k=1}^{n} a_{k}}$
\item $\varphi_{115} :\equiv \mathrm{NSeq}(\mathsf{s}) \land \forall k \in \mathbb{N}\,(\mathsf{s}(k) = S(k)) \to \mathrm{Seq}(a) \to \varphi_{9}$
\item $\varphi_{116} :\equiv n \in \mathbb{N} \to \Sigma a \circ \mathsf{s}(n) + -1 \odot \Sigma a(n) = {\textstyle\sum_{k=1}^{S(n)} a_{k}} + -1 \odot \Sigma a(n)$
\item $\varphi_{117} :\equiv \forall n\,(\mathrm{Seq}(a) \land \mathrm{Conv}(\Sigma a,s)) \to \forall n \in \mathbb{N}\,(\Sigma a \circ \mathsf{s} \oplus -1 \odot \Sigma a(n) = a \circ \mathsf{s}(n))$
\item $\varphi_{118} :\equiv n \in \mathbb{N} \to \Sigma a \circ \mathsf{s} \oplus -1 \odot \Sigma a(n) = {\textstyle\sum_{k=1}^{n} a_{k}} + a_{S(n)} + -{\textstyle\sum_{k=1}^{n} a_{k}}$
\item $\varphi_{119} :\equiv n \in \mathbb{N} \to \Sigma a \circ \mathsf{s} \oplus -1 \odot \Sigma a(n) = a_{S(n)} + {\textstyle\sum_{k=1}^{n} a_{k}} + -{\textstyle\sum_{k=1}^{n} a_{k}}$
\item $\varphi_{120} :\equiv \forall x\,\forall y\,(x \in \mathbb{R} \to y \in \mathbb{R} \to x + y - y = x) \to \forall y\,(a_{S(n)} \in \mathbb{R} \to y \in \mathbb{R} \to a_{S(n)} + y - y = a_{S(n)})$
\item $\varphi_{121} :\equiv n \in \mathbb{N} \to \Sigma a \circ \mathsf{s} \oplus -1 \odot \Sigma a(n) = \Sigma a \circ \mathsf{s}(n) + -1 \odot \Sigma a(n)$
\item $\varphi_{122} :\equiv \forall n \in \mathbb{N}\,({\textstyle\sum_{k=1}^{S(n)} a_{k}} = {\textstyle\sum_{k=1}^{n} a_{k}} + a_{S(n)}) \to n \in \mathbb{N} \to {\textstyle\sum_{k=1}^{S(n)} a_{k}} = {\textstyle\sum_{k=1}^{n} a_{k}} + a_{S(n)}$
\item $\varphi_{123} :\equiv \mathrm{Conv}(\Sigma a,s) \to n \in \mathbb{N} \to \Sigma a \circ \mathsf{s} \oplus -1 \odot \Sigma a(n) = {\textstyle\sum_{k=1}^{S(n)} a_{k}} + -{\textstyle\sum_{k=1}^{n} a_{k}}$
\item $\varphi_{124} :\equiv \forall \beta\,(\mathrm{Conv}(\Sigma a \circ \mathsf{s},s) \to \mathrm{Conv}(-1 \odot \Sigma a,\beta) \to \mathrm{Conv}(\Sigma a \circ \mathsf{s} \oplus -1 \odot \Sigma a,s + \beta))$
\item $\varphi_{125} :\equiv \forall y\,\forall z\,(\Sigma a \circ \mathsf{s} \oplus -1 \odot \Sigma a(n) = y \to y = z \to \Sigma a \circ \mathsf{s} \oplus -1 \odot \Sigma a(n) = z)$
\item $\varphi_{126} :\equiv \forall n \in \mathbb{N}\,(\Sigma a \circ \mathsf{s} \oplus -1 \odot \Sigma a(n) = \Sigma a \circ \mathsf{s}(n) + -1 \odot \Sigma a(n))$
\item $\varphi_{127} :\equiv \mathrm{Conv}(\Sigma a \circ \mathsf{s} \oplus -1 \odot \Sigma a,s + -1 \cdot s) \to \mathrm{Conv}(\Sigma a \circ \mathsf{s} \oplus -1 \odot \Sigma a,0)$
\item $\varphi_{128} :\equiv a_{S(n)} + {\textstyle\sum_{k=1}^{n} a_{k}} - {\textstyle\sum_{k=1}^{n} a_{k}} = a_{S(n)} \to \Sigma a \circ \mathsf{s} \oplus -1 \odot \Sigma a(n) = a_{S(n)}$
\item $\varphi_{129} :\equiv (\mathrm{Seq}(a) \to \mathrm{Conv}(\Sigma a,s) \to n \in \mathbb{N} \to n \in \mathbb{N}) \to \varphi_{29}$
\item $\varphi_{130} :\equiv \forall x\,\forall y\,(x \in \mathbb{R} \to y \in \mathbb{R} \to x + y = y + x) \to \forall y\,({\textstyle\sum_{k=1}^{n} a_{k}} \in \mathbb{R} \to y \in \mathbb{R} \to {\textstyle\sum_{k=1}^{n} a_{k}} + y = y + {\textstyle\sum_{k=1}^{n} a_{k}})$
\item $\varphi_{131} :\equiv \forall k \in \mathbb{N}\,(\Sigma a \circ \mathsf{s}(k) = {\textstyle\sum_{k=1}^{\mathsf{s}(k)} a_{k}}) \to n \in \mathbb{N} \to \Sigma a \circ \mathsf{s}(n) = {\textstyle\sum_{k=1}^{\mathsf{s}(n)} a_{k}}$
\item $\varphi_{132} :\equiv \varphi_{36} \to \mathrm{Seq}(a) \to \mathrm{Conv}(\Sigma a,s) \to \mathrm{Conv}(a \circ \mathsf{s},0)$
\item $\varphi_{133} :\equiv {\textstyle\sum_{k=1}^{S(n)} a_{k}} + -{\textstyle\sum_{k=1}^{n} a_{k}} = z + -{\textstyle\sum_{k=1}^{n} a_{k}} \to {\textstyle\sum_{k=1}^{S(n)} a_{k}} + -{\textstyle\sum_{k=1}^{n} a_{k}} = w + -{\textstyle\sum_{k=1}^{n} a_{k}}$
\item $\varphi_{134} :\equiv \forall n \in \mathbb{N}\,(-1 \odot \Sigma a(n) = -1 \cdot {\textstyle\sum_{k=1}^{n} a_{k}}) \to n \in \mathbb{N} \to -1 \odot \Sigma a(n) = -1 \cdot {\textstyle\sum_{k=1}^{n} a_{k}}$
\item $\varphi_{135} :\equiv (\mathrm{Seq}(a) \to \mathrm{Conv}(\Sigma a,s) \to n \in \mathbb{N} \to n \in \mathbb{N}) \to \varphi_{50}$
\item $\varphi_{136} :\equiv {\textstyle\sum_{k=1}^{S(n)} a_{k}} + -1 \cdot {\textstyle\sum_{k=1}^{n} a_{k}} = {\textstyle\sum_{k=1}^{S(n)} a_{k}} + z \to {\textstyle\sum_{k=1}^{S(n)} a_{k}} + -1 \cdot {\textstyle\sum_{k=1}^{n} a_{k}} = {\textstyle\sum_{k=1}^{S(n)} a_{k}} + w$
\item $\varphi_{137} :\equiv (\mathrm{Seq}(a) \to \mathrm{Conv}(\Sigma a,s) \to n \in \mathbb{N} \to n \in \mathbb{N}) \to \varphi_{57}$
\item $\varphi_{138} :\equiv (\mathrm{Seq}(a) \to \mathrm{Conv}(\Sigma a,s) \to \mathrm{NSeq}(\mathsf{s})) \to \mathrm{Seq}(a) \to \varphi_{63}$
\item $\varphi_{139} :\equiv \Sigma a \circ \mathsf{s}(n) = {\textstyle\sum_{k=1}^{\mathsf{s}(n)} a_{k}} \to {\textstyle\sum_{k=1}^{\mathsf{s}(n)} a_{k}} = {\textstyle\sum_{k=1}^{S(n)} a_{k}} \to \Sigma a \circ \mathsf{s}(n) = {\textstyle\sum_{k=1}^{S(n)} a_{k}}$
\item $\varphi_{140} :\equiv \forall z\,(\Sigma a \circ \mathsf{s} \oplus -1 \odot \Sigma a(n) = a_{S(n)} \to a_{S(n)} = z \to \Sigma a \circ \mathsf{s} \oplus -1 \odot \Sigma a(n) = z)$
\item $\varphi_{141} :\equiv (\mathrm{Seq}(a) \to \mathrm{Conv}(\Sigma a,s) \to n \in \mathbb{N} \to n \in \mathbb{N}) \to \mathrm{Seq}(a) \to \varphi_{40}$
\item $\varphi_{142} :\equiv \varphi_{45} \to \mathrm{Seq}(a) \to \mathrm{Conv}(\Sigma a,s) \to s + -1 \cdot s = s + -s$
\item $\varphi_{143} :\equiv {\textstyle\sum_{k=1}^{S(n)} a_{k}} + -1 \odot \Sigma a(n) = {\textstyle\sum_{k=1}^{S(n)} a_{k}} + z \to {\textstyle\sum_{k=1}^{S(n)} a_{k}} + -1 \odot \Sigma a(n) = {\textstyle\sum_{k=1}^{S(n)} a_{k}} + w$
\item $\varphi_{144} :\equiv (\mathrm{Seq}(a) \to \varphi_{63}) \to \mathrm{Seq}(a) \to \mathrm{Conv}(\Sigma a,s) \to \mathrm{Seq}(a \circ \mathsf{s})$
\item $\varphi_{145} :\equiv (\mathrm{Seq}(a) \to \mathrm{Conv}(\Sigma a,s) \to n \in \mathbb{N} \to {\textstyle\sum_{k=1}^{n} a_{k}} \in \mathbb{R}) \to \varphi_{41}$
\item $\varphi_{146} :\equiv (\mathrm{Seq}(a) \to \mathrm{Conv}(a \circ \mathsf{s},0) \to \mathrm{Conv}(a,0)) \to (\mathrm{Seq}(a) \to \mathrm{Seq}(a)) \to \varphi_{51}$
\item $\varphi_{147} :\equiv \forall w\,(s + -1 \cdot s = w \to \mathrm{Conv}(\Sigma a \circ \mathsf{s} \oplus -1 \odot \Sigma a,s + -1 \cdot s) \to \mathrm{Conv}(\Sigma a \circ \mathsf{s} \oplus -1 \odot \Sigma a,w))$
\item $\varphi_{148} :\equiv {\textstyle\sum_{k=1}^{n} a_{k}} + a_{S(n)} + -{\textstyle\sum_{k=1}^{n} a_{k}} = z + -{\textstyle\sum_{k=1}^{n} a_{k}} \to {\textstyle\sum_{k=1}^{n} a_{k}} + a_{S(n)} + -{\textstyle\sum_{k=1}^{n} a_{k}} = w + -{\textstyle\sum_{k=1}^{n} a_{k}}$
\item $\varphi_{149} :\equiv \forall z\,\forall w\,(z = w \to a_{\mathsf{s}(n)} = a_{z} \to a_{\mathsf{s}(n)} = a_{w}) \to \forall w\,(\mathsf{s}(n) = w \to a_{\mathsf{s}(n)} = a_{\mathsf{s}(n)} \to a_{\mathsf{s}(n)} = a_{w})$
\item $\varphi_{150} :\equiv (\mathrm{Seq}(a) \to \mathrm{Conv}(\Sigma a,s) \to n \in \mathbb{N} \to a_{S(n)} \in \mathbb{R}) \to \mathrm{Seq}(a) \to \varphi_{59}$
\item $\varphi_{151} :\equiv \forall b\,(\mathrm{Seq}(\Sigma a \circ \mathsf{s}) \to \mathrm{Seq}(b) \to \mathrm{Seq}(\Sigma a \circ \mathsf{s} \oplus b) \land \forall n \in \mathbb{N}\,(\Sigma a \circ \mathsf{s} \oplus b(n) = \Sigma a \circ \mathsf{s}(n) + b_{n}))$
\item $\varphi_{152} :\equiv (\mathsf{s}(n) = S(n) \to a_{\mathsf{s}(n)} = a_{\mathsf{s}(n)} \to a_{\mathsf{s}(n)} = a_{S(n)}) \to \varphi_{24}$
\item $\varphi_{153} :\equiv (\mathrm{Seq}(a) \to \mathrm{Conv}(\Sigma a,s) \to n \in \mathbb{N} \to \mathrm{Seq}(a)) \to \varphi_{109}$
\item $\varphi_{154} :\equiv (\mathrm{Seq}(a) \to \mathrm{Conv}(\Sigma a,s) \to n \in \mathbb{N} \to {\textstyle\sum_{k=1}^{n} a_{k}} \in \mathbb{R}) \to \mathrm{Seq}(a) \to \varphi_{58}$
\item $\varphi_{155} :\equiv (\mathrm{Seq}(a) \to \mathrm{Conv}(\Sigma a,s) \to n \in \mathbb{N} \to \mathrm{NSeq}(\mathsf{s})) \to \mathrm{Seq}(a) \to \mathrm{Conv}(\Sigma a,s) \to \varphi_{47}$
\item $\varphi_{156} :\equiv \forall b\,\forall \alpha\,(\mathrm{Seq}(b) \to \forall n \in \mathbb{N}\,(\Sigma a \circ \mathsf{s} \oplus -1 \odot \Sigma a(n) = b_{n}) \to \mathrm{Conv}(\Sigma a \circ \mathsf{s} \oplus -1 \odot \Sigma a,\alpha) \to \mathrm{Conv}(b,\alpha))$
\item $\varphi_{157} :\equiv (\mathrm{Seq}(a) \to \mathrm{Conv}(\Sigma a,s) \to n \in \mathbb{N} \to \mathrm{Seq}(\Sigma a)) \to \varphi_{108}$
\item $\varphi_{158} :\equiv \forall n \in \mathbb{N}\,(\Sigma a \circ \mathsf{s} \oplus -1 \odot \Sigma a(n) = a \circ \mathsf{s}(n)) \to \mathrm{Conv}(\Sigma a \circ \mathsf{s} \oplus -1 \odot \Sigma a,0) \to \mathrm{Conv}(a \circ \mathsf{s},0)$
\item $\varphi_{159} :\equiv \forall w\,({\textstyle\sum_{k=1}^{S(n)} a_{k}} = w \to {\textstyle\sum_{k=1}^{S(n)} a_{k}} + -{\textstyle\sum_{k=1}^{n} a_{k}} = {\textstyle\sum_{k=1}^{S(n)} a_{k}} + -{\textstyle\sum_{k=1}^{n} a_{k}} \to {\textstyle\sum_{k=1}^{S(n)} a_{k}} + -{\textstyle\sum_{k=1}^{n} a_{k}} = w + -{\textstyle\sum_{k=1}^{n} a_{k}})$
\item $\varphi_{160} :\equiv {\textstyle\sum_{k=1}^{S(n)} a_{k}} + -{\textstyle\sum_{k=1}^{n} a_{k}} = {\textstyle\sum_{k=1}^{S(n)} a_{k}} + -{\textstyle\sum_{k=1}^{n} a_{k}} \to {\textstyle\sum_{k=1}^{S(n)} a_{k}} + -{\textstyle\sum_{k=1}^{n} a_{k}} = {\textstyle\sum_{k=1}^{n} a_{k}} + a_{S(n)} + -{\textstyle\sum_{k=1}^{n} a_{k}}$
\item $\varphi_{161} :\equiv \forall w\,(\mathsf{s}(n) = w \to a_{\mathsf{s}(n)} = a_{\mathsf{s}(n)} \to a_{\mathsf{s}(n)} = a_{w}) \to \mathsf{s}(n) = S(n) \to a_{\mathsf{s}(n)} = a_{\mathsf{s}(n)} \to a_{\mathsf{s}(n)} = a_{S(n)}$
\item $\varphi_{162} :\equiv \forall y\,(a_{S(n)} + {\textstyle\sum_{k=1}^{n} a_{k}} - y = a_{S(n)} + {\textstyle\sum_{k=1}^{n} a_{k}} + -y) \to a_{S(n)} + {\textstyle\sum_{k=1}^{n} a_{k}} - {\textstyle\sum_{k=1}^{n} a_{k}} = a_{S(n)} + {\textstyle\sum_{k=1}^{n} a_{k}} + -{\textstyle\sum_{k=1}^{n} a_{k}}$
\item $\varphi_{163} :\equiv {\textstyle\sum_{k=1}^{S(n)} a_{k}} + -1 \cdot {\textstyle\sum_{k=1}^{n} a_{k}} = {\textstyle\sum_{k=1}^{S(n)} a_{k}} + -1 \cdot {\textstyle\sum_{k=1}^{n} a_{k}} \to {\textstyle\sum_{k=1}^{S(n)} a_{k}} + -1 \cdot {\textstyle\sum_{k=1}^{n} a_{k}} = {\textstyle\sum_{k=1}^{S(n)} a_{k}} + -{\textstyle\sum_{k=1}^{n} a_{k}}$
\item $\varphi_{164} :\equiv \Sigma a \circ \mathsf{s}(n) + -1 \odot \Sigma a(n) = z + -1 \odot \Sigma a(n) \to \Sigma a \circ \mathsf{s}(n) + -1 \odot \Sigma a(n) = w + -1 \odot \Sigma a(n)$
\item $\varphi_{165} :\equiv (\mathrm{Seq}(a) \to \mathrm{Conv}(\Sigma a,s) \to n \in \mathbb{N} \to \mathrm{NSeq}(\mathsf{s})) \to \mathrm{Seq}(a) \to \mathrm{Conv}(\Sigma a,s) \to \varphi_{78}$
\item $\varphi_{166} :\equiv \forall y\,(a_{S(n)} \in \mathbb{R} \to y \in \mathbb{R} \to a_{S(n)} + y - y = a_{S(n)}) \to \varphi_{96}$
\item $\varphi_{167} :\equiv {\textstyle\sum_{k=1}^{S(n)} a_{k}} + -1 \cdot {\textstyle\sum_{k=1}^{n} a_{k}} = {\textstyle\sum_{k=1}^{S(n)} a_{k}} + -{\textstyle\sum_{k=1}^{n} a_{k}} \to \Sigma a \circ \mathsf{s} \oplus -1 \odot \Sigma a(n) = {\textstyle\sum_{k=1}^{S(n)} a_{k}} + -{\textstyle\sum_{k=1}^{n} a_{k}}$
\item $\varphi_{168} :\equiv \forall y\,\forall z\,(\Sigma a \circ \mathsf{s}(n) = y \to y = z \to \Sigma a \circ \mathsf{s}(n) = z) \to \varphi_{104}$
\item $\varphi_{169} :\equiv (\mathrm{Seq}(a) \to \mathrm{Conv}(\Sigma a,s) \to n \in \mathbb{N} \to \mathrm{Seq}(\Sigma a)) \to \mathrm{Seq}(a) \to \mathrm{Conv}(\Sigma a,s) \to n \in \mathbb{N} \to \varphi_{68}$
\item $\varphi_{170} :\equiv \forall y\,({\textstyle\sum_{k=1}^{n} a_{k}} \in \mathbb{R} \to y \in \mathbb{R} \to {\textstyle\sum_{k=1}^{n} a_{k}} + y = y + {\textstyle\sum_{k=1}^{n} a_{k}}) \to \varphi_{97}$
\item $\varphi_{171} :\equiv \forall w\,(-1 \cdot {\textstyle\sum_{k=1}^{n} a_{k}} = w \to {\textstyle\sum_{k=1}^{S(n)} a_{k}} + -1 \cdot {\textstyle\sum_{k=1}^{n} a_{k}} = {\textstyle\sum_{k=1}^{S(n)} a_{k}} + -1 \cdot {\textstyle\sum_{k=1}^{n} a_{k}} \to {\textstyle\sum_{k=1}^{S(n)} a_{k}} + -1 \cdot {\textstyle\sum_{k=1}^{n} a_{k}} = {\textstyle\sum_{k=1}^{S(n)} a_{k}} + w)$
\item $\varphi_{172} :\equiv \forall \alpha\,(\mathrm{Seq}(a \circ \mathsf{s}) \to \forall n \in \mathbb{N}\,(\Sigma a \circ \mathsf{s} \oplus -1 \odot \Sigma a(n) = a \circ \mathsf{s}(n)) \to \mathrm{Conv}(\Sigma a \circ \mathsf{s} \oplus -1 \odot \Sigma a,\alpha) \to \mathrm{Conv}(a \circ \mathsf{s},\alpha))$
\item $\varphi_{173} :\equiv {\textstyle\sum_{k=1}^{S(n)} a_{k}} + -1 \odot \Sigma a(n) = {\textstyle\sum_{k=1}^{S(n)} a_{k}} + -1 \odot \Sigma a(n) \to {\textstyle\sum_{k=1}^{S(n)} a_{k}} + -1 \odot \Sigma a(n) = {\textstyle\sum_{k=1}^{S(n)} a_{k}} + -1 \cdot {\textstyle\sum_{k=1}^{n} a_{k}}$
\item $\varphi_{174} :\equiv (\mathrm{Seq}(a) \to \mathrm{Conv}(\Sigma a,s) \to n \in \mathbb{N} \to \mathrm{Seq}(a)) \to \mathrm{Seq}(a) \to \mathrm{Conv}(\Sigma a,s) \to n \in \mathbb{N} \to \varphi_{98}$
\item $\varphi_{175} :\equiv {\textstyle\sum_{k=1}^{S(n)} a_{k}} + -1 \odot \Sigma a(n) = {\textstyle\sum_{k=1}^{S(n)} a_{k}} + -1 \cdot {\textstyle\sum_{k=1}^{n} a_{k}} \to \Sigma a \circ \mathsf{s} \oplus -1 \odot \Sigma a(n) = {\textstyle\sum_{k=1}^{S(n)} a_{k}} + -1 \cdot {\textstyle\sum_{k=1}^{n} a_{k}}$
\item $\varphi_{176} :\equiv (\mathsf{s}(n) = S(n) \to {\textstyle\sum_{k=1}^{\mathsf{s}(n)} a_{k}} = {\textstyle\sum_{k=1}^{\mathsf{s}(n)} a_{k}} \to {\textstyle\sum_{k=1}^{\mathsf{s}(n)} a_{k}} = {\textstyle\sum_{k=1}^{S(n)} a_{k}}) \to \varphi_{37}$
\item $\varphi_{177} :\equiv \forall z\,(\Sigma a \circ \mathsf{s} \oplus -1 \odot \Sigma a(n) = {\textstyle\sum_{k=1}^{S(n)} a_{k}} + -{\textstyle\sum_{k=1}^{n} a_{k}} \to {\textstyle\sum_{k=1}^{S(n)} a_{k}} + -{\textstyle\sum_{k=1}^{n} a_{k}} = z \to \Sigma a \circ \mathsf{s} \oplus -1 \odot \Sigma a(n) = z)$
\item $\varphi_{178} :\equiv (\mathrm{Seq}(a) \to \mathrm{Conv}(\Sigma a,s) \to n \in \mathbb{N} \to n \in \mathbb{N}) \to \mathrm{Seq}(a) \to \mathrm{Conv}(\Sigma a,s) \to \varphi_{121}$
\item $\varphi_{179} :\equiv (\mathrm{Seq}(a) \to \mathrm{Conv}(\Sigma a,s) \to n \in \mathbb{N} \to \varphi_{98}) \to \mathrm{Seq}(a) \to \mathrm{Conv}(\Sigma a,s) \to n \in \mathbb{N} \to \mathrm{Seq}(\Sigma a)$
\item $\varphi_{180} :\equiv {\textstyle\sum_{k=1}^{S(n)} a_{k}} + -{\textstyle\sum_{k=1}^{n} a_{k}} = {\textstyle\sum_{k=1}^{n} a_{k}} + a_{S(n)} + -{\textstyle\sum_{k=1}^{n} a_{k}} \to \Sigma a \circ \mathsf{s} \oplus -1 \odot \Sigma a(n) = {\textstyle\sum_{k=1}^{n} a_{k}} + a_{S(n)} + -{\textstyle\sum_{k=1}^{n} a_{k}}$
\item $\varphi_{181} :\equiv a_{S(n)} + {\textstyle\sum_{k=1}^{n} a_{k}} - {\textstyle\sum_{k=1}^{n} a_{k}} = a_{S(n)} + {\textstyle\sum_{k=1}^{n} a_{k}} + -{\textstyle\sum_{k=1}^{n} a_{k}} \to a_{S(n)} + {\textstyle\sum_{k=1}^{n} a_{k}} + -{\textstyle\sum_{k=1}^{n} a_{k}} = a_{S(n)} + {\textstyle\sum_{k=1}^{n} a_{k}} - {\textstyle\sum_{k=1}^{n} a_{k}}$
\item $\varphi_{182} :\equiv a_{S(n)} + {\textstyle\sum_{k=1}^{n} a_{k}} + -{\textstyle\sum_{k=1}^{n} a_{k}} = a_{S(n)} + {\textstyle\sum_{k=1}^{n} a_{k}} - {\textstyle\sum_{k=1}^{n} a_{k}} \to \Sigma a \circ \mathsf{s} \oplus -1 \odot \Sigma a(n) = a_{S(n)} + {\textstyle\sum_{k=1}^{n} a_{k}} - {\textstyle\sum_{k=1}^{n} a_{k}}$
\item $\varphi_{183} :\equiv \forall w\,(-1 \odot \Sigma a(n) = w \to {\textstyle\sum_{k=1}^{S(n)} a_{k}} + -1 \odot \Sigma a(n) = {\textstyle\sum_{k=1}^{S(n)} a_{k}} + -1 \odot \Sigma a(n) \to {\textstyle\sum_{k=1}^{S(n)} a_{k}} + -1 \odot \Sigma a(n) = {\textstyle\sum_{k=1}^{S(n)} a_{k}} + w)$
\item $\varphi_{184} :\equiv \forall z\,\forall w\,(z = w \to {\textstyle\sum_{k=1}^{\mathsf{s}(n)} a_{k}} = {\textstyle\sum_{k=1}^{z} a_{k}} \to {\textstyle\sum_{k=1}^{\mathsf{s}(n)} a_{k}} = {\textstyle\sum_{k=1}^{w} a_{k}}) \to \forall w\,(\mathsf{s}(n) = w \to {\textstyle\sum_{k=1}^{\mathsf{s}(n)} a_{k}} = {\textstyle\sum_{k=1}^{\mathsf{s}(n)} a_{k}} \to {\textstyle\sum_{k=1}^{\mathsf{s}(n)} a_{k}} = {\textstyle\sum_{k=1}^{w} a_{k}})$
\item $\varphi_{185} :\equiv \mathrm{Conv}(\Sigma a,s) \to n \in \mathbb{N} \to \mathrm{Seq}(\Sigma a \circ \mathsf{s} \oplus -1 \odot \Sigma a) \land \varphi_{126}$
\item $\varphi_{186} :\equiv \forall z\,(\Sigma a \circ \mathsf{s} \oplus -1 \odot \Sigma a(n) = {\textstyle\sum_{k=1}^{S(n)} a_{k}} + -1 \cdot {\textstyle\sum_{k=1}^{n} a_{k}} \to {\textstyle\sum_{k=1}^{S(n)} a_{k}} + -1 \cdot {\textstyle\sum_{k=1}^{n} a_{k}} = z \to \Sigma a \circ \mathsf{s} \oplus -1 \odot \Sigma a(n) = z)$
\item $\varphi_{187} :\equiv \Sigma a \circ \mathsf{s}(n) + -1 \odot \Sigma a(n) = {\textstyle\sum_{k=1}^{S(n)} a_{k}} + -1 \odot \Sigma a(n) \to \Sigma a \circ \mathsf{s} \oplus -1 \odot \Sigma a(n) = {\textstyle\sum_{k=1}^{S(n)} a_{k}} + -1 \odot \Sigma a(n)$
\item $\varphi_{188} :\equiv (\mathrm{Seq}(a) \to \mathrm{Conv}(\Sigma a,s) \to n \in \mathbb{N} \to \mathrm{Seq}(a)) \to \varphi_{155}$
\item $\varphi_{189} :\equiv (\mathrm{Seq}(a) \to \mathrm{Conv}(\Sigma a,s) \to s + -1 \cdot s = s + -s) \to (\mathrm{Seq}(a) \to \mathrm{Conv}(\Sigma a,s) \to s + -s = 0) \to \varphi_{107}$
\item $\varphi_{190} :\equiv \forall z\,(\Sigma a \circ \mathsf{s} \oplus -1 \odot \Sigma a(n) = a_{S(n)} + {\textstyle\sum_{k=1}^{n} a_{k}} - {\textstyle\sum_{k=1}^{n} a_{k}} \to a_{S(n)} + {\textstyle\sum_{k=1}^{n} a_{k}} - {\textstyle\sum_{k=1}^{n} a_{k}} = z \to \Sigma a \circ \mathsf{s} \oplus -1 \odot \Sigma a(n) = z)$
\item $\varphi_{191} :\equiv \varphi_{85} \to \Sigma a \circ \mathsf{s} \oplus -1 \odot \Sigma a(n) = a_{S(n)} + {\textstyle\sum_{k=1}^{n} a_{k}} + -{\textstyle\sum_{k=1}^{n} a_{k}}$
\item $\varphi_{192} :\equiv \forall z\,(\Sigma a \circ \mathsf{s} \oplus -1 \odot \Sigma a(n) = {\textstyle\sum_{k=1}^{S(n)} a_{k}} + -1 \odot \Sigma a(n) \to {\textstyle\sum_{k=1}^{S(n)} a_{k}} + -1 \odot \Sigma a(n) = z \to \Sigma a \circ \mathsf{s} \oplus -1 \odot \Sigma a(n) = z)$
\item $\varphi_{193} :\equiv (\mathrm{Seq}(a) \to \mathrm{Conv}(\Sigma a,s) \to \mathrm{Conv}(\Sigma a \circ \mathsf{s},s)) \to \varphi_{14} \to \mathrm{Seq}(a) \to \varphi_{53}$
\item $\varphi_{194} :\equiv \varphi_{110} \to \Sigma a \circ \mathsf{s}(n) + -1 \odot \Sigma a(n) = {\textstyle\sum_{k=1}^{S(n)} a_{k}} + -1 \odot \Sigma a(n)$
\item $\varphi_{195} :\equiv \forall w\,({\textstyle\sum_{k=1}^{n} a_{k}} + a_{S(n)} = w \to \varphi_{84} \to {\textstyle\sum_{k=1}^{n} a_{k}} + a_{S(n)} + -{\textstyle\sum_{k=1}^{n} a_{k}} = w + -{\textstyle\sum_{k=1}^{n} a_{k}})$
\item $\varphi_{196} :\equiv \forall z\,(\Sigma a \circ \mathsf{s} \oplus -1 \odot \Sigma a(n) = a_{S(n)} + {\textstyle\sum_{k=1}^{n} a_{k}} + -{\textstyle\sum_{k=1}^{n} a_{k}} \to a_{S(n)} + {\textstyle\sum_{k=1}^{n} a_{k}} + -{\textstyle\sum_{k=1}^{n} a_{k}} = z \to \Sigma a \circ \mathsf{s} \oplus -1 \odot \Sigma a(n) = z)$
\item $\varphi_{197} :\equiv \forall z\,(\Sigma a \circ \mathsf{s} \oplus -1 \odot \Sigma a(n) = {\textstyle\sum_{k=1}^{n} a_{k}} + a_{S(n)} + -{\textstyle\sum_{k=1}^{n} a_{k}} \to {\textstyle\sum_{k=1}^{n} a_{k}} + a_{S(n)} + -{\textstyle\sum_{k=1}^{n} a_{k}} = z \to \Sigma a \circ \mathsf{s} \oplus -1 \odot \Sigma a(n) = z)$
\item $\varphi_{198} :\equiv (\mathrm{Seq}(a) \to \mathrm{Conv}(\Sigma a,s) \to n \in \mathbb{N} \to a_{S(n)} \in \mathbb{R}) \to \varphi_{154}$
\item $\varphi_{199} :\equiv (\mathrm{Seq}(a) \to \mathrm{Conv}(\Sigma a,s) \to n \in \mathbb{N} \to {\textstyle\sum_{k=1}^{n} a_{k}} \in \mathbb{R}) \to \varphi_{150}$
\item $\varphi_{200} :\equiv \forall w\,(\mathsf{s}(n) = w \to {\textstyle\sum_{k=1}^{\mathsf{s}(n)} a_{k}} = {\textstyle\sum_{k=1}^{\mathsf{s}(n)} a_{k}} \to {\textstyle\sum_{k=1}^{\mathsf{s}(n)} a_{k}} = {\textstyle\sum_{k=1}^{w} a_{k}}) \to \mathsf{s}(n) = S(n) \to {\textstyle\sum_{k=1}^{\mathsf{s}(n)} a_{k}} = {\textstyle\sum_{k=1}^{\mathsf{s}(n)} a_{k}} \to {\textstyle\sum_{k=1}^{\mathsf{s}(n)} a_{k}} = {\textstyle\sum_{k=1}^{S(n)} a_{k}}$
\item $\varphi_{201} :\equiv (\mathrm{Seq}(a) \to \mathrm{Conv}(\Sigma a,s) \to -1 \cdot s = -s) \to s + -1 \cdot s = s + -1 \cdot s \to \varphi_{142}$
\item $\varphi_{202} :\equiv \forall z\,(\Sigma a \circ \mathsf{s} \oplus -1 \odot \Sigma a(n) = \Sigma a \circ \mathsf{s}(n) + -1 \odot \Sigma a(n) \to \Sigma a \circ \mathsf{s}(n) + -1 \odot \Sigma a(n) = z \to \Sigma a \circ \mathsf{s} \oplus -1 \odot \Sigma a(n) = z)$
\item $\varphi_{203} :\equiv \mathrm{Seq}(\Sigma a \circ \mathsf{s}) \to \mathrm{Seq}(-1 \odot \Sigma a) \to \mathrm{Seq}(\Sigma a \circ \mathsf{s} \oplus -1 \odot \Sigma a) \land \varphi_{126}$
\item $\varphi_{204} :\equiv (\mathrm{Seq}(a) \to \mathrm{Conv}(\Sigma a,s) \to n \in \mathbb{N} \to \mathrm{Seq}(\Sigma a)) \to \varphi_{165}$
\item $\varphi_{205} :\equiv \forall w\,(\Sigma a \circ \mathsf{s}(n) = w \to \varphi_{110} \to \Sigma a \circ \mathsf{s}(n) + -1 \odot \Sigma a(n) = w + -1 \odot \Sigma a(n))$
\item $\varphi_{206} :\equiv (\mathrm{Seq}(a) \to \mathrm{Conv}(\Sigma a,s) \to n \in \mathbb{N} \to \varphi_{68}) \to \mathrm{Seq}(a) \to \varphi_{74}$
\item $\varphi_{207} :\equiv (\mathrm{Seq}(a) \to \mathrm{Conv}(\Sigma a,s) \to n \in \mathbb{N} \to -1 \in \mathbb{R}) \to \varphi_{169}$
\item $\varphi_{208} :\equiv (\mathrm{Seq}(a) \to \mathrm{Conv}(\Sigma a,s) \to n \in \mathbb{N} \to \varphi_{98}) \to \mathrm{Seq}(a) \to \varphi_{64}$
\item $\varphi_{209} :\equiv (\mathrm{Seq}(a) \to \varphi_{52}) \to (\mathrm{Seq}(a) \to \mathrm{Conv}(\Sigma a,s) \to \mathrm{Seq}(a)) \to \varphi_{138}$
\item $\varphi_{210} :\equiv (\Sigma a \circ \mathsf{s} \oplus -1 \odot \Sigma a(n) = a_{S(n)} \to \varphi_{95}) \to \mathrm{Seq}(a) \to \varphi_{79}$
\item $\varphi_{211} :\equiv (\mathrm{Seq}(a) \to \mathrm{Conv}(\Sigma a,s) \to \mathrm{Seq}(a \circ \mathsf{s})) \to (\mathrm{Seq}(a) \to \varphi_{92}) \to \varphi_{132}$
\item $\varphi_{212} :\equiv \varphi_{181} \to a_{S(n)} + {\textstyle\sum_{k=1}^{n} a_{k}} + -{\textstyle\sum_{k=1}^{n} a_{k}} = a_{S(n)} + {\textstyle\sum_{k=1}^{n} a_{k}} - {\textstyle\sum_{k=1}^{n} a_{k}}$
\item $\varphi_{213} :\equiv \Sigma a \circ \mathsf{s} \oplus -1 \odot \Sigma a(n) = a_{S(n)} + {\textstyle\sum_{k=1}^{n} a_{k}} + -{\textstyle\sum_{k=1}^{n} a_{k}} \to \varphi_{182}$
\item $\varphi_{214} :\equiv \Sigma a \circ \mathsf{s} \oplus -1 \odot \Sigma a(n) = {\textstyle\sum_{k=1}^{n} a_{k}} + a_{S(n)} + -{\textstyle\sum_{k=1}^{n} a_{k}} \to \varphi_{191}$
\item $\varphi_{215} :\equiv \Sigma a \circ \mathsf{s} \oplus -1 \odot \Sigma a(n) = \Sigma a \circ \mathsf{s}(n) + -1 \odot \Sigma a(n) \to \varphi_{187}$
\item $\varphi_{216} :\equiv (\Sigma a \circ \mathsf{s} \oplus -1 \odot \Sigma a(n) = a_{S(n)} + {\textstyle\sum_{k=1}^{n} a_{k}} - {\textstyle\sum_{k=1}^{n} a_{k}} \to \varphi_{128}) \to \mathrm{Seq}(a) \to \varphi_{70}$
\item $\varphi_{217} :\equiv ({\textstyle\sum_{k=1}^{S(n)} a_{k}} = {\textstyle\sum_{k=1}^{n} a_{k}} + a_{S(n)} \to \varphi_{160}) \to \mathrm{Seq}(a) \to \mathrm{Conv}(\Sigma a,s) \to \varphi_{86}$
\item $\varphi_{218} :\equiv (\mathrm{Seq}(a) \to \mathrm{Conv}(\Sigma a,s) \to s + -1 \cdot s = 0) \to (\mathrm{Seq}(a) \to \varphi_{53}) \to (s + -1 \cdot s = 0 \to \varphi_{127}) \to \varphi_{36}$
\item $\varphi_{219} :\equiv (-1 \odot \Sigma a(n) = -1 \cdot {\textstyle\sum_{k=1}^{n} a_{k}} \to \varphi_{173}) \to \mathrm{Seq}(a) \to \mathrm{Conv}(\Sigma a,s) \to \varphi_{91}$
\item $\varphi_{220} :\equiv (\mathrm{Seq}(a) \to \varphi_{185}) \to \mathrm{Seq}(a) \to \mathrm{Conv}(\Sigma a,s) \to n \in \mathbb{N} \to \varphi_{126}$
\item $\varphi_{221} :\equiv a_{S(n)} + {\textstyle\sum_{k=1}^{n} a_{k}} - {\textstyle\sum_{k=1}^{n} a_{k}} = a_{S(n)} + {\textstyle\sum_{k=1}^{n} a_{k}} + -{\textstyle\sum_{k=1}^{n} a_{k}} \to \varphi_{212}$
\item $\varphi_{222} :\equiv (\Sigma a \circ \mathsf{s}(n) = {\textstyle\sum_{k=1}^{S(n)} a_{k}} \to \varphi_{194}) \to \mathrm{Seq}(a) \to \mathrm{Conv}(\Sigma a,s) \to \varphi_{116}$
\item $\varphi_{223} :\equiv (\Sigma a \circ \mathsf{s} \oplus -1 \odot \Sigma a(n) = {\textstyle\sum_{k=1}^{S(n)} a_{k}} + -1 \cdot {\textstyle\sum_{k=1}^{n} a_{k}} \to \varphi_{167}) \to \mathrm{Seq}(a) \to \varphi_{123}$
\item $\varphi_{224} :\equiv ({\textstyle\sum_{k=1}^{n} a_{k}} + a_{S(n)} = a_{S(n)} + {\textstyle\sum_{k=1}^{n} a_{k}} \to \varphi_{84} \to \varphi_{85}) \to \mathrm{Seq}(a) \to \mathrm{Conv}(\Sigma a,s) \to n \in \mathbb{N} \to \varphi_{85}$
\item $\varphi_{225} :\equiv (\Sigma a \circ \mathsf{s} \oplus -1 \odot \Sigma a(n) = {\textstyle\sum_{k=1}^{S(n)} a_{k}} + -1 \odot \Sigma a(n) \to \varphi_{175}) \to \mathrm{Seq}(a) \to \mathrm{Conv}(\Sigma a,s) \to \varphi_{103}$
\item $\varphi_{226} :\equiv (\Sigma a \circ \mathsf{s} \oplus -1 \odot \Sigma a(n) = {\textstyle\sum_{k=1}^{S(n)} a_{k}} + -{\textstyle\sum_{k=1}^{n} a_{k}} \to \varphi_{180}) \to \mathrm{Seq}(a) \to \mathrm{Conv}(\Sigma a,s) \to \varphi_{118}$
\item $\varphi_{227} :\equiv \varphi_{190} \to \Sigma a \circ \mathsf{s} \oplus -1 \odot \Sigma a(n) = a_{S(n)} + {\textstyle\sum_{k=1}^{n} a_{k}} - {\textstyle\sum_{k=1}^{n} a_{k}} \to \varphi_{128}$
\item $\varphi_{228} :\equiv {\textstyle\sum_{k=1}^{S(n)} a_{k}} + -1 \cdot {\textstyle\sum_{k=1}^{n} a_{k}} = {\textstyle\sum_{k=1}^{S(n)} a_{k}} + -1 \cdot {\textstyle\sum_{k=1}^{n} a_{k}} \to (-1 \cdot {\textstyle\sum_{k=1}^{n} a_{k}} = -{\textstyle\sum_{k=1}^{n} a_{k}} \to \varphi_{163}) \to \mathrm{Seq}(a) \to \varphi_{114}$
\item $\varphi_{229} :\equiv \varphi_{192} \to \Sigma a \circ \mathsf{s} \oplus -1 \odot \Sigma a(n) = {\textstyle\sum_{k=1}^{S(n)} a_{k}} + -1 \odot \Sigma a(n) \to \varphi_{175}$
\item $\varphi_{230} :\equiv a_{S(n)} + {\textstyle\sum_{k=1}^{n} a_{k}} + -{\textstyle\sum_{k=1}^{n} a_{k}} = a_{S(n)} + {\textstyle\sum_{k=1}^{n} a_{k}} - {\textstyle\sum_{k=1}^{n} a_{k}} \to \varphi_{213} \to \mathrm{Seq}(a) \to \mathrm{Conv}(\Sigma a,s) \to \varphi_{111}$
\item $\varphi_{231} :\equiv (\mathrm{Seq}(a) \to \varphi_{40}) \to {\textstyle\sum_{k=1}^{S(n)} a_{k}} + -1 \odot \Sigma a(n) = {\textstyle\sum_{k=1}^{S(n)} a_{k}} + -1 \odot \Sigma a(n) \to \varphi_{219}$
\item $\varphi_{232} :\equiv (\mathrm{Seq}(a) \to \mathrm{Conv}(\Sigma a,s) \to n \in \mathbb{N} \to \varphi_{126}) \to (\varphi_{126} \to \varphi_{121}) \to \mathrm{Seq}(a) \to \mathrm{Conv}(\Sigma a,s) \to n \in \mathbb{N} \to \varphi_{121}$
\item $\varphi_{233} :\equiv (\mathrm{Seq}(a) \to \mathrm{Conv}(\Sigma a,s) \to \varphi_{111}) \to (\mathrm{Seq}(a) \to \varphi_{58}) \to \varphi_{216}$
\item $\varphi_{234} :\equiv (\mathrm{Seq}(a) \to \mathrm{Conv}(\Sigma a,s) \to \varphi_{116}) \to \varphi_{215} \to \mathrm{Seq}(a) \to \mathrm{Conv}(\Sigma a,s) \to \varphi_{112}$
\item $\varphi_{235} :\equiv (\mathrm{Seq}(a) \to \mathrm{Conv}(\Sigma a,s) \to n \in \mathbb{N} \to \varphi_{85}) \to \varphi_{214} \to \mathrm{Seq}(a) \to \mathrm{Conv}(\Sigma a,s) \to \varphi_{119}$
\item $\varphi_{236} :\equiv (\mathrm{Seq}(a) \to \mathrm{Conv}(\Sigma a,s) \to \varphi_{103}) \to (\mathrm{Seq}(a) \to \varphi_{114}) \to \varphi_{223}$
\item $\varphi_{237} :\equiv (\mathrm{Seq}(a) \to \varphi_{123}) \to (\mathrm{Seq}(a) \to \mathrm{Conv}(\Sigma a,s) \to \varphi_{86}) \to \varphi_{226}$
\item $\varphi_{238} :\equiv (\mathrm{Seq}(a) \to \mathrm{Conv}(\Sigma a,s) \to \varphi_{112}) \to (\mathrm{Seq}(a) \to \mathrm{Conv}(\Sigma a,s) \to \varphi_{91}) \to \varphi_{225}$
\end{abbrevlist}
\end{fullproof}

% 自動生成 (prover/build.js): GeoPartial  ノード数 787
\begin{fullproof}{[Q3]}{$\forall r \in \mathbb{R}\,\forall n \in \mathbb{N}\,((1 - r) \cdot {\textstyle\sum_{k=1}^{n} {{r}^{k}}} = r - {{r}^{S(n)}})$}{787}\phantomsection\label{tree:GeoPartial}
\noindent$\Pi_{1}$:
\begin{ptree}
\AXC{$\forall n\,(n \in \mathbb{N} \leftrightarrow n \in \omega \land n \neq \emptyset)$\lbl{Def$\mathbb{N}$}}
\AXC{$\varphi_{50}$\lbl{A2}}
\BIC{$n \in \mathbb{N} \leftrightarrow n \in \omega \land n \neq \emptyset$}
\AXC{$r \in \mathbb{R} \to n \in \mathbb{N} \to n \in \mathbb{N}$\lbl{A1}}
\AXC{$\varphi_{127}$\lbl{A1}}
\BIC{$(n \in \mathbb{N} \leftrightarrow n \in \omega \land n \neq \emptyset) \to r \in \mathbb{R} \to n \in \mathbb{N} \to n \in \omega$}
\BIC{$r \in \mathbb{R} \to n \in \mathbb{N} \to n \in \omega$}
\end{ptree}
\noindent$\Pi_{2}$:
\begin{ptree}
\AXC{$\varphi_{20}$\lbl{A4}}
\AXC{$\forall y\,(\varphi_{6}) \to \varphi_{23}$\lbl{A3}}
\AXC{$\varphi_{6}$\lbl{A2}}
\UIC{$\forall y\,(\varphi_{6})$}
\AXC{$\forall y\,(\varphi_{6}) \to (\forall y\,(\varphi_{6}) \to \varphi_{23}) \to \varphi_{20} \to \varphi_{11}$\lbl{A1}}
\BIC{$(\forall y\,(\varphi_{6}) \to \varphi_{23}) \to \varphi_{20} \to \varphi_{11}$}
\BIC{$\varphi_{20} \to \varphi_{11}$}
\BIC{$\varphi_{11}$}
\end{ptree}
\noindent$\Pi_{3}$:
\begin{ptree}
\AXC{$\varphi_{21}$\lbl{A4}}
\AXC{$\forall n\,(\varphi_{7}) \to \varphi_{24}$\lbl{A3}}
\AXC{$\varphi_{7}$\lbl{A2}}
\UIC{$\forall n\,(\varphi_{7})$}
\AXC{$\forall n\,(\varphi_{7}) \to (\forall n\,(\varphi_{7}) \to \varphi_{24}) \to \varphi_{21} \to \varphi_{12}$\lbl{A1}}
\BIC{$(\forall n\,(\varphi_{7}) \to \varphi_{24}) \to \varphi_{21} \to \varphi_{12}$}
\BIC{$\varphi_{21} \to \varphi_{12}$}
\BIC{$\varphi_{12}$}
\end{ptree}
\noindent$\Pi_{4}$:
\begin{ptree}
\AXC{$y = z \to x = y \to x = z$\lbl{A6}}
\AXC{$(y = z \to x = y \to x = z) \to x = y \to y = z \to x = z$\lbl{A1}}
\BIC{$x = y \to y = z \to x = z$}
\UIC{$\forall z\,(x = y \to y = z \to x = z)$}
\UIC{$\forall y\,\forall z\,(x = y \to y = z \to x = z)$}
\UIC{$\forall x\,\forall y\,\forall z\,(x = y \to y = z \to x = z)$}
\AXC{$\forall x\,\forall y\,\forall z\,(x = y \to y = z \to x = z) \to \varphi_{56}$\lbl{A2}}
\BIC{$\varphi_{56}$}
\end{ptree}
\noindent$\Pi_{5}$:
\begin{ptree}
\AXC{$\stackrel{\Pi_{4}}{\vdots}$}\noLine
\UIC{$\varphi_{56}$}
\AXC{$\varphi_{56} \to \varphi_{145}$\lbl{A2}}
\BIC{$\varphi_{145}$}
\AXC{$\varphi_{145} \to (1 - r) \cdot {\textstyle\sum_{k=1}^{S(n)} {{r}^{k}}} = r - {{r}^{S(n)}} \cdot r \to \varphi_{120}$\lbl{A2}}
\BIC{$(1 - r) \cdot {\textstyle\sum_{k=1}^{S(n)} {{r}^{k}}} = r - {{r}^{S(n)}} \cdot r \to \varphi_{120}$}
\end{ptree}
\noindent$\Pi_{6}$:
\begin{ptree}
\AXC{$z = w \to r - {{r}^{S(n)}} \cdot r = r - z \to r - {{r}^{S(n)}} \cdot r = r - w$\lbl{A6}}
\UIC{$\varphi_{41}$}
\UIC{$\forall z\,(\varphi_{41})$}
\AXC{$\forall z\,(\varphi_{41}) \to \varphi_{129}$\lbl{A2}}
\BIC{$\varphi_{129}$}
\AXC{$\varphi_{129} \to {{r}^{S(n)}} \cdot r = {{r}^{S(S(n))}} \to \varphi_{110}$\lbl{A2}}
\BIC{${{r}^{S(n)}} \cdot r = {{r}^{S(S(n))}} \to \varphi_{110}$}
\end{ptree}
\noindent$\Pi_{7}$:
\begin{ptree}
\AXC{$x = x$\lbl{A5}}
\AXC{$x = y \to x = x \to y = x$\lbl{A6}}
\AXC{$(x = y \to x = x \to y = x) \to x = x \to x = y \to y = x$\lbl{A1}}
\BIC{$x = x \to x = y \to y = x$}
\BIC{$x = y \to y = x$}
\UIC{$\forall y\,(x = y \to y = x)$}
\UIC{$\forall x\,\forall y\,(x = y \to y = x)$}
\end{ptree}
\noindent$\Pi_{8}$:
\begin{ptree}
\AXC{$\forall n\,(\varphi_{66}) \to \varphi_{66}$\lbl{A2}}
\AXC{$r \in \mathbb{R} \to \forall n\,(\varphi_{66}) \to \forall n\,(\varphi_{66})$\lbl{A1}}
\AXC{$\varphi_{276}$\lbl{A1}}
\BIC{$(\forall n\,(\varphi_{66}) \to \varphi_{66}) \to r \in \mathbb{R} \to \forall n\,(\varphi_{66}) \to \varphi_{66}$}
\BIC{$r \in \mathbb{R} \to \forall n\,(\varphi_{66}) \to \varphi_{66}$}
\AXC{$(r \in \mathbb{R} \to \forall n\,(\varphi_{66}) \to \varphi_{66}) \to \varphi_{232}$\lbl{A1}}
\BIC{$\varphi_{232}$}
\end{ptree}
\noindent$\Pi_{9}$:
\begin{ptree}
\AXC{$r \in \mathbb{R} \to \forall n\,(\varphi_{66}) \to n \in B_{51} \to n \in B_{51}$\lbl{A1}}
\AXC{$\stackrel{\Pi_{8}}{\vdots}$}\noLine
\UIC{$\varphi_{232}$}
\BIC{$r \in \mathbb{R} \to \forall n\,(\varphi_{66}) \to \varphi_{67}$}
\end{ptree}
\noindent$\Pi_{10}$:
\begin{ptree}
\AXC{$\stackrel{\Pi_{9}}{\vdots}$}\noLine
\UIC{$r \in \mathbb{R} \to \forall n\,(\varphi_{66}) \to \varphi_{67}$}
\AXC{$\varphi_{233}$\lbl{A1}}
\BIC{$r \in \mathbb{R} \to \forall n\,(\varphi_{66}) \to n \in B_{51} \to n \in \omega$}
\end{ptree}
\noindent$\Pi_{11}$:
\begin{ptree}
\AXC{$\varphi_{180}$\lbl{A1}}
\AXC{$\stackrel{\Pi_{10}}{\vdots}$}\noLine
\UIC{$r \in \mathbb{R} \to \forall n\,(\varphi_{66}) \to n \in B_{51} \to n \in \omega$}
\AXC{$\forall n\,(n \in \mathbb{N} \leftrightarrow n \in \omega \land n \neq \emptyset)$\lbl{Def$\mathbb{N}$}}
\AXC{$\varphi_{50}$\lbl{A2}}
\BIC{$n \in \mathbb{N} \leftrightarrow n \in \omega \land n \neq \emptyset$}
\AXC{$\varphi_{261}$\lbl{A1}}
\BIC{$(r \in \mathbb{R} \to \forall n\,(\varphi_{66}) \to n \in B_{51} \to n \in \omega) \to \varphi_{180} \to \varphi_{177}$}
\BIC{$\varphi_{180} \to \varphi_{177}$}
\BIC{$\varphi_{177}$}
\end{ptree}
\noindent$\Pi_{12}$:
\begin{ptree}
\AXC{$\forall n \in \mathbb{N}\,(S(n) \in \mathbb{N})$\lbl{[N2]}}
\AXC{$\forall n \in \mathbb{N}\,(S(n) \in \mathbb{N}) \to n \in \mathbb{N} \to S(n) \in \mathbb{N}$\lbl{A2}}
\BIC{$n \in \mathbb{N} \to S(n) \in \mathbb{N}$}
\end{ptree}
\noindent$\Pi_{13}$:
\begin{ptree}
\AXC{$\stackrel{\Pi_{11}}{\vdots}$}\noLine
\UIC{$\varphi_{177}$}
\AXC{$\stackrel{\Pi_{12}}{\vdots}$}\noLine
\UIC{$n \in \mathbb{N} \to S(n) \in \mathbb{N}$}
\AXC{$(n \in \mathbb{N} \to S(n) \in \mathbb{N}) \to \varphi_{177} \to r \in \mathbb{R} \to \varphi_{167}$\lbl{A1}}
\BIC{$\varphi_{177} \to r \in \mathbb{R} \to \varphi_{167}$}
\BIC{$r \in \mathbb{R} \to \varphi_{167}$}
\end{ptree}
\noindent$\Pi_{14}$:
\begin{ptree}
\AXC{$\varphi_{179}$\lbl{A1}}
\AXC{$\forall r \in \mathbb{R}\,(\varphi_{58})$\lbl{[S24]}}
\AXC{$\forall r \in \mathbb{R}\,(\varphi_{58}) \to r \in \mathbb{R} \to \varphi_{58}$\lbl{A2}}
\BIC{$r \in \mathbb{R} \to \varphi_{58}$}
\AXC{$(r \in \mathbb{R} \to \varphi_{58}) \to \varphi_{246}$\lbl{A1}}
\BIC{$\varphi_{246}$}
\BIC{$r \in \mathbb{R} \to \forall n\,(\varphi_{66}) \to n \in B_{51} \to S(n) \in \mathbb{N} \to n \neq \emptyset \to \varphi_{58}$}
\AXC{$\varphi_{255}$\lbl{A1}}
\BIC{$r \in \mathbb{R} \to \forall n\,(\varphi_{66}) \to \varphi_{72}$}
\AXC{$\varphi_{265}$\lbl{A1}}
\BIC{$\varphi_{116} \to r \in \mathbb{R} \to \forall n\,(\varphi_{66}) \to n \in B_{51} \to \varphi_{61}$}
\end{ptree}
\noindent$\Pi_{15}$:
\begin{ptree}
\AXC{$\stackrel{\Pi_{13}}{\vdots}$}\noLine
\UIC{$r \in \mathbb{R} \to \varphi_{167}$}
\AXC{$\varphi_{116}$\lbl{A2}}
\AXC{$\stackrel{\Pi_{14}}{\vdots}$}\noLine
\UIC{$\varphi_{116} \to r \in \mathbb{R} \to \forall n\,(\varphi_{66}) \to n \in B_{51} \to \varphi_{61}$}
\BIC{$r \in \mathbb{R} \to \forall n\,(\varphi_{66}) \to n \in B_{51} \to \varphi_{61}$}
\AXC{$\varphi_{270}$\lbl{A1}}
\BIC{$(r \in \mathbb{R} \to \varphi_{167}) \to r \in \mathbb{R} \to \forall n\,(\varphi_{66}) \to \varphi_{53}$}
\BIC{$r \in \mathbb{R} \to \forall n\,(\varphi_{66}) \to \varphi_{53}$}
\AXC{$(r \in \mathbb{R} \to \forall n\,(\varphi_{66}) \to \varphi_{53}) \to \varphi_{221}$\lbl{A1}}
\BIC{$\varphi_{221}$}
\end{ptree}
\noindent$\Pi_{16}$:
\begin{ptree}
\AXC{$\stackrel{\Pi_{7}}{\vdots}$}\noLine
\UIC{$\forall x\,\forall y\,(x = y \to y = x)$}
\AXC{$\varphi_{27}$\lbl{A2}}
\BIC{$\forall y\,({{r}^{S(S(n))}} = y \to y = {{r}^{S(S(n))}})$}
\AXC{$\varphi_{150}$\lbl{A2}}
\BIC{${{r}^{S(S(n))}} = {{r}^{S(n)}} \cdot r \to {{r}^{S(n)}} \cdot r = {{r}^{S(S(n))}}$}
\AXC{$\stackrel{\Pi_{15}}{\vdots}$}\noLine
\UIC{$\varphi_{221}$}
\BIC{$r \in \mathbb{R} \to \forall n\,(\varphi_{66}) \to \varphi_{54}$}
\AXC{$(r \in \mathbb{R} \to \forall n\,(\varphi_{66}) \to \varphi_{54}) \to \varphi_{250}$\lbl{A1}}
\BIC{$\varphi_{250}$}
\end{ptree}
\noindent$\Pi_{17}$:
\begin{ptree}
\AXC{$\stackrel{\Pi_{6}}{\vdots}$}\noLine
\UIC{${{r}^{S(n)}} \cdot r = {{r}^{S(S(n))}} \to \varphi_{110}$}
\AXC{$x = x$\lbl{A5}}
\UIC{$\forall x\,(x = x)$}
\AXC{$\forall x\,(x = x) \to r - {{r}^{S(n)}} \cdot r = r - {{r}^{S(n)}} \cdot r$\lbl{A2}}
\BIC{$r - {{r}^{S(n)}} \cdot r = r - {{r}^{S(n)}} \cdot r$}
\AXC{$\stackrel{\Pi_{16}}{\vdots}$}\noLine
\UIC{$\varphi_{250}$}
\BIC{$({{r}^{S(n)}} \cdot r = {{r}^{S(S(n))}} \to \varphi_{110}) \to r \in \mathbb{R} \to \forall n\,(\varphi_{66}) \to \varphi_{77}$}
\BIC{$r \in \mathbb{R} \to \forall n\,(\varphi_{66}) \to \varphi_{77}$}
\end{ptree}
\noindent$\Pi_{18}$:
\begin{ptree}
\AXC{$y = z \to x = y \to x = z$\lbl{A6}}
\AXC{$(y = z \to x = y \to x = z) \to x = y \to y = z \to x = z$\lbl{A1}}
\BIC{$x = y \to y = z \to x = z$}
\UIC{$\forall z\,(x = y \to y = z \to x = z)$}
\UIC{$\forall y\,\forall z\,(x = y \to y = z \to x = z)$}
\UIC{$\forall x\,\forall y\,\forall z\,(x = y \to y = z \to x = z)$}
\AXC{$\forall x\,\forall y\,\forall z\,(x = y \to y = z \to x = z) \to \varphi_{56}$\lbl{A2}}
\BIC{$\varphi_{56}$}
\end{ptree}
\noindent$\Pi_{19}$:
\begin{ptree}
\AXC{$\stackrel{\Pi_{7}}{\vdots}$}\noLine
\UIC{$\forall x\,\forall y\,(x = y \to y = x)$}
\AXC{$\varphi_{47}$\lbl{A2}}
\BIC{$\forall y\,(r - {{r}^{S(n)}} \cdot r = y \to y = r - {{r}^{S(n)}} \cdot r)$}
\AXC{$\forall y\,(r - {{r}^{S(n)}} \cdot r = y \to y = r - {{r}^{S(n)}} \cdot r) \to \varphi_{197}$\lbl{A2}}
\BIC{$\varphi_{197}$}
\end{ptree}
\noindent$\Pi_{20}$:
\begin{ptree}
\AXC{$\varphi_{179}$\lbl{A1}}
\AXC{$\forall r \in \mathbb{R}\,(\varphi_{58})$\lbl{[S24]}}
\AXC{$\forall r \in \mathbb{R}\,(\varphi_{58}) \to r \in \mathbb{R} \to \varphi_{58}$\lbl{A2}}
\BIC{$r \in \mathbb{R} \to \varphi_{58}$}
\AXC{$(r \in \mathbb{R} \to \varphi_{58}) \to \varphi_{246}$\lbl{A1}}
\BIC{$\varphi_{246}$}
\BIC{$r \in \mathbb{R} \to \forall n\,(\varphi_{66}) \to n \in B_{51} \to S(n) \in \mathbb{N} \to n \neq \emptyset \to \varphi_{58}$}
\AXC{$\varphi_{247}$\lbl{A1}}
\BIC{$r \in \mathbb{R} \to \forall n\,(\varphi_{66}) \to \varphi_{5}$}
\end{ptree}
\noindent$\Pi_{21}$:
\begin{ptree}
\AXC{$\stackrel{\Pi_{20}}{\vdots}$}\noLine
\UIC{$r \in \mathbb{R} \to \forall n\,(\varphi_{66}) \to \varphi_{5}$}
\AXC{$\forall a\,\forall n\,(\mathrm{Seq}(a) \to n \in \mathbb{N} \to a_{n} \in \mathbb{R})$\lbl{[S1]}}
\AXC{$\varphi_{89}$\lbl{A2}}
\BIC{$\forall n\,(\mathrm{Seq}({(r)^{\bullet}}) \to n \in \mathbb{N} \to {{r}^{n}} \in \mathbb{R})$}
\AXC{$\varphi_{115}$\lbl{A2}}
\BIC{$\mathrm{Seq}({(r)^{\bullet}}) \to S(n) \in \mathbb{N} \to {{r}^{S(n)}} \in \mathbb{R}$}
\AXC{$\varphi_{269}$\lbl{A1}}
\BIC{$(r \in \mathbb{R} \to \forall n\,(\varphi_{66}) \to \varphi_{5}) \to (r \in \mathbb{R} \to \varphi_{167}) \to r \in \mathbb{R} \to \varphi_{181}$}
\BIC{$(r \in \mathbb{R} \to \varphi_{167}) \to r \in \mathbb{R} \to \varphi_{181}$}
\end{ptree}
\noindent$\Pi_{22}$:
\begin{ptree}
\AXC{$\stackrel{\Pi_{13}}{\vdots}$}\noLine
\UIC{$r \in \mathbb{R} \to \varphi_{167}$}
\AXC{$\stackrel{\Pi_{21}}{\vdots}$}\noLine
\UIC{$(r \in \mathbb{R} \to \varphi_{167}) \to r \in \mathbb{R} \to \varphi_{181}$}
\BIC{$r \in \mathbb{R} \to \varphi_{181}$}
\end{ptree}
\noindent$\Pi_{23}$:
\begin{ptree}
\AXC{$\varphi_{179}$\lbl{A1}}
\AXC{$\stackrel{\Pi_{22}}{\vdots}$}\noLine
\UIC{$r \in \mathbb{R} \to \varphi_{181}$}
\AXC{$\forall x\,\forall y\,(x \in \mathbb{R} \to y \in \mathbb{R} \to x \cdot y \in \mathbb{R})$\lbl{(R1)}}
\AXC{$\varphi_{140}$\lbl{A2}}
\BIC{$\forall y\,({{r}^{S(n)}} \in \mathbb{R} \to y \in \mathbb{R} \to {{r}^{S(n)}} \cdot y \in \mathbb{R})$}
\AXC{$\varphi_{165}$\lbl{A2}}
\BIC{${{r}^{S(n)}} \in \mathbb{R} \to r \in \mathbb{R} \to {{r}^{S(n)}} \cdot r \in \mathbb{R}$}
\AXC{$\varphi_{278}$\lbl{A1}}
\BIC{$(r \in \mathbb{R} \to \varphi_{181}) \to \varphi_{179} \to r \in \mathbb{R} \to \varphi_{185}$}
\BIC{$\varphi_{179} \to r \in \mathbb{R} \to \varphi_{185}$}
\BIC{$r \in \mathbb{R} \to \varphi_{185}$}
\end{ptree}
\noindent$\Pi_{24}$:
\begin{ptree}
\AXC{$\varphi_{179}$\lbl{A1}}
\AXC{$\varphi_{35}$\lbl{[F16]}}
\AXC{$\varphi_{35} \to \varphi_{32}$\lbl{A2}}
\BIC{$\varphi_{32}$}
\AXC{$\varphi_{32} \to \varphi_{141}$\lbl{A2}}
\BIC{$\varphi_{141}$}
\AXC{$\varphi_{141} \to r \in \mathbb{R} \to {{r}^{S(n)}} \cdot r \in \mathbb{R} \to \varphi_{133}$\lbl{A2}}
\BIC{$r \in \mathbb{R} \to {{r}^{S(n)}} \cdot r \in \mathbb{R} \to \varphi_{133}$}
\AXC{$\varphi_{299}$\lbl{A1}}
\BIC{$\varphi_{179} \to (r \in \mathbb{R} \to \varphi_{185}) \to \varphi_{252}$}
\BIC{$(r \in \mathbb{R} \to \varphi_{185}) \to \varphi_{252}$}
\end{ptree}
\noindent$\Pi_{25}$:
\begin{ptree}
\AXC{$\stackrel{\Pi_{19}}{\vdots}$}\noLine
\UIC{$\varphi_{197}$}
\AXC{$\stackrel{\Pi_{22}}{\vdots}$}\noLine
\UIC{$r \in \mathbb{R} \to \varphi_{181}$}
\AXC{$\stackrel{\Pi_{23}}{\vdots}$}\noLine
\UIC{$r \in \mathbb{R} \to \varphi_{185}$}
\AXC{$\stackrel{\Pi_{24}}{\vdots}$}\noLine
\UIC{$(r \in \mathbb{R} \to \varphi_{185}) \to \varphi_{252}$}
\BIC{$\varphi_{252}$}
\BIC{$r \in \mathbb{R} \to \forall n\,(\varphi_{66}) \to n \in B_{51} \to S(n) \in \mathbb{N} \to \varphi_{111}$}
\AXC{$\varphi_{288}$\lbl{A1}}
\BIC{$\varphi_{197} \to r \in \mathbb{R} \to \forall n\,(\varphi_{66}) \to n \in B_{51} \to S(n) \in \mathbb{N} \to \varphi_{112}$}
\BIC{$r \in \mathbb{R} \to \forall n\,(\varphi_{66}) \to n \in B_{51} \to S(n) \in \mathbb{N} \to \varphi_{112}$}
\end{ptree}
\noindent$\Pi_{26}$:
\begin{ptree}
\AXC{$y = z \to x = y \to x = z$\lbl{A6}}
\AXC{$(y = z \to x = y \to x = z) \to x = y \to y = z \to x = z$\lbl{A1}}
\BIC{$x = y \to y = z \to x = z$}
\UIC{$\forall z\,(x = y \to y = z \to x = z)$}
\UIC{$\forall y\,\forall z\,(x = y \to y = z \to x = z)$}
\UIC{$\forall x\,\forall y\,\forall z\,(x = y \to y = z \to x = z)$}
\AXC{$\forall x\,\forall y\,\forall z\,(x = y \to y = z \to x = z) \to \varphi_{56}$\lbl{A2}}
\BIC{$\varphi_{56}$}
\end{ptree}
\noindent$\Pi_{27}$:
\begin{ptree}
\AXC{$z = w \to \varphi_{135} \to \varphi_{134}$\lbl{A6}}
\UIC{$\forall w\,(z = w \to \varphi_{135} \to \varphi_{134})$}
\UIC{$\forall z\,\forall w\,(z = w \to \varphi_{135} \to \varphi_{134})$}
\AXC{$\forall z\,\forall w\,(z = w \to \varphi_{135} \to \varphi_{134}) \to \forall w\,({{r}^{S(n)}} \cdot 1 = w \to \varphi_{156} \to \varphi_{134})$\lbl{A2}}
\BIC{$\forall w\,({{r}^{S(n)}} \cdot 1 = w \to \varphi_{156} \to \varphi_{134})$}
\AXC{$\varphi_{283}$\lbl{A2}}
\BIC{${{r}^{S(n)}} \cdot 1 = {{r}^{S(n)}} \to \varphi_{156} \to \varphi_{152}$}
\end{ptree}
\noindent$\Pi_{28}$:
\begin{ptree}
\AXC{$\stackrel{\Pi_{22}}{\vdots}$}\noLine
\UIC{$r \in \mathbb{R} \to \varphi_{181}$}
\AXC{$\forall x \in \mathbb{R}\,(x \cdot 1 = x)$\lbl{(R8)}}
\AXC{$\varphi_{76}$\lbl{A2}}
\BIC{${{r}^{S(n)}} \in \mathbb{R} \to {{r}^{S(n)}} \cdot 1 = {{r}^{S(n)}}$}
\AXC{$\varphi_{254}$\lbl{A1}}
\BIC{$(r \in \mathbb{R} \to \varphi_{181}) \to r \in \mathbb{R} \to \varphi_{190}$}
\BIC{$r \in \mathbb{R} \to \varphi_{190}$}
\AXC{$(r \in \mathbb{R} \to \varphi_{190}) \to \varphi_{156} \to \varphi_{281}$\lbl{A1}}
\BIC{$\varphi_{156} \to \varphi_{281}$}
\end{ptree}
\noindent$\Pi_{29}$:
\begin{ptree}
\AXC{$y = z \to x = y \to x = z$\lbl{A6}}
\AXC{$(y = z \to x = y \to x = z) \to x = y \to y = z \to x = z$\lbl{A1}}
\BIC{$x = y \to y = z \to x = z$}
\UIC{$\forall z\,(x = y \to y = z \to x = z)$}
\UIC{$\forall y\,\forall z\,(x = y \to y = z \to x = z)$}
\UIC{$\forall x\,\forall y\,\forall z\,(x = y \to y = z \to x = z)$}
\AXC{$\forall x\,\forall y\,\forall z\,(x = y \to y = z \to x = z) \to \varphi_{56}$\lbl{A2}}
\BIC{$\varphi_{56}$}
\end{ptree}
\noindent$\Pi_{30}$:
\begin{ptree}
\AXC{$z = w \to \varphi_{173}$\lbl{A6}}
\UIC{$\forall w\,(z = w \to \varphi_{173})$}
\UIC{$\forall z\,\forall w\,(z = w \to \varphi_{173})$}
\AXC{$\forall z\,\forall w\,(z = w \to \varphi_{173}) \to \varphi_{206}$\lbl{A2}}
\BIC{$\varphi_{206}$}
\AXC{$\varphi_{206} \to {{r}^{S(n)}} \cdot (1 - r) = {{r}^{S(n)}} \cdot 1 - {{r}^{S(n)}} \cdot r \to \varphi_{98} \to \varphi_{136}$\lbl{A2}}
\BIC{${{r}^{S(n)}} \cdot (1 - r) = {{r}^{S(n)}} \cdot 1 - {{r}^{S(n)}} \cdot r \to \varphi_{98} \to \varphi_{136}$}
\end{ptree}
\noindent$\Pi_{31}$:
\begin{ptree}
\AXC{$\stackrel{\Pi_{22}}{\vdots}$}\noLine
\UIC{$r \in \mathbb{R} \to \varphi_{181}$}
\AXC{$\varphi_{42}$\lbl{[F17]}}
\AXC{$\varphi_{42} \to \varphi_{149}$\lbl{A2}}
\BIC{$\varphi_{149}$}
\AXC{$\varphi_{149} \to \varphi_{146}$\lbl{A2}}
\BIC{$\varphi_{146}$}
\AXC{$\varphi_{146} \to {{r}^{S(n)}} \in \mathbb{R} \to \varphi_{91}$\lbl{A2}}
\BIC{${{r}^{S(n)}} \in \mathbb{R} \to \varphi_{91}$}
\AXC{$({{r}^{S(n)}} \in \mathbb{R} \to \varphi_{91}) \to \varphi_{275}$\lbl{A1}}
\BIC{$\varphi_{275}$}
\BIC{$1 \in \mathbb{R} \to \varphi_{179} \to r \in \mathbb{R} \to \forall n\,(\varphi_{66}) \to n \in B_{51} \to \varphi_{100}$}
\end{ptree}
\noindent$\Pi_{32}$:
\begin{ptree}
\AXC{$\varphi_{179}$\lbl{A1}}
\AXC{$1 \in \mathbb{R}$\lbl{(R1)}}
\AXC{$\stackrel{\Pi_{31}}{\vdots}$}\noLine
\UIC{$1 \in \mathbb{R} \to \varphi_{179} \to r \in \mathbb{R} \to \forall n\,(\varphi_{66}) \to n \in B_{51} \to \varphi_{100}$}
\BIC{$\varphi_{179} \to r \in \mathbb{R} \to \forall n\,(\varphi_{66}) \to n \in B_{51} \to \varphi_{100}$}
\BIC{$r \in \mathbb{R} \to \forall n\,(\varphi_{66}) \to n \in B_{51} \to \varphi_{100}$}
\AXC{$\varphi_{279}$\lbl{A1}}
\BIC{$\varphi_{164} \to r \in \mathbb{R} \to \forall n\,(\varphi_{66}) \to n \in B_{51} \to \varphi_{101}$}
\end{ptree}
\noindent$\Pi_{33}$:
\begin{ptree}
\AXC{$\stackrel{\Pi_{7}}{\vdots}$}\noLine
\UIC{$\forall x\,\forall y\,(x = y \to y = x)$}
\AXC{$\forall x\,\forall y\,(x = y \to y = x) \to \varphi_{93}$\lbl{A2}}
\BIC{$\varphi_{93}$}
\AXC{$\varphi_{93} \to \varphi_{164}$\lbl{A2}}
\BIC{$\varphi_{164}$}
\AXC{$\stackrel{\Pi_{32}}{\vdots}$}\noLine
\UIC{$\varphi_{164} \to r \in \mathbb{R} \to \forall n\,(\varphi_{66}) \to n \in B_{51} \to \varphi_{101}$}
\BIC{$r \in \mathbb{R} \to \forall n\,(\varphi_{66}) \to n \in B_{51} \to \varphi_{101}$}
\end{ptree}
\noindent$\Pi_{34}$:
\begin{ptree}
\AXC{$\stackrel{\Pi_{33}}{\vdots}$}\noLine
\UIC{$r \in \mathbb{R} \to \forall n\,(\varphi_{66}) \to n \in B_{51} \to \varphi_{101}$}
\AXC{$(r \in \mathbb{R} \to \forall n\,(\varphi_{66}) \to n \in B_{51} \to \varphi_{101}) \to \varphi_{98} \to \varphi_{274}$\lbl{A1}}
\BIC{$\varphi_{98} \to \varphi_{274}$}
\end{ptree}
\noindent$\Pi_{35}$:
\begin{ptree}
\AXC{$\stackrel{\Pi_{30}}{\vdots}$}\noLine
\UIC{${{r}^{S(n)}} \cdot (1 - r) = {{r}^{S(n)}} \cdot 1 - {{r}^{S(n)}} \cdot r \to \varphi_{98} \to \varphi_{136}$}
\AXC{$x = x$\lbl{A5}}
\UIC{$\forall x\,(x = x)$}
\AXC{$\forall x\,(x = x) \to \varphi_{98}$\lbl{A2}}
\BIC{$\varphi_{98}$}
\AXC{$\stackrel{\Pi_{34}}{\vdots}$}\noLine
\UIC{$\varphi_{98} \to \varphi_{274}$}
\BIC{$\varphi_{274}$}
\BIC{$r \in \mathbb{R} \to \forall n\,(\varphi_{66}) \to n \in B_{51} \to S(n) \in \mathbb{N} \to n \neq \emptyset \to \varphi_{136}$}
\end{ptree}
\noindent$\Pi_{36}$:
\begin{ptree}
\AXC{$y = z \to x = y \to x = z$\lbl{A6}}
\AXC{$(y = z \to x = y \to x = z) \to x = y \to y = z \to x = z$\lbl{A1}}
\BIC{$x = y \to y = z \to x = z$}
\UIC{$\forall z\,(x = y \to y = z \to x = z)$}
\UIC{$\forall y\,\forall z\,(x = y \to y = z \to x = z)$}
\UIC{$\forall x\,\forall y\,\forall z\,(x = y \to y = z \to x = z)$}
\AXC{$\forall x\,\forall y\,\forall z\,(x = y \to y = z \to x = z) \to \varphi_{56}$\lbl{A2}}
\BIC{$\varphi_{56}$}
\end{ptree}
\noindent$\Pi_{37}$:
\begin{ptree}
\AXC{$z = w \to \varphi_{174}$\lbl{A6}}
\UIC{$\forall w\,(z = w \to \varphi_{174})$}
\UIC{$\forall z\,\forall w\,(z = w \to \varphi_{174})$}
\AXC{$\forall z\,\forall w\,(z = w \to \varphi_{174}) \to \varphi_{207}$\lbl{A2}}
\BIC{$\varphi_{207}$}
\AXC{$\varphi_{207} \to (1 - r) \cdot {{r}^{S(n)}} = {{r}^{S(n)}} \cdot (1 - r) \to \varphi_{99} \to \varphi_{102}$\lbl{A2}}
\BIC{$(1 - r) \cdot {{r}^{S(n)}} = {{r}^{S(n)}} \cdot (1 - r) \to \varphi_{99} \to \varphi_{102}$}
\end{ptree}
\noindent$\Pi_{38}$:
\begin{ptree}
\AXC{$\forall x\,\forall y\,(x \in \mathbb{R} \to y \in \mathbb{R} \to x - y \in \mathbb{R})$\lbl{[F1]}}
\AXC{$\varphi_{84}$\lbl{A2}}
\BIC{$\forall y\,(1 \in \mathbb{R} \to y \in \mathbb{R} \to 1 - y \in \mathbb{R})$}
\AXC{$\varphi_{68}$\lbl{A2}}
\BIC{$1 \in \mathbb{R} \to r \in \mathbb{R} \to 1 - r \in \mathbb{R}$}
\AXC{$(1 \in \mathbb{R} \to r \in \mathbb{R} \to 1 - r \in \mathbb{R}) \to 1 \in \mathbb{R} \to \varphi_{179} \to \varphi_{183}$\lbl{A1}}
\BIC{$1 \in \mathbb{R} \to \varphi_{179} \to \varphi_{183}$}
\end{ptree}
\noindent$\Pi_{39}$:
\begin{ptree}
\AXC{$\varphi_{179}$\lbl{A1}}
\AXC{$1 \in \mathbb{R}$\lbl{(R1)}}
\AXC{$\stackrel{\Pi_{38}}{\vdots}$}\noLine
\UIC{$1 \in \mathbb{R} \to \varphi_{179} \to \varphi_{183}$}
\BIC{$\varphi_{179} \to \varphi_{183}$}
\BIC{$\varphi_{183}$}
\end{ptree}
\noindent$\Pi_{40}$:
\begin{ptree}
\AXC{$\forall x\,\forall y\,(x \in \mathbb{R} \to y \in \mathbb{R} \to x \cdot y = y \cdot x)$\lbl{(R7)}}
\AXC{$\forall x\,\forall y\,(x \in \mathbb{R} \to y \in \mathbb{R} \to x \cdot y = y \cdot x) \to \varphi_{13}$\lbl{A2}}
\BIC{$\varphi_{13}$}
\AXC{$\varphi_{13} \to \varphi_{90}$\lbl{A2}}
\BIC{$\varphi_{90}$}
\AXC{$\varphi_{285}$\lbl{A1}}
\BIC{$\varphi_{183} \to (r \in \mathbb{R} \to \varphi_{181}) \to r \in \mathbb{R} \to \forall n\,(\varphi_{66}) \to n \in B_{51} \to \varphi_{55}$}
\end{ptree}
\noindent$\Pi_{41}$:
\begin{ptree}
\AXC{$\stackrel{\Pi_{22}}{\vdots}$}\noLine
\UIC{$r \in \mathbb{R} \to \varphi_{181}$}
\AXC{$\stackrel{\Pi_{39}}{\vdots}$}\noLine
\UIC{$\varphi_{183}$}
\AXC{$\stackrel{\Pi_{40}}{\vdots}$}\noLine
\UIC{$\varphi_{183} \to (r \in \mathbb{R} \to \varphi_{181}) \to r \in \mathbb{R} \to \forall n\,(\varphi_{66}) \to n \in B_{51} \to \varphi_{55}$}
\BIC{$(r \in \mathbb{R} \to \varphi_{181}) \to r \in \mathbb{R} \to \forall n\,(\varphi_{66}) \to n \in B_{51} \to \varphi_{55}$}
\BIC{$r \in \mathbb{R} \to \forall n\,(\varphi_{66}) \to n \in B_{51} \to \varphi_{55}$}
\end{ptree}
\noindent$\Pi_{42}$:
\begin{ptree}
\AXC{$\stackrel{\Pi_{41}}{\vdots}$}\noLine
\UIC{$r \in \mathbb{R} \to \forall n\,(\varphi_{66}) \to n \in B_{51} \to \varphi_{55}$}
\AXC{$(r \in \mathbb{R} \to \forall n\,(\varphi_{66}) \to n \in B_{51} \to \varphi_{55}) \to \varphi_{99} \to \varphi_{263}$\lbl{A1}}
\BIC{$\varphi_{99} \to \varphi_{263}$}
\end{ptree}
\noindent$\Pi_{43}$:
\begin{ptree}
\AXC{$\stackrel{\Pi_{37}}{\vdots}$}\noLine
\UIC{$(1 - r) \cdot {{r}^{S(n)}} = {{r}^{S(n)}} \cdot (1 - r) \to \varphi_{99} \to \varphi_{102}$}
\AXC{$x = x$\lbl{A5}}
\UIC{$\forall x\,(x = x)$}
\AXC{$\forall x\,(x = x) \to \varphi_{99}$\lbl{A2}}
\BIC{$\varphi_{99}$}
\AXC{$\stackrel{\Pi_{42}}{\vdots}$}\noLine
\UIC{$\varphi_{99} \to \varphi_{263}$}
\BIC{$\varphi_{263}$}
\BIC{$r \in \mathbb{R} \to \forall n\,(\varphi_{66}) \to n \in B_{51} \to S(n) \in \mathbb{N} \to n \neq \emptyset \to \varphi_{102}$}
\end{ptree}
\noindent$\Pi_{44}$:
\begin{ptree}
\AXC{$y = z \to x = y \to x = z$\lbl{A6}}
\AXC{$(y = z \to x = y \to x = z) \to x = y \to y = z \to x = z$\lbl{A1}}
\BIC{$x = y \to y = z \to x = z$}
\UIC{$\forall z\,(x = y \to y = z \to x = z)$}
\UIC{$\forall y\,\forall z\,(x = y \to y = z \to x = z)$}
\UIC{$\forall x\,\forall y\,\forall z\,(x = y \to y = z \to x = z)$}
\AXC{$\forall x\,\forall y\,\forall z\,(x = y \to y = z \to x = z) \to \varphi_{56}$\lbl{A2}}
\BIC{$\varphi_{56}$}
\end{ptree}
\noindent$\Pi_{45}$:
\begin{ptree}
\AXC{$\stackrel{\Pi_{9}}{\vdots}$}\noLine
\UIC{$r \in \mathbb{R} \to \forall n\,(\varphi_{66}) \to \varphi_{67}$}
\AXC{$(r \in \mathbb{R} \to \forall n\,(\varphi_{66}) \to \varphi_{67}) \to r \in \mathbb{R} \to \varphi_{187}$\lbl{A1}}
\BIC{$r \in \mathbb{R} \to \varphi_{187}$}
\AXC{$(r \in \mathbb{R} \to \varphi_{187}) \to \varphi_{177} \to r \in \mathbb{R} \to \forall n\,(\varphi_{66}) \to \varphi_{78}$\lbl{A1}}
\BIC{$\varphi_{177} \to r \in \mathbb{R} \to \forall n\,(\varphi_{66}) \to \varphi_{78}$}
\end{ptree}
\noindent$\Pi_{46}$:
\begin{ptree}
\AXC{$x = x$\lbl{A5}}
\UIC{$\forall x\,(x = x)$}
\AXC{$\forall x\,(x = x) \to \varphi_{132}$\lbl{A2}}
\BIC{$\varphi_{132}$}
\AXC{$\stackrel{\Pi_{11}}{\vdots}$}\noLine
\UIC{$\varphi_{177}$}
\AXC{$\stackrel{\Pi_{45}}{\vdots}$}\noLine
\UIC{$\varphi_{177} \to r \in \mathbb{R} \to \forall n\,(\varphi_{66}) \to \varphi_{78}$}
\BIC{$r \in \mathbb{R} \to \forall n\,(\varphi_{66}) \to \varphi_{78}$}
\AXC{$(r \in \mathbb{R} \to \forall n\,(\varphi_{66}) \to \varphi_{78}) \to \varphi_{132} \to \varphi_{268}$\lbl{A1}}
\BIC{$\varphi_{132} \to \varphi_{268}$}
\BIC{$\varphi_{268}$}
\end{ptree}
\noindent$\Pi_{47}$:
\begin{ptree}
\AXC{$y = z \to x = y \to x = z$\lbl{A6}}
\AXC{$(y = z \to x = y \to x = z) \to x = y \to y = z \to x = z$\lbl{A1}}
\BIC{$x = y \to y = z \to x = z$}
\UIC{$\forall z\,(x = y \to y = z \to x = z)$}
\UIC{$\forall y\,\forall z\,(x = y \to y = z \to x = z)$}
\UIC{$\forall x\,\forall y\,\forall z\,(x = y \to y = z \to x = z)$}
\AXC{$\forall x\,\forall y\,\forall z\,(x = y \to y = z \to x = z) \to \varphi_{56}$\lbl{A2}}
\BIC{$\varphi_{56}$}
\end{ptree}
\noindent$\Pi_{48}$:
\begin{ptree}
\AXC{$\varphi_{179}$\lbl{A1}}
\AXC{$\forall r \in \mathbb{R}\,(\varphi_{58})$\lbl{[S24]}}
\AXC{$\forall r \in \mathbb{R}\,(\varphi_{58}) \to r \in \mathbb{R} \to \varphi_{58}$\lbl{A2}}
\BIC{$r \in \mathbb{R} \to \varphi_{58}$}
\AXC{$(r \in \mathbb{R} \to \varphi_{58}) \to \varphi_{246}$\lbl{A1}}
\BIC{$\varphi_{246}$}
\BIC{$r \in \mathbb{R} \to \forall n\,(\varphi_{66}) \to n \in B_{51} \to S(n) \in \mathbb{N} \to n \neq \emptyset \to \varphi_{58}$}
\AXC{$\varphi_{247}$\lbl{A1}}
\BIC{$r \in \mathbb{R} \to \forall n\,(\varphi_{66}) \to \varphi_{5}$}
\end{ptree}
\noindent$\Pi_{49}$:
\begin{ptree}
\AXC{$\stackrel{\Pi_{48}}{\vdots}$}\noLine
\UIC{$r \in \mathbb{R} \to \forall n\,(\varphi_{66}) \to \varphi_{5}$}
\AXC{$\varphi_{139}$\lbl{[S25]}}
\AXC{$\varphi_{139} \to \mathrm{Seq}({(r)^{\bullet}}) \to \varphi_{153}$\lbl{A2}}
\BIC{$\mathrm{Seq}({(r)^{\bullet}}) \to \varphi_{153}$}
\AXC{$(\mathrm{Seq}({(r)^{\bullet}}) \to \varphi_{153}) \to \varphi_{259}$\lbl{A1}}
\BIC{$\varphi_{259}$}
\BIC{$r \in \mathbb{R} \to \forall n\,(\varphi_{66}) \to n \in B_{51} \to S(n) \in \mathbb{N} \to n \neq \emptyset \to \varphi_{153}$}
\end{ptree}
\noindent$\Pi_{50}$:
\begin{ptree}
\AXC{$\forall a\,\forall n\,(\mathrm{Seq}(a) \to n \in \mathbb{N} \to a_{n} \in \mathbb{R})$\lbl{[S1]}}
\AXC{$\forall a\,\forall n\,(\mathrm{Seq}(a) \to n \in \mathbb{N} \to a_{n} \in \mathbb{R}) \to \varphi_{25}$\lbl{A2}}
\BIC{$\varphi_{25}$}
\end{ptree}
\noindent$\Pi_{51}$:
\begin{ptree}
\AXC{$\stackrel{\Pi_{50}}{\vdots}$}\noLine
\UIC{$\varphi_{25}$}
\AXC{$\varphi_{25} \to \mathrm{Seq}(\Sigma {(r)^{\bullet}}) \to n \in \mathbb{N} \to {\textstyle\sum_{k=1}^{n} {{r}^{k}}} \in \mathbb{R}$\lbl{A2}}
\BIC{$\mathrm{Seq}(\Sigma {(r)^{\bullet}}) \to n \in \mathbb{N} \to {\textstyle\sum_{k=1}^{n} {{r}^{k}}} \in \mathbb{R}$}
\AXC{$\varphi_{273}$\lbl{A1}}
\BIC{$(r \in \mathbb{R} \to \forall n\,(\varphi_{66}) \to \varphi_{15}) \to \varphi_{177} \to r \in \mathbb{R} \to \varphi_{182}$}
\end{ptree}
\noindent$\Pi_{52}$:
\begin{ptree}
\AXC{$\stackrel{\Pi_{11}}{\vdots}$}\noLine
\UIC{$\varphi_{177}$}
\AXC{$\stackrel{\Pi_{49}}{\vdots}$}\noLine
\UIC{$r \in \mathbb{R} \to \forall n\,(\varphi_{66}) \to n \in B_{51} \to S(n) \in \mathbb{N} \to n \neq \emptyset \to \varphi_{153}$}
\AXC{$\varphi_{260}$\lbl{A1}}
\BIC{$r \in \mathbb{R} \to \forall n\,(\varphi_{66}) \to \varphi_{15}$}
\AXC{$\stackrel{\Pi_{51}}{\vdots}$}\noLine
\UIC{$(r \in \mathbb{R} \to \forall n\,(\varphi_{66}) \to \varphi_{15}) \to \varphi_{177} \to r \in \mathbb{R} \to \varphi_{182}$}
\BIC{$\varphi_{177} \to r \in \mathbb{R} \to \varphi_{182}$}
\BIC{$r \in \mathbb{R} \to \varphi_{182}$}
\end{ptree}
\noindent$\Pi_{53}$:
\begin{ptree}
\AXC{$\stackrel{\Pi_{39}}{\vdots}$}\noLine
\UIC{$\varphi_{183}$}
\AXC{$\varphi_{43}$\lbl{(R10)}}
\AXC{$\varphi_{43} \to \varphi_{82}$\lbl{A2}}
\BIC{$\varphi_{82}$}
\AXC{$\varphi_{82} \to \varphi_{161}$\lbl{A2}}
\BIC{$\varphi_{161}$}
\AXC{$\varphi_{161} \to 1 - r \in \mathbb{R} \to {\textstyle\sum_{k=1}^{n} {{r}^{k}}} \in \mathbb{R} \to {{r}^{S(n)}} \in \mathbb{R} \to \varphi_{118}$\lbl{A2}}
\BIC{$1 - r \in \mathbb{R} \to {\textstyle\sum_{k=1}^{n} {{r}^{k}}} \in \mathbb{R} \to {{r}^{S(n)}} \in \mathbb{R} \to \varphi_{118}$}
\AXC{$\varphi_{301}$\lbl{A1}}
\BIC{$\varphi_{183} \to (r \in \mathbb{R} \to \varphi_{182}) \to \varphi_{256}$}
\BIC{$(r \in \mathbb{R} \to \varphi_{182}) \to \varphi_{256}$}
\end{ptree}
\noindent$\Pi_{54}$:
\begin{ptree}
\AXC{$\varphi_{179}$\lbl{A1}}
\AXC{$\forall r \in \mathbb{R}\,(\varphi_{58})$\lbl{[S24]}}
\AXC{$\forall r \in \mathbb{R}\,(\varphi_{58}) \to r \in \mathbb{R} \to \varphi_{58}$\lbl{A2}}
\BIC{$r \in \mathbb{R} \to \varphi_{58}$}
\AXC{$(r \in \mathbb{R} \to \varphi_{58}) \to \varphi_{246}$\lbl{A1}}
\BIC{$\varphi_{246}$}
\BIC{$r \in \mathbb{R} \to \forall n\,(\varphi_{66}) \to n \in B_{51} \to S(n) \in \mathbb{N} \to n \neq \emptyset \to \varphi_{58}$}
\AXC{$\varphi_{247}$\lbl{A1}}
\BIC{$r \in \mathbb{R} \to \forall n\,(\varphi_{66}) \to \varphi_{5}$}
\end{ptree}
\noindent$\Pi_{55}$:
\begin{ptree}
\AXC{$\stackrel{\Pi_{54}}{\vdots}$}\noLine
\UIC{$r \in \mathbb{R} \to \forall n\,(\varphi_{66}) \to \varphi_{5}$}
\AXC{$\varphi_{139}$\lbl{[S25]}}
\AXC{$\varphi_{139} \to \mathrm{Seq}({(r)^{\bullet}}) \to \varphi_{153}$\lbl{A2}}
\BIC{$\mathrm{Seq}({(r)^{\bullet}}) \to \varphi_{153}$}
\AXC{$(\mathrm{Seq}({(r)^{\bullet}}) \to \varphi_{153}) \to \varphi_{259}$\lbl{A1}}
\BIC{$\varphi_{259}$}
\BIC{$r \in \mathbb{R} \to \forall n\,(\varphi_{66}) \to n \in B_{51} \to S(n) \in \mathbb{N} \to n \neq \emptyset \to \varphi_{153}$}
\end{ptree}
\noindent$\Pi_{56}$:
\begin{ptree}
\AXC{$\stackrel{\Pi_{11}}{\vdots}$}\noLine
\UIC{$\varphi_{177}$}
\AXC{$\varphi_{169}$\lbl{A2}}
\AXC{$\stackrel{\Pi_{55}}{\vdots}$}\noLine
\UIC{$r \in \mathbb{R} \to \forall n\,(\varphi_{66}) \to n \in B_{51} \to S(n) \in \mathbb{N} \to n \neq \emptyset \to \varphi_{153}$}
\AXC{$\varphi_{264}$\lbl{A1}}
\BIC{$r \in \mathbb{R} \to \forall n\,(\varphi_{66}) \to n \in B_{51} \to \varphi_{113}$}
\AXC{$\varphi_{282}$\lbl{A1}}
\BIC{$\varphi_{169} \to r \in \mathbb{R} \to \forall n\,(\varphi_{66}) \to n \in B_{51} \to \varphi_{104}$}
\BIC{$r \in \mathbb{R} \to \forall n\,(\varphi_{66}) \to n \in B_{51} \to \varphi_{104}$}
\AXC{$\varphi_{280}$\lbl{A1}}
\BIC{$\varphi_{177} \to r \in \mathbb{R} \to \forall n\,(\varphi_{66}) \to \varphi_{105}$}
\BIC{$r \in \mathbb{R} \to \forall n\,(\varphi_{66}) \to \varphi_{105}$}
\end{ptree}
\noindent$\Pi_{57}$:
\begin{ptree}
\AXC{$x = x$\lbl{A5}}
\UIC{$\forall x\,(x = x)$}
\AXC{$\forall x\,(x = x) \to (1 - r) \cdot {\textstyle\sum_{k=1}^{S(n)} {{r}^{k}}} = (1 - r) \cdot {\textstyle\sum_{k=1}^{S(n)} {{r}^{k}}}$\lbl{A2}}
\BIC{$(1 - r) \cdot {\textstyle\sum_{k=1}^{S(n)} {{r}^{k}}} = (1 - r) \cdot {\textstyle\sum_{k=1}^{S(n)} {{r}^{k}}}$}
\AXC{$\stackrel{\Pi_{56}}{\vdots}$}\noLine
\UIC{$r \in \mathbb{R} \to \forall n\,(\varphi_{66}) \to \varphi_{105}$}
\AXC{$\varphi_{291}$\lbl{A1}}
\BIC{$(1 - r) \cdot {\textstyle\sum_{k=1}^{S(n)} {{r}^{k}}} = (1 - r) \cdot {\textstyle\sum_{k=1}^{S(n)} {{r}^{k}}} \to \varphi_{258}$}
\BIC{$\varphi_{258}$}
\end{ptree}
\noindent$\Pi_{58}$:
\begin{ptree}
\AXC{$z = w \to \varphi_{74}$\lbl{A6}}
\UIC{$\forall w\,(z = w \to \varphi_{74})$}
\UIC{$\forall z\,\forall w\,(z = w \to \varphi_{74})$}
\AXC{$\forall z\,\forall w\,(z = w \to \varphi_{74}) \to \varphi_{158}$\lbl{A2}}
\BIC{$\varphi_{158}$}
\AXC{$\varphi_{158} \to {\textstyle\sum_{k=1}^{S(n)} {{r}^{k}}} = {\textstyle\sum_{k=1}^{n} {{r}^{k}}} + {{r}^{S(n)}} \to \varphi_{166}$\lbl{A2}}
\BIC{${\textstyle\sum_{k=1}^{S(n)} {{r}^{k}}} = {\textstyle\sum_{k=1}^{n} {{r}^{k}}} + {{r}^{S(n)}} \to \varphi_{166}$}
\AXC{$\stackrel{\Pi_{57}}{\vdots}$}\noLine
\UIC{$\varphi_{258}$}
\BIC{$r \in \mathbb{R} \to \forall n\,(\varphi_{66}) \to n \in B_{51} \to S(n) \in \mathbb{N} \to \varphi_{83}$}
\end{ptree}
\noindent$\Pi_{59}$:
\begin{ptree}
\AXC{$\stackrel{\Pi_{58}}{\vdots}$}\noLine
\UIC{$r \in \mathbb{R} \to \forall n\,(\varphi_{66}) \to n \in B_{51} \to S(n) \in \mathbb{N} \to \varphi_{83}$}
\AXC{$(r \in \mathbb{R} \to \forall n\,(\varphi_{66}) \to n \in B_{51} \to S(n) \in \mathbb{N} \to \varphi_{83}) \to \varphi_{295}$\lbl{A1}}
\BIC{$\varphi_{295}$}
\end{ptree}
\noindent$\Pi_{60}$:
\begin{ptree}
\AXC{$\stackrel{\Pi_{47}}{\vdots}$}\noLine
\UIC{$\varphi_{56}$}
\AXC{$\varphi_{56} \to \varphi_{189}$\lbl{A2}}
\BIC{$\varphi_{189}$}
\AXC{$\varphi_{189} \to \varphi_{224}$\lbl{A2}}
\BIC{$\varphi_{224}$}
\AXC{$\stackrel{\Pi_{22}}{\vdots}$}\noLine
\UIC{$r \in \mathbb{R} \to \varphi_{181}$}
\AXC{$\stackrel{\Pi_{52}}{\vdots}$}\noLine
\UIC{$r \in \mathbb{R} \to \varphi_{182}$}
\AXC{$\stackrel{\Pi_{53}}{\vdots}$}\noLine
\UIC{$(r \in \mathbb{R} \to \varphi_{182}) \to \varphi_{256}$}
\BIC{$\varphi_{256}$}
\BIC{$r \in \mathbb{R} \to \forall n\,(\varphi_{66}) \to n \in B_{51} \to S(n) \in \mathbb{N} \to n \neq \emptyset \to \varphi_{118}$}
\AXC{$\stackrel{\Pi_{59}}{\vdots}$}\noLine
\UIC{$\varphi_{295}$}
\BIC{$\varphi_{266}$}
\BIC{$r \in \mathbb{R} \to \forall n\,(\varphi_{66}) \to n \in B_{51} \to S(n) \in \mathbb{N} \to n \neq \emptyset \to \varphi_{73}$}
\end{ptree}
\noindent$\Pi_{61}$:
\begin{ptree}
\AXC{$\stackrel{\Pi_{60}}{\vdots}$}\noLine
\UIC{$r \in \mathbb{R} \to \forall n\,(\varphi_{66}) \to n \in B_{51} \to S(n) \in \mathbb{N} \to n \neq \emptyset \to \varphi_{73}$}
\AXC{$\varphi_{303}$\lbl{A1}}
\BIC{$\varphi_{294}$}
\end{ptree}
\noindent$\Pi_{62}$:
\begin{ptree}
\AXC{$\stackrel{\Pi_{44}}{\vdots}$}\noLine
\UIC{$\varphi_{56}$}
\AXC{$\varphi_{56} \to \varphi_{202}$\lbl{A2}}
\BIC{$\varphi_{202}$}
\AXC{$\varphi_{202} \to \varphi_{225}$\lbl{A2}}
\BIC{$\varphi_{225}$}
\AXC{$z = w \to \varphi_{71} \to \varphi_{70}$\lbl{A6}}
\UIC{$\forall w\,(z = w \to \varphi_{71} \to \varphi_{70})$}
\UIC{$\forall z\,\forall w\,(z = w \to \varphi_{71} \to \varphi_{70})$}
\AXC{$\varphi_{257}$\lbl{A2}}
\BIC{$\forall w\,((1 - r) \cdot {\textstyle\sum_{k=1}^{n} {{r}^{k}}} = w \to \varphi_{132} \to \varphi_{70})$}
\AXC{$\varphi_{267}$\lbl{A2}}
\BIC{$(1 - r) \cdot {\textstyle\sum_{k=1}^{n} {{r}^{k}}} = r - {{r}^{S(n)}} \to \varphi_{132} \to \varphi_{121}$}
\AXC{$\stackrel{\Pi_{46}}{\vdots}$}\noLine
\UIC{$\varphi_{268}$}
\BIC{$r \in \mathbb{R} \to \forall n\,(\varphi_{66}) \to n \in B_{51} \to S(n) \in \mathbb{N} \to n \neq \emptyset \to \varphi_{121}$}
\AXC{$\stackrel{\Pi_{61}}{\vdots}$}\noLine
\UIC{$\varphi_{294}$}
\BIC{$\varphi_{225} \to r \in \mathbb{R} \to \forall n\,(\varphi_{66}) \to n \in B_{51} \to S(n) \in \mathbb{N} \to \varphi_{87}$}
\BIC{$r \in \mathbb{R} \to \forall n\,(\varphi_{66}) \to n \in B_{51} \to S(n) \in \mathbb{N} \to \varphi_{87}$}
\end{ptree}
\noindent$\Pi_{63}$:
\begin{ptree}
\AXC{$\stackrel{\Pi_{62}}{\vdots}$}\noLine
\UIC{$r \in \mathbb{R} \to \forall n\,(\varphi_{66}) \to n \in B_{51} \to S(n) \in \mathbb{N} \to \varphi_{87}$}
\AXC{$(r \in \mathbb{R} \to \forall n\,(\varphi_{66}) \to n \in B_{51} \to S(n) \in \mathbb{N} \to \varphi_{87}) \to \varphi_{292}$\lbl{A1}}
\BIC{$\varphi_{292}$}
\end{ptree}
\noindent$\Pi_{64}$:
\begin{ptree}
\AXC{$\stackrel{\Pi_{36}}{\vdots}$}\noLine
\UIC{$\varphi_{56}$}
\AXC{$\varphi_{56} \to \varphi_{193}$\lbl{A2}}
\BIC{$\varphi_{193}$}
\AXC{$\varphi_{193} \to \varphi_{220}$\lbl{A2}}
\BIC{$\varphi_{220}$}
\AXC{$\stackrel{\Pi_{43}}{\vdots}$}\noLine
\UIC{$r \in \mathbb{R} \to \forall n\,(\varphi_{66}) \to n \in B_{51} \to S(n) \in \mathbb{N} \to n \neq \emptyset \to \varphi_{102}$}
\AXC{$\stackrel{\Pi_{63}}{\vdots}$}\noLine
\UIC{$\varphi_{292}$}
\BIC{$\varphi_{220} \to r \in \mathbb{R} \to \forall n\,(\varphi_{66}) \to n \in B_{51} \to S(n) \in \mathbb{N} \to \varphi_{88}$}
\BIC{$r \in \mathbb{R} \to \forall n\,(\varphi_{66}) \to n \in B_{51} \to S(n) \in \mathbb{N} \to \varphi_{88}$}
\end{ptree}
\noindent$\Pi_{65}$:
\begin{ptree}
\AXC{$\stackrel{\Pi_{64}}{\vdots}$}\noLine
\UIC{$r \in \mathbb{R} \to \forall n\,(\varphi_{66}) \to n \in B_{51} \to S(n) \in \mathbb{N} \to \varphi_{88}$}
\AXC{$(r \in \mathbb{R} \to \forall n\,(\varphi_{66}) \to n \in B_{51} \to S(n) \in \mathbb{N} \to \varphi_{88}) \to \varphi_{297}$\lbl{A1}}
\BIC{$\varphi_{297}$}
\end{ptree}
\noindent$\Pi_{66}$:
\begin{ptree}
\AXC{$\stackrel{\Pi_{29}}{\vdots}$}\noLine
\UIC{$\varphi_{56}$}
\AXC{$\varphi_{56} \to \varphi_{192}$\lbl{A2}}
\BIC{$\varphi_{192}$}
\AXC{$\varphi_{192} \to \varphi_{227}$\lbl{A2}}
\BIC{$\varphi_{227}$}
\AXC{$\stackrel{\Pi_{35}}{\vdots}$}\noLine
\UIC{$r \in \mathbb{R} \to \forall n\,(\varphi_{66}) \to n \in B_{51} \to S(n) \in \mathbb{N} \to n \neq \emptyset \to \varphi_{136}$}
\AXC{$\stackrel{\Pi_{65}}{\vdots}$}\noLine
\UIC{$\varphi_{297}$}
\BIC{$\varphi_{271}$}
\BIC{$r \in \mathbb{R} \to \forall n\,(\varphi_{66}) \to n \in B_{51} \to S(n) \in \mathbb{N} \to n \neq \emptyset \to \varphi_{103}$}
\AXC{$\varphi_{304}$\lbl{A1}}
\BIC{$\varphi_{298}$}
\end{ptree}
\noindent$\Pi_{67}$:
\begin{ptree}
\AXC{$\stackrel{\Pi_{27}}{\vdots}$}\noLine
\UIC{${{r}^{S(n)}} \cdot 1 = {{r}^{S(n)}} \to \varphi_{156} \to \varphi_{152}$}
\AXC{$x = x$\lbl{A5}}
\UIC{$\forall x\,(x = x)$}
\AXC{$\forall x\,(x = x) \to \varphi_{156}$\lbl{A2}}
\BIC{$\varphi_{156}$}
\AXC{$\stackrel{\Pi_{28}}{\vdots}$}\noLine
\UIC{$\varphi_{156} \to \varphi_{281}$}
\BIC{$\varphi_{281}$}
\BIC{$r \in \mathbb{R} \to \forall n\,(\varphi_{66}) \to n \in B_{51} \to S(n) \in \mathbb{N} \to n \neq \emptyset \to \varphi_{152}$}
\AXC{$\stackrel{\Pi_{66}}{\vdots}$}\noLine
\UIC{$\varphi_{298}$}
\BIC{$\varphi_{277}$}
\end{ptree}
\noindent$\Pi_{68}$:
\begin{ptree}
\AXC{$\stackrel{\Pi_{25}}{\vdots}$}\noLine
\UIC{$r \in \mathbb{R} \to \forall n\,(\varphi_{66}) \to n \in B_{51} \to S(n) \in \mathbb{N} \to \varphi_{112}$}
\AXC{$\stackrel{\Pi_{26}}{\vdots}$}\noLine
\UIC{$\varphi_{56}$}
\AXC{$\varphi_{56} \to \varphi_{210}$\lbl{A2}}
\BIC{$\varphi_{210}$}
\AXC{$\varphi_{210} \to \varphi_{103} \to \varphi_{152} \to \varphi_{86}$\lbl{A2}}
\BIC{$\varphi_{103} \to \varphi_{152} \to \varphi_{86}$}
\AXC{$\stackrel{\Pi_{67}}{\vdots}$}\noLine
\UIC{$\varphi_{277}$}
\BIC{$r \in \mathbb{R} \to \forall n\,(\varphi_{66}) \to n \in B_{51} \to S(n) \in \mathbb{N} \to n \neq \emptyset \to \varphi_{86}$}
\AXC{$\varphi_{300}$\lbl{A1}}
\BIC{$\varphi_{290}$}
\BIC{$(\varphi_{86} \to \varphi_{170}) \to r \in \mathbb{R} \to \forall n\,(\varphi_{66}) \to n \in B_{51} \to \varphi_{65}$}
\end{ptree}
\noindent$\Pi_{69}$:
\begin{ptree}
\AXC{$\stackrel{\Pi_{18}}{\vdots}$}\noLine
\UIC{$\varphi_{56}$}
\AXC{$\varphi_{56} \to \varphi_{203}$\lbl{A2}}
\BIC{$\varphi_{203}$}
\AXC{$\varphi_{203} \to \varphi_{86} \to \varphi_{170}$\lbl{A2}}
\BIC{$\varphi_{86} \to \varphi_{170}$}
\AXC{$\stackrel{\Pi_{68}}{\vdots}$}\noLine
\UIC{$(\varphi_{86} \to \varphi_{170}) \to r \in \mathbb{R} \to \forall n\,(\varphi_{66}) \to n \in B_{51} \to \varphi_{65}$}
\BIC{$r \in \mathbb{R} \to \forall n\,(\varphi_{66}) \to n \in B_{51} \to \varphi_{65}$}
\AXC{$\varphi_{296}$\lbl{A1}}
\BIC{$(r \in \mathbb{R} \to \forall n\,(\varphi_{66}) \to \varphi_{77}) \to \varphi_{243}$}
\end{ptree}
\noindent$\Pi_{70}$:
\begin{ptree}
\AXC{$\stackrel{\Pi_{5}}{\vdots}$}\noLine
\UIC{$(1 - r) \cdot {\textstyle\sum_{k=1}^{S(n)} {{r}^{k}}} = r - {{r}^{S(n)}} \cdot r \to \varphi_{120}$}
\AXC{$\stackrel{\Pi_{17}}{\vdots}$}\noLine
\UIC{$r \in \mathbb{R} \to \forall n\,(\varphi_{66}) \to \varphi_{77}$}
\AXC{$\stackrel{\Pi_{69}}{\vdots}$}\noLine
\UIC{$(r \in \mathbb{R} \to \forall n\,(\varphi_{66}) \to \varphi_{77}) \to \varphi_{243}$}
\BIC{$\varphi_{243}$}
\BIC{$r \in \mathbb{R} \to \forall n\,(\varphi_{66}) \to n \in B_{51} \to \varphi_{64}$}
\end{ptree}
\noindent$\Pi_{71}$:
\begin{ptree}
\AXC{$y = z \to x = y \to x = z$\lbl{A6}}
\AXC{$(y = z \to x = y \to x = z) \to x = y \to y = z \to x = z$\lbl{A1}}
\BIC{$x = y \to y = z \to x = z$}
\UIC{$\forall z\,(x = y \to y = z \to x = z)$}
\UIC{$\forall y\,\forall z\,(x = y \to y = z \to x = z)$}
\UIC{$\forall x\,\forall y\,\forall z\,(x = y \to y = z \to x = z)$}
\AXC{$\forall x\,\forall y\,\forall z\,(x = y \to y = z \to x = z) \to \varphi_{57}$\lbl{A2}}
\BIC{$\varphi_{57}$}
\end{ptree}
\noindent$\Pi_{72}$:
\begin{ptree}
\AXC{$z = w \to r - {{r}^{1_{\omega}}} \cdot r = r - z \to r - {{r}^{1_{\omega}}} \cdot r = r - w$\lbl{A6}}
\UIC{$\varphi_{44}$}
\UIC{$\forall z\,(\varphi_{44})$}
\AXC{$\forall z\,(\varphi_{44}) \to \varphi_{130}$\lbl{A2}}
\BIC{$\varphi_{130}$}
\AXC{$\varphi_{130} \to {{r}^{1_{\omega}}} \cdot r = {{r}^{S(1_{\omega})}} \to \varphi_{114}$\lbl{A2}}
\BIC{${{r}^{1_{\omega}}} \cdot r = {{r}^{S(1_{\omega})}} \to \varphi_{114}$}
\end{ptree}
\noindent$\Pi_{73}$:
\begin{ptree}
\AXC{$\forall r \in \mathbb{R}\,(\varphi_{58})$\lbl{[S24]}}
\AXC{$\forall r \in \mathbb{R}\,(\varphi_{58}) \to r \in \mathbb{R} \to \varphi_{58}$\lbl{A2}}
\BIC{$r \in \mathbb{R} \to \varphi_{58}$}
\AXC{$(r \in \mathbb{R} \to \varphi_{58}) \to r \in \mathbb{R} \to \forall n \in \mathbb{N}\,({{r}^{S(n)}} = {{r}^{n}} \cdot r)$\lbl{A1}}
\BIC{$r \in \mathbb{R} \to \forall n \in \mathbb{N}\,({{r}^{S(n)}} = {{r}^{n}} \cdot r)$}
\end{ptree}
\noindent$\Pi_{74}$:
\begin{ptree}
\AXC{$1_{\omega} \in \mathbb{N}$\lbl{[N1]}}
\AXC{$\varphi_{117}$\lbl{A2}}
\AXC{$\stackrel{\Pi_{73}}{\vdots}$}\noLine
\UIC{$r \in \mathbb{R} \to \forall n \in \mathbb{N}\,({{r}^{S(n)}} = {{r}^{n}} \cdot r)$}
\AXC{$\varphi_{216}$\lbl{A1}}
\BIC{$\varphi_{117} \to r \in \mathbb{R} \to 1_{\omega} \in \mathbb{N} \to {{r}^{S(1_{\omega})}} = {{r}^{1_{\omega}}} \cdot r$}
\BIC{$r \in \mathbb{R} \to 1_{\omega} \in \mathbb{N} \to {{r}^{S(1_{\omega})}} = {{r}^{1_{\omega}}} \cdot r$}
\AXC{$\varphi_{157}$\lbl{A1}}
\BIC{$1_{\omega} \in \mathbb{N} \to r \in \mathbb{R} \to {{r}^{S(1_{\omega})}} = {{r}^{1_{\omega}}} \cdot r$}
\BIC{$r \in \mathbb{R} \to {{r}^{S(1_{\omega})}} = {{r}^{1_{\omega}}} \cdot r$}
\end{ptree}
\noindent$\Pi_{75}$:
\begin{ptree}
\AXC{$\stackrel{\Pi_{74}}{\vdots}$}\noLine
\UIC{$r \in \mathbb{R} \to {{r}^{S(1_{\omega})}} = {{r}^{1_{\omega}}} \cdot r$}
\AXC{$(r \in \mathbb{R} \to {{r}^{S(1_{\omega})}} = {{r}^{1_{\omega}}} \cdot r) \to \varphi_{154}$\lbl{A1}}
\BIC{$\varphi_{154}$}
\end{ptree}
\noindent$\Pi_{76}$:
\begin{ptree}
\AXC{$\stackrel{\Pi_{7}}{\vdots}$}\noLine
\UIC{$\forall x\,\forall y\,(x = y \to y = x)$}
\AXC{$\varphi_{28}$\lbl{A2}}
\BIC{$\forall y\,({{r}^{S(1_{\omega})}} = y \to y = {{r}^{S(1_{\omega})}})$}
\AXC{$\varphi_{151}$\lbl{A2}}
\BIC{${{r}^{S(1_{\omega})}} = {{r}^{1_{\omega}}} \cdot r \to {{r}^{1_{\omega}}} \cdot r = {{r}^{S(1_{\omega})}}$}
\AXC{$\stackrel{\Pi_{75}}{\vdots}$}\noLine
\UIC{$\varphi_{154}$}
\BIC{$r \in \mathbb{R} \to {{r}^{1_{\omega}}} \cdot r = {{r}^{S(1_{\omega})}}$}
\AXC{$\varphi_{229}$\lbl{A1}}
\BIC{$r - {{r}^{1_{\omega}}} \cdot r = r - {{r}^{1_{\omega}}} \cdot r \to \varphi_{208}$}
\end{ptree}
\noindent$\Pi_{77}$:
\begin{ptree}
\AXC{$x = x$\lbl{A5}}
\UIC{$\forall x\,(x = x)$}
\AXC{$\forall x\,(x = x) \to r - {{r}^{1_{\omega}}} \cdot r = r - {{r}^{1_{\omega}}} \cdot r$\lbl{A2}}
\BIC{$r - {{r}^{1_{\omega}}} \cdot r = r - {{r}^{1_{\omega}}} \cdot r$}
\AXC{$\stackrel{\Pi_{76}}{\vdots}$}\noLine
\UIC{$r - {{r}^{1_{\omega}}} \cdot r = r - {{r}^{1_{\omega}}} \cdot r \to \varphi_{208}$}
\BIC{$\varphi_{208}$}
\end{ptree}
\noindent$\Pi_{78}$:
\begin{ptree}
\AXC{$y = z \to x = y \to x = z$\lbl{A6}}
\AXC{$(y = z \to x = y \to x = z) \to x = y \to y = z \to x = z$\lbl{A1}}
\BIC{$x = y \to y = z \to x = z$}
\UIC{$\forall z\,(x = y \to y = z \to x = z)$}
\UIC{$\forall y\,\forall z\,(x = y \to y = z \to x = z)$}
\UIC{$\forall x\,\forall y\,\forall z\,(x = y \to y = z \to x = z)$}
\AXC{$\forall x\,\forall y\,\forall z\,(x = y \to y = z \to x = z) \to \varphi_{57}$\lbl{A2}}
\BIC{$\varphi_{57}$}
\end{ptree}
\noindent$\Pi_{79}$:
\begin{ptree}
\AXC{$z = w \to r - r \cdot r = r - z \cdot r \to r - r \cdot r = r - w \cdot r$\lbl{A6}}
\UIC{$\varphi_{16}$}
\UIC{$\forall z\,(\varphi_{16})$}
\AXC{$\forall z\,(\varphi_{16}) \to \varphi_{17}$\lbl{A2}}
\BIC{$\varphi_{17}$}
\AXC{$\varphi_{17} \to \varphi_{60}$\lbl{A2}}
\BIC{$\varphi_{60}$}
\end{ptree}
\noindent$\Pi_{80}$:
\begin{ptree}
\AXC{$\forall r \in \mathbb{R}\,(\varphi_{58})$\lbl{[S24]}}
\AXC{$\forall r \in \mathbb{R}\,(\varphi_{58}) \to r \in \mathbb{R} \to \varphi_{58}$\lbl{A2}}
\BIC{$r \in \mathbb{R} \to \varphi_{58}$}
\AXC{$(r \in \mathbb{R} \to \varphi_{58}) \to r \in \mathbb{R} \to {{r}^{1_{\omega}}} = r$\lbl{A1}}
\BIC{$r \in \mathbb{R} \to {{r}^{1_{\omega}}} = r$}
\AXC{$\varphi_{109}$\lbl{A1}}
\BIC{$({{r}^{1_{\omega}}} = r \to r = {{r}^{1_{\omega}}}) \to r \in \mathbb{R} \to r = {{r}^{1_{\omega}}}$}
\end{ptree}
\noindent$\Pi_{81}$:
\begin{ptree}
\AXC{$\stackrel{\Pi_{7}}{\vdots}$}\noLine
\UIC{$\forall x\,\forall y\,(x = y \to y = x)$}
\AXC{$\varphi_{18}$\lbl{A2}}
\BIC{$\forall y\,({{r}^{1_{\omega}}} = y \to y = {{r}^{1_{\omega}}})$}
\AXC{$\varphi_{51}$\lbl{A2}}
\BIC{${{r}^{1_{\omega}}} = r \to r = {{r}^{1_{\omega}}}$}
\AXC{$\stackrel{\Pi_{80}}{\vdots}$}\noLine
\UIC{$({{r}^{1_{\omega}}} = r \to r = {{r}^{1_{\omega}}}) \to r \in \mathbb{R} \to r = {{r}^{1_{\omega}}}$}
\BIC{$r \in \mathbb{R} \to r = {{r}^{1_{\omega}}}$}
\end{ptree}
\noindent$\Pi_{82}$:
\begin{ptree}
\AXC{$x = x$\lbl{A5}}
\UIC{$\forall x\,(x = x)$}
\AXC{$\forall x\,(x = x) \to r - r \cdot r = r - r \cdot r$\lbl{A2}}
\BIC{$r - r \cdot r = r - r \cdot r$}
\AXC{$\stackrel{\Pi_{81}}{\vdots}$}\noLine
\UIC{$r \in \mathbb{R} \to r = {{r}^{1_{\omega}}}$}
\AXC{$(r \in \mathbb{R} \to r = {{r}^{1_{\omega}}}) \to \varphi_{184}$\lbl{A1}}
\BIC{$\varphi_{184}$}
\BIC{$\varphi_{60} \to r \in \mathbb{R} \to r - r \cdot r = r - {{r}^{1_{\omega}}} \cdot r$}
\end{ptree}
\noindent$\Pi_{83}$:
\begin{ptree}
\AXC{$y = z \to x = y \to x = z$\lbl{A6}}
\AXC{$(y = z \to x = y \to x = z) \to x = y \to y = z \to x = z$\lbl{A1}}
\BIC{$x = y \to y = z \to x = z$}
\UIC{$\forall z\,(x = y \to y = z \to x = z)$}
\UIC{$\forall y\,\forall z\,(x = y \to y = z \to x = z)$}
\UIC{$\forall x\,\forall y\,\forall z\,(x = y \to y = z \to x = z)$}
\AXC{$\forall x\,\forall y\,\forall z\,(x = y \to y = z \to x = z) \to \varphi_{57}$\lbl{A2}}
\BIC{$\varphi_{57}$}
\end{ptree}
\noindent$\Pi_{84}$:
\begin{ptree}
\AXC{$\forall x \in \mathbb{R}\,(x \cdot 1 = x)$\lbl{(R8)}}
\AXC{$\forall x \in \mathbb{R}\,(x \cdot 1 = x) \to r \in \mathbb{R} \to r \cdot 1 = r$\lbl{A2}}
\BIC{$r \in \mathbb{R} \to r \cdot 1 = r$}
\AXC{$(r \in \mathbb{R} \to r \cdot 1 = r) \to \varphi_{175}$\lbl{A1}}
\BIC{$\varphi_{175}$}
\end{ptree}
\noindent$\Pi_{85}$:
\begin{ptree}
\AXC{$x = x$\lbl{A5}}
\UIC{$\forall x\,(x = x)$}
\AXC{$\forall x\,(x = x) \to r \cdot 1 - r \cdot r = r \cdot 1 - r \cdot r$\lbl{A2}}
\BIC{$r \cdot 1 - r \cdot r = r \cdot 1 - r \cdot r$}
\AXC{$\stackrel{\Pi_{84}}{\vdots}$}\noLine
\UIC{$\varphi_{175}$}
\BIC{$\varphi_{40} \to r \in \mathbb{R} \to r \cdot 1 - r \cdot r = r - r \cdot r$}
\end{ptree}
\noindent$\Pi_{86}$:
\begin{ptree}
\AXC{$\varphi_{29}$\lbl{A6}}
\UIC{$\forall w\,(\varphi_{29})$}
\UIC{$\forall z\,\forall w\,(\varphi_{29})$}
\AXC{$\forall z\,\forall w\,(\varphi_{29}) \to \varphi_{48}$\lbl{A2}}
\BIC{$\varphi_{48}$}
\AXC{$\varphi_{48} \to \varphi_{40}$\lbl{A2}}
\BIC{$\varphi_{40}$}
\AXC{$\stackrel{\Pi_{85}}{\vdots}$}\noLine
\UIC{$\varphi_{40} \to r \in \mathbb{R} \to r \cdot 1 - r \cdot r = r - r \cdot r$}
\BIC{$r \in \mathbb{R} \to r \cdot 1 - r \cdot r = r - r \cdot r$}
\end{ptree}
\noindent$\Pi_{87}$:
\begin{ptree}
\AXC{$y = z \to x = y \to x = z$\lbl{A6}}
\AXC{$(y = z \to x = y \to x = z) \to x = y \to y = z \to x = z$\lbl{A1}}
\BIC{$x = y \to y = z \to x = z$}
\UIC{$\forall z\,(x = y \to y = z \to x = z)$}
\UIC{$\forall y\,\forall z\,(x = y \to y = z \to x = z)$}
\UIC{$\forall x\,\forall y\,\forall z\,(x = y \to y = z \to x = z)$}
\AXC{$\forall x\,\forall y\,\forall z\,(x = y \to y = z \to x = z) \to \varphi_{57}$\lbl{A2}}
\BIC{$\varphi_{57}$}
\end{ptree}
\noindent$\Pi_{88}$:
\begin{ptree}
\AXC{$\stackrel{\Pi_{7}}{\vdots}$}\noLine
\UIC{$\forall x\,\forall y\,(x = y \to y = x)$}
\AXC{$\varphi_{22}$\lbl{A2}}
\BIC{$\forall y\,(r \cdot 1 - r \cdot r = y \to y = r \cdot 1 - r \cdot r)$}
\AXC{$\forall y\,(r \cdot 1 - r \cdot r = y \to y = r \cdot 1 - r \cdot r) \to \varphi_{19}$\lbl{A2}}
\BIC{$\varphi_{19}$}
\end{ptree}
\noindent$\Pi_{89}$:
\begin{ptree}
\AXC{$1 \in \mathbb{R}$\lbl{(R1)}}
\AXC{$r \in \mathbb{R} \to r \in \mathbb{R}$\lbl{A1}}
\AXC{$\varphi_{42}$\lbl{[F17]}}
\AXC{$\varphi_{42} \to \varphi_{37}$\lbl{A2}}
\BIC{$\varphi_{37}$}
\AXC{$\varphi_{37} \to \varphi_{34}$\lbl{A2}}
\BIC{$\varphi_{34}$}
\AXC{$\varphi_{160}$\lbl{A2}}
\BIC{$r \in \mathbb{R} \to 1 \in \mathbb{R} \to r \in \mathbb{R} \to r \cdot 1 - r \cdot r = r \cdot (1 - r)$}
\AXC{$\varphi_{196}$\lbl{A1}}
\BIC{$(r \in \mathbb{R} \to r \in \mathbb{R}) \to 1 \in \mathbb{R} \to \varphi_{31}$}
\BIC{$1 \in \mathbb{R} \to \varphi_{31}$}
\BIC{$\varphi_{31}$}
\end{ptree}
\noindent$\Pi_{90}$:
\begin{ptree}
\AXC{$y = z \to x = y \to x = z$\lbl{A6}}
\AXC{$(y = z \to x = y \to x = z) \to x = y \to y = z \to x = z$\lbl{A1}}
\BIC{$x = y \to y = z \to x = z$}
\UIC{$\forall z\,(x = y \to y = z \to x = z)$}
\UIC{$\forall y\,\forall z\,(x = y \to y = z \to x = z)$}
\UIC{$\forall x\,\forall y\,\forall z\,(x = y \to y = z \to x = z)$}
\AXC{$\forall x\,\forall y\,\forall z\,(x = y \to y = z \to x = z) \to \varphi_{57}$\lbl{A2}}
\BIC{$\varphi_{57}$}
\end{ptree}
\noindent$\Pi_{91}$:
\begin{ptree}
\AXC{$1 \in \mathbb{R}$\lbl{(R1)}}
\AXC{$\forall x\,\forall y\,(x \in \mathbb{R} \to y \in \mathbb{R} \to x - y \in \mathbb{R})$\lbl{[F1]}}
\AXC{$\varphi_{84}$\lbl{A2}}
\BIC{$\forall y\,(1 \in \mathbb{R} \to y \in \mathbb{R} \to 1 - y \in \mathbb{R})$}
\AXC{$\varphi_{68}$\lbl{A2}}
\BIC{$1 \in \mathbb{R} \to r \in \mathbb{R} \to 1 - r \in \mathbb{R}$}
\AXC{$\varphi_{124}$\lbl{A1}}
\BIC{$1 \in \mathbb{R} \to (r \in \mathbb{R} \to r \in \mathbb{R}) \to r \in \mathbb{R} \to 1 - r \in \mathbb{R}$}
\BIC{$(r \in \mathbb{R} \to r \in \mathbb{R}) \to r \in \mathbb{R} \to 1 - r \in \mathbb{R}$}
\end{ptree}
\noindent$\Pi_{92}$:
\begin{ptree}
\AXC{$\forall x\,\forall y\,(x \in \mathbb{R} \to y \in \mathbb{R} \to x \cdot y = y \cdot x)$\lbl{(R7)}}
\AXC{$\forall x\,\forall y\,(x \in \mathbb{R} \to y \in \mathbb{R} \to x \cdot y = y \cdot x) \to \varphi_{13}$\lbl{A2}}
\BIC{$\varphi_{13}$}
\end{ptree}
\noindent$\Pi_{93}$:
\begin{ptree}
\AXC{$\stackrel{\Pi_{92}}{\vdots}$}\noLine
\UIC{$\varphi_{13}$}
\AXC{$\varphi_{13} \to 1 - r \in \mathbb{R} \to r \in \mathbb{R} \to (1 - r) \cdot r = r \cdot (1 - r)$\lbl{A2}}
\BIC{$1 - r \in \mathbb{R} \to r \in \mathbb{R} \to (1 - r) \cdot r = r \cdot (1 - r)$}
\AXC{$\varphi_{172}$\lbl{A1}}
\BIC{$(r \in \mathbb{R} \to 1 - r \in \mathbb{R}) \to \varphi_{26}$}
\end{ptree}
\noindent$\Pi_{94}$:
\begin{ptree}
\AXC{$r \in \mathbb{R} \to r \in \mathbb{R}$\lbl{A1}}
\AXC{$\stackrel{\Pi_{91}}{\vdots}$}\noLine
\UIC{$(r \in \mathbb{R} \to r \in \mathbb{R}) \to r \in \mathbb{R} \to 1 - r \in \mathbb{R}$}
\BIC{$r \in \mathbb{R} \to 1 - r \in \mathbb{R}$}
\AXC{$\stackrel{\Pi_{93}}{\vdots}$}\noLine
\UIC{$(r \in \mathbb{R} \to 1 - r \in \mathbb{R}) \to \varphi_{26}$}
\BIC{$\varphi_{26}$}
\end{ptree}
\noindent$\Pi_{95}$:
\begin{ptree}
\AXC{$y = z \to x = y \to x = z$\lbl{A6}}
\AXC{$(y = z \to x = y \to x = z) \to x = y \to y = z \to x = z$\lbl{A1}}
\BIC{$x = y \to y = z \to x = z$}
\UIC{$\forall z\,(x = y \to y = z \to x = z)$}
\UIC{$\forall y\,\forall z\,(x = y \to y = z \to x = z)$}
\UIC{$\forall x\,\forall y\,\forall z\,(x = y \to y = z \to x = z)$}
\AXC{$\forall x\,\forall y\,\forall z\,(x = y \to y = z \to x = z) \to \varphi_{57}$\lbl{A2}}
\BIC{$\varphi_{57}$}
\end{ptree}
\noindent$\Pi_{96}$:
\begin{ptree}
\AXC{$\forall r \in \mathbb{R}\,(\varphi_{58})$\lbl{[S24]}}
\AXC{$\forall r \in \mathbb{R}\,(\varphi_{58}) \to r \in \mathbb{R} \to \varphi_{58}$\lbl{A2}}
\BIC{$r \in \mathbb{R} \to \varphi_{58}$}
\AXC{$(r \in \mathbb{R} \to \varphi_{58}) \to r \in \mathbb{R} \to {{r}^{1_{\omega}}} = r$\lbl{A1}}
\BIC{$r \in \mathbb{R} \to {{r}^{1_{\omega}}} = r$}
\AXC{$\varphi_{217}$\lbl{A1}}
\BIC{$(1 - r) \cdot {{r}^{1_{\omega}}} = (1 - r) \cdot {{r}^{1_{\omega}}} \to \varphi_{186}$}
\end{ptree}
\noindent$\Pi_{97}$:
\begin{ptree}
\AXC{$x = x$\lbl{A5}}
\UIC{$\forall x\,(x = x)$}
\AXC{$\forall x\,(x = x) \to (1 - r) \cdot {{r}^{1_{\omega}}} = (1 - r) \cdot {{r}^{1_{\omega}}}$\lbl{A2}}
\BIC{$(1 - r) \cdot {{r}^{1_{\omega}}} = (1 - r) \cdot {{r}^{1_{\omega}}}$}
\AXC{$\stackrel{\Pi_{96}}{\vdots}$}\noLine
\UIC{$(1 - r) \cdot {{r}^{1_{\omega}}} = (1 - r) \cdot {{r}^{1_{\omega}}} \to \varphi_{186}$}
\BIC{$\varphi_{186}$}
\end{ptree}
\noindent$\Pi_{98}$:
\begin{ptree}
\AXC{$z = w \to \varphi_{46}$\lbl{A6}}
\UIC{$\forall w\,(z = w \to \varphi_{46})$}
\UIC{$\forall z\,\forall w\,(z = w \to \varphi_{46})$}
\AXC{$\forall z\,\forall w\,(z = w \to \varphi_{46}) \to \varphi_{131}$\lbl{A2}}
\BIC{$\varphi_{131}$}
\AXC{$\varphi_{131} \to {{r}^{1_{\omega}}} = r \to \varphi_{80}$\lbl{A2}}
\BIC{${{r}^{1_{\omega}}} = r \to \varphi_{80}$}
\AXC{$\stackrel{\Pi_{97}}{\vdots}$}\noLine
\UIC{$\varphi_{186}$}
\BIC{$r \in \mathbb{R} \to (1 - r) \cdot {{r}^{1_{\omega}}} = (1 - r) \cdot r$}
\end{ptree}
\noindent$\Pi_{99}$:
\begin{ptree}
\AXC{$\forall r \in \mathbb{R}\,(\varphi_{58})$\lbl{[S24]}}
\AXC{$\forall r \in \mathbb{R}\,(\varphi_{58}) \to r \in \mathbb{R} \to \varphi_{58}$\lbl{A2}}
\BIC{$r \in \mathbb{R} \to \varphi_{58}$}
\AXC{$(r \in \mathbb{R} \to \varphi_{58}) \to r \in \mathbb{R} \to \mathrm{Seq}({(r)^{\bullet}})$\lbl{A1}}
\BIC{$r \in \mathbb{R} \to \mathrm{Seq}({(r)^{\bullet}})$}
\end{ptree}
\noindent$\Pi_{100}$:
\begin{ptree}
\AXC{$\stackrel{\Pi_{99}}{\vdots}$}\noLine
\UIC{$r \in \mathbb{R} \to \mathrm{Seq}({(r)^{\bullet}})$}
\AXC{$\varphi_{139}$\lbl{[S25]}}
\AXC{$\varphi_{139} \to \mathrm{Seq}({(r)^{\bullet}}) \to \varphi_{153}$\lbl{A2}}
\BIC{$\mathrm{Seq}({(r)^{\bullet}}) \to \varphi_{153}$}
\AXC{$\varphi_{230}$\lbl{A1}}
\BIC{$(r \in \mathbb{R} \to \mathrm{Seq}({(r)^{\bullet}})) \to r \in \mathbb{R} \to \varphi_{153}$}
\BIC{$r \in \mathbb{R} \to \varphi_{153}$}
\end{ptree}
\noindent$\Pi_{101}$:
\begin{ptree}
\AXC{$\stackrel{\Pi_{100}}{\vdots}$}\noLine
\UIC{$r \in \mathbb{R} \to \varphi_{153}$}
\AXC{$(r \in \mathbb{R} \to \varphi_{153}) \to r \in \mathbb{R} \to {\textstyle\sum_{k=1}^{1_{\omega}} {{r}^{k}}} = {{r}^{1_{\omega}}}$\lbl{A1}}
\BIC{$r \in \mathbb{R} \to {\textstyle\sum_{k=1}^{1_{\omega}} {{r}^{k}}} = {{r}^{1_{\omega}}}$}
\AXC{$\varphi_{240}$\lbl{A1}}
\BIC{$(1 - r) \cdot {\textstyle\sum_{k=1}^{1_{\omega}} {{r}^{k}}} = (1 - r) \cdot {\textstyle\sum_{k=1}^{1_{\omega}} {{r}^{k}}} \to \varphi_{214}$}
\end{ptree}
\noindent$\Pi_{102}$:
\begin{ptree}
\AXC{$x = x$\lbl{A5}}
\UIC{$\forall x\,(x = x)$}
\AXC{$\forall x\,(x = x) \to (1 - r) \cdot {\textstyle\sum_{k=1}^{1_{\omega}} {{r}^{k}}} = (1 - r) \cdot {\textstyle\sum_{k=1}^{1_{\omega}} {{r}^{k}}}$\lbl{A2}}
\BIC{$(1 - r) \cdot {\textstyle\sum_{k=1}^{1_{\omega}} {{r}^{k}}} = (1 - r) \cdot {\textstyle\sum_{k=1}^{1_{\omega}} {{r}^{k}}}$}
\AXC{$\stackrel{\Pi_{101}}{\vdots}$}\noLine
\UIC{$(1 - r) \cdot {\textstyle\sum_{k=1}^{1_{\omega}} {{r}^{k}}} = (1 - r) \cdot {\textstyle\sum_{k=1}^{1_{\omega}} {{r}^{k}}} \to \varphi_{214}$}
\BIC{$\varphi_{214}$}
\end{ptree}
\noindent$\Pi_{103}$:
\begin{ptree}
\AXC{$z = w \to \varphi_{75}$\lbl{A6}}
\UIC{$\forall w\,(z = w \to \varphi_{75})$}
\UIC{$\forall z\,\forall w\,(z = w \to \varphi_{75})$}
\AXC{$\forall z\,\forall w\,(z = w \to \varphi_{75}) \to \varphi_{159}$\lbl{A2}}
\BIC{$\varphi_{159}$}
\AXC{$\varphi_{159} \to {\textstyle\sum_{k=1}^{1_{\omega}} {{r}^{k}}} = {{r}^{1_{\omega}}} \to \varphi_{143}$\lbl{A2}}
\BIC{${\textstyle\sum_{k=1}^{1_{\omega}} {{r}^{k}}} = {{r}^{1_{\omega}}} \to \varphi_{143}$}
\AXC{$\stackrel{\Pi_{102}}{\vdots}$}\noLine
\UIC{$\varphi_{214}$}
\BIC{$r \in \mathbb{R} \to (1 - r) \cdot {\textstyle\sum_{k=1}^{1_{\omega}} {{r}^{k}}} = (1 - r) \cdot {{r}^{1_{\omega}}}$}
\AXC{$\varphi_{235}$\lbl{A1}}
\BIC{$(r \in \mathbb{R} \to (1 - r) \cdot {{r}^{1_{\omega}}} = (1 - r) \cdot r) \to \varphi_{204}$}
\end{ptree}
\noindent$\Pi_{104}$:
\begin{ptree}
\AXC{$\stackrel{\Pi_{98}}{\vdots}$}\noLine
\UIC{$r \in \mathbb{R} \to (1 - r) \cdot {{r}^{1_{\omega}}} = (1 - r) \cdot r$}
\AXC{$\stackrel{\Pi_{103}}{\vdots}$}\noLine
\UIC{$(r \in \mathbb{R} \to (1 - r) \cdot {{r}^{1_{\omega}}} = (1 - r) \cdot r) \to \varphi_{204}$}
\BIC{$\varphi_{204}$}
\end{ptree}
\noindent$\Pi_{105}$:
\begin{ptree}
\AXC{$\stackrel{\Pi_{95}}{\vdots}$}\noLine
\UIC{$\varphi_{57}$}
\AXC{$\varphi_{57} \to \varphi_{148}$\lbl{A2}}
\BIC{$\varphi_{148}$}
\AXC{$\varphi_{148} \to (1 - r) \cdot {\textstyle\sum_{k=1}^{1_{\omega}} {{r}^{k}}} = (1 - r) \cdot {{r}^{1_{\omega}}} \to \varphi_{62}$\lbl{A2}}
\BIC{$(1 - r) \cdot {\textstyle\sum_{k=1}^{1_{\omega}} {{r}^{k}}} = (1 - r) \cdot {{r}^{1_{\omega}}} \to \varphi_{62}$}
\AXC{$\stackrel{\Pi_{104}}{\vdots}$}\noLine
\UIC{$\varphi_{204}$}
\BIC{$r \in \mathbb{R} \to (1 - r) \cdot {\textstyle\sum_{k=1}^{1_{\omega}} {{r}^{k}}} = (1 - r) \cdot r$}
\AXC{$\varphi_{223}$\lbl{A1}}
\BIC{$(r \in \mathbb{R} \to (1 - r) \cdot r = r \cdot (1 - r)) \to \varphi_{194}$}
\end{ptree}
\noindent$\Pi_{106}$:
\begin{ptree}
\AXC{$r \in \mathbb{R} \to r \in \mathbb{R}$\lbl{A1}}
\AXC{$\stackrel{\Pi_{94}}{\vdots}$}\noLine
\UIC{$\varphi_{26}$}
\BIC{$r \in \mathbb{R} \to (1 - r) \cdot r = r \cdot (1 - r)$}
\AXC{$\stackrel{\Pi_{105}}{\vdots}$}\noLine
\UIC{$(r \in \mathbb{R} \to (1 - r) \cdot r = r \cdot (1 - r)) \to \varphi_{194}$}
\BIC{$\varphi_{194}$}
\end{ptree}
\noindent$\Pi_{107}$:
\begin{ptree}
\AXC{$\stackrel{\Pi_{90}}{\vdots}$}\noLine
\UIC{$\varphi_{57}$}
\AXC{$\varphi_{57} \to \varphi_{96}$\lbl{A2}}
\BIC{$\varphi_{96}$}
\AXC{$\varphi_{96} \to (1 - r) \cdot {\textstyle\sum_{k=1}^{1_{\omega}} {{r}^{k}}} = (1 - r) \cdot r \to \varphi_{33}$\lbl{A2}}
\BIC{$(1 - r) \cdot {\textstyle\sum_{k=1}^{1_{\omega}} {{r}^{k}}} = (1 - r) \cdot r \to \varphi_{33}$}
\AXC{$\stackrel{\Pi_{106}}{\vdots}$}\noLine
\UIC{$\varphi_{194}$}
\BIC{$r \in \mathbb{R} \to (1 - r) \cdot {\textstyle\sum_{k=1}^{1_{\omega}} {{r}^{k}}} = r \cdot (1 - r)$}
\AXC{$\varphi_{226}$\lbl{A1}}
\BIC{$(r \in \mathbb{R} \to r \cdot (1 - r) = r \cdot 1 - r \cdot r) \to \varphi_{201}$}
\end{ptree}
\noindent$\Pi_{108}$:
\begin{ptree}
\AXC{$\stackrel{\Pi_{88}}{\vdots}$}\noLine
\UIC{$\varphi_{19}$}
\AXC{$r \in \mathbb{R} \to r \in \mathbb{R}$\lbl{A1}}
\AXC{$\stackrel{\Pi_{89}}{\vdots}$}\noLine
\UIC{$\varphi_{31}$}
\BIC{$r \in \mathbb{R} \to r \cdot 1 - r \cdot r = r \cdot (1 - r)$}
\AXC{$\varphi_{168}$\lbl{A1}}
\BIC{$\varphi_{19} \to r \in \mathbb{R} \to r \cdot (1 - r) = r \cdot 1 - r \cdot r$}
\BIC{$r \in \mathbb{R} \to r \cdot (1 - r) = r \cdot 1 - r \cdot r$}
\AXC{$\stackrel{\Pi_{107}}{\vdots}$}\noLine
\UIC{$(r \in \mathbb{R} \to r \cdot (1 - r) = r \cdot 1 - r \cdot r) \to \varphi_{201}$}
\BIC{$\varphi_{201}$}
\end{ptree}
\noindent$\Pi_{109}$:
\begin{ptree}
\AXC{$\stackrel{\Pi_{87}}{\vdots}$}\noLine
\UIC{$\varphi_{57}$}
\AXC{$\varphi_{57} \to \varphi_{95}$\lbl{A2}}
\BIC{$\varphi_{95}$}
\AXC{$\varphi_{95} \to (1 - r) \cdot {\textstyle\sum_{k=1}^{1_{\omega}} {{r}^{k}}} = r \cdot (1 - r) \to \varphi_{49}$\lbl{A2}}
\BIC{$(1 - r) \cdot {\textstyle\sum_{k=1}^{1_{\omega}} {{r}^{k}}} = r \cdot (1 - r) \to \varphi_{49}$}
\AXC{$\stackrel{\Pi_{108}}{\vdots}$}\noLine
\UIC{$\varphi_{201}$}
\BIC{$r \in \mathbb{R} \to (1 - r) \cdot {\textstyle\sum_{k=1}^{1_{\omega}} {{r}^{k}}} = r \cdot 1 - r \cdot r$}
\end{ptree}
\noindent$\Pi_{110}$:
\begin{ptree}
\AXC{$\stackrel{\Pi_{109}}{\vdots}$}\noLine
\UIC{$r \in \mathbb{R} \to (1 - r) \cdot {\textstyle\sum_{k=1}^{1_{\omega}} {{r}^{k}}} = r \cdot 1 - r \cdot r$}
\AXC{$(r \in \mathbb{R} \to (1 - r) \cdot {\textstyle\sum_{k=1}^{1_{\omega}} {{r}^{k}}} = r \cdot 1 - r \cdot r) \to \varphi_{213}$\lbl{A1}}
\BIC{$\varphi_{213}$}
\end{ptree}
\noindent$\Pi_{111}$:
\begin{ptree}
\AXC{$\stackrel{\Pi_{83}}{\vdots}$}\noLine
\UIC{$\varphi_{57}$}
\AXC{$\varphi_{57} \to \varphi_{123}$\lbl{A2}}
\BIC{$\varphi_{123}$}
\AXC{$\varphi_{123} \to \varphi_{144}$\lbl{A2}}
\BIC{$\varphi_{144}$}
\AXC{$\stackrel{\Pi_{86}}{\vdots}$}\noLine
\UIC{$r \in \mathbb{R} \to r \cdot 1 - r \cdot r = r - r \cdot r$}
\AXC{$\stackrel{\Pi_{110}}{\vdots}$}\noLine
\UIC{$\varphi_{213}$}
\BIC{$\varphi_{144} \to r \in \mathbb{R} \to (1 - r) \cdot {\textstyle\sum_{k=1}^{1_{\omega}} {{r}^{k}}} = r - r \cdot r$}
\BIC{$r \in \mathbb{R} \to (1 - r) \cdot {\textstyle\sum_{k=1}^{1_{\omega}} {{r}^{k}}} = r - r \cdot r$}
\AXC{$\varphi_{236}$\lbl{A1}}
\BIC{$(r \in \mathbb{R} \to r - r \cdot r = r - {{r}^{1_{\omega}}} \cdot r) \to \varphi_{211}$}
\end{ptree}
\noindent$\Pi_{112}$:
\begin{ptree}
\AXC{$\stackrel{\Pi_{79}}{\vdots}$}\noLine
\UIC{$\varphi_{60}$}
\AXC{$\stackrel{\Pi_{82}}{\vdots}$}\noLine
\UIC{$\varphi_{60} \to r \in \mathbb{R} \to r - r \cdot r = r - {{r}^{1_{\omega}}} \cdot r$}
\BIC{$r \in \mathbb{R} \to r - r \cdot r = r - {{r}^{1_{\omega}}} \cdot r$}
\AXC{$\stackrel{\Pi_{111}}{\vdots}$}\noLine
\UIC{$(r \in \mathbb{R} \to r - r \cdot r = r - {{r}^{1_{\omega}}} \cdot r) \to \varphi_{211}$}
\BIC{$\varphi_{211}$}
\end{ptree}
\noindent$\Pi_{113}$:
\begin{ptree}
\AXC{$\stackrel{\Pi_{78}}{\vdots}$}\noLine
\UIC{$\varphi_{57}$}
\AXC{$\varphi_{57} \to \varphi_{94}$\lbl{A2}}
\BIC{$\varphi_{94}$}
\AXC{$\varphi_{94} \to (1 - r) \cdot {\textstyle\sum_{k=1}^{1_{\omega}} {{r}^{k}}} = r - r \cdot r \to \varphi_{92}$\lbl{A2}}
\BIC{$(1 - r) \cdot {\textstyle\sum_{k=1}^{1_{\omega}} {{r}^{k}}} = r - r \cdot r \to \varphi_{92}$}
\AXC{$\stackrel{\Pi_{112}}{\vdots}$}\noLine
\UIC{$\varphi_{211}$}
\BIC{$r \in \mathbb{R} \to (1 - r) \cdot {\textstyle\sum_{k=1}^{1_{\omega}} {{r}^{k}}} = r - {{r}^{1_{\omega}}} \cdot r$}
\AXC{$\varphi_{244}$\lbl{A1}}
\BIC{$(r \in \mathbb{R} \to r - {{r}^{1_{\omega}}} \cdot r = r - {{r}^{S(1_{\omega})}}) \to \varphi_{215}$}
\end{ptree}
\noindent$\Pi_{114}$:
\begin{ptree}
\AXC{$\stackrel{\Pi_{72}}{\vdots}$}\noLine
\UIC{${{r}^{1_{\omega}}} \cdot r = {{r}^{S(1_{\omega})}} \to \varphi_{114}$}
\AXC{$\stackrel{\Pi_{77}}{\vdots}$}\noLine
\UIC{$\varphi_{208}$}
\BIC{$r \in \mathbb{R} \to r - {{r}^{1_{\omega}}} \cdot r = r - {{r}^{S(1_{\omega})}}$}
\AXC{$\stackrel{\Pi_{113}}{\vdots}$}\noLine
\UIC{$(r \in \mathbb{R} \to r - {{r}^{1_{\omega}}} \cdot r = r - {{r}^{S(1_{\omega})}}) \to \varphi_{215}$}
\BIC{$\varphi_{215}$}
\end{ptree}
\noindent$\Pi_{115}$:
\begin{ptree}
\AXC{$\stackrel{\Pi_{71}}{\vdots}$}\noLine
\UIC{$\varphi_{57}$}
\AXC{$\varphi_{57} \to \varphi_{147}$\lbl{A2}}
\BIC{$\varphi_{147}$}
\AXC{$\varphi_{147} \to (1 - r) \cdot {\textstyle\sum_{k=1}^{1_{\omega}} {{r}^{k}}} = r - {{r}^{1_{\omega}}} \cdot r \to \varphi_{122}$\lbl{A2}}
\BIC{$(1 - r) \cdot {\textstyle\sum_{k=1}^{1_{\omega}} {{r}^{k}}} = r - {{r}^{1_{\omega}}} \cdot r \to \varphi_{122}$}
\AXC{$\stackrel{\Pi_{114}}{\vdots}$}\noLine
\UIC{$\varphi_{215}$}
\BIC{$r \in \mathbb{R} \to (1 - r) \cdot {\textstyle\sum_{k=1}^{1_{\omega}} {{r}^{k}}} = r - {{r}^{S(1_{\omega})}}$}
\end{ptree}
\noindent$\Pi_{116}$:
\begin{ptree}
\AXC{$y = z \to x = y \to x = z$\lbl{A6}}
\AXC{$(y = z \to x = y \to x = z) \to x = y \to y = z \to x = z$\lbl{A1}}
\BIC{$x = y \to y = z \to x = z$}
\UIC{$\forall z\,(x = y \to y = z \to x = z)$}
\UIC{$\forall y\,\forall z\,(x = y \to y = z \to x = z)$}
\UIC{$\forall x\,\forall y\,\forall z\,(x = y \to y = z \to x = z)$}
\AXC{$\varphi_{8}$\lbl{A2}}
\BIC{$\forall y\,\forall z\,(S(n) = y \to y = z \to S(n) = z)$}
\AXC{$\varphi_{38}$\lbl{A2}}
\BIC{$\forall z\,(S(n) = S(\emptyset) \to S(\emptyset) = z \to S(n) = z)$}
\end{ptree}
\noindent$\Pi_{117}$:
\begin{ptree}
\AXC{$\stackrel{\Pi_{116}}{\vdots}$}\noLine
\UIC{$\forall z\,(S(n) = S(\emptyset) \to S(\emptyset) = z \to S(n) = z)$}
\AXC{$\forall z\,(S(n) = S(\emptyset) \to S(\emptyset) = z \to S(n) = z) \to \varphi_{4}$\lbl{A2}}
\BIC{$\varphi_{4}$}
\end{ptree}
\noindent$\Pi_{118}$:
\begin{ptree}
\AXC{$\stackrel{\Pi_{7}}{\vdots}$}\noLine
\UIC{$\forall x\,\forall y\,(x = y \to y = x)$}
\AXC{$\varphi_{2}$\lbl{A2}}
\BIC{$\forall y\,(1_{\omega} = y \to y = 1_{\omega})$}
\AXC{$\varphi_{10}$\lbl{A2}}
\BIC{$1_{\omega} = S(\emptyset) \to S(\emptyset) = 1_{\omega}$}
\AXC{$1_{\omega} = S(\emptyset)$\lbl{Def$1_\omega$}}
\AXC{$1_{\omega} = S(\emptyset) \to \varphi_{9}$\lbl{A1}}
\BIC{$\varphi_{9}$}
\BIC{$S(\emptyset) = 1_{\omega}$}
\end{ptree}
\noindent$\Pi_{119}$:
\begin{ptree}
\AXC{$x = x$\lbl{A5}}
\UIC{$\forall x\,(x = x)$}
\AXC{$\forall x\,(x = x) \to S(n) = S(n)$\lbl{A2}}
\BIC{$S(n) = S(n)$}
\AXC{$\varphi_{178}$\lbl{A1}}
\AXC{$\varphi_{178} \to S(n) = S(n) \to \varphi_{205}$\lbl{A1}}
\BIC{$S(n) = S(n) \to \varphi_{205}$}
\BIC{$\varphi_{205}$}
\end{ptree}
\noindent$\Pi_{120}$:
\begin{ptree}
\AXC{$z = w \to S(n) = S(z) \to S(n) = S(w)$\lbl{A6}}
\UIC{$\forall w\,(z = w \to S(n) = S(z) \to S(n) = S(w))$}
\UIC{$\forall z\,\forall w\,(z = w \to S(n) = S(z) \to S(n) = S(w))$}
\AXC{$\varphi_{30}$\lbl{A2}}
\BIC{$\forall w\,(n = w \to S(n) = S(n) \to S(n) = S(w))$}
\AXC{$\varphi_{39}$\lbl{A2}}
\BIC{$n = \emptyset \to S(n) = S(n) \to S(n) = S(\emptyset)$}
\AXC{$\stackrel{\Pi_{119}}{\vdots}$}\noLine
\UIC{$\varphi_{205}$}
\BIC{$r \in \mathbb{R} \to \varphi_{171}$}
\end{ptree}
\noindent$\Pi_{121}$:
\begin{ptree}
\AXC{$\stackrel{\Pi_{117}}{\vdots}$}\noLine
\UIC{$\varphi_{4}$}
\AXC{$\stackrel{\Pi_{118}}{\vdots}$}\noLine
\UIC{$S(\emptyset) = 1_{\omega}$}
\AXC{$\stackrel{\Pi_{120}}{\vdots}$}\noLine
\UIC{$r \in \mathbb{R} \to \varphi_{171}$}
\AXC{$(r \in \mathbb{R} \to \varphi_{171}) \to S(\emptyset) = 1_{\omega} \to \varphi_{4} \to r \in \mathbb{R} \to \varphi_{162}$\lbl{A1}}
\BIC{$S(\emptyset) = 1_{\omega} \to \varphi_{4} \to r \in \mathbb{R} \to \varphi_{162}$}
\BIC{$\varphi_{4} \to r \in \mathbb{R} \to \varphi_{162}$}
\BIC{$r \in \mathbb{R} \to \varphi_{162}$}
\AXC{$\varphi_{238}$\lbl{A1}}
\BIC{$(S(n) = 1_{\omega} \to 1_{\omega} = S(n)) \to r \in \mathbb{R} \to \varphi_{163}$}
\end{ptree}
\noindent$\Pi_{122}$:
\begin{ptree}
\AXC{$\stackrel{\Pi_{7}}{\vdots}$}\noLine
\UIC{$\forall x\,\forall y\,(x = y \to y = x)$}
\AXC{$\varphi_{1}$\lbl{A2}}
\BIC{$\forall y\,(S(n) = y \to y = S(n))$}
\AXC{$\varphi_{3}$\lbl{A2}}
\BIC{$S(n) = 1_{\omega} \to 1_{\omega} = S(n)$}
\AXC{$\stackrel{\Pi_{121}}{\vdots}$}\noLine
\UIC{$(S(n) = 1_{\omega} \to 1_{\omega} = S(n)) \to r \in \mathbb{R} \to \varphi_{163}$}
\BIC{$r \in \mathbb{R} \to \varphi_{163}$}
\AXC{$\varphi_{272}$\lbl{A1}}
\BIC{$(r \in \mathbb{R} \to (1 - r) \cdot {\textstyle\sum_{k=1}^{1_{\omega}} {{r}^{k}}} = r - {{r}^{S(1_{\omega})}}) \to \varphi_{237}$}
\end{ptree}
\noindent$\Pi_{123}$:
\begin{ptree}
\AXC{$\stackrel{\Pi_{115}}{\vdots}$}\noLine
\UIC{$r \in \mathbb{R} \to (1 - r) \cdot {\textstyle\sum_{k=1}^{1_{\omega}} {{r}^{k}}} = r - {{r}^{S(1_{\omega})}}$}
\AXC{$\stackrel{\Pi_{122}}{\vdots}$}\noLine
\UIC{$(r \in \mathbb{R} \to (1 - r) \cdot {\textstyle\sum_{k=1}^{1_{\omega}} {{r}^{k}}} = r - {{r}^{S(1_{\omega})}}) \to \varphi_{237}$}
\BIC{$\varphi_{237}$}
\end{ptree}
\noindent$\Pi_{124}$:
\begin{ptree}
\AXC{$z = w \to \varphi_{107}$\lbl{A6}}
\UIC{$\forall w\,(z = w \to \varphi_{107})$}
\UIC{$\forall z\,\forall w\,(z = w \to \varphi_{107})$}
\AXC{$\forall z\,\forall w\,(z = w \to \varphi_{107}) \to \varphi_{137}$\lbl{A2}}
\BIC{$\varphi_{137}$}
\AXC{$\varphi_{137} \to 1_{\omega} = S(n) \to \varphi_{128}$\lbl{A2}}
\BIC{$1_{\omega} = S(n) \to \varphi_{128}$}
\AXC{$\stackrel{\Pi_{123}}{\vdots}$}\noLine
\UIC{$\varphi_{237}$}
\BIC{$r \in \mathbb{R} \to \forall n\,(\varphi_{66}) \to n \in B_{51} \to \varphi_{63}$}
\end{ptree}
\noindent$\Pi_{125}$:
\begin{ptree}
\AXC{$\stackrel{\Pi_{70}}{\vdots}$}\noLine
\UIC{$r \in \mathbb{R} \to \forall n\,(\varphi_{66}) \to n \in B_{51} \to \varphi_{64}$}
\AXC{$\stackrel{\Pi_{124}}{\vdots}$}\noLine
\UIC{$r \in \mathbb{R} \to \forall n\,(\varphi_{66}) \to n \in B_{51} \to \varphi_{63}$}
\AXC{$(r \in \mathbb{R} \to \forall n\,(\varphi_{66}) \to n \in B_{51} \to \varphi_{63}) \to \varphi_{249}$\lbl{A1}}
\BIC{$\varphi_{249}$}
\BIC{$r \in \mathbb{R} \to \forall n\,(\varphi_{66}) \to \varphi_{59}$}
\end{ptree}
\noindent$\Pi_{126}$:
\begin{ptree}
\AXC{$\forall x\,(\mathrm{Ind}(x) \leftrightarrow \emptyset \in x \land \forall y \in x\,(S(y) \in x))$\lbl{DefInd}}
\AXC{$\varphi_{138}$\lbl{A2}}
\BIC{$\mathrm{Ind}(\omega) \leftrightarrow \emptyset \in \omega \land \forall y \in \omega\,(S(y) \in \omega)$}
\AXC{$\mathrm{Ind}(\omega)$\lbl{補題(ωは帰納的)}}
\AXC{$\mathrm{Ind}(\omega) \to \varphi_{212}$\lbl{A1}}
\BIC{$\varphi_{212}$}
\BIC{$r \in \mathbb{R} \to \forall n\,(\varphi_{66}) \to \emptyset \in \omega \land \forall y \in \omega\,(S(y) \in \omega)$}
\end{ptree}
\noindent$\Pi_{127}$:
\begin{ptree}
\AXC{$\forall y \in \omega\,(S(y) \in \omega) \to n \in \omega \to S(n) \in \omega$\lbl{A2}}
\AXC{$\stackrel{\Pi_{126}}{\vdots}$}\noLine
\UIC{$r \in \mathbb{R} \to \forall n\,(\varphi_{66}) \to \emptyset \in \omega \land \forall y \in \omega\,(S(y) \in \omega)$}
\AXC{$\varphi_{228}$\lbl{A1}}
\BIC{$r \in \mathbb{R} \to \forall n\,(\varphi_{66}) \to n \in B_{51} \to \forall y \in \omega\,(S(y) \in \omega)$}
\AXC{$\varphi_{241}$\lbl{A1}}
\BIC{$\varphi_{209}$}
\BIC{$r \in \mathbb{R} \to \forall n\,(\varphi_{66}) \to n \in B_{51} \to n \in \omega \to S(n) \in \omega$}
\end{ptree}
\noindent$\Pi_{128}$:
\begin{ptree}
\AXC{$\stackrel{\Pi_{10}}{\vdots}$}\noLine
\UIC{$r \in \mathbb{R} \to \forall n\,(\varphi_{66}) \to n \in B_{51} \to n \in \omega$}
\AXC{$\stackrel{\Pi_{127}}{\vdots}$}\noLine
\UIC{$r \in \mathbb{R} \to \forall n\,(\varphi_{66}) \to n \in B_{51} \to n \in \omega \to S(n) \in \omega$}
\AXC{$\varphi_{253}$\lbl{A1}}
\BIC{$\varphi_{219}$}
\BIC{$r \in \mathbb{R} \to \forall n\,(\varphi_{66}) \to n \in B_{51} \to S(n) \in \omega$}
\end{ptree}
\noindent$\Pi_{129}$:
\begin{ptree}
\AXC{$\forall n\,(\varphi_{66}) \to \varphi_{108}$\lbl{A2}}
\AXC{$r \in \mathbb{R} \to \forall n\,(\varphi_{66}) \to \forall n\,(\varphi_{66})$\lbl{A1}}
\AXC{$\varphi_{284}$\lbl{A1}}
\BIC{$(\forall n\,(\varphi_{66}) \to \varphi_{108}) \to r \in \mathbb{R} \to \forall n\,(\varphi_{66}) \to \varphi_{108}$}
\BIC{$r \in \mathbb{R} \to \forall n\,(\varphi_{66}) \to \varphi_{108}$}
\AXC{$\varphi_{289}$\lbl{A1}}
\BIC{$(r \in \mathbb{R} \to \forall n\,(\varphi_{66}) \to n \in B_{51} \to S(n) \in \omega) \to \varphi_{231}$}
\end{ptree}
\noindent$\Pi_{130}$:
\begin{ptree}
\AXC{$\stackrel{\Pi_{125}}{\vdots}$}\noLine
\UIC{$r \in \mathbb{R} \to \forall n\,(\varphi_{66}) \to \varphi_{59}$}
\AXC{$\stackrel{\Pi_{128}}{\vdots}$}\noLine
\UIC{$r \in \mathbb{R} \to \forall n\,(\varphi_{66}) \to n \in B_{51} \to S(n) \in \omega$}
\AXC{$\stackrel{\Pi_{129}}{\vdots}$}\noLine
\UIC{$(r \in \mathbb{R} \to \forall n\,(\varphi_{66}) \to n \in B_{51} \to S(n) \in \omega) \to \varphi_{231}$}
\BIC{$\varphi_{231}$}
\BIC{$r \in \mathbb{R} \to \forall n\,(\varphi_{66}) \to n \in B_{51} \to S(n) \in B_{51}$}
\end{ptree}
\noindent$\Pi_{131}$:
\begin{ptree}
\AXC{$\varphi_{222}$\lbl{A3}}
\AXC{$\stackrel{\Pi_{130}}{\vdots}$}\noLine
\UIC{$r \in \mathbb{R} \to \forall n\,(\varphi_{66}) \to n \in B_{51} \to S(n) \in B_{51}$}
\AXC{$\varphi_{218}$\lbl{A1}}
\BIC{$r \in \mathbb{R} \land \forall n\,(\varphi_{66}) \to n \in B_{51} \to S(n) \in B_{51}$}
\UIC{$\forall n\,(r \in \mathbb{R} \land \forall n\,(\varphi_{66}) \to n \in B_{51} \to S(n) \in B_{51})$}
\AXC{$\varphi_{293}$\lbl{A1}}
\BIC{$\varphi_{222} \to \varphi_{239}$}
\BIC{$\varphi_{239}$}
\end{ptree}
\noindent$\Pi_{132}$:
\begin{ptree}
\AXC{$\forall n\,(n \in \mathbb{N} \leftrightarrow n \in \omega \land n \neq \emptyset)$\lbl{Def$\mathbb{N}$}}
\AXC{$\varphi_{106}$\lbl{A2}}
\BIC{$\emptyset \in \mathbb{N} \leftrightarrow \emptyset \in \omega \land \emptyset \neq \emptyset$}
\AXC{$(\emptyset \in \mathbb{N} \leftrightarrow \emptyset \in \omega \land \emptyset \neq \emptyset) \to \emptyset = \emptyset \to \varphi_{85}$\lbl{A1}}
\BIC{$\emptyset = \emptyset \to \varphi_{85}$}
\end{ptree}
\noindent$\Pi_{133}$:
\begin{ptree}
\AXC{$\forall n\,(\varphi_{66}) \to \varphi_{142}$\lbl{A2}}
\AXC{$r \in \mathbb{R} \to \forall n\,(\varphi_{66}) \to \forall n\,(\varphi_{66})$\lbl{A1}}
\AXC{$\varphi_{287}$\lbl{A1}}
\BIC{$(\forall n\,(\varphi_{66}) \to \varphi_{142}) \to r \in \mathbb{R} \to \forall n\,(\varphi_{66}) \to \varphi_{142}$}
\BIC{$r \in \mathbb{R} \to \forall n\,(\varphi_{66}) \to \varphi_{142}$}
\AXC{$(r \in \mathbb{R} \to \forall n\,(\varphi_{66}) \to \varphi_{142}) \to \varphi_{245}$\lbl{A1}}
\BIC{$\varphi_{245}$}
\end{ptree}
\noindent$\Pi_{134}$:
\begin{ptree}
\AXC{$x = x$\lbl{A5}}
\UIC{$\forall x\,(x = x)$}
\AXC{$\forall x\,(x = x) \to \emptyset = \emptyset$\lbl{A2}}
\BIC{$\emptyset = \emptyset$}
\AXC{$\stackrel{\Pi_{132}}{\vdots}$}\noLine
\UIC{$\emptyset = \emptyset \to \varphi_{85}$}
\BIC{$\varphi_{85}$}
\AXC{$\stackrel{\Pi_{126}}{\vdots}$}\noLine
\UIC{$r \in \mathbb{R} \to \forall n\,(\varphi_{66}) \to \emptyset \in \omega \land \forall y \in \omega\,(S(y) \in \omega)$}
\AXC{$\stackrel{\Pi_{133}}{\vdots}$}\noLine
\UIC{$\varphi_{245}$}
\BIC{$\varphi_{85} \to r \in \mathbb{R} \to \forall n\,(\varphi_{66}) \to \emptyset \in B_{51}$}
\BIC{$r \in \mathbb{R} \to \forall n\,(\varphi_{66}) \to \emptyset \in B_{51}$}
\end{ptree}
\noindent$\Pi_{135}$:
\begin{ptree}
\AXC{$\forall x\,(\mathrm{Ind}(x) \leftrightarrow \emptyset \in x \land \forall y \in x\,(S(y) \in x))$\lbl{DefInd}}
\AXC{$\forall x\,(\mathrm{Ind}(x) \leftrightarrow \emptyset \in x \land \forall y \in x\,(S(y) \in x)) \to \varphi_{14}$\lbl{A2}}
\BIC{$\varphi_{14}$}
\end{ptree}
\noindent$\Pi_{136}$:
\begin{ptree}
\AXC{$\stackrel{\Pi_{134}}{\vdots}$}\noLine
\UIC{$r \in \mathbb{R} \to \forall n\,(\varphi_{66}) \to \emptyset \in B_{51}$}
\AXC{$\stackrel{\Pi_{135}}{\vdots}$}\noLine
\UIC{$\varphi_{14}$}
\AXC{$\varphi_{14} \to (r \in \mathbb{R} \to \forall n\,(\varphi_{66}) \to \emptyset \in B_{51}) \to \varphi_{242}$\lbl{A1}}
\BIC{$(r \in \mathbb{R} \to \forall n\,(\varphi_{66}) \to \emptyset \in B_{51}) \to \varphi_{242}$}
\BIC{$\varphi_{242}$}
\end{ptree}
\noindent$\Pi_{137}$:
\begin{ptree}
\AXC{$\stackrel{\Pi_{2}}{\vdots}$}\noLine
\UIC{$\varphi_{11}$}
\AXC{$\stackrel{\Pi_{3}}{\vdots}$}\noLine
\UIC{$\varphi_{12}$}
\AXC{$r \in \mathbb{R} \land \forall n\,(\varphi_{66}) \to \forall n\,(r \in \mathbb{R} \land \forall n\,(\varphi_{66}))$\lbl{A4}}
\AXC{$\stackrel{\Pi_{131}}{\vdots}$}\noLine
\UIC{$\varphi_{239}$}
\BIC{$r \in \mathbb{R} \to \forall n\,(\varphi_{66}) \to \forall n \in B_{51}\,(S(n) \in B_{51})$}
\AXC{$\stackrel{\Pi_{136}}{\vdots}$}\noLine
\UIC{$\varphi_{242}$}
\BIC{$\varphi_{12} \to \varphi_{11} \to r \in \mathbb{R} \to \forall n\,(\varphi_{66}) \to \mathrm{Ind}(B_{51})$}
\BIC{$\varphi_{11} \to r \in \mathbb{R} \to \forall n\,(\varphi_{66}) \to \mathrm{Ind}(B_{51})$}
\BIC{$r \in \mathbb{R} \to \forall n\,(\varphi_{66}) \to \mathrm{Ind}(B_{51})$}
\end{ptree}
\noindent$\Pi_{138}$:
\begin{ptree}
\AXC{$r \in \mathbb{R} \to \forall n\,(\varphi_{66}) \to n \in \omega \to n \in \omega$\lbl{A1}}
\AXC{$(r \in \mathbb{R} \to \forall n\,(\varphi_{66}) \to n \in \omega \to n \in \omega) \to \varphi_{195}$\lbl{A1}}
\BIC{$\varphi_{195}$}
\end{ptree}
\noindent$\Pi_{139}$:
\begin{ptree}
\AXC{$\forall n\,(n \in \omega \leftrightarrow \forall x\,(\mathrm{Ind}(x) \to n \in x))$\lbl{Def$\omega$}}
\AXC{$\varphi_{36}$\lbl{A2}}
\BIC{$n \in \omega \leftrightarrow \forall x\,(\mathrm{Ind}(x) \to n \in x)$}
\AXC{$\stackrel{\Pi_{138}}{\vdots}$}\noLine
\UIC{$\varphi_{195}$}
\BIC{$r \in \mathbb{R} \to \forall n\,(\varphi_{66}) \to n \in \omega \to \forall x\,(\mathrm{Ind}(x) \to n \in x)$}
\AXC{$\varphi_{234}$\lbl{A1}}
\BIC{$\varphi_{199}$}
\end{ptree}
\noindent$\Pi_{140}$:
\begin{ptree}
\AXC{$\stackrel{\Pi_{137}}{\vdots}$}\noLine
\UIC{$r \in \mathbb{R} \to \forall n\,(\varphi_{66}) \to \mathrm{Ind}(B_{51})$}
\AXC{$\forall x\,(\mathrm{Ind}(x) \to n \in x) \to \mathrm{Ind}(B_{51}) \to n \in B_{51}$\lbl{A2}}
\AXC{$\stackrel{\Pi_{139}}{\vdots}$}\noLine
\UIC{$\varphi_{199}$}
\BIC{$r \in \mathbb{R} \to \forall n\,(\varphi_{66}) \to n \in \omega \to \mathrm{Ind}(B_{51}) \to n \in B_{51}$}
\AXC{$\varphi_{286}$\lbl{A1}}
\BIC{$\varphi_{262}$}
\BIC{$(r \in \mathbb{R} \to \forall n\,(\varphi_{66}) \to \varphi_{66}) \to r \in \mathbb{R} \to \varphi_{188}$}
\end{ptree}
\noindent$\Pi_{141}$:
\begin{ptree}
\AXC{$\forall n\,(\varphi_{66}) \to \varphi_{66}$\lbl{A2}}
\AXC{$r \in \mathbb{R} \to \forall n\,(\varphi_{66}) \to \forall n\,(\varphi_{66})$\lbl{A1}}
\AXC{$\varphi_{276}$\lbl{A1}}
\BIC{$(\forall n\,(\varphi_{66}) \to \varphi_{66}) \to r \in \mathbb{R} \to \forall n\,(\varphi_{66}) \to \varphi_{66}$}
\BIC{$r \in \mathbb{R} \to \forall n\,(\varphi_{66}) \to \varphi_{66}$}
\AXC{$\stackrel{\Pi_{140}}{\vdots}$}\noLine
\UIC{$(r \in \mathbb{R} \to \forall n\,(\varphi_{66}) \to \varphi_{66}) \to r \in \mathbb{R} \to \varphi_{188}$}
\BIC{$r \in \mathbb{R} \to \varphi_{188}$}
\end{ptree}
\noindent$\Pi_{142}$:
\begin{ptree}
\AXC{$\forall n\,(\varphi_{198}) \to \varphi_{200}$\lbl{A3}}
\AXC{$\stackrel{\Pi_{141}}{\vdots}$}\noLine
\UIC{$r \in \mathbb{R} \to \varphi_{188}$}
\AXC{$(r \in \mathbb{R} \to \varphi_{188}) \to \varphi_{198}$\lbl{A1}}
\BIC{$\varphi_{198}$}
\UIC{$\forall n\,(\varphi_{198})$}
\AXC{$\forall n\,(\varphi_{198}) \to (\forall n\,(\varphi_{198}) \to \varphi_{200}) \to \varphi_{251}$\lbl{A1}}
\BIC{$(\forall n\,(\varphi_{198}) \to \varphi_{200}) \to \varphi_{251}$}
\BIC{$\varphi_{251}$}
\end{ptree}
\noindent$\Pi_{143}$:
\begin{ptree}
\AXC{$r \in \mathbb{R} \land \forall n\,(\varphi_{66}) \to \forall n\,(r \in \mathbb{R} \land \forall n\,(\varphi_{66}))$\lbl{A4}}
\AXC{$\stackrel{\Pi_{142}}{\vdots}$}\noLine
\UIC{$\varphi_{251}$}
\BIC{$r \in \mathbb{R} \to \varphi_{191}$}
\AXC{$(r \in \mathbb{R} \to \varphi_{191}) \to \varphi_{81} \to \neg \forall n\,(\varphi_{66})$\lbl{A1}}
\BIC{$\varphi_{81} \to \neg \forall n\,(\varphi_{66})$}
\UIC{$\forall B_{51}\,(\varphi_{81} \to \neg \forall n\,(\varphi_{66}))$}
\end{ptree}
\noindent$\Pi_{144}$:
\begin{ptree}
\AXC{$\varphi_{125}$\lbl{Sep}}
\AXC{$\varphi_{125} \to \varphi_{97}$\lbl{A2}}
\BIC{$\varphi_{97}$}
\AXC{$\varphi_{97} \to \exists B_{51}\,\forall n\,(\varphi_{66})$\lbl{A2}}
\BIC{$\exists B_{51}\,\forall n\,(\varphi_{66})$}
\AXC{$\exists B_{51}\,\forall n\,(\varphi_{66}) \to r \in \mathbb{R} \to \exists B_{51}\,\forall n\,(\varphi_{66})$\lbl{A1}}
\BIC{$r \in \mathbb{R} \to \exists B_{51}\,\forall n\,(\varphi_{66})$}
\AXC{$\varphi_{302}$\lbl{A1}}
\BIC{$\forall B_{51}\,(\varphi_{81} \to \neg \forall n\,(\varphi_{66})) \to \varphi_{248} \to (\varphi_{81} \to \forall B_{51}\,(\varphi_{81})) \to \varphi_{79}$}
\end{ptree}
\noindent$\Pi_{145}$:
\begin{ptree}
\AXC{$\varphi_{81} \to \forall B_{51}\,(\varphi_{81})$\lbl{A4}}
\AXC{$\varphi_{248}$\lbl{A3}}
\AXC{$\stackrel{\Pi_{143}}{\vdots}$}\noLine
\UIC{$\forall B_{51}\,(\varphi_{81} \to \neg \forall n\,(\varphi_{66}))$}
\AXC{$\stackrel{\Pi_{144}}{\vdots}$}\noLine
\UIC{$\forall B_{51}\,(\varphi_{81} \to \neg \forall n\,(\varphi_{66})) \to \varphi_{248} \to (\varphi_{81} \to \forall B_{51}\,(\varphi_{81})) \to \varphi_{79}$}
\BIC{$\varphi_{248} \to (\varphi_{81} \to \forall B_{51}\,(\varphi_{81})) \to \varphi_{79}$}
\BIC{$(\varphi_{81} \to \forall B_{51}\,(\varphi_{81})) \to \varphi_{79}$}
\BIC{$\varphi_{79}$}
\end{ptree}
\noindent$\Pi_{146}$:
\begin{ptree}
\AXC{$\stackrel{\Pi_{1}}{\vdots}$}\noLine
\UIC{$r \in \mathbb{R} \to n \in \mathbb{N} \to n \in \omega$}
\AXC{$\varphi_{176}$\lbl{A2}}
\AXC{$\stackrel{\Pi_{145}}{\vdots}$}\noLine
\UIC{$\varphi_{79}$}
\AXC{$\varphi_{79} \to \varphi_{176} \to \varphi_{69}$\lbl{A1}}
\BIC{$\varphi_{176} \to \varphi_{69}$}
\BIC{$\varphi_{69}$}
\AXC{$\varphi_{69} \to (r \in \mathbb{R} \to n \in \mathbb{N} \to n \in \omega) \to \varphi_{126}$\lbl{A1}}
\BIC{$(r \in \mathbb{R} \to n \in \mathbb{N} \to n \in \omega) \to \varphi_{126}$}
\BIC{$\varphi_{126}$}
\end{ptree}
\noindent 主木:
\begin{ptree}
\AXC{$r \in \mathbb{R} \to \forall n\,(r \in \mathbb{R})$\lbl{A4}}
\AXC{$\varphi_{45} \to \varphi_{52}$\lbl{A3}}
\AXC{$r \in \mathbb{R} \to n \in \mathbb{N} \to n \in \mathbb{N}$\lbl{A1}}
\AXC{$\stackrel{\Pi_{146}}{\vdots}$}\noLine
\UIC{$\varphi_{126}$}
\BIC{$r \in \mathbb{R} \to n \in \mathbb{N} \to (1 - r) \cdot {\textstyle\sum_{k=1}^{n} {{r}^{k}}} = r - {{r}^{S(n)}}$}
\UIC{$\varphi_{45}$}
\AXC{$\varphi_{45} \to (\varphi_{45} \to \varphi_{52}) \to \varphi_{119}$\lbl{A1}}
\BIC{$(\varphi_{45} \to \varphi_{52}) \to \varphi_{119}$}
\BIC{$\varphi_{119}$}
\BIC{$r \in \mathbb{R} \to \forall n \in \mathbb{N}\,((1 - r) \cdot {\textstyle\sum_{k=1}^{n} {{r}^{k}}} = r - {{r}^{S(n)}})$}
\UIC{$\forall r \in \mathbb{R}\,\forall n \in \mathbb{N}\,((1 - r) \cdot {\textstyle\sum_{k=1}^{n} {{r}^{k}}} = r - {{r}^{S(n)}})$}
\end{ptree}
\begin{abbrevlist}
\item $\varphi_{1} :\equiv \forall x\,\forall y\,(x = y \to y = x) \to \forall y\,(S(n) = y \to y = S(n))$
\item $\varphi_{2} :\equiv \forall x\,\forall y\,(x = y \to y = x) \to \forall y\,(1_{\omega} = y \to y = 1_{\omega})$
\item $\varphi_{3} :\equiv \forall y\,(S(n) = y \to y = S(n)) \to S(n) = 1_{\omega} \to 1_{\omega} = S(n)$
\item $\varphi_{4} :\equiv S(n) = S(\emptyset) \to S(\emptyset) = 1_{\omega} \to S(n) = 1_{\omega}$
\item $\varphi_{5} :\equiv n \in B_{51} \to S(n) \in \mathbb{N} \to n \neq \emptyset \to \mathrm{Seq}({(r)^{\bullet}})$
\item $\varphi_{6} :\equiv \forall n \in B_{51}\,(S(n) \in B_{51}) \to y \in B_{51} \to S(y) \in B_{51}$
\item $\varphi_{7} :\equiv \forall y \in B_{51}\,(S(y) \in B_{51}) \to n \in B_{51} \to S(n) \in B_{51}$
\item $\varphi_{8} :\equiv \forall x\,\forall y\,\forall z\,(x = y \to y = z \to x = z) \to \forall y\,\forall z\,(S(n) = y \to y = z \to S(n) = z)$
\item $\varphi_{9} :\equiv (1_{\omega} = S(\emptyset) \to S(\emptyset) = 1_{\omega}) \to S(\emptyset) = 1_{\omega}$
\item $\varphi_{10} :\equiv \forall y\,(1_{\omega} = y \to y = 1_{\omega}) \to 1_{\omega} = S(\emptyset) \to S(\emptyset) = 1_{\omega}$
\item $\varphi_{11} :\equiv \forall n \in B_{51}\,(S(n) \in B_{51}) \to \forall y \in B_{51}\,(S(y) \in B_{51})$
\item $\varphi_{12} :\equiv \forall y \in B_{51}\,(S(y) \in B_{51}) \to \forall n \in B_{51}\,(S(n) \in B_{51})$
\item $\varphi_{13} :\equiv \forall y\,(1 - r \in \mathbb{R} \to y \in \mathbb{R} \to (1 - r) \cdot y = y \cdot (1 - r))$
\item $\varphi_{14} :\equiv \mathrm{Ind}(B_{51}) \leftrightarrow \emptyset \in B_{51} \land \forall y \in B_{51}\,(S(y) \in B_{51})$
\item $\varphi_{15} :\equiv n \in B_{51} \to S(n) \in \mathbb{N} \to n \neq \emptyset \to \mathrm{Seq}(\Sigma {(r)^{\bullet}})$
\item $\varphi_{16} :\equiv \forall w\,(z = w \to r - r \cdot r = r - z \cdot r \to r - r \cdot r = r - w \cdot r)$
\item $\varphi_{17} :\equiv \forall w\,(r = w \to r - r \cdot r = r - r \cdot r \to r - r \cdot r = r - w \cdot r)$
\item $\varphi_{18} :\equiv \forall x\,\forall y\,(x = y \to y = x) \to \forall y\,({{r}^{1_{\omega}}} = y \to y = {{r}^{1_{\omega}}})$
\item $\varphi_{19} :\equiv r \cdot 1 - r \cdot r = r \cdot (1 - r) \to r \cdot (1 - r) = r \cdot 1 - r \cdot r$
\item $\varphi_{20} :\equiv \forall n \in B_{51}\,(S(n) \in B_{51}) \to \forall y\,\forall n \in B_{51}\,(S(n) \in B_{51})$
\item $\varphi_{21} :\equiv \forall y \in B_{51}\,(S(y) \in B_{51}) \to \forall n\,\forall y \in B_{51}\,(S(y) \in B_{51})$
\item $\varphi_{22} :\equiv \forall x\,\forall y\,(x = y \to y = x) \to \forall y\,(r \cdot 1 - r \cdot r = y \to y = r \cdot 1 - r \cdot r)$
\item $\varphi_{23} :\equiv \forall y\,\forall n \in B_{51}\,(S(n) \in B_{51}) \to \forall y \in B_{51}\,(S(y) \in B_{51})$
\item $\varphi_{24} :\equiv \forall n\,\forall y \in B_{51}\,(S(y) \in B_{51}) \to \forall n \in B_{51}\,(S(n) \in B_{51})$
\item $\varphi_{25} :\equiv \forall n\,(\mathrm{Seq}(\Sigma {(r)^{\bullet}}) \to n \in \mathbb{N} \to {\textstyle\sum_{k=1}^{n} {{r}^{k}}} \in \mathbb{R})$
\item $\varphi_{26} :\equiv (r \in \mathbb{R} \to r \in \mathbb{R}) \to r \in \mathbb{R} \to (1 - r) \cdot r = r \cdot (1 - r)$
\item $\varphi_{27} :\equiv \forall x\,\forall y\,(x = y \to y = x) \to \forall y\,({{r}^{S(S(n))}} = y \to y = {{r}^{S(S(n))}})$
\item $\varphi_{28} :\equiv \forall x\,\forall y\,(x = y \to y = x) \to \forall y\,({{r}^{S(1_{\omega})}} = y \to y = {{r}^{S(1_{\omega})}})$
\item $\varphi_{29} :\equiv z = w \to r \cdot 1 - r \cdot r = z - r \cdot r \to r \cdot 1 - r \cdot r = w - r \cdot r$
\item $\varphi_{30} :\equiv \forall z\,\forall w\,(z = w \to S(n) = S(z) \to S(n) = S(w)) \to \forall w\,(n = w \to S(n) = S(n) \to S(n) = S(w))$
\item $\varphi_{31} :\equiv (r \in \mathbb{R} \to r \in \mathbb{R}) \to r \in \mathbb{R} \to r \cdot 1 - r \cdot r = r \cdot (1 - r)$
\item $\varphi_{32} :\equiv \forall b\,\forall x\,(r \in \mathbb{R} \to b \in \mathbb{R} \to x \in \mathbb{R} \to r - b = r - x + (x - b))$
\item $\varphi_{33} :\equiv (1 - r) \cdot r = r \cdot (1 - r) \to (1 - r) \cdot {\textstyle\sum_{k=1}^{1_{\omega}} {{r}^{k}}} = r \cdot (1 - r)$
\item $\varphi_{34} :\equiv \forall a\,(r \in \mathbb{R} \to 1 \in \mathbb{R} \to a \in \mathbb{R} \to r \cdot 1 - r \cdot a = r \cdot (1 - a))$
\item $\varphi_{35} :\equiv \forall a\,\forall b\,\forall x\,(a \in \mathbb{R} \to b \in \mathbb{R} \to x \in \mathbb{R} \to a - b = a - x + (x - b))$
\item $\varphi_{36} :\equiv \forall n\,(n \in \omega \leftrightarrow \forall x\,(\mathrm{Ind}(x) \to n \in x)) \to (n \in \omega \leftrightarrow \forall x\,(\mathrm{Ind}(x) \to n \in x))$
\item $\varphi_{37} :\equiv \forall x\,\forall a\,(r \in \mathbb{R} \to x \in \mathbb{R} \to a \in \mathbb{R} \to r \cdot x - r \cdot a = r \cdot (x - a))$
\item $\varphi_{38} :\equiv \forall y\,\forall z\,(S(n) = y \to y = z \to S(n) = z) \to \forall z\,(S(n) = S(\emptyset) \to S(\emptyset) = z \to S(n) = z)$
\item $\varphi_{39} :\equiv \forall w\,(n = w \to S(n) = S(n) \to S(n) = S(w)) \to n = \emptyset \to S(n) = S(n) \to S(n) = S(\emptyset)$
\item $\varphi_{40} :\equiv r \cdot 1 = r \to r \cdot 1 - r \cdot r = r \cdot 1 - r \cdot r \to r \cdot 1 - r \cdot r = r - r \cdot r$
\item $\varphi_{41} :\equiv \forall w\,(z = w \to r - {{r}^{S(n)}} \cdot r = r - z \to r - {{r}^{S(n)}} \cdot r = r - w)$
\item $\varphi_{42} :\equiv \forall c\,\forall x\,\forall a\,(c \in \mathbb{R} \to x \in \mathbb{R} \to a \in \mathbb{R} \to c \cdot x - c \cdot a = c \cdot (x - a))$
\item $\varphi_{43} :\equiv \forall x\,\forall y\,\forall z\,(x \in \mathbb{R} \to y \in \mathbb{R} \to z \in \mathbb{R} \to x \cdot (y + z) = x \cdot y + x \cdot z)$
\item $\varphi_{44} :\equiv \forall w\,(z = w \to r - {{r}^{1_{\omega}}} \cdot r = r - z \to r - {{r}^{1_{\omega}}} \cdot r = r - w)$
\item $\varphi_{45} :\equiv \forall n\,(r \in \mathbb{R} \to n \in \mathbb{N} \to (1 - r) \cdot {\textstyle\sum_{k=1}^{n} {{r}^{k}}} = r - {{r}^{S(n)}})$
\item $\varphi_{46} :\equiv (1 - r) \cdot {{r}^{1_{\omega}}} = (1 - r) \cdot z \to (1 - r) \cdot {{r}^{1_{\omega}}} = (1 - r) \cdot w$
\item $\varphi_{47} :\equiv \forall x\,\forall y\,(x = y \to y = x) \to \forall y\,(r - {{r}^{S(n)}} \cdot r = y \to y = r - {{r}^{S(n)}} \cdot r)$
\item $\varphi_{48} :\equiv \forall w\,(r \cdot 1 = w \to r \cdot 1 - r \cdot r = r \cdot 1 - r \cdot r \to r \cdot 1 - r \cdot r = w - r \cdot r)$
\item $\varphi_{49} :\equiv r \cdot (1 - r) = r \cdot 1 - r \cdot r \to (1 - r) \cdot {\textstyle\sum_{k=1}^{1_{\omega}} {{r}^{k}}} = r \cdot 1 - r \cdot r$
\item $\varphi_{50} :\equiv \forall n\,(n \in \mathbb{N} \leftrightarrow n \in \omega \land n \neq \emptyset) \to (n \in \mathbb{N} \leftrightarrow n \in \omega \land n \neq \emptyset)$
\item $\varphi_{51} :\equiv \forall y\,({{r}^{1_{\omega}}} = y \to y = {{r}^{1_{\omega}}}) \to {{r}^{1_{\omega}}} = r \to r = {{r}^{1_{\omega}}}$
\item $\varphi_{52} :\equiv \forall n\,(r \in \mathbb{R}) \to \forall n \in \mathbb{N}\,((1 - r) \cdot {\textstyle\sum_{k=1}^{n} {{r}^{k}}} = r - {{r}^{S(n)}})$
\item $\varphi_{53} :\equiv n \in B_{51} \to S(n) \in \mathbb{N} \to n \neq \emptyset \to {{r}^{S(S(n))}} = {{r}^{S(n)}} \cdot r$
\item $\varphi_{54} :\equiv n \in B_{51} \to S(n) \in \mathbb{N} \to n \neq \emptyset \to {{r}^{S(n)}} \cdot r = {{r}^{S(S(n))}}$
\item $\varphi_{55} :\equiv S(n) \in \mathbb{N} \to n \neq \emptyset \to (1 - r) \cdot {{r}^{S(n)}} = {{r}^{S(n)}} \cdot (1 - r)$
\item $\varphi_{56} :\equiv \forall y\,\forall z\,((1 - r) \cdot {\textstyle\sum_{k=1}^{S(n)} {{r}^{k}}} = y \to y = z \to (1 - r) \cdot {\textstyle\sum_{k=1}^{S(n)} {{r}^{k}}} = z)$
\item $\varphi_{57} :\equiv \forall y\,\forall z\,((1 - r) \cdot {\textstyle\sum_{k=1}^{1_{\omega}} {{r}^{k}}} = y \to y = z \to (1 - r) \cdot {\textstyle\sum_{k=1}^{1_{\omega}} {{r}^{k}}} = z)$
\item $\varphi_{58} :\equiv \mathrm{Seq}({(r)^{\bullet}}) \land {{r}^{1_{\omega}}} = r \land \forall n \in \mathbb{N}\,({{r}^{S(n)}} = {{r}^{n}} \cdot r)$
\item $\varphi_{59} :\equiv n \in B_{51} \to S(n) \in \mathbb{N} \to (1 - r) \cdot {\textstyle\sum_{k=1}^{S(n)} {{r}^{k}}} = r - {{r}^{S(S(n))}}$
\item $\varphi_{60} :\equiv r = {{r}^{1_{\omega}}} \to r - r \cdot r = r - r \cdot r \to r - r \cdot r = r - {{r}^{1_{\omega}}} \cdot r$
\item $\varphi_{61} :\equiv S(n) \in \mathbb{N} \to n \neq \emptyset \to S(n) \in \mathbb{N} \to {{r}^{S(S(n))}} = {{r}^{S(n)}} \cdot r$
\item $\varphi_{62} :\equiv (1 - r) \cdot {{r}^{1_{\omega}}} = (1 - r) \cdot r \to (1 - r) \cdot {\textstyle\sum_{k=1}^{1_{\omega}} {{r}^{k}}} = (1 - r) \cdot r$
\item $\varphi_{63} :\equiv S(n) \in \mathbb{N} \to n = \emptyset \to (1 - r) \cdot {\textstyle\sum_{k=1}^{S(n)} {{r}^{k}}} = r - {{r}^{S(S(n))}}$
\item $\varphi_{64} :\equiv S(n) \in \mathbb{N} \to n \neq \emptyset \to (1 - r) \cdot {\textstyle\sum_{k=1}^{S(n)} {{r}^{k}}} = r - {{r}^{S(S(n))}}$
\item $\varphi_{65} :\equiv S(n) \in \mathbb{N} \to n \neq \emptyset \to (1 - r) \cdot {\textstyle\sum_{k=1}^{S(n)} {{r}^{k}}} = r - {{r}^{S(n)}} \cdot r$
\item $\varphi_{66} :\equiv n \in B_{51} \leftrightarrow n \in \omega \land (n \in \mathbb{N} \to (1 - r) \cdot {\textstyle\sum_{k=1}^{n} {{r}^{k}}} = r - {{r}^{S(n)}})$
\item $\varphi_{67} :\equiv n \in B_{51} \to n \in \omega \land (n \in \mathbb{N} \to (1 - r) \cdot {\textstyle\sum_{k=1}^{n} {{r}^{k}}} = r - {{r}^{S(n)}})$
\item $\varphi_{68} :\equiv \forall y\,(1 \in \mathbb{R} \to y \in \mathbb{R} \to 1 - y \in \mathbb{R}) \to 1 \in \mathbb{R} \to r \in \mathbb{R} \to 1 - r \in \mathbb{R}$
\item $\varphi_{69} :\equiv r \in \mathbb{R} \to n \in \omega \to n \in \mathbb{N} \to (1 - r) \cdot {\textstyle\sum_{k=1}^{n} {{r}^{k}}} = r - {{r}^{S(n)}}$
\item $\varphi_{70} :\equiv (1 - r) \cdot {\textstyle\sum_{k=1}^{n} {{r}^{k}}} + (1 - r) \cdot {{r}^{S(n)}} = w + (1 - r) \cdot {{r}^{S(n)}}$
\item $\varphi_{71} :\equiv (1 - r) \cdot {\textstyle\sum_{k=1}^{n} {{r}^{k}}} + (1 - r) \cdot {{r}^{S(n)}} = z + (1 - r) \cdot {{r}^{S(n)}}$
\item $\varphi_{72} :\equiv n \in B_{51} \to S(n) \in \mathbb{N} \to n \neq \emptyset \to \forall n \in \mathbb{N}\,({{r}^{S(n)}} = {{r}^{n}} \cdot r)$
\item $\varphi_{73} :\equiv (1 - r) \cdot {\textstyle\sum_{k=1}^{S(n)} {{r}^{k}}} = (1 - r) \cdot {\textstyle\sum_{k=1}^{n} {{r}^{k}}} + (1 - r) \cdot {{r}^{S(n)}}$
\item $\varphi_{74} :\equiv (1 - r) \cdot {\textstyle\sum_{k=1}^{S(n)} {{r}^{k}}} = (1 - r) \cdot z \to (1 - r) \cdot {\textstyle\sum_{k=1}^{S(n)} {{r}^{k}}} = (1 - r) \cdot w$
\item $\varphi_{75} :\equiv (1 - r) \cdot {\textstyle\sum_{k=1}^{1_{\omega}} {{r}^{k}}} = (1 - r) \cdot z \to (1 - r) \cdot {\textstyle\sum_{k=1}^{1_{\omega}} {{r}^{k}}} = (1 - r) \cdot w$
\item $\varphi_{76} :\equiv \forall x \in \mathbb{R}\,(x \cdot 1 = x) \to {{r}^{S(n)}} \in \mathbb{R} \to {{r}^{S(n)}} \cdot 1 = {{r}^{S(n)}}$
\item $\varphi_{77} :\equiv n \in B_{51} \to S(n) \in \mathbb{N} \to n \neq \emptyset \to r - {{r}^{S(n)}} \cdot r = r - {{r}^{S(S(n))}}$
\item $\varphi_{78} :\equiv n \in B_{51} \to S(n) \in \mathbb{N} \to n \neq \emptyset \to (1 - r) \cdot {\textstyle\sum_{k=1}^{n} {{r}^{k}}} = r - {{r}^{S(n)}}$
\item $\varphi_{79} :\equiv r \in \mathbb{R} \to \forall n \in \omega\,(n \in \mathbb{N} \to (1 - r) \cdot {\textstyle\sum_{k=1}^{n} {{r}^{k}}} = r - {{r}^{S(n)}})$
\item $\varphi_{80} :\equiv (1 - r) \cdot {{r}^{1_{\omega}}} = (1 - r) \cdot {{r}^{1_{\omega}}} \to (1 - r) \cdot {{r}^{1_{\omega}}} = (1 - r) \cdot r$
\item $\varphi_{81} :\equiv r \in \mathbb{R} \land \neg \forall n \in \omega\,(n \in \mathbb{N} \to (1 - r) \cdot {\textstyle\sum_{k=1}^{n} {{r}^{k}}} = r - {{r}^{S(n)}})$
\item $\varphi_{82} :\equiv \forall y\,\forall z\,(1 - r \in \mathbb{R} \to y \in \mathbb{R} \to z \in \mathbb{R} \to (1 - r) \cdot (y + z) = (1 - r) \cdot y + (1 - r) \cdot z)$
\item $\varphi_{83} :\equiv n \neq \emptyset \to (1 - r) \cdot {\textstyle\sum_{k=1}^{S(n)} {{r}^{k}}} = (1 - r) \cdot ({\textstyle\sum_{k=1}^{n} {{r}^{k}}} + {{r}^{S(n)}})$
\item $\varphi_{84} :\equiv \forall x\,\forall y\,(x \in \mathbb{R} \to y \in \mathbb{R} \to x - y \in \mathbb{R}) \to \forall y\,(1 \in \mathbb{R} \to y \in \mathbb{R} \to 1 - y \in \mathbb{R})$
\item $\varphi_{85} :\equiv r \in \mathbb{R} \to \emptyset \in \mathbb{N} \to (1 - r) \cdot {\textstyle\sum_{k=1}^{\emptyset} {{r}^{k}}} = r - {{r}^{S(\emptyset)}}$
\item $\varphi_{86} :\equiv (1 - r) \cdot {\textstyle\sum_{k=1}^{S(n)} {{r}^{k}}} = r - {{r}^{S(n)}} + ({{r}^{S(n)}} - {{r}^{S(n)}} \cdot r)$
\item $\varphi_{87} :\equiv n \neq \emptyset \to (1 - r) \cdot {\textstyle\sum_{k=1}^{S(n)} {{r}^{k}}} = r - {{r}^{S(n)}} + (1 - r) \cdot {{r}^{S(n)}}$
\item $\varphi_{88} :\equiv n \neq \emptyset \to (1 - r) \cdot {\textstyle\sum_{k=1}^{S(n)} {{r}^{k}}} = r - {{r}^{S(n)}} + {{r}^{S(n)}} \cdot (1 - r)$
\item $\varphi_{89} :\equiv \forall a\,\forall n\,(\mathrm{Seq}(a) \to n \in \mathbb{N} \to a_{n} \in \mathbb{R}) \to \forall n\,(\mathrm{Seq}({(r)^{\bullet}}) \to n \in \mathbb{N} \to {{r}^{n}} \in \mathbb{R})$
\item $\varphi_{90} :\equiv 1 - r \in \mathbb{R} \to {{r}^{S(n)}} \in \mathbb{R} \to (1 - r) \cdot {{r}^{S(n)}} = {{r}^{S(n)}} \cdot (1 - r)$
\item $\varphi_{91} :\equiv 1 \in \mathbb{R} \to r \in \mathbb{R} \to {{r}^{S(n)}} \cdot 1 - {{r}^{S(n)}} \cdot r = {{r}^{S(n)}} \cdot (1 - r)$
\item $\varphi_{92} :\equiv r - r \cdot r = r - {{r}^{1_{\omega}}} \cdot r \to (1 - r) \cdot {\textstyle\sum_{k=1}^{1_{\omega}} {{r}^{k}}} = r - {{r}^{1_{\omega}}} \cdot r$
\item $\varphi_{93} :\equiv \forall y\,({{r}^{S(n)}} \cdot 1 - {{r}^{S(n)}} \cdot r = y \to y = {{r}^{S(n)}} \cdot 1 - {{r}^{S(n)}} \cdot r)$
\item $\varphi_{94} :\equiv \forall z\,((1 - r) \cdot {\textstyle\sum_{k=1}^{1_{\omega}} {{r}^{k}}} = r - r \cdot r \to r - r \cdot r = z \to (1 - r) \cdot {\textstyle\sum_{k=1}^{1_{\omega}} {{r}^{k}}} = z)$
\item $\varphi_{95} :\equiv \forall z\,((1 - r) \cdot {\textstyle\sum_{k=1}^{1_{\omega}} {{r}^{k}}} = r \cdot (1 - r) \to r \cdot (1 - r) = z \to (1 - r) \cdot {\textstyle\sum_{k=1}^{1_{\omega}} {{r}^{k}}} = z)$
\item $\varphi_{96} :\equiv \forall z\,((1 - r) \cdot {\textstyle\sum_{k=1}^{1_{\omega}} {{r}^{k}}} = (1 - r) \cdot r \to (1 - r) \cdot r = z \to (1 - r) \cdot {\textstyle\sum_{k=1}^{1_{\omega}} {{r}^{k}}} = z)$
\item $\varphi_{97} :\equiv \forall A_{52}\,\exists B_{51}\,\forall n\,(n \in B_{51} \leftrightarrow n \in A_{52} \land (n \in \mathbb{N} \to (1 - r) \cdot {\textstyle\sum_{k=1}^{n} {{r}^{k}}} = r - {{r}^{S(n)}}))$
\item $\varphi_{98} :\equiv r - {{r}^{S(n)}} + {{r}^{S(n)}} \cdot (1 - r) = r - {{r}^{S(n)}} + {{r}^{S(n)}} \cdot (1 - r)$
\item $\varphi_{99} :\equiv r - {{r}^{S(n)}} + (1 - r) \cdot {{r}^{S(n)}} = r - {{r}^{S(n)}} + (1 - r) \cdot {{r}^{S(n)}}$
\item $\varphi_{100} :\equiv S(n) \in \mathbb{N} \to n \neq \emptyset \to {{r}^{S(n)}} \cdot 1 - {{r}^{S(n)}} \cdot r = {{r}^{S(n)}} \cdot (1 - r)$
\item $\varphi_{101} :\equiv S(n) \in \mathbb{N} \to n \neq \emptyset \to {{r}^{S(n)}} \cdot (1 - r) = {{r}^{S(n)}} \cdot 1 - {{r}^{S(n)}} \cdot r$
\item $\varphi_{102} :\equiv r - {{r}^{S(n)}} + (1 - r) \cdot {{r}^{S(n)}} = r - {{r}^{S(n)}} + {{r}^{S(n)}} \cdot (1 - r)$
\item $\varphi_{103} :\equiv (1 - r) \cdot {\textstyle\sum_{k=1}^{S(n)} {{r}^{k}}} = r - {{r}^{S(n)}} + ({{r}^{S(n)}} \cdot 1 - {{r}^{S(n)}} \cdot r)$
\item $\varphi_{104} :\equiv S(n) \in \mathbb{N} \to n \neq \emptyset \to n \in \mathbb{N} \to {\textstyle\sum_{k=1}^{S(n)} {{r}^{k}}} = {\textstyle\sum_{k=1}^{n} {{r}^{k}}} + {{r}^{S(n)}}$
\item $\varphi_{105} :\equiv n \in B_{51} \to S(n) \in \mathbb{N} \to n \neq \emptyset \to {\textstyle\sum_{k=1}^{S(n)} {{r}^{k}}} = {\textstyle\sum_{k=1}^{n} {{r}^{k}}} + {{r}^{S(n)}}$
\item $\varphi_{106} :\equiv \forall n\,(n \in \mathbb{N} \leftrightarrow n \in \omega \land n \neq \emptyset) \to (\emptyset \in \mathbb{N} \leftrightarrow \emptyset \in \omega \land \emptyset \neq \emptyset)$
\item $\varphi_{107} :\equiv (1 - r) \cdot {\textstyle\sum_{k=1}^{z} {{r}^{k}}} = r - {{r}^{S(z)}} \to (1 - r) \cdot {\textstyle\sum_{k=1}^{w} {{r}^{k}}} = r - {{r}^{S(w)}}$
\item $\varphi_{108} :\equiv S(n) \in B_{51} \leftrightarrow S(n) \in \omega \land (S(n) \in \mathbb{N} \to (1 - r) \cdot {\textstyle\sum_{k=1}^{S(n)} {{r}^{k}}} = r - {{r}^{S(S(n))}})$
\item $\varphi_{109} :\equiv (r \in \mathbb{R} \to {{r}^{1_{\omega}}} = r) \to ({{r}^{1_{\omega}}} = r \to r = {{r}^{1_{\omega}}}) \to r \in \mathbb{R} \to r = {{r}^{1_{\omega}}}$
\item $\varphi_{110} :\equiv r - {{r}^{S(n)}} \cdot r = r - {{r}^{S(n)}} \cdot r \to r - {{r}^{S(n)}} \cdot r = r - {{r}^{S(S(n))}}$
\item $\varphi_{111} :\equiv n \neq \emptyset \to r - {{r}^{S(n)}} \cdot r = r - {{r}^{S(n)}} + ({{r}^{S(n)}} - {{r}^{S(n)}} \cdot r)$
\item $\varphi_{112} :\equiv n \neq \emptyset \to r - {{r}^{S(n)}} + ({{r}^{S(n)}} - {{r}^{S(n)}} \cdot r) = r - {{r}^{S(n)}} \cdot r$
\item $\varphi_{113} :\equiv S(n) \in \mathbb{N} \to n \neq \emptyset \to \forall n \in \mathbb{N}\,({\textstyle\sum_{k=1}^{S(n)} {{r}^{k}}} = {\textstyle\sum_{k=1}^{n} {{r}^{k}}} + {{r}^{S(n)}})$
\item $\varphi_{114} :\equiv r - {{r}^{1_{\omega}}} \cdot r = r - {{r}^{1_{\omega}}} \cdot r \to r - {{r}^{1_{\omega}}} \cdot r = r - {{r}^{S(1_{\omega})}}$
\item $\varphi_{115} :\equiv \forall n\,(\mathrm{Seq}({(r)^{\bullet}}) \to n \in \mathbb{N} \to {{r}^{n}} \in \mathbb{R}) \to \mathrm{Seq}({(r)^{\bullet}}) \to S(n) \in \mathbb{N} \to {{r}^{S(n)}} \in \mathbb{R}$
\item $\varphi_{116} :\equiv \forall n \in \mathbb{N}\,({{r}^{S(n)}} = {{r}^{n}} \cdot r) \to S(n) \in \mathbb{N} \to {{r}^{S(S(n))}} = {{r}^{S(n)}} \cdot r$
\item $\varphi_{117} :\equiv \forall n \in \mathbb{N}\,({{r}^{S(n)}} = {{r}^{n}} \cdot r) \to 1_{\omega} \in \mathbb{N} \to {{r}^{S(1_{\omega})}} = {{r}^{1_{\omega}}} \cdot r$
\item $\varphi_{118} :\equiv (1 - r) \cdot ({\textstyle\sum_{k=1}^{n} {{r}^{k}}} + {{r}^{S(n)}}) = (1 - r) \cdot {\textstyle\sum_{k=1}^{n} {{r}^{k}}} + (1 - r) \cdot {{r}^{S(n)}}$
\item $\varphi_{119} :\equiv (r \in \mathbb{R} \to \forall n\,(r \in \mathbb{R})) \to r \in \mathbb{R} \to \forall n \in \mathbb{N}\,((1 - r) \cdot {\textstyle\sum_{k=1}^{n} {{r}^{k}}} = r - {{r}^{S(n)}})$
\item $\varphi_{120} :\equiv r - {{r}^{S(n)}} \cdot r = r - {{r}^{S(S(n))}} \to (1 - r) \cdot {\textstyle\sum_{k=1}^{S(n)} {{r}^{k}}} = r - {{r}^{S(S(n))}}$
\item $\varphi_{121} :\equiv (1 - r) \cdot {\textstyle\sum_{k=1}^{n} {{r}^{k}}} + (1 - r) \cdot {{r}^{S(n)}} = r - {{r}^{S(n)}} + (1 - r) \cdot {{r}^{S(n)}}$
\item $\varphi_{122} :\equiv r - {{r}^{1_{\omega}}} \cdot r = r - {{r}^{S(1_{\omega})}} \to (1 - r) \cdot {\textstyle\sum_{k=1}^{1_{\omega}} {{r}^{k}}} = r - {{r}^{S(1_{\omega})}}$
\item $\varphi_{123} :\equiv \forall z\,((1 - r) \cdot {\textstyle\sum_{k=1}^{1_{\omega}} {{r}^{k}}} = r \cdot 1 - r \cdot r \to r \cdot 1 - r \cdot r = z \to (1 - r) \cdot {\textstyle\sum_{k=1}^{1_{\omega}} {{r}^{k}}} = z)$
\item $\varphi_{124} :\equiv (1 \in \mathbb{R} \to r \in \mathbb{R} \to 1 - r \in \mathbb{R}) \to 1 \in \mathbb{R} \to (r \in \mathbb{R} \to r \in \mathbb{R}) \to r \in \mathbb{R} \to 1 - r \in \mathbb{R}$
\item $\varphi_{125} :\equiv \forall z_{53}\,\forall A_{52}\,\exists B_{51}\,\forall n\,(n \in B_{51} \leftrightarrow n \in A_{52} \land (n \in \mathbb{N} \to (1 - z_{53}) \cdot {\textstyle\sum_{k=1}^{n} {{z_{53}}^{k}}} = z_{53} - {{z_{53}}^{S(n)}}))$
\item $\varphi_{126} :\equiv (r \in \mathbb{R} \to n \in \mathbb{N} \to n \in \mathbb{N}) \to r \in \mathbb{R} \to n \in \mathbb{N} \to (1 - r) \cdot {\textstyle\sum_{k=1}^{n} {{r}^{k}}} = r - {{r}^{S(n)}}$
\item $\varphi_{127} :\equiv (r \in \mathbb{R} \to n \in \mathbb{N} \to n \in \mathbb{N}) \to (n \in \mathbb{N} \leftrightarrow n \in \omega \land n \neq \emptyset) \to r \in \mathbb{R} \to n \in \mathbb{N} \to n \in \omega$
\item $\varphi_{128} :\equiv (1 - r) \cdot {\textstyle\sum_{k=1}^{1_{\omega}} {{r}^{k}}} = r - {{r}^{S(1_{\omega})}} \to (1 - r) \cdot {\textstyle\sum_{k=1}^{S(n)} {{r}^{k}}} = r - {{r}^{S(S(n))}}$
\item $\varphi_{129} :\equiv \forall w\,({{r}^{S(n)}} \cdot r = w \to r - {{r}^{S(n)}} \cdot r = r - {{r}^{S(n)}} \cdot r \to r - {{r}^{S(n)}} \cdot r = r - w)$
\item $\varphi_{130} :\equiv \forall w\,({{r}^{1_{\omega}}} \cdot r = w \to r - {{r}^{1_{\omega}}} \cdot r = r - {{r}^{1_{\omega}}} \cdot r \to r - {{r}^{1_{\omega}}} \cdot r = r - w)$
\item $\varphi_{131} :\equiv \forall w\,({{r}^{1_{\omega}}} = w \to (1 - r) \cdot {{r}^{1_{\omega}}} = (1 - r) \cdot {{r}^{1_{\omega}}} \to (1 - r) \cdot {{r}^{1_{\omega}}} = (1 - r) \cdot w)$
\item $\varphi_{132} :\equiv (1 - r) \cdot {\textstyle\sum_{k=1}^{n} {{r}^{k}}} + (1 - r) \cdot {{r}^{S(n)}} = (1 - r) \cdot {\textstyle\sum_{k=1}^{n} {{r}^{k}}} + (1 - r) \cdot {{r}^{S(n)}}$
\item $\varphi_{133} :\equiv {{r}^{S(n)}} \in \mathbb{R} \to r - {{r}^{S(n)}} \cdot r = r - {{r}^{S(n)}} + ({{r}^{S(n)}} - {{r}^{S(n)}} \cdot r)$
\item $\varphi_{134} :\equiv r - {{r}^{S(n)}} + ({{r}^{S(n)}} \cdot 1 - {{r}^{S(n)}} \cdot r) = r - {{r}^{S(n)}} + (w - {{r}^{S(n)}} \cdot r)$
\item $\varphi_{135} :\equiv r - {{r}^{S(n)}} + ({{r}^{S(n)}} \cdot 1 - {{r}^{S(n)}} \cdot r) = r - {{r}^{S(n)}} + (z - {{r}^{S(n)}} \cdot r)$
\item $\varphi_{136} :\equiv r - {{r}^{S(n)}} + {{r}^{S(n)}} \cdot (1 - r) = r - {{r}^{S(n)}} + ({{r}^{S(n)}} \cdot 1 - {{r}^{S(n)}} \cdot r)$
\item $\varphi_{137} :\equiv \forall w\,(1_{\omega} = w \to (1 - r) \cdot {\textstyle\sum_{k=1}^{1_{\omega}} {{r}^{k}}} = r - {{r}^{S(1_{\omega})}} \to (1 - r) \cdot {\textstyle\sum_{k=1}^{w} {{r}^{k}}} = r - {{r}^{S(w)}})$
\item $\varphi_{138} :\equiv \forall x\,(\mathrm{Ind}(x) \leftrightarrow \emptyset \in x \land \forall y \in x\,(S(y) \in x)) \to (\mathrm{Ind}(\omega) \leftrightarrow \emptyset \in \omega \land \forall y \in \omega\,(S(y) \in \omega))$
\item $\varphi_{139} :\equiv \forall a\,(\mathrm{Seq}(a) \to \mathrm{Seq}(\Sigma a) \land {\textstyle\sum_{k=1}^{1_{\omega}} a_{k}} = a_{1_{\omega}} \land \forall n \in \mathbb{N}\,({\textstyle\sum_{k=1}^{S(n)} a_{k}} = {\textstyle\sum_{k=1}^{n} a_{k}} + a_{S(n)}))$
\item $\varphi_{140} :\equiv \forall x\,\forall y\,(x \in \mathbb{R} \to y \in \mathbb{R} \to x \cdot y \in \mathbb{R}) \to \forall y\,({{r}^{S(n)}} \in \mathbb{R} \to y \in \mathbb{R} \to {{r}^{S(n)}} \cdot y \in \mathbb{R})$
\item $\varphi_{141} :\equiv \forall x\,(r \in \mathbb{R} \to {{r}^{S(n)}} \cdot r \in \mathbb{R} \to x \in \mathbb{R} \to r - {{r}^{S(n)}} \cdot r = r - x + (x - {{r}^{S(n)}} \cdot r))$
\item $\varphi_{142} :\equiv \emptyset \in B_{51} \leftrightarrow \emptyset \in \omega \land (\emptyset \in \mathbb{N} \to (1 - r) \cdot {\textstyle\sum_{k=1}^{\emptyset} {{r}^{k}}} = r - {{r}^{S(\emptyset)}})$
\item $\varphi_{143} :\equiv (1 - r) \cdot {\textstyle\sum_{k=1}^{1_{\omega}} {{r}^{k}}} = (1 - r) \cdot {\textstyle\sum_{k=1}^{1_{\omega}} {{r}^{k}}} \to (1 - r) \cdot {\textstyle\sum_{k=1}^{1_{\omega}} {{r}^{k}}} = (1 - r) \cdot {{r}^{1_{\omega}}}$
\item $\varphi_{144} :\equiv (1 - r) \cdot {\textstyle\sum_{k=1}^{1_{\omega}} {{r}^{k}}} = r \cdot 1 - r \cdot r \to r \cdot 1 - r \cdot r = r - r \cdot r \to (1 - r) \cdot {\textstyle\sum_{k=1}^{1_{\omega}} {{r}^{k}}} = r - r \cdot r$
\item $\varphi_{145} :\equiv \forall z\,((1 - r) \cdot {\textstyle\sum_{k=1}^{S(n)} {{r}^{k}}} = r - {{r}^{S(n)}} \cdot r \to r - {{r}^{S(n)}} \cdot r = z \to (1 - r) \cdot {\textstyle\sum_{k=1}^{S(n)} {{r}^{k}}} = z)$
\item $\varphi_{146} :\equiv \forall a\,({{r}^{S(n)}} \in \mathbb{R} \to 1 \in \mathbb{R} \to a \in \mathbb{R} \to {{r}^{S(n)}} \cdot 1 - {{r}^{S(n)}} \cdot a = {{r}^{S(n)}} \cdot (1 - a))$
\item $\varphi_{147} :\equiv \forall z\,((1 - r) \cdot {\textstyle\sum_{k=1}^{1_{\omega}} {{r}^{k}}} = r - {{r}^{1_{\omega}}} \cdot r \to r - {{r}^{1_{\omega}}} \cdot r = z \to (1 - r) \cdot {\textstyle\sum_{k=1}^{1_{\omega}} {{r}^{k}}} = z)$
\item $\varphi_{148} :\equiv \forall z\,((1 - r) \cdot {\textstyle\sum_{k=1}^{1_{\omega}} {{r}^{k}}} = (1 - r) \cdot {{r}^{1_{\omega}}} \to (1 - r) \cdot {{r}^{1_{\omega}}} = z \to (1 - r) \cdot {\textstyle\sum_{k=1}^{1_{\omega}} {{r}^{k}}} = z)$
\item $\varphi_{149} :\equiv \forall x\,\forall a\,({{r}^{S(n)}} \in \mathbb{R} \to x \in \mathbb{R} \to a \in \mathbb{R} \to {{r}^{S(n)}} \cdot x - {{r}^{S(n)}} \cdot a = {{r}^{S(n)}} \cdot (x - a))$
\item $\varphi_{150} :\equiv \forall y\,({{r}^{S(S(n))}} = y \to y = {{r}^{S(S(n))}}) \to {{r}^{S(S(n))}} = {{r}^{S(n)}} \cdot r \to {{r}^{S(n)}} \cdot r = {{r}^{S(S(n))}}$
\item $\varphi_{151} :\equiv \forall y\,({{r}^{S(1_{\omega})}} = y \to y = {{r}^{S(1_{\omega})}}) \to {{r}^{S(1_{\omega})}} = {{r}^{1_{\omega}}} \cdot r \to {{r}^{1_{\omega}}} \cdot r = {{r}^{S(1_{\omega})}}$
\item $\varphi_{152} :\equiv r - {{r}^{S(n)}} + ({{r}^{S(n)}} \cdot 1 - {{r}^{S(n)}} \cdot r) = r - {{r}^{S(n)}} + ({{r}^{S(n)}} - {{r}^{S(n)}} \cdot r)$
\item $\varphi_{153} :\equiv \mathrm{Seq}(\Sigma {(r)^{\bullet}}) \land {\textstyle\sum_{k=1}^{1_{\omega}} {{r}^{k}}} = {{r}^{1_{\omega}}} \land \forall n \in \mathbb{N}\,({\textstyle\sum_{k=1}^{S(n)} {{r}^{k}}} = {\textstyle\sum_{k=1}^{n} {{r}^{k}}} + {{r}^{S(n)}})$
\item $\varphi_{154} :\equiv ({{r}^{S(1_{\omega})}} = {{r}^{1_{\omega}}} \cdot r \to {{r}^{1_{\omega}}} \cdot r = {{r}^{S(1_{\omega})}}) \to r \in \mathbb{R} \to {{r}^{1_{\omega}}} \cdot r = {{r}^{S(1_{\omega})}}$
\item $\varphi_{155} :\equiv \forall n\,(r \in \mathbb{R} \land \forall n\,(\varphi_{66})) \to \forall n \in B_{51}\,(S(n) \in B_{51})$
\item $\varphi_{156} :\equiv r - {{r}^{S(n)}} + ({{r}^{S(n)}} \cdot 1 - {{r}^{S(n)}} \cdot r) = r - {{r}^{S(n)}} + ({{r}^{S(n)}} \cdot 1 - {{r}^{S(n)}} \cdot r)$
\item $\varphi_{157} :\equiv (r \in \mathbb{R} \to 1_{\omega} \in \mathbb{N} \to {{r}^{S(1_{\omega})}} = {{r}^{1_{\omega}}} \cdot r) \to 1_{\omega} \in \mathbb{N} \to r \in \mathbb{R} \to {{r}^{S(1_{\omega})}} = {{r}^{1_{\omega}}} \cdot r$
\item $\varphi_{158} :\equiv \forall w\,({\textstyle\sum_{k=1}^{S(n)} {{r}^{k}}} = w \to (1 - r) \cdot {\textstyle\sum_{k=1}^{S(n)} {{r}^{k}}} = (1 - r) \cdot {\textstyle\sum_{k=1}^{S(n)} {{r}^{k}}} \to (1 - r) \cdot {\textstyle\sum_{k=1}^{S(n)} {{r}^{k}}} = (1 - r) \cdot w)$
\item $\varphi_{159} :\equiv \forall w\,({\textstyle\sum_{k=1}^{1_{\omega}} {{r}^{k}}} = w \to (1 - r) \cdot {\textstyle\sum_{k=1}^{1_{\omega}} {{r}^{k}}} = (1 - r) \cdot {\textstyle\sum_{k=1}^{1_{\omega}} {{r}^{k}}} \to (1 - r) \cdot {\textstyle\sum_{k=1}^{1_{\omega}} {{r}^{k}}} = (1 - r) \cdot w)$
\item $\varphi_{160} :\equiv \varphi_{34} \to r \in \mathbb{R} \to 1 \in \mathbb{R} \to r \in \mathbb{R} \to r \cdot 1 - r \cdot r = r \cdot (1 - r)$
\item $\varphi_{161} :\equiv \forall z\,(1 - r \in \mathbb{R} \to {\textstyle\sum_{k=1}^{n} {{r}^{k}}} \in \mathbb{R} \to z \in \mathbb{R} \to (1 - r) \cdot ({\textstyle\sum_{k=1}^{n} {{r}^{k}}} + z) = (1 - r) \cdot {\textstyle\sum_{k=1}^{n} {{r}^{k}}} + (1 - r) \cdot z)$
\item $\varphi_{162} :\equiv \forall n\,(\varphi_{66}) \to n \in B_{51} \to S(n) \in \mathbb{N} \to n = \emptyset \to S(n) = 1_{\omega}$
\item $\varphi_{163} :\equiv \forall n\,(\varphi_{66}) \to n \in B_{51} \to S(n) \in \mathbb{N} \to n = \emptyset \to 1_{\omega} = S(n)$
\item $\varphi_{164} :\equiv {{r}^{S(n)}} \cdot 1 - {{r}^{S(n)}} \cdot r = {{r}^{S(n)}} \cdot (1 - r) \to {{r}^{S(n)}} \cdot (1 - r) = {{r}^{S(n)}} \cdot 1 - {{r}^{S(n)}} \cdot r$
\item $\varphi_{165} :\equiv \forall y\,({{r}^{S(n)}} \in \mathbb{R} \to y \in \mathbb{R} \to {{r}^{S(n)}} \cdot y \in \mathbb{R}) \to {{r}^{S(n)}} \in \mathbb{R} \to r \in \mathbb{R} \to {{r}^{S(n)}} \cdot r \in \mathbb{R}$
\item $\varphi_{166} :\equiv (1 - r) \cdot {\textstyle\sum_{k=1}^{S(n)} {{r}^{k}}} = (1 - r) \cdot {\textstyle\sum_{k=1}^{S(n)} {{r}^{k}}} \to (1 - r) \cdot {\textstyle\sum_{k=1}^{S(n)} {{r}^{k}}} = (1 - r) \cdot ({\textstyle\sum_{k=1}^{n} {{r}^{k}}} + {{r}^{S(n)}})$
\item $\varphi_{167} :\equiv \forall n\,(\varphi_{66}) \to n \in B_{51} \to S(n) \in \mathbb{N} \to n \neq \emptyset \to S(n) \in \mathbb{N}$
\item $\varphi_{168} :\equiv (r \in \mathbb{R} \to r \cdot 1 - r \cdot r = r \cdot (1 - r)) \to \varphi_{19} \to r \in \mathbb{R} \to r \cdot (1 - r) = r \cdot 1 - r \cdot r$
\item $\varphi_{169} :\equiv \forall n \in \mathbb{N}\,({\textstyle\sum_{k=1}^{S(n)} {{r}^{k}}} = {\textstyle\sum_{k=1}^{n} {{r}^{k}}} + {{r}^{S(n)}}) \to n \in \mathbb{N} \to {\textstyle\sum_{k=1}^{S(n)} {{r}^{k}}} = {\textstyle\sum_{k=1}^{n} {{r}^{k}}} + {{r}^{S(n)}}$
\item $\varphi_{170} :\equiv r - {{r}^{S(n)}} + ({{r}^{S(n)}} - {{r}^{S(n)}} \cdot r) = r - {{r}^{S(n)}} \cdot r \to (1 - r) \cdot {\textstyle\sum_{k=1}^{S(n)} {{r}^{k}}} = r - {{r}^{S(n)}} \cdot r$
\item $\varphi_{171} :\equiv \forall n\,(\varphi_{66}) \to n \in B_{51} \to S(n) \in \mathbb{N} \to n = \emptyset \to S(n) = S(\emptyset)$
\item $\varphi_{172} :\equiv (1 - r \in \mathbb{R} \to r \in \mathbb{R} \to (1 - r) \cdot r = r \cdot (1 - r)) \to (r \in \mathbb{R} \to 1 - r \in \mathbb{R}) \to \varphi_{26}$
\item $\varphi_{173} :\equiv r - {{r}^{S(n)}} + {{r}^{S(n)}} \cdot (1 - r) = r - {{r}^{S(n)}} + z \to r - {{r}^{S(n)}} + {{r}^{S(n)}} \cdot (1 - r) = r - {{r}^{S(n)}} + w$
\item $\varphi_{174} :\equiv r - {{r}^{S(n)}} + (1 - r) \cdot {{r}^{S(n)}} = r - {{r}^{S(n)}} + z \to r - {{r}^{S(n)}} + (1 - r) \cdot {{r}^{S(n)}} = r - {{r}^{S(n)}} + w$
\item $\varphi_{175} :\equiv r \cdot 1 - r \cdot r = r \cdot 1 - r \cdot r \to \varphi_{40} \to r \in \mathbb{R} \to r \cdot 1 - r \cdot r = r - r \cdot r$
\item $\varphi_{176} :\equiv \forall n \in \omega\,(n \in \mathbb{N} \to (1 - r) \cdot {\textstyle\sum_{k=1}^{n} {{r}^{k}}} = r - {{r}^{S(n)}}) \to n \in \omega \to n \in \mathbb{N} \to (1 - r) \cdot {\textstyle\sum_{k=1}^{n} {{r}^{k}}} = r - {{r}^{S(n)}}$
\item $\varphi_{177} :\equiv r \in \mathbb{R} \to \forall n\,(\varphi_{66}) \to n \in B_{51} \to S(n) \in \mathbb{N} \to n \neq \emptyset \to n \in \mathbb{N}$
\item $\varphi_{178} :\equiv r \in \mathbb{R} \to \forall n\,(\varphi_{66}) \to n \in B_{51} \to S(n) \in \mathbb{N} \to n = \emptyset \to n = \emptyset$
\item $\varphi_{179} :\equiv r \in \mathbb{R} \to \forall n\,(\varphi_{66}) \to n \in B_{51} \to S(n) \in \mathbb{N} \to n \neq \emptyset \to r \in \mathbb{R}$
\item $\varphi_{180} :\equiv r \in \mathbb{R} \to \forall n\,(\varphi_{66}) \to n \in B_{51} \to S(n) \in \mathbb{N} \to n \neq \emptyset \to n \neq \emptyset$
\item $\varphi_{181} :\equiv \forall n\,(\varphi_{66}) \to n \in B_{51} \to S(n) \in \mathbb{N} \to n \neq \emptyset \to {{r}^{S(n)}} \in \mathbb{R}$
\item $\varphi_{182} :\equiv \forall n\,(\varphi_{66}) \to n \in B_{51} \to S(n) \in \mathbb{N} \to n \neq \emptyset \to {\textstyle\sum_{k=1}^{n} {{r}^{k}}} \in \mathbb{R}$
\item $\varphi_{183} :\equiv r \in \mathbb{R} \to \forall n\,(\varphi_{66}) \to n \in B_{51} \to S(n) \in \mathbb{N} \to n \neq \emptyset \to 1 - r \in \mathbb{R}$
\item $\varphi_{184} :\equiv r - r \cdot r = r - r \cdot r \to \varphi_{60} \to r \in \mathbb{R} \to r - r \cdot r = r - {{r}^{1_{\omega}}} \cdot r$
\item $\varphi_{185} :\equiv \forall n\,(\varphi_{66}) \to n \in B_{51} \to S(n) \in \mathbb{N} \to n \neq \emptyset \to {{r}^{S(n)}} \cdot r \in \mathbb{R}$
\item $\varphi_{186} :\equiv ({{r}^{1_{\omega}}} = r \to \varphi_{80}) \to r \in \mathbb{R} \to (1 - r) \cdot {{r}^{1_{\omega}}} = (1 - r) \cdot r$
\item $\varphi_{187} :\equiv \forall n\,(\varphi_{66}) \to n \in B_{51} \to n \in \mathbb{N} \to (1 - r) \cdot {\textstyle\sum_{k=1}^{n} {{r}^{k}}} = r - {{r}^{S(n)}}$
\item $\varphi_{188} :\equiv \forall n\,(\varphi_{66}) \to n \in \omega \to n \in \mathbb{N} \to (1 - r) \cdot {\textstyle\sum_{k=1}^{n} {{r}^{k}}} = r - {{r}^{S(n)}}$
\item $\varphi_{189} :\equiv \forall z\,((1 - r) \cdot {\textstyle\sum_{k=1}^{S(n)} {{r}^{k}}} = (1 - r) \cdot ({\textstyle\sum_{k=1}^{n} {{r}^{k}}} + {{r}^{S(n)}}) \to (1 - r) \cdot ({\textstyle\sum_{k=1}^{n} {{r}^{k}}} + {{r}^{S(n)}}) = z \to (1 - r) \cdot {\textstyle\sum_{k=1}^{S(n)} {{r}^{k}}} = z)$
\item $\varphi_{190} :\equiv \forall n\,(\varphi_{66}) \to n \in B_{51} \to S(n) \in \mathbb{N} \to n \neq \emptyset \to {{r}^{S(n)}} \cdot 1 = {{r}^{S(n)}}$
\item $\varphi_{191} :\equiv \forall n\,(\varphi_{66}) \to \forall n \in \omega\,(n \in \mathbb{N} \to (1 - r) \cdot {\textstyle\sum_{k=1}^{n} {{r}^{k}}} = r - {{r}^{S(n)}})$
\item $\varphi_{192} :\equiv \forall z\,((1 - r) \cdot {\textstyle\sum_{k=1}^{S(n)} {{r}^{k}}} = r - {{r}^{S(n)}} + {{r}^{S(n)}} \cdot (1 - r) \to r - {{r}^{S(n)}} + {{r}^{S(n)}} \cdot (1 - r) = z \to (1 - r) \cdot {\textstyle\sum_{k=1}^{S(n)} {{r}^{k}}} = z)$
\item $\varphi_{193} :\equiv \forall z\,((1 - r) \cdot {\textstyle\sum_{k=1}^{S(n)} {{r}^{k}}} = r - {{r}^{S(n)}} + (1 - r) \cdot {{r}^{S(n)}} \to r - {{r}^{S(n)}} + (1 - r) \cdot {{r}^{S(n)}} = z \to (1 - r) \cdot {\textstyle\sum_{k=1}^{S(n)} {{r}^{k}}} = z)$
\item $\varphi_{194} :\equiv ((1 - r) \cdot {\textstyle\sum_{k=1}^{1_{\omega}} {{r}^{k}}} = (1 - r) \cdot r \to \varphi_{33}) \to r \in \mathbb{R} \to (1 - r) \cdot {\textstyle\sum_{k=1}^{1_{\omega}} {{r}^{k}}} = r \cdot (1 - r)$
\item $\varphi_{195} :\equiv (n \in \omega \leftrightarrow \forall x\,(\mathrm{Ind}(x) \to n \in x)) \to r \in \mathbb{R} \to \forall n\,(\varphi_{66}) \to n \in \omega \to \forall x\,(\mathrm{Ind}(x) \to n \in x)$
\item $\varphi_{196} :\equiv (r \in \mathbb{R} \to 1 \in \mathbb{R} \to r \in \mathbb{R} \to r \cdot 1 - r \cdot r = r \cdot (1 - r)) \to (r \in \mathbb{R} \to r \in \mathbb{R}) \to 1 \in \mathbb{R} \to \varphi_{31}$
\item $\varphi_{197} :\equiv r - {{r}^{S(n)}} \cdot r = r - {{r}^{S(n)}} + ({{r}^{S(n)}} - {{r}^{S(n)}} \cdot r) \to r - {{r}^{S(n)}} + ({{r}^{S(n)}} - {{r}^{S(n)}} \cdot r) = r - {{r}^{S(n)}} \cdot r$
\item $\varphi_{198} :\equiv r \in \mathbb{R} \land \forall n\,(\varphi_{66}) \to n \in \omega \to n \in \mathbb{N} \to (1 - r) \cdot {\textstyle\sum_{k=1}^{n} {{r}^{k}}} = r - {{r}^{S(n)}}$
\item $\varphi_{199} :\equiv (\forall x\,(\mathrm{Ind}(x) \to n \in x) \to \mathrm{Ind}(B_{51}) \to n \in B_{51}) \to r \in \mathbb{R} \to \forall n\,(\varphi_{66}) \to n \in \omega \to \mathrm{Ind}(B_{51}) \to n \in B_{51}$
\item $\varphi_{200} :\equiv \forall n\,(r \in \mathbb{R} \land \forall n\,(\varphi_{66})) \to \forall n \in \omega\,(n \in \mathbb{N} \to (1 - r) \cdot {\textstyle\sum_{k=1}^{n} {{r}^{k}}} = r - {{r}^{S(n)}})$
\item $\varphi_{201} :\equiv ((1 - r) \cdot {\textstyle\sum_{k=1}^{1_{\omega}} {{r}^{k}}} = r \cdot (1 - r) \to \varphi_{49}) \to r \in \mathbb{R} \to (1 - r) \cdot {\textstyle\sum_{k=1}^{1_{\omega}} {{r}^{k}}} = r \cdot 1 - r \cdot r$
\item $\varphi_{202} :\equiv \forall z\,(\varphi_{73} \to (1 - r) \cdot {\textstyle\sum_{k=1}^{n} {{r}^{k}}} + (1 - r) \cdot {{r}^{S(n)}} = z \to (1 - r) \cdot {\textstyle\sum_{k=1}^{S(n)} {{r}^{k}}} = z)$
\item $\varphi_{203} :\equiv \forall z\,(\varphi_{86} \to r - {{r}^{S(n)}} + ({{r}^{S(n)}} - {{r}^{S(n)}} \cdot r) = z \to (1 - r) \cdot {\textstyle\sum_{k=1}^{S(n)} {{r}^{k}}} = z)$
\item $\varphi_{204} :\equiv ((1 - r) \cdot {\textstyle\sum_{k=1}^{1_{\omega}} {{r}^{k}}} = (1 - r) \cdot {{r}^{1_{\omega}}} \to \varphi_{62}) \to r \in \mathbb{R} \to (1 - r) \cdot {\textstyle\sum_{k=1}^{1_{\omega}} {{r}^{k}}} = (1 - r) \cdot r$
\item $\varphi_{205} :\equiv (n = \emptyset \to S(n) = S(n) \to S(n) = S(\emptyset)) \to r \in \mathbb{R} \to \varphi_{171}$
\item $\varphi_{206} :\equiv \forall w\,({{r}^{S(n)}} \cdot (1 - r) = w \to \varphi_{98} \to r - {{r}^{S(n)}} + {{r}^{S(n)}} \cdot (1 - r) = r - {{r}^{S(n)}} + w)$
\item $\varphi_{207} :\equiv \forall w\,((1 - r) \cdot {{r}^{S(n)}} = w \to \varphi_{99} \to r - {{r}^{S(n)}} + (1 - r) \cdot {{r}^{S(n)}} = r - {{r}^{S(n)}} + w)$
\item $\varphi_{208} :\equiv ({{r}^{1_{\omega}}} \cdot r = {{r}^{S(1_{\omega})}} \to \varphi_{114}) \to r \in \mathbb{R} \to r - {{r}^{1_{\omega}}} \cdot r = r - {{r}^{S(1_{\omega})}}$
\item $\varphi_{209} :\equiv (\forall y \in \omega\,(S(y) \in \omega) \to n \in \omega \to S(n) \in \omega) \to r \in \mathbb{R} \to \forall n\,(\varphi_{66}) \to n \in B_{51} \to n \in \omega \to S(n) \in \omega$
\item $\varphi_{210} :\equiv \forall z\,(\varphi_{103} \to r - {{r}^{S(n)}} + ({{r}^{S(n)}} \cdot 1 - {{r}^{S(n)}} \cdot r) = z \to (1 - r) \cdot {\textstyle\sum_{k=1}^{S(n)} {{r}^{k}}} = z)$
\item $\varphi_{211} :\equiv ((1 - r) \cdot {\textstyle\sum_{k=1}^{1_{\omega}} {{r}^{k}}} = r - r \cdot r \to \varphi_{92}) \to r \in \mathbb{R} \to (1 - r) \cdot {\textstyle\sum_{k=1}^{1_{\omega}} {{r}^{k}}} = r - {{r}^{1_{\omega}}} \cdot r$
\item $\varphi_{212} :\equiv (\mathrm{Ind}(\omega) \leftrightarrow \emptyset \in \omega \land \forall y \in \omega\,(S(y) \in \omega)) \to r \in \mathbb{R} \to \forall n\,(\varphi_{66}) \to \emptyset \in \omega \land \forall y \in \omega\,(S(y) \in \omega)$
\item $\varphi_{213} :\equiv (r \in \mathbb{R} \to r \cdot 1 - r \cdot r = r - r \cdot r) \to \varphi_{144} \to r \in \mathbb{R} \to (1 - r) \cdot {\textstyle\sum_{k=1}^{1_{\omega}} {{r}^{k}}} = r - r \cdot r$
\item $\varphi_{214} :\equiv ({\textstyle\sum_{k=1}^{1_{\omega}} {{r}^{k}}} = {{r}^{1_{\omega}}} \to \varphi_{143}) \to r \in \mathbb{R} \to (1 - r) \cdot {\textstyle\sum_{k=1}^{1_{\omega}} {{r}^{k}}} = (1 - r) \cdot {{r}^{1_{\omega}}}$
\item $\varphi_{215} :\equiv ((1 - r) \cdot {\textstyle\sum_{k=1}^{1_{\omega}} {{r}^{k}}} = r - {{r}^{1_{\omega}}} \cdot r \to \varphi_{122}) \to r \in \mathbb{R} \to (1 - r) \cdot {\textstyle\sum_{k=1}^{1_{\omega}} {{r}^{k}}} = r - {{r}^{S(1_{\omega})}}$
\item $\varphi_{216} :\equiv (r \in \mathbb{R} \to \forall n \in \mathbb{N}\,({{r}^{S(n)}} = {{r}^{n}} \cdot r)) \to \varphi_{117} \to r \in \mathbb{R} \to 1_{\omega} \in \mathbb{N} \to {{r}^{S(1_{\omega})}} = {{r}^{1_{\omega}}} \cdot r$
\item $\varphi_{217} :\equiv (r \in \mathbb{R} \to {{r}^{1_{\omega}}} = r) \to (1 - r) \cdot {{r}^{1_{\omega}}} = (1 - r) \cdot {{r}^{1_{\omega}}} \to \varphi_{186}$
\item $\varphi_{218} :\equiv (r \in \mathbb{R} \to \forall n\,(\varphi_{66}) \to n \in B_{51} \to S(n) \in B_{51}) \to r \in \mathbb{R} \land \forall n\,(\varphi_{66}) \to n \in B_{51} \to S(n) \in B_{51}$
\item $\varphi_{219} :\equiv (r \in \mathbb{R} \to \forall n\,(\varphi_{66}) \to n \in B_{51} \to n \in \omega) \to r \in \mathbb{R} \to \forall n\,(\varphi_{66}) \to n \in B_{51} \to S(n) \in \omega$
\item $\varphi_{220} :\equiv (1 - r) \cdot {\textstyle\sum_{k=1}^{S(n)} {{r}^{k}}} = r - {{r}^{S(n)}} + (1 - r) \cdot {{r}^{S(n)}} \to \varphi_{102} \to (1 - r) \cdot {\textstyle\sum_{k=1}^{S(n)} {{r}^{k}}} = r - {{r}^{S(n)}} + {{r}^{S(n)}} \cdot (1 - r)$
\item $\varphi_{221} :\equiv ({{r}^{S(S(n))}} = {{r}^{S(n)}} \cdot r \to {{r}^{S(n)}} \cdot r = {{r}^{S(S(n))}}) \to r \in \mathbb{R} \to \forall n\,(\varphi_{66}) \to \varphi_{54}$
\item $\varphi_{222} :\equiv \forall n\,(r \in \mathbb{R} \land \forall n\,(\varphi_{66}) \to n \in B_{51} \to S(n) \in B_{51}) \to \varphi_{155}$
\item $\varphi_{223} :\equiv (r \in \mathbb{R} \to (1 - r) \cdot {\textstyle\sum_{k=1}^{1_{\omega}} {{r}^{k}}} = (1 - r) \cdot r) \to (r \in \mathbb{R} \to (1 - r) \cdot r = r \cdot (1 - r)) \to \varphi_{194}$
\item $\varphi_{224} :\equiv (1 - r) \cdot {\textstyle\sum_{k=1}^{S(n)} {{r}^{k}}} = (1 - r) \cdot ({\textstyle\sum_{k=1}^{n} {{r}^{k}}} + {{r}^{S(n)}}) \to \varphi_{118} \to \varphi_{73}$
\item $\varphi_{225} :\equiv \varphi_{73} \to \varphi_{121} \to (1 - r) \cdot {\textstyle\sum_{k=1}^{S(n)} {{r}^{k}}} = r - {{r}^{S(n)}} + (1 - r) \cdot {{r}^{S(n)}}$
\item $\varphi_{226} :\equiv (r \in \mathbb{R} \to (1 - r) \cdot {\textstyle\sum_{k=1}^{1_{\omega}} {{r}^{k}}} = r \cdot (1 - r)) \to (r \in \mathbb{R} \to r \cdot (1 - r) = r \cdot 1 - r \cdot r) \to \varphi_{201}$
\item $\varphi_{227} :\equiv (1 - r) \cdot {\textstyle\sum_{k=1}^{S(n)} {{r}^{k}}} = r - {{r}^{S(n)}} + {{r}^{S(n)}} \cdot (1 - r) \to \varphi_{136} \to \varphi_{103}$
\item $\varphi_{228} :\equiv (r \in \mathbb{R} \to \forall n\,(\varphi_{66}) \to \emptyset \in \omega \land \forall y \in \omega\,(S(y) \in \omega)) \to r \in \mathbb{R} \to \forall n\,(\varphi_{66}) \to n \in B_{51} \to \forall y \in \omega\,(S(y) \in \omega)$
\item $\varphi_{229} :\equiv (r \in \mathbb{R} \to {{r}^{1_{\omega}}} \cdot r = {{r}^{S(1_{\omega})}}) \to r - {{r}^{1_{\omega}}} \cdot r = r - {{r}^{1_{\omega}}} \cdot r \to \varphi_{208}$
\item $\varphi_{230} :\equiv (\mathrm{Seq}({(r)^{\bullet}}) \to \varphi_{153}) \to (r \in \mathbb{R} \to \mathrm{Seq}({(r)^{\bullet}})) \to r \in \mathbb{R} \to \varphi_{153}$
\item $\varphi_{231} :\equiv (r \in \mathbb{R} \to \forall n\,(\varphi_{66}) \to \varphi_{59}) \to r \in \mathbb{R} \to \forall n\,(\varphi_{66}) \to n \in B_{51} \to S(n) \in B_{51}$
\item $\varphi_{232} :\equiv (r \in \mathbb{R} \to \forall n\,(\varphi_{66}) \to n \in B_{51} \to n \in B_{51}) \to r \in \mathbb{R} \to \forall n\,(\varphi_{66}) \to \varphi_{67}$
\item $\varphi_{233} :\equiv (r \in \mathbb{R} \to \forall n\,(\varphi_{66}) \to \varphi_{67}) \to r \in \mathbb{R} \to \forall n\,(\varphi_{66}) \to n \in B_{51} \to n \in \omega$
\item $\varphi_{234} :\equiv (r \in \mathbb{R} \to \forall n\,(\varphi_{66}) \to n \in \omega \to \forall x\,(\mathrm{Ind}(x) \to n \in x)) \to \varphi_{199}$
\item $\varphi_{235} :\equiv (r \in \mathbb{R} \to (1 - r) \cdot {\textstyle\sum_{k=1}^{1_{\omega}} {{r}^{k}}} = (1 - r) \cdot {{r}^{1_{\omega}}}) \to (r \in \mathbb{R} \to (1 - r) \cdot {{r}^{1_{\omega}}} = (1 - r) \cdot r) \to \varphi_{204}$
\item $\varphi_{236} :\equiv (r \in \mathbb{R} \to (1 - r) \cdot {\textstyle\sum_{k=1}^{1_{\omega}} {{r}^{k}}} = r - r \cdot r) \to (r \in \mathbb{R} \to r - r \cdot r = r - {{r}^{1_{\omega}}} \cdot r) \to \varphi_{211}$
\item $\varphi_{237} :\equiv (1_{\omega} = S(n) \to \varphi_{128}) \to r \in \mathbb{R} \to \forall n\,(\varphi_{66}) \to n \in B_{51} \to \varphi_{63}$
\item $\varphi_{238} :\equiv (r \in \mathbb{R} \to \varphi_{162}) \to (S(n) = 1_{\omega} \to 1_{\omega} = S(n)) \to r \in \mathbb{R} \to \varphi_{163}$
\item $\varphi_{239} :\equiv (r \in \mathbb{R} \land \forall n\,(\varphi_{66}) \to \forall n\,(r \in \mathbb{R} \land \forall n\,(\varphi_{66}))) \to r \in \mathbb{R} \to \forall n\,(\varphi_{66}) \to \forall n \in B_{51}\,(S(n) \in B_{51})$
\item $\varphi_{240} :\equiv (r \in \mathbb{R} \to {\textstyle\sum_{k=1}^{1_{\omega}} {{r}^{k}}} = {{r}^{1_{\omega}}}) \to (1 - r) \cdot {\textstyle\sum_{k=1}^{1_{\omega}} {{r}^{k}}} = (1 - r) \cdot {\textstyle\sum_{k=1}^{1_{\omega}} {{r}^{k}}} \to \varphi_{214}$
\item $\varphi_{241} :\equiv (r \in \mathbb{R} \to \forall n\,(\varphi_{66}) \to n \in B_{51} \to \forall y \in \omega\,(S(y) \in \omega)) \to \varphi_{209}$
\item $\varphi_{242} :\equiv (r \in \mathbb{R} \to \forall n\,(\varphi_{66}) \to \forall n \in B_{51}\,(S(n) \in B_{51})) \to \varphi_{12} \to \varphi_{11} \to r \in \mathbb{R} \to \forall n\,(\varphi_{66}) \to \mathrm{Ind}(B_{51})$
\item $\varphi_{243} :\equiv ((1 - r) \cdot {\textstyle\sum_{k=1}^{S(n)} {{r}^{k}}} = r - {{r}^{S(n)}} \cdot r \to \varphi_{120}) \to r \in \mathbb{R} \to \forall n\,(\varphi_{66}) \to n \in B_{51} \to \varphi_{64}$
\item $\varphi_{244} :\equiv (r \in \mathbb{R} \to (1 - r) \cdot {\textstyle\sum_{k=1}^{1_{\omega}} {{r}^{k}}} = r - {{r}^{1_{\omega}}} \cdot r) \to (r \in \mathbb{R} \to r - {{r}^{1_{\omega}}} \cdot r = r - {{r}^{S(1_{\omega})}}) \to \varphi_{215}$
\item $\varphi_{245} :\equiv (r \in \mathbb{R} \to \forall n\,(\varphi_{66}) \to \emptyset \in \omega \land \forall y \in \omega\,(S(y) \in \omega)) \to \varphi_{85} \to r \in \mathbb{R} \to \forall n\,(\varphi_{66}) \to \emptyset \in B_{51}$
\item $\varphi_{246} :\equiv \varphi_{179} \to r \in \mathbb{R} \to \forall n\,(\varphi_{66}) \to n \in B_{51} \to S(n) \in \mathbb{N} \to n \neq \emptyset \to \varphi_{58}$
\item $\varphi_{247} :\equiv (r \in \mathbb{R} \to \forall n\,(\varphi_{66}) \to n \in B_{51} \to S(n) \in \mathbb{N} \to n \neq \emptyset \to \varphi_{58}) \to r \in \mathbb{R} \to \forall n\,(\varphi_{66}) \to \varphi_{5}$
\item $\varphi_{248} :\equiv \forall B_{51}\,(\varphi_{81} \to \neg \forall n\,(\varphi_{66})) \to \forall B_{51}\,(\varphi_{81}) \to \forall B_{51}\,(\neg \forall n\,(\varphi_{66}))$
\item $\varphi_{249} :\equiv (r \in \mathbb{R} \to \forall n\,(\varphi_{66}) \to n \in B_{51} \to \varphi_{64}) \to r \in \mathbb{R} \to \forall n\,(\varphi_{66}) \to \varphi_{59}$
\item $\varphi_{250} :\equiv r - {{r}^{S(n)}} \cdot r = r - {{r}^{S(n)}} \cdot r \to ({{r}^{S(n)}} \cdot r = {{r}^{S(S(n))}} \to \varphi_{110}) \to r \in \mathbb{R} \to \forall n\,(\varphi_{66}) \to \varphi_{77}$
\item $\varphi_{251} :\equiv (r \in \mathbb{R} \land \forall n\,(\varphi_{66}) \to \forall n\,(r \in \mathbb{R} \land \forall n\,(\varphi_{66}))) \to r \in \mathbb{R} \to \varphi_{191}$
\item $\varphi_{252} :\equiv (r \in \mathbb{R} \to \varphi_{181}) \to r \in \mathbb{R} \to \forall n\,(\varphi_{66}) \to n \in B_{51} \to S(n) \in \mathbb{N} \to \varphi_{111}$
\item $\varphi_{253} :\equiv (r \in \mathbb{R} \to \forall n\,(\varphi_{66}) \to n \in B_{51} \to n \in \omega \to S(n) \in \omega) \to \varphi_{219}$
\item $\varphi_{254} :\equiv ({{r}^{S(n)}} \in \mathbb{R} \to {{r}^{S(n)}} \cdot 1 = {{r}^{S(n)}}) \to (r \in \mathbb{R} \to \varphi_{181}) \to r \in \mathbb{R} \to \varphi_{190}$
\item $\varphi_{255} :\equiv (r \in \mathbb{R} \to \forall n\,(\varphi_{66}) \to n \in B_{51} \to S(n) \in \mathbb{N} \to n \neq \emptyset \to \varphi_{58}) \to r \in \mathbb{R} \to \forall n\,(\varphi_{66}) \to \varphi_{72}$
\item $\varphi_{256} :\equiv (r \in \mathbb{R} \to \varphi_{181}) \to r \in \mathbb{R} \to \forall n\,(\varphi_{66}) \to n \in B_{51} \to S(n) \in \mathbb{N} \to n \neq \emptyset \to \varphi_{118}$
\item $\varphi_{257} :\equiv \forall z\,\forall w\,(z = w \to \varphi_{71} \to \varphi_{70}) \to \forall w\,((1 - r) \cdot {\textstyle\sum_{k=1}^{n} {{r}^{k}}} = w \to \varphi_{132} \to \varphi_{70})$
\item $\varphi_{258} :\equiv ({\textstyle\sum_{k=1}^{S(n)} {{r}^{k}}} = {\textstyle\sum_{k=1}^{n} {{r}^{k}}} + {{r}^{S(n)}} \to \varphi_{166}) \to r \in \mathbb{R} \to \forall n\,(\varphi_{66}) \to n \in B_{51} \to S(n) \in \mathbb{N} \to \varphi_{83}$
\item $\varphi_{259} :\equiv (r \in \mathbb{R} \to \forall n\,(\varphi_{66}) \to \varphi_{5}) \to r \in \mathbb{R} \to \forall n\,(\varphi_{66}) \to n \in B_{51} \to S(n) \in \mathbb{N} \to n \neq \emptyset \to \varphi_{153}$
\item $\varphi_{260} :\equiv (r \in \mathbb{R} \to \forall n\,(\varphi_{66}) \to n \in B_{51} \to S(n) \in \mathbb{N} \to n \neq \emptyset \to \varphi_{153}) \to r \in \mathbb{R} \to \forall n\,(\varphi_{66}) \to \varphi_{15}$
\item $\varphi_{261} :\equiv (n \in \mathbb{N} \leftrightarrow n \in \omega \land n \neq \emptyset) \to (r \in \mathbb{R} \to \forall n\,(\varphi_{66}) \to n \in B_{51} \to n \in \omega) \to \varphi_{180} \to \varphi_{177}$
\item $\varphi_{262} :\equiv (r \in \mathbb{R} \to \forall n\,(\varphi_{66}) \to \mathrm{Ind}(B_{51})) \to (r \in \mathbb{R} \to \forall n\,(\varphi_{66}) \to \varphi_{66}) \to r \in \mathbb{R} \to \varphi_{188}$
\item $\varphi_{263} :\equiv ((1 - r) \cdot {{r}^{S(n)}} = {{r}^{S(n)}} \cdot (1 - r) \to \varphi_{99} \to \varphi_{102}) \to r \in \mathbb{R} \to \forall n\,(\varphi_{66}) \to n \in B_{51} \to S(n) \in \mathbb{N} \to n \neq \emptyset \to \varphi_{102}$
\item $\varphi_{264} :\equiv (r \in \mathbb{R} \to \forall n\,(\varphi_{66}) \to n \in B_{51} \to S(n) \in \mathbb{N} \to n \neq \emptyset \to \varphi_{153}) \to r \in \mathbb{R} \to \forall n\,(\varphi_{66}) \to n \in B_{51} \to \varphi_{113}$
\item $\varphi_{265} :\equiv (r \in \mathbb{R} \to \forall n\,(\varphi_{66}) \to \varphi_{72}) \to \varphi_{116} \to r \in \mathbb{R} \to \forall n\,(\varphi_{66}) \to n \in B_{51} \to \varphi_{61}$
\item $\varphi_{266} :\equiv \varphi_{224} \to r \in \mathbb{R} \to \forall n\,(\varphi_{66}) \to n \in B_{51} \to S(n) \in \mathbb{N} \to n \neq \emptyset \to \varphi_{73}$
\item $\varphi_{267} :\equiv \forall w\,((1 - r) \cdot {\textstyle\sum_{k=1}^{n} {{r}^{k}}} = w \to \varphi_{132} \to \varphi_{70}) \to (1 - r) \cdot {\textstyle\sum_{k=1}^{n} {{r}^{k}}} = r - {{r}^{S(n)}} \to \varphi_{132} \to \varphi_{121}$
\item $\varphi_{268} :\equiv ((1 - r) \cdot {\textstyle\sum_{k=1}^{n} {{r}^{k}}} = r - {{r}^{S(n)}} \to \varphi_{132} \to \varphi_{121}) \to r \in \mathbb{R} \to \forall n\,(\varphi_{66}) \to n \in B_{51} \to S(n) \in \mathbb{N} \to n \neq \emptyset \to \varphi_{121}$
\item $\varphi_{269} :\equiv (\mathrm{Seq}({(r)^{\bullet}}) \to S(n) \in \mathbb{N} \to {{r}^{S(n)}} \in \mathbb{R}) \to (r \in \mathbb{R} \to \forall n\,(\varphi_{66}) \to \varphi_{5}) \to (r \in \mathbb{R} \to \varphi_{167}) \to r \in \mathbb{R} \to \varphi_{181}$
\item $\varphi_{270} :\equiv (r \in \mathbb{R} \to \forall n\,(\varphi_{66}) \to n \in B_{51} \to \varphi_{61}) \to (r \in \mathbb{R} \to \varphi_{167}) \to r \in \mathbb{R} \to \forall n\,(\varphi_{66}) \to \varphi_{53}$
\item $\varphi_{271} :\equiv \varphi_{227} \to r \in \mathbb{R} \to \forall n\,(\varphi_{66}) \to n \in B_{51} \to S(n) \in \mathbb{N} \to n \neq \emptyset \to \varphi_{103}$
\item $\varphi_{272} :\equiv (r \in \mathbb{R} \to \varphi_{163}) \to (r \in \mathbb{R} \to (1 - r) \cdot {\textstyle\sum_{k=1}^{1_{\omega}} {{r}^{k}}} = r - {{r}^{S(1_{\omega})}}) \to \varphi_{237}$
\item $\varphi_{273} :\equiv (\mathrm{Seq}(\Sigma {(r)^{\bullet}}) \to n \in \mathbb{N} \to {\textstyle\sum_{k=1}^{n} {{r}^{k}}} \in \mathbb{R}) \to (r \in \mathbb{R} \to \forall n\,(\varphi_{66}) \to \varphi_{15}) \to \varphi_{177} \to r \in \mathbb{R} \to \varphi_{182}$
\item $\varphi_{274} :\equiv ({{r}^{S(n)}} \cdot (1 - r) = {{r}^{S(n)}} \cdot 1 - {{r}^{S(n)}} \cdot r \to \varphi_{98} \to \varphi_{136}) \to r \in \mathbb{R} \to \forall n\,(\varphi_{66}) \to n \in B_{51} \to S(n) \in \mathbb{N} \to n \neq \emptyset \to \varphi_{136}$
\item $\varphi_{275} :\equiv (r \in \mathbb{R} \to \varphi_{181}) \to 1 \in \mathbb{R} \to \varphi_{179} \to r \in \mathbb{R} \to \forall n\,(\varphi_{66}) \to n \in B_{51} \to \varphi_{100}$
\item $\varphi_{276} :\equiv (r \in \mathbb{R} \to \forall n\,(\varphi_{66}) \to \forall n\,(\varphi_{66})) \to (\forall n\,(\varphi_{66}) \to \varphi_{66}) \to r \in \mathbb{R} \to \forall n\,(\varphi_{66}) \to \varphi_{66}$
\item $\varphi_{277} :\equiv (\varphi_{103} \to \varphi_{152} \to \varphi_{86}) \to r \in \mathbb{R} \to \forall n\,(\varphi_{66}) \to n \in B_{51} \to S(n) \in \mathbb{N} \to n \neq \emptyset \to \varphi_{86}$
\item $\varphi_{278} :\equiv ({{r}^{S(n)}} \in \mathbb{R} \to r \in \mathbb{R} \to {{r}^{S(n)}} \cdot r \in \mathbb{R}) \to (r \in \mathbb{R} \to \varphi_{181}) \to \varphi_{179} \to r \in \mathbb{R} \to \varphi_{185}$
\item $\varphi_{279} :\equiv (r \in \mathbb{R} \to \forall n\,(\varphi_{66}) \to n \in B_{51} \to \varphi_{100}) \to \varphi_{164} \to r \in \mathbb{R} \to \forall n\,(\varphi_{66}) \to n \in B_{51} \to \varphi_{101}$
\item $\varphi_{280} :\equiv (r \in \mathbb{R} \to \forall n\,(\varphi_{66}) \to n \in B_{51} \to \varphi_{104}) \to \varphi_{177} \to r \in \mathbb{R} \to \forall n\,(\varphi_{66}) \to \varphi_{105}$
\item $\varphi_{281} :\equiv ({{r}^{S(n)}} \cdot 1 = {{r}^{S(n)}} \to \varphi_{156} \to \varphi_{152}) \to r \in \mathbb{R} \to \forall n\,(\varphi_{66}) \to n \in B_{51} \to S(n) \in \mathbb{N} \to n \neq \emptyset \to \varphi_{152}$
\item $\varphi_{282} :\equiv (r \in \mathbb{R} \to \forall n\,(\varphi_{66}) \to n \in B_{51} \to \varphi_{113}) \to \varphi_{169} \to r \in \mathbb{R} \to \forall n\,(\varphi_{66}) \to n \in B_{51} \to \varphi_{104}$
\item $\varphi_{283} :\equiv \forall w\,({{r}^{S(n)}} \cdot 1 = w \to \varphi_{156} \to \varphi_{134}) \to {{r}^{S(n)}} \cdot 1 = {{r}^{S(n)}} \to \varphi_{156} \to \varphi_{152}$
\item $\varphi_{284} :\equiv (r \in \mathbb{R} \to \forall n\,(\varphi_{66}) \to \forall n\,(\varphi_{66})) \to (\forall n\,(\varphi_{66}) \to \varphi_{108}) \to r \in \mathbb{R} \to \forall n\,(\varphi_{66}) \to \varphi_{108}$
\item $\varphi_{285} :\equiv \varphi_{90} \to \varphi_{183} \to (r \in \mathbb{R} \to \varphi_{181}) \to r \in \mathbb{R} \to \forall n\,(\varphi_{66}) \to n \in B_{51} \to \varphi_{55}$
\item $\varphi_{286} :\equiv (r \in \mathbb{R} \to \forall n\,(\varphi_{66}) \to n \in \omega \to \mathrm{Ind}(B_{51}) \to n \in B_{51}) \to \varphi_{262}$
\item $\varphi_{287} :\equiv (r \in \mathbb{R} \to \forall n\,(\varphi_{66}) \to \forall n\,(\varphi_{66})) \to (\forall n\,(\varphi_{66}) \to \varphi_{142}) \to r \in \mathbb{R} \to \forall n\,(\varphi_{66}) \to \varphi_{142}$
\item $\varphi_{288} :\equiv (r \in \mathbb{R} \to \forall n\,(\varphi_{66}) \to n \in B_{51} \to S(n) \in \mathbb{N} \to \varphi_{111}) \to \varphi_{197} \to r \in \mathbb{R} \to \forall n\,(\varphi_{66}) \to n \in B_{51} \to S(n) \in \mathbb{N} \to \varphi_{112}$
\item $\varphi_{289} :\equiv (r \in \mathbb{R} \to \forall n\,(\varphi_{66}) \to \varphi_{108}) \to (r \in \mathbb{R} \to \forall n\,(\varphi_{66}) \to n \in B_{51} \to S(n) \in \omega) \to \varphi_{231}$
\item $\varphi_{290} :\equiv (r \in \mathbb{R} \to \forall n\,(\varphi_{66}) \to n \in B_{51} \to S(n) \in \mathbb{N} \to \varphi_{112}) \to (\varphi_{86} \to \varphi_{170}) \to r \in \mathbb{R} \to \forall n\,(\varphi_{66}) \to n \in B_{51} \to \varphi_{65}$
\item $\varphi_{291} :\equiv (r \in \mathbb{R} \to \forall n\,(\varphi_{66}) \to \varphi_{105}) \to (1 - r) \cdot {\textstyle\sum_{k=1}^{S(n)} {{r}^{k}}} = (1 - r) \cdot {\textstyle\sum_{k=1}^{S(n)} {{r}^{k}}} \to \varphi_{258}$
\item $\varphi_{292} :\equiv (r \in \mathbb{R} \to \forall n\,(\varphi_{66}) \to n \in B_{51} \to S(n) \in \mathbb{N} \to n \neq \emptyset \to \varphi_{102}) \to \varphi_{220} \to r \in \mathbb{R} \to \forall n\,(\varphi_{66}) \to n \in B_{51} \to S(n) \in \mathbb{N} \to \varphi_{88}$
\item $\varphi_{293} :\equiv \forall n\,(r \in \mathbb{R} \land \forall n\,(\varphi_{66}) \to n \in B_{51} \to S(n) \in B_{51}) \to \varphi_{222} \to \varphi_{239}$
\item $\varphi_{294} :\equiv (r \in \mathbb{R} \to \forall n\,(\varphi_{66}) \to n \in B_{51} \to S(n) \in \mathbb{N} \to n \neq \emptyset \to \varphi_{121}) \to \varphi_{225} \to r \in \mathbb{R} \to \forall n\,(\varphi_{66}) \to n \in B_{51} \to S(n) \in \mathbb{N} \to \varphi_{87}$
\item $\varphi_{295} :\equiv (r \in \mathbb{R} \to \forall n\,(\varphi_{66}) \to n \in B_{51} \to S(n) \in \mathbb{N} \to n \neq \emptyset \to \varphi_{118}) \to \varphi_{266}$
\item $\varphi_{296} :\equiv (r \in \mathbb{R} \to \forall n\,(\varphi_{66}) \to n \in B_{51} \to \varphi_{65}) \to (r \in \mathbb{R} \to \forall n\,(\varphi_{66}) \to \varphi_{77}) \to \varphi_{243}$
\item $\varphi_{297} :\equiv (r \in \mathbb{R} \to \forall n\,(\varphi_{66}) \to n \in B_{51} \to S(n) \in \mathbb{N} \to n \neq \emptyset \to \varphi_{136}) \to \varphi_{271}$
\item $\varphi_{298} :\equiv (r \in \mathbb{R} \to \forall n\,(\varphi_{66}) \to n \in B_{51} \to S(n) \in \mathbb{N} \to n \neq \emptyset \to \varphi_{152}) \to \varphi_{277}$
\item $\varphi_{299} :\equiv (r \in \mathbb{R} \to {{r}^{S(n)}} \cdot r \in \mathbb{R} \to \varphi_{133}) \to \varphi_{179} \to (r \in \mathbb{R} \to \varphi_{185}) \to \varphi_{252}$
\item $\varphi_{300} :\equiv (r \in \mathbb{R} \to \forall n\,(\varphi_{66}) \to n \in B_{51} \to S(n) \in \mathbb{N} \to n \neq \emptyset \to \varphi_{86}) \to \varphi_{290}$
\item $\varphi_{301} :\equiv (1 - r \in \mathbb{R} \to {\textstyle\sum_{k=1}^{n} {{r}^{k}}} \in \mathbb{R} \to {{r}^{S(n)}} \in \mathbb{R} \to \varphi_{118}) \to \varphi_{183} \to (r \in \mathbb{R} \to \varphi_{182}) \to \varphi_{256}$
\item $\varphi_{302} :\equiv (r \in \mathbb{R} \to \exists B_{51}\,\forall n\,(\varphi_{66})) \to \forall B_{51}\,(\varphi_{81} \to \neg \forall n\,(\varphi_{66})) \to \varphi_{248} \to (\varphi_{81} \to \forall B_{51}\,(\varphi_{81})) \to \varphi_{79}$
\item $\varphi_{303} :\equiv (r \in \mathbb{R} \to \forall n\,(\varphi_{66}) \to n \in B_{51} \to S(n) \in \mathbb{N} \to n \neq \emptyset \to \varphi_{73}) \to \varphi_{294}$
\item $\varphi_{304} :\equiv (r \in \mathbb{R} \to \forall n\,(\varphi_{66}) \to n \in B_{51} \to S(n) \in \mathbb{N} \to n \neq \emptyset \to \varphi_{103}) \to \varphi_{298}$
\end{abbrevlist}
\end{fullproof}

% 自動生成 (prover/build.js): GeoSer  ノード数 884
\begin{fullproof}{命題4.3(1)}{$\forall r \in \mathbb{R}\,(\lvert r \rvert < 1 \to \mathrm{Conv}(\Sigma {(r)^{\bullet}},r \cdot (1 - r)^{-1}))$}{884}\phantomsection\label{tree:GeoSer}
\noindent$\Pi_{1}$:
\begin{ptree}
\AXC{$\varphi_{5}$\lbl{[C18]}}
\AXC{$\varphi_{5} \to \varphi_{7}$\lbl{A2}}
\BIC{$\varphi_{7}$}
\AXC{$\varphi_{7} \to \varphi_{11}$\lbl{A2}}
\BIC{$\varphi_{11}$}
\AXC{$\varphi_{11} \to \varphi_{6}$\lbl{A2}}
\BIC{$\varphi_{6}$}
\AXC{$\varphi_{6} \to \varphi_{130}$\lbl{A1}}
\BIC{$\varphi_{130}$}
\end{ptree}
\noindent$\Pi_{2}$:
\begin{ptree}
\AXC{$1 \in \mathbb{R}$\lbl{(R1)}}
\AXC{$\forall x \in \mathbb{R}\,(-x \in \mathbb{R})$\lbl{(R1)}}
\AXC{$\forall x \in \mathbb{R}\,(-x \in \mathbb{R}) \to 1 \in \mathbb{R} \to -1 \in \mathbb{R}$\lbl{A2}}
\BIC{$1 \in \mathbb{R} \to -1 \in \mathbb{R}$}
\AXC{$\varphi_{73}$\lbl{A1}}
\BIC{$1 \in \mathbb{R} \to r \in \mathbb{R} \to \lvert r \rvert < 1 \to -1 \in \mathbb{R}$}
\BIC{$r \in \mathbb{R} \to \lvert r \rvert < 1 \to -1 \in \mathbb{R}$}
\end{ptree}
\noindent$\Pi_{3}$:
\begin{ptree}
\AXC{$\varphi_{8}$\lbl{定理2.1(2)}}
\AXC{$\varphi_{8} \to \varphi_{33}$\lbl{A2}}
\BIC{$\varphi_{33}$}
\AXC{$\varphi_{33} \to \varphi_{105}$\lbl{A2}}
\BIC{$\varphi_{105}$}
\AXC{$\varphi_{105} \to -1 \in \mathbb{R} \to \varphi_{43}$\lbl{A2}}
\BIC{$-1 \in \mathbb{R} \to \varphi_{43}$}
\end{ptree}
\noindent$\Pi_{4}$:
\begin{ptree}
\AXC{$\stackrel{\Pi_{2}}{\vdots}$}\noLine
\UIC{$r \in \mathbb{R} \to \lvert r \rvert < 1 \to -1 \in \mathbb{R}$}
\AXC{$\stackrel{\Pi_{3}}{\vdots}$}\noLine
\UIC{$-1 \in \mathbb{R} \to \varphi_{43}$}
\AXC{$(-1 \in \mathbb{R} \to \varphi_{43}) \to (r \in \mathbb{R} \to \lvert r \rvert < 1 \to -1 \in \mathbb{R}) \to \varphi_{158}$\lbl{A1}}
\BIC{$(r \in \mathbb{R} \to \lvert r \rvert < 1 \to -1 \in \mathbb{R}) \to \varphi_{158}$}
\BIC{$\varphi_{158}$}
\end{ptree}
\noindent$\Pi_{5}$:
\begin{ptree}
\AXC{$\mathrm{SubIdx}(\mathsf{s})$\lbl{[C28]}}
\AXC{$\forall r \in \mathbb{R}\,(\lvert r \rvert < 1 \to \mathrm{Conv}({(r)^{\bullet}},0))$\lbl{命題3.1(3)}}
\AXC{$\varphi_{72}$\lbl{A2}}
\BIC{$r \in \mathbb{R} \to \lvert r \rvert < 1 \to \mathrm{Conv}({(r)^{\bullet}},0)$}
\AXC{$\stackrel{\Pi_{1}}{\vdots}$}\noLine
\UIC{$\varphi_{130}$}
\BIC{$\varphi_{23}$}
\BIC{$r \in \mathbb{R} \to \lvert r \rvert < 1 \to \mathrm{Conv}({(r)^{\bullet}} \circ \mathsf{s},0)$}
\AXC{$\stackrel{\Pi_{4}}{\vdots}$}\noLine
\UIC{$\varphi_{158}$}
\BIC{$\varphi_{39}$}
\end{ptree}
\noindent$\Pi_{6}$:
\begin{ptree}
\AXC{$r \in \mathbb{R} \to \lvert r \rvert < 1 \to r \in \mathbb{R}$\lbl{A1}}
\AXC{$\forall c \in \mathbb{R}\,(\mathrm{Conv}(\overline{c},c))$\lbl{[G4]}}
\AXC{$\varphi_{13}$\lbl{A2}}
\BIC{$r \in \mathbb{R} \to \mathrm{Conv}(\overline{r},r)$}
\AXC{$(r \in \mathbb{R} \to \mathrm{Conv}(\overline{r},r)) \to \varphi_{59}$\lbl{A1}}
\BIC{$\varphi_{59}$}
\BIC{$r \in \mathbb{R} \to \lvert r \rvert < 1 \to \mathrm{Conv}(\overline{r},r)$}
\end{ptree}
\noindent$\Pi_{7}$:
\begin{ptree}
\AXC{$\varphi_{20}$\lbl{定理2.1(1)}}
\AXC{$\varphi_{20} \to \varphi_{40}$\lbl{A2}}
\BIC{$\varphi_{40}$}
\AXC{$\varphi_{40} \to \varphi_{27}$\lbl{A2}}
\BIC{$\varphi_{27}$}
\AXC{$\varphi_{27} \to \varphi_{148}$\lbl{A2}}
\BIC{$\varphi_{148}$}
\AXC{$\varphi_{148} \to \mathrm{Conv}(\overline{r},r) \to \varphi_{144}$\lbl{A2}}
\BIC{$\mathrm{Conv}(\overline{r},r) \to \varphi_{144}$}
\end{ptree}
\noindent$\Pi_{8}$:
\begin{ptree}
\AXC{$\stackrel{\Pi_{5}}{\vdots}$}\noLine
\UIC{$\varphi_{39}$}
\AXC{$\stackrel{\Pi_{6}}{\vdots}$}\noLine
\UIC{$r \in \mathbb{R} \to \lvert r \rvert < 1 \to \mathrm{Conv}(\overline{r},r)$}
\AXC{$\stackrel{\Pi_{7}}{\vdots}$}\noLine
\UIC{$\mathrm{Conv}(\overline{r},r) \to \varphi_{144}$}
\AXC{$(\mathrm{Conv}(\overline{r},r) \to \varphi_{144}) \to \varphi_{224}$\lbl{A1}}
\BIC{$\varphi_{224}$}
\BIC{$\varphi_{39} \to r \in \mathbb{R} \to \varphi_{57}$}
\BIC{$r \in \mathbb{R} \to \varphi_{57}$}
\end{ptree}
\noindent$\Pi_{9}$:
\begin{ptree}
\AXC{$x = x$\lbl{A5}}
\AXC{$x = y \to x = x \to y = x$\lbl{A6}}
\AXC{$(x = y \to x = x \to y = x) \to x = x \to x = y \to y = x$\lbl{A1}}
\BIC{$x = x \to x = y \to y = x$}
\BIC{$x = y \to y = x$}
\UIC{$\forall y\,(x = y \to y = x)$}
\UIC{$\forall x\,\forall y\,(x = y \to y = x)$}
\end{ptree}
\noindent$\Pi_{10}$:
\begin{ptree}
\AXC{$\stackrel{\Pi_{9}}{\vdots}$}\noLine
\UIC{$\forall x\,\forall y\,(x = y \to y = x)$}
\AXC{$\forall x\,\forall y\,(x = y \to y = x) \to \forall y\,(1 - r = y \to y = 1 - r)$\lbl{A2}}
\BIC{$\forall y\,(1 - r = y \to y = 1 - r)$}
\end{ptree}
\noindent$\Pi_{11}$:
\begin{ptree}
\AXC{$\stackrel{\Pi_{10}}{\vdots}$}\noLine
\UIC{$\forall y\,(1 - r = y \to y = 1 - r)$}
\AXC{$\forall y\,(1 - r = y \to y = 1 - r) \to 1 - r = 0 \to 0 = 1 - r$\lbl{A2}}
\BIC{$1 - r = 0 \to 0 = 1 - r$}
\end{ptree}
\noindent$\Pi_{12}$:
\begin{ptree}
\AXC{$\forall x \in \mathbb{R}\,(-\lvert x \rvert \leq x \land x \leq \lvert x \rvert)$\lbl{[A8]}}
\AXC{$\varphi_{54}$\lbl{A2}}
\BIC{$r \in \mathbb{R} \to -\lvert r \rvert \leq r \land r \leq \lvert r \rvert$}
\AXC{$(r \in \mathbb{R} \to -\lvert r \rvert \leq r \land r \leq \lvert r \rvert) \to \varphi_{92}$\lbl{A1}}
\BIC{$\varphi_{92}$}
\end{ptree}
\noindent$\Pi_{13}$:
\begin{ptree}
\AXC{$\forall x\,\forall y\,\forall z\,(x \leq y \to y < z \to x < z)$\lbl{[O4]}}
\AXC{$\varphi_{9}$\lbl{A2}}
\BIC{$\forall y\,\forall z\,(r \leq y \to y < z \to r < z)$}
\AXC{$\varphi_{19}$\lbl{A2}}
\BIC{$\forall z\,(r \leq \lvert r \rvert \to \lvert r \rvert < z \to r < z)$}
\AXC{$\varphi_{28}$\lbl{A2}}
\BIC{$r \leq \lvert r \rvert \to \lvert r \rvert < 1 \to r < 1$}
\AXC{$\varphi_{177}$\lbl{A1}}
\BIC{$(r \in \mathbb{R} \to \lvert r \rvert < 1 \to r \leq \lvert r \rvert) \to \varphi_{41}$}
\end{ptree}
\noindent$\Pi_{14}$:
\begin{ptree}
\AXC{$r \in \mathbb{R} \to \lvert r \rvert < 1 \to r \in \mathbb{R}$\lbl{A1}}
\AXC{$\stackrel{\Pi_{12}}{\vdots}$}\noLine
\UIC{$\varphi_{92}$}
\BIC{$r \in \mathbb{R} \to \lvert r \rvert < 1 \to -\lvert r \rvert \leq r \land r \leq \lvert r \rvert$}
\AXC{$\varphi_{86}$\lbl{A1}}
\BIC{$r \in \mathbb{R} \to \lvert r \rvert < 1 \to r \leq \lvert r \rvert$}
\AXC{$\stackrel{\Pi_{13}}{\vdots}$}\noLine
\UIC{$(r \in \mathbb{R} \to \lvert r \rvert < 1 \to r \leq \lvert r \rvert) \to \varphi_{41}$}
\BIC{$\varphi_{41}$}
\end{ptree}
\noindent$\Pi_{15}$:
\begin{ptree}
\AXC{$r \in \mathbb{R} \to \lvert r \rvert < 1 \to r \in \mathbb{R}$\lbl{A1}}
\AXC{$\forall x\,\forall y\,(x \in \mathbb{R} \to y \in \mathbb{R} \to x < y \to 0 < y - x)$\lbl{[F24]}}
\AXC{$\varphi_{118}$\lbl{A2}}
\BIC{$\forall y\,(r \in \mathbb{R} \to y \in \mathbb{R} \to r < y \to 0 < y - r)$}
\AXC{$\varphi_{106}$\lbl{A2}}
\BIC{$r \in \mathbb{R} \to 1 \in \mathbb{R} \to r < 1 \to 0 < 1 - r$}
\AXC{$\varphi_{198}$\lbl{A1}}
\BIC{$(r \in \mathbb{R} \to \lvert r \rvert < 1 \to r \in \mathbb{R}) \to 1 \in \mathbb{R} \to \varphi_{42}$}
\BIC{$1 \in \mathbb{R} \to \varphi_{42}$}
\end{ptree}
\noindent$\Pi_{16}$:
\begin{ptree}
\AXC{$\forall x\,\forall y\,(x < y \to x \neq y)$\lbl{[F26]}}
\AXC{$\forall x\,\forall y\,(x < y \to x \neq y) \to \forall y\,(0 < y \to 0 \neq y)$\lbl{A2}}
\BIC{$\forall y\,(0 < y \to 0 \neq y)$}
\end{ptree}
\noindent$\Pi_{17}$:
\begin{ptree}
\AXC{$\stackrel{\Pi_{16}}{\vdots}$}\noLine
\UIC{$\forall y\,(0 < y \to 0 \neq y)$}
\AXC{$\forall y\,(0 < y \to 0 \neq y) \to 0 < 1 - r \to 0 \neq 1 - r$\lbl{A2}}
\BIC{$0 < 1 - r \to 0 \neq 1 - r$}
\AXC{$(0 < 1 - r \to 0 \neq 1 - r) \to \varphi_{47}$\lbl{A1}}
\BIC{$\varphi_{47}$}
\end{ptree}
\noindent$\Pi_{18}$:
\begin{ptree}
\AXC{$r \in \mathbb{R} \to \lvert r \rvert < 1 \to \lvert r \rvert < 1$\lbl{A1}}
\AXC{$\stackrel{\Pi_{14}}{\vdots}$}\noLine
\UIC{$\varphi_{41}$}
\BIC{$r \in \mathbb{R} \to \lvert r \rvert < 1 \to r < 1$}
\AXC{$1 \in \mathbb{R}$\lbl{(R1)}}
\AXC{$\stackrel{\Pi_{15}}{\vdots}$}\noLine
\UIC{$1 \in \mathbb{R} \to \varphi_{42}$}
\BIC{$\varphi_{42}$}
\BIC{$r \in \mathbb{R} \to \lvert r \rvert < 1 \to 0 < 1 - r$}
\AXC{$\stackrel{\Pi_{17}}{\vdots}$}\noLine
\UIC{$\varphi_{47}$}
\BIC{$r \in \mathbb{R} \to \lvert r \rvert < 1 \to 0 \neq 1 - r$}
\AXC{$\varphi_{104}$\lbl{A1}}
\BIC{$(1 - r = 0 \to 0 = 1 - r) \to r \in \mathbb{R} \to \lvert r \rvert < 1 \to 1 - r \neq 0$}
\end{ptree}
\noindent$\Pi_{19}$:
\begin{ptree}
\AXC{$\stackrel{\Pi_{11}}{\vdots}$}\noLine
\UIC{$1 - r = 0 \to 0 = 1 - r$}
\AXC{$\stackrel{\Pi_{18}}{\vdots}$}\noLine
\UIC{$(1 - r = 0 \to 0 = 1 - r) \to r \in \mathbb{R} \to \lvert r \rvert < 1 \to 1 - r \neq 0$}
\BIC{$r \in \mathbb{R} \to \lvert r \rvert < 1 \to 1 - r \neq 0$}
\end{ptree}
\noindent$\Pi_{20}$:
\begin{ptree}
\AXC{$\forall x\,\forall y\,(x \in \mathbb{R} \to y \in \mathbb{R} \to x - y \in \mathbb{R})$\lbl{[F1]}}
\AXC{$\varphi_{99}$\lbl{A2}}
\BIC{$\forall y\,(1 \in \mathbb{R} \to y \in \mathbb{R} \to 1 - y \in \mathbb{R})$}
\AXC{$\varphi_{82}$\lbl{A2}}
\BIC{$1 \in \mathbb{R} \to r \in \mathbb{R} \to 1 - r \in \mathbb{R}$}
\AXC{$(1 \in \mathbb{R} \to r \in \mathbb{R} \to 1 - r \in \mathbb{R}) \to 1 \in \mathbb{R} \to \varphi_{60}$\lbl{A1}}
\BIC{$1 \in \mathbb{R} \to \varphi_{60}$}
\end{ptree}
\noindent$\Pi_{21}$:
\begin{ptree}
\AXC{$r \in \mathbb{R} \to \lvert r \rvert < 1 \to r \in \mathbb{R}$\lbl{A1}}
\AXC{$1 \in \mathbb{R}$\lbl{(R1)}}
\AXC{$\stackrel{\Pi_{20}}{\vdots}$}\noLine
\UIC{$1 \in \mathbb{R} \to \varphi_{60}$}
\BIC{$\varphi_{60}$}
\BIC{$r \in \mathbb{R} \to \lvert r \rvert < 1 \to 1 - r \in \mathbb{R}$}
\end{ptree}
\noindent$\Pi_{22}$:
\begin{ptree}
\AXC{$\stackrel{\Pi_{21}}{\vdots}$}\noLine
\UIC{$r \in \mathbb{R} \to \lvert r \rvert < 1 \to 1 - r \in \mathbb{R}$}
\AXC{$\forall x \in \mathbb{R}\,(x \neq 0 \to x^{-1} \in \mathbb{R})$\lbl{(R1)}}
\AXC{$\varphi_{93}$\lbl{A2}}
\BIC{$1 - r \in \mathbb{R} \to 1 - r \neq 0 \to (1 - r)^{-1} \in \mathbb{R}$}
\AXC{$\varphi_{203}$\lbl{A1}}
\BIC{$(r \in \mathbb{R} \to \lvert r \rvert < 1 \to 1 - r \in \mathbb{R}) \to \varphi_{74}$}
\BIC{$\varphi_{74}$}
\end{ptree}
\noindent$\Pi_{23}$:
\begin{ptree}
\AXC{$\stackrel{\Pi_{19}}{\vdots}$}\noLine
\UIC{$r \in \mathbb{R} \to \lvert r \rvert < 1 \to 1 - r \neq 0$}
\AXC{$\stackrel{\Pi_{22}}{\vdots}$}\noLine
\UIC{$\varphi_{74}$}
\BIC{$r \in \mathbb{R} \to \lvert r \rvert < 1 \to (1 - r)^{-1} \in \mathbb{R}$}
\end{ptree}
\noindent$\Pi_{24}$:
\begin{ptree}
\AXC{$\varphi_{8}$\lbl{定理2.1(2)}}
\AXC{$\varphi_{8} \to \varphi_{75}$\lbl{A2}}
\BIC{$\varphi_{75}$}
\AXC{$\varphi_{75} \to \varphi_{190}$\lbl{A2}}
\BIC{$\varphi_{190}$}
\AXC{$\varphi_{190} \to (1 - r)^{-1} \in \mathbb{R} \to \varphi_{193}$\lbl{A2}}
\BIC{$(1 - r)^{-1} \in \mathbb{R} \to \varphi_{193}$}
\AXC{$((1 - r)^{-1} \in \mathbb{R} \to \varphi_{193}) \to \varphi_{247}$\lbl{A1}}
\BIC{$\varphi_{247}$}
\end{ptree}
\noindent$\Pi_{25}$:
\begin{ptree}
\AXC{$\stackrel{\Pi_{8}}{\vdots}$}\noLine
\UIC{$r \in \mathbb{R} \to \varphi_{57}$}
\AXC{$\stackrel{\Pi_{23}}{\vdots}$}\noLine
\UIC{$r \in \mathbb{R} \to \lvert r \rvert < 1 \to (1 - r)^{-1} \in \mathbb{R}$}
\AXC{$\stackrel{\Pi_{24}}{\vdots}$}\noLine
\UIC{$\varphi_{247}$}
\BIC{$(r \in \mathbb{R} \to \varphi_{57}) \to r \in \mathbb{R} \to \lvert r \rvert < 1 \to \varphi_{114}$}
\BIC{$r \in \mathbb{R} \to \lvert r \rvert < 1 \to \varphi_{114}$}
\end{ptree}
\noindent$\Pi_{26}$:
\begin{ptree}
\AXC{$y = z \to x = y \to x = z$\lbl{A6}}
\AXC{$(y = z \to x = y \to x = z) \to x = y \to y = z \to x = z$\lbl{A1}}
\BIC{$x = y \to y = z \to x = z$}
\UIC{$\forall z\,(x = y \to y = z \to x = z)$}
\UIC{$\forall y\,\forall z\,(x = y \to y = z \to x = z)$}
\UIC{$\forall x\,\forall y\,\forall z\,(x = y \to y = z \to x = z)$}
\AXC{$\forall x\,\forall y\,\forall z\,(x = y \to y = z \to x = z) \to \varphi_{185}$\lbl{A2}}
\BIC{$\varphi_{185}$}
\end{ptree}
\noindent$\Pi_{27}$:
\begin{ptree}
\AXC{$\forall r \in \mathbb{R}\,(\varphi_{58})$\lbl{[S24]}}
\AXC{$\forall r \in \mathbb{R}\,(\varphi_{58}) \to r \in \mathbb{R} \to \varphi_{58}$\lbl{A2}}
\BIC{$r \in \mathbb{R} \to \varphi_{58}$}
\AXC{$(r \in \mathbb{R} \to \varphi_{58}) \to \varphi_{194}$\lbl{A1}}
\BIC{$\varphi_{194}$}
\end{ptree}
\noindent$\Pi_{28}$:
\begin{ptree}
\AXC{$\varphi_{149}$\lbl{[S25]}}
\AXC{$\varphi_{149} \to \mathrm{Seq}({(r)^{\bullet}}) \to \varphi_{164}$\lbl{A2}}
\BIC{$\mathrm{Seq}({(r)^{\bullet}}) \to \varphi_{164}$}
\AXC{$(\mathrm{Seq}({(r)^{\bullet}}) \to \varphi_{164}) \to \varphi_{228}$\lbl{A1}}
\BIC{$\varphi_{228}$}
\end{ptree}
\noindent$\Pi_{29}$:
\begin{ptree}
\AXC{$r \in \mathbb{R} \to \lvert r \rvert < 1 \to n \in \mathbb{N} \to r \in \mathbb{R}$\lbl{A1}}
\AXC{$\stackrel{\Pi_{27}}{\vdots}$}\noLine
\UIC{$\varphi_{194}$}
\BIC{$r \in \mathbb{R} \to \lvert r \rvert < 1 \to n \in \mathbb{N} \to \varphi_{58}$}
\AXC{$\varphi_{195}$\lbl{A1}}
\BIC{$r \in \mathbb{R} \to \lvert r \rvert < 1 \to n \in \mathbb{N} \to \mathrm{Seq}({(r)^{\bullet}})$}
\AXC{$\stackrel{\Pi_{28}}{\vdots}$}\noLine
\UIC{$\varphi_{228}$}
\BIC{$r \in \mathbb{R} \to \lvert r \rvert < 1 \to n \in \mathbb{N} \to \varphi_{164}$}
\AXC{$\varphi_{232}$\lbl{A1}}
\BIC{$r \in \mathbb{R} \to \lvert r \rvert < 1 \to n \in \mathbb{N} \to \mathrm{Seq}(\Sigma {(r)^{\bullet}})$}
\end{ptree}
\noindent$\Pi_{30}$:
\begin{ptree}
\AXC{$\forall a\,\forall n\,(\mathrm{Seq}(a) \to n \in \mathbb{N} \to a_{n} \in \mathbb{R})$\lbl{[S1]}}
\AXC{$\forall a\,\forall n\,(\mathrm{Seq}(a) \to n \in \mathbb{N} \to a_{n} \in \mathbb{R}) \to \varphi_{12}$\lbl{A2}}
\BIC{$\varphi_{12}$}
\end{ptree}
\noindent$\Pi_{31}$:
\begin{ptree}
\AXC{$\stackrel{\Pi_{30}}{\vdots}$}\noLine
\UIC{$\varphi_{12}$}
\AXC{$\varphi_{12} \to \mathrm{Seq}(\Sigma {(r)^{\bullet}}) \to n \in \mathbb{N} \to {\textstyle\sum_{k=1}^{n} {{r}^{k}}} \in \mathbb{R}$\lbl{A2}}
\BIC{$\mathrm{Seq}(\Sigma {(r)^{\bullet}}) \to n \in \mathbb{N} \to {\textstyle\sum_{k=1}^{n} {{r}^{k}}} \in \mathbb{R}$}
\AXC{$\varphi_{235}$\lbl{A1}}
\BIC{$\varphi_{197}$}
\end{ptree}
\noindent$\Pi_{32}$:
\begin{ptree}
\AXC{$r \in \mathbb{R} \to \lvert r \rvert < 1 \to n \in \mathbb{N} \to n \in \mathbb{N}$\lbl{A1}}
\AXC{$\stackrel{\Pi_{29}}{\vdots}$}\noLine
\UIC{$r \in \mathbb{R} \to \lvert r \rvert < 1 \to n \in \mathbb{N} \to \mathrm{Seq}(\Sigma {(r)^{\bullet}})$}
\AXC{$\stackrel{\Pi_{31}}{\vdots}$}\noLine
\UIC{$\varphi_{197}$}
\BIC{$\varphi_{142}$}
\BIC{$r \in \mathbb{R} \to \lvert r \rvert < 1 \to n \in \mathbb{N} \to {\textstyle\sum_{k=1}^{n} {{r}^{k}}} \in \mathbb{R}$}
\end{ptree}
\noindent$\Pi_{33}$:
\begin{ptree}
\AXC{$\forall x \in \mathbb{R}\,(x \cdot 1 = x)$\lbl{(R8)}}
\AXC{$\varphi_{111}$\lbl{A2}}
\BIC{${\textstyle\sum_{k=1}^{n} {{r}^{k}}} \in \mathbb{R} \to {\textstyle\sum_{k=1}^{n} {{r}^{k}}} \cdot 1 = {\textstyle\sum_{k=1}^{n} {{r}^{k}}}$}
\AXC{$({\textstyle\sum_{k=1}^{n} {{r}^{k}}} \in \mathbb{R} \to {\textstyle\sum_{k=1}^{n} {{r}^{k}}} \cdot 1 = {\textstyle\sum_{k=1}^{n} {{r}^{k}}}) \to \varphi_{181}$\lbl{A1}}
\BIC{$\varphi_{181}$}
\end{ptree}
\noindent$\Pi_{34}$:
\begin{ptree}
\AXC{$y = z \to x = y \to x = z$\lbl{A6}}
\AXC{$(y = z \to x = y \to x = z) \to x = y \to y = z \to x = z$\lbl{A1}}
\BIC{$x = y \to y = z \to x = z$}
\UIC{$\forall z\,(x = y \to y = z \to x = z)$}
\UIC{$\forall y\,\forall z\,(x = y \to y = z \to x = z)$}
\UIC{$\forall x\,\forall y\,\forall z\,(x = y \to y = z \to x = z)$}
\AXC{$\forall x\,\forall y\,\forall z\,(x = y \to y = z \to x = z) \to \varphi_{185}$\lbl{A2}}
\BIC{$\varphi_{185}$}
\end{ptree}
\noindent$\Pi_{35}$:
\begin{ptree}
\AXC{$\stackrel{\Pi_{34}}{\vdots}$}\noLine
\UIC{$\varphi_{185}$}
\AXC{$\varphi_{185} \to \varphi_{213}$\lbl{A2}}
\BIC{$\varphi_{213}$}
\AXC{$\varphi_{213} \to \varphi_{97} \to 1 \cdot {\textstyle\sum_{k=1}^{n} {{r}^{k}}} = {\textstyle\sum_{k=1}^{n} {{r}^{k}}} \cdot 1 \to \varphi_{96}$\lbl{A2}}
\BIC{$\varphi_{97} \to 1 \cdot {\textstyle\sum_{k=1}^{n} {{r}^{k}}} = {\textstyle\sum_{k=1}^{n} {{r}^{k}}} \cdot 1 \to \varphi_{96}$}
\end{ptree}
\noindent$\Pi_{36}$:
\begin{ptree}
\AXC{$1 \in \mathbb{R}$\lbl{(R1)}}
\AXC{$\forall x\,\forall y\,(x \in \mathbb{R} \to y \in \mathbb{R} \to x \cdot y = y \cdot x)$\lbl{(R7)}}
\AXC{$\varphi_{89}$\lbl{A2}}
\BIC{$\forall y\,(1 \in \mathbb{R} \to y \in \mathbb{R} \to 1 \cdot y = y \cdot 1)$}
\AXC{$\varphi_{166}$\lbl{A2}}
\BIC{$1 \in \mathbb{R} \to {\textstyle\sum_{k=1}^{n} {{r}^{k}}} \in \mathbb{R} \to 1 \cdot {\textstyle\sum_{k=1}^{n} {{r}^{k}}} = {\textstyle\sum_{k=1}^{n} {{r}^{k}}} \cdot 1$}
\AXC{$\varphi_{244}$\lbl{A1}}
\BIC{$1 \in \mathbb{R} \to \varphi_{184}$}
\BIC{$\varphi_{184}$}
\end{ptree}
\noindent$\Pi_{37}$:
\begin{ptree}
\AXC{$y = z \to x = y \to x = z$\lbl{A6}}
\AXC{$(y = z \to x = y \to x = z) \to x = y \to y = z \to x = z$\lbl{A1}}
\BIC{$x = y \to y = z \to x = z$}
\UIC{$\forall z\,(x = y \to y = z \to x = z)$}
\UIC{$\forall y\,\forall z\,(x = y \to y = z \to x = z)$}
\UIC{$\forall x\,\forall y\,\forall z\,(x = y \to y = z \to x = z)$}
\AXC{$\forall x\,\forall y\,\forall z\,(x = y \to y = z \to x = z) \to \varphi_{185}$\lbl{A2}}
\BIC{$\varphi_{185}$}
\end{ptree}
\noindent$\Pi_{38}$:
\begin{ptree}
\AXC{$z = w \to \varphi_{170}$\lbl{A6}}
\UIC{$\forall w\,(z = w \to \varphi_{170})$}
\UIC{$\forall z\,\forall w\,(z = w \to \varphi_{170})$}
\AXC{$\forall z\,\forall w\,(z = w \to \varphi_{170}) \to \varphi_{206}$\lbl{A2}}
\BIC{$\varphi_{206}$}
\AXC{$\varphi_{206} \to \varphi_{202}$\lbl{A2}}
\BIC{$\varphi_{202}$}
\end{ptree}
\noindent$\Pi_{39}$:
\begin{ptree}
\AXC{$\forall x\,\forall y\,(x \in \mathbb{R} \to y \in \mathbb{R} \to x - y \in \mathbb{R})$\lbl{[F1]}}
\AXC{$\varphi_{99}$\lbl{A2}}
\BIC{$\forall y\,(1 \in \mathbb{R} \to y \in \mathbb{R} \to 1 - y \in \mathbb{R})$}
\AXC{$\varphi_{82}$\lbl{A2}}
\BIC{$1 \in \mathbb{R} \to r \in \mathbb{R} \to 1 - r \in \mathbb{R}$}
\AXC{$(1 \in \mathbb{R} \to r \in \mathbb{R} \to 1 - r \in \mathbb{R}) \to 1 \in \mathbb{R} \to \varphi_{122}$\lbl{A1}}
\BIC{$1 \in \mathbb{R} \to \varphi_{122}$}
\end{ptree}
\noindent$\Pi_{40}$:
\begin{ptree}
\AXC{$r \in \mathbb{R} \to \lvert r \rvert < 1 \to n \in \mathbb{N} \to r \in \mathbb{R}$\lbl{A1}}
\AXC{$1 \in \mathbb{R}$\lbl{(R1)}}
\AXC{$\stackrel{\Pi_{39}}{\vdots}$}\noLine
\UIC{$1 \in \mathbb{R} \to \varphi_{122}$}
\BIC{$\varphi_{122}$}
\BIC{$r \in \mathbb{R} \to \lvert r \rvert < 1 \to n \in \mathbb{N} \to 1 - r \in \mathbb{R}$}
\end{ptree}
\noindent$\Pi_{41}$:
\begin{ptree}
\AXC{$\stackrel{\Pi_{40}}{\vdots}$}\noLine
\UIC{$r \in \mathbb{R} \to \lvert r \rvert < 1 \to n \in \mathbb{N} \to 1 - r \in \mathbb{R}$}
\AXC{$\forall x \in \mathbb{R}\,(x \neq 0 \to x^{-1} \cdot x = 1)$\lbl{[E22]}}
\AXC{$\varphi_{95}$\lbl{A2}}
\BIC{$1 - r \in \mathbb{R} \to 1 - r \neq 0 \to (1 - r)^{-1} \cdot (1 - r) = 1$}
\AXC{$\varphi_{223}$\lbl{A1}}
\BIC{$\varphi_{186}$}
\BIC{$(r \in \mathbb{R} \to \lvert r \rvert < 1 \to 1 - r \neq 0) \to \varphi_{14}$}
\end{ptree}
\noindent$\Pi_{42}$:
\begin{ptree}
\AXC{$\stackrel{\Pi_{19}}{\vdots}$}\noLine
\UIC{$r \in \mathbb{R} \to \lvert r \rvert < 1 \to 1 - r \neq 0$}
\AXC{$\stackrel{\Pi_{41}}{\vdots}$}\noLine
\UIC{$(r \in \mathbb{R} \to \lvert r \rvert < 1 \to 1 - r \neq 0) \to \varphi_{14}$}
\BIC{$\varphi_{14}$}
\AXC{$\varphi_{14} \to \varphi_{77} \to \varphi_{202} \to r \in \mathbb{R} \to \varphi_{98}$\lbl{A1}}
\BIC{$\varphi_{77} \to \varphi_{202} \to r \in \mathbb{R} \to \varphi_{98}$}
\end{ptree}
\noindent$\Pi_{43}$:
\begin{ptree}
\AXC{$y = z \to x = y \to x = z$\lbl{A6}}
\AXC{$(y = z \to x = y \to x = z) \to x = y \to y = z \to x = z$\lbl{A1}}
\BIC{$x = y \to y = z \to x = z$}
\UIC{$\forall z\,(x = y \to y = z \to x = z)$}
\UIC{$\forall y\,\forall z\,(x = y \to y = z \to x = z)$}
\UIC{$\forall x\,\forall y\,\forall z\,(x = y \to y = z \to x = z)$}
\AXC{$\forall x\,\forall y\,\forall z\,(x = y \to y = z \to x = z) \to \varphi_{185}$\lbl{A2}}
\BIC{$\varphi_{185}$}
\end{ptree}
\noindent$\Pi_{44}$:
\begin{ptree}
\AXC{$\stackrel{\Pi_{9}}{\vdots}$}\noLine
\UIC{$\forall x\,\forall y\,(x = y \to y = x)$}
\AXC{$\forall x\,\forall y\,(x = y \to y = x) \to \varphi_{101}$\lbl{A2}}
\BIC{$\varphi_{101}$}
\AXC{$\varphi_{101} \to \varphi_{79} \to \varphi_{78}$\lbl{A2}}
\BIC{$\varphi_{79} \to \varphi_{78}$}
\end{ptree}
\noindent$\Pi_{45}$:
\begin{ptree}
\AXC{$\stackrel{\Pi_{40}}{\vdots}$}\noLine
\UIC{$r \in \mathbb{R} \to \lvert r \rvert < 1 \to n \in \mathbb{N} \to 1 - r \in \mathbb{R}$}
\AXC{$\forall x \in \mathbb{R}\,(x \neq 0 \to x^{-1} \in \mathbb{R})$\lbl{(R1)}}
\AXC{$\varphi_{93}$\lbl{A2}}
\BIC{$1 - r \in \mathbb{R} \to 1 - r \neq 0 \to (1 - r)^{-1} \in \mathbb{R}$}
\AXC{$\varphi_{218}$\lbl{A1}}
\BIC{$(r \in \mathbb{R} \to \lvert r \rvert < 1 \to n \in \mathbb{N} \to 1 - r \in \mathbb{R}) \to \varphi_{107}$}
\BIC{$\varphi_{107}$}
\end{ptree}
\noindent$\Pi_{46}$:
\begin{ptree}
\AXC{$\stackrel{\Pi_{19}}{\vdots}$}\noLine
\UIC{$r \in \mathbb{R} \to \lvert r \rvert < 1 \to 1 - r \neq 0$}
\AXC{$\stackrel{\Pi_{45}}{\vdots}$}\noLine
\UIC{$\varphi_{107}$}
\BIC{$r \in \mathbb{R} \to \lvert r \rvert < 1 \to n \in \mathbb{N} \to (1 - r)^{-1} \in \mathbb{R}$}
\end{ptree}
\noindent$\Pi_{47}$:
\begin{ptree}
\AXC{$\varphi_{34}$\lbl{(R6)}}
\AXC{$\varphi_{34} \to \varphi_{108}$\lbl{A2}}
\BIC{$\varphi_{108}$}
\AXC{$\varphi_{108} \to \varphi_{133}$\lbl{A2}}
\BIC{$\varphi_{133}$}
\AXC{$\varphi_{133} \to (1 - r)^{-1} \in \mathbb{R} \to 1 - r \in \mathbb{R} \to {\textstyle\sum_{k=1}^{n} {{r}^{k}}} \in \mathbb{R} \to \varphi_{79}$\lbl{A2}}
\BIC{$(1 - r)^{-1} \in \mathbb{R} \to 1 - r \in \mathbb{R} \to {\textstyle\sum_{k=1}^{n} {{r}^{k}}} \in \mathbb{R} \to \varphi_{79}$}
\AXC{$\varphi_{282}$\lbl{A1}}
\BIC{$\varphi_{257}$}
\end{ptree}
\noindent$\Pi_{48}$:
\begin{ptree}
\AXC{$\stackrel{\Pi_{40}}{\vdots}$}\noLine
\UIC{$r \in \mathbb{R} \to \lvert r \rvert < 1 \to n \in \mathbb{N} \to 1 - r \in \mathbb{R}$}
\AXC{$\stackrel{\Pi_{46}}{\vdots}$}\noLine
\UIC{$r \in \mathbb{R} \to \lvert r \rvert < 1 \to n \in \mathbb{N} \to (1 - r)^{-1} \in \mathbb{R}$}
\AXC{$\stackrel{\Pi_{47}}{\vdots}$}\noLine
\UIC{$\varphi_{257}$}
\BIC{$(r \in \mathbb{R} \to \lvert r \rvert < 1 \to n \in \mathbb{N} \to 1 - r \in \mathbb{R}) \to \varphi_{209}$}
\BIC{$\varphi_{209}$}
\end{ptree}
\noindent$\Pi_{49}$:
\begin{ptree}
\AXC{$\stackrel{\Pi_{44}}{\vdots}$}\noLine
\UIC{$\varphi_{79} \to \varphi_{78}$}
\AXC{$\stackrel{\Pi_{32}}{\vdots}$}\noLine
\UIC{$r \in \mathbb{R} \to \lvert r \rvert < 1 \to n \in \mathbb{N} \to {\textstyle\sum_{k=1}^{n} {{r}^{k}}} \in \mathbb{R}$}
\AXC{$\stackrel{\Pi_{48}}{\vdots}$}\noLine
\UIC{$\varphi_{209}$}
\BIC{$r \in \mathbb{R} \to \lvert r \rvert < 1 \to n \in \mathbb{N} \to \varphi_{79}$}
\AXC{$\varphi_{276}$\lbl{A1}}
\BIC{$(\varphi_{79} \to \varphi_{78}) \to r \in \mathbb{R} \to \lvert r \rvert < 1 \to n \in \mathbb{N} \to \varphi_{78}$}
\BIC{$r \in \mathbb{R} \to \lvert r \rvert < 1 \to n \in \mathbb{N} \to \varphi_{78}$}
\end{ptree}
\noindent$\Pi_{50}$:
\begin{ptree}
\AXC{$y = z \to x = y \to x = z$\lbl{A6}}
\AXC{$(y = z \to x = y \to x = z) \to x = y \to y = z \to x = z$\lbl{A1}}
\BIC{$x = y \to y = z \to x = z$}
\UIC{$\forall z\,(x = y \to y = z \to x = z)$}
\UIC{$\forall y\,\forall z\,(x = y \to y = z \to x = z)$}
\UIC{$\forall x\,\forall y\,\forall z\,(x = y \to y = z \to x = z)$}
\AXC{$\forall x\,\forall y\,\forall z\,(x = y \to y = z \to x = z) \to \varphi_{185}$\lbl{A2}}
\BIC{$\varphi_{185}$}
\end{ptree}
\noindent$\Pi_{51}$:
\begin{ptree}
\AXC{$z = w \to \varphi_{128}$\lbl{A6}}
\UIC{$\forall w\,(z = w \to \varphi_{128})$}
\UIC{$\forall z\,\forall w\,(z = w \to \varphi_{128})$}
\AXC{$\forall z\,\forall w\,(z = w \to \varphi_{128}) \to \varphi_{183}$\lbl{A2}}
\BIC{$\varphi_{183}$}
\AXC{$\varphi_{183} \to r - {{r}^{S(n)}} = (1 - r) \cdot {\textstyle\sum_{k=1}^{n} {{r}^{k}}} \to \varphi_{35} \to \varphi_{52}$\lbl{A2}}
\BIC{$r - {{r}^{S(n)}} = (1 - r) \cdot {\textstyle\sum_{k=1}^{n} {{r}^{k}}} \to \varphi_{35} \to \varphi_{52}$}
\end{ptree}
\noindent$\Pi_{52}$:
\begin{ptree}
\AXC{$r \in \mathbb{R} \to \lvert r \rvert < 1 \to n \in \mathbb{N} \to r \in \mathbb{R}$\lbl{A1}}
\AXC{$\forall r \in \mathbb{R}\,\forall n \in \mathbb{N}\,((1 - r) \cdot {\textstyle\sum_{k=1}^{n} {{r}^{k}}} = r - {{r}^{S(n)}})$\lbl{[Q3]}}
\AXC{$\varphi_{188}$\lbl{A2}}
\BIC{$r \in \mathbb{R} \to \forall n \in \mathbb{N}\,((1 - r) \cdot {\textstyle\sum_{k=1}^{n} {{r}^{k}}} = r - {{r}^{S(n)}})$}
\AXC{$\varphi_{234}$\lbl{A1}}
\BIC{$(r \in \mathbb{R} \to \lvert r \rvert < 1 \to n \in \mathbb{N} \to r \in \mathbb{R}) \to r \in \mathbb{R} \to \varphi_{91}$}
\BIC{$r \in \mathbb{R} \to \varphi_{91}$}
\end{ptree}
\noindent$\Pi_{53}$:
\begin{ptree}
\AXC{$r \in \mathbb{R} \to \lvert r \rvert < 1 \to n \in \mathbb{N} \to n \in \mathbb{N}$\lbl{A1}}
\AXC{$\varphi_{159}$\lbl{A2}}
\AXC{$\stackrel{\Pi_{52}}{\vdots}$}\noLine
\UIC{$r \in \mathbb{R} \to \varphi_{91}$}
\AXC{$(r \in \mathbb{R} \to \varphi_{91}) \to \varphi_{159} \to r \in \mathbb{R} \to \varphi_{80}$\lbl{A1}}
\BIC{$\varphi_{159} \to r \in \mathbb{R} \to \varphi_{80}$}
\BIC{$r \in \mathbb{R} \to \varphi_{80}$}
\AXC{$\varphi_{239}$\lbl{A1}}
\BIC{$(r \in \mathbb{R} \to \lvert r \rvert < 1 \to n \in \mathbb{N} \to n \in \mathbb{N}) \to \varphi_{83}$}
\BIC{$\varphi_{83}$}
\end{ptree}
\noindent$\Pi_{54}$:
\begin{ptree}
\AXC{$\stackrel{\Pi_{9}}{\vdots}$}\noLine
\UIC{$\forall x\,\forall y\,(x = y \to y = x)$}
\AXC{$\varphi_{63}$\lbl{A2}}
\BIC{$\forall y\,((1 - r) \cdot {\textstyle\sum_{k=1}^{n} {{r}^{k}}} = y \to y = (1 - r) \cdot {\textstyle\sum_{k=1}^{n} {{r}^{k}}})$}
\AXC{$\varphi_{192}$\lbl{A2}}
\BIC{$\varphi_{115}$}
\AXC{$\stackrel{\Pi_{53}}{\vdots}$}\noLine
\UIC{$\varphi_{83}$}
\AXC{$\varphi_{83} \to \varphi_{115} \to \varphi_{84}$\lbl{A1}}
\BIC{$\varphi_{115} \to \varphi_{84}$}
\BIC{$\varphi_{84}$}
\AXC{$\varphi_{84} \to \varphi_{35} \to \varphi_{258}$\lbl{A1}}
\BIC{$\varphi_{35} \to \varphi_{258}$}
\end{ptree}
\noindent$\Pi_{55}$:
\begin{ptree}
\AXC{$\stackrel{\Pi_{51}}{\vdots}$}\noLine
\UIC{$r - {{r}^{S(n)}} = (1 - r) \cdot {\textstyle\sum_{k=1}^{n} {{r}^{k}}} \to \varphi_{35} \to \varphi_{52}$}
\AXC{$x = x$\lbl{A5}}
\UIC{$\forall x\,(x = x)$}
\AXC{$\forall x\,(x = x) \to \varphi_{35}$\lbl{A2}}
\BIC{$\varphi_{35}$}
\AXC{$\stackrel{\Pi_{54}}{\vdots}$}\noLine
\UIC{$\varphi_{35} \to \varphi_{258}$}
\BIC{$\varphi_{258}$}
\BIC{$r \in \mathbb{R} \to \lvert r \rvert < 1 \to n \in \mathbb{N} \to \varphi_{52}$}
\end{ptree}
\noindent$\Pi_{56}$:
\begin{ptree}
\AXC{$y = z \to x = y \to x = z$\lbl{A6}}
\AXC{$(y = z \to x = y \to x = z) \to x = y \to y = z \to x = z$\lbl{A1}}
\BIC{$x = y \to y = z \to x = z$}
\UIC{$\forall z\,(x = y \to y = z \to x = z)$}
\UIC{$\forall y\,\forall z\,(x = y \to y = z \to x = z)$}
\UIC{$\forall x\,\forall y\,\forall z\,(x = y \to y = z \to x = z)$}
\AXC{$\forall x\,\forall y\,\forall z\,(x = y \to y = z \to x = z) \to \varphi_{185}$\lbl{A2}}
\BIC{$\varphi_{185}$}
\end{ptree}
\noindent$\Pi_{57}$:
\begin{ptree}
\AXC{$z = w \to \varphi_{145}$\lbl{A6}}
\UIC{$\forall w\,(z = w \to \varphi_{145})$}
\UIC{$\forall z\,\forall w\,(z = w \to \varphi_{145})$}
\AXC{$\forall z\,\forall w\,(z = w \to \varphi_{145}) \to \varphi_{196}$\lbl{A2}}
\BIC{$\varphi_{196}$}
\AXC{$\varphi_{196} \to r + -{{r}^{S(n)}} = r - {{r}^{S(n)}} \to \varphi_{55} \to \varphi_{46}$\lbl{A2}}
\BIC{$r + -{{r}^{S(n)}} = r - {{r}^{S(n)}} \to \varphi_{55} \to \varphi_{46}$}
\end{ptree}
\noindent$\Pi_{58}$:
\begin{ptree}
\AXC{$\forall x\,\forall y\,(x - y = x + -y)$\lbl{Def$-$}}
\AXC{$\forall x\,\forall y\,(x - y = x + -y) \to \forall y\,(r - y = r + -y)$\lbl{A2}}
\BIC{$\forall y\,(r - y = r + -y)$}
\end{ptree}
\noindent$\Pi_{59}$:
\begin{ptree}
\AXC{$\stackrel{\Pi_{58}}{\vdots}$}\noLine
\UIC{$\forall y\,(r - y = r + -y)$}
\AXC{$\forall y\,(r - y = r + -y) \to r - {{r}^{S(n)}} = r + -{{r}^{S(n)}}$\lbl{A2}}
\BIC{$r - {{r}^{S(n)}} = r + -{{r}^{S(n)}}$}
\AXC{$r - {{r}^{S(n)}} = r + -{{r}^{S(n)}} \to \varphi_{175}$\lbl{A1}}
\BIC{$\varphi_{175}$}
\end{ptree}
\noindent$\Pi_{60}$:
\begin{ptree}
\AXC{$\stackrel{\Pi_{9}}{\vdots}$}\noLine
\UIC{$\forall x\,\forall y\,(x = y \to y = x)$}
\AXC{$\varphi_{29}$\lbl{A2}}
\BIC{$\forall y\,(r - {{r}^{S(n)}} = y \to y = r - {{r}^{S(n)}})$}
\AXC{$\varphi_{179}$\lbl{A2}}
\BIC{$r - {{r}^{S(n)}} = r + -{{r}^{S(n)}} \to r + -{{r}^{S(n)}} = r - {{r}^{S(n)}}$}
\AXC{$\stackrel{\Pi_{59}}{\vdots}$}\noLine
\UIC{$\varphi_{175}$}
\BIC{$r + -{{r}^{S(n)}} = r - {{r}^{S(n)}}$}
\AXC{$r + -{{r}^{S(n)}} = r - {{r}^{S(n)}} \to \varphi_{55} \to \varphi_{259}$\lbl{A1}}
\BIC{$\varphi_{55} \to \varphi_{259}$}
\end{ptree}
\noindent$\Pi_{61}$:
\begin{ptree}
\AXC{$\stackrel{\Pi_{57}}{\vdots}$}\noLine
\UIC{$r + -{{r}^{S(n)}} = r - {{r}^{S(n)}} \to \varphi_{55} \to \varphi_{46}$}
\AXC{$x = x$\lbl{A5}}
\UIC{$\forall x\,(x = x)$}
\AXC{$\forall x\,(x = x) \to \varphi_{55}$\lbl{A2}}
\BIC{$\varphi_{55}$}
\AXC{$\stackrel{\Pi_{60}}{\vdots}$}\noLine
\UIC{$\varphi_{55} \to \varphi_{259}$}
\BIC{$\varphi_{259}$}
\BIC{$r \in \mathbb{R} \to \lvert r \rvert < 1 \to n \in \mathbb{N} \to \varphi_{46}$}
\end{ptree}
\noindent$\Pi_{62}$:
\begin{ptree}
\AXC{$y = z \to x = y \to x = z$\lbl{A6}}
\AXC{$(y = z \to x = y \to x = z) \to x = y \to y = z \to x = z$\lbl{A1}}
\BIC{$x = y \to y = z \to x = z$}
\UIC{$\forall z\,(x = y \to y = z \to x = z)$}
\UIC{$\forall y\,\forall z\,(x = y \to y = z \to x = z)$}
\UIC{$\forall x\,\forall y\,\forall z\,(x = y \to y = z \to x = z)$}
\AXC{$\forall x\,\forall y\,\forall z\,(x = y \to y = z \to x = z) \to \varphi_{185}$\lbl{A2}}
\BIC{$\varphi_{185}$}
\end{ptree}
\noindent$\Pi_{63}$:
\begin{ptree}
\AXC{$\forall n \in \mathbb{N}\,(S(n) \in \mathbb{N})$\lbl{[N2]}}
\AXC{$\forall n \in \mathbb{N}\,(S(n) \in \mathbb{N}) \to n \in \mathbb{N} \to S(n) \in \mathbb{N}$\lbl{A2}}
\BIC{$n \in \mathbb{N} \to S(n) \in \mathbb{N}$}
\AXC{$(n \in \mathbb{N} \to S(n) \in \mathbb{N}) \to \varphi_{116}$\lbl{A1}}
\BIC{$\varphi_{116}$}
\end{ptree}
\noindent$\Pi_{64}$:
\begin{ptree}
\AXC{$\forall r \in \mathbb{R}\,(\varphi_{58})$\lbl{[S24]}}
\AXC{$\forall r \in \mathbb{R}\,(\varphi_{58}) \to r \in \mathbb{R} \to \varphi_{58}$\lbl{A2}}
\BIC{$r \in \mathbb{R} \to \varphi_{58}$}
\AXC{$(r \in \mathbb{R} \to \varphi_{58}) \to \varphi_{194}$\lbl{A1}}
\BIC{$\varphi_{194}$}
\end{ptree}
\noindent$\Pi_{65}$:
\begin{ptree}
\AXC{$\forall a\,\forall n\,(\mathrm{Seq}(a) \to n \in \mathbb{N} \to a_{n} \in \mathbb{R})$\lbl{[S1]}}
\AXC{$\varphi_{102}$\lbl{A2}}
\BIC{$\forall n\,(\mathrm{Seq}({(r)^{\bullet}}) \to n \in \mathbb{N} \to {{r}^{n}} \in \mathbb{R})$}
\AXC{$\varphi_{119}$\lbl{A2}}
\BIC{$\mathrm{Seq}({(r)^{\bullet}}) \to S(n) \in \mathbb{N} \to {{r}^{S(n)}} \in \mathbb{R}$}
\AXC{$\varphi_{229}$\lbl{A1}}
\BIC{$\varphi_{191}$}
\end{ptree}
\noindent$\Pi_{66}$:
\begin{ptree}
\AXC{$r \in \mathbb{R} \to \lvert r \rvert < 1 \to n \in \mathbb{N} \to r \in \mathbb{R}$\lbl{A1}}
\AXC{$\stackrel{\Pi_{64}}{\vdots}$}\noLine
\UIC{$\varphi_{194}$}
\BIC{$r \in \mathbb{R} \to \lvert r \rvert < 1 \to n \in \mathbb{N} \to \varphi_{58}$}
\AXC{$\varphi_{195}$\lbl{A1}}
\BIC{$r \in \mathbb{R} \to \lvert r \rvert < 1 \to n \in \mathbb{N} \to \mathrm{Seq}({(r)^{\bullet}})$}
\AXC{$\stackrel{\Pi_{65}}{\vdots}$}\noLine
\UIC{$\varphi_{191}$}
\BIC{$\varphi_{143}$}
\end{ptree}
\noindent$\Pi_{67}$:
\begin{ptree}
\AXC{$\forall x \in \mathbb{R}\,(-1 \cdot x = -x)$\lbl{[E50]}}
\AXC{$\varphi_{134}$\lbl{A2}}
\BIC{${{r}^{S(n)}} \in \mathbb{R} \to -1 \cdot {{r}^{S(n)}} = -{{r}^{S(n)}}$}
\AXC{$({{r}^{S(n)}} \in \mathbb{R} \to -1 \cdot {{r}^{S(n)}} = -{{r}^{S(n)}}) \to \varphi_{182}$\lbl{A1}}
\BIC{$\varphi_{182}$}
\end{ptree}
\noindent$\Pi_{68}$:
\begin{ptree}
\AXC{$r \in \mathbb{R} \to \lvert r \rvert < 1 \to n \in \mathbb{N} \to n \in \mathbb{N}$\lbl{A1}}
\AXC{$\stackrel{\Pi_{63}}{\vdots}$}\noLine
\UIC{$\varphi_{116}$}
\BIC{$r \in \mathbb{R} \to \lvert r \rvert < 1 \to n \in \mathbb{N} \to S(n) \in \mathbb{N}$}
\AXC{$\stackrel{\Pi_{66}}{\vdots}$}\noLine
\UIC{$\varphi_{143}$}
\BIC{$r \in \mathbb{R} \to \lvert r \rvert < 1 \to n \in \mathbb{N} \to {{r}^{S(n)}} \in \mathbb{R}$}
\AXC{$\stackrel{\Pi_{67}}{\vdots}$}\noLine
\UIC{$\varphi_{182}$}
\BIC{$r \in \mathbb{R} \to \lvert r \rvert < 1 \to n \in \mathbb{N} \to -1 \cdot {{r}^{S(n)}} = -{{r}^{S(n)}}$}
\AXC{$\varphi_{286}$\lbl{A1}}
\BIC{$\varphi_{94} \to \varphi_{264}$}
\end{ptree}
\noindent$\Pi_{69}$:
\begin{ptree}
\AXC{$z = w \to \varphi_{25} \to \varphi_{24}$\lbl{A6}}
\UIC{$\forall w\,(z = w \to \varphi_{25} \to \varphi_{24})$}
\UIC{$\forall z\,\forall w\,(z = w \to \varphi_{25} \to \varphi_{24})$}
\AXC{$\varphi_{263}$\lbl{A2}}
\BIC{$\forall w\,(-1 \cdot {{r}^{S(n)}} = w \to \varphi_{94} \to \varphi_{24})$}
\AXC{$\varphi_{275}$\lbl{A2}}
\BIC{$-1 \cdot {{r}^{S(n)}} = -{{r}^{S(n)}} \to \varphi_{94} \to \varphi_{76}$}
\AXC{$x = x$\lbl{A5}}
\UIC{$\forall x\,(x = x)$}
\AXC{$\forall x\,(x = x) \to \varphi_{94}$\lbl{A2}}
\BIC{$\varphi_{94}$}
\AXC{$\stackrel{\Pi_{68}}{\vdots}$}\noLine
\UIC{$\varphi_{94} \to \varphi_{264}$}
\BIC{$\varphi_{264}$}
\BIC{$r \in \mathbb{R} \to \lvert r \rvert < 1 \to n \in \mathbb{N} \to \varphi_{76}$}
\end{ptree}
\noindent$\Pi_{70}$:
\begin{ptree}
\AXC{$y = z \to x = y \to x = z$\lbl{A6}}
\AXC{$(y = z \to x = y \to x = z) \to x = y \to y = z \to x = z$\lbl{A1}}
\BIC{$x = y \to y = z \to x = z$}
\UIC{$\forall z\,(x = y \to y = z \to x = z)$}
\UIC{$\forall y\,\forall z\,(x = y \to y = z \to x = z)$}
\UIC{$\forall x\,\forall y\,\forall z\,(x = y \to y = z \to x = z)$}
\AXC{$\forall x\,\forall y\,\forall z\,(x = y \to y = z \to x = z) \to \varphi_{185}$\lbl{A2}}
\BIC{$\varphi_{185}$}
\end{ptree}
\noindent$\Pi_{71}$:
\begin{ptree}
\AXC{$y = z \to x = y \to x = z$\lbl{A6}}
\AXC{$(y = z \to x = y \to x = z) \to x = y \to y = z \to x = z$\lbl{A1}}
\BIC{$x = y \to y = z \to x = z$}
\UIC{$\forall z\,(x = y \to y = z \to x = z)$}
\UIC{$\forall y\,\forall z\,(x = y \to y = z \to x = z)$}
\UIC{$\forall x\,\forall y\,\forall z\,(x = y \to y = z \to x = z)$}
\AXC{$\forall x\,\forall y\,\forall z\,(x = y \to y = z \to x = z) \to \varphi_{21}$\lbl{A2}}
\BIC{$\varphi_{21}$}
\end{ptree}
\noindent$\Pi_{72}$:
\begin{ptree}
\AXC{$\mathrm{NSeq}(\mathsf{s}) \land \forall k \in \mathbb{N}\,(\mathsf{s}(k) = S(k))$\lbl{[S14]}}
\AXC{$\mathrm{NSeq}(\mathsf{s}) \land \forall k \in \mathbb{N}\,(\mathsf{s}(k) = S(k)) \to \varphi_{31}$\lbl{A1}}
\BIC{$\varphi_{31}$}
\AXC{$\varphi_{31} \to \varphi_{161}$\lbl{A1}}
\BIC{$\varphi_{161}$}
\end{ptree}
\noindent$\Pi_{73}$:
\begin{ptree}
\AXC{$r \in \mathbb{R} \to \lvert r \rvert < 1 \to n \in \mathbb{N} \to n \in \mathbb{N}$\lbl{A1}}
\AXC{$\varphi_{15}$\lbl{A2}}
\AXC{$\stackrel{\Pi_{72}}{\vdots}$}\noLine
\UIC{$\varphi_{161}$}
\BIC{$r \in \mathbb{R} \to \lvert r \rvert < 1 \to n \in \mathbb{N} \to n \in \mathbb{N} \to \mathsf{s}(n) = S(n)$}
\AXC{$\varphi_{199}$\lbl{A1}}
\BIC{$\varphi_{123}$}
\BIC{$r \in \mathbb{R} \to \lvert r \rvert < 1 \to n \in \mathbb{N} \to \mathsf{s}(n) = S(n)$}
\AXC{$\varphi_{254}$\lbl{A1}}
\BIC{${{r}^{\mathsf{s}(n)}} = {{r}^{\mathsf{s}(n)}} \to \varphi_{204}$}
\end{ptree}
\noindent$\Pi_{74}$:
\begin{ptree}
\AXC{$x = x$\lbl{A5}}
\UIC{$\forall x\,(x = x)$}
\AXC{$\forall x\,(x = x) \to {{r}^{\mathsf{s}(n)}} = {{r}^{\mathsf{s}(n)}}$\lbl{A2}}
\BIC{${{r}^{\mathsf{s}(n)}} = {{r}^{\mathsf{s}(n)}}$}
\AXC{$\stackrel{\Pi_{73}}{\vdots}$}\noLine
\UIC{${{r}^{\mathsf{s}(n)}} = {{r}^{\mathsf{s}(n)}} \to \varphi_{204}$}
\BIC{$\varphi_{204}$}
\end{ptree}
\noindent$\Pi_{75}$:
\begin{ptree}
\AXC{$\forall r \in \mathbb{R}\,(\varphi_{58})$\lbl{[S24]}}
\AXC{$\forall r \in \mathbb{R}\,(\varphi_{58}) \to r \in \mathbb{R} \to \varphi_{58}$\lbl{A2}}
\BIC{$r \in \mathbb{R} \to \varphi_{58}$}
\AXC{$(r \in \mathbb{R} \to \varphi_{58}) \to \varphi_{194}$\lbl{A1}}
\BIC{$\varphi_{194}$}
\end{ptree}
\noindent$\Pi_{76}$:
\begin{ptree}
\AXC{$\varphi_{87}$\lbl{[S12]}}
\AXC{$\varphi_{87} \to \varphi_{117}$\lbl{A2}}
\BIC{$\varphi_{117}$}
\AXC{$\varphi_{117} \to \mathrm{Seq}({(r)^{\bullet}}) \to \mathrm{NSeq}(\mathsf{s}) \to \varphi_{71}$\lbl{A2}}
\BIC{$\mathrm{Seq}({(r)^{\bullet}}) \to \mathrm{NSeq}(\mathsf{s}) \to \varphi_{71}$}
\end{ptree}
\noindent$\Pi_{77}$:
\begin{ptree}
\AXC{$\stackrel{\Pi_{76}}{\vdots}$}\noLine
\UIC{$\mathrm{Seq}({(r)^{\bullet}}) \to \mathrm{NSeq}(\mathsf{s}) \to \varphi_{71}$}
\AXC{$(\mathrm{Seq}({(r)^{\bullet}}) \to \mathrm{NSeq}(\mathsf{s}) \to \varphi_{71}) \to \varphi_{231}$\lbl{A1}}
\BIC{$\varphi_{231}$}
\end{ptree}
\noindent$\Pi_{78}$:
\begin{ptree}
\AXC{$r \in \mathbb{R} \to \lvert r \rvert < 1 \to n \in \mathbb{N} \to r \in \mathbb{R}$\lbl{A1}}
\AXC{$\stackrel{\Pi_{75}}{\vdots}$}\noLine
\UIC{$\varphi_{194}$}
\BIC{$r \in \mathbb{R} \to \lvert r \rvert < 1 \to n \in \mathbb{N} \to \varphi_{58}$}
\AXC{$\varphi_{195}$\lbl{A1}}
\BIC{$r \in \mathbb{R} \to \lvert r \rvert < 1 \to n \in \mathbb{N} \to \mathrm{Seq}({(r)^{\bullet}})$}
\AXC{$\stackrel{\Pi_{77}}{\vdots}$}\noLine
\UIC{$\varphi_{231}$}
\BIC{$\varphi_{200}$}
\end{ptree}
\noindent$\Pi_{79}$:
\begin{ptree}
\AXC{$\mathrm{NSeq}(\mathsf{s}) \land \forall k \in \mathbb{N}\,(\mathsf{s}(k) = S(k))$\lbl{[S14]}}
\AXC{$\varphi_{103}$\lbl{A1}}
\BIC{$r \in \mathbb{R} \to \lvert r \rvert < 1 \to n \in \mathbb{N} \to \mathrm{NSeq}(\mathsf{s})$}
\AXC{$\stackrel{\Pi_{78}}{\vdots}$}\noLine
\UIC{$\varphi_{200}$}
\BIC{$r \in \mathbb{R} \to \lvert r \rvert < 1 \to n \in \mathbb{N} \to \varphi_{71}$}
\AXC{$(r \in \mathbb{R} \to \lvert r \rvert < 1 \to n \in \mathbb{N} \to \varphi_{71}) \to r \in \mathbb{R} \to \varphi_{81}$\lbl{A1}}
\BIC{$r \in \mathbb{R} \to \varphi_{81}$}
\end{ptree}
\noindent$\Pi_{80}$:
\begin{ptree}
\AXC{$r \in \mathbb{R} \to \lvert r \rvert < 1 \to n \in \mathbb{N} \to n \in \mathbb{N}$\lbl{A1}}
\AXC{$\varphi_{150}$\lbl{A2}}
\AXC{$\stackrel{\Pi_{79}}{\vdots}$}\noLine
\UIC{$r \in \mathbb{R} \to \varphi_{81}$}
\AXC{$(r \in \mathbb{R} \to \varphi_{81}) \to \varphi_{150} \to r \in \mathbb{R} \to \varphi_{64}$\lbl{A1}}
\BIC{$\varphi_{150} \to r \in \mathbb{R} \to \varphi_{64}$}
\BIC{$r \in \mathbb{R} \to \varphi_{64}$}
\AXC{$\varphi_{236}$\lbl{A1}}
\BIC{$(r \in \mathbb{R} \to \lvert r \rvert < 1 \to n \in \mathbb{N} \to n \in \mathbb{N}) \to \varphi_{67}$}
\BIC{$\varphi_{67}$}
\end{ptree}
\noindent$\Pi_{81}$:
\begin{ptree}
\AXC{$z = w \to {{r}^{\mathsf{s}(n)}} = {{r}^{z}} \to {{r}^{\mathsf{s}(n)}} = {{r}^{w}}$\lbl{A6}}
\UIC{$\forall w\,(z = w \to {{r}^{\mathsf{s}(n)}} = {{r}^{z}} \to {{r}^{\mathsf{s}(n)}} = {{r}^{w}})$}
\UIC{$\forall z\,\forall w\,(z = w \to {{r}^{\mathsf{s}(n)}} = {{r}^{z}} \to {{r}^{\mathsf{s}(n)}} = {{r}^{w}})$}
\AXC{$\varphi_{208}$\lbl{A2}}
\BIC{$\forall w\,(\mathsf{s}(n) = w \to {{r}^{\mathsf{s}(n)}} = {{r}^{\mathsf{s}(n)}} \to {{r}^{\mathsf{s}(n)}} = {{r}^{w}})$}
\AXC{$\varphi_{220}$\lbl{A2}}
\BIC{$\mathsf{s}(n) = S(n) \to {{r}^{\mathsf{s}(n)}} = {{r}^{\mathsf{s}(n)}} \to {{r}^{\mathsf{s}(n)}} = {{r}^{S(n)}}$}
\AXC{$\stackrel{\Pi_{74}}{\vdots}$}\noLine
\UIC{$\varphi_{204}$}
\BIC{$r \in \mathbb{R} \to \lvert r \rvert < 1 \to n \in \mathbb{N} \to {{r}^{\mathsf{s}(n)}} = {{r}^{S(n)}}$}
\AXC{$\stackrel{\Pi_{80}}{\vdots}$}\noLine
\UIC{$\varphi_{67}$}
\AXC{$\varphi_{67} \to \varphi_{252}$\lbl{A1}}
\BIC{$\varphi_{252}$}
\BIC{$\varphi_{160} \to \varphi_{49}$}
\end{ptree}
\noindent$\Pi_{82}$:
\begin{ptree}
\AXC{$x = x$\lbl{A5}}
\UIC{$\forall x\,(x = x)$}
\AXC{$\forall x\,(x = x) \to \varphi_{126}$\lbl{A2}}
\BIC{$\varphi_{126}$}
\AXC{$\stackrel{\Pi_{71}}{\vdots}$}\noLine
\UIC{$\varphi_{21}$}
\AXC{$\varphi_{21} \to \varphi_{120}$\lbl{A2}}
\BIC{$\varphi_{120}$}
\AXC{$\varphi_{120} \to \varphi_{160}$\lbl{A2}}
\BIC{$\varphi_{160}$}
\AXC{$\stackrel{\Pi_{81}}{\vdots}$}\noLine
\UIC{$\varphi_{160} \to \varphi_{49}$}
\BIC{$\varphi_{49}$}
\AXC{$\varphi_{49} \to \varphi_{126} \to \varphi_{269}$\lbl{A1}}
\BIC{$\varphi_{126} \to \varphi_{269}$}
\BIC{$\varphi_{269}$}
\end{ptree}
\noindent$\Pi_{83}$:
\begin{ptree}
\AXC{$y = z \to x = y \to x = z$\lbl{A6}}
\AXC{$(y = z \to x = y \to x = z) \to x = y \to y = z \to x = z$\lbl{A1}}
\BIC{$x = y \to y = z \to x = z$}
\UIC{$\forall z\,(x = y \to y = z \to x = z)$}
\UIC{$\forall y\,\forall z\,(x = y \to y = z \to x = z)$}
\UIC{$\forall x\,\forall y\,\forall z\,(x = y \to y = z \to x = z)$}
\AXC{$\forall x\,\forall y\,\forall z\,(x = y \to y = z \to x = z) \to \varphi_{185}$\lbl{A2}}
\BIC{$\varphi_{185}$}
\end{ptree}
\noindent$\Pi_{84}$:
\begin{ptree}
\AXC{$\forall r \in \mathbb{R}\,(\varphi_{58})$\lbl{[S24]}}
\AXC{$\forall r \in \mathbb{R}\,(\varphi_{58}) \to r \in \mathbb{R} \to \varphi_{58}$\lbl{A2}}
\BIC{$r \in \mathbb{R} \to \varphi_{58}$}
\AXC{$(r \in \mathbb{R} \to \varphi_{58}) \to \varphi_{194}$\lbl{A1}}
\BIC{$\varphi_{194}$}
\end{ptree}
\noindent$\Pi_{85}$:
\begin{ptree}
\AXC{$\varphi_{87}$\lbl{[S12]}}
\AXC{$\varphi_{87} \to \varphi_{117}$\lbl{A2}}
\BIC{$\varphi_{117}$}
\AXC{$\varphi_{117} \to \mathrm{Seq}({(r)^{\bullet}}) \to \mathrm{NSeq}(\mathsf{s}) \to \varphi_{71}$\lbl{A2}}
\BIC{$\mathrm{Seq}({(r)^{\bullet}}) \to \mathrm{NSeq}(\mathsf{s}) \to \varphi_{71}$}
\end{ptree}
\noindent$\Pi_{86}$:
\begin{ptree}
\AXC{$\stackrel{\Pi_{85}}{\vdots}$}\noLine
\UIC{$\mathrm{Seq}({(r)^{\bullet}}) \to \mathrm{NSeq}(\mathsf{s}) \to \varphi_{71}$}
\AXC{$(\mathrm{Seq}({(r)^{\bullet}}) \to \mathrm{NSeq}(\mathsf{s}) \to \varphi_{71}) \to \varphi_{231}$\lbl{A1}}
\BIC{$\varphi_{231}$}
\end{ptree}
\noindent$\Pi_{87}$:
\begin{ptree}
\AXC{$r \in \mathbb{R} \to \lvert r \rvert < 1 \to n \in \mathbb{N} \to r \in \mathbb{R}$\lbl{A1}}
\AXC{$\stackrel{\Pi_{84}}{\vdots}$}\noLine
\UIC{$\varphi_{194}$}
\BIC{$r \in \mathbb{R} \to \lvert r \rvert < 1 \to n \in \mathbb{N} \to \varphi_{58}$}
\AXC{$\varphi_{195}$\lbl{A1}}
\BIC{$r \in \mathbb{R} \to \lvert r \rvert < 1 \to n \in \mathbb{N} \to \mathrm{Seq}({(r)^{\bullet}})$}
\AXC{$\stackrel{\Pi_{86}}{\vdots}$}\noLine
\UIC{$\varphi_{231}$}
\BIC{$\varphi_{200}$}
\end{ptree}
\noindent$\Pi_{88}$:
\begin{ptree}
\AXC{$\mathrm{NSeq}(\mathsf{s}) \land \forall k \in \mathbb{N}\,(\mathsf{s}(k) = S(k))$\lbl{[S14]}}
\AXC{$\varphi_{103}$\lbl{A1}}
\BIC{$r \in \mathbb{R} \to \lvert r \rvert < 1 \to n \in \mathbb{N} \to \mathrm{NSeq}(\mathsf{s})$}
\AXC{$\stackrel{\Pi_{87}}{\vdots}$}\noLine
\UIC{$\varphi_{200}$}
\BIC{$r \in \mathbb{R} \to \lvert r \rvert < 1 \to n \in \mathbb{N} \to \varphi_{71}$}
\AXC{$(r \in \mathbb{R} \to \lvert r \rvert < 1 \to n \in \mathbb{N} \to \varphi_{71}) \to \varphi_{10}$\lbl{A1}}
\BIC{$\varphi_{10}$}
\end{ptree}
\noindent$\Pi_{89}$:
\begin{ptree}
\AXC{$1 \in \mathbb{R}$\lbl{(R1)}}
\AXC{$\forall x \in \mathbb{R}\,(-x \in \mathbb{R})$\lbl{(R1)}}
\AXC{$\forall x \in \mathbb{R}\,(-x \in \mathbb{R}) \to 1 \in \mathbb{R} \to -1 \in \mathbb{R}$\lbl{A2}}
\BIC{$1 \in \mathbb{R} \to -1 \in \mathbb{R}$}
\AXC{$\varphi_{109}$\lbl{A1}}
\BIC{$1 \in \mathbb{R} \to r \in \mathbb{R} \to \lvert r \rvert < 1 \to n \in \mathbb{N} \to -1 \in \mathbb{R}$}
\BIC{$r \in \mathbb{R} \to \lvert r \rvert < 1 \to n \in \mathbb{N} \to -1 \in \mathbb{R}$}
\end{ptree}
\noindent$\Pi_{90}$:
\begin{ptree}
\AXC{$\stackrel{\Pi_{89}}{\vdots}$}\noLine
\UIC{$r \in \mathbb{R} \to \lvert r \rvert < 1 \to n \in \mathbb{N} \to -1 \in \mathbb{R}$}
\AXC{$\varphi_{90}$\lbl{[S3]}}
\AXC{$\varphi_{90} \to \varphi_{127}$\lbl{A2}}
\BIC{$\varphi_{127}$}
\AXC{$\varphi_{127} \to -1 \in \mathbb{R} \to \varphi_{178}$\lbl{A2}}
\BIC{$-1 \in \mathbb{R} \to \varphi_{178}$}
\AXC{$(-1 \in \mathbb{R} \to \varphi_{178}) \to \varphi_{249}$\lbl{A1}}
\BIC{$\varphi_{249}$}
\BIC{$\varphi_{10} \to r \in \mathbb{R} \to \lvert r \rvert < 1 \to \varphi_{165}$}
\end{ptree}
\noindent$\Pi_{91}$:
\begin{ptree}
\AXC{$\stackrel{\Pi_{88}}{\vdots}$}\noLine
\UIC{$\varphi_{10}$}
\AXC{$\stackrel{\Pi_{90}}{\vdots}$}\noLine
\UIC{$\varphi_{10} \to r \in \mathbb{R} \to \lvert r \rvert < 1 \to \varphi_{165}$}
\BIC{$r \in \mathbb{R} \to \lvert r \rvert < 1 \to \varphi_{165}$}
\AXC{$\varphi_{251}$\lbl{A1}}
\BIC{$r \in \mathbb{R} \to \lvert r \rvert < 1 \to n \in \mathbb{N} \to \forall n \in \mathbb{N}\,(\varphi_{32})$}
\AXC{$(r \in \mathbb{R} \to \lvert r \rvert < 1 \to n \in \mathbb{N} \to \forall n \in \mathbb{N}\,(\varphi_{32})) \to \varphi_{255}$\lbl{A1}}
\BIC{$\varphi_{255}$}
\end{ptree}
\noindent$\Pi_{92}$:
\begin{ptree}
\AXC{$\forall n \in \mathbb{N}\,(\varphi_{32}) \to n \in \mathbb{N} \to \varphi_{32}$\lbl{A2}}
\AXC{$\stackrel{\Pi_{91}}{\vdots}$}\noLine
\UIC{$\varphi_{255}$}
\BIC{$r \in \mathbb{R} \to \lvert r \rvert < 1 \to n \in \mathbb{N} \to n \in \mathbb{N} \to \varphi_{32}$}
\AXC{$(r \in \mathbb{R} \to \lvert r \rvert < 1 \to n \in \mathbb{N} \to n \in \mathbb{N} \to \varphi_{32}) \to \varphi_{189}$\lbl{A1}}
\BIC{$\varphi_{189}$}
\end{ptree}
\noindent$\Pi_{93}$:
\begin{ptree}
\AXC{$x = x$\lbl{A5}}
\UIC{$\forall x\,(x = x)$}
\AXC{$\forall x\,(x = x) \to \varphi_{137}$\lbl{A2}}
\BIC{$\varphi_{137}$}
\AXC{$r \in \mathbb{R} \to \lvert r \rvert < 1 \to n \in \mathbb{N} \to n \in \mathbb{N}$\lbl{A1}}
\AXC{$\stackrel{\Pi_{92}}{\vdots}$}\noLine
\UIC{$\varphi_{189}$}
\BIC{$r \in \mathbb{R} \to \lvert r \rvert < 1 \to n \in \mathbb{N} \to \varphi_{32}$}
\AXC{$\varphi_{295}$\lbl{A1}}
\BIC{$\varphi_{137} \to (\varphi_{32} \to \varphi_{137} \to \varphi_{136}) \to r \in \mathbb{R} \to \lvert r \rvert < 1 \to n \in \mathbb{N} \to \varphi_{136}$}
\BIC{$(\varphi_{32} \to \varphi_{137} \to \varphi_{136}) \to r \in \mathbb{R} \to \lvert r \rvert < 1 \to n \in \mathbb{N} \to \varphi_{136}$}
\end{ptree}
\noindent$\Pi_{94}$:
\begin{ptree}
\AXC{$z = w \to \varphi_{51} \to \varphi_{50}$\lbl{A6}}
\UIC{$\forall w\,(z = w \to \varphi_{51} \to \varphi_{50})$}
\UIC{$\forall z\,\forall w\,(z = w \to \varphi_{51} \to \varphi_{50})$}
\AXC{$\varphi_{271}$\lbl{A2}}
\BIC{$\forall w\,(-1 \odot {(r)^{\bullet}} \circ \mathsf{s}(n) = w \to \varphi_{137} \to \varphi_{50})$}
\AXC{$\varphi_{287}$\lbl{A2}}
\BIC{$\varphi_{32} \to \varphi_{137} \to \varphi_{136}$}
\AXC{$\stackrel{\Pi_{93}}{\vdots}$}\noLine
\UIC{$(\varphi_{32} \to \varphi_{137} \to \varphi_{136}) \to r \in \mathbb{R} \to \lvert r \rvert < 1 \to n \in \mathbb{N} \to \varphi_{136}$}
\BIC{$r \in \mathbb{R} \to \lvert r \rvert < 1 \to n \in \mathbb{N} \to \varphi_{136}$}
\end{ptree}
\noindent$\Pi_{95}$:
\begin{ptree}
\AXC{$y = z \to x = y \to x = z$\lbl{A6}}
\AXC{$(y = z \to x = y \to x = z) \to x = y \to y = z \to x = z$\lbl{A1}}
\BIC{$x = y \to y = z \to x = z$}
\UIC{$\forall z\,(x = y \to y = z \to x = z)$}
\UIC{$\forall y\,\forall z\,(x = y \to y = z \to x = z)$}
\UIC{$\forall x\,\forall y\,\forall z\,(x = y \to y = z \to x = z)$}
\AXC{$\forall x\,\forall y\,\forall z\,(x = y \to y = z \to x = z) \to \varphi_{185}$\lbl{A2}}
\BIC{$\varphi_{185}$}
\end{ptree}
\noindent$\Pi_{96}$:
\begin{ptree}
\AXC{$z = w \to \varphi_{156} \to \varphi_{155}$\lbl{A6}}
\UIC{$\forall w\,(z = w \to \varphi_{156} \to \varphi_{155})$}
\UIC{$\forall z\,\forall w\,(z = w \to \varphi_{156} \to \varphi_{155})$}
\AXC{$\forall z\,\forall w\,(z = w \to \varphi_{156} \to \varphi_{155}) \to \forall w\,(\overline{r}(n) = w \to \varphi_{172} \to \varphi_{155})$\lbl{A2}}
\BIC{$\forall w\,(\overline{r}(n) = w \to \varphi_{172} \to \varphi_{155})$}
\end{ptree}
\noindent$\Pi_{97}$:
\begin{ptree}
\AXC{$\stackrel{\Pi_{96}}{\vdots}$}\noLine
\UIC{$\forall w\,(\overline{r}(n) = w \to \varphi_{172} \to \varphi_{155})$}
\AXC{$\forall w\,(\overline{r}(n) = w \to \varphi_{172} \to \varphi_{155}) \to \overline{r}(n) = r \to \varphi_{172} \to \varphi_{154}$\lbl{A2}}
\BIC{$\overline{r}(n) = r \to \varphi_{172} \to \varphi_{154}$}
\end{ptree}
\noindent$\Pi_{98}$:
\begin{ptree}
\AXC{$r \in \mathbb{R} \to \lvert r \rvert < 1 \to n \in \mathbb{N} \to r \in \mathbb{R}$\lbl{A1}}
\AXC{$\forall c \in \mathbb{R}\,(\mathrm{Seq}(\overline{c}) \land \forall n \in \mathbb{N}\,(\overline{c}(n) = c))$\lbl{[S10]}}
\AXC{$\varphi_{152}$\lbl{A2}}
\BIC{$r \in \mathbb{R} \to \mathrm{Seq}(\overline{r}) \land \forall n \in \mathbb{N}\,(\overline{r}(n) = r)$}
\AXC{$\varphi_{214}$\lbl{A1}}
\BIC{$(r \in \mathbb{R} \to \lvert r \rvert < 1 \to n \in \mathbb{N} \to r \in \mathbb{R}) \to r \in \mathbb{R} \to \varphi_{37}$}
\BIC{$r \in \mathbb{R} \to \varphi_{37}$}
\end{ptree}
\noindent$\Pi_{99}$:
\begin{ptree}
\AXC{$\forall n \in \mathbb{N}\,(\overline{r}(n) = r) \to n \in \mathbb{N} \to \overline{r}(n) = r$\lbl{A2}}
\AXC{$\stackrel{\Pi_{98}}{\vdots}$}\noLine
\UIC{$r \in \mathbb{R} \to \varphi_{37}$}
\AXC{$\varphi_{187}$\lbl{A1}}
\BIC{$r \in \mathbb{R} \to \lvert r \rvert < 1 \to n \in \mathbb{N} \to \forall n \in \mathbb{N}\,(\overline{r}(n) = r)$}
\AXC{$\varphi_{216}$\lbl{A1}}
\BIC{$\varphi_{162}$}
\BIC{$r \in \mathbb{R} \to \lvert r \rvert < 1 \to n \in \mathbb{N} \to n \in \mathbb{N} \to \overline{r}(n) = r$}
\end{ptree}
\noindent$\Pi_{100}$:
\begin{ptree}
\AXC{$r \in \mathbb{R} \to \lvert r \rvert < 1 \to n \in \mathbb{N} \to n \in \mathbb{N}$\lbl{A1}}
\AXC{$\stackrel{\Pi_{99}}{\vdots}$}\noLine
\UIC{$r \in \mathbb{R} \to \lvert r \rvert < 1 \to n \in \mathbb{N} \to n \in \mathbb{N} \to \overline{r}(n) = r$}
\AXC{$\varphi_{201}$\lbl{A1}}
\BIC{$\varphi_{124}$}
\BIC{$r \in \mathbb{R} \to \lvert r \rvert < 1 \to n \in \mathbb{N} \to \overline{r}(n) = r$}
\end{ptree}
\noindent$\Pi_{101}$:
\begin{ptree}
\AXC{$\stackrel{\Pi_{100}}{\vdots}$}\noLine
\UIC{$r \in \mathbb{R} \to \lvert r \rvert < 1 \to n \in \mathbb{N} \to \overline{r}(n) = r$}
\AXC{$(r \in \mathbb{R} \to \lvert r \rvert < 1 \to n \in \mathbb{N} \to \overline{r}(n) = r) \to \varphi_{172} \to \varphi_{279}$\lbl{A1}}
\BIC{$\varphi_{172} \to \varphi_{279}$}
\end{ptree}
\noindent$\Pi_{102}$:
\begin{ptree}
\AXC{$\stackrel{\Pi_{97}}{\vdots}$}\noLine
\UIC{$\overline{r}(n) = r \to \varphi_{172} \to \varphi_{154}$}
\AXC{$x = x$\lbl{A5}}
\UIC{$\forall x\,(x = x)$}
\AXC{$\forall x\,(x = x) \to \varphi_{172}$\lbl{A2}}
\BIC{$\varphi_{172}$}
\AXC{$\stackrel{\Pi_{101}}{\vdots}$}\noLine
\UIC{$\varphi_{172} \to \varphi_{279}$}
\BIC{$\varphi_{279}$}
\BIC{$r \in \mathbb{R} \to \lvert r \rvert < 1 \to n \in \mathbb{N} \to \varphi_{154}$}
\end{ptree}
\noindent$\Pi_{103}$:
\begin{ptree}
\AXC{$y = z \to x = y \to x = z$\lbl{A6}}
\AXC{$(y = z \to x = y \to x = z) \to x = y \to y = z \to x = z$\lbl{A1}}
\BIC{$x = y \to y = z \to x = z$}
\UIC{$\forall z\,(x = y \to y = z \to x = z)$}
\UIC{$\forall y\,\forall z\,(x = y \to y = z \to x = z)$}
\UIC{$\forall x\,\forall y\,\forall z\,(x = y \to y = z \to x = z)$}
\AXC{$\forall x\,\forall y\,\forall z\,(x = y \to y = z \to x = z) \to \varphi_{185}$\lbl{A2}}
\BIC{$\varphi_{185}$}
\end{ptree}
\noindent$\Pi_{104}$:
\begin{ptree}
\AXC{$\forall r \in \mathbb{R}\,(\varphi_{58})$\lbl{[S24]}}
\AXC{$\forall r \in \mathbb{R}\,(\varphi_{58}) \to r \in \mathbb{R} \to \varphi_{58}$\lbl{A2}}
\BIC{$r \in \mathbb{R} \to \varphi_{58}$}
\AXC{$(r \in \mathbb{R} \to \varphi_{58}) \to \varphi_{194}$\lbl{A1}}
\BIC{$\varphi_{194}$}
\end{ptree}
\noindent$\Pi_{105}$:
\begin{ptree}
\AXC{$\varphi_{87}$\lbl{[S12]}}
\AXC{$\varphi_{87} \to \varphi_{117}$\lbl{A2}}
\BIC{$\varphi_{117}$}
\AXC{$\varphi_{117} \to \mathrm{Seq}({(r)^{\bullet}}) \to \mathrm{NSeq}(\mathsf{s}) \to \varphi_{71}$\lbl{A2}}
\BIC{$\mathrm{Seq}({(r)^{\bullet}}) \to \mathrm{NSeq}(\mathsf{s}) \to \varphi_{71}$}
\end{ptree}
\noindent$\Pi_{106}$:
\begin{ptree}
\AXC{$\stackrel{\Pi_{105}}{\vdots}$}\noLine
\UIC{$\mathrm{Seq}({(r)^{\bullet}}) \to \mathrm{NSeq}(\mathsf{s}) \to \varphi_{71}$}
\AXC{$(\mathrm{Seq}({(r)^{\bullet}}) \to \mathrm{NSeq}(\mathsf{s}) \to \varphi_{71}) \to \varphi_{231}$\lbl{A1}}
\BIC{$\varphi_{231}$}
\end{ptree}
\noindent$\Pi_{107}$:
\begin{ptree}
\AXC{$r \in \mathbb{R} \to \lvert r \rvert < 1 \to n \in \mathbb{N} \to r \in \mathbb{R}$\lbl{A1}}
\AXC{$\stackrel{\Pi_{104}}{\vdots}$}\noLine
\UIC{$\varphi_{194}$}
\BIC{$r \in \mathbb{R} \to \lvert r \rvert < 1 \to n \in \mathbb{N} \to \varphi_{58}$}
\AXC{$\varphi_{195}$\lbl{A1}}
\BIC{$r \in \mathbb{R} \to \lvert r \rvert < 1 \to n \in \mathbb{N} \to \mathrm{Seq}({(r)^{\bullet}})$}
\AXC{$\stackrel{\Pi_{106}}{\vdots}$}\noLine
\UIC{$\varphi_{231}$}
\BIC{$\varphi_{200}$}
\end{ptree}
\noindent$\Pi_{108}$:
\begin{ptree}
\AXC{$\mathrm{NSeq}(\mathsf{s}) \land \forall k \in \mathbb{N}\,(\mathsf{s}(k) = S(k))$\lbl{[S14]}}
\AXC{$\varphi_{103}$\lbl{A1}}
\BIC{$r \in \mathbb{R} \to \lvert r \rvert < 1 \to n \in \mathbb{N} \to \mathrm{NSeq}(\mathsf{s})$}
\AXC{$\stackrel{\Pi_{107}}{\vdots}$}\noLine
\UIC{$\varphi_{200}$}
\BIC{$r \in \mathbb{R} \to \lvert r \rvert < 1 \to n \in \mathbb{N} \to \varphi_{71}$}
\AXC{$(r \in \mathbb{R} \to \lvert r \rvert < 1 \to n \in \mathbb{N} \to \varphi_{71}) \to \varphi_{10}$\lbl{A1}}
\BIC{$\varphi_{10}$}
\end{ptree}
\noindent$\Pi_{109}$:
\begin{ptree}
\AXC{$\stackrel{\Pi_{89}}{\vdots}$}\noLine
\UIC{$r \in \mathbb{R} \to \lvert r \rvert < 1 \to n \in \mathbb{N} \to -1 \in \mathbb{R}$}
\AXC{$\varphi_{90}$\lbl{[S3]}}
\AXC{$\varphi_{90} \to \varphi_{127}$\lbl{A2}}
\BIC{$\varphi_{127}$}
\AXC{$\varphi_{127} \to -1 \in \mathbb{R} \to \varphi_{178}$\lbl{A2}}
\BIC{$-1 \in \mathbb{R} \to \varphi_{178}$}
\AXC{$(-1 \in \mathbb{R} \to \varphi_{178}) \to \varphi_{249}$\lbl{A1}}
\BIC{$\varphi_{249}$}
\BIC{$\varphi_{10} \to r \in \mathbb{R} \to \lvert r \rvert < 1 \to \varphi_{165}$}
\end{ptree}
\noindent$\Pi_{110}$:
\begin{ptree}
\AXC{$\stackrel{\Pi_{108}}{\vdots}$}\noLine
\UIC{$\varphi_{10}$}
\AXC{$\stackrel{\Pi_{109}}{\vdots}$}\noLine
\UIC{$\varphi_{10} \to r \in \mathbb{R} \to \lvert r \rvert < 1 \to \varphi_{165}$}
\BIC{$r \in \mathbb{R} \to \lvert r \rvert < 1 \to \varphi_{165}$}
\AXC{$(r \in \mathbb{R} \to \lvert r \rvert < 1 \to \varphi_{165}) \to \varphi_{38}$\lbl{A1}}
\BIC{$\varphi_{38}$}
\end{ptree}
\noindent$\Pi_{111}$:
\begin{ptree}
\AXC{$r \in \mathbb{R} \to \lvert r \rvert < 1 \to n \in \mathbb{N} \to r \in \mathbb{R}$\lbl{A1}}
\AXC{$\forall c \in \mathbb{R}\,(\mathrm{Seq}(\overline{c}) \land \forall n \in \mathbb{N}\,(\overline{c}(n) = c))$\lbl{[S10]}}
\AXC{$\varphi_{152}$\lbl{A2}}
\BIC{$r \in \mathbb{R} \to \mathrm{Seq}(\overline{r}) \land \forall n \in \mathbb{N}\,(\overline{r}(n) = r)$}
\AXC{$\varphi_{214}$\lbl{A1}}
\BIC{$(r \in \mathbb{R} \to \lvert r \rvert < 1 \to n \in \mathbb{N} \to r \in \mathbb{R}) \to r \in \mathbb{R} \to \varphi_{37}$}
\BIC{$r \in \mathbb{R} \to \varphi_{37}$}
\end{ptree}
\noindent$\Pi_{112}$:
\begin{ptree}
\AXC{$\varphi_{100}$\lbl{[S2]}}
\AXC{$\varphi_{100} \to \varphi_{141}$\lbl{A2}}
\BIC{$\varphi_{141}$}
\AXC{$\varphi_{141} \to \mathrm{Seq}(\overline{r}) \to \varphi_{221}$\lbl{A2}}
\BIC{$\mathrm{Seq}(\overline{r}) \to \varphi_{221}$}
\AXC{$(\mathrm{Seq}(\overline{r}) \to \varphi_{221}) \to \varphi_{262}$\lbl{A1}}
\BIC{$\varphi_{262}$}
\end{ptree}
\noindent$\Pi_{113}$:
\begin{ptree}
\AXC{$\stackrel{\Pi_{110}}{\vdots}$}\noLine
\UIC{$\varphi_{38}$}
\AXC{$\stackrel{\Pi_{111}}{\vdots}$}\noLine
\UIC{$r \in \mathbb{R} \to \varphi_{37}$}
\AXC{$\varphi_{173}$\lbl{A1}}
\BIC{$r \in \mathbb{R} \to \lvert r \rvert < 1 \to n \in \mathbb{N} \to \mathrm{Seq}(\overline{r})$}
\AXC{$\stackrel{\Pi_{112}}{\vdots}$}\noLine
\UIC{$\varphi_{262}$}
\BIC{$\varphi_{38} \to r \in \mathbb{R} \to \lvert r \rvert < 1 \to \varphi_{205}$}
\BIC{$r \in \mathbb{R} \to \lvert r \rvert < 1 \to \varphi_{205}$}
\AXC{$\varphi_{268}$\lbl{A1}}
\BIC{$r \in \mathbb{R} \to \lvert r \rvert < 1 \to n \in \mathbb{N} \to \forall n \in \mathbb{N}\,(\varphi_{125})$}
\end{ptree}
\noindent$\Pi_{114}$:
\begin{ptree}
\AXC{$\stackrel{\Pi_{113}}{\vdots}$}\noLine
\UIC{$r \in \mathbb{R} \to \lvert r \rvert < 1 \to n \in \mathbb{N} \to \forall n \in \mathbb{N}\,(\varphi_{125})$}
\AXC{$(r \in \mathbb{R} \to \lvert r \rvert < 1 \to n \in \mathbb{N} \to \forall n \in \mathbb{N}\,(\varphi_{125})) \to \varphi_{270}$\lbl{A1}}
\BIC{$\varphi_{270}$}
\end{ptree}
\noindent$\Pi_{115}$:
\begin{ptree}
\AXC{$\forall n \in \mathbb{N}\,(\varphi_{125}) \to n \in \mathbb{N} \to \varphi_{125}$\lbl{A2}}
\AXC{$\stackrel{\Pi_{114}}{\vdots}$}\noLine
\UIC{$\varphi_{270}$}
\BIC{$r \in \mathbb{R} \to \lvert r \rvert < 1 \to n \in \mathbb{N} \to n \in \mathbb{N} \to \varphi_{125}$}
\AXC{$(r \in \mathbb{R} \to \lvert r \rvert < 1 \to n \in \mathbb{N} \to n \in \mathbb{N} \to \varphi_{125}) \to \varphi_{210}$\lbl{A1}}
\BIC{$\varphi_{210}$}
\end{ptree}
\noindent$\Pi_{116}$:
\begin{ptree}
\AXC{$x = x$\lbl{A5}}
\UIC{$\forall x\,(x = x)$}
\AXC{$\forall x\,(x = x) \to \varphi_{163}$\lbl{A2}}
\BIC{$\varphi_{163}$}
\AXC{$r \in \mathbb{R} \to \lvert r \rvert < 1 \to n \in \mathbb{N} \to n \in \mathbb{N}$\lbl{A1}}
\AXC{$\stackrel{\Pi_{115}}{\vdots}$}\noLine
\UIC{$\varphi_{210}$}
\BIC{$r \in \mathbb{R} \to \lvert r \rvert < 1 \to n \in \mathbb{N} \to \varphi_{125}$}
\AXC{$\varphi_{299}$\lbl{A1}}
\BIC{$\varphi_{163} \to (\varphi_{125} \to \varphi_{163} \to \varphi_{168}) \to r \in \mathbb{R} \to \lvert r \rvert < 1 \to n \in \mathbb{N} \to \varphi_{168}$}
\BIC{$(\varphi_{125} \to \varphi_{163} \to \varphi_{168}) \to r \in \mathbb{R} \to \lvert r \rvert < 1 \to n \in \mathbb{N} \to \varphi_{168}$}
\end{ptree}
\noindent$\Pi_{117}$:
\begin{ptree}
\AXC{$z = w \to \varphi_{62} \to \varphi_{61}$\lbl{A6}}
\UIC{$\forall w\,(z = w \to \varphi_{62} \to \varphi_{61})$}
\UIC{$\forall z\,\forall w\,(z = w \to \varphi_{62} \to \varphi_{61})$}
\AXC{$\varphi_{280}$\lbl{A2}}
\BIC{$\forall w\,(\overline{r} \oplus -1 \odot {(r)^{\bullet}} \circ \mathsf{s}(n) = w \to \varphi_{163} \to \varphi_{61})$}
\AXC{$\varphi_{293}$\lbl{A2}}
\BIC{$\varphi_{125} \to \varphi_{163} \to \varphi_{168}$}
\AXC{$\stackrel{\Pi_{116}}{\vdots}$}\noLine
\UIC{$(\varphi_{125} \to \varphi_{163} \to \varphi_{168}) \to r \in \mathbb{R} \to \lvert r \rvert < 1 \to n \in \mathbb{N} \to \varphi_{168}$}
\BIC{$r \in \mathbb{R} \to \lvert r \rvert < 1 \to n \in \mathbb{N} \to \varphi_{168}$}
\end{ptree}
\noindent$\Pi_{118}$:
\begin{ptree}
\AXC{$\forall r \in \mathbb{R}\,(\varphi_{58})$\lbl{[S24]}}
\AXC{$\forall r \in \mathbb{R}\,(\varphi_{58}) \to r \in \mathbb{R} \to \varphi_{58}$\lbl{A2}}
\BIC{$r \in \mathbb{R} \to \varphi_{58}$}
\AXC{$(r \in \mathbb{R} \to \varphi_{58}) \to \varphi_{194}$\lbl{A1}}
\BIC{$\varphi_{194}$}
\end{ptree}
\noindent$\Pi_{119}$:
\begin{ptree}
\AXC{$\varphi_{87}$\lbl{[S12]}}
\AXC{$\varphi_{87} \to \varphi_{117}$\lbl{A2}}
\BIC{$\varphi_{117}$}
\AXC{$\varphi_{117} \to \mathrm{Seq}({(r)^{\bullet}}) \to \mathrm{NSeq}(\mathsf{s}) \to \varphi_{71}$\lbl{A2}}
\BIC{$\mathrm{Seq}({(r)^{\bullet}}) \to \mathrm{NSeq}(\mathsf{s}) \to \varphi_{71}$}
\end{ptree}
\noindent$\Pi_{120}$:
\begin{ptree}
\AXC{$\stackrel{\Pi_{119}}{\vdots}$}\noLine
\UIC{$\mathrm{Seq}({(r)^{\bullet}}) \to \mathrm{NSeq}(\mathsf{s}) \to \varphi_{71}$}
\AXC{$(\mathrm{Seq}({(r)^{\bullet}}) \to \mathrm{NSeq}(\mathsf{s}) \to \varphi_{71}) \to \varphi_{231}$\lbl{A1}}
\BIC{$\varphi_{231}$}
\end{ptree}
\noindent$\Pi_{121}$:
\begin{ptree}
\AXC{$r \in \mathbb{R} \to \lvert r \rvert < 1 \to n \in \mathbb{N} \to r \in \mathbb{R}$\lbl{A1}}
\AXC{$\stackrel{\Pi_{118}}{\vdots}$}\noLine
\UIC{$\varphi_{194}$}
\BIC{$r \in \mathbb{R} \to \lvert r \rvert < 1 \to n \in \mathbb{N} \to \varphi_{58}$}
\AXC{$\varphi_{195}$\lbl{A1}}
\BIC{$r \in \mathbb{R} \to \lvert r \rvert < 1 \to n \in \mathbb{N} \to \mathrm{Seq}({(r)^{\bullet}})$}
\AXC{$\stackrel{\Pi_{120}}{\vdots}$}\noLine
\UIC{$\varphi_{231}$}
\BIC{$\varphi_{200}$}
\end{ptree}
\noindent$\Pi_{122}$:
\begin{ptree}
\AXC{$\mathrm{NSeq}(\mathsf{s}) \land \forall k \in \mathbb{N}\,(\mathsf{s}(k) = S(k))$\lbl{[S14]}}
\AXC{$\varphi_{103}$\lbl{A1}}
\BIC{$r \in \mathbb{R} \to \lvert r \rvert < 1 \to n \in \mathbb{N} \to \mathrm{NSeq}(\mathsf{s})$}
\AXC{$\stackrel{\Pi_{121}}{\vdots}$}\noLine
\UIC{$\varphi_{200}$}
\BIC{$r \in \mathbb{R} \to \lvert r \rvert < 1 \to n \in \mathbb{N} \to \varphi_{71}$}
\AXC{$(r \in \mathbb{R} \to \lvert r \rvert < 1 \to n \in \mathbb{N} \to \varphi_{71}) \to \varphi_{10}$\lbl{A1}}
\BIC{$\varphi_{10}$}
\end{ptree}
\noindent$\Pi_{123}$:
\begin{ptree}
\AXC{$\stackrel{\Pi_{89}}{\vdots}$}\noLine
\UIC{$r \in \mathbb{R} \to \lvert r \rvert < 1 \to n \in \mathbb{N} \to -1 \in \mathbb{R}$}
\AXC{$\varphi_{90}$\lbl{[S3]}}
\AXC{$\varphi_{90} \to \varphi_{127}$\lbl{A2}}
\BIC{$\varphi_{127}$}
\AXC{$\varphi_{127} \to -1 \in \mathbb{R} \to \varphi_{178}$\lbl{A2}}
\BIC{$-1 \in \mathbb{R} \to \varphi_{178}$}
\AXC{$(-1 \in \mathbb{R} \to \varphi_{178}) \to \varphi_{249}$\lbl{A1}}
\BIC{$\varphi_{249}$}
\BIC{$\varphi_{10} \to r \in \mathbb{R} \to \lvert r \rvert < 1 \to \varphi_{165}$}
\end{ptree}
\noindent$\Pi_{124}$:
\begin{ptree}
\AXC{$\stackrel{\Pi_{122}}{\vdots}$}\noLine
\UIC{$\varphi_{10}$}
\AXC{$\stackrel{\Pi_{123}}{\vdots}$}\noLine
\UIC{$\varphi_{10} \to r \in \mathbb{R} \to \lvert r \rvert < 1 \to \varphi_{165}$}
\BIC{$r \in \mathbb{R} \to \lvert r \rvert < 1 \to \varphi_{165}$}
\AXC{$(r \in \mathbb{R} \to \lvert r \rvert < 1 \to \varphi_{165}) \to \varphi_{38}$\lbl{A1}}
\BIC{$\varphi_{38}$}
\end{ptree}
\noindent$\Pi_{125}$:
\begin{ptree}
\AXC{$r \in \mathbb{R} \to \lvert r \rvert < 1 \to n \in \mathbb{N} \to r \in \mathbb{R}$\lbl{A1}}
\AXC{$\forall c \in \mathbb{R}\,(\mathrm{Seq}(\overline{c}) \land \forall n \in \mathbb{N}\,(\overline{c}(n) = c))$\lbl{[S10]}}
\AXC{$\varphi_{152}$\lbl{A2}}
\BIC{$r \in \mathbb{R} \to \mathrm{Seq}(\overline{r}) \land \forall n \in \mathbb{N}\,(\overline{r}(n) = r)$}
\AXC{$\varphi_{214}$\lbl{A1}}
\BIC{$(r \in \mathbb{R} \to \lvert r \rvert < 1 \to n \in \mathbb{N} \to r \in \mathbb{R}) \to r \in \mathbb{R} \to \varphi_{37}$}
\BIC{$r \in \mathbb{R} \to \varphi_{37}$}
\end{ptree}
\noindent$\Pi_{126}$:
\begin{ptree}
\AXC{$\varphi_{100}$\lbl{[S2]}}
\AXC{$\varphi_{100} \to \varphi_{141}$\lbl{A2}}
\BIC{$\varphi_{141}$}
\AXC{$\varphi_{141} \to \mathrm{Seq}(\overline{r}) \to \varphi_{221}$\lbl{A2}}
\BIC{$\mathrm{Seq}(\overline{r}) \to \varphi_{221}$}
\AXC{$(\mathrm{Seq}(\overline{r}) \to \varphi_{221}) \to \varphi_{262}$\lbl{A1}}
\BIC{$\varphi_{262}$}
\end{ptree}
\noindent$\Pi_{127}$:
\begin{ptree}
\AXC{$\stackrel{\Pi_{124}}{\vdots}$}\noLine
\UIC{$\varphi_{38}$}
\AXC{$\stackrel{\Pi_{125}}{\vdots}$}\noLine
\UIC{$r \in \mathbb{R} \to \varphi_{37}$}
\AXC{$\varphi_{173}$\lbl{A1}}
\BIC{$r \in \mathbb{R} \to \lvert r \rvert < 1 \to n \in \mathbb{N} \to \mathrm{Seq}(\overline{r})$}
\AXC{$\stackrel{\Pi_{126}}{\vdots}$}\noLine
\UIC{$\varphi_{262}$}
\BIC{$\varphi_{38} \to r \in \mathbb{R} \to \lvert r \rvert < 1 \to \varphi_{205}$}
\BIC{$r \in \mathbb{R} \to \lvert r \rvert < 1 \to \varphi_{205}$}
\AXC{$(r \in \mathbb{R} \to \lvert r \rvert < 1 \to \varphi_{205}) \to r \in \mathbb{R} \to \varphi_{45}$\lbl{A1}}
\BIC{$r \in \mathbb{R} \to \varphi_{45}$}
\end{ptree}
\noindent$\Pi_{128}$:
\begin{ptree}
\AXC{$\stackrel{\Pi_{46}}{\vdots}$}\noLine
\UIC{$r \in \mathbb{R} \to \lvert r \rvert < 1 \to n \in \mathbb{N} \to (1 - r)^{-1} \in \mathbb{R}$}
\AXC{$\varphi_{90}$\lbl{[S3]}}
\AXC{$\varphi_{90} \to \varphi_{157}$\lbl{A2}}
\BIC{$\varphi_{157}$}
\AXC{$\varphi_{157} \to \varphi_{256}$\lbl{A2}}
\BIC{$\varphi_{256}$}
\AXC{$\varphi_{256} \to \varphi_{272}$\lbl{A1}}
\BIC{$\varphi_{272}$}
\BIC{$(r \in \mathbb{R} \to \varphi_{45}) \to r \in \mathbb{R} \to \lvert r \rvert < 1 \to n \in \mathbb{N} \to \varphi_{225}$}
\end{ptree}
\noindent$\Pi_{129}$:
\begin{ptree}
\AXC{$\stackrel{\Pi_{127}}{\vdots}$}\noLine
\UIC{$r \in \mathbb{R} \to \varphi_{45}$}
\AXC{$\stackrel{\Pi_{128}}{\vdots}$}\noLine
\UIC{$(r \in \mathbb{R} \to \varphi_{45}) \to r \in \mathbb{R} \to \lvert r \rvert < 1 \to n \in \mathbb{N} \to \varphi_{225}$}
\BIC{$r \in \mathbb{R} \to \lvert r \rvert < 1 \to n \in \mathbb{N} \to \varphi_{225}$}
\AXC{$\varphi_{278}$\lbl{A1}}
\BIC{$r \in \mathbb{R} \to \lvert r \rvert < 1 \to n \in \mathbb{N} \to \forall n \in \mathbb{N}\,(\varphi_{167})$}
\end{ptree}
\noindent$\Pi_{130}$:
\begin{ptree}
\AXC{$\stackrel{\Pi_{129}}{\vdots}$}\noLine
\UIC{$r \in \mathbb{R} \to \lvert r \rvert < 1 \to n \in \mathbb{N} \to \forall n \in \mathbb{N}\,(\varphi_{167})$}
\AXC{$(r \in \mathbb{R} \to \lvert r \rvert < 1 \to n \in \mathbb{N} \to \forall n \in \mathbb{N}\,(\varphi_{167})) \to \varphi_{284}$\lbl{A1}}
\BIC{$\varphi_{284}$}
\end{ptree}
\noindent$\Pi_{131}$:
\begin{ptree}
\AXC{$\forall n \in \mathbb{N}\,(\varphi_{167}) \to n \in \mathbb{N} \to \varphi_{167}$\lbl{A2}}
\AXC{$\stackrel{\Pi_{130}}{\vdots}$}\noLine
\UIC{$\varphi_{284}$}
\BIC{$r \in \mathbb{R} \to \lvert r \rvert < 1 \to n \in \mathbb{N} \to n \in \mathbb{N} \to \varphi_{167}$}
\AXC{$(r \in \mathbb{R} \to \lvert r \rvert < 1 \to n \in \mathbb{N} \to n \in \mathbb{N} \to \varphi_{167}) \to \varphi_{230}$\lbl{A1}}
\BIC{$\varphi_{230}$}
\end{ptree}
\noindent$\Pi_{132}$:
\begin{ptree}
\AXC{$r \in \mathbb{R} \to \lvert r \rvert < 1 \to n \in \mathbb{N} \to n \in \mathbb{N}$\lbl{A1}}
\AXC{$\stackrel{\Pi_{131}}{\vdots}$}\noLine
\UIC{$\varphi_{230}$}
\BIC{$r \in \mathbb{R} \to \lvert r \rvert < 1 \to n \in \mathbb{N} \to \varphi_{167}$}
\AXC{$(r \in \mathbb{R} \to \lvert r \rvert < 1 \to n \in \mathbb{N} \to \varphi_{167}) \to \varphi_{298}$\lbl{A1}}
\BIC{$\varphi_{298}$}
\end{ptree}
\noindent$\Pi_{133}$:
\begin{ptree}
\AXC{$\stackrel{\Pi_{103}}{\vdots}$}\noLine
\UIC{$\varphi_{185}$}
\AXC{$\varphi_{185} \to \varphi_{250}$\lbl{A2}}
\BIC{$\varphi_{250}$}
\AXC{$\varphi_{250} \to \varphi_{167} \to \varphi_{168} \to \varphi_{174}$\lbl{A2}}
\BIC{$\varphi_{167} \to \varphi_{168} \to \varphi_{174}$}
\AXC{$\stackrel{\Pi_{117}}{\vdots}$}\noLine
\UIC{$r \in \mathbb{R} \to \lvert r \rvert < 1 \to n \in \mathbb{N} \to \varphi_{168}$}
\AXC{$\stackrel{\Pi_{132}}{\vdots}$}\noLine
\UIC{$\varphi_{298}$}
\BIC{$(\varphi_{167} \to \varphi_{168} \to \varphi_{174}) \to r \in \mathbb{R} \to \lvert r \rvert < 1 \to n \in \mathbb{N} \to \varphi_{174}$}
\BIC{$r \in \mathbb{R} \to \lvert r \rvert < 1 \to n \in \mathbb{N} \to \varphi_{174}$}
\end{ptree}
\noindent$\Pi_{134}$:
\begin{ptree}
\AXC{$\stackrel{\Pi_{102}}{\vdots}$}\noLine
\UIC{$r \in \mathbb{R} \to \lvert r \rvert < 1 \to n \in \mathbb{N} \to \varphi_{154}$}
\AXC{$\stackrel{\Pi_{133}}{\vdots}$}\noLine
\UIC{$r \in \mathbb{R} \to \lvert r \rvert < 1 \to n \in \mathbb{N} \to \varphi_{174}$}
\AXC{$(r \in \mathbb{R} \to \lvert r \rvert < 1 \to n \in \mathbb{N} \to \varphi_{174}) \to \varphi_{297}$\lbl{A1}}
\BIC{$\varphi_{297}$}
\BIC{$(\varphi_{174} \to \varphi_{154} \to \varphi_{153}) \to r \in \mathbb{R} \to \lvert r \rvert < 1 \to n \in \mathbb{N} \to \varphi_{153}$}
\end{ptree}
\noindent$\Pi_{135}$:
\begin{ptree}
\AXC{$\stackrel{\Pi_{95}}{\vdots}$}\noLine
\UIC{$\varphi_{185}$}
\AXC{$\varphi_{185} \to \varphi_{253}$\lbl{A2}}
\BIC{$\varphi_{253}$}
\AXC{$\varphi_{253} \to \varphi_{174} \to \varphi_{154} \to \varphi_{153}$\lbl{A2}}
\BIC{$\varphi_{174} \to \varphi_{154} \to \varphi_{153}$}
\AXC{$\stackrel{\Pi_{134}}{\vdots}$}\noLine
\UIC{$(\varphi_{174} \to \varphi_{154} \to \varphi_{153}) \to r \in \mathbb{R} \to \lvert r \rvert < 1 \to n \in \mathbb{N} \to \varphi_{153}$}
\BIC{$r \in \mathbb{R} \to \lvert r \rvert < 1 \to n \in \mathbb{N} \to \varphi_{153}$}
\end{ptree}
\noindent$\Pi_{136}$:
\begin{ptree}
\AXC{$\stackrel{\Pi_{94}}{\vdots}$}\noLine
\UIC{$r \in \mathbb{R} \to \lvert r \rvert < 1 \to n \in \mathbb{N} \to \varphi_{136}$}
\AXC{$\stackrel{\Pi_{135}}{\vdots}$}\noLine
\UIC{$r \in \mathbb{R} \to \lvert r \rvert < 1 \to n \in \mathbb{N} \to \varphi_{153}$}
\AXC{$(r \in \mathbb{R} \to \lvert r \rvert < 1 \to n \in \mathbb{N} \to \varphi_{153}) \to \varphi_{296}$\lbl{A1}}
\BIC{$\varphi_{296}$}
\BIC{$(\varphi_{153} \to \varphi_{136} \to \varphi_{151}) \to r \in \mathbb{R} \to \lvert r \rvert < 1 \to n \in \mathbb{N} \to \varphi_{151}$}
\end{ptree}
\noindent$\Pi_{137}$:
\begin{ptree}
\AXC{$\stackrel{\Pi_{83}}{\vdots}$}\noLine
\UIC{$\varphi_{185}$}
\AXC{$\varphi_{185} \to \varphi_{245}$\lbl{A2}}
\BIC{$\varphi_{245}$}
\AXC{$\varphi_{245} \to \varphi_{153} \to \varphi_{136} \to \varphi_{151}$\lbl{A2}}
\BIC{$\varphi_{153} \to \varphi_{136} \to \varphi_{151}$}
\AXC{$\stackrel{\Pi_{136}}{\vdots}$}\noLine
\UIC{$(\varphi_{153} \to \varphi_{136} \to \varphi_{151}) \to r \in \mathbb{R} \to \lvert r \rvert < 1 \to n \in \mathbb{N} \to \varphi_{151}$}
\BIC{$r \in \mathbb{R} \to \lvert r \rvert < 1 \to n \in \mathbb{N} \to \varphi_{151}$}
\end{ptree}
\noindent$\Pi_{138}$:
\begin{ptree}
\AXC{$\stackrel{\Pi_{137}}{\vdots}$}\noLine
\UIC{$r \in \mathbb{R} \to \lvert r \rvert < 1 \to n \in \mathbb{N} \to \varphi_{151}$}
\AXC{$(r \in \mathbb{R} \to \lvert r \rvert < 1 \to n \in \mathbb{N} \to \varphi_{151}) \to \varphi_{292}$\lbl{A1}}
\BIC{$\varphi_{292}$}
\end{ptree}
\noindent$\Pi_{139}$:
\begin{ptree}
\AXC{$z = w \to \varphi_{70} \to \varphi_{69}$\lbl{A6}}
\UIC{$\forall w\,(z = w \to \varphi_{70} \to \varphi_{69})$}
\UIC{$\forall z\,\forall w\,(z = w \to \varphi_{70} \to \varphi_{69})$}
\AXC{$\varphi_{273}$\lbl{A2}}
\BIC{$\forall w\,({(r)^{\bullet}} \circ \mathsf{s}(n) = w \to \varphi_{126} \to \varphi_{69})$}
\AXC{$\varphi_{283}$\lbl{A2}}
\BIC{${(r)^{\bullet}} \circ \mathsf{s}(n) = {{r}^{S(n)}} \to \varphi_{126} \to \varphi_{112}$}
\AXC{$\stackrel{\Pi_{82}}{\vdots}$}\noLine
\UIC{$\varphi_{269}$}
\BIC{$r \in \mathbb{R} \to \lvert r \rvert < 1 \to n \in \mathbb{N} \to \varphi_{112}$}
\AXC{$\stackrel{\Pi_{138}}{\vdots}$}\noLine
\UIC{$\varphi_{292}$}
\BIC{$(\varphi_{151} \to \varphi_{112} \to \varphi_{140}) \to r \in \mathbb{R} \to \lvert r \rvert < 1 \to n \in \mathbb{N} \to \varphi_{140}$}
\end{ptree}
\noindent$\Pi_{140}$:
\begin{ptree}
\AXC{$\stackrel{\Pi_{70}}{\vdots}$}\noLine
\UIC{$\varphi_{185}$}
\AXC{$\varphi_{185} \to \varphi_{243}$\lbl{A2}}
\BIC{$\varphi_{243}$}
\AXC{$\varphi_{243} \to \varphi_{151} \to \varphi_{112} \to \varphi_{140}$\lbl{A2}}
\BIC{$\varphi_{151} \to \varphi_{112} \to \varphi_{140}$}
\AXC{$\stackrel{\Pi_{139}}{\vdots}$}\noLine
\UIC{$(\varphi_{151} \to \varphi_{112} \to \varphi_{140}) \to r \in \mathbb{R} \to \lvert r \rvert < 1 \to n \in \mathbb{N} \to \varphi_{140}$}
\BIC{$r \in \mathbb{R} \to \lvert r \rvert < 1 \to n \in \mathbb{N} \to \varphi_{140}$}
\end{ptree}
\noindent$\Pi_{141}$:
\begin{ptree}
\AXC{$\stackrel{\Pi_{69}}{\vdots}$}\noLine
\UIC{$r \in \mathbb{R} \to \lvert r \rvert < 1 \to n \in \mathbb{N} \to \varphi_{76}$}
\AXC{$\stackrel{\Pi_{140}}{\vdots}$}\noLine
\UIC{$r \in \mathbb{R} \to \lvert r \rvert < 1 \to n \in \mathbb{N} \to \varphi_{140}$}
\AXC{$(r \in \mathbb{R} \to \lvert r \rvert < 1 \to n \in \mathbb{N} \to \varphi_{140}) \to \varphi_{290}$\lbl{A1}}
\BIC{$\varphi_{290}$}
\BIC{$(\varphi_{140} \to \varphi_{76} \to \varphi_{135}) \to r \in \mathbb{R} \to \lvert r \rvert < 1 \to n \in \mathbb{N} \to \varphi_{135}$}
\end{ptree}
\noindent$\Pi_{142}$:
\begin{ptree}
\AXC{$\stackrel{\Pi_{62}}{\vdots}$}\noLine
\UIC{$\varphi_{185}$}
\AXC{$\varphi_{185} \to \varphi_{240}$\lbl{A2}}
\BIC{$\varphi_{240}$}
\AXC{$\varphi_{240} \to \varphi_{140} \to \varphi_{76} \to \varphi_{135}$\lbl{A2}}
\BIC{$\varphi_{140} \to \varphi_{76} \to \varphi_{135}$}
\AXC{$\stackrel{\Pi_{141}}{\vdots}$}\noLine
\UIC{$(\varphi_{140} \to \varphi_{76} \to \varphi_{135}) \to r \in \mathbb{R} \to \lvert r \rvert < 1 \to n \in \mathbb{N} \to \varphi_{135}$}
\BIC{$r \in \mathbb{R} \to \lvert r \rvert < 1 \to n \in \mathbb{N} \to \varphi_{135}$}
\end{ptree}
\noindent$\Pi_{143}$:
\begin{ptree}
\AXC{$\stackrel{\Pi_{61}}{\vdots}$}\noLine
\UIC{$r \in \mathbb{R} \to \lvert r \rvert < 1 \to n \in \mathbb{N} \to \varphi_{46}$}
\AXC{$\stackrel{\Pi_{142}}{\vdots}$}\noLine
\UIC{$r \in \mathbb{R} \to \lvert r \rvert < 1 \to n \in \mathbb{N} \to \varphi_{135}$}
\AXC{$(r \in \mathbb{R} \to \lvert r \rvert < 1 \to n \in \mathbb{N} \to \varphi_{135}) \to \varphi_{288}$\lbl{A1}}
\BIC{$\varphi_{288}$}
\BIC{$(\varphi_{135} \to \varphi_{46} \to \varphi_{121}) \to r \in \mathbb{R} \to \lvert r \rvert < 1 \to n \in \mathbb{N} \to \varphi_{121}$}
\end{ptree}
\noindent$\Pi_{144}$:
\begin{ptree}
\AXC{$\stackrel{\Pi_{56}}{\vdots}$}\noLine
\UIC{$\varphi_{185}$}
\AXC{$\varphi_{185} \to \varphi_{233}$\lbl{A2}}
\BIC{$\varphi_{233}$}
\AXC{$\varphi_{233} \to \varphi_{135} \to \varphi_{46} \to \varphi_{121}$\lbl{A2}}
\BIC{$\varphi_{135} \to \varphi_{46} \to \varphi_{121}$}
\AXC{$\stackrel{\Pi_{143}}{\vdots}$}\noLine
\UIC{$(\varphi_{135} \to \varphi_{46} \to \varphi_{121}) \to r \in \mathbb{R} \to \lvert r \rvert < 1 \to n \in \mathbb{N} \to \varphi_{121}$}
\BIC{$r \in \mathbb{R} \to \lvert r \rvert < 1 \to n \in \mathbb{N} \to \varphi_{121}$}
\end{ptree}
\noindent$\Pi_{145}$:
\begin{ptree}
\AXC{$\stackrel{\Pi_{55}}{\vdots}$}\noLine
\UIC{$r \in \mathbb{R} \to \lvert r \rvert < 1 \to n \in \mathbb{N} \to \varphi_{52}$}
\AXC{$\stackrel{\Pi_{144}}{\vdots}$}\noLine
\UIC{$r \in \mathbb{R} \to \lvert r \rvert < 1 \to n \in \mathbb{N} \to \varphi_{121}$}
\AXC{$(r \in \mathbb{R} \to \lvert r \rvert < 1 \to n \in \mathbb{N} \to \varphi_{121}) \to \varphi_{289}$\lbl{A1}}
\BIC{$\varphi_{289}$}
\BIC{$(\varphi_{121} \to \varphi_{52} \to \varphi_{139}) \to r \in \mathbb{R} \to \lvert r \rvert < 1 \to n \in \mathbb{N} \to \varphi_{139}$}
\end{ptree}
\noindent$\Pi_{146}$:
\begin{ptree}
\AXC{$\stackrel{\Pi_{50}}{\vdots}$}\noLine
\UIC{$\varphi_{185}$}
\AXC{$\varphi_{185} \to \varphi_{226}$\lbl{A2}}
\BIC{$\varphi_{226}$}
\AXC{$\varphi_{226} \to \varphi_{121} \to \varphi_{52} \to \varphi_{139}$\lbl{A2}}
\BIC{$\varphi_{121} \to \varphi_{52} \to \varphi_{139}$}
\AXC{$\stackrel{\Pi_{145}}{\vdots}$}\noLine
\UIC{$(\varphi_{121} \to \varphi_{52} \to \varphi_{139}) \to r \in \mathbb{R} \to \lvert r \rvert < 1 \to n \in \mathbb{N} \to \varphi_{139}$}
\BIC{$r \in \mathbb{R} \to \lvert r \rvert < 1 \to n \in \mathbb{N} \to \varphi_{139}$}
\end{ptree}
\noindent$\Pi_{147}$:
\begin{ptree}
\AXC{$\stackrel{\Pi_{49}}{\vdots}$}\noLine
\UIC{$r \in \mathbb{R} \to \lvert r \rvert < 1 \to n \in \mathbb{N} \to \varphi_{78}$}
\AXC{$\stackrel{\Pi_{146}}{\vdots}$}\noLine
\UIC{$r \in \mathbb{R} \to \lvert r \rvert < 1 \to n \in \mathbb{N} \to \varphi_{139}$}
\AXC{$(r \in \mathbb{R} \to \lvert r \rvert < 1 \to n \in \mathbb{N} \to \varphi_{139}) \to \varphi_{291}$\lbl{A1}}
\BIC{$\varphi_{291}$}
\BIC{$(\varphi_{139} \to \varphi_{78} \to \varphi_{138}) \to r \in \mathbb{R} \to \lvert r \rvert < 1 \to n \in \mathbb{N} \to \varphi_{138}$}
\end{ptree}
\noindent$\Pi_{148}$:
\begin{ptree}
\AXC{$\stackrel{\Pi_{43}}{\vdots}$}\noLine
\UIC{$\varphi_{185}$}
\AXC{$\varphi_{185} \to \varphi_{238}$\lbl{A2}}
\BIC{$\varphi_{238}$}
\AXC{$\varphi_{238} \to \varphi_{139} \to \varphi_{78} \to \varphi_{138}$\lbl{A2}}
\BIC{$\varphi_{139} \to \varphi_{78} \to \varphi_{138}$}
\AXC{$\stackrel{\Pi_{147}}{\vdots}$}\noLine
\UIC{$(\varphi_{139} \to \varphi_{78} \to \varphi_{138}) \to r \in \mathbb{R} \to \lvert r \rvert < 1 \to n \in \mathbb{N} \to \varphi_{138}$}
\BIC{$r \in \mathbb{R} \to \lvert r \rvert < 1 \to n \in \mathbb{N} \to \varphi_{138}$}
\end{ptree}
\noindent$\Pi_{149}$:
\begin{ptree}
\AXC{$\stackrel{\Pi_{148}}{\vdots}$}\noLine
\UIC{$r \in \mathbb{R} \to \lvert r \rvert < 1 \to n \in \mathbb{N} \to \varphi_{138}$}
\AXC{$(r \in \mathbb{R} \to \lvert r \rvert < 1 \to n \in \mathbb{N} \to \varphi_{138}) \to (r \in \mathbb{R} \to \varphi_{98}) \to \varphi_{274}$\lbl{A1}}
\BIC{$(r \in \mathbb{R} \to \varphi_{98}) \to \varphi_{274}$}
\end{ptree}
\noindent$\Pi_{150}$:
\begin{ptree}
\AXC{$\stackrel{\Pi_{38}}{\vdots}$}\noLine
\UIC{$\varphi_{202}$}
\AXC{$x = x$\lbl{A5}}
\UIC{$\forall x\,(x = x)$}
\AXC{$\forall x\,(x = x) \to \varphi_{77}$\lbl{A2}}
\BIC{$\varphi_{77}$}
\AXC{$\stackrel{\Pi_{42}}{\vdots}$}\noLine
\UIC{$\varphi_{77} \to \varphi_{202} \to r \in \mathbb{R} \to \varphi_{98}$}
\BIC{$\varphi_{202} \to r \in \mathbb{R} \to \varphi_{98}$}
\BIC{$r \in \mathbb{R} \to \varphi_{98}$}
\AXC{$\stackrel{\Pi_{149}}{\vdots}$}\noLine
\UIC{$(r \in \mathbb{R} \to \varphi_{98}) \to \varphi_{274}$}
\BIC{$\varphi_{274}$}
\end{ptree}
\noindent$\Pi_{151}$:
\begin{ptree}
\AXC{$\stackrel{\Pi_{37}}{\vdots}$}\noLine
\UIC{$\varphi_{185}$}
\AXC{$\varphi_{185} \to \varphi_{237}$\lbl{A2}}
\BIC{$\varphi_{237}$}
\AXC{$\varphi_{285}$\lbl{A2}}
\BIC{$\varphi_{138} \to (1 - r)^{-1} \cdot (1 - r) \cdot {\textstyle\sum_{k=1}^{n} {{r}^{k}}} = 1 \cdot {\textstyle\sum_{k=1}^{n} {{r}^{k}}} \to \varphi_{97}$}
\AXC{$\stackrel{\Pi_{150}}{\vdots}$}\noLine
\UIC{$\varphi_{274}$}
\BIC{$r \in \mathbb{R} \to \lvert r \rvert < 1 \to n \in \mathbb{N} \to \varphi_{97}$}
\AXC{$(r \in \mathbb{R} \to \lvert r \rvert < 1 \to n \in \mathbb{N} \to \varphi_{97}) \to \varphi_{281}$\lbl{A1}}
\BIC{$\varphi_{281}$}
\end{ptree}
\noindent$\Pi_{152}$:
\begin{ptree}
\AXC{$\stackrel{\Pi_{35}}{\vdots}$}\noLine
\UIC{$\varphi_{97} \to 1 \cdot {\textstyle\sum_{k=1}^{n} {{r}^{k}}} = {\textstyle\sum_{k=1}^{n} {{r}^{k}}} \cdot 1 \to \varphi_{96}$}
\AXC{$\stackrel{\Pi_{32}}{\vdots}$}\noLine
\UIC{$r \in \mathbb{R} \to \lvert r \rvert < 1 \to n \in \mathbb{N} \to {\textstyle\sum_{k=1}^{n} {{r}^{k}}} \in \mathbb{R}$}
\AXC{$\stackrel{\Pi_{36}}{\vdots}$}\noLine
\UIC{$\varphi_{184}$}
\BIC{$r \in \mathbb{R} \to \lvert r \rvert < 1 \to n \in \mathbb{N} \to 1 \cdot {\textstyle\sum_{k=1}^{n} {{r}^{k}}} = {\textstyle\sum_{k=1}^{n} {{r}^{k}}} \cdot 1$}
\AXC{$\stackrel{\Pi_{151}}{\vdots}$}\noLine
\UIC{$\varphi_{281}$}
\BIC{$\varphi_{267}$}
\BIC{$r \in \mathbb{R} \to \lvert r \rvert < 1 \to n \in \mathbb{N} \to \varphi_{96}$}
\end{ptree}
\noindent$\Pi_{153}$:
\begin{ptree}
\AXC{$\stackrel{\Pi_{152}}{\vdots}$}\noLine
\UIC{$r \in \mathbb{R} \to \lvert r \rvert < 1 \to n \in \mathbb{N} \to \varphi_{96}$}
\AXC{$(r \in \mathbb{R} \to \lvert r \rvert < 1 \to n \in \mathbb{N} \to \varphi_{96}) \to \varphi_{277}$\lbl{A1}}
\BIC{$\varphi_{277}$}
\end{ptree}
\noindent$\Pi_{154}$:
\begin{ptree}
\AXC{$\stackrel{\Pi_{32}}{\vdots}$}\noLine
\UIC{$r \in \mathbb{R} \to \lvert r \rvert < 1 \to n \in \mathbb{N} \to {\textstyle\sum_{k=1}^{n} {{r}^{k}}} \in \mathbb{R}$}
\AXC{$\stackrel{\Pi_{33}}{\vdots}$}\noLine
\UIC{$\varphi_{181}$}
\BIC{$r \in \mathbb{R} \to \lvert r \rvert < 1 \to n \in \mathbb{N} \to {\textstyle\sum_{k=1}^{n} {{r}^{k}}} \cdot 1 = {\textstyle\sum_{k=1}^{n} {{r}^{k}}}$}
\AXC{$\stackrel{\Pi_{153}}{\vdots}$}\noLine
\UIC{$\varphi_{277}$}
\BIC{$\varphi_{266}$}
\end{ptree}
\noindent$\Pi_{155}$:
\begin{ptree}
\AXC{$\stackrel{\Pi_{26}}{\vdots}$}\noLine
\UIC{$\varphi_{185}$}
\AXC{$\varphi_{185} \to \varphi_{212}$\lbl{A2}}
\BIC{$\varphi_{212}$}
\AXC{$\varphi_{212} \to \varphi_{96} \to {\textstyle\sum_{k=1}^{n} {{r}^{k}}} \cdot 1 = {\textstyle\sum_{k=1}^{n} {{r}^{k}}} \to \varphi_{85}$\lbl{A2}}
\BIC{$\varphi_{96} \to {\textstyle\sum_{k=1}^{n} {{r}^{k}}} \cdot 1 = {\textstyle\sum_{k=1}^{n} {{r}^{k}}} \to \varphi_{85}$}
\AXC{$\stackrel{\Pi_{154}}{\vdots}$}\noLine
\UIC{$\varphi_{266}$}
\BIC{$r \in \mathbb{R} \to \lvert r \rvert < 1 \to n \in \mathbb{N} \to \varphi_{85}$}
\AXC{$\varphi_{242}$\lbl{A1}}
\BIC{$r \in \mathbb{R} \land \lvert r \rvert < 1 \to n \in \mathbb{N} \to \varphi_{85}$}
\UIC{$\forall n\,(r \in \mathbb{R} \land \lvert r \rvert < 1 \to n \in \mathbb{N} \to \varphi_{85})$}
\end{ptree}
\noindent$\Pi_{156}$:
\begin{ptree}
\AXC{$\varphi_{246}$\lbl{A3}}
\AXC{$\stackrel{\Pi_{155}}{\vdots}$}\noLine
\UIC{$\forall n\,(r \in \mathbb{R} \land \lvert r \rvert < 1 \to n \in \mathbb{N} \to \varphi_{85})$}
\AXC{$\forall n\,(r \in \mathbb{R} \land \lvert r \rvert < 1 \to n \in \mathbb{N} \to \varphi_{85}) \to \varphi_{246} \to \varphi_{207}$\lbl{A1}}
\BIC{$\varphi_{246} \to \varphi_{207}$}
\BIC{$\varphi_{207}$}
\end{ptree}
\noindent$\Pi_{157}$:
\begin{ptree}
\AXC{$r \in \mathbb{R} \to \lvert r \rvert < 1 \to r \in \mathbb{R}$\lbl{A1}}
\AXC{$\forall r \in \mathbb{R}\,(\varphi_{58})$\lbl{[S24]}}
\AXC{$\forall r \in \mathbb{R}\,(\varphi_{58}) \to r \in \mathbb{R} \to \varphi_{58}$\lbl{A2}}
\BIC{$r \in \mathbb{R} \to \varphi_{58}$}
\AXC{$\varphi_{241}$\lbl{A1}}
\BIC{$(r \in \mathbb{R} \to \lvert r \rvert < 1 \to r \in \mathbb{R}) \to r \in \mathbb{R} \to \lvert r \rvert < 1 \to \varphi_{58}$}
\BIC{$r \in \mathbb{R} \to \lvert r \rvert < 1 \to \varphi_{58}$}
\AXC{$\varphi_{180}$\lbl{A1}}
\BIC{$r \in \mathbb{R} \to \lvert r \rvert < 1 \to \mathrm{Seq}({(r)^{\bullet}})$}
\end{ptree}
\noindent$\Pi_{158}$:
\begin{ptree}
\AXC{$\varphi_{149}$\lbl{[S25]}}
\AXC{$\varphi_{149} \to \mathrm{Seq}({(r)^{\bullet}}) \to \varphi_{164}$\lbl{A2}}
\BIC{$\mathrm{Seq}({(r)^{\bullet}}) \to \varphi_{164}$}
\AXC{$(\mathrm{Seq}({(r)^{\bullet}}) \to \varphi_{164}) \to \varphi_{215}$\lbl{A1}}
\BIC{$\varphi_{215}$}
\end{ptree}
\noindent$\Pi_{159}$:
\begin{ptree}
\AXC{$\varphi_{48}$\lbl{[C13]}}
\AXC{$\varphi_{48} \to \varphi_{211}$\lbl{A2}}
\BIC{$\varphi_{211}$}
\AXC{$\varphi_{211} \to \varphi_{227}$\lbl{A2}}
\BIC{$\varphi_{227}$}
\AXC{$\varphi_{227} \to \mathrm{Seq}(\Sigma {(r)^{\bullet}}) \to \forall n \in \mathbb{N}\,(\varphi_{85}) \to \varphi_{114} \to \varphi_{3}$\lbl{A2}}
\BIC{$\mathrm{Seq}(\Sigma {(r)^{\bullet}}) \to \forall n \in \mathbb{N}\,(\varphi_{85}) \to \varphi_{114} \to \varphi_{3}$}
\AXC{$\varphi_{294}$\lbl{A1}}
\BIC{$(r \in \mathbb{R} \to \lvert r \rvert < 1 \to \mathrm{Seq}(\Sigma {(r)^{\bullet}})) \to \varphi_{261}$}
\end{ptree}
\noindent$\Pi_{160}$:
\begin{ptree}
\AXC{$\stackrel{\Pi_{157}}{\vdots}$}\noLine
\UIC{$r \in \mathbb{R} \to \lvert r \rvert < 1 \to \mathrm{Seq}({(r)^{\bullet}})$}
\AXC{$\stackrel{\Pi_{158}}{\vdots}$}\noLine
\UIC{$\varphi_{215}$}
\BIC{$r \in \mathbb{R} \to \lvert r \rvert < 1 \to \varphi_{164}$}
\AXC{$\varphi_{217}$\lbl{A1}}
\BIC{$r \in \mathbb{R} \to \lvert r \rvert < 1 \to \mathrm{Seq}(\Sigma {(r)^{\bullet}})$}
\AXC{$\stackrel{\Pi_{159}}{\vdots}$}\noLine
\UIC{$(r \in \mathbb{R} \to \lvert r \rvert < 1 \to \mathrm{Seq}(\Sigma {(r)^{\bullet}})) \to \varphi_{261}$}
\BIC{$\varphi_{261}$}
\end{ptree}
\noindent$\Pi_{161}$:
\begin{ptree}
\AXC{$\stackrel{\Pi_{25}}{\vdots}$}\noLine
\UIC{$r \in \mathbb{R} \to \lvert r \rvert < 1 \to \varphi_{114}$}
\AXC{$r \in \mathbb{R} \land \lvert r \rvert < 1 \to \forall n\,(r \in \mathbb{R} \land \lvert r \rvert < 1)$\lbl{A4}}
\AXC{$\stackrel{\Pi_{156}}{\vdots}$}\noLine
\UIC{$\varphi_{207}$}
\BIC{$r \in \mathbb{R} \to \lvert r \rvert < 1 \to \forall n \in \mathbb{N}\,(\varphi_{85})$}
\AXC{$\stackrel{\Pi_{160}}{\vdots}$}\noLine
\UIC{$\varphi_{261}$}
\BIC{$(r \in \mathbb{R} \to \lvert r \rvert < 1 \to \varphi_{114}) \to r \in \mathbb{R} \to \lvert r \rvert < 1 \to \varphi_{3}$}
\BIC{$r \in \mathbb{R} \to \lvert r \rvert < 1 \to \varphi_{3}$}
\end{ptree}
\noindent$\Pi_{162}$:
\begin{ptree}
\AXC{$y = z \to x = y \to x = z$\lbl{A6}}
\AXC{$(y = z \to x = y \to x = z) \to x = y \to y = z \to x = z$\lbl{A1}}
\BIC{$x = y \to y = z \to x = z$}
\UIC{$\forall z\,(x = y \to y = z \to x = z)$}
\UIC{$\forall y\,\forall z\,(x = y \to y = z \to x = z)$}
\UIC{$\forall x\,\forall y\,\forall z\,(x = y \to y = z \to x = z)$}
\AXC{$\forall x\,\forall y\,\forall z\,(x = y \to y = z \to x = z) \to \varphi_{56}$\lbl{A2}}
\BIC{$\varphi_{56}$}
\end{ptree}
\noindent$\Pi_{163}$:
\begin{ptree}
\AXC{$\forall x\,\forall y\,(x \in \mathbb{R} \to y \in \mathbb{R} \to x \cdot y = y \cdot x)$\lbl{(R7)}}
\AXC{$\forall x\,\forall y\,(x \in \mathbb{R} \to y \in \mathbb{R} \to x \cdot y = y \cdot x) \to \varphi_{36}$\lbl{A2}}
\BIC{$\varphi_{36}$}
\AXC{$\varphi_{36} \to \varphi_{30}$\lbl{A2}}
\BIC{$\varphi_{30}$}
\end{ptree}
\noindent$\Pi_{164}$:
\begin{ptree}
\AXC{$\stackrel{\Pi_{23}}{\vdots}$}\noLine
\UIC{$r \in \mathbb{R} \to \lvert r \rvert < 1 \to (1 - r)^{-1} \in \mathbb{R}$}
\AXC{$\stackrel{\Pi_{163}}{\vdots}$}\noLine
\UIC{$\varphi_{30}$}
\AXC{$\varphi_{30} \to (r \in \mathbb{R} \to \lvert r \rvert < 1 \to (1 - r)^{-1} \in \mathbb{R}) \to \varphi_{113}$\lbl{A1}}
\BIC{$(r \in \mathbb{R} \to \lvert r \rvert < 1 \to (1 - r)^{-1} \in \mathbb{R}) \to \varphi_{113}$}
\BIC{$\varphi_{113}$}
\end{ptree}
\noindent$\Pi_{165}$:
\begin{ptree}
\AXC{$y = z \to x = y \to x = z$\lbl{A6}}
\AXC{$(y = z \to x = y \to x = z) \to x = y \to y = z \to x = z$\lbl{A1}}
\BIC{$x = y \to y = z \to x = z$}
\UIC{$\forall z\,(x = y \to y = z \to x = z)$}
\UIC{$\forall y\,\forall z\,(x = y \to y = z \to x = z)$}
\UIC{$\forall x\,\forall y\,\forall z\,(x = y \to y = z \to x = z)$}
\AXC{$\forall x\,\forall y\,\forall z\,(x = y \to y = z \to x = z) \to \varphi_{56}$\lbl{A2}}
\BIC{$\varphi_{56}$}
\end{ptree}
\noindent$\Pi_{166}$:
\begin{ptree}
\AXC{$z = w \to \varphi_{53}$\lbl{A6}}
\UIC{$\forall w\,(z = w \to \varphi_{53})$}
\UIC{$\forall z\,\forall w\,(z = w \to \varphi_{53})$}
\AXC{$\forall z\,\forall w\,(z = w \to \varphi_{53}) \to \varphi_{110}$\lbl{A2}}
\BIC{$\varphi_{110}$}
\AXC{$\varphi_{110} \to r + 0 = r \to \varphi_{66}$\lbl{A2}}
\BIC{$r + 0 = r \to \varphi_{66}$}
\end{ptree}
\noindent$\Pi_{167}$:
\begin{ptree}
\AXC{$\forall x \in \mathbb{R}\,(x + 0 = x)$\lbl{(R4)}}
\AXC{$\forall x \in \mathbb{R}\,(x + 0 = x) \to r \in \mathbb{R} \to r + 0 = r$\lbl{A2}}
\BIC{$r \in \mathbb{R} \to r + 0 = r$}
\AXC{$(r \in \mathbb{R} \to r + 0 = r) \to \varphi_{44}$\lbl{A1}}
\BIC{$\varphi_{44}$}
\end{ptree}
\noindent$\Pi_{168}$:
\begin{ptree}
\AXC{$r \in \mathbb{R} \to \lvert r \rvert < 1 \to r \in \mathbb{R}$\lbl{A1}}
\AXC{$\stackrel{\Pi_{167}}{\vdots}$}\noLine
\UIC{$\varphi_{44}$}
\BIC{$r \in \mathbb{R} \to \lvert r \rvert < 1 \to r + 0 = r$}
\AXC{$(r \in \mathbb{R} \to \lvert r \rvert < 1 \to r + 0 = r) \to \varphi_{219}$\lbl{A1}}
\BIC{$\varphi_{219}$}
\end{ptree}
\noindent$\Pi_{169}$:
\begin{ptree}
\AXC{$\stackrel{\Pi_{166}}{\vdots}$}\noLine
\UIC{$r + 0 = r \to \varphi_{66}$}
\AXC{$x = x$\lbl{A5}}
\UIC{$\forall x\,(x = x)$}
\AXC{$\varphi_{4}$\lbl{A2}}
\BIC{$(1 - r)^{-1} \cdot (r + 0) = (1 - r)^{-1} \cdot (r + 0)$}
\AXC{$\stackrel{\Pi_{168}}{\vdots}$}\noLine
\UIC{$\varphi_{219}$}
\BIC{$(r + 0 = r \to \varphi_{66}) \to \varphi_{22}$}
\BIC{$\varphi_{22}$}
\end{ptree}
\noindent$\Pi_{170}$:
\begin{ptree}
\AXC{$\forall x \in \mathbb{R}\,(x \cdot 0 = 0)$\lbl{[F8]}}
\AXC{$\forall x \in \mathbb{R}\,(x \cdot 0 = 0) \to -1 \in \mathbb{R} \to -1 \cdot 0 = 0$\lbl{A2}}
\BIC{$-1 \in \mathbb{R} \to -1 \cdot 0 = 0$}
\AXC{$(-1 \in \mathbb{R} \to -1 \cdot 0 = 0) \to \varphi_{68}$\lbl{A1}}
\BIC{$\varphi_{68}$}
\end{ptree}
\noindent$\Pi_{171}$:
\begin{ptree}
\AXC{$x = x$\lbl{A5}}
\UIC{$\forall x\,(x = x)$}
\AXC{$\forall x\,(x = x) \to \varphi_{18}$\lbl{A2}}
\BIC{$\varphi_{18}$}
\AXC{$\stackrel{\Pi_{2}}{\vdots}$}\noLine
\UIC{$r \in \mathbb{R} \to \lvert r \rvert < 1 \to -1 \in \mathbb{R}$}
\AXC{$\stackrel{\Pi_{170}}{\vdots}$}\noLine
\UIC{$\varphi_{68}$}
\BIC{$r \in \mathbb{R} \to \lvert r \rvert < 1 \to -1 \cdot 0 = 0$}
\AXC{$\varphi_{260}$\lbl{A1}}
\BIC{$\varphi_{18} \to (-1 \cdot 0 = 0 \to \varphi_{147}) \to r \in \mathbb{R} \to \varphi_{26}$}
\BIC{$(-1 \cdot 0 = 0 \to \varphi_{147}) \to r \in \mathbb{R} \to \varphi_{26}$}
\end{ptree}
\noindent$\Pi_{172}$:
\begin{ptree}
\AXC{$z = w \to \varphi_{131}$\lbl{A6}}
\UIC{$\forall w\,(z = w \to \varphi_{131})$}
\UIC{$\forall z\,\forall w\,(z = w \to \varphi_{131})$}
\AXC{$\forall z\,\forall w\,(z = w \to \varphi_{131}) \to \varphi_{171}$\lbl{A2}}
\BIC{$\varphi_{171}$}
\AXC{$\varphi_{171} \to -1 \cdot 0 = 0 \to \varphi_{147}$\lbl{A2}}
\BIC{$-1 \cdot 0 = 0 \to \varphi_{147}$}
\AXC{$\stackrel{\Pi_{171}}{\vdots}$}\noLine
\UIC{$(-1 \cdot 0 = 0 \to \varphi_{147}) \to r \in \mathbb{R} \to \varphi_{26}$}
\BIC{$r \in \mathbb{R} \to \varphi_{26}$}
\end{ptree}
\noindent$\Pi_{173}$:
\begin{ptree}
\AXC{$\stackrel{\Pi_{165}}{\vdots}$}\noLine
\UIC{$\varphi_{56}$}
\AXC{$\varphi_{56} \to \varphi_{146}$\lbl{A2}}
\BIC{$\varphi_{146}$}
\AXC{$\varphi_{146} \to \varphi_{176}$\lbl{A2}}
\BIC{$\varphi_{176}$}
\AXC{$\stackrel{\Pi_{169}}{\vdots}$}\noLine
\UIC{$\varphi_{22}$}
\AXC{$\stackrel{\Pi_{172}}{\vdots}$}\noLine
\UIC{$r \in \mathbb{R} \to \varphi_{26}$}
\AXC{$(r \in \mathbb{R} \to \varphi_{26}) \to \varphi_{22} \to \varphi_{176} \to r \in \mathbb{R} \to \varphi_{16}$\lbl{A1}}
\BIC{$\varphi_{22} \to \varphi_{176} \to r \in \mathbb{R} \to \varphi_{16}$}
\BIC{$\varphi_{176} \to r \in \mathbb{R} \to \varphi_{16}$}
\BIC{$r \in \mathbb{R} \to \varphi_{16}$}
\end{ptree}
\noindent$\Pi_{174}$:
\begin{ptree}
\AXC{$r \in \mathbb{R} \to \lvert r \rvert < 1 \to r \in \mathbb{R}$\lbl{A1}}
\AXC{$\stackrel{\Pi_{164}}{\vdots}$}\noLine
\UIC{$\varphi_{113}$}
\BIC{$r \in \mathbb{R} \to \lvert r \rvert < 1 \to (1 - r)^{-1} \cdot r = r \cdot (1 - r)^{-1}$}
\AXC{$\stackrel{\Pi_{173}}{\vdots}$}\noLine
\UIC{$r \in \mathbb{R} \to \varphi_{16}$}
\AXC{$(r \in \mathbb{R} \to \varphi_{16}) \to \varphi_{248}$\lbl{A1}}
\BIC{$\varphi_{248}$}
\BIC{$\varphi_{222}$}
\end{ptree}
\noindent$\Pi_{175}$:
\begin{ptree}
\AXC{$\stackrel{\Pi_{162}}{\vdots}$}\noLine
\UIC{$\varphi_{56}$}
\AXC{$\varphi_{56} \to \varphi_{129}$\lbl{A2}}
\BIC{$\varphi_{129}$}
\AXC{$\varphi_{129} \to (1 - r)^{-1} \cdot (r + -1 \cdot 0) = (1 - r)^{-1} \cdot r \to \varphi_{65}$\lbl{A2}}
\BIC{$(1 - r)^{-1} \cdot (r + -1 \cdot 0) = (1 - r)^{-1} \cdot r \to \varphi_{65}$}
\AXC{$\stackrel{\Pi_{174}}{\vdots}$}\noLine
\UIC{$\varphi_{222}$}
\BIC{$r \in \mathbb{R} \to \varphi_{17}$}
\AXC{$\varphi_{265}$\lbl{A1}}
\BIC{$(r \in \mathbb{R} \to \lvert r \rvert < 1 \to \varphi_{3}) \to \varphi_{169} \to r \in \mathbb{R} \to \varphi_{2}$}
\end{ptree}
\noindent$\Pi_{176}$:
\begin{ptree}
\AXC{$\stackrel{\Pi_{161}}{\vdots}$}\noLine
\UIC{$r \in \mathbb{R} \to \lvert r \rvert < 1 \to \varphi_{3}$}
\AXC{$\stackrel{\Pi_{175}}{\vdots}$}\noLine
\UIC{$(r \in \mathbb{R} \to \lvert r \rvert < 1 \to \varphi_{3}) \to \varphi_{169} \to r \in \mathbb{R} \to \varphi_{2}$}
\BIC{$\varphi_{169} \to r \in \mathbb{R} \to \varphi_{2}$}
\end{ptree}
\noindent 主木:
\begin{ptree}
\AXC{$z = w \to \varphi_{1}$\lbl{A6}}
\UIC{$\forall w\,(z = w \to \varphi_{1})$}
\UIC{$\forall z\,\forall w\,(z = w \to \varphi_{1})$}
\AXC{$\forall z\,\forall w\,(z = w \to \varphi_{1}) \to \varphi_{132}$\lbl{A2}}
\BIC{$\varphi_{132}$}
\AXC{$\varphi_{132} \to \varphi_{169}$\lbl{A2}}
\BIC{$\varphi_{169}$}
\AXC{$\stackrel{\Pi_{176}}{\vdots}$}\noLine
\UIC{$\varphi_{169} \to r \in \mathbb{R} \to \varphi_{2}$}
\BIC{$r \in \mathbb{R} \to \varphi_{2}$}
\UIC{$\forall r \in \mathbb{R}\,(\varphi_{2})$}
\end{ptree}
\begin{abbrevlist}
\item $\varphi_{1} :\equiv \mathrm{Conv}(\Sigma {(r)^{\bullet}},z) \to \mathrm{Conv}(\Sigma {(r)^{\bullet}},w)$
\item $\varphi_{2} :\equiv \lvert r \rvert < 1 \to \mathrm{Conv}(\Sigma {(r)^{\bullet}},r \cdot (1 - r)^{-1})$
\item $\varphi_{3} :\equiv \mathrm{Conv}(\Sigma {(r)^{\bullet}},(1 - r)^{-1} \cdot (r + -1 \cdot 0))$
\item $\varphi_{4} :\equiv \forall x\,(x = x) \to (1 - r)^{-1} \cdot (r + 0) = (1 - r)^{-1} \cdot (r + 0)$
\item $\varphi_{5} :\equiv \forall a\,\forall s\,\forall \alpha\,(\mathrm{Conv}(a,\alpha) \to \mathrm{SubIdx}(s) \to \mathrm{Conv}(a \circ s,\alpha))$
\item $\varphi_{6} :\equiv \mathrm{Conv}({(r)^{\bullet}},0) \to \mathrm{SubIdx}(\mathsf{s}) \to \mathrm{Conv}({(r)^{\bullet}} \circ \mathsf{s},0)$
\item $\varphi_{7} :\equiv \forall s\,\forall \alpha\,(\mathrm{Conv}({(r)^{\bullet}},\alpha) \to \mathrm{SubIdx}(s) \to \mathrm{Conv}({(r)^{\bullet}} \circ s,\alpha))$
\item $\varphi_{8} :\equiv \forall c\,\forall a\,\forall \alpha\,(c \in \mathbb{R} \to \mathrm{Conv}(a,\alpha) \to \mathrm{Conv}(c \odot a,c \cdot \alpha))$
\item $\varphi_{9} :\equiv \forall x\,\forall y\,\forall z\,(x \leq y \to y < z \to x < z) \to \forall y\,\forall z\,(r \leq y \to y < z \to r < z)$
\item $\varphi_{10} :\equiv r \in \mathbb{R} \to \lvert r \rvert < 1 \to n \in \mathbb{N} \to \mathrm{Seq}({(r)^{\bullet}} \circ \mathsf{s})$
\item $\varphi_{11} :\equiv \forall \alpha\,(\mathrm{Conv}({(r)^{\bullet}},\alpha) \to \mathrm{SubIdx}(\mathsf{s}) \to \mathrm{Conv}({(r)^{\bullet}} \circ \mathsf{s},\alpha))$
\item $\varphi_{12} :\equiv \forall n\,(\mathrm{Seq}(\Sigma {(r)^{\bullet}}) \to n \in \mathbb{N} \to {\textstyle\sum_{k=1}^{n} {{r}^{k}}} \in \mathbb{R})$
\item $\varphi_{13} :\equiv \forall c \in \mathbb{R}\,(\mathrm{Conv}(\overline{c},c)) \to r \in \mathbb{R} \to \mathrm{Conv}(\overline{r},r)$
\item $\varphi_{14} :\equiv r \in \mathbb{R} \to \lvert r \rvert < 1 \to n \in \mathbb{N} \to (1 - r)^{-1} \cdot (1 - r) = 1$
\item $\varphi_{15} :\equiv \forall k \in \mathbb{N}\,(\mathsf{s}(k) = S(k)) \to n \in \mathbb{N} \to \mathsf{s}(n) = S(n)$
\item $\varphi_{16} :\equiv \lvert r \rvert < 1 \to (1 - r)^{-1} \cdot (r + -1 \cdot 0) = (1 - r)^{-1} \cdot r$
\item $\varphi_{17} :\equiv \lvert r \rvert < 1 \to (1 - r)^{-1} \cdot (r + -1 \cdot 0) = r \cdot (1 - r)^{-1}$
\item $\varphi_{18} :\equiv (1 - r)^{-1} \cdot (r + -1 \cdot 0) = (1 - r)^{-1} \cdot (r + -1 \cdot 0)$
\item $\varphi_{19} :\equiv \forall y\,\forall z\,(r \leq y \to y < z \to r < z) \to \forall z\,(r \leq \lvert r \rvert \to \lvert r \rvert < z \to r < z)$
\item $\varphi_{20} :\equiv \forall a\,\forall \alpha\,\forall b\,\forall \beta\,(\mathrm{Conv}(a,\alpha) \to \mathrm{Conv}(b,\beta) \to \mathrm{Conv}(a \oplus b,\alpha + \beta))$
\item $\varphi_{21} :\equiv \forall y\,\forall z\,({(r)^{\bullet}} \circ \mathsf{s}(n) = y \to y = z \to {(r)^{\bullet}} \circ \mathsf{s}(n) = z)$
\item $\varphi_{22} :\equiv r \in \mathbb{R} \to \lvert r \rvert < 1 \to (1 - r)^{-1} \cdot (r + 0) = (1 - r)^{-1} \cdot r$
\item $\varphi_{23} :\equiv \mathrm{SubIdx}(\mathsf{s}) \to r \in \mathbb{R} \to \lvert r \rvert < 1 \to \mathrm{Conv}({(r)^{\bullet}} \circ \mathsf{s},0)$
\item $\varphi_{24} :\equiv (1 - r)^{-1} \cdot (r + -1 \cdot {{r}^{S(n)}}) = (1 - r)^{-1} \cdot (r + w)$
\item $\varphi_{25} :\equiv (1 - r)^{-1} \cdot (r + -1 \cdot {{r}^{S(n)}}) = (1 - r)^{-1} \cdot (r + z)$
\item $\varphi_{26} :\equiv \lvert r \rvert < 1 \to (1 - r)^{-1} \cdot (r + -1 \cdot 0) = (1 - r)^{-1} \cdot (r + 0)$
\item $\varphi_{27} :\equiv \forall b\,\forall \beta\,(\mathrm{Conv}(\overline{r},r) \to \mathrm{Conv}(b,\beta) \to \mathrm{Conv}(\overline{r} \oplus b,r + \beta))$
\item $\varphi_{28} :\equiv \forall z\,(r \leq \lvert r \rvert \to \lvert r \rvert < z \to r < z) \to r \leq \lvert r \rvert \to \lvert r \rvert < 1 \to r < 1$
\item $\varphi_{29} :\equiv \forall x\,\forall y\,(x = y \to y = x) \to \forall y\,(r - {{r}^{S(n)}} = y \to y = r - {{r}^{S(n)}})$
\item $\varphi_{30} :\equiv (1 - r)^{-1} \in \mathbb{R} \to r \in \mathbb{R} \to (1 - r)^{-1} \cdot r = r \cdot (1 - r)^{-1}$
\item $\varphi_{31} :\equiv r \in \mathbb{R} \to \lvert r \rvert < 1 \to n \in \mathbb{N} \to \forall k \in \mathbb{N}\,(\mathsf{s}(k) = S(k))$
\item $\varphi_{32} :\equiv -1 \odot {(r)^{\bullet}} \circ \mathsf{s}(n) = -1 \cdot {(r)^{\bullet}} \circ \mathsf{s}(n)$
\item $\varphi_{33} :\equiv \forall a\,\forall \alpha\,(-1 \in \mathbb{R} \to \mathrm{Conv}(a,\alpha) \to \mathrm{Conv}(-1 \odot a,-1 \cdot \alpha))$
\item $\varphi_{34} :\equiv \forall x\,\forall y\,\forall z\,(x \in \mathbb{R} \to y \in \mathbb{R} \to z \in \mathbb{R} \to x \cdot y \cdot z = x \cdot (y \cdot z))$
\item $\varphi_{35} :\equiv (1 - r)^{-1} \cdot (r - {{r}^{S(n)}}) = (1 - r)^{-1} \cdot (r - {{r}^{S(n)}})$
\item $\varphi_{36} :\equiv \forall y\,((1 - r)^{-1} \in \mathbb{R} \to y \in \mathbb{R} \to (1 - r)^{-1} \cdot y = y \cdot (1 - r)^{-1})$
\item $\varphi_{37} :\equiv \lvert r \rvert < 1 \to n \in \mathbb{N} \to \mathrm{Seq}(\overline{r}) \land \forall n \in \mathbb{N}\,(\overline{r}(n) = r)$
\item $\varphi_{38} :\equiv r \in \mathbb{R} \to \lvert r \rvert < 1 \to n \in \mathbb{N} \to \mathrm{Seq}(-1 \odot {(r)^{\bullet}} \circ \mathsf{s})$
\item $\varphi_{39} :\equiv r \in \mathbb{R} \to \lvert r \rvert < 1 \to \mathrm{Conv}(-1 \odot {(r)^{\bullet}} \circ \mathsf{s},-1 \cdot 0)$
\item $\varphi_{40} :\equiv \forall \alpha\,\forall b\,\forall \beta\,(\mathrm{Conv}(\overline{r},\alpha) \to \mathrm{Conv}(b,\beta) \to \mathrm{Conv}(\overline{r} \oplus b,\alpha + \beta))$
\item $\varphi_{41} :\equiv (r \in \mathbb{R} \to \lvert r \rvert < 1 \to \lvert r \rvert < 1) \to r \in \mathbb{R} \to \lvert r \rvert < 1 \to r < 1$
\item $\varphi_{42} :\equiv (r \in \mathbb{R} \to \lvert r \rvert < 1 \to r < 1) \to r \in \mathbb{R} \to \lvert r \rvert < 1 \to 0 < 1 - r$
\item $\varphi_{43} :\equiv \mathrm{Conv}({(r)^{\bullet}} \circ \mathsf{s},0) \to \mathrm{Conv}(-1 \odot {(r)^{\bullet}} \circ \mathsf{s},-1 \cdot 0)$
\item $\varphi_{44} :\equiv (r \in \mathbb{R} \to \lvert r \rvert < 1 \to r \in \mathbb{R}) \to r \in \mathbb{R} \to \lvert r \rvert < 1 \to r + 0 = r$
\item $\varphi_{45} :\equiv \lvert r \rvert < 1 \to n \in \mathbb{N} \to \mathrm{Seq}(\overline{r} \oplus -1 \odot {(r)^{\bullet}} \circ \mathsf{s})$
\item $\varphi_{46} :\equiv (1 - r)^{-1} \cdot (r + -{{r}^{S(n)}}) = (1 - r)^{-1} \cdot (r - {{r}^{S(n)}})$
\item $\varphi_{47} :\equiv (r \in \mathbb{R} \to \lvert r \rvert < 1 \to 0 < 1 - r) \to r \in \mathbb{R} \to \lvert r \rvert < 1 \to 0 \neq 1 - r$
\item $\varphi_{48} :\equiv \forall a\,\forall b\,\forall \alpha\,(\mathrm{Seq}(b) \to \forall n \in \mathbb{N}\,(a_{n} = b_{n}) \to \mathrm{Conv}(a,\alpha) \to \mathrm{Conv}(b,\alpha))$
\item $\varphi_{49} :\equiv r \in \mathbb{R} \to \lvert r \rvert < 1 \to n \in \mathbb{N} \to {(r)^{\bullet}} \circ \mathsf{s}(n) = {{r}^{S(n)}}$
\item $\varphi_{50} :\equiv (1 - r)^{-1} \cdot (r + -1 \odot {(r)^{\bullet}} \circ \mathsf{s}(n)) = (1 - r)^{-1} \cdot (r + w)$
\item $\varphi_{51} :\equiv (1 - r)^{-1} \cdot (r + -1 \odot {(r)^{\bullet}} \circ \mathsf{s}(n)) = (1 - r)^{-1} \cdot (r + z)$
\item $\varphi_{52} :\equiv (1 - r)^{-1} \cdot (r - {{r}^{S(n)}}) = (1 - r)^{-1} \cdot ((1 - r) \cdot {\textstyle\sum_{k=1}^{n} {{r}^{k}}})$
\item $\varphi_{53} :\equiv (1 - r)^{-1} \cdot (r + 0) = (1 - r)^{-1} \cdot z \to (1 - r)^{-1} \cdot (r + 0) = (1 - r)^{-1} \cdot w$
\item $\varphi_{54} :\equiv \forall x \in \mathbb{R}\,(-\lvert x \rvert \leq x \land x \leq \lvert x \rvert) \to r \in \mathbb{R} \to -\lvert r \rvert \leq r \land r \leq \lvert r \rvert$
\item $\varphi_{55} :\equiv (1 - r)^{-1} \cdot (r + -{{r}^{S(n)}}) = (1 - r)^{-1} \cdot (r + -{{r}^{S(n)}})$
\item $\varphi_{56} :\equiv \forall y\,\forall z\,((1 - r)^{-1} \cdot (r + -1 \cdot 0) = y \to y = z \to (1 - r)^{-1} \cdot (r + -1 \cdot 0) = z)$
\item $\varphi_{57} :\equiv \lvert r \rvert < 1 \to \mathrm{Conv}(\overline{r} \oplus -1 \odot {(r)^{\bullet}} \circ \mathsf{s},r + -1 \cdot 0)$
\item $\varphi_{58} :\equiv \mathrm{Seq}({(r)^{\bullet}}) \land {{r}^{1_{\omega}}} = r \land \forall n \in \mathbb{N}\,({{r}^{S(n)}} = {{r}^{n}} \cdot r)$
\item $\varphi_{59} :\equiv (r \in \mathbb{R} \to \lvert r \rvert < 1 \to r \in \mathbb{R}) \to r \in \mathbb{R} \to \lvert r \rvert < 1 \to \mathrm{Conv}(\overline{r},r)$
\item $\varphi_{60} :\equiv (r \in \mathbb{R} \to \lvert r \rvert < 1 \to r \in \mathbb{R}) \to r \in \mathbb{R} \to \lvert r \rvert < 1 \to 1 - r \in \mathbb{R}$
\item $\varphi_{61} :\equiv (1 - r)^{-1} \cdot \overline{r} \oplus -1 \odot {(r)^{\bullet}} \circ \mathsf{s}(n) = (1 - r)^{-1} \cdot w$
\item $\varphi_{62} :\equiv (1 - r)^{-1} \cdot \overline{r} \oplus -1 \odot {(r)^{\bullet}} \circ \mathsf{s}(n) = (1 - r)^{-1} \cdot z$
\item $\varphi_{63} :\equiv \forall x\,\forall y\,(x = y \to y = x) \to \forall y\,((1 - r) \cdot {\textstyle\sum_{k=1}^{n} {{r}^{k}}} = y \to y = (1 - r) \cdot {\textstyle\sum_{k=1}^{n} {{r}^{k}}})$
\item $\varphi_{64} :\equiv \lvert r \rvert < 1 \to n \in \mathbb{N} \to n \in \mathbb{N} \to {(r)^{\bullet}} \circ \mathsf{s}(n) = {{r}^{\mathsf{s}(n)}}$
\item $\varphi_{65} :\equiv (1 - r)^{-1} \cdot r = r \cdot (1 - r)^{-1} \to (1 - r)^{-1} \cdot (r + -1 \cdot 0) = r \cdot (1 - r)^{-1}$
\item $\varphi_{66} :\equiv (1 - r)^{-1} \cdot (r + 0) = (1 - r)^{-1} \cdot (r + 0) \to (1 - r)^{-1} \cdot (r + 0) = (1 - r)^{-1} \cdot r$
\item $\varphi_{67} :\equiv r \in \mathbb{R} \to \lvert r \rvert < 1 \to n \in \mathbb{N} \to {(r)^{\bullet}} \circ \mathsf{s}(n) = {{r}^{\mathsf{s}(n)}}$
\item $\varphi_{68} :\equiv (r \in \mathbb{R} \to \lvert r \rvert < 1 \to -1 \in \mathbb{R}) \to r \in \mathbb{R} \to \lvert r \rvert < 1 \to -1 \cdot 0 = 0$
\item $\varphi_{69} :\equiv (1 - r)^{-1} \cdot (r + -1 \cdot {(r)^{\bullet}} \circ \mathsf{s}(n)) = (1 - r)^{-1} \cdot (r + -1 \cdot w)$
\item $\varphi_{70} :\equiv (1 - r)^{-1} \cdot (r + -1 \cdot {(r)^{\bullet}} \circ \mathsf{s}(n)) = (1 - r)^{-1} \cdot (r + -1 \cdot z)$
\item $\varphi_{71} :\equiv \mathrm{Seq}({(r)^{\bullet}} \circ \mathsf{s}) \land \forall k \in \mathbb{N}\,({(r)^{\bullet}} \circ \mathsf{s}(k) = {{r}^{\mathsf{s}(k)}})$
\item $\varphi_{72} :\equiv \forall r \in \mathbb{R}\,(\lvert r \rvert < 1 \to \mathrm{Conv}({(r)^{\bullet}},0)) \to r \in \mathbb{R} \to \lvert r \rvert < 1 \to \mathrm{Conv}({(r)^{\bullet}},0)$
\item $\varphi_{73} :\equiv (1 \in \mathbb{R} \to -1 \in \mathbb{R}) \to 1 \in \mathbb{R} \to r \in \mathbb{R} \to \lvert r \rvert < 1 \to -1 \in \mathbb{R}$
\item $\varphi_{74} :\equiv (r \in \mathbb{R} \to \lvert r \rvert < 1 \to 1 - r \neq 0) \to r \in \mathbb{R} \to \lvert r \rvert < 1 \to (1 - r)^{-1} \in \mathbb{R}$
\item $\varphi_{75} :\equiv \forall a\,\forall \alpha\,((1 - r)^{-1} \in \mathbb{R} \to \mathrm{Conv}(a,\alpha) \to \mathrm{Conv}((1 - r)^{-1} \odot a,(1 - r)^{-1} \cdot \alpha))$
\item $\varphi_{76} :\equiv (1 - r)^{-1} \cdot (r + -1 \cdot {{r}^{S(n)}}) = (1 - r)^{-1} \cdot (r + -{{r}^{S(n)}})$
\item $\varphi_{77} :\equiv (1 - r)^{-1} \cdot (1 - r) \cdot {\textstyle\sum_{k=1}^{n} {{r}^{k}}} = (1 - r)^{-1} \cdot (1 - r) \cdot {\textstyle\sum_{k=1}^{n} {{r}^{k}}}$
\item $\varphi_{78} :\equiv (1 - r)^{-1} \cdot ((1 - r) \cdot {\textstyle\sum_{k=1}^{n} {{r}^{k}}}) = (1 - r)^{-1} \cdot (1 - r) \cdot {\textstyle\sum_{k=1}^{n} {{r}^{k}}}$
\item $\varphi_{79} :\equiv (1 - r)^{-1} \cdot (1 - r) \cdot {\textstyle\sum_{k=1}^{n} {{r}^{k}}} = (1 - r)^{-1} \cdot ((1 - r) \cdot {\textstyle\sum_{k=1}^{n} {{r}^{k}}})$
\item $\varphi_{80} :\equiv \lvert r \rvert < 1 \to n \in \mathbb{N} \to n \in \mathbb{N} \to (1 - r) \cdot {\textstyle\sum_{k=1}^{n} {{r}^{k}}} = r - {{r}^{S(n)}}$
\item $\varphi_{81} :\equiv \lvert r \rvert < 1 \to n \in \mathbb{N} \to \forall k \in \mathbb{N}\,({(r)^{\bullet}} \circ \mathsf{s}(k) = {{r}^{\mathsf{s}(k)}})$
\item $\varphi_{82} :\equiv \forall y\,(1 \in \mathbb{R} \to y \in \mathbb{R} \to 1 - y \in \mathbb{R}) \to 1 \in \mathbb{R} \to r \in \mathbb{R} \to 1 - r \in \mathbb{R}$
\item $\varphi_{83} :\equiv r \in \mathbb{R} \to \lvert r \rvert < 1 \to n \in \mathbb{N} \to (1 - r) \cdot {\textstyle\sum_{k=1}^{n} {{r}^{k}}} = r - {{r}^{S(n)}}$
\item $\varphi_{84} :\equiv r \in \mathbb{R} \to \lvert r \rvert < 1 \to n \in \mathbb{N} \to r - {{r}^{S(n)}} = (1 - r) \cdot {\textstyle\sum_{k=1}^{n} {{r}^{k}}}$
\item $\varphi_{85} :\equiv (1 - r)^{-1} \odot \overline{r} \oplus -1 \odot {(r)^{\bullet}} \circ \mathsf{s}(n) = {\textstyle\sum_{k=1}^{n} {{r}^{k}}}$
\item $\varphi_{86} :\equiv (r \in \mathbb{R} \to \lvert r \rvert < 1 \to -\lvert r \rvert \leq r \land r \leq \lvert r \rvert) \to r \in \mathbb{R} \to \lvert r \rvert < 1 \to r \leq \lvert r \rvert$
\item $\varphi_{87} :\equiv \forall a\,\forall s\,(\mathrm{Seq}(a) \to \mathrm{NSeq}(s) \to \mathrm{Seq}(a \circ s) \land \forall k \in \mathbb{N}\,(a \circ s(k) = a_{s_{k}}))$
\item $\varphi_{88} :\equiv (1 - r)^{-1} \cdot (r + 0) = (1 - r)^{-1} \cdot r \to (1 - r)^{-1} \cdot (r + -1 \cdot 0) = (1 - r)^{-1} \cdot r$
\item $\varphi_{89} :\equiv \forall x\,\forall y\,(x \in \mathbb{R} \to y \in \mathbb{R} \to x \cdot y = y \cdot x) \to \forall y\,(1 \in \mathbb{R} \to y \in \mathbb{R} \to 1 \cdot y = y \cdot 1)$
\item $\varphi_{90} :\equiv \forall c\,\forall a\,(c \in \mathbb{R} \to \mathrm{Seq}(a) \to \mathrm{Seq}(c \odot a) \land \forall n \in \mathbb{N}\,(c \odot a(n) = c \cdot a_{n}))$
\item $\varphi_{91} :\equiv \lvert r \rvert < 1 \to n \in \mathbb{N} \to \forall n \in \mathbb{N}\,((1 - r) \cdot {\textstyle\sum_{k=1}^{n} {{r}^{k}}} = r - {{r}^{S(n)}})$
\item $\varphi_{92} :\equiv (r \in \mathbb{R} \to \lvert r \rvert < 1 \to r \in \mathbb{R}) \to r \in \mathbb{R} \to \lvert r \rvert < 1 \to -\lvert r \rvert \leq r \land r \leq \lvert r \rvert$
\item $\varphi_{93} :\equiv \forall x \in \mathbb{R}\,(x \neq 0 \to x^{-1} \in \mathbb{R}) \to 1 - r \in \mathbb{R} \to 1 - r \neq 0 \to (1 - r)^{-1} \in \mathbb{R}$
\item $\varphi_{94} :\equiv (1 - r)^{-1} \cdot (r + -1 \cdot {{r}^{S(n)}}) = (1 - r)^{-1} \cdot (r + -1 \cdot {{r}^{S(n)}})$
\item $\varphi_{95} :\equiv \forall x \in \mathbb{R}\,(x \neq 0 \to x^{-1} \cdot x = 1) \to 1 - r \in \mathbb{R} \to 1 - r \neq 0 \to (1 - r)^{-1} \cdot (1 - r) = 1$
\item $\varphi_{96} :\equiv (1 - r)^{-1} \odot \overline{r} \oplus -1 \odot {(r)^{\bullet}} \circ \mathsf{s}(n) = {\textstyle\sum_{k=1}^{n} {{r}^{k}}} \cdot 1$
\item $\varphi_{97} :\equiv (1 - r)^{-1} \odot \overline{r} \oplus -1 \odot {(r)^{\bullet}} \circ \mathsf{s}(n) = 1 \cdot {\textstyle\sum_{k=1}^{n} {{r}^{k}}}$
\item $\varphi_{98} :\equiv \lvert r \rvert < 1 \to n \in \mathbb{N} \to (1 - r)^{-1} \cdot (1 - r) \cdot {\textstyle\sum_{k=1}^{n} {{r}^{k}}} = 1 \cdot {\textstyle\sum_{k=1}^{n} {{r}^{k}}}$
\item $\varphi_{99} :\equiv \forall x\,\forall y\,(x \in \mathbb{R} \to y \in \mathbb{R} \to x - y \in \mathbb{R}) \to \forall y\,(1 \in \mathbb{R} \to y \in \mathbb{R} \to 1 - y \in \mathbb{R})$
\item $\varphi_{100} :\equiv \forall a\,\forall b\,(\mathrm{Seq}(a) \to \mathrm{Seq}(b) \to \mathrm{Seq}(a \oplus b) \land \forall n \in \mathbb{N}\,(a \oplus b(n) = a_{n} + b_{n}))$
\item $\varphi_{101} :\equiv \forall y\,((1 - r)^{-1} \cdot (1 - r) \cdot {\textstyle\sum_{k=1}^{n} {{r}^{k}}} = y \to y = (1 - r)^{-1} \cdot (1 - r) \cdot {\textstyle\sum_{k=1}^{n} {{r}^{k}}})$
\item $\varphi_{102} :\equiv \forall a\,\forall n\,(\mathrm{Seq}(a) \to n \in \mathbb{N} \to a_{n} \in \mathbb{R}) \to \forall n\,(\mathrm{Seq}({(r)^{\bullet}}) \to n \in \mathbb{N} \to {{r}^{n}} \in \mathbb{R})$
\item $\varphi_{103} :\equiv \mathrm{NSeq}(\mathsf{s}) \land \forall k \in \mathbb{N}\,(\mathsf{s}(k) = S(k)) \to r \in \mathbb{R} \to \lvert r \rvert < 1 \to n \in \mathbb{N} \to \mathrm{NSeq}(\mathsf{s})$
\item $\varphi_{104} :\equiv (r \in \mathbb{R} \to \lvert r \rvert < 1 \to 0 \neq 1 - r) \to (1 - r = 0 \to 0 = 1 - r) \to r \in \mathbb{R} \to \lvert r \rvert < 1 \to 1 - r \neq 0$
\item $\varphi_{105} :\equiv \forall \alpha\,(-1 \in \mathbb{R} \to \mathrm{Conv}({(r)^{\bullet}} \circ \mathsf{s},\alpha) \to \mathrm{Conv}(-1 \odot {(r)^{\bullet}} \circ \mathsf{s},-1 \cdot \alpha))$
\item $\varphi_{106} :\equiv \forall y\,(r \in \mathbb{R} \to y \in \mathbb{R} \to r < y \to 0 < y - r) \to r \in \mathbb{R} \to 1 \in \mathbb{R} \to r < 1 \to 0 < 1 - r$
\item $\varphi_{107} :\equiv (r \in \mathbb{R} \to \lvert r \rvert < 1 \to 1 - r \neq 0) \to r \in \mathbb{R} \to \lvert r \rvert < 1 \to n \in \mathbb{N} \to (1 - r)^{-1} \in \mathbb{R}$
\item $\varphi_{108} :\equiv \forall y\,\forall z\,((1 - r)^{-1} \in \mathbb{R} \to y \in \mathbb{R} \to z \in \mathbb{R} \to (1 - r)^{-1} \cdot y \cdot z = (1 - r)^{-1} \cdot (y \cdot z))$
\item $\varphi_{109} :\equiv (1 \in \mathbb{R} \to -1 \in \mathbb{R}) \to 1 \in \mathbb{R} \to r \in \mathbb{R} \to \lvert r \rvert < 1 \to n \in \mathbb{N} \to -1 \in \mathbb{R}$
\item $\varphi_{110} :\equiv \forall w\,(r + 0 = w \to (1 - r)^{-1} \cdot (r + 0) = (1 - r)^{-1} \cdot (r + 0) \to (1 - r)^{-1} \cdot (r + 0) = (1 - r)^{-1} \cdot w)$
\item $\varphi_{111} :\equiv \forall x \in \mathbb{R}\,(x \cdot 1 = x) \to {\textstyle\sum_{k=1}^{n} {{r}^{k}}} \in \mathbb{R} \to {\textstyle\sum_{k=1}^{n} {{r}^{k}}} \cdot 1 = {\textstyle\sum_{k=1}^{n} {{r}^{k}}}$
\item $\varphi_{112} :\equiv (1 - r)^{-1} \cdot (r + -1 \cdot {(r)^{\bullet}} \circ \mathsf{s}(n)) = (1 - r)^{-1} \cdot (r + -1 \cdot {{r}^{S(n)}})$
\item $\varphi_{113} :\equiv (r \in \mathbb{R} \to \lvert r \rvert < 1 \to r \in \mathbb{R}) \to r \in \mathbb{R} \to \lvert r \rvert < 1 \to (1 - r)^{-1} \cdot r = r \cdot (1 - r)^{-1}$
\item $\varphi_{114} :\equiv \mathrm{Conv}((1 - r)^{-1} \odot \overline{r} \oplus -1 \odot {(r)^{\bullet}} \circ \mathsf{s},(1 - r)^{-1} \cdot (r + -1 \cdot 0))$
\item $\varphi_{115} :\equiv (1 - r) \cdot {\textstyle\sum_{k=1}^{n} {{r}^{k}}} = r - {{r}^{S(n)}} \to r - {{r}^{S(n)}} = (1 - r) \cdot {\textstyle\sum_{k=1}^{n} {{r}^{k}}}$
\item $\varphi_{116} :\equiv (r \in \mathbb{R} \to \lvert r \rvert < 1 \to n \in \mathbb{N} \to n \in \mathbb{N}) \to r \in \mathbb{R} \to \lvert r \rvert < 1 \to n \in \mathbb{N} \to S(n) \in \mathbb{N}$
\item $\varphi_{117} :\equiv \forall s\,(\mathrm{Seq}({(r)^{\bullet}}) \to \mathrm{NSeq}(s) \to \mathrm{Seq}({(r)^{\bullet}} \circ s) \land \forall k \in \mathbb{N}\,({(r)^{\bullet}} \circ s(k) = {{r}^{s_{k}}}))$
\item $\varphi_{118} :\equiv \forall x\,\forall y\,(x \in \mathbb{R} \to y \in \mathbb{R} \to x < y \to 0 < y - x) \to \forall y\,(r \in \mathbb{R} \to y \in \mathbb{R} \to r < y \to 0 < y - r)$
\item $\varphi_{119} :\equiv \forall n\,(\mathrm{Seq}({(r)^{\bullet}}) \to n \in \mathbb{N} \to {{r}^{n}} \in \mathbb{R}) \to \mathrm{Seq}({(r)^{\bullet}}) \to S(n) \in \mathbb{N} \to {{r}^{S(n)}} \in \mathbb{R}$
\item $\varphi_{120} :\equiv \forall z\,({(r)^{\bullet}} \circ \mathsf{s}(n) = {{r}^{\mathsf{s}(n)}} \to {{r}^{\mathsf{s}(n)}} = z \to {(r)^{\bullet}} \circ \mathsf{s}(n) = z)$
\item $\varphi_{121} :\equiv (1 - r)^{-1} \odot \overline{r} \oplus -1 \odot {(r)^{\bullet}} \circ \mathsf{s}(n) = (1 - r)^{-1} \cdot (r - {{r}^{S(n)}})$
\item $\varphi_{122} :\equiv (r \in \mathbb{R} \to \lvert r \rvert < 1 \to n \in \mathbb{N} \to r \in \mathbb{R}) \to r \in \mathbb{R} \to \lvert r \rvert < 1 \to n \in \mathbb{N} \to 1 - r \in \mathbb{R}$
\item $\varphi_{123} :\equiv (r \in \mathbb{R} \to \lvert r \rvert < 1 \to n \in \mathbb{N} \to n \in \mathbb{N}) \to r \in \mathbb{R} \to \lvert r \rvert < 1 \to n \in \mathbb{N} \to \mathsf{s}(n) = S(n)$
\item $\varphi_{124} :\equiv (r \in \mathbb{R} \to \lvert r \rvert < 1 \to n \in \mathbb{N} \to n \in \mathbb{N}) \to r \in \mathbb{R} \to \lvert r \rvert < 1 \to n \in \mathbb{N} \to \overline{r}(n) = r$
\item $\varphi_{125} :\equiv \overline{r} \oplus -1 \odot {(r)^{\bullet}} \circ \mathsf{s}(n) = \overline{r}(n) + -1 \odot {(r)^{\bullet}} \circ \mathsf{s}(n)$
\item $\varphi_{126} :\equiv (1 - r)^{-1} \cdot (r + -1 \cdot {(r)^{\bullet}} \circ \mathsf{s}(n)) = (1 - r)^{-1} \cdot (r + -1 \cdot {(r)^{\bullet}} \circ \mathsf{s}(n))$
\item $\varphi_{127} :\equiv \forall a\,(-1 \in \mathbb{R} \to \mathrm{Seq}(a) \to \mathrm{Seq}(-1 \odot a) \land \forall n \in \mathbb{N}\,(-1 \odot a(n) = -1 \cdot a_{n}))$
\item $\varphi_{128} :\equiv (1 - r)^{-1} \cdot (r - {{r}^{S(n)}}) = (1 - r)^{-1} \cdot z \to (1 - r)^{-1} \cdot (r - {{r}^{S(n)}}) = (1 - r)^{-1} \cdot w$
\item $\varphi_{129} :\equiv \forall z\,((1 - r)^{-1} \cdot (r + -1 \cdot 0) = (1 - r)^{-1} \cdot r \to (1 - r)^{-1} \cdot r = z \to (1 - r)^{-1} \cdot (r + -1 \cdot 0) = z)$
\item $\varphi_{130} :\equiv (r \in \mathbb{R} \to \lvert r \rvert < 1 \to \mathrm{Conv}({(r)^{\bullet}},0)) \to \varphi_{23}$
\item $\varphi_{131} :\equiv (1 - r)^{-1} \cdot (r + -1 \cdot 0) = (1 - r)^{-1} \cdot (r + z) \to (1 - r)^{-1} \cdot (r + -1 \cdot 0) = (1 - r)^{-1} \cdot (r + w)$
\item $\varphi_{132} :\equiv \forall w\,((1 - r)^{-1} \cdot (r + -1 \cdot 0) = w \to \varphi_{3} \to \mathrm{Conv}(\Sigma {(r)^{\bullet}},w))$
\item $\varphi_{133} :\equiv \forall z\,((1 - r)^{-1} \in \mathbb{R} \to 1 - r \in \mathbb{R} \to z \in \mathbb{R} \to (1 - r)^{-1} \cdot (1 - r) \cdot z = (1 - r)^{-1} \cdot ((1 - r) \cdot z))$
\item $\varphi_{134} :\equiv \forall x \in \mathbb{R}\,(-1 \cdot x = -x) \to {{r}^{S(n)}} \in \mathbb{R} \to -1 \cdot {{r}^{S(n)}} = -{{r}^{S(n)}}$
\item $\varphi_{135} :\equiv (1 - r)^{-1} \odot \overline{r} \oplus -1 \odot {(r)^{\bullet}} \circ \mathsf{s}(n) = (1 - r)^{-1} \cdot (r + -{{r}^{S(n)}})$
\item $\varphi_{136} :\equiv (1 - r)^{-1} \cdot (r + -1 \odot {(r)^{\bullet}} \circ \mathsf{s}(n)) = (1 - r)^{-1} \cdot (r + -1 \cdot {(r)^{\bullet}} \circ \mathsf{s}(n))$
\item $\varphi_{137} :\equiv (1 - r)^{-1} \cdot (r + -1 \odot {(r)^{\bullet}} \circ \mathsf{s}(n)) = (1 - r)^{-1} \cdot (r + -1 \odot {(r)^{\bullet}} \circ \mathsf{s}(n))$
\item $\varphi_{138} :\equiv (1 - r)^{-1} \odot \overline{r} \oplus -1 \odot {(r)^{\bullet}} \circ \mathsf{s}(n) = (1 - r)^{-1} \cdot (1 - r) \cdot {\textstyle\sum_{k=1}^{n} {{r}^{k}}}$
\item $\varphi_{139} :\equiv (1 - r)^{-1} \odot \overline{r} \oplus -1 \odot {(r)^{\bullet}} \circ \mathsf{s}(n) = (1 - r)^{-1} \cdot ((1 - r) \cdot {\textstyle\sum_{k=1}^{n} {{r}^{k}}})$
\item $\varphi_{140} :\equiv (1 - r)^{-1} \odot \overline{r} \oplus -1 \odot {(r)^{\bullet}} \circ \mathsf{s}(n) = (1 - r)^{-1} \cdot (r + -1 \cdot {{r}^{S(n)}})$
\item $\varphi_{141} :\equiv \forall b\,(\mathrm{Seq}(\overline{r}) \to \mathrm{Seq}(b) \to \mathrm{Seq}(\overline{r} \oplus b) \land \forall n \in \mathbb{N}\,(\overline{r} \oplus b(n) = \overline{r}(n) + b_{n}))$
\item $\varphi_{142} :\equiv (r \in \mathbb{R} \to \lvert r \rvert < 1 \to n \in \mathbb{N} \to n \in \mathbb{N}) \to r \in \mathbb{R} \to \lvert r \rvert < 1 \to n \in \mathbb{N} \to {\textstyle\sum_{k=1}^{n} {{r}^{k}}} \in \mathbb{R}$
\item $\varphi_{143} :\equiv (r \in \mathbb{R} \to \lvert r \rvert < 1 \to n \in \mathbb{N} \to S(n) \in \mathbb{N}) \to r \in \mathbb{R} \to \lvert r \rvert < 1 \to n \in \mathbb{N} \to {{r}^{S(n)}} \in \mathbb{R}$
\item $\varphi_{144} :\equiv \mathrm{Conv}(-1 \odot {(r)^{\bullet}} \circ \mathsf{s},-1 \cdot 0) \to \mathrm{Conv}(\overline{r} \oplus -1 \odot {(r)^{\bullet}} \circ \mathsf{s},r + -1 \cdot 0)$
\item $\varphi_{145} :\equiv (1 - r)^{-1} \cdot (r + -{{r}^{S(n)}}) = (1 - r)^{-1} \cdot z \to (1 - r)^{-1} \cdot (r + -{{r}^{S(n)}}) = (1 - r)^{-1} \cdot w$
\item $\varphi_{146} :\equiv \forall z\,((1 - r)^{-1} \cdot (r + -1 \cdot 0) = (1 - r)^{-1} \cdot (r + 0) \to (1 - r)^{-1} \cdot (r + 0) = z \to (1 - r)^{-1} \cdot (r + -1 \cdot 0) = z)$
\item $\varphi_{147} :\equiv \varphi_{18} \to (1 - r)^{-1} \cdot (r + -1 \cdot 0) = (1 - r)^{-1} \cdot (r + 0)$
\item $\varphi_{148} :\equiv \forall \beta\,(\mathrm{Conv}(\overline{r},r) \to \mathrm{Conv}(-1 \odot {(r)^{\bullet}} \circ \mathsf{s},\beta) \to \mathrm{Conv}(\overline{r} \oplus -1 \odot {(r)^{\bullet}} \circ \mathsf{s},r + \beta))$
\item $\varphi_{149} :\equiv \forall a\,(\mathrm{Seq}(a) \to \mathrm{Seq}(\Sigma a) \land {\textstyle\sum_{k=1}^{1_{\omega}} a_{k}} = a_{1_{\omega}} \land \forall n \in \mathbb{N}\,({\textstyle\sum_{k=1}^{S(n)} a_{k}} = {\textstyle\sum_{k=1}^{n} a_{k}} + a_{S(n)}))$
\item $\varphi_{150} :\equiv \forall k \in \mathbb{N}\,({(r)^{\bullet}} \circ \mathsf{s}(k) = {{r}^{\mathsf{s}(k)}}) \to n \in \mathbb{N} \to {(r)^{\bullet}} \circ \mathsf{s}(n) = {{r}^{\mathsf{s}(n)}}$
\item $\varphi_{151} :\equiv (1 - r)^{-1} \odot \overline{r} \oplus -1 \odot {(r)^{\bullet}} \circ \mathsf{s}(n) = (1 - r)^{-1} \cdot (r + -1 \cdot {(r)^{\bullet}} \circ \mathsf{s}(n))$
\item $\varphi_{152} :\equiv \forall c \in \mathbb{R}\,(\mathrm{Seq}(\overline{c}) \land \forall n \in \mathbb{N}\,(\overline{c}(n) = c)) \to r \in \mathbb{R} \to \mathrm{Seq}(\overline{r}) \land \forall n \in \mathbb{N}\,(\overline{r}(n) = r)$
\item $\varphi_{153} :\equiv (1 - r)^{-1} \odot \overline{r} \oplus -1 \odot {(r)^{\bullet}} \circ \mathsf{s}(n) = (1 - r)^{-1} \cdot (r + -1 \odot {(r)^{\bullet}} \circ \mathsf{s}(n))$
\item $\varphi_{154} :\equiv (1 - r)^{-1} \cdot (\overline{r}(n) + -1 \odot {(r)^{\bullet}} \circ \mathsf{s}(n)) = (1 - r)^{-1} \cdot (r + -1 \odot {(r)^{\bullet}} \circ \mathsf{s}(n))$
\item $\varphi_{155} :\equiv (1 - r)^{-1} \cdot (\overline{r}(n) + -1 \odot {(r)^{\bullet}} \circ \mathsf{s}(n)) = (1 - r)^{-1} \cdot (w + -1 \odot {(r)^{\bullet}} \circ \mathsf{s}(n))$
\item $\varphi_{156} :\equiv (1 - r)^{-1} \cdot (\overline{r}(n) + -1 \odot {(r)^{\bullet}} \circ \mathsf{s}(n)) = (1 - r)^{-1} \cdot (z + -1 \odot {(r)^{\bullet}} \circ \mathsf{s}(n))$
\item $\varphi_{157} :\equiv \forall a\,((1 - r)^{-1} \in \mathbb{R} \to \mathrm{Seq}(a) \to \mathrm{Seq}((1 - r)^{-1} \odot a) \land \forall n \in \mathbb{N}\,((1 - r)^{-1} \odot a(n) = (1 - r)^{-1} \cdot a_{n}))$
\item $\varphi_{158} :\equiv (r \in \mathbb{R} \to \lvert r \rvert < 1 \to \mathrm{Conv}({(r)^{\bullet}} \circ \mathsf{s},0)) \to \varphi_{39}$
\item $\varphi_{159} :\equiv \forall n \in \mathbb{N}\,((1 - r) \cdot {\textstyle\sum_{k=1}^{n} {{r}^{k}}} = r - {{r}^{S(n)}}) \to n \in \mathbb{N} \to (1 - r) \cdot {\textstyle\sum_{k=1}^{n} {{r}^{k}}} = r - {{r}^{S(n)}}$
\item $\varphi_{160} :\equiv {(r)^{\bullet}} \circ \mathsf{s}(n) = {{r}^{\mathsf{s}(n)}} \to {{r}^{\mathsf{s}(n)}} = {{r}^{S(n)}} \to {(r)^{\bullet}} \circ \mathsf{s}(n) = {{r}^{S(n)}}$
\item $\varphi_{161} :\equiv \varphi_{15} \to r \in \mathbb{R} \to \lvert r \rvert < 1 \to n \in \mathbb{N} \to n \in \mathbb{N} \to \mathsf{s}(n) = S(n)$
\item $\varphi_{162} :\equiv (\forall n \in \mathbb{N}\,(\overline{r}(n) = r) \to n \in \mathbb{N} \to \overline{r}(n) = r) \to r \in \mathbb{R} \to \lvert r \rvert < 1 \to n \in \mathbb{N} \to n \in \mathbb{N} \to \overline{r}(n) = r$
\item $\varphi_{163} :\equiv (1 - r)^{-1} \cdot \overline{r} \oplus -1 \odot {(r)^{\bullet}} \circ \mathsf{s}(n) = (1 - r)^{-1} \cdot \overline{r} \oplus -1 \odot {(r)^{\bullet}} \circ \mathsf{s}(n)$
\item $\varphi_{164} :\equiv \mathrm{Seq}(\Sigma {(r)^{\bullet}}) \land {\textstyle\sum_{k=1}^{1_{\omega}} {{r}^{k}}} = {{r}^{1_{\omega}}} \land \forall n \in \mathbb{N}\,({\textstyle\sum_{k=1}^{S(n)} {{r}^{k}}} = {\textstyle\sum_{k=1}^{n} {{r}^{k}}} + {{r}^{S(n)}})$
\item $\varphi_{165} :\equiv n \in \mathbb{N} \to \mathrm{Seq}(-1 \odot {(r)^{\bullet}} \circ \mathsf{s}) \land \forall n \in \mathbb{N}\,(\varphi_{32})$
\item $\varphi_{166} :\equiv \forall y\,(1 \in \mathbb{R} \to y \in \mathbb{R} \to 1 \cdot y = y \cdot 1) \to 1 \in \mathbb{R} \to {\textstyle\sum_{k=1}^{n} {{r}^{k}}} \in \mathbb{R} \to 1 \cdot {\textstyle\sum_{k=1}^{n} {{r}^{k}}} = {\textstyle\sum_{k=1}^{n} {{r}^{k}}} \cdot 1$
\item $\varphi_{167} :\equiv (1 - r)^{-1} \odot \overline{r} \oplus -1 \odot {(r)^{\bullet}} \circ \mathsf{s}(n) = (1 - r)^{-1} \cdot \overline{r} \oplus -1 \odot {(r)^{\bullet}} \circ \mathsf{s}(n)$
\item $\varphi_{168} :\equiv (1 - r)^{-1} \cdot \overline{r} \oplus -1 \odot {(r)^{\bullet}} \circ \mathsf{s}(n) = (1 - r)^{-1} \cdot (\overline{r}(n) + -1 \odot {(r)^{\bullet}} \circ \mathsf{s}(n))$
\item $\varphi_{169} :\equiv (1 - r)^{-1} \cdot (r + -1 \cdot 0) = r \cdot (1 - r)^{-1} \to \varphi_{3} \to \mathrm{Conv}(\Sigma {(r)^{\bullet}},r \cdot (1 - r)^{-1})$
\item $\varphi_{170} :\equiv (1 - r)^{-1} \cdot (1 - r) \cdot {\textstyle\sum_{k=1}^{n} {{r}^{k}}} = z \cdot {\textstyle\sum_{k=1}^{n} {{r}^{k}}} \to (1 - r)^{-1} \cdot (1 - r) \cdot {\textstyle\sum_{k=1}^{n} {{r}^{k}}} = w \cdot {\textstyle\sum_{k=1}^{n} {{r}^{k}}}$
\item $\varphi_{171} :\equiv \forall w\,(-1 \cdot 0 = w \to \varphi_{18} \to (1 - r)^{-1} \cdot (r + -1 \cdot 0) = (1 - r)^{-1} \cdot (r + w))$
\item $\varphi_{172} :\equiv (1 - r)^{-1} \cdot (\overline{r}(n) + -1 \odot {(r)^{\bullet}} \circ \mathsf{s}(n)) = (1 - r)^{-1} \cdot (\overline{r}(n) + -1 \odot {(r)^{\bullet}} \circ \mathsf{s}(n))$
\item $\varphi_{173} :\equiv (r \in \mathbb{R} \to \varphi_{37}) \to r \in \mathbb{R} \to \lvert r \rvert < 1 \to n \in \mathbb{N} \to \mathrm{Seq}(\overline{r})$
\item $\varphi_{174} :\equiv (1 - r)^{-1} \odot \overline{r} \oplus -1 \odot {(r)^{\bullet}} \circ \mathsf{s}(n) = (1 - r)^{-1} \cdot (\overline{r}(n) + -1 \odot {(r)^{\bullet}} \circ \mathsf{s}(n))$
\item $\varphi_{175} :\equiv (r - {{r}^{S(n)}} = r + -{{r}^{S(n)}} \to r + -{{r}^{S(n)}} = r - {{r}^{S(n)}}) \to r + -{{r}^{S(n)}} = r - {{r}^{S(n)}}$
\item $\varphi_{176} :\equiv (1 - r)^{-1} \cdot (r + -1 \cdot 0) = (1 - r)^{-1} \cdot (r + 0) \to \varphi_{88}$
\item $\varphi_{177} :\equiv (r \leq \lvert r \rvert \to \lvert r \rvert < 1 \to r < 1) \to (r \in \mathbb{R} \to \lvert r \rvert < 1 \to r \leq \lvert r \rvert) \to \varphi_{41}$
\item $\varphi_{178} :\equiv \mathrm{Seq}({(r)^{\bullet}} \circ \mathsf{s}) \to \mathrm{Seq}(-1 \odot {(r)^{\bullet}} \circ \mathsf{s}) \land \forall n \in \mathbb{N}\,(\varphi_{32})$
\item $\varphi_{179} :\equiv \forall y\,(r - {{r}^{S(n)}} = y \to y = r - {{r}^{S(n)}}) \to r - {{r}^{S(n)}} = r + -{{r}^{S(n)}} \to r + -{{r}^{S(n)}} = r - {{r}^{S(n)}}$
\item $\varphi_{180} :\equiv (r \in \mathbb{R} \to \lvert r \rvert < 1 \to \varphi_{58}) \to r \in \mathbb{R} \to \lvert r \rvert < 1 \to \mathrm{Seq}({(r)^{\bullet}})$
\item $\varphi_{181} :\equiv (r \in \mathbb{R} \to \lvert r \rvert < 1 \to n \in \mathbb{N} \to {\textstyle\sum_{k=1}^{n} {{r}^{k}}} \in \mathbb{R}) \to r \in \mathbb{R} \to \lvert r \rvert < 1 \to n \in \mathbb{N} \to {\textstyle\sum_{k=1}^{n} {{r}^{k}}} \cdot 1 = {\textstyle\sum_{k=1}^{n} {{r}^{k}}}$
\item $\varphi_{182} :\equiv (r \in \mathbb{R} \to \lvert r \rvert < 1 \to n \in \mathbb{N} \to {{r}^{S(n)}} \in \mathbb{R}) \to r \in \mathbb{R} \to \lvert r \rvert < 1 \to n \in \mathbb{N} \to -1 \cdot {{r}^{S(n)}} = -{{r}^{S(n)}}$
\item $\varphi_{183} :\equiv \forall w\,(r - {{r}^{S(n)}} = w \to \varphi_{35} \to (1 - r)^{-1} \cdot (r - {{r}^{S(n)}}) = (1 - r)^{-1} \cdot w)$
\item $\varphi_{184} :\equiv (r \in \mathbb{R} \to \lvert r \rvert < 1 \to n \in \mathbb{N} \to {\textstyle\sum_{k=1}^{n} {{r}^{k}}} \in \mathbb{R}) \to r \in \mathbb{R} \to \lvert r \rvert < 1 \to n \in \mathbb{N} \to 1 \cdot {\textstyle\sum_{k=1}^{n} {{r}^{k}}} = {\textstyle\sum_{k=1}^{n} {{r}^{k}}} \cdot 1$
\item $\varphi_{185} :\equiv \forall y\,\forall z\,((1 - r)^{-1} \odot \overline{r} \oplus -1 \odot {(r)^{\bullet}} \circ \mathsf{s}(n) = y \to y = z \to (1 - r)^{-1} \odot \overline{r} \oplus -1 \odot {(r)^{\bullet}} \circ \mathsf{s}(n) = z)$
\item $\varphi_{186} :\equiv (r \in \mathbb{R} \to \lvert r \rvert < 1 \to n \in \mathbb{N} \to 1 - r \in \mathbb{R}) \to (r \in \mathbb{R} \to \lvert r \rvert < 1 \to 1 - r \neq 0) \to \varphi_{14}$
\item $\varphi_{187} :\equiv (r \in \mathbb{R} \to \varphi_{37}) \to r \in \mathbb{R} \to \lvert r \rvert < 1 \to n \in \mathbb{N} \to \forall n \in \mathbb{N}\,(\overline{r}(n) = r)$
\item $\varphi_{188} :\equiv \forall r \in \mathbb{R}\,\forall n \in \mathbb{N}\,((1 - r) \cdot {\textstyle\sum_{k=1}^{n} {{r}^{k}}} = r - {{r}^{S(n)}}) \to r \in \mathbb{R} \to \forall n \in \mathbb{N}\,((1 - r) \cdot {\textstyle\sum_{k=1}^{n} {{r}^{k}}} = r - {{r}^{S(n)}})$
\item $\varphi_{189} :\equiv (r \in \mathbb{R} \to \lvert r \rvert < 1 \to n \in \mathbb{N} \to n \in \mathbb{N}) \to r \in \mathbb{R} \to \lvert r \rvert < 1 \to n \in \mathbb{N} \to \varphi_{32}$
\item $\varphi_{190} :\equiv \forall \alpha\,((1 - r)^{-1} \in \mathbb{R} \to \mathrm{Conv}(\overline{r} \oplus -1 \odot {(r)^{\bullet}} \circ \mathsf{s},\alpha) \to \mathrm{Conv}((1 - r)^{-1} \odot \overline{r} \oplus -1 \odot {(r)^{\bullet}} \circ \mathsf{s},(1 - r)^{-1} \cdot \alpha))$
\item $\varphi_{191} :\equiv (r \in \mathbb{R} \to \lvert r \rvert < 1 \to n \in \mathbb{N} \to \mathrm{Seq}({(r)^{\bullet}})) \to \varphi_{143}$
\item $\varphi_{192} :\equiv \forall y\,((1 - r) \cdot {\textstyle\sum_{k=1}^{n} {{r}^{k}}} = y \to y = (1 - r) \cdot {\textstyle\sum_{k=1}^{n} {{r}^{k}}}) \to \varphi_{115}$
\item $\varphi_{193} :\equiv \mathrm{Conv}(\overline{r} \oplus -1 \odot {(r)^{\bullet}} \circ \mathsf{s},r + -1 \cdot 0) \to \varphi_{114}$
\item $\varphi_{194} :\equiv (r \in \mathbb{R} \to \lvert r \rvert < 1 \to n \in \mathbb{N} \to r \in \mathbb{R}) \to r \in \mathbb{R} \to \lvert r \rvert < 1 \to n \in \mathbb{N} \to \varphi_{58}$
\item $\varphi_{195} :\equiv (r \in \mathbb{R} \to \lvert r \rvert < 1 \to n \in \mathbb{N} \to \varphi_{58}) \to r \in \mathbb{R} \to \lvert r \rvert < 1 \to n \in \mathbb{N} \to \mathrm{Seq}({(r)^{\bullet}})$
\item $\varphi_{196} :\equiv \forall w\,(r + -{{r}^{S(n)}} = w \to \varphi_{55} \to (1 - r)^{-1} \cdot (r + -{{r}^{S(n)}}) = (1 - r)^{-1} \cdot w)$
\item $\varphi_{197} :\equiv (r \in \mathbb{R} \to \lvert r \rvert < 1 \to n \in \mathbb{N} \to \mathrm{Seq}(\Sigma {(r)^{\bullet}})) \to \varphi_{142}$
\item $\varphi_{198} :\equiv (r \in \mathbb{R} \to 1 \in \mathbb{R} \to r < 1 \to 0 < 1 - r) \to (r \in \mathbb{R} \to \lvert r \rvert < 1 \to r \in \mathbb{R}) \to 1 \in \mathbb{R} \to \varphi_{42}$
\item $\varphi_{199} :\equiv (r \in \mathbb{R} \to \lvert r \rvert < 1 \to n \in \mathbb{N} \to n \in \mathbb{N} \to \mathsf{s}(n) = S(n)) \to \varphi_{123}$
\item $\varphi_{200} :\equiv (r \in \mathbb{R} \to \lvert r \rvert < 1 \to n \in \mathbb{N} \to \mathrm{NSeq}(\mathsf{s})) \to r \in \mathbb{R} \to \lvert r \rvert < 1 \to n \in \mathbb{N} \to \varphi_{71}$
\item $\varphi_{201} :\equiv (r \in \mathbb{R} \to \lvert r \rvert < 1 \to n \in \mathbb{N} \to n \in \mathbb{N} \to \overline{r}(n) = r) \to \varphi_{124}$
\item $\varphi_{202} :\equiv (1 - r)^{-1} \cdot (1 - r) = 1 \to \varphi_{77} \to (1 - r)^{-1} \cdot (1 - r) \cdot {\textstyle\sum_{k=1}^{n} {{r}^{k}}} = 1 \cdot {\textstyle\sum_{k=1}^{n} {{r}^{k}}}$
\item $\varphi_{203} :\equiv (1 - r \in \mathbb{R} \to 1 - r \neq 0 \to (1 - r)^{-1} \in \mathbb{R}) \to (r \in \mathbb{R} \to \lvert r \rvert < 1 \to 1 - r \in \mathbb{R}) \to \varphi_{74}$
\item $\varphi_{204} :\equiv (\mathsf{s}(n) = S(n) \to {{r}^{\mathsf{s}(n)}} = {{r}^{\mathsf{s}(n)}} \to {{r}^{\mathsf{s}(n)}} = {{r}^{S(n)}}) \to r \in \mathbb{R} \to \lvert r \rvert < 1 \to n \in \mathbb{N} \to {{r}^{\mathsf{s}(n)}} = {{r}^{S(n)}}$
\item $\varphi_{205} :\equiv n \in \mathbb{N} \to \mathrm{Seq}(\overline{r} \oplus -1 \odot {(r)^{\bullet}} \circ \mathsf{s}) \land \forall n \in \mathbb{N}\,(\varphi_{125})$
\item $\varphi_{206} :\equiv \forall w\,((1 - r)^{-1} \cdot (1 - r) = w \to \varphi_{77} \to (1 - r)^{-1} \cdot (1 - r) \cdot {\textstyle\sum_{k=1}^{n} {{r}^{k}}} = w \cdot {\textstyle\sum_{k=1}^{n} {{r}^{k}}})$
\item $\varphi_{207} :\equiv (r \in \mathbb{R} \land \lvert r \rvert < 1 \to \forall n\,(r \in \mathbb{R} \land \lvert r \rvert < 1)) \to r \in \mathbb{R} \to \lvert r \rvert < 1 \to \forall n \in \mathbb{N}\,(\varphi_{85})$
\item $\varphi_{208} :\equiv \forall z\,\forall w\,(z = w \to {{r}^{\mathsf{s}(n)}} = {{r}^{z}} \to {{r}^{\mathsf{s}(n)}} = {{r}^{w}}) \to \forall w\,(\mathsf{s}(n) = w \to {{r}^{\mathsf{s}(n)}} = {{r}^{\mathsf{s}(n)}} \to {{r}^{\mathsf{s}(n)}} = {{r}^{w}})$
\item $\varphi_{209} :\equiv (r \in \mathbb{R} \to \lvert r \rvert < 1 \to n \in \mathbb{N} \to {\textstyle\sum_{k=1}^{n} {{r}^{k}}} \in \mathbb{R}) \to r \in \mathbb{R} \to \lvert r \rvert < 1 \to n \in \mathbb{N} \to \varphi_{79}$
\item $\varphi_{210} :\equiv (r \in \mathbb{R} \to \lvert r \rvert < 1 \to n \in \mathbb{N} \to n \in \mathbb{N}) \to r \in \mathbb{R} \to \lvert r \rvert < 1 \to n \in \mathbb{N} \to \varphi_{125}$
\item $\varphi_{211} :\equiv \forall b\,\forall \alpha\,(\mathrm{Seq}(b) \to \forall n \in \mathbb{N}\,((1 - r)^{-1} \odot \overline{r} \oplus -1 \odot {(r)^{\bullet}} \circ \mathsf{s}(n) = b_{n}) \to \mathrm{Conv}((1 - r)^{-1} \odot \overline{r} \oplus -1 \odot {(r)^{\bullet}} \circ \mathsf{s},\alpha) \to \mathrm{Conv}(b,\alpha))$
\item $\varphi_{212} :\equiv \forall z\,(\varphi_{96} \to {\textstyle\sum_{k=1}^{n} {{r}^{k}}} \cdot 1 = z \to (1 - r)^{-1} \odot \overline{r} \oplus -1 \odot {(r)^{\bullet}} \circ \mathsf{s}(n) = z)$
\item $\varphi_{213} :\equiv \forall z\,(\varphi_{97} \to 1 \cdot {\textstyle\sum_{k=1}^{n} {{r}^{k}}} = z \to (1 - r)^{-1} \odot \overline{r} \oplus -1 \odot {(r)^{\bullet}} \circ \mathsf{s}(n) = z)$
\item $\varphi_{214} :\equiv (r \in \mathbb{R} \to \mathrm{Seq}(\overline{r}) \land \forall n \in \mathbb{N}\,(\overline{r}(n) = r)) \to (r \in \mathbb{R} \to \lvert r \rvert < 1 \to n \in \mathbb{N} \to r \in \mathbb{R}) \to r \in \mathbb{R} \to \varphi_{37}$
\item $\varphi_{215} :\equiv (r \in \mathbb{R} \to \lvert r \rvert < 1 \to \mathrm{Seq}({(r)^{\bullet}})) \to r \in \mathbb{R} \to \lvert r \rvert < 1 \to \varphi_{164}$
\item $\varphi_{216} :\equiv (r \in \mathbb{R} \to \lvert r \rvert < 1 \to n \in \mathbb{N} \to \forall n \in \mathbb{N}\,(\overline{r}(n) = r)) \to \varphi_{162}$
\item $\varphi_{217} :\equiv (r \in \mathbb{R} \to \lvert r \rvert < 1 \to \varphi_{164}) \to r \in \mathbb{R} \to \lvert r \rvert < 1 \to \mathrm{Seq}(\Sigma {(r)^{\bullet}})$
\item $\varphi_{218} :\equiv (1 - r \in \mathbb{R} \to 1 - r \neq 0 \to (1 - r)^{-1} \in \mathbb{R}) \to (r \in \mathbb{R} \to \lvert r \rvert < 1 \to n \in \mathbb{N} \to 1 - r \in \mathbb{R}) \to \varphi_{107}$
\item $\varphi_{219} :\equiv (1 - r)^{-1} \cdot (r + 0) = (1 - r)^{-1} \cdot (r + 0) \to (r + 0 = r \to \varphi_{66}) \to \varphi_{22}$
\item $\varphi_{220} :\equiv \forall w\,(\mathsf{s}(n) = w \to {{r}^{\mathsf{s}(n)}} = {{r}^{\mathsf{s}(n)}} \to {{r}^{\mathsf{s}(n)}} = {{r}^{w}}) \to \mathsf{s}(n) = S(n) \to {{r}^{\mathsf{s}(n)}} = {{r}^{\mathsf{s}(n)}} \to {{r}^{\mathsf{s}(n)}} = {{r}^{S(n)}}$
\item $\varphi_{221} :\equiv \mathrm{Seq}(-1 \odot {(r)^{\bullet}} \circ \mathsf{s}) \to \mathrm{Seq}(\overline{r} \oplus -1 \odot {(r)^{\bullet}} \circ \mathsf{s}) \land \forall n \in \mathbb{N}\,(\varphi_{125})$
\item $\varphi_{222} :\equiv ((1 - r)^{-1} \cdot (r + -1 \cdot 0) = (1 - r)^{-1} \cdot r \to \varphi_{65}) \to r \in \mathbb{R} \to \varphi_{17}$
\item $\varphi_{223} :\equiv (1 - r \in \mathbb{R} \to 1 - r \neq 0 \to (1 - r)^{-1} \cdot (1 - r) = 1) \to \varphi_{186}$
\item $\varphi_{224} :\equiv (r \in \mathbb{R} \to \lvert r \rvert < 1 \to \mathrm{Conv}(\overline{r},r)) \to \varphi_{39} \to r \in \mathbb{R} \to \varphi_{57}$
\item $\varphi_{225} :\equiv \mathrm{Seq}((1 - r)^{-1} \odot \overline{r} \oplus -1 \odot {(r)^{\bullet}} \circ \mathsf{s}) \land \forall n \in \mathbb{N}\,(\varphi_{167})$
\item $\varphi_{226} :\equiv \forall z\,(\varphi_{121} \to (1 - r)^{-1} \cdot (r - {{r}^{S(n)}}) = z \to (1 - r)^{-1} \odot \overline{r} \oplus -1 \odot {(r)^{\bullet}} \circ \mathsf{s}(n) = z)$
\item $\varphi_{227} :\equiv \forall \alpha\,(\mathrm{Seq}(\Sigma {(r)^{\bullet}}) \to \forall n \in \mathbb{N}\,(\varphi_{85}) \to \mathrm{Conv}((1 - r)^{-1} \odot \overline{r} \oplus -1 \odot {(r)^{\bullet}} \circ \mathsf{s},\alpha) \to \mathrm{Conv}(\Sigma {(r)^{\bullet}},\alpha))$
\item $\varphi_{228} :\equiv (r \in \mathbb{R} \to \lvert r \rvert < 1 \to n \in \mathbb{N} \to \mathrm{Seq}({(r)^{\bullet}})) \to r \in \mathbb{R} \to \lvert r \rvert < 1 \to n \in \mathbb{N} \to \varphi_{164}$
\item $\varphi_{229} :\equiv (\mathrm{Seq}({(r)^{\bullet}}) \to S(n) \in \mathbb{N} \to {{r}^{S(n)}} \in \mathbb{R}) \to \varphi_{191}$
\item $\varphi_{230} :\equiv (r \in \mathbb{R} \to \lvert r \rvert < 1 \to n \in \mathbb{N} \to n \in \mathbb{N}) \to r \in \mathbb{R} \to \lvert r \rvert < 1 \to n \in \mathbb{N} \to \varphi_{167}$
\item $\varphi_{231} :\equiv (r \in \mathbb{R} \to \lvert r \rvert < 1 \to n \in \mathbb{N} \to \mathrm{Seq}({(r)^{\bullet}})) \to \varphi_{200}$
\item $\varphi_{232} :\equiv (r \in \mathbb{R} \to \lvert r \rvert < 1 \to n \in \mathbb{N} \to \varphi_{164}) \to r \in \mathbb{R} \to \lvert r \rvert < 1 \to n \in \mathbb{N} \to \mathrm{Seq}(\Sigma {(r)^{\bullet}})$
\item $\varphi_{233} :\equiv \forall z\,(\varphi_{135} \to (1 - r)^{-1} \cdot (r + -{{r}^{S(n)}}) = z \to (1 - r)^{-1} \odot \overline{r} \oplus -1 \odot {(r)^{\bullet}} \circ \mathsf{s}(n) = z)$
\item $\varphi_{234} :\equiv (r \in \mathbb{R} \to \forall n \in \mathbb{N}\,((1 - r) \cdot {\textstyle\sum_{k=1}^{n} {{r}^{k}}} = r - {{r}^{S(n)}})) \to (r \in \mathbb{R} \to \lvert r \rvert < 1 \to n \in \mathbb{N} \to r \in \mathbb{R}) \to r \in \mathbb{R} \to \varphi_{91}$
\item $\varphi_{235} :\equiv (\mathrm{Seq}(\Sigma {(r)^{\bullet}}) \to n \in \mathbb{N} \to {\textstyle\sum_{k=1}^{n} {{r}^{k}}} \in \mathbb{R}) \to \varphi_{197}$
\item $\varphi_{236} :\equiv (r \in \mathbb{R} \to \varphi_{64}) \to (r \in \mathbb{R} \to \lvert r \rvert < 1 \to n \in \mathbb{N} \to n \in \mathbb{N}) \to \varphi_{67}$
\item $\varphi_{237} :\equiv \forall z\,(\varphi_{138} \to (1 - r)^{-1} \cdot (1 - r) \cdot {\textstyle\sum_{k=1}^{n} {{r}^{k}}} = z \to (1 - r)^{-1} \odot \overline{r} \oplus -1 \odot {(r)^{\bullet}} \circ \mathsf{s}(n) = z)$
\item $\varphi_{238} :\equiv \forall z\,(\varphi_{139} \to (1 - r)^{-1} \cdot ((1 - r) \cdot {\textstyle\sum_{k=1}^{n} {{r}^{k}}}) = z \to (1 - r)^{-1} \odot \overline{r} \oplus -1 \odot {(r)^{\bullet}} \circ \mathsf{s}(n) = z)$
\item $\varphi_{239} :\equiv (r \in \mathbb{R} \to \varphi_{80}) \to (r \in \mathbb{R} \to \lvert r \rvert < 1 \to n \in \mathbb{N} \to n \in \mathbb{N}) \to \varphi_{83}$
\item $\varphi_{240} :\equiv \forall z\,(\varphi_{140} \to (1 - r)^{-1} \cdot (r + -1 \cdot {{r}^{S(n)}}) = z \to (1 - r)^{-1} \odot \overline{r} \oplus -1 \odot {(r)^{\bullet}} \circ \mathsf{s}(n) = z)$
\item $\varphi_{241} :\equiv (r \in \mathbb{R} \to \varphi_{58}) \to (r \in \mathbb{R} \to \lvert r \rvert < 1 \to r \in \mathbb{R}) \to r \in \mathbb{R} \to \lvert r \rvert < 1 \to \varphi_{58}$
\item $\varphi_{242} :\equiv (r \in \mathbb{R} \to \lvert r \rvert < 1 \to n \in \mathbb{N} \to \varphi_{85}) \to r \in \mathbb{R} \land \lvert r \rvert < 1 \to n \in \mathbb{N} \to \varphi_{85}$
\item $\varphi_{243} :\equiv \forall z\,(\varphi_{151} \to (1 - r)^{-1} \cdot (r + -1 \cdot {(r)^{\bullet}} \circ \mathsf{s}(n)) = z \to (1 - r)^{-1} \odot \overline{r} \oplus -1 \odot {(r)^{\bullet}} \circ \mathsf{s}(n) = z)$
\item $\varphi_{244} :\equiv (1 \in \mathbb{R} \to {\textstyle\sum_{k=1}^{n} {{r}^{k}}} \in \mathbb{R} \to 1 \cdot {\textstyle\sum_{k=1}^{n} {{r}^{k}}} = {\textstyle\sum_{k=1}^{n} {{r}^{k}}} \cdot 1) \to 1 \in \mathbb{R} \to \varphi_{184}$
\item $\varphi_{245} :\equiv \forall z\,(\varphi_{153} \to (1 - r)^{-1} \cdot (r + -1 \odot {(r)^{\bullet}} \circ \mathsf{s}(n)) = z \to (1 - r)^{-1} \odot \overline{r} \oplus -1 \odot {(r)^{\bullet}} \circ \mathsf{s}(n) = z)$
\item $\varphi_{246} :\equiv \forall n\,(r \in \mathbb{R} \land \lvert r \rvert < 1 \to n \in \mathbb{N} \to \varphi_{85}) \to \forall n\,(r \in \mathbb{R} \land \lvert r \rvert < 1) \to \forall n \in \mathbb{N}\,(\varphi_{85})$
\item $\varphi_{247} :\equiv (r \in \mathbb{R} \to \lvert r \rvert < 1 \to (1 - r)^{-1} \in \mathbb{R}) \to (r \in \mathbb{R} \to \varphi_{57}) \to r \in \mathbb{R} \to \lvert r \rvert < 1 \to \varphi_{114}$
\item $\varphi_{248} :\equiv (r \in \mathbb{R} \to \lvert r \rvert < 1 \to (1 - r)^{-1} \cdot r = r \cdot (1 - r)^{-1}) \to \varphi_{222}$
\item $\varphi_{249} :\equiv (r \in \mathbb{R} \to \lvert r \rvert < 1 \to n \in \mathbb{N} \to -1 \in \mathbb{R}) \to \varphi_{10} \to r \in \mathbb{R} \to \lvert r \rvert < 1 \to \varphi_{165}$
\item $\varphi_{250} :\equiv \forall z\,(\varphi_{167} \to (1 - r)^{-1} \cdot \overline{r} \oplus -1 \odot {(r)^{\bullet}} \circ \mathsf{s}(n) = z \to (1 - r)^{-1} \odot \overline{r} \oplus -1 \odot {(r)^{\bullet}} \circ \mathsf{s}(n) = z)$
\item $\varphi_{251} :\equiv (r \in \mathbb{R} \to \lvert r \rvert < 1 \to \varphi_{165}) \to r \in \mathbb{R} \to \lvert r \rvert < 1 \to n \in \mathbb{N} \to \forall n \in \mathbb{N}\,(\varphi_{32})$
\item $\varphi_{252} :\equiv (r \in \mathbb{R} \to \lvert r \rvert < 1 \to n \in \mathbb{N} \to {{r}^{\mathsf{s}(n)}} = {{r}^{S(n)}}) \to \varphi_{160} \to \varphi_{49}$
\item $\varphi_{253} :\equiv \forall z\,(\varphi_{174} \to (1 - r)^{-1} \cdot (\overline{r}(n) + -1 \odot {(r)^{\bullet}} \circ \mathsf{s}(n)) = z \to (1 - r)^{-1} \odot \overline{r} \oplus -1 \odot {(r)^{\bullet}} \circ \mathsf{s}(n) = z)$
\item $\varphi_{254} :\equiv (r \in \mathbb{R} \to \lvert r \rvert < 1 \to n \in \mathbb{N} \to \mathsf{s}(n) = S(n)) \to {{r}^{\mathsf{s}(n)}} = {{r}^{\mathsf{s}(n)}} \to \varphi_{204}$
\item $\varphi_{255} :\equiv (\forall n \in \mathbb{N}\,(\varphi_{32}) \to n \in \mathbb{N} \to \varphi_{32}) \to r \in \mathbb{R} \to \lvert r \rvert < 1 \to n \in \mathbb{N} \to n \in \mathbb{N} \to \varphi_{32}$
\item $\varphi_{256} :\equiv (1 - r)^{-1} \in \mathbb{R} \to \mathrm{Seq}(\overline{r} \oplus -1 \odot {(r)^{\bullet}} \circ \mathsf{s}) \to \varphi_{225}$
\item $\varphi_{257} :\equiv (r \in \mathbb{R} \to \lvert r \rvert < 1 \to n \in \mathbb{N} \to (1 - r)^{-1} \in \mathbb{R}) \to (r \in \mathbb{R} \to \lvert r \rvert < 1 \to n \in \mathbb{N} \to 1 - r \in \mathbb{R}) \to \varphi_{209}$
\item $\varphi_{258} :\equiv (r - {{r}^{S(n)}} = (1 - r) \cdot {\textstyle\sum_{k=1}^{n} {{r}^{k}}} \to \varphi_{35} \to \varphi_{52}) \to r \in \mathbb{R} \to \lvert r \rvert < 1 \to n \in \mathbb{N} \to \varphi_{52}$
\item $\varphi_{259} :\equiv (r + -{{r}^{S(n)}} = r - {{r}^{S(n)}} \to \varphi_{55} \to \varphi_{46}) \to r \in \mathbb{R} \to \lvert r \rvert < 1 \to n \in \mathbb{N} \to \varphi_{46}$
\item $\varphi_{260} :\equiv (r \in \mathbb{R} \to \lvert r \rvert < 1 \to -1 \cdot 0 = 0) \to \varphi_{18} \to (-1 \cdot 0 = 0 \to \varphi_{147}) \to r \in \mathbb{R} \to \varphi_{26}$
\item $\varphi_{261} :\equiv (r \in \mathbb{R} \to \lvert r \rvert < 1 \to \forall n \in \mathbb{N}\,(\varphi_{85})) \to (r \in \mathbb{R} \to \lvert r \rvert < 1 \to \varphi_{114}) \to r \in \mathbb{R} \to \lvert r \rvert < 1 \to \varphi_{3}$
\item $\varphi_{262} :\equiv (r \in \mathbb{R} \to \lvert r \rvert < 1 \to n \in \mathbb{N} \to \mathrm{Seq}(\overline{r})) \to \varphi_{38} \to r \in \mathbb{R} \to \lvert r \rvert < 1 \to \varphi_{205}$
\item $\varphi_{263} :\equiv \forall z\,\forall w\,(z = w \to \varphi_{25} \to \varphi_{24}) \to \forall w\,(-1 \cdot {{r}^{S(n)}} = w \to \varphi_{94} \to \varphi_{24})$
\item $\varphi_{264} :\equiv (-1 \cdot {{r}^{S(n)}} = -{{r}^{S(n)}} \to \varphi_{94} \to \varphi_{76}) \to r \in \mathbb{R} \to \lvert r \rvert < 1 \to n \in \mathbb{N} \to \varphi_{76}$
\item $\varphi_{265} :\equiv (r \in \mathbb{R} \to \varphi_{17}) \to (r \in \mathbb{R} \to \lvert r \rvert < 1 \to \varphi_{3}) \to \varphi_{169} \to r \in \mathbb{R} \to \varphi_{2}$
\item $\varphi_{266} :\equiv (\varphi_{96} \to {\textstyle\sum_{k=1}^{n} {{r}^{k}}} \cdot 1 = {\textstyle\sum_{k=1}^{n} {{r}^{k}}} \to \varphi_{85}) \to r \in \mathbb{R} \to \lvert r \rvert < 1 \to n \in \mathbb{N} \to \varphi_{85}$
\item $\varphi_{267} :\equiv (\varphi_{97} \to 1 \cdot {\textstyle\sum_{k=1}^{n} {{r}^{k}}} = {\textstyle\sum_{k=1}^{n} {{r}^{k}}} \cdot 1 \to \varphi_{96}) \to r \in \mathbb{R} \to \lvert r \rvert < 1 \to n \in \mathbb{N} \to \varphi_{96}$
\item $\varphi_{268} :\equiv (r \in \mathbb{R} \to \lvert r \rvert < 1 \to \varphi_{205}) \to r \in \mathbb{R} \to \lvert r \rvert < 1 \to n \in \mathbb{N} \to \forall n \in \mathbb{N}\,(\varphi_{125})$
\item $\varphi_{269} :\equiv ({(r)^{\bullet}} \circ \mathsf{s}(n) = {{r}^{S(n)}} \to \varphi_{126} \to \varphi_{112}) \to r \in \mathbb{R} \to \lvert r \rvert < 1 \to n \in \mathbb{N} \to \varphi_{112}$
\item $\varphi_{270} :\equiv (\forall n \in \mathbb{N}\,(\varphi_{125}) \to n \in \mathbb{N} \to \varphi_{125}) \to r \in \mathbb{R} \to \lvert r \rvert < 1 \to n \in \mathbb{N} \to n \in \mathbb{N} \to \varphi_{125}$
\item $\varphi_{271} :\equiv \forall z\,\forall w\,(z = w \to \varphi_{51} \to \varphi_{50}) \to \forall w\,(-1 \odot {(r)^{\bullet}} \circ \mathsf{s}(n) = w \to \varphi_{137} \to \varphi_{50})$
\item $\varphi_{272} :\equiv (r \in \mathbb{R} \to \lvert r \rvert < 1 \to n \in \mathbb{N} \to (1 - r)^{-1} \in \mathbb{R}) \to (r \in \mathbb{R} \to \varphi_{45}) \to r \in \mathbb{R} \to \lvert r \rvert < 1 \to n \in \mathbb{N} \to \varphi_{225}$
\item $\varphi_{273} :\equiv \forall z\,\forall w\,(z = w \to \varphi_{70} \to \varphi_{69}) \to \forall w\,({(r)^{\bullet}} \circ \mathsf{s}(n) = w \to \varphi_{126} \to \varphi_{69})$
\item $\varphi_{274} :\equiv (\varphi_{138} \to (1 - r)^{-1} \cdot (1 - r) \cdot {\textstyle\sum_{k=1}^{n} {{r}^{k}}} = 1 \cdot {\textstyle\sum_{k=1}^{n} {{r}^{k}}} \to \varphi_{97}) \to r \in \mathbb{R} \to \lvert r \rvert < 1 \to n \in \mathbb{N} \to \varphi_{97}$
\item $\varphi_{275} :\equiv \forall w\,(-1 \cdot {{r}^{S(n)}} = w \to \varphi_{94} \to \varphi_{24}) \to -1 \cdot {{r}^{S(n)}} = -{{r}^{S(n)}} \to \varphi_{94} \to \varphi_{76}$
\item $\varphi_{276} :\equiv (r \in \mathbb{R} \to \lvert r \rvert < 1 \to n \in \mathbb{N} \to \varphi_{79}) \to (\varphi_{79} \to \varphi_{78}) \to r \in \mathbb{R} \to \lvert r \rvert < 1 \to n \in \mathbb{N} \to \varphi_{78}$
\item $\varphi_{277} :\equiv (r \in \mathbb{R} \to \lvert r \rvert < 1 \to n \in \mathbb{N} \to {\textstyle\sum_{k=1}^{n} {{r}^{k}}} \cdot 1 = {\textstyle\sum_{k=1}^{n} {{r}^{k}}}) \to \varphi_{266}$
\item $\varphi_{278} :\equiv (r \in \mathbb{R} \to \lvert r \rvert < 1 \to n \in \mathbb{N} \to \varphi_{225}) \to r \in \mathbb{R} \to \lvert r \rvert < 1 \to n \in \mathbb{N} \to \forall n \in \mathbb{N}\,(\varphi_{167})$
\item $\varphi_{279} :\equiv (\overline{r}(n) = r \to \varphi_{172} \to \varphi_{154}) \to r \in \mathbb{R} \to \lvert r \rvert < 1 \to n \in \mathbb{N} \to \varphi_{154}$
\item $\varphi_{280} :\equiv \forall z\,\forall w\,(z = w \to \varphi_{62} \to \varphi_{61}) \to \forall w\,(\overline{r} \oplus -1 \odot {(r)^{\bullet}} \circ \mathsf{s}(n) = w \to \varphi_{163} \to \varphi_{61})$
\item $\varphi_{281} :\equiv (r \in \mathbb{R} \to \lvert r \rvert < 1 \to n \in \mathbb{N} \to 1 \cdot {\textstyle\sum_{k=1}^{n} {{r}^{k}}} = {\textstyle\sum_{k=1}^{n} {{r}^{k}}} \cdot 1) \to \varphi_{267}$
\item $\varphi_{282} :\equiv ((1 - r)^{-1} \in \mathbb{R} \to 1 - r \in \mathbb{R} \to {\textstyle\sum_{k=1}^{n} {{r}^{k}}} \in \mathbb{R} \to \varphi_{79}) \to \varphi_{257}$
\item $\varphi_{283} :\equiv \forall w\,({(r)^{\bullet}} \circ \mathsf{s}(n) = w \to \varphi_{126} \to \varphi_{69}) \to {(r)^{\bullet}} \circ \mathsf{s}(n) = {{r}^{S(n)}} \to \varphi_{126} \to \varphi_{112}$
\item $\varphi_{284} :\equiv (\forall n \in \mathbb{N}\,(\varphi_{167}) \to n \in \mathbb{N} \to \varphi_{167}) \to r \in \mathbb{R} \to \lvert r \rvert < 1 \to n \in \mathbb{N} \to n \in \mathbb{N} \to \varphi_{167}$
\item $\varphi_{285} :\equiv \varphi_{237} \to \varphi_{138} \to (1 - r)^{-1} \cdot (1 - r) \cdot {\textstyle\sum_{k=1}^{n} {{r}^{k}}} = 1 \cdot {\textstyle\sum_{k=1}^{n} {{r}^{k}}} \to \varphi_{97}$
\item $\varphi_{286} :\equiv (r \in \mathbb{R} \to \lvert r \rvert < 1 \to n \in \mathbb{N} \to -1 \cdot {{r}^{S(n)}} = -{{r}^{S(n)}}) \to \varphi_{94} \to \varphi_{264}$
\item $\varphi_{287} :\equiv \forall w\,(-1 \odot {(r)^{\bullet}} \circ \mathsf{s}(n) = w \to \varphi_{137} \to \varphi_{50}) \to \varphi_{32} \to \varphi_{137} \to \varphi_{136}$
\item $\varphi_{288} :\equiv (r \in \mathbb{R} \to \lvert r \rvert < 1 \to n \in \mathbb{N} \to \varphi_{46}) \to (\varphi_{135} \to \varphi_{46} \to \varphi_{121}) \to r \in \mathbb{R} \to \lvert r \rvert < 1 \to n \in \mathbb{N} \to \varphi_{121}$
\item $\varphi_{289} :\equiv (r \in \mathbb{R} \to \lvert r \rvert < 1 \to n \in \mathbb{N} \to \varphi_{52}) \to (\varphi_{121} \to \varphi_{52} \to \varphi_{139}) \to r \in \mathbb{R} \to \lvert r \rvert < 1 \to n \in \mathbb{N} \to \varphi_{139}$
\item $\varphi_{290} :\equiv (r \in \mathbb{R} \to \lvert r \rvert < 1 \to n \in \mathbb{N} \to \varphi_{76}) \to (\varphi_{140} \to \varphi_{76} \to \varphi_{135}) \to r \in \mathbb{R} \to \lvert r \rvert < 1 \to n \in \mathbb{N} \to \varphi_{135}$
\item $\varphi_{291} :\equiv (r \in \mathbb{R} \to \lvert r \rvert < 1 \to n \in \mathbb{N} \to \varphi_{78}) \to (\varphi_{139} \to \varphi_{78} \to \varphi_{138}) \to r \in \mathbb{R} \to \lvert r \rvert < 1 \to n \in \mathbb{N} \to \varphi_{138}$
\item $\varphi_{292} :\equiv (r \in \mathbb{R} \to \lvert r \rvert < 1 \to n \in \mathbb{N} \to \varphi_{112}) \to (\varphi_{151} \to \varphi_{112} \to \varphi_{140}) \to r \in \mathbb{R} \to \lvert r \rvert < 1 \to n \in \mathbb{N} \to \varphi_{140}$
\item $\varphi_{293} :\equiv \forall w\,(\overline{r} \oplus -1 \odot {(r)^{\bullet}} \circ \mathsf{s}(n) = w \to \varphi_{163} \to \varphi_{61}) \to \varphi_{125} \to \varphi_{163} \to \varphi_{168}$
\item $\varphi_{294} :\equiv (\mathrm{Seq}(\Sigma {(r)^{\bullet}}) \to \forall n \in \mathbb{N}\,(\varphi_{85}) \to \varphi_{114} \to \varphi_{3}) \to (r \in \mathbb{R} \to \lvert r \rvert < 1 \to \mathrm{Seq}(\Sigma {(r)^{\bullet}})) \to \varphi_{261}$
\item $\varphi_{295} :\equiv (r \in \mathbb{R} \to \lvert r \rvert < 1 \to n \in \mathbb{N} \to \varphi_{32}) \to \varphi_{137} \to (\varphi_{32} \to \varphi_{137} \to \varphi_{136}) \to r \in \mathbb{R} \to \lvert r \rvert < 1 \to n \in \mathbb{N} \to \varphi_{136}$
\item $\varphi_{296} :\equiv (r \in \mathbb{R} \to \lvert r \rvert < 1 \to n \in \mathbb{N} \to \varphi_{136}) \to (\varphi_{153} \to \varphi_{136} \to \varphi_{151}) \to r \in \mathbb{R} \to \lvert r \rvert < 1 \to n \in \mathbb{N} \to \varphi_{151}$
\item $\varphi_{297} :\equiv (r \in \mathbb{R} \to \lvert r \rvert < 1 \to n \in \mathbb{N} \to \varphi_{154}) \to (\varphi_{174} \to \varphi_{154} \to \varphi_{153}) \to r \in \mathbb{R} \to \lvert r \rvert < 1 \to n \in \mathbb{N} \to \varphi_{153}$
\item $\varphi_{298} :\equiv (r \in \mathbb{R} \to \lvert r \rvert < 1 \to n \in \mathbb{N} \to \varphi_{168}) \to (\varphi_{167} \to \varphi_{168} \to \varphi_{174}) \to r \in \mathbb{R} \to \lvert r \rvert < 1 \to n \in \mathbb{N} \to \varphi_{174}$
\item $\varphi_{299} :\equiv (r \in \mathbb{R} \to \lvert r \rvert < 1 \to n \in \mathbb{N} \to \varphi_{125}) \to \varphi_{163} \to (\varphi_{125} \to \varphi_{163} \to \varphi_{168}) \to r \in \mathbb{R} \to \lvert r \rvert < 1 \to n \in \mathbb{N} \to \varphi_{168}$
\end{abbrevlist}
\end{fullproof}

% 自動生成 (prover/build.js): PowGeOne  ノード数 389
\begin{fullproof}{[P8]}{$\forall r \in \mathbb{R}\,(1 \leq r \to \forall n \in \mathbb{N}\,(1 \leq {{r}^{n}}))$}{389}\phantomsection\label{tree:PowGeOne}
\noindent$\Pi_{1}$:
\begin{ptree}
\AXC{$r \in \mathbb{R} \to 1 \leq r \to n \in \mathbb{N} \to n \in \mathbb{N}$\lbl{A1}}
\AXC{$(r \in \mathbb{R} \to 1 \leq r \to n \in \mathbb{N} \to n \in \mathbb{N}) \to \varphi_{44}$\lbl{A1}}
\BIC{$\varphi_{44}$}
\end{ptree}
\noindent$\Pi_{2}$:
\begin{ptree}
\AXC{$\varphi_{82}$\lbl{A2}}
\AXC{$r \in \mathbb{R} \to 1 \leq r \to \varphi_{88}$\lbl{A1}}
\AXC{$(r \in \mathbb{R} \to 1 \leq r \to \varphi_{88}) \to \varphi_{82} \to r \in \mathbb{R} \to 1 \leq r \to \varphi_{82}$\lbl{A1}}
\BIC{$\varphi_{82} \to r \in \mathbb{R} \to 1 \leq r \to \varphi_{82}$}
\BIC{$r \in \mathbb{R} \to 1 \leq r \to \varphi_{82}$}
\end{ptree}
\noindent$\Pi_{3}$:
\begin{ptree}
\AXC{$\varphi_{82}$\lbl{A2}}
\AXC{$r \in \mathbb{R} \to 1 \leq r \to \varphi_{88}$\lbl{A1}}
\AXC{$(r \in \mathbb{R} \to 1 \leq r \to \varphi_{88}) \to \varphi_{82} \to r \in \mathbb{R} \to 1 \leq r \to \varphi_{82}$\lbl{A1}}
\BIC{$\varphi_{82} \to r \in \mathbb{R} \to 1 \leq r \to \varphi_{82}$}
\BIC{$r \in \mathbb{R} \to 1 \leq r \to \varphi_{82}$}
\AXC{$\varphi_{142}$\lbl{A1}}
\BIC{$(r \in \mathbb{R} \to 1 \leq r \to \varphi_{39}) \to r \in \mathbb{R} \to 1 \leq r \to \varphi_{85}$}
\end{ptree}
\noindent$\Pi_{4}$:
\begin{ptree}
\AXC{$r \in \mathbb{R} \to 1 \leq r \to \varphi_{39}$\lbl{A1}}
\AXC{$\stackrel{\Pi_{3}}{\vdots}$}\noLine
\UIC{$(r \in \mathbb{R} \to 1 \leq r \to \varphi_{39}) \to r \in \mathbb{R} \to 1 \leq r \to \varphi_{85}$}
\BIC{$r \in \mathbb{R} \to 1 \leq r \to \varphi_{85}$}
\end{ptree}
\noindent$\Pi_{5}$:
\begin{ptree}
\AXC{$\stackrel{\Pi_{4}}{\vdots}$}\noLine
\UIC{$r \in \mathbb{R} \to 1 \leq r \to \varphi_{85}$}
\AXC{$(r \in \mathbb{R} \to 1 \leq r \to \varphi_{85}) \to r \in \mathbb{R} \to 1 \leq r \to \varphi_{42}$\lbl{A1}}
\BIC{$r \in \mathbb{R} \to 1 \leq r \to \varphi_{42}$}
\end{ptree}
\noindent$\Pi_{6}$:
\begin{ptree}
\AXC{$\stackrel{\Pi_{5}}{\vdots}$}\noLine
\UIC{$r \in \mathbb{R} \to 1 \leq r \to \varphi_{42}$}
\AXC{$\forall n\,(n \in \mathbb{N} \leftrightarrow n \in \omega \land n \neq \emptyset)$\lbl{Def$\mathbb{N}$}}
\AXC{$\varphi_{40}$\lbl{A2}}
\BIC{$n \in \mathbb{N} \leftrightarrow n \in \omega \land n \neq \emptyset$}
\AXC{$(n \in \mathbb{N} \leftrightarrow n \in \omega \land n \neq \emptyset) \to \varphi_{143}$\lbl{A1}}
\BIC{$\varphi_{143}$}
\BIC{$(r \in \mathbb{R} \to 1 \leq r \to \varphi_{84}) \to r \in \mathbb{R} \to 1 \leq r \to \varphi_{79}$}
\end{ptree}
\noindent$\Pi_{7}$:
\begin{ptree}
\AXC{$r \in \mathbb{R} \to 1 \leq r \to \varphi_{84}$\lbl{A1}}
\AXC{$\stackrel{\Pi_{6}}{\vdots}$}\noLine
\UIC{$(r \in \mathbb{R} \to 1 \leq r \to \varphi_{84}) \to r \in \mathbb{R} \to 1 \leq r \to \varphi_{79}$}
\BIC{$r \in \mathbb{R} \to 1 \leq r \to \varphi_{79}$}
\end{ptree}
\noindent$\Pi_{8}$:
\begin{ptree}
\AXC{$\stackrel{\Pi_{7}}{\vdots}$}\noLine
\UIC{$r \in \mathbb{R} \to 1 \leq r \to \varphi_{79}$}
\AXC{$\stackrel{\Pi_{4}}{\vdots}$}\noLine
\UIC{$r \in \mathbb{R} \to 1 \leq r \to \varphi_{85}$}
\AXC{$(r \in \mathbb{R} \to 1 \leq r \to \varphi_{85}) \to r \in \mathbb{R} \to 1 \leq r \to \varphi_{62}$\lbl{A1}}
\BIC{$r \in \mathbb{R} \to 1 \leq r \to \varphi_{62}$}
\AXC{$\varphi_{144}$\lbl{A1}}
\BIC{$(r \in \mathbb{R} \to 1 \leq r \to \varphi_{79}) \to r \in \mathbb{R} \to 1 \leq r \to \varphi_{89}$}
\BIC{$r \in \mathbb{R} \to 1 \leq r \to \varphi_{89}$}
\end{ptree}
\noindent$\Pi_{9}$:
\begin{ptree}
\AXC{$r \in \mathbb{R} \to 1 \leq r \to \varphi_{81}$\lbl{A1}}
\AXC{$\forall r \in \mathbb{R}\,(\varphi_{43})$\lbl{[S24]}}
\AXC{$\forall r \in \mathbb{R}\,(\varphi_{43}) \to r \in \mathbb{R} \to \varphi_{43}$\lbl{A2}}
\BIC{$r \in \mathbb{R} \to \varphi_{43}$}
\AXC{$\varphi_{149}$\lbl{A1}}
\BIC{$(r \in \mathbb{R} \to 1 \leq r \to \varphi_{81}) \to r \in \mathbb{R} \to 1 \leq r \to \varphi_{123}$}
\BIC{$r \in \mathbb{R} \to 1 \leq r \to \varphi_{123}$}
\end{ptree}
\noindent$\Pi_{10}$:
\begin{ptree}
\AXC{$\forall a\,\forall n\,(\mathrm{Seq}(a) \to n \in \mathbb{N} \to a_{n} \in \mathbb{R})$\lbl{[S1]}}
\AXC{$\varphi_{60}$\lbl{A2}}
\BIC{$\forall n\,(\mathrm{Seq}({(r)^{\bullet}}) \to n \in \mathbb{N} \to {{r}^{n}} \in \mathbb{R})$}
\AXC{$\varphi_{63}$\lbl{A2}}
\BIC{$\mathrm{Seq}({(r)^{\bullet}}) \to n \in \mathbb{N} \to {{r}^{n}} \in \mathbb{R}$}
\end{ptree}
\noindent$\Pi_{11}$:
\begin{ptree}
\AXC{$\stackrel{\Pi_{10}}{\vdots}$}\noLine
\UIC{$\mathrm{Seq}({(r)^{\bullet}}) \to n \in \mathbb{N} \to {{r}^{n}} \in \mathbb{R}$}
\AXC{$(\mathrm{Seq}({(r)^{\bullet}}) \to n \in \mathbb{N} \to {{r}^{n}} \in \mathbb{R}) \to \varphi_{146}$\lbl{A1}}
\BIC{$\varphi_{146}$}
\end{ptree}
\noindent$\Pi_{12}$:
\begin{ptree}
\AXC{$\stackrel{\Pi_{7}}{\vdots}$}\noLine
\UIC{$r \in \mathbb{R} \to 1 \leq r \to \varphi_{79}$}
\AXC{$\stackrel{\Pi_{9}}{\vdots}$}\noLine
\UIC{$r \in \mathbb{R} \to 1 \leq r \to \varphi_{123}$}
\AXC{$(r \in \mathbb{R} \to 1 \leq r \to \varphi_{123}) \to r \in \mathbb{R} \to \varphi_{91}$\lbl{A1}}
\BIC{$r \in \mathbb{R} \to \varphi_{91}$}
\AXC{$\stackrel{\Pi_{11}}{\vdots}$}\noLine
\UIC{$\varphi_{146}$}
\BIC{$(r \in \mathbb{R} \to 1 \leq r \to \varphi_{79}) \to r \in \mathbb{R} \to 1 \leq r \to \varphi_{98}$}
\BIC{$r \in \mathbb{R} \to 1 \leq r \to \varphi_{98}$}
\end{ptree}
\noindent$\Pi_{13}$:
\begin{ptree}
\AXC{$\varphi_{68}$\lbl{[E12]}}
\AXC{$\varphi_{68} \to \varphi_{64}$\lbl{A2}}
\BIC{$\varphi_{64}$}
\AXC{$\varphi_{64} \to \varphi_{105}$\lbl{A2}}
\BIC{$\varphi_{105}$}
\AXC{$\varphi_{105} \to \varphi_{101}$\lbl{A2}}
\BIC{$\varphi_{101}$}
\AXC{$\varphi_{101} \to 1 \in \mathbb{R} \to {{r}^{n}} \in \mathbb{R} \to \varphi_{49}$\lbl{A2}}
\BIC{$1 \in \mathbb{R} \to {{r}^{n}} \in \mathbb{R} \to \varphi_{49}$}
\AXC{$\varphi_{156}$\lbl{A1}}
\BIC{$1 \in \mathbb{R} \to (r \in \mathbb{R} \to 1 \leq r \to \varphi_{98}) \to 1 \in \mathbb{R} \to \varphi_{155}$}
\end{ptree}
\noindent$\Pi_{14}$:
\begin{ptree}
\AXC{$1 \in \mathbb{R}$\lbl{(R1)}}
\AXC{$\stackrel{\Pi_{12}}{\vdots}$}\noLine
\UIC{$r \in \mathbb{R} \to 1 \leq r \to \varphi_{98}$}
\AXC{$1 \in \mathbb{R}$\lbl{(R1)}}
\AXC{$\stackrel{\Pi_{13}}{\vdots}$}\noLine
\UIC{$1 \in \mathbb{R} \to (r \in \mathbb{R} \to 1 \leq r \to \varphi_{98}) \to 1 \in \mathbb{R} \to \varphi_{155}$}
\BIC{$(r \in \mathbb{R} \to 1 \leq r \to \varphi_{98}) \to 1 \in \mathbb{R} \to \varphi_{155}$}
\BIC{$1 \in \mathbb{R} \to \varphi_{155}$}
\BIC{$\varphi_{155}$}
\end{ptree}
\noindent$\Pi_{15}$:
\begin{ptree}
\AXC{$0 < 1$\lbl{[F28]}}
\AXC{$0 < 1 \to r \in \mathbb{R} \to 1 \leq r \to \varphi_{73}$\lbl{A1}}
\BIC{$r \in \mathbb{R} \to 1 \leq r \to \varphi_{73}$}
\AXC{$r \in \mathbb{R} \to 1 \leq r \to \varphi_{81}$\lbl{A1}}
\AXC{$\stackrel{\Pi_{14}}{\vdots}$}\noLine
\UIC{$\varphi_{155}$}
\BIC{$(r \in \mathbb{R} \to 1 \leq r \to \varphi_{73}) \to (r \in \mathbb{R} \to 1 \leq r \to \varphi_{73}) \to \varphi_{137}$}
\BIC{$(r \in \mathbb{R} \to 1 \leq r \to \varphi_{73}) \to \varphi_{137}$}
\end{ptree}
\noindent$\Pi_{16}$:
\begin{ptree}
\AXC{$\stackrel{\Pi_{8}}{\vdots}$}\noLine
\UIC{$r \in \mathbb{R} \to 1 \leq r \to \varphi_{89}$}
\AXC{$0 < 1$\lbl{[F28]}}
\AXC{$0 < 1 \to r \in \mathbb{R} \to 1 \leq r \to \varphi_{73}$\lbl{A1}}
\BIC{$r \in \mathbb{R} \to 1 \leq r \to \varphi_{73}$}
\AXC{$\stackrel{\Pi_{15}}{\vdots}$}\noLine
\UIC{$(r \in \mathbb{R} \to 1 \leq r \to \varphi_{73}) \to \varphi_{137}$}
\BIC{$\varphi_{137}$}
\BIC{$(r \in \mathbb{R} \to 1 \leq r \to 1 \leq r) \to r \in \mathbb{R} \to 1 \leq r \to \varphi_{100}$}
\end{ptree}
\noindent$\Pi_{17}$:
\begin{ptree}
\AXC{$r \in \mathbb{R} \to 1 \leq r \to 1 \leq r$\lbl{A1}}
\AXC{$\stackrel{\Pi_{16}}{\vdots}$}\noLine
\UIC{$(r \in \mathbb{R} \to 1 \leq r \to 1 \leq r) \to r \in \mathbb{R} \to 1 \leq r \to \varphi_{100}$}
\BIC{$r \in \mathbb{R} \to 1 \leq r \to \varphi_{100}$}
\end{ptree}
\noindent$\Pi_{18}$:
\begin{ptree}
\AXC{$\forall x \in \mathbb{R}\,(x \cdot 1 = x)$\lbl{(R8)}}
\AXC{$\forall x \in \mathbb{R}\,(x \cdot 1 = x) \to 1 \in \mathbb{R} \to 1 \cdot 1 = 1$\lbl{A2}}
\BIC{$1 \in \mathbb{R} \to 1 \cdot 1 = 1$}
\end{ptree}
\noindent$\Pi_{19}$:
\begin{ptree}
\AXC{$1 \in \mathbb{R}$\lbl{(R1)}}
\AXC{$\stackrel{\Pi_{18}}{\vdots}$}\noLine
\UIC{$1 \in \mathbb{R} \to 1 \cdot 1 = 1$}
\AXC{$(1 \in \mathbb{R} \to 1 \cdot 1 = 1) \to 1 \in \mathbb{R} \to r \in \mathbb{R} \to 1 \leq r \to \varphi_{76}$\lbl{A1}}
\BIC{$1 \in \mathbb{R} \to r \in \mathbb{R} \to 1 \leq r \to \varphi_{76}$}
\BIC{$r \in \mathbb{R} \to 1 \leq r \to \varphi_{76}$}
\end{ptree}
\noindent$\Pi_{20}$:
\begin{ptree}
\AXC{$\stackrel{\Pi_{17}}{\vdots}$}\noLine
\UIC{$r \in \mathbb{R} \to 1 \leq r \to \varphi_{100}$}
\AXC{$\stackrel{\Pi_{19}}{\vdots}$}\noLine
\UIC{$r \in \mathbb{R} \to 1 \leq r \to \varphi_{76}$}
\AXC{$(r \in \mathbb{R} \to 1 \leq r \to \varphi_{76}) \to \varphi_{147}$\lbl{A1}}
\BIC{$\varphi_{147}$}
\BIC{$(1 \cdot 1 = 1 \to 1 \cdot 1 < {{r}^{n}} \cdot r \to 1 < {{r}^{n}} \cdot r) \to \varphi_{127}$}
\end{ptree}
\noindent$\Pi_{21}$:
\begin{ptree}
\AXC{$z = w \to z < {{r}^{n}} \cdot r \to w < {{r}^{n}} \cdot r$\lbl{A6}}
\UIC{$\forall w\,(z = w \to z < {{r}^{n}} \cdot r \to w < {{r}^{n}} \cdot r)$}
\UIC{$\forall z\,\forall w\,(z = w \to z < {{r}^{n}} \cdot r \to w < {{r}^{n}} \cdot r)$}
\AXC{$\varphi_{104}$\lbl{A2}}
\BIC{$\forall w\,(1 \cdot 1 = w \to 1 \cdot 1 < {{r}^{n}} \cdot r \to w < {{r}^{n}} \cdot r)$}
\AXC{$\varphi_{107}$\lbl{A2}}
\BIC{$1 \cdot 1 = 1 \to 1 \cdot 1 < {{r}^{n}} \cdot r \to 1 < {{r}^{n}} \cdot r$}
\AXC{$\stackrel{\Pi_{20}}{\vdots}$}\noLine
\UIC{$(1 \cdot 1 = 1 \to 1 \cdot 1 < {{r}^{n}} \cdot r \to 1 < {{r}^{n}} \cdot r) \to \varphi_{127}$}
\BIC{$\varphi_{127}$}
\end{ptree}
\noindent$\Pi_{22}$:
\begin{ptree}
\AXC{$x = x$\lbl{A5}}
\AXC{$x = y \to x = x \to y = x$\lbl{A6}}
\AXC{$(x = y \to x = x \to y = x) \to x = x \to x = y \to y = x$\lbl{A1}}
\BIC{$x = x \to x = y \to y = x$}
\BIC{$x = y \to y = x$}
\UIC{$\forall y\,(x = y \to y = x)$}
\UIC{$\forall x\,\forall y\,(x = y \to y = x)$}
\end{ptree}
\noindent$\Pi_{23}$:
\begin{ptree}
\AXC{$r \in \mathbb{R} \to 1 \leq r \to \varphi_{81}$\lbl{A1}}
\AXC{$\forall r \in \mathbb{R}\,(\varphi_{43})$\lbl{[S24]}}
\AXC{$\forall r \in \mathbb{R}\,(\varphi_{43}) \to r \in \mathbb{R} \to \varphi_{43}$\lbl{A2}}
\BIC{$r \in \mathbb{R} \to \varphi_{43}$}
\AXC{$\varphi_{149}$\lbl{A1}}
\BIC{$(r \in \mathbb{R} \to 1 \leq r \to \varphi_{81}) \to r \in \mathbb{R} \to 1 \leq r \to \varphi_{123}$}
\BIC{$r \in \mathbb{R} \to 1 \leq r \to \varphi_{123}$}
\end{ptree}
\noindent$\Pi_{24}$:
\begin{ptree}
\AXC{$\varphi_{61}$\lbl{A2}}
\AXC{$\stackrel{\Pi_{23}}{\vdots}$}\noLine
\UIC{$r \in \mathbb{R} \to 1 \leq r \to \varphi_{123}$}
\AXC{$(r \in \mathbb{R} \to 1 \leq r \to \varphi_{123}) \to r \in \mathbb{R} \to \varphi_{118}$\lbl{A1}}
\BIC{$r \in \mathbb{R} \to \varphi_{118}$}
\AXC{$(r \in \mathbb{R} \to \varphi_{118}) \to \varphi_{61} \to r \in \mathbb{R} \to \varphi_{117}$\lbl{A1}}
\BIC{$\varphi_{61} \to r \in \mathbb{R} \to \varphi_{117}$}
\BIC{$r \in \mathbb{R} \to \varphi_{117}$}
\end{ptree}
\noindent$\Pi_{25}$:
\begin{ptree}
\AXC{$\stackrel{\Pi_{7}}{\vdots}$}\noLine
\UIC{$r \in \mathbb{R} \to 1 \leq r \to \varphi_{79}$}
\AXC{$\stackrel{\Pi_{24}}{\vdots}$}\noLine
\UIC{$r \in \mathbb{R} \to \varphi_{117}$}
\AXC{$\varphi_{151}$\lbl{A1}}
\BIC{$(r \in \mathbb{R} \to 1 \leq r \to \varphi_{79}) \to r \in \mathbb{R} \to 1 \leq r \to \varphi_{108}$}
\BIC{$r \in \mathbb{R} \to 1 \leq r \to \varphi_{108}$}
\AXC{$(r \in \mathbb{R} \to 1 \leq r \to \varphi_{108}) \to \varphi_{132}$\lbl{A1}}
\BIC{$\varphi_{132}$}
\end{ptree}
\noindent$\Pi_{26}$:
\begin{ptree}
\AXC{$\stackrel{\Pi_{22}}{\vdots}$}\noLine
\UIC{$\forall x\,\forall y\,(x = y \to y = x)$}
\AXC{$\varphi_{16}$\lbl{A2}}
\BIC{$\forall y\,({{r}^{S(n)}} = y \to y = {{r}^{S(n)}})$}
\AXC{$\varphi_{87}$\lbl{A2}}
\BIC{${{r}^{S(n)}} = {{r}^{n}} \cdot r \to {{r}^{n}} \cdot r = {{r}^{S(n)}}$}
\AXC{$\stackrel{\Pi_{25}}{\vdots}$}\noLine
\UIC{$\varphi_{132}$}
\BIC{$r \in \mathbb{R} \to 1 \leq r \to \varphi_{109}$}
\AXC{$(r \in \mathbb{R} \to 1 \leq r \to \varphi_{109}) \to \varphi_{150}$\lbl{A1}}
\BIC{$\varphi_{150}$}
\end{ptree}
\noindent$\Pi_{27}$:
\begin{ptree}
\AXC{$z = w \to z = {{r}^{n}} \cdot r \to w = {{r}^{n}} \cdot r$\lbl{A6}}
\UIC{$\forall w\,(z = w \to z = {{r}^{n}} \cdot r \to w = {{r}^{n}} \cdot r)$}
\UIC{$\forall z\,\forall w\,(z = w \to z = {{r}^{n}} \cdot r \to w = {{r}^{n}} \cdot r)$}
\AXC{$\varphi_{93}$\lbl{A2}}
\BIC{$\forall w\,(1 \cdot 1 = w \to 1 \cdot 1 = {{r}^{n}} \cdot r \to w = {{r}^{n}} \cdot r)$}
\AXC{$\varphi_{96}$\lbl{A2}}
\BIC{$1 \cdot 1 = 1 \to 1 \cdot 1 = {{r}^{n}} \cdot r \to 1 = {{r}^{n}} \cdot r$}
\AXC{$\stackrel{\Pi_{21}}{\vdots}$}\noLine
\UIC{$\varphi_{127}$}
\BIC{$r \in \mathbb{R} \to 1 \leq r \to \varphi_{95}$}
\AXC{$\stackrel{\Pi_{26}}{\vdots}$}\noLine
\UIC{$\varphi_{150}$}
\BIC{$({{r}^{n}} \cdot r = {{r}^{S(n)}} \to 1 < {{r}^{n}} \cdot r \to 1 < {{r}^{S(n)}}) \to \varphi_{130}$}
\end{ptree}
\noindent$\Pi_{28}$:
\begin{ptree}
\AXC{$z = w \to 1 < z \to 1 < w$\lbl{A6}}
\UIC{$\forall w\,(z = w \to 1 < z \to 1 < w)$}
\UIC{$\forall z\,\forall w\,(z = w \to 1 < z \to 1 < w)$}
\AXC{$\varphi_{51}$\lbl{A2}}
\BIC{$\forall w\,({{r}^{n}} \cdot r = w \to 1 < {{r}^{n}} \cdot r \to 1 < w)$}
\AXC{$\varphi_{111}$\lbl{A2}}
\BIC{${{r}^{n}} \cdot r = {{r}^{S(n)}} \to 1 < {{r}^{n}} \cdot r \to 1 < {{r}^{S(n)}}$}
\AXC{$\stackrel{\Pi_{27}}{\vdots}$}\noLine
\UIC{$({{r}^{n}} \cdot r = {{r}^{S(n)}} \to 1 < {{r}^{n}} \cdot r \to 1 < {{r}^{S(n)}}) \to \varphi_{130}$}
\BIC{$\varphi_{130}$}
\end{ptree}
\noindent$\Pi_{29}$:
\begin{ptree}
\AXC{$r \in \mathbb{R} \to 1 \leq r \to r \in \mathbb{R}$\lbl{A1}}
\AXC{$\forall r \in \mathbb{R}\,(\varphi_{43})$\lbl{[S24]}}
\AXC{$\forall r \in \mathbb{R}\,(\varphi_{43}) \to r \in \mathbb{R} \to \varphi_{43}$\lbl{A2}}
\BIC{$r \in \mathbb{R} \to \varphi_{43}$}
\AXC{$\varphi_{133}$\lbl{A1}}
\BIC{$(r \in \mathbb{R} \to 1 \leq r \to r \in \mathbb{R}) \to r \in \mathbb{R} \to 1 \leq r \to \varphi_{43}$}
\BIC{$r \in \mathbb{R} \to 1 \leq r \to \varphi_{43}$}
\end{ptree}
\noindent$\Pi_{30}$:
\begin{ptree}
\AXC{$\stackrel{\Pi_{29}}{\vdots}$}\noLine
\UIC{$r \in \mathbb{R} \to 1 \leq r \to \varphi_{43}$}
\AXC{$(r \in \mathbb{R} \to 1 \leq r \to \varphi_{43}) \to r \in \mathbb{R} \to 1 \leq r \to {{r}^{1_{\omega}}} = r$\lbl{A1}}
\BIC{$r \in \mathbb{R} \to 1 \leq r \to {{r}^{1_{\omega}}} = r$}
\AXC{$(r \in \mathbb{R} \to 1 \leq r \to {{r}^{1_{\omega}}} = r) \to \varphi_{38}$\lbl{A1}}
\BIC{$\varphi_{38}$}
\end{ptree}
\noindent$\Pi_{31}$:
\begin{ptree}
\AXC{$\stackrel{\Pi_{22}}{\vdots}$}\noLine
\UIC{$\forall x\,\forall y\,(x = y \to y = x)$}
\AXC{$\varphi_{17}$\lbl{A2}}
\BIC{$\forall y\,({{r}^{1_{\omega}}} = y \to y = {{r}^{1_{\omega}}})$}
\AXC{$\varphi_{41}$\lbl{A2}}
\BIC{${{r}^{1_{\omega}}} = r \to r = {{r}^{1_{\omega}}}$}
\AXC{$\stackrel{\Pi_{30}}{\vdots}$}\noLine
\UIC{$\varphi_{38}$}
\BIC{$r \in \mathbb{R} \to 1 \leq r \to r = {{r}^{1_{\omega}}}$}
\AXC{$(r \in \mathbb{R} \to 1 \leq r \to r = {{r}^{1_{\omega}}}) \to \varphi_{121}$\lbl{A1}}
\BIC{$\varphi_{121}$}
\end{ptree}
\noindent$\Pi_{32}$:
\begin{ptree}
\AXC{$z = w \to 1 < z \to 1 < w$\lbl{A6}}
\UIC{$\forall w\,(z = w \to 1 < z \to 1 < w)$}
\UIC{$\forall z\,\forall w\,(z = w \to 1 < z \to 1 < w)$}
\AXC{$\varphi_{9}$\lbl{A2}}
\BIC{$\forall w\,(r = w \to 1 < r \to 1 < w)$}
\AXC{$\varphi_{37}$\lbl{A2}}
\BIC{$r = {{r}^{1_{\omega}}} \to 1 < r \to 1 < {{r}^{1_{\omega}}}$}
\AXC{$r \in \mathbb{R} \to 1 \leq r \to 1 \leq r$\lbl{A1}}
\AXC{$\stackrel{\Pi_{31}}{\vdots}$}\noLine
\UIC{$\varphi_{121}$}
\BIC{$(r = {{r}^{1_{\omega}}} \to 1 < r \to 1 < {{r}^{1_{\omega}}}) \to \varphi_{45}$}
\BIC{$\varphi_{45}$}
\end{ptree}
\noindent$\Pi_{33}$:
\begin{ptree}
\AXC{$y = z \to x = y \to x = z$\lbl{A6}}
\AXC{$(y = z \to x = y \to x = z) \to x = y \to y = z \to x = z$\lbl{A1}}
\BIC{$x = y \to y = z \to x = z$}
\UIC{$\forall z\,(x = y \to y = z \to x = z)$}
\UIC{$\forall y\,\forall z\,(x = y \to y = z \to x = z)$}
\UIC{$\forall x\,\forall y\,\forall z\,(x = y \to y = z \to x = z)$}
\AXC{$\varphi_{10}$\lbl{A2}}
\BIC{$\forall y\,\forall z\,(S(n) = y \to y = z \to S(n) = z)$}
\AXC{$\varphi_{33}$\lbl{A2}}
\BIC{$\forall z\,(S(n) = S(\emptyset) \to S(\emptyset) = z \to S(n) = z)$}
\end{ptree}
\noindent$\Pi_{34}$:
\begin{ptree}
\AXC{$\stackrel{\Pi_{33}}{\vdots}$}\noLine
\UIC{$\forall z\,(S(n) = S(\emptyset) \to S(\emptyset) = z \to S(n) = z)$}
\AXC{$\forall z\,(S(n) = S(\emptyset) \to S(\emptyset) = z \to S(n) = z) \to \varphi_{5}$\lbl{A2}}
\BIC{$\varphi_{5}$}
\end{ptree}
\noindent$\Pi_{35}$:
\begin{ptree}
\AXC{$\stackrel{\Pi_{22}}{\vdots}$}\noLine
\UIC{$\forall x\,\forall y\,(x = y \to y = x)$}
\AXC{$\varphi_{2}$\lbl{A2}}
\BIC{$\forall y\,(1_{\omega} = y \to y = 1_{\omega})$}
\AXC{$\varphi_{12}$\lbl{A2}}
\BIC{$1_{\omega} = S(\emptyset) \to S(\emptyset) = 1_{\omega}$}
\AXC{$1_{\omega} = S(\emptyset)$\lbl{Def$1_\omega$}}
\AXC{$1_{\omega} = S(\emptyset) \to \varphi_{11}$\lbl{A1}}
\BIC{$\varphi_{11}$}
\BIC{$S(\emptyset) = 1_{\omega}$}
\end{ptree}
\noindent$\Pi_{36}$:
\begin{ptree}
\AXC{$r \in \mathbb{R} \to 1 \leq r \to \varphi_{80}$\lbl{A1}}
\AXC{$(r \in \mathbb{R} \to 1 \leq r \to \varphi_{80}) \to S(n) = S(n) \to \varphi_{122}$\lbl{A1}}
\BIC{$S(n) = S(n) \to \varphi_{122}$}
\end{ptree}
\noindent$\Pi_{37}$:
\begin{ptree}
\AXC{$x = x$\lbl{A5}}
\UIC{$\forall x\,(x = x)$}
\AXC{$\forall x\,(x = x) \to S(n) = S(n)$\lbl{A2}}
\BIC{$S(n) = S(n)$}
\AXC{$\stackrel{\Pi_{36}}{\vdots}$}\noLine
\UIC{$S(n) = S(n) \to \varphi_{122}$}
\BIC{$\varphi_{122}$}
\end{ptree}
\noindent$\Pi_{38}$:
\begin{ptree}
\AXC{$z = w \to S(n) = S(z) \to S(n) = S(w)$\lbl{A6}}
\UIC{$\forall w\,(z = w \to S(n) = S(z) \to S(n) = S(w))$}
\UIC{$\forall z\,\forall w\,(z = w \to S(n) = S(z) \to S(n) = S(w))$}
\AXC{$\varphi_{26}$\lbl{A2}}
\BIC{$\forall w\,(n = w \to S(n) = S(n) \to S(n) = S(w))$}
\AXC{$\varphi_{34}$\lbl{A2}}
\BIC{$n = \emptyset \to S(n) = S(n) \to S(n) = S(\emptyset)$}
\AXC{$\stackrel{\Pi_{37}}{\vdots}$}\noLine
\UIC{$\varphi_{122}$}
\BIC{$r \in \mathbb{R} \to 1 \leq r \to \varphi_{90}$}
\AXC{$\varphi_{139}$\lbl{A1}}
\BIC{$S(\emptyset) = 1_{\omega} \to \varphi_{5} \to r \in \mathbb{R} \to 1 \leq r \to \varphi_{77}$}
\end{ptree}
\noindent$\Pi_{39}$:
\begin{ptree}
\AXC{$\stackrel{\Pi_{34}}{\vdots}$}\noLine
\UIC{$\varphi_{5}$}
\AXC{$\stackrel{\Pi_{35}}{\vdots}$}\noLine
\UIC{$S(\emptyset) = 1_{\omega}$}
\AXC{$\stackrel{\Pi_{38}}{\vdots}$}\noLine
\UIC{$S(\emptyset) = 1_{\omega} \to \varphi_{5} \to r \in \mathbb{R} \to 1 \leq r \to \varphi_{77}$}
\BIC{$\varphi_{5} \to r \in \mathbb{R} \to 1 \leq r \to \varphi_{77}$}
\BIC{$r \in \mathbb{R} \to 1 \leq r \to \varphi_{77}$}
\AXC{$(r \in \mathbb{R} \to 1 \leq r \to \varphi_{77}) \to \varphi_{116}$\lbl{A1}}
\BIC{$\varphi_{116}$}
\end{ptree}
\noindent$\Pi_{40}$:
\begin{ptree}
\AXC{$\stackrel{\Pi_{22}}{\vdots}$}\noLine
\UIC{$\forall x\,\forall y\,(x = y \to y = x)$}
\AXC{$\varphi_{1}$\lbl{A2}}
\BIC{$\forall y\,(S(n) = y \to y = S(n))$}
\AXC{$\varphi_{4}$\lbl{A2}}
\BIC{$S(n) = 1_{\omega} \to 1_{\omega} = S(n)$}
\AXC{$\stackrel{\Pi_{39}}{\vdots}$}\noLine
\UIC{$\varphi_{116}$}
\BIC{$r \in \mathbb{R} \to 1 \leq r \to \varphi_{78}$}
\AXC{$(r \in \mathbb{R} \to 1 \leq r \to \varphi_{78}) \to \varphi_{136}$\lbl{A1}}
\BIC{$\varphi_{136}$}
\end{ptree}
\noindent$\Pi_{41}$:
\begin{ptree}
\AXC{$z = w \to 1 = z \to 1 = w$\lbl{A6}}
\UIC{$\forall w\,(z = w \to 1 = z \to 1 = w)$}
\UIC{$\forall z\,\forall w\,(z = w \to 1 = z \to 1 = w)$}
\AXC{$\varphi_{3}$\lbl{A2}}
\BIC{$\forall w\,(r = w \to 1 = r \to 1 = w)$}
\AXC{$\varphi_{24}$\lbl{A2}}
\BIC{$r = {{r}^{1_{\omega}}} \to 1 = r \to 1 = {{r}^{1_{\omega}}}$}
\AXC{$\stackrel{\Pi_{32}}{\vdots}$}\noLine
\UIC{$\varphi_{45}$}
\BIC{$r \in \mathbb{R} \to 1 \leq r \to 1 \leq {{r}^{1_{\omega}}}$}
\AXC{$\stackrel{\Pi_{40}}{\vdots}$}\noLine
\UIC{$\varphi_{136}$}
\BIC{$(1_{\omega} = S(n) \to 1 < {{r}^{1_{\omega}}} \to 1 < {{r}^{S(n)}}) \to \varphi_{125}$}
\end{ptree}
\noindent$\Pi_{42}$:
\begin{ptree}
\AXC{$z = w \to 1 < {{r}^{z}} \to 1 < {{r}^{w}}$\lbl{A6}}
\UIC{$\forall w\,(z = w \to 1 < {{r}^{z}} \to 1 < {{r}^{w}})$}
\UIC{$\forall z\,\forall w\,(z = w \to 1 < {{r}^{z}} \to 1 < {{r}^{w}})$}
\AXC{$\varphi_{75}$\lbl{A2}}
\BIC{$\forall w\,(1_{\omega} = w \to 1 < {{r}^{1_{\omega}}} \to 1 < {{r}^{w}})$}
\AXC{$\varphi_{83}$\lbl{A2}}
\BIC{$1_{\omega} = S(n) \to 1 < {{r}^{1_{\omega}}} \to 1 < {{r}^{S(n)}}$}
\AXC{$\stackrel{\Pi_{41}}{\vdots}$}\noLine
\UIC{$(1_{\omega} = S(n) \to 1 < {{r}^{1_{\omega}}} \to 1 < {{r}^{S(n)}}) \to \varphi_{125}$}
\BIC{$\varphi_{125}$}
\end{ptree}
\noindent$\Pi_{43}$:
\begin{ptree}
\AXC{$z = w \to 1 = {{r}^{z}} \to 1 = {{r}^{w}}$\lbl{A6}}
\UIC{$\forall w\,(z = w \to 1 = {{r}^{z}} \to 1 = {{r}^{w}})$}
\UIC{$\forall z\,\forall w\,(z = w \to 1 = {{r}^{z}} \to 1 = {{r}^{w}})$}
\AXC{$\varphi_{57}$\lbl{A2}}
\BIC{$\forall w\,(1_{\omega} = w \to 1 = {{r}^{1_{\omega}}} \to 1 = {{r}^{w}})$}
\AXC{$\varphi_{65}$\lbl{A2}}
\BIC{$1_{\omega} = S(n) \to 1 = {{r}^{1_{\omega}}} \to 1 = {{r}^{S(n)}}$}
\AXC{$\stackrel{\Pi_{42}}{\vdots}$}\noLine
\UIC{$\varphi_{125}$}
\BIC{$r \in \mathbb{R} \to 1 \leq r \to \varphi_{92}$}
\AXC{$\varphi_{145}$\lbl{A1}}
\BIC{$(r \in \mathbb{R} \to 1 \leq r \to \varphi_{94}) \to r \in \mathbb{R} \to 1 \leq r \to \varphi_{74}$}
\end{ptree}
\noindent$\Pi_{44}$:
\begin{ptree}
\AXC{$z = w \to 1 = z \to 1 = w$\lbl{A6}}
\UIC{$\forall w\,(z = w \to 1 = z \to 1 = w)$}
\UIC{$\forall z\,\forall w\,(z = w \to 1 = z \to 1 = w)$}
\AXC{$\varphi_{32}$\lbl{A2}}
\BIC{$\forall w\,({{r}^{n}} \cdot r = w \to 1 = {{r}^{n}} \cdot r \to 1 = w)$}
\AXC{$\varphi_{103}$\lbl{A2}}
\BIC{${{r}^{n}} \cdot r = {{r}^{S(n)}} \to 1 = {{r}^{n}} \cdot r \to 1 = {{r}^{S(n)}}$}
\AXC{$\stackrel{\Pi_{28}}{\vdots}$}\noLine
\UIC{$\varphi_{130}$}
\BIC{$r \in \mathbb{R} \to 1 \leq r \to \varphi_{94}$}
\AXC{$\stackrel{\Pi_{43}}{\vdots}$}\noLine
\UIC{$(r \in \mathbb{R} \to 1 \leq r \to \varphi_{94}) \to r \in \mathbb{R} \to 1 \leq r \to \varphi_{74}$}
\BIC{$r \in \mathbb{R} \to 1 \leq r \to \varphi_{74}$}
\end{ptree}
\noindent$\Pi_{45}$:
\begin{ptree}
\AXC{$\forall x\,(\mathrm{Ind}(x) \leftrightarrow \emptyset \in x \land \forall y \in x\,(S(y) \in x))$\lbl{DefInd}}
\AXC{$\varphi_{97}$\lbl{A2}}
\BIC{$\mathrm{Ind}(\omega) \leftrightarrow \emptyset \in \omega \land \forall y \in \omega\,(S(y) \in \omega)$}
\AXC{$\mathrm{Ind}(\omega)$\lbl{補題(ωは帰納的)}}
\AXC{$\mathrm{Ind}(\omega) \to \varphi_{128}$\lbl{A1}}
\BIC{$\varphi_{128}$}
\BIC{$r \in \mathbb{R} \to 1 \leq r \to \varphi_{86}$}
\end{ptree}
\noindent$\Pi_{46}$:
\begin{ptree}
\AXC{$\stackrel{\Pi_{45}}{\vdots}$}\noLine
\UIC{$r \in \mathbb{R} \to 1 \leq r \to \varphi_{86}$}
\AXC{$(r \in \mathbb{R} \to 1 \leq r \to \varphi_{86}) \to r \in \mathbb{R} \to 1 \leq r \to \varphi_{70}$\lbl{A1}}
\BIC{$r \in \mathbb{R} \to 1 \leq r \to \varphi_{70}$}
\AXC{$(r \in \mathbb{R} \to 1 \leq r \to \varphi_{70}) \to \varphi_{124}$\lbl{A1}}
\BIC{$\varphi_{124}$}
\end{ptree}
\noindent$\Pi_{47}$:
\begin{ptree}
\AXC{$\stackrel{\Pi_{5}}{\vdots}$}\noLine
\UIC{$r \in \mathbb{R} \to 1 \leq r \to \varphi_{42}$}
\AXC{$\forall y \in \omega\,(S(y) \in \omega) \to n \in \omega \to S(n) \in \omega$\lbl{A2}}
\AXC{$\stackrel{\Pi_{46}}{\vdots}$}\noLine
\UIC{$\varphi_{124}$}
\BIC{$r \in \mathbb{R} \to 1 \leq r \to \varphi_{66}$}
\AXC{$\varphi_{138}$\lbl{A1}}
\BIC{$(r \in \mathbb{R} \to 1 \leq r \to \varphi_{42}) \to r \in \mathbb{R} \to 1 \leq r \to \varphi_{48}$}
\BIC{$r \in \mathbb{R} \to 1 \leq r \to \varphi_{48}$}
\end{ptree}
\noindent$\Pi_{48}$:
\begin{ptree}
\AXC{$\forall n\,(n \in B_{54} \leftrightarrow n \in \omega \land (n \in \mathbb{N} \to 1 \leq {{r}^{n}})) \to \varphi_{25}$\lbl{A2}}
\AXC{$r \in \mathbb{R} \to 1 \leq r \to \varphi_{88}$\lbl{A1}}
\AXC{$(r \in \mathbb{R} \to 1 \leq r \to \varphi_{88}) \to \varphi_{135}$\lbl{A1}}
\BIC{$\varphi_{135}$}
\BIC{$r \in \mathbb{R} \to \varphi_{106}$}
\end{ptree}
\noindent$\Pi_{49}$:
\begin{ptree}
\AXC{$\stackrel{\Pi_{44}}{\vdots}$}\noLine
\UIC{$r \in \mathbb{R} \to 1 \leq r \to \varphi_{74}$}
\AXC{$\stackrel{\Pi_{47}}{\vdots}$}\noLine
\UIC{$r \in \mathbb{R} \to 1 \leq r \to \varphi_{48}$}
\AXC{$\stackrel{\Pi_{48}}{\vdots}$}\noLine
\UIC{$r \in \mathbb{R} \to \varphi_{106}$}
\AXC{$(r \in \mathbb{R} \to \varphi_{106}) \to \varphi_{140}$\lbl{A1}}
\BIC{$\varphi_{140}$}
\BIC{$(r \in \mathbb{R} \to 1 \leq r \to \varphi_{74}) \to r \in \mathbb{R} \to 1 \leq r \to \varphi_{46}$}
\BIC{$r \in \mathbb{R} \to 1 \leq r \to \varphi_{46}$}
\end{ptree}
\noindent$\Pi_{50}$:
\begin{ptree}
\AXC{$\varphi_{131}$\lbl{A3}}
\AXC{$\stackrel{\Pi_{49}}{\vdots}$}\noLine
\UIC{$r \in \mathbb{R} \to 1 \leq r \to \varphi_{46}$}
\AXC{$(r \in \mathbb{R} \to 1 \leq r \to \varphi_{46}) \to \varphi_{35} \to n \in B_{54} \to S(n) \in B_{54}$\lbl{A1}}
\BIC{$\varphi_{35} \to n \in B_{54} \to S(n) \in B_{54}$}
\UIC{$\forall n\,(\varphi_{35} \to n \in B_{54} \to S(n) \in B_{54})$}
\AXC{$\varphi_{152}$\lbl{A1}}
\BIC{$\varphi_{131} \to (\varphi_{35} \to \forall n\,(\varphi_{35})) \to r \in \mathbb{R} \to 1 \leq r \to \varphi_{50}$}
\BIC{$(\varphi_{35} \to \forall n\,(\varphi_{35})) \to r \in \mathbb{R} \to 1 \leq r \to \varphi_{50}$}
\end{ptree}
\noindent$\Pi_{51}$:
\begin{ptree}
\AXC{$x = x$\lbl{A5}}
\UIC{$\forall x\,(x = x)$}
\AXC{$\forall x\,(x = x) \to \emptyset = \emptyset$\lbl{A2}}
\BIC{$\emptyset = \emptyset$}
\AXC{$\forall n\,(n \in \mathbb{N} \leftrightarrow n \in \omega \land n \neq \emptyset)$\lbl{Def$\mathbb{N}$}}
\AXC{$\varphi_{69}$\lbl{A2}}
\BIC{$\emptyset \in \mathbb{N} \leftrightarrow \emptyset \in \omega \land \emptyset \neq \emptyset$}
\AXC{$\varphi_{112}$\lbl{A1}}
\BIC{$\emptyset = \emptyset \to r \in \mathbb{R} \to 1 \leq r \to \emptyset \in \mathbb{N} \to 1 \leq {{r}^{\emptyset}}$}
\BIC{$r \in \mathbb{R} \to 1 \leq r \to \emptyset \in \mathbb{N} \to 1 \leq {{r}^{\emptyset}}$}
\end{ptree}
\noindent$\Pi_{52}$:
\begin{ptree}
\AXC{$\varphi_{110}$\lbl{A2}}
\AXC{$r \in \mathbb{R} \to 1 \leq r \to \varphi_{88}$\lbl{A1}}
\AXC{$(r \in \mathbb{R} \to 1 \leq r \to \varphi_{88}) \to \varphi_{110} \to r \in \mathbb{R} \to 1 \leq r \to \varphi_{110}$\lbl{A1}}
\BIC{$\varphi_{110} \to r \in \mathbb{R} \to 1 \leq r \to \varphi_{110}$}
\BIC{$r \in \mathbb{R} \to 1 \leq r \to \varphi_{110}$}
\end{ptree}
\noindent$\Pi_{53}$:
\begin{ptree}
\AXC{$\stackrel{\Pi_{45}}{\vdots}$}\noLine
\UIC{$r \in \mathbb{R} \to 1 \leq r \to \varphi_{86}$}
\AXC{$\stackrel{\Pi_{52}}{\vdots}$}\noLine
\UIC{$r \in \mathbb{R} \to 1 \leq r \to \varphi_{110}$}
\AXC{$(r \in \mathbb{R} \to 1 \leq r \to \varphi_{110}) \to (r \in \mathbb{R} \to 1 \leq r \to \varphi_{86}) \to \varphi_{119}$\lbl{A1}}
\BIC{$(r \in \mathbb{R} \to 1 \leq r \to \varphi_{86}) \to \varphi_{119}$}
\BIC{$\varphi_{119}$}
\end{ptree}
\noindent$\Pi_{54}$:
\begin{ptree}
\AXC{$\forall x\,(\mathrm{Ind}(x) \leftrightarrow \emptyset \in x \land \forall y \in x\,(S(y) \in x))$\lbl{DefInd}}
\AXC{$\forall x\,(\mathrm{Ind}(x) \leftrightarrow \emptyset \in x \land \forall y \in x\,(S(y) \in x)) \to \varphi_{15}$\lbl{A2}}
\BIC{$\varphi_{15}$}
\AXC{$\varphi_{15} \to \varphi_{148}$\lbl{A1}}
\BIC{$\varphi_{148}$}
\end{ptree}
\noindent$\Pi_{55}$:
\begin{ptree}
\AXC{$\stackrel{\Pi_{51}}{\vdots}$}\noLine
\UIC{$r \in \mathbb{R} \to 1 \leq r \to \emptyset \in \mathbb{N} \to 1 \leq {{r}^{\emptyset}}$}
\AXC{$\stackrel{\Pi_{53}}{\vdots}$}\noLine
\UIC{$\varphi_{119}$}
\BIC{$r \in \mathbb{R} \to 1 \leq r \to \varphi_{36}$}
\AXC{$\stackrel{\Pi_{54}}{\vdots}$}\noLine
\UIC{$\varphi_{148}$}
\BIC{$(r \in \mathbb{R} \to 1 \leq r \to \varphi_{50}) \to \varphi_{14} \to \varphi_{13} \to r \in \mathbb{R} \to 1 \leq r \to \varphi_{27}$}
\end{ptree}
\noindent$\Pi_{56}$:
\begin{ptree}
\AXC{$\varphi_{35} \to \forall n\,(\varphi_{35})$\lbl{A4}}
\AXC{$\stackrel{\Pi_{50}}{\vdots}$}\noLine
\UIC{$(\varphi_{35} \to \forall n\,(\varphi_{35})) \to r \in \mathbb{R} \to 1 \leq r \to \varphi_{50}$}
\BIC{$r \in \mathbb{R} \to 1 \leq r \to \varphi_{50}$}
\AXC{$\stackrel{\Pi_{55}}{\vdots}$}\noLine
\UIC{$(r \in \mathbb{R} \to 1 \leq r \to \varphi_{50}) \to \varphi_{14} \to \varphi_{13} \to r \in \mathbb{R} \to 1 \leq r \to \varphi_{27}$}
\BIC{$\varphi_{14} \to \varphi_{13} \to r \in \mathbb{R} \to 1 \leq r \to \varphi_{27}$}
\end{ptree}
\noindent$\Pi_{57}$:
\begin{ptree}
\AXC{$\varphi_{19}$\lbl{A4}}
\AXC{$\forall n\,(\varphi_{8}) \to \varphi_{21}$\lbl{A3}}
\AXC{$\varphi_{8}$\lbl{A2}}
\UIC{$\forall n\,(\varphi_{8})$}
\AXC{$\forall n\,(\varphi_{8}) \to (\forall n\,(\varphi_{8}) \to \varphi_{21}) \to \varphi_{19} \to \varphi_{14}$\lbl{A1}}
\BIC{$(\forall n\,(\varphi_{8}) \to \varphi_{21}) \to \varphi_{19} \to \varphi_{14}$}
\BIC{$\varphi_{19} \to \varphi_{14}$}
\BIC{$\varphi_{14}$}
\AXC{$\stackrel{\Pi_{56}}{\vdots}$}\noLine
\UIC{$\varphi_{14} \to \varphi_{13} \to r \in \mathbb{R} \to 1 \leq r \to \varphi_{27}$}
\BIC{$\varphi_{13} \to r \in \mathbb{R} \to 1 \leq r \to \varphi_{27}$}
\end{ptree}
\noindent$\Pi_{58}$:
\begin{ptree}
\AXC{$\varphi_{18}$\lbl{A4}}
\AXC{$\forall y\,(\varphi_{7}) \to \varphi_{20}$\lbl{A3}}
\AXC{$\varphi_{7}$\lbl{A2}}
\UIC{$\forall y\,(\varphi_{7})$}
\AXC{$\forall y\,(\varphi_{7}) \to (\forall y\,(\varphi_{7}) \to \varphi_{20}) \to \varphi_{18} \to \varphi_{13}$\lbl{A1}}
\BIC{$(\forall y\,(\varphi_{7}) \to \varphi_{20}) \to \varphi_{18} \to \varphi_{13}$}
\BIC{$\varphi_{18} \to \varphi_{13}$}
\BIC{$\varphi_{13}$}
\AXC{$\stackrel{\Pi_{57}}{\vdots}$}\noLine
\UIC{$\varphi_{13} \to r \in \mathbb{R} \to 1 \leq r \to \varphi_{27}$}
\BIC{$r \in \mathbb{R} \to 1 \leq r \to \varphi_{27}$}
\end{ptree}
\noindent$\Pi_{59}$:
\begin{ptree}
\AXC{$\forall n\,(n \in \omega \leftrightarrow \forall x\,(\mathrm{Ind}(x) \to n \in x))$\lbl{Def$\omega$}}
\AXC{$\varphi_{31}$\lbl{A2}}
\BIC{$n \in \omega \leftrightarrow \forall x\,(\mathrm{Ind}(x) \to n \in x)$}
\AXC{$r \in \mathbb{R} \to 1 \leq r \to \varphi_{47}$\lbl{A1}}
\AXC{$(r \in \mathbb{R} \to 1 \leq r \to \varphi_{47}) \to \varphi_{114}$\lbl{A1}}
\BIC{$\varphi_{114}$}
\BIC{$r \in \mathbb{R} \to 1 \leq r \to \varphi_{56}$}
\end{ptree}
\noindent$\Pi_{60}$:
\begin{ptree}
\AXC{$\forall x\,(\mathrm{Ind}(x) \to n \in x) \to \mathrm{Ind}(B_{54}) \to n \in B_{54}$\lbl{A2}}
\AXC{$\stackrel{\Pi_{59}}{\vdots}$}\noLine
\UIC{$r \in \mathbb{R} \to 1 \leq r \to \varphi_{56}$}
\AXC{$(r \in \mathbb{R} \to 1 \leq r \to \varphi_{56}) \to \varphi_{120}$\lbl{A1}}
\BIC{$\varphi_{120}$}
\BIC{$r \in \mathbb{R} \to 1 \leq r \to \varphi_{58}$}
\end{ptree}
\noindent$\Pi_{61}$:
\begin{ptree}
\AXC{$\stackrel{\Pi_{2}}{\vdots}$}\noLine
\UIC{$r \in \mathbb{R} \to 1 \leq r \to \varphi_{82}$}
\AXC{$\stackrel{\Pi_{58}}{\vdots}$}\noLine
\UIC{$r \in \mathbb{R} \to 1 \leq r \to \varphi_{27}$}
\AXC{$\stackrel{\Pi_{60}}{\vdots}$}\noLine
\UIC{$r \in \mathbb{R} \to 1 \leq r \to \varphi_{58}$}
\AXC{$(r \in \mathbb{R} \to 1 \leq r \to \varphi_{58}) \to \varphi_{141}$\lbl{A1}}
\BIC{$\varphi_{141}$}
\BIC{$(r \in \mathbb{R} \to 1 \leq r \to \varphi_{82}) \to r \in \mathbb{R} \to 1 \leq r \to \varphi_{67}$}
\BIC{$r \in \mathbb{R} \to 1 \leq r \to \varphi_{67}$}
\end{ptree}
\noindent$\Pi_{62}$:
\begin{ptree}
\AXC{$\varphi_{134}$\lbl{A3}}
\AXC{$\stackrel{\Pi_{61}}{\vdots}$}\noLine
\UIC{$r \in \mathbb{R} \to 1 \leq r \to \varphi_{67}$}
\AXC{$(r \in \mathbb{R} \to 1 \leq r \to \varphi_{67}) \to \varphi_{35} \to n \in \omega \to n \in \mathbb{N} \to 1 \leq {{r}^{n}}$\lbl{A1}}
\BIC{$\varphi_{35} \to n \in \omega \to n \in \mathbb{N} \to 1 \leq {{r}^{n}}$}
\UIC{$\forall n\,(\varphi_{35} \to n \in \omega \to n \in \mathbb{N} \to 1 \leq {{r}^{n}})$}
\AXC{$\varphi_{153}$\lbl{A1}}
\BIC{$\varphi_{134} \to (\varphi_{35} \to \forall n\,(\varphi_{35})) \to r \in \mathbb{R} \to 1 \leq r \to \varphi_{71}$}
\BIC{$(\varphi_{35} \to \forall n\,(\varphi_{35})) \to r \in \mathbb{R} \to 1 \leq r \to \varphi_{71}$}
\end{ptree}
\noindent$\Pi_{63}$:
\begin{ptree}
\AXC{$\varphi_{35} \to \forall n\,(\varphi_{35})$\lbl{A4}}
\AXC{$\stackrel{\Pi_{62}}{\vdots}$}\noLine
\UIC{$(\varphi_{35} \to \forall n\,(\varphi_{35})) \to r \in \mathbb{R} \to 1 \leq r \to \varphi_{71}$}
\BIC{$r \in \mathbb{R} \to 1 \leq r \to \varphi_{71}$}
\AXC{$(r \in \mathbb{R} \to 1 \leq r \to \varphi_{71}) \to \varphi_{99}$\lbl{A1}}
\BIC{$\varphi_{99}$}
\UIC{$\forall B_{54}\,(\varphi_{99})$}
\end{ptree}
\noindent$\Pi_{64}$:
\begin{ptree}
\AXC{$\forall B_{54}\,(\varphi_{99}) \to \varphi_{29} \to \varphi_{22}$\lbl{A3}}
\AXC{$\stackrel{\Pi_{63}}{\vdots}$}\noLine
\UIC{$\forall B_{54}\,(\varphi_{99})$}
\AXC{$\varphi_{30}$\lbl{Sep}}
\AXC{$\varphi_{30} \to \varphi_{23}$\lbl{A2}}
\BIC{$\varphi_{23}$}
\AXC{$\varphi_{102}$\lbl{A2}}
\BIC{$\exists B_{54}\,\forall n\,(n \in B_{54} \leftrightarrow n \in \omega \land (n \in \mathbb{N} \to 1 \leq {{r}^{n}}))$}
\AXC{$\varphi_{113}$\lbl{A1}}
\BIC{$r \in \mathbb{R} \to \varphi_{28}$}
\AXC{$\varphi_{154}$\lbl{A1}}
\BIC{$\forall B_{54}\,(\varphi_{99}) \to (\forall B_{54}\,(\varphi_{99}) \to \varphi_{29} \to \varphi_{22}) \to \varphi_{129}$}
\BIC{$(\forall B_{54}\,(\varphi_{99}) \to \varphi_{29} \to \varphi_{22}) \to \varphi_{129}$}
\BIC{$\varphi_{129}$}
\end{ptree}
\noindent$\Pi_{65}$:
\begin{ptree}
\AXC{$r \in \mathbb{R} \land 1 \leq r \land \neg \forall n \in \omega\,(n \in \mathbb{N} \to 1 \leq {{r}^{n}}) \to \varphi_{29}$\lbl{A4}}
\AXC{$\stackrel{\Pi_{64}}{\vdots}$}\noLine
\UIC{$\varphi_{129}$}
\BIC{$r \in \mathbb{R} \to 1 \leq r \to \forall n \in \omega\,(n \in \mathbb{N} \to 1 \leq {{r}^{n}})$}
\AXC{$(r \in \mathbb{R} \to 1 \leq r \to \forall n \in \omega\,(n \in \mathbb{N} \to 1 \leq {{r}^{n}})) \to \varphi_{115}$\lbl{A1}}
\BIC{$\varphi_{115}$}
\end{ptree}
\noindent$\Pi_{66}$:
\begin{ptree}
\AXC{$\forall n \in \omega\,(n \in \mathbb{N} \to 1 \leq {{r}^{n}}) \to n \in \omega \to n \in \mathbb{N} \to 1 \leq {{r}^{n}}$\lbl{A2}}
\AXC{$\stackrel{\Pi_{65}}{\vdots}$}\noLine
\UIC{$\varphi_{115}$}
\BIC{$r \in \mathbb{R} \to 1 \leq r \to n \in \omega \to n \in \mathbb{N} \to 1 \leq {{r}^{n}}$}
\AXC{$\varphi_{126}$\lbl{A1}}
\BIC{$(r \in \mathbb{R} \to 1 \leq r \to n \in \mathbb{N} \to n \in \omega) \to \varphi_{53}$}
\end{ptree}
\noindent$\Pi_{67}$:
\begin{ptree}
\AXC{$\forall n\,(n \in \mathbb{N} \leftrightarrow n \in \omega \land n \neq \emptyset)$\lbl{Def$\mathbb{N}$}}
\AXC{$\varphi_{40}$\lbl{A2}}
\BIC{$n \in \mathbb{N} \leftrightarrow n \in \omega \land n \neq \emptyset$}
\AXC{$\stackrel{\Pi_{1}}{\vdots}$}\noLine
\UIC{$\varphi_{44}$}
\BIC{$r \in \mathbb{R} \to 1 \leq r \to n \in \mathbb{N} \to n \in \omega$}
\AXC{$\stackrel{\Pi_{66}}{\vdots}$}\noLine
\UIC{$(r \in \mathbb{R} \to 1 \leq r \to n \in \mathbb{N} \to n \in \omega) \to \varphi_{53}$}
\BIC{$\varphi_{53}$}
\end{ptree}
\noindent$\Pi_{68}$:
\begin{ptree}
\AXC{$r \in \mathbb{R} \to 1 \leq r \to n \in \mathbb{N} \to n \in \mathbb{N}$\lbl{A1}}
\AXC{$\stackrel{\Pi_{67}}{\vdots}$}\noLine
\UIC{$\varphi_{53}$}
\BIC{$r \in \mathbb{R} \to 1 \leq r \to n \in \mathbb{N} \to 1 \leq {{r}^{n}}$}
\AXC{$\varphi_{59}$\lbl{A1}}
\BIC{$r \in \mathbb{R} \land 1 \leq r \to n \in \mathbb{N} \to 1 \leq {{r}^{n}}$}
\UIC{$\forall n\,(r \in \mathbb{R} \land 1 \leq r \to n \in \mathbb{N} \to 1 \leq {{r}^{n}})$}
\AXC{$\forall n\,(r \in \mathbb{R} \land 1 \leq r \to n \in \mathbb{N} \to 1 \leq {{r}^{n}}) \to \varphi_{72} \to \varphi_{54}$\lbl{A1}}
\BIC{$\varphi_{72} \to \varphi_{54}$}
\end{ptree}
\noindent 主木:
\begin{ptree}
\AXC{$r \in \mathbb{R} \land 1 \leq r \to \forall n\,(r \in \mathbb{R} \land 1 \leq r)$\lbl{A4}}
\AXC{$\varphi_{72}$\lbl{A3}}
\AXC{$\stackrel{\Pi_{68}}{\vdots}$}\noLine
\UIC{$\varphi_{72} \to \varphi_{54}$}
\BIC{$\varphi_{54}$}
\BIC{$r \in \mathbb{R} \to 1 \leq r \to \forall n \in \mathbb{N}\,(1 \leq {{r}^{n}})$}
\UIC{$\forall r \in \mathbb{R}\,(1 \leq r \to \forall n \in \mathbb{N}\,(1 \leq {{r}^{n}}))$}
\end{ptree}
\begin{abbrevlist}
\item $\varphi_{1} :\equiv \forall x\,\forall y\,(x = y \to y = x) \to \forall y\,(S(n) = y \to y = S(n))$
\item $\varphi_{2} :\equiv \forall x\,\forall y\,(x = y \to y = x) \to \forall y\,(1_{\omega} = y \to y = 1_{\omega})$
\item $\varphi_{3} :\equiv \forall z\,\forall w\,(z = w \to 1 = z \to 1 = w) \to \forall w\,(r = w \to 1 = r \to 1 = w)$
\item $\varphi_{4} :\equiv \forall y\,(S(n) = y \to y = S(n)) \to S(n) = 1_{\omega} \to 1_{\omega} = S(n)$
\item $\varphi_{5} :\equiv S(n) = S(\emptyset) \to S(\emptyset) = 1_{\omega} \to S(n) = 1_{\omega}$
\item $\varphi_{6} :\equiv n \in B_{54} \to S(n) \in \mathbb{N} \to n \neq \emptyset \to \mathrm{Seq}({(r)^{\bullet}})$
\item $\varphi_{7} :\equiv \forall n \in B_{54}\,(S(n) \in B_{54}) \to y \in B_{54} \to S(y) \in B_{54}$
\item $\varphi_{8} :\equiv \forall y \in B_{54}\,(S(y) \in B_{54}) \to n \in B_{54} \to S(n) \in B_{54}$
\item $\varphi_{9} :\equiv \forall z\,\forall w\,(z = w \to 1 < z \to 1 < w) \to \forall w\,(r = w \to 1 < r \to 1 < w)$
\item $\varphi_{10} :\equiv \forall x\,\forall y\,\forall z\,(x = y \to y = z \to x = z) \to \forall y\,\forall z\,(S(n) = y \to y = z \to S(n) = z)$
\item $\varphi_{11} :\equiv (1_{\omega} = S(\emptyset) \to S(\emptyset) = 1_{\omega}) \to S(\emptyset) = 1_{\omega}$
\item $\varphi_{12} :\equiv \forall y\,(1_{\omega} = y \to y = 1_{\omega}) \to 1_{\omega} = S(\emptyset) \to S(\emptyset) = 1_{\omega}$
\item $\varphi_{13} :\equiv \forall n \in B_{54}\,(S(n) \in B_{54}) \to \forall y \in B_{54}\,(S(y) \in B_{54})$
\item $\varphi_{14} :\equiv \forall y \in B_{54}\,(S(y) \in B_{54}) \to \forall n \in B_{54}\,(S(n) \in B_{54})$
\item $\varphi_{15} :\equiv \mathrm{Ind}(B_{54}) \leftrightarrow \emptyset \in B_{54} \land \forall y \in B_{54}\,(S(y) \in B_{54})$
\item $\varphi_{16} :\equiv \forall x\,\forall y\,(x = y \to y = x) \to \forall y\,({{r}^{S(n)}} = y \to y = {{r}^{S(n)}})$
\item $\varphi_{17} :\equiv \forall x\,\forall y\,(x = y \to y = x) \to \forall y\,({{r}^{1_{\omega}}} = y \to y = {{r}^{1_{\omega}}})$
\item $\varphi_{18} :\equiv \forall n \in B_{54}\,(S(n) \in B_{54}) \to \forall y\,\forall n \in B_{54}\,(S(n) \in B_{54})$
\item $\varphi_{19} :\equiv \forall y \in B_{54}\,(S(y) \in B_{54}) \to \forall n\,\forall y \in B_{54}\,(S(y) \in B_{54})$
\item $\varphi_{20} :\equiv \forall y\,\forall n \in B_{54}\,(S(n) \in B_{54}) \to \forall y \in B_{54}\,(S(y) \in B_{54})$
\item $\varphi_{21} :\equiv \forall n\,\forall y \in B_{54}\,(S(y) \in B_{54}) \to \forall n \in B_{54}\,(S(n) \in B_{54})$
\item $\varphi_{22} :\equiv \forall B_{54}\,(\neg \forall n\,(n \in B_{54} \leftrightarrow n \in \omega \land (n \in \mathbb{N} \to 1 \leq {{r}^{n}})))$
\item $\varphi_{23} :\equiv \forall A_{55}\,\exists B_{54}\,\forall n\,(n \in B_{54} \leftrightarrow n \in A_{55} \land (n \in \mathbb{N} \to 1 \leq {{r}^{n}}))$
\item $\varphi_{24} :\equiv \forall w\,(r = w \to 1 = r \to 1 = w) \to r = {{r}^{1_{\omega}}} \to 1 = r \to 1 = {{r}^{1_{\omega}}}$
\item $\varphi_{25} :\equiv S(n) \in B_{54} \leftrightarrow S(n) \in \omega \land (S(n) \in \mathbb{N} \to 1 \leq {{r}^{S(n)}})$
\item $\varphi_{26} :\equiv \forall z\,\forall w\,(z = w \to S(n) = S(z) \to S(n) = S(w)) \to \forall w\,(n = w \to S(n) = S(n) \to S(n) = S(w))$
\item $\varphi_{27} :\equiv \forall n\,(n \in B_{54} \leftrightarrow n \in \omega \land (n \in \mathbb{N} \to 1 \leq {{r}^{n}})) \to \mathrm{Ind}(B_{54})$
\item $\varphi_{28} :\equiv 1 \leq r \to \exists B_{54}\,\forall n\,(n \in B_{54} \leftrightarrow n \in \omega \land (n \in \mathbb{N} \to 1 \leq {{r}^{n}}))$
\item $\varphi_{29} :\equiv \forall B_{54}\,(r \in \mathbb{R} \land 1 \leq r \land \neg \forall n \in \omega\,(n \in \mathbb{N} \to 1 \leq {{r}^{n}}))$
\item $\varphi_{30} :\equiv \forall z_{56}\,\forall A_{55}\,\exists B_{54}\,\forall n\,(n \in B_{54} \leftrightarrow n \in A_{55} \land (n \in \mathbb{N} \to 1 \leq {{z_{56}}^{n}}))$
\item $\varphi_{31} :\equiv \forall n\,(n \in \omega \leftrightarrow \forall x\,(\mathrm{Ind}(x) \to n \in x)) \to (n \in \omega \leftrightarrow \forall x\,(\mathrm{Ind}(x) \to n \in x))$
\item $\varphi_{32} :\equiv \forall z\,\forall w\,(z = w \to 1 = z \to 1 = w) \to \forall w\,({{r}^{n}} \cdot r = w \to 1 = {{r}^{n}} \cdot r \to 1 = w)$
\item $\varphi_{33} :\equiv \forall y\,\forall z\,(S(n) = y \to y = z \to S(n) = z) \to \forall z\,(S(n) = S(\emptyset) \to S(\emptyset) = z \to S(n) = z)$
\item $\varphi_{34} :\equiv \forall w\,(n = w \to S(n) = S(n) \to S(n) = S(w)) \to n = \emptyset \to S(n) = S(n) \to S(n) = S(\emptyset)$
\item $\varphi_{35} :\equiv r \in \mathbb{R} \land 1 \leq r \land \forall n\,(n \in B_{54} \leftrightarrow n \in \omega \land (n \in \mathbb{N} \to 1 \leq {{r}^{n}}))$
\item $\varphi_{36} :\equiv \forall n\,(n \in B_{54} \leftrightarrow n \in \omega \land (n \in \mathbb{N} \to 1 \leq {{r}^{n}})) \to \emptyset \in B_{54}$
\item $\varphi_{37} :\equiv \forall w\,(r = w \to 1 < r \to 1 < w) \to r = {{r}^{1_{\omega}}} \to 1 < r \to 1 < {{r}^{1_{\omega}}}$
\item $\varphi_{38} :\equiv ({{r}^{1_{\omega}}} = r \to r = {{r}^{1_{\omega}}}) \to r \in \mathbb{R} \to 1 \leq r \to r = {{r}^{1_{\omega}}}$
\item $\varphi_{39} :\equiv \forall n\,(n \in B_{54} \leftrightarrow n \in \omega \land (n \in \mathbb{N} \to 1 \leq {{r}^{n}})) \to n \in B_{54} \to n \in B_{54}$
\item $\varphi_{40} :\equiv \forall n\,(n \in \mathbb{N} \leftrightarrow n \in \omega \land n \neq \emptyset) \to (n \in \mathbb{N} \leftrightarrow n \in \omega \land n \neq \emptyset)$
\item $\varphi_{41} :\equiv \forall y\,({{r}^{1_{\omega}}} = y \to y = {{r}^{1_{\omega}}}) \to {{r}^{1_{\omega}}} = r \to r = {{r}^{1_{\omega}}}$
\item $\varphi_{42} :\equiv \forall n\,(n \in B_{54} \leftrightarrow n \in \omega \land (n \in \mathbb{N} \to 1 \leq {{r}^{n}})) \to n \in B_{54} \to n \in \omega$
\item $\varphi_{43} :\equiv \mathrm{Seq}({(r)^{\bullet}}) \land {{r}^{1_{\omega}}} = r \land \forall n \in \mathbb{N}\,({{r}^{S(n)}} = {{r}^{n}} \cdot r)$
\item $\varphi_{44} :\equiv (n \in \mathbb{N} \leftrightarrow n \in \omega \land n \neq \emptyset) \to r \in \mathbb{R} \to 1 \leq r \to n \in \mathbb{N} \to n \in \omega$
\item $\varphi_{45} :\equiv (r = {{r}^{1_{\omega}}} \to 1 = r \to 1 = {{r}^{1_{\omega}}}) \to r \in \mathbb{R} \to 1 \leq r \to 1 \leq {{r}^{1_{\omega}}}$
\item $\varphi_{46} :\equiv \forall n\,(n \in B_{54} \leftrightarrow n \in \omega \land (n \in \mathbb{N} \to 1 \leq {{r}^{n}})) \to n \in B_{54} \to S(n) \in B_{54}$
\item $\varphi_{47} :\equiv \forall n\,(n \in B_{54} \leftrightarrow n \in \omega \land (n \in \mathbb{N} \to 1 \leq {{r}^{n}})) \to n \in \omega \to n \in \omega$
\item $\varphi_{48} :\equiv \forall n\,(n \in B_{54} \leftrightarrow n \in \omega \land (n \in \mathbb{N} \to 1 \leq {{r}^{n}})) \to n \in B_{54} \to S(n) \in \omega$
\item $\varphi_{49} :\equiv 1 \in \mathbb{R} \to r \in \mathbb{R} \to 0 \leq 1 \to 0 \leq 1 \to 1 \leq {{r}^{n}} \to 1 \leq r \to 1 \cdot 1 \leq {{r}^{n}} \cdot r$
\item $\varphi_{50} :\equiv \forall n\,(n \in B_{54} \leftrightarrow n \in \omega \land (n \in \mathbb{N} \to 1 \leq {{r}^{n}})) \to \forall n \in B_{54}\,(S(n) \in B_{54})$
\item $\varphi_{51} :\equiv \forall z\,\forall w\,(z = w \to 1 < z \to 1 < w) \to \forall w\,({{r}^{n}} \cdot r = w \to 1 < {{r}^{n}} \cdot r \to 1 < w)$
\item $\varphi_{52} :\equiv n \in B_{54} \to S(n) \in \mathbb{N} \to n \neq \emptyset \to n \in \mathbb{N} \to {{r}^{S(n)}} = {{r}^{n}} \cdot r$
\item $\varphi_{53} :\equiv (r \in \mathbb{R} \to 1 \leq r \to n \in \mathbb{N} \to n \in \mathbb{N}) \to r \in \mathbb{R} \to 1 \leq r \to n \in \mathbb{N} \to 1 \leq {{r}^{n}}$
\item $\varphi_{54} :\equiv (r \in \mathbb{R} \land 1 \leq r \to \forall n\,(r \in \mathbb{R} \land 1 \leq r)) \to r \in \mathbb{R} \to 1 \leq r \to \forall n \in \mathbb{N}\,(1 \leq {{r}^{n}})$
\item $\varphi_{55} :\equiv n \in B_{54} \to S(n) \in \mathbb{N} \to n \neq \emptyset \to \forall n \in \mathbb{N}\,({{r}^{S(n)}} = {{r}^{n}} \cdot r)$
\item $\varphi_{56} :\equiv \forall n\,(n \in B_{54} \leftrightarrow n \in \omega \land (n \in \mathbb{N} \to 1 \leq {{r}^{n}})) \to n \in \omega \to \forall x\,(\mathrm{Ind}(x) \to n \in x)$
\item $\varphi_{57} :\equiv \forall z\,\forall w\,(z = w \to 1 = {{r}^{z}} \to 1 = {{r}^{w}}) \to \forall w\,(1_{\omega} = w \to 1 = {{r}^{1_{\omega}}} \to 1 = {{r}^{w}})$
\item $\varphi_{58} :\equiv \forall n\,(n \in B_{54} \leftrightarrow n \in \omega \land (n \in \mathbb{N} \to 1 \leq {{r}^{n}})) \to n \in \omega \to \mathrm{Ind}(B_{54}) \to n \in B_{54}$
\item $\varphi_{59} :\equiv (r \in \mathbb{R} \to 1 \leq r \to n \in \mathbb{N} \to 1 \leq {{r}^{n}}) \to r \in \mathbb{R} \land 1 \leq r \to n \in \mathbb{N} \to 1 \leq {{r}^{n}}$
\item $\varphi_{60} :\equiv \forall a\,\forall n\,(\mathrm{Seq}(a) \to n \in \mathbb{N} \to a_{n} \in \mathbb{R}) \to \forall n\,(\mathrm{Seq}({(r)^{\bullet}}) \to n \in \mathbb{N} \to {{r}^{n}} \in \mathbb{R})$
\item $\varphi_{61} :\equiv \forall n \in \mathbb{N}\,({{r}^{S(n)}} = {{r}^{n}} \cdot r) \to n \in \mathbb{N} \to {{r}^{S(n)}} = {{r}^{n}} \cdot r$
\item $\varphi_{62} :\equiv \forall n\,(n \in B_{54} \leftrightarrow n \in \omega \land (n \in \mathbb{N} \to 1 \leq {{r}^{n}})) \to n \in B_{54} \to n \in \mathbb{N} \to 1 \leq {{r}^{n}}$
\item $\varphi_{63} :\equiv \forall n\,(\mathrm{Seq}({(r)^{\bullet}}) \to n \in \mathbb{N} \to {{r}^{n}} \in \mathbb{R}) \to \mathrm{Seq}({(r)^{\bullet}}) \to n \in \mathbb{N} \to {{r}^{n}} \in \mathbb{R}$
\item $\varphi_{64} :\equiv \forall y\,\forall u\,\forall v\,(1 \in \mathbb{R} \to y \in \mathbb{R} \to u \in \mathbb{R} \to v \in \mathbb{R} \to 0 \leq 1 \to 0 \leq u \to 1 \leq y \to u \leq v \to 1 \cdot u \leq y \cdot v)$
\item $\varphi_{65} :\equiv \forall w\,(1_{\omega} = w \to 1 = {{r}^{1_{\omega}}} \to 1 = {{r}^{w}}) \to 1_{\omega} = S(n) \to 1 = {{r}^{1_{\omega}}} \to 1 = {{r}^{S(n)}}$
\item $\varphi_{66} :\equiv \forall n\,(n \in B_{54} \leftrightarrow n \in \omega \land (n \in \mathbb{N} \to 1 \leq {{r}^{n}})) \to n \in B_{54} \to n \in \omega \to S(n) \in \omega$
\item $\varphi_{67} :\equiv \forall n\,(n \in B_{54} \leftrightarrow n \in \omega \land (n \in \mathbb{N} \to 1 \leq {{r}^{n}})) \to n \in \omega \to n \in \mathbb{N} \to 1 \leq {{r}^{n}}$
\item $\varphi_{68} :\equiv \forall x\,\forall y\,\forall u\,\forall v\,(x \in \mathbb{R} \to y \in \mathbb{R} \to u \in \mathbb{R} \to v \in \mathbb{R} \to 0 \leq x \to 0 \leq u \to x \leq y \to u \leq v \to x \cdot u \leq y \cdot v)$
\item $\varphi_{69} :\equiv \forall n\,(n \in \mathbb{N} \leftrightarrow n \in \omega \land n \neq \emptyset) \to (\emptyset \in \mathbb{N} \leftrightarrow \emptyset \in \omega \land \emptyset \neq \emptyset)$
\item $\varphi_{70} :\equiv \forall n\,(n \in B_{54} \leftrightarrow n \in \omega \land (n \in \mathbb{N} \to 1 \leq {{r}^{n}})) \to n \in B_{54} \to \forall y \in \omega\,(S(y) \in \omega)$
\item $\varphi_{71} :\equiv \forall n\,(n \in B_{54} \leftrightarrow n \in \omega \land (n \in \mathbb{N} \to 1 \leq {{r}^{n}})) \to \forall n \in \omega\,(n \in \mathbb{N} \to 1 \leq {{r}^{n}})$
\item $\varphi_{72} :\equiv \forall n\,(r \in \mathbb{R} \land 1 \leq r \to n \in \mathbb{N} \to 1 \leq {{r}^{n}}) \to \forall n\,(r \in \mathbb{R} \land 1 \leq r) \to \forall n \in \mathbb{N}\,(1 \leq {{r}^{n}})$
\item $\varphi_{73} :\equiv \forall n\,(n \in B_{54} \leftrightarrow n \in \omega \land (n \in \mathbb{N} \to 1 \leq {{r}^{n}})) \to n \in B_{54} \to S(n) \in \mathbb{N} \to n \neq \emptyset \to 0 \leq 1$
\item $\varphi_{74} :\equiv \forall n\,(n \in B_{54} \leftrightarrow n \in \omega \land (n \in \mathbb{N} \to 1 \leq {{r}^{n}})) \to n \in B_{54} \to S(n) \in \mathbb{N} \to 1 \leq {{r}^{S(n)}}$
\item $\varphi_{75} :\equiv \forall z\,\forall w\,(z = w \to 1 < {{r}^{z}} \to 1 < {{r}^{w}}) \to \forall w\,(1_{\omega} = w \to 1 < {{r}^{1_{\omega}}} \to 1 < {{r}^{w}})$
\item $\varphi_{76} :\equiv \forall n\,(n \in B_{54} \leftrightarrow n \in \omega \land (n \in \mathbb{N} \to 1 \leq {{r}^{n}})) \to n \in B_{54} \to S(n) \in \mathbb{N} \to n \neq \emptyset \to 1 \cdot 1 = 1$
\item $\varphi_{77} :\equiv \forall n\,(n \in B_{54} \leftrightarrow n \in \omega \land (n \in \mathbb{N} \to 1 \leq {{r}^{n}})) \to n \in B_{54} \to S(n) \in \mathbb{N} \to n = \emptyset \to S(n) = 1_{\omega}$
\item $\varphi_{78} :\equiv \forall n\,(n \in B_{54} \leftrightarrow n \in \omega \land (n \in \mathbb{N} \to 1 \leq {{r}^{n}})) \to n \in B_{54} \to S(n) \in \mathbb{N} \to n = \emptyset \to 1_{\omega} = S(n)$
\item $\varphi_{79} :\equiv \forall n\,(n \in B_{54} \leftrightarrow n \in \omega \land (n \in \mathbb{N} \to 1 \leq {{r}^{n}})) \to n \in B_{54} \to S(n) \in \mathbb{N} \to n \neq \emptyset \to n \in \mathbb{N}$
\item $\varphi_{80} :\equiv \forall n\,(n \in B_{54} \leftrightarrow n \in \omega \land (n \in \mathbb{N} \to 1 \leq {{r}^{n}})) \to n \in B_{54} \to S(n) \in \mathbb{N} \to n = \emptyset \to n = \emptyset$
\item $\varphi_{81} :\equiv \forall n\,(n \in B_{54} \leftrightarrow n \in \omega \land (n \in \mathbb{N} \to 1 \leq {{r}^{n}})) \to n \in B_{54} \to S(n) \in \mathbb{N} \to n \neq \emptyset \to r \in \mathbb{R}$
\item $\varphi_{82} :\equiv \forall n\,(n \in B_{54} \leftrightarrow n \in \omega \land (n \in \mathbb{N} \to 1 \leq {{r}^{n}})) \to (n \in B_{54} \leftrightarrow n \in \omega \land (n \in \mathbb{N} \to 1 \leq {{r}^{n}}))$
\item $\varphi_{83} :\equiv \forall w\,(1_{\omega} = w \to 1 < {{r}^{1_{\omega}}} \to 1 < {{r}^{w}}) \to 1_{\omega} = S(n) \to 1 < {{r}^{1_{\omega}}} \to 1 < {{r}^{S(n)}}$
\item $\varphi_{84} :\equiv \forall n\,(n \in B_{54} \leftrightarrow n \in \omega \land (n \in \mathbb{N} \to 1 \leq {{r}^{n}})) \to n \in B_{54} \to S(n) \in \mathbb{N} \to n \neq \emptyset \to n \neq \emptyset$
\item $\varphi_{85} :\equiv \forall n\,(n \in B_{54} \leftrightarrow n \in \omega \land (n \in \mathbb{N} \to 1 \leq {{r}^{n}})) \to n \in B_{54} \to n \in \omega \land (n \in \mathbb{N} \to 1 \leq {{r}^{n}})$
\item $\varphi_{86} :\equiv \forall n\,(n \in B_{54} \leftrightarrow n \in \omega \land (n \in \mathbb{N} \to 1 \leq {{r}^{n}})) \to \emptyset \in \omega \land \forall y \in \omega\,(S(y) \in \omega)$
\item $\varphi_{87} :\equiv \forall y\,({{r}^{S(n)}} = y \to y = {{r}^{S(n)}}) \to {{r}^{S(n)}} = {{r}^{n}} \cdot r \to {{r}^{n}} \cdot r = {{r}^{S(n)}}$
\item $\varphi_{88} :\equiv \forall n\,(n \in B_{54} \leftrightarrow n \in \omega \land (n \in \mathbb{N} \to 1 \leq {{r}^{n}})) \to \forall n\,(n \in B_{54} \leftrightarrow n \in \omega \land (n \in \mathbb{N} \to 1 \leq {{r}^{n}}))$
\item $\varphi_{89} :\equiv \forall n\,(n \in B_{54} \leftrightarrow n \in \omega \land (n \in \mathbb{N} \to 1 \leq {{r}^{n}})) \to n \in B_{54} \to S(n) \in \mathbb{N} \to n \neq \emptyset \to 1 \leq {{r}^{n}}$
\item $\varphi_{90} :\equiv \forall n\,(n \in B_{54} \leftrightarrow n \in \omega \land (n \in \mathbb{N} \to 1 \leq {{r}^{n}})) \to n \in B_{54} \to S(n) \in \mathbb{N} \to n = \emptyset \to S(n) = S(\emptyset)$
\item $\varphi_{91} :\equiv 1 \leq r \to \forall n\,(n \in B_{54} \leftrightarrow n \in \omega \land (n \in \mathbb{N} \to 1 \leq {{r}^{n}})) \to \varphi_{6}$
\item $\varphi_{92} :\equiv \forall n\,(n \in B_{54} \leftrightarrow n \in \omega \land (n \in \mathbb{N} \to 1 \leq {{r}^{n}})) \to n \in B_{54} \to S(n) \in \mathbb{N} \to n = \emptyset \to 1 \leq {{r}^{S(n)}}$
\item $\varphi_{93} :\equiv \forall z\,\forall w\,(z = w \to z = {{r}^{n}} \cdot r \to w = {{r}^{n}} \cdot r) \to \forall w\,(1 \cdot 1 = w \to 1 \cdot 1 = {{r}^{n}} \cdot r \to w = {{r}^{n}} \cdot r)$
\item $\varphi_{94} :\equiv \forall n\,(n \in B_{54} \leftrightarrow n \in \omega \land (n \in \mathbb{N} \to 1 \leq {{r}^{n}})) \to n \in B_{54} \to S(n) \in \mathbb{N} \to n \neq \emptyset \to 1 \leq {{r}^{S(n)}}$
\item $\varphi_{95} :\equiv \forall n\,(n \in B_{54} \leftrightarrow n \in \omega \land (n \in \mathbb{N} \to 1 \leq {{r}^{n}})) \to n \in B_{54} \to S(n) \in \mathbb{N} \to n \neq \emptyset \to 1 \leq {{r}^{n}} \cdot r$
\item $\varphi_{96} :\equiv \forall w\,(1 \cdot 1 = w \to 1 \cdot 1 = {{r}^{n}} \cdot r \to w = {{r}^{n}} \cdot r) \to 1 \cdot 1 = 1 \to 1 \cdot 1 = {{r}^{n}} \cdot r \to 1 = {{r}^{n}} \cdot r$
\item $\varphi_{97} :\equiv \forall x\,(\mathrm{Ind}(x) \leftrightarrow \emptyset \in x \land \forall y \in x\,(S(y) \in x)) \to (\mathrm{Ind}(\omega) \leftrightarrow \emptyset \in \omega \land \forall y \in \omega\,(S(y) \in \omega))$
\item $\varphi_{98} :\equiv \forall n\,(n \in B_{54} \leftrightarrow n \in \omega \land (n \in \mathbb{N} \to 1 \leq {{r}^{n}})) \to n \in B_{54} \to S(n) \in \mathbb{N} \to n \neq \emptyset \to {{r}^{n}} \in \mathbb{R}$
\item $\varphi_{99} :\equiv r \in \mathbb{R} \land 1 \leq r \land \neg \forall n \in \omega\,(n \in \mathbb{N} \to 1 \leq {{r}^{n}}) \to \neg \forall n\,(n \in B_{54} \leftrightarrow n \in \omega \land (n \in \mathbb{N} \to 1 \leq {{r}^{n}}))$
\item $\varphi_{100} :\equiv \forall n\,(n \in B_{54} \leftrightarrow n \in \omega \land (n \in \mathbb{N} \to 1 \leq {{r}^{n}})) \to n \in B_{54} \to S(n) \in \mathbb{N} \to n \neq \emptyset \to 1 \cdot 1 \leq {{r}^{n}} \cdot r$
\item $\varphi_{101} :\equiv \forall v\,(1 \in \mathbb{R} \to {{r}^{n}} \in \mathbb{R} \to 1 \in \mathbb{R} \to v \in \mathbb{R} \to 0 \leq 1 \to 0 \leq 1 \to 1 \leq {{r}^{n}} \to 1 \leq v \to 1 \cdot 1 \leq {{r}^{n}} \cdot v)$
\item $\varphi_{102} :\equiv \varphi_{23} \to \exists B_{54}\,\forall n\,(n \in B_{54} \leftrightarrow n \in \omega \land (n \in \mathbb{N} \to 1 \leq {{r}^{n}}))$
\item $\varphi_{103} :\equiv \forall w\,({{r}^{n}} \cdot r = w \to 1 = {{r}^{n}} \cdot r \to 1 = w) \to {{r}^{n}} \cdot r = {{r}^{S(n)}} \to 1 = {{r}^{n}} \cdot r \to 1 = {{r}^{S(n)}}$
\item $\varphi_{104} :\equiv \forall z\,\forall w\,(z = w \to z < {{r}^{n}} \cdot r \to w < {{r}^{n}} \cdot r) \to \forall w\,(1 \cdot 1 = w \to 1 \cdot 1 < {{r}^{n}} \cdot r \to w < {{r}^{n}} \cdot r)$
\item $\varphi_{105} :\equiv \forall u\,\forall v\,(1 \in \mathbb{R} \to {{r}^{n}} \in \mathbb{R} \to u \in \mathbb{R} \to v \in \mathbb{R} \to 0 \leq 1 \to 0 \leq u \to 1 \leq {{r}^{n}} \to u \leq v \to 1 \cdot u \leq {{r}^{n}} \cdot v)$
\item $\varphi_{106} :\equiv 1 \leq r \to \forall n\,(n \in B_{54} \leftrightarrow n \in \omega \land (n \in \mathbb{N} \to 1 \leq {{r}^{n}})) \to \varphi_{25}$
\item $\varphi_{107} :\equiv \forall w\,(1 \cdot 1 = w \to 1 \cdot 1 < {{r}^{n}} \cdot r \to w < {{r}^{n}} \cdot r) \to 1 \cdot 1 = 1 \to 1 \cdot 1 < {{r}^{n}} \cdot r \to 1 < {{r}^{n}} \cdot r$
\item $\varphi_{108} :\equiv \forall n\,(n \in B_{54} \leftrightarrow n \in \omega \land (n \in \mathbb{N} \to 1 \leq {{r}^{n}})) \to n \in B_{54} \to S(n) \in \mathbb{N} \to n \neq \emptyset \to {{r}^{S(n)}} = {{r}^{n}} \cdot r$
\item $\varphi_{109} :\equiv \forall n\,(n \in B_{54} \leftrightarrow n \in \omega \land (n \in \mathbb{N} \to 1 \leq {{r}^{n}})) \to n \in B_{54} \to S(n) \in \mathbb{N} \to n \neq \emptyset \to {{r}^{n}} \cdot r = {{r}^{S(n)}}$
\item $\varphi_{110} :\equiv \forall n\,(n \in B_{54} \leftrightarrow n \in \omega \land (n \in \mathbb{N} \to 1 \leq {{r}^{n}})) \to (\emptyset \in B_{54} \leftrightarrow \emptyset \in \omega \land (\emptyset \in \mathbb{N} \to 1 \leq {{r}^{\emptyset}}))$
\item $\varphi_{111} :\equiv \forall w\,({{r}^{n}} \cdot r = w \to 1 < {{r}^{n}} \cdot r \to 1 < w) \to {{r}^{n}} \cdot r = {{r}^{S(n)}} \to 1 < {{r}^{n}} \cdot r \to 1 < {{r}^{S(n)}}$
\item $\varphi_{112} :\equiv (\emptyset \in \mathbb{N} \leftrightarrow \emptyset \in \omega \land \emptyset \neq \emptyset) \to \emptyset = \emptyset \to r \in \mathbb{R} \to 1 \leq r \to \emptyset \in \mathbb{N} \to 1 \leq {{r}^{\emptyset}}$
\item $\varphi_{113} :\equiv \exists B_{54}\,\forall n\,(n \in B_{54} \leftrightarrow n \in \omega \land (n \in \mathbb{N} \to 1 \leq {{r}^{n}})) \to r \in \mathbb{R} \to \varphi_{28}$
\item $\varphi_{114} :\equiv (n \in \omega \leftrightarrow \forall x\,(\mathrm{Ind}(x) \to n \in x)) \to r \in \mathbb{R} \to 1 \leq r \to \varphi_{56}$
\item $\varphi_{115} :\equiv (\forall n \in \omega\,(n \in \mathbb{N} \to 1 \leq {{r}^{n}}) \to n \in \omega \to n \in \mathbb{N} \to 1 \leq {{r}^{n}}) \to r \in \mathbb{R} \to 1 \leq r \to n \in \omega \to n \in \mathbb{N} \to 1 \leq {{r}^{n}}$
\item $\varphi_{116} :\equiv (S(n) = 1_{\omega} \to 1_{\omega} = S(n)) \to r \in \mathbb{R} \to 1 \leq r \to \varphi_{78}$
\item $\varphi_{117} :\equiv 1 \leq r \to \forall n\,(n \in B_{54} \leftrightarrow n \in \omega \land (n \in \mathbb{N} \to 1 \leq {{r}^{n}})) \to \varphi_{52}$
\item $\varphi_{118} :\equiv 1 \leq r \to \forall n\,(n \in B_{54} \leftrightarrow n \in \omega \land (n \in \mathbb{N} \to 1 \leq {{r}^{n}})) \to \varphi_{55}$
\item $\varphi_{119} :\equiv (r \in \mathbb{R} \to 1 \leq r \to \emptyset \in \mathbb{N} \to 1 \leq {{r}^{\emptyset}}) \to r \in \mathbb{R} \to 1 \leq r \to \varphi_{36}$
\item $\varphi_{120} :\equiv (\forall x\,(\mathrm{Ind}(x) \to n \in x) \to \mathrm{Ind}(B_{54}) \to n \in B_{54}) \to r \in \mathbb{R} \to 1 \leq r \to \varphi_{58}$
\item $\varphi_{121} :\equiv (r \in \mathbb{R} \to 1 \leq r \to 1 \leq r) \to (r = {{r}^{1_{\omega}}} \to 1 < r \to 1 < {{r}^{1_{\omega}}}) \to \varphi_{45}$
\item $\varphi_{122} :\equiv (n = \emptyset \to S(n) = S(n) \to S(n) = S(\emptyset)) \to r \in \mathbb{R} \to 1 \leq r \to \varphi_{90}$
\item $\varphi_{123} :\equiv \forall n\,(n \in B_{54} \leftrightarrow n \in \omega \land (n \in \mathbb{N} \to 1 \leq {{r}^{n}})) \to n \in B_{54} \to S(n) \in \mathbb{N} \to n \neq \emptyset \to \varphi_{43}$
\item $\varphi_{124} :\equiv (\forall y \in \omega\,(S(y) \in \omega) \to n \in \omega \to S(n) \in \omega) \to r \in \mathbb{R} \to 1 \leq r \to \varphi_{66}$
\item $\varphi_{125} :\equiv (1_{\omega} = S(n) \to 1 = {{r}^{1_{\omega}}} \to 1 = {{r}^{S(n)}}) \to r \in \mathbb{R} \to 1 \leq r \to \varphi_{92}$
\item $\varphi_{126} :\equiv (r \in \mathbb{R} \to 1 \leq r \to n \in \omega \to n \in \mathbb{N} \to 1 \leq {{r}^{n}}) \to (r \in \mathbb{R} \to 1 \leq r \to n \in \mathbb{N} \to n \in \omega) \to \varphi_{53}$
\item $\varphi_{127} :\equiv (1 \cdot 1 = 1 \to 1 \cdot 1 = {{r}^{n}} \cdot r \to 1 = {{r}^{n}} \cdot r) \to r \in \mathbb{R} \to 1 \leq r \to \varphi_{95}$
\item $\varphi_{128} :\equiv (\mathrm{Ind}(\omega) \leftrightarrow \emptyset \in \omega \land \forall y \in \omega\,(S(y) \in \omega)) \to r \in \mathbb{R} \to 1 \leq r \to \varphi_{86}$
\item $\varphi_{129} :\equiv (r \in \mathbb{R} \land 1 \leq r \land \neg \forall n \in \omega\,(n \in \mathbb{N} \to 1 \leq {{r}^{n}}) \to \varphi_{29}) \to r \in \mathbb{R} \to 1 \leq r \to \forall n \in \omega\,(n \in \mathbb{N} \to 1 \leq {{r}^{n}})$
\item $\varphi_{130} :\equiv ({{r}^{n}} \cdot r = {{r}^{S(n)}} \to 1 = {{r}^{n}} \cdot r \to 1 = {{r}^{S(n)}}) \to r \in \mathbb{R} \to 1 \leq r \to \varphi_{94}$
\item $\varphi_{131} :\equiv \forall n\,(\varphi_{35} \to n \in B_{54} \to S(n) \in B_{54}) \to \forall n\,(\varphi_{35}) \to \forall n \in B_{54}\,(S(n) \in B_{54})$
\item $\varphi_{132} :\equiv ({{r}^{S(n)}} = {{r}^{n}} \cdot r \to {{r}^{n}} \cdot r = {{r}^{S(n)}}) \to r \in \mathbb{R} \to 1 \leq r \to \varphi_{109}$
\item $\varphi_{133} :\equiv (r \in \mathbb{R} \to \varphi_{43}) \to (r \in \mathbb{R} \to 1 \leq r \to r \in \mathbb{R}) \to r \in \mathbb{R} \to 1 \leq r \to \varphi_{43}$
\item $\varphi_{134} :\equiv \forall n\,(\varphi_{35} \to n \in \omega \to n \in \mathbb{N} \to 1 \leq {{r}^{n}}) \to \forall n\,(\varphi_{35}) \to \forall n \in \omega\,(n \in \mathbb{N} \to 1 \leq {{r}^{n}})$
\item $\varphi_{135} :\equiv (\forall n\,(n \in B_{54} \leftrightarrow n \in \omega \land (n \in \mathbb{N} \to 1 \leq {{r}^{n}})) \to \varphi_{25}) \to r \in \mathbb{R} \to \varphi_{106}$
\item $\varphi_{136} :\equiv (r \in \mathbb{R} \to 1 \leq r \to 1 \leq {{r}^{1_{\omega}}}) \to (1_{\omega} = S(n) \to 1 < {{r}^{1_{\omega}}} \to 1 < {{r}^{S(n)}}) \to \varphi_{125}$
\item $\varphi_{137} :\equiv (r \in \mathbb{R} \to 1 \leq r \to \varphi_{89}) \to (r \in \mathbb{R} \to 1 \leq r \to 1 \leq r) \to r \in \mathbb{R} \to 1 \leq r \to \varphi_{100}$
\item $\varphi_{138} :\equiv (r \in \mathbb{R} \to 1 \leq r \to \varphi_{66}) \to (r \in \mathbb{R} \to 1 \leq r \to \varphi_{42}) \to r \in \mathbb{R} \to 1 \leq r \to \varphi_{48}$
\item $\varphi_{139} :\equiv (r \in \mathbb{R} \to 1 \leq r \to \varphi_{90}) \to S(\emptyset) = 1_{\omega} \to \varphi_{5} \to r \in \mathbb{R} \to 1 \leq r \to \varphi_{77}$
\item $\varphi_{140} :\equiv (r \in \mathbb{R} \to 1 \leq r \to \varphi_{48}) \to (r \in \mathbb{R} \to 1 \leq r \to \varphi_{74}) \to r \in \mathbb{R} \to 1 \leq r \to \varphi_{46}$
\item $\varphi_{141} :\equiv (r \in \mathbb{R} \to 1 \leq r \to \varphi_{27}) \to (r \in \mathbb{R} \to 1 \leq r \to \varphi_{82}) \to r \in \mathbb{R} \to 1 \leq r \to \varphi_{67}$
\item $\varphi_{142} :\equiv (r \in \mathbb{R} \to 1 \leq r \to \varphi_{82}) \to (r \in \mathbb{R} \to 1 \leq r \to \varphi_{39}) \to r \in \mathbb{R} \to 1 \leq r \to \varphi_{85}$
\item $\varphi_{143} :\equiv (r \in \mathbb{R} \to 1 \leq r \to \varphi_{42}) \to (r \in \mathbb{R} \to 1 \leq r \to \varphi_{84}) \to r \in \mathbb{R} \to 1 \leq r \to \varphi_{79}$
\item $\varphi_{144} :\equiv (r \in \mathbb{R} \to 1 \leq r \to \varphi_{62}) \to (r \in \mathbb{R} \to 1 \leq r \to \varphi_{79}) \to r \in \mathbb{R} \to 1 \leq r \to \varphi_{89}$
\item $\varphi_{145} :\equiv (r \in \mathbb{R} \to 1 \leq r \to \varphi_{92}) \to (r \in \mathbb{R} \to 1 \leq r \to \varphi_{94}) \to r \in \mathbb{R} \to 1 \leq r \to \varphi_{74}$
\item $\varphi_{146} :\equiv (r \in \mathbb{R} \to \varphi_{91}) \to (r \in \mathbb{R} \to 1 \leq r \to \varphi_{79}) \to r \in \mathbb{R} \to 1 \leq r \to \varphi_{98}$
\item $\varphi_{147} :\equiv (r \in \mathbb{R} \to 1 \leq r \to \varphi_{100}) \to (1 \cdot 1 = 1 \to 1 \cdot 1 < {{r}^{n}} \cdot r \to 1 < {{r}^{n}} \cdot r) \to \varphi_{127}$
\item $\varphi_{148} :\equiv (r \in \mathbb{R} \to 1 \leq r \to \varphi_{36}) \to (r \in \mathbb{R} \to 1 \leq r \to \varphi_{50}) \to \varphi_{14} \to \varphi_{13} \to r \in \mathbb{R} \to 1 \leq r \to \varphi_{27}$
\item $\varphi_{149} :\equiv (r \in \mathbb{R} \to \varphi_{43}) \to (r \in \mathbb{R} \to 1 \leq r \to \varphi_{81}) \to r \in \mathbb{R} \to 1 \leq r \to \varphi_{123}$
\item $\varphi_{150} :\equiv (r \in \mathbb{R} \to 1 \leq r \to \varphi_{95}) \to ({{r}^{n}} \cdot r = {{r}^{S(n)}} \to 1 < {{r}^{n}} \cdot r \to 1 < {{r}^{S(n)}}) \to \varphi_{130}$
\item $\varphi_{151} :\equiv (r \in \mathbb{R} \to \varphi_{117}) \to (r \in \mathbb{R} \to 1 \leq r \to \varphi_{79}) \to r \in \mathbb{R} \to 1 \leq r \to \varphi_{108}$
\item $\varphi_{152} :\equiv \forall n\,(\varphi_{35} \to n \in B_{54} \to S(n) \in B_{54}) \to \varphi_{131} \to (\varphi_{35} \to \forall n\,(\varphi_{35})) \to r \in \mathbb{R} \to 1 \leq r \to \varphi_{50}$
\item $\varphi_{153} :\equiv \forall n\,(\varphi_{35} \to n \in \omega \to n \in \mathbb{N} \to 1 \leq {{r}^{n}}) \to \varphi_{134} \to (\varphi_{35} \to \forall n\,(\varphi_{35})) \to r \in \mathbb{R} \to 1 \leq r \to \varphi_{71}$
\item $\varphi_{154} :\equiv (r \in \mathbb{R} \to \varphi_{28}) \to \forall B_{54}\,(\varphi_{99}) \to (\forall B_{54}\,(\varphi_{99}) \to \varphi_{29} \to \varphi_{22}) \to \varphi_{129}$
\item $\varphi_{155} :\equiv (r \in \mathbb{R} \to 1 \leq r \to \varphi_{81}) \to (r \in \mathbb{R} \to 1 \leq r \to \varphi_{73}) \to (r \in \mathbb{R} \to 1 \leq r \to \varphi_{73}) \to \varphi_{137}$
\item $\varphi_{156} :\equiv (1 \in \mathbb{R} \to {{r}^{n}} \in \mathbb{R} \to \varphi_{49}) \to 1 \in \mathbb{R} \to (r \in \mathbb{R} \to 1 \leq r \to \varphi_{98}) \to 1 \in \mathbb{R} \to \varphi_{155}$
\end{abbrevlist}
\end{fullproof}

% 自動生成 (prover/build.js): GeoSerDvg  ノード数 294
\begin{fullproof}{命題4.3(2)}{$\forall r \in \mathbb{R}\,(1 \leq \lvert r \rvert \to \neg \exists s\,(\mathrm{Conv}(\Sigma {(r)^{\bullet}},s)))$}{294}\phantomsection\label{tree:GeoSerDvg}
\noindent$\Pi_{1}$:
\begin{ptree}
\AXC{$r \in \mathbb{R} \to 1 \leq \lvert r \rvert \to \varphi_{109}$\lbl{A1}}
\AXC{$(r \in \mathbb{R} \to 1 \leq \lvert r \rvert \to \varphi_{109}) \to r \in \mathbb{R} \to 1 \leq \lvert r \rvert \to \varphi_{69}$\lbl{A1}}
\BIC{$r \in \mathbb{R} \to 1 \leq \lvert r \rvert \to \varphi_{69}$}
\end{ptree}
\noindent$\Pi_{2}$:
\begin{ptree}
\AXC{$r \in \mathbb{R} \to 1 \leq \lvert r \rvert \to \varphi_{72}$\lbl{A1}}
\AXC{$\forall x \in \mathbb{R}\,(\lvert x \rvert \in \mathbb{R})$\lbl{[A5]}}
\AXC{$\forall x \in \mathbb{R}\,(\lvert x \rvert \in \mathbb{R}) \to r \in \mathbb{R} \to \lvert r \rvert \in \mathbb{R}$\lbl{A2}}
\BIC{$r \in \mathbb{R} \to \lvert r \rvert \in \mathbb{R}$}
\AXC{$\varphi_{126}$\lbl{A1}}
\BIC{$(r \in \mathbb{R} \to 1 \leq \lvert r \rvert \to \varphi_{72}) \to r \in \mathbb{R} \to 1 \leq \lvert r \rvert \to \varphi_{77}$}
\BIC{$r \in \mathbb{R} \to 1 \leq \lvert r \rvert \to \varphi_{77}$}
\end{ptree}
\noindent$\Pi_{3}$:
\begin{ptree}
\AXC{$\forall r \in \mathbb{R}\,(1 \leq r \to \forall n \in \mathbb{N}\,(1 \leq {{r}^{n}}))$\lbl{[P8]}}
\AXC{$\varphi_{61}$\lbl{A2}}
\BIC{$\lvert r \rvert \in \mathbb{R} \to 1 \leq \lvert r \rvert \to \forall n \in \mathbb{N}\,(1 \leq {{\lvert r \rvert}^{n}})$}
\AXC{$(\lvert r \rvert \in \mathbb{R} \to 1 \leq \lvert r \rvert \to \forall n \in \mathbb{N}\,(1 \leq {{\lvert r \rvert}^{n}})) \to \varphi_{127}$\lbl{A1}}
\BIC{$\varphi_{127}$}
\end{ptree}
\noindent$\Pi_{4}$:
\begin{ptree}
\AXC{$r \in \mathbb{R} \to 1 \leq \lvert r \rvert \to 1 \leq \lvert r \rvert$\lbl{A1}}
\AXC{$\stackrel{\Pi_{2}}{\vdots}$}\noLine
\UIC{$r \in \mathbb{R} \to 1 \leq \lvert r \rvert \to \varphi_{77}$}
\AXC{$\stackrel{\Pi_{3}}{\vdots}$}\noLine
\UIC{$\varphi_{127}$}
\BIC{$(r \in \mathbb{R} \to 1 \leq \lvert r \rvert \to 1 \leq \lvert r \rvert) \to \varphi_{106}$}
\BIC{$\varphi_{106}$}
\AXC{$\varphi_{132}$\lbl{A1}}
\BIC{$(\forall n \in \mathbb{N}\,(1 \leq {{\lvert r \rvert}^{n}}) \to N \in \mathbb{N} \to 1 \leq {{\lvert r \rvert}^{N}}) \to \varphi_{104}$}
\end{ptree}
\noindent$\Pi_{5}$:
\begin{ptree}
\AXC{$\forall n \in \mathbb{N}\,(1 \leq {{\lvert r \rvert}^{n}}) \to N \in \mathbb{N} \to 1 \leq {{\lvert r \rvert}^{N}}$\lbl{A2}}
\AXC{$\stackrel{\Pi_{4}}{\vdots}$}\noLine
\UIC{$(\forall n \in \mathbb{N}\,(1 \leq {{\lvert r \rvert}^{n}}) \to N \in \mathbb{N} \to 1 \leq {{\lvert r \rvert}^{N}}) \to \varphi_{104}$}
\BIC{$\varphi_{104}$}
\end{ptree}
\noindent$\Pi_{6}$:
\begin{ptree}
\AXC{$\stackrel{\Pi_{1}}{\vdots}$}\noLine
\UIC{$r \in \mathbb{R} \to 1 \leq \lvert r \rvert \to \varphi_{69}$}
\AXC{$\stackrel{\Pi_{5}}{\vdots}$}\noLine
\UIC{$\varphi_{104}$}
\AXC{$\varphi_{104} \to (r \in \mathbb{R} \to 1 \leq \lvert r \rvert \to \varphi_{69}) \to r \in \mathbb{R} \to 1 \leq \lvert r \rvert \to \varphi_{81}$\lbl{A1}}
\BIC{$(r \in \mathbb{R} \to 1 \leq \lvert r \rvert \to \varphi_{69}) \to r \in \mathbb{R} \to 1 \leq \lvert r \rvert \to \varphi_{81}$}
\BIC{$r \in \mathbb{R} \to 1 \leq \lvert r \rvert \to \varphi_{81}$}
\end{ptree}
\noindent$\Pi_{7}$:
\begin{ptree}
\AXC{$\forall x\,\forall y\,(x \leq y \to \neg y < x)$\lbl{[O7]}}
\AXC{$\forall x\,\forall y\,(x \leq y \to \neg y < x) \to \forall y\,(1 \leq y \to \neg y < 1)$\lbl{A2}}
\BIC{$\forall y\,(1 \leq y \to \neg y < 1)$}
\end{ptree}
\noindent$\Pi_{8}$:
\begin{ptree}
\AXC{$\stackrel{\Pi_{6}}{\vdots}$}\noLine
\UIC{$r \in \mathbb{R} \to 1 \leq \lvert r \rvert \to \varphi_{81}$}
\AXC{$\stackrel{\Pi_{7}}{\vdots}$}\noLine
\UIC{$\forall y\,(1 \leq y \to \neg y < 1)$}
\AXC{$\forall y\,(1 \leq y \to \neg y < 1) \to 1 \leq {{\lvert r \rvert}^{N}} \to \neg {{\lvert r \rvert}^{N}} < 1$\lbl{A2}}
\BIC{$1 \leq {{\lvert r \rvert}^{N}} \to \neg {{\lvert r \rvert}^{N}} < 1$}
\AXC{$\varphi_{129}$\lbl{A1}}
\BIC{$(r \in \mathbb{R} \to 1 \leq \lvert r \rvert \to \varphi_{81}) \to r \in \mathbb{R} \to 1 \leq \lvert r \rvert \to \varphi_{88}$}
\BIC{$r \in \mathbb{R} \to 1 \leq \lvert r \rvert \to \varphi_{88}$}
\end{ptree}
\noindent$\Pi_{9}$:
\begin{ptree}
\AXC{$z = w \to z < 1 \to w < 1$\lbl{A6}}
\UIC{$\forall w\,(z = w \to z < 1 \to w < 1)$}
\UIC{$\forall z\,\forall w\,(z = w \to z < 1 \to w < 1)$}
\AXC{$\forall z\,\forall w\,(z = w \to z < 1 \to w < 1) \to \varphi_{32}$\lbl{A2}}
\BIC{$\varphi_{32}$}
\AXC{$\varphi_{32} \to \varphi_{51}$\lbl{A2}}
\BIC{$\varphi_{51}$}
\end{ptree}
\noindent$\Pi_{10}$:
\begin{ptree}
\AXC{$x = x$\lbl{A5}}
\UIC{$\forall x\,(x = x)$}
\AXC{$\forall x\,(x = x) \to N = N$\lbl{A2}}
\BIC{$N = N$}
\AXC{$N = N \to r \in \mathbb{R} \to 1 \leq \lvert r \rvert \to \varphi_{63}$\lbl{A1}}
\BIC{$r \in \mathbb{R} \to 1 \leq \lvert r \rvert \to \varphi_{63}$}
\end{ptree}
\noindent$\Pi_{11}$:
\begin{ptree}
\AXC{$r \in \mathbb{R} \to 1 \leq \lvert r \rvert \to \varphi_{109}$\lbl{A1}}
\AXC{$(r \in \mathbb{R} \to 1 \leq \lvert r \rvert \to \varphi_{109}) \to r \in \mathbb{R} \to 1 \leq \lvert r \rvert \to \varphi_{103}$\lbl{A1}}
\BIC{$r \in \mathbb{R} \to 1 \leq \lvert r \rvert \to \varphi_{103}$}
\AXC{$(r \in \mathbb{R} \to 1 \leq \lvert r \rvert \to \varphi_{103}) \to \varphi_{53} \to \varphi_{117}$\lbl{A1}}
\BIC{$\varphi_{53} \to \varphi_{117}$}
\end{ptree}
\noindent$\Pi_{12}$:
\begin{ptree}
\AXC{$\stackrel{\Pi_{10}}{\vdots}$}\noLine
\UIC{$r \in \mathbb{R} \to 1 \leq \lvert r \rvert \to \varphi_{63}$}
\AXC{$\stackrel{\Pi_{1}}{\vdots}$}\noLine
\UIC{$r \in \mathbb{R} \to 1 \leq \lvert r \rvert \to \varphi_{69}$}
\AXC{$\varphi_{53}$\lbl{A2}}
\AXC{$\stackrel{\Pi_{11}}{\vdots}$}\noLine
\UIC{$\varphi_{53} \to \varphi_{117}$}
\BIC{$\varphi_{117}$}
\AXC{$\varphi_{117} \to \varphi_{135}$\lbl{A1}}
\BIC{$\varphi_{135}$}
\BIC{$(r \in \mathbb{R} \to 1 \leq \lvert r \rvert \to \varphi_{63}) \to \varphi_{105}$}
\BIC{$\varphi_{105}$}
\end{ptree}
\noindent$\Pi_{13}$:
\begin{ptree}
\AXC{$y = z \to x = y \to x = z$\lbl{A6}}
\AXC{$(y = z \to x = y \to x = z) \to x = y \to y = z \to x = z$\lbl{A1}}
\BIC{$x = y \to y = z \to x = z$}
\UIC{$\forall z\,(x = y \to y = z \to x = z)$}
\UIC{$\forall y\,\forall z\,(x = y \to y = z \to x = z)$}
\UIC{$\forall x\,\forall y\,\forall z\,(x = y \to y = z \to x = z)$}
\AXC{$\forall x\,\forall y\,\forall z\,(x = y \to y = z \to x = z) \to \varphi_{31}$\lbl{A2}}
\BIC{$\varphi_{31}$}
\end{ptree}
\noindent$\Pi_{14}$:
\begin{ptree}
\AXC{$\stackrel{\Pi_{13}}{\vdots}$}\noLine
\UIC{$\varphi_{31}$}
\AXC{$\varphi_{31} \to \varphi_{56}$\lbl{A2}}
\BIC{$\varphi_{56}$}
\AXC{$\varphi_{56} \to \lvert \lvert {(r)^{\bullet}} \rvert(N) - 0 \rvert = \lvert {{\lvert r \rvert}^{N}} \rvert \to \varphi_{38}$\lbl{A2}}
\BIC{$\lvert \lvert {(r)^{\bullet}} \rvert(N) - 0 \rvert = \lvert {{\lvert r \rvert}^{N}} \rvert \to \varphi_{38}$}
\end{ptree}
\noindent$\Pi_{15}$:
\begin{ptree}
\AXC{$\forall x\,\forall y\,\forall z\,(x \leq y \to y \leq z \to x \leq z)$\lbl{[O5]}}
\AXC{$\varphi_{13}$\lbl{A2}}
\BIC{$\forall y\,\forall z\,(0 \leq y \to y \leq z \to 0 \leq z)$}
\AXC{$\varphi_{7}$\lbl{A2}}
\BIC{$\forall z\,(0 \leq 1 \to 1 \leq z \to 0 \leq z)$}
\AXC{$\varphi_{23}$\lbl{A2}}
\BIC{$0 \leq 1 \to 1 \leq {{\lvert r \rvert}^{N}} \to 0 \leq {{\lvert r \rvert}^{N}}$}
\AXC{$(0 \leq 1 \to 1 \leq {{\lvert r \rvert}^{N}} \to 0 \leq {{\lvert r \rvert}^{N}}) \to \varphi_{134}$\lbl{A1}}
\BIC{$\varphi_{134}$}
\end{ptree}
\noindent$\Pi_{16}$:
\begin{ptree}
\AXC{$\stackrel{\Pi_{6}}{\vdots}$}\noLine
\UIC{$r \in \mathbb{R} \to 1 \leq \lvert r \rvert \to \varphi_{81}$}
\AXC{$0 < 1$\lbl{[F28]}}
\AXC{$0 < 1 \to r \in \mathbb{R} \to 1 \leq \lvert r \rvert \to \varphi_{64}$\lbl{A1}}
\BIC{$r \in \mathbb{R} \to 1 \leq \lvert r \rvert \to \varphi_{64}$}
\AXC{$\stackrel{\Pi_{15}}{\vdots}$}\noLine
\UIC{$\varphi_{134}$}
\BIC{$(r \in \mathbb{R} \to 1 \leq \lvert r \rvert \to \varphi_{81}) \to r \in \mathbb{R} \to 1 \leq \lvert r \rvert \to \varphi_{82}$}
\BIC{$r \in \mathbb{R} \to 1 \leq \lvert r \rvert \to \varphi_{82}$}
\end{ptree}
\noindent$\Pi_{17}$:
\begin{ptree}
\AXC{$y = z \to x = y \to x = z$\lbl{A6}}
\AXC{$(y = z \to x = y \to x = z) \to x = y \to y = z \to x = z$\lbl{A1}}
\BIC{$x = y \to y = z \to x = z$}
\UIC{$\forall z\,(x = y \to y = z \to x = z)$}
\UIC{$\forall y\,\forall z\,(x = y \to y = z \to x = z)$}
\UIC{$\forall x\,\forall y\,\forall z\,(x = y \to y = z \to x = z)$}
\AXC{$\forall x\,\forall y\,\forall z\,(x = y \to y = z \to x = z) \to \varphi_{31}$\lbl{A2}}
\BIC{$\varphi_{31}$}
\end{ptree}
\noindent$\Pi_{18}$:
\begin{ptree}
\AXC{$\stackrel{\Pi_{17}}{\vdots}$}\noLine
\UIC{$\varphi_{31}$}
\AXC{$\varphi_{31} \to \varphi_{57}$\lbl{A2}}
\BIC{$\varphi_{57}$}
\AXC{$\varphi_{57} \to \lvert \lvert {(r)^{\bullet}} \rvert(N) - 0 \rvert = \lvert \lvert {{r}^{N}} \rvert \rvert \to \varphi_{43}$\lbl{A2}}
\BIC{$\lvert \lvert {(r)^{\bullet}} \rvert(N) - 0 \rvert = \lvert \lvert {{r}^{N}} \rvert \rvert \to \varphi_{43}$}
\end{ptree}
\noindent$\Pi_{19}$:
\begin{ptree}
\AXC{$r \in \mathbb{R} \to 1 \leq \lvert r \rvert \to \varphi_{72}$\lbl{A1}}
\AXC{$\forall r \in \mathbb{R}\,\forall n \in \mathbb{N}\,(\lvert {{r}^{n}} \rvert = {{\lvert r \rvert}^{n}})$\lbl{[P1]}}
\AXC{$\varphi_{79}$\lbl{A2}}
\BIC{$r \in \mathbb{R} \to \forall n \in \mathbb{N}\,(\lvert {{r}^{n}} \rvert = {{\lvert r \rvert}^{n}})$}
\AXC{$\varphi_{131}$\lbl{A1}}
\BIC{$(r \in \mathbb{R} \to 1 \leq \lvert r \rvert \to \varphi_{72}) \to \varphi_{116}$}
\BIC{$\varphi_{116}$}
\AXC{$\varphi_{116} \to \varphi_{44} \to \varphi_{114}$\lbl{A1}}
\BIC{$\varphi_{44} \to \varphi_{114}$}
\end{ptree}
\noindent$\Pi_{20}$:
\begin{ptree}
\AXC{$\stackrel{\Pi_{1}}{\vdots}$}\noLine
\UIC{$r \in \mathbb{R} \to 1 \leq \lvert r \rvert \to \varphi_{69}$}
\AXC{$\varphi_{44}$\lbl{A2}}
\AXC{$\stackrel{\Pi_{19}}{\vdots}$}\noLine
\UIC{$\varphi_{44} \to \varphi_{114}$}
\BIC{$\varphi_{114}$}
\AXC{$\varphi_{114} \to (r \in \mathbb{R} \to 1 \leq \lvert r \rvert \to \varphi_{69}) \to r \in \mathbb{R} \to 1 \leq \lvert r \rvert \to \varphi_{96}$\lbl{A1}}
\BIC{$(r \in \mathbb{R} \to 1 \leq \lvert r \rvert \to \varphi_{69}) \to r \in \mathbb{R} \to 1 \leq \lvert r \rvert \to \varphi_{96}$}
\BIC{$r \in \mathbb{R} \to 1 \leq \lvert r \rvert \to \varphi_{96}$}
\AXC{$\varphi_{141}$\lbl{A1}}
\BIC{$\lvert \lvert {{r}^{N}} \rvert \rvert = \lvert \lvert {{r}^{N}} \rvert \rvert \to \varphi_{76} \to \varphi_{110}$}
\end{ptree}
\noindent$\Pi_{21}$:
\begin{ptree}
\AXC{$x = x$\lbl{A5}}
\UIC{$\forall x\,(x = x)$}
\AXC{$\forall x\,(x = x) \to \lvert \lvert {{r}^{N}} \rvert \rvert = \lvert \lvert {{r}^{N}} \rvert \rvert$\lbl{A2}}
\BIC{$\lvert \lvert {{r}^{N}} \rvert \rvert = \lvert \lvert {{r}^{N}} \rvert \rvert$}
\AXC{$\stackrel{\Pi_{20}}{\vdots}$}\noLine
\UIC{$\lvert \lvert {{r}^{N}} \rvert \rvert = \lvert \lvert {{r}^{N}} \rvert \rvert \to \varphi_{76} \to \varphi_{110}$}
\BIC{$\varphi_{76} \to \varphi_{110}$}
\end{ptree}
\noindent$\Pi_{22}$:
\begin{ptree}
\AXC{$z = w \to \lvert \lvert {{r}^{N}} \rvert \rvert = \lvert z \rvert \to \lvert \lvert {{r}^{N}} \rvert \rvert = \lvert w \rvert$\lbl{A6}}
\UIC{$\forall w\,(z = w \to \lvert \lvert {{r}^{N}} \rvert \rvert = \lvert z \rvert \to \lvert \lvert {{r}^{N}} \rvert \rvert = \lvert w \rvert)$}
\UIC{$\varphi_{28}$}
\AXC{$\varphi_{28} \to \varphi_{47}$\lbl{A2}}
\BIC{$\varphi_{47}$}
\AXC{$\varphi_{47} \to \varphi_{76}$\lbl{A2}}
\BIC{$\varphi_{76}$}
\AXC{$\stackrel{\Pi_{21}}{\vdots}$}\noLine
\UIC{$\varphi_{76} \to \varphi_{110}$}
\BIC{$\varphi_{110}$}
\end{ptree}
\noindent$\Pi_{23}$:
\begin{ptree}
\AXC{$y = z \to x = y \to x = z$\lbl{A6}}
\AXC{$(y = z \to x = y \to x = z) \to x = y \to y = z \to x = z$\lbl{A1}}
\BIC{$x = y \to y = z \to x = z$}
\UIC{$\forall z\,(x = y \to y = z \to x = z)$}
\UIC{$\forall y\,\forall z\,(x = y \to y = z \to x = z)$}
\UIC{$\forall x\,\forall y\,\forall z\,(x = y \to y = z \to x = z)$}
\AXC{$\forall x\,\forall y\,\forall z\,(x = y \to y = z \to x = z) \to \varphi_{31}$\lbl{A2}}
\BIC{$\varphi_{31}$}
\end{ptree}
\noindent$\Pi_{24}$:
\begin{ptree}
\AXC{$r \in \mathbb{R} \to 1 \leq \lvert r \rvert \to \varphi_{72}$\lbl{A1}}
\AXC{$\forall r \in \mathbb{R}\,(\varphi_{34})$\lbl{[S24]}}
\AXC{$\forall r \in \mathbb{R}\,(\varphi_{34}) \to r \in \mathbb{R} \to \varphi_{34}$\lbl{A2}}
\BIC{$r \in \mathbb{R} \to \varphi_{34}$}
\AXC{$\varphi_{136}$\lbl{A1}}
\BIC{$(r \in \mathbb{R} \to 1 \leq \lvert r \rvert \to \varphi_{72}) \to r \in \mathbb{R} \to 1 \leq \lvert r \rvert \to \varphi_{121}$}
\BIC{$r \in \mathbb{R} \to 1 \leq \lvert r \rvert \to \varphi_{121}$}
\end{ptree}
\noindent$\Pi_{25}$:
\begin{ptree}
\AXC{$\stackrel{\Pi_{24}}{\vdots}$}\noLine
\UIC{$r \in \mathbb{R} \to 1 \leq \lvert r \rvert \to \varphi_{121}$}
\AXC{$(r \in \mathbb{R} \to 1 \leq \lvert r \rvert \to \varphi_{121}) \to \varphi_{95}$\lbl{A1}}
\BIC{$\varphi_{95}$}
\end{ptree}
\noindent$\Pi_{26}$:
\begin{ptree}
\AXC{$\stackrel{\Pi_{25}}{\vdots}$}\noLine
\UIC{$\varphi_{95}$}
\AXC{$\varphi_{25}$\lbl{[S11]}}
\AXC{$\varphi_{25} \to \mathrm{Seq}({(r)^{\bullet}}) \to \varphi_{26}$\lbl{A2}}
\BIC{$\mathrm{Seq}({(r)^{\bullet}}) \to \varphi_{26}$}
\AXC{$\varphi_{133}$\lbl{A1}}
\BIC{$\varphi_{95} \to r \in \mathbb{R} \to 1 \leq \lvert r \rvert \to \varphi_{113}$}
\BIC{$r \in \mathbb{R} \to 1 \leq \lvert r \rvert \to \varphi_{113}$}
\AXC{$(r \in \mathbb{R} \to 1 \leq \lvert r \rvert \to \varphi_{113}) \to \varphi_{118}$\lbl{A1}}
\BIC{$\varphi_{118}$}
\end{ptree}
\noindent$\Pi_{27}$:
\begin{ptree}
\AXC{$\stackrel{\Pi_{1}}{\vdots}$}\noLine
\UIC{$r \in \mathbb{R} \to 1 \leq \lvert r \rvert \to \varphi_{69}$}
\AXC{$\varphi_{48}$\lbl{A2}}
\AXC{$\stackrel{\Pi_{26}}{\vdots}$}\noLine
\UIC{$\varphi_{118}$}
\AXC{$\varphi_{118} \to \varphi_{48} \to \varphi_{115}$\lbl{A1}}
\BIC{$\varphi_{48} \to \varphi_{115}$}
\BIC{$\varphi_{115}$}
\AXC{$\varphi_{115} \to (r \in \mathbb{R} \to 1 \leq \lvert r \rvert \to \varphi_{69}) \to \varphi_{107}$\lbl{A1}}
\BIC{$(r \in \mathbb{R} \to 1 \leq \lvert r \rvert \to \varphi_{69}) \to \varphi_{107}$}
\BIC{$\varphi_{107}$}
\end{ptree}
\noindent$\Pi_{28}$:
\begin{ptree}
\AXC{$x = x$\lbl{A5}}
\UIC{$\forall x\,(x = x)$}
\AXC{$\forall x\,(x = x) \to \varphi_{10}$\lbl{A2}}
\BIC{$\varphi_{10}$}
\AXC{$\stackrel{\Pi_{27}}{\vdots}$}\noLine
\UIC{$\varphi_{107}$}
\AXC{$\varphi_{107} \to \varphi_{10} \to \varphi_{80} \to \varphi_{111}$\lbl{A1}}
\BIC{$\varphi_{10} \to \varphi_{80} \to \varphi_{111}$}
\BIC{$\varphi_{80} \to \varphi_{111}$}
\end{ptree}
\noindent$\Pi_{29}$:
\begin{ptree}
\AXC{$z = w \to \varphi_{21}$\lbl{A6}}
\UIC{$\forall w\,(z = w \to \varphi_{21})$}
\UIC{$\forall z\,\forall w\,(z = w \to \varphi_{21})$}
\AXC{$\forall z\,\forall w\,(z = w \to \varphi_{21}) \to \varphi_{50}$\lbl{A2}}
\BIC{$\varphi_{50}$}
\AXC{$\varphi_{50} \to \varphi_{80}$\lbl{A2}}
\BIC{$\varphi_{80}$}
\AXC{$\stackrel{\Pi_{28}}{\vdots}$}\noLine
\UIC{$\varphi_{80} \to \varphi_{111}$}
\BIC{$\varphi_{111}$}
\end{ptree}
\noindent$\Pi_{30}$:
\begin{ptree}
\AXC{$r \in \mathbb{R} \to 1 \leq \lvert r \rvert \to \varphi_{72}$\lbl{A1}}
\AXC{$\forall r \in \mathbb{R}\,(\varphi_{34})$\lbl{[S24]}}
\AXC{$\forall r \in \mathbb{R}\,(\varphi_{34}) \to r \in \mathbb{R} \to \varphi_{34}$\lbl{A2}}
\BIC{$r \in \mathbb{R} \to \varphi_{34}$}
\AXC{$\varphi_{136}$\lbl{A1}}
\BIC{$(r \in \mathbb{R} \to 1 \leq \lvert r \rvert \to \varphi_{72}) \to r \in \mathbb{R} \to 1 \leq \lvert r \rvert \to \varphi_{121}$}
\BIC{$r \in \mathbb{R} \to 1 \leq \lvert r \rvert \to \varphi_{121}$}
\end{ptree}
\noindent$\Pi_{31}$:
\begin{ptree}
\AXC{$\stackrel{\Pi_{30}}{\vdots}$}\noLine
\UIC{$r \in \mathbb{R} \to 1 \leq \lvert r \rvert \to \varphi_{121}$}
\AXC{$(r \in \mathbb{R} \to 1 \leq \lvert r \rvert \to \varphi_{121}) \to \varphi_{95}$\lbl{A1}}
\BIC{$\varphi_{95}$}
\end{ptree}
\noindent$\Pi_{32}$:
\begin{ptree}
\AXC{$\stackrel{\Pi_{31}}{\vdots}$}\noLine
\UIC{$\varphi_{95}$}
\AXC{$\varphi_{25}$\lbl{[S11]}}
\AXC{$\varphi_{25} \to \mathrm{Seq}({(r)^{\bullet}}) \to \varphi_{26}$\lbl{A2}}
\BIC{$\mathrm{Seq}({(r)^{\bullet}}) \to \varphi_{26}$}
\AXC{$\varphi_{133}$\lbl{A1}}
\BIC{$\varphi_{95} \to r \in \mathbb{R} \to 1 \leq \lvert r \rvert \to \varphi_{113}$}
\BIC{$r \in \mathbb{R} \to 1 \leq \lvert r \rvert \to \varphi_{113}$}
\AXC{$(r \in \mathbb{R} \to 1 \leq \lvert r \rvert \to \varphi_{113}) \to \varphi_{97}$\lbl{A1}}
\BIC{$\varphi_{97}$}
\end{ptree}
\noindent$\Pi_{33}$:
\begin{ptree}
\AXC{$\forall a\,\forall n\,(\mathrm{Seq}(a) \to n \in \mathbb{N} \to a_{n} \in \mathbb{R})$\lbl{[S1]}}
\AXC{$\forall a\,\forall n\,(\mathrm{Seq}(a) \to n \in \mathbb{N} \to a_{n} \in \mathbb{R}) \to \varphi_{20}$\lbl{A2}}
\BIC{$\varphi_{20}$}
\AXC{$\varphi_{20} \to \varphi_{19}$\lbl{A2}}
\BIC{$\varphi_{19}$}
\end{ptree}
\noindent$\Pi_{34}$:
\begin{ptree}
\AXC{$\forall x \in \mathbb{R}\,(x - 0 = x)$\lbl{[E28]}}
\AXC{$\forall x \in \mathbb{R}\,(x - 0 = x) \to \varphi_{27}$\lbl{A2}}
\BIC{$\varphi_{27}$}
\AXC{$\varphi_{27} \to \varphi_{101} \to r \in \mathbb{R} \to 1 \leq \lvert r \rvert \to \varphi_{99}$\lbl{A1}}
\BIC{$\varphi_{101} \to r \in \mathbb{R} \to 1 \leq \lvert r \rvert \to \varphi_{99}$}
\end{ptree}
\noindent$\Pi_{35}$:
\begin{ptree}
\AXC{$\stackrel{\Pi_{1}}{\vdots}$}\noLine
\UIC{$r \in \mathbb{R} \to 1 \leq \lvert r \rvert \to \varphi_{69}$}
\AXC{$\stackrel{\Pi_{32}}{\vdots}$}\noLine
\UIC{$\varphi_{97}$}
\AXC{$\stackrel{\Pi_{33}}{\vdots}$}\noLine
\UIC{$\varphi_{19}$}
\AXC{$\varphi_{19} \to \varphi_{97} \to (r \in \mathbb{R} \to 1 \leq \lvert r \rvert \to \varphi_{69}) \to \varphi_{101}$\lbl{A1}}
\BIC{$\varphi_{97} \to (r \in \mathbb{R} \to 1 \leq \lvert r \rvert \to \varphi_{69}) \to \varphi_{101}$}
\BIC{$(r \in \mathbb{R} \to 1 \leq \lvert r \rvert \to \varphi_{69}) \to \varphi_{101}$}
\BIC{$\varphi_{101}$}
\AXC{$\stackrel{\Pi_{34}}{\vdots}$}\noLine
\UIC{$\varphi_{101} \to r \in \mathbb{R} \to 1 \leq \lvert r \rvert \to \varphi_{99}$}
\BIC{$r \in \mathbb{R} \to 1 \leq \lvert r \rvert \to \varphi_{99}$}
\end{ptree}
\noindent$\Pi_{36}$:
\begin{ptree}
\AXC{$x = x$\lbl{A5}}
\UIC{$\forall x\,(x = x)$}
\AXC{$\forall x\,(x = x) \to \varphi_{17}$\lbl{A2}}
\BIC{$\varphi_{17}$}
\AXC{$\stackrel{\Pi_{35}}{\vdots}$}\noLine
\UIC{$r \in \mathbb{R} \to 1 \leq \lvert r \rvert \to \varphi_{99}$}
\AXC{$\varphi_{142}$\lbl{A1}}
\BIC{$\varphi_{17} \to \varphi_{98} \to r \in \mathbb{R} \to 1 \leq \lvert r \rvert \to \varphi_{102}$}
\BIC{$\varphi_{98} \to r \in \mathbb{R} \to 1 \leq \lvert r \rvert \to \varphi_{102}$}
\end{ptree}
\noindent$\Pi_{37}$:
\begin{ptree}
\AXC{$z = w \to \varphi_{30}$\lbl{A6}}
\UIC{$\forall w\,(z = w \to \varphi_{30})$}
\UIC{$\forall z\,\forall w\,(z = w \to \varphi_{30})$}
\AXC{$\forall z\,\forall w\,(z = w \to \varphi_{30}) \to \varphi_{71}$\lbl{A2}}
\BIC{$\varphi_{71}$}
\AXC{$\varphi_{71} \to \varphi_{98}$\lbl{A2}}
\BIC{$\varphi_{98}$}
\AXC{$\stackrel{\Pi_{36}}{\vdots}$}\noLine
\UIC{$\varphi_{98} \to r \in \mathbb{R} \to 1 \leq \lvert r \rvert \to \varphi_{102}$}
\BIC{$r \in \mathbb{R} \to 1 \leq \lvert r \rvert \to \varphi_{102}$}
\end{ptree}
\noindent$\Pi_{38}$:
\begin{ptree}
\AXC{$\stackrel{\Pi_{29}}{\vdots}$}\noLine
\UIC{$\varphi_{111}$}
\AXC{$\stackrel{\Pi_{37}}{\vdots}$}\noLine
\UIC{$r \in \mathbb{R} \to 1 \leq \lvert r \rvert \to \varphi_{102}$}
\AXC{$(r \in \mathbb{R} \to 1 \leq \lvert r \rvert \to \varphi_{102}) \to \varphi_{111} \to (\varphi_{16} \to \varphi_{45}) \to \varphi_{119}$\lbl{A1}}
\BIC{$\varphi_{111} \to (\varphi_{16} \to \varphi_{45}) \to \varphi_{119}$}
\BIC{$(\varphi_{16} \to \varphi_{45}) \to \varphi_{119}$}
\end{ptree}
\noindent$\Pi_{39}$:
\begin{ptree}
\AXC{$\stackrel{\Pi_{22}}{\vdots}$}\noLine
\UIC{$\varphi_{110}$}
\AXC{$\stackrel{\Pi_{23}}{\vdots}$}\noLine
\UIC{$\varphi_{31}$}
\AXC{$\varphi_{31} \to \varphi_{59}$\lbl{A2}}
\BIC{$\varphi_{59}$}
\AXC{$\varphi_{59} \to \varphi_{16} \to \varphi_{45}$\lbl{A2}}
\BIC{$\varphi_{16} \to \varphi_{45}$}
\AXC{$\stackrel{\Pi_{38}}{\vdots}$}\noLine
\UIC{$(\varphi_{16} \to \varphi_{45}) \to \varphi_{119}$}
\BIC{$\varphi_{119}$}
\AXC{$\varphi_{119} \to \varphi_{139}$\lbl{A1}}
\BIC{$\varphi_{139}$}
\BIC{$(\lvert \lvert {(r)^{\bullet}} \rvert(N) - 0 \rvert = \lvert \lvert {{r}^{N}} \rvert \rvert \to \varphi_{43}) \to \varphi_{120}$}
\end{ptree}
\noindent$\Pi_{40}$:
\begin{ptree}
\AXC{$\stackrel{\Pi_{18}}{\vdots}$}\noLine
\UIC{$\lvert \lvert {(r)^{\bullet}} \rvert(N) - 0 \rvert = \lvert \lvert {{r}^{N}} \rvert \rvert \to \varphi_{43}$}
\AXC{$\stackrel{\Pi_{39}}{\vdots}$}\noLine
\UIC{$(\lvert \lvert {(r)^{\bullet}} \rvert(N) - 0 \rvert = \lvert \lvert {{r}^{N}} \rvert \rvert \to \varphi_{43}) \to \varphi_{120}$}
\BIC{$\varphi_{120}$}
\AXC{$\varphi_{120} \to \varphi_{138}$\lbl{A1}}
\BIC{$\varphi_{138}$}
\end{ptree}
\noindent$\Pi_{41}$:
\begin{ptree}
\AXC{$\stackrel{\Pi_{16}}{\vdots}$}\noLine
\UIC{$r \in \mathbb{R} \to 1 \leq \lvert r \rvert \to \varphi_{82}$}
\AXC{$\forall x\,(0 \leq x \to \lvert x \rvert = x)$\lbl{[A1]}}
\AXC{$\varphi_{35}$\lbl{A2}}
\BIC{$0 \leq {{\lvert r \rvert}^{N}} \to \lvert {{\lvert r \rvert}^{N}} \rvert = {{\lvert r \rvert}^{N}}$}
\AXC{$\varphi_{130}$\lbl{A1}}
\BIC{$(r \in \mathbb{R} \to 1 \leq \lvert r \rvert \to \varphi_{82}) \to \varphi_{108}$}
\BIC{$\varphi_{108}$}
\AXC{$\stackrel{\Pi_{40}}{\vdots}$}\noLine
\UIC{$\varphi_{138}$}
\BIC{$(\lvert \lvert {(r)^{\bullet}} \rvert(N) - 0 \rvert = \lvert {{\lvert r \rvert}^{N}} \rvert \to \varphi_{38}) \to \varphi_{112}$}
\end{ptree}
\noindent$\Pi_{42}$:
\begin{ptree}
\AXC{$\stackrel{\Pi_{14}}{\vdots}$}\noLine
\UIC{$\lvert \lvert {(r)^{\bullet}} \rvert(N) - 0 \rvert = \lvert {{\lvert r \rvert}^{N}} \rvert \to \varphi_{38}$}
\AXC{$\stackrel{\Pi_{41}}{\vdots}$}\noLine
\UIC{$(\lvert \lvert {(r)^{\bullet}} \rvert(N) - 0 \rvert = \lvert {{\lvert r \rvert}^{N}} \rvert \to \varphi_{38}) \to \varphi_{112}$}
\BIC{$\varphi_{112}$}
\end{ptree}
\noindent$\Pi_{43}$:
\begin{ptree}
\AXC{$\stackrel{\Pi_{8}}{\vdots}$}\noLine
\UIC{$r \in \mathbb{R} \to 1 \leq \lvert r \rvert \to \varphi_{88}$}
\AXC{$\stackrel{\Pi_{9}}{\vdots}$}\noLine
\UIC{$\varphi_{51}$}
\AXC{$\stackrel{\Pi_{12}}{\vdots}$}\noLine
\UIC{$\varphi_{105}$}
\AXC{$\stackrel{\Pi_{42}}{\vdots}$}\noLine
\UIC{$\varphi_{112}$}
\AXC{$\varphi_{112} \to \varphi_{105} \to \varphi_{51} \to r \in \mathbb{R} \to 1 \leq \lvert r \rvert \to \varphi_{87}$\lbl{A1}}
\BIC{$\varphi_{105} \to \varphi_{51} \to r \in \mathbb{R} \to 1 \leq \lvert r \rvert \to \varphi_{87}$}
\BIC{$\varphi_{51} \to r \in \mathbb{R} \to 1 \leq \lvert r \rvert \to \varphi_{87}$}
\BIC{$r \in \mathbb{R} \to 1 \leq \lvert r \rvert \to \varphi_{87}$}
\AXC{$\varphi_{137}$\lbl{A1}}
\BIC{$(r \in \mathbb{R} \to 1 \leq \lvert r \rvert \to \varphi_{88}) \to \varphi_{100}$}
\BIC{$\varphi_{100}$}
\end{ptree}
\noindent$\Pi_{44}$:
\begin{ptree}
\AXC{$r \in \mathbb{R} \to 1 \leq \lvert r \rvert \to \varphi_{18}$\lbl{A1}}
\AXC{$\forall r \in \mathbb{R}\,(\varphi_{34})$\lbl{[S24]}}
\AXC{$\forall r \in \mathbb{R}\,(\varphi_{34}) \to r \in \mathbb{R} \to \varphi_{34}$\lbl{A2}}
\BIC{$r \in \mathbb{R} \to \varphi_{34}$}
\AXC{$\varphi_{125}$\lbl{A1}}
\BIC{$(r \in \mathbb{R} \to 1 \leq \lvert r \rvert \to \varphi_{18}) \to r \in \mathbb{R} \to 1 \leq \lvert r \rvert \to \varphi_{74}$}
\BIC{$r \in \mathbb{R} \to 1 \leq \lvert r \rvert \to \varphi_{74}$}
\end{ptree}
\noindent$\Pi_{45}$:
\begin{ptree}
\AXC{$\stackrel{\Pi_{44}}{\vdots}$}\noLine
\UIC{$r \in \mathbb{R} \to 1 \leq \lvert r \rvert \to \varphi_{74}$}
\AXC{$(r \in \mathbb{R} \to 1 \leq \lvert r \rvert \to \varphi_{74}) \to \varphi_{33}$\lbl{A1}}
\BIC{$\varphi_{33}$}
\end{ptree}
\noindent$\Pi_{46}$:
\begin{ptree}
\AXC{$\varphi_{8}$\lbl{定理4.7(1)}}
\AXC{$\varphi_{8} \to \forall s\,(\mathrm{Seq}({(r)^{\bullet}}) \to \varphi_{2})$\lbl{A2}}
\BIC{$\forall s\,(\mathrm{Seq}({(r)^{\bullet}}) \to \varphi_{2})$}
\AXC{$\varphi_{49}$\lbl{A2}}
\BIC{$\mathrm{Seq}({(r)^{\bullet}}) \to \varphi_{2}$}
\AXC{$(\mathrm{Seq}({(r)^{\bullet}}) \to \varphi_{2}) \to \varphi_{33} \to \varphi_{39} \to \varphi_{37}$\lbl{A1}}
\BIC{$\varphi_{33} \to \varphi_{39} \to \varphi_{37}$}
\end{ptree}
\noindent$\Pi_{47}$:
\begin{ptree}
\AXC{$\stackrel{\Pi_{45}}{\vdots}$}\noLine
\UIC{$\varphi_{33}$}
\AXC{$\varphi_{4}$\lbl{定理2.15(2)}}
\AXC{$\varphi_{4} \to \mathrm{Seq}({(r)^{\bullet}}) \to \varphi_{3}$\lbl{A2}}
\BIC{$\mathrm{Seq}({(r)^{\bullet}}) \to \varphi_{3}$}
\AXC{$\varphi_{124}$\lbl{A1}}
\BIC{$\varphi_{33} \to r \in \mathbb{R} \to 1 \leq \lvert r \rvert \to \varphi_{36}$}
\BIC{$r \in \mathbb{R} \to 1 \leq \lvert r \rvert \to \varphi_{36}$}
\AXC{$(r \in \mathbb{R} \to 1 \leq \lvert r \rvert \to \varphi_{36}) \to \varphi_{37} \to \varphi_{40}$\lbl{A1}}
\BIC{$\varphi_{37} \to \varphi_{40}$}
\end{ptree}
\noindent$\Pi_{48}$:
\begin{ptree}
\AXC{$\varphi_{55}$\lbl{[C3]}}
\AXC{$\varphi_{55} \to \varphi_{75}$\lbl{A2}}
\BIC{$\varphi_{75}$}
\AXC{$\varphi_{75} \to \varphi_{67}$\lbl{A2}}
\BIC{$\varphi_{67}$}
\AXC{$\varphi_{123}$\lbl{A2}}
\BIC{$\mathrm{Conv}(\lvert {(r)^{\bullet}} \rvert,0) \to 1 \in \mathbb{R} \to 0 < 1 \to \varphi_{24}$}
\AXC{$\varphi_{128}$\lbl{A1}}
\BIC{$\varphi_{40} \to 1 \in \mathbb{R} \to 0 < 1 \to r \in \mathbb{R} \to 1 \leq \lvert r \rvert \to \varphi_{60}$}
\end{ptree}
\noindent$\Pi_{49}$:
\begin{ptree}
\AXC{$1 \in \mathbb{R}$\lbl{(R1)}}
\AXC{$\varphi_{39}$\lbl{A1}}
\AXC{$\stackrel{\Pi_{45}}{\vdots}$}\noLine
\UIC{$\varphi_{33}$}
\AXC{$\stackrel{\Pi_{46}}{\vdots}$}\noLine
\UIC{$\varphi_{33} \to \varphi_{39} \to \varphi_{37}$}
\BIC{$\varphi_{39} \to \varphi_{37}$}
\BIC{$\varphi_{37}$}
\AXC{$\stackrel{\Pi_{47}}{\vdots}$}\noLine
\UIC{$\varphi_{37} \to \varphi_{40}$}
\BIC{$\varphi_{40}$}
\AXC{$\stackrel{\Pi_{48}}{\vdots}$}\noLine
\UIC{$\varphi_{40} \to 1 \in \mathbb{R} \to 0 < 1 \to r \in \mathbb{R} \to 1 \leq \lvert r \rvert \to \varphi_{60}$}
\BIC{$1 \in \mathbb{R} \to 0 < 1 \to r \in \mathbb{R} \to 1 \leq \lvert r \rvert \to \varphi_{60}$}
\BIC{$0 < 1 \to r \in \mathbb{R} \to 1 \leq \lvert r \rvert \to \varphi_{60}$}
\end{ptree}
\noindent$\Pi_{50}$:
\begin{ptree}
\AXC{$0 < 1$\lbl{[F28]}}
\AXC{$\stackrel{\Pi_{49}}{\vdots}$}\noLine
\UIC{$0 < 1 \to r \in \mathbb{R} \to 1 \leq \lvert r \rvert \to \varphi_{60}$}
\BIC{$r \in \mathbb{R} \to 1 \leq \lvert r \rvert \to \varphi_{60}$}
\AXC{$(r \in \mathbb{R} \to 1 \leq \lvert r \rvert \to \varphi_{60}) \to \forall N\,(\varphi_{42} \to \neg \varphi_{22}) \to \varphi_{143}$\lbl{A1}}
\BIC{$\forall N\,(\varphi_{42} \to \neg \varphi_{22}) \to \varphi_{143}$}
\end{ptree}
\noindent$\Pi_{51}$:
\begin{ptree}
\AXC{$\varphi_{42} \to \forall N\,(\varphi_{42})$\lbl{A4}}
\AXC{$\forall N\,(\varphi_{42} \to \neg \varphi_{22}) \to \forall N\,(\varphi_{42}) \to \forall N\,(\neg \varphi_{22})$\lbl{A3}}
\AXC{$\stackrel{\Pi_{43}}{\vdots}$}\noLine
\UIC{$\varphi_{100}$}
\AXC{$\varphi_{100} \to \varphi_{42} \to \neg \varphi_{22}$\lbl{A1}}
\BIC{$\varphi_{42} \to \neg \varphi_{22}$}
\UIC{$\forall N\,(\varphi_{42} \to \neg \varphi_{22})$}
\AXC{$\stackrel{\Pi_{50}}{\vdots}$}\noLine
\UIC{$\forall N\,(\varphi_{42} \to \neg \varphi_{22}) \to \varphi_{143}$}
\BIC{$\varphi_{143}$}
\BIC{$(\varphi_{42} \to \forall N\,(\varphi_{42})) \to \varphi_{41}$}
\BIC{$\varphi_{41}$}
\end{ptree}
\noindent$\Pi_{52}$:
\begin{ptree}
\AXC{$\stackrel{\Pi_{51}}{\vdots}$}\noLine
\UIC{$\varphi_{41}$}
\AXC{$\varphi_{41} \to \varphi_{29} \to \neg \mathrm{Conv}(\Sigma {(r)^{\bullet}},s)$\lbl{A1}}
\BIC{$\varphi_{29} \to \neg \mathrm{Conv}(\Sigma {(r)^{\bullet}},s)$}
\UIC{$\forall s\,(\varphi_{29} \to \neg \mathrm{Conv}(\Sigma {(r)^{\bullet}},s))$}
\AXC{$r \in \mathbb{R} \to 1 \leq \lvert r \rvert \to \varphi_{12}$\lbl{A1}}
\AXC{$(r \in \mathbb{R} \to 1 \leq \lvert r \rvert \to \varphi_{12}) \to \varphi_{140}$\lbl{A1}}
\BIC{$\varphi_{140}$}
\BIC{$\varphi_{122} \to (\varphi_{29} \to \forall s\,(\varphi_{29})) \to r \in \mathbb{R} \to 1 \leq \lvert r \rvert \to \varphi_{14}$}
\end{ptree}
\noindent 主木:
\begin{ptree}
\AXC{$\varphi_{29} \to \forall s\,(\varphi_{29})$\lbl{A4}}
\AXC{$\varphi_{122}$\lbl{A3}}
\AXC{$\stackrel{\Pi_{52}}{\vdots}$}\noLine
\UIC{$\varphi_{122} \to (\varphi_{29} \to \forall s\,(\varphi_{29})) \to r \in \mathbb{R} \to 1 \leq \lvert r \rvert \to \varphi_{14}$}
\BIC{$(\varphi_{29} \to \forall s\,(\varphi_{29})) \to r \in \mathbb{R} \to 1 \leq \lvert r \rvert \to \varphi_{14}$}
\BIC{$r \in \mathbb{R} \to 1 \leq \lvert r \rvert \to \varphi_{14}$}
\AXC{$(r \in \mathbb{R} \to 1 \leq \lvert r \rvert \to \varphi_{14}) \to \varphi_{11}$\lbl{A1}}
\BIC{$\varphi_{11}$}
\UIC{$\varphi_{15}$}
\end{ptree}
\begin{abbrevlist}
\item $\varphi_{1} :\equiv \mathrm{Conv}(\Sigma {(r)^{\bullet}},s) \to \mathrm{Seq}({(r)^{\bullet}})$
\item $\varphi_{2} :\equiv \mathrm{Conv}(\Sigma {(r)^{\bullet}},s) \to \mathrm{Conv}({(r)^{\bullet}},0)$
\item $\varphi_{3} :\equiv \mathrm{Conv}({(r)^{\bullet}},0) \leftrightarrow \mathrm{Conv}(\lvert {(r)^{\bullet}} \rvert,0)$
\item $\varphi_{4} :\equiv \forall a\,(\mathrm{Seq}(a) \to (\mathrm{Conv}(a,0) \leftrightarrow \mathrm{Conv}(\lvert a \rvert,0)))$
\item $\varphi_{5} :\equiv \mathrm{Conv}(\Sigma {(r)^{\bullet}},s) \to \mathrm{Conv}(\Sigma {(r)^{\bullet}},s)$
\item $\varphi_{6} :\equiv \mathrm{Conv}(\Sigma {(r)^{\bullet}},s) \to \mathrm{Conv}(\lvert {(r)^{\bullet}} \rvert,0)$
\item $\varphi_{7} :\equiv \forall y\,\forall z\,(0 \leq y \to y \leq z \to 0 \leq z) \to \forall z\,(0 \leq 1 \to 1 \leq z \to 0 \leq z)$
\item $\varphi_{8} :\equiv \forall a\,\forall s\,(\mathrm{Seq}(a) \to \mathrm{Conv}(\Sigma a,s) \to \mathrm{Conv}(a,0))$
\item $\varphi_{9} :\equiv \mathrm{Conv}(\Sigma {(r)^{\bullet}},s) \to \neg \exists s\,(\mathrm{Conv}(\Sigma {(r)^{\bullet}},s))$
\item $\varphi_{10} :\equiv \lvert \lvert {(r)^{\bullet}} \rvert(N) \rvert = \lvert \lvert {(r)^{\bullet}} \rvert(N) \rvert$
\item $\varphi_{11} :\equiv r \in \mathbb{R} \to 1 \leq \lvert r \rvert \to \neg \exists s\,(\mathrm{Conv}(\Sigma {(r)^{\bullet}},s))$
\item $\varphi_{12} :\equiv \exists s\,(\mathrm{Conv}(\Sigma {(r)^{\bullet}},s)) \to \exists s\,(\mathrm{Conv}(\Sigma {(r)^{\bullet}},s))$
\item $\varphi_{13} :\equiv \forall x\,\forall y\,\forall z\,(x \leq y \to y \leq z \to x \leq z) \to \forall y\,\forall z\,(0 \leq y \to y \leq z \to 0 \leq z)$
\item $\varphi_{14} :\equiv \exists s\,(\mathrm{Conv}(\Sigma {(r)^{\bullet}},s)) \to \neg \exists s\,(\mathrm{Conv}(\Sigma {(r)^{\bullet}},s))$
\item $\varphi_{15} :\equiv \forall r \in \mathbb{R}\,(1 \leq \lvert r \rvert \to \neg \exists s\,(\mathrm{Conv}(\Sigma {(r)^{\bullet}},s)))$
\item $\varphi_{16} :\equiv \lvert \lvert {(r)^{\bullet}} \rvert(N) - 0 \rvert = \lvert \lvert {(r)^{\bullet}} \rvert(N) \rvert$
\item $\varphi_{17} :\equiv \lvert \lvert {(r)^{\bullet}} \rvert(N) - 0 \rvert = \lvert \lvert {(r)^{\bullet}} \rvert(N) - 0 \rvert$
\item $\varphi_{18} :\equiv \exists s\,(\mathrm{Conv}(\Sigma {(r)^{\bullet}},s)) \to \mathrm{Conv}(\Sigma {(r)^{\bullet}},s) \to r \in \mathbb{R}$
\item $\varphi_{19} :\equiv \mathrm{Seq}(\lvert {(r)^{\bullet}} \rvert) \to N \in \mathbb{N} \to \lvert {(r)^{\bullet}} \rvert(N) \in \mathbb{R}$
\item $\varphi_{20} :\equiv \forall n\,(\mathrm{Seq}(\lvert {(r)^{\bullet}} \rvert) \to n \in \mathbb{N} \to \lvert {(r)^{\bullet}} \rvert(n) \in \mathbb{R})$
\item $\varphi_{21} :\equiv \lvert \lvert {(r)^{\bullet}} \rvert(N) \rvert = \lvert z \rvert \to \lvert \lvert {(r)^{\bullet}} \rvert(N) \rvert = \lvert w \rvert$
\item $\varphi_{22} :\equiv N \in \mathbb{N} \land \forall n \in \mathbb{N}\,(N \leq n \to \lvert \lvert {(r)^{\bullet}} \rvert(n) - 0 \rvert < 1)$
\item $\varphi_{23} :\equiv \forall z\,(0 \leq 1 \to 1 \leq z \to 0 \leq z) \to 0 \leq 1 \to 1 \leq {{\lvert r \rvert}^{N}} \to 0 \leq {{\lvert r \rvert}^{N}}$
\item $\varphi_{24} :\equiv \exists N \in \mathbb{N}\,\forall n \in \mathbb{N}\,(N \leq n \to \lvert \lvert {(r)^{\bullet}} \rvert(n) - 0 \rvert < 1)$
\item $\varphi_{25} :\equiv \forall a\,(\mathrm{Seq}(a) \to \mathrm{Seq}(\lvert a \rvert) \land \forall n \in \mathbb{N}\,(\lvert a \rvert(n) = \lvert a_{n} \rvert))$
\item $\varphi_{26} :\equiv \mathrm{Seq}(\lvert {(r)^{\bullet}} \rvert) \land \forall n \in \mathbb{N}\,(\lvert {(r)^{\bullet}} \rvert(n) = \lvert {{r}^{n}} \rvert)$
\item $\varphi_{27} :\equiv \lvert {(r)^{\bullet}} \rvert(N) \in \mathbb{R} \to \lvert {(r)^{\bullet}} \rvert(N) - 0 = \lvert {(r)^{\bullet}} \rvert(N)$
\item $\varphi_{28} :\equiv \forall z\,\forall w\,(z = w \to \lvert \lvert {{r}^{N}} \rvert \rvert = \lvert z \rvert \to \lvert \lvert {{r}^{N}} \rvert \rvert = \lvert w \rvert)$
\item $\varphi_{29} :\equiv r \in \mathbb{R} \land 1 \leq \lvert r \rvert \land \exists s\,(\mathrm{Conv}(\Sigma {(r)^{\bullet}},s)) \land \neg \neg \exists s\,(\mathrm{Conv}(\Sigma {(r)^{\bullet}},s))$
\item $\varphi_{30} :\equiv \lvert \lvert {(r)^{\bullet}} \rvert(N) - 0 \rvert = \lvert z \rvert \to \lvert \lvert {(r)^{\bullet}} \rvert(N) - 0 \rvert = \lvert w \rvert$
\item $\varphi_{31} :\equiv \forall y\,\forall z\,(\lvert \lvert {(r)^{\bullet}} \rvert(N) - 0 \rvert = y \to y = z \to \lvert \lvert {(r)^{\bullet}} \rvert(N) - 0 \rvert = z)$
\item $\varphi_{32} :\equiv \forall w\,(\lvert \lvert {(r)^{\bullet}} \rvert(N) - 0 \rvert = w \to \lvert \lvert {(r)^{\bullet}} \rvert(N) - 0 \rvert < 1 \to w < 1)$
\item $\varphi_{33} :\equiv r \in \mathbb{R} \to 1 \leq \lvert r \rvert \to \exists s\,(\mathrm{Conv}(\Sigma {(r)^{\bullet}},s)) \to \varphi_{1}$
\item $\varphi_{34} :\equiv \mathrm{Seq}({(r)^{\bullet}}) \land {{r}^{1_{\omega}}} = r \land \forall n \in \mathbb{N}\,({{r}^{S(n)}} = {{r}^{n}} \cdot r)$
\item $\varphi_{35} :\equiv \forall x\,(0 \leq x \to \lvert x \rvert = x) \to 0 \leq {{\lvert r \rvert}^{N}} \to \lvert {{\lvert r \rvert}^{N}} \rvert = {{\lvert r \rvert}^{N}}$
\item $\varphi_{36} :\equiv \exists s\,(\mathrm{Conv}(\Sigma {(r)^{\bullet}},s)) \to \mathrm{Conv}(\Sigma {(r)^{\bullet}},s) \to \varphi_{3}$
\item $\varphi_{37} :\equiv r \in \mathbb{R} \to 1 \leq \lvert r \rvert \to \exists s\,(\mathrm{Conv}(\Sigma {(r)^{\bullet}},s)) \to \varphi_{2}$
\item $\varphi_{38} :\equiv \lvert {{\lvert r \rvert}^{N}} \rvert = {{\lvert r \rvert}^{N}} \to \lvert \lvert {(r)^{\bullet}} \rvert(N) - 0 \rvert = {{\lvert r \rvert}^{N}}$
\item $\varphi_{39} :\equiv r \in \mathbb{R} \to 1 \leq \lvert r \rvert \to \exists s\,(\mathrm{Conv}(\Sigma {(r)^{\bullet}},s)) \to \varphi_{5}$
\item $\varphi_{40} :\equiv r \in \mathbb{R} \to 1 \leq \lvert r \rvert \to \exists s\,(\mathrm{Conv}(\Sigma {(r)^{\bullet}},s)) \to \varphi_{6}$
\item $\varphi_{41} :\equiv r \in \mathbb{R} \to 1 \leq \lvert r \rvert \to \exists s\,(\mathrm{Conv}(\Sigma {(r)^{\bullet}},s)) \to \varphi_{9}$
\item $\varphi_{42} :\equiv r \in \mathbb{R} \land 1 \leq \lvert r \rvert \land \exists s\,(\mathrm{Conv}(\Sigma {(r)^{\bullet}},s)) \land \mathrm{Conv}(\Sigma {(r)^{\bullet}},s) \land \neg \neg \exists s\,(\mathrm{Conv}(\Sigma {(r)^{\bullet}},s))$
\item $\varphi_{43} :\equiv \lvert \lvert {{r}^{N}} \rvert \rvert = \lvert {{\lvert r \rvert}^{N}} \rvert \to \lvert \lvert {(r)^{\bullet}} \rvert(N) - 0 \rvert = \lvert {{\lvert r \rvert}^{N}} \rvert$
\item $\varphi_{44} :\equiv \forall n \in \mathbb{N}\,(\lvert {{r}^{n}} \rvert = {{\lvert r \rvert}^{n}}) \to N \in \mathbb{N} \to \lvert {{r}^{N}} \rvert = {{\lvert r \rvert}^{N}}$
\item $\varphi_{45} :\equiv \lvert \lvert {(r)^{\bullet}} \rvert(N) \rvert = \lvert \lvert {{r}^{N}} \rvert \rvert \to \lvert \lvert {(r)^{\bullet}} \rvert(N) - 0 \rvert = \lvert \lvert {{r}^{N}} \rvert \rvert$
\item $\varphi_{46} :\equiv \mathrm{Conv}(\Sigma {(r)^{\bullet}},s) \to \varphi_{22} \to \mathrm{Seq}({(r)^{\bullet}})$
\item $\varphi_{47} :\equiv \forall w\,(\lvert {{r}^{N}} \rvert = w \to \lvert \lvert {{r}^{N}} \rvert \rvert = \lvert \lvert {{r}^{N}} \rvert \rvert \to \lvert \lvert {{r}^{N}} \rvert \rvert = \lvert w \rvert)$
\item $\varphi_{48} :\equiv \forall n \in \mathbb{N}\,(\lvert {(r)^{\bullet}} \rvert(n) = \lvert {{r}^{n}} \rvert) \to N \in \mathbb{N} \to \lvert {(r)^{\bullet}} \rvert(N) = \lvert {{r}^{N}} \rvert$
\item $\varphi_{49} :\equiv \forall s\,(\mathrm{Seq}({(r)^{\bullet}}) \to \varphi_{2}) \to \mathrm{Seq}({(r)^{\bullet}}) \to \varphi_{2}$
\item $\varphi_{50} :\equiv \forall w\,(\lvert {(r)^{\bullet}} \rvert(N) = w \to \varphi_{10} \to \lvert \lvert {(r)^{\bullet}} \rvert(N) \rvert = \lvert w \rvert)$
\item $\varphi_{51} :\equiv \lvert \lvert {(r)^{\bullet}} \rvert(N) - 0 \rvert = {{\lvert r \rvert}^{N}} \to \lvert \lvert {(r)^{\bullet}} \rvert(N) - 0 \rvert < 1 \to {{\lvert r \rvert}^{N}} < 1$
\item $\varphi_{52} :\equiv \mathrm{Conv}(\Sigma {(r)^{\bullet}},s) \to \varphi_{22} \to \mathrm{Seq}(\lvert {(r)^{\bullet}} \rvert)$
\item $\varphi_{53} :\equiv \forall n \in \mathbb{N}\,(N \leq n \to \lvert \lvert {(r)^{\bullet}} \rvert(n) - 0 \rvert < 1) \to N \in \mathbb{N} \to N \leq N \to \lvert \lvert {(r)^{\bullet}} \rvert(N) - 0 \rvert < 1$
\item $\varphi_{54} :\equiv \varphi_{22} \to \lvert {(r)^{\bullet}} \rvert(N) - 0 = \lvert {(r)^{\bullet}} \rvert(N)$
\item $\varphi_{55} :\equiv \forall a\,\forall \alpha\,\forall \varepsilon\,(\mathrm{Conv}(a,\alpha) \to \varepsilon \in \mathbb{R} \to 0 < \varepsilon \to \exists N \in \mathbb{N}\,\forall n \in \mathbb{N}\,(N \leq n \to \lvert a_{n} - \alpha \rvert < \varepsilon))$
\item $\varphi_{56} :\equiv \forall z\,(\lvert \lvert {(r)^{\bullet}} \rvert(N) - 0 \rvert = \lvert {{\lvert r \rvert}^{N}} \rvert \to \lvert {{\lvert r \rvert}^{N}} \rvert = z \to \lvert \lvert {(r)^{\bullet}} \rvert(N) - 0 \rvert = z)$
\item $\varphi_{57} :\equiv \forall z\,(\lvert \lvert {(r)^{\bullet}} \rvert(N) - 0 \rvert = \lvert \lvert {{r}^{N}} \rvert \rvert \to \lvert \lvert {{r}^{N}} \rvert \rvert = z \to \lvert \lvert {(r)^{\bullet}} \rvert(N) - 0 \rvert = z)$
\item $\varphi_{58} :\equiv \mathrm{Conv}(\Sigma {(r)^{\bullet}},s) \to \varphi_{22} \to \neg \exists s\,(\mathrm{Conv}(\Sigma {(r)^{\bullet}},s))$
\item $\varphi_{59} :\equiv \forall z\,(\varphi_{16} \to \lvert \lvert {(r)^{\bullet}} \rvert(N) \rvert = z \to \lvert \lvert {(r)^{\bullet}} \rvert(N) - 0 \rvert = z)$
\item $\varphi_{60} :\equiv \exists s\,(\mathrm{Conv}(\Sigma {(r)^{\bullet}},s)) \to \mathrm{Conv}(\Sigma {(r)^{\bullet}},s) \to \varphi_{24}$
\item $\varphi_{61} :\equiv \forall r \in \mathbb{R}\,(1 \leq r \to \forall n \in \mathbb{N}\,(1 \leq {{r}^{n}})) \to \lvert r \rvert \in \mathbb{R} \to 1 \leq \lvert r \rvert \to \forall n \in \mathbb{N}\,(1 \leq {{\lvert r \rvert}^{n}})$
\item $\varphi_{62} :\equiv \mathrm{Conv}(\Sigma {(r)^{\bullet}},s) \to \varphi_{22} \to \lvert {(r)^{\bullet}} \rvert(N) \in \mathbb{R}$
\item $\varphi_{63} :\equiv \exists s\,(\mathrm{Conv}(\Sigma {(r)^{\bullet}},s)) \to \mathrm{Conv}(\Sigma {(r)^{\bullet}},s) \to \varphi_{22} \to N \leq N$
\item $\varphi_{64} :\equiv \exists s\,(\mathrm{Conv}(\Sigma {(r)^{\bullet}},s)) \to \mathrm{Conv}(\Sigma {(r)^{\bullet}},s) \to \varphi_{22} \to 0 \leq 1$
\item $\varphi_{65} :\equiv \mathrm{Conv}(\Sigma {(r)^{\bullet}},s) \to \varphi_{22} \to N \in \mathbb{N} \to 1 \leq {{\lvert r \rvert}^{N}}$
\item $\varphi_{66} :\equiv \varphi_{22} \to \forall n \in \mathbb{N}\,(N \leq n \to \lvert \lvert {(r)^{\bullet}} \rvert(n) - 0 \rvert < 1)$
\item $\varphi_{67} :\equiv \forall \varepsilon\,(\mathrm{Conv}(\lvert {(r)^{\bullet}} \rvert,0) \to \varepsilon \in \mathbb{R} \to 0 < \varepsilon \to \exists N \in \mathbb{N}\,\forall n \in \mathbb{N}\,(N \leq n \to \lvert \lvert {(r)^{\bullet}} \rvert(n) - 0 \rvert < \varepsilon))$
\item $\varphi_{68} :\equiv \mathrm{Conv}(\Sigma {(r)^{\bullet}},s) \to \varphi_{22} \to \lvert \lvert {(r)^{\bullet}} \rvert(N) - 0 \rvert < 1$
\item $\varphi_{69} :\equiv \exists s\,(\mathrm{Conv}(\Sigma {(r)^{\bullet}},s)) \to \mathrm{Conv}(\Sigma {(r)^{\bullet}},s) \to \varphi_{22} \to N \in \mathbb{N}$
\item $\varphi_{70} :\equiv \mathrm{Conv}(\Sigma {(r)^{\bullet}},s) \to \varphi_{22} \to \forall n \in \mathbb{N}\,(1 \leq {{\lvert r \rvert}^{n}})$
\item $\varphi_{71} :\equiv \forall w\,(\lvert {(r)^{\bullet}} \rvert(N) - 0 = w \to \varphi_{17} \to \lvert \lvert {(r)^{\bullet}} \rvert(N) - 0 \rvert = \lvert w \rvert)$
\item $\varphi_{72} :\equiv \exists s\,(\mathrm{Conv}(\Sigma {(r)^{\bullet}},s)) \to \mathrm{Conv}(\Sigma {(r)^{\bullet}},s) \to \varphi_{22} \to r \in \mathbb{R}$
\item $\varphi_{73} :\equiv \mathrm{Conv}(\Sigma {(r)^{\bullet}},s) \to \varphi_{22} \to \lvert {(r)^{\bullet}} \rvert(N) = \lvert {{r}^{N}} \rvert$
\item $\varphi_{74} :\equiv \exists s\,(\mathrm{Conv}(\Sigma {(r)^{\bullet}},s)) \to \mathrm{Conv}(\Sigma {(r)^{\bullet}},s) \to \varphi_{34}$
\item $\varphi_{75} :\equiv \forall \alpha\,\forall \varepsilon\,(\mathrm{Conv}(\lvert {(r)^{\bullet}} \rvert,\alpha) \to \varepsilon \in \mathbb{R} \to 0 < \varepsilon \to \exists N \in \mathbb{N}\,\forall n \in \mathbb{N}\,(N \leq n \to \lvert \lvert {(r)^{\bullet}} \rvert(n) - \alpha \rvert < \varepsilon))$
\item $\varphi_{76} :\equiv \lvert {{r}^{N}} \rvert = {{\lvert r \rvert}^{N}} \to \lvert \lvert {{r}^{N}} \rvert \rvert = \lvert \lvert {{r}^{N}} \rvert \rvert \to \lvert \lvert {{r}^{N}} \rvert \rvert = \lvert {{\lvert r \rvert}^{N}} \rvert$
\item $\varphi_{77} :\equiv \exists s\,(\mathrm{Conv}(\Sigma {(r)^{\bullet}},s)) \to \mathrm{Conv}(\Sigma {(r)^{\bullet}},s) \to \varphi_{22} \to \lvert r \rvert \in \mathbb{R}$
\item $\varphi_{78} :\equiv \mathrm{Conv}(\Sigma {(r)^{\bullet}},s) \to \varphi_{22} \to \lvert {{\lvert r \rvert}^{N}} \rvert = {{\lvert r \rvert}^{N}}$
\item $\varphi_{79} :\equiv \forall r \in \mathbb{R}\,\forall n \in \mathbb{N}\,(\lvert {{r}^{n}} \rvert = {{\lvert r \rvert}^{n}}) \to r \in \mathbb{R} \to \forall n \in \mathbb{N}\,(\lvert {{r}^{n}} \rvert = {{\lvert r \rvert}^{n}})$
\item $\varphi_{80} :\equiv \lvert {(r)^{\bullet}} \rvert(N) = \lvert {{r}^{N}} \rvert \to \varphi_{10} \to \lvert \lvert {(r)^{\bullet}} \rvert(N) \rvert = \lvert \lvert {{r}^{N}} \rvert \rvert$
\item $\varphi_{81} :\equiv \exists s\,(\mathrm{Conv}(\Sigma {(r)^{\bullet}},s)) \to \mathrm{Conv}(\Sigma {(r)^{\bullet}},s) \to \varphi_{22} \to 1 \leq {{\lvert r \rvert}^{N}}$
\item $\varphi_{82} :\equiv \exists s\,(\mathrm{Conv}(\Sigma {(r)^{\bullet}},s)) \to \mathrm{Conv}(\Sigma {(r)^{\bullet}},s) \to \varphi_{22} \to 0 \leq {{\lvert r \rvert}^{N}}$
\item $\varphi_{83} :\equiv \mathrm{Conv}(\Sigma {(r)^{\bullet}},s) \to \varphi_{22} \to \lvert \lvert {{r}^{N}} \rvert \rvert = \lvert {{\lvert r \rvert}^{N}} \rvert$
\item $\varphi_{84} :\equiv \mathrm{Conv}(\Sigma {(r)^{\bullet}},s) \to \varphi_{22} \to \lvert \lvert {(r)^{\bullet}} \rvert(N) \rvert = \lvert \lvert {{r}^{N}} \rvert \rvert$
\item $\varphi_{85} :\equiv \mathrm{Conv}(\Sigma {(r)^{\bullet}},s) \to \varphi_{22} \to \lvert \lvert {(r)^{\bullet}} \rvert(N) - 0 \rvert = {{\lvert r \rvert}^{N}}$
\item $\varphi_{86} :\equiv \mathrm{Conv}(\Sigma {(r)^{\bullet}},s) \to \varphi_{22} \to N \in \mathbb{N} \to \lvert {{r}^{N}} \rvert = {{\lvert r \rvert}^{N}}$
\item $\varphi_{87} :\equiv \exists s\,(\mathrm{Conv}(\Sigma {(r)^{\bullet}},s)) \to \mathrm{Conv}(\Sigma {(r)^{\bullet}},s) \to \varphi_{22} \to {{\lvert r \rvert}^{N}} < 1$
\item $\varphi_{88} :\equiv \exists s\,(\mathrm{Conv}(\Sigma {(r)^{\bullet}},s)) \to \mathrm{Conv}(\Sigma {(r)^{\bullet}},s) \to \varphi_{22} \to \neg {{\lvert r \rvert}^{N}} < 1$
\item $\varphi_{89} :\equiv \mathrm{Conv}(\Sigma {(r)^{\bullet}},s) \to \varphi_{22} \to N \in \mathbb{N} \to \lvert {(r)^{\bullet}} \rvert(N) = \lvert {{r}^{N}} \rvert$
\item $\varphi_{90} :\equiv \mathrm{Conv}(\Sigma {(r)^{\bullet}},s) \to \varphi_{22} \to \forall n \in \mathbb{N}\,(\lvert {{r}^{n}} \rvert = {{\lvert r \rvert}^{n}})$
\item $\varphi_{91} :\equiv \mathrm{Conv}(\Sigma {(r)^{\bullet}},s) \to \varphi_{22} \to N \in \mathbb{N} \to N \leq N \to \lvert \lvert {(r)^{\bullet}} \rvert(N) - 0 \rvert < 1$
\item $\varphi_{92} :\equiv \mathrm{Conv}(\Sigma {(r)^{\bullet}},s) \to \varphi_{22} \to \forall n \in \mathbb{N}\,(\lvert {(r)^{\bullet}} \rvert(n) = \lvert {{r}^{n}} \rvert)$
\item $\varphi_{93} :\equiv \mathrm{Conv}(\Sigma {(r)^{\bullet}},s) \to \varphi_{22} \to \lvert \lvert {(r)^{\bullet}} \rvert(N) - 0 \rvert = \lvert \lvert {{r}^{N}} \rvert \rvert$
\item $\varphi_{94} :\equiv \mathrm{Conv}(\Sigma {(r)^{\bullet}},s) \to \varphi_{22} \to \lvert \lvert {(r)^{\bullet}} \rvert(N) - 0 \rvert = \lvert {{\lvert r \rvert}^{N}} \rvert$
\item $\varphi_{95} :\equiv r \in \mathbb{R} \to 1 \leq \lvert r \rvert \to \exists s\,(\mathrm{Conv}(\Sigma {(r)^{\bullet}},s)) \to \varphi_{46}$
\item $\varphi_{96} :\equiv \exists s\,(\mathrm{Conv}(\Sigma {(r)^{\bullet}},s)) \to \mathrm{Conv}(\Sigma {(r)^{\bullet}},s) \to \varphi_{22} \to \lvert {{r}^{N}} \rvert = {{\lvert r \rvert}^{N}}$
\item $\varphi_{97} :\equiv r \in \mathbb{R} \to 1 \leq \lvert r \rvert \to \exists s\,(\mathrm{Conv}(\Sigma {(r)^{\bullet}},s)) \to \varphi_{52}$
\item $\varphi_{98} :\equiv \lvert {(r)^{\bullet}} \rvert(N) - 0 = \lvert {(r)^{\bullet}} \rvert(N) \to \varphi_{17} \to \varphi_{16}$
\item $\varphi_{99} :\equiv \exists s\,(\mathrm{Conv}(\Sigma {(r)^{\bullet}},s)) \to \mathrm{Conv}(\Sigma {(r)^{\bullet}},s) \to \varphi_{54}$
\item $\varphi_{100} :\equiv r \in \mathbb{R} \to 1 \leq \lvert r \rvert \to \exists s\,(\mathrm{Conv}(\Sigma {(r)^{\bullet}},s)) \to \varphi_{58}$
\item $\varphi_{101} :\equiv r \in \mathbb{R} \to 1 \leq \lvert r \rvert \to \exists s\,(\mathrm{Conv}(\Sigma {(r)^{\bullet}},s)) \to \varphi_{62}$
\item $\varphi_{102} :\equiv \exists s\,(\mathrm{Conv}(\Sigma {(r)^{\bullet}},s)) \to \mathrm{Conv}(\Sigma {(r)^{\bullet}},s) \to \varphi_{22} \to \varphi_{16}$
\item $\varphi_{103} :\equiv \exists s\,(\mathrm{Conv}(\Sigma {(r)^{\bullet}},s)) \to \mathrm{Conv}(\Sigma {(r)^{\bullet}},s) \to \varphi_{66}$
\item $\varphi_{104} :\equiv r \in \mathbb{R} \to 1 \leq \lvert r \rvert \to \exists s\,(\mathrm{Conv}(\Sigma {(r)^{\bullet}},s)) \to \varphi_{65}$
\item $\varphi_{105} :\equiv r \in \mathbb{R} \to 1 \leq \lvert r \rvert \to \exists s\,(\mathrm{Conv}(\Sigma {(r)^{\bullet}},s)) \to \varphi_{68}$
\item $\varphi_{106} :\equiv r \in \mathbb{R} \to 1 \leq \lvert r \rvert \to \exists s\,(\mathrm{Conv}(\Sigma {(r)^{\bullet}},s)) \to \varphi_{70}$
\item $\varphi_{107} :\equiv r \in \mathbb{R} \to 1 \leq \lvert r \rvert \to \exists s\,(\mathrm{Conv}(\Sigma {(r)^{\bullet}},s)) \to \varphi_{73}$
\item $\varphi_{108} :\equiv r \in \mathbb{R} \to 1 \leq \lvert r \rvert \to \exists s\,(\mathrm{Conv}(\Sigma {(r)^{\bullet}},s)) \to \varphi_{78}$
\item $\varphi_{109} :\equiv \exists s\,(\mathrm{Conv}(\Sigma {(r)^{\bullet}},s)) \to \mathrm{Conv}(\Sigma {(r)^{\bullet}},s) \to \varphi_{22} \to \varphi_{22}$
\item $\varphi_{110} :\equiv r \in \mathbb{R} \to 1 \leq \lvert r \rvert \to \exists s\,(\mathrm{Conv}(\Sigma {(r)^{\bullet}},s)) \to \varphi_{83}$
\item $\varphi_{111} :\equiv r \in \mathbb{R} \to 1 \leq \lvert r \rvert \to \exists s\,(\mathrm{Conv}(\Sigma {(r)^{\bullet}},s)) \to \varphi_{84}$
\item $\varphi_{112} :\equiv r \in \mathbb{R} \to 1 \leq \lvert r \rvert \to \exists s\,(\mathrm{Conv}(\Sigma {(r)^{\bullet}},s)) \to \varphi_{85}$
\item $\varphi_{113} :\equiv \exists s\,(\mathrm{Conv}(\Sigma {(r)^{\bullet}},s)) \to \mathrm{Conv}(\Sigma {(r)^{\bullet}},s) \to \varphi_{22} \to \varphi_{26}$
\item $\varphi_{114} :\equiv r \in \mathbb{R} \to 1 \leq \lvert r \rvert \to \exists s\,(\mathrm{Conv}(\Sigma {(r)^{\bullet}},s)) \to \varphi_{86}$
\item $\varphi_{115} :\equiv r \in \mathbb{R} \to 1 \leq \lvert r \rvert \to \exists s\,(\mathrm{Conv}(\Sigma {(r)^{\bullet}},s)) \to \varphi_{89}$
\item $\varphi_{116} :\equiv r \in \mathbb{R} \to 1 \leq \lvert r \rvert \to \exists s\,(\mathrm{Conv}(\Sigma {(r)^{\bullet}},s)) \to \varphi_{90}$
\item $\varphi_{117} :\equiv r \in \mathbb{R} \to 1 \leq \lvert r \rvert \to \exists s\,(\mathrm{Conv}(\Sigma {(r)^{\bullet}},s)) \to \varphi_{91}$
\item $\varphi_{118} :\equiv r \in \mathbb{R} \to 1 \leq \lvert r \rvert \to \exists s\,(\mathrm{Conv}(\Sigma {(r)^{\bullet}},s)) \to \varphi_{92}$
\item $\varphi_{119} :\equiv r \in \mathbb{R} \to 1 \leq \lvert r \rvert \to \exists s\,(\mathrm{Conv}(\Sigma {(r)^{\bullet}},s)) \to \varphi_{93}$
\item $\varphi_{120} :\equiv r \in \mathbb{R} \to 1 \leq \lvert r \rvert \to \exists s\,(\mathrm{Conv}(\Sigma {(r)^{\bullet}},s)) \to \varphi_{94}$
\item $\varphi_{121} :\equiv \exists s\,(\mathrm{Conv}(\Sigma {(r)^{\bullet}},s)) \to \mathrm{Conv}(\Sigma {(r)^{\bullet}},s) \to \varphi_{22} \to \varphi_{34}$
\item $\varphi_{122} :\equiv \forall s\,(\varphi_{29} \to \neg \mathrm{Conv}(\Sigma {(r)^{\bullet}},s)) \to \forall s\,(\varphi_{29}) \to \forall s\,(\neg \mathrm{Conv}(\Sigma {(r)^{\bullet}},s))$
\item $\varphi_{123} :\equiv \varphi_{67} \to \mathrm{Conv}(\lvert {(r)^{\bullet}} \rvert,0) \to 1 \in \mathbb{R} \to 0 < 1 \to \varphi_{24}$
\item $\varphi_{124} :\equiv (\mathrm{Seq}({(r)^{\bullet}}) \to \varphi_{3}) \to \varphi_{33} \to r \in \mathbb{R} \to 1 \leq \lvert r \rvert \to \varphi_{36}$
\item $\varphi_{125} :\equiv (r \in \mathbb{R} \to \varphi_{34}) \to (r \in \mathbb{R} \to 1 \leq \lvert r \rvert \to \varphi_{18}) \to r \in \mathbb{R} \to 1 \leq \lvert r \rvert \to \varphi_{74}$
\item $\varphi_{126} :\equiv (r \in \mathbb{R} \to \lvert r \rvert \in \mathbb{R}) \to (r \in \mathbb{R} \to 1 \leq \lvert r \rvert \to \varphi_{72}) \to r \in \mathbb{R} \to 1 \leq \lvert r \rvert \to \varphi_{77}$
\item $\varphi_{127} :\equiv (r \in \mathbb{R} \to 1 \leq \lvert r \rvert \to \varphi_{77}) \to (r \in \mathbb{R} \to 1 \leq \lvert r \rvert \to 1 \leq \lvert r \rvert) \to \varphi_{106}$
\item $\varphi_{128} :\equiv (\mathrm{Conv}(\lvert {(r)^{\bullet}} \rvert,0) \to 1 \in \mathbb{R} \to 0 < 1 \to \varphi_{24}) \to \varphi_{40} \to 1 \in \mathbb{R} \to 0 < 1 \to r \in \mathbb{R} \to 1 \leq \lvert r \rvert \to \varphi_{60}$
\item $\varphi_{129} :\equiv (1 \leq {{\lvert r \rvert}^{N}} \to \neg {{\lvert r \rvert}^{N}} < 1) \to (r \in \mathbb{R} \to 1 \leq \lvert r \rvert \to \varphi_{81}) \to r \in \mathbb{R} \to 1 \leq \lvert r \rvert \to \varphi_{88}$
\item $\varphi_{130} :\equiv (0 \leq {{\lvert r \rvert}^{N}} \to \lvert {{\lvert r \rvert}^{N}} \rvert = {{\lvert r \rvert}^{N}}) \to (r \in \mathbb{R} \to 1 \leq \lvert r \rvert \to \varphi_{82}) \to \varphi_{108}$
\item $\varphi_{131} :\equiv (r \in \mathbb{R} \to \forall n \in \mathbb{N}\,(\lvert {{r}^{n}} \rvert = {{\lvert r \rvert}^{n}})) \to (r \in \mathbb{R} \to 1 \leq \lvert r \rvert \to \varphi_{72}) \to \varphi_{116}$
\item $\varphi_{132} :\equiv \varphi_{106} \to (\forall n \in \mathbb{N}\,(1 \leq {{\lvert r \rvert}^{n}}) \to N \in \mathbb{N} \to 1 \leq {{\lvert r \rvert}^{N}}) \to \varphi_{104}$
\item $\varphi_{133} :\equiv (\mathrm{Seq}({(r)^{\bullet}}) \to \varphi_{26}) \to \varphi_{95} \to r \in \mathbb{R} \to 1 \leq \lvert r \rvert \to \varphi_{113}$
\item $\varphi_{134} :\equiv (r \in \mathbb{R} \to 1 \leq \lvert r \rvert \to \varphi_{64}) \to (r \in \mathbb{R} \to 1 \leq \lvert r \rvert \to \varphi_{81}) \to r \in \mathbb{R} \to 1 \leq \lvert r \rvert \to \varphi_{82}$
\item $\varphi_{135} :\equiv (r \in \mathbb{R} \to 1 \leq \lvert r \rvert \to \varphi_{69}) \to (r \in \mathbb{R} \to 1 \leq \lvert r \rvert \to \varphi_{63}) \to \varphi_{105}$
\item $\varphi_{136} :\equiv (r \in \mathbb{R} \to \varphi_{34}) \to (r \in \mathbb{R} \to 1 \leq \lvert r \rvert \to \varphi_{72}) \to r \in \mathbb{R} \to 1 \leq \lvert r \rvert \to \varphi_{121}$
\item $\varphi_{137} :\equiv (r \in \mathbb{R} \to 1 \leq \lvert r \rvert \to \varphi_{87}) \to (r \in \mathbb{R} \to 1 \leq \lvert r \rvert \to \varphi_{88}) \to \varphi_{100}$
\item $\varphi_{138} :\equiv \varphi_{108} \to (\lvert \lvert {(r)^{\bullet}} \rvert(N) - 0 \rvert = \lvert {{\lvert r \rvert}^{N}} \rvert \to \varphi_{38}) \to \varphi_{112}$
\item $\varphi_{139} :\equiv \varphi_{110} \to (\lvert \lvert {(r)^{\bullet}} \rvert(N) - 0 \rvert = \lvert \lvert {{r}^{N}} \rvert \rvert \to \varphi_{43}) \to \varphi_{120}$
\item $\varphi_{140} :\equiv \forall s\,(\varphi_{29} \to \neg \mathrm{Conv}(\Sigma {(r)^{\bullet}},s)) \to \varphi_{122} \to (\varphi_{29} \to \forall s\,(\varphi_{29})) \to r \in \mathbb{R} \to 1 \leq \lvert r \rvert \to \varphi_{14}$
\item $\varphi_{141} :\equiv (r \in \mathbb{R} \to 1 \leq \lvert r \rvert \to \varphi_{96}) \to \lvert \lvert {{r}^{N}} \rvert \rvert = \lvert \lvert {{r}^{N}} \rvert \rvert \to \varphi_{76} \to \varphi_{110}$
\item $\varphi_{142} :\equiv (r \in \mathbb{R} \to 1 \leq \lvert r \rvert \to \varphi_{99}) \to \varphi_{17} \to \varphi_{98} \to r \in \mathbb{R} \to 1 \leq \lvert r \rvert \to \varphi_{102}$
\item $\varphi_{143} :\equiv (\forall N\,(\varphi_{42} \to \neg \varphi_{22}) \to \forall N\,(\varphi_{42}) \to \forall N\,(\neg \varphi_{22})) \to (\varphi_{42} \to \forall N\,(\varphi_{42})) \to \varphi_{41}$
\end{abbrevlist}
\end{fullproof}

\section{第3章 5節・6節: 単調収束定理, 区間縮小法, Bolzano–Weierstrass, Cauchy 列}
% 自動生成 (prover/build.js): MonoIncLe  ノード数 428
\begin{fullproof}{[C20]}{$\forall a\,\forall m\,\forall n\,(\mathrm{MonoInc}(a) \to m \in \mathbb{N} \to n \in \mathbb{N} \to m \leq n \to a_{m} \leq a_{n})$}{428}\phantomsection\label{tree:MonoIncLe}
\noindent$\Pi_{1}$:
\begin{ptree}
\AXC{$\mathrm{MonoInc}(a) \to m \in \mathbb{N} \to n \in \mathbb{N} \to n \in \mathbb{N}$\lbl{A1}}
\AXC{$(\mathrm{MonoInc}(a) \to m \in \mathbb{N} \to n \in \mathbb{N} \to n \in \mathbb{N}) \to \varphi_{72}$\lbl{A1}}
\BIC{$\varphi_{72}$}
\end{ptree}
\noindent$\Pi_{2}$:
\begin{ptree}
\AXC{$\varphi_{30}$\lbl{A4}}
\AXC{$\forall y\,(\varphi_{15}) \to \varphi_{34}$\lbl{A3}}
\AXC{$\varphi_{15}$\lbl{A2}}
\UIC{$\forall y\,(\varphi_{15})$}
\AXC{$\forall y\,(\varphi_{15}) \to (\forall y\,(\varphi_{15}) \to \varphi_{34}) \to \varphi_{30} \to \varphi_{25}$\lbl{A1}}
\BIC{$(\forall y\,(\varphi_{15}) \to \varphi_{34}) \to \varphi_{30} \to \varphi_{25}$}
\BIC{$\varphi_{30} \to \varphi_{25}$}
\BIC{$\varphi_{25}$}
\end{ptree}
\noindent$\Pi_{3}$:
\begin{ptree}
\AXC{$\varphi_{31}$\lbl{A4}}
\AXC{$\forall n\,(\varphi_{16}) \to \varphi_{35}$\lbl{A3}}
\AXC{$\varphi_{16}$\lbl{A2}}
\UIC{$\forall n\,(\varphi_{16})$}
\AXC{$\forall n\,(\varphi_{16}) \to (\forall n\,(\varphi_{16}) \to \varphi_{35}) \to \varphi_{31} \to \varphi_{26}$\lbl{A1}}
\BIC{$(\forall n\,(\varphi_{16}) \to \varphi_{35}) \to \varphi_{31} \to \varphi_{26}$}
\BIC{$\varphi_{31} \to \varphi_{26}$}
\BIC{$\varphi_{26}$}
\end{ptree}
\noindent$\Pi_{4}$:
\begin{ptree}
\AXC{$x = x$\lbl{A5}}
\UIC{$\forall x\,(x = x)$}
\AXC{$\forall x\,(x = x) \to a_{m} = a_{m}$\lbl{A2}}
\BIC{$a_{m} = a_{m}$}
\AXC{$a_{m} = a_{m} \to \mathrm{MonoInc}(a) \to m \in \mathbb{N} \to \varphi_{136}$\lbl{A1}}
\BIC{$\mathrm{MonoInc}(a) \to m \in \mathbb{N} \to \varphi_{136}$}
\end{ptree}
\noindent$\Pi_{5}$:
\begin{ptree}
\AXC{$\stackrel{\Pi_{4}}{\vdots}$}\noLine
\UIC{$\mathrm{MonoInc}(a) \to m \in \mathbb{N} \to \varphi_{136}$}
\AXC{$\mathrm{MonoInc}(a) \to m \in \mathbb{N} \to \varphi_{129}$\lbl{A1}}
\AXC{$(\mathrm{MonoInc}(a) \to m \in \mathbb{N} \to \varphi_{129}) \to \varphi_{185}$\lbl{A1}}
\BIC{$\varphi_{185}$}
\BIC{$(m = S(n) \to a_{m} < a_{m} \to a_{m} < a_{S(n)}) \to \varphi_{153}$}
\end{ptree}
\noindent$\Pi_{6}$:
\begin{ptree}
\AXC{$z = w \to a_{m} < a_{z} \to a_{m} < a_{w}$\lbl{A6}}
\UIC{$\forall w\,(z = w \to a_{m} < a_{z} \to a_{m} < a_{w})$}
\UIC{$\forall z\,\forall w\,(z = w \to a_{m} < a_{z} \to a_{m} < a_{w})$}
\AXC{$\varphi_{95}$\lbl{A2}}
\BIC{$\forall w\,(m = w \to a_{m} < a_{m} \to a_{m} < a_{w})$}
\AXC{$\varphi_{96}$\lbl{A2}}
\BIC{$m = S(n) \to a_{m} < a_{m} \to a_{m} < a_{S(n)}$}
\AXC{$\stackrel{\Pi_{5}}{\vdots}$}\noLine
\UIC{$(m = S(n) \to a_{m} < a_{m} \to a_{m} < a_{S(n)}) \to \varphi_{153}$}
\BIC{$\varphi_{153}$}
\end{ptree}
\noindent$\Pi_{7}$:
\begin{ptree}
\AXC{$\varphi_{122}$\lbl{A2}}
\AXC{$\mathrm{MonoInc}(a) \to m \in \mathbb{N} \to \varphi_{127}$\lbl{A1}}
\AXC{$\varphi_{183}$\lbl{A1}}
\BIC{$\varphi_{122} \to \mathrm{MonoInc}(a) \to m \in \mathbb{N} \to \varphi_{122}$}
\BIC{$\mathrm{MonoInc}(a) \to m \in \mathbb{N} \to \varphi_{122}$}
\AXC{$(\mathrm{MonoInc}(a) \to m \in \mathbb{N} \to \varphi_{122}) \to \varphi_{160}$\lbl{A1}}
\BIC{$\varphi_{160}$}
\end{ptree}
\noindent$\Pi_{8}$:
\begin{ptree}
\AXC{$\mathrm{MonoInc}(a) \to m \in \mathbb{N} \to \varphi_{73}$\lbl{A1}}
\AXC{$\stackrel{\Pi_{7}}{\vdots}$}\noLine
\UIC{$\varphi_{160}$}
\BIC{$\mathrm{MonoInc}(a) \to m \in \mathbb{N} \to \varphi_{123}$}
\end{ptree}
\noindent$\Pi_{9}$:
\begin{ptree}
\AXC{$\stackrel{\Pi_{8}}{\vdots}$}\noLine
\UIC{$\mathrm{MonoInc}(a) \to m \in \mathbb{N} \to \varphi_{123}$}
\AXC{$\varphi_{161}$\lbl{A1}}
\BIC{$\mathrm{MonoInc}(a) \to m \in \mathbb{N} \to \varphi_{76}$}
\end{ptree}
\noindent$\Pi_{10}$:
\begin{ptree}
\AXC{$\stackrel{\Pi_{9}}{\vdots}$}\noLine
\UIC{$\mathrm{MonoInc}(a) \to m \in \mathbb{N} \to \varphi_{76}$}
\AXC{$\forall n\,(n \in \mathbb{N} \leftrightarrow n \in \omega \land n \neq \emptyset)$\lbl{Def$\mathbb{N}$}}
\AXC{$\varphi_{67}$\lbl{A2}}
\BIC{$n \in \mathbb{N} \leftrightarrow n \in \omega \land n \neq \emptyset$}
\AXC{$\varphi_{184}$\lbl{A1}}
\BIC{$(\mathrm{MonoInc}(a) \to m \in \mathbb{N} \to \varphi_{76}) \to \varphi_{164}$}
\BIC{$\varphi_{164}$}
\end{ptree}
\noindent$\Pi_{11}$:
\begin{ptree}
\AXC{$\mathrm{MonoInc}(a) \to m \in \mathbb{N} \to \varphi_{111}$\lbl{A1}}
\AXC{$\stackrel{\Pi_{10}}{\vdots}$}\noLine
\UIC{$\varphi_{164}$}
\BIC{$\mathrm{MonoInc}(a) \to m \in \mathbb{N} \to \varphi_{109}$}
\end{ptree}
\noindent$\Pi_{12}$:
\begin{ptree}
\AXC{$\forall a\,(\varphi_{27})$\lbl{DefMonoInc}}
\AXC{$\forall a\,(\varphi_{27}) \to \varphi_{27}$\lbl{A2}}
\BIC{$\varphi_{27}$}
\AXC{$\mathrm{MonoInc}(a) \to \mathrm{MonoInc}(a)$\lbl{A1}}
\AXC{$\varphi_{158}$\lbl{A1}}
\BIC{$\varphi_{27} \to \mathrm{MonoInc}(a) \to m \in \mathbb{N} \to \varphi_{140}$}
\BIC{$\mathrm{MonoInc}(a) \to m \in \mathbb{N} \to \varphi_{140}$}
\end{ptree}
\noindent$\Pi_{13}$:
\begin{ptree}
\AXC{$\stackrel{\Pi_{12}}{\vdots}$}\noLine
\UIC{$\mathrm{MonoInc}(a) \to m \in \mathbb{N} \to \varphi_{140}$}
\AXC{$(\mathrm{MonoInc}(a) \to m \in \mathbb{N} \to \varphi_{140}) \to \varphi_{157}$\lbl{A1}}
\BIC{$\varphi_{157}$}
\end{ptree}
\noindent$\Pi_{14}$:
\begin{ptree}
\AXC{$\forall n \in \mathbb{N}\,(a_{n} \leq a_{S(n)}) \to n \in \mathbb{N} \to a_{n} \leq a_{S(n)}$\lbl{A2}}
\AXC{$\stackrel{\Pi_{13}}{\vdots}$}\noLine
\UIC{$\varphi_{157}$}
\BIC{$\mathrm{MonoInc}(a) \to m \in \mathbb{N} \to \varphi_{139}$}
\AXC{$(\mathrm{MonoInc}(a) \to m \in \mathbb{N} \to \varphi_{139}) \to \varphi_{172}$\lbl{A1}}
\BIC{$\varphi_{172}$}
\end{ptree}
\noindent$\Pi_{15}$:
\begin{ptree}
\AXC{$\stackrel{\Pi_{8}}{\vdots}$}\noLine
\UIC{$\mathrm{MonoInc}(a) \to m \in \mathbb{N} \to \varphi_{123}$}
\AXC{$\varphi_{166}$\lbl{A1}}
\BIC{$\mathrm{MonoInc}(a) \to m \in \mathbb{N} \to \varphi_{103}$}
\AXC{$(\mathrm{MonoInc}(a) \to m \in \mathbb{N} \to \varphi_{103}) \to \varphi_{169}$\lbl{A1}}
\BIC{$\varphi_{169}$}
\end{ptree}
\noindent$\Pi_{16}$:
\begin{ptree}
\AXC{$\stackrel{\Pi_{11}}{\vdots}$}\noLine
\UIC{$\mathrm{MonoInc}(a) \to m \in \mathbb{N} \to \varphi_{109}$}
\AXC{$\stackrel{\Pi_{15}}{\vdots}$}\noLine
\UIC{$\varphi_{169}$}
\BIC{$\mathrm{MonoInc}(a) \to m \in \mathbb{N} \to \varphi_{128}$}
\AXC{$(\mathrm{MonoInc}(a) \to m \in \mathbb{N} \to \varphi_{128}) \to \varphi_{174}$\lbl{A1}}
\BIC{$\varphi_{174}$}
\end{ptree}
\noindent$\Pi_{17}$:
\begin{ptree}
\AXC{$\forall x\,\forall y\,\forall z\,(x \leq y \to y \leq z \to x \leq z)$\lbl{[O5]}}
\AXC{$\varphi_{36}$\lbl{A2}}
\BIC{$\forall y\,\forall z\,(a_{m} \leq y \to y \leq z \to a_{m} \leq z)$}
\AXC{$\varphi_{71}$\lbl{A2}}
\BIC{$\forall z\,(a_{m} \leq a_{n} \to a_{n} \leq z \to a_{m} \leq z)$}
\AXC{$\varphi_{102}$\lbl{A2}}
\BIC{$a_{m} \leq a_{n} \to a_{n} \leq a_{S(n)} \to a_{m} \leq a_{S(n)}$}
\AXC{$\varphi_{193}$\lbl{A1}}
\BIC{$(\mathrm{MonoInc}(a) \to m \in \mathbb{N} \to \varphi_{134}) \to \varphi_{175}$}
\end{ptree}
\noindent$\Pi_{18}$:
\begin{ptree}
\AXC{$\mathrm{MonoInc}(a) \to m \in \mathbb{N} \to \varphi_{121}$\lbl{A1}}
\AXC{$\stackrel{\Pi_{16}}{\vdots}$}\noLine
\UIC{$\varphi_{174}$}
\BIC{$\mathrm{MonoInc}(a) \to m \in \mathbb{N} \to \varphi_{134}$}
\AXC{$\stackrel{\Pi_{17}}{\vdots}$}\noLine
\UIC{$(\mathrm{MonoInc}(a) \to m \in \mathbb{N} \to \varphi_{134}) \to \varphi_{175}$}
\BIC{$\varphi_{175}$}
\end{ptree}
\noindent$\Pi_{19}$:
\begin{ptree}
\AXC{$\forall n \in \mathbb{N}\,(n \in \omega)$\lbl{[N7']}}
\AXC{$\forall n \in \mathbb{N}\,(n \in \omega) \to n \in \mathbb{N} \to n \in \omega$\lbl{A2}}
\BIC{$n \in \mathbb{N} \to n \in \omega$}
\AXC{$(n \in \mathbb{N} \to n \in \omega) \to \varphi_{168}$\lbl{A1}}
\BIC{$\varphi_{168}$}
\end{ptree}
\noindent$\Pi_{20}$:
\begin{ptree}
\AXC{$\forall n \in \mathbb{N}\,(n \in \omega)$\lbl{[N7']}}
\AXC{$\forall n \in \mathbb{N}\,(n \in \omega) \to m \in \mathbb{N} \to m \in \omega$\lbl{A2}}
\BIC{$m \in \mathbb{N} \to m \in \omega$}
\AXC{$(m \in \mathbb{N} \to m \in \omega) \to \varphi_{148}$\lbl{A1}}
\BIC{$\varphi_{148}$}
\end{ptree}
\noindent$\Pi_{21}$:
\begin{ptree}
\AXC{$\varphi_{22}$\lbl{[N9]}}
\AXC{$\varphi_{97}$\lbl{A2}}
\BIC{$\forall n\,(m \in \omega \to n \in \omega \to m \leq S(n) \to m \leq n \lor m = S(n))$}
\AXC{$\varphi_{90}$\lbl{A2}}
\BIC{$m \in \omega \to n \in \omega \to m \leq S(n) \to m \leq n \lor m = S(n)$}
\AXC{$(m \in \omega \to n \in \omega \to m \leq S(n) \to m \leq n \lor m = S(n)) \to \varphi_{195}$\lbl{A1}}
\BIC{$\varphi_{195}$}
\end{ptree}
\noindent$\Pi_{22}$:
\begin{ptree}
\AXC{$\stackrel{\Pi_{11}}{\vdots}$}\noLine
\UIC{$\mathrm{MonoInc}(a) \to m \in \mathbb{N} \to \varphi_{109}$}
\AXC{$\stackrel{\Pi_{19}}{\vdots}$}\noLine
\UIC{$\varphi_{168}$}
\BIC{$\mathrm{MonoInc}(a) \to m \in \mathbb{N} \to \varphi_{124}$}
\AXC{$\mathrm{MonoInc}(a) \to m \in \mathbb{N} \to m \in \mathbb{N}$\lbl{A1}}
\AXC{$\stackrel{\Pi_{20}}{\vdots}$}\noLine
\UIC{$\varphi_{148}$}
\BIC{$\mathrm{MonoInc}(a) \to m \in \mathbb{N} \to \varphi_{125}$}
\AXC{$\stackrel{\Pi_{21}}{\vdots}$}\noLine
\UIC{$\varphi_{195}$}
\BIC{$(\mathrm{MonoInc}(a) \to m \in \mathbb{N} \to \varphi_{124}) \to \varphi_{170}$}
\BIC{$\varphi_{170}$}
\end{ptree}
\noindent$\Pi_{23}$:
\begin{ptree}
\AXC{$\stackrel{\Pi_{11}}{\vdots}$}\noLine
\UIC{$\mathrm{MonoInc}(a) \to m \in \mathbb{N} \to \varphi_{109}$}
\AXC{$\stackrel{\Pi_{14}}{\vdots}$}\noLine
\UIC{$\varphi_{172}$}
\BIC{$\mathrm{MonoInc}(a) \to m \in \mathbb{N} \to \varphi_{131}$}
\AXC{$\stackrel{\Pi_{18}}{\vdots}$}\noLine
\UIC{$\varphi_{175}$}
\BIC{$\mathrm{MonoInc}(a) \to m \in \mathbb{N} \to \varphi_{137}$}
\AXC{$\mathrm{MonoInc}(a) \to m \in \mathbb{N} \to \varphi_{118}$\lbl{A1}}
\AXC{$\stackrel{\Pi_{22}}{\vdots}$}\noLine
\UIC{$\varphi_{170}$}
\BIC{$\mathrm{MonoInc}(a) \to m \in \mathbb{N} \to \varphi_{126}$}
\AXC{$\varphi_{196}$\lbl{A1}}
\BIC{$(\mathrm{MonoInc}(a) \to m \in \mathbb{N} \to \varphi_{137}) \to \varphi_{176}$}
\BIC{$\varphi_{176}$}
\end{ptree}
\noindent$\Pi_{24}$:
\begin{ptree}
\AXC{$x = x$\lbl{A5}}
\AXC{$x = y \to x = x \to y = x$\lbl{A6}}
\AXC{$(x = y \to x = x \to y = x) \to x = x \to x = y \to y = x$\lbl{A1}}
\BIC{$x = x \to x = y \to y = x$}
\BIC{$x = y \to y = x$}
\UIC{$\forall y\,(x = y \to y = x)$}
\UIC{$\forall x\,\forall y\,(x = y \to y = x)$}
\end{ptree}
\noindent$\Pi_{25}$:
\begin{ptree}
\AXC{$y = z \to x = y \to x = z$\lbl{A6}}
\AXC{$(y = z \to x = y \to x = z) \to x = y \to y = z \to x = z$\lbl{A1}}
\BIC{$x = y \to y = z \to x = z$}
\UIC{$\forall z\,(x = y \to y = z \to x = z)$}
\UIC{$\forall y\,\forall z\,(x = y \to y = z \to x = z)$}
\UIC{$\forall x\,\forall y\,\forall z\,(x = y \to y = z \to x = z)$}
\AXC{$\varphi_{17}$\lbl{A2}}
\BIC{$\forall y\,\forall z\,(S(n) = y \to y = z \to S(n) = z)$}
\AXC{$\varphi_{58}$\lbl{A2}}
\BIC{$\forall z\,(S(n) = S(\emptyset) \to S(\emptyset) = z \to S(n) = z)$}
\end{ptree}
\noindent$\Pi_{26}$:
\begin{ptree}
\AXC{$\stackrel{\Pi_{25}}{\vdots}$}\noLine
\UIC{$\forall z\,(S(n) = S(\emptyset) \to S(\emptyset) = z \to S(n) = z)$}
\AXC{$\forall z\,(S(n) = S(\emptyset) \to S(\emptyset) = z \to S(n) = z) \to \varphi_{10}$\lbl{A2}}
\BIC{$\varphi_{10}$}
\end{ptree}
\noindent$\Pi_{27}$:
\begin{ptree}
\AXC{$\stackrel{\Pi_{24}}{\vdots}$}\noLine
\UIC{$\forall x\,\forall y\,(x = y \to y = x)$}
\AXC{$\varphi_{1}$\lbl{A2}}
\BIC{$\forall y\,(1_{\omega} = y \to y = 1_{\omega})$}
\AXC{$\varphi_{23}$\lbl{A2}}
\BIC{$1_{\omega} = S(\emptyset) \to S(\emptyset) = 1_{\omega}$}
\AXC{$1_{\omega} = S(\emptyset)$\lbl{Def$1_\omega$}}
\AXC{$1_{\omega} = S(\emptyset) \to \varphi_{18}$\lbl{A1}}
\BIC{$\varphi_{18}$}
\BIC{$S(\emptyset) = 1_{\omega}$}
\end{ptree}
\noindent$\Pi_{28}$:
\begin{ptree}
\AXC{$x = x$\lbl{A5}}
\UIC{$\forall x\,(x = x)$}
\AXC{$\forall x\,(x = x) \to S(n) = S(n)$\lbl{A2}}
\BIC{$S(n) = S(n)$}
\AXC{$\mathrm{MonoInc}(a) \to m \in \mathbb{N} \to \varphi_{110}$\lbl{A1}}
\AXC{$\varphi_{179}$\lbl{A1}}
\BIC{$S(n) = S(n) \to \varphi_{149}$}
\BIC{$\varphi_{149}$}
\end{ptree}
\noindent$\Pi_{29}$:
\begin{ptree}
\AXC{$z = w \to S(n) = S(z) \to S(n) = S(w)$\lbl{A6}}
\UIC{$\forall w\,(z = w \to S(n) = S(z) \to S(n) = S(w))$}
\UIC{$\forall z\,\forall w\,(z = w \to S(n) = S(z) \to S(n) = S(w))$}
\AXC{$\varphi_{43}$\lbl{A2}}
\BIC{$\forall w\,(n = w \to S(n) = S(n) \to S(n) = S(w))$}
\AXC{$\varphi_{59}$\lbl{A2}}
\BIC{$n = \emptyset \to S(n) = S(n) \to S(n) = S(\emptyset)$}
\AXC{$\stackrel{\Pi_{28}}{\vdots}$}\noLine
\UIC{$\varphi_{149}$}
\BIC{$\mathrm{MonoInc}(a) \to m \in \mathbb{N} \to \varphi_{119}$}
\AXC{$\varphi_{180}$\lbl{A1}}
\BIC{$S(\emptyset) = 1_{\omega} \to \varphi_{10} \to \mathrm{MonoInc}(a) \to m \in \mathbb{N} \to \varphi_{106}$}
\end{ptree}
\noindent$\Pi_{30}$:
\begin{ptree}
\AXC{$\stackrel{\Pi_{26}}{\vdots}$}\noLine
\UIC{$\varphi_{10}$}
\AXC{$\stackrel{\Pi_{27}}{\vdots}$}\noLine
\UIC{$S(\emptyset) = 1_{\omega}$}
\AXC{$\stackrel{\Pi_{29}}{\vdots}$}\noLine
\UIC{$S(\emptyset) = 1_{\omega} \to \varphi_{10} \to \mathrm{MonoInc}(a) \to m \in \mathbb{N} \to \varphi_{106}$}
\BIC{$\varphi_{10} \to \mathrm{MonoInc}(a) \to m \in \mathbb{N} \to \varphi_{106}$}
\BIC{$\mathrm{MonoInc}(a) \to m \in \mathbb{N} \to \varphi_{106}$}
\end{ptree}
\noindent$\Pi_{31}$:
\begin{ptree}
\AXC{$\stackrel{\Pi_{30}}{\vdots}$}\noLine
\UIC{$\mathrm{MonoInc}(a) \to m \in \mathbb{N} \to \varphi_{106}$}
\AXC{$(\mathrm{MonoInc}(a) \to m \in \mathbb{N} \to \varphi_{106}) \to \varphi_{145}$\lbl{A1}}
\BIC{$\varphi_{145}$}
\end{ptree}
\noindent$\Pi_{32}$:
\begin{ptree}
\AXC{$\stackrel{\Pi_{24}}{\vdots}$}\noLine
\UIC{$\forall x\,\forall y\,(x = y \to y = x)$}
\AXC{$\varphi_{2}$\lbl{A2}}
\BIC{$\forall y\,(S(n) = y \to y = S(n))$}
\AXC{$\varphi_{7}$\lbl{A2}}
\BIC{$S(n) = 1_{\omega} \to 1_{\omega} = S(n)$}
\AXC{$\stackrel{\Pi_{31}}{\vdots}$}\noLine
\UIC{$\varphi_{145}$}
\BIC{$\mathrm{MonoInc}(a) \to m \in \mathbb{N} \to \varphi_{107}$}
\end{ptree}
\noindent$\Pi_{33}$:
\begin{ptree}
\AXC{$\stackrel{\Pi_{32}}{\vdots}$}\noLine
\UIC{$\mathrm{MonoInc}(a) \to m \in \mathbb{N} \to \varphi_{107}$}
\AXC{$(\mathrm{MonoInc}(a) \to m \in \mathbb{N} \to \varphi_{107}) \to \varphi_{146}$\lbl{A1}}
\BIC{$\varphi_{146}$}
\end{ptree}
\noindent$\Pi_{34}$:
\begin{ptree}
\AXC{$\forall n \in \mathbb{N}\,(1_{\omega} \leq n)$\lbl{[N12]}}
\AXC{$\forall n \in \mathbb{N}\,(1_{\omega} \leq n) \to m \in \mathbb{N} \to 1_{\omega} \leq m$\lbl{A2}}
\BIC{$m \in \mathbb{N} \to 1_{\omega} \leq m$}
\AXC{$\varphi_{85}$\lbl{A1}}
\BIC{$(\mathrm{MonoInc}(a) \to m \in \mathbb{N} \to m \in \mathbb{N}) \to \varphi_{3}$}
\end{ptree}
\noindent$\Pi_{35}$:
\begin{ptree}
\AXC{$\varphi_{4}$\lbl{A1}}
\AXC{$\forall x\,\forall y\,(x \leq y \to y \leq x \to x = y)$\lbl{[O6]}}
\AXC{$\varphi_{8}$\lbl{A2}}
\BIC{$\forall y\,(m \leq y \to y \leq m \to m = y)$}
\AXC{$\varphi_{9}$\lbl{A2}}
\BIC{$m \leq 1_{\omega} \to 1_{\omega} \leq m \to m = 1_{\omega}$}
\AXC{$\varphi_{141}$\lbl{A1}}
\BIC{$\varphi_{4} \to \varphi_{3} \to \varphi_{5}$}
\BIC{$\varphi_{3} \to \varphi_{5}$}
\end{ptree}
\noindent$\Pi_{36}$:
\begin{ptree}
\AXC{$\mathrm{MonoInc}(a) \to m \in \mathbb{N} \to m \in \mathbb{N}$\lbl{A1}}
\AXC{$\stackrel{\Pi_{34}}{\vdots}$}\noLine
\UIC{$(\mathrm{MonoInc}(a) \to m \in \mathbb{N} \to m \in \mathbb{N}) \to \varphi_{3}$}
\BIC{$\varphi_{3}$}
\AXC{$\stackrel{\Pi_{35}}{\vdots}$}\noLine
\UIC{$\varphi_{3} \to \varphi_{5}$}
\BIC{$\varphi_{5}$}
\AXC{$\varphi_{5} \to \varphi_{152}$\lbl{A1}}
\BIC{$\varphi_{152}$}
\end{ptree}
\noindent$\Pi_{37}$:
\begin{ptree}
\AXC{$x = x$\lbl{A5}}
\UIC{$\forall x\,(x = x)$}
\AXC{$\forall x\,(x = x) \to a_{m} = a_{m}$\lbl{A2}}
\BIC{$a_{m} = a_{m}$}
\AXC{$\varphi_{53}$\lbl{A1}}
\BIC{$\mathrm{MonoInc}(a) \to m \in \mathbb{N} \to m \leq 1_{\omega} \to a_{m} \leq a_{m}$}
\AXC{$\stackrel{\Pi_{36}}{\vdots}$}\noLine
\UIC{$\varphi_{152}$}
\BIC{$\varphi_{144}$}
\end{ptree}
\noindent$\Pi_{38}$:
\begin{ptree}
\AXC{$z = w \to a_{m} < a_{z} \to a_{m} < a_{w}$\lbl{A6}}
\UIC{$\forall w\,(z = w \to a_{m} < a_{z} \to a_{m} < a_{w})$}
\UIC{$\forall z\,\forall w\,(z = w \to a_{m} < a_{z} \to a_{m} < a_{w})$}
\AXC{$\varphi_{95}$\lbl{A2}}
\BIC{$\forall w\,(m = w \to a_{m} < a_{m} \to a_{m} < a_{w})$}
\AXC{$\varphi_{100}$\lbl{A2}}
\BIC{$m = 1_{\omega} \to a_{m} < a_{m} \to a_{m} < a_{1_{\omega}}$}
\AXC{$\stackrel{\Pi_{37}}{\vdots}$}\noLine
\UIC{$\varphi_{144}$}
\BIC{$(m = 1_{\omega} \to a_{m} = a_{m} \to a_{m} = a_{1_{\omega}}) \to \varphi_{21}$}
\end{ptree}
\noindent$\Pi_{39}$:
\begin{ptree}
\AXC{$z = w \to a_{m} = a_{z} \to a_{m} = a_{w}$\lbl{A6}}
\UIC{$\forall w\,(z = w \to a_{m} = a_{z} \to a_{m} = a_{w})$}
\UIC{$\forall z\,\forall w\,(z = w \to a_{m} = a_{z} \to a_{m} = a_{w})$}
\AXC{$\varphi_{82}$\lbl{A2}}
\BIC{$\forall w\,(m = w \to a_{m} = a_{m} \to a_{m} = a_{w})$}
\AXC{$\varphi_{84}$\lbl{A2}}
\BIC{$m = 1_{\omega} \to a_{m} = a_{m} \to a_{m} = a_{1_{\omega}}$}
\AXC{$\stackrel{\Pi_{38}}{\vdots}$}\noLine
\UIC{$(m = 1_{\omega} \to a_{m} = a_{m} \to a_{m} = a_{1_{\omega}}) \to \varphi_{21}$}
\BIC{$\varphi_{21}$}
\end{ptree}
\noindent$\Pi_{40}$:
\begin{ptree}
\AXC{$\stackrel{\Pi_{39}}{\vdots}$}\noLine
\UIC{$\varphi_{21}$}
\AXC{$\stackrel{\Pi_{32}}{\vdots}$}\noLine
\UIC{$\mathrm{MonoInc}(a) \to m \in \mathbb{N} \to \varphi_{107}$}
\AXC{$(\mathrm{MonoInc}(a) \to m \in \mathbb{N} \to \varphi_{107}) \to \varphi_{21} \to \varphi_{198}$\lbl{A1}}
\BIC{$\varphi_{21} \to \varphi_{198}$}
\BIC{$\varphi_{198}$}
\end{ptree}
\noindent$\Pi_{41}$:
\begin{ptree}
\AXC{$z = w \to m < z \to m < w$\lbl{A6}}
\UIC{$\forall w\,(z = w \to m < z \to m < w)$}
\UIC{$\forall z\,\forall w\,(z = w \to m < z \to m < w)$}
\AXC{$\varphi_{33}$\lbl{A2}}
\BIC{$\forall w\,(1_{\omega} = w \to m < 1_{\omega} \to m < w)$}
\AXC{$\varphi_{45}$\lbl{A2}}
\BIC{$1_{\omega} = S(n) \to m < 1_{\omega} \to m < S(n)$}
\AXC{$\stackrel{\Pi_{40}}{\vdots}$}\noLine
\UIC{$\varphi_{198}$}
\BIC{$(1_{\omega} = S(n) \to m = 1_{\omega} \to m = S(n)) \to \varphi_{194}$}
\end{ptree}
\noindent$\Pi_{42}$:
\begin{ptree}
\AXC{$z = w \to m = z \to m = w$\lbl{A6}}
\UIC{$\forall w\,(z = w \to m = z \to m = w)$}
\UIC{$\forall z\,\forall w\,(z = w \to m = z \to m = w)$}
\AXC{$\varphi_{12}$\lbl{A2}}
\BIC{$\forall w\,(1_{\omega} = w \to m = 1_{\omega} \to m = w)$}
\AXC{$\varphi_{20}$\lbl{A2}}
\BIC{$1_{\omega} = S(n) \to m = 1_{\omega} \to m = S(n)$}
\AXC{$\stackrel{\Pi_{41}}{\vdots}$}\noLine
\UIC{$(1_{\omega} = S(n) \to m = 1_{\omega} \to m = S(n)) \to \varphi_{194}$}
\BIC{$\varphi_{194}$}
\end{ptree}
\noindent$\Pi_{43}$:
\begin{ptree}
\AXC{$z = w \to a_{m} < a_{z} \to a_{m} < a_{w}$\lbl{A6}}
\UIC{$\forall w\,(z = w \to a_{m} < a_{z} \to a_{m} < a_{w})$}
\UIC{$\forall z\,\forall w\,(z = w \to a_{m} < a_{z} \to a_{m} < a_{w})$}
\AXC{$\varphi_{105}$\lbl{A2}}
\BIC{$\forall w\,(1_{\omega} = w \to a_{m} < a_{1_{\omega}} \to a_{m} < a_{w})$}
\AXC{$\varphi_{116}$\lbl{A2}}
\BIC{$1_{\omega} = S(n) \to a_{m} < a_{1_{\omega}} \to a_{m} < a_{S(n)}$}
\AXC{$\stackrel{\Pi_{42}}{\vdots}$}\noLine
\UIC{$\varphi_{194}$}
\BIC{$\varphi_{190}$}
\end{ptree}
\noindent$\Pi_{44}$:
\begin{ptree}
\AXC{$z = w \to a_{m} = a_{z} \to a_{m} = a_{w}$\lbl{A6}}
\UIC{$\forall w\,(z = w \to a_{m} = a_{z} \to a_{m} = a_{w})$}
\UIC{$\forall z\,\forall w\,(z = w \to a_{m} = a_{z} \to a_{m} = a_{w})$}
\AXC{$\varphi_{92}$\lbl{A2}}
\BIC{$\forall w\,(1_{\omega} = w \to a_{m} = a_{1_{\omega}} \to a_{m} = a_{w})$}
\AXC{$\varphi_{99}$\lbl{A2}}
\BIC{$1_{\omega} = S(n) \to a_{m} = a_{1_{\omega}} \to a_{m} = a_{S(n)}$}
\AXC{$\stackrel{\Pi_{43}}{\vdots}$}\noLine
\UIC{$\varphi_{190}$}
\BIC{$\varphi_{187}$}
\end{ptree}
\noindent$\Pi_{45}$:
\begin{ptree}
\AXC{$z = w \to m < z \to m < w$\lbl{A6}}
\UIC{$\forall w\,(z = w \to m < z \to m < w)$}
\UIC{$\forall z\,\forall w\,(z = w \to m < z \to m < w)$}
\AXC{$\varphi_{32}$\lbl{A2}}
\BIC{$\forall w\,(S(n) = w \to m < S(n) \to m < w)$}
\AXC{$\varphi_{44}$\lbl{A2}}
\BIC{$S(n) = 1_{\omega} \to m < S(n) \to m < 1_{\omega}$}
\AXC{$\stackrel{\Pi_{44}}{\vdots}$}\noLine
\UIC{$\varphi_{187}$}
\BIC{$(S(n) = 1_{\omega} \to m = S(n) \to m = 1_{\omega}) \to \varphi_{182}$}
\end{ptree}
\noindent$\Pi_{46}$:
\begin{ptree}
\AXC{$z = w \to m = z \to m = w$\lbl{A6}}
\UIC{$\forall w\,(z = w \to m = z \to m = w)$}
\UIC{$\forall z\,\forall w\,(z = w \to m = z \to m = w)$}
\AXC{$\varphi_{11}$\lbl{A2}}
\BIC{$\forall w\,(S(n) = w \to m = S(n) \to m = w)$}
\AXC{$\varphi_{19}$\lbl{A2}}
\BIC{$S(n) = 1_{\omega} \to m = S(n) \to m = 1_{\omega}$}
\AXC{$\stackrel{\Pi_{45}}{\vdots}$}\noLine
\UIC{$(S(n) = 1_{\omega} \to m = S(n) \to m = 1_{\omega}) \to \varphi_{182}$}
\BIC{$\varphi_{182}$}
\end{ptree}
\noindent$\Pi_{47}$:
\begin{ptree}
\AXC{$z = w \to a_{m} < a_{z} \to a_{m} < a_{w}$\lbl{A6}}
\UIC{$\forall w\,(z = w \to a_{m} < a_{z} \to a_{m} < a_{w})$}
\UIC{$\forall z\,\forall w\,(z = w \to a_{m} < a_{z} \to a_{m} < a_{w})$}
\AXC{$\varphi_{104}$\lbl{A2}}
\BIC{$\forall w\,(S(n) = w \to a_{m} < a_{S(n)} \to a_{m} < a_{w})$}
\AXC{$\varphi_{115}$\lbl{A2}}
\BIC{$S(n) = 1_{\omega} \to a_{m} < a_{S(n)} \to a_{m} < a_{1_{\omega}}$}
\AXC{$\stackrel{\Pi_{46}}{\vdots}$}\noLine
\UIC{$\varphi_{182}$}
\BIC{$\varphi_{181}$}
\end{ptree}
\noindent$\Pi_{48}$:
\begin{ptree}
\AXC{$z = w \to a_{m} = a_{z} \to a_{m} = a_{w}$\lbl{A6}}
\UIC{$\forall w\,(z = w \to a_{m} = a_{z} \to a_{m} = a_{w})$}
\UIC{$\forall z\,\forall w\,(z = w \to a_{m} = a_{z} \to a_{m} = a_{w})$}
\AXC{$\varphi_{91}$\lbl{A2}}
\BIC{$\forall w\,(S(n) = w \to a_{m} = a_{S(n)} \to a_{m} = a_{w})$}
\AXC{$\varphi_{98}$\lbl{A2}}
\BIC{$S(n) = 1_{\omega} \to a_{m} = a_{S(n)} \to a_{m} = a_{1_{\omega}}$}
\AXC{$\stackrel{\Pi_{47}}{\vdots}$}\noLine
\UIC{$\varphi_{181}$}
\BIC{$\varphi_{171}$}
\end{ptree}
\noindent$\Pi_{49}$:
\begin{ptree}
\AXC{$\stackrel{\Pi_{24}}{\vdots}$}\noLine
\UIC{$\forall x\,\forall y\,(x = y \to y = x)$}
\AXC{$\varphi_{1}$\lbl{A2}}
\BIC{$\forall y\,(1_{\omega} = y \to y = 1_{\omega})$}
\AXC{$\varphi_{6}$\lbl{A2}}
\BIC{$1_{\omega} = S(n) \to S(n) = 1_{\omega}$}
\AXC{$\stackrel{\Pi_{33}}{\vdots}$}\noLine
\UIC{$\varphi_{146}$}
\BIC{$\mathrm{MonoInc}(a) \to m \in \mathbb{N} \to \varphi_{106}$}
\AXC{$\stackrel{\Pi_{48}}{\vdots}$}\noLine
\UIC{$\varphi_{171}$}
\BIC{$\mathrm{MonoInc}(a) \to m \in \mathbb{N} \to \varphi_{130}$}
\AXC{$(\mathrm{MonoInc}(a) \to m \in \mathbb{N} \to \varphi_{130}) \to \varphi_{173}$\lbl{A1}}
\BIC{$\varphi_{173}$}
\end{ptree}
\noindent$\Pi_{50}$:
\begin{ptree}
\AXC{$\forall x\,(\mathrm{Ind}(x) \leftrightarrow \emptyset \in x \land \forall y \in x\,(S(y) \in x))$\lbl{DefInd}}
\AXC{$\varphi_{114}$\lbl{A2}}
\BIC{$\mathrm{Ind}(\omega) \leftrightarrow \emptyset \in \omega \land \forall y \in \omega\,(S(y) \in \omega)$}
\AXC{$\mathrm{Ind}(\omega)$\lbl{補題(ωは帰納的)}}
\AXC{$\mathrm{Ind}(\omega) \to \varphi_{151}$\lbl{A1}}
\BIC{$\varphi_{151}$}
\BIC{$\mathrm{MonoInc}(a) \to m \in \mathbb{N} \to \varphi_{112}$}
\end{ptree}
\noindent$\Pi_{51}$:
\begin{ptree}
\AXC{$\stackrel{\Pi_{50}}{\vdots}$}\noLine
\UIC{$\mathrm{MonoInc}(a) \to m \in \mathbb{N} \to \varphi_{112}$}
\AXC{$\varphi_{162}$\lbl{A1}}
\BIC{$\mathrm{MonoInc}(a) \to m \in \mathbb{N} \to \varphi_{101}$}
\AXC{$(\mathrm{MonoInc}(a) \to m \in \mathbb{N} \to \varphi_{101}) \to \varphi_{150}$\lbl{A1}}
\BIC{$\varphi_{150}$}
\end{ptree}
\noindent$\Pi_{52}$:
\begin{ptree}
\AXC{$\forall y \in \omega\,(S(y) \in \omega) \to n \in \omega \to S(n) \in \omega$\lbl{A2}}
\AXC{$\stackrel{\Pi_{51}}{\vdots}$}\noLine
\UIC{$\varphi_{150}$}
\BIC{$\mathrm{MonoInc}(a) \to m \in \mathbb{N} \to \varphi_{93}$}
\AXC{$(\mathrm{MonoInc}(a) \to m \in \mathbb{N} \to \varphi_{93}) \to \varphi_{155}$\lbl{A1}}
\BIC{$\varphi_{155}$}
\end{ptree}
\noindent$\Pi_{53}$:
\begin{ptree}
\AXC{$\stackrel{\Pi_{9}}{\vdots}$}\noLine
\UIC{$\mathrm{MonoInc}(a) \to m \in \mathbb{N} \to \varphi_{76}$}
\AXC{$\stackrel{\Pi_{52}}{\vdots}$}\noLine
\UIC{$\varphi_{155}$}
\BIC{$\mathrm{MonoInc}(a) \to m \in \mathbb{N} \to \varphi_{79}$}
\AXC{$\varphi_{135}$\lbl{A2}}
\AXC{$\mathrm{MonoInc}(a) \to m \in \mathbb{N} \to \varphi_{127}$\lbl{A1}}
\AXC{$\varphi_{186}$\lbl{A1}}
\BIC{$\varphi_{135} \to \mathrm{MonoInc}(a) \to m \in \mathbb{N} \to \varphi_{135}$}
\BIC{$\mathrm{MonoInc}(a) \to m \in \mathbb{N} \to \varphi_{135}$}
\AXC{$\varphi_{192}$\lbl{A1}}
\BIC{$(\mathrm{MonoInc}(a) \to m \in \mathbb{N} \to \varphi_{79}) \to \varphi_{159}$}
\BIC{$\varphi_{159}$}
\end{ptree}
\noindent$\Pi_{54}$:
\begin{ptree}
\AXC{$z = w \to a_{m} = a_{z} \to a_{m} = a_{w}$\lbl{A6}}
\UIC{$\forall w\,(z = w \to a_{m} = a_{z} \to a_{m} = a_{w})$}
\UIC{$\forall z\,\forall w\,(z = w \to a_{m} = a_{z} \to a_{m} = a_{w})$}
\AXC{$\varphi_{82}$\lbl{A2}}
\BIC{$\forall w\,(m = w \to a_{m} = a_{m} \to a_{m} = a_{w})$}
\AXC{$\varphi_{83}$\lbl{A2}}
\BIC{$m = S(n) \to a_{m} = a_{m} \to a_{m} = a_{S(n)}$}
\AXC{$\stackrel{\Pi_{6}}{\vdots}$}\noLine
\UIC{$\varphi_{153}$}
\BIC{$\mathrm{MonoInc}(a) \to m \in \mathbb{N} \to \varphi_{138}$}
\AXC{$\stackrel{\Pi_{23}}{\vdots}$}\noLine
\UIC{$\varphi_{176}$}
\BIC{$\mathrm{MonoInc}(a) \to m \in \mathbb{N} \to \varphi_{132}$}
\AXC{$\stackrel{\Pi_{49}}{\vdots}$}\noLine
\UIC{$\varphi_{173}$}
\BIC{$\mathrm{MonoInc}(a) \to m \in \mathbb{N} \to \varphi_{120}$}
\AXC{$\stackrel{\Pi_{53}}{\vdots}$}\noLine
\UIC{$\varphi_{159}$}
\BIC{$\mathrm{MonoInc}(a) \to m \in \mathbb{N} \to \varphi_{77}$}
\AXC{$\varphi_{154}$\lbl{A1}}
\BIC{$\varphi_{74} \to n \in B_{18} \to S(n) \in B_{18}$}
\UIC{$\forall n\,(\varphi_{74} \to n \in B_{18} \to S(n) \in B_{18})$}
\end{ptree}
\noindent$\Pi_{55}$:
\begin{ptree}
\AXC{$\forall n\,(n \in \mathbb{N} \leftrightarrow n \in \omega \land n \neq \emptyset)$\lbl{Def$\mathbb{N}$}}
\AXC{$\varphi_{87}$\lbl{A2}}
\BIC{$\emptyset \in \mathbb{N} \leftrightarrow \emptyset \in \omega \land \emptyset \neq \emptyset$}
\AXC{$(\emptyset \in \mathbb{N} \leftrightarrow \emptyset \in \omega \land \emptyset \neq \emptyset) \to \emptyset = \emptyset \to \varphi_{68}$\lbl{A1}}
\BIC{$\emptyset = \emptyset \to \varphi_{68}$}
\end{ptree}
\noindent$\Pi_{56}$:
\begin{ptree}
\AXC{$\varphi_{143}$\lbl{A2}}
\AXC{$\mathrm{MonoInc}(a) \to m \in \mathbb{N} \to \varphi_{127}$\lbl{A1}}
\AXC{$\varphi_{188}$\lbl{A1}}
\BIC{$\varphi_{143} \to \mathrm{MonoInc}(a) \to m \in \mathbb{N} \to \varphi_{143}$}
\BIC{$\mathrm{MonoInc}(a) \to m \in \mathbb{N} \to \varphi_{143}$}
\AXC{$(\mathrm{MonoInc}(a) \to m \in \mathbb{N} \to \varphi_{143}) \to \varphi_{177}$\lbl{A1}}
\BIC{$\varphi_{177}$}
\end{ptree}
\noindent$\Pi_{57}$:
\begin{ptree}
\AXC{$x = x$\lbl{A5}}
\UIC{$\forall x\,(x = x)$}
\AXC{$\forall x\,(x = x) \to \emptyset = \emptyset$\lbl{A2}}
\BIC{$\emptyset = \emptyset$}
\AXC{$\stackrel{\Pi_{55}}{\vdots}$}\noLine
\UIC{$\emptyset = \emptyset \to \varphi_{68}$}
\BIC{$\varphi_{68}$}
\AXC{$\stackrel{\Pi_{50}}{\vdots}$}\noLine
\UIC{$\mathrm{MonoInc}(a) \to m \in \mathbb{N} \to \varphi_{112}$}
\AXC{$\stackrel{\Pi_{56}}{\vdots}$}\noLine
\UIC{$\varphi_{177}$}
\BIC{$\varphi_{68} \to \mathrm{MonoInc}(a) \to m \in \mathbb{N} \to \varphi_{70}$}
\BIC{$\mathrm{MonoInc}(a) \to m \in \mathbb{N} \to \varphi_{70}$}
\end{ptree}
\noindent$\Pi_{58}$:
\begin{ptree}
\AXC{$\forall x\,(\mathrm{Ind}(x) \leftrightarrow \emptyset \in x \land \forall y \in x\,(S(y) \in x))$\lbl{DefInd}}
\AXC{$\forall x\,(\mathrm{Ind}(x) \leftrightarrow \emptyset \in x \land \forall y \in x\,(S(y) \in x)) \to \varphi_{28}$\lbl{A2}}
\BIC{$\varphi_{28}$}
\end{ptree}
\noindent$\Pi_{59}$:
\begin{ptree}
\AXC{$\stackrel{\Pi_{57}}{\vdots}$}\noLine
\UIC{$\mathrm{MonoInc}(a) \to m \in \mathbb{N} \to \varphi_{70}$}
\AXC{$\stackrel{\Pi_{58}}{\vdots}$}\noLine
\UIC{$\varphi_{28}$}
\AXC{$\varphi_{28} \to (\mathrm{MonoInc}(a) \to m \in \mathbb{N} \to \varphi_{70}) \to \varphi_{178}$\lbl{A1}}
\BIC{$(\mathrm{MonoInc}(a) \to m \in \mathbb{N} \to \varphi_{70}) \to \varphi_{178}$}
\BIC{$\varphi_{178}$}
\end{ptree}
\noindent$\Pi_{60}$:
\begin{ptree}
\AXC{$\varphi_{74} \to \forall n\,(\varphi_{74})$\lbl{A4}}
\AXC{$\varphi_{156}$\lbl{A3}}
\AXC{$\stackrel{\Pi_{54}}{\vdots}$}\noLine
\UIC{$\forall n\,(\varphi_{74} \to n \in B_{18} \to S(n) \in B_{18})$}
\AXC{$\varphi_{197}$\lbl{A1}}
\BIC{$\varphi_{156} \to (\varphi_{74} \to \forall n\,(\varphi_{74})) \to \mathrm{MonoInc}(a) \to m \in \mathbb{N} \to \varphi_{80}$}
\BIC{$(\varphi_{74} \to \forall n\,(\varphi_{74})) \to \mathrm{MonoInc}(a) \to m \in \mathbb{N} \to \varphi_{80}$}
\BIC{$\mathrm{MonoInc}(a) \to m \in \mathbb{N} \to \varphi_{80}$}
\AXC{$\stackrel{\Pi_{59}}{\vdots}$}\noLine
\UIC{$\varphi_{178}$}
\BIC{$\varphi_{26} \to \varphi_{25} \to \mathrm{MonoInc}(a) \to m \in \mathbb{N} \to \varphi_{61}$}
\end{ptree}
\noindent$\Pi_{61}$:
\begin{ptree}
\AXC{$\mathrm{MonoInc}(a) \to m \in \mathbb{N} \to \varphi_{78}$\lbl{A1}}
\AXC{$(\mathrm{MonoInc}(a) \to m \in \mathbb{N} \to \varphi_{78}) \to \varphi_{142}$\lbl{A1}}
\BIC{$\varphi_{142}$}
\end{ptree}
\noindent$\Pi_{62}$:
\begin{ptree}
\AXC{$\forall n\,(n \in \omega \leftrightarrow \forall x\,(\mathrm{Ind}(x) \to n \in x))$\lbl{Def$\omega$}}
\AXC{$\varphi_{54}$\lbl{A2}}
\BIC{$n \in \omega \leftrightarrow \forall x\,(\mathrm{Ind}(x) \to n \in x)$}
\AXC{$\stackrel{\Pi_{61}}{\vdots}$}\noLine
\UIC{$\varphi_{142}$}
\BIC{$\mathrm{MonoInc}(a) \to m \in \mathbb{N} \to \varphi_{86}$}
\AXC{$(\mathrm{MonoInc}(a) \to m \in \mathbb{N} \to \varphi_{86}) \to \varphi_{147}$\lbl{A1}}
\BIC{$\varphi_{147}$}
\end{ptree}
\noindent$\Pi_{63}$:
\begin{ptree}
\AXC{$\forall x\,(\mathrm{Ind}(x) \to n \in x) \to \mathrm{Ind}(B_{18}) \to n \in B_{18}$\lbl{A2}}
\AXC{$\stackrel{\Pi_{62}}{\vdots}$}\noLine
\UIC{$\varphi_{147}$}
\BIC{$\mathrm{MonoInc}(a) \to m \in \mathbb{N} \to \varphi_{88}$}
\AXC{$\varphi_{189}$\lbl{A1}}
\BIC{$(\mathrm{MonoInc}(a) \to m \in \mathbb{N} \to \varphi_{61}) \to \varphi_{167}$}
\end{ptree}
\noindent$\Pi_{64}$:
\begin{ptree}
\AXC{$\stackrel{\Pi_{2}}{\vdots}$}\noLine
\UIC{$\varphi_{25}$}
\AXC{$\stackrel{\Pi_{3}}{\vdots}$}\noLine
\UIC{$\varphi_{26}$}
\AXC{$\stackrel{\Pi_{60}}{\vdots}$}\noLine
\UIC{$\varphi_{26} \to \varphi_{25} \to \mathrm{MonoInc}(a) \to m \in \mathbb{N} \to \varphi_{61}$}
\BIC{$\varphi_{25} \to \mathrm{MonoInc}(a) \to m \in \mathbb{N} \to \varphi_{61}$}
\BIC{$\mathrm{MonoInc}(a) \to m \in \mathbb{N} \to \varphi_{61}$}
\AXC{$\stackrel{\Pi_{63}}{\vdots}$}\noLine
\UIC{$(\mathrm{MonoInc}(a) \to m \in \mathbb{N} \to \varphi_{61}) \to \varphi_{167}$}
\BIC{$\varphi_{167}$}
\end{ptree}
\noindent$\Pi_{65}$:
\begin{ptree}
\AXC{$\varphi_{165}$\lbl{A3}}
\AXC{$\varphi_{122}$\lbl{A2}}
\AXC{$\mathrm{MonoInc}(a) \to m \in \mathbb{N} \to \varphi_{127}$\lbl{A1}}
\AXC{$\varphi_{183}$\lbl{A1}}
\BIC{$\varphi_{122} \to \mathrm{MonoInc}(a) \to m \in \mathbb{N} \to \varphi_{122}$}
\BIC{$\mathrm{MonoInc}(a) \to m \in \mathbb{N} \to \varphi_{122}$}
\AXC{$\stackrel{\Pi_{64}}{\vdots}$}\noLine
\UIC{$\varphi_{167}$}
\BIC{$\mathrm{MonoInc}(a) \to m \in \mathbb{N} \to \varphi_{108}$}
\AXC{$\varphi_{163}$\lbl{A1}}
\BIC{$\varphi_{74} \to n \in \omega \to n \in \mathbb{N} \to m \leq n \to a_{m} \leq a_{n}$}
\UIC{$\forall n\,(\varphi_{74} \to n \in \omega \to n \in \mathbb{N} \to m \leq n \to a_{m} \leq a_{n})$}
\AXC{$\varphi_{199}$\lbl{A1}}
\BIC{$\varphi_{165} \to (\varphi_{74} \to \forall n\,(\varphi_{74})) \to \mathrm{MonoInc}(a) \to m \in \mathbb{N} \to \varphi_{113}$}
\BIC{$(\varphi_{74} \to \forall n\,(\varphi_{74})) \to \mathrm{MonoInc}(a) \to m \in \mathbb{N} \to \varphi_{113}$}
\end{ptree}
\noindent$\Pi_{66}$:
\begin{ptree}
\AXC{$\varphi_{74} \to \forall n\,(\varphi_{74})$\lbl{A4}}
\AXC{$\stackrel{\Pi_{65}}{\vdots}$}\noLine
\UIC{$(\varphi_{74} \to \forall n\,(\varphi_{74})) \to \mathrm{MonoInc}(a) \to m \in \mathbb{N} \to \varphi_{113}$}
\BIC{$\mathrm{MonoInc}(a) \to m \in \mathbb{N} \to \varphi_{113}$}
\AXC{$(\mathrm{MonoInc}(a) \to m \in \mathbb{N} \to \varphi_{113}) \to \varphi_{66} \to \varphi_{42}$\lbl{A1}}
\BIC{$\varphi_{66} \to \varphi_{42}$}
\UIC{$\forall B_{18}\,(\varphi_{66} \to \varphi_{42})$}
\end{ptree}
\noindent$\Pi_{67}$:
\begin{ptree}
\AXC{$\varphi_{75}$\lbl{Sep}}
\AXC{$\varphi_{75} \to \varphi_{69}$\lbl{A2}}
\BIC{$\varphi_{69}$}
\AXC{$\varphi_{69} \to \varphi_{55}$\lbl{A2}}
\BIC{$\varphi_{55}$}
\AXC{$\varphi_{55} \to \varphi_{48}$\lbl{A2}}
\BIC{$\varphi_{48}$}
\AXC{$\varphi_{48} \to \mathrm{MonoInc}(a) \to m \in \mathbb{N} \to \varphi_{48}$\lbl{A1}}
\BIC{$\mathrm{MonoInc}(a) \to m \in \mathbb{N} \to \varphi_{48}$}
\AXC{$\varphi_{200}$\lbl{A1}}
\BIC{$\forall B_{18}\,(\varphi_{66} \to \varphi_{42}) \to \varphi_{191}$}
\end{ptree}
\noindent$\Pi_{68}$:
\begin{ptree}
\AXC{$\varphi_{66} \to \forall B_{18}\,(\varphi_{66})$\lbl{A4}}
\AXC{$\forall B_{18}\,(\varphi_{66} \to \varphi_{42}) \to \forall B_{18}\,(\varphi_{66}) \to \forall B_{18}\,(\varphi_{42})$\lbl{A3}}
\AXC{$\stackrel{\Pi_{66}}{\vdots}$}\noLine
\UIC{$\forall B_{18}\,(\varphi_{66} \to \varphi_{42})$}
\AXC{$\stackrel{\Pi_{67}}{\vdots}$}\noLine
\UIC{$\forall B_{18}\,(\varphi_{66} \to \varphi_{42}) \to \varphi_{191}$}
\BIC{$\varphi_{191}$}
\BIC{$(\varphi_{66} \to \forall B_{18}\,(\varphi_{66})) \to \varphi_{65}$}
\BIC{$\varphi_{65}$}
\end{ptree}
\noindent$\Pi_{69}$:
\begin{ptree}
\AXC{$\varphi_{89}$\lbl{A2}}
\AXC{$\stackrel{\Pi_{68}}{\vdots}$}\noLine
\UIC{$\varphi_{65}$}
\AXC{$\varphi_{65} \to \varphi_{89} \to \varphi_{60}$\lbl{A1}}
\BIC{$\varphi_{89} \to \varphi_{60}$}
\BIC{$\varphi_{60}$}
\AXC{$\varphi_{60} \to (\mathrm{MonoInc}(a) \to m \in \mathbb{N} \to n \in \mathbb{N} \to n \in \omega) \to \varphi_{94}$\lbl{A1}}
\BIC{$(\mathrm{MonoInc}(a) \to m \in \mathbb{N} \to n \in \mathbb{N} \to n \in \omega) \to \varphi_{94}$}
\end{ptree}
\noindent$\Pi_{70}$:
\begin{ptree}
\AXC{$\forall n\,(n \in \mathbb{N} \leftrightarrow n \in \omega \land n \neq \emptyset)$\lbl{Def$\mathbb{N}$}}
\AXC{$\varphi_{67}$\lbl{A2}}
\BIC{$n \in \mathbb{N} \leftrightarrow n \in \omega \land n \neq \emptyset$}
\AXC{$\stackrel{\Pi_{1}}{\vdots}$}\noLine
\UIC{$\varphi_{72}$}
\BIC{$\mathrm{MonoInc}(a) \to m \in \mathbb{N} \to n \in \mathbb{N} \to n \in \omega$}
\AXC{$\stackrel{\Pi_{69}}{\vdots}$}\noLine
\UIC{$(\mathrm{MonoInc}(a) \to m \in \mathbb{N} \to n \in \mathbb{N} \to n \in \omega) \to \varphi_{94}$}
\BIC{$\varphi_{94}$}
\end{ptree}
\noindent$\Pi_{71}$:
\begin{ptree}
\AXC{$\varphi_{38} \to \varphi_{46}$\lbl{A3}}
\AXC{$\mathrm{MonoInc}(a) \to m \in \mathbb{N} \to n \in \mathbb{N} \to n \in \mathbb{N}$\lbl{A1}}
\AXC{$\stackrel{\Pi_{70}}{\vdots}$}\noLine
\UIC{$\varphi_{94}$}
\BIC{$\mathrm{MonoInc}(a) \to m \in \mathbb{N} \to n \in \mathbb{N} \to m \leq n \to a_{m} \leq a_{n}$}
\AXC{$\varphi_{117}$\lbl{A1}}
\BIC{$\mathrm{MonoInc}(a) \land m \in \mathbb{N} \to n \in \mathbb{N} \to m \leq n \to a_{m} \leq a_{n}$}
\UIC{$\varphi_{38}$}
\AXC{$\varphi_{38} \to (\varphi_{38} \to \varphi_{46}) \to \varphi_{13} \to \varphi_{39}$\lbl{A1}}
\BIC{$(\varphi_{38} \to \varphi_{46}) \to \varphi_{13} \to \varphi_{39}$}
\BIC{$\varphi_{13} \to \varphi_{39}$}
\end{ptree}
\noindent 主木:
\begin{ptree}
\AXC{$\varphi_{62}$\lbl{A2}}
\AXC{$\varphi_{13}$\lbl{A4}}
\AXC{$\stackrel{\Pi_{71}}{\vdots}$}\noLine
\UIC{$\varphi_{13} \to \varphi_{39}$}
\BIC{$\varphi_{39}$}
\AXC{$\varphi_{39} \to \varphi_{133}$\lbl{A1}}
\BIC{$\varphi_{133}$}
\BIC{$\mathrm{MonoInc}(a) \to m \in \mathbb{N} \to n \in \mathbb{N} \to m \leq n \to a_{m} \leq a_{n}$}
\UIC{$\varphi_{40}$}
\UIC{$\forall m\,(\varphi_{40})$}
\UIC{$\forall a\,\forall m\,(\varphi_{40})$}
\end{ptree}
\begin{abbrevlist}
\item $\varphi_{1} :\equiv \forall x\,\forall y\,(x = y \to y = x) \to \forall y\,(1_{\omega} = y \to y = 1_{\omega})$
\item $\varphi_{2} :\equiv \forall x\,\forall y\,(x = y \to y = x) \to \forall y\,(S(n) = y \to y = S(n))$
\item $\varphi_{3} :\equiv \mathrm{MonoInc}(a) \to m \in \mathbb{N} \to m \leq 1_{\omega} \to 1_{\omega} \leq m$
\item $\varphi_{4} :\equiv \mathrm{MonoInc}(a) \to m \in \mathbb{N} \to m \leq 1_{\omega} \to m \leq 1_{\omega}$
\item $\varphi_{5} :\equiv \mathrm{MonoInc}(a) \to m \in \mathbb{N} \to m \leq 1_{\omega} \to m = 1_{\omega}$
\item $\varphi_{6} :\equiv \forall y\,(1_{\omega} = y \to y = 1_{\omega}) \to 1_{\omega} = S(n) \to S(n) = 1_{\omega}$
\item $\varphi_{7} :\equiv \forall y\,(S(n) = y \to y = S(n)) \to S(n) = 1_{\omega} \to 1_{\omega} = S(n)$
\item $\varphi_{8} :\equiv \forall x\,\forall y\,(x \leq y \to y \leq x \to x = y) \to \forall y\,(m \leq y \to y \leq m \to m = y)$
\item $\varphi_{9} :\equiv \forall y\,(m \leq y \to y \leq m \to m = y) \to m \leq 1_{\omega} \to 1_{\omega} \leq m \to m = 1_{\omega}$
\item $\varphi_{10} :\equiv S(n) = S(\emptyset) \to S(\emptyset) = 1_{\omega} \to S(n) = 1_{\omega}$
\item $\varphi_{11} :\equiv \forall z\,\forall w\,(z = w \to m = z \to m = w) \to \forall w\,(S(n) = w \to m = S(n) \to m = w)$
\item $\varphi_{12} :\equiv \forall z\,\forall w\,(z = w \to m = z \to m = w) \to \forall w\,(1_{\omega} = w \to m = 1_{\omega} \to m = w)$
\item $\varphi_{13} :\equiv \mathrm{MonoInc}(a) \land m \in \mathbb{N} \to \forall n\,(\mathrm{MonoInc}(a) \land m \in \mathbb{N})$
\item $\varphi_{14} :\equiv S(n) \in \mathbb{N} \to n \neq \emptyset \to m \leq S(n) \to m = S(n) \to m = S(n)$
\item $\varphi_{15} :\equiv \forall n \in B_{18}\,(S(n) \in B_{18}) \to y \in B_{18} \to S(y) \in B_{18}$
\item $\varphi_{16} :\equiv \forall y \in B_{18}\,(S(y) \in B_{18}) \to n \in B_{18} \to S(n) \in B_{18}$
\item $\varphi_{17} :\equiv \forall x\,\forall y\,\forall z\,(x = y \to y = z \to x = z) \to \forall y\,\forall z\,(S(n) = y \to y = z \to S(n) = z)$
\item $\varphi_{18} :\equiv (1_{\omega} = S(\emptyset) \to S(\emptyset) = 1_{\omega}) \to S(\emptyset) = 1_{\omega}$
\item $\varphi_{19} :\equiv \forall w\,(S(n) = w \to m = S(n) \to m = w) \to S(n) = 1_{\omega} \to m = S(n) \to m = 1_{\omega}$
\item $\varphi_{20} :\equiv \forall w\,(1_{\omega} = w \to m = 1_{\omega} \to m = w) \to 1_{\omega} = S(n) \to m = 1_{\omega} \to m = S(n)$
\item $\varphi_{21} :\equiv \mathrm{MonoInc}(a) \to m \in \mathbb{N} \to m \leq 1_{\omega} \to a_{m} \leq a_{1_{\omega}}$
\item $\varphi_{22} :\equiv \forall m\,\forall n\,(m \in \omega \to n \in \omega \to m \leq S(n) \to m \leq n \lor m = S(n))$
\item $\varphi_{23} :\equiv \forall y\,(1_{\omega} = y \to y = 1_{\omega}) \to 1_{\omega} = S(\emptyset) \to S(\emptyset) = 1_{\omega}$
\item $\varphi_{24} :\equiv n \in B_{18} \to S(n) \in \mathbb{N} \to n \neq \emptyset \to m \leq S(n) \to m \leq S(n)$
\item $\varphi_{25} :\equiv \forall n \in B_{18}\,(S(n) \in B_{18}) \to \forall y \in B_{18}\,(S(y) \in B_{18})$
\item $\varphi_{26} :\equiv \forall y \in B_{18}\,(S(y) \in B_{18}) \to \forall n \in B_{18}\,(S(n) \in B_{18})$
\item $\varphi_{27} :\equiv \mathrm{MonoInc}(a) \leftrightarrow \mathrm{Seq}(a) \land \forall n \in \mathbb{N}\,(a_{n} \leq a_{S(n)})$
\item $\varphi_{28} :\equiv \mathrm{Ind}(B_{18}) \leftrightarrow \emptyset \in B_{18} \land \forall y \in B_{18}\,(S(y) \in B_{18})$
\item $\varphi_{29} :\equiv n \in B_{18} \to S(n) \in \mathbb{N} \to n \neq \emptyset \to m \leq S(n) \to m \leq n \to m \leq n$
\item $\varphi_{30} :\equiv \forall n \in B_{18}\,(S(n) \in B_{18}) \to \forall y\,\forall n \in B_{18}\,(S(n) \in B_{18})$
\item $\varphi_{31} :\equiv \forall y \in B_{18}\,(S(y) \in B_{18}) \to \forall n\,\forall y \in B_{18}\,(S(y) \in B_{18})$
\item $\varphi_{32} :\equiv \forall z\,\forall w\,(z = w \to m < z \to m < w) \to \forall w\,(S(n) = w \to m < S(n) \to m < w)$
\item $\varphi_{33} :\equiv \forall z\,\forall w\,(z = w \to m < z \to m < w) \to \forall w\,(1_{\omega} = w \to m < 1_{\omega} \to m < w)$
\item $\varphi_{34} :\equiv \forall y\,\forall n \in B_{18}\,(S(n) \in B_{18}) \to \forall y \in B_{18}\,(S(y) \in B_{18})$
\item $\varphi_{35} :\equiv \forall n\,\forall y \in B_{18}\,(S(y) \in B_{18}) \to \forall n \in B_{18}\,(S(n) \in B_{18})$
\item $\varphi_{36} :\equiv \forall x\,\forall y\,\forall z\,(x \leq y \to y \leq z \to x \leq z) \to \forall y\,\forall z\,(a_{m} \leq y \to y \leq z \to a_{m} \leq z)$
\item $\varphi_{37} :\equiv n \in B_{18} \to S(n) \in \mathbb{N} \to n \neq \emptyset \to m \leq S(n) \to m \leq n \lor m = S(n)$
\item $\varphi_{38} :\equiv \forall n\,(\mathrm{MonoInc}(a) \land m \in \mathbb{N} \to n \in \mathbb{N} \to m \leq n \to a_{m} \leq a_{n})$
\item $\varphi_{39} :\equiv \mathrm{MonoInc}(a) \to m \in \mathbb{N} \to \forall n \in \mathbb{N}\,(m \leq n \to a_{m} \leq a_{n})$
\item $\varphi_{40} :\equiv \forall n\,(\mathrm{MonoInc}(a) \to m \in \mathbb{N} \to n \in \mathbb{N} \to m \leq n \to a_{m} \leq a_{n})$
\item $\varphi_{41} :\equiv S(n) \in \mathbb{N} \to n \neq \emptyset \to m \leq S(n) \to m = S(n) \to a_{m} \leq a_{m}$
\item $\varphi_{42} :\equiv \neg \forall n\,(n \in B_{18} \leftrightarrow n \in \omega \land (n \in \mathbb{N} \to m \leq n \to a_{m} \leq a_{n}))$
\item $\varphi_{43} :\equiv \forall z\,\forall w\,(z = w \to S(n) = S(z) \to S(n) = S(w)) \to \forall w\,(n = w \to S(n) = S(n) \to S(n) = S(w))$
\item $\varphi_{44} :\equiv \forall w\,(S(n) = w \to m < S(n) \to m < w) \to S(n) = 1_{\omega} \to m < S(n) \to m < 1_{\omega}$
\item $\varphi_{45} :\equiv \forall w\,(1_{\omega} = w \to m < 1_{\omega} \to m < w) \to 1_{\omega} = S(n) \to m < 1_{\omega} \to m < S(n)$
\item $\varphi_{46} :\equiv \forall n\,(\mathrm{MonoInc}(a) \land m \in \mathbb{N}) \to \forall n \in \mathbb{N}\,(m \leq n \to a_{m} \leq a_{n})$
\item $\varphi_{47} :\equiv S(n) \in \mathbb{N} \to n \neq \emptyset \to m \leq S(n) \to m = S(n) \to a_{m} \leq a_{S(n)}$
\item $\varphi_{48} :\equiv \exists B_{18}\,\forall n\,(n \in B_{18} \leftrightarrow n \in \omega \land (n \in \mathbb{N} \to m \leq n \to a_{m} \leq a_{n}))$
\item $\varphi_{49} :\equiv n \in B_{18} \to S(n) \in \mathbb{N} \to n = \emptyset \to m \leq S(n) \to a_{m} \leq a_{S(n)}$
\item $\varphi_{50} :\equiv S(n) \in \mathbb{N} \to n \neq \emptyset \to m \leq S(n) \to n \in \mathbb{N} \to a_{n} \leq a_{S(n)}$
\item $\varphi_{51} :\equiv n \in B_{18} \to S(n) \in \mathbb{N} \to n \neq \emptyset \to m \leq S(n) \to a_{n} \leq a_{S(n)}$
\item $\varphi_{52} :\equiv n \in B_{18} \to S(n) \in \mathbb{N} \to n \neq \emptyset \to m \leq S(n) \to a_{m} \leq a_{S(n)}$
\item $\varphi_{53} :\equiv a_{m} = a_{m} \to \mathrm{MonoInc}(a) \to m \in \mathbb{N} \to m \leq 1_{\omega} \to a_{m} \leq a_{m}$
\item $\varphi_{54} :\equiv \forall n\,(n \in \omega \leftrightarrow \forall x\,(\mathrm{Ind}(x) \to n \in x)) \to (n \in \omega \leftrightarrow \forall x\,(\mathrm{Ind}(x) \to n \in x))$
\item $\varphi_{55} :\equiv \forall A_{19}\,\exists B_{18}\,\forall n\,(n \in B_{18} \leftrightarrow n \in A_{19} \land (n \in \mathbb{N} \to m \leq n \to a_{m} \leq a_{n}))$
\item $\varphi_{56} :\equiv S(n) \in \mathbb{N} \to n \neq \emptyset \to m \leq S(n) \to \forall n \in \mathbb{N}\,(a_{n} \leq a_{S(n)})$
\item $\varphi_{57} :\equiv n \in B_{18} \to S(n) \in \mathbb{N} \to n \neq \emptyset \to m \leq S(n) \to m \leq n \to a_{m} \leq a_{n}$
\item $\varphi_{58} :\equiv \forall y\,\forall z\,(S(n) = y \to y = z \to S(n) = z) \to \forall z\,(S(n) = S(\emptyset) \to S(\emptyset) = z \to S(n) = z)$
\item $\varphi_{59} :\equiv \forall w\,(n = w \to S(n) = S(n) \to S(n) = S(w)) \to n = \emptyset \to S(n) = S(n) \to S(n) = S(\emptyset)$
\item $\varphi_{60} :\equiv \mathrm{MonoInc}(a) \to m \in \mathbb{N} \to n \in \omega \to n \in \mathbb{N} \to m \leq n \to a_{m} \leq a_{n}$
\item $\varphi_{61} :\equiv \forall n\,(n \in B_{18} \leftrightarrow n \in \omega \land (n \in \mathbb{N} \to m \leq n \to a_{m} \leq a_{n})) \to \mathrm{Ind}(B_{18})$
\item $\varphi_{62} :\equiv \forall n \in \mathbb{N}\,(m \leq n \to a_{m} \leq a_{n}) \to n \in \mathbb{N} \to m \leq n \to a_{m} \leq a_{n}$
\item $\varphi_{63} :\equiv n \in B_{18} \to S(n) \in \mathbb{N} \to n \neq \emptyset \to m \leq S(n) \to m \leq n \to a_{m} \leq a_{S(n)}$
\item $\varphi_{64} :\equiv S(n) \in B_{18} \leftrightarrow S(n) \in \omega \land (S(n) \in \mathbb{N} \to m \leq S(n) \to a_{m} \leq a_{S(n)})$
\item $\varphi_{65} :\equiv \mathrm{MonoInc}(a) \to m \in \mathbb{N} \to \forall n \in \omega\,(n \in \mathbb{N} \to m \leq n \to a_{m} \leq a_{n})$
\item $\varphi_{66} :\equiv \mathrm{MonoInc}(a) \land m \in \mathbb{N} \land \neg \forall n \in \omega\,(n \in \mathbb{N} \to m \leq n \to a_{m} \leq a_{n})$
\item $\varphi_{67} :\equiv \forall n\,(n \in \mathbb{N} \leftrightarrow n \in \omega \land n \neq \emptyset) \to (n \in \mathbb{N} \leftrightarrow n \in \omega \land n \neq \emptyset)$
\item $\varphi_{68} :\equiv \mathrm{MonoInc}(a) \to m \in \mathbb{N} \to \emptyset \in \mathbb{N} \to m \leq \emptyset \to a_{m} \leq a_{\emptyset}$
\item $\varphi_{69} :\equiv \forall z_{21}\,\forall A_{19}\,\exists B_{18}\,\forall n\,(n \in B_{18} \leftrightarrow n \in A_{19} \land (n \in \mathbb{N} \to m \leq n \to z_{21}(m) \leq z_{21}(n)))$
\item $\varphi_{70} :\equiv \forall n\,(n \in B_{18} \leftrightarrow n \in \omega \land (n \in \mathbb{N} \to m \leq n \to a_{m} \leq a_{n})) \to \emptyset \in B_{18}$
\item $\varphi_{71} :\equiv \forall y\,\forall z\,(a_{m} \leq y \to y \leq z \to a_{m} \leq z) \to \forall z\,(a_{m} \leq a_{n} \to a_{n} \leq z \to a_{m} \leq z)$
\item $\varphi_{72} :\equiv (n \in \mathbb{N} \leftrightarrow n \in \omega \land n \neq \emptyset) \to \mathrm{MonoInc}(a) \to m \in \mathbb{N} \to n \in \mathbb{N} \to n \in \omega$
\item $\varphi_{73} :\equiv \forall n\,(n \in B_{18} \leftrightarrow n \in \omega \land (n \in \mathbb{N} \to m \leq n \to a_{m} \leq a_{n})) \to n \in B_{18} \to n \in B_{18}$
\item $\varphi_{74} :\equiv \mathrm{MonoInc}(a) \land m \in \mathbb{N} \land \forall n\,(n \in B_{18} \leftrightarrow n \in \omega \land (n \in \mathbb{N} \to m \leq n \to a_{m} \leq a_{n}))$
\item $\varphi_{75} :\equiv \forall z_{20}\,\forall z_{21}\,\forall A_{19}\,\exists B_{18}\,\forall n\,(n \in B_{18} \leftrightarrow n \in A_{19} \land (n \in \mathbb{N} \to z_{20} \leq n \to z_{21}(z_{20}) \leq z_{21}(n)))$
\item $\varphi_{76} :\equiv \forall n\,(n \in B_{18} \leftrightarrow n \in \omega \land (n \in \mathbb{N} \to m \leq n \to a_{m} \leq a_{n})) \to n \in B_{18} \to n \in \omega$
\item $\varphi_{77} :\equiv \forall n\,(n \in B_{18} \leftrightarrow n \in \omega \land (n \in \mathbb{N} \to m \leq n \to a_{m} \leq a_{n})) \to n \in B_{18} \to S(n) \in B_{18}$
\item $\varphi_{78} :\equiv \forall n\,(n \in B_{18} \leftrightarrow n \in \omega \land (n \in \mathbb{N} \to m \leq n \to a_{m} \leq a_{n})) \to n \in \omega \to n \in \omega$
\item $\varphi_{79} :\equiv \forall n\,(n \in B_{18} \leftrightarrow n \in \omega \land (n \in \mathbb{N} \to m \leq n \to a_{m} \leq a_{n})) \to n \in B_{18} \to S(n) \in \omega$
\item $\varphi_{80} :\equiv \forall n\,(n \in B_{18} \leftrightarrow n \in \omega \land (n \in \mathbb{N} \to m \leq n \to a_{m} \leq a_{n})) \to \forall n \in B_{18}\,(S(n) \in B_{18})$
\item $\varphi_{81} :\equiv \emptyset \in B_{18} \leftrightarrow \emptyset \in \omega \land (\emptyset \in \mathbb{N} \to m \leq \emptyset \to a_{m} \leq a_{\emptyset})$
\item $\varphi_{82} :\equiv \forall z\,\forall w\,(z = w \to a_{m} = a_{z} \to a_{m} = a_{w}) \to \forall w\,(m = w \to a_{m} = a_{m} \to a_{m} = a_{w})$
\item $\varphi_{83} :\equiv \forall w\,(m = w \to a_{m} = a_{m} \to a_{m} = a_{w}) \to m = S(n) \to a_{m} = a_{m} \to a_{m} = a_{S(n)}$
\item $\varphi_{84} :\equiv \forall w\,(m = w \to a_{m} = a_{m} \to a_{m} = a_{w}) \to m = 1_{\omega} \to a_{m} = a_{m} \to a_{m} = a_{1_{\omega}}$
\item $\varphi_{85} :\equiv (m \in \mathbb{N} \to 1_{\omega} \leq m) \to (\mathrm{MonoInc}(a) \to m \in \mathbb{N} \to m \in \mathbb{N}) \to \varphi_{3}$
\item $\varphi_{86} :\equiv \forall n\,(n \in B_{18} \leftrightarrow n \in \omega \land (n \in \mathbb{N} \to m \leq n \to a_{m} \leq a_{n})) \to n \in \omega \to \forall x\,(\mathrm{Ind}(x) \to n \in x)$
\item $\varphi_{87} :\equiv \forall n\,(n \in \mathbb{N} \leftrightarrow n \in \omega \land n \neq \emptyset) \to (\emptyset \in \mathbb{N} \leftrightarrow \emptyset \in \omega \land \emptyset \neq \emptyset)$
\item $\varphi_{88} :\equiv \forall n\,(n \in B_{18} \leftrightarrow n \in \omega \land (n \in \mathbb{N} \to m \leq n \to a_{m} \leq a_{n})) \to n \in \omega \to \mathrm{Ind}(B_{18}) \to n \in B_{18}$
\item $\varphi_{89} :\equiv \forall n \in \omega\,(n \in \mathbb{N} \to m \leq n \to a_{m} \leq a_{n}) \to n \in \omega \to n \in \mathbb{N} \to m \leq n \to a_{m} \leq a_{n}$
\item $\varphi_{90} :\equiv \forall n\,(m \in \omega \to n \in \omega \to m \leq S(n) \to m \leq n \lor m = S(n)) \to m \in \omega \to n \in \omega \to m \leq S(n) \to m \leq n \lor m = S(n)$
\item $\varphi_{91} :\equiv \forall z\,\forall w\,(z = w \to a_{m} = a_{z} \to a_{m} = a_{w}) \to \forall w\,(S(n) = w \to a_{m} = a_{S(n)} \to a_{m} = a_{w})$
\item $\varphi_{92} :\equiv \forall z\,\forall w\,(z = w \to a_{m} = a_{z} \to a_{m} = a_{w}) \to \forall w\,(1_{\omega} = w \to a_{m} = a_{1_{\omega}} \to a_{m} = a_{w})$
\item $\varphi_{93} :\equiv \forall n\,(n \in B_{18} \leftrightarrow n \in \omega \land (n \in \mathbb{N} \to m \leq n \to a_{m} \leq a_{n})) \to n \in B_{18} \to n \in \omega \to S(n) \in \omega$
\item $\varphi_{94} :\equiv (\mathrm{MonoInc}(a) \to m \in \mathbb{N} \to n \in \mathbb{N} \to n \in \mathbb{N}) \to \mathrm{MonoInc}(a) \to m \in \mathbb{N} \to n \in \mathbb{N} \to m \leq n \to a_{m} \leq a_{n}$
\item $\varphi_{95} :\equiv \forall z\,\forall w\,(z = w \to a_{m} < a_{z} \to a_{m} < a_{w}) \to \forall w\,(m = w \to a_{m} < a_{m} \to a_{m} < a_{w})$
\item $\varphi_{96} :\equiv \forall w\,(m = w \to a_{m} < a_{m} \to a_{m} < a_{w}) \to m = S(n) \to a_{m} < a_{m} \to a_{m} < a_{S(n)}$
\item $\varphi_{97} :\equiv \varphi_{22} \to \forall n\,(m \in \omega \to n \in \omega \to m \leq S(n) \to m \leq n \lor m = S(n))$
\item $\varphi_{98} :\equiv \forall w\,(S(n) = w \to a_{m} = a_{S(n)} \to a_{m} = a_{w}) \to S(n) = 1_{\omega} \to a_{m} = a_{S(n)} \to a_{m} = a_{1_{\omega}}$
\item $\varphi_{99} :\equiv \forall w\,(1_{\omega} = w \to a_{m} = a_{1_{\omega}} \to a_{m} = a_{w}) \to 1_{\omega} = S(n) \to a_{m} = a_{1_{\omega}} \to a_{m} = a_{S(n)}$
\item $\varphi_{100} :\equiv \forall w\,(m = w \to a_{m} < a_{m} \to a_{m} < a_{w}) \to m = 1_{\omega} \to a_{m} < a_{m} \to a_{m} < a_{1_{\omega}}$
\item $\varphi_{101} :\equiv \forall n\,(n \in B_{18} \leftrightarrow n \in \omega \land (n \in \mathbb{N} \to m \leq n \to a_{m} \leq a_{n})) \to n \in B_{18} \to \forall y \in \omega\,(S(y) \in \omega)$
\item $\varphi_{102} :\equiv \forall z\,(a_{m} \leq a_{n} \to a_{n} \leq z \to a_{m} \leq z) \to a_{m} \leq a_{n} \to a_{n} \leq a_{S(n)} \to a_{m} \leq a_{S(n)}$
\item $\varphi_{103} :\equiv \forall n\,(n \in B_{18} \leftrightarrow n \in \omega \land (n \in \mathbb{N} \to m \leq n \to a_{m} \leq a_{n})) \to n \in B_{18} \to n \in \mathbb{N} \to m \leq n \to a_{m} \leq a_{n}$
\item $\varphi_{104} :\equiv \forall z\,\forall w\,(z = w \to a_{m} < a_{z} \to a_{m} < a_{w}) \to \forall w\,(S(n) = w \to a_{m} < a_{S(n)} \to a_{m} < a_{w})$
\item $\varphi_{105} :\equiv \forall z\,\forall w\,(z = w \to a_{m} < a_{z} \to a_{m} < a_{w}) \to \forall w\,(1_{\omega} = w \to a_{m} < a_{1_{\omega}} \to a_{m} < a_{w})$
\item $\varphi_{106} :\equiv \forall n\,(n \in B_{18} \leftrightarrow n \in \omega \land (n \in \mathbb{N} \to m \leq n \to a_{m} \leq a_{n})) \to n \in B_{18} \to S(n) \in \mathbb{N} \to n = \emptyset \to S(n) = 1_{\omega}$
\item $\varphi_{107} :\equiv \forall n\,(n \in B_{18} \leftrightarrow n \in \omega \land (n \in \mathbb{N} \to m \leq n \to a_{m} \leq a_{n})) \to n \in B_{18} \to S(n) \in \mathbb{N} \to n = \emptyset \to 1_{\omega} = S(n)$
\item $\varphi_{108} :\equiv \forall n\,(n \in B_{18} \leftrightarrow n \in \omega \land (n \in \mathbb{N} \to m \leq n \to a_{m} \leq a_{n})) \to n \in \omega \to n \in \mathbb{N} \to m \leq n \to a_{m} \leq a_{n}$
\item $\varphi_{109} :\equiv \forall n\,(n \in B_{18} \leftrightarrow n \in \omega \land (n \in \mathbb{N} \to m \leq n \to a_{m} \leq a_{n})) \to n \in B_{18} \to S(n) \in \mathbb{N} \to n \neq \emptyset \to n \in \mathbb{N}$
\item $\varphi_{110} :\equiv \forall n\,(n \in B_{18} \leftrightarrow n \in \omega \land (n \in \mathbb{N} \to m \leq n \to a_{m} \leq a_{n})) \to n \in B_{18} \to S(n) \in \mathbb{N} \to n = \emptyset \to n = \emptyset$
\item $\varphi_{111} :\equiv \forall n\,(n \in B_{18} \leftrightarrow n \in \omega \land (n \in \mathbb{N} \to m \leq n \to a_{m} \leq a_{n})) \to n \in B_{18} \to S(n) \in \mathbb{N} \to n \neq \emptyset \to n \neq \emptyset$
\item $\varphi_{112} :\equiv \forall n\,(n \in B_{18} \leftrightarrow n \in \omega \land (n \in \mathbb{N} \to m \leq n \to a_{m} \leq a_{n})) \to \emptyset \in \omega \land \forall y \in \omega\,(S(y) \in \omega)$
\item $\varphi_{113} :\equiv \forall n\,(n \in B_{18} \leftrightarrow n \in \omega \land (n \in \mathbb{N} \to m \leq n \to a_{m} \leq a_{n})) \to \forall n \in \omega\,(n \in \mathbb{N} \to m \leq n \to a_{m} \leq a_{n})$
\item $\varphi_{114} :\equiv \forall x\,(\mathrm{Ind}(x) \leftrightarrow \emptyset \in x \land \forall y \in x\,(S(y) \in x)) \to (\mathrm{Ind}(\omega) \leftrightarrow \emptyset \in \omega \land \forall y \in \omega\,(S(y) \in \omega))$
\item $\varphi_{115} :\equiv \forall w\,(S(n) = w \to a_{m} < a_{S(n)} \to a_{m} < a_{w}) \to S(n) = 1_{\omega} \to a_{m} < a_{S(n)} \to a_{m} < a_{1_{\omega}}$
\item $\varphi_{116} :\equiv \forall w\,(1_{\omega} = w \to a_{m} < a_{1_{\omega}} \to a_{m} < a_{w}) \to 1_{\omega} = S(n) \to a_{m} < a_{1_{\omega}} \to a_{m} < a_{S(n)}$
\item $\varphi_{117} :\equiv (\mathrm{MonoInc}(a) \to m \in \mathbb{N} \to n \in \mathbb{N} \to m \leq n \to a_{m} \leq a_{n}) \to \mathrm{MonoInc}(a) \land m \in \mathbb{N} \to n \in \mathbb{N} \to m \leq n \to a_{m} \leq a_{n}$
\item $\varphi_{118} :\equiv \forall n\,(n \in B_{18} \leftrightarrow n \in \omega \land (n \in \mathbb{N} \to m \leq n \to a_{m} \leq a_{n})) \to \varphi_{24}$
\item $\varphi_{119} :\equiv \forall n\,(n \in B_{18} \leftrightarrow n \in \omega \land (n \in \mathbb{N} \to m \leq n \to a_{m} \leq a_{n})) \to n \in B_{18} \to S(n) \in \mathbb{N} \to n = \emptyset \to S(n) = S(\emptyset)$
\item $\varphi_{120} :\equiv \forall n\,(n \in B_{18} \leftrightarrow n \in \omega \land (n \in \mathbb{N} \to m \leq n \to a_{m} \leq a_{n})) \to n \in B_{18} \to S(n) \in \mathbb{N} \to m \leq S(n) \to a_{m} \leq a_{S(n)}$
\item $\varphi_{121} :\equiv \forall n\,(n \in B_{18} \leftrightarrow n \in \omega \land (n \in \mathbb{N} \to m \leq n \to a_{m} \leq a_{n})) \to \varphi_{29}$
\item $\varphi_{122} :\equiv \forall n\,(n \in B_{18} \leftrightarrow n \in \omega \land (n \in \mathbb{N} \to m \leq n \to a_{m} \leq a_{n})) \to (n \in B_{18} \leftrightarrow n \in \omega \land (n \in \mathbb{N} \to m \leq n \to a_{m} \leq a_{n}))$
\item $\varphi_{123} :\equiv \forall n\,(n \in B_{18} \leftrightarrow n \in \omega \land (n \in \mathbb{N} \to m \leq n \to a_{m} \leq a_{n})) \to n \in B_{18} \to n \in \omega \land (n \in \mathbb{N} \to m \leq n \to a_{m} \leq a_{n})$
\item $\varphi_{124} :\equiv \forall n\,(n \in B_{18} \leftrightarrow n \in \omega \land (n \in \mathbb{N} \to m \leq n \to a_{m} \leq a_{n})) \to n \in B_{18} \to S(n) \in \mathbb{N} \to n \neq \emptyset \to m \leq S(n) \to n \in \omega$
\item $\varphi_{125} :\equiv \forall n\,(n \in B_{18} \leftrightarrow n \in \omega \land (n \in \mathbb{N} \to m \leq n \to a_{m} \leq a_{n})) \to n \in B_{18} \to S(n) \in \mathbb{N} \to n \neq \emptyset \to m \leq S(n) \to m \in \omega$
\item $\varphi_{126} :\equiv \forall n\,(n \in B_{18} \leftrightarrow n \in \omega \land (n \in \mathbb{N} \to m \leq n \to a_{m} \leq a_{n})) \to \varphi_{37}$
\item $\varphi_{127} :\equiv \forall n\,(n \in B_{18} \leftrightarrow n \in \omega \land (n \in \mathbb{N} \to m \leq n \to a_{m} \leq a_{n})) \to \forall n\,(n \in B_{18} \leftrightarrow n \in \omega \land (n \in \mathbb{N} \to m \leq n \to a_{m} \leq a_{n}))$
\item $\varphi_{128} :\equiv \forall n\,(n \in B_{18} \leftrightarrow n \in \omega \land (n \in \mathbb{N} \to m \leq n \to a_{m} \leq a_{n})) \to n \in B_{18} \to S(n) \in \mathbb{N} \to n \neq \emptyset \to m \leq n \to a_{m} \leq a_{n}$
\item $\varphi_{129} :\equiv \forall n\,(n \in B_{18} \leftrightarrow n \in \omega \land (n \in \mathbb{N} \to m \leq n \to a_{m} \leq a_{n})) \to n \in B_{18} \to \varphi_{14}$
\item $\varphi_{130} :\equiv \forall n\,(n \in B_{18} \leftrightarrow n \in \omega \land (n \in \mathbb{N} \to m \leq n \to a_{m} \leq a_{n})) \to \varphi_{49}$
\item $\varphi_{131} :\equiv \forall n\,(n \in B_{18} \leftrightarrow n \in \omega \land (n \in \mathbb{N} \to m \leq n \to a_{m} \leq a_{n})) \to \varphi_{51}$
\item $\varphi_{132} :\equiv \forall n\,(n \in B_{18} \leftrightarrow n \in \omega \land (n \in \mathbb{N} \to m \leq n \to a_{m} \leq a_{n})) \to \varphi_{52}$
\item $\varphi_{133} :\equiv \varphi_{62} \to \mathrm{MonoInc}(a) \to m \in \mathbb{N} \to n \in \mathbb{N} \to m \leq n \to a_{m} \leq a_{n}$
\item $\varphi_{134} :\equiv \forall n\,(n \in B_{18} \leftrightarrow n \in \omega \land (n \in \mathbb{N} \to m \leq n \to a_{m} \leq a_{n})) \to \varphi_{57}$
\item $\varphi_{135} :\equiv \forall n\,(n \in B_{18} \leftrightarrow n \in \omega \land (n \in \mathbb{N} \to m \leq n \to a_{m} \leq a_{n})) \to \varphi_{64}$
\item $\varphi_{136} :\equiv \forall n\,(n \in B_{18} \leftrightarrow n \in \omega \land (n \in \mathbb{N} \to m \leq n \to a_{m} \leq a_{n})) \to n \in B_{18} \to \varphi_{41}$
\item $\varphi_{137} :\equiv \forall n\,(n \in B_{18} \leftrightarrow n \in \omega \land (n \in \mathbb{N} \to m \leq n \to a_{m} \leq a_{n})) \to \varphi_{63}$
\item $\varphi_{138} :\equiv \forall n\,(n \in B_{18} \leftrightarrow n \in \omega \land (n \in \mathbb{N} \to m \leq n \to a_{m} \leq a_{n})) \to n \in B_{18} \to \varphi_{47}$
\item $\varphi_{139} :\equiv \forall n\,(n \in B_{18} \leftrightarrow n \in \omega \land (n \in \mathbb{N} \to m \leq n \to a_{m} \leq a_{n})) \to n \in B_{18} \to \varphi_{50}$
\item $\varphi_{140} :\equiv \forall n\,(n \in B_{18} \leftrightarrow n \in \omega \land (n \in \mathbb{N} \to m \leq n \to a_{m} \leq a_{n})) \to n \in B_{18} \to \varphi_{56}$
\item $\varphi_{141} :\equiv (m \leq 1_{\omega} \to 1_{\omega} \leq m \to m = 1_{\omega}) \to \varphi_{4} \to \varphi_{3} \to \varphi_{5}$
\item $\varphi_{142} :\equiv (n \in \omega \leftrightarrow \forall x\,(\mathrm{Ind}(x) \to n \in x)) \to \mathrm{MonoInc}(a) \to m \in \mathbb{N} \to \varphi_{86}$
\item $\varphi_{143} :\equiv \forall n\,(n \in B_{18} \leftrightarrow n \in \omega \land (n \in \mathbb{N} \to m \leq n \to a_{m} \leq a_{n})) \to \varphi_{81}$
\item $\varphi_{144} :\equiv (m = 1_{\omega} \to a_{m} < a_{m} \to a_{m} < a_{1_{\omega}}) \to (m = 1_{\omega} \to a_{m} = a_{m} \to a_{m} = a_{1_{\omega}}) \to \varphi_{21}$
\item $\varphi_{145} :\equiv (S(n) = 1_{\omega} \to 1_{\omega} = S(n)) \to \mathrm{MonoInc}(a) \to m \in \mathbb{N} \to \varphi_{107}$
\item $\varphi_{146} :\equiv (1_{\omega} = S(n) \to S(n) = 1_{\omega}) \to \mathrm{MonoInc}(a) \to m \in \mathbb{N} \to \varphi_{106}$
\item $\varphi_{147} :\equiv (\forall x\,(\mathrm{Ind}(x) \to n \in x) \to \mathrm{Ind}(B_{18}) \to n \in B_{18}) \to \mathrm{MonoInc}(a) \to m \in \mathbb{N} \to \varphi_{88}$
\item $\varphi_{148} :\equiv (\mathrm{MonoInc}(a) \to m \in \mathbb{N} \to m \in \mathbb{N}) \to \mathrm{MonoInc}(a) \to m \in \mathbb{N} \to \varphi_{125}$
\item $\varphi_{149} :\equiv (n = \emptyset \to S(n) = S(n) \to S(n) = S(\emptyset)) \to \mathrm{MonoInc}(a) \to m \in \mathbb{N} \to \varphi_{119}$
\item $\varphi_{150} :\equiv (\forall y \in \omega\,(S(y) \in \omega) \to n \in \omega \to S(n) \in \omega) \to \mathrm{MonoInc}(a) \to m \in \mathbb{N} \to \varphi_{93}$
\item $\varphi_{151} :\equiv (\mathrm{Ind}(\omega) \leftrightarrow \emptyset \in \omega \land \forall y \in \omega\,(S(y) \in \omega)) \to \mathrm{MonoInc}(a) \to m \in \mathbb{N} \to \varphi_{112}$
\item $\varphi_{152} :\equiv (\mathrm{MonoInc}(a) \to m \in \mathbb{N} \to m \leq 1_{\omega} \to a_{m} \leq a_{m}) \to \varphi_{144}$
\item $\varphi_{153} :\equiv (m = S(n) \to a_{m} = a_{m} \to a_{m} = a_{S(n)}) \to \mathrm{MonoInc}(a) \to m \in \mathbb{N} \to \varphi_{138}$
\item $\varphi_{154} :\equiv (\mathrm{MonoInc}(a) \to m \in \mathbb{N} \to \varphi_{77}) \to \varphi_{74} \to n \in B_{18} \to S(n) \in B_{18}$
\item $\varphi_{155} :\equiv (\mathrm{MonoInc}(a) \to m \in \mathbb{N} \to \varphi_{76}) \to \mathrm{MonoInc}(a) \to m \in \mathbb{N} \to \varphi_{79}$
\item $\varphi_{156} :\equiv \forall n\,(\varphi_{74} \to n \in B_{18} \to S(n) \in B_{18}) \to \forall n\,(\varphi_{74}) \to \forall n \in B_{18}\,(S(n) \in B_{18})$
\item $\varphi_{157} :\equiv (\forall n \in \mathbb{N}\,(a_{n} \leq a_{S(n)}) \to n \in \mathbb{N} \to a_{n} \leq a_{S(n)}) \to \mathrm{MonoInc}(a) \to m \in \mathbb{N} \to \varphi_{139}$
\item $\varphi_{158} :\equiv (\mathrm{MonoInc}(a) \to \mathrm{MonoInc}(a)) \to \varphi_{27} \to \mathrm{MonoInc}(a) \to m \in \mathbb{N} \to \varphi_{140}$
\item $\varphi_{159} :\equiv (\mathrm{MonoInc}(a) \to m \in \mathbb{N} \to \varphi_{120}) \to \mathrm{MonoInc}(a) \to m \in \mathbb{N} \to \varphi_{77}$
\item $\varphi_{160} :\equiv (\mathrm{MonoInc}(a) \to m \in \mathbb{N} \to \varphi_{73}) \to \mathrm{MonoInc}(a) \to m \in \mathbb{N} \to \varphi_{123}$
\item $\varphi_{161} :\equiv (\mathrm{MonoInc}(a) \to m \in \mathbb{N} \to \varphi_{123}) \to \mathrm{MonoInc}(a) \to m \in \mathbb{N} \to \varphi_{76}$
\item $\varphi_{162} :\equiv (\mathrm{MonoInc}(a) \to m \in \mathbb{N} \to \varphi_{112}) \to \mathrm{MonoInc}(a) \to m \in \mathbb{N} \to \varphi_{101}$
\item $\varphi_{163} :\equiv (\mathrm{MonoInc}(a) \to m \in \mathbb{N} \to \varphi_{108}) \to \varphi_{74} \to n \in \omega \to n \in \mathbb{N} \to m \leq n \to a_{m} \leq a_{n}$
\item $\varphi_{164} :\equiv (\mathrm{MonoInc}(a) \to m \in \mathbb{N} \to \varphi_{111}) \to \mathrm{MonoInc}(a) \to m \in \mathbb{N} \to \varphi_{109}$
\item $\varphi_{165} :\equiv \forall n\,(\varphi_{74} \to n \in \omega \to n \in \mathbb{N} \to m \leq n \to a_{m} \leq a_{n}) \to \forall n\,(\varphi_{74}) \to \forall n \in \omega\,(n \in \mathbb{N} \to m \leq n \to a_{m} \leq a_{n})$
\item $\varphi_{166} :\equiv (\mathrm{MonoInc}(a) \to m \in \mathbb{N} \to \varphi_{123}) \to \mathrm{MonoInc}(a) \to m \in \mathbb{N} \to \varphi_{103}$
\item $\varphi_{167} :\equiv (\mathrm{MonoInc}(a) \to m \in \mathbb{N} \to \varphi_{122}) \to \mathrm{MonoInc}(a) \to m \in \mathbb{N} \to \varphi_{108}$
\item $\varphi_{168} :\equiv (\mathrm{MonoInc}(a) \to m \in \mathbb{N} \to \varphi_{109}) \to \mathrm{MonoInc}(a) \to m \in \mathbb{N} \to \varphi_{124}$
\item $\varphi_{169} :\equiv (\mathrm{MonoInc}(a) \to m \in \mathbb{N} \to \varphi_{109}) \to \mathrm{MonoInc}(a) \to m \in \mathbb{N} \to \varphi_{128}$
\item $\varphi_{170} :\equiv (\mathrm{MonoInc}(a) \to m \in \mathbb{N} \to \varphi_{118}) \to \mathrm{MonoInc}(a) \to m \in \mathbb{N} \to \varphi_{126}$
\item $\varphi_{171} :\equiv (\mathrm{MonoInc}(a) \to m \in \mathbb{N} \to \varphi_{106}) \to \mathrm{MonoInc}(a) \to m \in \mathbb{N} \to \varphi_{130}$
\item $\varphi_{172} :\equiv (\mathrm{MonoInc}(a) \to m \in \mathbb{N} \to \varphi_{109}) \to \mathrm{MonoInc}(a) \to m \in \mathbb{N} \to \varphi_{131}$
\item $\varphi_{173} :\equiv (\mathrm{MonoInc}(a) \to m \in \mathbb{N} \to \varphi_{132}) \to \mathrm{MonoInc}(a) \to m \in \mathbb{N} \to \varphi_{120}$
\item $\varphi_{174} :\equiv (\mathrm{MonoInc}(a) \to m \in \mathbb{N} \to \varphi_{121}) \to \mathrm{MonoInc}(a) \to m \in \mathbb{N} \to \varphi_{134}$
\item $\varphi_{175} :\equiv (\mathrm{MonoInc}(a) \to m \in \mathbb{N} \to \varphi_{131}) \to \mathrm{MonoInc}(a) \to m \in \mathbb{N} \to \varphi_{137}$
\item $\varphi_{176} :\equiv (\mathrm{MonoInc}(a) \to m \in \mathbb{N} \to \varphi_{138}) \to \mathrm{MonoInc}(a) \to m \in \mathbb{N} \to \varphi_{132}$
\item $\varphi_{177} :\equiv (\mathrm{MonoInc}(a) \to m \in \mathbb{N} \to \varphi_{112}) \to \varphi_{68} \to \mathrm{MonoInc}(a) \to m \in \mathbb{N} \to \varphi_{70}$
\item $\varphi_{178} :\equiv (\mathrm{MonoInc}(a) \to m \in \mathbb{N} \to \varphi_{80}) \to \varphi_{26} \to \varphi_{25} \to \mathrm{MonoInc}(a) \to m \in \mathbb{N} \to \varphi_{61}$
\item $\varphi_{179} :\equiv (\mathrm{MonoInc}(a) \to m \in \mathbb{N} \to \varphi_{110}) \to S(n) = S(n) \to \varphi_{149}$
\item $\varphi_{180} :\equiv (\mathrm{MonoInc}(a) \to m \in \mathbb{N} \to \varphi_{119}) \to S(\emptyset) = 1_{\omega} \to \varphi_{10} \to \mathrm{MonoInc}(a) \to m \in \mathbb{N} \to \varphi_{106}$
\item $\varphi_{181} :\equiv (S(n) = 1_{\omega} \to a_{m} = a_{S(n)} \to a_{m} = a_{1_{\omega}}) \to \varphi_{171}$
\item $\varphi_{182} :\equiv (S(n) = 1_{\omega} \to a_{m} < a_{S(n)} \to a_{m} < a_{1_{\omega}}) \to \varphi_{181}$
\item $\varphi_{183} :\equiv (\mathrm{MonoInc}(a) \to m \in \mathbb{N} \to \varphi_{127}) \to \varphi_{122} \to \mathrm{MonoInc}(a) \to m \in \mathbb{N} \to \varphi_{122}$
\item $\varphi_{184} :\equiv (n \in \mathbb{N} \leftrightarrow n \in \omega \land n \neq \emptyset) \to (\mathrm{MonoInc}(a) \to m \in \mathbb{N} \to \varphi_{76}) \to \varphi_{164}$
\item $\varphi_{185} :\equiv (\mathrm{MonoInc}(a) \to m \in \mathbb{N} \to \varphi_{136}) \to (m = S(n) \to a_{m} < a_{m} \to a_{m} < a_{S(n)}) \to \varphi_{153}$
\item $\varphi_{186} :\equiv (\mathrm{MonoInc}(a) \to m \in \mathbb{N} \to \varphi_{127}) \to \varphi_{135} \to \mathrm{MonoInc}(a) \to m \in \mathbb{N} \to \varphi_{135}$
\item $\varphi_{187} :\equiv (S(n) = 1_{\omega} \to m < S(n) \to m < 1_{\omega}) \to (S(n) = 1_{\omega} \to m = S(n) \to m = 1_{\omega}) \to \varphi_{182}$
\item $\varphi_{188} :\equiv (\mathrm{MonoInc}(a) \to m \in \mathbb{N} \to \varphi_{127}) \to \varphi_{143} \to \mathrm{MonoInc}(a) \to m \in \mathbb{N} \to \varphi_{143}$
\item $\varphi_{189} :\equiv (\mathrm{MonoInc}(a) \to m \in \mathbb{N} \to \varphi_{88}) \to (\mathrm{MonoInc}(a) \to m \in \mathbb{N} \to \varphi_{61}) \to \varphi_{167}$
\item $\varphi_{190} :\equiv (1_{\omega} = S(n) \to a_{m} = a_{1_{\omega}} \to a_{m} = a_{S(n)}) \to \varphi_{187}$
\item $\varphi_{191} :\equiv (\forall B_{18}\,(\varphi_{66} \to \varphi_{42}) \to \forall B_{18}\,(\varphi_{66}) \to \forall B_{18}\,(\varphi_{42})) \to (\varphi_{66} \to \forall B_{18}\,(\varphi_{66})) \to \varphi_{65}$
\item $\varphi_{192} :\equiv (\mathrm{MonoInc}(a) \to m \in \mathbb{N} \to \varphi_{135}) \to (\mathrm{MonoInc}(a) \to m \in \mathbb{N} \to \varphi_{79}) \to \varphi_{159}$
\item $\varphi_{193} :\equiv (a_{m} \leq a_{n} \to a_{n} \leq a_{S(n)} \to a_{m} \leq a_{S(n)}) \to (\mathrm{MonoInc}(a) \to m \in \mathbb{N} \to \varphi_{134}) \to \varphi_{175}$
\item $\varphi_{194} :\equiv (1_{\omega} = S(n) \to a_{m} < a_{1_{\omega}} \to a_{m} < a_{S(n)}) \to \varphi_{190}$
\item $\varphi_{195} :\equiv (\mathrm{MonoInc}(a) \to m \in \mathbb{N} \to \varphi_{125}) \to (\mathrm{MonoInc}(a) \to m \in \mathbb{N} \to \varphi_{124}) \to \varphi_{170}$
\item $\varphi_{196} :\equiv (\mathrm{MonoInc}(a) \to m \in \mathbb{N} \to \varphi_{126}) \to (\mathrm{MonoInc}(a) \to m \in \mathbb{N} \to \varphi_{137}) \to \varphi_{176}$
\item $\varphi_{197} :\equiv \forall n\,(\varphi_{74} \to n \in B_{18} \to S(n) \in B_{18}) \to \varphi_{156} \to (\varphi_{74} \to \forall n\,(\varphi_{74})) \to \mathrm{MonoInc}(a) \to m \in \mathbb{N} \to \varphi_{80}$
\item $\varphi_{198} :\equiv (1_{\omega} = S(n) \to m < 1_{\omega} \to m < S(n)) \to (1_{\omega} = S(n) \to m = 1_{\omega} \to m = S(n)) \to \varphi_{194}$
\item $\varphi_{199} :\equiv \forall n\,(\varphi_{74} \to n \in \omega \to n \in \mathbb{N} \to m \leq n \to a_{m} \leq a_{n}) \to \varphi_{165} \to (\varphi_{74} \to \forall n\,(\varphi_{74})) \to \mathrm{MonoInc}(a) \to m \in \mathbb{N} \to \varphi_{113}$
\item $\varphi_{200} :\equiv (\mathrm{MonoInc}(a) \to m \in \mathbb{N} \to \varphi_{48}) \to \forall B_{18}\,(\varphi_{66} \to \varphi_{42}) \to \varphi_{191}$
\end{abbrevlist}
\end{fullproof}

% 自動生成 (prover/build.js): MonoDecLe  ノード数 428
\begin{fullproof}{[C21]}{$\forall a\,\forall m\,\forall n\,(\mathrm{MonoDec}(a) \to m \in \mathbb{N} \to n \in \mathbb{N} \to m \leq n \to a_{n} \leq a_{m})$}{428}\phantomsection\label{tree:MonoDecLe}
\noindent$\Pi_{1}$:
\begin{ptree}
\AXC{$\mathrm{MonoDec}(a) \to m \in \mathbb{N} \to n \in \mathbb{N} \to n \in \mathbb{N}$\lbl{A1}}
\AXC{$(\mathrm{MonoDec}(a) \to m \in \mathbb{N} \to n \in \mathbb{N} \to n \in \mathbb{N}) \to \varphi_{71}$\lbl{A1}}
\BIC{$\varphi_{71}$}
\end{ptree}
\noindent$\Pi_{2}$:
\begin{ptree}
\AXC{$\varphi_{30}$\lbl{A4}}
\AXC{$\forall y\,(\varphi_{15}) \to \varphi_{34}$\lbl{A3}}
\AXC{$\varphi_{15}$\lbl{A2}}
\UIC{$\forall y\,(\varphi_{15})$}
\AXC{$\forall y\,(\varphi_{15}) \to (\forall y\,(\varphi_{15}) \to \varphi_{34}) \to \varphi_{30} \to \varphi_{25}$\lbl{A1}}
\BIC{$(\forall y\,(\varphi_{15}) \to \varphi_{34}) \to \varphi_{30} \to \varphi_{25}$}
\BIC{$\varphi_{30} \to \varphi_{25}$}
\BIC{$\varphi_{25}$}
\end{ptree}
\noindent$\Pi_{3}$:
\begin{ptree}
\AXC{$\varphi_{31}$\lbl{A4}}
\AXC{$\forall n\,(\varphi_{16}) \to \varphi_{35}$\lbl{A3}}
\AXC{$\varphi_{16}$\lbl{A2}}
\UIC{$\forall n\,(\varphi_{16})$}
\AXC{$\forall n\,(\varphi_{16}) \to (\forall n\,(\varphi_{16}) \to \varphi_{35}) \to \varphi_{31} \to \varphi_{26}$\lbl{A1}}
\BIC{$(\forall n\,(\varphi_{16}) \to \varphi_{35}) \to \varphi_{31} \to \varphi_{26}$}
\BIC{$\varphi_{31} \to \varphi_{26}$}
\BIC{$\varphi_{26}$}
\end{ptree}
\noindent$\Pi_{4}$:
\begin{ptree}
\AXC{$x = x$\lbl{A5}}
\UIC{$\forall x\,(x = x)$}
\AXC{$\forall x\,(x = x) \to a_{m} = a_{m}$\lbl{A2}}
\BIC{$a_{m} = a_{m}$}
\AXC{$a_{m} = a_{m} \to \mathrm{MonoDec}(a) \to m \in \mathbb{N} \to \varphi_{136}$\lbl{A1}}
\BIC{$\mathrm{MonoDec}(a) \to m \in \mathbb{N} \to \varphi_{136}$}
\end{ptree}
\noindent$\Pi_{5}$:
\begin{ptree}
\AXC{$\stackrel{\Pi_{4}}{\vdots}$}\noLine
\UIC{$\mathrm{MonoDec}(a) \to m \in \mathbb{N} \to \varphi_{136}$}
\AXC{$\mathrm{MonoDec}(a) \to m \in \mathbb{N} \to \varphi_{129}$\lbl{A1}}
\AXC{$(\mathrm{MonoDec}(a) \to m \in \mathbb{N} \to \varphi_{129}) \to \varphi_{185}$\lbl{A1}}
\BIC{$\varphi_{185}$}
\BIC{$(m = S(n) \to a_{m} < a_{m} \to a_{S(n)} < a_{m}) \to \varphi_{153}$}
\end{ptree}
\noindent$\Pi_{6}$:
\begin{ptree}
\AXC{$z = w \to a_{z} < a_{m} \to a_{w} < a_{m}$\lbl{A6}}
\UIC{$\forall w\,(z = w \to a_{z} < a_{m} \to a_{w} < a_{m})$}
\UIC{$\forall z\,\forall w\,(z = w \to a_{z} < a_{m} \to a_{w} < a_{m})$}
\AXC{$\varphi_{95}$\lbl{A2}}
\BIC{$\forall w\,(m = w \to a_{m} < a_{m} \to a_{w} < a_{m})$}
\AXC{$\varphi_{96}$\lbl{A2}}
\BIC{$m = S(n) \to a_{m} < a_{m} \to a_{S(n)} < a_{m}$}
\AXC{$\stackrel{\Pi_{5}}{\vdots}$}\noLine
\UIC{$(m = S(n) \to a_{m} < a_{m} \to a_{S(n)} < a_{m}) \to \varphi_{153}$}
\BIC{$\varphi_{153}$}
\end{ptree}
\noindent$\Pi_{7}$:
\begin{ptree}
\AXC{$\varphi_{122}$\lbl{A2}}
\AXC{$\mathrm{MonoDec}(a) \to m \in \mathbb{N} \to \varphi_{127}$\lbl{A1}}
\AXC{$\varphi_{183}$\lbl{A1}}
\BIC{$\varphi_{122} \to \mathrm{MonoDec}(a) \to m \in \mathbb{N} \to \varphi_{122}$}
\BIC{$\mathrm{MonoDec}(a) \to m \in \mathbb{N} \to \varphi_{122}$}
\AXC{$(\mathrm{MonoDec}(a) \to m \in \mathbb{N} \to \varphi_{122}) \to \varphi_{160}$\lbl{A1}}
\BIC{$\varphi_{160}$}
\end{ptree}
\noindent$\Pi_{8}$:
\begin{ptree}
\AXC{$\mathrm{MonoDec}(a) \to m \in \mathbb{N} \to \varphi_{72}$\lbl{A1}}
\AXC{$\stackrel{\Pi_{7}}{\vdots}$}\noLine
\UIC{$\varphi_{160}$}
\BIC{$\mathrm{MonoDec}(a) \to m \in \mathbb{N} \to \varphi_{123}$}
\end{ptree}
\noindent$\Pi_{9}$:
\begin{ptree}
\AXC{$\stackrel{\Pi_{8}}{\vdots}$}\noLine
\UIC{$\mathrm{MonoDec}(a) \to m \in \mathbb{N} \to \varphi_{123}$}
\AXC{$\varphi_{161}$\lbl{A1}}
\BIC{$\mathrm{MonoDec}(a) \to m \in \mathbb{N} \to \varphi_{75}$}
\end{ptree}
\noindent$\Pi_{10}$:
\begin{ptree}
\AXC{$\stackrel{\Pi_{9}}{\vdots}$}\noLine
\UIC{$\mathrm{MonoDec}(a) \to m \in \mathbb{N} \to \varphi_{75}$}
\AXC{$\forall n\,(n \in \mathbb{N} \leftrightarrow n \in \omega \land n \neq \emptyset)$\lbl{Def$\mathbb{N}$}}
\AXC{$\varphi_{67}$\lbl{A2}}
\BIC{$n \in \mathbb{N} \leftrightarrow n \in \omega \land n \neq \emptyset$}
\AXC{$\varphi_{184}$\lbl{A1}}
\BIC{$(\mathrm{MonoDec}(a) \to m \in \mathbb{N} \to \varphi_{75}) \to \varphi_{164}$}
\BIC{$\varphi_{164}$}
\end{ptree}
\noindent$\Pi_{11}$:
\begin{ptree}
\AXC{$\mathrm{MonoDec}(a) \to m \in \mathbb{N} \to \varphi_{111}$\lbl{A1}}
\AXC{$\stackrel{\Pi_{10}}{\vdots}$}\noLine
\UIC{$\varphi_{164}$}
\BIC{$\mathrm{MonoDec}(a) \to m \in \mathbb{N} \to \varphi_{108}$}
\end{ptree}
\noindent$\Pi_{12}$:
\begin{ptree}
\AXC{$\stackrel{\Pi_{8}}{\vdots}$}\noLine
\UIC{$\mathrm{MonoDec}(a) \to m \in \mathbb{N} \to \varphi_{123}$}
\AXC{$\varphi_{166}$\lbl{A1}}
\BIC{$\mathrm{MonoDec}(a) \to m \in \mathbb{N} \to \varphi_{102}$}
\AXC{$(\mathrm{MonoDec}(a) \to m \in \mathbb{N} \to \varphi_{102}) \to \varphi_{169}$\lbl{A1}}
\BIC{$\varphi_{169}$}
\end{ptree}
\noindent$\Pi_{13}$:
\begin{ptree}
\AXC{$\stackrel{\Pi_{11}}{\vdots}$}\noLine
\UIC{$\mathrm{MonoDec}(a) \to m \in \mathbb{N} \to \varphi_{108}$}
\AXC{$\stackrel{\Pi_{12}}{\vdots}$}\noLine
\UIC{$\varphi_{169}$}
\BIC{$\mathrm{MonoDec}(a) \to m \in \mathbb{N} \to \varphi_{128}$}
\AXC{$(\mathrm{MonoDec}(a) \to m \in \mathbb{N} \to \varphi_{128}) \to \varphi_{174}$\lbl{A1}}
\BIC{$\varphi_{174}$}
\end{ptree}
\noindent$\Pi_{14}$:
\begin{ptree}
\AXC{$\forall a\,(\varphi_{27})$\lbl{DefMonoDec}}
\AXC{$\forall a\,(\varphi_{27}) \to \varphi_{27}$\lbl{A2}}
\BIC{$\varphi_{27}$}
\AXC{$\mathrm{MonoDec}(a) \to \mathrm{MonoDec}(a)$\lbl{A1}}
\AXC{$\varphi_{158}$\lbl{A1}}
\BIC{$\varphi_{27} \to \mathrm{MonoDec}(a) \to m \in \mathbb{N} \to \varphi_{140}$}
\BIC{$\mathrm{MonoDec}(a) \to m \in \mathbb{N} \to \varphi_{140}$}
\end{ptree}
\noindent$\Pi_{15}$:
\begin{ptree}
\AXC{$\stackrel{\Pi_{14}}{\vdots}$}\noLine
\UIC{$\mathrm{MonoDec}(a) \to m \in \mathbb{N} \to \varphi_{140}$}
\AXC{$(\mathrm{MonoDec}(a) \to m \in \mathbb{N} \to \varphi_{140}) \to \varphi_{157}$\lbl{A1}}
\BIC{$\varphi_{157}$}
\end{ptree}
\noindent$\Pi_{16}$:
\begin{ptree}
\AXC{$\forall n \in \mathbb{N}\,(a_{S(n)} \leq a_{n}) \to n \in \mathbb{N} \to a_{S(n)} \leq a_{n}$\lbl{A2}}
\AXC{$\stackrel{\Pi_{15}}{\vdots}$}\noLine
\UIC{$\varphi_{157}$}
\BIC{$\mathrm{MonoDec}(a) \to m \in \mathbb{N} \to \varphi_{139}$}
\AXC{$(\mathrm{MonoDec}(a) \to m \in \mathbb{N} \to \varphi_{139}) \to \varphi_{172}$\lbl{A1}}
\BIC{$\varphi_{172}$}
\end{ptree}
\noindent$\Pi_{17}$:
\begin{ptree}
\AXC{$\forall x\,\forall y\,\forall z\,(x \leq y \to y \leq z \to x \leq z)$\lbl{[O5]}}
\AXC{$\varphi_{47}$\lbl{A2}}
\BIC{$\forall y\,\forall z\,(a_{S(n)} \leq y \to y \leq z \to a_{S(n)} \leq z)$}
\AXC{$\varphi_{80}$\lbl{A2}}
\BIC{$\forall z\,(a_{S(n)} \leq a_{n} \to a_{n} \leq z \to a_{S(n)} \leq z)$}
\AXC{$\varphi_{110}$\lbl{A2}}
\BIC{$a_{S(n)} \leq a_{n} \to a_{n} \leq a_{m} \to a_{S(n)} \leq a_{m}$}
\AXC{$\varphi_{193}$\lbl{A1}}
\BIC{$(\mathrm{MonoDec}(a) \to m \in \mathbb{N} \to \varphi_{131}) \to \varphi_{175}$}
\end{ptree}
\noindent$\Pi_{18}$:
\begin{ptree}
\AXC{$\stackrel{\Pi_{11}}{\vdots}$}\noLine
\UIC{$\mathrm{MonoDec}(a) \to m \in \mathbb{N} \to \varphi_{108}$}
\AXC{$\stackrel{\Pi_{16}}{\vdots}$}\noLine
\UIC{$\varphi_{172}$}
\BIC{$\mathrm{MonoDec}(a) \to m \in \mathbb{N} \to \varphi_{131}$}
\AXC{$\stackrel{\Pi_{17}}{\vdots}$}\noLine
\UIC{$(\mathrm{MonoDec}(a) \to m \in \mathbb{N} \to \varphi_{131}) \to \varphi_{175}$}
\BIC{$\varphi_{175}$}
\end{ptree}
\noindent$\Pi_{19}$:
\begin{ptree}
\AXC{$\forall n \in \mathbb{N}\,(n \in \omega)$\lbl{[N7']}}
\AXC{$\forall n \in \mathbb{N}\,(n \in \omega) \to n \in \mathbb{N} \to n \in \omega$\lbl{A2}}
\BIC{$n \in \mathbb{N} \to n \in \omega$}
\AXC{$(n \in \mathbb{N} \to n \in \omega) \to \varphi_{168}$\lbl{A1}}
\BIC{$\varphi_{168}$}
\end{ptree}
\noindent$\Pi_{20}$:
\begin{ptree}
\AXC{$\forall n \in \mathbb{N}\,(n \in \omega)$\lbl{[N7']}}
\AXC{$\forall n \in \mathbb{N}\,(n \in \omega) \to m \in \mathbb{N} \to m \in \omega$\lbl{A2}}
\BIC{$m \in \mathbb{N} \to m \in \omega$}
\AXC{$(m \in \mathbb{N} \to m \in \omega) \to \varphi_{148}$\lbl{A1}}
\BIC{$\varphi_{148}$}
\end{ptree}
\noindent$\Pi_{21}$:
\begin{ptree}
\AXC{$\varphi_{22}$\lbl{[N9]}}
\AXC{$\varphi_{97}$\lbl{A2}}
\BIC{$\forall n\,(m \in \omega \to n \in \omega \to m \leq S(n) \to m \leq n \lor m = S(n))$}
\AXC{$\varphi_{90}$\lbl{A2}}
\BIC{$m \in \omega \to n \in \omega \to m \leq S(n) \to m \leq n \lor m = S(n)$}
\AXC{$(m \in \omega \to n \in \omega \to m \leq S(n) \to m \leq n \lor m = S(n)) \to \varphi_{195}$\lbl{A1}}
\BIC{$\varphi_{195}$}
\end{ptree}
\noindent$\Pi_{22}$:
\begin{ptree}
\AXC{$\stackrel{\Pi_{11}}{\vdots}$}\noLine
\UIC{$\mathrm{MonoDec}(a) \to m \in \mathbb{N} \to \varphi_{108}$}
\AXC{$\stackrel{\Pi_{19}}{\vdots}$}\noLine
\UIC{$\varphi_{168}$}
\BIC{$\mathrm{MonoDec}(a) \to m \in \mathbb{N} \to \varphi_{124}$}
\AXC{$\mathrm{MonoDec}(a) \to m \in \mathbb{N} \to m \in \mathbb{N}$\lbl{A1}}
\AXC{$\stackrel{\Pi_{20}}{\vdots}$}\noLine
\UIC{$\varphi_{148}$}
\BIC{$\mathrm{MonoDec}(a) \to m \in \mathbb{N} \to \varphi_{125}$}
\AXC{$\stackrel{\Pi_{21}}{\vdots}$}\noLine
\UIC{$\varphi_{195}$}
\BIC{$(\mathrm{MonoDec}(a) \to m \in \mathbb{N} \to \varphi_{124}) \to \varphi_{170}$}
\BIC{$\varphi_{170}$}
\end{ptree}
\noindent$\Pi_{23}$:
\begin{ptree}
\AXC{$\mathrm{MonoDec}(a) \to m \in \mathbb{N} \to \varphi_{121}$\lbl{A1}}
\AXC{$\stackrel{\Pi_{13}}{\vdots}$}\noLine
\UIC{$\varphi_{174}$}
\BIC{$\mathrm{MonoDec}(a) \to m \in \mathbb{N} \to \varphi_{134}$}
\AXC{$\stackrel{\Pi_{18}}{\vdots}$}\noLine
\UIC{$\varphi_{175}$}
\BIC{$\mathrm{MonoDec}(a) \to m \in \mathbb{N} \to \varphi_{137}$}
\AXC{$\mathrm{MonoDec}(a) \to m \in \mathbb{N} \to \varphi_{118}$\lbl{A1}}
\AXC{$\stackrel{\Pi_{22}}{\vdots}$}\noLine
\UIC{$\varphi_{170}$}
\BIC{$\mathrm{MonoDec}(a) \to m \in \mathbb{N} \to \varphi_{126}$}
\AXC{$\varphi_{196}$\lbl{A1}}
\BIC{$(\mathrm{MonoDec}(a) \to m \in \mathbb{N} \to \varphi_{137}) \to \varphi_{176}$}
\BIC{$\varphi_{176}$}
\end{ptree}
\noindent$\Pi_{24}$:
\begin{ptree}
\AXC{$x = x$\lbl{A5}}
\AXC{$x = y \to x = x \to y = x$\lbl{A6}}
\AXC{$(x = y \to x = x \to y = x) \to x = x \to x = y \to y = x$\lbl{A1}}
\BIC{$x = x \to x = y \to y = x$}
\BIC{$x = y \to y = x$}
\UIC{$\forall y\,(x = y \to y = x)$}
\UIC{$\forall x\,\forall y\,(x = y \to y = x)$}
\end{ptree}
\noindent$\Pi_{25}$:
\begin{ptree}
\AXC{$y = z \to x = y \to x = z$\lbl{A6}}
\AXC{$(y = z \to x = y \to x = z) \to x = y \to y = z \to x = z$\lbl{A1}}
\BIC{$x = y \to y = z \to x = z$}
\UIC{$\forall z\,(x = y \to y = z \to x = z)$}
\UIC{$\forall y\,\forall z\,(x = y \to y = z \to x = z)$}
\UIC{$\forall x\,\forall y\,\forall z\,(x = y \to y = z \to x = z)$}
\AXC{$\varphi_{17}$\lbl{A2}}
\BIC{$\forall y\,\forall z\,(S(n) = y \to y = z \to S(n) = z)$}
\AXC{$\varphi_{58}$\lbl{A2}}
\BIC{$\forall z\,(S(n) = S(\emptyset) \to S(\emptyset) = z \to S(n) = z)$}
\end{ptree}
\noindent$\Pi_{26}$:
\begin{ptree}
\AXC{$\stackrel{\Pi_{25}}{\vdots}$}\noLine
\UIC{$\forall z\,(S(n) = S(\emptyset) \to S(\emptyset) = z \to S(n) = z)$}
\AXC{$\forall z\,(S(n) = S(\emptyset) \to S(\emptyset) = z \to S(n) = z) \to \varphi_{10}$\lbl{A2}}
\BIC{$\varphi_{10}$}
\end{ptree}
\noindent$\Pi_{27}$:
\begin{ptree}
\AXC{$\stackrel{\Pi_{24}}{\vdots}$}\noLine
\UIC{$\forall x\,\forall y\,(x = y \to y = x)$}
\AXC{$\varphi_{1}$\lbl{A2}}
\BIC{$\forall y\,(1_{\omega} = y \to y = 1_{\omega})$}
\AXC{$\varphi_{23}$\lbl{A2}}
\BIC{$1_{\omega} = S(\emptyset) \to S(\emptyset) = 1_{\omega}$}
\AXC{$1_{\omega} = S(\emptyset)$\lbl{Def$1_\omega$}}
\AXC{$1_{\omega} = S(\emptyset) \to \varphi_{18}$\lbl{A1}}
\BIC{$\varphi_{18}$}
\BIC{$S(\emptyset) = 1_{\omega}$}
\end{ptree}
\noindent$\Pi_{28}$:
\begin{ptree}
\AXC{$x = x$\lbl{A5}}
\UIC{$\forall x\,(x = x)$}
\AXC{$\forall x\,(x = x) \to S(n) = S(n)$\lbl{A2}}
\BIC{$S(n) = S(n)$}
\AXC{$\mathrm{MonoDec}(a) \to m \in \mathbb{N} \to \varphi_{109}$\lbl{A1}}
\AXC{$\varphi_{179}$\lbl{A1}}
\BIC{$S(n) = S(n) \to \varphi_{149}$}
\BIC{$\varphi_{149}$}
\end{ptree}
\noindent$\Pi_{29}$:
\begin{ptree}
\AXC{$z = w \to S(n) = S(z) \to S(n) = S(w)$\lbl{A6}}
\UIC{$\forall w\,(z = w \to S(n) = S(z) \to S(n) = S(w))$}
\UIC{$\forall z\,\forall w\,(z = w \to S(n) = S(z) \to S(n) = S(w))$}
\AXC{$\varphi_{42}$\lbl{A2}}
\BIC{$\forall w\,(n = w \to S(n) = S(n) \to S(n) = S(w))$}
\AXC{$\varphi_{59}$\lbl{A2}}
\BIC{$n = \emptyset \to S(n) = S(n) \to S(n) = S(\emptyset)$}
\AXC{$\stackrel{\Pi_{28}}{\vdots}$}\noLine
\UIC{$\varphi_{149}$}
\BIC{$\mathrm{MonoDec}(a) \to m \in \mathbb{N} \to \varphi_{119}$}
\AXC{$\varphi_{180}$\lbl{A1}}
\BIC{$S(\emptyset) = 1_{\omega} \to \varphi_{10} \to \mathrm{MonoDec}(a) \to m \in \mathbb{N} \to \varphi_{105}$}
\end{ptree}
\noindent$\Pi_{30}$:
\begin{ptree}
\AXC{$\stackrel{\Pi_{26}}{\vdots}$}\noLine
\UIC{$\varphi_{10}$}
\AXC{$\stackrel{\Pi_{27}}{\vdots}$}\noLine
\UIC{$S(\emptyset) = 1_{\omega}$}
\AXC{$\stackrel{\Pi_{29}}{\vdots}$}\noLine
\UIC{$S(\emptyset) = 1_{\omega} \to \varphi_{10} \to \mathrm{MonoDec}(a) \to m \in \mathbb{N} \to \varphi_{105}$}
\BIC{$\varphi_{10} \to \mathrm{MonoDec}(a) \to m \in \mathbb{N} \to \varphi_{105}$}
\BIC{$\mathrm{MonoDec}(a) \to m \in \mathbb{N} \to \varphi_{105}$}
\end{ptree}
\noindent$\Pi_{31}$:
\begin{ptree}
\AXC{$\stackrel{\Pi_{30}}{\vdots}$}\noLine
\UIC{$\mathrm{MonoDec}(a) \to m \in \mathbb{N} \to \varphi_{105}$}
\AXC{$(\mathrm{MonoDec}(a) \to m \in \mathbb{N} \to \varphi_{105}) \to \varphi_{145}$\lbl{A1}}
\BIC{$\varphi_{145}$}
\end{ptree}
\noindent$\Pi_{32}$:
\begin{ptree}
\AXC{$\stackrel{\Pi_{24}}{\vdots}$}\noLine
\UIC{$\forall x\,\forall y\,(x = y \to y = x)$}
\AXC{$\varphi_{2}$\lbl{A2}}
\BIC{$\forall y\,(S(n) = y \to y = S(n))$}
\AXC{$\varphi_{7}$\lbl{A2}}
\BIC{$S(n) = 1_{\omega} \to 1_{\omega} = S(n)$}
\AXC{$\stackrel{\Pi_{31}}{\vdots}$}\noLine
\UIC{$\varphi_{145}$}
\BIC{$\mathrm{MonoDec}(a) \to m \in \mathbb{N} \to \varphi_{106}$}
\end{ptree}
\noindent$\Pi_{33}$:
\begin{ptree}
\AXC{$\stackrel{\Pi_{32}}{\vdots}$}\noLine
\UIC{$\mathrm{MonoDec}(a) \to m \in \mathbb{N} \to \varphi_{106}$}
\AXC{$(\mathrm{MonoDec}(a) \to m \in \mathbb{N} \to \varphi_{106}) \to \varphi_{146}$\lbl{A1}}
\BIC{$\varphi_{146}$}
\end{ptree}
\noindent$\Pi_{34}$:
\begin{ptree}
\AXC{$\forall n \in \mathbb{N}\,(1_{\omega} \leq n)$\lbl{[N12]}}
\AXC{$\forall n \in \mathbb{N}\,(1_{\omega} \leq n) \to m \in \mathbb{N} \to 1_{\omega} \leq m$\lbl{A2}}
\BIC{$m \in \mathbb{N} \to 1_{\omega} \leq m$}
\AXC{$\varphi_{85}$\lbl{A1}}
\BIC{$(\mathrm{MonoDec}(a) \to m \in \mathbb{N} \to m \in \mathbb{N}) \to \varphi_{3}$}
\end{ptree}
\noindent$\Pi_{35}$:
\begin{ptree}
\AXC{$\varphi_{4}$\lbl{A1}}
\AXC{$\forall x\,\forall y\,(x \leq y \to y \leq x \to x = y)$\lbl{[O6]}}
\AXC{$\varphi_{8}$\lbl{A2}}
\BIC{$\forall y\,(m \leq y \to y \leq m \to m = y)$}
\AXC{$\varphi_{9}$\lbl{A2}}
\BIC{$m \leq 1_{\omega} \to 1_{\omega} \leq m \to m = 1_{\omega}$}
\AXC{$\varphi_{141}$\lbl{A1}}
\BIC{$\varphi_{4} \to \varphi_{3} \to \varphi_{5}$}
\BIC{$\varphi_{3} \to \varphi_{5}$}
\end{ptree}
\noindent$\Pi_{36}$:
\begin{ptree}
\AXC{$\mathrm{MonoDec}(a) \to m \in \mathbb{N} \to m \in \mathbb{N}$\lbl{A1}}
\AXC{$\stackrel{\Pi_{34}}{\vdots}$}\noLine
\UIC{$(\mathrm{MonoDec}(a) \to m \in \mathbb{N} \to m \in \mathbb{N}) \to \varphi_{3}$}
\BIC{$\varphi_{3}$}
\AXC{$\stackrel{\Pi_{35}}{\vdots}$}\noLine
\UIC{$\varphi_{3} \to \varphi_{5}$}
\BIC{$\varphi_{5}$}
\AXC{$\varphi_{5} \to \varphi_{152}$\lbl{A1}}
\BIC{$\varphi_{152}$}
\end{ptree}
\noindent$\Pi_{37}$:
\begin{ptree}
\AXC{$x = x$\lbl{A5}}
\UIC{$\forall x\,(x = x)$}
\AXC{$\forall x\,(x = x) \to a_{m} = a_{m}$\lbl{A2}}
\BIC{$a_{m} = a_{m}$}
\AXC{$\varphi_{53}$\lbl{A1}}
\BIC{$\mathrm{MonoDec}(a) \to m \in \mathbb{N} \to m \leq 1_{\omega} \to a_{m} \leq a_{m}$}
\AXC{$\stackrel{\Pi_{36}}{\vdots}$}\noLine
\UIC{$\varphi_{152}$}
\BIC{$\varphi_{144}$}
\end{ptree}
\noindent$\Pi_{38}$:
\begin{ptree}
\AXC{$z = w \to a_{z} < a_{m} \to a_{w} < a_{m}$\lbl{A6}}
\UIC{$\forall w\,(z = w \to a_{z} < a_{m} \to a_{w} < a_{m})$}
\UIC{$\forall z\,\forall w\,(z = w \to a_{z} < a_{m} \to a_{w} < a_{m})$}
\AXC{$\varphi_{95}$\lbl{A2}}
\BIC{$\forall w\,(m = w \to a_{m} < a_{m} \to a_{w} < a_{m})$}
\AXC{$\varphi_{100}$\lbl{A2}}
\BIC{$m = 1_{\omega} \to a_{m} < a_{m} \to a_{1_{\omega}} < a_{m}$}
\AXC{$\stackrel{\Pi_{37}}{\vdots}$}\noLine
\UIC{$\varphi_{144}$}
\BIC{$(m = 1_{\omega} \to a_{m} = a_{m} \to a_{1_{\omega}} = a_{m}) \to \varphi_{21}$}
\end{ptree}
\noindent$\Pi_{39}$:
\begin{ptree}
\AXC{$z = w \to a_{z} = a_{m} \to a_{w} = a_{m}$\lbl{A6}}
\UIC{$\forall w\,(z = w \to a_{z} = a_{m} \to a_{w} = a_{m})$}
\UIC{$\forall z\,\forall w\,(z = w \to a_{z} = a_{m} \to a_{w} = a_{m})$}
\AXC{$\varphi_{82}$\lbl{A2}}
\BIC{$\forall w\,(m = w \to a_{m} = a_{m} \to a_{w} = a_{m})$}
\AXC{$\varphi_{84}$\lbl{A2}}
\BIC{$m = 1_{\omega} \to a_{m} = a_{m} \to a_{1_{\omega}} = a_{m}$}
\AXC{$\stackrel{\Pi_{38}}{\vdots}$}\noLine
\UIC{$(m = 1_{\omega} \to a_{m} = a_{m} \to a_{1_{\omega}} = a_{m}) \to \varphi_{21}$}
\BIC{$\varphi_{21}$}
\end{ptree}
\noindent$\Pi_{40}$:
\begin{ptree}
\AXC{$\stackrel{\Pi_{39}}{\vdots}$}\noLine
\UIC{$\varphi_{21}$}
\AXC{$\stackrel{\Pi_{32}}{\vdots}$}\noLine
\UIC{$\mathrm{MonoDec}(a) \to m \in \mathbb{N} \to \varphi_{106}$}
\AXC{$(\mathrm{MonoDec}(a) \to m \in \mathbb{N} \to \varphi_{106}) \to \varphi_{21} \to \varphi_{198}$\lbl{A1}}
\BIC{$\varphi_{21} \to \varphi_{198}$}
\BIC{$\varphi_{198}$}
\end{ptree}
\noindent$\Pi_{41}$:
\begin{ptree}
\AXC{$z = w \to m < z \to m < w$\lbl{A6}}
\UIC{$\forall w\,(z = w \to m < z \to m < w)$}
\UIC{$\forall z\,\forall w\,(z = w \to m < z \to m < w)$}
\AXC{$\varphi_{33}$\lbl{A2}}
\BIC{$\forall w\,(1_{\omega} = w \to m < 1_{\omega} \to m < w)$}
\AXC{$\varphi_{44}$\lbl{A2}}
\BIC{$1_{\omega} = S(n) \to m < 1_{\omega} \to m < S(n)$}
\AXC{$\stackrel{\Pi_{40}}{\vdots}$}\noLine
\UIC{$\varphi_{198}$}
\BIC{$(1_{\omega} = S(n) \to m = 1_{\omega} \to m = S(n)) \to \varphi_{194}$}
\end{ptree}
\noindent$\Pi_{42}$:
\begin{ptree}
\AXC{$z = w \to m = z \to m = w$\lbl{A6}}
\UIC{$\forall w\,(z = w \to m = z \to m = w)$}
\UIC{$\forall z\,\forall w\,(z = w \to m = z \to m = w)$}
\AXC{$\varphi_{12}$\lbl{A2}}
\BIC{$\forall w\,(1_{\omega} = w \to m = 1_{\omega} \to m = w)$}
\AXC{$\varphi_{20}$\lbl{A2}}
\BIC{$1_{\omega} = S(n) \to m = 1_{\omega} \to m = S(n)$}
\AXC{$\stackrel{\Pi_{41}}{\vdots}$}\noLine
\UIC{$(1_{\omega} = S(n) \to m = 1_{\omega} \to m = S(n)) \to \varphi_{194}$}
\BIC{$\varphi_{194}$}
\end{ptree}
\noindent$\Pi_{43}$:
\begin{ptree}
\AXC{$z = w \to a_{z} < a_{m} \to a_{w} < a_{m}$\lbl{A6}}
\UIC{$\forall w\,(z = w \to a_{z} < a_{m} \to a_{w} < a_{m})$}
\UIC{$\forall z\,\forall w\,(z = w \to a_{z} < a_{m} \to a_{w} < a_{m})$}
\AXC{$\varphi_{104}$\lbl{A2}}
\BIC{$\forall w\,(1_{\omega} = w \to a_{1_{\omega}} < a_{m} \to a_{w} < a_{m})$}
\AXC{$\varphi_{116}$\lbl{A2}}
\BIC{$1_{\omega} = S(n) \to a_{1_{\omega}} < a_{m} \to a_{S(n)} < a_{m}$}
\AXC{$\stackrel{\Pi_{42}}{\vdots}$}\noLine
\UIC{$\varphi_{194}$}
\BIC{$\varphi_{190}$}
\end{ptree}
\noindent$\Pi_{44}$:
\begin{ptree}
\AXC{$z = w \to a_{z} = a_{m} \to a_{w} = a_{m}$\lbl{A6}}
\UIC{$\forall w\,(z = w \to a_{z} = a_{m} \to a_{w} = a_{m})$}
\UIC{$\forall z\,\forall w\,(z = w \to a_{z} = a_{m} \to a_{w} = a_{m})$}
\AXC{$\varphi_{92}$\lbl{A2}}
\BIC{$\forall w\,(1_{\omega} = w \to a_{1_{\omega}} = a_{m} \to a_{w} = a_{m})$}
\AXC{$\varphi_{99}$\lbl{A2}}
\BIC{$1_{\omega} = S(n) \to a_{1_{\omega}} = a_{m} \to a_{S(n)} = a_{m}$}
\AXC{$\stackrel{\Pi_{43}}{\vdots}$}\noLine
\UIC{$\varphi_{190}$}
\BIC{$\varphi_{187}$}
\end{ptree}
\noindent$\Pi_{45}$:
\begin{ptree}
\AXC{$z = w \to m < z \to m < w$\lbl{A6}}
\UIC{$\forall w\,(z = w \to m < z \to m < w)$}
\UIC{$\forall z\,\forall w\,(z = w \to m < z \to m < w)$}
\AXC{$\varphi_{32}$\lbl{A2}}
\BIC{$\forall w\,(S(n) = w \to m < S(n) \to m < w)$}
\AXC{$\varphi_{43}$\lbl{A2}}
\BIC{$S(n) = 1_{\omega} \to m < S(n) \to m < 1_{\omega}$}
\AXC{$\stackrel{\Pi_{44}}{\vdots}$}\noLine
\UIC{$\varphi_{187}$}
\BIC{$(S(n) = 1_{\omega} \to m = S(n) \to m = 1_{\omega}) \to \varphi_{182}$}
\end{ptree}
\noindent$\Pi_{46}$:
\begin{ptree}
\AXC{$z = w \to m = z \to m = w$\lbl{A6}}
\UIC{$\forall w\,(z = w \to m = z \to m = w)$}
\UIC{$\forall z\,\forall w\,(z = w \to m = z \to m = w)$}
\AXC{$\varphi_{11}$\lbl{A2}}
\BIC{$\forall w\,(S(n) = w \to m = S(n) \to m = w)$}
\AXC{$\varphi_{19}$\lbl{A2}}
\BIC{$S(n) = 1_{\omega} \to m = S(n) \to m = 1_{\omega}$}
\AXC{$\stackrel{\Pi_{45}}{\vdots}$}\noLine
\UIC{$(S(n) = 1_{\omega} \to m = S(n) \to m = 1_{\omega}) \to \varphi_{182}$}
\BIC{$\varphi_{182}$}
\end{ptree}
\noindent$\Pi_{47}$:
\begin{ptree}
\AXC{$z = w \to a_{z} < a_{m} \to a_{w} < a_{m}$\lbl{A6}}
\UIC{$\forall w\,(z = w \to a_{z} < a_{m} \to a_{w} < a_{m})$}
\UIC{$\forall z\,\forall w\,(z = w \to a_{z} < a_{m} \to a_{w} < a_{m})$}
\AXC{$\varphi_{103}$\lbl{A2}}
\BIC{$\forall w\,(S(n) = w \to a_{S(n)} < a_{m} \to a_{w} < a_{m})$}
\AXC{$\varphi_{115}$\lbl{A2}}
\BIC{$S(n) = 1_{\omega} \to a_{S(n)} < a_{m} \to a_{1_{\omega}} < a_{m}$}
\AXC{$\stackrel{\Pi_{46}}{\vdots}$}\noLine
\UIC{$\varphi_{182}$}
\BIC{$\varphi_{181}$}
\end{ptree}
\noindent$\Pi_{48}$:
\begin{ptree}
\AXC{$z = w \to a_{z} = a_{m} \to a_{w} = a_{m}$\lbl{A6}}
\UIC{$\forall w\,(z = w \to a_{z} = a_{m} \to a_{w} = a_{m})$}
\UIC{$\forall z\,\forall w\,(z = w \to a_{z} = a_{m} \to a_{w} = a_{m})$}
\AXC{$\varphi_{91}$\lbl{A2}}
\BIC{$\forall w\,(S(n) = w \to a_{S(n)} = a_{m} \to a_{w} = a_{m})$}
\AXC{$\varphi_{98}$\lbl{A2}}
\BIC{$S(n) = 1_{\omega} \to a_{S(n)} = a_{m} \to a_{1_{\omega}} = a_{m}$}
\AXC{$\stackrel{\Pi_{47}}{\vdots}$}\noLine
\UIC{$\varphi_{181}$}
\BIC{$\varphi_{171}$}
\end{ptree}
\noindent$\Pi_{49}$:
\begin{ptree}
\AXC{$\stackrel{\Pi_{24}}{\vdots}$}\noLine
\UIC{$\forall x\,\forall y\,(x = y \to y = x)$}
\AXC{$\varphi_{1}$\lbl{A2}}
\BIC{$\forall y\,(1_{\omega} = y \to y = 1_{\omega})$}
\AXC{$\varphi_{6}$\lbl{A2}}
\BIC{$1_{\omega} = S(n) \to S(n) = 1_{\omega}$}
\AXC{$\stackrel{\Pi_{33}}{\vdots}$}\noLine
\UIC{$\varphi_{146}$}
\BIC{$\mathrm{MonoDec}(a) \to m \in \mathbb{N} \to \varphi_{105}$}
\AXC{$\stackrel{\Pi_{48}}{\vdots}$}\noLine
\UIC{$\varphi_{171}$}
\BIC{$\mathrm{MonoDec}(a) \to m \in \mathbb{N} \to \varphi_{130}$}
\AXC{$(\mathrm{MonoDec}(a) \to m \in \mathbb{N} \to \varphi_{130}) \to \varphi_{173}$\lbl{A1}}
\BIC{$\varphi_{173}$}
\end{ptree}
\noindent$\Pi_{50}$:
\begin{ptree}
\AXC{$\forall x\,(\mathrm{Ind}(x) \leftrightarrow \emptyset \in x \land \forall y \in x\,(S(y) \in x))$\lbl{DefInd}}
\AXC{$\varphi_{114}$\lbl{A2}}
\BIC{$\mathrm{Ind}(\omega) \leftrightarrow \emptyset \in \omega \land \forall y \in \omega\,(S(y) \in \omega)$}
\AXC{$\mathrm{Ind}(\omega)$\lbl{補題(ωは帰納的)}}
\AXC{$\mathrm{Ind}(\omega) \to \varphi_{151}$\lbl{A1}}
\BIC{$\varphi_{151}$}
\BIC{$\mathrm{MonoDec}(a) \to m \in \mathbb{N} \to \varphi_{112}$}
\end{ptree}
\noindent$\Pi_{51}$:
\begin{ptree}
\AXC{$\stackrel{\Pi_{50}}{\vdots}$}\noLine
\UIC{$\mathrm{MonoDec}(a) \to m \in \mathbb{N} \to \varphi_{112}$}
\AXC{$\varphi_{162}$\lbl{A1}}
\BIC{$\mathrm{MonoDec}(a) \to m \in \mathbb{N} \to \varphi_{101}$}
\AXC{$(\mathrm{MonoDec}(a) \to m \in \mathbb{N} \to \varphi_{101}) \to \varphi_{150}$\lbl{A1}}
\BIC{$\varphi_{150}$}
\end{ptree}
\noindent$\Pi_{52}$:
\begin{ptree}
\AXC{$\forall y \in \omega\,(S(y) \in \omega) \to n \in \omega \to S(n) \in \omega$\lbl{A2}}
\AXC{$\stackrel{\Pi_{51}}{\vdots}$}\noLine
\UIC{$\varphi_{150}$}
\BIC{$\mathrm{MonoDec}(a) \to m \in \mathbb{N} \to \varphi_{93}$}
\AXC{$(\mathrm{MonoDec}(a) \to m \in \mathbb{N} \to \varphi_{93}) \to \varphi_{155}$\lbl{A1}}
\BIC{$\varphi_{155}$}
\end{ptree}
\noindent$\Pi_{53}$:
\begin{ptree}
\AXC{$\stackrel{\Pi_{9}}{\vdots}$}\noLine
\UIC{$\mathrm{MonoDec}(a) \to m \in \mathbb{N} \to \varphi_{75}$}
\AXC{$\stackrel{\Pi_{52}}{\vdots}$}\noLine
\UIC{$\varphi_{155}$}
\BIC{$\mathrm{MonoDec}(a) \to m \in \mathbb{N} \to \varphi_{78}$}
\AXC{$\varphi_{135}$\lbl{A2}}
\AXC{$\mathrm{MonoDec}(a) \to m \in \mathbb{N} \to \varphi_{127}$\lbl{A1}}
\AXC{$\varphi_{186}$\lbl{A1}}
\BIC{$\varphi_{135} \to \mathrm{MonoDec}(a) \to m \in \mathbb{N} \to \varphi_{135}$}
\BIC{$\mathrm{MonoDec}(a) \to m \in \mathbb{N} \to \varphi_{135}$}
\AXC{$\varphi_{192}$\lbl{A1}}
\BIC{$(\mathrm{MonoDec}(a) \to m \in \mathbb{N} \to \varphi_{78}) \to \varphi_{159}$}
\BIC{$\varphi_{159}$}
\end{ptree}
\noindent$\Pi_{54}$:
\begin{ptree}
\AXC{$z = w \to a_{z} = a_{m} \to a_{w} = a_{m}$\lbl{A6}}
\UIC{$\forall w\,(z = w \to a_{z} = a_{m} \to a_{w} = a_{m})$}
\UIC{$\forall z\,\forall w\,(z = w \to a_{z} = a_{m} \to a_{w} = a_{m})$}
\AXC{$\varphi_{82}$\lbl{A2}}
\BIC{$\forall w\,(m = w \to a_{m} = a_{m} \to a_{w} = a_{m})$}
\AXC{$\varphi_{83}$\lbl{A2}}
\BIC{$m = S(n) \to a_{m} = a_{m} \to a_{S(n)} = a_{m}$}
\AXC{$\stackrel{\Pi_{6}}{\vdots}$}\noLine
\UIC{$\varphi_{153}$}
\BIC{$\mathrm{MonoDec}(a) \to m \in \mathbb{N} \to \varphi_{138}$}
\AXC{$\stackrel{\Pi_{23}}{\vdots}$}\noLine
\UIC{$\varphi_{176}$}
\BIC{$\mathrm{MonoDec}(a) \to m \in \mathbb{N} \to \varphi_{132}$}
\AXC{$\stackrel{\Pi_{49}}{\vdots}$}\noLine
\UIC{$\varphi_{173}$}
\BIC{$\mathrm{MonoDec}(a) \to m \in \mathbb{N} \to \varphi_{120}$}
\AXC{$\stackrel{\Pi_{53}}{\vdots}$}\noLine
\UIC{$\varphi_{159}$}
\BIC{$\mathrm{MonoDec}(a) \to m \in \mathbb{N} \to \varphi_{76}$}
\AXC{$\varphi_{154}$\lbl{A1}}
\BIC{$\varphi_{73} \to n \in B_{22} \to S(n) \in B_{22}$}
\UIC{$\forall n\,(\varphi_{73} \to n \in B_{22} \to S(n) \in B_{22})$}
\end{ptree}
\noindent$\Pi_{55}$:
\begin{ptree}
\AXC{$\forall n\,(n \in \mathbb{N} \leftrightarrow n \in \omega \land n \neq \emptyset)$\lbl{Def$\mathbb{N}$}}
\AXC{$\varphi_{87}$\lbl{A2}}
\BIC{$\emptyset \in \mathbb{N} \leftrightarrow \emptyset \in \omega \land \emptyset \neq \emptyset$}
\AXC{$(\emptyset \in \mathbb{N} \leftrightarrow \emptyset \in \omega \land \emptyset \neq \emptyset) \to \emptyset = \emptyset \to \varphi_{68}$\lbl{A1}}
\BIC{$\emptyset = \emptyset \to \varphi_{68}$}
\end{ptree}
\noindent$\Pi_{56}$:
\begin{ptree}
\AXC{$\varphi_{143}$\lbl{A2}}
\AXC{$\mathrm{MonoDec}(a) \to m \in \mathbb{N} \to \varphi_{127}$\lbl{A1}}
\AXC{$\varphi_{188}$\lbl{A1}}
\BIC{$\varphi_{143} \to \mathrm{MonoDec}(a) \to m \in \mathbb{N} \to \varphi_{143}$}
\BIC{$\mathrm{MonoDec}(a) \to m \in \mathbb{N} \to \varphi_{143}$}
\AXC{$(\mathrm{MonoDec}(a) \to m \in \mathbb{N} \to \varphi_{143}) \to \varphi_{177}$\lbl{A1}}
\BIC{$\varphi_{177}$}
\end{ptree}
\noindent$\Pi_{57}$:
\begin{ptree}
\AXC{$x = x$\lbl{A5}}
\UIC{$\forall x\,(x = x)$}
\AXC{$\forall x\,(x = x) \to \emptyset = \emptyset$\lbl{A2}}
\BIC{$\emptyset = \emptyset$}
\AXC{$\stackrel{\Pi_{55}}{\vdots}$}\noLine
\UIC{$\emptyset = \emptyset \to \varphi_{68}$}
\BIC{$\varphi_{68}$}
\AXC{$\stackrel{\Pi_{50}}{\vdots}$}\noLine
\UIC{$\mathrm{MonoDec}(a) \to m \in \mathbb{N} \to \varphi_{112}$}
\AXC{$\stackrel{\Pi_{56}}{\vdots}$}\noLine
\UIC{$\varphi_{177}$}
\BIC{$\varphi_{68} \to \mathrm{MonoDec}(a) \to m \in \mathbb{N} \to \varphi_{70}$}
\BIC{$\mathrm{MonoDec}(a) \to m \in \mathbb{N} \to \varphi_{70}$}
\end{ptree}
\noindent$\Pi_{58}$:
\begin{ptree}
\AXC{$\forall x\,(\mathrm{Ind}(x) \leftrightarrow \emptyset \in x \land \forall y \in x\,(S(y) \in x))$\lbl{DefInd}}
\AXC{$\forall x\,(\mathrm{Ind}(x) \leftrightarrow \emptyset \in x \land \forall y \in x\,(S(y) \in x)) \to \varphi_{28}$\lbl{A2}}
\BIC{$\varphi_{28}$}
\end{ptree}
\noindent$\Pi_{59}$:
\begin{ptree}
\AXC{$\stackrel{\Pi_{57}}{\vdots}$}\noLine
\UIC{$\mathrm{MonoDec}(a) \to m \in \mathbb{N} \to \varphi_{70}$}
\AXC{$\stackrel{\Pi_{58}}{\vdots}$}\noLine
\UIC{$\varphi_{28}$}
\AXC{$\varphi_{28} \to (\mathrm{MonoDec}(a) \to m \in \mathbb{N} \to \varphi_{70}) \to \varphi_{178}$\lbl{A1}}
\BIC{$(\mathrm{MonoDec}(a) \to m \in \mathbb{N} \to \varphi_{70}) \to \varphi_{178}$}
\BIC{$\varphi_{178}$}
\end{ptree}
\noindent$\Pi_{60}$:
\begin{ptree}
\AXC{$\varphi_{73} \to \forall n\,(\varphi_{73})$\lbl{A4}}
\AXC{$\varphi_{156}$\lbl{A3}}
\AXC{$\stackrel{\Pi_{54}}{\vdots}$}\noLine
\UIC{$\forall n\,(\varphi_{73} \to n \in B_{22} \to S(n) \in B_{22})$}
\AXC{$\varphi_{197}$\lbl{A1}}
\BIC{$\varphi_{156} \to (\varphi_{73} \to \forall n\,(\varphi_{73})) \to \mathrm{MonoDec}(a) \to m \in \mathbb{N} \to \varphi_{79}$}
\BIC{$(\varphi_{73} \to \forall n\,(\varphi_{73})) \to \mathrm{MonoDec}(a) \to m \in \mathbb{N} \to \varphi_{79}$}
\BIC{$\mathrm{MonoDec}(a) \to m \in \mathbb{N} \to \varphi_{79}$}
\AXC{$\stackrel{\Pi_{59}}{\vdots}$}\noLine
\UIC{$\varphi_{178}$}
\BIC{$\varphi_{26} \to \varphi_{25} \to \mathrm{MonoDec}(a) \to m \in \mathbb{N} \to \varphi_{61}$}
\end{ptree}
\noindent$\Pi_{61}$:
\begin{ptree}
\AXC{$\mathrm{MonoDec}(a) \to m \in \mathbb{N} \to \varphi_{77}$\lbl{A1}}
\AXC{$(\mathrm{MonoDec}(a) \to m \in \mathbb{N} \to \varphi_{77}) \to \varphi_{142}$\lbl{A1}}
\BIC{$\varphi_{142}$}
\end{ptree}
\noindent$\Pi_{62}$:
\begin{ptree}
\AXC{$\forall n\,(n \in \omega \leftrightarrow \forall x\,(\mathrm{Ind}(x) \to n \in x))$\lbl{Def$\omega$}}
\AXC{$\varphi_{54}$\lbl{A2}}
\BIC{$n \in \omega \leftrightarrow \forall x\,(\mathrm{Ind}(x) \to n \in x)$}
\AXC{$\stackrel{\Pi_{61}}{\vdots}$}\noLine
\UIC{$\varphi_{142}$}
\BIC{$\mathrm{MonoDec}(a) \to m \in \mathbb{N} \to \varphi_{86}$}
\AXC{$(\mathrm{MonoDec}(a) \to m \in \mathbb{N} \to \varphi_{86}) \to \varphi_{147}$\lbl{A1}}
\BIC{$\varphi_{147}$}
\end{ptree}
\noindent$\Pi_{63}$:
\begin{ptree}
\AXC{$\forall x\,(\mathrm{Ind}(x) \to n \in x) \to \mathrm{Ind}(B_{22}) \to n \in B_{22}$\lbl{A2}}
\AXC{$\stackrel{\Pi_{62}}{\vdots}$}\noLine
\UIC{$\varphi_{147}$}
\BIC{$\mathrm{MonoDec}(a) \to m \in \mathbb{N} \to \varphi_{88}$}
\AXC{$\varphi_{189}$\lbl{A1}}
\BIC{$(\mathrm{MonoDec}(a) \to m \in \mathbb{N} \to \varphi_{61}) \to \varphi_{167}$}
\end{ptree}
\noindent$\Pi_{64}$:
\begin{ptree}
\AXC{$\stackrel{\Pi_{2}}{\vdots}$}\noLine
\UIC{$\varphi_{25}$}
\AXC{$\stackrel{\Pi_{3}}{\vdots}$}\noLine
\UIC{$\varphi_{26}$}
\AXC{$\stackrel{\Pi_{60}}{\vdots}$}\noLine
\UIC{$\varphi_{26} \to \varphi_{25} \to \mathrm{MonoDec}(a) \to m \in \mathbb{N} \to \varphi_{61}$}
\BIC{$\varphi_{25} \to \mathrm{MonoDec}(a) \to m \in \mathbb{N} \to \varphi_{61}$}
\BIC{$\mathrm{MonoDec}(a) \to m \in \mathbb{N} \to \varphi_{61}$}
\AXC{$\stackrel{\Pi_{63}}{\vdots}$}\noLine
\UIC{$(\mathrm{MonoDec}(a) \to m \in \mathbb{N} \to \varphi_{61}) \to \varphi_{167}$}
\BIC{$\varphi_{167}$}
\end{ptree}
\noindent$\Pi_{65}$:
\begin{ptree}
\AXC{$\varphi_{165}$\lbl{A3}}
\AXC{$\varphi_{122}$\lbl{A2}}
\AXC{$\mathrm{MonoDec}(a) \to m \in \mathbb{N} \to \varphi_{127}$\lbl{A1}}
\AXC{$\varphi_{183}$\lbl{A1}}
\BIC{$\varphi_{122} \to \mathrm{MonoDec}(a) \to m \in \mathbb{N} \to \varphi_{122}$}
\BIC{$\mathrm{MonoDec}(a) \to m \in \mathbb{N} \to \varphi_{122}$}
\AXC{$\stackrel{\Pi_{64}}{\vdots}$}\noLine
\UIC{$\varphi_{167}$}
\BIC{$\mathrm{MonoDec}(a) \to m \in \mathbb{N} \to \varphi_{107}$}
\AXC{$\varphi_{163}$\lbl{A1}}
\BIC{$\varphi_{73} \to n \in \omega \to n \in \mathbb{N} \to m \leq n \to a_{n} \leq a_{m}$}
\UIC{$\forall n\,(\varphi_{73} \to n \in \omega \to n \in \mathbb{N} \to m \leq n \to a_{n} \leq a_{m})$}
\AXC{$\varphi_{199}$\lbl{A1}}
\BIC{$\varphi_{165} \to (\varphi_{73} \to \forall n\,(\varphi_{73})) \to \mathrm{MonoDec}(a) \to m \in \mathbb{N} \to \varphi_{113}$}
\BIC{$(\varphi_{73} \to \forall n\,(\varphi_{73})) \to \mathrm{MonoDec}(a) \to m \in \mathbb{N} \to \varphi_{113}$}
\end{ptree}
\noindent$\Pi_{66}$:
\begin{ptree}
\AXC{$\varphi_{73} \to \forall n\,(\varphi_{73})$\lbl{A4}}
\AXC{$\stackrel{\Pi_{65}}{\vdots}$}\noLine
\UIC{$(\varphi_{73} \to \forall n\,(\varphi_{73})) \to \mathrm{MonoDec}(a) \to m \in \mathbb{N} \to \varphi_{113}$}
\BIC{$\mathrm{MonoDec}(a) \to m \in \mathbb{N} \to \varphi_{113}$}
\AXC{$(\mathrm{MonoDec}(a) \to m \in \mathbb{N} \to \varphi_{113}) \to \varphi_{66} \to \varphi_{41}$\lbl{A1}}
\BIC{$\varphi_{66} \to \varphi_{41}$}
\UIC{$\forall B_{22}\,(\varphi_{66} \to \varphi_{41})$}
\end{ptree}
\noindent$\Pi_{67}$:
\begin{ptree}
\AXC{$\varphi_{74}$\lbl{Sep}}
\AXC{$\varphi_{74} \to \varphi_{69}$\lbl{A2}}
\BIC{$\varphi_{69}$}
\AXC{$\varphi_{69} \to \varphi_{55}$\lbl{A2}}
\BIC{$\varphi_{55}$}
\AXC{$\varphi_{55} \to \varphi_{48}$\lbl{A2}}
\BIC{$\varphi_{48}$}
\AXC{$\varphi_{48} \to \mathrm{MonoDec}(a) \to m \in \mathbb{N} \to \varphi_{48}$\lbl{A1}}
\BIC{$\mathrm{MonoDec}(a) \to m \in \mathbb{N} \to \varphi_{48}$}
\AXC{$\varphi_{200}$\lbl{A1}}
\BIC{$\forall B_{22}\,(\varphi_{66} \to \varphi_{41}) \to \varphi_{191}$}
\end{ptree}
\noindent$\Pi_{68}$:
\begin{ptree}
\AXC{$\varphi_{66} \to \forall B_{22}\,(\varphi_{66})$\lbl{A4}}
\AXC{$\forall B_{22}\,(\varphi_{66} \to \varphi_{41}) \to \forall B_{22}\,(\varphi_{66}) \to \forall B_{22}\,(\varphi_{41})$\lbl{A3}}
\AXC{$\stackrel{\Pi_{66}}{\vdots}$}\noLine
\UIC{$\forall B_{22}\,(\varphi_{66} \to \varphi_{41})$}
\AXC{$\stackrel{\Pi_{67}}{\vdots}$}\noLine
\UIC{$\forall B_{22}\,(\varphi_{66} \to \varphi_{41}) \to \varphi_{191}$}
\BIC{$\varphi_{191}$}
\BIC{$(\varphi_{66} \to \forall B_{22}\,(\varphi_{66})) \to \varphi_{65}$}
\BIC{$\varphi_{65}$}
\end{ptree}
\noindent$\Pi_{69}$:
\begin{ptree}
\AXC{$\varphi_{89}$\lbl{A2}}
\AXC{$\stackrel{\Pi_{68}}{\vdots}$}\noLine
\UIC{$\varphi_{65}$}
\AXC{$\varphi_{65} \to \varphi_{89} \to \varphi_{60}$\lbl{A1}}
\BIC{$\varphi_{89} \to \varphi_{60}$}
\BIC{$\varphi_{60}$}
\AXC{$\varphi_{60} \to (\mathrm{MonoDec}(a) \to m \in \mathbb{N} \to n \in \mathbb{N} \to n \in \omega) \to \varphi_{94}$\lbl{A1}}
\BIC{$(\mathrm{MonoDec}(a) \to m \in \mathbb{N} \to n \in \mathbb{N} \to n \in \omega) \to \varphi_{94}$}
\end{ptree}
\noindent$\Pi_{70}$:
\begin{ptree}
\AXC{$\forall n\,(n \in \mathbb{N} \leftrightarrow n \in \omega \land n \neq \emptyset)$\lbl{Def$\mathbb{N}$}}
\AXC{$\varphi_{67}$\lbl{A2}}
\BIC{$n \in \mathbb{N} \leftrightarrow n \in \omega \land n \neq \emptyset$}
\AXC{$\stackrel{\Pi_{1}}{\vdots}$}\noLine
\UIC{$\varphi_{71}$}
\BIC{$\mathrm{MonoDec}(a) \to m \in \mathbb{N} \to n \in \mathbb{N} \to n \in \omega$}
\AXC{$\stackrel{\Pi_{69}}{\vdots}$}\noLine
\UIC{$(\mathrm{MonoDec}(a) \to m \in \mathbb{N} \to n \in \mathbb{N} \to n \in \omega) \to \varphi_{94}$}
\BIC{$\varphi_{94}$}
\end{ptree}
\noindent$\Pi_{71}$:
\begin{ptree}
\AXC{$\varphi_{37} \to \varphi_{45}$\lbl{A3}}
\AXC{$\mathrm{MonoDec}(a) \to m \in \mathbb{N} \to n \in \mathbb{N} \to n \in \mathbb{N}$\lbl{A1}}
\AXC{$\stackrel{\Pi_{70}}{\vdots}$}\noLine
\UIC{$\varphi_{94}$}
\BIC{$\mathrm{MonoDec}(a) \to m \in \mathbb{N} \to n \in \mathbb{N} \to m \leq n \to a_{n} \leq a_{m}$}
\AXC{$\varphi_{117}$\lbl{A1}}
\BIC{$\mathrm{MonoDec}(a) \land m \in \mathbb{N} \to n \in \mathbb{N} \to m \leq n \to a_{n} \leq a_{m}$}
\UIC{$\varphi_{37}$}
\AXC{$\varphi_{37} \to (\varphi_{37} \to \varphi_{45}) \to \varphi_{13} \to \varphi_{38}$\lbl{A1}}
\BIC{$(\varphi_{37} \to \varphi_{45}) \to \varphi_{13} \to \varphi_{38}$}
\BIC{$\varphi_{13} \to \varphi_{38}$}
\end{ptree}
\noindent 主木:
\begin{ptree}
\AXC{$\varphi_{62}$\lbl{A2}}
\AXC{$\varphi_{13}$\lbl{A4}}
\AXC{$\stackrel{\Pi_{71}}{\vdots}$}\noLine
\UIC{$\varphi_{13} \to \varphi_{38}$}
\BIC{$\varphi_{38}$}
\AXC{$\varphi_{38} \to \varphi_{133}$\lbl{A1}}
\BIC{$\varphi_{133}$}
\BIC{$\mathrm{MonoDec}(a) \to m \in \mathbb{N} \to n \in \mathbb{N} \to m \leq n \to a_{n} \leq a_{m}$}
\UIC{$\varphi_{39}$}
\UIC{$\forall m\,(\varphi_{39})$}
\UIC{$\forall a\,\forall m\,(\varphi_{39})$}
\end{ptree}
\begin{abbrevlist}
\item $\varphi_{1} :\equiv \forall x\,\forall y\,(x = y \to y = x) \to \forall y\,(1_{\omega} = y \to y = 1_{\omega})$
\item $\varphi_{2} :\equiv \forall x\,\forall y\,(x = y \to y = x) \to \forall y\,(S(n) = y \to y = S(n))$
\item $\varphi_{3} :\equiv \mathrm{MonoDec}(a) \to m \in \mathbb{N} \to m \leq 1_{\omega} \to 1_{\omega} \leq m$
\item $\varphi_{4} :\equiv \mathrm{MonoDec}(a) \to m \in \mathbb{N} \to m \leq 1_{\omega} \to m \leq 1_{\omega}$
\item $\varphi_{5} :\equiv \mathrm{MonoDec}(a) \to m \in \mathbb{N} \to m \leq 1_{\omega} \to m = 1_{\omega}$
\item $\varphi_{6} :\equiv \forall y\,(1_{\omega} = y \to y = 1_{\omega}) \to 1_{\omega} = S(n) \to S(n) = 1_{\omega}$
\item $\varphi_{7} :\equiv \forall y\,(S(n) = y \to y = S(n)) \to S(n) = 1_{\omega} \to 1_{\omega} = S(n)$
\item $\varphi_{8} :\equiv \forall x\,\forall y\,(x \leq y \to y \leq x \to x = y) \to \forall y\,(m \leq y \to y \leq m \to m = y)$
\item $\varphi_{9} :\equiv \forall y\,(m \leq y \to y \leq m \to m = y) \to m \leq 1_{\omega} \to 1_{\omega} \leq m \to m = 1_{\omega}$
\item $\varphi_{10} :\equiv S(n) = S(\emptyset) \to S(\emptyset) = 1_{\omega} \to S(n) = 1_{\omega}$
\item $\varphi_{11} :\equiv \forall z\,\forall w\,(z = w \to m = z \to m = w) \to \forall w\,(S(n) = w \to m = S(n) \to m = w)$
\item $\varphi_{12} :\equiv \forall z\,\forall w\,(z = w \to m = z \to m = w) \to \forall w\,(1_{\omega} = w \to m = 1_{\omega} \to m = w)$
\item $\varphi_{13} :\equiv \mathrm{MonoDec}(a) \land m \in \mathbb{N} \to \forall n\,(\mathrm{MonoDec}(a) \land m \in \mathbb{N})$
\item $\varphi_{14} :\equiv S(n) \in \mathbb{N} \to n \neq \emptyset \to m \leq S(n) \to m = S(n) \to m = S(n)$
\item $\varphi_{15} :\equiv \forall n \in B_{22}\,(S(n) \in B_{22}) \to y \in B_{22} \to S(y) \in B_{22}$
\item $\varphi_{16} :\equiv \forall y \in B_{22}\,(S(y) \in B_{22}) \to n \in B_{22} \to S(n) \in B_{22}$
\item $\varphi_{17} :\equiv \forall x\,\forall y\,\forall z\,(x = y \to y = z \to x = z) \to \forall y\,\forall z\,(S(n) = y \to y = z \to S(n) = z)$
\item $\varphi_{18} :\equiv (1_{\omega} = S(\emptyset) \to S(\emptyset) = 1_{\omega}) \to S(\emptyset) = 1_{\omega}$
\item $\varphi_{19} :\equiv \forall w\,(S(n) = w \to m = S(n) \to m = w) \to S(n) = 1_{\omega} \to m = S(n) \to m = 1_{\omega}$
\item $\varphi_{20} :\equiv \forall w\,(1_{\omega} = w \to m = 1_{\omega} \to m = w) \to 1_{\omega} = S(n) \to m = 1_{\omega} \to m = S(n)$
\item $\varphi_{21} :\equiv \mathrm{MonoDec}(a) \to m \in \mathbb{N} \to m \leq 1_{\omega} \to a_{1_{\omega}} \leq a_{m}$
\item $\varphi_{22} :\equiv \forall m\,\forall n\,(m \in \omega \to n \in \omega \to m \leq S(n) \to m \leq n \lor m = S(n))$
\item $\varphi_{23} :\equiv \forall y\,(1_{\omega} = y \to y = 1_{\omega}) \to 1_{\omega} = S(\emptyset) \to S(\emptyset) = 1_{\omega}$
\item $\varphi_{24} :\equiv n \in B_{22} \to S(n) \in \mathbb{N} \to n \neq \emptyset \to m \leq S(n) \to m \leq S(n)$
\item $\varphi_{25} :\equiv \forall n \in B_{22}\,(S(n) \in B_{22}) \to \forall y \in B_{22}\,(S(y) \in B_{22})$
\item $\varphi_{26} :\equiv \forall y \in B_{22}\,(S(y) \in B_{22}) \to \forall n \in B_{22}\,(S(n) \in B_{22})$
\item $\varphi_{27} :\equiv \mathrm{MonoDec}(a) \leftrightarrow \mathrm{Seq}(a) \land \forall n \in \mathbb{N}\,(a_{S(n)} \leq a_{n})$
\item $\varphi_{28} :\equiv \mathrm{Ind}(B_{22}) \leftrightarrow \emptyset \in B_{22} \land \forall y \in B_{22}\,(S(y) \in B_{22})$
\item $\varphi_{29} :\equiv n \in B_{22} \to S(n) \in \mathbb{N} \to n \neq \emptyset \to m \leq S(n) \to m \leq n \to m \leq n$
\item $\varphi_{30} :\equiv \forall n \in B_{22}\,(S(n) \in B_{22}) \to \forall y\,\forall n \in B_{22}\,(S(n) \in B_{22})$
\item $\varphi_{31} :\equiv \forall y \in B_{22}\,(S(y) \in B_{22}) \to \forall n\,\forall y \in B_{22}\,(S(y) \in B_{22})$
\item $\varphi_{32} :\equiv \forall z\,\forall w\,(z = w \to m < z \to m < w) \to \forall w\,(S(n) = w \to m < S(n) \to m < w)$
\item $\varphi_{33} :\equiv \forall z\,\forall w\,(z = w \to m < z \to m < w) \to \forall w\,(1_{\omega} = w \to m < 1_{\omega} \to m < w)$
\item $\varphi_{34} :\equiv \forall y\,\forall n \in B_{22}\,(S(n) \in B_{22}) \to \forall y \in B_{22}\,(S(y) \in B_{22})$
\item $\varphi_{35} :\equiv \forall n\,\forall y \in B_{22}\,(S(y) \in B_{22}) \to \forall n \in B_{22}\,(S(n) \in B_{22})$
\item $\varphi_{36} :\equiv n \in B_{22} \to S(n) \in \mathbb{N} \to n \neq \emptyset \to m \leq S(n) \to m \leq n \lor m = S(n)$
\item $\varphi_{37} :\equiv \forall n\,(\mathrm{MonoDec}(a) \land m \in \mathbb{N} \to n \in \mathbb{N} \to m \leq n \to a_{n} \leq a_{m})$
\item $\varphi_{38} :\equiv \mathrm{MonoDec}(a) \to m \in \mathbb{N} \to \forall n \in \mathbb{N}\,(m \leq n \to a_{n} \leq a_{m})$
\item $\varphi_{39} :\equiv \forall n\,(\mathrm{MonoDec}(a) \to m \in \mathbb{N} \to n \in \mathbb{N} \to m \leq n \to a_{n} \leq a_{m})$
\item $\varphi_{40} :\equiv S(n) \in \mathbb{N} \to n \neq \emptyset \to m \leq S(n) \to m = S(n) \to a_{m} \leq a_{m}$
\item $\varphi_{41} :\equiv \neg \forall n\,(n \in B_{22} \leftrightarrow n \in \omega \land (n \in \mathbb{N} \to m \leq n \to a_{n} \leq a_{m}))$
\item $\varphi_{42} :\equiv \forall z\,\forall w\,(z = w \to S(n) = S(z) \to S(n) = S(w)) \to \forall w\,(n = w \to S(n) = S(n) \to S(n) = S(w))$
\item $\varphi_{43} :\equiv \forall w\,(S(n) = w \to m < S(n) \to m < w) \to S(n) = 1_{\omega} \to m < S(n) \to m < 1_{\omega}$
\item $\varphi_{44} :\equiv \forall w\,(1_{\omega} = w \to m < 1_{\omega} \to m < w) \to 1_{\omega} = S(n) \to m < 1_{\omega} \to m < S(n)$
\item $\varphi_{45} :\equiv \forall n\,(\mathrm{MonoDec}(a) \land m \in \mathbb{N}) \to \forall n \in \mathbb{N}\,(m \leq n \to a_{n} \leq a_{m})$
\item $\varphi_{46} :\equiv S(n) \in \mathbb{N} \to n \neq \emptyset \to m \leq S(n) \to m = S(n) \to a_{S(n)} \leq a_{m}$
\item $\varphi_{47} :\equiv \forall x\,\forall y\,\forall z\,(x \leq y \to y \leq z \to x \leq z) \to \forall y\,\forall z\,(a_{S(n)} \leq y \to y \leq z \to a_{S(n)} \leq z)$
\item $\varphi_{48} :\equiv \exists B_{22}\,\forall n\,(n \in B_{22} \leftrightarrow n \in \omega \land (n \in \mathbb{N} \to m \leq n \to a_{n} \leq a_{m}))$
\item $\varphi_{49} :\equiv n \in B_{22} \to S(n) \in \mathbb{N} \to n = \emptyset \to m \leq S(n) \to a_{S(n)} \leq a_{m}$
\item $\varphi_{50} :\equiv S(n) \in \mathbb{N} \to n \neq \emptyset \to m \leq S(n) \to n \in \mathbb{N} \to a_{S(n)} \leq a_{n}$
\item $\varphi_{51} :\equiv n \in B_{22} \to S(n) \in \mathbb{N} \to n \neq \emptyset \to m \leq S(n) \to a_{S(n)} \leq a_{n}$
\item $\varphi_{52} :\equiv n \in B_{22} \to S(n) \in \mathbb{N} \to n \neq \emptyset \to m \leq S(n) \to a_{S(n)} \leq a_{m}$
\item $\varphi_{53} :\equiv a_{m} = a_{m} \to \mathrm{MonoDec}(a) \to m \in \mathbb{N} \to m \leq 1_{\omega} \to a_{m} \leq a_{m}$
\item $\varphi_{54} :\equiv \forall n\,(n \in \omega \leftrightarrow \forall x\,(\mathrm{Ind}(x) \to n \in x)) \to (n \in \omega \leftrightarrow \forall x\,(\mathrm{Ind}(x) \to n \in x))$
\item $\varphi_{55} :\equiv \forall A_{23}\,\exists B_{22}\,\forall n\,(n \in B_{22} \leftrightarrow n \in A_{23} \land (n \in \mathbb{N} \to m \leq n \to a_{n} \leq a_{m}))$
\item $\varphi_{56} :\equiv n \in B_{22} \to S(n) \in \mathbb{N} \to n \neq \emptyset \to m \leq S(n) \to m \leq n \to a_{n} \leq a_{m}$
\item $\varphi_{57} :\equiv S(n) \in \mathbb{N} \to n \neq \emptyset \to m \leq S(n) \to \forall n \in \mathbb{N}\,(a_{S(n)} \leq a_{n})$
\item $\varphi_{58} :\equiv \forall y\,\forall z\,(S(n) = y \to y = z \to S(n) = z) \to \forall z\,(S(n) = S(\emptyset) \to S(\emptyset) = z \to S(n) = z)$
\item $\varphi_{59} :\equiv \forall w\,(n = w \to S(n) = S(n) \to S(n) = S(w)) \to n = \emptyset \to S(n) = S(n) \to S(n) = S(\emptyset)$
\item $\varphi_{60} :\equiv \mathrm{MonoDec}(a) \to m \in \mathbb{N} \to n \in \omega \to n \in \mathbb{N} \to m \leq n \to a_{n} \leq a_{m}$
\item $\varphi_{61} :\equiv \forall n\,(n \in B_{22} \leftrightarrow n \in \omega \land (n \in \mathbb{N} \to m \leq n \to a_{n} \leq a_{m})) \to \mathrm{Ind}(B_{22})$
\item $\varphi_{62} :\equiv \forall n \in \mathbb{N}\,(m \leq n \to a_{n} \leq a_{m}) \to n \in \mathbb{N} \to m \leq n \to a_{n} \leq a_{m}$
\item $\varphi_{63} :\equiv n \in B_{22} \to S(n) \in \mathbb{N} \to n \neq \emptyset \to m \leq S(n) \to m \leq n \to a_{S(n)} \leq a_{m}$
\item $\varphi_{64} :\equiv S(n) \in B_{22} \leftrightarrow S(n) \in \omega \land (S(n) \in \mathbb{N} \to m \leq S(n) \to a_{S(n)} \leq a_{m})$
\item $\varphi_{65} :\equiv \mathrm{MonoDec}(a) \to m \in \mathbb{N} \to \forall n \in \omega\,(n \in \mathbb{N} \to m \leq n \to a_{n} \leq a_{m})$
\item $\varphi_{66} :\equiv \mathrm{MonoDec}(a) \land m \in \mathbb{N} \land \neg \forall n \in \omega\,(n \in \mathbb{N} \to m \leq n \to a_{n} \leq a_{m})$
\item $\varphi_{67} :\equiv \forall n\,(n \in \mathbb{N} \leftrightarrow n \in \omega \land n \neq \emptyset) \to (n \in \mathbb{N} \leftrightarrow n \in \omega \land n \neq \emptyset)$
\item $\varphi_{68} :\equiv \mathrm{MonoDec}(a) \to m \in \mathbb{N} \to \emptyset \in \mathbb{N} \to m \leq \emptyset \to a_{\emptyset} \leq a_{m}$
\item $\varphi_{69} :\equiv \forall z_{25}\,\forall A_{23}\,\exists B_{22}\,\forall n\,(n \in B_{22} \leftrightarrow n \in A_{23} \land (n \in \mathbb{N} \to m \leq n \to z_{25}(n) \leq z_{25}(m)))$
\item $\varphi_{70} :\equiv \forall n\,(n \in B_{22} \leftrightarrow n \in \omega \land (n \in \mathbb{N} \to m \leq n \to a_{n} \leq a_{m})) \to \emptyset \in B_{22}$
\item $\varphi_{71} :\equiv (n \in \mathbb{N} \leftrightarrow n \in \omega \land n \neq \emptyset) \to \mathrm{MonoDec}(a) \to m \in \mathbb{N} \to n \in \mathbb{N} \to n \in \omega$
\item $\varphi_{72} :\equiv \forall n\,(n \in B_{22} \leftrightarrow n \in \omega \land (n \in \mathbb{N} \to m \leq n \to a_{n} \leq a_{m})) \to n \in B_{22} \to n \in B_{22}$
\item $\varphi_{73} :\equiv \mathrm{MonoDec}(a) \land m \in \mathbb{N} \land \forall n\,(n \in B_{22} \leftrightarrow n \in \omega \land (n \in \mathbb{N} \to m \leq n \to a_{n} \leq a_{m}))$
\item $\varphi_{74} :\equiv \forall z_{24}\,\forall z_{25}\,\forall A_{23}\,\exists B_{22}\,\forall n\,(n \in B_{22} \leftrightarrow n \in A_{23} \land (n \in \mathbb{N} \to z_{24} \leq n \to z_{25}(n) \leq z_{25}(z_{24})))$
\item $\varphi_{75} :\equiv \forall n\,(n \in B_{22} \leftrightarrow n \in \omega \land (n \in \mathbb{N} \to m \leq n \to a_{n} \leq a_{m})) \to n \in B_{22} \to n \in \omega$
\item $\varphi_{76} :\equiv \forall n\,(n \in B_{22} \leftrightarrow n \in \omega \land (n \in \mathbb{N} \to m \leq n \to a_{n} \leq a_{m})) \to n \in B_{22} \to S(n) \in B_{22}$
\item $\varphi_{77} :\equiv \forall n\,(n \in B_{22} \leftrightarrow n \in \omega \land (n \in \mathbb{N} \to m \leq n \to a_{n} \leq a_{m})) \to n \in \omega \to n \in \omega$
\item $\varphi_{78} :\equiv \forall n\,(n \in B_{22} \leftrightarrow n \in \omega \land (n \in \mathbb{N} \to m \leq n \to a_{n} \leq a_{m})) \to n \in B_{22} \to S(n) \in \omega$
\item $\varphi_{79} :\equiv \forall n\,(n \in B_{22} \leftrightarrow n \in \omega \land (n \in \mathbb{N} \to m \leq n \to a_{n} \leq a_{m})) \to \forall n \in B_{22}\,(S(n) \in B_{22})$
\item $\varphi_{80} :\equiv \forall y\,\forall z\,(a_{S(n)} \leq y \to y \leq z \to a_{S(n)} \leq z) \to \forall z\,(a_{S(n)} \leq a_{n} \to a_{n} \leq z \to a_{S(n)} \leq z)$
\item $\varphi_{81} :\equiv \emptyset \in B_{22} \leftrightarrow \emptyset \in \omega \land (\emptyset \in \mathbb{N} \to m \leq \emptyset \to a_{\emptyset} \leq a_{m})$
\item $\varphi_{82} :\equiv \forall z\,\forall w\,(z = w \to a_{z} = a_{m} \to a_{w} = a_{m}) \to \forall w\,(m = w \to a_{m} = a_{m} \to a_{w} = a_{m})$
\item $\varphi_{83} :\equiv \forall w\,(m = w \to a_{m} = a_{m} \to a_{w} = a_{m}) \to m = S(n) \to a_{m} = a_{m} \to a_{S(n)} = a_{m}$
\item $\varphi_{84} :\equiv \forall w\,(m = w \to a_{m} = a_{m} \to a_{w} = a_{m}) \to m = 1_{\omega} \to a_{m} = a_{m} \to a_{1_{\omega}} = a_{m}$
\item $\varphi_{85} :\equiv (m \in \mathbb{N} \to 1_{\omega} \leq m) \to (\mathrm{MonoDec}(a) \to m \in \mathbb{N} \to m \in \mathbb{N}) \to \varphi_{3}$
\item $\varphi_{86} :\equiv \forall n\,(n \in B_{22} \leftrightarrow n \in \omega \land (n \in \mathbb{N} \to m \leq n \to a_{n} \leq a_{m})) \to n \in \omega \to \forall x\,(\mathrm{Ind}(x) \to n \in x)$
\item $\varphi_{87} :\equiv \forall n\,(n \in \mathbb{N} \leftrightarrow n \in \omega \land n \neq \emptyset) \to (\emptyset \in \mathbb{N} \leftrightarrow \emptyset \in \omega \land \emptyset \neq \emptyset)$
\item $\varphi_{88} :\equiv \forall n\,(n \in B_{22} \leftrightarrow n \in \omega \land (n \in \mathbb{N} \to m \leq n \to a_{n} \leq a_{m})) \to n \in \omega \to \mathrm{Ind}(B_{22}) \to n \in B_{22}$
\item $\varphi_{89} :\equiv \forall n \in \omega\,(n \in \mathbb{N} \to m \leq n \to a_{n} \leq a_{m}) \to n \in \omega \to n \in \mathbb{N} \to m \leq n \to a_{n} \leq a_{m}$
\item $\varphi_{90} :\equiv \forall n\,(m \in \omega \to n \in \omega \to m \leq S(n) \to m \leq n \lor m = S(n)) \to m \in \omega \to n \in \omega \to m \leq S(n) \to m \leq n \lor m = S(n)$
\item $\varphi_{91} :\equiv \forall z\,\forall w\,(z = w \to a_{z} = a_{m} \to a_{w} = a_{m}) \to \forall w\,(S(n) = w \to a_{S(n)} = a_{m} \to a_{w} = a_{m})$
\item $\varphi_{92} :\equiv \forall z\,\forall w\,(z = w \to a_{z} = a_{m} \to a_{w} = a_{m}) \to \forall w\,(1_{\omega} = w \to a_{1_{\omega}} = a_{m} \to a_{w} = a_{m})$
\item $\varphi_{93} :\equiv \forall n\,(n \in B_{22} \leftrightarrow n \in \omega \land (n \in \mathbb{N} \to m \leq n \to a_{n} \leq a_{m})) \to n \in B_{22} \to n \in \omega \to S(n) \in \omega$
\item $\varphi_{94} :\equiv (\mathrm{MonoDec}(a) \to m \in \mathbb{N} \to n \in \mathbb{N} \to n \in \mathbb{N}) \to \mathrm{MonoDec}(a) \to m \in \mathbb{N} \to n \in \mathbb{N} \to m \leq n \to a_{n} \leq a_{m}$
\item $\varphi_{95} :\equiv \forall z\,\forall w\,(z = w \to a_{z} < a_{m} \to a_{w} < a_{m}) \to \forall w\,(m = w \to a_{m} < a_{m} \to a_{w} < a_{m})$
\item $\varphi_{96} :\equiv \forall w\,(m = w \to a_{m} < a_{m} \to a_{w} < a_{m}) \to m = S(n) \to a_{m} < a_{m} \to a_{S(n)} < a_{m}$
\item $\varphi_{97} :\equiv \varphi_{22} \to \forall n\,(m \in \omega \to n \in \omega \to m \leq S(n) \to m \leq n \lor m = S(n))$
\item $\varphi_{98} :\equiv \forall w\,(S(n) = w \to a_{S(n)} = a_{m} \to a_{w} = a_{m}) \to S(n) = 1_{\omega} \to a_{S(n)} = a_{m} \to a_{1_{\omega}} = a_{m}$
\item $\varphi_{99} :\equiv \forall w\,(1_{\omega} = w \to a_{1_{\omega}} = a_{m} \to a_{w} = a_{m}) \to 1_{\omega} = S(n) \to a_{1_{\omega}} = a_{m} \to a_{S(n)} = a_{m}$
\item $\varphi_{100} :\equiv \forall w\,(m = w \to a_{m} < a_{m} \to a_{w} < a_{m}) \to m = 1_{\omega} \to a_{m} < a_{m} \to a_{1_{\omega}} < a_{m}$
\item $\varphi_{101} :\equiv \forall n\,(n \in B_{22} \leftrightarrow n \in \omega \land (n \in \mathbb{N} \to m \leq n \to a_{n} \leq a_{m})) \to n \in B_{22} \to \forall y \in \omega\,(S(y) \in \omega)$
\item $\varphi_{102} :\equiv \forall n\,(n \in B_{22} \leftrightarrow n \in \omega \land (n \in \mathbb{N} \to m \leq n \to a_{n} \leq a_{m})) \to n \in B_{22} \to n \in \mathbb{N} \to m \leq n \to a_{n} \leq a_{m}$
\item $\varphi_{103} :\equiv \forall z\,\forall w\,(z = w \to a_{z} < a_{m} \to a_{w} < a_{m}) \to \forall w\,(S(n) = w \to a_{S(n)} < a_{m} \to a_{w} < a_{m})$
\item $\varphi_{104} :\equiv \forall z\,\forall w\,(z = w \to a_{z} < a_{m} \to a_{w} < a_{m}) \to \forall w\,(1_{\omega} = w \to a_{1_{\omega}} < a_{m} \to a_{w} < a_{m})$
\item $\varphi_{105} :\equiv \forall n\,(n \in B_{22} \leftrightarrow n \in \omega \land (n \in \mathbb{N} \to m \leq n \to a_{n} \leq a_{m})) \to n \in B_{22} \to S(n) \in \mathbb{N} \to n = \emptyset \to S(n) = 1_{\omega}$
\item $\varphi_{106} :\equiv \forall n\,(n \in B_{22} \leftrightarrow n \in \omega \land (n \in \mathbb{N} \to m \leq n \to a_{n} \leq a_{m})) \to n \in B_{22} \to S(n) \in \mathbb{N} \to n = \emptyset \to 1_{\omega} = S(n)$
\item $\varphi_{107} :\equiv \forall n\,(n \in B_{22} \leftrightarrow n \in \omega \land (n \in \mathbb{N} \to m \leq n \to a_{n} \leq a_{m})) \to n \in \omega \to n \in \mathbb{N} \to m \leq n \to a_{n} \leq a_{m}$
\item $\varphi_{108} :\equiv \forall n\,(n \in B_{22} \leftrightarrow n \in \omega \land (n \in \mathbb{N} \to m \leq n \to a_{n} \leq a_{m})) \to n \in B_{22} \to S(n) \in \mathbb{N} \to n \neq \emptyset \to n \in \mathbb{N}$
\item $\varphi_{109} :\equiv \forall n\,(n \in B_{22} \leftrightarrow n \in \omega \land (n \in \mathbb{N} \to m \leq n \to a_{n} \leq a_{m})) \to n \in B_{22} \to S(n) \in \mathbb{N} \to n = \emptyset \to n = \emptyset$
\item $\varphi_{110} :\equiv \forall z\,(a_{S(n)} \leq a_{n} \to a_{n} \leq z \to a_{S(n)} \leq z) \to a_{S(n)} \leq a_{n} \to a_{n} \leq a_{m} \to a_{S(n)} \leq a_{m}$
\item $\varphi_{111} :\equiv \forall n\,(n \in B_{22} \leftrightarrow n \in \omega \land (n \in \mathbb{N} \to m \leq n \to a_{n} \leq a_{m})) \to n \in B_{22} \to S(n) \in \mathbb{N} \to n \neq \emptyset \to n \neq \emptyset$
\item $\varphi_{112} :\equiv \forall n\,(n \in B_{22} \leftrightarrow n \in \omega \land (n \in \mathbb{N} \to m \leq n \to a_{n} \leq a_{m})) \to \emptyset \in \omega \land \forall y \in \omega\,(S(y) \in \omega)$
\item $\varphi_{113} :\equiv \forall n\,(n \in B_{22} \leftrightarrow n \in \omega \land (n \in \mathbb{N} \to m \leq n \to a_{n} \leq a_{m})) \to \forall n \in \omega\,(n \in \mathbb{N} \to m \leq n \to a_{n} \leq a_{m})$
\item $\varphi_{114} :\equiv \forall x\,(\mathrm{Ind}(x) \leftrightarrow \emptyset \in x \land \forall y \in x\,(S(y) \in x)) \to (\mathrm{Ind}(\omega) \leftrightarrow \emptyset \in \omega \land \forall y \in \omega\,(S(y) \in \omega))$
\item $\varphi_{115} :\equiv \forall w\,(S(n) = w \to a_{S(n)} < a_{m} \to a_{w} < a_{m}) \to S(n) = 1_{\omega} \to a_{S(n)} < a_{m} \to a_{1_{\omega}} < a_{m}$
\item $\varphi_{116} :\equiv \forall w\,(1_{\omega} = w \to a_{1_{\omega}} < a_{m} \to a_{w} < a_{m}) \to 1_{\omega} = S(n) \to a_{1_{\omega}} < a_{m} \to a_{S(n)} < a_{m}$
\item $\varphi_{117} :\equiv (\mathrm{MonoDec}(a) \to m \in \mathbb{N} \to n \in \mathbb{N} \to m \leq n \to a_{n} \leq a_{m}) \to \mathrm{MonoDec}(a) \land m \in \mathbb{N} \to n \in \mathbb{N} \to m \leq n \to a_{n} \leq a_{m}$
\item $\varphi_{118} :\equiv \forall n\,(n \in B_{22} \leftrightarrow n \in \omega \land (n \in \mathbb{N} \to m \leq n \to a_{n} \leq a_{m})) \to \varphi_{24}$
\item $\varphi_{119} :\equiv \forall n\,(n \in B_{22} \leftrightarrow n \in \omega \land (n \in \mathbb{N} \to m \leq n \to a_{n} \leq a_{m})) \to n \in B_{22} \to S(n) \in \mathbb{N} \to n = \emptyset \to S(n) = S(\emptyset)$
\item $\varphi_{120} :\equiv \forall n\,(n \in B_{22} \leftrightarrow n \in \omega \land (n \in \mathbb{N} \to m \leq n \to a_{n} \leq a_{m})) \to n \in B_{22} \to S(n) \in \mathbb{N} \to m \leq S(n) \to a_{S(n)} \leq a_{m}$
\item $\varphi_{121} :\equiv \forall n\,(n \in B_{22} \leftrightarrow n \in \omega \land (n \in \mathbb{N} \to m \leq n \to a_{n} \leq a_{m})) \to \varphi_{29}$
\item $\varphi_{122} :\equiv \forall n\,(n \in B_{22} \leftrightarrow n \in \omega \land (n \in \mathbb{N} \to m \leq n \to a_{n} \leq a_{m})) \to (n \in B_{22} \leftrightarrow n \in \omega \land (n \in \mathbb{N} \to m \leq n \to a_{n} \leq a_{m}))$
\item $\varphi_{123} :\equiv \forall n\,(n \in B_{22} \leftrightarrow n \in \omega \land (n \in \mathbb{N} \to m \leq n \to a_{n} \leq a_{m})) \to n \in B_{22} \to n \in \omega \land (n \in \mathbb{N} \to m \leq n \to a_{n} \leq a_{m})$
\item $\varphi_{124} :\equiv \forall n\,(n \in B_{22} \leftrightarrow n \in \omega \land (n \in \mathbb{N} \to m \leq n \to a_{n} \leq a_{m})) \to n \in B_{22} \to S(n) \in \mathbb{N} \to n \neq \emptyset \to m \leq S(n) \to n \in \omega$
\item $\varphi_{125} :\equiv \forall n\,(n \in B_{22} \leftrightarrow n \in \omega \land (n \in \mathbb{N} \to m \leq n \to a_{n} \leq a_{m})) \to n \in B_{22} \to S(n) \in \mathbb{N} \to n \neq \emptyset \to m \leq S(n) \to m \in \omega$
\item $\varphi_{126} :\equiv \forall n\,(n \in B_{22} \leftrightarrow n \in \omega \land (n \in \mathbb{N} \to m \leq n \to a_{n} \leq a_{m})) \to \varphi_{36}$
\item $\varphi_{127} :\equiv \forall n\,(n \in B_{22} \leftrightarrow n \in \omega \land (n \in \mathbb{N} \to m \leq n \to a_{n} \leq a_{m})) \to \forall n\,(n \in B_{22} \leftrightarrow n \in \omega \land (n \in \mathbb{N} \to m \leq n \to a_{n} \leq a_{m}))$
\item $\varphi_{128} :\equiv \forall n\,(n \in B_{22} \leftrightarrow n \in \omega \land (n \in \mathbb{N} \to m \leq n \to a_{n} \leq a_{m})) \to n \in B_{22} \to S(n) \in \mathbb{N} \to n \neq \emptyset \to m \leq n \to a_{n} \leq a_{m}$
\item $\varphi_{129} :\equiv \forall n\,(n \in B_{22} \leftrightarrow n \in \omega \land (n \in \mathbb{N} \to m \leq n \to a_{n} \leq a_{m})) \to n \in B_{22} \to \varphi_{14}$
\item $\varphi_{130} :\equiv \forall n\,(n \in B_{22} \leftrightarrow n \in \omega \land (n \in \mathbb{N} \to m \leq n \to a_{n} \leq a_{m})) \to \varphi_{49}$
\item $\varphi_{131} :\equiv \forall n\,(n \in B_{22} \leftrightarrow n \in \omega \land (n \in \mathbb{N} \to m \leq n \to a_{n} \leq a_{m})) \to \varphi_{51}$
\item $\varphi_{132} :\equiv \forall n\,(n \in B_{22} \leftrightarrow n \in \omega \land (n \in \mathbb{N} \to m \leq n \to a_{n} \leq a_{m})) \to \varphi_{52}$
\item $\varphi_{133} :\equiv \varphi_{62} \to \mathrm{MonoDec}(a) \to m \in \mathbb{N} \to n \in \mathbb{N} \to m \leq n \to a_{n} \leq a_{m}$
\item $\varphi_{134} :\equiv \forall n\,(n \in B_{22} \leftrightarrow n \in \omega \land (n \in \mathbb{N} \to m \leq n \to a_{n} \leq a_{m})) \to \varphi_{56}$
\item $\varphi_{135} :\equiv \forall n\,(n \in B_{22} \leftrightarrow n \in \omega \land (n \in \mathbb{N} \to m \leq n \to a_{n} \leq a_{m})) \to \varphi_{64}$
\item $\varphi_{136} :\equiv \forall n\,(n \in B_{22} \leftrightarrow n \in \omega \land (n \in \mathbb{N} \to m \leq n \to a_{n} \leq a_{m})) \to n \in B_{22} \to \varphi_{40}$
\item $\varphi_{137} :\equiv \forall n\,(n \in B_{22} \leftrightarrow n \in \omega \land (n \in \mathbb{N} \to m \leq n \to a_{n} \leq a_{m})) \to \varphi_{63}$
\item $\varphi_{138} :\equiv \forall n\,(n \in B_{22} \leftrightarrow n \in \omega \land (n \in \mathbb{N} \to m \leq n \to a_{n} \leq a_{m})) \to n \in B_{22} \to \varphi_{46}$
\item $\varphi_{139} :\equiv \forall n\,(n \in B_{22} \leftrightarrow n \in \omega \land (n \in \mathbb{N} \to m \leq n \to a_{n} \leq a_{m})) \to n \in B_{22} \to \varphi_{50}$
\item $\varphi_{140} :\equiv \forall n\,(n \in B_{22} \leftrightarrow n \in \omega \land (n \in \mathbb{N} \to m \leq n \to a_{n} \leq a_{m})) \to n \in B_{22} \to \varphi_{57}$
\item $\varphi_{141} :\equiv (m \leq 1_{\omega} \to 1_{\omega} \leq m \to m = 1_{\omega}) \to \varphi_{4} \to \varphi_{3} \to \varphi_{5}$
\item $\varphi_{142} :\equiv (n \in \omega \leftrightarrow \forall x\,(\mathrm{Ind}(x) \to n \in x)) \to \mathrm{MonoDec}(a) \to m \in \mathbb{N} \to \varphi_{86}$
\item $\varphi_{143} :\equiv \forall n\,(n \in B_{22} \leftrightarrow n \in \omega \land (n \in \mathbb{N} \to m \leq n \to a_{n} \leq a_{m})) \to \varphi_{81}$
\item $\varphi_{144} :\equiv (m = 1_{\omega} \to a_{m} < a_{m} \to a_{1_{\omega}} < a_{m}) \to (m = 1_{\omega} \to a_{m} = a_{m} \to a_{1_{\omega}} = a_{m}) \to \varphi_{21}$
\item $\varphi_{145} :\equiv (S(n) = 1_{\omega} \to 1_{\omega} = S(n)) \to \mathrm{MonoDec}(a) \to m \in \mathbb{N} \to \varphi_{106}$
\item $\varphi_{146} :\equiv (1_{\omega} = S(n) \to S(n) = 1_{\omega}) \to \mathrm{MonoDec}(a) \to m \in \mathbb{N} \to \varphi_{105}$
\item $\varphi_{147} :\equiv (\forall x\,(\mathrm{Ind}(x) \to n \in x) \to \mathrm{Ind}(B_{22}) \to n \in B_{22}) \to \mathrm{MonoDec}(a) \to m \in \mathbb{N} \to \varphi_{88}$
\item $\varphi_{148} :\equiv (\mathrm{MonoDec}(a) \to m \in \mathbb{N} \to m \in \mathbb{N}) \to \mathrm{MonoDec}(a) \to m \in \mathbb{N} \to \varphi_{125}$
\item $\varphi_{149} :\equiv (n = \emptyset \to S(n) = S(n) \to S(n) = S(\emptyset)) \to \mathrm{MonoDec}(a) \to m \in \mathbb{N} \to \varphi_{119}$
\item $\varphi_{150} :\equiv (\forall y \in \omega\,(S(y) \in \omega) \to n \in \omega \to S(n) \in \omega) \to \mathrm{MonoDec}(a) \to m \in \mathbb{N} \to \varphi_{93}$
\item $\varphi_{151} :\equiv (\mathrm{Ind}(\omega) \leftrightarrow \emptyset \in \omega \land \forall y \in \omega\,(S(y) \in \omega)) \to \mathrm{MonoDec}(a) \to m \in \mathbb{N} \to \varphi_{112}$
\item $\varphi_{152} :\equiv (\mathrm{MonoDec}(a) \to m \in \mathbb{N} \to m \leq 1_{\omega} \to a_{m} \leq a_{m}) \to \varphi_{144}$
\item $\varphi_{153} :\equiv (m = S(n) \to a_{m} = a_{m} \to a_{S(n)} = a_{m}) \to \mathrm{MonoDec}(a) \to m \in \mathbb{N} \to \varphi_{138}$
\item $\varphi_{154} :\equiv (\mathrm{MonoDec}(a) \to m \in \mathbb{N} \to \varphi_{76}) \to \varphi_{73} \to n \in B_{22} \to S(n) \in B_{22}$
\item $\varphi_{155} :\equiv (\mathrm{MonoDec}(a) \to m \in \mathbb{N} \to \varphi_{75}) \to \mathrm{MonoDec}(a) \to m \in \mathbb{N} \to \varphi_{78}$
\item $\varphi_{156} :\equiv \forall n\,(\varphi_{73} \to n \in B_{22} \to S(n) \in B_{22}) \to \forall n\,(\varphi_{73}) \to \forall n \in B_{22}\,(S(n) \in B_{22})$
\item $\varphi_{157} :\equiv (\forall n \in \mathbb{N}\,(a_{S(n)} \leq a_{n}) \to n \in \mathbb{N} \to a_{S(n)} \leq a_{n}) \to \mathrm{MonoDec}(a) \to m \in \mathbb{N} \to \varphi_{139}$
\item $\varphi_{158} :\equiv (\mathrm{MonoDec}(a) \to \mathrm{MonoDec}(a)) \to \varphi_{27} \to \mathrm{MonoDec}(a) \to m \in \mathbb{N} \to \varphi_{140}$
\item $\varphi_{159} :\equiv (\mathrm{MonoDec}(a) \to m \in \mathbb{N} \to \varphi_{120}) \to \mathrm{MonoDec}(a) \to m \in \mathbb{N} \to \varphi_{76}$
\item $\varphi_{160} :\equiv (\mathrm{MonoDec}(a) \to m \in \mathbb{N} \to \varphi_{72}) \to \mathrm{MonoDec}(a) \to m \in \mathbb{N} \to \varphi_{123}$
\item $\varphi_{161} :\equiv (\mathrm{MonoDec}(a) \to m \in \mathbb{N} \to \varphi_{123}) \to \mathrm{MonoDec}(a) \to m \in \mathbb{N} \to \varphi_{75}$
\item $\varphi_{162} :\equiv (\mathrm{MonoDec}(a) \to m \in \mathbb{N} \to \varphi_{112}) \to \mathrm{MonoDec}(a) \to m \in \mathbb{N} \to \varphi_{101}$
\item $\varphi_{163} :\equiv (\mathrm{MonoDec}(a) \to m \in \mathbb{N} \to \varphi_{107}) \to \varphi_{73} \to n \in \omega \to n \in \mathbb{N} \to m \leq n \to a_{n} \leq a_{m}$
\item $\varphi_{164} :\equiv (\mathrm{MonoDec}(a) \to m \in \mathbb{N} \to \varphi_{111}) \to \mathrm{MonoDec}(a) \to m \in \mathbb{N} \to \varphi_{108}$
\item $\varphi_{165} :\equiv \forall n\,(\varphi_{73} \to n \in \omega \to n \in \mathbb{N} \to m \leq n \to a_{n} \leq a_{m}) \to \forall n\,(\varphi_{73}) \to \forall n \in \omega\,(n \in \mathbb{N} \to m \leq n \to a_{n} \leq a_{m})$
\item $\varphi_{166} :\equiv (\mathrm{MonoDec}(a) \to m \in \mathbb{N} \to \varphi_{123}) \to \mathrm{MonoDec}(a) \to m \in \mathbb{N} \to \varphi_{102}$
\item $\varphi_{167} :\equiv (\mathrm{MonoDec}(a) \to m \in \mathbb{N} \to \varphi_{122}) \to \mathrm{MonoDec}(a) \to m \in \mathbb{N} \to \varphi_{107}$
\item $\varphi_{168} :\equiv (\mathrm{MonoDec}(a) \to m \in \mathbb{N} \to \varphi_{108}) \to \mathrm{MonoDec}(a) \to m \in \mathbb{N} \to \varphi_{124}$
\item $\varphi_{169} :\equiv (\mathrm{MonoDec}(a) \to m \in \mathbb{N} \to \varphi_{108}) \to \mathrm{MonoDec}(a) \to m \in \mathbb{N} \to \varphi_{128}$
\item $\varphi_{170} :\equiv (\mathrm{MonoDec}(a) \to m \in \mathbb{N} \to \varphi_{118}) \to \mathrm{MonoDec}(a) \to m \in \mathbb{N} \to \varphi_{126}$
\item $\varphi_{171} :\equiv (\mathrm{MonoDec}(a) \to m \in \mathbb{N} \to \varphi_{105}) \to \mathrm{MonoDec}(a) \to m \in \mathbb{N} \to \varphi_{130}$
\item $\varphi_{172} :\equiv (\mathrm{MonoDec}(a) \to m \in \mathbb{N} \to \varphi_{108}) \to \mathrm{MonoDec}(a) \to m \in \mathbb{N} \to \varphi_{131}$
\item $\varphi_{173} :\equiv (\mathrm{MonoDec}(a) \to m \in \mathbb{N} \to \varphi_{132}) \to \mathrm{MonoDec}(a) \to m \in \mathbb{N} \to \varphi_{120}$
\item $\varphi_{174} :\equiv (\mathrm{MonoDec}(a) \to m \in \mathbb{N} \to \varphi_{121}) \to \mathrm{MonoDec}(a) \to m \in \mathbb{N} \to \varphi_{134}$
\item $\varphi_{175} :\equiv (\mathrm{MonoDec}(a) \to m \in \mathbb{N} \to \varphi_{134}) \to \mathrm{MonoDec}(a) \to m \in \mathbb{N} \to \varphi_{137}$
\item $\varphi_{176} :\equiv (\mathrm{MonoDec}(a) \to m \in \mathbb{N} \to \varphi_{138}) \to \mathrm{MonoDec}(a) \to m \in \mathbb{N} \to \varphi_{132}$
\item $\varphi_{177} :\equiv (\mathrm{MonoDec}(a) \to m \in \mathbb{N} \to \varphi_{112}) \to \varphi_{68} \to \mathrm{MonoDec}(a) \to m \in \mathbb{N} \to \varphi_{70}$
\item $\varphi_{178} :\equiv (\mathrm{MonoDec}(a) \to m \in \mathbb{N} \to \varphi_{79}) \to \varphi_{26} \to \varphi_{25} \to \mathrm{MonoDec}(a) \to m \in \mathbb{N} \to \varphi_{61}$
\item $\varphi_{179} :\equiv (\mathrm{MonoDec}(a) \to m \in \mathbb{N} \to \varphi_{109}) \to S(n) = S(n) \to \varphi_{149}$
\item $\varphi_{180} :\equiv (\mathrm{MonoDec}(a) \to m \in \mathbb{N} \to \varphi_{119}) \to S(\emptyset) = 1_{\omega} \to \varphi_{10} \to \mathrm{MonoDec}(a) \to m \in \mathbb{N} \to \varphi_{105}$
\item $\varphi_{181} :\equiv (S(n) = 1_{\omega} \to a_{S(n)} = a_{m} \to a_{1_{\omega}} = a_{m}) \to \varphi_{171}$
\item $\varphi_{182} :\equiv (S(n) = 1_{\omega} \to a_{S(n)} < a_{m} \to a_{1_{\omega}} < a_{m}) \to \varphi_{181}$
\item $\varphi_{183} :\equiv (\mathrm{MonoDec}(a) \to m \in \mathbb{N} \to \varphi_{127}) \to \varphi_{122} \to \mathrm{MonoDec}(a) \to m \in \mathbb{N} \to \varphi_{122}$
\item $\varphi_{184} :\equiv (n \in \mathbb{N} \leftrightarrow n \in \omega \land n \neq \emptyset) \to (\mathrm{MonoDec}(a) \to m \in \mathbb{N} \to \varphi_{75}) \to \varphi_{164}$
\item $\varphi_{185} :\equiv (\mathrm{MonoDec}(a) \to m \in \mathbb{N} \to \varphi_{136}) \to (m = S(n) \to a_{m} < a_{m} \to a_{S(n)} < a_{m}) \to \varphi_{153}$
\item $\varphi_{186} :\equiv (\mathrm{MonoDec}(a) \to m \in \mathbb{N} \to \varphi_{127}) \to \varphi_{135} \to \mathrm{MonoDec}(a) \to m \in \mathbb{N} \to \varphi_{135}$
\item $\varphi_{187} :\equiv (S(n) = 1_{\omega} \to m < S(n) \to m < 1_{\omega}) \to (S(n) = 1_{\omega} \to m = S(n) \to m = 1_{\omega}) \to \varphi_{182}$
\item $\varphi_{188} :\equiv (\mathrm{MonoDec}(a) \to m \in \mathbb{N} \to \varphi_{127}) \to \varphi_{143} \to \mathrm{MonoDec}(a) \to m \in \mathbb{N} \to \varphi_{143}$
\item $\varphi_{189} :\equiv (\mathrm{MonoDec}(a) \to m \in \mathbb{N} \to \varphi_{88}) \to (\mathrm{MonoDec}(a) \to m \in \mathbb{N} \to \varphi_{61}) \to \varphi_{167}$
\item $\varphi_{190} :\equiv (1_{\omega} = S(n) \to a_{1_{\omega}} = a_{m} \to a_{S(n)} = a_{m}) \to \varphi_{187}$
\item $\varphi_{191} :\equiv (\forall B_{22}\,(\varphi_{66} \to \varphi_{41}) \to \forall B_{22}\,(\varphi_{66}) \to \forall B_{22}\,(\varphi_{41})) \to (\varphi_{66} \to \forall B_{22}\,(\varphi_{66})) \to \varphi_{65}$
\item $\varphi_{192} :\equiv (\mathrm{MonoDec}(a) \to m \in \mathbb{N} \to \varphi_{135}) \to (\mathrm{MonoDec}(a) \to m \in \mathbb{N} \to \varphi_{78}) \to \varphi_{159}$
\item $\varphi_{193} :\equiv (a_{S(n)} \leq a_{n} \to a_{n} \leq a_{m} \to a_{S(n)} \leq a_{m}) \to (\mathrm{MonoDec}(a) \to m \in \mathbb{N} \to \varphi_{131}) \to \varphi_{175}$
\item $\varphi_{194} :\equiv (1_{\omega} = S(n) \to a_{1_{\omega}} < a_{m} \to a_{S(n)} < a_{m}) \to \varphi_{190}$
\item $\varphi_{195} :\equiv (\mathrm{MonoDec}(a) \to m \in \mathbb{N} \to \varphi_{125}) \to (\mathrm{MonoDec}(a) \to m \in \mathbb{N} \to \varphi_{124}) \to \varphi_{170}$
\item $\varphi_{196} :\equiv (\mathrm{MonoDec}(a) \to m \in \mathbb{N} \to \varphi_{126}) \to (\mathrm{MonoDec}(a) \to m \in \mathbb{N} \to \varphi_{137}) \to \varphi_{176}$
\item $\varphi_{197} :\equiv \forall n\,(\varphi_{73} \to n \in B_{22} \to S(n) \in B_{22}) \to \varphi_{156} \to (\varphi_{73} \to \forall n\,(\varphi_{73})) \to \mathrm{MonoDec}(a) \to m \in \mathbb{N} \to \varphi_{79}$
\item $\varphi_{198} :\equiv (1_{\omega} = S(n) \to m < 1_{\omega} \to m < S(n)) \to (1_{\omega} = S(n) \to m = 1_{\omega} \to m = S(n)) \to \varphi_{194}$
\item $\varphi_{199} :\equiv \forall n\,(\varphi_{73} \to n \in \omega \to n \in \mathbb{N} \to m \leq n \to a_{n} \leq a_{m}) \to \varphi_{165} \to (\varphi_{73} \to \forall n\,(\varphi_{73})) \to \mathrm{MonoDec}(a) \to m \in \mathbb{N} \to \varphi_{113}$
\item $\varphi_{200} :\equiv (\mathrm{MonoDec}(a) \to m \in \mathbb{N} \to \varphi_{48}) \to \forall B_{22}\,(\varphi_{66} \to \varphi_{41}) \to \varphi_{191}$
\end{abbrevlist}
\end{fullproof}

% 自動生成 (prover/build.js): SubLtSwap  ノード数 224
\begin{fullproof}{[E31]}{$\forall x\,\forall e\,\forall y\,(x \in \mathbb{R} \to e \in \mathbb{R} \to y \in \mathbb{R} \to x - e < y \to -e < y - x)$}{224}\phantomsection\label{tree:SubLtSwap}
\noindent$\Pi_{1}$:
\begin{ptree}
\AXC{$z = w \to z < y + -x \to w < y + -x$\lbl{A6}}
\UIC{$\forall w\,(z = w \to z < y + -x \to w < y + -x)$}
\UIC{$\forall z\,\forall w\,(z = w \to z < y + -x \to w < y + -x)$}
\AXC{$\forall z\,\forall w\,(z = w \to z < y + -x \to w < y + -x) \to \varphi_{11}$\lbl{A2}}
\BIC{$\varphi_{11}$}
\AXC{$\varphi_{57}$\lbl{A2}}
\BIC{$x - e + -x = -e \to x - e + -x < y + -x \to -e < y + -x$}
\end{ptree}
\noindent$\Pi_{2}$:
\begin{ptree}
\AXC{$\forall x \in \mathbb{R}\,(-x \in \mathbb{R})$\lbl{(R1)}}
\AXC{$\forall x \in \mathbb{R}\,(-x \in \mathbb{R}) \to x \in \mathbb{R} \to -x \in \mathbb{R}$\lbl{A2}}
\BIC{$x \in \mathbb{R} \to -x \in \mathbb{R}$}
\AXC{$(x \in \mathbb{R} \to -x \in \mathbb{R}) \to \varphi_{50}$\lbl{A1}}
\BIC{$\varphi_{50}$}
\end{ptree}
\noindent$\Pi_{3}$:
\begin{ptree}
\AXC{$x \in \mathbb{R} \to e \in \mathbb{R} \to y \in \mathbb{R} \to x - e < y \to x \in \mathbb{R}$\lbl{A1}}
\AXC{$\stackrel{\Pi_{2}}{\vdots}$}\noLine
\UIC{$\varphi_{50}$}
\BIC{$x \in \mathbb{R} \to e \in \mathbb{R} \to y \in \mathbb{R} \to x - e < y \to -x \in \mathbb{R}$}
\end{ptree}
\noindent$\Pi_{4}$:
\begin{ptree}
\AXC{$\forall x\,\forall y\,(x \in \mathbb{R} \to y \in \mathbb{R} \to x - y \in \mathbb{R})$\lbl{[F1]}}
\AXC{$\varphi_{39}$\lbl{A2}}
\BIC{$\forall y\,(x \in \mathbb{R} \to y \in \mathbb{R} \to x - y \in \mathbb{R})$}
\AXC{$\varphi_{29}$\lbl{A2}}
\BIC{$x \in \mathbb{R} \to e \in \mathbb{R} \to x - e \in \mathbb{R}$}
\AXC{$\varphi_{64}$\lbl{A1}}
\BIC{$(x \in \mathbb{R} \to e \in \mathbb{R} \to y \in \mathbb{R} \to x - e < y \to x \in \mathbb{R}) \to \varphi_{51}$}
\end{ptree}
\noindent$\Pi_{5}$:
\begin{ptree}
\AXC{$x \in \mathbb{R} \to e \in \mathbb{R} \to y \in \mathbb{R} \to x - e < y \to x \in \mathbb{R}$\lbl{A1}}
\AXC{$\stackrel{\Pi_{4}}{\vdots}$}\noLine
\UIC{$(x \in \mathbb{R} \to e \in \mathbb{R} \to y \in \mathbb{R} \to x - e < y \to x \in \mathbb{R}) \to \varphi_{51}$}
\BIC{$\varphi_{51}$}
\end{ptree}
\noindent$\Pi_{6}$:
\begin{ptree}
\AXC{$\varphi_{9}$\lbl{(R14)}}
\AXC{$\varphi_{9} \to \forall y\,(\varphi_{22})$\lbl{A2}}
\BIC{$\forall y\,(\varphi_{22})$}
\AXC{$\forall y\,(\varphi_{22}) \to \varphi_{22}$\lbl{A2}}
\BIC{$\varphi_{22}$}
\AXC{$\varphi_{22} \to \varphi_{40}$\lbl{A2}}
\BIC{$\varphi_{40}$}
\AXC{$\varphi_{40} \to \varphi_{72}$\lbl{A1}}
\BIC{$\varphi_{72}$}
\end{ptree}
\noindent$\Pi_{7}$:
\begin{ptree}
\AXC{$x \in \mathbb{R} \to e \in \mathbb{R} \to y \in \mathbb{R} \to x - e < y \to y \in \mathbb{R}$\lbl{A1}}
\AXC{$x \in \mathbb{R} \to e \in \mathbb{R} \to y \in \mathbb{R} \to x - e < y \to e \in \mathbb{R}$\lbl{A1}}
\AXC{$\stackrel{\Pi_{5}}{\vdots}$}\noLine
\UIC{$\varphi_{51}$}
\BIC{$x \in \mathbb{R} \to e \in \mathbb{R} \to y \in \mathbb{R} \to x - e < y \to x - e \in \mathbb{R}$}
\AXC{$\stackrel{\Pi_{6}}{\vdots}$}\noLine
\UIC{$\varphi_{72}$}
\BIC{$(x \in \mathbb{R} \to e \in \mathbb{R} \to y \in \mathbb{R} \to x - e < y \to y \in \mathbb{R}) \to \varphi_{61}$}
\BIC{$\varphi_{61}$}
\end{ptree}
\noindent$\Pi_{8}$:
\begin{ptree}
\AXC{$x \in \mathbb{R} \to e \in \mathbb{R} \to y \in \mathbb{R} \to x - e < y \to x - e < y$\lbl{A1}}
\AXC{$\stackrel{\Pi_{3}}{\vdots}$}\noLine
\UIC{$x \in \mathbb{R} \to e \in \mathbb{R} \to y \in \mathbb{R} \to x - e < y \to -x \in \mathbb{R}$}
\AXC{$\stackrel{\Pi_{7}}{\vdots}$}\noLine
\UIC{$\varphi_{61}$}
\BIC{$(x \in \mathbb{R} \to e \in \mathbb{R} \to y \in \mathbb{R} \to x - e < y \to x - e < y) \to \varphi_{26}$}
\BIC{$\varphi_{26}$}
\end{ptree}
\noindent$\Pi_{9}$:
\begin{ptree}
\AXC{$y = z \to x = y \to x = z$\lbl{A6}}
\AXC{$(y = z \to x = y \to x = z) \to x = y \to y = z \to x = z$\lbl{A1}}
\BIC{$x = y \to y = z \to x = z$}
\UIC{$\forall z\,(x = y \to y = z \to x = z)$}
\UIC{$\forall y\,\forall z\,(x = y \to y = z \to x = z)$}
\UIC{$\forall x\,\forall y\,\forall z\,(x = y \to y = z \to x = z)$}
\AXC{$\varphi_{14}$\lbl{A2}}
\BIC{$\forall y\,\forall z\,(x - e + -x = y \to y = z \to x - e + -x = z)$}
\end{ptree}
\noindent$\Pi_{10}$:
\begin{ptree}
\AXC{$\stackrel{\Pi_{9}}{\vdots}$}\noLine
\UIC{$\forall y\,\forall z\,(x - e + -x = y \to y = z \to x - e + -x = z)$}
\AXC{$\forall y\,\forall z\,(x - e + -x = y \to y = z \to x - e + -x = z) \to \varphi_{5}$\lbl{A2}}
\BIC{$\varphi_{5}$}
\AXC{$\varphi_{54}$\lbl{A2}}
\BIC{$x - e + -x = -e + 0 \to -e + 0 = -e \to x - e + -x = -e$}
\end{ptree}
\noindent$\Pi_{11}$:
\begin{ptree}
\AXC{$\forall x \in \mathbb{R}\,(-x \in \mathbb{R})$\lbl{(R1)}}
\AXC{$\forall x \in \mathbb{R}\,(-x \in \mathbb{R}) \to e \in \mathbb{R} \to -e \in \mathbb{R}$\lbl{A2}}
\BIC{$e \in \mathbb{R} \to -e \in \mathbb{R}$}
\AXC{$(e \in \mathbb{R} \to -e \in \mathbb{R}) \to \varphi_{52}$\lbl{A1}}
\BIC{$\varphi_{52}$}
\end{ptree}
\noindent$\Pi_{12}$:
\begin{ptree}
\AXC{$x \in \mathbb{R} \to e \in \mathbb{R} \to y \in \mathbb{R} \to x - e < y \to e \in \mathbb{R}$\lbl{A1}}
\AXC{$\stackrel{\Pi_{11}}{\vdots}$}\noLine
\UIC{$\varphi_{52}$}
\BIC{$x \in \mathbb{R} \to e \in \mathbb{R} \to y \in \mathbb{R} \to x - e < y \to -e \in \mathbb{R}$}
\end{ptree}
\noindent$\Pi_{13}$:
\begin{ptree}
\AXC{$\forall x \in \mathbb{R}\,(x + 0 = x)$\lbl{(R4)}}
\AXC{$\forall x \in \mathbb{R}\,(x + 0 = x) \to -e \in \mathbb{R} \to -e + 0 = -e$\lbl{A2}}
\BIC{$-e \in \mathbb{R} \to -e + 0 = -e$}
\AXC{$(-e \in \mathbb{R} \to -e + 0 = -e) \to \varphi_{53}$\lbl{A1}}
\BIC{$\varphi_{53}$}
\end{ptree}
\noindent$\Pi_{14}$:
\begin{ptree}
\AXC{$y = z \to x = y \to x = z$\lbl{A6}}
\AXC{$(y = z \to x = y \to x = z) \to x = y \to y = z \to x = z$\lbl{A1}}
\BIC{$x = y \to y = z \to x = z$}
\UIC{$\forall z\,(x = y \to y = z \to x = z)$}
\UIC{$\forall y\,\forall z\,(x = y \to y = z \to x = z)$}
\UIC{$\forall x\,\forall y\,\forall z\,(x = y \to y = z \to x = z)$}
\AXC{$\varphi_{14}$\lbl{A2}}
\BIC{$\forall y\,\forall z\,(x - e + -x = y \to y = z \to x - e + -x = z)$}
\end{ptree}
\noindent$\Pi_{15}$:
\begin{ptree}
\AXC{$\stackrel{\Pi_{14}}{\vdots}$}\noLine
\UIC{$\forall y\,\forall z\,(x - e + -x = y \to y = z \to x - e + -x = z)$}
\AXC{$\forall y\,\forall z\,(x - e + -x = y \to y = z \to x - e + -x = z) \to \varphi_{30}$\lbl{A2}}
\BIC{$\varphi_{30}$}
\AXC{$\varphi_{30} \to \varphi_{45}$\lbl{A2}}
\BIC{$\varphi_{45}$}
\end{ptree}
\noindent$\Pi_{16}$:
\begin{ptree}
\AXC{$\forall x \in \mathbb{R}\,(x + -x = 0)$\lbl{(R5)}}
\AXC{$\forall x \in \mathbb{R}\,(x + -x = 0) \to x \in \mathbb{R} \to x + -x = 0$\lbl{A2}}
\BIC{$x \in \mathbb{R} \to x + -x = 0$}
\AXC{$(x \in \mathbb{R} \to x + -x = 0) \to \varphi_{48}$\lbl{A1}}
\BIC{$\varphi_{48}$}
\end{ptree}
\noindent$\Pi_{17}$:
\begin{ptree}
\AXC{$x \in \mathbb{R} \to e \in \mathbb{R} \to y \in \mathbb{R} \to x - e < y \to x \in \mathbb{R}$\lbl{A1}}
\AXC{$\stackrel{\Pi_{16}}{\vdots}$}\noLine
\UIC{$\varphi_{48}$}
\BIC{$x \in \mathbb{R} \to e \in \mathbb{R} \to y \in \mathbb{R} \to x - e < y \to x + -x = 0$}
\AXC{$\varphi_{67}$\lbl{A1}}
\BIC{$\varphi_{63}$}
\end{ptree}
\noindent$\Pi_{18}$:
\begin{ptree}
\AXC{$x = x$\lbl{A5}}
\UIC{$\forall x\,(x = x)$}
\AXC{$\forall x\,(x = x) \to -e + (x + -x) = -e + (x + -x)$\lbl{A2}}
\BIC{$-e + (x + -x) = -e + (x + -x)$}
\AXC{$\stackrel{\Pi_{17}}{\vdots}$}\noLine
\UIC{$\varphi_{63}$}
\BIC{$(x + -x = 0 \to \varphi_{18}) \to \varphi_{28}$}
\end{ptree}
\noindent$\Pi_{19}$:
\begin{ptree}
\AXC{$\varphi_{12}$\lbl{A6}}
\UIC{$\forall w\,(\varphi_{12})$}
\UIC{$\forall z\,\forall w\,(\varphi_{12})$}
\AXC{$\forall z\,\forall w\,(\varphi_{12}) \to \varphi_{42}$\lbl{A2}}
\BIC{$\varphi_{42}$}
\AXC{$\varphi_{42} \to x + -x = 0 \to \varphi_{18}$\lbl{A2}}
\BIC{$x + -x = 0 \to \varphi_{18}$}
\AXC{$\stackrel{\Pi_{18}}{\vdots}$}\noLine
\UIC{$(x + -x = 0 \to \varphi_{18}) \to \varphi_{28}$}
\BIC{$\varphi_{28}$}
\end{ptree}
\noindent$\Pi_{20}$:
\begin{ptree}
\AXC{$y = z \to x = y \to x = z$\lbl{A6}}
\AXC{$(y = z \to x = y \to x = z) \to x = y \to y = z \to x = z$\lbl{A1}}
\BIC{$x = y \to y = z \to x = z$}
\UIC{$\forall z\,(x = y \to y = z \to x = z)$}
\UIC{$\forall y\,\forall z\,(x = y \to y = z \to x = z)$}
\UIC{$\forall x\,\forall y\,\forall z\,(x = y \to y = z \to x = z)$}
\AXC{$\varphi_{14}$\lbl{A2}}
\BIC{$\forall y\,\forall z\,(x - e + -x = y \to y = z \to x - e + -x = z)$}
\end{ptree}
\noindent$\Pi_{21}$:
\begin{ptree}
\AXC{$\stackrel{\Pi_{20}}{\vdots}$}\noLine
\UIC{$\forall y\,\forall z\,(x - e + -x = y \to y = z \to x - e + -x = z)$}
\AXC{$\forall y\,\forall z\,(x - e + -x = y \to y = z \to x - e + -x = z) \to \varphi_{31}$\lbl{A2}}
\BIC{$\varphi_{31}$}
\end{ptree}
\noindent$\Pi_{22}$:
\begin{ptree}
\AXC{$\varphi_{7}$\lbl{(R2)}}
\AXC{$\varphi_{7} \to \varphi_{23}$\lbl{A2}}
\BIC{$\varphi_{23}$}
\AXC{$\varphi_{23} \to \varphi_{19}$\lbl{A2}}
\BIC{$\varphi_{19}$}
\AXC{$\varphi_{19} \to \varphi_{37}$\lbl{A2}}
\BIC{$\varphi_{37}$}
\AXC{$\varphi_{37} \to \varphi_{68}$\lbl{A1}}
\BIC{$\varphi_{68}$}
\end{ptree}
\noindent$\Pi_{23}$:
\begin{ptree}
\AXC{$x \in \mathbb{R} \to e \in \mathbb{R} \to y \in \mathbb{R} \to x - e < y \to x \in \mathbb{R}$\lbl{A1}}
\AXC{$\stackrel{\Pi_{12}}{\vdots}$}\noLine
\UIC{$x \in \mathbb{R} \to e \in \mathbb{R} \to y \in \mathbb{R} \to x - e < y \to -e \in \mathbb{R}$}
\AXC{$\stackrel{\Pi_{22}}{\vdots}$}\noLine
\UIC{$\varphi_{68}$}
\BIC{$(x \in \mathbb{R} \to e \in \mathbb{R} \to y \in \mathbb{R} \to x - e < y \to x \in \mathbb{R}) \to \varphi_{59}$}
\BIC{$\varphi_{59}$}
\end{ptree}
\noindent$\Pi_{24}$:
\begin{ptree}
\AXC{$y = z \to x = y \to x = z$\lbl{A6}}
\AXC{$(y = z \to x = y \to x = z) \to x = y \to y = z \to x = z$\lbl{A1}}
\BIC{$x = y \to y = z \to x = z$}
\UIC{$\forall z\,(x = y \to y = z \to x = z)$}
\UIC{$\forall y\,\forall z\,(x = y \to y = z \to x = z)$}
\UIC{$\forall x\,\forall y\,\forall z\,(x = y \to y = z \to x = z)$}
\AXC{$\varphi_{14}$\lbl{A2}}
\BIC{$\forall y\,\forall z\,(x - e + -x = y \to y = z \to x - e + -x = z)$}
\end{ptree}
\noindent$\Pi_{25}$:
\begin{ptree}
\AXC{$\stackrel{\Pi_{24}}{\vdots}$}\noLine
\UIC{$\forall y\,\forall z\,(x - e + -x = y \to y = z \to x - e + -x = z)$}
\AXC{$\forall y\,\forall z\,(x - e + -x = y \to y = z \to x - e + -x = z) \to \varphi_{32}$\lbl{A2}}
\BIC{$\varphi_{32}$}
\AXC{$\varphi_{32} \to \varphi_{47}$\lbl{A2}}
\BIC{$\varphi_{47}$}
\end{ptree}
\noindent$\Pi_{26}$:
\begin{ptree}
\AXC{$z = w \to x + -e + -x = z + -x \to x + -e + -x = w + -x$\lbl{A6}}
\UIC{$\varphi_{15}$}
\UIC{$\forall z\,(\varphi_{15})$}
\AXC{$\forall z\,(\varphi_{15}) \to \varphi_{43}$\lbl{A2}}
\BIC{$\varphi_{43}$}
\AXC{$\varphi_{43} \to x + -e = -e + x \to \varphi_{33}$\lbl{A2}}
\BIC{$x + -e = -e + x \to \varphi_{33}$}
\end{ptree}
\noindent$\Pi_{27}$:
\begin{ptree}
\AXC{$\forall x\,\forall y\,(x \in \mathbb{R} \to y \in \mathbb{R} \to x + y = y + x)$\lbl{(R3)}}
\AXC{$\varphi_{34}$\lbl{A2}}
\BIC{$\forall y\,(x \in \mathbb{R} \to y \in \mathbb{R} \to x + y = y + x)$}
\AXC{$\varphi_{44}$\lbl{A2}}
\BIC{$x \in \mathbb{R} \to -e \in \mathbb{R} \to x + -e = -e + x$}
\AXC{$\varphi_{66}$\lbl{A1}}
\BIC{$(x \in \mathbb{R} \to e \in \mathbb{R} \to y \in \mathbb{R} \to x - e < y \to x \in \mathbb{R}) \to \varphi_{55}$}
\end{ptree}
\noindent$\Pi_{28}$:
\begin{ptree}
\AXC{$x \in \mathbb{R} \to e \in \mathbb{R} \to y \in \mathbb{R} \to x - e < y \to x \in \mathbb{R}$\lbl{A1}}
\AXC{$\stackrel{\Pi_{27}}{\vdots}$}\noLine
\UIC{$(x \in \mathbb{R} \to e \in \mathbb{R} \to y \in \mathbb{R} \to x - e < y \to x \in \mathbb{R}) \to \varphi_{55}$}
\BIC{$\varphi_{55}$}
\end{ptree}
\noindent$\Pi_{29}$:
\begin{ptree}
\AXC{$\stackrel{\Pi_{12}}{\vdots}$}\noLine
\UIC{$x \in \mathbb{R} \to e \in \mathbb{R} \to y \in \mathbb{R} \to x - e < y \to -e \in \mathbb{R}$}
\AXC{$\stackrel{\Pi_{28}}{\vdots}$}\noLine
\UIC{$\varphi_{55}$}
\BIC{$x \in \mathbb{R} \to e \in \mathbb{R} \to y \in \mathbb{R} \to x - e < y \to x + -e = -e + x$}
\AXC{$\varphi_{71}$\lbl{A1}}
\BIC{$\varphi_{65}$}
\end{ptree}
\noindent$\Pi_{30}$:
\begin{ptree}
\AXC{$\stackrel{\Pi_{26}}{\vdots}$}\noLine
\UIC{$x + -e = -e + x \to \varphi_{33}$}
\AXC{$x = x$\lbl{A5}}
\UIC{$\forall x\,(x = x)$}
\AXC{$\forall x\,(x = x) \to x + -e + -x = x + -e + -x$\lbl{A2}}
\BIC{$x + -e + -x = x + -e + -x$}
\AXC{$\stackrel{\Pi_{29}}{\vdots}$}\noLine
\UIC{$\varphi_{65}$}
\BIC{$(x + -e = -e + x \to \varphi_{33}) \to x \in \mathbb{R} \to \varphi_{25}$}
\BIC{$x \in \mathbb{R} \to \varphi_{25}$}
\end{ptree}
\noindent$\Pi_{31}$:
\begin{ptree}
\AXC{$\forall x\,\forall y\,(x - y = x + -y)$\lbl{Def$-$}}
\AXC{$\forall x\,\forall y\,(x - y = x + -y) \to \forall y\,(x - y = x + -y)$\lbl{A2}}
\BIC{$\forall y\,(x - y = x + -y)$}
\end{ptree}
\noindent$\Pi_{32}$:
\begin{ptree}
\AXC{$\stackrel{\Pi_{31}}{\vdots}$}\noLine
\UIC{$\forall y\,(x - y = x + -y)$}
\AXC{$\forall y\,(x - y = x + -y) \to x - e = x + -e$\lbl{A2}}
\BIC{$x - e = x + -e$}
\AXC{$x - e = x + -e \to x - e + -x = x - e + -x \to \varphi_{41} \to \varphi_{35}$\lbl{A1}}
\BIC{$x - e + -x = x - e + -x \to \varphi_{41} \to \varphi_{35}$}
\end{ptree}
\noindent$\Pi_{33}$:
\begin{ptree}
\AXC{$x = x$\lbl{A5}}
\UIC{$\forall x\,(x = x)$}
\AXC{$\forall x\,(x = x) \to x - e + -x = x - e + -x$\lbl{A2}}
\BIC{$x - e + -x = x - e + -x$}
\AXC{$\stackrel{\Pi_{32}}{\vdots}$}\noLine
\UIC{$x - e + -x = x - e + -x \to \varphi_{41} \to \varphi_{35}$}
\BIC{$\varphi_{41} \to \varphi_{35}$}
\end{ptree}
\noindent$\Pi_{34}$:
\begin{ptree}
\AXC{$z = w \to x - e + -x = z + -x \to x - e + -x = w + -x$\lbl{A6}}
\UIC{$\varphi_{6}$}
\UIC{$\forall z\,(\varphi_{6})$}
\AXC{$\forall z\,(\varphi_{6}) \to \varphi_{16}$\lbl{A2}}
\BIC{$\varphi_{16}$}
\AXC{$\varphi_{16} \to \varphi_{41}$\lbl{A2}}
\BIC{$\varphi_{41}$}
\AXC{$\stackrel{\Pi_{33}}{\vdots}$}\noLine
\UIC{$\varphi_{41} \to \varphi_{35}$}
\BIC{$\varphi_{35}$}
\end{ptree}
\noindent$\Pi_{35}$:
\begin{ptree}
\AXC{$\stackrel{\Pi_{25}}{\vdots}$}\noLine
\UIC{$\varphi_{47}$}
\AXC{$\stackrel{\Pi_{30}}{\vdots}$}\noLine
\UIC{$x \in \mathbb{R} \to \varphi_{25}$}
\AXC{$\stackrel{\Pi_{34}}{\vdots}$}\noLine
\UIC{$\varphi_{35}$}
\AXC{$\varphi_{35} \to (x \in \mathbb{R} \to \varphi_{25}) \to \varphi_{47} \to \varphi_{36}$\lbl{A1}}
\BIC{$(x \in \mathbb{R} \to \varphi_{25}) \to \varphi_{47} \to \varphi_{36}$}
\BIC{$\varphi_{47} \to \varphi_{36}$}
\BIC{$\varphi_{36}$}
\AXC{$\varphi_{36} \to \varphi_{69}$\lbl{A1}}
\BIC{$\varphi_{69}$}
\end{ptree}
\noindent$\Pi_{36}$:
\begin{ptree}
\AXC{$\stackrel{\Pi_{21}}{\vdots}$}\noLine
\UIC{$\varphi_{31}$}
\AXC{$\varphi_{31} \to x - e + -x = -e + x + -x \to \varphi_{27}$\lbl{A2}}
\BIC{$x - e + -x = -e + x + -x \to \varphi_{27}$}
\AXC{$\stackrel{\Pi_{3}}{\vdots}$}\noLine
\UIC{$x \in \mathbb{R} \to e \in \mathbb{R} \to y \in \mathbb{R} \to x - e < y \to -x \in \mathbb{R}$}
\AXC{$\stackrel{\Pi_{23}}{\vdots}$}\noLine
\UIC{$\varphi_{59}$}
\BIC{$x \in \mathbb{R} \to \varphi_{24}$}
\AXC{$\stackrel{\Pi_{35}}{\vdots}$}\noLine
\UIC{$\varphi_{69}$}
\BIC{$(x - e + -x = -e + x + -x \to \varphi_{27}) \to x \in \mathbb{R} \to \varphi_{17}$}
\BIC{$x \in \mathbb{R} \to \varphi_{17}$}
\end{ptree}
\noindent$\Pi_{37}$:
\begin{ptree}
\AXC{$\stackrel{\Pi_{15}}{\vdots}$}\noLine
\UIC{$\varphi_{45}$}
\AXC{$\stackrel{\Pi_{19}}{\vdots}$}\noLine
\UIC{$\varphi_{28}$}
\AXC{$\stackrel{\Pi_{36}}{\vdots}$}\noLine
\UIC{$x \in \mathbb{R} \to \varphi_{17}$}
\AXC{$(x \in \mathbb{R} \to \varphi_{17}) \to \varphi_{28} \to \varphi_{45} \to \varphi_{20}$\lbl{A1}}
\BIC{$\varphi_{28} \to \varphi_{45} \to \varphi_{20}$}
\BIC{$\varphi_{45} \to \varphi_{20}$}
\BIC{$\varphi_{20}$}
\AXC{$\varphi_{20} \to \varphi_{62}$\lbl{A1}}
\BIC{$\varphi_{62}$}
\end{ptree}
\noindent$\Pi_{38}$:
\begin{ptree}
\AXC{$\stackrel{\Pi_{12}}{\vdots}$}\noLine
\UIC{$x \in \mathbb{R} \to e \in \mathbb{R} \to y \in \mathbb{R} \to x - e < y \to -e \in \mathbb{R}$}
\AXC{$\stackrel{\Pi_{13}}{\vdots}$}\noLine
\UIC{$\varphi_{53}$}
\BIC{$x \in \mathbb{R} \to e \in \mathbb{R} \to y \in \mathbb{R} \to x - e < y \to -e + 0 = -e$}
\AXC{$\stackrel{\Pi_{37}}{\vdots}$}\noLine
\UIC{$\varphi_{62}$}
\BIC{$\varphi_{56}$}
\end{ptree}
\noindent$\Pi_{39}$:
\begin{ptree}
\AXC{$\stackrel{\Pi_{8}}{\vdots}$}\noLine
\UIC{$\varphi_{26}$}
\AXC{$\stackrel{\Pi_{10}}{\vdots}$}\noLine
\UIC{$x - e + -x = -e + 0 \to -e + 0 = -e \to x - e + -x = -e$}
\AXC{$\stackrel{\Pi_{38}}{\vdots}$}\noLine
\UIC{$\varphi_{56}$}
\BIC{$x \in \mathbb{R} \to e \in \mathbb{R} \to y \in \mathbb{R} \to x - e < y \to x - e + -x = -e$}
\AXC{$\varphi_{70}$\lbl{A1}}
\BIC{$\varphi_{26} \to \varphi_{58}$}
\BIC{$\varphi_{58}$}
\end{ptree}
\noindent$\Pi_{40}$:
\begin{ptree}
\AXC{$x = x$\lbl{A5}}
\AXC{$x = y \to x = x \to y = x$\lbl{A6}}
\AXC{$(x = y \to x = x \to y = x) \to x = x \to x = y \to y = x$\lbl{A1}}
\BIC{$x = x \to x = y \to y = x$}
\BIC{$x = y \to y = x$}
\UIC{$\forall y\,(x = y \to y = x)$}
\UIC{$\forall x\,\forall y\,(x = y \to y = x)$}
\end{ptree}
\noindent$\Pi_{41}$:
\begin{ptree}
\AXC{$\stackrel{\Pi_{40}}{\vdots}$}\noLine
\UIC{$\forall x\,\forall y\,(x = y \to y = x)$}
\AXC{$\forall x\,\forall y\,(x = y \to y = x) \to \forall y\,(v_{705} = y \to y = v_{705})$\lbl{A2}}
\BIC{$\forall y\,(v_{705} = y \to y = v_{705})$}
\AXC{$\varphi_{1}$\lbl{A2}}
\BIC{$v_{705} = v_{706} \to v_{706} = v_{705}$}
\UIC{$\forall v_{706}\,(v_{705} = v_{706} \to v_{706} = v_{705})$}
\UIC{$\forall v_{705}\,\forall v_{706}\,(v_{705} = v_{706} \to v_{706} = v_{705})$}
\AXC{$\varphi_{4}$\lbl{A2}}
\BIC{$\forall v_{706}\,(y - x = v_{706} \to v_{706} = y - x)$}
\AXC{$\varphi_{10}$\lbl{A2}}
\BIC{$y - x = y + -x \to y + -x = y - x$}
\end{ptree}
\noindent$\Pi_{42}$:
\begin{ptree}
\AXC{$\forall x\,\forall y\,(x - y = x + -y)$\lbl{Def$-$}}
\AXC{$\varphi_{2}$\lbl{A2}}
\BIC{$\forall y\,(v_{703} - y = v_{703} + -y)$}
\AXC{$\varphi_{3}$\lbl{A2}}
\BIC{$v_{703} - v_{704} = v_{703} + -v_{704}$}
\UIC{$\forall v_{704}\,(v_{703} - v_{704} = v_{703} + -v_{704})$}
\UIC{$\forall v_{703}\,\forall v_{704}\,(v_{703} - v_{704} = v_{703} + -v_{704})$}
\AXC{$\varphi_{13}$\lbl{A2}}
\BIC{$\forall v_{704}\,(y - v_{704} = y + -v_{704})$}
\AXC{$\forall v_{704}\,(y - v_{704} = y + -v_{704}) \to y - x = y + -x$\lbl{A2}}
\BIC{$y - x = y + -x$}
\end{ptree}
\noindent$\Pi_{43}$:
\begin{ptree}
\AXC{$\stackrel{\Pi_{41}}{\vdots}$}\noLine
\UIC{$y - x = y + -x \to y + -x = y - x$}
\AXC{$\stackrel{\Pi_{42}}{\vdots}$}\noLine
\UIC{$y - x = y + -x$}
\AXC{$\varphi_{21}$\lbl{A1}}
\BIC{$(y - x = y + -x \to y + -x = y - x) \to y + -x = y - x$}
\BIC{$y + -x = y - x$}
\AXC{$y + -x = y - x \to \varphi_{60}$\lbl{A1}}
\BIC{$\varphi_{60}$}
\end{ptree}
\noindent$\Pi_{44}$:
\begin{ptree}
\AXC{$\stackrel{\Pi_{1}}{\vdots}$}\noLine
\UIC{$x - e + -x = -e \to x - e + -x < y + -x \to -e < y + -x$}
\AXC{$\stackrel{\Pi_{39}}{\vdots}$}\noLine
\UIC{$\varphi_{58}$}
\BIC{$x \in \mathbb{R} \to e \in \mathbb{R} \to y \in \mathbb{R} \to x - e < y \to -e < y + -x$}
\AXC{$\stackrel{\Pi_{43}}{\vdots}$}\noLine
\UIC{$\varphi_{60}$}
\BIC{$\varphi_{49}$}
\end{ptree}
\noindent 主木:
\begin{ptree}
\AXC{$z = w \to -e < z \to -e < w$\lbl{A6}}
\UIC{$\forall w\,(z = w \to -e < z \to -e < w)$}
\UIC{$\forall z\,\forall w\,(z = w \to -e < z \to -e < w)$}
\AXC{$\varphi_{38}$\lbl{A2}}
\BIC{$\forall w\,(y + -x = w \to -e < y + -x \to -e < w)$}
\AXC{$\varphi_{46}$\lbl{A2}}
\BIC{$y + -x = y - x \to -e < y + -x \to -e < y - x$}
\AXC{$\stackrel{\Pi_{44}}{\vdots}$}\noLine
\UIC{$\varphi_{49}$}
\BIC{$x \in \mathbb{R} \to e \in \mathbb{R} \to y \in \mathbb{R} \to x - e < y \to -e < y - x$}
\UIC{$\varphi_{8}$}
\UIC{$\forall e\,(\varphi_{8})$}
\UIC{$\forall x\,\forall e\,(\varphi_{8})$}
\end{ptree}
\begin{abbrevlist}
\item $\varphi_{1} :\equiv \forall y\,(v_{705} = y \to y = v_{705}) \to v_{705} = v_{706} \to v_{706} = v_{705}$
\item $\varphi_{2} :\equiv \forall x\,\forall y\,(x - y = x + -y) \to \forall y\,(v_{703} - y = v_{703} + -y)$
\item $\varphi_{3} :\equiv \forall y\,(v_{703} - y = v_{703} + -y) \to v_{703} - v_{704} = v_{703} + -v_{704}$
\item $\varphi_{4} :\equiv \forall v_{705}\,\forall v_{706}\,(v_{705} = v_{706} \to v_{706} = v_{705}) \to \forall v_{706}\,(y - x = v_{706} \to v_{706} = y - x)$
\item $\varphi_{5} :\equiv \forall z\,(x - e + -x = -e + 0 \to -e + 0 = z \to x - e + -x = z)$
\item $\varphi_{6} :\equiv \forall w\,(z = w \to x - e + -x = z + -x \to x - e + -x = w + -x)$
\item $\varphi_{7} :\equiv \forall x\,\forall y\,\forall z\,(x \in \mathbb{R} \to y \in \mathbb{R} \to z \in \mathbb{R} \to x + y + z = x + (y + z))$
\item $\varphi_{8} :\equiv \forall y\,(x \in \mathbb{R} \to e \in \mathbb{R} \to y \in \mathbb{R} \to x - e < y \to -e < y - x)$
\item $\varphi_{9} :\equiv \forall x\,\forall y\,\forall z\,(x \in \mathbb{R} \to y \in \mathbb{R} \to z \in \mathbb{R} \to x < y \to x + z < y + z)$
\item $\varphi_{10} :\equiv \forall v_{706}\,(y - x = v_{706} \to v_{706} = y - x) \to y - x = y + -x \to y + -x = y - x$
\item $\varphi_{11} :\equiv \forall w\,(x - e + -x = w \to x - e + -x < y + -x \to w < y + -x)$
\item $\varphi_{12} :\equiv z = w \to -e + (x + -x) = -e + z \to -e + (x + -x) = -e + w$
\item $\varphi_{13} :\equiv \forall v_{703}\,\forall v_{704}\,(v_{703} - v_{704} = v_{703} + -v_{704}) \to \forall v_{704}\,(y - v_{704} = y + -v_{704})$
\item $\varphi_{14} :\equiv \forall x\,\forall y\,\forall z\,(x = y \to y = z \to x = z) \to \forall y\,\forall z\,(x - e + -x = y \to y = z \to x - e + -x = z)$
\item $\varphi_{15} :\equiv \forall w\,(z = w \to x + -e + -x = z + -x \to x + -e + -x = w + -x)$
\item $\varphi_{16} :\equiv \forall w\,(x - e = w \to x - e + -x = x - e + -x \to x - e + -x = w + -x)$
\item $\varphi_{17} :\equiv e \in \mathbb{R} \to y \in \mathbb{R} \to x - e < y \to x - e + -x = -e + (x + -x)$
\item $\varphi_{18} :\equiv -e + (x + -x) = -e + (x + -x) \to -e + (x + -x) = -e + 0$
\item $\varphi_{19} :\equiv \forall z\,(-e \in \mathbb{R} \to x \in \mathbb{R} \to z \in \mathbb{R} \to -e + x + z = -e + (x + z))$
\item $\varphi_{20} :\equiv x \in \mathbb{R} \to e \in \mathbb{R} \to y \in \mathbb{R} \to x - e < y \to x - e + -x = -e + 0$
\item $\varphi_{21} :\equiv y - x = y + -x \to (y - x = y + -x \to y + -x = y - x) \to y + -x = y - x$
\item $\varphi_{22} :\equiv \forall z\,(x - e \in \mathbb{R} \to y \in \mathbb{R} \to z \in \mathbb{R} \to x - e < y \to x - e + z < y + z)$
\item $\varphi_{23} :\equiv \forall y\,\forall z\,(-e \in \mathbb{R} \to y \in \mathbb{R} \to z \in \mathbb{R} \to -e + y + z = -e + (y + z))$
\item $\varphi_{24} :\equiv e \in \mathbb{R} \to y \in \mathbb{R} \to x - e < y \to -e + x + -x = -e + (x + -x)$
\item $\varphi_{25} :\equiv e \in \mathbb{R} \to y \in \mathbb{R} \to x - e < y \to x + -e + -x = -e + x + -x$
\item $\varphi_{26} :\equiv x \in \mathbb{R} \to e \in \mathbb{R} \to y \in \mathbb{R} \to x - e < y \to x - e + -x < y + -x$
\item $\varphi_{27} :\equiv -e + x + -x = -e + (x + -x) \to x - e + -x = -e + (x + -x)$
\item $\varphi_{28} :\equiv x \in \mathbb{R} \to e \in \mathbb{R} \to y \in \mathbb{R} \to x - e < y \to -e + (x + -x) = -e + 0$
\item $\varphi_{29} :\equiv \forall y\,(x \in \mathbb{R} \to y \in \mathbb{R} \to x - y \in \mathbb{R}) \to x \in \mathbb{R} \to e \in \mathbb{R} \to x - e \in \mathbb{R}$
\item $\varphi_{30} :\equiv \forall z\,(x - e + -x = -e + (x + -x) \to -e + (x + -x) = z \to x - e + -x = z)$
\item $\varphi_{31} :\equiv \forall z\,(x - e + -x = -e + x + -x \to -e + x + -x = z \to x - e + -x = z)$
\item $\varphi_{32} :\equiv \forall z\,(x - e + -x = x + -e + -x \to x + -e + -x = z \to x - e + -x = z)$
\item $\varphi_{33} :\equiv x + -e + -x = x + -e + -x \to x + -e + -x = -e + x + -x$
\item $\varphi_{34} :\equiv \forall x\,\forall y\,(x \in \mathbb{R} \to y \in \mathbb{R} \to x + y = y + x) \to \forall y\,(x \in \mathbb{R} \to y \in \mathbb{R} \to x + y = y + x)$
\item $\varphi_{35} :\equiv x \in \mathbb{R} \to e \in \mathbb{R} \to y \in \mathbb{R} \to x - e < y \to x - e + -x = x + -e + -x$
\item $\varphi_{36} :\equiv x \in \mathbb{R} \to e \in \mathbb{R} \to y \in \mathbb{R} \to x - e < y \to x - e + -x = -e + x + -x$
\item $\varphi_{37} :\equiv -e \in \mathbb{R} \to x \in \mathbb{R} \to -x \in \mathbb{R} \to -e + x + -x = -e + (x + -x)$
\item $\varphi_{38} :\equiv \forall z\,\forall w\,(z = w \to -e < z \to -e < w) \to \forall w\,(y + -x = w \to -e < y + -x \to -e < w)$
\item $\varphi_{39} :\equiv \forall x\,\forall y\,(x \in \mathbb{R} \to y \in \mathbb{R} \to x - y \in \mathbb{R}) \to \forall y\,(x \in \mathbb{R} \to y \in \mathbb{R} \to x - y \in \mathbb{R})$
\item $\varphi_{40} :\equiv x - e \in \mathbb{R} \to y \in \mathbb{R} \to -x \in \mathbb{R} \to x - e < y \to x - e + -x < y + -x$
\item $\varphi_{41} :\equiv x - e = x + -e \to x - e + -x = x - e + -x \to x - e + -x = x + -e + -x$
\item $\varphi_{42} :\equiv \forall w\,(x + -x = w \to -e + (x + -x) = -e + (x + -x) \to -e + (x + -x) = -e + w)$
\item $\varphi_{43} :\equiv \forall w\,(x + -e = w \to x + -e + -x = x + -e + -x \to x + -e + -x = w + -x)$
\item $\varphi_{44} :\equiv \forall y\,(x \in \mathbb{R} \to y \in \mathbb{R} \to x + y = y + x) \to x \in \mathbb{R} \to -e \in \mathbb{R} \to x + -e = -e + x$
\item $\varphi_{45} :\equiv x - e + -x = -e + (x + -x) \to -e + (x + -x) = -e + 0 \to x - e + -x = -e + 0$
\item $\varphi_{46} :\equiv \forall w\,(y + -x = w \to -e < y + -x \to -e < w) \to y + -x = y - x \to -e < y + -x \to -e < y - x$
\item $\varphi_{47} :\equiv x - e + -x = x + -e + -x \to x + -e + -x = -e + x + -x \to x - e + -x = -e + x + -x$
\item $\varphi_{48} :\equiv (x \in \mathbb{R} \to e \in \mathbb{R} \to y \in \mathbb{R} \to x - e < y \to x \in \mathbb{R}) \to x \in \mathbb{R} \to e \in \mathbb{R} \to y \in \mathbb{R} \to x - e < y \to x + -x = 0$
\item $\varphi_{49} :\equiv (y + -x = y - x \to -e < y + -x \to -e < y - x) \to x \in \mathbb{R} \to e \in \mathbb{R} \to y \in \mathbb{R} \to x - e < y \to -e < y - x$
\item $\varphi_{50} :\equiv (x \in \mathbb{R} \to e \in \mathbb{R} \to y \in \mathbb{R} \to x - e < y \to x \in \mathbb{R}) \to x \in \mathbb{R} \to e \in \mathbb{R} \to y \in \mathbb{R} \to x - e < y \to -x \in \mathbb{R}$
\item $\varphi_{51} :\equiv (x \in \mathbb{R} \to e \in \mathbb{R} \to y \in \mathbb{R} \to x - e < y \to e \in \mathbb{R}) \to x \in \mathbb{R} \to e \in \mathbb{R} \to y \in \mathbb{R} \to x - e < y \to x - e \in \mathbb{R}$
\item $\varphi_{52} :\equiv (x \in \mathbb{R} \to e \in \mathbb{R} \to y \in \mathbb{R} \to x - e < y \to e \in \mathbb{R}) \to x \in \mathbb{R} \to e \in \mathbb{R} \to y \in \mathbb{R} \to x - e < y \to -e \in \mathbb{R}$
\item $\varphi_{53} :\equiv (x \in \mathbb{R} \to e \in \mathbb{R} \to y \in \mathbb{R} \to x - e < y \to -e \in \mathbb{R}) \to x \in \mathbb{R} \to e \in \mathbb{R} \to y \in \mathbb{R} \to x - e < y \to -e + 0 = -e$
\item $\varphi_{54} :\equiv \varphi_{5} \to x - e + -x = -e + 0 \to -e + 0 = -e \to x - e + -x = -e$
\item $\varphi_{55} :\equiv (x \in \mathbb{R} \to e \in \mathbb{R} \to y \in \mathbb{R} \to x - e < y \to -e \in \mathbb{R}) \to x \in \mathbb{R} \to e \in \mathbb{R} \to y \in \mathbb{R} \to x - e < y \to x + -e = -e + x$
\item $\varphi_{56} :\equiv (x - e + -x = -e + 0 \to -e + 0 = -e \to x - e + -x = -e) \to x \in \mathbb{R} \to e \in \mathbb{R} \to y \in \mathbb{R} \to x - e < y \to x - e + -x = -e$
\item $\varphi_{57} :\equiv \varphi_{11} \to x - e + -x = -e \to x - e + -x < y + -x \to -e < y + -x$
\item $\varphi_{58} :\equiv (x - e + -x = -e \to x - e + -x < y + -x \to -e < y + -x) \to x \in \mathbb{R} \to e \in \mathbb{R} \to y \in \mathbb{R} \to x - e < y \to -e < y + -x$
\item $\varphi_{59} :\equiv (x \in \mathbb{R} \to e \in \mathbb{R} \to y \in \mathbb{R} \to x - e < y \to -x \in \mathbb{R}) \to x \in \mathbb{R} \to \varphi_{24}$
\item $\varphi_{60} :\equiv (x \in \mathbb{R} \to e \in \mathbb{R} \to y \in \mathbb{R} \to x - e < y \to -e < y + -x) \to \varphi_{49}$
\item $\varphi_{61} :\equiv (x \in \mathbb{R} \to e \in \mathbb{R} \to y \in \mathbb{R} \to x - e < y \to -x \in \mathbb{R}) \to (x \in \mathbb{R} \to e \in \mathbb{R} \to y \in \mathbb{R} \to x - e < y \to x - e < y) \to \varphi_{26}$
\item $\varphi_{62} :\equiv (x \in \mathbb{R} \to e \in \mathbb{R} \to y \in \mathbb{R} \to x - e < y \to -e + 0 = -e) \to \varphi_{56}$
\item $\varphi_{63} :\equiv -e + (x + -x) = -e + (x + -x) \to (x + -x = 0 \to \varphi_{18}) \to \varphi_{28}$
\item $\varphi_{64} :\equiv (x \in \mathbb{R} \to e \in \mathbb{R} \to x - e \in \mathbb{R}) \to (x \in \mathbb{R} \to e \in \mathbb{R} \to y \in \mathbb{R} \to x - e < y \to x \in \mathbb{R}) \to \varphi_{51}$
\item $\varphi_{65} :\equiv x + -e + -x = x + -e + -x \to (x + -e = -e + x \to \varphi_{33}) \to x \in \mathbb{R} \to \varphi_{25}$
\item $\varphi_{66} :\equiv (x \in \mathbb{R} \to -e \in \mathbb{R} \to x + -e = -e + x) \to (x \in \mathbb{R} \to e \in \mathbb{R} \to y \in \mathbb{R} \to x - e < y \to x \in \mathbb{R}) \to \varphi_{55}$
\item $\varphi_{67} :\equiv (x \in \mathbb{R} \to e \in \mathbb{R} \to y \in \mathbb{R} \to x - e < y \to x + -x = 0) \to \varphi_{63}$
\item $\varphi_{68} :\equiv (x \in \mathbb{R} \to e \in \mathbb{R} \to y \in \mathbb{R} \to x - e < y \to -e \in \mathbb{R}) \to (x \in \mathbb{R} \to e \in \mathbb{R} \to y \in \mathbb{R} \to x - e < y \to x \in \mathbb{R}) \to \varphi_{59}$
\item $\varphi_{69} :\equiv (x \in \mathbb{R} \to \varphi_{24}) \to (x - e + -x = -e + x + -x \to \varphi_{27}) \to x \in \mathbb{R} \to \varphi_{17}$
\item $\varphi_{70} :\equiv (x \in \mathbb{R} \to e \in \mathbb{R} \to y \in \mathbb{R} \to x - e < y \to x - e + -x = -e) \to \varphi_{26} \to \varphi_{58}$
\item $\varphi_{71} :\equiv (x \in \mathbb{R} \to e \in \mathbb{R} \to y \in \mathbb{R} \to x - e < y \to x + -e = -e + x) \to \varphi_{65}$
\item $\varphi_{72} :\equiv (x \in \mathbb{R} \to e \in \mathbb{R} \to y \in \mathbb{R} \to x - e < y \to x - e \in \mathbb{R}) \to (x \in \mathbb{R} \to e \in \mathbb{R} \to y \in \mathbb{R} \to x - e < y \to y \in \mathbb{R}) \to \varphi_{61}$
\end{abbrevlist}
\end{fullproof}

% 自動生成 (prover/build.js): LtAddSwap  ノード数 201
\begin{fullproof}{[E32]}{$\forall x\,\forall e\,\forall y\,(x \in \mathbb{R} \to e \in \mathbb{R} \to y \in \mathbb{R} \to y < x + e \to y - x < e)$}{201}\phantomsection\label{tree:LtAddSwap}
\noindent$\Pi_{1}$:
\begin{ptree}
\AXC{$z = w \to y + -x < z \to y + -x < w$\lbl{A6}}
\UIC{$\forall w\,(z = w \to y + -x < z \to y + -x < w)$}
\UIC{$\forall z\,\forall w\,(z = w \to y + -x < z \to y + -x < w)$}
\AXC{$\forall z\,\forall w\,(z = w \to y + -x < z \to y + -x < w) \to \varphi_{17}$\lbl{A2}}
\BIC{$\varphi_{17}$}
\end{ptree}
\noindent$\Pi_{2}$:
\begin{ptree}
\AXC{$\forall x \in \mathbb{R}\,(-x \in \mathbb{R})$\lbl{(R1)}}
\AXC{$\forall x \in \mathbb{R}\,(-x \in \mathbb{R}) \to x \in \mathbb{R} \to -x \in \mathbb{R}$\lbl{A2}}
\BIC{$x \in \mathbb{R} \to -x \in \mathbb{R}$}
\AXC{$(x \in \mathbb{R} \to -x \in \mathbb{R}) \to \varphi_{51}$\lbl{A1}}
\BIC{$\varphi_{51}$}
\end{ptree}
\noindent$\Pi_{3}$:
\begin{ptree}
\AXC{$x \in \mathbb{R} \to e \in \mathbb{R} \to y \in \mathbb{R} \to y < x + e \to x \in \mathbb{R}$\lbl{A1}}
\AXC{$\stackrel{\Pi_{2}}{\vdots}$}\noLine
\UIC{$\varphi_{51}$}
\BIC{$x \in \mathbb{R} \to e \in \mathbb{R} \to y \in \mathbb{R} \to y < x + e \to -x \in \mathbb{R}$}
\end{ptree}
\noindent$\Pi_{4}$:
\begin{ptree}
\AXC{$\forall x\,\forall y\,(x \in \mathbb{R} \to y \in \mathbb{R} \to x + y \in \mathbb{R})$\lbl{(R1)}}
\AXC{$\varphi_{39}$\lbl{A2}}
\BIC{$\forall y\,(x \in \mathbb{R} \to y \in \mathbb{R} \to x + y \in \mathbb{R})$}
\AXC{$\varphi_{37}$\lbl{A2}}
\BIC{$x \in \mathbb{R} \to e \in \mathbb{R} \to x + e \in \mathbb{R}$}
\AXC{$\varphi_{58}$\lbl{A1}}
\BIC{$(x \in \mathbb{R} \to e \in \mathbb{R} \to y \in \mathbb{R} \to y < x + e \to x \in \mathbb{R}) \to \varphi_{52}$}
\end{ptree}
\noindent$\Pi_{5}$:
\begin{ptree}
\AXC{$x \in \mathbb{R} \to e \in \mathbb{R} \to y \in \mathbb{R} \to y < x + e \to x \in \mathbb{R}$\lbl{A1}}
\AXC{$\stackrel{\Pi_{4}}{\vdots}$}\noLine
\UIC{$(x \in \mathbb{R} \to e \in \mathbb{R} \to y \in \mathbb{R} \to y < x + e \to x \in \mathbb{R}) \to \varphi_{52}$}
\BIC{$\varphi_{52}$}
\end{ptree}
\noindent$\Pi_{6}$:
\begin{ptree}
\AXC{$\varphi_{12}$\lbl{(R14)}}
\AXC{$\varphi_{12} \to \varphi_{29}$\lbl{A2}}
\BIC{$\varphi_{29}$}
\AXC{$\varphi_{29} \to \varphi_{33}$\lbl{A2}}
\BIC{$\varphi_{33}$}
\AXC{$\varphi_{33} \to v_{707} \in \mathbb{R} \to \varphi_{20}$\lbl{A2}}
\BIC{$v_{707} \in \mathbb{R} \to \varphi_{20}$}
\UIC{$\forall v_{709}\,(v_{707} \in \mathbb{R} \to \varphi_{20})$}
\UIC{$\forall v_{708}\,\forall v_{709}\,(v_{707} \in \mathbb{R} \to \varphi_{20})$}
\UIC{$\forall v_{707}\,\forall v_{708}\,\forall v_{709}\,(v_{707} \in \mathbb{R} \to \varphi_{20})$}
\AXC{$\forall v_{707}\,\forall v_{708}\,\forall v_{709}\,(v_{707} \in \mathbb{R} \to \varphi_{20}) \to \varphi_{40}$\lbl{A2}}
\BIC{$\varphi_{40}$}
\end{ptree}
\noindent$\Pi_{7}$:
\begin{ptree}
\AXC{$x \in \mathbb{R} \to e \in \mathbb{R} \to y \in \mathbb{R} \to y < x + e \to y \in \mathbb{R}$\lbl{A1}}
\AXC{$\stackrel{\Pi_{6}}{\vdots}$}\noLine
\UIC{$\varphi_{40}$}
\AXC{$\varphi_{40} \to \varphi_{41}$\lbl{A2}}
\BIC{$\varphi_{41}$}
\AXC{$\varphi_{41} \to \varphi_{42}$\lbl{A2}}
\BIC{$\varphi_{42}$}
\AXC{$\varphi_{66}$\lbl{A1}}
\BIC{$(x \in \mathbb{R} \to e \in \mathbb{R} \to y \in \mathbb{R} \to y < x + e \to y \in \mathbb{R}) \to \varphi_{64}$}
\BIC{$\varphi_{64}$}
\end{ptree}
\noindent$\Pi_{8}$:
\begin{ptree}
\AXC{$\stackrel{\Pi_{3}}{\vdots}$}\noLine
\UIC{$x \in \mathbb{R} \to e \in \mathbb{R} \to y \in \mathbb{R} \to y < x + e \to -x \in \mathbb{R}$}
\AXC{$x \in \mathbb{R} \to e \in \mathbb{R} \to y \in \mathbb{R} \to y < x + e \to e \in \mathbb{R}$\lbl{A1}}
\AXC{$\stackrel{\Pi_{5}}{\vdots}$}\noLine
\UIC{$\varphi_{52}$}
\BIC{$x \in \mathbb{R} \to e \in \mathbb{R} \to y \in \mathbb{R} \to y < x + e \to x + e \in \mathbb{R}$}
\AXC{$\stackrel{\Pi_{7}}{\vdots}$}\noLine
\UIC{$\varphi_{64}$}
\BIC{$\varphi_{56}$}
\BIC{$(x \in \mathbb{R} \to e \in \mathbb{R} \to y \in \mathbb{R} \to y < x + e \to y < x + e) \to \varphi_{31}$}
\end{ptree}
\noindent$\Pi_{9}$:
\begin{ptree}
\AXC{$x \in \mathbb{R} \to e \in \mathbb{R} \to y \in \mathbb{R} \to y < x + e \to y < x + e$\lbl{A1}}
\AXC{$\stackrel{\Pi_{8}}{\vdots}$}\noLine
\UIC{$(x \in \mathbb{R} \to e \in \mathbb{R} \to y \in \mathbb{R} \to y < x + e \to y < x + e) \to \varphi_{31}$}
\BIC{$\varphi_{31}$}
\end{ptree}
\noindent$\Pi_{10}$:
\begin{ptree}
\AXC{$y = z \to x = y \to x = z$\lbl{A6}}
\AXC{$(y = z \to x = y \to x = z) \to x = y \to y = z \to x = z$\lbl{A1}}
\BIC{$x = y \to y = z \to x = z$}
\UIC{$\forall z\,(x = y \to y = z \to x = z)$}
\UIC{$\forall y\,\forall z\,(x = y \to y = z \to x = z)$}
\UIC{$\forall x\,\forall y\,\forall z\,(x = y \to y = z \to x = z)$}
\AXC{$\varphi_{24}$\lbl{A2}}
\BIC{$\forall y\,\forall z\,(x + e + -x = y \to y = z \to x + e + -x = z)$}
\AXC{$\varphi_{44}$\lbl{A2}}
\BIC{$\forall z\,(x + e + -x = e + 0 \to e + 0 = z \to x + e + -x = z)$}
\AXC{$\varphi_{45}$\lbl{A2}}
\BIC{$x + e + -x = e + 0 \to e + 0 = e \to x + e + -x = e$}
\end{ptree}
\noindent$\Pi_{11}$:
\begin{ptree}
\AXC{$\forall x \in \mathbb{R}\,(x + 0 = x)$\lbl{(R4)}}
\AXC{$\forall x \in \mathbb{R}\,(x + 0 = x) \to e \in \mathbb{R} \to e + 0 = e$\lbl{A2}}
\BIC{$e \in \mathbb{R} \to e + 0 = e$}
\AXC{$(e \in \mathbb{R} \to e + 0 = e) \to \varphi_{47}$\lbl{A1}}
\BIC{$\varphi_{47}$}
\end{ptree}
\noindent$\Pi_{12}$:
\begin{ptree}
\AXC{$y = z \to x = y \to x = z$\lbl{A6}}
\AXC{$(y = z \to x = y \to x = z) \to x = y \to y = z \to x = z$\lbl{A1}}
\BIC{$x = y \to y = z \to x = z$}
\UIC{$\forall z\,(x = y \to y = z \to x = z)$}
\UIC{$\forall y\,\forall z\,(x = y \to y = z \to x = z)$}
\UIC{$\forall x\,\forall y\,\forall z\,(x = y \to y = z \to x = z)$}
\AXC{$\varphi_{24}$\lbl{A2}}
\BIC{$\forall y\,\forall z\,(x + e + -x = y \to y = z \to x + e + -x = z)$}
\end{ptree}
\noindent$\Pi_{13}$:
\begin{ptree}
\AXC{$\stackrel{\Pi_{12}}{\vdots}$}\noLine
\UIC{$\forall y\,\forall z\,(x + e + -x = y \to y = z \to x + e + -x = z)$}
\AXC{$\forall y\,\forall z\,(x + e + -x = y \to y = z \to x + e + -x = z) \to \varphi_{25}$\lbl{A2}}
\BIC{$\varphi_{25}$}
\AXC{$\varphi_{25} \to \varphi_{34}$\lbl{A2}}
\BIC{$\varphi_{34}$}
\end{ptree}
\noindent$\Pi_{14}$:
\begin{ptree}
\AXC{$\forall x \in \mathbb{R}\,(x + -x = 0)$\lbl{(R5)}}
\AXC{$\forall x \in \mathbb{R}\,(x + -x = 0) \to x \in \mathbb{R} \to x + -x = 0$\lbl{A2}}
\BIC{$x \in \mathbb{R} \to x + -x = 0$}
\AXC{$(x \in \mathbb{R} \to x + -x = 0) \to \varphi_{49}$\lbl{A1}}
\BIC{$\varphi_{49}$}
\end{ptree}
\noindent$\Pi_{15}$:
\begin{ptree}
\AXC{$x \in \mathbb{R} \to e \in \mathbb{R} \to y \in \mathbb{R} \to y < x + e \to x \in \mathbb{R}$\lbl{A1}}
\AXC{$\stackrel{\Pi_{14}}{\vdots}$}\noLine
\UIC{$\varphi_{49}$}
\BIC{$x \in \mathbb{R} \to e \in \mathbb{R} \to y \in \mathbb{R} \to y < x + e \to x + -x = 0$}
\AXC{$\varphi_{59}$\lbl{A1}}
\BIC{$e + (x + -x) = e + (x + -x) \to \varphi_{18} \to \varphi_{21}$}
\end{ptree}
\noindent$\Pi_{16}$:
\begin{ptree}
\AXC{$x = x$\lbl{A5}}
\UIC{$\forall x\,(x = x)$}
\AXC{$\forall x\,(x = x) \to e + (x + -x) = e + (x + -x)$\lbl{A2}}
\BIC{$e + (x + -x) = e + (x + -x)$}
\AXC{$\stackrel{\Pi_{15}}{\vdots}$}\noLine
\UIC{$e + (x + -x) = e + (x + -x) \to \varphi_{18} \to \varphi_{21}$}
\BIC{$\varphi_{18} \to \varphi_{21}$}
\end{ptree}
\noindent$\Pi_{17}$:
\begin{ptree}
\AXC{$z = w \to e + (x + -x) = e + z \to e + (x + -x) = e + w$\lbl{A6}}
\UIC{$\varphi_{3}$}
\UIC{$\forall z\,(\varphi_{3})$}
\AXC{$\forall z\,(\varphi_{3}) \to \varphi_{26}$\lbl{A2}}
\BIC{$\varphi_{26}$}
\AXC{$\varphi_{26} \to \varphi_{18}$\lbl{A2}}
\BIC{$\varphi_{18}$}
\AXC{$\stackrel{\Pi_{16}}{\vdots}$}\noLine
\UIC{$\varphi_{18} \to \varphi_{21}$}
\BIC{$\varphi_{21}$}
\end{ptree}
\noindent$\Pi_{18}$:
\begin{ptree}
\AXC{$y = z \to x = y \to x = z$\lbl{A6}}
\AXC{$(y = z \to x = y \to x = z) \to x = y \to y = z \to x = z$\lbl{A1}}
\BIC{$x = y \to y = z \to x = z$}
\UIC{$\forall z\,(x = y \to y = z \to x = z)$}
\UIC{$\forall y\,\forall z\,(x = y \to y = z \to x = z)$}
\UIC{$\forall x\,\forall y\,\forall z\,(x = y \to y = z \to x = z)$}
\AXC{$\varphi_{24}$\lbl{A2}}
\BIC{$\forall y\,\forall z\,(x + e + -x = y \to y = z \to x + e + -x = z)$}
\end{ptree}
\noindent$\Pi_{19}$:
\begin{ptree}
\AXC{$\stackrel{\Pi_{18}}{\vdots}$}\noLine
\UIC{$\forall y\,\forall z\,(x + e + -x = y \to y = z \to x + e + -x = z)$}
\AXC{$\forall y\,\forall z\,(x + e + -x = y \to y = z \to x + e + -x = z) \to \varphi_{27}$\lbl{A2}}
\BIC{$\varphi_{27}$}
\end{ptree}
\noindent$\Pi_{20}$:
\begin{ptree}
\AXC{$\varphi_{11}$\lbl{(R2)}}
\AXC{$\varphi_{11} \to \varphi_{7}$\lbl{A2}}
\BIC{$\varphi_{7}$}
\AXC{$\varphi_{7} \to \varphi_{5}$\lbl{A2}}
\BIC{$\varphi_{5}$}
\AXC{$\varphi_{5} \to \varphi_{23}$\lbl{A2}}
\BIC{$\varphi_{23}$}
\AXC{$\varphi_{23} \to \varphi_{63}$\lbl{A1}}
\BIC{$\varphi_{63}$}
\end{ptree}
\noindent$\Pi_{21}$:
\begin{ptree}
\AXC{$x \in \mathbb{R} \to e \in \mathbb{R} \to y \in \mathbb{R} \to y < x + e \to x \in \mathbb{R}$\lbl{A1}}
\AXC{$x \in \mathbb{R} \to e \in \mathbb{R} \to y \in \mathbb{R} \to y < x + e \to e \in \mathbb{R}$\lbl{A1}}
\AXC{$\stackrel{\Pi_{20}}{\vdots}$}\noLine
\UIC{$\varphi_{63}$}
\BIC{$(x \in \mathbb{R} \to e \in \mathbb{R} \to y \in \mathbb{R} \to y < x + e \to x \in \mathbb{R}) \to \varphi_{54}$}
\BIC{$\varphi_{54}$}
\end{ptree}
\noindent$\Pi_{22}$:
\begin{ptree}
\AXC{$\forall x\,\forall y\,(x \in \mathbb{R} \to y \in \mathbb{R} \to x + y = y + x)$\lbl{(R3)}}
\AXC{$\varphi_{38}$\lbl{A2}}
\BIC{$\forall y\,(x \in \mathbb{R} \to y \in \mathbb{R} \to x + y = y + x)$}
\AXC{$\varphi_{32}$\lbl{A2}}
\BIC{$x \in \mathbb{R} \to e \in \mathbb{R} \to x + e = e + x$}
\AXC{$\varphi_{57}$\lbl{A1}}
\BIC{$(x \in \mathbb{R} \to e \in \mathbb{R} \to y \in \mathbb{R} \to y < x + e \to x \in \mathbb{R}) \to \varphi_{50}$}
\end{ptree}
\noindent$\Pi_{23}$:
\begin{ptree}
\AXC{$x \in \mathbb{R} \to e \in \mathbb{R} \to y \in \mathbb{R} \to y < x + e \to x \in \mathbb{R}$\lbl{A1}}
\AXC{$\stackrel{\Pi_{22}}{\vdots}$}\noLine
\UIC{$(x \in \mathbb{R} \to e \in \mathbb{R} \to y \in \mathbb{R} \to y < x + e \to x \in \mathbb{R}) \to \varphi_{50}$}
\BIC{$\varphi_{50}$}
\end{ptree}
\noindent$\Pi_{24}$:
\begin{ptree}
\AXC{$x \in \mathbb{R} \to e \in \mathbb{R} \to y \in \mathbb{R} \to y < x + e \to e \in \mathbb{R}$\lbl{A1}}
\AXC{$\stackrel{\Pi_{23}}{\vdots}$}\noLine
\UIC{$\varphi_{50}$}
\BIC{$x \in \mathbb{R} \to e \in \mathbb{R} \to y \in \mathbb{R} \to y < x + e \to x + e = e + x$}
\AXC{$\varphi_{61}$\lbl{A1}}
\BIC{$x + e + -x = x + e + -x \to \varphi_{35} \to \varphi_{36}$}
\end{ptree}
\noindent$\Pi_{25}$:
\begin{ptree}
\AXC{$x = x$\lbl{A5}}
\UIC{$\forall x\,(x = x)$}
\AXC{$\forall x\,(x = x) \to x + e + -x = x + e + -x$\lbl{A2}}
\BIC{$x + e + -x = x + e + -x$}
\AXC{$\stackrel{\Pi_{24}}{\vdots}$}\noLine
\UIC{$x + e + -x = x + e + -x \to \varphi_{35} \to \varphi_{36}$}
\BIC{$\varphi_{35} \to \varphi_{36}$}
\end{ptree}
\noindent$\Pi_{26}$:
\begin{ptree}
\AXC{$z = w \to x + e + -x = z + -x \to x + e + -x = w + -x$\lbl{A6}}
\UIC{$\varphi_{8}$}
\UIC{$\forall z\,(\varphi_{8})$}
\AXC{$\forall z\,(\varphi_{8}) \to \varphi_{28}$\lbl{A2}}
\BIC{$\varphi_{28}$}
\AXC{$\varphi_{28} \to \varphi_{35}$\lbl{A2}}
\BIC{$\varphi_{35}$}
\AXC{$\stackrel{\Pi_{25}}{\vdots}$}\noLine
\UIC{$\varphi_{35} \to \varphi_{36}$}
\BIC{$\varphi_{36}$}
\AXC{$\varphi_{36} \to \varphi_{62}$\lbl{A1}}
\BIC{$\varphi_{62}$}
\end{ptree}
\noindent$\Pi_{27}$:
\begin{ptree}
\AXC{$\stackrel{\Pi_{19}}{\vdots}$}\noLine
\UIC{$\varphi_{27}$}
\AXC{$\varphi_{27} \to x + e + -x = e + x + -x \to \varphi_{9}$\lbl{A2}}
\BIC{$x + e + -x = e + x + -x \to \varphi_{9}$}
\AXC{$\stackrel{\Pi_{3}}{\vdots}$}\noLine
\UIC{$x \in \mathbb{R} \to e \in \mathbb{R} \to y \in \mathbb{R} \to y < x + e \to -x \in \mathbb{R}$}
\AXC{$\stackrel{\Pi_{21}}{\vdots}$}\noLine
\UIC{$\varphi_{54}$}
\BIC{$x \in \mathbb{R} \to \varphi_{14}$}
\AXC{$\stackrel{\Pi_{26}}{\vdots}$}\noLine
\UIC{$\varphi_{62}$}
\BIC{$(x + e + -x = e + x + -x \to \varphi_{9}) \to x \in \mathbb{R} \to \varphi_{15}$}
\BIC{$x \in \mathbb{R} \to \varphi_{15}$}
\end{ptree}
\noindent$\Pi_{28}$:
\begin{ptree}
\AXC{$\stackrel{\Pi_{13}}{\vdots}$}\noLine
\UIC{$\varphi_{34}$}
\AXC{$\stackrel{\Pi_{17}}{\vdots}$}\noLine
\UIC{$\varphi_{21}$}
\AXC{$\stackrel{\Pi_{27}}{\vdots}$}\noLine
\UIC{$x \in \mathbb{R} \to \varphi_{15}$}
\AXC{$(x \in \mathbb{R} \to \varphi_{15}) \to \varphi_{21} \to \varphi_{34} \to \varphi_{22}$\lbl{A1}}
\BIC{$\varphi_{21} \to \varphi_{34} \to \varphi_{22}$}
\BIC{$\varphi_{34} \to \varphi_{22}$}
\BIC{$\varphi_{22}$}
\AXC{$\varphi_{60}$\lbl{A1}}
\BIC{$(x \in \mathbb{R} \to e \in \mathbb{R} \to y \in \mathbb{R} \to y < x + e \to e + 0 = e) \to \varphi_{48}$}
\end{ptree}
\noindent$\Pi_{29}$:
\begin{ptree}
\AXC{$x \in \mathbb{R} \to e \in \mathbb{R} \to y \in \mathbb{R} \to y < x + e \to e \in \mathbb{R}$\lbl{A1}}
\AXC{$\stackrel{\Pi_{11}}{\vdots}$}\noLine
\UIC{$\varphi_{47}$}
\BIC{$x \in \mathbb{R} \to e \in \mathbb{R} \to y \in \mathbb{R} \to y < x + e \to e + 0 = e$}
\AXC{$\stackrel{\Pi_{28}}{\vdots}$}\noLine
\UIC{$(x \in \mathbb{R} \to e \in \mathbb{R} \to y \in \mathbb{R} \to y < x + e \to e + 0 = e) \to \varphi_{48}$}
\BIC{$\varphi_{48}$}
\end{ptree}
\noindent$\Pi_{30}$:
\begin{ptree}
\AXC{$\stackrel{\Pi_{9}}{\vdots}$}\noLine
\UIC{$\varphi_{31}$}
\AXC{$\stackrel{\Pi_{10}}{\vdots}$}\noLine
\UIC{$x + e + -x = e + 0 \to e + 0 = e \to x + e + -x = e$}
\AXC{$\stackrel{\Pi_{29}}{\vdots}$}\noLine
\UIC{$\varphi_{48}$}
\BIC{$x \in \mathbb{R} \to e \in \mathbb{R} \to y \in \mathbb{R} \to y < x + e \to x + e + -x = e$}
\AXC{$\varphi_{65}$\lbl{A1}}
\BIC{$\varphi_{31} \to \varphi_{53}$}
\BIC{$\varphi_{53}$}
\end{ptree}
\noindent$\Pi_{31}$:
\begin{ptree}
\AXC{$x = x$\lbl{A5}}
\AXC{$x = y \to x = x \to y = x$\lbl{A6}}
\AXC{$(x = y \to x = x \to y = x) \to x = x \to x = y \to y = x$\lbl{A1}}
\BIC{$x = x \to x = y \to y = x$}
\BIC{$x = y \to y = x$}
\UIC{$\forall y\,(x = y \to y = x)$}
\UIC{$\forall x\,\forall y\,(x = y \to y = x)$}
\end{ptree}
\noindent$\Pi_{32}$:
\begin{ptree}
\AXC{$\stackrel{\Pi_{31}}{\vdots}$}\noLine
\UIC{$\forall x\,\forall y\,(x = y \to y = x)$}
\AXC{$\forall x\,\forall y\,(x = y \to y = x) \to \forall y\,(v_{712} = y \to y = v_{712})$\lbl{A2}}
\BIC{$\forall y\,(v_{712} = y \to y = v_{712})$}
\AXC{$\varphi_{1}$\lbl{A2}}
\BIC{$v_{712} = v_{713} \to v_{713} = v_{712}$}
\UIC{$\forall v_{713}\,(v_{712} = v_{713} \to v_{713} = v_{712})$}
\UIC{$\forall v_{712}\,\forall v_{713}\,(v_{712} = v_{713} \to v_{713} = v_{712})$}
\AXC{$\varphi_{6}$\lbl{A2}}
\BIC{$\forall v_{713}\,(y - x = v_{713} \to v_{713} = y - x)$}
\AXC{$\varphi_{13}$\lbl{A2}}
\BIC{$y - x = y + -x \to y + -x = y - x$}
\end{ptree}
\noindent$\Pi_{33}$:
\begin{ptree}
\AXC{$\forall x\,\forall y\,(x - y = x + -y)$\lbl{Def$-$}}
\AXC{$\varphi_{2}$\lbl{A2}}
\BIC{$\forall y\,(v_{710} - y = v_{710} + -y)$}
\AXC{$\varphi_{4}$\lbl{A2}}
\BIC{$v_{710} - v_{711} = v_{710} + -v_{711}$}
\UIC{$\forall v_{711}\,(v_{710} - v_{711} = v_{710} + -v_{711})$}
\UIC{$\forall v_{710}\,\forall v_{711}\,(v_{710} - v_{711} = v_{710} + -v_{711})$}
\AXC{$\varphi_{19}$\lbl{A2}}
\BIC{$\forall v_{711}\,(y - v_{711} = y + -v_{711})$}
\AXC{$\forall v_{711}\,(y - v_{711} = y + -v_{711}) \to y - x = y + -x$\lbl{A2}}
\BIC{$y - x = y + -x$}
\end{ptree}
\noindent$\Pi_{34}$:
\begin{ptree}
\AXC{$\stackrel{\Pi_{32}}{\vdots}$}\noLine
\UIC{$y - x = y + -x \to y + -x = y - x$}
\AXC{$\stackrel{\Pi_{33}}{\vdots}$}\noLine
\UIC{$y - x = y + -x$}
\AXC{$\varphi_{30}$\lbl{A1}}
\BIC{$(y - x = y + -x \to y + -x = y - x) \to y + -x = y - x$}
\BIC{$y + -x = y - x$}
\AXC{$y + -x = y - x \to \varphi_{55}$\lbl{A1}}
\BIC{$\varphi_{55}$}
\end{ptree}
\noindent$\Pi_{35}$:
\begin{ptree}
\AXC{$\stackrel{\Pi_{1}}{\vdots}$}\noLine
\UIC{$\varphi_{17}$}
\AXC{$\varphi_{17} \to x + e + -x = e \to y + -x < x + e + -x \to y + -x < e$\lbl{A2}}
\BIC{$x + e + -x = e \to y + -x < x + e + -x \to y + -x < e$}
\AXC{$\stackrel{\Pi_{30}}{\vdots}$}\noLine
\UIC{$\varphi_{53}$}
\BIC{$x \in \mathbb{R} \to e \in \mathbb{R} \to y \in \mathbb{R} \to y < x + e \to y + -x < e$}
\AXC{$\stackrel{\Pi_{34}}{\vdots}$}\noLine
\UIC{$\varphi_{55}$}
\BIC{$\varphi_{46}$}
\end{ptree}
\noindent 主木:
\begin{ptree}
\AXC{$z = w \to z < e \to w < e$\lbl{A6}}
\UIC{$\forall w\,(z = w \to z < e \to w < e)$}
\UIC{$\forall z\,\forall w\,(z = w \to z < e \to w < e)$}
\AXC{$\varphi_{16}$\lbl{A2}}
\BIC{$\forall w\,(y + -x = w \to y + -x < e \to w < e)$}
\AXC{$\varphi_{43}$\lbl{A2}}
\BIC{$y + -x = y - x \to y + -x < e \to y - x < e$}
\AXC{$\stackrel{\Pi_{35}}{\vdots}$}\noLine
\UIC{$\varphi_{46}$}
\BIC{$x \in \mathbb{R} \to e \in \mathbb{R} \to y \in \mathbb{R} \to y < x + e \to y - x < e$}
\UIC{$\forall y\,(x \in \mathbb{R} \to e \in \mathbb{R} \to y \in \mathbb{R} \to y < x + e \to y - x < e)$}
\UIC{$\varphi_{10}$}
\UIC{$\forall x\,(\varphi_{10})$}
\end{ptree}
\begin{abbrevlist}
\item $\varphi_{1} :\equiv \forall y\,(v_{712} = y \to y = v_{712}) \to v_{712} = v_{713} \to v_{713} = v_{712}$
\item $\varphi_{2} :\equiv \forall x\,\forall y\,(x - y = x + -y) \to \forall y\,(v_{710} - y = v_{710} + -y)$
\item $\varphi_{3} :\equiv \forall w\,(z = w \to e + (x + -x) = e + z \to e + (x + -x) = e + w)$
\item $\varphi_{4} :\equiv \forall y\,(v_{710} - y = v_{710} + -y) \to v_{710} - v_{711} = v_{710} + -v_{711}$
\item $\varphi_{5} :\equiv \forall z\,(e \in \mathbb{R} \to x \in \mathbb{R} \to z \in \mathbb{R} \to e + x + z = e + (x + z))$
\item $\varphi_{6} :\equiv \forall v_{712}\,\forall v_{713}\,(v_{712} = v_{713} \to v_{713} = v_{712}) \to \forall v_{713}\,(y - x = v_{713} \to v_{713} = y - x)$
\item $\varphi_{7} :\equiv \forall y\,\forall z\,(e \in \mathbb{R} \to y \in \mathbb{R} \to z \in \mathbb{R} \to e + y + z = e + (y + z))$
\item $\varphi_{8} :\equiv \forall w\,(z = w \to x + e + -x = z + -x \to x + e + -x = w + -x)$
\item $\varphi_{9} :\equiv e + x + -x = e + (x + -x) \to x + e + -x = e + (x + -x)$
\item $\varphi_{10} :\equiv \forall e\,\forall y\,(x \in \mathbb{R} \to e \in \mathbb{R} \to y \in \mathbb{R} \to y < x + e \to y - x < e)$
\item $\varphi_{11} :\equiv \forall x\,\forall y\,\forall z\,(x \in \mathbb{R} \to y \in \mathbb{R} \to z \in \mathbb{R} \to x + y + z = x + (y + z))$
\item $\varphi_{12} :\equiv \forall x\,\forall y\,\forall z\,(x \in \mathbb{R} \to y \in \mathbb{R} \to z \in \mathbb{R} \to x < y \to x + z < y + z)$
\item $\varphi_{13} :\equiv \forall v_{713}\,(y - x = v_{713} \to v_{713} = y - x) \to y - x = y + -x \to y + -x = y - x$
\item $\varphi_{14} :\equiv e \in \mathbb{R} \to y \in \mathbb{R} \to y < x + e \to e + x + -x = e + (x + -x)$
\item $\varphi_{15} :\equiv e \in \mathbb{R} \to y \in \mathbb{R} \to y < x + e \to x + e + -x = e + (x + -x)$
\item $\varphi_{16} :\equiv \forall z\,\forall w\,(z = w \to z < e \to w < e) \to \forall w\,(y + -x = w \to y + -x < e \to w < e)$
\item $\varphi_{17} :\equiv \forall w\,(x + e + -x = w \to y + -x < x + e + -x \to y + -x < w)$
\item $\varphi_{18} :\equiv x + -x = 0 \to e + (x + -x) = e + (x + -x) \to e + (x + -x) = e + 0$
\item $\varphi_{19} :\equiv \forall v_{710}\,\forall v_{711}\,(v_{710} - v_{711} = v_{710} + -v_{711}) \to \forall v_{711}\,(y - v_{711} = y + -v_{711})$
\item $\varphi_{20} :\equiv v_{708} \in \mathbb{R} \to v_{709} \in \mathbb{R} \to v_{707} < v_{708} \to v_{707} + v_{709} < v_{708} + v_{709}$
\item $\varphi_{21} :\equiv x \in \mathbb{R} \to e \in \mathbb{R} \to y \in \mathbb{R} \to y < x + e \to e + (x + -x) = e + 0$
\item $\varphi_{22} :\equiv x \in \mathbb{R} \to e \in \mathbb{R} \to y \in \mathbb{R} \to y < x + e \to x + e + -x = e + 0$
\item $\varphi_{23} :\equiv e \in \mathbb{R} \to x \in \mathbb{R} \to -x \in \mathbb{R} \to e + x + -x = e + (x + -x)$
\item $\varphi_{24} :\equiv \forall x\,\forall y\,\forall z\,(x = y \to y = z \to x = z) \to \forall y\,\forall z\,(x + e + -x = y \to y = z \to x + e + -x = z)$
\item $\varphi_{25} :\equiv \forall z\,(x + e + -x = e + (x + -x) \to e + (x + -x) = z \to x + e + -x = z)$
\item $\varphi_{26} :\equiv \forall w\,(x + -x = w \to e + (x + -x) = e + (x + -x) \to e + (x + -x) = e + w)$
\item $\varphi_{27} :\equiv \forall z\,(x + e + -x = e + x + -x \to e + x + -x = z \to x + e + -x = z)$
\item $\varphi_{28} :\equiv \forall w\,(x + e = w \to x + e + -x = x + e + -x \to x + e + -x = w + -x)$
\item $\varphi_{29} :\equiv \forall y\,\forall z\,(v_{707} \in \mathbb{R} \to y \in \mathbb{R} \to z \in \mathbb{R} \to v_{707} < y \to v_{707} + z < y + z)$
\item $\varphi_{30} :\equiv y - x = y + -x \to (y - x = y + -x \to y + -x = y - x) \to y + -x = y - x$
\item $\varphi_{31} :\equiv x \in \mathbb{R} \to e \in \mathbb{R} \to y \in \mathbb{R} \to y < x + e \to y + -x < x + e + -x$
\item $\varphi_{32} :\equiv \forall y\,(x \in \mathbb{R} \to y \in \mathbb{R} \to x + y = y + x) \to x \in \mathbb{R} \to e \in \mathbb{R} \to x + e = e + x$
\item $\varphi_{33} :\equiv \forall z\,(v_{707} \in \mathbb{R} \to v_{708} \in \mathbb{R} \to z \in \mathbb{R} \to v_{707} < v_{708} \to v_{707} + z < v_{708} + z)$
\item $\varphi_{34} :\equiv x + e + -x = e + (x + -x) \to e + (x + -x) = e + 0 \to x + e + -x = e + 0$
\item $\varphi_{35} :\equiv x + e = e + x \to x + e + -x = x + e + -x \to x + e + -x = e + x + -x$
\item $\varphi_{36} :\equiv x \in \mathbb{R} \to e \in \mathbb{R} \to y \in \mathbb{R} \to y < x + e \to x + e + -x = e + x + -x$
\item $\varphi_{37} :\equiv \forall y\,(x \in \mathbb{R} \to y \in \mathbb{R} \to x + y \in \mathbb{R}) \to x \in \mathbb{R} \to e \in \mathbb{R} \to x + e \in \mathbb{R}$
\item $\varphi_{38} :\equiv \forall x\,\forall y\,(x \in \mathbb{R} \to y \in \mathbb{R} \to x + y = y + x) \to \forall y\,(x \in \mathbb{R} \to y \in \mathbb{R} \to x + y = y + x)$
\item $\varphi_{39} :\equiv \forall x\,\forall y\,(x \in \mathbb{R} \to y \in \mathbb{R} \to x + y \in \mathbb{R}) \to \forall y\,(x \in \mathbb{R} \to y \in \mathbb{R} \to x + y \in \mathbb{R})$
\item $\varphi_{40} :\equiv \forall v_{708}\,\forall v_{709}\,(y \in \mathbb{R} \to v_{708} \in \mathbb{R} \to v_{709} \in \mathbb{R} \to y < v_{708} \to y + v_{709} < v_{708} + v_{709})$
\item $\varphi_{41} :\equiv \forall v_{709}\,(y \in \mathbb{R} \to x + e \in \mathbb{R} \to v_{709} \in \mathbb{R} \to y < x + e \to y + v_{709} < x + e + v_{709})$
\item $\varphi_{42} :\equiv y \in \mathbb{R} \to x + e \in \mathbb{R} \to -x \in \mathbb{R} \to y < x + e \to y + -x < x + e + -x$
\item $\varphi_{43} :\equiv \forall w\,(y + -x = w \to y + -x < e \to w < e) \to y + -x = y - x \to y + -x < e \to y - x < e$
\item $\varphi_{44} :\equiv \forall y\,\forall z\,(x + e + -x = y \to y = z \to x + e + -x = z) \to \forall z\,(x + e + -x = e + 0 \to e + 0 = z \to x + e + -x = z)$
\item $\varphi_{45} :\equiv \forall z\,(x + e + -x = e + 0 \to e + 0 = z \to x + e + -x = z) \to x + e + -x = e + 0 \to e + 0 = e \to x + e + -x = e$
\item $\varphi_{46} :\equiv (y + -x = y - x \to y + -x < e \to y - x < e) \to x \in \mathbb{R} \to e \in \mathbb{R} \to y \in \mathbb{R} \to y < x + e \to y - x < e$
\item $\varphi_{47} :\equiv (x \in \mathbb{R} \to e \in \mathbb{R} \to y \in \mathbb{R} \to y < x + e \to e \in \mathbb{R}) \to x \in \mathbb{R} \to e \in \mathbb{R} \to y \in \mathbb{R} \to y < x + e \to e + 0 = e$
\item $\varphi_{48} :\equiv (x + e + -x = e + 0 \to e + 0 = e \to x + e + -x = e) \to x \in \mathbb{R} \to e \in \mathbb{R} \to y \in \mathbb{R} \to y < x + e \to x + e + -x = e$
\item $\varphi_{49} :\equiv (x \in \mathbb{R} \to e \in \mathbb{R} \to y \in \mathbb{R} \to y < x + e \to x \in \mathbb{R}) \to x \in \mathbb{R} \to e \in \mathbb{R} \to y \in \mathbb{R} \to y < x + e \to x + -x = 0$
\item $\varphi_{50} :\equiv (x \in \mathbb{R} \to e \in \mathbb{R} \to y \in \mathbb{R} \to y < x + e \to e \in \mathbb{R}) \to x \in \mathbb{R} \to e \in \mathbb{R} \to y \in \mathbb{R} \to y < x + e \to x + e = e + x$
\item $\varphi_{51} :\equiv (x \in \mathbb{R} \to e \in \mathbb{R} \to y \in \mathbb{R} \to y < x + e \to x \in \mathbb{R}) \to x \in \mathbb{R} \to e \in \mathbb{R} \to y \in \mathbb{R} \to y < x + e \to -x \in \mathbb{R}$
\item $\varphi_{52} :\equiv (x \in \mathbb{R} \to e \in \mathbb{R} \to y \in \mathbb{R} \to y < x + e \to e \in \mathbb{R}) \to x \in \mathbb{R} \to e \in \mathbb{R} \to y \in \mathbb{R} \to y < x + e \to x + e \in \mathbb{R}$
\item $\varphi_{53} :\equiv (x + e + -x = e \to y + -x < x + e + -x \to y + -x < e) \to x \in \mathbb{R} \to e \in \mathbb{R} \to y \in \mathbb{R} \to y < x + e \to y + -x < e$
\item $\varphi_{54} :\equiv (x \in \mathbb{R} \to e \in \mathbb{R} \to y \in \mathbb{R} \to y < x + e \to -x \in \mathbb{R}) \to x \in \mathbb{R} \to \varphi_{14}$
\item $\varphi_{55} :\equiv (x \in \mathbb{R} \to e \in \mathbb{R} \to y \in \mathbb{R} \to y < x + e \to y + -x < e) \to \varphi_{46}$
\item $\varphi_{56} :\equiv (x \in \mathbb{R} \to e \in \mathbb{R} \to y \in \mathbb{R} \to y < x + e \to -x \in \mathbb{R}) \to (x \in \mathbb{R} \to e \in \mathbb{R} \to y \in \mathbb{R} \to y < x + e \to y < x + e) \to \varphi_{31}$
\item $\varphi_{57} :\equiv (x \in \mathbb{R} \to e \in \mathbb{R} \to x + e = e + x) \to (x \in \mathbb{R} \to e \in \mathbb{R} \to y \in \mathbb{R} \to y < x + e \to x \in \mathbb{R}) \to \varphi_{50}$
\item $\varphi_{58} :\equiv (x \in \mathbb{R} \to e \in \mathbb{R} \to x + e \in \mathbb{R}) \to (x \in \mathbb{R} \to e \in \mathbb{R} \to y \in \mathbb{R} \to y < x + e \to x \in \mathbb{R}) \to \varphi_{52}$
\item $\varphi_{59} :\equiv (x \in \mathbb{R} \to e \in \mathbb{R} \to y \in \mathbb{R} \to y < x + e \to x + -x = 0) \to e + (x + -x) = e + (x + -x) \to \varphi_{18} \to \varphi_{21}$
\item $\varphi_{60} :\equiv \varphi_{22} \to (x \in \mathbb{R} \to e \in \mathbb{R} \to y \in \mathbb{R} \to y < x + e \to e + 0 = e) \to \varphi_{48}$
\item $\varphi_{61} :\equiv (x \in \mathbb{R} \to e \in \mathbb{R} \to y \in \mathbb{R} \to y < x + e \to x + e = e + x) \to x + e + -x = x + e + -x \to \varphi_{35} \to \varphi_{36}$
\item $\varphi_{62} :\equiv (x \in \mathbb{R} \to \varphi_{14}) \to (x + e + -x = e + x + -x \to \varphi_{9}) \to x \in \mathbb{R} \to \varphi_{15}$
\item $\varphi_{63} :\equiv (x \in \mathbb{R} \to e \in \mathbb{R} \to y \in \mathbb{R} \to y < x + e \to e \in \mathbb{R}) \to (x \in \mathbb{R} \to e \in \mathbb{R} \to y \in \mathbb{R} \to y < x + e \to x \in \mathbb{R}) \to \varphi_{54}$
\item $\varphi_{64} :\equiv (x \in \mathbb{R} \to e \in \mathbb{R} \to y \in \mathbb{R} \to y < x + e \to x + e \in \mathbb{R}) \to \varphi_{56}$
\item $\varphi_{65} :\equiv (x \in \mathbb{R} \to e \in \mathbb{R} \to y \in \mathbb{R} \to y < x + e \to x + e + -x = e) \to \varphi_{31} \to \varphi_{53}$
\item $\varphi_{66} :\equiv \varphi_{42} \to (x \in \mathbb{R} \to e \in \mathbb{R} \to y \in \mathbb{R} \to y < x + e \to y \in \mathbb{R}) \to \varphi_{64}$
\end{abbrevlist}
\end{fullproof}

% 自動生成 (prover/build.js): BandAbs  ノード数 72
\begin{fullproof}{[E33]}{$\forall x\,\forall e\,\forall y\,(x \in \mathbb{R} \to e \in \mathbb{R} \to y \in \mathbb{R} \to x - e < y \to y < x + e \to \lvert y - x \rvert < e)$}{72}\phantomsection\label{tree:BandAbs}
\noindent$\Pi_{1}$:
\begin{ptree}
\AXC{$\varphi_{6}$\lbl{A1}}
\AXC{$\varphi_{7}$\lbl{A1}}
\AXC{$\varphi_{8}$\lbl{A1}}
\AXC{$\forall x\,(\varphi_{3})$\lbl{[E32]}}
\AXC{$\forall x\,(\varphi_{3}) \to \varphi_{3}$\lbl{A2}}
\BIC{$\varphi_{3}$}
\AXC{$\varphi_{21}$\lbl{A2}}
\BIC{$\forall y\,(x \in \mathbb{R} \to e \in \mathbb{R} \to y \in \mathbb{R} \to y < x + e \to y - x < e)$}
\AXC{$\varphi_{20}$\lbl{A2}}
\BIC{$x \in \mathbb{R} \to e \in \mathbb{R} \to y \in \mathbb{R} \to y < x + e \to y - x < e$}
\AXC{$\varphi_{23}$\lbl{A1}}
\BIC{$\varphi_{8} \to \varphi_{7} \to \varphi_{6} \to \varphi_{9} \to \varphi_{10}$}
\BIC{$\varphi_{7} \to \varphi_{6} \to \varphi_{9} \to \varphi_{10}$}
\BIC{$\varphi_{6} \to \varphi_{9} \to \varphi_{10}$}
\BIC{$\varphi_{9} \to \varphi_{10}$}
\end{ptree}
\noindent$\Pi_{2}$:
\begin{ptree}
\AXC{$\forall x\,\forall e\,(\varphi_{4})$\lbl{[E31]}}
\AXC{$\forall x\,\forall e\,(\varphi_{4}) \to \forall e\,(\varphi_{4})$\lbl{A2}}
\BIC{$\forall e\,(\varphi_{4})$}
\AXC{$\forall e\,(\varphi_{4}) \to \varphi_{4}$\lbl{A2}}
\BIC{$\varphi_{4}$}
\end{ptree}
\noindent$\Pi_{3}$:
\begin{ptree}
\AXC{$\varphi_{6}$\lbl{A1}}
\AXC{$\varphi_{7}$\lbl{A1}}
\AXC{$\varphi_{8}$\lbl{A1}}
\AXC{$\stackrel{\Pi_{2}}{\vdots}$}\noLine
\UIC{$\varphi_{4}$}
\AXC{$\varphi_{4} \to x \in \mathbb{R} \to e \in \mathbb{R} \to y \in \mathbb{R} \to x - e < y \to -e < y - x$\lbl{A2}}
\BIC{$x \in \mathbb{R} \to e \in \mathbb{R} \to y \in \mathbb{R} \to x - e < y \to -e < y - x$}
\AXC{$\varphi_{24}$\lbl{A1}}
\BIC{$\varphi_{8} \to \varphi_{7} \to \varphi_{6} \to \varphi_{11} \to \varphi_{14}$}
\BIC{$\varphi_{7} \to \varphi_{6} \to \varphi_{11} \to \varphi_{14}$}
\BIC{$\varphi_{6} \to \varphi_{11} \to \varphi_{14}$}
\BIC{$\varphi_{11} \to \varphi_{14}$}
\end{ptree}
\noindent$\Pi_{4}$:
\begin{ptree}
\AXC{$\forall x\,\forall y\,(x \in \mathbb{R} \to y \in \mathbb{R} \to x - y \in \mathbb{R})$\lbl{[F1]}}
\AXC{$\varphi_{18}$\lbl{A2}}
\BIC{$\forall y\,(v_{716} \in \mathbb{R} \to y \in \mathbb{R} \to v_{716} - y \in \mathbb{R})$}
\AXC{$\varphi_{19}$\lbl{A2}}
\BIC{$v_{716} \in \mathbb{R} \to v_{717} \in \mathbb{R} \to v_{716} - v_{717} \in \mathbb{R}$}
\UIC{$\forall v_{717}\,(v_{716} \in \mathbb{R} \to v_{717} \in \mathbb{R} \to v_{716} - v_{717} \in \mathbb{R})$}
\UIC{$\varphi_{1}$}
\AXC{$\varphi_{1} \to \forall v_{717}\,(y \in \mathbb{R} \to v_{717} \in \mathbb{R} \to y - v_{717} \in \mathbb{R})$\lbl{A2}}
\BIC{$\forall v_{717}\,(y \in \mathbb{R} \to v_{717} \in \mathbb{R} \to y - v_{717} \in \mathbb{R})$}
\end{ptree}
\noindent$\Pi_{5}$:
\begin{ptree}
\AXC{$\varphi_{6}$\lbl{A1}}
\AXC{$\stackrel{\Pi_{4}}{\vdots}$}\noLine
\UIC{$\forall v_{717}\,(y \in \mathbb{R} \to v_{717} \in \mathbb{R} \to y - v_{717} \in \mathbb{R})$}
\AXC{$\varphi_{16}$\lbl{A2}}
\BIC{$y \in \mathbb{R} \to x \in \mathbb{R} \to y - x \in \mathbb{R}$}
\AXC{$(y \in \mathbb{R} \to x \in \mathbb{R} \to y - x \in \mathbb{R}) \to \varphi_{6} \to \varphi_{8} \to \varphi_{13}$\lbl{A1}}
\BIC{$\varphi_{6} \to \varphi_{8} \to \varphi_{13}$}
\BIC{$\varphi_{8} \to \varphi_{13}$}
\end{ptree}
\noindent$\Pi_{6}$:
\begin{ptree}
\AXC{$\forall x\,\forall y\,(x \in \mathbb{R} \to y \in \mathbb{R} \to -y < x \to x < y \to \lvert x \rvert < y)$\lbl{[A12]}}
\AXC{$\varphi_{22}$\lbl{A2}}
\BIC{$\varphi_{5}$}
\AXC{$\varphi_{5} \to v_{714} \in \mathbb{R} \to \varphi_{2}$\lbl{A2}}
\BIC{$v_{714} \in \mathbb{R} \to \varphi_{2}$}
\UIC{$\forall v_{715}\,(v_{714} \in \mathbb{R} \to \varphi_{2})$}
\UIC{$\forall v_{714}\,\forall v_{715}\,(v_{714} \in \mathbb{R} \to \varphi_{2})$}
\end{ptree}
\noindent$\Pi_{7}$:
\begin{ptree}
\AXC{$\stackrel{\Pi_{6}}{\vdots}$}\noLine
\UIC{$\forall v_{714}\,\forall v_{715}\,(v_{714} \in \mathbb{R} \to \varphi_{2})$}
\AXC{$\forall v_{714}\,\forall v_{715}\,(v_{714} \in \mathbb{R} \to \varphi_{2}) \to \varphi_{17}$\lbl{A2}}
\BIC{$\varphi_{17}$}
\AXC{$\varphi_{17} \to \varphi_{12}$\lbl{A2}}
\BIC{$\varphi_{12}$}
\AXC{$\varphi_{12} \to \varphi_{13} \to \varphi_{7} \to \varphi_{14} \to \varphi_{10} \to \varphi_{15}$\lbl{A1}}
\BIC{$\varphi_{13} \to \varphi_{7} \to \varphi_{14} \to \varphi_{10} \to \varphi_{15}$}
\end{ptree}
\noindent 主木:
\begin{ptree}
\AXC{$\varphi_{9}$\lbl{A1}}
\AXC{$\stackrel{\Pi_{1}}{\vdots}$}\noLine
\UIC{$\varphi_{9} \to \varphi_{10}$}
\BIC{$\varphi_{10}$}
\AXC{$\varphi_{11}$\lbl{A1}}
\AXC{$\stackrel{\Pi_{3}}{\vdots}$}\noLine
\UIC{$\varphi_{11} \to \varphi_{14}$}
\BIC{$\varphi_{14}$}
\AXC{$\varphi_{7}$\lbl{A1}}
\AXC{$\varphi_{8}$\lbl{A1}}
\AXC{$\stackrel{\Pi_{5}}{\vdots}$}\noLine
\UIC{$\varphi_{8} \to \varphi_{13}$}
\BIC{$\varphi_{13}$}
\AXC{$\stackrel{\Pi_{7}}{\vdots}$}\noLine
\UIC{$\varphi_{13} \to \varphi_{7} \to \varphi_{14} \to \varphi_{10} \to \varphi_{15}$}
\BIC{$\varphi_{7} \to \varphi_{14} \to \varphi_{10} \to \varphi_{15}$}
\BIC{$\varphi_{14} \to \varphi_{10} \to \varphi_{15}$}
\BIC{$\varphi_{10} \to \varphi_{15}$}
\BIC{$\varphi_{15}$}
\UIC{$\forall y\,(\varphi_{15})$}
\UIC{$\forall e\,\forall y\,(\varphi_{15})$}
\UIC{$\forall x\,\forall e\,\forall y\,(\varphi_{15})$}
\end{ptree}
\begin{abbrevlist}
\item $\varphi_{1} :\equiv \forall v_{716}\,\forall v_{717}\,(v_{716} \in \mathbb{R} \to v_{717} \in \mathbb{R} \to v_{716} - v_{717} \in \mathbb{R})$
\item $\varphi_{2} :\equiv v_{715} \in \mathbb{R} \to -v_{715} < v_{714} \to v_{714} < v_{715} \to \lvert v_{714} \rvert < v_{715}$
\item $\varphi_{3} :\equiv \forall e\,\forall y\,(x \in \mathbb{R} \to e \in \mathbb{R} \to y \in \mathbb{R} \to y < x + e \to y - x < e)$
\item $\varphi_{4} :\equiv \forall y\,(x \in \mathbb{R} \to e \in \mathbb{R} \to y \in \mathbb{R} \to x - e < y \to -e < y - x)$
\item $\varphi_{5} :\equiv \forall y\,(v_{714} \in \mathbb{R} \to y \in \mathbb{R} \to -y < v_{714} \to v_{714} < y \to \lvert v_{714} \rvert < y)$
\item $\varphi_{6} :\equiv x \in \mathbb{R} \to e \in \mathbb{R} \to y \in \mathbb{R} \to x - e < y \to y < x + e \to y \in \mathbb{R}$
\item $\varphi_{7} :\equiv x \in \mathbb{R} \to e \in \mathbb{R} \to y \in \mathbb{R} \to x - e < y \to y < x + e \to e \in \mathbb{R}$
\item $\varphi_{8} :\equiv x \in \mathbb{R} \to e \in \mathbb{R} \to y \in \mathbb{R} \to x - e < y \to y < x + e \to x \in \mathbb{R}$
\item $\varphi_{9} :\equiv x \in \mathbb{R} \to e \in \mathbb{R} \to y \in \mathbb{R} \to x - e < y \to y < x + e \to y < x + e$
\item $\varphi_{10} :\equiv x \in \mathbb{R} \to e \in \mathbb{R} \to y \in \mathbb{R} \to x - e < y \to y < x + e \to y - x < e$
\item $\varphi_{11} :\equiv x \in \mathbb{R} \to e \in \mathbb{R} \to y \in \mathbb{R} \to x - e < y \to y < x + e \to x - e < y$
\item $\varphi_{12} :\equiv y - x \in \mathbb{R} \to e \in \mathbb{R} \to -e < y - x \to y - x < e \to \lvert y - x \rvert < e$
\item $\varphi_{13} :\equiv x \in \mathbb{R} \to e \in \mathbb{R} \to y \in \mathbb{R} \to x - e < y \to y < x + e \to y - x \in \mathbb{R}$
\item $\varphi_{14} :\equiv x \in \mathbb{R} \to e \in \mathbb{R} \to y \in \mathbb{R} \to x - e < y \to y < x + e \to -e < y - x$
\item $\varphi_{15} :\equiv x \in \mathbb{R} \to e \in \mathbb{R} \to y \in \mathbb{R} \to x - e < y \to y < x + e \to \lvert y - x \rvert < e$
\item $\varphi_{16} :\equiv \forall v_{717}\,(y \in \mathbb{R} \to v_{717} \in \mathbb{R} \to y - v_{717} \in \mathbb{R}) \to y \in \mathbb{R} \to x \in \mathbb{R} \to y - x \in \mathbb{R}$
\item $\varphi_{17} :\equiv \forall v_{715}\,(y - x \in \mathbb{R} \to v_{715} \in \mathbb{R} \to -v_{715} < y - x \to y - x < v_{715} \to \lvert y - x \rvert < v_{715})$
\item $\varphi_{18} :\equiv \forall x\,\forall y\,(x \in \mathbb{R} \to y \in \mathbb{R} \to x - y \in \mathbb{R}) \to \forall y\,(v_{716} \in \mathbb{R} \to y \in \mathbb{R} \to v_{716} - y \in \mathbb{R})$
\item $\varphi_{19} :\equiv \forall y\,(v_{716} \in \mathbb{R} \to y \in \mathbb{R} \to v_{716} - y \in \mathbb{R}) \to v_{716} \in \mathbb{R} \to v_{717} \in \mathbb{R} \to v_{716} - v_{717} \in \mathbb{R}$
\item $\varphi_{20} :\equiv \forall y\,(x \in \mathbb{R} \to e \in \mathbb{R} \to y \in \mathbb{R} \to y < x + e \to y - x < e) \to x \in \mathbb{R} \to e \in \mathbb{R} \to y \in \mathbb{R} \to y < x + e \to y - x < e$
\item $\varphi_{21} :\equiv \varphi_{3} \to \forall y\,(x \in \mathbb{R} \to e \in \mathbb{R} \to y \in \mathbb{R} \to y < x + e \to y - x < e)$
\item $\varphi_{22} :\equiv \forall x\,\forall y\,(x \in \mathbb{R} \to y \in \mathbb{R} \to -y < x \to x < y \to \lvert x \rvert < y) \to \varphi_{5}$
\item $\varphi_{23} :\equiv (x \in \mathbb{R} \to e \in \mathbb{R} \to y \in \mathbb{R} \to y < x + e \to y - x < e) \to \varphi_{8} \to \varphi_{7} \to \varphi_{6} \to \varphi_{9} \to \varphi_{10}$
\item $\varphi_{24} :\equiv (x \in \mathbb{R} \to e \in \mathbb{R} \to y \in \mathbb{R} \to x - e < y \to -e < y - x) \to \varphi_{8} \to \varphi_{7} \to \varphi_{6} \to \varphi_{11} \to \varphi_{14}$
\end{abbrevlist}
\end{fullproof}

% 自動生成 (prover/build.js): SubLtSelf  ノード数 47
\begin{fullproof}{[E34]}{$\forall x\,\forall e\,(x \in \mathbb{R} \to e \in \mathbb{R} \to 0 < e \to x - e < x)$}{47}\phantomsection\label{tree:SubLtSelf}
\noindent$\Pi_{1}$:
\begin{ptree}
\AXC{$x \in \mathbb{R} \to e \in \mathbb{R} \to 0 < e \to x \in \mathbb{R}$\lbl{A1}}
\AXC{$\forall x\,\forall y\,(x \in \mathbb{R} \to y \in \mathbb{R} \to x - y \in \mathbb{R})$\lbl{[F1]}}
\AXC{$\varphi_{5}$\lbl{A2}}
\BIC{$\forall y\,(x \in \mathbb{R} \to y \in \mathbb{R} \to x - y \in \mathbb{R})$}
\AXC{$\varphi_{2}$\lbl{A2}}
\BIC{$x \in \mathbb{R} \to e \in \mathbb{R} \to x - e \in \mathbb{R}$}
\AXC{$\varphi_{14}$\lbl{A1}}
\BIC{$(x \in \mathbb{R} \to e \in \mathbb{R} \to 0 < e \to x \in \mathbb{R}) \to \varphi_{7}$}
\BIC{$\varphi_{7}$}
\end{ptree}
\noindent$\Pi_{2}$:
\begin{ptree}
\AXC{$\forall x\,\forall y\,(x \in \mathbb{R} \to y \in \mathbb{R} \to 0 < y \to x < x + y)$\lbl{[E3]}}
\AXC{$\varphi_{11}$\lbl{A2}}
\BIC{$\forall y\,(x - e \in \mathbb{R} \to y \in \mathbb{R} \to 0 < y \to x - e < x - e + y)$}
\AXC{$\varphi_{12}$\lbl{A2}}
\BIC{$x - e \in \mathbb{R} \to e \in \mathbb{R} \to 0 < e \to x - e < x - e + e$}
\end{ptree}
\noindent$\Pi_{3}$:
\begin{ptree}
\AXC{$\stackrel{\Pi_{2}}{\vdots}$}\noLine
\UIC{$x - e \in \mathbb{R} \to e \in \mathbb{R} \to 0 < e \to x - e < x - e + e$}
\AXC{$(x - e \in \mathbb{R} \to e \in \mathbb{R} \to 0 < e \to x - e < x - e + e) \to \varphi_{15}$\lbl{A1}}
\BIC{$\varphi_{15}$}
\end{ptree}
\noindent$\Pi_{4}$:
\begin{ptree}
\AXC{$x \in \mathbb{R} \to e \in \mathbb{R} \to 0 < e \to e \in \mathbb{R}$\lbl{A1}}
\AXC{$x \in \mathbb{R} \to e \in \mathbb{R} \to 0 < e \to e \in \mathbb{R}$\lbl{A1}}
\AXC{$\stackrel{\Pi_{1}}{\vdots}$}\noLine
\UIC{$\varphi_{7}$}
\BIC{$x \in \mathbb{R} \to e \in \mathbb{R} \to 0 < e \to x - e \in \mathbb{R}$}
\AXC{$\stackrel{\Pi_{3}}{\vdots}$}\noLine
\UIC{$\varphi_{15}$}
\BIC{$(x \in \mathbb{R} \to e \in \mathbb{R} \to 0 < e \to e \in \mathbb{R}) \to \varphi_{8}$}
\BIC{$\varphi_{8}$}
\end{ptree}
\noindent$\Pi_{5}$:
\begin{ptree}
\AXC{$x \in \mathbb{R} \to e \in \mathbb{R} \to 0 < e \to x \in \mathbb{R}$\lbl{A1}}
\AXC{$\forall x\,\forall y\,(x \in \mathbb{R} \to y \in \mathbb{R} \to x - y + y = x)$\lbl{[F12]}}
\AXC{$\varphi_{3}$\lbl{A2}}
\BIC{$\forall y\,(x \in \mathbb{R} \to y \in \mathbb{R} \to x - y + y = x)$}
\AXC{$\varphi_{1}$\lbl{A2}}
\BIC{$x \in \mathbb{R} \to e \in \mathbb{R} \to x - e + e = x$}
\AXC{$\varphi_{13}$\lbl{A1}}
\BIC{$(x \in \mathbb{R} \to e \in \mathbb{R} \to 0 < e \to x \in \mathbb{R}) \to \varphi_{6}$}
\BIC{$\varphi_{6}$}
\end{ptree}
\noindent$\Pi_{6}$:
\begin{ptree}
\AXC{$x \in \mathbb{R} \to e \in \mathbb{R} \to 0 < e \to 0 < e$\lbl{A1}}
\AXC{$\stackrel{\Pi_{4}}{\vdots}$}\noLine
\UIC{$\varphi_{8}$}
\BIC{$x \in \mathbb{R} \to e \in \mathbb{R} \to 0 < e \to x - e < x - e + e$}
\AXC{$x \in \mathbb{R} \to e \in \mathbb{R} \to 0 < e \to e \in \mathbb{R}$\lbl{A1}}
\AXC{$\stackrel{\Pi_{5}}{\vdots}$}\noLine
\UIC{$\varphi_{6}$}
\BIC{$x \in \mathbb{R} \to e \in \mathbb{R} \to 0 < e \to x - e + e = x$}
\AXC{$\varphi_{16}$\lbl{A1}}
\BIC{$(x \in \mathbb{R} \to e \in \mathbb{R} \to 0 < e \to x - e < x - e + e) \to \varphi_{9}$}
\BIC{$\varphi_{9}$}
\end{ptree}
\noindent 主木:
\begin{ptree}
\AXC{$z = w \to x - e < z \to x - e < w$\lbl{A6}}
\UIC{$\forall w\,(z = w \to x - e < z \to x - e < w)$}
\UIC{$\forall z\,\forall w\,(z = w \to x - e < z \to x - e < w)$}
\AXC{$\varphi_{4}$\lbl{A2}}
\BIC{$\forall w\,(x - e + e = w \to x - e < x - e + e \to x - e < w)$}
\AXC{$\varphi_{10}$\lbl{A2}}
\BIC{$x - e + e = x \to x - e < x - e + e \to x - e < x$}
\AXC{$\stackrel{\Pi_{6}}{\vdots}$}\noLine
\UIC{$\varphi_{9}$}
\BIC{$x \in \mathbb{R} \to e \in \mathbb{R} \to 0 < e \to x - e < x$}
\UIC{$\forall e\,(x \in \mathbb{R} \to e \in \mathbb{R} \to 0 < e \to x - e < x)$}
\UIC{$\forall x\,\forall e\,(x \in \mathbb{R} \to e \in \mathbb{R} \to 0 < e \to x - e < x)$}
\end{ptree}
\begin{abbrevlist}
\item $\varphi_{1} :\equiv \forall y\,(x \in \mathbb{R} \to y \in \mathbb{R} \to x - y + y = x) \to x \in \mathbb{R} \to e \in \mathbb{R} \to x - e + e = x$
\item $\varphi_{2} :\equiv \forall y\,(x \in \mathbb{R} \to y \in \mathbb{R} \to x - y \in \mathbb{R}) \to x \in \mathbb{R} \to e \in \mathbb{R} \to x - e \in \mathbb{R}$
\item $\varphi_{3} :\equiv \forall x\,\forall y\,(x \in \mathbb{R} \to y \in \mathbb{R} \to x - y + y = x) \to \forall y\,(x \in \mathbb{R} \to y \in \mathbb{R} \to x - y + y = x)$
\item $\varphi_{4} :\equiv \forall z\,\forall w\,(z = w \to x - e < z \to x - e < w) \to \forall w\,(x - e + e = w \to x - e < x - e + e \to x - e < w)$
\item $\varphi_{5} :\equiv \forall x\,\forall y\,(x \in \mathbb{R} \to y \in \mathbb{R} \to x - y \in \mathbb{R}) \to \forall y\,(x \in \mathbb{R} \to y \in \mathbb{R} \to x - y \in \mathbb{R})$
\item $\varphi_{6} :\equiv (x \in \mathbb{R} \to e \in \mathbb{R} \to 0 < e \to e \in \mathbb{R}) \to x \in \mathbb{R} \to e \in \mathbb{R} \to 0 < e \to x - e + e = x$
\item $\varphi_{7} :\equiv (x \in \mathbb{R} \to e \in \mathbb{R} \to 0 < e \to e \in \mathbb{R}) \to x \in \mathbb{R} \to e \in \mathbb{R} \to 0 < e \to x - e \in \mathbb{R}$
\item $\varphi_{8} :\equiv (x \in \mathbb{R} \to e \in \mathbb{R} \to 0 < e \to 0 < e) \to x \in \mathbb{R} \to e \in \mathbb{R} \to 0 < e \to x - e < x - e + e$
\item $\varphi_{9} :\equiv (x - e + e = x \to x - e < x - e + e \to x - e < x) \to x \in \mathbb{R} \to e \in \mathbb{R} \to 0 < e \to x - e < x$
\item $\varphi_{10} :\equiv \forall w\,(x - e + e = w \to x - e < x - e + e \to x - e < w) \to x - e + e = x \to x - e < x - e + e \to x - e < x$
\item $\varphi_{11} :\equiv \forall x\,\forall y\,(x \in \mathbb{R} \to y \in \mathbb{R} \to 0 < y \to x < x + y) \to \forall y\,(x - e \in \mathbb{R} \to y \in \mathbb{R} \to 0 < y \to x - e < x - e + y)$
\item $\varphi_{12} :\equiv \forall y\,(x - e \in \mathbb{R} \to y \in \mathbb{R} \to 0 < y \to x - e < x - e + y) \to x - e \in \mathbb{R} \to e \in \mathbb{R} \to 0 < e \to x - e < x - e + e$
\item $\varphi_{13} :\equiv (x \in \mathbb{R} \to e \in \mathbb{R} \to x - e + e = x) \to (x \in \mathbb{R} \to e \in \mathbb{R} \to 0 < e \to x \in \mathbb{R}) \to \varphi_{6}$
\item $\varphi_{14} :\equiv (x \in \mathbb{R} \to e \in \mathbb{R} \to x - e \in \mathbb{R}) \to (x \in \mathbb{R} \to e \in \mathbb{R} \to 0 < e \to x \in \mathbb{R}) \to \varphi_{7}$
\item $\varphi_{15} :\equiv (x \in \mathbb{R} \to e \in \mathbb{R} \to 0 < e \to x - e \in \mathbb{R}) \to (x \in \mathbb{R} \to e \in \mathbb{R} \to 0 < e \to e \in \mathbb{R}) \to \varphi_{8}$
\item $\varphi_{16} :\equiv (x \in \mathbb{R} \to e \in \mathbb{R} \to 0 < e \to x - e + e = x) \to (x \in \mathbb{R} \to e \in \mathbb{R} \to 0 < e \to x - e < x - e + e) \to \varphi_{9}$
\end{abbrevlist}
\end{fullproof}

% 自動生成 (prover/build.js): SupApprox  ノード数 139
\begin{fullproof}{[C22]}{$\forall s\,\forall E\,\forall e\,(E \subseteq \mathbb{R} \to \mathrm{IsSup}(s,E) \to e \in \mathbb{R} \to 0 < e \to \exists x \in E\,(s - e < x))$}{139}\phantomsection\label{tree:SupApprox}
\noindent$\Pi_{1}$:
\begin{ptree}
\AXC{$E \subseteq \mathbb{R} \to \mathrm{IsSup}(s,E) \to \varphi_{17}$\lbl{A1}}
\AXC{$(E \subseteq \mathbb{R} \to \mathrm{IsSup}(s,E) \to \varphi_{17}) \to E \subseteq \mathbb{R} \to \varphi_{23}$\lbl{A1}}
\BIC{$E \subseteq \mathbb{R} \to \varphi_{23}$}
\end{ptree}
\noindent$\Pi_{2}$:
\begin{ptree}
\AXC{$E \subseteq \mathbb{R} \to E \subseteq \mathbb{R}$\lbl{A1}}
\AXC{$\forall X\,\forall Y\,\forall x\,(X \subseteq Y \to x \in X \to x \in Y)$\lbl{[B5]}}
\AXC{$\varphi_{28}$\lbl{A2}}
\BIC{$\forall Y\,\forall x\,(E \subseteq Y \to x \in E \to x \in Y)$}
\AXC{$\varphi_{29}$\lbl{A2}}
\BIC{$\forall x\,(E \subseteq \mathbb{R} \to x \in E \to x \in \mathbb{R})$}
\AXC{$\varphi_{30}$\lbl{A2}}
\BIC{$E \subseteq \mathbb{R} \to x \in E \to x \in \mathbb{R}$}
\AXC{$\varphi_{52}$\lbl{A1}}
\BIC{$(E \subseteq \mathbb{R} \to E \subseteq \mathbb{R}) \to (E \subseteq \mathbb{R} \to \varphi_{20}) \to E \subseteq \mathbb{R} \to \varphi_{24}$}
\BIC{$(E \subseteq \mathbb{R} \to \varphi_{20}) \to E \subseteq \mathbb{R} \to \varphi_{24}$}
\end{ptree}
\noindent$\Pi_{3}$:
\begin{ptree}
\AXC{$E \subseteq \mathbb{R} \to \mathrm{IsSup}(s,E) \to \varphi_{17}$\lbl{A1}}
\AXC{$(E \subseteq \mathbb{R} \to \mathrm{IsSup}(s,E) \to \varphi_{17}) \to E \subseteq \mathbb{R} \to \varphi_{20}$\lbl{A1}}
\BIC{$E \subseteq \mathbb{R} \to \varphi_{20}$}
\AXC{$\stackrel{\Pi_{2}}{\vdots}$}\noLine
\UIC{$(E \subseteq \mathbb{R} \to \varphi_{20}) \to E \subseteq \mathbb{R} \to \varphi_{24}$}
\BIC{$E \subseteq \mathbb{R} \to \varphi_{24}$}
\end{ptree}
\noindent$\Pi_{4}$:
\begin{ptree}
\AXC{$\forall s\,(\varphi_{2})$\lbl{DefSup}}
\AXC{$\forall s\,(\varphi_{2}) \to \varphi_{2}$\lbl{A2}}
\BIC{$\varphi_{2}$}
\AXC{$\varphi_{2} \to \varphi_{1}$\lbl{A2}}
\BIC{$\varphi_{1}$}
\AXC{$E \subseteq \mathbb{R} \to \mathrm{IsSup}(s,E) \to \mathrm{IsSup}(s,E)$\lbl{A1}}
\AXC{$\varphi_{47}$\lbl{A1}}
\BIC{$\varphi_{1} \to \varphi_{37}$}
\BIC{$\varphi_{37}$}
\end{ptree}
\noindent$\Pi_{5}$:
\begin{ptree}
\AXC{$\forall a\,\forall X\,(\mathrm{UB}(a,X) \leftrightarrow a \in \mathbb{R} \land \forall x \in X\,(x \leq a))$\lbl{DefUB}}
\AXC{$\varphi_{34}$\lbl{A2}}
\BIC{$\forall X\,(\mathrm{UB}(s,X) \leftrightarrow s \in \mathbb{R} \land \forall x \in X\,(x \leq s))$}
\AXC{$\varphi_{32}$\lbl{A2}}
\BIC{$\mathrm{UB}(s,E) \leftrightarrow s \in \mathbb{R} \land \forall x \in E\,(x \leq s)$}
\AXC{$\stackrel{\Pi_{4}}{\vdots}$}\noLine
\UIC{$\varphi_{37}$}
\BIC{$E \subseteq \mathbb{R} \to \mathrm{IsSup}(s,E) \to e \in \mathbb{R} \to 0 < e \to s \in \mathbb{R}$}
\end{ptree}
\noindent$\Pi_{6}$:
\begin{ptree}
\AXC{$\forall x\,\forall y\,(x \in \mathbb{R} \to y \in \mathbb{R} \to x - y \in \mathbb{R})$\lbl{[F1]}}
\AXC{$\varphi_{35}$\lbl{A2}}
\BIC{$\forall y\,(s \in \mathbb{R} \to y \in \mathbb{R} \to s - y \in \mathbb{R})$}
\AXC{$\varphi_{33}$\lbl{A2}}
\BIC{$s \in \mathbb{R} \to e \in \mathbb{R} \to s - e \in \mathbb{R}$}
\AXC{$(s \in \mathbb{R} \to e \in \mathbb{R} \to s - e \in \mathbb{R}) \to \varphi_{50}$\lbl{A1}}
\BIC{$\varphi_{50}$}
\end{ptree}
\noindent$\Pi_{7}$:
\begin{ptree}
\AXC{$E \subseteq \mathbb{R} \to \mathrm{IsSup}(s,E) \to e \in \mathbb{R} \to e \in \mathbb{R}$\lbl{A1}}
\AXC{$\stackrel{\Pi_{5}}{\vdots}$}\noLine
\UIC{$E \subseteq \mathbb{R} \to \mathrm{IsSup}(s,E) \to e \in \mathbb{R} \to 0 < e \to s \in \mathbb{R}$}
\AXC{$\stackrel{\Pi_{6}}{\vdots}$}\noLine
\UIC{$\varphi_{50}$}
\BIC{$\varphi_{43}$}
\BIC{$E \subseteq \mathbb{R} \to \varphi_{27}$}
\end{ptree}
\noindent$\Pi_{8}$:
\begin{ptree}
\AXC{$\forall x\,\forall y\,(x \in \mathbb{R} \to y \in \mathbb{R} \to \neg y \leq x \to x < y)$\lbl{[A2]}}
\AXC{$\varphi_{39}$\lbl{A2}}
\BIC{$\forall y\,(s - e \in \mathbb{R} \to y \in \mathbb{R} \to \neg y \leq s - e \to s - e < y)$}
\AXC{$\varphi_{40}$\lbl{A2}}
\BIC{$s - e \in \mathbb{R} \to x \in \mathbb{R} \to \neg x \leq s - e \to s - e < x$}
\end{ptree}
\noindent$\Pi_{9}$:
\begin{ptree}
\AXC{$\stackrel{\Pi_{7}}{\vdots}$}\noLine
\UIC{$E \subseteq \mathbb{R} \to \varphi_{27}$}
\AXC{$\stackrel{\Pi_{8}}{\vdots}$}\noLine
\UIC{$s - e \in \mathbb{R} \to x \in \mathbb{R} \to \neg x \leq s - e \to s - e < x$}
\AXC{$(s - e \in \mathbb{R} \to x \in \mathbb{R} \to \neg x \leq s - e \to s - e < x) \to \varphi_{56}$\lbl{A1}}
\BIC{$\varphi_{56}$}
\BIC{$(E \subseteq \mathbb{R} \to \varphi_{24}) \to (E \subseteq \mathbb{R} \to \varphi_{23}) \to E \subseteq \mathbb{R} \to \varphi_{25}$}
\end{ptree}
\noindent$\Pi_{10}$:
\begin{ptree}
\AXC{$\stackrel{\Pi_{1}}{\vdots}$}\noLine
\UIC{$E \subseteq \mathbb{R} \to \varphi_{23}$}
\AXC{$\stackrel{\Pi_{3}}{\vdots}$}\noLine
\UIC{$E \subseteq \mathbb{R} \to \varphi_{24}$}
\AXC{$\stackrel{\Pi_{9}}{\vdots}$}\noLine
\UIC{$(E \subseteq \mathbb{R} \to \varphi_{24}) \to (E \subseteq \mathbb{R} \to \varphi_{23}) \to E \subseteq \mathbb{R} \to \varphi_{25}$}
\BIC{$(E \subseteq \mathbb{R} \to \varphi_{23}) \to E \subseteq \mathbb{R} \to \varphi_{25}$}
\BIC{$E \subseteq \mathbb{R} \to \varphi_{25}$}
\end{ptree}
\noindent$\Pi_{11}$:
\begin{ptree}
\AXC{$E \subseteq \mathbb{R} \to \mathrm{IsSup}(s,E) \to \varphi_{17}$\lbl{A1}}
\AXC{$(E \subseteq \mathbb{R} \to \mathrm{IsSup}(s,E) \to \varphi_{17}) \to E \subseteq \mathbb{R} \to \varphi_{20}$\lbl{A1}}
\BIC{$E \subseteq \mathbb{R} \to \varphi_{20}$}
\AXC{$\varphi_{55}$\lbl{A1}}
\BIC{$(E \subseteq \mathbb{R} \to \varphi_{25}) \to E \subseteq \mathbb{R} \to \mathrm{IsSup}(s,E) \to \varphi_{22}$}
\end{ptree}
\noindent$\Pi_{12}$:
\begin{ptree}
\AXC{$\stackrel{\Pi_{10}}{\vdots}$}\noLine
\UIC{$E \subseteq \mathbb{R} \to \varphi_{25}$}
\AXC{$\stackrel{\Pi_{11}}{\vdots}$}\noLine
\UIC{$(E \subseteq \mathbb{R} \to \varphi_{25}) \to E \subseteq \mathbb{R} \to \mathrm{IsSup}(s,E) \to \varphi_{22}$}
\BIC{$E \subseteq \mathbb{R} \to \mathrm{IsSup}(s,E) \to \varphi_{22}$}
\AXC{$(E \subseteq \mathbb{R} \to \mathrm{IsSup}(s,E) \to \varphi_{22}) \to \varphi_{45}$\lbl{A1}}
\BIC{$\varphi_{45}$}
\end{ptree}
\noindent$\Pi_{13}$:
\begin{ptree}
\AXC{$\forall x\,(\neg (x \in E \land s - e < x)) \to \neg (x \in E \land s - e < x)$\lbl{A2}}
\AXC{$\stackrel{\Pi_{12}}{\vdots}$}\noLine
\UIC{$\varphi_{45}$}
\BIC{$E \subseteq \mathbb{R} \to \mathrm{IsSup}(s,E) \to \varphi_{26}$}
\AXC{$\varphi_{48}$\lbl{A1}}
\BIC{$\varphi_{31} \to \neg \neg (x \in E \to x \leq s - e)$}
\UIC{$\forall x\,(\varphi_{31} \to \neg \neg (x \in E \to x \leq s - e))$}
\end{ptree}
\noindent$\Pi_{14}$:
\begin{ptree}
\AXC{$\varphi_{4}$\lbl{A2}}
\AXC{$\varphi_{4} \to \forall x\,(\neg \neg (x \in E \to x \leq s - e)) \to x \in E \to x \leq s - e$\lbl{A1}}
\BIC{$\forall x\,(\neg \neg (x \in E \to x \leq s - e)) \to x \in E \to x \leq s - e$}
\UIC{$\varphi_{5}$}
\end{ptree}
\noindent$\Pi_{15}$:
\begin{ptree}
\AXC{$\forall x\,\forall y\,(x \in \mathbb{R} \to y \in \mathbb{R} \to x - y \in \mathbb{R})$\lbl{[F1]}}
\AXC{$\varphi_{35}$\lbl{A2}}
\BIC{$\forall y\,(s \in \mathbb{R} \to y \in \mathbb{R} \to s - y \in \mathbb{R})$}
\AXC{$\varphi_{33}$\lbl{A2}}
\BIC{$s \in \mathbb{R} \to e \in \mathbb{R} \to s - e \in \mathbb{R}$}
\AXC{$(s \in \mathbb{R} \to e \in \mathbb{R} \to s - e \in \mathbb{R}) \to \varphi_{46}$\lbl{A1}}
\BIC{$\varphi_{46}$}
\end{ptree}
\noindent$\Pi_{16}$:
\begin{ptree}
\AXC{$E \subseteq \mathbb{R} \to \mathrm{IsSup}(s,E) \to e \in \mathbb{R} \to e \in \mathbb{R}$\lbl{A1}}
\AXC{$\stackrel{\Pi_{5}}{\vdots}$}\noLine
\UIC{$E \subseteq \mathbb{R} \to \mathrm{IsSup}(s,E) \to e \in \mathbb{R} \to 0 < e \to s \in \mathbb{R}$}
\AXC{$\stackrel{\Pi_{15}}{\vdots}$}\noLine
\UIC{$\varphi_{46}$}
\BIC{$(E \subseteq \mathbb{R} \to \mathrm{IsSup}(s,E) \to e \in \mathbb{R} \to e \in \mathbb{R}) \to \varphi_{19}$}
\BIC{$\varphi_{19}$}
\end{ptree}
\noindent$\Pi_{17}$:
\begin{ptree}
\AXC{$\forall a\,\forall X\,(\mathrm{UB}(a,X) \leftrightarrow a \in \mathbb{R} \land \forall x \in X\,(x \leq a))$\lbl{DefUB}}
\AXC{$\forall a\,\forall X\,(\mathrm{UB}(a,X) \leftrightarrow a \in \mathbb{R} \land \forall x \in X\,(x \leq a)) \to \varphi_{10}$\lbl{A2}}
\BIC{$\varphi_{10}$}
\AXC{$\varphi_{41}$\lbl{A2}}
\BIC{$\mathrm{UB}(s - e,E) \leftrightarrow s - e \in \mathbb{R} \land \forall x \in E\,(x \leq s - e)$}
\end{ptree}
\noindent$\Pi_{18}$:
\begin{ptree}
\AXC{$\forall s\,(\varphi_{2})$\lbl{DefSup}}
\AXC{$\forall s\,(\varphi_{2}) \to \varphi_{2}$\lbl{A2}}
\BIC{$\varphi_{2}$}
\AXC{$\varphi_{2} \to \varphi_{1}$\lbl{A2}}
\BIC{$\varphi_{1}$}
\AXC{$E \subseteq \mathbb{R} \to \mathrm{IsSup}(s,E) \to \mathrm{IsSup}(s,E)$\lbl{A1}}
\AXC{$\varphi_{44}$\lbl{A1}}
\BIC{$\varphi_{1} \to E \subseteq \mathbb{R} \to \varphi_{7}$}
\BIC{$E \subseteq \mathbb{R} \to \varphi_{7}$}
\end{ptree}
\noindent$\Pi_{19}$:
\begin{ptree}
\AXC{$E \subseteq \mathbb{R} \to \mathrm{IsSup}(s,E) \to \varphi_{9}$\lbl{A1}}
\AXC{$\varphi_{3}$\lbl{A2}}
\AXC{$\stackrel{\Pi_{18}}{\vdots}$}\noLine
\UIC{$E \subseteq \mathbb{R} \to \varphi_{7}$}
\AXC{$(E \subseteq \mathbb{R} \to \varphi_{7}) \to \varphi_{3} \to E \subseteq \mathbb{R} \to \varphi_{13}$\lbl{A1}}
\BIC{$\varphi_{3} \to E \subseteq \mathbb{R} \to \varphi_{13}$}
\BIC{$E \subseteq \mathbb{R} \to \varphi_{13}$}
\AXC{$\varphi_{53}$\lbl{A1}}
\BIC{$(E \subseteq \mathbb{R} \to \mathrm{IsSup}(s,E) \to \varphi_{9}) \to E \subseteq \mathbb{R} \to \varphi_{15}$}
\BIC{$E \subseteq \mathbb{R} \to \varphi_{15}$}
\end{ptree}
\noindent$\Pi_{20}$:
\begin{ptree}
\AXC{$\forall x\,\forall y\,(x \leq y \to \neg y < x)$\lbl{[O7]}}
\AXC{$\forall x\,\forall y\,(x \leq y \to \neg y < x) \to \forall y\,(s \leq y \to \neg y < s)$\lbl{A2}}
\BIC{$\forall y\,(s \leq y \to \neg y < s)$}
\end{ptree}
\noindent$\Pi_{21}$:
\begin{ptree}
\AXC{$\stackrel{\Pi_{20}}{\vdots}$}\noLine
\UIC{$\forall y\,(s \leq y \to \neg y < s)$}
\AXC{$\forall y\,(s \leq y \to \neg y < s) \to s \leq s - e \to \neg s - e < s$\lbl{A2}}
\BIC{$s \leq s - e \to \neg s - e < s$}
\end{ptree}
\noindent$\Pi_{22}$:
\begin{ptree}
\AXC{$\stackrel{\Pi_{19}}{\vdots}$}\noLine
\UIC{$E \subseteq \mathbb{R} \to \varphi_{15}$}
\AXC{$\stackrel{\Pi_{21}}{\vdots}$}\noLine
\UIC{$s \leq s - e \to \neg s - e < s$}
\AXC{$(s \leq s - e \to \neg s - e < s) \to (E \subseteq \mathbb{R} \to \varphi_{15}) \to E \subseteq \mathbb{R} \to \varphi_{21}$\lbl{A1}}
\BIC{$(E \subseteq \mathbb{R} \to \varphi_{15}) \to E \subseteq \mathbb{R} \to \varphi_{21}$}
\BIC{$E \subseteq \mathbb{R} \to \varphi_{21}$}
\end{ptree}
\noindent$\Pi_{23}$:
\begin{ptree}
\AXC{$\forall x\,\forall e\,(x \in \mathbb{R} \to e \in \mathbb{R} \to 0 < e \to x - e < x)$\lbl{[E34]}}
\AXC{$\varphi_{38}$\lbl{A2}}
\BIC{$\forall e\,(s \in \mathbb{R} \to e \in \mathbb{R} \to 0 < e \to s - e < s)$}
\AXC{$\varphi_{36}$\lbl{A2}}
\BIC{$s \in \mathbb{R} \to e \in \mathbb{R} \to 0 < e \to s - e < s$}
\AXC{$(s \in \mathbb{R} \to e \in \mathbb{R} \to 0 < e \to s - e < s) \to \varphi_{54}$\lbl{A1}}
\BIC{$\varphi_{54}$}
\end{ptree}
\noindent$\Pi_{24}$:
\begin{ptree}
\AXC{$E \subseteq \mathbb{R} \to \mathrm{IsSup}(s,E) \to e \in \mathbb{R} \to e \in \mathbb{R}$\lbl{A1}}
\AXC{$\stackrel{\Pi_{5}}{\vdots}$}\noLine
\UIC{$E \subseteq \mathbb{R} \to \mathrm{IsSup}(s,E) \to e \in \mathbb{R} \to 0 < e \to s \in \mathbb{R}$}
\AXC{$\stackrel{\Pi_{23}}{\vdots}$}\noLine
\UIC{$\varphi_{54}$}
\BIC{$(E \subseteq \mathbb{R} \to \mathrm{IsSup}(s,E) \to e \in \mathbb{R} \to e \in \mathbb{R}) \to \varphi_{42}$}
\BIC{$\varphi_{42}$}
\end{ptree}
\noindent$\Pi_{25}$:
\begin{ptree}
\AXC{$E \subseteq \mathbb{R} \to \mathrm{IsSup}(s,E) \to e \in \mathbb{R} \to 0 < e \to 0 < e$\lbl{A1}}
\AXC{$\stackrel{\Pi_{24}}{\vdots}$}\noLine
\UIC{$\varphi_{42}$}
\BIC{$E \subseteq \mathbb{R} \to \varphi_{18}$}
\AXC{$(E \subseteq \mathbb{R} \to \varphi_{18}) \to (E \subseteq \mathbb{R} \to \varphi_{21}) \to \varphi_{16}$\lbl{A1}}
\BIC{$(E \subseteq \mathbb{R} \to \varphi_{21}) \to \varphi_{16}$}
\end{ptree}
\noindent$\Pi_{26}$:
\begin{ptree}
\AXC{$\stackrel{\Pi_{16}}{\vdots}$}\noLine
\UIC{$\varphi_{19}$}
\AXC{$\stackrel{\Pi_{17}}{\vdots}$}\noLine
\UIC{$\mathrm{UB}(s - e,E) \leftrightarrow s - e \in \mathbb{R} \land \forall x \in E\,(x \leq s - e)$}
\AXC{$\stackrel{\Pi_{22}}{\vdots}$}\noLine
\UIC{$E \subseteq \mathbb{R} \to \varphi_{21}$}
\AXC{$\stackrel{\Pi_{25}}{\vdots}$}\noLine
\UIC{$(E \subseteq \mathbb{R} \to \varphi_{21}) \to \varphi_{16}$}
\BIC{$\varphi_{16}$}
\AXC{$\varphi_{16} \to \varphi_{49}$\lbl{A1}}
\BIC{$\varphi_{49}$}
\BIC{$\varphi_{19} \to E \subseteq \mathbb{R} \to \varphi_{11}$}
\BIC{$E \subseteq \mathbb{R} \to \varphi_{11}$}
\end{ptree}
\noindent$\Pi_{27}$:
\begin{ptree}
\AXC{$\varphi_{8}$\lbl{A4}}
\AXC{$\varphi_{5} \to \varphi_{6}$\lbl{A3}}
\AXC{$\stackrel{\Pi_{14}}{\vdots}$}\noLine
\UIC{$\varphi_{5}$}
\AXC{$\stackrel{\Pi_{26}}{\vdots}$}\noLine
\UIC{$E \subseteq \mathbb{R} \to \varphi_{11}$}
\AXC{$(E \subseteq \mathbb{R} \to \varphi_{11}) \to \varphi_{5} \to (\varphi_{5} \to \varphi_{6}) \to \varphi_{8} \to E \subseteq \mathbb{R} \to \varphi_{12}$\lbl{A1}}
\BIC{$\varphi_{5} \to (\varphi_{5} \to \varphi_{6}) \to \varphi_{8} \to E \subseteq \mathbb{R} \to \varphi_{12}$}
\BIC{$(\varphi_{5} \to \varphi_{6}) \to \varphi_{8} \to E \subseteq \mathbb{R} \to \varphi_{12}$}
\BIC{$\varphi_{8} \to E \subseteq \mathbb{R} \to \varphi_{12}$}
\BIC{$E \subseteq \mathbb{R} \to \varphi_{12}$}
\end{ptree}
\noindent 主木:
\begin{ptree}
\AXC{$\varphi_{31} \to \forall x\,(\varphi_{31})$\lbl{A4}}
\AXC{$\varphi_{51}$\lbl{A3}}
\AXC{$\stackrel{\Pi_{13}}{\vdots}$}\noLine
\UIC{$\forall x\,(\varphi_{31} \to \neg \neg (x \in E \to x \leq s - e))$}
\AXC{$\stackrel{\Pi_{27}}{\vdots}$}\noLine
\UIC{$E \subseteq \mathbb{R} \to \varphi_{12}$}
\AXC{$(E \subseteq \mathbb{R} \to \varphi_{12}) \to \varphi_{57}$\lbl{A1}}
\BIC{$\varphi_{57}$}
\BIC{$\varphi_{51} \to (\varphi_{31} \to \forall x\,(\varphi_{31})) \to E \subseteq \mathbb{R} \to \varphi_{14}$}
\BIC{$(\varphi_{31} \to \forall x\,(\varphi_{31})) \to E \subseteq \mathbb{R} \to \varphi_{14}$}
\BIC{$E \subseteq \mathbb{R} \to \varphi_{14}$}
\UIC{$\forall e\,(E \subseteq \mathbb{R} \to \varphi_{14})$}
\UIC{$\forall E\,\forall e\,(E \subseteq \mathbb{R} \to \varphi_{14})$}
\UIC{$\forall s\,\forall E\,\forall e\,(E \subseteq \mathbb{R} \to \varphi_{14})$}
\end{ptree}
\begin{abbrevlist}
\item $\varphi_{1} :\equiv \mathrm{IsSup}(s,E) \leftrightarrow \mathrm{UB}(s,E) \land \forall w\,(\mathrm{UB}(w,E) \to s \leq w)$
\item $\varphi_{2} :\equiv \forall X\,(\mathrm{IsSup}(s,X) \leftrightarrow \mathrm{UB}(s,X) \land \forall w\,(\mathrm{UB}(w,X) \to s \leq w))$
\item $\varphi_{3} :\equiv \forall w\,(\mathrm{UB}(w,E) \to s \leq w) \to \mathrm{UB}(s - e,E) \to s \leq s - e$
\item $\varphi_{4} :\equiv \forall x\,(\neg \neg (x \in E \to x \leq s - e)) \to \neg \neg (x \in E \to x \leq s - e)$
\item $\varphi_{5} :\equiv \forall x\,(\forall x\,(\neg \neg (x \in E \to x \leq s - e)) \to x \in E \to x \leq s - e)$
\item $\varphi_{6} :\equiv \forall x\,\forall x\,(\neg \neg (x \in E \to x \leq s - e)) \to \forall x \in E\,(x \leq s - e)$
\item $\varphi_{7} :\equiv \mathrm{IsSup}(s,E) \to e \in \mathbb{R} \to 0 < e \to \forall w\,(\mathrm{UB}(w,E) \to s \leq w)$
\item $\varphi_{8} :\equiv \forall x\,(\neg \neg (x \in E \to x \leq s - e)) \to \forall x\,\forall x\,(\neg \neg (x \in E \to x \leq s - e))$
\item $\varphi_{9} :\equiv e \in \mathbb{R} \to 0 < e \to \neg \neg \mathrm{UB}(s - e,E) \to \neg \neg \mathrm{UB}(s - e,E)$
\item $\varphi_{10} :\equiv \forall X\,(\mathrm{UB}(s - e,X) \leftrightarrow s - e \in \mathbb{R} \land \forall x \in X\,(x \leq s - e))$
\item $\varphi_{11} :\equiv \mathrm{IsSup}(s,E) \to e \in \mathbb{R} \to 0 < e \to \neg \forall x \in E\,(x \leq s - e)$
\item $\varphi_{12} :\equiv \mathrm{IsSup}(s,E) \to e \in \mathbb{R} \to 0 < e \to \exists x\,(\neg (x \in E \to x \leq s - e))$
\item $\varphi_{13} :\equiv \mathrm{IsSup}(s,E) \to e \in \mathbb{R} \to 0 < e \to \mathrm{UB}(s - e,E) \to s \leq s - e$
\item $\varphi_{14} :\equiv \mathrm{IsSup}(s,E) \to e \in \mathbb{R} \to 0 < e \to \exists x \in E\,(s - e < x)$
\item $\varphi_{15} :\equiv \mathrm{IsSup}(s,E) \to e \in \mathbb{R} \to 0 < e \to \neg \neg \mathrm{UB}(s - e,E) \to s \leq s - e$
\item $\varphi_{16} :\equiv E \subseteq \mathbb{R} \to \mathrm{IsSup}(s,E) \to e \in \mathbb{R} \to 0 < e \to \neg \mathrm{UB}(s - e,E)$
\item $\varphi_{17} :\equiv e \in \mathbb{R} \to 0 < e \to \neg (x \in E \to x \leq s - e) \to \neg (x \in E \to x \leq s - e)$
\item $\varphi_{18} :\equiv \mathrm{IsSup}(s,E) \to e \in \mathbb{R} \to 0 < e \to \neg \neg \mathrm{UB}(s - e,E) \to s - e < s$
\item $\varphi_{19} :\equiv E \subseteq \mathbb{R} \to \mathrm{IsSup}(s,E) \to e \in \mathbb{R} \to 0 < e \to s - e \in \mathbb{R}$
\item $\varphi_{20} :\equiv \mathrm{IsSup}(s,E) \to e \in \mathbb{R} \to 0 < e \to \neg (x \in E \to x \leq s - e) \to x \in E$
\item $\varphi_{21} :\equiv \mathrm{IsSup}(s,E) \to e \in \mathbb{R} \to 0 < e \to \neg \neg \mathrm{UB}(s - e,E) \to \neg s - e < s$
\item $\varphi_{22} :\equiv e \in \mathbb{R} \to 0 < e \to \neg (x \in E \to x \leq s - e) \to x \in E \land s - e < x$
\item $\varphi_{23} :\equiv \mathrm{IsSup}(s,E) \to e \in \mathbb{R} \to 0 < e \to \neg (x \in E \to x \leq s - e) \to \neg x \leq s - e$
\item $\varphi_{24} :\equiv \mathrm{IsSup}(s,E) \to e \in \mathbb{R} \to 0 < e \to \neg (x \in E \to x \leq s - e) \to x \in \mathbb{R}$
\item $\varphi_{25} :\equiv \mathrm{IsSup}(s,E) \to e \in \mathbb{R} \to 0 < e \to \neg (x \in E \to x \leq s - e) \to s - e < x$
\item $\varphi_{26} :\equiv e \in \mathbb{R} \to 0 < e \to \neg (x \in E \to x \leq s - e) \to \exists x \in E\,(s - e < x)$
\item $\varphi_{27} :\equiv \mathrm{IsSup}(s,E) \to e \in \mathbb{R} \to 0 < e \to \neg (x \in E \to x \leq s - e) \to s - e \in \mathbb{R}$
\item $\varphi_{28} :\equiv \forall X\,\forall Y\,\forall x\,(X \subseteq Y \to x \in X \to x \in Y) \to \forall Y\,\forall x\,(E \subseteq Y \to x \in E \to x \in Y)$
\item $\varphi_{29} :\equiv \forall Y\,\forall x\,(E \subseteq Y \to x \in E \to x \in Y) \to \forall x\,(E \subseteq \mathbb{R} \to x \in E \to x \in \mathbb{R})$
\item $\varphi_{30} :\equiv \forall x\,(E \subseteq \mathbb{R} \to x \in E \to x \in \mathbb{R}) \to E \subseteq \mathbb{R} \to x \in E \to x \in \mathbb{R}$
\item $\varphi_{31} :\equiv E \subseteq \mathbb{R} \land \mathrm{IsSup}(s,E) \land e \in \mathbb{R} \land 0 < e \land \neg \exists x \in E\,(s - e < x)$
\item $\varphi_{32} :\equiv \forall X\,(\mathrm{UB}(s,X) \leftrightarrow s \in \mathbb{R} \land \forall x \in X\,(x \leq s)) \to (\mathrm{UB}(s,E) \leftrightarrow s \in \mathbb{R} \land \forall x \in E\,(x \leq s))$
\item $\varphi_{33} :\equiv \forall y\,(s \in \mathbb{R} \to y \in \mathbb{R} \to s - y \in \mathbb{R}) \to s \in \mathbb{R} \to e \in \mathbb{R} \to s - e \in \mathbb{R}$
\item $\varphi_{34} :\equiv \forall a\,\forall X\,(\mathrm{UB}(a,X) \leftrightarrow a \in \mathbb{R} \land \forall x \in X\,(x \leq a)) \to \forall X\,(\mathrm{UB}(s,X) \leftrightarrow s \in \mathbb{R} \land \forall x \in X\,(x \leq s))$
\item $\varphi_{35} :\equiv \forall x\,\forall y\,(x \in \mathbb{R} \to y \in \mathbb{R} \to x - y \in \mathbb{R}) \to \forall y\,(s \in \mathbb{R} \to y \in \mathbb{R} \to s - y \in \mathbb{R})$
\item $\varphi_{36} :\equiv \forall e\,(s \in \mathbb{R} \to e \in \mathbb{R} \to 0 < e \to s - e < s) \to s \in \mathbb{R} \to e \in \mathbb{R} \to 0 < e \to s - e < s$
\item $\varphi_{37} :\equiv (\mathrm{UB}(s,E) \leftrightarrow s \in \mathbb{R} \land \forall x \in E\,(x \leq s)) \to E \subseteq \mathbb{R} \to \mathrm{IsSup}(s,E) \to e \in \mathbb{R} \to 0 < e \to s \in \mathbb{R}$
\item $\varphi_{38} :\equiv \forall x\,\forall e\,(x \in \mathbb{R} \to e \in \mathbb{R} \to 0 < e \to x - e < x) \to \forall e\,(s \in \mathbb{R} \to e \in \mathbb{R} \to 0 < e \to s - e < s)$
\item $\varphi_{39} :\equiv \forall x\,\forall y\,(x \in \mathbb{R} \to y \in \mathbb{R} \to \neg y \leq x \to x < y) \to \forall y\,(s - e \in \mathbb{R} \to y \in \mathbb{R} \to \neg y \leq s - e \to s - e < y)$
\item $\varphi_{40} :\equiv \forall y\,(s - e \in \mathbb{R} \to y \in \mathbb{R} \to \neg y \leq s - e \to s - e < y) \to s - e \in \mathbb{R} \to x \in \mathbb{R} \to \neg x \leq s - e \to s - e < x$
\item $\varphi_{41} :\equiv \varphi_{10} \to (\mathrm{UB}(s - e,E) \leftrightarrow s - e \in \mathbb{R} \land \forall x \in E\,(x \leq s - e))$
\item $\varphi_{42} :\equiv (E \subseteq \mathbb{R} \to \mathrm{IsSup}(s,E) \to e \in \mathbb{R} \to 0 < e \to 0 < e) \to E \subseteq \mathbb{R} \to \varphi_{18}$
\item $\varphi_{43} :\equiv (E \subseteq \mathbb{R} \to \mathrm{IsSup}(s,E) \to e \in \mathbb{R} \to e \in \mathbb{R}) \to E \subseteq \mathbb{R} \to \varphi_{27}$
\item $\varphi_{44} :\equiv (E \subseteq \mathbb{R} \to \mathrm{IsSup}(s,E) \to \mathrm{IsSup}(s,E)) \to \varphi_{1} \to E \subseteq \mathbb{R} \to \varphi_{7}$
\item $\varphi_{45} :\equiv (\forall x\,(\neg (x \in E \land s - e < x)) \to \neg (x \in E \land s - e < x)) \to E \subseteq \mathbb{R} \to \mathrm{IsSup}(s,E) \to \varphi_{26}$
\item $\varphi_{46} :\equiv (E \subseteq \mathbb{R} \to \mathrm{IsSup}(s,E) \to e \in \mathbb{R} \to 0 < e \to s \in \mathbb{R}) \to (E \subseteq \mathbb{R} \to \mathrm{IsSup}(s,E) \to e \in \mathbb{R} \to e \in \mathbb{R}) \to \varphi_{19}$
\item $\varphi_{47} :\equiv (E \subseteq \mathbb{R} \to \mathrm{IsSup}(s,E) \to \mathrm{IsSup}(s,E)) \to \varphi_{1} \to \varphi_{37}$
\item $\varphi_{48} :\equiv (E \subseteq \mathbb{R} \to \mathrm{IsSup}(s,E) \to \varphi_{26}) \to \varphi_{31} \to \neg \neg (x \in E \to x \leq s - e)$
\item $\varphi_{49} :\equiv (\mathrm{UB}(s - e,E) \leftrightarrow s - e \in \mathbb{R} \land \forall x \in E\,(x \leq s - e)) \to \varphi_{19} \to E \subseteq \mathbb{R} \to \varphi_{11}$
\item $\varphi_{50} :\equiv (E \subseteq \mathbb{R} \to \mathrm{IsSup}(s,E) \to e \in \mathbb{R} \to 0 < e \to s \in \mathbb{R}) \to \varphi_{43}$
\item $\varphi_{51} :\equiv \forall x\,(\varphi_{31} \to \neg \neg (x \in E \to x \leq s - e)) \to \forall x\,(\varphi_{31}) \to \forall x\,(\neg \neg (x \in E \to x \leq s - e))$
\item $\varphi_{52} :\equiv (E \subseteq \mathbb{R} \to x \in E \to x \in \mathbb{R}) \to (E \subseteq \mathbb{R} \to E \subseteq \mathbb{R}) \to (E \subseteq \mathbb{R} \to \varphi_{20}) \to E \subseteq \mathbb{R} \to \varphi_{24}$
\item $\varphi_{53} :\equiv (E \subseteq \mathbb{R} \to \varphi_{13}) \to (E \subseteq \mathbb{R} \to \mathrm{IsSup}(s,E) \to \varphi_{9}) \to E \subseteq \mathbb{R} \to \varphi_{15}$
\item $\varphi_{54} :\equiv (E \subseteq \mathbb{R} \to \mathrm{IsSup}(s,E) \to e \in \mathbb{R} \to 0 < e \to s \in \mathbb{R}) \to (E \subseteq \mathbb{R} \to \mathrm{IsSup}(s,E) \to e \in \mathbb{R} \to e \in \mathbb{R}) \to \varphi_{42}$
\item $\varphi_{55} :\equiv (E \subseteq \mathbb{R} \to \varphi_{20}) \to (E \subseteq \mathbb{R} \to \varphi_{25}) \to E \subseteq \mathbb{R} \to \mathrm{IsSup}(s,E) \to \varphi_{22}$
\item $\varphi_{56} :\equiv (E \subseteq \mathbb{R} \to \varphi_{27}) \to (E \subseteq \mathbb{R} \to \varphi_{24}) \to (E \subseteq \mathbb{R} \to \varphi_{23}) \to E \subseteq \mathbb{R} \to \varphi_{25}$
\item $\varphi_{57} :\equiv \forall x\,(\varphi_{31} \to \neg \neg (x \in E \to x \leq s - e)) \to \varphi_{51} \to (\varphi_{31} \to \forall x\,(\varphi_{31})) \to E \subseteq \mathbb{R} \to \varphi_{14}$
\end{abbrevlist}
\end{fullproof}

% 自動生成 (prover/build.js): InfApprox  ノード数 139
\begin{fullproof}{[C23]}{$\forall s\,\forall E\,\forall e\,(E \subseteq \mathbb{R} \to \mathrm{IsInf}(s,E) \to e \in \mathbb{R} \to 0 < e \to \exists x \in E\,(x < s + e))$}{139}\phantomsection\label{tree:InfApprox}
\noindent$\Pi_{1}$:
\begin{ptree}
\AXC{$E \subseteq \mathbb{R} \to \mathrm{IsInf}(s,E) \to \varphi_{18}$\lbl{A1}}
\AXC{$(E \subseteq \mathbb{R} \to \mathrm{IsInf}(s,E) \to \varphi_{18}) \to E \subseteq \mathbb{R} \to \varphi_{24}$\lbl{A1}}
\BIC{$E \subseteq \mathbb{R} \to \varphi_{24}$}
\end{ptree}
\noindent$\Pi_{2}$:
\begin{ptree}
\AXC{$\forall s\,(\varphi_{2})$\lbl{DefInf}}
\AXC{$\forall s\,(\varphi_{2}) \to \varphi_{2}$\lbl{A2}}
\BIC{$\varphi_{2}$}
\AXC{$\varphi_{2} \to \varphi_{1}$\lbl{A2}}
\BIC{$\varphi_{1}$}
\AXC{$E \subseteq \mathbb{R} \to \mathrm{IsInf}(s,E) \to \mathrm{IsInf}(s,E)$\lbl{A1}}
\AXC{$\varphi_{48}$\lbl{A1}}
\BIC{$\varphi_{1} \to \varphi_{40}$}
\BIC{$\varphi_{40}$}
\end{ptree}
\noindent$\Pi_{3}$:
\begin{ptree}
\AXC{$\forall a\,\forall X\,(\mathrm{LB}(a,X) \leftrightarrow a \in \mathbb{R} \land \forall x \in X\,(a \leq x))$\lbl{DefLB}}
\AXC{$\varphi_{35}$\lbl{A2}}
\BIC{$\forall X\,(\mathrm{LB}(s,X) \leftrightarrow s \in \mathbb{R} \land \forall x \in X\,(s \leq x))$}
\AXC{$\varphi_{33}$\lbl{A2}}
\BIC{$\mathrm{LB}(s,E) \leftrightarrow s \in \mathbb{R} \land \forall x \in E\,(s \leq x)$}
\AXC{$\stackrel{\Pi_{2}}{\vdots}$}\noLine
\UIC{$\varphi_{40}$}
\BIC{$E \subseteq \mathbb{R} \to \mathrm{IsInf}(s,E) \to e \in \mathbb{R} \to 0 < e \to s \in \mathbb{R}$}
\end{ptree}
\noindent$\Pi_{4}$:
\begin{ptree}
\AXC{$\forall x\,\forall y\,(x \in \mathbb{R} \to y \in \mathbb{R} \to x + y \in \mathbb{R})$\lbl{(R1)}}
\AXC{$\varphi_{37}$\lbl{A2}}
\BIC{$\forall y\,(s \in \mathbb{R} \to y \in \mathbb{R} \to s + y \in \mathbb{R})$}
\AXC{$\varphi_{34}$\lbl{A2}}
\BIC{$s \in \mathbb{R} \to e \in \mathbb{R} \to s + e \in \mathbb{R}$}
\AXC{$(s \in \mathbb{R} \to e \in \mathbb{R} \to s + e \in \mathbb{R}) \to \varphi_{51}$\lbl{A1}}
\BIC{$\varphi_{51}$}
\end{ptree}
\noindent$\Pi_{5}$:
\begin{ptree}
\AXC{$E \subseteq \mathbb{R} \to \mathrm{IsInf}(s,E) \to e \in \mathbb{R} \to e \in \mathbb{R}$\lbl{A1}}
\AXC{$\stackrel{\Pi_{3}}{\vdots}$}\noLine
\UIC{$E \subseteq \mathbb{R} \to \mathrm{IsInf}(s,E) \to e \in \mathbb{R} \to 0 < e \to s \in \mathbb{R}$}
\AXC{$\stackrel{\Pi_{4}}{\vdots}$}\noLine
\UIC{$\varphi_{51}$}
\BIC{$\varphi_{44}$}
\BIC{$E \subseteq \mathbb{R} \to \varphi_{28}$}
\end{ptree}
\noindent$\Pi_{6}$:
\begin{ptree}
\AXC{$E \subseteq \mathbb{R} \to E \subseteq \mathbb{R}$\lbl{A1}}
\AXC{$\forall X\,\forall Y\,\forall x\,(X \subseteq Y \to x \in X \to x \in Y)$\lbl{[B5]}}
\AXC{$\varphi_{29}$\lbl{A2}}
\BIC{$\forall Y\,\forall x\,(E \subseteq Y \to x \in E \to x \in Y)$}
\AXC{$\varphi_{30}$\lbl{A2}}
\BIC{$\forall x\,(E \subseteq \mathbb{R} \to x \in E \to x \in \mathbb{R})$}
\AXC{$\varphi_{31}$\lbl{A2}}
\BIC{$E \subseteq \mathbb{R} \to x \in E \to x \in \mathbb{R}$}
\AXC{$\varphi_{53}$\lbl{A1}}
\BIC{$(E \subseteq \mathbb{R} \to E \subseteq \mathbb{R}) \to (E \subseteq \mathbb{R} \to \varphi_{21}) \to E \subseteq \mathbb{R} \to \varphi_{25}$}
\BIC{$(E \subseteq \mathbb{R} \to \varphi_{21}) \to E \subseteq \mathbb{R} \to \varphi_{25}$}
\end{ptree}
\noindent$\Pi_{7}$:
\begin{ptree}
\AXC{$E \subseteq \mathbb{R} \to \mathrm{IsInf}(s,E) \to \varphi_{18}$\lbl{A1}}
\AXC{$(E \subseteq \mathbb{R} \to \mathrm{IsInf}(s,E) \to \varphi_{18}) \to E \subseteq \mathbb{R} \to \varphi_{21}$\lbl{A1}}
\BIC{$E \subseteq \mathbb{R} \to \varphi_{21}$}
\AXC{$\stackrel{\Pi_{6}}{\vdots}$}\noLine
\UIC{$(E \subseteq \mathbb{R} \to \varphi_{21}) \to E \subseteq \mathbb{R} \to \varphi_{25}$}
\BIC{$E \subseteq \mathbb{R} \to \varphi_{25}$}
\end{ptree}
\noindent$\Pi_{8}$:
\begin{ptree}
\AXC{$\forall x\,\forall y\,(x \in \mathbb{R} \to y \in \mathbb{R} \to \neg y \leq x \to x < y)$\lbl{[A2]}}
\AXC{$\varphi_{36}$\lbl{A2}}
\BIC{$\forall y\,(x \in \mathbb{R} \to y \in \mathbb{R} \to \neg y \leq x \to x < y)$}
\AXC{$\varphi_{39}$\lbl{A2}}
\BIC{$x \in \mathbb{R} \to s + e \in \mathbb{R} \to \neg s + e \leq x \to x < s + e$}
\end{ptree}
\noindent$\Pi_{9}$:
\begin{ptree}
\AXC{$\stackrel{\Pi_{7}}{\vdots}$}\noLine
\UIC{$E \subseteq \mathbb{R} \to \varphi_{25}$}
\AXC{$\stackrel{\Pi_{8}}{\vdots}$}\noLine
\UIC{$x \in \mathbb{R} \to s + e \in \mathbb{R} \to \neg s + e \leq x \to x < s + e$}
\AXC{$(x \in \mathbb{R} \to s + e \in \mathbb{R} \to \neg s + e \leq x \to x < s + e) \to \varphi_{57}$\lbl{A1}}
\BIC{$\varphi_{57}$}
\BIC{$(E \subseteq \mathbb{R} \to \varphi_{28}) \to (E \subseteq \mathbb{R} \to \varphi_{24}) \to E \subseteq \mathbb{R} \to \varphi_{26}$}
\end{ptree}
\noindent$\Pi_{10}$:
\begin{ptree}
\AXC{$\stackrel{\Pi_{1}}{\vdots}$}\noLine
\UIC{$E \subseteq \mathbb{R} \to \varphi_{24}$}
\AXC{$\stackrel{\Pi_{5}}{\vdots}$}\noLine
\UIC{$E \subseteq \mathbb{R} \to \varphi_{28}$}
\AXC{$\stackrel{\Pi_{9}}{\vdots}$}\noLine
\UIC{$(E \subseteq \mathbb{R} \to \varphi_{28}) \to (E \subseteq \mathbb{R} \to \varphi_{24}) \to E \subseteq \mathbb{R} \to \varphi_{26}$}
\BIC{$(E \subseteq \mathbb{R} \to \varphi_{24}) \to E \subseteq \mathbb{R} \to \varphi_{26}$}
\BIC{$E \subseteq \mathbb{R} \to \varphi_{26}$}
\end{ptree}
\noindent$\Pi_{11}$:
\begin{ptree}
\AXC{$E \subseteq \mathbb{R} \to \mathrm{IsInf}(s,E) \to \varphi_{18}$\lbl{A1}}
\AXC{$(E \subseteq \mathbb{R} \to \mathrm{IsInf}(s,E) \to \varphi_{18}) \to E \subseteq \mathbb{R} \to \varphi_{21}$\lbl{A1}}
\BIC{$E \subseteq \mathbb{R} \to \varphi_{21}$}
\AXC{$\varphi_{56}$\lbl{A1}}
\BIC{$(E \subseteq \mathbb{R} \to \varphi_{26}) \to E \subseteq \mathbb{R} \to \mathrm{IsInf}(s,E) \to \varphi_{23}$}
\end{ptree}
\noindent$\Pi_{12}$:
\begin{ptree}
\AXC{$\stackrel{\Pi_{10}}{\vdots}$}\noLine
\UIC{$E \subseteq \mathbb{R} \to \varphi_{26}$}
\AXC{$\stackrel{\Pi_{11}}{\vdots}$}\noLine
\UIC{$(E \subseteq \mathbb{R} \to \varphi_{26}) \to E \subseteq \mathbb{R} \to \mathrm{IsInf}(s,E) \to \varphi_{23}$}
\BIC{$E \subseteq \mathbb{R} \to \mathrm{IsInf}(s,E) \to \varphi_{23}$}
\AXC{$(E \subseteq \mathbb{R} \to \mathrm{IsInf}(s,E) \to \varphi_{23}) \to \varphi_{46}$\lbl{A1}}
\BIC{$\varphi_{46}$}
\end{ptree}
\noindent$\Pi_{13}$:
\begin{ptree}
\AXC{$\forall x\,(\neg (x \in E \land x < s + e)) \to \neg (x \in E \land x < s + e)$\lbl{A2}}
\AXC{$\stackrel{\Pi_{12}}{\vdots}$}\noLine
\UIC{$\varphi_{46}$}
\BIC{$E \subseteq \mathbb{R} \to \mathrm{IsInf}(s,E) \to \varphi_{27}$}
\AXC{$\varphi_{49}$\lbl{A1}}
\BIC{$\varphi_{32} \to \neg \neg (x \in E \to s + e \leq x)$}
\UIC{$\forall x\,(\varphi_{32} \to \neg \neg (x \in E \to s + e \leq x))$}
\end{ptree}
\noindent$\Pi_{14}$:
\begin{ptree}
\AXC{$\varphi_{4}$\lbl{A2}}
\AXC{$\varphi_{4} \to \forall x\,(\neg \neg (x \in E \to s + e \leq x)) \to x \in E \to s + e \leq x$\lbl{A1}}
\BIC{$\forall x\,(\neg \neg (x \in E \to s + e \leq x)) \to x \in E \to s + e \leq x$}
\UIC{$\varphi_{5}$}
\end{ptree}
\noindent$\Pi_{15}$:
\begin{ptree}
\AXC{$\forall x\,\forall y\,(x \in \mathbb{R} \to y \in \mathbb{R} \to x + y \in \mathbb{R})$\lbl{(R1)}}
\AXC{$\varphi_{37}$\lbl{A2}}
\BIC{$\forall y\,(s \in \mathbb{R} \to y \in \mathbb{R} \to s + y \in \mathbb{R})$}
\AXC{$\varphi_{34}$\lbl{A2}}
\BIC{$s \in \mathbb{R} \to e \in \mathbb{R} \to s + e \in \mathbb{R}$}
\AXC{$(s \in \mathbb{R} \to e \in \mathbb{R} \to s + e \in \mathbb{R}) \to \varphi_{47}$\lbl{A1}}
\BIC{$\varphi_{47}$}
\end{ptree}
\noindent$\Pi_{16}$:
\begin{ptree}
\AXC{$E \subseteq \mathbb{R} \to \mathrm{IsInf}(s,E) \to e \in \mathbb{R} \to e \in \mathbb{R}$\lbl{A1}}
\AXC{$\stackrel{\Pi_{3}}{\vdots}$}\noLine
\UIC{$E \subseteq \mathbb{R} \to \mathrm{IsInf}(s,E) \to e \in \mathbb{R} \to 0 < e \to s \in \mathbb{R}$}
\AXC{$\stackrel{\Pi_{15}}{\vdots}$}\noLine
\UIC{$\varphi_{47}$}
\BIC{$(E \subseteq \mathbb{R} \to \mathrm{IsInf}(s,E) \to e \in \mathbb{R} \to e \in \mathbb{R}) \to \varphi_{20}$}
\BIC{$\varphi_{20}$}
\end{ptree}
\noindent$\Pi_{17}$:
\begin{ptree}
\AXC{$\forall a\,\forall X\,(\mathrm{LB}(a,X) \leftrightarrow a \in \mathbb{R} \land \forall x \in X\,(a \leq x))$\lbl{DefLB}}
\AXC{$\forall a\,\forall X\,(\mathrm{LB}(a,X) \leftrightarrow a \in \mathbb{R} \land \forall x \in X\,(a \leq x)) \to \varphi_{11}$\lbl{A2}}
\BIC{$\varphi_{11}$}
\AXC{$\varphi_{42}$\lbl{A2}}
\BIC{$\mathrm{LB}(s + e,E) \leftrightarrow s + e \in \mathbb{R} \land \forall x \in E\,(s + e \leq x)$}
\end{ptree}
\noindent$\Pi_{18}$:
\begin{ptree}
\AXC{$\forall s\,(\varphi_{2})$\lbl{DefInf}}
\AXC{$\forall s\,(\varphi_{2}) \to \varphi_{2}$\lbl{A2}}
\BIC{$\varphi_{2}$}
\AXC{$\varphi_{2} \to \varphi_{1}$\lbl{A2}}
\BIC{$\varphi_{1}$}
\AXC{$E \subseteq \mathbb{R} \to \mathrm{IsInf}(s,E) \to \mathrm{IsInf}(s,E)$\lbl{A1}}
\AXC{$\varphi_{45}$\lbl{A1}}
\BIC{$\varphi_{1} \to E \subseteq \mathbb{R} \to \varphi_{8}$}
\BIC{$E \subseteq \mathbb{R} \to \varphi_{8}$}
\end{ptree}
\noindent$\Pi_{19}$:
\begin{ptree}
\AXC{$E \subseteq \mathbb{R} \to \mathrm{IsInf}(s,E) \to \varphi_{10}$\lbl{A1}}
\AXC{$\varphi_{3}$\lbl{A2}}
\AXC{$\stackrel{\Pi_{18}}{\vdots}$}\noLine
\UIC{$E \subseteq \mathbb{R} \to \varphi_{8}$}
\AXC{$(E \subseteq \mathbb{R} \to \varphi_{8}) \to \varphi_{3} \to E \subseteq \mathbb{R} \to \varphi_{14}$\lbl{A1}}
\BIC{$\varphi_{3} \to E \subseteq \mathbb{R} \to \varphi_{14}$}
\BIC{$E \subseteq \mathbb{R} \to \varphi_{14}$}
\AXC{$\varphi_{54}$\lbl{A1}}
\BIC{$(E \subseteq \mathbb{R} \to \mathrm{IsInf}(s,E) \to \varphi_{10}) \to E \subseteq \mathbb{R} \to \varphi_{16}$}
\BIC{$E \subseteq \mathbb{R} \to \varphi_{16}$}
\end{ptree}
\noindent$\Pi_{20}$:
\begin{ptree}
\AXC{$\forall x\,\forall y\,(x \leq y \to \neg y < x)$\lbl{[O7]}}
\AXC{$\varphi_{6}$\lbl{A2}}
\BIC{$\forall y\,(s + e \leq y \to \neg y < s + e)$}
\AXC{$\forall y\,(s + e \leq y \to \neg y < s + e) \to s + e \leq s \to \neg s < s + e$\lbl{A2}}
\BIC{$s + e \leq s \to \neg s < s + e$}
\end{ptree}
\noindent$\Pi_{21}$:
\begin{ptree}
\AXC{$\stackrel{\Pi_{19}}{\vdots}$}\noLine
\UIC{$E \subseteq \mathbb{R} \to \varphi_{16}$}
\AXC{$\stackrel{\Pi_{20}}{\vdots}$}\noLine
\UIC{$s + e \leq s \to \neg s < s + e$}
\AXC{$(s + e \leq s \to \neg s < s + e) \to (E \subseteq \mathbb{R} \to \varphi_{16}) \to E \subseteq \mathbb{R} \to \varphi_{22}$\lbl{A1}}
\BIC{$(E \subseteq \mathbb{R} \to \varphi_{16}) \to E \subseteq \mathbb{R} \to \varphi_{22}$}
\BIC{$E \subseteq \mathbb{R} \to \varphi_{22}$}
\end{ptree}
\noindent$\Pi_{22}$:
\begin{ptree}
\AXC{$\forall x\,\forall y\,(x \in \mathbb{R} \to y \in \mathbb{R} \to 0 < y \to x < x + y)$\lbl{[E3]}}
\AXC{$\varphi_{41}$\lbl{A2}}
\BIC{$\forall y\,(s \in \mathbb{R} \to y \in \mathbb{R} \to 0 < y \to s < s + y)$}
\AXC{$\varphi_{38}$\lbl{A2}}
\BIC{$s \in \mathbb{R} \to e \in \mathbb{R} \to 0 < e \to s < s + e$}
\AXC{$(s \in \mathbb{R} \to e \in \mathbb{R} \to 0 < e \to s < s + e) \to \varphi_{55}$\lbl{A1}}
\BIC{$\varphi_{55}$}
\end{ptree}
\noindent$\Pi_{23}$:
\begin{ptree}
\AXC{$E \subseteq \mathbb{R} \to \mathrm{IsInf}(s,E) \to e \in \mathbb{R} \to e \in \mathbb{R}$\lbl{A1}}
\AXC{$\stackrel{\Pi_{3}}{\vdots}$}\noLine
\UIC{$E \subseteq \mathbb{R} \to \mathrm{IsInf}(s,E) \to e \in \mathbb{R} \to 0 < e \to s \in \mathbb{R}$}
\AXC{$\stackrel{\Pi_{22}}{\vdots}$}\noLine
\UIC{$\varphi_{55}$}
\BIC{$(E \subseteq \mathbb{R} \to \mathrm{IsInf}(s,E) \to e \in \mathbb{R} \to e \in \mathbb{R}) \to \varphi_{43}$}
\BIC{$\varphi_{43}$}
\end{ptree}
\noindent$\Pi_{24}$:
\begin{ptree}
\AXC{$E \subseteq \mathbb{R} \to \mathrm{IsInf}(s,E) \to e \in \mathbb{R} \to 0 < e \to 0 < e$\lbl{A1}}
\AXC{$\stackrel{\Pi_{23}}{\vdots}$}\noLine
\UIC{$\varphi_{43}$}
\BIC{$E \subseteq \mathbb{R} \to \varphi_{19}$}
\AXC{$(E \subseteq \mathbb{R} \to \varphi_{19}) \to (E \subseteq \mathbb{R} \to \varphi_{22}) \to \varphi_{17}$\lbl{A1}}
\BIC{$(E \subseteq \mathbb{R} \to \varphi_{22}) \to \varphi_{17}$}
\end{ptree}
\noindent$\Pi_{25}$:
\begin{ptree}
\AXC{$\stackrel{\Pi_{16}}{\vdots}$}\noLine
\UIC{$\varphi_{20}$}
\AXC{$\stackrel{\Pi_{17}}{\vdots}$}\noLine
\UIC{$\mathrm{LB}(s + e,E) \leftrightarrow s + e \in \mathbb{R} \land \forall x \in E\,(s + e \leq x)$}
\AXC{$\stackrel{\Pi_{21}}{\vdots}$}\noLine
\UIC{$E \subseteq \mathbb{R} \to \varphi_{22}$}
\AXC{$\stackrel{\Pi_{24}}{\vdots}$}\noLine
\UIC{$(E \subseteq \mathbb{R} \to \varphi_{22}) \to \varphi_{17}$}
\BIC{$\varphi_{17}$}
\AXC{$\varphi_{17} \to \varphi_{50}$\lbl{A1}}
\BIC{$\varphi_{50}$}
\BIC{$\varphi_{20} \to E \subseteq \mathbb{R} \to \varphi_{12}$}
\BIC{$E \subseteq \mathbb{R} \to \varphi_{12}$}
\end{ptree}
\noindent$\Pi_{26}$:
\begin{ptree}
\AXC{$\varphi_{9}$\lbl{A4}}
\AXC{$\varphi_{5} \to \varphi_{7}$\lbl{A3}}
\AXC{$\stackrel{\Pi_{14}}{\vdots}$}\noLine
\UIC{$\varphi_{5}$}
\AXC{$\stackrel{\Pi_{25}}{\vdots}$}\noLine
\UIC{$E \subseteq \mathbb{R} \to \varphi_{12}$}
\AXC{$(E \subseteq \mathbb{R} \to \varphi_{12}) \to \varphi_{5} \to (\varphi_{5} \to \varphi_{7}) \to \varphi_{9} \to E \subseteq \mathbb{R} \to \varphi_{13}$\lbl{A1}}
\BIC{$\varphi_{5} \to (\varphi_{5} \to \varphi_{7}) \to \varphi_{9} \to E \subseteq \mathbb{R} \to \varphi_{13}$}
\BIC{$(\varphi_{5} \to \varphi_{7}) \to \varphi_{9} \to E \subseteq \mathbb{R} \to \varphi_{13}$}
\BIC{$\varphi_{9} \to E \subseteq \mathbb{R} \to \varphi_{13}$}
\BIC{$E \subseteq \mathbb{R} \to \varphi_{13}$}
\end{ptree}
\noindent 主木:
\begin{ptree}
\AXC{$\varphi_{32} \to \forall x\,(\varphi_{32})$\lbl{A4}}
\AXC{$\varphi_{52}$\lbl{A3}}
\AXC{$\stackrel{\Pi_{13}}{\vdots}$}\noLine
\UIC{$\forall x\,(\varphi_{32} \to \neg \neg (x \in E \to s + e \leq x))$}
\AXC{$\stackrel{\Pi_{26}}{\vdots}$}\noLine
\UIC{$E \subseteq \mathbb{R} \to \varphi_{13}$}
\AXC{$(E \subseteq \mathbb{R} \to \varphi_{13}) \to \varphi_{58}$\lbl{A1}}
\BIC{$\varphi_{58}$}
\BIC{$\varphi_{52} \to (\varphi_{32} \to \forall x\,(\varphi_{32})) \to E \subseteq \mathbb{R} \to \varphi_{15}$}
\BIC{$(\varphi_{32} \to \forall x\,(\varphi_{32})) \to E \subseteq \mathbb{R} \to \varphi_{15}$}
\BIC{$E \subseteq \mathbb{R} \to \varphi_{15}$}
\UIC{$\forall e\,(E \subseteq \mathbb{R} \to \varphi_{15})$}
\UIC{$\forall E\,\forall e\,(E \subseteq \mathbb{R} \to \varphi_{15})$}
\UIC{$\forall s\,\forall E\,\forall e\,(E \subseteq \mathbb{R} \to \varphi_{15})$}
\end{ptree}
\begin{abbrevlist}
\item $\varphi_{1} :\equiv \mathrm{IsInf}(s,E) \leftrightarrow \mathrm{LB}(s,E) \land \forall w\,(\mathrm{LB}(w,E) \to w \leq s)$
\item $\varphi_{2} :\equiv \forall X\,(\mathrm{IsInf}(s,X) \leftrightarrow \mathrm{LB}(s,X) \land \forall w\,(\mathrm{LB}(w,X) \to w \leq s))$
\item $\varphi_{3} :\equiv \forall w\,(\mathrm{LB}(w,E) \to w \leq s) \to \mathrm{LB}(s + e,E) \to s + e \leq s$
\item $\varphi_{4} :\equiv \forall x\,(\neg \neg (x \in E \to s + e \leq x)) \to \neg \neg (x \in E \to s + e \leq x)$
\item $\varphi_{5} :\equiv \forall x\,(\forall x\,(\neg \neg (x \in E \to s + e \leq x)) \to x \in E \to s + e \leq x)$
\item $\varphi_{6} :\equiv \forall x\,\forall y\,(x \leq y \to \neg y < x) \to \forall y\,(s + e \leq y \to \neg y < s + e)$
\item $\varphi_{7} :\equiv \forall x\,\forall x\,(\neg \neg (x \in E \to s + e \leq x)) \to \forall x \in E\,(s + e \leq x)$
\item $\varphi_{8} :\equiv \mathrm{IsInf}(s,E) \to e \in \mathbb{R} \to 0 < e \to \forall w\,(\mathrm{LB}(w,E) \to w \leq s)$
\item $\varphi_{9} :\equiv \forall x\,(\neg \neg (x \in E \to s + e \leq x)) \to \forall x\,\forall x\,(\neg \neg (x \in E \to s + e \leq x))$
\item $\varphi_{10} :\equiv e \in \mathbb{R} \to 0 < e \to \neg \neg \mathrm{LB}(s + e,E) \to \neg \neg \mathrm{LB}(s + e,E)$
\item $\varphi_{11} :\equiv \forall X\,(\mathrm{LB}(s + e,X) \leftrightarrow s + e \in \mathbb{R} \land \forall x \in X\,(s + e \leq x))$
\item $\varphi_{12} :\equiv \mathrm{IsInf}(s,E) \to e \in \mathbb{R} \to 0 < e \to \neg \forall x \in E\,(s + e \leq x)$
\item $\varphi_{13} :\equiv \mathrm{IsInf}(s,E) \to e \in \mathbb{R} \to 0 < e \to \exists x\,(\neg (x \in E \to s + e \leq x))$
\item $\varphi_{14} :\equiv \mathrm{IsInf}(s,E) \to e \in \mathbb{R} \to 0 < e \to \mathrm{LB}(s + e,E) \to s + e \leq s$
\item $\varphi_{15} :\equiv \mathrm{IsInf}(s,E) \to e \in \mathbb{R} \to 0 < e \to \exists x \in E\,(x < s + e)$
\item $\varphi_{16} :\equiv \mathrm{IsInf}(s,E) \to e \in \mathbb{R} \to 0 < e \to \neg \neg \mathrm{LB}(s + e,E) \to s + e \leq s$
\item $\varphi_{17} :\equiv E \subseteq \mathbb{R} \to \mathrm{IsInf}(s,E) \to e \in \mathbb{R} \to 0 < e \to \neg \mathrm{LB}(s + e,E)$
\item $\varphi_{18} :\equiv e \in \mathbb{R} \to 0 < e \to \neg (x \in E \to s + e \leq x) \to \neg (x \in E \to s + e \leq x)$
\item $\varphi_{19} :\equiv \mathrm{IsInf}(s,E) \to e \in \mathbb{R} \to 0 < e \to \neg \neg \mathrm{LB}(s + e,E) \to s < s + e$
\item $\varphi_{20} :\equiv E \subseteq \mathbb{R} \to \mathrm{IsInf}(s,E) \to e \in \mathbb{R} \to 0 < e \to s + e \in \mathbb{R}$
\item $\varphi_{21} :\equiv \mathrm{IsInf}(s,E) \to e \in \mathbb{R} \to 0 < e \to \neg (x \in E \to s + e \leq x) \to x \in E$
\item $\varphi_{22} :\equiv \mathrm{IsInf}(s,E) \to e \in \mathbb{R} \to 0 < e \to \neg \neg \mathrm{LB}(s + e,E) \to \neg s < s + e$
\item $\varphi_{23} :\equiv e \in \mathbb{R} \to 0 < e \to \neg (x \in E \to s + e \leq x) \to x \in E \land x < s + e$
\item $\varphi_{24} :\equiv \mathrm{IsInf}(s,E) \to e \in \mathbb{R} \to 0 < e \to \neg (x \in E \to s + e \leq x) \to \neg s + e \leq x$
\item $\varphi_{25} :\equiv \mathrm{IsInf}(s,E) \to e \in \mathbb{R} \to 0 < e \to \neg (x \in E \to s + e \leq x) \to x \in \mathbb{R}$
\item $\varphi_{26} :\equiv \mathrm{IsInf}(s,E) \to e \in \mathbb{R} \to 0 < e \to \neg (x \in E \to s + e \leq x) \to x < s + e$
\item $\varphi_{27} :\equiv e \in \mathbb{R} \to 0 < e \to \neg (x \in E \to s + e \leq x) \to \exists x \in E\,(x < s + e)$
\item $\varphi_{28} :\equiv \mathrm{IsInf}(s,E) \to e \in \mathbb{R} \to 0 < e \to \neg (x \in E \to s + e \leq x) \to s + e \in \mathbb{R}$
\item $\varphi_{29} :\equiv \forall X\,\forall Y\,\forall x\,(X \subseteq Y \to x \in X \to x \in Y) \to \forall Y\,\forall x\,(E \subseteq Y \to x \in E \to x \in Y)$
\item $\varphi_{30} :\equiv \forall Y\,\forall x\,(E \subseteq Y \to x \in E \to x \in Y) \to \forall x\,(E \subseteq \mathbb{R} \to x \in E \to x \in \mathbb{R})$
\item $\varphi_{31} :\equiv \forall x\,(E \subseteq \mathbb{R} \to x \in E \to x \in \mathbb{R}) \to E \subseteq \mathbb{R} \to x \in E \to x \in \mathbb{R}$
\item $\varphi_{32} :\equiv E \subseteq \mathbb{R} \land \mathrm{IsInf}(s,E) \land e \in \mathbb{R} \land 0 < e \land \neg \exists x \in E\,(x < s + e)$
\item $\varphi_{33} :\equiv \forall X\,(\mathrm{LB}(s,X) \leftrightarrow s \in \mathbb{R} \land \forall x \in X\,(s \leq x)) \to (\mathrm{LB}(s,E) \leftrightarrow s \in \mathbb{R} \land \forall x \in E\,(s \leq x))$
\item $\varphi_{34} :\equiv \forall y\,(s \in \mathbb{R} \to y \in \mathbb{R} \to s + y \in \mathbb{R}) \to s \in \mathbb{R} \to e \in \mathbb{R} \to s + e \in \mathbb{R}$
\item $\varphi_{35} :\equiv \forall a\,\forall X\,(\mathrm{LB}(a,X) \leftrightarrow a \in \mathbb{R} \land \forall x \in X\,(a \leq x)) \to \forall X\,(\mathrm{LB}(s,X) \leftrightarrow s \in \mathbb{R} \land \forall x \in X\,(s \leq x))$
\item $\varphi_{36} :\equiv \forall x\,\forall y\,(x \in \mathbb{R} \to y \in \mathbb{R} \to \neg y \leq x \to x < y) \to \forall y\,(x \in \mathbb{R} \to y \in \mathbb{R} \to \neg y \leq x \to x < y)$
\item $\varphi_{37} :\equiv \forall x\,\forall y\,(x \in \mathbb{R} \to y \in \mathbb{R} \to x + y \in \mathbb{R}) \to \forall y\,(s \in \mathbb{R} \to y \in \mathbb{R} \to s + y \in \mathbb{R})$
\item $\varphi_{38} :\equiv \forall y\,(s \in \mathbb{R} \to y \in \mathbb{R} \to 0 < y \to s < s + y) \to s \in \mathbb{R} \to e \in \mathbb{R} \to 0 < e \to s < s + e$
\item $\varphi_{39} :\equiv \forall y\,(x \in \mathbb{R} \to y \in \mathbb{R} \to \neg y \leq x \to x < y) \to x \in \mathbb{R} \to s + e \in \mathbb{R} \to \neg s + e \leq x \to x < s + e$
\item $\varphi_{40} :\equiv (\mathrm{LB}(s,E) \leftrightarrow s \in \mathbb{R} \land \forall x \in E\,(s \leq x)) \to E \subseteq \mathbb{R} \to \mathrm{IsInf}(s,E) \to e \in \mathbb{R} \to 0 < e \to s \in \mathbb{R}$
\item $\varphi_{41} :\equiv \forall x\,\forall y\,(x \in \mathbb{R} \to y \in \mathbb{R} \to 0 < y \to x < x + y) \to \forall y\,(s \in \mathbb{R} \to y \in \mathbb{R} \to 0 < y \to s < s + y)$
\item $\varphi_{42} :\equiv \varphi_{11} \to (\mathrm{LB}(s + e,E) \leftrightarrow s + e \in \mathbb{R} \land \forall x \in E\,(s + e \leq x))$
\item $\varphi_{43} :\equiv (E \subseteq \mathbb{R} \to \mathrm{IsInf}(s,E) \to e \in \mathbb{R} \to 0 < e \to 0 < e) \to E \subseteq \mathbb{R} \to \varphi_{19}$
\item $\varphi_{44} :\equiv (E \subseteq \mathbb{R} \to \mathrm{IsInf}(s,E) \to e \in \mathbb{R} \to e \in \mathbb{R}) \to E \subseteq \mathbb{R} \to \varphi_{28}$
\item $\varphi_{45} :\equiv (E \subseteq \mathbb{R} \to \mathrm{IsInf}(s,E) \to \mathrm{IsInf}(s,E)) \to \varphi_{1} \to E \subseteq \mathbb{R} \to \varphi_{8}$
\item $\varphi_{46} :\equiv (\forall x\,(\neg (x \in E \land x < s + e)) \to \neg (x \in E \land x < s + e)) \to E \subseteq \mathbb{R} \to \mathrm{IsInf}(s,E) \to \varphi_{27}$
\item $\varphi_{47} :\equiv (E \subseteq \mathbb{R} \to \mathrm{IsInf}(s,E) \to e \in \mathbb{R} \to 0 < e \to s \in \mathbb{R}) \to (E \subseteq \mathbb{R} \to \mathrm{IsInf}(s,E) \to e \in \mathbb{R} \to e \in \mathbb{R}) \to \varphi_{20}$
\item $\varphi_{48} :\equiv (E \subseteq \mathbb{R} \to \mathrm{IsInf}(s,E) \to \mathrm{IsInf}(s,E)) \to \varphi_{1} \to \varphi_{40}$
\item $\varphi_{49} :\equiv (E \subseteq \mathbb{R} \to \mathrm{IsInf}(s,E) \to \varphi_{27}) \to \varphi_{32} \to \neg \neg (x \in E \to s + e \leq x)$
\item $\varphi_{50} :\equiv (\mathrm{LB}(s + e,E) \leftrightarrow s + e \in \mathbb{R} \land \forall x \in E\,(s + e \leq x)) \to \varphi_{20} \to E \subseteq \mathbb{R} \to \varphi_{12}$
\item $\varphi_{51} :\equiv (E \subseteq \mathbb{R} \to \mathrm{IsInf}(s,E) \to e \in \mathbb{R} \to 0 < e \to s \in \mathbb{R}) \to \varphi_{44}$
\item $\varphi_{52} :\equiv \forall x\,(\varphi_{32} \to \neg \neg (x \in E \to s + e \leq x)) \to \forall x\,(\varphi_{32}) \to \forall x\,(\neg \neg (x \in E \to s + e \leq x))$
\item $\varphi_{53} :\equiv (E \subseteq \mathbb{R} \to x \in E \to x \in \mathbb{R}) \to (E \subseteq \mathbb{R} \to E \subseteq \mathbb{R}) \to (E \subseteq \mathbb{R} \to \varphi_{21}) \to E \subseteq \mathbb{R} \to \varphi_{25}$
\item $\varphi_{54} :\equiv (E \subseteq \mathbb{R} \to \varphi_{14}) \to (E \subseteq \mathbb{R} \to \mathrm{IsInf}(s,E) \to \varphi_{10}) \to E \subseteq \mathbb{R} \to \varphi_{16}$
\item $\varphi_{55} :\equiv (E \subseteq \mathbb{R} \to \mathrm{IsInf}(s,E) \to e \in \mathbb{R} \to 0 < e \to s \in \mathbb{R}) \to (E \subseteq \mathbb{R} \to \mathrm{IsInf}(s,E) \to e \in \mathbb{R} \to e \in \mathbb{R}) \to \varphi_{43}$
\item $\varphi_{56} :\equiv (E \subseteq \mathbb{R} \to \varphi_{21}) \to (E \subseteq \mathbb{R} \to \varphi_{26}) \to E \subseteq \mathbb{R} \to \mathrm{IsInf}(s,E) \to \varphi_{23}$
\item $\varphi_{57} :\equiv (E \subseteq \mathbb{R} \to \varphi_{25}) \to (E \subseteq \mathbb{R} \to \varphi_{28}) \to (E \subseteq \mathbb{R} \to \varphi_{24}) \to E \subseteq \mathbb{R} \to \varphi_{26}$
\item $\varphi_{58} :\equiv \forall x\,(\varphi_{32} \to \neg \neg (x \in E \to s + e \leq x)) \to \varphi_{52} \to (\varphi_{32} \to \forall x\,(\varphi_{32})) \to E \subseteq \mathbb{R} \to \varphi_{15}$
\end{abbrevlist}
\end{fullproof}

% 自動生成 (prover/build.js): MonoConv  ノード数 443
\begin{fullproof}{定理5.3(1)}{$\forall a\,\forall M\,(\mathrm{MonoInc}(a) \to M \in \mathbb{R} \to \forall n \in \mathbb{N}\,(a_{n} \leq M) \to \exists \alpha\,(\mathrm{Conv}(a,\alpha)))$}{443}\phantomsection\label{tree:MonoConv}
\noindent$\Pi_{1}$:
\begin{ptree}
\AXC{$\forall a\,\forall X\,(\mathrm{UB}(a,X) \leftrightarrow a \in \mathbb{R} \land \forall x \in X\,(x \leq a))$\lbl{DefUB}}
\AXC{$\varphi_{107}$\lbl{A2}}
\BIC{$\forall X\,(\mathrm{UB}(s,X) \leftrightarrow s \in \mathbb{R} \land \forall x \in X\,(x \leq s))$}
\AXC{$\forall X\,(\mathrm{UB}(s,X) \leftrightarrow s \in \mathbb{R} \land \forall x \in X\,(x \leq s)) \to \varphi_{7}$\lbl{A2}}
\BIC{$\varphi_{7}$}
\end{ptree}
\noindent$\Pi_{2}$:
\begin{ptree}
\AXC{$\mathrm{MonoInc}(a) \to M \in \mathbb{R} \to \forall n \in \mathbb{N}\,(a_{n} \leq M) \to \varphi_{2}$\lbl{A1}}
\AXC{$\varphi_{215}$\lbl{A1}}
\BIC{$\varphi_{9} \to \varphi_{7} \to \mathrm{MonoInc}(a) \to M \in \mathbb{R} \to \varphi_{16}$}
\end{ptree}
\noindent$\Pi_{3}$:
\begin{ptree}
\AXC{$\stackrel{\Pi_{1}}{\vdots}$}\noLine
\UIC{$\varphi_{7}$}
\AXC{$\forall s\,(\varphi_{4})$\lbl{DefSup}}
\AXC{$\forall s\,(\varphi_{4}) \to \varphi_{4}$\lbl{A2}}
\BIC{$\varphi_{4}$}
\AXC{$\varphi_{4} \to \varphi_{9}$\lbl{A2}}
\BIC{$\varphi_{9}$}
\AXC{$\stackrel{\Pi_{2}}{\vdots}$}\noLine
\UIC{$\varphi_{9} \to \varphi_{7} \to \mathrm{MonoInc}(a) \to M \in \mathbb{R} \to \varphi_{16}$}
\BIC{$\varphi_{7} \to \mathrm{MonoInc}(a) \to M \in \mathbb{R} \to \varphi_{16}$}
\BIC{$\mathrm{MonoInc}(a) \to M \in \mathbb{R} \to \varphi_{16}$}
\end{ptree}
\noindent$\Pi_{4}$:
\begin{ptree}
\AXC{$\forall x\,\forall y\,(x \in \mathbb{R} \to y \in \mathbb{R} \to 0 < y \to x < x + y)$\lbl{[E3]}}
\AXC{$\varphi_{133}$\lbl{A2}}
\BIC{$\forall y\,(s \in \mathbb{R} \to y \in \mathbb{R} \to 0 < y \to s < s + y)$}
\AXC{$\varphi_{134}$\lbl{A2}}
\BIC{$s \in \mathbb{R} \to \varepsilon \in \mathbb{R} \to 0 < \varepsilon \to s < s + \varepsilon$}
\AXC{$(s \in \mathbb{R} \to \varepsilon \in \mathbb{R} \to 0 < \varepsilon \to s < s + \varepsilon) \to \varphi_{256}$\lbl{A1}}
\BIC{$\varphi_{256}$}
\end{ptree}
\noindent$\Pi_{5}$:
\begin{ptree}
\AXC{$\mathrm{MonoInc}(a) \to M \in \mathbb{R} \to \varphi_{89}$\lbl{A1}}
\AXC{$\mathrm{MonoInc}(a) \to M \in \mathbb{R} \to \varphi_{69}$\lbl{A1}}
\AXC{$\stackrel{\Pi_{3}}{\vdots}$}\noLine
\UIC{$\mathrm{MonoInc}(a) \to M \in \mathbb{R} \to \varphi_{16}$}
\AXC{$\stackrel{\Pi_{4}}{\vdots}$}\noLine
\UIC{$\varphi_{256}$}
\BIC{$(\mathrm{MonoInc}(a) \to M \in \mathbb{R} \to \varphi_{69}) \to \varphi_{224}$}
\BIC{$\varphi_{224}$}
\BIC{$\mathrm{MonoInc}(a) \to M \in \mathbb{R} \to \forall n \in \mathbb{N}\,(a_{n} \leq M) \to \varphi_{173}$}
\end{ptree}
\noindent$\Pi_{6}$:
\begin{ptree}
\AXC{$x = x$\lbl{A5}}
\UIC{$\forall x\,(x = x)$}
\AXC{$\forall x\,(x = x) \to a_{n} = a_{n}$\lbl{A2}}
\BIC{$a_{n} = a_{n}$}
\AXC{$\mathrm{MonoInc}(a) \to M \in \mathbb{R} \to \forall n \in \mathbb{N}\,(a_{n} \leq M) \to \varphi_{160}$\lbl{A1}}
\AXC{$\varphi_{251}$\lbl{A1}}
\BIC{$a_{n} = a_{n} \to \mathrm{MonoInc}(a) \to M \in \mathbb{R} \to \varphi_{197}$}
\BIC{$\mathrm{MonoInc}(a) \to M \in \mathbb{R} \to \varphi_{197}$}
\end{ptree}
\noindent$\Pi_{7}$:
\begin{ptree}
\AXC{$\mathrm{MonoInc}(a) \to \mathrm{MonoInc}(a)$\lbl{A1}}
\AXC{$(\mathrm{MonoInc}(a) \to \mathrm{MonoInc}(a)) \to \varphi_{18} \to \varphi_{24}$\lbl{A1}}
\BIC{$\varphi_{18} \to \varphi_{24}$}
\end{ptree}
\noindent$\Pi_{8}$:
\begin{ptree}
\AXC{$\forall a\,(\varphi_{18})$\lbl{DefMonoInc}}
\AXC{$\forall a\,(\varphi_{18}) \to \varphi_{18}$\lbl{A2}}
\BIC{$\varphi_{18}$}
\AXC{$\stackrel{\Pi_{7}}{\vdots}$}\noLine
\UIC{$\varphi_{18} \to \varphi_{24}$}
\BIC{$\varphi_{24}$}
\end{ptree}
\noindent$\Pi_{9}$:
\begin{ptree}
\AXC{$\stackrel{\Pi_{8}}{\vdots}$}\noLine
\UIC{$\varphi_{24}$}
\AXC{$\forall a\,(\mathrm{Seq}(a) \to \varphi_{32})$\lbl{[S16]}}
\AXC{$\forall a\,(\mathrm{Seq}(a) \to \varphi_{32}) \to \mathrm{Seq}(a) \to \varphi_{32}$\lbl{A2}}
\BIC{$\mathrm{Seq}(a) \to \varphi_{32}$}
\AXC{$(\mathrm{Seq}(a) \to \varphi_{32}) \to \varphi_{184}$\lbl{A1}}
\BIC{$\varphi_{184}$}
\BIC{$\mathrm{MonoInc}(a) \to M \in \mathbb{R} \to \forall n \in \mathbb{N}\,(a_{n} \leq M) \to \varphi_{32}$}
\end{ptree}
\noindent$\Pi_{10}$:
\begin{ptree}
\AXC{$\varphi_{128}$\lbl{A2}}
\AXC{$\stackrel{\Pi_{9}}{\vdots}$}\noLine
\UIC{$\mathrm{MonoInc}(a) \to M \in \mathbb{R} \to \forall n \in \mathbb{N}\,(a_{n} \leq M) \to \varphi_{32}$}
\AXC{$\varphi_{228}$\lbl{A1}}
\BIC{$\mathrm{MonoInc}(a) \to M \in \mathbb{R} \to \varphi_{205}$}
\AXC{$\varphi_{261}$\lbl{A1}}
\BIC{$\varphi_{128} \to \mathrm{MonoInc}(a) \to M \in \mathbb{R} \to \varphi_{210}$}
\BIC{$\mathrm{MonoInc}(a) \to M \in \mathbb{R} \to \varphi_{210}$}
\end{ptree}
\noindent$\Pi_{11}$:
\begin{ptree}
\AXC{$\stackrel{\Pi_{10}}{\vdots}$}\noLine
\UIC{$\mathrm{MonoInc}(a) \to M \in \mathbb{R} \to \varphi_{210}$}
\AXC{$(\mathrm{MonoInc}(a) \to M \in \mathbb{R} \to \varphi_{210}) \to \varphi_{250}$\lbl{A1}}
\BIC{$\varphi_{250}$}
\end{ptree}
\noindent$\Pi_{12}$:
\begin{ptree}
\AXC{$\forall a\,\forall X\,(\mathrm{UB}(a,X) \leftrightarrow a \in \mathbb{R} \land \forall x \in X\,(x \leq a))$\lbl{DefUB}}
\AXC{$\varphi_{107}$\lbl{A2}}
\BIC{$\forall X\,(\mathrm{UB}(s,X) \leftrightarrow s \in \mathbb{R} \land \forall x \in X\,(x \leq s))$}
\AXC{$\forall X\,(\mathrm{UB}(s,X) \leftrightarrow s \in \mathbb{R} \land \forall x \in X\,(x \leq s)) \to \varphi_{7}$\lbl{A2}}
\BIC{$\varphi_{7}$}
\end{ptree}
\noindent$\Pi_{13}$:
\begin{ptree}
\AXC{$\mathrm{MonoInc}(a) \to M \in \mathbb{R} \to \forall n \in \mathbb{N}\,(a_{n} \leq M) \to \varphi_{2}$\lbl{A1}}
\AXC{$\varphi_{219}$\lbl{A1}}
\BIC{$\varphi_{9} \to \varphi_{7} \to \mathrm{MonoInc}(a) \to M \in \mathbb{R} \to \varphi_{45}$}
\end{ptree}
\noindent$\Pi_{14}$:
\begin{ptree}
\AXC{$\stackrel{\Pi_{12}}{\vdots}$}\noLine
\UIC{$\varphi_{7}$}
\AXC{$\forall s\,(\varphi_{4})$\lbl{DefSup}}
\AXC{$\forall s\,(\varphi_{4}) \to \varphi_{4}$\lbl{A2}}
\BIC{$\varphi_{4}$}
\AXC{$\varphi_{4} \to \varphi_{9}$\lbl{A2}}
\BIC{$\varphi_{9}$}
\AXC{$\stackrel{\Pi_{13}}{\vdots}$}\noLine
\UIC{$\varphi_{9} \to \varphi_{7} \to \mathrm{MonoInc}(a) \to M \in \mathbb{R} \to \varphi_{45}$}
\BIC{$\varphi_{7} \to \mathrm{MonoInc}(a) \to M \in \mathbb{R} \to \varphi_{45}$}
\BIC{$\mathrm{MonoInc}(a) \to M \in \mathbb{R} \to \varphi_{45}$}
\end{ptree}
\noindent$\Pi_{15}$:
\begin{ptree}
\AXC{$\stackrel{\Pi_{14}}{\vdots}$}\noLine
\UIC{$\mathrm{MonoInc}(a) \to M \in \mathbb{R} \to \varphi_{45}$}
\AXC{$(\mathrm{MonoInc}(a) \to M \in \mathbb{R} \to \varphi_{45}) \to \varphi_{164}$\lbl{A1}}
\BIC{$\varphi_{164}$}
\end{ptree}
\noindent$\Pi_{16}$:
\begin{ptree}
\AXC{$\forall x \in \mathrm{ran}(a)\,(x \leq s) \to a_{n} \in \mathrm{ran}(a) \to a_{n} \leq s$\lbl{A2}}
\AXC{$\stackrel{\Pi_{15}}{\vdots}$}\noLine
\UIC{$\varphi_{164}$}
\BIC{$\mathrm{MonoInc}(a) \to M \in \mathbb{R} \to \varphi_{82}$}
\AXC{$(\mathrm{MonoInc}(a) \to M \in \mathbb{R} \to \varphi_{82}) \to \varphi_{241}$\lbl{A1}}
\BIC{$\varphi_{241}$}
\end{ptree}
\noindent$\Pi_{17}$:
\begin{ptree}
\AXC{$\forall x\,\forall y\,\forall z\,(x \leq y \to y < z \to x < z)$\lbl{[O4]}}
\AXC{$\varphi_{71}$\lbl{A2}}
\BIC{$\forall y\,\forall z\,(a_{n} \leq y \to y < z \to a_{n} < z)$}
\AXC{$\varphi_{91}$\lbl{A2}}
\BIC{$\forall z\,(a_{n} \leq s \to s < z \to a_{n} < z)$}
\AXC{$\varphi_{108}$\lbl{A2}}
\BIC{$a_{n} \leq s \to s < s + \varepsilon \to a_{n} < s + \varepsilon$}
\AXC{$(a_{n} \leq s \to s < s + \varepsilon \to a_{n} < s + \varepsilon) \to \varphi_{265}$\lbl{A1}}
\BIC{$\varphi_{265}$}
\end{ptree}
\noindent$\Pi_{18}$:
\begin{ptree}
\AXC{$\varphi_{41}$\lbl{A2}}
\AXC{$\stackrel{\Pi_{6}}{\vdots}$}\noLine
\UIC{$\mathrm{MonoInc}(a) \to M \in \mathbb{R} \to \varphi_{197}$}
\AXC{$\varphi_{259}$\lbl{A1}}
\BIC{$\varphi_{41} \to \mathrm{MonoInc}(a) \to M \in \mathbb{R} \to \varphi_{198}$}
\BIC{$\mathrm{MonoInc}(a) \to M \in \mathbb{R} \to \varphi_{198}$}
\AXC{$\stackrel{\Pi_{11}}{\vdots}$}\noLine
\UIC{$\varphi_{250}$}
\BIC{$\mathrm{MonoInc}(a) \to M \in \mathbb{R} \to \forall n \in \mathbb{N}\,(a_{n} \leq M) \to \varphi_{177}$}
\AXC{$\stackrel{\Pi_{16}}{\vdots}$}\noLine
\UIC{$\varphi_{241}$}
\BIC{$\mathrm{MonoInc}(a) \to M \in \mathbb{R} \to \forall n \in \mathbb{N}\,(a_{n} \leq M) \to \varphi_{170}$}
\AXC{$\stackrel{\Pi_{17}}{\vdots}$}\noLine
\UIC{$\varphi_{265}$}
\BIC{$\varphi_{243}$}
\end{ptree}
\noindent$\Pi_{19}$:
\begin{ptree}
\AXC{$\forall a\,(\varphi_{42})$\lbl{[C20]}}
\AXC{$\forall a\,(\varphi_{42}) \to \varphi_{42}$\lbl{A2}}
\BIC{$\varphi_{42}$}
\AXC{$\varphi_{42} \to \varphi_{34}$\lbl{A2}}
\BIC{$\varphi_{34}$}
\AXC{$\varphi_{140}$\lbl{A2}}
\BIC{$\mathrm{MonoInc}(a) \to N \in \mathbb{N} \to n \in \mathbb{N} \to N \leq n \to a_{N} \leq a_{n}$}
\AXC{$\varphi_{271}$\lbl{A1}}
\BIC{$(\mathrm{MonoInc}(a) \to \mathrm{MonoInc}(a)) \to \varphi_{269}$}
\end{ptree}
\noindent$\Pi_{20}$:
\begin{ptree}
\AXC{$\mathrm{MonoInc}(a) \to \mathrm{MonoInc}(a)$\lbl{A1}}
\AXC{$\stackrel{\Pi_{19}}{\vdots}$}\noLine
\UIC{$(\mathrm{MonoInc}(a) \to \mathrm{MonoInc}(a)) \to \varphi_{269}$}
\BIC{$\varphi_{269}$}
\end{ptree}
\noindent$\Pi_{21}$:
\begin{ptree}
\AXC{$\mathrm{MonoInc}(a) \to M \in \mathbb{R} \to \forall n \in \mathbb{N}\,(a_{n} \leq M) \to \varphi_{174}$\lbl{A1}}
\AXC{$\varphi_{240}$\lbl{A1}}
\BIC{$\mathrm{MonoInc}(a) \to M \in \mathbb{R} \to \forall n \in \mathbb{N}\,(a_{n} \leq M) \to \varphi_{166}$}
\AXC{$\stackrel{\Pi_{20}}{\vdots}$}\noLine
\UIC{$\varphi_{269}$}
\BIC{$\varphi_{263}$}
\end{ptree}
\noindent$\Pi_{22}$:
\begin{ptree}
\AXC{$\mathrm{MonoInc}(a) \to M \in \mathbb{R} \to \forall n \in \mathbb{N}\,(a_{n} \leq M) \to \varphi_{161}$\lbl{A1}}
\AXC{$\mathrm{MonoInc}(a) \to M \in \mathbb{R} \to \forall n \in \mathbb{N}\,(a_{n} \leq M) \to \varphi_{160}$\lbl{A1}}
\AXC{$\stackrel{\Pi_{21}}{\vdots}$}\noLine
\UIC{$\varphi_{263}$}
\BIC{$\varphi_{237}$}
\BIC{$\mathrm{MonoInc}(a) \to M \in \mathbb{R} \to \forall n \in \mathbb{N}\,(a_{n} \leq M) \to \varphi_{175}$}
\end{ptree}
\noindent$\Pi_{23}$:
\begin{ptree}
\AXC{$\forall x\,\forall y\,\forall z\,(x < y \to y \leq z \to x < z)$\lbl{[O3]}}
\AXC{$\varphi_{72}$\lbl{A2}}
\BIC{$\forall y\,\forall z\,(s - \varepsilon < y \to y \leq z \to s - \varepsilon < z)$}
\AXC{$\varphi_{123}$\lbl{A2}}
\BIC{$\forall z\,(s - \varepsilon < a_{N} \to a_{N} \leq z \to s - \varepsilon < z)$}
\AXC{$\varphi_{138}$\lbl{A2}}
\BIC{$s - \varepsilon < a_{N} \to a_{N} \leq a_{n} \to s - \varepsilon < a_{n}$}
\end{ptree}
\noindent$\Pi_{24}$:
\begin{ptree}
\AXC{$\stackrel{\Pi_{23}}{\vdots}$}\noLine
\UIC{$s - \varepsilon < a_{N} \to a_{N} \leq a_{n} \to s - \varepsilon < a_{n}$}
\AXC{$(s - \varepsilon < a_{N} \to a_{N} \leq a_{n} \to s - \varepsilon < a_{n}) \to \varphi_{266}$\lbl{A1}}
\BIC{$\varphi_{266}$}
\end{ptree}
\noindent$\Pi_{25}$:
\begin{ptree}
\AXC{$\mathrm{MonoInc}(a) \to M \in \mathbb{R} \to \forall n \in \mathbb{N}\,(a_{n} \leq M) \to \varphi_{174}$\lbl{A1}}
\AXC{$\varphi_{244}$\lbl{A1}}
\BIC{$\mathrm{MonoInc}(a) \to M \in \mathbb{R} \to \forall n \in \mathbb{N}\,(a_{n} \leq M) \to \varphi_{180}$}
\AXC{$\stackrel{\Pi_{24}}{\vdots}$}\noLine
\UIC{$\varphi_{266}$}
\BIC{$\varphi_{245}$}
\end{ptree}
\noindent$\Pi_{26}$:
\begin{ptree}
\AXC{$\forall a\,\forall n\,(\mathrm{Seq}(a) \to n \in \mathbb{N} \to a_{n} \in \mathbb{R})$\lbl{[S1]}}
\AXC{$\varphi_{103}$\lbl{A2}}
\BIC{$\forall n\,(\mathrm{Seq}(a) \to n \in \mathbb{N} \to a_{n} \in \mathbb{R})$}
\AXC{$\varphi_{92}$\lbl{A2}}
\BIC{$\mathrm{Seq}(a) \to n \in \mathbb{N} \to a_{n} \in \mathbb{R}$}
\AXC{$(\mathrm{Seq}(a) \to n \in \mathbb{N} \to a_{n} \in \mathbb{R}) \to \varphi_{24} \to \varphi_{238}$\lbl{A1}}
\BIC{$\varphi_{24} \to \varphi_{238}$}
\end{ptree}
\noindent$\Pi_{27}$:
\begin{ptree}
\AXC{$\varphi_{113}$\lbl{[E33]}}
\AXC{$\varphi_{113} \to \varphi_{105}$\lbl{A2}}
\BIC{$\varphi_{105}$}
\AXC{$\varphi_{105} \to \varphi_{121}$\lbl{A2}}
\BIC{$\varphi_{121}$}
\AXC{$\varphi_{121} \to s \in \mathbb{R} \to \varphi_{136}$\lbl{A2}}
\BIC{$s \in \mathbb{R} \to \varphi_{136}$}
\AXC{$(s \in \mathbb{R} \to \varphi_{136}) \to \varphi_{273}$\lbl{A1}}
\BIC{$\varphi_{273}$}
\end{ptree}
\noindent$\Pi_{28}$:
\begin{ptree}
\AXC{$\mathrm{MonoInc}(a) \to M \in \mathbb{R} \to \varphi_{69}$\lbl{A1}}
\AXC{$\stackrel{\Pi_{3}}{\vdots}$}\noLine
\UIC{$\mathrm{MonoInc}(a) \to M \in \mathbb{R} \to \varphi_{16}$}
\AXC{$\stackrel{\Pi_{27}}{\vdots}$}\noLine
\UIC{$\varphi_{273}$}
\BIC{$(\mathrm{MonoInc}(a) \to M \in \mathbb{R} \to \varphi_{69}) \to \varphi_{270}$}
\BIC{$\varphi_{270}$}
\end{ptree}
\noindent$\Pi_{29}$:
\begin{ptree}
\AXC{$\mathrm{MonoInc}(a) \to M \in \mathbb{R} \to \forall n \in \mathbb{N}\,(a_{n} \leq M) \to \varphi_{160}$\lbl{A1}}
\AXC{$\stackrel{\Pi_{8}}{\vdots}$}\noLine
\UIC{$\varphi_{24}$}
\AXC{$\stackrel{\Pi_{26}}{\vdots}$}\noLine
\UIC{$\varphi_{24} \to \varphi_{238}$}
\BIC{$\varphi_{238}$}
\BIC{$\mathrm{MonoInc}(a) \to M \in \mathbb{R} \to \forall n \in \mathbb{N}\,(a_{n} \leq M) \to \varphi_{176}$}
\AXC{$\stackrel{\Pi_{28}}{\vdots}$}\noLine
\UIC{$\varphi_{270}$}
\BIC{$\varphi_{267}$}
\end{ptree}
\noindent$\Pi_{30}$:
\begin{ptree}
\AXC{$\stackrel{\Pi_{22}}{\vdots}$}\noLine
\UIC{$\mathrm{MonoInc}(a) \to M \in \mathbb{R} \to \forall n \in \mathbb{N}\,(a_{n} \leq M) \to \varphi_{175}$}
\AXC{$\stackrel{\Pi_{25}}{\vdots}$}\noLine
\UIC{$\varphi_{245}$}
\BIC{$\mathrm{MonoInc}(a) \to M \in \mathbb{R} \to \forall n \in \mathbb{N}\,(a_{n} \leq M) \to \varphi_{181}$}
\AXC{$\stackrel{\Pi_{29}}{\vdots}$}\noLine
\UIC{$\varphi_{267}$}
\BIC{$\varphi_{247}$}
\end{ptree}
\noindent$\Pi_{31}$:
\begin{ptree}
\AXC{$\varphi_{252}$\lbl{A3}}
\AXC{$\stackrel{\Pi_{5}}{\vdots}$}\noLine
\UIC{$\mathrm{MonoInc}(a) \to M \in \mathbb{R} \to \forall n \in \mathbb{N}\,(a_{n} \leq M) \to \varphi_{173}$}
\AXC{$\stackrel{\Pi_{18}}{\vdots}$}\noLine
\UIC{$\varphi_{243}$}
\BIC{$\mathrm{MonoInc}(a) \to M \in \mathbb{R} \to \forall n \in \mathbb{N}\,(a_{n} \leq M) \to \varphi_{179}$}
\AXC{$\stackrel{\Pi_{30}}{\vdots}$}\noLine
\UIC{$\varphi_{247}$}
\BIC{$\mathrm{MonoInc}(a) \to M \in \mathbb{R} \to \forall n \in \mathbb{N}\,(a_{n} \leq M) \to \varphi_{185}$}
\AXC{$\varphi_{249}$\lbl{A1}}
\BIC{$\varphi_{188} \to n \in \mathbb{N} \to N \leq n \to \lvert a_{n} - s \rvert < \varepsilon$}
\UIC{$\forall n\,(\varphi_{188} \to n \in \mathbb{N} \to N \leq n \to \lvert a_{n} - s \rvert < \varepsilon)$}
\AXC{$\varphi_{275}$\lbl{A1}}
\BIC{$\varphi_{252} \to (\varphi_{188} \to \forall n\,(\varphi_{188})) \to \mathrm{MonoInc}(a) \to M \in \mathbb{R} \to \varphi_{196}$}
\BIC{$(\varphi_{188} \to \forall n\,(\varphi_{188})) \to \mathrm{MonoInc}(a) \to M \in \mathbb{R} \to \varphi_{196}$}
\end{ptree}
\noindent$\Pi_{32}$:
\begin{ptree}
\AXC{$\mathrm{MonoInc}(a) \to M \in \mathbb{R} \to \forall n \in \mathbb{N}\,(a_{n} \leq M) \to \varphi_{174}$\lbl{A1}}
\AXC{$\varphi_{233}$\lbl{A1}}
\BIC{$\mathrm{MonoInc}(a) \to M \in \mathbb{R} \to \varphi_{178}$}
\AXC{$(\mathrm{MonoInc}(a) \to M \in \mathbb{R} \to \varphi_{178}) \to \varphi_{253}$\lbl{A1}}
\BIC{$\varphi_{253}$}
\end{ptree}
\noindent$\Pi_{33}$:
\begin{ptree}
\AXC{$\varphi_{188} \to \forall n\,(\varphi_{188})$\lbl{A4}}
\AXC{$\stackrel{\Pi_{31}}{\vdots}$}\noLine
\UIC{$(\varphi_{188} \to \forall n\,(\varphi_{188})) \to \mathrm{MonoInc}(a) \to M \in \mathbb{R} \to \varphi_{196}$}
\BIC{$\mathrm{MonoInc}(a) \to M \in \mathbb{R} \to \varphi_{196}$}
\AXC{$\stackrel{\Pi_{32}}{\vdots}$}\noLine
\UIC{$\varphi_{253}$}
\BIC{$\mathrm{MonoInc}(a) \to M \in \mathbb{R} \to \varphi_{204}$}
\end{ptree}
\noindent$\Pi_{34}$:
\begin{ptree}
\AXC{$\stackrel{\Pi_{33}}{\vdots}$}\noLine
\UIC{$\mathrm{MonoInc}(a) \to M \in \mathbb{R} \to \varphi_{204}$}
\AXC{$(\mathrm{MonoInc}(a) \to M \in \mathbb{R} \to \varphi_{204}) \to \varphi_{229}$\lbl{A1}}
\BIC{$\varphi_{229}$}
\end{ptree}
\noindent$\Pi_{35}$:
\begin{ptree}
\AXC{$\varphi_{31} \to \neg (N \in \mathbb{N} \land \forall n \in \mathbb{N}\,(N \leq n \to \lvert a_{n} - s \rvert < \varepsilon))$\lbl{A2}}
\AXC{$\stackrel{\Pi_{34}}{\vdots}$}\noLine
\UIC{$\varphi_{229}$}
\BIC{$\mathrm{MonoInc}(a) \to M \in \mathbb{R} \to \varphi_{206}$}
\AXC{$\varphi_{255}$\lbl{A1}}
\BIC{$\varphi_{200} \to \neg (N \in \mathbb{N} \land s - \varepsilon < a_{N})$}
\UIC{$\forall N\,(\varphi_{200} \to \neg (N \in \mathbb{N} \land s - \varepsilon < a_{N}))$}
\end{ptree}
\noindent$\Pi_{36}$:
\begin{ptree}
\AXC{$\mathrm{MonoInc}(a) \to M \in \mathbb{R} \to \varphi_{143}$\lbl{A1}}
\AXC{$\varphi_{223}$\lbl{A1}}
\BIC{$\mathrm{MonoInc}(a) \to M \in \mathbb{R} \to \varphi_{148}$}
\AXC{$\mathrm{MonoInc}(a) \to M \in \mathbb{R} \to \varphi_{159}$\lbl{A1}}
\AXC{$\varphi_{226}$\lbl{A1}}
\BIC{$\mathrm{MonoInc}(a) \to M \in \mathbb{R} \to \varphi_{146}$}
\AXC{$\varphi_{260}$\lbl{A1}}
\BIC{$(\mathrm{MonoInc}(a) \to M \in \mathbb{R} \to \varphi_{148}) \to \varphi_{207}$}
\BIC{$\varphi_{207}$}
\end{ptree}
\noindent$\Pi_{37}$:
\begin{ptree}
\AXC{$\mathrm{MonoInc}(a) \to M \in \mathbb{R} \to \varphi_{159}$\lbl{A1}}
\AXC{$\varphi_{225}$\lbl{A1}}
\BIC{$\mathrm{MonoInc}(a) \to M \in \mathbb{R} \to \varphi_{145}$}
\AXC{$(\mathrm{MonoInc}(a) \to M \in \mathbb{R} \to \varphi_{145}) \to \varphi_{227}$\lbl{A1}}
\BIC{$\varphi_{227}$}
\end{ptree}
\noindent$\Pi_{38}$:
\begin{ptree}
\AXC{$\varphi_{73}$\lbl{A2}}
\AXC{$z = w \to s - \varepsilon < z \to s - \varepsilon < w$\lbl{A6}}
\UIC{$\forall w\,(z = w \to s - \varepsilon < z \to s - \varepsilon < w)$}
\UIC{$\forall z\,\forall w\,(z = w \to s - \varepsilon < z \to s - \varepsilon < w)$}
\AXC{$\varphi_{93}$\lbl{A2}}
\BIC{$\forall w\,(x = w \to s - \varepsilon < x \to s - \varepsilon < w)$}
\AXC{$\varphi_{114}$\lbl{A2}}
\BIC{$x = a_{i_{0}} \to s - \varepsilon < x \to s - \varepsilon < a_{i_{0}}$}
\AXC{$\stackrel{\Pi_{36}}{\vdots}$}\noLine
\UIC{$\varphi_{207}$}
\BIC{$\mathrm{MonoInc}(a) \to M \in \mathbb{R} \to \varphi_{155}$}
\AXC{$\stackrel{\Pi_{37}}{\vdots}$}\noLine
\UIC{$\varphi_{227}$}
\BIC{$\mathrm{MonoInc}(a) \to M \in \mathbb{R} \to \varphi_{167}$}
\AXC{$\varphi_{246}$\lbl{A1}}
\BIC{$\varphi_{73} \to \mathrm{MonoInc}(a) \to M \in \mathbb{R} \to \varphi_{172}$}
\BIC{$\mathrm{MonoInc}(a) \to M \in \mathbb{R} \to \varphi_{172}$}
\end{ptree}
\noindent$\Pi_{39}$:
\begin{ptree}
\AXC{$\varphi_{200} \to \forall N\,(\varphi_{200})$\lbl{A4}}
\AXC{$\varphi_{257}$\lbl{A3}}
\AXC{$\stackrel{\Pi_{35}}{\vdots}$}\noLine
\UIC{$\forall N\,(\varphi_{200} \to \neg (N \in \mathbb{N} \land s - \varepsilon < a_{N}))$}
\AXC{$\stackrel{\Pi_{38}}{\vdots}$}\noLine
\UIC{$\mathrm{MonoInc}(a) \to M \in \mathbb{R} \to \varphi_{172}$}
\AXC{$\varphi_{277}$\lbl{A1}}
\BIC{$\forall N\,(\varphi_{200} \to \neg (N \in \mathbb{N} \land s - \varepsilon < a_{N})) \to \varphi_{257} \to \varphi_{264}$}
\BIC{$\varphi_{257} \to \varphi_{264}$}
\BIC{$\varphi_{264}$}
\BIC{$\mathrm{MonoInc}(a) \to M \in \mathbb{R} \to \forall n \in \mathbb{N}\,(a_{n} \leq M) \to \varphi_{165}$}
\end{ptree}
\noindent$\Pi_{40}$:
\begin{ptree}
\AXC{$\stackrel{\Pi_{39}}{\vdots}$}\noLine
\UIC{$\mathrm{MonoInc}(a) \to M \in \mathbb{R} \to \forall n \in \mathbb{N}\,(a_{n} \leq M) \to \varphi_{165}$}
\AXC{$\varphi_{236}$\lbl{A1}}
\BIC{$\varphi_{190} \to \neg (i_{0} \in \mathbb{N} \land x = a_{i_{0}})$}
\UIC{$\forall i_{0}\,(\varphi_{190} \to \neg (i_{0} \in \mathbb{N} \land x = a_{i_{0}}))$}
\end{ptree}
\noindent$\Pi_{41}$:
\begin{ptree}
\AXC{$\varphi_{100}$\lbl{A2}}
\AXC{$\stackrel{\Pi_{9}}{\vdots}$}\noLine
\UIC{$\mathrm{MonoInc}(a) \to M \in \mathbb{R} \to \forall n \in \mathbb{N}\,(a_{n} \leq M) \to \varphi_{32}$}
\AXC{$\varphi_{218}$\lbl{A1}}
\BIC{$\mathrm{MonoInc}(a) \to M \in \mathbb{R} \to \varphi_{151}$}
\AXC{$\varphi_{235}$\lbl{A1}}
\BIC{$\varphi_{100} \to \mathrm{MonoInc}(a) \to M \in \mathbb{R} \to \varphi_{149}$}
\BIC{$\mathrm{MonoInc}(a) \to M \in \mathbb{R} \to \varphi_{149}$}
\AXC{$\varphi_{248}$\lbl{A1}}
\BIC{$(\mathrm{MonoInc}(a) \to \varphi_{141}) \to \mathrm{MonoInc}(a) \to \varphi_{147}$}
\end{ptree}
\noindent$\Pi_{42}$:
\begin{ptree}
\AXC{$\mathrm{MonoInc}(a) \to M \in \mathbb{R} \to \varphi_{143}$\lbl{A1}}
\AXC{$\varphi_{220}$\lbl{A1}}
\BIC{$\mathrm{MonoInc}(a) \to \varphi_{141}$}
\AXC{$\stackrel{\Pi_{41}}{\vdots}$}\noLine
\UIC{$(\mathrm{MonoInc}(a) \to \varphi_{141}) \to \mathrm{MonoInc}(a) \to \varphi_{147}$}
\BIC{$\mathrm{MonoInc}(a) \to \varphi_{147}$}
\end{ptree}
\noindent$\Pi_{43}$:
\begin{ptree}
\AXC{$\varphi_{242}$\lbl{A3}}
\AXC{$\stackrel{\Pi_{40}}{\vdots}$}\noLine
\UIC{$\forall i_{0}\,(\varphi_{190} \to \neg (i_{0} \in \mathbb{N} \land x = a_{i_{0}}))$}
\AXC{$\stackrel{\Pi_{42}}{\vdots}$}\noLine
\UIC{$\mathrm{MonoInc}(a) \to \varphi_{147}$}
\AXC{$(\mathrm{MonoInc}(a) \to \varphi_{147}) \to \varphi_{276}$\lbl{A1}}
\BIC{$\varphi_{276}$}
\BIC{$\varphi_{242} \to (\varphi_{190} \to \forall i_{0}\,(\varphi_{190})) \to \varphi_{239}$}
\BIC{$(\varphi_{190} \to \forall i_{0}\,(\varphi_{190})) \to \varphi_{239}$}
\end{ptree}
\noindent$\Pi_{44}$:
\begin{ptree}
\AXC{$\varphi_{36}$\lbl{A4}}
\AXC{$\forall j\,(\varphi_{11}) \to \varphi_{27}$\lbl{A3}}
\AXC{$\varphi_{11}$\lbl{A2}}
\UIC{$\forall j\,(\varphi_{11})$}
\AXC{$\varphi_{190} \to \forall i_{0}\,(\varphi_{190})$\lbl{A4}}
\AXC{$\stackrel{\Pi_{43}}{\vdots}$}\noLine
\UIC{$(\varphi_{190} \to \forall i_{0}\,(\varphi_{190})) \to \varphi_{239}$}
\BIC{$\varphi_{239}$}
\BIC{$(\forall j\,(\varphi_{11}) \to \varphi_{27}) \to \varphi_{36} \to \mathrm{MonoInc}(a) \to M \in \mathbb{R} \to \varphi_{168}$}
\BIC{$\varphi_{36} \to \mathrm{MonoInc}(a) \to M \in \mathbb{R} \to \varphi_{168}$}
\BIC{$\mathrm{MonoInc}(a) \to M \in \mathbb{R} \to \varphi_{168}$}
\AXC{$\varphi_{230}$\lbl{A1}}
\BIC{$\varphi_{162} \to \neg (x \in \mathrm{ran}(a) \land s - \varepsilon < x)$}
\UIC{$\forall x\,(\varphi_{162} \to \neg (x \in \mathrm{ran}(a) \land s - \varepsilon < x))$}
\end{ptree}
\noindent$\Pi_{45}$:
\begin{ptree}
\AXC{$\stackrel{\Pi_{9}}{\vdots}$}\noLine
\UIC{$\mathrm{MonoInc}(a) \to M \in \mathbb{R} \to \forall n \in \mathbb{N}\,(a_{n} \leq M) \to \varphi_{32}$}
\AXC{$\varphi_{189}$\lbl{A1}}
\BIC{$\varphi_{59}$}
\end{ptree}
\noindent$\Pi_{46}$:
\begin{ptree}
\AXC{$\forall s\,(\varphi_{86})$\lbl{[C22]}}
\AXC{$\forall s\,(\varphi_{86}) \to \varphi_{86}$\lbl{A2}}
\BIC{$\varphi_{86}$}
\AXC{$\varphi_{86} \to \varphi_{112}$\lbl{A2}}
\BIC{$\varphi_{112}$}
\AXC{$\varphi_{112} \to \mathrm{ran}(a) \subseteq \mathbb{R} \to \varphi_{75}$\lbl{A2}}
\BIC{$\mathrm{ran}(a) \subseteq \mathbb{R} \to \varphi_{75}$}
\end{ptree}
\noindent$\Pi_{47}$:
\begin{ptree}
\AXC{$\stackrel{\Pi_{45}}{\vdots}$}\noLine
\UIC{$\varphi_{59}$}
\AXC{$\stackrel{\Pi_{46}}{\vdots}$}\noLine
\UIC{$\mathrm{ran}(a) \subseteq \mathbb{R} \to \varphi_{75}$}
\AXC{$(\mathrm{ran}(a) \subseteq \mathbb{R} \to \varphi_{75}) \to \varphi_{59} \to \varphi_{234}$\lbl{A1}}
\BIC{$\varphi_{59} \to \varphi_{234}$}
\BIC{$\varphi_{234}$}
\end{ptree}
\noindent$\Pi_{48}$:
\begin{ptree}
\AXC{$\mathrm{MonoInc}(a) \to M \in \mathbb{R} \to \varphi_{69}$\lbl{A1}}
\AXC{$\mathrm{MonoInc}(a) \to M \in \mathbb{R} \to \forall n \in \mathbb{N}\,(a_{n} \leq M) \to \varphi_{2}$\lbl{A1}}
\AXC{$\stackrel{\Pi_{47}}{\vdots}$}\noLine
\UIC{$\varphi_{234}$}
\BIC{$(\mathrm{MonoInc}(a) \to M \in \mathbb{R} \to \varphi_{69}) \to \varphi_{211}$}
\BIC{$\varphi_{211}$}
\end{ptree}
\noindent$\Pi_{49}$:
\begin{ptree}
\AXC{$\varphi_{231}$\lbl{A3}}
\AXC{$\stackrel{\Pi_{44}}{\vdots}$}\noLine
\UIC{$\forall x\,(\varphi_{162} \to \neg (x \in \mathrm{ran}(a) \land s - \varepsilon < x))$}
\AXC{$\mathrm{MonoInc}(a) \to M \in \mathbb{R} \to \varphi_{89}$\lbl{A1}}
\AXC{$\stackrel{\Pi_{48}}{\vdots}$}\noLine
\UIC{$\varphi_{211}$}
\BIC{$\mathrm{MonoInc}(a) \to M \in \mathbb{R} \to \forall n \in \mathbb{N}\,(a_{n} \leq M) \to \varphi_{75}$}
\AXC{$\varphi_{274}$\lbl{A1}}
\BIC{$\varphi_{272}$}
\BIC{$\varphi_{231} \to (\varphi_{162} \to \forall x\,(\varphi_{162})) \to \mathrm{MonoInc}(a) \to \varphi_{152}$}
\BIC{$(\varphi_{162} \to \forall x\,(\varphi_{162})) \to \mathrm{MonoInc}(a) \to \varphi_{152}$}
\end{ptree}
\noindent$\Pi_{50}$:
\begin{ptree}
\AXC{$\varphi_{162} \to \forall x\,(\varphi_{162})$\lbl{A4}}
\AXC{$\stackrel{\Pi_{49}}{\vdots}$}\noLine
\UIC{$(\varphi_{162} \to \forall x\,(\varphi_{162})) \to \mathrm{MonoInc}(a) \to \varphi_{152}$}
\BIC{$\mathrm{MonoInc}(a) \to \varphi_{152}$}
\AXC{$(\mathrm{MonoInc}(a) \to \varphi_{152}) \to \varphi_{56} \to \varphi_{102}$\lbl{A1}}
\BIC{$\varphi_{56} \to \varphi_{102}$}
\UIC{$\forall \varepsilon\,(\varphi_{56} \to \varphi_{102})$}
\end{ptree}
\noindent$\Pi_{51}$:
\begin{ptree}
\AXC{$\varphi_{56} \to \forall \varepsilon\,(\varphi_{56})$\lbl{A4}}
\AXC{$\forall \varepsilon\,(\varphi_{56} \to \varphi_{102}) \to \forall \varepsilon\,(\varphi_{56}) \to \varphi_{115}$\lbl{A3}}
\AXC{$\stackrel{\Pi_{50}}{\vdots}$}\noLine
\UIC{$\forall \varepsilon\,(\varphi_{56} \to \varphi_{102})$}
\AXC{$\forall \varepsilon\,(\varphi_{56} \to \varphi_{102}) \to \varphi_{262}$\lbl{A1}}
\BIC{$\varphi_{262}$}
\BIC{$(\varphi_{56} \to \forall \varepsilon\,(\varphi_{56})) \to \mathrm{MonoInc}(a) \to \varphi_{156}$}
\BIC{$\mathrm{MonoInc}(a) \to \varphi_{156}$}
\end{ptree}
\noindent$\Pi_{52}$:
\begin{ptree}
\AXC{$\stackrel{\Pi_{8}}{\vdots}$}\noLine
\UIC{$\varphi_{24}$}
\AXC{$\forall a\,(\varphi_{144})$\lbl{DefConv}}
\AXC{$\forall a\,(\varphi_{144}) \to \varphi_{144}$\lbl{A2}}
\BIC{$\varphi_{144}$}
\AXC{$\varphi_{144} \to (\mathrm{Conv}(a,s) \leftrightarrow \mathrm{Seq}(a) \land s \in \mathbb{R} \land \varphi_{115})$\lbl{A2}}
\BIC{$\mathrm{Conv}(a,s) \leftrightarrow \mathrm{Seq}(a) \land s \in \mathbb{R} \land \varphi_{115}$}
\AXC{$\varphi_{258}$\lbl{A1}}
\BIC{$\varphi_{24} \to (\mathrm{MonoInc}(a) \to M \in \mathbb{R} \to \varphi_{16}) \to \varphi_{212}$}
\BIC{$(\mathrm{MonoInc}(a) \to M \in \mathbb{R} \to \varphi_{16}) \to \varphi_{212}$}
\end{ptree}
\noindent$\Pi_{53}$:
\begin{ptree}
\AXC{$\stackrel{\Pi_{51}}{\vdots}$}\noLine
\UIC{$\mathrm{MonoInc}(a) \to \varphi_{156}$}
\AXC{$\stackrel{\Pi_{3}}{\vdots}$}\noLine
\UIC{$\mathrm{MonoInc}(a) \to M \in \mathbb{R} \to \varphi_{16}$}
\AXC{$\stackrel{\Pi_{52}}{\vdots}$}\noLine
\UIC{$(\mathrm{MonoInc}(a) \to M \in \mathbb{R} \to \varphi_{16}) \to \varphi_{212}$}
\BIC{$\varphi_{212}$}
\BIC{$\mathrm{MonoInc}(a) \to M \in \mathbb{R} \to \varphi_{13}$}
\end{ptree}
\noindent$\Pi_{54}$:
\begin{ptree}
\AXC{$\stackrel{\Pi_{53}}{\vdots}$}\noLine
\UIC{$\mathrm{MonoInc}(a) \to M \in \mathbb{R} \to \varphi_{13}$}
\AXC{$(\mathrm{MonoInc}(a) \to M \in \mathbb{R} \to \varphi_{13}) \to \varphi_{137}$\lbl{A1}}
\BIC{$\varphi_{137}$}
\end{ptree}
\noindent$\Pi_{55}$:
\begin{ptree}
\AXC{$\forall n \in \mathbb{N}\,(a_{n} \leq M) \to n \in \mathbb{N} \to a_{n} \leq M$\lbl{A2}}
\AXC{$\varphi_{87}$\lbl{A1}}
\AXC{$\varphi_{87} \to (\forall n \in \mathbb{N}\,(a_{n} \leq M) \to n \in \mathbb{N} \to a_{n} \leq M) \to \varphi_{77}$\lbl{A1}}
\BIC{$(\forall n \in \mathbb{N}\,(a_{n} \leq M) \to n \in \mathbb{N} \to a_{n} \leq M) \to \varphi_{77}$}
\BIC{$\varphi_{77}$}
\AXC{$\varphi_{77} \to \varphi_{195}$\lbl{A1}}
\BIC{$\varphi_{195}$}
\end{ptree}
\noindent$\Pi_{56}$:
\begin{ptree}
\AXC{$x = x$\lbl{A5}}
\AXC{$x = y \to x = x \to y = x$\lbl{A6}}
\AXC{$(x = y \to x = x \to y = x) \to x = x \to x = y \to y = x$\lbl{A1}}
\BIC{$x = x \to x = y \to y = x$}
\BIC{$x = y \to y = x$}
\UIC{$\forall y\,(x = y \to y = x)$}
\UIC{$\forall x\,\forall y\,(x = y \to y = x)$}
\end{ptree}
\noindent$\Pi_{57}$:
\begin{ptree}
\AXC{$\stackrel{\Pi_{56}}{\vdots}$}\noLine
\UIC{$\forall x\,\forall y\,(x = y \to y = x)$}
\AXC{$\forall x\,\forall y\,(x = y \to y = x) \to \forall y\,(v_{718} = y \to y = v_{718})$\lbl{A2}}
\BIC{$\forall y\,(v_{718} = y \to y = v_{718})$}
\AXC{$\varphi_{5}$\lbl{A2}}
\BIC{$v_{718} = v_{719} \to v_{719} = v_{718}$}
\UIC{$\forall v_{719}\,(v_{718} = v_{719} \to v_{719} = v_{718})$}
\UIC{$\forall v_{718}\,\forall v_{719}\,(v_{718} = v_{719} \to v_{719} = v_{718})$}
\AXC{$\varphi_{19}$\lbl{A2}}
\BIC{$\forall v_{719}\,(y = v_{719} \to v_{719} = y)$}
\end{ptree}
\noindent$\Pi_{58}$:
\begin{ptree}
\AXC{$\stackrel{\Pi_{57}}{\vdots}$}\noLine
\UIC{$\forall v_{719}\,(y = v_{719} \to v_{719} = y)$}
\AXC{$\forall v_{719}\,(y = v_{719} \to v_{719} = y) \to y = a_{n} \to a_{n} = y$\lbl{A2}}
\BIC{$y = a_{n} \to a_{n} = y$}
\end{ptree}
\noindent$\Pi_{59}$:
\begin{ptree}
\AXC{$\stackrel{\Pi_{58}}{\vdots}$}\noLine
\UIC{$y = a_{n} \to a_{n} = y$}
\AXC{$\mathrm{MonoInc}(a) \to M \in \mathbb{R} \to \varphi_{111}$\lbl{A1}}
\AXC{$\varphi_{202}$\lbl{A1}}
\BIC{$\mathrm{MonoInc}(a) \to M \in \mathbb{R} \to \varphi_{84}$}
\AXC{$\varphi_{209}$\lbl{A1}}
\BIC{$(y = a_{n} \to a_{n} = y) \to \mathrm{MonoInc}(a) \to M \in \mathbb{R} \to \varphi_{85}$}
\BIC{$\mathrm{MonoInc}(a) \to M \in \mathbb{R} \to \varphi_{85}$}
\end{ptree}
\noindent$\Pi_{60}$:
\begin{ptree}
\AXC{$\stackrel{\Pi_{59}}{\vdots}$}\noLine
\UIC{$\mathrm{MonoInc}(a) \to M \in \mathbb{R} \to \varphi_{85}$}
\AXC{$(\mathrm{MonoInc}(a) \to M \in \mathbb{R} \to \varphi_{85}) \to \varphi_{214}$\lbl{A1}}
\BIC{$\varphi_{214}$}
\end{ptree}
\noindent$\Pi_{61}$:
\begin{ptree}
\AXC{$\mathrm{MonoInc}(a) \to M \in \mathbb{R} \to \varphi_{111}$\lbl{A1}}
\AXC{$\varphi_{201}$\lbl{A1}}
\BIC{$\mathrm{MonoInc}(a) \to M \in \mathbb{R} \to \varphi_{79}$}
\AXC{$\stackrel{\Pi_{55}}{\vdots}$}\noLine
\UIC{$\varphi_{195}$}
\BIC{$\mathrm{MonoInc}(a) \to M \in \mathbb{R} \to \varphi_{83}$}
\AXC{$\stackrel{\Pi_{60}}{\vdots}$}\noLine
\UIC{$\varphi_{214}$}
\BIC{$(a_{n} = y \to a_{n} < M \to y < M) \to \varphi_{142}$}
\end{ptree}
\noindent$\Pi_{62}$:
\begin{ptree}
\AXC{$z = w \to z < M \to w < M$\lbl{A6}}
\UIC{$\forall w\,(z = w \to z < M \to w < M)$}
\UIC{$\forall z\,\forall w\,(z = w \to z < M \to w < M)$}
\AXC{$\varphi_{48}$\lbl{A2}}
\BIC{$\forall w\,(a_{n} = w \to a_{n} < M \to w < M)$}
\AXC{$\varphi_{76}$\lbl{A2}}
\BIC{$a_{n} = y \to a_{n} < M \to y < M$}
\AXC{$\stackrel{\Pi_{61}}{\vdots}$}\noLine
\UIC{$(a_{n} = y \to a_{n} < M \to y < M) \to \varphi_{142}$}
\BIC{$\varphi_{142}$}
\end{ptree}
\noindent$\Pi_{63}$:
\begin{ptree}
\AXC{$z = w \to z = M \to w = M$\lbl{A6}}
\UIC{$\forall w\,(z = w \to z = M \to w = M)$}
\UIC{$\forall z\,\forall w\,(z = w \to z = M \to w = M)$}
\AXC{$\varphi_{15}$\lbl{A2}}
\BIC{$\forall w\,(a_{n} = w \to a_{n} = M \to w = M)$}
\AXC{$\varphi_{37}$\lbl{A2}}
\BIC{$a_{n} = y \to a_{n} = M \to y = M$}
\AXC{$\stackrel{\Pi_{62}}{\vdots}$}\noLine
\UIC{$\varphi_{142}$}
\BIC{$\mathrm{MonoInc}(a) \to M \in \mathbb{R} \to \varphi_{64}$}
\AXC{$\varphi_{193}$\lbl{A1}}
\BIC{$\varphi_{70} \to \neg (n \in \mathbb{N} \land y = a_{n})$}
\UIC{$\forall n\,(\varphi_{70} \to \neg (n \in \mathbb{N} \land y = a_{n}))$}
\end{ptree}
\noindent$\Pi_{64}$:
\begin{ptree}
\AXC{$\varphi_{101}$\lbl{A2}}
\AXC{$\stackrel{\Pi_{9}}{\vdots}$}\noLine
\UIC{$\mathrm{MonoInc}(a) \to M \in \mathbb{R} \to \forall n \in \mathbb{N}\,(a_{n} \leq M) \to \varphi_{32}$}
\AXC{$\varphi_{208}$\lbl{A1}}
\BIC{$\mathrm{MonoInc}(a) \to M \in \mathbb{R} \to \forall n \in \mathbb{N}\,(a_{n} \leq M) \to \varphi_{23}$}
\AXC{$\varphi_{221}$\lbl{A1}}
\BIC{$\varphi_{192}$}
\BIC{$\mathrm{MonoInc}(a) \to M \in \mathbb{R} \to \forall n \in \mathbb{N}\,(a_{n} \leq M) \to \varphi_{14}$}
\AXC{$\varphi_{217}$\lbl{A1}}
\BIC{$\varphi_{183}$}
\end{ptree}
\noindent$\Pi_{65}$:
\begin{ptree}
\AXC{$\stackrel{\Pi_{63}}{\vdots}$}\noLine
\UIC{$\forall n\,(\varphi_{70} \to \neg (n \in \mathbb{N} \land y = a_{n}))$}
\AXC{$\mathrm{MonoInc}(a) \to \varphi_{53}$\lbl{A1}}
\AXC{$\stackrel{\Pi_{64}}{\vdots}$}\noLine
\UIC{$\varphi_{183}$}
\BIC{$\mathrm{MonoInc}(a) \to M \in \mathbb{R} \to \varphi_{54}$}
\AXC{$(\mathrm{MonoInc}(a) \to M \in \mathbb{R} \to \varphi_{54}) \to \varphi_{268}$\lbl{A1}}
\BIC{$\varphi_{268}$}
\BIC{$\varphi_{194} \to (\varphi_{70} \to \forall n\,(\varphi_{70})) \to \varphi_{12} \to (\varphi_{12} \to \varphi_{22}) \to \varphi_{21} \to \varphi_{66}$}
\end{ptree}
\noindent$\Pi_{66}$:
\begin{ptree}
\AXC{$\varphi_{70} \to \forall n\,(\varphi_{70})$\lbl{A4}}
\AXC{$\varphi_{194}$\lbl{A3}}
\AXC{$\stackrel{\Pi_{65}}{\vdots}$}\noLine
\UIC{$\varphi_{194} \to (\varphi_{70} \to \forall n\,(\varphi_{70})) \to \varphi_{12} \to (\varphi_{12} \to \varphi_{22}) \to \varphi_{21} \to \varphi_{66}$}
\BIC{$(\varphi_{70} \to \forall n\,(\varphi_{70})) \to \varphi_{12} \to (\varphi_{12} \to \varphi_{22}) \to \varphi_{21} \to \varphi_{66}$}
\BIC{$\varphi_{12} \to (\varphi_{12} \to \varphi_{22}) \to \varphi_{21} \to \varphi_{66}$}
\end{ptree}
\noindent$\Pi_{67}$:
\begin{ptree}
\AXC{$\varphi_{21}$\lbl{A4}}
\AXC{$\varphi_{12} \to \varphi_{22}$\lbl{A3}}
\AXC{$\forall n\,(\neg (n \in \mathbb{N} \land y = a_{n})) \to \neg (j \in \mathbb{N} \land y = a_{j})$\lbl{A2}}
\UIC{$\varphi_{12}$}
\AXC{$\stackrel{\Pi_{66}}{\vdots}$}\noLine
\UIC{$\varphi_{12} \to (\varphi_{12} \to \varphi_{22}) \to \varphi_{21} \to \varphi_{66}$}
\BIC{$(\varphi_{12} \to \varphi_{22}) \to \varphi_{21} \to \varphi_{66}$}
\BIC{$\varphi_{21} \to \varphi_{66}$}
\BIC{$\varphi_{66}$}
\end{ptree}
\noindent$\Pi_{68}$:
\begin{ptree}
\AXC{$\forall y\,(\varphi_{67}) \to \varphi_{10} \to \forall y \in \mathrm{ran}(a)\,(y \leq M)$\lbl{A3}}
\AXC{$\stackrel{\Pi_{67}}{\vdots}$}\noLine
\UIC{$\varphi_{66}$}
\AXC{$\varphi_{66} \to \varphi_{67}$\lbl{A1}}
\BIC{$\varphi_{67}$}
\UIC{$\forall y\,(\varphi_{67})$}
\AXC{$\varphi_{232}$\lbl{A1}}
\BIC{$(\forall y\,(\varphi_{67}) \to \varphi_{10} \to \forall y \in \mathrm{ran}(a)\,(y \leq M)) \to \varphi_{186}$}
\BIC{$\varphi_{186}$}
\end{ptree}
\noindent$\Pi_{69}$:
\begin{ptree}
\AXC{$\forall a\,\forall X\,(\mathrm{UB}(a,X) \leftrightarrow a \in \mathbb{R} \land \forall x \in X\,(x \leq a))$\lbl{DefUB}}
\AXC{$\varphi_{109}$\lbl{A2}}
\BIC{$\forall X\,(\mathrm{UB}(M,X) \leftrightarrow M \in \mathbb{R} \land \forall x \in X\,(x \leq M))$}
\AXC{$\forall X\,(\mathrm{UB}(M,X) \leftrightarrow M \in \mathbb{R} \land \forall x \in X\,(x \leq M)) \to \varphi_{8}$\lbl{A2}}
\BIC{$\varphi_{8}$}
\end{ptree}
\noindent$\Pi_{70}$:
\begin{ptree}
\AXC{$\mathrm{MonoInc}(a) \to M \in \mathbb{R} \to M \in \mathbb{R}$\lbl{A1}}
\AXC{$\stackrel{\Pi_{69}}{\vdots}$}\noLine
\UIC{$\varphi_{8}$}
\AXC{$\varphi_{8} \to \varphi_{222}$\lbl{A1}}
\BIC{$\varphi_{222}$}
\BIC{$(\mathrm{MonoInc}(a) \to \varphi_{33}) \to \varphi_{40} \to \varphi_{39} \to \varphi_{43}$}
\end{ptree}
\noindent$\Pi_{71}$:
\begin{ptree}
\AXC{$\mathrm{MonoInc}(a) \land M \in \mathbb{R} \land \forall n \in \mathbb{N}\,(a_{n} \leq M) \to \varphi_{10}$\lbl{A4}}
\AXC{$\stackrel{\Pi_{68}}{\vdots}$}\noLine
\UIC{$\varphi_{186}$}
\BIC{$\mathrm{MonoInc}(a) \to \varphi_{33}$}
\AXC{$\stackrel{\Pi_{70}}{\vdots}$}\noLine
\UIC{$(\mathrm{MonoInc}(a) \to \varphi_{33}) \to \varphi_{40} \to \varphi_{39} \to \varphi_{43}$}
\BIC{$\varphi_{40} \to \varphi_{39} \to \varphi_{43}$}
\end{ptree}
\noindent$\Pi_{72}$:
\begin{ptree}
\AXC{$\varphi_{47}$\lbl{A4}}
\AXC{$\forall y\,(\varphi_{30}) \to \varphi_{52}$\lbl{A3}}
\AXC{$\varphi_{30}$\lbl{A2}}
\UIC{$\forall y\,(\varphi_{30})$}
\AXC{$\forall y\,(\varphi_{30}) \to (\forall y\,(\varphi_{30}) \to \varphi_{52}) \to \varphi_{47} \to \varphi_{40}$\lbl{A1}}
\BIC{$(\forall y\,(\varphi_{30}) \to \varphi_{52}) \to \varphi_{47} \to \varphi_{40}$}
\BIC{$\varphi_{47} \to \varphi_{40}$}
\BIC{$\varphi_{40}$}
\AXC{$\stackrel{\Pi_{71}}{\vdots}$}\noLine
\UIC{$\varphi_{40} \to \varphi_{39} \to \varphi_{43}$}
\BIC{$\varphi_{39} \to \varphi_{43}$}
\end{ptree}
\noindent$\Pi_{73}$:
\begin{ptree}
\AXC{$\varphi_{46}$\lbl{A4}}
\AXC{$\forall x\,(\varphi_{29}) \to \varphi_{51}$\lbl{A3}}
\AXC{$\varphi_{29}$\lbl{A2}}
\UIC{$\forall x\,(\varphi_{29})$}
\AXC{$\forall x\,(\varphi_{29}) \to (\forall x\,(\varphi_{29}) \to \varphi_{51}) \to \varphi_{46} \to \varphi_{39}$\lbl{A1}}
\BIC{$(\forall x\,(\varphi_{29}) \to \varphi_{51}) \to \varphi_{46} \to \varphi_{39}$}
\BIC{$\varphi_{46} \to \varphi_{39}$}
\BIC{$\varphi_{39}$}
\AXC{$\stackrel{\Pi_{72}}{\vdots}$}\noLine
\UIC{$\varphi_{39} \to \varphi_{43}$}
\BIC{$\varphi_{43}$}
\end{ptree}
\noindent$\Pi_{74}$:
\begin{ptree}
\AXC{$x = x$\lbl{A5}}
\UIC{$\forall x\,(x = x)$}
\AXC{$\forall x\,(x = x) \to a_{1_{\omega}} = a_{1_{\omega}}$\lbl{A2}}
\BIC{$a_{1_{\omega}} = a_{1_{\omega}}$}
\AXC{$1_{\omega} \in \mathbb{N}$\lbl{[N1]}}
\AXC{$1_{\omega} \in \mathbb{N} \to \varphi_{139}$\lbl{A1}}
\BIC{$\varphi_{139}$}
\BIC{$\mathrm{MonoInc}(a) \to M \in \mathbb{R} \to \varphi_{44}$}
\end{ptree}
\noindent$\Pi_{75}$:
\begin{ptree}
\AXC{$\stackrel{\Pi_{9}}{\vdots}$}\noLine
\UIC{$\mathrm{MonoInc}(a) \to M \in \mathbb{R} \to \forall n \in \mathbb{N}\,(a_{n} \leq M) \to \varphi_{32}$}
\AXC{$\varphi_{199}$\lbl{A1}}
\BIC{$\mathrm{MonoInc}(a) \to M \in \mathbb{R} \to \varphi_{65}$}
\AXC{$(\mathrm{MonoInc}(a) \to M \in \mathbb{R} \to \varphi_{65}) \to \varphi_{203}$\lbl{A1}}
\BIC{$\varphi_{203}$}
\end{ptree}
\noindent$\Pi_{76}$:
\begin{ptree}
\AXC{$\forall y\,(y \in \mathrm{ran}(a) \leftrightarrow \exists j \in \mathbb{N}\,(y = a_{j})) \to \varphi_{26}$\lbl{A2}}
\AXC{$\stackrel{\Pi_{75}}{\vdots}$}\noLine
\UIC{$\varphi_{203}$}
\BIC{$\mathrm{MonoInc}(a) \to M \in \mathbb{R} \to \forall n \in \mathbb{N}\,(a_{n} \leq M) \to \varphi_{26}$}
\AXC{$\varphi_{216}$\lbl{A1}}
\BIC{$(\mathrm{MonoInc}(a) \to \varphi_{78}) \to \mathrm{MonoInc}(a) \to \varphi_{35}$}
\end{ptree}
\noindent$\Pi_{77}$:
\begin{ptree}
\AXC{$\varphi_{80}$\lbl{A2}}
\AXC{$\stackrel{\Pi_{74}}{\vdots}$}\noLine
\UIC{$\mathrm{MonoInc}(a) \to M \in \mathbb{R} \to \varphi_{44}$}
\AXC{$\varphi_{213}$\lbl{A1}}
\BIC{$\varphi_{80} \to \mathrm{MonoInc}(a) \to \varphi_{78}$}
\BIC{$\mathrm{MonoInc}(a) \to \varphi_{78}$}
\AXC{$\stackrel{\Pi_{76}}{\vdots}$}\noLine
\UIC{$(\mathrm{MonoInc}(a) \to \varphi_{78}) \to \mathrm{MonoInc}(a) \to \varphi_{35}$}
\BIC{$\mathrm{MonoInc}(a) \to \varphi_{35}$}
\end{ptree}
\noindent$\Pi_{78}$:
\begin{ptree}
\AXC{$\stackrel{\Pi_{45}}{\vdots}$}\noLine
\UIC{$\varphi_{59}$}
\AXC{$\varphi_{20}$\lbl{(R16)}}
\AXC{$\varphi_{20} \to \mathrm{ran}(a) \subseteq \mathbb{R} \to \mathrm{ran}(a) \neq \emptyset \to \varphi_{3}$\lbl{A2}}
\BIC{$\mathrm{ran}(a) \subseteq \mathbb{R} \to \mathrm{ran}(a) \neq \emptyset \to \varphi_{3}$}
\AXC{$(\mathrm{ran}(a) \subseteq \mathbb{R} \to \mathrm{ran}(a) \neq \emptyset \to \varphi_{3}) \to \varphi_{59} \to \varphi_{60} \to \varphi_{150}$\lbl{A1}}
\BIC{$\varphi_{59} \to \varphi_{60} \to \varphi_{150}$}
\BIC{$\varphi_{60} \to \varphi_{150}$}
\end{ptree}
\noindent$\Pi_{79}$:
\begin{ptree}
\AXC{$\stackrel{\Pi_{77}}{\vdots}$}\noLine
\UIC{$\mathrm{MonoInc}(a) \to \varphi_{35}$}
\AXC{$\forall X\,\forall x \in X\,(X \neq \emptyset)$\lbl{[B2]}}
\AXC{$\varphi_{38}$\lbl{A2}}
\BIC{$\forall x \in \mathrm{ran}(a)\,(\mathrm{ran}(a) \neq \emptyset)$}
\AXC{$\varphi_{81}$\lbl{A2}}
\BIC{$a_{1_{\omega}} \in \mathrm{ran}(a) \to \mathrm{ran}(a) \neq \emptyset$}
\AXC{$\varphi_{191}$\lbl{A1}}
\BIC{$(\mathrm{MonoInc}(a) \to \varphi_{35}) \to \varphi_{60}$}
\BIC{$\varphi_{60}$}
\AXC{$\stackrel{\Pi_{78}}{\vdots}$}\noLine
\UIC{$\varphi_{60} \to \varphi_{150}$}
\BIC{$\varphi_{150}$}
\end{ptree}
\noindent$\Pi_{80}$:
\begin{ptree}
\AXC{$\varphi_{1}$\lbl{A2}}
\AXC{$\stackrel{\Pi_{73}}{\vdots}$}\noLine
\UIC{$\varphi_{43}$}
\AXC{$\varphi_{43} \to \varphi_{1} \to \mathrm{MonoInc}(a) \to \varphi_{17}$\lbl{A1}}
\BIC{$\varphi_{1} \to \mathrm{MonoInc}(a) \to \varphi_{17}$}
\BIC{$\mathrm{MonoInc}(a) \to \varphi_{17}$}
\AXC{$\stackrel{\Pi_{79}}{\vdots}$}\noLine
\UIC{$\varphi_{150}$}
\BIC{$\mathrm{MonoInc}(a) \to M \in \mathbb{R} \to \varphi_{6}$}
\AXC{$(\mathrm{MonoInc}(a) \to M \in \mathbb{R} \to \varphi_{6}) \to \varphi_{254}$\lbl{A1}}
\BIC{$\varphi_{254}$}
\end{ptree}
\noindent 主木:
\begin{ptree}
\AXC{$\varphi_{55} \to \forall s\,(\varphi_{55})$\lbl{A4}}
\AXC{$\varphi_{187}$\lbl{A3}}
\AXC{$\forall \alpha\,(\neg \mathrm{Conv}(a,\alpha)) \to \neg \mathrm{Conv}(a,s)$\lbl{A2}}
\AXC{$\stackrel{\Pi_{54}}{\vdots}$}\noLine
\UIC{$\varphi_{137}$}
\BIC{$\mathrm{MonoInc}(a) \to M \in \mathbb{R} \to \varphi_{25}$}
\AXC{$\varphi_{182}$\lbl{A1}}
\BIC{$\varphi_{55} \to \neg \mathrm{IsSup}(s,\mathrm{ran}(a))$}
\UIC{$\forall s\,(\varphi_{55} \to \neg \mathrm{IsSup}(s,\mathrm{ran}(a)))$}
\AXC{$\stackrel{\Pi_{80}}{\vdots}$}\noLine
\UIC{$\varphi_{254}$}
\BIC{$\varphi_{187} \to (\varphi_{55} \to \forall s\,(\varphi_{55})) \to \varphi_{49}$}
\BIC{$(\varphi_{55} \to \forall s\,(\varphi_{55})) \to \varphi_{49}$}
\BIC{$\varphi_{49}$}
\UIC{$\forall M\,(\varphi_{49})$}
\UIC{$\forall a\,\forall M\,(\varphi_{49})$}
\end{ptree}
\begin{abbrevlist}
\item $\varphi_{1} :\equiv \forall b\,(\neg \mathrm{UB}(b,\mathrm{ran}(a))) \to \neg \mathrm{UB}(M,\mathrm{ran}(a))$
\item $\varphi_{2} :\equiv \mathrm{IsSup}(s,\mathrm{ran}(a)) \to \mathrm{IsSup}(s,\mathrm{ran}(a))$
\item $\varphi_{3} :\equiv \exists b\,(\mathrm{UB}(b,\mathrm{ran}(a))) \to \exists s\,(\mathrm{IsSup}(s,\mathrm{ran}(a)))$
\item $\varphi_{4} :\equiv \forall X\,(\mathrm{IsSup}(s,X) \leftrightarrow \mathrm{UB}(s,X) \land \forall w\,(\mathrm{UB}(w,X) \to s \leq w))$
\item $\varphi_{5} :\equiv \forall y\,(v_{718} = y \to y = v_{718}) \to v_{718} = v_{719} \to v_{719} = v_{718}$
\item $\varphi_{6} :\equiv \forall n \in \mathbb{N}\,(a_{n} \leq M) \to \exists s\,(\mathrm{IsSup}(s,\mathrm{ran}(a)))$
\item $\varphi_{7} :\equiv \mathrm{UB}(s,\mathrm{ran}(a)) \leftrightarrow s \in \mathbb{R} \land \forall x \in \mathrm{ran}(a)\,(x \leq s)$
\item $\varphi_{8} :\equiv \mathrm{UB}(M,\mathrm{ran}(a)) \leftrightarrow M \in \mathbb{R} \land \forall x \in \mathrm{ran}(a)\,(x \leq M)$
\item $\varphi_{9} :\equiv \mathrm{IsSup}(s,\mathrm{ran}(a)) \leftrightarrow \mathrm{UB}(s,\mathrm{ran}(a)) \land \forall w\,(\mathrm{UB}(w,\mathrm{ran}(a)) \to s \leq w)$
\item $\varphi_{10} :\equiv \forall y\,(\mathrm{MonoInc}(a) \land M \in \mathbb{R} \land \forall n \in \mathbb{N}\,(a_{n} \leq M))$
\item $\varphi_{11} :\equiv \forall i_{0}\,(\neg (i_{0} \in \mathbb{N} \land x = a_{i_{0}})) \to \neg (j \in \mathbb{N} \land x = a_{j})$
\item $\varphi_{12} :\equiv \forall j\,(\forall n\,(\neg (n \in \mathbb{N} \land y = a_{n})) \to \neg (j \in \mathbb{N} \land y = a_{j}))$
\item $\varphi_{13} :\equiv \forall n \in \mathbb{N}\,(a_{n} \leq M) \to \mathrm{IsSup}(s,\mathrm{ran}(a)) \to \mathrm{Conv}(a,s)$
\item $\varphi_{14} :\equiv y \in \mathrm{ran}(a) \to (y \in \mathrm{ran}(a) \leftrightarrow \exists j \in \mathbb{N}\,(y = a_{j}))$
\item $\varphi_{15} :\equiv \forall z\,\forall w\,(z = w \to z = M \to w = M) \to \forall w\,(a_{n} = w \to a_{n} = M \to w = M)$
\item $\varphi_{16} :\equiv \forall n \in \mathbb{N}\,(a_{n} \leq M) \to \mathrm{IsSup}(s,\mathrm{ran}(a)) \to s \in \mathbb{R}$
\item $\varphi_{17} :\equiv M \in \mathbb{R} \to \forall n \in \mathbb{N}\,(a_{n} \leq M) \to \exists b\,(\mathrm{UB}(b,\mathrm{ran}(a)))$
\item $\varphi_{18} :\equiv \mathrm{MonoInc}(a) \leftrightarrow \mathrm{Seq}(a) \land \forall n \in \mathbb{N}\,(a_{n} \leq a_{S(n)})$
\item $\varphi_{19} :\equiv \forall v_{718}\,\forall v_{719}\,(v_{718} = v_{719} \to v_{719} = v_{718}) \to \forall v_{719}\,(y = v_{719} \to v_{719} = y)$
\item $\varphi_{20} :\equiv \forall X\,(X \subseteq \mathbb{R} \to X \neq \emptyset \to \exists b\,(\mathrm{UB}(b,X)) \to \exists s\,(\mathrm{IsSup}(s,X)))$
\item $\varphi_{21} :\equiv \forall n\,(\neg (n \in \mathbb{N} \land y = a_{n})) \to \forall j\,\forall n\,(\neg (n \in \mathbb{N} \land y = a_{n}))$
\item $\varphi_{22} :\equiv \forall j\,\forall n\,(\neg (n \in \mathbb{N} \land y = a_{n})) \to \forall j\,(\neg (j \in \mathbb{N} \land y = a_{j}))$
\item $\varphi_{23} :\equiv y \in \mathrm{ran}(a) \to \forall y\,(y \in \mathrm{ran}(a) \leftrightarrow \exists j \in \mathbb{N}\,(y = a_{j}))$
\item $\varphi_{24} :\equiv \mathrm{MonoInc}(a) \to M \in \mathbb{R} \to \forall n \in \mathbb{N}\,(a_{n} \leq M) \to \mathrm{Seq}(a)$
\item $\varphi_{25} :\equiv \forall n \in \mathbb{N}\,(a_{n} \leq M) \to \mathrm{IsSup}(s,\mathrm{ran}(a)) \to \exists \alpha\,(\mathrm{Conv}(a,\alpha))$
\item $\varphi_{26} :\equiv a_{1_{\omega}} \in \mathrm{ran}(a) \leftrightarrow \exists j \in \mathbb{N}\,(a_{1_{\omega}} = a_{j})$
\item $\varphi_{27} :\equiv \forall j\,\forall i_{0}\,(\neg (i_{0} \in \mathbb{N} \land x = a_{i_{0}})) \to \forall j\,(\neg (j \in \mathbb{N} \land x = a_{j}))$
\item $\varphi_{28} :\equiv n \in \mathbb{N} \to N \leq n \to \forall y\,(y \in \mathrm{ran}(a) \leftrightarrow \exists j \in \mathbb{N}\,(y = a_{j}))$
\item $\varphi_{29} :\equiv \forall y \in \mathrm{ran}(a)\,(\neg y < M \to y = M) \to x \in \mathrm{ran}(a) \to \neg x < M \to x = M$
\item $\varphi_{30} :\equiv \forall x \in \mathrm{ran}(a)\,(\neg x < M \to x = M) \to y \in \mathrm{ran}(a) \to \neg y < M \to y = M$
\item $\varphi_{31} :\equiv \forall N\,(\neg (N \in \mathbb{N} \land \forall n \in \mathbb{N}\,(N \leq n \to \lvert a_{n} - s \rvert < \varepsilon)))$
\item $\varphi_{32} :\equiv \mathrm{ran}(a) \subseteq \mathbb{R} \land \forall y\,(y \in \mathrm{ran}(a) \leftrightarrow \exists j \in \mathbb{N}\,(y = a_{j}))$
\item $\varphi_{33} :\equiv M \in \mathbb{R} \to \forall n \in \mathbb{N}\,(a_{n} \leq M) \to \forall y \in \mathrm{ran}(a)\,(y \leq M)$
\item $\varphi_{34} :\equiv \forall n\,(\mathrm{MonoInc}(a) \to N \in \mathbb{N} \to n \in \mathbb{N} \to N \leq n \to a_{N} \leq a_{n})$
\item $\varphi_{35} :\equiv M \in \mathbb{R} \to \forall n \in \mathbb{N}\,(a_{n} \leq M) \to a_{1_{\omega}} \in \mathrm{ran}(a)$
\item $\varphi_{36} :\equiv \forall i_{0}\,(\neg (i_{0} \in \mathbb{N} \land x = a_{i_{0}})) \to \forall j\,\forall i_{0}\,(\neg (i_{0} \in \mathbb{N} \land x = a_{i_{0}}))$
\item $\varphi_{37} :\equiv \forall w\,(a_{n} = w \to a_{n} = M \to w = M) \to a_{n} = y \to a_{n} = M \to y = M$
\item $\varphi_{38} :\equiv \forall X\,\forall x \in X\,(X \neq \emptyset) \to \forall x \in \mathrm{ran}(a)\,(\mathrm{ran}(a) \neq \emptyset)$
\item $\varphi_{39} :\equiv \forall y \in \mathrm{ran}(a)\,(\neg y < M \to y = M) \to \forall x \in \mathrm{ran}(a)\,(\neg x < M \to x = M)$
\item $\varphi_{40} :\equiv \forall x \in \mathrm{ran}(a)\,(\neg x < M \to x = M) \to \forall y \in \mathrm{ran}(a)\,(\neg y < M \to y = M)$
\item $\varphi_{41} :\equiv \forall j\,(\neg (j \in \mathbb{N} \land a_{n} = a_{j})) \to \neg (n \in \mathbb{N} \land a_{n} = a_{n})$
\item $\varphi_{42} :\equiv \forall m\,\forall n\,(\mathrm{MonoInc}(a) \to m \in \mathbb{N} \to n \in \mathbb{N} \to m \leq n \to a_{m} \leq a_{n})$
\item $\varphi_{43} :\equiv \mathrm{MonoInc}(a) \to M \in \mathbb{R} \to \forall n \in \mathbb{N}\,(a_{n} \leq M) \to \mathrm{UB}(M,\mathrm{ran}(a))$
\item $\varphi_{44} :\equiv \forall n \in \mathbb{N}\,(a_{n} \leq M) \to 1_{\omega} \in \mathbb{N} \land a_{1_{\omega}} = a_{1_{\omega}}$
\item $\varphi_{45} :\equiv \forall n \in \mathbb{N}\,(a_{n} \leq M) \to \mathrm{IsSup}(s,\mathrm{ran}(a)) \to \forall x \in \mathrm{ran}(a)\,(x \leq s)$
\item $\varphi_{46} :\equiv \forall y \in \mathrm{ran}(a)\,(\neg y < M \to y = M) \to \forall x\,\forall y \in \mathrm{ran}(a)\,(\neg y < M \to y = M)$
\item $\varphi_{47} :\equiv \forall x \in \mathrm{ran}(a)\,(\neg x < M \to x = M) \to \forall y\,\forall x \in \mathrm{ran}(a)\,(\neg x < M \to x = M)$
\item $\varphi_{48} :\equiv \forall z\,\forall w\,(z = w \to z < M \to w < M) \to \forall w\,(a_{n} = w \to a_{n} < M \to w < M)$
\item $\varphi_{49} :\equiv \mathrm{MonoInc}(a) \to M \in \mathbb{R} \to \forall n \in \mathbb{N}\,(a_{n} \leq M) \to \exists \alpha\,(\mathrm{Conv}(a,\alpha))$
\item $\varphi_{50} :\equiv x \in \mathrm{ran}(a) \land s - \varepsilon < x \to i_{0} \in \mathbb{N} \land x = a_{i_{0}} \to i_{0} \in \mathbb{N}$
\item $\varphi_{51} :\equiv \forall x\,\forall y \in \mathrm{ran}(a)\,(\neg y < M \to y = M) \to \forall x \in \mathrm{ran}(a)\,(\neg x < M \to x = M)$
\item $\varphi_{52} :\equiv \forall y\,\forall x \in \mathrm{ran}(a)\,(\neg x < M \to x = M) \to \forall y \in \mathrm{ran}(a)\,(\neg y < M \to y = M)$
\item $\varphi_{53} :\equiv M \in \mathbb{R} \to \forall n \in \mathbb{N}\,(a_{n} \leq M) \to y \in \mathrm{ran}(a) \to y \in \mathrm{ran}(a)$
\item $\varphi_{54} :\equiv \forall n \in \mathbb{N}\,(a_{n} \leq M) \to y \in \mathrm{ran}(a) \to \exists j \in \mathbb{N}\,(y = a_{j})$
\item $\varphi_{55} :\equiv \mathrm{MonoInc}(a) \land M \in \mathbb{R} \land \forall n \in \mathbb{N}\,(a_{n} \leq M) \land \neg \exists \alpha\,(\mathrm{Conv}(a,\alpha))$
\item $\varphi_{56} :\equiv \mathrm{MonoInc}(a) \land M \in \mathbb{R} \land \forall n \in \mathbb{N}\,(a_{n} \leq M) \land \mathrm{IsSup}(s,\mathrm{ran}(a))$
\item $\varphi_{57} :\equiv x \in \mathrm{ran}(a) \land s - \varepsilon < x \to i_{0} \in \mathbb{N} \land x = a_{i_{0}} \to x = a_{i_{0}}$
\item $\varphi_{58} :\equiv \varepsilon \in \mathbb{R} \to 0 < \varepsilon \to x \in \mathrm{ran}(a) \land s - \varepsilon < x \to x \in \mathrm{ran}(a)$
\item $\varphi_{59} :\equiv \mathrm{MonoInc}(a) \to M \in \mathbb{R} \to \forall n \in \mathbb{N}\,(a_{n} \leq M) \to \mathrm{ran}(a) \subseteq \mathbb{R}$
\item $\varphi_{60} :\equiv \mathrm{MonoInc}(a) \to M \in \mathbb{R} \to \forall n \in \mathbb{N}\,(a_{n} \leq M) \to \mathrm{ran}(a) \neq \emptyset$
\item $\varphi_{61} :\equiv 0 < \varepsilon \to x \in \mathrm{ran}(a) \land s - \varepsilon < x \to x \in \mathrm{ran}(a) \land s - \varepsilon < x$
\item $\varphi_{62} :\equiv x \in \mathrm{ran}(a) \land s - \varepsilon < x \to i_{0} \in \mathbb{N} \land x = a_{i_{0}} \to s - \varepsilon < x$
\item $\varphi_{63} :\equiv x \in \mathrm{ran}(a) \land s - \varepsilon < x \to (x \in \mathrm{ran}(a) \leftrightarrow \exists j \in \mathbb{N}\,(x = a_{j}))$
\item $\varphi_{64} :\equiv \forall n \in \mathbb{N}\,(a_{n} \leq M) \to y \in \mathrm{ran}(a) \to n \in \mathbb{N} \land y = a_{n} \to y \leq M$
\item $\varphi_{65} :\equiv \forall n \in \mathbb{N}\,(a_{n} \leq M) \to \forall y\,(y \in \mathrm{ran}(a) \leftrightarrow \exists j \in \mathbb{N}\,(y = a_{j}))$
\item $\varphi_{66} :\equiv \mathrm{MonoInc}(a) \to M \in \mathbb{R} \to \forall n \in \mathbb{N}\,(a_{n} \leq M) \to y \in \mathrm{ran}(a) \to y \leq M$
\item $\varphi_{67} :\equiv \mathrm{MonoInc}(a) \land M \in \mathbb{R} \land \forall n \in \mathbb{N}\,(a_{n} \leq M) \to y \in \mathrm{ran}(a) \to y \leq M$
\item $\varphi_{68} :\equiv n \in \mathbb{N} \to N \leq n \to (a_{n} \in \mathrm{ran}(a) \leftrightarrow \exists j \in \mathbb{N}\,(a_{n} = a_{j}))$
\item $\varphi_{69} :\equiv \forall n \in \mathbb{N}\,(a_{n} \leq M) \to \mathrm{IsSup}(s,\mathrm{ran}(a)) \to \varepsilon \in \mathbb{R} \to \varepsilon \in \mathbb{R}$
\item $\varphi_{70} :\equiv \mathrm{MonoInc}(a) \land M \in \mathbb{R} \land \forall n \in \mathbb{N}\,(a_{n} \leq M) \land y \in \mathrm{ran}(a) \land \neg y \leq M$
\item $\varphi_{71} :\equiv \forall x\,\forall y\,\forall z\,(x \leq y \to y < z \to x < z) \to \forall y\,\forall z\,(a_{n} \leq y \to y < z \to a_{n} < z)$
\item $\varphi_{72} :\equiv \forall x\,\forall y\,\forall z\,(x < y \to y \leq z \to x < z) \to \forall y\,\forall z\,(s - \varepsilon < y \to y \leq z \to s - \varepsilon < z)$
\item $\varphi_{73} :\equiv \forall N\,(\neg (N \in \mathbb{N} \land s - \varepsilon < a_{N})) \to \neg (i_{0} \in \mathbb{N} \land s - \varepsilon < a_{i_{0}})$
\item $\varphi_{74} :\equiv x \in \mathrm{ran}(a) \land s - \varepsilon < x \to \forall y\,(y \in \mathrm{ran}(a) \leftrightarrow \exists j \in \mathbb{N}\,(y = a_{j}))$
\item $\varphi_{75} :\equiv \mathrm{IsSup}(s,\mathrm{ran}(a)) \to \varepsilon \in \mathbb{R} \to 0 < \varepsilon \to \exists x \in \mathrm{ran}(a)\,(s - \varepsilon < x)$
\item $\varphi_{76} :\equiv \forall w\,(a_{n} = w \to a_{n} < M \to w < M) \to a_{n} = y \to a_{n} < M \to y < M$
\item $\varphi_{77} :\equiv \mathrm{MonoInc}(a) \to M \in \mathbb{R} \to \forall n \in \mathbb{N}\,(a_{n} \leq M) \to n \in \mathbb{N} \to a_{n} \leq M$
\item $\varphi_{78} :\equiv M \in \mathbb{R} \to \forall n \in \mathbb{N}\,(a_{n} \leq M) \to \exists j \in \mathbb{N}\,(a_{1_{\omega}} = a_{j})$
\item $\varphi_{79} :\equiv \forall n \in \mathbb{N}\,(a_{n} \leq M) \to y \in \mathrm{ran}(a) \to n \in \mathbb{N} \land y = a_{n} \to n \in \mathbb{N}$
\item $\varphi_{80} :\equiv \forall j\,(\neg (j \in \mathbb{N} \land a_{1_{\omega}} = a_{j})) \to \neg (1_{\omega} \in \mathbb{N} \land a_{1_{\omega}} = a_{1_{\omega}})$
\item $\varphi_{81} :\equiv \forall x \in \mathrm{ran}(a)\,(\mathrm{ran}(a) \neq \emptyset) \to a_{1_{\omega}} \in \mathrm{ran}(a) \to \mathrm{ran}(a) \neq \emptyset$
\item $\varphi_{82} :\equiv \forall n \in \mathbb{N}\,(a_{n} \leq M) \to \mathrm{IsSup}(s,\mathrm{ran}(a)) \to a_{n} \in \mathrm{ran}(a) \to a_{n} \leq s$
\item $\varphi_{83} :\equiv \forall n \in \mathbb{N}\,(a_{n} \leq M) \to y \in \mathrm{ran}(a) \to n \in \mathbb{N} \land y = a_{n} \to a_{n} \leq M$
\item $\varphi_{84} :\equiv \forall n \in \mathbb{N}\,(a_{n} \leq M) \to y \in \mathrm{ran}(a) \to n \in \mathbb{N} \land y = a_{n} \to y = a_{n}$
\item $\varphi_{85} :\equiv \forall n \in \mathbb{N}\,(a_{n} \leq M) \to y \in \mathrm{ran}(a) \to n \in \mathbb{N} \land y = a_{n} \to a_{n} = y$
\item $\varphi_{86} :\equiv \forall E\,\forall e\,(E \subseteq \mathbb{R} \to \mathrm{IsSup}(s,E) \to e \in \mathbb{R} \to 0 < e \to \exists x \in E\,(s - e < x))$
\item $\varphi_{87} :\equiv \mathrm{MonoInc}(a) \to M \in \mathbb{R} \to \forall n \in \mathbb{N}\,(a_{n} \leq M) \to \forall n \in \mathbb{N}\,(a_{n} \leq M)$
\item $\varphi_{88} :\equiv x \in \mathrm{ran}(a) \land s - \varepsilon < x \to i_{0} \in \mathbb{N} \land x = a_{i_{0}} \to s - \varepsilon < a_{i_{0}}$
\item $\varphi_{89} :\equiv \forall n \in \mathbb{N}\,(a_{n} \leq M) \to \mathrm{IsSup}(s,\mathrm{ran}(a)) \to \varepsilon \in \mathbb{R} \to 0 < \varepsilon \to 0 < \varepsilon$
\item $\varphi_{90} :\equiv N \in \mathbb{N} \land s - \varepsilon < a_{N} \to \forall n \in \mathbb{N}\,(N \leq n \to \lvert a_{n} - s \rvert < \varepsilon)$
\item $\varphi_{91} :\equiv \forall y\,\forall z\,(a_{n} \leq y \to y < z \to a_{n} < z) \to \forall z\,(a_{n} \leq s \to s < z \to a_{n} < z)$
\item $\varphi_{92} :\equiv \forall n\,(\mathrm{Seq}(a) \to n \in \mathbb{N} \to a_{n} \in \mathbb{R}) \to \mathrm{Seq}(a) \to n \in \mathbb{N} \to a_{n} \in \mathbb{R}$
\item $\varphi_{93} :\equiv \forall z\,\forall w\,(z = w \to s - \varepsilon < z \to s - \varepsilon < w) \to \forall w\,(x = w \to s - \varepsilon < x \to s - \varepsilon < w)$
\item $\varphi_{94} :\equiv N \in \mathbb{N} \land s - \varepsilon < a_{N} \to n \in \mathbb{N} \to N \leq n \to n \in \mathbb{N} \land a_{n} = a_{n}$
\item $\varphi_{95} :\equiv x \in \mathrm{ran}(a) \land s - \varepsilon < x \to i_{0} \in \mathbb{N} \land x = a_{i_{0}} \to i_{0} \in \mathbb{N} \land x = a_{i_{0}}$
\item $\varphi_{96} :\equiv \varepsilon \in \mathbb{R} \to 0 < \varepsilon \to x \in \mathrm{ran}(a) \land s - \varepsilon < x \to \exists j \in \mathbb{N}\,(x = a_{j})$
\item $\varphi_{97} :\equiv i_{0} \in \mathbb{N} \land x = a_{i_{0}} \to N \in \mathbb{N} \land s - \varepsilon < a_{N} \to n \in \mathbb{N} \to n \in \mathbb{N}$
\item $\varphi_{98} :\equiv i_{0} \in \mathbb{N} \land x = a_{i_{0}} \to N \in \mathbb{N} \land s - \varepsilon < a_{N} \to n \in \mathbb{N} \to N \leq n \to N \leq n$
\item $\varphi_{99} :\equiv N \in \mathbb{N} \land s - \varepsilon < a_{N} \to n \in \mathbb{N} \to N \leq n \to \exists j \in \mathbb{N}\,(a_{n} = a_{j})$
\item $\varphi_{100} :\equiv \forall y\,(y \in \mathrm{ran}(a) \leftrightarrow \exists j \in \mathbb{N}\,(y = a_{j})) \to (x \in \mathrm{ran}(a) \leftrightarrow \exists j \in \mathbb{N}\,(x = a_{j}))$
\item $\varphi_{101} :\equiv \forall y\,(y \in \mathrm{ran}(a) \leftrightarrow \exists j \in \mathbb{N}\,(y = a_{j})) \to (y \in \mathrm{ran}(a) \leftrightarrow \exists j \in \mathbb{N}\,(y = a_{j}))$
\item $\varphi_{102} :\equiv \varepsilon \in \mathbb{R} \to 0 < \varepsilon \to \exists N \in \mathbb{N}\,\forall n \in \mathbb{N}\,(N \leq n \to \lvert a_{n} - s \rvert < \varepsilon)$
\item $\varphi_{103} :\equiv \forall a\,\forall n\,(\mathrm{Seq}(a) \to n \in \mathbb{N} \to a_{n} \in \mathbb{R}) \to \forall n\,(\mathrm{Seq}(a) \to n \in \mathbb{N} \to a_{n} \in \mathbb{R})$
\item $\varphi_{104} :\equiv i_{0} \in \mathbb{N} \land x = a_{i_{0}} \to \exists N \in \mathbb{N}\,\forall n \in \mathbb{N}\,(N \leq n \to \lvert a_{n} - s \rvert < \varepsilon)$
\item $\varphi_{105} :\equiv \forall e\,\forall y\,(s \in \mathbb{R} \to e \in \mathbb{R} \to y \in \mathbb{R} \to s - e < y \to y < s + e \to \lvert y - s \rvert < e)$
\item $\varphi_{106} :\equiv i_{0} \in \mathbb{N} \land x = a_{i_{0}} \to N \in \mathbb{N} \land s - \varepsilon < a_{N} \to n \in \mathbb{N} \to N \leq n \to N \in \mathbb{N}$
\item $\varphi_{107} :\equiv \forall a\,\forall X\,(\mathrm{UB}(a,X) \leftrightarrow a \in \mathbb{R} \land \forall x \in X\,(x \leq a)) \to \forall X\,(\mathrm{UB}(s,X) \leftrightarrow s \in \mathbb{R} \land \forall x \in X\,(x \leq s))$
\item $\varphi_{108} :\equiv \forall z\,(a_{n} \leq s \to s < z \to a_{n} < z) \to a_{n} \leq s \to s < s + \varepsilon \to a_{n} < s + \varepsilon$
\item $\varphi_{109} :\equiv \forall a\,\forall X\,(\mathrm{UB}(a,X) \leftrightarrow a \in \mathbb{R} \land \forall x \in X\,(x \leq a)) \to \forall X\,(\mathrm{UB}(M,X) \leftrightarrow M \in \mathbb{R} \land \forall x \in X\,(x \leq M))$
\item $\varphi_{110} :\equiv i_{0} \in \mathbb{N} \land x = a_{i_{0}} \to N \in \mathbb{N} \land s - \varepsilon < a_{N} \to n \in \mathbb{N} \to N \leq n \to a_{n} \leq s$
\item $\varphi_{111} :\equiv \forall n \in \mathbb{N}\,(a_{n} \leq M) \to y \in \mathrm{ran}(a) \to n \in \mathbb{N} \land y = a_{n} \to n \in \mathbb{N} \land y = a_{n}$
\item $\varphi_{112} :\equiv \forall e\,(\mathrm{ran}(a) \subseteq \mathbb{R} \to \mathrm{IsSup}(s,\mathrm{ran}(a)) \to e \in \mathbb{R} \to 0 < e \to \exists x \in \mathrm{ran}(a)\,(s - e < x))$
\item $\varphi_{113} :\equiv \forall x\,\forall e\,\forall y\,(x \in \mathbb{R} \to e \in \mathbb{R} \to y \in \mathbb{R} \to x - e < y \to y < x + e \to \lvert y - x \rvert < e)$
\item $\varphi_{114} :\equiv \forall w\,(x = w \to s - \varepsilon < x \to s - \varepsilon < w) \to x = a_{i_{0}} \to s - \varepsilon < x \to s - \varepsilon < a_{i_{0}}$
\item $\varphi_{115} :\equiv \forall \varepsilon \in \mathbb{R}\,(0 < \varepsilon \to \exists N \in \mathbb{N}\,\forall n \in \mathbb{N}\,(N \leq n \to \lvert a_{n} - s \rvert < \varepsilon))$
\item $\varphi_{116} :\equiv x \in \mathrm{ran}(a) \land s - \varepsilon < x \to i_{0} \in \mathbb{N} \land x = a_{i_{0}} \to i_{0} \in \mathbb{N} \land s - \varepsilon < a_{i_{0}}$
\item $\varphi_{117} :\equiv x \in \mathrm{ran}(a) \land s - \varepsilon < x \to \exists N \in \mathbb{N}\,\forall n \in \mathbb{N}\,(N \leq n \to \lvert a_{n} - s \rvert < \varepsilon)$
\item $\varphi_{118} :\equiv i_{0} \in \mathbb{N} \land x = a_{i_{0}} \to N \in \mathbb{N} \land s - \varepsilon < a_{N} \to n \in \mathbb{N} \to N \leq n \to s < s + \varepsilon$
\item $\varphi_{119} :\equiv N \in \mathbb{N} \land s - \varepsilon < a_{N} \to N \in \mathbb{N} \land \forall n \in \mathbb{N}\,(N \leq n \to \lvert a_{n} - s \rvert < \varepsilon)$
\item $\varphi_{120} :\equiv x \in \mathrm{ran}(a) \land s - \varepsilon < x \to i_{0} \in \mathbb{N} \land x = a_{i_{0}} \to \exists N \in \mathbb{N}\,(s - \varepsilon < a_{N})$
\item $\varphi_{121} :\equiv \forall y\,(s \in \mathbb{R} \to \varepsilon \in \mathbb{R} \to y \in \mathbb{R} \to s - \varepsilon < y \to y < s + \varepsilon \to \lvert y - s \rvert < \varepsilon)$
\item $\varphi_{122} :\equiv i_{0} \in \mathbb{N} \land x = a_{i_{0}} \to N \in \mathbb{N} \land s - \varepsilon < a_{N} \to N \in \mathbb{N} \land s - \varepsilon < a_{N}$
\item $\varphi_{123} :\equiv \forall y\,\forall z\,(s - \varepsilon < y \to y \leq z \to s - \varepsilon < z) \to \forall z\,(s - \varepsilon < a_{N} \to a_{N} \leq z \to s - \varepsilon < z)$
\item $\varphi_{124} :\equiv N \in \mathbb{N} \land s - \varepsilon < a_{N} \to \exists N \in \mathbb{N}\,\forall n \in \mathbb{N}\,(N \leq n \to \lvert a_{n} - s \rvert < \varepsilon)$
\item $\varphi_{125} :\equiv i_{0} \in \mathbb{N} \land x = a_{i_{0}} \to N \in \mathbb{N} \land s - \varepsilon < a_{N} \to n \in \mathbb{N} \to N \leq n \to a_{N} \leq a_{n}$
\item $\varphi_{126} :\equiv i_{0} \in \mathbb{N} \land x = a_{i_{0}} \to N \in \mathbb{N} \land s - \varepsilon < a_{N} \to n \in \mathbb{N} \to N \leq n \to a_{n} \in \mathbb{R}$
\item $\varphi_{127} :\equiv i_{0} \in \mathbb{N} \land x = a_{i_{0}} \to N \in \mathbb{N} \land s - \varepsilon < a_{N} \to n \in \mathbb{N} \to N \leq n \to a_{n} \in \mathrm{ran}(a)$
\item $\varphi_{128} :\equiv \forall y\,(y \in \mathrm{ran}(a) \leftrightarrow \exists j \in \mathbb{N}\,(y = a_{j})) \to (a_{n} \in \mathrm{ran}(a) \leftrightarrow \exists j \in \mathbb{N}\,(a_{n} = a_{j}))$
\item $\varphi_{129} :\equiv i_{0} \in \mathbb{N} \land x = a_{i_{0}} \to N \in \mathbb{N} \land s - \varepsilon < a_{N} \to n \in \mathbb{N} \to N \leq n \to a_{n} < s + \varepsilon$
\item $\varphi_{130} :\equiv i_{0} \in \mathbb{N} \land x = a_{i_{0}} \to N \in \mathbb{N} \land s - \varepsilon < a_{N} \to n \in \mathbb{N} \to N \leq n \to s - \varepsilon < a_{N}$
\item $\varphi_{131} :\equiv i_{0} \in \mathbb{N} \land x = a_{i_{0}} \to N \in \mathbb{N} \land s - \varepsilon < a_{N} \to n \in \mathbb{N} \to N \leq n \to s - \varepsilon < a_{n}$
\item $\varphi_{132} :\equiv x \in \mathrm{ran}(a) \land s - \varepsilon < x \to i_{0} \in \mathbb{N} \land x = a_{i_{0}} \to N \in \mathbb{N} \land s - \varepsilon < a_{N} \to N \in \mathbb{N}$
\item $\varphi_{133} :\equiv \forall x\,\forall y\,(x \in \mathbb{R} \to y \in \mathbb{R} \to 0 < y \to x < x + y) \to \forall y\,(s \in \mathbb{R} \to y \in \mathbb{R} \to 0 < y \to s < s + y)$
\item $\varphi_{134} :\equiv \forall y\,(s \in \mathbb{R} \to y \in \mathbb{R} \to 0 < y \to s < s + y) \to s \in \mathbb{R} \to \varepsilon \in \mathbb{R} \to 0 < \varepsilon \to s < s + \varepsilon$
\item $\varphi_{135} :\equiv i_{0} \in \mathbb{N} \land x = a_{i_{0}} \to N \in \mathbb{N} \land s - \varepsilon < a_{N} \to n \in \mathbb{N} \to N \leq n \to \lvert a_{n} - s \rvert < \varepsilon$
\item $\varphi_{136} :\equiv \varepsilon \in \mathbb{R} \to a_{n} \in \mathbb{R} \to s - \varepsilon < a_{n} \to a_{n} < s + \varepsilon \to \lvert a_{n} - s \rvert < \varepsilon$
\item $\varphi_{137} :\equiv (\forall \alpha\,(\neg \mathrm{Conv}(a,\alpha)) \to \neg \mathrm{Conv}(a,s)) \to \mathrm{MonoInc}(a) \to M \in \mathbb{R} \to \varphi_{25}$
\item $\varphi_{138} :\equiv \forall z\,(s - \varepsilon < a_{N} \to a_{N} \leq z \to s - \varepsilon < z) \to s - \varepsilon < a_{N} \to a_{N} \leq a_{n} \to s - \varepsilon < a_{n}$
\item $\varphi_{139} :\equiv a_{1_{\omega}} = a_{1_{\omega}} \to \mathrm{MonoInc}(a) \to M \in \mathbb{R} \to \varphi_{44}$
\item $\varphi_{140} :\equiv \varphi_{34} \to \mathrm{MonoInc}(a) \to N \in \mathbb{N} \to n \in \mathbb{N} \to N \leq n \to a_{N} \leq a_{n}$
\item $\varphi_{141} :\equiv M \in \mathbb{R} \to \forall n \in \mathbb{N}\,(a_{n} \leq M) \to \mathrm{IsSup}(s,\mathrm{ran}(a)) \to \varphi_{58}$
\item $\varphi_{142} :\equiv (a_{n} = y \to a_{n} = M \to y = M) \to \mathrm{MonoInc}(a) \to M \in \mathbb{R} \to \varphi_{64}$
\item $\varphi_{143} :\equiv \forall n \in \mathbb{N}\,(a_{n} \leq M) \to \mathrm{IsSup}(s,\mathrm{ran}(a)) \to \varepsilon \in \mathbb{R} \to \varphi_{61}$
\item $\varphi_{144} :\equiv \forall \alpha\,(\mathrm{Conv}(a,\alpha) \leftrightarrow \mathrm{Seq}(a) \land \alpha \in \mathbb{R} \land \forall \varepsilon \in \mathbb{R}\,(0 < \varepsilon \to \exists N \in \mathbb{N}\,\forall n \in \mathbb{N}\,(N \leq n \to \lvert a_{n} - \alpha \rvert < \varepsilon)))$
\item $\varphi_{145} :\equiv \forall n \in \mathbb{N}\,(a_{n} \leq M) \to \mathrm{IsSup}(s,\mathrm{ran}(a)) \to \varepsilon \in \mathbb{R} \to 0 < \varepsilon \to \varphi_{50}$
\item $\varphi_{146} :\equiv \forall n \in \mathbb{N}\,(a_{n} \leq M) \to \mathrm{IsSup}(s,\mathrm{ran}(a)) \to \varepsilon \in \mathbb{R} \to 0 < \varepsilon \to \varphi_{57}$
\item $\varphi_{147} :\equiv M \in \mathbb{R} \to \forall n \in \mathbb{N}\,(a_{n} \leq M) \to \mathrm{IsSup}(s,\mathrm{ran}(a)) \to \varphi_{96}$
\item $\varphi_{148} :\equiv \forall n \in \mathbb{N}\,(a_{n} \leq M) \to \mathrm{IsSup}(s,\mathrm{ran}(a)) \to \varepsilon \in \mathbb{R} \to 0 < \varepsilon \to \varphi_{62}$
\item $\varphi_{149} :\equiv \forall n \in \mathbb{N}\,(a_{n} \leq M) \to \mathrm{IsSup}(s,\mathrm{ran}(a)) \to \varepsilon \in \mathbb{R} \to 0 < \varepsilon \to \varphi_{63}$
\item $\varphi_{150} :\equiv (\mathrm{MonoInc}(a) \to \varphi_{17}) \to \mathrm{MonoInc}(a) \to M \in \mathbb{R} \to \varphi_{6}$
\item $\varphi_{151} :\equiv \forall n \in \mathbb{N}\,(a_{n} \leq M) \to \mathrm{IsSup}(s,\mathrm{ran}(a)) \to \varepsilon \in \mathbb{R} \to 0 < \varepsilon \to \varphi_{74}$
\item $\varphi_{152} :\equiv M \in \mathbb{R} \to \forall n \in \mathbb{N}\,(a_{n} \leq M) \to \mathrm{IsSup}(s,\mathrm{ran}(a)) \to \varphi_{102}$
\item $\varphi_{153} :\equiv 0 < \varepsilon \to x \in \mathrm{ran}(a) \land s - \varepsilon < x \to i_{0} \in \mathbb{N} \land x = a_{i_{0}} \to \varphi_{90}$
\item $\varphi_{154} :\equiv 0 < \varepsilon \to x \in \mathrm{ran}(a) \land s - \varepsilon < x \to i_{0} \in \mathbb{N} \land x = a_{i_{0}} \to \varphi_{94}$
\item $\varphi_{155} :\equiv \forall n \in \mathbb{N}\,(a_{n} \leq M) \to \mathrm{IsSup}(s,\mathrm{ran}(a)) \to \varepsilon \in \mathbb{R} \to 0 < \varepsilon \to \varphi_{88}$
\item $\varphi_{156} :\equiv M \in \mathbb{R} \to \forall n \in \mathbb{N}\,(a_{n} \leq M) \to \mathrm{IsSup}(s,\mathrm{ran}(a)) \to \varphi_{115}$
\item $\varphi_{157} :\equiv 0 < \varepsilon \to x \in \mathrm{ran}(a) \land s - \varepsilon < x \to i_{0} \in \mathbb{N} \land x = a_{i_{0}} \to \varphi_{99}$
\item $\varphi_{158} :\equiv x \in \mathrm{ran}(a) \land s - \varepsilon < x \to i_{0} \in \mathbb{N} \land x = a_{i_{0}} \to N \in \mathbb{N} \land s - \varepsilon < a_{N} \to \varphi_{28}$
\item $\varphi_{159} :\equiv \forall n \in \mathbb{N}\,(a_{n} \leq M) \to \mathrm{IsSup}(s,\mathrm{ran}(a)) \to \varepsilon \in \mathbb{R} \to 0 < \varepsilon \to \varphi_{95}$
\item $\varphi_{160} :\equiv \mathrm{IsSup}(s,\mathrm{ran}(a)) \to \varepsilon \in \mathbb{R} \to 0 < \varepsilon \to x \in \mathrm{ran}(a) \land s - \varepsilon < x \to \varphi_{97}$
\item $\varphi_{161} :\equiv \mathrm{IsSup}(s,\mathrm{ran}(a)) \to \varepsilon \in \mathbb{R} \to 0 < \varepsilon \to x \in \mathrm{ran}(a) \land s - \varepsilon < x \to \varphi_{98}$
\item $\varphi_{162} :\equiv \varphi_{56} \land \varepsilon \in \mathbb{R} \land 0 < \varepsilon \land \neg \exists N \in \mathbb{N}\,\forall n \in \mathbb{N}\,(N \leq n \to \lvert a_{n} - s \rvert < \varepsilon)$
\item $\varphi_{163} :\equiv 0 < \varepsilon \to x \in \mathrm{ran}(a) \land s - \varepsilon < x \to i_{0} \in \mathbb{N} \land x = a_{i_{0}} \to \varphi_{119}$
\item $\varphi_{164} :\equiv (\forall x \in \mathrm{ran}(a)\,(x \leq s) \to a_{n} \in \mathrm{ran}(a) \to a_{n} \leq s) \to \mathrm{MonoInc}(a) \to M \in \mathbb{R} \to \varphi_{82}$
\item $\varphi_{165} :\equiv \mathrm{IsSup}(s,\mathrm{ran}(a)) \to \varepsilon \in \mathbb{R} \to 0 < \varepsilon \to x \in \mathrm{ran}(a) \land s - \varepsilon < x \to \varphi_{104}$
\item $\varphi_{166} :\equiv \mathrm{IsSup}(s,\mathrm{ran}(a)) \to \varepsilon \in \mathbb{R} \to 0 < \varepsilon \to x \in \mathrm{ran}(a) \land s - \varepsilon < x \to \varphi_{106}$
\item $\varphi_{167} :\equiv \forall n \in \mathbb{N}\,(a_{n} \leq M) \to \mathrm{IsSup}(s,\mathrm{ran}(a)) \to \varepsilon \in \mathbb{R} \to 0 < \varepsilon \to \varphi_{116}$
\item $\varphi_{168} :\equiv \forall n \in \mathbb{N}\,(a_{n} \leq M) \to \mathrm{IsSup}(s,\mathrm{ran}(a)) \to \varepsilon \in \mathbb{R} \to 0 < \varepsilon \to \varphi_{117}$
\item $\varphi_{169} :\equiv x \in \mathrm{ran}(a) \land s - \varepsilon < x \to i_{0} \in \mathbb{N} \land x = a_{i_{0}} \to N \in \mathbb{N} \land s - \varepsilon < a_{N} \to \varphi_{68}$
\item $\varphi_{170} :\equiv \mathrm{IsSup}(s,\mathrm{ran}(a)) \to \varepsilon \in \mathbb{R} \to 0 < \varepsilon \to x \in \mathrm{ran}(a) \land s - \varepsilon < x \to \varphi_{110}$
\item $\varphi_{171} :\equiv 0 < \varepsilon \to x \in \mathrm{ran}(a) \land s - \varepsilon < x \to i_{0} \in \mathbb{N} \land x = a_{i_{0}} \to \varphi_{124}$
\item $\varphi_{172} :\equiv \forall n \in \mathbb{N}\,(a_{n} \leq M) \to \mathrm{IsSup}(s,\mathrm{ran}(a)) \to \varepsilon \in \mathbb{R} \to 0 < \varepsilon \to \varphi_{120}$
\item $\varphi_{173} :\equiv \mathrm{IsSup}(s,\mathrm{ran}(a)) \to \varepsilon \in \mathbb{R} \to 0 < \varepsilon \to x \in \mathrm{ran}(a) \land s - \varepsilon < x \to \varphi_{118}$
\item $\varphi_{174} :\equiv \mathrm{IsSup}(s,\mathrm{ran}(a)) \to \varepsilon \in \mathbb{R} \to 0 < \varepsilon \to x \in \mathrm{ran}(a) \land s - \varepsilon < x \to \varphi_{122}$
\item $\varphi_{175} :\equiv \mathrm{IsSup}(s,\mathrm{ran}(a)) \to \varepsilon \in \mathbb{R} \to 0 < \varepsilon \to x \in \mathrm{ran}(a) \land s - \varepsilon < x \to \varphi_{125}$
\item $\varphi_{176} :\equiv \mathrm{IsSup}(s,\mathrm{ran}(a)) \to \varepsilon \in \mathbb{R} \to 0 < \varepsilon \to x \in \mathrm{ran}(a) \land s - \varepsilon < x \to \varphi_{126}$
\item $\varphi_{177} :\equiv \mathrm{IsSup}(s,\mathrm{ran}(a)) \to \varepsilon \in \mathbb{R} \to 0 < \varepsilon \to x \in \mathrm{ran}(a) \land s - \varepsilon < x \to \varphi_{127}$
\item $\varphi_{178} :\equiv \forall n \in \mathbb{N}\,(a_{n} \leq M) \to \mathrm{IsSup}(s,\mathrm{ran}(a)) \to \varepsilon \in \mathbb{R} \to 0 < \varepsilon \to \varphi_{132}$
\item $\varphi_{179} :\equiv \mathrm{IsSup}(s,\mathrm{ran}(a)) \to \varepsilon \in \mathbb{R} \to 0 < \varepsilon \to x \in \mathrm{ran}(a) \land s - \varepsilon < x \to \varphi_{129}$
\item $\varphi_{180} :\equiv \mathrm{IsSup}(s,\mathrm{ran}(a)) \to \varepsilon \in \mathbb{R} \to 0 < \varepsilon \to x \in \mathrm{ran}(a) \land s - \varepsilon < x \to \varphi_{130}$
\item $\varphi_{181} :\equiv \mathrm{IsSup}(s,\mathrm{ran}(a)) \to \varepsilon \in \mathbb{R} \to 0 < \varepsilon \to x \in \mathrm{ran}(a) \land s - \varepsilon < x \to \varphi_{131}$
\item $\varphi_{182} :\equiv (\mathrm{MonoInc}(a) \to M \in \mathbb{R} \to \varphi_{25}) \to \varphi_{55} \to \neg \mathrm{IsSup}(s,\mathrm{ran}(a))$
\item $\varphi_{183} :\equiv (\mathrm{MonoInc}(a) \to \varphi_{53}) \to \mathrm{MonoInc}(a) \to M \in \mathbb{R} \to \varphi_{54}$
\item $\varphi_{184} :\equiv \varphi_{24} \to \mathrm{MonoInc}(a) \to M \in \mathbb{R} \to \forall n \in \mathbb{N}\,(a_{n} \leq M) \to \varphi_{32}$
\item $\varphi_{185} :\equiv \mathrm{IsSup}(s,\mathrm{ran}(a)) \to \varepsilon \in \mathbb{R} \to 0 < \varepsilon \to x \in \mathrm{ran}(a) \land s - \varepsilon < x \to \varphi_{135}$
\item $\varphi_{186} :\equiv (\mathrm{MonoInc}(a) \land M \in \mathbb{R} \land \forall n \in \mathbb{N}\,(a_{n} \leq M) \to \varphi_{10}) \to \mathrm{MonoInc}(a) \to \varphi_{33}$
\item $\varphi_{187} :\equiv \forall s\,(\varphi_{55} \to \neg \mathrm{IsSup}(s,\mathrm{ran}(a))) \to \forall s\,(\varphi_{55}) \to \forall s\,(\neg \mathrm{IsSup}(s,\mathrm{ran}(a)))$
\item $\varphi_{188} :\equiv \varphi_{56} \land \varepsilon \in \mathbb{R} \land 0 < \varepsilon \land (x \in \mathrm{ran}(a) \land s - \varepsilon < x) \land (i_{0} \in \mathbb{N} \land x = a_{i_{0}}) \land (N \in \mathbb{N} \land s - \varepsilon < a_{N})$
\item $\varphi_{189} :\equiv (\mathrm{MonoInc}(a) \to M \in \mathbb{R} \to \forall n \in \mathbb{N}\,(a_{n} \leq M) \to \varphi_{32}) \to \varphi_{59}$
\item $\varphi_{190} :\equiv \varphi_{56} \land \varepsilon \in \mathbb{R} \land 0 < \varepsilon \land (x \in \mathrm{ran}(a) \land s - \varepsilon < x) \land \neg \exists N \in \mathbb{N}\,\forall n \in \mathbb{N}\,(N \leq n \to \lvert a_{n} - s \rvert < \varepsilon)$
\item $\varphi_{191} :\equiv (a_{1_{\omega}} \in \mathrm{ran}(a) \to \mathrm{ran}(a) \neq \emptyset) \to (\mathrm{MonoInc}(a) \to \varphi_{35}) \to \varphi_{60}$
\item $\varphi_{192} :\equiv \varphi_{101} \to \mathrm{MonoInc}(a) \to M \in \mathbb{R} \to \forall n \in \mathbb{N}\,(a_{n} \leq M) \to \varphi_{14}$
\item $\varphi_{193} :\equiv (\mathrm{MonoInc}(a) \to M \in \mathbb{R} \to \varphi_{64}) \to \varphi_{70} \to \neg (n \in \mathbb{N} \land y = a_{n})$
\item $\varphi_{194} :\equiv \forall n\,(\varphi_{70} \to \neg (n \in \mathbb{N} \land y = a_{n})) \to \forall n\,(\varphi_{70}) \to \forall n\,(\neg (n \in \mathbb{N} \land y = a_{n}))$
\item $\varphi_{195} :\equiv (\mathrm{MonoInc}(a) \to M \in \mathbb{R} \to \varphi_{79}) \to \mathrm{MonoInc}(a) \to M \in \mathbb{R} \to \varphi_{83}$
\item $\varphi_{196} :\equiv \forall n \in \mathbb{N}\,(a_{n} \leq M) \to \mathrm{IsSup}(s,\mathrm{ran}(a)) \to \varepsilon \in \mathbb{R} \to \varphi_{153}$
\item $\varphi_{197} :\equiv \forall n \in \mathbb{N}\,(a_{n} \leq M) \to \mathrm{IsSup}(s,\mathrm{ran}(a)) \to \varepsilon \in \mathbb{R} \to \varphi_{154}$
\item $\varphi_{198} :\equiv \forall n \in \mathbb{N}\,(a_{n} \leq M) \to \mathrm{IsSup}(s,\mathrm{ran}(a)) \to \varepsilon \in \mathbb{R} \to \varphi_{157}$
\item $\varphi_{199} :\equiv (\mathrm{MonoInc}(a) \to M \in \mathbb{R} \to \forall n \in \mathbb{N}\,(a_{n} \leq M) \to \varphi_{32}) \to \mathrm{MonoInc}(a) \to M \in \mathbb{R} \to \varphi_{65}$
\item $\varphi_{200} :\equiv \varphi_{56} \land \varepsilon \in \mathbb{R} \land 0 < \varepsilon \land (x \in \mathrm{ran}(a) \land s - \varepsilon < x) \land (i_{0} \in \mathbb{N} \land x = a_{i_{0}}) \land \neg \exists N \in \mathbb{N}\,\forall n \in \mathbb{N}\,(N \leq n \to \lvert a_{n} - s \rvert < \varepsilon)$
\item $\varphi_{201} :\equiv (\mathrm{MonoInc}(a) \to M \in \mathbb{R} \to \varphi_{111}) \to \mathrm{MonoInc}(a) \to M \in \mathbb{R} \to \varphi_{79}$
\item $\varphi_{202} :\equiv (\mathrm{MonoInc}(a) \to M \in \mathbb{R} \to \varphi_{111}) \to \mathrm{MonoInc}(a) \to M \in \mathbb{R} \to \varphi_{84}$
\item $\varphi_{203} :\equiv (\forall y\,(y \in \mathrm{ran}(a) \leftrightarrow \exists j \in \mathbb{N}\,(y = a_{j})) \to \varphi_{26}) \to \mathrm{MonoInc}(a) \to M \in \mathbb{R} \to \forall n \in \mathbb{N}\,(a_{n} \leq M) \to \varphi_{26}$
\item $\varphi_{204} :\equiv \forall n \in \mathbb{N}\,(a_{n} \leq M) \to \mathrm{IsSup}(s,\mathrm{ran}(a)) \to \varepsilon \in \mathbb{R} \to \varphi_{163}$
\item $\varphi_{205} :\equiv \forall n \in \mathbb{N}\,(a_{n} \leq M) \to \mathrm{IsSup}(s,\mathrm{ran}(a)) \to \varepsilon \in \mathbb{R} \to 0 < \varepsilon \to \varphi_{158}$
\item $\varphi_{206} :\equiv \forall n \in \mathbb{N}\,(a_{n} \leq M) \to \mathrm{IsSup}(s,\mathrm{ran}(a)) \to \varepsilon \in \mathbb{R} \to \varphi_{171}$
\item $\varphi_{207} :\equiv (x = a_{i_{0}} \to s - \varepsilon < x \to s - \varepsilon < a_{i_{0}}) \to \mathrm{MonoInc}(a) \to M \in \mathbb{R} \to \varphi_{155}$
\item $\varphi_{208} :\equiv (\mathrm{MonoInc}(a) \to M \in \mathbb{R} \to \forall n \in \mathbb{N}\,(a_{n} \leq M) \to \varphi_{32}) \to \mathrm{MonoInc}(a) \to M \in \mathbb{R} \to \forall n \in \mathbb{N}\,(a_{n} \leq M) \to \varphi_{23}$
\item $\varphi_{209} :\equiv (\mathrm{MonoInc}(a) \to M \in \mathbb{R} \to \varphi_{84}) \to (y = a_{n} \to a_{n} = y) \to \mathrm{MonoInc}(a) \to M \in \mathbb{R} \to \varphi_{85}$
\item $\varphi_{210} :\equiv \forall n \in \mathbb{N}\,(a_{n} \leq M) \to \mathrm{IsSup}(s,\mathrm{ran}(a)) \to \varepsilon \in \mathbb{R} \to 0 < \varepsilon \to \varphi_{169}$
\item $\varphi_{211} :\equiv (\mathrm{MonoInc}(a) \to M \in \mathbb{R} \to \varphi_{89}) \to \mathrm{MonoInc}(a) \to M \in \mathbb{R} \to \forall n \in \mathbb{N}\,(a_{n} \leq M) \to \varphi_{75}$
\item $\varphi_{212} :\equiv (\mathrm{MonoInc}(a) \to \varphi_{156}) \to \mathrm{MonoInc}(a) \to M \in \mathbb{R} \to \varphi_{13}$
\item $\varphi_{213} :\equiv (\mathrm{MonoInc}(a) \to M \in \mathbb{R} \to \varphi_{44}) \to \varphi_{80} \to \mathrm{MonoInc}(a) \to \varphi_{78}$
\item $\varphi_{214} :\equiv (\mathrm{MonoInc}(a) \to M \in \mathbb{R} \to \varphi_{83}) \to (a_{n} = y \to a_{n} < M \to y < M) \to \varphi_{142}$
\item $\varphi_{215} :\equiv (\mathrm{MonoInc}(a) \to M \in \mathbb{R} \to \forall n \in \mathbb{N}\,(a_{n} \leq M) \to \varphi_{2}) \to \varphi_{9} \to \varphi_{7} \to \mathrm{MonoInc}(a) \to M \in \mathbb{R} \to \varphi_{16}$
\item $\varphi_{216} :\equiv (\mathrm{MonoInc}(a) \to M \in \mathbb{R} \to \forall n \in \mathbb{N}\,(a_{n} \leq M) \to \varphi_{26}) \to (\mathrm{MonoInc}(a) \to \varphi_{78}) \to \mathrm{MonoInc}(a) \to \varphi_{35}$
\item $\varphi_{217} :\equiv (\mathrm{MonoInc}(a) \to M \in \mathbb{R} \to \forall n \in \mathbb{N}\,(a_{n} \leq M) \to \varphi_{14}) \to \varphi_{183}$
\item $\varphi_{218} :\equiv (\mathrm{MonoInc}(a) \to M \in \mathbb{R} \to \forall n \in \mathbb{N}\,(a_{n} \leq M) \to \varphi_{32}) \to \mathrm{MonoInc}(a) \to M \in \mathbb{R} \to \varphi_{151}$
\item $\varphi_{219} :\equiv (\mathrm{MonoInc}(a) \to M \in \mathbb{R} \to \forall n \in \mathbb{N}\,(a_{n} \leq M) \to \varphi_{2}) \to \varphi_{9} \to \varphi_{7} \to \mathrm{MonoInc}(a) \to M \in \mathbb{R} \to \varphi_{45}$
\item $\varphi_{220} :\equiv (\mathrm{MonoInc}(a) \to M \in \mathbb{R} \to \varphi_{143}) \to \mathrm{MonoInc}(a) \to \varphi_{141}$
\item $\varphi_{221} :\equiv (\mathrm{MonoInc}(a) \to M \in \mathbb{R} \to \forall n \in \mathbb{N}\,(a_{n} \leq M) \to \varphi_{23}) \to \varphi_{192}$
\item $\varphi_{222} :\equiv (\mathrm{MonoInc}(a) \to M \in \mathbb{R} \to M \in \mathbb{R}) \to (\mathrm{MonoInc}(a) \to \varphi_{33}) \to \varphi_{40} \to \varphi_{39} \to \varphi_{43}$
\item $\varphi_{223} :\equiv (\mathrm{MonoInc}(a) \to M \in \mathbb{R} \to \varphi_{143}) \to \mathrm{MonoInc}(a) \to M \in \mathbb{R} \to \varphi_{148}$
\item $\varphi_{224} :\equiv (\mathrm{MonoInc}(a) \to M \in \mathbb{R} \to \varphi_{89}) \to \mathrm{MonoInc}(a) \to M \in \mathbb{R} \to \forall n \in \mathbb{N}\,(a_{n} \leq M) \to \varphi_{173}$
\item $\varphi_{225} :\equiv (\mathrm{MonoInc}(a) \to M \in \mathbb{R} \to \varphi_{159}) \to \mathrm{MonoInc}(a) \to M \in \mathbb{R} \to \varphi_{145}$
\item $\varphi_{226} :\equiv (\mathrm{MonoInc}(a) \to M \in \mathbb{R} \to \varphi_{159}) \to \mathrm{MonoInc}(a) \to M \in \mathbb{R} \to \varphi_{146}$
\item $\varphi_{227} :\equiv (\mathrm{MonoInc}(a) \to M \in \mathbb{R} \to \varphi_{155}) \to \mathrm{MonoInc}(a) \to M \in \mathbb{R} \to \varphi_{167}$
\item $\varphi_{228} :\equiv (\mathrm{MonoInc}(a) \to M \in \mathbb{R} \to \forall n \in \mathbb{N}\,(a_{n} \leq M) \to \varphi_{32}) \to \mathrm{MonoInc}(a) \to M \in \mathbb{R} \to \varphi_{205}$
\item $\varphi_{229} :\equiv (\varphi_{31} \to \neg (N \in \mathbb{N} \land \forall n \in \mathbb{N}\,(N \leq n \to \lvert a_{n} - s \rvert < \varepsilon))) \to \mathrm{MonoInc}(a) \to M \in \mathbb{R} \to \varphi_{206}$
\item $\varphi_{230} :\equiv (\mathrm{MonoInc}(a) \to M \in \mathbb{R} \to \varphi_{168}) \to \varphi_{162} \to \neg (x \in \mathrm{ran}(a) \land s - \varepsilon < x)$
\item $\varphi_{231} :\equiv \forall x\,(\varphi_{162} \to \neg (x \in \mathrm{ran}(a) \land s - \varepsilon < x)) \to \forall x\,(\varphi_{162}) \to \forall x\,(\neg (x \in \mathrm{ran}(a) \land s - \varepsilon < x))$
\item $\varphi_{232} :\equiv \forall y\,(\varphi_{67}) \to (\forall y\,(\varphi_{67}) \to \varphi_{10} \to \forall y \in \mathrm{ran}(a)\,(y \leq M)) \to \varphi_{186}$
\item $\varphi_{233} :\equiv (\mathrm{MonoInc}(a) \to M \in \mathbb{R} \to \forall n \in \mathbb{N}\,(a_{n} \leq M) \to \varphi_{174}) \to \mathrm{MonoInc}(a) \to M \in \mathbb{R} \to \varphi_{178}$
\item $\varphi_{234} :\equiv (\mathrm{MonoInc}(a) \to M \in \mathbb{R} \to \forall n \in \mathbb{N}\,(a_{n} \leq M) \to \varphi_{2}) \to (\mathrm{MonoInc}(a) \to M \in \mathbb{R} \to \varphi_{69}) \to \varphi_{211}$
\item $\varphi_{235} :\equiv (\mathrm{MonoInc}(a) \to M \in \mathbb{R} \to \varphi_{151}) \to \varphi_{100} \to \mathrm{MonoInc}(a) \to M \in \mathbb{R} \to \varphi_{149}$
\item $\varphi_{236} :\equiv (\mathrm{MonoInc}(a) \to M \in \mathbb{R} \to \forall n \in \mathbb{N}\,(a_{n} \leq M) \to \varphi_{165}) \to \varphi_{190} \to \neg (i_{0} \in \mathbb{N} \land x = a_{i_{0}})$
\item $\varphi_{237} :\equiv (\mathrm{MonoInc}(a) \to M \in \mathbb{R} \to \forall n \in \mathbb{N}\,(a_{n} \leq M) \to \varphi_{161}) \to \mathrm{MonoInc}(a) \to M \in \mathbb{R} \to \forall n \in \mathbb{N}\,(a_{n} \leq M) \to \varphi_{175}$
\item $\varphi_{238} :\equiv (\mathrm{MonoInc}(a) \to M \in \mathbb{R} \to \forall n \in \mathbb{N}\,(a_{n} \leq M) \to \varphi_{160}) \to \mathrm{MonoInc}(a) \to M \in \mathbb{R} \to \forall n \in \mathbb{N}\,(a_{n} \leq M) \to \varphi_{176}$
\item $\varphi_{239} :\equiv \forall j\,(\varphi_{11}) \to (\forall j\,(\varphi_{11}) \to \varphi_{27}) \to \varphi_{36} \to \mathrm{MonoInc}(a) \to M \in \mathbb{R} \to \varphi_{168}$
\item $\varphi_{240} :\equiv (\mathrm{MonoInc}(a) \to M \in \mathbb{R} \to \forall n \in \mathbb{N}\,(a_{n} \leq M) \to \varphi_{174}) \to \mathrm{MonoInc}(a) \to M \in \mathbb{R} \to \forall n \in \mathbb{N}\,(a_{n} \leq M) \to \varphi_{166}$
\item $\varphi_{241} :\equiv (\mathrm{MonoInc}(a) \to M \in \mathbb{R} \to \forall n \in \mathbb{N}\,(a_{n} \leq M) \to \varphi_{177}) \to \mathrm{MonoInc}(a) \to M \in \mathbb{R} \to \forall n \in \mathbb{N}\,(a_{n} \leq M) \to \varphi_{170}$
\item $\varphi_{242} :\equiv \forall i_{0}\,(\varphi_{190} \to \neg (i_{0} \in \mathbb{N} \land x = a_{i_{0}})) \to \forall i_{0}\,(\varphi_{190}) \to \forall i_{0}\,(\neg (i_{0} \in \mathbb{N} \land x = a_{i_{0}}))$
\item $\varphi_{243} :\equiv (\mathrm{MonoInc}(a) \to M \in \mathbb{R} \to \forall n \in \mathbb{N}\,(a_{n} \leq M) \to \varphi_{173}) \to \mathrm{MonoInc}(a) \to M \in \mathbb{R} \to \forall n \in \mathbb{N}\,(a_{n} \leq M) \to \varphi_{179}$
\item $\varphi_{244} :\equiv (\mathrm{MonoInc}(a) \to M \in \mathbb{R} \to \forall n \in \mathbb{N}\,(a_{n} \leq M) \to \varphi_{174}) \to \mathrm{MonoInc}(a) \to M \in \mathbb{R} \to \forall n \in \mathbb{N}\,(a_{n} \leq M) \to \varphi_{180}$
\item $\varphi_{245} :\equiv (\mathrm{MonoInc}(a) \to M \in \mathbb{R} \to \forall n \in \mathbb{N}\,(a_{n} \leq M) \to \varphi_{175}) \to \mathrm{MonoInc}(a) \to M \in \mathbb{R} \to \forall n \in \mathbb{N}\,(a_{n} \leq M) \to \varphi_{181}$
\item $\varphi_{246} :\equiv (\mathrm{MonoInc}(a) \to M \in \mathbb{R} \to \varphi_{167}) \to \varphi_{73} \to \mathrm{MonoInc}(a) \to M \in \mathbb{R} \to \varphi_{172}$
\item $\varphi_{247} :\equiv (\mathrm{MonoInc}(a) \to M \in \mathbb{R} \to \forall n \in \mathbb{N}\,(a_{n} \leq M) \to \varphi_{179}) \to \mathrm{MonoInc}(a) \to M \in \mathbb{R} \to \forall n \in \mathbb{N}\,(a_{n} \leq M) \to \varphi_{185}$
\item $\varphi_{248} :\equiv (\mathrm{MonoInc}(a) \to M \in \mathbb{R} \to \varphi_{149}) \to (\mathrm{MonoInc}(a) \to \varphi_{141}) \to \mathrm{MonoInc}(a) \to \varphi_{147}$
\item $\varphi_{249} :\equiv (\mathrm{MonoInc}(a) \to M \in \mathbb{R} \to \forall n \in \mathbb{N}\,(a_{n} \leq M) \to \varphi_{185}) \to \varphi_{188} \to n \in \mathbb{N} \to N \leq n \to \lvert a_{n} - s \rvert < \varepsilon$
\item $\varphi_{250} :\equiv (\mathrm{MonoInc}(a) \to M \in \mathbb{R} \to \varphi_{198}) \to \mathrm{MonoInc}(a) \to M \in \mathbb{R} \to \forall n \in \mathbb{N}\,(a_{n} \leq M) \to \varphi_{177}$
\item $\varphi_{251} :\equiv (\mathrm{MonoInc}(a) \to M \in \mathbb{R} \to \forall n \in \mathbb{N}\,(a_{n} \leq M) \to \varphi_{160}) \to a_{n} = a_{n} \to \mathrm{MonoInc}(a) \to M \in \mathbb{R} \to \varphi_{197}$
\item $\varphi_{252} :\equiv \forall n\,(\varphi_{188} \to n \in \mathbb{N} \to N \leq n \to \lvert a_{n} - s \rvert < \varepsilon) \to \forall n\,(\varphi_{188}) \to \forall n \in \mathbb{N}\,(N \leq n \to \lvert a_{n} - s \rvert < \varepsilon)$
\item $\varphi_{253} :\equiv (\mathrm{MonoInc}(a) \to M \in \mathbb{R} \to \varphi_{196}) \to \mathrm{MonoInc}(a) \to M \in \mathbb{R} \to \varphi_{204}$
\item $\varphi_{254} :\equiv \forall s\,(\varphi_{55} \to \neg \mathrm{IsSup}(s,\mathrm{ran}(a))) \to \varphi_{187} \to (\varphi_{55} \to \forall s\,(\varphi_{55})) \to \varphi_{49}$
\item $\varphi_{255} :\equiv (\mathrm{MonoInc}(a) \to M \in \mathbb{R} \to \varphi_{206}) \to \varphi_{200} \to \neg (N \in \mathbb{N} \land s - \varepsilon < a_{N})$
\item $\varphi_{256} :\equiv (\mathrm{MonoInc}(a) \to M \in \mathbb{R} \to \varphi_{16}) \to (\mathrm{MonoInc}(a) \to M \in \mathbb{R} \to \varphi_{69}) \to \varphi_{224}$
\item $\varphi_{257} :\equiv \forall N\,(\varphi_{200} \to \neg (N \in \mathbb{N} \land s - \varepsilon < a_{N})) \to \forall N\,(\varphi_{200}) \to \forall N\,(\neg (N \in \mathbb{N} \land s - \varepsilon < a_{N}))$
\item $\varphi_{258} :\equiv (\mathrm{Conv}(a,s) \leftrightarrow \mathrm{Seq}(a) \land s \in \mathbb{R} \land \varphi_{115}) \to \varphi_{24} \to (\mathrm{MonoInc}(a) \to M \in \mathbb{R} \to \varphi_{16}) \to \varphi_{212}$
\item $\varphi_{259} :\equiv (\mathrm{MonoInc}(a) \to M \in \mathbb{R} \to \varphi_{197}) \to \varphi_{41} \to \mathrm{MonoInc}(a) \to M \in \mathbb{R} \to \varphi_{198}$
\item $\varphi_{260} :\equiv (\mathrm{MonoInc}(a) \to M \in \mathbb{R} \to \varphi_{146}) \to (\mathrm{MonoInc}(a) \to M \in \mathbb{R} \to \varphi_{148}) \to \varphi_{207}$
\item $\varphi_{261} :\equiv (\mathrm{MonoInc}(a) \to M \in \mathbb{R} \to \varphi_{205}) \to \varphi_{128} \to \mathrm{MonoInc}(a) \to M \in \mathbb{R} \to \varphi_{210}$
\item $\varphi_{262} :\equiv (\forall \varepsilon\,(\varphi_{56} \to \varphi_{102}) \to \forall \varepsilon\,(\varphi_{56}) \to \varphi_{115}) \to (\varphi_{56} \to \forall \varepsilon\,(\varphi_{56})) \to \mathrm{MonoInc}(a) \to \varphi_{156}$
\item $\varphi_{263} :\equiv (\mathrm{MonoInc}(a) \to M \in \mathbb{R} \to \forall n \in \mathbb{N}\,(a_{n} \leq M) \to \varphi_{160}) \to \varphi_{237}$
\item $\varphi_{264} :\equiv (\varphi_{200} \to \forall N\,(\varphi_{200})) \to \mathrm{MonoInc}(a) \to M \in \mathbb{R} \to \forall n \in \mathbb{N}\,(a_{n} \leq M) \to \varphi_{165}$
\item $\varphi_{265} :\equiv (\mathrm{MonoInc}(a) \to M \in \mathbb{R} \to \forall n \in \mathbb{N}\,(a_{n} \leq M) \to \varphi_{170}) \to \varphi_{243}$
\item $\varphi_{266} :\equiv (\mathrm{MonoInc}(a) \to M \in \mathbb{R} \to \forall n \in \mathbb{N}\,(a_{n} \leq M) \to \varphi_{180}) \to \varphi_{245}$
\item $\varphi_{267} :\equiv (\mathrm{MonoInc}(a) \to M \in \mathbb{R} \to \forall n \in \mathbb{N}\,(a_{n} \leq M) \to \varphi_{181}) \to \varphi_{247}$
\item $\varphi_{268} :\equiv \forall n\,(\varphi_{70} \to \neg (n \in \mathbb{N} \land y = a_{n})) \to \varphi_{194} \to (\varphi_{70} \to \forall n\,(\varphi_{70})) \to \varphi_{12} \to (\varphi_{12} \to \varphi_{22}) \to \varphi_{21} \to \varphi_{66}$
\item $\varphi_{269} :\equiv (\mathrm{MonoInc}(a) \to M \in \mathbb{R} \to \forall n \in \mathbb{N}\,(a_{n} \leq M) \to \varphi_{166}) \to \varphi_{263}$
\item $\varphi_{270} :\equiv (\mathrm{MonoInc}(a) \to M \in \mathbb{R} \to \forall n \in \mathbb{N}\,(a_{n} \leq M) \to \varphi_{176}) \to \varphi_{267}$
\item $\varphi_{271} :\equiv (\mathrm{MonoInc}(a) \to N \in \mathbb{N} \to n \in \mathbb{N} \to N \leq n \to a_{N} \leq a_{n}) \to (\mathrm{MonoInc}(a) \to \mathrm{MonoInc}(a)) \to \varphi_{269}$
\item $\varphi_{272} :\equiv \forall x\,(\varphi_{162} \to \neg (x \in \mathrm{ran}(a) \land s - \varepsilon < x)) \to \varphi_{231} \to (\varphi_{162} \to \forall x\,(\varphi_{162})) \to \mathrm{MonoInc}(a) \to \varphi_{152}$
\item $\varphi_{273} :\equiv (\mathrm{MonoInc}(a) \to M \in \mathbb{R} \to \varphi_{16}) \to (\mathrm{MonoInc}(a) \to M \in \mathbb{R} \to \varphi_{69}) \to \varphi_{270}$
\item $\varphi_{274} :\equiv (\mathrm{MonoInc}(a) \to M \in \mathbb{R} \to \forall n \in \mathbb{N}\,(a_{n} \leq M) \to \varphi_{75}) \to \varphi_{272}$
\item $\varphi_{275} :\equiv \forall n\,(\varphi_{188} \to n \in \mathbb{N} \to N \leq n \to \lvert a_{n} - s \rvert < \varepsilon) \to \varphi_{252} \to (\varphi_{188} \to \forall n\,(\varphi_{188})) \to \mathrm{MonoInc}(a) \to M \in \mathbb{R} \to \varphi_{196}$
\item $\varphi_{276} :\equiv \forall i_{0}\,(\varphi_{190} \to \neg (i_{0} \in \mathbb{N} \land x = a_{i_{0}})) \to \varphi_{242} \to (\varphi_{190} \to \forall i_{0}\,(\varphi_{190})) \to \varphi_{239}$
\item $\varphi_{277} :\equiv (\mathrm{MonoInc}(a) \to M \in \mathbb{R} \to \varphi_{172}) \to \forall N\,(\varphi_{200} \to \neg (N \in \mathbb{N} \land s - \varepsilon < a_{N})) \to \varphi_{257} \to \varphi_{264}$
\end{abbrevlist}
\end{fullproof}

% 自動生成 (prover/build.js): MonoConvDec  ノード数 438
\begin{fullproof}{定理5.3(2)}{$\forall a\,\forall M\,(\mathrm{MonoDec}(a) \to M \in \mathbb{R} \to \forall n \in \mathbb{N}\,(M \leq a_{n}) \to \exists \alpha\,(\mathrm{Conv}(a,\alpha)))$}{438}\phantomsection\label{tree:MonoConvDec}
\noindent$\Pi_{1}$:
\begin{ptree}
\AXC{$\mathrm{MonoDec}(a) \to M \in \mathbb{R} \to \forall n \in \mathbb{N}\,(M \leq a_{n}) \to \varphi_{174}$\lbl{A1}}
\AXC{$\varphi_{244}$\lbl{A1}}
\BIC{$\mathrm{MonoDec}(a) \to M \in \mathbb{R} \to \forall n \in \mathbb{N}\,(M \leq a_{n}) \to \varphi_{179}$}
\end{ptree}
\noindent$\Pi_{2}$:
\begin{ptree}
\AXC{$\forall a\,(\varphi_{41})$\lbl{[C21]}}
\AXC{$\forall a\,(\varphi_{41}) \to \varphi_{41}$\lbl{A2}}
\BIC{$\varphi_{41}$}
\AXC{$\varphi_{41} \to \varphi_{34}$\lbl{A2}}
\BIC{$\varphi_{34}$}
\AXC{$\varphi_{140}$\lbl{A2}}
\BIC{$\mathrm{MonoDec}(a) \to N \in \mathbb{N} \to n \in \mathbb{N} \to N \leq n \to a_{n} \leq a_{N}$}
\AXC{$\varphi_{271}$\lbl{A1}}
\BIC{$(\mathrm{MonoDec}(a) \to \mathrm{MonoDec}(a)) \to \varphi_{269}$}
\end{ptree}
\noindent$\Pi_{3}$:
\begin{ptree}
\AXC{$\mathrm{MonoDec}(a) \to \mathrm{MonoDec}(a)$\lbl{A1}}
\AXC{$\stackrel{\Pi_{2}}{\vdots}$}\noLine
\UIC{$(\mathrm{MonoDec}(a) \to \mathrm{MonoDec}(a)) \to \varphi_{269}$}
\BIC{$\varphi_{269}$}
\end{ptree}
\noindent$\Pi_{4}$:
\begin{ptree}
\AXC{$\mathrm{MonoDec}(a) \to M \in \mathbb{R} \to \forall n \in \mathbb{N}\,(M \leq a_{n}) \to \varphi_{174}$\lbl{A1}}
\AXC{$\varphi_{240}$\lbl{A1}}
\BIC{$\mathrm{MonoDec}(a) \to M \in \mathbb{R} \to \forall n \in \mathbb{N}\,(M \leq a_{n}) \to \varphi_{166}$}
\AXC{$\stackrel{\Pi_{3}}{\vdots}$}\noLine
\UIC{$\varphi_{269}$}
\BIC{$\varphi_{263}$}
\end{ptree}
\noindent$\Pi_{5}$:
\begin{ptree}
\AXC{$\mathrm{MonoDec}(a) \to M \in \mathbb{R} \to \forall n \in \mathbb{N}\,(M \leq a_{n}) \to \varphi_{160}$\lbl{A1}}
\AXC{$\mathrm{MonoDec}(a) \to M \in \mathbb{R} \to \forall n \in \mathbb{N}\,(M \leq a_{n}) \to \varphi_{161}$\lbl{A1}}
\AXC{$\stackrel{\Pi_{4}}{\vdots}$}\noLine
\UIC{$\varphi_{263}$}
\BIC{$\varphi_{237}$}
\BIC{$\mathrm{MonoDec}(a) \to M \in \mathbb{R} \to \forall n \in \mathbb{N}\,(M \leq a_{n}) \to \varphi_{175}$}
\end{ptree}
\noindent$\Pi_{6}$:
\begin{ptree}
\AXC{$\forall x\,\forall y\,\forall z\,(x \leq y \to y < z \to x < z)$\lbl{[O4]}}
\AXC{$\varphi_{71}$\lbl{A2}}
\BIC{$\forall y\,\forall z\,(a_{n} \leq y \to y < z \to a_{n} < z)$}
\AXC{$\varphi_{123}$\lbl{A2}}
\BIC{$\forall z\,(a_{n} \leq a_{N} \to a_{N} < z \to a_{n} < z)$}
\AXC{$\varphi_{138}$\lbl{A2}}
\BIC{$a_{n} \leq a_{N} \to a_{N} < s + \varepsilon \to a_{n} < s + \varepsilon$}
\end{ptree}
\noindent$\Pi_{7}$:
\begin{ptree}
\AXC{$\stackrel{\Pi_{6}}{\vdots}$}\noLine
\UIC{$a_{n} \leq a_{N} \to a_{N} < s + \varepsilon \to a_{n} < s + \varepsilon$}
\AXC{$(a_{n} \leq a_{N} \to a_{N} < s + \varepsilon \to a_{n} < s + \varepsilon) \to \varphi_{266}$\lbl{A1}}
\BIC{$\varphi_{266}$}
\end{ptree}
\noindent$\Pi_{8}$:
\begin{ptree}
\AXC{$\stackrel{\Pi_{1}}{\vdots}$}\noLine
\UIC{$\mathrm{MonoDec}(a) \to M \in \mathbb{R} \to \forall n \in \mathbb{N}\,(M \leq a_{n}) \to \varphi_{179}$}
\AXC{$\stackrel{\Pi_{5}}{\vdots}$}\noLine
\UIC{$\mathrm{MonoDec}(a) \to M \in \mathbb{R} \to \forall n \in \mathbb{N}\,(M \leq a_{n}) \to \varphi_{175}$}
\AXC{$\stackrel{\Pi_{7}}{\vdots}$}\noLine
\UIC{$\varphi_{266}$}
\BIC{$\varphi_{245}$}
\BIC{$\mathrm{MonoDec}(a) \to M \in \mathbb{R} \to \forall n \in \mathbb{N}\,(M \leq a_{n}) \to \varphi_{180}$}
\end{ptree}
\noindent$\Pi_{9}$:
\begin{ptree}
\AXC{$x = x$\lbl{A5}}
\UIC{$\forall x\,(x = x)$}
\AXC{$\forall x\,(x = x) \to a_{n} = a_{n}$\lbl{A2}}
\BIC{$a_{n} = a_{n}$}
\AXC{$\mathrm{MonoDec}(a) \to M \in \mathbb{R} \to \forall n \in \mathbb{N}\,(M \leq a_{n}) \to \varphi_{161}$\lbl{A1}}
\AXC{$\varphi_{251}$\lbl{A1}}
\BIC{$a_{n} = a_{n} \to \mathrm{MonoDec}(a) \to M \in \mathbb{R} \to \varphi_{197}$}
\BIC{$\mathrm{MonoDec}(a) \to M \in \mathbb{R} \to \varphi_{197}$}
\end{ptree}
\noindent$\Pi_{10}$:
\begin{ptree}
\AXC{$\mathrm{MonoDec}(a) \to \mathrm{MonoDec}(a)$\lbl{A1}}
\AXC{$(\mathrm{MonoDec}(a) \to \mathrm{MonoDec}(a)) \to \varphi_{18} \to \varphi_{24}$\lbl{A1}}
\BIC{$\varphi_{18} \to \varphi_{24}$}
\end{ptree}
\noindent$\Pi_{11}$:
\begin{ptree}
\AXC{$\forall a\,(\varphi_{18})$\lbl{DefMonoDec}}
\AXC{$\forall a\,(\varphi_{18}) \to \varphi_{18}$\lbl{A2}}
\BIC{$\varphi_{18}$}
\AXC{$\stackrel{\Pi_{10}}{\vdots}$}\noLine
\UIC{$\varphi_{18} \to \varphi_{24}$}
\BIC{$\varphi_{24}$}
\end{ptree}
\noindent$\Pi_{12}$:
\begin{ptree}
\AXC{$\stackrel{\Pi_{11}}{\vdots}$}\noLine
\UIC{$\varphi_{24}$}
\AXC{$\forall a\,(\mathrm{Seq}(a) \to \varphi_{32})$\lbl{[S16]}}
\AXC{$\forall a\,(\mathrm{Seq}(a) \to \varphi_{32}) \to \mathrm{Seq}(a) \to \varphi_{32}$\lbl{A2}}
\BIC{$\mathrm{Seq}(a) \to \varphi_{32}$}
\AXC{$(\mathrm{Seq}(a) \to \varphi_{32}) \to \varphi_{184}$\lbl{A1}}
\BIC{$\varphi_{184}$}
\BIC{$\mathrm{MonoDec}(a) \to M \in \mathbb{R} \to \forall n \in \mathbb{N}\,(M \leq a_{n}) \to \varphi_{32}$}
\end{ptree}
\noindent$\Pi_{13}$:
\begin{ptree}
\AXC{$\varphi_{128}$\lbl{A2}}
\AXC{$\stackrel{\Pi_{12}}{\vdots}$}\noLine
\UIC{$\mathrm{MonoDec}(a) \to M \in \mathbb{R} \to \forall n \in \mathbb{N}\,(M \leq a_{n}) \to \varphi_{32}$}
\AXC{$\varphi_{228}$\lbl{A1}}
\BIC{$\mathrm{MonoDec}(a) \to M \in \mathbb{R} \to \varphi_{205}$}
\AXC{$\varphi_{261}$\lbl{A1}}
\BIC{$\varphi_{128} \to \mathrm{MonoDec}(a) \to M \in \mathbb{R} \to \varphi_{210}$}
\BIC{$\mathrm{MonoDec}(a) \to M \in \mathbb{R} \to \varphi_{210}$}
\end{ptree}
\noindent$\Pi_{14}$:
\begin{ptree}
\AXC{$\stackrel{\Pi_{13}}{\vdots}$}\noLine
\UIC{$\mathrm{MonoDec}(a) \to M \in \mathbb{R} \to \varphi_{210}$}
\AXC{$(\mathrm{MonoDec}(a) \to M \in \mathbb{R} \to \varphi_{210}) \to \varphi_{250}$\lbl{A1}}
\BIC{$\varphi_{250}$}
\end{ptree}
\noindent$\Pi_{15}$:
\begin{ptree}
\AXC{$\forall a\,\forall X\,(\mathrm{LB}(a,X) \leftrightarrow a \in \mathbb{R} \land \forall x \in X\,(a \leq x))$\lbl{DefLB}}
\AXC{$\varphi_{107}$\lbl{A2}}
\BIC{$\forall X\,(\mathrm{LB}(s,X) \leftrightarrow s \in \mathbb{R} \land \forall x \in X\,(s \leq x))$}
\AXC{$\forall X\,(\mathrm{LB}(s,X) \leftrightarrow s \in \mathbb{R} \land \forall x \in X\,(s \leq x)) \to \varphi_{7}$\lbl{A2}}
\BIC{$\varphi_{7}$}
\end{ptree}
\noindent$\Pi_{16}$:
\begin{ptree}
\AXC{$\mathrm{MonoDec}(a) \to M \in \mathbb{R} \to \forall n \in \mathbb{N}\,(M \leq a_{n}) \to \varphi_{2}$\lbl{A1}}
\AXC{$\varphi_{219}$\lbl{A1}}
\BIC{$\varphi_{9} \to \varphi_{7} \to \mathrm{MonoDec}(a) \to M \in \mathbb{R} \to \varphi_{45}$}
\end{ptree}
\noindent$\Pi_{17}$:
\begin{ptree}
\AXC{$\stackrel{\Pi_{15}}{\vdots}$}\noLine
\UIC{$\varphi_{7}$}
\AXC{$\forall s\,(\varphi_{4})$\lbl{DefInf}}
\AXC{$\forall s\,(\varphi_{4}) \to \varphi_{4}$\lbl{A2}}
\BIC{$\varphi_{4}$}
\AXC{$\varphi_{4} \to \varphi_{9}$\lbl{A2}}
\BIC{$\varphi_{9}$}
\AXC{$\stackrel{\Pi_{16}}{\vdots}$}\noLine
\UIC{$\varphi_{9} \to \varphi_{7} \to \mathrm{MonoDec}(a) \to M \in \mathbb{R} \to \varphi_{45}$}
\BIC{$\varphi_{7} \to \mathrm{MonoDec}(a) \to M \in \mathbb{R} \to \varphi_{45}$}
\BIC{$\mathrm{MonoDec}(a) \to M \in \mathbb{R} \to \varphi_{45}$}
\end{ptree}
\noindent$\Pi_{18}$:
\begin{ptree}
\AXC{$\stackrel{\Pi_{17}}{\vdots}$}\noLine
\UIC{$\mathrm{MonoDec}(a) \to M \in \mathbb{R} \to \varphi_{45}$}
\AXC{$(\mathrm{MonoDec}(a) \to M \in \mathbb{R} \to \varphi_{45}) \to \varphi_{164}$\lbl{A1}}
\BIC{$\varphi_{164}$}
\end{ptree}
\noindent$\Pi_{19}$:
\begin{ptree}
\AXC{$\forall x \in \mathrm{ran}(a)\,(s \leq x) \to a_{n} \in \mathrm{ran}(a) \to s \leq a_{n}$\lbl{A2}}
\AXC{$\stackrel{\Pi_{18}}{\vdots}$}\noLine
\UIC{$\varphi_{164}$}
\BIC{$\mathrm{MonoDec}(a) \to M \in \mathbb{R} \to \varphi_{82}$}
\AXC{$(\mathrm{MonoDec}(a) \to M \in \mathbb{R} \to \varphi_{82}) \to \varphi_{241}$\lbl{A1}}
\BIC{$\varphi_{241}$}
\end{ptree}
\noindent$\Pi_{20}$:
\begin{ptree}
\AXC{$\varphi_{42}$\lbl{A2}}
\AXC{$\stackrel{\Pi_{9}}{\vdots}$}\noLine
\UIC{$\mathrm{MonoDec}(a) \to M \in \mathbb{R} \to \varphi_{197}$}
\AXC{$\varphi_{259}$\lbl{A1}}
\BIC{$\varphi_{42} \to \mathrm{MonoDec}(a) \to M \in \mathbb{R} \to \varphi_{198}$}
\BIC{$\mathrm{MonoDec}(a) \to M \in \mathbb{R} \to \varphi_{198}$}
\AXC{$\stackrel{\Pi_{14}}{\vdots}$}\noLine
\UIC{$\varphi_{250}$}
\BIC{$\mathrm{MonoDec}(a) \to M \in \mathbb{R} \to \forall n \in \mathbb{N}\,(M \leq a_{n}) \to \varphi_{177}$}
\AXC{$\stackrel{\Pi_{19}}{\vdots}$}\noLine
\UIC{$\varphi_{241}$}
\BIC{$\mathrm{MonoDec}(a) \to M \in \mathbb{R} \to \forall n \in \mathbb{N}\,(M \leq a_{n}) \to \varphi_{170}$}
\end{ptree}
\noindent$\Pi_{21}$:
\begin{ptree}
\AXC{$\forall a\,\forall X\,(\mathrm{LB}(a,X) \leftrightarrow a \in \mathbb{R} \land \forall x \in X\,(a \leq x))$\lbl{DefLB}}
\AXC{$\varphi_{107}$\lbl{A2}}
\BIC{$\forall X\,(\mathrm{LB}(s,X) \leftrightarrow s \in \mathbb{R} \land \forall x \in X\,(s \leq x))$}
\AXC{$\forall X\,(\mathrm{LB}(s,X) \leftrightarrow s \in \mathbb{R} \land \forall x \in X\,(s \leq x)) \to \varphi_{7}$\lbl{A2}}
\BIC{$\varphi_{7}$}
\end{ptree}
\noindent$\Pi_{22}$:
\begin{ptree}
\AXC{$\mathrm{MonoDec}(a) \to M \in \mathbb{R} \to \forall n \in \mathbb{N}\,(M \leq a_{n}) \to \varphi_{2}$\lbl{A1}}
\AXC{$\varphi_{215}$\lbl{A1}}
\BIC{$\varphi_{9} \to \varphi_{7} \to \mathrm{MonoDec}(a) \to M \in \mathbb{R} \to \varphi_{16}$}
\end{ptree}
\noindent$\Pi_{23}$:
\begin{ptree}
\AXC{$\stackrel{\Pi_{21}}{\vdots}$}\noLine
\UIC{$\varphi_{7}$}
\AXC{$\forall s\,(\varphi_{4})$\lbl{DefInf}}
\AXC{$\forall s\,(\varphi_{4}) \to \varphi_{4}$\lbl{A2}}
\BIC{$\varphi_{4}$}
\AXC{$\varphi_{4} \to \varphi_{9}$\lbl{A2}}
\BIC{$\varphi_{9}$}
\AXC{$\stackrel{\Pi_{22}}{\vdots}$}\noLine
\UIC{$\varphi_{9} \to \varphi_{7} \to \mathrm{MonoDec}(a) \to M \in \mathbb{R} \to \varphi_{16}$}
\BIC{$\varphi_{7} \to \mathrm{MonoDec}(a) \to M \in \mathbb{R} \to \varphi_{16}$}
\BIC{$\mathrm{MonoDec}(a) \to M \in \mathbb{R} \to \varphi_{16}$}
\end{ptree}
\noindent$\Pi_{24}$:
\begin{ptree}
\AXC{$\forall x\,\forall e\,(x \in \mathbb{R} \to e \in \mathbb{R} \to 0 < e \to x - e < x)$\lbl{[E34]}}
\AXC{$\varphi_{133}$\lbl{A2}}
\BIC{$\forall e\,(s \in \mathbb{R} \to e \in \mathbb{R} \to 0 < e \to s - e < s)$}
\AXC{$\varphi_{134}$\lbl{A2}}
\BIC{$s \in \mathbb{R} \to \varepsilon \in \mathbb{R} \to 0 < \varepsilon \to s - \varepsilon < s$}
\AXC{$(s \in \mathbb{R} \to \varepsilon \in \mathbb{R} \to 0 < \varepsilon \to s - \varepsilon < s) \to \varphi_{256}$\lbl{A1}}
\BIC{$\varphi_{256}$}
\end{ptree}
\noindent$\Pi_{25}$:
\begin{ptree}
\AXC{$\mathrm{MonoDec}(a) \to M \in \mathbb{R} \to \varphi_{89}$\lbl{A1}}
\AXC{$\mathrm{MonoDec}(a) \to M \in \mathbb{R} \to \varphi_{69}$\lbl{A1}}
\AXC{$\stackrel{\Pi_{23}}{\vdots}$}\noLine
\UIC{$\mathrm{MonoDec}(a) \to M \in \mathbb{R} \to \varphi_{16}$}
\AXC{$\stackrel{\Pi_{24}}{\vdots}$}\noLine
\UIC{$\varphi_{256}$}
\BIC{$(\mathrm{MonoDec}(a) \to M \in \mathbb{R} \to \varphi_{69}) \to \varphi_{224}$}
\BIC{$\varphi_{224}$}
\BIC{$\mathrm{MonoDec}(a) \to M \in \mathbb{R} \to \forall n \in \mathbb{N}\,(M \leq a_{n}) \to \varphi_{173}$}
\end{ptree}
\noindent$\Pi_{26}$:
\begin{ptree}
\AXC{$\forall x\,\forall y\,\forall z\,(x < y \to y \leq z \to x < z)$\lbl{[O3]}}
\AXC{$\varphi_{72}$\lbl{A2}}
\BIC{$\forall y\,\forall z\,(s - \varepsilon < y \to y \leq z \to s - \varepsilon < z)$}
\AXC{$\varphi_{91}$\lbl{A2}}
\BIC{$\forall z\,(s - \varepsilon < s \to s \leq z \to s - \varepsilon < z)$}
\AXC{$\varphi_{108}$\lbl{A2}}
\BIC{$s - \varepsilon < s \to s \leq a_{n} \to s - \varepsilon < a_{n}$}
\AXC{$(s - \varepsilon < s \to s \leq a_{n} \to s - \varepsilon < a_{n}) \to \varphi_{265}$\lbl{A1}}
\BIC{$\varphi_{265}$}
\end{ptree}
\noindent$\Pi_{27}$:
\begin{ptree}
\AXC{$\stackrel{\Pi_{20}}{\vdots}$}\noLine
\UIC{$\mathrm{MonoDec}(a) \to M \in \mathbb{R} \to \forall n \in \mathbb{N}\,(M \leq a_{n}) \to \varphi_{170}$}
\AXC{$\stackrel{\Pi_{25}}{\vdots}$}\noLine
\UIC{$\mathrm{MonoDec}(a) \to M \in \mathbb{R} \to \forall n \in \mathbb{N}\,(M \leq a_{n}) \to \varphi_{173}$}
\AXC{$\stackrel{\Pi_{26}}{\vdots}$}\noLine
\UIC{$\varphi_{265}$}
\BIC{$\varphi_{242}$}
\BIC{$\mathrm{MonoDec}(a) \to M \in \mathbb{R} \to \forall n \in \mathbb{N}\,(M \leq a_{n}) \to \varphi_{181}$}
\end{ptree}
\noindent$\Pi_{28}$:
\begin{ptree}
\AXC{$\forall a\,\forall n\,(\mathrm{Seq}(a) \to n \in \mathbb{N} \to a_{n} \in \mathbb{R})$\lbl{[S1]}}
\AXC{$\varphi_{103}$\lbl{A2}}
\BIC{$\forall n\,(\mathrm{Seq}(a) \to n \in \mathbb{N} \to a_{n} \in \mathbb{R})$}
\AXC{$\varphi_{92}$\lbl{A2}}
\BIC{$\mathrm{Seq}(a) \to n \in \mathbb{N} \to a_{n} \in \mathbb{R}$}
\AXC{$(\mathrm{Seq}(a) \to n \in \mathbb{N} \to a_{n} \in \mathbb{R}) \to \varphi_{24} \to \varphi_{238}$\lbl{A1}}
\BIC{$\varphi_{24} \to \varphi_{238}$}
\end{ptree}
\noindent$\Pi_{29}$:
\begin{ptree}
\AXC{$\varphi_{113}$\lbl{[E33]}}
\AXC{$\varphi_{113} \to \varphi_{105}$\lbl{A2}}
\BIC{$\varphi_{105}$}
\AXC{$\varphi_{105} \to \varphi_{121}$\lbl{A2}}
\BIC{$\varphi_{121}$}
\AXC{$\varphi_{121} \to s \in \mathbb{R} \to \varphi_{136}$\lbl{A2}}
\BIC{$s \in \mathbb{R} \to \varphi_{136}$}
\AXC{$(s \in \mathbb{R} \to \varphi_{136}) \to \varphi_{273}$\lbl{A1}}
\BIC{$\varphi_{273}$}
\end{ptree}
\noindent$\Pi_{30}$:
\begin{ptree}
\AXC{$\mathrm{MonoDec}(a) \to M \in \mathbb{R} \to \varphi_{69}$\lbl{A1}}
\AXC{$\stackrel{\Pi_{23}}{\vdots}$}\noLine
\UIC{$\mathrm{MonoDec}(a) \to M \in \mathbb{R} \to \varphi_{16}$}
\AXC{$\stackrel{\Pi_{29}}{\vdots}$}\noLine
\UIC{$\varphi_{273}$}
\BIC{$(\mathrm{MonoDec}(a) \to M \in \mathbb{R} \to \varphi_{69}) \to \varphi_{270}$}
\BIC{$\varphi_{270}$}
\end{ptree}
\noindent$\Pi_{31}$:
\begin{ptree}
\AXC{$\mathrm{MonoDec}(a) \to M \in \mathbb{R} \to \forall n \in \mathbb{N}\,(M \leq a_{n}) \to \varphi_{161}$\lbl{A1}}
\AXC{$\stackrel{\Pi_{11}}{\vdots}$}\noLine
\UIC{$\varphi_{24}$}
\AXC{$\stackrel{\Pi_{28}}{\vdots}$}\noLine
\UIC{$\varphi_{24} \to \varphi_{238}$}
\BIC{$\varphi_{238}$}
\BIC{$\mathrm{MonoDec}(a) \to M \in \mathbb{R} \to \forall n \in \mathbb{N}\,(M \leq a_{n}) \to \varphi_{176}$}
\AXC{$\stackrel{\Pi_{30}}{\vdots}$}\noLine
\UIC{$\varphi_{270}$}
\BIC{$\varphi_{267}$}
\end{ptree}
\noindent$\Pi_{32}$:
\begin{ptree}
\AXC{$\stackrel{\Pi_{8}}{\vdots}$}\noLine
\UIC{$\mathrm{MonoDec}(a) \to M \in \mathbb{R} \to \forall n \in \mathbb{N}\,(M \leq a_{n}) \to \varphi_{180}$}
\AXC{$\stackrel{\Pi_{27}}{\vdots}$}\noLine
\UIC{$\mathrm{MonoDec}(a) \to M \in \mathbb{R} \to \forall n \in \mathbb{N}\,(M \leq a_{n}) \to \varphi_{181}$}
\AXC{$\stackrel{\Pi_{31}}{\vdots}$}\noLine
\UIC{$\varphi_{267}$}
\BIC{$\varphi_{247}$}
\BIC{$\mathrm{MonoDec}(a) \to M \in \mathbb{R} \to \forall n \in \mathbb{N}\,(M \leq a_{n}) \to \varphi_{185}$}
\end{ptree}
\noindent$\Pi_{33}$:
\begin{ptree}
\AXC{$\varphi_{188} \to \forall n\,(\varphi_{188})$\lbl{A4}}
\AXC{$\varphi_{252}$\lbl{A3}}
\AXC{$\stackrel{\Pi_{32}}{\vdots}$}\noLine
\UIC{$\mathrm{MonoDec}(a) \to M \in \mathbb{R} \to \forall n \in \mathbb{N}\,(M \leq a_{n}) \to \varphi_{185}$}
\AXC{$\varphi_{249}$\lbl{A1}}
\BIC{$\varphi_{188} \to n \in \mathbb{N} \to N \leq n \to \lvert a_{n} - s \rvert < \varepsilon$}
\UIC{$\forall n\,(\varphi_{188} \to n \in \mathbb{N} \to N \leq n \to \lvert a_{n} - s \rvert < \varepsilon)$}
\AXC{$\varphi_{275}$\lbl{A1}}
\BIC{$\varphi_{252} \to (\varphi_{188} \to \forall n\,(\varphi_{188})) \to \mathrm{MonoDec}(a) \to M \in \mathbb{R} \to \varphi_{196}$}
\BIC{$(\varphi_{188} \to \forall n\,(\varphi_{188})) \to \mathrm{MonoDec}(a) \to M \in \mathbb{R} \to \varphi_{196}$}
\BIC{$\mathrm{MonoDec}(a) \to M \in \mathbb{R} \to \varphi_{196}$}
\end{ptree}
\noindent$\Pi_{34}$:
\begin{ptree}
\AXC{$\mathrm{MonoDec}(a) \to M \in \mathbb{R} \to \forall n \in \mathbb{N}\,(M \leq a_{n}) \to \varphi_{174}$\lbl{A1}}
\AXC{$\varphi_{233}$\lbl{A1}}
\BIC{$\mathrm{MonoDec}(a) \to M \in \mathbb{R} \to \varphi_{178}$}
\AXC{$(\mathrm{MonoDec}(a) \to M \in \mathbb{R} \to \varphi_{178}) \to \varphi_{253}$\lbl{A1}}
\BIC{$\varphi_{253}$}
\end{ptree}
\noindent$\Pi_{35}$:
\begin{ptree}
\AXC{$\stackrel{\Pi_{33}}{\vdots}$}\noLine
\UIC{$\mathrm{MonoDec}(a) \to M \in \mathbb{R} \to \varphi_{196}$}
\AXC{$\stackrel{\Pi_{34}}{\vdots}$}\noLine
\UIC{$\varphi_{253}$}
\BIC{$\mathrm{MonoDec}(a) \to M \in \mathbb{R} \to \varphi_{204}$}
\AXC{$(\mathrm{MonoDec}(a) \to M \in \mathbb{R} \to \varphi_{204}) \to \varphi_{229}$\lbl{A1}}
\BIC{$\varphi_{229}$}
\end{ptree}
\noindent$\Pi_{36}$:
\begin{ptree}
\AXC{$\varphi_{31} \to \neg (N \in \mathbb{N} \land \forall n \in \mathbb{N}\,(N \leq n \to \lvert a_{n} - s \rvert < \varepsilon))$\lbl{A2}}
\AXC{$\stackrel{\Pi_{35}}{\vdots}$}\noLine
\UIC{$\varphi_{229}$}
\BIC{$\mathrm{MonoDec}(a) \to M \in \mathbb{R} \to \varphi_{206}$}
\AXC{$\varphi_{255}$\lbl{A1}}
\BIC{$\varphi_{200} \to \neg (N \in \mathbb{N} \land a_{N} < s + \varepsilon)$}
\UIC{$\forall N\,(\varphi_{200} \to \neg (N \in \mathbb{N} \land a_{N} < s + \varepsilon))$}
\end{ptree}
\noindent$\Pi_{37}$:
\begin{ptree}
\AXC{$\mathrm{MonoDec}(a) \to M \in \mathbb{R} \to \varphi_{143}$\lbl{A1}}
\AXC{$\varphi_{223}$\lbl{A1}}
\BIC{$\mathrm{MonoDec}(a) \to M \in \mathbb{R} \to \varphi_{148}$}
\AXC{$\mathrm{MonoDec}(a) \to M \in \mathbb{R} \to \varphi_{159}$\lbl{A1}}
\AXC{$\varphi_{226}$\lbl{A1}}
\BIC{$\mathrm{MonoDec}(a) \to M \in \mathbb{R} \to \varphi_{146}$}
\AXC{$\varphi_{260}$\lbl{A1}}
\BIC{$(\mathrm{MonoDec}(a) \to M \in \mathbb{R} \to \varphi_{148}) \to \varphi_{207}$}
\BIC{$\varphi_{207}$}
\end{ptree}
\noindent$\Pi_{38}$:
\begin{ptree}
\AXC{$\mathrm{MonoDec}(a) \to M \in \mathbb{R} \to \varphi_{159}$\lbl{A1}}
\AXC{$\varphi_{225}$\lbl{A1}}
\BIC{$\mathrm{MonoDec}(a) \to M \in \mathbb{R} \to \varphi_{145}$}
\AXC{$(\mathrm{MonoDec}(a) \to M \in \mathbb{R} \to \varphi_{145}) \to \varphi_{227}$\lbl{A1}}
\BIC{$\varphi_{227}$}
\end{ptree}
\noindent$\Pi_{39}$:
\begin{ptree}
\AXC{$\varphi_{73}$\lbl{A2}}
\AXC{$z = w \to z < s + \varepsilon \to w < s + \varepsilon$\lbl{A6}}
\UIC{$\forall w\,(z = w \to z < s + \varepsilon \to w < s + \varepsilon)$}
\UIC{$\forall z\,\forall w\,(z = w \to z < s + \varepsilon \to w < s + \varepsilon)$}
\AXC{$\varphi_{93}$\lbl{A2}}
\BIC{$\forall w\,(x = w \to x < s + \varepsilon \to w < s + \varepsilon)$}
\AXC{$\varphi_{114}$\lbl{A2}}
\BIC{$x = a_{i_{0}} \to x < s + \varepsilon \to a_{i_{0}} < s + \varepsilon$}
\AXC{$\stackrel{\Pi_{37}}{\vdots}$}\noLine
\UIC{$\varphi_{207}$}
\BIC{$\mathrm{MonoDec}(a) \to M \in \mathbb{R} \to \varphi_{155}$}
\AXC{$\stackrel{\Pi_{38}}{\vdots}$}\noLine
\UIC{$\varphi_{227}$}
\BIC{$\mathrm{MonoDec}(a) \to M \in \mathbb{R} \to \varphi_{167}$}
\AXC{$\varphi_{246}$\lbl{A1}}
\BIC{$\varphi_{73} \to \mathrm{MonoDec}(a) \to M \in \mathbb{R} \to \varphi_{172}$}
\BIC{$\mathrm{MonoDec}(a) \to M \in \mathbb{R} \to \varphi_{172}$}
\end{ptree}
\noindent$\Pi_{40}$:
\begin{ptree}
\AXC{$\varphi_{200} \to \forall N\,(\varphi_{200})$\lbl{A4}}
\AXC{$\varphi_{257}$\lbl{A3}}
\AXC{$\stackrel{\Pi_{36}}{\vdots}$}\noLine
\UIC{$\forall N\,(\varphi_{200} \to \neg (N \in \mathbb{N} \land a_{N} < s + \varepsilon))$}
\AXC{$\stackrel{\Pi_{39}}{\vdots}$}\noLine
\UIC{$\mathrm{MonoDec}(a) \to M \in \mathbb{R} \to \varphi_{172}$}
\AXC{$\varphi_{277}$\lbl{A1}}
\BIC{$\forall N\,(\varphi_{200} \to \neg (N \in \mathbb{N} \land a_{N} < s + \varepsilon)) \to \varphi_{257} \to \varphi_{264}$}
\BIC{$\varphi_{257} \to \varphi_{264}$}
\BIC{$\varphi_{264}$}
\BIC{$\mathrm{MonoDec}(a) \to M \in \mathbb{R} \to \forall n \in \mathbb{N}\,(M \leq a_{n}) \to \varphi_{165}$}
\end{ptree}
\noindent$\Pi_{41}$:
\begin{ptree}
\AXC{$\stackrel{\Pi_{40}}{\vdots}$}\noLine
\UIC{$\mathrm{MonoDec}(a) \to M \in \mathbb{R} \to \forall n \in \mathbb{N}\,(M \leq a_{n}) \to \varphi_{165}$}
\AXC{$\varphi_{236}$\lbl{A1}}
\BIC{$\varphi_{190} \to \neg (i_{0} \in \mathbb{N} \land x = a_{i_{0}})$}
\UIC{$\forall i_{0}\,(\varphi_{190} \to \neg (i_{0} \in \mathbb{N} \land x = a_{i_{0}}))$}
\end{ptree}
\noindent$\Pi_{42}$:
\begin{ptree}
\AXC{$\varphi_{100}$\lbl{A2}}
\AXC{$\stackrel{\Pi_{12}}{\vdots}$}\noLine
\UIC{$\mathrm{MonoDec}(a) \to M \in \mathbb{R} \to \forall n \in \mathbb{N}\,(M \leq a_{n}) \to \varphi_{32}$}
\AXC{$\varphi_{218}$\lbl{A1}}
\BIC{$\mathrm{MonoDec}(a) \to M \in \mathbb{R} \to \varphi_{151}$}
\AXC{$\varphi_{235}$\lbl{A1}}
\BIC{$\varphi_{100} \to \mathrm{MonoDec}(a) \to M \in \mathbb{R} \to \varphi_{149}$}
\BIC{$\mathrm{MonoDec}(a) \to M \in \mathbb{R} \to \varphi_{149}$}
\AXC{$\varphi_{248}$\lbl{A1}}
\BIC{$(\mathrm{MonoDec}(a) \to \varphi_{141}) \to \mathrm{MonoDec}(a) \to \varphi_{147}$}
\end{ptree}
\noindent$\Pi_{43}$:
\begin{ptree}
\AXC{$\mathrm{MonoDec}(a) \to M \in \mathbb{R} \to \varphi_{143}$\lbl{A1}}
\AXC{$\varphi_{220}$\lbl{A1}}
\BIC{$\mathrm{MonoDec}(a) \to \varphi_{141}$}
\AXC{$\stackrel{\Pi_{42}}{\vdots}$}\noLine
\UIC{$(\mathrm{MonoDec}(a) \to \varphi_{141}) \to \mathrm{MonoDec}(a) \to \varphi_{147}$}
\BIC{$\mathrm{MonoDec}(a) \to \varphi_{147}$}
\end{ptree}
\noindent$\Pi_{44}$:
\begin{ptree}
\AXC{$\varphi_{243}$\lbl{A3}}
\AXC{$\stackrel{\Pi_{41}}{\vdots}$}\noLine
\UIC{$\forall i_{0}\,(\varphi_{190} \to \neg (i_{0} \in \mathbb{N} \land x = a_{i_{0}}))$}
\AXC{$\stackrel{\Pi_{43}}{\vdots}$}\noLine
\UIC{$\mathrm{MonoDec}(a) \to \varphi_{147}$}
\AXC{$(\mathrm{MonoDec}(a) \to \varphi_{147}) \to \varphi_{276}$\lbl{A1}}
\BIC{$\varphi_{276}$}
\BIC{$\varphi_{243} \to (\varphi_{190} \to \forall i_{0}\,(\varphi_{190})) \to \varphi_{239}$}
\BIC{$(\varphi_{190} \to \forall i_{0}\,(\varphi_{190})) \to \varphi_{239}$}
\end{ptree}
\noindent$\Pi_{45}$:
\begin{ptree}
\AXC{$\varphi_{36}$\lbl{A4}}
\AXC{$\forall j\,(\varphi_{11}) \to \varphi_{27}$\lbl{A3}}
\AXC{$\varphi_{11}$\lbl{A2}}
\UIC{$\forall j\,(\varphi_{11})$}
\AXC{$\varphi_{190} \to \forall i_{0}\,(\varphi_{190})$\lbl{A4}}
\AXC{$\stackrel{\Pi_{44}}{\vdots}$}\noLine
\UIC{$(\varphi_{190} \to \forall i_{0}\,(\varphi_{190})) \to \varphi_{239}$}
\BIC{$\varphi_{239}$}
\BIC{$(\forall j\,(\varphi_{11}) \to \varphi_{27}) \to \varphi_{36} \to \mathrm{MonoDec}(a) \to M \in \mathbb{R} \to \varphi_{168}$}
\BIC{$\varphi_{36} \to \mathrm{MonoDec}(a) \to M \in \mathbb{R} \to \varphi_{168}$}
\BIC{$\mathrm{MonoDec}(a) \to M \in \mathbb{R} \to \varphi_{168}$}
\AXC{$\varphi_{230}$\lbl{A1}}
\BIC{$\varphi_{162} \to \neg (x \in \mathrm{ran}(a) \land x < s + \varepsilon)$}
\UIC{$\forall x\,(\varphi_{162} \to \neg (x \in \mathrm{ran}(a) \land x < s + \varepsilon))$}
\end{ptree}
\noindent$\Pi_{46}$:
\begin{ptree}
\AXC{$\stackrel{\Pi_{12}}{\vdots}$}\noLine
\UIC{$\mathrm{MonoDec}(a) \to M \in \mathbb{R} \to \forall n \in \mathbb{N}\,(M \leq a_{n}) \to \varphi_{32}$}
\AXC{$\varphi_{189}$\lbl{A1}}
\BIC{$\varphi_{59}$}
\end{ptree}
\noindent$\Pi_{47}$:
\begin{ptree}
\AXC{$\forall s\,(\varphi_{86})$\lbl{[C23]}}
\AXC{$\forall s\,(\varphi_{86}) \to \varphi_{86}$\lbl{A2}}
\BIC{$\varphi_{86}$}
\AXC{$\varphi_{86} \to \varphi_{112}$\lbl{A2}}
\BIC{$\varphi_{112}$}
\AXC{$\varphi_{112} \to \mathrm{ran}(a) \subseteq \mathbb{R} \to \varphi_{75}$\lbl{A2}}
\BIC{$\mathrm{ran}(a) \subseteq \mathbb{R} \to \varphi_{75}$}
\end{ptree}
\noindent$\Pi_{48}$:
\begin{ptree}
\AXC{$\stackrel{\Pi_{46}}{\vdots}$}\noLine
\UIC{$\varphi_{59}$}
\AXC{$\stackrel{\Pi_{47}}{\vdots}$}\noLine
\UIC{$\mathrm{ran}(a) \subseteq \mathbb{R} \to \varphi_{75}$}
\AXC{$(\mathrm{ran}(a) \subseteq \mathbb{R} \to \varphi_{75}) \to \varphi_{59} \to \varphi_{234}$\lbl{A1}}
\BIC{$\varphi_{59} \to \varphi_{234}$}
\BIC{$\varphi_{234}$}
\end{ptree}
\noindent$\Pi_{49}$:
\begin{ptree}
\AXC{$\mathrm{MonoDec}(a) \to M \in \mathbb{R} \to \varphi_{69}$\lbl{A1}}
\AXC{$\mathrm{MonoDec}(a) \to M \in \mathbb{R} \to \forall n \in \mathbb{N}\,(M \leq a_{n}) \to \varphi_{2}$\lbl{A1}}
\AXC{$\stackrel{\Pi_{48}}{\vdots}$}\noLine
\UIC{$\varphi_{234}$}
\BIC{$(\mathrm{MonoDec}(a) \to M \in \mathbb{R} \to \varphi_{69}) \to \varphi_{211}$}
\BIC{$\varphi_{211}$}
\end{ptree}
\noindent$\Pi_{50}$:
\begin{ptree}
\AXC{$\varphi_{231}$\lbl{A3}}
\AXC{$\stackrel{\Pi_{45}}{\vdots}$}\noLine
\UIC{$\forall x\,(\varphi_{162} \to \neg (x \in \mathrm{ran}(a) \land x < s + \varepsilon))$}
\AXC{$\mathrm{MonoDec}(a) \to M \in \mathbb{R} \to \varphi_{89}$\lbl{A1}}
\AXC{$\stackrel{\Pi_{49}}{\vdots}$}\noLine
\UIC{$\varphi_{211}$}
\BIC{$\mathrm{MonoDec}(a) \to M \in \mathbb{R} \to \forall n \in \mathbb{N}\,(M \leq a_{n}) \to \varphi_{75}$}
\AXC{$\varphi_{274}$\lbl{A1}}
\BIC{$\varphi_{272}$}
\BIC{$\varphi_{231} \to (\varphi_{162} \to \forall x\,(\varphi_{162})) \to \mathrm{MonoDec}(a) \to \varphi_{152}$}
\BIC{$(\varphi_{162} \to \forall x\,(\varphi_{162})) \to \mathrm{MonoDec}(a) \to \varphi_{152}$}
\end{ptree}
\noindent$\Pi_{51}$:
\begin{ptree}
\AXC{$\varphi_{162} \to \forall x\,(\varphi_{162})$\lbl{A4}}
\AXC{$\stackrel{\Pi_{50}}{\vdots}$}\noLine
\UIC{$(\varphi_{162} \to \forall x\,(\varphi_{162})) \to \mathrm{MonoDec}(a) \to \varphi_{152}$}
\BIC{$\mathrm{MonoDec}(a) \to \varphi_{152}$}
\AXC{$(\mathrm{MonoDec}(a) \to \varphi_{152}) \to \varphi_{56} \to \varphi_{102}$\lbl{A1}}
\BIC{$\varphi_{56} \to \varphi_{102}$}
\UIC{$\forall \varepsilon\,(\varphi_{56} \to \varphi_{102})$}
\end{ptree}
\noindent$\Pi_{52}$:
\begin{ptree}
\AXC{$\varphi_{56} \to \forall \varepsilon\,(\varphi_{56})$\lbl{A4}}
\AXC{$\forall \varepsilon\,(\varphi_{56} \to \varphi_{102}) \to \forall \varepsilon\,(\varphi_{56}) \to \varphi_{115}$\lbl{A3}}
\AXC{$\stackrel{\Pi_{51}}{\vdots}$}\noLine
\UIC{$\forall \varepsilon\,(\varphi_{56} \to \varphi_{102})$}
\AXC{$\forall \varepsilon\,(\varphi_{56} \to \varphi_{102}) \to \varphi_{262}$\lbl{A1}}
\BIC{$\varphi_{262}$}
\BIC{$(\varphi_{56} \to \forall \varepsilon\,(\varphi_{56})) \to \mathrm{MonoDec}(a) \to \varphi_{156}$}
\BIC{$\mathrm{MonoDec}(a) \to \varphi_{156}$}
\end{ptree}
\noindent$\Pi_{53}$:
\begin{ptree}
\AXC{$\stackrel{\Pi_{11}}{\vdots}$}\noLine
\UIC{$\varphi_{24}$}
\AXC{$\forall a\,(\varphi_{144})$\lbl{DefConv}}
\AXC{$\forall a\,(\varphi_{144}) \to \varphi_{144}$\lbl{A2}}
\BIC{$\varphi_{144}$}
\AXC{$\varphi_{144} \to (\mathrm{Conv}(a,s) \leftrightarrow \mathrm{Seq}(a) \land s \in \mathbb{R} \land \varphi_{115})$\lbl{A2}}
\BIC{$\mathrm{Conv}(a,s) \leftrightarrow \mathrm{Seq}(a) \land s \in \mathbb{R} \land \varphi_{115}$}
\AXC{$\varphi_{258}$\lbl{A1}}
\BIC{$\varphi_{24} \to (\mathrm{MonoDec}(a) \to M \in \mathbb{R} \to \varphi_{16}) \to \varphi_{212}$}
\BIC{$(\mathrm{MonoDec}(a) \to M \in \mathbb{R} \to \varphi_{16}) \to \varphi_{212}$}
\end{ptree}
\noindent$\Pi_{54}$:
\begin{ptree}
\AXC{$\stackrel{\Pi_{52}}{\vdots}$}\noLine
\UIC{$\mathrm{MonoDec}(a) \to \varphi_{156}$}
\AXC{$\stackrel{\Pi_{23}}{\vdots}$}\noLine
\UIC{$\mathrm{MonoDec}(a) \to M \in \mathbb{R} \to \varphi_{16}$}
\AXC{$\stackrel{\Pi_{53}}{\vdots}$}\noLine
\UIC{$(\mathrm{MonoDec}(a) \to M \in \mathbb{R} \to \varphi_{16}) \to \varphi_{212}$}
\BIC{$\varphi_{212}$}
\BIC{$\mathrm{MonoDec}(a) \to M \in \mathbb{R} \to \varphi_{13}$}
\end{ptree}
\noindent$\Pi_{55}$:
\begin{ptree}
\AXC{$\stackrel{\Pi_{54}}{\vdots}$}\noLine
\UIC{$\mathrm{MonoDec}(a) \to M \in \mathbb{R} \to \varphi_{13}$}
\AXC{$(\mathrm{MonoDec}(a) \to M \in \mathbb{R} \to \varphi_{13}) \to \varphi_{137}$\lbl{A1}}
\BIC{$\varphi_{137}$}
\end{ptree}
\noindent$\Pi_{56}$:
\begin{ptree}
\AXC{$\forall n \in \mathbb{N}\,(M \leq a_{n}) \to n \in \mathbb{N} \to M \leq a_{n}$\lbl{A2}}
\AXC{$\varphi_{87}$\lbl{A1}}
\AXC{$\varphi_{87} \to (\forall n \in \mathbb{N}\,(M \leq a_{n}) \to n \in \mathbb{N} \to M \leq a_{n}) \to \varphi_{77}$\lbl{A1}}
\BIC{$(\forall n \in \mathbb{N}\,(M \leq a_{n}) \to n \in \mathbb{N} \to M \leq a_{n}) \to \varphi_{77}$}
\BIC{$\varphi_{77}$}
\AXC{$\varphi_{77} \to \varphi_{195}$\lbl{A1}}
\BIC{$\varphi_{195}$}
\end{ptree}
\noindent$\Pi_{57}$:
\begin{ptree}
\AXC{$x = x$\lbl{A5}}
\AXC{$x = y \to x = x \to y = x$\lbl{A6}}
\AXC{$(x = y \to x = x \to y = x) \to x = x \to x = y \to y = x$\lbl{A1}}
\BIC{$x = x \to x = y \to y = x$}
\BIC{$x = y \to y = x$}
\UIC{$\forall y\,(x = y \to y = x)$}
\UIC{$\forall x\,\forall y\,(x = y \to y = x)$}
\end{ptree}
\noindent$\Pi_{58}$:
\begin{ptree}
\AXC{$\stackrel{\Pi_{57}}{\vdots}$}\noLine
\UIC{$\forall x\,\forall y\,(x = y \to y = x)$}
\AXC{$\forall x\,\forall y\,(x = y \to y = x) \to \forall y\,(v_{720} = y \to y = v_{720})$\lbl{A2}}
\BIC{$\forall y\,(v_{720} = y \to y = v_{720})$}
\AXC{$\varphi_{5}$\lbl{A2}}
\BIC{$v_{720} = v_{721} \to v_{721} = v_{720}$}
\UIC{$\forall v_{721}\,(v_{720} = v_{721} \to v_{721} = v_{720})$}
\UIC{$\forall v_{720}\,\forall v_{721}\,(v_{720} = v_{721} \to v_{721} = v_{720})$}
\AXC{$\varphi_{19}$\lbl{A2}}
\BIC{$\forall v_{721}\,(y = v_{721} \to v_{721} = y)$}
\end{ptree}
\noindent$\Pi_{59}$:
\begin{ptree}
\AXC{$\stackrel{\Pi_{58}}{\vdots}$}\noLine
\UIC{$\forall v_{721}\,(y = v_{721} \to v_{721} = y)$}
\AXC{$\forall v_{721}\,(y = v_{721} \to v_{721} = y) \to y = a_{n} \to a_{n} = y$\lbl{A2}}
\BIC{$y = a_{n} \to a_{n} = y$}
\end{ptree}
\noindent$\Pi_{60}$:
\begin{ptree}
\AXC{$\stackrel{\Pi_{59}}{\vdots}$}\noLine
\UIC{$y = a_{n} \to a_{n} = y$}
\AXC{$\mathrm{MonoDec}(a) \to M \in \mathbb{R} \to \varphi_{111}$\lbl{A1}}
\AXC{$\varphi_{202}$\lbl{A1}}
\BIC{$\mathrm{MonoDec}(a) \to M \in \mathbb{R} \to \varphi_{84}$}
\AXC{$\varphi_{209}$\lbl{A1}}
\BIC{$(y = a_{n} \to a_{n} = y) \to \mathrm{MonoDec}(a) \to M \in \mathbb{R} \to \varphi_{85}$}
\BIC{$\mathrm{MonoDec}(a) \to M \in \mathbb{R} \to \varphi_{85}$}
\end{ptree}
\noindent$\Pi_{61}$:
\begin{ptree}
\AXC{$\stackrel{\Pi_{60}}{\vdots}$}\noLine
\UIC{$\mathrm{MonoDec}(a) \to M \in \mathbb{R} \to \varphi_{85}$}
\AXC{$(\mathrm{MonoDec}(a) \to M \in \mathbb{R} \to \varphi_{85}) \to \varphi_{214}$\lbl{A1}}
\BIC{$\varphi_{214}$}
\end{ptree}
\noindent$\Pi_{62}$:
\begin{ptree}
\AXC{$\mathrm{MonoDec}(a) \to M \in \mathbb{R} \to \varphi_{111}$\lbl{A1}}
\AXC{$\varphi_{201}$\lbl{A1}}
\BIC{$\mathrm{MonoDec}(a) \to M \in \mathbb{R} \to \varphi_{79}$}
\AXC{$\stackrel{\Pi_{56}}{\vdots}$}\noLine
\UIC{$\varphi_{195}$}
\BIC{$\mathrm{MonoDec}(a) \to M \in \mathbb{R} \to \varphi_{83}$}
\AXC{$\stackrel{\Pi_{61}}{\vdots}$}\noLine
\UIC{$\varphi_{214}$}
\BIC{$(a_{n} = y \to M < a_{n} \to M < y) \to \varphi_{142}$}
\end{ptree}
\noindent$\Pi_{63}$:
\begin{ptree}
\AXC{$z = w \to M < z \to M < w$\lbl{A6}}
\UIC{$\forall w\,(z = w \to M < z \to M < w)$}
\UIC{$\forall z\,\forall w\,(z = w \to M < z \to M < w)$}
\AXC{$\varphi_{48}$\lbl{A2}}
\BIC{$\forall w\,(a_{n} = w \to M < a_{n} \to M < w)$}
\AXC{$\varphi_{76}$\lbl{A2}}
\BIC{$a_{n} = y \to M < a_{n} \to M < y$}
\AXC{$\stackrel{\Pi_{62}}{\vdots}$}\noLine
\UIC{$(a_{n} = y \to M < a_{n} \to M < y) \to \varphi_{142}$}
\BIC{$\varphi_{142}$}
\end{ptree}
\noindent$\Pi_{64}$:
\begin{ptree}
\AXC{$z = w \to M = z \to M = w$\lbl{A6}}
\UIC{$\forall w\,(z = w \to M = z \to M = w)$}
\UIC{$\forall z\,\forall w\,(z = w \to M = z \to M = w)$}
\AXC{$\varphi_{15}$\lbl{A2}}
\BIC{$\forall w\,(a_{n} = w \to M = a_{n} \to M = w)$}
\AXC{$\varphi_{37}$\lbl{A2}}
\BIC{$a_{n} = y \to M = a_{n} \to M = y$}
\AXC{$\stackrel{\Pi_{63}}{\vdots}$}\noLine
\UIC{$\varphi_{142}$}
\BIC{$\mathrm{MonoDec}(a) \to M \in \mathbb{R} \to \varphi_{64}$}
\AXC{$\varphi_{193}$\lbl{A1}}
\BIC{$\varphi_{70} \to \neg (n \in \mathbb{N} \land y = a_{n})$}
\UIC{$\forall n\,(\varphi_{70} \to \neg (n \in \mathbb{N} \land y = a_{n}))$}
\end{ptree}
\noindent$\Pi_{65}$:
\begin{ptree}
\AXC{$\varphi_{101}$\lbl{A2}}
\AXC{$\stackrel{\Pi_{12}}{\vdots}$}\noLine
\UIC{$\mathrm{MonoDec}(a) \to M \in \mathbb{R} \to \forall n \in \mathbb{N}\,(M \leq a_{n}) \to \varphi_{32}$}
\AXC{$\varphi_{208}$\lbl{A1}}
\BIC{$\mathrm{MonoDec}(a) \to M \in \mathbb{R} \to \forall n \in \mathbb{N}\,(M \leq a_{n}) \to \varphi_{23}$}
\AXC{$\varphi_{221}$\lbl{A1}}
\BIC{$\varphi_{192}$}
\BIC{$\mathrm{MonoDec}(a) \to M \in \mathbb{R} \to \forall n \in \mathbb{N}\,(M \leq a_{n}) \to \varphi_{14}$}
\AXC{$\varphi_{217}$\lbl{A1}}
\BIC{$\varphi_{183}$}
\end{ptree}
\noindent$\Pi_{66}$:
\begin{ptree}
\AXC{$\stackrel{\Pi_{64}}{\vdots}$}\noLine
\UIC{$\forall n\,(\varphi_{70} \to \neg (n \in \mathbb{N} \land y = a_{n}))$}
\AXC{$\mathrm{MonoDec}(a) \to \varphi_{53}$\lbl{A1}}
\AXC{$\stackrel{\Pi_{65}}{\vdots}$}\noLine
\UIC{$\varphi_{183}$}
\BIC{$\mathrm{MonoDec}(a) \to M \in \mathbb{R} \to \varphi_{54}$}
\AXC{$(\mathrm{MonoDec}(a) \to M \in \mathbb{R} \to \varphi_{54}) \to \varphi_{268}$\lbl{A1}}
\BIC{$\varphi_{268}$}
\BIC{$\varphi_{194} \to (\varphi_{70} \to \forall n\,(\varphi_{70})) \to \varphi_{12} \to (\varphi_{12} \to \varphi_{22}) \to \varphi_{21} \to \varphi_{66}$}
\end{ptree}
\noindent$\Pi_{67}$:
\begin{ptree}
\AXC{$\varphi_{70} \to \forall n\,(\varphi_{70})$\lbl{A4}}
\AXC{$\varphi_{194}$\lbl{A3}}
\AXC{$\stackrel{\Pi_{66}}{\vdots}$}\noLine
\UIC{$\varphi_{194} \to (\varphi_{70} \to \forall n\,(\varphi_{70})) \to \varphi_{12} \to (\varphi_{12} \to \varphi_{22}) \to \varphi_{21} \to \varphi_{66}$}
\BIC{$(\varphi_{70} \to \forall n\,(\varphi_{70})) \to \varphi_{12} \to (\varphi_{12} \to \varphi_{22}) \to \varphi_{21} \to \varphi_{66}$}
\BIC{$\varphi_{12} \to (\varphi_{12} \to \varphi_{22}) \to \varphi_{21} \to \varphi_{66}$}
\end{ptree}
\noindent$\Pi_{68}$:
\begin{ptree}
\AXC{$\varphi_{21}$\lbl{A4}}
\AXC{$\varphi_{12} \to \varphi_{22}$\lbl{A3}}
\AXC{$\forall n\,(\neg (n \in \mathbb{N} \land y = a_{n})) \to \neg (j \in \mathbb{N} \land y = a_{j})$\lbl{A2}}
\UIC{$\varphi_{12}$}
\AXC{$\stackrel{\Pi_{67}}{\vdots}$}\noLine
\UIC{$\varphi_{12} \to (\varphi_{12} \to \varphi_{22}) \to \varphi_{21} \to \varphi_{66}$}
\BIC{$(\varphi_{12} \to \varphi_{22}) \to \varphi_{21} \to \varphi_{66}$}
\BIC{$\varphi_{21} \to \varphi_{66}$}
\BIC{$\varphi_{66}$}
\end{ptree}
\noindent$\Pi_{69}$:
\begin{ptree}
\AXC{$\forall y\,(\varphi_{67}) \to \varphi_{10} \to \forall y \in \mathrm{ran}(a)\,(M \leq y)$\lbl{A3}}
\AXC{$\stackrel{\Pi_{68}}{\vdots}$}\noLine
\UIC{$\varphi_{66}$}
\AXC{$\varphi_{66} \to \varphi_{67}$\lbl{A1}}
\BIC{$\varphi_{67}$}
\UIC{$\forall y\,(\varphi_{67})$}
\AXC{$\varphi_{232}$\lbl{A1}}
\BIC{$(\forall y\,(\varphi_{67}) \to \varphi_{10} \to \forall y \in \mathrm{ran}(a)\,(M \leq y)) \to \varphi_{186}$}
\BIC{$\varphi_{186}$}
\end{ptree}
\noindent$\Pi_{70}$:
\begin{ptree}
\AXC{$\forall a\,\forall X\,(\mathrm{LB}(a,X) \leftrightarrow a \in \mathbb{R} \land \forall x \in X\,(a \leq x))$\lbl{DefLB}}
\AXC{$\varphi_{109}$\lbl{A2}}
\BIC{$\forall X\,(\mathrm{LB}(M,X) \leftrightarrow M \in \mathbb{R} \land \forall x \in X\,(M \leq x))$}
\AXC{$\forall X\,(\mathrm{LB}(M,X) \leftrightarrow M \in \mathbb{R} \land \forall x \in X\,(M \leq x)) \to \varphi_{8}$\lbl{A2}}
\BIC{$\varphi_{8}$}
\end{ptree}
\noindent$\Pi_{71}$:
\begin{ptree}
\AXC{$\mathrm{MonoDec}(a) \to M \in \mathbb{R} \to M \in \mathbb{R}$\lbl{A1}}
\AXC{$\stackrel{\Pi_{70}}{\vdots}$}\noLine
\UIC{$\varphi_{8}$}
\AXC{$\varphi_{8} \to \varphi_{222}$\lbl{A1}}
\BIC{$\varphi_{222}$}
\BIC{$(\mathrm{MonoDec}(a) \to \varphi_{33}) \to \varphi_{40} \to \varphi_{39} \to \varphi_{43}$}
\end{ptree}
\noindent$\Pi_{72}$:
\begin{ptree}
\AXC{$\mathrm{MonoDec}(a) \land M \in \mathbb{R} \land \forall n \in \mathbb{N}\,(M \leq a_{n}) \to \varphi_{10}$\lbl{A4}}
\AXC{$\stackrel{\Pi_{69}}{\vdots}$}\noLine
\UIC{$\varphi_{186}$}
\BIC{$\mathrm{MonoDec}(a) \to \varphi_{33}$}
\AXC{$\stackrel{\Pi_{71}}{\vdots}$}\noLine
\UIC{$(\mathrm{MonoDec}(a) \to \varphi_{33}) \to \varphi_{40} \to \varphi_{39} \to \varphi_{43}$}
\BIC{$\varphi_{40} \to \varphi_{39} \to \varphi_{43}$}
\end{ptree}
\noindent$\Pi_{73}$:
\begin{ptree}
\AXC{$\varphi_{47}$\lbl{A4}}
\AXC{$\forall y\,(\varphi_{30}) \to \varphi_{52}$\lbl{A3}}
\AXC{$\varphi_{30}$\lbl{A2}}
\UIC{$\forall y\,(\varphi_{30})$}
\AXC{$\forall y\,(\varphi_{30}) \to (\forall y\,(\varphi_{30}) \to \varphi_{52}) \to \varphi_{47} \to \varphi_{40}$\lbl{A1}}
\BIC{$(\forall y\,(\varphi_{30}) \to \varphi_{52}) \to \varphi_{47} \to \varphi_{40}$}
\BIC{$\varphi_{47} \to \varphi_{40}$}
\BIC{$\varphi_{40}$}
\AXC{$\stackrel{\Pi_{72}}{\vdots}$}\noLine
\UIC{$\varphi_{40} \to \varphi_{39} \to \varphi_{43}$}
\BIC{$\varphi_{39} \to \varphi_{43}$}
\end{ptree}
\noindent$\Pi_{74}$:
\begin{ptree}
\AXC{$\varphi_{46}$\lbl{A4}}
\AXC{$\forall x\,(\varphi_{29}) \to \varphi_{51}$\lbl{A3}}
\AXC{$\varphi_{29}$\lbl{A2}}
\UIC{$\forall x\,(\varphi_{29})$}
\AXC{$\forall x\,(\varphi_{29}) \to (\forall x\,(\varphi_{29}) \to \varphi_{51}) \to \varphi_{46} \to \varphi_{39}$\lbl{A1}}
\BIC{$(\forall x\,(\varphi_{29}) \to \varphi_{51}) \to \varphi_{46} \to \varphi_{39}$}
\BIC{$\varphi_{46} \to \varphi_{39}$}
\BIC{$\varphi_{39}$}
\AXC{$\stackrel{\Pi_{73}}{\vdots}$}\noLine
\UIC{$\varphi_{39} \to \varphi_{43}$}
\BIC{$\varphi_{43}$}
\end{ptree}
\noindent$\Pi_{75}$:
\begin{ptree}
\AXC{$x = x$\lbl{A5}}
\UIC{$\forall x\,(x = x)$}
\AXC{$\forall x\,(x = x) \to a_{1_{\omega}} = a_{1_{\omega}}$\lbl{A2}}
\BIC{$a_{1_{\omega}} = a_{1_{\omega}}$}
\AXC{$1_{\omega} \in \mathbb{N}$\lbl{[N1]}}
\AXC{$1_{\omega} \in \mathbb{N} \to \varphi_{139}$\lbl{A1}}
\BIC{$\varphi_{139}$}
\BIC{$\mathrm{MonoDec}(a) \to M \in \mathbb{R} \to \varphi_{44}$}
\end{ptree}
\noindent$\Pi_{76}$:
\begin{ptree}
\AXC{$\stackrel{\Pi_{12}}{\vdots}$}\noLine
\UIC{$\mathrm{MonoDec}(a) \to M \in \mathbb{R} \to \forall n \in \mathbb{N}\,(M \leq a_{n}) \to \varphi_{32}$}
\AXC{$\varphi_{199}$\lbl{A1}}
\BIC{$\mathrm{MonoDec}(a) \to M \in \mathbb{R} \to \varphi_{65}$}
\AXC{$(\mathrm{MonoDec}(a) \to M \in \mathbb{R} \to \varphi_{65}) \to \varphi_{203}$\lbl{A1}}
\BIC{$\varphi_{203}$}
\end{ptree}
\noindent$\Pi_{77}$:
\begin{ptree}
\AXC{$\forall y\,(y \in \mathrm{ran}(a) \leftrightarrow \exists j \in \mathbb{N}\,(y = a_{j})) \to \varphi_{26}$\lbl{A2}}
\AXC{$\stackrel{\Pi_{76}}{\vdots}$}\noLine
\UIC{$\varphi_{203}$}
\BIC{$\mathrm{MonoDec}(a) \to M \in \mathbb{R} \to \forall n \in \mathbb{N}\,(M \leq a_{n}) \to \varphi_{26}$}
\AXC{$\varphi_{216}$\lbl{A1}}
\BIC{$(\mathrm{MonoDec}(a) \to \varphi_{78}) \to \mathrm{MonoDec}(a) \to \varphi_{35}$}
\end{ptree}
\noindent$\Pi_{78}$:
\begin{ptree}
\AXC{$\varphi_{80}$\lbl{A2}}
\AXC{$\stackrel{\Pi_{75}}{\vdots}$}\noLine
\UIC{$\mathrm{MonoDec}(a) \to M \in \mathbb{R} \to \varphi_{44}$}
\AXC{$\varphi_{213}$\lbl{A1}}
\BIC{$\varphi_{80} \to \mathrm{MonoDec}(a) \to \varphi_{78}$}
\BIC{$\mathrm{MonoDec}(a) \to \varphi_{78}$}
\AXC{$\stackrel{\Pi_{77}}{\vdots}$}\noLine
\UIC{$(\mathrm{MonoDec}(a) \to \varphi_{78}) \to \mathrm{MonoDec}(a) \to \varphi_{35}$}
\BIC{$\mathrm{MonoDec}(a) \to \varphi_{35}$}
\end{ptree}
\noindent$\Pi_{79}$:
\begin{ptree}
\AXC{$\stackrel{\Pi_{46}}{\vdots}$}\noLine
\UIC{$\varphi_{59}$}
\AXC{$\varphi_{20}$\lbl{公理3.1(2)}}
\AXC{$\varphi_{20} \to \mathrm{ran}(a) \subseteq \mathbb{R} \to \mathrm{ran}(a) \neq \emptyset \to \varphi_{3}$\lbl{A2}}
\BIC{$\mathrm{ran}(a) \subseteq \mathbb{R} \to \mathrm{ran}(a) \neq \emptyset \to \varphi_{3}$}
\AXC{$(\mathrm{ran}(a) \subseteq \mathbb{R} \to \mathrm{ran}(a) \neq \emptyset \to \varphi_{3}) \to \varphi_{59} \to \varphi_{60} \to \varphi_{150}$\lbl{A1}}
\BIC{$\varphi_{59} \to \varphi_{60} \to \varphi_{150}$}
\BIC{$\varphi_{60} \to \varphi_{150}$}
\end{ptree}
\noindent$\Pi_{80}$:
\begin{ptree}
\AXC{$\stackrel{\Pi_{78}}{\vdots}$}\noLine
\UIC{$\mathrm{MonoDec}(a) \to \varphi_{35}$}
\AXC{$\forall X\,\forall x \in X\,(X \neq \emptyset)$\lbl{[B2]}}
\AXC{$\varphi_{38}$\lbl{A2}}
\BIC{$\forall x \in \mathrm{ran}(a)\,(\mathrm{ran}(a) \neq \emptyset)$}
\AXC{$\varphi_{81}$\lbl{A2}}
\BIC{$a_{1_{\omega}} \in \mathrm{ran}(a) \to \mathrm{ran}(a) \neq \emptyset$}
\AXC{$\varphi_{191}$\lbl{A1}}
\BIC{$(\mathrm{MonoDec}(a) \to \varphi_{35}) \to \varphi_{60}$}
\BIC{$\varphi_{60}$}
\AXC{$\stackrel{\Pi_{79}}{\vdots}$}\noLine
\UIC{$\varphi_{60} \to \varphi_{150}$}
\BIC{$\varphi_{150}$}
\end{ptree}
\noindent$\Pi_{81}$:
\begin{ptree}
\AXC{$\varphi_{1}$\lbl{A2}}
\AXC{$\stackrel{\Pi_{74}}{\vdots}$}\noLine
\UIC{$\varphi_{43}$}
\AXC{$\varphi_{43} \to \varphi_{1} \to \mathrm{MonoDec}(a) \to \varphi_{17}$\lbl{A1}}
\BIC{$\varphi_{1} \to \mathrm{MonoDec}(a) \to \varphi_{17}$}
\BIC{$\mathrm{MonoDec}(a) \to \varphi_{17}$}
\AXC{$\stackrel{\Pi_{80}}{\vdots}$}\noLine
\UIC{$\varphi_{150}$}
\BIC{$\mathrm{MonoDec}(a) \to M \in \mathbb{R} \to \varphi_{6}$}
\AXC{$(\mathrm{MonoDec}(a) \to M \in \mathbb{R} \to \varphi_{6}) \to \varphi_{254}$\lbl{A1}}
\BIC{$\varphi_{254}$}
\end{ptree}
\noindent 主木:
\begin{ptree}
\AXC{$\varphi_{55} \to \forall s\,(\varphi_{55})$\lbl{A4}}
\AXC{$\varphi_{187}$\lbl{A3}}
\AXC{$\forall \alpha\,(\neg \mathrm{Conv}(a,\alpha)) \to \neg \mathrm{Conv}(a,s)$\lbl{A2}}
\AXC{$\stackrel{\Pi_{55}}{\vdots}$}\noLine
\UIC{$\varphi_{137}$}
\BIC{$\mathrm{MonoDec}(a) \to M \in \mathbb{R} \to \varphi_{25}$}
\AXC{$\varphi_{182}$\lbl{A1}}
\BIC{$\varphi_{55} \to \neg \mathrm{IsInf}(s,\mathrm{ran}(a))$}
\UIC{$\forall s\,(\varphi_{55} \to \neg \mathrm{IsInf}(s,\mathrm{ran}(a)))$}
\AXC{$\stackrel{\Pi_{81}}{\vdots}$}\noLine
\UIC{$\varphi_{254}$}
\BIC{$\varphi_{187} \to (\varphi_{55} \to \forall s\,(\varphi_{55})) \to \varphi_{49}$}
\BIC{$(\varphi_{55} \to \forall s\,(\varphi_{55})) \to \varphi_{49}$}
\BIC{$\varphi_{49}$}
\UIC{$\forall M\,(\varphi_{49})$}
\UIC{$\forall a\,\forall M\,(\varphi_{49})$}
\end{ptree}
\begin{abbrevlist}
\item $\varphi_{1} :\equiv \forall b\,(\neg \mathrm{LB}(b,\mathrm{ran}(a))) \to \neg \mathrm{LB}(M,\mathrm{ran}(a))$
\item $\varphi_{2} :\equiv \mathrm{IsInf}(s,\mathrm{ran}(a)) \to \mathrm{IsInf}(s,\mathrm{ran}(a))$
\item $\varphi_{3} :\equiv \exists b\,(\mathrm{LB}(b,\mathrm{ran}(a))) \to \exists s\,(\mathrm{IsInf}(s,\mathrm{ran}(a)))$
\item $\varphi_{4} :\equiv \forall X\,(\mathrm{IsInf}(s,X) \leftrightarrow \mathrm{LB}(s,X) \land \forall w\,(\mathrm{LB}(w,X) \to w \leq s))$
\item $\varphi_{5} :\equiv \forall y\,(v_{720} = y \to y = v_{720}) \to v_{720} = v_{721} \to v_{721} = v_{720}$
\item $\varphi_{6} :\equiv \forall n \in \mathbb{N}\,(M \leq a_{n}) \to \exists s\,(\mathrm{IsInf}(s,\mathrm{ran}(a)))$
\item $\varphi_{7} :\equiv \mathrm{LB}(s,\mathrm{ran}(a)) \leftrightarrow s \in \mathbb{R} \land \forall x \in \mathrm{ran}(a)\,(s \leq x)$
\item $\varphi_{8} :\equiv \mathrm{LB}(M,\mathrm{ran}(a)) \leftrightarrow M \in \mathbb{R} \land \forall x \in \mathrm{ran}(a)\,(M \leq x)$
\item $\varphi_{9} :\equiv \mathrm{IsInf}(s,\mathrm{ran}(a)) \leftrightarrow \mathrm{LB}(s,\mathrm{ran}(a)) \land \forall w\,(\mathrm{LB}(w,\mathrm{ran}(a)) \to w \leq s)$
\item $\varphi_{10} :\equiv \forall y\,(\mathrm{MonoDec}(a) \land M \in \mathbb{R} \land \forall n \in \mathbb{N}\,(M \leq a_{n}))$
\item $\varphi_{11} :\equiv \forall i_{0}\,(\neg (i_{0} \in \mathbb{N} \land x = a_{i_{0}})) \to \neg (j \in \mathbb{N} \land x = a_{j})$
\item $\varphi_{12} :\equiv \forall j\,(\forall n\,(\neg (n \in \mathbb{N} \land y = a_{n})) \to \neg (j \in \mathbb{N} \land y = a_{j}))$
\item $\varphi_{13} :\equiv \forall n \in \mathbb{N}\,(M \leq a_{n}) \to \mathrm{IsInf}(s,\mathrm{ran}(a)) \to \mathrm{Conv}(a,s)$
\item $\varphi_{14} :\equiv y \in \mathrm{ran}(a) \to (y \in \mathrm{ran}(a) \leftrightarrow \exists j \in \mathbb{N}\,(y = a_{j}))$
\item $\varphi_{15} :\equiv \forall z\,\forall w\,(z = w \to M = z \to M = w) \to \forall w\,(a_{n} = w \to M = a_{n} \to M = w)$
\item $\varphi_{16} :\equiv \forall n \in \mathbb{N}\,(M \leq a_{n}) \to \mathrm{IsInf}(s,\mathrm{ran}(a)) \to s \in \mathbb{R}$
\item $\varphi_{17} :\equiv M \in \mathbb{R} \to \forall n \in \mathbb{N}\,(M \leq a_{n}) \to \exists b\,(\mathrm{LB}(b,\mathrm{ran}(a)))$
\item $\varphi_{18} :\equiv \mathrm{MonoDec}(a) \leftrightarrow \mathrm{Seq}(a) \land \forall n \in \mathbb{N}\,(a_{S(n)} \leq a_{n})$
\item $\varphi_{19} :\equiv \forall v_{720}\,\forall v_{721}\,(v_{720} = v_{721} \to v_{721} = v_{720}) \to \forall v_{721}\,(y = v_{721} \to v_{721} = y)$
\item $\varphi_{20} :\equiv \forall X\,(X \subseteq \mathbb{R} \to X \neq \emptyset \to \exists b\,(\mathrm{LB}(b,X)) \to \exists s\,(\mathrm{IsInf}(s,X)))$
\item $\varphi_{21} :\equiv \forall n\,(\neg (n \in \mathbb{N} \land y = a_{n})) \to \forall j\,\forall n\,(\neg (n \in \mathbb{N} \land y = a_{n}))$
\item $\varphi_{22} :\equiv \forall j\,\forall n\,(\neg (n \in \mathbb{N} \land y = a_{n})) \to \forall j\,(\neg (j \in \mathbb{N} \land y = a_{j}))$
\item $\varphi_{23} :\equiv y \in \mathrm{ran}(a) \to \forall y\,(y \in \mathrm{ran}(a) \leftrightarrow \exists j \in \mathbb{N}\,(y = a_{j}))$
\item $\varphi_{24} :\equiv \mathrm{MonoDec}(a) \to M \in \mathbb{R} \to \forall n \in \mathbb{N}\,(M \leq a_{n}) \to \mathrm{Seq}(a)$
\item $\varphi_{25} :\equiv \forall n \in \mathbb{N}\,(M \leq a_{n}) \to \mathrm{IsInf}(s,\mathrm{ran}(a)) \to \exists \alpha\,(\mathrm{Conv}(a,\alpha))$
\item $\varphi_{26} :\equiv a_{1_{\omega}} \in \mathrm{ran}(a) \leftrightarrow \exists j \in \mathbb{N}\,(a_{1_{\omega}} = a_{j})$
\item $\varphi_{27} :\equiv \forall j\,\forall i_{0}\,(\neg (i_{0} \in \mathbb{N} \land x = a_{i_{0}})) \to \forall j\,(\neg (j \in \mathbb{N} \land x = a_{j}))$
\item $\varphi_{28} :\equiv n \in \mathbb{N} \to N \leq n \to \forall y\,(y \in \mathrm{ran}(a) \leftrightarrow \exists j \in \mathbb{N}\,(y = a_{j}))$
\item $\varphi_{29} :\equiv \forall y \in \mathrm{ran}(a)\,(\neg M < y \to M = y) \to x \in \mathrm{ran}(a) \to \neg M < x \to M = x$
\item $\varphi_{30} :\equiv \forall x \in \mathrm{ran}(a)\,(\neg M < x \to M = x) \to y \in \mathrm{ran}(a) \to \neg M < y \to M = y$
\item $\varphi_{31} :\equiv \forall N\,(\neg (N \in \mathbb{N} \land \forall n \in \mathbb{N}\,(N \leq n \to \lvert a_{n} - s \rvert < \varepsilon)))$
\item $\varphi_{32} :\equiv \mathrm{ran}(a) \subseteq \mathbb{R} \land \forall y\,(y \in \mathrm{ran}(a) \leftrightarrow \exists j \in \mathbb{N}\,(y = a_{j}))$
\item $\varphi_{33} :\equiv M \in \mathbb{R} \to \forall n \in \mathbb{N}\,(M \leq a_{n}) \to \forall y \in \mathrm{ran}(a)\,(M \leq y)$
\item $\varphi_{34} :\equiv \forall n\,(\mathrm{MonoDec}(a) \to N \in \mathbb{N} \to n \in \mathbb{N} \to N \leq n \to a_{n} \leq a_{N})$
\item $\varphi_{35} :\equiv M \in \mathbb{R} \to \forall n \in \mathbb{N}\,(M \leq a_{n}) \to a_{1_{\omega}} \in \mathrm{ran}(a)$
\item $\varphi_{36} :\equiv \forall i_{0}\,(\neg (i_{0} \in \mathbb{N} \land x = a_{i_{0}})) \to \forall j\,\forall i_{0}\,(\neg (i_{0} \in \mathbb{N} \land x = a_{i_{0}}))$
\item $\varphi_{37} :\equiv \forall w\,(a_{n} = w \to M = a_{n} \to M = w) \to a_{n} = y \to M = a_{n} \to M = y$
\item $\varphi_{38} :\equiv \forall X\,\forall x \in X\,(X \neq \emptyset) \to \forall x \in \mathrm{ran}(a)\,(\mathrm{ran}(a) \neq \emptyset)$
\item $\varphi_{39} :\equiv \forall y \in \mathrm{ran}(a)\,(\neg M < y \to M = y) \to \forall x \in \mathrm{ran}(a)\,(\neg M < x \to M = x)$
\item $\varphi_{40} :\equiv \forall x \in \mathrm{ran}(a)\,(\neg M < x \to M = x) \to \forall y \in \mathrm{ran}(a)\,(\neg M < y \to M = y)$
\item $\varphi_{41} :\equiv \forall m\,\forall n\,(\mathrm{MonoDec}(a) \to m \in \mathbb{N} \to n \in \mathbb{N} \to m \leq n \to a_{n} \leq a_{m})$
\item $\varphi_{42} :\equiv \forall j\,(\neg (j \in \mathbb{N} \land a_{n} = a_{j})) \to \neg (n \in \mathbb{N} \land a_{n} = a_{n})$
\item $\varphi_{43} :\equiv \mathrm{MonoDec}(a) \to M \in \mathbb{R} \to \forall n \in \mathbb{N}\,(M \leq a_{n}) \to \mathrm{LB}(M,\mathrm{ran}(a))$
\item $\varphi_{44} :\equiv \forall n \in \mathbb{N}\,(M \leq a_{n}) \to 1_{\omega} \in \mathbb{N} \land a_{1_{\omega}} = a_{1_{\omega}}$
\item $\varphi_{45} :\equiv \forall n \in \mathbb{N}\,(M \leq a_{n}) \to \mathrm{IsInf}(s,\mathrm{ran}(a)) \to \forall x \in \mathrm{ran}(a)\,(s \leq x)$
\item $\varphi_{46} :\equiv \forall y \in \mathrm{ran}(a)\,(\neg M < y \to M = y) \to \forall x\,\forall y \in \mathrm{ran}(a)\,(\neg M < y \to M = y)$
\item $\varphi_{47} :\equiv \forall x \in \mathrm{ran}(a)\,(\neg M < x \to M = x) \to \forall y\,\forall x \in \mathrm{ran}(a)\,(\neg M < x \to M = x)$
\item $\varphi_{48} :\equiv \forall z\,\forall w\,(z = w \to M < z \to M < w) \to \forall w\,(a_{n} = w \to M < a_{n} \to M < w)$
\item $\varphi_{49} :\equiv \mathrm{MonoDec}(a) \to M \in \mathbb{R} \to \forall n \in \mathbb{N}\,(M \leq a_{n}) \to \exists \alpha\,(\mathrm{Conv}(a,\alpha))$
\item $\varphi_{50} :\equiv x \in \mathrm{ran}(a) \land x < s + \varepsilon \to i_{0} \in \mathbb{N} \land x = a_{i_{0}} \to i_{0} \in \mathbb{N}$
\item $\varphi_{51} :\equiv \forall x\,\forall y \in \mathrm{ran}(a)\,(\neg M < y \to M = y) \to \forall x \in \mathrm{ran}(a)\,(\neg M < x \to M = x)$
\item $\varphi_{52} :\equiv \forall y\,\forall x \in \mathrm{ran}(a)\,(\neg M < x \to M = x) \to \forall y \in \mathrm{ran}(a)\,(\neg M < y \to M = y)$
\item $\varphi_{53} :\equiv M \in \mathbb{R} \to \forall n \in \mathbb{N}\,(M \leq a_{n}) \to y \in \mathrm{ran}(a) \to y \in \mathrm{ran}(a)$
\item $\varphi_{54} :\equiv \forall n \in \mathbb{N}\,(M \leq a_{n}) \to y \in \mathrm{ran}(a) \to \exists j \in \mathbb{N}\,(y = a_{j})$
\item $\varphi_{55} :\equiv \mathrm{MonoDec}(a) \land M \in \mathbb{R} \land \forall n \in \mathbb{N}\,(M \leq a_{n}) \land \neg \exists \alpha\,(\mathrm{Conv}(a,\alpha))$
\item $\varphi_{56} :\equiv \mathrm{MonoDec}(a) \land M \in \mathbb{R} \land \forall n \in \mathbb{N}\,(M \leq a_{n}) \land \mathrm{IsInf}(s,\mathrm{ran}(a))$
\item $\varphi_{57} :\equiv x \in \mathrm{ran}(a) \land x < s + \varepsilon \to i_{0} \in \mathbb{N} \land x = a_{i_{0}} \to x = a_{i_{0}}$
\item $\varphi_{58} :\equiv \varepsilon \in \mathbb{R} \to 0 < \varepsilon \to x \in \mathrm{ran}(a) \land x < s + \varepsilon \to x \in \mathrm{ran}(a)$
\item $\varphi_{59} :\equiv \mathrm{MonoDec}(a) \to M \in \mathbb{R} \to \forall n \in \mathbb{N}\,(M \leq a_{n}) \to \mathrm{ran}(a) \subseteq \mathbb{R}$
\item $\varphi_{60} :\equiv \mathrm{MonoDec}(a) \to M \in \mathbb{R} \to \forall n \in \mathbb{N}\,(M \leq a_{n}) \to \mathrm{ran}(a) \neq \emptyset$
\item $\varphi_{61} :\equiv 0 < \varepsilon \to x \in \mathrm{ran}(a) \land x < s + \varepsilon \to x \in \mathrm{ran}(a) \land x < s + \varepsilon$
\item $\varphi_{62} :\equiv x \in \mathrm{ran}(a) \land x < s + \varepsilon \to i_{0} \in \mathbb{N} \land x = a_{i_{0}} \to x < s + \varepsilon$
\item $\varphi_{63} :\equiv x \in \mathrm{ran}(a) \land x < s + \varepsilon \to (x \in \mathrm{ran}(a) \leftrightarrow \exists j \in \mathbb{N}\,(x = a_{j}))$
\item $\varphi_{64} :\equiv \forall n \in \mathbb{N}\,(M \leq a_{n}) \to y \in \mathrm{ran}(a) \to n \in \mathbb{N} \land y = a_{n} \to M \leq y$
\item $\varphi_{65} :\equiv \forall n \in \mathbb{N}\,(M \leq a_{n}) \to \forall y\,(y \in \mathrm{ran}(a) \leftrightarrow \exists j \in \mathbb{N}\,(y = a_{j}))$
\item $\varphi_{66} :\equiv \mathrm{MonoDec}(a) \to M \in \mathbb{R} \to \forall n \in \mathbb{N}\,(M \leq a_{n}) \to y \in \mathrm{ran}(a) \to M \leq y$
\item $\varphi_{67} :\equiv \mathrm{MonoDec}(a) \land M \in \mathbb{R} \land \forall n \in \mathbb{N}\,(M \leq a_{n}) \to y \in \mathrm{ran}(a) \to M \leq y$
\item $\varphi_{68} :\equiv n \in \mathbb{N} \to N \leq n \to (a_{n} \in \mathrm{ran}(a) \leftrightarrow \exists j \in \mathbb{N}\,(a_{n} = a_{j}))$
\item $\varphi_{69} :\equiv \forall n \in \mathbb{N}\,(M \leq a_{n}) \to \mathrm{IsInf}(s,\mathrm{ran}(a)) \to \varepsilon \in \mathbb{R} \to \varepsilon \in \mathbb{R}$
\item $\varphi_{70} :\equiv \mathrm{MonoDec}(a) \land M \in \mathbb{R} \land \forall n \in \mathbb{N}\,(M \leq a_{n}) \land y \in \mathrm{ran}(a) \land \neg M \leq y$
\item $\varphi_{71} :\equiv \forall x\,\forall y\,\forall z\,(x \leq y \to y < z \to x < z) \to \forall y\,\forall z\,(a_{n} \leq y \to y < z \to a_{n} < z)$
\item $\varphi_{72} :\equiv \forall x\,\forall y\,\forall z\,(x < y \to y \leq z \to x < z) \to \forall y\,\forall z\,(s - \varepsilon < y \to y \leq z \to s - \varepsilon < z)$
\item $\varphi_{73} :\equiv \forall N\,(\neg (N \in \mathbb{N} \land a_{N} < s + \varepsilon)) \to \neg (i_{0} \in \mathbb{N} \land a_{i_{0}} < s + \varepsilon)$
\item $\varphi_{74} :\equiv x \in \mathrm{ran}(a) \land x < s + \varepsilon \to \forall y\,(y \in \mathrm{ran}(a) \leftrightarrow \exists j \in \mathbb{N}\,(y = a_{j}))$
\item $\varphi_{75} :\equiv \mathrm{IsInf}(s,\mathrm{ran}(a)) \to \varepsilon \in \mathbb{R} \to 0 < \varepsilon \to \exists x \in \mathrm{ran}(a)\,(x < s + \varepsilon)$
\item $\varphi_{76} :\equiv \forall w\,(a_{n} = w \to M < a_{n} \to M < w) \to a_{n} = y \to M < a_{n} \to M < y$
\item $\varphi_{77} :\equiv \mathrm{MonoDec}(a) \to M \in \mathbb{R} \to \forall n \in \mathbb{N}\,(M \leq a_{n}) \to n \in \mathbb{N} \to M \leq a_{n}$
\item $\varphi_{78} :\equiv M \in \mathbb{R} \to \forall n \in \mathbb{N}\,(M \leq a_{n}) \to \exists j \in \mathbb{N}\,(a_{1_{\omega}} = a_{j})$
\item $\varphi_{79} :\equiv \forall n \in \mathbb{N}\,(M \leq a_{n}) \to y \in \mathrm{ran}(a) \to n \in \mathbb{N} \land y = a_{n} \to n \in \mathbb{N}$
\item $\varphi_{80} :\equiv \forall j\,(\neg (j \in \mathbb{N} \land a_{1_{\omega}} = a_{j})) \to \neg (1_{\omega} \in \mathbb{N} \land a_{1_{\omega}} = a_{1_{\omega}})$
\item $\varphi_{81} :\equiv \forall x \in \mathrm{ran}(a)\,(\mathrm{ran}(a) \neq \emptyset) \to a_{1_{\omega}} \in \mathrm{ran}(a) \to \mathrm{ran}(a) \neq \emptyset$
\item $\varphi_{82} :\equiv \forall n \in \mathbb{N}\,(M \leq a_{n}) \to \mathrm{IsInf}(s,\mathrm{ran}(a)) \to a_{n} \in \mathrm{ran}(a) \to s \leq a_{n}$
\item $\varphi_{83} :\equiv \forall n \in \mathbb{N}\,(M \leq a_{n}) \to y \in \mathrm{ran}(a) \to n \in \mathbb{N} \land y = a_{n} \to M \leq a_{n}$
\item $\varphi_{84} :\equiv \forall n \in \mathbb{N}\,(M \leq a_{n}) \to y \in \mathrm{ran}(a) \to n \in \mathbb{N} \land y = a_{n} \to y = a_{n}$
\item $\varphi_{85} :\equiv \forall n \in \mathbb{N}\,(M \leq a_{n}) \to y \in \mathrm{ran}(a) \to n \in \mathbb{N} \land y = a_{n} \to a_{n} = y$
\item $\varphi_{86} :\equiv \forall E\,\forall e\,(E \subseteq \mathbb{R} \to \mathrm{IsInf}(s,E) \to e \in \mathbb{R} \to 0 < e \to \exists x \in E\,(x < s + e))$
\item $\varphi_{87} :\equiv \mathrm{MonoDec}(a) \to M \in \mathbb{R} \to \forall n \in \mathbb{N}\,(M \leq a_{n}) \to \forall n \in \mathbb{N}\,(M \leq a_{n})$
\item $\varphi_{88} :\equiv x \in \mathrm{ran}(a) \land x < s + \varepsilon \to i_{0} \in \mathbb{N} \land x = a_{i_{0}} \to a_{i_{0}} < s + \varepsilon$
\item $\varphi_{89} :\equiv \forall n \in \mathbb{N}\,(M \leq a_{n}) \to \mathrm{IsInf}(s,\mathrm{ran}(a)) \to \varepsilon \in \mathbb{R} \to 0 < \varepsilon \to 0 < \varepsilon$
\item $\varphi_{90} :\equiv N \in \mathbb{N} \land a_{N} < s + \varepsilon \to \forall n \in \mathbb{N}\,(N \leq n \to \lvert a_{n} - s \rvert < \varepsilon)$
\item $\varphi_{91} :\equiv \forall y\,\forall z\,(s - \varepsilon < y \to y \leq z \to s - \varepsilon < z) \to \forall z\,(s - \varepsilon < s \to s \leq z \to s - \varepsilon < z)$
\item $\varphi_{92} :\equiv \forall n\,(\mathrm{Seq}(a) \to n \in \mathbb{N} \to a_{n} \in \mathbb{R}) \to \mathrm{Seq}(a) \to n \in \mathbb{N} \to a_{n} \in \mathbb{R}$
\item $\varphi_{93} :\equiv \forall z\,\forall w\,(z = w \to z < s + \varepsilon \to w < s + \varepsilon) \to \forall w\,(x = w \to x < s + \varepsilon \to w < s + \varepsilon)$
\item $\varphi_{94} :\equiv N \in \mathbb{N} \land a_{N} < s + \varepsilon \to n \in \mathbb{N} \to N \leq n \to n \in \mathbb{N} \land a_{n} = a_{n}$
\item $\varphi_{95} :\equiv x \in \mathrm{ran}(a) \land x < s + \varepsilon \to i_{0} \in \mathbb{N} \land x = a_{i_{0}} \to i_{0} \in \mathbb{N} \land x = a_{i_{0}}$
\item $\varphi_{96} :\equiv \varepsilon \in \mathbb{R} \to 0 < \varepsilon \to x \in \mathrm{ran}(a) \land x < s + \varepsilon \to \exists j \in \mathbb{N}\,(x = a_{j})$
\item $\varphi_{97} :\equiv i_{0} \in \mathbb{N} \land x = a_{i_{0}} \to N \in \mathbb{N} \land a_{N} < s + \varepsilon \to n \in \mathbb{N} \to N \leq n \to N \leq n$
\item $\varphi_{98} :\equiv i_{0} \in \mathbb{N} \land x = a_{i_{0}} \to N \in \mathbb{N} \land a_{N} < s + \varepsilon \to n \in \mathbb{N} \to n \in \mathbb{N}$
\item $\varphi_{99} :\equiv N \in \mathbb{N} \land a_{N} < s + \varepsilon \to n \in \mathbb{N} \to N \leq n \to \exists j \in \mathbb{N}\,(a_{n} = a_{j})$
\item $\varphi_{100} :\equiv \forall y\,(y \in \mathrm{ran}(a) \leftrightarrow \exists j \in \mathbb{N}\,(y = a_{j})) \to (x \in \mathrm{ran}(a) \leftrightarrow \exists j \in \mathbb{N}\,(x = a_{j}))$
\item $\varphi_{101} :\equiv \forall y\,(y \in \mathrm{ran}(a) \leftrightarrow \exists j \in \mathbb{N}\,(y = a_{j})) \to (y \in \mathrm{ran}(a) \leftrightarrow \exists j \in \mathbb{N}\,(y = a_{j}))$
\item $\varphi_{102} :\equiv \varepsilon \in \mathbb{R} \to 0 < \varepsilon \to \exists N \in \mathbb{N}\,\forall n \in \mathbb{N}\,(N \leq n \to \lvert a_{n} - s \rvert < \varepsilon)$
\item $\varphi_{103} :\equiv \forall a\,\forall n\,(\mathrm{Seq}(a) \to n \in \mathbb{N} \to a_{n} \in \mathbb{R}) \to \forall n\,(\mathrm{Seq}(a) \to n \in \mathbb{N} \to a_{n} \in \mathbb{R})$
\item $\varphi_{104} :\equiv i_{0} \in \mathbb{N} \land x = a_{i_{0}} \to \exists N \in \mathbb{N}\,\forall n \in \mathbb{N}\,(N \leq n \to \lvert a_{n} - s \rvert < \varepsilon)$
\item $\varphi_{105} :\equiv \forall e\,\forall y\,(s \in \mathbb{R} \to e \in \mathbb{R} \to y \in \mathbb{R} \to s - e < y \to y < s + e \to \lvert y - s \rvert < e)$
\item $\varphi_{106} :\equiv i_{0} \in \mathbb{N} \land x = a_{i_{0}} \to N \in \mathbb{N} \land a_{N} < s + \varepsilon \to n \in \mathbb{N} \to N \leq n \to N \in \mathbb{N}$
\item $\varphi_{107} :\equiv \forall a\,\forall X\,(\mathrm{LB}(a,X) \leftrightarrow a \in \mathbb{R} \land \forall x \in X\,(a \leq x)) \to \forall X\,(\mathrm{LB}(s,X) \leftrightarrow s \in \mathbb{R} \land \forall x \in X\,(s \leq x))$
\item $\varphi_{108} :\equiv \forall z\,(s - \varepsilon < s \to s \leq z \to s - \varepsilon < z) \to s - \varepsilon < s \to s \leq a_{n} \to s - \varepsilon < a_{n}$
\item $\varphi_{109} :\equiv \forall a\,\forall X\,(\mathrm{LB}(a,X) \leftrightarrow a \in \mathbb{R} \land \forall x \in X\,(a \leq x)) \to \forall X\,(\mathrm{LB}(M,X) \leftrightarrow M \in \mathbb{R} \land \forall x \in X\,(M \leq x))$
\item $\varphi_{110} :\equiv i_{0} \in \mathbb{N} \land x = a_{i_{0}} \to N \in \mathbb{N} \land a_{N} < s + \varepsilon \to n \in \mathbb{N} \to N \leq n \to s \leq a_{n}$
\item $\varphi_{111} :\equiv \forall n \in \mathbb{N}\,(M \leq a_{n}) \to y \in \mathrm{ran}(a) \to n \in \mathbb{N} \land y = a_{n} \to n \in \mathbb{N} \land y = a_{n}$
\item $\varphi_{112} :\equiv \forall e\,(\mathrm{ran}(a) \subseteq \mathbb{R} \to \mathrm{IsInf}(s,\mathrm{ran}(a)) \to e \in \mathbb{R} \to 0 < e \to \exists x \in \mathrm{ran}(a)\,(x < s + e))$
\item $\varphi_{113} :\equiv \forall x\,\forall e\,\forall y\,(x \in \mathbb{R} \to e \in \mathbb{R} \to y \in \mathbb{R} \to x - e < y \to y < x + e \to \lvert y - x \rvert < e)$
\item $\varphi_{114} :\equiv \forall w\,(x = w \to x < s + \varepsilon \to w < s + \varepsilon) \to x = a_{i_{0}} \to x < s + \varepsilon \to a_{i_{0}} < s + \varepsilon$
\item $\varphi_{115} :\equiv \forall \varepsilon \in \mathbb{R}\,(0 < \varepsilon \to \exists N \in \mathbb{N}\,\forall n \in \mathbb{N}\,(N \leq n \to \lvert a_{n} - s \rvert < \varepsilon))$
\item $\varphi_{116} :\equiv x \in \mathrm{ran}(a) \land x < s + \varepsilon \to i_{0} \in \mathbb{N} \land x = a_{i_{0}} \to i_{0} \in \mathbb{N} \land a_{i_{0}} < s + \varepsilon$
\item $\varphi_{117} :\equiv x \in \mathrm{ran}(a) \land x < s + \varepsilon \to \exists N \in \mathbb{N}\,\forall n \in \mathbb{N}\,(N \leq n \to \lvert a_{n} - s \rvert < \varepsilon)$
\item $\varphi_{118} :\equiv i_{0} \in \mathbb{N} \land x = a_{i_{0}} \to N \in \mathbb{N} \land a_{N} < s + \varepsilon \to n \in \mathbb{N} \to N \leq n \to s - \varepsilon < s$
\item $\varphi_{119} :\equiv N \in \mathbb{N} \land a_{N} < s + \varepsilon \to N \in \mathbb{N} \land \forall n \in \mathbb{N}\,(N \leq n \to \lvert a_{n} - s \rvert < \varepsilon)$
\item $\varphi_{120} :\equiv x \in \mathrm{ran}(a) \land x < s + \varepsilon \to i_{0} \in \mathbb{N} \land x = a_{i_{0}} \to \exists N \in \mathbb{N}\,(a_{N} < s + \varepsilon)$
\item $\varphi_{121} :\equiv \forall y\,(s \in \mathbb{R} \to \varepsilon \in \mathbb{R} \to y \in \mathbb{R} \to s - \varepsilon < y \to y < s + \varepsilon \to \lvert y - s \rvert < \varepsilon)$
\item $\varphi_{122} :\equiv i_{0} \in \mathbb{N} \land x = a_{i_{0}} \to N \in \mathbb{N} \land a_{N} < s + \varepsilon \to N \in \mathbb{N} \land a_{N} < s + \varepsilon$
\item $\varphi_{123} :\equiv \forall y\,\forall z\,(a_{n} \leq y \to y < z \to a_{n} < z) \to \forall z\,(a_{n} \leq a_{N} \to a_{N} < z \to a_{n} < z)$
\item $\varphi_{124} :\equiv N \in \mathbb{N} \land a_{N} < s + \varepsilon \to \exists N \in \mathbb{N}\,\forall n \in \mathbb{N}\,(N \leq n \to \lvert a_{n} - s \rvert < \varepsilon)$
\item $\varphi_{125} :\equiv i_{0} \in \mathbb{N} \land x = a_{i_{0}} \to N \in \mathbb{N} \land a_{N} < s + \varepsilon \to n \in \mathbb{N} \to N \leq n \to a_{n} \leq a_{N}$
\item $\varphi_{126} :\equiv i_{0} \in \mathbb{N} \land x = a_{i_{0}} \to N \in \mathbb{N} \land a_{N} < s + \varepsilon \to n \in \mathbb{N} \to N \leq n \to a_{n} \in \mathbb{R}$
\item $\varphi_{127} :\equiv i_{0} \in \mathbb{N} \land x = a_{i_{0}} \to N \in \mathbb{N} \land a_{N} < s + \varepsilon \to n \in \mathbb{N} \to N \leq n \to a_{n} \in \mathrm{ran}(a)$
\item $\varphi_{128} :\equiv \forall y\,(y \in \mathrm{ran}(a) \leftrightarrow \exists j \in \mathbb{N}\,(y = a_{j})) \to (a_{n} \in \mathrm{ran}(a) \leftrightarrow \exists j \in \mathbb{N}\,(a_{n} = a_{j}))$
\item $\varphi_{129} :\equiv i_{0} \in \mathbb{N} \land x = a_{i_{0}} \to N \in \mathbb{N} \land a_{N} < s + \varepsilon \to n \in \mathbb{N} \to N \leq n \to a_{N} < s + \varepsilon$
\item $\varphi_{130} :\equiv i_{0} \in \mathbb{N} \land x = a_{i_{0}} \to N \in \mathbb{N} \land a_{N} < s + \varepsilon \to n \in \mathbb{N} \to N \leq n \to a_{n} < s + \varepsilon$
\item $\varphi_{131} :\equiv i_{0} \in \mathbb{N} \land x = a_{i_{0}} \to N \in \mathbb{N} \land a_{N} < s + \varepsilon \to n \in \mathbb{N} \to N \leq n \to s - \varepsilon < a_{n}$
\item $\varphi_{132} :\equiv x \in \mathrm{ran}(a) \land x < s + \varepsilon \to i_{0} \in \mathbb{N} \land x = a_{i_{0}} \to N \in \mathbb{N} \land a_{N} < s + \varepsilon \to N \in \mathbb{N}$
\item $\varphi_{133} :\equiv \forall x\,\forall e\,(x \in \mathbb{R} \to e \in \mathbb{R} \to 0 < e \to x - e < x) \to \forall e\,(s \in \mathbb{R} \to e \in \mathbb{R} \to 0 < e \to s - e < s)$
\item $\varphi_{134} :\equiv \forall e\,(s \in \mathbb{R} \to e \in \mathbb{R} \to 0 < e \to s - e < s) \to s \in \mathbb{R} \to \varepsilon \in \mathbb{R} \to 0 < \varepsilon \to s - \varepsilon < s$
\item $\varphi_{135} :\equiv i_{0} \in \mathbb{N} \land x = a_{i_{0}} \to N \in \mathbb{N} \land a_{N} < s + \varepsilon \to n \in \mathbb{N} \to N \leq n \to \lvert a_{n} - s \rvert < \varepsilon$
\item $\varphi_{136} :\equiv \varepsilon \in \mathbb{R} \to a_{n} \in \mathbb{R} \to s - \varepsilon < a_{n} \to a_{n} < s + \varepsilon \to \lvert a_{n} - s \rvert < \varepsilon$
\item $\varphi_{137} :\equiv (\forall \alpha\,(\neg \mathrm{Conv}(a,\alpha)) \to \neg \mathrm{Conv}(a,s)) \to \mathrm{MonoDec}(a) \to M \in \mathbb{R} \to \varphi_{25}$
\item $\varphi_{138} :\equiv \forall z\,(a_{n} \leq a_{N} \to a_{N} < z \to a_{n} < z) \to a_{n} \leq a_{N} \to a_{N} < s + \varepsilon \to a_{n} < s + \varepsilon$
\item $\varphi_{139} :\equiv a_{1_{\omega}} = a_{1_{\omega}} \to \mathrm{MonoDec}(a) \to M \in \mathbb{R} \to \varphi_{44}$
\item $\varphi_{140} :\equiv \varphi_{34} \to \mathrm{MonoDec}(a) \to N \in \mathbb{N} \to n \in \mathbb{N} \to N \leq n \to a_{n} \leq a_{N}$
\item $\varphi_{141} :\equiv M \in \mathbb{R} \to \forall n \in \mathbb{N}\,(M \leq a_{n}) \to \mathrm{IsInf}(s,\mathrm{ran}(a)) \to \varphi_{58}$
\item $\varphi_{142} :\equiv (a_{n} = y \to M = a_{n} \to M = y) \to \mathrm{MonoDec}(a) \to M \in \mathbb{R} \to \varphi_{64}$
\item $\varphi_{143} :\equiv \forall n \in \mathbb{N}\,(M \leq a_{n}) \to \mathrm{IsInf}(s,\mathrm{ran}(a)) \to \varepsilon \in \mathbb{R} \to \varphi_{61}$
\item $\varphi_{144} :\equiv \forall \alpha\,(\mathrm{Conv}(a,\alpha) \leftrightarrow \mathrm{Seq}(a) \land \alpha \in \mathbb{R} \land \forall \varepsilon \in \mathbb{R}\,(0 < \varepsilon \to \exists N \in \mathbb{N}\,\forall n \in \mathbb{N}\,(N \leq n \to \lvert a_{n} - \alpha \rvert < \varepsilon)))$
\item $\varphi_{145} :\equiv \forall n \in \mathbb{N}\,(M \leq a_{n}) \to \mathrm{IsInf}(s,\mathrm{ran}(a)) \to \varepsilon \in \mathbb{R} \to 0 < \varepsilon \to \varphi_{50}$
\item $\varphi_{146} :\equiv \forall n \in \mathbb{N}\,(M \leq a_{n}) \to \mathrm{IsInf}(s,\mathrm{ran}(a)) \to \varepsilon \in \mathbb{R} \to 0 < \varepsilon \to \varphi_{57}$
\item $\varphi_{147} :\equiv M \in \mathbb{R} \to \forall n \in \mathbb{N}\,(M \leq a_{n}) \to \mathrm{IsInf}(s,\mathrm{ran}(a)) \to \varphi_{96}$
\item $\varphi_{148} :\equiv \forall n \in \mathbb{N}\,(M \leq a_{n}) \to \mathrm{IsInf}(s,\mathrm{ran}(a)) \to \varepsilon \in \mathbb{R} \to 0 < \varepsilon \to \varphi_{62}$
\item $\varphi_{149} :\equiv \forall n \in \mathbb{N}\,(M \leq a_{n}) \to \mathrm{IsInf}(s,\mathrm{ran}(a)) \to \varepsilon \in \mathbb{R} \to 0 < \varepsilon \to \varphi_{63}$
\item $\varphi_{150} :\equiv (\mathrm{MonoDec}(a) \to \varphi_{17}) \to \mathrm{MonoDec}(a) \to M \in \mathbb{R} \to \varphi_{6}$
\item $\varphi_{151} :\equiv \forall n \in \mathbb{N}\,(M \leq a_{n}) \to \mathrm{IsInf}(s,\mathrm{ran}(a)) \to \varepsilon \in \mathbb{R} \to 0 < \varepsilon \to \varphi_{74}$
\item $\varphi_{152} :\equiv M \in \mathbb{R} \to \forall n \in \mathbb{N}\,(M \leq a_{n}) \to \mathrm{IsInf}(s,\mathrm{ran}(a)) \to \varphi_{102}$
\item $\varphi_{153} :\equiv 0 < \varepsilon \to x \in \mathrm{ran}(a) \land x < s + \varepsilon \to i_{0} \in \mathbb{N} \land x = a_{i_{0}} \to \varphi_{90}$
\item $\varphi_{154} :\equiv 0 < \varepsilon \to x \in \mathrm{ran}(a) \land x < s + \varepsilon \to i_{0} \in \mathbb{N} \land x = a_{i_{0}} \to \varphi_{94}$
\item $\varphi_{155} :\equiv \forall n \in \mathbb{N}\,(M \leq a_{n}) \to \mathrm{IsInf}(s,\mathrm{ran}(a)) \to \varepsilon \in \mathbb{R} \to 0 < \varepsilon \to \varphi_{88}$
\item $\varphi_{156} :\equiv M \in \mathbb{R} \to \forall n \in \mathbb{N}\,(M \leq a_{n}) \to \mathrm{IsInf}(s,\mathrm{ran}(a)) \to \varphi_{115}$
\item $\varphi_{157} :\equiv 0 < \varepsilon \to x \in \mathrm{ran}(a) \land x < s + \varepsilon \to i_{0} \in \mathbb{N} \land x = a_{i_{0}} \to \varphi_{99}$
\item $\varphi_{158} :\equiv x \in \mathrm{ran}(a) \land x < s + \varepsilon \to i_{0} \in \mathbb{N} \land x = a_{i_{0}} \to N \in \mathbb{N} \land a_{N} < s + \varepsilon \to \varphi_{28}$
\item $\varphi_{159} :\equiv \forall n \in \mathbb{N}\,(M \leq a_{n}) \to \mathrm{IsInf}(s,\mathrm{ran}(a)) \to \varepsilon \in \mathbb{R} \to 0 < \varepsilon \to \varphi_{95}$
\item $\varphi_{160} :\equiv \mathrm{IsInf}(s,\mathrm{ran}(a)) \to \varepsilon \in \mathbb{R} \to 0 < \varepsilon \to x \in \mathrm{ran}(a) \land x < s + \varepsilon \to \varphi_{97}$
\item $\varphi_{161} :\equiv \mathrm{IsInf}(s,\mathrm{ran}(a)) \to \varepsilon \in \mathbb{R} \to 0 < \varepsilon \to x \in \mathrm{ran}(a) \land x < s + \varepsilon \to \varphi_{98}$
\item $\varphi_{162} :\equiv \varphi_{56} \land \varepsilon \in \mathbb{R} \land 0 < \varepsilon \land \neg \exists N \in \mathbb{N}\,\forall n \in \mathbb{N}\,(N \leq n \to \lvert a_{n} - s \rvert < \varepsilon)$
\item $\varphi_{163} :\equiv 0 < \varepsilon \to x \in \mathrm{ran}(a) \land x < s + \varepsilon \to i_{0} \in \mathbb{N} \land x = a_{i_{0}} \to \varphi_{119}$
\item $\varphi_{164} :\equiv (\forall x \in \mathrm{ran}(a)\,(s \leq x) \to a_{n} \in \mathrm{ran}(a) \to s \leq a_{n}) \to \mathrm{MonoDec}(a) \to M \in \mathbb{R} \to \varphi_{82}$
\item $\varphi_{165} :\equiv \mathrm{IsInf}(s,\mathrm{ran}(a)) \to \varepsilon \in \mathbb{R} \to 0 < \varepsilon \to x \in \mathrm{ran}(a) \land x < s + \varepsilon \to \varphi_{104}$
\item $\varphi_{166} :\equiv \mathrm{IsInf}(s,\mathrm{ran}(a)) \to \varepsilon \in \mathbb{R} \to 0 < \varepsilon \to x \in \mathrm{ran}(a) \land x < s + \varepsilon \to \varphi_{106}$
\item $\varphi_{167} :\equiv \forall n \in \mathbb{N}\,(M \leq a_{n}) \to \mathrm{IsInf}(s,\mathrm{ran}(a)) \to \varepsilon \in \mathbb{R} \to 0 < \varepsilon \to \varphi_{116}$
\item $\varphi_{168} :\equiv \forall n \in \mathbb{N}\,(M \leq a_{n}) \to \mathrm{IsInf}(s,\mathrm{ran}(a)) \to \varepsilon \in \mathbb{R} \to 0 < \varepsilon \to \varphi_{117}$
\item $\varphi_{169} :\equiv x \in \mathrm{ran}(a) \land x < s + \varepsilon \to i_{0} \in \mathbb{N} \land x = a_{i_{0}} \to N \in \mathbb{N} \land a_{N} < s + \varepsilon \to \varphi_{68}$
\item $\varphi_{170} :\equiv \mathrm{IsInf}(s,\mathrm{ran}(a)) \to \varepsilon \in \mathbb{R} \to 0 < \varepsilon \to x \in \mathrm{ran}(a) \land x < s + \varepsilon \to \varphi_{110}$
\item $\varphi_{171} :\equiv 0 < \varepsilon \to x \in \mathrm{ran}(a) \land x < s + \varepsilon \to i_{0} \in \mathbb{N} \land x = a_{i_{0}} \to \varphi_{124}$
\item $\varphi_{172} :\equiv \forall n \in \mathbb{N}\,(M \leq a_{n}) \to \mathrm{IsInf}(s,\mathrm{ran}(a)) \to \varepsilon \in \mathbb{R} \to 0 < \varepsilon \to \varphi_{120}$
\item $\varphi_{173} :\equiv \mathrm{IsInf}(s,\mathrm{ran}(a)) \to \varepsilon \in \mathbb{R} \to 0 < \varepsilon \to x \in \mathrm{ran}(a) \land x < s + \varepsilon \to \varphi_{118}$
\item $\varphi_{174} :\equiv \mathrm{IsInf}(s,\mathrm{ran}(a)) \to \varepsilon \in \mathbb{R} \to 0 < \varepsilon \to x \in \mathrm{ran}(a) \land x < s + \varepsilon \to \varphi_{122}$
\item $\varphi_{175} :\equiv \mathrm{IsInf}(s,\mathrm{ran}(a)) \to \varepsilon \in \mathbb{R} \to 0 < \varepsilon \to x \in \mathrm{ran}(a) \land x < s + \varepsilon \to \varphi_{125}$
\item $\varphi_{176} :\equiv \mathrm{IsInf}(s,\mathrm{ran}(a)) \to \varepsilon \in \mathbb{R} \to 0 < \varepsilon \to x \in \mathrm{ran}(a) \land x < s + \varepsilon \to \varphi_{126}$
\item $\varphi_{177} :\equiv \mathrm{IsInf}(s,\mathrm{ran}(a)) \to \varepsilon \in \mathbb{R} \to 0 < \varepsilon \to x \in \mathrm{ran}(a) \land x < s + \varepsilon \to \varphi_{127}$
\item $\varphi_{178} :\equiv \forall n \in \mathbb{N}\,(M \leq a_{n}) \to \mathrm{IsInf}(s,\mathrm{ran}(a)) \to \varepsilon \in \mathbb{R} \to 0 < \varepsilon \to \varphi_{132}$
\item $\varphi_{179} :\equiv \mathrm{IsInf}(s,\mathrm{ran}(a)) \to \varepsilon \in \mathbb{R} \to 0 < \varepsilon \to x \in \mathrm{ran}(a) \land x < s + \varepsilon \to \varphi_{129}$
\item $\varphi_{180} :\equiv \mathrm{IsInf}(s,\mathrm{ran}(a)) \to \varepsilon \in \mathbb{R} \to 0 < \varepsilon \to x \in \mathrm{ran}(a) \land x < s + \varepsilon \to \varphi_{130}$
\item $\varphi_{181} :\equiv \mathrm{IsInf}(s,\mathrm{ran}(a)) \to \varepsilon \in \mathbb{R} \to 0 < \varepsilon \to x \in \mathrm{ran}(a) \land x < s + \varepsilon \to \varphi_{131}$
\item $\varphi_{182} :\equiv (\mathrm{MonoDec}(a) \to M \in \mathbb{R} \to \varphi_{25}) \to \varphi_{55} \to \neg \mathrm{IsInf}(s,\mathrm{ran}(a))$
\item $\varphi_{183} :\equiv (\mathrm{MonoDec}(a) \to \varphi_{53}) \to \mathrm{MonoDec}(a) \to M \in \mathbb{R} \to \varphi_{54}$
\item $\varphi_{184} :\equiv \varphi_{24} \to \mathrm{MonoDec}(a) \to M \in \mathbb{R} \to \forall n \in \mathbb{N}\,(M \leq a_{n}) \to \varphi_{32}$
\item $\varphi_{185} :\equiv \mathrm{IsInf}(s,\mathrm{ran}(a)) \to \varepsilon \in \mathbb{R} \to 0 < \varepsilon \to x \in \mathrm{ran}(a) \land x < s + \varepsilon \to \varphi_{135}$
\item $\varphi_{186} :\equiv (\mathrm{MonoDec}(a) \land M \in \mathbb{R} \land \forall n \in \mathbb{N}\,(M \leq a_{n}) \to \varphi_{10}) \to \mathrm{MonoDec}(a) \to \varphi_{33}$
\item $\varphi_{187} :\equiv \forall s\,(\varphi_{55} \to \neg \mathrm{IsInf}(s,\mathrm{ran}(a))) \to \forall s\,(\varphi_{55}) \to \forall s\,(\neg \mathrm{IsInf}(s,\mathrm{ran}(a)))$
\item $\varphi_{188} :\equiv \varphi_{56} \land \varepsilon \in \mathbb{R} \land 0 < \varepsilon \land (x \in \mathrm{ran}(a) \land x < s + \varepsilon) \land (i_{0} \in \mathbb{N} \land x = a_{i_{0}}) \land (N \in \mathbb{N} \land a_{N} < s + \varepsilon)$
\item $\varphi_{189} :\equiv (\mathrm{MonoDec}(a) \to M \in \mathbb{R} \to \forall n \in \mathbb{N}\,(M \leq a_{n}) \to \varphi_{32}) \to \varphi_{59}$
\item $\varphi_{190} :\equiv \varphi_{56} \land \varepsilon \in \mathbb{R} \land 0 < \varepsilon \land (x \in \mathrm{ran}(a) \land x < s + \varepsilon) \land \neg \exists N \in \mathbb{N}\,\forall n \in \mathbb{N}\,(N \leq n \to \lvert a_{n} - s \rvert < \varepsilon)$
\item $\varphi_{191} :\equiv (a_{1_{\omega}} \in \mathrm{ran}(a) \to \mathrm{ran}(a) \neq \emptyset) \to (\mathrm{MonoDec}(a) \to \varphi_{35}) \to \varphi_{60}$
\item $\varphi_{192} :\equiv \varphi_{101} \to \mathrm{MonoDec}(a) \to M \in \mathbb{R} \to \forall n \in \mathbb{N}\,(M \leq a_{n}) \to \varphi_{14}$
\item $\varphi_{193} :\equiv (\mathrm{MonoDec}(a) \to M \in \mathbb{R} \to \varphi_{64}) \to \varphi_{70} \to \neg (n \in \mathbb{N} \land y = a_{n})$
\item $\varphi_{194} :\equiv \forall n\,(\varphi_{70} \to \neg (n \in \mathbb{N} \land y = a_{n})) \to \forall n\,(\varphi_{70}) \to \forall n\,(\neg (n \in \mathbb{N} \land y = a_{n}))$
\item $\varphi_{195} :\equiv (\mathrm{MonoDec}(a) \to M \in \mathbb{R} \to \varphi_{79}) \to \mathrm{MonoDec}(a) \to M \in \mathbb{R} \to \varphi_{83}$
\item $\varphi_{196} :\equiv \forall n \in \mathbb{N}\,(M \leq a_{n}) \to \mathrm{IsInf}(s,\mathrm{ran}(a)) \to \varepsilon \in \mathbb{R} \to \varphi_{153}$
\item $\varphi_{197} :\equiv \forall n \in \mathbb{N}\,(M \leq a_{n}) \to \mathrm{IsInf}(s,\mathrm{ran}(a)) \to \varepsilon \in \mathbb{R} \to \varphi_{154}$
\item $\varphi_{198} :\equiv \forall n \in \mathbb{N}\,(M \leq a_{n}) \to \mathrm{IsInf}(s,\mathrm{ran}(a)) \to \varepsilon \in \mathbb{R} \to \varphi_{157}$
\item $\varphi_{199} :\equiv (\mathrm{MonoDec}(a) \to M \in \mathbb{R} \to \forall n \in \mathbb{N}\,(M \leq a_{n}) \to \varphi_{32}) \to \mathrm{MonoDec}(a) \to M \in \mathbb{R} \to \varphi_{65}$
\item $\varphi_{200} :\equiv \varphi_{56} \land \varepsilon \in \mathbb{R} \land 0 < \varepsilon \land (x \in \mathrm{ran}(a) \land x < s + \varepsilon) \land (i_{0} \in \mathbb{N} \land x = a_{i_{0}}) \land \neg \exists N \in \mathbb{N}\,\forall n \in \mathbb{N}\,(N \leq n \to \lvert a_{n} - s \rvert < \varepsilon)$
\item $\varphi_{201} :\equiv (\mathrm{MonoDec}(a) \to M \in \mathbb{R} \to \varphi_{111}) \to \mathrm{MonoDec}(a) \to M \in \mathbb{R} \to \varphi_{79}$
\item $\varphi_{202} :\equiv (\mathrm{MonoDec}(a) \to M \in \mathbb{R} \to \varphi_{111}) \to \mathrm{MonoDec}(a) \to M \in \mathbb{R} \to \varphi_{84}$
\item $\varphi_{203} :\equiv (\forall y\,(y \in \mathrm{ran}(a) \leftrightarrow \exists j \in \mathbb{N}\,(y = a_{j})) \to \varphi_{26}) \to \mathrm{MonoDec}(a) \to M \in \mathbb{R} \to \forall n \in \mathbb{N}\,(M \leq a_{n}) \to \varphi_{26}$
\item $\varphi_{204} :\equiv \forall n \in \mathbb{N}\,(M \leq a_{n}) \to \mathrm{IsInf}(s,\mathrm{ran}(a)) \to \varepsilon \in \mathbb{R} \to \varphi_{163}$
\item $\varphi_{205} :\equiv \forall n \in \mathbb{N}\,(M \leq a_{n}) \to \mathrm{IsInf}(s,\mathrm{ran}(a)) \to \varepsilon \in \mathbb{R} \to 0 < \varepsilon \to \varphi_{158}$
\item $\varphi_{206} :\equiv \forall n \in \mathbb{N}\,(M \leq a_{n}) \to \mathrm{IsInf}(s,\mathrm{ran}(a)) \to \varepsilon \in \mathbb{R} \to \varphi_{171}$
\item $\varphi_{207} :\equiv (x = a_{i_{0}} \to x < s + \varepsilon \to a_{i_{0}} < s + \varepsilon) \to \mathrm{MonoDec}(a) \to M \in \mathbb{R} \to \varphi_{155}$
\item $\varphi_{208} :\equiv (\mathrm{MonoDec}(a) \to M \in \mathbb{R} \to \forall n \in \mathbb{N}\,(M \leq a_{n}) \to \varphi_{32}) \to \mathrm{MonoDec}(a) \to M \in \mathbb{R} \to \forall n \in \mathbb{N}\,(M \leq a_{n}) \to \varphi_{23}$
\item $\varphi_{209} :\equiv (\mathrm{MonoDec}(a) \to M \in \mathbb{R} \to \varphi_{84}) \to (y = a_{n} \to a_{n} = y) \to \mathrm{MonoDec}(a) \to M \in \mathbb{R} \to \varphi_{85}$
\item $\varphi_{210} :\equiv \forall n \in \mathbb{N}\,(M \leq a_{n}) \to \mathrm{IsInf}(s,\mathrm{ran}(a)) \to \varepsilon \in \mathbb{R} \to 0 < \varepsilon \to \varphi_{169}$
\item $\varphi_{211} :\equiv (\mathrm{MonoDec}(a) \to M \in \mathbb{R} \to \varphi_{89}) \to \mathrm{MonoDec}(a) \to M \in \mathbb{R} \to \forall n \in \mathbb{N}\,(M \leq a_{n}) \to \varphi_{75}$
\item $\varphi_{212} :\equiv (\mathrm{MonoDec}(a) \to \varphi_{156}) \to \mathrm{MonoDec}(a) \to M \in \mathbb{R} \to \varphi_{13}$
\item $\varphi_{213} :\equiv (\mathrm{MonoDec}(a) \to M \in \mathbb{R} \to \varphi_{44}) \to \varphi_{80} \to \mathrm{MonoDec}(a) \to \varphi_{78}$
\item $\varphi_{214} :\equiv (\mathrm{MonoDec}(a) \to M \in \mathbb{R} \to \varphi_{83}) \to (a_{n} = y \to M < a_{n} \to M < y) \to \varphi_{142}$
\item $\varphi_{215} :\equiv (\mathrm{MonoDec}(a) \to M \in \mathbb{R} \to \forall n \in \mathbb{N}\,(M \leq a_{n}) \to \varphi_{2}) \to \varphi_{9} \to \varphi_{7} \to \mathrm{MonoDec}(a) \to M \in \mathbb{R} \to \varphi_{16}$
\item $\varphi_{216} :\equiv (\mathrm{MonoDec}(a) \to M \in \mathbb{R} \to \forall n \in \mathbb{N}\,(M \leq a_{n}) \to \varphi_{26}) \to (\mathrm{MonoDec}(a) \to \varphi_{78}) \to \mathrm{MonoDec}(a) \to \varphi_{35}$
\item $\varphi_{217} :\equiv (\mathrm{MonoDec}(a) \to M \in \mathbb{R} \to \forall n \in \mathbb{N}\,(M \leq a_{n}) \to \varphi_{14}) \to \varphi_{183}$
\item $\varphi_{218} :\equiv (\mathrm{MonoDec}(a) \to M \in \mathbb{R} \to \forall n \in \mathbb{N}\,(M \leq a_{n}) \to \varphi_{32}) \to \mathrm{MonoDec}(a) \to M \in \mathbb{R} \to \varphi_{151}$
\item $\varphi_{219} :\equiv (\mathrm{MonoDec}(a) \to M \in \mathbb{R} \to \forall n \in \mathbb{N}\,(M \leq a_{n}) \to \varphi_{2}) \to \varphi_{9} \to \varphi_{7} \to \mathrm{MonoDec}(a) \to M \in \mathbb{R} \to \varphi_{45}$
\item $\varphi_{220} :\equiv (\mathrm{MonoDec}(a) \to M \in \mathbb{R} \to \varphi_{143}) \to \mathrm{MonoDec}(a) \to \varphi_{141}$
\item $\varphi_{221} :\equiv (\mathrm{MonoDec}(a) \to M \in \mathbb{R} \to \forall n \in \mathbb{N}\,(M \leq a_{n}) \to \varphi_{23}) \to \varphi_{192}$
\item $\varphi_{222} :\equiv (\mathrm{MonoDec}(a) \to M \in \mathbb{R} \to M \in \mathbb{R}) \to (\mathrm{MonoDec}(a) \to \varphi_{33}) \to \varphi_{40} \to \varphi_{39} \to \varphi_{43}$
\item $\varphi_{223} :\equiv (\mathrm{MonoDec}(a) \to M \in \mathbb{R} \to \varphi_{143}) \to \mathrm{MonoDec}(a) \to M \in \mathbb{R} \to \varphi_{148}$
\item $\varphi_{224} :\equiv (\mathrm{MonoDec}(a) \to M \in \mathbb{R} \to \varphi_{89}) \to \mathrm{MonoDec}(a) \to M \in \mathbb{R} \to \forall n \in \mathbb{N}\,(M \leq a_{n}) \to \varphi_{173}$
\item $\varphi_{225} :\equiv (\mathrm{MonoDec}(a) \to M \in \mathbb{R} \to \varphi_{159}) \to \mathrm{MonoDec}(a) \to M \in \mathbb{R} \to \varphi_{145}$
\item $\varphi_{226} :\equiv (\mathrm{MonoDec}(a) \to M \in \mathbb{R} \to \varphi_{159}) \to \mathrm{MonoDec}(a) \to M \in \mathbb{R} \to \varphi_{146}$
\item $\varphi_{227} :\equiv (\mathrm{MonoDec}(a) \to M \in \mathbb{R} \to \varphi_{155}) \to \mathrm{MonoDec}(a) \to M \in \mathbb{R} \to \varphi_{167}$
\item $\varphi_{228} :\equiv (\mathrm{MonoDec}(a) \to M \in \mathbb{R} \to \forall n \in \mathbb{N}\,(M \leq a_{n}) \to \varphi_{32}) \to \mathrm{MonoDec}(a) \to M \in \mathbb{R} \to \varphi_{205}$
\item $\varphi_{229} :\equiv (\varphi_{31} \to \neg (N \in \mathbb{N} \land \forall n \in \mathbb{N}\,(N \leq n \to \lvert a_{n} - s \rvert < \varepsilon))) \to \mathrm{MonoDec}(a) \to M \in \mathbb{R} \to \varphi_{206}$
\item $\varphi_{230} :\equiv (\mathrm{MonoDec}(a) \to M \in \mathbb{R} \to \varphi_{168}) \to \varphi_{162} \to \neg (x \in \mathrm{ran}(a) \land x < s + \varepsilon)$
\item $\varphi_{231} :\equiv \forall x\,(\varphi_{162} \to \neg (x \in \mathrm{ran}(a) \land x < s + \varepsilon)) \to \forall x\,(\varphi_{162}) \to \forall x\,(\neg (x \in \mathrm{ran}(a) \land x < s + \varepsilon))$
\item $\varphi_{232} :\equiv \forall y\,(\varphi_{67}) \to (\forall y\,(\varphi_{67}) \to \varphi_{10} \to \forall y \in \mathrm{ran}(a)\,(M \leq y)) \to \varphi_{186}$
\item $\varphi_{233} :\equiv (\mathrm{MonoDec}(a) \to M \in \mathbb{R} \to \forall n \in \mathbb{N}\,(M \leq a_{n}) \to \varphi_{174}) \to \mathrm{MonoDec}(a) \to M \in \mathbb{R} \to \varphi_{178}$
\item $\varphi_{234} :\equiv (\mathrm{MonoDec}(a) \to M \in \mathbb{R} \to \forall n \in \mathbb{N}\,(M \leq a_{n}) \to \varphi_{2}) \to (\mathrm{MonoDec}(a) \to M \in \mathbb{R} \to \varphi_{69}) \to \varphi_{211}$
\item $\varphi_{235} :\equiv (\mathrm{MonoDec}(a) \to M \in \mathbb{R} \to \varphi_{151}) \to \varphi_{100} \to \mathrm{MonoDec}(a) \to M \in \mathbb{R} \to \varphi_{149}$
\item $\varphi_{236} :\equiv (\mathrm{MonoDec}(a) \to M \in \mathbb{R} \to \forall n \in \mathbb{N}\,(M \leq a_{n}) \to \varphi_{165}) \to \varphi_{190} \to \neg (i_{0} \in \mathbb{N} \land x = a_{i_{0}})$
\item $\varphi_{237} :\equiv (\mathrm{MonoDec}(a) \to M \in \mathbb{R} \to \forall n \in \mathbb{N}\,(M \leq a_{n}) \to \varphi_{160}) \to \mathrm{MonoDec}(a) \to M \in \mathbb{R} \to \forall n \in \mathbb{N}\,(M \leq a_{n}) \to \varphi_{175}$
\item $\varphi_{238} :\equiv (\mathrm{MonoDec}(a) \to M \in \mathbb{R} \to \forall n \in \mathbb{N}\,(M \leq a_{n}) \to \varphi_{161}) \to \mathrm{MonoDec}(a) \to M \in \mathbb{R} \to \forall n \in \mathbb{N}\,(M \leq a_{n}) \to \varphi_{176}$
\item $\varphi_{239} :\equiv \forall j\,(\varphi_{11}) \to (\forall j\,(\varphi_{11}) \to \varphi_{27}) \to \varphi_{36} \to \mathrm{MonoDec}(a) \to M \in \mathbb{R} \to \varphi_{168}$
\item $\varphi_{240} :\equiv (\mathrm{MonoDec}(a) \to M \in \mathbb{R} \to \forall n \in \mathbb{N}\,(M \leq a_{n}) \to \varphi_{174}) \to \mathrm{MonoDec}(a) \to M \in \mathbb{R} \to \forall n \in \mathbb{N}\,(M \leq a_{n}) \to \varphi_{166}$
\item $\varphi_{241} :\equiv (\mathrm{MonoDec}(a) \to M \in \mathbb{R} \to \forall n \in \mathbb{N}\,(M \leq a_{n}) \to \varphi_{177}) \to \mathrm{MonoDec}(a) \to M \in \mathbb{R} \to \forall n \in \mathbb{N}\,(M \leq a_{n}) \to \varphi_{170}$
\item $\varphi_{242} :\equiv (\mathrm{MonoDec}(a) \to M \in \mathbb{R} \to \forall n \in \mathbb{N}\,(M \leq a_{n}) \to \varphi_{170}) \to \mathrm{MonoDec}(a) \to M \in \mathbb{R} \to \forall n \in \mathbb{N}\,(M \leq a_{n}) \to \varphi_{181}$
\item $\varphi_{243} :\equiv \forall i_{0}\,(\varphi_{190} \to \neg (i_{0} \in \mathbb{N} \land x = a_{i_{0}})) \to \forall i_{0}\,(\varphi_{190}) \to \forall i_{0}\,(\neg (i_{0} \in \mathbb{N} \land x = a_{i_{0}}))$
\item $\varphi_{244} :\equiv (\mathrm{MonoDec}(a) \to M \in \mathbb{R} \to \forall n \in \mathbb{N}\,(M \leq a_{n}) \to \varphi_{174}) \to \mathrm{MonoDec}(a) \to M \in \mathbb{R} \to \forall n \in \mathbb{N}\,(M \leq a_{n}) \to \varphi_{179}$
\item $\varphi_{245} :\equiv (\mathrm{MonoDec}(a) \to M \in \mathbb{R} \to \forall n \in \mathbb{N}\,(M \leq a_{n}) \to \varphi_{179}) \to \mathrm{MonoDec}(a) \to M \in \mathbb{R} \to \forall n \in \mathbb{N}\,(M \leq a_{n}) \to \varphi_{180}$
\item $\varphi_{246} :\equiv (\mathrm{MonoDec}(a) \to M \in \mathbb{R} \to \varphi_{167}) \to \varphi_{73} \to \mathrm{MonoDec}(a) \to M \in \mathbb{R} \to \varphi_{172}$
\item $\varphi_{247} :\equiv (\mathrm{MonoDec}(a) \to M \in \mathbb{R} \to \forall n \in \mathbb{N}\,(M \leq a_{n}) \to \varphi_{180}) \to \mathrm{MonoDec}(a) \to M \in \mathbb{R} \to \forall n \in \mathbb{N}\,(M \leq a_{n}) \to \varphi_{185}$
\item $\varphi_{248} :\equiv (\mathrm{MonoDec}(a) \to M \in \mathbb{R} \to \varphi_{149}) \to (\mathrm{MonoDec}(a) \to \varphi_{141}) \to \mathrm{MonoDec}(a) \to \varphi_{147}$
\item $\varphi_{249} :\equiv (\mathrm{MonoDec}(a) \to M \in \mathbb{R} \to \forall n \in \mathbb{N}\,(M \leq a_{n}) \to \varphi_{185}) \to \varphi_{188} \to n \in \mathbb{N} \to N \leq n \to \lvert a_{n} - s \rvert < \varepsilon$
\item $\varphi_{250} :\equiv (\mathrm{MonoDec}(a) \to M \in \mathbb{R} \to \varphi_{198}) \to \mathrm{MonoDec}(a) \to M \in \mathbb{R} \to \forall n \in \mathbb{N}\,(M \leq a_{n}) \to \varphi_{177}$
\item $\varphi_{251} :\equiv (\mathrm{MonoDec}(a) \to M \in \mathbb{R} \to \forall n \in \mathbb{N}\,(M \leq a_{n}) \to \varphi_{161}) \to a_{n} = a_{n} \to \mathrm{MonoDec}(a) \to M \in \mathbb{R} \to \varphi_{197}$
\item $\varphi_{252} :\equiv \forall n\,(\varphi_{188} \to n \in \mathbb{N} \to N \leq n \to \lvert a_{n} - s \rvert < \varepsilon) \to \forall n\,(\varphi_{188}) \to \forall n \in \mathbb{N}\,(N \leq n \to \lvert a_{n} - s \rvert < \varepsilon)$
\item $\varphi_{253} :\equiv (\mathrm{MonoDec}(a) \to M \in \mathbb{R} \to \varphi_{196}) \to \mathrm{MonoDec}(a) \to M \in \mathbb{R} \to \varphi_{204}$
\item $\varphi_{254} :\equiv \forall s\,(\varphi_{55} \to \neg \mathrm{IsInf}(s,\mathrm{ran}(a))) \to \varphi_{187} \to (\varphi_{55} \to \forall s\,(\varphi_{55})) \to \varphi_{49}$
\item $\varphi_{255} :\equiv (\mathrm{MonoDec}(a) \to M \in \mathbb{R} \to \varphi_{206}) \to \varphi_{200} \to \neg (N \in \mathbb{N} \land a_{N} < s + \varepsilon)$
\item $\varphi_{256} :\equiv (\mathrm{MonoDec}(a) \to M \in \mathbb{R} \to \varphi_{16}) \to (\mathrm{MonoDec}(a) \to M \in \mathbb{R} \to \varphi_{69}) \to \varphi_{224}$
\item $\varphi_{257} :\equiv \forall N\,(\varphi_{200} \to \neg (N \in \mathbb{N} \land a_{N} < s + \varepsilon)) \to \forall N\,(\varphi_{200}) \to \forall N\,(\neg (N \in \mathbb{N} \land a_{N} < s + \varepsilon))$
\item $\varphi_{258} :\equiv (\mathrm{Conv}(a,s) \leftrightarrow \mathrm{Seq}(a) \land s \in \mathbb{R} \land \varphi_{115}) \to \varphi_{24} \to (\mathrm{MonoDec}(a) \to M \in \mathbb{R} \to \varphi_{16}) \to \varphi_{212}$
\item $\varphi_{259} :\equiv (\mathrm{MonoDec}(a) \to M \in \mathbb{R} \to \varphi_{197}) \to \varphi_{42} \to \mathrm{MonoDec}(a) \to M \in \mathbb{R} \to \varphi_{198}$
\item $\varphi_{260} :\equiv (\mathrm{MonoDec}(a) \to M \in \mathbb{R} \to \varphi_{146}) \to (\mathrm{MonoDec}(a) \to M \in \mathbb{R} \to \varphi_{148}) \to \varphi_{207}$
\item $\varphi_{261} :\equiv (\mathrm{MonoDec}(a) \to M \in \mathbb{R} \to \varphi_{205}) \to \varphi_{128} \to \mathrm{MonoDec}(a) \to M \in \mathbb{R} \to \varphi_{210}$
\item $\varphi_{262} :\equiv (\forall \varepsilon\,(\varphi_{56} \to \varphi_{102}) \to \forall \varepsilon\,(\varphi_{56}) \to \varphi_{115}) \to (\varphi_{56} \to \forall \varepsilon\,(\varphi_{56})) \to \mathrm{MonoDec}(a) \to \varphi_{156}$
\item $\varphi_{263} :\equiv (\mathrm{MonoDec}(a) \to M \in \mathbb{R} \to \forall n \in \mathbb{N}\,(M \leq a_{n}) \to \varphi_{161}) \to \varphi_{237}$
\item $\varphi_{264} :\equiv (\varphi_{200} \to \forall N\,(\varphi_{200})) \to \mathrm{MonoDec}(a) \to M \in \mathbb{R} \to \forall n \in \mathbb{N}\,(M \leq a_{n}) \to \varphi_{165}$
\item $\varphi_{265} :\equiv (\mathrm{MonoDec}(a) \to M \in \mathbb{R} \to \forall n \in \mathbb{N}\,(M \leq a_{n}) \to \varphi_{173}) \to \varphi_{242}$
\item $\varphi_{266} :\equiv (\mathrm{MonoDec}(a) \to M \in \mathbb{R} \to \forall n \in \mathbb{N}\,(M \leq a_{n}) \to \varphi_{175}) \to \varphi_{245}$
\item $\varphi_{267} :\equiv (\mathrm{MonoDec}(a) \to M \in \mathbb{R} \to \forall n \in \mathbb{N}\,(M \leq a_{n}) \to \varphi_{181}) \to \varphi_{247}$
\item $\varphi_{268} :\equiv \forall n\,(\varphi_{70} \to \neg (n \in \mathbb{N} \land y = a_{n})) \to \varphi_{194} \to (\varphi_{70} \to \forall n\,(\varphi_{70})) \to \varphi_{12} \to (\varphi_{12} \to \varphi_{22}) \to \varphi_{21} \to \varphi_{66}$
\item $\varphi_{269} :\equiv (\mathrm{MonoDec}(a) \to M \in \mathbb{R} \to \forall n \in \mathbb{N}\,(M \leq a_{n}) \to \varphi_{166}) \to \varphi_{263}$
\item $\varphi_{270} :\equiv (\mathrm{MonoDec}(a) \to M \in \mathbb{R} \to \forall n \in \mathbb{N}\,(M \leq a_{n}) \to \varphi_{176}) \to \varphi_{267}$
\item $\varphi_{271} :\equiv (\mathrm{MonoDec}(a) \to N \in \mathbb{N} \to n \in \mathbb{N} \to N \leq n \to a_{n} \leq a_{N}) \to (\mathrm{MonoDec}(a) \to \mathrm{MonoDec}(a)) \to \varphi_{269}$
\item $\varphi_{272} :\equiv \forall x\,(\varphi_{162} \to \neg (x \in \mathrm{ran}(a) \land x < s + \varepsilon)) \to \varphi_{231} \to (\varphi_{162} \to \forall x\,(\varphi_{162})) \to \mathrm{MonoDec}(a) \to \varphi_{152}$
\item $\varphi_{273} :\equiv (\mathrm{MonoDec}(a) \to M \in \mathbb{R} \to \varphi_{16}) \to (\mathrm{MonoDec}(a) \to M \in \mathbb{R} \to \varphi_{69}) \to \varphi_{270}$
\item $\varphi_{274} :\equiv (\mathrm{MonoDec}(a) \to M \in \mathbb{R} \to \forall n \in \mathbb{N}\,(M \leq a_{n}) \to \varphi_{75}) \to \varphi_{272}$
\item $\varphi_{275} :\equiv \forall n\,(\varphi_{188} \to n \in \mathbb{N} \to N \leq n \to \lvert a_{n} - s \rvert < \varepsilon) \to \varphi_{252} \to (\varphi_{188} \to \forall n\,(\varphi_{188})) \to \mathrm{MonoDec}(a) \to M \in \mathbb{R} \to \varphi_{196}$
\item $\varphi_{276} :\equiv \forall i_{0}\,(\varphi_{190} \to \neg (i_{0} \in \mathbb{N} \land x = a_{i_{0}})) \to \varphi_{243} \to (\varphi_{190} \to \forall i_{0}\,(\varphi_{190})) \to \varphi_{239}$
\item $\varphi_{277} :\equiv (\mathrm{MonoDec}(a) \to M \in \mathbb{R} \to \varphi_{172}) \to \forall N\,(\varphi_{200} \to \neg (N \in \mathbb{N} \land a_{N} < s + \varepsilon)) \to \varphi_{257} \to \varphi_{264}$
\end{abbrevlist}
\end{fullproof}

% 自動生成 (prover/build.js): NestedInt  ノード数 646
\begin{fullproof}{定理5.10}{$\forall a\,\forall b\,(\mathrm{MonoInc}(a) \to \mathrm{MonoDec}(b) \to \forall n \in \mathbb{N}\,(a_{n} \leq b_{n}) \to \exists x \in \mathbb{R}\,\forall n \in \mathbb{N}\,(a_{n} \leq x \land x \leq b_{n}))$}{646}\phantomsection\label{tree:NestedInt}
\noindent$\Pi_{1}$:
\begin{ptree}
\AXC{$\varphi_{146}$\lbl{A4}}
\AXC{$\forall x\,(\varphi_{139}) \to \varphi_{149}$\lbl{A3}}
\AXC{$\varphi_{139}$\lbl{A2}}
\UIC{$\forall x\,(\varphi_{139})$}
\AXC{$\forall x\,(\varphi_{139}) \to (\forall x\,(\varphi_{139}) \to \varphi_{149}) \to \varphi_{146} \to \varphi_{143}$\lbl{A1}}
\BIC{$(\forall x\,(\varphi_{139}) \to \varphi_{149}) \to \varphi_{146} \to \varphi_{143}$}
\BIC{$\varphi_{146} \to \varphi_{143}$}
\BIC{$\varphi_{143}$}
\end{ptree}
\noindent$\Pi_{2}$:
\begin{ptree}
\AXC{$\varphi_{147}$\lbl{A4}}
\AXC{$\forall y\,(\varphi_{140}) \to \varphi_{150}$\lbl{A3}}
\AXC{$\varphi_{140}$\lbl{A2}}
\UIC{$\forall y\,(\varphi_{140})$}
\AXC{$\forall y\,(\varphi_{140}) \to (\forall y\,(\varphi_{140}) \to \varphi_{150}) \to \varphi_{147} \to \varphi_{144}$\lbl{A1}}
\BIC{$(\forall y\,(\varphi_{140}) \to \varphi_{150}) \to \varphi_{147} \to \varphi_{144}$}
\BIC{$\varphi_{147} \to \varphi_{144}$}
\BIC{$\varphi_{144}$}
\end{ptree}
\noindent$\Pi_{3}$:
\begin{ptree}
\AXC{$\varphi_{2}$\lbl{A1}}
\AXC{$\varphi_{62}$\lbl{[C21]}}
\AXC{$\varphi_{62} \to \forall m\,(\varphi_{44})$\lbl{A2}}
\BIC{$\forall m\,(\varphi_{44})$}
\AXC{$\forall m\,(\varphi_{44}) \to \varphi_{44}$\lbl{A2}}
\BIC{$\varphi_{44}$}
\AXC{$\varphi_{187}$\lbl{A2}}
\BIC{$\mathrm{MonoDec}(b) \to m \in \mathbb{N} \to n \in \mathbb{N} \to m \leq n \to b_{n} \leq b_{m}$}
\AXC{$\varphi_{258}$\lbl{A1}}
\BIC{$\varphi_{2} \to \varphi_{255}$}
\BIC{$\varphi_{255}$}
\end{ptree}
\noindent$\Pi_{4}$:
\begin{ptree}
\AXC{$\varphi_{196}$\lbl{A1}}
\AXC{$\varphi_{196} \to \mathrm{MonoInc}(a) \to \mathrm{MonoDec}(b) \to \varphi_{151}$\lbl{A1}}
\BIC{$\mathrm{MonoInc}(a) \to \mathrm{MonoDec}(b) \to \varphi_{151}$}
\AXC{$\mathrm{MonoInc}(a) \to \varphi_{69}$\lbl{A1}}
\AXC{$\stackrel{\Pi_{3}}{\vdots}$}\noLine
\UIC{$\varphi_{255}$}
\BIC{$(\mathrm{MonoInc}(a) \to \mathrm{MonoDec}(b) \to \varphi_{151}) \to \varphi_{183} \to \varphi_{203}$}
\BIC{$\varphi_{183} \to \varphi_{203}$}
\end{ptree}
\noindent$\Pi_{5}$:
\begin{ptree}
\AXC{$\forall n \in \mathbb{N}\,(a_{n} \leq b_{n}) \to n \in \mathbb{N} \to a_{n} \leq b_{n}$\lbl{A2}}
\AXC{$\mathrm{MonoInc}(a) \to \varphi_{118}$\lbl{A1}}
\AXC{$(\mathrm{MonoInc}(a) \to \varphi_{118}) \to \varphi_{217}$\lbl{A1}}
\BIC{$\varphi_{217}$}
\BIC{$\mathrm{MonoInc}(a) \to \varphi_{107}$}
\AXC{$\varphi_{251}$\lbl{A1}}
\BIC{$(\mathrm{MonoInc}(a) \to \mathrm{MonoDec}(b) \to \varphi_{151}) \to \varphi_{204}$}
\end{ptree}
\noindent$\Pi_{6}$:
\begin{ptree}
\AXC{$\varphi_{196}$\lbl{A1}}
\AXC{$\varphi_{196} \to \mathrm{MonoInc}(a) \to \mathrm{MonoDec}(b) \to \varphi_{151}$\lbl{A1}}
\BIC{$\mathrm{MonoInc}(a) \to \mathrm{MonoDec}(b) \to \varphi_{151}$}
\AXC{$\stackrel{\Pi_{5}}{\vdots}$}\noLine
\UIC{$(\mathrm{MonoInc}(a) \to \mathrm{MonoDec}(b) \to \varphi_{151}) \to \varphi_{204}$}
\BIC{$\varphi_{204}$}
\end{ptree}
\noindent$\Pi_{7}$:
\begin{ptree}
\AXC{$\forall x\,\forall y\,\forall z\,(x \leq y \to y \leq z \to x \leq z)$\lbl{[O5]}}
\AXC{$\varphi_{41}$\lbl{A2}}
\BIC{$\forall y\,\forall z\,(a_{n} \leq y \to y \leq z \to a_{n} \leq z)$}
\AXC{$\varphi_{127}$\lbl{A2}}
\BIC{$\forall z\,(a_{n} \leq b_{n} \to b_{n} \leq z \to a_{n} \leq z)$}
\AXC{$\varphi_{165}$\lbl{A2}}
\BIC{$a_{n} \leq b_{n} \to b_{n} \leq b_{m} \to a_{n} \leq b_{m}$}
\end{ptree}
\noindent$\Pi_{8}$:
\begin{ptree}
\AXC{$\stackrel{\Pi_{6}}{\vdots}$}\noLine
\UIC{$\varphi_{204}$}
\AXC{$\stackrel{\Pi_{7}}{\vdots}$}\noLine
\UIC{$a_{n} \leq b_{n} \to b_{n} \leq b_{m} \to a_{n} \leq b_{m}$}
\AXC{$(a_{n} \leq b_{n} \to b_{n} \leq b_{m} \to a_{n} \leq b_{m}) \to \varphi_{204} \to \varphi_{203} \to \varphi_{205}$\lbl{A1}}
\BIC{$\varphi_{204} \to \varphi_{203} \to \varphi_{205}$}
\BIC{$\varphi_{203} \to \varphi_{205}$}
\end{ptree}
\noindent$\Pi_{9}$:
\begin{ptree}
\AXC{$\forall n \in \mathbb{N}\,(a_{n} \leq b_{n}) \to m \in \mathbb{N} \to a_{m} \leq b_{m}$\lbl{A2}}
\AXC{$\mathrm{MonoInc}(a) \to \varphi_{118}$\lbl{A1}}
\AXC{$(\mathrm{MonoInc}(a) \to \varphi_{118}) \to \varphi_{218}$\lbl{A1}}
\BIC{$\varphi_{218}$}
\BIC{$\mathrm{MonoInc}(a) \to \varphi_{108}$}
\AXC{$\varphi_{246}$\lbl{A1}}
\BIC{$(\mathrm{MonoInc}(a) \to \varphi_{69}) \to \varphi_{198}$}
\end{ptree}
\noindent$\Pi_{10}$:
\begin{ptree}
\AXC{$\varphi_{63}$\lbl{[C20]}}
\AXC{$\varphi_{63} \to \varphi_{92}$\lbl{A2}}
\BIC{$\varphi_{92}$}
\AXC{$\varphi_{92} \to \varphi_{121}$\lbl{A2}}
\BIC{$\varphi_{121}$}
\AXC{$\varphi_{121} \to \mathrm{MonoInc}(v_{724}) \to v_{725} \in \mathbb{N} \to \varphi_{26}$\lbl{A2}}
\BIC{$\mathrm{MonoInc}(v_{724}) \to v_{725} \in \mathbb{N} \to \varphi_{26}$}
\UIC{$\forall v_{726}\,(\mathrm{MonoInc}(v_{724}) \to v_{725} \in \mathbb{N} \to \varphi_{26})$}
\UIC{$\varphi_{159}$}
\UIC{$\forall v_{724}\,(\varphi_{159})$}
\AXC{$\forall v_{724}\,(\varphi_{159}) \to \varphi_{141}$\lbl{A2}}
\BIC{$\varphi_{141}$}
\end{ptree}
\noindent$\Pi_{11}$:
\begin{ptree}
\AXC{$\mathrm{MonoInc}(a) \to \mathrm{MonoInc}(a)$\lbl{A1}}
\AXC{$\stackrel{\Pi_{10}}{\vdots}$}\noLine
\UIC{$\varphi_{141}$}
\AXC{$\varphi_{141} \to \varphi_{93}$\lbl{A2}}
\BIC{$\varphi_{93}$}
\AXC{$\varphi_{202}$\lbl{A2}}
\BIC{$\mathrm{MonoInc}(a) \to n \in \mathbb{N} \to m \in \mathbb{N} \to n \leq m \to a_{n} \leq a_{m}$}
\AXC{$\varphi_{257}$\lbl{A1}}
\BIC{$(\mathrm{MonoInc}(a) \to \mathrm{MonoInc}(a)) \to \varphi_{254}$}
\BIC{$\varphi_{254}$}
\end{ptree}
\noindent$\Pi_{12}$:
\begin{ptree}
\AXC{$\varphi_{180}$\lbl{A1}}
\AXC{$\mathrm{MonoInc}(a) \to \varphi_{69}$\lbl{A1}}
\AXC{$\varphi_{196}$\lbl{A1}}
\AXC{$\varphi_{196} \to \mathrm{MonoInc}(a) \to \mathrm{MonoDec}(b) \to \varphi_{151}$\lbl{A1}}
\BIC{$\mathrm{MonoInc}(a) \to \mathrm{MonoDec}(b) \to \varphi_{151}$}
\AXC{$\stackrel{\Pi_{11}}{\vdots}$}\noLine
\UIC{$\varphi_{254}$}
\BIC{$(\mathrm{MonoInc}(a) \to \varphi_{69}) \to \varphi_{180} \to \varphi_{199}$}
\BIC{$\varphi_{180} \to \varphi_{199}$}
\BIC{$\varphi_{199}$}
\end{ptree}
\noindent$\Pi_{13}$:
\begin{ptree}
\AXC{$\forall x\,\forall y\,\forall z\,(x \leq y \to y \leq z \to x \leq z)$\lbl{[O5]}}
\AXC{$\varphi_{41}$\lbl{A2}}
\BIC{$\forall y\,\forall z\,(a_{n} \leq y \to y \leq z \to a_{n} \leq z)$}
\AXC{$\varphi_{128}$\lbl{A2}}
\BIC{$\forall z\,(a_{n} \leq a_{m} \to a_{m} \leq z \to a_{n} \leq z)$}
\AXC{$\varphi_{166}$\lbl{A2}}
\BIC{$a_{n} \leq a_{m} \to a_{m} \leq b_{m} \to a_{n} \leq b_{m}$}
\end{ptree}
\noindent$\Pi_{14}$:
\begin{ptree}
\AXC{$\stackrel{\Pi_{12}}{\vdots}$}\noLine
\UIC{$\varphi_{199}$}
\AXC{$\stackrel{\Pi_{13}}{\vdots}$}\noLine
\UIC{$a_{n} \leq a_{m} \to a_{m} \leq b_{m} \to a_{n} \leq b_{m}$}
\AXC{$(a_{n} \leq a_{m} \to a_{m} \leq b_{m} \to a_{n} \leq b_{m}) \to \varphi_{199} \to \varphi_{198} \to \varphi_{200}$\lbl{A1}}
\BIC{$\varphi_{199} \to \varphi_{198} \to \varphi_{200}$}
\BIC{$\varphi_{198} \to \varphi_{200}$}
\end{ptree}
\noindent$\Pi_{15}$:
\begin{ptree}
\AXC{$\mathrm{MonoInc}(a) \to \varphi_{69}$\lbl{A1}}
\AXC{$\stackrel{\Pi_{9}}{\vdots}$}\noLine
\UIC{$(\mathrm{MonoInc}(a) \to \varphi_{69}) \to \varphi_{198}$}
\BIC{$\varphi_{198}$}
\AXC{$\stackrel{\Pi_{14}}{\vdots}$}\noLine
\UIC{$\varphi_{198} \to \varphi_{200}$}
\BIC{$\varphi_{200}$}
\end{ptree}
\noindent$\Pi_{16}$:
\begin{ptree}
\AXC{$\forall n \in \mathbb{N}\,(n \in \omega)$\lbl{[N7']}}
\AXC{$\forall n \in \mathbb{N}\,(n \in \omega) \to m \in \mathbb{N} \to m \in \omega$\lbl{A2}}
\BIC{$m \in \mathbb{N} \to m \in \omega$}
\AXC{$(m \in \mathbb{N} \to m \in \omega) \to \varphi_{230}$\lbl{A1}}
\BIC{$\varphi_{230}$}
\end{ptree}
\noindent$\Pi_{17}$:
\begin{ptree}
\AXC{$\forall n \in \mathbb{N}\,(n \in \omega)$\lbl{[N7']}}
\AXC{$\forall n \in \mathbb{N}\,(n \in \omega) \to n \in \mathbb{N} \to n \in \omega$\lbl{A2}}
\BIC{$n \in \mathbb{N} \to n \in \omega$}
\AXC{$(n \in \mathbb{N} \to n \in \omega) \to \varphi_{236}$\lbl{A1}}
\BIC{$\varphi_{236}$}
\end{ptree}
\noindent$\Pi_{18}$:
\begin{ptree}
\AXC{$\varphi_{196}$\lbl{A1}}
\AXC{$\varphi_{196} \to \mathrm{MonoInc}(a) \to \mathrm{MonoDec}(b) \to \varphi_{151}$\lbl{A1}}
\BIC{$\mathrm{MonoInc}(a) \to \mathrm{MonoDec}(b) \to \varphi_{151}$}
\AXC{$\stackrel{\Pi_{17}}{\vdots}$}\noLine
\UIC{$\varphi_{236}$}
\BIC{$\mathrm{MonoInc}(a) \to \mathrm{MonoDec}(b) \to \varphi_{158}$}
\end{ptree}
\noindent$\Pi_{19}$:
\begin{ptree}
\AXC{$\forall m\,\forall n\,(m \in \omega \to n \in \omega \to m \leq n \lor n < m)$\lbl{[N11]}}
\AXC{$\varphi_{162}$\lbl{A2}}
\BIC{$\forall n\,(v_{722} \in \omega \to n \in \omega \to v_{722} \leq n \lor n < v_{722})$}
\AXC{$\varphi_{172}$\lbl{A2}}
\BIC{$v_{722} \in \omega \to v_{723} \in \omega \to v_{722} \leq v_{723} \lor v_{723} < v_{722}$}
\UIC{$\varphi_{47}$}
\UIC{$\forall v_{722}\,(\varphi_{47})$}
\AXC{$\varphi_{189}$\lbl{A2}}
\BIC{$\forall v_{723}\,(n \in \omega \to v_{723} \in \omega \to n \leq v_{723} \lor v_{723} < n)$}
\AXC{$\varphi_{155}$\lbl{A2}}
\BIC{$n \in \omega \to m \in \omega \to n \leq m \lor m < n$}
\end{ptree}
\noindent$\Pi_{20}$:
\begin{ptree}
\AXC{$\stackrel{\Pi_{18}}{\vdots}$}\noLine
\UIC{$\mathrm{MonoInc}(a) \to \mathrm{MonoDec}(b) \to \varphi_{158}$}
\AXC{$\stackrel{\Pi_{19}}{\vdots}$}\noLine
\UIC{$n \in \omega \to m \in \omega \to n \leq m \lor m < n$}
\AXC{$(n \in \omega \to m \in \omega \to n \leq m \lor m < n) \to \varphi_{252}$\lbl{A1}}
\BIC{$\varphi_{252}$}
\BIC{$(\mathrm{MonoInc}(a) \to \mathrm{MonoDec}(b) \to \varphi_{157}) \to \varphi_{184}$}
\end{ptree}
\noindent$\Pi_{21}$:
\begin{ptree}
\AXC{$\mathrm{MonoInc}(a) \to \varphi_{69}$\lbl{A1}}
\AXC{$\stackrel{\Pi_{16}}{\vdots}$}\noLine
\UIC{$\varphi_{230}$}
\BIC{$\mathrm{MonoInc}(a) \to \mathrm{MonoDec}(b) \to \varphi_{157}$}
\AXC{$\stackrel{\Pi_{20}}{\vdots}$}\noLine
\UIC{$(\mathrm{MonoInc}(a) \to \mathrm{MonoDec}(b) \to \varphi_{157}) \to \varphi_{184}$}
\BIC{$\varphi_{184}$}
\end{ptree}
\noindent$\Pi_{22}$:
\begin{ptree}
\AXC{$\stackrel{\Pi_{15}}{\vdots}$}\noLine
\UIC{$\varphi_{200}$}
\AXC{$\stackrel{\Pi_{21}}{\vdots}$}\noLine
\UIC{$\varphi_{184}$}
\AXC{$\varphi_{184} \to \varphi_{200} \to \varphi_{205} \to \mathrm{MonoInc}(a) \to \mathrm{MonoDec}(b) \to \varphi_{164}$\lbl{A1}}
\BIC{$\varphi_{200} \to \varphi_{205} \to \mathrm{MonoInc}(a) \to \mathrm{MonoDec}(b) \to \varphi_{164}$}
\BIC{$\varphi_{205} \to \mathrm{MonoInc}(a) \to \mathrm{MonoDec}(b) \to \varphi_{164}$}
\end{ptree}
\noindent$\Pi_{23}$:
\begin{ptree}
\AXC{$\varphi_{188}$\lbl{A1}}
\AXC{$\varphi_{188} \to \varphi_{183}$\lbl{A1}}
\BIC{$\varphi_{183}$}
\AXC{$\stackrel{\Pi_{4}}{\vdots}$}\noLine
\UIC{$\varphi_{183} \to \varphi_{203}$}
\BIC{$\varphi_{203}$}
\AXC{$\stackrel{\Pi_{8}}{\vdots}$}\noLine
\UIC{$\varphi_{203} \to \varphi_{205}$}
\BIC{$\varphi_{205}$}
\AXC{$\stackrel{\Pi_{22}}{\vdots}$}\noLine
\UIC{$\varphi_{205} \to \mathrm{MonoInc}(a) \to \mathrm{MonoDec}(b) \to \varphi_{164}$}
\BIC{$\mathrm{MonoInc}(a) \to \mathrm{MonoDec}(b) \to \varphi_{164}$}
\end{ptree}
\noindent$\Pi_{24}$:
\begin{ptree}
\AXC{$x = x$\lbl{A5}}
\AXC{$x = y \to x = x \to y = x$\lbl{A6}}
\AXC{$(x = y \to x = x \to y = x) \to x = x \to x = y \to y = x$\lbl{A1}}
\BIC{$x = x \to x = y \to y = x$}
\BIC{$x = y \to y = x$}
\UIC{$\forall y\,(x = y \to y = x)$}
\UIC{$\forall x\,\forall y\,(x = y \to y = x)$}
\end{ptree}
\noindent$\Pi_{25}$:
\begin{ptree}
\AXC{$\stackrel{\Pi_{24}}{\vdots}$}\noLine
\UIC{$\forall x\,\forall y\,(x = y \to y = x)$}
\AXC{$\forall x\,\forall y\,(x = y \to y = x) \to \forall y\,(v_{727} = y \to y = v_{727})$\lbl{A2}}
\BIC{$\forall y\,(v_{727} = y \to y = v_{727})$}
\AXC{$\varphi_{6}$\lbl{A2}}
\BIC{$v_{727} = v_{728} \to v_{728} = v_{727}$}
\UIC{$\forall v_{728}\,(v_{727} = v_{728} \to v_{728} = v_{727})$}
\UIC{$\forall v_{727}\,\forall v_{728}\,(v_{727} = v_{728} \to v_{728} = v_{727})$}
\AXC{$\varphi_{28}$\lbl{A2}}
\BIC{$\forall v_{728}\,(y = v_{728} \to v_{728} = y)$}
\end{ptree}
\noindent$\Pi_{26}$:
\begin{ptree}
\AXC{$\varphi_{196}$\lbl{A1}}
\AXC{$\varphi_{196} \to \mathrm{MonoInc}(a) \to \mathrm{MonoDec}(b) \to \varphi_{152}$\lbl{A1}}
\BIC{$\mathrm{MonoInc}(a) \to \mathrm{MonoDec}(b) \to \varphi_{152}$}
\AXC{$(\mathrm{MonoInc}(a) \to \mathrm{MonoDec}(b) \to \varphi_{152}) \to \varphi_{214}$\lbl{A1}}
\BIC{$\varphi_{214}$}
\end{ptree}
\noindent$\Pi_{27}$:
\begin{ptree}
\AXC{$\stackrel{\Pi_{25}}{\vdots}$}\noLine
\UIC{$\forall v_{728}\,(y = v_{728} \to v_{728} = y)$}
\AXC{$\forall v_{728}\,(y = v_{728} \to v_{728} = y) \to y = a_{n} \to a_{n} = y$\lbl{A2}}
\BIC{$y = a_{n} \to a_{n} = y$}
\AXC{$\stackrel{\Pi_{26}}{\vdots}$}\noLine
\UIC{$\varphi_{214}$}
\BIC{$\mathrm{MonoInc}(a) \to \mathrm{MonoDec}(b) \to \varphi_{153}$}
\end{ptree}
\noindent$\Pi_{28}$:
\begin{ptree}
\AXC{$\stackrel{\Pi_{27}}{\vdots}$}\noLine
\UIC{$\mathrm{MonoInc}(a) \to \mathrm{MonoDec}(b) \to \varphi_{153}$}
\AXC{$(\mathrm{MonoInc}(a) \to \mathrm{MonoDec}(b) \to \varphi_{153}) \to \varphi_{249}$\lbl{A1}}
\BIC{$\varphi_{249}$}
\end{ptree}
\noindent$\Pi_{29}$:
\begin{ptree}
\AXC{$z = w \to z < b_{m} \to w < b_{m}$\lbl{A6}}
\UIC{$\forall w\,(z = w \to z < b_{m} \to w < b_{m})$}
\UIC{$\forall z\,\forall w\,(z = w \to z < b_{m} \to w < b_{m})$}
\AXC{$\varphi_{148}$\lbl{A2}}
\BIC{$\forall w\,(a_{n} = w \to a_{n} < b_{m} \to w < b_{m})$}
\AXC{$\varphi_{163}$\lbl{A2}}
\BIC{$a_{n} = y \to a_{n} < b_{m} \to y < b_{m}$}
\AXC{$\stackrel{\Pi_{23}}{\vdots}$}\noLine
\UIC{$\mathrm{MonoInc}(a) \to \mathrm{MonoDec}(b) \to \varphi_{164}$}
\AXC{$\stackrel{\Pi_{28}}{\vdots}$}\noLine
\UIC{$\varphi_{249}$}
\BIC{$(a_{n} = y \to a_{n} < b_{m} \to y < b_{m}) \to \varphi_{221}$}
\BIC{$\varphi_{221}$}
\end{ptree}
\noindent$\Pi_{30}$:
\begin{ptree}
\AXC{$z = w \to z = b_{m} \to w = b_{m}$\lbl{A6}}
\UIC{$\forall w\,(z = w \to z = b_{m} \to w = b_{m})$}
\UIC{$\forall z\,\forall w\,(z = w \to z = b_{m} \to w = b_{m})$}
\AXC{$\varphi_{126}$\lbl{A2}}
\BIC{$\forall w\,(a_{n} = w \to a_{n} = b_{m} \to w = b_{m})$}
\AXC{$\varphi_{142}$\lbl{A2}}
\BIC{$a_{n} = y \to a_{n} = b_{m} \to y = b_{m}$}
\AXC{$\stackrel{\Pi_{29}}{\vdots}$}\noLine
\UIC{$\varphi_{221}$}
\BIC{$\mathrm{MonoInc}(a) \to \mathrm{MonoDec}(b) \to \varphi_{154}$}
\AXC{$(\mathrm{MonoInc}(a) \to \mathrm{MonoDec}(b) \to \varphi_{154}) \to \varphi_{185}$\lbl{A1}}
\BIC{$\varphi_{185}$}
\UIC{$\forall n\,(\varphi_{185})$}
\end{ptree}
\noindent$\Pi_{31}$:
\begin{ptree}
\AXC{$\mathrm{MonoInc}(a) \to \mathrm{MonoInc}(a)$\lbl{A1}}
\AXC{$(\mathrm{MonoInc}(a) \to \mathrm{MonoInc}(a)) \to \varphi_{29} \to \varphi_{58}$\lbl{A1}}
\BIC{$\varphi_{29} \to \varphi_{58}$}
\end{ptree}
\noindent$\Pi_{32}$:
\begin{ptree}
\AXC{$\forall a\,(\mathrm{Seq}(a) \to \varphi_{42})$\lbl{[S16]}}
\AXC{$\forall a\,(\mathrm{Seq}(a) \to \varphi_{42}) \to \mathrm{Seq}(a) \to \varphi_{42}$\lbl{A2}}
\BIC{$\mathrm{Seq}(a) \to \varphi_{42}$}
\AXC{$(\mathrm{Seq}(a) \to \varphi_{42}) \to \varphi_{58} \to \varphi_{181}$\lbl{A1}}
\BIC{$\varphi_{58} \to \varphi_{181}$}
\end{ptree}
\noindent$\Pi_{33}$:
\begin{ptree}
\AXC{$\forall a\,(\varphi_{29})$\lbl{DefMonoInc}}
\AXC{$\forall a\,(\varphi_{29}) \to \varphi_{29}$\lbl{A2}}
\BIC{$\varphi_{29}$}
\AXC{$\stackrel{\Pi_{31}}{\vdots}$}\noLine
\UIC{$\varphi_{29} \to \varphi_{58}$}
\BIC{$\varphi_{58}$}
\AXC{$\stackrel{\Pi_{32}}{\vdots}$}\noLine
\UIC{$\varphi_{58} \to \varphi_{181}$}
\BIC{$\varphi_{181}$}
\end{ptree}
\noindent$\Pi_{34}$:
\begin{ptree}
\AXC{$\mathrm{MonoInc}(a) \to \mathrm{MonoDec}(b) \to \varphi_{90}$\lbl{A1}}
\AXC{$\varphi_{134}$\lbl{A2}}
\AXC{$\stackrel{\Pi_{33}}{\vdots}$}\noLine
\UIC{$\varphi_{181}$}
\AXC{$\varphi_{181} \to \varphi_{192}$\lbl{A1}}
\BIC{$\varphi_{192}$}
\AXC{$\varphi_{192} \to \varphi_{134} \to \varphi_{190}$\lbl{A1}}
\BIC{$\varphi_{134} \to \varphi_{190}$}
\BIC{$\varphi_{190}$}
\AXC{$\varphi_{190} \to \varphi_{231}$\lbl{A1}}
\BIC{$\varphi_{231}$}
\BIC{$\mathrm{MonoInc}(a) \to \mathrm{MonoDec}(b) \to \varphi_{135}$}
\end{ptree}
\noindent$\Pi_{35}$:
\begin{ptree}
\AXC{$\varphi_{226}$\lbl{A4}}
\AXC{$\forall n\,(\varphi_{185}) \to \varphi_{193}$\lbl{A3}}
\AXC{$\stackrel{\Pi_{30}}{\vdots}$}\noLine
\UIC{$\forall n\,(\varphi_{185})$}
\AXC{$\stackrel{\Pi_{34}}{\vdots}$}\noLine
\UIC{$\mathrm{MonoInc}(a) \to \mathrm{MonoDec}(b) \to \varphi_{135}$}
\AXC{$\varphi_{270}$\lbl{A1}}
\BIC{$\forall n\,(\varphi_{185}) \to (\forall n\,(\varphi_{185}) \to \varphi_{193}) \to \varphi_{226} \to \varphi_{23} \to \varphi_{241}$}
\BIC{$(\forall n\,(\varphi_{185}) \to \varphi_{193}) \to \varphi_{226} \to \varphi_{23} \to \varphi_{241}$}
\BIC{$\varphi_{226} \to \varphi_{23} \to \varphi_{241}$}
\BIC{$\varphi_{23} \to \varphi_{241}$}
\end{ptree}
\noindent$\Pi_{36}$:
\begin{ptree}
\AXC{$\varphi_{33}$\lbl{A4}}
\AXC{$\varphi_{23} \to \varphi_{35}$\lbl{A3}}
\AXC{$\forall n\,(\neg (n \in \mathbb{N} \land y = a_{n})) \to \neg (j \in \mathbb{N} \land y = a_{j})$\lbl{A2}}
\UIC{$\varphi_{23}$}
\AXC{$\stackrel{\Pi_{35}}{\vdots}$}\noLine
\UIC{$\varphi_{23} \to \varphi_{241}$}
\BIC{$\varphi_{241}$}
\BIC{$\varphi_{33} \to \mathrm{MonoInc}(a) \to \mathrm{MonoDec}(b) \to \varphi_{74}$}
\BIC{$\mathrm{MonoInc}(a) \to \mathrm{MonoDec}(b) \to \varphi_{74}$}
\AXC{$\varphi_{225}$\lbl{A1}}
\BIC{$\varphi_{34} \land m \in \mathbb{N} \to y \in \mathrm{ran}(a) \to y \leq b_{m}$}
\UIC{$\forall y\,(\varphi_{34} \land m \in \mathbb{N} \to y \in \mathrm{ran}(a) \to y \leq b_{m})$}
\end{ptree}
\noindent$\Pi_{37}$:
\begin{ptree}
\AXC{$\varphi_{229}$\lbl{A3}}
\AXC{$\stackrel{\Pi_{36}}{\vdots}$}\noLine
\UIC{$\forall y\,(\varphi_{34} \land m \in \mathbb{N} \to y \in \mathrm{ran}(a) \to y \leq b_{m})$}
\AXC{$\forall y\,(\varphi_{34} \land m \in \mathbb{N} \to y \in \mathrm{ran}(a) \to y \leq b_{m}) \to \varphi_{229} \to \varphi_{237}$\lbl{A1}}
\BIC{$\varphi_{229} \to \varphi_{237}$}
\BIC{$\varphi_{237}$}
\end{ptree}
\noindent$\Pi_{38}$:
\begin{ptree}
\AXC{$\forall a\,\forall n\,(\mathrm{Seq}(a) \to n \in \mathbb{N} \to a_{n} \in \mathbb{R})$\lbl{[S1]}}
\AXC{$\varphi_{136}$\lbl{A2}}
\BIC{$\forall n\,(\mathrm{Seq}(b) \to n \in \mathbb{N} \to b_{n} \in \mathbb{R})$}
\AXC{$\varphi_{124}$\lbl{A2}}
\BIC{$\mathrm{Seq}(b) \to m \in \mathbb{N} \to b_{m} \in \mathbb{R}$}
\AXC{$\varphi_{243}$\lbl{A1}}
\BIC{$\varphi_{59} \to (\mathrm{MonoInc}(a) \to \varphi_{69}) \to \mathrm{MonoInc}(a) \to \varphi_{109}$}
\end{ptree}
\noindent$\Pi_{39}$:
\begin{ptree}
\AXC{$\varphi_{38}$\lbl{DefMonoDec}}
\AXC{$\varphi_{38} \to \varphi_{30}$\lbl{A2}}
\BIC{$\varphi_{30}$}
\AXC{$\varphi_{2}$\lbl{A1}}
\AXC{$\varphi_{2} \to \varphi_{30} \to \varphi_{59}$\lbl{A1}}
\BIC{$\varphi_{30} \to \varphi_{59}$}
\BIC{$\varphi_{59}$}
\AXC{$\stackrel{\Pi_{38}}{\vdots}$}\noLine
\UIC{$\varphi_{59} \to (\mathrm{MonoInc}(a) \to \varphi_{69}) \to \mathrm{MonoInc}(a) \to \varphi_{109}$}
\BIC{$(\mathrm{MonoInc}(a) \to \varphi_{69}) \to \mathrm{MonoInc}(a) \to \varphi_{109}$}
\end{ptree}
\noindent$\Pi_{40}$:
\begin{ptree}
\AXC{$\forall a\,\forall X\,(\mathrm{UB}(a,X) \leftrightarrow a \in \mathbb{R} \land \forall x \in X\,(x \leq a))$\lbl{DefUB}}
\AXC{$\varphi_{168}$\lbl{A2}}
\BIC{$\forall X\,(\mathrm{UB}(b_{m},X) \leftrightarrow b_{m} \in \mathbb{R} \land \forall x \in X\,(x \leq b_{m}))$}
\AXC{$\varphi_{201}$\lbl{A2}}
\BIC{$\varphi_{72}$}
\AXC{$\varphi_{72} \to (\mathrm{MonoInc}(a) \to \varphi_{109}) \to \varphi_{253}$\lbl{A1}}
\BIC{$(\mathrm{MonoInc}(a) \to \varphi_{109}) \to \varphi_{253}$}
\end{ptree}
\noindent$\Pi_{41}$:
\begin{ptree}
\AXC{$\mathrm{MonoInc}(a) \to \varphi_{69}$\lbl{A1}}
\AXC{$\stackrel{\Pi_{39}}{\vdots}$}\noLine
\UIC{$(\mathrm{MonoInc}(a) \to \varphi_{69}) \to \mathrm{MonoInc}(a) \to \varphi_{109}$}
\BIC{$\mathrm{MonoInc}(a) \to \varphi_{109}$}
\AXC{$\stackrel{\Pi_{40}}{\vdots}$}\noLine
\UIC{$(\mathrm{MonoInc}(a) \to \varphi_{109}) \to \varphi_{253}$}
\BIC{$\varphi_{253}$}
\end{ptree}
\noindent$\Pi_{42}$:
\begin{ptree}
\AXC{$\stackrel{\Pi_{1}}{\vdots}$}\noLine
\UIC{$\varphi_{143}$}
\AXC{$\stackrel{\Pi_{2}}{\vdots}$}\noLine
\UIC{$\varphi_{144}$}
\AXC{$\varphi_{34} \land m \in \mathbb{N} \to \forall y\,(\varphi_{34} \land m \in \mathbb{N})$\lbl{A4}}
\AXC{$\stackrel{\Pi_{37}}{\vdots}$}\noLine
\UIC{$\varphi_{237}$}
\BIC{$\mathrm{MonoInc}(a) \to \mathrm{MonoDec}(b) \to \varphi_{97}$}
\AXC{$\stackrel{\Pi_{41}}{\vdots}$}\noLine
\UIC{$\varphi_{253}$}
\BIC{$\varphi_{144} \to \varphi_{143} \to \mathrm{MonoInc}(a) \to \mathrm{MonoDec}(b) \to \varphi_{56}$}
\BIC{$\varphi_{143} \to \mathrm{MonoInc}(a) \to \mathrm{MonoDec}(b) \to \varphi_{56}$}
\BIC{$\mathrm{MonoInc}(a) \to \mathrm{MonoDec}(b) \to \varphi_{56}$}
\AXC{$\varphi_{223}$\lbl{A1}}
\BIC{$\varphi_{34} \to m \in \mathbb{N} \to \mathrm{UB}(b_{m},\mathrm{ran}(a))$}
\UIC{$\forall m\,(\varphi_{34} \to m \in \mathbb{N} \to \mathrm{UB}(b_{m},\mathrm{ran}(a)))$}
\end{ptree}
\noindent$\Pi_{43}$:
\begin{ptree}
\AXC{$\varphi_{224}$\lbl{A3}}
\AXC{$\stackrel{\Pi_{42}}{\vdots}$}\noLine
\UIC{$\forall m\,(\varphi_{34} \to m \in \mathbb{N} \to \mathrm{UB}(b_{m},\mathrm{ran}(a)))$}
\AXC{$\forall m\,(\varphi_{34} \to m \in \mathbb{N} \to \mathrm{UB}(b_{m},\mathrm{ran}(a))) \to \varphi_{224} \to \varphi_{233}$\lbl{A1}}
\BIC{$\varphi_{224} \to \varphi_{233}$}
\BIC{$\varphi_{233}$}
\end{ptree}
\noindent$\Pi_{44}$:
\begin{ptree}
\AXC{$\varphi_{34} \to \forall m\,(\varphi_{34})$\lbl{A4}}
\AXC{$\stackrel{\Pi_{43}}{\vdots}$}\noLine
\UIC{$\varphi_{233}$}
\BIC{$\mathrm{MonoInc}(a) \to \mathrm{MonoDec}(b) \to \varphi_{70}$}
\end{ptree}
\noindent$\Pi_{45}$:
\begin{ptree}
\AXC{$\varphi_{25}$\lbl{A2}}
\AXC{$\varphi_{11}$\lbl{DefSup}}
\AXC{$\varphi_{11} \to \varphi_{7}$\lbl{A2}}
\BIC{$\varphi_{7}$}
\AXC{$\varphi_{7} \to \varphi_{18}$\lbl{A2}}
\BIC{$\varphi_{18}$}
\AXC{$\varphi_{145}$\lbl{A1}}
\AXC{$\varphi_{145} \to \varphi_{18} \to \varphi_{160}$\lbl{A1}}
\BIC{$\varphi_{18} \to \varphi_{160}$}
\BIC{$\varphi_{160}$}
\AXC{$\varphi_{160} \to \varphi_{25} \to \varphi_{169}$\lbl{A1}}
\BIC{$\varphi_{25} \to \varphi_{169}$}
\BIC{$\varphi_{169}$}
\end{ptree}
\noindent$\Pi_{46}$:
\begin{ptree}
\AXC{$x = x$\lbl{A5}}
\UIC{$\forall x\,(x = x)$}
\AXC{$\forall x\,(x = x) \to a_{n} = a_{n}$\lbl{A2}}
\BIC{$a_{n} = a_{n}$}
\AXC{$\mathrm{MonoInc}(a) \to \mathrm{MonoDec}(b) \to \varphi_{88}$\lbl{A1}}
\AXC{$\varphi_{235}$\lbl{A1}}
\BIC{$a_{n} = a_{n} \to \varphi_{177}$}
\BIC{$\varphi_{177}$}
\end{ptree}
\noindent$\Pi_{47}$:
\begin{ptree}
\AXC{$\varphi_{161}$\lbl{A2}}
\AXC{$\stackrel{\Pi_{33}}{\vdots}$}\noLine
\UIC{$\varphi_{181}$}
\AXC{$\varphi_{181} \to \varphi_{195}$\lbl{A1}}
\BIC{$\varphi_{195}$}
\AXC{$\varphi_{195} \to \varphi_{161} \to \varphi_{207}$\lbl{A1}}
\BIC{$\varphi_{161} \to \varphi_{207}$}
\BIC{$\varphi_{207}$}
\AXC{$\varphi_{207} \to \varphi_{182} \to \mathrm{MonoInc}(a) \to \mathrm{MonoDec}(b) \to \varphi_{132}$\lbl{A1}}
\BIC{$\varphi_{182} \to \mathrm{MonoInc}(a) \to \mathrm{MonoDec}(b) \to \varphi_{132}$}
\end{ptree}
\noindent$\Pi_{48}$:
\begin{ptree}
\AXC{$\varphi_{64}$\lbl{A4}}
\AXC{$\forall u_{2}\,(\varphi_{53}) \to \varphi_{83}$\lbl{A3}}
\AXC{$\varphi_{53}$\lbl{A2}}
\UIC{$\forall u_{2}\,(\varphi_{53})$}
\AXC{$\forall u_{2}\,(\varphi_{53}) \to (\forall u_{2}\,(\varphi_{53}) \to \varphi_{83}) \to \varphi_{64} \to \varphi_{65}$\lbl{A1}}
\BIC{$(\forall u_{2}\,(\varphi_{53}) \to \varphi_{83}) \to \varphi_{64} \to \varphi_{65}$}
\BIC{$\varphi_{64} \to \varphi_{65}$}
\BIC{$\varphi_{65}$}
\end{ptree}
\noindent$\Pi_{49}$:
\begin{ptree}
\AXC{$\varphi_{98}$\lbl{A4}}
\AXC{$\forall y\,(\varphi_{55}) \to \varphi_{75}$\lbl{A3}}
\AXC{$\varphi_{55}$\lbl{A2}}
\UIC{$\forall y\,(\varphi_{55})$}
\AXC{$\forall y\,(\varphi_{55}) \to (\forall y\,(\varphi_{55}) \to \varphi_{75}) \to \varphi_{98} \to \varphi_{66}$\lbl{A1}}
\BIC{$(\forall y\,(\varphi_{55}) \to \varphi_{75}) \to \varphi_{98} \to \varphi_{66}$}
\BIC{$\varphi_{98} \to \varphi_{66}$}
\BIC{$\varphi_{66}$}
\end{ptree}
\noindent$\Pi_{50}$:
\begin{ptree}
\AXC{$\varphi_{129}$\lbl{A4}}
\AXC{$\forall x\,(\varphi_{85}) \to \varphi_{119}$\lbl{A3}}
\AXC{$\varphi_{85}$\lbl{A2}}
\UIC{$\forall x\,(\varphi_{85})$}
\AXC{$\forall x\,(\varphi_{85}) \to (\forall x\,(\varphi_{85}) \to \varphi_{119}) \to \varphi_{129} \to \varphi_{99}$\lbl{A1}}
\BIC{$(\forall x\,(\varphi_{85}) \to \varphi_{119}) \to \varphi_{129} \to \varphi_{99}$}
\BIC{$\varphi_{129} \to \varphi_{99}$}
\BIC{$\varphi_{99}$}
\end{ptree}
\noindent$\Pi_{51}$:
\begin{ptree}
\AXC{$\varphi_{100}$\lbl{A4}}
\AXC{$\forall u_{2}\,(\varphi_{78}) \to \varphi_{122}$\lbl{A3}}
\AXC{$\varphi_{78}$\lbl{A2}}
\UIC{$\forall u_{2}\,(\varphi_{78})$}
\AXC{$\forall u_{2}\,(\varphi_{78}) \to (\forall u_{2}\,(\varphi_{78}) \to \varphi_{122}) \to \varphi_{100} \to \varphi_{101}$\lbl{A1}}
\BIC{$(\forall u_{2}\,(\varphi_{78}) \to \varphi_{122}) \to \varphi_{100} \to \varphi_{101}$}
\BIC{$\varphi_{100} \to \varphi_{101}$}
\BIC{$\varphi_{101}$}
\end{ptree}
\noindent$\Pi_{52}$:
\begin{ptree}
\AXC{$\stackrel{\Pi_{50}}{\vdots}$}\noLine
\UIC{$\varphi_{99}$}
\AXC{$\stackrel{\Pi_{51}}{\vdots}$}\noLine
\UIC{$\varphi_{101}$}
\AXC{$\varphi_{79}$\lbl{A1}}
\AXC{$\varphi_{79} \to \varphi_{101} \to \varphi_{99} \to \varphi_{102}$\lbl{A1}}
\BIC{$\varphi_{101} \to \varphi_{99} \to \varphi_{102}$}
\BIC{$\varphi_{99} \to \varphi_{102}$}
\BIC{$\varphi_{102}$}
\AXC{$\varphi_{102} \to \varphi_{116}$\lbl{A1}}
\BIC{$\varphi_{116}$}
\end{ptree}
\noindent$\Pi_{53}$:
\begin{ptree}
\AXC{$\stackrel{\Pi_{52}}{\vdots}$}\noLine
\UIC{$\varphi_{116}$}
\AXC{$v_{731} \in \mathbb{R} \leftrightarrow v_{731} \in \mathbb{R}$\lbl{A1}}
\AXC{$(v_{731} \in \mathbb{R} \leftrightarrow v_{731} \in \mathbb{R}) \to \varphi_{116} \to \varphi_{170}$\lbl{A1}}
\BIC{$\varphi_{116} \to \varphi_{170}$}
\BIC{$\varphi_{170}$}
\AXC{$\varphi_{170} \to \varphi_{173}$\lbl{A1}}
\BIC{$\varphi_{173}$}
\end{ptree}
\noindent$\Pi_{54}$:
\begin{ptree}
\AXC{$\mathrm{UB}(v_{731},v_{732}) \leftrightarrow \mathrm{UB}(v_{731},v_{732})$\lbl{A1}}
\AXC{$\stackrel{\Pi_{53}}{\vdots}$}\noLine
\UIC{$\varphi_{173}$}
\AXC{$\varphi_{173} \to (\mathrm{UB}(v_{731},v_{732}) \leftrightarrow \mathrm{UB}(v_{731},v_{732})) \to \varphi_{208}$\lbl{A1}}
\BIC{$(\mathrm{UB}(v_{731},v_{732}) \leftrightarrow \mathrm{UB}(v_{731},v_{732})) \to \varphi_{208}$}
\BIC{$\varphi_{208}$}
\end{ptree}
\noindent$\Pi_{55}$:
\begin{ptree}
\AXC{$\varphi_{129}$\lbl{A4}}
\AXC{$\forall x\,(\varphi_{85}) \to \varphi_{119}$\lbl{A3}}
\AXC{$\varphi_{85}$\lbl{A2}}
\UIC{$\forall x\,(\varphi_{85})$}
\AXC{$\forall x\,(\varphi_{85}) \to (\forall x\,(\varphi_{85}) \to \varphi_{119}) \to \varphi_{129} \to \varphi_{99}$\lbl{A1}}
\BIC{$(\forall x\,(\varphi_{85}) \to \varphi_{119}) \to \varphi_{129} \to \varphi_{99}$}
\BIC{$\varphi_{129} \to \varphi_{99}$}
\BIC{$\varphi_{99}$}
\end{ptree}
\noindent$\Pi_{56}$:
\begin{ptree}
\AXC{$\varphi_{100}$\lbl{A4}}
\AXC{$\forall u_{2}\,(\varphi_{78}) \to \varphi_{122}$\lbl{A3}}
\AXC{$\varphi_{78}$\lbl{A2}}
\UIC{$\forall u_{2}\,(\varphi_{78})$}
\AXC{$\forall u_{2}\,(\varphi_{78}) \to (\forall u_{2}\,(\varphi_{78}) \to \varphi_{122}) \to \varphi_{100} \to \varphi_{101}$\lbl{A1}}
\BIC{$(\forall u_{2}\,(\varphi_{78}) \to \varphi_{122}) \to \varphi_{100} \to \varphi_{101}$}
\BIC{$\varphi_{100} \to \varphi_{101}$}
\BIC{$\varphi_{101}$}
\end{ptree}
\noindent$\Pi_{57}$:
\begin{ptree}
\AXC{$\stackrel{\Pi_{55}}{\vdots}$}\noLine
\UIC{$\varphi_{99}$}
\AXC{$\stackrel{\Pi_{56}}{\vdots}$}\noLine
\UIC{$\varphi_{101}$}
\AXC{$\varphi_{79}$\lbl{A1}}
\AXC{$\varphi_{79} \to \varphi_{101} \to \varphi_{99} \to \varphi_{102}$\lbl{A1}}
\BIC{$\varphi_{101} \to \varphi_{99} \to \varphi_{102}$}
\BIC{$\varphi_{99} \to \varphi_{102}$}
\BIC{$\varphi_{102}$}
\AXC{$\varphi_{102} \to \varphi_{116}$\lbl{A1}}
\BIC{$\varphi_{116}$}
\end{ptree}
\noindent$\Pi_{58}$:
\begin{ptree}
\AXC{$\stackrel{\Pi_{57}}{\vdots}$}\noLine
\UIC{$\varphi_{116}$}
\AXC{$v_{731} \in \mathbb{R} \leftrightarrow v_{731} \in \mathbb{R}$\lbl{A1}}
\AXC{$(v_{731} \in \mathbb{R} \leftrightarrow v_{731} \in \mathbb{R}) \to \varphi_{116} \to \varphi_{170}$\lbl{A1}}
\BIC{$\varphi_{116} \to \varphi_{170}$}
\BIC{$\varphi_{170}$}
\AXC{$\varphi_{170} \to \varphi_{173}$\lbl{A1}}
\BIC{$\varphi_{173}$}
\end{ptree}
\noindent$\Pi_{59}$:
\begin{ptree}
\AXC{$\stackrel{\Pi_{58}}{\vdots}$}\noLine
\UIC{$\varphi_{173}$}
\AXC{$\mathrm{UB}(v_{731},v_{732}) \leftrightarrow \mathrm{UB}(v_{731},v_{732})$\lbl{A1}}
\AXC{$(\mathrm{UB}(v_{731},v_{732}) \leftrightarrow \mathrm{UB}(v_{731},v_{732})) \to \varphi_{173} \to \varphi_{209}$\lbl{A1}}
\BIC{$\varphi_{173} \to \varphi_{209}$}
\BIC{$\varphi_{209}$}
\end{ptree}
\noindent$\Pi_{60}$:
\begin{ptree}
\AXC{$\stackrel{\Pi_{54}}{\vdots}$}\noLine
\UIC{$\varphi_{208}$}
\AXC{$\varphi_{208} \to \varphi_{212}$\lbl{A1}}
\BIC{$\varphi_{212}$}
\AXC{$\stackrel{\Pi_{59}}{\vdots}$}\noLine
\UIC{$\varphi_{209}$}
\AXC{$\varphi_{209} \to \varphi_{212} \to \varphi_{244}$\lbl{A1}}
\BIC{$\varphi_{212} \to \varphi_{244}$}
\BIC{$\varphi_{244}$}
\AXC{$\varphi_{244} \to (\varphi_{19} \leftrightarrow \varphi_{215})$\lbl{A1}}
\BIC{$\varphi_{19} \leftrightarrow \varphi_{215}$}
\end{ptree}
\noindent$\Pi_{61}$:
\begin{ptree}
\AXC{$\stackrel{\Pi_{60}}{\vdots}$}\noLine
\UIC{$\varphi_{19} \leftrightarrow \varphi_{215}$}
\AXC{$(\varphi_{19} \leftrightarrow \varphi_{215}) \to \varphi_{215} \to \varphi_{19}$\lbl{A1}}
\BIC{$\varphi_{215} \to \varphi_{19}$}
\UIC{$\forall v_{732}\,(\varphi_{215} \to \varphi_{19})$}
\AXC{$\stackrel{\Pi_{60}}{\vdots}$}\noLine
\UIC{$\varphi_{19} \leftrightarrow \varphi_{215}$}
\AXC{$(\varphi_{19} \leftrightarrow \varphi_{215}) \to \varphi_{19} \to \varphi_{215}$\lbl{A1}}
\BIC{$\varphi_{19} \to \varphi_{215}$}
\UIC{$\forall v_{732}\,(\varphi_{19} \to \varphi_{215})$}
\AXC{$\forall v_{732}\,(\varphi_{19} \to \varphi_{215}) \to \varphi_{271}$\lbl{A1}}
\BIC{$\varphi_{271}$}
\BIC{$(\forall v_{732}\,(\varphi_{19} \to \varphi_{215}) \to \forall v_{732}\,(\varphi_{19}) \to \forall v_{732}\,(\varphi_{215})) \to \varphi_{267}$}
\end{ptree}
\noindent$\Pi_{62}$:
\begin{ptree}
\AXC{$\forall v_{732}\,(\varphi_{19} \to \varphi_{215}) \to \forall v_{732}\,(\varphi_{19}) \to \forall v_{732}\,(\varphi_{215})$\lbl{A3}}
\AXC{$\stackrel{\Pi_{61}}{\vdots}$}\noLine
\UIC{$(\forall v_{732}\,(\varphi_{19} \to \varphi_{215}) \to \forall v_{732}\,(\varphi_{19}) \to \forall v_{732}\,(\varphi_{215})) \to \varphi_{267}$}
\BIC{$\varphi_{267}$}
\end{ptree}
\noindent$\Pi_{63}$:
\begin{ptree}
\AXC{$\forall v_{732}\,(\varphi_{215} \to \varphi_{19}) \to \forall v_{732}\,(\varphi_{215}) \to \forall v_{732}\,(\varphi_{19})$\lbl{A3}}
\AXC{$\stackrel{\Pi_{62}}{\vdots}$}\noLine
\UIC{$\varphi_{267}$}
\BIC{$\forall v_{732}\,(\varphi_{19}) \leftrightarrow \forall v_{732}\,(\varphi_{215})$}
\end{ptree}
\noindent$\Pi_{64}$:
\begin{ptree}
\AXC{$\stackrel{\Pi_{63}}{\vdots}$}\noLine
\UIC{$\forall v_{732}\,(\varphi_{19}) \leftrightarrow \forall v_{732}\,(\varphi_{215})$}
\AXC{$\varphi_{260}$\lbl{A1}}
\BIC{$\forall v_{732}\,(\varphi_{19}) \to \forall v_{732}\,(\varphi_{215})$}
\UIC{$\forall v_{731}\,(\forall v_{732}\,(\varphi_{19}) \to \forall v_{732}\,(\varphi_{215}))$}
\AXC{$\forall v_{731}\,(\forall v_{732}\,(\varphi_{19}) \to \forall v_{732}\,(\varphi_{215})) \to \varphi_{273}$\lbl{A1}}
\BIC{$\varphi_{273}$}
\end{ptree}
\noindent$\Pi_{65}$:
\begin{ptree}
\AXC{$\varphi_{263}$\lbl{A3}}
\AXC{$\varphi_{264}$\lbl{A3}}
\AXC{$\stackrel{\Pi_{63}}{\vdots}$}\noLine
\UIC{$\forall v_{732}\,(\varphi_{19}) \leftrightarrow \forall v_{732}\,(\varphi_{215})$}
\AXC{$\varphi_{259}$\lbl{A1}}
\BIC{$\forall v_{732}\,(\varphi_{215}) \to \forall v_{732}\,(\varphi_{19})$}
\UIC{$\forall v_{731}\,(\forall v_{732}\,(\varphi_{215}) \to \forall v_{732}\,(\varphi_{19}))$}
\AXC{$\stackrel{\Pi_{64}}{\vdots}$}\noLine
\UIC{$\varphi_{273}$}
\BIC{$\varphi_{264} \to \varphi_{263} \to (\forall v_{731}\,\forall v_{732}\,(\varphi_{19}) \leftrightarrow \varphi_{67})$}
\BIC{$\varphi_{263} \to (\forall v_{731}\,\forall v_{732}\,(\varphi_{19}) \leftrightarrow \varphi_{67})$}
\BIC{$\forall v_{731}\,\forall v_{732}\,(\varphi_{19}) \leftrightarrow \varphi_{67}$}
\end{ptree}
\noindent$\Pi_{66}$:
\begin{ptree}
\AXC{$\forall a\,\forall X\,(\mathrm{UB}(a,X) \leftrightarrow a \in \mathbb{R} \land \forall x \in X\,(x \leq a))$\lbl{DefUB}}
\AXC{$\forall a\,\forall X\,(\mathrm{UB}(a,X) \leftrightarrow a \in \mathbb{R} \land \forall x \in X\,(x \leq a)) \to \varphi_{13}$\lbl{A2}}
\BIC{$\varphi_{13}$}
\end{ptree}
\noindent$\Pi_{67}$:
\begin{ptree}
\AXC{$\stackrel{\Pi_{65}}{\vdots}$}\noLine
\UIC{$\forall v_{731}\,\forall v_{732}\,(\varphi_{19}) \leftrightarrow \varphi_{67}$}
\AXC{$\stackrel{\Pi_{66}}{\vdots}$}\noLine
\UIC{$\varphi_{13}$}
\AXC{$\varphi_{13} \to \varphi_{19}$\lbl{A2}}
\BIC{$\varphi_{19}$}
\UIC{$\forall v_{732}\,(\varphi_{19})$}
\UIC{$\forall v_{731}\,\forall v_{732}\,(\varphi_{19})$}
\AXC{$\forall v_{731}\,\forall v_{732}\,(\varphi_{19}) \to (\forall v_{731}\,\forall v_{732}\,(\varphi_{19}) \leftrightarrow \varphi_{67}) \to \varphi_{67}$\lbl{A1}}
\BIC{$(\forall v_{731}\,\forall v_{732}\,(\varphi_{19}) \leftrightarrow \varphi_{67}) \to \varphi_{67}$}
\BIC{$\varphi_{67}$}
\end{ptree}
\noindent$\Pi_{68}$:
\begin{ptree}
\AXC{$\stackrel{\Pi_{67}}{\vdots}$}\noLine
\UIC{$\varphi_{67}$}
\AXC{$\varphi_{67} \to \varphi_{20}$\lbl{A2}}
\BIC{$\varphi_{20}$}
\AXC{$\varphi_{20} \to \varphi_{15}$\lbl{A2}}
\BIC{$\varphi_{15}$}
\AXC{$\varphi_{11}$\lbl{DefSup}}
\AXC{$\varphi_{11} \to \varphi_{7}$\lbl{A2}}
\BIC{$\varphi_{7}$}
\AXC{$\varphi_{7} \to \varphi_{18}$\lbl{A2}}
\BIC{$\varphi_{18}$}
\AXC{$\varphi_{145}$\lbl{A1}}
\AXC{$\varphi_{145} \to \varphi_{250}$\lbl{A1}}
\BIC{$\varphi_{250}$}
\BIC{$\varphi_{15} \to \varphi_{66} \to \varphi_{65} \to \mathrm{MonoInc}(a) \to \mathrm{MonoDec}(b) \to \varphi_{91}$}
\BIC{$\varphi_{66} \to \varphi_{65} \to \mathrm{MonoInc}(a) \to \mathrm{MonoDec}(b) \to \varphi_{91}$}
\end{ptree}
\noindent$\Pi_{69}$:
\begin{ptree}
\AXC{$\stackrel{\Pi_{48}}{\vdots}$}\noLine
\UIC{$\varphi_{65}$}
\AXC{$\stackrel{\Pi_{49}}{\vdots}$}\noLine
\UIC{$\varphi_{66}$}
\AXC{$\stackrel{\Pi_{68}}{\vdots}$}\noLine
\UIC{$\varphi_{66} \to \varphi_{65} \to \mathrm{MonoInc}(a) \to \mathrm{MonoDec}(b) \to \varphi_{91}$}
\BIC{$\varphi_{65} \to \mathrm{MonoInc}(a) \to \mathrm{MonoDec}(b) \to \varphi_{91}$}
\BIC{$\mathrm{MonoInc}(a) \to \mathrm{MonoDec}(b) \to \varphi_{91}$}
\end{ptree}
\noindent$\Pi_{70}$:
\begin{ptree}
\AXC{$\stackrel{\Pi_{69}}{\vdots}$}\noLine
\UIC{$\mathrm{MonoInc}(a) \to \mathrm{MonoDec}(b) \to \varphi_{91}$}
\AXC{$(\mathrm{MonoInc}(a) \to \mathrm{MonoDec}(b) \to \varphi_{91}) \to \varphi_{222}$\lbl{A1}}
\BIC{$\varphi_{222}$}
\end{ptree}
\noindent$\Pi_{71}$:
\begin{ptree}
\AXC{$\forall y \in \mathrm{ran}(a)\,(y \leq x) \to a_{n} \in \mathrm{ran}(a) \to a_{n} \leq x$\lbl{A2}}
\AXC{$\stackrel{\Pi_{70}}{\vdots}$}\noLine
\UIC{$\varphi_{222}$}
\BIC{$\mathrm{MonoInc}(a) \to \mathrm{MonoDec}(b) \to \varphi_{133}$}
\AXC{$(\mathrm{MonoInc}(a) \to \mathrm{MonoDec}(b) \to \varphi_{133}) \to \varphi_{256}$\lbl{A1}}
\BIC{$\varphi_{256}$}
\end{ptree}
\noindent$\Pi_{72}$:
\begin{ptree}
\AXC{$\stackrel{\Pi_{45}}{\vdots}$}\noLine
\UIC{$\varphi_{169}$}
\AXC{$\varphi_{54}$\lbl{A2}}
\AXC{$\stackrel{\Pi_{46}}{\vdots}$}\noLine
\UIC{$\varphi_{177}$}
\AXC{$\varphi_{177} \to \varphi_{54} \to \varphi_{182}$\lbl{A1}}
\BIC{$\varphi_{54} \to \varphi_{182}$}
\BIC{$\varphi_{182}$}
\AXC{$\stackrel{\Pi_{47}}{\vdots}$}\noLine
\UIC{$\varphi_{182} \to \mathrm{MonoInc}(a) \to \mathrm{MonoDec}(b) \to \varphi_{132}$}
\BIC{$\mathrm{MonoInc}(a) \to \mathrm{MonoDec}(b) \to \varphi_{132}$}
\AXC{$\stackrel{\Pi_{71}}{\vdots}$}\noLine
\UIC{$\varphi_{256}$}
\BIC{$\varphi_{169} \to (\mathrm{MonoInc}(a) \to \mathrm{MonoDec}(b) \to \varphi_{57}) \to \varphi_{232}$}
\BIC{$(\mathrm{MonoInc}(a) \to \mathrm{MonoDec}(b) \to \varphi_{57}) \to \varphi_{232}$}
\end{ptree}
\noindent$\Pi_{73}$:
\begin{ptree}
\AXC{$\varphi_{61}$\lbl{A2}}
\AXC{$\stackrel{\Pi_{44}}{\vdots}$}\noLine
\UIC{$\mathrm{MonoInc}(a) \to \mathrm{MonoDec}(b) \to \varphi_{70}$}
\AXC{$\varphi_{239}$\lbl{A1}}
\BIC{$\varphi_{61} \to \mathrm{MonoInc}(a) \to \mathrm{MonoDec}(b) \to \varphi_{57}$}
\BIC{$\mathrm{MonoInc}(a) \to \mathrm{MonoDec}(b) \to \varphi_{57}$}
\AXC{$\stackrel{\Pi_{72}}{\vdots}$}\noLine
\UIC{$(\mathrm{MonoInc}(a) \to \mathrm{MonoDec}(b) \to \varphi_{57}) \to \varphi_{232}$}
\BIC{$\varphi_{232}$}
\end{ptree}
\noindent$\Pi_{74}$:
\begin{ptree}
\AXC{$\mathrm{MonoInc}(a) \to \mathrm{MonoDec}(b) \to \varphi_{88}$\lbl{A1}}
\AXC{$\stackrel{\Pi_{73}}{\vdots}$}\noLine
\UIC{$\varphi_{232}$}
\BIC{$\mathrm{MonoInc}(a) \to \mathrm{MonoDec}(b) \to \varphi_{137}$}
\AXC{$(\mathrm{MonoInc}(a) \to \mathrm{MonoDec}(b) \to \varphi_{137}) \to \varphi_{175}$\lbl{A1}}
\BIC{$\varphi_{175}$}
\UIC{$\forall n\,(\varphi_{175})$}
\end{ptree}
\noindent$\Pi_{75}$:
\begin{ptree}
\AXC{$\varphi_{130}$\lbl{A4}}
\AXC{$\forall x\,(\varphi_{86}) \to \varphi_{120}$\lbl{A3}}
\AXC{$\varphi_{86}$\lbl{A2}}
\UIC{$\forall x\,(\varphi_{86})$}
\AXC{$\forall x\,(\varphi_{86}) \to (\forall x\,(\varphi_{86}) \to \varphi_{120}) \to \varphi_{130} \to \varphi_{103}$\lbl{A1}}
\BIC{$(\forall x\,(\varphi_{86}) \to \varphi_{120}) \to \varphi_{130} \to \varphi_{103}$}
\BIC{$\varphi_{130} \to \varphi_{103}$}
\BIC{$\varphi_{103}$}
\end{ptree}
\noindent$\Pi_{76}$:
\begin{ptree}
\AXC{$\varphi_{104}$\lbl{A4}}
\AXC{$\forall u_{1}\,(\varphi_{80}) \to \varphi_{123}$\lbl{A3}}
\AXC{$\varphi_{80}$\lbl{A2}}
\UIC{$\forall u_{1}\,(\varphi_{80})$}
\AXC{$\forall u_{1}\,(\varphi_{80}) \to (\forall u_{1}\,(\varphi_{80}) \to \varphi_{123}) \to \varphi_{104} \to \varphi_{105}$\lbl{A1}}
\BIC{$(\forall u_{1}\,(\varphi_{80}) \to \varphi_{123}) \to \varphi_{104} \to \varphi_{105}$}
\BIC{$\varphi_{104} \to \varphi_{105}$}
\BIC{$\varphi_{105}$}
\end{ptree}
\noindent$\Pi_{77}$:
\begin{ptree}
\AXC{$\stackrel{\Pi_{75}}{\vdots}$}\noLine
\UIC{$\varphi_{103}$}
\AXC{$\stackrel{\Pi_{76}}{\vdots}$}\noLine
\UIC{$\varphi_{105}$}
\AXC{$\varphi_{81}$\lbl{A1}}
\AXC{$\varphi_{81} \to \varphi_{105} \to \varphi_{103} \to \varphi_{106}$\lbl{A1}}
\BIC{$\varphi_{105} \to \varphi_{103} \to \varphi_{106}$}
\BIC{$\varphi_{103} \to \varphi_{106}$}
\BIC{$\varphi_{106}$}
\AXC{$\varphi_{106} \to \varphi_{117}$\lbl{A1}}
\BIC{$\varphi_{117}$}
\end{ptree}
\noindent$\Pi_{78}$:
\begin{ptree}
\AXC{$\stackrel{\Pi_{77}}{\vdots}$}\noLine
\UIC{$\varphi_{117}$}
\AXC{$v_{729} \in \mathbb{R} \leftrightarrow v_{729} \in \mathbb{R}$\lbl{A1}}
\AXC{$(v_{729} \in \mathbb{R} \leftrightarrow v_{729} \in \mathbb{R}) \to \varphi_{117} \to \varphi_{171}$\lbl{A1}}
\BIC{$\varphi_{117} \to \varphi_{171}$}
\BIC{$\varphi_{171}$}
\AXC{$\varphi_{171} \to \varphi_{174}$\lbl{A1}}
\BIC{$\varphi_{174}$}
\end{ptree}
\noindent$\Pi_{79}$:
\begin{ptree}
\AXC{$\mathrm{UB}(v_{729},v_{730}) \leftrightarrow \mathrm{UB}(v_{729},v_{730})$\lbl{A1}}
\AXC{$\stackrel{\Pi_{78}}{\vdots}$}\noLine
\UIC{$\varphi_{174}$}
\AXC{$\varphi_{174} \to (\mathrm{UB}(v_{729},v_{730}) \leftrightarrow \mathrm{UB}(v_{729},v_{730})) \to \varphi_{210}$\lbl{A1}}
\BIC{$(\mathrm{UB}(v_{729},v_{730}) \leftrightarrow \mathrm{UB}(v_{729},v_{730})) \to \varphi_{210}$}
\BIC{$\varphi_{210}$}
\end{ptree}
\noindent$\Pi_{80}$:
\begin{ptree}
\AXC{$\varphi_{130}$\lbl{A4}}
\AXC{$\forall x\,(\varphi_{86}) \to \varphi_{120}$\lbl{A3}}
\AXC{$\varphi_{86}$\lbl{A2}}
\UIC{$\forall x\,(\varphi_{86})$}
\AXC{$\forall x\,(\varphi_{86}) \to (\forall x\,(\varphi_{86}) \to \varphi_{120}) \to \varphi_{130} \to \varphi_{103}$\lbl{A1}}
\BIC{$(\forall x\,(\varphi_{86}) \to \varphi_{120}) \to \varphi_{130} \to \varphi_{103}$}
\BIC{$\varphi_{130} \to \varphi_{103}$}
\BIC{$\varphi_{103}$}
\end{ptree}
\noindent$\Pi_{81}$:
\begin{ptree}
\AXC{$\varphi_{104}$\lbl{A4}}
\AXC{$\forall u_{1}\,(\varphi_{80}) \to \varphi_{123}$\lbl{A3}}
\AXC{$\varphi_{80}$\lbl{A2}}
\UIC{$\forall u_{1}\,(\varphi_{80})$}
\AXC{$\forall u_{1}\,(\varphi_{80}) \to (\forall u_{1}\,(\varphi_{80}) \to \varphi_{123}) \to \varphi_{104} \to \varphi_{105}$\lbl{A1}}
\BIC{$(\forall u_{1}\,(\varphi_{80}) \to \varphi_{123}) \to \varphi_{104} \to \varphi_{105}$}
\BIC{$\varphi_{104} \to \varphi_{105}$}
\BIC{$\varphi_{105}$}
\end{ptree}
\noindent$\Pi_{82}$:
\begin{ptree}
\AXC{$\stackrel{\Pi_{80}}{\vdots}$}\noLine
\UIC{$\varphi_{103}$}
\AXC{$\stackrel{\Pi_{81}}{\vdots}$}\noLine
\UIC{$\varphi_{105}$}
\AXC{$\varphi_{81}$\lbl{A1}}
\AXC{$\varphi_{81} \to \varphi_{105} \to \varphi_{103} \to \varphi_{106}$\lbl{A1}}
\BIC{$\varphi_{105} \to \varphi_{103} \to \varphi_{106}$}
\BIC{$\varphi_{103} \to \varphi_{106}$}
\BIC{$\varphi_{106}$}
\AXC{$\varphi_{106} \to \varphi_{117}$\lbl{A1}}
\BIC{$\varphi_{117}$}
\end{ptree}
\noindent$\Pi_{83}$:
\begin{ptree}
\AXC{$\stackrel{\Pi_{82}}{\vdots}$}\noLine
\UIC{$\varphi_{117}$}
\AXC{$v_{729} \in \mathbb{R} \leftrightarrow v_{729} \in \mathbb{R}$\lbl{A1}}
\AXC{$(v_{729} \in \mathbb{R} \leftrightarrow v_{729} \in \mathbb{R}) \to \varphi_{117} \to \varphi_{171}$\lbl{A1}}
\BIC{$\varphi_{117} \to \varphi_{171}$}
\BIC{$\varphi_{171}$}
\AXC{$\varphi_{171} \to \varphi_{174}$\lbl{A1}}
\BIC{$\varphi_{174}$}
\end{ptree}
\noindent$\Pi_{84}$:
\begin{ptree}
\AXC{$\stackrel{\Pi_{83}}{\vdots}$}\noLine
\UIC{$\varphi_{174}$}
\AXC{$\mathrm{UB}(v_{729},v_{730}) \leftrightarrow \mathrm{UB}(v_{729},v_{730})$\lbl{A1}}
\AXC{$(\mathrm{UB}(v_{729},v_{730}) \leftrightarrow \mathrm{UB}(v_{729},v_{730})) \to \varphi_{174} \to \varphi_{211}$\lbl{A1}}
\BIC{$\varphi_{174} \to \varphi_{211}$}
\BIC{$\varphi_{211}$}
\end{ptree}
\noindent$\Pi_{85}$:
\begin{ptree}
\AXC{$\stackrel{\Pi_{79}}{\vdots}$}\noLine
\UIC{$\varphi_{210}$}
\AXC{$\varphi_{210} \to \varphi_{213}$\lbl{A1}}
\BIC{$\varphi_{213}$}
\AXC{$\stackrel{\Pi_{84}}{\vdots}$}\noLine
\UIC{$\varphi_{211}$}
\AXC{$\varphi_{211} \to \varphi_{213} \to \varphi_{245}$\lbl{A1}}
\BIC{$\varphi_{213} \to \varphi_{245}$}
\BIC{$\varphi_{245}$}
\AXC{$\varphi_{245} \to (\varphi_{21} \leftrightarrow \varphi_{216})$\lbl{A1}}
\BIC{$\varphi_{21} \leftrightarrow \varphi_{216}$}
\end{ptree}
\noindent$\Pi_{86}$:
\begin{ptree}
\AXC{$\stackrel{\Pi_{85}}{\vdots}$}\noLine
\UIC{$\varphi_{21} \leftrightarrow \varphi_{216}$}
\AXC{$(\varphi_{21} \leftrightarrow \varphi_{216}) \to \varphi_{216} \to \varphi_{21}$\lbl{A1}}
\BIC{$\varphi_{216} \to \varphi_{21}$}
\UIC{$\forall v_{730}\,(\varphi_{216} \to \varphi_{21})$}
\AXC{$\stackrel{\Pi_{85}}{\vdots}$}\noLine
\UIC{$\varphi_{21} \leftrightarrow \varphi_{216}$}
\AXC{$(\varphi_{21} \leftrightarrow \varphi_{216}) \to \varphi_{21} \to \varphi_{216}$\lbl{A1}}
\BIC{$\varphi_{21} \to \varphi_{216}$}
\UIC{$\forall v_{730}\,(\varphi_{21} \to \varphi_{216})$}
\AXC{$\forall v_{730}\,(\varphi_{21} \to \varphi_{216}) \to \varphi_{272}$\lbl{A1}}
\BIC{$\varphi_{272}$}
\BIC{$(\forall v_{730}\,(\varphi_{21} \to \varphi_{216}) \to \forall v_{730}\,(\varphi_{21}) \to \forall v_{730}\,(\varphi_{216})) \to \varphi_{268}$}
\end{ptree}
\noindent$\Pi_{87}$:
\begin{ptree}
\AXC{$\forall v_{730}\,(\varphi_{21} \to \varphi_{216}) \to \forall v_{730}\,(\varphi_{21}) \to \forall v_{730}\,(\varphi_{216})$\lbl{A3}}
\AXC{$\stackrel{\Pi_{86}}{\vdots}$}\noLine
\UIC{$(\forall v_{730}\,(\varphi_{21} \to \varphi_{216}) \to \forall v_{730}\,(\varphi_{21}) \to \forall v_{730}\,(\varphi_{216})) \to \varphi_{268}$}
\BIC{$\varphi_{268}$}
\end{ptree}
\noindent$\Pi_{88}$:
\begin{ptree}
\AXC{$\forall v_{730}\,(\varphi_{216} \to \varphi_{21}) \to \forall v_{730}\,(\varphi_{216}) \to \forall v_{730}\,(\varphi_{21})$\lbl{A3}}
\AXC{$\stackrel{\Pi_{87}}{\vdots}$}\noLine
\UIC{$\varphi_{268}$}
\BIC{$\forall v_{730}\,(\varphi_{21}) \leftrightarrow \forall v_{730}\,(\varphi_{216})$}
\end{ptree}
\noindent$\Pi_{89}$:
\begin{ptree}
\AXC{$\stackrel{\Pi_{88}}{\vdots}$}\noLine
\UIC{$\forall v_{730}\,(\varphi_{21}) \leftrightarrow \forall v_{730}\,(\varphi_{216})$}
\AXC{$\varphi_{262}$\lbl{A1}}
\BIC{$\forall v_{730}\,(\varphi_{21}) \to \forall v_{730}\,(\varphi_{216})$}
\UIC{$\forall v_{729}\,(\forall v_{730}\,(\varphi_{21}) \to \forall v_{730}\,(\varphi_{216}))$}
\AXC{$\forall v_{729}\,(\forall v_{730}\,(\varphi_{21}) \to \forall v_{730}\,(\varphi_{216})) \to \varphi_{274}$\lbl{A1}}
\BIC{$\varphi_{274}$}
\end{ptree}
\noindent$\Pi_{90}$:
\begin{ptree}
\AXC{$\varphi_{265}$\lbl{A3}}
\AXC{$\varphi_{266}$\lbl{A3}}
\AXC{$\stackrel{\Pi_{88}}{\vdots}$}\noLine
\UIC{$\forall v_{730}\,(\varphi_{21}) \leftrightarrow \forall v_{730}\,(\varphi_{216})$}
\AXC{$\varphi_{261}$\lbl{A1}}
\BIC{$\forall v_{730}\,(\varphi_{216}) \to \forall v_{730}\,(\varphi_{21})$}
\UIC{$\forall v_{729}\,(\forall v_{730}\,(\varphi_{216}) \to \forall v_{730}\,(\varphi_{21}))$}
\AXC{$\stackrel{\Pi_{89}}{\vdots}$}\noLine
\UIC{$\varphi_{274}$}
\BIC{$\varphi_{266} \to \varphi_{265} \to (\forall v_{729}\,\forall v_{730}\,(\varphi_{21}) \leftrightarrow \varphi_{68})$}
\BIC{$\varphi_{265} \to (\forall v_{729}\,\forall v_{730}\,(\varphi_{21}) \leftrightarrow \varphi_{68})$}
\BIC{$\forall v_{729}\,\forall v_{730}\,(\varphi_{21}) \leftrightarrow \varphi_{68}$}
\end{ptree}
\noindent$\Pi_{91}$:
\begin{ptree}
\AXC{$\forall a\,\forall X\,(\mathrm{UB}(a,X) \leftrightarrow a \in \mathbb{R} \land \forall x \in X\,(x \leq a))$\lbl{DefUB}}
\AXC{$\forall a\,\forall X\,(\mathrm{UB}(a,X) \leftrightarrow a \in \mathbb{R} \land \forall x \in X\,(x \leq a)) \to \varphi_{14}$\lbl{A2}}
\BIC{$\varphi_{14}$}
\end{ptree}
\noindent$\Pi_{92}$:
\begin{ptree}
\AXC{$\stackrel{\Pi_{90}}{\vdots}$}\noLine
\UIC{$\forall v_{729}\,\forall v_{730}\,(\varphi_{21}) \leftrightarrow \varphi_{68}$}
\AXC{$\stackrel{\Pi_{91}}{\vdots}$}\noLine
\UIC{$\varphi_{14}$}
\AXC{$\varphi_{14} \to \varphi_{21}$\lbl{A2}}
\BIC{$\varphi_{21}$}
\UIC{$\forall v_{730}\,(\varphi_{21})$}
\UIC{$\forall v_{729}\,\forall v_{730}\,(\varphi_{21})$}
\AXC{$\forall v_{729}\,\forall v_{730}\,(\varphi_{21}) \to (\forall v_{729}\,\forall v_{730}\,(\varphi_{21}) \leftrightarrow \varphi_{68}) \to \varphi_{68}$\lbl{A1}}
\BIC{$(\forall v_{729}\,\forall v_{730}\,(\varphi_{21}) \leftrightarrow \varphi_{68}) \to \varphi_{68}$}
\BIC{$\varphi_{68}$}
\end{ptree}
\noindent$\Pi_{93}$:
\begin{ptree}
\AXC{$\varphi_{11}$\lbl{DefSup}}
\AXC{$\varphi_{11} \to \varphi_{7}$\lbl{A2}}
\BIC{$\varphi_{7}$}
\AXC{$\varphi_{7} \to \varphi_{18}$\lbl{A2}}
\BIC{$\varphi_{18}$}
\AXC{$\varphi_{145}$\lbl{A1}}
\AXC{$\varphi_{145} \to \varphi_{18} \to \varphi_{16} \to \mathrm{MonoInc}(a) \to \mathrm{MonoDec}(b) \to \varphi_{45}$\lbl{A1}}
\BIC{$\varphi_{18} \to \varphi_{16} \to \mathrm{MonoInc}(a) \to \mathrm{MonoDec}(b) \to \varphi_{45}$}
\BIC{$\varphi_{16} \to \mathrm{MonoInc}(a) \to \mathrm{MonoDec}(b) \to \varphi_{45}$}
\end{ptree}
\noindent$\Pi_{94}$:
\begin{ptree}
\AXC{$\stackrel{\Pi_{92}}{\vdots}$}\noLine
\UIC{$\varphi_{68}$}
\AXC{$\varphi_{68} \to \varphi_{22}$\lbl{A2}}
\BIC{$\varphi_{22}$}
\AXC{$\varphi_{22} \to \varphi_{16}$\lbl{A2}}
\BIC{$\varphi_{16}$}
\AXC{$\stackrel{\Pi_{93}}{\vdots}$}\noLine
\UIC{$\varphi_{16} \to \mathrm{MonoInc}(a) \to \mathrm{MonoDec}(b) \to \varphi_{45}$}
\BIC{$\mathrm{MonoInc}(a) \to \mathrm{MonoDec}(b) \to \varphi_{45}$}
\end{ptree}
\noindent$\Pi_{95}$:
\begin{ptree}
\AXC{$\stackrel{\Pi_{94}}{\vdots}$}\noLine
\UIC{$\mathrm{MonoInc}(a) \to \mathrm{MonoDec}(b) \to \varphi_{45}$}
\AXC{$(\mathrm{MonoInc}(a) \to \mathrm{MonoDec}(b) \to \varphi_{45}) \to \varphi_{176} \to \varphi_{191}$\lbl{A1}}
\BIC{$\varphi_{176} \to \varphi_{191}$}
\end{ptree}
\noindent$\Pi_{96}$:
\begin{ptree}
\AXC{$\varphi_{219}$\lbl{A4}}
\AXC{$\forall n\,(\varphi_{175}) \to \varphi_{178}$\lbl{A3}}
\AXC{$\stackrel{\Pi_{74}}{\vdots}$}\noLine
\UIC{$\forall n\,(\varphi_{175})$}
\AXC{$\forall n\,(\varphi_{175}) \to (\forall n\,(\varphi_{175}) \to \varphi_{178}) \to \varphi_{219} \to \varphi_{176}$\lbl{A1}}
\BIC{$(\forall n\,(\varphi_{175}) \to \varphi_{178}) \to \varphi_{219} \to \varphi_{176}$}
\BIC{$\varphi_{219} \to \varphi_{176}$}
\BIC{$\varphi_{176}$}
\AXC{$\stackrel{\Pi_{95}}{\vdots}$}\noLine
\UIC{$\varphi_{176} \to \varphi_{191}$}
\BIC{$\varphi_{191}$}
\end{ptree}
\noindent$\Pi_{97}$:
\begin{ptree}
\AXC{$\varphi_{87}$\lbl{A2}}
\AXC{$\stackrel{\Pi_{44}}{\vdots}$}\noLine
\UIC{$\mathrm{MonoInc}(a) \to \mathrm{MonoDec}(b) \to \varphi_{70}$}
\AXC{$\varphi_{240}$\lbl{A1}}
\BIC{$\varphi_{87} \to \mathrm{MonoInc}(a) \to \mathrm{MonoDec}(b) \to \varphi_{77}$}
\BIC{$\mathrm{MonoInc}(a) \to \mathrm{MonoDec}(b) \to \varphi_{77}$}
\AXC{$(\mathrm{MonoInc}(a) \to \mathrm{MonoDec}(b) \to \varphi_{77}) \to \varphi_{156}$\lbl{A1}}
\BIC{$\varphi_{156}$}
\end{ptree}
\noindent$\Pi_{98}$:
\begin{ptree}
\AXC{$x = x$\lbl{A5}}
\UIC{$\forall x\,(x = x)$}
\AXC{$\forall x\,(x = x) \to a_{1_{\omega}} = a_{1_{\omega}}$\lbl{A2}}
\BIC{$a_{1_{\omega}} = a_{1_{\omega}}$}
\AXC{$1_{\omega} \in \mathbb{N}$\lbl{[N1]}}
\AXC{$1_{\omega} \in \mathbb{N} \to \varphi_{186}$\lbl{A1}}
\BIC{$\varphi_{186}$}
\BIC{$\mathrm{MonoInc}(a) \to \mathrm{MonoDec}(b) \to \varphi_{84}$}
\end{ptree}
\noindent$\Pi_{99}$:
\begin{ptree}
\AXC{$\varphi_{110}$\lbl{A2}}
\AXC{$\stackrel{\Pi_{98}}{\vdots}$}\noLine
\UIC{$\mathrm{MonoInc}(a) \to \mathrm{MonoDec}(b) \to \varphi_{84}$}
\AXC{$\varphi_{242}$\lbl{A1}}
\BIC{$\varphi_{110} \to \mathrm{MonoInc}(a) \to \varphi_{125}$}
\BIC{$\mathrm{MonoInc}(a) \to \varphi_{125}$}
\end{ptree}
\noindent$\Pi_{100}$:
\begin{ptree}
\AXC{$\forall y\,(y \in \mathrm{ran}(a) \leftrightarrow \exists j \in \mathbb{N}\,(y = a_{j})) \to \varphi_{39}$\lbl{A2}}
\AXC{$\stackrel{\Pi_{33}}{\vdots}$}\noLine
\UIC{$\varphi_{181}$}
\AXC{$\varphi_{181} \to \mathrm{MonoInc}(a) \to \mathrm{MonoDec}(b) \to \varphi_{115}$\lbl{A1}}
\BIC{$\mathrm{MonoInc}(a) \to \mathrm{MonoDec}(b) \to \varphi_{115}$}
\AXC{$\varphi_{248}$\lbl{A1}}
\BIC{$(\forall y\,(y \in \mathrm{ran}(a) \leftrightarrow \exists j \in \mathbb{N}\,(y = a_{j})) \to \varphi_{39}) \to \varphi_{179}$}
\BIC{$\varphi_{179}$}
\end{ptree}
\noindent$\Pi_{101}$:
\begin{ptree}
\AXC{$\stackrel{\Pi_{99}}{\vdots}$}\noLine
\UIC{$\mathrm{MonoInc}(a) \to \varphi_{125}$}
\AXC{$\stackrel{\Pi_{100}}{\vdots}$}\noLine
\UIC{$\varphi_{179}$}
\AXC{$\varphi_{179} \to (\mathrm{MonoInc}(a) \to \varphi_{125}) \to \mathrm{MonoInc}(a) \to \varphi_{73}$\lbl{A1}}
\BIC{$(\mathrm{MonoInc}(a) \to \varphi_{125}) \to \mathrm{MonoInc}(a) \to \varphi_{73}$}
\BIC{$\mathrm{MonoInc}(a) \to \varphi_{73}$}
\end{ptree}
\noindent$\Pi_{102}$:
\begin{ptree}
\AXC{$\stackrel{\Pi_{101}}{\vdots}$}\noLine
\UIC{$\mathrm{MonoInc}(a) \to \varphi_{73}$}
\AXC{$\forall X\,\forall x \in X\,(X \neq \emptyset)$\lbl{[B2]}}
\AXC{$\varphi_{48}$\lbl{A2}}
\BIC{$\forall x \in \mathrm{ran}(a)\,(\mathrm{ran}(a) \neq \emptyset)$}
\AXC{$\varphi_{111}$\lbl{A2}}
\BIC{$a_{1_{\omega}} \in \mathrm{ran}(a) \to \mathrm{ran}(a) \neq \emptyset$}
\AXC{$\varphi_{228}$\lbl{A1}}
\BIC{$(\mathrm{MonoInc}(a) \to \varphi_{73}) \to \mathrm{MonoInc}(a) \to \varphi_{51}$}
\BIC{$\mathrm{MonoInc}(a) \to \varphi_{51}$}
\end{ptree}
\noindent$\Pi_{103}$:
\begin{ptree}
\AXC{$\varphi_{31}$\lbl{(R16)}}
\AXC{$\varphi_{31} \to \mathrm{ran}(a) \subseteq \mathbb{R} \to \mathrm{ran}(a) \neq \emptyset \to \varphi_{3}$\lbl{A2}}
\BIC{$\mathrm{ran}(a) \subseteq \mathbb{R} \to \mathrm{ran}(a) \neq \emptyset \to \varphi_{3}$}
\AXC{$(\mathrm{ran}(a) \subseteq \mathbb{R} \to \mathrm{ran}(a) \neq \emptyset \to \varphi_{3}) \to \varphi_{247}$\lbl{A1}}
\BIC{$\varphi_{247}$}
\end{ptree}
\noindent$\Pi_{104}$:
\begin{ptree}
\AXC{$\stackrel{\Pi_{102}}{\vdots}$}\noLine
\UIC{$\mathrm{MonoInc}(a) \to \varphi_{51}$}
\AXC{$\stackrel{\Pi_{33}}{\vdots}$}\noLine
\UIC{$\varphi_{181}$}
\AXC{$\varphi_{181} \to \mathrm{MonoInc}(a) \to \varphi_{52}$\lbl{A1}}
\BIC{$\mathrm{MonoInc}(a) \to \varphi_{52}$}
\AXC{$\stackrel{\Pi_{103}}{\vdots}$}\noLine
\UIC{$\varphi_{247}$}
\BIC{$(\mathrm{MonoInc}(a) \to \varphi_{51}) \to \varphi_{220}$}
\BIC{$\varphi_{220}$}
\end{ptree}
\noindent$\Pi_{105}$:
\begin{ptree}
\AXC{$\varphi_{8}$\lbl{A2}}
\AXC{$1_{\omega} \in \mathbb{N}$\lbl{[N1]}}
\AXC{$\stackrel{\Pi_{97}}{\vdots}$}\noLine
\UIC{$\varphi_{156}$}
\BIC{$\mathrm{MonoInc}(a) \to \mathrm{MonoDec}(b) \to \varphi_{32}$}
\AXC{$\varphi_{227}$\lbl{A1}}
\BIC{$\varphi_{8} \to \mathrm{MonoInc}(a) \to \varphi_{46}$}
\BIC{$\mathrm{MonoInc}(a) \to \varphi_{46}$}
\AXC{$\stackrel{\Pi_{104}}{\vdots}$}\noLine
\UIC{$\varphi_{220}$}
\BIC{$\mathrm{MonoInc}(a) \to \mathrm{MonoDec}(b) \to \varphi_{17}$}
\AXC{$\varphi_{269}$\lbl{A1}}
\BIC{$\forall x\,(\varphi_{197}) \to (\forall x\,(\varphi_{197}) \to \varphi_{206}) \to \varphi_{234} \to \forall s\,(\varphi_{4}) \to \varphi_{238}$}
\end{ptree}
\noindent$\Pi_{106}$:
\begin{ptree}
\AXC{$\varphi_{167}$\lbl{A2}}
\AXC{$\stackrel{\Pi_{96}}{\vdots}$}\noLine
\UIC{$\varphi_{191}$}
\AXC{$\varphi_{191} \to \varphi_{167} \to \varphi_{194}$\lbl{A1}}
\BIC{$\varphi_{167} \to \varphi_{194}$}
\BIC{$\varphi_{194}$}
\AXC{$\varphi_{194} \to \varphi_{197}$\lbl{A1}}
\BIC{$\varphi_{197}$}
\UIC{$\forall x\,(\varphi_{197})$}
\AXC{$\stackrel{\Pi_{105}}{\vdots}$}\noLine
\UIC{$\forall x\,(\varphi_{197}) \to (\forall x\,(\varphi_{197}) \to \varphi_{206}) \to \varphi_{234} \to \forall s\,(\varphi_{4}) \to \varphi_{238}$}
\BIC{$(\forall x\,(\varphi_{197}) \to \varphi_{206}) \to \varphi_{234} \to \forall s\,(\varphi_{4}) \to \varphi_{238}$}
\end{ptree}
\noindent 主木:
\begin{ptree}
\AXC{$\varphi_{9}$\lbl{A4}}
\AXC{$\forall s\,(\varphi_{4}) \to \varphi_{10}$\lbl{A3}}
\AXC{$\varphi_{4}$\lbl{A2}}
\UIC{$\forall s\,(\varphi_{4})$}
\AXC{$\varphi_{234}$\lbl{A4}}
\AXC{$\forall x\,(\varphi_{197}) \to \varphi_{206}$\lbl{A3}}
\AXC{$\stackrel{\Pi_{106}}{\vdots}$}\noLine
\UIC{$(\forall x\,(\varphi_{197}) \to \varphi_{206}) \to \varphi_{234} \to \forall s\,(\varphi_{4}) \to \varphi_{238}$}
\BIC{$\varphi_{234} \to \forall s\,(\varphi_{4}) \to \varphi_{238}$}
\BIC{$\forall s\,(\varphi_{4}) \to \varphi_{238}$}
\BIC{$\varphi_{238}$}
\BIC{$\varphi_{9} \to \mathrm{MonoInc}(a) \to \mathrm{MonoDec}(b) \to \varphi_{138}$}
\BIC{$\mathrm{MonoInc}(a) \to \mathrm{MonoDec}(b) \to \varphi_{138}$}
\UIC{$\forall b\,(\mathrm{MonoInc}(a) \to \mathrm{MonoDec}(b) \to \varphi_{138})$}
\UIC{$\forall a\,\forall b\,(\mathrm{MonoInc}(a) \to \mathrm{MonoDec}(b) \to \varphi_{138})$}
\end{ptree}
\begin{abbrevlist}
\item $\varphi_{1} :\equiv \mathrm{IsSup}(x,\mathrm{ran}(a)) \to \mathrm{IsSup}(x,\mathrm{ran}(a))$
\item $\varphi_{2} :\equiv \mathrm{MonoInc}(a) \to \mathrm{MonoDec}(b) \to \mathrm{MonoDec}(b)$
\item $\varphi_{3} :\equiv \exists b\,(\mathrm{UB}(b,\mathrm{ran}(a))) \to \exists s\,(\mathrm{IsSup}(s,\mathrm{ran}(a)))$
\item $\varphi_{4} :\equiv \forall x\,(\neg \mathrm{IsSup}(x,\mathrm{ran}(a))) \to \neg \mathrm{IsSup}(s,\mathrm{ran}(a))$
\item $\varphi_{5} :\equiv \mathrm{IsSup}(x,\mathrm{ran}(a)) \to \forall w\,(\mathrm{UB}(w,\mathrm{ran}(a)) \to x \leq w)$
\item $\varphi_{6} :\equiv \forall y\,(v_{727} = y \to y = v_{727}) \to v_{727} = v_{728} \to v_{728} = v_{727}$
\item $\varphi_{7} :\equiv \forall X\,(\mathrm{IsSup}(x,X) \leftrightarrow \mathrm{UB}(x,X) \land \forall w\,(\mathrm{UB}(w,X) \to x \leq w))$
\item $\varphi_{8} :\equiv \forall b\,(\neg \mathrm{UB}(b,\mathrm{ran}(a))) \to \neg \mathrm{UB}(b_{1_{\omega}},\mathrm{ran}(a))$
\item $\varphi_{9} :\equiv \forall x\,(\neg \mathrm{IsSup}(x,\mathrm{ran}(a))) \to \forall s\,\forall x\,(\neg \mathrm{IsSup}(x,\mathrm{ran}(a)))$
\item $\varphi_{10} :\equiv \forall s\,\forall x\,(\neg \mathrm{IsSup}(x,\mathrm{ran}(a))) \to \forall s\,(\neg \mathrm{IsSup}(s,\mathrm{ran}(a)))$
\item $\varphi_{11} :\equiv \forall s\,\forall X\,(\mathrm{IsSup}(s,X) \leftrightarrow \mathrm{UB}(s,X) \land \forall w\,(\mathrm{UB}(w,X) \to s \leq w))$
\item $\varphi_{12} :\equiv \mathrm{IsSup}(x,\mathrm{ran}(a)) \to \mathrm{UB}(b_{n},\mathrm{ran}(a)) \to x \leq b_{n}$
\item $\varphi_{13} :\equiv \forall X\,(\mathrm{UB}(v_{731},X) \leftrightarrow v_{731} \in \mathbb{R} \land \forall x \in X\,(x \leq v_{731}))$
\item $\varphi_{14} :\equiv \forall X\,(\mathrm{UB}(v_{729},X) \leftrightarrow v_{729} \in \mathbb{R} \land \forall x \in X\,(x \leq v_{729}))$
\item $\varphi_{15} :\equiv \mathrm{UB}(x,\mathrm{ran}(a)) \leftrightarrow x \in \mathbb{R} \land \forall u_{2} \in \mathrm{ran}(a)\,(u_{2} \leq x)$
\item $\varphi_{16} :\equiv \mathrm{UB}(x,\mathrm{ran}(a)) \leftrightarrow x \in \mathbb{R} \land \forall u_{1} \in \mathrm{ran}(a)\,(u_{1} \leq x)$
\item $\varphi_{17} :\equiv \forall n \in \mathbb{N}\,(a_{n} \leq b_{n}) \to \exists s\,(\mathrm{IsSup}(s,\mathrm{ran}(a)))$
\item $\varphi_{18} :\equiv \mathrm{IsSup}(x,\mathrm{ran}(a)) \leftrightarrow \mathrm{UB}(x,\mathrm{ran}(a)) \land \forall w\,(\mathrm{UB}(w,\mathrm{ran}(a)) \to x \leq w)$
\item $\varphi_{19} :\equiv \mathrm{UB}(v_{731},v_{732}) \leftrightarrow v_{731} \in \mathbb{R} \land \forall x \in v_{732}\,(x \leq v_{731})$
\item $\varphi_{20} :\equiv \forall v_{732}\,(\mathrm{UB}(x,v_{732}) \leftrightarrow x \in \mathbb{R} \land \forall u_{2} \in v_{732}\,(u_{2} \leq x))$
\item $\varphi_{21} :\equiv \mathrm{UB}(v_{729},v_{730}) \leftrightarrow v_{729} \in \mathbb{R} \land \forall x \in v_{730}\,(x \leq v_{729})$
\item $\varphi_{22} :\equiv \forall v_{730}\,(\mathrm{UB}(x,v_{730}) \leftrightarrow x \in \mathbb{R} \land \forall u_{1} \in v_{730}\,(u_{1} \leq x))$
\item $\varphi_{23} :\equiv \forall j\,(\forall n\,(\neg (n \in \mathbb{N} \land y = a_{n})) \to \neg (j \in \mathbb{N} \land y = a_{j}))$
\item $\varphi_{24} :\equiv y \in \mathrm{ran}(a) \to (y \in \mathrm{ran}(a) \leftrightarrow \exists j \in \mathbb{N}\,(y = a_{j}))$
\item $\varphi_{25} :\equiv \forall w\,(\mathrm{UB}(w,\mathrm{ran}(a)) \to x \leq w) \to \mathrm{UB}(b_{n},\mathrm{ran}(a)) \to x \leq b_{n}$
\item $\varphi_{26} :\equiv v_{726} \in \mathbb{N} \to v_{725} \leq v_{726} \to v_{724}(v_{725}) \leq v_{724}(v_{726})$
\item $\varphi_{27} :\equiv \mathrm{IsSup}(x,\mathrm{ran}(a)) \to \forall n \in \mathbb{N}\,(a_{n} \leq x \land x \leq b_{n})$
\item $\varphi_{28} :\equiv \forall v_{727}\,\forall v_{728}\,(v_{727} = v_{728} \to v_{728} = v_{727}) \to \forall v_{728}\,(y = v_{728} \to v_{728} = y)$
\item $\varphi_{29} :\equiv \mathrm{MonoInc}(a) \leftrightarrow \mathrm{Seq}(a) \land \forall n \in \mathbb{N}\,(a_{n} \leq a_{S(n)})$
\item $\varphi_{30} :\equiv \mathrm{MonoDec}(b) \leftrightarrow \mathrm{Seq}(b) \land \forall n \in \mathbb{N}\,(b_{S(n)} \leq b_{n})$
\item $\varphi_{31} :\equiv \forall X\,(X \subseteq \mathbb{R} \to X \neq \emptyset \to \exists b\,(\mathrm{UB}(b,X)) \to \exists s\,(\mathrm{IsSup}(s,X)))$
\item $\varphi_{32} :\equiv \forall n \in \mathbb{N}\,(a_{n} \leq b_{n}) \to \mathrm{UB}(b_{1_{\omega}},\mathrm{ran}(a))$
\item $\varphi_{33} :\equiv \forall n\,(\neg (n \in \mathbb{N} \land y = a_{n})) \to \forall j\,\forall n\,(\neg (n \in \mathbb{N} \land y = a_{n}))$
\item $\varphi_{34} :\equiv \mathrm{MonoInc}(a) \land \mathrm{MonoDec}(b) \land \forall n \in \mathbb{N}\,(a_{n} \leq b_{n})$
\item $\varphi_{35} :\equiv \forall j\,\forall n\,(\neg (n \in \mathbb{N} \land y = a_{n})) \to \forall j\,(\neg (j \in \mathbb{N} \land y = a_{j}))$
\item $\varphi_{36} :\equiv y \in \mathrm{ran}(a) \to \forall y\,(y \in \mathrm{ran}(a) \leftrightarrow \exists j \in \mathbb{N}\,(y = a_{j}))$
\item $\varphi_{37} :\equiv \mathrm{IsSup}(x,\mathrm{ran}(a)) \to n \in \mathbb{N} \to n \in \mathbb{N} \land a_{n} = a_{n}$
\item $\varphi_{38} :\equiv \forall a\,(\mathrm{MonoDec}(a) \leftrightarrow \mathrm{Seq}(a) \land \forall n \in \mathbb{N}\,(a_{S(n)} \leq a_{n}))$
\item $\varphi_{39} :\equiv a_{1_{\omega}} \in \mathrm{ran}(a) \leftrightarrow \exists j \in \mathbb{N}\,(a_{1_{\omega}} = a_{j})$
\item $\varphi_{40} :\equiv m \in \mathbb{N} \to y \in \mathrm{ran}(a) \to n \in \mathbb{N} \land y = a_{n} \to n \leq m \to n \leq m$
\item $\varphi_{41} :\equiv \forall x\,\forall y\,\forall z\,(x \leq y \to y \leq z \to x \leq z) \to \forall y\,\forall z\,(a_{n} \leq y \to y \leq z \to a_{n} \leq z)$
\item $\varphi_{42} :\equiv \mathrm{ran}(a) \subseteq \mathbb{R} \land \forall y\,(y \in \mathrm{ran}(a) \leftrightarrow \exists j \in \mathbb{N}\,(y = a_{j}))$
\item $\varphi_{43} :\equiv \mathrm{IsSup}(x,\mathrm{ran}(a)) \to n \in \mathbb{N} \to \exists j \in \mathbb{N}\,(a_{n} = a_{j})$
\item $\varphi_{44} :\equiv \forall n\,(\mathrm{MonoDec}(b) \to m \in \mathbb{N} \to n \in \mathbb{N} \to m \leq n \to b_{n} \leq b_{m})$
\item $\varphi_{45} :\equiv \forall n \in \mathbb{N}\,(a_{n} \leq b_{n}) \to \mathrm{IsSup}(x,\mathrm{ran}(a)) \to x \in \mathbb{R}$
\item $\varphi_{46} :\equiv \mathrm{MonoDec}(b) \to \forall n \in \mathbb{N}\,(a_{n} \leq b_{n}) \to \exists b\,(\mathrm{UB}(b,\mathrm{ran}(a)))$
\item $\varphi_{47} :\equiv \forall v_{723}\,(v_{722} \in \omega \to v_{723} \in \omega \to v_{722} \leq v_{723} \lor v_{723} < v_{722})$
\item $\varphi_{48} :\equiv \forall X\,\forall x \in X\,(X \neq \emptyset) \to \forall x \in \mathrm{ran}(a)\,(\mathrm{ran}(a) \neq \emptyset)$
\item $\varphi_{49} :\equiv m \in \mathbb{N} \to y \in \mathrm{ran}(a) \to n \in \mathbb{N} \land y = a_{n} \to m < n \to m \leq n$
\item $\varphi_{50} :\equiv m \in \mathbb{N} \to y \in \mathrm{ran}(a) \to n \in \mathbb{N} \land y = a_{n} \to n \leq m \lor m < n$
\item $\varphi_{51} :\equiv \mathrm{MonoDec}(b) \to \forall n \in \mathbb{N}\,(a_{n} \leq b_{n}) \to \mathrm{ran}(a) \neq \emptyset$
\item $\varphi_{52} :\equiv \mathrm{MonoDec}(b) \to \forall n \in \mathbb{N}\,(a_{n} \leq b_{n}) \to \mathrm{ran}(a) \subseteq \mathbb{R}$
\item $\varphi_{53} :\equiv \forall y \in \mathrm{ran}(a)\,(\neg y < x \to y = x) \to u_{2} \in \mathrm{ran}(a) \to \neg u_{2} < x \to u_{2} = x$
\item $\varphi_{54} :\equiv \forall j\,(\neg (j \in \mathbb{N} \land a_{n} = a_{j})) \to \neg (n \in \mathbb{N} \land a_{n} = a_{n})$
\item $\varphi_{55} :\equiv \forall u_{2} \in \mathrm{ran}(a)\,(\neg u_{2} < x \to u_{2} = x) \to y \in \mathrm{ran}(a) \to \neg y < x \to y = x$
\item $\varphi_{56} :\equiv \forall n \in \mathbb{N}\,(a_{n} \leq b_{n}) \to m \in \mathbb{N} \to \mathrm{UB}(b_{m},\mathrm{ran}(a))$
\item $\varphi_{57} :\equiv \forall n \in \mathbb{N}\,(a_{n} \leq b_{n}) \to n \in \mathbb{N} \to \mathrm{UB}(b_{n},\mathrm{ran}(a))$
\item $\varphi_{58} :\equiv \mathrm{MonoInc}(a) \to \mathrm{MonoDec}(b) \to \forall n \in \mathbb{N}\,(a_{n} \leq b_{n}) \to \mathrm{Seq}(a)$
\item $\varphi_{59} :\equiv \mathrm{MonoInc}(a) \to \mathrm{MonoDec}(b) \to \forall n \in \mathbb{N}\,(a_{n} \leq b_{n}) \to \mathrm{Seq}(b)$
\item $\varphi_{60} :\equiv m \in \mathbb{N} \to y \in \mathrm{ran}(a) \to n \in \mathbb{N} \land y = a_{n} \to m < n \to m < n$
\item $\varphi_{61} :\equiv \forall m \in \mathbb{N}\,(\mathrm{UB}(b_{m},\mathrm{ran}(a))) \to n \in \mathbb{N} \to \mathrm{UB}(b_{n},\mathrm{ran}(a))$
\item $\varphi_{62} :\equiv \forall a\,\forall m\,\forall n\,(\mathrm{MonoDec}(a) \to m \in \mathbb{N} \to n \in \mathbb{N} \to m \leq n \to a_{n} \leq a_{m})$
\item $\varphi_{63} :\equiv \forall a\,\forall m\,\forall n\,(\mathrm{MonoInc}(a) \to m \in \mathbb{N} \to n \in \mathbb{N} \to m \leq n \to a_{m} \leq a_{n})$
\item $\varphi_{64} :\equiv \forall y \in \mathrm{ran}(a)\,(\neg y < x \to y = x) \to \forall u_{2}\,\forall y \in \mathrm{ran}(a)\,(\neg y < x \to y = x)$
\item $\varphi_{65} :\equiv \forall y \in \mathrm{ran}(a)\,(\neg y < x \to y = x) \to \forall u_{2} \in \mathrm{ran}(a)\,(\neg u_{2} < x \to u_{2} = x)$
\item $\varphi_{66} :\equiv \forall u_{2} \in \mathrm{ran}(a)\,(\neg u_{2} < x \to u_{2} = x) \to \forall y \in \mathrm{ran}(a)\,(\neg y < x \to y = x)$
\item $\varphi_{67} :\equiv \forall v_{731}\,\forall v_{732}\,(\varphi_{215})$
\item $\varphi_{68} :\equiv \forall v_{729}\,\forall v_{730}\,(\varphi_{216})$
\item $\varphi_{69} :\equiv \mathrm{MonoDec}(b) \to \forall n \in \mathbb{N}\,(a_{n} \leq b_{n}) \to m \in \mathbb{N} \to m \in \mathbb{N}$
\item $\varphi_{70} :\equiv \forall n \in \mathbb{N}\,(a_{n} \leq b_{n}) \to \forall m \in \mathbb{N}\,(\mathrm{UB}(b_{m},\mathrm{ran}(a)))$
\item $\varphi_{71} :\equiv \mathrm{IsSup}(x,\mathrm{ran}(a)) \to x \in \mathbb{R} \land \forall n \in \mathbb{N}\,(a_{n} \leq x \land x \leq b_{n})$
\item $\varphi_{72} :\equiv \mathrm{UB}(b_{m},\mathrm{ran}(a)) \leftrightarrow b_{m} \in \mathbb{R} \land \forall x \in \mathrm{ran}(a)\,(x \leq b_{m})$
\item $\varphi_{73} :\equiv \mathrm{MonoDec}(b) \to \forall n \in \mathbb{N}\,(a_{n} \leq b_{n}) \to a_{1_{\omega}} \in \mathrm{ran}(a)$
\item $\varphi_{74} :\equiv \forall n \in \mathbb{N}\,(a_{n} \leq b_{n}) \to m \in \mathbb{N} \to y \in \mathrm{ran}(a) \to y \leq b_{m}$
\item $\varphi_{75} :\equiv \forall y\,\forall u_{2} \in \mathrm{ran}(a)\,(\neg u_{2} < x \to u_{2} = x) \to \forall y \in \mathrm{ran}(a)\,(\neg y < x \to y = x)$
\item $\varphi_{76} :\equiv \mathrm{IsSup}(x,\mathrm{ran}(a)) \to \exists x \in \mathbb{R}\,\forall n \in \mathbb{N}\,(a_{n} \leq x \land x \leq b_{n})$
\item $\varphi_{77} :\equiv \forall n \in \mathbb{N}\,(a_{n} \leq b_{n}) \to 1_{\omega} \in \mathbb{N} \to \mathrm{UB}(b_{1_{\omega}},\mathrm{ran}(a))$
\item $\varphi_{78} :\equiv \forall x \in v_{732}\,(\neg x < v_{731} \to x = v_{731}) \to u_{2} \in v_{732} \to \neg u_{2} < v_{731} \to u_{2} = v_{731}$
\item $\varphi_{79} :\equiv \forall x \in v_{732}\,(\neg x < v_{731} \to x = v_{731}) \leftrightarrow \forall x \in v_{732}\,(\neg x < v_{731} \to x = v_{731})$
\item $\varphi_{80} :\equiv \forall x \in v_{730}\,(\neg x < v_{729} \to x = v_{729}) \to u_{1} \in v_{730} \to \neg u_{1} < v_{729} \to u_{1} = v_{729}$
\item $\varphi_{81} :\equiv \forall x \in v_{730}\,(\neg x < v_{729} \to x = v_{729}) \leftrightarrow \forall x \in v_{730}\,(\neg x < v_{729} \to x = v_{729})$
\item $\varphi_{82} :\equiv \mathrm{IsSup}(x,\mathrm{ran}(a)) \to n \in \mathbb{N} \to \forall y\,(y \in \mathrm{ran}(a) \leftrightarrow \exists j \in \mathbb{N}\,(y = a_{j}))$
\item $\varphi_{83} :\equiv \forall u_{2}\,\forall y \in \mathrm{ran}(a)\,(\neg y < x \to y = x) \to \forall u_{2} \in \mathrm{ran}(a)\,(\neg u_{2} < x \to u_{2} = x)$
\item $\varphi_{84} :\equiv \forall n \in \mathbb{N}\,(a_{n} \leq b_{n}) \to 1_{\omega} \in \mathbb{N} \land a_{1_{\omega}} = a_{1_{\omega}}$
\item $\varphi_{85} :\equiv \forall u_{2} \in v_{732}\,(\neg u_{2} < v_{731} \to u_{2} = v_{731}) \to x \in v_{732} \to \neg x < v_{731} \to x = v_{731}$
\item $\varphi_{86} :\equiv \forall u_{1} \in v_{730}\,(\neg u_{1} < v_{729} \to u_{1} = v_{729}) \to x \in v_{730} \to \neg x < v_{729} \to x = v_{729}$
\item $\varphi_{87} :\equiv \forall m \in \mathbb{N}\,(\mathrm{UB}(b_{m},\mathrm{ran}(a))) \to 1_{\omega} \in \mathbb{N} \to \mathrm{UB}(b_{1_{\omega}},\mathrm{ran}(a))$
\item $\varphi_{88} :\equiv \forall n \in \mathbb{N}\,(a_{n} \leq b_{n}) \to \mathrm{IsSup}(x,\mathrm{ran}(a)) \to n \in \mathbb{N} \to n \in \mathbb{N}$
\item $\varphi_{89} :\equiv m \in \mathbb{N} \to y \in \mathrm{ran}(a) \to n \in \mathbb{N} \land y = a_{n} \to n \in \mathbb{N} \land y = a_{n}$
\item $\varphi_{90} :\equiv \forall n \in \mathbb{N}\,(a_{n} \leq b_{n}) \to m \in \mathbb{N} \to y \in \mathrm{ran}(a) \to y \in \mathrm{ran}(a)$
\item $\varphi_{91} :\equiv \forall n \in \mathbb{N}\,(a_{n} \leq b_{n}) \to \mathrm{IsSup}(x,\mathrm{ran}(a)) \to \forall y \in \mathrm{ran}(a)\,(y \leq x)$
\item $\varphi_{92} :\equiv \forall m\,\forall n\,(\mathrm{MonoInc}(v_{724}) \to m \in \mathbb{N} \to n \in \mathbb{N} \to m \leq n \to v_{724}(m) \leq v_{724}(n))$
\item $\varphi_{93} :\equiv \forall v_{726}\,(\mathrm{MonoInc}(a) \to n \in \mathbb{N} \to v_{726} \in \mathbb{N} \to n \leq v_{726} \to a_{n} \leq a_{v_{726}})$
\item $\varphi_{94} :\equiv m \in \mathbb{N} \to y \in \mathrm{ran}(a) \to n \in \mathbb{N} \land y = a_{n} \to n \leq m \to a_{m} \leq b_{m}$
\item $\varphi_{95} :\equiv m \in \mathbb{N} \to y \in \mathrm{ran}(a) \to n \in \mathbb{N} \land y = a_{n} \to n \leq m \to a_{n} \leq a_{m}$
\item $\varphi_{96} :\equiv m \in \mathbb{N} \to y \in \mathrm{ran}(a) \to n \in \mathbb{N} \land y = a_{n} \to n \leq m \to a_{n} \leq b_{m}$
\item $\varphi_{97} :\equiv \forall n \in \mathbb{N}\,(a_{n} \leq b_{n}) \to m \in \mathbb{N} \to \forall y \in \mathrm{ran}(a)\,(y \leq b_{m})$
\item $\varphi_{98} :\equiv \forall u_{2} \in \mathrm{ran}(a)\,(\neg u_{2} < x \to u_{2} = x) \to \forall y\,\forall u_{2} \in \mathrm{ran}(a)\,(\neg u_{2} < x \to u_{2} = x)$
\item $\varphi_{99} :\equiv \forall u_{2} \in v_{732}\,(\neg u_{2} < v_{731} \to u_{2} = v_{731}) \to \forall x \in v_{732}\,(\neg x < v_{731} \to x = v_{731})$
\item $\varphi_{100} :\equiv \forall x \in v_{732}\,(\neg x < v_{731} \to x = v_{731}) \to \forall u_{2}\,\forall x \in v_{732}\,(\neg x < v_{731} \to x = v_{731})$
\item $\varphi_{101} :\equiv \forall x \in v_{732}\,(\neg x < v_{731} \to x = v_{731}) \to \forall u_{2} \in v_{732}\,(\neg u_{2} < v_{731} \to u_{2} = v_{731})$
\item $\varphi_{102} :\equiv \forall x \in v_{732}\,(\neg x < v_{731} \to x = v_{731}) \leftrightarrow \forall u_{2} \in v_{732}\,(\neg u_{2} < v_{731} \to u_{2} = v_{731})$
\item $\varphi_{103} :\equiv \forall u_{1} \in v_{730}\,(\neg u_{1} < v_{729} \to u_{1} = v_{729}) \to \forall x \in v_{730}\,(\neg x < v_{729} \to x = v_{729})$
\item $\varphi_{104} :\equiv \forall x \in v_{730}\,(\neg x < v_{729} \to x = v_{729}) \to \forall u_{1}\,\forall x \in v_{730}\,(\neg x < v_{729} \to x = v_{729})$
\item $\varphi_{105} :\equiv \forall x \in v_{730}\,(\neg x < v_{729} \to x = v_{729}) \to \forall u_{1} \in v_{730}\,(\neg u_{1} < v_{729} \to u_{1} = v_{729})$
\item $\varphi_{106} :\equiv \forall x \in v_{730}\,(\neg x < v_{729} \to x = v_{729}) \leftrightarrow \forall u_{1} \in v_{730}\,(\neg u_{1} < v_{729} \to u_{1} = v_{729})$
\item $\varphi_{107} :\equiv \mathrm{MonoDec}(b) \to \forall n \in \mathbb{N}\,(a_{n} \leq b_{n}) \to n \in \mathbb{N} \to a_{n} \leq b_{n}$
\item $\varphi_{108} :\equiv \mathrm{MonoDec}(b) \to \forall n \in \mathbb{N}\,(a_{n} \leq b_{n}) \to m \in \mathbb{N} \to a_{m} \leq b_{m}$
\item $\varphi_{109} :\equiv \mathrm{MonoDec}(b) \to \forall n \in \mathbb{N}\,(a_{n} \leq b_{n}) \to m \in \mathbb{N} \to b_{m} \in \mathbb{R}$
\item $\varphi_{110} :\equiv \forall j\,(\neg (j \in \mathbb{N} \land a_{1_{\omega}} = a_{j})) \to \neg (1_{\omega} \in \mathbb{N} \land a_{1_{\omega}} = a_{1_{\omega}})$
\item $\varphi_{111} :\equiv \forall x \in \mathrm{ran}(a)\,(\mathrm{ran}(a) \neq \emptyset) \to a_{1_{\omega}} \in \mathrm{ran}(a) \to \mathrm{ran}(a) \neq \emptyset$
\item $\varphi_{112} :\equiv m \in \mathbb{N} \to y \in \mathrm{ran}(a) \to n \in \mathbb{N} \land y = a_{n} \to m < n \to b_{n} \leq b_{m}$
\item $\varphi_{113} :\equiv m \in \mathbb{N} \to y \in \mathrm{ran}(a) \to n \in \mathbb{N} \land y = a_{n} \to m < n \to a_{n} \leq b_{n}$
\item $\varphi_{114} :\equiv m \in \mathbb{N} \to y \in \mathrm{ran}(a) \to n \in \mathbb{N} \land y = a_{n} \to m < n \to a_{n} \leq b_{m}$
\item $\varphi_{115} :\equiv \forall n \in \mathbb{N}\,(a_{n} \leq b_{n}) \to \forall y\,(y \in \mathrm{ran}(a) \leftrightarrow \exists j \in \mathbb{N}\,(y = a_{j}))$
\item $\varphi_{116} :\equiv \neg \forall x \in v_{732}\,(\neg x < v_{731} \to x = v_{731}) \leftrightarrow \neg \forall u_{2} \in v_{732}\,(\neg u_{2} < v_{731} \to u_{2} = v_{731})$
\item $\varphi_{117} :\equiv \neg \forall x \in v_{730}\,(\neg x < v_{729} \to x = v_{729}) \leftrightarrow \neg \forall u_{1} \in v_{730}\,(\neg u_{1} < v_{729} \to u_{1} = v_{729})$
\item $\varphi_{118} :\equiv \mathrm{MonoDec}(b) \to \forall n \in \mathbb{N}\,(a_{n} \leq b_{n}) \to \forall n \in \mathbb{N}\,(a_{n} \leq b_{n})$
\item $\varphi_{119} :\equiv \forall x\,\forall u_{2} \in v_{732}\,(\neg u_{2} < v_{731} \to u_{2} = v_{731}) \to \forall x \in v_{732}\,(\neg x < v_{731} \to x = v_{731})$
\item $\varphi_{120} :\equiv \forall x\,\forall u_{1} \in v_{730}\,(\neg u_{1} < v_{729} \to u_{1} = v_{729}) \to \forall x \in v_{730}\,(\neg x < v_{729} \to x = v_{729})$
\item $\varphi_{121} :\equiv \forall n\,(\mathrm{MonoInc}(v_{724}) \to v_{725} \in \mathbb{N} \to n \in \mathbb{N} \to v_{725} \leq n \to v_{724}(v_{725}) \leq v_{724}(n))$
\item $\varphi_{122} :\equiv \forall u_{2}\,\forall x \in v_{732}\,(\neg x < v_{731} \to x = v_{731}) \to \forall u_{2} \in v_{732}\,(\neg u_{2} < v_{731} \to u_{2} = v_{731})$
\item $\varphi_{123} :\equiv \forall u_{1}\,\forall x \in v_{730}\,(\neg x < v_{729} \to x = v_{729}) \to \forall u_{1} \in v_{730}\,(\neg u_{1} < v_{729} \to u_{1} = v_{729})$
\item $\varphi_{124} :\equiv \forall n\,(\mathrm{Seq}(b) \to n \in \mathbb{N} \to b_{n} \in \mathbb{R}) \to \mathrm{Seq}(b) \to m \in \mathbb{N} \to b_{m} \in \mathbb{R}$
\item $\varphi_{125} :\equiv \mathrm{MonoDec}(b) \to \forall n \in \mathbb{N}\,(a_{n} \leq b_{n}) \to \exists j \in \mathbb{N}\,(a_{1_{\omega}} = a_{j})$
\item $\varphi_{126} :\equiv \forall z\,\forall w\,(z = w \to z = b_{m} \to w = b_{m}) \to \forall w\,(a_{n} = w \to a_{n} = b_{m} \to w = b_{m})$
\item $\varphi_{127} :\equiv \forall y\,\forall z\,(a_{n} \leq y \to y \leq z \to a_{n} \leq z) \to \forall z\,(a_{n} \leq b_{n} \to b_{n} \leq z \to a_{n} \leq z)$
\item $\varphi_{128} :\equiv \forall y\,\forall z\,(a_{n} \leq y \to y \leq z \to a_{n} \leq z) \to \forall z\,(a_{n} \leq a_{m} \to a_{m} \leq z \to a_{n} \leq z)$
\item $\varphi_{129} :\equiv \forall u_{2} \in v_{732}\,(\neg u_{2} < v_{731} \to u_{2} = v_{731}) \to \forall x\,\forall u_{2} \in v_{732}\,(\neg u_{2} < v_{731} \to u_{2} = v_{731})$
\item $\varphi_{130} :\equiv \forall u_{1} \in v_{730}\,(\neg u_{1} < v_{729} \to u_{1} = v_{729}) \to \forall x\,\forall u_{1} \in v_{730}\,(\neg u_{1} < v_{729} \to u_{1} = v_{729})$
\item $\varphi_{131} :\equiv \mathrm{IsSup}(x,\mathrm{ran}(a)) \to n \in \mathbb{N} \to (a_{n} \in \mathrm{ran}(a) \leftrightarrow \exists j \in \mathbb{N}\,(a_{n} = a_{j}))$
\item $\varphi_{132} :\equiv \forall n \in \mathbb{N}\,(a_{n} \leq b_{n}) \to \mathrm{IsSup}(x,\mathrm{ran}(a)) \to n \in \mathbb{N} \to a_{n} \in \mathrm{ran}(a)$
\item $\varphi_{133} :\equiv \forall n \in \mathbb{N}\,(a_{n} \leq b_{n}) \to \mathrm{IsSup}(x,\mathrm{ran}(a)) \to a_{n} \in \mathrm{ran}(a) \to a_{n} \leq x$
\item $\varphi_{134} :\equiv \forall y\,(y \in \mathrm{ran}(a) \leftrightarrow \exists j \in \mathbb{N}\,(y = a_{j})) \to (y \in \mathrm{ran}(a) \leftrightarrow \exists j \in \mathbb{N}\,(y = a_{j}))$
\item $\varphi_{135} :\equiv \forall n \in \mathbb{N}\,(a_{n} \leq b_{n}) \to m \in \mathbb{N} \to y \in \mathrm{ran}(a) \to \exists j \in \mathbb{N}\,(y = a_{j})$
\item $\varphi_{136} :\equiv \forall a\,\forall n\,(\mathrm{Seq}(a) \to n \in \mathbb{N} \to a_{n} \in \mathbb{R}) \to \forall n\,(\mathrm{Seq}(b) \to n \in \mathbb{N} \to b_{n} \in \mathbb{R})$
\item $\varphi_{137} :\equiv \forall n \in \mathbb{N}\,(a_{n} \leq b_{n}) \to \mathrm{IsSup}(x,\mathrm{ran}(a)) \to n \in \mathbb{N} \to a_{n} \leq x \land x \leq b_{n}$
\item $\varphi_{138} :\equiv \forall n \in \mathbb{N}\,(a_{n} \leq b_{n}) \to \exists x \in \mathbb{R}\,\forall n \in \mathbb{N}\,(a_{n} \leq x \land x \leq b_{n})$
\item $\varphi_{139} :\equiv \forall y \in \mathrm{ran}(a)\,(\neg y < b_{m} \to y = b_{m}) \to x \in \mathrm{ran}(a) \to \neg x < b_{m} \to x = b_{m}$
\item $\varphi_{140} :\equiv \forall x \in \mathrm{ran}(a)\,(\neg x < b_{m} \to x = b_{m}) \to y \in \mathrm{ran}(a) \to \neg y < b_{m} \to y = b_{m}$
\item $\varphi_{141} :\equiv \forall v_{725}\,\forall v_{726}\,(\mathrm{MonoInc}(a) \to v_{725} \in \mathbb{N} \to v_{726} \in \mathbb{N} \to v_{725} \leq v_{726} \to a_{v_{725}} \leq a_{v_{726}})$
\item $\varphi_{142} :\equiv \forall w\,(a_{n} = w \to a_{n} = b_{m} \to w = b_{m}) \to a_{n} = y \to a_{n} = b_{m} \to y = b_{m}$
\item $\varphi_{143} :\equiv \forall y \in \mathrm{ran}(a)\,(\neg y < b_{m} \to y = b_{m}) \to \forall x \in \mathrm{ran}(a)\,(\neg x < b_{m} \to x = b_{m})$
\item $\varphi_{144} :\equiv \forall x \in \mathrm{ran}(a)\,(\neg x < b_{m} \to x = b_{m}) \to \forall y \in \mathrm{ran}(a)\,(\neg y < b_{m} \to y = b_{m})$
\item $\varphi_{145} :\equiv \mathrm{MonoInc}(a) \to \mathrm{MonoDec}(b) \to \forall n \in \mathbb{N}\,(a_{n} \leq b_{n}) \to \varphi_{1}$
\item $\varphi_{146} :\equiv \forall y \in \mathrm{ran}(a)\,(\neg y < b_{m} \to y = b_{m}) \to \forall x\,\forall y \in \mathrm{ran}(a)\,(\neg y < b_{m} \to y = b_{m})$
\item $\varphi_{147} :\equiv \forall x \in \mathrm{ran}(a)\,(\neg x < b_{m} \to x = b_{m}) \to \forall y\,\forall x \in \mathrm{ran}(a)\,(\neg x < b_{m} \to x = b_{m})$
\item $\varphi_{148} :\equiv \forall z\,\forall w\,(z = w \to z < b_{m} \to w < b_{m}) \to \forall w\,(a_{n} = w \to a_{n} < b_{m} \to w < b_{m})$
\item $\varphi_{149} :\equiv \forall x\,\forall y \in \mathrm{ran}(a)\,(\neg y < b_{m} \to y = b_{m}) \to \forall x \in \mathrm{ran}(a)\,(\neg x < b_{m} \to x = b_{m})$
\item $\varphi_{150} :\equiv \forall y\,\forall x \in \mathrm{ran}(a)\,(\neg x < b_{m} \to x = b_{m}) \to \forall y \in \mathrm{ran}(a)\,(\neg y < b_{m} \to y = b_{m})$
\item $\varphi_{151} :\equiv \forall n \in \mathbb{N}\,(a_{n} \leq b_{n}) \to m \in \mathbb{N} \to y \in \mathrm{ran}(a) \to n \in \mathbb{N} \land y = a_{n} \to n \in \mathbb{N}$
\item $\varphi_{152} :\equiv \forall n \in \mathbb{N}\,(a_{n} \leq b_{n}) \to m \in \mathbb{N} \to y \in \mathrm{ran}(a) \to n \in \mathbb{N} \land y = a_{n} \to y = a_{n}$
\item $\varphi_{153} :\equiv \forall n \in \mathbb{N}\,(a_{n} \leq b_{n}) \to m \in \mathbb{N} \to y \in \mathrm{ran}(a) \to n \in \mathbb{N} \land y = a_{n} \to a_{n} = y$
\item $\varphi_{154} :\equiv \forall n \in \mathbb{N}\,(a_{n} \leq b_{n}) \to m \in \mathbb{N} \to y \in \mathrm{ran}(a) \to n \in \mathbb{N} \land y = a_{n} \to y \leq b_{m}$
\item $\varphi_{155} :\equiv \forall v_{723}\,(n \in \omega \to v_{723} \in \omega \to n \leq v_{723} \lor v_{723} < n) \to n \in \omega \to m \in \omega \to n \leq m \lor m < n$
\item $\varphi_{156} :\equiv 1_{\omega} \in \mathbb{N} \to \mathrm{MonoInc}(a) \to \mathrm{MonoDec}(b) \to \varphi_{32}$
\item $\varphi_{157} :\equiv \forall n \in \mathbb{N}\,(a_{n} \leq b_{n}) \to m \in \mathbb{N} \to y \in \mathrm{ran}(a) \to n \in \mathbb{N} \land y = a_{n} \to m \in \omega$
\item $\varphi_{158} :\equiv \forall n \in \mathbb{N}\,(a_{n} \leq b_{n}) \to m \in \mathbb{N} \to y \in \mathrm{ran}(a) \to n \in \mathbb{N} \land y = a_{n} \to n \in \omega$
\item $\varphi_{159} :\equiv \forall v_{725}\,\forall v_{726}\,(\mathrm{MonoInc}(v_{724}) \to v_{725} \in \mathbb{N} \to \varphi_{26})$
\item $\varphi_{160} :\equiv \mathrm{MonoInc}(a) \to \mathrm{MonoDec}(b) \to \forall n \in \mathbb{N}\,(a_{n} \leq b_{n}) \to \varphi_{5}$
\item $\varphi_{161} :\equiv \forall y\,(y \in \mathrm{ran}(a) \leftrightarrow \exists j \in \mathbb{N}\,(y = a_{j})) \to (a_{n} \in \mathrm{ran}(a) \leftrightarrow \exists j \in \mathbb{N}\,(a_{n} = a_{j}))$
\item $\varphi_{162} :\equiv \forall m\,\forall n\,(m \in \omega \to n \in \omega \to m \leq n \lor n < m) \to \forall n\,(v_{722} \in \omega \to n \in \omega \to v_{722} \leq n \lor n < v_{722})$
\item $\varphi_{163} :\equiv \forall w\,(a_{n} = w \to a_{n} < b_{m} \to w < b_{m}) \to a_{n} = y \to a_{n} < b_{m} \to y < b_{m}$
\item $\varphi_{164} :\equiv \forall n \in \mathbb{N}\,(a_{n} \leq b_{n}) \to m \in \mathbb{N} \to y \in \mathrm{ran}(a) \to n \in \mathbb{N} \land y = a_{n} \to a_{n} \leq b_{m}$
\item $\varphi_{165} :\equiv \forall z\,(a_{n} \leq b_{n} \to b_{n} \leq z \to a_{n} \leq z) \to a_{n} \leq b_{n} \to b_{n} \leq b_{m} \to a_{n} \leq b_{m}$
\item $\varphi_{166} :\equiv \forall z\,(a_{n} \leq a_{m} \to a_{m} \leq z \to a_{n} \leq z) \to a_{n} \leq a_{m} \to a_{m} \leq b_{m} \to a_{n} \leq b_{m}$
\item $\varphi_{167} :\equiv \forall x\,(\neg (x \in \mathbb{R} \land \forall n \in \mathbb{N}\,(a_{n} \leq x \land x \leq b_{n}))) \to \neg (x \in \mathbb{R} \land \forall n \in \mathbb{N}\,(a_{n} \leq x \land x \leq b_{n}))$
\item $\varphi_{168} :\equiv \forall a\,\forall X\,(\mathrm{UB}(a,X) \leftrightarrow a \in \mathbb{R} \land \forall x \in X\,(x \leq a)) \to \forall X\,(\mathrm{UB}(b_{m},X) \leftrightarrow b_{m} \in \mathbb{R} \land \forall x \in X\,(x \leq b_{m}))$
\item $\varphi_{169} :\equiv \mathrm{MonoInc}(a) \to \mathrm{MonoDec}(b) \to \forall n \in \mathbb{N}\,(a_{n} \leq b_{n}) \to \varphi_{12}$
\item $\varphi_{170} :\equiv v_{731} \in \mathbb{R} \to \neg \forall x \in v_{732}\,(\neg x < v_{731} \to x = v_{731}) \leftrightarrow v_{731} \in \mathbb{R} \to \neg \forall u_{2} \in v_{732}\,(\neg u_{2} < v_{731} \to u_{2} = v_{731})$
\item $\varphi_{171} :\equiv v_{729} \in \mathbb{R} \to \neg \forall x \in v_{730}\,(\neg x < v_{729} \to x = v_{729}) \leftrightarrow v_{729} \in \mathbb{R} \to \neg \forall u_{1} \in v_{730}\,(\neg u_{1} < v_{729} \to u_{1} = v_{729})$
\item $\varphi_{172} :\equiv \forall n\,(v_{722} \in \omega \to n \in \omega \to v_{722} \leq n \lor n < v_{722}) \to v_{722} \in \omega \to v_{723} \in \omega \to v_{722} \leq v_{723} \lor v_{723} < v_{722}$
\item $\varphi_{173} :\equiv \neg (v_{731} \in \mathbb{R} \to \neg \forall x \in v_{732}\,(\neg x < v_{731} \to x = v_{731})) \leftrightarrow \neg (v_{731} \in \mathbb{R} \to \neg \forall u_{2} \in v_{732}\,(\neg u_{2} < v_{731} \to u_{2} = v_{731}))$
\item $\varphi_{174} :\equiv \neg (v_{729} \in \mathbb{R} \to \neg \forall x \in v_{730}\,(\neg x < v_{729} \to x = v_{729})) \leftrightarrow \neg (v_{729} \in \mathbb{R} \to \neg \forall u_{1} \in v_{730}\,(\neg u_{1} < v_{729} \to u_{1} = v_{729}))$
\item $\varphi_{175} :\equiv \varphi_{34} \land \mathrm{IsSup}(x,\mathrm{ran}(a)) \to n \in \mathbb{N} \to a_{n} \leq x \land x \leq b_{n}$
\item $\varphi_{176} :\equiv \mathrm{MonoInc}(a) \to \mathrm{MonoDec}(b) \to \forall n \in \mathbb{N}\,(a_{n} \leq b_{n}) \to \varphi_{27}$
\item $\varphi_{177} :\equiv \mathrm{MonoInc}(a) \to \mathrm{MonoDec}(b) \to \forall n \in \mathbb{N}\,(a_{n} \leq b_{n}) \to \varphi_{37}$
\item $\varphi_{178} :\equiv \forall n\,(\varphi_{34} \land \mathrm{IsSup}(x,\mathrm{ran}(a))) \to \forall n \in \mathbb{N}\,(a_{n} \leq x \land x \leq b_{n})$
\item $\varphi_{179} :\equiv \mathrm{MonoInc}(a) \to \mathrm{MonoDec}(b) \to \forall n \in \mathbb{N}\,(a_{n} \leq b_{n}) \to \varphi_{39}$
\item $\varphi_{180} :\equiv \mathrm{MonoInc}(a) \to \mathrm{MonoDec}(b) \to \forall n \in \mathbb{N}\,(a_{n} \leq b_{n}) \to \varphi_{40}$
\item $\varphi_{181} :\equiv \mathrm{MonoInc}(a) \to \mathrm{MonoDec}(b) \to \forall n \in \mathbb{N}\,(a_{n} \leq b_{n}) \to \varphi_{42}$
\item $\varphi_{182} :\equiv \mathrm{MonoInc}(a) \to \mathrm{MonoDec}(b) \to \forall n \in \mathbb{N}\,(a_{n} \leq b_{n}) \to \varphi_{43}$
\item $\varphi_{183} :\equiv \mathrm{MonoInc}(a) \to \mathrm{MonoDec}(b) \to \forall n \in \mathbb{N}\,(a_{n} \leq b_{n}) \to \varphi_{49}$
\item $\varphi_{184} :\equiv \mathrm{MonoInc}(a) \to \mathrm{MonoDec}(b) \to \forall n \in \mathbb{N}\,(a_{n} \leq b_{n}) \to \varphi_{50}$
\item $\varphi_{185} :\equiv \varphi_{34} \land m \in \mathbb{N} \land y \in \mathrm{ran}(a) \land \neg y \leq b_{m} \to \neg (n \in \mathbb{N} \land y = a_{n})$
\item $\varphi_{186} :\equiv a_{1_{\omega}} = a_{1_{\omega}} \to \mathrm{MonoInc}(a) \to \mathrm{MonoDec}(b) \to \varphi_{84}$
\item $\varphi_{187} :\equiv \varphi_{44} \to \mathrm{MonoDec}(b) \to m \in \mathbb{N} \to n \in \mathbb{N} \to m \leq n \to b_{n} \leq b_{m}$
\item $\varphi_{188} :\equiv \mathrm{MonoInc}(a) \to \mathrm{MonoDec}(b) \to \forall n \in \mathbb{N}\,(a_{n} \leq b_{n}) \to \varphi_{60}$
\item $\varphi_{189} :\equiv \forall v_{722}\,(\varphi_{47}) \to \forall v_{723}\,(n \in \omega \to v_{723} \in \omega \to n \leq v_{723} \lor v_{723} < n)$
\item $\varphi_{190} :\equiv \mathrm{MonoInc}(a) \to \mathrm{MonoDec}(b) \to \forall n \in \mathbb{N}\,(a_{n} \leq b_{n}) \to m \in \mathbb{N} \to \varphi_{24}$
\item $\varphi_{191} :\equiv \mathrm{MonoInc}(a) \to \mathrm{MonoDec}(b) \to \forall n \in \mathbb{N}\,(a_{n} \leq b_{n}) \to \varphi_{71}$
\item $\varphi_{192} :\equiv \mathrm{MonoInc}(a) \to \mathrm{MonoDec}(b) \to \forall n \in \mathbb{N}\,(a_{n} \leq b_{n}) \to m \in \mathbb{N} \to \varphi_{36}$
\item $\varphi_{193} :\equiv \forall n\,(\varphi_{34} \land m \in \mathbb{N} \land y \in \mathrm{ran}(a) \land \neg y \leq b_{m}) \to \forall n\,(\neg (n \in \mathbb{N} \land y = a_{n}))$
\item $\varphi_{194} :\equiv \mathrm{MonoInc}(a) \to \mathrm{MonoDec}(b) \to \forall n \in \mathbb{N}\,(a_{n} \leq b_{n}) \to \varphi_{76}$
\item $\varphi_{195} :\equiv \mathrm{MonoInc}(a) \to \mathrm{MonoDec}(b) \to \forall n \in \mathbb{N}\,(a_{n} \leq b_{n}) \to \varphi_{82}$
\item $\varphi_{196} :\equiv \mathrm{MonoInc}(a) \to \mathrm{MonoDec}(b) \to \forall n \in \mathbb{N}\,(a_{n} \leq b_{n}) \to \varphi_{89}$
\item $\varphi_{197} :\equiv \varphi_{34} \land \neg \exists x \in \mathbb{R}\,\forall n \in \mathbb{N}\,(a_{n} \leq x \land x \leq b_{n}) \to \neg \mathrm{IsSup}(x,\mathrm{ran}(a))$
\item $\varphi_{198} :\equiv \mathrm{MonoInc}(a) \to \mathrm{MonoDec}(b) \to \forall n \in \mathbb{N}\,(a_{n} \leq b_{n}) \to \varphi_{94}$
\item $\varphi_{199} :\equiv \mathrm{MonoInc}(a) \to \mathrm{MonoDec}(b) \to \forall n \in \mathbb{N}\,(a_{n} \leq b_{n}) \to \varphi_{95}$
\item $\varphi_{200} :\equiv \mathrm{MonoInc}(a) \to \mathrm{MonoDec}(b) \to \forall n \in \mathbb{N}\,(a_{n} \leq b_{n}) \to \varphi_{96}$
\item $\varphi_{201} :\equiv \forall X\,(\mathrm{UB}(b_{m},X) \leftrightarrow b_{m} \in \mathbb{R} \land \forall x \in X\,(x \leq b_{m})) \to \varphi_{72}$
\item $\varphi_{202} :\equiv \varphi_{93} \to \mathrm{MonoInc}(a) \to n \in \mathbb{N} \to m \in \mathbb{N} \to n \leq m \to a_{n} \leq a_{m}$
\item $\varphi_{203} :\equiv \mathrm{MonoInc}(a) \to \mathrm{MonoDec}(b) \to \forall n \in \mathbb{N}\,(a_{n} \leq b_{n}) \to \varphi_{112}$
\item $\varphi_{204} :\equiv \mathrm{MonoInc}(a) \to \mathrm{MonoDec}(b) \to \forall n \in \mathbb{N}\,(a_{n} \leq b_{n}) \to \varphi_{113}$
\item $\varphi_{205} :\equiv \mathrm{MonoInc}(a) \to \mathrm{MonoDec}(b) \to \forall n \in \mathbb{N}\,(a_{n} \leq b_{n}) \to \varphi_{114}$
\item $\varphi_{206} :\equiv \forall x\,(\varphi_{34} \land \neg \exists x \in \mathbb{R}\,\forall n \in \mathbb{N}\,(a_{n} \leq x \land x \leq b_{n})) \to \forall x\,(\neg \mathrm{IsSup}(x,\mathrm{ran}(a)))$
\item $\varphi_{207} :\equiv \mathrm{MonoInc}(a) \to \mathrm{MonoDec}(b) \to \forall n \in \mathbb{N}\,(a_{n} \leq b_{n}) \to \varphi_{131}$
\item $\varphi_{208} :\equiv \neg (v_{731} \in \mathbb{R} \to \neg \forall x \in v_{732}\,(\neg x < v_{731} \to x = v_{731})) \to \mathrm{UB}(v_{731},v_{732}) \leftrightarrow \neg (v_{731} \in \mathbb{R} \to \neg \forall u_{2} \in v_{732}\,(\neg u_{2} < v_{731} \to u_{2} = v_{731})) \to \mathrm{UB}(v_{731},v_{732})$
\item $\varphi_{209} :\equiv \mathrm{UB}(v_{731},v_{732}) \to \neg (v_{731} \in \mathbb{R} \to \neg \forall x \in v_{732}\,(\neg x < v_{731} \to x = v_{731})) \leftrightarrow \mathrm{UB}(v_{731},v_{732}) \to \neg (v_{731} \in \mathbb{R} \to \neg \forall u_{2} \in v_{732}\,(\neg u_{2} < v_{731} \to u_{2} = v_{731}))$
\item $\varphi_{210} :\equiv \neg (v_{729} \in \mathbb{R} \to \neg \forall x \in v_{730}\,(\neg x < v_{729} \to x = v_{729})) \to \mathrm{UB}(v_{729},v_{730}) \leftrightarrow \neg (v_{729} \in \mathbb{R} \to \neg \forall u_{1} \in v_{730}\,(\neg u_{1} < v_{729} \to u_{1} = v_{729})) \to \mathrm{UB}(v_{729},v_{730})$
\item $\varphi_{211} :\equiv \mathrm{UB}(v_{729},v_{730}) \to \neg (v_{729} \in \mathbb{R} \to \neg \forall x \in v_{730}\,(\neg x < v_{729} \to x = v_{729})) \leftrightarrow \mathrm{UB}(v_{729},v_{730}) \to \neg (v_{729} \in \mathbb{R} \to \neg \forall u_{1} \in v_{730}\,(\neg u_{1} < v_{729} \to u_{1} = v_{729}))$
\item $\varphi_{212} :\equiv \neg (\neg (v_{731} \in \mathbb{R} \to \neg \forall x \in v_{732}\,(\neg x < v_{731} \to x = v_{731})) \to \mathrm{UB}(v_{731},v_{732})) \leftrightarrow \neg (\neg (v_{731} \in \mathbb{R} \to \neg \forall u_{2} \in v_{732}\,(\neg u_{2} < v_{731} \to u_{2} = v_{731})) \to \mathrm{UB}(v_{731},v_{732}))$
\item $\varphi_{213} :\equiv \neg (\neg (v_{729} \in \mathbb{R} \to \neg \forall x \in v_{730}\,(\neg x < v_{729} \to x = v_{729})) \to \mathrm{UB}(v_{729},v_{730})) \leftrightarrow \neg (\neg (v_{729} \in \mathbb{R} \to \neg \forall u_{1} \in v_{730}\,(\neg u_{1} < v_{729} \to u_{1} = v_{729})) \to \mathrm{UB}(v_{729},v_{730}))$
\item $\varphi_{214} :\equiv (y = a_{n} \to a_{n} = y) \to \mathrm{MonoInc}(a) \to \mathrm{MonoDec}(b) \to \varphi_{153}$
\item $\varphi_{215} :\equiv \neg ((\mathrm{UB}(v_{731},v_{732}) \to \neg (v_{731} \in \mathbb{R} \to \neg \forall u_{2} \in v_{732}\,(\neg u_{2} < v_{731} \to u_{2} = v_{731}))) \to \neg (\neg (v_{731} \in \mathbb{R} \to \neg \forall u_{2} \in v_{732}\,(\neg u_{2} < v_{731} \to u_{2} = v_{731})) \to \mathrm{UB}(v_{731},v_{732})))$
\item $\varphi_{216} :\equiv \neg ((\mathrm{UB}(v_{729},v_{730}) \to \neg (v_{729} \in \mathbb{R} \to \neg \forall u_{1} \in v_{730}\,(\neg u_{1} < v_{729} \to u_{1} = v_{729}))) \to \neg (\neg (v_{729} \in \mathbb{R} \to \neg \forall u_{1} \in v_{730}\,(\neg u_{1} < v_{729} \to u_{1} = v_{729})) \to \mathrm{UB}(v_{729},v_{730})))$
\item $\varphi_{217} :\equiv (\forall n \in \mathbb{N}\,(a_{n} \leq b_{n}) \to n \in \mathbb{N} \to a_{n} \leq b_{n}) \to \mathrm{MonoInc}(a) \to \varphi_{107}$
\item $\varphi_{218} :\equiv (\forall n \in \mathbb{N}\,(a_{n} \leq b_{n}) \to m \in \mathbb{N} \to a_{m} \leq b_{m}) \to \mathrm{MonoInc}(a) \to \varphi_{108}$
\item $\varphi_{219} :\equiv \varphi_{34} \land \mathrm{IsSup}(x,\mathrm{ran}(a)) \to \forall n\,(\varphi_{34} \land \mathrm{IsSup}(x,\mathrm{ran}(a)))$
\item $\varphi_{220} :\equiv (\mathrm{MonoInc}(a) \to \varphi_{46}) \to \mathrm{MonoInc}(a) \to \mathrm{MonoDec}(b) \to \varphi_{17}$
\item $\varphi_{221} :\equiv (a_{n} = y \to a_{n} = b_{m} \to y = b_{m}) \to \mathrm{MonoInc}(a) \to \mathrm{MonoDec}(b) \to \varphi_{154}$
\item $\varphi_{222} :\equiv (\forall y \in \mathrm{ran}(a)\,(y \leq x) \to a_{n} \in \mathrm{ran}(a) \to a_{n} \leq x) \to \mathrm{MonoInc}(a) \to \mathrm{MonoDec}(b) \to \varphi_{133}$
\item $\varphi_{223} :\equiv (\mathrm{MonoInc}(a) \to \mathrm{MonoDec}(b) \to \varphi_{56}) \to \varphi_{34} \to m \in \mathbb{N} \to \mathrm{UB}(b_{m},\mathrm{ran}(a))$
\item $\varphi_{224} :\equiv \forall m\,(\varphi_{34} \to m \in \mathbb{N} \to \mathrm{UB}(b_{m},\mathrm{ran}(a))) \to \forall m\,(\varphi_{34}) \to \forall m \in \mathbb{N}\,(\mathrm{UB}(b_{m},\mathrm{ran}(a)))$
\item $\varphi_{225} :\equiv (\mathrm{MonoInc}(a) \to \mathrm{MonoDec}(b) \to \varphi_{74}) \to \varphi_{34} \land m \in \mathbb{N} \to y \in \mathrm{ran}(a) \to y \leq b_{m}$
\item $\varphi_{226} :\equiv \varphi_{34} \land m \in \mathbb{N} \land y \in \mathrm{ran}(a) \land \neg y \leq b_{m} \to \forall n\,(\varphi_{34} \land m \in \mathbb{N} \land y \in \mathrm{ran}(a) \land \neg y \leq b_{m})$
\item $\varphi_{227} :\equiv (\mathrm{MonoInc}(a) \to \mathrm{MonoDec}(b) \to \varphi_{32}) \to \varphi_{8} \to \mathrm{MonoInc}(a) \to \varphi_{46}$
\item $\varphi_{228} :\equiv (a_{1_{\omega}} \in \mathrm{ran}(a) \to \mathrm{ran}(a) \neq \emptyset) \to (\mathrm{MonoInc}(a) \to \varphi_{73}) \to \mathrm{MonoInc}(a) \to \varphi_{51}$
\item $\varphi_{229} :\equiv \forall y\,(\varphi_{34} \land m \in \mathbb{N} \to y \in \mathrm{ran}(a) \to y \leq b_{m}) \to \forall y\,(\varphi_{34} \land m \in \mathbb{N}) \to \forall y \in \mathrm{ran}(a)\,(y \leq b_{m})$
\item $\varphi_{230} :\equiv (\mathrm{MonoInc}(a) \to \varphi_{69}) \to \mathrm{MonoInc}(a) \to \mathrm{MonoDec}(b) \to \varphi_{157}$
\item $\varphi_{231} :\equiv (\mathrm{MonoInc}(a) \to \mathrm{MonoDec}(b) \to \varphi_{90}) \to \mathrm{MonoInc}(a) \to \mathrm{MonoDec}(b) \to \varphi_{135}$
\item $\varphi_{232} :\equiv (\mathrm{MonoInc}(a) \to \mathrm{MonoDec}(b) \to \varphi_{88}) \to \mathrm{MonoInc}(a) \to \mathrm{MonoDec}(b) \to \varphi_{137}$
\item $\varphi_{233} :\equiv (\varphi_{34} \to \forall m\,(\varphi_{34})) \to \mathrm{MonoInc}(a) \to \mathrm{MonoDec}(b) \to \varphi_{70}$
\item $\varphi_{234} :\equiv \varphi_{34} \land \neg \exists x \in \mathbb{R}\,\forall n \in \mathbb{N}\,(a_{n} \leq x \land x \leq b_{n}) \to \forall x\,(\varphi_{34} \land \neg \exists x \in \mathbb{R}\,\forall n \in \mathbb{N}\,(a_{n} \leq x \land x \leq b_{n}))$
\item $\varphi_{235} :\equiv (\mathrm{MonoInc}(a) \to \mathrm{MonoDec}(b) \to \varphi_{88}) \to a_{n} = a_{n} \to \varphi_{177}$
\item $\varphi_{236} :\equiv (\mathrm{MonoInc}(a) \to \mathrm{MonoDec}(b) \to \varphi_{151}) \to \mathrm{MonoInc}(a) \to \mathrm{MonoDec}(b) \to \varphi_{158}$
\item $\varphi_{237} :\equiv (\varphi_{34} \land m \in \mathbb{N} \to \forall y\,(\varphi_{34} \land m \in \mathbb{N})) \to \mathrm{MonoInc}(a) \to \mathrm{MonoDec}(b) \to \varphi_{97}$
\item $\varphi_{238} :\equiv (\forall s\,(\varphi_{4}) \to \varphi_{10}) \to \varphi_{9} \to \mathrm{MonoInc}(a) \to \mathrm{MonoDec}(b) \to \varphi_{138}$
\item $\varphi_{239} :\equiv (\mathrm{MonoInc}(a) \to \mathrm{MonoDec}(b) \to \varphi_{70}) \to \varphi_{61} \to \mathrm{MonoInc}(a) \to \mathrm{MonoDec}(b) \to \varphi_{57}$
\item $\varphi_{240} :\equiv (\mathrm{MonoInc}(a) \to \mathrm{MonoDec}(b) \to \varphi_{70}) \to \varphi_{87} \to \mathrm{MonoInc}(a) \to \mathrm{MonoDec}(b) \to \varphi_{77}$
\item $\varphi_{241} :\equiv (\varphi_{23} \to \varphi_{35}) \to \varphi_{33} \to \mathrm{MonoInc}(a) \to \mathrm{MonoDec}(b) \to \varphi_{74}$
\item $\varphi_{242} :\equiv (\mathrm{MonoInc}(a) \to \mathrm{MonoDec}(b) \to \varphi_{84}) \to \varphi_{110} \to \mathrm{MonoInc}(a) \to \varphi_{125}$
\item $\varphi_{243} :\equiv (\mathrm{Seq}(b) \to m \in \mathbb{N} \to b_{m} \in \mathbb{R}) \to \varphi_{59} \to (\mathrm{MonoInc}(a) \to \varphi_{69}) \to \mathrm{MonoInc}(a) \to \varphi_{109}$
\item $\varphi_{244} :\equiv (\mathrm{UB}(v_{731},v_{732}) \to \neg (v_{731} \in \mathbb{R} \to \neg \forall x \in v_{732}\,(\neg x < v_{731} \to x = v_{731}))) \to \neg (\neg (v_{731} \in \mathbb{R} \to \neg \forall x \in v_{732}\,(\neg x < v_{731} \to x = v_{731})) \to \mathrm{UB}(v_{731},v_{732})) \leftrightarrow (\mathrm{UB}(v_{731},v_{732}) \to \neg (v_{731} \in \mathbb{R} \to \neg \forall u_{2} \in v_{732}\,(\neg u_{2} < v_{731} \to u_{2} = v_{731}))) \to \neg (\neg (v_{731} \in \mathbb{R} \to \neg \forall u_{2} \in v_{732}\,(\neg u_{2} < v_{731} \to u_{2} = v_{731})) \to \mathrm{UB}(v_{731},v_{732}))$
\item $\varphi_{245} :\equiv (\mathrm{UB}(v_{729},v_{730}) \to \neg (v_{729} \in \mathbb{R} \to \neg \forall x \in v_{730}\,(\neg x < v_{729} \to x = v_{729}))) \to \neg (\neg (v_{729} \in \mathbb{R} \to \neg \forall x \in v_{730}\,(\neg x < v_{729} \to x = v_{729})) \to \mathrm{UB}(v_{729},v_{730})) \leftrightarrow (\mathrm{UB}(v_{729},v_{730}) \to \neg (v_{729} \in \mathbb{R} \to \neg \forall u_{1} \in v_{730}\,(\neg u_{1} < v_{729} \to u_{1} = v_{729}))) \to \neg (\neg (v_{729} \in \mathbb{R} \to \neg \forall u_{1} \in v_{730}\,(\neg u_{1} < v_{729} \to u_{1} = v_{729})) \to \mathrm{UB}(v_{729},v_{730}))$
\item $\varphi_{246} :\equiv (\mathrm{MonoInc}(a) \to \varphi_{108}) \to (\mathrm{MonoInc}(a) \to \varphi_{69}) \to \varphi_{198}$
\item $\varphi_{247} :\equiv (\mathrm{MonoInc}(a) \to \varphi_{52}) \to (\mathrm{MonoInc}(a) \to \varphi_{51}) \to \varphi_{220}$
\item $\varphi_{248} :\equiv (\mathrm{MonoInc}(a) \to \mathrm{MonoDec}(b) \to \varphi_{115}) \to (\forall y\,(y \in \mathrm{ran}(a) \leftrightarrow \exists j \in \mathbb{N}\,(y = a_{j})) \to \varphi_{39}) \to \varphi_{179}$
\item $\varphi_{249} :\equiv (\mathrm{MonoInc}(a) \to \mathrm{MonoDec}(b) \to \varphi_{164}) \to (a_{n} = y \to a_{n} < b_{m} \to y < b_{m}) \to \varphi_{221}$
\item $\varphi_{250} :\equiv \varphi_{18} \to \varphi_{15} \to \varphi_{66} \to \varphi_{65} \to \mathrm{MonoInc}(a) \to \mathrm{MonoDec}(b) \to \varphi_{91}$
\item $\varphi_{251} :\equiv (\mathrm{MonoInc}(a) \to \varphi_{107}) \to (\mathrm{MonoInc}(a) \to \mathrm{MonoDec}(b) \to \varphi_{151}) \to \varphi_{204}$
\item $\varphi_{252} :\equiv (\mathrm{MonoInc}(a) \to \mathrm{MonoDec}(b) \to \varphi_{158}) \to (\mathrm{MonoInc}(a) \to \mathrm{MonoDec}(b) \to \varphi_{157}) \to \varphi_{184}$
\item $\varphi_{253} :\equiv (\mathrm{MonoInc}(a) \to \mathrm{MonoDec}(b) \to \varphi_{97}) \to \varphi_{144} \to \varphi_{143} \to \mathrm{MonoInc}(a) \to \mathrm{MonoDec}(b) \to \varphi_{56}$
\item $\varphi_{254} :\equiv (\mathrm{MonoInc}(a) \to \mathrm{MonoDec}(b) \to \varphi_{151}) \to (\mathrm{MonoInc}(a) \to \varphi_{69}) \to \varphi_{180} \to \varphi_{199}$
\item $\varphi_{255} :\equiv (\mathrm{MonoInc}(a) \to \varphi_{69}) \to (\mathrm{MonoInc}(a) \to \mathrm{MonoDec}(b) \to \varphi_{151}) \to \varphi_{183} \to \varphi_{203}$
\item $\varphi_{256} :\equiv (\mathrm{MonoInc}(a) \to \mathrm{MonoDec}(b) \to \varphi_{132}) \to \varphi_{169} \to (\mathrm{MonoInc}(a) \to \mathrm{MonoDec}(b) \to \varphi_{57}) \to \varphi_{232}$
\item $\varphi_{257} :\equiv (\mathrm{MonoInc}(a) \to n \in \mathbb{N} \to m \in \mathbb{N} \to n \leq m \to a_{n} \leq a_{m}) \to (\mathrm{MonoInc}(a) \to \mathrm{MonoInc}(a)) \to \varphi_{254}$
\item $\varphi_{258} :\equiv (\mathrm{MonoDec}(b) \to m \in \mathbb{N} \to n \in \mathbb{N} \to m \leq n \to b_{n} \leq b_{m}) \to \varphi_{2} \to \varphi_{255}$
\item $\varphi_{259} :\equiv (\forall v_{732}\,(\varphi_{19}) \leftrightarrow \forall v_{732}\,(\varphi_{215})) \to \forall v_{732}\,(\varphi_{215}) \to \forall v_{732}\,(\varphi_{19})$
\item $\varphi_{260} :\equiv (\forall v_{732}\,(\varphi_{19}) \leftrightarrow \forall v_{732}\,(\varphi_{215})) \to \forall v_{732}\,(\varphi_{19}) \to \forall v_{732}\,(\varphi_{215})$
\item $\varphi_{261} :\equiv (\forall v_{730}\,(\varphi_{21}) \leftrightarrow \forall v_{730}\,(\varphi_{216})) \to \forall v_{730}\,(\varphi_{216}) \to \forall v_{730}\,(\varphi_{21})$
\item $\varphi_{262} :\equiv (\forall v_{730}\,(\varphi_{21}) \leftrightarrow \forall v_{730}\,(\varphi_{216})) \to \forall v_{730}\,(\varphi_{21}) \to \forall v_{730}\,(\varphi_{216})$
\item $\varphi_{263} :\equiv \forall v_{731}\,(\forall v_{732}\,(\varphi_{215}) \to \forall v_{732}\,(\varphi_{19})) \to \varphi_{67} \to \forall v_{731}\,\forall v_{732}\,(\varphi_{19})$
\item $\varphi_{264} :\equiv \forall v_{731}\,(\forall v_{732}\,(\varphi_{19}) \to \forall v_{732}\,(\varphi_{215})) \to \forall v_{731}\,\forall v_{732}\,(\varphi_{19}) \to \varphi_{67}$
\item $\varphi_{265} :\equiv \forall v_{729}\,(\forall v_{730}\,(\varphi_{216}) \to \forall v_{730}\,(\varphi_{21})) \to \varphi_{68} \to \forall v_{729}\,\forall v_{730}\,(\varphi_{21})$
\item $\varphi_{266} :\equiv \forall v_{729}\,(\forall v_{730}\,(\varphi_{21}) \to \forall v_{730}\,(\varphi_{216})) \to \forall v_{729}\,\forall v_{730}\,(\varphi_{21}) \to \varphi_{68}$
\item $\varphi_{267} :\equiv (\forall v_{732}\,(\varphi_{215} \to \varphi_{19}) \to \forall v_{732}\,(\varphi_{215}) \to \forall v_{732}\,(\varphi_{19})) \to (\forall v_{732}\,(\varphi_{19}) \leftrightarrow \forall v_{732}\,(\varphi_{215}))$
\item $\varphi_{268} :\equiv (\forall v_{730}\,(\varphi_{216} \to \varphi_{21}) \to \forall v_{730}\,(\varphi_{216}) \to \forall v_{730}\,(\varphi_{21})) \to (\forall v_{730}\,(\varphi_{21}) \leftrightarrow \forall v_{730}\,(\varphi_{216}))$
\item $\varphi_{269} :\equiv (\mathrm{MonoInc}(a) \to \mathrm{MonoDec}(b) \to \varphi_{17}) \to \forall x\,(\varphi_{197}) \to (\forall x\,(\varphi_{197}) \to \varphi_{206}) \to \varphi_{234} \to \forall s\,(\varphi_{4}) \to \varphi_{238}$
\item $\varphi_{270} :\equiv (\mathrm{MonoInc}(a) \to \mathrm{MonoDec}(b) \to \varphi_{135}) \to \forall n\,(\varphi_{185}) \to (\forall n\,(\varphi_{185}) \to \varphi_{193}) \to \varphi_{226} \to \varphi_{23} \to \varphi_{241}$
\item $\varphi_{271} :\equiv \forall v_{732}\,(\varphi_{215} \to \varphi_{19}) \to (\forall v_{732}\,(\varphi_{19} \to \varphi_{215}) \to \forall v_{732}\,(\varphi_{19}) \to \forall v_{732}\,(\varphi_{215})) \to \varphi_{267}$
\item $\varphi_{272} :\equiv \forall v_{730}\,(\varphi_{216} \to \varphi_{21}) \to (\forall v_{730}\,(\varphi_{21} \to \varphi_{216}) \to \forall v_{730}\,(\varphi_{21}) \to \forall v_{730}\,(\varphi_{216})) \to \varphi_{268}$
\item $\varphi_{273} :\equiv \forall v_{731}\,(\forall v_{732}\,(\varphi_{215}) \to \forall v_{732}\,(\varphi_{19})) \to \varphi_{264} \to \varphi_{263} \to (\forall v_{731}\,\forall v_{732}\,(\varphi_{19}) \leftrightarrow \varphi_{67})$
\item $\varphi_{274} :\equiv \forall v_{729}\,(\forall v_{730}\,(\varphi_{216}) \to \forall v_{730}\,(\varphi_{21})) \to \varphi_{266} \to \varphi_{265} \to (\forall v_{729}\,\forall v_{730}\,(\varphi_{21}) \leftrightarrow \varphi_{68})$
\end{abbrevlist}
\end{fullproof}

\section{第3章 6節: Bolzano–Weierstrass の定理, Cauchy 列と実数の完備性}
% 自動生成 (prover/build.js): ArchInv  ノード数 167
\begin{fullproof}{[C24]}{$\forall \varepsilon \in \mathbb{R}\,(0 < \varepsilon \to \exists K \in \mathbb{N}\,(\underline{K}^{-1} < \varepsilon))$}{167}\phantomsection\label{tree:ArchInv}
\noindent$\Pi_{1}$:
\begin{ptree}
\AXC{$\forall n \in \mathbb{N}\,(\underline{n} \in \mathbb{R})$\lbl{[N4]}}
\AXC{$\forall n \in \mathbb{N}\,(\underline{n} \in \mathbb{R}) \to K \in \mathbb{N} \to \underline{K} \in \mathbb{R}$\lbl{A2}}
\BIC{$K \in \mathbb{N} \to \underline{K} \in \mathbb{R}$}
\AXC{$(K \in \mathbb{N} \to \underline{K} \in \mathbb{R}) \to \varphi_{38}$\lbl{A1}}
\BIC{$\varphi_{38}$}
\end{ptree}
\noindent$\Pi_{2}$:
\begin{ptree}
\AXC{$\varphi_{17}$\lbl{A1}}
\AXC{$\varphi_{17} \to \varepsilon \in \mathbb{R} \to 0 < \varepsilon \to K \in \mathbb{N} \land 1 < \underline{K} \cdot \varepsilon \to K \in \mathbb{N}$\lbl{A1}}
\BIC{$\varepsilon \in \mathbb{R} \to 0 < \varepsilon \to K \in \mathbb{N} \land 1 < \underline{K} \cdot \varepsilon \to K \in \mathbb{N}$}
\AXC{$\stackrel{\Pi_{1}}{\vdots}$}\noLine
\UIC{$\varphi_{38}$}
\BIC{$\varepsilon \in \mathbb{R} \to 0 < \varepsilon \to K \in \mathbb{N} \land 1 < \underline{K} \cdot \varepsilon \to \underline{K} \in \mathbb{R}$}
\end{ptree}
\noindent$\Pi_{3}$:
\begin{ptree}
\AXC{$\forall x\,\forall y\,(x \in \mathbb{R} \to y \in \mathbb{R} \to x \cdot y = y \cdot x)$\lbl{(R7)}}
\AXC{$\varphi_{30}$\lbl{A2}}
\BIC{$\forall y\,(\underline{K} \in \mathbb{R} \to y \in \mathbb{R} \to \underline{K} \cdot y = y \cdot \underline{K})$}
\AXC{$\varphi_{33}$\lbl{A2}}
\BIC{$\underline{K} \in \mathbb{R} \to \varepsilon \in \mathbb{R} \to \underline{K} \cdot \varepsilon = \varepsilon \cdot \underline{K}$}
\AXC{$\varphi_{53}$\lbl{A1}}
\BIC{$(\varepsilon \in \mathbb{R} \to 0 < \varepsilon \to K \in \mathbb{N} \land 1 < \underline{K} \cdot \varepsilon \to \underline{K} \in \mathbb{R}) \to \varphi_{42}$}
\end{ptree}
\noindent$\Pi_{4}$:
\begin{ptree}
\AXC{$\stackrel{\Pi_{2}}{\vdots}$}\noLine
\UIC{$\varepsilon \in \mathbb{R} \to 0 < \varepsilon \to K \in \mathbb{N} \land 1 < \underline{K} \cdot \varepsilon \to \underline{K} \in \mathbb{R}$}
\AXC{$\stackrel{\Pi_{3}}{\vdots}$}\noLine
\UIC{$(\varepsilon \in \mathbb{R} \to 0 < \varepsilon \to K \in \mathbb{N} \land 1 < \underline{K} \cdot \varepsilon \to \underline{K} \in \mathbb{R}) \to \varphi_{42}$}
\BIC{$\varphi_{42}$}
\end{ptree}
\noindent$\Pi_{5}$:
\begin{ptree}
\AXC{$\varepsilon \in \mathbb{R} \to 0 < \varepsilon \to K \in \mathbb{N} \land 1 < \underline{K} \cdot \varepsilon \to \varepsilon \in \mathbb{R}$\lbl{A1}}
\AXC{$\stackrel{\Pi_{4}}{\vdots}$}\noLine
\UIC{$\varphi_{42}$}
\BIC{$\varepsilon \in \mathbb{R} \to 0 < \varepsilon \to K \in \mathbb{N} \land 1 < \underline{K} \cdot \varepsilon \to \underline{K} \cdot \varepsilon = \varepsilon \cdot \underline{K}$}
\AXC{$\varphi_{55}$\lbl{A1}}
\BIC{$\varphi_{48}$}
\end{ptree}
\noindent$\Pi_{6}$:
\begin{ptree}
\AXC{$\varphi_{17}$\lbl{A1}}
\AXC{$\varphi_{17} \to \varepsilon \in \mathbb{R} \to 0 < \varepsilon \to K \in \mathbb{N} \land 1 < \underline{K} \cdot \varepsilon \to 1 < \underline{K} \cdot \varepsilon$\lbl{A1}}
\BIC{$\varepsilon \in \mathbb{R} \to 0 < \varepsilon \to K \in \mathbb{N} \land 1 < \underline{K} \cdot \varepsilon \to 1 < \underline{K} \cdot \varepsilon$}
\AXC{$\stackrel{\Pi_{5}}{\vdots}$}\noLine
\UIC{$\varphi_{48}$}
\BIC{$\varphi_{40}$}
\end{ptree}
\noindent$\Pi_{7}$:
\begin{ptree}
\AXC{$\forall n \in \mathbb{N}\,(0 < \underline{n})$\lbl{[N15]}}
\AXC{$\forall n \in \mathbb{N}\,(0 < \underline{n}) \to K \in \mathbb{N} \to 0 < \underline{K}$\lbl{A2}}
\BIC{$K \in \mathbb{N} \to 0 < \underline{K}$}
\AXC{$(K \in \mathbb{N} \to 0 < \underline{K}) \to \varphi_{36}$\lbl{A1}}
\BIC{$\varphi_{36}$}
\end{ptree}
\noindent$\Pi_{8}$:
\begin{ptree}
\AXC{$\varphi_{17}$\lbl{A1}}
\AXC{$\varphi_{17} \to \varepsilon \in \mathbb{R} \to 0 < \varepsilon \to K \in \mathbb{N} \land 1 < \underline{K} \cdot \varepsilon \to K \in \mathbb{N}$\lbl{A1}}
\BIC{$\varepsilon \in \mathbb{R} \to 0 < \varepsilon \to K \in \mathbb{N} \land 1 < \underline{K} \cdot \varepsilon \to K \in \mathbb{N}$}
\AXC{$\stackrel{\Pi_{7}}{\vdots}$}\noLine
\UIC{$\varphi_{36}$}
\BIC{$\varepsilon \in \mathbb{R} \to 0 < \varepsilon \to K \in \mathbb{N} \land 1 < \underline{K} \cdot \varepsilon \to 0 < \underline{K}$}
\end{ptree}
\noindent$\Pi_{9}$:
\begin{ptree}
\AXC{$1 \in \mathbb{R}$\lbl{(R1)}}
\AXC{$\varphi_{13}$\lbl{[E24]}}
\AXC{$\varphi_{13} \to \varphi_{11}$\lbl{A2}}
\BIC{$\varphi_{11}$}
\AXC{$\varphi_{11} \to \varphi_{29}$\lbl{A2}}
\BIC{$\varphi_{29}$}
\AXC{$\varphi_{29} \to 1 \in \mathbb{R} \to \varphi_{25}$\lbl{A2}}
\BIC{$1 \in \mathbb{R} \to \varphi_{25}$}
\AXC{$(1 \in \mathbb{R} \to \varphi_{25}) \to \varphi_{56}$\lbl{A1}}
\BIC{$\varphi_{56}$}
\BIC{$(\varepsilon \in \mathbb{R} \to 0 < \varepsilon \to K \in \mathbb{N} \land 1 < \underline{K} \cdot \varepsilon \to \underline{K} \in \mathbb{R}) \to \varphi_{54}$}
\end{ptree}
\noindent$\Pi_{10}$:
\begin{ptree}
\AXC{$\stackrel{\Pi_{2}}{\vdots}$}\noLine
\UIC{$\varepsilon \in \mathbb{R} \to 0 < \varepsilon \to K \in \mathbb{N} \land 1 < \underline{K} \cdot \varepsilon \to \underline{K} \in \mathbb{R}$}
\AXC{$\stackrel{\Pi_{9}}{\vdots}$}\noLine
\UIC{$(\varepsilon \in \mathbb{R} \to 0 < \varepsilon \to K \in \mathbb{N} \land 1 < \underline{K} \cdot \varepsilon \to \underline{K} \in \mathbb{R}) \to \varphi_{54}$}
\BIC{$\varphi_{54}$}
\end{ptree}
\noindent$\Pi_{11}$:
\begin{ptree}
\AXC{$\stackrel{\Pi_{8}}{\vdots}$}\noLine
\UIC{$\varepsilon \in \mathbb{R} \to 0 < \varepsilon \to K \in \mathbb{N} \land 1 < \underline{K} \cdot \varepsilon \to 0 < \underline{K}$}
\AXC{$\varepsilon \in \mathbb{R} \to 0 < \varepsilon \to K \in \mathbb{N} \land 1 < \underline{K} \cdot \varepsilon \to \varepsilon \in \mathbb{R}$\lbl{A1}}
\AXC{$\stackrel{\Pi_{10}}{\vdots}$}\noLine
\UIC{$\varphi_{54}$}
\BIC{$(\varepsilon \in \mathbb{R} \to 0 < \varepsilon \to K \in \mathbb{N} \land 1 < \underline{K} \cdot \varepsilon \to 0 < \underline{K}) \to \varphi_{43}$}
\BIC{$\varphi_{43}$}
\end{ptree}
\noindent$\Pi_{12}$:
\begin{ptree}
\AXC{$y = z \to x = y \to x = z$\lbl{A6}}
\AXC{$(y = z \to x = y \to x = z) \to x = y \to y = z \to x = z$\lbl{A1}}
\BIC{$x = y \to y = z \to x = z$}
\UIC{$\forall z\,(x = y \to y = z \to x = z)$}
\UIC{$\forall y\,\forall z\,(x = y \to y = z \to x = z)$}
\UIC{$\forall x\,\forall y\,\forall z\,(x = y \to y = z \to x = z)$}
\AXC{$\varphi_{6}$\lbl{A2}}
\BIC{$\forall y\,\forall z\,(1 \cdot \underline{K}^{-1} = y \to y = z \to 1 \cdot \underline{K}^{-1} = z)$}
\end{ptree}
\noindent$\Pi_{13}$:
\begin{ptree}
\AXC{$\stackrel{\Pi_{12}}{\vdots}$}\noLine
\UIC{$\forall y\,\forall z\,(1 \cdot \underline{K}^{-1} = y \to y = z \to 1 \cdot \underline{K}^{-1} = z)$}
\AXC{$\forall y\,\forall z\,(1 \cdot \underline{K}^{-1} = y \to y = z \to 1 \cdot \underline{K}^{-1} = z) \to \varphi_{8}$\lbl{A2}}
\BIC{$\varphi_{8}$}
\AXC{$\varphi_{8} \to \varphi_{27}$\lbl{A2}}
\BIC{$\varphi_{27}$}
\end{ptree}
\noindent$\Pi_{14}$:
\begin{ptree}
\AXC{$x = x$\lbl{A5}}
\AXC{$x = y \to x = x \to y = x$\lbl{A6}}
\AXC{$(x = y \to x = x \to y = x) \to x = x \to x = y \to y = x$\lbl{A1}}
\BIC{$x = x \to x = y \to y = x$}
\BIC{$x = y \to y = x$}
\UIC{$\forall y\,(x = y \to y = x)$}
\UIC{$\forall x\,\forall y\,(x = y \to y = x)$}
\end{ptree}
\noindent$\Pi_{15}$:
\begin{ptree}
\AXC{$\stackrel{\Pi_{14}}{\vdots}$}\noLine
\UIC{$\forall x\,\forall y\,(x = y \to y = x)$}
\AXC{$\forall x\,\forall y\,(x = y \to y = x) \to \forall y\,(\underline{K} = y \to y = \underline{K})$\lbl{A2}}
\BIC{$\forall y\,(\underline{K} = y \to y = \underline{K})$}
\end{ptree}
\noindent$\Pi_{16}$:
\begin{ptree}
\AXC{$\forall x\,\forall y\,(x < y \to x \neq y)$\lbl{[F26]}}
\AXC{$\forall x\,\forall y\,(x < y \to x \neq y) \to \forall y\,(0 < y \to 0 \neq y)$\lbl{A2}}
\BIC{$\forall y\,(0 < y \to 0 \neq y)$}
\AXC{$\forall y\,(0 < y \to 0 \neq y) \to 0 < \underline{K} \to 0 \neq \underline{K}$\lbl{A2}}
\BIC{$0 < \underline{K} \to 0 \neq \underline{K}$}
\end{ptree}
\noindent$\Pi_{17}$:
\begin{ptree}
\AXC{$\stackrel{\Pi_{8}}{\vdots}$}\noLine
\UIC{$\varepsilon \in \mathbb{R} \to 0 < \varepsilon \to K \in \mathbb{N} \land 1 < \underline{K} \cdot \varepsilon \to 0 < \underline{K}$}
\AXC{$\stackrel{\Pi_{16}}{\vdots}$}\noLine
\UIC{$0 < \underline{K} \to 0 \neq \underline{K}$}
\AXC{$(0 < \underline{K} \to 0 \neq \underline{K}) \to \varphi_{37}$\lbl{A1}}
\BIC{$\varphi_{37}$}
\BIC{$\varepsilon \in \mathbb{R} \to 0 < \varepsilon \to K \in \mathbb{N} \land 1 < \underline{K} \cdot \varepsilon \to 0 \neq \underline{K}$}
\end{ptree}
\noindent$\Pi_{18}$:
\begin{ptree}
\AXC{$\stackrel{\Pi_{17}}{\vdots}$}\noLine
\UIC{$\varepsilon \in \mathbb{R} \to 0 < \varepsilon \to K \in \mathbb{N} \land 1 < \underline{K} \cdot \varepsilon \to 0 \neq \underline{K}$}
\AXC{$(\varepsilon \in \mathbb{R} \to 0 < \varepsilon \to K \in \mathbb{N} \land 1 < \underline{K} \cdot \varepsilon \to 0 \neq \underline{K}) \to \varphi_{20}$\lbl{A1}}
\BIC{$\varphi_{20}$}
\end{ptree}
\noindent$\Pi_{19}$:
\begin{ptree}
\AXC{$\stackrel{\Pi_{2}}{\vdots}$}\noLine
\UIC{$\varepsilon \in \mathbb{R} \to 0 < \varepsilon \to K \in \mathbb{N} \land 1 < \underline{K} \cdot \varepsilon \to \underline{K} \in \mathbb{R}$}
\AXC{$\forall x \in \mathbb{R}\,(x \neq 0 \to x^{-1} \in \mathbb{R})$\lbl{(R1)}}
\AXC{$\varphi_{21}$\lbl{A2}}
\BIC{$\underline{K} \in \mathbb{R} \to \underline{K} \neq 0 \to \underline{K}^{-1} \in \mathbb{R}$}
\AXC{$\varphi_{51}$\lbl{A1}}
\BIC{$(\varepsilon \in \mathbb{R} \to 0 < \varepsilon \to K \in \mathbb{N} \land 1 < \underline{K} \cdot \varepsilon \to \underline{K} \in \mathbb{R}) \to \varphi_{39}$}
\BIC{$\varphi_{39}$}
\end{ptree}
\noindent$\Pi_{20}$:
\begin{ptree}
\AXC{$\stackrel{\Pi_{15}}{\vdots}$}\noLine
\UIC{$\forall y\,(\underline{K} = y \to y = \underline{K})$}
\AXC{$\forall y\,(\underline{K} = y \to y = \underline{K}) \to \underline{K} = 0 \to 0 = \underline{K}$\lbl{A2}}
\BIC{$\underline{K} = 0 \to 0 = \underline{K}$}
\AXC{$\stackrel{\Pi_{18}}{\vdots}$}\noLine
\UIC{$\varphi_{20}$}
\BIC{$\varepsilon \in \mathbb{R} \to 0 < \varepsilon \to K \in \mathbb{N} \land 1 < \underline{K} \cdot \varepsilon \to \underline{K} \neq 0$}
\AXC{$\stackrel{\Pi_{19}}{\vdots}$}\noLine
\UIC{$\varphi_{39}$}
\BIC{$\varepsilon \in \mathbb{R} \to 0 < \varepsilon \to K \in \mathbb{N} \land 1 < \underline{K} \cdot \varepsilon \to \underline{K}^{-1} \in \mathbb{R}$}
\end{ptree}
\noindent$\Pi_{21}$:
\begin{ptree}
\AXC{$\forall x \in \mathbb{R}\,(x \cdot 1 = x)$\lbl{(R8)}}
\AXC{$\varphi_{7}$\lbl{A2}}
\BIC{$\underline{K}^{-1} \in \mathbb{R} \to \underline{K}^{-1} \cdot 1 = \underline{K}^{-1}$}
\AXC{$(\underline{K}^{-1} \in \mathbb{R} \to \underline{K}^{-1} \cdot 1 = \underline{K}^{-1}) \to \varphi_{45}$\lbl{A1}}
\BIC{$\varphi_{45}$}
\end{ptree}
\noindent$\Pi_{22}$:
\begin{ptree}
\AXC{$1 \in \mathbb{R}$\lbl{(R1)}}
\AXC{$\forall x\,\forall y\,(x \in \mathbb{R} \to y \in \mathbb{R} \to x \cdot y = y \cdot x)$\lbl{(R7)}}
\AXC{$\varphi_{15}$\lbl{A2}}
\BIC{$\forall y\,(1 \in \mathbb{R} \to y \in \mathbb{R} \to 1 \cdot y = y \cdot 1)$}
\AXC{$\varphi_{31}$\lbl{A2}}
\BIC{$1 \in \mathbb{R} \to \underline{K}^{-1} \in \mathbb{R} \to 1 \cdot \underline{K}^{-1} = \underline{K}^{-1} \cdot 1$}
\AXC{$\varphi_{50}$\lbl{A1}}
\BIC{$1 \in \mathbb{R} \to \varphi_{46}$}
\BIC{$\varphi_{46}$}
\end{ptree}
\noindent$\Pi_{23}$:
\begin{ptree}
\AXC{$\stackrel{\Pi_{20}}{\vdots}$}\noLine
\UIC{$\varepsilon \in \mathbb{R} \to 0 < \varepsilon \to K \in \mathbb{N} \land 1 < \underline{K} \cdot \varepsilon \to \underline{K}^{-1} \in \mathbb{R}$}
\AXC{$\stackrel{\Pi_{22}}{\vdots}$}\noLine
\UIC{$\varphi_{46}$}
\BIC{$\varphi_{23}$}
\AXC{$\varphi_{23} \to \varphi_{14} \to \varphi_{27} \to \varphi_{16}$\lbl{A1}}
\BIC{$\varphi_{14} \to \varphi_{27} \to \varphi_{16}$}
\end{ptree}
\noindent$\Pi_{24}$:
\begin{ptree}
\AXC{$\stackrel{\Pi_{13}}{\vdots}$}\noLine
\UIC{$\varphi_{27}$}
\AXC{$\stackrel{\Pi_{20}}{\vdots}$}\noLine
\UIC{$\varepsilon \in \mathbb{R} \to 0 < \varepsilon \to K \in \mathbb{N} \land 1 < \underline{K} \cdot \varepsilon \to \underline{K}^{-1} \in \mathbb{R}$}
\AXC{$\stackrel{\Pi_{21}}{\vdots}$}\noLine
\UIC{$\varphi_{45}$}
\BIC{$\varphi_{14}$}
\AXC{$\stackrel{\Pi_{23}}{\vdots}$}\noLine
\UIC{$\varphi_{14} \to \varphi_{27} \to \varphi_{16}$}
\BIC{$\varphi_{27} \to \varphi_{16}$}
\BIC{$\varphi_{16}$}
\AXC{$\varphi_{16} \to \varphi_{49}$\lbl{A1}}
\BIC{$\varphi_{49}$}
\end{ptree}
\noindent$\Pi_{25}$:
\begin{ptree}
\AXC{$z = w \to 1 < z \to 1 < w$\lbl{A6}}
\UIC{$\forall w\,(z = w \to 1 < z \to 1 < w)$}
\UIC{$\forall z\,\forall w\,(z = w \to 1 < z \to 1 < w)$}
\AXC{$\varphi_{5}$\lbl{A2}}
\BIC{$\forall w\,(\underline{K} \cdot \varepsilon = w \to 1 < \underline{K} \cdot \varepsilon \to 1 < w)$}
\AXC{$\varphi_{32}$\lbl{A2}}
\BIC{$\underline{K} \cdot \varepsilon = \varepsilon \cdot \underline{K} \to 1 < \underline{K} \cdot \varepsilon \to 1 < \varepsilon \cdot \underline{K}$}
\AXC{$\stackrel{\Pi_{6}}{\vdots}$}\noLine
\UIC{$\varphi_{40}$}
\BIC{$\varepsilon \in \mathbb{R} \to 0 < \varepsilon \to K \in \mathbb{N} \land 1 < \underline{K} \cdot \varepsilon \to 1 < \varepsilon \cdot \underline{K}$}
\AXC{$\stackrel{\Pi_{11}}{\vdots}$}\noLine
\UIC{$\varphi_{43}$}
\BIC{$\varepsilon \in \mathbb{R} \to 0 < \varepsilon \to K \in \mathbb{N} \land 1 < \underline{K} \cdot \varepsilon \to 1 \cdot \underline{K}^{-1} < \varepsilon$}
\AXC{$\stackrel{\Pi_{24}}{\vdots}$}\noLine
\UIC{$\varphi_{49}$}
\BIC{$\varphi_{41}$}
\end{ptree}
\noindent$\Pi_{26}$:
\begin{ptree}
\AXC{$\varphi_{17}$\lbl{A1}}
\AXC{$\varphi_{17} \to \varepsilon \in \mathbb{R} \to 0 < \varepsilon \to K \in \mathbb{N} \land 1 < \underline{K} \cdot \varepsilon \to K \in \mathbb{N}$\lbl{A1}}
\BIC{$\varepsilon \in \mathbb{R} \to 0 < \varepsilon \to K \in \mathbb{N} \land 1 < \underline{K} \cdot \varepsilon \to K \in \mathbb{N}$}
\AXC{$(\varepsilon \in \mathbb{R} \to 0 < \varepsilon \to K \in \mathbb{N} \land 1 < \underline{K} \cdot \varepsilon \to K \in \mathbb{N}) \to \varphi_{44}$\lbl{A1}}
\BIC{$\varphi_{44}$}
\end{ptree}
\noindent$\Pi_{27}$:
\begin{ptree}
\AXC{$\varphi_{2}$\lbl{A2}}
\AXC{$z = w \to z < \varepsilon \to w < \varepsilon$\lbl{A6}}
\UIC{$\forall w\,(z = w \to z < \varepsilon \to w < \varepsilon)$}
\UIC{$\forall z\,\forall w\,(z = w \to z < \varepsilon \to w < \varepsilon)$}
\AXC{$\varphi_{19}$\lbl{A2}}
\BIC{$\forall w\,(1 \cdot \underline{K}^{-1} = w \to 1 \cdot \underline{K}^{-1} < \varepsilon \to w < \varepsilon)$}
\AXC{$\varphi_{35}$\lbl{A2}}
\BIC{$1 \cdot \underline{K}^{-1} = \underline{K}^{-1} \to 1 \cdot \underline{K}^{-1} < \varepsilon \to \underline{K}^{-1} < \varepsilon$}
\AXC{$\stackrel{\Pi_{25}}{\vdots}$}\noLine
\UIC{$\varphi_{41}$}
\BIC{$\varepsilon \in \mathbb{R} \to 0 < \varepsilon \to K \in \mathbb{N} \land 1 < \underline{K} \cdot \varepsilon \to \underline{K}^{-1} < \varepsilon$}
\AXC{$\stackrel{\Pi_{26}}{\vdots}$}\noLine
\UIC{$\varphi_{44}$}
\BIC{$\varphi_{18}$}
\AXC{$\varphi_{18} \to \varphi_{2} \to \varphi_{22}$\lbl{A1}}
\BIC{$\varphi_{2} \to \varphi_{22}$}
\BIC{$\varphi_{22}$}
\end{ptree}
\noindent$\Pi_{28}$:
\begin{ptree}
\AXC{$1 \in \mathbb{R}$\lbl{(R1)}}
\AXC{$\varepsilon \in \mathbb{R} \to 0 < \varepsilon \to \varepsilon \in \mathbb{R}$\lbl{A1}}
\AXC{$\varphi_{9}$\lbl{命題3.2}}
\AXC{$\varphi_{9} \to \varphi_{12}$\lbl{A2}}
\BIC{$\varphi_{12}$}
\AXC{$\varphi_{12} \to \varphi_{10}$\lbl{A2}}
\BIC{$\varphi_{10}$}
\AXC{$\varphi_{10} \to \varphi_{47}$\lbl{A1}}
\BIC{$\varphi_{47}$}
\BIC{$1 \in \mathbb{R} \to (\varepsilon \in \mathbb{R} \to 0 < \varepsilon \to 0 < \varepsilon) \to \varphi_{24}$}
\BIC{$(\varepsilon \in \mathbb{R} \to 0 < \varepsilon \to 0 < \varepsilon) \to \varphi_{24}$}
\end{ptree}
\noindent$\Pi_{29}$:
\begin{ptree}
\AXC{$0 < 1$\lbl{[F28]}}
\AXC{$0 < 1 \to \varepsilon \in \mathbb{R} \to 0 < \varepsilon \to 0 < 1$\lbl{A1}}
\BIC{$\varepsilon \in \mathbb{R} \to 0 < \varepsilon \to 0 < 1$}
\AXC{$\varepsilon \in \mathbb{R} \to 0 < \varepsilon \to 0 < \varepsilon$\lbl{A1}}
\AXC{$\stackrel{\Pi_{28}}{\vdots}$}\noLine
\UIC{$(\varepsilon \in \mathbb{R} \to 0 < \varepsilon \to 0 < \varepsilon) \to \varphi_{24}$}
\BIC{$\varphi_{24}$}
\BIC{$\varepsilon \in \mathbb{R} \to 0 < \varepsilon \to \exists n \in \mathbb{N}\,(1 < \underline{n} \cdot \varepsilon)$}
\end{ptree}
\noindent$\Pi_{30}$:
\begin{ptree}
\AXC{$\varphi_{34}$\lbl{A4}}
\AXC{$\forall K\,(\varphi_{26}) \to \varphi_{28}$\lbl{A3}}
\AXC{$\stackrel{\Pi_{27}}{\vdots}$}\noLine
\UIC{$\varphi_{22}$}
\AXC{$\varphi_{22} \to \varphi_{26}$\lbl{A1}}
\BIC{$\varphi_{26}$}
\UIC{$\forall K\,(\varphi_{26})$}
\AXC{$\stackrel{\Pi_{29}}{\vdots}$}\noLine
\UIC{$\varepsilon \in \mathbb{R} \to 0 < \varepsilon \to \exists n \in \mathbb{N}\,(1 < \underline{n} \cdot \varepsilon)$}
\AXC{$\varphi_{57}$\lbl{A1}}
\BIC{$\forall K\,(\varphi_{26}) \to (\forall K\,(\varphi_{26}) \to \varphi_{28}) \to \varphi_{34} \to \forall n\,(\varphi_{1}) \to \varphi_{52}$}
\BIC{$(\forall K\,(\varphi_{26}) \to \varphi_{28}) \to \varphi_{34} \to \forall n\,(\varphi_{1}) \to \varphi_{52}$}
\BIC{$\varphi_{34} \to \forall n\,(\varphi_{1}) \to \varphi_{52}$}
\BIC{$\forall n\,(\varphi_{1}) \to \varphi_{52}$}
\end{ptree}
\noindent 主木:
\begin{ptree}
\AXC{$\varphi_{3}$\lbl{A4}}
\AXC{$\forall n\,(\varphi_{1}) \to \varphi_{4}$\lbl{A3}}
\AXC{$\varphi_{1}$\lbl{A2}}
\UIC{$\forall n\,(\varphi_{1})$}
\AXC{$\stackrel{\Pi_{30}}{\vdots}$}\noLine
\UIC{$\forall n\,(\varphi_{1}) \to \varphi_{52}$}
\BIC{$\varphi_{52}$}
\BIC{$\varphi_{3} \to \varepsilon \in \mathbb{R} \to 0 < \varepsilon \to \exists K \in \mathbb{N}\,(\underline{K}^{-1} < \varepsilon)$}
\BIC{$\varepsilon \in \mathbb{R} \to 0 < \varepsilon \to \exists K \in \mathbb{N}\,(\underline{K}^{-1} < \varepsilon)$}
\UIC{$\forall \varepsilon \in \mathbb{R}\,(0 < \varepsilon \to \exists K \in \mathbb{N}\,(\underline{K}^{-1} < \varepsilon))$}
\end{ptree}
\begin{abbrevlist}
\item $\varphi_{1} :\equiv \forall K\,(\neg (K \in \mathbb{N} \land 1 < \underline{K} \cdot \varepsilon)) \to \neg (n \in \mathbb{N} \land 1 < \underline{n} \cdot \varepsilon)$
\item $\varphi_{2} :\equiv \forall K\,(\neg (K \in \mathbb{N} \land \underline{K}^{-1} < \varepsilon)) \to \neg (K \in \mathbb{N} \land \underline{K}^{-1} < \varepsilon)$
\item $\varphi_{3} :\equiv \forall K\,(\neg (K \in \mathbb{N} \land 1 < \underline{K} \cdot \varepsilon)) \to \forall n\,\forall K\,(\neg (K \in \mathbb{N} \land 1 < \underline{K} \cdot \varepsilon))$
\item $\varphi_{4} :\equiv \forall n\,\forall K\,(\neg (K \in \mathbb{N} \land 1 < \underline{K} \cdot \varepsilon)) \to \forall n\,(\neg (n \in \mathbb{N} \land 1 < \underline{n} \cdot \varepsilon))$
\item $\varphi_{5} :\equiv \forall z\,\forall w\,(z = w \to 1 < z \to 1 < w) \to \forall w\,(\underline{K} \cdot \varepsilon = w \to 1 < \underline{K} \cdot \varepsilon \to 1 < w)$
\item $\varphi_{6} :\equiv \forall x\,\forall y\,\forall z\,(x = y \to y = z \to x = z) \to \forall y\,\forall z\,(1 \cdot \underline{K}^{-1} = y \to y = z \to 1 \cdot \underline{K}^{-1} = z)$
\item $\varphi_{7} :\equiv \forall x \in \mathbb{R}\,(x \cdot 1 = x) \to \underline{K}^{-1} \in \mathbb{R} \to \underline{K}^{-1} \cdot 1 = \underline{K}^{-1}$
\item $\varphi_{8} :\equiv \forall z\,(1 \cdot \underline{K}^{-1} = \underline{K}^{-1} \cdot 1 \to \underline{K}^{-1} \cdot 1 = z \to 1 \cdot \underline{K}^{-1} = z)$
\item $\varphi_{9} :\equiv \forall a\,\forall b\,(a \in \mathbb{R} \to b \in \mathbb{R} \to 0 < a \to 0 < b \to \exists n \in \mathbb{N}\,(b < \underline{n} \cdot a))$
\item $\varphi_{10} :\equiv \varepsilon \in \mathbb{R} \to 1 \in \mathbb{R} \to 0 < \varepsilon \to 0 < 1 \to \exists n \in \mathbb{N}\,(1 < \underline{n} \cdot \varepsilon)$
\item $\varphi_{11} :\equiv \forall y\,\forall z\,(1 \in \mathbb{R} \to y \in \mathbb{R} \to z \in \mathbb{R} \to 0 < y \to 1 < z \cdot y \to 1 \cdot y^{-1} < z)$
\item $\varphi_{12} :\equiv \forall b\,(\varepsilon \in \mathbb{R} \to b \in \mathbb{R} \to 0 < \varepsilon \to 0 < b \to \exists n \in \mathbb{N}\,(b < \underline{n} \cdot \varepsilon))$
\item $\varphi_{13} :\equiv \forall x\,\forall y\,\forall z\,(x \in \mathbb{R} \to y \in \mathbb{R} \to z \in \mathbb{R} \to 0 < y \to x < z \cdot y \to x \cdot y^{-1} < z)$
\item $\varphi_{14} :\equiv \varepsilon \in \mathbb{R} \to 0 < \varepsilon \to K \in \mathbb{N} \land 1 < \underline{K} \cdot \varepsilon \to \underline{K}^{-1} \cdot 1 = \underline{K}^{-1}$
\item $\varphi_{15} :\equiv \forall x\,\forall y\,(x \in \mathbb{R} \to y \in \mathbb{R} \to x \cdot y = y \cdot x) \to \forall y\,(1 \in \mathbb{R} \to y \in \mathbb{R} \to 1 \cdot y = y \cdot 1)$
\item $\varphi_{16} :\equiv \varepsilon \in \mathbb{R} \to 0 < \varepsilon \to K \in \mathbb{N} \land 1 < \underline{K} \cdot \varepsilon \to 1 \cdot \underline{K}^{-1} = \underline{K}^{-1}$
\item $\varphi_{17} :\equiv \varepsilon \in \mathbb{R} \to 0 < \varepsilon \to K \in \mathbb{N} \land 1 < \underline{K} \cdot \varepsilon \to K \in \mathbb{N} \land 1 < \underline{K} \cdot \varepsilon$
\item $\varphi_{18} :\equiv \varepsilon \in \mathbb{R} \to 0 < \varepsilon \to K \in \mathbb{N} \land 1 < \underline{K} \cdot \varepsilon \to K \in \mathbb{N} \land \underline{K}^{-1} < \varepsilon$
\item $\varphi_{19} :\equiv \forall z\,\forall w\,(z = w \to z < \varepsilon \to w < \varepsilon) \to \forall w\,(1 \cdot \underline{K}^{-1} = w \to 1 \cdot \underline{K}^{-1} < \varepsilon \to w < \varepsilon)$
\item $\varphi_{20} :\equiv (\underline{K} = 0 \to 0 = \underline{K}) \to \varepsilon \in \mathbb{R} \to 0 < \varepsilon \to K \in \mathbb{N} \land 1 < \underline{K} \cdot \varepsilon \to \underline{K} \neq 0$
\item $\varphi_{21} :\equiv \forall x \in \mathbb{R}\,(x \neq 0 \to x^{-1} \in \mathbb{R}) \to \underline{K} \in \mathbb{R} \to \underline{K} \neq 0 \to \underline{K}^{-1} \in \mathbb{R}$
\item $\varphi_{22} :\equiv \varepsilon \in \mathbb{R} \to 0 < \varepsilon \to K \in \mathbb{N} \land 1 < \underline{K} \cdot \varepsilon \to \exists K \in \mathbb{N}\,(\underline{K}^{-1} < \varepsilon)$
\item $\varphi_{23} :\equiv \varepsilon \in \mathbb{R} \to 0 < \varepsilon \to K \in \mathbb{N} \land 1 < \underline{K} \cdot \varepsilon \to 1 \cdot \underline{K}^{-1} = \underline{K}^{-1} \cdot 1$
\item $\varphi_{24} :\equiv (\varepsilon \in \mathbb{R} \to 0 < \varepsilon \to 0 < 1) \to \varepsilon \in \mathbb{R} \to 0 < \varepsilon \to \exists n \in \mathbb{N}\,(1 < \underline{n} \cdot \varepsilon)$
\item $\varphi_{25} :\equiv \underline{K} \in \mathbb{R} \to \varepsilon \in \mathbb{R} \to 0 < \underline{K} \to 1 < \varepsilon \cdot \underline{K} \to 1 \cdot \underline{K}^{-1} < \varepsilon$
\item $\varphi_{26} :\equiv \varepsilon \in \mathbb{R} \land 0 < \varepsilon \land \neg \exists K \in \mathbb{N}\,(\underline{K}^{-1} < \varepsilon) \to \neg (K \in \mathbb{N} \land 1 < \underline{K} \cdot \varepsilon)$
\item $\varphi_{27} :\equiv 1 \cdot \underline{K}^{-1} = \underline{K}^{-1} \cdot 1 \to \underline{K}^{-1} \cdot 1 = \underline{K}^{-1} \to 1 \cdot \underline{K}^{-1} = \underline{K}^{-1}$
\item $\varphi_{28} :\equiv \forall K\,(\varepsilon \in \mathbb{R} \land 0 < \varepsilon \land \neg \exists K \in \mathbb{N}\,(\underline{K}^{-1} < \varepsilon)) \to \forall K\,(\neg (K \in \mathbb{N} \land 1 < \underline{K} \cdot \varepsilon))$
\item $\varphi_{29} :\equiv \forall z\,(1 \in \mathbb{R} \to \underline{K} \in \mathbb{R} \to z \in \mathbb{R} \to 0 < \underline{K} \to 1 < z \cdot \underline{K} \to 1 \cdot \underline{K}^{-1} < z)$
\item $\varphi_{30} :\equiv \forall x\,\forall y\,(x \in \mathbb{R} \to y \in \mathbb{R} \to x \cdot y = y \cdot x) \to \forall y\,(\underline{K} \in \mathbb{R} \to y \in \mathbb{R} \to \underline{K} \cdot y = y \cdot \underline{K})$
\item $\varphi_{31} :\equiv \forall y\,(1 \in \mathbb{R} \to y \in \mathbb{R} \to 1 \cdot y = y \cdot 1) \to 1 \in \mathbb{R} \to \underline{K}^{-1} \in \mathbb{R} \to 1 \cdot \underline{K}^{-1} = \underline{K}^{-1} \cdot 1$
\item $\varphi_{32} :\equiv \forall w\,(\underline{K} \cdot \varepsilon = w \to 1 < \underline{K} \cdot \varepsilon \to 1 < w) \to \underline{K} \cdot \varepsilon = \varepsilon \cdot \underline{K} \to 1 < \underline{K} \cdot \varepsilon \to 1 < \varepsilon \cdot \underline{K}$
\item $\varphi_{33} :\equiv \forall y\,(\underline{K} \in \mathbb{R} \to y \in \mathbb{R} \to \underline{K} \cdot y = y \cdot \underline{K}) \to \underline{K} \in \mathbb{R} \to \varepsilon \in \mathbb{R} \to \underline{K} \cdot \varepsilon = \varepsilon \cdot \underline{K}$
\item $\varphi_{34} :\equiv \varepsilon \in \mathbb{R} \land 0 < \varepsilon \land \neg \exists K \in \mathbb{N}\,(\underline{K}^{-1} < \varepsilon) \to \forall K\,(\varepsilon \in \mathbb{R} \land 0 < \varepsilon \land \neg \exists K \in \mathbb{N}\,(\underline{K}^{-1} < \varepsilon))$
\item $\varphi_{35} :\equiv \forall w\,(1 \cdot \underline{K}^{-1} = w \to 1 \cdot \underline{K}^{-1} < \varepsilon \to w < \varepsilon) \to 1 \cdot \underline{K}^{-1} = \underline{K}^{-1} \to 1 \cdot \underline{K}^{-1} < \varepsilon \to \underline{K}^{-1} < \varepsilon$
\item $\varphi_{36} :\equiv (\varepsilon \in \mathbb{R} \to 0 < \varepsilon \to K \in \mathbb{N} \land 1 < \underline{K} \cdot \varepsilon \to K \in \mathbb{N}) \to \varepsilon \in \mathbb{R} \to 0 < \varepsilon \to K \in \mathbb{N} \land 1 < \underline{K} \cdot \varepsilon \to 0 < \underline{K}$
\item $\varphi_{37} :\equiv (\varepsilon \in \mathbb{R} \to 0 < \varepsilon \to K \in \mathbb{N} \land 1 < \underline{K} \cdot \varepsilon \to 0 < \underline{K}) \to \varepsilon \in \mathbb{R} \to 0 < \varepsilon \to K \in \mathbb{N} \land 1 < \underline{K} \cdot \varepsilon \to 0 \neq \underline{K}$
\item $\varphi_{38} :\equiv (\varepsilon \in \mathbb{R} \to 0 < \varepsilon \to K \in \mathbb{N} \land 1 < \underline{K} \cdot \varepsilon \to K \in \mathbb{N}) \to \varepsilon \in \mathbb{R} \to 0 < \varepsilon \to K \in \mathbb{N} \land 1 < \underline{K} \cdot \varepsilon \to \underline{K} \in \mathbb{R}$
\item $\varphi_{39} :\equiv (\varepsilon \in \mathbb{R} \to 0 < \varepsilon \to K \in \mathbb{N} \land 1 < \underline{K} \cdot \varepsilon \to \underline{K} \neq 0) \to \varepsilon \in \mathbb{R} \to 0 < \varepsilon \to K \in \mathbb{N} \land 1 < \underline{K} \cdot \varepsilon \to \underline{K}^{-1} \in \mathbb{R}$
\item $\varphi_{40} :\equiv (\underline{K} \cdot \varepsilon = \varepsilon \cdot \underline{K} \to 1 < \underline{K} \cdot \varepsilon \to 1 < \varepsilon \cdot \underline{K}) \to \varepsilon \in \mathbb{R} \to 0 < \varepsilon \to K \in \mathbb{N} \land 1 < \underline{K} \cdot \varepsilon \to 1 < \varepsilon \cdot \underline{K}$
\item $\varphi_{41} :\equiv (1 \cdot \underline{K}^{-1} = \underline{K}^{-1} \to 1 \cdot \underline{K}^{-1} < \varepsilon \to \underline{K}^{-1} < \varepsilon) \to \varepsilon \in \mathbb{R} \to 0 < \varepsilon \to K \in \mathbb{N} \land 1 < \underline{K} \cdot \varepsilon \to \underline{K}^{-1} < \varepsilon$
\item $\varphi_{42} :\equiv (\varepsilon \in \mathbb{R} \to 0 < \varepsilon \to K \in \mathbb{N} \land 1 < \underline{K} \cdot \varepsilon \to \varepsilon \in \mathbb{R}) \to \varepsilon \in \mathbb{R} \to 0 < \varepsilon \to K \in \mathbb{N} \land 1 < \underline{K} \cdot \varepsilon \to \underline{K} \cdot \varepsilon = \varepsilon \cdot \underline{K}$
\item $\varphi_{43} :\equiv (\varepsilon \in \mathbb{R} \to 0 < \varepsilon \to K \in \mathbb{N} \land 1 < \underline{K} \cdot \varepsilon \to 1 < \varepsilon \cdot \underline{K}) \to \varepsilon \in \mathbb{R} \to 0 < \varepsilon \to K \in \mathbb{N} \land 1 < \underline{K} \cdot \varepsilon \to 1 \cdot \underline{K}^{-1} < \varepsilon$
\item $\varphi_{44} :\equiv (\varepsilon \in \mathbb{R} \to 0 < \varepsilon \to K \in \mathbb{N} \land 1 < \underline{K} \cdot \varepsilon \to \underline{K}^{-1} < \varepsilon) \to \varphi_{18}$
\item $\varphi_{45} :\equiv (\varepsilon \in \mathbb{R} \to 0 < \varepsilon \to K \in \mathbb{N} \land 1 < \underline{K} \cdot \varepsilon \to \underline{K}^{-1} \in \mathbb{R}) \to \varphi_{14}$
\item $\varphi_{46} :\equiv (\varepsilon \in \mathbb{R} \to 0 < \varepsilon \to K \in \mathbb{N} \land 1 < \underline{K} \cdot \varepsilon \to \underline{K}^{-1} \in \mathbb{R}) \to \varphi_{23}$
\item $\varphi_{47} :\equiv (\varepsilon \in \mathbb{R} \to 0 < \varepsilon \to \varepsilon \in \mathbb{R}) \to 1 \in \mathbb{R} \to (\varepsilon \in \mathbb{R} \to 0 < \varepsilon \to 0 < \varepsilon) \to \varphi_{24}$
\item $\varphi_{48} :\equiv (\varepsilon \in \mathbb{R} \to 0 < \varepsilon \to K \in \mathbb{N} \land 1 < \underline{K} \cdot \varepsilon \to 1 < \underline{K} \cdot \varepsilon) \to \varphi_{40}$
\item $\varphi_{49} :\equiv (\varepsilon \in \mathbb{R} \to 0 < \varepsilon \to K \in \mathbb{N} \land 1 < \underline{K} \cdot \varepsilon \to 1 \cdot \underline{K}^{-1} < \varepsilon) \to \varphi_{41}$
\item $\varphi_{50} :\equiv (1 \in \mathbb{R} \to \underline{K}^{-1} \in \mathbb{R} \to 1 \cdot \underline{K}^{-1} = \underline{K}^{-1} \cdot 1) \to 1 \in \mathbb{R} \to \varphi_{46}$
\item $\varphi_{51} :\equiv (\underline{K} \in \mathbb{R} \to \underline{K} \neq 0 \to \underline{K}^{-1} \in \mathbb{R}) \to (\varepsilon \in \mathbb{R} \to 0 < \varepsilon \to K \in \mathbb{N} \land 1 < \underline{K} \cdot \varepsilon \to \underline{K} \in \mathbb{R}) \to \varphi_{39}$
\item $\varphi_{52} :\equiv (\forall n\,(\varphi_{1}) \to \varphi_{4}) \to \varphi_{3} \to \varepsilon \in \mathbb{R} \to 0 < \varepsilon \to \exists K \in \mathbb{N}\,(\underline{K}^{-1} < \varepsilon)$
\item $\varphi_{53} :\equiv (\underline{K} \in \mathbb{R} \to \varepsilon \in \mathbb{R} \to \underline{K} \cdot \varepsilon = \varepsilon \cdot \underline{K}) \to (\varepsilon \in \mathbb{R} \to 0 < \varepsilon \to K \in \mathbb{N} \land 1 < \underline{K} \cdot \varepsilon \to \underline{K} \in \mathbb{R}) \to \varphi_{42}$
\item $\varphi_{54} :\equiv (\varepsilon \in \mathbb{R} \to 0 < \varepsilon \to K \in \mathbb{N} \land 1 < \underline{K} \cdot \varepsilon \to \varepsilon \in \mathbb{R}) \to (\varepsilon \in \mathbb{R} \to 0 < \varepsilon \to K \in \mathbb{N} \land 1 < \underline{K} \cdot \varepsilon \to 0 < \underline{K}) \to \varphi_{43}$
\item $\varphi_{55} :\equiv (\varepsilon \in \mathbb{R} \to 0 < \varepsilon \to K \in \mathbb{N} \land 1 < \underline{K} \cdot \varepsilon \to \underline{K} \cdot \varepsilon = \varepsilon \cdot \underline{K}) \to \varphi_{48}$
\item $\varphi_{56} :\equiv 1 \in \mathbb{R} \to (\varepsilon \in \mathbb{R} \to 0 < \varepsilon \to K \in \mathbb{N} \land 1 < \underline{K} \cdot \varepsilon \to \underline{K} \in \mathbb{R}) \to \varphi_{54}$
\item $\varphi_{57} :\equiv (\varepsilon \in \mathbb{R} \to 0 < \varepsilon \to \exists n \in \mathbb{N}\,(1 < \underline{n} \cdot \varepsilon)) \to \forall K\,(\varphi_{26}) \to (\forall K\,(\varphi_{26}) \to \varphi_{28}) \to \varphi_{34} \to \forall n\,(\varphi_{1}) \to \varphi_{52}$
\end{abbrevlist}
\end{fullproof}

% 自動生成 (prover/build.js): IotaInvMono  ノード数 38
\begin{fullproof}{[C25]}{$\forall K\,\forall n\,(K \in \mathbb{N} \to n \in \mathbb{N} \to K \leq n \to \underline{n}^{-1} \leq \underline{K}^{-1})$}{38}\phantomsection\label{tree:IotaInvMono}
\noindent$\Pi_{1}$:
\begin{ptree}
\AXC{$K \in \mathbb{N} \to K \in \mathbb{N}$\lbl{A1}}
\AXC{$\forall n \in \mathbb{N}\,(0 < \underline{n})$\lbl{[N15]}}
\AXC{$\forall n \in \mathbb{N}\,(0 < \underline{n}) \to K \in \mathbb{N} \to 0 < \underline{K}$\lbl{A2}}
\BIC{$K \in \mathbb{N} \to 0 < \underline{K}$}
\AXC{$\varphi_{5}$\lbl{A1}}
\BIC{$(K \in \mathbb{N} \to K \in \mathbb{N}) \to K \in \mathbb{N} \to n \in \mathbb{N} \to K \leq n \to 0 < \underline{K}$}
\BIC{$K \in \mathbb{N} \to n \in \mathbb{N} \to K \leq n \to 0 < \underline{K}$}
\end{ptree}
\noindent$\Pi_{2}$:
\begin{ptree}
\AXC{$\forall n \in \mathbb{N}\,(\underline{n} \in \mathbb{R})$\lbl{[N4]}}
\AXC{$\forall n \in \mathbb{N}\,(\underline{n} \in \mathbb{R}) \to n \in \mathbb{N} \to \underline{n} \in \mathbb{R}$\lbl{A2}}
\BIC{$n \in \mathbb{N} \to \underline{n} \in \mathbb{R}$}
\AXC{$(n \in \mathbb{N} \to \underline{n} \in \mathbb{R}) \to \varphi_{4}$\lbl{A1}}
\BIC{$\varphi_{4}$}
\end{ptree}
\noindent$\Pi_{3}$:
\begin{ptree}
\AXC{$\forall n \in \mathbb{N}\,(\underline{n} \in \mathbb{R})$\lbl{[N4]}}
\AXC{$\forall n \in \mathbb{N}\,(\underline{n} \in \mathbb{R}) \to K \in \mathbb{N} \to \underline{K} \in \mathbb{R}$\lbl{A2}}
\BIC{$K \in \mathbb{N} \to \underline{K} \in \mathbb{R}$}
\AXC{$(K \in \mathbb{N} \to \underline{K} \in \mathbb{R}) \to \varphi_{2}$\lbl{A1}}
\BIC{$\varphi_{2}$}
\end{ptree}
\noindent$\Pi_{4}$:
\begin{ptree}
\AXC{$\varphi_{1}$\lbl{[E20]}}
\AXC{$\varphi_{1} \to \varphi_{3}$\lbl{A2}}
\BIC{$\varphi_{3}$}
\AXC{$\varphi_{3} \to \varphi_{7}$\lbl{A2}}
\BIC{$\varphi_{7}$}
\AXC{$\varphi_{7} \to (K \in \mathbb{N} \to n \in \mathbb{N} \to K \leq n \to \underline{K} \in \mathbb{R}) \to \varphi_{10}$\lbl{A1}}
\BIC{$(K \in \mathbb{N} \to n \in \mathbb{N} \to K \leq n \to \underline{K} \in \mathbb{R}) \to \varphi_{10}$}
\end{ptree}
\noindent$\Pi_{5}$:
\begin{ptree}
\AXC{$K \in \mathbb{N} \to K \in \mathbb{N}$\lbl{A1}}
\AXC{$\stackrel{\Pi_{3}}{\vdots}$}\noLine
\UIC{$\varphi_{2}$}
\BIC{$K \in \mathbb{N} \to n \in \mathbb{N} \to K \leq n \to \underline{K} \in \mathbb{R}$}
\AXC{$\stackrel{\Pi_{4}}{\vdots}$}\noLine
\UIC{$(K \in \mathbb{N} \to n \in \mathbb{N} \to K \leq n \to \underline{K} \in \mathbb{R}) \to \varphi_{10}$}
\BIC{$\varphi_{10}$}
\end{ptree}
\noindent$\Pi_{6}$:
\begin{ptree}
\AXC{$\stackrel{\Pi_{1}}{\vdots}$}\noLine
\UIC{$K \in \mathbb{N} \to n \in \mathbb{N} \to K \leq n \to 0 < \underline{K}$}
\AXC{$K \in \mathbb{N} \to n \in \mathbb{N} \to n \in \mathbb{N}$\lbl{A1}}
\AXC{$\stackrel{\Pi_{2}}{\vdots}$}\noLine
\UIC{$\varphi_{4}$}
\BIC{$K \in \mathbb{N} \to n \in \mathbb{N} \to K \leq n \to \underline{n} \in \mathbb{R}$}
\AXC{$\stackrel{\Pi_{5}}{\vdots}$}\noLine
\UIC{$\varphi_{10}$}
\BIC{$(K \in \mathbb{N} \to n \in \mathbb{N} \to K \leq n \to 0 < \underline{K}) \to \varphi_{9}$}
\BIC{$\varphi_{9}$}
\end{ptree}
\noindent 主木:
\begin{ptree}
\AXC{$\forall m\,\forall n\,(m \in \mathbb{N} \to n \in \mathbb{N} \to m \leq n \to \underline{m} \leq \underline{n})$\lbl{[N16]}}
\AXC{$\varphi_{8}$\lbl{A2}}
\BIC{$\forall n\,(K \in \mathbb{N} \to n \in \mathbb{N} \to K \leq n \to \underline{K} \leq \underline{n})$}
\AXC{$\varphi_{6}$\lbl{A2}}
\BIC{$K \in \mathbb{N} \to n \in \mathbb{N} \to K \leq n \to \underline{K} \leq \underline{n}$}
\AXC{$\stackrel{\Pi_{6}}{\vdots}$}\noLine
\UIC{$\varphi_{9}$}
\BIC{$K \in \mathbb{N} \to n \in \mathbb{N} \to K \leq n \to \underline{n}^{-1} \leq \underline{K}^{-1}$}
\UIC{$\forall n\,(K \in \mathbb{N} \to n \in \mathbb{N} \to K \leq n \to \underline{n}^{-1} \leq \underline{K}^{-1})$}
\UIC{$\forall K\,\forall n\,(K \in \mathbb{N} \to n \in \mathbb{N} \to K \leq n \to \underline{n}^{-1} \leq \underline{K}^{-1})$}
\end{ptree}
\begin{abbrevlist}
\item $\varphi_{1} :\equiv \forall x\,\forall y\,(x \in \mathbb{R} \to y \in \mathbb{R} \to 0 < x \to x \leq y \to y^{-1} \leq x^{-1})$
\item $\varphi_{2} :\equiv (K \in \mathbb{N} \to K \in \mathbb{N}) \to K \in \mathbb{N} \to n \in \mathbb{N} \to K \leq n \to \underline{K} \in \mathbb{R}$
\item $\varphi_{3} :\equiv \forall y\,(\underline{K} \in \mathbb{R} \to y \in \mathbb{R} \to 0 < \underline{K} \to \underline{K} \leq y \to y^{-1} \leq \underline{K}^{-1})$
\item $\varphi_{4} :\equiv (K \in \mathbb{N} \to n \in \mathbb{N} \to n \in \mathbb{N}) \to K \in \mathbb{N} \to n \in \mathbb{N} \to K \leq n \to \underline{n} \in \mathbb{R}$
\item $\varphi_{5} :\equiv (K \in \mathbb{N} \to 0 < \underline{K}) \to (K \in \mathbb{N} \to K \in \mathbb{N}) \to K \in \mathbb{N} \to n \in \mathbb{N} \to K \leq n \to 0 < \underline{K}$
\item $\varphi_{6} :\equiv \forall n\,(K \in \mathbb{N} \to n \in \mathbb{N} \to K \leq n \to \underline{K} \leq \underline{n}) \to K \in \mathbb{N} \to n \in \mathbb{N} \to K \leq n \to \underline{K} \leq \underline{n}$
\item $\varphi_{7} :\equiv \underline{K} \in \mathbb{R} \to \underline{n} \in \mathbb{R} \to 0 < \underline{K} \to \underline{K} \leq \underline{n} \to \underline{n}^{-1} \leq \underline{K}^{-1}$
\item $\varphi_{8} :\equiv \forall m\,\forall n\,(m \in \mathbb{N} \to n \in \mathbb{N} \to m \leq n \to \underline{m} \leq \underline{n}) \to \forall n\,(K \in \mathbb{N} \to n \in \mathbb{N} \to K \leq n \to \underline{K} \leq \underline{n})$
\item $\varphi_{9} :\equiv (K \in \mathbb{N} \to n \in \mathbb{N} \to K \leq n \to \underline{K} \leq \underline{n}) \to K \in \mathbb{N} \to n \in \mathbb{N} \to K \leq n \to \underline{n}^{-1} \leq \underline{K}^{-1}$
\item $\varphi_{10} :\equiv (K \in \mathbb{N} \to n \in \mathbb{N} \to K \leq n \to \underline{n} \in \mathbb{R}) \to (K \in \mathbb{N} \to n \in \mathbb{N} \to K \leq n \to 0 < \underline{K}) \to \varphi_{9}$
\end{abbrevlist}
\end{fullproof}

% 自動生成 (prover/build.js): BW  ノード数 941
\begin{fullproof}{定理6.8}{$\forall a\,(\mathrm{SeqBdd}(a) \to \exists b\,(\mathrm{SubSeq}(b,a) \land \mathrm{Cgt}(b)))$}{941}\phantomsection\label{tree:BW}
\noindent$\Pi_{1}$:
\begin{ptree}
\AXC{$\mathrm{SeqBdd}(a) \to \varphi_{344}$\lbl{A1}}
\AXC{$(\mathrm{SeqBdd}(a) \to \varphi_{344}) \to \mathrm{SeqBdd}(a) \to \varphi_{350}$\lbl{A1}}
\BIC{$\mathrm{SeqBdd}(a) \to \varphi_{350}$}
\end{ptree}
\noindent$\Pi_{2}$:
\begin{ptree}
\AXC{$\forall K\,\forall n\,(K \in \mathbb{N} \to n \in \mathbb{N} \to K \leq n \to \underline{n}^{-1} \leq \underline{K}^{-1})$\lbl{[C25]}}
\AXC{$\varphi_{165}$\lbl{A2}}
\BIC{$\forall n\,(K \in \mathbb{N} \to n \in \mathbb{N} \to K \leq n \to \underline{n}^{-1} \leq \underline{K}^{-1})$}
\AXC{$\varphi_{159}$\lbl{A2}}
\BIC{$K \in \mathbb{N} \to n \in \mathbb{N} \to K \leq n \to \underline{n}^{-1} \leq \underline{K}^{-1}$}
\AXC{$\varphi_{482}$\lbl{A1}}
\BIC{$(\mathrm{SeqBdd}(a) \to \varphi_{334}) \to \varphi_{470}$}
\end{ptree}
\noindent$\Pi_{3}$:
\begin{ptree}
\AXC{$\mathrm{SeqBdd}(a) \to \varphi_{344}$\lbl{A1}}
\AXC{$(\mathrm{SeqBdd}(a) \to \varphi_{344}) \to \mathrm{SeqBdd}(a) \to \varphi_{334}$\lbl{A1}}
\BIC{$\mathrm{SeqBdd}(a) \to \varphi_{334}$}
\AXC{$\stackrel{\Pi_{2}}{\vdots}$}\noLine
\UIC{$(\mathrm{SeqBdd}(a) \to \varphi_{334}) \to \varphi_{470}$}
\BIC{$\varphi_{470}$}
\end{ptree}
\noindent$\Pi_{4}$:
\begin{ptree}
\AXC{$\mathrm{SeqBdd}(a) \to \varphi_{332}$\lbl{A1}}
\AXC{$(\mathrm{SeqBdd}(a) \to \varphi_{332}) \to \mathrm{SeqBdd}(a) \to \varphi_{323}$\lbl{A1}}
\BIC{$\mathrm{SeqBdd}(a) \to \varphi_{323}$}
\end{ptree}
\noindent$\Pi_{5}$:
\begin{ptree}
\AXC{$\varphi_{154}$\lbl{A2}}
\AXC{$\stackrel{\Pi_{4}}{\vdots}$}\noLine
\UIC{$\mathrm{SeqBdd}(a) \to \varphi_{323}$}
\AXC{$(\mathrm{SeqBdd}(a) \to \varphi_{323}) \to \varphi_{154} \to \mathrm{SeqBdd}(a) \to \varphi_{321}$\lbl{A1}}
\BIC{$\varphi_{154} \to \mathrm{SeqBdd}(a) \to \varphi_{321}$}
\BIC{$\mathrm{SeqBdd}(a) \to \varphi_{321}$}
\AXC{$\varphi_{468}$\lbl{A1}}
\BIC{$(\mathrm{SeqBdd}(a) \to \varphi_{341}) \to \mathrm{SeqBdd}(a) \to \varphi_{370}$}
\end{ptree}
\noindent$\Pi_{6}$:
\begin{ptree}
\AXC{$\forall x\,\forall y\,\forall z\,(x < y \to y \leq z \to x < z)$\lbl{[O3]}}
\AXC{$\forall x\,\forall y\,\forall z\,(x < y \to y \leq z \to x < z) \to \varphi_{56}$\lbl{A2}}
\BIC{$\varphi_{56}$}
\AXC{$\varphi_{56} \to \varphi_{121}$\lbl{A2}}
\BIC{$\varphi_{121}$}
\AXC{$\varphi_{121} \to \varphi_{157}$\lbl{A2}}
\BIC{$\varphi_{157}$}
\end{ptree}
\noindent$\Pi_{7}$:
\begin{ptree}
\AXC{$\mathrm{SeqBdd}(a) \to \varphi_{341}$\lbl{A1}}
\AXC{$\stackrel{\Pi_{5}}{\vdots}$}\noLine
\UIC{$(\mathrm{SeqBdd}(a) \to \varphi_{341}) \to \mathrm{SeqBdd}(a) \to \varphi_{370}$}
\BIC{$\mathrm{SeqBdd}(a) \to \varphi_{370}$}
\AXC{$\stackrel{\Pi_{6}}{\vdots}$}\noLine
\UIC{$\varphi_{157}$}
\AXC{$\varphi_{157} \to \varphi_{474}$\lbl{A1}}
\BIC{$\varphi_{474}$}
\BIC{$(\mathrm{SeqBdd}(a) \to \varphi_{355}) \to \mathrm{SeqBdd}(a) \to \varphi_{371}$}
\end{ptree}
\noindent$\Pi_{8}$:
\begin{ptree}
\AXC{$\mathrm{SeqBdd}(a) \to \varphi_{340}$\lbl{A1}}
\AXC{$\mathrm{SeqBdd}(a) \to \varphi_{341}$\lbl{A1}}
\AXC{$\stackrel{\Pi_{3}}{\vdots}$}\noLine
\UIC{$\varphi_{470}$}
\BIC{$(\mathrm{SeqBdd}(a) \to \varphi_{340}) \to \mathrm{SeqBdd}(a) \to \varphi_{355}$}
\BIC{$\mathrm{SeqBdd}(a) \to \varphi_{355}$}
\AXC{$\stackrel{\Pi_{7}}{\vdots}$}\noLine
\UIC{$(\mathrm{SeqBdd}(a) \to \varphi_{355}) \to \mathrm{SeqBdd}(a) \to \varphi_{371}$}
\BIC{$\mathrm{SeqBdd}(a) \to \varphi_{371}$}
\end{ptree}
\noindent$\Pi_{9}$:
\begin{ptree}
\AXC{$\forall x\,\forall y\,\forall z\,(x < y \to y < z \to x < z)$\lbl{[O2]}}
\AXC{$\forall x\,\forall y\,\forall z\,(x < y \to y < z \to x < z) \to \varphi_{65}$\lbl{A2}}
\BIC{$\varphi_{65}$}
\AXC{$\varphi_{65} \to \varphi_{131}$\lbl{A2}}
\BIC{$\varphi_{131}$}
\AXC{$\varphi_{131} \to \varphi_{133}$\lbl{A2}}
\BIC{$\varphi_{133}$}
\end{ptree}
\noindent$\Pi_{10}$:
\begin{ptree}
\AXC{$\stackrel{\Pi_{1}}{\vdots}$}\noLine
\UIC{$\mathrm{SeqBdd}(a) \to \varphi_{350}$}
\AXC{$\stackrel{\Pi_{8}}{\vdots}$}\noLine
\UIC{$\mathrm{SeqBdd}(a) \to \varphi_{371}$}
\AXC{$\stackrel{\Pi_{9}}{\vdots}$}\noLine
\UIC{$\varphi_{133}$}
\AXC{$\varphi_{133} \to \varphi_{473}$\lbl{A1}}
\BIC{$\varphi_{473}$}
\BIC{$(\mathrm{SeqBdd}(a) \to \varphi_{350}) \to \mathrm{SeqBdd}(a) \to \varphi_{360}$}
\BIC{$\mathrm{SeqBdd}(a) \to \varphi_{360}$}
\end{ptree}
\noindent$\Pi_{11}$:
\begin{ptree}
\AXC{$x = x$\lbl{A5}}
\AXC{$x = y \to x = x \to y = x$\lbl{A6}}
\AXC{$(x = y \to x = x \to y = x) \to x = x \to x = y \to y = x$\lbl{A1}}
\BIC{$x = x \to x = y \to y = x$}
\BIC{$x = y \to y = x$}
\UIC{$\forall y\,(x = y \to y = x)$}
\UIC{$\forall x\,\forall y\,(x = y \to y = x)$}
\end{ptree}
\noindent$\Pi_{12}$:
\begin{ptree}
\AXC{$\mathrm{SeqBdd}(a) \to \varphi_{332}$\lbl{A1}}
\AXC{$(\mathrm{SeqBdd}(a) \to \varphi_{332}) \to \mathrm{SeqBdd}(a) \to \varphi_{297}$\lbl{A1}}
\BIC{$\mathrm{SeqBdd}(a) \to \varphi_{297}$}
\end{ptree}
\noindent$\Pi_{13}$:
\begin{ptree}
\AXC{$\forall s\,(\varphi_{20})$\lbl{DefSubIdx}}
\AXC{$\forall s\,(\varphi_{20}) \to \varphi_{20}$\lbl{A2}}
\BIC{$\varphi_{20}$}
\AXC{$\stackrel{\Pi_{12}}{\vdots}$}\noLine
\UIC{$\mathrm{SeqBdd}(a) \to \varphi_{297}$}
\AXC{$(\mathrm{SeqBdd}(a) \to \varphi_{297}) \to \varphi_{20} \to \mathrm{SeqBdd}(a) \to \varphi_{294}$\lbl{A1}}
\BIC{$\varphi_{20} \to \mathrm{SeqBdd}(a) \to \varphi_{294}$}
\BIC{$\mathrm{SeqBdd}(a) \to \varphi_{294}$}
\end{ptree}
\noindent$\Pi_{14}$:
\begin{ptree}
\AXC{$\mathrm{SeqBdd}(a) \to \mathrm{SeqBdd}(a)$\lbl{A1}}
\AXC{$(\mathrm{SeqBdd}(a) \to \mathrm{SeqBdd}(a)) \to \varphi_{35} \to \varphi_{36}$\lbl{A1}}
\BIC{$\varphi_{35} \to \varphi_{36}$}
\end{ptree}
\noindent$\Pi_{15}$:
\begin{ptree}
\AXC{$\forall a\,(\varphi_{35})$\lbl{DefSeqBdd}}
\AXC{$\forall a\,(\varphi_{35}) \to \varphi_{35}$\lbl{A2}}
\BIC{$\varphi_{35}$}
\AXC{$\stackrel{\Pi_{14}}{\vdots}$}\noLine
\UIC{$\varphi_{35} \to \varphi_{36}$}
\BIC{$\varphi_{36}$}
\end{ptree}
\noindent$\Pi_{16}$:
\begin{ptree}
\AXC{$\stackrel{\Pi_{15}}{\vdots}$}\noLine
\UIC{$\varphi_{36}$}
\AXC{$\varphi_{36} \to \mathrm{SeqBdd}(a) \to \mathrm{Seq}(a)$\lbl{A1}}
\BIC{$\mathrm{SeqBdd}(a) \to \mathrm{Seq}(a)$}
\end{ptree}
\noindent$\Pi_{17}$:
\begin{ptree}
\AXC{$\forall a\,\forall s\,(\mathrm{Seq}(a) \to \varphi_{55})$\lbl{[S12]}}
\AXC{$\forall a\,\forall s\,(\mathrm{Seq}(a) \to \varphi_{55}) \to \forall s\,(\mathrm{Seq}(a) \to \varphi_{55})$\lbl{A2}}
\BIC{$\forall s\,(\mathrm{Seq}(a) \to \varphi_{55})$}
\end{ptree}
\noindent$\Pi_{18}$:
\begin{ptree}
\AXC{$\stackrel{\Pi_{17}}{\vdots}$}\noLine
\UIC{$\forall s\,(\mathrm{Seq}(a) \to \varphi_{55})$}
\AXC{$\forall s\,(\mathrm{Seq}(a) \to \varphi_{55}) \to \mathrm{Seq}(a) \to \varphi_{55}$\lbl{A2}}
\BIC{$\mathrm{Seq}(a) \to \varphi_{55}$}
\AXC{$(\mathrm{Seq}(a) \to \varphi_{55}) \to \varphi_{426}$\lbl{A1}}
\BIC{$\varphi_{426}$}
\end{ptree}
\noindent$\Pi_{19}$:
\begin{ptree}
\AXC{$\stackrel{\Pi_{13}}{\vdots}$}\noLine
\UIC{$\mathrm{SeqBdd}(a) \to \varphi_{294}$}
\AXC{$\stackrel{\Pi_{16}}{\vdots}$}\noLine
\UIC{$\mathrm{SeqBdd}(a) \to \mathrm{Seq}(a)$}
\AXC{$\stackrel{\Pi_{18}}{\vdots}$}\noLine
\UIC{$\varphi_{426}$}
\BIC{$(\mathrm{SeqBdd}(a) \to \varphi_{294}) \to \mathrm{SeqBdd}(a) \to \varphi_{330}$}
\BIC{$\mathrm{SeqBdd}(a) \to \varphi_{330}$}
\end{ptree}
\noindent$\Pi_{20}$:
\begin{ptree}
\AXC{$\stackrel{\Pi_{19}}{\vdots}$}\noLine
\UIC{$\mathrm{SeqBdd}(a) \to \varphi_{330}$}
\AXC{$(\mathrm{SeqBdd}(a) \to \varphi_{330}) \to \mathrm{SeqBdd}(a) \to \varphi_{317}$\lbl{A1}}
\BIC{$\mathrm{SeqBdd}(a) \to \varphi_{317}$}
\end{ptree}
\noindent$\Pi_{21}$:
\begin{ptree}
\AXC{$\varphi_{122}$\lbl{A2}}
\AXC{$\stackrel{\Pi_{20}}{\vdots}$}\noLine
\UIC{$\mathrm{SeqBdd}(a) \to \varphi_{317}$}
\AXC{$(\mathrm{SeqBdd}(a) \to \varphi_{317}) \to \varphi_{122} \to \mathrm{SeqBdd}(a) \to \varphi_{315}$\lbl{A1}}
\BIC{$\varphi_{122} \to \mathrm{SeqBdd}(a) \to \varphi_{315}$}
\BIC{$\mathrm{SeqBdd}(a) \to \varphi_{315}$}
\AXC{$\varphi_{466}$\lbl{A1}}
\BIC{$(\mathrm{SeqBdd}(a) \to \varphi_{341}) \to \mathrm{SeqBdd}(a) \to \varphi_{361}$}
\end{ptree}
\noindent$\Pi_{22}$:
\begin{ptree}
\AXC{$\mathrm{SeqBdd}(a) \to \varphi_{341}$\lbl{A1}}
\AXC{$\stackrel{\Pi_{21}}{\vdots}$}\noLine
\UIC{$(\mathrm{SeqBdd}(a) \to \varphi_{341}) \to \mathrm{SeqBdd}(a) \to \varphi_{361}$}
\BIC{$\mathrm{SeqBdd}(a) \to \varphi_{361}$}
\AXC{$(\mathrm{SeqBdd}(a) \to \varphi_{361}) \to \varphi_{404}$\lbl{A1}}
\BIC{$\varphi_{404}$}
\end{ptree}
\noindent$\Pi_{23}$:
\begin{ptree}
\AXC{$\stackrel{\Pi_{11}}{\vdots}$}\noLine
\UIC{$\forall x\,\forall y\,(x = y \to y = x)$}
\AXC{$\varphi_{19}$\lbl{A2}}
\BIC{$\forall y\,(a \circ s(n) = y \to y = a \circ s(n))$}
\AXC{$\varphi_{155}$\lbl{A2}}
\BIC{$a \circ s(n) = a_{s_{n}} \to a_{s_{n}} = a \circ s(n)$}
\AXC{$\stackrel{\Pi_{22}}{\vdots}$}\noLine
\UIC{$\varphi_{404}$}
\BIC{$\mathrm{SeqBdd}(a) \to \varphi_{362}$}
\AXC{$\varphi_{478}$\lbl{A1}}
\BIC{$(\mathrm{SeqBdd}(a) \to \varphi_{360}) \to \varphi_{140} \to \mathrm{SeqBdd}(a) \to \varphi_{363}$}
\end{ptree}
\noindent$\Pi_{24}$:
\begin{ptree}
\AXC{$\stackrel{\Pi_{10}}{\vdots}$}\noLine
\UIC{$\mathrm{SeqBdd}(a) \to \varphi_{360}$}
\AXC{$\stackrel{\Pi_{23}}{\vdots}$}\noLine
\UIC{$(\mathrm{SeqBdd}(a) \to \varphi_{360}) \to \varphi_{140} \to \mathrm{SeqBdd}(a) \to \varphi_{363}$}
\BIC{$\varphi_{140} \to \mathrm{SeqBdd}(a) \to \varphi_{363}$}
\end{ptree}
\noindent$\Pi_{25}$:
\begin{ptree}
\AXC{$z = w \to \lvert z - c \rvert < \varepsilon \to \lvert w - c \rvert < \varepsilon$\lbl{A6}}
\UIC{$\forall w\,(z = w \to \lvert z - c \rvert < \varepsilon \to \lvert w - c \rvert < \varepsilon)$}
\UIC{$\forall z\,\forall w\,(z = w \to \lvert z - c \rvert < \varepsilon \to \lvert w - c \rvert < \varepsilon)$}
\AXC{$\varphi_{179}$\lbl{A2}}
\BIC{$\forall w\,(a_{s_{n}} = w \to \lvert a_{s_{n}} - c \rvert < \varepsilon \to \lvert w - c \rvert < \varepsilon)$}
\AXC{$\varphi_{257}$\lbl{A2}}
\BIC{$\varphi_{140}$}
\AXC{$\stackrel{\Pi_{24}}{\vdots}$}\noLine
\UIC{$\varphi_{140} \to \mathrm{SeqBdd}(a) \to \varphi_{363}$}
\BIC{$\mathrm{SeqBdd}(a) \to \varphi_{363}$}
\end{ptree}
\noindent$\Pi_{26}$:
\begin{ptree}
\AXC{$\stackrel{\Pi_{25}}{\vdots}$}\noLine
\UIC{$\mathrm{SeqBdd}(a) \to \varphi_{363}$}
\AXC{$(\mathrm{SeqBdd}(a) \to \varphi_{363}) \to \varphi_{374}$\lbl{A1}}
\BIC{$\varphi_{374}$}
\UIC{$\forall n\,(\varphi_{374})$}
\AXC{$\forall n\,(\varphi_{374}) \to (\forall n\,(\varphi_{374}) \to \varphi_{377}) \to \varphi_{429} \to \mathrm{SeqBdd}(a) \to \varphi_{366}$\lbl{A1}}
\BIC{$(\forall n\,(\varphi_{374}) \to \varphi_{377}) \to \varphi_{429} \to \mathrm{SeqBdd}(a) \to \varphi_{366}$}
\end{ptree}
\noindent$\Pi_{27}$:
\begin{ptree}
\AXC{$\varphi_{429}$\lbl{A4}}
\AXC{$\forall n\,(\varphi_{374}) \to \varphi_{377}$\lbl{A3}}
\AXC{$\stackrel{\Pi_{26}}{\vdots}$}\noLine
\UIC{$(\forall n\,(\varphi_{374}) \to \varphi_{377}) \to \varphi_{429} \to \mathrm{SeqBdd}(a) \to \varphi_{366}$}
\BIC{$\varphi_{429} \to \mathrm{SeqBdd}(a) \to \varphi_{366}$}
\BIC{$\mathrm{SeqBdd}(a) \to \varphi_{366}$}
\end{ptree}
\noindent$\Pi_{28}$:
\begin{ptree}
\AXC{$\mathrm{SeqBdd}(a) \to \varphi_{344}$\lbl{A1}}
\AXC{$(\mathrm{SeqBdd}(a) \to \varphi_{344}) \to \mathrm{SeqBdd}(a) \to \varphi_{334}$\lbl{A1}}
\BIC{$\mathrm{SeqBdd}(a) \to \varphi_{334}$}
\AXC{$\varphi_{471}$\lbl{A1}}
\BIC{$(\mathrm{SeqBdd}(a) \to \varphi_{366}) \to \mathrm{SeqBdd}(a) \to \varphi_{375}$}
\end{ptree}
\noindent$\Pi_{29}$:
\begin{ptree}
\AXC{$\varphi_{46} \to \varphi_{37}$\lbl{A2}}
\AXC{$\stackrel{\Pi_{27}}{\vdots}$}\noLine
\UIC{$\mathrm{SeqBdd}(a) \to \varphi_{366}$}
\AXC{$\stackrel{\Pi_{28}}{\vdots}$}\noLine
\UIC{$(\mathrm{SeqBdd}(a) \to \varphi_{366}) \to \mathrm{SeqBdd}(a) \to \varphi_{375}$}
\BIC{$\mathrm{SeqBdd}(a) \to \varphi_{375}$}
\AXC{$\varphi_{460}$\lbl{A1}}
\BIC{$(\varphi_{46} \to \varphi_{37}) \to \mathrm{SeqBdd}(a) \to \varphi_{378}$}
\BIC{$\mathrm{SeqBdd}(a) \to \varphi_{378}$}
\end{ptree}
\noindent$\Pi_{30}$:
\begin{ptree}
\AXC{$\forall \varepsilon \in \mathbb{R}\,(0 < \varepsilon \to \exists K \in \mathbb{N}\,(\underline{K}^{-1} < \varepsilon))$\lbl{[C24]}}
\AXC{$\varphi_{174}$\lbl{A2}}
\BIC{$\varepsilon \in \mathbb{R} \to 0 < \varepsilon \to \exists K \in \mathbb{N}\,(\underline{K}^{-1} < \varepsilon)$}
\AXC{$(\varepsilon \in \mathbb{R} \to 0 < \varepsilon \to \exists K \in \mathbb{N}\,(\underline{K}^{-1} < \varepsilon)) \to \varphi_{458}$\lbl{A1}}
\BIC{$\varphi_{458}$}
\end{ptree}
\noindent$\Pi_{31}$:
\begin{ptree}
\AXC{$\mathrm{SeqBdd}(a) \to \varphi_{313}$\lbl{A1}}
\AXC{$\mathrm{SeqBdd}(a) \to \varphi_{311}$\lbl{A1}}
\AXC{$\stackrel{\Pi_{30}}{\vdots}$}\noLine
\UIC{$\varphi_{458}$}
\BIC{$(\mathrm{SeqBdd}(a) \to \varphi_{313}) \to \mathrm{SeqBdd}(a) \to \varphi_{326}$}
\BIC{$\mathrm{SeqBdd}(a) \to \varphi_{326}$}
\AXC{$(\mathrm{SeqBdd}(a) \to \varphi_{326}) \to \varphi_{491}$\lbl{A1}}
\BIC{$\varphi_{491}$}
\end{ptree}
\noindent$\Pi_{32}$:
\begin{ptree}
\AXC{$\varphi_{448}$\lbl{A3}}
\AXC{$\stackrel{\Pi_{29}}{\vdots}$}\noLine
\UIC{$\mathrm{SeqBdd}(a) \to \varphi_{378}$}
\AXC{$(\mathrm{SeqBdd}(a) \to \varphi_{378}) \to \varphi_{357} \to \neg (K \in \mathbb{N} \land \underline{K}^{-1} < \varepsilon)$\lbl{A1}}
\BIC{$\varphi_{357} \to \neg (K \in \mathbb{N} \land \underline{K}^{-1} < \varepsilon)$}
\UIC{$\forall K\,(\varphi_{357} \to \neg (K \in \mathbb{N} \land \underline{K}^{-1} < \varepsilon))$}
\AXC{$\stackrel{\Pi_{31}}{\vdots}$}\noLine
\UIC{$\varphi_{491}$}
\BIC{$\varphi_{448} \to (\varphi_{357} \to \forall K\,(\varphi_{357})) \to \mathrm{SeqBdd}(a) \to \varphi_{347}$}
\BIC{$(\varphi_{357} \to \forall K\,(\varphi_{357})) \to \mathrm{SeqBdd}(a) \to \varphi_{347}$}
\end{ptree}
\noindent$\Pi_{33}$:
\begin{ptree}
\AXC{$\varphi_{357} \to \forall K\,(\varphi_{357})$\lbl{A4}}
\AXC{$\stackrel{\Pi_{32}}{\vdots}$}\noLine
\UIC{$(\varphi_{357} \to \forall K\,(\varphi_{357})) \to \mathrm{SeqBdd}(a) \to \varphi_{347}$}
\BIC{$\mathrm{SeqBdd}(a) \to \varphi_{347}$}
\AXC{$(\mathrm{SeqBdd}(a) \to \varphi_{347}) \to \varphi_{296} \to \varepsilon \in \mathbb{R} \to \varphi_{70}$\lbl{A1}}
\BIC{$\varphi_{296} \to \varepsilon \in \mathbb{R} \to \varphi_{70}$}
\UIC{$\forall \varepsilon\,(\varphi_{296} \to \varepsilon \in \mathbb{R} \to \varphi_{70})$}
\end{ptree}
\noindent$\Pi_{34}$:
\begin{ptree}
\AXC{$\forall \varepsilon\,(\varphi_{296} \to \varepsilon \in \mathbb{R} \to \varphi_{70}) \to \forall \varepsilon\,(\varphi_{296}) \to \forall \varepsilon \in \mathbb{R}\,(\varphi_{70})$\lbl{A3}}
\AXC{$\stackrel{\Pi_{33}}{\vdots}$}\noLine
\UIC{$\forall \varepsilon\,(\varphi_{296} \to \varepsilon \in \mathbb{R} \to \varphi_{70})$}
\AXC{$\forall \varepsilon\,(\varphi_{296} \to \varepsilon \in \mathbb{R} \to \varphi_{70}) \to \varphi_{485}$\lbl{A1}}
\BIC{$\varphi_{485}$}
\BIC{$(\varphi_{296} \to \forall \varepsilon\,(\varphi_{296})) \to \mathrm{SeqBdd}(a) \to \varphi_{351}$}
\end{ptree}
\noindent$\Pi_{35}$:
\begin{ptree}
\AXC{$\varphi_{10}$\lbl{DefSup}}
\AXC{$\varphi_{10} \to \varphi_{9}$\lbl{A2}}
\BIC{$\varphi_{9}$}
\AXC{$\varphi_{9} \to \varphi_{8}$\lbl{A2}}
\BIC{$\varphi_{8}$}
\AXC{$\mathrm{SeqBdd}(a) \to \varphi_{212}$\lbl{A1}}
\AXC{$(\mathrm{SeqBdd}(a) \to \varphi_{212}) \to \varphi_{8} \to \varphi_{267}$\lbl{A1}}
\BIC{$\varphi_{8} \to \varphi_{267}$}
\BIC{$\varphi_{267}$}
\end{ptree}
\noindent$\Pi_{36}$:
\begin{ptree}
\AXC{$\forall a\,\forall X\,(\mathrm{UB}(a,X) \leftrightarrow a \in \mathbb{R} \land \forall x \in X\,(x \leq a))$\lbl{DefUB}}
\AXC{$\varphi_{106}$\lbl{A2}}
\BIC{$\forall X\,(\mathrm{UB}(c,X) \leftrightarrow c \in \mathbb{R} \land \forall x \in X\,(x \leq c))$}
\AXC{$\varphi_{84}$\lbl{A2}}
\BIC{$\mathrm{UB}(c,E) \leftrightarrow c \in \mathbb{R} \land \forall x \in E\,(x \leq c)$}
\AXC{$\stackrel{\Pi_{35}}{\vdots}$}\noLine
\UIC{$\varphi_{267}$}
\BIC{$\mathrm{SeqBdd}(a) \to \varphi_{213}$}
\end{ptree}
\noindent$\Pi_{37}$:
\begin{ptree}
\AXC{$\varphi_{197}$\lbl{DefConv}}
\AXC{$\varphi_{197} \to \varphi_{229}$\lbl{A2}}
\BIC{$\varphi_{229}$}
\AXC{$\varphi_{229} \to \varphi_{221}$\lbl{A2}}
\BIC{$\varphi_{221}$}
\AXC{$\varphi_{221} \to (\mathrm{SeqBdd}(a) \to \varphi_{300}) \to \varphi_{449}$\lbl{A1}}
\BIC{$(\mathrm{SeqBdd}(a) \to \varphi_{300}) \to \varphi_{449}$}
\end{ptree}
\noindent$\Pi_{38}$:
\begin{ptree}
\AXC{$\stackrel{\Pi_{19}}{\vdots}$}\noLine
\UIC{$\mathrm{SeqBdd}(a) \to \varphi_{330}$}
\AXC{$(\mathrm{SeqBdd}(a) \to \varphi_{330}) \to \mathrm{SeqBdd}(a) \to \varphi_{300}$\lbl{A1}}
\BIC{$\mathrm{SeqBdd}(a) \to \varphi_{300}$}
\AXC{$\stackrel{\Pi_{37}}{\vdots}$}\noLine
\UIC{$(\mathrm{SeqBdd}(a) \to \varphi_{300}) \to \varphi_{449}$}
\BIC{$\varphi_{449}$}
\end{ptree}
\noindent$\Pi_{39}$:
\begin{ptree}
\AXC{$\varphi_{296} \to \forall \varepsilon\,(\varphi_{296})$\lbl{A4}}
\AXC{$\stackrel{\Pi_{34}}{\vdots}$}\noLine
\UIC{$(\varphi_{296} \to \forall \varepsilon\,(\varphi_{296})) \to \mathrm{SeqBdd}(a) \to \varphi_{351}$}
\BIC{$\mathrm{SeqBdd}(a) \to \varphi_{351}$}
\AXC{$\stackrel{\Pi_{36}}{\vdots}$}\noLine
\UIC{$\mathrm{SeqBdd}(a) \to \varphi_{213}$}
\AXC{$\stackrel{\Pi_{38}}{\vdots}$}\noLine
\UIC{$\varphi_{449}$}
\BIC{$(\mathrm{SeqBdd}(a) \to \varphi_{351}) \to \mathrm{SeqBdd}(a) \to \varphi_{304}$}
\BIC{$\mathrm{SeqBdd}(a) \to \varphi_{304}$}
\end{ptree}
\noindent$\Pi_{40}$:
\begin{ptree}
\AXC{$\forall \alpha\,(\neg \mathrm{Conv}(a \circ s,\alpha)) \to \neg \mathrm{Conv}(a \circ s,c)$\lbl{A2}}
\AXC{$\stackrel{\Pi_{39}}{\vdots}$}\noLine
\UIC{$\mathrm{SeqBdd}(a) \to \varphi_{304}$}
\AXC{$(\mathrm{SeqBdd}(a) \to \varphi_{304}) \to \varphi_{339}$\lbl{A1}}
\BIC{$\varphi_{339}$}
\BIC{$\mathrm{SeqBdd}(a) \to \varphi_{306}$}
\end{ptree}
\noindent$\Pi_{41}$:
\begin{ptree}
\AXC{$\stackrel{\Pi_{40}}{\vdots}$}\noLine
\UIC{$\mathrm{SeqBdd}(a) \to \varphi_{306}$}
\AXC{$\forall a\,(\mathrm{Cgt}(a) \leftrightarrow \exists \alpha\,(\mathrm{Conv}(a,\alpha)))$\lbl{DefCgt}}
\AXC{$\varphi_{38}$\lbl{A2}}
\BIC{$\mathrm{Cgt}(a \circ s) \leftrightarrow \exists \alpha\,(\mathrm{Conv}(a \circ s,\alpha))$}
\AXC{$\varphi_{421}$\lbl{A1}}
\BIC{$(\mathrm{SeqBdd}(a) \to \varphi_{306}) \to \mathrm{SeqBdd}(a) \to \varphi_{301}$}
\BIC{$\mathrm{SeqBdd}(a) \to \varphi_{301}$}
\end{ptree}
\noindent$\Pi_{42}$:
\begin{ptree}
\AXC{$\forall u_{3}\,(\forall t\,(\varphi_{25}) \to \varphi_{32}) \to \forall u_{3}\,\forall t\,(\varphi_{25}) \to \forall u_{3}\,(\varphi_{32})$\lbl{A3}}
\AXC{$\forall t\,(\varphi_{25}) \to \varphi_{32}$\lbl{A2}}
\UIC{$\forall u_{3}\,(\forall t\,(\varphi_{25}) \to \varphi_{32})$}
\AXC{$\forall u_{3}\,(\forall t\,(\varphi_{25}) \to \varphi_{32}) \to \varphi_{430}$\lbl{A1}}
\BIC{$\varphi_{430}$}
\BIC{$(\forall t\,(\varphi_{25}) \to \forall u_{3}\,\forall t\,(\varphi_{25})) \to \forall t\,(\varphi_{25}) \to \forall u_{3}\,(\varphi_{32})$}
\end{ptree}
\noindent$\Pi_{43}$:
\begin{ptree}
\AXC{$\forall t\,(\varphi_{25}) \to \forall u_{3}\,\forall t\,(\varphi_{25})$\lbl{A4}}
\AXC{$\stackrel{\Pi_{42}}{\vdots}$}\noLine
\UIC{$(\forall t\,(\varphi_{25}) \to \forall u_{3}\,\forall t\,(\varphi_{25})) \to \forall t\,(\varphi_{25}) \to \forall u_{3}\,(\varphi_{32})$}
\BIC{$\forall t\,(\varphi_{25}) \to \forall u_{3}\,(\varphi_{32})$}
\end{ptree}
\noindent$\Pi_{44}$:
\begin{ptree}
\AXC{$\forall t\,(\forall u_{3}\,(\varphi_{32}) \to \varphi_{25}) \to \forall t\,\forall u_{3}\,(\varphi_{32}) \to \forall t\,(\varphi_{25})$\lbl{A3}}
\AXC{$\forall u_{3}\,(\varphi_{32}) \to \varphi_{25}$\lbl{A2}}
\UIC{$\forall t\,(\forall u_{3}\,(\varphi_{32}) \to \varphi_{25})$}
\AXC{$\forall t\,(\forall u_{3}\,(\varphi_{32}) \to \varphi_{25}) \to \varphi_{432}$\lbl{A1}}
\BIC{$\varphi_{432}$}
\BIC{$(\forall u_{3}\,(\varphi_{32}) \to \forall t\,\forall u_{3}\,(\varphi_{32})) \to \forall u_{3}\,(\varphi_{32}) \to \forall t\,(\varphi_{25})$}
\end{ptree}
\noindent$\Pi_{45}$:
\begin{ptree}
\AXC{$\stackrel{\Pi_{20}}{\vdots}$}\noLine
\UIC{$\mathrm{SeqBdd}(a) \to \varphi_{317}$}
\AXC{$\stackrel{\Pi_{12}}{\vdots}$}\noLine
\UIC{$\mathrm{SeqBdd}(a) \to \varphi_{297}$}
\AXC{$\varphi_{456}$\lbl{A1}}
\BIC{$(\mathrm{SeqBdd}(a) \to \varphi_{317}) \to \mathrm{SeqBdd}(a) \to \varphi_{324}$}
\BIC{$\mathrm{SeqBdd}(a) \to \varphi_{324}$}
\end{ptree}
\noindent$\Pi_{46}$:
\begin{ptree}
\AXC{$\stackrel{\Pi_{19}}{\vdots}$}\noLine
\UIC{$\mathrm{SeqBdd}(a) \to \varphi_{330}$}
\AXC{$(\mathrm{SeqBdd}(a) \to \varphi_{330}) \to \mathrm{SeqBdd}(a) \to \varphi_{300}$\lbl{A1}}
\BIC{$\mathrm{SeqBdd}(a) \to \varphi_{300}$}
\end{ptree}
\noindent$\Pi_{47}$:
\begin{ptree}
\AXC{$\forall s\,(\forall u_{3}\,(\varphi_{27}) \to \varphi_{22}) \to \forall s\,\forall u_{3}\,(\varphi_{27}) \to \forall s\,(\varphi_{22})$\lbl{A3}}
\AXC{$\forall u_{3}\,(\varphi_{27}) \to \varphi_{22}$\lbl{A2}}
\UIC{$\forall s\,(\forall u_{3}\,(\varphi_{27}) \to \varphi_{22})$}
\AXC{$\forall s\,(\forall u_{3}\,(\varphi_{27}) \to \varphi_{22}) \to \varphi_{428}$\lbl{A1}}
\BIC{$\varphi_{428}$}
\BIC{$(\forall u_{3}\,(\varphi_{27}) \to \forall s\,\forall u_{3}\,(\varphi_{27})) \to \forall u_{3}\,(\varphi_{27}) \to \forall s\,(\varphi_{22})$}
\end{ptree}
\noindent$\Pi_{48}$:
\begin{ptree}
\AXC{$\forall u_{3}\,(\varphi_{27}) \to \forall s\,\forall u_{3}\,(\varphi_{27})$\lbl{A4}}
\AXC{$\stackrel{\Pi_{47}}{\vdots}$}\noLine
\UIC{$(\forall u_{3}\,(\varphi_{27}) \to \forall s\,\forall u_{3}\,(\varphi_{27})) \to \forall u_{3}\,(\varphi_{27}) \to \forall s\,(\varphi_{22})$}
\BIC{$\forall u_{3}\,(\varphi_{27}) \to \forall s\,(\varphi_{22})$}
\end{ptree}
\noindent$\Pi_{49}$:
\begin{ptree}
\AXC{$\forall u_{3}\,(\forall s\,(\varphi_{22}) \to \varphi_{27}) \to \forall u_{3}\,\forall s\,(\varphi_{22}) \to \forall u_{3}\,(\varphi_{27})$\lbl{A3}}
\AXC{$\forall s\,(\varphi_{22}) \to \varphi_{27}$\lbl{A2}}
\UIC{$\forall u_{3}\,(\forall s\,(\varphi_{22}) \to \varphi_{27})$}
\AXC{$\forall u_{3}\,(\forall s\,(\varphi_{22}) \to \varphi_{27}) \to \varphi_{427}$\lbl{A1}}
\BIC{$\varphi_{427}$}
\BIC{$(\forall s\,(\varphi_{22}) \to \forall u_{3}\,\forall s\,(\varphi_{22})) \to \forall s\,(\varphi_{22}) \to \forall u_{3}\,(\varphi_{27})$}
\end{ptree}
\noindent$\Pi_{50}$:
\begin{ptree}
\AXC{$\forall s\,(\varphi_{22}) \to \forall u_{3}\,\forall s\,(\varphi_{22})$\lbl{A4}}
\AXC{$\stackrel{\Pi_{49}}{\vdots}$}\noLine
\UIC{$(\forall s\,(\varphi_{22}) \to \forall u_{3}\,\forall s\,(\varphi_{22})) \to \forall s\,(\varphi_{22}) \to \forall u_{3}\,(\varphi_{27})$}
\BIC{$\forall s\,(\varphi_{22}) \to \forall u_{3}\,(\varphi_{27})$}
\end{ptree}
\noindent$\Pi_{51}$:
\begin{ptree}
\AXC{$\stackrel{\Pi_{50}}{\vdots}$}\noLine
\UIC{$\forall s\,(\varphi_{22}) \to \forall u_{3}\,(\varphi_{27})$}
\AXC{$\forall s\,(\varphi_{22}) \leftrightarrow \forall s\,(\varphi_{22})$\lbl{A1}}
\AXC{$(\forall s\,(\varphi_{22}) \leftrightarrow \forall s\,(\varphi_{22})) \to \varphi_{408}$\lbl{A1}}
\BIC{$\varphi_{408}$}
\BIC{$(\forall u_{3}\,(\varphi_{27}) \to \forall s\,(\varphi_{22})) \to (\forall s\,(\varphi_{22}) \leftrightarrow \forall u_{3}\,(\varphi_{27}))$}
\end{ptree}
\noindent$\Pi_{52}$:
\begin{ptree}
\AXC{$\stackrel{\Pi_{48}}{\vdots}$}\noLine
\UIC{$\forall u_{3}\,(\varphi_{27}) \to \forall s\,(\varphi_{22})$}
\AXC{$\stackrel{\Pi_{51}}{\vdots}$}\noLine
\UIC{$(\forall u_{3}\,(\varphi_{27}) \to \forall s\,(\varphi_{22})) \to (\forall s\,(\varphi_{22}) \leftrightarrow \forall u_{3}\,(\varphi_{27}))$}
\BIC{$\forall s\,(\varphi_{22}) \leftrightarrow \forall u_{3}\,(\varphi_{27})$}
\AXC{$\varphi_{328}$\lbl{A1}}
\BIC{$\neg \forall s\,(\varphi_{22}) \leftrightarrow \neg \forall u_{3}\,(\varphi_{27})$}
\AXC{$\varphi_{331}$\lbl{A1}}
\BIC{$\neg \neg \forall s\,(\varphi_{22}) \leftrightarrow \neg \neg \forall u_{3}\,(\varphi_{27})$}
\end{ptree}
\noindent$\Pi_{53}$:
\begin{ptree}
\AXC{$\stackrel{\Pi_{52}}{\vdots}$}\noLine
\UIC{$\neg \neg \forall s\,(\varphi_{22}) \leftrightarrow \neg \neg \forall u_{3}\,(\varphi_{27})$}
\AXC{$\mathrm{Seq}(v_{733}) \leftrightarrow \mathrm{Seq}(v_{733})$\lbl{A1}}
\AXC{$\varphi_{369}$\lbl{A1}}
\BIC{$(\neg \neg \forall s\,(\varphi_{22}) \leftrightarrow \neg \neg \forall u_{3}\,(\varphi_{27})) \to \varphi_{231}$}
\BIC{$\varphi_{231}$}
\AXC{$\varphi_{231} \to \varphi_{234}$\lbl{A1}}
\BIC{$\varphi_{234}$}
\AXC{$\varphi_{234} \to \varphi_{7} \to \varphi_{269}$\lbl{A1}}
\BIC{$\varphi_{7} \to \varphi_{269}$}
\end{ptree}
\noindent$\Pi_{54}$:
\begin{ptree}
\AXC{$\forall s\,(\forall u_{3}\,(\varphi_{27}) \to \varphi_{22}) \to \forall s\,\forall u_{3}\,(\varphi_{27}) \to \forall s\,(\varphi_{22})$\lbl{A3}}
\AXC{$\forall u_{3}\,(\varphi_{27}) \to \varphi_{22}$\lbl{A2}}
\UIC{$\forall s\,(\forall u_{3}\,(\varphi_{27}) \to \varphi_{22})$}
\AXC{$\forall s\,(\forall u_{3}\,(\varphi_{27}) \to \varphi_{22}) \to \varphi_{428}$\lbl{A1}}
\BIC{$\varphi_{428}$}
\BIC{$(\forall u_{3}\,(\varphi_{27}) \to \forall s\,\forall u_{3}\,(\varphi_{27})) \to \forall u_{3}\,(\varphi_{27}) \to \forall s\,(\varphi_{22})$}
\end{ptree}
\noindent$\Pi_{55}$:
\begin{ptree}
\AXC{$\forall u_{3}\,(\varphi_{27}) \to \forall s\,\forall u_{3}\,(\varphi_{27})$\lbl{A4}}
\AXC{$\stackrel{\Pi_{54}}{\vdots}$}\noLine
\UIC{$(\forall u_{3}\,(\varphi_{27}) \to \forall s\,\forall u_{3}\,(\varphi_{27})) \to \forall u_{3}\,(\varphi_{27}) \to \forall s\,(\varphi_{22})$}
\BIC{$\forall u_{3}\,(\varphi_{27}) \to \forall s\,(\varphi_{22})$}
\end{ptree}
\noindent$\Pi_{56}$:
\begin{ptree}
\AXC{$\forall u_{3}\,(\forall s\,(\varphi_{22}) \to \varphi_{27}) \to \forall u_{3}\,\forall s\,(\varphi_{22}) \to \forall u_{3}\,(\varphi_{27})$\lbl{A3}}
\AXC{$\forall s\,(\varphi_{22}) \to \varphi_{27}$\lbl{A2}}
\UIC{$\forall u_{3}\,(\forall s\,(\varphi_{22}) \to \varphi_{27})$}
\AXC{$\forall u_{3}\,(\forall s\,(\varphi_{22}) \to \varphi_{27}) \to \varphi_{427}$\lbl{A1}}
\BIC{$\varphi_{427}$}
\BIC{$(\forall s\,(\varphi_{22}) \to \forall u_{3}\,\forall s\,(\varphi_{22})) \to \forall s\,(\varphi_{22}) \to \forall u_{3}\,(\varphi_{27})$}
\end{ptree}
\noindent$\Pi_{57}$:
\begin{ptree}
\AXC{$\forall s\,(\varphi_{22}) \to \forall u_{3}\,\forall s\,(\varphi_{22})$\lbl{A4}}
\AXC{$\stackrel{\Pi_{56}}{\vdots}$}\noLine
\UIC{$(\forall s\,(\varphi_{22}) \to \forall u_{3}\,\forall s\,(\varphi_{22})) \to \forall s\,(\varphi_{22}) \to \forall u_{3}\,(\varphi_{27})$}
\BIC{$\forall s\,(\varphi_{22}) \to \forall u_{3}\,(\varphi_{27})$}
\end{ptree}
\noindent$\Pi_{58}$:
\begin{ptree}
\AXC{$\stackrel{\Pi_{57}}{\vdots}$}\noLine
\UIC{$\forall s\,(\varphi_{22}) \to \forall u_{3}\,(\varphi_{27})$}
\AXC{$\forall s\,(\varphi_{22}) \leftrightarrow \forall s\,(\varphi_{22})$\lbl{A1}}
\AXC{$(\forall s\,(\varphi_{22}) \leftrightarrow \forall s\,(\varphi_{22})) \to \varphi_{408}$\lbl{A1}}
\BIC{$\varphi_{408}$}
\BIC{$(\forall u_{3}\,(\varphi_{27}) \to \forall s\,(\varphi_{22})) \to (\forall s\,(\varphi_{22}) \leftrightarrow \forall u_{3}\,(\varphi_{27}))$}
\end{ptree}
\noindent$\Pi_{59}$:
\begin{ptree}
\AXC{$\stackrel{\Pi_{55}}{\vdots}$}\noLine
\UIC{$\forall u_{3}\,(\varphi_{27}) \to \forall s\,(\varphi_{22})$}
\AXC{$\stackrel{\Pi_{58}}{\vdots}$}\noLine
\UIC{$(\forall u_{3}\,(\varphi_{27}) \to \forall s\,(\varphi_{22})) \to (\forall s\,(\varphi_{22}) \leftrightarrow \forall u_{3}\,(\varphi_{27}))$}
\BIC{$\forall s\,(\varphi_{22}) \leftrightarrow \forall u_{3}\,(\varphi_{27})$}
\AXC{$\varphi_{328}$\lbl{A1}}
\BIC{$\neg \forall s\,(\varphi_{22}) \leftrightarrow \neg \forall u_{3}\,(\varphi_{27})$}
\AXC{$\varphi_{331}$\lbl{A1}}
\BIC{$\neg \neg \forall s\,(\varphi_{22}) \leftrightarrow \neg \neg \forall u_{3}\,(\varphi_{27})$}
\end{ptree}
\noindent$\Pi_{60}$:
\begin{ptree}
\AXC{$\stackrel{\Pi_{59}}{\vdots}$}\noLine
\UIC{$\neg \neg \forall s\,(\varphi_{22}) \leftrightarrow \neg \neg \forall u_{3}\,(\varphi_{27})$}
\AXC{$\mathrm{Seq}(v_{733}) \leftrightarrow \mathrm{Seq}(v_{733})$\lbl{A1}}
\AXC{$\varphi_{369}$\lbl{A1}}
\BIC{$(\neg \neg \forall s\,(\varphi_{22}) \leftrightarrow \neg \neg \forall u_{3}\,(\varphi_{27})) \to \varphi_{231}$}
\BIC{$\varphi_{231}$}
\AXC{$\varphi_{231} \to \varphi_{234}$\lbl{A1}}
\BIC{$\varphi_{234}$}
\end{ptree}
\noindent$\Pi_{61}$:
\begin{ptree}
\AXC{$\varphi_{7}$\lbl{A1}}
\AXC{$\stackrel{\Pi_{53}}{\vdots}$}\noLine
\UIC{$\varphi_{7} \to \varphi_{269}$}
\BIC{$\varphi_{269}$}
\AXC{$\varphi_{269} \to \varphi_{275}$\lbl{A1}}
\BIC{$\varphi_{275}$}
\AXC{$\stackrel{\Pi_{60}}{\vdots}$}\noLine
\UIC{$\varphi_{234}$}
\AXC{$\varphi_{7}$\lbl{A1}}
\AXC{$\varphi_{7} \to \varphi_{234} \to \varphi_{270}$\lbl{A1}}
\BIC{$\varphi_{234} \to \varphi_{270}$}
\BIC{$\varphi_{270}$}
\AXC{$\varphi_{270} \to \varphi_{275} \to \varphi_{411}$\lbl{A1}}
\BIC{$\varphi_{275} \to \varphi_{411}$}
\BIC{$\varphi_{411}$}
\end{ptree}
\noindent$\Pi_{62}$:
\begin{ptree}
\AXC{$\stackrel{\Pi_{61}}{\vdots}$}\noLine
\UIC{$\varphi_{411}$}
\AXC{$\varphi_{411} \to (\varphi_{113} \leftrightarrow \varphi_{285})$\lbl{A1}}
\BIC{$\varphi_{113} \leftrightarrow \varphi_{285}$}
\end{ptree}
\noindent$\Pi_{63}$:
\begin{ptree}
\AXC{$\stackrel{\Pi_{62}}{\vdots}$}\noLine
\UIC{$\varphi_{113} \leftrightarrow \varphi_{285}$}
\AXC{$(\varphi_{113} \leftrightarrow \varphi_{285}) \to \varphi_{285} \to \varphi_{113}$\lbl{A1}}
\BIC{$\varphi_{285} \to \varphi_{113}$}
\UIC{$\forall v_{734}\,(\varphi_{285} \to \varphi_{113})$}
\AXC{$\stackrel{\Pi_{62}}{\vdots}$}\noLine
\UIC{$\varphi_{113} \leftrightarrow \varphi_{285}$}
\AXC{$(\varphi_{113} \leftrightarrow \varphi_{285}) \to \varphi_{113} \to \varphi_{285}$\lbl{A1}}
\BIC{$\varphi_{113} \to \varphi_{285}$}
\UIC{$\forall v_{734}\,(\varphi_{113} \to \varphi_{285})$}
\AXC{$\forall v_{734}\,(\varphi_{113} \to \varphi_{285}) \to \varphi_{493}$\lbl{A1}}
\BIC{$\varphi_{493}$}
\BIC{$(\forall v_{734}\,(\varphi_{113} \to \varphi_{285}) \to \forall v_{734}\,(\varphi_{113}) \to \forall v_{734}\,(\varphi_{285})) \to \varphi_{483}$}
\end{ptree}
\noindent$\Pi_{64}$:
\begin{ptree}
\AXC{$\forall v_{734}\,(\varphi_{113} \to \varphi_{285}) \to \forall v_{734}\,(\varphi_{113}) \to \forall v_{734}\,(\varphi_{285})$\lbl{A3}}
\AXC{$\stackrel{\Pi_{63}}{\vdots}$}\noLine
\UIC{$(\forall v_{734}\,(\varphi_{113} \to \varphi_{285}) \to \forall v_{734}\,(\varphi_{113}) \to \forall v_{734}\,(\varphi_{285})) \to \varphi_{483}$}
\BIC{$\varphi_{483}$}
\end{ptree}
\noindent$\Pi_{65}$:
\begin{ptree}
\AXC{$\forall v_{734}\,(\varphi_{285} \to \varphi_{113}) \to \forall v_{734}\,(\varphi_{285}) \to \forall v_{734}\,(\varphi_{113})$\lbl{A3}}
\AXC{$\stackrel{\Pi_{64}}{\vdots}$}\noLine
\UIC{$\varphi_{483}$}
\BIC{$\forall v_{734}\,(\varphi_{113}) \leftrightarrow \forall v_{734}\,(\varphi_{285})$}
\end{ptree}
\noindent$\Pi_{66}$:
\begin{ptree}
\AXC{$\stackrel{\Pi_{65}}{\vdots}$}\noLine
\UIC{$\forall v_{734}\,(\varphi_{113}) \leftrightarrow \forall v_{734}\,(\varphi_{285})$}
\AXC{$\varphi_{462}$\lbl{A1}}
\BIC{$\forall v_{734}\,(\varphi_{113}) \to \forall v_{734}\,(\varphi_{285})$}
\UIC{$\forall v_{733}\,(\forall v_{734}\,(\varphi_{113}) \to \forall v_{734}\,(\varphi_{285}))$}
\AXC{$\forall v_{733}\,(\forall v_{734}\,(\varphi_{113}) \to \forall v_{734}\,(\varphi_{285})) \to \varphi_{494}$\lbl{A1}}
\BIC{$\varphi_{494}$}
\end{ptree}
\noindent$\Pi_{67}$:
\begin{ptree}
\AXC{$\varphi_{83}$\lbl{DefSubSeq}}
\AXC{$\varphi_{83} \to \varphi_{101}$\lbl{A2}}
\BIC{$\varphi_{101}$}
\AXC{$\varphi_{101} \to \varphi_{113}$\lbl{A2}}
\BIC{$\varphi_{113}$}
\UIC{$\forall v_{734}\,(\varphi_{113})$}
\UIC{$\forall v_{733}\,\forall v_{734}\,(\varphi_{113})$}
\AXC{$\forall v_{733}\,\forall v_{734}\,(\varphi_{113}) \to (\forall v_{733}\,\forall v_{734}\,(\varphi_{113}) \leftrightarrow \varphi_{147}) \to \varphi_{147}$\lbl{A1}}
\BIC{$(\forall v_{733}\,\forall v_{734}\,(\varphi_{113}) \leftrightarrow \varphi_{147}) \to \varphi_{147}$}
\end{ptree}
\noindent$\Pi_{68}$:
\begin{ptree}
\AXC{$\varphi_{464}$\lbl{A3}}
\AXC{$\varphi_{465}$\lbl{A3}}
\AXC{$\stackrel{\Pi_{65}}{\vdots}$}\noLine
\UIC{$\forall v_{734}\,(\varphi_{113}) \leftrightarrow \forall v_{734}\,(\varphi_{285})$}
\AXC{$\varphi_{461}$\lbl{A1}}
\BIC{$\forall v_{734}\,(\varphi_{285}) \to \forall v_{734}\,(\varphi_{113})$}
\UIC{$\forall v_{733}\,(\forall v_{734}\,(\varphi_{285}) \to \forall v_{734}\,(\varphi_{113}))$}
\AXC{$\stackrel{\Pi_{66}}{\vdots}$}\noLine
\UIC{$\varphi_{494}$}
\BIC{$\varphi_{465} \to \varphi_{464} \to (\forall v_{733}\,\forall v_{734}\,(\varphi_{113}) \leftrightarrow \varphi_{147})$}
\BIC{$\varphi_{464} \to (\forall v_{733}\,\forall v_{734}\,(\varphi_{113}) \leftrightarrow \varphi_{147})$}
\BIC{$\forall v_{733}\,\forall v_{734}\,(\varphi_{113}) \leftrightarrow \varphi_{147}$}
\AXC{$\stackrel{\Pi_{67}}{\vdots}$}\noLine
\UIC{$(\forall v_{733}\,\forall v_{734}\,(\varphi_{113}) \leftrightarrow \varphi_{147}) \to \varphi_{147}$}
\BIC{$\varphi_{147}$}
\end{ptree}
\noindent$\Pi_{69}$:
\begin{ptree}
\AXC{$\stackrel{\Pi_{46}}{\vdots}$}\noLine
\UIC{$\mathrm{SeqBdd}(a) \to \varphi_{300}$}
\AXC{$\stackrel{\Pi_{68}}{\vdots}$}\noLine
\UIC{$\varphi_{147}$}
\AXC{$\varphi_{147} \to \varphi_{160}$\lbl{A2}}
\BIC{$\varphi_{160}$}
\AXC{$\varphi_{160} \to \varphi_{143}$\lbl{A2}}
\BIC{$\varphi_{143}$}
\AXC{$\varphi_{479}$\lbl{A1}}
\BIC{$(\mathrm{SeqBdd}(a) \to \varphi_{300}) \to (\mathrm{SeqBdd}(a) \to \varphi_{327}) \to \varphi_{420}$}
\BIC{$(\mathrm{SeqBdd}(a) \to \varphi_{327}) \to \varphi_{420}$}
\end{ptree}
\noindent$\Pi_{70}$:
\begin{ptree}
\AXC{$\forall t\,(\varphi_{25}) \to \varphi_{21}$\lbl{A2}}
\AXC{$\stackrel{\Pi_{45}}{\vdots}$}\noLine
\UIC{$\mathrm{SeqBdd}(a) \to \varphi_{324}$}
\AXC{$\varphi_{447}$\lbl{A1}}
\BIC{$(\forall t\,(\varphi_{25}) \to \varphi_{21}) \to \mathrm{SeqBdd}(a) \to \varphi_{327}$}
\BIC{$\mathrm{SeqBdd}(a) \to \varphi_{327}$}
\AXC{$\stackrel{\Pi_{69}}{\vdots}$}\noLine
\UIC{$(\mathrm{SeqBdd}(a) \to \varphi_{327}) \to \varphi_{420}$}
\BIC{$\varphi_{420}$}
\end{ptree}
\noindent$\Pi_{71}$:
\begin{ptree}
\AXC{$\forall u_{3}\,(\varphi_{32}) \to \forall t\,\forall u_{3}\,(\varphi_{32})$\lbl{A4}}
\AXC{$\stackrel{\Pi_{44}}{\vdots}$}\noLine
\UIC{$(\forall u_{3}\,(\varphi_{32}) \to \forall t\,\forall u_{3}\,(\varphi_{32})) \to \forall u_{3}\,(\varphi_{32}) \to \forall t\,(\varphi_{25})$}
\BIC{$\forall u_{3}\,(\varphi_{32}) \to \forall t\,(\varphi_{25})$}
\AXC{$\stackrel{\Pi_{70}}{\vdots}$}\noLine
\UIC{$\varphi_{420}$}
\BIC{$(\forall t\,(\varphi_{25}) \to \forall u_{3}\,(\varphi_{32})) \to \mathrm{SeqBdd}(a) \to \varphi_{305}$}
\end{ptree}
\noindent$\Pi_{72}$:
\begin{ptree}
\AXC{$\stackrel{\Pi_{41}}{\vdots}$}\noLine
\UIC{$\mathrm{SeqBdd}(a) \to \varphi_{301}$}
\AXC{$\stackrel{\Pi_{43}}{\vdots}$}\noLine
\UIC{$\forall t\,(\varphi_{25}) \to \forall u_{3}\,(\varphi_{32})$}
\AXC{$\stackrel{\Pi_{71}}{\vdots}$}\noLine
\UIC{$(\forall t\,(\varphi_{25}) \to \forall u_{3}\,(\varphi_{32})) \to \mathrm{SeqBdd}(a) \to \varphi_{305}$}
\BIC{$\mathrm{SeqBdd}(a) \to \varphi_{305}$}
\AXC{$\varphi_{451}$\lbl{A1}}
\BIC{$(\mathrm{SeqBdd}(a) \to \varphi_{301}) \to \mathrm{SeqBdd}(a) \to \varphi_{314}$}
\BIC{$\mathrm{SeqBdd}(a) \to \varphi_{314}$}
\end{ptree}
\noindent$\Pi_{73}$:
\begin{ptree}
\AXC{$\varphi_{30}$\lbl{A2}}
\AXC{$\stackrel{\Pi_{72}}{\vdots}$}\noLine
\UIC{$\mathrm{SeqBdd}(a) \to \varphi_{314}$}
\AXC{$(\mathrm{SeqBdd}(a) \to \varphi_{314}) \to \varphi_{30} \to \mathrm{SeqBdd}(a) \to \varphi_{307}$\lbl{A1}}
\BIC{$\varphi_{30} \to \mathrm{SeqBdd}(a) \to \varphi_{307}$}
\BIC{$\mathrm{SeqBdd}(a) \to \varphi_{307}$}
\end{ptree}
\noindent$\Pi_{74}$:
\begin{ptree}
\AXC{$\forall m\,\forall n\,(n \leq m \cup n)$\lbl{[C8]}}
\AXC{$\forall m\,\forall n\,(n \leq m \cup n) \to \forall n\,(n \leq N \cup n)$\lbl{A2}}
\BIC{$\forall n\,(n \leq N \cup n)$}
\AXC{$\forall n\,(n \leq N \cup n) \to N_{1} \leq N \cup N_{1}$\lbl{A2}}
\BIC{$N_{1} \leq N \cup N_{1}$}
\end{ptree}
\noindent$\Pi_{75}$:
\begin{ptree}
\AXC{$\forall x\,\forall y\,\forall z\,(x \leq y \to y \leq z \to x \leq z)$\lbl{[O5]}}
\AXC{$\varphi_{13}$\lbl{A2}}
\BIC{$\forall y\,\forall z\,(N_{1} \leq y \to y \leq z \to N_{1} \leq z)$}
\AXC{$\varphi_{24}$\lbl{A2}}
\BIC{$\forall z\,(N_{1} \leq N \cup N_{1} \to N \cup N_{1} \leq z \to N_{1} \leq z)$}
\AXC{$\varphi_{49}$\lbl{A2}}
\BIC{$N_{1} \leq N \cup N_{1} \to N \cup N_{1} \leq n \to N_{1} \leq n$}
\end{ptree}
\noindent$\Pi_{76}$:
\begin{ptree}
\AXC{$\stackrel{\Pi_{74}}{\vdots}$}\noLine
\UIC{$N_{1} \leq N \cup N_{1}$}
\AXC{$\stackrel{\Pi_{75}}{\vdots}$}\noLine
\UIC{$N_{1} \leq N \cup N_{1} \to N \cup N_{1} \leq n \to N_{1} \leq n$}
\AXC{$(N_{1} \leq N \cup N_{1} \to N \cup N_{1} \leq n \to N_{1} \leq n) \to \varphi_{441}$\lbl{A1}}
\BIC{$\varphi_{441}$}
\BIC{$(\mathrm{SeqBdd}(a) \to \varphi_{356}) \to \mathrm{SeqBdd}(a) \to \varphi_{349}$}
\end{ptree}
\noindent$\Pi_{77}$:
\begin{ptree}
\AXC{$\mathrm{SeqBdd}(a) \to \varphi_{333}$\lbl{A1}}
\AXC{$(\mathrm{SeqBdd}(a) \to \varphi_{333}) \to \mathrm{SeqBdd}(a) \to \varphi_{322}$\lbl{A1}}
\BIC{$\mathrm{SeqBdd}(a) \to \varphi_{322}$}
\AXC{$(\mathrm{SeqBdd}(a) \to \varphi_{322}) \to \varphi_{134} \to \varphi_{329}$\lbl{A1}}
\BIC{$\varphi_{134} \to \varphi_{329}$}
\end{ptree}
\noindent$\Pi_{78}$:
\begin{ptree}
\AXC{$\varphi_{387}$\lbl{A1}}
\AXC{$\varphi_{387} \to \mathrm{SeqBdd}(a) \to \varphi_{353}$\lbl{A1}}
\BIC{$\mathrm{SeqBdd}(a) \to \varphi_{353}$}
\AXC{$\varphi_{134}$\lbl{A2}}
\AXC{$\stackrel{\Pi_{77}}{\vdots}$}\noLine
\UIC{$\varphi_{134} \to \varphi_{329}$}
\BIC{$\varphi_{329}$}
\AXC{$\varphi_{481}$\lbl{A1}}
\BIC{$(\mathrm{SeqBdd}(a) \to \varphi_{353}) \to (\mathrm{SeqBdd}(a) \to \varphi_{349}) \to \varphi_{380}$}
\BIC{$(\mathrm{SeqBdd}(a) \to \varphi_{349}) \to \varphi_{380}$}
\end{ptree}
\noindent$\Pi_{79}$:
\begin{ptree}
\AXC{$\varphi_{387}$\lbl{A1}}
\AXC{$\varphi_{387} \to \mathrm{SeqBdd}(a) \to \varphi_{356}$\lbl{A1}}
\BIC{$\mathrm{SeqBdd}(a) \to \varphi_{356}$}
\AXC{$\stackrel{\Pi_{76}}{\vdots}$}\noLine
\UIC{$(\mathrm{SeqBdd}(a) \to \varphi_{356}) \to \mathrm{SeqBdd}(a) \to \varphi_{349}$}
\BIC{$\mathrm{SeqBdd}(a) \to \varphi_{349}$}
\AXC{$\stackrel{\Pi_{78}}{\vdots}$}\noLine
\UIC{$(\mathrm{SeqBdd}(a) \to \varphi_{349}) \to \varphi_{380}$}
\BIC{$\varphi_{380}$}
\end{ptree}
\noindent$\Pi_{80}$:
\begin{ptree}
\AXC{$\stackrel{\Pi_{11}}{\vdots}$}\noLine
\UIC{$\forall x\,\forall y\,(x = y \to y = x)$}
\AXC{$\forall x\,\forall y\,(x = y \to y = x) \to \forall y\,(\underline{k} = y \to y = \underline{k})$\lbl{A2}}
\BIC{$\forall y\,(\underline{k} = y \to y = \underline{k})$}
\end{ptree}
\noindent$\Pi_{81}$:
\begin{ptree}
\AXC{$\mathrm{SeqBdd}(a) \to \varphi_{226}$\lbl{A1}}
\AXC{$\forall n \in \mathbb{N}\,(0 < \underline{n})$\lbl{[N15]}}
\AXC{$\forall n \in \mathbb{N}\,(0 < \underline{n}) \to k \in \mathbb{N} \to 0 < \underline{k}$\lbl{A2}}
\BIC{$k \in \mathbb{N} \to 0 < \underline{k}$}
\AXC{$\varphi_{395}$\lbl{A1}}
\BIC{$(\mathrm{SeqBdd}(a) \to \varphi_{226}) \to \mathrm{SeqBdd}(a) \to \varphi_{230}$}
\BIC{$\mathrm{SeqBdd}(a) \to \varphi_{230}$}
\end{ptree}
\noindent$\Pi_{82}$:
\begin{ptree}
\AXC{$\forall x\,\forall y\,(x < y \to x \neq y)$\lbl{[F26]}}
\AXC{$\forall x\,\forall y\,(x < y \to x \neq y) \to \forall y\,(0 < y \to 0 \neq y)$\lbl{A2}}
\BIC{$\forall y\,(0 < y \to 0 \neq y)$}
\AXC{$\forall y\,(0 < y \to 0 \neq y) \to 0 < \underline{k} \to 0 \neq \underline{k}$\lbl{A2}}
\BIC{$0 < \underline{k} \to 0 \neq \underline{k}$}
\end{ptree}
\noindent$\Pi_{83}$:
\begin{ptree}
\AXC{$\stackrel{\Pi_{81}}{\vdots}$}\noLine
\UIC{$\mathrm{SeqBdd}(a) \to \varphi_{230}$}
\AXC{$\stackrel{\Pi_{82}}{\vdots}$}\noLine
\UIC{$0 < \underline{k} \to 0 \neq \underline{k}$}
\AXC{$\varphi_{396}$\lbl{A1}}
\BIC{$(\mathrm{SeqBdd}(a) \to \varphi_{230}) \to \mathrm{SeqBdd}(a) \to \varphi_{227}$}
\BIC{$\mathrm{SeqBdd}(a) \to \varphi_{227}$}
\AXC{$\varphi_{390}$\lbl{A1}}
\BIC{$(\underline{k} = 0 \to 0 = \underline{k}) \to \mathrm{SeqBdd}(a) \to \varphi_{228}$}
\end{ptree}
\noindent$\Pi_{84}$:
\begin{ptree}
\AXC{$\stackrel{\Pi_{80}}{\vdots}$}\noLine
\UIC{$\forall y\,(\underline{k} = y \to y = \underline{k})$}
\AXC{$\forall y\,(\underline{k} = y \to y = \underline{k}) \to \underline{k} = 0 \to 0 = \underline{k}$\lbl{A2}}
\BIC{$\underline{k} = 0 \to 0 = \underline{k}$}
\AXC{$\stackrel{\Pi_{83}}{\vdots}$}\noLine
\UIC{$(\underline{k} = 0 \to 0 = \underline{k}) \to \mathrm{SeqBdd}(a) \to \varphi_{228}$}
\BIC{$\mathrm{SeqBdd}(a) \to \varphi_{228}$}
\end{ptree}
\noindent$\Pi_{85}$:
\begin{ptree}
\AXC{$\mathrm{SeqBdd}(a) \to \varphi_{226}$\lbl{A1}}
\AXC{$\forall n \in \mathbb{N}\,(\underline{n} \in \mathbb{R})$\lbl{[N4]}}
\AXC{$\forall n \in \mathbb{N}\,(\underline{n} \in \mathbb{R}) \to k \in \mathbb{N} \to \underline{k} \in \mathbb{R}$\lbl{A2}}
\BIC{$k \in \mathbb{N} \to \underline{k} \in \mathbb{R}$}
\AXC{$\varphi_{400}$\lbl{A1}}
\BIC{$(\mathrm{SeqBdd}(a) \to \varphi_{226}) \to \mathrm{SeqBdd}(a) \to \varphi_{236}$}
\BIC{$\mathrm{SeqBdd}(a) \to \varphi_{236}$}
\end{ptree}
\noindent$\Pi_{86}$:
\begin{ptree}
\AXC{$\stackrel{\Pi_{85}}{\vdots}$}\noLine
\UIC{$\mathrm{SeqBdd}(a) \to \varphi_{236}$}
\AXC{$\forall x \in \mathbb{R}\,(x \neq 0 \to x^{-1} \in \mathbb{R})$\lbl{(R1)}}
\AXC{$\varphi_{123}$\lbl{A2}}
\BIC{$\underline{k} \in \mathbb{R} \to \underline{k} \neq 0 \to \underline{k}^{-1} \in \mathbb{R}$}
\AXC{$(\underline{k} \in \mathbb{R} \to \underline{k} \neq 0 \to \underline{k}^{-1} \in \mathbb{R}) \to \varphi_{425}$\lbl{A1}}
\BIC{$\varphi_{425}$}
\BIC{$(\mathrm{SeqBdd}(a) \to \varphi_{228}) \to \mathrm{SeqBdd}(a) \to \varphi_{242}$}
\end{ptree}
\noindent$\Pi_{87}$:
\begin{ptree}
\AXC{$\stackrel{\Pi_{84}}{\vdots}$}\noLine
\UIC{$\mathrm{SeqBdd}(a) \to \varphi_{228}$}
\AXC{$\stackrel{\Pi_{86}}{\vdots}$}\noLine
\UIC{$(\mathrm{SeqBdd}(a) \to \varphi_{228}) \to \mathrm{SeqBdd}(a) \to \varphi_{242}$}
\BIC{$\mathrm{SeqBdd}(a) \to \varphi_{242}$}
\end{ptree}
\noindent$\Pi_{88}$:
\begin{ptree}
\AXC{$\stackrel{\Pi_{36}}{\vdots}$}\noLine
\UIC{$\mathrm{SeqBdd}(a) \to \varphi_{213}$}
\AXC{$\forall x\,\forall y\,(x \in \mathbb{R} \to y \in \mathbb{R} \to x + y \in \mathbb{R})$\lbl{(R1)}}
\AXC{$\varphi_{124}$\lbl{A2}}
\BIC{$\forall y\,(c \in \mathbb{R} \to y \in \mathbb{R} \to c + y \in \mathbb{R})$}
\AXC{$\varphi_{152}$\lbl{A2}}
\BIC{$c \in \mathbb{R} \to \underline{k}^{-1} \in \mathbb{R} \to c + \underline{k}^{-1} \in \mathbb{R}$}
\AXC{$\varphi_{453}$\lbl{A1}}
\BIC{$(\mathrm{SeqBdd}(a) \to \varphi_{213}) \to (\mathrm{SeqBdd}(a) \to \varphi_{242}) \to \varphi_{379}$}
\BIC{$(\mathrm{SeqBdd}(a) \to \varphi_{242}) \to \varphi_{379}$}
\end{ptree}
\noindent$\Pi_{89}$:
\begin{ptree}
\AXC{$\forall a\,\forall n\,(\mathrm{Seq}(a) \to n \in \mathbb{N} \to a_{n} \in \mathbb{R})$\lbl{[S1]}}
\AXC{$\varphi_{96}$\lbl{A2}}
\BIC{$\forall n\,(\mathrm{Seq}(a) \to n \in \mathbb{N} \to a_{n} \in \mathbb{R})$}
\AXC{$\varphi_{75}$\lbl{A2}}
\BIC{$\mathrm{Seq}(a) \to n \in \mathbb{N} \to a_{n} \in \mathbb{R}$}
\AXC{$(\mathrm{Seq}(a) \to n \in \mathbb{N} \to a_{n} \in \mathbb{R}) \to \varphi_{444}$\lbl{A1}}
\BIC{$\varphi_{444}$}
\end{ptree}
\noindent$\Pi_{90}$:
\begin{ptree}
\AXC{$\varphi_{387}$\lbl{A1}}
\AXC{$\varphi_{387} \to \mathrm{SeqBdd}(a) \to \varphi_{353}$\lbl{A1}}
\BIC{$\mathrm{SeqBdd}(a) \to \varphi_{353}$}
\AXC{$\stackrel{\Pi_{16}}{\vdots}$}\noLine
\UIC{$\mathrm{SeqBdd}(a) \to \mathrm{Seq}(a)$}
\AXC{$\stackrel{\Pi_{89}}{\vdots}$}\noLine
\UIC{$\varphi_{444}$}
\BIC{$(\mathrm{SeqBdd}(a) \to \varphi_{353}) \to \mathrm{SeqBdd}(a) \to \varphi_{359}$}
\BIC{$\mathrm{SeqBdd}(a) \to \varphi_{359}$}
\end{ptree}
\noindent$\Pi_{91}$:
\begin{ptree}
\AXC{$\forall x\,\forall y\,(x \in \mathbb{R} \to y \in \mathbb{R} \to \neg y \leq x \to x < y)$\lbl{[A2]}}
\AXC{$\varphi_{153}$\lbl{A2}}
\BIC{$\forall y\,(a_{n} \in \mathbb{R} \to y \in \mathbb{R} \to \neg y \leq a_{n} \to a_{n} < y)$}
\AXC{$\forall y\,(a_{n} \in \mathbb{R} \to y \in \mathbb{R} \to \neg y \leq a_{n} \to a_{n} < y) \to \varphi_{149}$\lbl{A2}}
\BIC{$\varphi_{149}$}
\end{ptree}
\noindent$\Pi_{92}$:
\begin{ptree}
\AXC{$\stackrel{\Pi_{90}}{\vdots}$}\noLine
\UIC{$\mathrm{SeqBdd}(a) \to \varphi_{359}$}
\AXC{$\stackrel{\Pi_{91}}{\vdots}$}\noLine
\UIC{$\varphi_{149}$}
\AXC{$\varphi_{149} \to (\mathrm{SeqBdd}(a) \to \varphi_{359}) \to \varphi_{379} \to \varphi_{380} \to \varphi_{381}$\lbl{A1}}
\BIC{$(\mathrm{SeqBdd}(a) \to \varphi_{359}) \to \varphi_{379} \to \varphi_{380} \to \varphi_{381}$}
\BIC{$\varphi_{379} \to \varphi_{380} \to \varphi_{381}$}
\end{ptree}
\noindent$\Pi_{93}$:
\begin{ptree}
\AXC{$\mathrm{SeqBdd}(a) \to \varphi_{336}$\lbl{A1}}
\AXC{$(\mathrm{SeqBdd}(a) \to \varphi_{336}) \to \mathrm{SeqBdd}(a) \to \varphi_{367}$\lbl{A1}}
\BIC{$\mathrm{SeqBdd}(a) \to \varphi_{367}$}
\end{ptree}
\noindent$\Pi_{94}$:
\begin{ptree}
\AXC{$\forall x\,\forall y\,\forall z\,(x < y \to y \leq z \to x < z)$\lbl{[O3]}}
\AXC{$\varphi_{114}$\lbl{A2}}
\BIC{$\forall y\,\forall z\,(c - \underline{k}^{-1} < y \to y \leq z \to c - \underline{k}^{-1} < z)$}
\AXC{$\varphi_{162}$\lbl{A2}}
\BIC{$\forall z\,(c - \underline{k}^{-1} < x \to x \leq z \to c - \underline{k}^{-1} < z)$}
\AXC{$\varphi_{169}$\lbl{A2}}
\BIC{$c - \underline{k}^{-1} < x \to x \leq a_{n} \to c - \underline{k}^{-1} < a_{n}$}
\AXC{$\varphi_{477}$\lbl{A1}}
\BIC{$(\mathrm{SeqBdd}(a) \to \varphi_{367}) \to (\mathrm{SeqBdd}(a) \to \varphi_{354}) \to \varphi_{382}$}
\end{ptree}
\noindent$\Pi_{95}$:
\begin{ptree}
\AXC{$\varphi_{387}$\lbl{A1}}
\AXC{$\varphi_{387} \to \mathrm{SeqBdd}(a) \to \varphi_{354}$\lbl{A1}}
\BIC{$\mathrm{SeqBdd}(a) \to \varphi_{354}$}
\AXC{$\stackrel{\Pi_{93}}{\vdots}$}\noLine
\UIC{$\mathrm{SeqBdd}(a) \to \varphi_{367}$}
\AXC{$\stackrel{\Pi_{94}}{\vdots}$}\noLine
\UIC{$(\mathrm{SeqBdd}(a) \to \varphi_{367}) \to (\mathrm{SeqBdd}(a) \to \varphi_{354}) \to \varphi_{382}$}
\BIC{$(\mathrm{SeqBdd}(a) \to \varphi_{354}) \to \varphi_{382}$}
\BIC{$\varphi_{382}$}
\end{ptree}
\noindent$\Pi_{96}$:
\begin{ptree}
\AXC{$\varphi_{115}$\lbl{[E33]}}
\AXC{$\varphi_{115} \to \varphi_{102}$\lbl{A2}}
\BIC{$\varphi_{102}$}
\AXC{$\varphi_{102} \to \varphi_{180}$\lbl{A2}}
\BIC{$\varphi_{180}$}
\AXC{$\varphi_{180} \to c \in \mathbb{R} \to \underline{k}^{-1} \in \mathbb{R} \to a_{n} \in \mathbb{R} \to \varphi_{137}$\lbl{A2}}
\BIC{$c \in \mathbb{R} \to \underline{k}^{-1} \in \mathbb{R} \to a_{n} \in \mathbb{R} \to \varphi_{137}$}
\AXC{$\varphi_{488}$\lbl{A1}}
\BIC{$(\mathrm{SeqBdd}(a) \to \varphi_{213}) \to \varphi_{487}$}
\end{ptree}
\noindent$\Pi_{97}$:
\begin{ptree}
\AXC{$\stackrel{\Pi_{95}}{\vdots}$}\noLine
\UIC{$\varphi_{382}$}
\AXC{$\stackrel{\Pi_{90}}{\vdots}$}\noLine
\UIC{$\mathrm{SeqBdd}(a) \to \varphi_{359}$}
\AXC{$\stackrel{\Pi_{87}}{\vdots}$}\noLine
\UIC{$\mathrm{SeqBdd}(a) \to \varphi_{242}$}
\AXC{$\stackrel{\Pi_{36}}{\vdots}$}\noLine
\UIC{$\mathrm{SeqBdd}(a) \to \varphi_{213}$}
\AXC{$\stackrel{\Pi_{96}}{\vdots}$}\noLine
\UIC{$(\mathrm{SeqBdd}(a) \to \varphi_{213}) \to \varphi_{487}$}
\BIC{$\varphi_{487}$}
\BIC{$(\mathrm{SeqBdd}(a) \to \varphi_{359}) \to \varphi_{382} \to \varphi_{381} \to \varphi_{385}$}
\BIC{$\varphi_{382} \to \varphi_{381} \to \varphi_{385}$}
\BIC{$\varphi_{381} \to \varphi_{385}$}
\end{ptree}
\noindent$\Pi_{98}$:
\begin{ptree}
\AXC{$\stackrel{\Pi_{79}}{\vdots}$}\noLine
\UIC{$\varphi_{380}$}
\AXC{$\stackrel{\Pi_{87}}{\vdots}$}\noLine
\UIC{$\mathrm{SeqBdd}(a) \to \varphi_{242}$}
\AXC{$\stackrel{\Pi_{88}}{\vdots}$}\noLine
\UIC{$(\mathrm{SeqBdd}(a) \to \varphi_{242}) \to \varphi_{379}$}
\BIC{$\varphi_{379}$}
\AXC{$\stackrel{\Pi_{92}}{\vdots}$}\noLine
\UIC{$\varphi_{379} \to \varphi_{380} \to \varphi_{381}$}
\BIC{$\varphi_{380} \to \varphi_{381}$}
\BIC{$\varphi_{381}$}
\AXC{$\stackrel{\Pi_{97}}{\vdots}$}\noLine
\UIC{$\varphi_{381} \to \varphi_{385}$}
\BIC{$\varphi_{385}$}
\end{ptree}
\noindent$\Pi_{99}$:
\begin{ptree}
\AXC{$\forall x\,\forall y\,\forall z\,(x \leq y \to y \leq z \to x \leq z)$\lbl{[O5]}}
\AXC{$\varphi_{11}$\lbl{A2}}
\BIC{$\forall y\,\forall z\,(N \leq y \to y \leq z \to N \leq z)$}
\AXC{$\varphi_{17}$\lbl{A2}}
\BIC{$\forall z\,(N \leq N \cup N_{1} \to N \cup N_{1} \leq z \to N \leq z)$}
\AXC{$\varphi_{39}$\lbl{A2}}
\BIC{$N \leq N \cup N_{1} \to N \cup N_{1} \leq n \to N \leq n$}
\AXC{$(N \leq N \cup N_{1} \to N \cup N_{1} \leq n \to N \leq n) \to \varphi_{440}$\lbl{A1}}
\BIC{$\varphi_{440}$}
\end{ptree}
\noindent$\Pi_{100}$:
\begin{ptree}
\AXC{$\forall m\,\forall n\,(m \leq m \cup n)$\lbl{[C7]}}
\AXC{$\forall m\,\forall n\,(m \leq m \cup n) \to \forall n\,(N \leq N \cup n)$\lbl{A2}}
\BIC{$\forall n\,(N \leq N \cup n)$}
\AXC{$\forall n\,(N \leq N \cup n) \to N \leq N \cup N_{1}$\lbl{A2}}
\BIC{$N \leq N \cup N_{1}$}
\AXC{$\stackrel{\Pi_{99}}{\vdots}$}\noLine
\UIC{$\varphi_{440}$}
\BIC{$(\mathrm{SeqBdd}(a) \to \varphi_{356}) \to \mathrm{SeqBdd}(a) \to \varphi_{348}$}
\end{ptree}
\noindent$\Pi_{101}$:
\begin{ptree}
\AXC{$\varphi_{387}$\lbl{A1}}
\AXC{$\varphi_{387} \to \mathrm{SeqBdd}(a) \to \varphi_{356}$\lbl{A1}}
\BIC{$\mathrm{SeqBdd}(a) \to \varphi_{356}$}
\AXC{$\stackrel{\Pi_{100}}{\vdots}$}\noLine
\UIC{$(\mathrm{SeqBdd}(a) \to \varphi_{356}) \to \mathrm{SeqBdd}(a) \to \varphi_{348}$}
\BIC{$\mathrm{SeqBdd}(a) \to \varphi_{348}$}
\end{ptree}
\noindent$\Pi_{102}$:
\begin{ptree}
\AXC{$\stackrel{\Pi_{98}}{\vdots}$}\noLine
\UIC{$\varphi_{385}$}
\AXC{$\stackrel{\Pi_{101}}{\vdots}$}\noLine
\UIC{$\mathrm{SeqBdd}(a) \to \varphi_{348}$}
\AXC{$\varphi_{387}$\lbl{A1}}
\AXC{$\varphi_{387} \to \mathrm{SeqBdd}(a) \to \varphi_{353}$\lbl{A1}}
\BIC{$\mathrm{SeqBdd}(a) \to \varphi_{353}$}
\AXC{$\varphi_{484}$\lbl{A1}}
\BIC{$(\mathrm{SeqBdd}(a) \to \varphi_{348}) \to \varphi_{385} \to \varphi_{389}$}
\BIC{$\varphi_{385} \to \varphi_{389}$}
\BIC{$\varphi_{389}$}
\AXC{$\varphi_{389} \to \varphi_{156} \to \varphi_{391}$\lbl{A1}}
\BIC{$\varphi_{156} \to \varphi_{391}$}
\end{ptree}
\noindent$\Pi_{103}$:
\begin{ptree}
\AXC{$\varphi_{337}$\lbl{A1}}
\AXC{$\forall m\,\forall n\,(m \in \mathbb{N} \to n \in \mathbb{N} \to m \cup n \in \mathbb{N})$\lbl{[N3]}}
\AXC{$\varphi_{116}$\lbl{A2}}
\BIC{$\forall n\,(N \in \mathbb{N} \to n \in \mathbb{N} \to N \cup n \in \mathbb{N})$}
\AXC{$\varphi_{97}$\lbl{A2}}
\BIC{$N \in \mathbb{N} \to N_{1} \in \mathbb{N} \to N \cup N_{1} \in \mathbb{N}$}
\AXC{$\varphi_{463}$\lbl{A1}}
\BIC{$\varphi_{337} \to (\mathrm{SeqBdd}(a) \to \varphi_{303}) \to \mathrm{SeqBdd}(a) \to \varphi_{335}$}
\BIC{$(\mathrm{SeqBdd}(a) \to \varphi_{303}) \to \mathrm{SeqBdd}(a) \to \varphi_{335}$}
\end{ptree}
\noindent$\Pi_{104}$:
\begin{ptree}
\AXC{$\mathrm{SeqBdd}(a) \to \varphi_{333}$\lbl{A1}}
\AXC{$(\mathrm{SeqBdd}(a) \to \varphi_{333}) \to \mathrm{SeqBdd}(a) \to \varphi_{303}$\lbl{A1}}
\BIC{$\mathrm{SeqBdd}(a) \to \varphi_{303}$}
\AXC{$\stackrel{\Pi_{103}}{\vdots}$}\noLine
\UIC{$(\mathrm{SeqBdd}(a) \to \varphi_{303}) \to \mathrm{SeqBdd}(a) \to \varphi_{335}$}
\BIC{$\mathrm{SeqBdd}(a) \to \varphi_{335}$}
\end{ptree}
\noindent$\Pi_{105}$:
\begin{ptree}
\AXC{$\forall x\,(\varphi_{42}) \to \varphi_{42}$\lbl{A2}}
\AXC{$\varphi_{278}$\lbl{A1}}
\AXC{$\varphi_{278} \to (\forall x\,(\varphi_{42}) \to \varphi_{42}) \to \varphi_{271}$\lbl{A1}}
\BIC{$(\forall x\,(\varphi_{42}) \to \varphi_{42}) \to \varphi_{271}$}
\BIC{$\varphi_{271}$}
\AXC{$\varphi_{271} \to \varphi_{325} \to \mathrm{SeqBdd}(a) \to \varphi_{342}$\lbl{A1}}
\BIC{$\varphi_{325} \to \mathrm{SeqBdd}(a) \to \varphi_{342}$}
\end{ptree}
\noindent$\Pi_{106}$:
\begin{ptree}
\AXC{$\mathrm{SeqBdd}(a) \to \varphi_{336}$\lbl{A1}}
\AXC{$(\mathrm{SeqBdd}(a) \to \varphi_{336}) \to \varphi_{325}$\lbl{A1}}
\BIC{$\varphi_{325}$}
\AXC{$\stackrel{\Pi_{105}}{\vdots}$}\noLine
\UIC{$\varphi_{325} \to \mathrm{SeqBdd}(a) \to \varphi_{342}$}
\BIC{$\mathrm{SeqBdd}(a) \to \varphi_{342}$}
\end{ptree}
\noindent$\Pi_{107}$:
\begin{ptree}
\AXC{$\varphi_{150}$\lbl{A2}}
\AXC{$\stackrel{\Pi_{106}}{\vdots}$}\noLine
\UIC{$\mathrm{SeqBdd}(a) \to \varphi_{342}$}
\AXC{$(\mathrm{SeqBdd}(a) \to \varphi_{342}) \to \varphi_{150} \to \mathrm{SeqBdd}(a) \to \varphi_{352}$\lbl{A1}}
\BIC{$\varphi_{150} \to \mathrm{SeqBdd}(a) \to \varphi_{352}$}
\BIC{$\mathrm{SeqBdd}(a) \to \varphi_{352}$}
\AXC{$\varphi_{469}$\lbl{A1}}
\BIC{$(\mathrm{SeqBdd}(a) \to \varphi_{335}) \to \mathrm{SeqBdd}(a) \to \varphi_{346}$}
\end{ptree}
\noindent$\Pi_{108}$:
\begin{ptree}
\AXC{$\stackrel{\Pi_{104}}{\vdots}$}\noLine
\UIC{$\mathrm{SeqBdd}(a) \to \varphi_{335}$}
\AXC{$\stackrel{\Pi_{107}}{\vdots}$}\noLine
\UIC{$(\mathrm{SeqBdd}(a) \to \varphi_{335}) \to \mathrm{SeqBdd}(a) \to \varphi_{346}$}
\BIC{$\mathrm{SeqBdd}(a) \to \varphi_{346}$}
\AXC{$(\mathrm{SeqBdd}(a) \to \varphi_{346}) \to \varphi_{492}$\lbl{A1}}
\BIC{$\varphi_{492}$}
\end{ptree}
\noindent$\Pi_{109}$:
\begin{ptree}
\AXC{$\varphi_{373} \to \forall n\,(\varphi_{373})$\lbl{A4}}
\AXC{$\varphi_{454}$\lbl{A3}}
\AXC{$\varphi_{156}$\lbl{A2}}
\AXC{$\stackrel{\Pi_{102}}{\vdots}$}\noLine
\UIC{$\varphi_{156} \to \varphi_{391}$}
\BIC{$\varphi_{391}$}
\AXC{$\varphi_{391} \to \varphi_{373} \to \neg (n \in \mathbb{N} \land N \cup N_{1} \leq n \land x \leq a_{n})$\lbl{A1}}
\BIC{$\varphi_{373} \to \neg (n \in \mathbb{N} \land N \cup N_{1} \leq n \land x \leq a_{n})$}
\UIC{$\forall n\,(\varphi_{373} \to \neg (n \in \mathbb{N} \land N \cup N_{1} \leq n \land x \leq a_{n}))$}
\AXC{$\stackrel{\Pi_{108}}{\vdots}$}\noLine
\UIC{$\varphi_{492}$}
\BIC{$\varphi_{454} \to (\varphi_{373} \to \forall n\,(\varphi_{373})) \to \mathrm{SeqBdd}(a) \to \varphi_{358}$}
\BIC{$(\varphi_{373} \to \forall n\,(\varphi_{373})) \to \mathrm{SeqBdd}(a) \to \varphi_{358}$}
\BIC{$\mathrm{SeqBdd}(a) \to \varphi_{358}$}
\end{ptree}
\noindent$\Pi_{110}$:
\begin{ptree}
\AXC{$\forall N\,(\varphi_{372}) \to \varphi_{376}$\lbl{A3}}
\AXC{$\stackrel{\Pi_{109}}{\vdots}$}\noLine
\UIC{$\mathrm{SeqBdd}(a) \to \varphi_{358}$}
\AXC{$(\mathrm{SeqBdd}(a) \to \varphi_{358}) \to \varphi_{372}$\lbl{A1}}
\BIC{$\varphi_{372}$}
\UIC{$\forall N\,(\varphi_{372})$}
\AXC{$\forall N\,(\varphi_{372}) \to \varphi_{486}$\lbl{A1}}
\BIC{$\varphi_{486}$}
\BIC{$(\varphi_{318} \to \forall N\,(\varphi_{318})) \to \mathrm{SeqBdd}(a) \to \varphi_{364}$}
\end{ptree}
\noindent$\Pi_{111}$:
\begin{ptree}
\AXC{$\stackrel{\Pi_{85}}{\vdots}$}\noLine
\UIC{$\mathrm{SeqBdd}(a) \to \varphi_{236}$}
\AXC{$\forall x \in \mathbb{R}\,(0 < x \to 0 < x^{-1})$\lbl{[F34]}}
\AXC{$\varphi_{107}$\lbl{A2}}
\BIC{$\underline{k} \in \mathbb{R} \to 0 < \underline{k} \to 0 < \underline{k}^{-1}$}
\AXC{$(\underline{k} \in \mathbb{R} \to 0 < \underline{k} \to 0 < \underline{k}^{-1}) \to \varphi_{423}$\lbl{A1}}
\BIC{$\varphi_{423}$}
\BIC{$(\mathrm{SeqBdd}(a) \to \varphi_{230}) \to \mathrm{SeqBdd}(a) \to \varphi_{237}$}
\end{ptree}
\noindent$\Pi_{112}$:
\begin{ptree}
\AXC{$\stackrel{\Pi_{81}}{\vdots}$}\noLine
\UIC{$\mathrm{SeqBdd}(a) \to \varphi_{230}$}
\AXC{$\stackrel{\Pi_{111}}{\vdots}$}\noLine
\UIC{$(\mathrm{SeqBdd}(a) \to \varphi_{230}) \to \mathrm{SeqBdd}(a) \to \varphi_{237}$}
\BIC{$\mathrm{SeqBdd}(a) \to \varphi_{237}$}
\end{ptree}
\noindent$\Pi_{113}$:
\begin{ptree}
\AXC{$\forall x\,(\varphi_{42}) \to \varphi_{43}$\lbl{A2}}
\AXC{$\varphi_{278}$\lbl{A1}}
\AXC{$\varphi_{278} \to (\forall x\,(\varphi_{42}) \to \varphi_{43}) \to \varphi_{272}$\lbl{A1}}
\BIC{$(\forall x\,(\varphi_{42}) \to \varphi_{43}) \to \varphi_{272}$}
\BIC{$\varphi_{272}$}
\AXC{$\varphi_{272} \to (\mathrm{SeqBdd}(a) \to \varphi_{201}) \to \mathrm{SeqBdd}(a) \to \varphi_{206}$\lbl{A1}}
\BIC{$(\mathrm{SeqBdd}(a) \to \varphi_{201}) \to \mathrm{SeqBdd}(a) \to \varphi_{206}$}
\end{ptree}
\noindent$\Pi_{114}$:
\begin{ptree}
\AXC{$\mathrm{SeqBdd}(a) \to \varphi_{201}$\lbl{A1}}
\AXC{$\stackrel{\Pi_{113}}{\vdots}$}\noLine
\UIC{$(\mathrm{SeqBdd}(a) \to \varphi_{201}) \to \mathrm{SeqBdd}(a) \to \varphi_{206}$}
\BIC{$\mathrm{SeqBdd}(a) \to \varphi_{206}$}
\AXC{$(\mathrm{SeqBdd}(a) \to \varphi_{206}) \to \varphi_{182} \to u \in E \to u \in \mathbb{R}$\lbl{A1}}
\BIC{$\varphi_{182} \to u \in E \to u \in \mathbb{R}$}
\UIC{$\forall u\,(\varphi_{182} \to u \in E \to u \in \mathbb{R})$}
\end{ptree}
\noindent$\Pi_{115}$:
\begin{ptree}
\AXC{$\varphi_{182} \to \forall u\,(\varphi_{182})$\lbl{A4}}
\AXC{$\varphi_{365}$\lbl{A3}}
\AXC{$\stackrel{\Pi_{114}}{\vdots}$}\noLine
\UIC{$\forall u\,(\varphi_{182} \to u \in E \to u \in \mathbb{R})$}
\AXC{$\varphi_{467}$\lbl{A1}}
\BIC{$\varphi_{365} \to (\varphi_{182} \to \forall u\,(\varphi_{182})) \to \mathrm{SeqBdd}(a) \to \varphi_{215}$}
\BIC{$(\varphi_{182} \to \forall u\,(\varphi_{182})) \to \mathrm{SeqBdd}(a) \to \varphi_{215}$}
\BIC{$\mathrm{SeqBdd}(a) \to \varphi_{215}$}
\end{ptree}
\noindent$\Pi_{116}$:
\begin{ptree}
\AXC{$\stackrel{\Pi_{115}}{\vdots}$}\noLine
\UIC{$\mathrm{SeqBdd}(a) \to \varphi_{215}$}
\AXC{$\forall x\,\forall y\,(x \subseteq y \leftrightarrow \forall u \in x\,(u \in y))$\lbl{Def$\subseteq$}}
\AXC{$\varphi_{52}$\lbl{A2}}
\BIC{$\forall y\,(E \subseteq y \leftrightarrow \forall u \in E\,(u \in y))$}
\AXC{$\varphi_{53}$\lbl{A2}}
\BIC{$E \subseteq \mathbb{R} \leftrightarrow \forall u \in E\,(u \in \mathbb{R})$}
\AXC{$\varphi_{386}$\lbl{A1}}
\BIC{$(\mathrm{SeqBdd}(a) \to \varphi_{215}) \to \varphi_{210}$}
\BIC{$\varphi_{210}$}
\end{ptree}
\noindent$\Pi_{117}$:
\begin{ptree}
\AXC{$\stackrel{\Pi_{116}}{\vdots}$}\noLine
\UIC{$\varphi_{210}$}
\AXC{$\varphi_{76}$\lbl{[C22]}}
\AXC{$\varphi_{76} \to \forall E\,(\varphi_{61})$\lbl{A2}}
\BIC{$\forall E\,(\varphi_{61})$}
\AXC{$\forall E\,(\varphi_{61}) \to \varphi_{61}$\lbl{A2}}
\BIC{$\varphi_{61}$}
\AXC{$\varphi_{61} \to E \subseteq \mathbb{R} \to \varphi_{117}$\lbl{A2}}
\BIC{$E \subseteq \mathbb{R} \to \varphi_{117}$}
\AXC{$(E \subseteq \mathbb{R} \to \varphi_{117}) \to \varphi_{210} \to \varphi_{457}$\lbl{A1}}
\BIC{$\varphi_{210} \to \varphi_{457}$}
\BIC{$\varphi_{457}$}
\end{ptree}
\noindent$\Pi_{118}$:
\begin{ptree}
\AXC{$\stackrel{\Pi_{112}}{\vdots}$}\noLine
\UIC{$\mathrm{SeqBdd}(a) \to \varphi_{237}$}
\AXC{$\stackrel{\Pi_{87}}{\vdots}$}\noLine
\UIC{$\mathrm{SeqBdd}(a) \to \varphi_{242}$}
\AXC{$\mathrm{SeqBdd}(a) \to \varphi_{212}$\lbl{A1}}
\AXC{$\stackrel{\Pi_{117}}{\vdots}$}\noLine
\UIC{$\varphi_{457}$}
\BIC{$(\mathrm{SeqBdd}(a) \to \varphi_{242}) \to (\mathrm{SeqBdd}(a) \to \varphi_{237}) \to \varphi_{319}$}
\BIC{$(\mathrm{SeqBdd}(a) \to \varphi_{237}) \to \varphi_{319}$}
\BIC{$\varphi_{319}$}
\AXC{$\varphi_{319} \to \varphi_{490}$\lbl{A1}}
\BIC{$\varphi_{490}$}
\end{ptree}
\noindent$\Pi_{119}$:
\begin{ptree}
\AXC{$\varphi_{343} \to \forall x\,(\varphi_{343})$\lbl{A4}}
\AXC{$\varphi_{446}$\lbl{A3}}
\AXC{$\varphi_{318} \to \forall N\,(\varphi_{318})$\lbl{A4}}
\AXC{$\stackrel{\Pi_{110}}{\vdots}$}\noLine
\UIC{$(\varphi_{318} \to \forall N\,(\varphi_{318})) \to \mathrm{SeqBdd}(a) \to \varphi_{364}$}
\BIC{$\mathrm{SeqBdd}(a) \to \varphi_{364}$}
\AXC{$\varphi_{443}$\lbl{A1}}
\BIC{$\varphi_{343} \to \neg (x \in E \land c - \underline{k}^{-1} < x)$}
\UIC{$\forall x\,(\varphi_{343} \to \neg (x \in E \land c - \underline{k}^{-1} < x))$}
\AXC{$\stackrel{\Pi_{118}}{\vdots}$}\noLine
\UIC{$\varphi_{490}$}
\BIC{$\varphi_{446} \to (\varphi_{343} \to \forall x\,(\varphi_{343})) \to \mathrm{SeqBdd}(a) \to \varphi_{338}$}
\BIC{$(\varphi_{343} \to \forall x\,(\varphi_{343})) \to \mathrm{SeqBdd}(a) \to \varphi_{338}$}
\BIC{$\mathrm{SeqBdd}(a) \to \varphi_{338}$}
\end{ptree}
\noindent$\Pi_{120}$:
\begin{ptree}
\AXC{$\stackrel{\Pi_{36}}{\vdots}$}\noLine
\UIC{$\mathrm{SeqBdd}(a) \to \varphi_{213}$}
\AXC{$\forall x\,\forall y\,(x \in \mathbb{R} \to y \in \mathbb{R} \to x + y \in \mathbb{R})$\lbl{(R1)}}
\AXC{$\varphi_{124}$\lbl{A2}}
\BIC{$\forall y\,(c \in \mathbb{R} \to y \in \mathbb{R} \to c + y \in \mathbb{R})$}
\AXC{$\varphi_{152}$\lbl{A2}}
\BIC{$c \in \mathbb{R} \to \underline{k}^{-1} \in \mathbb{R} \to c + \underline{k}^{-1} \in \mathbb{R}$}
\AXC{$\varphi_{435}$\lbl{A1}}
\BIC{$(\mathrm{SeqBdd}(a) \to \varphi_{213}) \to \varphi_{394}$}
\BIC{$\varphi_{394}$}
\end{ptree}
\noindent$\Pi_{121}$:
\begin{ptree}
\AXC{$\varphi_{10}$\lbl{DefSup}}
\AXC{$\varphi_{10} \to \varphi_{9}$\lbl{A2}}
\BIC{$\varphi_{9}$}
\AXC{$\varphi_{9} \to \varphi_{8}$\lbl{A2}}
\BIC{$\varphi_{8}$}
\AXC{$\mathrm{SeqBdd}(a) \to \varphi_{212}$\lbl{A1}}
\AXC{$(\mathrm{SeqBdd}(a) \to \varphi_{212}) \to \varphi_{8} \to \varphi_{276}$\lbl{A1}}
\BIC{$\varphi_{8} \to \varphi_{276}$}
\BIC{$\varphi_{276}$}
\end{ptree}
\noindent$\Pi_{122}$:
\begin{ptree}
\AXC{$\forall a\,\forall X\,(\mathrm{UB}(a,X) \leftrightarrow a \in \mathbb{R} \land \forall x \in X\,(x \leq a))$\lbl{DefUB}}
\AXC{$\varphi_{106}$\lbl{A2}}
\BIC{$\forall X\,(\mathrm{UB}(c,X) \leftrightarrow c \in \mathbb{R} \land \forall x \in X\,(x \leq c))$}
\AXC{$\varphi_{84}$\lbl{A2}}
\BIC{$\mathrm{UB}(c,E) \leftrightarrow c \in \mathbb{R} \land \forall x \in E\,(x \leq c)$}
\AXC{$\stackrel{\Pi_{121}}{\vdots}$}\noLine
\UIC{$\varphi_{276}$}
\BIC{$\mathrm{SeqBdd}(a) \to \varphi_{222}$}
\end{ptree}
\noindent$\Pi_{123}$:
\begin{ptree}
\AXC{$\forall x \in E\,(x \leq c) \to c + \underline{k}^{-1} \in E \to c + \underline{k}^{-1} \leq c$\lbl{A2}}
\AXC{$\stackrel{\Pi_{122}}{\vdots}$}\noLine
\UIC{$\mathrm{SeqBdd}(a) \to \varphi_{222}$}
\AXC{$(\mathrm{SeqBdd}(a) \to \varphi_{222}) \to \varphi_{320}$\lbl{A1}}
\BIC{$\varphi_{320}$}
\BIC{$\mathrm{SeqBdd}(a) \to M \in \mathbb{R} \land \forall n \in \mathbb{N}\,(\lvert a_{n} \rvert \leq M) \to \varphi_{191}$}
\end{ptree}
\noindent$\Pi_{124}$:
\begin{ptree}
\AXC{$\forall x\,\forall y\,(x \leq y \to \neg y < x)$\lbl{[O7]}}
\AXC{$\varphi_{45}$\lbl{A2}}
\BIC{$\forall y\,(c + \underline{k}^{-1} \leq y \to \neg y < c + \underline{k}^{-1})$}
\AXC{$\varphi_{129}$\lbl{A2}}
\BIC{$c + \underline{k}^{-1} \leq c \to \neg c < c + \underline{k}^{-1}$}
\AXC{$(c + \underline{k}^{-1} \leq c \to \neg c < c + \underline{k}^{-1}) \to \varphi_{279} \to \varphi_{289}$\lbl{A1}}
\BIC{$\varphi_{279} \to \varphi_{289}$}
\end{ptree}
\noindent$\Pi_{125}$:
\begin{ptree}
\AXC{$\varphi_{293}$\lbl{A1}}
\AXC{$\varphi_{293} \to \varphi_{288}$\lbl{A1}}
\BIC{$\varphi_{288}$}
\AXC{$\stackrel{\Pi_{123}}{\vdots}$}\noLine
\UIC{$\mathrm{SeqBdd}(a) \to M \in \mathbb{R} \land \forall n \in \mathbb{N}\,(\lvert a_{n} \rvert \leq M) \to \varphi_{191}$}
\AXC{$\varphi_{436}$\lbl{A1}}
\BIC{$\varphi_{288} \to \varphi_{279}$}
\BIC{$\varphi_{279}$}
\AXC{$\stackrel{\Pi_{124}}{\vdots}$}\noLine
\UIC{$\varphi_{279} \to \varphi_{289}$}
\BIC{$\varphi_{289}$}
\end{ptree}
\noindent$\Pi_{126}$:
\begin{ptree}
\AXC{$\stackrel{\Pi_{36}}{\vdots}$}\noLine
\UIC{$\mathrm{SeqBdd}(a) \to \varphi_{213}$}
\AXC{$\forall x\,\forall y\,(x \in \mathbb{R} \to y \in \mathbb{R} \to 0 < y \to x < x + y)$\lbl{[E3]}}
\AXC{$\varphi_{144}$\lbl{A2}}
\BIC{$\forall y\,(c \in \mathbb{R} \to y \in \mathbb{R} \to 0 < y \to c < c + y)$}
\AXC{$\varphi_{183}$\lbl{A2}}
\BIC{$c \in \mathbb{R} \to \underline{k}^{-1} \in \mathbb{R} \to 0 < \underline{k}^{-1} \to c < c + \underline{k}^{-1}$}
\AXC{$\varphi_{459}$\lbl{A1}}
\BIC{$\varphi_{452}$}
\BIC{$(\mathrm{SeqBdd}(a) \to \varphi_{242}) \to (\mathrm{SeqBdd}(a) \to \varphi_{237}) \to \varphi_{286}$}
\end{ptree}
\noindent$\Pi_{127}$:
\begin{ptree}
\AXC{$\stackrel{\Pi_{112}}{\vdots}$}\noLine
\UIC{$\mathrm{SeqBdd}(a) \to \varphi_{237}$}
\AXC{$\stackrel{\Pi_{87}}{\vdots}$}\noLine
\UIC{$\mathrm{SeqBdd}(a) \to \varphi_{242}$}
\AXC{$\stackrel{\Pi_{126}}{\vdots}$}\noLine
\UIC{$(\mathrm{SeqBdd}(a) \to \varphi_{242}) \to (\mathrm{SeqBdd}(a) \to \varphi_{237}) \to \varphi_{286}$}
\BIC{$(\mathrm{SeqBdd}(a) \to \varphi_{237}) \to \varphi_{286}$}
\BIC{$\varphi_{286}$}
\end{ptree}
\noindent$\Pi_{128}$:
\begin{ptree}
\AXC{$\stackrel{\Pi_{125}}{\vdots}$}\noLine
\UIC{$\varphi_{289}$}
\AXC{$\stackrel{\Pi_{127}}{\vdots}$}\noLine
\UIC{$\varphi_{286}$}
\AXC{$\varphi_{286} \to \varphi_{289} \to \mathrm{SeqBdd}(a) \to \varphi_{248}$\lbl{A1}}
\BIC{$\varphi_{289} \to \mathrm{SeqBdd}(a) \to \varphi_{248}$}
\BIC{$\mathrm{SeqBdd}(a) \to \varphi_{248}$}
\AXC{$(\mathrm{SeqBdd}(a) \to \varphi_{248}) \to \varphi_{316} \to \varphi_{410}$\lbl{A1}}
\BIC{$\varphi_{316} \to \varphi_{410}$}
\end{ptree}
\noindent$\Pi_{129}$:
\begin{ptree}
\AXC{$\forall x\,(\varphi_{42}) \to \varphi_{161}$\lbl{A2}}
\AXC{$\varphi_{278}$\lbl{A1}}
\AXC{$\varphi_{278} \to (\forall x\,(\varphi_{42}) \to \varphi_{161}) \to \varphi_{316}$\lbl{A1}}
\BIC{$(\forall x\,(\varphi_{42}) \to \varphi_{161}) \to \varphi_{316}$}
\BIC{$\varphi_{316}$}
\AXC{$\stackrel{\Pi_{128}}{\vdots}$}\noLine
\UIC{$\varphi_{316} \to \varphi_{410}$}
\BIC{$\varphi_{410}$}
\end{ptree}
\noindent$\Pi_{130}$:
\begin{ptree}
\AXC{$\stackrel{\Pi_{87}}{\vdots}$}\noLine
\UIC{$\mathrm{SeqBdd}(a) \to \varphi_{242}$}
\AXC{$\stackrel{\Pi_{120}}{\vdots}$}\noLine
\UIC{$\varphi_{394}$}
\BIC{$\mathrm{SeqBdd}(a) \to M \in \mathbb{R} \land \forall n \in \mathbb{N}\,(\lvert a_{n} \rvert \leq M) \to \varphi_{172}$}
\AXC{$\stackrel{\Pi_{129}}{\vdots}$}\noLine
\UIC{$\varphi_{410}$}
\BIC{$\mathrm{SeqBdd}(a) \to \varphi_{290}$}
\AXC{$(\mathrm{SeqBdd}(a) \to \varphi_{290}) \to \forall N\,(\varphi_{170}) \to \varphi_{437}$\lbl{A1}}
\BIC{$\forall N\,(\varphi_{170}) \to \varphi_{437}$}
\end{ptree}
\noindent$\Pi_{131}$:
\begin{ptree}
\AXC{$\forall N\,(\varphi_{23}) \to \forall N\,\forall N\,(\varphi_{23})$\lbl{A4}}
\AXC{$\forall N\,(\varphi_{170}) \to \varphi_{188}$\lbl{A3}}
\AXC{$\forall N\,(\varphi_{23}) \to \varphi_{23}$\lbl{A2}}
\AXC{$(\forall N\,(\varphi_{23}) \to \varphi_{23}) \to \varphi_{170}$\lbl{A1}}
\BIC{$\varphi_{170}$}
\UIC{$\forall N\,(\varphi_{170})$}
\AXC{$\stackrel{\Pi_{130}}{\vdots}$}\noLine
\UIC{$\forall N\,(\varphi_{170}) \to \varphi_{437}$}
\BIC{$\varphi_{437}$}
\BIC{$(\forall N\,(\varphi_{23}) \to \forall N\,\forall N\,(\varphi_{23})) \to \mathrm{SeqBdd}(a) \to \varphi_{291}$}
\BIC{$\mathrm{SeqBdd}(a) \to \varphi_{291}$}
\end{ptree}
\noindent$\Pi_{132}$:
\begin{ptree}
\AXC{$\stackrel{\Pi_{131}}{\vdots}$}\noLine
\UIC{$\mathrm{SeqBdd}(a) \to \varphi_{291}$}
\AXC{$(\mathrm{SeqBdd}(a) \to \varphi_{291}) \to \forall N_{1}\,(\varphi_{309} \to \neg \varphi_{28}) \to \varphi_{489}$\lbl{A1}}
\BIC{$\forall N_{1}\,(\varphi_{309} \to \neg \varphi_{28}) \to \varphi_{489}$}
\end{ptree}
\noindent$\Pi_{133}$:
\begin{ptree}
\AXC{$\stackrel{\Pi_{119}}{\vdots}$}\noLine
\UIC{$\mathrm{SeqBdd}(a) \to \varphi_{338}$}
\AXC{$(\mathrm{SeqBdd}(a) \to \varphi_{338}) \to \varphi_{309} \to \neg \varphi_{28}$\lbl{A1}}
\BIC{$\varphi_{309} \to \neg \varphi_{28}$}
\UIC{$\forall N_{1}\,(\varphi_{309} \to \neg \varphi_{28})$}
\AXC{$\stackrel{\Pi_{132}}{\vdots}$}\noLine
\UIC{$\forall N_{1}\,(\varphi_{309} \to \neg \varphi_{28}) \to \varphi_{489}$}
\BIC{$\varphi_{489}$}
\end{ptree}
\noindent$\Pi_{134}$:
\begin{ptree}
\AXC{$\forall N_{1}\,(\neg \varphi_{28}) \to \varphi_{23}$\lbl{A2}}
\UIC{$\forall N\,(\forall N_{1}\,(\neg \varphi_{28}) \to \varphi_{23})$}
\AXC{$\varphi_{309} \to \forall N_{1}\,(\varphi_{309})$\lbl{A4}}
\AXC{$\forall N_{1}\,(\varphi_{309} \to \neg \varphi_{28}) \to \forall N_{1}\,(\varphi_{309}) \to \forall N_{1}\,(\neg \varphi_{28})$\lbl{A3}}
\AXC{$\stackrel{\Pi_{133}}{\vdots}$}\noLine
\UIC{$\varphi_{489}$}
\BIC{$(\varphi_{309} \to \forall N_{1}\,(\varphi_{309})) \to \forall N\,(\forall N_{1}\,(\neg \varphi_{28}) \to \varphi_{23}) \to \varphi_{439}$}
\BIC{$\forall N\,(\forall N_{1}\,(\neg \varphi_{28}) \to \varphi_{23}) \to \varphi_{439}$}
\BIC{$\varphi_{439}$}
\end{ptree}
\noindent$\Pi_{135}$:
\begin{ptree}
\AXC{$\forall N_{1}\,(\neg \varphi_{28}) \to \forall N\,\forall N_{1}\,(\neg \varphi_{28})$\lbl{A4}}
\AXC{$\forall N\,(\forall N_{1}\,(\neg \varphi_{28}) \to \varphi_{23}) \to \forall N\,\forall N_{1}\,(\neg \varphi_{28}) \to \forall N\,(\varphi_{23})$\lbl{A3}}
\AXC{$\stackrel{\Pi_{134}}{\vdots}$}\noLine
\UIC{$\varphi_{439}$}
\BIC{$\varphi_{398}$}
\BIC{$\mathrm{SeqBdd}(a) \to \varphi_{298}$}
\end{ptree}
\noindent$\Pi_{136}$:
\begin{ptree}
\AXC{$\stackrel{\Pi_{135}}{\vdots}$}\noLine
\UIC{$\mathrm{SeqBdd}(a) \to \varphi_{298}$}
\AXC{$(\mathrm{SeqBdd}(a) \to \varphi_{298}) \to \varphi_{182} \land \mathrm{IsSup}(c,E) \to \varphi_{82}$\lbl{A1}}
\BIC{$\varphi_{182} \land \mathrm{IsSup}(c,E) \to \varphi_{82}$}
\UIC{$\forall k\,(\varphi_{182} \land \mathrm{IsSup}(c,E) \to \varphi_{82})$}
\AXC{$\forall k\,(\varphi_{182} \land \mathrm{IsSup}(c,E) \to \varphi_{82}) \to \varphi_{415} \to \varphi_{424}$\lbl{A1}}
\BIC{$\varphi_{415} \to \varphi_{424}$}
\end{ptree}
\noindent$\Pi_{137}$:
\begin{ptree}
\AXC{$\varphi_{182} \land \mathrm{IsSup}(c,E) \to \forall k\,(\varphi_{182} \land \mathrm{IsSup}(c,E))$\lbl{A4}}
\AXC{$\varphi_{415}$\lbl{A3}}
\AXC{$\stackrel{\Pi_{136}}{\vdots}$}\noLine
\UIC{$\varphi_{415} \to \varphi_{424}$}
\BIC{$\varphi_{424}$}
\BIC{$\mathrm{SeqBdd}(a) \to \varphi_{302}$}
\end{ptree}
\noindent$\Pi_{138}$:
\begin{ptree}
\AXC{$\forall a\,\forall c\,(\mathrm{Seq}(a) \to \varphi_{93} \to \varphi_{47})$\lbl{[S28]}}
\AXC{$\varphi_{388}$\lbl{A2}}
\BIC{$\forall c\,(\mathrm{Seq}(a) \to \varphi_{93} \to \varphi_{47})$}
\AXC{$\forall c\,(\mathrm{Seq}(a) \to \varphi_{93} \to \varphi_{47}) \to \mathrm{Seq}(a) \to \varphi_{93} \to \varphi_{47}$\lbl{A2}}
\BIC{$\mathrm{Seq}(a) \to \varphi_{93} \to \varphi_{47}$}
\end{ptree}
\noindent$\Pi_{139}$:
\begin{ptree}
\AXC{$\stackrel{\Pi_{16}}{\vdots}$}\noLine
\UIC{$\mathrm{SeqBdd}(a) \to \mathrm{Seq}(a)$}
\AXC{$\stackrel{\Pi_{138}}{\vdots}$}\noLine
\UIC{$\mathrm{Seq}(a) \to \varphi_{93} \to \varphi_{47}$}
\AXC{$(\mathrm{Seq}(a) \to \varphi_{93} \to \varphi_{47}) \to \varphi_{416}$\lbl{A1}}
\BIC{$\varphi_{416}$}
\BIC{$(\mathrm{SeqBdd}(a) \to \varphi_{302}) \to \mathrm{SeqBdd}(a) \to \varphi_{283}$}
\end{ptree}
\noindent$\Pi_{140}$:
\begin{ptree}
\AXC{$\stackrel{\Pi_{137}}{\vdots}$}\noLine
\UIC{$\mathrm{SeqBdd}(a) \to \varphi_{302}$}
\AXC{$\stackrel{\Pi_{139}}{\vdots}$}\noLine
\UIC{$(\mathrm{SeqBdd}(a) \to \varphi_{302}) \to \mathrm{SeqBdd}(a) \to \varphi_{283}$}
\BIC{$\mathrm{SeqBdd}(a) \to \varphi_{283}$}
\AXC{$(\mathrm{SeqBdd}(a) \to \varphi_{283}) \to \forall s\,(\varphi_{247} \to \varphi_{41}) \to \varphi_{475}$\lbl{A1}}
\BIC{$\forall s\,(\varphi_{247} \to \varphi_{41}) \to \varphi_{475}$}
\end{ptree}
\noindent$\Pi_{141}$:
\begin{ptree}
\AXC{$\stackrel{\Pi_{73}}{\vdots}$}\noLine
\UIC{$\mathrm{SeqBdd}(a) \to \varphi_{307}$}
\AXC{$(\mathrm{SeqBdd}(a) \to \varphi_{307}) \to \varphi_{247} \to \varphi_{41}$\lbl{A1}}
\BIC{$\varphi_{247} \to \varphi_{41}$}
\UIC{$\forall s\,(\varphi_{247} \to \varphi_{41})$}
\AXC{$\stackrel{\Pi_{140}}{\vdots}$}\noLine
\UIC{$\forall s\,(\varphi_{247} \to \varphi_{41}) \to \varphi_{475}$}
\BIC{$\varphi_{475}$}
\end{ptree}
\noindent$\Pi_{142}$:
\begin{ptree}
\AXC{$\varphi_{247} \to \forall s\,(\varphi_{247})$\lbl{A4}}
\AXC{$\forall s\,(\varphi_{247} \to \varphi_{41}) \to \forall s\,(\varphi_{247}) \to \forall s\,(\varphi_{41})$\lbl{A3}}
\AXC{$\stackrel{\Pi_{141}}{\vdots}$}\noLine
\UIC{$\varphi_{475}$}
\BIC{$(\varphi_{247} \to \forall s\,(\varphi_{247})) \to \varphi_{245}$}
\BIC{$\varphi_{245}$}
\AXC{$\varphi_{245} \to \varphi_{249}$\lbl{A1}}
\BIC{$\varphi_{249}$}
\UIC{$\forall c\,(\varphi_{249})$}
\end{ptree}
\noindent$\Pi_{143}$:
\begin{ptree}
\AXC{$\forall n \in \mathbb{N}\,(\lvert a_{n} \rvert \leq M) \to n \in \mathbb{N} \to \lvert a_{n} \rvert \leq M$\lbl{A2}}
\AXC{$\mathrm{SeqBdd}(a) \to \varphi_{105}$\lbl{A1}}
\AXC{$(\mathrm{SeqBdd}(a) \to \varphi_{105}) \to \varphi_{284}$\lbl{A1}}
\BIC{$\varphi_{284}$}
\AXC{$\varphi_{417}$\lbl{A1}}
\BIC{$(\forall n \in \mathbb{N}\,(\lvert a_{n} \rvert \leq M) \to n \in \mathbb{N} \to \lvert a_{n} \rvert \leq M) \to \varphi_{277}$}
\BIC{$\varphi_{277}$}
\end{ptree}
\noindent$\Pi_{144}$:
\begin{ptree}
\AXC{$\varphi_{287}$\lbl{A1}}
\AXC{$\varphi_{405}$\lbl{A1}}
\BIC{$\mathrm{SeqBdd}(a) \to M \in \mathbb{R} \land \forall n \in \mathbb{N}\,(\lvert a_{n} \rvert \leq M) \to \varphi_{175}$}
\AXC{$\stackrel{\Pi_{143}}{\vdots}$}\noLine
\UIC{$\varphi_{277}$}
\AXC{$\varphi_{277} \to \varphi_{401}$\lbl{A1}}
\BIC{$\varphi_{401}$}
\BIC{$\mathrm{SeqBdd}(a) \to M \in \mathbb{R} \land \forall n \in \mathbb{N}\,(\lvert a_{n} \rvert \leq M) \to \varphi_{192}$}
\end{ptree}
\noindent$\Pi_{145}$:
\begin{ptree}
\AXC{$\stackrel{\Pi_{16}}{\vdots}$}\noLine
\UIC{$\mathrm{SeqBdd}(a) \to \mathrm{Seq}(a)$}
\AXC{$\forall a\,\forall n\,(\mathrm{Seq}(a) \to n \in \mathbb{N} \to a_{n} \in \mathbb{R})$\lbl{[S1]}}
\AXC{$\varphi_{96}$\lbl{A2}}
\BIC{$\forall n\,(\mathrm{Seq}(a) \to n \in \mathbb{N} \to a_{n} \in \mathbb{R})$}
\AXC{$\varphi_{75}$\lbl{A2}}
\BIC{$\mathrm{Seq}(a) \to n \in \mathbb{N} \to a_{n} \in \mathbb{R}$}
\AXC{$\varphi_{412}$\lbl{A1}}
\BIC{$(\mathrm{SeqBdd}(a) \to \mathrm{Seq}(a)) \to \varphi_{402}$}
\BIC{$\varphi_{402}$}
\end{ptree}
\noindent$\Pi_{146}$:
\begin{ptree}
\AXC{$\forall x\,\forall y\,(x \in \mathbb{R} \to y \in \mathbb{R} \to \lvert x \rvert \leq y \to -y \leq x \land x \leq y)$\lbl{[A9]}}
\AXC{$\varphi_{207}$\lbl{A2}}
\BIC{$\varphi_{77}$}
\end{ptree}
\noindent$\Pi_{147}$:
\begin{ptree}
\AXC{$\stackrel{\Pi_{146}}{\vdots}$}\noLine
\UIC{$\varphi_{77}$}
\AXC{$\varphi_{77} \to a_{n} \in \mathbb{R} \to M \in \mathbb{R} \to \lvert a_{n} \rvert \leq M \to -M \leq a_{n} \land a_{n} \leq M$\lbl{A2}}
\BIC{$a_{n} \in \mathbb{R} \to M \in \mathbb{R} \to \lvert a_{n} \rvert \leq M \to -M \leq a_{n} \land a_{n} \leq M$}
\AXC{$\varphi_{455}$\lbl{A1}}
\BIC{$\varphi_{445}$}
\end{ptree}
\noindent$\Pi_{148}$:
\begin{ptree}
\AXC{$\varphi_{287}$\lbl{A1}}
\AXC{$\varphi_{405}$\lbl{A1}}
\BIC{$\mathrm{SeqBdd}(a) \to M \in \mathbb{R} \land \forall n \in \mathbb{N}\,(\lvert a_{n} \rvert \leq M) \to \varphi_{175}$}
\AXC{$\stackrel{\Pi_{145}}{\vdots}$}\noLine
\UIC{$\varphi_{402}$}
\BIC{$\mathrm{SeqBdd}(a) \to M \in \mathbb{R} \land \forall n \in \mathbb{N}\,(\lvert a_{n} \rvert \leq M) \to \varphi_{199}$}
\AXC{$\stackrel{\Pi_{147}}{\vdots}$}\noLine
\UIC{$\varphi_{445}$}
\BIC{$\varphi_{40} \to \varphi_{407}$}
\end{ptree}
\noindent$\Pi_{149}$:
\begin{ptree}
\AXC{$\mathrm{SeqBdd}(a) \to \varphi_{105}$\lbl{A1}}
\AXC{$(\mathrm{SeqBdd}(a) \to \varphi_{105}) \to \varphi_{40}$\lbl{A1}}
\BIC{$\varphi_{40}$}
\AXC{$\stackrel{\Pi_{148}}{\vdots}$}\noLine
\UIC{$\varphi_{40} \to \varphi_{407}$}
\BIC{$\varphi_{407}$}
\end{ptree}
\noindent$\Pi_{150}$:
\begin{ptree}
\AXC{$\stackrel{\Pi_{144}}{\vdots}$}\noLine
\UIC{$\mathrm{SeqBdd}(a) \to M \in \mathbb{R} \land \forall n \in \mathbb{N}\,(\lvert a_{n} \rvert \leq M) \to \varphi_{192}$}
\AXC{$\stackrel{\Pi_{149}}{\vdots}$}\noLine
\UIC{$\varphi_{407}$}
\BIC{$\varphi_{280}$}
\AXC{$\varphi_{403}$\lbl{A1}}
\BIC{$\mathrm{SeqBdd}(a) \to M \in \mathbb{R} \land \forall n \in \mathbb{N}\,(\lvert a_{n} \rvert \leq M) \to \varphi_{177}$}
\end{ptree}
\noindent$\Pi_{151}$:
\begin{ptree}
\AXC{$\forall x\,\forall y\,\forall z\,(x \leq y \to y \leq z \to x \leq z)$\lbl{[O5]}}
\AXC{$\varphi_{12}$\lbl{A2}}
\BIC{$\forall y\,\forall z\,(x \leq y \to y \leq z \to x \leq z)$}
\AXC{$\varphi_{15}$\lbl{A2}}
\BIC{$\forall z\,(x \leq a_{n} \to a_{n} \leq z \to x \leq z)$}
\AXC{$\varphi_{31}$\lbl{A2}}
\BIC{$x \leq a_{n} \to a_{n} \leq M \to x \leq M$}
\AXC{$(x \leq a_{n} \to a_{n} \leq M \to x \leq M) \to \varphi_{431}$\lbl{A1}}
\BIC{$\varphi_{431}$}
\end{ptree}
\noindent$\Pi_{152}$:
\begin{ptree}
\AXC{$\stackrel{\Pi_{150}}{\vdots}$}\noLine
\UIC{$\mathrm{SeqBdd}(a) \to M \in \mathbb{R} \land \forall n \in \mathbb{N}\,(\lvert a_{n} \rvert \leq M) \to \varphi_{177}$}
\AXC{$\varphi_{287}$\lbl{A1}}
\AXC{$\varphi_{406}$\lbl{A1}}
\BIC{$\mathrm{SeqBdd}(a) \to M \in \mathbb{R} \land \forall n \in \mathbb{N}\,(\lvert a_{n} \rvert \leq M) \to \varphi_{178}$}
\AXC{$\stackrel{\Pi_{151}}{\vdots}$}\noLine
\UIC{$\varphi_{431}$}
\BIC{$\varphi_{397}$}
\BIC{$\mathrm{SeqBdd}(a) \to M \in \mathbb{R} \land \forall n \in \mathbb{N}\,(\lvert a_{n} \rvert \leq M) \to \varphi_{166}$}
\end{ptree}
\noindent$\Pi_{153}$:
\begin{ptree}
\AXC{$\mathrm{SeqBdd}(a) \to \varphi_{202}$\lbl{A1}}
\AXC{$\forall x\,(\varphi_{42}) \to \varphi_{42}$\lbl{A2}}
\AXC{$\varphi_{278}$\lbl{A1}}
\AXC{$\varphi_{278} \to (\forall x\,(\varphi_{42}) \to \varphi_{42}) \to \varphi_{271}$\lbl{A1}}
\BIC{$(\forall x\,(\varphi_{42}) \to \varphi_{42}) \to \varphi_{271}$}
\BIC{$\varphi_{271}$}
\AXC{$\varphi_{271} \to \varphi_{384}$\lbl{A1}}
\BIC{$\varphi_{384}$}
\BIC{$\mathrm{SeqBdd}(a) \to M \in \mathbb{R} \land \forall n \in \mathbb{N}\,(\lvert a_{n} \rvert \leq M) \to \varphi_{185}$}
\AXC{$\varphi_{418}$\lbl{A1}}
\BIC{$\varphi_{135} \to \varphi_{264}$}
\end{ptree}
\noindent$\Pi_{154}$:
\begin{ptree}
\AXC{$1_{\omega} \in \mathbb{N}$\lbl{[N1]}}
\AXC{$\varphi_{135}$\lbl{A2}}
\AXC{$\stackrel{\Pi_{153}}{\vdots}$}\noLine
\UIC{$\varphi_{135} \to \varphi_{264}$}
\BIC{$\varphi_{264}$}
\AXC{$\varphi_{264} \to \varphi_{265}$\lbl{A1}}
\BIC{$\varphi_{265}$}
\BIC{$\mathrm{SeqBdd}(a) \to M \in \mathbb{R} \land \forall n \in \mathbb{N}\,(\lvert a_{n} \rvert \leq M) \to \varphi_{164}$}
\AXC{$\varphi_{480}$\lbl{A1}}
\BIC{$\forall n\,(\varphi_{258}) \to (\forall n\,(\varphi_{258}) \to \varphi_{261}) \to \varphi_{414}$}
\end{ptree}
\noindent$\Pi_{155}$:
\begin{ptree}
\AXC{$\stackrel{\Pi_{152}}{\vdots}$}\noLine
\UIC{$\mathrm{SeqBdd}(a) \to M \in \mathbb{R} \land \forall n \in \mathbb{N}\,(\lvert a_{n} \rvert \leq M) \to \varphi_{166}$}
\AXC{$\varphi_{392}$\lbl{A1}}
\BIC{$\varphi_{258}$}
\UIC{$\forall n\,(\varphi_{258})$}
\AXC{$\stackrel{\Pi_{154}}{\vdots}$}\noLine
\UIC{$\forall n\,(\varphi_{258}) \to (\forall n\,(\varphi_{258}) \to \varphi_{261}) \to \varphi_{414}$}
\BIC{$(\forall n\,(\varphi_{258}) \to \varphi_{261}) \to \varphi_{414}$}
\end{ptree}
\noindent$\Pi_{156}$:
\begin{ptree}
\AXC{$\varphi_{182} \land x \in E \land \neg x \leq M \to \forall n\,(\varphi_{182} \land x \in E \land \neg x \leq M)$\lbl{A4}}
\AXC{$\forall n\,(\varphi_{258}) \to \varphi_{261}$\lbl{A3}}
\AXC{$\stackrel{\Pi_{155}}{\vdots}$}\noLine
\UIC{$(\forall n\,(\varphi_{258}) \to \varphi_{261}) \to \varphi_{414}$}
\BIC{$\varphi_{414}$}
\BIC{$\mathrm{SeqBdd}(a) \to \varphi_{193}$}
\end{ptree}
\noindent$\Pi_{157}$:
\begin{ptree}
\AXC{$\stackrel{\Pi_{156}}{\vdots}$}\noLine
\UIC{$\mathrm{SeqBdd}(a) \to \varphi_{193}$}
\AXC{$(\mathrm{SeqBdd}(a) \to \varphi_{193}) \to \varphi_{182} \to x \in E \to x \leq M$\lbl{A1}}
\BIC{$\varphi_{182} \to x \in E \to x \leq M$}
\UIC{$\forall x\,(\varphi_{182} \to x \in E \to x \leq M)$}
\AXC{$\forall x\,(\varphi_{182} \to x \in E \to x \leq M) \to \varphi_{450}$\lbl{A1}}
\BIC{$\varphi_{450}$}
\end{ptree}
\noindent$\Pi_{158}$:
\begin{ptree}
\AXC{$\varphi_{182} \to \forall x\,(\varphi_{182})$\lbl{A4}}
\AXC{$\forall x\,(\varphi_{182} \to x \in E \to x \leq M) \to \forall x\,(\varphi_{182}) \to \forall x \in E\,(x \leq M)$\lbl{A3}}
\AXC{$\stackrel{\Pi_{157}}{\vdots}$}\noLine
\UIC{$\varphi_{450}$}
\BIC{$(\varphi_{182} \to \forall x\,(\varphi_{182})) \to \mathrm{SeqBdd}(a) \to \varphi_{200}$}
\BIC{$\mathrm{SeqBdd}(a) \to \varphi_{200}$}
\end{ptree}
\noindent$\Pi_{159}$:
\begin{ptree}
\AXC{$\forall a\,\forall X\,(\mathrm{UB}(a,X) \leftrightarrow a \in \mathbb{R} \land \forall x \in X\,(x \leq a))$\lbl{DefUB}}
\AXC{$\varphi_{108}$\lbl{A2}}
\BIC{$\forall X\,(\mathrm{UB}(M,X) \leftrightarrow M \in \mathbb{R} \land \forall x \in X\,(x \leq M))$}
\AXC{$\varphi_{85}$\lbl{A2}}
\BIC{$\mathrm{UB}(M,E) \leftrightarrow M \in \mathbb{R} \land \forall x \in E\,(x \leq M)$}
\AXC{$\varphi_{409}$\lbl{A1}}
\BIC{$\varphi_{40} \to (\mathrm{SeqBdd}(a) \to \varphi_{200}) \to \mathrm{SeqBdd}(a) \to \varphi_{176}$}
\end{ptree}
\noindent$\Pi_{160}$:
\begin{ptree}
\AXC{$\mathrm{SeqBdd}(a) \to \varphi_{105}$\lbl{A1}}
\AXC{$(\mathrm{SeqBdd}(a) \to \varphi_{105}) \to \varphi_{40}$\lbl{A1}}
\BIC{$\varphi_{40}$}
\AXC{$\stackrel{\Pi_{159}}{\vdots}$}\noLine
\UIC{$\varphi_{40} \to (\mathrm{SeqBdd}(a) \to \varphi_{200}) \to \mathrm{SeqBdd}(a) \to \varphi_{176}$}
\BIC{$(\mathrm{SeqBdd}(a) \to \varphi_{200}) \to \mathrm{SeqBdd}(a) \to \varphi_{176}$}
\end{ptree}
\noindent$\Pi_{161}$:
\begin{ptree}
\AXC{$\stackrel{\Pi_{158}}{\vdots}$}\noLine
\UIC{$\mathrm{SeqBdd}(a) \to \varphi_{200}$}
\AXC{$\stackrel{\Pi_{160}}{\vdots}$}\noLine
\UIC{$(\mathrm{SeqBdd}(a) \to \varphi_{200}) \to \mathrm{SeqBdd}(a) \to \varphi_{176}$}
\BIC{$\mathrm{SeqBdd}(a) \to \varphi_{176}$}
\AXC{$(\mathrm{SeqBdd}(a) \to \varphi_{176}) \to \varphi_{243}$\lbl{A1}}
\BIC{$\varphi_{243}$}
\end{ptree}
\noindent$\Pi_{162}$:
\begin{ptree}
\AXC{$\mathrm{SeqBdd}(a) \to \varphi_{105}$\lbl{A1}}
\AXC{$(\mathrm{SeqBdd}(a) \to \varphi_{105}) \to \mathrm{SeqBdd}(a) \to \varphi_{238}$\lbl{A1}}
\BIC{$\mathrm{SeqBdd}(a) \to \varphi_{238}$}
\AXC{$(\mathrm{SeqBdd}(a) \to \varphi_{238}) \to \varphi_{312}$\lbl{A1}}
\BIC{$\varphi_{312}$}
\end{ptree}
\noindent$\Pi_{163}$:
\begin{ptree}
\AXC{$\mathrm{SeqBdd}(a) \to \varphi_{209}$\lbl{A1}}
\AXC{$\forall n \in \mathbb{N}\,(\lvert a_{n} \rvert \leq M) \to N \in \mathbb{N} \to \lvert a_{N} \rvert \leq M$\lbl{A2}}
\AXC{$\stackrel{\Pi_{162}}{\vdots}$}\noLine
\UIC{$\varphi_{312}$}
\BIC{$\mathrm{SeqBdd}(a) \to \varphi_{232}$}
\AXC{$\varphi_{419}$\lbl{A1}}
\BIC{$(\mathrm{SeqBdd}(a) \to \varphi_{209}) \to \mathrm{SeqBdd}(a) \to \varphi_{220}$}
\BIC{$\mathrm{SeqBdd}(a) \to \varphi_{220}$}
\end{ptree}
\noindent$\Pi_{164}$:
\begin{ptree}
\AXC{$\forall a\,\forall n\,(\mathrm{Seq}(a) \to n \in \mathbb{N} \to a_{n} \in \mathbb{R})$\lbl{[S1]}}
\AXC{$\varphi_{96}$\lbl{A2}}
\BIC{$\forall n\,(\mathrm{Seq}(a) \to n \in \mathbb{N} \to a_{n} \in \mathbb{R})$}
\AXC{$\varphi_{78}$\lbl{A2}}
\BIC{$\mathrm{Seq}(a) \to N \in \mathbb{N} \to a_{N} \in \mathbb{R}$}
\AXC{$(\mathrm{Seq}(a) \to N \in \mathbb{N} \to a_{N} \in \mathbb{R}) \to \varphi_{383}$\lbl{A1}}
\BIC{$\varphi_{383}$}
\end{ptree}
\noindent$\Pi_{165}$:
\begin{ptree}
\AXC{$\forall x\,\forall y\,(x \in \mathbb{R} \to y \in \mathbb{R} \to \lvert x \rvert \leq y \to -y \leq x \land x \leq y)$\lbl{[A9]}}
\AXC{$\varphi_{208}$\lbl{A2}}
\BIC{$\varphi_{79}$}
\end{ptree}
\noindent$\Pi_{166}$:
\begin{ptree}
\AXC{$\stackrel{\Pi_{165}}{\vdots}$}\noLine
\UIC{$\varphi_{79}$}
\AXC{$\varphi_{79} \to a_{N} \in \mathbb{R} \to M \in \mathbb{R} \to \lvert a_{N} \rvert \leq M \to -M \leq a_{N} \land a_{N} \leq M$\lbl{A2}}
\BIC{$a_{N} \in \mathbb{R} \to M \in \mathbb{R} \to \lvert a_{N} \rvert \leq M \to -M \leq a_{N} \land a_{N} \leq M$}
\AXC{$\varphi_{442}$\lbl{A1}}
\BIC{$\varphi_{433}$}
\end{ptree}
\noindent$\Pi_{167}$:
\begin{ptree}
\AXC{$\mathrm{SeqBdd}(a) \to \varphi_{209}$\lbl{A1}}
\AXC{$\stackrel{\Pi_{16}}{\vdots}$}\noLine
\UIC{$\mathrm{SeqBdd}(a) \to \mathrm{Seq}(a)$}
\AXC{$\stackrel{\Pi_{164}}{\vdots}$}\noLine
\UIC{$\varphi_{383}$}
\BIC{$(\mathrm{SeqBdd}(a) \to \varphi_{209}) \to \mathrm{SeqBdd}(a) \to \varphi_{223}$}
\BIC{$\mathrm{SeqBdd}(a) \to \varphi_{223}$}
\AXC{$\stackrel{\Pi_{166}}{\vdots}$}\noLine
\UIC{$\varphi_{433}$}
\BIC{$\varphi_{40} \to (\mathrm{SeqBdd}(a) \to \varphi_{220}) \to \mathrm{SeqBdd}(a) \to \varphi_{233}$}
\end{ptree}
\noindent$\Pi_{168}$:
\begin{ptree}
\AXC{$\mathrm{SeqBdd}(a) \to \varphi_{105}$\lbl{A1}}
\AXC{$(\mathrm{SeqBdd}(a) \to \varphi_{105}) \to \varphi_{40}$\lbl{A1}}
\BIC{$\varphi_{40}$}
\AXC{$\stackrel{\Pi_{167}}{\vdots}$}\noLine
\UIC{$\varphi_{40} \to (\mathrm{SeqBdd}(a) \to \varphi_{220}) \to \mathrm{SeqBdd}(a) \to \varphi_{233}$}
\BIC{$(\mathrm{SeqBdd}(a) \to \varphi_{220}) \to \mathrm{SeqBdd}(a) \to \varphi_{233}$}
\end{ptree}
\noindent$\Pi_{169}$:
\begin{ptree}
\AXC{$x = x$\lbl{A5}}
\UIC{$\forall x\,(x = x)$}
\AXC{$\forall x\,(x = x) \to N = N$\lbl{A2}}
\BIC{$N = N$}
\AXC{$\mathrm{SeqBdd}(a) \to \varphi_{209}$\lbl{A1}}
\AXC{$\varphi_{422}$\lbl{A1}}
\BIC{$N = N \to (\mathrm{SeqBdd}(a) \to \varphi_{233}) \to \mathrm{SeqBdd}(a) \to \varphi_{240}$}
\BIC{$(\mathrm{SeqBdd}(a) \to \varphi_{233}) \to \mathrm{SeqBdd}(a) \to \varphi_{240}$}
\end{ptree}
\noindent$\Pi_{170}$:
\begin{ptree}
\AXC{$\stackrel{\Pi_{163}}{\vdots}$}\noLine
\UIC{$\mathrm{SeqBdd}(a) \to \varphi_{220}$}
\AXC{$\stackrel{\Pi_{168}}{\vdots}$}\noLine
\UIC{$(\mathrm{SeqBdd}(a) \to \varphi_{220}) \to \mathrm{SeqBdd}(a) \to \varphi_{233}$}
\BIC{$\mathrm{SeqBdd}(a) \to \varphi_{233}$}
\AXC{$\stackrel{\Pi_{169}}{\vdots}$}\noLine
\UIC{$(\mathrm{SeqBdd}(a) \to \varphi_{233}) \to \mathrm{SeqBdd}(a) \to \varphi_{240}$}
\BIC{$\mathrm{SeqBdd}(a) \to \varphi_{240}$}
\end{ptree}
\noindent$\Pi_{171}$:
\begin{ptree}
\AXC{$\varphi_{54}$\lbl{A2}}
\AXC{$\stackrel{\Pi_{170}}{\vdots}$}\noLine
\UIC{$\mathrm{SeqBdd}(a) \to \varphi_{240}$}
\AXC{$(\mathrm{SeqBdd}(a) \to \varphi_{240}) \to \varphi_{54} \to \mathrm{SeqBdd}(a) \to \varphi_{246}$\lbl{A1}}
\BIC{$\varphi_{54} \to \mathrm{SeqBdd}(a) \to \varphi_{246}$}
\BIC{$\mathrm{SeqBdd}(a) \to \varphi_{246}$}
\AXC{$\varphi_{393}$\lbl{A1}}
\BIC{$\varphi_{182} \to N \in \mathbb{N} \to \exists n\,(n \in \mathbb{N} \land N \leq n \land -M \leq a_{n})$}
\UIC{$\forall N\,(\varphi_{182} \to N \in \mathbb{N} \to \exists n\,(n \in \mathbb{N} \land N \leq n \land -M \leq a_{n}))$}
\AXC{$\varphi_{476}$\lbl{A1}}
\BIC{$\varphi_{399} \to (\varphi_{182} \to \forall N\,(\varphi_{182})) \to \mathrm{SeqBdd}(a) \to \varphi_{251}$}
\end{ptree}
\noindent$\Pi_{172}$:
\begin{ptree}
\AXC{$\forall x \in \mathbb{R}\,(-x \in \mathbb{R})$\lbl{(R1)}}
\AXC{$\forall x \in \mathbb{R}\,(-x \in \mathbb{R}) \to M \in \mathbb{R} \to -M \in \mathbb{R}$\lbl{A2}}
\BIC{$M \in \mathbb{R} \to -M \in \mathbb{R}$}
\AXC{$(M \in \mathbb{R} \to -M \in \mathbb{R}) \to \varphi_{40} \to \varphi_{219}$\lbl{A1}}
\BIC{$\varphi_{40} \to \varphi_{219}$}
\end{ptree}
\noindent$\Pi_{173}$:
\begin{ptree}
\AXC{$\forall x\,(\varphi_{42}) \to \varphi_{89}$\lbl{A2}}
\AXC{$\varphi_{278}$\lbl{A1}}
\AXC{$\varphi_{278} \to (\forall x\,(\varphi_{42}) \to \varphi_{89}) \to \varphi_{299}$\lbl{A1}}
\BIC{$(\forall x\,(\varphi_{42}) \to \varphi_{89}) \to \varphi_{299}$}
\BIC{$\varphi_{299}$}
\AXC{$\varphi_{299} \to \varphi_{219} \to (\mathrm{SeqBdd}(a) \to \varphi_{251}) \to \varphi_{217}$\lbl{A1}}
\BIC{$\varphi_{219} \to (\mathrm{SeqBdd}(a) \to \varphi_{251}) \to \varphi_{217}$}
\end{ptree}
\noindent$\Pi_{174}$:
\begin{ptree}
\AXC{$\mathrm{SeqBdd}(a) \to \varphi_{105}$\lbl{A1}}
\AXC{$(\mathrm{SeqBdd}(a) \to \varphi_{105}) \to \varphi_{40}$\lbl{A1}}
\BIC{$\varphi_{40}$}
\AXC{$\stackrel{\Pi_{172}}{\vdots}$}\noLine
\UIC{$\varphi_{40} \to \varphi_{219}$}
\BIC{$\varphi_{219}$}
\AXC{$\stackrel{\Pi_{173}}{\vdots}$}\noLine
\UIC{$\varphi_{219} \to (\mathrm{SeqBdd}(a) \to \varphi_{251}) \to \varphi_{217}$}
\BIC{$(\mathrm{SeqBdd}(a) \to \varphi_{251}) \to \varphi_{217}$}
\end{ptree}
\noindent$\Pi_{175}$:
\begin{ptree}
\AXC{$\varphi_{182} \to \forall N\,(\varphi_{182})$\lbl{A4}}
\AXC{$\varphi_{399}$\lbl{A3}}
\AXC{$\stackrel{\Pi_{171}}{\vdots}$}\noLine
\UIC{$\varphi_{399} \to (\varphi_{182} \to \forall N\,(\varphi_{182})) \to \mathrm{SeqBdd}(a) \to \varphi_{251}$}
\BIC{$(\varphi_{182} \to \forall N\,(\varphi_{182})) \to \mathrm{SeqBdd}(a) \to \varphi_{251}$}
\BIC{$\mathrm{SeqBdd}(a) \to \varphi_{251}$}
\AXC{$\stackrel{\Pi_{174}}{\vdots}$}\noLine
\UIC{$(\mathrm{SeqBdd}(a) \to \varphi_{251}) \to \varphi_{217}$}
\BIC{$\varphi_{217}$}
\end{ptree}
\noindent$\Pi_{176}$:
\begin{ptree}
\AXC{$\forall X\,\forall x \in X\,(X \neq \emptyset)$\lbl{[B2]}}
\AXC{$\forall X\,\forall x \in X\,(X \neq \emptyset) \to \forall x \in E\,(E \neq \emptyset)$\lbl{A2}}
\BIC{$\forall x \in E\,(E \neq \emptyset)$}
\AXC{$\forall x \in E\,(E \neq \emptyset) \to -M \in E \to E \neq \emptyset$\lbl{A2}}
\BIC{$-M \in E \to E \neq \emptyset$}
\end{ptree}
\noindent$\Pi_{177}$:
\begin{ptree}
\AXC{$\stackrel{\Pi_{175}}{\vdots}$}\noLine
\UIC{$\varphi_{217}$}
\AXC{$\stackrel{\Pi_{176}}{\vdots}$}\noLine
\UIC{$-M \in E \to E \neq \emptyset$}
\AXC{$(-M \in E \to E \neq \emptyset) \to \varphi_{217} \to \varphi_{211}$\lbl{A1}}
\BIC{$\varphi_{217} \to \varphi_{211}$}
\BIC{$\varphi_{211}$}
\end{ptree}
\noindent$\Pi_{178}$:
\begin{ptree}
\AXC{$\stackrel{\Pi_{177}}{\vdots}$}\noLine
\UIC{$\varphi_{211}$}
\AXC{$\stackrel{\Pi_{116}}{\vdots}$}\noLine
\UIC{$\varphi_{210}$}
\AXC{$\varphi_{16}$\lbl{(R16)}}
\AXC{$\varphi_{16} \to \varphi_{14}$\lbl{A2}}
\BIC{$\varphi_{14}$}
\AXC{$\varphi_{438}$\lbl{A1}}
\BIC{$\varphi_{210} \to \varphi_{211} \to (\mathrm{SeqBdd}(a) \to \varphi_{186}) \to \mathrm{SeqBdd}(a) \to \varphi_{194}$}
\BIC{$\varphi_{211} \to (\mathrm{SeqBdd}(a) \to \varphi_{186}) \to \mathrm{SeqBdd}(a) \to \varphi_{194}$}
\BIC{$(\mathrm{SeqBdd}(a) \to \varphi_{186}) \to \mathrm{SeqBdd}(a) \to \varphi_{194}$}
\end{ptree}
\noindent$\Pi_{179}$:
\begin{ptree}
\AXC{$\forall b\,(\neg \mathrm{UB}(b,E)) \to \neg \mathrm{UB}(M,E)$\lbl{A2}}
\AXC{$\stackrel{\Pi_{161}}{\vdots}$}\noLine
\UIC{$\varphi_{243}$}
\BIC{$\mathrm{SeqBdd}(a) \to \varphi_{186}$}
\AXC{$\stackrel{\Pi_{178}}{\vdots}$}\noLine
\UIC{$(\mathrm{SeqBdd}(a) \to \varphi_{186}) \to \mathrm{SeqBdd}(a) \to \varphi_{194}$}
\BIC{$\mathrm{SeqBdd}(a) \to \varphi_{194}$}
\end{ptree}
\noindent$\Pi_{180}$:
\begin{ptree}
\AXC{$\varphi_{368}$\lbl{A4}}
\AXC{$\forall c\,(\varphi_{249}) \to \varphi_{253}$\lbl{A3}}
\AXC{$\stackrel{\Pi_{142}}{\vdots}$}\noLine
\UIC{$\forall c\,(\varphi_{249})$}
\AXC{$\stackrel{\Pi_{179}}{\vdots}$}\noLine
\UIC{$\mathrm{SeqBdd}(a) \to \varphi_{194}$}
\AXC{$(\mathrm{SeqBdd}(a) \to \varphi_{194}) \to \forall c\,(\varphi_{249}) \to \varphi_{472}$\lbl{A1}}
\BIC{$\forall c\,(\varphi_{249}) \to \varphi_{472}$}
\BIC{$\varphi_{472}$}
\BIC{$\varphi_{368} \to \varphi_{1} \to (\varphi_{1} \to \varphi_{3}) \to \varphi_{2} \to \mathrm{SeqBdd}(a) \to \varphi_{216}$}
\BIC{$\varphi_{1} \to (\varphi_{1} \to \varphi_{3}) \to \varphi_{2} \to \mathrm{SeqBdd}(a) \to \varphi_{216}$}
\end{ptree}
\noindent$\Pi_{181}$:
\begin{ptree}
\AXC{$\varphi_{2}$\lbl{A4}}
\AXC{$\varphi_{1} \to \varphi_{3}$\lbl{A3}}
\AXC{$\forall c\,(\neg \mathrm{IsSup}(c,E)) \to \neg \mathrm{IsSup}(s,E)$\lbl{A2}}
\UIC{$\varphi_{1}$}
\AXC{$\stackrel{\Pi_{180}}{\vdots}$}\noLine
\UIC{$\varphi_{1} \to (\varphi_{1} \to \varphi_{3}) \to \varphi_{2} \to \mathrm{SeqBdd}(a) \to \varphi_{216}$}
\BIC{$(\varphi_{1} \to \varphi_{3}) \to \varphi_{2} \to \mathrm{SeqBdd}(a) \to \varphi_{216}$}
\BIC{$\varphi_{2} \to \mathrm{SeqBdd}(a) \to \varphi_{216}$}
\BIC{$\mathrm{SeqBdd}(a) \to \varphi_{216}$}
\end{ptree}
\noindent$\Pi_{182}$:
\begin{ptree}
\AXC{$\forall a\,(\varphi_{57})$\lbl{Sep}}
\AXC{$\forall a\,(\varphi_{57}) \to \varphi_{57}$\lbl{A2}}
\BIC{$\varphi_{57}$}
\AXC{$\varphi_{57} \to \exists E\,\forall x\,(\varphi_{42})$\lbl{A2}}
\BIC{$\exists E\,\forall x\,(\varphi_{42})$}
\AXC{$\exists E\,\forall x\,(\varphi_{42}) \to \varphi_{189}$\lbl{A1}}
\BIC{$\varphi_{189}$}
\end{ptree}
\noindent$\Pi_{183}$:
\begin{ptree}
\AXC{$\stackrel{\Pi_{181}}{\vdots}$}\noLine
\UIC{$\mathrm{SeqBdd}(a) \to \varphi_{216}$}
\AXC{$(\mathrm{SeqBdd}(a) \to \varphi_{216}) \to \varphi_{80} \to \neg \forall x\,(\varphi_{42})$\lbl{A1}}
\BIC{$\varphi_{80} \to \neg \forall x\,(\varphi_{42})$}
\UIC{$\forall E\,(\varphi_{80} \to \neg \forall x\,(\varphi_{42}))$}
\AXC{$\stackrel{\Pi_{182}}{\vdots}$}\noLine
\UIC{$\varphi_{189}$}
\AXC{$\varphi_{189} \to \forall E\,(\varphi_{80} \to \neg \forall x\,(\varphi_{42})) \to \varphi_{434}$\lbl{A1}}
\BIC{$\forall E\,(\varphi_{80} \to \neg \forall x\,(\varphi_{42})) \to \varphi_{434}$}
\BIC{$\varphi_{434}$}
\end{ptree}
\noindent$\Pi_{184}$:
\begin{ptree}
\AXC{$\varphi_{80} \to \forall E\,(\varphi_{80})$\lbl{A4}}
\AXC{$\forall E\,(\varphi_{80} \to \neg \forall x\,(\varphi_{42})) \to \forall E\,(\varphi_{80}) \to \forall E\,(\neg \forall x\,(\varphi_{42}))$\lbl{A3}}
\AXC{$\stackrel{\Pi_{183}}{\vdots}$}\noLine
\UIC{$\varphi_{434}$}
\BIC{$(\varphi_{80} \to \forall E\,(\varphi_{80})) \to \mathrm{SeqBdd}(a) \to \varphi_{44}$}
\BIC{$\mathrm{SeqBdd}(a) \to \varphi_{44}$}
\AXC{$\varphi_{241}$\lbl{A1}}
\BIC{$\varphi_{5} \to \neg (M \in \mathbb{R} \land \forall n \in \mathbb{N}\,(\lvert a_{n} \rvert \leq M))$}
\UIC{$\forall M\,(\varphi_{5} \to \neg (M \in \mathbb{R} \land \forall n \in \mathbb{N}\,(\lvert a_{n} \rvert \leq M)))$}
\end{ptree}
\noindent$\Pi_{185}$:
\begin{ptree}
\AXC{$\stackrel{\Pi_{15}}{\vdots}$}\noLine
\UIC{$\varphi_{36}$}
\AXC{$\varphi_{36} \to \mathrm{SeqBdd}(a) \to \exists M \in \mathbb{R}\,\forall n \in \mathbb{N}\,(\lvert a_{n} \rvert \leq M)$\lbl{A1}}
\BIC{$\mathrm{SeqBdd}(a) \to \exists M \in \mathbb{R}\,\forall n \in \mathbb{N}\,(\lvert a_{n} \rvert \leq M)$}
\AXC{$\varphi_{413}$\lbl{A1}}
\BIC{$\forall M\,(\varphi_{5} \to \neg (M \in \mathbb{R} \land \forall n \in \mathbb{N}\,(\lvert a_{n} \rvert \leq M))) \to \varphi_{345}$}
\end{ptree}
\noindent$\Pi_{186}$:
\begin{ptree}
\AXC{$\stackrel{\Pi_{184}}{\vdots}$}\noLine
\UIC{$\forall M\,(\varphi_{5} \to \neg (M \in \mathbb{R} \land \forall n \in \mathbb{N}\,(\lvert a_{n} \rvert \leq M)))$}
\AXC{$\stackrel{\Pi_{185}}{\vdots}$}\noLine
\UIC{$\forall M\,(\varphi_{5} \to \neg (M \in \mathbb{R} \land \forall n \in \mathbb{N}\,(\lvert a_{n} \rvert \leq M))) \to \varphi_{345}$}
\BIC{$\varphi_{345}$}
\end{ptree}
\noindent 主木:
\begin{ptree}
\AXC{$\varphi_{5} \to \forall M\,(\varphi_{5})$\lbl{A4}}
\AXC{$\forall M\,(\varphi_{5} \to \neg (M \in \mathbb{R} \land \forall n \in \mathbb{N}\,(\lvert a_{n} \rvert \leq M))) \to \varphi_{103}$\lbl{A3}}
\AXC{$\stackrel{\Pi_{186}}{\vdots}$}\noLine
\UIC{$\varphi_{345}$}
\BIC{$(\varphi_{5} \to \forall M\,(\varphi_{5})) \to \varphi_{4}$}
\BIC{$\varphi_{4}$}
\UIC{$\forall a\,(\varphi_{4})$}
\end{ptree}
\begin{abbrevlist}
\item $\varphi_{1} :\equiv \forall s\,(\forall c\,(\neg \mathrm{IsSup}(c,E)) \to \neg \mathrm{IsSup}(s,E))$
\item $\varphi_{2} :\equiv \forall c\,(\neg \mathrm{IsSup}(c,E)) \to \forall s\,\forall c\,(\neg \mathrm{IsSup}(c,E))$
\item $\varphi_{3} :\equiv \forall s\,\forall c\,(\neg \mathrm{IsSup}(c,E)) \to \forall s\,(\neg \mathrm{IsSup}(s,E))$
\item $\varphi_{4} :\equiv \mathrm{SeqBdd}(a) \to \exists b\,(\mathrm{SubSeq}(b,a) \land \mathrm{Cgt}(b))$
\item $\varphi_{5} :\equiv \mathrm{SeqBdd}(a) \land \neg \exists b\,(\mathrm{SubSeq}(b,a) \land \mathrm{Cgt}(b))$
\item $\varphi_{6} :\equiv \mathrm{IsSup}(c,E) \to \exists b\,(\mathrm{SubSeq}(b,a) \land \mathrm{Cgt}(b))$
\item $\varphi_{7} :\equiv \mathrm{SubSeq}(v_{733},v_{734}) \leftrightarrow \mathrm{SubSeq}(v_{733},v_{734})$
\item $\varphi_{8} :\equiv \mathrm{IsSup}(c,E) \leftrightarrow \mathrm{UB}(c,E) \land \forall w\,(\mathrm{UB}(w,E) \to c \leq w)$
\item $\varphi_{9} :\equiv \forall X\,(\mathrm{IsSup}(c,X) \leftrightarrow \mathrm{UB}(c,X) \land \forall w\,(\mathrm{UB}(w,X) \to c \leq w))$
\item $\varphi_{10} :\equiv \forall s\,\forall X\,(\mathrm{IsSup}(s,X) \leftrightarrow \mathrm{UB}(s,X) \land \forall w\,(\mathrm{UB}(w,X) \to s \leq w))$
\item $\varphi_{11} :\equiv \forall x\,\forall y\,\forall z\,(x \leq y \to y \leq z \to x \leq z) \to \forall y\,\forall z\,(N \leq y \to y \leq z \to N \leq z)$
\item $\varphi_{12} :\equiv \forall x\,\forall y\,\forall z\,(x \leq y \to y \leq z \to x \leq z) \to \forall y\,\forall z\,(x \leq y \to y \leq z \to x \leq z)$
\item $\varphi_{13} :\equiv \forall x\,\forall y\,\forall z\,(x \leq y \to y \leq z \to x \leq z) \to \forall y\,\forall z\,(N_{1} \leq y \to y \leq z \to N_{1} \leq z)$
\item $\varphi_{14} :\equiv E \subseteq \mathbb{R} \to E \neq \emptyset \to \exists b\,(\mathrm{UB}(b,E)) \to \exists s\,(\mathrm{IsSup}(s,E))$
\item $\varphi_{15} :\equiv \forall y\,\forall z\,(x \leq y \to y \leq z \to x \leq z) \to \forall z\,(x \leq a_{n} \to a_{n} \leq z \to x \leq z)$
\item $\varphi_{16} :\equiv \forall X\,(X \subseteq \mathbb{R} \to X \neq \emptyset \to \exists b\,(\mathrm{UB}(b,X)) \to \exists s\,(\mathrm{IsSup}(s,X)))$
\item $\varphi_{17} :\equiv \forall y\,\forall z\,(N \leq y \to y \leq z \to N \leq z) \to \forall z\,(N \leq N \cup N_{1} \to N \cup N_{1} \leq z \to N \leq z)$
\item $\varphi_{18} :\equiv x \in E \to 1_{\omega} \in \mathbb{N} \to \exists n\,(n \in \mathbb{N} \land 1_{\omega} \leq n \land x \leq a_{n})$
\item $\varphi_{19} :\equiv \forall x\,\forall y\,(x = y \to y = x) \to \forall y\,(a \circ s(n) = y \to y = a \circ s(n))$
\item $\varphi_{20} :\equiv \mathrm{SubIdx}(s) \leftrightarrow \mathrm{NSeq}(s) \land \forall k \in \mathbb{N}\,(s_{k} < s_{S(k)})$
\item $\varphi_{21} :\equiv \neg (\mathrm{SubIdx}(s) \land \forall k \in \mathbb{N}\,(a \circ s(k) = a_{s_{k}}))$
\item $\varphi_{22} :\equiv \neg \neg (\mathrm{SubIdx}(s) \to \neg \forall k \in \mathbb{N}\,(v_{733}(k) = v_{734}(s_{k})))$
\item $\varphi_{23} :\equiv \neg \neg (N \in \mathbb{N} \to \exists n\,(n \in \mathbb{N} \land N \leq n \land c + \underline{k}^{-1} \leq a_{n}))$
\item $\varphi_{24} :\equiv \forall y\,\forall z\,(N_{1} \leq y \to y \leq z \to N_{1} \leq z) \to \forall z\,(N_{1} \leq N \cup N_{1} \to N \cup N_{1} \leq z \to N_{1} \leq z)$
\item $\varphi_{25} :\equiv \neg \neg (\mathrm{SubIdx}(t) \to \neg \forall k \in \mathbb{N}\,(a \circ s(k) = a_{t(k)}))$
\item $\varphi_{26} :\equiv \exists t\,(\mathrm{SubIdx}(t) \land \forall k \in \mathbb{N}\,(a \circ s(k) = a_{t(k)}))$
\item $\varphi_{27} :\equiv \neg \neg (\mathrm{SubIdx}(u_{3}) \to \neg \forall k \in \mathbb{N}\,(v_{733}(k) = v_{734}(u_{3}(k))))$
\item $\varphi_{28} :\equiv \neg (N_{1} \in \mathbb{N} \to \exists n\,(n \in \mathbb{N} \land N_{1} \leq n \land c + \underline{k}^{-1} \leq a_{n}))$
\item $\varphi_{29} :\equiv \exists N\,(\neg (N \in \mathbb{N} \to \exists n\,(n \in \mathbb{N} \land N \leq n \land c + \underline{k}^{-1} \leq a_{n})))$
\item $\varphi_{30} :\equiv \forall b\,(\neg (\mathrm{SubSeq}(b,a) \land \mathrm{Cgt}(b))) \to \neg (\mathrm{SubSeq}(a \circ s,a) \land \mathrm{Cgt}(a \circ s))$
\item $\varphi_{31} :\equiv \forall z\,(x \leq a_{n} \to a_{n} \leq z \to x \leq z) \to x \leq a_{n} \to a_{n} \leq M \to x \leq M$
\item $\varphi_{32} :\equiv \neg \neg (\mathrm{SubIdx}(u_{3}) \to \neg \forall k \in \mathbb{N}\,(a \circ s(k) = a_{u_{3}(k)}))$
\item $\varphi_{33} :\equiv n \in \mathbb{N} \land 1_{\omega} \leq n \land x \leq a_{n} \to \forall n \in \mathbb{N}\,(\lvert a_{n} \rvert \leq M)$
\item $\varphi_{34} :\equiv n \in \mathbb{N} \land 1_{\omega} \leq n \land x \leq a_{n} \to n \in \mathbb{N} \land 1_{\omega} \leq n \land x \leq a_{n}$
\item $\varphi_{35} :\equiv \mathrm{SeqBdd}(a) \leftrightarrow \mathrm{Seq}(a) \land \exists M \in \mathbb{R}\,\forall n \in \mathbb{N}\,(\lvert a_{n} \rvert \leq M)$
\item $\varphi_{36} :\equiv \mathrm{SeqBdd}(a) \to \mathrm{Seq}(a) \land \exists M \in \mathbb{R}\,\forall n \in \mathbb{N}\,(\lvert a_{n} \rvert \leq M)$
\item $\varphi_{37} :\equiv \neg (K \in \mathbb{N} \land \forall n \in \mathbb{N}\,(K \leq n \to \lvert a \circ s(n) - c \rvert < \varepsilon))$
\item $\varphi_{38} :\equiv \forall a\,(\mathrm{Cgt}(a) \leftrightarrow \exists \alpha\,(\mathrm{Conv}(a,\alpha))) \to (\mathrm{Cgt}(a \circ s) \leftrightarrow \exists \alpha\,(\mathrm{Conv}(a \circ s,\alpha)))$
\item $\varphi_{39} :\equiv \forall z\,(N \leq N \cup N_{1} \to N \cup N_{1} \leq z \to N \leq z) \to N \leq N \cup N_{1} \to N \cup N_{1} \leq n \to N \leq n$
\item $\varphi_{40} :\equiv \mathrm{SeqBdd}(a) \to M \in \mathbb{R} \land \forall n \in \mathbb{N}\,(\lvert a_{n} \rvert \leq M) \to M \in \mathbb{R}$
\item $\varphi_{41} :\equiv \neg (\mathrm{SubIdx}(s) \land \forall k \in \mathbb{N}\,(\lvert a_{s_{k}} - c \rvert < \underline{k}^{-1}))$
\item $\varphi_{42} :\equiv x \in E \leftrightarrow x \in \mathbb{R} \land \forall N \in \mathbb{N}\,\exists n\,(n \in \mathbb{N} \land N \leq n \land x \leq a_{n})$
\item $\varphi_{43} :\equiv u \in E \leftrightarrow u \in \mathbb{R} \land \forall N \in \mathbb{N}\,\exists n\,(n \in \mathbb{N} \land N \leq n \land u \leq a_{n})$
\item $\varphi_{44} :\equiv M \in \mathbb{R} \land \forall n \in \mathbb{N}\,(\lvert a_{n} \rvert \leq M) \to \exists b\,(\mathrm{SubSeq}(b,a) \land \mathrm{Cgt}(b))$
\item $\varphi_{45} :\equiv \forall x\,\forall y\,(x \leq y \to \neg y < x) \to \forall y\,(c + \underline{k}^{-1} \leq y \to \neg y < c + \underline{k}^{-1})$
\item $\varphi_{46} :\equiv \forall N\,(\neg (N \in \mathbb{N} \land \forall n \in \mathbb{N}\,(N \leq n \to \lvert a \circ s(n) - c \rvert < \varepsilon)))$
\item $\varphi_{47} :\equiv \exists s\,(\mathrm{SubIdx}(s) \land \forall k \in \mathbb{N}\,(\lvert a_{s_{k}} - c \rvert < \underline{k}^{-1}))$
\item $\varphi_{48} :\equiv x \in E \to n \in \mathbb{N} \land 1_{\omega} \leq n \land x \leq a_{n} \to n \in \mathbb{N} \to \lvert a_{n} \rvert \leq M$
\item $\varphi_{49} :\equiv \forall z\,(N_{1} \leq N \cup N_{1} \to N \cup N_{1} \leq z \to N_{1} \leq z) \to N_{1} \leq N \cup N_{1} \to N \cup N_{1} \leq n \to N_{1} \leq n$
\item $\varphi_{50} :\equiv \mathrm{IsSup}(c,E) \to k \in \mathbb{N} \to \neg c + \underline{k}^{-1} \notin E \to c + \underline{k}^{-1} \leq c$
\item $\varphi_{51} :\equiv x \in E \to n \in \mathbb{N} \land 1_{\omega} \leq n \land x \leq a_{n} \to -M \leq a_{n} \land a_{n} \leq M$
\item $\varphi_{52} :\equiv \forall x\,\forall y\,(x \subseteq y \leftrightarrow \forall u \in x\,(u \in y)) \to \forall y\,(E \subseteq y \leftrightarrow \forall u \in E\,(u \in y))$
\item $\varphi_{53} :\equiv \forall y\,(E \subseteq y \leftrightarrow \forall u \in E\,(u \in y)) \to (E \subseteq \mathbb{R} \leftrightarrow \forall u \in E\,(u \in \mathbb{R}))$
\item $\varphi_{54} :\equiv \forall n\,(\neg (n \in \mathbb{N} \land N \leq n \land -M \leq a_{n})) \to \neg (N \in \mathbb{N} \land N \leq N \land -M \leq a_{N})$
\item $\varphi_{55} :\equiv \mathrm{NSeq}(s) \to \mathrm{Seq}(a \circ s) \land \forall k \in \mathbb{N}\,(a \circ s(k) = a_{s_{k}})$
\item $\varphi_{56} :\equiv \forall y\,\forall z\,(\lvert a_{s_{n}} - c \rvert < y \to y \leq z \to \lvert a_{s_{n}} - c \rvert < z)$
\item $\varphi_{57} :\equiv \forall A\,\exists E\,\forall x\,(x \in E \leftrightarrow x \in A \land \forall N \in \mathbb{N}\,\exists n\,(n \in \mathbb{N} \land N \leq n \land x \leq a_{n}))$
\item $\varphi_{58} :\equiv n \in \mathbb{N} \land N \cup N_{1} \leq n \land x \leq a_{n} \to n \in \mathbb{N} \land N \cup N_{1} \leq n \land x \leq a_{n}$
\item $\varphi_{59} :\equiv \mathrm{IsSup}(c,E) \to k \in \mathbb{N} \to \neg c + \underline{k}^{-1} \notin E \to c < c + \underline{k}^{-1}$
\item $\varphi_{60} :\equiv k \in \mathbb{N} \to \neg \forall N \in \mathbb{N}\,\exists n\,(n \in \mathbb{N} \land N \leq n \land c + \underline{k}^{-1} \leq a_{n})$
\item $\varphi_{61} :\equiv \forall e\,(E \subseteq \mathbb{R} \to \mathrm{IsSup}(c,E) \to e \in \mathbb{R} \to 0 < e \to \exists x \in E\,(c - e < x))$
\item $\varphi_{62} :\equiv \mathrm{IsSup}(c,E) \to k \in \mathbb{N} \to \neg c + \underline{k}^{-1} \notin E \to c + \underline{k}^{-1} \in E$
\item $\varphi_{63} :\equiv \mathrm{IsSup}(c,E) \to k \in \mathbb{N} \to \neg c + \underline{k}^{-1} \notin E \to \neg c < c + \underline{k}^{-1}$
\item $\varphi_{64} :\equiv N \in \mathbb{N} \to n \in \mathbb{N} \land N \cup N_{1} \leq n \land x \leq a_{n} \to c + \underline{k}^{-1} \in \mathbb{R}$
\item $\varphi_{65} :\equiv \forall y\,\forall z\,(\lvert a_{s_{n}} - c \rvert < y \to y < z \to \lvert a_{s_{n}} - c \rvert < z)$
\item $\varphi_{66} :\equiv \mathrm{SubIdx}(s) \land \forall k \in \mathbb{N}\,(\lvert a_{s_{k}} - c \rvert < \underline{k}^{-1}) \to \mathrm{NSeq}(s)$
\item $\varphi_{67} :\equiv N \in \mathbb{N} \to n \in \mathbb{N} \land N \cup N_{1} \leq n \land x \leq a_{n} \to \neg c + \underline{k}^{-1} \leq a_{n}$
\item $\varphi_{68} :\equiv \mathrm{IsSup}(c,E) \to k \in \mathbb{N} \to \neg c + \underline{k}^{-1} \notin E \to \neg c + \underline{k}^{-1} \notin E$
\item $\varphi_{69} :\equiv \mathrm{SubIdx}(s) \land \forall k \in \mathbb{N}\,(\lvert a_{s_{k}} - c \rvert < \underline{k}^{-1}) \to \mathrm{SubIdx}(s)$
\item $\varphi_{70} :\equiv 0 < \varepsilon \to \exists N \in \mathbb{N}\,\forall n \in \mathbb{N}\,(N \leq n \to \lvert a \circ s(n) - c \rvert < \varepsilon)$
\item $\varphi_{71} :\equiv N \in \mathbb{N} \to n \in \mathbb{N} \land N \cup N_{1} \leq n \land x \leq a_{n} \to a_{n} < c + \underline{k}^{-1}$
\item $\varphi_{72} :\equiv N \in \mathbb{N} \to n \in \mathbb{N} \land N \cup N_{1} \leq n \land x \leq a_{n} \to c - \underline{k}^{-1} < a_{n}$
\item $\varphi_{73} :\equiv \varepsilon \in \mathbb{R} \to 0 < \varepsilon \to K \in \mathbb{N} \land \underline{K}^{-1} < \varepsilon \to n \in \mathbb{N} \to K \leq n \to K \leq n$
\item $\varphi_{74} :\equiv x \in E \land c - \underline{k}^{-1} < x \to \forall N \in \mathbb{N}\,\exists n\,(n \in \mathbb{N} \land N \leq n \land x \leq a_{n})$
\item $\varphi_{75} :\equiv \forall n\,(\mathrm{Seq}(a) \to n \in \mathbb{N} \to a_{n} \in \mathbb{R}) \to \mathrm{Seq}(a) \to n \in \mathbb{N} \to a_{n} \in \mathbb{R}$
\item $\varphi_{76} :\equiv \forall s\,\forall E\,\forall e\,(E \subseteq \mathbb{R} \to \mathrm{IsSup}(s,E) \to e \in \mathbb{R} \to 0 < e \to \exists x \in E\,(s - e < x))$
\item $\varphi_{77} :\equiv \forall y\,(a_{n} \in \mathbb{R} \to y \in \mathbb{R} \to \lvert a_{n} \rvert \leq y \to -y \leq a_{n} \land a_{n} \leq y)$
\item $\varphi_{78} :\equiv \forall n\,(\mathrm{Seq}(a) \to n \in \mathbb{N} \to a_{n} \in \mathbb{R}) \to \mathrm{Seq}(a) \to N \in \mathbb{N} \to a_{N} \in \mathbb{R}$
\item $\varphi_{79} :\equiv \forall y\,(a_{N} \in \mathbb{R} \to y \in \mathbb{R} \to \lvert a_{N} \rvert \leq y \to -y \leq a_{N} \land a_{N} \leq y)$
\item $\varphi_{80} :\equiv \mathrm{SeqBdd}(a) \land (M \in \mathbb{R} \land \forall n \in \mathbb{N}\,(\lvert a_{n} \rvert \leq M)) \land \neg \exists b\,(\mathrm{SubSeq}(b,a) \land \mathrm{Cgt}(b))$
\item $\varphi_{81} :\equiv 0 < \varepsilon \to K \in \mathbb{N} \land \underline{K}^{-1} < \varepsilon \to n \in \mathbb{N} \to K \leq n \to \underline{K}^{-1} < \varepsilon$
\item $\varphi_{82} :\equiv k \in \mathbb{N} \to \forall N \in \mathbb{N}\,\exists n\,(n \in \mathbb{N} \land N \leq n \land \lvert a_{n} - c \rvert < \underline{k}^{-1})$
\item $\varphi_{83} :\equiv \forall b\,\forall a\,(\mathrm{SubSeq}(b,a) \leftrightarrow \mathrm{Seq}(b) \land \exists s\,(\mathrm{SubIdx}(s) \land \forall k \in \mathbb{N}\,(b_{k} = a_{s_{k}})))$
\item $\varphi_{84} :\equiv \forall X\,(\mathrm{UB}(c,X) \leftrightarrow c \in \mathbb{R} \land \forall x \in X\,(x \leq c)) \to (\mathrm{UB}(c,E) \leftrightarrow c \in \mathbb{R} \land \forall x \in E\,(x \leq c))$
\item $\varphi_{85} :\equiv \forall X\,(\mathrm{UB}(M,X) \leftrightarrow M \in \mathbb{R} \land \forall x \in X\,(x \leq M)) \to (\mathrm{UB}(M,E) \leftrightarrow M \in \mathbb{R} \land \forall x \in E\,(x \leq M))$
\item $\varphi_{86} :\equiv K \in \mathbb{N} \land \underline{K}^{-1} < \varepsilon \to n \in \mathbb{N} \to K \leq n \to \lvert a_{s_{n}} - c \rvert < \varepsilon$
\item $\varphi_{87} :\equiv K \in \mathbb{N} \land \underline{K}^{-1} < \varepsilon \to n \in \mathbb{N} \to K \leq n \to a \circ s(n) = a_{s_{n}}$
\item $\varphi_{88} :\equiv K \in \mathbb{N} \land \underline{K}^{-1} < \varepsilon \to n \in \mathbb{N} \to K \leq n \to a_{s_{n}} = a \circ s(n)$
\item $\varphi_{89} :\equiv -M \in E \leftrightarrow -M \in \mathbb{R} \land \forall N \in \mathbb{N}\,\exists n\,(n \in \mathbb{N} \land N \leq n \land -M \leq a_{n})$
\item $\varphi_{90} :\equiv K \in \mathbb{N} \land \underline{K}^{-1} < \varepsilon \to n \in \mathbb{N} \to K \leq n \to \lvert a \circ s(n) - c \rvert < \varepsilon$
\item $\varphi_{91} :\equiv \mathrm{SubIdx}(s) \land \forall k \in \mathbb{N}\,(\lvert a_{s_{k}} - c \rvert < \underline{k}^{-1}) \to \mathrm{Seq}(a \circ s)$
\item $\varphi_{92} :\equiv \mathrm{SubIdx}(s) \land \forall k \in \mathbb{N}\,(\lvert a_{s_{k}} - c \rvert < \underline{k}^{-1}) \to \mathrm{Cgt}(a \circ s)$
\item $\varphi_{93} :\equiv \forall k \in \mathbb{N}\,\forall N \in \mathbb{N}\,\exists n\,(n \in \mathbb{N} \land N \leq n \land \lvert a_{n} - c \rvert < \underline{k}^{-1})$
\item $\varphi_{94} :\equiv N \in \mathbb{N} \to n \in \mathbb{N} \land N \cup N_{1} \leq n \land x \leq a_{n} \to \lvert a_{n} - c \rvert < \underline{k}^{-1}$
\item $\varphi_{95} :\equiv x \in E \land c - \underline{k}^{-1} < x \to N \in \mathbb{N} \to \exists n\,(n \in \mathbb{N} \land N \cup N_{1} \leq n \land x \leq a_{n})$
\item $\varphi_{96} :\equiv \forall a\,\forall n\,(\mathrm{Seq}(a) \to n \in \mathbb{N} \to a_{n} \in \mathbb{R}) \to \forall n\,(\mathrm{Seq}(a) \to n \in \mathbb{N} \to a_{n} \in \mathbb{R})$
\item $\varphi_{97} :\equiv \forall n\,(N \in \mathbb{N} \to n \in \mathbb{N} \to N \cup n \in \mathbb{N}) \to N \in \mathbb{N} \to N_{1} \in \mathbb{N} \to N \cup N_{1} \in \mathbb{N}$
\item $\varphi_{98} :\equiv 0 < \varepsilon \to K \in \mathbb{N} \land \underline{K}^{-1} < \varepsilon \to n \in \mathbb{N} \to K \leq n \to \underline{n}^{-1} \leq \underline{K}^{-1}$
\item $\varphi_{99} :\equiv K \in \mathbb{N} \land \underline{K}^{-1} < \varepsilon \to \forall n \in \mathbb{N}\,(K \leq n \to \lvert a \circ s(n) - c \rvert < \varepsilon)$
\item $\varphi_{100} :\equiv \mathrm{SubIdx}(s) \land \forall k \in \mathbb{N}\,(\lvert a_{s_{k}} - c \rvert < \underline{k}^{-1}) \to \mathrm{Conv}(a \circ s,c)$
\item $\varphi_{101} :\equiv \forall a\,(\mathrm{SubSeq}(v_{733},a) \leftrightarrow \mathrm{Seq}(v_{733}) \land \exists s\,(\mathrm{SubIdx}(s) \land \forall k \in \mathbb{N}\,(v_{733}(k) = a_{s_{k}})))$
\item $\varphi_{102} :\equiv \forall e\,\forall y\,(c \in \mathbb{R} \to e \in \mathbb{R} \to y \in \mathbb{R} \to c - e < y \to y < c + e \to \lvert y - c \rvert < e)$
\item $\varphi_{103} :\equiv \forall M\,(\varphi_{5}) \to \forall M\,(\neg (M \in \mathbb{R} \land \forall n \in \mathbb{N}\,(\lvert a_{n} \rvert \leq M)))$
\item $\varphi_{104} :\equiv x \in E \land c - \underline{k}^{-1} < x \to N \in \mathbb{N} \to n \in \mathbb{N} \land N \cup N_{1} \leq n \land x \leq a_{n} \to N \leq n$
\item $\varphi_{105} :\equiv M \in \mathbb{R} \land \forall n \in \mathbb{N}\,(\lvert a_{n} \rvert \leq M) \to M \in \mathbb{R} \land \forall n \in \mathbb{N}\,(\lvert a_{n} \rvert \leq M)$
\item $\varphi_{106} :\equiv \forall a\,\forall X\,(\mathrm{UB}(a,X) \leftrightarrow a \in \mathbb{R} \land \forall x \in X\,(x \leq a)) \to \forall X\,(\mathrm{UB}(c,X) \leftrightarrow c \in \mathbb{R} \land \forall x \in X\,(x \leq c))$
\item $\varphi_{107} :\equiv \forall x \in \mathbb{R}\,(0 < x \to 0 < x^{-1}) \to \underline{k} \in \mathbb{R} \to 0 < \underline{k} \to 0 < \underline{k}^{-1}$
\item $\varphi_{108} :\equiv \forall a\,\forall X\,(\mathrm{UB}(a,X) \leftrightarrow a \in \mathbb{R} \land \forall x \in X\,(x \leq a)) \to \forall X\,(\mathrm{UB}(M,X) \leftrightarrow M \in \mathbb{R} \land \forall x \in X\,(x \leq M))$
\item $\varphi_{109} :\equiv \mathrm{SubIdx}(s) \land \forall k \in \mathbb{N}\,(\lvert a_{s_{k}} - c \rvert < \underline{k}^{-1}) \to \mathrm{SubSeq}(a \circ s,a)$
\item $\varphi_{110} :\equiv x \in E \land c - \underline{k}^{-1} < x \to N \in \mathbb{N} \to n \in \mathbb{N} \land N \cup N_{1} \leq n \land x \leq a_{n} \to N_{1} \leq n$
\item $\varphi_{111} :\equiv \varepsilon \in \mathbb{R} \to 0 < \varepsilon \to K \in \mathbb{N} \land \underline{K}^{-1} < \varepsilon \to K \in \mathbb{N} \land \underline{K}^{-1} < \varepsilon$
\item $\varphi_{112} :\equiv n \in \mathbb{N} \land N \cup N_{1} \leq n \land x \leq a_{n} \to n \in \mathbb{N} \land N \leq n \land \lvert a_{n} - c \rvert < \underline{k}^{-1}$
\item $\varphi_{113} :\equiv \mathrm{SubSeq}(v_{733},v_{734}) \leftrightarrow \mathrm{Seq}(v_{733}) \land \exists s\,(\mathrm{SubIdx}(s) \land \forall k \in \mathbb{N}\,(v_{733}(k) = v_{734}(s_{k})))$
\item $\varphi_{114} :\equiv \forall x\,\forall y\,\forall z\,(x < y \to y \leq z \to x < z) \to \forall y\,\forall z\,(c - \underline{k}^{-1} < y \to y \leq z \to c - \underline{k}^{-1} < z)$
\item $\varphi_{115} :\equiv \forall x\,\forall e\,\forall y\,(x \in \mathbb{R} \to e \in \mathbb{R} \to y \in \mathbb{R} \to x - e < y \to y < x + e \to \lvert y - x \rvert < e)$
\item $\varphi_{116} :\equiv \forall m\,\forall n\,(m \in \mathbb{N} \to n \in \mathbb{N} \to m \cup n \in \mathbb{N}) \to \forall n\,(N \in \mathbb{N} \to n \in \mathbb{N} \to N \cup n \in \mathbb{N})$
\item $\varphi_{117} :\equiv \mathrm{IsSup}(c,E) \to \underline{k}^{-1} \in \mathbb{R} \to 0 < \underline{k}^{-1} \to \exists x \in E\,(c - \underline{k}^{-1} < x)$
\item $\varphi_{118} :\equiv K \in \mathbb{N} \land \underline{K}^{-1} < \varepsilon \to n \in \mathbb{N} \to K \leq n \to \lvert a_{s_{n}} - c \rvert < \underline{n}^{-1}$
\item $\varphi_{119} :\equiv K \in \mathbb{N} \land \underline{K}^{-1} < \varepsilon \to n \in \mathbb{N} \to K \leq n \to \lvert a_{s_{n}} - c \rvert < \underline{K}^{-1}$
\item $\varphi_{120} :\equiv \mathrm{SubIdx}(s) \land \forall k \in \mathbb{N}\,(\lvert a_{s_{k}} - c \rvert < \underline{k}^{-1}) \to \exists \alpha\,(\mathrm{Conv}(a \circ s,\alpha))$
\item $\varphi_{121} :\equiv \forall z\,(\lvert a_{s_{n}} - c \rvert < \underline{n}^{-1} \to \underline{n}^{-1} \leq z \to \lvert a_{s_{n}} - c \rvert < z)$
\item $\varphi_{122} :\equiv \forall k \in \mathbb{N}\,(a \circ s(k) = a_{s_{k}}) \to n \in \mathbb{N} \to a \circ s(n) = a_{s_{n}}$
\item $\varphi_{123} :\equiv \forall x \in \mathbb{R}\,(x \neq 0 \to x^{-1} \in \mathbb{R}) \to \underline{k} \in \mathbb{R} \to \underline{k} \neq 0 \to \underline{k}^{-1} \in \mathbb{R}$
\item $\varphi_{124} :\equiv \forall x\,\forall y\,(x \in \mathbb{R} \to y \in \mathbb{R} \to x + y \in \mathbb{R}) \to \forall y\,(c \in \mathbb{R} \to y \in \mathbb{R} \to c + y \in \mathbb{R})$
\item $\varphi_{125} :\equiv \mathrm{SubIdx}(s) \land \forall k \in \mathbb{N}\,(\lvert a_{s_{k}} - c \rvert < \underline{k}^{-1}) \to \exists b\,(\mathrm{SubSeq}(b,a) \land \mathrm{Cgt}(b))$
\item $\varphi_{126} :\equiv n \in \mathbb{N} \land N \cup N_{1} \leq n \land x \leq a_{n} \to \exists n\,(n \in \mathbb{N} \land N \leq n \land \lvert a_{n} - c \rvert < \underline{k}^{-1})$
\item $\varphi_{127} :\equiv x \in E \land c - \underline{k}^{-1} < x \to N \cup N_{1} \in \mathbb{N} \to \exists n\,(n \in \mathbb{N} \land N \cup N_{1} \leq n \land x \leq a_{n})$
\item $\varphi_{128} :\equiv x \in E \land c - \underline{k}^{-1} < x \to N \in \mathbb{N} \to n \in \mathbb{N} \land N \cup N_{1} \leq n \land x \leq a_{n} \to n \in \mathbb{N}$
\item $\varphi_{129} :\equiv \forall y\,(c + \underline{k}^{-1} \leq y \to \neg y < c + \underline{k}^{-1}) \to c + \underline{k}^{-1} \leq c \to \neg c < c + \underline{k}^{-1}$
\item $\varphi_{130} :\equiv x \in E \land c - \underline{k}^{-1} < x \to N \in \mathbb{N} \to n \in \mathbb{N} \land N \cup N_{1} \leq n \land x \leq a_{n} \to x \leq a_{n}$
\item $\varphi_{131} :\equiv \forall z\,(\lvert a_{s_{n}} - c \rvert < \underline{K}^{-1} \to \underline{K}^{-1} < z \to \lvert a_{s_{n}} - c \rvert < z)$
\item $\varphi_{132} :\equiv x \in E \land c - \underline{k}^{-1} < x \to N \in \mathbb{N} \to n \in \mathbb{N} \land N \cup N_{1} \leq n \land x \leq a_{n} \to N \cup N_{1} \leq n$
\item $\varphi_{133} :\equiv \lvert a_{s_{n}} - c \rvert < \underline{K}^{-1} \to \underline{K}^{-1} < \varepsilon \to \lvert a_{s_{n}} - c \rvert < \varepsilon$
\item $\varphi_{134} :\equiv \forall n\,(\neg (n \in \mathbb{N} \land N_{1} \leq n \land c + \underline{k}^{-1} \leq a_{n})) \to \neg (n \in \mathbb{N} \land N_{1} \leq n \land c + \underline{k}^{-1} \leq a_{n})$
\item $\varphi_{135} :\equiv \forall N \in \mathbb{N}\,\exists n\,(n \in \mathbb{N} \land N \leq n \land x \leq a_{n}) \to 1_{\omega} \in \mathbb{N} \to \exists n\,(n \in \mathbb{N} \land 1_{\omega} \leq n \land x \leq a_{n})$
\item $\varphi_{136} :\equiv \mathrm{SubIdx}(s) \land \forall k \in \mathbb{N}\,(\lvert a_{s_{k}} - c \rvert < \underline{k}^{-1}) \to \varepsilon \in \mathbb{R} \to \varepsilon \in \mathbb{R}$
\item $\varphi_{137} :\equiv c - \underline{k}^{-1} < a_{n} \to a_{n} < c + \underline{k}^{-1} \to \lvert a_{n} - c \rvert < \underline{k}^{-1}$
\item $\varphi_{138} :\equiv K \in \mathbb{N} \land \underline{K}^{-1} < \varepsilon \to K \in \mathbb{N} \land \forall n \in \mathbb{N}\,(K \leq n \to \lvert a \circ s(n) - c \rvert < \varepsilon)$
\item $\varphi_{139} :\equiv x \in E \land c - \underline{k}^{-1} < x \to N \in \mathbb{N} \to \exists n\,(n \in \mathbb{N} \land N \leq n \land \lvert a_{n} - c \rvert < \underline{k}^{-1})$
\item $\varphi_{140} :\equiv a_{s_{n}} = a \circ s(n) \to \lvert a_{s_{n}} - c \rvert < \varepsilon \to \lvert a \circ s(n) - c \rvert < \varepsilon$
\item $\varphi_{141} :\equiv K \in \mathbb{N} \land \underline{K}^{-1} < \varepsilon \to \exists N \in \mathbb{N}\,\forall n \in \mathbb{N}\,(N \leq n \to \lvert a \circ s(n) - c \rvert < \varepsilon)$
\item $\varphi_{142} :\equiv x \in E \land c - \underline{k}^{-1} < x \to N \in \mathbb{N} \to n \in \mathbb{N} \land N \cup N_{1} \leq n \land x \leq a_{n} \to a_{n} \in \mathbb{R}$
\item $\varphi_{143} :\equiv \mathrm{SubSeq}(a \circ s,a) \leftrightarrow \mathrm{Seq}(a \circ s) \land \exists u_{3}\,(\mathrm{SubIdx}(u_{3}) \land \forall k \in \mathbb{N}\,(a \circ s(k) = a_{u_{3}(k)}))$
\item $\varphi_{144} :\equiv \forall x\,\forall y\,(x \in \mathbb{R} \to y \in \mathbb{R} \to 0 < y \to x < x + y) \to \forall y\,(c \in \mathbb{R} \to y \in \mathbb{R} \to 0 < y \to c < c + y)$
\item $\varphi_{145} :\equiv \mathrm{SubIdx}(s) \land \forall k \in \mathbb{N}\,(\lvert a_{s_{k}} - c \rvert < \underline{k}^{-1}) \to \varepsilon \in \mathbb{R} \to 0 < \varepsilon \to 0 < \varepsilon$
\item $\varphi_{146} :\equiv x \in E \land c - \underline{k}^{-1} < x \to \forall N \in \mathbb{N}\,\exists n\,(n \in \mathbb{N} \land N \leq n \land \lvert a_{n} - c \rvert < \underline{k}^{-1})$
\item $\varphi_{147} :\equiv \forall v_{733}\,\forall v_{734}\,(\varphi_{285})$
\item $\varphi_{148} :\equiv \mathrm{SubIdx}(s) \land \forall k \in \mathbb{N}\,(\lvert a_{s_{k}} - c \rvert < \underline{k}^{-1}) \to \mathrm{SubSeq}(a \circ s,a) \land \mathrm{Cgt}(a \circ s)$
\item $\varphi_{149} :\equiv a_{n} \in \mathbb{R} \to c + \underline{k}^{-1} \in \mathbb{R} \to \neg c + \underline{k}^{-1} \leq a_{n} \to a_{n} < c + \underline{k}^{-1}$
\item $\varphi_{150} :\equiv \forall N \in \mathbb{N}\,\exists n\,(n \in \mathbb{N} \land N \leq n \land x \leq a_{n}) \to N \cup N_{1} \in \mathbb{N} \to \exists n\,(n \in \mathbb{N} \land N \cup N_{1} \leq n \land x \leq a_{n})$
\item $\varphi_{151} :\equiv x \in E \land c - \underline{k}^{-1} < x \to N \in \mathbb{N} \to n \in \mathbb{N} \land N \cup N_{1} \leq n \land x \leq a_{n} \to c - \underline{k}^{-1} < x$
\item $\varphi_{152} :\equiv \forall y\,(c \in \mathbb{R} \to y \in \mathbb{R} \to c + y \in \mathbb{R}) \to c \in \mathbb{R} \to \underline{k}^{-1} \in \mathbb{R} \to c + \underline{k}^{-1} \in \mathbb{R}$
\item $\varphi_{153} :\equiv \forall x\,\forall y\,(x \in \mathbb{R} \to y \in \mathbb{R} \to \neg y \leq x \to x < y) \to \forall y\,(a_{n} \in \mathbb{R} \to y \in \mathbb{R} \to \neg y \leq a_{n} \to a_{n} < y)$
\item $\varphi_{154} :\equiv \forall k \in \mathbb{N}\,(\lvert a_{s_{k}} - c \rvert < \underline{k}^{-1}) \to n \in \mathbb{N} \to \lvert a_{s_{n}} - c \rvert < \underline{n}^{-1}$
\item $\varphi_{155} :\equiv \forall y\,(a \circ s(n) = y \to y = a \circ s(n)) \to a \circ s(n) = a_{s_{n}} \to a_{s_{n}} = a \circ s(n)$
\item $\varphi_{156} :\equiv \forall n\,(\neg (n \in \mathbb{N} \land N \leq n \land \lvert a_{n} - c \rvert < \underline{k}^{-1})) \to \neg (n \in \mathbb{N} \land N \leq n \land \lvert a_{n} - c \rvert < \underline{k}^{-1})$
\item $\varphi_{157} :\equiv \lvert a_{s_{n}} - c \rvert < \underline{n}^{-1} \to \underline{n}^{-1} \leq \underline{K}^{-1} \to \lvert a_{s_{n}} - c \rvert < \underline{K}^{-1}$
\item $\varphi_{158} :\equiv \mathrm{SubIdx}(s) \land \forall k \in \mathbb{N}\,(\lvert a_{s_{k}} - c \rvert < \underline{k}^{-1}) \to n \in \mathbb{N} \to a \circ s(n) = a_{s_{n}}$
\item $\varphi_{159} :\equiv \forall n\,(K \in \mathbb{N} \to n \in \mathbb{N} \to K \leq n \to \underline{n}^{-1} \leq \underline{K}^{-1}) \to K \in \mathbb{N} \to n \in \mathbb{N} \to K \leq n \to \underline{n}^{-1} \leq \underline{K}^{-1}$
\item $\varphi_{160} :\equiv \forall v_{734}\,(\mathrm{SubSeq}(a \circ s,v_{734}) \leftrightarrow \mathrm{Seq}(a \circ s) \land \exists u_{3}\,(\mathrm{SubIdx}(u_{3}) \land \forall k \in \mathbb{N}\,(a \circ s(k) = v_{734}(u_{3}(k)))))$
\item $\varphi_{161} :\equiv c + \underline{k}^{-1} \in E \leftrightarrow c + \underline{k}^{-1} \in \mathbb{R} \land \forall N \in \mathbb{N}\,\exists n\,(n \in \mathbb{N} \land N \leq n \land c + \underline{k}^{-1} \leq a_{n})$
\item $\varphi_{162} :\equiv \forall y\,\forall z\,(c - \underline{k}^{-1} < y \to y \leq z \to c - \underline{k}^{-1} < z) \to \forall z\,(c - \underline{k}^{-1} < x \to x \leq z \to c - \underline{k}^{-1} < z)$
\item $\varphi_{163} :\equiv \mathrm{SubIdx}(s) \land \forall k \in \mathbb{N}\,(\lvert a_{s_{k}} - c \rvert < \underline{k}^{-1}) \to \forall k \in \mathbb{N}\,(a \circ s(k) = a_{s_{k}})$
\item $\varphi_{164} :\equiv \forall x\,(\varphi_{42}) \to x \in E \to \exists n\,(n \in \mathbb{N} \land 1_{\omega} \leq n \land x \leq a_{n})$
\item $\varphi_{165} :\equiv \forall K\,\forall n\,(K \in \mathbb{N} \to n \in \mathbb{N} \to K \leq n \to \underline{n}^{-1} \leq \underline{K}^{-1}) \to \forall n\,(K \in \mathbb{N} \to n \in \mathbb{N} \to K \leq n \to \underline{n}^{-1} \leq \underline{K}^{-1})$
\item $\varphi_{166} :\equiv \forall x\,(\varphi_{42}) \to x \in E \to n \in \mathbb{N} \land 1_{\omega} \leq n \land x \leq a_{n} \to x \leq M$
\item $\varphi_{167} :\equiv \mathrm{SubIdx}(s) \land \forall k \in \mathbb{N}\,(\lvert a_{s_{k}} - c \rvert < \underline{k}^{-1}) \to n \in \mathbb{N} \to \lvert a_{s_{n}} - c \rvert < \underline{n}^{-1}$
\item $\varphi_{168} :\equiv \mathrm{IsSup}(c,E) \to k \in \mathbb{N} \to \varphi_{28} \to \exists x \in E\,(c - \underline{k}^{-1} < x)$
\item $\varphi_{169} :\equiv \forall z\,(c - \underline{k}^{-1} < x \to x \leq z \to c - \underline{k}^{-1} < z) \to c - \underline{k}^{-1} < x \to x \leq a_{n} \to c - \underline{k}^{-1} < a_{n}$
\item $\varphi_{170} :\equiv \forall N\,(\varphi_{23}) \to N \in \mathbb{N} \to \exists n\,(n \in \mathbb{N} \land N \leq n \land c + \underline{k}^{-1} \leq a_{n})$
\item $\varphi_{171} :\equiv k \in \mathbb{N} \to \varphi_{28} \to \forall n\,(\neg (n \in \mathbb{N} \land N_{1} \leq n \land c + \underline{k}^{-1} \leq a_{n}))$
\item $\varphi_{172} :\equiv \forall x\,(\varphi_{42}) \to \mathrm{IsSup}(c,E) \to k \in \mathbb{N} \to c + \underline{k}^{-1} \in \mathbb{R}$
\item $\varphi_{173} :\equiv \mathrm{SubIdx}(s) \land \forall k \in \mathbb{N}\,(\lvert a_{s_{k}} - c \rvert < \underline{k}^{-1}) \to \forall k \in \mathbb{N}\,(\lvert a_{s_{k}} - c \rvert < \underline{k}^{-1})$
\item $\varphi_{174} :\equiv \forall \varepsilon \in \mathbb{R}\,(0 < \varepsilon \to \exists K \in \mathbb{N}\,(\underline{K}^{-1} < \varepsilon)) \to \varepsilon \in \mathbb{R} \to 0 < \varepsilon \to \exists K \in \mathbb{N}\,(\underline{K}^{-1} < \varepsilon)$
\item $\varphi_{175} :\equiv \forall x\,(\varphi_{42}) \to x \in E \to n \in \mathbb{N} \land 1_{\omega} \leq n \land x \leq a_{n} \to n \in \mathbb{N}$
\item $\varphi_{176} :\equiv M \in \mathbb{R} \land \forall n \in \mathbb{N}\,(\lvert a_{n} \rvert \leq M) \to \forall x\,(\varphi_{42}) \to \mathrm{UB}(M,E)$
\item $\varphi_{177} :\equiv \forall x\,(\varphi_{42}) \to x \in E \to n \in \mathbb{N} \land 1_{\omega} \leq n \land x \leq a_{n} \to a_{n} \leq M$
\item $\varphi_{178} :\equiv \forall x\,(\varphi_{42}) \to x \in E \to n \in \mathbb{N} \land 1_{\omega} \leq n \land x \leq a_{n} \to x \leq a_{n}$
\item $\varphi_{179} :\equiv \forall z\,\forall w\,(z = w \to \lvert z - c \rvert < \varepsilon \to \lvert w - c \rvert < \varepsilon) \to \forall w\,(a_{s_{n}} = w \to \lvert a_{s_{n}} - c \rvert < \varepsilon \to \lvert w - c \rvert < \varepsilon)$
\item $\varphi_{180} :\equiv \forall y\,(c \in \mathbb{R} \to \underline{k}^{-1} \in \mathbb{R} \to y \in \mathbb{R} \to c - \underline{k}^{-1} < y \to y < c + \underline{k}^{-1} \to \lvert y - c \rvert < \underline{k}^{-1})$
\item $\varphi_{181} :\equiv \mathrm{SubIdx}(s) \land \forall k \in \mathbb{N}\,(\lvert a_{s_{k}} - c \rvert < \underline{k}^{-1}) \to \mathrm{SubIdx}(s) \land \forall k \in \mathbb{N}\,(a \circ s(k) = a_{s_{k}})$
\item $\varphi_{182} :\equiv \mathrm{SeqBdd}(a) \land (M \in \mathbb{R} \land \forall n \in \mathbb{N}\,(\lvert a_{n} \rvert \leq M)) \land \forall x\,(\varphi_{42})$
\item $\varphi_{183} :\equiv \forall y\,(c \in \mathbb{R} \to y \in \mathbb{R} \to 0 < y \to c < c + y) \to c \in \mathbb{R} \to \underline{k}^{-1} \in \mathbb{R} \to 0 < \underline{k}^{-1} \to c < c + \underline{k}^{-1}$
\item $\varphi_{184} :\equiv \mathrm{SubIdx}(s) \land \forall k \in \mathbb{N}\,(\lvert a_{s_{k}} - c \rvert < \underline{k}^{-1}) \to \varepsilon \in \mathbb{R} \to 0 < \varepsilon \to \exists K \in \mathbb{N}\,(\underline{K}^{-1} < \varepsilon)$
\item $\varphi_{185} :\equiv \forall x\,(\varphi_{42}) \to x \in E \to \forall N \in \mathbb{N}\,\exists n\,(n \in \mathbb{N} \land N \leq n \land x \leq a_{n})$
\item $\varphi_{186} :\equiv M \in \mathbb{R} \land \forall n \in \mathbb{N}\,(\lvert a_{n} \rvert \leq M) \to \forall x\,(\varphi_{42}) \to \exists b\,(\mathrm{UB}(b,E))$
\item $\varphi_{187} :\equiv \mathrm{IsSup}(c,E) \to k \in \mathbb{N} \to \varphi_{28} \to x \in E \land c - \underline{k}^{-1} < x \to x \in E$
\item $\varphi_{188} :\equiv \forall N\,\forall N\,(\varphi_{23}) \to \forall N \in \mathbb{N}\,\exists n\,(n \in \mathbb{N} \land N \leq n \land c + \underline{k}^{-1} \leq a_{n})$
\item $\varphi_{189} :\equiv \mathrm{SeqBdd}(a) \to M \in \mathbb{R} \land \forall n \in \mathbb{N}\,(\lvert a_{n} \rvert \leq M) \to \exists E\,\forall x\,(\varphi_{42})$
\item $\varphi_{190} :\equiv \mathrm{SubIdx}(s) \land \forall k \in \mathbb{N}\,(\lvert a_{s_{k}} - c \rvert < \underline{k}^{-1}) \to \varphi_{26}$
\item $\varphi_{191} :\equiv \forall x\,(\varphi_{42}) \to \mathrm{IsSup}(c,E) \to c + \underline{k}^{-1} \in E \to c + \underline{k}^{-1} \leq c$
\item $\varphi_{192} :\equiv \forall x\,(\varphi_{42}) \to x \in E \to n \in \mathbb{N} \land 1_{\omega} \leq n \land x \leq a_{n} \to \lvert a_{n} \rvert \leq M$
\item $\varphi_{193} :\equiv M \in \mathbb{R} \land \forall n \in \mathbb{N}\,(\lvert a_{n} \rvert \leq M) \to \forall x\,(\varphi_{42}) \to x \in E \to x \leq M$
\item $\varphi_{194} :\equiv M \in \mathbb{R} \land \forall n \in \mathbb{N}\,(\lvert a_{n} \rvert \leq M) \to \forall x\,(\varphi_{42}) \to \exists s\,(\mathrm{IsSup}(s,E))$
\item $\varphi_{195} :\equiv \mathrm{SubIdx}(s) \land \forall k \in \mathbb{N}\,(\lvert a_{s_{k}} - c \rvert < \underline{k}^{-1}) \to \mathrm{Seq}(a \circ s) \land \forall k \in \mathbb{N}\,(a \circ s(k) = a_{s_{k}})$
\item $\varphi_{196} :\equiv \varphi_{28} \to \forall N \in \mathbb{N}\,\exists n\,(n \in \mathbb{N} \land N \leq n \land \lvert a_{n} - c \rvert < \underline{k}^{-1})$
\item $\varphi_{197} :\equiv \forall a\,\forall \alpha\,(\mathrm{Conv}(a,\alpha) \leftrightarrow \mathrm{Seq}(a) \land \alpha \in \mathbb{R} \land \forall \varepsilon \in \mathbb{R}\,(0 < \varepsilon \to \exists N \in \mathbb{N}\,\forall n \in \mathbb{N}\,(N \leq n \to \lvert a_{n} - \alpha \rvert < \varepsilon)))$
\item $\varphi_{198} :\equiv \mathrm{IsSup}(c,E) \to k \in \mathbb{N} \to \varphi_{28} \to \neg (n \in \mathbb{N} \land N_{1} \leq n \land c + \underline{k}^{-1} \leq a_{n})$
\item $\varphi_{199} :\equiv \forall x\,(\varphi_{42}) \to x \in E \to n \in \mathbb{N} \land 1_{\omega} \leq n \land x \leq a_{n} \to a_{n} \in \mathbb{R}$
\item $\varphi_{200} :\equiv M \in \mathbb{R} \land \forall n \in \mathbb{N}\,(\lvert a_{n} \rvert \leq M) \to \forall x\,(\varphi_{42}) \to \forall x \in E\,(x \leq M)$
\item $\varphi_{201} :\equiv M \in \mathbb{R} \land \forall n \in \mathbb{N}\,(\lvert a_{n} \rvert \leq M) \to \forall x\,(\varphi_{42}) \to u \in E \to u \in E$
\item $\varphi_{202} :\equiv M \in \mathbb{R} \land \forall n \in \mathbb{N}\,(\lvert a_{n} \rvert \leq M) \to \forall x\,(\varphi_{42}) \to x \in E \to x \in E$
\item $\varphi_{203} :\equiv \mathrm{SubIdx}(s) \land \forall k \in \mathbb{N}\,(\lvert a_{s_{k}} - c \rvert < \underline{k}^{-1}) \to \mathrm{SubIdx}(s) \land \forall k \in \mathbb{N}\,(\lvert a_{s_{k}} - c \rvert < \underline{k}^{-1})$
\item $\varphi_{204} :\equiv \mathrm{SubIdx}(s) \land \forall k \in \mathbb{N}\,(\lvert a_{s_{k}} - c \rvert < \underline{k}^{-1}) \to \varepsilon \in \mathbb{R} \to 0 < \varepsilon \to K \in \mathbb{N} \land \underline{K}^{-1} < \varepsilon \to K \in \mathbb{N}$
\item $\varphi_{205} :\equiv k \in \mathbb{N} \to \varphi_{28} \to x \in E \land c - \underline{k}^{-1} < x \to N \in \mathbb{N} \to N \cup N_{1} \in \mathbb{N}$
\item $\varphi_{206} :\equiv M \in \mathbb{R} \land \forall n \in \mathbb{N}\,(\lvert a_{n} \rvert \leq M) \to \forall x\,(\varphi_{42}) \to u \in E \to u \in \mathbb{R}$
\item $\varphi_{207} :\equiv \forall x\,\forall y\,(x \in \mathbb{R} \to y \in \mathbb{R} \to \lvert x \rvert \leq y \to -y \leq x \land x \leq y) \to \varphi_{77}$
\item $\varphi_{208} :\equiv \forall x\,\forall y\,(x \in \mathbb{R} \to y \in \mathbb{R} \to \lvert x \rvert \leq y \to -y \leq x \land x \leq y) \to \varphi_{79}$
\item $\varphi_{209} :\equiv M \in \mathbb{R} \land \forall n \in \mathbb{N}\,(\lvert a_{n} \rvert \leq M) \to \forall x\,(\varphi_{42}) \to N \in \mathbb{N} \to N \in \mathbb{N}$
\item $\varphi_{210} :\equiv \mathrm{SeqBdd}(a) \to M \in \mathbb{R} \land \forall n \in \mathbb{N}\,(\lvert a_{n} \rvert \leq M) \to \forall x\,(\varphi_{42}) \to E \subseteq \mathbb{R}$
\item $\varphi_{211} :\equiv \mathrm{SeqBdd}(a) \to M \in \mathbb{R} \land \forall n \in \mathbb{N}\,(\lvert a_{n} \rvert \leq M) \to \forall x\,(\varphi_{42}) \to E \neq \emptyset$
\item $\varphi_{212} :\equiv M \in \mathbb{R} \land \forall n \in \mathbb{N}\,(\lvert a_{n} \rvert \leq M) \to \forall x\,(\varphi_{42}) \to \mathrm{IsSup}(c,E) \to \mathrm{IsSup}(c,E)$
\item $\varphi_{213} :\equiv M \in \mathbb{R} \land \forall n \in \mathbb{N}\,(\lvert a_{n} \rvert \leq M) \to \forall x\,(\varphi_{42}) \to \mathrm{IsSup}(c,E) \to c \in \mathbb{R}$
\item $\varphi_{214} :\equiv k \in \mathbb{N} \to \varphi_{28} \to x \in E \land c - \underline{k}^{-1} < x \to x \in E \land c - \underline{k}^{-1} < x$
\item $\varphi_{215} :\equiv M \in \mathbb{R} \land \forall n \in \mathbb{N}\,(\lvert a_{n} \rvert \leq M) \to \forall x\,(\varphi_{42}) \to \forall u \in E\,(u \in \mathbb{R})$
\item $\varphi_{216} :\equiv M \in \mathbb{R} \land \forall n \in \mathbb{N}\,(\lvert a_{n} \rvert \leq M) \to \forall x\,(\varphi_{42}) \to \exists b\,(\mathrm{SubSeq}(b,a) \land \mathrm{Cgt}(b))$
\item $\varphi_{217} :\equiv \mathrm{SeqBdd}(a) \to M \in \mathbb{R} \land \forall n \in \mathbb{N}\,(\lvert a_{n} \rvert \leq M) \to \forall x\,(\varphi_{42}) \to -M \in E$
\item $\varphi_{218} :\equiv \mathrm{IsSup}(c,E) \to k \in \mathbb{N} \to \varphi_{28} \to x \in E \land c - \underline{k}^{-1} < x \to N \in \mathbb{N} \to N \in \mathbb{N}$
\item $\varphi_{219} :\equiv \mathrm{SeqBdd}(a) \to M \in \mathbb{R} \land \forall n \in \mathbb{N}\,(\lvert a_{n} \rvert \leq M) \to \forall x\,(\varphi_{42}) \to -M \in \mathbb{R}$
\item $\varphi_{220} :\equiv M \in \mathbb{R} \land \forall n \in \mathbb{N}\,(\lvert a_{n} \rvert \leq M) \to \forall x\,(\varphi_{42}) \to N \in \mathbb{N} \to \lvert a_{N} \rvert \leq M$
\item $\varphi_{221} :\equiv \mathrm{Conv}(a \circ s,c) \leftrightarrow \mathrm{Seq}(a \circ s) \land c \in \mathbb{R} \land \forall \varepsilon \in \mathbb{R}\,(\varphi_{70})$
\item $\varphi_{222} :\equiv M \in \mathbb{R} \land \forall n \in \mathbb{N}\,(\lvert a_{n} \rvert \leq M) \to \forall x\,(\varphi_{42}) \to \mathrm{IsSup}(c,E) \to \forall x \in E\,(x \leq c)$
\item $\varphi_{223} :\equiv M \in \mathbb{R} \land \forall n \in \mathbb{N}\,(\lvert a_{n} \rvert \leq M) \to \forall x\,(\varphi_{42}) \to N \in \mathbb{N} \to a_{N} \in \mathbb{R}$
\item $\varphi_{224} :\equiv \mathrm{SubIdx}(s) \land \forall k \in \mathbb{N}\,(\lvert a_{s_{k}} - c \rvert < \underline{k}^{-1}) \to \varphi_{73}$
\item $\varphi_{225} :\equiv \mathrm{SubIdx}(s) \land \forall k \in \mathbb{N}\,(\lvert a_{s_{k}} - c \rvert < \underline{k}^{-1}) \to \varepsilon \in \mathbb{R} \to 0 < \varepsilon \to K \in \mathbb{N} \land \underline{K}^{-1} < \varepsilon \to n \in \mathbb{N} \to n \in \mathbb{N}$
\item $\varphi_{226} :\equiv M \in \mathbb{R} \land \forall n \in \mathbb{N}\,(\lvert a_{n} \rvert \leq M) \to \forall x\,(\varphi_{42}) \to \mathrm{IsSup}(c,E) \to k \in \mathbb{N} \to k \in \mathbb{N}$
\item $\varphi_{227} :\equiv M \in \mathbb{R} \land \forall n \in \mathbb{N}\,(\lvert a_{n} \rvert \leq M) \to \forall x\,(\varphi_{42}) \to \mathrm{IsSup}(c,E) \to k \in \mathbb{N} \to 0 \neq \underline{k}$
\item $\varphi_{228} :\equiv M \in \mathbb{R} \land \forall n \in \mathbb{N}\,(\lvert a_{n} \rvert \leq M) \to \forall x\,(\varphi_{42}) \to \mathrm{IsSup}(c,E) \to k \in \mathbb{N} \to \underline{k} \neq 0$
\item $\varphi_{229} :\equiv \forall \alpha\,(\mathrm{Conv}(a \circ s,\alpha) \leftrightarrow \mathrm{Seq}(a \circ s) \land \alpha \in \mathbb{R} \land \forall \varepsilon \in \mathbb{R}\,(0 < \varepsilon \to \exists N \in \mathbb{N}\,\forall n \in \mathbb{N}\,(N \leq n \to \lvert a \circ s(n) - \alpha \rvert < \varepsilon)))$
\item $\varphi_{230} :\equiv M \in \mathbb{R} \land \forall n \in \mathbb{N}\,(\lvert a_{n} \rvert \leq M) \to \forall x\,(\varphi_{42}) \to \mathrm{IsSup}(c,E) \to k \in \mathbb{N} \to 0 < \underline{k}$
\item $\varphi_{231} :\equiv \mathrm{Seq}(v_{733}) \to \neg \neg \forall s\,(\varphi_{22}) \leftrightarrow \mathrm{Seq}(v_{733}) \to \neg \neg \forall u_{3}\,(\varphi_{27})$
\item $\varphi_{232} :\equiv M \in \mathbb{R} \land \forall n \in \mathbb{N}\,(\lvert a_{n} \rvert \leq M) \to \forall x\,(\varphi_{42}) \to N \in \mathbb{N} \to N \in \mathbb{N} \to \lvert a_{N} \rvert \leq M$
\item $\varphi_{233} :\equiv M \in \mathbb{R} \land \forall n \in \mathbb{N}\,(\lvert a_{n} \rvert \leq M) \to \forall x\,(\varphi_{42}) \to N \in \mathbb{N} \to -M \leq a_{N} \land a_{N} \leq M$
\item $\varphi_{234} :\equiv \neg (\mathrm{Seq}(v_{733}) \to \neg \neg \forall s\,(\varphi_{22})) \leftrightarrow \neg (\mathrm{Seq}(v_{733}) \to \neg \neg \forall u_{3}\,(\varphi_{27}))$
\item $\varphi_{235} :\equiv \mathrm{SubIdx}(s) \land \forall k \in \mathbb{N}\,(\lvert a_{s_{k}} - c \rvert < \underline{k}^{-1}) \to \varphi_{111}$
\item $\varphi_{236} :\equiv M \in \mathbb{R} \land \forall n \in \mathbb{N}\,(\lvert a_{n} \rvert \leq M) \to \forall x\,(\varphi_{42}) \to \mathrm{IsSup}(c,E) \to k \in \mathbb{N} \to \underline{k} \in \mathbb{R}$
\item $\varphi_{237} :\equiv M \in \mathbb{R} \land \forall n \in \mathbb{N}\,(\lvert a_{n} \rvert \leq M) \to \forall x\,(\varphi_{42}) \to \mathrm{IsSup}(c,E) \to k \in \mathbb{N} \to 0 < \underline{k}^{-1}$
\item $\varphi_{238} :\equiv M \in \mathbb{R} \land \forall n \in \mathbb{N}\,(\lvert a_{n} \rvert \leq M) \to \forall x\,(\varphi_{42}) \to N \in \mathbb{N} \to \forall n \in \mathbb{N}\,(\lvert a_{n} \rvert \leq M)$
\item $\varphi_{239} :\equiv \mathrm{SubIdx}(s) \land \forall k \in \mathbb{N}\,(\lvert a_{s_{k}} - c \rvert < \underline{k}^{-1}) \to \varepsilon \in \mathbb{R} \to \varphi_{70}$
\item $\varphi_{240} :\equiv M \in \mathbb{R} \land \forall n \in \mathbb{N}\,(\lvert a_{n} \rvert \leq M) \to \forall x\,(\varphi_{42}) \to N \in \mathbb{N} \to N \in \mathbb{N} \land N \leq N \land -M \leq a_{N}$
\item $\varphi_{241} :\equiv (\mathrm{SeqBdd}(a) \to \varphi_{44}) \to \varphi_{5} \to \neg (M \in \mathbb{R} \land \forall n \in \mathbb{N}\,(\lvert a_{n} \rvert \leq M))$
\item $\varphi_{242} :\equiv M \in \mathbb{R} \land \forall n \in \mathbb{N}\,(\lvert a_{n} \rvert \leq M) \to \forall x\,(\varphi_{42}) \to \mathrm{IsSup}(c,E) \to k \in \mathbb{N} \to \underline{k}^{-1} \in \mathbb{R}$
\item $\varphi_{243} :\equiv (\forall b\,(\neg \mathrm{UB}(b,E)) \to \neg \mathrm{UB}(M,E)) \to \mathrm{SeqBdd}(a) \to \varphi_{186}$
\item $\varphi_{244} :\equiv \mathrm{SubIdx}(s) \land \forall k \in \mathbb{N}\,(\lvert a_{s_{k}} - c \rvert < \underline{k}^{-1}) \to \varepsilon \in \mathbb{R} \to \varphi_{81}$
\item $\varphi_{245} :\equiv \mathrm{SeqBdd}(a) \to M \in \mathbb{R} \land \forall n \in \mathbb{N}\,(\lvert a_{n} \rvert \leq M) \to \forall x\,(\varphi_{42}) \to \varphi_{6}$
\item $\varphi_{246} :\equiv M \in \mathbb{R} \land \forall n \in \mathbb{N}\,(\lvert a_{n} \rvert \leq M) \to \forall x\,(\varphi_{42}) \to N \in \mathbb{N} \to \exists n\,(n \in \mathbb{N} \land N \leq n \land -M \leq a_{n})$
\item $\varphi_{247} :\equiv \varphi_{182} \land \mathrm{IsSup}(c,E) \land \neg \exists b\,(\mathrm{SubSeq}(b,a) \land \mathrm{Cgt}(b))$
\item $\varphi_{248} :\equiv M \in \mathbb{R} \land \forall n \in \mathbb{N}\,(\lvert a_{n} \rvert \leq M) \to \forall x\,(\varphi_{42}) \to \mathrm{IsSup}(c,E) \to k \in \mathbb{N} \to c + \underline{k}^{-1} \notin E$
\item $\varphi_{249} :\equiv \varphi_{182} \land \neg \exists b\,(\mathrm{SubSeq}(b,a) \land \mathrm{Cgt}(b)) \to \neg \mathrm{IsSup}(c,E)$
\item $\varphi_{250} :\equiv \mathrm{SubIdx}(s) \land \forall k \in \mathbb{N}\,(\lvert a_{s_{k}} - c \rvert < \underline{k}^{-1}) \to \forall \varepsilon \in \mathbb{R}\,(\varphi_{70})$
\item $\varphi_{251} :\equiv M \in \mathbb{R} \land \forall n \in \mathbb{N}\,(\lvert a_{n} \rvert \leq M) \to \forall x\,(\varphi_{42}) \to \forall N \in \mathbb{N}\,\exists n\,(n \in \mathbb{N} \land N \leq n \land -M \leq a_{n})$
\item $\varphi_{252} :\equiv \mathrm{SubIdx}(s) \land \forall k \in \mathbb{N}\,(\lvert a_{s_{k}} - c \rvert < \underline{k}^{-1}) \to \varepsilon \in \mathbb{R} \to \varphi_{98}$
\item $\varphi_{253} :\equiv \forall c\,(\varphi_{182} \land \neg \exists b\,(\mathrm{SubSeq}(b,a) \land \mathrm{Cgt}(b))) \to \forall c\,(\neg \mathrm{IsSup}(c,E))$
\item $\varphi_{254} :\equiv \mathrm{SubIdx}(s) \land \forall k \in \mathbb{N}\,(\lvert a_{s_{k}} - c \rvert < \underline{k}^{-1}) \to \varepsilon \in \mathbb{R} \to 0 < \varepsilon \to \varphi_{86}$
\item $\varphi_{255} :\equiv \mathrm{SubIdx}(s) \land \forall k \in \mathbb{N}\,(\lvert a_{s_{k}} - c \rvert < \underline{k}^{-1}) \to \varepsilon \in \mathbb{R} \to 0 < \varepsilon \to \varphi_{87}$
\item $\varphi_{256} :\equiv \mathrm{SubIdx}(s) \land \forall k \in \mathbb{N}\,(\lvert a_{s_{k}} - c \rvert < \underline{k}^{-1}) \to \varepsilon \in \mathbb{R} \to 0 < \varepsilon \to \varphi_{88}$
\item $\varphi_{257} :\equiv \forall w\,(a_{s_{n}} = w \to \lvert a_{s_{n}} - c \rvert < \varepsilon \to \lvert w - c \rvert < \varepsilon) \to \varphi_{140}$
\item $\varphi_{258} :\equiv \varphi_{182} \land x \in E \land \neg x \leq M \to \neg (n \in \mathbb{N} \land 1_{\omega} \leq n \land x \leq a_{n})$
\item $\varphi_{259} :\equiv \mathrm{SubIdx}(s) \land \forall k \in \mathbb{N}\,(\lvert a_{s_{k}} - c \rvert < \underline{k}^{-1}) \to \varepsilon \in \mathbb{R} \to 0 < \varepsilon \to \varphi_{90}$
\item $\varphi_{260} :\equiv \mathrm{SubIdx}(s) \land \forall k \in \mathbb{N}\,(\lvert a_{s_{k}} - c \rvert < \underline{k}^{-1}) \to \varepsilon \in \mathbb{R} \to 0 < \varepsilon \to \varphi_{99}$
\item $\varphi_{261} :\equiv \forall n\,(\varphi_{182} \land x \in E \land \neg x \leq M) \to \forall n\,(\neg (n \in \mathbb{N} \land 1_{\omega} \leq n \land x \leq a_{n}))$
\item $\varphi_{262} :\equiv \mathrm{SubIdx}(s) \land \forall k \in \mathbb{N}\,(\lvert a_{s_{k}} - c \rvert < \underline{k}^{-1}) \to \varepsilon \in \mathbb{R} \to 0 < \varepsilon \to \varphi_{118}$
\item $\varphi_{263} :\equiv \mathrm{SubIdx}(s) \land \forall k \in \mathbb{N}\,(\lvert a_{s_{k}} - c \rvert < \underline{k}^{-1}) \to \varepsilon \in \mathbb{R} \to 0 < \varepsilon \to \varphi_{119}$
\item $\varphi_{264} :\equiv \mathrm{SeqBdd}(a) \to M \in \mathbb{R} \land \forall n \in \mathbb{N}\,(\lvert a_{n} \rvert \leq M) \to \forall x\,(\varphi_{42}) \to \varphi_{18}$
\item $\varphi_{265} :\equiv 1_{\omega} \in \mathbb{N} \to \mathrm{SeqBdd}(a) \to M \in \mathbb{R} \land \forall n \in \mathbb{N}\,(\lvert a_{n} \rvert \leq M) \to \varphi_{164}$
\item $\varphi_{266} :\equiv \mathrm{SubIdx}(s) \land \forall k \in \mathbb{N}\,(\lvert a_{s_{k}} - c \rvert < \underline{k}^{-1}) \to \varepsilon \in \mathbb{R} \to 0 < \varepsilon \to \varphi_{138}$
\item $\varphi_{267} :\equiv (\mathrm{UB}(c,E) \leftrightarrow c \in \mathbb{R} \land \forall x \in E\,(x \leq c)) \to \mathrm{SeqBdd}(a) \to \varphi_{213}$
\item $\varphi_{268} :\equiv \mathrm{SubIdx}(s) \land \forall k \in \mathbb{N}\,(\lvert a_{s_{k}} - c \rvert < \underline{k}^{-1}) \to \varepsilon \in \mathbb{R} \to 0 < \varepsilon \to \varphi_{141}$
\item $\varphi_{269} :\equiv \neg (\mathrm{Seq}(v_{733}) \to \neg \neg \forall s\,(\varphi_{22})) \to \mathrm{SubSeq}(v_{733},v_{734}) \leftrightarrow \neg (\mathrm{Seq}(v_{733}) \to \neg \neg \forall u_{3}\,(\varphi_{27})) \to \mathrm{SubSeq}(v_{733},v_{734})$
\item $\varphi_{270} :\equiv \mathrm{SubSeq}(v_{733},v_{734}) \to \neg (\mathrm{Seq}(v_{733}) \to \neg \neg \forall s\,(\varphi_{22})) \leftrightarrow \mathrm{SubSeq}(v_{733},v_{734}) \to \neg (\mathrm{Seq}(v_{733}) \to \neg \neg \forall u_{3}\,(\varphi_{27}))$
\item $\varphi_{271} :\equiv \mathrm{SeqBdd}(a) \to M \in \mathbb{R} \land \forall n \in \mathbb{N}\,(\lvert a_{n} \rvert \leq M) \to \forall x\,(\varphi_{42}) \to \varphi_{42}$
\item $\varphi_{272} :\equiv \mathrm{SeqBdd}(a) \to M \in \mathbb{R} \land \forall n \in \mathbb{N}\,(\lvert a_{n} \rvert \leq M) \to \forall x\,(\varphi_{42}) \to \varphi_{43}$
\item $\varphi_{273} :\equiv \mathrm{IsSup}(c,E) \to k \in \mathbb{N} \to \varphi_{28} \to x \in E \land c - \underline{k}^{-1} < x \to \varphi_{64}$
\item $\varphi_{274} :\equiv \mathrm{IsSup}(c,E) \to k \in \mathbb{N} \to \varphi_{28} \to x \in E \land c - \underline{k}^{-1} < x \to \varphi_{67}$
\item $\varphi_{275} :\equiv \neg (\neg (\mathrm{Seq}(v_{733}) \to \neg \neg \forall s\,(\varphi_{22})) \to \mathrm{SubSeq}(v_{733},v_{734})) \leftrightarrow \neg (\neg (\mathrm{Seq}(v_{733}) \to \neg \neg \forall u_{3}\,(\varphi_{27})) \to \mathrm{SubSeq}(v_{733},v_{734}))$
\item $\varphi_{276} :\equiv (\mathrm{UB}(c,E) \leftrightarrow c \in \mathbb{R} \land \forall x \in E\,(x \leq c)) \to \mathrm{SeqBdd}(a) \to \varphi_{222}$
\item $\varphi_{277} :\equiv \mathrm{SeqBdd}(a) \to M \in \mathbb{R} \land \forall n \in \mathbb{N}\,(\lvert a_{n} \rvert \leq M) \to \forall x\,(\varphi_{42}) \to \varphi_{48}$
\item $\varphi_{278} :\equiv \mathrm{SeqBdd}(a) \to M \in \mathbb{R} \land \forall n \in \mathbb{N}\,(\lvert a_{n} \rvert \leq M) \to \forall x\,(\varphi_{42}) \to \forall x\,(\varphi_{42})$
\item $\varphi_{279} :\equiv \mathrm{SeqBdd}(a) \to M \in \mathbb{R} \land \forall n \in \mathbb{N}\,(\lvert a_{n} \rvert \leq M) \to \forall x\,(\varphi_{42}) \to \varphi_{50}$
\item $\varphi_{280} :\equiv \mathrm{SeqBdd}(a) \to M \in \mathbb{R} \land \forall n \in \mathbb{N}\,(\lvert a_{n} \rvert \leq M) \to \forall x\,(\varphi_{42}) \to \varphi_{51}$
\item $\varphi_{281} :\equiv \mathrm{IsSup}(c,E) \to k \in \mathbb{N} \to \varphi_{28} \to x \in E \land c - \underline{k}^{-1} < x \to \varphi_{71}$
\item $\varphi_{282} :\equiv \mathrm{IsSup}(c,E) \to k \in \mathbb{N} \to \varphi_{28} \to x \in E \land c - \underline{k}^{-1} < x \to \varphi_{72}$
\item $\varphi_{283} :\equiv M \in \mathbb{R} \land \forall n \in \mathbb{N}\,(\lvert a_{n} \rvert \leq M) \to \forall x\,(\varphi_{42}) \to \mathrm{IsSup}(c,E) \to \varphi_{47}$
\item $\varphi_{284} :\equiv \mathrm{SeqBdd}(a) \to M \in \mathbb{R} \land \forall n \in \mathbb{N}\,(\lvert a_{n} \rvert \leq M) \to \forall x\,(\varphi_{42}) \to x \in E \to \varphi_{33}$
\item $\varphi_{285} :\equiv \neg ((\mathrm{SubSeq}(v_{733},v_{734}) \to \neg (\mathrm{Seq}(v_{733}) \to \neg \neg \forall u_{3}\,(\varphi_{27}))) \to \neg (\neg (\mathrm{Seq}(v_{733}) \to \neg \neg \forall u_{3}\,(\varphi_{27})) \to \mathrm{SubSeq}(v_{733},v_{734})))$
\item $\varphi_{286} :\equiv \mathrm{SeqBdd}(a) \to M \in \mathbb{R} \land \forall n \in \mathbb{N}\,(\lvert a_{n} \rvert \leq M) \to \forall x\,(\varphi_{42}) \to \varphi_{59}$
\item $\varphi_{287} :\equiv \mathrm{SeqBdd}(a) \to M \in \mathbb{R} \land \forall n \in \mathbb{N}\,(\lvert a_{n} \rvert \leq M) \to \forall x\,(\varphi_{42}) \to x \in E \to \varphi_{34}$
\item $\varphi_{288} :\equiv \mathrm{SeqBdd}(a) \to M \in \mathbb{R} \land \forall n \in \mathbb{N}\,(\lvert a_{n} \rvert \leq M) \to \forall x\,(\varphi_{42}) \to \varphi_{62}$
\item $\varphi_{289} :\equiv \mathrm{SeqBdd}(a) \to M \in \mathbb{R} \land \forall n \in \mathbb{N}\,(\lvert a_{n} \rvert \leq M) \to \forall x\,(\varphi_{42}) \to \varphi_{63}$
\item $\varphi_{290} :\equiv M \in \mathbb{R} \land \forall n \in \mathbb{N}\,(\lvert a_{n} \rvert \leq M) \to \forall x\,(\varphi_{42}) \to \mathrm{IsSup}(c,E) \to \varphi_{60}$
\item $\varphi_{291} :\equiv M \in \mathbb{R} \land \forall n \in \mathbb{N}\,(\lvert a_{n} \rvert \leq M) \to \forall x\,(\varphi_{42}) \to \mathrm{IsSup}(c,E) \to k \in \mathbb{N} \to \varphi_{29}$
\item $\varphi_{292} :\equiv \mathrm{IsSup}(c,E) \to k \in \mathbb{N} \to \varphi_{28} \to x \in E \land c - \underline{k}^{-1} < x \to \varphi_{94}$
\item $\varphi_{293} :\equiv \mathrm{SeqBdd}(a) \to M \in \mathbb{R} \land \forall n \in \mathbb{N}\,(\lvert a_{n} \rvert \leq M) \to \forall x\,(\varphi_{42}) \to \varphi_{68}$
\item $\varphi_{294} :\equiv M \in \mathbb{R} \land \forall n \in \mathbb{N}\,(\lvert a_{n} \rvert \leq M) \to \forall x\,(\varphi_{42}) \to \mathrm{IsSup}(c,E) \to \varphi_{66}$
\item $\varphi_{295} :\equiv \mathrm{IsSup}(c,E) \to k \in \mathbb{N} \to \varphi_{28} \to x \in E \land c - \underline{k}^{-1} < x \to N \in \mathbb{N} \to \varphi_{58}$
\item $\varphi_{296} :\equiv \varphi_{182} \land \mathrm{IsSup}(c,E) \land (\mathrm{SubIdx}(s) \land \forall k \in \mathbb{N}\,(\lvert a_{s_{k}} - c \rvert < \underline{k}^{-1}))$
\item $\varphi_{297} :\equiv M \in \mathbb{R} \land \forall n \in \mathbb{N}\,(\lvert a_{n} \rvert \leq M) \to \forall x\,(\varphi_{42}) \to \mathrm{IsSup}(c,E) \to \varphi_{69}$
\item $\varphi_{298} :\equiv M \in \mathbb{R} \land \forall n \in \mathbb{N}\,(\lvert a_{n} \rvert \leq M) \to \forall x\,(\varphi_{42}) \to \mathrm{IsSup}(c,E) \to \varphi_{82}$
\item $\varphi_{299} :\equiv \mathrm{SeqBdd}(a) \to M \in \mathbb{R} \land \forall n \in \mathbb{N}\,(\lvert a_{n} \rvert \leq M) \to \forall x\,(\varphi_{42}) \to \varphi_{89}$
\item $\varphi_{300} :\equiv M \in \mathbb{R} \land \forall n \in \mathbb{N}\,(\lvert a_{n} \rvert \leq M) \to \forall x\,(\varphi_{42}) \to \mathrm{IsSup}(c,E) \to \varphi_{91}$
\item $\varphi_{301} :\equiv M \in \mathbb{R} \land \forall n \in \mathbb{N}\,(\lvert a_{n} \rvert \leq M) \to \forall x\,(\varphi_{42}) \to \mathrm{IsSup}(c,E) \to \varphi_{92}$
\item $\varphi_{302} :\equiv M \in \mathbb{R} \land \forall n \in \mathbb{N}\,(\lvert a_{n} \rvert \leq M) \to \forall x\,(\varphi_{42}) \to \mathrm{IsSup}(c,E) \to \varphi_{93}$
\item $\varphi_{303} :\equiv M \in \mathbb{R} \land \forall n \in \mathbb{N}\,(\lvert a_{n} \rvert \leq M) \to \forall x\,(\varphi_{42}) \to \mathrm{IsSup}(c,E) \to k \in \mathbb{N} \to \varphi_{28} \to N_{1} \in \mathbb{N}$
\item $\varphi_{304} :\equiv M \in \mathbb{R} \land \forall n \in \mathbb{N}\,(\lvert a_{n} \rvert \leq M) \to \forall x\,(\varphi_{42}) \to \mathrm{IsSup}(c,E) \to \varphi_{100}$
\item $\varphi_{305} :\equiv M \in \mathbb{R} \land \forall n \in \mathbb{N}\,(\lvert a_{n} \rvert \leq M) \to \forall x\,(\varphi_{42}) \to \mathrm{IsSup}(c,E) \to \varphi_{109}$
\item $\varphi_{306} :\equiv M \in \mathbb{R} \land \forall n \in \mathbb{N}\,(\lvert a_{n} \rvert \leq M) \to \forall x\,(\varphi_{42}) \to \mathrm{IsSup}(c,E) \to \varphi_{120}$
\item $\varphi_{307} :\equiv M \in \mathbb{R} \land \forall n \in \mathbb{N}\,(\lvert a_{n} \rvert \leq M) \to \forall x\,(\varphi_{42}) \to \mathrm{IsSup}(c,E) \to \varphi_{125}$
\item $\varphi_{308} :\equiv \mathrm{IsSup}(c,E) \to k \in \mathbb{N} \to \varphi_{28} \to x \in E \land c - \underline{k}^{-1} < x \to N \in \mathbb{N} \to \varphi_{112}$
\item $\varphi_{309} :\equiv \varphi_{182} \land \mathrm{IsSup}(c,E) \land k \in \mathbb{N} \land \neg \forall N \in \mathbb{N}\,\exists n\,(n \in \mathbb{N} \land N \leq n \land \lvert a_{n} - c \rvert < \underline{k}^{-1})$
\item $\varphi_{310} :\equiv \mathrm{IsSup}(c,E) \to k \in \mathbb{N} \to \varphi_{28} \to x \in E \land c - \underline{k}^{-1} < x \to N \in \mathbb{N} \to \varphi_{126}$
\item $\varphi_{311} :\equiv M \in \mathbb{R} \land \forall n \in \mathbb{N}\,(\lvert a_{n} \rvert \leq M) \to \forall x\,(\varphi_{42}) \to \mathrm{IsSup}(c,E) \to \varphi_{136}$
\item $\varphi_{312} :\equiv (\forall n \in \mathbb{N}\,(\lvert a_{n} \rvert \leq M) \to N \in \mathbb{N} \to \lvert a_{N} \rvert \leq M) \to \mathrm{SeqBdd}(a) \to \varphi_{232}$
\item $\varphi_{313} :\equiv M \in \mathbb{R} \land \forall n \in \mathbb{N}\,(\lvert a_{n} \rvert \leq M) \to \forall x\,(\varphi_{42}) \to \mathrm{IsSup}(c,E) \to \varphi_{145}$
\item $\varphi_{314} :\equiv M \in \mathbb{R} \land \forall n \in \mathbb{N}\,(\lvert a_{n} \rvert \leq M) \to \forall x\,(\varphi_{42}) \to \mathrm{IsSup}(c,E) \to \varphi_{148}$
\item $\varphi_{315} :\equiv M \in \mathbb{R} \land \forall n \in \mathbb{N}\,(\lvert a_{n} \rvert \leq M) \to \forall x\,(\varphi_{42}) \to \mathrm{IsSup}(c,E) \to \varphi_{158}$
\item $\varphi_{316} :\equiv \mathrm{SeqBdd}(a) \to M \in \mathbb{R} \land \forall n \in \mathbb{N}\,(\lvert a_{n} \rvert \leq M) \to \forall x\,(\varphi_{42}) \to \varphi_{161}$
\item $\varphi_{317} :\equiv M \in \mathbb{R} \land \forall n \in \mathbb{N}\,(\lvert a_{n} \rvert \leq M) \to \forall x\,(\varphi_{42}) \to \mathrm{IsSup}(c,E) \to \varphi_{163}$
\item $\varphi_{318} :\equiv \varphi_{182} \land \mathrm{IsSup}(c,E) \land k \in \mathbb{N} \land \varphi_{28} \land (x \in E \land c - \underline{k}^{-1} < x)$
\item $\varphi_{319} :\equiv \mathrm{SeqBdd}(a) \to M \in \mathbb{R} \land \forall n \in \mathbb{N}\,(\lvert a_{n} \rvert \leq M) \to \forall x\,(\varphi_{42}) \to \varphi_{168}$
\item $\varphi_{320} :\equiv (\forall x \in E\,(x \leq c) \to c + \underline{k}^{-1} \in E \to c + \underline{k}^{-1} \leq c) \to \mathrm{SeqBdd}(a) \to M \in \mathbb{R} \land \forall n \in \mathbb{N}\,(\lvert a_{n} \rvert \leq M) \to \varphi_{191}$
\item $\varphi_{321} :\equiv M \in \mathbb{R} \land \forall n \in \mathbb{N}\,(\lvert a_{n} \rvert \leq M) \to \forall x\,(\varphi_{42}) \to \mathrm{IsSup}(c,E) \to \varphi_{167}$
\item $\varphi_{322} :\equiv M \in \mathbb{R} \land \forall n \in \mathbb{N}\,(\lvert a_{n} \rvert \leq M) \to \forall x\,(\varphi_{42}) \to \mathrm{IsSup}(c,E) \to \varphi_{171}$
\item $\varphi_{323} :\equiv M \in \mathbb{R} \land \forall n \in \mathbb{N}\,(\lvert a_{n} \rvert \leq M) \to \forall x\,(\varphi_{42}) \to \mathrm{IsSup}(c,E) \to \varphi_{173}$
\item $\varphi_{324} :\equiv M \in \mathbb{R} \land \forall n \in \mathbb{N}\,(\lvert a_{n} \rvert \leq M) \to \forall x\,(\varphi_{42}) \to \mathrm{IsSup}(c,E) \to \varphi_{181}$
\item $\varphi_{325} :\equiv \mathrm{SeqBdd}(a) \to M \in \mathbb{R} \land \forall n \in \mathbb{N}\,(\lvert a_{n} \rvert \leq M) \to \forall x\,(\varphi_{42}) \to \varphi_{187}$
\item $\varphi_{326} :\equiv M \in \mathbb{R} \land \forall n \in \mathbb{N}\,(\lvert a_{n} \rvert \leq M) \to \forall x\,(\varphi_{42}) \to \mathrm{IsSup}(c,E) \to \varphi_{184}$
\item $\varphi_{327} :\equiv M \in \mathbb{R} \land \forall n \in \mathbb{N}\,(\lvert a_{n} \rvert \leq M) \to \forall x\,(\varphi_{42}) \to \mathrm{IsSup}(c,E) \to \varphi_{190}$
\item $\varphi_{328} :\equiv (\forall s\,(\varphi_{22}) \leftrightarrow \forall u_{3}\,(\varphi_{27})) \to (\neg \forall s\,(\varphi_{22}) \leftrightarrow \neg \forall u_{3}\,(\varphi_{27}))$
\item $\varphi_{329} :\equiv \mathrm{SeqBdd}(a) \to M \in \mathbb{R} \land \forall n \in \mathbb{N}\,(\lvert a_{n} \rvert \leq M) \to \forall x\,(\varphi_{42}) \to \varphi_{198}$
\item $\varphi_{330} :\equiv M \in \mathbb{R} \land \forall n \in \mathbb{N}\,(\lvert a_{n} \rvert \leq M) \to \forall x\,(\varphi_{42}) \to \mathrm{IsSup}(c,E) \to \varphi_{195}$
\item $\varphi_{331} :\equiv (\neg \forall s\,(\varphi_{22}) \leftrightarrow \neg \forall u_{3}\,(\varphi_{27})) \to (\neg \neg \forall s\,(\varphi_{22}) \leftrightarrow \neg \neg \forall u_{3}\,(\varphi_{27}))$
\item $\varphi_{332} :\equiv M \in \mathbb{R} \land \forall n \in \mathbb{N}\,(\lvert a_{n} \rvert \leq M) \to \forall x\,(\varphi_{42}) \to \mathrm{IsSup}(c,E) \to \varphi_{203}$
\item $\varphi_{333} :\equiv M \in \mathbb{R} \land \forall n \in \mathbb{N}\,(\lvert a_{n} \rvert \leq M) \to \forall x\,(\varphi_{42}) \to \mathrm{IsSup}(c,E) \to k \in \mathbb{N} \to \varphi_{28} \to \varphi_{28}$
\item $\varphi_{334} :\equiv M \in \mathbb{R} \land \forall n \in \mathbb{N}\,(\lvert a_{n} \rvert \leq M) \to \forall x\,(\varphi_{42}) \to \mathrm{IsSup}(c,E) \to \varphi_{204}$
\item $\varphi_{335} :\equiv M \in \mathbb{R} \land \forall n \in \mathbb{N}\,(\lvert a_{n} \rvert \leq M) \to \forall x\,(\varphi_{42}) \to \mathrm{IsSup}(c,E) \to \varphi_{205}$
\item $\varphi_{336} :\equiv M \in \mathbb{R} \land \forall n \in \mathbb{N}\,(\lvert a_{n} \rvert \leq M) \to \forall x\,(\varphi_{42}) \to \mathrm{IsSup}(c,E) \to \varphi_{214}$
\item $\varphi_{337} :\equiv \mathrm{SeqBdd}(a) \to M \in \mathbb{R} \land \forall n \in \mathbb{N}\,(\lvert a_{n} \rvert \leq M) \to \forall x\,(\varphi_{42}) \to \varphi_{218}$
\item $\varphi_{338} :\equiv M \in \mathbb{R} \land \forall n \in \mathbb{N}\,(\lvert a_{n} \rvert \leq M) \to \forall x\,(\varphi_{42}) \to \mathrm{IsSup}(c,E) \to k \in \mathbb{N} \to \varphi_{196}$
\item $\varphi_{339} :\equiv (\forall \alpha\,(\neg \mathrm{Conv}(a \circ s,\alpha)) \to \neg \mathrm{Conv}(a \circ s,c)) \to \mathrm{SeqBdd}(a) \to \varphi_{306}$
\item $\varphi_{340} :\equiv M \in \mathbb{R} \land \forall n \in \mathbb{N}\,(\lvert a_{n} \rvert \leq M) \to \forall x\,(\varphi_{42}) \to \mathrm{IsSup}(c,E) \to \varphi_{224}$
\item $\varphi_{341} :\equiv M \in \mathbb{R} \land \forall n \in \mathbb{N}\,(\lvert a_{n} \rvert \leq M) \to \forall x\,(\varphi_{42}) \to \mathrm{IsSup}(c,E) \to \varphi_{225}$
\item $\varphi_{342} :\equiv M \in \mathbb{R} \land \forall n \in \mathbb{N}\,(\lvert a_{n} \rvert \leq M) \to \forall x\,(\varphi_{42}) \to \mathrm{IsSup}(c,E) \to k \in \mathbb{N} \to \varphi_{28} \to \varphi_{74}$
\item $\varphi_{343} :\equiv \varphi_{182} \land \mathrm{IsSup}(c,E) \land k \in \mathbb{N} \land \varphi_{28} \land \neg \forall N \in \mathbb{N}\,\exists n\,(n \in \mathbb{N} \land N \leq n \land \lvert a_{n} - c \rvert < \underline{k}^{-1})$
\item $\varphi_{344} :\equiv M \in \mathbb{R} \land \forall n \in \mathbb{N}\,(\lvert a_{n} \rvert \leq M) \to \forall x\,(\varphi_{42}) \to \mathrm{IsSup}(c,E) \to \varphi_{235}$
\item $\varphi_{345} :\equiv (\forall M\,(\varphi_{5} \to \neg (M \in \mathbb{R} \land \forall n \in \mathbb{N}\,(\lvert a_{n} \rvert \leq M))) \to \varphi_{103}) \to (\varphi_{5} \to \forall M\,(\varphi_{5})) \to \varphi_{4}$
\item $\varphi_{346} :\equiv M \in \mathbb{R} \land \forall n \in \mathbb{N}\,(\lvert a_{n} \rvert \leq M) \to \forall x\,(\varphi_{42}) \to \mathrm{IsSup}(c,E) \to k \in \mathbb{N} \to \varphi_{28} \to \varphi_{95}$
\item $\varphi_{347} :\equiv M \in \mathbb{R} \land \forall n \in \mathbb{N}\,(\lvert a_{n} \rvert \leq M) \to \forall x\,(\varphi_{42}) \to \mathrm{IsSup}(c,E) \to \varphi_{239}$
\item $\varphi_{348} :\equiv M \in \mathbb{R} \land \forall n \in \mathbb{N}\,(\lvert a_{n} \rvert \leq M) \to \forall x\,(\varphi_{42}) \to \mathrm{IsSup}(c,E) \to k \in \mathbb{N} \to \varphi_{28} \to \varphi_{104}$
\item $\varphi_{349} :\equiv M \in \mathbb{R} \land \forall n \in \mathbb{N}\,(\lvert a_{n} \rvert \leq M) \to \forall x\,(\varphi_{42}) \to \mathrm{IsSup}(c,E) \to k \in \mathbb{N} \to \varphi_{28} \to \varphi_{110}$
\item $\varphi_{350} :\equiv M \in \mathbb{R} \land \forall n \in \mathbb{N}\,(\lvert a_{n} \rvert \leq M) \to \forall x\,(\varphi_{42}) \to \mathrm{IsSup}(c,E) \to \varphi_{244}$
\item $\varphi_{351} :\equiv M \in \mathbb{R} \land \forall n \in \mathbb{N}\,(\lvert a_{n} \rvert \leq M) \to \forall x\,(\varphi_{42}) \to \mathrm{IsSup}(c,E) \to \varphi_{250}$
\item $\varphi_{352} :\equiv M \in \mathbb{R} \land \forall n \in \mathbb{N}\,(\lvert a_{n} \rvert \leq M) \to \forall x\,(\varphi_{42}) \to \mathrm{IsSup}(c,E) \to k \in \mathbb{N} \to \varphi_{28} \to \varphi_{127}$
\item $\varphi_{353} :\equiv M \in \mathbb{R} \land \forall n \in \mathbb{N}\,(\lvert a_{n} \rvert \leq M) \to \forall x\,(\varphi_{42}) \to \mathrm{IsSup}(c,E) \to k \in \mathbb{N} \to \varphi_{28} \to \varphi_{128}$
\item $\varphi_{354} :\equiv M \in \mathbb{R} \land \forall n \in \mathbb{N}\,(\lvert a_{n} \rvert \leq M) \to \forall x\,(\varphi_{42}) \to \mathrm{IsSup}(c,E) \to k \in \mathbb{N} \to \varphi_{28} \to \varphi_{130}$
\item $\varphi_{355} :\equiv M \in \mathbb{R} \land \forall n \in \mathbb{N}\,(\lvert a_{n} \rvert \leq M) \to \forall x\,(\varphi_{42}) \to \mathrm{IsSup}(c,E) \to \varphi_{252}$
\item $\varphi_{356} :\equiv M \in \mathbb{R} \land \forall n \in \mathbb{N}\,(\lvert a_{n} \rvert \leq M) \to \forall x\,(\varphi_{42}) \to \mathrm{IsSup}(c,E) \to k \in \mathbb{N} \to \varphi_{28} \to \varphi_{132}$
\item $\varphi_{357} :\equiv \varphi_{296} \land \varepsilon \in \mathbb{R} \land 0 < \varepsilon \land \neg \exists N \in \mathbb{N}\,\forall n \in \mathbb{N}\,(N \leq n \to \lvert a \circ s(n) - c \rvert < \varepsilon)$
\item $\varphi_{358} :\equiv M \in \mathbb{R} \land \forall n \in \mathbb{N}\,(\lvert a_{n} \rvert \leq M) \to \forall x\,(\varphi_{42}) \to \mathrm{IsSup}(c,E) \to k \in \mathbb{N} \to \varphi_{28} \to \varphi_{139}$
\item $\varphi_{359} :\equiv M \in \mathbb{R} \land \forall n \in \mathbb{N}\,(\lvert a_{n} \rvert \leq M) \to \forall x\,(\varphi_{42}) \to \mathrm{IsSup}(c,E) \to k \in \mathbb{N} \to \varphi_{28} \to \varphi_{142}$
\item $\varphi_{360} :\equiv M \in \mathbb{R} \land \forall n \in \mathbb{N}\,(\lvert a_{n} \rvert \leq M) \to \forall x\,(\varphi_{42}) \to \mathrm{IsSup}(c,E) \to \varphi_{254}$
\item $\varphi_{361} :\equiv M \in \mathbb{R} \land \forall n \in \mathbb{N}\,(\lvert a_{n} \rvert \leq M) \to \forall x\,(\varphi_{42}) \to \mathrm{IsSup}(c,E) \to \varphi_{255}$
\item $\varphi_{362} :\equiv M \in \mathbb{R} \land \forall n \in \mathbb{N}\,(\lvert a_{n} \rvert \leq M) \to \forall x\,(\varphi_{42}) \to \mathrm{IsSup}(c,E) \to \varphi_{256}$
\item $\varphi_{363} :\equiv M \in \mathbb{R} \land \forall n \in \mathbb{N}\,(\lvert a_{n} \rvert \leq M) \to \forall x\,(\varphi_{42}) \to \mathrm{IsSup}(c,E) \to \varphi_{259}$
\item $\varphi_{364} :\equiv M \in \mathbb{R} \land \forall n \in \mathbb{N}\,(\lvert a_{n} \rvert \leq M) \to \forall x\,(\varphi_{42}) \to \mathrm{IsSup}(c,E) \to k \in \mathbb{N} \to \varphi_{28} \to \varphi_{146}$
\item $\varphi_{365} :\equiv \forall u\,(\varphi_{182} \to u \in E \to u \in \mathbb{R}) \to \forall u\,(\varphi_{182}) \to \forall u \in E\,(u \in \mathbb{R})$
\item $\varphi_{366} :\equiv M \in \mathbb{R} \land \forall n \in \mathbb{N}\,(\lvert a_{n} \rvert \leq M) \to \forall x\,(\varphi_{42}) \to \mathrm{IsSup}(c,E) \to \varphi_{260}$
\item $\varphi_{367} :\equiv M \in \mathbb{R} \land \forall n \in \mathbb{N}\,(\lvert a_{n} \rvert \leq M) \to \forall x\,(\varphi_{42}) \to \mathrm{IsSup}(c,E) \to k \in \mathbb{N} \to \varphi_{28} \to \varphi_{151}$
\item $\varphi_{368} :\equiv \varphi_{182} \land \neg \exists b\,(\mathrm{SubSeq}(b,a) \land \mathrm{Cgt}(b)) \to \forall c\,(\varphi_{182} \land \neg \exists b\,(\mathrm{SubSeq}(b,a) \land \mathrm{Cgt}(b)))$
\item $\varphi_{369} :\equiv (\mathrm{Seq}(v_{733}) \leftrightarrow \mathrm{Seq}(v_{733})) \to (\neg \neg \forall s\,(\varphi_{22}) \leftrightarrow \neg \neg \forall u_{3}\,(\varphi_{27})) \to \varphi_{231}$
\item $\varphi_{370} :\equiv M \in \mathbb{R} \land \forall n \in \mathbb{N}\,(\lvert a_{n} \rvert \leq M) \to \forall x\,(\varphi_{42}) \to \mathrm{IsSup}(c,E) \to \varphi_{262}$
\item $\varphi_{371} :\equiv M \in \mathbb{R} \land \forall n \in \mathbb{N}\,(\lvert a_{n} \rvert \leq M) \to \forall x\,(\varphi_{42}) \to \mathrm{IsSup}(c,E) \to \varphi_{263}$
\item $\varphi_{372} :\equiv \varphi_{318} \to N \in \mathbb{N} \to \exists n\,(n \in \mathbb{N} \land N \leq n \land \lvert a_{n} - c \rvert < \underline{k}^{-1})$
\item $\varphi_{373} :\equiv \varphi_{318} \land N \in \mathbb{N} \land \neg \exists n\,(n \in \mathbb{N} \land N \leq n \land \lvert a_{n} - c \rvert < \underline{k}^{-1})$
\item $\varphi_{374} :\equiv \varphi_{296} \land \varepsilon \in \mathbb{R} \land 0 < \varepsilon \land (K \in \mathbb{N} \land \underline{K}^{-1} < \varepsilon) \to n \in \mathbb{N} \to K \leq n \to \lvert a \circ s(n) - c \rvert < \varepsilon$
\item $\varphi_{375} :\equiv M \in \mathbb{R} \land \forall n \in \mathbb{N}\,(\lvert a_{n} \rvert \leq M) \to \forall x\,(\varphi_{42}) \to \mathrm{IsSup}(c,E) \to \varphi_{266}$
\item $\varphi_{376} :\equiv \forall N\,(\varphi_{318}) \to \forall N \in \mathbb{N}\,\exists n\,(n \in \mathbb{N} \land N \leq n \land \lvert a_{n} - c \rvert < \underline{k}^{-1})$
\item $\varphi_{377} :\equiv \forall n\,(\varphi_{296} \land \varepsilon \in \mathbb{R} \land 0 < \varepsilon \land (K \in \mathbb{N} \land \underline{K}^{-1} < \varepsilon)) \to \forall n \in \mathbb{N}\,(K \leq n \to \lvert a \circ s(n) - c \rvert < \varepsilon)$
\item $\varphi_{378} :\equiv M \in \mathbb{R} \land \forall n \in \mathbb{N}\,(\lvert a_{n} \rvert \leq M) \to \forall x\,(\varphi_{42}) \to \mathrm{IsSup}(c,E) \to \varphi_{268}$
\item $\varphi_{379} :\equiv \mathrm{SeqBdd}(a) \to M \in \mathbb{R} \land \forall n \in \mathbb{N}\,(\lvert a_{n} \rvert \leq M) \to \forall x\,(\varphi_{42}) \to \varphi_{273}$
\item $\varphi_{380} :\equiv \mathrm{SeqBdd}(a) \to M \in \mathbb{R} \land \forall n \in \mathbb{N}\,(\lvert a_{n} \rvert \leq M) \to \forall x\,(\varphi_{42}) \to \varphi_{274}$
\item $\varphi_{381} :\equiv \mathrm{SeqBdd}(a) \to M \in \mathbb{R} \land \forall n \in \mathbb{N}\,(\lvert a_{n} \rvert \leq M) \to \forall x\,(\varphi_{42}) \to \varphi_{281}$
\item $\varphi_{382} :\equiv \mathrm{SeqBdd}(a) \to M \in \mathbb{R} \land \forall n \in \mathbb{N}\,(\lvert a_{n} \rvert \leq M) \to \forall x\,(\varphi_{42}) \to \varphi_{282}$
\item $\varphi_{383} :\equiv (\mathrm{SeqBdd}(a) \to \mathrm{Seq}(a)) \to (\mathrm{SeqBdd}(a) \to \varphi_{209}) \to \mathrm{SeqBdd}(a) \to \varphi_{223}$
\item $\varphi_{384} :\equiv (\mathrm{SeqBdd}(a) \to \varphi_{202}) \to \mathrm{SeqBdd}(a) \to M \in \mathbb{R} \land \forall n \in \mathbb{N}\,(\lvert a_{n} \rvert \leq M) \to \varphi_{185}$
\item $\varphi_{385} :\equiv \mathrm{SeqBdd}(a) \to M \in \mathbb{R} \land \forall n \in \mathbb{N}\,(\lvert a_{n} \rvert \leq M) \to \forall x\,(\varphi_{42}) \to \varphi_{292}$
\item $\varphi_{386} :\equiv (E \subseteq \mathbb{R} \leftrightarrow \forall u \in E\,(u \in \mathbb{R})) \to (\mathrm{SeqBdd}(a) \to \varphi_{215}) \to \varphi_{210}$
\item $\varphi_{387} :\equiv \mathrm{SeqBdd}(a) \to M \in \mathbb{R} \land \forall n \in \mathbb{N}\,(\lvert a_{n} \rvert \leq M) \to \forall x\,(\varphi_{42}) \to \varphi_{295}$
\item $\varphi_{388} :\equiv \forall a\,\forall c\,(\mathrm{Seq}(a) \to \varphi_{93} \to \varphi_{47}) \to \forall c\,(\mathrm{Seq}(a) \to \varphi_{93} \to \varphi_{47})$
\item $\varphi_{389} :\equiv \mathrm{SeqBdd}(a) \to M \in \mathbb{R} \land \forall n \in \mathbb{N}\,(\lvert a_{n} \rvert \leq M) \to \forall x\,(\varphi_{42}) \to \varphi_{308}$
\item $\varphi_{390} :\equiv (\mathrm{SeqBdd}(a) \to \varphi_{227}) \to (\underline{k} = 0 \to 0 = \underline{k}) \to \mathrm{SeqBdd}(a) \to \varphi_{228}$
\item $\varphi_{391} :\equiv \mathrm{SeqBdd}(a) \to M \in \mathbb{R} \land \forall n \in \mathbb{N}\,(\lvert a_{n} \rvert \leq M) \to \forall x\,(\varphi_{42}) \to \varphi_{310}$
\item $\varphi_{392} :\equiv (\mathrm{SeqBdd}(a) \to M \in \mathbb{R} \land \forall n \in \mathbb{N}\,(\lvert a_{n} \rvert \leq M) \to \varphi_{166}) \to \varphi_{258}$
\item $\varphi_{393} :\equiv (\mathrm{SeqBdd}(a) \to \varphi_{246}) \to \varphi_{182} \to N \in \mathbb{N} \to \exists n\,(n \in \mathbb{N} \land N \leq n \land -M \leq a_{n})$
\item $\varphi_{394} :\equiv (\mathrm{SeqBdd}(a) \to \varphi_{242}) \to \mathrm{SeqBdd}(a) \to M \in \mathbb{R} \land \forall n \in \mathbb{N}\,(\lvert a_{n} \rvert \leq M) \to \varphi_{172}$
\item $\varphi_{395} :\equiv (k \in \mathbb{N} \to 0 < \underline{k}) \to (\mathrm{SeqBdd}(a) \to \varphi_{226}) \to \mathrm{SeqBdd}(a) \to \varphi_{230}$
\item $\varphi_{396} :\equiv (0 < \underline{k} \to 0 \neq \underline{k}) \to (\mathrm{SeqBdd}(a) \to \varphi_{230}) \to \mathrm{SeqBdd}(a) \to \varphi_{227}$
\item $\varphi_{397} :\equiv (\mathrm{SeqBdd}(a) \to M \in \mathbb{R} \land \forall n \in \mathbb{N}\,(\lvert a_{n} \rvert \leq M) \to \varphi_{177}) \to \mathrm{SeqBdd}(a) \to M \in \mathbb{R} \land \forall n \in \mathbb{N}\,(\lvert a_{n} \rvert \leq M) \to \varphi_{166}$
\item $\varphi_{398} :\equiv (\forall N_{1}\,(\neg \varphi_{28}) \to \forall N\,\forall N_{1}\,(\neg \varphi_{28})) \to \mathrm{SeqBdd}(a) \to \varphi_{298}$
\item $\varphi_{399} :\equiv \forall N\,(\varphi_{182} \to N \in \mathbb{N} \to \exists n\,(n \in \mathbb{N} \land N \leq n \land -M \leq a_{n})) \to \forall N\,(\varphi_{182}) \to \forall N \in \mathbb{N}\,\exists n\,(n \in \mathbb{N} \land N \leq n \land -M \leq a_{n})$
\item $\varphi_{400} :\equiv (k \in \mathbb{N} \to \underline{k} \in \mathbb{R}) \to (\mathrm{SeqBdd}(a) \to \varphi_{226}) \to \mathrm{SeqBdd}(a) \to \varphi_{236}$
\item $\varphi_{401} :\equiv (\mathrm{SeqBdd}(a) \to M \in \mathbb{R} \land \forall n \in \mathbb{N}\,(\lvert a_{n} \rvert \leq M) \to \varphi_{175}) \to \mathrm{SeqBdd}(a) \to M \in \mathbb{R} \land \forall n \in \mathbb{N}\,(\lvert a_{n} \rvert \leq M) \to \varphi_{192}$
\item $\varphi_{402} :\equiv (\mathrm{SeqBdd}(a) \to M \in \mathbb{R} \land \forall n \in \mathbb{N}\,(\lvert a_{n} \rvert \leq M) \to \varphi_{175}) \to \mathrm{SeqBdd}(a) \to M \in \mathbb{R} \land \forall n \in \mathbb{N}\,(\lvert a_{n} \rvert \leq M) \to \varphi_{199}$
\item $\varphi_{403} :\equiv \varphi_{280} \to \mathrm{SeqBdd}(a) \to M \in \mathbb{R} \land \forall n \in \mathbb{N}\,(\lvert a_{n} \rvert \leq M) \to \varphi_{177}$
\item $\varphi_{404} :\equiv (a \circ s(n) = a_{s_{n}} \to a_{s_{n}} = a \circ s(n)) \to \mathrm{SeqBdd}(a) \to \varphi_{362}$
\item $\varphi_{405} :\equiv \varphi_{287} \to \mathrm{SeqBdd}(a) \to M \in \mathbb{R} \land \forall n \in \mathbb{N}\,(\lvert a_{n} \rvert \leq M) \to \varphi_{175}$
\item $\varphi_{406} :\equiv \varphi_{287} \to \mathrm{SeqBdd}(a) \to M \in \mathbb{R} \land \forall n \in \mathbb{N}\,(\lvert a_{n} \rvert \leq M) \to \varphi_{178}$
\item $\varphi_{407} :\equiv (\mathrm{SeqBdd}(a) \to M \in \mathbb{R} \land \forall n \in \mathbb{N}\,(\lvert a_{n} \rvert \leq M) \to \varphi_{192}) \to \varphi_{280}$
\item $\varphi_{408} :\equiv (\forall s\,(\varphi_{22}) \to \forall u_{3}\,(\varphi_{27})) \to (\forall u_{3}\,(\varphi_{27}) \to \forall s\,(\varphi_{22})) \to (\forall s\,(\varphi_{22}) \leftrightarrow \forall u_{3}\,(\varphi_{27}))$
\item $\varphi_{409} :\equiv (\mathrm{UB}(M,E) \leftrightarrow M \in \mathbb{R} \land \forall x \in E\,(x \leq M)) \to \varphi_{40} \to (\mathrm{SeqBdd}(a) \to \varphi_{200}) \to \mathrm{SeqBdd}(a) \to \varphi_{176}$
\item $\varphi_{410} :\equiv (\mathrm{SeqBdd}(a) \to M \in \mathbb{R} \land \forall n \in \mathbb{N}\,(\lvert a_{n} \rvert \leq M) \to \varphi_{172}) \to \mathrm{SeqBdd}(a) \to \varphi_{290}$
\item $\varphi_{411} :\equiv (\mathrm{SubSeq}(v_{733},v_{734}) \to \neg (\mathrm{Seq}(v_{733}) \to \neg \neg \forall s\,(\varphi_{22}))) \to \neg (\neg (\mathrm{Seq}(v_{733}) \to \neg \neg \forall s\,(\varphi_{22})) \to \mathrm{SubSeq}(v_{733},v_{734})) \leftrightarrow (\mathrm{SubSeq}(v_{733},v_{734}) \to \neg (\mathrm{Seq}(v_{733}) \to \neg \neg \forall u_{3}\,(\varphi_{27}))) \to \neg (\neg (\mathrm{Seq}(v_{733}) \to \neg \neg \forall u_{3}\,(\varphi_{27})) \to \mathrm{SubSeq}(v_{733},v_{734}))$
\item $\varphi_{412} :\equiv (\mathrm{Seq}(a) \to n \in \mathbb{N} \to a_{n} \in \mathbb{R}) \to (\mathrm{SeqBdd}(a) \to \mathrm{Seq}(a)) \to \varphi_{402}$
\item $\varphi_{413} :\equiv (\mathrm{SeqBdd}(a) \to \exists M \in \mathbb{R}\,\forall n \in \mathbb{N}\,(\lvert a_{n} \rvert \leq M)) \to \forall M\,(\varphi_{5} \to \neg (M \in \mathbb{R} \land \forall n \in \mathbb{N}\,(\lvert a_{n} \rvert \leq M))) \to \varphi_{345}$
\item $\varphi_{414} :\equiv (\varphi_{182} \land x \in E \land \neg x \leq M \to \forall n\,(\varphi_{182} \land x \in E \land \neg x \leq M)) \to \mathrm{SeqBdd}(a) \to \varphi_{193}$
\item $\varphi_{415} :\equiv \forall k\,(\varphi_{182} \land \mathrm{IsSup}(c,E) \to \varphi_{82}) \to \forall k\,(\varphi_{182} \land \mathrm{IsSup}(c,E)) \to \varphi_{93}$
\item $\varphi_{416} :\equiv (\mathrm{SeqBdd}(a) \to \mathrm{Seq}(a)) \to (\mathrm{SeqBdd}(a) \to \varphi_{302}) \to \mathrm{SeqBdd}(a) \to \varphi_{283}$
\item $\varphi_{417} :\equiv \varphi_{284} \to (\forall n \in \mathbb{N}\,(\lvert a_{n} \rvert \leq M) \to n \in \mathbb{N} \to \lvert a_{n} \rvert \leq M) \to \varphi_{277}$
\item $\varphi_{418} :\equiv (\mathrm{SeqBdd}(a) \to M \in \mathbb{R} \land \forall n \in \mathbb{N}\,(\lvert a_{n} \rvert \leq M) \to \varphi_{185}) \to \varphi_{135} \to \varphi_{264}$
\item $\varphi_{419} :\equiv (\mathrm{SeqBdd}(a) \to \varphi_{232}) \to (\mathrm{SeqBdd}(a) \to \varphi_{209}) \to \mathrm{SeqBdd}(a) \to \varphi_{220}$
\item $\varphi_{420} :\equiv (\forall u_{3}\,(\varphi_{32}) \to \forall t\,(\varphi_{25})) \to (\forall t\,(\varphi_{25}) \to \forall u_{3}\,(\varphi_{32})) \to \mathrm{SeqBdd}(a) \to \varphi_{305}$
\item $\varphi_{421} :\equiv (\mathrm{Cgt}(a \circ s) \leftrightarrow \exists \alpha\,(\mathrm{Conv}(a \circ s,\alpha))) \to (\mathrm{SeqBdd}(a) \to \varphi_{306}) \to \mathrm{SeqBdd}(a) \to \varphi_{301}$
\item $\varphi_{422} :\equiv (\mathrm{SeqBdd}(a) \to \varphi_{209}) \to N = N \to (\mathrm{SeqBdd}(a) \to \varphi_{233}) \to \mathrm{SeqBdd}(a) \to \varphi_{240}$
\item $\varphi_{423} :\equiv (\mathrm{SeqBdd}(a) \to \varphi_{236}) \to (\mathrm{SeqBdd}(a) \to \varphi_{230}) \to \mathrm{SeqBdd}(a) \to \varphi_{237}$
\item $\varphi_{424} :\equiv (\varphi_{182} \land \mathrm{IsSup}(c,E) \to \forall k\,(\varphi_{182} \land \mathrm{IsSup}(c,E))) \to \mathrm{SeqBdd}(a) \to \varphi_{302}$
\item $\varphi_{425} :\equiv (\mathrm{SeqBdd}(a) \to \varphi_{236}) \to (\mathrm{SeqBdd}(a) \to \varphi_{228}) \to \mathrm{SeqBdd}(a) \to \varphi_{242}$
\item $\varphi_{426} :\equiv (\mathrm{SeqBdd}(a) \to \mathrm{Seq}(a)) \to (\mathrm{SeqBdd}(a) \to \varphi_{294}) \to \mathrm{SeqBdd}(a) \to \varphi_{330}$
\item $\varphi_{427} :\equiv (\forall u_{3}\,(\forall s\,(\varphi_{22}) \to \varphi_{27}) \to \forall u_{3}\,\forall s\,(\varphi_{22}) \to \forall u_{3}\,(\varphi_{27})) \to (\forall s\,(\varphi_{22}) \to \forall u_{3}\,\forall s\,(\varphi_{22})) \to \forall s\,(\varphi_{22}) \to \forall u_{3}\,(\varphi_{27})$
\item $\varphi_{428} :\equiv (\forall s\,(\forall u_{3}\,(\varphi_{27}) \to \varphi_{22}) \to \forall s\,\forall u_{3}\,(\varphi_{27}) \to \forall s\,(\varphi_{22})) \to (\forall u_{3}\,(\varphi_{27}) \to \forall s\,\forall u_{3}\,(\varphi_{27})) \to \forall u_{3}\,(\varphi_{27}) \to \forall s\,(\varphi_{22})$
\item $\varphi_{429} :\equiv \varphi_{296} \land \varepsilon \in \mathbb{R} \land 0 < \varepsilon \land (K \in \mathbb{N} \land \underline{K}^{-1} < \varepsilon) \to \forall n\,(\varphi_{296} \land \varepsilon \in \mathbb{R} \land 0 < \varepsilon \land (K \in \mathbb{N} \land \underline{K}^{-1} < \varepsilon))$
\item $\varphi_{430} :\equiv (\forall u_{3}\,(\forall t\,(\varphi_{25}) \to \varphi_{32}) \to \forall u_{3}\,\forall t\,(\varphi_{25}) \to \forall u_{3}\,(\varphi_{32})) \to (\forall t\,(\varphi_{25}) \to \forall u_{3}\,\forall t\,(\varphi_{25})) \to \forall t\,(\varphi_{25}) \to \forall u_{3}\,(\varphi_{32})$
\item $\varphi_{431} :\equiv (\mathrm{SeqBdd}(a) \to M \in \mathbb{R} \land \forall n \in \mathbb{N}\,(\lvert a_{n} \rvert \leq M) \to \varphi_{178}) \to \varphi_{397}$
\item $\varphi_{432} :\equiv (\forall t\,(\forall u_{3}\,(\varphi_{32}) \to \varphi_{25}) \to \forall t\,\forall u_{3}\,(\varphi_{32}) \to \forall t\,(\varphi_{25})) \to (\forall u_{3}\,(\varphi_{32}) \to \forall t\,\forall u_{3}\,(\varphi_{32})) \to \forall u_{3}\,(\varphi_{32}) \to \forall t\,(\varphi_{25})$
\item $\varphi_{433} :\equiv (\mathrm{SeqBdd}(a) \to \varphi_{223}) \to \varphi_{40} \to (\mathrm{SeqBdd}(a) \to \varphi_{220}) \to \mathrm{SeqBdd}(a) \to \varphi_{233}$
\item $\varphi_{434} :\equiv (\forall E\,(\varphi_{80} \to \neg \forall x\,(\varphi_{42})) \to \forall E\,(\varphi_{80}) \to \forall E\,(\neg \forall x\,(\varphi_{42}))) \to (\varphi_{80} \to \forall E\,(\varphi_{80})) \to \mathrm{SeqBdd}(a) \to \varphi_{44}$
\item $\varphi_{435} :\equiv (c \in \mathbb{R} \to \underline{k}^{-1} \in \mathbb{R} \to c + \underline{k}^{-1} \in \mathbb{R}) \to (\mathrm{SeqBdd}(a) \to \varphi_{213}) \to \varphi_{394}$
\item $\varphi_{436} :\equiv (\mathrm{SeqBdd}(a) \to M \in \mathbb{R} \land \forall n \in \mathbb{N}\,(\lvert a_{n} \rvert \leq M) \to \varphi_{191}) \to \varphi_{288} \to \varphi_{279}$
\item $\varphi_{437} :\equiv (\forall N\,(\varphi_{170}) \to \varphi_{188}) \to (\forall N\,(\varphi_{23}) \to \forall N\,\forall N\,(\varphi_{23})) \to \mathrm{SeqBdd}(a) \to \varphi_{291}$
\item $\varphi_{438} :\equiv \varphi_{14} \to \varphi_{210} \to \varphi_{211} \to (\mathrm{SeqBdd}(a) \to \varphi_{186}) \to \mathrm{SeqBdd}(a) \to \varphi_{194}$
\item $\varphi_{439} :\equiv (\forall N\,(\forall N_{1}\,(\neg \varphi_{28}) \to \varphi_{23}) \to \forall N\,\forall N_{1}\,(\neg \varphi_{28}) \to \forall N\,(\varphi_{23})) \to \varphi_{398}$
\item $\varphi_{440} :\equiv N \leq N \cup N_{1} \to (\mathrm{SeqBdd}(a) \to \varphi_{356}) \to \mathrm{SeqBdd}(a) \to \varphi_{348}$
\item $\varphi_{441} :\equiv N_{1} \leq N \cup N_{1} \to (\mathrm{SeqBdd}(a) \to \varphi_{356}) \to \mathrm{SeqBdd}(a) \to \varphi_{349}$
\item $\varphi_{442} :\equiv (a_{N} \in \mathbb{R} \to M \in \mathbb{R} \to \lvert a_{N} \rvert \leq M \to -M \leq a_{N} \land a_{N} \leq M) \to \varphi_{433}$
\item $\varphi_{443} :\equiv (\mathrm{SeqBdd}(a) \to \varphi_{364}) \to \varphi_{343} \to \neg (x \in E \land c - \underline{k}^{-1} < x)$
\item $\varphi_{444} :\equiv (\mathrm{SeqBdd}(a) \to \mathrm{Seq}(a)) \to (\mathrm{SeqBdd}(a) \to \varphi_{353}) \to \mathrm{SeqBdd}(a) \to \varphi_{359}$
\item $\varphi_{445} :\equiv (\mathrm{SeqBdd}(a) \to M \in \mathbb{R} \land \forall n \in \mathbb{N}\,(\lvert a_{n} \rvert \leq M) \to \varphi_{199}) \to \varphi_{40} \to \varphi_{407}$
\item $\varphi_{446} :\equiv \forall x\,(\varphi_{343} \to \neg (x \in E \land c - \underline{k}^{-1} < x)) \to \forall x\,(\varphi_{343}) \to \forall x\,(\neg (x \in E \land c - \underline{k}^{-1} < x))$
\item $\varphi_{447} :\equiv (\mathrm{SeqBdd}(a) \to \varphi_{324}) \to (\forall t\,(\varphi_{25}) \to \varphi_{21}) \to \mathrm{SeqBdd}(a) \to \varphi_{327}$
\item $\varphi_{448} :\equiv \forall K\,(\varphi_{357} \to \neg (K \in \mathbb{N} \land \underline{K}^{-1} < \varepsilon)) \to \forall K\,(\varphi_{357}) \to \forall K\,(\neg (K \in \mathbb{N} \land \underline{K}^{-1} < \varepsilon))$
\item $\varphi_{449} :\equiv (\mathrm{SeqBdd}(a) \to \varphi_{213}) \to (\mathrm{SeqBdd}(a) \to \varphi_{351}) \to \mathrm{SeqBdd}(a) \to \varphi_{304}$
\item $\varphi_{450} :\equiv (\forall x\,(\varphi_{182} \to x \in E \to x \leq M) \to \forall x\,(\varphi_{182}) \to \forall x \in E\,(x \leq M)) \to (\varphi_{182} \to \forall x\,(\varphi_{182})) \to \mathrm{SeqBdd}(a) \to \varphi_{200}$
\item $\varphi_{451} :\equiv (\mathrm{SeqBdd}(a) \to \varphi_{305}) \to (\mathrm{SeqBdd}(a) \to \varphi_{301}) \to \mathrm{SeqBdd}(a) \to \varphi_{314}$
\item $\varphi_{452} :\equiv (\mathrm{SeqBdd}(a) \to \varphi_{213}) \to (\mathrm{SeqBdd}(a) \to \varphi_{242}) \to (\mathrm{SeqBdd}(a) \to \varphi_{237}) \to \varphi_{286}$
\item $\varphi_{453} :\equiv (c \in \mathbb{R} \to \underline{k}^{-1} \in \mathbb{R} \to c + \underline{k}^{-1} \in \mathbb{R}) \to (\mathrm{SeqBdd}(a) \to \varphi_{213}) \to (\mathrm{SeqBdd}(a) \to \varphi_{242}) \to \varphi_{379}$
\item $\varphi_{454} :\equiv \forall n\,(\varphi_{373} \to \neg (n \in \mathbb{N} \land N \cup N_{1} \leq n \land x \leq a_{n})) \to \forall n\,(\varphi_{373}) \to \forall n\,(\neg (n \in \mathbb{N} \land N \cup N_{1} \leq n \land x \leq a_{n}))$
\item $\varphi_{455} :\equiv (a_{n} \in \mathbb{R} \to M \in \mathbb{R} \to \lvert a_{n} \rvert \leq M \to -M \leq a_{n} \land a_{n} \leq M) \to \varphi_{445}$
\item $\varphi_{456} :\equiv (\mathrm{SeqBdd}(a) \to \varphi_{297}) \to (\mathrm{SeqBdd}(a) \to \varphi_{317}) \to \mathrm{SeqBdd}(a) \to \varphi_{324}$
\item $\varphi_{457} :\equiv (\mathrm{SeqBdd}(a) \to \varphi_{212}) \to (\mathrm{SeqBdd}(a) \to \varphi_{242}) \to (\mathrm{SeqBdd}(a) \to \varphi_{237}) \to \varphi_{319}$
\item $\varphi_{458} :\equiv (\mathrm{SeqBdd}(a) \to \varphi_{311}) \to (\mathrm{SeqBdd}(a) \to \varphi_{313}) \to \mathrm{SeqBdd}(a) \to \varphi_{326}$
\item $\varphi_{459} :\equiv (c \in \mathbb{R} \to \underline{k}^{-1} \in \mathbb{R} \to 0 < \underline{k}^{-1} \to c < c + \underline{k}^{-1}) \to \varphi_{452}$
\item $\varphi_{460} :\equiv (\mathrm{SeqBdd}(a) \to \varphi_{375}) \to (\varphi_{46} \to \varphi_{37}) \to \mathrm{SeqBdd}(a) \to \varphi_{378}$
\item $\varphi_{461} :\equiv (\forall v_{734}\,(\varphi_{113}) \leftrightarrow \forall v_{734}\,(\varphi_{285})) \to \forall v_{734}\,(\varphi_{285}) \to \forall v_{734}\,(\varphi_{113})$
\item $\varphi_{462} :\equiv (\forall v_{734}\,(\varphi_{113}) \leftrightarrow \forall v_{734}\,(\varphi_{285})) \to \forall v_{734}\,(\varphi_{113}) \to \forall v_{734}\,(\varphi_{285})$
\item $\varphi_{463} :\equiv (N \in \mathbb{N} \to N_{1} \in \mathbb{N} \to N \cup N_{1} \in \mathbb{N}) \to \varphi_{337} \to (\mathrm{SeqBdd}(a) \to \varphi_{303}) \to \mathrm{SeqBdd}(a) \to \varphi_{335}$
\item $\varphi_{464} :\equiv \forall v_{733}\,(\forall v_{734}\,(\varphi_{285}) \to \forall v_{734}\,(\varphi_{113})) \to \varphi_{147} \to \forall v_{733}\,\forall v_{734}\,(\varphi_{113})$
\item $\varphi_{465} :\equiv \forall v_{733}\,(\forall v_{734}\,(\varphi_{113}) \to \forall v_{734}\,(\varphi_{285})) \to \forall v_{733}\,\forall v_{734}\,(\varphi_{113}) \to \varphi_{147}$
\item $\varphi_{466} :\equiv (\mathrm{SeqBdd}(a) \to \varphi_{315}) \to (\mathrm{SeqBdd}(a) \to \varphi_{341}) \to \mathrm{SeqBdd}(a) \to \varphi_{361}$
\item $\varphi_{467} :\equiv \forall u\,(\varphi_{182} \to u \in E \to u \in \mathbb{R}) \to \varphi_{365} \to (\varphi_{182} \to \forall u\,(\varphi_{182})) \to \mathrm{SeqBdd}(a) \to \varphi_{215}$
\item $\varphi_{468} :\equiv (\mathrm{SeqBdd}(a) \to \varphi_{321}) \to (\mathrm{SeqBdd}(a) \to \varphi_{341}) \to \mathrm{SeqBdd}(a) \to \varphi_{370}$
\item $\varphi_{469} :\equiv (\mathrm{SeqBdd}(a) \to \varphi_{352}) \to (\mathrm{SeqBdd}(a) \to \varphi_{335}) \to \mathrm{SeqBdd}(a) \to \varphi_{346}$
\item $\varphi_{470} :\equiv (\mathrm{SeqBdd}(a) \to \varphi_{341}) \to (\mathrm{SeqBdd}(a) \to \varphi_{340}) \to \mathrm{SeqBdd}(a) \to \varphi_{355}$
\item $\varphi_{471} :\equiv (\mathrm{SeqBdd}(a) \to \varphi_{334}) \to (\mathrm{SeqBdd}(a) \to \varphi_{366}) \to \mathrm{SeqBdd}(a) \to \varphi_{375}$
\item $\varphi_{472} :\equiv (\forall c\,(\varphi_{249}) \to \varphi_{253}) \to \varphi_{368} \to \varphi_{1} \to (\varphi_{1} \to \varphi_{3}) \to \varphi_{2} \to \mathrm{SeqBdd}(a) \to \varphi_{216}$
\item $\varphi_{473} :\equiv (\mathrm{SeqBdd}(a) \to \varphi_{371}) \to (\mathrm{SeqBdd}(a) \to \varphi_{350}) \to \mathrm{SeqBdd}(a) \to \varphi_{360}$
\item $\varphi_{474} :\equiv (\mathrm{SeqBdd}(a) \to \varphi_{370}) \to (\mathrm{SeqBdd}(a) \to \varphi_{355}) \to \mathrm{SeqBdd}(a) \to \varphi_{371}$
\item $\varphi_{475} :\equiv (\forall s\,(\varphi_{247} \to \varphi_{41}) \to \forall s\,(\varphi_{247}) \to \forall s\,(\varphi_{41})) \to (\varphi_{247} \to \forall s\,(\varphi_{247})) \to \varphi_{245}$
\item $\varphi_{476} :\equiv \forall N\,(\varphi_{182} \to N \in \mathbb{N} \to \exists n\,(n \in \mathbb{N} \land N \leq n \land -M \leq a_{n})) \to \varphi_{399} \to (\varphi_{182} \to \forall N\,(\varphi_{182})) \to \mathrm{SeqBdd}(a) \to \varphi_{251}$
\item $\varphi_{477} :\equiv (c - \underline{k}^{-1} < x \to x \leq a_{n} \to c - \underline{k}^{-1} < a_{n}) \to (\mathrm{SeqBdd}(a) \to \varphi_{367}) \to (\mathrm{SeqBdd}(a) \to \varphi_{354}) \to \varphi_{382}$
\item $\varphi_{478} :\equiv (\mathrm{SeqBdd}(a) \to \varphi_{362}) \to (\mathrm{SeqBdd}(a) \to \varphi_{360}) \to \varphi_{140} \to \mathrm{SeqBdd}(a) \to \varphi_{363}$
\item $\varphi_{479} :\equiv \varphi_{143} \to (\mathrm{SeqBdd}(a) \to \varphi_{300}) \to (\mathrm{SeqBdd}(a) \to \varphi_{327}) \to \varphi_{420}$
\item $\varphi_{480} :\equiv (\mathrm{SeqBdd}(a) \to M \in \mathbb{R} \land \forall n \in \mathbb{N}\,(\lvert a_{n} \rvert \leq M) \to \varphi_{164}) \to \forall n\,(\varphi_{258}) \to (\forall n\,(\varphi_{258}) \to \varphi_{261}) \to \varphi_{414}$
\item $\varphi_{481} :\equiv \varphi_{329} \to (\mathrm{SeqBdd}(a) \to \varphi_{353}) \to (\mathrm{SeqBdd}(a) \to \varphi_{349}) \to \varphi_{380}$
\item $\varphi_{482} :\equiv (K \in \mathbb{N} \to n \in \mathbb{N} \to K \leq n \to \underline{n}^{-1} \leq \underline{K}^{-1}) \to (\mathrm{SeqBdd}(a) \to \varphi_{334}) \to \varphi_{470}$
\item $\varphi_{483} :\equiv (\forall v_{734}\,(\varphi_{285} \to \varphi_{113}) \to \forall v_{734}\,(\varphi_{285}) \to \forall v_{734}\,(\varphi_{113})) \to (\forall v_{734}\,(\varphi_{113}) \leftrightarrow \forall v_{734}\,(\varphi_{285}))$
\item $\varphi_{484} :\equiv (\mathrm{SeqBdd}(a) \to \varphi_{353}) \to (\mathrm{SeqBdd}(a) \to \varphi_{348}) \to \varphi_{385} \to \varphi_{389}$
\item $\varphi_{485} :\equiv (\forall \varepsilon\,(\varphi_{296} \to \varepsilon \in \mathbb{R} \to \varphi_{70}) \to \forall \varepsilon\,(\varphi_{296}) \to \forall \varepsilon \in \mathbb{R}\,(\varphi_{70})) \to (\varphi_{296} \to \forall \varepsilon\,(\varphi_{296})) \to \mathrm{SeqBdd}(a) \to \varphi_{351}$
\item $\varphi_{486} :\equiv (\forall N\,(\varphi_{372}) \to \varphi_{376}) \to (\varphi_{318} \to \forall N\,(\varphi_{318})) \to \mathrm{SeqBdd}(a) \to \varphi_{364}$
\item $\varphi_{487} :\equiv (\mathrm{SeqBdd}(a) \to \varphi_{242}) \to (\mathrm{SeqBdd}(a) \to \varphi_{359}) \to \varphi_{382} \to \varphi_{381} \to \varphi_{385}$
\item $\varphi_{488} :\equiv (c \in \mathbb{R} \to \underline{k}^{-1} \in \mathbb{R} \to a_{n} \in \mathbb{R} \to \varphi_{137}) \to (\mathrm{SeqBdd}(a) \to \varphi_{213}) \to \varphi_{487}$
\item $\varphi_{489} :\equiv (\forall N_{1}\,(\varphi_{309} \to \neg \varphi_{28}) \to \forall N_{1}\,(\varphi_{309}) \to \forall N_{1}\,(\neg \varphi_{28})) \to (\varphi_{309} \to \forall N_{1}\,(\varphi_{309})) \to \forall N\,(\forall N_{1}\,(\neg \varphi_{28}) \to \varphi_{23}) \to \varphi_{439}$
\item $\varphi_{490} :\equiv \forall x\,(\varphi_{343} \to \neg (x \in E \land c - \underline{k}^{-1} < x)) \to \varphi_{446} \to (\varphi_{343} \to \forall x\,(\varphi_{343})) \to \mathrm{SeqBdd}(a) \to \varphi_{338}$
\item $\varphi_{491} :\equiv \forall K\,(\varphi_{357} \to \neg (K \in \mathbb{N} \land \underline{K}^{-1} < \varepsilon)) \to \varphi_{448} \to (\varphi_{357} \to \forall K\,(\varphi_{357})) \to \mathrm{SeqBdd}(a) \to \varphi_{347}$
\item $\varphi_{492} :\equiv \forall n\,(\varphi_{373} \to \neg (n \in \mathbb{N} \land N \cup N_{1} \leq n \land x \leq a_{n})) \to \varphi_{454} \to (\varphi_{373} \to \forall n\,(\varphi_{373})) \to \mathrm{SeqBdd}(a) \to \varphi_{358}$
\item $\varphi_{493} :\equiv \forall v_{734}\,(\varphi_{285} \to \varphi_{113}) \to (\forall v_{734}\,(\varphi_{113} \to \varphi_{285}) \to \forall v_{734}\,(\varphi_{113}) \to \forall v_{734}\,(\varphi_{285})) \to \varphi_{483}$
\item $\varphi_{494} :\equiv \forall v_{733}\,(\forall v_{734}\,(\varphi_{285}) \to \forall v_{734}\,(\varphi_{113})) \to \varphi_{465} \to \varphi_{464} \to (\forall v_{733}\,\forall v_{734}\,(\varphi_{113}) \leftrightarrow \varphi_{147})$
\end{abbrevlist}
\end{fullproof}

% 自動生成 (prover/build.js): CauchyE  ノード数 19
\begin{fullproof}{[C26]}{$\forall a\,\forall \varepsilon\,(\mathrm{Cauchy}(a) \to \varepsilon \in \mathbb{R} \to 0 < \varepsilon \to \exists N \in \mathbb{N}\,\forall m\,\forall n\,(m \in \mathbb{N} \to n \in \mathbb{N} \to N \leq m \to N \leq n \to \lvert a_{m} - a_{n} \rvert < \varepsilon))$}{19}\phantomsection\label{tree:CauchyE}
\noindent$\Pi_{1}$:
\begin{ptree}
\AXC{$\mathrm{Cauchy}(a) \to \mathrm{Cauchy}(a)$\lbl{A1}}
\AXC{$(\mathrm{Cauchy}(a) \to \mathrm{Cauchy}(a)) \to \varphi_{6}$\lbl{A1}}
\BIC{$\varphi_{6}$}
\end{ptree}
\noindent$\Pi_{2}$:
\begin{ptree}
\AXC{$\forall \varepsilon \in \mathbb{R}\,(0 < \varepsilon \to \varphi_{1}) \to \varepsilon \in \mathbb{R} \to 0 < \varepsilon \to \varphi_{1}$\lbl{A2}}
\AXC{$\varphi_{2}$\lbl{DefCauchy}}
\AXC{$\varphi_{5}$\lbl{A2}}
\BIC{$\mathrm{Cauchy}(a) \leftrightarrow \mathrm{Seq}(a) \land \forall \varepsilon \in \mathbb{R}\,(0 < \varepsilon \to \varphi_{1})$}
\AXC{$\stackrel{\Pi_{1}}{\vdots}$}\noLine
\UIC{$\varphi_{6}$}
\BIC{$\varphi_{3}$}
\AXC{$\varphi_{3} \to \varphi_{8}$\lbl{A1}}
\BIC{$\varphi_{8}$}
\BIC{$\mathrm{Cauchy}(a) \to \varepsilon \in \mathbb{R} \to 0 < \varepsilon \to \varepsilon \in \mathbb{R} \to 0 < \varepsilon \to \varphi_{1}$}
\end{ptree}
\noindent$\Pi_{3}$:
\begin{ptree}
\AXC{$\mathrm{Cauchy}(a) \to \varepsilon \in \mathbb{R} \to \varepsilon \in \mathbb{R}$\lbl{A1}}
\AXC{$\stackrel{\Pi_{2}}{\vdots}$}\noLine
\UIC{$\mathrm{Cauchy}(a) \to \varepsilon \in \mathbb{R} \to 0 < \varepsilon \to \varepsilon \in \mathbb{R} \to 0 < \varepsilon \to \varphi_{1}$}
\AXC{$\varphi_{7}$\lbl{A1}}
\BIC{$(\mathrm{Cauchy}(a) \to \varepsilon \in \mathbb{R} \to \varepsilon \in \mathbb{R}) \to \varphi_{4}$}
\BIC{$\varphi_{4}$}
\end{ptree}
\noindent 主木:
\begin{ptree}
\AXC{$\mathrm{Cauchy}(a) \to \varepsilon \in \mathbb{R} \to 0 < \varepsilon \to 0 < \varepsilon$\lbl{A1}}
\AXC{$\stackrel{\Pi_{3}}{\vdots}$}\noLine
\UIC{$\varphi_{4}$}
\BIC{$\mathrm{Cauchy}(a) \to \varepsilon \in \mathbb{R} \to 0 < \varepsilon \to \varphi_{1}$}
\UIC{$\forall \varepsilon\,(\mathrm{Cauchy}(a) \to \varepsilon \in \mathbb{R} \to 0 < \varepsilon \to \varphi_{1})$}
\UIC{$\forall a\,\forall \varepsilon\,(\mathrm{Cauchy}(a) \to \varepsilon \in \mathbb{R} \to 0 < \varepsilon \to \varphi_{1})$}
\end{ptree}
\begin{abbrevlist}
\item $\varphi_{1} :\equiv \exists N \in \mathbb{N}\,\forall m\,\forall n\,(m \in \mathbb{N} \to n \in \mathbb{N} \to N \leq m \to N \leq n \to \lvert a_{m} - a_{n} \rvert < \varepsilon)$
\item $\varphi_{2} :\equiv \forall a\,(\mathrm{Cauchy}(a) \leftrightarrow \mathrm{Seq}(a) \land \forall \varepsilon \in \mathbb{R}\,(0 < \varepsilon \to \varphi_{1}))$
\item $\varphi_{3} :\equiv \mathrm{Cauchy}(a) \to \varepsilon \in \mathbb{R} \to 0 < \varepsilon \to \forall \varepsilon \in \mathbb{R}\,(0 < \varepsilon \to \varphi_{1})$
\item $\varphi_{4} :\equiv (\mathrm{Cauchy}(a) \to \varepsilon \in \mathbb{R} \to 0 < \varepsilon \to 0 < \varepsilon) \to \mathrm{Cauchy}(a) \to \varepsilon \in \mathbb{R} \to 0 < \varepsilon \to \varphi_{1}$
\item $\varphi_{5} :\equiv \varphi_{2} \to (\mathrm{Cauchy}(a) \leftrightarrow \mathrm{Seq}(a) \land \forall \varepsilon \in \mathbb{R}\,(0 < \varepsilon \to \varphi_{1}))$
\item $\varphi_{6} :\equiv (\mathrm{Cauchy}(a) \leftrightarrow \mathrm{Seq}(a) \land \forall \varepsilon \in \mathbb{R}\,(0 < \varepsilon \to \varphi_{1})) \to \varphi_{3}$
\item $\varphi_{7} :\equiv (\mathrm{Cauchy}(a) \to \varepsilon \in \mathbb{R} \to 0 < \varepsilon \to \varepsilon \in \mathbb{R} \to 0 < \varepsilon \to \varphi_{1}) \to (\mathrm{Cauchy}(a) \to \varepsilon \in \mathbb{R} \to \varepsilon \in \mathbb{R}) \to \varphi_{4}$
\item $\varphi_{8} :\equiv (\forall \varepsilon \in \mathbb{R}\,(0 < \varepsilon \to \varphi_{1}) \to \varepsilon \in \mathbb{R} \to 0 < \varepsilon \to \varphi_{1}) \to \mathrm{Cauchy}(a) \to \varepsilon \in \mathbb{R} \to 0 < \varepsilon \to \varepsilon \in \mathbb{R} \to 0 < \varepsilon \to \varphi_{1}$
\end{abbrevlist}
\end{fullproof}

% 自動生成 (prover/build.js): ErrChain  ノード数 163
\begin{fullproof}{[C27]}{$\forall x\,\forall y\,\forall z\,\forall e\,(x \in \mathbb{R} \to y \in \mathbb{R} \to z \in \mathbb{R} \to e \in \mathbb{R} \to \lvert x - y \rvert < e \to \lvert y - z \rvert < e \to \lvert x - z \rvert < e + e)$}{163}\phantomsection\label{tree:ErrChain}
\noindent$\Pi_{1}$:
\begin{ptree}
\AXC{$\forall x\,\forall y\,(x \in \mathbb{R} \to y \in \mathbb{R} \to x - y \in \mathbb{R})$\lbl{[F1]}}
\AXC{$\varphi_{41}$\lbl{A2}}
\BIC{$\forall y\,(v_{738} \in \mathbb{R} \to y \in \mathbb{R} \to v_{738} - y \in \mathbb{R})$}
\AXC{$\varphi_{45}$\lbl{A2}}
\BIC{$v_{738} \in \mathbb{R} \to v_{739} \in \mathbb{R} \to v_{738} - v_{739} \in \mathbb{R}$}
\UIC{$\forall v_{739}\,(v_{738} \in \mathbb{R} \to v_{739} \in \mathbb{R} \to v_{738} - v_{739} \in \mathbb{R})$}
\UIC{$\varphi_{4}$}
\AXC{$\varphi_{4} \to \forall v_{739}\,(y \in \mathbb{R} \to v_{739} \in \mathbb{R} \to y - v_{739} \in \mathbb{R})$\lbl{A2}}
\BIC{$\forall v_{739}\,(y \in \mathbb{R} \to v_{739} \in \mathbb{R} \to y - v_{739} \in \mathbb{R})$}
\end{ptree}
\noindent$\Pi_{2}$:
\begin{ptree}
\AXC{$x \in \mathbb{R} \to \varphi_{31}$\lbl{A1}}
\AXC{$x \in \mathbb{R} \to \varphi_{32}$\lbl{A1}}
\AXC{$\stackrel{\Pi_{1}}{\vdots}$}\noLine
\UIC{$\forall v_{739}\,(y \in \mathbb{R} \to v_{739} \in \mathbb{R} \to y - v_{739} \in \mathbb{R})$}
\AXC{$\varphi_{39}$\lbl{A2}}
\BIC{$y \in \mathbb{R} \to z \in \mathbb{R} \to y - z \in \mathbb{R}$}
\AXC{$\varphi_{66}$\lbl{A1}}
\BIC{$(x \in \mathbb{R} \to \varphi_{32}) \to (x \in \mathbb{R} \to \varphi_{31}) \to x \in \mathbb{R} \to y \in \mathbb{R} \to \varphi_{21}$}
\BIC{$(x \in \mathbb{R} \to \varphi_{31}) \to x \in \mathbb{R} \to y \in \mathbb{R} \to \varphi_{21}$}
\BIC{$x \in \mathbb{R} \to y \in \mathbb{R} \to \varphi_{21}$}
\end{ptree}
\noindent$\Pi_{3}$:
\begin{ptree}
\AXC{$\forall x \in \mathbb{R}\,(\lvert x \rvert \in \mathbb{R})$\lbl{[A5]}}
\AXC{$\forall x \in \mathbb{R}\,(\lvert x \rvert \in \mathbb{R}) \to y - z \in \mathbb{R} \to \lvert y - z \rvert \in \mathbb{R}$\lbl{A2}}
\BIC{$y - z \in \mathbb{R} \to \lvert y - z \rvert \in \mathbb{R}$}
\AXC{$\varphi_{62}$\lbl{A1}}
\BIC{$(x \in \mathbb{R} \to y \in \mathbb{R} \to \varphi_{21}) \to x \in \mathbb{R} \to y \in \mathbb{R} \to \varphi_{27}$}
\end{ptree}
\noindent$\Pi_{4}$:
\begin{ptree}
\AXC{$x \in \mathbb{R} \to \varphi_{33}$\lbl{A1}}
\AXC{$\forall x\,\forall y\,(x \in \mathbb{R} \to y \in \mathbb{R} \to x - y \in \mathbb{R})$\lbl{[F1]}}
\AXC{$\varphi_{37}$\lbl{A2}}
\BIC{$\forall y\,(x \in \mathbb{R} \to y \in \mathbb{R} \to x - y \in \mathbb{R})$}
\AXC{$\varphi_{34}$\lbl{A2}}
\BIC{$x \in \mathbb{R} \to y \in \mathbb{R} \to x - y \in \mathbb{R}$}
\AXC{$\varphi_{67}$\lbl{A1}}
\BIC{$(x \in \mathbb{R} \to \varphi_{33}) \to (x \in \mathbb{R} \to \varphi_{32}) \to x \in \mathbb{R} \to y \in \mathbb{R} \to \varphi_{22}$}
\BIC{$(x \in \mathbb{R} \to \varphi_{32}) \to x \in \mathbb{R} \to y \in \mathbb{R} \to \varphi_{22}$}
\end{ptree}
\noindent$\Pi_{5}$:
\begin{ptree}
\AXC{$x \in \mathbb{R} \to \varphi_{32}$\lbl{A1}}
\AXC{$\stackrel{\Pi_{4}}{\vdots}$}\noLine
\UIC{$(x \in \mathbb{R} \to \varphi_{32}) \to x \in \mathbb{R} \to y \in \mathbb{R} \to \varphi_{22}$}
\BIC{$x \in \mathbb{R} \to y \in \mathbb{R} \to \varphi_{22}$}
\end{ptree}
\noindent$\Pi_{6}$:
\begin{ptree}
\AXC{$\forall x \in \mathbb{R}\,(\lvert x \rvert \in \mathbb{R})$\lbl{[A5]}}
\AXC{$\forall x \in \mathbb{R}\,(\lvert x \rvert \in \mathbb{R}) \to x - y \in \mathbb{R} \to \lvert x - y \rvert \in \mathbb{R}$\lbl{A2}}
\BIC{$x - y \in \mathbb{R} \to \lvert x - y \rvert \in \mathbb{R}$}
\AXC{$\varphi_{63}$\lbl{A1}}
\BIC{$(x \in \mathbb{R} \to y \in \mathbb{R} \to \varphi_{22}) \to x \in \mathbb{R} \to y \in \mathbb{R} \to \varphi_{28}$}
\end{ptree}
\noindent$\Pi_{7}$:
\begin{ptree}
\AXC{$\stackrel{\Pi_{5}}{\vdots}$}\noLine
\UIC{$x \in \mathbb{R} \to y \in \mathbb{R} \to \varphi_{22}$}
\AXC{$\stackrel{\Pi_{6}}{\vdots}$}\noLine
\UIC{$(x \in \mathbb{R} \to y \in \mathbb{R} \to \varphi_{22}) \to x \in \mathbb{R} \to y \in \mathbb{R} \to \varphi_{28}$}
\BIC{$x \in \mathbb{R} \to y \in \mathbb{R} \to \varphi_{28}$}
\end{ptree}
\noindent$\Pi_{8}$:
\begin{ptree}
\AXC{$\varphi_{40}$\lbl{[E7]}}
\AXC{$\varphi_{40} \to \varphi_{44}$\lbl{A2}}
\BIC{$\varphi_{44}$}
\AXC{$\varphi_{44} \to \varphi_{47}$\lbl{A2}}
\BIC{$\varphi_{47}$}
\AXC{$\varphi_{47} \to \varphi_{49}$\lbl{A2}}
\BIC{$\varphi_{49}$}
\AXC{$\varphi_{59}$\lbl{A2}}
\BIC{$v_{740} \in \mathbb{R} \to v_{741} \in \mathbb{R} \to v_{742} \in \mathbb{R} \to v_{743} \in \mathbb{R} \to \varphi_{3}$}
\UIC{$\varphi_{50}$}
\UIC{$\forall v_{742}\,(\varphi_{50})$}
\UIC{$\forall v_{741}\,\forall v_{742}\,(\varphi_{50})$}
\UIC{$\forall v_{740}\,\forall v_{741}\,\forall v_{742}\,(\varphi_{50})$}
\AXC{$\forall v_{740}\,\forall v_{741}\,\forall v_{742}\,(\varphi_{50}) \to \varphi_{53}$\lbl{A2}}
\BIC{$\varphi_{53}$}
\end{ptree}
\noindent$\Pi_{9}$:
\begin{ptree}
\AXC{$x \in \mathbb{R} \to \varphi_{30}$\lbl{A1}}
\AXC{$\stackrel{\Pi_{7}}{\vdots}$}\noLine
\UIC{$x \in \mathbb{R} \to y \in \mathbb{R} \to \varphi_{28}$}
\AXC{$\stackrel{\Pi_{8}}{\vdots}$}\noLine
\UIC{$\varphi_{53}$}
\AXC{$\varphi_{53} \to \varphi_{52}$\lbl{A2}}
\BIC{$\varphi_{52}$}
\AXC{$\varphi_{52} \to \varphi_{54}$\lbl{A2}}
\BIC{$\varphi_{54}$}
\AXC{$\varphi_{64}$\lbl{A2}}
\BIC{$\lvert x - y \rvert \in \mathbb{R} \to e \in \mathbb{R} \to \lvert y - z \rvert \in \mathbb{R} \to e \in \mathbb{R} \to \varphi_{13}$}
\AXC{$\varphi_{71}$\lbl{A1}}
\BIC{$(x \in \mathbb{R} \to y \in \mathbb{R} \to \varphi_{28}) \to (x \in \mathbb{R} \to \varphi_{30}) \to \varphi_{70}$}
\BIC{$(x \in \mathbb{R} \to \varphi_{30}) \to \varphi_{70}$}
\BIC{$\varphi_{70}$}
\end{ptree}
\noindent$\Pi_{10}$:
\begin{ptree}
\AXC{$x \in \mathbb{R} \to \varphi_{30}$\lbl{A1}}
\AXC{$\stackrel{\Pi_{2}}{\vdots}$}\noLine
\UIC{$x \in \mathbb{R} \to y \in \mathbb{R} \to \varphi_{21}$}
\AXC{$\stackrel{\Pi_{3}}{\vdots}$}\noLine
\UIC{$(x \in \mathbb{R} \to y \in \mathbb{R} \to \varphi_{21}) \to x \in \mathbb{R} \to y \in \mathbb{R} \to \varphi_{27}$}
\BIC{$x \in \mathbb{R} \to y \in \mathbb{R} \to \varphi_{27}$}
\AXC{$\stackrel{\Pi_{9}}{\vdots}$}\noLine
\UIC{$\varphi_{70}$}
\BIC{$(x \in \mathbb{R} \to \varphi_{30}) \to (x \in \mathbb{R} \to y \in \mathbb{R} \to \varphi_{24}) \to \varphi_{57}$}
\BIC{$(x \in \mathbb{R} \to y \in \mathbb{R} \to \varphi_{24}) \to \varphi_{57}$}
\end{ptree}
\noindent$\Pi_{11}$:
\begin{ptree}
\AXC{$z_{1} = w \to \varphi_{16}$\lbl{A6}}
\UIC{$\forall w\,(z_{1} = w \to \varphi_{16})$}
\UIC{$\forall z_{1}\,\forall w\,(z_{1} = w \to \varphi_{16})$}
\AXC{$\forall z_{1}\,\forall w\,(z_{1} = w \to \varphi_{16}) \to \varphi_{48}$\lbl{A2}}
\BIC{$\varphi_{48}$}
\AXC{$\varphi_{48} \to x - y + (y - z) = x - z \to \varphi_{38}$\lbl{A2}}
\BIC{$x - y + (y - z) = x - z \to \varphi_{38}$}
\end{ptree}
\noindent$\Pi_{12}$:
\begin{ptree}
\AXC{$\forall a\,\forall b\,(a \in \mathbb{R} \to b \in \mathbb{R} \to \lvert a + b \rvert \leq \lvert a \rvert + \lvert b \rvert)$\lbl{命題1.2}}
\AXC{$\varphi_{51}$\lbl{A2}}
\BIC{$\varphi_{10}$}
\AXC{$\varphi_{10} \to \varphi_{26}$\lbl{A2}}
\BIC{$\varphi_{26}$}
\AXC{$\varphi_{26} \to (x \in \mathbb{R} \to y \in \mathbb{R} \to \varphi_{22}) \to \varphi_{60}$\lbl{A1}}
\BIC{$(x \in \mathbb{R} \to y \in \mathbb{R} \to \varphi_{22}) \to \varphi_{60}$}
\end{ptree}
\noindent$\Pi_{13}$:
\begin{ptree}
\AXC{$\stackrel{\Pi_{2}}{\vdots}$}\noLine
\UIC{$x \in \mathbb{R} \to y \in \mathbb{R} \to \varphi_{21}$}
\AXC{$\stackrel{\Pi_{5}}{\vdots}$}\noLine
\UIC{$x \in \mathbb{R} \to y \in \mathbb{R} \to \varphi_{22}$}
\AXC{$\stackrel{\Pi_{12}}{\vdots}$}\noLine
\UIC{$(x \in \mathbb{R} \to y \in \mathbb{R} \to \varphi_{22}) \to \varphi_{60}$}
\BIC{$\varphi_{60}$}
\BIC{$x \in \mathbb{R} \to y \in \mathbb{R} \to z \in \mathbb{R} \to e \in \mathbb{R} \to \lvert x - y \rvert < e \to \varphi_{7}$}
\end{ptree}
\noindent$\Pi_{14}$:
\begin{ptree}
\AXC{$x = x$\lbl{A5}}
\AXC{$x = y \to x = x \to y = x$\lbl{A6}}
\AXC{$(x = y \to x = x \to y = x) \to x = x \to x = y \to y = x$\lbl{A1}}
\BIC{$x = x \to x = y \to y = x$}
\BIC{$x = y \to y = x$}
\UIC{$\forall y\,(x = y \to y = x)$}
\UIC{$\forall x\,\forall y\,(x = y \to y = x)$}
\end{ptree}
\noindent$\Pi_{15}$:
\begin{ptree}
\AXC{$\varphi_{9}$\lbl{[F16]}}
\AXC{$\varphi_{9} \to \varphi_{18}$\lbl{A2}}
\BIC{$\varphi_{18}$}
\AXC{$\varphi_{18} \to \varphi_{25}$\lbl{A2}}
\BIC{$\varphi_{25}$}
\AXC{$\varphi_{25} \to v_{735} \in \mathbb{R} \to v_{736} \in \mathbb{R} \to \varphi_{2}$\lbl{A2}}
\BIC{$v_{735} \in \mathbb{R} \to v_{736} \in \mathbb{R} \to \varphi_{2}$}
\UIC{$\forall v_{737}\,(v_{735} \in \mathbb{R} \to v_{736} \in \mathbb{R} \to \varphi_{2})$}
\UIC{$\forall v_{736}\,\forall v_{737}\,(v_{735} \in \mathbb{R} \to v_{736} \in \mathbb{R} \to \varphi_{2})$}
\UIC{$\forall v_{735}\,\forall v_{736}\,\forall v_{737}\,(v_{735} \in \mathbb{R} \to v_{736} \in \mathbb{R} \to \varphi_{2})$}
\AXC{$\varphi_{56}$\lbl{A2}}
\BIC{$\varphi_{36}$}
\AXC{$\varphi_{36} \to \varphi_{19}$\lbl{A2}}
\BIC{$\varphi_{19}$}
\end{ptree}
\noindent$\Pi_{16}$:
\begin{ptree}
\AXC{$x \in \mathbb{R} \to \varphi_{31}$\lbl{A1}}
\AXC{$x \in \mathbb{R} \to \varphi_{33}$\lbl{A1}}
\AXC{$\stackrel{\Pi_{15}}{\vdots}$}\noLine
\UIC{$\varphi_{19}$}
\AXC{$\varphi_{19} \to x \in \mathbb{R} \to z \in \mathbb{R} \to y \in \mathbb{R} \to x - z = x - y + (y - z)$\lbl{A2}}
\BIC{$x \in \mathbb{R} \to z \in \mathbb{R} \to y \in \mathbb{R} \to x - z = x - y + (y - z)$}
\AXC{$\varphi_{68}$\lbl{A1}}
\BIC{$(x \in \mathbb{R} \to \varphi_{33}) \to \varphi_{65}$}
\BIC{$\varphi_{65}$}
\BIC{$(x \in \mathbb{R} \to \varphi_{32}) \to x \in \mathbb{R} \to y \in \mathbb{R} \to z \in \mathbb{R} \to \varphi_{14}$}
\end{ptree}
\noindent$\Pi_{17}$:
\begin{ptree}
\AXC{$x \in \mathbb{R} \to \varphi_{32}$\lbl{A1}}
\AXC{$\stackrel{\Pi_{16}}{\vdots}$}\noLine
\UIC{$(x \in \mathbb{R} \to \varphi_{32}) \to x \in \mathbb{R} \to y \in \mathbb{R} \to z \in \mathbb{R} \to \varphi_{14}$}
\BIC{$x \in \mathbb{R} \to y \in \mathbb{R} \to z \in \mathbb{R} \to \varphi_{14}$}
\AXC{$(x \in \mathbb{R} \to y \in \mathbb{R} \to z \in \mathbb{R} \to \varphi_{14}) \to \varphi_{55}$\lbl{A1}}
\BIC{$\varphi_{55}$}
\end{ptree}
\noindent$\Pi_{18}$:
\begin{ptree}
\AXC{$\stackrel{\Pi_{14}}{\vdots}$}\noLine
\UIC{$\forall x\,\forall y\,(x = y \to y = x)$}
\AXC{$\forall x\,\forall y\,(x = y \to y = x) \to \forall y\,(x - z = y \to y = x - z)$\lbl{A2}}
\BIC{$\forall y\,(x - z = y \to y = x - z)$}
\AXC{$\varphi_{12}$\lbl{A2}}
\BIC{$x - z = x - y + (y - z) \to x - y + (y - z) = x - z$}
\AXC{$\stackrel{\Pi_{17}}{\vdots}$}\noLine
\UIC{$\varphi_{55}$}
\BIC{$x \in \mathbb{R} \to y \in \mathbb{R} \to z \in \mathbb{R} \to \varphi_{15}$}
\end{ptree}
\noindent$\Pi_{19}$:
\begin{ptree}
\AXC{$\stackrel{\Pi_{18}}{\vdots}$}\noLine
\UIC{$x \in \mathbb{R} \to y \in \mathbb{R} \to z \in \mathbb{R} \to \varphi_{15}$}
\AXC{$(x \in \mathbb{R} \to y \in \mathbb{R} \to z \in \mathbb{R} \to \varphi_{15}) \to \varphi_{69}$\lbl{A1}}
\BIC{$\varphi_{69}$}
\end{ptree}
\noindent$\Pi_{20}$:
\begin{ptree}
\AXC{$\stackrel{\Pi_{11}}{\vdots}$}\noLine
\UIC{$x - y + (y - z) = x - z \to \varphi_{38}$}
\AXC{$\stackrel{\Pi_{13}}{\vdots}$}\noLine
\UIC{$x \in \mathbb{R} \to y \in \mathbb{R} \to z \in \mathbb{R} \to e \in \mathbb{R} \to \lvert x - y \rvert < e \to \varphi_{7}$}
\AXC{$\stackrel{\Pi_{19}}{\vdots}$}\noLine
\UIC{$\varphi_{69}$}
\BIC{$(x - y + (y - z) = x - z \to \varphi_{38}) \to \varphi_{61}$}
\BIC{$\varphi_{61}$}
\end{ptree}
\noindent$\Pi_{21}$:
\begin{ptree}
\AXC{$z_{1} = w \to \varphi_{8}$\lbl{A6}}
\UIC{$\forall w\,(z_{1} = w \to \varphi_{8})$}
\UIC{$\forall z_{1}\,\forall w\,(z_{1} = w \to \varphi_{8})$}
\AXC{$\forall z_{1}\,\forall w\,(z_{1} = w \to \varphi_{8}) \to \varphi_{46}$\lbl{A2}}
\BIC{$\varphi_{46}$}
\AXC{$\varphi_{46} \to x - y + (y - z) = x - z \to \varphi_{35}$\lbl{A2}}
\BIC{$x - y + (y - z) = x - z \to \varphi_{35}$}
\AXC{$\stackrel{\Pi_{20}}{\vdots}$}\noLine
\UIC{$\varphi_{61}$}
\BIC{$x \in \mathbb{R} \to y \in \mathbb{R} \to z \in \mathbb{R} \to e \in \mathbb{R} \to \varphi_{17}$}
\end{ptree}
\noindent$\Pi_{22}$:
\begin{ptree}
\AXC{$\forall x\,\forall y\,\forall z\,(x \leq y \to y < z \to x < z)$\lbl{[O4]}}
\AXC{$\varphi_{6}$\lbl{A2}}
\BIC{$\forall y\,\forall z\,(v_{744} \leq y \to y < z \to v_{744} < z)$}
\AXC{$\varphi_{11}$\lbl{A2}}
\BIC{$\forall z\,(v_{744} \leq v_{745} \to v_{745} < z \to v_{744} < z)$}
\AXC{$\varphi_{20}$\lbl{A2}}
\BIC{$v_{744} \leq v_{745} \to v_{745} < v_{746} \to v_{744} < v_{746}$}
\UIC{$\forall v_{746}\,(v_{744} \leq v_{745} \to v_{745} < v_{746} \to v_{744} < v_{746})$}
\UIC{$\varphi_{1}$}
\UIC{$\forall v_{744}\,(\varphi_{1})$}
\AXC{$\forall v_{744}\,(\varphi_{1}) \to \varphi_{5}$\lbl{A2}}
\BIC{$\varphi_{5}$}
\AXC{$\varphi_{5} \to \varphi_{43}$\lbl{A2}}
\BIC{$\varphi_{43}$}
\end{ptree}
\noindent$\Pi_{23}$:
\begin{ptree}
\AXC{$\stackrel{\Pi_{21}}{\vdots}$}\noLine
\UIC{$x \in \mathbb{R} \to y \in \mathbb{R} \to z \in \mathbb{R} \to e \in \mathbb{R} \to \varphi_{17}$}
\AXC{$\stackrel{\Pi_{22}}{\vdots}$}\noLine
\UIC{$\varphi_{43}$}
\AXC{$\varphi_{43} \to \varphi_{42}$\lbl{A2}}
\BIC{$\varphi_{42}$}
\AXC{$\varphi_{42} \to (x \in \mathbb{R} \to y \in \mathbb{R} \to z \in \mathbb{R} \to e \in \mathbb{R} \to \varphi_{17}) \to \varphi_{58}$\lbl{A1}}
\BIC{$(x \in \mathbb{R} \to y \in \mathbb{R} \to z \in \mathbb{R} \to e \in \mathbb{R} \to \varphi_{17}) \to \varphi_{58}$}
\BIC{$\varphi_{58}$}
\end{ptree}
\noindent 主木:
\begin{ptree}
\AXC{$x \in \mathbb{R} \to y \in \mathbb{R} \to \varphi_{23}$\lbl{A1}}
\AXC{$x \in \mathbb{R} \to y \in \mathbb{R} \to \varphi_{24}$\lbl{A1}}
\AXC{$\stackrel{\Pi_{10}}{\vdots}$}\noLine
\UIC{$(x \in \mathbb{R} \to y \in \mathbb{R} \to \varphi_{24}) \to \varphi_{57}$}
\BIC{$\varphi_{57}$}
\BIC{$x \in \mathbb{R} \to y \in \mathbb{R} \to z \in \mathbb{R} \to e \in \mathbb{R} \to \varphi_{13}$}
\AXC{$\stackrel{\Pi_{23}}{\vdots}$}\noLine
\UIC{$\varphi_{58}$}
\BIC{$x \in \mathbb{R} \to y \in \mathbb{R} \to \varphi_{29}$}
\UIC{$\forall e\,(x \in \mathbb{R} \to y \in \mathbb{R} \to \varphi_{29})$}
\UIC{$\forall z\,\forall e\,(x \in \mathbb{R} \to y \in \mathbb{R} \to \varphi_{29})$}
\UIC{$\forall y\,\forall z\,\forall e\,(x \in \mathbb{R} \to y \in \mathbb{R} \to \varphi_{29})$}
\UIC{$\forall x\,\forall y\,\forall z\,\forall e\,(x \in \mathbb{R} \to y \in \mathbb{R} \to \varphi_{29})$}
\end{ptree}
\begin{abbrevlist}
\item $\varphi_{1} :\equiv \forall v_{745}\,\forall v_{746}\,(v_{744} \leq v_{745} \to v_{745} < v_{746} \to v_{744} < v_{746})$
\item $\varphi_{2} :\equiv v_{737} \in \mathbb{R} \to v_{735} - v_{736} = v_{735} - v_{737} + (v_{737} - v_{736})$
\item $\varphi_{3} :\equiv v_{740} < v_{741} \to v_{742} < v_{743} \to v_{740} + v_{742} < v_{741} + v_{743}$
\item $\varphi_{4} :\equiv \forall v_{738}\,\forall v_{739}\,(v_{738} \in \mathbb{R} \to v_{739} \in \mathbb{R} \to v_{738} - v_{739} \in \mathbb{R})$
\item $\varphi_{5} :\equiv \forall v_{745}\,\forall v_{746}\,(\lvert x - z \rvert \leq v_{745} \to v_{745} < v_{746} \to \lvert x - z \rvert < v_{746})$
\item $\varphi_{6} :\equiv \forall x\,\forall y\,\forall z\,(x \leq y \to y < z \to x < z) \to \forall y\,\forall z\,(v_{744} \leq y \to y < z \to v_{744} < z)$
\item $\varphi_{7} :\equiv \lvert y - z \rvert < e \to \lvert x - y + (y - z) \rvert \leq \lvert x - y \rvert + \lvert y - z \rvert$
\item $\varphi_{8} :\equiv \lvert z_{1} \rvert = \lvert x - y \rvert + \lvert y - z \rvert \to \lvert w \rvert = \lvert x - y \rvert + \lvert y - z \rvert$
\item $\varphi_{9} :\equiv \forall a\,\forall b\,\forall x\,(a \in \mathbb{R} \to b \in \mathbb{R} \to x \in \mathbb{R} \to a - b = a - x + (x - b))$
\item $\varphi_{10} :\equiv \forall b\,(x - y \in \mathbb{R} \to b \in \mathbb{R} \to \lvert x - y + b \rvert \leq \lvert x - y \rvert + \lvert b \rvert)$
\item $\varphi_{11} :\equiv \forall y\,\forall z\,(v_{744} \leq y \to y < z \to v_{744} < z) \to \forall z\,(v_{744} \leq v_{745} \to v_{745} < z \to v_{744} < z)$
\item $\varphi_{12} :\equiv \forall y\,(x - z = y \to y = x - z) \to x - z = x - y + (y - z) \to x - y + (y - z) = x - z$
\item $\varphi_{13} :\equiv \lvert x - y \rvert < e \to \lvert y - z \rvert < e \to \lvert x - y \rvert + \lvert y - z \rvert < e + e$
\item $\varphi_{14} :\equiv e \in \mathbb{R} \to \lvert x - y \rvert < e \to \lvert y - z \rvert < e \to x - z = x - y + (y - z)$
\item $\varphi_{15} :\equiv e \in \mathbb{R} \to \lvert x - y \rvert < e \to \lvert y - z \rvert < e \to x - y + (y - z) = x - z$
\item $\varphi_{16} :\equiv \lvert z_{1} \rvert < \lvert x - y \rvert + \lvert y - z \rvert \to \lvert w \rvert < \lvert x - y \rvert + \lvert y - z \rvert$
\item $\varphi_{17} :\equiv \lvert x - y \rvert < e \to \lvert y - z \rvert < e \to \lvert x - z \rvert \leq \lvert x - y \rvert + \lvert y - z \rvert$
\item $\varphi_{18} :\equiv \forall b\,\forall x\,(v_{735} \in \mathbb{R} \to b \in \mathbb{R} \to x \in \mathbb{R} \to v_{735} - b = v_{735} - x + (x - b))$
\item $\varphi_{19} :\equiv \forall v_{737}\,(x \in \mathbb{R} \to z \in \mathbb{R} \to v_{737} \in \mathbb{R} \to x - z = x - v_{737} + (v_{737} - z))$
\item $\varphi_{20} :\equiv \forall z\,(v_{744} \leq v_{745} \to v_{745} < z \to v_{744} < z) \to v_{744} \leq v_{745} \to v_{745} < v_{746} \to v_{744} < v_{746}$
\item $\varphi_{21} :\equiv z \in \mathbb{R} \to e \in \mathbb{R} \to \lvert x - y \rvert < e \to \lvert y - z \rvert < e \to y - z \in \mathbb{R}$
\item $\varphi_{22} :\equiv z \in \mathbb{R} \to e \in \mathbb{R} \to \lvert x - y \rvert < e \to \lvert y - z \rvert < e \to x - y \in \mathbb{R}$
\item $\varphi_{23} :\equiv z \in \mathbb{R} \to e \in \mathbb{R} \to \lvert x - y \rvert < e \to \lvert y - z \rvert < e \to \lvert y - z \rvert < e$
\item $\varphi_{24} :\equiv z \in \mathbb{R} \to e \in \mathbb{R} \to \lvert x - y \rvert < e \to \lvert y - z \rvert < e \to \lvert x - y \rvert < e$
\item $\varphi_{25} :\equiv \forall x\,(v_{735} \in \mathbb{R} \to v_{736} \in \mathbb{R} \to x \in \mathbb{R} \to v_{735} - v_{736} = v_{735} - x + (x - v_{736}))$
\item $\varphi_{26} :\equiv x - y \in \mathbb{R} \to y - z \in \mathbb{R} \to \lvert x - y + (y - z) \rvert \leq \lvert x - y \rvert + \lvert y - z \rvert$
\item $\varphi_{27} :\equiv z \in \mathbb{R} \to e \in \mathbb{R} \to \lvert x - y \rvert < e \to \lvert y - z \rvert < e \to \lvert y - z \rvert \in \mathbb{R}$
\item $\varphi_{28} :\equiv z \in \mathbb{R} \to e \in \mathbb{R} \to \lvert x - y \rvert < e \to \lvert y - z \rvert < e \to \lvert x - y \rvert \in \mathbb{R}$
\item $\varphi_{29} :\equiv z \in \mathbb{R} \to e \in \mathbb{R} \to \lvert x - y \rvert < e \to \lvert y - z \rvert < e \to \lvert x - z \rvert < e + e$
\item $\varphi_{30} :\equiv y \in \mathbb{R} \to z \in \mathbb{R} \to e \in \mathbb{R} \to \lvert x - y \rvert < e \to \lvert y - z \rvert < e \to e \in \mathbb{R}$
\item $\varphi_{31} :\equiv y \in \mathbb{R} \to z \in \mathbb{R} \to e \in \mathbb{R} \to \lvert x - y \rvert < e \to \lvert y - z \rvert < e \to z \in \mathbb{R}$
\item $\varphi_{32} :\equiv y \in \mathbb{R} \to z \in \mathbb{R} \to e \in \mathbb{R} \to \lvert x - y \rvert < e \to \lvert y - z \rvert < e \to y \in \mathbb{R}$
\item $\varphi_{33} :\equiv y \in \mathbb{R} \to z \in \mathbb{R} \to e \in \mathbb{R} \to \lvert x - y \rvert < e \to \lvert y - z \rvert < e \to x \in \mathbb{R}$
\item $\varphi_{34} :\equiv \forall y\,(x \in \mathbb{R} \to y \in \mathbb{R} \to x - y \in \mathbb{R}) \to x \in \mathbb{R} \to y \in \mathbb{R} \to x - y \in \mathbb{R}$
\item $\varphi_{35} :\equiv \lvert x - y + (y - z) \rvert = \lvert x - y \rvert + \lvert y - z \rvert \to \lvert x - z \rvert = \lvert x - y \rvert + \lvert y - z \rvert$
\item $\varphi_{36} :\equiv \forall v_{736}\,\forall v_{737}\,(x \in \mathbb{R} \to v_{736} \in \mathbb{R} \to v_{737} \in \mathbb{R} \to x - v_{736} = x - v_{737} + (v_{737} - v_{736}))$
\item $\varphi_{37} :\equiv \forall x\,\forall y\,(x \in \mathbb{R} \to y \in \mathbb{R} \to x - y \in \mathbb{R}) \to \forall y\,(x \in \mathbb{R} \to y \in \mathbb{R} \to x - y \in \mathbb{R})$
\item $\varphi_{38} :\equiv \lvert x - y + (y - z) \rvert < \lvert x - y \rvert + \lvert y - z \rvert \to \lvert x - z \rvert < \lvert x - y \rvert + \lvert y - z \rvert$
\item $\varphi_{39} :\equiv \forall v_{739}\,(y \in \mathbb{R} \to v_{739} \in \mathbb{R} \to y - v_{739} \in \mathbb{R}) \to y \in \mathbb{R} \to z \in \mathbb{R} \to y - z \in \mathbb{R}$
\item $\varphi_{40} :\equiv \forall x\,\forall y\,\forall u\,\forall v\,(x \in \mathbb{R} \to y \in \mathbb{R} \to u \in \mathbb{R} \to v \in \mathbb{R} \to x < y \to u < v \to x + u < y + v)$
\item $\varphi_{41} :\equiv \forall x\,\forall y\,(x \in \mathbb{R} \to y \in \mathbb{R} \to x - y \in \mathbb{R}) \to \forall y\,(v_{738} \in \mathbb{R} \to y \in \mathbb{R} \to v_{738} - y \in \mathbb{R})$
\item $\varphi_{42} :\equiv \lvert x - z \rvert \leq \lvert x - y \rvert + \lvert y - z \rvert \to \lvert x - y \rvert + \lvert y - z \rvert < e + e \to \lvert x - z \rvert < e + e$
\item $\varphi_{43} :\equiv \forall v_{746}\,(\lvert x - z \rvert \leq \lvert x - y \rvert + \lvert y - z \rvert \to \lvert x - y \rvert + \lvert y - z \rvert < v_{746} \to \lvert x - z \rvert < v_{746})$
\item $\varphi_{44} :\equiv \forall y\,\forall u\,\forall v\,(v_{740} \in \mathbb{R} \to y \in \mathbb{R} \to u \in \mathbb{R} \to v \in \mathbb{R} \to v_{740} < y \to u < v \to v_{740} + u < y + v)$
\item $\varphi_{45} :\equiv \forall y\,(v_{738} \in \mathbb{R} \to y \in \mathbb{R} \to v_{738} - y \in \mathbb{R}) \to v_{738} \in \mathbb{R} \to v_{739} \in \mathbb{R} \to v_{738} - v_{739} \in \mathbb{R}$
\item $\varphi_{46} :\equiv \forall w\,(x - y + (y - z) = w \to \lvert x - y + (y - z) \rvert = \lvert x - y \rvert + \lvert y - z \rvert \to \lvert w \rvert = \lvert x - y \rvert + \lvert y - z \rvert)$
\item $\varphi_{47} :\equiv \forall u\,\forall v\,(v_{740} \in \mathbb{R} \to v_{741} \in \mathbb{R} \to u \in \mathbb{R} \to v \in \mathbb{R} \to v_{740} < v_{741} \to u < v \to v_{740} + u < v_{741} + v)$
\item $\varphi_{48} :\equiv \forall w\,(x - y + (y - z) = w \to \lvert x - y + (y - z) \rvert < \lvert x - y \rvert + \lvert y - z \rvert \to \lvert w \rvert < \lvert x - y \rvert + \lvert y - z \rvert)$
\item $\varphi_{49} :\equiv \forall v\,(v_{740} \in \mathbb{R} \to v_{741} \in \mathbb{R} \to v_{742} \in \mathbb{R} \to v \in \mathbb{R} \to v_{740} < v_{741} \to v_{742} < v \to v_{740} + v_{742} < v_{741} + v)$
\item $\varphi_{50} :\equiv \forall v_{743}\,(v_{740} \in \mathbb{R} \to v_{741} \in \mathbb{R} \to v_{742} \in \mathbb{R} \to v_{743} \in \mathbb{R} \to \varphi_{3})$
\item $\varphi_{51} :\equiv \forall a\,\forall b\,(a \in \mathbb{R} \to b \in \mathbb{R} \to \lvert a + b \rvert \leq \lvert a \rvert + \lvert b \rvert) \to \varphi_{10}$
\item $\varphi_{52} :\equiv \forall v_{742}\,\forall v_{743}\,(\lvert x - y \rvert \in \mathbb{R} \to e \in \mathbb{R} \to v_{742} \in \mathbb{R} \to v_{743} \in \mathbb{R} \to \lvert x - y \rvert < e \to v_{742} < v_{743} \to \lvert x - y \rvert + v_{742} < e + v_{743})$
\item $\varphi_{53} :\equiv \forall v_{741}\,\forall v_{742}\,\forall v_{743}\,(\lvert x - y \rvert \in \mathbb{R} \to v_{741} \in \mathbb{R} \to v_{742} \in \mathbb{R} \to v_{743} \in \mathbb{R} \to \lvert x - y \rvert < v_{741} \to v_{742} < v_{743} \to \lvert x - y \rvert + v_{742} < v_{741} + v_{743})$
\item $\varphi_{54} :\equiv \forall v_{743}\,(\lvert x - y \rvert \in \mathbb{R} \to e \in \mathbb{R} \to \lvert y - z \rvert \in \mathbb{R} \to v_{743} \in \mathbb{R} \to \lvert x - y \rvert < e \to \lvert y - z \rvert < v_{743} \to \lvert x - y \rvert + \lvert y - z \rvert < e + v_{743})$
\item $\varphi_{55} :\equiv (x - z = x - y + (y - z) \to x - y + (y - z) = x - z) \to x \in \mathbb{R} \to y \in \mathbb{R} \to z \in \mathbb{R} \to \varphi_{15}$
\item $\varphi_{56} :\equiv \forall v_{735}\,\forall v_{736}\,\forall v_{737}\,(v_{735} \in \mathbb{R} \to v_{736} \in \mathbb{R} \to \varphi_{2}) \to \varphi_{36}$
\item $\varphi_{57} :\equiv (x \in \mathbb{R} \to y \in \mathbb{R} \to \varphi_{23}) \to x \in \mathbb{R} \to y \in \mathbb{R} \to z \in \mathbb{R} \to e \in \mathbb{R} \to \varphi_{13}$
\item $\varphi_{58} :\equiv (x \in \mathbb{R} \to y \in \mathbb{R} \to z \in \mathbb{R} \to e \in \mathbb{R} \to \varphi_{13}) \to x \in \mathbb{R} \to y \in \mathbb{R} \to \varphi_{29}$
\item $\varphi_{59} :\equiv \varphi_{49} \to v_{740} \in \mathbb{R} \to v_{741} \in \mathbb{R} \to v_{742} \in \mathbb{R} \to v_{743} \in \mathbb{R} \to \varphi_{3}$
\item $\varphi_{60} :\equiv (x \in \mathbb{R} \to y \in \mathbb{R} \to \varphi_{21}) \to x \in \mathbb{R} \to y \in \mathbb{R} \to z \in \mathbb{R} \to e \in \mathbb{R} \to \lvert x - y \rvert < e \to \varphi_{7}$
\item $\varphi_{61} :\equiv (x - y + (y - z) = x - z \to \varphi_{35}) \to x \in \mathbb{R} \to y \in \mathbb{R} \to z \in \mathbb{R} \to e \in \mathbb{R} \to \varphi_{17}$
\item $\varphi_{62} :\equiv (y - z \in \mathbb{R} \to \lvert y - z \rvert \in \mathbb{R}) \to (x \in \mathbb{R} \to y \in \mathbb{R} \to \varphi_{21}) \to x \in \mathbb{R} \to y \in \mathbb{R} \to \varphi_{27}$
\item $\varphi_{63} :\equiv (x - y \in \mathbb{R} \to \lvert x - y \rvert \in \mathbb{R}) \to (x \in \mathbb{R} \to y \in \mathbb{R} \to \varphi_{22}) \to x \in \mathbb{R} \to y \in \mathbb{R} \to \varphi_{28}$
\item $\varphi_{64} :\equiv \varphi_{54} \to \lvert x - y \rvert \in \mathbb{R} \to e \in \mathbb{R} \to \lvert y - z \rvert \in \mathbb{R} \to e \in \mathbb{R} \to \varphi_{13}$
\item $\varphi_{65} :\equiv (x \in \mathbb{R} \to \varphi_{31}) \to (x \in \mathbb{R} \to \varphi_{32}) \to x \in \mathbb{R} \to y \in \mathbb{R} \to z \in \mathbb{R} \to \varphi_{14}$
\item $\varphi_{66} :\equiv (y \in \mathbb{R} \to z \in \mathbb{R} \to y - z \in \mathbb{R}) \to (x \in \mathbb{R} \to \varphi_{32}) \to (x \in \mathbb{R} \to \varphi_{31}) \to x \in \mathbb{R} \to y \in \mathbb{R} \to \varphi_{21}$
\item $\varphi_{67} :\equiv (x \in \mathbb{R} \to y \in \mathbb{R} \to x - y \in \mathbb{R}) \to (x \in \mathbb{R} \to \varphi_{33}) \to (x \in \mathbb{R} \to \varphi_{32}) \to x \in \mathbb{R} \to y \in \mathbb{R} \to \varphi_{22}$
\item $\varphi_{68} :\equiv (x \in \mathbb{R} \to z \in \mathbb{R} \to y \in \mathbb{R} \to x - z = x - y + (y - z)) \to (x \in \mathbb{R} \to \varphi_{33}) \to \varphi_{65}$
\item $\varphi_{69} :\equiv (x \in \mathbb{R} \to y \in \mathbb{R} \to z \in \mathbb{R} \to e \in \mathbb{R} \to \lvert x - y \rvert < e \to \varphi_{7}) \to (x - y + (y - z) = x - z \to \varphi_{38}) \to \varphi_{61}$
\item $\varphi_{70} :\equiv (x \in \mathbb{R} \to y \in \mathbb{R} \to \varphi_{27}) \to (x \in \mathbb{R} \to \varphi_{30}) \to (x \in \mathbb{R} \to y \in \mathbb{R} \to \varphi_{24}) \to \varphi_{57}$
\item $\varphi_{71} :\equiv (\lvert x - y \rvert \in \mathbb{R} \to e \in \mathbb{R} \to \lvert y - z \rvert \in \mathbb{R} \to e \in \mathbb{R} \to \varphi_{13}) \to (x \in \mathbb{R} \to y \in \mathbb{R} \to \varphi_{28}) \to (x \in \mathbb{R} \to \varphi_{30}) \to \varphi_{70}$
\end{abbrevlist}
\end{fullproof}

% 自動生成 (prover/build.js): CauchyBdd  ノード数 412
\begin{fullproof}{補題6.10}{$\forall a\,(\mathrm{Cauchy}(a) \to \mathrm{SeqBdd}(a))$}{412}\phantomsection\label{tree:CauchyBdd}
\noindent$\Pi_{1}$:
\begin{ptree}
\AXC{$\forall a\,(\varphi_{94})$\lbl{DefCauchy}}
\AXC{$\forall a\,(\varphi_{94}) \to \varphi_{94}$\lbl{A2}}
\BIC{$\varphi_{94}$}
\AXC{$\mathrm{Cauchy}(a) \to \mathrm{Cauchy}(a)$\lbl{A1}}
\AXC{$\varphi_{108}$\lbl{A1}}
\BIC{$\varphi_{94} \to \mathrm{Cauchy}(a) \to \mathrm{Seq}(a)$}
\BIC{$\mathrm{Cauchy}(a) \to \mathrm{Seq}(a)$}
\end{ptree}
\noindent$\Pi_{2}$:
\begin{ptree}
\AXC{$\stackrel{\Pi_{1}}{\vdots}$}\noLine
\UIC{$\mathrm{Cauchy}(a) \to \mathrm{Seq}(a)$}
\AXC{$\forall a\,\forall n\,(\mathrm{Seq}(a) \to n \in \mathbb{N} \to a_{n} \in \mathbb{R})$\lbl{[S1]}}
\AXC{$\varphi_{36}$\lbl{A2}}
\BIC{$\forall n\,(\mathrm{Seq}(a) \to n \in \mathbb{N} \to a_{n} \in \mathbb{R})$}
\AXC{$\varphi_{25}$\lbl{A2}}
\BIC{$\mathrm{Seq}(a) \to N \in \mathbb{N} \to a_{N} \in \mathbb{R}$}
\AXC{$\varphi_{133}$\lbl{A1}}
\BIC{$(\mathrm{Cauchy}(a) \to \mathrm{Seq}(a)) \to \varphi_{126}$}
\BIC{$\varphi_{126}$}
\end{ptree}
\noindent$\Pi_{3}$:
\begin{ptree}
\AXC{$\mathrm{Cauchy}(a) \to N \in \mathbb{N} \land \varphi_{10} \to N \in \mathbb{N} \land \varphi_{10}$\lbl{A1}}
\AXC{$\varphi_{128}$\lbl{A1}}
\BIC{$\mathrm{Cauchy}(a) \to N \in \mathbb{N} \land \varphi_{10} \to N \in \mathbb{N}$}
\AXC{$\stackrel{\Pi_{2}}{\vdots}$}\noLine
\UIC{$\varphi_{126}$}
\BIC{$\mathrm{Cauchy}(a) \to N \in \mathbb{N} \land \varphi_{10} \to \varphi_{13}$}
\end{ptree}
\noindent$\Pi_{4}$:
\begin{ptree}
\AXC{$\forall x \in \mathbb{R}\,(0 \leq \lvert x \rvert)$\lbl{[A6]}}
\AXC{$\forall x \in \mathbb{R}\,(0 \leq \lvert x \rvert) \to a_{N} \in \mathbb{R} \to 0 \leq \lvert a_{N} \rvert$\lbl{A2}}
\BIC{$a_{N} \in \mathbb{R} \to 0 \leq \lvert a_{N} \rvert$}
\AXC{$(a_{N} \in \mathbb{R} \to 0 \leq \lvert a_{N} \rvert) \to \varphi_{134}$\lbl{A1}}
\BIC{$\varphi_{134}$}
\end{ptree}
\noindent$\Pi_{5}$:
\begin{ptree}
\AXC{$0 < 1$\lbl{[F28]}}
\AXC{$0 < 1 \to \mathrm{Cauchy}(a) \to 0 < 1$\lbl{A1}}
\BIC{$\mathrm{Cauchy}(a) \to 0 < 1$}
\AXC{$(\mathrm{Cauchy}(a) \to 0 < 1) \to \mathrm{Cauchy}(a) \to \varphi_{98}$\lbl{A1}}
\BIC{$\mathrm{Cauchy}(a) \to \varphi_{98}$}
\end{ptree}
\noindent$\Pi_{6}$:
\begin{ptree}
\AXC{$\forall x \in \mathbb{R}\,(\lvert x \rvert \in \mathbb{R})$\lbl{[A5]}}
\AXC{$\forall x \in \mathbb{R}\,(\lvert x \rvert \in \mathbb{R}) \to a_{N} \in \mathbb{R} \to \lvert a_{N} \rvert \in \mathbb{R}$\lbl{A2}}
\BIC{$a_{N} \in \mathbb{R} \to \lvert a_{N} \rvert \in \mathbb{R}$}
\AXC{$(a_{N} \in \mathbb{R} \to \lvert a_{N} \rvert \in \mathbb{R}) \to \varphi_{137}$\lbl{A1}}
\BIC{$\varphi_{137}$}
\end{ptree}
\noindent$\Pi_{7}$:
\begin{ptree}
\AXC{$\stackrel{\Pi_{3}}{\vdots}$}\noLine
\UIC{$\mathrm{Cauchy}(a) \to N \in \mathbb{N} \land \varphi_{10} \to \varphi_{13}$}
\AXC{$\stackrel{\Pi_{6}}{\vdots}$}\noLine
\UIC{$\varphi_{137}$}
\BIC{$\mathrm{Cauchy}(a) \to N \in \mathbb{N} \land \varphi_{10} \to \varphi_{21}$}
\end{ptree}
\noindent$\Pi_{8}$:
\begin{ptree}
\AXC{$\forall x\,\forall y\,(x \in \mathbb{R} \to y \in \mathbb{R} \to 0 \leq x \to 0 \leq y \to 0 \leq x + y)$\lbl{[E6]}}
\AXC{$\varphi_{58}$\lbl{A2}}
\BIC{$\forall y\,(1 \in \mathbb{R} \to y \in \mathbb{R} \to 0 \leq 1 \to 0 \leq y \to 0 \leq 1 + y)$}
\AXC{$\forall y\,(1 \in \mathbb{R} \to y \in \mathbb{R} \to 0 \leq 1 \to 0 \leq y \to 0 \leq 1 + y) \to \varphi_{26}$\lbl{A2}}
\BIC{$\varphi_{26}$}
\end{ptree}
\noindent$\Pi_{9}$:
\begin{ptree}
\AXC{$\stackrel{\Pi_{5}}{\vdots}$}\noLine
\UIC{$\mathrm{Cauchy}(a) \to \varphi_{98}$}
\AXC{$\stackrel{\Pi_{7}}{\vdots}$}\noLine
\UIC{$\mathrm{Cauchy}(a) \to N \in \mathbb{N} \land \varphi_{10} \to \varphi_{21}$}
\AXC{$1 \in \mathbb{R}$\lbl{(R1)}}
\AXC{$\stackrel{\Pi_{8}}{\vdots}$}\noLine
\UIC{$\varphi_{26}$}
\AXC{$\varphi_{26} \to 1 \in \mathbb{R} \to \varphi_{188}$\lbl{A1}}
\BIC{$1 \in \mathbb{R} \to \varphi_{188}$}
\BIC{$\varphi_{188}$}
\BIC{$(\mathrm{Cauchy}(a) \to \varphi_{98}) \to \varphi_{136}$}
\BIC{$\varphi_{136}$}
\end{ptree}
\noindent$\Pi_{10}$:
\begin{ptree}
\AXC{$\stackrel{\Pi_{1}}{\vdots}$}\noLine
\UIC{$\mathrm{Cauchy}(a) \to \mathrm{Seq}(a)$}
\AXC{$\forall a\,\forall n\,(\mathrm{Seq}(a) \to n \in \mathbb{N} \to a_{n} \in \mathbb{R})$\lbl{[S1]}}
\AXC{$\varphi_{36}$\lbl{A2}}
\BIC{$\forall n\,(\mathrm{Seq}(a) \to n \in \mathbb{N} \to a_{n} \in \mathbb{R})$}
\AXC{$\varphi_{25}$\lbl{A2}}
\BIC{$\mathrm{Seq}(a) \to N \in \mathbb{N} \to a_{N} \in \mathbb{R}$}
\AXC{$\varphi_{144}$\lbl{A1}}
\BIC{$(\mathrm{Cauchy}(a) \to \mathrm{Seq}(a)) \to \varphi_{130}$}
\BIC{$\varphi_{130}$}
\end{ptree}
\noindent$\Pi_{11}$:
\begin{ptree}
\AXC{$\forall x \in \mathbb{R}\,(\lvert x \rvert \in \mathbb{R})$\lbl{[A5]}}
\AXC{$\forall x \in \mathbb{R}\,(\lvert x \rvert \in \mathbb{R}) \to a_{N} \in \mathbb{R} \to \lvert a_{N} \rvert \in \mathbb{R}$\lbl{A2}}
\BIC{$a_{N} \in \mathbb{R} \to \lvert a_{N} \rvert \in \mathbb{R}$}
\AXC{$(a_{N} \in \mathbb{R} \to \lvert a_{N} \rvert \in \mathbb{R}) \to \varphi_{156}$\lbl{A1}}
\BIC{$\varphi_{156}$}
\end{ptree}
\noindent$\Pi_{12}$:
\begin{ptree}
\AXC{$1 \in \mathbb{R}$\lbl{(R1)}}
\AXC{$\forall x\,\forall y\,(x \in \mathbb{R} \to y \in \mathbb{R} \to x + y \in \mathbb{R})$\lbl{(R1)}}
\AXC{$\varphi_{44}$\lbl{A2}}
\BIC{$\forall y\,(1 \in \mathbb{R} \to y \in \mathbb{R} \to 1 + y \in \mathbb{R})$}
\AXC{$\varphi_{72}$\lbl{A2}}
\BIC{$1 \in \mathbb{R} \to \lvert a_{N} \rvert \in \mathbb{R} \to 1 + \lvert a_{N} \rvert \in \mathbb{R}$}
\AXC{$\varphi_{177}$\lbl{A1}}
\BIC{$1 \in \mathbb{R} \to \varphi_{159}$}
\BIC{$\varphi_{159}$}
\end{ptree}
\noindent$\Pi_{13}$:
\begin{ptree}
\AXC{$\mathrm{Cauchy}(a) \to N \in \mathbb{N} \land \varphi_{10} \to N \in \mathbb{N} \land \varphi_{10}$\lbl{A1}}
\AXC{$\varphi_{128}$\lbl{A1}}
\BIC{$\mathrm{Cauchy}(a) \to N \in \mathbb{N} \land \varphi_{10} \to N \in \mathbb{N}$}
\AXC{$\stackrel{\Pi_{10}}{\vdots}$}\noLine
\UIC{$\varphi_{130}$}
\BIC{$\mathrm{Cauchy}(a) \to N \in \mathbb{N} \land \varphi_{10} \to \varphi_{49}$}
\AXC{$\stackrel{\Pi_{11}}{\vdots}$}\noLine
\UIC{$\varphi_{156}$}
\BIC{$\mathrm{Cauchy}(a) \to N \in \mathbb{N} \land \varphi_{10} \to \varphi_{53}$}
\AXC{$\stackrel{\Pi_{12}}{\vdots}$}\noLine
\UIC{$\varphi_{159}$}
\BIC{$\mathrm{Cauchy}(a) \to N \in \mathbb{N} \land \varphi_{10} \to \varphi_{62}$}
\end{ptree}
\noindent$\Pi_{14}$:
\begin{ptree}
\AXC{$\forall x\,\forall y\,(x \in \mathbb{R} \to y \in \mathbb{R} \to 0 \leq y \to x \leq x + y)$\lbl{[E1]}}
\AXC{$\varphi_{40}$\lbl{A2}}
\BIC{$\forall y\,(M \in \mathbb{R} \to y \in \mathbb{R} \to 0 \leq y \to M \leq M + y)$}
\AXC{$\forall y\,(M \in \mathbb{R} \to y \in \mathbb{R} \to 0 \leq y \to M \leq M + y) \to M \in \mathbb{R} \to \varphi_{16}$\lbl{A2}}
\BIC{$M \in \mathbb{R} \to \varphi_{16}$}
\end{ptree}
\noindent$\Pi_{15}$:
\begin{ptree}
\AXC{$\mathrm{Cauchy}(a) \to N \in \mathbb{N} \land \varphi_{10} \to \varphi_{76}$\lbl{A1}}
\AXC{$\varphi_{152}$\lbl{A1}}
\BIC{$\mathrm{Cauchy}(a) \to N \in \mathbb{N} \land \varphi_{10} \to \varphi_{5}$}
\AXC{$\stackrel{\Pi_{14}}{\vdots}$}\noLine
\UIC{$M \in \mathbb{R} \to \varphi_{16}$}
\AXC{$(M \in \mathbb{R} \to \varphi_{16}) \to \varphi_{190}$\lbl{A1}}
\BIC{$\varphi_{190}$}
\BIC{$(\mathrm{Cauchy}(a) \to N \in \mathbb{N} \land \varphi_{10} \to \varphi_{62}) \to \varphi_{153}$}
\end{ptree}
\noindent$\Pi_{16}$:
\begin{ptree}
\AXC{$\stackrel{\Pi_{13}}{\vdots}$}\noLine
\UIC{$\mathrm{Cauchy}(a) \to N \in \mathbb{N} \land \varphi_{10} \to \varphi_{62}$}
\AXC{$\stackrel{\Pi_{15}}{\vdots}$}\noLine
\UIC{$(\mathrm{Cauchy}(a) \to N \in \mathbb{N} \land \varphi_{10} \to \varphi_{62}) \to \varphi_{153}$}
\BIC{$\varphi_{153}$}
\end{ptree}
\noindent$\Pi_{17}$:
\begin{ptree}
\AXC{$\varphi_{8}$\lbl{A2}}
\AXC{$\mathrm{Cauchy}(a) \to N \in \mathbb{N} \land \varphi_{10} \to \varphi_{76}$\lbl{A1}}
\AXC{$\varphi_{168}$\lbl{A1}}
\BIC{$\mathrm{Cauchy}(a) \to N \in \mathbb{N} \land \varphi_{10} \to \varphi_{79}$}
\AXC{$\varphi_{179}$\lbl{A1}}
\BIC{$\varphi_{8} \to \mathrm{Cauchy}(a) \to N \in \mathbb{N} \land \varphi_{10} \to \varphi_{74}$}
\BIC{$\mathrm{Cauchy}(a) \to N \in \mathbb{N} \land \varphi_{10} \to \varphi_{74}$}
\AXC{$\varphi_{192}$\lbl{A1}}
\BIC{$(\mathrm{Cauchy}(a) \to N \in \mathbb{N} \land \varphi_{10} \to \varphi_{18}) \to \varphi_{150}$}
\end{ptree}
\noindent$\Pi_{18}$:
\begin{ptree}
\AXC{$\mathrm{Cauchy}(a) \to N \in \mathbb{N} \land \varphi_{10} \to \varphi_{32}$\lbl{A1}}
\AXC{$\mathrm{Cauchy}(a) \to N \in \mathbb{N} \land \varphi_{10} \to \varphi_{18}$\lbl{A1}}
\AXC{$\stackrel{\Pi_{17}}{\vdots}$}\noLine
\UIC{$(\mathrm{Cauchy}(a) \to N \in \mathbb{N} \land \varphi_{10} \to \varphi_{18}) \to \varphi_{150}$}
\BIC{$\varphi_{150}$}
\BIC{$\mathrm{Cauchy}(a) \to N \in \mathbb{N} \land \varphi_{10} \to \varphi_{47}$}
\end{ptree}
\noindent$\Pi_{19}$:
\begin{ptree}
\AXC{$\forall x\,\forall y\,\forall z\,(x \leq y \to y \leq z \to x \leq z)$\lbl{[O5]}}
\AXC{$\varphi_{9}$\lbl{A2}}
\BIC{$\forall y\,\forall z\,(\lvert a_{n} \rvert \leq y \to y \leq z \to \lvert a_{n} \rvert \leq z)$}
\AXC{$\varphi_{37}$\lbl{A2}}
\BIC{$\forall z\,(\lvert a_{n} \rvert \leq M \to M \leq z \to \lvert a_{n} \rvert \leq z)$}
\AXC{$\forall z\,(\lvert a_{n} \rvert \leq M \to M \leq z \to \lvert a_{n} \rvert \leq z) \to \varphi_{30}$\lbl{A2}}
\BIC{$\varphi_{30}$}
\end{ptree}
\noindent$\Pi_{20}$:
\begin{ptree}
\AXC{$\stackrel{\Pi_{18}}{\vdots}$}\noLine
\UIC{$\mathrm{Cauchy}(a) \to N \in \mathbb{N} \land \varphi_{10} \to \varphi_{47}$}
\AXC{$\stackrel{\Pi_{19}}{\vdots}$}\noLine
\UIC{$\varphi_{30}$}
\AXC{$\varphi_{30} \to (\mathrm{Cauchy}(a) \to N \in \mathbb{N} \land \varphi_{10} \to \varphi_{47}) \to \varphi_{165}$\lbl{A1}}
\BIC{$(\mathrm{Cauchy}(a) \to N \in \mathbb{N} \land \varphi_{10} \to \varphi_{47}) \to \varphi_{165}$}
\BIC{$\varphi_{165}$}
\end{ptree}
\noindent$\Pi_{21}$:
\begin{ptree}
\AXC{$\stackrel{\Pi_{1}}{\vdots}$}\noLine
\UIC{$\mathrm{Cauchy}(a) \to \mathrm{Seq}(a)$}
\AXC{$\forall a\,\forall n\,(\mathrm{Seq}(a) \to n \in \mathbb{N} \to a_{n} \in \mathbb{R})$\lbl{[S1]}}
\AXC{$\varphi_{36}$\lbl{A2}}
\BIC{$\forall n\,(\mathrm{Seq}(a) \to n \in \mathbb{N} \to a_{n} \in \mathbb{R})$}
\AXC{$\varphi_{25}$\lbl{A2}}
\BIC{$\mathrm{Seq}(a) \to N \in \mathbb{N} \to a_{N} \in \mathbb{R}$}
\AXC{$\varphi_{142}$\lbl{A1}}
\BIC{$(\mathrm{Cauchy}(a) \to \mathrm{Seq}(a)) \to \varphi_{129}$}
\BIC{$\varphi_{129}$}
\end{ptree}
\noindent$\Pi_{22}$:
\begin{ptree}
\AXC{$\mathrm{Cauchy}(a) \to N \in \mathbb{N} \land \varphi_{10} \to N \in \mathbb{N} \land \varphi_{10}$\lbl{A1}}
\AXC{$\varphi_{128}$\lbl{A1}}
\BIC{$\mathrm{Cauchy}(a) \to N \in \mathbb{N} \land \varphi_{10} \to N \in \mathbb{N}$}
\AXC{$\stackrel{\Pi_{21}}{\vdots}$}\noLine
\UIC{$\varphi_{129}$}
\BIC{$\mathrm{Cauchy}(a) \to N \in \mathbb{N} \land \varphi_{10} \to \varphi_{42}$}
\end{ptree}
\noindent$\Pi_{23}$:
\begin{ptree}
\AXC{$\forall x \in \mathbb{R}\,(\lvert x \rvert \in \mathbb{R})$\lbl{[A5]}}
\AXC{$\forall x \in \mathbb{R}\,(\lvert x \rvert \in \mathbb{R}) \to a_{N} \in \mathbb{R} \to \lvert a_{N} \rvert \in \mathbb{R}$\lbl{A2}}
\BIC{$a_{N} \in \mathbb{R} \to \lvert a_{N} \rvert \in \mathbb{R}$}
\AXC{$(a_{N} \in \mathbb{R} \to \lvert a_{N} \rvert \in \mathbb{R}) \to \varphi_{155}$\lbl{A1}}
\BIC{$\varphi_{155}$}
\end{ptree}
\noindent$\Pi_{24}$:
\begin{ptree}
\AXC{$\stackrel{\Pi_{22}}{\vdots}$}\noLine
\UIC{$\mathrm{Cauchy}(a) \to N \in \mathbb{N} \land \varphi_{10} \to \varphi_{42}$}
\AXC{$\stackrel{\Pi_{23}}{\vdots}$}\noLine
\UIC{$\varphi_{155}$}
\BIC{$\mathrm{Cauchy}(a) \to N \in \mathbb{N} \land \varphi_{10} \to \varphi_{50}$}
\end{ptree}
\noindent$\Pi_{25}$:
\begin{ptree}
\AXC{$1 \in \mathbb{R}$\lbl{(R1)}}
\AXC{$\forall x\,\forall y\,(x \in \mathbb{R} \to y \in \mathbb{R} \to x + y \in \mathbb{R})$\lbl{(R1)}}
\AXC{$\varphi_{44}$\lbl{A2}}
\BIC{$\forall y\,(1 \in \mathbb{R} \to y \in \mathbb{R} \to 1 + y \in \mathbb{R})$}
\AXC{$\varphi_{72}$\lbl{A2}}
\BIC{$1 \in \mathbb{R} \to \lvert a_{N} \rvert \in \mathbb{R} \to 1 + \lvert a_{N} \rvert \in \mathbb{R}$}
\AXC{$\varphi_{176}$\lbl{A1}}
\BIC{$1 \in \mathbb{R} \to \varphi_{158}$}
\BIC{$\varphi_{158}$}
\end{ptree}
\noindent$\Pi_{26}$:
\begin{ptree}
\AXC{$\forall x\,\forall y\,(x \in \mathbb{R} \to y \in \mathbb{R} \to 0 \leq y \to x \leq y + x)$\lbl{[E2]}}
\AXC{$\forall x\,\forall y\,(x \in \mathbb{R} \to y \in \mathbb{R} \to 0 \leq y \to x \leq y + x) \to \varphi_{45}$\lbl{A2}}
\BIC{$\varphi_{45}$}
\AXC{$\varphi_{45} \to \varphi_{39}$\lbl{A2}}
\BIC{$\varphi_{39}$}
\end{ptree}
\noindent$\Pi_{27}$:
\begin{ptree}
\AXC{$\stackrel{\Pi_{24}}{\vdots}$}\noLine
\UIC{$\mathrm{Cauchy}(a) \to N \in \mathbb{N} \land \varphi_{10} \to \varphi_{50}$}
\AXC{$\stackrel{\Pi_{25}}{\vdots}$}\noLine
\UIC{$\varphi_{158}$}
\BIC{$\mathrm{Cauchy}(a) \to N \in \mathbb{N} \land \varphi_{10} \to \varphi_{56}$}
\AXC{$\stackrel{\Pi_{26}}{\vdots}$}\noLine
\UIC{$\varphi_{39}$}
\AXC{$\varphi_{39} \to \varphi_{189}$\lbl{A1}}
\BIC{$\varphi_{189}$}
\BIC{$(\mathrm{Cauchy}(a) \to N \in \mathbb{N} \land \varphi_{10} \to \varphi_{5}) \to \varphi_{151}$}
\end{ptree}
\noindent$\Pi_{28}$:
\begin{ptree}
\AXC{$\mathrm{Cauchy}(a) \to N \in \mathbb{N} \land \varphi_{10} \to \varphi_{76}$\lbl{A1}}
\AXC{$\varphi_{152}$\lbl{A1}}
\BIC{$\mathrm{Cauchy}(a) \to N \in \mathbb{N} \land \varphi_{10} \to \varphi_{5}$}
\AXC{$\stackrel{\Pi_{27}}{\vdots}$}\noLine
\UIC{$(\mathrm{Cauchy}(a) \to N \in \mathbb{N} \land \varphi_{10} \to \varphi_{5}) \to \varphi_{151}$}
\BIC{$\varphi_{151}$}
\end{ptree}
\noindent$\Pi_{29}$:
\begin{ptree}
\AXC{$x = x$\lbl{A5}}
\UIC{$\forall x\,(x = x)$}
\AXC{$\forall x\,(x = x) \to \lvert a_{N} \rvert = \lvert a_{N} \rvert$\lbl{A2}}
\BIC{$\lvert a_{N} \rvert = \lvert a_{N} \rvert$}
\AXC{$\lvert a_{N} \rvert = \lvert a_{N} \rvert \to \mathrm{Cauchy}(a) \to N \in \mathbb{N} \land \varphi_{10} \to \varphi_{57}$\lbl{A1}}
\BIC{$\mathrm{Cauchy}(a) \to N \in \mathbb{N} \land \varphi_{10} \to \varphi_{57}$}
\end{ptree}
\noindent$\Pi_{30}$:
\begin{ptree}
\AXC{$\varphi_{10} \to \varphi_{17}$\lbl{A2}}
\AXC{$\mathrm{Cauchy}(a) \to N \in \mathbb{N} \land \varphi_{10} \to N \in \mathbb{N} \land \varphi_{10}$\lbl{A1}}
\AXC{$\varphi_{138}$\lbl{A1}}
\BIC{$\mathrm{Cauchy}(a) \to N \in \mathbb{N} \land \varphi_{10} \to \varphi_{10}$}
\AXC{$\varphi_{180}$\lbl{A1}}
\BIC{$(\varphi_{10} \to \varphi_{17}) \to \mathrm{Cauchy}(a) \to N \in \mathbb{N} \land \varphi_{10} \to \varphi_{17}$}
\BIC{$\mathrm{Cauchy}(a) \to N \in \mathbb{N} \land \varphi_{10} \to \varphi_{17}$}
\end{ptree}
\noindent$\Pi_{31}$:
\begin{ptree}
\AXC{$\varphi_{17} \to v_{747} \in \mathbb{N} \to \varphi_{6}$\lbl{A2}}
\AXC{$\stackrel{\Pi_{30}}{\vdots}$}\noLine
\UIC{$\mathrm{Cauchy}(a) \to N \in \mathbb{N} \land \varphi_{10} \to \varphi_{17}$}
\AXC{$(\mathrm{Cauchy}(a) \to N \in \mathbb{N} \land \varphi_{10} \to \varphi_{17}) \to \varphi_{135}$\lbl{A1}}
\BIC{$\varphi_{135}$}
\BIC{$\mathrm{Cauchy}(a) \to N \in \mathbb{N} \land \varphi_{10} \to v_{747} \in \mathbb{N} \to \varphi_{6}$}
\AXC{$\varphi_{145}$\lbl{A1}}
\BIC{$\mathrm{Cauchy}(a) \land (N \in \mathbb{N} \land \varphi_{10}) \to v_{747} \in \mathbb{N} \to \varphi_{6}$}
\UIC{$\varphi_{109}$}
\end{ptree}
\noindent$\Pi_{32}$:
\begin{ptree}
\AXC{$\varphi_{117}$\lbl{A4}}
\AXC{$\varphi_{109} \to \varphi_{115}$\lbl{A3}}
\AXC{$\stackrel{\Pi_{31}}{\vdots}$}\noLine
\UIC{$\varphi_{109}$}
\AXC{$\varphi_{109} \to (\varphi_{109} \to \varphi_{115}) \to \varphi_{117} \to \varphi_{110}$\lbl{A1}}
\BIC{$(\varphi_{109} \to \varphi_{115}) \to \varphi_{117} \to \varphi_{110}$}
\BIC{$\varphi_{117} \to \varphi_{110}$}
\BIC{$\varphi_{110}$}
\AXC{$\varphi_{110} \to \varphi_{111}$\lbl{A1}}
\BIC{$\varphi_{111}$}
\UIC{$\forall v_{747}\,(\varphi_{111})$}
\end{ptree}
\noindent$\Pi_{33}$:
\begin{ptree}
\AXC{$\varphi_{116}$\lbl{A4}}
\AXC{$\forall v_{747}\,(\varphi_{111}) \to \varphi_{118}$\lbl{A3}}
\AXC{$\stackrel{\Pi_{32}}{\vdots}$}\noLine
\UIC{$\forall v_{747}\,(\varphi_{111})$}
\AXC{$\forall v_{747}\,(\varphi_{111}) \to (\forall v_{747}\,(\varphi_{111}) \to \varphi_{118}) \to \varphi_{116} \to \varphi_{114}$\lbl{A1}}
\BIC{$(\forall v_{747}\,(\varphi_{111}) \to \varphi_{118}) \to \varphi_{116} \to \varphi_{114}$}
\BIC{$\varphi_{116} \to \varphi_{114}$}
\BIC{$\varphi_{114}$}
\AXC{$\varphi_{114} \to \varphi_{140}$\lbl{A1}}
\BIC{$\varphi_{140}$}
\end{ptree}
\noindent$\Pi_{34}$:
\begin{ptree}
\AXC{$\forall v_{747}\,\forall v_{748}\,(v_{747} \in \mathbb{N} \to \varphi_{6}) \to \varphi_{22}$\lbl{A2}}
\AXC{$\stackrel{\Pi_{33}}{\vdots}$}\noLine
\UIC{$\varphi_{140}$}
\BIC{$\mathrm{Cauchy}(a) \to N \in \mathbb{N} \land \varphi_{10} \to \varphi_{22}$}
\AXC{$(\mathrm{Cauchy}(a) \to N \in \mathbb{N} \land \varphi_{10} \to \varphi_{22}) \to \varphi_{131}$\lbl{A1}}
\BIC{$\varphi_{131}$}
\end{ptree}
\noindent$\Pi_{35}$:
\begin{ptree}
\AXC{$\mathrm{Cauchy}(a) \to N \in \mathbb{N} \land \varphi_{10} \to \varphi_{18}$\lbl{A1}}
\AXC{$\varphi_{22} \to n \in \mathbb{N} \to N \in \mathbb{N} \to N \leq n \to N \leq N \to \lvert a_{n} - a_{N} \rvert < 1$\lbl{A2}}
\AXC{$\stackrel{\Pi_{34}}{\vdots}$}\noLine
\UIC{$\varphi_{131}$}
\BIC{$\mathrm{Cauchy}(a) \to \varphi_{100}$}
\AXC{$\varphi_{197}$\lbl{A1}}
\BIC{$(\mathrm{Cauchy}(a) \to N \in \mathbb{N} \land \varphi_{10} \to \varphi_{18}) \to \varphi_{185}$}
\BIC{$\varphi_{185}$}
\end{ptree}
\noindent$\Pi_{36}$:
\begin{ptree}
\AXC{$\mathrm{Cauchy}(a) \to N \in \mathbb{N} \land \varphi_{10} \to \varphi_{20}$\lbl{A1}}
\AXC{$\mathrm{Cauchy}(a) \to N \in \mathbb{N} \land \varphi_{10} \to N \in \mathbb{N} \land \varphi_{10}$\lbl{A1}}
\AXC{$\varphi_{128}$\lbl{A1}}
\BIC{$\mathrm{Cauchy}(a) \to N \in \mathbb{N} \land \varphi_{10} \to N \in \mathbb{N}$}
\AXC{$\stackrel{\Pi_{35}}{\vdots}$}\noLine
\UIC{$\varphi_{185}$}
\BIC{$(\mathrm{Cauchy}(a) \to N \in \mathbb{N} \land \varphi_{10} \to \varphi_{20}) \to \varphi_{154}$}
\BIC{$\varphi_{154}$}
\end{ptree}
\noindent$\Pi_{37}$:
\begin{ptree}
\AXC{$x = x$\lbl{A5}}
\UIC{$\forall x\,(x = x)$}
\AXC{$\forall x\,(x = x) \to N = N$\lbl{A2}}
\BIC{$N = N$}
\AXC{$N = N \to \mathrm{Cauchy}(a) \to N \in \mathbb{N} \land \varphi_{10} \to \varphi_{19}$\lbl{A1}}
\BIC{$\mathrm{Cauchy}(a) \to N \in \mathbb{N} \land \varphi_{10} \to \varphi_{19}$}
\AXC{$\stackrel{\Pi_{36}}{\vdots}$}\noLine
\UIC{$\varphi_{154}$}
\BIC{$\mathrm{Cauchy}(a) \to N \in \mathbb{N} \land \varphi_{10} \to \varphi_{63}$}
\end{ptree}
\noindent$\Pi_{38}$:
\begin{ptree}
\AXC{$\stackrel{\Pi_{1}}{\vdots}$}\noLine
\UIC{$\mathrm{Cauchy}(a) \to \mathrm{Seq}(a)$}
\AXC{$\forall a\,\forall n\,(\mathrm{Seq}(a) \to n \in \mathbb{N} \to a_{n} \in \mathbb{R})$\lbl{[S1]}}
\AXC{$\varphi_{36}$\lbl{A2}}
\BIC{$\forall n\,(\mathrm{Seq}(a) \to n \in \mathbb{N} \to a_{n} \in \mathbb{R})$}
\AXC{$\varphi_{27}$\lbl{A2}}
\BIC{$\mathrm{Seq}(a) \to n \in \mathbb{N} \to a_{n} \in \mathbb{R}$}
\AXC{$\varphi_{170}$\lbl{A1}}
\BIC{$(\mathrm{Cauchy}(a) \to \mathrm{Seq}(a)) \to \varphi_{147}$}
\BIC{$\varphi_{147}$}
\end{ptree}
\noindent$\Pi_{39}$:
\begin{ptree}
\AXC{$\mathrm{Cauchy}(a) \to N \in \mathbb{N} \land \varphi_{10} \to \varphi_{18}$\lbl{A1}}
\AXC{$\stackrel{\Pi_{38}}{\vdots}$}\noLine
\UIC{$\varphi_{147}$}
\BIC{$\mathrm{Cauchy}(a) \to N \in \mathbb{N} \land \varphi_{10} \to \varphi_{43}$}
\end{ptree}
\noindent$\Pi_{40}$:
\begin{ptree}
\AXC{$\stackrel{\Pi_{39}}{\vdots}$}\noLine
\UIC{$\mathrm{Cauchy}(a) \to N \in \mathbb{N} \land \varphi_{10} \to \varphi_{43}$}
\AXC{$\forall x\,\forall y\,(x \in \mathbb{R} \to y \in \mathbb{R} \to x - y \in \mathbb{R})$\lbl{[F1]}}
\AXC{$\varphi_{66}$\lbl{A2}}
\BIC{$\forall y\,(a_{n} \in \mathbb{R} \to y \in \mathbb{R} \to a_{n} - y \in \mathbb{R})$}
\AXC{$\varphi_{87}$\lbl{A2}}
\BIC{$a_{n} \in \mathbb{R} \to a_{N} \in \mathbb{R} \to a_{n} - a_{N} \in \mathbb{R}$}
\AXC{$\varphi_{184}$\lbl{A1}}
\BIC{$(\mathrm{Cauchy}(a) \to N \in \mathbb{N} \land \varphi_{10} \to \varphi_{43}) \to \varphi_{157}$}
\BIC{$\varphi_{157}$}
\end{ptree}
\noindent$\Pi_{41}$:
\begin{ptree}
\AXC{$\stackrel{\Pi_{22}}{\vdots}$}\noLine
\UIC{$\mathrm{Cauchy}(a) \to N \in \mathbb{N} \land \varphi_{10} \to \varphi_{42}$}
\AXC{$\stackrel{\Pi_{40}}{\vdots}$}\noLine
\UIC{$\varphi_{157}$}
\BIC{$\mathrm{Cauchy}(a) \to N \in \mathbb{N} \land \varphi_{10} \to \varphi_{61}$}
\end{ptree}
\noindent$\Pi_{42}$:
\begin{ptree}
\AXC{$\forall x \in \mathbb{R}\,(\lvert x \rvert \in \mathbb{R})$\lbl{[A5]}}
\AXC{$\varphi_{52}$\lbl{A2}}
\BIC{$a_{n} - a_{N} \in \mathbb{R} \to \lvert a_{n} - a_{N} \rvert \in \mathbb{R}$}
\AXC{$(a_{n} - a_{N} \in \mathbb{R} \to \lvert a_{n} - a_{N} \rvert \in \mathbb{R}) \to \varphi_{162}$\lbl{A1}}
\BIC{$\varphi_{162}$}
\end{ptree}
\noindent$\Pi_{43}$:
\begin{ptree}
\AXC{$\varphi_{48}$\lbl{[F20]}}
\AXC{$\varphi_{48} \to \varphi_{103}$\lbl{A2}}
\BIC{$\varphi_{103}$}
\AXC{$\varphi_{103} \to \varphi_{102}$\lbl{A2}}
\BIC{$\varphi_{102}$}
\AXC{$\varphi_{102} \to \varphi_{113}$\lbl{A2}}
\BIC{$\varphi_{113}$}
\AXC{$\varphi_{113} \to \varphi_{121}$\lbl{A2}}
\BIC{$\varphi_{121}$}
\AXC{$\varphi_{121} \to \varphi_{199}$\lbl{A1}}
\BIC{$\varphi_{199}$}
\end{ptree}
\noindent$\Pi_{44}$:
\begin{ptree}
\AXC{$\stackrel{\Pi_{24}}{\vdots}$}\noLine
\UIC{$\mathrm{Cauchy}(a) \to N \in \mathbb{N} \land \varphi_{10} \to \varphi_{50}$}
\AXC{$1 \in \mathbb{R}$\lbl{(R1)}}
\AXC{$\stackrel{\Pi_{41}}{\vdots}$}\noLine
\UIC{$\mathrm{Cauchy}(a) \to N \in \mathbb{N} \land \varphi_{10} \to \varphi_{61}$}
\AXC{$\stackrel{\Pi_{42}}{\vdots}$}\noLine
\UIC{$\varphi_{162}$}
\BIC{$\mathrm{Cauchy}(a) \to N \in \mathbb{N} \land \varphi_{10} \to \varphi_{68}$}
\AXC{$\stackrel{\Pi_{43}}{\vdots}$}\noLine
\UIC{$\varphi_{199}$}
\BIC{$1 \in \mathbb{R} \to (\mathrm{Cauchy}(a) \to N \in \mathbb{N} \land \varphi_{10} \to \varphi_{50}) \to \varphi_{195}$}
\BIC{$(\mathrm{Cauchy}(a) \to N \in \mathbb{N} \land \varphi_{10} \to \varphi_{50}) \to \varphi_{195}$}
\BIC{$\varphi_{195}$}
\end{ptree}
\noindent$\Pi_{45}$:
\begin{ptree}
\AXC{$\stackrel{\Pi_{29}}{\vdots}$}\noLine
\UIC{$\mathrm{Cauchy}(a) \to N \in \mathbb{N} \land \varphi_{10} \to \varphi_{57}$}
\AXC{$\stackrel{\Pi_{37}}{\vdots}$}\noLine
\UIC{$\mathrm{Cauchy}(a) \to N \in \mathbb{N} \land \varphi_{10} \to \varphi_{63}$}
\AXC{$\stackrel{\Pi_{24}}{\vdots}$}\noLine
\UIC{$\mathrm{Cauchy}(a) \to N \in \mathbb{N} \land \varphi_{10} \to \varphi_{50}$}
\AXC{$\stackrel{\Pi_{44}}{\vdots}$}\noLine
\UIC{$\varphi_{195}$}
\BIC{$(\mathrm{Cauchy}(a) \to N \in \mathbb{N} \land \varphi_{10} \to \varphi_{63}) \to \varphi_{171}$}
\BIC{$\varphi_{171}$}
\BIC{$\mathrm{Cauchy}(a) \to N \in \mathbb{N} \land \varphi_{10} \to \varphi_{93}$}
\end{ptree}
\noindent$\Pi_{46}$:
\begin{ptree}
\AXC{$z = w \to \varphi_{64}$\lbl{A6}}
\UIC{$\forall w\,(z = w \to \varphi_{64})$}
\UIC{$\forall z\,\forall w\,(z = w \to \varphi_{64})$}
\AXC{$\forall z\,\forall w\,(z = w \to \varphi_{64}) \to \varphi_{107}$\lbl{A2}}
\BIC{$\varphi_{107}$}
\AXC{$\varphi_{107} \to a_{n} - a_{N} + a_{N} = a_{n} \to \varphi_{95}$\lbl{A2}}
\BIC{$a_{n} - a_{N} + a_{N} = a_{n} \to \varphi_{95}$}
\end{ptree}
\noindent$\Pi_{47}$:
\begin{ptree}
\AXC{$\forall a\,\forall b\,(a \in \mathbb{R} \to b \in \mathbb{R} \to \lvert a + b \rvert \leq \lvert a \rvert + \lvert b \rvert)$\lbl{命題1.2}}
\AXC{$\varphi_{106}$\lbl{A2}}
\BIC{$\varphi_{78}$}
\AXC{$\varphi_{78} \to \varphi_{89}$\lbl{A2}}
\BIC{$\varphi_{89}$}
\end{ptree}
\noindent$\Pi_{48}$:
\begin{ptree}
\AXC{$\stackrel{\Pi_{41}}{\vdots}$}\noLine
\UIC{$\mathrm{Cauchy}(a) \to N \in \mathbb{N} \land \varphi_{10} \to \varphi_{61}$}
\AXC{$\stackrel{\Pi_{47}}{\vdots}$}\noLine
\UIC{$\varphi_{89}$}
\AXC{$\varphi_{89} \to (\mathrm{Cauchy}(a) \to N \in \mathbb{N} \land \varphi_{10} \to \varphi_{61}) \to \varphi_{173}$\lbl{A1}}
\BIC{$(\mathrm{Cauchy}(a) \to N \in \mathbb{N} \land \varphi_{10} \to \varphi_{61}) \to \varphi_{173}$}
\BIC{$\varphi_{173}$}
\end{ptree}
\noindent$\Pi_{49}$:
\begin{ptree}
\AXC{$\stackrel{\Pi_{39}}{\vdots}$}\noLine
\UIC{$\mathrm{Cauchy}(a) \to N \in \mathbb{N} \land \varphi_{10} \to \varphi_{43}$}
\AXC{$\forall x\,\forall y\,(x \in \mathbb{R} \to y \in \mathbb{R} \to x - y + y = x)$\lbl{[F12]}}
\AXC{$\varphi_{70}$\lbl{A2}}
\BIC{$\forall y\,(a_{n} \in \mathbb{R} \to y \in \mathbb{R} \to a_{n} - y + y = a_{n})$}
\AXC{$\varphi_{97}$\lbl{A2}}
\BIC{$a_{n} \in \mathbb{R} \to a_{N} \in \mathbb{R} \to a_{n} - a_{N} + a_{N} = a_{n}$}
\AXC{$\varphi_{187}$\lbl{A1}}
\BIC{$(\mathrm{Cauchy}(a) \to N \in \mathbb{N} \land \varphi_{10} \to \varphi_{43}) \to \varphi_{160}$}
\BIC{$\varphi_{160}$}
\end{ptree}
\noindent$\Pi_{50}$:
\begin{ptree}
\AXC{$\stackrel{\Pi_{22}}{\vdots}$}\noLine
\UIC{$\mathrm{Cauchy}(a) \to N \in \mathbb{N} \land \varphi_{10} \to \varphi_{42}$}
\AXC{$\stackrel{\Pi_{49}}{\vdots}$}\noLine
\UIC{$\varphi_{160}$}
\BIC{$\mathrm{Cauchy}(a) \to N \in \mathbb{N} \land \varphi_{10} \to \varphi_{82}$}
\AXC{$(\mathrm{Cauchy}(a) \to N \in \mathbb{N} \land \varphi_{10} \to \varphi_{82}) \to \varphi_{196}$\lbl{A1}}
\BIC{$\varphi_{196}$}
\end{ptree}
\noindent$\Pi_{51}$:
\begin{ptree}
\AXC{$\stackrel{\Pi_{46}}{\vdots}$}\noLine
\UIC{$a_{n} - a_{N} + a_{N} = a_{n} \to \varphi_{95}$}
\AXC{$\stackrel{\Pi_{22}}{\vdots}$}\noLine
\UIC{$\mathrm{Cauchy}(a) \to N \in \mathbb{N} \land \varphi_{10} \to \varphi_{42}$}
\AXC{$\stackrel{\Pi_{48}}{\vdots}$}\noLine
\UIC{$\varphi_{173}$}
\BIC{$\mathrm{Cauchy}(a) \to \varphi_{123}$}
\AXC{$\stackrel{\Pi_{50}}{\vdots}$}\noLine
\UIC{$\varphi_{196}$}
\BIC{$(a_{n} - a_{N} + a_{N} = a_{n} \to \varphi_{95}) \to \varphi_{163}$}
\BIC{$\varphi_{163}$}
\end{ptree}
\noindent$\Pi_{52}$:
\begin{ptree}
\AXC{$z = w \to \varphi_{54}$\lbl{A6}}
\UIC{$\forall w\,(z = w \to \varphi_{54})$}
\UIC{$\forall z\,\forall w\,(z = w \to \varphi_{54})$}
\AXC{$\forall z\,\forall w\,(z = w \to \varphi_{54}) \to \varphi_{105}$\lbl{A2}}
\BIC{$\varphi_{105}$}
\AXC{$\varphi_{105} \to a_{n} - a_{N} + a_{N} = a_{n} \to \varphi_{92}$\lbl{A2}}
\BIC{$a_{n} - a_{N} + a_{N} = a_{n} \to \varphi_{92}$}
\AXC{$\stackrel{\Pi_{51}}{\vdots}$}\noLine
\UIC{$\varphi_{163}$}
\BIC{$\mathrm{Cauchy}(a) \to N \in \mathbb{N} \land \varphi_{10} \to \varphi_{91}$}
\end{ptree}
\noindent$\Pi_{53}$:
\begin{ptree}
\AXC{$\forall x\,\forall y\,\forall z\,(x \leq y \to y < z \to x < z)$\lbl{[O4]}}
\AXC{$\varphi_{31}$\lbl{A2}}
\BIC{$\forall y\,\forall z\,(\lvert a_{n} \rvert \leq y \to y < z \to \lvert a_{n} \rvert < z)$}
\AXC{$\forall y\,\forall z\,(\lvert a_{n} \rvert \leq y \to y < z \to \lvert a_{n} \rvert < z) \to \varphi_{88}$\lbl{A2}}
\BIC{$\varphi_{88}$}
\end{ptree}
\noindent$\Pi_{54}$:
\begin{ptree}
\AXC{$\stackrel{\Pi_{52}}{\vdots}$}\noLine
\UIC{$\mathrm{Cauchy}(a) \to N \in \mathbb{N} \land \varphi_{10} \to \varphi_{91}$}
\AXC{$\stackrel{\Pi_{53}}{\vdots}$}\noLine
\UIC{$\varphi_{88}$}
\AXC{$\varphi_{88} \to \lvert a_{n} \rvert \leq \lvert a_{n} - a_{N} \rvert + \lvert a_{N} \rvert \to \varphi_{55}$\lbl{A2}}
\BIC{$\lvert a_{n} \rvert \leq \lvert a_{n} - a_{N} \rvert + \lvert a_{N} \rvert \to \varphi_{55}$}
\AXC{$\varphi_{194}$\lbl{A1}}
\BIC{$(\mathrm{Cauchy}(a) \to N \in \mathbb{N} \land \varphi_{10} \to \varphi_{91}) \to \varphi_{174}$}
\BIC{$\varphi_{174}$}
\end{ptree}
\noindent$\Pi_{55}$:
\begin{ptree}
\AXC{$\forall x\,\forall y\,\forall z\,(x \leq y \to y \leq z \to x \leq z)$\lbl{[O5]}}
\AXC{$\varphi_{9}$\lbl{A2}}
\BIC{$\forall y\,\forall z\,(\lvert a_{n} \rvert \leq y \to y \leq z \to \lvert a_{n} \rvert \leq z)$}
\AXC{$\forall y\,\forall z\,(\lvert a_{n} \rvert \leq y \to y \leq z \to \lvert a_{n} \rvert \leq z) \to \varphi_{14}$\lbl{A2}}
\BIC{$\varphi_{14}$}
\end{ptree}
\noindent$\Pi_{56}$:
\begin{ptree}
\AXC{$\stackrel{\Pi_{45}}{\vdots}$}\noLine
\UIC{$\mathrm{Cauchy}(a) \to N \in \mathbb{N} \land \varphi_{10} \to \varphi_{93}$}
\AXC{$\stackrel{\Pi_{54}}{\vdots}$}\noLine
\UIC{$\varphi_{174}$}
\BIC{$\mathrm{Cauchy}(a) \to N \in \mathbb{N} \land \varphi_{10} \to \varphi_{71}$}
\AXC{$\varphi_{164}$\lbl{A1}}
\BIC{$\mathrm{Cauchy}(a) \to N \in \mathbb{N} \land \varphi_{10} \to \varphi_{67}$}
\AXC{$\stackrel{\Pi_{55}}{\vdots}$}\noLine
\UIC{$\varphi_{14}$}
\AXC{$\varphi_{14} \to \lvert a_{n} \rvert \leq 1 + \lvert a_{N} \rvert \to \varphi_{28}$\lbl{A2}}
\BIC{$\lvert a_{n} \rvert \leq 1 + \lvert a_{N} \rvert \to \varphi_{28}$}
\AXC{$\varphi_{191}$\lbl{A1}}
\BIC{$(\mathrm{Cauchy}(a) \to N \in \mathbb{N} \land \varphi_{10} \to \varphi_{67}) \to \varphi_{169}$}
\BIC{$\varphi_{169}$}
\end{ptree}
\noindent$\Pi_{57}$:
\begin{ptree}
\AXC{$\forall n \in \mathbb{N}\,(n \in \omega)$\lbl{[N7']}}
\AXC{$\forall n \in \mathbb{N}\,(n \in \omega) \to n \in \mathbb{N} \to n \in \omega$\lbl{A2}}
\BIC{$n \in \mathbb{N} \to n \in \omega$}
\AXC{$(n \in \mathbb{N} \to n \in \omega) \to \varphi_{139}$\lbl{A1}}
\BIC{$\varphi_{139}$}
\end{ptree}
\noindent$\Pi_{58}$:
\begin{ptree}
\AXC{$\forall n \in \mathbb{N}\,(n \in \omega)$\lbl{[N7']}}
\AXC{$\forall n \in \mathbb{N}\,(n \in \omega) \to N \in \mathbb{N} \to N \in \omega$\lbl{A2}}
\BIC{$N \in \mathbb{N} \to N \in \omega$}
\AXC{$(N \in \mathbb{N} \to N \in \omega) \to \varphi_{127}$\lbl{A1}}
\BIC{$\varphi_{127}$}
\end{ptree}
\noindent$\Pi_{59}$:
\begin{ptree}
\AXC{$\forall m\,\forall n\,(m \in \omega \to n \in \omega \to m \leq n \lor n < m)$\lbl{[N11]}}
\AXC{$\varphi_{41}$\lbl{A2}}
\BIC{$\forall n\,(N \in \omega \to n \in \omega \to N \leq n \lor n < N)$}
\AXC{$\varphi_{34}$\lbl{A2}}
\BIC{$N \in \omega \to n \in \omega \to N \leq n \lor n < N$}
\AXC{$\varphi_{182}$\lbl{A1}}
\BIC{$(\mathrm{Cauchy}(a) \to N \in \mathbb{N} \land \varphi_{10} \to \varphi_{24}) \to \varphi_{141}$}
\end{ptree}
\noindent$\Pi_{60}$:
\begin{ptree}
\AXC{$\mathrm{Cauchy}(a) \to N \in \mathbb{N} \land \varphi_{10} \to N \in \mathbb{N} \land \varphi_{10}$\lbl{A1}}
\AXC{$\varphi_{128}$\lbl{A1}}
\BIC{$\mathrm{Cauchy}(a) \to N \in \mathbb{N} \land \varphi_{10} \to N \in \mathbb{N}$}
\AXC{$\stackrel{\Pi_{58}}{\vdots}$}\noLine
\UIC{$\varphi_{127}$}
\BIC{$\mathrm{Cauchy}(a) \to N \in \mathbb{N} \land \varphi_{10} \to \varphi_{24}$}
\AXC{$\stackrel{\Pi_{59}}{\vdots}$}\noLine
\UIC{$(\mathrm{Cauchy}(a) \to N \in \mathbb{N} \land \varphi_{10} \to \varphi_{24}) \to \varphi_{141}$}
\BIC{$\varphi_{141}$}
\end{ptree}
\noindent$\Pi_{61}$:
\begin{ptree}
\AXC{$\mathrm{Cauchy}(a) \to N \in \mathbb{N} \land \varphi_{10} \to \varphi_{18}$\lbl{A1}}
\AXC{$\stackrel{\Pi_{57}}{\vdots}$}\noLine
\UIC{$\varphi_{139}$}
\BIC{$\mathrm{Cauchy}(a) \to N \in \mathbb{N} \land \varphi_{10} \to \varphi_{23}$}
\AXC{$\stackrel{\Pi_{60}}{\vdots}$}\noLine
\UIC{$\varphi_{141}$}
\BIC{$\mathrm{Cauchy}(a) \to N \in \mathbb{N} \land \varphi_{10} \to \varphi_{29}$}
\AXC{$\varphi_{193}$\lbl{A1}}
\BIC{$(\mathrm{Cauchy}(a) \to N \in \mathbb{N} \land \varphi_{10} \to \varphi_{77}) \to \varphi_{167}$}
\end{ptree}
\noindent$\Pi_{62}$:
\begin{ptree}
\AXC{$\mathrm{Cauchy}(a) \to N \in \mathbb{N} \land \varphi_{10} \to \varphi_{76}$\lbl{A1}}
\AXC{$\varphi_{148}$\lbl{A1}}
\BIC{$\mathrm{Cauchy}(a) \to \varphi_{99}$}
\AXC{$\stackrel{\Pi_{28}}{\vdots}$}\noLine
\UIC{$\varphi_{151}$}
\BIC{$\mathrm{Cauchy}(a) \to N \in \mathbb{N} \land \varphi_{10} \to \varphi_{83}$}
\AXC{$\stackrel{\Pi_{56}}{\vdots}$}\noLine
\UIC{$\varphi_{169}$}
\BIC{$\mathrm{Cauchy}(a) \to N \in \mathbb{N} \land \varphi_{10} \to \varphi_{77}$}
\AXC{$\stackrel{\Pi_{61}}{\vdots}$}\noLine
\UIC{$(\mathrm{Cauchy}(a) \to N \in \mathbb{N} \land \varphi_{10} \to \varphi_{77}) \to \varphi_{167}$}
\BIC{$\varphi_{167}$}
\end{ptree}
\noindent$\Pi_{63}$:
\begin{ptree}
\AXC{$\stackrel{\Pi_{3}}{\vdots}$}\noLine
\UIC{$\mathrm{Cauchy}(a) \to N \in \mathbb{N} \land \varphi_{10} \to \varphi_{13}$}
\AXC{$\stackrel{\Pi_{4}}{\vdots}$}\noLine
\UIC{$\varphi_{134}$}
\BIC{$\mathrm{Cauchy}(a) \to N \in \mathbb{N} \land \varphi_{10} \to \varphi_{11}$}
\AXC{$\stackrel{\Pi_{9}}{\vdots}$}\noLine
\UIC{$\varphi_{136}$}
\BIC{$\mathrm{Cauchy}(a) \to N \in \mathbb{N} \land \varphi_{10} \to \varphi_{15}$}
\AXC{$\stackrel{\Pi_{16}}{\vdots}$}\noLine
\UIC{$\varphi_{153}$}
\BIC{$\mathrm{Cauchy}(a) \to N \in \mathbb{N} \land \varphi_{10} \to \varphi_{59}$}
\AXC{$\stackrel{\Pi_{20}}{\vdots}$}\noLine
\UIC{$\varphi_{165}$}
\BIC{$\mathrm{Cauchy}(a) \to N \in \mathbb{N} \land \varphi_{10} \to \varphi_{81}$}
\AXC{$\stackrel{\Pi_{62}}{\vdots}$}\noLine
\UIC{$\varphi_{167}$}
\BIC{$\mathrm{Cauchy}(a) \to N \in \mathbb{N} \land \varphi_{10} \to \varphi_{65}$}
\AXC{$\varphi_{161}$\lbl{A1}}
\BIC{$\varphi_{101} \to n \in \mathbb{N} \to \lvert a_{n} \rvert \leq M + (1 + \lvert a_{N} \rvert)$}
\UIC{$\forall n\,(\varphi_{101} \to n \in \mathbb{N} \to \lvert a_{n} \rvert \leq M + (1 + \lvert a_{N} \rvert))$}
\AXC{$\varphi_{198}$\lbl{A1}}
\BIC{$\varphi_{166} \to (\varphi_{101} \to \forall n\,(\varphi_{101})) \to \mathrm{Cauchy}(a) \to N \in \mathbb{N} \land \varphi_{10} \to \varphi_{69}$}
\end{ptree}
\noindent$\Pi_{64}$:
\begin{ptree}
\AXC{$1 \in \mathbb{R}$\lbl{(R1)}}
\AXC{$\forall x\,\forall y\,(x \in \mathbb{R} \to y \in \mathbb{R} \to x + y \in \mathbb{R})$\lbl{(R1)}}
\AXC{$\varphi_{44}$\lbl{A2}}
\BIC{$\forall y\,(1 \in \mathbb{R} \to y \in \mathbb{R} \to 1 + y \in \mathbb{R})$}
\AXC{$\varphi_{72}$\lbl{A2}}
\BIC{$1 \in \mathbb{R} \to \lvert a_{N} \rvert \in \mathbb{R} \to 1 + \lvert a_{N} \rvert \in \mathbb{R}$}
\AXC{$\varphi_{172}$\lbl{A1}}
\BIC{$1 \in \mathbb{R} \to \varphi_{143}$}
\BIC{$\varphi_{143}$}
\end{ptree}
\noindent$\Pi_{65}$:
\begin{ptree}
\AXC{$\forall x\,\forall y\,(x \in \mathbb{R} \to y \in \mathbb{R} \to x + y \in \mathbb{R})$\lbl{(R1)}}
\AXC{$\varphi_{46}$\lbl{A2}}
\BIC{$\forall y\,(M \in \mathbb{R} \to y \in \mathbb{R} \to M + y \in \mathbb{R})$}
\AXC{$\varphi_{85}$\lbl{A2}}
\BIC{$M \in \mathbb{R} \to 1 + \lvert a_{N} \rvert \in \mathbb{R} \to M + (1 + \lvert a_{N} \rvert) \in \mathbb{R}$}
\AXC{$\varphi_{183}$\lbl{A1}}
\BIC{$(\mathrm{Cauchy}(a) \to N \in \mathbb{N} \land \varphi_{10} \to \varphi_{5}) \to \varphi_{149}$}
\end{ptree}
\noindent$\Pi_{66}$:
\begin{ptree}
\AXC{$\mathrm{Cauchy}(a) \to N \in \mathbb{N} \land \varphi_{10} \to \varphi_{76}$\lbl{A1}}
\AXC{$\varphi_{152}$\lbl{A1}}
\BIC{$\mathrm{Cauchy}(a) \to N \in \mathbb{N} \land \varphi_{10} \to \varphi_{5}$}
\AXC{$\stackrel{\Pi_{65}}{\vdots}$}\noLine
\UIC{$(\mathrm{Cauchy}(a) \to N \in \mathbb{N} \land \varphi_{10} \to \varphi_{5}) \to \varphi_{149}$}
\BIC{$\varphi_{149}$}
\end{ptree}
\noindent$\Pi_{67}$:
\begin{ptree}
\AXC{$\stackrel{\Pi_{7}}{\vdots}$}\noLine
\UIC{$\mathrm{Cauchy}(a) \to N \in \mathbb{N} \land \varphi_{10} \to \varphi_{21}$}
\AXC{$\stackrel{\Pi_{64}}{\vdots}$}\noLine
\UIC{$\varphi_{143}$}
\BIC{$\mathrm{Cauchy}(a) \to N \in \mathbb{N} \land \varphi_{10} \to \varphi_{33}$}
\AXC{$\stackrel{\Pi_{66}}{\vdots}$}\noLine
\UIC{$\varphi_{149}$}
\BIC{$\mathrm{Cauchy}(a) \to N \in \mathbb{N} \land \varphi_{10} \to \varphi_{38}$}
\AXC{$(\mathrm{Cauchy}(a) \to N \in \mathbb{N} \land \varphi_{10} \to \varphi_{38}) \to \varphi_{175}$\lbl{A1}}
\BIC{$\varphi_{175}$}
\end{ptree}
\noindent$\Pi_{68}$:
\begin{ptree}
\AXC{$\varphi_{101} \to \forall n\,(\varphi_{101})$\lbl{A4}}
\AXC{$\varphi_{166}$\lbl{A3}}
\AXC{$\stackrel{\Pi_{63}}{\vdots}$}\noLine
\UIC{$\varphi_{166} \to (\varphi_{101} \to \forall n\,(\varphi_{101})) \to \mathrm{Cauchy}(a) \to N \in \mathbb{N} \land \varphi_{10} \to \varphi_{69}$}
\BIC{$(\varphi_{101} \to \forall n\,(\varphi_{101})) \to \mathrm{Cauchy}(a) \to N \in \mathbb{N} \land \varphi_{10} \to \varphi_{69}$}
\BIC{$\mathrm{Cauchy}(a) \to N \in \mathbb{N} \land \varphi_{10} \to \varphi_{69}$}
\AXC{$\stackrel{\Pi_{67}}{\vdots}$}\noLine
\UIC{$\varphi_{175}$}
\BIC{$\mathrm{Cauchy}(a) \to \varphi_{122}$}
\end{ptree}
\noindent$\Pi_{69}$:
\begin{ptree}
\AXC{$\forall M\,(\neg (M \in \mathbb{R} \land \forall n \in \mathbb{N}\,(\lvert a_{n} \rvert \leq M))) \to \neg \varphi_{35}$\lbl{A2}}
\AXC{$\stackrel{\Pi_{68}}{\vdots}$}\noLine
\UIC{$\mathrm{Cauchy}(a) \to \varphi_{122}$}
\AXC{$(\mathrm{Cauchy}(a) \to \varphi_{122}) \to \varphi_{132}$\lbl{A1}}
\BIC{$\varphi_{132}$}
\BIC{$\mathrm{Cauchy}(a) \to N \in \mathbb{N} \land \varphi_{10} \to \varphi_{60}$}
\end{ptree}
\noindent$\Pi_{70}$:
\begin{ptree}
\AXC{$\stackrel{\Pi_{1}}{\vdots}$}\noLine
\UIC{$\mathrm{Cauchy}(a) \to \mathrm{Seq}(a)$}
\AXC{$\forall a\,(\varphi_{3})$\lbl{DefSeqBdd}}
\AXC{$\forall a\,(\varphi_{3}) \to \varphi_{3}$\lbl{A2}}
\BIC{$\varphi_{3}$}
\AXC{$\varphi_{3} \to (\mathrm{Cauchy}(a) \to \mathrm{Seq}(a)) \to \varphi_{146}$\lbl{A1}}
\BIC{$(\mathrm{Cauchy}(a) \to \mathrm{Seq}(a)) \to \varphi_{146}$}
\BIC{$\varphi_{146}$}
\end{ptree}
\noindent$\Pi_{71}$:
\begin{ptree}
\AXC{$\forall a\,(\mathrm{Seq}(a) \to \varphi_{12})$\lbl{[C14]}}
\AXC{$\forall a\,(\mathrm{Seq}(a) \to \varphi_{12}) \to \mathrm{Seq}(a) \to \varphi_{12}$\lbl{A2}}
\BIC{$\mathrm{Seq}(a) \to \varphi_{12}$}
\AXC{$(\mathrm{Seq}(a) \to \varphi_{12}) \to \varphi_{112}$\lbl{A1}}
\BIC{$\varphi_{112}$}
\end{ptree}
\noindent$\Pi_{72}$:
\begin{ptree}
\AXC{$\varphi_{12} \to \varphi_{7}$\lbl{A2}}
\AXC{$\stackrel{\Pi_{1}}{\vdots}$}\noLine
\UIC{$\mathrm{Cauchy}(a) \to \mathrm{Seq}(a)$}
\AXC{$\stackrel{\Pi_{71}}{\vdots}$}\noLine
\UIC{$\varphi_{112}$}
\BIC{$\mathrm{Cauchy}(a) \to N \in \mathbb{N} \land \varphi_{10} \to \varphi_{12}$}
\AXC{$\varphi_{178}$\lbl{A1}}
\BIC{$(\varphi_{12} \to \varphi_{7}) \to \mathrm{Cauchy}(a) \to N \in \mathbb{N} \land \varphi_{10} \to \varphi_{7}$}
\BIC{$\mathrm{Cauchy}(a) \to N \in \mathbb{N} \land \varphi_{10} \to \varphi_{7}$}
\AXC{$(\mathrm{Cauchy}(a) \to N \in \mathbb{N} \land \varphi_{10} \to \varphi_{7}) \to \varphi_{124}$\lbl{A1}}
\BIC{$\varphi_{124}$}
\end{ptree}
\noindent$\Pi_{73}$:
\begin{ptree}
\AXC{$\mathrm{Cauchy}(a) \to N \in \mathbb{N} \land \varphi_{10} \to N \in \mathbb{N} \land \varphi_{10}$\lbl{A1}}
\AXC{$\varphi_{128}$\lbl{A1}}
\BIC{$\mathrm{Cauchy}(a) \to N \in \mathbb{N} \land \varphi_{10} \to N \in \mathbb{N}$}
\AXC{$\stackrel{\Pi_{72}}{\vdots}$}\noLine
\UIC{$\varphi_{124}$}
\BIC{$\mathrm{Cauchy}(a) \to \varphi_{96}$}
\AXC{$(\mathrm{Cauchy}(a) \to \varphi_{96}) \to \forall M\,(\varphi_{104}) \to \varphi_{186}$\lbl{A1}}
\BIC{$\forall M\,(\varphi_{104}) \to \varphi_{186}$}
\end{ptree}
\noindent$\Pi_{74}$:
\begin{ptree}
\AXC{$\stackrel{\Pi_{69}}{\vdots}$}\noLine
\UIC{$\mathrm{Cauchy}(a) \to N \in \mathbb{N} \land \varphi_{10} \to \varphi_{60}$}
\AXC{$\stackrel{\Pi_{70}}{\vdots}$}\noLine
\UIC{$\varphi_{146}$}
\BIC{$\mathrm{Cauchy}(a) \to N \in \mathbb{N} \land \varphi_{10} \to \varphi_{4}$}
\AXC{$(\mathrm{Cauchy}(a) \to N \in \mathbb{N} \land \varphi_{10} \to \varphi_{4}) \to \varphi_{104}$\lbl{A1}}
\BIC{$\varphi_{104}$}
\UIC{$\forall M\,(\varphi_{104})$}
\AXC{$\stackrel{\Pi_{73}}{\vdots}$}\noLine
\UIC{$\forall M\,(\varphi_{104}) \to \varphi_{186}$}
\BIC{$\varphi_{186}$}
\end{ptree}
\noindent$\Pi_{75}$:
\begin{ptree}
\AXC{$\varphi_{120}$\lbl{A4}}
\AXC{$\forall M\,(\varphi_{104}) \to \varphi_{73} \to \varphi_{2}$\lbl{A3}}
\AXC{$\stackrel{\Pi_{74}}{\vdots}$}\noLine
\UIC{$\varphi_{186}$}
\BIC{$\varphi_{120} \to \mathrm{Cauchy}(a) \to N \in \mathbb{N} \land \varphi_{10} \to \mathrm{SeqBdd}(a)$}
\BIC{$\mathrm{Cauchy}(a) \to N \in \mathbb{N} \land \varphi_{10} \to \mathrm{SeqBdd}(a)$}
\AXC{$\varphi_{119}$\lbl{A1}}
\BIC{$\mathrm{Cauchy}(a) \land \neg \mathrm{SeqBdd}(a) \to \neg (N \in \mathbb{N} \land \varphi_{10})$}
\UIC{$\varphi_{75}$}
\end{ptree}
\noindent$\Pi_{76}$:
\begin{ptree}
\AXC{$\forall a\,(\varphi_{90})$\lbl{[C26]}}
\AXC{$\forall a\,(\varphi_{90}) \to \varphi_{90}$\lbl{A2}}
\BIC{$\varphi_{90}$}
\AXC{$\varphi_{90} \to \mathrm{Cauchy}(a) \to 1 \in \mathbb{R} \to 0 < 1 \to \exists N \in \mathbb{N}\,(\varphi_{10})$\lbl{A2}}
\BIC{$\mathrm{Cauchy}(a) \to 1 \in \mathbb{R} \to 0 < 1 \to \exists N \in \mathbb{N}\,(\varphi_{10})$}
\AXC{$\varphi_{125}$\lbl{A1}}
\BIC{$(\mathrm{Cauchy}(a) \to \mathrm{Cauchy}(a)) \to 1 \in \mathbb{R} \to \varphi_{84}$}
\end{ptree}
\noindent$\Pi_{77}$:
\begin{ptree}
\AXC{$1 \in \mathbb{R}$\lbl{(R1)}}
\AXC{$\mathrm{Cauchy}(a) \to \mathrm{Cauchy}(a)$\lbl{A1}}
\AXC{$\stackrel{\Pi_{76}}{\vdots}$}\noLine
\UIC{$(\mathrm{Cauchy}(a) \to \mathrm{Cauchy}(a)) \to 1 \in \mathbb{R} \to \varphi_{84}$}
\BIC{$1 \in \mathbb{R} \to \varphi_{84}$}
\BIC{$\varphi_{84}$}
\end{ptree}
\noindent$\Pi_{78}$:
\begin{ptree}
\AXC{$0 < 1$\lbl{[F28]}}
\AXC{$0 < 1 \to \mathrm{Cauchy}(a) \to 0 < 1$\lbl{A1}}
\BIC{$\mathrm{Cauchy}(a) \to 0 < 1$}
\AXC{$\stackrel{\Pi_{77}}{\vdots}$}\noLine
\UIC{$\varphi_{84}$}
\BIC{$\mathrm{Cauchy}(a) \to \exists N \in \mathbb{N}\,(\varphi_{10})$}
\AXC{$\varphi_{181}$\lbl{A1}}
\BIC{$\varphi_{75} \to (\varphi_{75} \to \varphi_{80}) \to \varphi_{1} \to \mathrm{Cauchy}(a) \to \mathrm{SeqBdd}(a)$}
\end{ptree}
\noindent 主木:
\begin{ptree}
\AXC{$\varphi_{1}$\lbl{A4}}
\AXC{$\varphi_{75} \to \varphi_{80}$\lbl{A3}}
\AXC{$\stackrel{\Pi_{75}}{\vdots}$}\noLine
\UIC{$\varphi_{75}$}
\AXC{$\stackrel{\Pi_{78}}{\vdots}$}\noLine
\UIC{$\varphi_{75} \to (\varphi_{75} \to \varphi_{80}) \to \varphi_{1} \to \mathrm{Cauchy}(a) \to \mathrm{SeqBdd}(a)$}
\BIC{$(\varphi_{75} \to \varphi_{80}) \to \varphi_{1} \to \mathrm{Cauchy}(a) \to \mathrm{SeqBdd}(a)$}
\BIC{$\varphi_{1} \to \mathrm{Cauchy}(a) \to \mathrm{SeqBdd}(a)$}
\BIC{$\mathrm{Cauchy}(a) \to \mathrm{SeqBdd}(a)$}
\UIC{$\forall a\,(\mathrm{Cauchy}(a) \to \mathrm{SeqBdd}(a))$}
\end{ptree}
\begin{abbrevlist}
\item $\varphi_{1} :\equiv \mathrm{Cauchy}(a) \land \neg \mathrm{SeqBdd}(a) \to \forall N\,(\mathrm{Cauchy}(a) \land \neg \mathrm{SeqBdd}(a))$
\item $\varphi_{2} :\equiv \forall M\,(\neg (M \in \mathbb{R} \land 0 \leq M \land \forall n \in \mathbb{N}\,(n \leq N \to \lvert a_{n} \rvert \leq M)))$
\item $\varphi_{3} :\equiv \mathrm{SeqBdd}(a) \leftrightarrow \mathrm{Seq}(a) \land \exists M \in \mathbb{R}\,\forall n \in \mathbb{N}\,(\lvert a_{n} \rvert \leq M)$
\item $\varphi_{4} :\equiv M \in \mathbb{R} \land 0 \leq M \land \forall n \in \mathbb{N}\,(n \leq N \to \lvert a_{n} \rvert \leq M) \to \mathrm{SeqBdd}(a)$
\item $\varphi_{5} :\equiv M \in \mathbb{R} \land 0 \leq M \land \forall n \in \mathbb{N}\,(n \leq N \to \lvert a_{n} \rvert \leq M) \to M \in \mathbb{R}$
\item $\varphi_{6} :\equiv v_{748} \in \mathbb{N} \to N \leq v_{747} \to N \leq v_{748} \to \lvert a_{v_{747}} - a_{v_{748}} \rvert < 1$
\item $\varphi_{7} :\equiv N \in \mathbb{N} \to \exists M\,(M \in \mathbb{R} \land 0 \leq M \land \forall n \in \mathbb{N}\,(n \leq N \to \lvert a_{n} \rvert \leq M))$
\item $\varphi_{8} :\equiv \forall n \in \mathbb{N}\,(n \leq N \to \lvert a_{n} \rvert \leq M) \to n \in \mathbb{N} \to n \leq N \to \lvert a_{n} \rvert \leq M$
\item $\varphi_{9} :\equiv \forall x\,\forall y\,\forall z\,(x \leq y \to y \leq z \to x \leq z) \to \forall y\,\forall z\,(\lvert a_{n} \rvert \leq y \to y \leq z \to \lvert a_{n} \rvert \leq z)$
\item $\varphi_{10} :\equiv \forall m\,\forall n\,(m \in \mathbb{N} \to n \in \mathbb{N} \to N \leq m \to N \leq n \to \lvert a_{m} - a_{n} \rvert < 1)$
\item $\varphi_{11} :\equiv M \in \mathbb{R} \land 0 \leq M \land \forall n \in \mathbb{N}\,(n \leq N \to \lvert a_{n} \rvert \leq M) \to 0 \leq \lvert a_{N} \rvert$
\item $\varphi_{12} :\equiv \forall N \in \mathbb{N}\,\exists M\,(M \in \mathbb{R} \land 0 \leq M \land \forall n \in \mathbb{N}\,(n \leq N \to \lvert a_{n} \rvert \leq M))$
\item $\varphi_{13} :\equiv M \in \mathbb{R} \land 0 \leq M \land \forall n \in \mathbb{N}\,(n \leq N \to \lvert a_{n} \rvert \leq M) \to a_{N} \in \mathbb{R}$
\item $\varphi_{14} :\equiv \forall z\,(\lvert a_{n} \rvert \leq 1 + \lvert a_{N} \rvert \to 1 + \lvert a_{N} \rvert \leq z \to \lvert a_{n} \rvert \leq z)$
\item $\varphi_{15} :\equiv M \in \mathbb{R} \land 0 \leq M \land \forall n \in \mathbb{N}\,(n \leq N \to \lvert a_{n} \rvert \leq M) \to 0 \leq 1 + \lvert a_{N} \rvert$
\item $\varphi_{16} :\equiv 1 + \lvert a_{N} \rvert \in \mathbb{R} \to 0 \leq 1 + \lvert a_{N} \rvert \to M \leq M + (1 + \lvert a_{N} \rvert)$
\item $\varphi_{17} :\equiv \forall n\,(v_{747} \in \mathbb{N} \to n \in \mathbb{N} \to N \leq v_{747} \to N \leq n \to \lvert a_{v_{747}} - a_{n} \rvert < 1)$
\item $\varphi_{18} :\equiv M \in \mathbb{R} \land 0 \leq M \land \forall n \in \mathbb{N}\,(n \leq N \to \lvert a_{n} \rvert \leq M) \to n \in \mathbb{N} \to n \in \mathbb{N}$
\item $\varphi_{19} :\equiv M \in \mathbb{R} \land 0 \leq M \land \forall n \in \mathbb{N}\,(n \leq N \to \lvert a_{n} \rvert \leq M) \to n \in \mathbb{N} \to N \leq n \to N \leq N$
\item $\varphi_{20} :\equiv M \in \mathbb{R} \land 0 \leq M \land \forall n \in \mathbb{N}\,(n \leq N \to \lvert a_{n} \rvert \leq M) \to n \in \mathbb{N} \to N \leq n \to N \leq n$
\item $\varphi_{21} :\equiv M \in \mathbb{R} \land 0 \leq M \land \forall n \in \mathbb{N}\,(n \leq N \to \lvert a_{n} \rvert \leq M) \to \lvert a_{N} \rvert \in \mathbb{R}$
\item $\varphi_{22} :\equiv \forall v_{748}\,(n \in \mathbb{N} \to v_{748} \in \mathbb{N} \to N \leq n \to N \leq v_{748} \to \lvert a_{n} - a_{v_{748}} \rvert < 1)$
\item $\varphi_{23} :\equiv M \in \mathbb{R} \land 0 \leq M \land \forall n \in \mathbb{N}\,(n \leq N \to \lvert a_{n} \rvert \leq M) \to n \in \mathbb{N} \to n \in \omega$
\item $\varphi_{24} :\equiv M \in \mathbb{R} \land 0 \leq M \land \forall n \in \mathbb{N}\,(n \leq N \to \lvert a_{n} \rvert \leq M) \to n \in \mathbb{N} \to N \in \omega$
\item $\varphi_{25} :\equiv \forall n\,(\mathrm{Seq}(a) \to n \in \mathbb{N} \to a_{n} \in \mathbb{R}) \to \mathrm{Seq}(a) \to N \in \mathbb{N} \to a_{N} \in \mathbb{R}$
\item $\varphi_{26} :\equiv 1 \in \mathbb{R} \to \lvert a_{N} \rvert \in \mathbb{R} \to 0 \leq 1 \to 0 \leq \lvert a_{N} \rvert \to 0 \leq 1 + \lvert a_{N} \rvert$
\item $\varphi_{27} :\equiv \forall n\,(\mathrm{Seq}(a) \to n \in \mathbb{N} \to a_{n} \in \mathbb{R}) \to \mathrm{Seq}(a) \to n \in \mathbb{N} \to a_{n} \in \mathbb{R}$
\item $\varphi_{28} :\equiv 1 + \lvert a_{N} \rvert \leq M + (1 + \lvert a_{N} \rvert) \to \lvert a_{n} \rvert \leq M + (1 + \lvert a_{N} \rvert)$
\item $\varphi_{29} :\equiv M \in \mathbb{R} \land 0 \leq M \land \forall n \in \mathbb{N}\,(n \leq N \to \lvert a_{n} \rvert \leq M) \to n \in \mathbb{N} \to N \leq n \lor n < N$
\item $\varphi_{30} :\equiv \lvert a_{n} \rvert \leq M \to M \leq M + (1 + \lvert a_{N} \rvert) \to \lvert a_{n} \rvert \leq M + (1 + \lvert a_{N} \rvert)$
\item $\varphi_{31} :\equiv \forall x\,\forall y\,\forall z\,(x \leq y \to y < z \to x < z) \to \forall y\,\forall z\,(\lvert a_{n} \rvert \leq y \to y < z \to \lvert a_{n} \rvert < z)$
\item $\varphi_{32} :\equiv M \in \mathbb{R} \land 0 \leq M \land \forall n \in \mathbb{N}\,(n \leq N \to \lvert a_{n} \rvert \leq M) \to n \in \mathbb{N} \to n < N \to n < N$
\item $\varphi_{33} :\equiv M \in \mathbb{R} \land 0 \leq M \land \forall n \in \mathbb{N}\,(n \leq N \to \lvert a_{n} \rvert \leq M) \to 1 + \lvert a_{N} \rvert \in \mathbb{R}$
\item $\varphi_{34} :\equiv \forall n\,(N \in \omega \to n \in \omega \to N \leq n \lor n < N) \to N \in \omega \to n \in \omega \to N \leq n \lor n < N$
\item $\varphi_{35} :\equiv M + (1 + \lvert a_{N} \rvert) \in \mathbb{R} \land \forall n \in \mathbb{N}\,(\lvert a_{n} \rvert \leq M + (1 + \lvert a_{N} \rvert))$
\item $\varphi_{36} :\equiv \forall a\,\forall n\,(\mathrm{Seq}(a) \to n \in \mathbb{N} \to a_{n} \in \mathbb{R}) \to \forall n\,(\mathrm{Seq}(a) \to n \in \mathbb{N} \to a_{n} \in \mathbb{R})$
\item $\varphi_{37} :\equiv \forall y\,\forall z\,(\lvert a_{n} \rvert \leq y \to y \leq z \to \lvert a_{n} \rvert \leq z) \to \forall z\,(\lvert a_{n} \rvert \leq M \to M \leq z \to \lvert a_{n} \rvert \leq z)$
\item $\varphi_{38} :\equiv M \in \mathbb{R} \land 0 \leq M \land \forall n \in \mathbb{N}\,(n \leq N \to \lvert a_{n} \rvert \leq M) \to M + (1 + \lvert a_{N} \rvert) \in \mathbb{R}$
\item $\varphi_{39} :\equiv 1 + \lvert a_{N} \rvert \in \mathbb{R} \to M \in \mathbb{R} \to 0 \leq M \to 1 + \lvert a_{N} \rvert \leq M + (1 + \lvert a_{N} \rvert)$
\item $\varphi_{40} :\equiv \forall x\,\forall y\,(x \in \mathbb{R} \to y \in \mathbb{R} \to 0 \leq y \to x \leq x + y) \to \forall y\,(M \in \mathbb{R} \to y \in \mathbb{R} \to 0 \leq y \to M \leq M + y)$
\item $\varphi_{41} :\equiv \forall m\,\forall n\,(m \in \omega \to n \in \omega \to m \leq n \lor n < m) \to \forall n\,(N \in \omega \to n \in \omega \to N \leq n \lor n < N)$
\item $\varphi_{42} :\equiv M \in \mathbb{R} \land 0 \leq M \land \forall n \in \mathbb{N}\,(n \leq N \to \lvert a_{n} \rvert \leq M) \to n \in \mathbb{N} \to N \leq n \to a_{N} \in \mathbb{R}$
\item $\varphi_{43} :\equiv M \in \mathbb{R} \land 0 \leq M \land \forall n \in \mathbb{N}\,(n \leq N \to \lvert a_{n} \rvert \leq M) \to n \in \mathbb{N} \to N \leq n \to a_{n} \in \mathbb{R}$
\item $\varphi_{44} :\equiv \forall x\,\forall y\,(x \in \mathbb{R} \to y \in \mathbb{R} \to x + y \in \mathbb{R}) \to \forall y\,(1 \in \mathbb{R} \to y \in \mathbb{R} \to 1 + y \in \mathbb{R})$
\item $\varphi_{45} :\equiv \forall y\,(1 + \lvert a_{N} \rvert \in \mathbb{R} \to y \in \mathbb{R} \to 0 \leq y \to 1 + \lvert a_{N} \rvert \leq y + (1 + \lvert a_{N} \rvert))$
\item $\varphi_{46} :\equiv \forall x\,\forall y\,(x \in \mathbb{R} \to y \in \mathbb{R} \to x + y \in \mathbb{R}) \to \forall y\,(M \in \mathbb{R} \to y \in \mathbb{R} \to M + y \in \mathbb{R})$
\item $\varphi_{47} :\equiv M \in \mathbb{R} \land 0 \leq M \land \forall n \in \mathbb{N}\,(n \leq N \to \lvert a_{n} \rvert \leq M) \to n \in \mathbb{N} \to n < N \to \lvert a_{n} \rvert \leq M$
\item $\varphi_{48} :\equiv \forall x\,\forall y\,\forall u\,\forall v\,(x \in \mathbb{R} \to y \in \mathbb{R} \to u \in \mathbb{R} \to v \in \mathbb{R} \to x < y \to u \leq v \to x + u < y + v)$
\item $\varphi_{49} :\equiv M \in \mathbb{R} \land 0 \leq M \land \forall n \in \mathbb{N}\,(n \leq N \to \lvert a_{n} \rvert \leq M) \to n \in \mathbb{N} \to n < N \to a_{N} \in \mathbb{R}$
\item $\varphi_{50} :\equiv M \in \mathbb{R} \land 0 \leq M \land \forall n \in \mathbb{N}\,(n \leq N \to \lvert a_{n} \rvert \leq M) \to n \in \mathbb{N} \to N \leq n \to \lvert a_{N} \rvert \in \mathbb{R}$
\item $\varphi_{51} :\equiv n \in \mathbb{N} \to N \leq n \to \lvert a_{n} - a_{N} + a_{N} \rvert \leq \lvert a_{n} - a_{N} \rvert + \lvert a_{N} \rvert$
\item $\varphi_{52} :\equiv \forall x \in \mathbb{R}\,(\lvert x \rvert \in \mathbb{R}) \to a_{n} - a_{N} \in \mathbb{R} \to \lvert a_{n} - a_{N} \rvert \in \mathbb{R}$
\item $\varphi_{53} :\equiv M \in \mathbb{R} \land 0 \leq M \land \forall n \in \mathbb{N}\,(n \leq N \to \lvert a_{n} \rvert \leq M) \to n \in \mathbb{N} \to n < N \to \lvert a_{N} \rvert \in \mathbb{R}$
\item $\varphi_{54} :\equiv \lvert z \rvert = \lvert a_{n} - a_{N} \rvert + \lvert a_{N} \rvert \to \lvert w \rvert = \lvert a_{n} - a_{N} \rvert + \lvert a_{N} \rvert$
\item $\varphi_{55} :\equiv \lvert a_{n} - a_{N} \rvert + \lvert a_{N} \rvert < 1 + \lvert a_{N} \rvert \to \lvert a_{n} \rvert < 1 + \lvert a_{N} \rvert$
\item $\varphi_{56} :\equiv M \in \mathbb{R} \land 0 \leq M \land \forall n \in \mathbb{N}\,(n \leq N \to \lvert a_{n} \rvert \leq M) \to n \in \mathbb{N} \to N \leq n \to 1 + \lvert a_{N} \rvert \in \mathbb{R}$
\item $\varphi_{57} :\equiv M \in \mathbb{R} \land 0 \leq M \land \forall n \in \mathbb{N}\,(n \leq N \to \lvert a_{n} \rvert \leq M) \to n \in \mathbb{N} \to N \leq n \to \lvert a_{N} \rvert \leq \lvert a_{N} \rvert$
\item $\varphi_{58} :\equiv \forall x\,\forall y\,(x \in \mathbb{R} \to y \in \mathbb{R} \to 0 \leq x \to 0 \leq y \to 0 \leq x + y) \to \forall y\,(1 \in \mathbb{R} \to y \in \mathbb{R} \to 0 \leq 1 \to 0 \leq y \to 0 \leq 1 + y)$
\item $\varphi_{59} :\equiv M \in \mathbb{R} \land 0 \leq M \land \forall n \in \mathbb{N}\,(n \leq N \to \lvert a_{n} \rvert \leq M) \to n \in \mathbb{N} \to n < N \to M \leq M + (1 + \lvert a_{N} \rvert)$
\item $\varphi_{60} :\equiv M \in \mathbb{R} \land 0 \leq M \land \forall n \in \mathbb{N}\,(n \leq N \to \lvert a_{n} \rvert \leq M) \to \exists M \in \mathbb{R}\,\forall n \in \mathbb{N}\,(\lvert a_{n} \rvert \leq M)$
\item $\varphi_{61} :\equiv M \in \mathbb{R} \land 0 \leq M \land \forall n \in \mathbb{N}\,(n \leq N \to \lvert a_{n} \rvert \leq M) \to n \in \mathbb{N} \to N \leq n \to a_{n} - a_{N} \in \mathbb{R}$
\item $\varphi_{62} :\equiv M \in \mathbb{R} \land 0 \leq M \land \forall n \in \mathbb{N}\,(n \leq N \to \lvert a_{n} \rvert \leq M) \to n \in \mathbb{N} \to n < N \to 1 + \lvert a_{N} \rvert \in \mathbb{R}$
\item $\varphi_{63} :\equiv M \in \mathbb{R} \land 0 \leq M \land \forall n \in \mathbb{N}\,(n \leq N \to \lvert a_{n} \rvert \leq M) \to n \in \mathbb{N} \to N \leq n \to \lvert a_{n} - a_{N} \rvert < 1$
\item $\varphi_{64} :\equiv \lvert z \rvert < \lvert a_{n} - a_{N} \rvert + \lvert a_{N} \rvert \to \lvert w \rvert < \lvert a_{n} - a_{N} \rvert + \lvert a_{N} \rvert$
\item $\varphi_{65} :\equiv M \in \mathbb{R} \land 0 \leq M \land \forall n \in \mathbb{N}\,(n \leq N \to \lvert a_{n} \rvert \leq M) \to n \in \mathbb{N} \to \lvert a_{n} \rvert \leq M + (1 + \lvert a_{N} \rvert)$
\item $\varphi_{66} :\equiv \forall x\,\forall y\,(x \in \mathbb{R} \to y \in \mathbb{R} \to x - y \in \mathbb{R}) \to \forall y\,(a_{n} \in \mathbb{R} \to y \in \mathbb{R} \to a_{n} - y \in \mathbb{R})$
\item $\varphi_{67} :\equiv M \in \mathbb{R} \land 0 \leq M \land \forall n \in \mathbb{N}\,(n \leq N \to \lvert a_{n} \rvert \leq M) \to n \in \mathbb{N} \to N \leq n \to \lvert a_{n} \rvert \leq 1 + \lvert a_{N} \rvert$
\item $\varphi_{68} :\equiv M \in \mathbb{R} \land 0 \leq M \land \forall n \in \mathbb{N}\,(n \leq N \to \lvert a_{n} \rvert \leq M) \to n \in \mathbb{N} \to N \leq n \to \lvert a_{n} - a_{N} \rvert \in \mathbb{R}$
\item $\varphi_{69} :\equiv M \in \mathbb{R} \land 0 \leq M \land \forall n \in \mathbb{N}\,(n \leq N \to \lvert a_{n} \rvert \leq M) \to \forall n \in \mathbb{N}\,(\lvert a_{n} \rvert \leq M + (1 + \lvert a_{N} \rvert))$
\item $\varphi_{70} :\equiv \forall x\,\forall y\,(x \in \mathbb{R} \to y \in \mathbb{R} \to x - y + y = x) \to \forall y\,(a_{n} \in \mathbb{R} \to y \in \mathbb{R} \to a_{n} - y + y = a_{n})$
\item $\varphi_{71} :\equiv M \in \mathbb{R} \land 0 \leq M \land \forall n \in \mathbb{N}\,(n \leq N \to \lvert a_{n} \rvert \leq M) \to n \in \mathbb{N} \to N \leq n \to \lvert a_{n} \rvert < 1 + \lvert a_{N} \rvert$
\item $\varphi_{72} :\equiv \forall y\,(1 \in \mathbb{R} \to y \in \mathbb{R} \to 1 + y \in \mathbb{R}) \to 1 \in \mathbb{R} \to \lvert a_{N} \rvert \in \mathbb{R} \to 1 + \lvert a_{N} \rvert \in \mathbb{R}$
\item $\varphi_{73} :\equiv \forall M\,(\mathrm{Cauchy}(a) \land (N \in \mathbb{N} \land \varphi_{10}) \land \neg \mathrm{SeqBdd}(a))$
\item $\varphi_{74} :\equiv M \in \mathbb{R} \land 0 \leq M \land \forall n \in \mathbb{N}\,(n \leq N \to \lvert a_{n} \rvert \leq M) \to n \in \mathbb{N} \to n < N \to n \in \mathbb{N} \to n \leq N \to \lvert a_{n} \rvert \leq M$
\item $\varphi_{75} :\equiv \forall N\,(\mathrm{Cauchy}(a) \land \neg \mathrm{SeqBdd}(a) \to \neg (N \in \mathbb{N} \land \varphi_{10}))$
\item $\varphi_{76} :\equiv M \in \mathbb{R} \land 0 \leq M \land \forall n \in \mathbb{N}\,(n \leq N \to \lvert a_{n} \rvert \leq M) \to M \in \mathbb{R} \land 0 \leq M \land \forall n \in \mathbb{N}\,(n \leq N \to \lvert a_{n} \rvert \leq M)$
\item $\varphi_{77} :\equiv M \in \mathbb{R} \land 0 \leq M \land \forall n \in \mathbb{N}\,(n \leq N \to \lvert a_{n} \rvert \leq M) \to n \in \mathbb{N} \to N \leq n \to \lvert a_{n} \rvert \leq M + (1 + \lvert a_{N} \rvert)$
\item $\varphi_{78} :\equiv \forall b\,(a_{n} - a_{N} \in \mathbb{R} \to b \in \mathbb{R} \to \lvert a_{n} - a_{N} + b \rvert \leq \lvert a_{n} - a_{N} \rvert + \lvert b \rvert)$
\item $\varphi_{79} :\equiv M \in \mathbb{R} \land 0 \leq M \land \forall n \in \mathbb{N}\,(n \leq N \to \lvert a_{n} \rvert \leq M) \to n \in \mathbb{N} \to n < N \to \forall n \in \mathbb{N}\,(n \leq N \to \lvert a_{n} \rvert \leq M)$
\item $\varphi_{80} :\equiv \forall N\,(\mathrm{Cauchy}(a) \land \neg \mathrm{SeqBdd}(a)) \to \forall N\,(\neg (N \in \mathbb{N} \land \varphi_{10}))$
\item $\varphi_{81} :\equiv M \in \mathbb{R} \land 0 \leq M \land \forall n \in \mathbb{N}\,(n \leq N \to \lvert a_{n} \rvert \leq M) \to n \in \mathbb{N} \to n < N \to \lvert a_{n} \rvert \leq M + (1 + \lvert a_{N} \rvert)$
\item $\varphi_{82} :\equiv M \in \mathbb{R} \land 0 \leq M \land \forall n \in \mathbb{N}\,(n \leq N \to \lvert a_{n} \rvert \leq M) \to n \in \mathbb{N} \to N \leq n \to a_{n} - a_{N} + a_{N} = a_{n}$
\item $\varphi_{83} :\equiv M \in \mathbb{R} \land 0 \leq M \land \forall n \in \mathbb{N}\,(n \leq N \to \lvert a_{n} \rvert \leq M) \to n \in \mathbb{N} \to N \leq n \to 1 + \lvert a_{N} \rvert \leq M + (1 + \lvert a_{N} \rvert)$
\item $\varphi_{84} :\equiv (\mathrm{Cauchy}(a) \to 0 < 1) \to \mathrm{Cauchy}(a) \to \exists N \in \mathbb{N}\,(\varphi_{10})$
\item $\varphi_{85} :\equiv \forall y\,(M \in \mathbb{R} \to y \in \mathbb{R} \to M + y \in \mathbb{R}) \to M \in \mathbb{R} \to 1 + \lvert a_{N} \rvert \in \mathbb{R} \to M + (1 + \lvert a_{N} \rvert) \in \mathbb{R}$
\item $\varphi_{86} :\equiv \lvert a_{n} - a_{N} \rvert < 1 \to \lvert a_{N} \rvert \leq \lvert a_{N} \rvert \to \lvert a_{n} - a_{N} \rvert + \lvert a_{N} \rvert < 1 + \lvert a_{N} \rvert$
\item $\varphi_{87} :\equiv \forall y\,(a_{n} \in \mathbb{R} \to y \in \mathbb{R} \to a_{n} - y \in \mathbb{R}) \to a_{n} \in \mathbb{R} \to a_{N} \in \mathbb{R} \to a_{n} - a_{N} \in \mathbb{R}$
\item $\varphi_{88} :\equiv \forall z\,(\lvert a_{n} \rvert \leq \lvert a_{n} - a_{N} \rvert + \lvert a_{N} \rvert \to \lvert a_{n} - a_{N} \rvert + \lvert a_{N} \rvert < z \to \lvert a_{n} \rvert < z)$
\item $\varphi_{89} :\equiv a_{n} - a_{N} \in \mathbb{R} \to a_{N} \in \mathbb{R} \to \lvert a_{n} - a_{N} + a_{N} \rvert \leq \lvert a_{n} - a_{N} \rvert + \lvert a_{N} \rvert$
\item $\varphi_{90} :\equiv \forall \varepsilon\,(\mathrm{Cauchy}(a) \to \varepsilon \in \mathbb{R} \to 0 < \varepsilon \to \exists N \in \mathbb{N}\,\forall m\,\forall n\,(m \in \mathbb{N} \to n \in \mathbb{N} \to N \leq m \to N \leq n \to \lvert a_{m} - a_{n} \rvert < \varepsilon))$
\item $\varphi_{91} :\equiv M \in \mathbb{R} \land 0 \leq M \land \forall n \in \mathbb{N}\,(n \leq N \to \lvert a_{n} \rvert \leq M) \to n \in \mathbb{N} \to N \leq n \to \lvert a_{n} \rvert \leq \lvert a_{n} - a_{N} \rvert + \lvert a_{N} \rvert$
\item $\varphi_{92} :\equiv \lvert a_{n} - a_{N} + a_{N} \rvert = \lvert a_{n} - a_{N} \rvert + \lvert a_{N} \rvert \to \lvert a_{n} \rvert = \lvert a_{n} - a_{N} \rvert + \lvert a_{N} \rvert$
\item $\varphi_{93} :\equiv M \in \mathbb{R} \land 0 \leq M \land \forall n \in \mathbb{N}\,(n \leq N \to \lvert a_{n} \rvert \leq M) \to n \in \mathbb{N} \to N \leq n \to \lvert a_{n} - a_{N} \rvert + \lvert a_{N} \rvert < 1 + \lvert a_{N} \rvert$
\item $\varphi_{94} :\equiv \mathrm{Cauchy}(a) \leftrightarrow \mathrm{Seq}(a) \land \forall \varepsilon \in \mathbb{R}\,(0 < \varepsilon \to \exists N \in \mathbb{N}\,\forall m\,\forall n\,(m \in \mathbb{N} \to n \in \mathbb{N} \to N \leq m \to N \leq n \to \lvert a_{m} - a_{n} \rvert < \varepsilon))$
\item $\varphi_{95} :\equiv \lvert a_{n} - a_{N} + a_{N} \rvert < \lvert a_{n} - a_{N} \rvert + \lvert a_{N} \rvert \to \lvert a_{n} \rvert < \lvert a_{n} - a_{N} \rvert + \lvert a_{N} \rvert$
\item $\varphi_{96} :\equiv N \in \mathbb{N} \land \varphi_{10} \to \exists M\,(M \in \mathbb{R} \land 0 \leq M \land \forall n \in \mathbb{N}\,(n \leq N \to \lvert a_{n} \rvert \leq M))$
\item $\varphi_{97} :\equiv \forall y\,(a_{n} \in \mathbb{R} \to y \in \mathbb{R} \to a_{n} - y + y = a_{n}) \to a_{n} \in \mathbb{R} \to a_{N} \in \mathbb{R} \to a_{n} - a_{N} + a_{N} = a_{n}$
\item $\varphi_{98} :\equiv N \in \mathbb{N} \land \varphi_{10} \to M \in \mathbb{R} \land 0 \leq M \land \forall n \in \mathbb{N}\,(n \leq N \to \lvert a_{n} \rvert \leq M) \to 0 \leq 1$
\item $\varphi_{99} :\equiv N \in \mathbb{N} \land \varphi_{10} \to M \in \mathbb{R} \land 0 \leq M \land \forall n \in \mathbb{N}\,(n \leq N \to \lvert a_{n} \rvert \leq M) \to 0 \leq M$
\item $\varphi_{100} :\equiv N \in \mathbb{N} \land \varphi_{10} \to n \in \mathbb{N} \to N \in \mathbb{N} \to N \leq n \to N \leq N \to \lvert a_{n} - a_{N} \rvert < 1$
\item $\varphi_{101} :\equiv \mathrm{Cauchy}(a) \land (N \in \mathbb{N} \land \varphi_{10}) \land (M \in \mathbb{R} \land 0 \leq M \land \forall n \in \mathbb{N}\,(n \leq N \to \lvert a_{n} \rvert \leq M))$
\item $\varphi_{102} :\equiv \forall u\,\forall v\,(\lvert a_{n} - a_{N} \rvert \in \mathbb{R} \to 1 \in \mathbb{R} \to u \in \mathbb{R} \to v \in \mathbb{R} \to \lvert a_{n} - a_{N} \rvert < 1 \to u \leq v \to \lvert a_{n} - a_{N} \rvert + u < 1 + v)$
\item $\varphi_{103} :\equiv \forall y\,\forall u\,\forall v\,(\lvert a_{n} - a_{N} \rvert \in \mathbb{R} \to y \in \mathbb{R} \to u \in \mathbb{R} \to v \in \mathbb{R} \to \lvert a_{n} - a_{N} \rvert < y \to u \leq v \to \lvert a_{n} - a_{N} \rvert + u < y + v)$
\item $\varphi_{104} :\equiv \mathrm{Cauchy}(a) \land (N \in \mathbb{N} \land \varphi_{10}) \land \neg \mathrm{SeqBdd}(a) \to \neg (M \in \mathbb{R} \land 0 \leq M \land \forall n \in \mathbb{N}\,(n \leq N \to \lvert a_{n} \rvert \leq M))$
\item $\varphi_{105} :\equiv \forall w\,(a_{n} - a_{N} + a_{N} = w \to \lvert a_{n} - a_{N} + a_{N} \rvert = \lvert a_{n} - a_{N} \rvert + \lvert a_{N} \rvert \to \lvert w \rvert = \lvert a_{n} - a_{N} \rvert + \lvert a_{N} \rvert)$
\item $\varphi_{106} :\equiv \forall a\,\forall b\,(a \in \mathbb{R} \to b \in \mathbb{R} \to \lvert a + b \rvert \leq \lvert a \rvert + \lvert b \rvert) \to \varphi_{78}$
\item $\varphi_{107} :\equiv \forall w\,(a_{n} - a_{N} + a_{N} = w \to \lvert a_{n} - a_{N} + a_{N} \rvert < \lvert a_{n} - a_{N} \rvert + \lvert a_{N} \rvert \to \lvert w \rvert < \lvert a_{n} - a_{N} \rvert + \lvert a_{N} \rvert)$
\item $\varphi_{108} :\equiv (\mathrm{Cauchy}(a) \to \mathrm{Cauchy}(a)) \to \varphi_{94} \to \mathrm{Cauchy}(a) \to \mathrm{Seq}(a)$
\item $\varphi_{109} :\equiv \forall v_{748}\,(\mathrm{Cauchy}(a) \land (N \in \mathbb{N} \land \varphi_{10}) \to v_{747} \in \mathbb{N} \to \varphi_{6})$
\item $\varphi_{110} :\equiv \mathrm{Cauchy}(a) \to N \in \mathbb{N} \land \varphi_{10} \to \forall v_{748}\,(v_{747} \in \mathbb{N} \to \varphi_{6})$
\item $\varphi_{111} :\equiv \mathrm{Cauchy}(a) \land (N \in \mathbb{N} \land \varphi_{10}) \to \forall v_{748}\,(v_{747} \in \mathbb{N} \to \varphi_{6})$
\item $\varphi_{112} :\equiv (\mathrm{Cauchy}(a) \to \mathrm{Seq}(a)) \to \mathrm{Cauchy}(a) \to N \in \mathbb{N} \land \varphi_{10} \to \varphi_{12}$
\item $\varphi_{113} :\equiv \forall v\,(\lvert a_{n} - a_{N} \rvert \in \mathbb{R} \to 1 \in \mathbb{R} \to \lvert a_{N} \rvert \in \mathbb{R} \to v \in \mathbb{R} \to \lvert a_{n} - a_{N} \rvert < 1 \to \lvert a_{N} \rvert \leq v \to \lvert a_{n} - a_{N} \rvert + \lvert a_{N} \rvert < 1 + v)$
\item $\varphi_{114} :\equiv \mathrm{Cauchy}(a) \to N \in \mathbb{N} \land \varphi_{10} \to \forall v_{747}\,\forall v_{748}\,(v_{747} \in \mathbb{N} \to \varphi_{6})$
\item $\varphi_{115} :\equiv \forall v_{748}\,(\mathrm{Cauchy}(a) \land (N \in \mathbb{N} \land \varphi_{10})) \to \forall v_{748}\,(v_{747} \in \mathbb{N} \to \varphi_{6})$
\item $\varphi_{116} :\equiv \mathrm{Cauchy}(a) \land (N \in \mathbb{N} \land \varphi_{10}) \to \forall v_{747}\,(\mathrm{Cauchy}(a) \land (N \in \mathbb{N} \land \varphi_{10}))$
\item $\varphi_{117} :\equiv \mathrm{Cauchy}(a) \land (N \in \mathbb{N} \land \varphi_{10}) \to \forall v_{748}\,(\mathrm{Cauchy}(a) \land (N \in \mathbb{N} \land \varphi_{10}))$
\item $\varphi_{118} :\equiv \forall v_{747}\,(\mathrm{Cauchy}(a) \land (N \in \mathbb{N} \land \varphi_{10})) \to \forall v_{747}\,\forall v_{748}\,(v_{747} \in \mathbb{N} \to \varphi_{6})$
\item $\varphi_{119} :\equiv (\mathrm{Cauchy}(a) \to N \in \mathbb{N} \land \varphi_{10} \to \mathrm{SeqBdd}(a)) \to \mathrm{Cauchy}(a) \land \neg \mathrm{SeqBdd}(a) \to \neg (N \in \mathbb{N} \land \varphi_{10})$
\item $\varphi_{120} :\equiv \mathrm{Cauchy}(a) \land (N \in \mathbb{N} \land \varphi_{10}) \land \neg \mathrm{SeqBdd}(a) \to \varphi_{73}$
\item $\varphi_{121} :\equiv \lvert a_{n} - a_{N} \rvert \in \mathbb{R} \to 1 \in \mathbb{R} \to \lvert a_{N} \rvert \in \mathbb{R} \to \lvert a_{N} \rvert \in \mathbb{R} \to \varphi_{86}$
\item $\varphi_{122} :\equiv N \in \mathbb{N} \land \varphi_{10} \to M \in \mathbb{R} \land 0 \leq M \land \forall n \in \mathbb{N}\,(n \leq N \to \lvert a_{n} \rvert \leq M) \to \varphi_{35}$
\item $\varphi_{123} :\equiv N \in \mathbb{N} \land \varphi_{10} \to M \in \mathbb{R} \land 0 \leq M \land \forall n \in \mathbb{N}\,(n \leq N \to \lvert a_{n} \rvert \leq M) \to \varphi_{51}$
\item $\varphi_{124} :\equiv (\mathrm{Cauchy}(a) \to N \in \mathbb{N} \land \varphi_{10} \to N \in \mathbb{N}) \to \mathrm{Cauchy}(a) \to \varphi_{96}$
\item $\varphi_{125} :\equiv (\mathrm{Cauchy}(a) \to 1 \in \mathbb{R} \to 0 < 1 \to \exists N \in \mathbb{N}\,(\varphi_{10})) \to (\mathrm{Cauchy}(a) \to \mathrm{Cauchy}(a)) \to 1 \in \mathbb{R} \to \varphi_{84}$
\item $\varphi_{126} :\equiv (\mathrm{Cauchy}(a) \to N \in \mathbb{N} \land \varphi_{10} \to N \in \mathbb{N}) \to \mathrm{Cauchy}(a) \to N \in \mathbb{N} \land \varphi_{10} \to \varphi_{13}$
\item $\varphi_{127} :\equiv (\mathrm{Cauchy}(a) \to N \in \mathbb{N} \land \varphi_{10} \to N \in \mathbb{N}) \to \mathrm{Cauchy}(a) \to N \in \mathbb{N} \land \varphi_{10} \to \varphi_{24}$
\item $\varphi_{128} :\equiv (\mathrm{Cauchy}(a) \to N \in \mathbb{N} \land \varphi_{10} \to N \in \mathbb{N} \land \varphi_{10}) \to \mathrm{Cauchy}(a) \to N \in \mathbb{N} \land \varphi_{10} \to N \in \mathbb{N}$
\item $\varphi_{129} :\equiv (\mathrm{Cauchy}(a) \to N \in \mathbb{N} \land \varphi_{10} \to N \in \mathbb{N}) \to \mathrm{Cauchy}(a) \to N \in \mathbb{N} \land \varphi_{10} \to \varphi_{42}$
\item $\varphi_{130} :\equiv (\mathrm{Cauchy}(a) \to N \in \mathbb{N} \land \varphi_{10} \to N \in \mathbb{N}) \to \mathrm{Cauchy}(a) \to N \in \mathbb{N} \land \varphi_{10} \to \varphi_{49}$
\item $\varphi_{131} :\equiv (\varphi_{22} \to n \in \mathbb{N} \to N \in \mathbb{N} \to N \leq n \to N \leq N \to \lvert a_{n} - a_{N} \rvert < 1) \to \mathrm{Cauchy}(a) \to \varphi_{100}$
\item $\varphi_{132} :\equiv (\forall M\,(\neg (M \in \mathbb{R} \land \forall n \in \mathbb{N}\,(\lvert a_{n} \rvert \leq M))) \to \neg \varphi_{35}) \to \mathrm{Cauchy}(a) \to N \in \mathbb{N} \land \varphi_{10} \to \varphi_{60}$
\item $\varphi_{133} :\equiv (\mathrm{Seq}(a) \to N \in \mathbb{N} \to a_{N} \in \mathbb{R}) \to (\mathrm{Cauchy}(a) \to \mathrm{Seq}(a)) \to \varphi_{126}$
\item $\varphi_{134} :\equiv (\mathrm{Cauchy}(a) \to N \in \mathbb{N} \land \varphi_{10} \to \varphi_{13}) \to \mathrm{Cauchy}(a) \to N \in \mathbb{N} \land \varphi_{10} \to \varphi_{11}$
\item $\varphi_{135} :\equiv (\varphi_{17} \to v_{747} \in \mathbb{N} \to \varphi_{6}) \to \mathrm{Cauchy}(a) \to N \in \mathbb{N} \land \varphi_{10} \to v_{747} \in \mathbb{N} \to \varphi_{6}$
\item $\varphi_{136} :\equiv (\mathrm{Cauchy}(a) \to N \in \mathbb{N} \land \varphi_{10} \to \varphi_{11}) \to \mathrm{Cauchy}(a) \to N \in \mathbb{N} \land \varphi_{10} \to \varphi_{15}$
\item $\varphi_{137} :\equiv (\mathrm{Cauchy}(a) \to N \in \mathbb{N} \land \varphi_{10} \to \varphi_{13}) \to \mathrm{Cauchy}(a) \to N \in \mathbb{N} \land \varphi_{10} \to \varphi_{21}$
\item $\varphi_{138} :\equiv (\mathrm{Cauchy}(a) \to N \in \mathbb{N} \land \varphi_{10} \to N \in \mathbb{N} \land \varphi_{10}) \to \mathrm{Cauchy}(a) \to N \in \mathbb{N} \land \varphi_{10} \to \varphi_{10}$
\item $\varphi_{139} :\equiv (\mathrm{Cauchy}(a) \to N \in \mathbb{N} \land \varphi_{10} \to \varphi_{18}) \to \mathrm{Cauchy}(a) \to N \in \mathbb{N} \land \varphi_{10} \to \varphi_{23}$
\item $\varphi_{140} :\equiv (\forall v_{747}\,\forall v_{748}\,(v_{747} \in \mathbb{N} \to \varphi_{6}) \to \varphi_{22}) \to \mathrm{Cauchy}(a) \to N \in \mathbb{N} \land \varphi_{10} \to \varphi_{22}$
\item $\varphi_{141} :\equiv (\mathrm{Cauchy}(a) \to N \in \mathbb{N} \land \varphi_{10} \to \varphi_{23}) \to \mathrm{Cauchy}(a) \to N \in \mathbb{N} \land \varphi_{10} \to \varphi_{29}$
\item $\varphi_{142} :\equiv (\mathrm{Seq}(a) \to N \in \mathbb{N} \to a_{N} \in \mathbb{R}) \to (\mathrm{Cauchy}(a) \to \mathrm{Seq}(a)) \to \varphi_{129}$
\item $\varphi_{143} :\equiv (\mathrm{Cauchy}(a) \to N \in \mathbb{N} \land \varphi_{10} \to \varphi_{21}) \to \mathrm{Cauchy}(a) \to N \in \mathbb{N} \land \varphi_{10} \to \varphi_{33}$
\item $\varphi_{144} :\equiv (\mathrm{Seq}(a) \to N \in \mathbb{N} \to a_{N} \in \mathbb{R}) \to (\mathrm{Cauchy}(a) \to \mathrm{Seq}(a)) \to \varphi_{130}$
\item $\varphi_{145} :\equiv (\mathrm{Cauchy}(a) \to N \in \mathbb{N} \land \varphi_{10} \to v_{747} \in \mathbb{N} \to \varphi_{6}) \to \mathrm{Cauchy}(a) \land (N \in \mathbb{N} \land \varphi_{10}) \to v_{747} \in \mathbb{N} \to \varphi_{6}$
\item $\varphi_{146} :\equiv (\mathrm{Cauchy}(a) \to N \in \mathbb{N} \land \varphi_{10} \to \varphi_{60}) \to \mathrm{Cauchy}(a) \to N \in \mathbb{N} \land \varphi_{10} \to \varphi_{4}$
\item $\varphi_{147} :\equiv (\mathrm{Cauchy}(a) \to N \in \mathbb{N} \land \varphi_{10} \to \varphi_{18}) \to \mathrm{Cauchy}(a) \to N \in \mathbb{N} \land \varphi_{10} \to \varphi_{43}$
\item $\varphi_{148} :\equiv (\mathrm{Cauchy}(a) \to N \in \mathbb{N} \land \varphi_{10} \to \varphi_{76}) \to \mathrm{Cauchy}(a) \to \varphi_{99}$
\item $\varphi_{149} :\equiv (\mathrm{Cauchy}(a) \to N \in \mathbb{N} \land \varphi_{10} \to \varphi_{33}) \to \mathrm{Cauchy}(a) \to N \in \mathbb{N} \land \varphi_{10} \to \varphi_{38}$
\item $\varphi_{150} :\equiv (\mathrm{Cauchy}(a) \to N \in \mathbb{N} \land \varphi_{10} \to \varphi_{32}) \to \mathrm{Cauchy}(a) \to N \in \mathbb{N} \land \varphi_{10} \to \varphi_{47}$
\item $\varphi_{151} :\equiv (\mathrm{Cauchy}(a) \to \varphi_{99}) \to \mathrm{Cauchy}(a) \to N \in \mathbb{N} \land \varphi_{10} \to \varphi_{83}$
\item $\varphi_{152} :\equiv (\mathrm{Cauchy}(a) \to N \in \mathbb{N} \land \varphi_{10} \to \varphi_{76}) \to \mathrm{Cauchy}(a) \to N \in \mathbb{N} \land \varphi_{10} \to \varphi_{5}$
\item $\varphi_{153} :\equiv (\mathrm{Cauchy}(a) \to N \in \mathbb{N} \land \varphi_{10} \to \varphi_{15}) \to \mathrm{Cauchy}(a) \to N \in \mathbb{N} \land \varphi_{10} \to \varphi_{59}$
\item $\varphi_{154} :\equiv (\mathrm{Cauchy}(a) \to N \in \mathbb{N} \land \varphi_{10} \to \varphi_{19}) \to \mathrm{Cauchy}(a) \to N \in \mathbb{N} \land \varphi_{10} \to \varphi_{63}$
\item $\varphi_{155} :\equiv (\mathrm{Cauchy}(a) \to N \in \mathbb{N} \land \varphi_{10} \to \varphi_{42}) \to \mathrm{Cauchy}(a) \to N \in \mathbb{N} \land \varphi_{10} \to \varphi_{50}$
\item $\varphi_{156} :\equiv (\mathrm{Cauchy}(a) \to N \in \mathbb{N} \land \varphi_{10} \to \varphi_{49}) \to \mathrm{Cauchy}(a) \to N \in \mathbb{N} \land \varphi_{10} \to \varphi_{53}$
\item $\varphi_{157} :\equiv (\mathrm{Cauchy}(a) \to N \in \mathbb{N} \land \varphi_{10} \to \varphi_{42}) \to \mathrm{Cauchy}(a) \to N \in \mathbb{N} \land \varphi_{10} \to \varphi_{61}$
\item $\varphi_{158} :\equiv (\mathrm{Cauchy}(a) \to N \in \mathbb{N} \land \varphi_{10} \to \varphi_{50}) \to \mathrm{Cauchy}(a) \to N \in \mathbb{N} \land \varphi_{10} \to \varphi_{56}$
\item $\varphi_{159} :\equiv (\mathrm{Cauchy}(a) \to N \in \mathbb{N} \land \varphi_{10} \to \varphi_{53}) \to \mathrm{Cauchy}(a) \to N \in \mathbb{N} \land \varphi_{10} \to \varphi_{62}$
\item $\varphi_{160} :\equiv (\mathrm{Cauchy}(a) \to N \in \mathbb{N} \land \varphi_{10} \to \varphi_{42}) \to \mathrm{Cauchy}(a) \to N \in \mathbb{N} \land \varphi_{10} \to \varphi_{82}$
\item $\varphi_{161} :\equiv (\mathrm{Cauchy}(a) \to N \in \mathbb{N} \land \varphi_{10} \to \varphi_{65}) \to \varphi_{101} \to n \in \mathbb{N} \to \lvert a_{n} \rvert \leq M + (1 + \lvert a_{N} \rvert)$
\item $\varphi_{162} :\equiv (\mathrm{Cauchy}(a) \to N \in \mathbb{N} \land \varphi_{10} \to \varphi_{61}) \to \mathrm{Cauchy}(a) \to N \in \mathbb{N} \land \varphi_{10} \to \varphi_{68}$
\item $\varphi_{163} :\equiv (a_{n} - a_{N} + a_{N} = a_{n} \to \varphi_{92}) \to \mathrm{Cauchy}(a) \to N \in \mathbb{N} \land \varphi_{10} \to \varphi_{91}$
\item $\varphi_{164} :\equiv (\mathrm{Cauchy}(a) \to N \in \mathbb{N} \land \varphi_{10} \to \varphi_{71}) \to \mathrm{Cauchy}(a) \to N \in \mathbb{N} \land \varphi_{10} \to \varphi_{67}$
\item $\varphi_{165} :\equiv (\mathrm{Cauchy}(a) \to N \in \mathbb{N} \land \varphi_{10} \to \varphi_{59}) \to \mathrm{Cauchy}(a) \to N \in \mathbb{N} \land \varphi_{10} \to \varphi_{81}$
\item $\varphi_{166} :\equiv \forall n\,(\varphi_{101} \to n \in \mathbb{N} \to \lvert a_{n} \rvert \leq M + (1 + \lvert a_{N} \rvert)) \to \forall n\,(\varphi_{101}) \to \forall n \in \mathbb{N}\,(\lvert a_{n} \rvert \leq M + (1 + \lvert a_{N} \rvert))$
\item $\varphi_{167} :\equiv (\mathrm{Cauchy}(a) \to N \in \mathbb{N} \land \varphi_{10} \to \varphi_{81}) \to \mathrm{Cauchy}(a) \to N \in \mathbb{N} \land \varphi_{10} \to \varphi_{65}$
\item $\varphi_{168} :\equiv (\mathrm{Cauchy}(a) \to N \in \mathbb{N} \land \varphi_{10} \to \varphi_{76}) \to \mathrm{Cauchy}(a) \to N \in \mathbb{N} \land \varphi_{10} \to \varphi_{79}$
\item $\varphi_{169} :\equiv (\mathrm{Cauchy}(a) \to N \in \mathbb{N} \land \varphi_{10} \to \varphi_{83}) \to \mathrm{Cauchy}(a) \to N \in \mathbb{N} \land \varphi_{10} \to \varphi_{77}$
\item $\varphi_{170} :\equiv (\mathrm{Seq}(a) \to n \in \mathbb{N} \to a_{n} \in \mathbb{R}) \to (\mathrm{Cauchy}(a) \to \mathrm{Seq}(a)) \to \varphi_{147}$
\item $\varphi_{171} :\equiv (\mathrm{Cauchy}(a) \to N \in \mathbb{N} \land \varphi_{10} \to \varphi_{57}) \to \mathrm{Cauchy}(a) \to N \in \mathbb{N} \land \varphi_{10} \to \varphi_{93}$
\item $\varphi_{172} :\equiv (1 \in \mathbb{R} \to \lvert a_{N} \rvert \in \mathbb{R} \to 1 + \lvert a_{N} \rvert \in \mathbb{R}) \to 1 \in \mathbb{R} \to \varphi_{143}$
\item $\varphi_{173} :\equiv (\mathrm{Cauchy}(a) \to N \in \mathbb{N} \land \varphi_{10} \to \varphi_{42}) \to \mathrm{Cauchy}(a) \to \varphi_{123}$
\item $\varphi_{174} :\equiv (\mathrm{Cauchy}(a) \to N \in \mathbb{N} \land \varphi_{10} \to \varphi_{93}) \to \mathrm{Cauchy}(a) \to N \in \mathbb{N} \land \varphi_{10} \to \varphi_{71}$
\item $\varphi_{175} :\equiv (\mathrm{Cauchy}(a) \to N \in \mathbb{N} \land \varphi_{10} \to \varphi_{69}) \to \mathrm{Cauchy}(a) \to \varphi_{122}$
\item $\varphi_{176} :\equiv (1 \in \mathbb{R} \to \lvert a_{N} \rvert \in \mathbb{R} \to 1 + \lvert a_{N} \rvert \in \mathbb{R}) \to 1 \in \mathbb{R} \to \varphi_{158}$
\item $\varphi_{177} :\equiv (1 \in \mathbb{R} \to \lvert a_{N} \rvert \in \mathbb{R} \to 1 + \lvert a_{N} \rvert \in \mathbb{R}) \to 1 \in \mathbb{R} \to \varphi_{159}$
\item $\varphi_{178} :\equiv (\mathrm{Cauchy}(a) \to N \in \mathbb{N} \land \varphi_{10} \to \varphi_{12}) \to (\varphi_{12} \to \varphi_{7}) \to \mathrm{Cauchy}(a) \to N \in \mathbb{N} \land \varphi_{10} \to \varphi_{7}$
\item $\varphi_{179} :\equiv (\mathrm{Cauchy}(a) \to N \in \mathbb{N} \land \varphi_{10} \to \varphi_{79}) \to \varphi_{8} \to \mathrm{Cauchy}(a) \to N \in \mathbb{N} \land \varphi_{10} \to \varphi_{74}$
\item $\varphi_{180} :\equiv (\mathrm{Cauchy}(a) \to N \in \mathbb{N} \land \varphi_{10} \to \varphi_{10}) \to (\varphi_{10} \to \varphi_{17}) \to \mathrm{Cauchy}(a) \to N \in \mathbb{N} \land \varphi_{10} \to \varphi_{17}$
\item $\varphi_{181} :\equiv (\mathrm{Cauchy}(a) \to \exists N \in \mathbb{N}\,(\varphi_{10})) \to \varphi_{75} \to (\varphi_{75} \to \varphi_{80}) \to \varphi_{1} \to \mathrm{Cauchy}(a) \to \mathrm{SeqBdd}(a)$
\item $\varphi_{182} :\equiv (N \in \omega \to n \in \omega \to N \leq n \lor n < N) \to (\mathrm{Cauchy}(a) \to N \in \mathbb{N} \land \varphi_{10} \to \varphi_{24}) \to \varphi_{141}$
\item $\varphi_{183} :\equiv (M \in \mathbb{R} \to 1 + \lvert a_{N} \rvert \in \mathbb{R} \to M + (1 + \lvert a_{N} \rvert) \in \mathbb{R}) \to (\mathrm{Cauchy}(a) \to N \in \mathbb{N} \land \varphi_{10} \to \varphi_{5}) \to \varphi_{149}$
\item $\varphi_{184} :\equiv (a_{n} \in \mathbb{R} \to a_{N} \in \mathbb{R} \to a_{n} - a_{N} \in \mathbb{R}) \to (\mathrm{Cauchy}(a) \to N \in \mathbb{N} \land \varphi_{10} \to \varphi_{43}) \to \varphi_{157}$
\item $\varphi_{185} :\equiv (\mathrm{Cauchy}(a) \to N \in \mathbb{N} \land \varphi_{10} \to N \in \mathbb{N}) \to (\mathrm{Cauchy}(a) \to N \in \mathbb{N} \land \varphi_{10} \to \varphi_{20}) \to \varphi_{154}$
\item $\varphi_{186} :\equiv (\forall M\,(\varphi_{104}) \to \varphi_{73} \to \varphi_{2}) \to \varphi_{120} \to \mathrm{Cauchy}(a) \to N \in \mathbb{N} \land \varphi_{10} \to \mathrm{SeqBdd}(a)$
\item $\varphi_{187} :\equiv (a_{n} \in \mathbb{R} \to a_{N} \in \mathbb{R} \to a_{n} - a_{N} + a_{N} = a_{n}) \to (\mathrm{Cauchy}(a) \to N \in \mathbb{N} \land \varphi_{10} \to \varphi_{43}) \to \varphi_{160}$
\item $\varphi_{188} :\equiv (\mathrm{Cauchy}(a) \to N \in \mathbb{N} \land \varphi_{10} \to \varphi_{21}) \to (\mathrm{Cauchy}(a) \to \varphi_{98}) \to \varphi_{136}$
\item $\varphi_{189} :\equiv (\mathrm{Cauchy}(a) \to N \in \mathbb{N} \land \varphi_{10} \to \varphi_{56}) \to (\mathrm{Cauchy}(a) \to N \in \mathbb{N} \land \varphi_{10} \to \varphi_{5}) \to \varphi_{151}$
\item $\varphi_{190} :\equiv (\mathrm{Cauchy}(a) \to N \in \mathbb{N} \land \varphi_{10} \to \varphi_{5}) \to (\mathrm{Cauchy}(a) \to N \in \mathbb{N} \land \varphi_{10} \to \varphi_{62}) \to \varphi_{153}$
\item $\varphi_{191} :\equiv (\lvert a_{n} \rvert \leq 1 + \lvert a_{N} \rvert \to \varphi_{28}) \to (\mathrm{Cauchy}(a) \to N \in \mathbb{N} \land \varphi_{10} \to \varphi_{67}) \to \varphi_{169}$
\item $\varphi_{192} :\equiv (\mathrm{Cauchy}(a) \to N \in \mathbb{N} \land \varphi_{10} \to \varphi_{74}) \to (\mathrm{Cauchy}(a) \to N \in \mathbb{N} \land \varphi_{10} \to \varphi_{18}) \to \varphi_{150}$
\item $\varphi_{193} :\equiv (\mathrm{Cauchy}(a) \to N \in \mathbb{N} \land \varphi_{10} \to \varphi_{29}) \to (\mathrm{Cauchy}(a) \to N \in \mathbb{N} \land \varphi_{10} \to \varphi_{77}) \to \varphi_{167}$
\item $\varphi_{194} :\equiv (\lvert a_{n} \rvert \leq \lvert a_{n} - a_{N} \rvert + \lvert a_{N} \rvert \to \varphi_{55}) \to (\mathrm{Cauchy}(a) \to N \in \mathbb{N} \land \varphi_{10} \to \varphi_{91}) \to \varphi_{174}$
\item $\varphi_{195} :\equiv (\mathrm{Cauchy}(a) \to N \in \mathbb{N} \land \varphi_{10} \to \varphi_{50}) \to (\mathrm{Cauchy}(a) \to N \in \mathbb{N} \land \varphi_{10} \to \varphi_{63}) \to \varphi_{171}$
\item $\varphi_{196} :\equiv (\mathrm{Cauchy}(a) \to \varphi_{123}) \to (a_{n} - a_{N} + a_{N} = a_{n} \to \varphi_{95}) \to \varphi_{163}$
\item $\varphi_{197} :\equiv (\mathrm{Cauchy}(a) \to \varphi_{100}) \to (\mathrm{Cauchy}(a) \to N \in \mathbb{N} \land \varphi_{10} \to \varphi_{18}) \to \varphi_{185}$
\item $\varphi_{198} :\equiv \forall n\,(\varphi_{101} \to n \in \mathbb{N} \to \lvert a_{n} \rvert \leq M + (1 + \lvert a_{N} \rvert)) \to \varphi_{166} \to (\varphi_{101} \to \forall n\,(\varphi_{101})) \to \mathrm{Cauchy}(a) \to N \in \mathbb{N} \land \varphi_{10} \to \varphi_{69}$
\item $\varphi_{199} :\equiv (\mathrm{Cauchy}(a) \to N \in \mathbb{N} \land \varphi_{10} \to \varphi_{68}) \to 1 \in \mathbb{R} \to (\mathrm{Cauchy}(a) \to N \in \mathbb{N} \land \varphi_{10} \to \varphi_{50}) \to \varphi_{195}$
\end{abbrevlist}
\end{fullproof}

% 自動生成 (prover/build.js): CauchyComplete  ノード数 657
\begin{fullproof}{定理6.11}{$\forall a\,(\mathrm{Cgt}(a) \leftrightarrow \mathrm{Cauchy}(a))$}{657}\phantomsection\label{tree:CauchyComplete}
\noindent$\Pi_{1}$:
\begin{ptree}
\AXC{$z = w \to \lvert a_{n} - c \rvert < z \to \lvert a_{n} - c \rvert < w$\lbl{A6}}
\UIC{$\forall w\,(z = w \to \lvert a_{n} - c \rvert < z \to \lvert a_{n} - c \rvert < w)$}
\UIC{$\forall z\,\forall w\,(z = w \to \lvert a_{n} - c \rvert < z \to \lvert a_{n} - c \rvert < w)$}
\AXC{$\forall z\,\forall w\,(z = w \to \lvert a_{n} - c \rvert < z \to \lvert a_{n} - c \rvert < w) \to \varphi_{122}$\lbl{A2}}
\BIC{$\varphi_{122}$}
\AXC{$\varphi_{122} \to \varphi_{123}$\lbl{A2}}
\BIC{$\varphi_{123}$}
\end{ptree}
\noindent$\Pi_{2}$:
\begin{ptree}
\AXC{$z = w \to \lvert z - c \rvert < \varepsilon \cdot 2^{-1} \to \lvert w - c \rvert < \varepsilon \cdot 2^{-1}$\lbl{A6}}
\UIC{$\varphi_{46}$}
\UIC{$\forall z\,(\varphi_{46})$}
\AXC{$\forall z\,(\varphi_{46}) \to \varphi_{100}$\lbl{A2}}
\BIC{$\varphi_{100}$}
\AXC{$\varphi_{100} \to b_{N_{1} \cup N_{2}} = a_{s_{N_{1} \cup N_{2}}} \to \varphi_{96}$\lbl{A2}}
\BIC{$b_{N_{1} \cup N_{2}} = a_{s_{N_{1} \cup N_{2}}} \to \varphi_{96}$}
\end{ptree}
\noindent$\Pi_{3}$:
\begin{ptree}
\AXC{$\forall m\,\forall n\,(n \leq m \cup n)$\lbl{[C8]}}
\AXC{$\forall m\,\forall n\,(n \leq m \cup n) \to \forall n\,(n \leq N_{1} \cup n)$\lbl{A2}}
\BIC{$\forall n\,(n \leq N_{1} \cup n)$}
\AXC{$\forall n\,(n \leq N_{1} \cup n) \to N_{2} \leq N_{1} \cup N_{2}$\lbl{A2}}
\BIC{$N_{2} \leq N_{1} \cup N_{2}$}
\end{ptree}
\noindent$\Pi_{4}$:
\begin{ptree}
\AXC{$\forall m\,\forall n\,(m \in \mathbb{N} \to n \in \mathbb{N} \to m \cup n \in \mathbb{N})$\lbl{[N3]}}
\AXC{$\varphi_{87}$\lbl{A2}}
\BIC{$\forall n\,(N_{1} \in \mathbb{N} \to n \in \mathbb{N} \to N_{1} \cup n \in \mathbb{N})$}
\AXC{$\varphi_{81}$\lbl{A2}}
\BIC{$N_{1} \in \mathbb{N} \to N_{2} \in \mathbb{N} \to N_{1} \cup N_{2} \in \mathbb{N}$}
\AXC{$(N_{1} \in \mathbb{N} \to N_{2} \in \mathbb{N} \to N_{1} \cup N_{2} \in \mathbb{N}) \to \varphi_{269} \to \varphi_{268} \to \varphi_{288}$\lbl{A1}}
\BIC{$\varphi_{269} \to \varphi_{268} \to \varphi_{288}$}
\end{ptree}
\noindent$\Pi_{5}$:
\begin{ptree}
\AXC{$\varphi_{310}$\lbl{A1}}
\AXC{$\varphi_{310} \to \varphi_{268}$\lbl{A1}}
\BIC{$\varphi_{268}$}
\AXC{$\varphi_{280}$\lbl{A1}}
\AXC{$\varphi_{280} \to \varphi_{269}$\lbl{A1}}
\BIC{$\varphi_{269}$}
\AXC{$\stackrel{\Pi_{4}}{\vdots}$}\noLine
\UIC{$\varphi_{269} \to \varphi_{268} \to \varphi_{288}$}
\BIC{$\varphi_{268} \to \varphi_{288}$}
\BIC{$\varphi_{288}$}
\end{ptree}
\noindent$\Pi_{6}$:
\begin{ptree}
\AXC{$\forall n \in \mathbb{N}\,(N_{2} \leq n \to \lvert b_{n} - c \rvert < \varepsilon \cdot 2^{-1}) \to \varphi_{65}$\lbl{A2}}
\AXC{$\varphi_{310}$\lbl{A1}}
\AXC{$\varphi_{310} \to \varphi_{313}$\lbl{A1}}
\BIC{$\varphi_{313}$}
\AXC{$\varphi_{362}$\lbl{A1}}
\BIC{$(\forall n \in \mathbb{N}\,(N_{2} \leq n \to \lvert b_{n} - c \rvert < \varepsilon \cdot 2^{-1}) \to \varphi_{65}) \to \varphi_{324}$}
\BIC{$\varphi_{324}$}
\end{ptree}
\noindent$\Pi_{7}$:
\begin{ptree}
\AXC{$\stackrel{\Pi_{3}}{\vdots}$}\noLine
\UIC{$N_{2} \leq N_{1} \cup N_{2}$}
\AXC{$\stackrel{\Pi_{5}}{\vdots}$}\noLine
\UIC{$\varphi_{288}$}
\AXC{$\stackrel{\Pi_{6}}{\vdots}$}\noLine
\UIC{$\varphi_{324}$}
\AXC{$\varphi_{324} \to \varphi_{288} \to N_{2} \leq N_{1} \cup N_{2} \to \varphi_{308}$\lbl{A1}}
\BIC{$\varphi_{288} \to N_{2} \leq N_{1} \cup N_{2} \to \varphi_{308}$}
\BIC{$N_{2} \leq N_{1} \cup N_{2} \to \varphi_{308}$}
\BIC{$\varphi_{308}$}
\end{ptree}
\noindent$\Pi_{8}$:
\begin{ptree}
\AXC{$\stackrel{\Pi_{5}}{\vdots}$}\noLine
\UIC{$\varphi_{288}$}
\AXC{$\varphi_{106}$\lbl{A2}}
\AXC{$\varphi_{173}$\lbl{A1}}
\AXC{$\varphi_{173} \to \varphi_{304}$\lbl{A1}}
\BIC{$\varphi_{304}$}
\AXC{$\varphi_{304} \to \varphi_{106} \to \varphi_{317}$\lbl{A1}}
\BIC{$\varphi_{106} \to \varphi_{317}$}
\BIC{$\varphi_{317}$}
\AXC{$\varphi_{317} \to \varphi_{288} \to \varphi_{306}$\lbl{A1}}
\BIC{$\varphi_{288} \to \varphi_{306}$}
\BIC{$\varphi_{306}$}
\end{ptree}
\noindent$\Pi_{9}$:
\begin{ptree}
\AXC{$\stackrel{\Pi_{2}}{\vdots}$}\noLine
\UIC{$b_{N_{1} \cup N_{2}} = a_{s_{N_{1} \cup N_{2}}} \to \varphi_{96}$}
\AXC{$\stackrel{\Pi_{7}}{\vdots}$}\noLine
\UIC{$\varphi_{308}$}
\AXC{$\stackrel{\Pi_{8}}{\vdots}$}\noLine
\UIC{$\varphi_{306}$}
\AXC{$\varphi_{306} \to \varphi_{308} \to (b_{N_{1} \cup N_{2}} = a_{s_{N_{1} \cup N_{2}}} \to \varphi_{96}) \to \varphi_{311}$\lbl{A1}}
\BIC{$\varphi_{308} \to (b_{N_{1} \cup N_{2}} = a_{s_{N_{1} \cup N_{2}}} \to \varphi_{96}) \to \varphi_{311}$}
\BIC{$(b_{N_{1} \cup N_{2}} = a_{s_{N_{1} \cup N_{2}}} \to \varphi_{96}) \to \varphi_{311}$}
\BIC{$\varphi_{311}$}
\end{ptree}
\noindent$\Pi_{10}$:
\begin{ptree}
\AXC{$\forall s\,(\mathrm{SubIdx}(s) \to \forall k \in \mathbb{N}\,(k \leq s_{k}))$\lbl{[C17]}}
\AXC{$\varphi_{64}$\lbl{A2}}
\BIC{$\mathrm{SubIdx}(s) \to \forall k \in \mathbb{N}\,(k \leq s_{k})$}
\AXC{$(\mathrm{SubIdx}(s) \to \forall k \in \mathbb{N}\,(k \leq s_{k})) \to \varphi_{113} \to \varphi_{135}$\lbl{A1}}
\BIC{$\varphi_{113} \to \varphi_{135}$}
\end{ptree}
\noindent$\Pi_{11}$:
\begin{ptree}
\AXC{$\stackrel{\Pi_{5}}{\vdots}$}\noLine
\UIC{$\varphi_{288}$}
\AXC{$\varphi_{62}$\lbl{A2}}
\AXC{$\varphi_{173}$\lbl{A1}}
\AXC{$\varphi_{173} \to \varphi_{113}$\lbl{A1}}
\BIC{$\varphi_{113}$}
\AXC{$\stackrel{\Pi_{10}}{\vdots}$}\noLine
\UIC{$\varphi_{113} \to \varphi_{135}$}
\BIC{$\varphi_{135}$}
\AXC{$\varphi_{135} \to \varphi_{62} \to \varphi_{170}$\lbl{A1}}
\BIC{$\varphi_{62} \to \varphi_{170}$}
\BIC{$\varphi_{170}$}
\AXC{$\varphi_{170} \to \varphi_{288} \to \varphi_{295}$\lbl{A1}}
\BIC{$\varphi_{288} \to \varphi_{295}$}
\BIC{$\varphi_{295}$}
\end{ptree}
\noindent$\Pi_{12}$:
\begin{ptree}
\AXC{$\forall m\,\forall n\,(m \leq m \cup n)$\lbl{[C7]}}
\AXC{$\forall m\,\forall n\,(m \leq m \cup n) \to \forall n\,(N_{1} \leq N_{1} \cup n)$\lbl{A2}}
\BIC{$\forall n\,(N_{1} \leq N_{1} \cup n)$}
\AXC{$\forall n\,(N_{1} \leq N_{1} \cup n) \to N_{1} \leq N_{1} \cup N_{2}$\lbl{A2}}
\BIC{$N_{1} \leq N_{1} \cup N_{2}$}
\end{ptree}
\noindent$\Pi_{13}$:
\begin{ptree}
\AXC{$\forall x\,\forall y\,\forall z\,(x \leq y \to y \leq z \to x \leq z)$\lbl{[O5]}}
\AXC{$\varphi_{22}$\lbl{A2}}
\BIC{$\forall y\,\forall z\,(N_{1} \leq y \to y \leq z \to N_{1} \leq z)$}
\AXC{$\varphi_{41}$\lbl{A2}}
\BIC{$\forall z\,(N_{1} \leq N_{1} \cup N_{2} \to N_{1} \cup N_{2} \leq z \to N_{1} \leq z)$}
\AXC{$\forall z\,(N_{1} \leq N_{1} \cup N_{2} \to N_{1} \cup N_{2} \leq z \to N_{1} \leq z) \to \varphi_{44}$\lbl{A2}}
\BIC{$\varphi_{44}$}
\end{ptree}
\noindent$\Pi_{14}$:
\begin{ptree}
\AXC{$\forall s\,(\varphi_{30})$\lbl{DefSubIdx}}
\AXC{$\forall s\,(\varphi_{30}) \to \varphi_{30}$\lbl{A2}}
\BIC{$\varphi_{30}$}
\AXC{$\varphi_{173}$\lbl{A1}}
\AXC{$\varphi_{173} \to \varphi_{113}$\lbl{A1}}
\BIC{$\varphi_{113}$}
\AXC{$\varphi_{113} \to \varphi_{30} \to \varphi_{111}$\lbl{A1}}
\BIC{$\varphi_{30} \to \varphi_{111}$}
\BIC{$\varphi_{111}$}
\end{ptree}
\noindent$\Pi_{15}$:
\begin{ptree}
\AXC{$\stackrel{\Pi_{5}}{\vdots}$}\noLine
\UIC{$\varphi_{288}$}
\AXC{$\stackrel{\Pi_{14}}{\vdots}$}\noLine
\UIC{$\varphi_{111}$}
\AXC{$\forall s\,\forall k\,(\mathrm{NSeq}(s) \to k \in \mathbb{N} \to s_{k} \in \mathbb{N})$\lbl{[S13]}}
\AXC{$\varphi_{77}$\lbl{A2}}
\BIC{$\forall k\,(\mathrm{NSeq}(s) \to k \in \mathbb{N} \to s_{k} \in \mathbb{N})$}
\AXC{$\varphi_{101}$\lbl{A2}}
\BIC{$\mathrm{NSeq}(s) \to N_{1} \cup N_{2} \in \mathbb{N} \to s_{N_{1} \cup N_{2}} \in \mathbb{N}$}
\AXC{$\varphi_{361}$\lbl{A1}}
\BIC{$\varphi_{111} \to \varphi_{288} \to \varphi_{292}$}
\BIC{$\varphi_{288} \to \varphi_{292}$}
\BIC{$\varphi_{292}$}
\end{ptree}
\noindent$\Pi_{16}$:
\begin{ptree}
\AXC{$\varphi_{75} \to \varphi_{83}$\lbl{A2}}
\AXC{$\varphi_{280}$\lbl{A1}}
\AXC{$\varphi_{280} \to \varphi_{318}$\lbl{A1}}
\BIC{$\varphi_{318}$}
\AXC{$\varphi_{318} \to (\varphi_{75} \to \varphi_{83}) \to \varphi_{319}$\lbl{A1}}
\BIC{$(\varphi_{75} \to \varphi_{83}) \to \varphi_{319}$}
\BIC{$\varphi_{319}$}
\AXC{$\varphi_{319} \to (\varphi_{83} \to v_{753} \in \mathbb{N} \to v_{754} \in \mathbb{N} \to \varphi_{48}) \to \varphi_{322}$\lbl{A1}}
\BIC{$(\varphi_{83} \to v_{753} \in \mathbb{N} \to v_{754} \in \mathbb{N} \to \varphi_{48}) \to \varphi_{322}$}
\end{ptree}
\noindent$\Pi_{17}$:
\begin{ptree}
\AXC{$\varphi_{83} \to v_{753} \in \mathbb{N} \to v_{754} \in \mathbb{N} \to \varphi_{48}$\lbl{A2}}
\AXC{$\stackrel{\Pi_{16}}{\vdots}$}\noLine
\UIC{$(\varphi_{83} \to v_{753} \in \mathbb{N} \to v_{754} \in \mathbb{N} \to \varphi_{48}) \to \varphi_{322}$}
\BIC{$\varphi_{322}$}
\AXC{$\varphi_{322} \to \varphi_{259} \to v_{753} \in \mathbb{N} \to v_{754} \in \mathbb{N} \to \varphi_{48}$\lbl{A1}}
\BIC{$\varphi_{259} \to v_{753} \in \mathbb{N} \to v_{754} \in \mathbb{N} \to \varphi_{48}$}
\UIC{$\forall v_{754}\,(\varphi_{259} \to v_{753} \in \mathbb{N} \to v_{754} \in \mathbb{N} \to \varphi_{48})$}
\end{ptree}
\noindent$\Pi_{18}$:
\begin{ptree}
\AXC{$\varphi_{259} \to \forall v_{754}\,(\varphi_{259})$\lbl{A4}}
\AXC{$\varphi_{359}$\lbl{A3}}
\AXC{$\stackrel{\Pi_{17}}{\vdots}$}\noLine
\UIC{$\forall v_{754}\,(\varphi_{259} \to v_{753} \in \mathbb{N} \to v_{754} \in \mathbb{N} \to \varphi_{48})$}
\AXC{$\varphi_{371}$\lbl{A1}}
\BIC{$\varphi_{359} \to (\varphi_{259} \to \forall v_{754}\,(\varphi_{259})) \to \varphi_{327}$}
\BIC{$(\varphi_{259} \to \forall v_{754}\,(\varphi_{259})) \to \varphi_{327}$}
\BIC{$\varphi_{327}$}
\end{ptree}
\noindent$\Pi_{19}$:
\begin{ptree}
\AXC{$\varphi_{259} \to \forall v_{753}\,(\varphi_{259})$\lbl{A4}}
\AXC{$\varphi_{360}$\lbl{A3}}
\AXC{$\stackrel{\Pi_{18}}{\vdots}$}\noLine
\UIC{$\varphi_{327}$}
\AXC{$\varphi_{327} \to \varphi_{259} \to \forall v_{754}\,(v_{753} \in \mathbb{N} \to v_{754} \in \mathbb{N} \to \varphi_{48})$\lbl{A1}}
\BIC{$\varphi_{259} \to \forall v_{754}\,(v_{753} \in \mathbb{N} \to v_{754} \in \mathbb{N} \to \varphi_{48})$}
\UIC{$\forall v_{753}\,(\varphi_{259} \to \forall v_{754}\,(v_{753} \in \mathbb{N} \to v_{754} \in \mathbb{N} \to \varphi_{48}))$}
\AXC{$\varphi_{373}$\lbl{A1}}
\BIC{$\varphi_{360} \to (\varphi_{259} \to \forall v_{753}\,(\varphi_{259})) \to \varphi_{329}$}
\BIC{$(\varphi_{259} \to \forall v_{753}\,(\varphi_{259})) \to \varphi_{329}$}
\BIC{$\varphi_{329}$}
\end{ptree}
\noindent$\Pi_{20}$:
\begin{ptree}
\AXC{$\forall v_{753}\,\forall v_{754}\,(v_{753} \in \mathbb{N} \to v_{754} \in \mathbb{N} \to \varphi_{48}) \to \varphi_{89}$\lbl{A2}}
\AXC{$\stackrel{\Pi_{19}}{\vdots}$}\noLine
\UIC{$\varphi_{329}$}
\AXC{$\varphi_{329} \to \varphi_{349}$\lbl{A1}}
\BIC{$\varphi_{349}$}
\BIC{$\varphi_{321}$}
\end{ptree}
\noindent$\Pi_{21}$:
\begin{ptree}
\AXC{$\varphi_{89} \to n \in \mathbb{N} \to s_{N_{1} \cup N_{2}} \in \mathbb{N} \to N_{1} \leq n \to \varphi_{66}$\lbl{A2}}
\AXC{$\stackrel{\Pi_{20}}{\vdots}$}\noLine
\UIC{$\varphi_{321}$}
\AXC{$\varphi_{321} \to (\varphi_{89} \to n \in \mathbb{N} \to s_{N_{1} \cup N_{2}} \in \mathbb{N} \to N_{1} \leq n \to \varphi_{66}) \to \varphi_{334}$\lbl{A1}}
\BIC{$(\varphi_{89} \to n \in \mathbb{N} \to s_{N_{1} \cup N_{2}} \in \mathbb{N} \to N_{1} \leq n \to \varphi_{66}) \to \varphi_{334}$}
\BIC{$\varphi_{334}$}
\end{ptree}
\noindent$\Pi_{22}$:
\begin{ptree}
\AXC{$\varphi_{276}$\lbl{A1}}
\AXC{$\stackrel{\Pi_{15}}{\vdots}$}\noLine
\UIC{$\varphi_{292}$}
\AXC{$\varphi_{275}$\lbl{A1}}
\AXC{$\stackrel{\Pi_{21}}{\vdots}$}\noLine
\UIC{$\varphi_{334}$}
\AXC{$\varphi_{334} \to \varphi_{275} \to \varphi_{292} \to \varphi_{276} \to \varphi_{289} \to \varphi_{314}$\lbl{A1}}
\BIC{$\varphi_{275} \to \varphi_{292} \to \varphi_{276} \to \varphi_{289} \to \varphi_{314}$}
\BIC{$\varphi_{292} \to \varphi_{276} \to \varphi_{289} \to \varphi_{314}$}
\BIC{$\varphi_{276} \to \varphi_{289} \to \varphi_{314}$}
\BIC{$\varphi_{289} \to \varphi_{314}$}
\end{ptree}
\noindent$\Pi_{23}$:
\begin{ptree}
\AXC{$\stackrel{\Pi_{11}}{\vdots}$}\noLine
\UIC{$\varphi_{295}$}
\AXC{$\stackrel{\Pi_{12}}{\vdots}$}\noLine
\UIC{$N_{1} \leq N_{1} \cup N_{2}$}
\AXC{$\stackrel{\Pi_{13}}{\vdots}$}\noLine
\UIC{$\varphi_{44}$}
\AXC{$\varphi_{44} \to N_{1} \leq N_{1} \cup N_{2} \to \varphi_{295} \to \varphi_{289}$\lbl{A1}}
\BIC{$N_{1} \leq N_{1} \cup N_{2} \to \varphi_{295} \to \varphi_{289}$}
\BIC{$\varphi_{295} \to \varphi_{289}$}
\BIC{$\varphi_{289}$}
\AXC{$\stackrel{\Pi_{22}}{\vdots}$}\noLine
\UIC{$\varphi_{289} \to \varphi_{314}$}
\BIC{$\varphi_{314}$}
\end{ptree}
\noindent$\Pi_{24}$:
\begin{ptree}
\AXC{$2^{-1} \in \mathbb{R}$\lbl{[F31]}}
\AXC{$\varphi_{137}$\lbl{A1}}
\AXC{$\forall x\,\forall y\,(x \in \mathbb{R} \to y \in \mathbb{R} \to x \cdot y \in \mathbb{R})$\lbl{(R1)}}
\AXC{$\varphi_{94}$\lbl{A2}}
\BIC{$\forall y\,(\varepsilon \in \mathbb{R} \to y \in \mathbb{R} \to \varepsilon \cdot y \in \mathbb{R})$}
\AXC{$\varphi_{108}$\lbl{A2}}
\BIC{$\varepsilon \in \mathbb{R} \to 2^{-1} \in \mathbb{R} \to \varepsilon \cdot 2^{-1} \in \mathbb{R}$}
\AXC{$\varphi_{302}$\lbl{A1}}
\BIC{$\varphi_{137} \to 2^{-1} \in \mathbb{R} \to \varphi_{165}$}
\BIC{$2^{-1} \in \mathbb{R} \to \varphi_{165}$}
\BIC{$\varphi_{165}$}
\end{ptree}
\noindent$\Pi_{25}$:
\begin{ptree}
\AXC{$\forall a\,\forall \alpha\,(\mathrm{Conv}(a,\alpha) \to \alpha \in \mathbb{R})$\lbl{[C2]}}
\AXC{$\varphi_{34}$\lbl{A2}}
\BIC{$\forall \alpha\,(\mathrm{Conv}(b,\alpha) \to \alpha \in \mathbb{R})$}
\AXC{$\varphi_{26}$\lbl{A2}}
\BIC{$\mathrm{Conv}(b,c) \to c \in \mathbb{R}$}
\AXC{$(\mathrm{Conv}(b,c) \to c \in \mathbb{R}) \to \varphi_{192}$\lbl{A1}}
\BIC{$\varphi_{192}$}
\end{ptree}
\noindent$\Pi_{26}$:
\begin{ptree}
\AXC{$\mathrm{Cauchy}(a) \to \varphi_{15}$\lbl{A1}}
\AXC{$\stackrel{\Pi_{25}}{\vdots}$}\noLine
\UIC{$\varphi_{192}$}
\BIC{$\mathrm{Cauchy}(a) \to \mathrm{SubSeq}(b,a) \land \mathrm{Cgt}(b) \to \varphi_{61}$}
\end{ptree}
\noindent$\Pi_{27}$:
\begin{ptree}
\AXC{$\forall a\,(\varphi_{156})$\lbl{DefCauchy}}
\AXC{$\forall a\,(\varphi_{156}) \to \varphi_{156}$\lbl{A2}}
\BIC{$\varphi_{156}$}
\AXC{$\mathrm{Cauchy}(a) \to \mathrm{Cauchy}(a)$\lbl{A1}}
\AXC{$\varphi_{216}$\lbl{A1}}
\BIC{$\varphi_{156} \to \mathrm{Cauchy}(a) \to \mathrm{Seq}(a)$}
\BIC{$\mathrm{Cauchy}(a) \to \mathrm{Seq}(a)$}
\end{ptree}
\noindent$\Pi_{28}$:
\begin{ptree}
\AXC{$\stackrel{\Pi_{27}}{\vdots}$}\noLine
\UIC{$\mathrm{Cauchy}(a) \to \mathrm{Seq}(a)$}
\AXC{$\forall a\,\forall n\,(\mathrm{Seq}(a) \to n \in \mathbb{N} \to a_{n} \in \mathbb{R})$\lbl{[S1]}}
\AXC{$\varphi_{78}$\lbl{A2}}
\BIC{$\forall n\,(\mathrm{Seq}(a) \to n \in \mathbb{N} \to a_{n} \in \mathbb{R})$}
\AXC{$\varphi_{120}$\lbl{A2}}
\BIC{$\mathrm{Seq}(a) \to s_{N_{1} \cup N_{2}} \in \mathbb{N} \to a_{s_{N_{1} \cup N_{2}}} \in \mathbb{R}$}
\AXC{$\varphi_{358}$\lbl{A1}}
\BIC{$(\mathrm{Cauchy}(a) \to \mathrm{Seq}(a)) \to \varphi_{292} \to \varphi_{299}$}
\BIC{$\varphi_{292} \to \varphi_{299}$}
\end{ptree}
\noindent$\Pi_{29}$:
\begin{ptree}
\AXC{$\stackrel{\Pi_{27}}{\vdots}$}\noLine
\UIC{$\mathrm{Cauchy}(a) \to \mathrm{Seq}(a)$}
\AXC{$\forall a\,\forall n\,(\mathrm{Seq}(a) \to n \in \mathbb{N} \to a_{n} \in \mathbb{R})$\lbl{[S1]}}
\AXC{$\varphi_{78}$\lbl{A2}}
\BIC{$\forall n\,(\mathrm{Seq}(a) \to n \in \mathbb{N} \to a_{n} \in \mathbb{R})$}
\AXC{$\varphi_{68}$\lbl{A2}}
\BIC{$\mathrm{Seq}(a) \to n \in \mathbb{N} \to a_{n} \in \mathbb{R}$}
\AXC{$\varphi_{355}$\lbl{A1}}
\BIC{$(\mathrm{Cauchy}(a) \to \mathrm{Seq}(a)) \to \varphi_{275} \to \varphi_{287}$}
\BIC{$\varphi_{275} \to \varphi_{287}$}
\end{ptree}
\noindent$\Pi_{30}$:
\begin{ptree}
\AXC{$\varphi_{131}$\lbl{[C27]}}
\AXC{$\varphi_{131} \to \varphi_{157}$\lbl{A2}}
\BIC{$\varphi_{157}$}
\AXC{$\varphi_{157} \to \varphi_{223}$\lbl{A2}}
\BIC{$\varphi_{223}$}
\AXC{$\varphi_{223} \to \varphi_{222}$\lbl{A2}}
\BIC{$\varphi_{222}$}
\AXC{$\varphi_{222} \to a_{n} \in \mathbb{R} \to a_{s_{N_{1} \cup N_{2}}} \in \mathbb{R} \to c \in \mathbb{R} \to \varphi_{218}$\lbl{A2}}
\BIC{$a_{n} \in \mathbb{R} \to a_{s_{N_{1} \cup N_{2}}} \in \mathbb{R} \to c \in \mathbb{R} \to \varphi_{218}$}
\end{ptree}
\noindent$\Pi_{31}$:
\begin{ptree}
\AXC{$\stackrel{\Pi_{15}}{\vdots}$}\noLine
\UIC{$\varphi_{292}$}
\AXC{$\stackrel{\Pi_{28}}{\vdots}$}\noLine
\UIC{$\varphi_{292} \to \varphi_{299}$}
\BIC{$\varphi_{299}$}
\AXC{$\varphi_{275}$\lbl{A1}}
\AXC{$\stackrel{\Pi_{29}}{\vdots}$}\noLine
\UIC{$\varphi_{275} \to \varphi_{287}$}
\BIC{$\varphi_{287}$}
\AXC{$\stackrel{\Pi_{30}}{\vdots}$}\noLine
\UIC{$a_{n} \in \mathbb{R} \to a_{s_{N_{1} \cup N_{2}}} \in \mathbb{R} \to c \in \mathbb{R} \to \varphi_{218}$}
\AXC{$\varphi_{374}$\lbl{A1}}
\BIC{$\varphi_{287} \to \varphi_{299} \to \varphi_{367}$}
\BIC{$\varphi_{299} \to \varphi_{367}$}
\BIC{$\varphi_{367}$}
\end{ptree}
\noindent$\Pi_{32}$:
\begin{ptree}
\AXC{$\stackrel{\Pi_{9}}{\vdots}$}\noLine
\UIC{$\varphi_{311}$}
\AXC{$\stackrel{\Pi_{23}}{\vdots}$}\noLine
\UIC{$\varphi_{314}$}
\AXC{$\stackrel{\Pi_{24}}{\vdots}$}\noLine
\UIC{$\varphi_{165}$}
\AXC{$\stackrel{\Pi_{26}}{\vdots}$}\noLine
\UIC{$\mathrm{Cauchy}(a) \to \mathrm{SubSeq}(b,a) \land \mathrm{Cgt}(b) \to \varphi_{61}$}
\AXC{$\stackrel{\Pi_{31}}{\vdots}$}\noLine
\UIC{$\varphi_{367}$}
\BIC{$\varphi_{165} \to \varphi_{314} \to \varphi_{311} \to \varphi_{312}$}
\BIC{$\varphi_{314} \to \varphi_{311} \to \varphi_{312}$}
\BIC{$\varphi_{311} \to \varphi_{312}$}
\BIC{$\varphi_{312}$}
\end{ptree}
\noindent$\Pi_{33}$:
\begin{ptree}
\AXC{$\varphi_{283}$\lbl{A1}}
\AXC{$\forall e \in \mathbb{R}\,(e \cdot 2^{-1} + e \cdot 2^{-1} = e)$\lbl{[F36]}}
\AXC{$\varphi_{90}$\lbl{A2}}
\BIC{$\varepsilon \in \mathbb{R} \to \varepsilon \cdot 2^{-1} + \varepsilon \cdot 2^{-1} = \varepsilon$}
\AXC{$(\varepsilon \in \mathbb{R} \to \varepsilon \cdot 2^{-1} + \varepsilon \cdot 2^{-1} = \varepsilon) \to \varphi_{283} \to \varphi_{300}$\lbl{A1}}
\BIC{$\varphi_{283} \to \varphi_{300}$}
\BIC{$\varphi_{300}$}
\end{ptree}
\noindent$\Pi_{34}$:
\begin{ptree}
\AXC{$\stackrel{\Pi_{1}}{\vdots}$}\noLine
\UIC{$\varphi_{123}$}
\AXC{$\stackrel{\Pi_{32}}{\vdots}$}\noLine
\UIC{$\varphi_{312}$}
\AXC{$\stackrel{\Pi_{33}}{\vdots}$}\noLine
\UIC{$\varphi_{300}$}
\AXC{$\varphi_{300} \to \varphi_{312} \to \varphi_{123} \to \varphi_{293}$\lbl{A1}}
\BIC{$\varphi_{312} \to \varphi_{123} \to \varphi_{293}$}
\BIC{$\varphi_{123} \to \varphi_{293}$}
\BIC{$\varphi_{293}$}
\AXC{$\varphi_{293} \to \varphi_{259} \to n \in \mathbb{N} \to N_{1} \leq n \to \lvert a_{n} - c \rvert < \varepsilon$\lbl{A1}}
\BIC{$\varphi_{259} \to n \in \mathbb{N} \to N_{1} \leq n \to \lvert a_{n} - c \rvert < \varepsilon$}
\UIC{$\forall n\,(\varphi_{259} \to n \in \mathbb{N} \to N_{1} \leq n \to \lvert a_{n} - c \rvert < \varepsilon)$}
\AXC{$\varphi_{370}$\lbl{A1}}
\BIC{$\varphi_{352} \to (\varphi_{259} \to \forall n\,(\varphi_{259})) \to \varphi_{296}$}
\end{ptree}
\noindent$\Pi_{35}$:
\begin{ptree}
\AXC{$\varphi_{259} \to \forall n\,(\varphi_{259})$\lbl{A4}}
\AXC{$\varphi_{352}$\lbl{A3}}
\AXC{$\stackrel{\Pi_{34}}{\vdots}$}\noLine
\UIC{$\varphi_{352} \to (\varphi_{259} \to \forall n\,(\varphi_{259})) \to \varphi_{296}$}
\BIC{$(\varphi_{259} \to \forall n\,(\varphi_{259})) \to \varphi_{296}$}
\BIC{$\varphi_{296}$}
\AXC{$\varphi_{280}$\lbl{A1}}
\AXC{$\varphi_{280} \to \varphi_{269}$\lbl{A1}}
\BIC{$\varphi_{269}$}
\AXC{$\varphi_{269} \to \varphi_{296} \to \varphi_{305}$\lbl{A1}}
\BIC{$\varphi_{296} \to \varphi_{305}$}
\BIC{$\varphi_{305}$}
\end{ptree}
\noindent$\Pi_{36}$:
\begin{ptree}
\AXC{$\varphi_{137}$\lbl{A1}}
\AXC{$\forall e \in \mathbb{R}\,(0 < e \to 0 < e \cdot 2^{-1})$\lbl{[F35]}}
\AXC{$\varphi_{79}$\lbl{A2}}
\BIC{$\varepsilon \in \mathbb{R} \to 0 < \varepsilon \to 0 < \varepsilon \cdot 2^{-1}$}
\AXC{$(\varepsilon \in \mathbb{R} \to 0 < \varepsilon \to 0 < \varepsilon \cdot 2^{-1}) \to \varphi_{137} \to \varphi_{149} \to \varphi_{162}$\lbl{A1}}
\BIC{$\varphi_{137} \to \varphi_{149} \to \varphi_{162}$}
\BIC{$\varphi_{149} \to \varphi_{162}$}
\end{ptree}
\noindent$\Pi_{37}$:
\begin{ptree}
\AXC{$\varphi_{149}$\lbl{A1}}
\AXC{$\stackrel{\Pi_{36}}{\vdots}$}\noLine
\UIC{$\varphi_{149} \to \varphi_{162}$}
\BIC{$\varphi_{162}$}
\end{ptree}
\noindent$\Pi_{38}$:
\begin{ptree}
\AXC{$\varphi_{119}$\lbl{[C3]}}
\AXC{$\varphi_{119} \to \varphi_{114}$\lbl{A2}}
\BIC{$\varphi_{114}$}
\AXC{$\varphi_{114} \to \varphi_{104}$\lbl{A2}}
\BIC{$\varphi_{104}$}
\AXC{$\varphi_{104} \to \mathrm{Conv}(b,c) \to \varepsilon \cdot 2^{-1} \in \mathbb{R} \to \varphi_{84}$\lbl{A2}}
\BIC{$\mathrm{Conv}(b,c) \to \varepsilon \cdot 2^{-1} \in \mathbb{R} \to \varphi_{84}$}
\AXC{$\varphi_{357}$\lbl{A1}}
\BIC{$(\mathrm{Cauchy}(a) \to \varphi_{15}) \to \varphi_{165} \to \varphi_{162} \to \varphi_{262}$}
\end{ptree}
\noindent$\Pi_{39}$:
\begin{ptree}
\AXC{$\stackrel{\Pi_{37}}{\vdots}$}\noLine
\UIC{$\varphi_{162}$}
\AXC{$\stackrel{\Pi_{24}}{\vdots}$}\noLine
\UIC{$\varphi_{165}$}
\AXC{$\mathrm{Cauchy}(a) \to \varphi_{15}$\lbl{A1}}
\AXC{$\stackrel{\Pi_{38}}{\vdots}$}\noLine
\UIC{$(\mathrm{Cauchy}(a) \to \varphi_{15}) \to \varphi_{165} \to \varphi_{162} \to \varphi_{262}$}
\BIC{$\varphi_{165} \to \varphi_{162} \to \varphi_{262}$}
\BIC{$\varphi_{162} \to \varphi_{262}$}
\BIC{$\varphi_{262}$}
\AXC{$\varphi_{262} \to \forall N_{2}\,(\varphi_{255} \to \varphi_{55}) \to \varphi_{372}$\lbl{A1}}
\BIC{$\forall N_{2}\,(\varphi_{255} \to \varphi_{55}) \to \varphi_{372}$}
\end{ptree}
\noindent$\Pi_{40}$:
\begin{ptree}
\AXC{$\varphi_{134}$\lbl{A2}}
\AXC{$\stackrel{\Pi_{35}}{\vdots}$}\noLine
\UIC{$\varphi_{305}$}
\AXC{$\varphi_{305} \to \varphi_{134} \to \varphi_{307}$\lbl{A1}}
\BIC{$\varphi_{134} \to \varphi_{307}$}
\BIC{$\varphi_{307}$}
\AXC{$\varphi_{307} \to \varphi_{255} \to \varphi_{55}$\lbl{A1}}
\BIC{$\varphi_{255} \to \varphi_{55}$}
\UIC{$\forall N_{2}\,(\varphi_{255} \to \varphi_{55})$}
\AXC{$\stackrel{\Pi_{39}}{\vdots}$}\noLine
\UIC{$\forall N_{2}\,(\varphi_{255} \to \varphi_{55}) \to \varphi_{372}$}
\BIC{$\varphi_{372}$}
\end{ptree}
\noindent$\Pi_{41}$:
\begin{ptree}
\AXC{$\forall N_{2}\,(\varphi_{55}) \to \varphi_{51}$\lbl{A2}}
\UIC{$\forall N\,(\forall N_{2}\,(\varphi_{55}) \to \varphi_{51})$}
\AXC{$\varphi_{255} \to \forall N_{2}\,(\varphi_{255})$\lbl{A4}}
\AXC{$\forall N_{2}\,(\varphi_{255} \to \varphi_{55}) \to \forall N_{2}\,(\varphi_{255}) \to \forall N_{2}\,(\varphi_{55})$\lbl{A3}}
\AXC{$\stackrel{\Pi_{40}}{\vdots}$}\noLine
\UIC{$\varphi_{372}$}
\BIC{$(\varphi_{255} \to \forall N_{2}\,(\varphi_{255})) \to \forall N\,(\forall N_{2}\,(\varphi_{55}) \to \varphi_{51}) \to \varphi_{356}$}
\BIC{$\forall N\,(\forall N_{2}\,(\varphi_{55}) \to \varphi_{51}) \to \varphi_{356}$}
\BIC{$\varphi_{356}$}
\end{ptree}
\noindent$\Pi_{42}$:
\begin{ptree}
\AXC{$\forall N_{2}\,(\varphi_{55}) \to \forall N\,\forall N_{2}\,(\varphi_{55})$\lbl{A4}}
\AXC{$\forall N\,(\forall N_{2}\,(\varphi_{55}) \to \varphi_{51}) \to \forall N\,\forall N_{2}\,(\varphi_{55}) \to \forall N\,(\varphi_{51})$\lbl{A3}}
\AXC{$\stackrel{\Pi_{41}}{\vdots}$}\noLine
\UIC{$\varphi_{356}$}
\BIC{$(\forall N_{2}\,(\varphi_{55}) \to \forall N\,\forall N_{2}\,(\varphi_{55})) \to \varphi_{254}$}
\BIC{$\varphi_{254}$}
\end{ptree}
\noindent$\Pi_{43}$:
\begin{ptree}
\AXC{$\forall a\,(\varphi_{143})$\lbl{[C26]}}
\AXC{$\forall a\,(\varphi_{143}) \to \varphi_{143}$\lbl{A2}}
\BIC{$\varphi_{143}$}
\AXC{$\varphi_{143} \to \mathrm{Cauchy}(a) \to \varepsilon \cdot 2^{-1} \in \mathbb{R} \to 0 < \varepsilon \cdot 2^{-1} \to \varphi_{99}$\lbl{A2}}
\BIC{$\mathrm{Cauchy}(a) \to \varepsilon \cdot 2^{-1} \in \mathbb{R} \to 0 < \varepsilon \cdot 2^{-1} \to \varphi_{99}$}
\AXC{$\varphi_{353}$\lbl{A1}}
\BIC{$(\mathrm{Cauchy}(a) \to \mathrm{Cauchy}(a)) \to \varphi_{165} \to \varphi_{162} \to \varphi_{231}$}
\end{ptree}
\noindent$\Pi_{44}$:
\begin{ptree}
\AXC{$\stackrel{\Pi_{37}}{\vdots}$}\noLine
\UIC{$\varphi_{162}$}
\AXC{$\stackrel{\Pi_{24}}{\vdots}$}\noLine
\UIC{$\varphi_{165}$}
\AXC{$\mathrm{Cauchy}(a) \to \mathrm{Cauchy}(a)$\lbl{A1}}
\AXC{$\stackrel{\Pi_{43}}{\vdots}$}\noLine
\UIC{$(\mathrm{Cauchy}(a) \to \mathrm{Cauchy}(a)) \to \varphi_{165} \to \varphi_{162} \to \varphi_{231}$}
\BIC{$\varphi_{165} \to \varphi_{162} \to \varphi_{231}$}
\BIC{$\varphi_{162} \to \varphi_{231}$}
\BIC{$\varphi_{231}$}
\end{ptree}
\noindent$\Pi_{45}$:
\begin{ptree}
\AXC{$\varphi_{348}$\lbl{A3}}
\AXC{$\stackrel{\Pi_{42}}{\vdots}$}\noLine
\UIC{$\varphi_{254}$}
\AXC{$\varphi_{254} \to \varphi_{217} \to \neg (N_{1} \in \mathbb{N} \land \varphi_{75})$\lbl{A1}}
\BIC{$\varphi_{217} \to \neg (N_{1} \in \mathbb{N} \land \varphi_{75})$}
\UIC{$\forall N_{1}\,(\varphi_{217} \to \neg (N_{1} \in \mathbb{N} \land \varphi_{75}))$}
\AXC{$\stackrel{\Pi_{44}}{\vdots}$}\noLine
\UIC{$\varphi_{231}$}
\AXC{$\varphi_{231} \to \forall N_{1}\,(\varphi_{217} \to \neg (N_{1} \in \mathbb{N} \land \varphi_{75})) \to \varphi_{348} \to \varphi_{368}$\lbl{A1}}
\BIC{$\forall N_{1}\,(\varphi_{217} \to \neg (N_{1} \in \mathbb{N} \land \varphi_{75})) \to \varphi_{348} \to \varphi_{368}$}
\BIC{$\varphi_{348} \to \varphi_{368}$}
\BIC{$\varphi_{368}$}
\end{ptree}
\noindent$\Pi_{46}$:
\begin{ptree}
\AXC{$\varphi_{225}$\lbl{A4}}
\AXC{$\varphi_{331}$\lbl{A3}}
\AXC{$\forall N_{1}\,(\neg (N_{1} \in \mathbb{N} \land \varphi_{75})) \to \varphi_{93}$\lbl{A2}}
\UIC{$\forall N\,(\forall N_{1}\,(\neg (N_{1} \in \mathbb{N} \land \varphi_{75})) \to \varphi_{93})$}
\AXC{$\varphi_{217} \to \forall N_{1}\,(\varphi_{217})$\lbl{A4}}
\AXC{$\stackrel{\Pi_{45}}{\vdots}$}\noLine
\UIC{$\varphi_{368}$}
\BIC{$\forall N\,(\forall N_{1}\,(\neg (N_{1} \in \mathbb{N} \land \varphi_{75})) \to \varphi_{93}) \to \varphi_{331} \to \varphi_{326}$}
\BIC{$\varphi_{331} \to \varphi_{326}$}
\BIC{$\varphi_{326}$}
\BIC{$\mathrm{Cauchy}(a) \to \mathrm{SubSeq}(b,a) \land \mathrm{Cgt}(b) \to \varphi_{176}$}
\AXC{$\varphi_{315}$\lbl{A1}}
\BIC{$\varphi_{95} \to \varphi_{76}$}
\UIC{$\forall \varepsilon\,(\varphi_{95} \to \varphi_{76})$}
\end{ptree}
\noindent$\Pi_{47}$:
\begin{ptree}
\AXC{$\forall \varepsilon\,(\varphi_{95} \to \varphi_{76}) \to \forall \varepsilon\,(\varphi_{95}) \to \varphi_{82}$\lbl{A3}}
\AXC{$\stackrel{\Pi_{46}}{\vdots}$}\noLine
\UIC{$\forall \varepsilon\,(\varphi_{95} \to \varphi_{76})$}
\AXC{$\forall \varepsilon\,(\varphi_{95} \to \varphi_{76}) \to (\forall \varepsilon\,(\varphi_{95} \to \varphi_{76}) \to \forall \varepsilon\,(\varphi_{95}) \to \varphi_{82}) \to \varphi_{323}$\lbl{A1}}
\BIC{$(\forall \varepsilon\,(\varphi_{95} \to \varphi_{76}) \to \forall \varepsilon\,(\varphi_{95}) \to \varphi_{82}) \to \varphi_{323}$}
\BIC{$\varphi_{323}$}
\end{ptree}
\noindent$\Pi_{48}$:
\begin{ptree}
\AXC{$\forall a\,(\varphi_{144})$\lbl{DefConv}}
\AXC{$\forall a\,(\varphi_{144}) \to \varphi_{144}$\lbl{A2}}
\BIC{$\varphi_{144}$}
\AXC{$\varphi_{144} \to (\mathrm{Conv}(a,c) \leftrightarrow \mathrm{Seq}(a) \land c \in \mathbb{R} \land \varphi_{82})$\lbl{A2}}
\BIC{$\mathrm{Conv}(a,c) \leftrightarrow \mathrm{Seq}(a) \land c \in \mathbb{R} \land \varphi_{82}$}
\AXC{$\varphi_{345}$\lbl{A1}}
\BIC{$(\mathrm{Cauchy}(a) \to \mathrm{Seq}(a)) \to \varphi_{332}$}
\end{ptree}
\noindent$\Pi_{49}$:
\begin{ptree}
\AXC{$\stackrel{\Pi_{27}}{\vdots}$}\noLine
\UIC{$\mathrm{Cauchy}(a) \to \mathrm{Seq}(a)$}
\AXC{$\stackrel{\Pi_{48}}{\vdots}$}\noLine
\UIC{$(\mathrm{Cauchy}(a) \to \mathrm{Seq}(a)) \to \varphi_{332}$}
\BIC{$\varphi_{332}$}
\end{ptree}
\noindent$\Pi_{50}$:
\begin{ptree}
\AXC{$\varphi_{95} \to \forall \varepsilon\,(\varphi_{95})$\lbl{A4}}
\AXC{$\stackrel{\Pi_{47}}{\vdots}$}\noLine
\UIC{$\varphi_{323}$}
\BIC{$\mathrm{Cauchy}(a) \to \mathrm{SubSeq}(b,a) \land \mathrm{Cgt}(b) \to \varphi_{180}$}
\AXC{$\stackrel{\Pi_{26}}{\vdots}$}\noLine
\UIC{$\mathrm{Cauchy}(a) \to \mathrm{SubSeq}(b,a) \land \mathrm{Cgt}(b) \to \varphi_{61}$}
\AXC{$\stackrel{\Pi_{49}}{\vdots}$}\noLine
\UIC{$\varphi_{332}$}
\BIC{$\varphi_{265}$}
\BIC{$\varphi_{115}$}
\end{ptree}
\noindent$\Pi_{51}$:
\begin{ptree}
\AXC{$\forall a\,(\mathrm{Cgt}(a) \leftrightarrow \exists \alpha\,(\mathrm{Conv}(a,\alpha)))$\lbl{DefCgt}}
\AXC{$\varphi_{23}$\lbl{A2}}
\BIC{$\mathrm{Cgt}(a) \leftrightarrow \exists \alpha\,(\mathrm{Conv}(a,\alpha))$}
\AXC{$(\mathrm{Cgt}(a) \leftrightarrow \exists \alpha\,(\mathrm{Conv}(a,\alpha))) \to \varphi_{121} \to \varphi_{110}$\lbl{A1}}
\BIC{$\varphi_{121} \to \varphi_{110}$}
\end{ptree}
\noindent$\Pi_{52}$:
\begin{ptree}
\AXC{$\forall \alpha\,(\neg \mathrm{Conv}(a,\alpha)) \to \neg \mathrm{Conv}(a,c)$\lbl{A2}}
\AXC{$\stackrel{\Pi_{50}}{\vdots}$}\noLine
\UIC{$\varphi_{115}$}
\AXC{$\varphi_{247}$\lbl{A1}}
\BIC{$(\forall \alpha\,(\neg \mathrm{Conv}(a,\alpha)) \to \neg \mathrm{Conv}(a,c)) \to \varphi_{121}$}
\BIC{$\varphi_{121}$}
\AXC{$\stackrel{\Pi_{51}}{\vdots}$}\noLine
\UIC{$\varphi_{121} \to \varphi_{110}$}
\BIC{$\varphi_{110}$}
\end{ptree}
\noindent$\Pi_{53}$:
\begin{ptree}
\AXC{$\forall b\,\forall a\,(\mathrm{SubSeq}(b,a) \leftrightarrow \varphi_{40})$\lbl{DefSubSeq}}
\AXC{$\varphi_{188}$\lbl{A2}}
\BIC{$\forall a\,(\mathrm{SubSeq}(b,a) \leftrightarrow \varphi_{40})$}
\AXC{$\varphi_{183}$\lbl{A2}}
\BIC{$\mathrm{SubSeq}(b,a) \leftrightarrow \varphi_{40}$}
\end{ptree}
\noindent$\Pi_{54}$:
\begin{ptree}
\AXC{$\stackrel{\Pi_{53}}{\vdots}$}\noLine
\UIC{$\mathrm{SubSeq}(b,a) \leftrightarrow \varphi_{40}$}
\AXC{$\mathrm{Cauchy}(a) \to \varphi_{13}$\lbl{A1}}
\AXC{$(\mathrm{Cauchy}(a) \to \varphi_{13}) \to \varphi_{205}$\lbl{A1}}
\BIC{$\varphi_{205}$}
\BIC{$\mathrm{Cauchy}(a) \to \mathrm{SubSeq}(b,a) \land \mathrm{Cgt}(b) \to \varphi_{40}$}
\AXC{$\varphi_{220}$\lbl{A1}}
\BIC{$\mathrm{Cauchy}(a) \to \mathrm{SubSeq}(b,a) \land \mathrm{Cgt}(b) \to \varphi_{43}$}
\AXC{$\varphi_{351}$\lbl{A1}}
\BIC{$\varphi_{117} \to \varphi_{325}$}
\end{ptree}
\noindent$\Pi_{55}$:
\begin{ptree}
\AXC{$\varphi_{27} \to \forall s\,(\varphi_{27})$\lbl{A4}}
\AXC{$\varphi_{117} \to \forall s\,(\varphi_{27}) \to \varphi_{29}$\lbl{A3}}
\AXC{$\stackrel{\Pi_{52}}{\vdots}$}\noLine
\UIC{$\varphi_{110}$}
\AXC{$\varphi_{228}$\lbl{A1}}
\BIC{$\varphi_{27} \to \neg (\mathrm{SubIdx}(s) \land \forall k \in \mathbb{N}\,(b_{k} = a_{s_{k}}))$}
\UIC{$\varphi_{117}$}
\AXC{$\stackrel{\Pi_{54}}{\vdots}$}\noLine
\UIC{$\varphi_{117} \to \varphi_{325}$}
\BIC{$\varphi_{325}$}
\BIC{$(\varphi_{27} \to \forall s\,(\varphi_{27})) \to \mathrm{Cauchy}(a) \to \varphi_{9}$}
\BIC{$\mathrm{Cauchy}(a) \to \varphi_{9}$}
\end{ptree}
\noindent$\Pi_{56}$:
\begin{ptree}
\AXC{$\stackrel{\Pi_{55}}{\vdots}$}\noLine
\UIC{$\mathrm{Cauchy}(a) \to \varphi_{9}$}
\AXC{$(\mathrm{Cauchy}(a) \to \varphi_{9}) \to \varphi_{11} \to \neg \mathrm{Conv}(b,c)$\lbl{A1}}
\BIC{$\varphi_{11} \to \neg \mathrm{Conv}(b,c)$}
\UIC{$\forall c\,(\varphi_{11} \to \neg \mathrm{Conv}(b,c))$}
\end{ptree}
\noindent$\Pi_{57}$:
\begin{ptree}
\AXC{$\forall a\,(\mathrm{Cgt}(a) \leftrightarrow \exists \alpha\,(\mathrm{Conv}(a,\alpha)))$\lbl{DefCgt}}
\AXC{$\varphi_{24}$\lbl{A2}}
\BIC{$\mathrm{Cgt}(b) \leftrightarrow \exists \alpha\,(\mathrm{Conv}(b,\alpha))$}
\AXC{$\mathrm{Cauchy}(a) \to \varphi_{13}$\lbl{A1}}
\AXC{$(\mathrm{Cauchy}(a) \to \varphi_{13}) \to \varphi_{57}$\lbl{A1}}
\BIC{$\varphi_{57}$}
\BIC{$\mathrm{Cauchy}(a) \to \varphi_{6}$}
\end{ptree}
\noindent$\Pi_{58}$:
\begin{ptree}
\AXC{$\varphi_{11} \to \forall c\,(\varphi_{11})$\lbl{A4}}
\AXC{$\varphi_{130}$\lbl{A3}}
\AXC{$\stackrel{\Pi_{56}}{\vdots}$}\noLine
\UIC{$\forall c\,(\varphi_{11} \to \neg \mathrm{Conv}(b,c))$}
\AXC{$\stackrel{\Pi_{57}}{\vdots}$}\noLine
\UIC{$\mathrm{Cauchy}(a) \to \varphi_{6}$}
\AXC{$\varphi_{343}$\lbl{A1}}
\BIC{$\forall c\,(\varphi_{11} \to \neg \mathrm{Conv}(b,c)) \to \varphi_{130} \to (\varphi_{11} \to \forall c\,(\varphi_{11})) \to \varphi_{209}$}
\BIC{$\varphi_{130} \to (\varphi_{11} \to \forall c\,(\varphi_{11})) \to \varphi_{209}$}
\BIC{$(\varphi_{11} \to \forall c\,(\varphi_{11})) \to \varphi_{209}$}
\BIC{$\varphi_{209}$}
\end{ptree}
\noindent$\Pi_{59}$:
\begin{ptree}
\AXC{$\forall a\,(\varphi_{2})$\lbl{定理6.8}}
\AXC{$\forall a\,(\varphi_{2}) \to \varphi_{2}$\lbl{A2}}
\BIC{$\varphi_{2}$}
\AXC{$\varphi_{2} \to (\mathrm{Cauchy}(a) \to \mathrm{SeqBdd}(a)) \to \varphi_{3}$\lbl{A1}}
\BIC{$(\mathrm{Cauchy}(a) \to \mathrm{SeqBdd}(a)) \to \varphi_{3}$}
\end{ptree}
\noindent$\Pi_{60}$:
\begin{ptree}
\AXC{$\forall a\,(\mathrm{Cauchy}(a) \to \mathrm{SeqBdd}(a))$\lbl{補題6.10}}
\AXC{$\varphi_{19}$\lbl{A2}}
\BIC{$\mathrm{Cauchy}(a) \to \mathrm{SeqBdd}(a)$}
\AXC{$\stackrel{\Pi_{59}}{\vdots}$}\noLine
\UIC{$(\mathrm{Cauchy}(a) \to \mathrm{SeqBdd}(a)) \to \varphi_{3}$}
\BIC{$\varphi_{3}$}
\AXC{$\varphi_{3} \to \varphi_{226}$\lbl{A1}}
\BIC{$\varphi_{226}$}
\end{ptree}
\noindent$\Pi_{61}$:
\begin{ptree}
\AXC{$\varphi_{1}$\lbl{A4}}
\AXC{$\varphi_{50}$\lbl{A3}}
\AXC{$\forall c\,(\neg \mathrm{Conv}(b,c)) \to \neg \mathrm{Conv}(b,\alpha)$\lbl{A2}}
\UIC{$\forall \alpha\,(\forall c\,(\neg \mathrm{Conv}(b,c)) \to \neg \mathrm{Conv}(b,\alpha))$}
\AXC{$\stackrel{\Pi_{58}}{\vdots}$}\noLine
\UIC{$\varphi_{209}$}
\BIC{$\varphi_{50} \to \varphi_{1} \to \varphi_{8}$}
\BIC{$\varphi_{1} \to \varphi_{8}$}
\BIC{$\varphi_{8}$}
\AXC{$\varphi_{8} \to \varphi_{12}$\lbl{A1}}
\BIC{$\varphi_{12}$}
\UIC{$\forall b\,(\varphi_{12})$}
\AXC{$\stackrel{\Pi_{60}}{\vdots}$}\noLine
\UIC{$\varphi_{226}$}
\BIC{$(\forall b\,(\varphi_{12}) \to \varphi_{21}) \to \varphi_{14} \to \mathrm{Cauchy}(a) \to \mathrm{Cgt}(a)$}
\end{ptree}
\noindent$\Pi_{62}$:
\begin{ptree}
\AXC{$z = w \to \lvert a_{m} - a_{n} \rvert < z \to \lvert a_{m} - a_{n} \rvert < w$\lbl{A6}}
\UIC{$\forall w\,(z = w \to \lvert a_{m} - a_{n} \rvert < z \to \lvert a_{m} - a_{n} \rvert < w)$}
\UIC{$\forall z\,\forall w\,(z = w \to \lvert a_{m} - a_{n} \rvert < z \to \lvert a_{m} - a_{n} \rvert < w)$}
\AXC{$\varphi_{224}$\lbl{A2}}
\BIC{$\varphi_{138}$}
\AXC{$\varphi_{138} \to \varphi_{139}$\lbl{A2}}
\BIC{$\varphi_{139}$}
\end{ptree}
\noindent$\Pi_{63}$:
\begin{ptree}
\AXC{$z = w \to z < \varepsilon \cdot 2^{-1} \to w < \varepsilon \cdot 2^{-1}$\lbl{A6}}
\UIC{$\forall w\,(z = w \to z < \varepsilon \cdot 2^{-1} \to w < \varepsilon \cdot 2^{-1})$}
\UIC{$\forall z\,\forall w\,(z = w \to z < \varepsilon \cdot 2^{-1} \to w < \varepsilon \cdot 2^{-1})$}
\AXC{$\forall z\,\forall w\,(z = w \to z < \varepsilon \cdot 2^{-1} \to w < \varepsilon \cdot 2^{-1}) \to \varphi_{73}$\lbl{A2}}
\BIC{$\varphi_{73}$}
\AXC{$\varphi_{73} \to \varphi_{124}$\lbl{A2}}
\BIC{$\varphi_{124}$}
\end{ptree}
\noindent$\Pi_{64}$:
\begin{ptree}
\AXC{$\varphi_{92} \to \varphi_{85}$\lbl{A2}}
\AXC{$\mathrm{Cgt}(a) \to \varphi_{246}$\lbl{A1}}
\AXC{$(\mathrm{Cgt}(a) \to \varphi_{246}) \to \varphi_{245}$\lbl{A1}}
\BIC{$\varphi_{245}$}
\AXC{$\varphi_{245} \to (\varphi_{92} \to \varphi_{85}) \to \varphi_{242}$\lbl{A1}}
\BIC{$(\varphi_{92} \to \varphi_{85}) \to \varphi_{242}$}
\BIC{$\varphi_{242}$}
\end{ptree}
\noindent$\Pi_{65}$:
\begin{ptree}
\AXC{$\mathrm{Cgt}(a) \to \mathrm{Conv}(a,\alpha) \to \varepsilon \in \mathbb{R} \to \varphi_{184}$\lbl{A1}}
\AXC{$\mathrm{Cgt}(a) \to \varphi_{211}$\lbl{A1}}
\AXC{$\stackrel{\Pi_{64}}{\vdots}$}\noLine
\UIC{$\varphi_{242}$}
\AXC{$\varphi_{242} \to (\mathrm{Cgt}(a) \to \varphi_{211}) \to \varphi_{335}$\lbl{A1}}
\BIC{$(\mathrm{Cgt}(a) \to \varphi_{211}) \to \varphi_{335}$}
\BIC{$\varphi_{335}$}
\BIC{$\mathrm{Cgt}(a) \to \mathrm{Conv}(a,\alpha) \to \varepsilon \in \mathbb{R} \to 0 < \varepsilon \to \varphi_{212}$}
\end{ptree}
\noindent$\Pi_{66}$:
\begin{ptree}
\AXC{$\forall a\,\forall \alpha\,(\mathrm{Conv}(a,\alpha) \to \alpha \in \mathbb{R})$\lbl{[C2]}}
\AXC{$\varphi_{35}$\lbl{A2}}
\BIC{$\forall \alpha\,(\mathrm{Conv}(a,\alpha) \to \alpha \in \mathbb{R})$}
\AXC{$\varphi_{28}$\lbl{A2}}
\BIC{$\mathrm{Conv}(a,\alpha) \to \alpha \in \mathbb{R}$}
\AXC{$(\mathrm{Conv}(a,\alpha) \to \alpha \in \mathbb{R}) \to \varphi_{47}$\lbl{A1}}
\BIC{$\varphi_{47}$}
\end{ptree}
\noindent$\Pi_{67}$:
\begin{ptree}
\AXC{$\mathrm{Cgt}(a) \to \mathrm{Conv}(a,\alpha) \to \mathrm{Conv}(a,\alpha)$\lbl{A1}}
\AXC{$\stackrel{\Pi_{66}}{\vdots}$}\noLine
\UIC{$\varphi_{47}$}
\BIC{$\mathrm{Cgt}(a) \to \mathrm{Conv}(a,\alpha) \to \alpha \in \mathbb{R}$}
\end{ptree}
\noindent$\Pi_{68}$:
\begin{ptree}
\AXC{$\forall a\,\forall \alpha\,(\mathrm{Conv}(a,\alpha) \to \mathrm{Seq}(a))$\lbl{[C1]}}
\AXC{$\varphi_{25}$\lbl{A2}}
\BIC{$\forall \alpha\,(\mathrm{Conv}(a,\alpha) \to \mathrm{Seq}(a))$}
\AXC{$\varphi_{17}$\lbl{A2}}
\BIC{$\mathrm{Conv}(a,\alpha) \to \mathrm{Seq}(a)$}
\AXC{$(\mathrm{Conv}(a,\alpha) \to \mathrm{Seq}(a)) \to \varphi_{37}$\lbl{A1}}
\BIC{$\varphi_{37}$}
\end{ptree}
\noindent$\Pi_{69}$:
\begin{ptree}
\AXC{$\mathrm{Cgt}(a) \to \mathrm{Conv}(a,\alpha) \to \mathrm{Conv}(a,\alpha)$\lbl{A1}}
\AXC{$\stackrel{\Pi_{68}}{\vdots}$}\noLine
\UIC{$\varphi_{37}$}
\BIC{$\mathrm{Cgt}(a) \to \mathrm{Conv}(a,\alpha) \to \mathrm{Seq}(a)$}
\end{ptree}
\noindent$\Pi_{70}$:
\begin{ptree}
\AXC{$\forall a\,\forall n\,(\mathrm{Seq}(a) \to n \in \mathbb{N} \to a_{n} \in \mathbb{R})$\lbl{[S1]}}
\AXC{$\varphi_{78}$\lbl{A2}}
\BIC{$\forall n\,(\mathrm{Seq}(a) \to n \in \mathbb{N} \to a_{n} \in \mathbb{R})$}
\AXC{$\varphi_{68}$\lbl{A2}}
\BIC{$\mathrm{Seq}(a) \to n \in \mathbb{N} \to a_{n} \in \mathbb{R}$}
\AXC{$\varphi_{338}$\lbl{A1}}
\BIC{$(\mathrm{Cgt}(a) \to \mathrm{Conv}(a,\alpha) \to \mathrm{Seq}(a)) \to \varphi_{328}$}
\end{ptree}
\noindent$\Pi_{71}$:
\begin{ptree}
\AXC{$\mathrm{Cgt}(a) \to \varphi_{211}$\lbl{A1}}
\AXC{$\stackrel{\Pi_{69}}{\vdots}$}\noLine
\UIC{$\mathrm{Cgt}(a) \to \mathrm{Conv}(a,\alpha) \to \mathrm{Seq}(a)$}
\AXC{$\stackrel{\Pi_{70}}{\vdots}$}\noLine
\UIC{$(\mathrm{Cgt}(a) \to \mathrm{Conv}(a,\alpha) \to \mathrm{Seq}(a)) \to \varphi_{328}$}
\BIC{$\varphi_{328}$}
\BIC{$\mathrm{Cgt}(a) \to \mathrm{Conv}(a,\alpha) \to \varepsilon \in \mathbb{R} \to \varphi_{200}$}
\end{ptree}
\noindent$\Pi_{72}$:
\begin{ptree}
\AXC{$\forall x\,\forall y\,(x \in \mathbb{R} \to y \in \mathbb{R} \to \lvert x - y \rvert = \lvert y - x \rvert)$\lbl{[A13]}}
\AXC{$\varphi_{140}$\lbl{A2}}
\BIC{$\forall y\,(a_{n} \in \mathbb{R} \to y \in \mathbb{R} \to \lvert a_{n} - y \rvert = \lvert y - a_{n} \rvert)$}
\AXC{$\varphi_{166}$\lbl{A2}}
\BIC{$a_{n} \in \mathbb{R} \to \alpha \in \mathbb{R} \to \lvert a_{n} - \alpha \rvert = \lvert \alpha - a_{n} \rvert$}
\end{ptree}
\noindent$\Pi_{73}$:
\begin{ptree}
\AXC{$\stackrel{\Pi_{72}}{\vdots}$}\noLine
\UIC{$a_{n} \in \mathbb{R} \to \alpha \in \mathbb{R} \to \lvert a_{n} - \alpha \rvert = \lvert \alpha - a_{n} \rvert$}
\AXC{$(a_{n} \in \mathbb{R} \to \alpha \in \mathbb{R} \to \lvert a_{n} - \alpha \rvert = \lvert \alpha - a_{n} \rvert) \to \varphi_{339}$\lbl{A1}}
\BIC{$\varphi_{339}$}
\end{ptree}
\noindent$\Pi_{74}$:
\begin{ptree}
\AXC{$\stackrel{\Pi_{67}}{\vdots}$}\noLine
\UIC{$\mathrm{Cgt}(a) \to \mathrm{Conv}(a,\alpha) \to \alpha \in \mathbb{R}$}
\AXC{$\stackrel{\Pi_{71}}{\vdots}$}\noLine
\UIC{$\mathrm{Cgt}(a) \to \mathrm{Conv}(a,\alpha) \to \varepsilon \in \mathbb{R} \to \varphi_{200}$}
\AXC{$\stackrel{\Pi_{73}}{\vdots}$}\noLine
\UIC{$\varphi_{339}$}
\BIC{$(\mathrm{Cgt}(a) \to \mathrm{Conv}(a,\alpha) \to \alpha \in \mathbb{R}) \to \varphi_{235}$}
\BIC{$\varphi_{235}$}
\AXC{$\varphi_{235} \to \varphi_{346}$\lbl{A1}}
\BIC{$\varphi_{346}$}
\end{ptree}
\noindent$\Pi_{75}$:
\begin{ptree}
\AXC{$z = w \to z < \varepsilon \cdot 2^{-1} \to w < \varepsilon \cdot 2^{-1}$\lbl{A6}}
\UIC{$\forall w\,(z = w \to z < \varepsilon \cdot 2^{-1} \to w < \varepsilon \cdot 2^{-1})$}
\UIC{$\forall z\,\forall w\,(z = w \to z < \varepsilon \cdot 2^{-1} \to w < \varepsilon \cdot 2^{-1})$}
\AXC{$\forall z\,\forall w\,(z = w \to z < \varepsilon \cdot 2^{-1} \to w < \varepsilon \cdot 2^{-1}) \to \varphi_{74}$\lbl{A2}}
\BIC{$\varphi_{74}$}
\AXC{$\varphi_{74} \to \varphi_{125}$\lbl{A2}}
\BIC{$\varphi_{125}$}
\end{ptree}
\noindent$\Pi_{76}$:
\begin{ptree}
\AXC{$\varphi_{92} \to \varphi_{86}$\lbl{A2}}
\AXC{$\mathrm{Cgt}(a) \to \varphi_{246}$\lbl{A1}}
\AXC{$(\mathrm{Cgt}(a) \to \varphi_{246}) \to \varphi_{245}$\lbl{A1}}
\BIC{$\varphi_{245}$}
\AXC{$\varphi_{245} \to (\varphi_{92} \to \varphi_{86}) \to \varphi_{243}$\lbl{A1}}
\BIC{$(\varphi_{92} \to \varphi_{86}) \to \varphi_{243}$}
\BIC{$\varphi_{243}$}
\end{ptree}
\noindent$\Pi_{77}$:
\begin{ptree}
\AXC{$\forall a\,\forall n\,(\mathrm{Seq}(a) \to n \in \mathbb{N} \to a_{n} \in \mathbb{R})$\lbl{[S1]}}
\AXC{$\varphi_{78}$\lbl{A2}}
\BIC{$\forall n\,(\mathrm{Seq}(a) \to n \in \mathbb{N} \to a_{n} \in \mathbb{R})$}
\AXC{$\varphi_{69}$\lbl{A2}}
\BIC{$\mathrm{Seq}(a) \to m \in \mathbb{N} \to a_{m} \in \mathbb{R}$}
\AXC{$\varphi_{337}$\lbl{A1}}
\BIC{$(\mathrm{Cgt}(a) \to \mathrm{Conv}(a,\alpha) \to \mathrm{Seq}(a)) \to \varphi_{320}$}
\end{ptree}
\noindent$\Pi_{78}$:
\begin{ptree}
\AXC{$\mathrm{Cgt}(a) \to \varphi_{197}$\lbl{A1}}
\AXC{$\stackrel{\Pi_{69}}{\vdots}$}\noLine
\UIC{$\mathrm{Cgt}(a) \to \mathrm{Conv}(a,\alpha) \to \mathrm{Seq}(a)$}
\AXC{$\stackrel{\Pi_{77}}{\vdots}$}\noLine
\UIC{$(\mathrm{Cgt}(a) \to \mathrm{Conv}(a,\alpha) \to \mathrm{Seq}(a)) \to \varphi_{320}$}
\BIC{$\varphi_{320}$}
\BIC{$\mathrm{Cgt}(a) \to \mathrm{Conv}(a,\alpha) \to \varepsilon \in \mathbb{R} \to \varphi_{201}$}
\end{ptree}
\noindent$\Pi_{79}$:
\begin{ptree}
\AXC{$\forall x\,\forall y\,(x \in \mathbb{R} \to y \in \mathbb{R} \to \lvert x - y \rvert = \lvert y - x \rvert)$\lbl{[A13]}}
\AXC{$\varphi_{141}$\lbl{A2}}
\BIC{$\forall y\,(a_{m} \in \mathbb{R} \to y \in \mathbb{R} \to \lvert a_{m} - y \rvert = \lvert y - a_{m} \rvert)$}
\AXC{$\varphi_{167}$\lbl{A2}}
\BIC{$a_{m} \in \mathbb{R} \to \alpha \in \mathbb{R} \to \lvert a_{m} - \alpha \rvert = \lvert \alpha - a_{m} \rvert$}
\end{ptree}
\noindent$\Pi_{80}$:
\begin{ptree}
\AXC{$\stackrel{\Pi_{79}}{\vdots}$}\noLine
\UIC{$a_{m} \in \mathbb{R} \to \alpha \in \mathbb{R} \to \lvert a_{m} - \alpha \rvert = \lvert \alpha - a_{m} \rvert$}
\AXC{$(a_{m} \in \mathbb{R} \to \alpha \in \mathbb{R} \to \lvert a_{m} - \alpha \rvert = \lvert \alpha - a_{m} \rvert) \to \varphi_{340}$\lbl{A1}}
\BIC{$\varphi_{340}$}
\end{ptree}
\noindent$\Pi_{81}$:
\begin{ptree}
\AXC{$\stackrel{\Pi_{67}}{\vdots}$}\noLine
\UIC{$\mathrm{Cgt}(a) \to \mathrm{Conv}(a,\alpha) \to \alpha \in \mathbb{R}$}
\AXC{$\stackrel{\Pi_{78}}{\vdots}$}\noLine
\UIC{$\mathrm{Cgt}(a) \to \mathrm{Conv}(a,\alpha) \to \varepsilon \in \mathbb{R} \to \varphi_{201}$}
\AXC{$\stackrel{\Pi_{80}}{\vdots}$}\noLine
\UIC{$\varphi_{340}$}
\BIC{$(\mathrm{Cgt}(a) \to \mathrm{Conv}(a,\alpha) \to \alpha \in \mathbb{R}) \to \varphi_{236}$}
\BIC{$\varphi_{236}$}
\AXC{$\varphi_{236} \to \varphi_{347}$\lbl{A1}}
\BIC{$\varphi_{347}$}
\end{ptree}
\noindent$\Pi_{82}$:
\begin{ptree}
\AXC{$\mathrm{Cgt}(a) \to \mathrm{Conv}(a,\alpha) \to \varphi_{193}$\lbl{A1}}
\AXC{$\mathrm{Cgt}(a) \to \varphi_{197}$\lbl{A1}}
\AXC{$\stackrel{\Pi_{76}}{\vdots}$}\noLine
\UIC{$\varphi_{243}$}
\AXC{$\varphi_{243} \to (\mathrm{Cgt}(a) \to \varphi_{197}) \to \varphi_{333}$\lbl{A1}}
\BIC{$(\mathrm{Cgt}(a) \to \varphi_{197}) \to \varphi_{333}$}
\BIC{$\varphi_{333}$}
\BIC{$\mathrm{Cgt}(a) \to \mathrm{Conv}(a,\alpha) \to \varepsilon \in \mathbb{R} \to 0 < \varepsilon \to \varphi_{214}$}
\AXC{$\stackrel{\Pi_{81}}{\vdots}$}\noLine
\UIC{$\varphi_{347}$}
\BIC{$\varphi_{125} \to \mathrm{Cgt}(a) \to \mathrm{Conv}(a,\alpha) \to \varepsilon \in \mathbb{R} \to 0 < \varepsilon \to \varphi_{215}$}
\end{ptree}
\noindent$\Pi_{83}$:
\begin{ptree}
\AXC{$\forall x\,\forall y\,(x \in \mathbb{R} \to y \in \mathbb{R} \to x \cdot y \in \mathbb{R})$\lbl{(R1)}}
\AXC{$\varphi_{94}$\lbl{A2}}
\BIC{$\forall y\,(\varepsilon \in \mathbb{R} \to y \in \mathbb{R} \to \varepsilon \cdot y \in \mathbb{R})$}
\AXC{$\varphi_{108}$\lbl{A2}}
\BIC{$\varepsilon \in \mathbb{R} \to 2^{-1} \in \mathbb{R} \to \varepsilon \cdot 2^{-1} \in \mathbb{R}$}
\AXC{$\varphi_{291}$\lbl{A1}}
\BIC{$(\mathrm{Cgt}(a) \to \mathrm{Conv}(a,\alpha) \to \varepsilon \in \mathbb{R} \to \varepsilon \in \mathbb{R}) \to \varphi_{237}$}
\end{ptree}
\noindent$\Pi_{84}$:
\begin{ptree}
\AXC{$\mathrm{Cgt}(a) \to \mathrm{Conv}(a,\alpha) \to \varepsilon \in \mathbb{R} \to \varepsilon \in \mathbb{R}$\lbl{A1}}
\AXC{$\stackrel{\Pi_{83}}{\vdots}$}\noLine
\UIC{$(\mathrm{Cgt}(a) \to \mathrm{Conv}(a,\alpha) \to \varepsilon \in \mathbb{R} \to \varepsilon \in \mathbb{R}) \to \varphi_{237}$}
\BIC{$\varphi_{237}$}
\end{ptree}
\noindent$\Pi_{85}$:
\begin{ptree}
\AXC{$\varphi_{158}$\lbl{[C10]}}
\AXC{$\varphi_{158} \to \varphi_{168}$\lbl{A2}}
\BIC{$\varphi_{168}$}
\AXC{$\varphi_{168} \to \varphi_{172}$\lbl{A2}}
\BIC{$\varphi_{172}$}
\AXC{$\varphi_{172} \to \varphi_{175}$\lbl{A2}}
\BIC{$\varphi_{175}$}
\AXC{$\varphi_{175} \to v_{749} \in \mathbb{R} \to \varphi_{153}$\lbl{A2}}
\BIC{$v_{749} \in \mathbb{R} \to \varphi_{153}$}
\UIC{$\forall v_{752}\,(v_{749} \in \mathbb{R} \to \varphi_{153})$}
\UIC{$\forall v_{751}\,\forall v_{752}\,(v_{749} \in \mathbb{R} \to \varphi_{153})$}
\UIC{$\forall v_{750}\,\forall v_{751}\,\forall v_{752}\,(v_{749} \in \mathbb{R} \to \varphi_{153})$}
\UIC{$\forall v_{749}\,\forall v_{750}\,\forall v_{751}\,\forall v_{752}\,(v_{749} \in \mathbb{R} \to \varphi_{153})$}
\end{ptree}
\noindent$\Pi_{86}$:
\begin{ptree}
\AXC{$\stackrel{\Pi_{85}}{\vdots}$}\noLine
\UIC{$\forall v_{749}\,\forall v_{750}\,\forall v_{751}\,\forall v_{752}\,(v_{749} \in \mathbb{R} \to \varphi_{153})$}
\AXC{$\forall v_{749}\,\forall v_{750}\,\forall v_{751}\,\forall v_{752}\,(v_{749} \in \mathbb{R} \to \varphi_{153}) \to \varphi_{186}$\lbl{A2}}
\BIC{$\varphi_{186}$}
\AXC{$\varphi_{186} \to \varphi_{195}$\lbl{A2}}
\BIC{$\varphi_{195}$}
\end{ptree}
\noindent$\Pi_{87}$:
\begin{ptree}
\AXC{$\stackrel{\Pi_{67}}{\vdots}$}\noLine
\UIC{$\mathrm{Cgt}(a) \to \mathrm{Conv}(a,\alpha) \to \alpha \in \mathbb{R}$}
\AXC{$\stackrel{\Pi_{86}}{\vdots}$}\noLine
\UIC{$\varphi_{195}$}
\AXC{$\varphi_{195} \to \varphi_{203}$\lbl{A2}}
\BIC{$\varphi_{203}$}
\AXC{$\varphi_{203} \to \alpha \in \mathbb{R} \to \varphi_{221}$\lbl{A2}}
\BIC{$\alpha \in \mathbb{R} \to \varphi_{221}$}
\AXC{$\varphi_{369}$\lbl{A1}}
\BIC{$(\mathrm{Cgt}(a) \to \mathrm{Conv}(a,\alpha) \to \alpha \in \mathbb{R}) \to \varphi_{366}$}
\BIC{$\varphi_{366}$}
\end{ptree}
\noindent$\Pi_{88}$:
\begin{ptree}
\AXC{$\stackrel{\Pi_{71}}{\vdots}$}\noLine
\UIC{$\mathrm{Cgt}(a) \to \mathrm{Conv}(a,\alpha) \to \varepsilon \in \mathbb{R} \to \varphi_{200}$}
\AXC{$\stackrel{\Pi_{78}}{\vdots}$}\noLine
\UIC{$\mathrm{Cgt}(a) \to \mathrm{Conv}(a,\alpha) \to \varepsilon \in \mathbb{R} \to \varphi_{201}$}
\AXC{$\stackrel{\Pi_{87}}{\vdots}$}\noLine
\UIC{$\varphi_{366}$}
\BIC{$(\mathrm{Cgt}(a) \to \mathrm{Conv}(a,\alpha) \to \varepsilon \in \mathbb{R} \to \varphi_{200}) \to \varphi_{363}$}
\BIC{$\varphi_{363}$}
\end{ptree}
\noindent$\Pi_{89}$:
\begin{ptree}
\AXC{$\stackrel{\Pi_{75}}{\vdots}$}\noLine
\UIC{$\varphi_{125}$}
\AXC{$\stackrel{\Pi_{82}}{\vdots}$}\noLine
\UIC{$\varphi_{125} \to \mathrm{Cgt}(a) \to \mathrm{Conv}(a,\alpha) \to \varepsilon \in \mathbb{R} \to 0 < \varepsilon \to \varphi_{215}$}
\BIC{$\mathrm{Cgt}(a) \to \mathrm{Conv}(a,\alpha) \to \varepsilon \in \mathbb{R} \to 0 < \varepsilon \to \varphi_{215}$}
\AXC{$2^{-1} \in \mathbb{R}$\lbl{[F31]}}
\AXC{$\stackrel{\Pi_{84}}{\vdots}$}\noLine
\UIC{$\varphi_{237}$}
\BIC{$\mathrm{Cgt}(a) \to \mathrm{Conv}(a,\alpha) \to \varepsilon \in \mathbb{R} \to 0 < \varepsilon \to \varphi_{191}$}
\AXC{$\stackrel{\Pi_{88}}{\vdots}$}\noLine
\UIC{$\varphi_{363}$}
\BIC{$\varphi_{354}$}
\BIC{$\varphi_{341}$}
\end{ptree}
\noindent$\Pi_{90}$:
\begin{ptree}
\AXC{$\mathrm{Cgt}(a) \to \mathrm{Conv}(a,\alpha) \to \varepsilon \in \mathbb{R} \to \varepsilon \in \mathbb{R}$\lbl{A1}}
\AXC{$\forall e \in \mathbb{R}\,(e \cdot 2^{-1} + e \cdot 2^{-1} = e)$\lbl{[F36]}}
\AXC{$\varphi_{90}$\lbl{A2}}
\BIC{$\varepsilon \in \mathbb{R} \to \varepsilon \cdot 2^{-1} + \varepsilon \cdot 2^{-1} = \varepsilon$}
\AXC{$\varphi_{284}$\lbl{A1}}
\BIC{$(\mathrm{Cgt}(a) \to \mathrm{Conv}(a,\alpha) \to \varepsilon \in \mathbb{R} \to \varepsilon \in \mathbb{R}) \to \varphi_{232}$}
\BIC{$\varphi_{232}$}
\end{ptree}
\noindent$\Pi_{91}$:
\begin{ptree}
\AXC{$\stackrel{\Pi_{63}}{\vdots}$}\noLine
\UIC{$\varphi_{124}$}
\AXC{$\stackrel{\Pi_{65}}{\vdots}$}\noLine
\UIC{$\mathrm{Cgt}(a) \to \mathrm{Conv}(a,\alpha) \to \varepsilon \in \mathbb{R} \to 0 < \varepsilon \to \varphi_{212}$}
\AXC{$\stackrel{\Pi_{74}}{\vdots}$}\noLine
\UIC{$\varphi_{346}$}
\BIC{$\varphi_{124} \to \mathrm{Cgt}(a) \to \mathrm{Conv}(a,\alpha) \to \varepsilon \in \mathbb{R} \to 0 < \varepsilon \to \varphi_{213}$}
\BIC{$\mathrm{Cgt}(a) \to \mathrm{Conv}(a,\alpha) \to \varepsilon \in \mathbb{R} \to 0 < \varepsilon \to \varphi_{213}$}
\AXC{$\stackrel{\Pi_{89}}{\vdots}$}\noLine
\UIC{$\varphi_{341}$}
\BIC{$\mathrm{Cgt}(a) \to \varphi_{241}$}
\AXC{$\stackrel{\Pi_{90}}{\vdots}$}\noLine
\UIC{$\varphi_{232}$}
\AXC{$\varphi_{232} \to \varphi_{350}$\lbl{A1}}
\BIC{$\varphi_{350}$}
\BIC{$\varphi_{139} \to \mathrm{Cgt}(a) \to \mathrm{Conv}(a,\alpha) \to \varepsilon \in \mathbb{R} \to 0 < \varepsilon \to \varphi_{207}$}
\end{ptree}
\noindent$\Pi_{92}$:
\begin{ptree}
\AXC{$\stackrel{\Pi_{62}}{\vdots}$}\noLine
\UIC{$\varphi_{139}$}
\AXC{$\stackrel{\Pi_{91}}{\vdots}$}\noLine
\UIC{$\varphi_{139} \to \mathrm{Cgt}(a) \to \mathrm{Conv}(a,\alpha) \to \varepsilon \in \mathbb{R} \to 0 < \varepsilon \to \varphi_{207}$}
\BIC{$\mathrm{Cgt}(a) \to \mathrm{Conv}(a,\alpha) \to \varepsilon \in \mathbb{R} \to 0 < \varepsilon \to \varphi_{207}$}
\AXC{$\varphi_{336}$\lbl{A1}}
\BIC{$\varphi_{179} \to m \in \mathbb{N} \to n \in \mathbb{N} \to N \leq m \to N \leq n \to \lvert a_{m} - a_{n} \rvert < \varepsilon$}
\UIC{$\varphi_{233}$}
\end{ptree}
\noindent$\Pi_{93}$:
\begin{ptree}
\AXC{$\varphi_{179} \to \forall n\,(\varphi_{179})$\lbl{A4}}
\AXC{$\varphi_{233} \to \forall n\,(\varphi_{179}) \to \varphi_{56}$\lbl{A3}}
\AXC{$\stackrel{\Pi_{92}}{\vdots}$}\noLine
\UIC{$\varphi_{233}$}
\AXC{$\varphi_{233} \to (\varphi_{233} \to \forall n\,(\varphi_{179}) \to \varphi_{56}) \to (\varphi_{179} \to \forall n\,(\varphi_{179})) \to \varphi_{234}$\lbl{A1}}
\BIC{$(\varphi_{233} \to \forall n\,(\varphi_{179}) \to \varphi_{56}) \to (\varphi_{179} \to \forall n\,(\varphi_{179})) \to \varphi_{234}$}
\BIC{$(\varphi_{179} \to \forall n\,(\varphi_{179})) \to \varphi_{234}$}
\BIC{$\varphi_{234}$}
\end{ptree}
\noindent$\Pi_{94}$:
\begin{ptree}
\AXC{$\varphi_{179} \to \forall m\,(\varphi_{179})$\lbl{A4}}
\AXC{$\forall m\,(\varphi_{179} \to \varphi_{56}) \to \forall m\,(\varphi_{179}) \to \forall m\,(\varphi_{56})$\lbl{A3}}
\AXC{$\stackrel{\Pi_{93}}{\vdots}$}\noLine
\UIC{$\varphi_{234}$}
\AXC{$\varphi_{234} \to \varphi_{179} \to \varphi_{56}$\lbl{A1}}
\BIC{$\varphi_{179} \to \varphi_{56}$}
\UIC{$\forall m\,(\varphi_{179} \to \varphi_{56})$}
\AXC{$\forall m\,(\varphi_{179} \to \varphi_{56}) \to \varphi_{364}$\lbl{A1}}
\BIC{$\varphi_{364}$}
\BIC{$(\varphi_{179} \to \forall m\,(\varphi_{179})) \to \mathrm{Cgt}(a) \to \varphi_{230}$}
\BIC{$\mathrm{Cgt}(a) \to \varphi_{230}$}
\end{ptree}
\noindent$\Pi_{95}$:
\begin{ptree}
\AXC{$\stackrel{\Pi_{94}}{\vdots}$}\noLine
\UIC{$\mathrm{Cgt}(a) \to \varphi_{230}$}
\AXC{$\mathrm{Cgt}(a) \to \varphi_{246}$\lbl{A1}}
\AXC{$(\mathrm{Cgt}(a) \to \varphi_{246}) \to \varphi_{189}$\lbl{A1}}
\BIC{$\varphi_{189}$}
\AXC{$\varphi_{189} \to (\mathrm{Cgt}(a) \to \varphi_{230}) \to \mathrm{Cgt}(a) \to \varphi_{238}$\lbl{A1}}
\BIC{$(\mathrm{Cgt}(a) \to \varphi_{230}) \to \mathrm{Cgt}(a) \to \varphi_{238}$}
\BIC{$\mathrm{Cgt}(a) \to \varphi_{238}$}
\end{ptree}
\noindent$\Pi_{96}$:
\begin{ptree}
\AXC{$\varphi_{208}$\lbl{A2}}
\AXC{$\stackrel{\Pi_{95}}{\vdots}$}\noLine
\UIC{$\mathrm{Cgt}(a) \to \varphi_{238}$}
\AXC{$(\mathrm{Cgt}(a) \to \varphi_{238}) \to \varphi_{208} \to \mathrm{Cgt}(a) \to \varphi_{239}$\lbl{A1}}
\BIC{$\varphi_{208} \to \mathrm{Cgt}(a) \to \varphi_{239}$}
\BIC{$\mathrm{Cgt}(a) \to \varphi_{239}$}
\AXC{$(\mathrm{Cgt}(a) \to \varphi_{239}) \to \varphi_{154} \to \neg (N \in \mathbb{N} \land \varphi_{92})$\lbl{A1}}
\BIC{$\varphi_{154} \to \neg (N \in \mathbb{N} \land \varphi_{92})$}
\UIC{$\forall N\,(\varphi_{154} \to \neg (N \in \mathbb{N} \land \varphi_{92}))$}
\end{ptree}
\noindent$\Pi_{97}$:
\begin{ptree}
\AXC{$\mathrm{Cgt}(a) \to \mathrm{Conv}(a,\alpha) \to \varepsilon \in \mathbb{R} \to \varepsilon \in \mathbb{R}$\lbl{A1}}
\AXC{$\forall e \in \mathbb{R}\,(0 < e \to 0 < e \cdot 2^{-1})$\lbl{[F35]}}
\AXC{$\varphi_{79}$\lbl{A2}}
\BIC{$\varepsilon \in \mathbb{R} \to 0 < \varepsilon \to 0 < \varepsilon \cdot 2^{-1}$}
\AXC{$\varphi_{229}$\lbl{A1}}
\BIC{$(\mathrm{Cgt}(a) \to \mathrm{Conv}(a,\alpha) \to \varepsilon \in \mathbb{R} \to \varepsilon \in \mathbb{R}) \to \varphi_{132}$}
\BIC{$\varphi_{132}$}
\end{ptree}
\noindent$\Pi_{98}$:
\begin{ptree}
\AXC{$\forall x\,\forall y\,(x \in \mathbb{R} \to y \in \mathbb{R} \to x \cdot y \in \mathbb{R})$\lbl{(R1)}}
\AXC{$\varphi_{94}$\lbl{A2}}
\BIC{$\forall y\,(\varepsilon \in \mathbb{R} \to y \in \mathbb{R} \to \varepsilon \cdot y \in \mathbb{R})$}
\AXC{$\varphi_{108}$\lbl{A2}}
\BIC{$\varepsilon \in \mathbb{R} \to 2^{-1} \in \mathbb{R} \to \varepsilon \cdot 2^{-1} \in \mathbb{R}$}
\AXC{$(\varepsilon \in \mathbb{R} \to 2^{-1} \in \mathbb{R} \to \varepsilon \cdot 2^{-1} \in \mathbb{R}) \to \varphi_{155}$\lbl{A1}}
\BIC{$\varphi_{155}$}
\end{ptree}
\noindent$\Pi_{99}$:
\begin{ptree}
\AXC{$2^{-1} \in \mathbb{R}$\lbl{[F31]}}
\AXC{$\mathrm{Cgt}(a) \to \mathrm{Conv}(a,\alpha) \to \varepsilon \in \mathbb{R} \to \varepsilon \in \mathbb{R}$\lbl{A1}}
\AXC{$\stackrel{\Pi_{98}}{\vdots}$}\noLine
\UIC{$\varphi_{155}$}
\BIC{$2^{-1} \in \mathbb{R} \to \varphi_{52}$}
\BIC{$\varphi_{52}$}
\end{ptree}
\noindent$\Pi_{100}$:
\begin{ptree}
\AXC{$\forall a\,\forall \alpha\,(\varphi_{177})$\lbl{[C9]}}
\AXC{$\forall a\,\forall \alpha\,(\varphi_{177}) \to \forall \alpha\,(\varphi_{177})$\lbl{A2}}
\BIC{$\forall \alpha\,(\varphi_{177})$}
\AXC{$\forall \alpha\,(\varphi_{177}) \to \varphi_{177}$\lbl{A2}}
\BIC{$\varphi_{177}$}
\AXC{$\varphi_{177} \to \varphi_{174}$\lbl{A2}}
\BIC{$\varphi_{174}$}
\AXC{$\varphi_{174} \to \varphi_{164}$\lbl{A2}}
\BIC{$\varphi_{164}$}
\end{ptree}
\noindent$\Pi_{101}$:
\begin{ptree}
\AXC{$\mathrm{Cgt}(a) \to \mathrm{Conv}(a,\alpha) \to \mathrm{Conv}(a,\alpha)$\lbl{A1}}
\AXC{$\stackrel{\Pi_{100}}{\vdots}$}\noLine
\UIC{$\varphi_{164}$}
\AXC{$\varphi_{164} \to \mathrm{Conv}(a,\alpha) \to \varphi_{194}$\lbl{A2}}
\BIC{$\mathrm{Conv}(a,\alpha) \to \varphi_{194}$}
\AXC{$\varphi_{344}$\lbl{A1}}
\BIC{$(\mathrm{Cgt}(a) \to \mathrm{Conv}(a,\alpha) \to \mathrm{Conv}(a,\alpha)) \to \varphi_{285}$}
\BIC{$\varphi_{285}$}
\end{ptree}
\noindent$\Pi_{102}$:
\begin{ptree}
\AXC{$\stackrel{\Pi_{99}}{\vdots}$}\noLine
\UIC{$\varphi_{52}$}
\AXC{$\mathrm{Cgt}(a) \to \mathrm{Conv}(a,\alpha) \to \mathrm{Conv}(a,\alpha)$\lbl{A1}}
\AXC{$\stackrel{\Pi_{101}}{\vdots}$}\noLine
\UIC{$\varphi_{285}$}
\BIC{$\varphi_{52} \to \varphi_{42} \to \varphi_{181}$}
\BIC{$\varphi_{42} \to \varphi_{181}$}
\end{ptree}
\noindent$\Pi_{103}$:
\begin{ptree}
\AXC{$\mathrm{Cgt}(a) \to \mathrm{Conv}(a,\alpha) \to \varepsilon \in \mathbb{R} \to 0 < \varepsilon \to 0 < \varepsilon$\lbl{A1}}
\AXC{$\stackrel{\Pi_{97}}{\vdots}$}\noLine
\UIC{$\varphi_{132}$}
\BIC{$\varphi_{42}$}
\AXC{$\stackrel{\Pi_{102}}{\vdots}$}\noLine
\UIC{$\varphi_{42} \to \varphi_{181}$}
\BIC{$\varphi_{181}$}
\AXC{$\varphi_{365}$\lbl{A1}}
\BIC{$\forall N\,(\varphi_{154} \to \neg (N \in \mathbb{N} \land \varphi_{92})) \to \varphi_{342} \to (\varphi_{154} \to \forall N\,(\varphi_{154})) \to \varphi_{151}$}
\end{ptree}
\noindent$\Pi_{104}$:
\begin{ptree}
\AXC{$\varphi_{154} \to \forall N\,(\varphi_{154})$\lbl{A4}}
\AXC{$\varphi_{342}$\lbl{A3}}
\AXC{$\stackrel{\Pi_{96}}{\vdots}$}\noLine
\UIC{$\forall N\,(\varphi_{154} \to \neg (N \in \mathbb{N} \land \varphi_{92}))$}
\AXC{$\stackrel{\Pi_{103}}{\vdots}$}\noLine
\UIC{$\forall N\,(\varphi_{154} \to \neg (N \in \mathbb{N} \land \varphi_{92})) \to \varphi_{342} \to (\varphi_{154} \to \forall N\,(\varphi_{154})) \to \varphi_{151}$}
\BIC{$\varphi_{342} \to (\varphi_{154} \to \forall N\,(\varphi_{154})) \to \varphi_{151}$}
\BIC{$(\varphi_{154} \to \forall N\,(\varphi_{154})) \to \varphi_{151}$}
\BIC{$\varphi_{151}$}
\end{ptree}
\noindent$\Pi_{105}$:
\begin{ptree}
\AXC{$\varphi_{18}$\lbl{A4}}
\AXC{$\forall \varepsilon\,(\varphi_{152}) \to \varphi_{163}$\lbl{A3}}
\AXC{$\stackrel{\Pi_{104}}{\vdots}$}\noLine
\UIC{$\varphi_{151}$}
\AXC{$\varphi_{151} \to \varphi_{152}$\lbl{A1}}
\BIC{$\varphi_{152}$}
\UIC{$\forall \varepsilon\,(\varphi_{152})$}
\AXC{$\forall \varepsilon\,(\varphi_{152}) \to (\forall \varepsilon\,(\varphi_{152}) \to \varphi_{163}) \to \varphi_{18} \to \mathrm{Cgt}(a) \to \varphi_{146}$\lbl{A1}}
\BIC{$(\forall \varepsilon\,(\varphi_{152}) \to \varphi_{163}) \to \varphi_{18} \to \mathrm{Cgt}(a) \to \varphi_{146}$}
\BIC{$\varphi_{18} \to \mathrm{Cgt}(a) \to \varphi_{146}$}
\BIC{$\mathrm{Cgt}(a) \to \varphi_{146}$}
\end{ptree}
\noindent$\Pi_{106}$:
\begin{ptree}
\AXC{$\forall a\,(\varphi_{156})$\lbl{DefCauchy}}
\AXC{$\forall a\,(\varphi_{156}) \to \varphi_{156}$\lbl{A2}}
\BIC{$\varphi_{156}$}
\AXC{$\varphi_{156} \to (\mathrm{Cgt}(a) \to \mathrm{Conv}(a,\alpha) \to \mathrm{Seq}(a)) \to \varphi_{202}$\lbl{A1}}
\BIC{$(\mathrm{Cgt}(a) \to \mathrm{Conv}(a,\alpha) \to \mathrm{Seq}(a)) \to \varphi_{202}$}
\end{ptree}
\noindent$\Pi_{107}$:
\begin{ptree}
\AXC{$\stackrel{\Pi_{105}}{\vdots}$}\noLine
\UIC{$\mathrm{Cgt}(a) \to \varphi_{146}$}
\AXC{$\stackrel{\Pi_{69}}{\vdots}$}\noLine
\UIC{$\mathrm{Cgt}(a) \to \mathrm{Conv}(a,\alpha) \to \mathrm{Seq}(a)$}
\AXC{$\stackrel{\Pi_{106}}{\vdots}$}\noLine
\UIC{$(\mathrm{Cgt}(a) \to \mathrm{Conv}(a,\alpha) \to \mathrm{Seq}(a)) \to \varphi_{202}$}
\BIC{$\varphi_{202}$}
\BIC{$\mathrm{Cgt}(a) \to \mathrm{Conv}(a,\alpha) \to \mathrm{Cauchy}(a)$}
\end{ptree}
\noindent$\Pi_{108}$:
\begin{ptree}
\AXC{$\mathrm{Cgt}(a) \to \mathrm{Cgt}(a)$\lbl{A1}}
\AXC{$(\mathrm{Cgt}(a) \to \mathrm{Cgt}(a)) \to \varphi_{20}$\lbl{A1}}
\BIC{$\varphi_{20}$}
\end{ptree}
\noindent$\Pi_{109}$:
\begin{ptree}
\AXC{$\forall a\,(\mathrm{Cgt}(a) \leftrightarrow \exists \alpha\,(\mathrm{Conv}(a,\alpha)))$\lbl{DefCgt}}
\AXC{$\varphi_{23}$\lbl{A2}}
\BIC{$\mathrm{Cgt}(a) \leftrightarrow \exists \alpha\,(\mathrm{Conv}(a,\alpha))$}
\AXC{$\stackrel{\Pi_{108}}{\vdots}$}\noLine
\UIC{$\varphi_{20}$}
\BIC{$\mathrm{Cgt}(a) \to \exists \alpha\,(\mathrm{Conv}(a,\alpha))$}
\AXC{$\varphi_{227}$\lbl{A1}}
\BIC{$\varphi_{5} \to (\varphi_{5} \to \varphi_{10}) \to \varphi_{16} \to \mathrm{Cgt}(a) \to \mathrm{Cauchy}(a)$}
\end{ptree}
\noindent$\Pi_{110}$:
\begin{ptree}
\AXC{$\stackrel{\Pi_{107}}{\vdots}$}\noLine
\UIC{$\mathrm{Cgt}(a) \to \mathrm{Conv}(a,\alpha) \to \mathrm{Cauchy}(a)$}
\AXC{$\varphi_{45}$\lbl{A1}}
\BIC{$\mathrm{Cgt}(a) \land \neg \mathrm{Cauchy}(a) \to \neg \mathrm{Conv}(a,\alpha)$}
\UIC{$\varphi_{5}$}
\AXC{$\stackrel{\Pi_{109}}{\vdots}$}\noLine
\UIC{$\varphi_{5} \to (\varphi_{5} \to \varphi_{10}) \to \varphi_{16} \to \mathrm{Cgt}(a) \to \mathrm{Cauchy}(a)$}
\BIC{$(\varphi_{5} \to \varphi_{10}) \to \varphi_{16} \to \mathrm{Cgt}(a) \to \mathrm{Cauchy}(a)$}
\end{ptree}
\noindent$\Pi_{111}$:
\begin{ptree}
\AXC{$\varphi_{16}$\lbl{A4}}
\AXC{$\varphi_{5} \to \varphi_{10}$\lbl{A3}}
\AXC{$\stackrel{\Pi_{110}}{\vdots}$}\noLine
\UIC{$(\varphi_{5} \to \varphi_{10}) \to \varphi_{16} \to \mathrm{Cgt}(a) \to \mathrm{Cauchy}(a)$}
\BIC{$\varphi_{16} \to \mathrm{Cgt}(a) \to \mathrm{Cauchy}(a)$}
\BIC{$\mathrm{Cgt}(a) \to \mathrm{Cauchy}(a)$}
\AXC{$(\mathrm{Cgt}(a) \to \mathrm{Cauchy}(a)) \to \varphi_{7}$\lbl{A1}}
\BIC{$\varphi_{7}$}
\end{ptree}
\noindent 主木:
\begin{ptree}
\AXC{$\varphi_{14}$\lbl{A4}}
\AXC{$\forall b\,(\varphi_{12}) \to \varphi_{21}$\lbl{A3}}
\AXC{$\stackrel{\Pi_{61}}{\vdots}$}\noLine
\UIC{$(\forall b\,(\varphi_{12}) \to \varphi_{21}) \to \varphi_{14} \to \mathrm{Cauchy}(a) \to \mathrm{Cgt}(a)$}
\BIC{$\varphi_{14} \to \mathrm{Cauchy}(a) \to \mathrm{Cgt}(a)$}
\BIC{$\mathrm{Cauchy}(a) \to \mathrm{Cgt}(a)$}
\AXC{$\stackrel{\Pi_{111}}{\vdots}$}\noLine
\UIC{$\varphi_{7}$}
\BIC{$\mathrm{Cgt}(a) \leftrightarrow \mathrm{Cauchy}(a)$}
\UIC{$\forall a\,(\mathrm{Cgt}(a) \leftrightarrow \mathrm{Cauchy}(a))$}
\end{ptree}
\begin{abbrevlist}
\item $\varphi_{1} :\equiv \forall c\,(\neg \mathrm{Conv}(b,c)) \to \forall \alpha\,\forall c\,(\neg \mathrm{Conv}(b,c))$
\item $\varphi_{2} :\equiv \mathrm{SeqBdd}(a) \to \exists b\,(\mathrm{SubSeq}(b,a) \land \mathrm{Cgt}(b))$
\item $\varphi_{3} :\equiv \mathrm{Cauchy}(a) \to \exists b\,(\mathrm{SubSeq}(b,a) \land \mathrm{Cgt}(b))$
\item $\varphi_{4} :\equiv \forall \alpha\,\forall c\,(\neg \mathrm{Conv}(b,c)) \to \forall \alpha\,(\neg \mathrm{Conv}(b,\alpha))$
\item $\varphi_{5} :\equiv \forall \alpha\,(\mathrm{Cgt}(a) \land \neg \mathrm{Cauchy}(a) \to \neg \mathrm{Conv}(a,\alpha))$
\item $\varphi_{6} :\equiv \mathrm{SubSeq}(b,a) \land \mathrm{Cgt}(b) \to \exists \alpha\,(\mathrm{Conv}(b,\alpha))$
\item $\varphi_{7} :\equiv (\mathrm{Cauchy}(a) \to \mathrm{Cgt}(a)) \to (\mathrm{Cgt}(a) \leftrightarrow \mathrm{Cauchy}(a))$
\item $\varphi_{8} :\equiv \mathrm{Cauchy}(a) \to \mathrm{SubSeq}(b,a) \land \mathrm{Cgt}(b) \to \mathrm{Cgt}(a)$
\item $\varphi_{9} :\equiv \mathrm{SubSeq}(b,a) \land \mathrm{Cgt}(b) \to \mathrm{Conv}(b,c) \to \mathrm{Cgt}(a)$
\item $\varphi_{10} :\equiv \forall \alpha\,(\mathrm{Cgt}(a) \land \neg \mathrm{Cauchy}(a)) \to \forall \alpha\,(\neg \mathrm{Conv}(a,\alpha))$
\item $\varphi_{11} :\equiv \mathrm{Cauchy}(a) \land (\mathrm{SubSeq}(b,a) \land \mathrm{Cgt}(b)) \land \neg \mathrm{Cgt}(a)$
\item $\varphi_{12} :\equiv \mathrm{Cauchy}(a) \land \neg \mathrm{Cgt}(a) \to \neg (\mathrm{SubSeq}(b,a) \land \mathrm{Cgt}(b))$
\item $\varphi_{13} :\equiv \mathrm{SubSeq}(b,a) \land \mathrm{Cgt}(b) \to \mathrm{SubSeq}(b,a) \land \mathrm{Cgt}(b)$
\item $\varphi_{14} :\equiv \mathrm{Cauchy}(a) \land \neg \mathrm{Cgt}(a) \to \forall b\,(\mathrm{Cauchy}(a) \land \neg \mathrm{Cgt}(a))$
\item $\varphi_{15} :\equiv \mathrm{SubSeq}(b,a) \land \mathrm{Cgt}(b) \to \mathrm{Conv}(b,c) \to \mathrm{Conv}(b,c)$
\item $\varphi_{16} :\equiv \mathrm{Cgt}(a) \land \neg \mathrm{Cauchy}(a) \to \forall \alpha\,(\mathrm{Cgt}(a) \land \neg \mathrm{Cauchy}(a))$
\item $\varphi_{17} :\equiv \forall \alpha\,(\mathrm{Conv}(a,\alpha) \to \mathrm{Seq}(a)) \to \mathrm{Conv}(a,\alpha) \to \mathrm{Seq}(a)$
\item $\varphi_{18} :\equiv \mathrm{Cgt}(a) \land \mathrm{Conv}(a,\alpha) \to \forall \varepsilon\,(\mathrm{Cgt}(a) \land \mathrm{Conv}(a,\alpha))$
\item $\varphi_{19} :\equiv \forall a\,(\mathrm{Cauchy}(a) \to \mathrm{SeqBdd}(a)) \to \mathrm{Cauchy}(a) \to \mathrm{SeqBdd}(a)$
\item $\varphi_{20} :\equiv (\mathrm{Cgt}(a) \leftrightarrow \exists \alpha\,(\mathrm{Conv}(a,\alpha))) \to \mathrm{Cgt}(a) \to \exists \alpha\,(\mathrm{Conv}(a,\alpha))$
\item $\varphi_{21} :\equiv \forall b\,(\mathrm{Cauchy}(a) \land \neg \mathrm{Cgt}(a)) \to \forall b\,(\neg (\mathrm{SubSeq}(b,a) \land \mathrm{Cgt}(b)))$
\item $\varphi_{22} :\equiv \forall x\,\forall y\,\forall z\,(x \leq y \to y \leq z \to x \leq z) \to \forall y\,\forall z\,(N_{1} \leq y \to y \leq z \to N_{1} \leq z)$
\item $\varphi_{23} :\equiv \forall a\,(\mathrm{Cgt}(a) \leftrightarrow \exists \alpha\,(\mathrm{Conv}(a,\alpha))) \to (\mathrm{Cgt}(a) \leftrightarrow \exists \alpha\,(\mathrm{Conv}(a,\alpha)))$
\item $\varphi_{24} :\equiv \forall a\,(\mathrm{Cgt}(a) \leftrightarrow \exists \alpha\,(\mathrm{Conv}(a,\alpha))) \to (\mathrm{Cgt}(b) \leftrightarrow \exists \alpha\,(\mathrm{Conv}(b,\alpha)))$
\item $\varphi_{25} :\equiv \forall a\,\forall \alpha\,(\mathrm{Conv}(a,\alpha) \to \mathrm{Seq}(a)) \to \forall \alpha\,(\mathrm{Conv}(a,\alpha) \to \mathrm{Seq}(a))$
\item $\varphi_{26} :\equiv \forall \alpha\,(\mathrm{Conv}(b,\alpha) \to \alpha \in \mathbb{R}) \to \mathrm{Conv}(b,c) \to c \in \mathbb{R}$
\item $\varphi_{27} :\equiv \mathrm{Cauchy}(a) \land (\mathrm{SubSeq}(b,a) \land \mathrm{Cgt}(b)) \land \mathrm{Conv}(b,c) \land \neg \mathrm{Cgt}(a)$
\item $\varphi_{28} :\equiv \forall \alpha\,(\mathrm{Conv}(a,\alpha) \to \alpha \in \mathbb{R}) \to \mathrm{Conv}(a,\alpha) \to \alpha \in \mathbb{R}$
\item $\varphi_{29} :\equiv \forall s\,(\neg (\mathrm{SubIdx}(s) \land \forall k \in \mathbb{N}\,(b_{k} = a_{s_{k}})))$
\item $\varphi_{30} :\equiv \mathrm{SubIdx}(s) \leftrightarrow \mathrm{NSeq}(s) \land \forall k \in \mathbb{N}\,(s_{k} < s_{S(k)})$
\item $\varphi_{31} :\equiv \mathrm{SubIdx}(s) \land \forall k \in \mathbb{N}\,(b_{k} = a_{s_{k}}) \to \mathrm{Cgt}(a)$
\item $\varphi_{32} :\equiv \mathrm{SubIdx}(s) \land \forall k \in \mathbb{N}\,(b_{k} = a_{s_{k}}) \to \mathrm{NSeq}(s)$
\item $\varphi_{33} :\equiv \lvert v_{749} - v_{751} \rvert < v_{752} \to \lvert v_{750} - v_{751} \rvert < v_{752} + v_{752}$
\item $\varphi_{34} :\equiv \forall a\,\forall \alpha\,(\mathrm{Conv}(a,\alpha) \to \alpha \in \mathbb{R}) \to \forall \alpha\,(\mathrm{Conv}(b,\alpha) \to \alpha \in \mathbb{R})$
\item $\varphi_{35} :\equiv \forall a\,\forall \alpha\,(\mathrm{Conv}(a,\alpha) \to \alpha \in \mathbb{R}) \to \forall \alpha\,(\mathrm{Conv}(a,\alpha) \to \alpha \in \mathbb{R})$
\item $\varphi_{36} :\equiv \forall N\,(\neg (N \in \mathbb{N} \land \forall n \in \mathbb{N}\,(N \leq n \to \lvert a_{n} - c \rvert < \varepsilon)))$
\item $\varphi_{37} :\equiv (\mathrm{Cgt}(a) \to \mathrm{Conv}(a,\alpha) \to \mathrm{Conv}(a,\alpha)) \to \mathrm{Cgt}(a) \to \mathrm{Conv}(a,\alpha) \to \mathrm{Seq}(a)$
\item $\varphi_{38} :\equiv \mathrm{SubIdx}(s) \land \forall k \in \mathbb{N}\,(b_{k} = a_{s_{k}}) \to \mathrm{SubIdx}(s)$
\item $\varphi_{39} :\equiv \mathrm{SubIdx}(s) \land \forall k \in \mathbb{N}\,(b_{k} = a_{s_{k}}) \to \mathrm{Conv}(a,c)$
\item $\varphi_{40} :\equiv \mathrm{Seq}(b) \land \exists s\,(\mathrm{SubIdx}(s) \land \forall k \in \mathbb{N}\,(b_{k} = a_{s_{k}}))$
\item $\varphi_{41} :\equiv \forall y\,\forall z\,(N_{1} \leq y \to y \leq z \to N_{1} \leq z) \to \forall z\,(N_{1} \leq N_{1} \cup N_{2} \to N_{1} \cup N_{2} \leq z \to N_{1} \leq z)$
\item $\varphi_{42} :\equiv \mathrm{Cgt}(a) \to \mathrm{Conv}(a,\alpha) \to \varepsilon \in \mathbb{R} \to 0 < \varepsilon \to 0 < \varepsilon \cdot 2^{-1}$
\item $\varphi_{43} :\equiv \mathrm{Conv}(b,c) \to \exists s\,(\mathrm{SubIdx}(s) \land \forall k \in \mathbb{N}\,(b_{k} = a_{s_{k}}))$
\item $\varphi_{44} :\equiv N_{1} \leq N_{1} \cup N_{2} \to N_{1} \cup N_{2} \leq s_{N_{1} \cup N_{2}} \to N_{1} \leq s_{N_{1} \cup N_{2}}$
\item $\varphi_{45} :\equiv (\mathrm{Cgt}(a) \to \mathrm{Conv}(a,\alpha) \to \mathrm{Cauchy}(a)) \to \mathrm{Cgt}(a) \land \neg \mathrm{Cauchy}(a) \to \neg \mathrm{Conv}(a,\alpha)$
\item $\varphi_{46} :\equiv \forall w\,(z = w \to \lvert z - c \rvert < \varepsilon \cdot 2^{-1} \to \lvert w - c \rvert < \varepsilon \cdot 2^{-1})$
\item $\varphi_{47} :\equiv (\mathrm{Cgt}(a) \to \mathrm{Conv}(a,\alpha) \to \mathrm{Conv}(a,\alpha)) \to \mathrm{Cgt}(a) \to \mathrm{Conv}(a,\alpha) \to \alpha \in \mathbb{R}$
\item $\varphi_{48} :\equiv N_{1} \leq v_{753} \to N_{1} \leq v_{754} \to \lvert a_{v_{753}} - a_{v_{754}} \rvert < \varepsilon \cdot 2^{-1}$
\item $\varphi_{49} :\equiv \mathrm{SubIdx}(s) \land \forall k \in \mathbb{N}\,(b_{k} = a_{s_{k}}) \to \exists \alpha\,(\mathrm{Conv}(a,\alpha))$
\item $\varphi_{50} :\equiv \forall \alpha\,(\forall c\,(\neg \mathrm{Conv}(b,c)) \to \neg \mathrm{Conv}(b,\alpha)) \to \varphi_{4}$
\item $\varphi_{51} :\equiv \neg (N \in \mathbb{N} \land \forall n \in \mathbb{N}\,(N \leq n \to \lvert b_{n} - c \rvert < \varepsilon \cdot 2^{-1}))$
\item $\varphi_{52} :\equiv \mathrm{Cgt}(a) \to \mathrm{Conv}(a,\alpha) \to \varepsilon \in \mathbb{R} \to 0 < \varepsilon \to \varepsilon \cdot 2^{-1} \in \mathbb{R}$
\item $\varphi_{53} :\equiv N_{1} \leq n \to N_{1} \cup N_{2} \in \mathbb{N} \to b_{N_{1} \cup N_{2}} = a_{s_{N_{1} \cup N_{2}}}$
\item $\varphi_{54} :\equiv m \in \mathbb{N} \to n \in \mathbb{N} \to N \leq m \to N \leq n \to \varepsilon \cdot 2^{-1} + \varepsilon \cdot 2^{-1} = \varepsilon$
\item $\varphi_{55} :\equiv \neg (N_{2} \in \mathbb{N} \land \forall n \in \mathbb{N}\,(N_{2} \leq n \to \lvert b_{n} - c \rvert < \varepsilon \cdot 2^{-1}))$
\item $\varphi_{56} :\equiv \forall n\,(m \in \mathbb{N} \to n \in \mathbb{N} \to N \leq m \to N \leq n \to \lvert a_{m} - a_{n} \rvert < \varepsilon)$
\item $\varphi_{57} :\equiv (\mathrm{Cgt}(b) \leftrightarrow \exists \alpha\,(\mathrm{Conv}(b,\alpha))) \to \mathrm{Cauchy}(a) \to \varphi_{6}$
\item $\varphi_{58} :\equiv m \in \mathbb{N} \to n \in \mathbb{N} \to N \leq m \to N \leq n \to \lvert a_{n} - \alpha \rvert = \lvert \alpha - a_{n} \rvert$
\item $\varphi_{59} :\equiv m \in \mathbb{N} \to n \in \mathbb{N} \to N \leq m \to N \leq n \to \lvert a_{m} - \alpha \rvert = \lvert \alpha - a_{m} \rvert$
\item $\varphi_{60} :\equiv n \in \mathbb{N} \to N_{1} \leq n \to \forall n \in \mathbb{N}\,(N_{2} \leq n \to \lvert b_{n} - c \rvert < \varepsilon \cdot 2^{-1})$
\item $\varphi_{61} :\equiv \mathrm{Conv}(b,c) \to \mathrm{SubIdx}(s) \land \forall k \in \mathbb{N}\,(b_{k} = a_{s_{k}}) \to c \in \mathbb{R}$
\item $\varphi_{62} :\equiv \forall k \in \mathbb{N}\,(k \leq s_{k}) \to N_{1} \cup N_{2} \in \mathbb{N} \to N_{1} \cup N_{2} \leq s_{N_{1} \cup N_{2}}$
\item $\varphi_{63} :\equiv \mathrm{SubIdx}(s) \land \forall k \in \mathbb{N}\,(b_{k} = a_{s_{k}}) \to \forall k \in \mathbb{N}\,(k \leq s_{k})$
\item $\varphi_{64} :\equiv \forall s\,(\mathrm{SubIdx}(s) \to \forall k \in \mathbb{N}\,(k \leq s_{k})) \to \mathrm{SubIdx}(s) \to \forall k \in \mathbb{N}\,(k \leq s_{k})$
\item $\varphi_{65} :\equiv N_{1} \cup N_{2} \in \mathbb{N} \to N_{2} \leq N_{1} \cup N_{2} \to \lvert b_{N_{1} \cup N_{2}} - c \rvert < \varepsilon \cdot 2^{-1}$
\item $\varphi_{66} :\equiv N_{1} \leq s_{N_{1} \cup N_{2}} \to \lvert a_{n} - a_{s_{N_{1} \cup N_{2}}} \rvert < \varepsilon \cdot 2^{-1}$
\item $\varphi_{67} :\equiv \mathrm{SubIdx}(s) \land \forall k \in \mathbb{N}\,(b_{k} = a_{s_{k}}) \to \varepsilon \in \mathbb{R} \to \varepsilon \in \mathbb{R}$
\item $\varphi_{68} :\equiv \forall n\,(\mathrm{Seq}(a) \to n \in \mathbb{N} \to a_{n} \in \mathbb{R}) \to \mathrm{Seq}(a) \to n \in \mathbb{N} \to a_{n} \in \mathbb{R}$
\item $\varphi_{69} :\equiv \forall n\,(\mathrm{Seq}(a) \to n \in \mathbb{N} \to a_{n} \in \mathbb{R}) \to \mathrm{Seq}(a) \to m \in \mathbb{N} \to a_{m} \in \mathbb{R}$
\item $\varphi_{70} :\equiv N_{2} \in \mathbb{N} \land \forall n \in \mathbb{N}\,(N_{2} \leq n \to \lvert b_{n} - c \rvert < \varepsilon \cdot 2^{-1}) \to N_{2} \in \mathbb{N}$
\item $\varphi_{71} :\equiv N_{2} \in \mathbb{N} \land \forall n \in \mathbb{N}\,(N_{2} \leq n \to \lvert b_{n} - c \rvert < \varepsilon \cdot 2^{-1}) \to N_{1} \in \mathbb{N}$
\item $\varphi_{72} :\equiv n \in \mathbb{N} \to N \leq m \to N \leq n \to \lvert a_{m} - a_{n} \rvert < \varepsilon \cdot 2^{-1} + \varepsilon \cdot 2^{-1}$
\item $\varphi_{73} :\equiv \forall w\,(\lvert a_{n} - \alpha \rvert = w \to \lvert a_{n} - \alpha \rvert < \varepsilon \cdot 2^{-1} \to w < \varepsilon \cdot 2^{-1})$
\item $\varphi_{74} :\equiv \forall w\,(\lvert a_{m} - \alpha \rvert = w \to \lvert a_{m} - \alpha \rvert < \varepsilon \cdot 2^{-1} \to w < \varepsilon \cdot 2^{-1})$
\item $\varphi_{75} :\equiv \forall m\,\forall n\,(m \in \mathbb{N} \to n \in \mathbb{N} \to N_{1} \leq m \to N_{1} \leq n \to \lvert a_{m} - a_{n} \rvert < \varepsilon \cdot 2^{-1})$
\item $\varphi_{76} :\equiv \varepsilon \in \mathbb{R} \to 0 < \varepsilon \to \exists N \in \mathbb{N}\,\forall n \in \mathbb{N}\,(N \leq n \to \lvert a_{n} - c \rvert < \varepsilon)$
\item $\varphi_{77} :\equiv \forall s\,\forall k\,(\mathrm{NSeq}(s) \to k \in \mathbb{N} \to s_{k} \in \mathbb{N}) \to \forall k\,(\mathrm{NSeq}(s) \to k \in \mathbb{N} \to s_{k} \in \mathbb{N})$
\item $\varphi_{78} :\equiv \forall a\,\forall n\,(\mathrm{Seq}(a) \to n \in \mathbb{N} \to a_{n} \in \mathbb{R}) \to \forall n\,(\mathrm{Seq}(a) \to n \in \mathbb{N} \to a_{n} \in \mathbb{R})$
\item $\varphi_{79} :\equiv \forall e \in \mathbb{R}\,(0 < e \to 0 < e \cdot 2^{-1}) \to \varepsilon \in \mathbb{R} \to 0 < \varepsilon \to 0 < \varepsilon \cdot 2^{-1}$
\item $\varphi_{80} :\equiv \mathrm{SubIdx}(s) \land \forall k \in \mathbb{N}\,(b_{k} = a_{s_{k}}) \to \varepsilon \in \mathbb{R} \to 0 < \varepsilon \to 0 < \varepsilon$
\item $\varphi_{81} :\equiv \forall n\,(N_{1} \in \mathbb{N} \to n \in \mathbb{N} \to N_{1} \cup n \in \mathbb{N}) \to N_{1} \in \mathbb{N} \to N_{2} \in \mathbb{N} \to N_{1} \cup N_{2} \in \mathbb{N}$
\item $\varphi_{82} :\equiv \forall \varepsilon \in \mathbb{R}\,(0 < \varepsilon \to \exists N \in \mathbb{N}\,\forall n \in \mathbb{N}\,(N \leq n \to \lvert a_{n} - c \rvert < \varepsilon))$
\item $\varphi_{83} :\equiv \forall n\,(v_{753} \in \mathbb{N} \to n \in \mathbb{N} \to N_{1} \leq v_{753} \to N_{1} \leq n \to \lvert a_{v_{753}} - a_{n} \rvert < \varepsilon \cdot 2^{-1})$
\item $\varphi_{84} :\equiv 0 < \varepsilon \cdot 2^{-1} \to \exists N \in \mathbb{N}\,\forall n \in \mathbb{N}\,(N \leq n \to \lvert b_{n} - c \rvert < \varepsilon \cdot 2^{-1})$
\item $\varphi_{85} :\equiv n \in \mathbb{N} \to N \leq n \to \lvert a_{n} - \alpha \rvert < \varepsilon \cdot 2^{-1} \land \lvert a_{n} - \alpha \rvert < \varepsilon \cdot 2^{-1}$
\item $\varphi_{86} :\equiv m \in \mathbb{N} \to N \leq m \to \lvert a_{m} - \alpha \rvert < \varepsilon \cdot 2^{-1} \land \lvert a_{m} - \alpha \rvert < \varepsilon \cdot 2^{-1}$
\item $\varphi_{87} :\equiv \forall m\,\forall n\,(m \in \mathbb{N} \to n \in \mathbb{N} \to m \cup n \in \mathbb{N}) \to \forall n\,(N_{1} \in \mathbb{N} \to n \in \mathbb{N} \to N_{1} \cup n \in \mathbb{N})$
\item $\varphi_{88} :\equiv N_{2} \in \mathbb{N} \land \forall n \in \mathbb{N}\,(N_{2} \leq n \to \lvert b_{n} - c \rvert < \varepsilon \cdot 2^{-1}) \to n \in \mathbb{N} \to n \in \mathbb{N}$
\item $\varphi_{89} :\equiv \forall v_{754}\,(n \in \mathbb{N} \to v_{754} \in \mathbb{N} \to N_{1} \leq n \to N_{1} \leq v_{754} \to \lvert a_{n} - a_{v_{754}} \rvert < \varepsilon \cdot 2^{-1})$
\item $\varphi_{90} :\equiv \forall e \in \mathbb{R}\,(e \cdot 2^{-1} + e \cdot 2^{-1} = e) \to \varepsilon \in \mathbb{R} \to \varepsilon \cdot 2^{-1} + \varepsilon \cdot 2^{-1} = \varepsilon$
\item $\varphi_{91} :\equiv N_{2} \in \mathbb{N} \land \forall n \in \mathbb{N}\,(N_{2} \leq n \to \lvert b_{n} - c \rvert < \varepsilon \cdot 2^{-1}) \to n \in \mathbb{N} \to N_{1} \leq n \to N_{1} \leq n$
\item $\varphi_{92} :\equiv \forall n \in \mathbb{N}\,(N \leq n \to \lvert a_{n} - \alpha \rvert < \varepsilon \cdot 2^{-1} \land \lvert a_{n} - \alpha \rvert < \varepsilon \cdot 2^{-1})$
\item $\varphi_{93} :\equiv \neg (N \in \mathbb{N} \land \forall m\,\forall n\,(m \in \mathbb{N} \to n \in \mathbb{N} \to N \leq m \to N \leq n \to \lvert a_{m} - a_{n} \rvert < \varepsilon \cdot 2^{-1}))$
\item $\varphi_{94} :\equiv \forall x\,\forall y\,(x \in \mathbb{R} \to y \in \mathbb{R} \to x \cdot y \in \mathbb{R}) \to \forall y\,(\varepsilon \in \mathbb{R} \to y \in \mathbb{R} \to \varepsilon \cdot y \in \mathbb{R})$
\item $\varphi_{95} :\equiv \mathrm{Cauchy}(a) \land (\mathrm{SubSeq}(b,a) \land \mathrm{Cgt}(b)) \land \mathrm{Conv}(b,c) \land (\mathrm{SubIdx}(s) \land \forall k \in \mathbb{N}\,(b_{k} = a_{s_{k}}))$
\item $\varphi_{96} :\equiv \lvert b_{N_{1} \cup N_{2}} - c \rvert < \varepsilon \cdot 2^{-1} \to \lvert a_{s_{N_{1} \cup N_{2}}} - c \rvert < \varepsilon \cdot 2^{-1}$
\item $\varphi_{97} :\equiv \mathrm{SubIdx}(s) \land \forall k \in \mathbb{N}\,(b_{k} = a_{s_{k}}) \to \varepsilon \in \mathbb{R} \to 0 < \varepsilon \to 0 < \varepsilon \cdot 2^{-1}$
\item $\varphi_{98} :\equiv \lvert \alpha - a_{n} \rvert < \varepsilon \cdot 2^{-1} \to \lvert a_{m} - a_{n} \rvert < \varepsilon \cdot 2^{-1} + \varepsilon \cdot 2^{-1}$
\item $\varphi_{99} :\equiv \exists N \in \mathbb{N}\,\forall m\,\forall n\,(m \in \mathbb{N} \to n \in \mathbb{N} \to N \leq m \to N \leq n \to \lvert a_{m} - a_{n} \rvert < \varepsilon \cdot 2^{-1})$
\item $\varphi_{100} :\equiv \forall w\,(b_{N_{1} \cup N_{2}} = w \to \lvert b_{N_{1} \cup N_{2}} - c \rvert < \varepsilon \cdot 2^{-1} \to \lvert w - c \rvert < \varepsilon \cdot 2^{-1})$
\item $\varphi_{101} :\equiv \forall k\,(\mathrm{NSeq}(s) \to k \in \mathbb{N} \to s_{k} \in \mathbb{N}) \to \mathrm{NSeq}(s) \to N_{1} \cup N_{2} \in \mathbb{N} \to s_{N_{1} \cup N_{2}} \in \mathbb{N}$
\item $\varphi_{102} :\equiv \mathrm{SubIdx}(s) \land \forall k \in \mathbb{N}\,(b_{k} = a_{s_{k}}) \to \varepsilon \in \mathbb{R} \to 0 < \varepsilon \to \varepsilon \cdot 2^{-1} \in \mathbb{R}$
\item $\varphi_{103} :\equiv \mathrm{SubIdx}(s) \land \forall k \in \mathbb{N}\,(b_{k} = a_{s_{k}}) \to N_{1} \cup N_{2} \in \mathbb{N} \to N_{1} \cup N_{2} \leq s_{N_{1} \cup N_{2}}$
\item $\varphi_{104} :\equiv \forall \varepsilon\,(\mathrm{Conv}(b,c) \to \varepsilon \in \mathbb{R} \to 0 < \varepsilon \to \exists N \in \mathbb{N}\,\forall n \in \mathbb{N}\,(N \leq n \to \lvert b_{n} - c \rvert < \varepsilon))$
\item $\varphi_{105} :\equiv N_{2} \in \mathbb{N} \land \forall n \in \mathbb{N}\,(N_{2} \leq n \to \lvert b_{n} - c \rvert < \varepsilon \cdot 2^{-1}) \to n \in \mathbb{N} \to N_{1} \leq n \to \varepsilon \in \mathbb{R}$
\item $\varphi_{106} :\equiv \forall k \in \mathbb{N}\,(b_{k} = a_{s_{k}}) \to N_{1} \cup N_{2} \in \mathbb{N} \to b_{N_{1} \cup N_{2}} = a_{s_{N_{1} \cup N_{2}}}$
\item $\varphi_{107} :\equiv \mathrm{SubIdx}(s) \land \forall k \in \mathbb{N}\,(b_{k} = a_{s_{k}}) \to \mathrm{SubIdx}(s) \land \forall k \in \mathbb{N}\,(b_{k} = a_{s_{k}})$
\item $\varphi_{108} :\equiv \forall y\,(\varepsilon \in \mathbb{R} \to y \in \mathbb{R} \to \varepsilon \cdot y \in \mathbb{R}) \to \varepsilon \in \mathbb{R} \to 2^{-1} \in \mathbb{R} \to \varepsilon \cdot 2^{-1} \in \mathbb{R}$
\item $\varphi_{109} :\equiv \lvert a_{s_{N_{1} \cup N_{2}}} - c \rvert < \varepsilon \cdot 2^{-1} \to \lvert a_{n} - c \rvert < \varepsilon \cdot 2^{-1} + \varepsilon \cdot 2^{-1}$
\item $\varphi_{110} :\equiv \mathrm{Cauchy}(a) \to \mathrm{SubSeq}(b,a) \land \mathrm{Cgt}(b) \to \mathrm{Conv}(b,c) \to \varphi_{31}$
\item $\varphi_{111} :\equiv \mathrm{Cauchy}(a) \to \mathrm{SubSeq}(b,a) \land \mathrm{Cgt}(b) \to \mathrm{Conv}(b,c) \to \varphi_{32}$
\item $\varphi_{112} :\equiv N_{2} \in \mathbb{N} \land \forall n \in \mathbb{N}\,(N_{2} \leq n \to \lvert b_{n} - c \rvert < \varepsilon \cdot 2^{-1}) \to n \in \mathbb{N} \to N_{1} \leq n \to a_{n} \in \mathbb{R}$
\item $\varphi_{113} :\equiv \mathrm{Cauchy}(a) \to \mathrm{SubSeq}(b,a) \land \mathrm{Cgt}(b) \to \mathrm{Conv}(b,c) \to \varphi_{38}$
\item $\varphi_{114} :\equiv \forall \alpha\,\forall \varepsilon\,(\mathrm{Conv}(b,\alpha) \to \varepsilon \in \mathbb{R} \to 0 < \varepsilon \to \exists N \in \mathbb{N}\,\forall n \in \mathbb{N}\,(N \leq n \to \lvert b_{n} - \alpha \rvert < \varepsilon))$
\item $\varphi_{115} :\equiv \mathrm{Cauchy}(a) \to \mathrm{SubSeq}(b,a) \land \mathrm{Cgt}(b) \to \mathrm{Conv}(b,c) \to \varphi_{39}$
\item $\varphi_{116} :\equiv N_{2} \in \mathbb{N} \land \forall n \in \mathbb{N}\,(N_{2} \leq n \to \lvert b_{n} - c \rvert < \varepsilon \cdot 2^{-1}) \to n \in \mathbb{N} \to N_{1} \leq n \to N_{1} \cup N_{2} \in \mathbb{N}$
\item $\varphi_{117} :\equiv \forall s\,(\varphi_{27} \to \neg (\mathrm{SubIdx}(s) \land \forall k \in \mathbb{N}\,(b_{k} = a_{s_{k}})))$
\item $\varphi_{118} :\equiv N_{2} \in \mathbb{N} \land \forall n \in \mathbb{N}\,(N_{2} \leq n \to \lvert b_{n} - c \rvert < \varepsilon \cdot 2^{-1}) \to n \in \mathbb{N} \to N_{1} \leq n \to N_{1} \leq s_{N_{1} \cup N_{2}}$
\item $\varphi_{119} :\equiv \forall a\,\forall \alpha\,\forall \varepsilon\,(\mathrm{Conv}(a,\alpha) \to \varepsilon \in \mathbb{R} \to 0 < \varepsilon \to \exists N \in \mathbb{N}\,\forall n \in \mathbb{N}\,(N \leq n \to \lvert a_{n} - \alpha \rvert < \varepsilon))$
\item $\varphi_{120} :\equiv \forall n\,(\mathrm{Seq}(a) \to n \in \mathbb{N} \to a_{n} \in \mathbb{R}) \to \mathrm{Seq}(a) \to s_{N_{1} \cup N_{2}} \in \mathbb{N} \to a_{s_{N_{1} \cup N_{2}}} \in \mathbb{R}$
\item $\varphi_{121} :\equiv \mathrm{Cauchy}(a) \to \mathrm{SubSeq}(b,a) \land \mathrm{Cgt}(b) \to \mathrm{Conv}(b,c) \to \varphi_{49}$
\item $\varphi_{122} :\equiv \forall w\,(\varepsilon \cdot 2^{-1} + \varepsilon \cdot 2^{-1} = w \to \lvert a_{n} - c \rvert < \varepsilon \cdot 2^{-1} + \varepsilon \cdot 2^{-1} \to \lvert a_{n} - c \rvert < w)$
\item $\varphi_{123} :\equiv \varepsilon \cdot 2^{-1} + \varepsilon \cdot 2^{-1} = \varepsilon \to \lvert a_{n} - c \rvert < \varepsilon \cdot 2^{-1} + \varepsilon \cdot 2^{-1} \to \lvert a_{n} - c \rvert < \varepsilon$
\item $\varphi_{124} :\equiv \lvert a_{n} - \alpha \rvert = \lvert \alpha - a_{n} \rvert \to \lvert a_{n} - \alpha \rvert < \varepsilon \cdot 2^{-1} \to \lvert \alpha - a_{n} \rvert < \varepsilon \cdot 2^{-1}$
\item $\varphi_{125} :\equiv \lvert a_{m} - \alpha \rvert = \lvert \alpha - a_{m} \rvert \to \lvert a_{m} - \alpha \rvert < \varepsilon \cdot 2^{-1} \to \lvert \alpha - a_{m} \rvert < \varepsilon \cdot 2^{-1}$
\item $\varphi_{126} :\equiv N_{2} \in \mathbb{N} \land \forall n \in \mathbb{N}\,(N_{2} \leq n \to \lvert b_{n} - c \rvert < \varepsilon \cdot 2^{-1}) \to n \in \mathbb{N} \to N_{1} \leq n \to s_{N_{1} \cup N_{2}} \in \mathbb{N}$
\item $\varphi_{127} :\equiv N_{2} \in \mathbb{N} \land \forall n \in \mathbb{N}\,(N_{2} \leq n \to \lvert b_{n} - c \rvert < \varepsilon \cdot 2^{-1}) \to n \in \mathbb{N} \to N_{1} \leq n \to \lvert a_{n} - c \rvert < \varepsilon$
\item $\varphi_{128} :\equiv N_{2} \in \mathbb{N} \land \forall n \in \mathbb{N}\,(N_{2} \leq n \to \lvert b_{n} - c \rvert < \varepsilon \cdot 2^{-1}) \to n \in \mathbb{N} \to N_{1} \leq n \to N_{1} \cup N_{2} \leq s_{N_{1} \cup N_{2}}$
\item $\varphi_{129} :\equiv N_{2} \in \mathbb{N} \land \forall n \in \mathbb{N}\,(N_{2} \leq n \to \lvert b_{n} - c \rvert < \varepsilon \cdot 2^{-1}) \to \forall n \in \mathbb{N}\,(N_{1} \leq n \to \lvert a_{n} - c \rvert < \varepsilon)$
\item $\varphi_{130} :\equiv \forall c\,(\varphi_{11} \to \neg \mathrm{Conv}(b,c)) \to \forall c\,(\varphi_{11}) \to \forall c\,(\neg \mathrm{Conv}(b,c))$
\item $\varphi_{131} :\equiv \forall x\,\forall y\,\forall z\,\forall e\,(x \in \mathbb{R} \to y \in \mathbb{R} \to z \in \mathbb{R} \to e \in \mathbb{R} \to \lvert x - y \rvert < e \to \lvert y - z \rvert < e \to \lvert x - z \rvert < e + e)$
\item $\varphi_{132} :\equiv (\mathrm{Cgt}(a) \to \mathrm{Conv}(a,\alpha) \to \varepsilon \in \mathbb{R} \to 0 < \varepsilon \to 0 < \varepsilon) \to \varphi_{42}$
\item $\varphi_{133} :\equiv N_{2} \in \mathbb{N} \land \forall n \in \mathbb{N}\,(N_{2} \leq n \to \lvert b_{n} - c \rvert < \varepsilon \cdot 2^{-1}) \to n \in \mathbb{N} \to N_{1} \leq n \to a_{s_{N_{1} \cup N_{2}}} \in \mathbb{R}$
\item $\varphi_{134} :\equiv \varphi_{36} \to \neg (N_{1} \in \mathbb{N} \land \forall n \in \mathbb{N}\,(N_{1} \leq n \to \lvert a_{n} - c \rvert < \varepsilon))$
\item $\varphi_{135} :\equiv \mathrm{Cauchy}(a) \to \mathrm{SubSeq}(b,a) \land \mathrm{Cgt}(b) \to \mathrm{Conv}(b,c) \to \varphi_{63}$
\item $\varphi_{136} :\equiv N_{2} \in \mathbb{N} \land \forall n \in \mathbb{N}\,(N_{2} \leq n \to \lvert b_{n} - c \rvert < \varepsilon \cdot 2^{-1}) \to n \in \mathbb{N} \to N_{1} \leq n \to \varepsilon \cdot 2^{-1} + \varepsilon \cdot 2^{-1} = \varepsilon$
\item $\varphi_{137} :\equiv \mathrm{Cauchy}(a) \to \mathrm{SubSeq}(b,a) \land \mathrm{Cgt}(b) \to \mathrm{Conv}(b,c) \to \varphi_{67}$
\item $\varphi_{138} :\equiv \forall w\,(\varepsilon \cdot 2^{-1} + \varepsilon \cdot 2^{-1} = w \to \lvert a_{m} - a_{n} \rvert < \varepsilon \cdot 2^{-1} + \varepsilon \cdot 2^{-1} \to \lvert a_{m} - a_{n} \rvert < w)$
\item $\varphi_{139} :\equiv \varepsilon \cdot 2^{-1} + \varepsilon \cdot 2^{-1} = \varepsilon \to \lvert a_{m} - a_{n} \rvert < \varepsilon \cdot 2^{-1} + \varepsilon \cdot 2^{-1} \to \lvert a_{m} - a_{n} \rvert < \varepsilon$
\item $\varphi_{140} :\equiv \forall x\,\forall y\,(x \in \mathbb{R} \to y \in \mathbb{R} \to \lvert x - y \rvert = \lvert y - x \rvert) \to \forall y\,(a_{n} \in \mathbb{R} \to y \in \mathbb{R} \to \lvert a_{n} - y \rvert = \lvert y - a_{n} \rvert)$
\item $\varphi_{141} :\equiv \forall x\,\forall y\,(x \in \mathbb{R} \to y \in \mathbb{R} \to \lvert x - y \rvert = \lvert y - x \rvert) \to \forall y\,(a_{m} \in \mathbb{R} \to y \in \mathbb{R} \to \lvert a_{m} - y \rvert = \lvert y - a_{m} \rvert)$
\item $\varphi_{142} :\equiv N_{2} \in \mathbb{N} \land \forall n \in \mathbb{N}\,(N_{2} \leq n \to \lvert b_{n} - c \rvert < \varepsilon \cdot 2^{-1}) \to n \in \mathbb{N} \to N_{1} \leq n \to \forall k \in \mathbb{N}\,(b_{k} = a_{s_{k}})$
\item $\varphi_{143} :\equiv \forall \varepsilon\,(\mathrm{Cauchy}(a) \to \varepsilon \in \mathbb{R} \to 0 < \varepsilon \to \exists N \in \mathbb{N}\,\forall m\,(\varphi_{56}))$
\item $\varphi_{144} :\equiv \forall \alpha\,(\mathrm{Conv}(a,\alpha) \leftrightarrow \mathrm{Seq}(a) \land \alpha \in \mathbb{R} \land \forall \varepsilon \in \mathbb{R}\,(0 < \varepsilon \to \exists N \in \mathbb{N}\,\forall n \in \mathbb{N}\,(N \leq n \to \lvert a_{n} - \alpha \rvert < \varepsilon)))$
\item $\varphi_{145} :\equiv N_{2} \in \mathbb{N} \land \forall n \in \mathbb{N}\,(N_{2} \leq n \to \lvert b_{n} - c \rvert < \varepsilon \cdot 2^{-1}) \to N_{1} \in \mathbb{N} \land \forall n \in \mathbb{N}\,(N_{1} \leq n \to \lvert a_{n} - c \rvert < \varepsilon)$
\item $\varphi_{146} :\equiv \mathrm{Conv}(a,\alpha) \to \forall \varepsilon \in \mathbb{R}\,(0 < \varepsilon \to \exists N \in \mathbb{N}\,\forall m\,(\varphi_{56}))$
\item $\varphi_{147} :\equiv N_{2} \in \mathbb{N} \land \forall n \in \mathbb{N}\,(N_{2} \leq n \to \lvert b_{n} - c \rvert < \varepsilon \cdot 2^{-1}) \to n \in \mathbb{N} \to N_{1} \leq n \to b_{N_{1} \cup N_{2}} = a_{s_{N_{1} \cup N_{2}}}$
\item $\varphi_{148} :\equiv N_{2} \in \mathbb{N} \land \forall n \in \mathbb{N}\,(N_{2} \leq n \to \lvert b_{n} - c \rvert < \varepsilon \cdot 2^{-1}) \to \exists N \in \mathbb{N}\,\forall n \in \mathbb{N}\,(N \leq n \to \lvert a_{n} - c \rvert < \varepsilon)$
\item $\varphi_{149} :\equiv \mathrm{Cauchy}(a) \to \mathrm{SubSeq}(b,a) \land \mathrm{Cgt}(b) \to \mathrm{Conv}(b,c) \to \varphi_{80}$
\item $\varphi_{150} :\equiv N_{2} \in \mathbb{N} \land \forall n \in \mathbb{N}\,(N_{2} \leq n \to \lvert b_{n} - c \rvert < \varepsilon \cdot 2^{-1}) \to n \in \mathbb{N} \to N_{1} \leq n \to \lvert b_{N_{1} \cup N_{2}} - c \rvert < \varepsilon \cdot 2^{-1}$
\item $\varphi_{151} :\equiv \mathrm{Cgt}(a) \to \mathrm{Conv}(a,\alpha) \to \varepsilon \in \mathbb{R} \to 0 < \varepsilon \to \exists N \in \mathbb{N}\,\forall m\,(\varphi_{56})$
\item $\varphi_{152} :\equiv \mathrm{Cgt}(a) \land \mathrm{Conv}(a,\alpha) \to \varepsilon \in \mathbb{R} \to 0 < \varepsilon \to \exists N \in \mathbb{N}\,\forall m\,(\varphi_{56})$
\item $\varphi_{153} :\equiv v_{750} \in \mathbb{R} \to v_{751} \in \mathbb{R} \to v_{752} \in \mathbb{R} \to \lvert v_{749} - v_{750} \rvert < v_{752} \to \varphi_{33}$
\item $\varphi_{154} :\equiv \mathrm{Cgt}(a) \land \mathrm{Conv}(a,\alpha) \land \varepsilon \in \mathbb{R} \land 0 < \varepsilon \land \neg \exists N \in \mathbb{N}\,\forall m\,(\varphi_{56})$
\item $\varphi_{155} :\equiv (\mathrm{Cgt}(a) \to \mathrm{Conv}(a,\alpha) \to \varepsilon \in \mathbb{R} \to \varepsilon \in \mathbb{R}) \to 2^{-1} \in \mathbb{R} \to \varphi_{52}$
\item $\varphi_{156} :\equiv \mathrm{Cauchy}(a) \leftrightarrow \mathrm{Seq}(a) \land \forall \varepsilon \in \mathbb{R}\,(0 < \varepsilon \to \exists N \in \mathbb{N}\,\forall m\,(\varphi_{56}))$
\item $\varphi_{157} :\equiv \forall y\,\forall z\,\forall e\,(a_{n} \in \mathbb{R} \to y \in \mathbb{R} \to z \in \mathbb{R} \to e \in \mathbb{R} \to \lvert a_{n} - y \rvert < e \to \lvert y - z \rvert < e \to \lvert a_{n} - z \rvert < e + e)$
\item $\varphi_{158} :\equiv \forall x\,\forall \alpha\,\forall \beta\,\forall \varepsilon\,(x \in \mathbb{R} \to \alpha \in \mathbb{R} \to \beta \in \mathbb{R} \to \varepsilon \in \mathbb{R} \to \lvert x - \alpha \rvert < \varepsilon \to \lvert x - \beta \rvert < \varepsilon \to \lvert \alpha - \beta \rvert < \varepsilon + \varepsilon)$
\item $\varphi_{159} :\equiv N_{2} \in \mathbb{N} \land \forall n \in \mathbb{N}\,(N_{2} \leq n \to \lvert b_{n} - c \rvert < \varepsilon \cdot 2^{-1}) \to N_{2} \in \mathbb{N} \land \forall n \in \mathbb{N}\,(N_{2} \leq n \to \lvert b_{n} - c \rvert < \varepsilon \cdot 2^{-1})$
\item $\varphi_{160} :\equiv N_{2} \in \mathbb{N} \land \forall n \in \mathbb{N}\,(N_{2} \leq n \to \lvert b_{n} - c \rvert < \varepsilon \cdot 2^{-1}) \to n \in \mathbb{N} \to N_{1} \leq n \to \lvert a_{s_{N_{1} \cup N_{2}}} - c \rvert < \varepsilon \cdot 2^{-1}$
\item $\varphi_{161} :\equiv N_{2} \in \mathbb{N} \land \forall n \in \mathbb{N}\,(N_{2} \leq n \to \lvert b_{n} - c \rvert < \varepsilon \cdot 2^{-1}) \to n \in \mathbb{N} \to N_{1} \leq n \to \lvert a_{n} - c \rvert < \varepsilon \cdot 2^{-1} + \varepsilon \cdot 2^{-1}$
\item $\varphi_{162} :\equiv \mathrm{Cauchy}(a) \to \mathrm{SubSeq}(b,a) \land \mathrm{Cgt}(b) \to \mathrm{Conv}(b,c) \to \varphi_{97}$
\item $\varphi_{163} :\equiv \forall \varepsilon\,(\mathrm{Cgt}(a) \land \mathrm{Conv}(a,\alpha)) \to \forall \varepsilon \in \mathbb{R}\,(0 < \varepsilon \to \exists N \in \mathbb{N}\,\forall m\,(\varphi_{56}))$
\item $\varphi_{164} :\equiv \forall \varepsilon\,(\mathrm{Conv}(a,\alpha) \to \mathrm{Conv}(a,\alpha) \to \varepsilon \in \mathbb{R} \to 0 < \varepsilon \to \exists N \in \mathbb{N}\,\forall n \in \mathbb{N}\,(N \leq n \to \lvert a_{n} - \alpha \rvert < \varepsilon \land \lvert a_{n} - \alpha \rvert < \varepsilon))$
\item $\varphi_{165} :\equiv \mathrm{Cauchy}(a) \to \mathrm{SubSeq}(b,a) \land \mathrm{Cgt}(b) \to \mathrm{Conv}(b,c) \to \varphi_{102}$
\item $\varphi_{166} :\equiv \forall y\,(a_{n} \in \mathbb{R} \to y \in \mathbb{R} \to \lvert a_{n} - y \rvert = \lvert y - a_{n} \rvert) \to a_{n} \in \mathbb{R} \to \alpha \in \mathbb{R} \to \lvert a_{n} - \alpha \rvert = \lvert \alpha - a_{n} \rvert$
\item $\varphi_{167} :\equiv \forall y\,(a_{m} \in \mathbb{R} \to y \in \mathbb{R} \to \lvert a_{m} - y \rvert = \lvert y - a_{m} \rvert) \to a_{m} \in \mathbb{R} \to \alpha \in \mathbb{R} \to \lvert a_{m} - \alpha \rvert = \lvert \alpha - a_{m} \rvert$
\item $\varphi_{168} :\equiv \forall \alpha\,\forall \beta\,\forall \varepsilon\,(v_{749} \in \mathbb{R} \to \alpha \in \mathbb{R} \to \beta \in \mathbb{R} \to \varepsilon \in \mathbb{R} \to \lvert v_{749} - \alpha \rvert < \varepsilon \to \lvert v_{749} - \beta \rvert < \varepsilon \to \lvert \alpha - \beta \rvert < \varepsilon + \varepsilon)$
\item $\varphi_{169} :\equiv N_{2} \in \mathbb{N} \land \forall n \in \mathbb{N}\,(N_{2} \leq n \to \lvert b_{n} - c \rvert < \varepsilon \cdot 2^{-1}) \to \varphi_{60}$
\item $\varphi_{170} :\equiv \mathrm{Cauchy}(a) \to \mathrm{SubSeq}(b,a) \land \mathrm{Cgt}(b) \to \mathrm{Conv}(b,c) \to \varphi_{103}$
\item $\varphi_{171} :\equiv N_{2} \in \mathbb{N} \land \forall n \in \mathbb{N}\,(N_{2} \leq n \to \lvert b_{n} - c \rvert < \varepsilon \cdot 2^{-1}) \to n \in \mathbb{N} \to N_{1} \leq n \to \lvert a_{n} - a_{s_{N_{1} \cup N_{2}}} \rvert < \varepsilon \cdot 2^{-1}$
\item $\varphi_{172} :\equiv \forall \beta\,\forall \varepsilon\,(v_{749} \in \mathbb{R} \to v_{750} \in \mathbb{R} \to \beta \in \mathbb{R} \to \varepsilon \in \mathbb{R} \to \lvert v_{749} - v_{750} \rvert < \varepsilon \to \lvert v_{749} - \beta \rvert < \varepsilon \to \lvert v_{750} - \beta \rvert < \varepsilon + \varepsilon)$
\item $\varphi_{173} :\equiv \mathrm{Cauchy}(a) \to \mathrm{SubSeq}(b,a) \land \mathrm{Cgt}(b) \to \mathrm{Conv}(b,c) \to \varphi_{107}$
\item $\varphi_{174} :\equiv \forall \beta\,\forall \varepsilon\,(\mathrm{Conv}(a,\alpha) \to \mathrm{Conv}(a,\beta) \to \varepsilon \in \mathbb{R} \to 0 < \varepsilon \to \exists N \in \mathbb{N}\,\forall n \in \mathbb{N}\,(N \leq n \to \lvert a_{n} - \alpha \rvert < \varepsilon \land \lvert a_{n} - \beta \rvert < \varepsilon))$
\item $\varphi_{175} :\equiv \forall \varepsilon\,(v_{749} \in \mathbb{R} \to v_{750} \in \mathbb{R} \to v_{751} \in \mathbb{R} \to \varepsilon \in \mathbb{R} \to \lvert v_{749} - v_{750} \rvert < \varepsilon \to \lvert v_{749} - v_{751} \rvert < \varepsilon \to \lvert v_{750} - v_{751} \rvert < \varepsilon + \varepsilon)$
\item $\varphi_{176} :\equiv \mathrm{Conv}(b,c) \to \mathrm{SubIdx}(s) \land \forall k \in \mathbb{N}\,(b_{k} = a_{s_{k}}) \to \varphi_{76}$
\item $\varphi_{177} :\equiv \forall b\,\forall \beta\,\forall \varepsilon\,(\mathrm{Conv}(a,\alpha) \to \mathrm{Conv}(b,\beta) \to \varepsilon \in \mathbb{R} \to 0 < \varepsilon \to \exists N \in \mathbb{N}\,\forall n \in \mathbb{N}\,(N \leq n \to \lvert a_{n} - \alpha \rvert < \varepsilon \land \lvert b_{n} - \beta \rvert < \varepsilon))$
\item $\varphi_{178} :\equiv N_{2} \in \mathbb{N} \land \forall n \in \mathbb{N}\,(N_{2} \leq n \to \lvert b_{n} - c \rvert < \varepsilon \cdot 2^{-1}) \to n \in \mathbb{N} \to \varphi_{53}$
\item $\varphi_{179} :\equiv \mathrm{Cgt}(a) \land \mathrm{Conv}(a,\alpha) \land \varepsilon \in \mathbb{R} \land 0 < \varepsilon \land (N \in \mathbb{N} \land \varphi_{92})$
\item $\varphi_{180} :\equiv \mathrm{Conv}(b,c) \to \mathrm{SubIdx}(s) \land \forall k \in \mathbb{N}\,(b_{k} = a_{s_{k}}) \to \varphi_{82}$
\item $\varphi_{181} :\equiv \mathrm{Cgt}(a) \to \mathrm{Conv}(a,\alpha) \to \varepsilon \in \mathbb{R} \to 0 < \varepsilon \to \exists N \in \mathbb{N}\,(\varphi_{92})$
\item $\varphi_{182} :\equiv N_{2} \in \mathbb{N} \land \forall n \in \mathbb{N}\,(N_{2} \leq n \to \lvert b_{n} - c \rvert < \varepsilon \cdot 2^{-1}) \to \varphi_{75}$
\item $\varphi_{183} :\equiv \forall a\,(\mathrm{SubSeq}(b,a) \leftrightarrow \varphi_{40}) \to (\mathrm{SubSeq}(b,a) \leftrightarrow \varphi_{40})$
\item $\varphi_{184} :\equiv 0 < \varepsilon \to N \in \mathbb{N} \land \varphi_{92} \to m \in \mathbb{N} \to n \in \mathbb{N} \to N \leq m \to N \leq n \to N \leq n$
\item $\varphi_{185} :\equiv N_{1} \in \mathbb{N} \land \varphi_{75} \to \exists N \in \mathbb{N}\,\forall n \in \mathbb{N}\,(N \leq n \to \lvert a_{n} - c \rvert < \varepsilon)$
\item $\varphi_{186} :\equiv \forall v_{750}\,\forall v_{751}\,\forall v_{752}\,(\alpha \in \mathbb{R} \to v_{750} \in \mathbb{R} \to v_{751} \in \mathbb{R} \to v_{752} \in \mathbb{R} \to \lvert \alpha - v_{750} \rvert < v_{752} \to \lvert \alpha - v_{751} \rvert < v_{752} \to \lvert v_{750} - v_{751} \rvert < v_{752} + v_{752})$
\item $\varphi_{187} :\equiv N_{2} \in \mathbb{N} \land \forall n \in \mathbb{N}\,(N_{2} \leq n \to \lvert b_{n} - c \rvert < \varepsilon \cdot 2^{-1}) \to \varphi_{83}$
\item $\varphi_{188} :\equiv \forall b\,\forall a\,(\mathrm{SubSeq}(b,a) \leftrightarrow \varphi_{40}) \to \forall a\,(\mathrm{SubSeq}(b,a) \leftrightarrow \varphi_{40})$
\item $\varphi_{189} :\equiv \mathrm{Cgt}(a) \to \mathrm{Conv}(a,\alpha) \to \varepsilon \in \mathbb{R} \to 0 < \varepsilon \to N \in \mathbb{N} \land \varphi_{92} \to N \in \mathbb{N}$
\item $\varphi_{190} :\equiv N_{2} \in \mathbb{N} \land \forall n \in \mathbb{N}\,(N_{2} \leq n \to \lvert b_{n} - c \rvert < \varepsilon \cdot 2^{-1}) \to \varphi_{89}$
\item $\varphi_{191} :\equiv N \in \mathbb{N} \land \varphi_{92} \to m \in \mathbb{N} \to n \in \mathbb{N} \to N \leq m \to N \leq n \to \varepsilon \cdot 2^{-1} \in \mathbb{R}$
\item $\varphi_{192} :\equiv (\mathrm{Cauchy}(a) \to \varphi_{15}) \to \mathrm{Cauchy}(a) \to \mathrm{SubSeq}(b,a) \land \mathrm{Cgt}(b) \to \varphi_{61}$
\item $\varphi_{193} :\equiv \varepsilon \in \mathbb{R} \to 0 < \varepsilon \to N \in \mathbb{N} \land \varphi_{92} \to m \in \mathbb{N} \to n \in \mathbb{N} \to N \leq m \to N \leq m$
\item $\varphi_{194} :\equiv \mathrm{Conv}(a,\alpha) \to \varepsilon \cdot 2^{-1} \in \mathbb{R} \to 0 < \varepsilon \cdot 2^{-1} \to \exists N \in \mathbb{N}\,(\varphi_{92})$
\item $\varphi_{195} :\equiv \forall v_{751}\,\forall v_{752}\,(\alpha \in \mathbb{R} \to a_{m} \in \mathbb{R} \to v_{751} \in \mathbb{R} \to v_{752} \in \mathbb{R} \to \lvert \alpha - a_{m} \rvert < v_{752} \to \lvert \alpha - v_{751} \rvert < v_{752} \to \lvert a_{m} - v_{751} \rvert < v_{752} + v_{752})$
\item $\varphi_{196} :\equiv N_{2} \in \mathbb{N} \land \forall n \in \mathbb{N}\,(N_{2} \leq n \to \lvert b_{n} - c \rvert < \varepsilon \cdot 2^{-1}) \to v_{753} \in \mathbb{N} \to v_{754} \in \mathbb{N} \to \varphi_{48}$
\item $\varphi_{197} :\equiv \mathrm{Conv}(a,\alpha) \to \varepsilon \in \mathbb{R} \to 0 < \varepsilon \to N \in \mathbb{N} \land \varphi_{92} \to m \in \mathbb{N} \to m \in \mathbb{N}$
\item $\varphi_{198} :\equiv N_{2} \in \mathbb{N} \land \forall n \in \mathbb{N}\,(N_{2} \leq n \to \lvert b_{n} - c \rvert < \varepsilon \cdot 2^{-1}) \to n \in \mathbb{N} \to N_{1} \leq n \to \varphi_{65}$
\item $\varphi_{199} :\equiv N_{1} \in \mathbb{N} \land \varphi_{75} \to \exists N \in \mathbb{N}\,\forall n \in \mathbb{N}\,(N \leq n \to \lvert b_{n} - c \rvert < \varepsilon \cdot 2^{-1})$
\item $\varphi_{200} :\equiv 0 < \varepsilon \to N \in \mathbb{N} \land \varphi_{92} \to m \in \mathbb{N} \to n \in \mathbb{N} \to N \leq m \to N \leq n \to a_{n} \in \mathbb{R}$
\item $\varphi_{201} :\equiv 0 < \varepsilon \to N \in \mathbb{N} \land \varphi_{92} \to m \in \mathbb{N} \to n \in \mathbb{N} \to N \leq m \to N \leq n \to a_{m} \in \mathbb{R}$
\item $\varphi_{202} :\equiv (\mathrm{Cgt}(a) \to \varphi_{146}) \to \mathrm{Cgt}(a) \to \mathrm{Conv}(a,\alpha) \to \mathrm{Cauchy}(a)$
\item $\varphi_{203} :\equiv \forall v_{752}\,(\alpha \in \mathbb{R} \to a_{m} \in \mathbb{R} \to a_{n} \in \mathbb{R} \to v_{752} \in \mathbb{R} \to \lvert \alpha - a_{m} \rvert < v_{752} \to \lvert \alpha - a_{n} \rvert < v_{752} \to \lvert a_{m} - a_{n} \rvert < v_{752} + v_{752})$
\item $\varphi_{204} :\equiv N_{2} \in \mathbb{N} \land \forall n \in \mathbb{N}\,(N_{2} \leq n \to \lvert b_{n} - c \rvert < \varepsilon \cdot 2^{-1}) \to \forall v_{754}\,(v_{753} \in \mathbb{N} \to v_{754} \in \mathbb{N} \to \varphi_{48})$
\item $\varphi_{205} :\equiv (\mathrm{SubSeq}(b,a) \leftrightarrow \varphi_{40}) \to \mathrm{Cauchy}(a) \to \mathrm{SubSeq}(b,a) \land \mathrm{Cgt}(b) \to \varphi_{40}$
\item $\varphi_{206} :\equiv \mathrm{SubIdx}(s) \land \forall k \in \mathbb{N}\,(b_{k} = a_{s_{k}}) \to \varepsilon \in \mathbb{R} \to 0 < \varepsilon \to \varphi_{99}$
\item $\varphi_{207} :\equiv N \in \mathbb{N} \land \varphi_{92} \to m \in \mathbb{N} \to n \in \mathbb{N} \to N \leq m \to N \leq n \to \lvert a_{m} - a_{n} \rvert < \varepsilon$
\item $\varphi_{208} :\equiv \forall N\,(\neg (N \in \mathbb{N} \land \forall m\,(\varphi_{56}))) \to \neg (N \in \mathbb{N} \land \forall m\,(\varphi_{56}))$
\item $\varphi_{209} :\equiv \forall \alpha\,(\forall c\,(\neg \mathrm{Conv}(b,c)) \to \neg \mathrm{Conv}(b,\alpha)) \to \varphi_{50} \to \varphi_{1} \to \varphi_{8}$
\item $\varphi_{210} :\equiv N_{2} \in \mathbb{N} \land \forall n \in \mathbb{N}\,(N_{2} \leq n \to \lvert b_{n} - c \rvert < \varepsilon \cdot 2^{-1}) \to \forall v_{753}\,\forall v_{754}\,(v_{753} \in \mathbb{N} \to v_{754} \in \mathbb{N} \to \varphi_{48})$
\item $\varphi_{211} :\equiv \mathrm{Conv}(a,\alpha) \to \varepsilon \in \mathbb{R} \to 0 < \varepsilon \to N \in \mathbb{N} \land \varphi_{92} \to m \in \mathbb{N} \to n \in \mathbb{N} \to n \in \mathbb{N}$
\item $\varphi_{212} :\equiv N \in \mathbb{N} \land \varphi_{92} \to m \in \mathbb{N} \to n \in \mathbb{N} \to N \leq m \to N \leq n \to \lvert a_{n} - \alpha \rvert < \varepsilon \cdot 2^{-1}$
\item $\varphi_{213} :\equiv N \in \mathbb{N} \land \varphi_{92} \to m \in \mathbb{N} \to n \in \mathbb{N} \to N \leq m \to N \leq n \to \lvert \alpha - a_{n} \rvert < \varepsilon \cdot 2^{-1}$
\item $\varphi_{214} :\equiv N \in \mathbb{N} \land \varphi_{92} \to m \in \mathbb{N} \to n \in \mathbb{N} \to N \leq m \to N \leq n \to \lvert a_{m} - \alpha \rvert < \varepsilon \cdot 2^{-1}$
\item $\varphi_{215} :\equiv N \in \mathbb{N} \land \varphi_{92} \to m \in \mathbb{N} \to n \in \mathbb{N} \to N \leq m \to N \leq n \to \lvert \alpha - a_{m} \rvert < \varepsilon \cdot 2^{-1}$
\item $\varphi_{216} :\equiv (\mathrm{Cauchy}(a) \to \mathrm{Cauchy}(a)) \to \varphi_{156} \to \mathrm{Cauchy}(a) \to \mathrm{Seq}(a)$
\item $\varphi_{217} :\equiv \varphi_{95} \land \varepsilon \in \mathbb{R} \land 0 < \varepsilon \land \neg \exists N \in \mathbb{N}\,\forall n \in \mathbb{N}\,(N \leq n \to \lvert a_{n} - c \rvert < \varepsilon)$
\item $\varphi_{218} :\equiv \varepsilon \cdot 2^{-1} \in \mathbb{R} \to \lvert a_{n} - a_{s_{N_{1} \cup N_{2}}} \rvert < \varepsilon \cdot 2^{-1} \to \varphi_{109}$
\item $\varphi_{219} :\equiv N_{2} \in \mathbb{N} \land \forall n \in \mathbb{N}\,(N_{2} \leq n \to \lvert b_{n} - c \rvert < \varepsilon \cdot 2^{-1}) \to n \in \mathbb{N} \to s_{N_{1} \cup N_{2}} \in \mathbb{N} \to N_{1} \leq n \to \varphi_{66}$
\item $\varphi_{220} :\equiv (\mathrm{Cauchy}(a) \to \mathrm{SubSeq}(b,a) \land \mathrm{Cgt}(b) \to \varphi_{40}) \to \mathrm{Cauchy}(a) \to \mathrm{SubSeq}(b,a) \land \mathrm{Cgt}(b) \to \varphi_{43}$
\item $\varphi_{221} :\equiv a_{m} \in \mathbb{R} \to a_{n} \in \mathbb{R} \to \varepsilon \cdot 2^{-1} \in \mathbb{R} \to \lvert \alpha - a_{m} \rvert < \varepsilon \cdot 2^{-1} \to \varphi_{98}$
\item $\varphi_{222} :\equiv \forall e\,(a_{n} \in \mathbb{R} \to a_{s_{N_{1} \cup N_{2}}} \in \mathbb{R} \to c \in \mathbb{R} \to e \in \mathbb{R} \to \lvert a_{n} - a_{s_{N_{1} \cup N_{2}}} \rvert < e \to \lvert a_{s_{N_{1} \cup N_{2}}} - c \rvert < e \to \lvert a_{n} - c \rvert < e + e)$
\item $\varphi_{223} :\equiv \forall z\,\forall e\,(a_{n} \in \mathbb{R} \to a_{s_{N_{1} \cup N_{2}}} \in \mathbb{R} \to z \in \mathbb{R} \to e \in \mathbb{R} \to \lvert a_{n} - a_{s_{N_{1} \cup N_{2}}} \rvert < e \to \lvert a_{s_{N_{1} \cup N_{2}}} - z \rvert < e \to \lvert a_{n} - z \rvert < e + e)$
\item $\varphi_{224} :\equiv \forall z\,\forall w\,(z = w \to \lvert a_{m} - a_{n} \rvert < z \to \lvert a_{m} - a_{n} \rvert < w) \to \varphi_{138}$
\item $\varphi_{225} :\equiv \forall N_{1}\,(\neg (N_{1} \in \mathbb{N} \land \varphi_{75})) \to \forall N\,\forall N_{1}\,(\neg (N_{1} \in \mathbb{N} \land \varphi_{75}))$
\item $\varphi_{226} :\equiv \forall b\,(\varphi_{12}) \to (\forall b\,(\varphi_{12}) \to \varphi_{21}) \to \varphi_{14} \to \mathrm{Cauchy}(a) \to \mathrm{Cgt}(a)$
\item $\varphi_{227} :\equiv (\mathrm{Cgt}(a) \to \exists \alpha\,(\mathrm{Conv}(a,\alpha))) \to \varphi_{5} \to (\varphi_{5} \to \varphi_{10}) \to \varphi_{16} \to \mathrm{Cgt}(a) \to \mathrm{Cauchy}(a)$
\item $\varphi_{228} :\equiv \varphi_{110} \to \varphi_{27} \to \neg (\mathrm{SubIdx}(s) \land \forall k \in \mathbb{N}\,(b_{k} = a_{s_{k}}))$
\item $\varphi_{229} :\equiv (\varepsilon \in \mathbb{R} \to 0 < \varepsilon \to 0 < \varepsilon \cdot 2^{-1}) \to (\mathrm{Cgt}(a) \to \mathrm{Conv}(a,\alpha) \to \varepsilon \in \mathbb{R} \to \varepsilon \in \mathbb{R}) \to \varphi_{132}$
\item $\varphi_{230} :\equiv \mathrm{Conv}(a,\alpha) \to \varepsilon \in \mathbb{R} \to 0 < \varepsilon \to N \in \mathbb{N} \land \varphi_{92} \to \forall m\,(\varphi_{56})$
\item $\varphi_{231} :\equiv \mathrm{Cauchy}(a) \to \mathrm{SubSeq}(b,a) \land \mathrm{Cgt}(b) \to \mathrm{Conv}(b,c) \to \varphi_{206}$
\item $\varphi_{232} :\equiv \mathrm{Cgt}(a) \to \mathrm{Conv}(a,\alpha) \to \varepsilon \in \mathbb{R} \to 0 < \varepsilon \to N \in \mathbb{N} \land \varphi_{92} \to \varphi_{54}$
\item $\varphi_{233} :\equiv \forall n\,(\varphi_{179} \to m \in \mathbb{N} \to n \in \mathbb{N} \to N \leq m \to N \leq n \to \lvert a_{m} - a_{n} \rvert < \varepsilon)$
\item $\varphi_{234} :\equiv \mathrm{Cgt}(a) \to \mathrm{Conv}(a,\alpha) \to \varepsilon \in \mathbb{R} \to 0 < \varepsilon \to N \in \mathbb{N} \land \varphi_{92} \to \varphi_{56}$
\item $\varphi_{235} :\equiv \mathrm{Cgt}(a) \to \mathrm{Conv}(a,\alpha) \to \varepsilon \in \mathbb{R} \to 0 < \varepsilon \to N \in \mathbb{N} \land \varphi_{92} \to \varphi_{58}$
\item $\varphi_{236} :\equiv \mathrm{Cgt}(a) \to \mathrm{Conv}(a,\alpha) \to \varepsilon \in \mathbb{R} \to 0 < \varepsilon \to N \in \mathbb{N} \land \varphi_{92} \to \varphi_{59}$
\item $\varphi_{237} :\equiv 2^{-1} \in \mathbb{R} \to \mathrm{Cgt}(a) \to \mathrm{Conv}(a,\alpha) \to \varepsilon \in \mathbb{R} \to 0 < \varepsilon \to \varphi_{191}$
\item $\varphi_{238} :\equiv \mathrm{Conv}(a,\alpha) \to \varepsilon \in \mathbb{R} \to 0 < \varepsilon \to N \in \mathbb{N} \land \varphi_{92} \to N \in \mathbb{N} \land \forall m\,(\varphi_{56})$
\item $\varphi_{239} :\equiv \mathrm{Conv}(a,\alpha) \to \varepsilon \in \mathbb{R} \to 0 < \varepsilon \to N \in \mathbb{N} \land \varphi_{92} \to \exists N \in \mathbb{N}\,\forall m\,(\varphi_{56})$
\item $\varphi_{240} :\equiv \mathrm{SubIdx}(s) \land \forall k \in \mathbb{N}\,(b_{k} = a_{s_{k}}) \to \varepsilon \in \mathbb{R} \to 0 < \varepsilon \to \varphi_{185}$
\item $\varphi_{241} :\equiv \mathrm{Conv}(a,\alpha) \to \varepsilon \in \mathbb{R} \to 0 < \varepsilon \to N \in \mathbb{N} \land \varphi_{92} \to m \in \mathbb{N} \to \varphi_{72}$
\item $\varphi_{242} :\equiv \mathrm{Cgt}(a) \to \mathrm{Conv}(a,\alpha) \to \varepsilon \in \mathbb{R} \to 0 < \varepsilon \to N \in \mathbb{N} \land \varphi_{92} \to \varphi_{85}$
\item $\varphi_{243} :\equiv \mathrm{Cgt}(a) \to \mathrm{Conv}(a,\alpha) \to \varepsilon \in \mathbb{R} \to 0 < \varepsilon \to N \in \mathbb{N} \land \varphi_{92} \to \varphi_{86}$
\item $\varphi_{244} :\equiv \mathrm{SubIdx}(s) \land \forall k \in \mathbb{N}\,(b_{k} = a_{s_{k}}) \to \varepsilon \in \mathbb{R} \to 0 < \varepsilon \to \varphi_{199}$
\item $\varphi_{245} :\equiv \mathrm{Cgt}(a) \to \mathrm{Conv}(a,\alpha) \to \varepsilon \in \mathbb{R} \to 0 < \varepsilon \to N \in \mathbb{N} \land \varphi_{92} \to \varphi_{92}$
\item $\varphi_{246} :\equiv \mathrm{Conv}(a,\alpha) \to \varepsilon \in \mathbb{R} \to 0 < \varepsilon \to N \in \mathbb{N} \land \varphi_{92} \to N \in \mathbb{N} \land \varphi_{92}$
\item $\varphi_{247} :\equiv \varphi_{115} \to (\forall \alpha\,(\neg \mathrm{Conv}(a,\alpha)) \to \neg \mathrm{Conv}(a,c)) \to \varphi_{121}$
\item $\varphi_{248} :\equiv \mathrm{SubIdx}(s) \land \forall k \in \mathbb{N}\,(b_{k} = a_{s_{k}}) \to \varepsilon \in \mathbb{R} \to 0 < \varepsilon \to N_{1} \in \mathbb{N} \land \varphi_{75} \to \varphi_{70}$
\item $\varphi_{249} :\equiv \mathrm{SubIdx}(s) \land \forall k \in \mathbb{N}\,(b_{k} = a_{s_{k}}) \to \varepsilon \in \mathbb{R} \to 0 < \varepsilon \to N_{1} \in \mathbb{N} \land \varphi_{75} \to \varphi_{71}$
\item $\varphi_{250} :\equiv \mathrm{SubIdx}(s) \land \forall k \in \mathbb{N}\,(b_{k} = a_{s_{k}}) \to \varepsilon \in \mathbb{R} \to 0 < \varepsilon \to N_{1} \in \mathbb{N} \land \varphi_{75} \to \varphi_{88}$
\item $\varphi_{251} :\equiv \mathrm{SubIdx}(s) \land \forall k \in \mathbb{N}\,(b_{k} = a_{s_{k}}) \to \varepsilon \in \mathbb{R} \to 0 < \varepsilon \to N_{1} \in \mathbb{N} \land \varphi_{75} \to \varphi_{91}$
\item $\varphi_{252} :\equiv \mathrm{SubIdx}(s) \land \forall k \in \mathbb{N}\,(b_{k} = a_{s_{k}}) \to \varepsilon \in \mathbb{R} \to 0 < \varepsilon \to N_{1} \in \mathbb{N} \land \varphi_{75} \to N_{1} \in \mathbb{N} \land \varphi_{75}$
\item $\varphi_{253} :\equiv \mathrm{SubIdx}(s) \land \forall k \in \mathbb{N}\,(b_{k} = a_{s_{k}}) \to \varepsilon \in \mathbb{R} \to 0 < \varepsilon \to N_{1} \in \mathbb{N} \land \varphi_{75} \to \varphi_{105}$
\item $\varphi_{254} :\equiv \mathrm{Cauchy}(a) \to \mathrm{SubSeq}(b,a) \land \mathrm{Cgt}(b) \to \mathrm{Conv}(b,c) \to \varphi_{240}$
\item $\varphi_{255} :\equiv \varphi_{95} \land \varepsilon \in \mathbb{R} \land 0 < \varepsilon \land (N_{1} \in \mathbb{N} \land \varphi_{75}) \land \neg \exists N \in \mathbb{N}\,\forall n \in \mathbb{N}\,(N \leq n \to \lvert a_{n} - c \rvert < \varepsilon)$
\item $\varphi_{256} :\equiv \mathrm{SubIdx}(s) \land \forall k \in \mathbb{N}\,(b_{k} = a_{s_{k}}) \to \varepsilon \in \mathbb{R} \to 0 < \varepsilon \to N_{1} \in \mathbb{N} \land \varphi_{75} \to \varphi_{112}$
\item $\varphi_{257} :\equiv \mathrm{SubIdx}(s) \land \forall k \in \mathbb{N}\,(b_{k} = a_{s_{k}}) \to \varepsilon \in \mathbb{R} \to 0 < \varepsilon \to N_{1} \in \mathbb{N} \land \varphi_{75} \to \varphi_{116}$
\item $\varphi_{258} :\equiv \mathrm{SubIdx}(s) \land \forall k \in \mathbb{N}\,(b_{k} = a_{s_{k}}) \to \varepsilon \in \mathbb{R} \to 0 < \varepsilon \to N_{1} \in \mathbb{N} \land \varphi_{75} \to \varphi_{118}$
\item $\varphi_{259} :\equiv \varphi_{95} \land \varepsilon \in \mathbb{R} \land 0 < \varepsilon \land (N_{1} \in \mathbb{N} \land \varphi_{75}) \land (N_{2} \in \mathbb{N} \land \forall n \in \mathbb{N}\,(N_{2} \leq n \to \lvert b_{n} - c \rvert < \varepsilon \cdot 2^{-1}))$
\item $\varphi_{260} :\equiv \mathrm{SubIdx}(s) \land \forall k \in \mathbb{N}\,(b_{k} = a_{s_{k}}) \to \varepsilon \in \mathbb{R} \to 0 < \varepsilon \to N_{1} \in \mathbb{N} \land \varphi_{75} \to \varphi_{126}$
\item $\varphi_{261} :\equiv \mathrm{SubIdx}(s) \land \forall k \in \mathbb{N}\,(b_{k} = a_{s_{k}}) \to \varepsilon \in \mathbb{R} \to 0 < \varepsilon \to N_{1} \in \mathbb{N} \land \varphi_{75} \to \varphi_{127}$
\item $\varphi_{262} :\equiv \mathrm{Cauchy}(a) \to \mathrm{SubSeq}(b,a) \land \mathrm{Cgt}(b) \to \mathrm{Conv}(b,c) \to \varphi_{244}$
\item $\varphi_{263} :\equiv \mathrm{SubIdx}(s) \land \forall k \in \mathbb{N}\,(b_{k} = a_{s_{k}}) \to \varepsilon \in \mathbb{R} \to 0 < \varepsilon \to N_{1} \in \mathbb{N} \land \varphi_{75} \to \varphi_{128}$
\item $\varphi_{264} :\equiv \mathrm{SubIdx}(s) \land \forall k \in \mathbb{N}\,(b_{k} = a_{s_{k}}) \to \varepsilon \in \mathbb{R} \to 0 < \varepsilon \to N_{1} \in \mathbb{N} \land \varphi_{75} \to \varphi_{129}$
\item $\varphi_{265} :\equiv (\mathrm{Cauchy}(a) \to \mathrm{SubSeq}(b,a) \land \mathrm{Cgt}(b) \to \varphi_{180}) \to \varphi_{115}$
\item $\varphi_{266} :\equiv \mathrm{SubIdx}(s) \land \forall k \in \mathbb{N}\,(b_{k} = a_{s_{k}}) \to \varepsilon \in \mathbb{R} \to 0 < \varepsilon \to N_{1} \in \mathbb{N} \land \varphi_{75} \to \varphi_{133}$
\item $\varphi_{267} :\equiv \mathrm{SubIdx}(s) \land \forall k \in \mathbb{N}\,(b_{k} = a_{s_{k}}) \to \varepsilon \in \mathbb{R} \to 0 < \varepsilon \to N_{1} \in \mathbb{N} \land \varphi_{75} \to \varphi_{136}$
\item $\varphi_{268} :\equiv \mathrm{Cauchy}(a) \to \mathrm{SubSeq}(b,a) \land \mathrm{Cgt}(b) \to \mathrm{Conv}(b,c) \to \varphi_{248}$
\item $\varphi_{269} :\equiv \mathrm{Cauchy}(a) \to \mathrm{SubSeq}(b,a) \land \mathrm{Cgt}(b) \to \mathrm{Conv}(b,c) \to \varphi_{249}$
\item $\varphi_{270} :\equiv \mathrm{SubIdx}(s) \land \forall k \in \mathbb{N}\,(b_{k} = a_{s_{k}}) \to \varepsilon \in \mathbb{R} \to 0 < \varepsilon \to N_{1} \in \mathbb{N} \land \varphi_{75} \to \varphi_{142}$
\item $\varphi_{271} :\equiv \mathrm{SubIdx}(s) \land \forall k \in \mathbb{N}\,(b_{k} = a_{s_{k}}) \to \varepsilon \in \mathbb{R} \to 0 < \varepsilon \to N_{1} \in \mathbb{N} \land \varphi_{75} \to \varphi_{145}$
\item $\varphi_{272} :\equiv \mathrm{SubIdx}(s) \land \forall k \in \mathbb{N}\,(b_{k} = a_{s_{k}}) \to \varepsilon \in \mathbb{R} \to 0 < \varepsilon \to N_{1} \in \mathbb{N} \land \varphi_{75} \to \varphi_{147}$
\item $\varphi_{273} :\equiv \mathrm{SubIdx}(s) \land \forall k \in \mathbb{N}\,(b_{k} = a_{s_{k}}) \to \varepsilon \in \mathbb{R} \to 0 < \varepsilon \to N_{1} \in \mathbb{N} \land \varphi_{75} \to \varphi_{148}$
\item $\varphi_{274} :\equiv \mathrm{SubIdx}(s) \land \forall k \in \mathbb{N}\,(b_{k} = a_{s_{k}}) \to \varepsilon \in \mathbb{R} \to 0 < \varepsilon \to N_{1} \in \mathbb{N} \land \varphi_{75} \to \varphi_{150}$
\item $\varphi_{275} :\equiv \mathrm{Cauchy}(a) \to \mathrm{SubSeq}(b,a) \land \mathrm{Cgt}(b) \to \mathrm{Conv}(b,c) \to \varphi_{250}$
\item $\varphi_{276} :\equiv \mathrm{Cauchy}(a) \to \mathrm{SubSeq}(b,a) \land \mathrm{Cgt}(b) \to \mathrm{Conv}(b,c) \to \varphi_{251}$
\item $\varphi_{277} :\equiv \mathrm{SubIdx}(s) \land \forall k \in \mathbb{N}\,(b_{k} = a_{s_{k}}) \to \varepsilon \in \mathbb{R} \to 0 < \varepsilon \to N_{1} \in \mathbb{N} \land \varphi_{75} \to \varphi_{159}$
\item $\varphi_{278} :\equiv \mathrm{SubIdx}(s) \land \forall k \in \mathbb{N}\,(b_{k} = a_{s_{k}}) \to \varepsilon \in \mathbb{R} \to 0 < \varepsilon \to N_{1} \in \mathbb{N} \land \varphi_{75} \to \varphi_{160}$
\item $\varphi_{279} :\equiv \mathrm{SubIdx}(s) \land \forall k \in \mathbb{N}\,(b_{k} = a_{s_{k}}) \to \varepsilon \in \mathbb{R} \to 0 < \varepsilon \to N_{1} \in \mathbb{N} \land \varphi_{75} \to \varphi_{161}$
\item $\varphi_{280} :\equiv \mathrm{Cauchy}(a) \to \mathrm{SubSeq}(b,a) \land \mathrm{Cgt}(b) \to \mathrm{Conv}(b,c) \to \varphi_{252}$
\item $\varphi_{281} :\equiv \mathrm{SubIdx}(s) \land \forall k \in \mathbb{N}\,(b_{k} = a_{s_{k}}) \to \varepsilon \in \mathbb{R} \to 0 < \varepsilon \to N_{1} \in \mathbb{N} \land \varphi_{75} \to \varphi_{169}$
\item $\varphi_{282} :\equiv \mathrm{SubIdx}(s) \land \forall k \in \mathbb{N}\,(b_{k} = a_{s_{k}}) \to \varepsilon \in \mathbb{R} \to 0 < \varepsilon \to N_{1} \in \mathbb{N} \land \varphi_{75} \to \varphi_{171}$
\item $\varphi_{283} :\equiv \mathrm{Cauchy}(a) \to \mathrm{SubSeq}(b,a) \land \mathrm{Cgt}(b) \to \mathrm{Conv}(b,c) \to \varphi_{253}$
\item $\varphi_{284} :\equiv (\varepsilon \in \mathbb{R} \to \varepsilon \cdot 2^{-1} + \varepsilon \cdot 2^{-1} = \varepsilon) \to (\mathrm{Cgt}(a) \to \mathrm{Conv}(a,\alpha) \to \varepsilon \in \mathbb{R} \to \varepsilon \in \mathbb{R}) \to \varphi_{232}$
\item $\varphi_{285} :\equiv (\mathrm{Cgt}(a) \to \mathrm{Conv}(a,\alpha) \to \mathrm{Conv}(a,\alpha)) \to \varphi_{52} \to \varphi_{42} \to \varphi_{181}$
\item $\varphi_{286} :\equiv \mathrm{SubIdx}(s) \land \forall k \in \mathbb{N}\,(b_{k} = a_{s_{k}}) \to \varepsilon \in \mathbb{R} \to 0 < \varepsilon \to N_{1} \in \mathbb{N} \land \varphi_{75} \to \varphi_{178}$
\item $\varphi_{287} :\equiv \mathrm{Cauchy}(a) \to \mathrm{SubSeq}(b,a) \land \mathrm{Cgt}(b) \to \mathrm{Conv}(b,c) \to \varphi_{256}$
\item $\varphi_{288} :\equiv \mathrm{Cauchy}(a) \to \mathrm{SubSeq}(b,a) \land \mathrm{Cgt}(b) \to \mathrm{Conv}(b,c) \to \varphi_{257}$
\item $\varphi_{289} :\equiv \mathrm{Cauchy}(a) \to \mathrm{SubSeq}(b,a) \land \mathrm{Cgt}(b) \to \mathrm{Conv}(b,c) \to \varphi_{258}$
\item $\varphi_{290} :\equiv \mathrm{SubIdx}(s) \land \forall k \in \mathbb{N}\,(b_{k} = a_{s_{k}}) \to \varepsilon \in \mathbb{R} \to 0 < \varepsilon \to N_{1} \in \mathbb{N} \land \varphi_{75} \to \varphi_{182}$
\item $\varphi_{291} :\equiv (\varepsilon \in \mathbb{R} \to 2^{-1} \in \mathbb{R} \to \varepsilon \cdot 2^{-1} \in \mathbb{R}) \to (\mathrm{Cgt}(a) \to \mathrm{Conv}(a,\alpha) \to \varepsilon \in \mathbb{R} \to \varepsilon \in \mathbb{R}) \to \varphi_{237}$
\item $\varphi_{292} :\equiv \mathrm{Cauchy}(a) \to \mathrm{SubSeq}(b,a) \land \mathrm{Cgt}(b) \to \mathrm{Conv}(b,c) \to \varphi_{260}$
\item $\varphi_{293} :\equiv \mathrm{Cauchy}(a) \to \mathrm{SubSeq}(b,a) \land \mathrm{Cgt}(b) \to \mathrm{Conv}(b,c) \to \varphi_{261}$
\item $\varphi_{294} :\equiv \mathrm{SubIdx}(s) \land \forall k \in \mathbb{N}\,(b_{k} = a_{s_{k}}) \to \varepsilon \in \mathbb{R} \to 0 < \varepsilon \to N_{1} \in \mathbb{N} \land \varphi_{75} \to \varphi_{187}$
\item $\varphi_{295} :\equiv \mathrm{Cauchy}(a) \to \mathrm{SubSeq}(b,a) \land \mathrm{Cgt}(b) \to \mathrm{Conv}(b,c) \to \varphi_{263}$
\item $\varphi_{296} :\equiv \mathrm{Cauchy}(a) \to \mathrm{SubSeq}(b,a) \land \mathrm{Cgt}(b) \to \mathrm{Conv}(b,c) \to \varphi_{264}$
\item $\varphi_{297} :\equiv \mathrm{SubIdx}(s) \land \forall k \in \mathbb{N}\,(b_{k} = a_{s_{k}}) \to \varepsilon \in \mathbb{R} \to 0 < \varepsilon \to N_{1} \in \mathbb{N} \land \varphi_{75} \to \varphi_{190}$
\item $\varphi_{298} :\equiv \mathrm{SubIdx}(s) \land \forall k \in \mathbb{N}\,(b_{k} = a_{s_{k}}) \to \varepsilon \in \mathbb{R} \to 0 < \varepsilon \to N_{1} \in \mathbb{N} \land \varphi_{75} \to \varphi_{196}$
\item $\varphi_{299} :\equiv \mathrm{Cauchy}(a) \to \mathrm{SubSeq}(b,a) \land \mathrm{Cgt}(b) \to \mathrm{Conv}(b,c) \to \varphi_{266}$
\item $\varphi_{300} :\equiv \mathrm{Cauchy}(a) \to \mathrm{SubSeq}(b,a) \land \mathrm{Cgt}(b) \to \mathrm{Conv}(b,c) \to \varphi_{267}$
\item $\varphi_{301} :\equiv \mathrm{SubIdx}(s) \land \forall k \in \mathbb{N}\,(b_{k} = a_{s_{k}}) \to \varepsilon \in \mathbb{R} \to 0 < \varepsilon \to N_{1} \in \mathbb{N} \land \varphi_{75} \to \varphi_{198}$
\item $\varphi_{302} :\equiv (\varepsilon \in \mathbb{R} \to 2^{-1} \in \mathbb{R} \to \varepsilon \cdot 2^{-1} \in \mathbb{R}) \to \varphi_{137} \to 2^{-1} \in \mathbb{R} \to \varphi_{165}$
\item $\varphi_{303} :\equiv \mathrm{SubIdx}(s) \land \forall k \in \mathbb{N}\,(b_{k} = a_{s_{k}}) \to \varepsilon \in \mathbb{R} \to 0 < \varepsilon \to N_{1} \in \mathbb{N} \land \varphi_{75} \to \varphi_{204}$
\item $\varphi_{304} :\equiv \mathrm{Cauchy}(a) \to \mathrm{SubSeq}(b,a) \land \mathrm{Cgt}(b) \to \mathrm{Conv}(b,c) \to \varphi_{270}$
\item $\varphi_{305} :\equiv \mathrm{Cauchy}(a) \to \mathrm{SubSeq}(b,a) \land \mathrm{Cgt}(b) \to \mathrm{Conv}(b,c) \to \varphi_{271}$
\item $\varphi_{306} :\equiv \mathrm{Cauchy}(a) \to \mathrm{SubSeq}(b,a) \land \mathrm{Cgt}(b) \to \mathrm{Conv}(b,c) \to \varphi_{272}$
\item $\varphi_{307} :\equiv \mathrm{Cauchy}(a) \to \mathrm{SubSeq}(b,a) \land \mathrm{Cgt}(b) \to \mathrm{Conv}(b,c) \to \varphi_{273}$
\item $\varphi_{308} :\equiv \mathrm{Cauchy}(a) \to \mathrm{SubSeq}(b,a) \land \mathrm{Cgt}(b) \to \mathrm{Conv}(b,c) \to \varphi_{274}$
\item $\varphi_{309} :\equiv \mathrm{SubIdx}(s) \land \forall k \in \mathbb{N}\,(b_{k} = a_{s_{k}}) \to \varepsilon \in \mathbb{R} \to 0 < \varepsilon \to N_{1} \in \mathbb{N} \land \varphi_{75} \to \varphi_{210}$
\item $\varphi_{310} :\equiv \mathrm{Cauchy}(a) \to \mathrm{SubSeq}(b,a) \land \mathrm{Cgt}(b) \to \mathrm{Conv}(b,c) \to \varphi_{277}$
\item $\varphi_{311} :\equiv \mathrm{Cauchy}(a) \to \mathrm{SubSeq}(b,a) \land \mathrm{Cgt}(b) \to \mathrm{Conv}(b,c) \to \varphi_{278}$
\item $\varphi_{312} :\equiv \mathrm{Cauchy}(a) \to \mathrm{SubSeq}(b,a) \land \mathrm{Cgt}(b) \to \mathrm{Conv}(b,c) \to \varphi_{279}$
\item $\varphi_{313} :\equiv \mathrm{Cauchy}(a) \to \mathrm{SubSeq}(b,a) \land \mathrm{Cgt}(b) \to \mathrm{Conv}(b,c) \to \varphi_{281}$
\item $\varphi_{314} :\equiv \mathrm{Cauchy}(a) \to \mathrm{SubSeq}(b,a) \land \mathrm{Cgt}(b) \to \mathrm{Conv}(b,c) \to \varphi_{282}$
\item $\varphi_{315} :\equiv (\mathrm{Cauchy}(a) \to \mathrm{SubSeq}(b,a) \land \mathrm{Cgt}(b) \to \varphi_{176}) \to \varphi_{95} \to \varphi_{76}$
\item $\varphi_{316} :\equiv \mathrm{SubIdx}(s) \land \forall k \in \mathbb{N}\,(b_{k} = a_{s_{k}}) \to \varepsilon \in \mathbb{R} \to 0 < \varepsilon \to N_{1} \in \mathbb{N} \land \varphi_{75} \to \varphi_{219}$
\item $\varphi_{317} :\equiv \mathrm{Cauchy}(a) \to \mathrm{SubSeq}(b,a) \land \mathrm{Cgt}(b) \to \mathrm{Conv}(b,c) \to \varphi_{286}$
\item $\varphi_{318} :\equiv \mathrm{Cauchy}(a) \to \mathrm{SubSeq}(b,a) \land \mathrm{Cgt}(b) \to \mathrm{Conv}(b,c) \to \varphi_{290}$
\item $\varphi_{319} :\equiv \mathrm{Cauchy}(a) \to \mathrm{SubSeq}(b,a) \land \mathrm{Cgt}(b) \to \mathrm{Conv}(b,c) \to \varphi_{294}$
\item $\varphi_{320} :\equiv (\mathrm{Cgt}(a) \to \varphi_{197}) \to \mathrm{Cgt}(a) \to \mathrm{Conv}(a,\alpha) \to \varepsilon \in \mathbb{R} \to \varphi_{201}$
\item $\varphi_{321} :\equiv \mathrm{Cauchy}(a) \to \mathrm{SubSeq}(b,a) \land \mathrm{Cgt}(b) \to \mathrm{Conv}(b,c) \to \varphi_{297}$
\item $\varphi_{322} :\equiv \mathrm{Cauchy}(a) \to \mathrm{SubSeq}(b,a) \land \mathrm{Cgt}(b) \to \mathrm{Conv}(b,c) \to \varphi_{298}$
\item $\varphi_{323} :\equiv (\varphi_{95} \to \forall \varepsilon\,(\varphi_{95})) \to \mathrm{Cauchy}(a) \to \mathrm{SubSeq}(b,a) \land \mathrm{Cgt}(b) \to \varphi_{180}$
\item $\varphi_{324} :\equiv \mathrm{Cauchy}(a) \to \mathrm{SubSeq}(b,a) \land \mathrm{Cgt}(b) \to \mathrm{Conv}(b,c) \to \varphi_{301}$
\item $\varphi_{325} :\equiv (\varphi_{117} \to \forall s\,(\varphi_{27}) \to \varphi_{29}) \to (\varphi_{27} \to \forall s\,(\varphi_{27})) \to \mathrm{Cauchy}(a) \to \varphi_{9}$
\item $\varphi_{326} :\equiv \varphi_{225} \to \mathrm{Cauchy}(a) \to \mathrm{SubSeq}(b,a) \land \mathrm{Cgt}(b) \to \varphi_{176}$
\item $\varphi_{327} :\equiv \mathrm{Cauchy}(a) \to \mathrm{SubSeq}(b,a) \land \mathrm{Cgt}(b) \to \mathrm{Conv}(b,c) \to \varphi_{303}$
\item $\varphi_{328} :\equiv (\mathrm{Cgt}(a) \to \varphi_{211}) \to \mathrm{Cgt}(a) \to \mathrm{Conv}(a,\alpha) \to \varepsilon \in \mathbb{R} \to \varphi_{200}$
\item $\varphi_{329} :\equiv \mathrm{Cauchy}(a) \to \mathrm{SubSeq}(b,a) \land \mathrm{Cgt}(b) \to \mathrm{Conv}(b,c) \to \varphi_{309}$
\item $\varphi_{330} :\equiv \forall v_{753}\,(\varphi_{259}) \to \forall v_{753}\,\forall v_{754}\,(v_{753} \in \mathbb{N} \to v_{754} \in \mathbb{N} \to \varphi_{48})$
\item $\varphi_{331} :\equiv \forall N\,(\forall N_{1}\,(\neg (N_{1} \in \mathbb{N} \land \varphi_{75})) \to \varphi_{93}) \to \forall N\,\forall N_{1}\,(\neg (N_{1} \in \mathbb{N} \land \varphi_{75})) \to \forall N\,(\varphi_{93})$
\item $\varphi_{332} :\equiv (\mathrm{Cauchy}(a) \to \mathrm{SubSeq}(b,a) \land \mathrm{Cgt}(b) \to \varphi_{61}) \to \varphi_{265}$
\item $\varphi_{333} :\equiv (\mathrm{Cgt}(a) \to \mathrm{Conv}(a,\alpha) \to \varphi_{193}) \to \mathrm{Cgt}(a) \to \mathrm{Conv}(a,\alpha) \to \varepsilon \in \mathbb{R} \to 0 < \varepsilon \to \varphi_{214}$
\item $\varphi_{334} :\equiv \mathrm{Cauchy}(a) \to \mathrm{SubSeq}(b,a) \land \mathrm{Cgt}(b) \to \mathrm{Conv}(b,c) \to \varphi_{316}$
\item $\varphi_{335} :\equiv (\mathrm{Cgt}(a) \to \mathrm{Conv}(a,\alpha) \to \varepsilon \in \mathbb{R} \to \varphi_{184}) \to \mathrm{Cgt}(a) \to \mathrm{Conv}(a,\alpha) \to \varepsilon \in \mathbb{R} \to 0 < \varepsilon \to \varphi_{212}$
\item $\varphi_{336} :\equiv (\mathrm{Cgt}(a) \to \mathrm{Conv}(a,\alpha) \to \varepsilon \in \mathbb{R} \to 0 < \varepsilon \to \varphi_{207}) \to \varphi_{179} \to m \in \mathbb{N} \to n \in \mathbb{N} \to N \leq m \to N \leq n \to \lvert a_{m} - a_{n} \rvert < \varepsilon$
\item $\varphi_{337} :\equiv (\mathrm{Seq}(a) \to m \in \mathbb{N} \to a_{m} \in \mathbb{R}) \to (\mathrm{Cgt}(a) \to \mathrm{Conv}(a,\alpha) \to \mathrm{Seq}(a)) \to \varphi_{320}$
\item $\varphi_{338} :\equiv (\mathrm{Seq}(a) \to n \in \mathbb{N} \to a_{n} \in \mathbb{R}) \to (\mathrm{Cgt}(a) \to \mathrm{Conv}(a,\alpha) \to \mathrm{Seq}(a)) \to \varphi_{328}$
\item $\varphi_{339} :\equiv (\mathrm{Cgt}(a) \to \mathrm{Conv}(a,\alpha) \to \varepsilon \in \mathbb{R} \to \varphi_{200}) \to (\mathrm{Cgt}(a) \to \mathrm{Conv}(a,\alpha) \to \alpha \in \mathbb{R}) \to \varphi_{235}$
\item $\varphi_{340} :\equiv (\mathrm{Cgt}(a) \to \mathrm{Conv}(a,\alpha) \to \varepsilon \in \mathbb{R} \to \varphi_{201}) \to (\mathrm{Cgt}(a) \to \mathrm{Conv}(a,\alpha) \to \alpha \in \mathbb{R}) \to \varphi_{236}$
\item $\varphi_{341} :\equiv (\mathrm{Cgt}(a) \to \mathrm{Conv}(a,\alpha) \to \varepsilon \in \mathbb{R} \to 0 < \varepsilon \to \varphi_{213}) \to \mathrm{Cgt}(a) \to \varphi_{241}$
\item $\varphi_{342} :\equiv \forall N\,(\varphi_{154} \to \neg (N \in \mathbb{N} \land \varphi_{92})) \to \forall N\,(\varphi_{154}) \to \forall N\,(\neg (N \in \mathbb{N} \land \varphi_{92}))$
\item $\varphi_{343} :\equiv (\mathrm{Cauchy}(a) \to \varphi_{6}) \to \forall c\,(\varphi_{11} \to \neg \mathrm{Conv}(b,c)) \to \varphi_{130} \to (\varphi_{11} \to \forall c\,(\varphi_{11})) \to \varphi_{209}$
\item $\varphi_{344} :\equiv (\mathrm{Conv}(a,\alpha) \to \varphi_{194}) \to (\mathrm{Cgt}(a) \to \mathrm{Conv}(a,\alpha) \to \mathrm{Conv}(a,\alpha)) \to \varphi_{285}$
\item $\varphi_{345} :\equiv (\mathrm{Conv}(a,c) \leftrightarrow \mathrm{Seq}(a) \land c \in \mathbb{R} \land \varphi_{82}) \to (\mathrm{Cauchy}(a) \to \mathrm{Seq}(a)) \to \varphi_{332}$
\item $\varphi_{346} :\equiv (\mathrm{Cgt}(a) \to \mathrm{Conv}(a,\alpha) \to \varepsilon \in \mathbb{R} \to 0 < \varepsilon \to \varphi_{212}) \to \varphi_{124} \to \mathrm{Cgt}(a) \to \mathrm{Conv}(a,\alpha) \to \varepsilon \in \mathbb{R} \to 0 < \varepsilon \to \varphi_{213}$
\item $\varphi_{347} :\equiv (\mathrm{Cgt}(a) \to \mathrm{Conv}(a,\alpha) \to \varepsilon \in \mathbb{R} \to 0 < \varepsilon \to \varphi_{214}) \to \varphi_{125} \to \mathrm{Cgt}(a) \to \mathrm{Conv}(a,\alpha) \to \varepsilon \in \mathbb{R} \to 0 < \varepsilon \to \varphi_{215}$
\item $\varphi_{348} :\equiv \forall N_{1}\,(\varphi_{217} \to \neg (N_{1} \in \mathbb{N} \land \varphi_{75})) \to \forall N_{1}\,(\varphi_{217}) \to \forall N_{1}\,(\neg (N_{1} \in \mathbb{N} \land \varphi_{75}))$
\item $\varphi_{349} :\equiv (\forall v_{753}\,\forall v_{754}\,(v_{753} \in \mathbb{N} \to v_{754} \in \mathbb{N} \to \varphi_{48}) \to \varphi_{89}) \to \varphi_{321}$
\item $\varphi_{350} :\equiv (\mathrm{Cgt}(a) \to \varphi_{241}) \to \varphi_{139} \to \mathrm{Cgt}(a) \to \mathrm{Conv}(a,\alpha) \to \varepsilon \in \mathbb{R} \to 0 < \varepsilon \to \varphi_{207}$
\item $\varphi_{351} :\equiv (\mathrm{Cauchy}(a) \to \mathrm{SubSeq}(b,a) \land \mathrm{Cgt}(b) \to \varphi_{43}) \to \varphi_{117} \to \varphi_{325}$
\item $\varphi_{352} :\equiv \forall n\,(\varphi_{259} \to n \in \mathbb{N} \to N_{1} \leq n \to \lvert a_{n} - c \rvert < \varepsilon) \to \forall n\,(\varphi_{259}) \to \forall n \in \mathbb{N}\,(N_{1} \leq n \to \lvert a_{n} - c \rvert < \varepsilon)$
\item $\varphi_{353} :\equiv (\mathrm{Cauchy}(a) \to \varepsilon \cdot 2^{-1} \in \mathbb{R} \to 0 < \varepsilon \cdot 2^{-1} \to \varphi_{99}) \to (\mathrm{Cauchy}(a) \to \mathrm{Cauchy}(a)) \to \varphi_{165} \to \varphi_{162} \to \varphi_{231}$
\item $\varphi_{354} :\equiv (\mathrm{Cgt}(a) \to \mathrm{Conv}(a,\alpha) \to \varepsilon \in \mathbb{R} \to 0 < \varepsilon \to \varphi_{215}) \to \varphi_{341}$
\item $\varphi_{355} :\equiv (\mathrm{Seq}(a) \to n \in \mathbb{N} \to a_{n} \in \mathbb{R}) \to (\mathrm{Cauchy}(a) \to \mathrm{Seq}(a)) \to \varphi_{275} \to \varphi_{287}$
\item $\varphi_{356} :\equiv (\forall N\,(\forall N_{2}\,(\varphi_{55}) \to \varphi_{51}) \to \forall N\,\forall N_{2}\,(\varphi_{55}) \to \forall N\,(\varphi_{51})) \to (\forall N_{2}\,(\varphi_{55}) \to \forall N\,\forall N_{2}\,(\varphi_{55})) \to \varphi_{254}$
\item $\varphi_{357} :\equiv (\mathrm{Conv}(b,c) \to \varepsilon \cdot 2^{-1} \in \mathbb{R} \to \varphi_{84}) \to (\mathrm{Cauchy}(a) \to \varphi_{15}) \to \varphi_{165} \to \varphi_{162} \to \varphi_{262}$
\item $\varphi_{358} :\equiv (\mathrm{Seq}(a) \to s_{N_{1} \cup N_{2}} \in \mathbb{N} \to a_{s_{N_{1} \cup N_{2}}} \in \mathbb{R}) \to (\mathrm{Cauchy}(a) \to \mathrm{Seq}(a)) \to \varphi_{292} \to \varphi_{299}$
\item $\varphi_{359} :\equiv \forall v_{754}\,(\varphi_{259} \to v_{753} \in \mathbb{N} \to v_{754} \in \mathbb{N} \to \varphi_{48}) \to \forall v_{754}\,(\varphi_{259}) \to \forall v_{754}\,(v_{753} \in \mathbb{N} \to v_{754} \in \mathbb{N} \to \varphi_{48})$
\item $\varphi_{360} :\equiv \forall v_{753}\,(\varphi_{259} \to \forall v_{754}\,(v_{753} \in \mathbb{N} \to v_{754} \in \mathbb{N} \to \varphi_{48})) \to \varphi_{330}$
\item $\varphi_{361} :\equiv (\mathrm{NSeq}(s) \to N_{1} \cup N_{2} \in \mathbb{N} \to s_{N_{1} \cup N_{2}} \in \mathbb{N}) \to \varphi_{111} \to \varphi_{288} \to \varphi_{292}$
\item $\varphi_{362} :\equiv \varphi_{313} \to (\forall n \in \mathbb{N}\,(N_{2} \leq n \to \lvert b_{n} - c \rvert < \varepsilon \cdot 2^{-1}) \to \varphi_{65}) \to \varphi_{324}$
\item $\varphi_{363} :\equiv (\mathrm{Cgt}(a) \to \mathrm{Conv}(a,\alpha) \to \varepsilon \in \mathbb{R} \to 0 < \varepsilon \to \varphi_{191}) \to \varphi_{354}$
\item $\varphi_{364} :\equiv (\forall m\,(\varphi_{179} \to \varphi_{56}) \to \forall m\,(\varphi_{179}) \to \forall m\,(\varphi_{56})) \to (\varphi_{179} \to \forall m\,(\varphi_{179})) \to \mathrm{Cgt}(a) \to \varphi_{230}$
\item $\varphi_{365} :\equiv \varphi_{181} \to \forall N\,(\varphi_{154} \to \neg (N \in \mathbb{N} \land \varphi_{92})) \to \varphi_{342} \to (\varphi_{154} \to \forall N\,(\varphi_{154})) \to \varphi_{151}$
\item $\varphi_{366} :\equiv (\mathrm{Cgt}(a) \to \mathrm{Conv}(a,\alpha) \to \varepsilon \in \mathbb{R} \to \varphi_{201}) \to (\mathrm{Cgt}(a) \to \mathrm{Conv}(a,\alpha) \to \varepsilon \in \mathbb{R} \to \varphi_{200}) \to \varphi_{363}$
\item $\varphi_{367} :\equiv (\mathrm{Cauchy}(a) \to \mathrm{SubSeq}(b,a) \land \mathrm{Cgt}(b) \to \varphi_{61}) \to \varphi_{165} \to \varphi_{314} \to \varphi_{311} \to \varphi_{312}$
\item $\varphi_{368} :\equiv (\varphi_{217} \to \forall N_{1}\,(\varphi_{217})) \to \forall N\,(\forall N_{1}\,(\neg (N_{1} \in \mathbb{N} \land \varphi_{75})) \to \varphi_{93}) \to \varphi_{331} \to \varphi_{326}$
\item $\varphi_{369} :\equiv (\alpha \in \mathbb{R} \to \varphi_{221}) \to (\mathrm{Cgt}(a) \to \mathrm{Conv}(a,\alpha) \to \alpha \in \mathbb{R}) \to \varphi_{366}$
\item $\varphi_{370} :\equiv \forall n\,(\varphi_{259} \to n \in \mathbb{N} \to N_{1} \leq n \to \lvert a_{n} - c \rvert < \varepsilon) \to \varphi_{352} \to (\varphi_{259} \to \forall n\,(\varphi_{259})) \to \varphi_{296}$
\item $\varphi_{371} :\equiv \forall v_{754}\,(\varphi_{259} \to v_{753} \in \mathbb{N} \to v_{754} \in \mathbb{N} \to \varphi_{48}) \to \varphi_{359} \to (\varphi_{259} \to \forall v_{754}\,(\varphi_{259})) \to \varphi_{327}$
\item $\varphi_{372} :\equiv (\forall N_{2}\,(\varphi_{255} \to \varphi_{55}) \to \forall N_{2}\,(\varphi_{255}) \to \forall N_{2}\,(\varphi_{55})) \to (\varphi_{255} \to \forall N_{2}\,(\varphi_{255})) \to \forall N\,(\forall N_{2}\,(\varphi_{55}) \to \varphi_{51}) \to \varphi_{356}$
\item $\varphi_{373} :\equiv \forall v_{753}\,(\varphi_{259} \to \forall v_{754}\,(v_{753} \in \mathbb{N} \to v_{754} \in \mathbb{N} \to \varphi_{48})) \to \varphi_{360} \to (\varphi_{259} \to \forall v_{753}\,(\varphi_{259})) \to \varphi_{329}$
\item $\varphi_{374} :\equiv (a_{n} \in \mathbb{R} \to a_{s_{N_{1} \cup N_{2}}} \in \mathbb{R} \to c \in \mathbb{R} \to \varphi_{218}) \to \varphi_{287} \to \varphi_{299} \to \varphi_{367}$
\end{abbrevlist}
\end{fullproof}

\chapter{テキスト第4章 関数の極限 の証明木}
本章の定理数 77, ノード数の合計 25656.

\section{第4章 1節: 関数の極限}
% 自動生成 (prover/build.js): FLimE  ノード数 25
\begin{fullproof}{[D1]}{$\forall f\,\forall c\,\forall \alpha\,\forall \varepsilon\,(\mathrm{FLim}(f,c,\alpha) \to \varepsilon \in \mathbb{R} \to 0 < \varepsilon \to \exists \delta\,(\delta \in \mathbb{R} \land 0 < \delta \land \forall x \in \mathrm{dom}(f)\,(0 < \lvert x - c \rvert \to \lvert x - c \rvert < \delta \to \lvert f(x) - \alpha \rvert < \varepsilon)))$}{25}\phantomsection\label{tree:FLimE}
\noindent$\Pi_{1}$:
\begin{ptree}
\AXC{$\forall f\,\forall c\,\forall \alpha\,(\varphi_{6})$\lbl{DefFLim}}
\AXC{$\forall f\,\forall c\,\forall \alpha\,(\varphi_{6}) \to \forall c\,\forall \alpha\,(\varphi_{6})$\lbl{A2}}
\BIC{$\forall c\,\forall \alpha\,(\varphi_{6})$}
\AXC{$\forall c\,\forall \alpha\,(\varphi_{6}) \to \forall \alpha\,(\varphi_{6})$\lbl{A2}}
\BIC{$\forall \alpha\,(\varphi_{6})$}
\AXC{$\forall \alpha\,(\varphi_{6}) \to \varphi_{6}$\lbl{A2}}
\BIC{$\varphi_{6}$}
\end{ptree}
\noindent$\Pi_{2}$:
\begin{ptree}
\AXC{$\stackrel{\Pi_{1}}{\vdots}$}\noLine
\UIC{$\varphi_{6}$}
\AXC{$\mathrm{FLim}(f,c,\alpha) \to \mathrm{FLim}(f,c,\alpha)$\lbl{A1}}
\AXC{$(\mathrm{FLim}(f,c,\alpha) \to \mathrm{FLim}(f,c,\alpha)) \to \varphi_{6} \to \varphi_{4}$\lbl{A1}}
\BIC{$\varphi_{6} \to \varphi_{4}$}
\BIC{$\varphi_{4}$}
\end{ptree}
\noindent$\Pi_{3}$:
\begin{ptree}
\AXC{$\forall \varepsilon \in \mathbb{R}\,(0 < \varepsilon \to \varphi_{1}) \to \varepsilon \in \mathbb{R} \to 0 < \varepsilon \to \varphi_{1}$\lbl{A2}}
\AXC{$\stackrel{\Pi_{2}}{\vdots}$}\noLine
\UIC{$\varphi_{4}$}
\AXC{$\varphi_{4} \to (\forall \varepsilon \in \mathbb{R}\,(0 < \varepsilon \to \varphi_{1}) \to \varepsilon \in \mathbb{R} \to 0 < \varepsilon \to \varphi_{1}) \to \varphi_{3}$\lbl{A1}}
\BIC{$(\forall \varepsilon \in \mathbb{R}\,(0 < \varepsilon \to \varphi_{1}) \to \varepsilon \in \mathbb{R} \to 0 < \varepsilon \to \varphi_{1}) \to \varphi_{3}$}
\BIC{$\varphi_{3}$}
\end{ptree}
\noindent$\Pi_{4}$:
\begin{ptree}
\AXC{$\mathrm{FLim}(f,c,\alpha) \to \varepsilon \in \mathbb{R} \to \varepsilon \in \mathbb{R}$\lbl{A1}}
\AXC{$\stackrel{\Pi_{3}}{\vdots}$}\noLine
\UIC{$\varphi_{3}$}
\AXC{$\varphi_{3} \to (\mathrm{FLim}(f,c,\alpha) \to \varepsilon \in \mathbb{R} \to \varepsilon \in \mathbb{R}) \to \varphi_{5}$\lbl{A1}}
\BIC{$(\mathrm{FLim}(f,c,\alpha) \to \varepsilon \in \mathbb{R} \to \varepsilon \in \mathbb{R}) \to \varphi_{5}$}
\BIC{$\varphi_{5}$}
\end{ptree}
\noindent 主木:
\begin{ptree}
\AXC{$\mathrm{FLim}(f,c,\alpha) \to \varepsilon \in \mathbb{R} \to 0 < \varepsilon \to 0 < \varepsilon$\lbl{A1}}
\AXC{$\stackrel{\Pi_{4}}{\vdots}$}\noLine
\UIC{$\varphi_{5}$}
\BIC{$\mathrm{FLim}(f,c,\alpha) \to \varepsilon \in \mathbb{R} \to 0 < \varepsilon \to \varphi_{1}$}
\UIC{$\forall \varepsilon\,(\mathrm{FLim}(f,c,\alpha) \to \varepsilon \in \mathbb{R} \to 0 < \varepsilon \to \varphi_{1})$}
\UIC{$\forall \alpha\,\forall \varepsilon\,(\mathrm{FLim}(f,c,\alpha) \to \varepsilon \in \mathbb{R} \to 0 < \varepsilon \to \varphi_{1})$}
\UIC{$\forall c\,\forall \alpha\,\forall \varepsilon\,(\mathrm{FLim}(f,c,\alpha) \to \varepsilon \in \mathbb{R} \to 0 < \varepsilon \to \varphi_{1})$}
\UIC{$\varphi_{2}$}
\end{ptree}
\begin{abbrevlist}
\item $\varphi_{1} :\equiv \exists \delta\,(\delta \in \mathbb{R} \land 0 < \delta \land \forall x \in \mathrm{dom}(f)\,(0 < \lvert x - c \rvert \to \lvert x - c \rvert < \delta \to \lvert f(x) - \alpha \rvert < \varepsilon))$
\item $\varphi_{2} :\equiv \forall f\,\forall c\,\forall \alpha\,\forall \varepsilon\,(\mathrm{FLim}(f,c,\alpha) \to \varepsilon \in \mathbb{R} \to 0 < \varepsilon \to \varphi_{1})$
\item $\varphi_{3} :\equiv \mathrm{FLim}(f,c,\alpha) \to \varepsilon \in \mathbb{R} \to 0 < \varepsilon \to \varepsilon \in \mathbb{R} \to 0 < \varepsilon \to \varphi_{1}$
\item $\varphi_{4} :\equiv \mathrm{FLim}(f,c,\alpha) \to \varepsilon \in \mathbb{R} \to 0 < \varepsilon \to \forall \varepsilon \in \mathbb{R}\,(0 < \varepsilon \to \varphi_{1})$
\item $\varphi_{5} :\equiv (\mathrm{FLim}(f,c,\alpha) \to \varepsilon \in \mathbb{R} \to 0 < \varepsilon \to 0 < \varepsilon) \to \mathrm{FLim}(f,c,\alpha) \to \varepsilon \in \mathbb{R} \to 0 < \varepsilon \to \varphi_{1}$
\item $\varphi_{6} :\equiv \mathrm{FLim}(f,c,\alpha) \leftrightarrow \mathrm{RFn}(f) \land c \in \mathbb{R} \land \alpha \in \mathbb{R} \land \mathrm{PNbhd}(f,c) \land \forall \varepsilon \in \mathbb{R}\,(0 < \varepsilon \to \varphi_{1})$
\end{abbrevlist}
\end{fullproof}

% 自動生成 (prover/build.js): FLimRFn  ノード数 14
\begin{fullproof}{[D2]}{$\forall f\,\forall c\,\forall \alpha\,(\mathrm{FLim}(f,c,\alpha) \to \mathrm{RFn}(f))$}{14}\phantomsection\label{tree:FLimRFn}
\noindent$\Pi_{1}$:
\begin{ptree}
\AXC{$\forall f\,\forall c\,\forall \alpha\,(\varphi_{1})$\lbl{DefFLim}}
\AXC{$\forall f\,\forall c\,\forall \alpha\,(\varphi_{1}) \to \forall c\,\forall \alpha\,(\varphi_{1})$\lbl{A2}}
\BIC{$\forall c\,\forall \alpha\,(\varphi_{1})$}
\AXC{$\forall c\,\forall \alpha\,(\varphi_{1}) \to \forall \alpha\,(\varphi_{1})$\lbl{A2}}
\BIC{$\forall \alpha\,(\varphi_{1})$}
\AXC{$\forall \alpha\,(\varphi_{1}) \to \varphi_{1}$\lbl{A2}}
\BIC{$\varphi_{1}$}
\end{ptree}
\noindent 主木:
\begin{ptree}
\AXC{$\stackrel{\Pi_{1}}{\vdots}$}\noLine
\UIC{$\varphi_{1}$}
\AXC{$\mathrm{FLim}(f,c,\alpha) \to \mathrm{FLim}(f,c,\alpha)$\lbl{A1}}
\AXC{$\varphi_{2}$\lbl{A1}}
\BIC{$\varphi_{1} \to \mathrm{FLim}(f,c,\alpha) \to \mathrm{RFn}(f)$}
\BIC{$\mathrm{FLim}(f,c,\alpha) \to \mathrm{RFn}(f)$}
\UIC{$\forall \alpha\,(\mathrm{FLim}(f,c,\alpha) \to \mathrm{RFn}(f))$}
\UIC{$\forall c\,\forall \alpha\,(\mathrm{FLim}(f,c,\alpha) \to \mathrm{RFn}(f))$}
\UIC{$\forall f\,\forall c\,\forall \alpha\,(\mathrm{FLim}(f,c,\alpha) \to \mathrm{RFn}(f))$}
\end{ptree}
\begin{abbrevlist}
\item $\varphi_{1} :\equiv \mathrm{FLim}(f,c,\alpha) \leftrightarrow \mathrm{RFn}(f) \land c \in \mathbb{R} \land \alpha \in \mathbb{R} \land \mathrm{PNbhd}(f,c) \land \forall \varepsilon \in \mathbb{R}\,(0 < \varepsilon \to \exists \delta\,(\delta \in \mathbb{R} \land 0 < \delta \land \forall x \in \mathrm{dom}(f)\,(0 < \lvert x - c \rvert \to \lvert x - c \rvert < \delta \to \lvert f(x) - \alpha \rvert < \varepsilon)))$
\item $\varphi_{2} :\equiv (\mathrm{FLim}(f,c,\alpha) \to \mathrm{FLim}(f,c,\alpha)) \to \varphi_{1} \to \mathrm{FLim}(f,c,\alpha) \to \mathrm{RFn}(f)$
\end{abbrevlist}
\end{fullproof}

% 自動生成 (prover/build.js): FLimCR  ノード数 14
\begin{fullproof}{[D3]}{$\forall f\,\forall c\,\forall \alpha\,(\mathrm{FLim}(f,c,\alpha) \to c \in \mathbb{R})$}{14}\phantomsection\label{tree:FLimCR}
\noindent$\Pi_{1}$:
\begin{ptree}
\AXC{$\forall f\,\forall c\,\forall \alpha\,(\varphi_{1})$\lbl{DefFLim}}
\AXC{$\forall f\,\forall c\,\forall \alpha\,(\varphi_{1}) \to \forall c\,\forall \alpha\,(\varphi_{1})$\lbl{A2}}
\BIC{$\forall c\,\forall \alpha\,(\varphi_{1})$}
\AXC{$\forall c\,\forall \alpha\,(\varphi_{1}) \to \forall \alpha\,(\varphi_{1})$\lbl{A2}}
\BIC{$\forall \alpha\,(\varphi_{1})$}
\AXC{$\forall \alpha\,(\varphi_{1}) \to \varphi_{1}$\lbl{A2}}
\BIC{$\varphi_{1}$}
\end{ptree}
\noindent 主木:
\begin{ptree}
\AXC{$\stackrel{\Pi_{1}}{\vdots}$}\noLine
\UIC{$\varphi_{1}$}
\AXC{$\mathrm{FLim}(f,c,\alpha) \to \mathrm{FLim}(f,c,\alpha)$\lbl{A1}}
\AXC{$\varphi_{2}$\lbl{A1}}
\BIC{$\varphi_{1} \to \mathrm{FLim}(f,c,\alpha) \to c \in \mathbb{R}$}
\BIC{$\mathrm{FLim}(f,c,\alpha) \to c \in \mathbb{R}$}
\UIC{$\forall \alpha\,(\mathrm{FLim}(f,c,\alpha) \to c \in \mathbb{R})$}
\UIC{$\forall c\,\forall \alpha\,(\mathrm{FLim}(f,c,\alpha) \to c \in \mathbb{R})$}
\UIC{$\forall f\,\forall c\,\forall \alpha\,(\mathrm{FLim}(f,c,\alpha) \to c \in \mathbb{R})$}
\end{ptree}
\begin{abbrevlist}
\item $\varphi_{1} :\equiv \mathrm{FLim}(f,c,\alpha) \leftrightarrow \mathrm{RFn}(f) \land c \in \mathbb{R} \land \alpha \in \mathbb{R} \land \mathrm{PNbhd}(f,c) \land \forall \varepsilon \in \mathbb{R}\,(0 < \varepsilon \to \exists \delta\,(\delta \in \mathbb{R} \land 0 < \delta \land \forall x \in \mathrm{dom}(f)\,(0 < \lvert x - c \rvert \to \lvert x - c \rvert < \delta \to \lvert f(x) - \alpha \rvert < \varepsilon)))$
\item $\varphi_{2} :\equiv (\mathrm{FLim}(f,c,\alpha) \to \mathrm{FLim}(f,c,\alpha)) \to \varphi_{1} \to \mathrm{FLim}(f,c,\alpha) \to c \in \mathbb{R}$
\end{abbrevlist}
\end{fullproof}

% 自動生成 (prover/build.js): FLimAR  ノード数 14
\begin{fullproof}{[D4]}{$\forall f\,\forall c\,\forall \alpha\,(\mathrm{FLim}(f,c,\alpha) \to \alpha \in \mathbb{R})$}{14}\phantomsection\label{tree:FLimAR}
\noindent$\Pi_{1}$:
\begin{ptree}
\AXC{$\forall f\,\forall c\,\forall \alpha\,(\varphi_{1})$\lbl{DefFLim}}
\AXC{$\forall f\,\forall c\,\forall \alpha\,(\varphi_{1}) \to \forall c\,\forall \alpha\,(\varphi_{1})$\lbl{A2}}
\BIC{$\forall c\,\forall \alpha\,(\varphi_{1})$}
\AXC{$\forall c\,\forall \alpha\,(\varphi_{1}) \to \forall \alpha\,(\varphi_{1})$\lbl{A2}}
\BIC{$\forall \alpha\,(\varphi_{1})$}
\AXC{$\forall \alpha\,(\varphi_{1}) \to \varphi_{1}$\lbl{A2}}
\BIC{$\varphi_{1}$}
\end{ptree}
\noindent 主木:
\begin{ptree}
\AXC{$\stackrel{\Pi_{1}}{\vdots}$}\noLine
\UIC{$\varphi_{1}$}
\AXC{$\mathrm{FLim}(f,c,\alpha) \to \mathrm{FLim}(f,c,\alpha)$\lbl{A1}}
\AXC{$\varphi_{2}$\lbl{A1}}
\BIC{$\varphi_{1} \to \mathrm{FLim}(f,c,\alpha) \to \alpha \in \mathbb{R}$}
\BIC{$\mathrm{FLim}(f,c,\alpha) \to \alpha \in \mathbb{R}$}
\UIC{$\forall \alpha\,(\mathrm{FLim}(f,c,\alpha) \to \alpha \in \mathbb{R})$}
\UIC{$\forall c\,\forall \alpha\,(\mathrm{FLim}(f,c,\alpha) \to \alpha \in \mathbb{R})$}
\UIC{$\forall f\,\forall c\,\forall \alpha\,(\mathrm{FLim}(f,c,\alpha) \to \alpha \in \mathbb{R})$}
\end{ptree}
\begin{abbrevlist}
\item $\varphi_{1} :\equiv \mathrm{FLim}(f,c,\alpha) \leftrightarrow \mathrm{RFn}(f) \land c \in \mathbb{R} \land \alpha \in \mathbb{R} \land \mathrm{PNbhd}(f,c) \land \forall \varepsilon \in \mathbb{R}\,(0 < \varepsilon \to \exists \delta\,(\delta \in \mathbb{R} \land 0 < \delta \land \forall x \in \mathrm{dom}(f)\,(0 < \lvert x - c \rvert \to \lvert x - c \rvert < \delta \to \lvert f(x) - \alpha \rvert < \varepsilon)))$
\item $\varphi_{2} :\equiv (\mathrm{FLim}(f,c,\alpha) \to \mathrm{FLim}(f,c,\alpha)) \to \varphi_{1} \to \mathrm{FLim}(f,c,\alpha) \to \alpha \in \mathbb{R}$
\end{abbrevlist}
\end{fullproof}

% 自動生成 (prover/build.js): FLimNb  ノード数 14
\begin{fullproof}{[D5]}{$\forall f\,\forall c\,\forall \alpha\,(\mathrm{FLim}(f,c,\alpha) \to \mathrm{PNbhd}(f,c))$}{14}\phantomsection\label{tree:FLimNb}
\noindent$\Pi_{1}$:
\begin{ptree}
\AXC{$\forall f\,\forall c\,\forall \alpha\,(\varphi_{1})$\lbl{DefFLim}}
\AXC{$\forall f\,\forall c\,\forall \alpha\,(\varphi_{1}) \to \forall c\,\forall \alpha\,(\varphi_{1})$\lbl{A2}}
\BIC{$\forall c\,\forall \alpha\,(\varphi_{1})$}
\AXC{$\forall c\,\forall \alpha\,(\varphi_{1}) \to \forall \alpha\,(\varphi_{1})$\lbl{A2}}
\BIC{$\forall \alpha\,(\varphi_{1})$}
\AXC{$\forall \alpha\,(\varphi_{1}) \to \varphi_{1}$\lbl{A2}}
\BIC{$\varphi_{1}$}
\end{ptree}
\noindent 主木:
\begin{ptree}
\AXC{$\stackrel{\Pi_{1}}{\vdots}$}\noLine
\UIC{$\varphi_{1}$}
\AXC{$\mathrm{FLim}(f,c,\alpha) \to \mathrm{FLim}(f,c,\alpha)$\lbl{A1}}
\AXC{$\varphi_{2}$\lbl{A1}}
\BIC{$\varphi_{1} \to \mathrm{FLim}(f,c,\alpha) \to \mathrm{PNbhd}(f,c)$}
\BIC{$\mathrm{FLim}(f,c,\alpha) \to \mathrm{PNbhd}(f,c)$}
\UIC{$\forall \alpha\,(\mathrm{FLim}(f,c,\alpha) \to \mathrm{PNbhd}(f,c))$}
\UIC{$\forall c\,\forall \alpha\,(\mathrm{FLim}(f,c,\alpha) \to \mathrm{PNbhd}(f,c))$}
\UIC{$\forall f\,\forall c\,\forall \alpha\,(\mathrm{FLim}(f,c,\alpha) \to \mathrm{PNbhd}(f,c))$}
\end{ptree}
\begin{abbrevlist}
\item $\varphi_{1} :\equiv \mathrm{FLim}(f,c,\alpha) \leftrightarrow \mathrm{RFn}(f) \land c \in \mathbb{R} \land \alpha \in \mathbb{R} \land \mathrm{PNbhd}(f,c) \land \forall \varepsilon \in \mathbb{R}\,(0 < \varepsilon \to \exists \delta\,(\delta \in \mathbb{R} \land 0 < \delta \land \forall x \in \mathrm{dom}(f)\,(0 < \lvert x - c \rvert \to \lvert x - c \rvert < \delta \to \lvert f(x) - \alpha \rvert < \varepsilon)))$
\item $\varphi_{2} :\equiv (\mathrm{FLim}(f,c,\alpha) \to \mathrm{FLim}(f,c,\alpha)) \to \varphi_{1} \to \mathrm{FLim}(f,c,\alpha) \to \mathrm{PNbhd}(f,c)$
\end{abbrevlist}
\end{fullproof}

% 自動生成 (prover/build.js): AbsZeroEq  ノード数 30
\begin{fullproof}{[D6]}{$\forall x \in \mathbb{R}\,(\lvert x \rvert = 0 \to x = 0)$}{30}\phantomsection\label{tree:AbsZeroEq}
\noindent$\Pi_{1}$:
\begin{ptree}
\AXC{$x \in \mathbb{R} \to \lvert x \rvert = 0 \to x \neq 0 \to x \in \mathbb{R}$\lbl{A1}}
\AXC{$\forall x \in \mathbb{R}\,(x \neq 0 \to 0 < \lvert x \rvert)$\lbl{[A14]}}
\AXC{$\varphi_{3}$\lbl{A2}}
\BIC{$x \in \mathbb{R} \to x \neq 0 \to 0 < \lvert x \rvert$}
\AXC{$\varphi_{8}$\lbl{A1}}
\BIC{$(x \in \mathbb{R} \to \lvert x \rvert = 0 \to x \neq 0 \to x \in \mathbb{R}) \to \varphi_{6}$}
\BIC{$\varphi_{6}$}
\end{ptree}
\noindent$\Pi_{2}$:
\begin{ptree}
\AXC{$x \in \mathbb{R} \to \lvert x \rvert = 0 \to x \neq 0 \to x \neq 0$\lbl{A1}}
\AXC{$\stackrel{\Pi_{1}}{\vdots}$}\noLine
\UIC{$\varphi_{6}$}
\BIC{$x \in \mathbb{R} \to \lvert x \rvert = 0 \to x \neq 0 \to 0 < \lvert x \rvert$}
\AXC{$x \in \mathbb{R} \to \lvert x \rvert = 0 \to \lvert x \rvert = 0$\lbl{A1}}
\AXC{$\varphi_{7}$\lbl{A1}}
\BIC{$(x \in \mathbb{R} \to \lvert x \rvert = 0 \to x \neq 0 \to 0 < \lvert x \rvert) \to \varphi_{5}$}
\BIC{$\varphi_{5}$}
\end{ptree}
\noindent$\Pi_{3}$:
\begin{ptree}
\AXC{$z = w \to 0 < z \to 0 < w$\lbl{A6}}
\UIC{$\forall w\,(z = w \to 0 < z \to 0 < w)$}
\UIC{$\forall z\,\forall w\,(z = w \to 0 < z \to 0 < w)$}
\AXC{$\varphi_{1}$\lbl{A2}}
\BIC{$\forall w\,(\lvert x \rvert = w \to 0 < \lvert x \rvert \to 0 < w)$}
\AXC{$\varphi_{2}$\lbl{A2}}
\BIC{$\lvert x \rvert = 0 \to 0 < \lvert x \rvert \to 0 < 0$}
\AXC{$\stackrel{\Pi_{2}}{\vdots}$}\noLine
\UIC{$\varphi_{5}$}
\BIC{$x \in \mathbb{R} \to \lvert x \rvert = 0 \to x \neq 0 \to 0 < 0$}
\AXC{$(x \in \mathbb{R} \to \lvert x \rvert = 0 \to x \neq 0 \to 0 < 0) \to \varphi_{4}$\lbl{A1}}
\BIC{$\varphi_{4}$}
\end{ptree}
\noindent 主木:
\begin{ptree}
\AXC{$\forall x\,(\neg x < x)$\lbl{[O1]}}
\AXC{$\forall x\,(\neg x < x) \to \neg 0 < 0$\lbl{A2}}
\BIC{$\neg 0 < 0$}
\AXC{$\neg 0 < 0 \to x \in \mathbb{R} \to \lvert x \rvert = 0 \to x \neq 0 \to \neg 0 < 0$\lbl{A1}}
\BIC{$x \in \mathbb{R} \to \lvert x \rvert = 0 \to x \neq 0 \to \neg 0 < 0$}
\AXC{$\stackrel{\Pi_{3}}{\vdots}$}\noLine
\UIC{$\varphi_{4}$}
\BIC{$x \in \mathbb{R} \to \lvert x \rvert = 0 \to x = 0$}
\UIC{$\forall x \in \mathbb{R}\,(\lvert x \rvert = 0 \to x = 0)$}
\end{ptree}
\begin{abbrevlist}
\item $\varphi_{1} :\equiv \forall z\,\forall w\,(z = w \to 0 < z \to 0 < w) \to \forall w\,(\lvert x \rvert = w \to 0 < \lvert x \rvert \to 0 < w)$
\item $\varphi_{2} :\equiv \forall w\,(\lvert x \rvert = w \to 0 < \lvert x \rvert \to 0 < w) \to \lvert x \rvert = 0 \to 0 < \lvert x \rvert \to 0 < 0$
\item $\varphi_{3} :\equiv \forall x \in \mathbb{R}\,(x \neq 0 \to 0 < \lvert x \rvert) \to x \in \mathbb{R} \to x \neq 0 \to 0 < \lvert x \rvert$
\item $\varphi_{4} :\equiv (x \in \mathbb{R} \to \lvert x \rvert = 0 \to x \neq 0 \to \neg 0 < 0) \to x \in \mathbb{R} \to \lvert x \rvert = 0 \to x = 0$
\item $\varphi_{5} :\equiv (\lvert x \rvert = 0 \to 0 < \lvert x \rvert \to 0 < 0) \to x \in \mathbb{R} \to \lvert x \rvert = 0 \to x \neq 0 \to 0 < 0$
\item $\varphi_{6} :\equiv (x \in \mathbb{R} \to \lvert x \rvert = 0 \to x \neq 0 \to x \neq 0) \to x \in \mathbb{R} \to \lvert x \rvert = 0 \to x \neq 0 \to 0 < \lvert x \rvert$
\item $\varphi_{7} :\equiv (x \in \mathbb{R} \to \lvert x \rvert = 0 \to \lvert x \rvert = 0) \to (x \in \mathbb{R} \to \lvert x \rvert = 0 \to x \neq 0 \to 0 < \lvert x \rvert) \to \varphi_{5}$
\item $\varphi_{8} :\equiv (x \in \mathbb{R} \to x \neq 0 \to 0 < \lvert x \rvert) \to (x \in \mathbb{R} \to \lvert x \rvert = 0 \to x \neq 0 \to x \in \mathbb{R}) \to \varphi_{6}$
\end{abbrevlist}
\end{fullproof}

% 自動生成 (prover/build.js): SubZeroEq  ノード数 103
\begin{fullproof}{[D7]}{$\forall x\,\forall y\,(x \in \mathbb{R} \to y \in \mathbb{R} \to x - y = 0 \to x = y)$}{103}\phantomsection\label{tree:SubZeroEq}
\noindent$\Pi_{1}$:
\begin{ptree}
\AXC{$y = z \to x = y \to x = z$\lbl{A6}}
\AXC{$(y = z \to x = y \to x = z) \to x = y \to y = z \to x = z$\lbl{A1}}
\BIC{$x = y \to y = z \to x = z$}
\UIC{$\forall z\,(x = y \to y = z \to x = z)$}
\UIC{$\forall y\,\forall z\,(x = y \to y = z \to x = z)$}
\UIC{$\forall x\,\forall y\,\forall z\,(x = y \to y = z \to x = z)$}
\AXC{$\varphi_{2}$\lbl{A2}}
\BIC{$\forall y\,\forall z\,(x = y \to y = z \to x = z)$}
\AXC{$\varphi_{3}$\lbl{A2}}
\BIC{$\forall z\,(x = y + 0 \to y + 0 = z \to x = z)$}
\AXC{$\varphi_{5}$\lbl{A2}}
\BIC{$x = y + 0 \to y + 0 = y \to x = y$}
\end{ptree}
\noindent$\Pi_{2}$:
\begin{ptree}
\AXC{$\forall x \in \mathbb{R}\,(x + 0 = x)$\lbl{(R4)}}
\AXC{$\forall x \in \mathbb{R}\,(x + 0 = x) \to y \in \mathbb{R} \to y + 0 = y$\lbl{A2}}
\BIC{$y \in \mathbb{R} \to y + 0 = y$}
\AXC{$(y \in \mathbb{R} \to y + 0 = y) \to \varphi_{20}$\lbl{A1}}
\BIC{$\varphi_{20}$}
\end{ptree}
\noindent$\Pi_{3}$:
\begin{ptree}
\AXC{$y = z \to x = y \to x = z$\lbl{A6}}
\AXC{$(y = z \to x = y \to x = z) \to x = y \to y = z \to x = z$\lbl{A1}}
\BIC{$x = y \to y = z \to x = z$}
\UIC{$\forall z\,(x = y \to y = z \to x = z)$}
\UIC{$\forall y\,\forall z\,(x = y \to y = z \to x = z)$}
\UIC{$\forall x\,\forall y\,\forall z\,(x = y \to y = z \to x = z)$}
\AXC{$\varphi_{2}$\lbl{A2}}
\BIC{$\forall y\,\forall z\,(x = y \to y = z \to x = z)$}
\AXC{$\varphi_{4}$\lbl{A2}}
\BIC{$\forall z\,(x = 0 + y \to 0 + y = z \to x = z)$}
\AXC{$\varphi_{8}$\lbl{A2}}
\BIC{$x = 0 + y \to 0 + y = y + 0 \to x = y + 0$}
\end{ptree}
\noindent$\Pi_{4}$:
\begin{ptree}
\AXC{$\forall x\,\forall y\,(x \in \mathbb{R} \to y \in \mathbb{R} \to x + y = y + x)$\lbl{(R3)}}
\AXC{$\varphi_{17}$\lbl{A2}}
\BIC{$\forall y\,(0 \in \mathbb{R} \to y \in \mathbb{R} \to 0 + y = y + 0)$}
\AXC{$\varphi_{14}$\lbl{A2}}
\BIC{$0 \in \mathbb{R} \to y \in \mathbb{R} \to 0 + y = y + 0$}
\AXC{$(0 \in \mathbb{R} \to y \in \mathbb{R} \to 0 + y = y + 0) \to 0 \in \mathbb{R} \to \varphi_{21}$\lbl{A1}}
\BIC{$0 \in \mathbb{R} \to \varphi_{21}$}
\end{ptree}
\noindent$\Pi_{5}$:
\begin{ptree}
\AXC{$x \in \mathbb{R} \to y \in \mathbb{R} \to x - y = 0 \to y \in \mathbb{R}$\lbl{A1}}
\AXC{$0 \in \mathbb{R}$\lbl{(R1)}}
\AXC{$\stackrel{\Pi_{4}}{\vdots}$}\noLine
\UIC{$0 \in \mathbb{R} \to \varphi_{21}$}
\BIC{$\varphi_{21}$}
\BIC{$x \in \mathbb{R} \to y \in \mathbb{R} \to x - y = 0 \to 0 + y = y + 0$}
\end{ptree}
\noindent$\Pi_{6}$:
\begin{ptree}
\AXC{$y = z \to x = y \to x = z$\lbl{A6}}
\AXC{$(y = z \to x = y \to x = z) \to x = y \to y = z \to x = z$\lbl{A1}}
\BIC{$x = y \to y = z \to x = z$}
\UIC{$\forall z\,(x = y \to y = z \to x = z)$}
\UIC{$\forall y\,\forall z\,(x = y \to y = z \to x = z)$}
\UIC{$\forall x\,\forall y\,\forall z\,(x = y \to y = z \to x = z)$}
\AXC{$\varphi_{2}$\lbl{A2}}
\BIC{$\forall y\,\forall z\,(x = y \to y = z \to x = z)$}
\AXC{$\varphi_{7}$\lbl{A2}}
\BIC{$\forall z\,(x = x - y + y \to x - y + y = z \to x = z)$}
\AXC{$\varphi_{16}$\lbl{A2}}
\BIC{$x = x - y + y \to x - y + y = 0 + y \to x = 0 + y$}
\end{ptree}
\noindent$\Pi_{7}$:
\begin{ptree}
\AXC{$z = w \to x - y + y = z + y \to x - y + y = w + y$\lbl{A6}}
\UIC{$\forall w\,(z = w \to x - y + y = z + y \to x - y + y = w + y)$}
\UIC{$\forall z\,\forall w\,(z = w \to x - y + y = z + y \to x - y + y = w + y)$}
\AXC{$\varphi_{23}$\lbl{A2}}
\BIC{$\varphi_{6}$}
\AXC{$\varphi_{6} \to x - y = 0 \to x - y + y = x - y + y \to x - y + y = 0 + y$\lbl{A2}}
\BIC{$x - y = 0 \to x - y + y = x - y + y \to x - y + y = 0 + y$}
\end{ptree}
\noindent$\Pi_{8}$:
\begin{ptree}
\AXC{$x = x$\lbl{A5}}
\UIC{$\forall x\,(x = x)$}
\AXC{$\forall x\,(x = x) \to x - y + y = x - y + y$\lbl{A2}}
\BIC{$x - y + y = x - y + y$}
\AXC{$x \in \mathbb{R} \to y \in \mathbb{R} \to x - y = 0 \to x - y = 0$\lbl{A1}}
\AXC{$\varphi_{27}$\lbl{A1}}
\BIC{$x - y + y = x - y + y \to \varphi_{24}$}
\BIC{$\varphi_{24}$}
\end{ptree}
\noindent$\Pi_{9}$:
\begin{ptree}
\AXC{$x = x$\lbl{A5}}
\AXC{$x = y \to x = x \to y = x$\lbl{A6}}
\AXC{$(x = y \to x = x \to y = x) \to x = x \to x = y \to y = x$\lbl{A1}}
\BIC{$x = x \to x = y \to y = x$}
\BIC{$x = y \to y = x$}
\UIC{$\forall y\,(x = y \to y = x)$}
\UIC{$\forall x\,\forall y\,(x = y \to y = x)$}
\end{ptree}
\noindent$\Pi_{10}$:
\begin{ptree}
\AXC{$\stackrel{\Pi_{9}}{\vdots}$}\noLine
\UIC{$\forall x\,\forall y\,(x = y \to y = x)$}
\AXC{$\forall x\,\forall y\,(x = y \to y = x) \to \forall y\,(v_{406} = y \to y = v_{406})$\lbl{A2}}
\BIC{$\forall y\,(v_{406} = y \to y = v_{406})$}
\AXC{$\varphi_{1}$\lbl{A2}}
\BIC{$v_{406} = v_{407} \to v_{407} = v_{406}$}
\UIC{$\forall v_{407}\,(v_{406} = v_{407} \to v_{407} = v_{406})$}
\UIC{$\forall v_{406}\,\forall v_{407}\,(v_{406} = v_{407} \to v_{407} = v_{406})$}
\AXC{$\varphi_{11}$\lbl{A2}}
\BIC{$\forall v_{407}\,(x - y + y = v_{407} \to v_{407} = x - y + y)$}
\AXC{$\varphi_{10}$\lbl{A2}}
\BIC{$x - y + y = x \to x = x - y + y$}
\end{ptree}
\noindent$\Pi_{11}$:
\begin{ptree}
\AXC{$\forall x\,\forall y\,(x \in \mathbb{R} \to y \in \mathbb{R} \to x - y + y = x)$\lbl{[F12]}}
\AXC{$\varphi_{18}$\lbl{A2}}
\BIC{$\forall y\,(x \in \mathbb{R} \to y \in \mathbb{R} \to x - y + y = x)$}
\AXC{$\varphi_{15}$\lbl{A2}}
\BIC{$x \in \mathbb{R} \to y \in \mathbb{R} \to x - y + y = x$}
\AXC{$\varphi_{28}$\lbl{A1}}
\BIC{$(x \in \mathbb{R} \to y \in \mathbb{R} \to x - y = 0 \to x \in \mathbb{R}) \to \varphi_{22}$}
\end{ptree}
\noindent$\Pi_{12}$:
\begin{ptree}
\AXC{$x \in \mathbb{R} \to y \in \mathbb{R} \to x - y = 0 \to x \in \mathbb{R}$\lbl{A1}}
\AXC{$\stackrel{\Pi_{11}}{\vdots}$}\noLine
\UIC{$(x \in \mathbb{R} \to y \in \mathbb{R} \to x - y = 0 \to x \in \mathbb{R}) \to \varphi_{22}$}
\BIC{$\varphi_{22}$}
\end{ptree}
\noindent$\Pi_{13}$:
\begin{ptree}
\AXC{$x \in \mathbb{R} \to y \in \mathbb{R} \to x - y = 0 \to y \in \mathbb{R}$\lbl{A1}}
\AXC{$\stackrel{\Pi_{12}}{\vdots}$}\noLine
\UIC{$\varphi_{22}$}
\BIC{$x \in \mathbb{R} \to y \in \mathbb{R} \to x - y = 0 \to x - y + y = x$}
\AXC{$(x \in \mathbb{R} \to y \in \mathbb{R} \to x - y = 0 \to x - y + y = x) \to \varphi_{12}$\lbl{A1}}
\BIC{$\varphi_{12}$}
\end{ptree}
\noindent$\Pi_{14}$:
\begin{ptree}
\AXC{$\stackrel{\Pi_{7}}{\vdots}$}\noLine
\UIC{$x - y = 0 \to x - y + y = x - y + y \to x - y + y = 0 + y$}
\AXC{$\stackrel{\Pi_{8}}{\vdots}$}\noLine
\UIC{$\varphi_{24}$}
\BIC{$x \in \mathbb{R} \to y \in \mathbb{R} \to x - y = 0 \to x - y + y = 0 + y$}
\AXC{$\stackrel{\Pi_{10}}{\vdots}$}\noLine
\UIC{$x - y + y = x \to x = x - y + y$}
\AXC{$\stackrel{\Pi_{13}}{\vdots}$}\noLine
\UIC{$\varphi_{12}$}
\BIC{$x \in \mathbb{R} \to y \in \mathbb{R} \to x - y = 0 \to x = x - y + y$}
\AXC{$\varphi_{29}$\lbl{A1}}
\BIC{$(x \in \mathbb{R} \to y \in \mathbb{R} \to x - y = 0 \to x - y + y = 0 + y) \to \varphi_{19}$}
\BIC{$\varphi_{19}$}
\end{ptree}
\noindent$\Pi_{15}$:
\begin{ptree}
\AXC{$\stackrel{\Pi_{5}}{\vdots}$}\noLine
\UIC{$x \in \mathbb{R} \to y \in \mathbb{R} \to x - y = 0 \to 0 + y = y + 0$}
\AXC{$\stackrel{\Pi_{6}}{\vdots}$}\noLine
\UIC{$x = x - y + y \to x - y + y = 0 + y \to x = 0 + y$}
\AXC{$\stackrel{\Pi_{14}}{\vdots}$}\noLine
\UIC{$\varphi_{19}$}
\BIC{$x \in \mathbb{R} \to y \in \mathbb{R} \to x - y = 0 \to x = 0 + y$}
\AXC{$\varphi_{26}$\lbl{A1}}
\BIC{$(x \in \mathbb{R} \to y \in \mathbb{R} \to x - y = 0 \to 0 + y = y + 0) \to \varphi_{13}$}
\BIC{$\varphi_{13}$}
\end{ptree}
\noindent$\Pi_{16}$:
\begin{ptree}
\AXC{$x \in \mathbb{R} \to y \in \mathbb{R} \to x - y = 0 \to y \in \mathbb{R}$\lbl{A1}}
\AXC{$\stackrel{\Pi_{2}}{\vdots}$}\noLine
\UIC{$\varphi_{20}$}
\BIC{$x \in \mathbb{R} \to y \in \mathbb{R} \to x - y = 0 \to y + 0 = y$}
\AXC{$\stackrel{\Pi_{3}}{\vdots}$}\noLine
\UIC{$x = 0 + y \to 0 + y = y + 0 \to x = y + 0$}
\AXC{$\stackrel{\Pi_{15}}{\vdots}$}\noLine
\UIC{$\varphi_{13}$}
\BIC{$x \in \mathbb{R} \to y \in \mathbb{R} \to x - y = 0 \to x = y + 0$}
\AXC{$\varphi_{25}$\lbl{A1}}
\BIC{$(x \in \mathbb{R} \to y \in \mathbb{R} \to x - y = 0 \to y + 0 = y) \to \varphi_{9}$}
\BIC{$\varphi_{9}$}
\end{ptree}
\noindent 主木:
\begin{ptree}
\AXC{$\stackrel{\Pi_{1}}{\vdots}$}\noLine
\UIC{$x = y + 0 \to y + 0 = y \to x = y$}
\AXC{$\stackrel{\Pi_{16}}{\vdots}$}\noLine
\UIC{$\varphi_{9}$}
\BIC{$x \in \mathbb{R} \to y \in \mathbb{R} \to x - y = 0 \to x = y$}
\UIC{$\forall y\,(x \in \mathbb{R} \to y \in \mathbb{R} \to x - y = 0 \to x = y)$}
\UIC{$\forall x\,\forall y\,(x \in \mathbb{R} \to y \in \mathbb{R} \to x - y = 0 \to x = y)$}
\end{ptree}
\begin{abbrevlist}
\item $\varphi_{1} :\equiv \forall y\,(v_{406} = y \to y = v_{406}) \to v_{406} = v_{407} \to v_{407} = v_{406}$
\item $\varphi_{2} :\equiv \forall x\,\forall y\,\forall z\,(x = y \to y = z \to x = z) \to \forall y\,\forall z\,(x = y \to y = z \to x = z)$
\item $\varphi_{3} :\equiv \forall y\,\forall z\,(x = y \to y = z \to x = z) \to \forall z\,(x = y + 0 \to y + 0 = z \to x = z)$
\item $\varphi_{4} :\equiv \forall y\,\forall z\,(x = y \to y = z \to x = z) \to \forall z\,(x = 0 + y \to 0 + y = z \to x = z)$
\item $\varphi_{5} :\equiv \forall z\,(x = y + 0 \to y + 0 = z \to x = z) \to x = y + 0 \to y + 0 = y \to x = y$
\item $\varphi_{6} :\equiv \forall w\,(x - y = w \to x - y + y = x - y + y \to x - y + y = w + y)$
\item $\varphi_{7} :\equiv \forall y\,\forall z\,(x = y \to y = z \to x = z) \to \forall z\,(x = x - y + y \to x - y + y = z \to x = z)$
\item $\varphi_{8} :\equiv \forall z\,(x = 0 + y \to 0 + y = z \to x = z) \to x = 0 + y \to 0 + y = y + 0 \to x = y + 0$
\item $\varphi_{9} :\equiv (x = y + 0 \to y + 0 = y \to x = y) \to x \in \mathbb{R} \to y \in \mathbb{R} \to x - y = 0 \to x = y$
\item $\varphi_{10} :\equiv \forall v_{407}\,(x - y + y = v_{407} \to v_{407} = x - y + y) \to x - y + y = x \to x = x - y + y$
\item $\varphi_{11} :\equiv \forall v_{406}\,\forall v_{407}\,(v_{406} = v_{407} \to v_{407} = v_{406}) \to \forall v_{407}\,(x - y + y = v_{407} \to v_{407} = x - y + y)$
\item $\varphi_{12} :\equiv (x - y + y = x \to x = x - y + y) \to x \in \mathbb{R} \to y \in \mathbb{R} \to x - y = 0 \to x = x - y + y$
\item $\varphi_{13} :\equiv (x = 0 + y \to 0 + y = y + 0 \to x = y + 0) \to x \in \mathbb{R} \to y \in \mathbb{R} \to x - y = 0 \to x = y + 0$
\item $\varphi_{14} :\equiv \forall y\,(0 \in \mathbb{R} \to y \in \mathbb{R} \to 0 + y = y + 0) \to 0 \in \mathbb{R} \to y \in \mathbb{R} \to 0 + y = y + 0$
\item $\varphi_{15} :\equiv \forall y\,(x \in \mathbb{R} \to y \in \mathbb{R} \to x - y + y = x) \to x \in \mathbb{R} \to y \in \mathbb{R} \to x - y + y = x$
\item $\varphi_{16} :\equiv \forall z\,(x = x - y + y \to x - y + y = z \to x = z) \to x = x - y + y \to x - y + y = 0 + y \to x = 0 + y$
\item $\varphi_{17} :\equiv \forall x\,\forall y\,(x \in \mathbb{R} \to y \in \mathbb{R} \to x + y = y + x) \to \forall y\,(0 \in \mathbb{R} \to y \in \mathbb{R} \to 0 + y = y + 0)$
\item $\varphi_{18} :\equiv \forall x\,\forall y\,(x \in \mathbb{R} \to y \in \mathbb{R} \to x - y + y = x) \to \forall y\,(x \in \mathbb{R} \to y \in \mathbb{R} \to x - y + y = x)$
\item $\varphi_{19} :\equiv (x = x - y + y \to x - y + y = 0 + y \to x = 0 + y) \to x \in \mathbb{R} \to y \in \mathbb{R} \to x - y = 0 \to x = 0 + y$
\item $\varphi_{20} :\equiv (x \in \mathbb{R} \to y \in \mathbb{R} \to x - y = 0 \to y \in \mathbb{R}) \to x \in \mathbb{R} \to y \in \mathbb{R} \to x - y = 0 \to y + 0 = y$
\item $\varphi_{21} :\equiv (x \in \mathbb{R} \to y \in \mathbb{R} \to x - y = 0 \to y \in \mathbb{R}) \to x \in \mathbb{R} \to y \in \mathbb{R} \to x - y = 0 \to 0 + y = y + 0$
\item $\varphi_{22} :\equiv (x \in \mathbb{R} \to y \in \mathbb{R} \to x - y = 0 \to y \in \mathbb{R}) \to x \in \mathbb{R} \to y \in \mathbb{R} \to x - y = 0 \to x - y + y = x$
\item $\varphi_{23} :\equiv \forall z\,\forall w\,(z = w \to x - y + y = z + y \to x - y + y = w + y) \to \varphi_{6}$
\item $\varphi_{24} :\equiv (x - y = 0 \to x - y + y = x - y + y \to x - y + y = 0 + y) \to x \in \mathbb{R} \to y \in \mathbb{R} \to x - y = 0 \to x - y + y = 0 + y$
\item $\varphi_{25} :\equiv (x \in \mathbb{R} \to y \in \mathbb{R} \to x - y = 0 \to x = y + 0) \to (x \in \mathbb{R} \to y \in \mathbb{R} \to x - y = 0 \to y + 0 = y) \to \varphi_{9}$
\item $\varphi_{26} :\equiv (x \in \mathbb{R} \to y \in \mathbb{R} \to x - y = 0 \to x = 0 + y) \to (x \in \mathbb{R} \to y \in \mathbb{R} \to x - y = 0 \to 0 + y = y + 0) \to \varphi_{13}$
\item $\varphi_{27} :\equiv (x \in \mathbb{R} \to y \in \mathbb{R} \to x - y = 0 \to x - y = 0) \to x - y + y = x - y + y \to \varphi_{24}$
\item $\varphi_{28} :\equiv (x \in \mathbb{R} \to y \in \mathbb{R} \to x - y + y = x) \to (x \in \mathbb{R} \to y \in \mathbb{R} \to x - y = 0 \to x \in \mathbb{R}) \to \varphi_{22}$
\item $\varphi_{29} :\equiv (x \in \mathbb{R} \to y \in \mathbb{R} \to x - y = 0 \to x = x - y + y) \to (x \in \mathbb{R} \to y \in \mathbb{R} \to x - y = 0 \to x - y + y = 0 + y) \to \varphi_{19}$
\end{abbrevlist}
\end{fullproof}

% 自動生成 (prover/build.js): MinPos  ノード数 61
\begin{fullproof}{[D8]}{$\forall p\,\forall q\,(p \in \mathbb{R} \to q \in \mathbb{R} \to 0 < p \to 0 < q \to \exists d\,(d \in \mathbb{R} \land 0 < d \land d \leq p \land d \leq q))$}{61}\phantomsection\label{tree:MinPos}
\noindent$\Pi_{1}$:
\begin{ptree}
\AXC{$p \in \mathbb{R} \to q \in \mathbb{R} \to 0 < p \to 0 < q \to 0 < q$\lbl{A1}}
\AXC{$(p \in \mathbb{R} \to q \in \mathbb{R} \to 0 < p \to 0 < q \to 0 < q) \to \varphi_{17}$\lbl{A1}}
\BIC{$\varphi_{17}$}
\end{ptree}
\noindent$\Pi_{2}$:
\begin{ptree}
\AXC{$x = x$\lbl{A5}}
\UIC{$\forall x\,(x = x)$}
\AXC{$\forall x\,(x = x) \to q = q$\lbl{A2}}
\BIC{$q = q$}
\AXC{$p \in \mathbb{R} \to q \in \mathbb{R} \to 0 < p \to 0 < q \to q < p \to q < p$\lbl{A1}}
\AXC{$\stackrel{\Pi_{1}}{\vdots}$}\noLine
\UIC{$\varphi_{17}$}
\BIC{$q = q \to \varphi_{14}$}
\BIC{$\varphi_{14}$}
\end{ptree}
\noindent$\Pi_{3}$:
\begin{ptree}
\AXC{$\varphi_{6}$\lbl{A2}}
\AXC{$p \in \mathbb{R} \to q \in \mathbb{R} \to 0 < p \to 0 < q \to q < p \to q \in \mathbb{R}$\lbl{A1}}
\AXC{$\stackrel{\Pi_{2}}{\vdots}$}\noLine
\UIC{$\varphi_{14}$}
\BIC{$p \in \mathbb{R} \to \varphi_{8}$}
\AXC{$(p \in \mathbb{R} \to \varphi_{8}) \to \varphi_{6} \to p \in \mathbb{R} \to q \in \mathbb{R} \to \varphi_{2}$\lbl{A1}}
\BIC{$\varphi_{6} \to p \in \mathbb{R} \to q \in \mathbb{R} \to \varphi_{2}$}
\BIC{$p \in \mathbb{R} \to q \in \mathbb{R} \to \varphi_{2}$}
\end{ptree}
\noindent$\Pi_{4}$:
\begin{ptree}
\AXC{$p \in \mathbb{R} \to q \in \mathbb{R} \to 0 < p \to 0 < q \to 0 < p$\lbl{A1}}
\AXC{$(p \in \mathbb{R} \to q \in \mathbb{R} \to 0 < p \to 0 < q \to 0 < p) \to \varphi_{16}$\lbl{A1}}
\BIC{$\varphi_{16}$}
\end{ptree}
\noindent$\Pi_{5}$:
\begin{ptree}
\AXC{$x = x$\lbl{A5}}
\UIC{$\forall x\,(x = x)$}
\AXC{$\forall x\,(x = x) \to p = p$\lbl{A2}}
\BIC{$p = p$}
\AXC{$p \in \mathbb{R} \to q \in \mathbb{R} \to 0 < p \to 0 < q \to p = q \to p = q$\lbl{A1}}
\AXC{$\stackrel{\Pi_{4}}{\vdots}$}\noLine
\UIC{$\varphi_{16}$}
\BIC{$p = p \to \varphi_{13}$}
\BIC{$\varphi_{13}$}
\end{ptree}
\noindent$\Pi_{6}$:
\begin{ptree}
\AXC{$\varphi_{7}$\lbl{A2}}
\AXC{$p \in \mathbb{R} \to q \in \mathbb{R} \to 0 < p \to 0 < q \to p = q \to p \in \mathbb{R}$\lbl{A1}}
\AXC{$\stackrel{\Pi_{5}}{\vdots}$}\noLine
\UIC{$\varphi_{13}$}
\BIC{$p \in \mathbb{R} \to \varphi_{5}$}
\AXC{$(p \in \mathbb{R} \to \varphi_{5}) \to \varphi_{7} \to p \in \mathbb{R} \to q \in \mathbb{R} \to \varphi_{1}$\lbl{A1}}
\BIC{$\varphi_{7} \to p \in \mathbb{R} \to q \in \mathbb{R} \to \varphi_{1}$}
\BIC{$p \in \mathbb{R} \to q \in \mathbb{R} \to \varphi_{1}$}
\end{ptree}
\noindent$\Pi_{7}$:
\begin{ptree}
\AXC{$p \in \mathbb{R} \to q \in \mathbb{R} \to 0 < p \to 0 < q \to 0 < p$\lbl{A1}}
\AXC{$(p \in \mathbb{R} \to q \in \mathbb{R} \to 0 < p \to 0 < q \to 0 < p) \to \varphi_{18}$\lbl{A1}}
\BIC{$\varphi_{18}$}
\end{ptree}
\noindent$\Pi_{8}$:
\begin{ptree}
\AXC{$x = x$\lbl{A5}}
\UIC{$\forall x\,(x = x)$}
\AXC{$\forall x\,(x = x) \to p = p$\lbl{A2}}
\BIC{$p = p$}
\AXC{$p \in \mathbb{R} \to q \in \mathbb{R} \to 0 < p \to 0 < q \to p < q \to p < q$\lbl{A1}}
\AXC{$\stackrel{\Pi_{7}}{\vdots}$}\noLine
\UIC{$\varphi_{18}$}
\BIC{$p = p \to \varphi_{15}$}
\BIC{$\varphi_{15}$}
\end{ptree}
\noindent$\Pi_{9}$:
\begin{ptree}
\AXC{$\varphi_{7}$\lbl{A2}}
\AXC{$p \in \mathbb{R} \to q \in \mathbb{R} \to 0 < p \to 0 < q \to p < q \to p \in \mathbb{R}$\lbl{A1}}
\AXC{$\stackrel{\Pi_{8}}{\vdots}$}\noLine
\UIC{$\varphi_{15}$}
\BIC{$p \in \mathbb{R} \to \varphi_{9}$}
\AXC{$(p \in \mathbb{R} \to \varphi_{9}) \to \varphi_{7} \to p \in \mathbb{R} \to q \in \mathbb{R} \to \varphi_{3}$\lbl{A1}}
\BIC{$\varphi_{7} \to p \in \mathbb{R} \to q \in \mathbb{R} \to \varphi_{3}$}
\BIC{$p \in \mathbb{R} \to q \in \mathbb{R} \to \varphi_{3}$}
\end{ptree}
\noindent$\Pi_{10}$:
\begin{ptree}
\AXC{$\forall x\,\forall y\,(x \in \mathbb{R} \to y \in \mathbb{R} \to x < y \lor x = y \lor y < x)$\lbl{(R13)}}
\AXC{$\varphi_{11}$\lbl{A2}}
\BIC{$\forall y\,(p \in \mathbb{R} \to y \in \mathbb{R} \to p < y \lor p = y \lor y < p)$}
\AXC{$\varphi_{10}$\lbl{A2}}
\BIC{$p \in \mathbb{R} \to q \in \mathbb{R} \to p < q \lor p = q \lor q < p$}
\AXC{$\varphi_{19}$\lbl{A1}}
\BIC{$(p \in \mathbb{R} \to q \in \mathbb{R} \to 0 < p \to 0 < q \to p \in \mathbb{R}) \to \varphi_{12}$}
\end{ptree}
\noindent$\Pi_{11}$:
\begin{ptree}
\AXC{$p \in \mathbb{R} \to q \in \mathbb{R} \to 0 < p \to 0 < q \to p \in \mathbb{R}$\lbl{A1}}
\AXC{$\stackrel{\Pi_{10}}{\vdots}$}\noLine
\UIC{$(p \in \mathbb{R} \to q \in \mathbb{R} \to 0 < p \to 0 < q \to p \in \mathbb{R}) \to \varphi_{12}$}
\BIC{$\varphi_{12}$}
\end{ptree}
\noindent$\Pi_{12}$:
\begin{ptree}
\AXC{$\stackrel{\Pi_{6}}{\vdots}$}\noLine
\UIC{$p \in \mathbb{R} \to q \in \mathbb{R} \to \varphi_{1}$}
\AXC{$\stackrel{\Pi_{9}}{\vdots}$}\noLine
\UIC{$p \in \mathbb{R} \to q \in \mathbb{R} \to \varphi_{3}$}
\AXC{$p \in \mathbb{R} \to q \in \mathbb{R} \to 0 < p \to 0 < q \to q \in \mathbb{R}$\lbl{A1}}
\AXC{$\stackrel{\Pi_{11}}{\vdots}$}\noLine
\UIC{$\varphi_{12}$}
\BIC{$p \in \mathbb{R} \to q \in \mathbb{R} \to 0 < p \to 0 < q \to p < q \lor p = q \lor q < p$}
\AXC{$\varphi_{21}$\lbl{A1}}
\BIC{$(p \in \mathbb{R} \to q \in \mathbb{R} \to \varphi_{3}) \to \varphi_{20}$}
\BIC{$\varphi_{20}$}
\BIC{$(p \in \mathbb{R} \to q \in \mathbb{R} \to \varphi_{2}) \to p \in \mathbb{R} \to \varphi_{4}$}
\end{ptree}
\noindent 主木:
\begin{ptree}
\AXC{$\stackrel{\Pi_{3}}{\vdots}$}\noLine
\UIC{$p \in \mathbb{R} \to q \in \mathbb{R} \to \varphi_{2}$}
\AXC{$\stackrel{\Pi_{12}}{\vdots}$}\noLine
\UIC{$(p \in \mathbb{R} \to q \in \mathbb{R} \to \varphi_{2}) \to p \in \mathbb{R} \to \varphi_{4}$}
\BIC{$p \in \mathbb{R} \to \varphi_{4}$}
\UIC{$\forall q\,(p \in \mathbb{R} \to \varphi_{4})$}
\UIC{$\forall p\,\forall q\,(p \in \mathbb{R} \to \varphi_{4})$}
\end{ptree}
\begin{abbrevlist}
\item $\varphi_{1} :\equiv 0 < p \to 0 < q \to p = q \to \exists d\,(d \in \mathbb{R} \land 0 < d \land d \leq p \land d \leq q)$
\item $\varphi_{2} :\equiv 0 < p \to 0 < q \to q < p \to \exists d\,(d \in \mathbb{R} \land 0 < d \land d \leq p \land d \leq q)$
\item $\varphi_{3} :\equiv 0 < p \to 0 < q \to p < q \to \exists d\,(d \in \mathbb{R} \land 0 < d \land d \leq p \land d \leq q)$
\item $\varphi_{4} :\equiv q \in \mathbb{R} \to 0 < p \to 0 < q \to \exists d\,(d \in \mathbb{R} \land 0 < d \land d \leq p \land d \leq q)$
\item $\varphi_{5} :\equiv q \in \mathbb{R} \to 0 < p \to 0 < q \to p = q \to p \in \mathbb{R} \land 0 < p \land p \leq p \land p \leq q$
\item $\varphi_{6} :\equiv \forall d\,(\neg (d \in \mathbb{R} \land 0 < d \land d \leq p \land d \leq q)) \to \neg (q \in \mathbb{R} \land 0 < q \land q \leq p \land q \leq q)$
\item $\varphi_{7} :\equiv \forall d\,(\neg (d \in \mathbb{R} \land 0 < d \land d \leq p \land d \leq q)) \to \neg (p \in \mathbb{R} \land 0 < p \land p \leq p \land p \leq q)$
\item $\varphi_{8} :\equiv q \in \mathbb{R} \to 0 < p \to 0 < q \to q < p \to q \in \mathbb{R} \land 0 < q \land q \leq p \land q \leq q$
\item $\varphi_{9} :\equiv q \in \mathbb{R} \to 0 < p \to 0 < q \to p < q \to p \in \mathbb{R} \land 0 < p \land p \leq p \land p \leq q$
\item $\varphi_{10} :\equiv \forall y\,(p \in \mathbb{R} \to y \in \mathbb{R} \to p < y \lor p = y \lor y < p) \to p \in \mathbb{R} \to q \in \mathbb{R} \to p < q \lor p = q \lor q < p$
\item $\varphi_{11} :\equiv \forall x\,\forall y\,(x \in \mathbb{R} \to y \in \mathbb{R} \to x < y \lor x = y \lor y < x) \to \forall y\,(p \in \mathbb{R} \to y \in \mathbb{R} \to p < y \lor p = y \lor y < p)$
\item $\varphi_{12} :\equiv (p \in \mathbb{R} \to q \in \mathbb{R} \to 0 < p \to 0 < q \to q \in \mathbb{R}) \to p \in \mathbb{R} \to q \in \mathbb{R} \to 0 < p \to 0 < q \to p < q \lor p = q \lor q < p$
\item $\varphi_{13} :\equiv (p \in \mathbb{R} \to q \in \mathbb{R} \to 0 < p \to 0 < q \to p = q \to p \in \mathbb{R}) \to p \in \mathbb{R} \to \varphi_{5}$
\item $\varphi_{14} :\equiv (p \in \mathbb{R} \to q \in \mathbb{R} \to 0 < p \to 0 < q \to q < p \to q \in \mathbb{R}) \to p \in \mathbb{R} \to \varphi_{8}$
\item $\varphi_{15} :\equiv (p \in \mathbb{R} \to q \in \mathbb{R} \to 0 < p \to 0 < q \to p < q \to p \in \mathbb{R}) \to p \in \mathbb{R} \to \varphi_{9}$
\item $\varphi_{16} :\equiv (p \in \mathbb{R} \to q \in \mathbb{R} \to 0 < p \to 0 < q \to p = q \to p = q) \to p = p \to \varphi_{13}$
\item $\varphi_{17} :\equiv (p \in \mathbb{R} \to q \in \mathbb{R} \to 0 < p \to 0 < q \to q < p \to q < p) \to q = q \to \varphi_{14}$
\item $\varphi_{18} :\equiv (p \in \mathbb{R} \to q \in \mathbb{R} \to 0 < p \to 0 < q \to p < q \to p < q) \to p = p \to \varphi_{15}$
\item $\varphi_{19} :\equiv (p \in \mathbb{R} \to q \in \mathbb{R} \to p < q \lor p = q \lor q < p) \to (p \in \mathbb{R} \to q \in \mathbb{R} \to 0 < p \to 0 < q \to p \in \mathbb{R}) \to \varphi_{12}$
\item $\varphi_{20} :\equiv (p \in \mathbb{R} \to q \in \mathbb{R} \to \varphi_{1}) \to (p \in \mathbb{R} \to q \in \mathbb{R} \to \varphi_{2}) \to p \in \mathbb{R} \to \varphi_{4}$
\item $\varphi_{21} :\equiv (p \in \mathbb{R} \to q \in \mathbb{R} \to 0 < p \to 0 < q \to p < q \lor p = q \lor q < p) \to (p \in \mathbb{R} \to q \in \mathbb{R} \to \varphi_{3}) \to \varphi_{20}$
\end{abbrevlist}
\end{fullproof}

% 自動生成 (prover/build.js): SeqCrit  ノード数 261
\begin{fullproof}{定理1.7(⇒)}{$\forall f\,\forall c\,\forall \alpha\,\forall a\,(\mathrm{FLim}(f,c,\alpha) \to \mathrm{Conv}(a,c) \to \forall n \in \mathbb{N}\,(a_{n} \in \mathrm{dom}(f) \land a_{n} \neq c) \to \mathrm{Conv}(f \circ a,\alpha))$}{261}\phantomsection\label{tree:SeqCrit}
\noindent$\Pi_{1}$:
\begin{ptree}
\AXC{$z = w \to \lvert z - \alpha \rvert < \varepsilon \to \lvert w - \alpha \rvert < \varepsilon$\lbl{A6}}
\UIC{$\forall w\,(z = w \to \lvert z - \alpha \rvert < \varepsilon \to \lvert w - \alpha \rvert < \varepsilon)$}
\UIC{$\forall z\,\forall w\,(z = w \to \lvert z - \alpha \rvert < \varepsilon \to \lvert w - \alpha \rvert < \varepsilon)$}
\AXC{$\forall z\,\forall w\,(z = w \to \lvert z - \alpha \rvert < \varepsilon \to \lvert w - \alpha \rvert < \varepsilon) \to \varphi_{23}$\lbl{A2}}
\BIC{$\varphi_{23}$}
\AXC{$\varphi_{23} \to \varphi_{55}$\lbl{A2}}
\BIC{$\varphi_{55}$}
\end{ptree}
\noindent$\Pi_{2}$:
\begin{ptree}
\AXC{$\varphi_{39}$\lbl{A2}}
\AXC{$\mathrm{FLim}(f,c,\alpha) \to \mathrm{Conv}(a,c) \to \varphi_{41}$\lbl{A1}}
\AXC{$\varphi_{122}$\lbl{A1}}
\BIC{$\varphi_{39} \to \mathrm{FLim}(f,c,\alpha) \to \mathrm{Conv}(a,c) \to \varphi_{39}$}
\BIC{$\mathrm{FLim}(f,c,\alpha) \to \mathrm{Conv}(a,c) \to \varphi_{39}$}
\AXC{$(\mathrm{FLim}(f,c,\alpha) \to \mathrm{Conv}(a,c) \to \varphi_{39}) \to \varphi_{140}$\lbl{A1}}
\BIC{$\varphi_{140}$}
\end{ptree}
\noindent$\Pi_{3}$:
\begin{ptree}
\AXC{$\mathrm{FLim}(f,c,\alpha) \to \mathrm{Conv}(a,c) \to \varphi_{96}$\lbl{A1}}
\AXC{$\stackrel{\Pi_{2}}{\vdots}$}\noLine
\UIC{$\varphi_{140}$}
\BIC{$\mathrm{FLim}(f,c,\alpha) \to \mathrm{Conv}(a,c) \to \varphi_{112}$}
\end{ptree}
\noindent$\Pi_{4}$:
\begin{ptree}
\AXC{$\forall f\,\forall c\,\forall \alpha\,(\mathrm{FLim}(f,c,\alpha) \to c \in \mathbb{R})$\lbl{[D3]}}
\AXC{$\varphi_{15}$\lbl{A2}}
\BIC{$\forall c\,\forall \alpha\,(\mathrm{FLim}(f,c,\alpha) \to c \in \mathbb{R})$}
\AXC{$\varphi_{10}$\lbl{A2}}
\BIC{$\forall \alpha\,(\mathrm{FLim}(f,c,\alpha) \to c \in \mathbb{R})$}
\AXC{$\varphi_{6}$\lbl{A2}}
\BIC{$\mathrm{FLim}(f,c,\alpha) \to c \in \mathbb{R}$}
\end{ptree}
\noindent$\Pi_{5}$:
\begin{ptree}
\AXC{$\mathrm{FLim}(f,c,\alpha) \to \mathrm{FLim}(f,c,\alpha)$\lbl{A1}}
\AXC{$\stackrel{\Pi_{4}}{\vdots}$}\noLine
\UIC{$\mathrm{FLim}(f,c,\alpha) \to c \in \mathbb{R}$}
\AXC{$(\mathrm{FLim}(f,c,\alpha) \to c \in \mathbb{R}) \to \varphi_{65}$\lbl{A1}}
\BIC{$\varphi_{65}$}
\BIC{$\mathrm{FLim}(f,c,\alpha) \to \varphi_{20}$}
\end{ptree}
\noindent$\Pi_{6}$:
\begin{ptree}
\AXC{$\varphi_{1}$\lbl{A1}}
\AXC{$\forall a\,\forall \alpha\,(\mathrm{Conv}(a,\alpha) \to \mathrm{Seq}(a))$\lbl{[C1]}}
\AXC{$\varphi_{4}$\lbl{A2}}
\BIC{$\forall \alpha\,(\mathrm{Conv}(a,\alpha) \to \mathrm{Seq}(a))$}
\AXC{$\varphi_{2}$\lbl{A2}}
\BIC{$\mathrm{Conv}(a,c) \to \mathrm{Seq}(a)$}
\AXC{$\varphi_{74}$\lbl{A1}}
\BIC{$\varphi_{1} \to \mathrm{FLim}(f,c,\alpha) \to \varphi_{13}$}
\BIC{$\mathrm{FLim}(f,c,\alpha) \to \varphi_{13}$}
\end{ptree}
\noindent$\Pi_{7}$:
\begin{ptree}
\AXC{$\stackrel{\Pi_{6}}{\vdots}$}\noLine
\UIC{$\mathrm{FLim}(f,c,\alpha) \to \varphi_{13}$}
\AXC{$\forall a\,\forall n\,(\mathrm{Seq}(a) \to n \in \mathbb{N} \to a_{n} \in \mathbb{R})$\lbl{[S1]}}
\AXC{$\varphi_{34}$\lbl{A2}}
\BIC{$\forall n\,(\mathrm{Seq}(a) \to n \in \mathbb{N} \to a_{n} \in \mathbb{R})$}
\AXC{$\varphi_{27}$\lbl{A2}}
\BIC{$\mathrm{Seq}(a) \to n \in \mathbb{N} \to a_{n} \in \mathbb{R}$}
\AXC{$\varphi_{152}$\lbl{A1}}
\BIC{$(\mathrm{FLim}(f,c,\alpha) \to \varphi_{13}) \to \varphi_{136}$}
\BIC{$\varphi_{136}$}
\end{ptree}
\noindent$\Pi_{8}$:
\begin{ptree}
\AXC{$\mathrm{FLim}(f,c,\alpha) \to \mathrm{Conv}(a,c) \to \varphi_{96}$\lbl{A1}}
\AXC{$\stackrel{\Pi_{7}}{\vdots}$}\noLine
\UIC{$\varphi_{136}$}
\BIC{$\mathrm{FLim}(f,c,\alpha) \to \mathrm{Conv}(a,c) \to \varphi_{103}$}
\end{ptree}
\noindent$\Pi_{9}$:
\begin{ptree}
\AXC{$\forall x\,\forall y\,(x \in \mathbb{R} \to y \in \mathbb{R} \to x \neq y \to x - y \neq 0)$\lbl{[A15]}}
\AXC{$\varphi_{63}$\lbl{A2}}
\BIC{$\forall y\,(a_{n} \in \mathbb{R} \to y \in \mathbb{R} \to a_{n} \neq y \to a_{n} - y \neq 0)$}
\AXC{$\varphi_{71}$\lbl{A2}}
\BIC{$a_{n} \in \mathbb{R} \to c \in \mathbb{R} \to a_{n} \neq c \to a_{n} - c \neq 0$}
\end{ptree}
\noindent$\Pi_{10}$:
\begin{ptree}
\AXC{$\stackrel{\Pi_{9}}{\vdots}$}\noLine
\UIC{$a_{n} \in \mathbb{R} \to c \in \mathbb{R} \to a_{n} \neq c \to a_{n} - c \neq 0$}
\AXC{$(a_{n} \in \mathbb{R} \to c \in \mathbb{R} \to a_{n} \neq c \to a_{n} - c \neq 0) \to \varphi_{160}$\lbl{A1}}
\BIC{$\varphi_{160}$}
\end{ptree}
\noindent$\Pi_{11}$:
\begin{ptree}
\AXC{$\stackrel{\Pi_{5}}{\vdots}$}\noLine
\UIC{$\mathrm{FLim}(f,c,\alpha) \to \varphi_{20}$}
\AXC{$\stackrel{\Pi_{8}}{\vdots}$}\noLine
\UIC{$\mathrm{FLim}(f,c,\alpha) \to \mathrm{Conv}(a,c) \to \varphi_{103}$}
\AXC{$\stackrel{\Pi_{10}}{\vdots}$}\noLine
\UIC{$\varphi_{160}$}
\BIC{$(\mathrm{FLim}(f,c,\alpha) \to \varphi_{20}) \to \varphi_{137}$}
\BIC{$\varphi_{137}$}
\end{ptree}
\noindent$\Pi_{12}$:
\begin{ptree}
\AXC{$\forall x\,\forall y\,(x \in \mathbb{R} \to y \in \mathbb{R} \to x - y \in \mathbb{R})$\lbl{[F1]}}
\AXC{$\varphi_{51}$\lbl{A2}}
\BIC{$\forall y\,(a_{n} \in \mathbb{R} \to y \in \mathbb{R} \to a_{n} - y \in \mathbb{R})$}
\AXC{$\varphi_{62}$\lbl{A2}}
\BIC{$a_{n} \in \mathbb{R} \to c \in \mathbb{R} \to a_{n} - c \in \mathbb{R}$}
\AXC{$\varphi_{154}$\lbl{A1}}
\BIC{$(\mathrm{FLim}(f,c,\alpha) \to \mathrm{Conv}(a,c) \to \varphi_{103}) \to \varphi_{126}$}
\end{ptree}
\noindent$\Pi_{13}$:
\begin{ptree}
\AXC{$\stackrel{\Pi_{5}}{\vdots}$}\noLine
\UIC{$\mathrm{FLim}(f,c,\alpha) \to \varphi_{20}$}
\AXC{$\stackrel{\Pi_{8}}{\vdots}$}\noLine
\UIC{$\mathrm{FLim}(f,c,\alpha) \to \mathrm{Conv}(a,c) \to \varphi_{103}$}
\AXC{$\stackrel{\Pi_{12}}{\vdots}$}\noLine
\UIC{$(\mathrm{FLim}(f,c,\alpha) \to \mathrm{Conv}(a,c) \to \varphi_{103}) \to \varphi_{126}$}
\BIC{$\varphi_{126}$}
\BIC{$\mathrm{FLim}(f,c,\alpha) \to \mathrm{Conv}(a,c) \to \varphi_{105}$}
\end{ptree}
\noindent$\Pi_{14}$:
\begin{ptree}
\AXC{$\stackrel{\Pi_{13}}{\vdots}$}\noLine
\UIC{$\mathrm{FLim}(f,c,\alpha) \to \mathrm{Conv}(a,c) \to \varphi_{105}$}
\AXC{$\forall x \in \mathbb{R}\,(x \neq 0 \to 0 < \lvert x \rvert)$\lbl{[A14]}}
\AXC{$\varphi_{47}$\lbl{A2}}
\BIC{$a_{n} - c \in \mathbb{R} \to a_{n} - c \neq 0 \to 0 < \lvert a_{n} - c \rvert$}
\AXC{$\varphi_{159}$\lbl{A1}}
\BIC{$(\mathrm{FLim}(f,c,\alpha) \to \mathrm{Conv}(a,c) \to \varphi_{105}) \to \varphi_{142}$}
\BIC{$\varphi_{142}$}
\end{ptree}
\noindent$\Pi_{15}$:
\begin{ptree}
\AXC{$\stackrel{\Pi_{3}}{\vdots}$}\noLine
\UIC{$\mathrm{FLim}(f,c,\alpha) \to \mathrm{Conv}(a,c) \to \varphi_{112}$}
\AXC{$\varphi_{143}$\lbl{A1}}
\BIC{$\mathrm{FLim}(f,c,\alpha) \to \mathrm{Conv}(a,c) \to \varphi_{100}$}
\AXC{$\stackrel{\Pi_{11}}{\vdots}$}\noLine
\UIC{$\varphi_{137}$}
\BIC{$\mathrm{FLim}(f,c,\alpha) \to \mathrm{Conv}(a,c) \to \varphi_{102}$}
\AXC{$\stackrel{\Pi_{14}}{\vdots}$}\noLine
\UIC{$\varphi_{142}$}
\BIC{$\mathrm{FLim}(f,c,\alpha) \to \mathrm{Conv}(a,c) \to \varphi_{107}$}
\end{ptree}
\noindent$\Pi_{16}$:
\begin{ptree}
\AXC{$\varphi_{29} \to \varphi_{52}$\lbl{A2}}
\AXC{$\mathrm{FLim}(f,c,\alpha) \to \mathrm{Conv}(a,c) \to \varphi_{110}$\lbl{A1}}
\AXC{$\varphi_{139}$\lbl{A1}}
\BIC{$\mathrm{FLim}(f,c,\alpha) \to \mathrm{Conv}(a,c) \to \varphi_{97}$}
\AXC{$\varphi_{155}$\lbl{A1}}
\BIC{$(\varphi_{29} \to \varphi_{52}) \to \mathrm{FLim}(f,c,\alpha) \to \mathrm{Conv}(a,c) \to \varphi_{111}$}
\BIC{$\mathrm{FLim}(f,c,\alpha) \to \mathrm{Conv}(a,c) \to \varphi_{111}$}
\end{ptree}
\noindent$\Pi_{17}$:
\begin{ptree}
\AXC{$\stackrel{\Pi_{16}}{\vdots}$}\noLine
\UIC{$\mathrm{FLim}(f,c,\alpha) \to \mathrm{Conv}(a,c) \to \varphi_{111}$}
\AXC{$(\mathrm{FLim}(f,c,\alpha) \to \mathrm{Conv}(a,c) \to \varphi_{111}) \to \varphi_{161}$\lbl{A1}}
\BIC{$\varphi_{161}$}
\end{ptree}
\noindent$\Pi_{18}$:
\begin{ptree}
\AXC{$\stackrel{\Pi_{15}}{\vdots}$}\noLine
\UIC{$\mathrm{FLim}(f,c,\alpha) \to \mathrm{Conv}(a,c) \to \varphi_{107}$}
\AXC{$\stackrel{\Pi_{3}}{\vdots}$}\noLine
\UIC{$\mathrm{FLim}(f,c,\alpha) \to \mathrm{Conv}(a,c) \to \varphi_{112}$}
\AXC{$\stackrel{\Pi_{17}}{\vdots}$}\noLine
\UIC{$\varphi_{161}$}
\BIC{$(\mathrm{FLim}(f,c,\alpha) \to \mathrm{Conv}(a,c) \to \varphi_{107}) \to \varphi_{144}$}
\BIC{$\varphi_{144}$}
\end{ptree}
\noindent$\Pi_{19}$:
\begin{ptree}
\AXC{$x = x$\lbl{A5}}
\AXC{$x = y \to x = x \to y = x$\lbl{A6}}
\AXC{$(x = y \to x = x \to y = x) \to x = x \to x = y \to y = x$\lbl{A1}}
\BIC{$x = x \to x = y \to y = x$}
\BIC{$x = y \to y = x$}
\UIC{$\forall y\,(x = y \to y = x)$}
\UIC{$\forall x\,\forall y\,(x = y \to y = x)$}
\AXC{$\varphi_{11}$\lbl{A2}}
\BIC{$\forall y\,(f \circ a(n) = y \to y = f \circ a(n))$}
\AXC{$\varphi_{69}$\lbl{A2}}
\BIC{$f \circ a(n) = f(a_{n}) \to f(a_{n}) = f \circ a(n)$}
\end{ptree}
\noindent$\Pi_{20}$:
\begin{ptree}
\AXC{$\varphi_{39}$\lbl{A2}}
\AXC{$\mathrm{FLim}(f,c,\alpha) \to \mathrm{Conv}(a,c) \to \varphi_{41}$\lbl{A1}}
\AXC{$\varphi_{122}$\lbl{A1}}
\BIC{$\varphi_{39} \to \mathrm{FLim}(f,c,\alpha) \to \mathrm{Conv}(a,c) \to \varphi_{39}$}
\BIC{$\mathrm{FLim}(f,c,\alpha) \to \mathrm{Conv}(a,c) \to \varphi_{39}$}
\AXC{$(\mathrm{FLim}(f,c,\alpha) \to \mathrm{Conv}(a,c) \to \varphi_{39}) \to \varphi_{86}$\lbl{A1}}
\BIC{$\varphi_{86}$}
\end{ptree}
\noindent$\Pi_{21}$:
\begin{ptree}
\AXC{$\mathrm{FLim}(f,c,\alpha) \to \mathrm{Conv}(a,c) \to \varphi_{18}$\lbl{A1}}
\AXC{$\stackrel{\Pi_{20}}{\vdots}$}\noLine
\UIC{$\varphi_{86}$}
\BIC{$\mathrm{FLim}(f,c,\alpha) \to \mathrm{Conv}(a,c) \to \varphi_{28}$}
\AXC{$\varphi_{87}$\lbl{A1}}
\BIC{$\varphi_{22} \to n \in \mathbb{N} \to a_{n} \in \mathrm{dom}(f)$}
\UIC{$\forall n\,(\varphi_{22} \to n \in \mathbb{N} \to a_{n} \in \mathrm{dom}(f))$}
\AXC{$\forall n\,(\varphi_{22} \to n \in \mathbb{N} \to a_{n} \in \mathrm{dom}(f)) \to \varphi_{88} \to \varphi_{98}$\lbl{A1}}
\BIC{$\varphi_{88} \to \varphi_{98}$}
\end{ptree}
\noindent$\Pi_{22}$:
\begin{ptree}
\AXC{$\forall f\,\forall c\,\forall \alpha\,(\mathrm{FLim}(f,c,\alpha) \to \mathrm{RFn}(f))$\lbl{[D2]}}
\AXC{$\varphi_{9}$\lbl{A2}}
\BIC{$\forall c\,\forall \alpha\,(\mathrm{FLim}(f,c,\alpha) \to \mathrm{RFn}(f))$}
\AXC{$\varphi_{5}$\lbl{A2}}
\BIC{$\forall \alpha\,(\mathrm{FLim}(f,c,\alpha) \to \mathrm{RFn}(f))$}
\AXC{$\varphi_{3}$\lbl{A2}}
\BIC{$\mathrm{FLim}(f,c,\alpha) \to \mathrm{RFn}(f)$}
\end{ptree}
\noindent$\Pi_{23}$:
\begin{ptree}
\AXC{$\stackrel{\Pi_{22}}{\vdots}$}\noLine
\UIC{$\mathrm{FLim}(f,c,\alpha) \to \mathrm{RFn}(f)$}
\AXC{$(\mathrm{FLim}(f,c,\alpha) \to \mathrm{RFn}(f)) \to \varphi_{60}$\lbl{A1}}
\BIC{$\varphi_{60}$}
\end{ptree}
\noindent$\Pi_{24}$:
\begin{ptree}
\AXC{$\forall f\,\forall a\,(\mathrm{RFn}(f) \to \mathrm{Seq}(a) \to \varphi_{50})$\lbl{[S8]}}
\AXC{$\varphi_{94}$\lbl{A2}}
\BIC{$\forall a\,(\mathrm{RFn}(f) \to \mathrm{Seq}(a) \to \varphi_{50})$}
\AXC{$\varphi_{91}$\lbl{A2}}
\BIC{$\mathrm{RFn}(f) \to \mathrm{Seq}(a) \to \varphi_{50}$}
\AXC{$(\mathrm{RFn}(f) \to \mathrm{Seq}(a) \to \varphi_{50}) \to \varphi_{127}$\lbl{A1}}
\BIC{$\varphi_{127}$}
\end{ptree}
\noindent$\Pi_{25}$:
\begin{ptree}
\AXC{$\stackrel{\Pi_{6}}{\vdots}$}\noLine
\UIC{$\mathrm{FLim}(f,c,\alpha) \to \varphi_{13}$}
\AXC{$\mathrm{FLim}(f,c,\alpha) \to \mathrm{FLim}(f,c,\alpha)$\lbl{A1}}
\AXC{$\stackrel{\Pi_{23}}{\vdots}$}\noLine
\UIC{$\varphi_{60}$}
\BIC{$\mathrm{FLim}(f,c,\alpha) \to \varphi_{14}$}
\AXC{$\stackrel{\Pi_{24}}{\vdots}$}\noLine
\UIC{$\varphi_{127}$}
\BIC{$(\mathrm{FLim}(f,c,\alpha) \to \varphi_{13}) \to \varphi_{95}$}
\BIC{$\varphi_{95}$}
\end{ptree}
\noindent$\Pi_{26}$:
\begin{ptree}
\AXC{$\varphi_{22} \to \forall n\,(\varphi_{22})$\lbl{A4}}
\AXC{$\varphi_{88}$\lbl{A3}}
\AXC{$\stackrel{\Pi_{21}}{\vdots}$}\noLine
\UIC{$\varphi_{88} \to \varphi_{98}$}
\BIC{$\varphi_{98}$}
\BIC{$\mathrm{FLim}(f,c,\alpha) \to \mathrm{Conv}(a,c) \to \varphi_{32}$}
\AXC{$\stackrel{\Pi_{25}}{\vdots}$}\noLine
\UIC{$\varphi_{95}$}
\BIC{$\mathrm{FLim}(f,c,\alpha) \to \mathrm{Conv}(a,c) \to \varphi_{68}$}
\end{ptree}
\noindent$\Pi_{27}$:
\begin{ptree}
\AXC{$\varphi_{42}$\lbl{A2}}
\AXC{$\stackrel{\Pi_{26}}{\vdots}$}\noLine
\UIC{$\mathrm{FLim}(f,c,\alpha) \to \mathrm{Conv}(a,c) \to \varphi_{68}$}
\AXC{$\varphi_{101}$\lbl{A1}}
\BIC{$\mathrm{FLim}(f,c,\alpha) \to \mathrm{Conv}(a,c) \to \varphi_{45}$}
\AXC{$\varphi_{124}$\lbl{A1}}
\BIC{$\varphi_{42} \to \mathrm{FLim}(f,c,\alpha) \to \mathrm{Conv}(a,c) \to \varphi_{40}$}
\BIC{$\mathrm{FLim}(f,c,\alpha) \to \mathrm{Conv}(a,c) \to \varphi_{40}$}
\end{ptree}
\noindent$\Pi_{28}$:
\begin{ptree}
\AXC{$\stackrel{\Pi_{27}}{\vdots}$}\noLine
\UIC{$\mathrm{FLim}(f,c,\alpha) \to \mathrm{Conv}(a,c) \to \varphi_{40}$}
\AXC{$(\mathrm{FLim}(f,c,\alpha) \to \mathrm{Conv}(a,c) \to \varphi_{40}) \to \varphi_{141}$\lbl{A1}}
\BIC{$\varphi_{141}$}
\end{ptree}
\noindent$\Pi_{29}$:
\begin{ptree}
\AXC{$\mathrm{FLim}(f,c,\alpha) \to \mathrm{Conv}(a,c) \to \varphi_{96}$\lbl{A1}}
\AXC{$\stackrel{\Pi_{28}}{\vdots}$}\noLine
\UIC{$\varphi_{141}$}
\BIC{$\mathrm{FLim}(f,c,\alpha) \to \mathrm{Conv}(a,c) \to \varphi_{114}$}
\AXC{$(\mathrm{FLim}(f,c,\alpha) \to \mathrm{Conv}(a,c) \to \varphi_{114}) \to \varphi_{125}$\lbl{A1}}
\BIC{$\varphi_{125}$}
\end{ptree}
\noindent$\Pi_{30}$:
\begin{ptree}
\AXC{$\stackrel{\Pi_{19}}{\vdots}$}\noLine
\UIC{$f \circ a(n) = f(a_{n}) \to f(a_{n}) = f \circ a(n)$}
\AXC{$\stackrel{\Pi_{29}}{\vdots}$}\noLine
\UIC{$\varphi_{125}$}
\BIC{$\mathrm{FLim}(f,c,\alpha) \to \mathrm{Conv}(a,c) \to \varphi_{115}$}
\AXC{$(\mathrm{FLim}(f,c,\alpha) \to \mathrm{Conv}(a,c) \to \varphi_{115}) \to \varphi_{153}$\lbl{A1}}
\BIC{$\varphi_{153}$}
\end{ptree}
\noindent$\Pi_{31}$:
\begin{ptree}
\AXC{$\stackrel{\Pi_{1}}{\vdots}$}\noLine
\UIC{$\varphi_{55}$}
\AXC{$\varphi_{36}$\lbl{A2}}
\AXC{$\mathrm{FLim}(f,c,\alpha) \to \mathrm{Conv}(a,c) \to \varphi_{117}$\lbl{A1}}
\AXC{$\varphi_{145}$\lbl{A1}}
\BIC{$\mathrm{FLim}(f,c,\alpha) \to \mathrm{Conv}(a,c) \to \varphi_{109}$}
\AXC{$\varphi_{151}$\lbl{A1}}
\BIC{$\varphi_{36} \to \mathrm{FLim}(f,c,\alpha) \to \mathrm{Conv}(a,c) \to \varphi_{108}$}
\BIC{$\mathrm{FLim}(f,c,\alpha) \to \mathrm{Conv}(a,c) \to \varphi_{108}$}
\AXC{$\stackrel{\Pi_{18}}{\vdots}$}\noLine
\UIC{$\varphi_{144}$}
\BIC{$\mathrm{FLim}(f,c,\alpha) \to \mathrm{Conv}(a,c) \to \varphi_{113}$}
\AXC{$\stackrel{\Pi_{30}}{\vdots}$}\noLine
\UIC{$\varphi_{153}$}
\BIC{$\varphi_{55} \to \mathrm{FLim}(f,c,\alpha) \to \mathrm{Conv}(a,c) \to \varphi_{116}$}
\BIC{$\mathrm{FLim}(f,c,\alpha) \to \mathrm{Conv}(a,c) \to \varphi_{116}$}
\end{ptree}
\noindent$\Pi_{32}$:
\begin{ptree}
\AXC{$\mathrm{FLim}(f,c,\alpha) \to \mathrm{Conv}(a,c) \to \varphi_{117}$\lbl{A1}}
\AXC{$\varphi_{138}$\lbl{A1}}
\BIC{$\mathrm{FLim}(f,c,\alpha) \to \mathrm{Conv}(a,c) \to \varphi_{92}$}
\AXC{$(\mathrm{FLim}(f,c,\alpha) \to \mathrm{Conv}(a,c) \to \varphi_{92}) \to \varphi_{148}$\lbl{A1}}
\BIC{$\varphi_{148}$}
\end{ptree}
\noindent$\Pi_{33}$:
\begin{ptree}
\AXC{$\varphi_{99} \to \forall n\,(\varphi_{99})$\lbl{A4}}
\AXC{$\varphi_{147}$\lbl{A3}}
\AXC{$\stackrel{\Pi_{31}}{\vdots}$}\noLine
\UIC{$\mathrm{FLim}(f,c,\alpha) \to \mathrm{Conv}(a,c) \to \varphi_{116}$}
\AXC{$\varphi_{146}$\lbl{A1}}
\BIC{$\varphi_{99} \to n \in \mathbb{N} \to N \leq n \to \lvert f \circ a(n) - \alpha \rvert < \varepsilon$}
\UIC{$\forall n\,(\varphi_{99} \to n \in \mathbb{N} \to N \leq n \to \lvert f \circ a(n) - \alpha \rvert < \varepsilon)$}
\AXC{$\varphi_{163}$\lbl{A1}}
\BIC{$\varphi_{147} \to \varphi_{158}$}
\BIC{$\varphi_{158}$}
\BIC{$\mathrm{FLim}(f,c,\alpha) \to \mathrm{Conv}(a,c) \to \varphi_{118}$}
\AXC{$\stackrel{\Pi_{32}}{\vdots}$}\noLine
\UIC{$\varphi_{148}$}
\BIC{$\mathrm{FLim}(f,c,\alpha) \to \mathrm{Conv}(a,c) \to \varphi_{119}$}
\end{ptree}
\noindent$\Pi_{34}$:
\begin{ptree}
\AXC{$\forall N\,(\varphi_{16}) \to \varphi_{16}$\lbl{A2}}
\AXC{$\stackrel{\Pi_{33}}{\vdots}$}\noLine
\UIC{$\mathrm{FLim}(f,c,\alpha) \to \mathrm{Conv}(a,c) \to \varphi_{119}$}
\AXC{$(\mathrm{FLim}(f,c,\alpha) \to \mathrm{Conv}(a,c) \to \varphi_{119}) \to \varphi_{130}$\lbl{A1}}
\BIC{$\varphi_{130}$}
\BIC{$\mathrm{FLim}(f,c,\alpha) \to \mathrm{Conv}(a,c) \to \varphi_{120}$}
\end{ptree}
\noindent$\Pi_{35}$:
\begin{ptree}
\AXC{$\varphi_{1}$\lbl{A1}}
\AXC{$\forall a\,(\varphi_{54})$\lbl{[C3]}}
\AXC{$\forall a\,(\varphi_{54}) \to \varphi_{54}$\lbl{A2}}
\BIC{$\varphi_{54}$}
\AXC{$\varphi_{54} \to \varphi_{48}$\lbl{A2}}
\BIC{$\varphi_{48}$}
\AXC{$\varphi_{48} \to \mathrm{Conv}(a,c) \to \varphi_{30}$\lbl{A2}}
\BIC{$\mathrm{Conv}(a,c) \to \varphi_{30}$}
\AXC{$\varphi_{156}$\lbl{A1}}
\BIC{$\varphi_{1} \to (\mathrm{FLim}(f,c,\alpha) \to \mathrm{Conv}(a,c) \to \varphi_{85}) \to \varphi_{131}$}
\BIC{$(\mathrm{FLim}(f,c,\alpha) \to \mathrm{Conv}(a,c) \to \varphi_{85}) \to \varphi_{131}$}
\end{ptree}
\noindent$\Pi_{36}$:
\begin{ptree}
\AXC{$\mathrm{FLim}(f,c,\alpha) \to \mathrm{Conv}(a,c) \to \varphi_{110}$\lbl{A1}}
\AXC{$\varphi_{133}$\lbl{A1}}
\BIC{$\mathrm{FLim}(f,c,\alpha) \to \mathrm{Conv}(a,c) \to \varphi_{85}$}
\AXC{$\stackrel{\Pi_{35}}{\vdots}$}\noLine
\UIC{$(\mathrm{FLim}(f,c,\alpha) \to \mathrm{Conv}(a,c) \to \varphi_{85}) \to \varphi_{131}$}
\BIC{$\varphi_{131}$}
\end{ptree}
\noindent$\Pi_{37}$:
\begin{ptree}
\AXC{$\mathrm{FLim}(f,c,\alpha) \to \mathrm{Conv}(a,c) \to \varphi_{110}$\lbl{A1}}
\AXC{$\varphi_{132}$\lbl{A1}}
\BIC{$\mathrm{FLim}(f,c,\alpha) \to \mathrm{Conv}(a,c) \to \varphi_{84}$}
\AXC{$\stackrel{\Pi_{36}}{\vdots}$}\noLine
\UIC{$\varphi_{131}$}
\BIC{$\mathrm{FLim}(f,c,\alpha) \to \mathrm{Conv}(a,c) \to \varphi_{90}$}
\AXC{$\varphi_{164}$\lbl{A1}}
\BIC{$\varphi_{123} \to (\varphi_{123} \to \forall N\,(\varphi_{106}) \to \varphi_{8}) \to \varphi_{157}$}
\end{ptree}
\noindent$\Pi_{38}$:
\begin{ptree}
\AXC{$\varphi_{123} \to \forall N\,(\varphi_{106}) \to \varphi_{8}$\lbl{A3}}
\AXC{$\stackrel{\Pi_{34}}{\vdots}$}\noLine
\UIC{$\mathrm{FLim}(f,c,\alpha) \to \mathrm{Conv}(a,c) \to \varphi_{120}$}
\AXC{$\varphi_{149}$\lbl{A1}}
\BIC{$\varphi_{106} \to \neg (N \in \mathbb{N} \land \forall n \in \mathbb{N}\,(N \leq n \to \lvert a_{n} - c \rvert < \delta))$}
\UIC{$\varphi_{123}$}
\AXC{$\stackrel{\Pi_{37}}{\vdots}$}\noLine
\UIC{$\varphi_{123} \to (\varphi_{123} \to \forall N\,(\varphi_{106}) \to \varphi_{8}) \to \varphi_{157}$}
\BIC{$(\varphi_{123} \to \forall N\,(\varphi_{106}) \to \varphi_{8}) \to \varphi_{157}$}
\BIC{$\varphi_{157}$}
\end{ptree}
\noindent$\Pi_{39}$:
\begin{ptree}
\AXC{$\forall f\,\forall c\,\forall \alpha\,\forall \varepsilon\,(\varphi_{75})$\lbl{[D1]}}
\AXC{$\forall f\,\forall c\,\forall \alpha\,\forall \varepsilon\,(\varphi_{75}) \to \forall c\,\forall \alpha\,\forall \varepsilon\,(\varphi_{75})$\lbl{A2}}
\BIC{$\forall c\,\forall \alpha\,\forall \varepsilon\,(\varphi_{75})$}
\AXC{$\forall c\,\forall \alpha\,\forall \varepsilon\,(\varphi_{75}) \to \forall \alpha\,\forall \varepsilon\,(\varphi_{75})$\lbl{A2}}
\BIC{$\forall \alpha\,\forall \varepsilon\,(\varphi_{75})$}
\end{ptree}
\noindent$\Pi_{40}$:
\begin{ptree}
\AXC{$\stackrel{\Pi_{39}}{\vdots}$}\noLine
\UIC{$\forall \alpha\,\forall \varepsilon\,(\varphi_{75})$}
\AXC{$\forall \alpha\,\forall \varepsilon\,(\varphi_{75}) \to \forall \varepsilon\,(\varphi_{75})$\lbl{A2}}
\BIC{$\forall \varepsilon\,(\varphi_{75})$}
\AXC{$\forall \varepsilon\,(\varphi_{75}) \to \varphi_{75}$\lbl{A2}}
\BIC{$\varphi_{75}$}
\AXC{$\varphi_{75} \to \varphi_{128}$\lbl{A1}}
\BIC{$\varphi_{128}$}
\end{ptree}
\noindent$\Pi_{41}$:
\begin{ptree}
\AXC{$\mathrm{FLim}(f,c,\alpha) \to \mathrm{Conv}(a,c) \to \varphi_{24}$\lbl{A1}}
\AXC{$\mathrm{FLim}(f,c,\alpha) \to \mathrm{FLim}(f,c,\alpha)$\lbl{A1}}
\AXC{$\stackrel{\Pi_{40}}{\vdots}$}\noLine
\UIC{$\varphi_{128}$}
\BIC{$(\mathrm{FLim}(f,c,\alpha) \to \mathrm{Conv}(a,c) \to \varphi_{24}) \to \varphi_{121}$}
\BIC{$\varphi_{121}$}
\end{ptree}
\noindent$\Pi_{42}$:
\begin{ptree}
\AXC{$\mathrm{FLim}(f,c,\alpha) \to \mathrm{Conv}(a,c) \to \varphi_{33}$\lbl{A1}}
\AXC{$\stackrel{\Pi_{41}}{\vdots}$}\noLine
\UIC{$\varphi_{121}$}
\BIC{$\mathrm{FLim}(f,c,\alpha) \to \mathrm{Conv}(a,c) \to \varphi_{83}$}
\AXC{$(\mathrm{FLim}(f,c,\alpha) \to \mathrm{Conv}(a,c) \to \varphi_{83}) \to \varphi_{162}$\lbl{A1}}
\BIC{$\varphi_{162}$}
\end{ptree}
\noindent$\Pi_{43}$:
\begin{ptree}
\AXC{$\varphi_{80} \to \forall \delta\,(\varphi_{80})$\lbl{A4}}
\AXC{$\varphi_{135}$\lbl{A3}}
\AXC{$\varphi_{106} \to \forall N\,(\varphi_{106})$\lbl{A4}}
\AXC{$\stackrel{\Pi_{38}}{\vdots}$}\noLine
\UIC{$\varphi_{157}$}
\BIC{$\mathrm{FLim}(f,c,\alpha) \to \mathrm{Conv}(a,c) \to \varphi_{93}$}
\AXC{$\varphi_{134}$\lbl{A1}}
\BIC{$\varphi_{80} \to \neg (\delta \in \mathbb{R} \land 0 < \delta \land \varphi_{29})$}
\UIC{$\forall \delta\,(\varphi_{80} \to \neg (\delta \in \mathbb{R} \land 0 < \delta \land \varphi_{29}))$}
\AXC{$\stackrel{\Pi_{42}}{\vdots}$}\noLine
\UIC{$\varphi_{162}$}
\BIC{$\varphi_{135} \to (\varphi_{80} \to \forall \delta\,(\varphi_{80})) \to \mathrm{FLim}(f,c,\alpha) \to \varphi_{77}$}
\BIC{$(\varphi_{80} \to \forall \delta\,(\varphi_{80})) \to \mathrm{FLim}(f,c,\alpha) \to \varphi_{77}$}
\BIC{$\mathrm{FLim}(f,c,\alpha) \to \varphi_{77}$}
\end{ptree}
\noindent$\Pi_{44}$:
\begin{ptree}
\AXC{$\stackrel{\Pi_{43}}{\vdots}$}\noLine
\UIC{$\mathrm{FLim}(f,c,\alpha) \to \varphi_{77}$}
\AXC{$(\mathrm{FLim}(f,c,\alpha) \to \varphi_{77}) \to \varphi_{22} \to \varepsilon \in \mathbb{R} \to \varphi_{31}$\lbl{A1}}
\BIC{$\varphi_{22} \to \varepsilon \in \mathbb{R} \to \varphi_{31}$}
\UIC{$\forall \varepsilon\,(\varphi_{22} \to \varepsilon \in \mathbb{R} \to \varphi_{31})$}
\AXC{$\forall \varepsilon\,(\varphi_{22} \to \varepsilon \in \mathbb{R} \to \varphi_{31}) \to \varphi_{150}$\lbl{A1}}
\BIC{$\varphi_{150}$}
\end{ptree}
\noindent$\Pi_{45}$:
\begin{ptree}
\AXC{$\forall f\,\forall c\,\forall \alpha\,(\mathrm{FLim}(f,c,\alpha) \to \alpha \in \mathbb{R})$\lbl{[D4]}}
\AXC{$\varphi_{17}$\lbl{A2}}
\BIC{$\forall c\,\forall \alpha\,(\mathrm{FLim}(f,c,\alpha) \to \alpha \in \mathbb{R})$}
\AXC{$\varphi_{12}$\lbl{A2}}
\BIC{$\forall \alpha\,(\mathrm{FLim}(f,c,\alpha) \to \alpha \in \mathbb{R})$}
\AXC{$\varphi_{7}$\lbl{A2}}
\BIC{$\mathrm{FLim}(f,c,\alpha) \to \alpha \in \mathbb{R}$}
\end{ptree}
\noindent$\Pi_{46}$:
\begin{ptree}
\AXC{$\mathrm{FLim}(f,c,\alpha) \to \mathrm{FLim}(f,c,\alpha)$\lbl{A1}}
\AXC{$\stackrel{\Pi_{45}}{\vdots}$}\noLine
\UIC{$\mathrm{FLim}(f,c,\alpha) \to \alpha \in \mathbb{R}$}
\AXC{$(\mathrm{FLim}(f,c,\alpha) \to \alpha \in \mathbb{R}) \to \varphi_{67}$\lbl{A1}}
\BIC{$\varphi_{67}$}
\BIC{$\mathrm{FLim}(f,c,\alpha) \to \varphi_{21}$}
\end{ptree}
\noindent$\Pi_{47}$:
\begin{ptree}
\AXC{$\varphi_{73}$\lbl{DefConv}}
\AXC{$\varphi_{73} \to \forall \alpha\,(\varphi_{76})$\lbl{A2}}
\BIC{$\forall \alpha\,(\varphi_{76})$}
\AXC{$\forall \alpha\,(\varphi_{76}) \to \varphi_{76}$\lbl{A2}}
\BIC{$\varphi_{76}$}
\AXC{$\varphi_{76} \to \varphi_{129}$\lbl{A1}}
\BIC{$\varphi_{129}$}
\end{ptree}
\noindent$\Pi_{48}$:
\begin{ptree}
\AXC{$\stackrel{\Pi_{46}}{\vdots}$}\noLine
\UIC{$\mathrm{FLim}(f,c,\alpha) \to \varphi_{21}$}
\AXC{$\stackrel{\Pi_{26}}{\vdots}$}\noLine
\UIC{$\mathrm{FLim}(f,c,\alpha) \to \mathrm{Conv}(a,c) \to \varphi_{68}$}
\AXC{$\varphi_{89}$\lbl{A1}}
\BIC{$\mathrm{FLim}(f,c,\alpha) \to \varphi_{25}$}
\AXC{$\stackrel{\Pi_{47}}{\vdots}$}\noLine
\UIC{$\varphi_{129}$}
\BIC{$(\mathrm{FLim}(f,c,\alpha) \to \varphi_{21}) \to \varphi_{104}$}
\BIC{$\varphi_{104}$}
\end{ptree}
\noindent 主木:
\begin{ptree}
\AXC{$\varphi_{22} \to \forall \varepsilon\,(\varphi_{22})$\lbl{A4}}
\AXC{$\forall \varepsilon\,(\varphi_{22} \to \varepsilon \in \mathbb{R} \to \varphi_{31}) \to \forall \varepsilon\,(\varphi_{22}) \to \forall \varepsilon \in \mathbb{R}\,(\varphi_{31})$\lbl{A3}}
\AXC{$\stackrel{\Pi_{44}}{\vdots}$}\noLine
\UIC{$\varphi_{150}$}
\BIC{$(\varphi_{22} \to \forall \varepsilon\,(\varphi_{22})) \to \mathrm{FLim}(f,c,\alpha) \to \varphi_{78}$}
\BIC{$\mathrm{FLim}(f,c,\alpha) \to \varphi_{78}$}
\AXC{$\stackrel{\Pi_{48}}{\vdots}$}\noLine
\UIC{$\varphi_{104}$}
\BIC{$\mathrm{FLim}(f,c,\alpha) \to \mathrm{Conv}(a,c) \to \varphi_{19}$}
\UIC{$\forall a\,(\mathrm{FLim}(f,c,\alpha) \to \mathrm{Conv}(a,c) \to \varphi_{19})$}
\UIC{$\forall \alpha\,\forall a\,(\mathrm{FLim}(f,c,\alpha) \to \mathrm{Conv}(a,c) \to \varphi_{19})$}
\UIC{$\varphi_{53}$}
\UIC{$\forall f\,(\varphi_{53})$}
\end{ptree}
\begin{abbrevlist}
\item $\varphi_{1} :\equiv \mathrm{FLim}(f,c,\alpha) \to \mathrm{Conv}(a,c) \to \mathrm{Conv}(a,c)$
\item $\varphi_{2} :\equiv \forall \alpha\,(\mathrm{Conv}(a,\alpha) \to \mathrm{Seq}(a)) \to \mathrm{Conv}(a,c) \to \mathrm{Seq}(a)$
\item $\varphi_{3} :\equiv \forall \alpha\,(\mathrm{FLim}(f,c,\alpha) \to \mathrm{RFn}(f)) \to \mathrm{FLim}(f,c,\alpha) \to \mathrm{RFn}(f)$
\item $\varphi_{4} :\equiv \forall a\,\forall \alpha\,(\mathrm{Conv}(a,\alpha) \to \mathrm{Seq}(a)) \to \forall \alpha\,(\mathrm{Conv}(a,\alpha) \to \mathrm{Seq}(a))$
\item $\varphi_{5} :\equiv \forall c\,\forall \alpha\,(\mathrm{FLim}(f,c,\alpha) \to \mathrm{RFn}(f)) \to \forall \alpha\,(\mathrm{FLim}(f,c,\alpha) \to \mathrm{RFn}(f))$
\item $\varphi_{6} :\equiv \forall \alpha\,(\mathrm{FLim}(f,c,\alpha) \to c \in \mathbb{R}) \to \mathrm{FLim}(f,c,\alpha) \to c \in \mathbb{R}$
\item $\varphi_{7} :\equiv \forall \alpha\,(\mathrm{FLim}(f,c,\alpha) \to \alpha \in \mathbb{R}) \to \mathrm{FLim}(f,c,\alpha) \to \alpha \in \mathbb{R}$
\item $\varphi_{8} :\equiv \forall N\,(\neg (N \in \mathbb{N} \land \forall n \in \mathbb{N}\,(N \leq n \to \lvert a_{n} - c \rvert < \delta)))$
\item $\varphi_{9} :\equiv \forall f\,\forall c\,\forall \alpha\,(\mathrm{FLim}(f,c,\alpha) \to \mathrm{RFn}(f)) \to \forall c\,\forall \alpha\,(\mathrm{FLim}(f,c,\alpha) \to \mathrm{RFn}(f))$
\item $\varphi_{10} :\equiv \forall c\,\forall \alpha\,(\mathrm{FLim}(f,c,\alpha) \to c \in \mathbb{R}) \to \forall \alpha\,(\mathrm{FLim}(f,c,\alpha) \to c \in \mathbb{R})$
\item $\varphi_{11} :\equiv \forall x\,\forall y\,(x = y \to y = x) \to \forall y\,(f \circ a(n) = y \to y = f \circ a(n))$
\item $\varphi_{12} :\equiv \forall c\,\forall \alpha\,(\mathrm{FLim}(f,c,\alpha) \to \alpha \in \mathbb{R}) \to \forall \alpha\,(\mathrm{FLim}(f,c,\alpha) \to \alpha \in \mathbb{R})$
\item $\varphi_{13} :\equiv \mathrm{Conv}(a,c) \to \forall n \in \mathbb{N}\,(a_{n} \in \mathrm{dom}(f) \land a_{n} \neq c) \to \mathrm{Seq}(a)$
\item $\varphi_{14} :\equiv \mathrm{Conv}(a,c) \to \forall n \in \mathbb{N}\,(a_{n} \in \mathrm{dom}(f) \land a_{n} \neq c) \to \mathrm{RFn}(f)$
\item $\varphi_{15} :\equiv \forall f\,\forall c\,\forall \alpha\,(\mathrm{FLim}(f,c,\alpha) \to c \in \mathbb{R}) \to \forall c\,\forall \alpha\,(\mathrm{FLim}(f,c,\alpha) \to c \in \mathbb{R})$
\item $\varphi_{16} :\equiv \neg (N \in \mathbb{N} \land \forall n \in \mathbb{N}\,(N \leq n \to \lvert f \circ a(n) - \alpha \rvert < \varepsilon))$
\item $\varphi_{17} :\equiv \forall f\,\forall c\,\forall \alpha\,(\mathrm{FLim}(f,c,\alpha) \to \alpha \in \mathbb{R}) \to \forall c\,\forall \alpha\,(\mathrm{FLim}(f,c,\alpha) \to \alpha \in \mathbb{R})$
\item $\varphi_{18} :\equiv \forall n \in \mathbb{N}\,(a_{n} \in \mathrm{dom}(f) \land a_{n} \neq c) \to n \in \mathbb{N} \to n \in \mathbb{N}$
\item $\varphi_{19} :\equiv \forall n \in \mathbb{N}\,(a_{n} \in \mathrm{dom}(f) \land a_{n} \neq c) \to \mathrm{Conv}(f \circ a,\alpha)$
\item $\varphi_{20} :\equiv \mathrm{Conv}(a,c) \to \forall n \in \mathbb{N}\,(a_{n} \in \mathrm{dom}(f) \land a_{n} \neq c) \to c \in \mathbb{R}$
\item $\varphi_{21} :\equiv \mathrm{Conv}(a,c) \to \forall n \in \mathbb{N}\,(a_{n} \in \mathrm{dom}(f) \land a_{n} \neq c) \to \alpha \in \mathbb{R}$
\item $\varphi_{22} :\equiv \mathrm{FLim}(f,c,\alpha) \land \mathrm{Conv}(a,c) \land \forall n \in \mathbb{N}\,(a_{n} \in \mathrm{dom}(f) \land a_{n} \neq c)$
\item $\varphi_{23} :\equiv \forall w\,(f(a_{n}) = w \to \lvert f(a_{n}) - \alpha \rvert < \varepsilon \to \lvert w - \alpha \rvert < \varepsilon)$
\item $\varphi_{24} :\equiv \forall n \in \mathbb{N}\,(a_{n} \in \mathrm{dom}(f) \land a_{n} \neq c) \to \varepsilon \in \mathbb{R} \to \varepsilon \in \mathbb{R}$
\item $\varphi_{25} :\equiv \mathrm{Conv}(a,c) \to \forall n \in \mathbb{N}\,(a_{n} \in \mathrm{dom}(f) \land a_{n} \neq c) \to \mathrm{Seq}(f \circ a)$
\item $\varphi_{26} :\equiv N \in \mathbb{N} \land \forall n \in \mathbb{N}\,(N \leq n \to \lvert a_{n} - c \rvert < \delta) \to n \in \mathbb{N} \to n \in \mathbb{N}$
\item $\varphi_{27} :\equiv \forall n\,(\mathrm{Seq}(a) \to n \in \mathbb{N} \to a_{n} \in \mathbb{R}) \to \mathrm{Seq}(a) \to n \in \mathbb{N} \to a_{n} \in \mathbb{R}$
\item $\varphi_{28} :\equiv \forall n \in \mathbb{N}\,(a_{n} \in \mathrm{dom}(f) \land a_{n} \neq c) \to n \in \mathbb{N} \to a_{n} \in \mathrm{dom}(f)$
\item $\varphi_{29} :\equiv \forall x \in \mathrm{dom}(f)\,(0 < \lvert x - c \rvert \to \lvert x - c \rvert < \delta \to \lvert f(x) - \alpha \rvert < \varepsilon)$
\item $\varphi_{30} :\equiv \delta \in \mathbb{R} \to 0 < \delta \to \exists N \in \mathbb{N}\,\forall n \in \mathbb{N}\,(N \leq n \to \lvert a_{n} - c \rvert < \delta)$
\item $\varphi_{31} :\equiv 0 < \varepsilon \to \exists N \in \mathbb{N}\,\forall n \in \mathbb{N}\,(N \leq n \to \lvert f \circ a(n) - \alpha \rvert < \varepsilon)$
\item $\varphi_{32} :\equiv \forall n \in \mathbb{N}\,(a_{n} \in \mathrm{dom}(f) \land a_{n} \neq c) \to \forall n \in \mathbb{N}\,(a_{n} \in \mathrm{dom}(f))$
\item $\varphi_{33} :\equiv \forall n \in \mathbb{N}\,(a_{n} \in \mathrm{dom}(f) \land a_{n} \neq c) \to \varepsilon \in \mathbb{R} \to 0 < \varepsilon \to 0 < \varepsilon$
\item $\varphi_{34} :\equiv \forall a\,\forall n\,(\mathrm{Seq}(a) \to n \in \mathbb{N} \to a_{n} \in \mathbb{R}) \to \forall n\,(\mathrm{Seq}(a) \to n \in \mathbb{N} \to a_{n} \in \mathbb{R})$
\item $\varphi_{35} :\equiv N \in \mathbb{N} \land \forall n \in \mathbb{N}\,(N \leq n \to \lvert a_{n} - c \rvert < \delta) \to n \in \mathbb{N} \to N \leq n \to a_{n} \neq c$
\item $\varphi_{36} :\equiv \forall n \in \mathbb{N}\,(N \leq n \to \lvert a_{n} - c \rvert < \delta) \to n \in \mathbb{N} \to N \leq n \to \lvert a_{n} - c \rvert < \delta$
\item $\varphi_{37} :\equiv N \in \mathbb{N} \land \forall n \in \mathbb{N}\,(N \leq n \to \lvert a_{n} - c \rvert < \delta) \to n \in \mathbb{N} \to N \leq n \to a_{n} - c \neq 0$
\item $\varphi_{38} :\equiv N \in \mathbb{N} \land \forall n \in \mathbb{N}\,(N \leq n \to \lvert a_{n} - c \rvert < \delta) \to n \in \mathbb{N} \to N \leq n \to a_{n} \in \mathbb{R}$
\item $\varphi_{39} :\equiv \forall n \in \mathbb{N}\,(a_{n} \in \mathrm{dom}(f) \land a_{n} \neq c) \to n \in \mathbb{N} \to a_{n} \in \mathrm{dom}(f) \land a_{n} \neq c$
\item $\varphi_{40} :\equiv \forall n \in \mathbb{N}\,(a_{n} \in \mathrm{dom}(f) \land a_{n} \neq c) \to n \in \mathbb{N} \to f \circ a(n) = f(a_{n})$
\item $\varphi_{41} :\equiv \forall n \in \mathbb{N}\,(a_{n} \in \mathrm{dom}(f) \land a_{n} \neq c) \to \forall n \in \mathbb{N}\,(a_{n} \in \mathrm{dom}(f) \land a_{n} \neq c)$
\item $\varphi_{42} :\equiv \forall n \in \mathbb{N}\,(f \circ a(n) = f(a_{n})) \to n \in \mathbb{N} \to f \circ a(n) = f(a_{n})$
\item $\varphi_{43} :\equiv N \in \mathbb{N} \land \forall n \in \mathbb{N}\,(N \leq n \to \lvert a_{n} - c \rvert < \delta) \to n \in \mathbb{N} \to N \leq n \to a_{n} - c \in \mathbb{R}$
\item $\varphi_{44} :\equiv N \in \mathbb{N} \land \forall n \in \mathbb{N}\,(N \leq n \to \lvert a_{n} - c \rvert < \delta) \to n \in \mathbb{N} \to N \leq n \to 0 < \lvert a_{n} - c \rvert$
\item $\varphi_{45} :\equiv \forall n \in \mathbb{N}\,(a_{n} \in \mathrm{dom}(f) \land a_{n} \neq c) \to \forall n \in \mathbb{N}\,(f \circ a(n) = f(a_{n}))$
\item $\varphi_{46} :\equiv N \in \mathbb{N} \land \forall n \in \mathbb{N}\,(N \leq n \to \lvert a_{n} - c \rvert < \delta) \to n \in \mathbb{N} \to N \leq n \to \lvert a_{n} - c \rvert < \delta$
\item $\varphi_{47} :\equiv \forall x \in \mathbb{R}\,(x \neq 0 \to 0 < \lvert x \rvert) \to a_{n} - c \in \mathbb{R} \to a_{n} - c \neq 0 \to 0 < \lvert a_{n} - c \rvert$
\item $\varphi_{48} :\equiv \forall \varepsilon\,(\mathrm{Conv}(a,c) \to \varepsilon \in \mathbb{R} \to 0 < \varepsilon \to \exists N \in \mathbb{N}\,\forall n \in \mathbb{N}\,(N \leq n \to \lvert a_{n} - c \rvert < \varepsilon))$
\item $\varphi_{49} :\equiv N \in \mathbb{N} \land \forall n \in \mathbb{N}\,(N \leq n \to \lvert a_{n} - c \rvert < \delta) \to \forall n \in \mathbb{N}\,(N \leq n \to \lvert a_{n} - c \rvert < \delta)$
\item $\varphi_{50} :\equiv \forall n \in \mathbb{N}\,(a_{n} \in \mathrm{dom}(f)) \to \mathrm{Seq}(f \circ a) \land \forall n \in \mathbb{N}\,(f \circ a(n) = f(a_{n}))$
\item $\varphi_{51} :\equiv \forall x\,\forall y\,(x \in \mathbb{R} \to y \in \mathbb{R} \to x - y \in \mathbb{R}) \to \forall y\,(a_{n} \in \mathbb{R} \to y \in \mathbb{R} \to a_{n} - y \in \mathbb{R})$
\item $\varphi_{52} :\equiv a_{n} \in \mathrm{dom}(f) \to 0 < \lvert a_{n} - c \rvert \to \lvert a_{n} - c \rvert < \delta \to \lvert f(a_{n}) - \alpha \rvert < \varepsilon$
\item $\varphi_{53} :\equiv \forall c\,\forall \alpha\,\forall a\,(\mathrm{FLim}(f,c,\alpha) \to \mathrm{Conv}(a,c) \to \varphi_{19})$
\item $\varphi_{54} :\equiv \forall \alpha\,\forall \varepsilon\,(\mathrm{Conv}(a,\alpha) \to \varepsilon \in \mathbb{R} \to 0 < \varepsilon \to \exists N \in \mathbb{N}\,\forall n \in \mathbb{N}\,(N \leq n \to \lvert a_{n} - \alpha \rvert < \varepsilon))$
\item $\varphi_{55} :\equiv f(a_{n}) = f \circ a(n) \to \lvert f(a_{n}) - \alpha \rvert < \varepsilon \to \lvert f \circ a(n) - \alpha \rvert < \varepsilon$
\item $\varphi_{56} :\equiv N \in \mathbb{N} \land \forall n \in \mathbb{N}\,(N \leq n \to \lvert a_{n} - c \rvert < \delta) \to n \in \mathbb{N} \to N \leq n \to a_{n} \in \mathrm{dom}(f) \land a_{n} \neq c$
\item $\varphi_{57} :\equiv N \in \mathbb{N} \land \forall n \in \mathbb{N}\,(N \leq n \to \lvert a_{n} - c \rvert < \delta) \to n \in \mathbb{N} \to N \leq n \to \lvert f(a_{n}) - \alpha \rvert < \varepsilon$
\item $\varphi_{58} :\equiv N \in \mathbb{N} \land \forall n \in \mathbb{N}\,(N \leq n \to \lvert a_{n} - c \rvert < \delta) \to n \in \mathbb{N} \to N \leq n \to f \circ a(n) = f(a_{n})$
\item $\varphi_{59} :\equiv N \in \mathbb{N} \land \forall n \in \mathbb{N}\,(N \leq n \to \lvert a_{n} - c \rvert < \delta) \to n \in \mathbb{N} \to N \leq n \to f(a_{n}) = f \circ a(n)$
\item $\varphi_{60} :\equiv (\mathrm{FLim}(f,c,\alpha) \to \mathrm{FLim}(f,c,\alpha)) \to \mathrm{FLim}(f,c,\alpha) \to \varphi_{14}$
\item $\varphi_{61} :\equiv N \in \mathbb{N} \land \forall n \in \mathbb{N}\,(N \leq n \to \lvert a_{n} - c \rvert < \delta) \to n \in \mathbb{N} \to N \leq n \to \lvert f \circ a(n) - \alpha \rvert < \varepsilon$
\item $\varphi_{62} :\equiv \forall y\,(a_{n} \in \mathbb{R} \to y \in \mathbb{R} \to a_{n} - y \in \mathbb{R}) \to a_{n} \in \mathbb{R} \to c \in \mathbb{R} \to a_{n} - c \in \mathbb{R}$
\item $\varphi_{63} :\equiv \forall x\,\forall y\,(x \in \mathbb{R} \to y \in \mathbb{R} \to x \neq y \to x - y \neq 0) \to \forall y\,(a_{n} \in \mathbb{R} \to y \in \mathbb{R} \to a_{n} \neq y \to a_{n} - y \neq 0)$
\item $\varphi_{64} :\equiv N \in \mathbb{N} \land \forall n \in \mathbb{N}\,(N \leq n \to \lvert a_{n} - c \rvert < \delta) \to N \in \mathbb{N} \land \forall n \in \mathbb{N}\,(N \leq n \to \lvert a_{n} - c \rvert < \delta)$
\item $\varphi_{65} :\equiv (\mathrm{FLim}(f,c,\alpha) \to \mathrm{FLim}(f,c,\alpha)) \to \mathrm{FLim}(f,c,\alpha) \to \varphi_{20}$
\item $\varphi_{66} :\equiv N \in \mathbb{N} \land \forall n \in \mathbb{N}\,(N \leq n \to \lvert a_{n} - c \rvert < \delta) \to \forall n \in \mathbb{N}\,(N \leq n \to \lvert f \circ a(n) - \alpha \rvert < \varepsilon)$
\item $\varphi_{67} :\equiv (\mathrm{FLim}(f,c,\alpha) \to \mathrm{FLim}(f,c,\alpha)) \to \mathrm{FLim}(f,c,\alpha) \to \varphi_{21}$
\item $\varphi_{68} :\equiv \forall n \in \mathbb{N}\,(a_{n} \in \mathrm{dom}(f) \land a_{n} \neq c) \to \mathrm{Seq}(f \circ a) \land \forall n \in \mathbb{N}\,(f \circ a(n) = f(a_{n}))$
\item $\varphi_{69} :\equiv \forall y\,(f \circ a(n) = y \to y = f \circ a(n)) \to f \circ a(n) = f(a_{n}) \to f(a_{n}) = f \circ a(n)$
\item $\varphi_{70} :\equiv N \in \mathbb{N} \land \forall n \in \mathbb{N}\,(N \leq n \to \lvert a_{n} - c \rvert < \delta) \to N \in \mathbb{N} \land \forall n \in \mathbb{N}\,(N \leq n \to \lvert f \circ a(n) - \alpha \rvert < \varepsilon)$
\item $\varphi_{71} :\equiv \forall y\,(a_{n} \in \mathbb{R} \to y \in \mathbb{R} \to a_{n} \neq y \to a_{n} - y \neq 0) \to a_{n} \in \mathbb{R} \to c \in \mathbb{R} \to a_{n} \neq c \to a_{n} - c \neq 0$
\item $\varphi_{72} :\equiv N \in \mathbb{N} \land \forall n \in \mathbb{N}\,(N \leq n \to \lvert a_{n} - c \rvert < \delta) \to \exists N \in \mathbb{N}\,\forall n \in \mathbb{N}\,(N \leq n \to \lvert f \circ a(n) - \alpha \rvert < \varepsilon)$
\item $\varphi_{73} :\equiv \forall a\,\forall \alpha\,(\mathrm{Conv}(a,\alpha) \leftrightarrow \mathrm{Seq}(a) \land \alpha \in \mathbb{R} \land \forall \varepsilon \in \mathbb{R}\,(0 < \varepsilon \to \exists N \in \mathbb{N}\,\forall n \in \mathbb{N}\,(N \leq n \to \lvert a_{n} - \alpha \rvert < \varepsilon)))$
\item $\varphi_{74} :\equiv (\mathrm{Conv}(a,c) \to \mathrm{Seq}(a)) \to \varphi_{1} \to \mathrm{FLim}(f,c,\alpha) \to \varphi_{13}$
\item $\varphi_{75} :\equiv \mathrm{FLim}(f,c,\alpha) \to \varepsilon \in \mathbb{R} \to 0 < \varepsilon \to \exists \delta\,(\delta \in \mathbb{R} \land 0 < \delta \land \varphi_{29})$
\item $\varphi_{76} :\equiv \mathrm{Conv}(f \circ a,\alpha) \leftrightarrow \mathrm{Seq}(f \circ a) \land \alpha \in \mathbb{R} \land \forall \varepsilon \in \mathbb{R}\,(\varphi_{31})$
\item $\varphi_{77} :\equiv \mathrm{Conv}(a,c) \to \forall n \in \mathbb{N}\,(a_{n} \in \mathrm{dom}(f) \land a_{n} \neq c) \to \varepsilon \in \mathbb{R} \to \varphi_{31}$
\item $\varphi_{78} :\equiv \mathrm{Conv}(a,c) \to \forall n \in \mathbb{N}\,(a_{n} \in \mathrm{dom}(f) \land a_{n} \neq c) \to \forall \varepsilon \in \mathbb{R}\,(\varphi_{31})$
\item $\varphi_{79} :\equiv \delta \in \mathbb{R} \land 0 < \delta \land \varphi_{29} \to \exists N \in \mathbb{N}\,\forall n \in \mathbb{N}\,(N \leq n \to \lvert a_{n} - c \rvert < \delta)$
\item $\varphi_{80} :\equiv \varphi_{22} \land \varepsilon \in \mathbb{R} \land 0 < \varepsilon \land \neg \exists N \in \mathbb{N}\,\forall n \in \mathbb{N}\,(N \leq n \to \lvert f \circ a(n) - \alpha \rvert < \varepsilon)$
\item $\varphi_{81} :\equiv \delta \in \mathbb{R} \land 0 < \delta \land \varphi_{29} \to N \in \mathbb{N} \land \forall n \in \mathbb{N}\,(N \leq n \to \lvert a_{n} - c \rvert < \delta) \to N \in \mathbb{N}$
\item $\varphi_{82} :\equiv \delta \in \mathbb{R} \land 0 < \delta \land \varphi_{29} \to \exists N \in \mathbb{N}\,\forall n \in \mathbb{N}\,(N \leq n \to \lvert f \circ a(n) - \alpha \rvert < \varepsilon)$
\item $\varphi_{83} :\equiv \forall n \in \mathbb{N}\,(a_{n} \in \mathrm{dom}(f) \land a_{n} \neq c) \to \varepsilon \in \mathbb{R} \to 0 < \varepsilon \to \exists \delta\,(\delta \in \mathbb{R} \land 0 < \delta \land \varphi_{29})$
\item $\varphi_{84} :\equiv \forall n \in \mathbb{N}\,(a_{n} \in \mathrm{dom}(f) \land a_{n} \neq c) \to \varepsilon \in \mathbb{R} \to 0 < \varepsilon \to \delta \in \mathbb{R} \land 0 < \delta \land \varphi_{29} \to 0 < \delta$
\item $\varphi_{85} :\equiv \forall n \in \mathbb{N}\,(a_{n} \in \mathrm{dom}(f) \land a_{n} \neq c) \to \varepsilon \in \mathbb{R} \to 0 < \varepsilon \to \delta \in \mathbb{R} \land 0 < \delta \land \varphi_{29} \to \delta \in \mathbb{R}$
\item $\varphi_{86} :\equiv (\mathrm{FLim}(f,c,\alpha) \to \mathrm{Conv}(a,c) \to \varphi_{18}) \to \mathrm{FLim}(f,c,\alpha) \to \mathrm{Conv}(a,c) \to \varphi_{28}$
\item $\varphi_{87} :\equiv (\mathrm{FLim}(f,c,\alpha) \to \mathrm{Conv}(a,c) \to \varphi_{28}) \to \varphi_{22} \to n \in \mathbb{N} \to a_{n} \in \mathrm{dom}(f)$
\item $\varphi_{88} :\equiv \forall n\,(\varphi_{22} \to n \in \mathbb{N} \to a_{n} \in \mathrm{dom}(f)) \to \forall n\,(\varphi_{22}) \to \forall n \in \mathbb{N}\,(a_{n} \in \mathrm{dom}(f))$
\item $\varphi_{89} :\equiv (\mathrm{FLim}(f,c,\alpha) \to \mathrm{Conv}(a,c) \to \varphi_{68}) \to \mathrm{FLim}(f,c,\alpha) \to \varphi_{25}$
\item $\varphi_{90} :\equiv \forall n \in \mathbb{N}\,(a_{n} \in \mathrm{dom}(f) \land a_{n} \neq c) \to \varepsilon \in \mathbb{R} \to 0 < \varepsilon \to \varphi_{79}$
\item $\varphi_{91} :\equiv \forall a\,(\mathrm{RFn}(f) \to \mathrm{Seq}(a) \to \varphi_{50}) \to \mathrm{RFn}(f) \to \mathrm{Seq}(a) \to \varphi_{50}$
\item $\varphi_{92} :\equiv \forall n \in \mathbb{N}\,(a_{n} \in \mathrm{dom}(f) \land a_{n} \neq c) \to \varepsilon \in \mathbb{R} \to 0 < \varepsilon \to \varphi_{81}$
\item $\varphi_{93} :\equiv \forall n \in \mathbb{N}\,(a_{n} \in \mathrm{dom}(f) \land a_{n} \neq c) \to \varepsilon \in \mathbb{R} \to 0 < \varepsilon \to \varphi_{82}$
\item $\varphi_{94} :\equiv \forall f\,\forall a\,(\mathrm{RFn}(f) \to \mathrm{Seq}(a) \to \varphi_{50}) \to \forall a\,(\mathrm{RFn}(f) \to \mathrm{Seq}(a) \to \varphi_{50})$
\item $\varphi_{95} :\equiv (\mathrm{FLim}(f,c,\alpha) \to \mathrm{Conv}(a,c) \to \varphi_{32}) \to \mathrm{FLim}(f,c,\alpha) \to \mathrm{Conv}(a,c) \to \varphi_{68}$
\item $\varphi_{96} :\equiv \forall n \in \mathbb{N}\,(a_{n} \in \mathrm{dom}(f) \land a_{n} \neq c) \to \varepsilon \in \mathbb{R} \to 0 < \varepsilon \to \delta \in \mathbb{R} \land 0 < \delta \land \varphi_{29} \to \varphi_{26}$
\item $\varphi_{97} :\equiv \forall n \in \mathbb{N}\,(a_{n} \in \mathrm{dom}(f) \land a_{n} \neq c) \to \varepsilon \in \mathbb{R} \to 0 < \varepsilon \to \delta \in \mathbb{R} \land 0 < \delta \land \varphi_{29} \to \varphi_{29}$
\item $\varphi_{98} :\equiv (\varphi_{22} \to \forall n\,(\varphi_{22})) \to \mathrm{FLim}(f,c,\alpha) \to \mathrm{Conv}(a,c) \to \varphi_{32}$
\item $\varphi_{99} :\equiv \varphi_{22} \land \varepsilon \in \mathbb{R} \land 0 < \varepsilon \land (\delta \in \mathbb{R} \land 0 < \delta \land \varphi_{29}) \land (N \in \mathbb{N} \land \forall n \in \mathbb{N}\,(N \leq n \to \lvert a_{n} - c \rvert < \delta))$
\item $\varphi_{100} :\equiv \forall n \in \mathbb{N}\,(a_{n} \in \mathrm{dom}(f) \land a_{n} \neq c) \to \varepsilon \in \mathbb{R} \to 0 < \varepsilon \to \delta \in \mathbb{R} \land 0 < \delta \land \varphi_{29} \to \varphi_{35}$
\item $\varphi_{101} :\equiv (\mathrm{FLim}(f,c,\alpha) \to \mathrm{Conv}(a,c) \to \varphi_{68}) \to \mathrm{FLim}(f,c,\alpha) \to \mathrm{Conv}(a,c) \to \varphi_{45}$
\item $\varphi_{102} :\equiv \forall n \in \mathbb{N}\,(a_{n} \in \mathrm{dom}(f) \land a_{n} \neq c) \to \varepsilon \in \mathbb{R} \to 0 < \varepsilon \to \delta \in \mathbb{R} \land 0 < \delta \land \varphi_{29} \to \varphi_{37}$
\item $\varphi_{103} :\equiv \forall n \in \mathbb{N}\,(a_{n} \in \mathrm{dom}(f) \land a_{n} \neq c) \to \varepsilon \in \mathbb{R} \to 0 < \varepsilon \to \delta \in \mathbb{R} \land 0 < \delta \land \varphi_{29} \to \varphi_{38}$
\item $\varphi_{104} :\equiv (\mathrm{FLim}(f,c,\alpha) \to \varphi_{78}) \to \mathrm{FLim}(f,c,\alpha) \to \mathrm{Conv}(a,c) \to \varphi_{19}$
\item $\varphi_{105} :\equiv \forall n \in \mathbb{N}\,(a_{n} \in \mathrm{dom}(f) \land a_{n} \neq c) \to \varepsilon \in \mathbb{R} \to 0 < \varepsilon \to \delta \in \mathbb{R} \land 0 < \delta \land \varphi_{29} \to \varphi_{43}$
\item $\varphi_{106} :\equiv \varphi_{22} \land \varepsilon \in \mathbb{R} \land 0 < \varepsilon \land (\delta \in \mathbb{R} \land 0 < \delta \land \varphi_{29}) \land \neg \exists N \in \mathbb{N}\,\forall n \in \mathbb{N}\,(N \leq n \to \lvert f \circ a(n) - \alpha \rvert < \varepsilon)$
\item $\varphi_{107} :\equiv \forall n \in \mathbb{N}\,(a_{n} \in \mathrm{dom}(f) \land a_{n} \neq c) \to \varepsilon \in \mathbb{R} \to 0 < \varepsilon \to \delta \in \mathbb{R} \land 0 < \delta \land \varphi_{29} \to \varphi_{44}$
\item $\varphi_{108} :\equiv \forall n \in \mathbb{N}\,(a_{n} \in \mathrm{dom}(f) \land a_{n} \neq c) \to \varepsilon \in \mathbb{R} \to 0 < \varepsilon \to \delta \in \mathbb{R} \land 0 < \delta \land \varphi_{29} \to \varphi_{46}$
\item $\varphi_{109} :\equiv \forall n \in \mathbb{N}\,(a_{n} \in \mathrm{dom}(f) \land a_{n} \neq c) \to \varepsilon \in \mathbb{R} \to 0 < \varepsilon \to \delta \in \mathbb{R} \land 0 < \delta \land \varphi_{29} \to \varphi_{49}$
\item $\varphi_{110} :\equiv \forall n \in \mathbb{N}\,(a_{n} \in \mathrm{dom}(f) \land a_{n} \neq c) \to \varepsilon \in \mathbb{R} \to 0 < \varepsilon \to \delta \in \mathbb{R} \land 0 < \delta \land \varphi_{29} \to \delta \in \mathbb{R} \land 0 < \delta \land \varphi_{29}$
\item $\varphi_{111} :\equiv \forall n \in \mathbb{N}\,(a_{n} \in \mathrm{dom}(f) \land a_{n} \neq c) \to \varepsilon \in \mathbb{R} \to 0 < \varepsilon \to \delta \in \mathbb{R} \land 0 < \delta \land \varphi_{29} \to \varphi_{52}$
\item $\varphi_{112} :\equiv \forall n \in \mathbb{N}\,(a_{n} \in \mathrm{dom}(f) \land a_{n} \neq c) \to \varepsilon \in \mathbb{R} \to 0 < \varepsilon \to \delta \in \mathbb{R} \land 0 < \delta \land \varphi_{29} \to \varphi_{56}$
\item $\varphi_{113} :\equiv \forall n \in \mathbb{N}\,(a_{n} \in \mathrm{dom}(f) \land a_{n} \neq c) \to \varepsilon \in \mathbb{R} \to 0 < \varepsilon \to \delta \in \mathbb{R} \land 0 < \delta \land \varphi_{29} \to \varphi_{57}$
\item $\varphi_{114} :\equiv \forall n \in \mathbb{N}\,(a_{n} \in \mathrm{dom}(f) \land a_{n} \neq c) \to \varepsilon \in \mathbb{R} \to 0 < \varepsilon \to \delta \in \mathbb{R} \land 0 < \delta \land \varphi_{29} \to \varphi_{58}$
\item $\varphi_{115} :\equiv \forall n \in \mathbb{N}\,(a_{n} \in \mathrm{dom}(f) \land a_{n} \neq c) \to \varepsilon \in \mathbb{R} \to 0 < \varepsilon \to \delta \in \mathbb{R} \land 0 < \delta \land \varphi_{29} \to \varphi_{59}$
\item $\varphi_{116} :\equiv \forall n \in \mathbb{N}\,(a_{n} \in \mathrm{dom}(f) \land a_{n} \neq c) \to \varepsilon \in \mathbb{R} \to 0 < \varepsilon \to \delta \in \mathbb{R} \land 0 < \delta \land \varphi_{29} \to \varphi_{61}$
\item $\varphi_{117} :\equiv \forall n \in \mathbb{N}\,(a_{n} \in \mathrm{dom}(f) \land a_{n} \neq c) \to \varepsilon \in \mathbb{R} \to 0 < \varepsilon \to \delta \in \mathbb{R} \land 0 < \delta \land \varphi_{29} \to \varphi_{64}$
\item $\varphi_{118} :\equiv \forall n \in \mathbb{N}\,(a_{n} \in \mathrm{dom}(f) \land a_{n} \neq c) \to \varepsilon \in \mathbb{R} \to 0 < \varepsilon \to \delta \in \mathbb{R} \land 0 < \delta \land \varphi_{29} \to \varphi_{66}$
\item $\varphi_{119} :\equiv \forall n \in \mathbb{N}\,(a_{n} \in \mathrm{dom}(f) \land a_{n} \neq c) \to \varepsilon \in \mathbb{R} \to 0 < \varepsilon \to \delta \in \mathbb{R} \land 0 < \delta \land \varphi_{29} \to \varphi_{70}$
\item $\varphi_{120} :\equiv \forall n \in \mathbb{N}\,(a_{n} \in \mathrm{dom}(f) \land a_{n} \neq c) \to \varepsilon \in \mathbb{R} \to 0 < \varepsilon \to \delta \in \mathbb{R} \land 0 < \delta \land \varphi_{29} \to \varphi_{72}$
\item $\varphi_{121} :\equiv (\mathrm{FLim}(f,c,\alpha) \to \mathrm{Conv}(a,c) \to \varphi_{33}) \to \mathrm{FLim}(f,c,\alpha) \to \mathrm{Conv}(a,c) \to \varphi_{83}$
\item $\varphi_{122} :\equiv (\mathrm{FLim}(f,c,\alpha) \to \mathrm{Conv}(a,c) \to \varphi_{41}) \to \varphi_{39} \to \mathrm{FLim}(f,c,\alpha) \to \mathrm{Conv}(a,c) \to \varphi_{39}$
\item $\varphi_{123} :\equiv \forall N\,(\varphi_{106} \to \neg (N \in \mathbb{N} \land \forall n \in \mathbb{N}\,(N \leq n \to \lvert a_{n} - c \rvert < \delta)))$
\item $\varphi_{124} :\equiv (\mathrm{FLim}(f,c,\alpha) \to \mathrm{Conv}(a,c) \to \varphi_{45}) \to \varphi_{42} \to \mathrm{FLim}(f,c,\alpha) \to \mathrm{Conv}(a,c) \to \varphi_{40}$
\item $\varphi_{125} :\equiv (f \circ a(n) = f(a_{n}) \to f(a_{n}) = f \circ a(n)) \to \mathrm{FLim}(f,c,\alpha) \to \mathrm{Conv}(a,c) \to \varphi_{115}$
\item $\varphi_{126} :\equiv (\mathrm{FLim}(f,c,\alpha) \to \varphi_{20}) \to \mathrm{FLim}(f,c,\alpha) \to \mathrm{Conv}(a,c) \to \varphi_{105}$
\item $\varphi_{127} :\equiv (\mathrm{FLim}(f,c,\alpha) \to \varphi_{14}) \to (\mathrm{FLim}(f,c,\alpha) \to \varphi_{13}) \to \varphi_{95}$
\item $\varphi_{128} :\equiv (\mathrm{FLim}(f,c,\alpha) \to \mathrm{FLim}(f,c,\alpha)) \to (\mathrm{FLim}(f,c,\alpha) \to \mathrm{Conv}(a,c) \to \varphi_{24}) \to \varphi_{121}$
\item $\varphi_{129} :\equiv (\mathrm{FLim}(f,c,\alpha) \to \varphi_{25}) \to (\mathrm{FLim}(f,c,\alpha) \to \varphi_{21}) \to \varphi_{104}$
\item $\varphi_{130} :\equiv (\forall N\,(\varphi_{16}) \to \varphi_{16}) \to \mathrm{FLim}(f,c,\alpha) \to \mathrm{Conv}(a,c) \to \varphi_{120}$
\item $\varphi_{131} :\equiv (\mathrm{FLim}(f,c,\alpha) \to \mathrm{Conv}(a,c) \to \varphi_{84}) \to \mathrm{FLim}(f,c,\alpha) \to \mathrm{Conv}(a,c) \to \varphi_{90}$
\item $\varphi_{132} :\equiv (\mathrm{FLim}(f,c,\alpha) \to \mathrm{Conv}(a,c) \to \varphi_{110}) \to \mathrm{FLim}(f,c,\alpha) \to \mathrm{Conv}(a,c) \to \varphi_{84}$
\item $\varphi_{133} :\equiv (\mathrm{FLim}(f,c,\alpha) \to \mathrm{Conv}(a,c) \to \varphi_{110}) \to \mathrm{FLim}(f,c,\alpha) \to \mathrm{Conv}(a,c) \to \varphi_{85}$
\item $\varphi_{134} :\equiv (\mathrm{FLim}(f,c,\alpha) \to \mathrm{Conv}(a,c) \to \varphi_{93}) \to \varphi_{80} \to \neg (\delta \in \mathbb{R} \land 0 < \delta \land \varphi_{29})$
\item $\varphi_{135} :\equiv \forall \delta\,(\varphi_{80} \to \neg (\delta \in \mathbb{R} \land 0 < \delta \land \varphi_{29})) \to \forall \delta\,(\varphi_{80}) \to \forall \delta\,(\neg (\delta \in \mathbb{R} \land 0 < \delta \land \varphi_{29}))$
\item $\varphi_{136} :\equiv (\mathrm{FLim}(f,c,\alpha) \to \mathrm{Conv}(a,c) \to \varphi_{96}) \to \mathrm{FLim}(f,c,\alpha) \to \mathrm{Conv}(a,c) \to \varphi_{103}$
\item $\varphi_{137} :\equiv (\mathrm{FLim}(f,c,\alpha) \to \mathrm{Conv}(a,c) \to \varphi_{100}) \to \mathrm{FLim}(f,c,\alpha) \to \mathrm{Conv}(a,c) \to \varphi_{102}$
\item $\varphi_{138} :\equiv (\mathrm{FLim}(f,c,\alpha) \to \mathrm{Conv}(a,c) \to \varphi_{117}) \to \mathrm{FLim}(f,c,\alpha) \to \mathrm{Conv}(a,c) \to \varphi_{92}$
\item $\varphi_{139} :\equiv (\mathrm{FLim}(f,c,\alpha) \to \mathrm{Conv}(a,c) \to \varphi_{110}) \to \mathrm{FLim}(f,c,\alpha) \to \mathrm{Conv}(a,c) \to \varphi_{97}$
\item $\varphi_{140} :\equiv (\mathrm{FLim}(f,c,\alpha) \to \mathrm{Conv}(a,c) \to \varphi_{96}) \to \mathrm{FLim}(f,c,\alpha) \to \mathrm{Conv}(a,c) \to \varphi_{112}$
\item $\varphi_{141} :\equiv (\mathrm{FLim}(f,c,\alpha) \to \mathrm{Conv}(a,c) \to \varphi_{96}) \to \mathrm{FLim}(f,c,\alpha) \to \mathrm{Conv}(a,c) \to \varphi_{114}$
\item $\varphi_{142} :\equiv (\mathrm{FLim}(f,c,\alpha) \to \mathrm{Conv}(a,c) \to \varphi_{102}) \to \mathrm{FLim}(f,c,\alpha) \to \mathrm{Conv}(a,c) \to \varphi_{107}$
\item $\varphi_{143} :\equiv (\mathrm{FLim}(f,c,\alpha) \to \mathrm{Conv}(a,c) \to \varphi_{112}) \to \mathrm{FLim}(f,c,\alpha) \to \mathrm{Conv}(a,c) \to \varphi_{100}$
\item $\varphi_{144} :\equiv (\mathrm{FLim}(f,c,\alpha) \to \mathrm{Conv}(a,c) \to \varphi_{108}) \to \mathrm{FLim}(f,c,\alpha) \to \mathrm{Conv}(a,c) \to \varphi_{113}$
\item $\varphi_{145} :\equiv (\mathrm{FLim}(f,c,\alpha) \to \mathrm{Conv}(a,c) \to \varphi_{117}) \to \mathrm{FLim}(f,c,\alpha) \to \mathrm{Conv}(a,c) \to \varphi_{109}$
\item $\varphi_{146} :\equiv (\mathrm{FLim}(f,c,\alpha) \to \mathrm{Conv}(a,c) \to \varphi_{116}) \to \varphi_{99} \to n \in \mathbb{N} \to N \leq n \to \lvert f \circ a(n) - \alpha \rvert < \varepsilon$
\item $\varphi_{147} :\equiv \forall n\,(\varphi_{99} \to n \in \mathbb{N} \to N \leq n \to \lvert f \circ a(n) - \alpha \rvert < \varepsilon) \to \forall n\,(\varphi_{99}) \to \forall n \in \mathbb{N}\,(N \leq n \to \lvert f \circ a(n) - \alpha \rvert < \varepsilon)$
\item $\varphi_{148} :\equiv (\mathrm{FLim}(f,c,\alpha) \to \mathrm{Conv}(a,c) \to \varphi_{118}) \to \mathrm{FLim}(f,c,\alpha) \to \mathrm{Conv}(a,c) \to \varphi_{119}$
\item $\varphi_{149} :\equiv (\mathrm{FLim}(f,c,\alpha) \to \mathrm{Conv}(a,c) \to \varphi_{120}) \to \varphi_{106} \to \neg (N \in \mathbb{N} \land \forall n \in \mathbb{N}\,(N \leq n \to \lvert a_{n} - c \rvert < \delta))$
\item $\varphi_{150} :\equiv (\forall \varepsilon\,(\varphi_{22} \to \varepsilon \in \mathbb{R} \to \varphi_{31}) \to \forall \varepsilon\,(\varphi_{22}) \to \forall \varepsilon \in \mathbb{R}\,(\varphi_{31})) \to (\varphi_{22} \to \forall \varepsilon\,(\varphi_{22})) \to \mathrm{FLim}(f,c,\alpha) \to \varphi_{78}$
\item $\varphi_{151} :\equiv (\mathrm{FLim}(f,c,\alpha) \to \mathrm{Conv}(a,c) \to \varphi_{109}) \to \varphi_{36} \to \mathrm{FLim}(f,c,\alpha) \to \mathrm{Conv}(a,c) \to \varphi_{108}$
\item $\varphi_{152} :\equiv (\mathrm{Seq}(a) \to n \in \mathbb{N} \to a_{n} \in \mathbb{R}) \to (\mathrm{FLim}(f,c,\alpha) \to \varphi_{13}) \to \varphi_{136}$
\item $\varphi_{153} :\equiv (\mathrm{FLim}(f,c,\alpha) \to \mathrm{Conv}(a,c) \to \varphi_{113}) \to \varphi_{55} \to \mathrm{FLim}(f,c,\alpha) \to \mathrm{Conv}(a,c) \to \varphi_{116}$
\item $\varphi_{154} :\equiv (a_{n} \in \mathbb{R} \to c \in \mathbb{R} \to a_{n} - c \in \mathbb{R}) \to (\mathrm{FLim}(f,c,\alpha) \to \mathrm{Conv}(a,c) \to \varphi_{103}) \to \varphi_{126}$
\item $\varphi_{155} :\equiv (\mathrm{FLim}(f,c,\alpha) \to \mathrm{Conv}(a,c) \to \varphi_{97}) \to (\varphi_{29} \to \varphi_{52}) \to \mathrm{FLim}(f,c,\alpha) \to \mathrm{Conv}(a,c) \to \varphi_{111}$
\item $\varphi_{156} :\equiv (\mathrm{Conv}(a,c) \to \varphi_{30}) \to \varphi_{1} \to (\mathrm{FLim}(f,c,\alpha) \to \mathrm{Conv}(a,c) \to \varphi_{85}) \to \varphi_{131}$
\item $\varphi_{157} :\equiv (\varphi_{106} \to \forall N\,(\varphi_{106})) \to \mathrm{FLim}(f,c,\alpha) \to \mathrm{Conv}(a,c) \to \varphi_{93}$
\item $\varphi_{158} :\equiv (\varphi_{99} \to \forall n\,(\varphi_{99})) \to \mathrm{FLim}(f,c,\alpha) \to \mathrm{Conv}(a,c) \to \varphi_{118}$
\item $\varphi_{159} :\equiv (a_{n} - c \in \mathbb{R} \to a_{n} - c \neq 0 \to 0 < \lvert a_{n} - c \rvert) \to (\mathrm{FLim}(f,c,\alpha) \to \mathrm{Conv}(a,c) \to \varphi_{105}) \to \varphi_{142}$
\item $\varphi_{160} :\equiv (\mathrm{FLim}(f,c,\alpha) \to \mathrm{Conv}(a,c) \to \varphi_{103}) \to (\mathrm{FLim}(f,c,\alpha) \to \varphi_{20}) \to \varphi_{137}$
\item $\varphi_{161} :\equiv (\mathrm{FLim}(f,c,\alpha) \to \mathrm{Conv}(a,c) \to \varphi_{112}) \to (\mathrm{FLim}(f,c,\alpha) \to \mathrm{Conv}(a,c) \to \varphi_{107}) \to \varphi_{144}$
\item $\varphi_{162} :\equiv \forall \delta\,(\varphi_{80} \to \neg (\delta \in \mathbb{R} \land 0 < \delta \land \varphi_{29})) \to \varphi_{135} \to (\varphi_{80} \to \forall \delta\,(\varphi_{80})) \to \mathrm{FLim}(f,c,\alpha) \to \varphi_{77}$
\item $\varphi_{163} :\equiv \forall n\,(\varphi_{99} \to n \in \mathbb{N} \to N \leq n \to \lvert f \circ a(n) - \alpha \rvert < \varepsilon) \to \varphi_{147} \to \varphi_{158}$
\item $\varphi_{164} :\equiv (\mathrm{FLim}(f,c,\alpha) \to \mathrm{Conv}(a,c) \to \varphi_{90}) \to \varphi_{123} \to (\varphi_{123} \to \forall N\,(\varphi_{106}) \to \varphi_{8}) \to \varphi_{157}$
\end{abbrevlist}
\end{fullproof}

% 自動生成 (prover/build.js): FLimComp  ノード数 446
\begin{fullproof}{定理1.11}{$\forall f\,\forall g\,\forall c\,\forall b\,(\mathrm{FLim}(f,c,b) \to \mathrm{FLim}(g,b,g(b)) \to b \in \mathrm{dom}(g) \to \mathrm{PNbhd}(g \circ f,c) \to \mathrm{FLim}(g \circ f,c,g(b)))$}{446}\phantomsection\label{tree:FLimComp}
\noindent$\Pi_{1}$:
\begin{ptree}
\AXC{$z = w \to \lvert z - g(b) \rvert < \varepsilon \to \lvert w - g(b) \rvert < \varepsilon$\lbl{A6}}
\UIC{$\forall w\,(z = w \to \lvert z - g(b) \rvert < \varepsilon \to \lvert w - g(b) \rvert < \varepsilon)$}
\UIC{$\varphi_{34}$}
\AXC{$\varphi_{34} \to \varphi_{83}$\lbl{A2}}
\BIC{$\varphi_{83}$}
\AXC{$\varphi_{83} \to g(f(x)) = g \circ f(x) \to \varphi_{71}$\lbl{A2}}
\BIC{$g(f(x)) = g \circ f(x) \to \varphi_{71}$}
\end{ptree}
\noindent$\Pi_{2}$:
\begin{ptree}
\AXC{$\forall f\,\forall c\,\forall \alpha\,(\mathrm{FLim}(f,c,\alpha) \to \mathrm{RFn}(f))$\lbl{[D2]}}
\AXC{$\varphi_{24}$\lbl{A2}}
\BIC{$\forall c\,\forall \alpha\,(\mathrm{FLim}(f,c,\alpha) \to \mathrm{RFn}(f))$}
\AXC{$\varphi_{12}$\lbl{A2}}
\BIC{$\forall \alpha\,(\mathrm{FLim}(f,c,\alpha) \to \mathrm{RFn}(f))$}
\AXC{$\varphi_{4}$\lbl{A2}}
\BIC{$\mathrm{FLim}(f,c,b) \to \mathrm{RFn}(f)$}
\end{ptree}
\noindent$\Pi_{3}$:
\begin{ptree}
\AXC{$\stackrel{\Pi_{2}}{\vdots}$}\noLine
\UIC{$\mathrm{FLim}(f,c,b) \to \mathrm{RFn}(f)$}
\AXC{$(\mathrm{FLim}(f,c,b) \to \mathrm{RFn}(f)) \to \varphi_{93}$\lbl{A1}}
\BIC{$\varphi_{93}$}
\end{ptree}
\noindent$\Pi_{4}$:
\begin{ptree}
\AXC{$\varphi_{6}$\lbl{A1}}
\AXC{$\forall f\,\forall c\,\forall \alpha\,(\mathrm{FLim}(f,c,\alpha) \to \mathrm{RFn}(f))$\lbl{[D2]}}
\AXC{$\varphi_{25}$\lbl{A2}}
\BIC{$\forall c\,\forall \alpha\,(\mathrm{FLim}(g,c,\alpha) \to \mathrm{RFn}(g))$}
\AXC{$\varphi_{13}$\lbl{A2}}
\BIC{$\forall \alpha\,(\mathrm{FLim}(g,b,\alpha) \to \mathrm{RFn}(g))$}
\AXC{$\varphi_{10}$\lbl{A2}}
\BIC{$\mathrm{FLim}(g,b,g(b)) \to \mathrm{RFn}(g)$}
\AXC{$\varphi_{121}$\lbl{A1}}
\BIC{$\varphi_{6} \to \mathrm{FLim}(f,c,b) \to \varphi_{29}$}
\BIC{$\mathrm{FLim}(f,c,b) \to \varphi_{29}$}
\end{ptree}
\noindent$\Pi_{5}$:
\begin{ptree}
\AXC{$\stackrel{\Pi_{4}}{\vdots}$}\noLine
\UIC{$\mathrm{FLim}(f,c,b) \to \varphi_{29}$}
\AXC{$\forall g\,\forall f\,(\varphi_{110})$\lbl{[S9]}}
\AXC{$\forall g\,\forall f\,(\varphi_{110}) \to \forall f\,(\varphi_{110})$\lbl{A2}}
\BIC{$\forall f\,(\varphi_{110})$}
\AXC{$\forall f\,(\varphi_{110}) \to \varphi_{110}$\lbl{A2}}
\BIC{$\varphi_{110}$}
\AXC{$\varphi_{110} \to \varphi_{209}$\lbl{A1}}
\BIC{$\varphi_{209}$}
\BIC{$(\mathrm{FLim}(f,c,b) \to \varphi_{28}) \to \varphi_{132}$}
\end{ptree}
\noindent$\Pi_{6}$:
\begin{ptree}
\AXC{$\mathrm{FLim}(f,c,b) \to \mathrm{FLim}(f,c,b)$\lbl{A1}}
\AXC{$\stackrel{\Pi_{3}}{\vdots}$}\noLine
\UIC{$\varphi_{93}$}
\BIC{$\mathrm{FLim}(f,c,b) \to \varphi_{28}$}
\AXC{$\stackrel{\Pi_{5}}{\vdots}$}\noLine
\UIC{$(\mathrm{FLim}(f,c,b) \to \varphi_{28}) \to \varphi_{132}$}
\BIC{$\varphi_{132}$}
\end{ptree}
\noindent$\Pi_{7}$:
\begin{ptree}
\AXC{$\stackrel{\Pi_{6}}{\vdots}$}\noLine
\UIC{$\varphi_{132}$}
\AXC{$\varphi_{132} \to \mathrm{FLim}(f,c,b) \to \mathrm{FLim}(g,b,g(b)) \to \varphi_{111}$\lbl{A1}}
\BIC{$\mathrm{FLim}(f,c,b) \to \mathrm{FLim}(g,b,g(b)) \to \varphi_{111}$}
\AXC{$\varphi_{254}$\lbl{A1}}
\BIC{$\varphi_{217}$}
\end{ptree}
\noindent$\Pi_{8}$:
\begin{ptree}
\AXC{$\varphi_{219}$\lbl{A1}}
\AXC{$\forall x \in \mathrm{dom}(g \circ f)\,(\varphi_{52}) \to x \in \mathrm{dom}(g \circ f) \to \varphi_{52}$\lbl{A2}}
\AXC{$\stackrel{\Pi_{7}}{\vdots}$}\noLine
\UIC{$\varphi_{217}$}
\BIC{$\mathrm{FLim}(f,c,b) \to \mathrm{FLim}(g,b,g(b)) \to \varphi_{109}$}
\AXC{$\varphi_{278}$\lbl{A1}}
\BIC{$\varphi_{219} \to \varphi_{242}$}
\BIC{$\varphi_{242}$}
\end{ptree}
\noindent$\Pi_{9}$:
\begin{ptree}
\AXC{$\varphi_{72} \to \varphi_{66}$\lbl{A2}}
\AXC{$\mathrm{FLim}(f,c,b) \to \mathrm{FLim}(g,b,g(b)) \to \varphi_{229}$\lbl{A1}}
\AXC{$\varphi_{270}$\lbl{A1}}
\BIC{$\mathrm{FLim}(f,c,b) \to \varphi_{226}$}
\AXC{$\varphi_{277}$\lbl{A1}}
\BIC{$(\varphi_{72} \to \varphi_{66}) \to \mathrm{FLim}(f,c,b) \to \varphi_{223}$}
\BIC{$\mathrm{FLim}(f,c,b) \to \varphi_{223}$}
\end{ptree}
\noindent$\Pi_{10}$:
\begin{ptree}
\AXC{$\varphi_{220}$\lbl{A1}}
\AXC{$\stackrel{\Pi_{8}}{\vdots}$}\noLine
\UIC{$\varphi_{242}$}
\AXC{$\stackrel{\Pi_{9}}{\vdots}$}\noLine
\UIC{$\mathrm{FLim}(f,c,b) \to \varphi_{223}$}
\AXC{$(\mathrm{FLim}(f,c,b) \to \varphi_{223}) \to \varphi_{242} \to \varphi_{281}$\lbl{A1}}
\BIC{$\varphi_{242} \to \varphi_{281}$}
\BIC{$\varphi_{281}$}
\BIC{$(\mathrm{FLim}(f,c,b) \to \varphi_{224}) \to \mathrm{FLim}(f,c,b) \to \varphi_{230}$}
\end{ptree}
\noindent$\Pi_{11}$:
\begin{ptree}
\AXC{$\mathrm{FLim}(f,c,b) \to \varphi_{224}$\lbl{A1}}
\AXC{$\stackrel{\Pi_{10}}{\vdots}$}\noLine
\UIC{$(\mathrm{FLim}(f,c,b) \to \varphi_{224}) \to \mathrm{FLim}(f,c,b) \to \varphi_{230}$}
\BIC{$\mathrm{FLim}(f,c,b) \to \varphi_{230}$}
\end{ptree}
\noindent$\Pi_{12}$:
\begin{ptree}
\AXC{$\varphi_{82} \to \varphi_{75}$\lbl{A2}}
\AXC{$\mathrm{FLim}(f,c,b) \to \varphi_{192}$\lbl{A1}}
\AXC{$(\mathrm{FLim}(f,c,b) \to \varphi_{192}) \to \varphi_{191}$\lbl{A1}}
\BIC{$\varphi_{191}$}
\AXC{$\varphi_{191} \to (\varphi_{82} \to \varphi_{75}) \to \varphi_{189}$\lbl{A1}}
\BIC{$(\varphi_{82} \to \varphi_{75}) \to \varphi_{189}$}
\BIC{$\varphi_{189}$}
\end{ptree}
\noindent$\Pi_{13}$:
\begin{ptree}
\AXC{$z = w \to z < \varepsilon \to w < \varepsilon$\lbl{A6}}
\UIC{$\forall w\,(z = w \to z < \varepsilon \to w < \varepsilon)$}
\UIC{$\forall z\,\forall w\,(z = w \to z < \varepsilon \to w < \varepsilon)$}
\AXC{$\varphi_{26}$\lbl{A2}}
\BIC{$\forall w\,(0 = w \to 0 < \varepsilon \to w < \varepsilon)$}
\AXC{$\forall w\,(0 = w \to 0 < \varepsilon \to w < \varepsilon) \to \varphi_{58}$\lbl{A2}}
\BIC{$\varphi_{58}$}
\end{ptree}
\noindent$\Pi_{14}$:
\begin{ptree}
\AXC{$x = x$\lbl{A5}}
\AXC{$x = y \to x = x \to y = x$\lbl{A6}}
\AXC{$(x = y \to x = x \to y = x) \to x = x \to x = y \to y = x$\lbl{A1}}
\BIC{$x = x \to x = y \to y = x$}
\BIC{$x = y \to y = x$}
\UIC{$\forall y\,(x = y \to y = x)$}
\UIC{$\forall x\,\forall y\,(x = y \to y = x)$}
\end{ptree}
\noindent$\Pi_{15}$:
\begin{ptree}
\AXC{$\stackrel{\Pi_{14}}{\vdots}$}\noLine
\UIC{$\forall x\,\forall y\,(x = y \to y = x)$}
\AXC{$\forall x\,\forall y\,(x = y \to y = x) \to \forall y\,(v_{420} = y \to y = v_{420})$\lbl{A2}}
\BIC{$\forall y\,(v_{420} = y \to y = v_{420})$}
\AXC{$\varphi_{1}$\lbl{A2}}
\BIC{$v_{420} = v_{421} \to v_{421} = v_{420}$}
\UIC{$\forall v_{421}\,(v_{420} = v_{421} \to v_{421} = v_{420})$}
\UIC{$\forall v_{420}\,\forall v_{421}\,(v_{420} = v_{421} \to v_{421} = v_{420})$}
\end{ptree}
\noindent$\Pi_{16}$:
\begin{ptree}
\AXC{$\stackrel{\Pi_{15}}{\vdots}$}\noLine
\UIC{$\forall v_{420}\,\forall v_{421}\,(v_{420} = v_{421} \to v_{421} = v_{420})$}
\AXC{$\forall v_{420}\,\forall v_{421}\,(v_{420} = v_{421} \to v_{421} = v_{420}) \to \varphi_{50}$\lbl{A2}}
\BIC{$\varphi_{50}$}
\end{ptree}
\noindent$\Pi_{17}$:
\begin{ptree}
\AXC{$y = z \to x = y \to x = z$\lbl{A6}}
\AXC{$(y = z \to x = y \to x = z) \to x = y \to y = z \to x = z$\lbl{A1}}
\BIC{$x = y \to y = z \to x = z$}
\UIC{$\forall z\,(x = y \to y = z \to x = z)$}
\UIC{$\forall y\,\forall z\,(x = y \to y = z \to x = z)$}
\UIC{$\forall x\,\forall y\,\forall z\,(x = y \to y = z \to x = z)$}
\AXC{$\varphi_{7}$\lbl{A2}}
\BIC{$\forall y\,\forall z\,(v_{417} = y \to y = z \to v_{417} = z)$}
\AXC{$\varphi_{19}$\lbl{A2}}
\BIC{$\forall z\,(v_{417} = v_{418} \to v_{418} = z \to v_{417} = z)$}
\AXC{$\varphi_{37}$\lbl{A2}}
\BIC{$v_{417} = v_{418} \to v_{418} = v_{419} \to v_{417} = v_{419}$}
\UIC{$\forall v_{419}\,(v_{417} = v_{418} \to v_{418} = v_{419} \to v_{417} = v_{419})$}
\UIC{$\varphi_{2}$}
\UIC{$\forall v_{417}\,(\varphi_{2})$}
\end{ptree}
\noindent$\Pi_{18}$:
\begin{ptree}
\AXC{$\forall x\,(0 \leq x \to \lvert x \rvert = x)$\lbl{[A1]}}
\AXC{$\forall x\,(0 \leq x \to \lvert x \rvert = x) \to 0 \leq 0 \to \lvert 0 \rvert = 0$\lbl{A2}}
\BIC{$0 \leq 0 \to \lvert 0 \rvert = 0$}
\AXC{$(0 \leq 0 \to \lvert 0 \rvert = 0) \to \varphi_{256}$\lbl{A1}}
\BIC{$\varphi_{256}$}
\end{ptree}
\noindent$\Pi_{19}$:
\begin{ptree}
\AXC{$y = z \to x = y \to x = z$\lbl{A6}}
\AXC{$(y = z \to x = y \to x = z) \to x = y \to y = z \to x = z$\lbl{A1}}
\BIC{$x = y \to y = z \to x = z$}
\UIC{$\forall z\,(x = y \to y = z \to x = z)$}
\UIC{$\forall y\,\forall z\,(x = y \to y = z \to x = z)$}
\UIC{$\forall x\,\forall y\,\forall z\,(x = y \to y = z \to x = z)$}
\AXC{$\varphi_{8}$\lbl{A2}}
\BIC{$\forall y\,\forall z\,(v_{414} = y \to y = z \to v_{414} = z)$}
\AXC{$\varphi_{20}$\lbl{A2}}
\BIC{$\forall z\,(v_{414} = v_{415} \to v_{415} = z \to v_{414} = z)$}
\AXC{$\varphi_{38}$\lbl{A2}}
\BIC{$v_{414} = v_{415} \to v_{415} = v_{416} \to v_{414} = v_{416}$}
\UIC{$\forall v_{416}\,(v_{414} = v_{415} \to v_{415} = v_{416} \to v_{414} = v_{416})$}
\UIC{$\varphi_{3}$}
\UIC{$\forall v_{414}\,(\varphi_{3})$}
\end{ptree}
\noindent$\Pi_{20}$:
\begin{ptree}
\AXC{$\stackrel{\Pi_{19}}{\vdots}$}\noLine
\UIC{$\forall v_{414}\,(\varphi_{3})$}
\AXC{$\forall v_{414}\,(\varphi_{3}) \to \varphi_{78}$\lbl{A2}}
\BIC{$\varphi_{78}$}
\AXC{$\varphi_{78} \to \varphi_{108}$\lbl{A2}}
\BIC{$\varphi_{108}$}
\AXC{$\varphi_{108} \to \lvert g(y) - g(b) \rvert = \lvert g(b) - g(b) \rvert \to \varphi_{41}$\lbl{A2}}
\BIC{$\lvert g(y) - g(b) \rvert = \lvert g(b) - g(b) \rvert \to \varphi_{41}$}
\end{ptree}
\noindent$\Pi_{21}$:
\begin{ptree}
\AXC{$z = w \to \varphi_{42}$\lbl{A6}}
\UIC{$\forall w\,(z = w \to \varphi_{42})$}
\UIC{$\forall z\,\forall w\,(z = w \to \varphi_{42})$}
\AXC{$\forall z\,\forall w\,(z = w \to \varphi_{42}) \to \varphi_{104}$\lbl{A2}}
\BIC{$\varphi_{104}$}
\AXC{$\varphi_{104} \to g(b) - g(b) = 0 \to \varphi_{76}$\lbl{A2}}
\BIC{$g(b) - g(b) = 0 \to \varphi_{76}$}
\end{ptree}
\noindent$\Pi_{22}$:
\begin{ptree}
\AXC{$\forall f\,\forall x\,(\mathrm{RFn}(f) \to x \in \mathrm{dom}(f) \to f(x) \in \mathbb{R})$\lbl{[S6]}}
\AXC{$\varphi_{84}$\lbl{A2}}
\BIC{$\forall x\,(\mathrm{RFn}(g) \to x \in \mathrm{dom}(g) \to g(x) \in \mathbb{R})$}
\AXC{$\varphi_{80}$\lbl{A2}}
\BIC{$\mathrm{RFn}(g) \to b \in \mathrm{dom}(g) \to g(b) \in \mathbb{R}$}
\end{ptree}
\noindent$\Pi_{23}$:
\begin{ptree}
\AXC{$\stackrel{\Pi_{22}}{\vdots}$}\noLine
\UIC{$\mathrm{RFn}(g) \to b \in \mathrm{dom}(g) \to g(b) \in \mathbb{R}$}
\AXC{$(\mathrm{RFn}(g) \to b \in \mathrm{dom}(g) \to g(b) \in \mathbb{R}) \to \varphi_{154}$\lbl{A1}}
\BIC{$\varphi_{154}$}
\end{ptree}
\noindent$\Pi_{24}$:
\begin{ptree}
\AXC{$\mathrm{FLim}(f,c,b) \to \varphi_{5}$\lbl{A1}}
\AXC{$\stackrel{\Pi_{4}}{\vdots}$}\noLine
\UIC{$\mathrm{FLim}(f,c,b) \to \varphi_{29}$}
\AXC{$\stackrel{\Pi_{23}}{\vdots}$}\noLine
\UIC{$\varphi_{154}$}
\BIC{$(\mathrm{FLim}(f,c,b) \to \varphi_{5}) \to \mathrm{FLim}(f,c,b) \to \varphi_{62}$}
\BIC{$\mathrm{FLim}(f,c,b) \to \varphi_{62}$}
\end{ptree}
\noindent$\Pi_{25}$:
\begin{ptree}
\AXC{$\stackrel{\Pi_{24}}{\vdots}$}\noLine
\UIC{$\mathrm{FLim}(f,c,b) \to \varphi_{62}$}
\AXC{$\forall x \in \mathbb{R}\,(x - x = 0)$\lbl{[F13]}}
\AXC{$\varphi_{51}$\lbl{A2}}
\BIC{$g(b) \in \mathbb{R} \to g(b) - g(b) = 0$}
\AXC{$\varphi_{236}$\lbl{A1}}
\BIC{$(\mathrm{FLim}(f,c,b) \to \varphi_{62}) \to \varphi_{183}$}
\BIC{$\varphi_{183}$}
\AXC{$\varphi_{183} \to \varphi_{247}$\lbl{A1}}
\BIC{$\varphi_{247}$}
\end{ptree}
\noindent$\Pi_{26}$:
\begin{ptree}
\AXC{$\stackrel{\Pi_{14}}{\vdots}$}\noLine
\UIC{$\forall x\,\forall y\,(x = y \to y = x)$}
\AXC{$\forall x\,\forall y\,(x = y \to y = x) \to \forall y\,(0 = y \to y = 0)$\lbl{A2}}
\BIC{$\forall y\,(0 = y \to y = 0)$}
\end{ptree}
\noindent$\Pi_{27}$:
\begin{ptree}
\AXC{$\stackrel{\Pi_{26}}{\vdots}$}\noLine
\UIC{$\forall y\,(0 = y \to y = 0)$}
\AXC{$\forall y\,(0 = y \to y = 0) \to 0 = \lvert y - b \rvert \to \lvert y - b \rvert = 0$\lbl{A2}}
\BIC{$0 = \lvert y - b \rvert \to \lvert y - b \rvert = 0$}
\end{ptree}
\noindent$\Pi_{28}$:
\begin{ptree}
\AXC{$\forall f\,\forall c\,\forall \alpha\,(\mathrm{FLim}(f,c,\alpha) \to \alpha \in \mathbb{R})$\lbl{[D4]}}
\AXC{$\varphi_{57}$\lbl{A2}}
\BIC{$\forall c\,\forall \alpha\,(\mathrm{FLim}(f,c,\alpha) \to \alpha \in \mathbb{R})$}
\AXC{$\varphi_{39}$\lbl{A2}}
\BIC{$\forall \alpha\,(\mathrm{FLim}(f,c,\alpha) \to \alpha \in \mathbb{R})$}
\AXC{$\varphi_{15}$\lbl{A2}}
\BIC{$\mathrm{FLim}(f,c,b) \to b \in \mathbb{R}$}
\end{ptree}
\noindent$\Pi_{29}$:
\begin{ptree}
\AXC{$\mathrm{FLim}(f,c,b) \to \mathrm{FLim}(f,c,b)$\lbl{A1}}
\AXC{$\stackrel{\Pi_{28}}{\vdots}$}\noLine
\UIC{$\mathrm{FLim}(f,c,b) \to b \in \mathbb{R}$}
\AXC{$(\mathrm{FLim}(f,c,b) \to b \in \mathbb{R}) \to \varphi_{99}$\lbl{A1}}
\BIC{$\varphi_{99}$}
\BIC{$\mathrm{FLim}(f,c,b) \to \varphi_{43}$}
\end{ptree}
\noindent$\Pi_{30}$:
\begin{ptree}
\AXC{$\stackrel{\Pi_{4}}{\vdots}$}\noLine
\UIC{$\mathrm{FLim}(f,c,b) \to \varphi_{29}$}
\AXC{$\forall f\,\forall x\,(\mathrm{RFn}(f) \to x \in \mathrm{dom}(f) \to x \in \mathbb{R})$\lbl{[S7]}}
\AXC{$\varphi_{68}$\lbl{A2}}
\BIC{$\forall x\,(\mathrm{RFn}(g) \to x \in \mathrm{dom}(g) \to x \in \mathbb{R})$}
\AXC{$\varphi_{61}$\lbl{A2}}
\BIC{$\mathrm{RFn}(g) \to y \in \mathrm{dom}(g) \to y \in \mathbb{R}$}
\AXC{$\varphi_{261}$\lbl{A1}}
\BIC{$(\mathrm{FLim}(f,c,b) \to \varphi_{29}) \to \varphi_{253}$}
\BIC{$\varphi_{253}$}
\end{ptree}
\noindent$\Pi_{31}$:
\begin{ptree}
\AXC{$\mathrm{FLim}(f,c,b) \to \varphi_{144}$\lbl{A1}}
\AXC{$\stackrel{\Pi_{30}}{\vdots}$}\noLine
\UIC{$\varphi_{253}$}
\BIC{$\mathrm{FLim}(f,c,b) \to \mathrm{FLim}(g,b,g(b)) \to \varphi_{140}$}
\end{ptree}
\noindent$\Pi_{32}$:
\begin{ptree}
\AXC{$\forall x\,\forall y\,(x \in \mathbb{R} \to y \in \mathbb{R} \to x - y \in \mathbb{R})$\lbl{[F1]}}
\AXC{$\varphi_{89}$\lbl{A2}}
\BIC{$\forall y\,(v_{410} \in \mathbb{R} \to y \in \mathbb{R} \to v_{410} - y \in \mathbb{R})$}
\AXC{$\varphi_{96}$\lbl{A2}}
\BIC{$v_{410} \in \mathbb{R} \to v_{411} \in \mathbb{R} \to v_{410} - v_{411} \in \mathbb{R}$}
\UIC{$\forall v_{411}\,(v_{410} \in \mathbb{R} \to v_{411} \in \mathbb{R} \to v_{410} - v_{411} \in \mathbb{R})$}
\UIC{$\varphi_{35}$}
\AXC{$\varphi_{35} \to \forall v_{411}\,(y \in \mathbb{R} \to v_{411} \in \mathbb{R} \to y - v_{411} \in \mathbb{R})$\lbl{A2}}
\BIC{$\forall v_{411}\,(y \in \mathbb{R} \to v_{411} \in \mathbb{R} \to y - v_{411} \in \mathbb{R})$}
\end{ptree}
\noindent$\Pi_{33}$:
\begin{ptree}
\AXC{$\stackrel{\Pi_{32}}{\vdots}$}\noLine
\UIC{$\forall v_{411}\,(y \in \mathbb{R} \to v_{411} \in \mathbb{R} \to y - v_{411} \in \mathbb{R})$}
\AXC{$\varphi_{87}$\lbl{A2}}
\BIC{$y \in \mathbb{R} \to b \in \mathbb{R} \to y - b \in \mathbb{R}$}
\AXC{$(y \in \mathbb{R} \to b \in \mathbb{R} \to y - b \in \mathbb{R}) \to \varphi_{262}$\lbl{A1}}
\BIC{$\varphi_{262}$}
\end{ptree}
\noindent$\Pi_{34}$:
\begin{ptree}
\AXC{$\forall x \in \mathbb{R}\,(\lvert x \rvert = 0 \to x = 0)$\lbl{[D6]}}
\AXC{$\varphi_{63}$\lbl{A2}}
\BIC{$y - b \in \mathbb{R} \to \lvert y - b \rvert = 0 \to y - b = 0$}
\AXC{$(y - b \in \mathbb{R} \to \lvert y - b \rvert = 0 \to y - b = 0) \to \varphi_{176} \to \varphi_{175} \to \varphi_{171}$\lbl{A1}}
\BIC{$\varphi_{176} \to \varphi_{175} \to \varphi_{171}$}
\end{ptree}
\noindent$\Pi_{35}$:
\begin{ptree}
\AXC{$\stackrel{\Pi_{29}}{\vdots}$}\noLine
\UIC{$\mathrm{FLim}(f,c,b) \to \varphi_{43}$}
\AXC{$\stackrel{\Pi_{31}}{\vdots}$}\noLine
\UIC{$\mathrm{FLim}(f,c,b) \to \mathrm{FLim}(g,b,g(b)) \to \varphi_{140}$}
\AXC{$\stackrel{\Pi_{33}}{\vdots}$}\noLine
\UIC{$\varphi_{262}$}
\BIC{$(\mathrm{FLim}(f,c,b) \to \varphi_{43}) \to \varphi_{176}$}
\BIC{$\varphi_{176}$}
\AXC{$\stackrel{\Pi_{34}}{\vdots}$}\noLine
\UIC{$\varphi_{176} \to \varphi_{175} \to \varphi_{171}$}
\BIC{$\varphi_{175} \to \varphi_{171}$}
\end{ptree}
\noindent$\Pi_{36}$:
\begin{ptree}
\AXC{$\stackrel{\Pi_{27}}{\vdots}$}\noLine
\UIC{$0 = \lvert y - b \rvert \to \lvert y - b \rvert = 0$}
\AXC{$\varphi_{174}$\lbl{A1}}
\AXC{$\varphi_{174} \to (0 = \lvert y - b \rvert \to \lvert y - b \rvert = 0) \to \varphi_{175}$\lbl{A1}}
\BIC{$(0 = \lvert y - b \rvert \to \lvert y - b \rvert = 0) \to \varphi_{175}$}
\BIC{$\varphi_{175}$}
\AXC{$\stackrel{\Pi_{35}}{\vdots}$}\noLine
\UIC{$\varphi_{175} \to \varphi_{171}$}
\BIC{$\varphi_{171}$}
\end{ptree}
\noindent$\Pi_{37}$:
\begin{ptree}
\AXC{$\forall x\,\forall y\,(x \in \mathbb{R} \to y \in \mathbb{R} \to x - y = 0 \to x = y)$\lbl{[D7]}}
\AXC{$\varphi_{92}$\lbl{A2}}
\BIC{$\forall y\,(v_{412} \in \mathbb{R} \to y \in \mathbb{R} \to v_{412} - y = 0 \to v_{412} = y)$}
\AXC{$\varphi_{103}$\lbl{A2}}
\BIC{$\varphi_{16}$}
\UIC{$\forall v_{413}\,(\varphi_{16})$}
\UIC{$\forall v_{412}\,\forall v_{413}\,(\varphi_{16})$}
\AXC{$\forall v_{412}\,\forall v_{413}\,(\varphi_{16}) \to \varphi_{11}$\lbl{A2}}
\BIC{$\varphi_{11}$}
\end{ptree}
\noindent$\Pi_{38}$:
\begin{ptree}
\AXC{$\stackrel{\Pi_{37}}{\vdots}$}\noLine
\UIC{$\varphi_{11}$}
\AXC{$\varphi_{11} \to y \in \mathbb{R} \to b \in \mathbb{R} \to y - b = 0 \to y = b$\lbl{A2}}
\BIC{$y \in \mathbb{R} \to b \in \mathbb{R} \to y - b = 0 \to y = b$}
\AXC{$(y \in \mathbb{R} \to b \in \mathbb{R} \to y - b = 0 \to y = b) \to \varphi_{273}$\lbl{A1}}
\BIC{$\varphi_{273}$}
\end{ptree}
\noindent$\Pi_{39}$:
\begin{ptree}
\AXC{$\stackrel{\Pi_{36}}{\vdots}$}\noLine
\UIC{$\varphi_{171}$}
\AXC{$\stackrel{\Pi_{29}}{\vdots}$}\noLine
\UIC{$\mathrm{FLim}(f,c,b) \to \varphi_{43}$}
\AXC{$\stackrel{\Pi_{31}}{\vdots}$}\noLine
\UIC{$\mathrm{FLim}(f,c,b) \to \mathrm{FLim}(g,b,g(b)) \to \varphi_{140}$}
\AXC{$\stackrel{\Pi_{38}}{\vdots}$}\noLine
\UIC{$\varphi_{273}$}
\BIC{$\varphi_{263}$}
\BIC{$\varphi_{171} \to \mathrm{FLim}(f,c,b) \to \mathrm{FLim}(g,b,g(b)) \to \varphi_{153}$}
\BIC{$\mathrm{FLim}(f,c,b) \to \mathrm{FLim}(g,b,g(b)) \to \varphi_{153}$}
\AXC{$\varphi_{267}$\lbl{A1}}
\BIC{$\lvert g(y) - g(b) \rvert = \lvert g(y) - g(b) \rvert \to \varphi_{240}$}
\end{ptree}
\noindent$\Pi_{40}$:
\begin{ptree}
\AXC{$x = x$\lbl{A5}}
\UIC{$\forall x\,(x = x)$}
\AXC{$\varphi_{22}$\lbl{A2}}
\BIC{$\lvert g(y) - g(b) \rvert = \lvert g(y) - g(b) \rvert$}
\AXC{$\stackrel{\Pi_{39}}{\vdots}$}\noLine
\UIC{$\lvert g(y) - g(b) \rvert = \lvert g(y) - g(b) \rvert \to \varphi_{240}$}
\BIC{$\varphi_{240}$}
\end{ptree}
\noindent$\Pi_{41}$:
\begin{ptree}
\AXC{$z = w \to \varphi_{94}$\lbl{A6}}
\UIC{$\forall w\,(z = w \to \varphi_{94})$}
\UIC{$\forall z\,\forall w\,(z = w \to \varphi_{94})$}
\AXC{$\forall z\,\forall w\,(z = w \to \varphi_{94}) \to \varphi_{105}$\lbl{A2}}
\BIC{$\varphi_{105}$}
\AXC{$\varphi_{105} \to y = b \to \varphi_{95}$\lbl{A2}}
\BIC{$y = b \to \varphi_{95}$}
\AXC{$\stackrel{\Pi_{40}}{\vdots}$}\noLine
\UIC{$\varphi_{240}$}
\BIC{$\mathrm{FLim}(f,c,b) \to \mathrm{FLim}(g,b,g(b)) \to \varphi_{178}$}
\AXC{$\varphi_{279}$\lbl{A1}}
\BIC{$(\mathrm{FLim}(f,c,b) \to \varphi_{179}) \to \varphi_{235}$}
\end{ptree}
\noindent$\Pi_{42}$:
\begin{ptree}
\AXC{$\stackrel{\Pi_{21}}{\vdots}$}\noLine
\UIC{$g(b) - g(b) = 0 \to \varphi_{76}$}
\AXC{$x = x$\lbl{A5}}
\UIC{$\forall x\,(x = x)$}
\AXC{$\varphi_{21}$\lbl{A2}}
\BIC{$\lvert g(b) - g(b) \rvert = \lvert g(b) - g(b) \rvert$}
\AXC{$\stackrel{\Pi_{25}}{\vdots}$}\noLine
\UIC{$\varphi_{247}$}
\BIC{$(g(b) - g(b) = 0 \to \varphi_{76}) \to \mathrm{FLim}(f,c,b) \to \varphi_{179}$}
\BIC{$\mathrm{FLim}(f,c,b) \to \varphi_{179}$}
\AXC{$\stackrel{\Pi_{41}}{\vdots}$}\noLine
\UIC{$(\mathrm{FLim}(f,c,b) \to \varphi_{179}) \to \varphi_{235}$}
\BIC{$\varphi_{235}$}
\end{ptree}
\noindent$\Pi_{43}$:
\begin{ptree}
\AXC{$\stackrel{\Pi_{20}}{\vdots}$}\noLine
\UIC{$\lvert g(y) - g(b) \rvert = \lvert g(b) - g(b) \rvert \to \varphi_{41}$}
\AXC{$\stackrel{\Pi_{42}}{\vdots}$}\noLine
\UIC{$\varphi_{235}$}
\BIC{$\mathrm{FLim}(f,c,b) \to \varphi_{180}$}
\AXC{$(\mathrm{FLim}(f,c,b) \to \varphi_{180}) \to \varphi_{170} \to \varphi_{59} \to \varphi_{185}$\lbl{A1}}
\BIC{$\varphi_{170} \to \varphi_{59} \to \varphi_{185}$}
\end{ptree}
\noindent$\Pi_{44}$:
\begin{ptree}
\AXC{$x = x$\lbl{A5}}
\UIC{$\forall x\,(x = x)$}
\AXC{$\forall x\,(x = x) \to 0 = 0$\lbl{A2}}
\BIC{$0 = 0$}
\AXC{$\varphi_{173}$\lbl{A1}}
\BIC{$\mathrm{FLim}(f,c,b) \to \mathrm{FLim}(g,b,g(b)) \to \varphi_{152}$}
\AXC{$\stackrel{\Pi_{18}}{\vdots}$}\noLine
\UIC{$\varphi_{256}$}
\BIC{$\varphi_{170}$}
\AXC{$\stackrel{\Pi_{43}}{\vdots}$}\noLine
\UIC{$\varphi_{170} \to \varphi_{59} \to \varphi_{185}$}
\BIC{$\varphi_{59} \to \varphi_{185}$}
\end{ptree}
\noindent$\Pi_{45}$:
\begin{ptree}
\AXC{$\stackrel{\Pi_{17}}{\vdots}$}\noLine
\UIC{$\forall v_{417}\,(\varphi_{2})$}
\AXC{$\forall v_{417}\,(\varphi_{2}) \to \varphi_{77}$\lbl{A2}}
\BIC{$\varphi_{77}$}
\AXC{$\varphi_{77} \to \varphi_{74}$\lbl{A2}}
\BIC{$\varphi_{74}$}
\AXC{$\varphi_{74} \to \varphi_{59}$\lbl{A2}}
\BIC{$\varphi_{59}$}
\AXC{$\stackrel{\Pi_{44}}{\vdots}$}\noLine
\UIC{$\varphi_{59} \to \varphi_{185}$}
\BIC{$\varphi_{185}$}
\AXC{$\varphi_{185} \to \varphi_{214}$\lbl{A1}}
\BIC{$\varphi_{214}$}
\end{ptree}
\noindent$\Pi_{46}$:
\begin{ptree}
\AXC{$\stackrel{\Pi_{16}}{\vdots}$}\noLine
\UIC{$\varphi_{50}$}
\AXC{$\varphi_{50} \to \lvert g(y) - g(b) \rvert = 0 \to 0 = \lvert g(y) - g(b) \rvert$\lbl{A2}}
\BIC{$\lvert g(y) - g(b) \rvert = 0 \to 0 = \lvert g(y) - g(b) \rvert$}
\AXC{$\stackrel{\Pi_{45}}{\vdots}$}\noLine
\UIC{$\varphi_{214}$}
\BIC{$\varphi_{186}$}
\AXC{$\varphi_{186} \to \varphi_{250}$\lbl{A1}}
\BIC{$\varphi_{250}$}
\end{ptree}
\noindent$\Pi_{47}$:
\begin{ptree}
\AXC{$\forall x\,\forall y\,(x \in \mathbb{R} \to y \in \mathbb{R} \to x - y \in \mathbb{R})$\lbl{[F1]}}
\AXC{$\varphi_{90}$\lbl{A2}}
\BIC{$\forall y\,(v_{408} \in \mathbb{R} \to y \in \mathbb{R} \to v_{408} - y \in \mathbb{R})$}
\AXC{$\varphi_{97}$\lbl{A2}}
\BIC{$v_{408} \in \mathbb{R} \to v_{409} \in \mathbb{R} \to v_{408} - v_{409} \in \mathbb{R}$}
\UIC{$\forall v_{409}\,(v_{408} \in \mathbb{R} \to v_{409} \in \mathbb{R} \to v_{408} - v_{409} \in \mathbb{R})$}
\UIC{$\varphi_{36}$}
\AXC{$\varphi_{36} \to \forall v_{409}\,(y \in \mathbb{R} \to v_{409} \in \mathbb{R} \to y - v_{409} \in \mathbb{R})$\lbl{A2}}
\BIC{$\forall v_{409}\,(y \in \mathbb{R} \to v_{409} \in \mathbb{R} \to y - v_{409} \in \mathbb{R})$}
\end{ptree}
\noindent$\Pi_{48}$:
\begin{ptree}
\AXC{$\stackrel{\Pi_{47}}{\vdots}$}\noLine
\UIC{$\forall v_{409}\,(y \in \mathbb{R} \to v_{409} \in \mathbb{R} \to y - v_{409} \in \mathbb{R})$}
\AXC{$\varphi_{88}$\lbl{A2}}
\BIC{$y \in \mathbb{R} \to b \in \mathbb{R} \to y - b \in \mathbb{R}$}
\AXC{$(y \in \mathbb{R} \to b \in \mathbb{R} \to y - b \in \mathbb{R}) \to \varphi_{260}$\lbl{A1}}
\BIC{$\varphi_{260}$}
\end{ptree}
\noindent$\Pi_{49}$:
\begin{ptree}
\AXC{$\forall x \in \mathbb{R}\,(0 \leq \lvert x \rvert)$\lbl{[A6]}}
\AXC{$\forall x \in \mathbb{R}\,(0 \leq \lvert x \rvert) \to y - b \in \mathbb{R} \to 0 \leq \lvert y - b \rvert$\lbl{A2}}
\BIC{$y - b \in \mathbb{R} \to 0 \leq \lvert y - b \rvert$}
\AXC{$(y - b \in \mathbb{R} \to 0 \leq \lvert y - b \rvert) \to \varphi_{255}$\lbl{A1}}
\BIC{$\varphi_{255}$}
\end{ptree}
\noindent$\Pi_{50}$:
\begin{ptree}
\AXC{$\stackrel{\Pi_{29}}{\vdots}$}\noLine
\UIC{$\mathrm{FLim}(f,c,b) \to \varphi_{43}$}
\AXC{$\stackrel{\Pi_{31}}{\vdots}$}\noLine
\UIC{$\mathrm{FLim}(f,c,b) \to \mathrm{FLim}(g,b,g(b)) \to \varphi_{140}$}
\AXC{$\stackrel{\Pi_{48}}{\vdots}$}\noLine
\UIC{$\varphi_{260}$}
\BIC{$\varphi_{211}$}
\BIC{$\mathrm{FLim}(f,c,b) \to \mathrm{FLim}(g,b,g(b)) \to \varphi_{148}$}
\AXC{$\stackrel{\Pi_{49}}{\vdots}$}\noLine
\UIC{$\varphi_{255}$}
\BIC{$\mathrm{FLim}(f,c,b) \to \mathrm{FLim}(g,b,g(b)) \to \varphi_{145}$}
\AXC{$\varphi_{283}$\lbl{A1}}
\BIC{$\varphi_{187} \to \varphi_{189} \to (\mathrm{FLim}(f,c,b) \to \varphi_{144}) \to \varphi_{258}$}
\end{ptree}
\noindent$\Pi_{51}$:
\begin{ptree}
\AXC{$\stackrel{\Pi_{13}}{\vdots}$}\noLine
\UIC{$\varphi_{58}$}
\AXC{$\mathrm{FLim}(f,c,b) \to \mathrm{FLim}(g,b,g(b)) \to \varphi_{55}$\lbl{A1}}
\AXC{$\stackrel{\Pi_{46}}{\vdots}$}\noLine
\UIC{$\varphi_{250}$}
\BIC{$\varphi_{58} \to \varphi_{187}$}
\BIC{$\varphi_{187}$}
\AXC{$\stackrel{\Pi_{50}}{\vdots}$}\noLine
\UIC{$\varphi_{187} \to \varphi_{189} \to (\mathrm{FLim}(f,c,b) \to \varphi_{144}) \to \varphi_{258}$}
\BIC{$\varphi_{189} \to (\mathrm{FLim}(f,c,b) \to \varphi_{144}) \to \varphi_{258}$}
\end{ptree}
\noindent$\Pi_{52}$:
\begin{ptree}
\AXC{$\mathrm{FLim}(f,c,b) \to \varphi_{144}$\lbl{A1}}
\AXC{$\stackrel{\Pi_{12}}{\vdots}$}\noLine
\UIC{$\varphi_{189}$}
\AXC{$\stackrel{\Pi_{51}}{\vdots}$}\noLine
\UIC{$\varphi_{189} \to (\mathrm{FLim}(f,c,b) \to \varphi_{144}) \to \varphi_{258}$}
\BIC{$(\mathrm{FLim}(f,c,b) \to \varphi_{144}) \to \varphi_{258}$}
\BIC{$\varphi_{258}$}
\end{ptree}
\noindent$\Pi_{53}$:
\begin{ptree}
\AXC{$\mathrm{FLim}(f,c,b) \to \mathrm{FLim}(g,b,g(b)) \to \varphi_{150}$\lbl{A1}}
\AXC{$\stackrel{\Pi_{52}}{\vdots}$}\noLine
\UIC{$\varphi_{258}$}
\BIC{$\varphi_{177}$}
\AXC{$\varphi_{177} \to \varphi_{47} \land \varepsilon \in \mathbb{R} \land 0 < \varepsilon \land (\mathit{dg} \in \mathbb{R} \land 0 < \mathit{dg} \land \varphi_{82}) \to \varphi_{45}$\lbl{A1}}
\BIC{$\varphi_{47} \land \varepsilon \in \mathbb{R} \land 0 < \varepsilon \land (\mathit{dg} \in \mathbb{R} \land 0 < \mathit{dg} \land \varphi_{82}) \to \varphi_{45}$}
\UIC{$\varphi_{181}$}
\end{ptree}
\noindent$\Pi_{54}$:
\begin{ptree}
\AXC{$\varphi_{54} \to \varphi_{81}$\lbl{A2}}
\AXC{$\varphi_{238}$\lbl{A4}}
\AXC{$\varphi_{181} \to \varphi_{184}$\lbl{A3}}
\AXC{$\stackrel{\Pi_{53}}{\vdots}$}\noLine
\UIC{$\varphi_{181}$}
\AXC{$\varphi_{181} \to (\varphi_{181} \to \varphi_{184}) \to \varphi_{238} \to \varphi_{182}$\lbl{A1}}
\BIC{$(\varphi_{181} \to \varphi_{184}) \to \varphi_{238} \to \varphi_{182}$}
\BIC{$\varphi_{238} \to \varphi_{182}$}
\BIC{$\varphi_{182}$}
\AXC{$\varphi_{182} \to (\varphi_{54} \to \varphi_{81}) \to \varphi_{190}$\lbl{A1}}
\BIC{$(\varphi_{54} \to \varphi_{81}) \to \varphi_{190}$}
\BIC{$\varphi_{190}$}
\end{ptree}
\noindent$\Pi_{55}$:
\begin{ptree}
\AXC{$\stackrel{\Pi_{8}}{\vdots}$}\noLine
\UIC{$\varphi_{242}$}
\AXC{$\varphi_{242} \to \mathrm{FLim}(f,c,b) \to \varphi_{231}$\lbl{A1}}
\BIC{$\mathrm{FLim}(f,c,b) \to \varphi_{231}$}
\AXC{$(\mathrm{FLim}(f,c,b) \to \varphi_{231}) \to \varphi_{252}$\lbl{A1}}
\BIC{$\varphi_{252}$}
\end{ptree}
\noindent$\Pi_{56}$:
\begin{ptree}
\AXC{$\stackrel{\Pi_{14}}{\vdots}$}\noLine
\UIC{$\forall x\,\forall y\,(x = y \to y = x)$}
\AXC{$\varphi_{23}$\lbl{A2}}
\BIC{$\forall y\,(g \circ f(x) = y \to y = g \circ f(x))$}
\AXC{$\varphi_{102}$\lbl{A2}}
\BIC{$g \circ f(x) = g(f(x)) \to g(f(x)) = g \circ f(x)$}
\AXC{$\stackrel{\Pi_{55}}{\vdots}$}\noLine
\UIC{$\varphi_{252}$}
\BIC{$\mathrm{FLim}(f,c,b) \to \varphi_{232}$}
\AXC{$(\mathrm{FLim}(f,c,b) \to \varphi_{232}) \to \varphi_{276}$\lbl{A1}}
\BIC{$\varphi_{276}$}
\end{ptree}
\noindent$\Pi_{57}$:
\begin{ptree}
\AXC{$\stackrel{\Pi_{1}}{\vdots}$}\noLine
\UIC{$g(f(x)) = g \circ f(x) \to \varphi_{71}$}
\AXC{$\stackrel{\Pi_{11}}{\vdots}$}\noLine
\UIC{$\mathrm{FLim}(f,c,b) \to \varphi_{230}$}
\AXC{$\stackrel{\Pi_{8}}{\vdots}$}\noLine
\UIC{$\varphi_{242}$}
\AXC{$\stackrel{\Pi_{54}}{\vdots}$}\noLine
\UIC{$\varphi_{190}$}
\AXC{$\varphi_{190} \to \varphi_{242} \to \varphi_{269}$\lbl{A1}}
\BIC{$\varphi_{242} \to \varphi_{269}$}
\BIC{$\varphi_{269}$}
\BIC{$\mathrm{FLim}(f,c,b) \to \mathrm{FLim}(g,b,g(b)) \to \varphi_{225}$}
\AXC{$\stackrel{\Pi_{56}}{\vdots}$}\noLine
\UIC{$\varphi_{276}$}
\BIC{$\varphi_{257}$}
\BIC{$\mathrm{FLim}(f,c,b) \to \mathrm{FLim}(g,b,g(b)) \to \varphi_{227}$}
\end{ptree}
\noindent$\Pi_{58}$:
\begin{ptree}
\AXC{$\varphi_{272}$\lbl{A3}}
\AXC{$\stackrel{\Pi_{57}}{\vdots}$}\noLine
\UIC{$\mathrm{FLim}(f,c,b) \to \mathrm{FLim}(g,b,g(b)) \to \varphi_{227}$}
\AXC{$\varphi_{271}$\lbl{A1}}
\BIC{$\varphi_{194} \to x \in \mathrm{dom}(g \circ f) \to \varphi_{69}$}
\UIC{$\forall x\,(\varphi_{194} \to x \in \mathrm{dom}(g \circ f) \to \varphi_{69})$}
\AXC{$\forall x\,(\varphi_{194} \to x \in \mathrm{dom}(g \circ f) \to \varphi_{69}) \to \varphi_{272} \to \varphi_{280}$\lbl{A1}}
\BIC{$\varphi_{272} \to \varphi_{280}$}
\BIC{$\varphi_{280}$}
\end{ptree}
\noindent$\Pi_{59}$:
\begin{ptree}
\AXC{$\varphi_{194} \to \forall x\,(\varphi_{194})$\lbl{A4}}
\AXC{$\stackrel{\Pi_{58}}{\vdots}$}\noLine
\UIC{$\varphi_{280}$}
\BIC{$\mathrm{FLim}(f,c,b) \to \mathrm{FLim}(g,b,g(b)) \to \varphi_{228}$}
\AXC{$\mathrm{FLim}(f,c,b) \to \mathrm{FLim}(g,b,g(b)) \to \varphi_{229}$\lbl{A1}}
\AXC{$\varphi_{282}$\lbl{A1}}
\BIC{$\varphi_{274}$}
\BIC{$\varphi_{243}$}
\end{ptree}
\noindent$\Pi_{60}$:
\begin{ptree}
\AXC{$\varphi_{114} \to \neg (\delta_{2} \in \mathbb{R} \land 0 < \delta_{2} \land \forall x \in \mathrm{dom}(g \circ f)\,(\varphi_{69}))$\lbl{A2}}
\AXC{$\stackrel{\Pi_{59}}{\vdots}$}\noLine
\UIC{$\varphi_{243}$}
\AXC{$\varphi_{243} \to \varphi_{266}$\lbl{A1}}
\BIC{$\varphi_{266}$}
\BIC{$\mathrm{FLim}(f,c,b) \to \varphi_{241}$}
\AXC{$(\mathrm{FLim}(f,c,b) \to \varphi_{241}) \to \varphi_{244}$\lbl{A1}}
\BIC{$\varphi_{244}$}
\UIC{$\forall \delta_{2}\,(\varphi_{244})$}
\end{ptree}
\noindent$\Pi_{61}$:
\begin{ptree}
\AXC{$\mathrm{FLim}(f,c,b) \to \varphi_{192}$\lbl{A1}}
\AXC{$(\mathrm{FLim}(f,c,b) \to \varphi_{192}) \to \varphi_{141}$\lbl{A1}}
\BIC{$\varphi_{141}$}
\end{ptree}
\noindent$\Pi_{62}$:
\begin{ptree}
\AXC{$\forall f\,\forall c\,(\varphi_{118})$\lbl{[D1]}}
\AXC{$\forall f\,\forall c\,(\varphi_{118}) \to \forall c\,(\varphi_{118})$\lbl{A2}}
\BIC{$\forall c\,(\varphi_{118})$}
\AXC{$\forall c\,(\varphi_{118}) \to \varphi_{118}$\lbl{A2}}
\BIC{$\varphi_{118}$}
\AXC{$\varphi_{118} \to \varphi_{117}$\lbl{A2}}
\BIC{$\varphi_{117}$}
\end{ptree}
\noindent$\Pi_{63}$:
\begin{ptree}
\AXC{$\stackrel{\Pi_{62}}{\vdots}$}\noLine
\UIC{$\varphi_{117}$}
\AXC{$\varphi_{117} \to \mathrm{FLim}(f,c,b) \to \mathit{dg} \in \mathbb{R} \to 0 < \mathit{dg} \to \varphi_{101}$\lbl{A2}}
\BIC{$\mathrm{FLim}(f,c,b) \to \mathit{dg} \in \mathbb{R} \to 0 < \mathit{dg} \to \varphi_{101}$}
\AXC{$(\mathrm{FLim}(f,c,b) \to \mathit{dg} \in \mathbb{R} \to 0 < \mathit{dg} \to \varphi_{101}) \to \varphi_{268}$\lbl{A1}}
\BIC{$\varphi_{268}$}
\end{ptree}
\noindent$\Pi_{64}$:
\begin{ptree}
\AXC{$\stackrel{\Pi_{61}}{\vdots}$}\noLine
\UIC{$\varphi_{141}$}
\AXC{$\mathrm{FLim}(f,c,b) \to \mathrm{FLim}(f,c,b)$\lbl{A1}}
\AXC{$\stackrel{\Pi_{63}}{\vdots}$}\noLine
\UIC{$\varphi_{268}$}
\BIC{$\varphi_{141} \to \varphi_{137} \to \varphi_{195}$}
\BIC{$\varphi_{137} \to \varphi_{195}$}
\end{ptree}
\noindent$\Pi_{65}$:
\begin{ptree}
\AXC{$\stackrel{\Pi_{60}}{\vdots}$}\noLine
\UIC{$\forall \delta_{2}\,(\varphi_{244})$}
\AXC{$\mathrm{FLim}(f,c,b) \to \varphi_{192}$\lbl{A1}}
\AXC{$(\mathrm{FLim}(f,c,b) \to \varphi_{192}) \to \varphi_{137}$\lbl{A1}}
\BIC{$\varphi_{137}$}
\AXC{$\stackrel{\Pi_{64}}{\vdots}$}\noLine
\UIC{$\varphi_{137} \to \varphi_{195}$}
\BIC{$\varphi_{195}$}
\AXC{$\varphi_{195} \to \varphi_{287}$\lbl{A1}}
\BIC{$\varphi_{287}$}
\BIC{$(\forall \delta_{2}\,(\varphi_{244}) \to \varphi_{216} \to \forall \delta_{2}\,(\neg (\delta_{2} \in \mathbb{R} \land 0 < \delta_{2} \land \varphi_{72}))) \to \varphi_{284}$}
\end{ptree}
\noindent$\Pi_{66}$:
\begin{ptree}
\AXC{$\forall \delta_{2}\,(\varphi_{244}) \to \varphi_{216} \to \forall \delta_{2}\,(\neg (\delta_{2} \in \mathbb{R} \land 0 < \delta_{2} \land \varphi_{72}))$\lbl{A3}}
\AXC{$\stackrel{\Pi_{65}}{\vdots}$}\noLine
\UIC{$(\forall \delta_{2}\,(\varphi_{244}) \to \varphi_{216} \to \forall \delta_{2}\,(\neg (\delta_{2} \in \mathbb{R} \land 0 < \delta_{2} \land \varphi_{72}))) \to \varphi_{284}$}
\BIC{$\varphi_{284}$}
\end{ptree}
\noindent$\Pi_{67}$:
\begin{ptree}
\AXC{$\varphi_{245}$\lbl{A3}}
\AXC{$\forall \delta_{2}\,(\neg (\delta_{2} \in \mathbb{R} \land 0 < \delta_{2} \land \varphi_{72})) \to \varphi_{98}$\lbl{A2}}
\UIC{$\forall \delta\,(\forall \delta_{2}\,(\neg (\delta_{2} \in \mathbb{R} \land 0 < \delta_{2} \land \varphi_{72})) \to \varphi_{98})$}
\AXC{$\varphi_{47} \land \varepsilon \in \mathbb{R} \land 0 < \varepsilon \land (\mathit{dg} \in \mathbb{R} \land 0 < \mathit{dg} \land \varphi_{82}) \land \neg \varphi_{113} \to \varphi_{216}$\lbl{A4}}
\AXC{$\stackrel{\Pi_{66}}{\vdots}$}\noLine
\UIC{$\varphi_{284}$}
\BIC{$\forall \delta\,(\forall \delta_{2}\,(\neg (\delta_{2} \in \mathbb{R} \land 0 < \delta_{2} \land \varphi_{72})) \to \varphi_{98}) \to \varphi_{245} \to \varphi_{156} \to \varphi_{210}$}
\BIC{$\varphi_{245} \to \varphi_{156} \to \varphi_{210}$}
\BIC{$\varphi_{156} \to \varphi_{210}$}
\end{ptree}
\noindent$\Pi_{68}$:
\begin{ptree}
\AXC{$\varphi_{6}$\lbl{A1}}
\AXC{$\forall f\,\forall c\,(\varphi_{118})$\lbl{[D1]}}
\AXC{$\forall f\,\forall c\,(\varphi_{118}) \to \varphi_{122}$\lbl{A2}}
\BIC{$\varphi_{122}$}
\AXC{$\varphi_{122} \to \varphi_{119}$\lbl{A2}}
\BIC{$\varphi_{119}$}
\AXC{$\varphi_{201}$\lbl{A2}}
\BIC{$\forall \varepsilon\,(\mathrm{FLim}(g,b,g(b)) \to \varepsilon \in \mathbb{R} \to 0 < \varepsilon \to \varphi_{107})$}
\AXC{$\varphi_{207}$\lbl{A2}}
\BIC{$\mathrm{FLim}(g,b,g(b)) \to \varepsilon \in \mathbb{R} \to 0 < \varepsilon \to \varphi_{107}$}
\AXC{$\varphi_{264}$\lbl{A1}}
\BIC{$\varphi_{6} \to \varphi_{237}$}
\BIC{$\varphi_{237}$}
\end{ptree}
\noindent$\Pi_{69}$:
\begin{ptree}
\AXC{$\mathrm{FLim}(f,c,b) \to \mathrm{FLim}(g,b,g(b)) \to \varphi_{55}$\lbl{A1}}
\AXC{$\mathrm{FLim}(f,c,b) \to \mathrm{FLim}(g,b,g(b)) \to \varphi_{30}$\lbl{A1}}
\AXC{$\stackrel{\Pi_{68}}{\vdots}$}\noLine
\UIC{$\varphi_{237}$}
\BIC{$\varphi_{196}$}
\BIC{$\mathrm{FLim}(f,c,b) \to \mathrm{FLim}(g,b,g(b)) \to \varphi_{131}$}
\end{ptree}
\noindent$\Pi_{70}$:
\begin{ptree}
\AXC{$\varphi_{156}$\lbl{A4}}
\AXC{$\stackrel{\Pi_{67}}{\vdots}$}\noLine
\UIC{$\varphi_{156} \to \varphi_{210}$}
\BIC{$\varphi_{210}$}
\AXC{$\varphi_{265}$\lbl{A1}}
\BIC{$\varphi_{47} \land \varepsilon \in \mathbb{R} \land 0 < \varepsilon \land \neg \varphi_{113} \to \neg (\mathit{dg} \in \mathbb{R} \land 0 < \mathit{dg} \land \varphi_{82})$}
\UIC{$\varphi_{215}$}
\AXC{$\stackrel{\Pi_{69}}{\vdots}$}\noLine
\UIC{$\mathrm{FLim}(f,c,b) \to \mathrm{FLim}(g,b,g(b)) \to \varphi_{131}$}
\AXC{$\varphi_{286}$\lbl{A1}}
\BIC{$\varphi_{215} \to \varphi_{285}$}
\BIC{$\varphi_{285}$}
\end{ptree}
\noindent$\Pi_{71}$:
\begin{ptree}
\AXC{$\varphi_{251}$\lbl{A3}}
\AXC{$\forall \mathit{dg}\,(\neg (\mathit{dg} \in \mathbb{R} \land 0 < \mathit{dg} \land \varphi_{82})) \to \varphi_{106}$\lbl{A2}}
\UIC{$\forall \delta\,(\forall \mathit{dg}\,(\neg (\mathit{dg} \in \mathbb{R} \land 0 < \mathit{dg} \land \varphi_{82})) \to \varphi_{106})$}
\AXC{$\varphi_{248}$\lbl{A4}}
\AXC{$\varphi_{215} \to \varphi_{218}$\lbl{A3}}
\AXC{$\stackrel{\Pi_{70}}{\vdots}$}\noLine
\UIC{$\varphi_{285}$}
\BIC{$\varphi_{248} \to \forall \delta\,(\forall \mathit{dg}\,(\neg (\mathit{dg} \in \mathbb{R} \land 0 < \mathit{dg} \land \varphi_{82})) \to \varphi_{106}) \to \varphi_{275}$}
\BIC{$\forall \delta\,(\forall \mathit{dg}\,(\neg (\mathit{dg} \in \mathbb{R} \land 0 < \mathit{dg} \land \varphi_{82})) \to \varphi_{106}) \to \varphi_{275}$}
\BIC{$\varphi_{275}$}
\BIC{$\varphi_{160} \to \mathrm{FLim}(f,c,b) \to \mathrm{FLim}(g,b,g(b)) \to \varphi_{133}$}
\end{ptree}
\noindent$\Pi_{72}$:
\begin{ptree}
\AXC{$\varphi_{160}$\lbl{A4}}
\AXC{$\stackrel{\Pi_{71}}{\vdots}$}\noLine
\UIC{$\varphi_{160} \to \mathrm{FLim}(f,c,b) \to \mathrm{FLim}(g,b,g(b)) \to \varphi_{133}$}
\BIC{$\mathrm{FLim}(f,c,b) \to \mathrm{FLim}(g,b,g(b)) \to \varphi_{133}$}
\AXC{$\varphi_{246}$\lbl{A1}}
\BIC{$\varphi_{47} \to \varepsilon \in \mathbb{R} \to 0 < \varepsilon \to \varphi_{113}$}
\UIC{$\forall \varepsilon\,(\varphi_{47} \to \varepsilon \in \mathbb{R} \to 0 < \varepsilon \to \varphi_{113})$}
\AXC{$\forall \varepsilon\,(\varphi_{47} \to \varepsilon \in \mathbb{R} \to 0 < \varepsilon \to \varphi_{113}) \to \varphi_{249} \to \varphi_{221}$\lbl{A1}}
\BIC{$\varphi_{249} \to \varphi_{221}$}
\end{ptree}
\noindent$\Pi_{73}$:
\begin{ptree}
\AXC{$\forall f\,\forall c\,\forall \alpha\,(\mathrm{FLim}(f,c,\alpha) \to c \in \mathbb{R})$\lbl{[D3]}}
\AXC{$\varphi_{53}$\lbl{A2}}
\BIC{$\forall c\,\forall \alpha\,(\mathrm{FLim}(f,c,\alpha) \to c \in \mathbb{R})$}
\AXC{$\varphi_{31}$\lbl{A2}}
\BIC{$\forall \alpha\,(\mathrm{FLim}(f,c,\alpha) \to c \in \mathbb{R})$}
\AXC{$\varphi_{14}$\lbl{A2}}
\BIC{$\mathrm{FLim}(f,c,b) \to c \in \mathbb{R}$}
\end{ptree}
\noindent$\Pi_{74}$:
\begin{ptree}
\AXC{$\mathrm{FLim}(f,c,b) \to \mathrm{FLim}(f,c,b)$\lbl{A1}}
\AXC{$\stackrel{\Pi_{73}}{\vdots}$}\noLine
\UIC{$\mathrm{FLim}(f,c,b) \to c \in \mathbb{R}$}
\AXC{$(\mathrm{FLim}(f,c,b) \to c \in \mathbb{R}) \to \varphi_{100}$\lbl{A1}}
\BIC{$\varphi_{100}$}
\BIC{$\mathrm{FLim}(f,c,b) \to \varphi_{46}$}
\end{ptree}
\noindent$\Pi_{75}$:
\begin{ptree}
\AXC{$\varphi_{135}$\lbl{DefFLim}}
\AXC{$\varphi_{135} \to \forall c\,(\varphi_{165})$\lbl{A2}}
\BIC{$\forall c\,(\varphi_{165})$}
\AXC{$\forall c\,(\varphi_{165}) \to \varphi_{165}$\lbl{A2}}
\BIC{$\varphi_{165}$}
\AXC{$\varphi_{165} \to \varphi_{169}$\lbl{A2}}
\BIC{$\varphi_{169}$}
\AXC{$\varphi_{169} \to (\mathrm{FLim}(f,c,b) \to \varphi_{56}) \to \varphi_{259}$\lbl{A1}}
\BIC{$(\mathrm{FLim}(f,c,b) \to \varphi_{56}) \to \varphi_{259}$}
\end{ptree}
\noindent$\Pi_{76}$:
\begin{ptree}
\AXC{$\stackrel{\Pi_{74}}{\vdots}$}\noLine
\UIC{$\mathrm{FLim}(f,c,b) \to \varphi_{46}$}
\AXC{$\stackrel{\Pi_{6}}{\vdots}$}\noLine
\UIC{$\varphi_{132}$}
\AXC{$\varphi_{132} \to \mathrm{FLim}(f,c,b) \to \varphi_{56}$\lbl{A1}}
\BIC{$\mathrm{FLim}(f,c,b) \to \varphi_{56}$}
\AXC{$\stackrel{\Pi_{75}}{\vdots}$}\noLine
\UIC{$(\mathrm{FLim}(f,c,b) \to \varphi_{56}) \to \varphi_{259}$}
\BIC{$\varphi_{259}$}
\BIC{$(\mathrm{FLim}(f,c,b) \to \varphi_{62}) \to \varphi_{239}$}
\end{ptree}
\noindent$\Pi_{77}$:
\begin{ptree}
\AXC{$\mathrm{FLim}(f,c,b) \to \mathrm{FLim}(g,b,g(b)) \to \varphi_{9}$\lbl{A1}}
\AXC{$\stackrel{\Pi_{24}}{\vdots}$}\noLine
\UIC{$\mathrm{FLim}(f,c,b) \to \varphi_{62}$}
\AXC{$\stackrel{\Pi_{76}}{\vdots}$}\noLine
\UIC{$(\mathrm{FLim}(f,c,b) \to \varphi_{62}) \to \varphi_{239}$}
\BIC{$\varphi_{239}$}
\BIC{$\varphi_{204}$}
\end{ptree}
\noindent 主木:
\begin{ptree}
\AXC{$\varphi_{47} \to \forall \varepsilon\,(\varphi_{47})$\lbl{A4}}
\AXC{$\varphi_{249}$\lbl{A3}}
\AXC{$\stackrel{\Pi_{72}}{\vdots}$}\noLine
\UIC{$\varphi_{249} \to \varphi_{221}$}
\BIC{$\varphi_{221}$}
\BIC{$\mathrm{FLim}(f,c,b) \to \mathrm{FLim}(g,b,g(b)) \to \varphi_{134}$}
\AXC{$\stackrel{\Pi_{77}}{\vdots}$}\noLine
\UIC{$\varphi_{204}$}
\BIC{$\mathrm{FLim}(f,c,b) \to \mathrm{FLim}(g,b,g(b)) \to \varphi_{27}$}
\UIC{$\varphi_{91}$}
\UIC{$\forall c\,(\varphi_{91})$}
\UIC{$\forall g\,\forall c\,(\varphi_{91})$}
\UIC{$\forall f\,\forall g\,\forall c\,(\varphi_{91})$}
\end{ptree}
\begin{abbrevlist}
\item $\varphi_{1} :\equiv \forall y\,(v_{420} = y \to y = v_{420}) \to v_{420} = v_{421} \to v_{421} = v_{420}$
\item $\varphi_{2} :\equiv \forall v_{418}\,\forall v_{419}\,(v_{417} = v_{418} \to v_{418} = v_{419} \to v_{417} = v_{419})$
\item $\varphi_{3} :\equiv \forall v_{415}\,\forall v_{416}\,(v_{414} = v_{415} \to v_{415} = v_{416} \to v_{414} = v_{416})$
\item $\varphi_{4} :\equiv \forall \alpha\,(\mathrm{FLim}(f,c,\alpha) \to \mathrm{RFn}(f)) \to \mathrm{FLim}(f,c,b) \to \mathrm{RFn}(f)$
\item $\varphi_{5} :\equiv \mathrm{FLim}(g,b,g(b)) \to b \in \mathrm{dom}(g) \to b \in \mathrm{dom}(g)$
\item $\varphi_{6} :\equiv \mathrm{FLim}(f,c,b) \to \mathrm{FLim}(g,b,g(b)) \to \mathrm{FLim}(g,b,g(b))$
\item $\varphi_{7} :\equiv \forall x\,\forall y\,\forall z\,(x = y \to y = z \to x = z) \to \forall y\,\forall z\,(v_{417} = y \to y = z \to v_{417} = z)$
\item $\varphi_{8} :\equiv \forall x\,\forall y\,\forall z\,(x = y \to y = z \to x = z) \to \forall y\,\forall z\,(v_{414} = y \to y = z \to v_{414} = z)$
\item $\varphi_{9} :\equiv b \in \mathrm{dom}(g) \to \mathrm{PNbhd}(g \circ f,c) \to \mathrm{PNbhd}(g \circ f,c)$
\item $\varphi_{10} :\equiv \forall \alpha\,(\mathrm{FLim}(g,b,\alpha) \to \mathrm{RFn}(g)) \to \mathrm{FLim}(g,b,g(b)) \to \mathrm{RFn}(g)$
\item $\varphi_{11} :\equiv \forall v_{413}\,(y \in \mathbb{R} \to v_{413} \in \mathbb{R} \to y - v_{413} = 0 \to y = v_{413})$
\item $\varphi_{12} :\equiv \forall c\,\forall \alpha\,(\mathrm{FLim}(f,c,\alpha) \to \mathrm{RFn}(f)) \to \forall \alpha\,(\mathrm{FLim}(f,c,\alpha) \to \mathrm{RFn}(f))$
\item $\varphi_{13} :\equiv \forall c\,\forall \alpha\,(\mathrm{FLim}(g,c,\alpha) \to \mathrm{RFn}(g)) \to \forall \alpha\,(\mathrm{FLim}(g,b,\alpha) \to \mathrm{RFn}(g))$
\item $\varphi_{14} :\equiv \forall \alpha\,(\mathrm{FLim}(f,c,\alpha) \to c \in \mathbb{R}) \to \mathrm{FLim}(f,c,b) \to c \in \mathbb{R}$
\item $\varphi_{15} :\equiv \forall \alpha\,(\mathrm{FLim}(f,c,\alpha) \to \alpha \in \mathbb{R}) \to \mathrm{FLim}(f,c,b) \to b \in \mathbb{R}$
\item $\varphi_{16} :\equiv v_{412} \in \mathbb{R} \to v_{413} \in \mathbb{R} \to v_{412} - v_{413} = 0 \to v_{412} = v_{413}$
\item $\varphi_{17} :\equiv y \in \mathrm{dom}(g) \to \lvert y - b \rvert < \mathit{dg} \to 0 = \lvert y - b \rvert \to \lvert 0 \rvert = 0$
\item $\varphi_{18} :\equiv y \in \mathrm{dom}(g) \to \lvert y - b \rvert < \mathit{dg} \to 0 = \lvert y - b \rvert \to y - b = 0$
\item $\varphi_{19} :\equiv \forall y\,\forall z\,(v_{417} = y \to y = z \to v_{417} = z) \to \forall z\,(v_{417} = v_{418} \to v_{418} = z \to v_{417} = z)$
\item $\varphi_{20} :\equiv \forall y\,\forall z\,(v_{414} = y \to y = z \to v_{414} = z) \to \forall z\,(v_{414} = v_{415} \to v_{415} = z \to v_{414} = z)$
\item $\varphi_{21} :\equiv \forall x\,(x = x) \to \lvert g(b) - g(b) \rvert = \lvert g(b) - g(b) \rvert$
\item $\varphi_{22} :\equiv \forall x\,(x = x) \to \lvert g(y) - g(b) \rvert = \lvert g(y) - g(b) \rvert$
\item $\varphi_{23} :\equiv \forall x\,\forall y\,(x = y \to y = x) \to \forall y\,(g \circ f(x) = y \to y = g \circ f(x))$
\item $\varphi_{24} :\equiv \forall f\,\forall c\,\forall \alpha\,(\mathrm{FLim}(f,c,\alpha) \to \mathrm{RFn}(f)) \to \forall c\,\forall \alpha\,(\mathrm{FLim}(f,c,\alpha) \to \mathrm{RFn}(f))$
\item $\varphi_{25} :\equiv \forall f\,\forall c\,\forall \alpha\,(\mathrm{FLim}(f,c,\alpha) \to \mathrm{RFn}(f)) \to \forall c\,\forall \alpha\,(\mathrm{FLim}(g,c,\alpha) \to \mathrm{RFn}(g))$
\item $\varphi_{26} :\equiv \forall z\,\forall w\,(z = w \to z < \varepsilon \to w < \varepsilon) \to \forall w\,(0 = w \to 0 < \varepsilon \to w < \varepsilon)$
\item $\varphi_{27} :\equiv b \in \mathrm{dom}(g) \to \mathrm{PNbhd}(g \circ f,c) \to \mathrm{FLim}(g \circ f,c,g(b))$
\item $\varphi_{28} :\equiv \mathrm{FLim}(g,b,g(b)) \to b \in \mathrm{dom}(g) \to \mathrm{PNbhd}(g \circ f,c) \to \mathrm{RFn}(f)$
\item $\varphi_{29} :\equiv \mathrm{FLim}(g,b,g(b)) \to b \in \mathrm{dom}(g) \to \mathrm{PNbhd}(g \circ f,c) \to \mathrm{RFn}(g)$
\item $\varphi_{30} :\equiv b \in \mathrm{dom}(g) \to \mathrm{PNbhd}(g \circ f,c) \to \varepsilon \in \mathbb{R} \to \varepsilon \in \mathbb{R}$
\item $\varphi_{31} :\equiv \forall c\,\forall \alpha\,(\mathrm{FLim}(f,c,\alpha) \to c \in \mathbb{R}) \to \forall \alpha\,(\mathrm{FLim}(f,c,\alpha) \to c \in \mathbb{R})$
\item $\varphi_{32} :\equiv y \in \mathrm{dom}(g) \to \lvert y - b \rvert < \mathit{dg} \to 0 = \lvert y - b \rvert \to 0 = \lvert y - b \rvert$
\item $\varphi_{33} :\equiv y \in \mathrm{dom}(g) \to \lvert y - b \rvert < \mathit{dg} \to 0 = \lvert y - b \rvert \to \lvert y - b \rvert = 0$
\item $\varphi_{34} :\equiv \forall z\,\forall w\,(z = w \to \lvert z - g(b) \rvert < \varepsilon \to \lvert w - g(b) \rvert < \varepsilon)$
\item $\varphi_{35} :\equiv \forall v_{410}\,\forall v_{411}\,(v_{410} \in \mathbb{R} \to v_{411} \in \mathbb{R} \to v_{410} - v_{411} \in \mathbb{R})$
\item $\varphi_{36} :\equiv \forall v_{408}\,\forall v_{409}\,(v_{408} \in \mathbb{R} \to v_{409} \in \mathbb{R} \to v_{408} - v_{409} \in \mathbb{R})$
\item $\varphi_{37} :\equiv \forall z\,(v_{417} = v_{418} \to v_{418} = z \to v_{417} = z) \to v_{417} = v_{418} \to v_{418} = v_{419} \to v_{417} = v_{419}$
\item $\varphi_{38} :\equiv \forall z\,(v_{414} = v_{415} \to v_{415} = z \to v_{414} = z) \to v_{414} = v_{415} \to v_{415} = v_{416} \to v_{414} = v_{416}$
\item $\varphi_{39} :\equiv \forall c\,\forall \alpha\,(\mathrm{FLim}(f,c,\alpha) \to \alpha \in \mathbb{R}) \to \forall \alpha\,(\mathrm{FLim}(f,c,\alpha) \to \alpha \in \mathbb{R})$
\item $\varphi_{40} :\equiv y \in \mathrm{dom}(g) \to \lvert y - b \rvert < \mathit{dg} \to 0 = \lvert y - b \rvert \to y - b \in \mathbb{R}$
\item $\varphi_{41} :\equiv \lvert g(b) - g(b) \rvert = \lvert 0 \rvert \to \lvert g(y) - g(b) \rvert = \lvert 0 \rvert$
\item $\varphi_{42} :\equiv \lvert g(b) - g(b) \rvert = \lvert z \rvert \to \lvert g(b) - g(b) \rvert = \lvert w \rvert$
\item $\varphi_{43} :\equiv \mathrm{FLim}(g,b,g(b)) \to b \in \mathrm{dom}(g) \to \mathrm{PNbhd}(g \circ f,c) \to b \in \mathbb{R}$
\item $\varphi_{44} :\equiv 0 = \lvert y - b \rvert \to \lvert g(y) - g(b) \rvert = \lvert g(b) - g(b) \rvert$
\item $\varphi_{45} :\equiv y \in \mathrm{dom}(g) \to \lvert y - b \rvert < \mathit{dg} \to \lvert g(y) - g(b) \rvert < \varepsilon$
\item $\varphi_{46} :\equiv \mathrm{FLim}(g,b,g(b)) \to b \in \mathrm{dom}(g) \to \mathrm{PNbhd}(g \circ f,c) \to c \in \mathbb{R}$
\item $\varphi_{47} :\equiv \mathrm{FLim}(f,c,b) \land \mathrm{FLim}(g,b,g(b)) \land b \in \mathrm{dom}(g) \land \mathrm{PNbhd}(g \circ f,c)$
\item $\varphi_{48} :\equiv \lvert y - b \rvert < \mathit{dg} \to 0 = \lvert y - b \rvert \to \lvert g(b) - g(b) \rvert = \lvert 0 \rvert$
\item $\varphi_{49} :\equiv \lvert y - b \rvert < \mathit{dg} \to 0 = \lvert y - b \rvert \to \lvert g(y) - g(b) \rvert = \lvert 0 \rvert$
\item $\varphi_{50} :\equiv \forall v_{421}\,(\lvert g(y) - g(b) \rvert = v_{421} \to v_{421} = \lvert g(y) - g(b) \rvert)$
\item $\varphi_{51} :\equiv \forall x \in \mathbb{R}\,(x - x = 0) \to g(b) \in \mathbb{R} \to g(b) - g(b) = 0$
\item $\varphi_{52} :\equiv x \in \mathrm{dom}(f) \land f(x) \in \mathrm{dom}(g) \land g \circ f(x) = g(f(x))$
\item $\varphi_{53} :\equiv \forall f\,\forall c\,\forall \alpha\,(\mathrm{FLim}(f,c,\alpha) \to c \in \mathbb{R}) \to \forall c\,\forall \alpha\,(\mathrm{FLim}(f,c,\alpha) \to c \in \mathbb{R})$
\item $\varphi_{54} :\equiv \forall y \in \mathrm{dom}(g)\,(\lvert y - b \rvert < \mathit{dg} \to \lvert g(y) - g(b) \rvert < \varepsilon)$
\item $\varphi_{55} :\equiv b \in \mathrm{dom}(g) \to \mathrm{PNbhd}(g \circ f,c) \to \varepsilon \in \mathbb{R} \to 0 < \varepsilon \to 0 < \varepsilon$
\item $\varphi_{56} :\equiv \mathrm{FLim}(g,b,g(b)) \to b \in \mathrm{dom}(g) \to \mathrm{PNbhd}(g \circ f,c) \to \mathrm{RFn}(g \circ f)$
\item $\varphi_{57} :\equiv \forall f\,\forall c\,\forall \alpha\,(\mathrm{FLim}(f,c,\alpha) \to \alpha \in \mathbb{R}) \to \forall c\,\forall \alpha\,(\mathrm{FLim}(f,c,\alpha) \to \alpha \in \mathbb{R})$
\item $\varphi_{58} :\equiv 0 = \lvert g(y) - g(b) \rvert \to 0 < \varepsilon \to \lvert g(y) - g(b) \rvert < \varepsilon$
\item $\varphi_{59} :\equiv \lvert g(y) - g(b) \rvert = \lvert 0 \rvert \to \lvert 0 \rvert = 0 \to \lvert g(y) - g(b) \rvert = 0$
\item $\varphi_{60} :\equiv y \in \mathrm{dom}(g) \to \lvert y - b \rvert < \mathit{dg} \to 0 = \lvert y - b \rvert \to g(b) - g(b) = 0$
\item $\varphi_{61} :\equiv \forall x\,(\mathrm{RFn}(g) \to x \in \mathrm{dom}(g) \to x \in \mathbb{R}) \to \mathrm{RFn}(g) \to y \in \mathrm{dom}(g) \to y \in \mathbb{R}$
\item $\varphi_{62} :\equiv \mathrm{FLim}(g,b,g(b)) \to b \in \mathrm{dom}(g) \to \mathrm{PNbhd}(g \circ f,c) \to g(b) \in \mathbb{R}$
\item $\varphi_{63} :\equiv \forall x \in \mathbb{R}\,(\lvert x \rvert = 0 \to x = 0) \to y - b \in \mathbb{R} \to \lvert y - b \rvert = 0 \to y - b = 0$
\item $\varphi_{64} :\equiv y \in \mathrm{dom}(g) \to \lvert y - b \rvert < \mathit{dg} \to 0 = \lvert y - b \rvert \to \lvert g(y) - g(b) \rvert = 0$
\item $\varphi_{65} :\equiv y \in \mathrm{dom}(g) \to \lvert y - b \rvert < \mathit{dg} \to 0 = \lvert y - b \rvert \to 0 = \lvert g(y) - g(b) \rvert$
\item $\varphi_{66} :\equiv x \in \mathrm{dom}(f) \to 0 < \lvert x - c \rvert \to \lvert x - c \rvert < \delta_{2} \to \lvert f(x) - b \rvert < \mathit{dg}$
\item $\varphi_{67} :\equiv 0 < \lvert x - c \rvert \to \lvert x - c \rvert < \delta_{2} \to \lvert g(f(x)) - g(b) \rvert < \varepsilon$
\item $\varphi_{68} :\equiv \forall f\,\forall x\,(\mathrm{RFn}(f) \to x \in \mathrm{dom}(f) \to x \in \mathbb{R}) \to \forall x\,(\mathrm{RFn}(g) \to x \in \mathrm{dom}(g) \to x \in \mathbb{R})$
\item $\varphi_{69} :\equiv 0 < \lvert x - c \rvert \to \lvert x - c \rvert < \delta_{2} \to \lvert g \circ f(x) - g(b) \rvert < \varepsilon$
\item $\varphi_{70} :\equiv x \in \mathrm{dom}(g \circ f) \to 0 < \lvert x - c \rvert \to \lvert x - c \rvert < \delta_{2} \to \lvert x - c \rvert < \delta_{2}$
\item $\varphi_{71} :\equiv \lvert g(f(x)) - g(b) \rvert < \varepsilon \to \lvert g \circ f(x) - g(b) \rvert < \varepsilon$
\item $\varphi_{72} :\equiv \forall x \in \mathrm{dom}(f)\,(0 < \lvert x - c \rvert \to \lvert x - c \rvert < \delta_{2} \to \lvert f(x) - b \rvert < \mathit{dg})$
\item $\varphi_{73} :\equiv y \in \mathrm{dom}(g) \to \lvert y - b \rvert < \mathit{dg} \to 0 = \lvert y - b \rvert \to \lvert g(y) - g(b) \rvert < \varepsilon$
\item $\varphi_{74} :\equiv \forall v_{419}\,(\lvert g(y) - g(b) \rvert = \lvert 0 \rvert \to \lvert 0 \rvert = v_{419} \to \lvert g(y) - g(b) \rvert = v_{419})$
\item $\varphi_{75} :\equiv y \in \mathrm{dom}(g) \to 0 < \lvert y - b \rvert \to \lvert y - b \rvert < \mathit{dg} \to \lvert g(y) - g(b) \rvert < \varepsilon$
\item $\varphi_{76} :\equiv \lvert g(b) - g(b) \rvert = \lvert g(b) - g(b) \rvert \to \lvert g(b) - g(b) \rvert = \lvert 0 \rvert$
\item $\varphi_{77} :\equiv \forall v_{418}\,\forall v_{419}\,(\lvert g(y) - g(b) \rvert = v_{418} \to v_{418} = v_{419} \to \lvert g(y) - g(b) \rvert = v_{419})$
\item $\varphi_{78} :\equiv \forall v_{415}\,\forall v_{416}\,(\lvert g(y) - g(b) \rvert = v_{415} \to v_{415} = v_{416} \to \lvert g(y) - g(b) \rvert = v_{416})$
\item $\varphi_{79} :\equiv x \in \mathrm{dom}(g \circ f) \to 0 < \lvert x - c \rvert \to \lvert x - c \rvert < \delta_{2} \to \lvert f(x) - b \rvert < \mathit{dg}$
\item $\varphi_{80} :\equiv \forall x\,(\mathrm{RFn}(g) \to x \in \mathrm{dom}(g) \to g(x) \in \mathbb{R}) \to \mathrm{RFn}(g) \to b \in \mathrm{dom}(g) \to g(b) \in \mathbb{R}$
\item $\varphi_{81} :\equiv f(x) \in \mathrm{dom}(g) \to \lvert f(x) - b \rvert < \mathit{dg} \to \lvert g(f(x)) - g(b) \rvert < \varepsilon$
\item $\varphi_{82} :\equiv \forall x \in \mathrm{dom}(g)\,(0 < \lvert x - b \rvert \to \lvert x - b \rvert < \mathit{dg} \to \lvert g(x) - g(b) \rvert < \varepsilon)$
\item $\varphi_{83} :\equiv \forall w\,(g(f(x)) = w \to \lvert g(f(x)) - g(b) \rvert < \varepsilon \to \lvert w - g(b) \rvert < \varepsilon)$
\item $\varphi_{84} :\equiv \forall f\,\forall x\,(\mathrm{RFn}(f) \to x \in \mathrm{dom}(f) \to f(x) \in \mathbb{R}) \to \forall x\,(\mathrm{RFn}(g) \to x \in \mathrm{dom}(g) \to g(x) \in \mathbb{R})$
\item $\varphi_{85} :\equiv x \in \mathrm{dom}(g \circ f) \to 0 < \lvert x - c \rvert \to \lvert x - c \rvert < \delta_{2} \to g \circ f(x) = g(f(x))$
\item $\varphi_{86} :\equiv x \in \mathrm{dom}(g \circ f) \to 0 < \lvert x - c \rvert \to \lvert x - c \rvert < \delta_{2} \to g(f(x)) = g \circ f(x)$
\item $\varphi_{87} :\equiv \forall v_{411}\,(y \in \mathbb{R} \to v_{411} \in \mathbb{R} \to y - v_{411} \in \mathbb{R}) \to y \in \mathbb{R} \to b \in \mathbb{R} \to y - b \in \mathbb{R}$
\item $\varphi_{88} :\equiv \forall v_{409}\,(y \in \mathbb{R} \to v_{409} \in \mathbb{R} \to y - v_{409} \in \mathbb{R}) \to y \in \mathbb{R} \to b \in \mathbb{R} \to y - b \in \mathbb{R}$
\item $\varphi_{89} :\equiv \forall x\,\forall y\,(x \in \mathbb{R} \to y \in \mathbb{R} \to x - y \in \mathbb{R}) \to \forall y\,(v_{410} \in \mathbb{R} \to y \in \mathbb{R} \to v_{410} - y \in \mathbb{R})$
\item $\varphi_{90} :\equiv \forall x\,\forall y\,(x \in \mathbb{R} \to y \in \mathbb{R} \to x - y \in \mathbb{R}) \to \forall y\,(v_{408} \in \mathbb{R} \to y \in \mathbb{R} \to v_{408} - y \in \mathbb{R})$
\item $\varphi_{91} :\equiv \forall b\,(\mathrm{FLim}(f,c,b) \to \mathrm{FLim}(g,b,g(b)) \to \varphi_{27})$
\item $\varphi_{92} :\equiv \forall x\,\forall y\,(x \in \mathbb{R} \to y \in \mathbb{R} \to x - y = 0 \to x = y) \to \forall y\,(v_{412} \in \mathbb{R} \to y \in \mathbb{R} \to v_{412} - y = 0 \to v_{412} = y)$
\item $\varphi_{93} :\equiv (\mathrm{FLim}(f,c,b) \to \mathrm{FLim}(f,c,b)) \to \mathrm{FLim}(f,c,b) \to \varphi_{28}$
\item $\varphi_{94} :\equiv \lvert g(y) - g(b) \rvert = \lvert g(z) - g(b) \rvert \to \lvert g(y) - g(b) \rvert = \lvert g(w) - g(b) \rvert$
\item $\varphi_{95} :\equiv \lvert g(y) - g(b) \rvert = \lvert g(y) - g(b) \rvert \to \lvert g(y) - g(b) \rvert = \lvert g(b) - g(b) \rvert$
\item $\varphi_{96} :\equiv \forall y\,(v_{410} \in \mathbb{R} \to y \in \mathbb{R} \to v_{410} - y \in \mathbb{R}) \to v_{410} \in \mathbb{R} \to v_{411} \in \mathbb{R} \to v_{410} - v_{411} \in \mathbb{R}$
\item $\varphi_{97} :\equiv \forall y\,(v_{408} \in \mathbb{R} \to y \in \mathbb{R} \to v_{408} - y \in \mathbb{R}) \to v_{408} \in \mathbb{R} \to v_{409} \in \mathbb{R} \to v_{408} - v_{409} \in \mathbb{R}$
\item $\varphi_{98} :\equiv \neg (\delta \in \mathbb{R} \land 0 < \delta \land \forall x \in \mathrm{dom}(f)\,(0 < \lvert x - c \rvert \to \lvert x - c \rvert < \delta \to \lvert f(x) - b \rvert < \mathit{dg}))$
\item $\varphi_{99} :\equiv (\mathrm{FLim}(f,c,b) \to \mathrm{FLim}(f,c,b)) \to \mathrm{FLim}(f,c,b) \to \varphi_{43}$
\item $\varphi_{100} :\equiv (\mathrm{FLim}(f,c,b) \to \mathrm{FLim}(f,c,b)) \to \mathrm{FLim}(f,c,b) \to \varphi_{46}$
\item $\varphi_{101} :\equiv \exists \delta\,(\delta \in \mathbb{R} \land 0 < \delta \land \forall x \in \mathrm{dom}(f)\,(0 < \lvert x - c \rvert \to \lvert x - c \rvert < \delta \to \lvert f(x) - b \rvert < \mathit{dg}))$
\item $\varphi_{102} :\equiv \forall y\,(g \circ f(x) = y \to y = g \circ f(x)) \to g \circ f(x) = g(f(x)) \to g(f(x)) = g \circ f(x)$
\item $\varphi_{103} :\equiv \forall y\,(v_{412} \in \mathbb{R} \to y \in \mathbb{R} \to v_{412} - y = 0 \to v_{412} = y) \to \varphi_{16}$
\item $\varphi_{104} :\equiv \forall w\,(g(b) - g(b) = w \to \lvert g(b) - g(b) \rvert = \lvert g(b) - g(b) \rvert \to \lvert g(b) - g(b) \rvert = \lvert w \rvert)$
\item $\varphi_{105} :\equiv \forall w\,(y = w \to \lvert g(y) - g(b) \rvert = \lvert g(y) - g(b) \rvert \to \lvert g(y) - g(b) \rvert = \lvert g(w) - g(b) \rvert)$
\item $\varphi_{106} :\equiv \neg (\delta \in \mathbb{R} \land 0 < \delta \land \forall x \in \mathrm{dom}(g)\,(0 < \lvert x - b \rvert \to \lvert x - b \rvert < \delta \to \lvert g(x) - g(b) \rvert < \varepsilon))$
\item $\varphi_{107} :\equiv \exists \delta\,(\delta \in \mathbb{R} \land 0 < \delta \land \forall x \in \mathrm{dom}(g)\,(0 < \lvert x - b \rvert \to \lvert x - b \rvert < \delta \to \lvert g(x) - g(b) \rvert < \varepsilon))$
\item $\varphi_{108} :\equiv \forall v_{416}\,(\lvert g(y) - g(b) \rvert = \lvert g(b) - g(b) \rvert \to \lvert g(b) - g(b) \rvert = v_{416} \to \lvert g(y) - g(b) \rvert = v_{416})$
\item $\varphi_{109} :\equiv b \in \mathrm{dom}(g) \to \mathrm{PNbhd}(g \circ f,c) \to x \in \mathrm{dom}(g \circ f) \to \varphi_{52}$
\item $\varphi_{110} :\equiv \mathrm{RFn}(g) \to \mathrm{RFn}(f) \to \mathrm{RFn}(g \circ f) \land \forall x \in \mathrm{dom}(g \circ f)\,(\varphi_{52})$
\item $\varphi_{111} :\equiv b \in \mathrm{dom}(g) \to \mathrm{PNbhd}(g \circ f,c) \to \forall x \in \mathrm{dom}(g \circ f)\,(\varphi_{52})$
\item $\varphi_{112} :\equiv \mathrm{PNbhd}(g \circ f,c) \to \mathrm{RFn}(g \circ f) \land \forall x \in \mathrm{dom}(g \circ f)\,(\varphi_{52})$
\item $\varphi_{113} :\equiv \exists \delta\,(\delta \in \mathbb{R} \land 0 < \delta \land \forall x \in \mathrm{dom}(g \circ f)\,(0 < \lvert x - c \rvert \to \lvert x - c \rvert < \delta \to \lvert g \circ f(x) - g(b) \rvert < \varepsilon))$
\item $\varphi_{114} :\equiv \forall \delta\,(\neg (\delta \in \mathbb{R} \land 0 < \delta \land \forall x \in \mathrm{dom}(g \circ f)\,(0 < \lvert x - c \rvert \to \lvert x - c \rvert < \delta \to \lvert g \circ f(x) - g(b) \rvert < \varepsilon)))$
\item $\varphi_{115} :\equiv \delta_{2} \in \mathbb{R} \land 0 < \delta_{2} \land \varphi_{72} \to x \in \mathrm{dom}(g \circ f) \to x \in \mathrm{dom}(g \circ f)$
\item $\varphi_{116} :\equiv 0 < \varepsilon \to \mathit{dg} \in \mathbb{R} \land 0 < \mathit{dg} \land \varphi_{82} \to y \in \mathrm{dom}(g) \to y \in \mathrm{dom}(g)$
\item $\varphi_{117} :\equiv \forall \varepsilon\,(\mathrm{FLim}(f,c,b) \to \varepsilon \in \mathbb{R} \to 0 < \varepsilon \to \exists \delta\,(\delta \in \mathbb{R} \land 0 < \delta \land \forall x \in \mathrm{dom}(f)\,(0 < \lvert x - c \rvert \to \lvert x - c \rvert < \delta \to \lvert f(x) - b \rvert < \varepsilon)))$
\item $\varphi_{118} :\equiv \forall \alpha\,\forall \varepsilon\,(\mathrm{FLim}(f,c,\alpha) \to \varepsilon \in \mathbb{R} \to 0 < \varepsilon \to \exists \delta\,(\delta \in \mathbb{R} \land 0 < \delta \land \forall x \in \mathrm{dom}(f)\,(0 < \lvert x - c \rvert \to \lvert x - c \rvert < \delta \to \lvert f(x) - \alpha \rvert < \varepsilon)))$
\item $\varphi_{119} :\equiv \forall \alpha\,\forall \varepsilon\,(\mathrm{FLim}(g,b,\alpha) \to \varepsilon \in \mathbb{R} \to 0 < \varepsilon \to \exists \delta\,(\delta \in \mathbb{R} \land 0 < \delta \land \forall x \in \mathrm{dom}(g)\,(0 < \lvert x - b \rvert \to \lvert x - b \rvert < \delta \to \lvert g(x) - \alpha \rvert < \varepsilon)))$
\item $\varphi_{120} :\equiv \mathit{dg} \in \mathbb{R} \land 0 < \mathit{dg} \land \varphi_{82} \to y \in \mathrm{dom}(g) \to \lvert y - b \rvert < \mathit{dg} \to y \in \mathbb{R}$
\item $\varphi_{121} :\equiv (\mathrm{FLim}(g,b,g(b)) \to \mathrm{RFn}(g)) \to \varphi_{6} \to \mathrm{FLim}(f,c,b) \to \varphi_{29}$
\item $\varphi_{122} :\equiv \forall c\,\forall \alpha\,\forall \varepsilon\,(\mathrm{FLim}(g,c,\alpha) \to \varepsilon \in \mathbb{R} \to 0 < \varepsilon \to \exists \delta\,(\delta \in \mathbb{R} \land 0 < \delta \land \forall x \in \mathrm{dom}(g)\,(0 < \lvert x - c \rvert \to \lvert x - c \rvert < \delta \to \lvert g(x) - \alpha \rvert < \varepsilon)))$
\item $\varphi_{123} :\equiv \mathit{dg} \in \mathbb{R} \land 0 < \mathit{dg} \land \varphi_{82} \to y \in \mathrm{dom}(g) \to \lvert y - b \rvert < \mathit{dg} \to 0 \leq \lvert y - b \rvert$
\item $\varphi_{124} :\equiv \mathit{dg} \in \mathbb{R} \land 0 < \mathit{dg} \land \varphi_{82} \to y \in \mathrm{dom}(g) \to \lvert y - b \rvert < \mathit{dg} \to y - b \in \mathbb{R}$
\item $\varphi_{125} :\equiv \mathit{dg} \in \mathbb{R} \land 0 < \mathit{dg} \land \varphi_{82} \to y \in \mathrm{dom}(g) \to \lvert y - b \rvert < \mathit{dg} \to \lvert y - b \rvert < \mathit{dg}$
\item $\varphi_{126} :\equiv \delta_{2} \in \mathbb{R} \land 0 < \delta_{2} \land \varphi_{72} \to x \in \mathrm{dom}(g \circ f) \to 0 < \lvert x - c \rvert \to 0 < \lvert x - c \rvert$
\item $\varphi_{127} :\equiv \mathit{dg} \in \mathbb{R} \land 0 < \mathit{dg} \land \varphi_{82} \to y \in \mathrm{dom}(g) \to \lvert y - b \rvert < \mathit{dg} \to 0 = \lvert y - b \rvert \to 0 \leq 0$
\item $\varphi_{128} :\equiv \mathit{dg} \in \mathbb{R} \land 0 < \mathit{dg} \land \varphi_{82} \to y \in \mathrm{dom}(g) \to \lvert y - b \rvert < \mathit{dg} \to 0 = \lvert y - b \rvert \to y = b$
\item $\varphi_{129} :\equiv \mathrm{PNbhd}(g \circ f,c) \to \varepsilon \in \mathbb{R} \to 0 < \varepsilon \to \mathit{dg} \in \mathbb{R} \land 0 < \mathit{dg} \land \varphi_{82} \to 0 < \mathit{dg}$
\item $\varphi_{130} :\equiv \mathrm{PNbhd}(g \circ f,c) \to \varepsilon \in \mathbb{R} \to 0 < \varepsilon \to \mathit{dg} \in \mathbb{R} \land 0 < \mathit{dg} \land \varphi_{82} \to \mathit{dg} \in \mathbb{R}$
\item $\varphi_{131} :\equiv b \in \mathrm{dom}(g) \to \mathrm{PNbhd}(g \circ f,c) \to \varepsilon \in \mathbb{R} \to 0 < \varepsilon \to \varphi_{107}$
\item $\varphi_{132} :\equiv \mathrm{FLim}(f,c,b) \to \mathrm{FLim}(g,b,g(b)) \to b \in \mathrm{dom}(g) \to \varphi_{112}$
\item $\varphi_{133} :\equiv b \in \mathrm{dom}(g) \to \mathrm{PNbhd}(g \circ f,c) \to \varepsilon \in \mathbb{R} \to 0 < \varepsilon \to \varphi_{113}$
\item $\varphi_{134} :\equiv b \in \mathrm{dom}(g) \to \mathrm{PNbhd}(g \circ f,c) \to \forall \varepsilon \in \mathbb{R}\,(0 < \varepsilon \to \varphi_{113})$
\item $\varphi_{135} :\equiv \forall f\,\forall c\,\forall \alpha\,(\mathrm{FLim}(f,c,\alpha) \leftrightarrow \mathrm{RFn}(f) \land c \in \mathbb{R} \land \alpha \in \mathbb{R} \land \mathrm{PNbhd}(f,c) \land \forall \varepsilon \in \mathbb{R}\,(0 < \varepsilon \to \exists \delta\,(\delta \in \mathbb{R} \land 0 < \delta \land \forall x \in \mathrm{dom}(f)\,(0 < \lvert x - c \rvert \to \lvert x - c \rvert < \delta \to \lvert f(x) - \alpha \rvert < \varepsilon))))$
\item $\varphi_{136} :\equiv \mathit{dg} \in \mathbb{R} \land 0 < \mathit{dg} \land \varphi_{82} \to y \in \mathrm{dom}(g) \to \lvert y - b \rvert < \mathit{dg} \to \varphi_{44}$
\item $\varphi_{137} :\equiv \mathrm{FLim}(f,c,b) \to \mathrm{FLim}(g,b,g(b)) \to b \in \mathrm{dom}(g) \to \varphi_{129}$
\item $\varphi_{138} :\equiv \mathrm{PNbhd}(g \circ f,c) \to \varepsilon \in \mathbb{R} \to 0 < \varepsilon \to \mathit{dg} \in \mathbb{R} \land 0 < \mathit{dg} \land \varphi_{82} \to \varphi_{17}$
\item $\varphi_{139} :\equiv \mathrm{PNbhd}(g \circ f,c) \to \varepsilon \in \mathbb{R} \to 0 < \varepsilon \to \mathit{dg} \in \mathbb{R} \land 0 < \mathit{dg} \land \varphi_{82} \to \varphi_{18}$
\item $\varphi_{140} :\equiv b \in \mathrm{dom}(g) \to \mathrm{PNbhd}(g \circ f,c) \to \varepsilon \in \mathbb{R} \to 0 < \varepsilon \to \varphi_{120}$
\item $\varphi_{141} :\equiv \mathrm{FLim}(f,c,b) \to \mathrm{FLim}(g,b,g(b)) \to b \in \mathrm{dom}(g) \to \varphi_{130}$
\item $\varphi_{142} :\equiv \varepsilon \in \mathbb{R} \to 0 < \varepsilon \to \mathit{dg} \in \mathbb{R} \land 0 < \mathit{dg} \land \varphi_{82} \to y \in \mathrm{dom}(g) \to \varphi_{48}$
\item $\varphi_{143} :\equiv \varepsilon \in \mathbb{R} \to 0 < \varepsilon \to \mathit{dg} \in \mathbb{R} \land 0 < \mathit{dg} \land \varphi_{82} \to y \in \mathrm{dom}(g) \to \varphi_{49}$
\item $\varphi_{144} :\equiv \mathrm{FLim}(g,b,g(b)) \to b \in \mathrm{dom}(g) \to \mathrm{PNbhd}(g \circ f,c) \to \varepsilon \in \mathbb{R} \to \varphi_{116}$
\item $\varphi_{145} :\equiv b \in \mathrm{dom}(g) \to \mathrm{PNbhd}(g \circ f,c) \to \varepsilon \in \mathbb{R} \to 0 < \varepsilon \to \varphi_{123}$
\item $\varphi_{146} :\equiv \mathrm{PNbhd}(g \circ f,c) \to \varepsilon \in \mathbb{R} \to 0 < \varepsilon \to \mathit{dg} \in \mathbb{R} \land 0 < \mathit{dg} \land \varphi_{82} \to \varphi_{32}$
\item $\varphi_{147} :\equiv \mathrm{PNbhd}(g \circ f,c) \to \varepsilon \in \mathbb{R} \to 0 < \varepsilon \to \mathit{dg} \in \mathbb{R} \land 0 < \mathit{dg} \land \varphi_{82} \to \varphi_{33}$
\item $\varphi_{148} :\equiv b \in \mathrm{dom}(g) \to \mathrm{PNbhd}(g \circ f,c) \to \varepsilon \in \mathbb{R} \to 0 < \varepsilon \to \varphi_{124}$
\item $\varphi_{149} :\equiv \mathrm{PNbhd}(g \circ f,c) \to \varepsilon \in \mathbb{R} \to 0 < \varepsilon \to \mathit{dg} \in \mathbb{R} \land 0 < \mathit{dg} \land \varphi_{82} \to \varphi_{40}$
\item $\varphi_{150} :\equiv b \in \mathrm{dom}(g) \to \mathrm{PNbhd}(g \circ f,c) \to \varepsilon \in \mathbb{R} \to 0 < \varepsilon \to \varphi_{125}$
\item $\varphi_{151} :\equiv \mathrm{PNbhd}(g \circ f,c) \to \varepsilon \in \mathbb{R} \to 0 < \varepsilon \to \mathit{dg} \in \mathbb{R} \land 0 < \mathit{dg} \land \varphi_{82} \to \varphi_{45}$
\item $\varphi_{152} :\equiv b \in \mathrm{dom}(g) \to \mathrm{PNbhd}(g \circ f,c) \to \varepsilon \in \mathbb{R} \to 0 < \varepsilon \to \varphi_{127}$
\item $\varphi_{153} :\equiv b \in \mathrm{dom}(g) \to \mathrm{PNbhd}(g \circ f,c) \to \varepsilon \in \mathbb{R} \to 0 < \varepsilon \to \varphi_{128}$
\item $\varphi_{154} :\equiv (\mathrm{FLim}(f,c,b) \to \varphi_{29}) \to (\mathrm{FLim}(f,c,b) \to \varphi_{5}) \to \mathrm{FLim}(f,c,b) \to \varphi_{62}$
\item $\varphi_{155} :\equiv \mathrm{PNbhd}(g \circ f,c) \to \varepsilon \in \mathbb{R} \to 0 < \varepsilon \to \mathit{dg} \in \mathbb{R} \land 0 < \mathit{dg} \land \varphi_{82} \to \varphi_{54}$
\item $\varphi_{156} :\equiv \forall \delta_{2}\,(\neg (\delta_{2} \in \mathbb{R} \land 0 < \delta_{2} \land \varphi_{72})) \to \forall \delta\,\forall \delta_{2}\,(\neg (\delta_{2} \in \mathbb{R} \land 0 < \delta_{2} \land \varphi_{72}))$
\item $\varphi_{157} :\equiv \mathrm{PNbhd}(g \circ f,c) \to \varepsilon \in \mathbb{R} \to 0 < \varepsilon \to \mathit{dg} \in \mathbb{R} \land 0 < \mathit{dg} \land \varphi_{82} \to \varphi_{60}$
\item $\varphi_{158} :\equiv \mathrm{PNbhd}(g \circ f,c) \to \varepsilon \in \mathbb{R} \to 0 < \varepsilon \to \mathit{dg} \in \mathbb{R} \land 0 < \mathit{dg} \land \varphi_{82} \to \varphi_{64}$
\item $\varphi_{159} :\equiv \mathrm{PNbhd}(g \circ f,c) \to \varepsilon \in \mathbb{R} \to 0 < \varepsilon \to \mathit{dg} \in \mathbb{R} \land 0 < \mathit{dg} \land \varphi_{82} \to \varphi_{65}$
\item $\varphi_{160} :\equiv \forall \mathit{dg}\,(\neg (\mathit{dg} \in \mathbb{R} \land 0 < \mathit{dg} \land \varphi_{82})) \to \forall \delta\,\forall \mathit{dg}\,(\neg (\mathit{dg} \in \mathbb{R} \land 0 < \mathit{dg} \land \varphi_{82}))$
\item $\varphi_{161} :\equiv \delta_{2} \in \mathbb{R} \land 0 < \delta_{2} \land \varphi_{72} \to x \in \mathrm{dom}(g \circ f) \to 0 < \lvert x - c \rvert \to \lvert x - c \rvert < \delta_{2} \to \varphi_{52}$
\item $\varphi_{162} :\equiv \delta_{2} \in \mathbb{R} \land 0 < \delta_{2} \land \varphi_{72} \to \delta_{2} \in \mathbb{R} \land 0 < \delta_{2} \land \forall x \in \mathrm{dom}(g \circ f)\,(\varphi_{69})$
\item $\varphi_{163} :\equiv \mathrm{PNbhd}(g \circ f,c) \to \varepsilon \in \mathbb{R} \to 0 < \varepsilon \to \mathit{dg} \in \mathbb{R} \land 0 < \mathit{dg} \land \varphi_{82} \to \varphi_{73}$
\item $\varphi_{164} :\equiv \mathrm{PNbhd}(g \circ f,c) \to \varepsilon \in \mathbb{R} \to 0 < \varepsilon \to \mathit{dg} \in \mathbb{R} \land 0 < \mathit{dg} \land \varphi_{82} \to \varphi_{75}$
\item $\varphi_{165} :\equiv \forall \alpha\,(\mathrm{FLim}(g \circ f,c,\alpha) \leftrightarrow \mathrm{RFn}(g \circ f) \land c \in \mathbb{R} \land \alpha \in \mathbb{R} \land \mathrm{PNbhd}(g \circ f,c) \land \forall \varepsilon \in \mathbb{R}\,(0 < \varepsilon \to \exists \delta\,(\delta \in \mathbb{R} \land 0 < \delta \land \forall x \in \mathrm{dom}(g \circ f)\,(0 < \lvert x - c \rvert \to \lvert x - c \rvert < \delta \to \lvert g \circ f(x) - \alpha \rvert < \varepsilon))))$
\item $\varphi_{166} :\equiv \mathrm{PNbhd}(g \circ f,c) \to \varepsilon \in \mathbb{R} \to 0 < \varepsilon \to \mathit{dg} \in \mathbb{R} \land 0 < \mathit{dg} \land \varphi_{82} \to \varphi_{81}$
\item $\varphi_{167} :\equiv \mathrm{PNbhd}(g \circ f,c) \to \varepsilon \in \mathbb{R} \to 0 < \varepsilon \to \mathit{dg} \in \mathbb{R} \land 0 < \mathit{dg} \land \varphi_{82} \to \varphi_{82}$
\item $\varphi_{168} :\equiv \varepsilon \in \mathbb{R} \to 0 < \varepsilon \to \mathit{dg} \in \mathbb{R} \land 0 < \mathit{dg} \land \varphi_{82} \to \mathit{dg} \in \mathbb{R} \land 0 < \mathit{dg} \land \varphi_{82}$
\item $\varphi_{169} :\equiv \mathrm{FLim}(g \circ f,c,g(b)) \leftrightarrow \mathrm{RFn}(g \circ f) \land c \in \mathbb{R} \land g(b) \in \mathbb{R} \land \mathrm{PNbhd}(g \circ f,c) \land \forall \varepsilon \in \mathbb{R}\,(0 < \varepsilon \to \varphi_{113})$
\item $\varphi_{170} :\equiv \mathrm{FLim}(f,c,b) \to \mathrm{FLim}(g,b,g(b)) \to b \in \mathrm{dom}(g) \to \varphi_{138}$
\item $\varphi_{171} :\equiv \mathrm{FLim}(f,c,b) \to \mathrm{FLim}(g,b,g(b)) \to b \in \mathrm{dom}(g) \to \varphi_{139}$
\item $\varphi_{172} :\equiv \mathrm{PNbhd}(g \circ f,c) \to \varepsilon \in \mathbb{R} \to 0 < \varepsilon \to \mathit{dg} \in \mathbb{R} \land 0 < \mathit{dg} \land \varphi_{82} \to \varphi_{101}$
\item $\varphi_{173} :\equiv 0 = 0 \to \mathrm{FLim}(f,c,b) \to \mathrm{FLim}(g,b,g(b)) \to \varphi_{152}$
\item $\varphi_{174} :\equiv \mathrm{FLim}(f,c,b) \to \mathrm{FLim}(g,b,g(b)) \to b \in \mathrm{dom}(g) \to \varphi_{146}$
\item $\varphi_{175} :\equiv \mathrm{FLim}(f,c,b) \to \mathrm{FLim}(g,b,g(b)) \to b \in \mathrm{dom}(g) \to \varphi_{147}$
\item $\varphi_{176} :\equiv \mathrm{FLim}(f,c,b) \to \mathrm{FLim}(g,b,g(b)) \to b \in \mathrm{dom}(g) \to \varphi_{149}$
\item $\varphi_{177} :\equiv \mathrm{FLim}(f,c,b) \to \mathrm{FLim}(g,b,g(b)) \to b \in \mathrm{dom}(g) \to \varphi_{151}$
\item $\varphi_{178} :\equiv b \in \mathrm{dom}(g) \to \mathrm{PNbhd}(g \circ f,c) \to \varepsilon \in \mathbb{R} \to 0 < \varepsilon \to \varphi_{136}$
\item $\varphi_{179} :\equiv \mathrm{FLim}(g,b,g(b)) \to b \in \mathrm{dom}(g) \to \mathrm{PNbhd}(g \circ f,c) \to \varphi_{142}$
\item $\varphi_{180} :\equiv \mathrm{FLim}(g,b,g(b)) \to b \in \mathrm{dom}(g) \to \mathrm{PNbhd}(g \circ f,c) \to \varphi_{143}$
\item $\varphi_{181} :\equiv \forall y\,(\varphi_{47} \land \varepsilon \in \mathbb{R} \land 0 < \varepsilon \land (\mathit{dg} \in \mathbb{R} \land 0 < \mathit{dg} \land \varphi_{82}) \to \varphi_{45})$
\item $\varphi_{182} :\equiv \mathrm{FLim}(f,c,b) \to \mathrm{FLim}(g,b,g(b)) \to b \in \mathrm{dom}(g) \to \varphi_{155}$
\item $\varphi_{183} :\equiv \mathrm{FLim}(f,c,b) \to \mathrm{FLim}(g,b,g(b)) \to b \in \mathrm{dom}(g) \to \varphi_{157}$
\item $\varphi_{184} :\equiv \forall y\,(\varphi_{47} \land \varepsilon \in \mathbb{R} \land 0 < \varepsilon \land (\mathit{dg} \in \mathbb{R} \land 0 < \mathit{dg} \land \varphi_{82})) \to \varphi_{54}$
\item $\varphi_{185} :\equiv \mathrm{FLim}(f,c,b) \to \mathrm{FLim}(g,b,g(b)) \to b \in \mathrm{dom}(g) \to \varphi_{158}$
\item $\varphi_{186} :\equiv \mathrm{FLim}(f,c,b) \to \mathrm{FLim}(g,b,g(b)) \to b \in \mathrm{dom}(g) \to \varphi_{159}$
\item $\varphi_{187} :\equiv \mathrm{FLim}(f,c,b) \to \mathrm{FLim}(g,b,g(b)) \to b \in \mathrm{dom}(g) \to \varphi_{163}$
\item $\varphi_{188} :\equiv \mathrm{PNbhd}(g \circ f,c) \to \varepsilon \in \mathbb{R} \to 0 < \varepsilon \to \mathit{dg} \in \mathbb{R} \land 0 < \mathit{dg} \land \varphi_{82} \to \varphi_{113}$
\item $\varphi_{189} :\equiv \mathrm{FLim}(f,c,b) \to \mathrm{FLim}(g,b,g(b)) \to b \in \mathrm{dom}(g) \to \varphi_{164}$
\item $\varphi_{190} :\equiv \mathrm{FLim}(f,c,b) \to \mathrm{FLim}(g,b,g(b)) \to b \in \mathrm{dom}(g) \to \varphi_{166}$
\item $\varphi_{191} :\equiv \mathrm{FLim}(f,c,b) \to \mathrm{FLim}(g,b,g(b)) \to b \in \mathrm{dom}(g) \to \varphi_{167}$
\item $\varphi_{192} :\equiv \mathrm{FLim}(g,b,g(b)) \to b \in \mathrm{dom}(g) \to \mathrm{PNbhd}(g \circ f,c) \to \varphi_{168}$
\item $\varphi_{193} :\equiv \mathrm{PNbhd}(g \circ f,c) \to \varepsilon \in \mathbb{R} \to 0 < \varepsilon \to \mathit{dg} \in \mathbb{R} \land 0 < \mathit{dg} \land \varphi_{82} \to \varphi_{115}$
\item $\varphi_{194} :\equiv \varphi_{47} \land \varepsilon \in \mathbb{R} \land 0 < \varepsilon \land (\mathit{dg} \in \mathbb{R} \land 0 < \mathit{dg} \land \varphi_{82}) \land (\delta_{2} \in \mathbb{R} \land 0 < \delta_{2} \land \varphi_{72})$
\item $\varphi_{195} :\equiv \mathrm{FLim}(f,c,b) \to \mathrm{FLim}(g,b,g(b)) \to b \in \mathrm{dom}(g) \to \varphi_{172}$
\item $\varphi_{196} :\equiv (\mathrm{FLim}(f,c,b) \to \mathrm{FLim}(g,b,g(b)) \to \varphi_{55}) \to \mathrm{FLim}(f,c,b) \to \mathrm{FLim}(g,b,g(b)) \to \varphi_{131}$
\item $\varphi_{197} :\equiv \mathit{dg} \in \mathbb{R} \land 0 < \mathit{dg} \land \varphi_{82} \to \delta_{2} \in \mathbb{R} \land 0 < \delta_{2} \land \varphi_{72} \to x \in \mathrm{dom}(g \circ f) \to \varphi_{67}$
\item $\varphi_{198} :\equiv \mathit{dg} \in \mathbb{R} \land 0 < \mathit{dg} \land \varphi_{82} \to \delta_{2} \in \mathbb{R} \land 0 < \delta_{2} \land \varphi_{72} \to x \in \mathrm{dom}(g \circ f) \to \varphi_{69}$
\item $\varphi_{199} :\equiv \mathrm{PNbhd}(g \circ f,c) \to \varepsilon \in \mathbb{R} \to 0 < \varepsilon \to \mathit{dg} \in \mathbb{R} \land 0 < \mathit{dg} \land \varphi_{82} \to \varphi_{126}$
\item $\varphi_{200} :\equiv \varepsilon \in \mathbb{R} \to 0 < \varepsilon \to \mathit{dg} \in \mathbb{R} \land 0 < \mathit{dg} \land \varphi_{82} \to \delta_{2} \in \mathbb{R} \land 0 < \delta_{2} \land \varphi_{72} \to \varphi_{66}$
\item $\varphi_{201} :\equiv \varphi_{119} \to \forall \varepsilon\,(\mathrm{FLim}(g,b,g(b)) \to \varepsilon \in \mathbb{R} \to 0 < \varepsilon \to \varphi_{107})$
\item $\varphi_{202} :\equiv \mathit{dg} \in \mathbb{R} \land 0 < \mathit{dg} \land \varphi_{82} \to \delta_{2} \in \mathbb{R} \land 0 < \delta_{2} \land \varphi_{72} \to \forall x \in \mathrm{dom}(g \circ f)\,(\varphi_{69})$
\item $\varphi_{203} :\equiv \varepsilon \in \mathbb{R} \to 0 < \varepsilon \to \mathit{dg} \in \mathbb{R} \land 0 < \mathit{dg} \land \varphi_{82} \to \delta_{2} \in \mathbb{R} \land 0 < \delta_{2} \land \varphi_{72} \to \varphi_{70}$
\item $\varphi_{204} :\equiv (\mathrm{FLim}(f,c,b) \to \mathrm{FLim}(g,b,g(b)) \to \varphi_{134}) \to \mathrm{FLim}(f,c,b) \to \mathrm{FLim}(g,b,g(b)) \to \varphi_{27}$
\item $\varphi_{205} :\equiv \mathit{dg} \in \mathbb{R} \land 0 < \mathit{dg} \land \varphi_{82} \to \delta_{2} \in \mathbb{R} \land 0 < \delta_{2} \land \varphi_{72} \to \delta_{2} \in \mathbb{R} \land 0 < \delta_{2} \land \varphi_{72}$
\item $\varphi_{206} :\equiv \varepsilon \in \mathbb{R} \to 0 < \varepsilon \to \mathit{dg} \in \mathbb{R} \land 0 < \mathit{dg} \land \varphi_{82} \to \delta_{2} \in \mathbb{R} \land 0 < \delta_{2} \land \varphi_{72} \to \varphi_{72}$
\item $\varphi_{207} :\equiv \forall \varepsilon\,(\mathrm{FLim}(g,b,g(b)) \to \varepsilon \in \mathbb{R} \to 0 < \varepsilon \to \varphi_{107}) \to \mathrm{FLim}(g,b,g(b)) \to \varepsilon \in \mathbb{R} \to 0 < \varepsilon \to \varphi_{107}$
\item $\varphi_{208} :\equiv \varepsilon \in \mathbb{R} \to 0 < \varepsilon \to \mathit{dg} \in \mathbb{R} \land 0 < \mathit{dg} \land \varphi_{82} \to \delta_{2} \in \mathbb{R} \land 0 < \delta_{2} \land \varphi_{72} \to \varphi_{79}$
\item $\varphi_{209} :\equiv (\mathrm{FLim}(f,c,b) \to \varphi_{29}) \to (\mathrm{FLim}(f,c,b) \to \varphi_{28}) \to \varphi_{132}$
\item $\varphi_{210} :\equiv \mathrm{FLim}(f,c,b) \to \mathrm{FLim}(g,b,g(b)) \to b \in \mathrm{dom}(g) \to \varphi_{188}$
\item $\varphi_{211} :\equiv (\mathrm{FLim}(f,c,b) \to \varphi_{43}) \to \mathrm{FLim}(f,c,b) \to \mathrm{FLim}(g,b,g(b)) \to \varphi_{148}$
\item $\varphi_{212} :\equiv \varepsilon \in \mathbb{R} \to 0 < \varepsilon \to \mathit{dg} \in \mathbb{R} \land 0 < \mathit{dg} \land \varphi_{82} \to \delta_{2} \in \mathbb{R} \land 0 < \delta_{2} \land \varphi_{72} \to \varphi_{85}$
\item $\varphi_{213} :\equiv \varepsilon \in \mathbb{R} \to 0 < \varepsilon \to \mathit{dg} \in \mathbb{R} \land 0 < \mathit{dg} \land \varphi_{82} \to \delta_{2} \in \mathbb{R} \land 0 < \delta_{2} \land \varphi_{72} \to \varphi_{86}$
\item $\varphi_{214} :\equiv (\lvert g(y) - g(b) \rvert = 0 \to 0 = \lvert g(y) - g(b) \rvert) \to \varphi_{186}$
\item $\varphi_{215} :\equiv \forall \mathit{dg}\,(\varphi_{47} \land \varepsilon \in \mathbb{R} \land 0 < \varepsilon \land \neg \varphi_{113} \to \neg (\mathit{dg} \in \mathbb{R} \land 0 < \mathit{dg} \land \varphi_{82}))$
\item $\varphi_{216} :\equiv \forall \delta_{2}\,(\varphi_{47} \land \varepsilon \in \mathbb{R} \land 0 < \varepsilon \land (\mathit{dg} \in \mathbb{R} \land 0 < \mathit{dg} \land \varphi_{82}) \land \neg \varphi_{113})$
\item $\varphi_{217} :\equiv (\forall x \in \mathrm{dom}(g \circ f)\,(\varphi_{52}) \to x \in \mathrm{dom}(g \circ f) \to \varphi_{52}) \to \mathrm{FLim}(f,c,b) \to \mathrm{FLim}(g,b,g(b)) \to \varphi_{109}$
\item $\varphi_{218} :\equiv \forall \mathit{dg}\,(\varphi_{47} \land \varepsilon \in \mathbb{R} \land 0 < \varepsilon \land \neg \varphi_{113}) \to \forall \mathit{dg}\,(\neg (\mathit{dg} \in \mathbb{R} \land 0 < \mathit{dg} \land \varphi_{82}))$
\item $\varphi_{219} :\equiv \mathrm{FLim}(f,c,b) \to \mathrm{FLim}(g,b,g(b)) \to b \in \mathrm{dom}(g) \to \varphi_{193}$
\item $\varphi_{220} :\equiv \mathrm{FLim}(f,c,b) \to \mathrm{FLim}(g,b,g(b)) \to b \in \mathrm{dom}(g) \to \varphi_{199}$
\item $\varphi_{221} :\equiv (\varphi_{47} \to \forall \varepsilon\,(\varphi_{47})) \to \mathrm{FLim}(f,c,b) \to \mathrm{FLim}(g,b,g(b)) \to \varphi_{134}$
\item $\varphi_{222} :\equiv \varepsilon \in \mathbb{R} \to 0 < \varepsilon \to \mathit{dg} \in \mathbb{R} \land 0 < \mathit{dg} \land \varphi_{82} \to \delta_{2} \in \mathbb{R} \land 0 < \delta_{2} \land \varphi_{72} \to \varphi_{113}$
\item $\varphi_{223} :\equiv \mathrm{FLim}(g,b,g(b)) \to b \in \mathrm{dom}(g) \to \mathrm{PNbhd}(g \circ f,c) \to \varphi_{200}$
\item $\varphi_{224} :\equiv \mathrm{FLim}(g,b,g(b)) \to b \in \mathrm{dom}(g) \to \mathrm{PNbhd}(g \circ f,c) \to \varphi_{203}$
\item $\varphi_{225} :\equiv b \in \mathrm{dom}(g) \to \mathrm{PNbhd}(g \circ f,c) \to \varepsilon \in \mathbb{R} \to 0 < \varepsilon \to \varphi_{197}$
\item $\varphi_{226} :\equiv \mathrm{FLim}(g,b,g(b)) \to b \in \mathrm{dom}(g) \to \mathrm{PNbhd}(g \circ f,c) \to \varphi_{206}$
\item $\varphi_{227} :\equiv b \in \mathrm{dom}(g) \to \mathrm{PNbhd}(g \circ f,c) \to \varepsilon \in \mathbb{R} \to 0 < \varepsilon \to \varphi_{198}$
\item $\varphi_{228} :\equiv b \in \mathrm{dom}(g) \to \mathrm{PNbhd}(g \circ f,c) \to \varepsilon \in \mathbb{R} \to 0 < \varepsilon \to \varphi_{202}$
\item $\varphi_{229} :\equiv b \in \mathrm{dom}(g) \to \mathrm{PNbhd}(g \circ f,c) \to \varepsilon \in \mathbb{R} \to 0 < \varepsilon \to \varphi_{205}$
\item $\varphi_{230} :\equiv \mathrm{FLim}(g,b,g(b)) \to b \in \mathrm{dom}(g) \to \mathrm{PNbhd}(g \circ f,c) \to \varphi_{208}$
\item $\varphi_{231} :\equiv \mathrm{FLim}(g,b,g(b)) \to b \in \mathrm{dom}(g) \to \mathrm{PNbhd}(g \circ f,c) \to \varphi_{212}$
\item $\varphi_{232} :\equiv \mathrm{FLim}(g,b,g(b)) \to b \in \mathrm{dom}(g) \to \mathrm{PNbhd}(g \circ f,c) \to \varphi_{213}$
\item $\varphi_{233} :\equiv \mathrm{PNbhd}(g \circ f,c) \to \varepsilon \in \mathbb{R} \to 0 < \varepsilon \to \mathit{dg} \in \mathbb{R} \land 0 < \mathit{dg} \land \varphi_{82} \to \varphi_{161}$
\item $\varphi_{234} :\equiv \mathrm{PNbhd}(g \circ f,c) \to \varepsilon \in \mathbb{R} \to 0 < \varepsilon \to \mathit{dg} \in \mathbb{R} \land 0 < \mathit{dg} \land \varphi_{82} \to \varphi_{162}$
\item $\varphi_{235} :\equiv (\lvert g(y) - g(b) \rvert = \lvert g(b) - g(b) \rvert \to \varphi_{41}) \to \mathrm{FLim}(f,c,b) \to \varphi_{180}$
\item $\varphi_{236} :\equiv (g(b) \in \mathbb{R} \to g(b) - g(b) = 0) \to (\mathrm{FLim}(f,c,b) \to \varphi_{62}) \to \varphi_{183}$
\item $\varphi_{237} :\equiv (\mathrm{FLim}(f,c,b) \to \mathrm{FLim}(g,b,g(b)) \to \varphi_{30}) \to \varphi_{196}$
\item $\varphi_{238} :\equiv \varphi_{47} \land \varepsilon \in \mathbb{R} \land 0 < \varepsilon \land (\mathit{dg} \in \mathbb{R} \land 0 < \mathit{dg} \land \varphi_{82}) \to \forall y\,(\varphi_{47} \land \varepsilon \in \mathbb{R} \land 0 < \varepsilon \land (\mathit{dg} \in \mathbb{R} \land 0 < \mathit{dg} \land \varphi_{82}))$
\item $\varphi_{239} :\equiv (\mathrm{FLim}(f,c,b) \to \mathrm{FLim}(g,b,g(b)) \to \varphi_{9}) \to \varphi_{204}$
\item $\varphi_{240} :\equiv (y = b \to \varphi_{95}) \to \mathrm{FLim}(f,c,b) \to \mathrm{FLim}(g,b,g(b)) \to \varphi_{178}$
\item $\varphi_{241} :\equiv \mathrm{FLim}(g,b,g(b)) \to b \in \mathrm{dom}(g) \to \mathrm{PNbhd}(g \circ f,c) \to \varphi_{222}$
\item $\varphi_{242} :\equiv \mathrm{FLim}(f,c,b) \to \mathrm{FLim}(g,b,g(b)) \to b \in \mathrm{dom}(g) \to \varphi_{233}$
\item $\varphi_{243} :\equiv \mathrm{FLim}(f,c,b) \to \mathrm{FLim}(g,b,g(b)) \to b \in \mathrm{dom}(g) \to \varphi_{234}$
\item $\varphi_{244} :\equiv \varphi_{47} \land \varepsilon \in \mathbb{R} \land 0 < \varepsilon \land (\mathit{dg} \in \mathbb{R} \land 0 < \mathit{dg} \land \varphi_{82}) \land \neg \varphi_{113} \to \neg (\delta_{2} \in \mathbb{R} \land 0 < \delta_{2} \land \varphi_{72})$
\item $\varphi_{245} :\equiv \forall \delta\,(\forall \delta_{2}\,(\neg (\delta_{2} \in \mathbb{R} \land 0 < \delta_{2} \land \varphi_{72})) \to \varphi_{98}) \to \forall \delta\,\forall \delta_{2}\,(\neg (\delta_{2} \in \mathbb{R} \land 0 < \delta_{2} \land \varphi_{72})) \to \forall \delta\,(\varphi_{98})$
\item $\varphi_{246} :\equiv (\mathrm{FLim}(f,c,b) \to \mathrm{FLim}(g,b,g(b)) \to \varphi_{133}) \to \varphi_{47} \to \varepsilon \in \mathbb{R} \to 0 < \varepsilon \to \varphi_{113}$
\item $\varphi_{247} :\equiv \lvert g(b) - g(b) \rvert = \lvert g(b) - g(b) \rvert \to (g(b) - g(b) = 0 \to \varphi_{76}) \to \mathrm{FLim}(f,c,b) \to \varphi_{179}$
\item $\varphi_{248} :\equiv \varphi_{47} \land \varepsilon \in \mathbb{R} \land 0 < \varepsilon \land \neg \varphi_{113} \to \forall \mathit{dg}\,(\varphi_{47} \land \varepsilon \in \mathbb{R} \land 0 < \varepsilon \land \neg \varphi_{113})$
\item $\varphi_{249} :\equiv \forall \varepsilon\,(\varphi_{47} \to \varepsilon \in \mathbb{R} \to 0 < \varepsilon \to \varphi_{113}) \to \forall \varepsilon\,(\varphi_{47}) \to \forall \varepsilon \in \mathbb{R}\,(0 < \varepsilon \to \varphi_{113})$
\item $\varphi_{250} :\equiv (\mathrm{FLim}(f,c,b) \to \mathrm{FLim}(g,b,g(b)) \to \varphi_{55}) \to \varphi_{58} \to \varphi_{187}$
\item $\varphi_{251} :\equiv \forall \delta\,(\forall \mathit{dg}\,(\neg (\mathit{dg} \in \mathbb{R} \land 0 < \mathit{dg} \land \varphi_{82})) \to \varphi_{106}) \to \forall \delta\,\forall \mathit{dg}\,(\neg (\mathit{dg} \in \mathbb{R} \land 0 < \mathit{dg} \land \varphi_{82})) \to \forall \delta\,(\varphi_{106})$
\item $\varphi_{252} :\equiv (g \circ f(x) = g(f(x)) \to g(f(x)) = g \circ f(x)) \to \mathrm{FLim}(f,c,b) \to \varphi_{232}$
\item $\varphi_{253} :\equiv (\mathrm{FLim}(f,c,b) \to \varphi_{144}) \to \mathrm{FLim}(f,c,b) \to \mathrm{FLim}(g,b,g(b)) \to \varphi_{140}$
\item $\varphi_{254} :\equiv (\mathrm{FLim}(f,c,b) \to \mathrm{FLim}(g,b,g(b)) \to \varphi_{111}) \to \varphi_{217}$
\item $\varphi_{255} :\equiv (\mathrm{FLim}(f,c,b) \to \mathrm{FLim}(g,b,g(b)) \to \varphi_{148}) \to \mathrm{FLim}(f,c,b) \to \mathrm{FLim}(g,b,g(b)) \to \varphi_{145}$
\item $\varphi_{256} :\equiv (\mathrm{FLim}(f,c,b) \to \mathrm{FLim}(g,b,g(b)) \to \varphi_{152}) \to \varphi_{170}$
\item $\varphi_{257} :\equiv (g(f(x)) = g \circ f(x) \to \varphi_{71}) \to \mathrm{FLim}(f,c,b) \to \mathrm{FLim}(g,b,g(b)) \to \varphi_{227}$
\item $\varphi_{258} :\equiv (\mathrm{FLim}(f,c,b) \to \mathrm{FLim}(g,b,g(b)) \to \varphi_{150}) \to \varphi_{177}$
\item $\varphi_{259} :\equiv (\mathrm{FLim}(f,c,b) \to \varphi_{46}) \to (\mathrm{FLim}(f,c,b) \to \varphi_{62}) \to \varphi_{239}$
\item $\varphi_{260} :\equiv (\mathrm{FLim}(f,c,b) \to \mathrm{FLim}(g,b,g(b)) \to \varphi_{140}) \to \varphi_{211}$
\item $\varphi_{261} :\equiv (\mathrm{RFn}(g) \to y \in \mathrm{dom}(g) \to y \in \mathbb{R}) \to (\mathrm{FLim}(f,c,b) \to \varphi_{29}) \to \varphi_{253}$
\item $\varphi_{262} :\equiv (\mathrm{FLim}(f,c,b) \to \mathrm{FLim}(g,b,g(b)) \to \varphi_{140}) \to (\mathrm{FLim}(f,c,b) \to \varphi_{43}) \to \varphi_{176}$
\item $\varphi_{263} :\equiv (\mathrm{FLim}(f,c,b) \to \varphi_{43}) \to \varphi_{171} \to \mathrm{FLim}(f,c,b) \to \mathrm{FLim}(g,b,g(b)) \to \varphi_{153}$
\item $\varphi_{264} :\equiv (\mathrm{FLim}(g,b,g(b)) \to \varepsilon \in \mathbb{R} \to 0 < \varepsilon \to \varphi_{107}) \to \varphi_{6} \to \varphi_{237}$
\item $\varphi_{265} :\equiv \varphi_{210} \to \varphi_{47} \land \varepsilon \in \mathbb{R} \land 0 < \varepsilon \land \neg \varphi_{113} \to \neg (\mathit{dg} \in \mathbb{R} \land 0 < \mathit{dg} \land \varphi_{82})$
\item $\varphi_{266} :\equiv (\varphi_{114} \to \neg (\delta_{2} \in \mathbb{R} \land 0 < \delta_{2} \land \forall x \in \mathrm{dom}(g \circ f)\,(\varphi_{69}))) \to \mathrm{FLim}(f,c,b) \to \varphi_{241}$
\item $\varphi_{267} :\equiv (\mathrm{FLim}(f,c,b) \to \mathrm{FLim}(g,b,g(b)) \to \varphi_{153}) \to \lvert g(y) - g(b) \rvert = \lvert g(y) - g(b) \rvert \to \varphi_{240}$
\item $\varphi_{268} :\equiv (\mathrm{FLim}(f,c,b) \to \mathrm{FLim}(f,c,b)) \to \varphi_{141} \to \varphi_{137} \to \varphi_{195}$
\item $\varphi_{269} :\equiv (\mathrm{FLim}(f,c,b) \to \varphi_{230}) \to \mathrm{FLim}(f,c,b) \to \mathrm{FLim}(g,b,g(b)) \to \varphi_{225}$
\item $\varphi_{270} :\equiv (\mathrm{FLim}(f,c,b) \to \mathrm{FLim}(g,b,g(b)) \to \varphi_{229}) \to \mathrm{FLim}(f,c,b) \to \varphi_{226}$
\item $\varphi_{271} :\equiv (\mathrm{FLim}(f,c,b) \to \mathrm{FLim}(g,b,g(b)) \to \varphi_{227}) \to \varphi_{194} \to x \in \mathrm{dom}(g \circ f) \to \varphi_{69}$
\item $\varphi_{272} :\equiv \forall x\,(\varphi_{194} \to x \in \mathrm{dom}(g \circ f) \to \varphi_{69}) \to \forall x\,(\varphi_{194}) \to \forall x \in \mathrm{dom}(g \circ f)\,(\varphi_{69})$
\item $\varphi_{273} :\equiv (\mathrm{FLim}(f,c,b) \to \mathrm{FLim}(g,b,g(b)) \to \varphi_{140}) \to \varphi_{263}$
\item $\varphi_{274} :\equiv (\mathrm{FLim}(f,c,b) \to \mathrm{FLim}(g,b,g(b)) \to \varphi_{228}) \to \varphi_{243}$
\item $\varphi_{275} :\equiv \varphi_{251} \to \varphi_{160} \to \mathrm{FLim}(f,c,b) \to \mathrm{FLim}(g,b,g(b)) \to \varphi_{133}$
\item $\varphi_{276} :\equiv (\mathrm{FLim}(f,c,b) \to \mathrm{FLim}(g,b,g(b)) \to \varphi_{225}) \to \varphi_{257}$
\item $\varphi_{277} :\equiv (\mathrm{FLim}(f,c,b) \to \varphi_{226}) \to (\varphi_{72} \to \varphi_{66}) \to \mathrm{FLim}(f,c,b) \to \varphi_{223}$
\item $\varphi_{278} :\equiv (\mathrm{FLim}(f,c,b) \to \mathrm{FLim}(g,b,g(b)) \to \varphi_{109}) \to \varphi_{219} \to \varphi_{242}$
\item $\varphi_{279} :\equiv (\mathrm{FLim}(f,c,b) \to \mathrm{FLim}(g,b,g(b)) \to \varphi_{178}) \to (\mathrm{FLim}(f,c,b) \to \varphi_{179}) \to \varphi_{235}$
\item $\varphi_{280} :\equiv (\varphi_{194} \to \forall x\,(\varphi_{194})) \to \mathrm{FLim}(f,c,b) \to \mathrm{FLim}(g,b,g(b)) \to \varphi_{228}$
\item $\varphi_{281} :\equiv \varphi_{220} \to (\mathrm{FLim}(f,c,b) \to \varphi_{224}) \to \mathrm{FLim}(f,c,b) \to \varphi_{230}$
\item $\varphi_{282} :\equiv (\mathrm{FLim}(f,c,b) \to \mathrm{FLim}(g,b,g(b)) \to \varphi_{229}) \to \varphi_{274}$
\item $\varphi_{283} :\equiv (\mathrm{FLim}(f,c,b) \to \mathrm{FLim}(g,b,g(b)) \to \varphi_{145}) \to \varphi_{187} \to \varphi_{189} \to (\mathrm{FLim}(f,c,b) \to \varphi_{144}) \to \varphi_{258}$
\item $\varphi_{284} :\equiv (\varphi_{47} \land \varepsilon \in \mathbb{R} \land 0 < \varepsilon \land (\mathit{dg} \in \mathbb{R} \land 0 < \mathit{dg} \land \varphi_{82}) \land \neg \varphi_{113} \to \varphi_{216}) \to \forall \delta\,(\forall \delta_{2}\,(\neg (\delta_{2} \in \mathbb{R} \land 0 < \delta_{2} \land \varphi_{72})) \to \varphi_{98}) \to \varphi_{245} \to \varphi_{156} \to \varphi_{210}$
\item $\varphi_{285} :\equiv (\varphi_{215} \to \varphi_{218}) \to \varphi_{248} \to \forall \delta\,(\forall \mathit{dg}\,(\neg (\mathit{dg} \in \mathbb{R} \land 0 < \mathit{dg} \land \varphi_{82})) \to \varphi_{106}) \to \varphi_{275}$
\item $\varphi_{286} :\equiv (\mathrm{FLim}(f,c,b) \to \mathrm{FLim}(g,b,g(b)) \to \varphi_{131}) \to \varphi_{215} \to \varphi_{285}$
\item $\varphi_{287} :\equiv \forall \delta_{2}\,(\varphi_{244}) \to (\forall \delta_{2}\,(\varphi_{244}) \to \varphi_{216} \to \forall \delta_{2}\,(\neg (\delta_{2} \in \mathbb{R} \land 0 < \delta_{2} \land \varphi_{72}))) \to \varphi_{284}$
\end{abbrevlist}
\end{fullproof}

% 自動生成 (prover/build.js): FSqueeze  ノード数 371
\begin{fullproof}{定理1.10}{$\forall f\,\forall g\,\forall h\,\forall c\,\forall \alpha\,\forall r\,(\mathrm{FLim}(f,c,\alpha) \to \mathrm{FLim}(g,c,\alpha) \to \mathrm{RFn}(h) \to \mathrm{PNbhd}(h,c) \to r \in \mathbb{R} \to 0 < r \to \forall x \in \mathrm{dom}(h)\,(0 < \lvert x - c \rvert \to \lvert x - c \rvert < r \to x \in \mathrm{dom}(f) \land x \in \mathrm{dom}(g) \land f(x) \leq h(x) \land h(x) \leq g(x)) \to \mathrm{FLim}(h,c,\alpha))$}{371}\phantomsection\label{tree:FSqueeze}
\noindent$\Pi_{1}$:
\begin{ptree}
\AXC{$\mathrm{FLim}(f,c,\alpha) \to \varphi_{230}$\lbl{A1}}
\AXC{$(\mathrm{FLim}(f,c,\alpha) \to \varphi_{230}) \to \mathrm{FLim}(f,c,\alpha) \to \varphi_{266}$\lbl{A1}}
\BIC{$\mathrm{FLim}(f,c,\alpha) \to \varphi_{266}$}
\end{ptree}
\noindent$\Pi_{2}$:
\begin{ptree}
\AXC{$\mathrm{FLim}(f,c,\alpha) \to \varphi_{260}$\lbl{A1}}
\AXC{$(\mathrm{FLim}(f,c,\alpha) \to \varphi_{260}) \to \mathrm{FLim}(f,c,\alpha) \to \varphi_{267}$\lbl{A1}}
\BIC{$\mathrm{FLim}(f,c,\alpha) \to \varphi_{267}$}
\end{ptree}
\noindent$\Pi_{3}$:
\begin{ptree}
\AXC{$\forall x\,\forall y\,\forall z\,(x < y \to y \leq z \to x < z)$\lbl{[O3]}}
\AXC{$\varphi_{40}$\lbl{A2}}
\BIC{$\forall y\,\forall z\,(\lvert x - c \rvert < y \to y \leq z \to \lvert x - c \rvert < z)$}
\AXC{$\varphi_{67}$\lbl{A2}}
\BIC{$\forall z\,(\lvert x - c \rvert < \delta \to \delta \leq z \to \lvert x - c \rvert < z)$}
\AXC{$\varphi_{61}$\lbl{A2}}
\BIC{$\lvert x - c \rvert < \delta \to \delta \leq d_{3} \to \lvert x - c \rvert < d_{3}$}
\end{ptree}
\noindent$\Pi_{4}$:
\begin{ptree}
\AXC{$\mathrm{FLim}(f,c,\alpha) \to \varphi_{273}$\lbl{A1}}
\AXC{$\stackrel{\Pi_{3}}{\vdots}$}\noLine
\UIC{$\lvert x - c \rvert < \delta \to \delta \leq d_{3} \to \lvert x - c \rvert < d_{3}$}
\AXC{$(\lvert x - c \rvert < \delta \to \delta \leq d_{3} \to \lvert x - c \rvert < d_{3}) \to \varphi_{323}$\lbl{A1}}
\BIC{$\varphi_{323}$}
\BIC{$(\mathrm{FLim}(f,c,\alpha) \to \varphi_{267}) \to \mathrm{FLim}(f,c,\alpha) \to \varphi_{274}$}
\end{ptree}
\noindent$\Pi_{5}$:
\begin{ptree}
\AXC{$\stackrel{\Pi_{2}}{\vdots}$}\noLine
\UIC{$\mathrm{FLim}(f,c,\alpha) \to \varphi_{267}$}
\AXC{$\stackrel{\Pi_{4}}{\vdots}$}\noLine
\UIC{$(\mathrm{FLim}(f,c,\alpha) \to \varphi_{267}) \to \mathrm{FLim}(f,c,\alpha) \to \varphi_{274}$}
\BIC{$\mathrm{FLim}(f,c,\alpha) \to \varphi_{274}$}
\end{ptree}
\noindent$\Pi_{6}$:
\begin{ptree}
\AXC{$\forall x\,\forall y\,\forall z\,(x < y \to y \leq z \to x < z)$\lbl{[O3]}}
\AXC{$\varphi_{40}$\lbl{A2}}
\BIC{$\forall y\,\forall z\,(\lvert x - c \rvert < y \to y \leq z \to \lvert x - c \rvert < z)$}
\AXC{$\varphi_{68}$\lbl{A2}}
\BIC{$\forall z\,(\lvert x - c \rvert < d_{3} \to d_{3} \leq z \to \lvert x - c \rvert < z)$}
\AXC{$\varphi_{62}$\lbl{A2}}
\BIC{$\lvert x - c \rvert < d_{3} \to d_{3} \leq d_{2} \to \lvert x - c \rvert < d_{2}$}
\end{ptree}
\noindent$\Pi_{7}$:
\begin{ptree}
\AXC{$\stackrel{\Pi_{6}}{\vdots}$}\noLine
\UIC{$\lvert x - c \rvert < d_{3} \to d_{3} \leq d_{2} \to \lvert x - c \rvert < d_{2}$}
\AXC{$(\lvert x - c \rvert < d_{3} \to d_{3} \leq d_{2} \to \lvert x - c \rvert < d_{2}) \to \varphi_{324}$\lbl{A1}}
\BIC{$\varphi_{324}$}
\end{ptree}
\noindent$\Pi_{8}$:
\begin{ptree}
\AXC{$\stackrel{\Pi_{1}}{\vdots}$}\noLine
\UIC{$\mathrm{FLim}(f,c,\alpha) \to \varphi_{266}$}
\AXC{$\stackrel{\Pi_{5}}{\vdots}$}\noLine
\UIC{$\mathrm{FLim}(f,c,\alpha) \to \varphi_{274}$}
\AXC{$\stackrel{\Pi_{7}}{\vdots}$}\noLine
\UIC{$\varphi_{324}$}
\BIC{$(\mathrm{FLim}(f,c,\alpha) \to \varphi_{266}) \to \mathrm{FLim}(f,c,\alpha) \to \varphi_{275}$}
\BIC{$\mathrm{FLim}(f,c,\alpha) \to \varphi_{275}$}
\end{ptree}
\noindent$\Pi_{9}$:
\begin{ptree}
\AXC{$\mathrm{FLim}(f,c,\alpha) \to \varphi_{260}$\lbl{A1}}
\AXC{$(\mathrm{FLim}(f,c,\alpha) \to \varphi_{260}) \to \mathrm{FLim}(f,c,\alpha) \to \varphi_{265}$\lbl{A1}}
\BIC{$\mathrm{FLim}(f,c,\alpha) \to \varphi_{265}$}
\end{ptree}
\noindent$\Pi_{10}$:
\begin{ptree}
\AXC{$\forall x\,\forall y\,\forall z\,(x < y \to y \leq z \to x < z)$\lbl{[O3]}}
\AXC{$\varphi_{40}$\lbl{A2}}
\BIC{$\forall y\,\forall z\,(\lvert x - c \rvert < y \to y \leq z \to \lvert x - c \rvert < z)$}
\AXC{$\varphi_{67}$\lbl{A2}}
\BIC{$\forall z\,(\lvert x - c \rvert < \delta \to \delta \leq z \to \lvert x - c \rvert < z)$}
\AXC{$\varphi_{60}$\lbl{A2}}
\BIC{$\lvert x - c \rvert < \delta \to \delta \leq r \to \lvert x - c \rvert < r$}
\end{ptree}
\noindent$\Pi_{11}$:
\begin{ptree}
\AXC{$\mathrm{FLim}(f,c,\alpha) \to \varphi_{273}$\lbl{A1}}
\AXC{$\stackrel{\Pi_{10}}{\vdots}$}\noLine
\UIC{$\lvert x - c \rvert < \delta \to \delta \leq r \to \lvert x - c \rvert < r$}
\AXC{$(\lvert x - c \rvert < \delta \to \delta \leq r \to \lvert x - c \rvert < r) \to \varphi_{322}$\lbl{A1}}
\BIC{$\varphi_{322}$}
\BIC{$(\mathrm{FLim}(f,c,\alpha) \to \varphi_{265}) \to \mathrm{FLim}(f,c,\alpha) \to \varphi_{272}$}
\end{ptree}
\noindent$\Pi_{12}$:
\begin{ptree}
\AXC{$\stackrel{\Pi_{9}}{\vdots}$}\noLine
\UIC{$\mathrm{FLim}(f,c,\alpha) \to \varphi_{265}$}
\AXC{$\stackrel{\Pi_{11}}{\vdots}$}\noLine
\UIC{$(\mathrm{FLim}(f,c,\alpha) \to \varphi_{265}) \to \mathrm{FLim}(f,c,\alpha) \to \varphi_{272}$}
\BIC{$\mathrm{FLim}(f,c,\alpha) \to \varphi_{272}$}
\end{ptree}
\noindent$\Pi_{13}$:
\begin{ptree}
\AXC{$\mathrm{FLim}(f,c,\alpha) \to \varphi_{257}$\lbl{A1}}
\AXC{$\varphi_{143}$\lbl{A2}}
\AXC{$\mathrm{FLim}(f,c,\alpha) \to \varphi_{157}$\lbl{A1}}
\AXC{$\varphi_{301}$\lbl{A1}}
\BIC{$\varphi_{143} \to \mathrm{FLim}(f,c,\alpha) \to \varphi_{154}$}
\BIC{$\mathrm{FLim}(f,c,\alpha) \to \varphi_{154}$}
\AXC{$\varphi_{335}$\lbl{A1}}
\BIC{$(\mathrm{FLim}(f,c,\alpha) \to \varphi_{257}) \to \varphi_{331}$}
\BIC{$\varphi_{331}$}
\end{ptree}
\noindent$\Pi_{14}$:
\begin{ptree}
\AXC{$\stackrel{\Pi_{12}}{\vdots}$}\noLine
\UIC{$\mathrm{FLim}(f,c,\alpha) \to \varphi_{272}$}
\AXC{$\mathrm{FLim}(f,c,\alpha) \to \varphi_{263}$\lbl{A1}}
\AXC{$\stackrel{\Pi_{13}}{\vdots}$}\noLine
\UIC{$\varphi_{331}$}
\BIC{$(\mathrm{FLim}(f,c,\alpha) \to \varphi_{272}) \to \mathrm{FLim}(f,c,\alpha) \to \varphi_{297}$}
\BIC{$\mathrm{FLim}(f,c,\alpha) \to \varphi_{297}$}
\end{ptree}
\noindent$\Pi_{15}$:
\begin{ptree}
\AXC{$\varphi_{48} \to \varphi_{41}$\lbl{A2}}
\AXC{$\mathrm{FLim}(f,c,\alpha) \to \varphi_{258}$\lbl{A1}}
\AXC{$(\mathrm{FLim}(f,c,\alpha) \to \varphi_{258}) \to \mathrm{FLim}(f,c,\alpha) \to \varphi_{288}$\lbl{A1}}
\BIC{$\mathrm{FLim}(f,c,\alpha) \to \varphi_{288}$}
\AXC{$\varphi_{315}$\lbl{A1}}
\BIC{$(\varphi_{48} \to \varphi_{41}) \to \mathrm{FLim}(f,c,\alpha) \to \varphi_{283}$}
\BIC{$\mathrm{FLim}(f,c,\alpha) \to \varphi_{283}$}
\end{ptree}
\noindent$\Pi_{16}$:
\begin{ptree}
\AXC{$\mathrm{FLim}(f,c,\alpha) \to \varphi_{263}$\lbl{A1}}
\AXC{$\stackrel{\Pi_{14}}{\vdots}$}\noLine
\UIC{$\mathrm{FLim}(f,c,\alpha) \to \varphi_{297}$}
\AXC{$\stackrel{\Pi_{15}}{\vdots}$}\noLine
\UIC{$\mathrm{FLim}(f,c,\alpha) \to \varphi_{283}$}
\AXC{$\varphi_{339}$\lbl{A1}}
\BIC{$(\mathrm{FLim}(f,c,\alpha) \to \varphi_{297}) \to \varphi_{325}$}
\BIC{$\varphi_{325}$}
\BIC{$(\mathrm{FLim}(f,c,\alpha) \to \varphi_{275}) \to \mathrm{FLim}(f,c,\alpha) \to \varphi_{284}$}
\end{ptree}
\noindent$\Pi_{17}$:
\begin{ptree}
\AXC{$\stackrel{\Pi_{8}}{\vdots}$}\noLine
\UIC{$\mathrm{FLim}(f,c,\alpha) \to \varphi_{275}$}
\AXC{$\stackrel{\Pi_{16}}{\vdots}$}\noLine
\UIC{$(\mathrm{FLim}(f,c,\alpha) \to \varphi_{275}) \to \mathrm{FLim}(f,c,\alpha) \to \varphi_{284}$}
\BIC{$\mathrm{FLim}(f,c,\alpha) \to \varphi_{284}$}
\end{ptree}
\noindent$\Pi_{18}$:
\begin{ptree}
\AXC{$\mathrm{FLim}(f,c,\alpha) \to \varphi_{230}$\lbl{A1}}
\AXC{$(\mathrm{FLim}(f,c,\alpha) \to \varphi_{230}) \to \mathrm{FLim}(f,c,\alpha) \to \varphi_{268}$\lbl{A1}}
\BIC{$\mathrm{FLim}(f,c,\alpha) \to \varphi_{268}$}
\end{ptree}
\noindent$\Pi_{19}$:
\begin{ptree}
\AXC{$\forall x\,\forall y\,\forall z\,(x < y \to y \leq z \to x < z)$\lbl{[O3]}}
\AXC{$\varphi_{40}$\lbl{A2}}
\BIC{$\forall y\,\forall z\,(\lvert x - c \rvert < y \to y \leq z \to \lvert x - c \rvert < z)$}
\AXC{$\varphi_{68}$\lbl{A2}}
\BIC{$\forall z\,(\lvert x - c \rvert < d_{3} \to d_{3} \leq z \to \lvert x - c \rvert < z)$}
\AXC{$\varphi_{63}$\lbl{A2}}
\BIC{$\lvert x - c \rvert < d_{3} \to d_{3} \leq d_{1} \to \lvert x - c \rvert < d_{1}$}
\end{ptree}
\noindent$\Pi_{20}$:
\begin{ptree}
\AXC{$\stackrel{\Pi_{19}}{\vdots}$}\noLine
\UIC{$\lvert x - c \rvert < d_{3} \to d_{3} \leq d_{1} \to \lvert x - c \rvert < d_{1}$}
\AXC{$(\lvert x - c \rvert < d_{3} \to d_{3} \leq d_{1} \to \lvert x - c \rvert < d_{1}) \to \varphi_{326}$\lbl{A1}}
\BIC{$\varphi_{326}$}
\end{ptree}
\noindent$\Pi_{21}$:
\begin{ptree}
\AXC{$\stackrel{\Pi_{18}}{\vdots}$}\noLine
\UIC{$\mathrm{FLim}(f,c,\alpha) \to \varphi_{268}$}
\AXC{$\stackrel{\Pi_{5}}{\vdots}$}\noLine
\UIC{$\mathrm{FLim}(f,c,\alpha) \to \varphi_{274}$}
\AXC{$\stackrel{\Pi_{20}}{\vdots}$}\noLine
\UIC{$\varphi_{326}$}
\BIC{$(\mathrm{FLim}(f,c,\alpha) \to \varphi_{268}) \to \mathrm{FLim}(f,c,\alpha) \to \varphi_{276}$}
\BIC{$\mathrm{FLim}(f,c,\alpha) \to \varphi_{276}$}
\end{ptree}
\noindent$\Pi_{22}$:
\begin{ptree}
\AXC{$\varphi_{49} \to \varphi_{43}$\lbl{A2}}
\AXC{$\mathrm{FLim}(f,c,\alpha) \to \varphi_{206}$\lbl{A1}}
\AXC{$(\mathrm{FLim}(f,c,\alpha) \to \varphi_{206}) \to \mathrm{FLim}(f,c,\alpha) \to \varphi_{289}$\lbl{A1}}
\BIC{$\mathrm{FLim}(f,c,\alpha) \to \varphi_{289}$}
\AXC{$\varphi_{316}$\lbl{A1}}
\BIC{$(\varphi_{49} \to \varphi_{43}) \to \mathrm{FLim}(f,c,\alpha) \to \varphi_{285}$}
\BIC{$\mathrm{FLim}(f,c,\alpha) \to \varphi_{285}$}
\end{ptree}
\noindent$\Pi_{23}$:
\begin{ptree}
\AXC{$\mathrm{FLim}(f,c,\alpha) \to \varphi_{263}$\lbl{A1}}
\AXC{$\stackrel{\Pi_{14}}{\vdots}$}\noLine
\UIC{$\mathrm{FLim}(f,c,\alpha) \to \varphi_{297}$}
\AXC{$\stackrel{\Pi_{22}}{\vdots}$}\noLine
\UIC{$\mathrm{FLim}(f,c,\alpha) \to \varphi_{285}$}
\AXC{$\varphi_{340}$\lbl{A1}}
\BIC{$(\mathrm{FLim}(f,c,\alpha) \to \varphi_{297}) \to \varphi_{327}$}
\BIC{$\varphi_{327}$}
\BIC{$(\mathrm{FLim}(f,c,\alpha) \to \varphi_{276}) \to \mathrm{FLim}(f,c,\alpha) \to \varphi_{286}$}
\end{ptree}
\noindent$\Pi_{24}$:
\begin{ptree}
\AXC{$\stackrel{\Pi_{14}}{\vdots}$}\noLine
\UIC{$\mathrm{FLim}(f,c,\alpha) \to \varphi_{297}$}
\AXC{$(\mathrm{FLim}(f,c,\alpha) \to \varphi_{297}) \to \mathrm{FLim}(f,c,\alpha) \to \varphi_{278}$\lbl{A1}}
\BIC{$\mathrm{FLim}(f,c,\alpha) \to \varphi_{278}$}
\end{ptree}
\noindent$\Pi_{25}$:
\begin{ptree}
\AXC{$\forall f\,\forall c\,\forall \alpha\,(\mathrm{FLim}(f,c,\alpha) \to \alpha \in \mathbb{R})$\lbl{[D4]}}
\AXC{$\varphi_{34}$\lbl{A2}}
\BIC{$\forall c\,\forall \alpha\,(\mathrm{FLim}(f,c,\alpha) \to \alpha \in \mathbb{R})$}
\AXC{$\varphi_{17}$\lbl{A2}}
\BIC{$\forall \alpha\,(\mathrm{FLim}(f,c,\alpha) \to \alpha \in \mathbb{R})$}
\AXC{$\varphi_{11}$\lbl{A2}}
\BIC{$\mathrm{FLim}(f,c,\alpha) \to \alpha \in \mathbb{R}$}
\end{ptree}
\noindent$\Pi_{26}$:
\begin{ptree}
\AXC{$\mathrm{FLim}(f,c,\alpha) \to \mathrm{FLim}(f,c,\alpha)$\lbl{A1}}
\AXC{$\stackrel{\Pi_{25}}{\vdots}$}\noLine
\UIC{$\mathrm{FLim}(f,c,\alpha) \to \alpha \in \mathbb{R}$}
\AXC{$(\mathrm{FLim}(f,c,\alpha) \to \alpha \in \mathbb{R}) \to \varphi_{141}$\lbl{A1}}
\BIC{$\varphi_{141}$}
\BIC{$\mathrm{FLim}(f,c,\alpha) \to \varphi_{119}$}
\end{ptree}
\noindent$\Pi_{27}$:
\begin{ptree}
\AXC{$\forall f\,\forall x\,(\mathrm{RFn}(f) \to x \in \mathrm{dom}(f) \to f(x) \in \mathbb{R})$\lbl{[S6]}}
\AXC{$\varphi_{64}$\lbl{A2}}
\BIC{$\forall x\,(\mathrm{RFn}(h) \to x \in \mathrm{dom}(h) \to h(x) \in \mathbb{R})$}
\AXC{$\varphi_{55}$\lbl{A2}}
\BIC{$\mathrm{RFn}(h) \to x \in \mathrm{dom}(h) \to h(x) \in \mathbb{R}$}
\end{ptree}
\noindent$\Pi_{28}$:
\begin{ptree}
\AXC{$\varphi_{2}$\lbl{A1}}
\AXC{$\stackrel{\Pi_{27}}{\vdots}$}\noLine
\UIC{$\mathrm{RFn}(h) \to x \in \mathrm{dom}(h) \to h(x) \in \mathbb{R}$}
\AXC{$(\mathrm{RFn}(h) \to x \in \mathrm{dom}(h) \to h(x) \in \mathbb{R}) \to \varphi_{309}$\lbl{A1}}
\BIC{$\varphi_{309}$}
\BIC{$(\mathrm{FLim}(f,c,\alpha) \to \varphi_{257}) \to \mathrm{FLim}(f,c,\alpha) \to \varphi_{279}$}
\end{ptree}
\noindent$\Pi_{29}$:
\begin{ptree}
\AXC{$\stackrel{\Pi_{14}}{\vdots}$}\noLine
\UIC{$\mathrm{FLim}(f,c,\alpha) \to \varphi_{297}$}
\AXC{$(\mathrm{FLim}(f,c,\alpha) \to \varphi_{297}) \to \mathrm{FLim}(f,c,\alpha) \to \varphi_{270}$\lbl{A1}}
\BIC{$\mathrm{FLim}(f,c,\alpha) \to \varphi_{270}$}
\end{ptree}
\noindent$\Pi_{30}$:
\begin{ptree}
\AXC{$\varphi_{1}$\lbl{A1}}
\AXC{$\forall f\,\forall c\,\forall \alpha\,(\mathrm{FLim}(f,c,\alpha) \to \mathrm{RFn}(f))$\lbl{[D2]}}
\AXC{$\varphi_{13}$\lbl{A2}}
\BIC{$\forall c\,\forall \alpha\,(\mathrm{FLim}(g,c,\alpha) \to \mathrm{RFn}(g))$}
\AXC{$\varphi_{7}$\lbl{A2}}
\BIC{$\forall \alpha\,(\mathrm{FLim}(g,c,\alpha) \to \mathrm{RFn}(g))$}
\AXC{$\varphi_{5}$\lbl{A2}}
\BIC{$\mathrm{FLim}(g,c,\alpha) \to \mathrm{RFn}(g)$}
\AXC{$\varphi_{152}$\lbl{A1}}
\BIC{$\varphi_{1} \to \mathrm{FLim}(f,c,\alpha) \to \varphi_{116}$}
\BIC{$\mathrm{FLim}(f,c,\alpha) \to \varphi_{116}$}
\end{ptree}
\noindent$\Pi_{31}$:
\begin{ptree}
\AXC{$\forall f\,\forall x\,(\mathrm{RFn}(f) \to x \in \mathrm{dom}(f) \to f(x) \in \mathbb{R})$\lbl{[S6]}}
\AXC{$\varphi_{65}$\lbl{A2}}
\BIC{$\forall x\,(\mathrm{RFn}(g) \to x \in \mathrm{dom}(g) \to g(x) \in \mathbb{R})$}
\AXC{$\varphi_{56}$\lbl{A2}}
\BIC{$\mathrm{RFn}(g) \to x \in \mathrm{dom}(g) \to g(x) \in \mathbb{R}$}
\end{ptree}
\noindent$\Pi_{32}$:
\begin{ptree}
\AXC{$\stackrel{\Pi_{31}}{\vdots}$}\noLine
\UIC{$\mathrm{RFn}(g) \to x \in \mathrm{dom}(g) \to g(x) \in \mathbb{R}$}
\AXC{$(\mathrm{RFn}(g) \to x \in \mathrm{dom}(g) \to g(x) \in \mathbb{R}) \to \varphi_{318}$\lbl{A1}}
\BIC{$\varphi_{318}$}
\end{ptree}
\noindent$\Pi_{33}$:
\begin{ptree}
\AXC{$\stackrel{\Pi_{29}}{\vdots}$}\noLine
\UIC{$\mathrm{FLim}(f,c,\alpha) \to \varphi_{270}$}
\AXC{$\stackrel{\Pi_{30}}{\vdots}$}\noLine
\UIC{$\mathrm{FLim}(f,c,\alpha) \to \varphi_{116}$}
\AXC{$\stackrel{\Pi_{32}}{\vdots}$}\noLine
\UIC{$\varphi_{318}$}
\BIC{$(\mathrm{FLim}(f,c,\alpha) \to \varphi_{270}) \to \mathrm{FLim}(f,c,\alpha) \to \varphi_{280}$}
\BIC{$\mathrm{FLim}(f,c,\alpha) \to \varphi_{280}$}
\end{ptree}
\noindent$\Pi_{34}$:
\begin{ptree}
\AXC{$\stackrel{\Pi_{14}}{\vdots}$}\noLine
\UIC{$\mathrm{FLim}(f,c,\alpha) \to \varphi_{297}$}
\AXC{$(\mathrm{FLim}(f,c,\alpha) \to \varphi_{297}) \to \mathrm{FLim}(f,c,\alpha) \to \varphi_{271}$\lbl{A1}}
\BIC{$\mathrm{FLim}(f,c,\alpha) \to \varphi_{271}$}
\end{ptree}
\noindent$\Pi_{35}$:
\begin{ptree}
\AXC{$\forall f\,\forall c\,\forall \alpha\,(\mathrm{FLim}(f,c,\alpha) \to \mathrm{RFn}(f))$\lbl{[D2]}}
\AXC{$\varphi_{14}$\lbl{A2}}
\BIC{$\forall c\,\forall \alpha\,(\mathrm{FLim}(f,c,\alpha) \to \mathrm{RFn}(f))$}
\AXC{$\varphi_{8}$\lbl{A2}}
\BIC{$\forall \alpha\,(\mathrm{FLim}(f,c,\alpha) \to \mathrm{RFn}(f))$}
\AXC{$\varphi_{6}$\lbl{A2}}
\BIC{$\mathrm{FLim}(f,c,\alpha) \to \mathrm{RFn}(f)$}
\end{ptree}
\noindent$\Pi_{36}$:
\begin{ptree}
\AXC{$\stackrel{\Pi_{35}}{\vdots}$}\noLine
\UIC{$\mathrm{FLim}(f,c,\alpha) \to \mathrm{RFn}(f)$}
\AXC{$(\mathrm{FLim}(f,c,\alpha) \to \mathrm{RFn}(f)) \to \varphi_{136}$\lbl{A1}}
\BIC{$\varphi_{136}$}
\end{ptree}
\noindent$\Pi_{37}$:
\begin{ptree}
\AXC{$\forall f\,\forall x\,(\mathrm{RFn}(f) \to x \in \mathrm{dom}(f) \to f(x) \in \mathbb{R})$\lbl{[S6]}}
\AXC{$\varphi_{66}$\lbl{A2}}
\BIC{$\forall x\,(\mathrm{RFn}(f) \to x \in \mathrm{dom}(f) \to f(x) \in \mathbb{R})$}
\AXC{$\varphi_{57}$\lbl{A2}}
\BIC{$\mathrm{RFn}(f) \to x \in \mathrm{dom}(f) \to f(x) \in \mathbb{R}$}
\end{ptree}
\noindent$\Pi_{38}$:
\begin{ptree}
\AXC{$\stackrel{\Pi_{37}}{\vdots}$}\noLine
\UIC{$\mathrm{RFn}(f) \to x \in \mathrm{dom}(f) \to f(x) \in \mathbb{R}$}
\AXC{$(\mathrm{RFn}(f) \to x \in \mathrm{dom}(f) \to f(x) \in \mathbb{R}) \to \varphi_{319}$\lbl{A1}}
\BIC{$\varphi_{319}$}
\end{ptree}
\noindent$\Pi_{39}$:
\begin{ptree}
\AXC{$\stackrel{\Pi_{34}}{\vdots}$}\noLine
\UIC{$\mathrm{FLim}(f,c,\alpha) \to \varphi_{271}$}
\AXC{$\mathrm{FLim}(f,c,\alpha) \to \mathrm{FLim}(f,c,\alpha)$\lbl{A1}}
\AXC{$\stackrel{\Pi_{36}}{\vdots}$}\noLine
\UIC{$\varphi_{136}$}
\BIC{$\mathrm{FLim}(f,c,\alpha) \to \varphi_{117}$}
\AXC{$\stackrel{\Pi_{38}}{\vdots}$}\noLine
\UIC{$\varphi_{319}$}
\BIC{$(\mathrm{FLim}(f,c,\alpha) \to \varphi_{271}) \to \mathrm{FLim}(f,c,\alpha) \to \varphi_{281}$}
\BIC{$\mathrm{FLim}(f,c,\alpha) \to \varphi_{281}$}
\end{ptree}
\noindent$\Pi_{40}$:
\begin{ptree}
\AXC{$\varphi_{107}$\lbl{[C12]}}
\AXC{$\varphi_{107} \to \varphi_{114}$\lbl{A2}}
\BIC{$\varphi_{114}$}
\AXC{$\varphi_{114} \to \varphi_{120}$\lbl{A2}}
\BIC{$\varphi_{120}$}
\AXC{$\varphi_{120} \to \varphi_{128}$\lbl{A2}}
\BIC{$\varphi_{128}$}
\AXC{$\varphi_{214}$\lbl{A2}}
\BIC{$\forall \varepsilon\,(f(x) \in \mathbb{R} \to g(x) \in \mathbb{R} \to h(x) \in \mathbb{R} \to \varphi_{98})$}
\AXC{$\varphi_{209}$\lbl{A2}}
\BIC{$f(x) \in \mathbb{R} \to g(x) \in \mathbb{R} \to h(x) \in \mathbb{R} \to \varphi_{98}$}
\AXC{$\varphi_{345}$\lbl{A1}}
\BIC{$(\mathrm{FLim}(f,c,\alpha) \to \varphi_{281}) \to \varphi_{344}$}
\end{ptree}
\noindent$\Pi_{41}$:
\begin{ptree}
\AXC{$\stackrel{\Pi_{33}}{\vdots}$}\noLine
\UIC{$\mathrm{FLim}(f,c,\alpha) \to \varphi_{280}$}
\AXC{$\stackrel{\Pi_{39}}{\vdots}$}\noLine
\UIC{$\mathrm{FLim}(f,c,\alpha) \to \varphi_{281}$}
\AXC{$\stackrel{\Pi_{40}}{\vdots}$}\noLine
\UIC{$(\mathrm{FLim}(f,c,\alpha) \to \varphi_{281}) \to \varphi_{344}$}
\BIC{$\varphi_{344}$}
\BIC{$(\mathrm{FLim}(f,c,\alpha) \to \varphi_{279}) \to \varphi_{341}$}
\end{ptree}
\noindent$\Pi_{42}$:
\begin{ptree}
\AXC{$\mathrm{FLim}(f,c,\alpha) \to \varphi_{257}$\lbl{A1}}
\AXC{$\stackrel{\Pi_{28}}{\vdots}$}\noLine
\UIC{$(\mathrm{FLim}(f,c,\alpha) \to \varphi_{257}) \to \mathrm{FLim}(f,c,\alpha) \to \varphi_{279}$}
\BIC{$\mathrm{FLim}(f,c,\alpha) \to \varphi_{279}$}
\AXC{$\stackrel{\Pi_{41}}{\vdots}$}\noLine
\UIC{$(\mathrm{FLim}(f,c,\alpha) \to \varphi_{279}) \to \varphi_{341}$}
\BIC{$\varphi_{341}$}
\end{ptree}
\noindent$\Pi_{43}$:
\begin{ptree}
\AXC{$\stackrel{\Pi_{24}}{\vdots}$}\noLine
\UIC{$\mathrm{FLim}(f,c,\alpha) \to \varphi_{278}$}
\AXC{$\mathrm{FLim}(f,c,\alpha) \to \varphi_{124}$\lbl{A1}}
\AXC{$\stackrel{\Pi_{26}}{\vdots}$}\noLine
\UIC{$\mathrm{FLim}(f,c,\alpha) \to \varphi_{119}$}
\AXC{$\stackrel{\Pi_{42}}{\vdots}$}\noLine
\UIC{$\varphi_{341}$}
\BIC{$(\mathrm{FLim}(f,c,\alpha) \to \varphi_{124}) \to \varphi_{338}$}
\BIC{$\varphi_{338}$}
\BIC{$(\mathrm{FLim}(f,c,\alpha) \to \varphi_{277}) \to \varphi_{330}$}
\end{ptree}
\noindent$\Pi_{44}$:
\begin{ptree}
\AXC{$\stackrel{\Pi_{14}}{\vdots}$}\noLine
\UIC{$\mathrm{FLim}(f,c,\alpha) \to \varphi_{297}$}
\AXC{$(\mathrm{FLim}(f,c,\alpha) \to \varphi_{297}) \to \mathrm{FLim}(f,c,\alpha) \to \varphi_{277}$\lbl{A1}}
\BIC{$\mathrm{FLim}(f,c,\alpha) \to \varphi_{277}$}
\AXC{$\stackrel{\Pi_{43}}{\vdots}$}\noLine
\UIC{$(\mathrm{FLim}(f,c,\alpha) \to \varphi_{277}) \to \varphi_{330}$}
\BIC{$\varphi_{330}$}
\end{ptree}
\noindent$\Pi_{45}$:
\begin{ptree}
\AXC{$\stackrel{\Pi_{21}}{\vdots}$}\noLine
\UIC{$\mathrm{FLim}(f,c,\alpha) \to \varphi_{276}$}
\AXC{$\stackrel{\Pi_{23}}{\vdots}$}\noLine
\UIC{$(\mathrm{FLim}(f,c,\alpha) \to \varphi_{276}) \to \mathrm{FLim}(f,c,\alpha) \to \varphi_{286}$}
\BIC{$\mathrm{FLim}(f,c,\alpha) \to \varphi_{286}$}
\AXC{$\stackrel{\Pi_{44}}{\vdots}$}\noLine
\UIC{$\varphi_{330}$}
\BIC{$(\mathrm{FLim}(f,c,\alpha) \to \varphi_{284}) \to \mathrm{FLim}(f,c,\alpha) \to \varphi_{287}$}
\end{ptree}
\noindent$\Pi_{46}$:
\begin{ptree}
\AXC{$\stackrel{\Pi_{17}}{\vdots}$}\noLine
\UIC{$\mathrm{FLim}(f,c,\alpha) \to \varphi_{284}$}
\AXC{$\stackrel{\Pi_{45}}{\vdots}$}\noLine
\UIC{$(\mathrm{FLim}(f,c,\alpha) \to \varphi_{284}) \to \mathrm{FLim}(f,c,\alpha) \to \varphi_{287}$}
\BIC{$\mathrm{FLim}(f,c,\alpha) \to \varphi_{287}$}
\AXC{$(\mathrm{FLim}(f,c,\alpha) \to \varphi_{287}) \to \varphi_{238} \to \varphi_{45}$\lbl{A1}}
\BIC{$\varphi_{238} \to \varphi_{45}$}
\UIC{$\forall x\,(\varphi_{238} \to \varphi_{45})$}
\end{ptree}
\noindent$\Pi_{47}$:
\begin{ptree}
\AXC{$\varphi_{238} \to \forall x\,(\varphi_{238})$\lbl{A4}}
\AXC{$\forall x\,(\varphi_{238} \to \varphi_{45}) \to \forall x\,(\varphi_{238}) \to \varphi_{47}$\lbl{A3}}
\AXC{$\stackrel{\Pi_{46}}{\vdots}$}\noLine
\UIC{$\forall x\,(\varphi_{238} \to \varphi_{45})$}
\AXC{$\forall x\,(\varphi_{238} \to \varphi_{45}) \to \varphi_{337}$\lbl{A1}}
\BIC{$\varphi_{337}$}
\BIC{$(\varphi_{238} \to \forall x\,(\varphi_{238})) \to \mathrm{FLim}(f,c,\alpha) \to \varphi_{290}$}
\BIC{$\mathrm{FLim}(f,c,\alpha) \to \varphi_{290}$}
\end{ptree}
\noindent$\Pi_{48}$:
\begin{ptree}
\AXC{$\varphi_{129}$\lbl{A2}}
\AXC{$\stackrel{\Pi_{47}}{\vdots}$}\noLine
\UIC{$\mathrm{FLim}(f,c,\alpha) \to \varphi_{290}$}
\AXC{$\mathrm{FLim}(f,c,\alpha) \to \varphi_{260}$\lbl{A1}}
\AXC{$\varphi_{329}$\lbl{A1}}
\BIC{$(\mathrm{FLim}(f,c,\alpha) \to \varphi_{290}) \to \mathrm{FLim}(f,c,\alpha) \to \varphi_{294}$}
\BIC{$\mathrm{FLim}(f,c,\alpha) \to \varphi_{294}$}
\AXC{$\varphi_{320}$\lbl{A1}}
\BIC{$\varphi_{129} \to \mathrm{FLim}(f,c,\alpha) \to \varphi_{295}$}
\BIC{$\mathrm{FLim}(f,c,\alpha) \to \varphi_{295}$}
\AXC{$\varphi_{312}$\lbl{A1}}
\BIC{$\varphi_{291} \to \neg (\delta \in \mathbb{R} \land 0 < \delta \land \delta \leq d_{3} \land \delta \leq r)$}
\UIC{$\forall \delta\,(\varphi_{291} \to \neg (\delta \in \mathbb{R} \land 0 < \delta \land \delta \leq d_{3} \land \delta \leq r))$}
\end{ptree}
\noindent$\Pi_{49}$:
\begin{ptree}
\AXC{$\mathrm{FLim}(f,c,\alpha) \to \varphi_{230}$\lbl{A1}}
\AXC{$(\mathrm{FLim}(f,c,\alpha) \to \varphi_{230}) \to \mathrm{FLim}(f,c,\alpha) \to \varphi_{223}$\lbl{A1}}
\BIC{$\mathrm{FLim}(f,c,\alpha) \to \varphi_{223}$}
\end{ptree}
\noindent$\Pi_{50}$:
\begin{ptree}
\AXC{$\mathrm{FLim}(f,c,\alpha) \to \varphi_{230}$\lbl{A1}}
\AXC{$(\mathrm{FLim}(f,c,\alpha) \to \varphi_{230}) \to \mathrm{FLim}(f,c,\alpha) \to \varphi_{225}$\lbl{A1}}
\BIC{$\mathrm{FLim}(f,c,\alpha) \to \varphi_{225}$}
\end{ptree}
\noindent$\Pi_{51}$:
\begin{ptree}
\AXC{$\stackrel{\Pi_{50}}{\vdots}$}\noLine
\UIC{$\mathrm{FLim}(f,c,\alpha) \to \varphi_{225}$}
\AXC{$\varphi_{50}$\lbl{[D8]}}
\AXC{$\varphi_{50} \to \varphi_{51}$\lbl{A2}}
\BIC{$\varphi_{51}$}
\AXC{$\varphi_{51} \to d_{3} \in \mathbb{R} \to \varphi_{22}$\lbl{A2}}
\BIC{$d_{3} \in \mathbb{R} \to \varphi_{22}$}
\AXC{$(d_{3} \in \mathbb{R} \to \varphi_{22}) \to \varphi_{321}$\lbl{A1}}
\BIC{$\varphi_{321}$}
\BIC{$(\mathrm{FLim}(f,c,\alpha) \to \varphi_{12}) \to \varphi_{306}$}
\end{ptree}
\noindent$\Pi_{52}$:
\begin{ptree}
\AXC{$\stackrel{\Pi_{49}}{\vdots}$}\noLine
\UIC{$\mathrm{FLim}(f,c,\alpha) \to \varphi_{223}$}
\AXC{$\mathrm{FLim}(f,c,\alpha) \to \varphi_{12}$\lbl{A1}}
\AXC{$\stackrel{\Pi_{51}}{\vdots}$}\noLine
\UIC{$(\mathrm{FLim}(f,c,\alpha) \to \varphi_{12}) \to \varphi_{306}$}
\BIC{$\varphi_{306}$}
\BIC{$(\mathrm{FLim}(f,c,\alpha) \to \varphi_{16}) \to \mathrm{FLim}(f,c,\alpha) \to \varphi_{229}$}
\end{ptree}
\noindent$\Pi_{53}$:
\begin{ptree}
\AXC{$\mathrm{FLim}(f,c,\alpha) \to \varphi_{16}$\lbl{A1}}
\AXC{$\stackrel{\Pi_{52}}{\vdots}$}\noLine
\UIC{$(\mathrm{FLim}(f,c,\alpha) \to \varphi_{16}) \to \mathrm{FLim}(f,c,\alpha) \to \varphi_{229}$}
\BIC{$\mathrm{FLim}(f,c,\alpha) \to \varphi_{229}$}
\AXC{$(\mathrm{FLim}(f,c,\alpha) \to \varphi_{229}) \to \varphi_{343}$\lbl{A1}}
\BIC{$\varphi_{343}$}
\end{ptree}
\noindent$\Pi_{54}$:
\begin{ptree}
\AXC{$\varphi_{37}$\lbl{A2}}
\UIC{$\forall d\,(\varphi_{37})$}
\AXC{$\varphi_{291} \to \forall \delta\,(\varphi_{291})$\lbl{A4}}
\AXC{$\varphi_{313}$\lbl{A3}}
\AXC{$\stackrel{\Pi_{48}}{\vdots}$}\noLine
\UIC{$\forall \delta\,(\varphi_{291} \to \neg (\delta \in \mathbb{R} \land 0 < \delta \land \delta \leq d_{3} \land \delta \leq r))$}
\AXC{$\stackrel{\Pi_{53}}{\vdots}$}\noLine
\UIC{$\varphi_{343}$}
\BIC{$\varphi_{313} \to (\varphi_{291} \to \forall \delta\,(\varphi_{291})) \to \varphi_{303}$}
\BIC{$(\varphi_{291} \to \forall \delta\,(\varphi_{291})) \to \varphi_{303}$}
\BIC{$\varphi_{303}$}
\BIC{$(\forall d\,(\varphi_{37}) \to \varphi_{46}) \to \varphi_{54} \to \mathrm{FLim}(f,c,\alpha) \to \varphi_{282}$}
\end{ptree}
\noindent$\Pi_{55}$:
\begin{ptree}
\AXC{$\varphi_{54}$\lbl{A4}}
\AXC{$\forall d\,(\varphi_{37}) \to \varphi_{46}$\lbl{A3}}
\AXC{$\stackrel{\Pi_{54}}{\vdots}$}\noLine
\UIC{$(\forall d\,(\varphi_{37}) \to \varphi_{46}) \to \varphi_{54} \to \mathrm{FLim}(f,c,\alpha) \to \varphi_{282}$}
\BIC{$\varphi_{54} \to \mathrm{FLim}(f,c,\alpha) \to \varphi_{282}$}
\BIC{$\mathrm{FLim}(f,c,\alpha) \to \varphi_{282}$}
\AXC{$\varphi_{310}$\lbl{A1}}
\BIC{$\varphi_{261} \to \neg (d_{3} \in \mathbb{R} \land 0 < d_{3} \land d_{3} \leq d_{1} \land d_{3} \leq d_{2})$}
\UIC{$\varphi_{292}$}
\end{ptree}
\noindent$\Pi_{56}$:
\begin{ptree}
\AXC{$\mathrm{FLim}(f,c,\alpha) \to \varphi_{258}$\lbl{A1}}
\AXC{$(\mathrm{FLim}(f,c,\alpha) \to \varphi_{258}) \to \mathrm{FLim}(f,c,\alpha) \to \varphi_{210}$\lbl{A1}}
\BIC{$\mathrm{FLim}(f,c,\alpha) \to \varphi_{210}$}
\end{ptree}
\noindent$\Pi_{57}$:
\begin{ptree}
\AXC{$\varphi_{50}$\lbl{[D8]}}
\AXC{$\varphi_{50} \to \varphi_{52}$\lbl{A2}}
\BIC{$\varphi_{52}$}
\AXC{$\varphi_{52} \to d_{1} \in \mathbb{R} \to d_{2} \in \mathbb{R} \to \varphi_{9}$\lbl{A2}}
\BIC{$d_{1} \in \mathbb{R} \to d_{2} \in \mathbb{R} \to \varphi_{9}$}
\AXC{$(d_{1} \in \mathbb{R} \to d_{2} \in \mathbb{R} \to \varphi_{9}) \to \varphi_{334}$\lbl{A1}}
\BIC{$\varphi_{334}$}
\end{ptree}
\noindent$\Pi_{58}$:
\begin{ptree}
\AXC{$\mathrm{FLim}(f,c,\alpha) \to \varphi_{206}$\lbl{A1}}
\AXC{$(\mathrm{FLim}(f,c,\alpha) \to \varphi_{206}) \to \mathrm{FLim}(f,c,\alpha) \to \varphi_{213}$\lbl{A1}}
\BIC{$\mathrm{FLim}(f,c,\alpha) \to \varphi_{213}$}
\AXC{$\stackrel{\Pi_{57}}{\vdots}$}\noLine
\UIC{$\varphi_{334}$}
\BIC{$(\mathrm{FLim}(f,c,\alpha) \to \varphi_{212}) \to \varphi_{317}$}
\end{ptree}
\noindent$\Pi_{59}$:
\begin{ptree}
\AXC{$\mathrm{FLim}(f,c,\alpha) \to \varphi_{258}$\lbl{A1}}
\AXC{$(\mathrm{FLim}(f,c,\alpha) \to \varphi_{258}) \to \mathrm{FLim}(f,c,\alpha) \to \varphi_{212}$\lbl{A1}}
\BIC{$\mathrm{FLim}(f,c,\alpha) \to \varphi_{212}$}
\AXC{$\stackrel{\Pi_{58}}{\vdots}$}\noLine
\UIC{$(\mathrm{FLim}(f,c,\alpha) \to \varphi_{212}) \to \varphi_{317}$}
\BIC{$\varphi_{317}$}
\end{ptree}
\noindent$\Pi_{60}$:
\begin{ptree}
\AXC{$\mathrm{FLim}(f,c,\alpha) \to \varphi_{206}$\lbl{A1}}
\AXC{$(\mathrm{FLim}(f,c,\alpha) \to \varphi_{206}) \to \mathrm{FLim}(f,c,\alpha) \to \varphi_{211}$\lbl{A1}}
\BIC{$\mathrm{FLim}(f,c,\alpha) \to \varphi_{211}$}
\AXC{$\stackrel{\Pi_{59}}{\vdots}$}\noLine
\UIC{$\varphi_{317}$}
\BIC{$(\mathrm{FLim}(f,c,\alpha) \to \varphi_{210}) \to \mathrm{FLim}(f,c,\alpha) \to \varphi_{221}$}
\end{ptree}
\noindent$\Pi_{61}$:
\begin{ptree}
\AXC{$\stackrel{\Pi_{56}}{\vdots}$}\noLine
\UIC{$\mathrm{FLim}(f,c,\alpha) \to \varphi_{210}$}
\AXC{$\stackrel{\Pi_{60}}{\vdots}$}\noLine
\UIC{$(\mathrm{FLim}(f,c,\alpha) \to \varphi_{210}) \to \mathrm{FLim}(f,c,\alpha) \to \varphi_{221}$}
\BIC{$\mathrm{FLim}(f,c,\alpha) \to \varphi_{221}$}
\end{ptree}
\noindent$\Pi_{62}$:
\begin{ptree}
\AXC{$\varphi_{292} \to \varphi_{293}$\lbl{A3}}
\AXC{$\stackrel{\Pi_{55}}{\vdots}$}\noLine
\UIC{$\varphi_{292}$}
\AXC{$\stackrel{\Pi_{61}}{\vdots}$}\noLine
\UIC{$\mathrm{FLim}(f,c,\alpha) \to \varphi_{221}$}
\AXC{$(\mathrm{FLim}(f,c,\alpha) \to \varphi_{221}) \to \varphi_{292} \to (\varphi_{292} \to \varphi_{293}) \to \varphi_{333}$\lbl{A1}}
\BIC{$\varphi_{292} \to (\varphi_{292} \to \varphi_{293}) \to \varphi_{333}$}
\BIC{$(\varphi_{292} \to \varphi_{293}) \to \varphi_{333}$}
\BIC{$\varphi_{333}$}
\end{ptree}
\noindent$\Pi_{63}$:
\begin{ptree}
\AXC{$\varphi_{112}$\lbl{[D1]}}
\AXC{$\varphi_{170}$\lbl{A2}}
\BIC{$\forall c\,\forall \alpha\,\forall \varepsilon\,(\mathrm{FLim}(g,c,\alpha) \to \varepsilon \in \mathbb{R} \to 0 < \varepsilon \to \varphi_{75})$}
\AXC{$\varphi_{167}$\lbl{A2}}
\BIC{$\forall \alpha\,\forall \varepsilon\,(\mathrm{FLim}(g,c,\alpha) \to \varepsilon \in \mathbb{R} \to 0 < \varepsilon \to \varphi_{75})$}
\AXC{$\varphi_{161}$\lbl{A2}}
\BIC{$\forall \varepsilon\,(\mathrm{FLim}(g,c,\alpha) \to \varepsilon \in \mathbb{R} \to 0 < \varepsilon \to \varphi_{75})$}
\AXC{$\varphi_{159}$\lbl{A2}}
\BIC{$\mathrm{FLim}(g,c,\alpha) \to \varepsilon \in \mathbb{R} \to 0 < \varepsilon \to \varphi_{75}$}
\end{ptree}
\noindent$\Pi_{64}$:
\begin{ptree}
\AXC{$\varphi_{1}$\lbl{A1}}
\AXC{$\stackrel{\Pi_{63}}{\vdots}$}\noLine
\UIC{$\mathrm{FLim}(g,c,\alpha) \to \varepsilon \in \mathbb{R} \to 0 < \varepsilon \to \varphi_{75}$}
\AXC{$(\mathrm{FLim}(g,c,\alpha) \to \varepsilon \in \mathbb{R} \to 0 < \varepsilon \to \varphi_{75}) \to \varphi_{1} \to \varphi_{302}$\lbl{A1}}
\BIC{$\varphi_{1} \to \varphi_{302}$}
\BIC{$\varphi_{302}$}
\end{ptree}
\noindent$\Pi_{65}$:
\begin{ptree}
\AXC{$\mathrm{FLim}(f,c,\alpha) \to \varphi_{127}$\lbl{A1}}
\AXC{$\mathrm{FLim}(f,c,\alpha) \to \varphi_{124}$\lbl{A1}}
\AXC{$\stackrel{\Pi_{64}}{\vdots}$}\noLine
\UIC{$\varphi_{302}$}
\BIC{$(\mathrm{FLim}(f,c,\alpha) \to \varphi_{127}) \to \mathrm{FLim}(f,c,\alpha) \to \varphi_{207}$}
\BIC{$\mathrm{FLim}(f,c,\alpha) \to \varphi_{207}$}
\AXC{$(\mathrm{FLim}(f,c,\alpha) \to \varphi_{207}) \to \varphi_{342}$\lbl{A1}}
\BIC{$\varphi_{342}$}
\end{ptree}
\noindent$\Pi_{66}$:
\begin{ptree}
\AXC{$\varphi_{308}$\lbl{A3}}
\AXC{$\varphi_{59}$\lbl{A4}}
\AXC{$\varphi_{113}$\lbl{A3}}
\AXC{$\varphi_{39}$\lbl{A2}}
\UIC{$\forall d\,(\varphi_{39})$}
\AXC{$\varphi_{261} \to \forall d_{3}\,(\varphi_{261})$\lbl{A4}}
\AXC{$\stackrel{\Pi_{62}}{\vdots}$}\noLine
\UIC{$\varphi_{333}$}
\BIC{$\forall d\,(\varphi_{39}) \to \varphi_{113} \to \varphi_{59} \to \mathrm{FLim}(f,c,\alpha) \to \varphi_{259}$}
\BIC{$\varphi_{113} \to \varphi_{59} \to \mathrm{FLim}(f,c,\alpha) \to \varphi_{259}$}
\BIC{$\varphi_{59} \to \mathrm{FLim}(f,c,\alpha) \to \varphi_{259}$}
\BIC{$\mathrm{FLim}(f,c,\alpha) \to \varphi_{259}$}
\AXC{$\varphi_{307}$\lbl{A1}}
\BIC{$\varphi_{216} \to \neg (d_{2} \in \mathbb{R} \land 0 < d_{2} \land \varphi_{48})$}
\UIC{$\forall d_{2}\,(\varphi_{216} \to \neg (d_{2} \in \mathbb{R} \land 0 < d_{2} \land \varphi_{48}))$}
\AXC{$\stackrel{\Pi_{65}}{\vdots}$}\noLine
\UIC{$\varphi_{342}$}
\BIC{$\varphi_{308} \to (\varphi_{216} \to \forall d_{2}\,(\varphi_{216})) \to \varphi_{314}$}
\BIC{$(\varphi_{216} \to \forall d_{2}\,(\varphi_{216})) \to \varphi_{314}$}
\end{ptree}
\noindent$\Pi_{67}$:
\begin{ptree}
\AXC{$\varphi_{132}$\lbl{A4}}
\AXC{$\varphi_{220}$\lbl{A3}}
\AXC{$\forall d_{2}\,(\neg (d_{2} \in \mathbb{R} \land 0 < d_{2} \land \varphi_{48})) \to \varphi_{74}$\lbl{A2}}
\UIC{$\forall \delta\,(\forall d_{2}\,(\neg (d_{2} \in \mathbb{R} \land 0 < d_{2} \land \varphi_{48})) \to \varphi_{74})$}
\AXC{$\varphi_{216} \to \forall d_{2}\,(\varphi_{216})$\lbl{A4}}
\AXC{$\stackrel{\Pi_{66}}{\vdots}$}\noLine
\UIC{$(\varphi_{216} \to \forall d_{2}\,(\varphi_{216})) \to \varphi_{314}$}
\BIC{$\varphi_{314}$}
\BIC{$\varphi_{220} \to \varphi_{132} \to \mathrm{FLim}(f,c,\alpha) \to \varphi_{208}$}
\BIC{$\varphi_{132} \to \mathrm{FLim}(f,c,\alpha) \to \varphi_{208}$}
\BIC{$\mathrm{FLim}(f,c,\alpha) \to \varphi_{208}$}
\end{ptree}
\noindent$\Pi_{68}$:
\begin{ptree}
\AXC{$\varphi_{112}$\lbl{[D1]}}
\AXC{$\varphi_{171}$\lbl{A2}}
\BIC{$\forall c\,\forall \alpha\,\forall \varepsilon\,(\mathrm{FLim}(f,c,\alpha) \to \varepsilon \in \mathbb{R} \to 0 < \varepsilon \to \varphi_{76})$}
\AXC{$\varphi_{168}$\lbl{A2}}
\BIC{$\forall \alpha\,\forall \varepsilon\,(\mathrm{FLim}(f,c,\alpha) \to \varepsilon \in \mathbb{R} \to 0 < \varepsilon \to \varphi_{76})$}
\AXC{$\varphi_{162}$\lbl{A2}}
\BIC{$\forall \varepsilon\,(\mathrm{FLim}(f,c,\alpha) \to \varepsilon \in \mathbb{R} \to 0 < \varepsilon \to \varphi_{76})$}
\AXC{$\varphi_{160}$\lbl{A2}}
\BIC{$\mathrm{FLim}(f,c,\alpha) \to \varepsilon \in \mathbb{R} \to 0 < \varepsilon \to \varphi_{76}$}
\AXC{$\varphi_{304}$\lbl{A1}}
\BIC{$(\mathrm{FLim}(f,c,\alpha) \to \mathrm{FLim}(f,c,\alpha)) \to \varphi_{300}$}
\end{ptree}
\noindent$\Pi_{69}$:
\begin{ptree}
\AXC{$\mathrm{FLim}(f,c,\alpha) \to \mathrm{FLim}(f,c,\alpha)$\lbl{A1}}
\AXC{$\stackrel{\Pi_{68}}{\vdots}$}\noLine
\UIC{$(\mathrm{FLim}(f,c,\alpha) \to \mathrm{FLim}(f,c,\alpha)) \to \varphi_{300}$}
\BIC{$\varphi_{300}$}
\end{ptree}
\noindent$\Pi_{70}$:
\begin{ptree}
\AXC{$\mathrm{FLim}(f,c,\alpha) \to \varphi_{127}$\lbl{A1}}
\AXC{$\mathrm{FLim}(f,c,\alpha) \to \varphi_{124}$\lbl{A1}}
\AXC{$\stackrel{\Pi_{69}}{\vdots}$}\noLine
\UIC{$\varphi_{300}$}
\BIC{$(\mathrm{FLim}(f,c,\alpha) \to \varphi_{127}) \to \mathrm{FLim}(f,c,\alpha) \to \varphi_{163}$}
\BIC{$\mathrm{FLim}(f,c,\alpha) \to \varphi_{163}$}
\end{ptree}
\noindent$\Pi_{71}$:
\begin{ptree}
\AXC{$\stackrel{\Pi_{70}}{\vdots}$}\noLine
\UIC{$\mathrm{FLim}(f,c,\alpha) \to \varphi_{163}$}
\AXC{$(\mathrm{FLim}(f,c,\alpha) \to \varphi_{163}) \to \forall d_{1}\,(\varphi_{217}) \to \varphi_{336}$\lbl{A1}}
\BIC{$\forall d_{1}\,(\varphi_{217}) \to \varphi_{336}$}
\end{ptree}
\noindent$\Pi_{72}$:
\begin{ptree}
\AXC{$\stackrel{\Pi_{67}}{\vdots}$}\noLine
\UIC{$\mathrm{FLim}(f,c,\alpha) \to \varphi_{208}$}
\AXC{$(\mathrm{FLim}(f,c,\alpha) \to \varphi_{208}) \to \varphi_{217}$\lbl{A1}}
\BIC{$\varphi_{217}$}
\UIC{$\forall d_{1}\,(\varphi_{217})$}
\AXC{$\stackrel{\Pi_{71}}{\vdots}$}\noLine
\UIC{$\forall d_{1}\,(\varphi_{217}) \to \varphi_{336}$}
\BIC{$\varphi_{336}$}
\end{ptree}
\noindent$\Pi_{73}$:
\begin{ptree}
\AXC{$\varphi_{121} \land \varepsilon \in \mathbb{R} \land 0 < \varepsilon \land \neg \exists \delta\,(\delta \in \mathbb{R} \land 0 < \delta \land \varphi_{47}) \to \varphi_{172}$\lbl{A4}}
\AXC{$\forall d_{1}\,(\varphi_{217}) \to \varphi_{172} \to \forall d_{1}\,(\neg (d_{1} \in \mathbb{R} \land 0 < d_{1} \land \varphi_{49}))$\lbl{A3}}
\AXC{$\stackrel{\Pi_{72}}{\vdots}$}\noLine
\UIC{$\varphi_{336}$}
\BIC{$\varphi_{332}$}
\BIC{$\varphi_{311}$}
\end{ptree}
\noindent$\Pi_{74}$:
\begin{ptree}
\AXC{$\varphi_{299}$\lbl{A3}}
\AXC{$\varphi_{131}$\lbl{A4}}
\AXC{$\varphi_{219}$\lbl{A3}}
\AXC{$\forall d_{1}\,(\neg (d_{1} \in \mathbb{R} \land 0 < d_{1} \land \varphi_{49})) \to \varphi_{73}$\lbl{A2}}
\UIC{$\forall \delta\,(\forall d_{1}\,(\neg (d_{1} \in \mathbb{R} \land 0 < d_{1} \land \varphi_{49})) \to \varphi_{73})$}
\AXC{$\stackrel{\Pi_{73}}{\vdots}$}\noLine
\UIC{$\varphi_{311}$}
\BIC{$\varphi_{219} \to \varphi_{131} \to \mathrm{FLim}(f,c,\alpha) \to \varphi_{164}$}
\BIC{$\varphi_{131} \to \mathrm{FLim}(f,c,\alpha) \to \varphi_{164}$}
\BIC{$\mathrm{FLim}(f,c,\alpha) \to \varphi_{164}$}
\AXC{$\varphi_{298}$\lbl{A1}}
\BIC{$\varphi_{121} \to \varepsilon \in \mathbb{R} \to 0 < \varepsilon \to \exists \delta\,(\delta \in \mathbb{R} \land 0 < \delta \land \varphi_{47})$}
\UIC{$\forall \varepsilon\,(\varphi_{121} \to \varepsilon \in \mathbb{R} \to 0 < \varepsilon \to \exists \delta\,(\delta \in \mathbb{R} \land 0 < \delta \land \varphi_{47}))$}
\AXC{$\varphi_{328}$\lbl{A1}}
\BIC{$\varphi_{299} \to (\varphi_{121} \to \forall \varepsilon\,(\varphi_{121})) \to \mathrm{FLim}(f,c,\alpha) \to \varphi_{166}$}
\BIC{$(\varphi_{121} \to \forall \varepsilon\,(\varphi_{121})) \to \mathrm{FLim}(f,c,\alpha) \to \varphi_{166}$}
\end{ptree}
\noindent$\Pi_{75}$:
\begin{ptree}
\AXC{$\forall f\,\forall c\,\forall \alpha\,(\mathrm{FLim}(f,c,\alpha) \to c \in \mathbb{R})$\lbl{[D3]}}
\AXC{$\varphi_{23}$\lbl{A2}}
\BIC{$\forall c\,\forall \alpha\,(\mathrm{FLim}(f,c,\alpha) \to c \in \mathbb{R})$}
\AXC{$\varphi_{15}$\lbl{A2}}
\BIC{$\forall \alpha\,(\mathrm{FLim}(f,c,\alpha) \to c \in \mathbb{R})$}
\AXC{$\varphi_{10}$\lbl{A2}}
\BIC{$\mathrm{FLim}(f,c,\alpha) \to c \in \mathbb{R}$}
\end{ptree}
\noindent$\Pi_{76}$:
\begin{ptree}
\AXC{$\mathrm{FLim}(f,c,\alpha) \to \mathrm{FLim}(f,c,\alpha)$\lbl{A1}}
\AXC{$\stackrel{\Pi_{75}}{\vdots}$}\noLine
\UIC{$\mathrm{FLim}(f,c,\alpha) \to c \in \mathbb{R}$}
\AXC{$(\mathrm{FLim}(f,c,\alpha) \to c \in \mathbb{R}) \to \varphi_{140}$\lbl{A1}}
\BIC{$\varphi_{140}$}
\BIC{$\mathrm{FLim}(f,c,\alpha) \to \varphi_{118}$}
\end{ptree}
\noindent$\Pi_{77}$:
\begin{ptree}
\AXC{$\varphi_{2}$\lbl{A1}}
\AXC{$\varphi_{125}$\lbl{DefFLim}}
\AXC{$\varphi_{125} \to \forall c\,\forall \alpha\,(\varphi_{123})$\lbl{A2}}
\BIC{$\forall c\,\forall \alpha\,(\varphi_{123})$}
\AXC{$\forall c\,\forall \alpha\,(\varphi_{123}) \to \forall \alpha\,(\varphi_{123})$\lbl{A2}}
\BIC{$\forall \alpha\,(\varphi_{123})$}
\AXC{$\forall \alpha\,(\varphi_{123}) \to \varphi_{123}$\lbl{A2}}
\BIC{$\varphi_{123}$}
\AXC{$\varphi_{123} \to \varphi_{2} \to \varphi_{305}$\lbl{A1}}
\BIC{$\varphi_{2} \to \varphi_{305}$}
\BIC{$\varphi_{305}$}
\end{ptree}
\noindent$\Pi_{78}$:
\begin{ptree}
\AXC{$\mathrm{FLim}(f,c,\alpha) \to \varphi_{4}$\lbl{A1}}
\AXC{$\stackrel{\Pi_{26}}{\vdots}$}\noLine
\UIC{$\mathrm{FLim}(f,c,\alpha) \to \varphi_{119}$}
\AXC{$\stackrel{\Pi_{76}}{\vdots}$}\noLine
\UIC{$\mathrm{FLim}(f,c,\alpha) \to \varphi_{118}$}
\AXC{$\stackrel{\Pi_{77}}{\vdots}$}\noLine
\UIC{$\varphi_{305}$}
\BIC{$(\mathrm{FLim}(f,c,\alpha) \to \varphi_{119}) \to \varphi_{296}$}
\BIC{$\varphi_{296}$}
\BIC{$(\mathrm{FLim}(f,c,\alpha) \to \varphi_{166}) \to \mathrm{FLim}(f,c,\alpha) \to \varphi_{122}$}
\end{ptree}
\noindent 主木:
\begin{ptree}
\AXC{$\varphi_{121} \to \forall \varepsilon\,(\varphi_{121})$\lbl{A4}}
\AXC{$\stackrel{\Pi_{74}}{\vdots}$}\noLine
\UIC{$(\varphi_{121} \to \forall \varepsilon\,(\varphi_{121})) \to \mathrm{FLim}(f,c,\alpha) \to \varphi_{166}$}
\BIC{$\mathrm{FLim}(f,c,\alpha) \to \varphi_{166}$}
\AXC{$\stackrel{\Pi_{78}}{\vdots}$}\noLine
\UIC{$(\mathrm{FLim}(f,c,\alpha) \to \varphi_{166}) \to \mathrm{FLim}(f,c,\alpha) \to \varphi_{122}$}
\BIC{$\mathrm{FLim}(f,c,\alpha) \to \varphi_{122}$}
\UIC{$\forall r\,(\mathrm{FLim}(f,c,\alpha) \to \varphi_{122})$}
\UIC{$\forall \alpha\,\forall r\,(\mathrm{FLim}(f,c,\alpha) \to \varphi_{122})$}
\UIC{$\forall c\,\forall \alpha\,\forall r\,(\mathrm{FLim}(f,c,\alpha) \to \varphi_{122})$}
\UIC{$\forall h\,\forall c\,\forall \alpha\,\forall r\,(\mathrm{FLim}(f,c,\alpha) \to \varphi_{122})$}
\UIC{$\forall g\,\forall h\,\forall c\,\forall \alpha\,\forall r\,(\mathrm{FLim}(f,c,\alpha) \to \varphi_{122})$}
\UIC{$\forall f\,\forall g\,\forall h\,\forall c\,\forall \alpha\,\forall r\,(\mathrm{FLim}(f,c,\alpha) \to \varphi_{122})$}
\end{ptree}
\begin{abbrevlist}
\item $\varphi_{1} :\equiv \mathrm{FLim}(f,c,\alpha) \to \mathrm{FLim}(g,c,\alpha) \to \mathrm{FLim}(g,c,\alpha)$
\item $\varphi_{2} :\equiv \mathrm{FLim}(f,c,\alpha) \to \mathrm{FLim}(g,c,\alpha) \to \mathrm{RFn}(h) \to \mathrm{RFn}(h)$
\item $\varphi_{3} :\equiv \forall d\,\forall d_{3}\,(\neg (d_{3} \in \mathbb{R} \land 0 < d_{3} \land d_{3} \leq d_{1} \land d_{3} \leq d_{2}))$
\item $\varphi_{4} :\equiv \mathrm{FLim}(g,c,\alpha) \to \mathrm{RFn}(h) \to \mathrm{PNbhd}(h,c) \to \mathrm{PNbhd}(h,c)$
\item $\varphi_{5} :\equiv \forall \alpha\,(\mathrm{FLim}(g,c,\alpha) \to \mathrm{RFn}(g)) \to \mathrm{FLim}(g,c,\alpha) \to \mathrm{RFn}(g)$
\item $\varphi_{6} :\equiv \forall \alpha\,(\mathrm{FLim}(f,c,\alpha) \to \mathrm{RFn}(f)) \to \mathrm{FLim}(f,c,\alpha) \to \mathrm{RFn}(f)$
\item $\varphi_{7} :\equiv \forall c\,\forall \alpha\,(\mathrm{FLim}(g,c,\alpha) \to \mathrm{RFn}(g)) \to \forall \alpha\,(\mathrm{FLim}(g,c,\alpha) \to \mathrm{RFn}(g))$
\item $\varphi_{8} :\equiv \forall c\,\forall \alpha\,(\mathrm{FLim}(f,c,\alpha) \to \mathrm{RFn}(f)) \to \forall \alpha\,(\mathrm{FLim}(f,c,\alpha) \to \mathrm{RFn}(f))$
\item $\varphi_{9} :\equiv 0 < d_{1} \to 0 < d_{2} \to \exists d\,(d \in \mathbb{R} \land 0 < d \land d \leq d_{1} \land d \leq d_{2})$
\item $\varphi_{10} :\equiv \forall \alpha\,(\mathrm{FLim}(f,c,\alpha) \to c \in \mathbb{R}) \to \mathrm{FLim}(f,c,\alpha) \to c \in \mathbb{R}$
\item $\varphi_{11} :\equiv \forall \alpha\,(\mathrm{FLim}(f,c,\alpha) \to \alpha \in \mathbb{R}) \to \mathrm{FLim}(f,c,\alpha) \to \alpha \in \mathbb{R}$
\item $\varphi_{12} :\equiv \mathrm{FLim}(g,c,\alpha) \to \mathrm{RFn}(h) \to \mathrm{PNbhd}(h,c) \to r \in \mathbb{R} \to r \in \mathbb{R}$
\item $\varphi_{13} :\equiv \forall f\,\forall c\,\forall \alpha\,(\mathrm{FLim}(f,c,\alpha) \to \mathrm{RFn}(f)) \to \forall c\,\forall \alpha\,(\mathrm{FLim}(g,c,\alpha) \to \mathrm{RFn}(g))$
\item $\varphi_{14} :\equiv \forall f\,\forall c\,\forall \alpha\,(\mathrm{FLim}(f,c,\alpha) \to \mathrm{RFn}(f)) \to \forall c\,\forall \alpha\,(\mathrm{FLim}(f,c,\alpha) \to \mathrm{RFn}(f))$
\item $\varphi_{15} :\equiv \forall c\,\forall \alpha\,(\mathrm{FLim}(f,c,\alpha) \to c \in \mathbb{R}) \to \forall \alpha\,(\mathrm{FLim}(f,c,\alpha) \to c \in \mathbb{R})$
\item $\varphi_{16} :\equiv \mathrm{FLim}(g,c,\alpha) \to \mathrm{RFn}(h) \to \mathrm{PNbhd}(h,c) \to r \in \mathbb{R} \to 0 < r \to 0 < r$
\item $\varphi_{17} :\equiv \forall c\,\forall \alpha\,(\mathrm{FLim}(f,c,\alpha) \to \alpha \in \mathbb{R}) \to \forall \alpha\,(\mathrm{FLim}(f,c,\alpha) \to \alpha \in \mathbb{R})$
\item $\varphi_{18} :\equiv \delta \in \mathbb{R} \land 0 < \delta \land \delta \leq d_{3} \land \delta \leq r \to x \in \mathrm{dom}(h) \to x \in \mathrm{dom}(h)$
\item $\varphi_{19} :\equiv x \in \mathrm{dom}(h) \to 0 < \lvert x - c \rvert \to \lvert x - c \rvert < \delta \to x \in \mathrm{dom}(g)$
\item $\varphi_{20} :\equiv x \in \mathrm{dom}(h) \to 0 < \lvert x - c \rvert \to \lvert x - c \rvert < \delta \to x \in \mathrm{dom}(f)$
\item $\varphi_{21} :\equiv x \in \mathrm{dom}(f) \land x \in \mathrm{dom}(g) \land f(x) \leq h(x) \land h(x) \leq g(x)$
\item $\varphi_{22} :\equiv r \in \mathbb{R} \to 0 < d_{3} \to 0 < r \to \exists d\,(d \in \mathbb{R} \land 0 < d \land d \leq d_{3} \land d \leq r)$
\item $\varphi_{23} :\equiv \forall f\,\forall c\,\forall \alpha\,(\mathrm{FLim}(f,c,\alpha) \to c \in \mathbb{R}) \to \forall c\,\forall \alpha\,(\mathrm{FLim}(f,c,\alpha) \to c \in \mathbb{R})$
\item $\varphi_{24} :\equiv x \in \mathrm{dom}(h) \to 0 < \lvert x - c \rvert \to \lvert x - c \rvert < \delta \to \lvert x - c \rvert < r$
\item $\varphi_{25} :\equiv x \in \mathrm{dom}(h) \to 0 < \lvert x - c \rvert \to \lvert x - c \rvert < \delta \to \lvert x - c \rvert < \delta$
\item $\varphi_{26} :\equiv x \in \mathrm{dom}(h) \to 0 < \lvert x - c \rvert \to \lvert x - c \rvert < \delta \to \lvert x - c \rvert < d_{3}$
\item $\varphi_{27} :\equiv x \in \mathrm{dom}(h) \to 0 < \lvert x - c \rvert \to \lvert x - c \rvert < \delta \to \lvert x - c \rvert < d_{2}$
\item $\varphi_{28} :\equiv x \in \mathrm{dom}(h) \to 0 < \lvert x - c \rvert \to \lvert x - c \rvert < \delta \to \lvert x - c \rvert < d_{1}$
\item $\varphi_{29} :\equiv x \in \mathrm{dom}(h) \to 0 < \lvert x - c \rvert \to \lvert x - c \rvert < \delta \to h(x) \leq g(x)$
\item $\varphi_{30} :\equiv x \in \mathrm{dom}(h) \to 0 < \lvert x - c \rvert \to \lvert x - c \rvert < \delta \to f(x) \leq h(x)$
\item $\varphi_{31} :\equiv x \in \mathrm{dom}(h) \to 0 < \lvert x - c \rvert \to \lvert x - c \rvert < \delta \to h(x) \in \mathbb{R}$
\item $\varphi_{32} :\equiv x \in \mathrm{dom}(h) \to 0 < \lvert x - c \rvert \to \lvert x - c \rvert < \delta \to g(x) \in \mathbb{R}$
\item $\varphi_{33} :\equiv x \in \mathrm{dom}(h) \to 0 < \lvert x - c \rvert \to \lvert x - c \rvert < \delta \to f(x) \in \mathbb{R}$
\item $\varphi_{34} :\equiv \forall f\,\forall c\,\forall \alpha\,(\mathrm{FLim}(f,c,\alpha) \to \alpha \in \mathbb{R}) \to \forall c\,\forall \alpha\,(\mathrm{FLim}(f,c,\alpha) \to \alpha \in \mathbb{R})$
\item $\varphi_{35} :\equiv \delta \in \mathbb{R} \land 0 < \delta \land \delta \leq d_{3} \land \delta \leq r \to \delta \in \mathbb{R} \land 0 < \delta \land \delta \leq d_{3} \land \delta \leq r$
\item $\varphi_{36} :\equiv d_{3} \in \mathbb{R} \land 0 < d_{3} \land d_{3} \leq d_{1} \land d_{3} \leq d_{2} \to \exists d\,(d \in \mathbb{R} \land 0 < d \land d \leq d_{3} \land d \leq r)$
\item $\varphi_{37} :\equiv \forall \delta\,(\neg (\delta \in \mathbb{R} \land 0 < \delta \land \delta \leq d_{3} \land \delta \leq r)) \to \neg (d \in \mathbb{R} \land 0 < d \land d \leq d_{3} \land d \leq r)$
\item $\varphi_{38} :\equiv d_{3} \in \mathbb{R} \land 0 < d_{3} \land d_{3} \leq d_{1} \land d_{3} \leq d_{2} \to d_{3} \in \mathbb{R} \land 0 < d_{3} \land d_{3} \leq d_{1} \land d_{3} \leq d_{2}$
\item $\varphi_{39} :\equiv \forall d_{3}\,(\neg (d_{3} \in \mathbb{R} \land 0 < d_{3} \land d_{3} \leq d_{1} \land d_{3} \leq d_{2})) \to \neg (d \in \mathbb{R} \land 0 < d \land d \leq d_{1} \land d \leq d_{2})$
\item $\varphi_{40} :\equiv \forall x\,\forall y\,\forall z\,(x < y \to y \leq z \to x < z) \to \forall y\,\forall z\,(\lvert x - c \rvert < y \to y \leq z \to \lvert x - c \rvert < z)$
\item $\varphi_{41} :\equiv x \in \mathrm{dom}(g) \to 0 < \lvert x - c \rvert \to \lvert x - c \rvert < d_{2} \to \lvert g(x) - \alpha \rvert < \varepsilon$
\item $\varphi_{42} :\equiv x \in \mathrm{dom}(h) \to 0 < \lvert x - c \rvert \to \lvert x - c \rvert < \delta \to \lvert g(x) - \alpha \rvert < \varepsilon$
\item $\varphi_{43} :\equiv x \in \mathrm{dom}(f) \to 0 < \lvert x - c \rvert \to \lvert x - c \rvert < d_{1} \to \lvert f(x) - \alpha \rvert < \varepsilon$
\item $\varphi_{44} :\equiv x \in \mathrm{dom}(h) \to 0 < \lvert x - c \rvert \to \lvert x - c \rvert < \delta \to \lvert f(x) - \alpha \rvert < \varepsilon$
\item $\varphi_{45} :\equiv x \in \mathrm{dom}(h) \to 0 < \lvert x - c \rvert \to \lvert x - c \rvert < \delta \to \lvert h(x) - \alpha \rvert < \varepsilon$
\item $\varphi_{46} :\equiv \forall d\,\forall \delta\,(\neg (\delta \in \mathbb{R} \land 0 < \delta \land \delta \leq d_{3} \land \delta \leq r)) \to \forall d\,(\neg (d \in \mathbb{R} \land 0 < d \land d \leq d_{3} \land d \leq r))$
\item $\varphi_{47} :\equiv \forall x \in \mathrm{dom}(h)\,(0 < \lvert x - c \rvert \to \lvert x - c \rvert < \delta \to \lvert h(x) - \alpha \rvert < \varepsilon)$
\item $\varphi_{48} :\equiv \forall x \in \mathrm{dom}(g)\,(0 < \lvert x - c \rvert \to \lvert x - c \rvert < d_{2} \to \lvert g(x) - \alpha \rvert < \varepsilon)$
\item $\varphi_{49} :\equiv \forall x \in \mathrm{dom}(f)\,(0 < \lvert x - c \rvert \to \lvert x - c \rvert < d_{1} \to \lvert f(x) - \alpha \rvert < \varepsilon)$
\item $\varphi_{50} :\equiv \forall p\,\forall q\,(p \in \mathbb{R} \to q \in \mathbb{R} \to 0 < p \to 0 < q \to \exists d\,(d \in \mathbb{R} \land 0 < d \land d \leq p \land d \leq q))$
\item $\varphi_{51} :\equiv \forall q\,(d_{3} \in \mathbb{R} \to q \in \mathbb{R} \to 0 < d_{3} \to 0 < q \to \exists d\,(d \in \mathbb{R} \land 0 < d \land d \leq d_{3} \land d \leq q))$
\item $\varphi_{52} :\equiv \forall q\,(d_{1} \in \mathbb{R} \to q \in \mathbb{R} \to 0 < d_{1} \to 0 < q \to \exists d\,(d \in \mathbb{R} \land 0 < d \land d \leq d_{1} \land d \leq q))$
\item $\varphi_{53} :\equiv \lvert f(x) - \alpha \rvert < \varepsilon \to \lvert g(x) - \alpha \rvert < \varepsilon \to \lvert h(x) - \alpha \rvert < \varepsilon$
\item $\varphi_{54} :\equiv \forall \delta\,(\neg (\delta \in \mathbb{R} \land 0 < \delta \land \delta \leq d_{3} \land \delta \leq r)) \to \forall d\,\forall \delta\,(\neg (\delta \in \mathbb{R} \land 0 < \delta \land \delta \leq d_{3} \land \delta \leq r))$
\item $\varphi_{55} :\equiv \forall x\,(\mathrm{RFn}(h) \to x \in \mathrm{dom}(h) \to h(x) \in \mathbb{R}) \to \mathrm{RFn}(h) \to x \in \mathrm{dom}(h) \to h(x) \in \mathbb{R}$
\item $\varphi_{56} :\equiv \forall x\,(\mathrm{RFn}(g) \to x \in \mathrm{dom}(g) \to g(x) \in \mathbb{R}) \to \mathrm{RFn}(g) \to x \in \mathrm{dom}(g) \to g(x) \in \mathbb{R}$
\item $\varphi_{57} :\equiv \forall x\,(\mathrm{RFn}(f) \to x \in \mathrm{dom}(f) \to f(x) \in \mathbb{R}) \to \mathrm{RFn}(f) \to x \in \mathrm{dom}(f) \to f(x) \in \mathbb{R}$
\item $\varphi_{58} :\equiv \delta \in \mathbb{R} \land 0 < \delta \land \delta \leq d_{3} \land \delta \leq r \to x \in \mathrm{dom}(h) \to 0 < \lvert x - c \rvert \to 0 < \lvert x - c \rvert$
\item $\varphi_{59} :\equiv \forall d_{3}\,(\neg (d_{3} \in \mathbb{R} \land 0 < d_{3} \land d_{3} \leq d_{1} \land d_{3} \leq d_{2})) \to \varphi_{3}$
\item $\varphi_{60} :\equiv \forall z\,(\lvert x - c \rvert < \delta \to \delta \leq z \to \lvert x - c \rvert < z) \to \lvert x - c \rvert < \delta \to \delta \leq r \to \lvert x - c \rvert < r$
\item $\varphi_{61} :\equiv \forall z\,(\lvert x - c \rvert < \delta \to \delta \leq z \to \lvert x - c \rvert < z) \to \lvert x - c \rvert < \delta \to \delta \leq d_{3} \to \lvert x - c \rvert < d_{3}$
\item $\varphi_{62} :\equiv \forall z\,(\lvert x - c \rvert < d_{3} \to d_{3} \leq z \to \lvert x - c \rvert < z) \to \lvert x - c \rvert < d_{3} \to d_{3} \leq d_{2} \to \lvert x - c \rvert < d_{2}$
\item $\varphi_{63} :\equiv \forall z\,(\lvert x - c \rvert < d_{3} \to d_{3} \leq z \to \lvert x - c \rvert < z) \to \lvert x - c \rvert < d_{3} \to d_{3} \leq d_{1} \to \lvert x - c \rvert < d_{1}$
\item $\varphi_{64} :\equiv \forall f\,\forall x\,(\mathrm{RFn}(f) \to x \in \mathrm{dom}(f) \to f(x) \in \mathbb{R}) \to \forall x\,(\mathrm{RFn}(h) \to x \in \mathrm{dom}(h) \to h(x) \in \mathbb{R})$
\item $\varphi_{65} :\equiv \forall f\,\forall x\,(\mathrm{RFn}(f) \to x \in \mathrm{dom}(f) \to f(x) \in \mathbb{R}) \to \forall x\,(\mathrm{RFn}(g) \to x \in \mathrm{dom}(g) \to g(x) \in \mathbb{R})$
\item $\varphi_{66} :\equiv \forall f\,\forall x\,(\mathrm{RFn}(f) \to x \in \mathrm{dom}(f) \to f(x) \in \mathbb{R}) \to \forall x\,(\mathrm{RFn}(f) \to x \in \mathrm{dom}(f) \to f(x) \in \mathbb{R})$
\item $\varphi_{67} :\equiv \forall y\,\forall z\,(\lvert x - c \rvert < y \to y \leq z \to \lvert x - c \rvert < z) \to \forall z\,(\lvert x - c \rvert < \delta \to \delta \leq z \to \lvert x - c \rvert < z)$
\item $\varphi_{68} :\equiv \forall y\,\forall z\,(\lvert x - c \rvert < y \to y \leq z \to \lvert x - c \rvert < z) \to \forall z\,(\lvert x - c \rvert < d_{3} \to d_{3} \leq z \to \lvert x - c \rvert < z)$
\item $\varphi_{69} :\equiv \delta \in \mathbb{R} \land 0 < \delta \land \delta \leq d_{3} \land \delta \leq r \to x \in \mathrm{dom}(h) \to 0 < \lvert x - c \rvert \to \lvert x - c \rvert < \delta \to \delta \leq r$
\item $\varphi_{70} :\equiv \delta \in \mathbb{R} \land 0 < \delta \land \delta \leq d_{3} \land \delta \leq r \to x \in \mathrm{dom}(h) \to 0 < \lvert x - c \rvert \to \lvert x - c \rvert < \delta \to d_{3} \leq d_{2}$
\item $\varphi_{71} :\equiv \delta \in \mathbb{R} \land 0 < \delta \land \delta \leq d_{3} \land \delta \leq r \to x \in \mathrm{dom}(h) \to 0 < \lvert x - c \rvert \to \lvert x - c \rvert < \delta \to \delta \leq d_{3}$
\item $\varphi_{72} :\equiv \delta \in \mathbb{R} \land 0 < \delta \land \delta \leq d_{3} \land \delta \leq r \to x \in \mathrm{dom}(h) \to 0 < \lvert x - c \rvert \to \lvert x - c \rvert < \delta \to d_{3} \leq d_{1}$
\item $\varphi_{73} :\equiv \neg (\delta \in \mathbb{R} \land 0 < \delta \land \forall x \in \mathrm{dom}(f)\,(0 < \lvert x - c \rvert \to \lvert x - c \rvert < \delta \to \lvert f(x) - \alpha \rvert < \varepsilon))$
\item $\varphi_{74} :\equiv \neg (\delta \in \mathbb{R} \land 0 < \delta \land \forall x \in \mathrm{dom}(g)\,(0 < \lvert x - c \rvert \to \lvert x - c \rvert < \delta \to \lvert g(x) - \alpha \rvert < \varepsilon))$
\item $\varphi_{75} :\equiv \exists \delta\,(\delta \in \mathbb{R} \land 0 < \delta \land \forall x \in \mathrm{dom}(g)\,(0 < \lvert x - c \rvert \to \lvert x - c \rvert < \delta \to \lvert g(x) - \alpha \rvert < \varepsilon))$
\item $\varphi_{76} :\equiv \exists \delta\,(\delta \in \mathbb{R} \land 0 < \delta \land \forall x \in \mathrm{dom}(f)\,(0 < \lvert x - c \rvert \to \lvert x - c \rvert < \delta \to \lvert f(x) - \alpha \rvert < \varepsilon))$
\item $\varphi_{77} :\equiv \forall x \in \mathrm{dom}(h)\,(0 < \lvert x - c \rvert \to \lvert x - c \rvert < r \to \varphi_{21}) \to \mathrm{RFn}(g)$
\item $\varphi_{78} :\equiv \forall x \in \mathrm{dom}(h)\,(0 < \lvert x - c \rvert \to \lvert x - c \rvert < r \to \varphi_{21}) \to \mathrm{RFn}(f)$
\item $\varphi_{79} :\equiv \forall x \in \mathrm{dom}(h)\,(0 < \lvert x - c \rvert \to \lvert x - c \rvert < r \to \varphi_{21}) \to c \in \mathbb{R}$
\item $\varphi_{80} :\equiv \forall x \in \mathrm{dom}(h)\,(0 < \lvert x - c \rvert \to \lvert x - c \rvert < r \to \varphi_{21}) \to \alpha \in \mathbb{R}$
\item $\varphi_{81} :\equiv \forall x \in \mathrm{dom}(h)\,(0 < \lvert x - c \rvert \to \lvert x - c \rvert < r \to \varphi_{21}) \to \mathrm{FLim}(h,c,\alpha)$
\item $\varphi_{82} :\equiv d_{3} \in \mathbb{R} \land 0 < d_{3} \land d_{3} \leq d_{1} \land d_{3} \leq d_{2} \to \delta \in \mathbb{R} \land 0 < \delta \land \delta \leq d_{3} \land \delta \leq r \to \varphi_{19}$
\item $\varphi_{83} :\equiv d_{3} \in \mathbb{R} \land 0 < d_{3} \land d_{3} \leq d_{1} \land d_{3} \leq d_{2} \to \delta \in \mathbb{R} \land 0 < \delta \land \delta \leq d_{3} \land \delta \leq r \to \varphi_{20}$
\item $\varphi_{84} :\equiv d_{2} \in \mathbb{R} \land 0 < d_{2} \land \varphi_{48} \to \exists d\,(d \in \mathbb{R} \land 0 < d \land d \leq d_{1} \land d \leq d_{2})$
\item $\varphi_{85} :\equiv d_{3} \in \mathbb{R} \land 0 < d_{3} \land d_{3} \leq d_{1} \land d_{3} \leq d_{2} \to \delta \in \mathbb{R} \land 0 < \delta \land \delta \leq d_{3} \land \delta \leq r \to \varphi_{24}$
\item $\varphi_{86} :\equiv d_{3} \in \mathbb{R} \land 0 < d_{3} \land d_{3} \leq d_{1} \land d_{3} \leq d_{2} \to \delta \in \mathbb{R} \land 0 < \delta \land \delta \leq d_{3} \land \delta \leq r \to \varphi_{25}$
\item $\varphi_{87} :\equiv d_{3} \in \mathbb{R} \land 0 < d_{3} \land d_{3} \leq d_{1} \land d_{3} \leq d_{2} \to \delta \in \mathbb{R} \land 0 < \delta \land \delta \leq d_{3} \land \delta \leq r \to \varphi_{26}$
\item $\varphi_{88} :\equiv d_{3} \in \mathbb{R} \land 0 < d_{3} \land d_{3} \leq d_{1} \land d_{3} \leq d_{2} \to \delta \in \mathbb{R} \land 0 < \delta \land \delta \leq d_{3} \land \delta \leq r \to \varphi_{27}$
\item $\varphi_{89} :\equiv d_{3} \in \mathbb{R} \land 0 < d_{3} \land d_{3} \leq d_{1} \land d_{3} \leq d_{2} \to \delta \in \mathbb{R} \land 0 < \delta \land \delta \leq d_{3} \land \delta \leq r \to \varphi_{28}$
\item $\varphi_{90} :\equiv d_{3} \in \mathbb{R} \land 0 < d_{3} \land d_{3} \leq d_{1} \land d_{3} \leq d_{2} \to \delta \in \mathbb{R} \land 0 < \delta \land \delta \leq d_{3} \land \delta \leq r \to \varphi_{29}$
\item $\varphi_{91} :\equiv d_{3} \in \mathbb{R} \land 0 < d_{3} \land d_{3} \leq d_{1} \land d_{3} \leq d_{2} \to \delta \in \mathbb{R} \land 0 < \delta \land \delta \leq d_{3} \land \delta \leq r \to \varphi_{30}$
\item $\varphi_{92} :\equiv d_{3} \in \mathbb{R} \land 0 < d_{3} \land d_{3} \leq d_{1} \land d_{3} \leq d_{2} \to \delta \in \mathbb{R} \land 0 < \delta \land \delta \leq d_{3} \land \delta \leq r \to \varphi_{31}$
\item $\varphi_{93} :\equiv d_{3} \in \mathbb{R} \land 0 < d_{3} \land d_{3} \leq d_{1} \land d_{3} \leq d_{2} \to \delta \in \mathbb{R} \land 0 < \delta \land \delta \leq d_{3} \land \delta \leq r \to \varphi_{32}$
\item $\varphi_{94} :\equiv d_{3} \in \mathbb{R} \land 0 < d_{3} \land d_{3} \leq d_{1} \land d_{3} \leq d_{2} \to \delta \in \mathbb{R} \land 0 < \delta \land \delta \leq d_{3} \land \delta \leq r \to \varphi_{33}$
\item $\varphi_{95} :\equiv \delta \in \mathbb{R} \land 0 < \delta \land \delta \leq d_{3} \land \delta \leq r \to \exists \delta\,(\delta \in \mathbb{R} \land 0 < \delta \land \varphi_{47})$
\item $\varphi_{96} :\equiv d_{3} \in \mathbb{R} \land 0 < d_{3} \land d_{3} \leq d_{1} \land d_{3} \leq d_{2} \to \exists \delta\,(\delta \in \mathbb{R} \land 0 < \delta \land \varphi_{47})$
\item $\varphi_{97} :\equiv \forall x \in \mathrm{dom}(h)\,(0 < \lvert x - c \rvert \to \lvert x - c \rvert < r \to \varphi_{21}) \to \varepsilon \in \mathbb{R} \to \varepsilon \in \mathbb{R}$
\item $\varphi_{98} :\equiv \alpha \in \mathbb{R} \to \varepsilon \in \mathbb{R} \to f(x) \leq h(x) \to h(x) \leq g(x) \to \varphi_{53}$
\item $\varphi_{99} :\equiv d_{2} \in \mathbb{R} \land 0 < d_{2} \land \varphi_{48} \to d_{3} \in \mathbb{R} \land 0 < d_{3} \land d_{3} \leq d_{1} \land d_{3} \leq d_{2} \to 0 < d_{3}$
\item $\varphi_{100} :\equiv d_{3} \in \mathbb{R} \land 0 < d_{3} \land d_{3} \leq d_{1} \land d_{3} \leq d_{2} \to \delta \in \mathbb{R} \land 0 < \delta \land \delta \leq d_{3} \land \delta \leq r \to \varphi_{41}$
\item $\varphi_{101} :\equiv d_{3} \in \mathbb{R} \land 0 < d_{3} \land d_{3} \leq d_{1} \land d_{3} \leq d_{2} \to \delta \in \mathbb{R} \land 0 < \delta \land \delta \leq d_{3} \land \delta \leq r \to \varphi_{42}$
\item $\varphi_{102} :\equiv d_{3} \in \mathbb{R} \land 0 < d_{3} \land d_{3} \leq d_{1} \land d_{3} \leq d_{2} \to \delta \in \mathbb{R} \land 0 < \delta \land \delta \leq d_{3} \land \delta \leq r \to \varphi_{43}$
\item $\varphi_{103} :\equiv d_{3} \in \mathbb{R} \land 0 < d_{3} \land d_{3} \leq d_{1} \land d_{3} \leq d_{2} \to \delta \in \mathbb{R} \land 0 < \delta \land \delta \leq d_{3} \land \delta \leq r \to \varphi_{44}$
\item $\varphi_{104} :\equiv d_{3} \in \mathbb{R} \land 0 < d_{3} \land d_{3} \leq d_{1} \land d_{3} \leq d_{2} \to \delta \in \mathbb{R} \land 0 < \delta \land \delta \leq d_{3} \land \delta \leq r \to \varphi_{45}$
\item $\varphi_{105} :\equiv \delta \in \mathbb{R} \land 0 < \delta \land \delta \leq d_{3} \land \delta \leq r \to x \in \mathrm{dom}(h) \to 0 < \lvert x - c \rvert \to \lvert x - c \rvert < \delta \to \varphi_{21}$
\item $\varphi_{106} :\equiv \forall x \in \mathrm{dom}(h)\,(0 < \lvert x - c \rvert \to \lvert x - c \rvert < r \to \varphi_{21}) \to \varepsilon \in \mathbb{R} \to 0 < \varepsilon \to 0 < \varepsilon$
\item $\varphi_{107} :\equiv \forall x\,\forall y\,\forall z\,\forall \alpha\,\forall \varepsilon\,(x \in \mathbb{R} \to y \in \mathbb{R} \to z \in \mathbb{R} \to \alpha \in \mathbb{R} \to \varepsilon \in \mathbb{R} \to x \leq z \to z \leq y \to \lvert x - \alpha \rvert < \varepsilon \to \lvert y - \alpha \rvert < \varepsilon \to \lvert z - \alpha \rvert < \varepsilon)$
\item $\varphi_{108} :\equiv d_{2} \in \mathbb{R} \land 0 < d_{2} \land \varphi_{48} \to d_{3} \in \mathbb{R} \land 0 < d_{3} \land d_{3} \leq d_{1} \land d_{3} \leq d_{2} \to d_{3} \in \mathbb{R}$
\item $\varphi_{109} :\equiv d_{3} \in \mathbb{R} \land 0 < d_{3} \land d_{3} \leq d_{1} \land d_{3} \leq d_{2} \to \delta \in \mathbb{R} \land 0 < \delta \land \delta \leq d_{3} \land \delta \leq r \to \varphi_{48}$
\item $\varphi_{110} :\equiv d_{3} \in \mathbb{R} \land 0 < d_{3} \land d_{3} \leq d_{1} \land d_{3} \leq d_{2} \to \delta \in \mathbb{R} \land 0 < \delta \land \delta \leq d_{3} \land \delta \leq r \to \varphi_{49}$
\item $\varphi_{111} :\equiv d_{3} \in \mathbb{R} \land 0 < d_{3} \land d_{3} \leq d_{1} \land d_{3} \leq d_{2} \to \delta \in \mathbb{R} \land 0 < \delta \land \delta \leq d_{3} \land \delta \leq r \to \varphi_{47}$
\item $\varphi_{112} :\equiv \forall f\,\forall c\,\forall \alpha\,\forall \varepsilon\,(\mathrm{FLim}(f,c,\alpha) \to \varepsilon \in \mathbb{R} \to 0 < \varepsilon \to \varphi_{76})$
\item $\varphi_{113} :\equiv \forall d\,(\varphi_{39}) \to \varphi_{3} \to \forall d\,(\neg (d \in \mathbb{R} \land 0 < d \land d \leq d_{1} \land d \leq d_{2}))$
\item $\varphi_{114} :\equiv \forall y\,\forall z\,\forall \alpha\,\forall \varepsilon\,(f(x) \in \mathbb{R} \to y \in \mathbb{R} \to z \in \mathbb{R} \to \alpha \in \mathbb{R} \to \varepsilon \in \mathbb{R} \to f(x) \leq z \to z \leq y \to \lvert f(x) - \alpha \rvert < \varepsilon \to \lvert y - \alpha \rvert < \varepsilon \to \lvert z - \alpha \rvert < \varepsilon)$
\item $\varphi_{115} :\equiv d_{3} \in \mathbb{R} \land 0 < d_{3} \land d_{3} \leq d_{1} \land d_{3} \leq d_{2} \to \delta \in \mathbb{R} \land 0 < \delta \land \delta \leq d_{3} \land \delta \leq r \to \delta \in \mathbb{R} \land 0 < \delta \land \varphi_{47}$
\item $\varphi_{116} :\equiv \mathrm{FLim}(g,c,\alpha) \to \mathrm{RFn}(h) \to \mathrm{PNbhd}(h,c) \to r \in \mathbb{R} \to 0 < r \to \varphi_{77}$
\item $\varphi_{117} :\equiv \mathrm{FLim}(g,c,\alpha) \to \mathrm{RFn}(h) \to \mathrm{PNbhd}(h,c) \to r \in \mathbb{R} \to 0 < r \to \varphi_{78}$
\item $\varphi_{118} :\equiv \mathrm{FLim}(g,c,\alpha) \to \mathrm{RFn}(h) \to \mathrm{PNbhd}(h,c) \to r \in \mathbb{R} \to 0 < r \to \varphi_{79}$
\item $\varphi_{119} :\equiv \mathrm{FLim}(g,c,\alpha) \to \mathrm{RFn}(h) \to \mathrm{PNbhd}(h,c) \to r \in \mathbb{R} \to 0 < r \to \varphi_{80}$
\item $\varphi_{120} :\equiv \forall z\,\forall \alpha\,\forall \varepsilon\,(f(x) \in \mathbb{R} \to g(x) \in \mathbb{R} \to z \in \mathbb{R} \to \alpha \in \mathbb{R} \to \varepsilon \in \mathbb{R} \to f(x) \leq z \to z \leq g(x) \to \lvert f(x) - \alpha \rvert < \varepsilon \to \lvert g(x) - \alpha \rvert < \varepsilon \to \lvert z - \alpha \rvert < \varepsilon)$
\item $\varphi_{121} :\equiv \mathrm{FLim}(f,c,\alpha) \land \mathrm{FLim}(g,c,\alpha) \land \mathrm{RFn}(h) \land \mathrm{PNbhd}(h,c) \land r \in \mathbb{R} \land 0 < r \land \forall x \in \mathrm{dom}(h)\,(0 < \lvert x - c \rvert \to \lvert x - c \rvert < r \to \varphi_{21})$
\item $\varphi_{122} :\equiv \mathrm{FLim}(g,c,\alpha) \to \mathrm{RFn}(h) \to \mathrm{PNbhd}(h,c) \to r \in \mathbb{R} \to 0 < r \to \varphi_{81}$
\item $\varphi_{123} :\equiv \mathrm{FLim}(h,c,\alpha) \leftrightarrow \mathrm{RFn}(h) \land c \in \mathbb{R} \land \alpha \in \mathbb{R} \land \mathrm{PNbhd}(h,c) \land \forall \varepsilon \in \mathbb{R}\,(0 < \varepsilon \to \exists \delta\,(\delta \in \mathbb{R} \land 0 < \delta \land \varphi_{47}))$
\item $\varphi_{124} :\equiv \mathrm{FLim}(g,c,\alpha) \to \mathrm{RFn}(h) \to \mathrm{PNbhd}(h,c) \to r \in \mathbb{R} \to 0 < r \to \varphi_{97}$
\item $\varphi_{125} :\equiv \forall f\,\forall c\,\forall \alpha\,(\mathrm{FLim}(f,c,\alpha) \leftrightarrow \mathrm{RFn}(f) \land c \in \mathbb{R} \land \alpha \in \mathbb{R} \land \mathrm{PNbhd}(f,c) \land \forall \varepsilon \in \mathbb{R}\,(0 < \varepsilon \to \varphi_{76}))$
\item $\varphi_{126} :\equiv d_{2} \in \mathbb{R} \land 0 < d_{2} \land \varphi_{48} \to d_{3} \in \mathbb{R} \land 0 < d_{3} \land d_{3} \leq d_{1} \land d_{3} \leq d_{2} \to \varphi_{18}$
\item $\varphi_{127} :\equiv \mathrm{FLim}(g,c,\alpha) \to \mathrm{RFn}(h) \to \mathrm{PNbhd}(h,c) \to r \in \mathbb{R} \to 0 < r \to \varphi_{106}$
\item $\varphi_{128} :\equiv \forall \alpha\,\forall \varepsilon\,(f(x) \in \mathbb{R} \to g(x) \in \mathbb{R} \to h(x) \in \mathbb{R} \to \varphi_{98})$
\item $\varphi_{129} :\equiv \forall \delta\,(\neg (\delta \in \mathbb{R} \land 0 < \delta \land \varphi_{47})) \to \neg (\delta \in \mathbb{R} \land 0 < \delta \land \varphi_{47})$
\item $\varphi_{130} :\equiv d_{2} \in \mathbb{R} \land 0 < d_{2} \land \varphi_{48} \to d_{3} \in \mathbb{R} \land 0 < d_{3} \land d_{3} \leq d_{1} \land d_{3} \leq d_{2} \to \varphi_{35}$
\item $\varphi_{131} :\equiv \forall d_{1}\,(\neg (d_{1} \in \mathbb{R} \land 0 < d_{1} \land \varphi_{49})) \to \forall \delta\,\forall d_{1}\,(\neg (d_{1} \in \mathbb{R} \land 0 < d_{1} \land \varphi_{49}))$
\item $\varphi_{132} :\equiv \forall d_{2}\,(\neg (d_{2} \in \mathbb{R} \land 0 < d_{2} \land \varphi_{48})) \to \forall \delta\,\forall d_{2}\,(\neg (d_{2} \in \mathbb{R} \land 0 < d_{2} \land \varphi_{48}))$
\item $\varphi_{133} :\equiv d_{1} \in \mathbb{R} \land 0 < d_{1} \land \varphi_{49} \to d_{2} \in \mathbb{R} \land 0 < d_{2} \land \varphi_{48} \to d_{2} \in \mathbb{R}$
\item $\varphi_{134} :\equiv d_{1} \in \mathbb{R} \land 0 < d_{1} \land \varphi_{49} \to d_{2} \in \mathbb{R} \land 0 < d_{2} \land \varphi_{48} \to d_{1} \in \mathbb{R}$
\item $\varphi_{135} :\equiv 0 < \varepsilon \to d_{1} \in \mathbb{R} \land 0 < d_{1} \land \varphi_{49} \to \exists \delta\,(\delta \in \mathbb{R} \land 0 < \delta \land \varphi_{47})$
\item $\varphi_{136} :\equiv (\mathrm{FLim}(f,c,\alpha) \to \mathrm{FLim}(f,c,\alpha)) \to \mathrm{FLim}(f,c,\alpha) \to \varphi_{117}$
\item $\varphi_{137} :\equiv 0 < \varepsilon \to d_{1} \in \mathbb{R} \land 0 < d_{1} \land \varphi_{49} \to d_{2} \in \mathbb{R} \land 0 < d_{2} \land \varphi_{48} \to 0 < d_{2}$
\item $\varphi_{138} :\equiv 0 < \varepsilon \to d_{1} \in \mathbb{R} \land 0 < d_{1} \land \varphi_{49} \to d_{2} \in \mathbb{R} \land 0 < d_{2} \land \varphi_{48} \to 0 < d_{1}$
\item $\varphi_{139} :\equiv d_{2} \in \mathbb{R} \land 0 < d_{2} \land \varphi_{48} \to d_{3} \in \mathbb{R} \land 0 < d_{3} \land d_{3} \leq d_{1} \land d_{3} \leq d_{2} \to \varphi_{58}$
\item $\varphi_{140} :\equiv (\mathrm{FLim}(f,c,\alpha) \to \mathrm{FLim}(f,c,\alpha)) \to \mathrm{FLim}(f,c,\alpha) \to \varphi_{118}$
\item $\varphi_{141} :\equiv (\mathrm{FLim}(f,c,\alpha) \to \mathrm{FLim}(f,c,\alpha)) \to \mathrm{FLim}(f,c,\alpha) \to \varphi_{119}$
\item $\varphi_{142} :\equiv \varepsilon \in \mathbb{R} \to 0 < \varepsilon \to d_{1} \in \mathbb{R} \land 0 < d_{1} \land \varphi_{49} \to d_{1} \in \mathbb{R} \land 0 < d_{1} \land \varphi_{49}$
\item $\varphi_{143} :\equiv \forall x \in \mathrm{dom}(h)\,(0 < \lvert x - c \rvert \to \lvert x - c \rvert < r \to \varphi_{21}) \to x \in \mathrm{dom}(h) \to 0 < \lvert x - c \rvert \to \lvert x - c \rvert < r \to \varphi_{21}$
\item $\varphi_{144} :\equiv d_{2} \in \mathbb{R} \land 0 < d_{2} \land \varphi_{48} \to d_{3} \in \mathbb{R} \land 0 < d_{3} \land d_{3} \leq d_{1} \land d_{3} \leq d_{2} \to \varphi_{69}$
\item $\varphi_{145} :\equiv d_{2} \in \mathbb{R} \land 0 < d_{2} \land \varphi_{48} \to d_{3} \in \mathbb{R} \land 0 < d_{3} \land d_{3} \leq d_{1} \land d_{3} \leq d_{2} \to \varphi_{70}$
\item $\varphi_{146} :\equiv d_{2} \in \mathbb{R} \land 0 < d_{2} \land \varphi_{48} \to d_{3} \in \mathbb{R} \land 0 < d_{3} \land d_{3} \leq d_{1} \land d_{3} \leq d_{2} \to \varphi_{71}$
\item $\varphi_{147} :\equiv d_{2} \in \mathbb{R} \land 0 < d_{2} \land \varphi_{48} \to d_{3} \in \mathbb{R} \land 0 < d_{3} \land d_{3} \leq d_{1} \land d_{3} \leq d_{2} \to \varphi_{72}$
\item $\varphi_{148} :\equiv \forall x \in \mathrm{dom}(h)\,(0 < \lvert x - c \rvert \to \lvert x - c \rvert < r \to \varphi_{21}) \to \forall x \in \mathrm{dom}(h)\,(0 < \lvert x - c \rvert \to \lvert x - c \rvert < r \to \varphi_{21})$
\item $\varphi_{149} :\equiv \forall x \in \mathrm{dom}(h)\,(0 < \lvert x - c \rvert \to \lvert x - c \rvert < r \to \varphi_{21}) \to \varepsilon \in \mathbb{R} \to 0 < \varepsilon \to \varphi_{76}$
\item $\varphi_{150} :\equiv \forall x \in \mathrm{dom}(h)\,(0 < \lvert x - c \rvert \to \lvert x - c \rvert < r \to \varphi_{21}) \to \varepsilon \in \mathbb{R} \to 0 < \varepsilon \to \exists \delta\,(\delta \in \mathbb{R} \land 0 < \delta \land \varphi_{47})$
\item $\varphi_{151} :\equiv \forall x \in \mathrm{dom}(h)\,(0 < \lvert x - c \rvert \to \lvert x - c \rvert < r \to \varphi_{21}) \to \forall \varepsilon \in \mathbb{R}\,(0 < \varepsilon \to \exists \delta\,(\delta \in \mathbb{R} \land 0 < \delta \land \varphi_{47}))$
\item $\varphi_{152} :\equiv (\mathrm{FLim}(g,c,\alpha) \to \mathrm{RFn}(g)) \to \varphi_{1} \to \mathrm{FLim}(f,c,\alpha) \to \varphi_{116}$
\item $\varphi_{153} :\equiv d_{2} \in \mathbb{R} \land 0 < d_{2} \land \varphi_{48} \to d_{3} \in \mathbb{R} \land 0 < d_{3} \land d_{3} \leq d_{1} \land d_{3} \leq d_{2} \to \varphi_{95}$
\item $\varphi_{154} :\equiv \mathrm{FLim}(g,c,\alpha) \to \mathrm{RFn}(h) \to \mathrm{PNbhd}(h,c) \to r \in \mathbb{R} \to 0 < r \to \varphi_{143}$
\item $\varphi_{155} :\equiv 0 < \varepsilon \to d_{1} \in \mathbb{R} \land 0 < d_{1} \land \varphi_{49} \to d_{2} \in \mathbb{R} \land 0 < d_{2} \land \varphi_{48} \to \varphi_{36}$
\item $\varphi_{156} :\equiv 0 < \varepsilon \to d_{1} \in \mathbb{R} \land 0 < d_{1} \land \varphi_{49} \to d_{2} \in \mathbb{R} \land 0 < d_{2} \land \varphi_{48} \to \varphi_{38}$
\item $\varphi_{157} :\equiv \mathrm{FLim}(g,c,\alpha) \to \mathrm{RFn}(h) \to \mathrm{PNbhd}(h,c) \to r \in \mathbb{R} \to 0 < r \to \varphi_{148}$
\item $\varphi_{158} :\equiv d_{2} \in \mathbb{R} \land 0 < d_{2} \land \varphi_{48} \to d_{3} \in \mathbb{R} \land 0 < d_{3} \land d_{3} \leq d_{1} \land d_{3} \leq d_{2} \to \varphi_{105}$
\item $\varphi_{159} :\equiv \forall \varepsilon\,(\mathrm{FLim}(g,c,\alpha) \to \varepsilon \in \mathbb{R} \to 0 < \varepsilon \to \varphi_{75}) \to \mathrm{FLim}(g,c,\alpha) \to \varepsilon \in \mathbb{R} \to 0 < \varepsilon \to \varphi_{75}$
\item $\varphi_{160} :\equiv \forall \varepsilon\,(\mathrm{FLim}(f,c,\alpha) \to \varepsilon \in \mathbb{R} \to 0 < \varepsilon \to \varphi_{76}) \to \mathrm{FLim}(f,c,\alpha) \to \varepsilon \in \mathbb{R} \to 0 < \varepsilon \to \varphi_{76}$
\item $\varphi_{161} :\equiv \forall \alpha\,\forall \varepsilon\,(\mathrm{FLim}(g,c,\alpha) \to \varepsilon \in \mathbb{R} \to 0 < \varepsilon \to \varphi_{75}) \to \forall \varepsilon\,(\mathrm{FLim}(g,c,\alpha) \to \varepsilon \in \mathbb{R} \to 0 < \varepsilon \to \varphi_{75})$
\item $\varphi_{162} :\equiv \forall \alpha\,\forall \varepsilon\,(\mathrm{FLim}(f,c,\alpha) \to \varepsilon \in \mathbb{R} \to 0 < \varepsilon \to \varphi_{76}) \to \forall \varepsilon\,(\mathrm{FLim}(f,c,\alpha) \to \varepsilon \in \mathbb{R} \to 0 < \varepsilon \to \varphi_{76})$
\item $\varphi_{163} :\equiv \mathrm{FLim}(g,c,\alpha) \to \mathrm{RFn}(h) \to \mathrm{PNbhd}(h,c) \to r \in \mathbb{R} \to 0 < r \to \varphi_{149}$
\item $\varphi_{164} :\equiv \mathrm{FLim}(g,c,\alpha) \to \mathrm{RFn}(h) \to \mathrm{PNbhd}(h,c) \to r \in \mathbb{R} \to 0 < r \to \varphi_{150}$
\item $\varphi_{165} :\equiv d_{1} \in \mathbb{R} \land 0 < d_{1} \land \varphi_{49} \to d_{2} \in \mathbb{R} \land 0 < d_{2} \land \varphi_{48} \to d_{2} \in \mathbb{R} \land 0 < d_{2} \land \varphi_{48}$
\item $\varphi_{166} :\equiv \mathrm{FLim}(g,c,\alpha) \to \mathrm{RFn}(h) \to \mathrm{PNbhd}(h,c) \to r \in \mathbb{R} \to 0 < r \to \varphi_{151}$
\item $\varphi_{167} :\equiv \forall c\,\forall \alpha\,\forall \varepsilon\,(\mathrm{FLim}(g,c,\alpha) \to \varepsilon \in \mathbb{R} \to 0 < \varepsilon \to \varphi_{75}) \to \forall \alpha\,\forall \varepsilon\,(\mathrm{FLim}(g,c,\alpha) \to \varepsilon \in \mathbb{R} \to 0 < \varepsilon \to \varphi_{75})$
\item $\varphi_{168} :\equiv \forall c\,\forall \alpha\,\forall \varepsilon\,(\mathrm{FLim}(f,c,\alpha) \to \varepsilon \in \mathbb{R} \to 0 < \varepsilon \to \varphi_{76}) \to \forall \alpha\,\forall \varepsilon\,(\mathrm{FLim}(f,c,\alpha) \to \varepsilon \in \mathbb{R} \to 0 < \varepsilon \to \varphi_{76})$
\item $\varphi_{169} :\equiv d_{1} \in \mathbb{R} \land 0 < d_{1} \land \varphi_{49} \to d_{2} \in \mathbb{R} \land 0 < d_{2} \land \varphi_{48} \to \exists \delta\,(\delta \in \mathbb{R} \land 0 < \delta \land \varphi_{47})$
\item $\varphi_{170} :\equiv \varphi_{112} \to \forall c\,\forall \alpha\,\forall \varepsilon\,(\mathrm{FLim}(g,c,\alpha) \to \varepsilon \in \mathbb{R} \to 0 < \varepsilon \to \varphi_{75})$
\item $\varphi_{171} :\equiv \varphi_{112} \to \forall c\,\forall \alpha\,\forall \varepsilon\,(\mathrm{FLim}(f,c,\alpha) \to \varepsilon \in \mathbb{R} \to 0 < \varepsilon \to \varphi_{76})$
\item $\varphi_{172} :\equiv \forall d_{1}\,(\varphi_{121} \land \varepsilon \in \mathbb{R} \land 0 < \varepsilon \land \neg \exists \delta\,(\delta \in \mathbb{R} \land 0 < \delta \land \varphi_{47}))$
\item $\varphi_{173} :\equiv \forall \varepsilon\,(\varphi_{121}) \to \forall \varepsilon \in \mathbb{R}\,(0 < \varepsilon \to \exists \delta\,(\delta \in \mathbb{R} \land 0 < \delta \land \varphi_{47}))$
\item $\varphi_{174} :\equiv \forall x \in \mathrm{dom}(h)\,(0 < \lvert x - c \rvert \to \lvert x - c \rvert < r \to \varphi_{21}) \to \varepsilon \in \mathbb{R} \to 0 < \varepsilon \to d_{1} \in \mathbb{R} \land 0 < d_{1} \land \varphi_{49} \to \varphi_{75}$
\item $\varphi_{175} :\equiv \forall x \in \mathrm{dom}(h)\,(0 < \lvert x - c \rvert \to \lvert x - c \rvert < r \to \varphi_{21}) \to \varepsilon \in \mathbb{R} \to \varphi_{135}$
\item $\varphi_{176} :\equiv 0 < \varepsilon \to d_{1} \in \mathbb{R} \land 0 < d_{1} \land \varphi_{49} \to d_{2} \in \mathbb{R} \land 0 < d_{2} \land \varphi_{48} \to \varphi_{82}$
\item $\varphi_{177} :\equiv 0 < \varepsilon \to d_{1} \in \mathbb{R} \land 0 < d_{1} \land \varphi_{49} \to d_{2} \in \mathbb{R} \land 0 < d_{2} \land \varphi_{48} \to \varphi_{83}$
\item $\varphi_{178} :\equiv 0 < r \to \forall x \in \mathrm{dom}(h)\,(0 < \lvert x - c \rvert \to \lvert x - c \rvert < r \to \varphi_{21}) \to \varphi_{142}$
\item $\varphi_{179} :\equiv \forall x \in \mathrm{dom}(h)\,(0 < \lvert x - c \rvert \to \lvert x - c \rvert < r \to \varphi_{21}) \to \varepsilon \in \mathbb{R} \to \varphi_{137}$
\item $\varphi_{180} :\equiv \forall x \in \mathrm{dom}(h)\,(0 < \lvert x - c \rvert \to \lvert x - c \rvert < r \to \varphi_{21}) \to \varepsilon \in \mathbb{R} \to \varphi_{138}$
\item $\varphi_{181} :\equiv 0 < \varepsilon \to d_{1} \in \mathbb{R} \land 0 < d_{1} \land \varphi_{49} \to d_{2} \in \mathbb{R} \land 0 < d_{2} \land \varphi_{48} \to \varphi_{85}$
\item $\varphi_{182} :\equiv 0 < \varepsilon \to d_{1} \in \mathbb{R} \land 0 < d_{1} \land \varphi_{49} \to d_{2} \in \mathbb{R} \land 0 < d_{2} \land \varphi_{48} \to \varphi_{86}$
\item $\varphi_{183} :\equiv 0 < \varepsilon \to d_{1} \in \mathbb{R} \land 0 < d_{1} \land \varphi_{49} \to d_{2} \in \mathbb{R} \land 0 < d_{2} \land \varphi_{48} \to \varphi_{87}$
\item $\varphi_{184} :\equiv 0 < \varepsilon \to d_{1} \in \mathbb{R} \land 0 < d_{1} \land \varphi_{49} \to d_{2} \in \mathbb{R} \land 0 < d_{2} \land \varphi_{48} \to \varphi_{88}$
\item $\varphi_{185} :\equiv 0 < \varepsilon \to d_{1} \in \mathbb{R} \land 0 < d_{1} \land \varphi_{49} \to d_{2} \in \mathbb{R} \land 0 < d_{2} \land \varphi_{48} \to \varphi_{89}$
\item $\varphi_{186} :\equiv 0 < \varepsilon \to d_{1} \in \mathbb{R} \land 0 < d_{1} \land \varphi_{49} \to d_{2} \in \mathbb{R} \land 0 < d_{2} \land \varphi_{48} \to \varphi_{90}$
\item $\varphi_{187} :\equiv 0 < \varepsilon \to d_{1} \in \mathbb{R} \land 0 < d_{1} \land \varphi_{49} \to d_{2} \in \mathbb{R} \land 0 < d_{2} \land \varphi_{48} \to \varphi_{91}$
\item $\varphi_{188} :\equiv 0 < \varepsilon \to d_{1} \in \mathbb{R} \land 0 < d_{1} \land \varphi_{49} \to d_{2} \in \mathbb{R} \land 0 < d_{2} \land \varphi_{48} \to \varphi_{92}$
\item $\varphi_{189} :\equiv 0 < \varepsilon \to d_{1} \in \mathbb{R} \land 0 < d_{1} \land \varphi_{49} \to d_{2} \in \mathbb{R} \land 0 < d_{2} \land \varphi_{48} \to \varphi_{93}$
\item $\varphi_{190} :\equiv 0 < \varepsilon \to d_{1} \in \mathbb{R} \land 0 < d_{1} \land \varphi_{49} \to d_{2} \in \mathbb{R} \land 0 < d_{2} \land \varphi_{48} \to \varphi_{94}$
\item $\varphi_{191} :\equiv 0 < \varepsilon \to d_{1} \in \mathbb{R} \land 0 < d_{1} \land \varphi_{49} \to d_{2} \in \mathbb{R} \land 0 < d_{2} \land \varphi_{48} \to \varphi_{96}$
\item $\varphi_{192} :\equiv \forall x \in \mathrm{dom}(h)\,(0 < \lvert x - c \rvert \to \lvert x - c \rvert < r \to \varphi_{21}) \to \varepsilon \in \mathbb{R} \to 0 < \varepsilon \to \varphi_{133}$
\item $\varphi_{193} :\equiv \forall x \in \mathrm{dom}(h)\,(0 < \lvert x - c \rvert \to \lvert x - c \rvert < r \to \varphi_{21}) \to \varepsilon \in \mathbb{R} \to 0 < \varepsilon \to \varphi_{134}$
\item $\varphi_{194} :\equiv 0 < \varepsilon \to d_{1} \in \mathbb{R} \land 0 < d_{1} \land \varphi_{49} \to d_{2} \in \mathbb{R} \land 0 < d_{2} \land \varphi_{48} \to \varphi_{100}$
\item $\varphi_{195} :\equiv 0 < \varepsilon \to d_{1} \in \mathbb{R} \land 0 < d_{1} \land \varphi_{49} \to d_{2} \in \mathbb{R} \land 0 < d_{2} \land \varphi_{48} \to \varphi_{101}$
\item $\varphi_{196} :\equiv 0 < \varepsilon \to d_{1} \in \mathbb{R} \land 0 < d_{1} \land \varphi_{49} \to d_{2} \in \mathbb{R} \land 0 < d_{2} \land \varphi_{48} \to \varphi_{102}$
\item $\varphi_{197} :\equiv 0 < \varepsilon \to d_{1} \in \mathbb{R} \land 0 < d_{1} \land \varphi_{49} \to d_{2} \in \mathbb{R} \land 0 < d_{2} \land \varphi_{48} \to \varphi_{103}$
\item $\varphi_{198} :\equiv 0 < \varepsilon \to d_{1} \in \mathbb{R} \land 0 < d_{1} \land \varphi_{49} \to d_{2} \in \mathbb{R} \land 0 < d_{2} \land \varphi_{48} \to \varphi_{104}$
\item $\varphi_{199} :\equiv 0 < \varepsilon \to d_{1} \in \mathbb{R} \land 0 < d_{1} \land \varphi_{49} \to d_{2} \in \mathbb{R} \land 0 < d_{2} \land \varphi_{48} \to \varphi_{109}$
\item $\varphi_{200} :\equiv 0 < \varepsilon \to d_{1} \in \mathbb{R} \land 0 < d_{1} \land \varphi_{49} \to d_{2} \in \mathbb{R} \land 0 < d_{2} \land \varphi_{48} \to \varphi_{110}$
\item $\varphi_{201} :\equiv 0 < \varepsilon \to d_{1} \in \mathbb{R} \land 0 < d_{1} \land \varphi_{49} \to d_{2} \in \mathbb{R} \land 0 < d_{2} \land \varphi_{48} \to \varphi_{111}$
\item $\varphi_{202} :\equiv \forall x \in \mathrm{dom}(h)\,(0 < \lvert x - c \rvert \to \lvert x - c \rvert < r \to \varphi_{21}) \to \varepsilon \in \mathbb{R} \to 0 < \varepsilon \to d_{1} \in \mathbb{R} \land 0 < d_{1} \land \varphi_{49} \to \varphi_{84}$
\item $\varphi_{203} :\equiv 0 < \varepsilon \to d_{1} \in \mathbb{R} \land 0 < d_{1} \land \varphi_{49} \to d_{2} \in \mathbb{R} \land 0 < d_{2} \land \varphi_{48} \to \varphi_{115}$
\item $\varphi_{204} :\equiv \forall x \in \mathrm{dom}(h)\,(0 < \lvert x - c \rvert \to \lvert x - c \rvert < r \to \varphi_{21}) \to \varepsilon \in \mathbb{R} \to 0 < \varepsilon \to d_{1} \in \mathbb{R} \land 0 < d_{1} \land \varphi_{49} \to \varphi_{99}$
\item $\varphi_{205} :\equiv \forall x \in \mathrm{dom}(h)\,(0 < \lvert x - c \rvert \to \lvert x - c \rvert < r \to \varphi_{21}) \to \varepsilon \in \mathbb{R} \to 0 < \varepsilon \to d_{1} \in \mathbb{R} \land 0 < d_{1} \land \varphi_{49} \to \varphi_{108}$
\item $\varphi_{206} :\equiv \mathrm{FLim}(g,c,\alpha) \to \mathrm{RFn}(h) \to \mathrm{PNbhd}(h,c) \to r \in \mathbb{R} \to \varphi_{178}$
\item $\varphi_{207} :\equiv \mathrm{FLim}(g,c,\alpha) \to \mathrm{RFn}(h) \to \mathrm{PNbhd}(h,c) \to r \in \mathbb{R} \to 0 < r \to \varphi_{174}$
\item $\varphi_{208} :\equiv \mathrm{FLim}(g,c,\alpha) \to \mathrm{RFn}(h) \to \mathrm{PNbhd}(h,c) \to r \in \mathbb{R} \to 0 < r \to \varphi_{175}$
\item $\varphi_{209} :\equiv \forall \varepsilon\,(f(x) \in \mathbb{R} \to g(x) \in \mathbb{R} \to h(x) \in \mathbb{R} \to \varphi_{98}) \to f(x) \in \mathbb{R} \to g(x) \in \mathbb{R} \to h(x) \in \mathbb{R} \to \varphi_{98}$
\item $\varphi_{210} :\equiv \mathrm{FLim}(g,c,\alpha) \to \mathrm{RFn}(h) \to \mathrm{PNbhd}(h,c) \to r \in \mathbb{R} \to 0 < r \to \varphi_{179}$
\item $\varphi_{211} :\equiv \mathrm{FLim}(g,c,\alpha) \to \mathrm{RFn}(h) \to \mathrm{PNbhd}(h,c) \to r \in \mathbb{R} \to 0 < r \to \varphi_{180}$
\item $\varphi_{212} :\equiv \mathrm{FLim}(g,c,\alpha) \to \mathrm{RFn}(h) \to \mathrm{PNbhd}(h,c) \to r \in \mathbb{R} \to 0 < r \to \varphi_{192}$
\item $\varphi_{213} :\equiv \mathrm{FLim}(g,c,\alpha) \to \mathrm{RFn}(h) \to \mathrm{PNbhd}(h,c) \to r \in \mathbb{R} \to 0 < r \to \varphi_{193}$
\item $\varphi_{214} :\equiv \varphi_{128} \to \forall \varepsilon\,(f(x) \in \mathbb{R} \to g(x) \in \mathbb{R} \to h(x) \in \mathbb{R} \to \varphi_{98})$
\item $\varphi_{215} :\equiv \forall x \in \mathrm{dom}(h)\,(0 < \lvert x - c \rvert \to \lvert x - c \rvert < r \to \varphi_{21}) \to \varepsilon \in \mathbb{R} \to \varphi_{155}$
\item $\varphi_{216} :\equiv \varphi_{121} \land \varepsilon \in \mathbb{R} \land 0 < \varepsilon \land (d_{1} \in \mathbb{R} \land 0 < d_{1} \land \varphi_{49}) \land \neg \exists \delta\,(\delta \in \mathbb{R} \land 0 < \delta \land \varphi_{47})$
\item $\varphi_{217} :\equiv \varphi_{121} \land \varepsilon \in \mathbb{R} \land 0 < \varepsilon \land \neg \exists \delta\,(\delta \in \mathbb{R} \land 0 < \delta \land \varphi_{47}) \to \neg (d_{1} \in \mathbb{R} \land 0 < d_{1} \land \varphi_{49})$
\item $\varphi_{218} :\equiv \forall x \in \mathrm{dom}(h)\,(0 < \lvert x - c \rvert \to \lvert x - c \rvert < r \to \varphi_{21}) \to \varepsilon \in \mathbb{R} \to \varphi_{156}$
\item $\varphi_{219} :\equiv \forall \delta\,(\forall d_{1}\,(\neg (d_{1} \in \mathbb{R} \land 0 < d_{1} \land \varphi_{49})) \to \varphi_{73}) \to \forall \delta\,\forall d_{1}\,(\neg (d_{1} \in \mathbb{R} \land 0 < d_{1} \land \varphi_{49})) \to \forall \delta\,(\varphi_{73})$
\item $\varphi_{220} :\equiv \forall \delta\,(\forall d_{2}\,(\neg (d_{2} \in \mathbb{R} \land 0 < d_{2} \land \varphi_{48})) \to \varphi_{74}) \to \forall \delta\,\forall d_{2}\,(\neg (d_{2} \in \mathbb{R} \land 0 < d_{2} \land \varphi_{48})) \to \forall \delta\,(\varphi_{74})$
\item $\varphi_{221} :\equiv \mathrm{FLim}(g,c,\alpha) \to \mathrm{RFn}(h) \to \mathrm{PNbhd}(h,c) \to r \in \mathbb{R} \to 0 < r \to \varphi_{202}$
\item $\varphi_{222} :\equiv \forall x \in \mathrm{dom}(h)\,(0 < \lvert x - c \rvert \to \lvert x - c \rvert < r \to \varphi_{21}) \to \varepsilon \in \mathbb{R} \to 0 < \varepsilon \to d_{1} \in \mathbb{R} \land 0 < d_{1} \land \varphi_{49} \to \varphi_{126}$
\item $\varphi_{223} :\equiv \mathrm{FLim}(g,c,\alpha) \to \mathrm{RFn}(h) \to \mathrm{PNbhd}(h,c) \to r \in \mathbb{R} \to 0 < r \to \varphi_{204}$
\item $\varphi_{224} :\equiv \forall x \in \mathrm{dom}(h)\,(0 < \lvert x - c \rvert \to \lvert x - c \rvert < r \to \varphi_{21}) \to \varepsilon \in \mathbb{R} \to 0 < \varepsilon \to \varphi_{165}$
\item $\varphi_{225} :\equiv \mathrm{FLim}(g,c,\alpha) \to \mathrm{RFn}(h) \to \mathrm{PNbhd}(h,c) \to r \in \mathbb{R} \to 0 < r \to \varphi_{205}$
\item $\varphi_{226} :\equiv \forall x \in \mathrm{dom}(h)\,(0 < \lvert x - c \rvert \to \lvert x - c \rvert < r \to \varphi_{21}) \to \varepsilon \in \mathbb{R} \to 0 < \varepsilon \to \varphi_{169}$
\item $\varphi_{227} :\equiv \forall x \in \mathrm{dom}(h)\,(0 < \lvert x - c \rvert \to \lvert x - c \rvert < r \to \varphi_{21}) \to \varepsilon \in \mathbb{R} \to 0 < \varepsilon \to d_{1} \in \mathbb{R} \land 0 < d_{1} \land \varphi_{49} \to \varphi_{130}$
\item $\varphi_{228} :\equiv \forall x \in \mathrm{dom}(h)\,(0 < \lvert x - c \rvert \to \lvert x - c \rvert < r \to \varphi_{21}) \to \varepsilon \in \mathbb{R} \to 0 < \varepsilon \to d_{1} \in \mathbb{R} \land 0 < d_{1} \land \varphi_{49} \to \varphi_{139}$
\item $\varphi_{229} :\equiv \mathrm{FLim}(g,c,\alpha) \to \mathrm{RFn}(h) \to \mathrm{PNbhd}(h,c) \to r \in \mathbb{R} \to 0 < r \to \varphi_{215}$
\item $\varphi_{230} :\equiv \mathrm{FLim}(g,c,\alpha) \to \mathrm{RFn}(h) \to \mathrm{PNbhd}(h,c) \to r \in \mathbb{R} \to 0 < r \to \varphi_{218}$
\item $\varphi_{231} :\equiv \forall x \in \mathrm{dom}(h)\,(0 < \lvert x - c \rvert \to \lvert x - c \rvert < r \to \varphi_{21}) \to \varepsilon \in \mathbb{R} \to 0 < \varepsilon \to d_{1} \in \mathbb{R} \land 0 < d_{1} \land \varphi_{49} \to \varphi_{144}$
\item $\varphi_{232} :\equiv \forall x \in \mathrm{dom}(h)\,(0 < \lvert x - c \rvert \to \lvert x - c \rvert < r \to \varphi_{21}) \to \varepsilon \in \mathbb{R} \to 0 < \varepsilon \to d_{1} \in \mathbb{R} \land 0 < d_{1} \land \varphi_{49} \to \varphi_{145}$
\item $\varphi_{233} :\equiv \forall x \in \mathrm{dom}(h)\,(0 < \lvert x - c \rvert \to \lvert x - c \rvert < r \to \varphi_{21}) \to \varepsilon \in \mathbb{R} \to 0 < \varepsilon \to d_{1} \in \mathbb{R} \land 0 < d_{1} \land \varphi_{49} \to \varphi_{146}$
\item $\varphi_{234} :\equiv \forall x \in \mathrm{dom}(h)\,(0 < \lvert x - c \rvert \to \lvert x - c \rvert < r \to \varphi_{21}) \to \varepsilon \in \mathbb{R} \to 0 < \varepsilon \to d_{1} \in \mathbb{R} \land 0 < d_{1} \land \varphi_{49} \to \varphi_{147}$
\item $\varphi_{235} :\equiv \forall x \in \mathrm{dom}(h)\,(0 < \lvert x - c \rvert \to \lvert x - c \rvert < r \to \varphi_{21}) \to \varepsilon \in \mathbb{R} \to \varphi_{176}$
\item $\varphi_{236} :\equiv \forall x \in \mathrm{dom}(h)\,(0 < \lvert x - c \rvert \to \lvert x - c \rvert < r \to \varphi_{21}) \to \varepsilon \in \mathbb{R} \to \varphi_{177}$
\item $\varphi_{237} :\equiv \forall x \in \mathrm{dom}(h)\,(0 < \lvert x - c \rvert \to \lvert x - c \rvert < r \to \varphi_{21}) \to \varepsilon \in \mathbb{R} \to \varphi_{181}$
\item $\varphi_{238} :\equiv \varphi_{121} \land \varepsilon \in \mathbb{R} \land 0 < \varepsilon \land (d_{1} \in \mathbb{R} \land 0 < d_{1} \land \varphi_{49}) \land (d_{2} \in \mathbb{R} \land 0 < d_{2} \land \varphi_{48}) \land (d_{3} \in \mathbb{R} \land 0 < d_{3} \land d_{3} \leq d_{1} \land d_{3} \leq d_{2}) \land (\delta \in \mathbb{R} \land 0 < \delta \land \delta \leq d_{3} \land \delta \leq r)$
\item $\varphi_{239} :\equiv \forall x \in \mathrm{dom}(h)\,(0 < \lvert x - c \rvert \to \lvert x - c \rvert < r \to \varphi_{21}) \to \varepsilon \in \mathbb{R} \to \varphi_{182}$
\item $\varphi_{240} :\equiv \forall x \in \mathrm{dom}(h)\,(0 < \lvert x - c \rvert \to \lvert x - c \rvert < r \to \varphi_{21}) \to \varepsilon \in \mathbb{R} \to \varphi_{183}$
\item $\varphi_{241} :\equiv \forall x \in \mathrm{dom}(h)\,(0 < \lvert x - c \rvert \to \lvert x - c \rvert < r \to \varphi_{21}) \to \varepsilon \in \mathbb{R} \to \varphi_{184}$
\item $\varphi_{242} :\equiv \forall x \in \mathrm{dom}(h)\,(0 < \lvert x - c \rvert \to \lvert x - c \rvert < r \to \varphi_{21}) \to \varepsilon \in \mathbb{R} \to \varphi_{185}$
\item $\varphi_{243} :\equiv \forall x \in \mathrm{dom}(h)\,(0 < \lvert x - c \rvert \to \lvert x - c \rvert < r \to \varphi_{21}) \to \varepsilon \in \mathbb{R} \to \varphi_{186}$
\item $\varphi_{244} :\equiv \forall x \in \mathrm{dom}(h)\,(0 < \lvert x - c \rvert \to \lvert x - c \rvert < r \to \varphi_{21}) \to \varepsilon \in \mathbb{R} \to \varphi_{187}$
\item $\varphi_{245} :\equiv \forall x \in \mathrm{dom}(h)\,(0 < \lvert x - c \rvert \to \lvert x - c \rvert < r \to \varphi_{21}) \to \varepsilon \in \mathbb{R} \to \varphi_{188}$
\item $\varphi_{246} :\equiv \forall x \in \mathrm{dom}(h)\,(0 < \lvert x - c \rvert \to \lvert x - c \rvert < r \to \varphi_{21}) \to \varepsilon \in \mathbb{R} \to \varphi_{189}$
\item $\varphi_{247} :\equiv \forall x \in \mathrm{dom}(h)\,(0 < \lvert x - c \rvert \to \lvert x - c \rvert < r \to \varphi_{21}) \to \varepsilon \in \mathbb{R} \to \varphi_{190}$
\item $\varphi_{248} :\equiv \forall x \in \mathrm{dom}(h)\,(0 < \lvert x - c \rvert \to \lvert x - c \rvert < r \to \varphi_{21}) \to \varepsilon \in \mathbb{R} \to \varphi_{191}$
\item $\varphi_{249} :\equiv \forall x \in \mathrm{dom}(h)\,(0 < \lvert x - c \rvert \to \lvert x - c \rvert < r \to \varphi_{21}) \to \varepsilon \in \mathbb{R} \to \varphi_{194}$
\item $\varphi_{250} :\equiv \forall x \in \mathrm{dom}(h)\,(0 < \lvert x - c \rvert \to \lvert x - c \rvert < r \to \varphi_{21}) \to \varepsilon \in \mathbb{R} \to \varphi_{195}$
\item $\varphi_{251} :\equiv \forall x \in \mathrm{dom}(h)\,(0 < \lvert x - c \rvert \to \lvert x - c \rvert < r \to \varphi_{21}) \to \varepsilon \in \mathbb{R} \to \varphi_{196}$
\item $\varphi_{252} :\equiv \forall x \in \mathrm{dom}(h)\,(0 < \lvert x - c \rvert \to \lvert x - c \rvert < r \to \varphi_{21}) \to \varepsilon \in \mathbb{R} \to \varphi_{197}$
\item $\varphi_{253} :\equiv \forall x \in \mathrm{dom}(h)\,(0 < \lvert x - c \rvert \to \lvert x - c \rvert < r \to \varphi_{21}) \to \varepsilon \in \mathbb{R} \to \varphi_{198}$
\item $\varphi_{254} :\equiv \forall x \in \mathrm{dom}(h)\,(0 < \lvert x - c \rvert \to \lvert x - c \rvert < r \to \varphi_{21}) \to \varepsilon \in \mathbb{R} \to \varphi_{199}$
\item $\varphi_{255} :\equiv \forall x \in \mathrm{dom}(h)\,(0 < \lvert x - c \rvert \to \lvert x - c \rvert < r \to \varphi_{21}) \to \varepsilon \in \mathbb{R} \to \varphi_{200}$
\item $\varphi_{256} :\equiv \forall x \in \mathrm{dom}(h)\,(0 < \lvert x - c \rvert \to \lvert x - c \rvert < r \to \varphi_{21}) \to \varepsilon \in \mathbb{R} \to \varphi_{201}$
\item $\varphi_{257} :\equiv \mathrm{FLim}(g,c,\alpha) \to \mathrm{RFn}(h) \to \mathrm{PNbhd}(h,c) \to r \in \mathbb{R} \to 0 < r \to \varphi_{222}$
\item $\varphi_{258} :\equiv \mathrm{FLim}(g,c,\alpha) \to \mathrm{RFn}(h) \to \mathrm{PNbhd}(h,c) \to r \in \mathbb{R} \to 0 < r \to \varphi_{224}$
\item $\varphi_{259} :\equiv \mathrm{FLim}(g,c,\alpha) \to \mathrm{RFn}(h) \to \mathrm{PNbhd}(h,c) \to r \in \mathbb{R} \to 0 < r \to \varphi_{226}$
\item $\varphi_{260} :\equiv \mathrm{FLim}(g,c,\alpha) \to \mathrm{RFn}(h) \to \mathrm{PNbhd}(h,c) \to r \in \mathbb{R} \to 0 < r \to \varphi_{227}$
\item $\varphi_{261} :\equiv \varphi_{121} \land \varepsilon \in \mathbb{R} \land 0 < \varepsilon \land (d_{1} \in \mathbb{R} \land 0 < d_{1} \land \varphi_{49}) \land (d_{2} \in \mathbb{R} \land 0 < d_{2} \land \varphi_{48}) \land \neg \exists \delta\,(\delta \in \mathbb{R} \land 0 < \delta \land \varphi_{47})$
\item $\varphi_{262} :\equiv \forall x \in \mathrm{dom}(h)\,(0 < \lvert x - c \rvert \to \lvert x - c \rvert < r \to \varphi_{21}) \to \varepsilon \in \mathbb{R} \to \varphi_{203}$
\item $\varphi_{263} :\equiv \mathrm{FLim}(g,c,\alpha) \to \mathrm{RFn}(h) \to \mathrm{PNbhd}(h,c) \to r \in \mathbb{R} \to 0 < r \to \varphi_{228}$
\item $\varphi_{264} :\equiv \forall x \in \mathrm{dom}(h)\,(0 < \lvert x - c \rvert \to \lvert x - c \rvert < r \to \varphi_{21}) \to \varepsilon \in \mathbb{R} \to 0 < \varepsilon \to d_{1} \in \mathbb{R} \land 0 < d_{1} \land \varphi_{49} \to \varphi_{153}$
\item $\varphi_{265} :\equiv \mathrm{FLim}(g,c,\alpha) \to \mathrm{RFn}(h) \to \mathrm{PNbhd}(h,c) \to r \in \mathbb{R} \to 0 < r \to \varphi_{231}$
\item $\varphi_{266} :\equiv \mathrm{FLim}(g,c,\alpha) \to \mathrm{RFn}(h) \to \mathrm{PNbhd}(h,c) \to r \in \mathbb{R} \to 0 < r \to \varphi_{232}$
\item $\varphi_{267} :\equiv \mathrm{FLim}(g,c,\alpha) \to \mathrm{RFn}(h) \to \mathrm{PNbhd}(h,c) \to r \in \mathbb{R} \to 0 < r \to \varphi_{233}$
\item $\varphi_{268} :\equiv \mathrm{FLim}(g,c,\alpha) \to \mathrm{RFn}(h) \to \mathrm{PNbhd}(h,c) \to r \in \mathbb{R} \to 0 < r \to \varphi_{234}$
\item $\varphi_{269} :\equiv \forall x \in \mathrm{dom}(h)\,(0 < \lvert x - c \rvert \to \lvert x - c \rvert < r \to \varphi_{21}) \to \varepsilon \in \mathbb{R} \to 0 < \varepsilon \to d_{1} \in \mathbb{R} \land 0 < d_{1} \land \varphi_{49} \to \varphi_{158}$
\item $\varphi_{270} :\equiv \mathrm{FLim}(g,c,\alpha) \to \mathrm{RFn}(h) \to \mathrm{PNbhd}(h,c) \to r \in \mathbb{R} \to 0 < r \to \varphi_{235}$
\item $\varphi_{271} :\equiv \mathrm{FLim}(g,c,\alpha) \to \mathrm{RFn}(h) \to \mathrm{PNbhd}(h,c) \to r \in \mathbb{R} \to 0 < r \to \varphi_{236}$
\item $\varphi_{272} :\equiv \mathrm{FLim}(g,c,\alpha) \to \mathrm{RFn}(h) \to \mathrm{PNbhd}(h,c) \to r \in \mathbb{R} \to 0 < r \to \varphi_{237}$
\item $\varphi_{273} :\equiv \mathrm{FLim}(g,c,\alpha) \to \mathrm{RFn}(h) \to \mathrm{PNbhd}(h,c) \to r \in \mathbb{R} \to 0 < r \to \varphi_{239}$
\item $\varphi_{274} :\equiv \mathrm{FLim}(g,c,\alpha) \to \mathrm{RFn}(h) \to \mathrm{PNbhd}(h,c) \to r \in \mathbb{R} \to 0 < r \to \varphi_{240}$
\item $\varphi_{275} :\equiv \mathrm{FLim}(g,c,\alpha) \to \mathrm{RFn}(h) \to \mathrm{PNbhd}(h,c) \to r \in \mathbb{R} \to 0 < r \to \varphi_{241}$
\item $\varphi_{276} :\equiv \mathrm{FLim}(g,c,\alpha) \to \mathrm{RFn}(h) \to \mathrm{PNbhd}(h,c) \to r \in \mathbb{R} \to 0 < r \to \varphi_{242}$
\item $\varphi_{277} :\equiv \mathrm{FLim}(g,c,\alpha) \to \mathrm{RFn}(h) \to \mathrm{PNbhd}(h,c) \to r \in \mathbb{R} \to 0 < r \to \varphi_{243}$
\item $\varphi_{278} :\equiv \mathrm{FLim}(g,c,\alpha) \to \mathrm{RFn}(h) \to \mathrm{PNbhd}(h,c) \to r \in \mathbb{R} \to 0 < r \to \varphi_{244}$
\item $\varphi_{279} :\equiv \mathrm{FLim}(g,c,\alpha) \to \mathrm{RFn}(h) \to \mathrm{PNbhd}(h,c) \to r \in \mathbb{R} \to 0 < r \to \varphi_{245}$
\item $\varphi_{280} :\equiv \mathrm{FLim}(g,c,\alpha) \to \mathrm{RFn}(h) \to \mathrm{PNbhd}(h,c) \to r \in \mathbb{R} \to 0 < r \to \varphi_{246}$
\item $\varphi_{281} :\equiv \mathrm{FLim}(g,c,\alpha) \to \mathrm{RFn}(h) \to \mathrm{PNbhd}(h,c) \to r \in \mathbb{R} \to 0 < r \to \varphi_{247}$
\item $\varphi_{282} :\equiv \mathrm{FLim}(g,c,\alpha) \to \mathrm{RFn}(h) \to \mathrm{PNbhd}(h,c) \to r \in \mathbb{R} \to 0 < r \to \varphi_{248}$
\item $\varphi_{283} :\equiv \mathrm{FLim}(g,c,\alpha) \to \mathrm{RFn}(h) \to \mathrm{PNbhd}(h,c) \to r \in \mathbb{R} \to 0 < r \to \varphi_{249}$
\item $\varphi_{284} :\equiv \mathrm{FLim}(g,c,\alpha) \to \mathrm{RFn}(h) \to \mathrm{PNbhd}(h,c) \to r \in \mathbb{R} \to 0 < r \to \varphi_{250}$
\item $\varphi_{285} :\equiv \mathrm{FLim}(g,c,\alpha) \to \mathrm{RFn}(h) \to \mathrm{PNbhd}(h,c) \to r \in \mathbb{R} \to 0 < r \to \varphi_{251}$
\item $\varphi_{286} :\equiv \mathrm{FLim}(g,c,\alpha) \to \mathrm{RFn}(h) \to \mathrm{PNbhd}(h,c) \to r \in \mathbb{R} \to 0 < r \to \varphi_{252}$
\item $\varphi_{287} :\equiv \mathrm{FLim}(g,c,\alpha) \to \mathrm{RFn}(h) \to \mathrm{PNbhd}(h,c) \to r \in \mathbb{R} \to 0 < r \to \varphi_{253}$
\item $\varphi_{288} :\equiv \mathrm{FLim}(g,c,\alpha) \to \mathrm{RFn}(h) \to \mathrm{PNbhd}(h,c) \to r \in \mathbb{R} \to 0 < r \to \varphi_{254}$
\item $\varphi_{289} :\equiv \mathrm{FLim}(g,c,\alpha) \to \mathrm{RFn}(h) \to \mathrm{PNbhd}(h,c) \to r \in \mathbb{R} \to 0 < r \to \varphi_{255}$
\item $\varphi_{290} :\equiv \mathrm{FLim}(g,c,\alpha) \to \mathrm{RFn}(h) \to \mathrm{PNbhd}(h,c) \to r \in \mathbb{R} \to 0 < r \to \varphi_{256}$
\item $\varphi_{291} :\equiv \varphi_{121} \land \varepsilon \in \mathbb{R} \land 0 < \varepsilon \land (d_{1} \in \mathbb{R} \land 0 < d_{1} \land \varphi_{49}) \land (d_{2} \in \mathbb{R} \land 0 < d_{2} \land \varphi_{48}) \land (d_{3} \in \mathbb{R} \land 0 < d_{3} \land d_{3} \leq d_{1} \land d_{3} \leq d_{2}) \land \neg \exists \delta\,(\delta \in \mathbb{R} \land 0 < \delta \land \varphi_{47})$
\item $\varphi_{292} :\equiv \forall d_{3}\,(\varphi_{261} \to \neg (d_{3} \in \mathbb{R} \land 0 < d_{3} \land d_{3} \leq d_{1} \land d_{3} \leq d_{2}))$
\item $\varphi_{293} :\equiv \forall d_{3}\,(\varphi_{261}) \to \forall d_{3}\,(\neg (d_{3} \in \mathbb{R} \land 0 < d_{3} \land d_{3} \leq d_{1} \land d_{3} \leq d_{2}))$
\item $\varphi_{294} :\equiv \mathrm{FLim}(g,c,\alpha) \to \mathrm{RFn}(h) \to \mathrm{PNbhd}(h,c) \to r \in \mathbb{R} \to 0 < r \to \varphi_{262}$
\item $\varphi_{295} :\equiv \mathrm{FLim}(g,c,\alpha) \to \mathrm{RFn}(h) \to \mathrm{PNbhd}(h,c) \to r \in \mathbb{R} \to 0 < r \to \varphi_{264}$
\item $\varphi_{296} :\equiv (\mathrm{FLim}(f,c,\alpha) \to \varphi_{4}) \to (\mathrm{FLim}(f,c,\alpha) \to \varphi_{166}) \to \mathrm{FLim}(f,c,\alpha) \to \varphi_{122}$
\item $\varphi_{297} :\equiv \mathrm{FLim}(g,c,\alpha) \to \mathrm{RFn}(h) \to \mathrm{PNbhd}(h,c) \to r \in \mathbb{R} \to 0 < r \to \varphi_{269}$
\item $\varphi_{298} :\equiv (\mathrm{FLim}(f,c,\alpha) \to \varphi_{164}) \to \varphi_{121} \to \varepsilon \in \mathbb{R} \to 0 < \varepsilon \to \exists \delta\,(\delta \in \mathbb{R} \land 0 < \delta \land \varphi_{47})$
\item $\varphi_{299} :\equiv \forall \varepsilon\,(\varphi_{121} \to \varepsilon \in \mathbb{R} \to 0 < \varepsilon \to \exists \delta\,(\delta \in \mathbb{R} \land 0 < \delta \land \varphi_{47})) \to \varphi_{173}$
\item $\varphi_{300} :\equiv (\mathrm{FLim}(f,c,\alpha) \to \varphi_{124}) \to (\mathrm{FLim}(f,c,\alpha) \to \varphi_{127}) \to \mathrm{FLim}(f,c,\alpha) \to \varphi_{163}$
\item $\varphi_{301} :\equiv (\mathrm{FLim}(f,c,\alpha) \to \varphi_{157}) \to \varphi_{143} \to \mathrm{FLim}(f,c,\alpha) \to \varphi_{154}$
\item $\varphi_{302} :\equiv (\mathrm{FLim}(f,c,\alpha) \to \varphi_{124}) \to (\mathrm{FLim}(f,c,\alpha) \to \varphi_{127}) \to \mathrm{FLim}(f,c,\alpha) \to \varphi_{207}$
\item $\varphi_{303} :\equiv \forall d\,(\varphi_{37}) \to (\forall d\,(\varphi_{37}) \to \varphi_{46}) \to \varphi_{54} \to \mathrm{FLim}(f,c,\alpha) \to \varphi_{282}$
\item $\varphi_{304} :\equiv (\mathrm{FLim}(f,c,\alpha) \to \varepsilon \in \mathbb{R} \to 0 < \varepsilon \to \varphi_{76}) \to (\mathrm{FLim}(f,c,\alpha) \to \mathrm{FLim}(f,c,\alpha)) \to \varphi_{300}$
\item $\varphi_{305} :\equiv (\mathrm{FLim}(f,c,\alpha) \to \varphi_{118}) \to (\mathrm{FLim}(f,c,\alpha) \to \varphi_{119}) \to \varphi_{296}$
\item $\varphi_{306} :\equiv (\mathrm{FLim}(f,c,\alpha) \to \varphi_{223}) \to (\mathrm{FLim}(f,c,\alpha) \to \varphi_{16}) \to \mathrm{FLim}(f,c,\alpha) \to \varphi_{229}$
\item $\varphi_{307} :\equiv (\mathrm{FLim}(f,c,\alpha) \to \varphi_{259}) \to \varphi_{216} \to \neg (d_{2} \in \mathbb{R} \land 0 < d_{2} \land \varphi_{48})$
\item $\varphi_{308} :\equiv \forall d_{2}\,(\varphi_{216} \to \neg (d_{2} \in \mathbb{R} \land 0 < d_{2} \land \varphi_{48})) \to \forall d_{2}\,(\varphi_{216}) \to \forall d_{2}\,(\neg (d_{2} \in \mathbb{R} \land 0 < d_{2} \land \varphi_{48}))$
\item $\varphi_{309} :\equiv \varphi_{2} \to (\mathrm{FLim}(f,c,\alpha) \to \varphi_{257}) \to \mathrm{FLim}(f,c,\alpha) \to \varphi_{279}$
\item $\varphi_{310} :\equiv (\mathrm{FLim}(f,c,\alpha) \to \varphi_{282}) \to \varphi_{261} \to \neg (d_{3} \in \mathbb{R} \land 0 < d_{3} \land d_{3} \leq d_{1} \land d_{3} \leq d_{2})$
\item $\varphi_{311} :\equiv \forall \delta\,(\forall d_{1}\,(\neg (d_{1} \in \mathbb{R} \land 0 < d_{1} \land \varphi_{49})) \to \varphi_{73}) \to \varphi_{219} \to \varphi_{131} \to \mathrm{FLim}(f,c,\alpha) \to \varphi_{164}$
\item $\varphi_{312} :\equiv (\mathrm{FLim}(f,c,\alpha) \to \varphi_{295}) \to \varphi_{291} \to \neg (\delta \in \mathbb{R} \land 0 < \delta \land \delta \leq d_{3} \land \delta \leq r)$
\item $\varphi_{313} :\equiv \forall \delta\,(\varphi_{291} \to \neg (\delta \in \mathbb{R} \land 0 < \delta \land \delta \leq d_{3} \land \delta \leq r)) \to \forall \delta\,(\varphi_{291}) \to \forall \delta\,(\neg (\delta \in \mathbb{R} \land 0 < \delta \land \delta \leq d_{3} \land \delta \leq r))$
\item $\varphi_{314} :\equiv \forall \delta\,(\forall d_{2}\,(\neg (d_{2} \in \mathbb{R} \land 0 < d_{2} \land \varphi_{48})) \to \varphi_{74}) \to \varphi_{220} \to \varphi_{132} \to \mathrm{FLim}(f,c,\alpha) \to \varphi_{208}$
\item $\varphi_{315} :\equiv (\mathrm{FLim}(f,c,\alpha) \to \varphi_{288}) \to (\varphi_{48} \to \varphi_{41}) \to \mathrm{FLim}(f,c,\alpha) \to \varphi_{283}$
\item $\varphi_{316} :\equiv (\mathrm{FLim}(f,c,\alpha) \to \varphi_{289}) \to (\varphi_{49} \to \varphi_{43}) \to \mathrm{FLim}(f,c,\alpha) \to \varphi_{285}$
\item $\varphi_{317} :\equiv (\mathrm{FLim}(f,c,\alpha) \to \varphi_{211}) \to (\mathrm{FLim}(f,c,\alpha) \to \varphi_{210}) \to \mathrm{FLim}(f,c,\alpha) \to \varphi_{221}$
\item $\varphi_{318} :\equiv (\mathrm{FLim}(f,c,\alpha) \to \varphi_{116}) \to (\mathrm{FLim}(f,c,\alpha) \to \varphi_{270}) \to \mathrm{FLim}(f,c,\alpha) \to \varphi_{280}$
\item $\varphi_{319} :\equiv (\mathrm{FLim}(f,c,\alpha) \to \varphi_{117}) \to (\mathrm{FLim}(f,c,\alpha) \to \varphi_{271}) \to \mathrm{FLim}(f,c,\alpha) \to \varphi_{281}$
\item $\varphi_{320} :\equiv (\mathrm{FLim}(f,c,\alpha) \to \varphi_{294}) \to \varphi_{129} \to \mathrm{FLim}(f,c,\alpha) \to \varphi_{295}$
\item $\varphi_{321} :\equiv (\mathrm{FLim}(f,c,\alpha) \to \varphi_{225}) \to (\mathrm{FLim}(f,c,\alpha) \to \varphi_{12}) \to \varphi_{306}$
\item $\varphi_{322} :\equiv (\mathrm{FLim}(f,c,\alpha) \to \varphi_{273}) \to (\mathrm{FLim}(f,c,\alpha) \to \varphi_{265}) \to \mathrm{FLim}(f,c,\alpha) \to \varphi_{272}$
\item $\varphi_{323} :\equiv (\mathrm{FLim}(f,c,\alpha) \to \varphi_{273}) \to (\mathrm{FLim}(f,c,\alpha) \to \varphi_{267}) \to \mathrm{FLim}(f,c,\alpha) \to \varphi_{274}$
\item $\varphi_{324} :\equiv (\mathrm{FLim}(f,c,\alpha) \to \varphi_{274}) \to (\mathrm{FLim}(f,c,\alpha) \to \varphi_{266}) \to \mathrm{FLim}(f,c,\alpha) \to \varphi_{275}$
\item $\varphi_{325} :\equiv (\mathrm{FLim}(f,c,\alpha) \to \varphi_{263}) \to (\mathrm{FLim}(f,c,\alpha) \to \varphi_{275}) \to \mathrm{FLim}(f,c,\alpha) \to \varphi_{284}$
\item $\varphi_{326} :\equiv (\mathrm{FLim}(f,c,\alpha) \to \varphi_{274}) \to (\mathrm{FLim}(f,c,\alpha) \to \varphi_{268}) \to \mathrm{FLim}(f,c,\alpha) \to \varphi_{276}$
\item $\varphi_{327} :\equiv (\mathrm{FLim}(f,c,\alpha) \to \varphi_{263}) \to (\mathrm{FLim}(f,c,\alpha) \to \varphi_{276}) \to \mathrm{FLim}(f,c,\alpha) \to \varphi_{286}$
\item $\varphi_{328} :\equiv \forall \varepsilon\,(\varphi_{121} \to \varepsilon \in \mathbb{R} \to 0 < \varepsilon \to \exists \delta\,(\delta \in \mathbb{R} \land 0 < \delta \land \varphi_{47})) \to \varphi_{299} \to (\varphi_{121} \to \forall \varepsilon\,(\varphi_{121})) \to \mathrm{FLim}(f,c,\alpha) \to \varphi_{166}$
\item $\varphi_{329} :\equiv (\mathrm{FLim}(f,c,\alpha) \to \varphi_{260}) \to (\mathrm{FLim}(f,c,\alpha) \to \varphi_{290}) \to \mathrm{FLim}(f,c,\alpha) \to \varphi_{294}$
\item $\varphi_{330} :\equiv (\mathrm{FLim}(f,c,\alpha) \to \varphi_{286}) \to (\mathrm{FLim}(f,c,\alpha) \to \varphi_{284}) \to \mathrm{FLim}(f,c,\alpha) \to \varphi_{287}$
\item $\varphi_{331} :\equiv (\mathrm{FLim}(f,c,\alpha) \to \varphi_{263}) \to (\mathrm{FLim}(f,c,\alpha) \to \varphi_{272}) \to \mathrm{FLim}(f,c,\alpha) \to \varphi_{297}$
\item $\varphi_{332} :\equiv (\varphi_{121} \land \varepsilon \in \mathbb{R} \land 0 < \varepsilon \land \neg \exists \delta\,(\delta \in \mathbb{R} \land 0 < \delta \land \varphi_{47}) \to \varphi_{172}) \to \varphi_{311}$
\item $\varphi_{333} :\equiv (\varphi_{261} \to \forall d_{3}\,(\varphi_{261})) \to \forall d\,(\varphi_{39}) \to \varphi_{113} \to \varphi_{59} \to \mathrm{FLim}(f,c,\alpha) \to \varphi_{259}$
\item $\varphi_{334} :\equiv (\mathrm{FLim}(f,c,\alpha) \to \varphi_{213}) \to (\mathrm{FLim}(f,c,\alpha) \to \varphi_{212}) \to \varphi_{317}$
\item $\varphi_{335} :\equiv (\mathrm{FLim}(f,c,\alpha) \to \varphi_{154}) \to (\mathrm{FLim}(f,c,\alpha) \to \varphi_{257}) \to \varphi_{331}$
\item $\varphi_{336} :\equiv (\forall d_{1}\,(\varphi_{217}) \to \varphi_{172} \to \forall d_{1}\,(\neg (d_{1} \in \mathbb{R} \land 0 < d_{1} \land \varphi_{49}))) \to \varphi_{332}$
\item $\varphi_{337} :\equiv (\forall x\,(\varphi_{238} \to \varphi_{45}) \to \forall x\,(\varphi_{238}) \to \varphi_{47}) \to (\varphi_{238} \to \forall x\,(\varphi_{238})) \to \mathrm{FLim}(f,c,\alpha) \to \varphi_{290}$
\item $\varphi_{338} :\equiv (\mathrm{FLim}(f,c,\alpha) \to \varphi_{278}) \to (\mathrm{FLim}(f,c,\alpha) \to \varphi_{277}) \to \varphi_{330}$
\item $\varphi_{339} :\equiv (\mathrm{FLim}(f,c,\alpha) \to \varphi_{283}) \to (\mathrm{FLim}(f,c,\alpha) \to \varphi_{297}) \to \varphi_{325}$
\item $\varphi_{340} :\equiv (\mathrm{FLim}(f,c,\alpha) \to \varphi_{285}) \to (\mathrm{FLim}(f,c,\alpha) \to \varphi_{297}) \to \varphi_{327}$
\item $\varphi_{341} :\equiv (\mathrm{FLim}(f,c,\alpha) \to \varphi_{119}) \to (\mathrm{FLim}(f,c,\alpha) \to \varphi_{124}) \to \varphi_{338}$
\item $\varphi_{342} :\equiv \forall d_{2}\,(\varphi_{216} \to \neg (d_{2} \in \mathbb{R} \land 0 < d_{2} \land \varphi_{48})) \to \varphi_{308} \to (\varphi_{216} \to \forall d_{2}\,(\varphi_{216})) \to \varphi_{314}$
\item $\varphi_{343} :\equiv \forall \delta\,(\varphi_{291} \to \neg (\delta \in \mathbb{R} \land 0 < \delta \land \delta \leq d_{3} \land \delta \leq r)) \to \varphi_{313} \to (\varphi_{291} \to \forall \delta\,(\varphi_{291})) \to \varphi_{303}$
\item $\varphi_{344} :\equiv (\mathrm{FLim}(f,c,\alpha) \to \varphi_{280}) \to (\mathrm{FLim}(f,c,\alpha) \to \varphi_{279}) \to \varphi_{341}$
\item $\varphi_{345} :\equiv (f(x) \in \mathbb{R} \to g(x) \in \mathbb{R} \to h(x) \in \mathbb{R} \to \varphi_{98}) \to (\mathrm{FLim}(f,c,\alpha) \to \varphi_{281}) \to \varphi_{344}$
\end{abbrevlist}
\end{fullproof}

\section{第4章 1節(続き): 関数の極限と数列の極限(定理1.7 ⇐), 極限の性質(定理1.8, 1.9), Cauchy の判定法(1.13), 片側極限}
% 自動生成 (prover/build.js): AddSub  ノード数 126
\begin{fullproof}{[E38]}{$\forall x\,\forall y\,(x \in \mathbb{R} \to y \in \mathbb{R} \to x + y - y = x)$}{126}\phantomsection\label{tree:AddSub}
\noindent$\Pi_{1}$:
\begin{ptree}
\AXC{$y = z \to x = y \to x = z$\lbl{A6}}
\AXC{$(y = z \to x = y \to x = z) \to x = y \to y = z \to x = z$\lbl{A1}}
\BIC{$x = y \to y = z \to x = z$}
\UIC{$\forall z\,(x = y \to y = z \to x = z)$}
\UIC{$\forall y\,\forall z\,(x = y \to y = z \to x = z)$}
\UIC{$\forall x\,\forall y\,\forall z\,(x = y \to y = z \to x = z)$}
\AXC{$\varphi_{4}$\lbl{A2}}
\BIC{$\forall y\,\forall z\,(v_{763} = y \to y = z \to v_{763} = z)$}
\AXC{$\varphi_{13}$\lbl{A2}}
\BIC{$\forall z\,(v_{763} = v_{764} \to v_{764} = z \to v_{763} = z)$}
\AXC{$\varphi_{18}$\lbl{A2}}
\BIC{$v_{763} = v_{764} \to v_{764} = v_{765} \to v_{763} = v_{765}$}
\UIC{$\forall v_{765}\,(v_{763} = v_{764} \to v_{764} = v_{765} \to v_{763} = v_{765})$}
\UIC{$\varphi_{1}$}
\UIC{$\forall v_{763}\,(\varphi_{1})$}
\end{ptree}
\noindent$\Pi_{2}$:
\begin{ptree}
\AXC{$\stackrel{\Pi_{1}}{\vdots}$}\noLine
\UIC{$\forall v_{763}\,(\varphi_{1})$}
\AXC{$\forall v_{763}\,(\varphi_{1}) \to \varphi_{10}$\lbl{A2}}
\BIC{$\varphi_{10}$}
\AXC{$\varphi_{10} \to \varphi_{8}$\lbl{A2}}
\BIC{$\varphi_{8}$}
\AXC{$\varphi_{8} \to x + y - y = x + 0 \to x + 0 = x \to x + y - y = x$\lbl{A2}}
\BIC{$x + y - y = x + 0 \to x + 0 = x \to x + y - y = x$}
\end{ptree}
\noindent$\Pi_{3}$:
\begin{ptree}
\AXC{$\forall x \in \mathbb{R}\,(x + 0 = x)$\lbl{(R4)}}
\AXC{$\forall x \in \mathbb{R}\,(x + 0 = x) \to x \in \mathbb{R} \to x + 0 = x$\lbl{A2}}
\BIC{$x \in \mathbb{R} \to x + 0 = x$}
\AXC{$(x \in \mathbb{R} \to x + 0 = x) \to \varphi_{22}$\lbl{A1}}
\BIC{$\varphi_{22}$}
\end{ptree}
\noindent$\Pi_{4}$:
\begin{ptree}
\AXC{$y = z \to x = y \to x = z$\lbl{A6}}
\AXC{$(y = z \to x = y \to x = z) \to x = y \to y = z \to x = z$\lbl{A1}}
\BIC{$x = y \to y = z \to x = z$}
\UIC{$\forall z\,(x = y \to y = z \to x = z)$}
\UIC{$\forall y\,\forall z\,(x = y \to y = z \to x = z)$}
\UIC{$\forall x\,\forall y\,\forall z\,(x = y \to y = z \to x = z)$}
\AXC{$\varphi_{5}$\lbl{A2}}
\BIC{$\forall y\,\forall z\,(v_{760} = y \to y = z \to v_{760} = z)$}
\AXC{$\varphi_{14}$\lbl{A2}}
\BIC{$\forall z\,(v_{760} = v_{761} \to v_{761} = z \to v_{760} = z)$}
\AXC{$\varphi_{19}$\lbl{A2}}
\BIC{$v_{760} = v_{761} \to v_{761} = v_{762} \to v_{760} = v_{762}$}
\UIC{$\forall v_{762}\,(v_{760} = v_{761} \to v_{761} = v_{762} \to v_{760} = v_{762})$}
\UIC{$\varphi_{2}$}
\UIC{$\forall v_{760}\,(\varphi_{2})$}
\end{ptree}
\noindent$\Pi_{5}$:
\begin{ptree}
\AXC{$z = w \to x + (y + -y) = x + z \to x + (y + -y) = x + w$\lbl{A6}}
\UIC{$\varphi_{9}$}
\UIC{$\forall z\,(\varphi_{9})$}
\AXC{$\forall z\,(\varphi_{9}) \to \varphi_{28}$\lbl{A2}}
\BIC{$\varphi_{28}$}
\AXC{$\varphi_{28} \to \varphi_{24}$\lbl{A2}}
\BIC{$\varphi_{24}$}
\end{ptree}
\noindent$\Pi_{6}$:
\begin{ptree}
\AXC{$\forall x \in \mathbb{R}\,(x + -x = 0)$\lbl{(R5)}}
\AXC{$\forall x \in \mathbb{R}\,(x + -x = 0) \to y \in \mathbb{R} \to y + -y = 0$\lbl{A2}}
\BIC{$y \in \mathbb{R} \to y + -y = 0$}
\AXC{$(y \in \mathbb{R} \to y + -y = 0) \to \varphi_{29}$\lbl{A1}}
\BIC{$\varphi_{29}$}
\end{ptree}
\noindent$\Pi_{7}$:
\begin{ptree}
\AXC{$x \in \mathbb{R} \to y \in \mathbb{R} \to y \in \mathbb{R}$\lbl{A1}}
\AXC{$\stackrel{\Pi_{6}}{\vdots}$}\noLine
\UIC{$\varphi_{29}$}
\BIC{$x \in \mathbb{R} \to y \in \mathbb{R} \to y + -y = 0$}
\AXC{$(x \in \mathbb{R} \to y \in \mathbb{R} \to y + -y = 0) \to \varphi_{35}$\lbl{A1}}
\BIC{$\varphi_{35}$}
\end{ptree}
\noindent$\Pi_{8}$:
\begin{ptree}
\AXC{$y = z \to x = y \to x = z$\lbl{A6}}
\AXC{$(y = z \to x = y \to x = z) \to x = y \to y = z \to x = z$\lbl{A1}}
\BIC{$x = y \to y = z \to x = z$}
\UIC{$\forall z\,(x = y \to y = z \to x = z)$}
\UIC{$\forall y\,\forall z\,(x = y \to y = z \to x = z)$}
\UIC{$\forall x\,\forall y\,\forall z\,(x = y \to y = z \to x = z)$}
\AXC{$\varphi_{6}$\lbl{A2}}
\BIC{$\forall y\,\forall z\,(v_{757} = y \to y = z \to v_{757} = z)$}
\AXC{$\varphi_{15}$\lbl{A2}}
\BIC{$\forall z\,(v_{757} = v_{758} \to v_{758} = z \to v_{757} = z)$}
\AXC{$\varphi_{20}$\lbl{A2}}
\BIC{$v_{757} = v_{758} \to v_{758} = v_{759} \to v_{757} = v_{759}$}
\UIC{$\forall v_{759}\,(v_{757} = v_{758} \to v_{758} = v_{759} \to v_{757} = v_{759})$}
\UIC{$\varphi_{3}$}
\UIC{$\forall v_{757}\,(\varphi_{3})$}
\end{ptree}
\noindent$\Pi_{9}$:
\begin{ptree}
\AXC{$\stackrel{\Pi_{8}}{\vdots}$}\noLine
\UIC{$\forall v_{757}\,(\varphi_{3})$}
\AXC{$\forall v_{757}\,(\varphi_{3}) \to \varphi_{12}$\lbl{A2}}
\BIC{$\varphi_{12}$}
\AXC{$\varphi_{12} \to \varphi_{26}$\lbl{A2}}
\BIC{$\varphi_{26}$}
\AXC{$\varphi_{26} \to \varphi_{33}$\lbl{A2}}
\BIC{$\varphi_{33}$}
\end{ptree}
\noindent$\Pi_{10}$:
\begin{ptree}
\AXC{$\forall x \in \mathbb{R}\,(-x \in \mathbb{R})$\lbl{(R1)}}
\AXC{$\forall x \in \mathbb{R}\,(-x \in \mathbb{R}) \to y \in \mathbb{R} \to -y \in \mathbb{R}$\lbl{A2}}
\BIC{$y \in \mathbb{R} \to -y \in \mathbb{R}$}
\AXC{$(y \in \mathbb{R} \to -y \in \mathbb{R}) \to \varphi_{30}$\lbl{A1}}
\BIC{$\varphi_{30}$}
\end{ptree}
\noindent$\Pi_{11}$:
\begin{ptree}
\AXC{$\forall x\,\forall y\,(\varphi_{17})$\lbl{(R2)}}
\AXC{$\forall x\,\forall y\,(\varphi_{17}) \to \forall y\,(\varphi_{17})$\lbl{A2}}
\BIC{$\forall y\,(\varphi_{17})$}
\AXC{$\forall y\,(\varphi_{17}) \to \varphi_{17}$\lbl{A2}}
\BIC{$\varphi_{17}$}
\AXC{$\varphi_{17} \to \varphi_{27}$\lbl{A2}}
\BIC{$\varphi_{27}$}
\AXC{$\varphi_{27} \to \varphi_{38}$\lbl{A1}}
\BIC{$\varphi_{38}$}
\end{ptree}
\noindent$\Pi_{12}$:
\begin{ptree}
\AXC{$x \in \mathbb{R} \to y \in \mathbb{R} \to y \in \mathbb{R}$\lbl{A1}}
\AXC{$\stackrel{\Pi_{10}}{\vdots}$}\noLine
\UIC{$\varphi_{30}$}
\BIC{$x \in \mathbb{R} \to y \in \mathbb{R} \to -y \in \mathbb{R}$}
\AXC{$x \in \mathbb{R} \to y \in \mathbb{R} \to y \in \mathbb{R}$\lbl{A1}}
\AXC{$x \in \mathbb{R} \to y \in \mathbb{R} \to x \in \mathbb{R}$\lbl{A1}}
\AXC{$\stackrel{\Pi_{11}}{\vdots}$}\noLine
\UIC{$\varphi_{38}$}
\BIC{$(x \in \mathbb{R} \to y \in \mathbb{R} \to y \in \mathbb{R}) \to \varphi_{34}$}
\BIC{$\varphi_{34}$}
\BIC{$x \in \mathbb{R} \to y \in \mathbb{R} \to x + y + -y = x + (y + -y)$}
\end{ptree}
\noindent$\Pi_{13}$:
\begin{ptree}
\AXC{$\forall x\,\forall y\,(x - y = x + -y)$\lbl{Def$-$}}
\AXC{$\varphi_{7}$\lbl{A2}}
\BIC{$\forall y\,(v_{755} - y = v_{755} + -y)$}
\AXC{$\varphi_{16}$\lbl{A2}}
\BIC{$v_{755} - v_{756} = v_{755} + -v_{756}$}
\UIC{$\forall v_{756}\,(v_{755} - v_{756} = v_{755} + -v_{756})$}
\UIC{$\forall v_{755}\,\forall v_{756}\,(v_{755} - v_{756} = v_{755} + -v_{756})$}
\AXC{$\varphi_{31}$\lbl{A2}}
\BIC{$\forall v_{756}\,(x + y - v_{756} = x + y + -v_{756})$}
\AXC{$\varphi_{21}$\lbl{A2}}
\BIC{$x + y - y = x + y + -y$}
\AXC{$x + y - y = x + y + -y \to \varphi_{39}$\lbl{A1}}
\BIC{$\varphi_{39}$}
\end{ptree}
\noindent$\Pi_{14}$:
\begin{ptree}
\AXC{$\stackrel{\Pi_{9}}{\vdots}$}\noLine
\UIC{$\varphi_{33}$}
\AXC{$\stackrel{\Pi_{12}}{\vdots}$}\noLine
\UIC{$x \in \mathbb{R} \to y \in \mathbb{R} \to x + y + -y = x + (y + -y)$}
\AXC{$\stackrel{\Pi_{13}}{\vdots}$}\noLine
\UIC{$\varphi_{39}$}
\BIC{$\varphi_{33} \to x \in \mathbb{R} \to y \in \mathbb{R} \to x + y - y = x + (y + -y)$}
\BIC{$x \in \mathbb{R} \to y \in \mathbb{R} \to x + y - y = x + (y + -y)$}
\AXC{$(x \in \mathbb{R} \to y \in \mathbb{R} \to x + y - y = x + (y + -y)) \to \varphi_{37}$\lbl{A1}}
\BIC{$\varphi_{37}$}
\end{ptree}
\noindent$\Pi_{15}$:
\begin{ptree}
\AXC{$\stackrel{\Pi_{5}}{\vdots}$}\noLine
\UIC{$\varphi_{24}$}
\AXC{$x = x$\lbl{A5}}
\UIC{$\forall x\,(x = x)$}
\AXC{$\forall x\,(x = x) \to x + (y + -y) = x + (y + -y)$\lbl{A2}}
\BIC{$x + (y + -y) = x + (y + -y)$}
\AXC{$\stackrel{\Pi_{7}}{\vdots}$}\noLine
\UIC{$\varphi_{35}$}
\BIC{$\varphi_{24} \to x \in \mathbb{R} \to y \in \mathbb{R} \to x + (y + -y) = x + 0$}
\BIC{$x \in \mathbb{R} \to y \in \mathbb{R} \to x + (y + -y) = x + 0$}
\AXC{$\stackrel{\Pi_{14}}{\vdots}$}\noLine
\UIC{$\varphi_{37}$}
\BIC{$\varphi_{23} \to x \in \mathbb{R} \to y \in \mathbb{R} \to x + y - y = x + 0$}
\end{ptree}
\noindent$\Pi_{16}$:
\begin{ptree}
\AXC{$\stackrel{\Pi_{4}}{\vdots}$}\noLine
\UIC{$\forall v_{760}\,(\varphi_{2})$}
\AXC{$\forall v_{760}\,(\varphi_{2}) \to \varphi_{11}$\lbl{A2}}
\BIC{$\varphi_{11}$}
\AXC{$\varphi_{11} \to \varphi_{25}$\lbl{A2}}
\BIC{$\varphi_{25}$}
\AXC{$\varphi_{25} \to \varphi_{23}$\lbl{A2}}
\BIC{$\varphi_{23}$}
\AXC{$\stackrel{\Pi_{15}}{\vdots}$}\noLine
\UIC{$\varphi_{23} \to x \in \mathbb{R} \to y \in \mathbb{R} \to x + y - y = x + 0$}
\BIC{$x \in \mathbb{R} \to y \in \mathbb{R} \to x + y - y = x + 0$}
\end{ptree}
\noindent$\Pi_{17}$:
\begin{ptree}
\AXC{$x \in \mathbb{R} \to y \in \mathbb{R} \to x \in \mathbb{R}$\lbl{A1}}
\AXC{$\stackrel{\Pi_{3}}{\vdots}$}\noLine
\UIC{$\varphi_{22}$}
\BIC{$x \in \mathbb{R} \to y \in \mathbb{R} \to x + 0 = x$}
\AXC{$\stackrel{\Pi_{16}}{\vdots}$}\noLine
\UIC{$x \in \mathbb{R} \to y \in \mathbb{R} \to x + y - y = x + 0$}
\AXC{$\varphi_{36}$\lbl{A1}}
\BIC{$(x \in \mathbb{R} \to y \in \mathbb{R} \to x + 0 = x) \to \varphi_{32}$}
\BIC{$\varphi_{32}$}
\end{ptree}
\noindent 主木:
\begin{ptree}
\AXC{$\stackrel{\Pi_{2}}{\vdots}$}\noLine
\UIC{$x + y - y = x + 0 \to x + 0 = x \to x + y - y = x$}
\AXC{$\stackrel{\Pi_{17}}{\vdots}$}\noLine
\UIC{$\varphi_{32}$}
\BIC{$x \in \mathbb{R} \to y \in \mathbb{R} \to x + y - y = x$}
\UIC{$\forall y\,(x \in \mathbb{R} \to y \in \mathbb{R} \to x + y - y = x)$}
\UIC{$\forall x\,\forall y\,(x \in \mathbb{R} \to y \in \mathbb{R} \to x + y - y = x)$}
\end{ptree}
\begin{abbrevlist}
\item $\varphi_{1} :\equiv \forall v_{764}\,\forall v_{765}\,(v_{763} = v_{764} \to v_{764} = v_{765} \to v_{763} = v_{765})$
\item $\varphi_{2} :\equiv \forall v_{761}\,\forall v_{762}\,(v_{760} = v_{761} \to v_{761} = v_{762} \to v_{760} = v_{762})$
\item $\varphi_{3} :\equiv \forall v_{758}\,\forall v_{759}\,(v_{757} = v_{758} \to v_{758} = v_{759} \to v_{757} = v_{759})$
\item $\varphi_{4} :\equiv \forall x\,\forall y\,\forall z\,(x = y \to y = z \to x = z) \to \forall y\,\forall z\,(v_{763} = y \to y = z \to v_{763} = z)$
\item $\varphi_{5} :\equiv \forall x\,\forall y\,\forall z\,(x = y \to y = z \to x = z) \to \forall y\,\forall z\,(v_{760} = y \to y = z \to v_{760} = z)$
\item $\varphi_{6} :\equiv \forall x\,\forall y\,\forall z\,(x = y \to y = z \to x = z) \to \forall y\,\forall z\,(v_{757} = y \to y = z \to v_{757} = z)$
\item $\varphi_{7} :\equiv \forall x\,\forall y\,(x - y = x + -y) \to \forall y\,(v_{755} - y = v_{755} + -y)$
\item $\varphi_{8} :\equiv \forall v_{765}\,(x + y - y = x + 0 \to x + 0 = v_{765} \to x + y - y = v_{765})$
\item $\varphi_{9} :\equiv \forall w\,(z = w \to x + (y + -y) = x + z \to x + (y + -y) = x + w)$
\item $\varphi_{10} :\equiv \forall v_{764}\,\forall v_{765}\,(x + y - y = v_{764} \to v_{764} = v_{765} \to x + y - y = v_{765})$
\item $\varphi_{11} :\equiv \forall v_{761}\,\forall v_{762}\,(x + y - y = v_{761} \to v_{761} = v_{762} \to x + y - y = v_{762})$
\item $\varphi_{12} :\equiv \forall v_{758}\,\forall v_{759}\,(x + y - y = v_{758} \to v_{758} = v_{759} \to x + y - y = v_{759})$
\item $\varphi_{13} :\equiv \forall y\,\forall z\,(v_{763} = y \to y = z \to v_{763} = z) \to \forall z\,(v_{763} = v_{764} \to v_{764} = z \to v_{763} = z)$
\item $\varphi_{14} :\equiv \forall y\,\forall z\,(v_{760} = y \to y = z \to v_{760} = z) \to \forall z\,(v_{760} = v_{761} \to v_{761} = z \to v_{760} = z)$
\item $\varphi_{15} :\equiv \forall y\,\forall z\,(v_{757} = y \to y = z \to v_{757} = z) \to \forall z\,(v_{757} = v_{758} \to v_{758} = z \to v_{757} = z)$
\item $\varphi_{16} :\equiv \forall y\,(v_{755} - y = v_{755} + -y) \to v_{755} - v_{756} = v_{755} + -v_{756}$
\item $\varphi_{17} :\equiv \forall z\,(x \in \mathbb{R} \to y \in \mathbb{R} \to z \in \mathbb{R} \to x + y + z = x + (y + z))$
\item $\varphi_{18} :\equiv \forall z\,(v_{763} = v_{764} \to v_{764} = z \to v_{763} = z) \to v_{763} = v_{764} \to v_{764} = v_{765} \to v_{763} = v_{765}$
\item $\varphi_{19} :\equiv \forall z\,(v_{760} = v_{761} \to v_{761} = z \to v_{760} = z) \to v_{760} = v_{761} \to v_{761} = v_{762} \to v_{760} = v_{762}$
\item $\varphi_{20} :\equiv \forall z\,(v_{757} = v_{758} \to v_{758} = z \to v_{757} = z) \to v_{757} = v_{758} \to v_{758} = v_{759} \to v_{757} = v_{759}$
\item $\varphi_{21} :\equiv \forall v_{756}\,(x + y - v_{756} = x + y + -v_{756}) \to x + y - y = x + y + -y$
\item $\varphi_{22} :\equiv (x \in \mathbb{R} \to y \in \mathbb{R} \to x \in \mathbb{R}) \to x \in \mathbb{R} \to y \in \mathbb{R} \to x + 0 = x$
\item $\varphi_{23} :\equiv x + y - y = x + (y + -y) \to x + (y + -y) = x + 0 \to x + y - y = x + 0$
\item $\varphi_{24} :\equiv y + -y = 0 \to x + (y + -y) = x + (y + -y) \to x + (y + -y) = x + 0$
\item $\varphi_{25} :\equiv \forall v_{762}\,(x + y - y = x + (y + -y) \to x + (y + -y) = v_{762} \to x + y - y = v_{762})$
\item $\varphi_{26} :\equiv \forall v_{759}\,(x + y - y = x + y + -y \to x + y + -y = v_{759} \to x + y - y = v_{759})$
\item $\varphi_{27} :\equiv x \in \mathbb{R} \to y \in \mathbb{R} \to -y \in \mathbb{R} \to x + y + -y = x + (y + -y)$
\item $\varphi_{28} :\equiv \forall w\,(y + -y = w \to x + (y + -y) = x + (y + -y) \to x + (y + -y) = x + w)$
\item $\varphi_{29} :\equiv (x \in \mathbb{R} \to y \in \mathbb{R} \to y \in \mathbb{R}) \to x \in \mathbb{R} \to y \in \mathbb{R} \to y + -y = 0$
\item $\varphi_{30} :\equiv (x \in \mathbb{R} \to y \in \mathbb{R} \to y \in \mathbb{R}) \to x \in \mathbb{R} \to y \in \mathbb{R} \to -y \in \mathbb{R}$
\item $\varphi_{31} :\equiv \forall v_{755}\,\forall v_{756}\,(v_{755} - v_{756} = v_{755} + -v_{756}) \to \forall v_{756}\,(x + y - v_{756} = x + y + -v_{756})$
\item $\varphi_{32} :\equiv (x + y - y = x + 0 \to x + 0 = x \to x + y - y = x) \to x \in \mathbb{R} \to y \in \mathbb{R} \to x + y - y = x$
\item $\varphi_{33} :\equiv x + y - y = x + y + -y \to x + y + -y = x + (y + -y) \to x + y - y = x + (y + -y)$
\item $\varphi_{34} :\equiv (x \in \mathbb{R} \to y \in \mathbb{R} \to -y \in \mathbb{R}) \to x \in \mathbb{R} \to y \in \mathbb{R} \to x + y + -y = x + (y + -y)$
\item $\varphi_{35} :\equiv x + (y + -y) = x + (y + -y) \to \varphi_{24} \to x \in \mathbb{R} \to y \in \mathbb{R} \to x + (y + -y) = x + 0$
\item $\varphi_{36} :\equiv (x \in \mathbb{R} \to y \in \mathbb{R} \to x + y - y = x + 0) \to (x \in \mathbb{R} \to y \in \mathbb{R} \to x + 0 = x) \to \varphi_{32}$
\item $\varphi_{37} :\equiv (x \in \mathbb{R} \to y \in \mathbb{R} \to x + (y + -y) = x + 0) \to \varphi_{23} \to x \in \mathbb{R} \to y \in \mathbb{R} \to x + y - y = x + 0$
\item $\varphi_{38} :\equiv (x \in \mathbb{R} \to y \in \mathbb{R} \to x \in \mathbb{R}) \to (x \in \mathbb{R} \to y \in \mathbb{R} \to y \in \mathbb{R}) \to \varphi_{34}$
\item $\varphi_{39} :\equiv (x \in \mathbb{R} \to y \in \mathbb{R} \to x + y + -y = x + (y + -y)) \to \varphi_{33} \to x \in \mathbb{R} \to y \in \mathbb{R} \to x + y - y = x + (y + -y)$
\end{abbrevlist}
\end{fullproof}

% 自動生成 (prover/build.js): Inv1  ノード数 41
\begin{fullproof}{[E35]}{$1^{-1} = 1$}{41}\phantomsection\label{tree:Inv1}
\noindent$\Pi_{1}$:
\begin{ptree}
\AXC{$x = x$\lbl{A5}}
\AXC{$x = y \to x = x \to y = x$\lbl{A6}}
\AXC{$(x = y \to x = x \to y = x) \to x = x \to x = y \to y = x$\lbl{A1}}
\BIC{$x = x \to x = y \to y = x$}
\BIC{$x = y \to y = x$}
\UIC{$\forall y\,(x = y \to y = x)$}
\UIC{$\forall x\,\forall y\,(x = y \to y = x)$}
\end{ptree}
\noindent$\Pi_{2}$:
\begin{ptree}
\AXC{$\stackrel{\Pi_{1}}{\vdots}$}\noLine
\UIC{$\forall x\,\forall y\,(x = y \to y = x)$}
\AXC{$\forall x\,\forall y\,(x = y \to y = x) \to \forall y\,(1 = y \to y = 1)$\lbl{A2}}
\BIC{$\forall y\,(1 = y \to y = 1)$}
\AXC{$\forall y\,(1 = y \to y = 1) \to 1 = 1^{-1} \to 1^{-1} = 1$\lbl{A2}}
\BIC{$1 = 1^{-1} \to 1^{-1} = 1$}
\end{ptree}
\noindent$\Pi_{3}$:
\begin{ptree}
\AXC{$\forall x \in \mathbb{R}\,(x \cdot 1 = x)$\lbl{(R8)}}
\AXC{$\forall x \in \mathbb{R}\,(x \cdot 1 = x) \to 1 \in \mathbb{R} \to 1 \cdot 1 = 1$\lbl{A2}}
\BIC{$1 \in \mathbb{R} \to 1 \cdot 1 = 1$}
\AXC{$(1 \in \mathbb{R} \to 1 \cdot 1 = 1) \to 1 \in \mathbb{R} \to 1 \cdot 1 = 1$\lbl{A1}}
\BIC{$1 \in \mathbb{R} \to 1 \cdot 1 = 1$}
\end{ptree}
\noindent$\Pi_{4}$:
\begin{ptree}
\AXC{$\stackrel{\Pi_{1}}{\vdots}$}\noLine
\UIC{$\forall x\,\forall y\,(x = y \to y = x)$}
\AXC{$\forall x\,\forall y\,(x = y \to y = x) \to \forall y\,(1 = y \to y = 1)$\lbl{A2}}
\BIC{$\forall y\,(1 = y \to y = 1)$}
\AXC{$\forall y\,(1 = y \to y = 1) \to 1 = 0 \to 0 = 1$\lbl{A2}}
\BIC{$1 = 0 \to 0 = 1$}
\end{ptree}
\noindent$\Pi_{5}$:
\begin{ptree}
\AXC{$1 \in \mathbb{R}$\lbl{(R1)}}
\AXC{$1 \in \mathbb{R}$\lbl{(R1)}}
\AXC{$\varphi_{1}$\lbl{[E13]}}
\AXC{$\varphi_{4}$\lbl{A2}}
\BIC{$\forall y\,(1 \in \mathbb{R} \to y \in \mathbb{R} \to 1 \neq 0 \to 1 \cdot y = 1 \to y = 1^{-1})$}
\AXC{$\varphi_{3}$\lbl{A2}}
\BIC{$1 \in \mathbb{R} \to 1 \in \mathbb{R} \to 1 \neq 0 \to 1 \cdot 1 = 1 \to 1 = 1^{-1}$}
\AXC{$\varphi_{2}$\lbl{A1}}
\BIC{$1 \in \mathbb{R} \to 1 \in \mathbb{R} \to 1 \neq 0 \to 1 \cdot 1 = 1 \to 1 = 1^{-1}$}
\BIC{$1 \in \mathbb{R} \to 1 \neq 0 \to 1 \cdot 1 = 1 \to 1 = 1^{-1}$}
\BIC{$1 \neq 0 \to 1 \cdot 1 = 1 \to 1 = 1^{-1}$}
\end{ptree}
\noindent$\Pi_{6}$:
\begin{ptree}
\AXC{$\stackrel{\Pi_{4}}{\vdots}$}\noLine
\UIC{$1 = 0 \to 0 = 1$}
\AXC{$0 \neq 1$\lbl{(R8)}}
\AXC{$0 \neq 1 \to (1 = 0 \to 0 = 1) \to 1 \neq 0$\lbl{A1}}
\BIC{$(1 = 0 \to 0 = 1) \to 1 \neq 0$}
\BIC{$1 \neq 0$}
\AXC{$\stackrel{\Pi_{5}}{\vdots}$}\noLine
\UIC{$1 \neq 0 \to 1 \cdot 1 = 1 \to 1 = 1^{-1}$}
\BIC{$1 \cdot 1 = 1 \to 1 = 1^{-1}$}
\end{ptree}
\noindent$\Pi_{7}$:
\begin{ptree}
\AXC{$1 \in \mathbb{R}$\lbl{(R1)}}
\AXC{$\stackrel{\Pi_{3}}{\vdots}$}\noLine
\UIC{$1 \in \mathbb{R} \to 1 \cdot 1 = 1$}
\BIC{$1 \cdot 1 = 1$}
\AXC{$\stackrel{\Pi_{6}}{\vdots}$}\noLine
\UIC{$1 \cdot 1 = 1 \to 1 = 1^{-1}$}
\BIC{$1 = 1^{-1}$}
\AXC{$1 = 1^{-1} \to (1 = 1^{-1} \to 1^{-1} = 1) \to 1^{-1} = 1$\lbl{A1}}
\BIC{$(1 = 1^{-1} \to 1^{-1} = 1) \to 1^{-1} = 1$}
\end{ptree}
\noindent 主木:
\begin{ptree}
\AXC{$\stackrel{\Pi_{2}}{\vdots}$}\noLine
\UIC{$1 = 1^{-1} \to 1^{-1} = 1$}
\AXC{$\stackrel{\Pi_{7}}{\vdots}$}\noLine
\UIC{$(1 = 1^{-1} \to 1^{-1} = 1) \to 1^{-1} = 1$}
\BIC{$1^{-1} = 1$}
\end{ptree}
\begin{abbrevlist}
\item $\varphi_{1} :\equiv \forall x\,\forall y\,(x \in \mathbb{R} \to y \in \mathbb{R} \to x \neq 0 \to x \cdot y = 1 \to y = x^{-1})$
\item $\varphi_{2} :\equiv (1 \in \mathbb{R} \to 1 \in \mathbb{R} \to 1 \neq 0 \to 1 \cdot 1 = 1 \to 1 = 1^{-1}) \to 1 \in \mathbb{R} \to 1 \in \mathbb{R} \to 1 \neq 0 \to 1 \cdot 1 = 1 \to 1 = 1^{-1}$
\item $\varphi_{3} :\equiv \forall y\,(1 \in \mathbb{R} \to y \in \mathbb{R} \to 1 \neq 0 \to 1 \cdot y = 1 \to y = 1^{-1}) \to 1 \in \mathbb{R} \to 1 \in \mathbb{R} \to 1 \neq 0 \to 1 \cdot 1 = 1 \to 1 = 1^{-1}$
\item $\varphi_{4} :\equiv \varphi_{1} \to \forall y\,(1 \in \mathbb{R} \to y \in \mathbb{R} \to 1 \neq 0 \to 1 \cdot y = 1 \to y = 1^{-1})$
\end{abbrevlist}
\end{fullproof}

% 自動生成 (prover/build.js): AbsSelf0  ノード数 69
\begin{fullproof}{[E36]}{$\forall x \in \mathbb{R}\,(0 < \lvert x - x \rvert \to x \neq x)$}{69}\phantomsection\label{tree:AbsSelf0}
\noindent$\Pi_{1}$:
\begin{ptree}
\AXC{$y = z \to x = y \to x = z$\lbl{A6}}
\AXC{$(y = z \to x = y \to x = z) \to x = y \to y = z \to x = z$\lbl{A1}}
\BIC{$x = y \to y = z \to x = z$}
\UIC{$\forall z\,(x = y \to y = z \to x = z)$}
\UIC{$\forall y\,\forall z\,(x = y \to y = z \to x = z)$}
\UIC{$\forall x\,\forall y\,\forall z\,(x = y \to y = z \to x = z)$}
\AXC{$\varphi_{2}$\lbl{A2}}
\BIC{$\forall y\,\forall z\,(\lvert x - x \rvert = y \to y = z \to \lvert x - x \rvert = z)$}
\AXC{$\varphi_{7}$\lbl{A2}}
\BIC{$\forall z\,(\lvert x - x \rvert = \lvert 0 \rvert \to \lvert 0 \rvert = z \to \lvert x - x \rvert = z)$}
\AXC{$\varphi_{9}$\lbl{A2}}
\BIC{$\lvert x - x \rvert = \lvert 0 \rvert \to \lvert 0 \rvert = 0 \to \lvert x - x \rvert = 0$}
\end{ptree}
\noindent$\Pi_{2}$:
\begin{ptree}
\AXC{$\forall x\,(0 \leq x \to \lvert x \rvert = x)$\lbl{[A1]}}
\AXC{$\forall x\,(0 \leq x \to \lvert x \rvert = x) \to 0 \leq 0 \to \lvert 0 \rvert = 0$\lbl{A2}}
\BIC{$0 \leq 0 \to \lvert 0 \rvert = 0$}
\AXC{$(0 \leq 0 \to \lvert 0 \rvert = 0) \to \varphi_{8}$\lbl{A1}}
\BIC{$\varphi_{8}$}
\end{ptree}
\noindent$\Pi_{3}$:
\begin{ptree}
\AXC{$x = x$\lbl{A5}}
\UIC{$\forall x\,(x = x)$}
\AXC{$\forall x\,(x = x) \to 0 = 0$\lbl{A2}}
\BIC{$0 = 0$}
\AXC{$0 = 0 \to x \in \mathbb{R} \to 0 < \lvert x - x \rvert \to \neg x \neq x \to 0 \leq 0$\lbl{A1}}
\BIC{$x \in \mathbb{R} \to 0 < \lvert x - x \rvert \to \neg x \neq x \to 0 \leq 0$}
\AXC{$\stackrel{\Pi_{2}}{\vdots}$}\noLine
\UIC{$\varphi_{8}$}
\BIC{$x \in \mathbb{R} \to 0 < \lvert x - x \rvert \to \neg x \neq x \to \lvert 0 \rvert = 0$}
\end{ptree}
\noindent$\Pi_{4}$:
\begin{ptree}
\AXC{$z = w \to \lvert x - x \rvert = \lvert z \rvert \to \lvert x - x \rvert = \lvert w \rvert$\lbl{A6}}
\UIC{$\forall w\,(z = w \to \lvert x - x \rvert = \lvert z \rvert \to \lvert x - x \rvert = \lvert w \rvert)$}
\UIC{$\forall z\,\forall w\,(z = w \to \lvert x - x \rvert = \lvert z \rvert \to \lvert x - x \rvert = \lvert w \rvert)$}
\AXC{$\varphi_{12}$\lbl{A2}}
\BIC{$\varphi_{1}$}
\AXC{$\varphi_{1} \to x - x = 0 \to \lvert x - x \rvert = \lvert x - x \rvert \to \lvert x - x \rvert = \lvert 0 \rvert$\lbl{A2}}
\BIC{$x - x = 0 \to \lvert x - x \rvert = \lvert x - x \rvert \to \lvert x - x \rvert = \lvert 0 \rvert$}
\end{ptree}
\noindent$\Pi_{5}$:
\begin{ptree}
\AXC{$\forall x \in \mathbb{R}\,(x - x = 0)$\lbl{[F13]}}
\AXC{$\forall x \in \mathbb{R}\,(x - x = 0) \to x \in \mathbb{R} \to x - x = 0$\lbl{A2}}
\BIC{$x \in \mathbb{R} \to x - x = 0$}
\AXC{$(x \in \mathbb{R} \to x - x = 0) \to \varphi_{10}$\lbl{A1}}
\BIC{$\varphi_{10}$}
\end{ptree}
\noindent$\Pi_{6}$:
\begin{ptree}
\AXC{$x \in \mathbb{R} \to 0 < \lvert x - x \rvert \to \neg x \neq x \to x \in \mathbb{R}$\lbl{A1}}
\AXC{$\stackrel{\Pi_{5}}{\vdots}$}\noLine
\UIC{$\varphi_{10}$}
\BIC{$x \in \mathbb{R} \to 0 < \lvert x - x \rvert \to \neg x \neq x \to x - x = 0$}
\AXC{$\varphi_{14}$\lbl{A1}}
\BIC{$\lvert x - x \rvert = \lvert x - x \rvert \to \varphi_{13}$}
\end{ptree}
\noindent$\Pi_{7}$:
\begin{ptree}
\AXC{$x = x$\lbl{A5}}
\UIC{$\forall x\,(x = x)$}
\AXC{$\forall x\,(x = x) \to \lvert x - x \rvert = \lvert x - x \rvert$\lbl{A2}}
\BIC{$\lvert x - x \rvert = \lvert x - x \rvert$}
\AXC{$\stackrel{\Pi_{6}}{\vdots}$}\noLine
\UIC{$\lvert x - x \rvert = \lvert x - x \rvert \to \varphi_{13}$}
\BIC{$\varphi_{13}$}
\end{ptree}
\noindent$\Pi_{8}$:
\begin{ptree}
\AXC{$\stackrel{\Pi_{3}}{\vdots}$}\noLine
\UIC{$x \in \mathbb{R} \to 0 < \lvert x - x \rvert \to \neg x \neq x \to \lvert 0 \rvert = 0$}
\AXC{$\stackrel{\Pi_{4}}{\vdots}$}\noLine
\UIC{$x - x = 0 \to \lvert x - x \rvert = \lvert x - x \rvert \to \lvert x - x \rvert = \lvert 0 \rvert$}
\AXC{$\stackrel{\Pi_{7}}{\vdots}$}\noLine
\UIC{$\varphi_{13}$}
\BIC{$x \in \mathbb{R} \to 0 < \lvert x - x \rvert \to \neg x \neq x \to \lvert x - x \rvert = \lvert 0 \rvert$}
\AXC{$\varphi_{16}$\lbl{A1}}
\BIC{$(x \in \mathbb{R} \to 0 < \lvert x - x \rvert \to \neg x \neq x \to \lvert 0 \rvert = 0) \to \varphi_{11}$}
\BIC{$\varphi_{11}$}
\end{ptree}
\noindent$\Pi_{9}$:
\begin{ptree}
\AXC{$x \in \mathbb{R} \to 0 < \lvert x - x \rvert \to 0 < \lvert x - x \rvert$\lbl{A1}}
\AXC{$\stackrel{\Pi_{1}}{\vdots}$}\noLine
\UIC{$\lvert x - x \rvert = \lvert 0 \rvert \to \lvert 0 \rvert = 0 \to \lvert x - x \rvert = 0$}
\AXC{$\stackrel{\Pi_{8}}{\vdots}$}\noLine
\UIC{$\varphi_{11}$}
\BIC{$x \in \mathbb{R} \to 0 < \lvert x - x \rvert \to \neg x \neq x \to \lvert x - x \rvert = 0$}
\AXC{$\varphi_{15}$\lbl{A1}}
\BIC{$(x \in \mathbb{R} \to 0 < \lvert x - x \rvert \to 0 < \lvert x - x \rvert) \to \varphi_{6}$}
\BIC{$\varphi_{6}$}
\end{ptree}
\noindent$\Pi_{10}$:
\begin{ptree}
\AXC{$z = w \to 0 < z \to 0 < w$\lbl{A6}}
\UIC{$\forall w\,(z = w \to 0 < z \to 0 < w)$}
\UIC{$\forall z\,\forall w\,(z = w \to 0 < z \to 0 < w)$}
\AXC{$\varphi_{3}$\lbl{A2}}
\BIC{$\forall w\,(\lvert x - x \rvert = w \to 0 < \lvert x - x \rvert \to 0 < w)$}
\AXC{$\varphi_{5}$\lbl{A2}}
\BIC{$\lvert x - x \rvert = 0 \to 0 < \lvert x - x \rvert \to 0 < 0$}
\AXC{$\stackrel{\Pi_{9}}{\vdots}$}\noLine
\UIC{$\varphi_{6}$}
\BIC{$x \in \mathbb{R} \to 0 < \lvert x - x \rvert \to \neg x \neq x \to 0 < 0$}
\AXC{$(x \in \mathbb{R} \to 0 < \lvert x - x \rvert \to \neg x \neq x \to 0 < 0) \to \varphi_{4}$\lbl{A1}}
\BIC{$\varphi_{4}$}
\end{ptree}
\noindent 主木:
\begin{ptree}
\AXC{$\forall x\,(\neg x < x)$\lbl{[O1]}}
\AXC{$\forall x\,(\neg x < x) \to \neg 0 < 0$\lbl{A2}}
\BIC{$\neg 0 < 0$}
\AXC{$\neg 0 < 0 \to x \in \mathbb{R} \to 0 < \lvert x - x \rvert \to \neg x \neq x \to \neg 0 < 0$\lbl{A1}}
\BIC{$x \in \mathbb{R} \to 0 < \lvert x - x \rvert \to \neg x \neq x \to \neg 0 < 0$}
\AXC{$\stackrel{\Pi_{10}}{\vdots}$}\noLine
\UIC{$\varphi_{4}$}
\BIC{$x \in \mathbb{R} \to 0 < \lvert x - x \rvert \to x \neq x$}
\UIC{$\forall x \in \mathbb{R}\,(0 < \lvert x - x \rvert \to x \neq x)$}
\end{ptree}
\begin{abbrevlist}
\item $\varphi_{1} :\equiv \forall w\,(x - x = w \to \lvert x - x \rvert = \lvert x - x \rvert \to \lvert x - x \rvert = \lvert w \rvert)$
\item $\varphi_{2} :\equiv \forall x\,\forall y\,\forall z\,(x = y \to y = z \to x = z) \to \forall y\,\forall z\,(\lvert x - x \rvert = y \to y = z \to \lvert x - x \rvert = z)$
\item $\varphi_{3} :\equiv \forall z\,\forall w\,(z = w \to 0 < z \to 0 < w) \to \forall w\,(\lvert x - x \rvert = w \to 0 < \lvert x - x \rvert \to 0 < w)$
\item $\varphi_{4} :\equiv (x \in \mathbb{R} \to 0 < \lvert x - x \rvert \to \neg x \neq x \to \neg 0 < 0) \to x \in \mathbb{R} \to 0 < \lvert x - x \rvert \to x \neq x$
\item $\varphi_{5} :\equiv \forall w\,(\lvert x - x \rvert = w \to 0 < \lvert x - x \rvert \to 0 < w) \to \lvert x - x \rvert = 0 \to 0 < \lvert x - x \rvert \to 0 < 0$
\item $\varphi_{6} :\equiv (\lvert x - x \rvert = 0 \to 0 < \lvert x - x \rvert \to 0 < 0) \to x \in \mathbb{R} \to 0 < \lvert x - x \rvert \to \neg x \neq x \to 0 < 0$
\item $\varphi_{7} :\equiv \forall y\,\forall z\,(\lvert x - x \rvert = y \to y = z \to \lvert x - x \rvert = z) \to \forall z\,(\lvert x - x \rvert = \lvert 0 \rvert \to \lvert 0 \rvert = z \to \lvert x - x \rvert = z)$
\item $\varphi_{8} :\equiv (x \in \mathbb{R} \to 0 < \lvert x - x \rvert \to \neg x \neq x \to 0 \leq 0) \to x \in \mathbb{R} \to 0 < \lvert x - x \rvert \to \neg x \neq x \to \lvert 0 \rvert = 0$
\item $\varphi_{9} :\equiv \forall z\,(\lvert x - x \rvert = \lvert 0 \rvert \to \lvert 0 \rvert = z \to \lvert x - x \rvert = z) \to \lvert x - x \rvert = \lvert 0 \rvert \to \lvert 0 \rvert = 0 \to \lvert x - x \rvert = 0$
\item $\varphi_{10} :\equiv (x \in \mathbb{R} \to 0 < \lvert x - x \rvert \to \neg x \neq x \to x \in \mathbb{R}) \to x \in \mathbb{R} \to 0 < \lvert x - x \rvert \to \neg x \neq x \to x - x = 0$
\item $\varphi_{11} :\equiv (\lvert x - x \rvert = \lvert 0 \rvert \to \lvert 0 \rvert = 0 \to \lvert x - x \rvert = 0) \to x \in \mathbb{R} \to 0 < \lvert x - x \rvert \to \neg x \neq x \to \lvert x - x \rvert = 0$
\item $\varphi_{12} :\equiv \forall z\,\forall w\,(z = w \to \lvert x - x \rvert = \lvert z \rvert \to \lvert x - x \rvert = \lvert w \rvert) \to \varphi_{1}$
\item $\varphi_{13} :\equiv (x - x = 0 \to \lvert x - x \rvert = \lvert x - x \rvert \to \lvert x - x \rvert = \lvert 0 \rvert) \to x \in \mathbb{R} \to 0 < \lvert x - x \rvert \to \neg x \neq x \to \lvert x - x \rvert = \lvert 0 \rvert$
\item $\varphi_{14} :\equiv (x \in \mathbb{R} \to 0 < \lvert x - x \rvert \to \neg x \neq x \to x - x = 0) \to \lvert x - x \rvert = \lvert x - x \rvert \to \varphi_{13}$
\item $\varphi_{15} :\equiv (x \in \mathbb{R} \to 0 < \lvert x - x \rvert \to \neg x \neq x \to \lvert x - x \rvert = 0) \to (x \in \mathbb{R} \to 0 < \lvert x - x \rvert \to 0 < \lvert x - x \rvert) \to \varphi_{6}$
\item $\varphi_{16} :\equiv (x \in \mathbb{R} \to 0 < \lvert x - x \rvert \to \neg x \neq x \to \lvert x - x \rvert = \lvert 0 \rvert) \to (x \in \mathbb{R} \to 0 < \lvert x - x \rvert \to \neg x \neq x \to \lvert 0 \rvert = 0) \to \varphi_{11}$
\end{abbrevlist}
\end{fullproof}

% 自動生成 (prover/build.js): PosAbsNe  ノード数 34
\begin{fullproof}{[E37]}{$\forall y\,\forall c\,(y \in \mathbb{R} \to c \in \mathbb{R} \to 0 < \lvert y - c \rvert \to y \neq c)$}{34}\phantomsection\label{tree:PosAbsNe}
\noindent$\Pi_{1}$:
\begin{ptree}
\AXC{$y \in \mathbb{R} \to c \in \mathbb{R} \to 0 < \lvert y - c \rvert \to \neg y \neq c \to \neg y \neq c$\lbl{A1}}
\AXC{$\varphi_{6}$\lbl{A1}}
\BIC{$y \in \mathbb{R} \to c \in \mathbb{R} \to 0 < \lvert y - c \rvert \to \neg y \neq c \to y = c$}
\AXC{$\varphi_{9}$\lbl{A1}}
\BIC{$(y \in \mathbb{R} \to c \in \mathbb{R} \to 0 < \lvert y - c \rvert \to 0 < \lvert y - c \rvert) \to \varphi_{5}$}
\end{ptree}
\noindent$\Pi_{2}$:
\begin{ptree}
\AXC{$y \in \mathbb{R} \to c \in \mathbb{R} \to 0 < \lvert y - c \rvert \to 0 < \lvert y - c \rvert$\lbl{A1}}
\AXC{$\stackrel{\Pi_{1}}{\vdots}$}\noLine
\UIC{$(y \in \mathbb{R} \to c \in \mathbb{R} \to 0 < \lvert y - c \rvert \to 0 < \lvert y - c \rvert) \to \varphi_{5}$}
\BIC{$\varphi_{5}$}
\end{ptree}
\noindent$\Pi_{3}$:
\begin{ptree}
\AXC{$z = w \to 0 < \lvert z - c \rvert \to 0 < \lvert w - c \rvert$\lbl{A6}}
\UIC{$\forall w\,(z = w \to 0 < \lvert z - c \rvert \to 0 < \lvert w - c \rvert)$}
\UIC{$\forall z\,\forall w\,(z = w \to 0 < \lvert z - c \rvert \to 0 < \lvert w - c \rvert)$}
\AXC{$\varphi_{3}$\lbl{A2}}
\BIC{$\forall w\,(y = w \to 0 < \lvert y - c \rvert \to 0 < \lvert w - c \rvert)$}
\AXC{$\varphi_{2}$\lbl{A2}}
\BIC{$y = c \to 0 < \lvert y - c \rvert \to 0 < \lvert c - c \rvert$}
\AXC{$\stackrel{\Pi_{2}}{\vdots}$}\noLine
\UIC{$\varphi_{5}$}
\BIC{$y \in \mathbb{R} \to c \in \mathbb{R} \to 0 < \lvert y - c \rvert \to \neg y \neq c \to 0 < \lvert c - c \rvert$}
\end{ptree}
\noindent$\Pi_{4}$:
\begin{ptree}
\AXC{$\forall x \in \mathbb{R}\,(0 < \lvert x - x \rvert \to x \neq x)$\lbl{[E36]}}
\AXC{$\varphi_{1}$\lbl{A2}}
\BIC{$c \in \mathbb{R} \to 0 < \lvert c - c \rvert \to c \neq c$}
\AXC{$(c \in \mathbb{R} \to 0 < \lvert c - c \rvert \to c \neq c) \to \varphi_{8}$\lbl{A1}}
\BIC{$\varphi_{8}$}
\end{ptree}
\noindent$\Pi_{5}$:
\begin{ptree}
\AXC{$\stackrel{\Pi_{3}}{\vdots}$}\noLine
\UIC{$y \in \mathbb{R} \to c \in \mathbb{R} \to 0 < \lvert y - c \rvert \to \neg y \neq c \to 0 < \lvert c - c \rvert$}
\AXC{$y \in \mathbb{R} \to c \in \mathbb{R} \to 0 < \lvert y - c \rvert \to \neg y \neq c \to c \in \mathbb{R}$\lbl{A1}}
\AXC{$\stackrel{\Pi_{4}}{\vdots}$}\noLine
\UIC{$\varphi_{8}$}
\BIC{$\varphi_{7}$}
\BIC{$y \in \mathbb{R} \to c \in \mathbb{R} \to 0 < \lvert y - c \rvert \to \neg y \neq c \to c \neq c$}
\end{ptree}
\noindent$\Pi_{6}$:
\begin{ptree}
\AXC{$\stackrel{\Pi_{5}}{\vdots}$}\noLine
\UIC{$y \in \mathbb{R} \to c \in \mathbb{R} \to 0 < \lvert y - c \rvert \to \neg y \neq c \to c \neq c$}
\AXC{$(y \in \mathbb{R} \to c \in \mathbb{R} \to 0 < \lvert y - c \rvert \to \neg y \neq c \to c \neq c) \to \varphi_{4}$\lbl{A1}}
\BIC{$\varphi_{4}$}
\end{ptree}
\noindent 主木:
\begin{ptree}
\AXC{$x = x$\lbl{A5}}
\UIC{$\forall x\,(x = x)$}
\AXC{$\forall x\,(x = x) \to c = c$\lbl{A2}}
\BIC{$c = c$}
\AXC{$c = c \to y \in \mathbb{R} \to c \in \mathbb{R} \to 0 < \lvert y - c \rvert \to \neg y \neq c \to c = c$\lbl{A1}}
\BIC{$y \in \mathbb{R} \to c \in \mathbb{R} \to 0 < \lvert y - c \rvert \to \neg y \neq c \to c = c$}
\AXC{$\stackrel{\Pi_{6}}{\vdots}$}\noLine
\UIC{$\varphi_{4}$}
\BIC{$y \in \mathbb{R} \to c \in \mathbb{R} \to 0 < \lvert y - c \rvert \to y \neq c$}
\UIC{$\forall c\,(y \in \mathbb{R} \to c \in \mathbb{R} \to 0 < \lvert y - c \rvert \to y \neq c)$}
\UIC{$\forall y\,\forall c\,(y \in \mathbb{R} \to c \in \mathbb{R} \to 0 < \lvert y - c \rvert \to y \neq c)$}
\end{ptree}
\begin{abbrevlist}
\item $\varphi_{1} :\equiv \forall x \in \mathbb{R}\,(0 < \lvert x - x \rvert \to x \neq x) \to c \in \mathbb{R} \to 0 < \lvert c - c \rvert \to c \neq c$
\item $\varphi_{2} :\equiv \forall w\,(y = w \to 0 < \lvert y - c \rvert \to 0 < \lvert w - c \rvert) \to y = c \to 0 < \lvert y - c \rvert \to 0 < \lvert c - c \rvert$
\item $\varphi_{3} :\equiv \forall z\,\forall w\,(z = w \to 0 < \lvert z - c \rvert \to 0 < \lvert w - c \rvert) \to \forall w\,(y = w \to 0 < \lvert y - c \rvert \to 0 < \lvert w - c \rvert)$
\item $\varphi_{4} :\equiv (y \in \mathbb{R} \to c \in \mathbb{R} \to 0 < \lvert y - c \rvert \to \neg y \neq c \to c = c) \to y \in \mathbb{R} \to c \in \mathbb{R} \to 0 < \lvert y - c \rvert \to y \neq c$
\item $\varphi_{5} :\equiv (y = c \to 0 < \lvert y - c \rvert \to 0 < \lvert c - c \rvert) \to y \in \mathbb{R} \to c \in \mathbb{R} \to 0 < \lvert y - c \rvert \to \neg y \neq c \to 0 < \lvert c - c \rvert$
\item $\varphi_{6} :\equiv (y \in \mathbb{R} \to c \in \mathbb{R} \to 0 < \lvert y - c \rvert \to \neg y \neq c \to \neg y \neq c) \to y \in \mathbb{R} \to c \in \mathbb{R} \to 0 < \lvert y - c \rvert \to \neg y \neq c \to y = c$
\item $\varphi_{7} :\equiv (y \in \mathbb{R} \to c \in \mathbb{R} \to 0 < \lvert y - c \rvert \to \neg y \neq c \to 0 < \lvert c - c \rvert) \to y \in \mathbb{R} \to c \in \mathbb{R} \to 0 < \lvert y - c \rvert \to \neg y \neq c \to c \neq c$
\item $\varphi_{8} :\equiv (y \in \mathbb{R} \to c \in \mathbb{R} \to 0 < \lvert y - c \rvert \to \neg y \neq c \to c \in \mathbb{R}) \to \varphi_{7}$
\item $\varphi_{9} :\equiv (y \in \mathbb{R} \to c \in \mathbb{R} \to 0 < \lvert y - c \rvert \to \neg y \neq c \to y = c) \to (y \in \mathbb{R} \to c \in \mathbb{R} \to 0 < \lvert y - c \rvert \to 0 < \lvert y - c \rvert) \to \varphi_{5}$
\end{abbrevlist}
\end{fullproof}

% 自動生成 (prover/build.js): ApproxSeq  ノード数 818
\begin{fullproof}{[H1]}{$\forall c\,\forall r\,(c \in \mathbb{R} \to r \in \mathbb{R} \to 0 < r \to \exists a\,(\mathrm{Seq}(a) \land \mathrm{Conv}(a,c) \land \forall n \in \mathbb{N}\,(0 < \lvert a_{n} - c \rvert \land \lvert a_{n} - c \rvert < r)))$}{818}\phantomsection\label{tree:ApproxSeq}
\noindent$\Pi_{1}$:
\begin{ptree}
\AXC{$z = w \to r \cdot (\underline{n} + 1)^{-1} < z \to r \cdot (\underline{n} + 1)^{-1} < w$\lbl{A6}}
\UIC{$\varphi_{7}$}
\UIC{$\forall z\,(\varphi_{7})$}
\AXC{$\forall z\,(\varphi_{7}) \to \varphi_{31}$\lbl{A2}}
\BIC{$\varphi_{31}$}
\AXC{$\varphi_{31} \to \varphi_{21}$\lbl{A2}}
\BIC{$\varphi_{21}$}
\end{ptree}
\noindent$\Pi_{2}$:
\begin{ptree}
\AXC{$z = w \to (\underline{n} + 1)^{-1} < z \to (\underline{n} + 1)^{-1} < w$\lbl{A6}}
\UIC{$\forall w\,(z = w \to (\underline{n} + 1)^{-1} < z \to (\underline{n} + 1)^{-1} < w)$}
\UIC{$\varphi_{3}$}
\AXC{$\varphi_{3} \to \varphi_{8}$\lbl{A2}}
\BIC{$\varphi_{8}$}
\AXC{$\varphi_{176}$\lbl{A2}}
\BIC{$1^{-1} = 1 \to (\underline{n} + 1)^{-1} < 1^{-1} \to (\underline{n} + 1)^{-1} < 1$}
\end{ptree}
\noindent$\Pi_{3}$:
\begin{ptree}
\AXC{$\forall n \in \mathbb{N}\,(0 < \underline{n})$\lbl{[N15]}}
\AXC{$\forall n \in \mathbb{N}\,(0 < \underline{n}) \to n \in \mathbb{N} \to 0 < \underline{n}$\lbl{A2}}
\BIC{$n \in \mathbb{N} \to 0 < \underline{n}$}
\AXC{$(n \in \mathbb{N} \to 0 < \underline{n}) \to \varphi_{153}$\lbl{A1}}
\BIC{$\varphi_{153}$}
\end{ptree}
\noindent$\Pi_{4}$:
\begin{ptree}
\AXC{$\forall n \in \mathbb{N}\,(\underline{n} \in \mathbb{R})$\lbl{[N4]}}
\AXC{$\forall n \in \mathbb{N}\,(\underline{n} \in \mathbb{R}) \to n \in \mathbb{N} \to \underline{n} \in \mathbb{R}$\lbl{A2}}
\BIC{$n \in \mathbb{N} \to \underline{n} \in \mathbb{R}$}
\AXC{$(n \in \mathbb{N} \to \underline{n} \in \mathbb{R}) \to \varphi_{163}$\lbl{A1}}
\BIC{$\varphi_{163}$}
\end{ptree}
\noindent$\Pi_{5}$:
\begin{ptree}
\AXC{$c \in \mathbb{R} \to r \in \mathbb{R} \to 0 < r \to n \in \mathbb{N} \to n \in \mathbb{N}$\lbl{A1}}
\AXC{$\stackrel{\Pi_{4}}{\vdots}$}\noLine
\UIC{$\varphi_{163}$}
\BIC{$c \in \mathbb{R} \to r \in \mathbb{R} \to 0 < r \to n \in \mathbb{N} \to \underline{n} \in \mathbb{R}$}
\end{ptree}
\noindent$\Pi_{6}$:
\begin{ptree}
\AXC{$\forall x\,\forall y\,(x \in \mathbb{R} \to y \in \mathbb{R} \to 0 < y \to x < x + y)$\lbl{[E3]}}
\AXC{$\varphi_{127}$\lbl{A2}}
\BIC{$\forall y\,(1 \in \mathbb{R} \to y \in \mathbb{R} \to 0 < y \to 1 < 1 + y)$}
\AXC{$\varphi_{147}$\lbl{A2}}
\BIC{$1 \in \mathbb{R} \to \underline{n} \in \mathbb{R} \to 0 < \underline{n} \to 1 < 1 + \underline{n}$}
\end{ptree}
\noindent$\Pi_{7}$:
\begin{ptree}
\AXC{$1 \in \mathbb{R}$\lbl{(R1)}}
\AXC{$\stackrel{\Pi_{6}}{\vdots}$}\noLine
\UIC{$1 \in \mathbb{R} \to \underline{n} \in \mathbb{R} \to 0 < \underline{n} \to 1 < 1 + \underline{n}$}
\AXC{$(1 \in \mathbb{R} \to \underline{n} \in \mathbb{R} \to 0 < \underline{n} \to 1 < 1 + \underline{n}) \to \varphi_{221}$\lbl{A1}}
\BIC{$\varphi_{221}$}
\BIC{$(c \in \mathbb{R} \to r \in \mathbb{R} \to 0 < r \to n \in \mathbb{N} \to \underline{n} \in \mathbb{R}) \to \varphi_{172}$}
\end{ptree}
\noindent$\Pi_{8}$:
\begin{ptree}
\AXC{$\stackrel{\Pi_{5}}{\vdots}$}\noLine
\UIC{$c \in \mathbb{R} \to r \in \mathbb{R} \to 0 < r \to n \in \mathbb{N} \to \underline{n} \in \mathbb{R}$}
\AXC{$\stackrel{\Pi_{7}}{\vdots}$}\noLine
\UIC{$(c \in \mathbb{R} \to r \in \mathbb{R} \to 0 < r \to n \in \mathbb{N} \to \underline{n} \in \mathbb{R}) \to \varphi_{172}$}
\BIC{$\varphi_{172}$}
\end{ptree}
\noindent$\Pi_{9}$:
\begin{ptree}
\AXC{$\forall x\,\forall y\,(x \in \mathbb{R} \to y \in \mathbb{R} \to x + y = y + x)$\lbl{(R3)}}
\AXC{$\varphi_{90}$\lbl{A2}}
\BIC{$\forall y\,(1 \in \mathbb{R} \to y \in \mathbb{R} \to 1 + y = y + 1)$}
\AXC{$\varphi_{118}$\lbl{A2}}
\BIC{$1 \in \mathbb{R} \to \underline{n} \in \mathbb{R} \to 1 + \underline{n} = \underline{n} + 1$}
\end{ptree}
\noindent$\Pi_{10}$:
\begin{ptree}
\AXC{$1 \in \mathbb{R}$\lbl{(R1)}}
\AXC{$\stackrel{\Pi_{9}}{\vdots}$}\noLine
\UIC{$1 \in \mathbb{R} \to \underline{n} \in \mathbb{R} \to 1 + \underline{n} = \underline{n} + 1$}
\AXC{$(1 \in \mathbb{R} \to \underline{n} \in \mathbb{R} \to 1 + \underline{n} = \underline{n} + 1) \to 1 \in \mathbb{R} \to \varphi_{184}$\lbl{A1}}
\BIC{$1 \in \mathbb{R} \to \varphi_{184}$}
\BIC{$\varphi_{184}$}
\end{ptree}
\noindent$\Pi_{11}$:
\begin{ptree}
\AXC{$\stackrel{\Pi_{5}}{\vdots}$}\noLine
\UIC{$c \in \mathbb{R} \to r \in \mathbb{R} \to 0 < r \to n \in \mathbb{N} \to \underline{n} \in \mathbb{R}$}
\AXC{$\stackrel{\Pi_{10}}{\vdots}$}\noLine
\UIC{$\varphi_{184}$}
\BIC{$c \in \mathbb{R} \to r \in \mathbb{R} \to 0 < r \to n \in \mathbb{N} \to 1 + \underline{n} = \underline{n} + 1$}
\AXC{$\varphi_{234}$\lbl{A1}}
\BIC{$(c \in \mathbb{R} \to r \in \mathbb{R} \to 0 < r \to n \in \mathbb{N} \to 1 < 1 + \underline{n}) \to \varphi_{177}$}
\end{ptree}
\noindent$\Pi_{12}$:
\begin{ptree}
\AXC{$c \in \mathbb{R} \to r \in \mathbb{R} \to 0 < r \to n \in \mathbb{N} \to n \in \mathbb{N}$\lbl{A1}}
\AXC{$\stackrel{\Pi_{3}}{\vdots}$}\noLine
\UIC{$\varphi_{153}$}
\BIC{$c \in \mathbb{R} \to r \in \mathbb{R} \to 0 < r \to n \in \mathbb{N} \to 0 < \underline{n}$}
\AXC{$\stackrel{\Pi_{8}}{\vdots}$}\noLine
\UIC{$\varphi_{172}$}
\BIC{$c \in \mathbb{R} \to r \in \mathbb{R} \to 0 < r \to n \in \mathbb{N} \to 1 < 1 + \underline{n}$}
\AXC{$\stackrel{\Pi_{11}}{\vdots}$}\noLine
\UIC{$(c \in \mathbb{R} \to r \in \mathbb{R} \to 0 < r \to n \in \mathbb{N} \to 1 < 1 + \underline{n}) \to \varphi_{177}$}
\BIC{$\varphi_{177}$}
\end{ptree}
\noindent$\Pi_{13}$:
\begin{ptree}
\AXC{$\forall x\,\forall y\,(x \in \mathbb{R} \to y \in \mathbb{R} \to x + y \in \mathbb{R})$\lbl{(R1)}}
\AXC{$\varphi_{130}$\lbl{A2}}
\BIC{$\forall y\,(\underline{n} \in \mathbb{R} \to y \in \mathbb{R} \to \underline{n} + y \in \mathbb{R})$}
\AXC{$\varphi_{137}$\lbl{A2}}
\BIC{$\underline{n} \in \mathbb{R} \to 1 \in \mathbb{R} \to \underline{n} + 1 \in \mathbb{R}$}
\AXC{$\varphi_{220}$\lbl{A1}}
\BIC{$(c \in \mathbb{R} \to r \in \mathbb{R} \to 0 < r \to n \in \mathbb{N} \to \underline{n} \in \mathbb{R}) \to \varphi_{54}$}
\end{ptree}
\noindent$\Pi_{14}$:
\begin{ptree}
\AXC{$1 \in \mathbb{R}$\lbl{(R1)}}
\AXC{$\stackrel{\Pi_{5}}{\vdots}$}\noLine
\UIC{$c \in \mathbb{R} \to r \in \mathbb{R} \to 0 < r \to n \in \mathbb{N} \to \underline{n} \in \mathbb{R}$}
\AXC{$\stackrel{\Pi_{13}}{\vdots}$}\noLine
\UIC{$(c \in \mathbb{R} \to r \in \mathbb{R} \to 0 < r \to n \in \mathbb{N} \to \underline{n} \in \mathbb{R}) \to \varphi_{54}$}
\BIC{$\varphi_{54}$}
\BIC{$c \in \mathbb{R} \to r \in \mathbb{R} \to 0 < r \to n \in \mathbb{N} \to \underline{n} + 1 \in \mathbb{R}$}
\end{ptree}
\noindent$\Pi_{15}$:
\begin{ptree}
\AXC{$1 \in \mathbb{R}$\lbl{(R1)}}
\AXC{$\varphi_{9}$\lbl{[E19]}}
\AXC{$\varphi_{178}$\lbl{A2}}
\BIC{$\forall y\,(1 \in \mathbb{R} \to y \in \mathbb{R} \to 0 < 1 \to 1 < y \to y^{-1} < 1^{-1})$}
\AXC{$\varphi_{197}$\lbl{A2}}
\BIC{$1 \in \mathbb{R} \to \varphi_{53}$}
\AXC{$(1 \in \mathbb{R} \to \varphi_{53}) \to \varphi_{240}$\lbl{A1}}
\BIC{$\varphi_{240}$}
\BIC{$(c \in \mathbb{R} \to r \in \mathbb{R} \to 0 < r \to n \in \mathbb{N} \to \underline{n} + 1 \in \mathbb{R}) \to \varphi_{217}$}
\end{ptree}
\noindent$\Pi_{16}$:
\begin{ptree}
\AXC{$\stackrel{\Pi_{14}}{\vdots}$}\noLine
\UIC{$c \in \mathbb{R} \to r \in \mathbb{R} \to 0 < r \to n \in \mathbb{N} \to \underline{n} + 1 \in \mathbb{R}$}
\AXC{$\stackrel{\Pi_{15}}{\vdots}$}\noLine
\UIC{$(c \in \mathbb{R} \to r \in \mathbb{R} \to 0 < r \to n \in \mathbb{N} \to \underline{n} + 1 \in \mathbb{R}) \to \varphi_{217}$}
\BIC{$\varphi_{217}$}
\end{ptree}
\noindent$\Pi_{17}$:
\begin{ptree}
\AXC{$0 < 1$\lbl{[F28]}}
\AXC{$0 < 1 \to c \in \mathbb{R} \to r \in \mathbb{R} \to 0 < r \to n \in \mathbb{N} \to 0 < 1$\lbl{A1}}
\BIC{$c \in \mathbb{R} \to r \in \mathbb{R} \to 0 < r \to n \in \mathbb{N} \to 0 < 1$}
\AXC{$\stackrel{\Pi_{16}}{\vdots}$}\noLine
\UIC{$\varphi_{217}$}
\BIC{$(c \in \mathbb{R} \to r \in \mathbb{R} \to 0 < r \to n \in \mathbb{N} \to 1 < \underline{n} + 1) \to \varphi_{26}$}
\end{ptree}
\noindent$\Pi_{18}$:
\begin{ptree}
\AXC{$z = w \to 1 < z \to 1 < w$\lbl{A6}}
\UIC{$\forall w\,(z = w \to 1 < z \to 1 < w)$}
\UIC{$\forall z\,\forall w\,(z = w \to 1 < z \to 1 < w)$}
\AXC{$\varphi_{32}$\lbl{A2}}
\BIC{$\forall w\,(1 + \underline{n} = w \to 1 < 1 + \underline{n} \to 1 < w)$}
\AXC{$\varphi_{141}$\lbl{A2}}
\BIC{$1 + \underline{n} = \underline{n} + 1 \to 1 < 1 + \underline{n} \to 1 < \underline{n} + 1$}
\AXC{$\stackrel{\Pi_{12}}{\vdots}$}\noLine
\UIC{$\varphi_{177}$}
\BIC{$c \in \mathbb{R} \to r \in \mathbb{R} \to 0 < r \to n \in \mathbb{N} \to 1 < \underline{n} + 1$}
\AXC{$\stackrel{\Pi_{17}}{\vdots}$}\noLine
\UIC{$(c \in \mathbb{R} \to r \in \mathbb{R} \to 0 < r \to n \in \mathbb{N} \to 1 < \underline{n} + 1) \to \varphi_{26}$}
\BIC{$\varphi_{26}$}
\end{ptree}
\noindent$\Pi_{19}$:
\begin{ptree}
\AXC{$\stackrel{\Pi_{2}}{\vdots}$}\noLine
\UIC{$1^{-1} = 1 \to (\underline{n} + 1)^{-1} < 1^{-1} \to (\underline{n} + 1)^{-1} < 1$}
\AXC{$\stackrel{\Pi_{18}}{\vdots}$}\noLine
\UIC{$\varphi_{26}$}
\AXC{$1^{-1} = 1$\lbl{[E35]}}
\AXC{$1^{-1} = 1 \to \varphi_{26} \to \varphi_{181}$\lbl{A1}}
\BIC{$\varphi_{26} \to \varphi_{181}$}
\BIC{$\varphi_{181}$}
\BIC{$c \in \mathbb{R} \to r \in \mathbb{R} \to 0 < r \to n \in \mathbb{N} \to (\underline{n} + 1)^{-1} < 1$}
\end{ptree}
\noindent$\Pi_{20}$:
\begin{ptree}
\AXC{$x = x$\lbl{A5}}
\AXC{$x = y \to x = x \to y = x$\lbl{A6}}
\AXC{$(x = y \to x = x \to y = x) \to x = x \to x = y \to y = x$\lbl{A1}}
\BIC{$x = x \to x = y \to y = x$}
\BIC{$x = y \to y = x$}
\UIC{$\forall y\,(x = y \to y = x)$}
\UIC{$\forall x\,\forall y\,(x = y \to y = x)$}
\end{ptree}
\noindent$\Pi_{21}$:
\begin{ptree}
\AXC{$\forall n \in \mathbb{N}\,(0 < \underline{n})$\lbl{[N15]}}
\AXC{$\forall n \in \mathbb{N}\,(0 < \underline{n}) \to n \in \mathbb{N} \to 0 < \underline{n}$\lbl{A2}}
\BIC{$n \in \mathbb{N} \to 0 < \underline{n}$}
\AXC{$(n \in \mathbb{N} \to 0 < \underline{n}) \to \varphi_{153}$\lbl{A1}}
\BIC{$\varphi_{153}$}
\end{ptree}
\noindent$\Pi_{22}$:
\begin{ptree}
\AXC{$\forall x\,\forall y\,(x \in \mathbb{R} \to y \in \mathbb{R} \to 0 < x \to 0 < y \to 0 < x + y)$\lbl{[E5]}}
\AXC{$\varphi_{186}$\lbl{A2}}
\BIC{$\forall y\,(\underline{n} \in \mathbb{R} \to y \in \mathbb{R} \to 0 < \underline{n} \to 0 < y \to 0 < \underline{n} + y)$}
\AXC{$\varphi_{192}$\lbl{A2}}
\BIC{$\underline{n} \in \mathbb{R} \to 1 \in \mathbb{R} \to 0 < \underline{n} \to 0 < 1 \to 0 < \underline{n} + 1$}
\AXC{$\varphi_{275}$\lbl{A1}}
\BIC{$(c \in \mathbb{R} \to r \in \mathbb{R} \to 0 < r \to n \in \mathbb{N} \to \underline{n} \in \mathbb{R}) \to \varphi_{215}$}
\end{ptree}
\noindent$\Pi_{23}$:
\begin{ptree}
\AXC{$1 \in \mathbb{R}$\lbl{(R1)}}
\AXC{$\stackrel{\Pi_{5}}{\vdots}$}\noLine
\UIC{$c \in \mathbb{R} \to r \in \mathbb{R} \to 0 < r \to n \in \mathbb{N} \to \underline{n} \in \mathbb{R}$}
\AXC{$\stackrel{\Pi_{22}}{\vdots}$}\noLine
\UIC{$(c \in \mathbb{R} \to r \in \mathbb{R} \to 0 < r \to n \in \mathbb{N} \to \underline{n} \in \mathbb{R}) \to \varphi_{215}$}
\BIC{$\varphi_{215}$}
\BIC{$(c \in \mathbb{R} \to r \in \mathbb{R} \to 0 < r \to n \in \mathbb{N} \to 0 < \underline{n}) \to \varphi_{156}$}
\end{ptree}
\noindent$\Pi_{24}$:
\begin{ptree}
\AXC{$c \in \mathbb{R} \to r \in \mathbb{R} \to 0 < r \to n \in \mathbb{N} \to n \in \mathbb{N}$\lbl{A1}}
\AXC{$\stackrel{\Pi_{21}}{\vdots}$}\noLine
\UIC{$\varphi_{153}$}
\BIC{$c \in \mathbb{R} \to r \in \mathbb{R} \to 0 < r \to n \in \mathbb{N} \to 0 < \underline{n}$}
\AXC{$\stackrel{\Pi_{23}}{\vdots}$}\noLine
\UIC{$(c \in \mathbb{R} \to r \in \mathbb{R} \to 0 < r \to n \in \mathbb{N} \to 0 < \underline{n}) \to \varphi_{156}$}
\BIC{$\varphi_{156}$}
\end{ptree}
\noindent$\Pi_{25}$:
\begin{ptree}
\AXC{$0 < 1$\lbl{[F28]}}
\AXC{$0 < 1 \to c \in \mathbb{R} \to r \in \mathbb{R} \to 0 < r \to n \in \mathbb{N} \to 0 < 1$\lbl{A1}}
\BIC{$c \in \mathbb{R} \to r \in \mathbb{R} \to 0 < r \to n \in \mathbb{N} \to 0 < 1$}
\AXC{$\stackrel{\Pi_{24}}{\vdots}$}\noLine
\UIC{$\varphi_{156}$}
\BIC{$c \in \mathbb{R} \to r \in \mathbb{R} \to 0 < r \to n \in \mathbb{N} \to 0 < \underline{n} + 1$}
\end{ptree}
\noindent$\Pi_{26}$:
\begin{ptree}
\AXC{$\forall x\,\forall y\,(x < y \to x \neq y)$\lbl{[F26]}}
\AXC{$\forall x\,\forall y\,(x < y \to x \neq y) \to \forall y\,(0 < y \to 0 \neq y)$\lbl{A2}}
\BIC{$\forall y\,(0 < y \to 0 \neq y)$}
\end{ptree}
\noindent$\Pi_{27}$:
\begin{ptree}
\AXC{$\stackrel{\Pi_{26}}{\vdots}$}\noLine
\UIC{$\forall y\,(0 < y \to 0 \neq y)$}
\AXC{$\forall y\,(0 < y \to 0 \neq y) \to 0 < \underline{n} + 1 \to 0 \neq \underline{n} + 1$\lbl{A2}}
\BIC{$0 < \underline{n} + 1 \to 0 \neq \underline{n} + 1$}
\AXC{$(0 < \underline{n} + 1 \to 0 \neq \underline{n} + 1) \to \varphi_{175}$\lbl{A1}}
\BIC{$\varphi_{175}$}
\end{ptree}
\noindent$\Pi_{28}$:
\begin{ptree}
\AXC{$\stackrel{\Pi_{25}}{\vdots}$}\noLine
\UIC{$c \in \mathbb{R} \to r \in \mathbb{R} \to 0 < r \to n \in \mathbb{N} \to 0 < \underline{n} + 1$}
\AXC{$\stackrel{\Pi_{27}}{\vdots}$}\noLine
\UIC{$\varphi_{175}$}
\BIC{$c \in \mathbb{R} \to r \in \mathbb{R} \to 0 < r \to n \in \mathbb{N} \to 0 \neq \underline{n} + 1$}
\AXC{$(c \in \mathbb{R} \to r \in \mathbb{R} \to 0 < r \to n \in \mathbb{N} \to 0 \neq \underline{n} + 1) \to \varphi_{106}$\lbl{A1}}
\BIC{$\varphi_{106}$}
\end{ptree}
\noindent$\Pi_{29}$:
\begin{ptree}
\AXC{$\stackrel{\Pi_{14}}{\vdots}$}\noLine
\UIC{$c \in \mathbb{R} \to r \in \mathbb{R} \to 0 < r \to n \in \mathbb{N} \to \underline{n} + 1 \in \mathbb{R}$}
\AXC{$\forall x \in \mathbb{R}\,(x \neq 0 \to x^{-1} \in \mathbb{R})$\lbl{(R1)}}
\AXC{$\varphi_{135}$\lbl{A2}}
\BIC{$\underline{n} + 1 \in \mathbb{R} \to \underline{n} + 1 \neq 0 \to (\underline{n} + 1)^{-1} \in \mathbb{R}$}
\AXC{$\varphi_{238}$\lbl{A1}}
\BIC{$(c \in \mathbb{R} \to r \in \mathbb{R} \to 0 < r \to n \in \mathbb{N} \to \underline{n} + 1 \in \mathbb{R}) \to \varphi_{185}$}
\BIC{$\varphi_{185}$}
\end{ptree}
\noindent$\Pi_{30}$:
\begin{ptree}
\AXC{$\stackrel{\Pi_{20}}{\vdots}$}\noLine
\UIC{$\forall x\,\forall y\,(x = y \to y = x)$}
\AXC{$\varphi_{1}$\lbl{A2}}
\BIC{$\forall y\,(\underline{n} + 1 = y \to y = \underline{n} + 1)$}
\AXC{$\varphi_{11}$\lbl{A2}}
\BIC{$\underline{n} + 1 = 0 \to 0 = \underline{n} + 1$}
\AXC{$\stackrel{\Pi_{28}}{\vdots}$}\noLine
\UIC{$\varphi_{106}$}
\BIC{$c \in \mathbb{R} \to r \in \mathbb{R} \to 0 < r \to n \in \mathbb{N} \to \underline{n} + 1 \neq 0$}
\AXC{$\stackrel{\Pi_{29}}{\vdots}$}\noLine
\UIC{$\varphi_{185}$}
\BIC{$c \in \mathbb{R} \to r \in \mathbb{R} \to 0 < r \to n \in \mathbb{N} \to (\underline{n} + 1)^{-1} \in \mathbb{R}$}
\end{ptree}
\noindent$\Pi_{31}$:
\begin{ptree}
\AXC{$\varphi_{58}$\lbl{[E8]}}
\AXC{$\varphi_{58} \to \varphi_{173}$\lbl{A2}}
\BIC{$\varphi_{173}$}
\AXC{$\varphi_{173} \to \varphi_{171}$\lbl{A2}}
\BIC{$\varphi_{171}$}
\AXC{$\varphi_{171} \to (\underline{n} + 1)^{-1} \in \mathbb{R} \to 1 \in \mathbb{R} \to \varphi_{47}$\lbl{A2}}
\BIC{$(\underline{n} + 1)^{-1} \in \mathbb{R} \to 1 \in \mathbb{R} \to \varphi_{47}$}
\end{ptree}
\noindent$\Pi_{32}$:
\begin{ptree}
\AXC{$\stackrel{\Pi_{31}}{\vdots}$}\noLine
\UIC{$(\underline{n} + 1)^{-1} \in \mathbb{R} \to 1 \in \mathbb{R} \to \varphi_{47}$}
\AXC{$((\underline{n} + 1)^{-1} \in \mathbb{R} \to 1 \in \mathbb{R} \to \varphi_{47}) \to \varphi_{278}$\lbl{A1}}
\BIC{$\varphi_{278}$}
\end{ptree}
\noindent$\Pi_{33}$:
\begin{ptree}
\AXC{$1 \in \mathbb{R}$\lbl{(R1)}}
\AXC{$\stackrel{\Pi_{30}}{\vdots}$}\noLine
\UIC{$c \in \mathbb{R} \to r \in \mathbb{R} \to 0 < r \to n \in \mathbb{N} \to (\underline{n} + 1)^{-1} \in \mathbb{R}$}
\AXC{$\stackrel{\Pi_{32}}{\vdots}$}\noLine
\UIC{$\varphi_{278}$}
\BIC{$1 \in \mathbb{R} \to \varphi_{224}$}
\BIC{$\varphi_{224}$}
\end{ptree}
\noindent$\Pi_{34}$:
\begin{ptree}
\AXC{$c \in \mathbb{R} \to r \in \mathbb{R} \to 0 < r \to 0 < r$\lbl{A1}}
\AXC{$c \in \mathbb{R} \to r \in \mathbb{R} \to 0 < r \to n \in \mathbb{N} \to r \in \mathbb{R}$\lbl{A1}}
\AXC{$\stackrel{\Pi_{33}}{\vdots}$}\noLine
\UIC{$\varphi_{224}$}
\BIC{$(c \in \mathbb{R} \to r \in \mathbb{R} \to 0 < r \to 0 < r) \to \varphi_{191}$}
\BIC{$\varphi_{191}$}
\end{ptree}
\noindent$\Pi_{35}$:
\begin{ptree}
\AXC{$\forall x \in \mathbb{R}\,(x \cdot 1 = x)$\lbl{(R8)}}
\AXC{$\forall x \in \mathbb{R}\,(x \cdot 1 = x) \to r \in \mathbb{R} \to r \cdot 1 = r$\lbl{A2}}
\BIC{$r \in \mathbb{R} \to r \cdot 1 = r$}
\AXC{$(r \in \mathbb{R} \to r \cdot 1 = r) \to \varphi_{145}$\lbl{A1}}
\BIC{$\varphi_{145}$}
\end{ptree}
\noindent$\Pi_{36}$:
\begin{ptree}
\AXC{$c \in \mathbb{R} \to r \in \mathbb{R} \to 0 < r \to n \in \mathbb{N} \to r \in \mathbb{R}$\lbl{A1}}
\AXC{$\stackrel{\Pi_{35}}{\vdots}$}\noLine
\UIC{$\varphi_{145}$}
\BIC{$c \in \mathbb{R} \to r \in \mathbb{R} \to 0 < r \to n \in \mathbb{N} \to r \cdot 1 = r$}
\AXC{$\varphi_{247}$\lbl{A1}}
\BIC{$\varphi_{46} \to \varphi_{21} \to \varphi_{27}$}
\end{ptree}
\noindent$\Pi_{37}$:
\begin{ptree}
\AXC{$\stackrel{\Pi_{20}}{\vdots}$}\noLine
\UIC{$\forall x\,\forall y\,(x = y \to y = x)$}
\AXC{$\varphi_{59}$\lbl{A2}}
\BIC{$\forall y\,(\lvert \pi_{c,r}(n) - c \rvert = y \to y = \lvert \pi_{c,r}(n) - c \rvert)$}
\AXC{$\forall y\,(\lvert \pi_{c,r}(n) - c \rvert = y \to y = \lvert \pi_{c,r}(n) - c \rvert) \to \varphi_{123}$\lbl{A2}}
\BIC{$\varphi_{123}$}
\end{ptree}
\noindent$\Pi_{38}$:
\begin{ptree}
\AXC{$y = z \to x = y \to x = z$\lbl{A6}}
\AXC{$(y = z \to x = y \to x = z) \to x = y \to y = z \to x = z$\lbl{A1}}
\BIC{$x = y \to y = z \to x = z$}
\UIC{$\forall z\,(x = y \to y = z \to x = z)$}
\UIC{$\forall y\,\forall z\,(x = y \to y = z \to x = z)$}
\UIC{$\forall x\,\forall y\,\forall z\,(x = y \to y = z \to x = z)$}
\AXC{$\forall x\,\forall y\,\forall z\,(x = y \to y = z \to x = z) \to \varphi_{22}$\lbl{A2}}
\BIC{$\varphi_{22}$}
\end{ptree}
\noindent$\Pi_{39}$:
\begin{ptree}
\AXC{$\stackrel{\Pi_{14}}{\vdots}$}\noLine
\UIC{$c \in \mathbb{R} \to r \in \mathbb{R} \to 0 < r \to n \in \mathbb{N} \to \underline{n} + 1 \in \mathbb{R}$}
\AXC{$\forall x \in \mathbb{R}\,(0 < x \to 0 < x^{-1})$\lbl{[F34]}}
\AXC{$\varphi_{128}$\lbl{A2}}
\BIC{$\underline{n} + 1 \in \mathbb{R} \to 0 < \underline{n} + 1 \to 0 < (\underline{n} + 1)^{-1}$}
\AXC{$\varphi_{235}$\lbl{A1}}
\BIC{$(c \in \mathbb{R} \to r \in \mathbb{R} \to 0 < r \to n \in \mathbb{N} \to \underline{n} + 1 \in \mathbb{R}) \to \varphi_{182}$}
\BIC{$\varphi_{182}$}
\end{ptree}
\noindent$\Pi_{40}$:
\begin{ptree}
\AXC{$\forall x\,\forall y\,(x \in \mathbb{R} \to y \in \mathbb{R} \to 0 < x \to 0 < y \to 0 < x \cdot y)$\lbl{(R15)}}
\AXC{$\varphi_{165}$\lbl{A2}}
\BIC{$\forall y\,(r \in \mathbb{R} \to y \in \mathbb{R} \to 0 < r \to 0 < y \to 0 < r \cdot y)$}
\AXC{$\varphi_{204}$\lbl{A2}}
\BIC{$r \in \mathbb{R} \to \varphi_{95}$}
\AXC{$\varphi_{289}$\lbl{A1}}
\BIC{$(c \in \mathbb{R} \to r \in \mathbb{R} \to 0 < r \to n \in \mathbb{N} \to r \in \mathbb{R}) \to \varphi_{231}$}
\end{ptree}
\noindent$\Pi_{41}$:
\begin{ptree}
\AXC{$c \in \mathbb{R} \to r \in \mathbb{R} \to 0 < r \to n \in \mathbb{N} \to r \in \mathbb{R}$\lbl{A1}}
\AXC{$\stackrel{\Pi_{40}}{\vdots}$}\noLine
\UIC{$(c \in \mathbb{R} \to r \in \mathbb{R} \to 0 < r \to n \in \mathbb{N} \to r \in \mathbb{R}) \to \varphi_{231}$}
\BIC{$\varphi_{231}$}
\end{ptree}
\noindent$\Pi_{42}$:
\begin{ptree}
\AXC{$c \in \mathbb{R} \to r \in \mathbb{R} \to 0 < r \to 0 < r$\lbl{A1}}
\AXC{$\stackrel{\Pi_{30}}{\vdots}$}\noLine
\UIC{$c \in \mathbb{R} \to r \in \mathbb{R} \to 0 < r \to n \in \mathbb{N} \to (\underline{n} + 1)^{-1} \in \mathbb{R}$}
\AXC{$\stackrel{\Pi_{41}}{\vdots}$}\noLine
\UIC{$\varphi_{231}$}
\BIC{$(c \in \mathbb{R} \to r \in \mathbb{R} \to 0 < r \to 0 < r) \to \varphi_{190}$}
\BIC{$\varphi_{190}$}
\end{ptree}
\noindent$\Pi_{43}$:
\begin{ptree}
\AXC{$\stackrel{\Pi_{25}}{\vdots}$}\noLine
\UIC{$c \in \mathbb{R} \to r \in \mathbb{R} \to 0 < r \to n \in \mathbb{N} \to 0 < \underline{n} + 1$}
\AXC{$\stackrel{\Pi_{39}}{\vdots}$}\noLine
\UIC{$\varphi_{182}$}
\BIC{$c \in \mathbb{R} \to r \in \mathbb{R} \to 0 < r \to n \in \mathbb{N} \to 0 < (\underline{n} + 1)^{-1}$}
\AXC{$\stackrel{\Pi_{42}}{\vdots}$}\noLine
\UIC{$\varphi_{190}$}
\BIC{$\varphi_{28}$}
\end{ptree}
\noindent$\Pi_{44}$:
\begin{ptree}
\AXC{$\forall x\,(0 \leq x \to \lvert x \rvert = x)$\lbl{[A1]}}
\AXC{$\forall x\,(0 \leq x \to \lvert x \rvert = x) \to \varphi_{38}$\lbl{A2}}
\BIC{$\varphi_{38}$}
\AXC{$\varphi_{38} \to \varphi_{15} \to c \in \mathbb{R} \to r \in \mathbb{R} \to \varphi_{16}$\lbl{A1}}
\BIC{$\varphi_{15} \to c \in \mathbb{R} \to r \in \mathbb{R} \to \varphi_{16}$}
\end{ptree}
\noindent$\Pi_{45}$:
\begin{ptree}
\AXC{$z = w \to \lvert \pi_{c,r}(n) - c \rvert = \lvert z \rvert \to \lvert \pi_{c,r}(n) - c \rvert = \lvert w \rvert$\lbl{A6}}
\UIC{$\varphi_{52}$}
\UIC{$\forall z\,(\varphi_{52})$}
\AXC{$\forall z\,(\varphi_{52}) \to \varphi_{166}$\lbl{A2}}
\BIC{$\varphi_{166}$}
\AXC{$\varphi_{166} \to \pi_{c,r}(n) - c = r \cdot (\underline{n} + 1)^{-1} \to \varphi_{140}$\lbl{A2}}
\BIC{$\pi_{c,r}(n) - c = r \cdot (\underline{n} + 1)^{-1} \to \varphi_{140}$}
\end{ptree}
\noindent$\Pi_{46}$:
\begin{ptree}
\AXC{$y = z \to x = y \to x = z$\lbl{A6}}
\AXC{$(y = z \to x = y \to x = z) \to x = y \to y = z \to x = z$\lbl{A1}}
\BIC{$x = y \to y = z \to x = z$}
\UIC{$\forall z\,(x = y \to y = z \to x = z)$}
\UIC{$\forall y\,\forall z\,(x = y \to y = z \to x = z)$}
\UIC{$\forall x\,\forall y\,\forall z\,(x = y \to y = z \to x = z)$}
\AXC{$\varphi_{91}$\lbl{A2}}
\BIC{$\forall y\,\forall z\,(\pi_{c,r}(n) - c = y \to y = z \to \pi_{c,r}(n) - c = z)$}
\end{ptree}
\noindent$\Pi_{47}$:
\begin{ptree}
\AXC{$\stackrel{\Pi_{46}}{\vdots}$}\noLine
\UIC{$\forall y\,\forall z\,(\pi_{c,r}(n) - c = y \to y = z \to \pi_{c,r}(n) - c = z)$}
\AXC{$\forall y\,\forall z\,(\pi_{c,r}(n) - c = y \to y = z \to \pi_{c,r}(n) - c = z) \to \varphi_{167}$\lbl{A2}}
\BIC{$\varphi_{167}$}
\end{ptree}
\noindent$\Pi_{48}$:
\begin{ptree}
\AXC{$\forall x\,\forall y\,(x \in \mathbb{R} \to y \in \mathbb{R} \to x \cdot y \in \mathbb{R})$\lbl{(R1)}}
\AXC{$\varphi_{103}$\lbl{A2}}
\BIC{$\forall y\,(r \in \mathbb{R} \to y \in \mathbb{R} \to r \cdot y \in \mathbb{R})$}
\AXC{$\varphi_{157}$\lbl{A2}}
\BIC{$r \in \mathbb{R} \to (\underline{n} + 1)^{-1} \in \mathbb{R} \to r \cdot (\underline{n} + 1)^{-1} \in \mathbb{R}$}
\AXC{$\varphi_{248}$\lbl{A1}}
\BIC{$(c \in \mathbb{R} \to r \in \mathbb{R} \to 0 < r \to n \in \mathbb{N} \to r \in \mathbb{R}) \to \varphi_{194}$}
\end{ptree}
\noindent$\Pi_{49}$:
\begin{ptree}
\AXC{$c \in \mathbb{R} \to r \in \mathbb{R} \to 0 < r \to n \in \mathbb{N} \to r \in \mathbb{R}$\lbl{A1}}
\AXC{$\stackrel{\Pi_{48}}{\vdots}$}\noLine
\UIC{$(c \in \mathbb{R} \to r \in \mathbb{R} \to 0 < r \to n \in \mathbb{N} \to r \in \mathbb{R}) \to \varphi_{194}$}
\BIC{$\varphi_{194}$}
\end{ptree}
\noindent$\Pi_{50}$:
\begin{ptree}
\AXC{$\stackrel{\Pi_{30}}{\vdots}$}\noLine
\UIC{$c \in \mathbb{R} \to r \in \mathbb{R} \to 0 < r \to n \in \mathbb{N} \to (\underline{n} + 1)^{-1} \in \mathbb{R}$}
\AXC{$\stackrel{\Pi_{49}}{\vdots}$}\noLine
\UIC{$\varphi_{194}$}
\BIC{$\varphi_{43}$}
\end{ptree}
\noindent$\Pi_{51}$:
\begin{ptree}
\AXC{$\forall x\,\forall y\,(x \in \mathbb{R} \to y \in \mathbb{R} \to x + y - y = x)$\lbl{[E38]}}
\AXC{$\forall x\,\forall y\,(x \in \mathbb{R} \to y \in \mathbb{R} \to x + y - y = x) \to \varphi_{129}$\lbl{A2}}
\BIC{$\varphi_{129}$}
\AXC{$\varphi_{129} \to \varphi_{119}$\lbl{A2}}
\BIC{$\varphi_{119}$}
\end{ptree}
\noindent$\Pi_{52}$:
\begin{ptree}
\AXC{$y = z \to x = y \to x = z$\lbl{A6}}
\AXC{$(y = z \to x = y \to x = z) \to x = y \to y = z \to x = z$\lbl{A1}}
\BIC{$x = y \to y = z \to x = z$}
\UIC{$\forall z\,(x = y \to y = z \to x = z)$}
\UIC{$\forall y\,\forall z\,(x = y \to y = z \to x = z)$}
\UIC{$\forall x\,\forall y\,\forall z\,(x = y \to y = z \to x = z)$}
\AXC{$\varphi_{91}$\lbl{A2}}
\BIC{$\forall y\,\forall z\,(\pi_{c,r}(n) - c = y \to y = z \to \pi_{c,r}(n) - c = z)$}
\end{ptree}
\noindent$\Pi_{53}$:
\begin{ptree}
\AXC{$\stackrel{\Pi_{52}}{\vdots}$}\noLine
\UIC{$\forall y\,\forall z\,(\pi_{c,r}(n) - c = y \to y = z \to \pi_{c,r}(n) - c = z)$}
\AXC{$\forall y\,\forall z\,(\pi_{c,r}(n) - c = y \to y = z \to \pi_{c,r}(n) - c = z) \to \varphi_{168}$\lbl{A2}}
\BIC{$\varphi_{168}$}
\end{ptree}
\noindent$\Pi_{54}$:
\begin{ptree}
\AXC{$\forall x\,\forall y\,(x \in \mathbb{R} \to y \in \mathbb{R} \to x + y = y + x)$\lbl{(R3)}}
\AXC{$\varphi_{92}$\lbl{A2}}
\BIC{$\forall y\,(c \in \mathbb{R} \to y \in \mathbb{R} \to c + y = y + c)$}
\AXC{$\forall y\,(c \in \mathbb{R} \to y \in \mathbb{R} \to c + y = y + c) \to c \in \mathbb{R} \to \varphi_{84}$\lbl{A2}}
\BIC{$c \in \mathbb{R} \to \varphi_{84}$}
\end{ptree}
\noindent$\Pi_{55}$:
\begin{ptree}
\AXC{$\stackrel{\Pi_{50}}{\vdots}$}\noLine
\UIC{$\varphi_{43}$}
\AXC{$c \in \mathbb{R} \to r \in \mathbb{R} \to 0 < r \to n \in \mathbb{N} \to c \in \mathbb{R}$\lbl{A1}}
\AXC{$\stackrel{\Pi_{54}}{\vdots}$}\noLine
\UIC{$c \in \mathbb{R} \to \varphi_{84}$}
\AXC{$(c \in \mathbb{R} \to \varphi_{84}) \to \varphi_{223}$\lbl{A1}}
\BIC{$\varphi_{223}$}
\BIC{$\varphi_{43} \to c \in \mathbb{R} \to r \in \mathbb{R} \to \varphi_{35}$}
\BIC{$c \in \mathbb{R} \to r \in \mathbb{R} \to \varphi_{35}$}
\end{ptree}
\noindent$\Pi_{56}$:
\begin{ptree}
\AXC{$x = x$\lbl{A5}}
\UIC{$\forall x\,(x = x)$}
\AXC{$\varphi_{29}$\lbl{A2}}
\BIC{$c + r \cdot (\underline{n} + 1)^{-1} - c = c + r \cdot (\underline{n} + 1)^{-1} - c$}
\AXC{$\stackrel{\Pi_{55}}{\vdots}$}\noLine
\UIC{$c \in \mathbb{R} \to r \in \mathbb{R} \to \varphi_{35}$}
\AXC{$(c \in \mathbb{R} \to r \in \mathbb{R} \to \varphi_{35}) \to \varphi_{284}$\lbl{A1}}
\BIC{$\varphi_{284}$}
\BIC{$\varphi_{259}$}
\end{ptree}
\noindent$\Pi_{57}$:
\begin{ptree}
\AXC{$z = w \to \varphi_{56}$\lbl{A6}}
\UIC{$\forall w\,(z = w \to \varphi_{56})$}
\UIC{$\forall z\,\forall w\,(z = w \to \varphi_{56})$}
\AXC{$\forall z\,\forall w\,(z = w \to \varphi_{56}) \to \varphi_{187}$\lbl{A2}}
\BIC{$\varphi_{187}$}
\AXC{$\varphi_{187} \to c + r \cdot (\underline{n} + 1)^{-1} = r \cdot (\underline{n} + 1)^{-1} + c \to \varphi_{180}$\lbl{A2}}
\BIC{$c + r \cdot (\underline{n} + 1)^{-1} = r \cdot (\underline{n} + 1)^{-1} + c \to \varphi_{180}$}
\AXC{$\stackrel{\Pi_{56}}{\vdots}$}\noLine
\UIC{$\varphi_{259}$}
\BIC{$c \in \mathbb{R} \to r \in \mathbb{R} \to 0 < r \to \varphi_{39}$}
\end{ptree}
\noindent$\Pi_{58}$:
\begin{ptree}
\AXC{$z = w \to \pi_{c,r}(n) - c = z - c \to \pi_{c,r}(n) - c = w - c$\lbl{A6}}
\UIC{$\forall w\,(z = w \to \pi_{c,r}(n) - c = z - c \to \pi_{c,r}(n) - c = w - c)$}
\UIC{$\varphi_{23}$}
\AXC{$\varphi_{23} \to \varphi_{120}$\lbl{A2}}
\BIC{$\varphi_{120}$}
\AXC{$\varphi_{120} \to \pi_{c,r}(n) = c + r \cdot (\underline{n} + 1)^{-1} \to \varphi_{126}$\lbl{A2}}
\BIC{$\pi_{c,r}(n) = c + r \cdot (\underline{n} + 1)^{-1} \to \varphi_{126}$}
\end{ptree}
\noindent$\Pi_{59}$:
\begin{ptree}
\AXC{$\forall c\,\forall r\,(c \in \mathbb{R} \to r \in \mathbb{R} \to \varphi_{45})$\lbl{[S22]}}
\AXC{$\varphi_{222}$\lbl{A2}}
\BIC{$\forall r\,(c \in \mathbb{R} \to r \in \mathbb{R} \to \varphi_{45})$}
\AXC{$\forall r\,(c \in \mathbb{R} \to r \in \mathbb{R} \to \varphi_{45}) \to c \in \mathbb{R} \to r \in \mathbb{R} \to \varphi_{45}$\lbl{A2}}
\BIC{$c \in \mathbb{R} \to r \in \mathbb{R} \to \varphi_{45}$}
\AXC{$\varphi_{264}$\lbl{A1}}
\BIC{$(c \in \mathbb{R} \to r \in \mathbb{R} \to 0 < r \to c \in \mathbb{R}) \to \varphi_{201}$}
\end{ptree}
\noindent$\Pi_{60}$:
\begin{ptree}
\AXC{$c \in \mathbb{R} \to r \in \mathbb{R} \to 0 < r \to r \in \mathbb{R}$\lbl{A1}}
\AXC{$c \in \mathbb{R} \to r \in \mathbb{R} \to 0 < r \to c \in \mathbb{R}$\lbl{A1}}
\AXC{$\stackrel{\Pi_{59}}{\vdots}$}\noLine
\UIC{$(c \in \mathbb{R} \to r \in \mathbb{R} \to 0 < r \to c \in \mathbb{R}) \to \varphi_{201}$}
\BIC{$\varphi_{201}$}
\BIC{$c \in \mathbb{R} \to r \in \mathbb{R} \to 0 < r \to \varphi_{45}$}
\end{ptree}
\noindent$\Pi_{61}$:
\begin{ptree}
\AXC{$\stackrel{\Pi_{60}}{\vdots}$}\noLine
\UIC{$c \in \mathbb{R} \to r \in \mathbb{R} \to 0 < r \to \varphi_{45}$}
\AXC{$(c \in \mathbb{R} \to r \in \mathbb{R} \to 0 < r \to \varphi_{45}) \to c \in \mathbb{R} \to \varphi_{60}$\lbl{A1}}
\BIC{$c \in \mathbb{R} \to \varphi_{60}$}
\end{ptree}
\noindent$\Pi_{62}$:
\begin{ptree}
\AXC{$\varphi_{158}$\lbl{A2}}
\AXC{$\stackrel{\Pi_{61}}{\vdots}$}\noLine
\UIC{$c \in \mathbb{R} \to \varphi_{60}$}
\AXC{$(c \in \mathbb{R} \to \varphi_{60}) \to \varphi_{158} \to c \in \mathbb{R} \to \varphi_{48}$\lbl{A1}}
\BIC{$\varphi_{158} \to c \in \mathbb{R} \to \varphi_{48}$}
\BIC{$c \in \mathbb{R} \to \varphi_{48}$}
\AXC{$(c \in \mathbb{R} \to \varphi_{48}) \to \varphi_{256}$\lbl{A1}}
\BIC{$\varphi_{256}$}
\end{ptree}
\noindent$\Pi_{63}$:
\begin{ptree}
\AXC{$\stackrel{\Pi_{58}}{\vdots}$}\noLine
\UIC{$\pi_{c,r}(n) = c + r \cdot (\underline{n} + 1)^{-1} \to \varphi_{126}$}
\AXC{$x = x$\lbl{A5}}
\UIC{$\forall x\,(x = x)$}
\AXC{$\forall x\,(x = x) \to \pi_{c,r}(n) - c = \pi_{c,r}(n) - c$\lbl{A2}}
\BIC{$\pi_{c,r}(n) - c = \pi_{c,r}(n) - c$}
\AXC{$\stackrel{\Pi_{62}}{\vdots}$}\noLine
\UIC{$\varphi_{256}$}
\BIC{$(\pi_{c,r}(n) = c + r \cdot (\underline{n} + 1)^{-1} \to \varphi_{126}) \to c \in \mathbb{R} \to \varphi_{85}$}
\BIC{$c \in \mathbb{R} \to \varphi_{85}$}
\end{ptree}
\noindent$\Pi_{64}$:
\begin{ptree}
\AXC{$\stackrel{\Pi_{57}}{\vdots}$}\noLine
\UIC{$c \in \mathbb{R} \to r \in \mathbb{R} \to 0 < r \to \varphi_{39}$}
\AXC{$\stackrel{\Pi_{63}}{\vdots}$}\noLine
\UIC{$c \in \mathbb{R} \to \varphi_{85}$}
\AXC{$(c \in \mathbb{R} \to \varphi_{85}) \to (c \in \mathbb{R} \to r \in \mathbb{R} \to 0 < r \to \varphi_{39}) \to \varphi_{243}$\lbl{A1}}
\BIC{$(c \in \mathbb{R} \to r \in \mathbb{R} \to 0 < r \to \varphi_{39}) \to \varphi_{243}$}
\BIC{$\varphi_{243}$}
\end{ptree}
\noindent$\Pi_{65}$:
\begin{ptree}
\AXC{$\stackrel{\Pi_{53}}{\vdots}$}\noLine
\UIC{$\varphi_{168}$}
\AXC{$\varphi_{168} \to \pi_{c,r}(n) - c = c + r \cdot (\underline{n} + 1)^{-1} - c \to \varphi_{170}$\lbl{A2}}
\BIC{$\pi_{c,r}(n) - c = c + r \cdot (\underline{n} + 1)^{-1} - c \to \varphi_{170}$}
\AXC{$\stackrel{\Pi_{64}}{\vdots}$}\noLine
\UIC{$\varphi_{243}$}
\BIC{$c \in \mathbb{R} \to \varphi_{86}$}
\AXC{$(c \in \mathbb{R} \to \varphi_{86}) \to (c \in \mathbb{R} \to r \in \mathbb{R} \to \varphi_{34}) \to \varphi_{225}$\lbl{A1}}
\BIC{$(c \in \mathbb{R} \to r \in \mathbb{R} \to \varphi_{34}) \to \varphi_{225}$}
\end{ptree}
\noindent$\Pi_{66}$:
\begin{ptree}
\AXC{$c \in \mathbb{R} \to r \in \mathbb{R} \to 0 < r \to n \in \mathbb{N} \to c \in \mathbb{R}$\lbl{A1}}
\AXC{$\stackrel{\Pi_{50}}{\vdots}$}\noLine
\UIC{$\varphi_{43}$}
\AXC{$\stackrel{\Pi_{51}}{\vdots}$}\noLine
\UIC{$\varphi_{119}$}
\AXC{$\varphi_{119} \to \varphi_{43} \to \varphi_{202}$\lbl{A1}}
\BIC{$\varphi_{43} \to \varphi_{202}$}
\BIC{$\varphi_{202}$}
\BIC{$c \in \mathbb{R} \to r \in \mathbb{R} \to \varphi_{34}$}
\AXC{$\stackrel{\Pi_{65}}{\vdots}$}\noLine
\UIC{$(c \in \mathbb{R} \to r \in \mathbb{R} \to \varphi_{34}) \to \varphi_{225}$}
\BIC{$\varphi_{225}$}
\end{ptree}
\noindent$\Pi_{67}$:
\begin{ptree}
\AXC{$\stackrel{\Pi_{47}}{\vdots}$}\noLine
\UIC{$\varphi_{167}$}
\AXC{$\varphi_{167} \to \pi_{c,r}(n) - c = r \cdot (\underline{n} + 1)^{-1} + c - c \to \varphi_{125}$\lbl{A2}}
\BIC{$\pi_{c,r}(n) - c = r \cdot (\underline{n} + 1)^{-1} + c - c \to \varphi_{125}$}
\AXC{$\stackrel{\Pi_{66}}{\vdots}$}\noLine
\UIC{$\varphi_{225}$}
\BIC{$c \in \mathbb{R} \to \varphi_{49}$}
\AXC{$(c \in \mathbb{R} \to \varphi_{49}) \to \varphi_{267}$\lbl{A1}}
\BIC{$\varphi_{267}$}
\end{ptree}
\noindent$\Pi_{68}$:
\begin{ptree}
\AXC{$\stackrel{\Pi_{45}}{\vdots}$}\noLine
\UIC{$\pi_{c,r}(n) - c = r \cdot (\underline{n} + 1)^{-1} \to \varphi_{140}$}
\AXC{$x = x$\lbl{A5}}
\UIC{$\forall x\,(x = x)$}
\AXC{$\forall x\,(x = x) \to \lvert \pi_{c,r}(n) - c \rvert = \lvert \pi_{c,r}(n) - c \rvert$\lbl{A2}}
\BIC{$\lvert \pi_{c,r}(n) - c \rvert = \lvert \pi_{c,r}(n) - c \rvert$}
\AXC{$\stackrel{\Pi_{67}}{\vdots}$}\noLine
\UIC{$\varphi_{267}$}
\BIC{$(\pi_{c,r}(n) - c = r \cdot (\underline{n} + 1)^{-1} \to \varphi_{140}) \to c \in \mathbb{R} \to \varphi_{87}$}
\BIC{$c \in \mathbb{R} \to \varphi_{87}$}
\end{ptree}
\noindent$\Pi_{69}$:
\begin{ptree}
\AXC{$\stackrel{\Pi_{43}}{\vdots}$}\noLine
\UIC{$\varphi_{28}$}
\AXC{$\varphi_{28} \to \varphi_{15}$\lbl{A1}}
\BIC{$\varphi_{15}$}
\AXC{$\stackrel{\Pi_{44}}{\vdots}$}\noLine
\UIC{$\varphi_{15} \to c \in \mathbb{R} \to r \in \mathbb{R} \to \varphi_{16}$}
\BIC{$c \in \mathbb{R} \to r \in \mathbb{R} \to \varphi_{16}$}
\AXC{$\stackrel{\Pi_{68}}{\vdots}$}\noLine
\UIC{$c \in \mathbb{R} \to \varphi_{87}$}
\AXC{$(c \in \mathbb{R} \to \varphi_{87}) \to (c \in \mathbb{R} \to r \in \mathbb{R} \to \varphi_{16}) \to \varphi_{227}$\lbl{A1}}
\BIC{$(c \in \mathbb{R} \to r \in \mathbb{R} \to \varphi_{16}) \to \varphi_{227}$}
\BIC{$\varphi_{227}$}
\end{ptree}
\noindent$\Pi_{70}$:
\begin{ptree}
\AXC{$\stackrel{\Pi_{38}}{\vdots}$}\noLine
\UIC{$\varphi_{22}$}
\AXC{$\varphi_{22} \to \varphi_{164}$\lbl{A2}}
\BIC{$\varphi_{164}$}
\AXC{$\varphi_{164} \to \lvert \pi_{c,r}(n) - c \rvert = \lvert r \cdot (\underline{n} + 1)^{-1} \rvert \to \varphi_{124}$\lbl{A2}}
\BIC{$\lvert \pi_{c,r}(n) - c \rvert = \lvert r \cdot (\underline{n} + 1)^{-1} \rvert \to \varphi_{124}$}
\AXC{$\stackrel{\Pi_{69}}{\vdots}$}\noLine
\UIC{$\varphi_{227}$}
\BIC{$c \in \mathbb{R} \to \varphi_{66}$}
\end{ptree}
\noindent$\Pi_{71}$:
\begin{ptree}
\AXC{$\stackrel{\Pi_{37}}{\vdots}$}\noLine
\UIC{$\varphi_{123}$}
\AXC{$\stackrel{\Pi_{70}}{\vdots}$}\noLine
\UIC{$c \in \mathbb{R} \to \varphi_{66}$}
\AXC{$(c \in \mathbb{R} \to \varphi_{66}) \to \varphi_{123} \to c \in \mathbb{R} \to \varphi_{67}$\lbl{A1}}
\BIC{$\varphi_{123} \to c \in \mathbb{R} \to \varphi_{67}$}
\BIC{$c \in \mathbb{R} \to \varphi_{67}$}
\AXC{$(c \in \mathbb{R} \to \varphi_{67}) \to \varphi_{27} \to \varphi_{218}$\lbl{A1}}
\BIC{$\varphi_{27} \to \varphi_{218}$}
\end{ptree}
\noindent$\Pi_{72}$:
\begin{ptree}
\AXC{$\stackrel{\Pi_{1}}{\vdots}$}\noLine
\UIC{$\varphi_{21}$}
\AXC{$\stackrel{\Pi_{19}}{\vdots}$}\noLine
\UIC{$c \in \mathbb{R} \to r \in \mathbb{R} \to 0 < r \to n \in \mathbb{N} \to (\underline{n} + 1)^{-1} < 1$}
\AXC{$\stackrel{\Pi_{34}}{\vdots}$}\noLine
\UIC{$\varphi_{191}$}
\BIC{$\varphi_{46}$}
\AXC{$\stackrel{\Pi_{36}}{\vdots}$}\noLine
\UIC{$\varphi_{46} \to \varphi_{21} \to \varphi_{27}$}
\BIC{$\varphi_{21} \to \varphi_{27}$}
\BIC{$\varphi_{27}$}
\AXC{$\stackrel{\Pi_{71}}{\vdots}$}\noLine
\UIC{$\varphi_{27} \to \varphi_{218}$}
\BIC{$\varphi_{218}$}
\end{ptree}
\noindent$\Pi_{73}$:
\begin{ptree}
\AXC{$\stackrel{\Pi_{20}}{\vdots}$}\noLine
\UIC{$\forall x\,\forall y\,(x = y \to y = x)$}
\AXC{$\varphi_{59}$\lbl{A2}}
\BIC{$\forall y\,(\lvert \pi_{c,r}(n) - c \rvert = y \to y = \lvert \pi_{c,r}(n) - c \rvert)$}
\AXC{$\forall y\,(\lvert \pi_{c,r}(n) - c \rvert = y \to y = \lvert \pi_{c,r}(n) - c \rvert) \to \varphi_{123}$\lbl{A2}}
\BIC{$\varphi_{123}$}
\end{ptree}
\noindent$\Pi_{74}$:
\begin{ptree}
\AXC{$\stackrel{\Pi_{43}}{\vdots}$}\noLine
\UIC{$\varphi_{28}$}
\AXC{$\stackrel{\Pi_{73}}{\vdots}$}\noLine
\UIC{$\varphi_{123}$}
\AXC{$\stackrel{\Pi_{70}}{\vdots}$}\noLine
\UIC{$c \in \mathbb{R} \to \varphi_{66}$}
\AXC{$(c \in \mathbb{R} \to \varphi_{66}) \to \varphi_{123} \to c \in \mathbb{R} \to \varphi_{67}$\lbl{A1}}
\BIC{$\varphi_{123} \to c \in \mathbb{R} \to \varphi_{67}$}
\BIC{$c \in \mathbb{R} \to \varphi_{67}$}
\AXC{$(c \in \mathbb{R} \to \varphi_{67}) \to \varphi_{28} \to \varphi_{219}$\lbl{A1}}
\BIC{$\varphi_{28} \to \varphi_{219}$}
\BIC{$\varphi_{219}$}
\end{ptree}
\noindent$\Pi_{75}$:
\begin{ptree}
\AXC{$z = w \to 0 < z \to 0 < w$\lbl{A6}}
\UIC{$\forall w\,(z = w \to 0 < z \to 0 < w)$}
\UIC{$\forall z\,\forall w\,(z = w \to 0 < z \to 0 < w)$}
\AXC{$\forall z\,\forall w\,(z = w \to 0 < z \to 0 < w) \to \varphi_{10}$\lbl{A2}}
\BIC{$\varphi_{10}$}
\AXC{$\varphi_{10} \to \varphi_{146}$\lbl{A2}}
\BIC{$\varphi_{146}$}
\AXC{$\stackrel{\Pi_{74}}{\vdots}$}\noLine
\UIC{$\varphi_{219}$}
\BIC{$c \in \mathbb{R} \to r \in \mathbb{R} \to 0 < r \to n \in \mathbb{N} \to 0 < \lvert \pi_{c,r}(n) - c \rvert$}
\AXC{$\varphi_{237}$\lbl{A1}}
\BIC{$\varphi_{216}$}
\end{ptree}
\noindent$\Pi_{76}$:
\begin{ptree}
\AXC{$z = w \to z < r \to w < r$\lbl{A6}}
\UIC{$\forall w\,(z = w \to z < r \to w < r)$}
\UIC{$\forall z\,\forall w\,(z = w \to z < r \to w < r)$}
\AXC{$\forall z\,\forall w\,(z = w \to z < r \to w < r) \to \varphi_{6}$\lbl{A2}}
\BIC{$\varphi_{6}$}
\AXC{$\varphi_{6} \to \varphi_{144}$\lbl{A2}}
\BIC{$\varphi_{144}$}
\AXC{$\stackrel{\Pi_{72}}{\vdots}$}\noLine
\UIC{$\varphi_{218}$}
\BIC{$c \in \mathbb{R} \to r \in \mathbb{R} \to 0 < r \to n \in \mathbb{N} \to \lvert \pi_{c,r}(n) - c \rvert < r$}
\AXC{$\stackrel{\Pi_{75}}{\vdots}$}\noLine
\UIC{$\varphi_{216}$}
\BIC{$c \in \mathbb{R} \to \varphi_{116}$}
\end{ptree}
\noindent$\Pi_{77}$:
\begin{ptree}
\AXC{$\forall n\,(\varphi_{142}) \to \varphi_{161}$\lbl{A3}}
\AXC{$\stackrel{\Pi_{76}}{\vdots}$}\noLine
\UIC{$c \in \mathbb{R} \to \varphi_{116}$}
\AXC{$(c \in \mathbb{R} \to \varphi_{116}) \to \varphi_{142}$\lbl{A1}}
\BIC{$\varphi_{142}$}
\UIC{$\forall n\,(\varphi_{142})$}
\AXC{$\forall n\,(\varphi_{142}) \to (\forall n\,(\varphi_{142}) \to \varphi_{161}) \to \varphi_{13} \to c \in \mathbb{R} \to r \in \mathbb{R} \to \varphi_{96}$\lbl{A1}}
\BIC{$(\forall n\,(\varphi_{142}) \to \varphi_{161}) \to \varphi_{13} \to c \in \mathbb{R} \to r \in \mathbb{R} \to \varphi_{96}$}
\BIC{$\varphi_{13} \to c \in \mathbb{R} \to r \in \mathbb{R} \to \varphi_{96}$}
\end{ptree}
\noindent$\Pi_{78}$:
\begin{ptree}
\AXC{$z = w \to z < \varepsilon \to w < \varepsilon$\lbl{A6}}
\UIC{$\forall w\,(z = w \to z < \varepsilon \to w < \varepsilon)$}
\UIC{$\forall z\,\forall w\,(z = w \to z < \varepsilon \to w < \varepsilon)$}
\AXC{$\forall z\,\forall w\,(z = w \to z < \varepsilon \to w < \varepsilon) \to \varphi_{14}$\lbl{A2}}
\BIC{$\varphi_{14}$}
\AXC{$\varphi_{14} \to \varphi_{155}$\lbl{A2}}
\BIC{$\varphi_{155}$}
\end{ptree}
\noindent$\Pi_{79}$:
\begin{ptree}
\AXC{$z = w \to r \cdot (\underline{n} + 1)^{-1} < z \to r \cdot (\underline{n} + 1)^{-1} < w$\lbl{A6}}
\UIC{$\varphi_{7}$}
\UIC{$\forall z\,(\varphi_{7})$}
\AXC{$\forall z\,(\varphi_{7}) \to \varphi_{115}$\lbl{A2}}
\BIC{$\varphi_{115}$}
\AXC{$\varphi_{115} \to r \cdot (\varepsilon \cdot r^{-1}) = \varepsilon \to \varphi_{36}$\lbl{A2}}
\BIC{$r \cdot (\varepsilon \cdot r^{-1}) = \varepsilon \to \varphi_{36}$}
\end{ptree}
\noindent$\Pi_{80}$:
\begin{ptree}
\AXC{$\forall K\,\forall n\,(K \in \mathbb{N} \to n \in \mathbb{N} \to K \leq n \to \underline{n}^{-1} \leq \underline{K}^{-1})$\lbl{[C25]}}
\AXC{$\varphi_{169}$\lbl{A2}}
\BIC{$\forall n\,(K \in \mathbb{N} \to n \in \mathbb{N} \to K \leq n \to \underline{n}^{-1} \leq \underline{K}^{-1})$}
\AXC{$\varphi_{148}$\lbl{A2}}
\BIC{$K \in \mathbb{N} \to n \in \mathbb{N} \to K \leq n \to \underline{n}^{-1} \leq \underline{K}^{-1}$}
\end{ptree}
\noindent$\Pi_{81}$:
\begin{ptree}
\AXC{$\stackrel{\Pi_{80}}{\vdots}$}\noLine
\UIC{$K \in \mathbb{N} \to n \in \mathbb{N} \to K \leq n \to \underline{n}^{-1} \leq \underline{K}^{-1}$}
\AXC{$(K \in \mathbb{N} \to n \in \mathbb{N} \to K \leq n \to \underline{n}^{-1} \leq \underline{K}^{-1}) \to \varphi_{302}$\lbl{A1}}
\BIC{$\varphi_{302}$}
\end{ptree}
\noindent$\Pi_{82}$:
\begin{ptree}
\AXC{$c \in \mathbb{R} \to r \in \mathbb{R} \to 0 < r \to \varphi_{98}$\lbl{A1}}
\AXC{$c \in \mathbb{R} \to r \in \mathbb{R} \to 0 < r \to \varepsilon \in \mathbb{R} \to 0 < \varepsilon \to \varphi_{72}$\lbl{A1}}
\AXC{$\varphi_{228}$\lbl{A1}}
\BIC{$c \in \mathbb{R} \to r \in \mathbb{R} \to \varphi_{80}$}
\AXC{$\stackrel{\Pi_{81}}{\vdots}$}\noLine
\UIC{$\varphi_{302}$}
\BIC{$(c \in \mathbb{R} \to r \in \mathbb{R} \to 0 < r \to \varphi_{98}) \to \varphi_{236}$}
\BIC{$\varphi_{236}$}
\end{ptree}
\noindent$\Pi_{83}$:
\begin{ptree}
\AXC{$\forall n \in \mathbb{N}\,(\underline{n} \in \mathbb{R})$\lbl{[N4]}}
\AXC{$\forall n \in \mathbb{N}\,(\underline{n} \in \mathbb{R}) \to n \in \mathbb{N} \to \underline{n} \in \mathbb{R}$\lbl{A2}}
\BIC{$n \in \mathbb{N} \to \underline{n} \in \mathbb{R}$}
\AXC{$(n \in \mathbb{N} \to \underline{n} \in \mathbb{R}) \to \varphi_{232}$\lbl{A1}}
\BIC{$\varphi_{232}$}
\end{ptree}
\noindent$\Pi_{84}$:
\begin{ptree}
\AXC{$c \in \mathbb{R} \to r \in \mathbb{R} \to 0 < r \to \varphi_{98}$\lbl{A1}}
\AXC{$\stackrel{\Pi_{83}}{\vdots}$}\noLine
\UIC{$\varphi_{232}$}
\BIC{$c \in \mathbb{R} \to r \in \mathbb{R} \to 0 < r \to \varepsilon \in \mathbb{R} \to \varphi_{88}$}
\end{ptree}
\noindent$\Pi_{85}$:
\begin{ptree}
\AXC{$\forall x \in \mathbb{R}\,(x < x + 1)$\lbl{[N13]}}
\AXC{$\forall x \in \mathbb{R}\,(x < x + 1) \to \underline{n} \in \mathbb{R} \to \underline{n} < \underline{n} + 1$\lbl{A2}}
\BIC{$\underline{n} \in \mathbb{R} \to \underline{n} < \underline{n} + 1$}
\AXC{$(\underline{n} \in \mathbb{R} \to \underline{n} < \underline{n} + 1) \to \varphi_{245}$\lbl{A1}}
\BIC{$\varphi_{245}$}
\end{ptree}
\noindent$\Pi_{86}$:
\begin{ptree}
\AXC{$\forall n \in \mathbb{N}\,(0 < \underline{n})$\lbl{[N15]}}
\AXC{$\forall n \in \mathbb{N}\,(0 < \underline{n}) \to n \in \mathbb{N} \to 0 < \underline{n}$\lbl{A2}}
\BIC{$n \in \mathbb{N} \to 0 < \underline{n}$}
\AXC{$(n \in \mathbb{N} \to 0 < \underline{n}) \to \varphi_{229}$\lbl{A1}}
\BIC{$\varphi_{229}$}
\end{ptree}
\noindent$\Pi_{87}$:
\begin{ptree}
\AXC{$c \in \mathbb{R} \to r \in \mathbb{R} \to 0 < r \to \varphi_{98}$\lbl{A1}}
\AXC{$\stackrel{\Pi_{86}}{\vdots}$}\noLine
\UIC{$\varphi_{229}$}
\BIC{$c \in \mathbb{R} \to r \in \mathbb{R} \to 0 < r \to \varepsilon \in \mathbb{R} \to \varphi_{74}$}
\end{ptree}
\noindent$\Pi_{88}$:
\begin{ptree}
\AXC{$\forall x\,\forall y\,(x \in \mathbb{R} \to y \in \mathbb{R} \to x + y \in \mathbb{R})$\lbl{(R1)}}
\AXC{$\varphi_{130}$\lbl{A2}}
\BIC{$\forall y\,(\underline{n} \in \mathbb{R} \to y \in \mathbb{R} \to \underline{n} + y \in \mathbb{R})$}
\AXC{$\varphi_{137}$\lbl{A2}}
\BIC{$\underline{n} \in \mathbb{R} \to 1 \in \mathbb{R} \to \underline{n} + 1 \in \mathbb{R}$}
\AXC{$(\underline{n} \in \mathbb{R} \to 1 \in \mathbb{R} \to \underline{n} + 1 \in \mathbb{R}) \to \varphi_{253}$\lbl{A1}}
\BIC{$\varphi_{253}$}
\end{ptree}
\noindent$\Pi_{89}$:
\begin{ptree}
\AXC{$1 \in \mathbb{R}$\lbl{(R1)}}
\AXC{$\stackrel{\Pi_{84}}{\vdots}$}\noLine
\UIC{$c \in \mathbb{R} \to r \in \mathbb{R} \to 0 < r \to \varepsilon \in \mathbb{R} \to \varphi_{88}$}
\AXC{$\stackrel{\Pi_{88}}{\vdots}$}\noLine
\UIC{$\varphi_{253}$}
\BIC{$1 \in \mathbb{R} \to c \in \mathbb{R} \to r \in \mathbb{R} \to 0 < r \to \varepsilon \in \mathbb{R} \to 0 < \varepsilon \to \varphi_{63}$}
\BIC{$c \in \mathbb{R} \to r \in \mathbb{R} \to 0 < r \to \varepsilon \in \mathbb{R} \to 0 < \varepsilon \to \varphi_{63}$}
\end{ptree}
\noindent$\Pi_{90}$:
\begin{ptree}
\AXC{$\stackrel{\Pi_{84}}{\vdots}$}\noLine
\UIC{$c \in \mathbb{R} \to r \in \mathbb{R} \to 0 < r \to \varepsilon \in \mathbb{R} \to \varphi_{88}$}
\AXC{$\varphi_{9}$\lbl{[E19]}}
\AXC{$\varphi_{9} \to \varphi_{77}$\lbl{A2}}
\BIC{$\varphi_{77}$}
\AXC{$\varphi_{77} \to \underline{n} \in \mathbb{R} \to \varphi_{112}$\lbl{A2}}
\BIC{$\underline{n} \in \mathbb{R} \to \varphi_{112}$}
\AXC{$\varphi_{315}$\lbl{A1}}
\BIC{$(c \in \mathbb{R} \to r \in \mathbb{R} \to 0 < r \to \varepsilon \in \mathbb{R} \to \varphi_{88}) \to \varphi_{308}$}
\BIC{$\varphi_{308}$}
\end{ptree}
\noindent$\Pi_{91}$:
\begin{ptree}
\AXC{$\stackrel{\Pi_{87}}{\vdots}$}\noLine
\UIC{$c \in \mathbb{R} \to r \in \mathbb{R} \to 0 < r \to \varepsilon \in \mathbb{R} \to \varphi_{74}$}
\AXC{$\stackrel{\Pi_{89}}{\vdots}$}\noLine
\UIC{$c \in \mathbb{R} \to r \in \mathbb{R} \to 0 < r \to \varepsilon \in \mathbb{R} \to 0 < \varepsilon \to \varphi_{63}$}
\AXC{$\stackrel{\Pi_{90}}{\vdots}$}\noLine
\UIC{$\varphi_{308}$}
\BIC{$(c \in \mathbb{R} \to r \in \mathbb{R} \to 0 < r \to \varepsilon \in \mathbb{R} \to \varphi_{74}) \to \varphi_{257}$}
\BIC{$\varphi_{257}$}
\end{ptree}
\noindent$\Pi_{92}$:
\begin{ptree}
\AXC{$\forall x\,\forall y\,\forall z\,(x < y \to y \leq z \to x < z)$\lbl{[O3]}}
\AXC{$\forall x\,\forall y\,\forall z\,(x < y \to y \leq z \to x < z) \to \varphi_{4}$\lbl{A2}}
\BIC{$\varphi_{4}$}
\AXC{$\varphi_{4} \to \varphi_{44}$\lbl{A2}}
\BIC{$\varphi_{44}$}
\AXC{$\varphi_{44} \to \varphi_{110}$\lbl{A2}}
\BIC{$\varphi_{110}$}
\end{ptree}
\noindent$\Pi_{93}$:
\begin{ptree}
\AXC{$\stackrel{\Pi_{92}}{\vdots}$}\noLine
\UIC{$\varphi_{110}$}
\AXC{$\varphi_{110} \to (c \in \mathbb{R} \to r \in \mathbb{R} \to 0 < r \to \varepsilon \in \mathbb{R} \to 0 < \varepsilon \to \varphi_{104}) \to \varphi_{261}$\lbl{A1}}
\BIC{$(c \in \mathbb{R} \to r \in \mathbb{R} \to 0 < r \to \varepsilon \in \mathbb{R} \to 0 < \varepsilon \to \varphi_{104}) \to \varphi_{261}$}
\end{ptree}
\noindent$\Pi_{94}$:
\begin{ptree}
\AXC{$\stackrel{\Pi_{84}}{\vdots}$}\noLine
\UIC{$c \in \mathbb{R} \to r \in \mathbb{R} \to 0 < r \to \varepsilon \in \mathbb{R} \to \varphi_{88}$}
\AXC{$\stackrel{\Pi_{85}}{\vdots}$}\noLine
\UIC{$\varphi_{245}$}
\BIC{$c \in \mathbb{R} \to r \in \mathbb{R} \to 0 < r \to \varepsilon \in \mathbb{R} \to 0 < \varepsilon \to \varphi_{73}$}
\AXC{$\stackrel{\Pi_{91}}{\vdots}$}\noLine
\UIC{$\varphi_{257}$}
\BIC{$c \in \mathbb{R} \to r \in \mathbb{R} \to 0 < r \to \varepsilon \in \mathbb{R} \to 0 < \varepsilon \to \varphi_{104}$}
\AXC{$\stackrel{\Pi_{93}}{\vdots}$}\noLine
\UIC{$(c \in \mathbb{R} \to r \in \mathbb{R} \to 0 < r \to \varepsilon \in \mathbb{R} \to 0 < \varepsilon \to \varphi_{104}) \to \varphi_{261}$}
\BIC{$\varphi_{261}$}
\end{ptree}
\noindent$\Pi_{95}$:
\begin{ptree}
\AXC{$\forall x\,\forall y\,\forall z\,(x < y \to y < z \to x < z)$\lbl{[O2]}}
\AXC{$\forall x\,\forall y\,\forall z\,(x < y \to y < z \to x < z) \to \varphi_{5}$\lbl{A2}}
\BIC{$\varphi_{5}$}
\AXC{$\varphi_{5} \to \varphi_{55}$\lbl{A2}}
\BIC{$\varphi_{55}$}
\AXC{$\varphi_{55} \to \varphi_{121}$\lbl{A2}}
\BIC{$\varphi_{121}$}
\end{ptree}
\noindent$\Pi_{96}$:
\begin{ptree}
\AXC{$\stackrel{\Pi_{95}}{\vdots}$}\noLine
\UIC{$\varphi_{121}$}
\AXC{$\varphi_{121} \to (c \in \mathbb{R} \to r \in \mathbb{R} \to 0 < r \to \varepsilon \in \mathbb{R} \to 0 < \varepsilon \to \varphi_{105}) \to \varphi_{263}$\lbl{A1}}
\BIC{$(c \in \mathbb{R} \to r \in \mathbb{R} \to 0 < r \to \varepsilon \in \mathbb{R} \to 0 < \varepsilon \to \varphi_{105}) \to \varphi_{263}$}
\end{ptree}
\noindent$\Pi_{97}$:
\begin{ptree}
\AXC{$c \in \mathbb{R} \to r \in \mathbb{R} \to 0 < r \to \varepsilon \in \mathbb{R} \to \varphi_{37}$\lbl{A1}}
\AXC{$\stackrel{\Pi_{82}}{\vdots}$}\noLine
\UIC{$\varphi_{236}$}
\BIC{$c \in \mathbb{R} \to r \in \mathbb{R} \to 0 < r \to \varepsilon \in \mathbb{R} \to 0 < \varepsilon \to \varphi_{79}$}
\AXC{$\stackrel{\Pi_{94}}{\vdots}$}\noLine
\UIC{$\varphi_{261}$}
\BIC{$c \in \mathbb{R} \to r \in \mathbb{R} \to 0 < r \to \varepsilon \in \mathbb{R} \to 0 < \varepsilon \to \varphi_{105}$}
\AXC{$\stackrel{\Pi_{96}}{\vdots}$}\noLine
\UIC{$(c \in \mathbb{R} \to r \in \mathbb{R} \to 0 < r \to \varepsilon \in \mathbb{R} \to 0 < \varepsilon \to \varphi_{105}) \to \varphi_{263}$}
\BIC{$\varphi_{263}$}
\end{ptree}
\noindent$\Pi_{98}$:
\begin{ptree}
\AXC{$\stackrel{\Pi_{20}}{\vdots}$}\noLine
\UIC{$\forall x\,\forall y\,(x = y \to y = x)$}
\AXC{$\forall x\,\forall y\,(x = y \to y = x) \to \forall y\,(r = y \to y = r)$\lbl{A2}}
\BIC{$\forall y\,(r = y \to y = r)$}
\AXC{$\forall y\,(r = y \to y = r) \to r = 0 \to 0 = r$\lbl{A2}}
\BIC{$r = 0 \to 0 = r$}
\end{ptree}
\noindent$\Pi_{99}$:
\begin{ptree}
\AXC{$\forall x\,\forall y\,(x < y \to x \neq y)$\lbl{[F26]}}
\AXC{$\forall x\,\forall y\,(x < y \to x \neq y) \to \forall y\,(0 < y \to 0 \neq y)$\lbl{A2}}
\BIC{$\forall y\,(0 < y \to 0 \neq y)$}
\AXC{$\forall y\,(0 < y \to 0 \neq y) \to 0 < r \to 0 \neq r$\lbl{A2}}
\BIC{$0 < r \to 0 \neq r$}
\end{ptree}
\noindent$\Pi_{100}$:
\begin{ptree}
\AXC{$\stackrel{\Pi_{98}}{\vdots}$}\noLine
\UIC{$r = 0 \to 0 = r$}
\AXC{$c \in \mathbb{R} \to r \in \mathbb{R} \to 0 < r \to 0 < r$\lbl{A1}}
\AXC{$\stackrel{\Pi_{99}}{\vdots}$}\noLine
\UIC{$0 < r \to 0 \neq r$}
\AXC{$(0 < r \to 0 \neq r) \to \varphi_{61}$\lbl{A1}}
\BIC{$\varphi_{61}$}
\BIC{$c \in \mathbb{R} \to r \in \mathbb{R} \to 0 < r \to 0 \neq r$}
\AXC{$\varphi_{101}$\lbl{A1}}
\BIC{$(r = 0 \to 0 = r) \to c \in \mathbb{R} \to r \in \mathbb{R} \to 0 < r \to r \neq 0$}
\BIC{$c \in \mathbb{R} \to r \in \mathbb{R} \to 0 < r \to r \neq 0$}
\end{ptree}
\noindent$\Pi_{101}$:
\begin{ptree}
\AXC{$c \in \mathbb{R} \to r \in \mathbb{R} \to 0 < r \to \varepsilon \in \mathbb{R} \to \varphi_{64}$\lbl{A1}}
\AXC{$\forall x \in \mathbb{R}\,(x \neq 0 \to x^{-1} \in \mathbb{R})$\lbl{(R1)}}
\AXC{$\varphi_{33}$\lbl{A2}}
\BIC{$r \in \mathbb{R} \to r \neq 0 \to r^{-1} \in \mathbb{R}$}
\AXC{$\varphi_{279}$\lbl{A1}}
\BIC{$(c \in \mathbb{R} \to r \in \mathbb{R} \to 0 < r \to \varepsilon \in \mathbb{R} \to \varphi_{64}) \to \varphi_{209}$}
\BIC{$\varphi_{209}$}
\end{ptree}
\noindent$\Pi_{102}$:
\begin{ptree}
\AXC{$\forall x\,\forall y\,(x \in \mathbb{R} \to y \in \mathbb{R} \to x \cdot y \in \mathbb{R})$\lbl{(R1)}}
\AXC{$\varphi_{111}$\lbl{A2}}
\BIC{$\forall y\,(\varepsilon \in \mathbb{R} \to y \in \mathbb{R} \to \varepsilon \cdot y \in \mathbb{R})$}
\AXC{$\varphi_{131}$\lbl{A2}}
\BIC{$\varepsilon \in \mathbb{R} \to r^{-1} \in \mathbb{R} \to \varepsilon \cdot r^{-1} \in \mathbb{R}$}
\AXC{$\varphi_{288}$\lbl{A1}}
\BIC{$(c \in \mathbb{R} \to r \in \mathbb{R} \to 0 < r \to \varepsilon \in \mathbb{R} \to \varepsilon \in \mathbb{R}) \to \varphi_{244}$}
\end{ptree}
\noindent$\Pi_{103}$:
\begin{ptree}
\AXC{$c \in \mathbb{R} \to r \in \mathbb{R} \to 0 < r \to \varepsilon \in \mathbb{R} \to \varepsilon \in \mathbb{R}$\lbl{A1}}
\AXC{$\stackrel{\Pi_{102}}{\vdots}$}\noLine
\UIC{$(c \in \mathbb{R} \to r \in \mathbb{R} \to 0 < r \to \varepsilon \in \mathbb{R} \to \varepsilon \in \mathbb{R}) \to \varphi_{244}$}
\BIC{$\varphi_{244}$}
\end{ptree}
\noindent$\Pi_{104}$:
\begin{ptree}
\AXC{$\stackrel{\Pi_{100}}{\vdots}$}\noLine
\UIC{$c \in \mathbb{R} \to r \in \mathbb{R} \to 0 < r \to r \neq 0$}
\AXC{$\stackrel{\Pi_{101}}{\vdots}$}\noLine
\UIC{$\varphi_{209}$}
\BIC{$c \in \mathbb{R} \to r \in \mathbb{R} \to 0 < r \to \varepsilon \in \mathbb{R} \to 0 < \varepsilon \to \varphi_{40}$}
\AXC{$\stackrel{\Pi_{103}}{\vdots}$}\noLine
\UIC{$\varphi_{244}$}
\BIC{$c \in \mathbb{R} \to r \in \mathbb{R} \to 0 < r \to \varepsilon \in \mathbb{R} \to 0 < \varepsilon \to \varphi_{68}$}
\end{ptree}
\noindent$\Pi_{105}$:
\begin{ptree}
\AXC{$\forall x\,\forall y\,(x \in \mathbb{R} \to y \in \mathbb{R} \to 0 < x \to 0 < y \to 0 < x + y)$\lbl{[E5]}}
\AXC{$\varphi_{186}$\lbl{A2}}
\BIC{$\forall y\,(\underline{n} \in \mathbb{R} \to y \in \mathbb{R} \to 0 < \underline{n} \to 0 < y \to 0 < \underline{n} + y)$}
\AXC{$\varphi_{192}$\lbl{A2}}
\BIC{$\underline{n} \in \mathbb{R} \to 1 \in \mathbb{R} \to 0 < \underline{n} \to 0 < 1 \to 0 < \underline{n} + 1$}
\end{ptree}
\noindent$\Pi_{106}$:
\begin{ptree}
\AXC{$\stackrel{\Pi_{105}}{\vdots}$}\noLine
\UIC{$\underline{n} \in \mathbb{R} \to 1 \in \mathbb{R} \to 0 < \underline{n} \to 0 < 1 \to 0 < \underline{n} + 1$}
\AXC{$(\underline{n} \in \mathbb{R} \to 1 \in \mathbb{R} \to 0 < \underline{n} \to 0 < 1 \to 0 < \underline{n} + 1) \to \varphi_{304}$\lbl{A1}}
\BIC{$\varphi_{304}$}
\end{ptree}
\noindent$\Pi_{107}$:
\begin{ptree}
\AXC{$\stackrel{\Pi_{87}}{\vdots}$}\noLine
\UIC{$c \in \mathbb{R} \to r \in \mathbb{R} \to 0 < r \to \varepsilon \in \mathbb{R} \to \varphi_{74}$}
\AXC{$1 \in \mathbb{R}$\lbl{(R1)}}
\AXC{$\stackrel{\Pi_{84}}{\vdots}$}\noLine
\UIC{$c \in \mathbb{R} \to r \in \mathbb{R} \to 0 < r \to \varepsilon \in \mathbb{R} \to \varphi_{88}$}
\AXC{$\stackrel{\Pi_{106}}{\vdots}$}\noLine
\UIC{$\varphi_{304}$}
\BIC{$1 \in \mathbb{R} \to (c \in \mathbb{R} \to r \in \mathbb{R} \to 0 < r \to \varepsilon \in \mathbb{R} \to \varphi_{74}) \to \varphi_{233}$}
\BIC{$(c \in \mathbb{R} \to r \in \mathbb{R} \to 0 < r \to \varepsilon \in \mathbb{R} \to \varphi_{74}) \to \varphi_{233}$}
\BIC{$\varphi_{233}$}
\end{ptree}
\noindent$\Pi_{108}$:
\begin{ptree}
\AXC{$0 < 1$\lbl{[F28]}}
\AXC{$0 < 1 \to c \in \mathbb{R} \to r \in \mathbb{R} \to 0 < r \to \varepsilon \in \mathbb{R} \to \varphi_{50}$\lbl{A1}}
\BIC{$c \in \mathbb{R} \to r \in \mathbb{R} \to 0 < r \to \varepsilon \in \mathbb{R} \to \varphi_{50}$}
\AXC{$\stackrel{\Pi_{107}}{\vdots}$}\noLine
\UIC{$\varphi_{233}$}
\BIC{$c \in \mathbb{R} \to r \in \mathbb{R} \to 0 < r \to \varepsilon \in \mathbb{R} \to 0 < \varepsilon \to \varphi_{51}$}
\end{ptree}
\noindent$\Pi_{109}$:
\begin{ptree}
\AXC{$\forall x\,\forall y\,(x < y \to x \neq y)$\lbl{[F26]}}
\AXC{$\forall x\,\forall y\,(x < y \to x \neq y) \to \forall y\,(0 < y \to 0 \neq y)$\lbl{A2}}
\BIC{$\forall y\,(0 < y \to 0 \neq y)$}
\end{ptree}
\noindent$\Pi_{110}$:
\begin{ptree}
\AXC{$\stackrel{\Pi_{109}}{\vdots}$}\noLine
\UIC{$\forall y\,(0 < y \to 0 \neq y)$}
\AXC{$\forall y\,(0 < y \to 0 \neq y) \to 0 < \underline{n} + 1 \to 0 \neq \underline{n} + 1$\lbl{A2}}
\BIC{$0 < \underline{n} + 1 \to 0 \neq \underline{n} + 1$}
\AXC{$(0 < \underline{n} + 1 \to 0 \neq \underline{n} + 1) \to \varphi_{241}$\lbl{A1}}
\BIC{$\varphi_{241}$}
\end{ptree}
\noindent$\Pi_{111}$:
\begin{ptree}
\AXC{$\stackrel{\Pi_{108}}{\vdots}$}\noLine
\UIC{$c \in \mathbb{R} \to r \in \mathbb{R} \to 0 < r \to \varepsilon \in \mathbb{R} \to 0 < \varepsilon \to \varphi_{51}$}
\AXC{$\stackrel{\Pi_{110}}{\vdots}$}\noLine
\UIC{$\varphi_{241}$}
\BIC{$c \in \mathbb{R} \to r \in \mathbb{R} \to 0 < r \to \varepsilon \in \mathbb{R} \to 0 < \varepsilon \to \varphi_{41}$}
\AXC{$(c \in \mathbb{R} \to r \in \mathbb{R} \to 0 < r \to \varepsilon \in \mathbb{R} \to 0 < \varepsilon \to \varphi_{41}) \to \varphi_{208}$\lbl{A1}}
\BIC{$\varphi_{208}$}
\end{ptree}
\noindent$\Pi_{112}$:
\begin{ptree}
\AXC{$\stackrel{\Pi_{89}}{\vdots}$}\noLine
\UIC{$c \in \mathbb{R} \to r \in \mathbb{R} \to 0 < r \to \varepsilon \in \mathbb{R} \to 0 < \varepsilon \to \varphi_{63}$}
\AXC{$\forall x \in \mathbb{R}\,(x \neq 0 \to x^{-1} \in \mathbb{R})$\lbl{(R1)}}
\AXC{$\varphi_{135}$\lbl{A2}}
\BIC{$\underline{n} + 1 \in \mathbb{R} \to \underline{n} + 1 \neq 0 \to (\underline{n} + 1)^{-1} \in \mathbb{R}$}
\AXC{$\varphi_{300}$\lbl{A1}}
\BIC{$(c \in \mathbb{R} \to r \in \mathbb{R} \to 0 < r \to \varepsilon \in \mathbb{R} \to 0 < \varepsilon \to \varphi_{63}) \to \varphi_{249}$}
\BIC{$\varphi_{249}$}
\end{ptree}
\noindent$\Pi_{113}$:
\begin{ptree}
\AXC{$\stackrel{\Pi_{20}}{\vdots}$}\noLine
\UIC{$\forall x\,\forall y\,(x = y \to y = x)$}
\AXC{$\varphi_{1}$\lbl{A2}}
\BIC{$\forall y\,(\underline{n} + 1 = y \to y = \underline{n} + 1)$}
\AXC{$\varphi_{11}$\lbl{A2}}
\BIC{$\underline{n} + 1 = 0 \to 0 = \underline{n} + 1$}
\AXC{$\stackrel{\Pi_{111}}{\vdots}$}\noLine
\UIC{$\varphi_{208}$}
\BIC{$c \in \mathbb{R} \to r \in \mathbb{R} \to 0 < r \to \varepsilon \in \mathbb{R} \to 0 < \varepsilon \to \varphi_{42}$}
\AXC{$\stackrel{\Pi_{112}}{\vdots}$}\noLine
\UIC{$\varphi_{249}$}
\BIC{$c \in \mathbb{R} \to r \in \mathbb{R} \to 0 < r \to \varepsilon \in \mathbb{R} \to 0 < \varepsilon \to \varphi_{81}$}
\end{ptree}
\noindent$\Pi_{114}$:
\begin{ptree}
\AXC{$\stackrel{\Pi_{113}}{\vdots}$}\noLine
\UIC{$c \in \mathbb{R} \to r \in \mathbb{R} \to 0 < r \to \varepsilon \in \mathbb{R} \to 0 < \varepsilon \to \varphi_{81}$}
\AXC{$\varphi_{58}$\lbl{[E8]}}
\AXC{$\varphi_{58} \to \varphi_{173}$\lbl{A2}}
\BIC{$\varphi_{173}$}
\AXC{$\varphi_{173} \to \varphi_{198}$\lbl{A2}}
\BIC{$\varphi_{198}$}
\AXC{$\varphi_{198} \to \varphi_{196}$\lbl{A2}}
\BIC{$\varphi_{196}$}
\AXC{$\varphi_{196} \to \varphi_{314}$\lbl{A1}}
\BIC{$\varphi_{314}$}
\BIC{$(c \in \mathbb{R} \to r \in \mathbb{R} \to 0 < r \to \varepsilon \in \mathbb{R} \to 0 < \varepsilon \to \varphi_{68}) \to \varphi_{299}$}
\end{ptree}
\noindent$\Pi_{115}$:
\begin{ptree}
\AXC{$\stackrel{\Pi_{104}}{\vdots}$}\noLine
\UIC{$c \in \mathbb{R} \to r \in \mathbb{R} \to 0 < r \to \varepsilon \in \mathbb{R} \to 0 < \varepsilon \to \varphi_{68}$}
\AXC{$\stackrel{\Pi_{114}}{\vdots}$}\noLine
\UIC{$(c \in \mathbb{R} \to r \in \mathbb{R} \to 0 < r \to \varepsilon \in \mathbb{R} \to 0 < \varepsilon \to \varphi_{68}) \to \varphi_{299}$}
\BIC{$\varphi_{299}$}
\end{ptree}
\noindent$\Pi_{116}$:
\begin{ptree}
\AXC{$c \in \mathbb{R} \to r \in \mathbb{R} \to 0 < r \to 0 < r$\lbl{A1}}
\AXC{$c \in \mathbb{R} \to r \in \mathbb{R} \to 0 < r \to \varepsilon \in \mathbb{R} \to \varphi_{64}$\lbl{A1}}
\AXC{$\stackrel{\Pi_{115}}{\vdots}$}\noLine
\UIC{$\varphi_{299}$}
\BIC{$(c \in \mathbb{R} \to r \in \mathbb{R} \to 0 < r \to 0 < r) \to \varphi_{269}$}
\BIC{$\varphi_{269}$}
\end{ptree}
\noindent$\Pi_{117}$:
\begin{ptree}
\AXC{$c \in \mathbb{R} \to r \in \mathbb{R} \to 0 < r \to \varepsilon \in \mathbb{R} \to 0 < \varepsilon \to \varphi_{72}$\lbl{A1}}
\AXC{$\varphi_{255}$\lbl{A1}}
\BIC{$c \in \mathbb{R} \to r \in \mathbb{R} \to 0 < r \to \varepsilon \in \mathbb{R} \to 0 < \varepsilon \to \varphi_{97}$}
\AXC{$\stackrel{\Pi_{97}}{\vdots}$}\noLine
\UIC{$\varphi_{263}$}
\BIC{$c \in \mathbb{R} \to r \in \mathbb{R} \to 0 < r \to \varepsilon \in \mathbb{R} \to 0 < \varepsilon \to \varphi_{107}$}
\AXC{$\stackrel{\Pi_{116}}{\vdots}$}\noLine
\UIC{$\varphi_{269}$}
\BIC{$c \in \mathbb{R} \to r \in \mathbb{R} \to 0 < r \to \varepsilon \in \mathbb{R} \to 0 < \varepsilon \to \varphi_{132}$}
\end{ptree}
\noindent$\Pi_{118}$:
\begin{ptree}
\AXC{$y = z \to x = y \to x = z$\lbl{A6}}
\AXC{$(y = z \to x = y \to x = z) \to x = y \to y = z \to x = z$\lbl{A1}}
\BIC{$x = y \to y = z \to x = z$}
\UIC{$\forall z\,(x = y \to y = z \to x = z)$}
\UIC{$\forall y\,\forall z\,(x = y \to y = z \to x = z)$}
\UIC{$\forall x\,\forall y\,\forall z\,(x = y \to y = z \to x = z)$}
\AXC{$\forall x\,\forall y\,\forall z\,(x = y \to y = z \to x = z) \to \varphi_{2}$\lbl{A2}}
\BIC{$\varphi_{2}$}
\end{ptree}
\noindent$\Pi_{119}$:
\begin{ptree}
\AXC{$\forall x\,\forall y\,(x \in \mathbb{R} \to y \in \mathbb{R} \to y \neq 0 \to x \cdot y^{-1} \cdot y = x)$\lbl{[E23]}}
\AXC{$\varphi_{159}$\lbl{A2}}
\BIC{$\forall y\,(\varepsilon \in \mathbb{R} \to y \in \mathbb{R} \to y \neq 0 \to \varepsilon \cdot y^{-1} \cdot y = \varepsilon)$}
\AXC{$\varphi_{160}$\lbl{A2}}
\BIC{$\varepsilon \in \mathbb{R} \to r \in \mathbb{R} \to r \neq 0 \to \varepsilon \cdot r^{-1} \cdot r = \varepsilon$}
\end{ptree}
\noindent$\Pi_{120}$:
\begin{ptree}
\AXC{$\stackrel{\Pi_{119}}{\vdots}$}\noLine
\UIC{$\varepsilon \in \mathbb{R} \to r \in \mathbb{R} \to r \neq 0 \to \varepsilon \cdot r^{-1} \cdot r = \varepsilon$}
\AXC{$(\varepsilon \in \mathbb{R} \to r \in \mathbb{R} \to r \neq 0 \to \varepsilon \cdot r^{-1} \cdot r = \varepsilon) \to \varphi_{287}$\lbl{A1}}
\BIC{$\varphi_{287}$}
\end{ptree}
\noindent$\Pi_{121}$:
\begin{ptree}
\AXC{$c \in \mathbb{R} \to r \in \mathbb{R} \to 0 < r \to \varepsilon \in \mathbb{R} \to \varphi_{64}$\lbl{A1}}
\AXC{$c \in \mathbb{R} \to r \in \mathbb{R} \to 0 < r \to \varepsilon \in \mathbb{R} \to \varepsilon \in \mathbb{R}$\lbl{A1}}
\AXC{$\stackrel{\Pi_{120}}{\vdots}$}\noLine
\UIC{$\varphi_{287}$}
\BIC{$(c \in \mathbb{R} \to r \in \mathbb{R} \to 0 < r \to \varepsilon \in \mathbb{R} \to \varphi_{64}) \to \varphi_{210}$}
\BIC{$\varphi_{210}$}
\end{ptree}
\noindent$\Pi_{122}$:
\begin{ptree}
\AXC{$\forall x\,\forall y\,(x \in \mathbb{R} \to y \in \mathbb{R} \to x \cdot y = y \cdot x)$\lbl{(R7)}}
\AXC{$\varphi_{93}$\lbl{A2}}
\BIC{$\forall y\,(r \in \mathbb{R} \to y \in \mathbb{R} \to r \cdot y = y \cdot r)$}
\AXC{$\forall y\,(r \in \mathbb{R} \to y \in \mathbb{R} \to r \cdot y = y \cdot r) \to \varphi_{30}$\lbl{A2}}
\BIC{$\varphi_{30}$}
\end{ptree}
\noindent$\Pi_{123}$:
\begin{ptree}
\AXC{$c \in \mathbb{R} \to r \in \mathbb{R} \to 0 < r \to \varepsilon \in \mathbb{R} \to \varphi_{64}$\lbl{A1}}
\AXC{$\stackrel{\Pi_{122}}{\vdots}$}\noLine
\UIC{$\varphi_{30}$}
\AXC{$\varphi_{30} \to (c \in \mathbb{R} \to r \in \mathbb{R} \to 0 < r \to \varepsilon \in \mathbb{R} \to \varphi_{64}) \to \varphi_{262}$\lbl{A1}}
\BIC{$(c \in \mathbb{R} \to r \in \mathbb{R} \to 0 < r \to \varepsilon \in \mathbb{R} \to \varphi_{64}) \to \varphi_{262}$}
\BIC{$\varphi_{262}$}
\end{ptree}
\noindent$\Pi_{124}$:
\begin{ptree}
\AXC{$\stackrel{\Pi_{104}}{\vdots}$}\noLine
\UIC{$c \in \mathbb{R} \to r \in \mathbb{R} \to 0 < r \to \varepsilon \in \mathbb{R} \to 0 < \varepsilon \to \varphi_{68}$}
\AXC{$\stackrel{\Pi_{123}}{\vdots}$}\noLine
\UIC{$\varphi_{262}$}
\BIC{$c \in \mathbb{R} \to r \in \mathbb{R} \to 0 < r \to \varepsilon \in \mathbb{R} \to 0 < \varepsilon \to \varphi_{113}$}
\AXC{$(c \in \mathbb{R} \to r \in \mathbb{R} \to 0 < r \to \varepsilon \in \mathbb{R} \to 0 < \varepsilon \to \varphi_{113}) \to \varphi_{286}$\lbl{A1}}
\BIC{$\varphi_{286}$}
\end{ptree}
\noindent$\Pi_{125}$:
\begin{ptree}
\AXC{$\stackrel{\Pi_{118}}{\vdots}$}\noLine
\UIC{$\varphi_{2}$}
\AXC{$\varphi_{2} \to \varphi_{75}$\lbl{A2}}
\BIC{$\varphi_{75}$}
\AXC{$\varphi_{75} \to \varphi_{78}$\lbl{A2}}
\BIC{$\varphi_{78}$}
\AXC{$\stackrel{\Pi_{100}}{\vdots}$}\noLine
\UIC{$c \in \mathbb{R} \to r \in \mathbb{R} \to 0 < r \to r \neq 0$}
\AXC{$\stackrel{\Pi_{121}}{\vdots}$}\noLine
\UIC{$\varphi_{210}$}
\BIC{$c \in \mathbb{R} \to r \in \mathbb{R} \to 0 < r \to \varepsilon \in \mathbb{R} \to 0 < \varepsilon \to \varphi_{69}$}
\AXC{$\stackrel{\Pi_{124}}{\vdots}$}\noLine
\UIC{$\varphi_{286}$}
\BIC{$\varphi_{78} \to c \in \mathbb{R} \to r \in \mathbb{R} \to 0 < r \to \varepsilon \in \mathbb{R} \to 0 < \varepsilon \to \varphi_{70}$}
\BIC{$c \in \mathbb{R} \to r \in \mathbb{R} \to 0 < r \to \varepsilon \in \mathbb{R} \to 0 < \varepsilon \to \varphi_{70}$}
\AXC{$\varphi_{306}$\lbl{A1}}
\BIC{$(c \in \mathbb{R} \to r \in \mathbb{R} \to 0 < r \to \varepsilon \in \mathbb{R} \to 0 < \varepsilon \to \varphi_{132}) \to \varphi_{226}$}
\end{ptree}
\noindent$\Pi_{126}$:
\begin{ptree}
\AXC{$\stackrel{\Pi_{117}}{\vdots}$}\noLine
\UIC{$c \in \mathbb{R} \to r \in \mathbb{R} \to 0 < r \to \varepsilon \in \mathbb{R} \to 0 < \varepsilon \to \varphi_{132}$}
\AXC{$\stackrel{\Pi_{125}}{\vdots}$}\noLine
\UIC{$(c \in \mathbb{R} \to r \in \mathbb{R} \to 0 < r \to \varepsilon \in \mathbb{R} \to 0 < \varepsilon \to \varphi_{132}) \to \varphi_{226}$}
\BIC{$\varphi_{226}$}
\end{ptree}
\noindent$\Pi_{127}$:
\begin{ptree}
\AXC{$\stackrel{\Pi_{20}}{\vdots}$}\noLine
\UIC{$\forall x\,\forall y\,(x = y \to y = x)$}
\AXC{$\varphi_{59}$\lbl{A2}}
\BIC{$\forall y\,(\lvert \pi_{c,r}(n) - c \rvert = y \to y = \lvert \pi_{c,r}(n) - c \rvert)$}
\AXC{$\forall y\,(\lvert \pi_{c,r}(n) - c \rvert = y \to y = \lvert \pi_{c,r}(n) - c \rvert) \to \varphi_{123}$\lbl{A2}}
\BIC{$\varphi_{123}$}
\end{ptree}
\noindent$\Pi_{128}$:
\begin{ptree}
\AXC{$y = z \to x = y \to x = z$\lbl{A6}}
\AXC{$(y = z \to x = y \to x = z) \to x = y \to y = z \to x = z$\lbl{A1}}
\BIC{$x = y \to y = z \to x = z$}
\UIC{$\forall z\,(x = y \to y = z \to x = z)$}
\UIC{$\forall y\,\forall z\,(x = y \to y = z \to x = z)$}
\UIC{$\forall x\,\forall y\,\forall z\,(x = y \to y = z \to x = z)$}
\AXC{$\forall x\,\forall y\,\forall z\,(x = y \to y = z \to x = z) \to \varphi_{22}$\lbl{A2}}
\BIC{$\varphi_{22}$}
\end{ptree}
\noindent$\Pi_{129}$:
\begin{ptree}
\AXC{$\stackrel{\Pi_{89}}{\vdots}$}\noLine
\UIC{$c \in \mathbb{R} \to r \in \mathbb{R} \to 0 < r \to \varepsilon \in \mathbb{R} \to 0 < \varepsilon \to \varphi_{63}$}
\AXC{$\forall x \in \mathbb{R}\,(0 < x \to 0 < x^{-1})$\lbl{[F34]}}
\AXC{$\varphi_{128}$\lbl{A2}}
\BIC{$\underline{n} + 1 \in \mathbb{R} \to 0 < \underline{n} + 1 \to 0 < (\underline{n} + 1)^{-1}$}
\AXC{$\varphi_{298}$\lbl{A1}}
\BIC{$(c \in \mathbb{R} \to r \in \mathbb{R} \to 0 < r \to \varepsilon \in \mathbb{R} \to 0 < \varepsilon \to \varphi_{63}) \to \varphi_{246}$}
\BIC{$\varphi_{246}$}
\end{ptree}
\noindent$\Pi_{130}$:
\begin{ptree}
\AXC{$\forall x\,\forall y\,(x \in \mathbb{R} \to y \in \mathbb{R} \to 0 < x \to 0 < y \to 0 < x \cdot y)$\lbl{(R15)}}
\AXC{$\varphi_{165}$\lbl{A2}}
\BIC{$\forall y\,(r \in \mathbb{R} \to y \in \mathbb{R} \to 0 < r \to 0 < y \to 0 < r \cdot y)$}
\AXC{$\varphi_{204}$\lbl{A2}}
\BIC{$r \in \mathbb{R} \to \varphi_{95}$}
\AXC{$\varphi_{313}$\lbl{A1}}
\BIC{$(c \in \mathbb{R} \to r \in \mathbb{R} \to 0 < r \to \varepsilon \in \mathbb{R} \to \varphi_{64}) \to \varphi_{295}$}
\end{ptree}
\noindent$\Pi_{131}$:
\begin{ptree}
\AXC{$c \in \mathbb{R} \to r \in \mathbb{R} \to 0 < r \to \varepsilon \in \mathbb{R} \to \varphi_{64}$\lbl{A1}}
\AXC{$\stackrel{\Pi_{130}}{\vdots}$}\noLine
\UIC{$(c \in \mathbb{R} \to r \in \mathbb{R} \to 0 < r \to \varepsilon \in \mathbb{R} \to \varphi_{64}) \to \varphi_{295}$}
\BIC{$\varphi_{295}$}
\end{ptree}
\noindent$\Pi_{132}$:
\begin{ptree}
\AXC{$c \in \mathbb{R} \to r \in \mathbb{R} \to 0 < r \to 0 < r$\lbl{A1}}
\AXC{$\stackrel{\Pi_{113}}{\vdots}$}\noLine
\UIC{$c \in \mathbb{R} \to r \in \mathbb{R} \to 0 < r \to \varepsilon \in \mathbb{R} \to 0 < \varepsilon \to \varphi_{81}$}
\AXC{$\stackrel{\Pi_{131}}{\vdots}$}\noLine
\UIC{$\varphi_{295}$}
\BIC{$(c \in \mathbb{R} \to r \in \mathbb{R} \to 0 < r \to 0 < r) \to \varphi_{252}$}
\BIC{$\varphi_{252}$}
\end{ptree}
\noindent$\Pi_{133}$:
\begin{ptree}
\AXC{$\stackrel{\Pi_{108}}{\vdots}$}\noLine
\UIC{$c \in \mathbb{R} \to r \in \mathbb{R} \to 0 < r \to \varepsilon \in \mathbb{R} \to 0 < \varepsilon \to \varphi_{51}$}
\AXC{$\stackrel{\Pi_{129}}{\vdots}$}\noLine
\UIC{$\varphi_{246}$}
\BIC{$c \in \mathbb{R} \to r \in \mathbb{R} \to 0 < r \to \varepsilon \in \mathbb{R} \to 0 < \varepsilon \to \varphi_{71}$}
\AXC{$\stackrel{\Pi_{132}}{\vdots}$}\noLine
\UIC{$\varphi_{252}$}
\BIC{$c \in \mathbb{R} \to r \in \mathbb{R} \to 0 < r \to \varepsilon \in \mathbb{R} \to 0 < \varepsilon \to \varphi_{89}$}
\AXC{$\varphi_{254}$\lbl{A1}}
\BIC{$c \in \mathbb{R} \to r \in \mathbb{R} \to 0 < r \to \varepsilon \in \mathbb{R} \to 0 < \varepsilon \to \varphi_{76}$}
\end{ptree}
\noindent$\Pi_{134}$:
\begin{ptree}
\AXC{$\forall x\,(0 \leq x \to \lvert x \rvert = x)$\lbl{[A1]}}
\AXC{$\forall x\,(0 \leq x \to \lvert x \rvert = x) \to \varphi_{38}$\lbl{A2}}
\BIC{$\varphi_{38}$}
\AXC{$\varphi_{38} \to \varphi_{266}$\lbl{A1}}
\BIC{$\varphi_{266}$}
\end{ptree}
\noindent$\Pi_{135}$:
\begin{ptree}
\AXC{$z = w \to \lvert \pi_{c,r}(n) - c \rvert = \lvert z \rvert \to \lvert \pi_{c,r}(n) - c \rvert = \lvert w \rvert$\lbl{A6}}
\UIC{$\varphi_{52}$}
\UIC{$\forall z\,(\varphi_{52})$}
\AXC{$\forall z\,(\varphi_{52}) \to \varphi_{166}$\lbl{A2}}
\BIC{$\varphi_{166}$}
\AXC{$\varphi_{166} \to \pi_{c,r}(n) - c = r \cdot (\underline{n} + 1)^{-1} \to \varphi_{140}$\lbl{A2}}
\BIC{$\pi_{c,r}(n) - c = r \cdot (\underline{n} + 1)^{-1} \to \varphi_{140}$}
\end{ptree}
\noindent$\Pi_{136}$:
\begin{ptree}
\AXC{$y = z \to x = y \to x = z$\lbl{A6}}
\AXC{$(y = z \to x = y \to x = z) \to x = y \to y = z \to x = z$\lbl{A1}}
\BIC{$x = y \to y = z \to x = z$}
\UIC{$\forall z\,(x = y \to y = z \to x = z)$}
\UIC{$\forall y\,\forall z\,(x = y \to y = z \to x = z)$}
\UIC{$\forall x\,\forall y\,\forall z\,(x = y \to y = z \to x = z)$}
\AXC{$\varphi_{91}$\lbl{A2}}
\BIC{$\forall y\,\forall z\,(\pi_{c,r}(n) - c = y \to y = z \to \pi_{c,r}(n) - c = z)$}
\end{ptree}
\noindent$\Pi_{137}$:
\begin{ptree}
\AXC{$\stackrel{\Pi_{136}}{\vdots}$}\noLine
\UIC{$\forall y\,\forall z\,(\pi_{c,r}(n) - c = y \to y = z \to \pi_{c,r}(n) - c = z)$}
\AXC{$\forall y\,\forall z\,(\pi_{c,r}(n) - c = y \to y = z \to \pi_{c,r}(n) - c = z) \to \varphi_{167}$\lbl{A2}}
\BIC{$\varphi_{167}$}
\end{ptree}
\noindent$\Pi_{138}$:
\begin{ptree}
\AXC{$\forall x\,\forall y\,(x \in \mathbb{R} \to y \in \mathbb{R} \to x \cdot y \in \mathbb{R})$\lbl{(R1)}}
\AXC{$\varphi_{103}$\lbl{A2}}
\BIC{$\forall y\,(r \in \mathbb{R} \to y \in \mathbb{R} \to r \cdot y \in \mathbb{R})$}
\AXC{$\varphi_{157}$\lbl{A2}}
\BIC{$r \in \mathbb{R} \to (\underline{n} + 1)^{-1} \in \mathbb{R} \to r \cdot (\underline{n} + 1)^{-1} \in \mathbb{R}$}
\AXC{$\varphi_{301}$\lbl{A1}}
\BIC{$(c \in \mathbb{R} \to r \in \mathbb{R} \to 0 < r \to \varepsilon \in \mathbb{R} \to \varphi_{64}) \to \varphi_{260}$}
\end{ptree}
\noindent$\Pi_{139}$:
\begin{ptree}
\AXC{$c \in \mathbb{R} \to r \in \mathbb{R} \to 0 < r \to \varepsilon \in \mathbb{R} \to \varphi_{64}$\lbl{A1}}
\AXC{$\stackrel{\Pi_{138}}{\vdots}$}\noLine
\UIC{$(c \in \mathbb{R} \to r \in \mathbb{R} \to 0 < r \to \varepsilon \in \mathbb{R} \to \varphi_{64}) \to \varphi_{260}$}
\BIC{$\varphi_{260}$}
\end{ptree}
\noindent$\Pi_{140}$:
\begin{ptree}
\AXC{$\stackrel{\Pi_{113}}{\vdots}$}\noLine
\UIC{$c \in \mathbb{R} \to r \in \mathbb{R} \to 0 < r \to \varepsilon \in \mathbb{R} \to 0 < \varepsilon \to \varphi_{81}$}
\AXC{$\stackrel{\Pi_{139}}{\vdots}$}\noLine
\UIC{$\varphi_{260}$}
\BIC{$c \in \mathbb{R} \to r \in \mathbb{R} \to 0 < r \to \varepsilon \in \mathbb{R} \to 0 < \varepsilon \to \varphi_{102}$}
\end{ptree}
\noindent$\Pi_{141}$:
\begin{ptree}
\AXC{$\forall x\,\forall y\,(x \in \mathbb{R} \to y \in \mathbb{R} \to x + y - y = x)$\lbl{[E38]}}
\AXC{$\forall x\,\forall y\,(x \in \mathbb{R} \to y \in \mathbb{R} \to x + y - y = x) \to \varphi_{129}$\lbl{A2}}
\BIC{$\varphi_{129}$}
\AXC{$\varphi_{129} \to \varphi_{119}$\lbl{A2}}
\BIC{$\varphi_{119}$}
\end{ptree}
\noindent$\Pi_{142}$:
\begin{ptree}
\AXC{$c \in \mathbb{R} \to r \in \mathbb{R} \to 0 < r \to \varepsilon \in \mathbb{R} \to \varphi_{65}$\lbl{A1}}
\AXC{$\stackrel{\Pi_{140}}{\vdots}$}\noLine
\UIC{$c \in \mathbb{R} \to r \in \mathbb{R} \to 0 < r \to \varepsilon \in \mathbb{R} \to 0 < \varepsilon \to \varphi_{102}$}
\AXC{$\stackrel{\Pi_{141}}{\vdots}$}\noLine
\UIC{$\varphi_{119}$}
\AXC{$\varphi_{119} \to \varphi_{294}$\lbl{A1}}
\BIC{$\varphi_{294}$}
\BIC{$(c \in \mathbb{R} \to r \in \mathbb{R} \to 0 < r \to \varepsilon \in \mathbb{R} \to \varphi_{65}) \to c \in \mathbb{R} \to \varphi_{205}$}
\BIC{$c \in \mathbb{R} \to \varphi_{205}$}
\end{ptree}
\noindent$\Pi_{143}$:
\begin{ptree}
\AXC{$y = z \to x = y \to x = z$\lbl{A6}}
\AXC{$(y = z \to x = y \to x = z) \to x = y \to y = z \to x = z$\lbl{A1}}
\BIC{$x = y \to y = z \to x = z$}
\UIC{$\forall z\,(x = y \to y = z \to x = z)$}
\UIC{$\forall y\,\forall z\,(x = y \to y = z \to x = z)$}
\UIC{$\forall x\,\forall y\,\forall z\,(x = y \to y = z \to x = z)$}
\AXC{$\varphi_{91}$\lbl{A2}}
\BIC{$\forall y\,\forall z\,(\pi_{c,r}(n) - c = y \to y = z \to \pi_{c,r}(n) - c = z)$}
\end{ptree}
\noindent$\Pi_{144}$:
\begin{ptree}
\AXC{$\stackrel{\Pi_{143}}{\vdots}$}\noLine
\UIC{$\forall y\,\forall z\,(\pi_{c,r}(n) - c = y \to y = z \to \pi_{c,r}(n) - c = z)$}
\AXC{$\forall y\,\forall z\,(\pi_{c,r}(n) - c = y \to y = z \to \pi_{c,r}(n) - c = z) \to \varphi_{168}$\lbl{A2}}
\BIC{$\varphi_{168}$}
\end{ptree}
\noindent$\Pi_{145}$:
\begin{ptree}
\AXC{$z = w \to \varphi_{56}$\lbl{A6}}
\UIC{$\forall w\,(z = w \to \varphi_{56})$}
\UIC{$\forall z\,\forall w\,(z = w \to \varphi_{56})$}
\AXC{$\forall z\,\forall w\,(z = w \to \varphi_{56}) \to \varphi_{187}$\lbl{A2}}
\BIC{$\varphi_{187}$}
\AXC{$\varphi_{187} \to c + r \cdot (\underline{n} + 1)^{-1} = r \cdot (\underline{n} + 1)^{-1} + c \to \varphi_{180}$\lbl{A2}}
\BIC{$c + r \cdot (\underline{n} + 1)^{-1} = r \cdot (\underline{n} + 1)^{-1} + c \to \varphi_{180}$}
\end{ptree}
\noindent$\Pi_{146}$:
\begin{ptree}
\AXC{$\forall x\,\forall y\,(x \in \mathbb{R} \to y \in \mathbb{R} \to x + y = y + x)$\lbl{(R3)}}
\AXC{$\varphi_{92}$\lbl{A2}}
\BIC{$\forall y\,(c \in \mathbb{R} \to y \in \mathbb{R} \to c + y = y + c)$}
\AXC{$\forall y\,(c \in \mathbb{R} \to y \in \mathbb{R} \to c + y = y + c) \to c \in \mathbb{R} \to \varphi_{84}$\lbl{A2}}
\BIC{$c \in \mathbb{R} \to \varphi_{84}$}
\end{ptree}
\noindent$\Pi_{147}$:
\begin{ptree}
\AXC{$c \in \mathbb{R} \to r \in \mathbb{R} \to 0 < r \to \varepsilon \in \mathbb{R} \to \varphi_{65}$\lbl{A1}}
\AXC{$\stackrel{\Pi_{146}}{\vdots}$}\noLine
\UIC{$c \in \mathbb{R} \to \varphi_{84}$}
\AXC{$\varphi_{307}$\lbl{A1}}
\BIC{$(c \in \mathbb{R} \to r \in \mathbb{R} \to 0 < r \to \varepsilon \in \mathbb{R} \to \varphi_{65}) \to \varphi_{273}$}
\BIC{$\varphi_{273}$}
\end{ptree}
\noindent$\Pi_{148}$:
\begin{ptree}
\AXC{$\stackrel{\Pi_{140}}{\vdots}$}\noLine
\UIC{$c \in \mathbb{R} \to r \in \mathbb{R} \to 0 < r \to \varepsilon \in \mathbb{R} \to 0 < \varepsilon \to \varphi_{102}$}
\AXC{$\stackrel{\Pi_{147}}{\vdots}$}\noLine
\UIC{$\varphi_{273}$}
\BIC{$c \in \mathbb{R} \to \varphi_{206}$}
\AXC{$(c \in \mathbb{R} \to \varphi_{206}) \to \varphi_{292}$\lbl{A1}}
\BIC{$\varphi_{292}$}
\end{ptree}
\noindent$\Pi_{149}$:
\begin{ptree}
\AXC{$\stackrel{\Pi_{145}}{\vdots}$}\noLine
\UIC{$c + r \cdot (\underline{n} + 1)^{-1} = r \cdot (\underline{n} + 1)^{-1} + c \to \varphi_{180}$}
\AXC{$x = x$\lbl{A5}}
\UIC{$\forall x\,(x = x)$}
\AXC{$\varphi_{29}$\lbl{A2}}
\BIC{$c + r \cdot (\underline{n} + 1)^{-1} - c = c + r \cdot (\underline{n} + 1)^{-1} - c$}
\AXC{$\stackrel{\Pi_{148}}{\vdots}$}\noLine
\UIC{$\varphi_{292}$}
\BIC{$\varphi_{285}$}
\BIC{$c \in \mathbb{R} \to r \in \mathbb{R} \to 0 < r \to \varphi_{193}$}
\end{ptree}
\noindent$\Pi_{150}$:
\begin{ptree}
\AXC{$\varphi_{158}$\lbl{A2}}
\AXC{$\stackrel{\Pi_{61}}{\vdots}$}\noLine
\UIC{$c \in \mathbb{R} \to \varphi_{60}$}
\AXC{$(c \in \mathbb{R} \to \varphi_{60}) \to \varphi_{158} \to c \in \mathbb{R} \to \varphi_{48}$\lbl{A1}}
\BIC{$\varphi_{158} \to c \in \mathbb{R} \to \varphi_{48}$}
\BIC{$c \in \mathbb{R} \to \varphi_{48}$}
\AXC{$(c \in \mathbb{R} \to \varphi_{48}) \to \varphi_{251}$\lbl{A1}}
\BIC{$\varphi_{251}$}
\end{ptree}
\noindent$\Pi_{151}$:
\begin{ptree}
\AXC{$x = x$\lbl{A5}}
\UIC{$\forall x\,(x = x)$}
\AXC{$\forall x\,(x = x) \to \pi_{c,r}(n) - c = \pi_{c,r}(n) - c$\lbl{A2}}
\BIC{$\pi_{c,r}(n) - c = \pi_{c,r}(n) - c$}
\AXC{$c \in \mathbb{R} \to r \in \mathbb{R} \to 0 < r \to \varphi_{98}$\lbl{A1}}
\AXC{$\stackrel{\Pi_{150}}{\vdots}$}\noLine
\UIC{$\varphi_{251}$}
\BIC{$c \in \mathbb{R} \to r \in \mathbb{R} \to 0 < r \to \varepsilon \in \mathbb{R} \to 0 < \varepsilon \to \varphi_{133}$}
\AXC{$\varphi_{303}$\lbl{A1}}
\BIC{$\pi_{c,r}(n) - c = \pi_{c,r}(n) - c \to \varphi_{270}$}
\BIC{$\varphi_{270}$}
\end{ptree}
\noindent$\Pi_{152}$:
\begin{ptree}
\AXC{$z = w \to \pi_{c,r}(n) - c = z - c \to \pi_{c,r}(n) - c = w - c$\lbl{A6}}
\UIC{$\forall w\,(z = w \to \pi_{c,r}(n) - c = z - c \to \pi_{c,r}(n) - c = w - c)$}
\UIC{$\varphi_{23}$}
\AXC{$\varphi_{23} \to \varphi_{120}$\lbl{A2}}
\BIC{$\varphi_{120}$}
\AXC{$\varphi_{120} \to \pi_{c,r}(n) = c + r \cdot (\underline{n} + 1)^{-1} \to \varphi_{126}$\lbl{A2}}
\BIC{$\pi_{c,r}(n) = c + r \cdot (\underline{n} + 1)^{-1} \to \varphi_{126}$}
\AXC{$\stackrel{\Pi_{151}}{\vdots}$}\noLine
\UIC{$\varphi_{270}$}
\BIC{$c \in \mathbb{R} \to r \in \mathbb{R} \to 0 < r \to \varepsilon \in \mathbb{R} \to 0 < \varepsilon \to \varphi_{150}$}
\end{ptree}
\noindent$\Pi_{153}$:
\begin{ptree}
\AXC{$\stackrel{\Pi_{149}}{\vdots}$}\noLine
\UIC{$c \in \mathbb{R} \to r \in \mathbb{R} \to 0 < r \to \varphi_{193}$}
\AXC{$\stackrel{\Pi_{152}}{\vdots}$}\noLine
\UIC{$c \in \mathbb{R} \to r \in \mathbb{R} \to 0 < r \to \varepsilon \in \mathbb{R} \to 0 < \varepsilon \to \varphi_{150}$}
\AXC{$\varphi_{312}$\lbl{A1}}
\BIC{$(c \in \mathbb{R} \to r \in \mathbb{R} \to 0 < r \to \varphi_{193}) \to \varphi_{281}$}
\BIC{$\varphi_{281}$}
\end{ptree}
\noindent$\Pi_{154}$:
\begin{ptree}
\AXC{$\stackrel{\Pi_{142}}{\vdots}$}\noLine
\UIC{$c \in \mathbb{R} \to \varphi_{205}$}
\AXC{$\stackrel{\Pi_{144}}{\vdots}$}\noLine
\UIC{$\varphi_{168}$}
\AXC{$\varphi_{168} \to \pi_{c,r}(n) - c = c + r \cdot (\underline{n} + 1)^{-1} - c \to \varphi_{170}$\lbl{A2}}
\BIC{$\pi_{c,r}(n) - c = c + r \cdot (\underline{n} + 1)^{-1} - c \to \varphi_{170}$}
\AXC{$\stackrel{\Pi_{153}}{\vdots}$}\noLine
\UIC{$\varphi_{281}$}
\BIC{$c \in \mathbb{R} \to r \in \mathbb{R} \to 0 < r \to \varepsilon \in \mathbb{R} \to 0 < \varepsilon \to \varphi_{151}$}
\AXC{$\varphi_{310}$\lbl{A1}}
\BIC{$(c \in \mathbb{R} \to \varphi_{205}) \to \varphi_{271}$}
\BIC{$\varphi_{271}$}
\end{ptree}
\noindent$\Pi_{155}$:
\begin{ptree}
\AXC{$\stackrel{\Pi_{137}}{\vdots}$}\noLine
\UIC{$\varphi_{167}$}
\AXC{$\varphi_{167} \to \pi_{c,r}(n) - c = r \cdot (\underline{n} + 1)^{-1} + c - c \to \varphi_{125}$\lbl{A2}}
\BIC{$\pi_{c,r}(n) - c = r \cdot (\underline{n} + 1)^{-1} + c - c \to \varphi_{125}$}
\AXC{$\stackrel{\Pi_{154}}{\vdots}$}\noLine
\UIC{$\varphi_{271}$}
\BIC{$c \in \mathbb{R} \to r \in \mathbb{R} \to 0 < r \to \varepsilon \in \mathbb{R} \to 0 < \varepsilon \to \varphi_{134}$}
\AXC{$\varphi_{305}$\lbl{A1}}
\BIC{$\lvert \pi_{c,r}(n) - c \rvert = \lvert \pi_{c,r}(n) - c \rvert \to \varphi_{277}$}
\end{ptree}
\noindent$\Pi_{156}$:
\begin{ptree}
\AXC{$x = x$\lbl{A5}}
\UIC{$\forall x\,(x = x)$}
\AXC{$\forall x\,(x = x) \to \lvert \pi_{c,r}(n) - c \rvert = \lvert \pi_{c,r}(n) - c \rvert$\lbl{A2}}
\BIC{$\lvert \pi_{c,r}(n) - c \rvert = \lvert \pi_{c,r}(n) - c \rvert$}
\AXC{$\stackrel{\Pi_{155}}{\vdots}$}\noLine
\UIC{$\lvert \pi_{c,r}(n) - c \rvert = \lvert \pi_{c,r}(n) - c \rvert \to \varphi_{277}$}
\BIC{$\varphi_{277}$}
\end{ptree}
\noindent$\Pi_{157}$:
\begin{ptree}
\AXC{$\stackrel{\Pi_{133}}{\vdots}$}\noLine
\UIC{$c \in \mathbb{R} \to r \in \mathbb{R} \to 0 < r \to \varepsilon \in \mathbb{R} \to 0 < \varepsilon \to \varphi_{76}$}
\AXC{$\stackrel{\Pi_{134}}{\vdots}$}\noLine
\UIC{$\varphi_{266}$}
\BIC{$c \in \mathbb{R} \to \varphi_{200}$}
\AXC{$\stackrel{\Pi_{135}}{\vdots}$}\noLine
\UIC{$\pi_{c,r}(n) - c = r \cdot (\underline{n} + 1)^{-1} \to \varphi_{140}$}
\AXC{$\stackrel{\Pi_{156}}{\vdots}$}\noLine
\UIC{$\varphi_{277}$}
\BIC{$c \in \mathbb{R} \to r \in \mathbb{R} \to 0 < r \to \varepsilon \in \mathbb{R} \to 0 < \varepsilon \to \varphi_{152}$}
\AXC{$\varphi_{311}$\lbl{A1}}
\BIC{$(c \in \mathbb{R} \to \varphi_{200}) \to \varphi_{276}$}
\BIC{$\varphi_{276}$}
\end{ptree}
\noindent$\Pi_{158}$:
\begin{ptree}
\AXC{$\stackrel{\Pi_{128}}{\vdots}$}\noLine
\UIC{$\varphi_{22}$}
\AXC{$\varphi_{22} \to \varphi_{164}$\lbl{A2}}
\BIC{$\varphi_{164}$}
\AXC{$\varphi_{164} \to \lvert \pi_{c,r}(n) - c \rvert = \lvert r \cdot (\underline{n} + 1)^{-1} \rvert \to \varphi_{124}$\lbl{A2}}
\BIC{$\lvert \pi_{c,r}(n) - c \rvert = \lvert r \cdot (\underline{n} + 1)^{-1} \rvert \to \varphi_{124}$}
\AXC{$\stackrel{\Pi_{157}}{\vdots}$}\noLine
\UIC{$\varphi_{276}$}
\BIC{$c \in \mathbb{R} \to r \in \mathbb{R} \to 0 < r \to \varepsilon \in \mathbb{R} \to 0 < \varepsilon \to \varphi_{138}$}
\AXC{$\varphi_{293}$\lbl{A1}}
\BIC{$\varphi_{123} \to c \in \mathbb{R} \to r \in \mathbb{R} \to 0 < r \to \varepsilon \in \mathbb{R} \to 0 < \varepsilon \to \varphi_{139}$}
\end{ptree}
\noindent$\Pi_{159}$:
\begin{ptree}
\AXC{$\stackrel{\Pi_{127}}{\vdots}$}\noLine
\UIC{$\varphi_{123}$}
\AXC{$\stackrel{\Pi_{158}}{\vdots}$}\noLine
\UIC{$\varphi_{123} \to c \in \mathbb{R} \to r \in \mathbb{R} \to 0 < r \to \varepsilon \in \mathbb{R} \to 0 < \varepsilon \to \varphi_{139}$}
\BIC{$c \in \mathbb{R} \to r \in \mathbb{R} \to 0 < r \to \varepsilon \in \mathbb{R} \to 0 < \varepsilon \to \varphi_{139}$}
\AXC{$(c \in \mathbb{R} \to r \in \mathbb{R} \to 0 < r \to \varepsilon \in \mathbb{R} \to 0 < \varepsilon \to \varphi_{139}) \to \varphi_{290}$\lbl{A1}}
\BIC{$\varphi_{290}$}
\end{ptree}
\noindent$\Pi_{160}$:
\begin{ptree}
\AXC{$\stackrel{\Pi_{78}}{\vdots}$}\noLine
\UIC{$\varphi_{155}$}
\AXC{$\stackrel{\Pi_{79}}{\vdots}$}\noLine
\UIC{$r \cdot (\varepsilon \cdot r^{-1}) = \varepsilon \to \varphi_{36}$}
\AXC{$\stackrel{\Pi_{126}}{\vdots}$}\noLine
\UIC{$\varphi_{226}$}
\BIC{$c \in \mathbb{R} \to r \in \mathbb{R} \to 0 < r \to \varepsilon \in \mathbb{R} \to 0 < \varepsilon \to \varphi_{99}$}
\AXC{$\stackrel{\Pi_{159}}{\vdots}$}\noLine
\UIC{$\varphi_{290}$}
\BIC{$\varphi_{155} \to c \in \mathbb{R} \to r \in \mathbb{R} \to 0 < r \to \varepsilon \in \mathbb{R} \to 0 < \varepsilon \to \varphi_{108}$}
\BIC{$c \in \mathbb{R} \to r \in \mathbb{R} \to 0 < r \to \varepsilon \in \mathbb{R} \to 0 < \varepsilon \to \varphi_{108}$}
\AXC{$\varphi_{265}$\lbl{A1}}
\BIC{$\varphi_{122} \to n \in \mathbb{N} \to K \leq n \to \lvert \pi_{c,r}(n) - c \rvert < \varepsilon$}
\UIC{$\forall n\,(\varphi_{122} \to n \in \mathbb{N} \to K \leq n \to \lvert \pi_{c,r}(n) - c \rvert < \varepsilon)$}
\end{ptree}
\noindent$\Pi_{161}$:
\begin{ptree}
\AXC{$\varphi_{268}$\lbl{A3}}
\AXC{$\stackrel{\Pi_{160}}{\vdots}$}\noLine
\UIC{$\forall n\,(\varphi_{122} \to n \in \mathbb{N} \to K \leq n \to \lvert \pi_{c,r}(n) - c \rvert < \varepsilon)$}
\AXC{$\forall n\,(\varphi_{122} \to n \in \mathbb{N} \to K \leq n \to \lvert \pi_{c,r}(n) - c \rvert < \varepsilon) \to \varphi_{268} \to \varphi_{283}$\lbl{A1}}
\BIC{$\varphi_{268} \to \varphi_{283}$}
\BIC{$\varphi_{283}$}
\end{ptree}
\noindent$\Pi_{162}$:
\begin{ptree}
\AXC{$c \in \mathbb{R} \to r \in \mathbb{R} \to 0 < r \to \varepsilon \in \mathbb{R} \to 0 < \varepsilon \to \varphi_{72}$\lbl{A1}}
\AXC{$\varphi_{228}$\lbl{A1}}
\BIC{$c \in \mathbb{R} \to r \in \mathbb{R} \to \varphi_{80}$}
\AXC{$(c \in \mathbb{R} \to r \in \mathbb{R} \to \varphi_{80}) \to \varphi_{274}$\lbl{A1}}
\BIC{$\varphi_{274}$}
\end{ptree}
\noindent$\Pi_{163}$:
\begin{ptree}
\AXC{$\varphi_{188}$\lbl{A2}}
\AXC{$\varphi_{122} \to \forall n\,(\varphi_{122})$\lbl{A4}}
\AXC{$\stackrel{\Pi_{161}}{\vdots}$}\noLine
\UIC{$\varphi_{283}$}
\BIC{$c \in \mathbb{R} \to r \in \mathbb{R} \to 0 < r \to \varepsilon \in \mathbb{R} \to 0 < \varepsilon \to \varphi_{114}$}
\AXC{$\stackrel{\Pi_{162}}{\vdots}$}\noLine
\UIC{$\varphi_{274}$}
\BIC{$c \in \mathbb{R} \to r \in \mathbb{R} \to 0 < r \to \varepsilon \in \mathbb{R} \to 0 < \varepsilon \to \varphi_{136}$}
\AXC{$\varphi_{296}$\lbl{A1}}
\BIC{$\varphi_{188} \to c \in \mathbb{R} \to r \in \mathbb{R} \to 0 < r \to \varepsilon \in \mathbb{R} \to 0 < \varepsilon \to \varphi_{143}$}
\BIC{$c \in \mathbb{R} \to r \in \mathbb{R} \to 0 < r \to \varepsilon \in \mathbb{R} \to 0 < \varepsilon \to \varphi_{143}$}
\AXC{$\varphi_{280}$\lbl{A1}}
\BIC{$\varphi_{179} \to \neg (K \in \mathbb{N} \land \underline{K}^{-1} < \varepsilon \cdot r^{-1})$}
\UIC{$\forall K\,(\varphi_{179} \to \neg (K \in \mathbb{N} \land \underline{K}^{-1} < \varepsilon \cdot r^{-1}))$}
\end{ptree}
\noindent$\Pi_{164}$:
\begin{ptree}
\AXC{$\forall x \in \mathbb{R}\,(0 < x \to 0 < x^{-1})$\lbl{[F34]}}
\AXC{$\varphi_{20}$\lbl{A2}}
\BIC{$r \in \mathbb{R} \to 0 < r \to 0 < r^{-1}$}
\AXC{$(r \in \mathbb{R} \to 0 < r \to 0 < r^{-1}) \to \varphi_{212}$\lbl{A1}}
\BIC{$\varphi_{212}$}
\end{ptree}
\noindent$\Pi_{165}$:
\begin{ptree}
\AXC{$\forall x \in \mathbb{R}\,(x \neq 0 \to x^{-1} \in \mathbb{R})$\lbl{(R1)}}
\AXC{$\varphi_{33}$\lbl{A2}}
\BIC{$r \in \mathbb{R} \to r \neq 0 \to r^{-1} \in \mathbb{R}$}
\AXC{$(r \in \mathbb{R} \to r \neq 0 \to r^{-1} \in \mathbb{R}) \to \varphi_{213}$\lbl{A1}}
\BIC{$\varphi_{213}$}
\end{ptree}
\noindent$\Pi_{166}$:
\begin{ptree}
\AXC{$\stackrel{\Pi_{100}}{\vdots}$}\noLine
\UIC{$c \in \mathbb{R} \to r \in \mathbb{R} \to 0 < r \to r \neq 0$}
\AXC{$c \in \mathbb{R} \to r \in \mathbb{R} \to 0 < r \to \varepsilon \in \mathbb{R} \to 0 < \varepsilon \to r \in \mathbb{R}$\lbl{A1}}
\AXC{$\stackrel{\Pi_{165}}{\vdots}$}\noLine
\UIC{$\varphi_{213}$}
\BIC{$\varphi_{154}$}
\BIC{$c \in \mathbb{R} \to r \in \mathbb{R} \to 0 < r \to \varepsilon \in \mathbb{R} \to 0 < \varepsilon \to r^{-1} \in \mathbb{R}$}
\end{ptree}
\noindent$\Pi_{167}$:
\begin{ptree}
\AXC{$\forall x\,\forall y\,(x \in \mathbb{R} \to y \in \mathbb{R} \to 0 < x \to 0 < y \to 0 < x \cdot y)$\lbl{(R15)}}
\AXC{$\varphi_{174}$\lbl{A2}}
\BIC{$\forall y\,(\varepsilon \in \mathbb{R} \to y \in \mathbb{R} \to 0 < \varepsilon \to 0 < y \to 0 < \varepsilon \cdot y)$}
\AXC{$\varphi_{189}$\lbl{A2}}
\BIC{$\varepsilon \in \mathbb{R} \to r^{-1} \in \mathbb{R} \to 0 < \varepsilon \to 0 < r^{-1} \to 0 < \varepsilon \cdot r^{-1}$}
\AXC{$\varphi_{291}$\lbl{A1}}
\BIC{$(c \in \mathbb{R} \to r \in \mathbb{R} \to 0 < r \to \varepsilon \in \mathbb{R} \to \varepsilon \in \mathbb{R}) \to \varphi_{258}$}
\end{ptree}
\noindent$\Pi_{168}$:
\begin{ptree}
\AXC{$c \in \mathbb{R} \to r \in \mathbb{R} \to 0 < r \to \varepsilon \in \mathbb{R} \to \varepsilon \in \mathbb{R}$\lbl{A1}}
\AXC{$\stackrel{\Pi_{167}}{\vdots}$}\noLine
\UIC{$(c \in \mathbb{R} \to r \in \mathbb{R} \to 0 < r \to \varepsilon \in \mathbb{R} \to \varepsilon \in \mathbb{R}) \to \varphi_{258}$}
\BIC{$\varphi_{258}$}
\end{ptree}
\noindent$\Pi_{169}$:
\begin{ptree}
\AXC{$c \in \mathbb{R} \to r \in \mathbb{R} \to 0 < r \to \varepsilon \in \mathbb{R} \to 0 < \varepsilon \to 0 < \varepsilon$\lbl{A1}}
\AXC{$\stackrel{\Pi_{166}}{\vdots}$}\noLine
\UIC{$c \in \mathbb{R} \to r \in \mathbb{R} \to 0 < r \to \varepsilon \in \mathbb{R} \to 0 < \varepsilon \to r^{-1} \in \mathbb{R}$}
\AXC{$\stackrel{\Pi_{168}}{\vdots}$}\noLine
\UIC{$\varphi_{258}$}
\BIC{$(c \in \mathbb{R} \to r \in \mathbb{R} \to 0 < r \to \varepsilon \in \mathbb{R} \to 0 < \varepsilon \to 0 < \varepsilon) \to \varphi_{195}$}
\BIC{$\varphi_{195}$}
\end{ptree}
\noindent$\Pi_{170}$:
\begin{ptree}
\AXC{$c \in \mathbb{R} \to r \in \mathbb{R} \to 0 < r \to 0 < r$\lbl{A1}}
\AXC{$c \in \mathbb{R} \to r \in \mathbb{R} \to 0 < r \to \varepsilon \in \mathbb{R} \to 0 < \varepsilon \to r \in \mathbb{R}$\lbl{A1}}
\AXC{$\stackrel{\Pi_{164}}{\vdots}$}\noLine
\UIC{$\varphi_{212}$}
\BIC{$\varphi_{149}$}
\BIC{$c \in \mathbb{R} \to r \in \mathbb{R} \to 0 < r \to \varepsilon \in \mathbb{R} \to 0 < \varepsilon \to 0 < r^{-1}$}
\AXC{$\stackrel{\Pi_{169}}{\vdots}$}\noLine
\UIC{$\varphi_{195}$}
\BIC{$c \in \mathbb{R} \to r \in \mathbb{R} \to 0 < r \to \varepsilon \in \mathbb{R} \to 0 < \varepsilon \to 0 < \varepsilon \cdot r^{-1}$}
\end{ptree}
\noindent$\Pi_{171}$:
\begin{ptree}
\AXC{$\forall x\,\forall y\,(x \in \mathbb{R} \to y \in \mathbb{R} \to x \cdot y \in \mathbb{R})$\lbl{(R1)}}
\AXC{$\varphi_{111}$\lbl{A2}}
\BIC{$\forall y\,(\varepsilon \in \mathbb{R} \to y \in \mathbb{R} \to \varepsilon \cdot y \in \mathbb{R})$}
\AXC{$\varphi_{131}$\lbl{A2}}
\BIC{$\varepsilon \in \mathbb{R} \to r^{-1} \in \mathbb{R} \to \varepsilon \cdot r^{-1} \in \mathbb{R}$}
\AXC{$\varphi_{242}$\lbl{A1}}
\BIC{$(c \in \mathbb{R} \to r \in \mathbb{R} \to 0 < r \to \varepsilon \in \mathbb{R} \to \varepsilon \in \mathbb{R}) \to \varphi_{203}$}
\end{ptree}
\noindent$\Pi_{172}$:
\begin{ptree}
\AXC{$c \in \mathbb{R} \to r \in \mathbb{R} \to 0 < r \to \varepsilon \in \mathbb{R} \to \varepsilon \in \mathbb{R}$\lbl{A1}}
\AXC{$\stackrel{\Pi_{171}}{\vdots}$}\noLine
\UIC{$(c \in \mathbb{R} \to r \in \mathbb{R} \to 0 < r \to \varepsilon \in \mathbb{R} \to \varepsilon \in \mathbb{R}) \to \varphi_{203}$}
\BIC{$\varphi_{203}$}
\end{ptree}
\noindent$\Pi_{173}$:
\begin{ptree}
\AXC{$\forall \varepsilon \in \mathbb{R}\,(0 < \varepsilon \to \exists K \in \mathbb{N}\,(\underline{K}^{-1} < \varepsilon))$\lbl{[C24]}}
\AXC{$\forall \varepsilon \in \mathbb{R}\,(0 < \varepsilon \to \exists K \in \mathbb{N}\,(\underline{K}^{-1} < \varepsilon)) \to \varphi_{94}$\lbl{A2}}
\BIC{$\varphi_{94}$}
\AXC{$\varphi_{94} \to \varphi_{62} \to \varphi_{211}$\lbl{A1}}
\BIC{$\varphi_{62} \to \varphi_{211}$}
\end{ptree}
\noindent$\Pi_{174}$:
\begin{ptree}
\AXC{$\stackrel{\Pi_{166}}{\vdots}$}\noLine
\UIC{$c \in \mathbb{R} \to r \in \mathbb{R} \to 0 < r \to \varepsilon \in \mathbb{R} \to 0 < \varepsilon \to r^{-1} \in \mathbb{R}$}
\AXC{$\stackrel{\Pi_{172}}{\vdots}$}\noLine
\UIC{$\varphi_{203}$}
\BIC{$\varphi_{62}$}
\AXC{$\stackrel{\Pi_{173}}{\vdots}$}\noLine
\UIC{$\varphi_{62} \to \varphi_{211}$}
\BIC{$\varphi_{211}$}
\end{ptree}
\noindent$\Pi_{175}$:
\begin{ptree}
\AXC{$\stackrel{\Pi_{170}}{\vdots}$}\noLine
\UIC{$c \in \mathbb{R} \to r \in \mathbb{R} \to 0 < r \to \varepsilon \in \mathbb{R} \to 0 < \varepsilon \to 0 < \varepsilon \cdot r^{-1}$}
\AXC{$\stackrel{\Pi_{174}}{\vdots}$}\noLine
\UIC{$\varphi_{211}$}
\BIC{$c \in \mathbb{R} \to \varphi_{100}$}
\AXC{$(c \in \mathbb{R} \to \varphi_{100}) \to \varphi_{316}$\lbl{A1}}
\BIC{$\varphi_{316}$}
\end{ptree}
\noindent$\Pi_{176}$:
\begin{ptree}
\AXC{$\varphi_{179} \to \forall K\,(\varphi_{179})$\lbl{A4}}
\AXC{$\varphi_{282}$\lbl{A3}}
\AXC{$\stackrel{\Pi_{163}}{\vdots}$}\noLine
\UIC{$\forall K\,(\varphi_{179} \to \neg (K \in \mathbb{N} \land \underline{K}^{-1} < \varepsilon \cdot r^{-1}))$}
\AXC{$\stackrel{\Pi_{175}}{\vdots}$}\noLine
\UIC{$\varphi_{316}$}
\BIC{$\varphi_{282} \to (\varphi_{179} \to \forall K\,(\varphi_{179})) \to c \in \mathbb{R} \to r \in \mathbb{R} \to 0 < r \to \varphi_{109}$}
\BIC{$(\varphi_{179} \to \forall K\,(\varphi_{179})) \to c \in \mathbb{R} \to r \in \mathbb{R} \to 0 < r \to \varphi_{109}$}
\BIC{$c \in \mathbb{R} \to r \in \mathbb{R} \to 0 < r \to \varphi_{109}$}
\AXC{$\varphi_{230}$\lbl{A1}}
\BIC{$c \in \mathbb{R} \land r \in \mathbb{R} \land 0 < r \to \varphi_{109}$}
\UIC{$\forall \varepsilon\,(c \in \mathbb{R} \land r \in \mathbb{R} \land 0 < r \to \varphi_{109})$}
\AXC{$\varphi_{309}$\lbl{A1}}
\BIC{$\varphi_{239} \to \varphi_{19} \to c \in \mathbb{R} \to r \in \mathbb{R} \to 0 < r \to \varphi_{117}$}
\end{ptree}
\noindent$\Pi_{177}$:
\begin{ptree}
\AXC{$\stackrel{\Pi_{60}}{\vdots}$}\noLine
\UIC{$c \in \mathbb{R} \to r \in \mathbb{R} \to 0 < r \to \varphi_{45}$}
\AXC{$\varphi_{207}$\lbl{A1}}
\BIC{$c \in \mathbb{R} \to r \in \mathbb{R} \to 0 < r \to \mathrm{Seq}(\pi_{c,r})$}
\end{ptree}
\noindent$\Pi_{178}$:
\begin{ptree}
\AXC{$\varphi_{183}$\lbl{DefConv}}
\AXC{$\varphi_{183} \to \varphi_{199}$\lbl{A2}}
\BIC{$\varphi_{199}$}
\AXC{$\varphi_{199} \to (\mathrm{Conv}(\pi_{c,r},c) \leftrightarrow \mathrm{Seq}(\pi_{c,r}) \land c \in \mathbb{R} \land \varphi_{117})$\lbl{A2}}
\BIC{$\mathrm{Conv}(\pi_{c,r},c) \leftrightarrow \mathrm{Seq}(\pi_{c,r}) \land c \in \mathbb{R} \land \varphi_{117}$}
\end{ptree}
\noindent$\Pi_{179}$:
\begin{ptree}
\AXC{$\stackrel{\Pi_{178}}{\vdots}$}\noLine
\UIC{$\mathrm{Conv}(\pi_{c,r},c) \leftrightarrow \mathrm{Seq}(\pi_{c,r}) \land c \in \mathbb{R} \land \varphi_{117}$}
\AXC{$(\mathrm{Conv}(\pi_{c,r},c) \leftrightarrow \mathrm{Seq}(\pi_{c,r}) \land c \in \mathbb{R} \land \varphi_{117}) \to \varphi_{250}$\lbl{A1}}
\BIC{$\varphi_{250}$}
\end{ptree}
\noindent$\Pi_{180}$:
\begin{ptree}
\AXC{$c \in \mathbb{R} \to r \in \mathbb{R} \to 0 < r \to c \in \mathbb{R}$\lbl{A1}}
\AXC{$\stackrel{\Pi_{177}}{\vdots}$}\noLine
\UIC{$c \in \mathbb{R} \to r \in \mathbb{R} \to 0 < r \to \mathrm{Seq}(\pi_{c,r})$}
\AXC{$\stackrel{\Pi_{179}}{\vdots}$}\noLine
\UIC{$\varphi_{250}$}
\BIC{$(c \in \mathbb{R} \to r \in \mathbb{R} \to 0 < r \to c \in \mathbb{R}) \to \varphi_{214}$}
\BIC{$\varphi_{214}$}
\end{ptree}
\noindent$\Pi_{181}$:
\begin{ptree}
\AXC{$\stackrel{\Pi_{177}}{\vdots}$}\noLine
\UIC{$c \in \mathbb{R} \to r \in \mathbb{R} \to 0 < r \to \mathrm{Seq}(\pi_{c,r})$}
\AXC{$(c \in \mathbb{R} \to r \in \mathbb{R} \to 0 < r \to \mathrm{Seq}(\pi_{c,r})) \to \varphi_{272}$\lbl{A1}}
\BIC{$\varphi_{272}$}
\end{ptree}
\noindent$\Pi_{182}$:
\begin{ptree}
\AXC{$\varphi_{19}$\lbl{A4}}
\AXC{$\varphi_{239}$\lbl{A3}}
\AXC{$\stackrel{\Pi_{176}}{\vdots}$}\noLine
\UIC{$\varphi_{239} \to \varphi_{19} \to c \in \mathbb{R} \to r \in \mathbb{R} \to 0 < r \to \varphi_{117}$}
\BIC{$\varphi_{19} \to c \in \mathbb{R} \to r \in \mathbb{R} \to 0 < r \to \varphi_{117}$}
\BIC{$c \in \mathbb{R} \to r \in \mathbb{R} \to 0 < r \to \varphi_{117}$}
\AXC{$\stackrel{\Pi_{180}}{\vdots}$}\noLine
\UIC{$\varphi_{214}$}
\BIC{$c \in \mathbb{R} \to r \in \mathbb{R} \to 0 < r \to \mathrm{Conv}(\pi_{c,r},c)$}
\AXC{$\stackrel{\Pi_{181}}{\vdots}$}\noLine
\UIC{$\varphi_{272}$}
\BIC{$(c \in \mathbb{R} \to r \in \mathbb{R} \to \varphi_{96}) \to c \in \mathbb{R} \to r \in \mathbb{R} \to 0 < r \to \varphi_{162}$}
\end{ptree}
\noindent$\Pi_{183}$:
\begin{ptree}
\AXC{$\varphi_{13}$\lbl{A4}}
\AXC{$\stackrel{\Pi_{77}}{\vdots}$}\noLine
\UIC{$\varphi_{13} \to c \in \mathbb{R} \to r \in \mathbb{R} \to \varphi_{96}$}
\BIC{$c \in \mathbb{R} \to r \in \mathbb{R} \to \varphi_{96}$}
\AXC{$\stackrel{\Pi_{182}}{\vdots}$}\noLine
\UIC{$(c \in \mathbb{R} \to r \in \mathbb{R} \to \varphi_{96}) \to c \in \mathbb{R} \to r \in \mathbb{R} \to 0 < r \to \varphi_{162}$}
\BIC{$c \in \mathbb{R} \to r \in \mathbb{R} \to 0 < r \to \varphi_{162}$}
\AXC{$\varphi_{297}$\lbl{A1}}
\BIC{$(\varphi_{83} \to \neg \varphi_{162}) \to c \in \mathbb{R} \to r \in \mathbb{R} \to 0 < r \to \varphi_{82}$}
\end{ptree}
\noindent 主木:
\begin{ptree}
\AXC{$\varphi_{83} \to \neg \varphi_{162}$\lbl{A2}}
\AXC{$\stackrel{\Pi_{183}}{\vdots}$}\noLine
\UIC{$(\varphi_{83} \to \neg \varphi_{162}) \to c \in \mathbb{R} \to r \in \mathbb{R} \to 0 < r \to \varphi_{82}$}
\BIC{$c \in \mathbb{R} \to r \in \mathbb{R} \to 0 < r \to \varphi_{82}$}
\UIC{$\forall r\,(c \in \mathbb{R} \to r \in \mathbb{R} \to 0 < r \to \varphi_{82})$}
\UIC{$\forall c\,\forall r\,(c \in \mathbb{R} \to r \in \mathbb{R} \to 0 < r \to \varphi_{82})$}
\end{ptree}
\begin{abbrevlist}
\item $\varphi_{1} :\equiv \forall x\,\forall y\,(x = y \to y = x) \to \forall y\,(\underline{n} + 1 = y \to y = \underline{n} + 1)$
\item $\varphi_{2} :\equiv \forall y\,\forall z\,(r \cdot (\varepsilon \cdot r^{-1}) = y \to y = z \to r \cdot (\varepsilon \cdot r^{-1}) = z)$
\item $\varphi_{3} :\equiv \forall z\,\forall w\,(z = w \to (\underline{n} + 1)^{-1} < z \to (\underline{n} + 1)^{-1} < w)$
\item $\varphi_{4} :\equiv \forall y\,\forall z\,((\underline{n} + 1)^{-1} < y \to y \leq z \to (\underline{n} + 1)^{-1} < z)$
\item $\varphi_{5} :\equiv \forall y\,\forall z\,((\underline{n} + 1)^{-1} < y \to y < z \to (\underline{n} + 1)^{-1} < z)$
\item $\varphi_{6} :\equiv \forall w\,(r \cdot (\underline{n} + 1)^{-1} = w \to r \cdot (\underline{n} + 1)^{-1} < r \to w < r)$
\item $\varphi_{7} :\equiv \forall w\,(z = w \to r \cdot (\underline{n} + 1)^{-1} < z \to r \cdot (\underline{n} + 1)^{-1} < w)$
\item $\varphi_{8} :\equiv \forall w\,(1^{-1} = w \to (\underline{n} + 1)^{-1} < 1^{-1} \to (\underline{n} + 1)^{-1} < w)$
\item $\varphi_{9} :\equiv \forall x\,\forall y\,(x \in \mathbb{R} \to y \in \mathbb{R} \to 0 < x \to x < y \to y^{-1} < x^{-1})$
\item $\varphi_{10} :\equiv \forall w\,(r \cdot (\underline{n} + 1)^{-1} = w \to 0 < r \cdot (\underline{n} + 1)^{-1} \to 0 < w)$
\item $\varphi_{11} :\equiv \forall y\,(\underline{n} + 1 = y \to y = \underline{n} + 1) \to \underline{n} + 1 = 0 \to 0 = \underline{n} + 1$
\item $\varphi_{12} :\equiv n \in \mathbb{N} \to K \leq n \to \lvert r \cdot (\underline{n} + 1)^{-1} \rvert = r \cdot (\underline{n} + 1)^{-1}$
\item $\varphi_{13} :\equiv c \in \mathbb{R} \land r \in \mathbb{R} \land 0 < r \to \forall n\,(c \in \mathbb{R} \land r \in \mathbb{R} \land 0 < r)$
\item $\varphi_{14} :\equiv \forall w\,(r \cdot (\underline{n} + 1)^{-1} = w \to r \cdot (\underline{n} + 1)^{-1} < \varepsilon \to w < \varepsilon)$
\item $\varphi_{15} :\equiv c \in \mathbb{R} \to r \in \mathbb{R} \to 0 < r \to n \in \mathbb{N} \to 0 \leq r \cdot (\underline{n} + 1)^{-1}$
\item $\varphi_{16} :\equiv 0 < r \to n \in \mathbb{N} \to \lvert r \cdot (\underline{n} + 1)^{-1} \rvert = r \cdot (\underline{n} + 1)^{-1}$
\item $\varphi_{17} :\equiv \forall N\,(\neg (N \in \mathbb{N} \land \forall n \in \mathbb{N}\,(N \leq n \to \lvert \pi_{c,r}(n) - c \rvert < \varepsilon)))$
\item $\varphi_{18} :\equiv K \leq n \to c + r \cdot (\underline{n} + 1)^{-1} - c = r \cdot (\underline{n} + 1)^{-1} + c - c$
\item $\varphi_{19} :\equiv c \in \mathbb{R} \land r \in \mathbb{R} \land 0 < r \to \forall \varepsilon\,(c \in \mathbb{R} \land r \in \mathbb{R} \land 0 < r)$
\item $\varphi_{20} :\equiv \forall x \in \mathbb{R}\,(0 < x \to 0 < x^{-1}) \to r \in \mathbb{R} \to 0 < r \to 0 < r^{-1}$
\item $\varphi_{21} :\equiv r \cdot 1 = r \to r \cdot (\underline{n} + 1)^{-1} < r \cdot 1 \to r \cdot (\underline{n} + 1)^{-1} < r$
\item $\varphi_{22} :\equiv \forall y\,\forall z\,(\lvert \pi_{c,r}(n) - c \rvert = y \to y = z \to \lvert \pi_{c,r}(n) - c \rvert = z)$
\item $\varphi_{23} :\equiv \forall z\,\forall w\,(z = w \to \pi_{c,r}(n) - c = z - c \to \pi_{c,r}(n) - c = w - c)$
\item $\varphi_{24} :\equiv n \in \mathbb{N} \to K \leq n \to r \cdot (\underline{n} + 1)^{-1} + c - c = r \cdot (\underline{n} + 1)^{-1}$
\item $\varphi_{25} :\equiv n \in \mathbb{N} \to K \leq n \to c + r \cdot (\underline{n} + 1)^{-1} = r \cdot (\underline{n} + 1)^{-1} + c$
\item $\varphi_{26} :\equiv c \in \mathbb{R} \to r \in \mathbb{R} \to 0 < r \to n \in \mathbb{N} \to (\underline{n} + 1)^{-1} < 1^{-1}$
\item $\varphi_{27} :\equiv c \in \mathbb{R} \to r \in \mathbb{R} \to 0 < r \to n \in \mathbb{N} \to r \cdot (\underline{n} + 1)^{-1} < r$
\item $\varphi_{28} :\equiv c \in \mathbb{R} \to r \in \mathbb{R} \to 0 < r \to n \in \mathbb{N} \to 0 < r \cdot (\underline{n} + 1)^{-1}$
\item $\varphi_{29} :\equiv \forall x\,(x = x) \to c + r \cdot (\underline{n} + 1)^{-1} - c = c + r \cdot (\underline{n} + 1)^{-1} - c$
\item $\varphi_{30} :\equiv r \in \mathbb{R} \to \varepsilon \cdot r^{-1} \in \mathbb{R} \to r \cdot (\varepsilon \cdot r^{-1}) = \varepsilon \cdot r^{-1} \cdot r$
\item $\varphi_{31} :\equiv \forall w\,(r \cdot 1 = w \to r \cdot (\underline{n} + 1)^{-1} < r \cdot 1 \to r \cdot (\underline{n} + 1)^{-1} < w)$
\item $\varphi_{32} :\equiv \forall z\,\forall w\,(z = w \to 1 < z \to 1 < w) \to \forall w\,(1 + \underline{n} = w \to 1 < 1 + \underline{n} \to 1 < w)$
\item $\varphi_{33} :\equiv \forall x \in \mathbb{R}\,(x \neq 0 \to x^{-1} \in \mathbb{R}) \to r \in \mathbb{R} \to r \neq 0 \to r^{-1} \in \mathbb{R}$
\item $\varphi_{34} :\equiv 0 < r \to n \in \mathbb{N} \to r \cdot (\underline{n} + 1)^{-1} + c - c = r \cdot (\underline{n} + 1)^{-1}$
\item $\varphi_{35} :\equiv 0 < r \to n \in \mathbb{N} \to c + r \cdot (\underline{n} + 1)^{-1} = r \cdot (\underline{n} + 1)^{-1} + c$
\item $\varphi_{36} :\equiv r \cdot (\underline{n} + 1)^{-1} < r \cdot (\varepsilon \cdot r^{-1}) \to r \cdot (\underline{n} + 1)^{-1} < \varepsilon$
\item $\varphi_{37} :\equiv 0 < \varepsilon \to K \in \mathbb{N} \land \underline{K}^{-1} < \varepsilon \cdot r^{-1} \to n \in \mathbb{N} \to K \leq n \to K \leq n$
\item $\varphi_{38} :\equiv 0 \leq r \cdot (\underline{n} + 1)^{-1} \to \lvert r \cdot (\underline{n} + 1)^{-1} \rvert = r \cdot (\underline{n} + 1)^{-1}$
\item $\varphi_{39} :\equiv n \in \mathbb{N} \to c + r \cdot (\underline{n} + 1)^{-1} - c = r \cdot (\underline{n} + 1)^{-1} + c - c$
\item $\varphi_{40} :\equiv K \in \mathbb{N} \land \underline{K}^{-1} < \varepsilon \cdot r^{-1} \to n \in \mathbb{N} \to K \leq n \to r^{-1} \in \mathbb{R}$
\item $\varphi_{41} :\equiv K \in \mathbb{N} \land \underline{K}^{-1} < \varepsilon \cdot r^{-1} \to n \in \mathbb{N} \to K \leq n \to 0 \neq \underline{n} + 1$
\item $\varphi_{42} :\equiv K \in \mathbb{N} \land \underline{K}^{-1} < \varepsilon \cdot r^{-1} \to n \in \mathbb{N} \to K \leq n \to \underline{n} + 1 \neq 0$
\item $\varphi_{43} :\equiv c \in \mathbb{R} \to r \in \mathbb{R} \to 0 < r \to n \in \mathbb{N} \to r \cdot (\underline{n} + 1)^{-1} \in \mathbb{R}$
\item $\varphi_{44} :\equiv \forall z\,((\underline{n} + 1)^{-1} < \underline{n}^{-1} \to \underline{n}^{-1} \leq z \to (\underline{n} + 1)^{-1} < z)$
\item $\varphi_{45} :\equiv \mathrm{Seq}(\pi_{c,r}) \land \forall n \in \mathbb{N}\,(\pi_{c,r}(n) = c + r \cdot (\underline{n} + 1)^{-1})$
\item $\varphi_{46} :\equiv c \in \mathbb{R} \to r \in \mathbb{R} \to 0 < r \to n \in \mathbb{N} \to r \cdot (\underline{n} + 1)^{-1} < r \cdot 1$
\item $\varphi_{47} :\equiv r \in \mathbb{R} \to 0 < r \to (\underline{n} + 1)^{-1} < 1 \to r \cdot (\underline{n} + 1)^{-1} < r \cdot 1$
\item $\varphi_{48} :\equiv r \in \mathbb{R} \to 0 < r \to n \in \mathbb{N} \to \pi_{c,r}(n) = c + r \cdot (\underline{n} + 1)^{-1}$
\item $\varphi_{49} :\equiv r \in \mathbb{R} \to 0 < r \to n \in \mathbb{N} \to \pi_{c,r}(n) - c = r \cdot (\underline{n} + 1)^{-1}$
\item $\varphi_{50} :\equiv 0 < \varepsilon \to K \in \mathbb{N} \land \underline{K}^{-1} < \varepsilon \cdot r^{-1} \to n \in \mathbb{N} \to K \leq n \to 0 < 1$
\item $\varphi_{51} :\equiv K \in \mathbb{N} \land \underline{K}^{-1} < \varepsilon \cdot r^{-1} \to n \in \mathbb{N} \to K \leq n \to 0 < \underline{n} + 1$
\item $\varphi_{52} :\equiv \forall w\,(z = w \to \lvert \pi_{c,r}(n) - c \rvert = \lvert z \rvert \to \lvert \pi_{c,r}(n) - c \rvert = \lvert w \rvert)$
\item $\varphi_{53} :\equiv \underline{n} + 1 \in \mathbb{R} \to 0 < 1 \to 1 < \underline{n} + 1 \to (\underline{n} + 1)^{-1} < 1^{-1}$
\item $\varphi_{54} :\equiv 1 \in \mathbb{R} \to c \in \mathbb{R} \to r \in \mathbb{R} \to 0 < r \to n \in \mathbb{N} \to \underline{n} + 1 \in \mathbb{R}$
\item $\varphi_{55} :\equiv \forall z\,((\underline{n} + 1)^{-1} < \underline{K}^{-1} \to \underline{K}^{-1} < z \to (\underline{n} + 1)^{-1} < z)$
\item $\varphi_{56} :\equiv c + r \cdot (\underline{n} + 1)^{-1} - c = z - c \to c + r \cdot (\underline{n} + 1)^{-1} - c = w - c$
\item $\varphi_{57} :\equiv (\underline{n} + 1)^{-1} < \varepsilon \cdot r^{-1} \to r \cdot (\underline{n} + 1)^{-1} < r \cdot (\varepsilon \cdot r^{-1})$
\item $\varphi_{58} :\equiv \forall x\,\forall y\,\forall z\,(x \in \mathbb{R} \to y \in \mathbb{R} \to z \in \mathbb{R} \to 0 < z \to x < y \to z \cdot x < z \cdot y)$
\item $\varphi_{59} :\equiv \forall x\,\forall y\,(x = y \to y = x) \to \forall y\,(\lvert \pi_{c,r}(n) - c \rvert = y \to y = \lvert \pi_{c,r}(n) - c \rvert)$
\item $\varphi_{60} :\equiv r \in \mathbb{R} \to 0 < r \to \forall n \in \mathbb{N}\,(\pi_{c,r}(n) = c + r \cdot (\underline{n} + 1)^{-1})$
\item $\varphi_{61} :\equiv (c \in \mathbb{R} \to r \in \mathbb{R} \to 0 < r \to 0 < r) \to c \in \mathbb{R} \to r \in \mathbb{R} \to 0 < r \to 0 \neq r$
\item $\varphi_{62} :\equiv c \in \mathbb{R} \to r \in \mathbb{R} \to 0 < r \to \varepsilon \in \mathbb{R} \to 0 < \varepsilon \to \varepsilon \cdot r^{-1} \in \mathbb{R}$
\item $\varphi_{63} :\equiv K \in \mathbb{N} \land \underline{K}^{-1} < \varepsilon \cdot r^{-1} \to n \in \mathbb{N} \to K \leq n \to \underline{n} + 1 \in \mathbb{R}$
\item $\varphi_{64} :\equiv 0 < \varepsilon \to K \in \mathbb{N} \land \underline{K}^{-1} < \varepsilon \cdot r^{-1} \to n \in \mathbb{N} \to K \leq n \to r \in \mathbb{R}$
\item $\varphi_{65} :\equiv 0 < \varepsilon \to K \in \mathbb{N} \land \underline{K}^{-1} < \varepsilon \cdot r^{-1} \to n \in \mathbb{N} \to K \leq n \to c \in \mathbb{R}$
\item $\varphi_{66} :\equiv r \in \mathbb{R} \to 0 < r \to n \in \mathbb{N} \to \lvert \pi_{c,r}(n) - c \rvert = r \cdot (\underline{n} + 1)^{-1}$
\item $\varphi_{67} :\equiv r \in \mathbb{R} \to 0 < r \to n \in \mathbb{N} \to r \cdot (\underline{n} + 1)^{-1} = \lvert \pi_{c,r}(n) - c \rvert$
\item $\varphi_{68} :\equiv K \in \mathbb{N} \land \underline{K}^{-1} < \varepsilon \cdot r^{-1} \to n \in \mathbb{N} \to K \leq n \to \varepsilon \cdot r^{-1} \in \mathbb{R}$
\item $\varphi_{69} :\equiv K \in \mathbb{N} \land \underline{K}^{-1} < \varepsilon \cdot r^{-1} \to n \in \mathbb{N} \to K \leq n \to \varepsilon \cdot r^{-1} \cdot r = \varepsilon$
\item $\varphi_{70} :\equiv K \in \mathbb{N} \land \underline{K}^{-1} < \varepsilon \cdot r^{-1} \to n \in \mathbb{N} \to K \leq n \to r \cdot (\varepsilon \cdot r^{-1}) = \varepsilon$
\item $\varphi_{71} :\equiv K \in \mathbb{N} \land \underline{K}^{-1} < \varepsilon \cdot r^{-1} \to n \in \mathbb{N} \to K \leq n \to 0 < (\underline{n} + 1)^{-1}$
\item $\varphi_{72} :\equiv K \in \mathbb{N} \land \underline{K}^{-1} < \varepsilon \cdot r^{-1} \to K \in \mathbb{N} \land \underline{K}^{-1} < \varepsilon \cdot r^{-1}$
\item $\varphi_{73} :\equiv K \in \mathbb{N} \land \underline{K}^{-1} < \varepsilon \cdot r^{-1} \to n \in \mathbb{N} \to K \leq n \to \underline{n} < \underline{n} + 1$
\item $\varphi_{74} :\equiv 0 < \varepsilon \to K \in \mathbb{N} \land \underline{K}^{-1} < \varepsilon \cdot r^{-1} \to n \in \mathbb{N} \to K \leq n \to 0 < \underline{n}$
\item $\varphi_{75} :\equiv \forall z\,(r \cdot (\varepsilon \cdot r^{-1}) = \varepsilon \cdot r^{-1} \cdot r \to \varepsilon \cdot r^{-1} \cdot r = z \to r \cdot (\varepsilon \cdot r^{-1}) = z)$
\item $\varphi_{76} :\equiv K \in \mathbb{N} \land \underline{K}^{-1} < \varepsilon \cdot r^{-1} \to n \in \mathbb{N} \to K \leq n \to 0 \leq r \cdot (\underline{n} + 1)^{-1}$
\item $\varphi_{77} :\equiv \forall y\,(\underline{n} \in \mathbb{R} \to y \in \mathbb{R} \to 0 < \underline{n} \to \underline{n} < y \to y^{-1} < \underline{n}^{-1})$
\item $\varphi_{78} :\equiv r \cdot (\varepsilon \cdot r^{-1}) = \varepsilon \cdot r^{-1} \cdot r \to \varepsilon \cdot r^{-1} \cdot r = \varepsilon \to r \cdot (\varepsilon \cdot r^{-1}) = \varepsilon$
\item $\varphi_{79} :\equiv K \in \mathbb{N} \land \underline{K}^{-1} < \varepsilon \cdot r^{-1} \to n \in \mathbb{N} \to K \leq n \to \underline{n}^{-1} \leq \underline{K}^{-1}$
\item $\varphi_{80} :\equiv 0 < r \to \varepsilon \in \mathbb{R} \to 0 < \varepsilon \to K \in \mathbb{N} \land \underline{K}^{-1} < \varepsilon \cdot r^{-1} \to K \in \mathbb{N}$
\item $\varphi_{81} :\equiv K \in \mathbb{N} \land \underline{K}^{-1} < \varepsilon \cdot r^{-1} \to n \in \mathbb{N} \to K \leq n \to (\underline{n} + 1)^{-1} \in \mathbb{R}$
\item $\varphi_{82} :\equiv \exists a\,(\mathrm{Seq}(a) \land \mathrm{Conv}(a,c) \land \forall n \in \mathbb{N}\,(0 < \lvert a_{n} - c \rvert \land \lvert a_{n} - c \rvert < r))$
\item $\varphi_{83} :\equiv \forall a\,(\neg (\mathrm{Seq}(a) \land \mathrm{Conv}(a,c) \land \forall n \in \mathbb{N}\,(0 < \lvert a_{n} - c \rvert \land \lvert a_{n} - c \rvert < r)))$
\item $\varphi_{84} :\equiv r \cdot (\underline{n} + 1)^{-1} \in \mathbb{R} \to c + r \cdot (\underline{n} + 1)^{-1} = r \cdot (\underline{n} + 1)^{-1} + c$
\item $\varphi_{85} :\equiv r \in \mathbb{R} \to 0 < r \to n \in \mathbb{N} \to \pi_{c,r}(n) - c = c + r \cdot (\underline{n} + 1)^{-1} - c$
\item $\varphi_{86} :\equiv r \in \mathbb{R} \to 0 < r \to n \in \mathbb{N} \to \pi_{c,r}(n) - c = r \cdot (\underline{n} + 1)^{-1} + c - c$
\item $\varphi_{87} :\equiv r \in \mathbb{R} \to 0 < r \to n \in \mathbb{N} \to \lvert \pi_{c,r}(n) - c \rvert = \lvert r \cdot (\underline{n} + 1)^{-1} \rvert$
\item $\varphi_{88} :\equiv 0 < \varepsilon \to K \in \mathbb{N} \land \underline{K}^{-1} < \varepsilon \cdot r^{-1} \to n \in \mathbb{N} \to K \leq n \to \underline{n} \in \mathbb{R}$
\item $\varphi_{89} :\equiv K \in \mathbb{N} \land \underline{K}^{-1} < \varepsilon \cdot r^{-1} \to n \in \mathbb{N} \to K \leq n \to 0 < r \cdot (\underline{n} + 1)^{-1}$
\item $\varphi_{90} :\equiv \forall x\,\forall y\,(x \in \mathbb{R} \to y \in \mathbb{R} \to x + y = y + x) \to \forall y\,(1 \in \mathbb{R} \to y \in \mathbb{R} \to 1 + y = y + 1)$
\item $\varphi_{91} :\equiv \forall x\,\forall y\,\forall z\,(x = y \to y = z \to x = z) \to \forall y\,\forall z\,(\pi_{c,r}(n) - c = y \to y = z \to \pi_{c,r}(n) - c = z)$
\item $\varphi_{92} :\equiv \forall x\,\forall y\,(x \in \mathbb{R} \to y \in \mathbb{R} \to x + y = y + x) \to \forall y\,(c \in \mathbb{R} \to y \in \mathbb{R} \to c + y = y + c)$
\item $\varphi_{93} :\equiv \forall x\,\forall y\,(x \in \mathbb{R} \to y \in \mathbb{R} \to x \cdot y = y \cdot x) \to \forall y\,(r \in \mathbb{R} \to y \in \mathbb{R} \to r \cdot y = y \cdot r)$
\item $\varphi_{94} :\equiv \varepsilon \cdot r^{-1} \in \mathbb{R} \to 0 < \varepsilon \cdot r^{-1} \to \exists K \in \mathbb{N}\,(\underline{K}^{-1} < \varepsilon \cdot r^{-1})$
\item $\varphi_{95} :\equiv (\underline{n} + 1)^{-1} \in \mathbb{R} \to 0 < r \to 0 < (\underline{n} + 1)^{-1} \to 0 < r \cdot (\underline{n} + 1)^{-1}$
\item $\varphi_{96} :\equiv 0 < r \to \forall n \in \mathbb{N}\,(0 < \lvert \pi_{c,r}(n) - c \rvert \land \lvert \pi_{c,r}(n) - c \rvert < r)$
\item $\varphi_{97} :\equiv K \in \mathbb{N} \land \underline{K}^{-1} < \varepsilon \cdot r^{-1} \to n \in \mathbb{N} \to K \leq n \to \underline{K}^{-1} < \varepsilon \cdot r^{-1}$
\item $\varphi_{98} :\equiv \varepsilon \in \mathbb{R} \to 0 < \varepsilon \to K \in \mathbb{N} \land \underline{K}^{-1} < \varepsilon \cdot r^{-1} \to n \in \mathbb{N} \to n \in \mathbb{N}$
\item $\varphi_{99} :\equiv K \in \mathbb{N} \land \underline{K}^{-1} < \varepsilon \cdot r^{-1} \to n \in \mathbb{N} \to K \leq n \to r \cdot (\underline{n} + 1)^{-1} < \varepsilon$
\item $\varphi_{100} :\equiv r \in \mathbb{R} \to 0 < r \to \varepsilon \in \mathbb{R} \to 0 < \varepsilon \to \exists K \in \mathbb{N}\,(\underline{K}^{-1} < \varepsilon \cdot r^{-1})$
\item $\varphi_{101} :\equiv (c \in \mathbb{R} \to r \in \mathbb{R} \to 0 < r \to 0 \neq r) \to (r = 0 \to 0 = r) \to c \in \mathbb{R} \to r \in \mathbb{R} \to 0 < r \to r \neq 0$
\item $\varphi_{102} :\equiv K \in \mathbb{N} \land \underline{K}^{-1} < \varepsilon \cdot r^{-1} \to n \in \mathbb{N} \to K \leq n \to r \cdot (\underline{n} + 1)^{-1} \in \mathbb{R}$
\item $\varphi_{103} :\equiv \forall x\,\forall y\,(x \in \mathbb{R} \to y \in \mathbb{R} \to x \cdot y \in \mathbb{R}) \to \forall y\,(r \in \mathbb{R} \to y \in \mathbb{R} \to r \cdot y \in \mathbb{R})$
\item $\varphi_{104} :\equiv K \in \mathbb{N} \land \underline{K}^{-1} < \varepsilon \cdot r^{-1} \to n \in \mathbb{N} \to K \leq n \to (\underline{n} + 1)^{-1} < \underline{n}^{-1}$
\item $\varphi_{105} :\equiv K \in \mathbb{N} \land \underline{K}^{-1} < \varepsilon \cdot r^{-1} \to n \in \mathbb{N} \to K \leq n \to (\underline{n} + 1)^{-1} < \underline{K}^{-1}$
\item $\varphi_{106} :\equiv (\underline{n} + 1 = 0 \to 0 = \underline{n} + 1) \to c \in \mathbb{R} \to r \in \mathbb{R} \to 0 < r \to n \in \mathbb{N} \to \underline{n} + 1 \neq 0$
\item $\varphi_{107} :\equiv K \in \mathbb{N} \land \underline{K}^{-1} < \varepsilon \cdot r^{-1} \to n \in \mathbb{N} \to K \leq n \to (\underline{n} + 1)^{-1} < \varepsilon \cdot r^{-1}$
\item $\varphi_{108} :\equiv K \in \mathbb{N} \land \underline{K}^{-1} < \varepsilon \cdot r^{-1} \to n \in \mathbb{N} \to K \leq n \to \lvert \pi_{c,r}(n) - c \rvert < \varepsilon$
\item $\varphi_{109} :\equiv \varepsilon \in \mathbb{R} \to 0 < \varepsilon \to \exists N \in \mathbb{N}\,\forall n \in \mathbb{N}\,(N \leq n \to \lvert \pi_{c,r}(n) - c \rvert < \varepsilon)$
\item $\varphi_{110} :\equiv (\underline{n} + 1)^{-1} < \underline{n}^{-1} \to \underline{n}^{-1} \leq \underline{K}^{-1} \to (\underline{n} + 1)^{-1} < \underline{K}^{-1}$
\item $\varphi_{111} :\equiv \forall x\,\forall y\,(x \in \mathbb{R} \to y \in \mathbb{R} \to x \cdot y \in \mathbb{R}) \to \forall y\,(\varepsilon \in \mathbb{R} \to y \in \mathbb{R} \to \varepsilon \cdot y \in \mathbb{R})$
\item $\varphi_{112} :\equiv \underline{n} + 1 \in \mathbb{R} \to 0 < \underline{n} \to \underline{n} < \underline{n} + 1 \to (\underline{n} + 1)^{-1} < \underline{n}^{-1}$
\item $\varphi_{113} :\equiv K \in \mathbb{N} \land \underline{K}^{-1} < \varepsilon \cdot r^{-1} \to n \in \mathbb{N} \to K \leq n \to r \cdot (\varepsilon \cdot r^{-1}) = \varepsilon \cdot r^{-1} \cdot r$
\item $\varphi_{114} :\equiv K \in \mathbb{N} \land \underline{K}^{-1} < \varepsilon \cdot r^{-1} \to \forall n \in \mathbb{N}\,(K \leq n \to \lvert \pi_{c,r}(n) - c \rvert < \varepsilon)$
\item $\varphi_{115} :\equiv \forall w\,(r \cdot (\varepsilon \cdot r^{-1}) = w \to r \cdot (\underline{n} + 1)^{-1} < r \cdot (\varepsilon \cdot r^{-1}) \to r \cdot (\underline{n} + 1)^{-1} < w)$
\item $\varphi_{116} :\equiv r \in \mathbb{R} \to 0 < r \to n \in \mathbb{N} \to 0 < \lvert \pi_{c,r}(n) - c \rvert \land \lvert \pi_{c,r}(n) - c \rvert < r$
\item $\varphi_{117} :\equiv \forall \varepsilon \in \mathbb{R}\,(0 < \varepsilon \to \exists N \in \mathbb{N}\,\forall n \in \mathbb{N}\,(N \leq n \to \lvert \pi_{c,r}(n) - c \rvert < \varepsilon))$
\item $\varphi_{118} :\equiv \forall y\,(1 \in \mathbb{R} \to y \in \mathbb{R} \to 1 + y = y + 1) \to 1 \in \mathbb{R} \to \underline{n} \in \mathbb{R} \to 1 + \underline{n} = \underline{n} + 1$
\item $\varphi_{119} :\equiv r \cdot (\underline{n} + 1)^{-1} \in \mathbb{R} \to c \in \mathbb{R} \to r \cdot (\underline{n} + 1)^{-1} + c - c = r \cdot (\underline{n} + 1)^{-1}$
\item $\varphi_{120} :\equiv \forall w\,(\pi_{c,r}(n) = w \to \pi_{c,r}(n) - c = \pi_{c,r}(n) - c \to \pi_{c,r}(n) - c = w - c)$
\item $\varphi_{121} :\equiv (\underline{n} + 1)^{-1} < \underline{K}^{-1} \to \underline{K}^{-1} < \varepsilon \cdot r^{-1} \to (\underline{n} + 1)^{-1} < \varepsilon \cdot r^{-1}$
\item $\varphi_{122} :\equiv c \in \mathbb{R} \land r \in \mathbb{R} \land 0 < r \land \varepsilon \in \mathbb{R} \land 0 < \varepsilon \land (K \in \mathbb{N} \land \underline{K}^{-1} < \varepsilon \cdot r^{-1})$
\item $\varphi_{123} :\equiv \lvert \pi_{c,r}(n) - c \rvert = r \cdot (\underline{n} + 1)^{-1} \to r \cdot (\underline{n} + 1)^{-1} = \lvert \pi_{c,r}(n) - c \rvert$
\item $\varphi_{124} :\equiv \lvert r \cdot (\underline{n} + 1)^{-1} \rvert = r \cdot (\underline{n} + 1)^{-1} \to \lvert \pi_{c,r}(n) - c \rvert = r \cdot (\underline{n} + 1)^{-1}$
\item $\varphi_{125} :\equiv r \cdot (\underline{n} + 1)^{-1} + c - c = r \cdot (\underline{n} + 1)^{-1} \to \pi_{c,r}(n) - c = r \cdot (\underline{n} + 1)^{-1}$
\item $\varphi_{126} :\equiv \pi_{c,r}(n) - c = \pi_{c,r}(n) - c \to \pi_{c,r}(n) - c = c + r \cdot (\underline{n} + 1)^{-1} - c$
\item $\varphi_{127} :\equiv \forall x\,\forall y\,(x \in \mathbb{R} \to y \in \mathbb{R} \to 0 < y \to x < x + y) \to \forall y\,(1 \in \mathbb{R} \to y \in \mathbb{R} \to 0 < y \to 1 < 1 + y)$
\item $\varphi_{128} :\equiv \forall x \in \mathbb{R}\,(0 < x \to 0 < x^{-1}) \to \underline{n} + 1 \in \mathbb{R} \to 0 < \underline{n} + 1 \to 0 < (\underline{n} + 1)^{-1}$
\item $\varphi_{129} :\equiv \forall y\,(r \cdot (\underline{n} + 1)^{-1} \in \mathbb{R} \to y \in \mathbb{R} \to r \cdot (\underline{n} + 1)^{-1} + y - y = r \cdot (\underline{n} + 1)^{-1})$
\item $\varphi_{130} :\equiv \forall x\,\forall y\,(x \in \mathbb{R} \to y \in \mathbb{R} \to x + y \in \mathbb{R}) \to \forall y\,(\underline{n} \in \mathbb{R} \to y \in \mathbb{R} \to \underline{n} + y \in \mathbb{R})$
\item $\varphi_{131} :\equiv \forall y\,(\varepsilon \in \mathbb{R} \to y \in \mathbb{R} \to \varepsilon \cdot y \in \mathbb{R}) \to \varepsilon \in \mathbb{R} \to r^{-1} \in \mathbb{R} \to \varepsilon \cdot r^{-1} \in \mathbb{R}$
\item $\varphi_{132} :\equiv K \in \mathbb{N} \land \underline{K}^{-1} < \varepsilon \cdot r^{-1} \to n \in \mathbb{N} \to K \leq n \to r \cdot (\underline{n} + 1)^{-1} < r \cdot (\varepsilon \cdot r^{-1})$
\item $\varphi_{133} :\equiv K \in \mathbb{N} \land \underline{K}^{-1} < \varepsilon \cdot r^{-1} \to n \in \mathbb{N} \to K \leq n \to \pi_{c,r}(n) = c + r \cdot (\underline{n} + 1)^{-1}$
\item $\varphi_{134} :\equiv K \in \mathbb{N} \land \underline{K}^{-1} < \varepsilon \cdot r^{-1} \to n \in \mathbb{N} \to K \leq n \to \pi_{c,r}(n) - c = r \cdot (\underline{n} + 1)^{-1}$
\item $\varphi_{135} :\equiv \forall x \in \mathbb{R}\,(x \neq 0 \to x^{-1} \in \mathbb{R}) \to \underline{n} + 1 \in \mathbb{R} \to \underline{n} + 1 \neq 0 \to (\underline{n} + 1)^{-1} \in \mathbb{R}$
\item $\varphi_{136} :\equiv K \in \mathbb{N} \land \underline{K}^{-1} < \varepsilon \cdot r^{-1} \to K \in \mathbb{N} \land \forall n \in \mathbb{N}\,(K \leq n \to \lvert \pi_{c,r}(n) - c \rvert < \varepsilon)$
\item $\varphi_{137} :\equiv \forall y\,(\underline{n} \in \mathbb{R} \to y \in \mathbb{R} \to \underline{n} + y \in \mathbb{R}) \to \underline{n} \in \mathbb{R} \to 1 \in \mathbb{R} \to \underline{n} + 1 \in \mathbb{R}$
\item $\varphi_{138} :\equiv K \in \mathbb{N} \land \underline{K}^{-1} < \varepsilon \cdot r^{-1} \to n \in \mathbb{N} \to K \leq n \to \lvert \pi_{c,r}(n) - c \rvert = r \cdot (\underline{n} + 1)^{-1}$
\item $\varphi_{139} :\equiv K \in \mathbb{N} \land \underline{K}^{-1} < \varepsilon \cdot r^{-1} \to n \in \mathbb{N} \to K \leq n \to r \cdot (\underline{n} + 1)^{-1} = \lvert \pi_{c,r}(n) - c \rvert$
\item $\varphi_{140} :\equiv \lvert \pi_{c,r}(n) - c \rvert = \lvert \pi_{c,r}(n) - c \rvert \to \lvert \pi_{c,r}(n) - c \rvert = \lvert r \cdot (\underline{n} + 1)^{-1} \rvert$
\item $\varphi_{141} :\equiv \forall w\,(1 + \underline{n} = w \to 1 < 1 + \underline{n} \to 1 < w) \to 1 + \underline{n} = \underline{n} + 1 \to 1 < 1 + \underline{n} \to 1 < \underline{n} + 1$
\item $\varphi_{142} :\equiv c \in \mathbb{R} \land r \in \mathbb{R} \land 0 < r \to n \in \mathbb{N} \to 0 < \lvert \pi_{c,r}(n) - c \rvert \land \lvert \pi_{c,r}(n) - c \rvert < r$
\item $\varphi_{143} :\equiv K \in \mathbb{N} \land \underline{K}^{-1} < \varepsilon \cdot r^{-1} \to \exists N \in \mathbb{N}\,\forall n \in \mathbb{N}\,(N \leq n \to \lvert \pi_{c,r}(n) - c \rvert < \varepsilon)$
\item $\varphi_{144} :\equiv r \cdot (\underline{n} + 1)^{-1} = \lvert \pi_{c,r}(n) - c \rvert \to r \cdot (\underline{n} + 1)^{-1} < r \to \lvert \pi_{c,r}(n) - c \rvert < r$
\item $\varphi_{145} :\equiv (c \in \mathbb{R} \to r \in \mathbb{R} \to 0 < r \to n \in \mathbb{N} \to r \in \mathbb{R}) \to c \in \mathbb{R} \to r \in \mathbb{R} \to 0 < r \to n \in \mathbb{N} \to r \cdot 1 = r$
\item $\varphi_{146} :\equiv r \cdot (\underline{n} + 1)^{-1} = \lvert \pi_{c,r}(n) - c \rvert \to 0 < r \cdot (\underline{n} + 1)^{-1} \to 0 < \lvert \pi_{c,r}(n) - c \rvert$
\item $\varphi_{147} :\equiv \forall y\,(1 \in \mathbb{R} \to y \in \mathbb{R} \to 0 < y \to 1 < 1 + y) \to 1 \in \mathbb{R} \to \underline{n} \in \mathbb{R} \to 0 < \underline{n} \to 1 < 1 + \underline{n}$
\item $\varphi_{148} :\equiv \forall n\,(K \in \mathbb{N} \to n \in \mathbb{N} \to K \leq n \to \underline{n}^{-1} \leq \underline{K}^{-1}) \to K \in \mathbb{N} \to n \in \mathbb{N} \to K \leq n \to \underline{n}^{-1} \leq \underline{K}^{-1}$
\item $\varphi_{149} :\equiv (c \in \mathbb{R} \to r \in \mathbb{R} \to 0 < r \to 0 < r) \to c \in \mathbb{R} \to r \in \mathbb{R} \to 0 < r \to \varepsilon \in \mathbb{R} \to 0 < \varepsilon \to 0 < r^{-1}$
\item $\varphi_{150} :\equiv K \in \mathbb{N} \land \underline{K}^{-1} < \varepsilon \cdot r^{-1} \to n \in \mathbb{N} \to K \leq n \to \pi_{c,r}(n) - c = c + r \cdot (\underline{n} + 1)^{-1} - c$
\item $\varphi_{151} :\equiv K \in \mathbb{N} \land \underline{K}^{-1} < \varepsilon \cdot r^{-1} \to n \in \mathbb{N} \to K \leq n \to \pi_{c,r}(n) - c = r \cdot (\underline{n} + 1)^{-1} + c - c$
\item $\varphi_{152} :\equiv K \in \mathbb{N} \land \underline{K}^{-1} < \varepsilon \cdot r^{-1} \to n \in \mathbb{N} \to K \leq n \to \lvert \pi_{c,r}(n) - c \rvert = \lvert r \cdot (\underline{n} + 1)^{-1} \rvert$
\item $\varphi_{153} :\equiv (c \in \mathbb{R} \to r \in \mathbb{R} \to 0 < r \to n \in \mathbb{N} \to n \in \mathbb{N}) \to c \in \mathbb{R} \to r \in \mathbb{R} \to 0 < r \to n \in \mathbb{N} \to 0 < \underline{n}$
\item $\varphi_{154} :\equiv (c \in \mathbb{R} \to r \in \mathbb{R} \to 0 < r \to r \neq 0) \to c \in \mathbb{R} \to r \in \mathbb{R} \to 0 < r \to \varepsilon \in \mathbb{R} \to 0 < \varepsilon \to r^{-1} \in \mathbb{R}$
\item $\varphi_{155} :\equiv r \cdot (\underline{n} + 1)^{-1} = \lvert \pi_{c,r}(n) - c \rvert \to r \cdot (\underline{n} + 1)^{-1} < \varepsilon \to \lvert \pi_{c,r}(n) - c \rvert < \varepsilon$
\item $\varphi_{156} :\equiv (c \in \mathbb{R} \to r \in \mathbb{R} \to 0 < r \to n \in \mathbb{N} \to 0 < 1) \to c \in \mathbb{R} \to r \in \mathbb{R} \to 0 < r \to n \in \mathbb{N} \to 0 < \underline{n} + 1$
\item $\varphi_{157} :\equiv \forall y\,(r \in \mathbb{R} \to y \in \mathbb{R} \to r \cdot y \in \mathbb{R}) \to r \in \mathbb{R} \to (\underline{n} + 1)^{-1} \in \mathbb{R} \to r \cdot (\underline{n} + 1)^{-1} \in \mathbb{R}$
\item $\varphi_{158} :\equiv \forall n \in \mathbb{N}\,(\pi_{c,r}(n) = c + r \cdot (\underline{n} + 1)^{-1}) \to n \in \mathbb{N} \to \pi_{c,r}(n) = c + r \cdot (\underline{n} + 1)^{-1}$
\item $\varphi_{159} :\equiv \forall x\,\forall y\,(x \in \mathbb{R} \to y \in \mathbb{R} \to y \neq 0 \to x \cdot y^{-1} \cdot y = x) \to \forall y\,(\varepsilon \in \mathbb{R} \to y \in \mathbb{R} \to y \neq 0 \to \varepsilon \cdot y^{-1} \cdot y = \varepsilon)$
\item $\varphi_{160} :\equiv \forall y\,(\varepsilon \in \mathbb{R} \to y \in \mathbb{R} \to y \neq 0 \to \varepsilon \cdot y^{-1} \cdot y = \varepsilon) \to \varepsilon \in \mathbb{R} \to r \in \mathbb{R} \to r \neq 0 \to \varepsilon \cdot r^{-1} \cdot r = \varepsilon$
\item $\varphi_{161} :\equiv \forall n\,(c \in \mathbb{R} \land r \in \mathbb{R} \land 0 < r) \to \forall n \in \mathbb{N}\,(0 < \lvert \pi_{c,r}(n) - c \rvert \land \lvert \pi_{c,r}(n) - c \rvert < r)$
\item $\varphi_{162} :\equiv \mathrm{Seq}(\pi_{c,r}) \land \mathrm{Conv}(\pi_{c,r},c) \land \forall n \in \mathbb{N}\,(0 < \lvert \pi_{c,r}(n) - c \rvert \land \lvert \pi_{c,r}(n) - c \rvert < r)$
\item $\varphi_{163} :\equiv (c \in \mathbb{R} \to r \in \mathbb{R} \to 0 < r \to n \in \mathbb{N} \to n \in \mathbb{N}) \to c \in \mathbb{R} \to r \in \mathbb{R} \to 0 < r \to n \in \mathbb{N} \to \underline{n} \in \mathbb{R}$
\item $\varphi_{164} :\equiv \forall z\,(\lvert \pi_{c,r}(n) - c \rvert = \lvert r \cdot (\underline{n} + 1)^{-1} \rvert \to \lvert r \cdot (\underline{n} + 1)^{-1} \rvert = z \to \lvert \pi_{c,r}(n) - c \rvert = z)$
\item $\varphi_{165} :\equiv \forall x\,\forall y\,(x \in \mathbb{R} \to y \in \mathbb{R} \to 0 < x \to 0 < y \to 0 < x \cdot y) \to \forall y\,(r \in \mathbb{R} \to y \in \mathbb{R} \to 0 < r \to 0 < y \to 0 < r \cdot y)$
\item $\varphi_{166} :\equiv \forall w\,(\pi_{c,r}(n) - c = w \to \lvert \pi_{c,r}(n) - c \rvert = \lvert \pi_{c,r}(n) - c \rvert \to \lvert \pi_{c,r}(n) - c \rvert = \lvert w \rvert)$
\item $\varphi_{167} :\equiv \forall z\,(\pi_{c,r}(n) - c = r \cdot (\underline{n} + 1)^{-1} + c - c \to r \cdot (\underline{n} + 1)^{-1} + c - c = z \to \pi_{c,r}(n) - c = z)$
\item $\varphi_{168} :\equiv \forall z\,(\pi_{c,r}(n) - c = c + r \cdot (\underline{n} + 1)^{-1} - c \to c + r \cdot (\underline{n} + 1)^{-1} - c = z \to \pi_{c,r}(n) - c = z)$
\item $\varphi_{169} :\equiv \forall K\,\forall n\,(K \in \mathbb{N} \to n \in \mathbb{N} \to K \leq n \to \underline{n}^{-1} \leq \underline{K}^{-1}) \to \forall n\,(K \in \mathbb{N} \to n \in \mathbb{N} \to K \leq n \to \underline{n}^{-1} \leq \underline{K}^{-1})$
\item $\varphi_{170} :\equiv c + r \cdot (\underline{n} + 1)^{-1} - c = r \cdot (\underline{n} + 1)^{-1} + c - c \to \pi_{c,r}(n) - c = r \cdot (\underline{n} + 1)^{-1} + c - c$
\item $\varphi_{171} :\equiv \forall z\,((\underline{n} + 1)^{-1} \in \mathbb{R} \to 1 \in \mathbb{R} \to z \in \mathbb{R} \to 0 < z \to (\underline{n} + 1)^{-1} < 1 \to z \cdot (\underline{n} + 1)^{-1} < z \cdot 1)$
\item $\varphi_{172} :\equiv (c \in \mathbb{R} \to r \in \mathbb{R} \to 0 < r \to n \in \mathbb{N} \to 0 < \underline{n}) \to c \in \mathbb{R} \to r \in \mathbb{R} \to 0 < r \to n \in \mathbb{N} \to 1 < 1 + \underline{n}$
\item $\varphi_{173} :\equiv \forall y\,\forall z\,((\underline{n} + 1)^{-1} \in \mathbb{R} \to y \in \mathbb{R} \to z \in \mathbb{R} \to 0 < z \to (\underline{n} + 1)^{-1} < y \to z \cdot (\underline{n} + 1)^{-1} < z \cdot y)$
\item $\varphi_{174} :\equiv \forall x\,\forall y\,(x \in \mathbb{R} \to y \in \mathbb{R} \to 0 < x \to 0 < y \to 0 < x \cdot y) \to \forall y\,(\varepsilon \in \mathbb{R} \to y \in \mathbb{R} \to 0 < \varepsilon \to 0 < y \to 0 < \varepsilon \cdot y)$
\item $\varphi_{175} :\equiv (c \in \mathbb{R} \to r \in \mathbb{R} \to 0 < r \to n \in \mathbb{N} \to 0 < \underline{n} + 1) \to c \in \mathbb{R} \to r \in \mathbb{R} \to 0 < r \to n \in \mathbb{N} \to 0 \neq \underline{n} + 1$
\item $\varphi_{176} :\equiv \varphi_{8} \to 1^{-1} = 1 \to (\underline{n} + 1)^{-1} < 1^{-1} \to (\underline{n} + 1)^{-1} < 1$
\item $\varphi_{177} :\equiv (1 + \underline{n} = \underline{n} + 1 \to 1 < 1 + \underline{n} \to 1 < \underline{n} + 1) \to c \in \mathbb{R} \to r \in \mathbb{R} \to 0 < r \to n \in \mathbb{N} \to 1 < \underline{n} + 1$
\item $\varphi_{178} :\equiv \varphi_{9} \to \forall y\,(1 \in \mathbb{R} \to y \in \mathbb{R} \to 0 < 1 \to 1 < y \to y^{-1} < 1^{-1})$
\item $\varphi_{179} :\equiv c \in \mathbb{R} \land r \in \mathbb{R} \land 0 < r \land \varepsilon \in \mathbb{R} \land 0 < \varepsilon \land \neg \exists N \in \mathbb{N}\,\forall n \in \mathbb{N}\,(N \leq n \to \lvert \pi_{c,r}(n) - c \rvert < \varepsilon)$
\item $\varphi_{180} :\equiv c + r \cdot (\underline{n} + 1)^{-1} - c = c + r \cdot (\underline{n} + 1)^{-1} - c \to c + r \cdot (\underline{n} + 1)^{-1} - c = r \cdot (\underline{n} + 1)^{-1} + c - c$
\item $\varphi_{181} :\equiv (1^{-1} = 1 \to (\underline{n} + 1)^{-1} < 1^{-1} \to (\underline{n} + 1)^{-1} < 1) \to c \in \mathbb{R} \to r \in \mathbb{R} \to 0 < r \to n \in \mathbb{N} \to (\underline{n} + 1)^{-1} < 1$
\item $\varphi_{182} :\equiv (c \in \mathbb{R} \to r \in \mathbb{R} \to 0 < r \to n \in \mathbb{N} \to 0 < \underline{n} + 1) \to c \in \mathbb{R} \to r \in \mathbb{R} \to 0 < r \to n \in \mathbb{N} \to 0 < (\underline{n} + 1)^{-1}$
\item $\varphi_{183} :\equiv \forall a\,\forall \alpha\,(\mathrm{Conv}(a,\alpha) \leftrightarrow \mathrm{Seq}(a) \land \alpha \in \mathbb{R} \land \forall \varepsilon \in \mathbb{R}\,(0 < \varepsilon \to \exists N \in \mathbb{N}\,\forall n \in \mathbb{N}\,(N \leq n \to \lvert a_{n} - \alpha \rvert < \varepsilon)))$
\item $\varphi_{184} :\equiv (c \in \mathbb{R} \to r \in \mathbb{R} \to 0 < r \to n \in \mathbb{N} \to \underline{n} \in \mathbb{R}) \to c \in \mathbb{R} \to r \in \mathbb{R} \to 0 < r \to n \in \mathbb{N} \to 1 + \underline{n} = \underline{n} + 1$
\item $\varphi_{185} :\equiv (c \in \mathbb{R} \to r \in \mathbb{R} \to 0 < r \to n \in \mathbb{N} \to \underline{n} + 1 \neq 0) \to c \in \mathbb{R} \to r \in \mathbb{R} \to 0 < r \to n \in \mathbb{N} \to (\underline{n} + 1)^{-1} \in \mathbb{R}$
\item $\varphi_{186} :\equiv \forall x\,\forall y\,(x \in \mathbb{R} \to y \in \mathbb{R} \to 0 < x \to 0 < y \to 0 < x + y) \to \forall y\,(\underline{n} \in \mathbb{R} \to y \in \mathbb{R} \to 0 < \underline{n} \to 0 < y \to 0 < \underline{n} + y)$
\item $\varphi_{187} :\equiv \forall w\,(c + r \cdot (\underline{n} + 1)^{-1} = w \to c + r \cdot (\underline{n} + 1)^{-1} - c = c + r \cdot (\underline{n} + 1)^{-1} - c \to c + r \cdot (\underline{n} + 1)^{-1} - c = w - c)$
\item $\varphi_{188} :\equiv \varphi_{17} \to \neg (K \in \mathbb{N} \land \forall n \in \mathbb{N}\,(K \leq n \to \lvert \pi_{c,r}(n) - c \rvert < \varepsilon))$
\item $\varphi_{189} :\equiv \forall y\,(\varepsilon \in \mathbb{R} \to y \in \mathbb{R} \to 0 < \varepsilon \to 0 < y \to 0 < \varepsilon \cdot y) \to \varepsilon \in \mathbb{R} \to r^{-1} \in \mathbb{R} \to 0 < \varepsilon \to 0 < r^{-1} \to 0 < \varepsilon \cdot r^{-1}$
\item $\varphi_{190} :\equiv (c \in \mathbb{R} \to r \in \mathbb{R} \to 0 < r \to n \in \mathbb{N} \to 0 < (\underline{n} + 1)^{-1}) \to \varphi_{28}$
\item $\varphi_{191} :\equiv (c \in \mathbb{R} \to r \in \mathbb{R} \to 0 < r \to n \in \mathbb{N} \to (\underline{n} + 1)^{-1} < 1) \to \varphi_{46}$
\item $\varphi_{192} :\equiv \forall y\,(\underline{n} \in \mathbb{R} \to y \in \mathbb{R} \to 0 < \underline{n} \to 0 < y \to 0 < \underline{n} + y) \to \underline{n} \in \mathbb{R} \to 1 \in \mathbb{R} \to 0 < \underline{n} \to 0 < 1 \to 0 < \underline{n} + 1$
\item $\varphi_{193} :\equiv \varepsilon \in \mathbb{R} \to 0 < \varepsilon \to K \in \mathbb{N} \land \underline{K}^{-1} < \varepsilon \cdot r^{-1} \to n \in \mathbb{N} \to \varphi_{18}$
\item $\varphi_{194} :\equiv (c \in \mathbb{R} \to r \in \mathbb{R} \to 0 < r \to n \in \mathbb{N} \to (\underline{n} + 1)^{-1} \in \mathbb{R}) \to \varphi_{43}$
\item $\varphi_{195} :\equiv (c \in \mathbb{R} \to r \in \mathbb{R} \to 0 < r \to \varepsilon \in \mathbb{R} \to 0 < \varepsilon \to 0 < r^{-1}) \to c \in \mathbb{R} \to r \in \mathbb{R} \to 0 < r \to \varepsilon \in \mathbb{R} \to 0 < \varepsilon \to 0 < \varepsilon \cdot r^{-1}$
\item $\varphi_{196} :\equiv (\underline{n} + 1)^{-1} \in \mathbb{R} \to \varepsilon \cdot r^{-1} \in \mathbb{R} \to r \in \mathbb{R} \to 0 < r \to \varphi_{57}$
\item $\varphi_{197} :\equiv \forall y\,(1 \in \mathbb{R} \to y \in \mathbb{R} \to 0 < 1 \to 1 < y \to y^{-1} < 1^{-1}) \to 1 \in \mathbb{R} \to \varphi_{53}$
\item $\varphi_{198} :\equiv \forall z\,((\underline{n} + 1)^{-1} \in \mathbb{R} \to \varepsilon \cdot r^{-1} \in \mathbb{R} \to z \in \mathbb{R} \to 0 < z \to (\underline{n} + 1)^{-1} < \varepsilon \cdot r^{-1} \to z \cdot (\underline{n} + 1)^{-1} < z \cdot (\varepsilon \cdot r^{-1}))$
\item $\varphi_{199} :\equiv \forall \alpha\,(\mathrm{Conv}(\pi_{c,r},\alpha) \leftrightarrow \mathrm{Seq}(\pi_{c,r}) \land \alpha \in \mathbb{R} \land \forall \varepsilon \in \mathbb{R}\,(0 < \varepsilon \to \exists N \in \mathbb{N}\,\forall n \in \mathbb{N}\,(N \leq n \to \lvert \pi_{c,r}(n) - \alpha \rvert < \varepsilon)))$
\item $\varphi_{200} :\equiv r \in \mathbb{R} \to 0 < r \to \varepsilon \in \mathbb{R} \to 0 < \varepsilon \to K \in \mathbb{N} \land \underline{K}^{-1} < \varepsilon \cdot r^{-1} \to \varphi_{12}$
\item $\varphi_{201} :\equiv (c \in \mathbb{R} \to r \in \mathbb{R} \to 0 < r \to r \in \mathbb{R}) \to c \in \mathbb{R} \to r \in \mathbb{R} \to 0 < r \to \varphi_{45}$
\item $\varphi_{202} :\equiv (c \in \mathbb{R} \to r \in \mathbb{R} \to 0 < r \to n \in \mathbb{N} \to c \in \mathbb{R}) \to c \in \mathbb{R} \to r \in \mathbb{R} \to \varphi_{34}$
\item $\varphi_{203} :\equiv (c \in \mathbb{R} \to r \in \mathbb{R} \to 0 < r \to \varepsilon \in \mathbb{R} \to 0 < \varepsilon \to r^{-1} \in \mathbb{R}) \to \varphi_{62}$
\item $\varphi_{204} :\equiv \forall y\,(r \in \mathbb{R} \to y \in \mathbb{R} \to 0 < r \to 0 < y \to 0 < r \cdot y) \to r \in \mathbb{R} \to \varphi_{95}$
\item $\varphi_{205} :\equiv r \in \mathbb{R} \to 0 < r \to \varepsilon \in \mathbb{R} \to 0 < \varepsilon \to K \in \mathbb{N} \land \underline{K}^{-1} < \varepsilon \cdot r^{-1} \to \varphi_{24}$
\item $\varphi_{206} :\equiv r \in \mathbb{R} \to 0 < r \to \varepsilon \in \mathbb{R} \to 0 < \varepsilon \to K \in \mathbb{N} \land \underline{K}^{-1} < \varepsilon \cdot r^{-1} \to \varphi_{25}$
\item $\varphi_{207} :\equiv (c \in \mathbb{R} \to r \in \mathbb{R} \to 0 < r \to \varphi_{45}) \to c \in \mathbb{R} \to r \in \mathbb{R} \to 0 < r \to \mathrm{Seq}(\pi_{c,r})$
\item $\varphi_{208} :\equiv (\underline{n} + 1 = 0 \to 0 = \underline{n} + 1) \to c \in \mathbb{R} \to r \in \mathbb{R} \to 0 < r \to \varepsilon \in \mathbb{R} \to 0 < \varepsilon \to \varphi_{42}$
\item $\varphi_{209} :\equiv (c \in \mathbb{R} \to r \in \mathbb{R} \to 0 < r \to r \neq 0) \to c \in \mathbb{R} \to r \in \mathbb{R} \to 0 < r \to \varepsilon \in \mathbb{R} \to 0 < \varepsilon \to \varphi_{40}$
\item $\varphi_{210} :\equiv (c \in \mathbb{R} \to r \in \mathbb{R} \to 0 < r \to r \neq 0) \to c \in \mathbb{R} \to r \in \mathbb{R} \to 0 < r \to \varepsilon \in \mathbb{R} \to 0 < \varepsilon \to \varphi_{69}$
\item $\varphi_{211} :\equiv (c \in \mathbb{R} \to r \in \mathbb{R} \to 0 < r \to \varepsilon \in \mathbb{R} \to 0 < \varepsilon \to 0 < \varepsilon \cdot r^{-1}) \to c \in \mathbb{R} \to \varphi_{100}$
\item $\varphi_{212} :\equiv (c \in \mathbb{R} \to r \in \mathbb{R} \to 0 < r \to \varepsilon \in \mathbb{R} \to 0 < \varepsilon \to r \in \mathbb{R}) \to \varphi_{149}$
\item $\varphi_{213} :\equiv (c \in \mathbb{R} \to r \in \mathbb{R} \to 0 < r \to \varepsilon \in \mathbb{R} \to 0 < \varepsilon \to r \in \mathbb{R}) \to \varphi_{154}$
\item $\varphi_{214} :\equiv (c \in \mathbb{R} \to r \in \mathbb{R} \to 0 < r \to \varphi_{117}) \to c \in \mathbb{R} \to r \in \mathbb{R} \to 0 < r \to \mathrm{Conv}(\pi_{c,r},c)$
\item $\varphi_{215} :\equiv 1 \in \mathbb{R} \to (c \in \mathbb{R} \to r \in \mathbb{R} \to 0 < r \to n \in \mathbb{N} \to 0 < \underline{n}) \to \varphi_{156}$
\item $\varphi_{216} :\equiv (c \in \mathbb{R} \to r \in \mathbb{R} \to 0 < r \to n \in \mathbb{N} \to \lvert \pi_{c,r}(n) - c \rvert < r) \to c \in \mathbb{R} \to \varphi_{116}$
\item $\varphi_{217} :\equiv (c \in \mathbb{R} \to r \in \mathbb{R} \to 0 < r \to n \in \mathbb{N} \to 0 < 1) \to (c \in \mathbb{R} \to r \in \mathbb{R} \to 0 < r \to n \in \mathbb{N} \to 1 < \underline{n} + 1) \to \varphi_{26}$
\item $\varphi_{218} :\equiv \varphi_{144} \to c \in \mathbb{R} \to r \in \mathbb{R} \to 0 < r \to n \in \mathbb{N} \to \lvert \pi_{c,r}(n) - c \rvert < r$
\item $\varphi_{219} :\equiv \varphi_{146} \to c \in \mathbb{R} \to r \in \mathbb{R} \to 0 < r \to n \in \mathbb{N} \to 0 < \lvert \pi_{c,r}(n) - c \rvert$
\item $\varphi_{220} :\equiv (\underline{n} \in \mathbb{R} \to 1 \in \mathbb{R} \to \underline{n} + 1 \in \mathbb{R}) \to (c \in \mathbb{R} \to r \in \mathbb{R} \to 0 < r \to n \in \mathbb{N} \to \underline{n} \in \mathbb{R}) \to \varphi_{54}$
\item $\varphi_{221} :\equiv 1 \in \mathbb{R} \to (c \in \mathbb{R} \to r \in \mathbb{R} \to 0 < r \to n \in \mathbb{N} \to \underline{n} \in \mathbb{R}) \to \varphi_{172}$
\item $\varphi_{222} :\equiv \forall c\,\forall r\,(c \in \mathbb{R} \to r \in \mathbb{R} \to \varphi_{45}) \to \forall r\,(c \in \mathbb{R} \to r \in \mathbb{R} \to \varphi_{45})$
\item $\varphi_{223} :\equiv (c \in \mathbb{R} \to r \in \mathbb{R} \to 0 < r \to n \in \mathbb{N} \to c \in \mathbb{R}) \to \varphi_{43} \to c \in \mathbb{R} \to r \in \mathbb{R} \to \varphi_{35}$
\item $\varphi_{224} :\equiv (c \in \mathbb{R} \to r \in \mathbb{R} \to 0 < r \to n \in \mathbb{N} \to r \in \mathbb{R}) \to (c \in \mathbb{R} \to r \in \mathbb{R} \to 0 < r \to 0 < r) \to \varphi_{191}$
\item $\varphi_{225} :\equiv (\pi_{c,r}(n) - c = r \cdot (\underline{n} + 1)^{-1} + c - c \to \varphi_{125}) \to c \in \mathbb{R} \to \varphi_{49}$
\item $\varphi_{226} :\equiv (r \cdot (\varepsilon \cdot r^{-1}) = \varepsilon \to \varphi_{36}) \to c \in \mathbb{R} \to r \in \mathbb{R} \to 0 < r \to \varepsilon \in \mathbb{R} \to 0 < \varepsilon \to \varphi_{99}$
\item $\varphi_{227} :\equiv (\lvert \pi_{c,r}(n) - c \rvert = \lvert r \cdot (\underline{n} + 1)^{-1} \rvert \to \varphi_{124}) \to c \in \mathbb{R} \to \varphi_{66}$
\item $\varphi_{228} :\equiv (c \in \mathbb{R} \to r \in \mathbb{R} \to 0 < r \to \varepsilon \in \mathbb{R} \to 0 < \varepsilon \to \varphi_{72}) \to c \in \mathbb{R} \to r \in \mathbb{R} \to \varphi_{80}$
\item $\varphi_{229} :\equiv (c \in \mathbb{R} \to r \in \mathbb{R} \to 0 < r \to \varphi_{98}) \to c \in \mathbb{R} \to r \in \mathbb{R} \to 0 < r \to \varepsilon \in \mathbb{R} \to \varphi_{74}$
\item $\varphi_{230} :\equiv (c \in \mathbb{R} \to r \in \mathbb{R} \to 0 < r \to \varphi_{109}) \to c \in \mathbb{R} \land r \in \mathbb{R} \land 0 < r \to \varphi_{109}$
\item $\varphi_{231} :\equiv (c \in \mathbb{R} \to r \in \mathbb{R} \to 0 < r \to n \in \mathbb{N} \to (\underline{n} + 1)^{-1} \in \mathbb{R}) \to (c \in \mathbb{R} \to r \in \mathbb{R} \to 0 < r \to 0 < r) \to \varphi_{190}$
\item $\varphi_{232} :\equiv (c \in \mathbb{R} \to r \in \mathbb{R} \to 0 < r \to \varphi_{98}) \to c \in \mathbb{R} \to r \in \mathbb{R} \to 0 < r \to \varepsilon \in \mathbb{R} \to \varphi_{88}$
\item $\varphi_{233} :\equiv (c \in \mathbb{R} \to r \in \mathbb{R} \to 0 < r \to \varepsilon \in \mathbb{R} \to \varphi_{50}) \to c \in \mathbb{R} \to r \in \mathbb{R} \to 0 < r \to \varepsilon \in \mathbb{R} \to 0 < \varepsilon \to \varphi_{51}$
\item $\varphi_{234} :\equiv (c \in \mathbb{R} \to r \in \mathbb{R} \to 0 < r \to n \in \mathbb{N} \to 1 + \underline{n} = \underline{n} + 1) \to (c \in \mathbb{R} \to r \in \mathbb{R} \to 0 < r \to n \in \mathbb{N} \to 1 < 1 + \underline{n}) \to \varphi_{177}$
\item $\varphi_{235} :\equiv (\underline{n} + 1 \in \mathbb{R} \to 0 < \underline{n} + 1 \to 0 < (\underline{n} + 1)^{-1}) \to (c \in \mathbb{R} \to r \in \mathbb{R} \to 0 < r \to n \in \mathbb{N} \to \underline{n} + 1 \in \mathbb{R}) \to \varphi_{182}$
\item $\varphi_{236} :\equiv (c \in \mathbb{R} \to r \in \mathbb{R} \to 0 < r \to \varepsilon \in \mathbb{R} \to \varphi_{37}) \to c \in \mathbb{R} \to r \in \mathbb{R} \to 0 < r \to \varepsilon \in \mathbb{R} \to 0 < \varepsilon \to \varphi_{79}$
\item $\varphi_{237} :\equiv (c \in \mathbb{R} \to r \in \mathbb{R} \to 0 < r \to n \in \mathbb{N} \to 0 < \lvert \pi_{c,r}(n) - c \rvert) \to \varphi_{216}$
\item $\varphi_{238} :\equiv (\underline{n} + 1 \in \mathbb{R} \to \underline{n} + 1 \neq 0 \to (\underline{n} + 1)^{-1} \in \mathbb{R}) \to (c \in \mathbb{R} \to r \in \mathbb{R} \to 0 < r \to n \in \mathbb{N} \to \underline{n} + 1 \in \mathbb{R}) \to \varphi_{185}$
\item $\varphi_{239} :\equiv \forall \varepsilon\,(c \in \mathbb{R} \land r \in \mathbb{R} \land 0 < r \to \varphi_{109}) \to \forall \varepsilon\,(c \in \mathbb{R} \land r \in \mathbb{R} \land 0 < r) \to \varphi_{117}$
\item $\varphi_{240} :\equiv 1 \in \mathbb{R} \to (c \in \mathbb{R} \to r \in \mathbb{R} \to 0 < r \to n \in \mathbb{N} \to \underline{n} + 1 \in \mathbb{R}) \to \varphi_{217}$
\item $\varphi_{241} :\equiv (c \in \mathbb{R} \to r \in \mathbb{R} \to 0 < r \to \varepsilon \in \mathbb{R} \to 0 < \varepsilon \to \varphi_{51}) \to c \in \mathbb{R} \to r \in \mathbb{R} \to 0 < r \to \varepsilon \in \mathbb{R} \to 0 < \varepsilon \to \varphi_{41}$
\item $\varphi_{242} :\equiv (\varepsilon \in \mathbb{R} \to r^{-1} \in \mathbb{R} \to \varepsilon \cdot r^{-1} \in \mathbb{R}) \to (c \in \mathbb{R} \to r \in \mathbb{R} \to 0 < r \to \varepsilon \in \mathbb{R} \to \varepsilon \in \mathbb{R}) \to \varphi_{203}$
\item $\varphi_{243} :\equiv (\pi_{c,r}(n) - c = c + r \cdot (\underline{n} + 1)^{-1} - c \to \varphi_{170}) \to c \in \mathbb{R} \to \varphi_{86}$
\item $\varphi_{244} :\equiv (c \in \mathbb{R} \to r \in \mathbb{R} \to 0 < r \to \varepsilon \in \mathbb{R} \to 0 < \varepsilon \to \varphi_{40}) \to c \in \mathbb{R} \to r \in \mathbb{R} \to 0 < r \to \varepsilon \in \mathbb{R} \to 0 < \varepsilon \to \varphi_{68}$
\item $\varphi_{245} :\equiv (c \in \mathbb{R} \to r \in \mathbb{R} \to 0 < r \to \varepsilon \in \mathbb{R} \to \varphi_{88}) \to c \in \mathbb{R} \to r \in \mathbb{R} \to 0 < r \to \varepsilon \in \mathbb{R} \to 0 < \varepsilon \to \varphi_{73}$
\item $\varphi_{246} :\equiv (c \in \mathbb{R} \to r \in \mathbb{R} \to 0 < r \to \varepsilon \in \mathbb{R} \to 0 < \varepsilon \to \varphi_{51}) \to c \in \mathbb{R} \to r \in \mathbb{R} \to 0 < r \to \varepsilon \in \mathbb{R} \to 0 < \varepsilon \to \varphi_{71}$
\item $\varphi_{247} :\equiv (c \in \mathbb{R} \to r \in \mathbb{R} \to 0 < r \to n \in \mathbb{N} \to r \cdot 1 = r) \to \varphi_{46} \to \varphi_{21} \to \varphi_{27}$
\item $\varphi_{248} :\equiv (r \in \mathbb{R} \to (\underline{n} + 1)^{-1} \in \mathbb{R} \to r \cdot (\underline{n} + 1)^{-1} \in \mathbb{R}) \to (c \in \mathbb{R} \to r \in \mathbb{R} \to 0 < r \to n \in \mathbb{N} \to r \in \mathbb{R}) \to \varphi_{194}$
\item $\varphi_{249} :\equiv (c \in \mathbb{R} \to r \in \mathbb{R} \to 0 < r \to \varepsilon \in \mathbb{R} \to 0 < \varepsilon \to \varphi_{42}) \to c \in \mathbb{R} \to r \in \mathbb{R} \to 0 < r \to \varepsilon \in \mathbb{R} \to 0 < \varepsilon \to \varphi_{81}$
\item $\varphi_{250} :\equiv (c \in \mathbb{R} \to r \in \mathbb{R} \to 0 < r \to \mathrm{Seq}(\pi_{c,r})) \to (c \in \mathbb{R} \to r \in \mathbb{R} \to 0 < r \to c \in \mathbb{R}) \to \varphi_{214}$
\item $\varphi_{251} :\equiv (c \in \mathbb{R} \to r \in \mathbb{R} \to 0 < r \to \varphi_{98}) \to c \in \mathbb{R} \to r \in \mathbb{R} \to 0 < r \to \varepsilon \in \mathbb{R} \to 0 < \varepsilon \to \varphi_{133}$
\item $\varphi_{252} :\equiv (c \in \mathbb{R} \to r \in \mathbb{R} \to 0 < r \to \varepsilon \in \mathbb{R} \to 0 < \varepsilon \to \varphi_{71}) \to c \in \mathbb{R} \to r \in \mathbb{R} \to 0 < r \to \varepsilon \in \mathbb{R} \to 0 < \varepsilon \to \varphi_{89}$
\item $\varphi_{253} :\equiv (c \in \mathbb{R} \to r \in \mathbb{R} \to 0 < r \to \varepsilon \in \mathbb{R} \to \varphi_{88}) \to 1 \in \mathbb{R} \to c \in \mathbb{R} \to r \in \mathbb{R} \to 0 < r \to \varepsilon \in \mathbb{R} \to 0 < \varepsilon \to \varphi_{63}$
\item $\varphi_{254} :\equiv (c \in \mathbb{R} \to r \in \mathbb{R} \to 0 < r \to \varepsilon \in \mathbb{R} \to 0 < \varepsilon \to \varphi_{89}) \to c \in \mathbb{R} \to r \in \mathbb{R} \to 0 < r \to \varepsilon \in \mathbb{R} \to 0 < \varepsilon \to \varphi_{76}$
\item $\varphi_{255} :\equiv (c \in \mathbb{R} \to r \in \mathbb{R} \to 0 < r \to \varepsilon \in \mathbb{R} \to 0 < \varepsilon \to \varphi_{72}) \to c \in \mathbb{R} \to r \in \mathbb{R} \to 0 < r \to \varepsilon \in \mathbb{R} \to 0 < \varepsilon \to \varphi_{97}$
\item $\varphi_{256} :\equiv \pi_{c,r}(n) - c = \pi_{c,r}(n) - c \to (\pi_{c,r}(n) = c + r \cdot (\underline{n} + 1)^{-1} \to \varphi_{126}) \to c \in \mathbb{R} \to \varphi_{85}$
\item $\varphi_{257} :\equiv (c \in \mathbb{R} \to r \in \mathbb{R} \to 0 < r \to \varepsilon \in \mathbb{R} \to 0 < \varepsilon \to \varphi_{73}) \to c \in \mathbb{R} \to r \in \mathbb{R} \to 0 < r \to \varepsilon \in \mathbb{R} \to 0 < \varepsilon \to \varphi_{104}$
\item $\varphi_{258} :\equiv (c \in \mathbb{R} \to r \in \mathbb{R} \to 0 < r \to \varepsilon \in \mathbb{R} \to 0 < \varepsilon \to r^{-1} \in \mathbb{R}) \to (c \in \mathbb{R} \to r \in \mathbb{R} \to 0 < r \to \varepsilon \in \mathbb{R} \to 0 < \varepsilon \to 0 < \varepsilon) \to \varphi_{195}$
\item $\varphi_{259} :\equiv (c + r \cdot (\underline{n} + 1)^{-1} = r \cdot (\underline{n} + 1)^{-1} + c \to \varphi_{180}) \to c \in \mathbb{R} \to r \in \mathbb{R} \to 0 < r \to \varphi_{39}$
\item $\varphi_{260} :\equiv (c \in \mathbb{R} \to r \in \mathbb{R} \to 0 < r \to \varepsilon \in \mathbb{R} \to 0 < \varepsilon \to \varphi_{81}) \to c \in \mathbb{R} \to r \in \mathbb{R} \to 0 < r \to \varepsilon \in \mathbb{R} \to 0 < \varepsilon \to \varphi_{102}$
\item $\varphi_{261} :\equiv (c \in \mathbb{R} \to r \in \mathbb{R} \to 0 < r \to \varepsilon \in \mathbb{R} \to 0 < \varepsilon \to \varphi_{79}) \to c \in \mathbb{R} \to r \in \mathbb{R} \to 0 < r \to \varepsilon \in \mathbb{R} \to 0 < \varepsilon \to \varphi_{105}$
\item $\varphi_{262} :\equiv (c \in \mathbb{R} \to r \in \mathbb{R} \to 0 < r \to \varepsilon \in \mathbb{R} \to 0 < \varepsilon \to \varphi_{68}) \to c \in \mathbb{R} \to r \in \mathbb{R} \to 0 < r \to \varepsilon \in \mathbb{R} \to 0 < \varepsilon \to \varphi_{113}$
\item $\varphi_{263} :\equiv (c \in \mathbb{R} \to r \in \mathbb{R} \to 0 < r \to \varepsilon \in \mathbb{R} \to 0 < \varepsilon \to \varphi_{97}) \to c \in \mathbb{R} \to r \in \mathbb{R} \to 0 < r \to \varepsilon \in \mathbb{R} \to 0 < \varepsilon \to \varphi_{107}$
\item $\varphi_{264} :\equiv (c \in \mathbb{R} \to r \in \mathbb{R} \to \varphi_{45}) \to (c \in \mathbb{R} \to r \in \mathbb{R} \to 0 < r \to c \in \mathbb{R}) \to \varphi_{201}$
\item $\varphi_{265} :\equiv (c \in \mathbb{R} \to r \in \mathbb{R} \to 0 < r \to \varepsilon \in \mathbb{R} \to 0 < \varepsilon \to \varphi_{108}) \to \varphi_{122} \to n \in \mathbb{N} \to K \leq n \to \lvert \pi_{c,r}(n) - c \rvert < \varepsilon$
\item $\varphi_{266} :\equiv (c \in \mathbb{R} \to r \in \mathbb{R} \to 0 < r \to \varepsilon \in \mathbb{R} \to 0 < \varepsilon \to \varphi_{76}) \to c \in \mathbb{R} \to \varphi_{200}$
\item $\varphi_{267} :\equiv \lvert \pi_{c,r}(n) - c \rvert = \lvert \pi_{c,r}(n) - c \rvert \to (\pi_{c,r}(n) - c = r \cdot (\underline{n} + 1)^{-1} \to \varphi_{140}) \to c \in \mathbb{R} \to \varphi_{87}$
\item $\varphi_{268} :\equiv \forall n\,(\varphi_{122} \to n \in \mathbb{N} \to K \leq n \to \lvert \pi_{c,r}(n) - c \rvert < \varepsilon) \to \forall n\,(\varphi_{122}) \to \forall n \in \mathbb{N}\,(K \leq n \to \lvert \pi_{c,r}(n) - c \rvert < \varepsilon)$
\item $\varphi_{269} :\equiv (c \in \mathbb{R} \to r \in \mathbb{R} \to 0 < r \to \varepsilon \in \mathbb{R} \to 0 < \varepsilon \to \varphi_{107}) \to c \in \mathbb{R} \to r \in \mathbb{R} \to 0 < r \to \varepsilon \in \mathbb{R} \to 0 < \varepsilon \to \varphi_{132}$
\item $\varphi_{270} :\equiv (\pi_{c,r}(n) = c + r \cdot (\underline{n} + 1)^{-1} \to \varphi_{126}) \to c \in \mathbb{R} \to r \in \mathbb{R} \to 0 < r \to \varepsilon \in \mathbb{R} \to 0 < \varepsilon \to \varphi_{150}$
\item $\varphi_{271} :\equiv (\pi_{c,r}(n) - c = r \cdot (\underline{n} + 1)^{-1} + c - c \to \varphi_{125}) \to c \in \mathbb{R} \to r \in \mathbb{R} \to 0 < r \to \varepsilon \in \mathbb{R} \to 0 < \varepsilon \to \varphi_{134}$
\item $\varphi_{272} :\equiv (c \in \mathbb{R} \to r \in \mathbb{R} \to 0 < r \to \mathrm{Conv}(\pi_{c,r},c)) \to (c \in \mathbb{R} \to r \in \mathbb{R} \to \varphi_{96}) \to c \in \mathbb{R} \to r \in \mathbb{R} \to 0 < r \to \varphi_{162}$
\item $\varphi_{273} :\equiv (c \in \mathbb{R} \to r \in \mathbb{R} \to 0 < r \to \varepsilon \in \mathbb{R} \to 0 < \varepsilon \to \varphi_{102}) \to c \in \mathbb{R} \to \varphi_{206}$
\item $\varphi_{274} :\equiv (c \in \mathbb{R} \to r \in \mathbb{R} \to 0 < r \to \varepsilon \in \mathbb{R} \to 0 < \varepsilon \to \varphi_{114}) \to c \in \mathbb{R} \to r \in \mathbb{R} \to 0 < r \to \varepsilon \in \mathbb{R} \to 0 < \varepsilon \to \varphi_{136}$
\item $\varphi_{275} :\equiv (\underline{n} \in \mathbb{R} \to 1 \in \mathbb{R} \to 0 < \underline{n} \to 0 < 1 \to 0 < \underline{n} + 1) \to (c \in \mathbb{R} \to r \in \mathbb{R} \to 0 < r \to n \in \mathbb{N} \to \underline{n} \in \mathbb{R}) \to \varphi_{215}$
\item $\varphi_{276} :\equiv (\lvert \pi_{c,r}(n) - c \rvert = \lvert r \cdot (\underline{n} + 1)^{-1} \rvert \to \varphi_{124}) \to c \in \mathbb{R} \to r \in \mathbb{R} \to 0 < r \to \varepsilon \in \mathbb{R} \to 0 < \varepsilon \to \varphi_{138}$
\item $\varphi_{277} :\equiv (\pi_{c,r}(n) - c = r \cdot (\underline{n} + 1)^{-1} \to \varphi_{140}) \to c \in \mathbb{R} \to r \in \mathbb{R} \to 0 < r \to \varepsilon \in \mathbb{R} \to 0 < \varepsilon \to \varphi_{152}$
\item $\varphi_{278} :\equiv (c \in \mathbb{R} \to r \in \mathbb{R} \to 0 < r \to n \in \mathbb{N} \to (\underline{n} + 1)^{-1} \in \mathbb{R}) \to 1 \in \mathbb{R} \to \varphi_{224}$
\item $\varphi_{279} :\equiv (r \in \mathbb{R} \to r \neq 0 \to r^{-1} \in \mathbb{R}) \to (c \in \mathbb{R} \to r \in \mathbb{R} \to 0 < r \to \varepsilon \in \mathbb{R} \to \varphi_{64}) \to \varphi_{209}$
\item $\varphi_{280} :\equiv (c \in \mathbb{R} \to r \in \mathbb{R} \to 0 < r \to \varepsilon \in \mathbb{R} \to 0 < \varepsilon \to \varphi_{143}) \to \varphi_{179} \to \neg (K \in \mathbb{N} \land \underline{K}^{-1} < \varepsilon \cdot r^{-1})$
\item $\varphi_{281} :\equiv (\pi_{c,r}(n) - c = c + r \cdot (\underline{n} + 1)^{-1} - c \to \varphi_{170}) \to c \in \mathbb{R} \to r \in \mathbb{R} \to 0 < r \to \varepsilon \in \mathbb{R} \to 0 < \varepsilon \to \varphi_{151}$
\item $\varphi_{282} :\equiv \forall K\,(\varphi_{179} \to \neg (K \in \mathbb{N} \land \underline{K}^{-1} < \varepsilon \cdot r^{-1})) \to \forall K\,(\varphi_{179}) \to \forall K\,(\neg (K \in \mathbb{N} \land \underline{K}^{-1} < \varepsilon \cdot r^{-1}))$
\item $\varphi_{283} :\equiv (\varphi_{122} \to \forall n\,(\varphi_{122})) \to c \in \mathbb{R} \to r \in \mathbb{R} \to 0 < r \to \varepsilon \in \mathbb{R} \to 0 < \varepsilon \to \varphi_{114}$
\item $\varphi_{284} :\equiv c + r \cdot (\underline{n} + 1)^{-1} - c = c + r \cdot (\underline{n} + 1)^{-1} - c \to \varphi_{259}$
\item $\varphi_{285} :\equiv (c + r \cdot (\underline{n} + 1)^{-1} = r \cdot (\underline{n} + 1)^{-1} + c \to \varphi_{180}) \to c \in \mathbb{R} \to r \in \mathbb{R} \to 0 < r \to \varphi_{193}$
\item $\varphi_{286} :\equiv (c \in \mathbb{R} \to r \in \mathbb{R} \to 0 < r \to \varepsilon \in \mathbb{R} \to 0 < \varepsilon \to \varphi_{69}) \to \varphi_{78} \to c \in \mathbb{R} \to r \in \mathbb{R} \to 0 < r \to \varepsilon \in \mathbb{R} \to 0 < \varepsilon \to \varphi_{70}$
\item $\varphi_{287} :\equiv (c \in \mathbb{R} \to r \in \mathbb{R} \to 0 < r \to \varepsilon \in \mathbb{R} \to \varepsilon \in \mathbb{R}) \to (c \in \mathbb{R} \to r \in \mathbb{R} \to 0 < r \to \varepsilon \in \mathbb{R} \to \varphi_{64}) \to \varphi_{210}$
\item $\varphi_{288} :\equiv (\varepsilon \in \mathbb{R} \to r^{-1} \in \mathbb{R} \to \varepsilon \cdot r^{-1} \in \mathbb{R}) \to (c \in \mathbb{R} \to r \in \mathbb{R} \to 0 < r \to \varepsilon \in \mathbb{R} \to \varepsilon \in \mathbb{R}) \to \varphi_{244}$
\item $\varphi_{289} :\equiv (r \in \mathbb{R} \to \varphi_{95}) \to (c \in \mathbb{R} \to r \in \mathbb{R} \to 0 < r \to n \in \mathbb{N} \to r \in \mathbb{R}) \to \varphi_{231}$
\item $\varphi_{290} :\equiv (c \in \mathbb{R} \to r \in \mathbb{R} \to 0 < r \to \varepsilon \in \mathbb{R} \to 0 < \varepsilon \to \varphi_{99}) \to \varphi_{155} \to c \in \mathbb{R} \to r \in \mathbb{R} \to 0 < r \to \varepsilon \in \mathbb{R} \to 0 < \varepsilon \to \varphi_{108}$
\item $\varphi_{291} :\equiv (\varepsilon \in \mathbb{R} \to r^{-1} \in \mathbb{R} \to 0 < \varepsilon \to 0 < r^{-1} \to 0 < \varepsilon \cdot r^{-1}) \to (c \in \mathbb{R} \to r \in \mathbb{R} \to 0 < r \to \varepsilon \in \mathbb{R} \to \varepsilon \in \mathbb{R}) \to \varphi_{258}$
\item $\varphi_{292} :\equiv c + r \cdot (\underline{n} + 1)^{-1} - c = c + r \cdot (\underline{n} + 1)^{-1} - c \to \varphi_{285}$
\item $\varphi_{293} :\equiv (c \in \mathbb{R} \to r \in \mathbb{R} \to 0 < r \to \varepsilon \in \mathbb{R} \to 0 < \varepsilon \to \varphi_{138}) \to \varphi_{123} \to c \in \mathbb{R} \to r \in \mathbb{R} \to 0 < r \to \varepsilon \in \mathbb{R} \to 0 < \varepsilon \to \varphi_{139}$
\item $\varphi_{294} :\equiv (c \in \mathbb{R} \to r \in \mathbb{R} \to 0 < r \to \varepsilon \in \mathbb{R} \to 0 < \varepsilon \to \varphi_{102}) \to (c \in \mathbb{R} \to r \in \mathbb{R} \to 0 < r \to \varepsilon \in \mathbb{R} \to \varphi_{65}) \to c \in \mathbb{R} \to \varphi_{205}$
\item $\varphi_{295} :\equiv (c \in \mathbb{R} \to r \in \mathbb{R} \to 0 < r \to \varepsilon \in \mathbb{R} \to 0 < \varepsilon \to \varphi_{81}) \to (c \in \mathbb{R} \to r \in \mathbb{R} \to 0 < r \to 0 < r) \to \varphi_{252}$
\item $\varphi_{296} :\equiv (c \in \mathbb{R} \to r \in \mathbb{R} \to 0 < r \to \varepsilon \in \mathbb{R} \to 0 < \varepsilon \to \varphi_{136}) \to \varphi_{188} \to c \in \mathbb{R} \to r \in \mathbb{R} \to 0 < r \to \varepsilon \in \mathbb{R} \to 0 < \varepsilon \to \varphi_{143}$
\item $\varphi_{297} :\equiv (c \in \mathbb{R} \to r \in \mathbb{R} \to 0 < r \to \varphi_{162}) \to (\varphi_{83} \to \neg \varphi_{162}) \to c \in \mathbb{R} \to r \in \mathbb{R} \to 0 < r \to \varphi_{82}$
\item $\varphi_{298} :\equiv (\underline{n} + 1 \in \mathbb{R} \to 0 < \underline{n} + 1 \to 0 < (\underline{n} + 1)^{-1}) \to (c \in \mathbb{R} \to r \in \mathbb{R} \to 0 < r \to \varepsilon \in \mathbb{R} \to 0 < \varepsilon \to \varphi_{63}) \to \varphi_{246}$
\item $\varphi_{299} :\equiv (c \in \mathbb{R} \to r \in \mathbb{R} \to 0 < r \to \varepsilon \in \mathbb{R} \to \varphi_{64}) \to (c \in \mathbb{R} \to r \in \mathbb{R} \to 0 < r \to 0 < r) \to \varphi_{269}$
\item $\varphi_{300} :\equiv (\underline{n} + 1 \in \mathbb{R} \to \underline{n} + 1 \neq 0 \to (\underline{n} + 1)^{-1} \in \mathbb{R}) \to (c \in \mathbb{R} \to r \in \mathbb{R} \to 0 < r \to \varepsilon \in \mathbb{R} \to 0 < \varepsilon \to \varphi_{63}) \to \varphi_{249}$
\item $\varphi_{301} :\equiv (r \in \mathbb{R} \to (\underline{n} + 1)^{-1} \in \mathbb{R} \to r \cdot (\underline{n} + 1)^{-1} \in \mathbb{R}) \to (c \in \mathbb{R} \to r \in \mathbb{R} \to 0 < r \to \varepsilon \in \mathbb{R} \to \varphi_{64}) \to \varphi_{260}$
\item $\varphi_{302} :\equiv (c \in \mathbb{R} \to r \in \mathbb{R} \to \varphi_{80}) \to (c \in \mathbb{R} \to r \in \mathbb{R} \to 0 < r \to \varphi_{98}) \to \varphi_{236}$
\item $\varphi_{303} :\equiv (c \in \mathbb{R} \to r \in \mathbb{R} \to 0 < r \to \varepsilon \in \mathbb{R} \to 0 < \varepsilon \to \varphi_{133}) \to \pi_{c,r}(n) - c = \pi_{c,r}(n) - c \to \varphi_{270}$
\item $\varphi_{304} :\equiv (c \in \mathbb{R} \to r \in \mathbb{R} \to 0 < r \to \varepsilon \in \mathbb{R} \to \varphi_{88}) \to 1 \in \mathbb{R} \to (c \in \mathbb{R} \to r \in \mathbb{R} \to 0 < r \to \varepsilon \in \mathbb{R} \to \varphi_{74}) \to \varphi_{233}$
\item $\varphi_{305} :\equiv (c \in \mathbb{R} \to r \in \mathbb{R} \to 0 < r \to \varepsilon \in \mathbb{R} \to 0 < \varepsilon \to \varphi_{134}) \to \lvert \pi_{c,r}(n) - c \rvert = \lvert \pi_{c,r}(n) - c \rvert \to \varphi_{277}$
\item $\varphi_{306} :\equiv (c \in \mathbb{R} \to r \in \mathbb{R} \to 0 < r \to \varepsilon \in \mathbb{R} \to 0 < \varepsilon \to \varphi_{70}) \to (c \in \mathbb{R} \to r \in \mathbb{R} \to 0 < r \to \varepsilon \in \mathbb{R} \to 0 < \varepsilon \to \varphi_{132}) \to \varphi_{226}$
\item $\varphi_{307} :\equiv (c \in \mathbb{R} \to \varphi_{84}) \to (c \in \mathbb{R} \to r \in \mathbb{R} \to 0 < r \to \varepsilon \in \mathbb{R} \to \varphi_{65}) \to \varphi_{273}$
\item $\varphi_{308} :\equiv (c \in \mathbb{R} \to r \in \mathbb{R} \to 0 < r \to \varepsilon \in \mathbb{R} \to 0 < \varepsilon \to \varphi_{63}) \to (c \in \mathbb{R} \to r \in \mathbb{R} \to 0 < r \to \varepsilon \in \mathbb{R} \to \varphi_{74}) \to \varphi_{257}$
\item $\varphi_{309} :\equiv \forall \varepsilon\,(c \in \mathbb{R} \land r \in \mathbb{R} \land 0 < r \to \varphi_{109}) \to \varphi_{239} \to \varphi_{19} \to c \in \mathbb{R} \to r \in \mathbb{R} \to 0 < r \to \varphi_{117}$
\item $\varphi_{310} :\equiv (c \in \mathbb{R} \to r \in \mathbb{R} \to 0 < r \to \varepsilon \in \mathbb{R} \to 0 < \varepsilon \to \varphi_{151}) \to (c \in \mathbb{R} \to \varphi_{205}) \to \varphi_{271}$
\item $\varphi_{311} :\equiv (c \in \mathbb{R} \to r \in \mathbb{R} \to 0 < r \to \varepsilon \in \mathbb{R} \to 0 < \varepsilon \to \varphi_{152}) \to (c \in \mathbb{R} \to \varphi_{200}) \to \varphi_{276}$
\item $\varphi_{312} :\equiv (c \in \mathbb{R} \to r \in \mathbb{R} \to 0 < r \to \varepsilon \in \mathbb{R} \to 0 < \varepsilon \to \varphi_{150}) \to (c \in \mathbb{R} \to r \in \mathbb{R} \to 0 < r \to \varphi_{193}) \to \varphi_{281}$
\item $\varphi_{313} :\equiv (r \in \mathbb{R} \to \varphi_{95}) \to (c \in \mathbb{R} \to r \in \mathbb{R} \to 0 < r \to \varepsilon \in \mathbb{R} \to \varphi_{64}) \to \varphi_{295}$
\item $\varphi_{314} :\equiv (c \in \mathbb{R} \to r \in \mathbb{R} \to 0 < r \to \varepsilon \in \mathbb{R} \to 0 < \varepsilon \to \varphi_{81}) \to (c \in \mathbb{R} \to r \in \mathbb{R} \to 0 < r \to \varepsilon \in \mathbb{R} \to 0 < \varepsilon \to \varphi_{68}) \to \varphi_{299}$
\item $\varphi_{315} :\equiv (\underline{n} \in \mathbb{R} \to \varphi_{112}) \to (c \in \mathbb{R} \to r \in \mathbb{R} \to 0 < r \to \varepsilon \in \mathbb{R} \to \varphi_{88}) \to \varphi_{308}$
\item $\varphi_{316} :\equiv \forall K\,(\varphi_{179} \to \neg (K \in \mathbb{N} \land \underline{K}^{-1} < \varepsilon \cdot r^{-1})) \to \varphi_{282} \to (\varphi_{179} \to \forall K\,(\varphi_{179})) \to c \in \mathbb{R} \to r \in \mathbb{R} \to 0 < r \to \varphi_{109}$
\end{abbrevlist}
\end{fullproof}

% 自動生成 (prover/build.js): SeqCritInv  ノード数 453
\begin{fullproof}{定理1.7(⇐)}{$\forall f\,\forall c\,\forall \alpha\,(\mathrm{RFn}(f) \to c \in \mathbb{R} \to \alpha \in \mathbb{R} \to \mathrm{PNbhd}(f,c) \to \forall a\,(\mathrm{Seq}(a) \to \mathrm{Conv}(a,c) \to \forall n \in \mathbb{N}\,(a_{n} \in \mathrm{dom}(f) \land a_{n} \neq c) \to \mathrm{Conv}(f \circ a,\alpha)) \to \mathrm{FLim}(f,c,\alpha))$}{453}\phantomsection\label{tree:SeqCritInv}
\noindent$\Pi_{1}$:
\begin{ptree}
\AXC{$\forall x\,(\neg x < x)$\lbl{[O1]}}
\AXC{$\forall x\,(\neg x < x) \to \neg 0 < 0$\lbl{A2}}
\BIC{$\neg 0 < 0$}
\AXC{$\varphi_{90}$\lbl{A1}}
\BIC{$\mathrm{RFn}(f) \to c \in \mathbb{R} \to \alpha \in \mathbb{R} \to \mathrm{PNbhd}(f,c) \to \varphi_{80}$}
\end{ptree}
\noindent$\Pi_{2}$:
\begin{ptree}
\AXC{$\mathrm{RFn}(f) \to c \in \mathbb{R} \to \alpha \in \mathbb{R} \to \mathrm{PNbhd}(f,c) \to \varphi_{205}$\lbl{A1}}
\AXC{$\varphi_{272}$\lbl{A1}}
\BIC{$\mathrm{RFn}(f) \to c \in \mathbb{R} \to \alpha \in \mathbb{R} \to \mathrm{PNbhd}(f,c) \to \varphi_{175}$}
\end{ptree}
\noindent$\Pi_{3}$:
\begin{ptree}
\AXC{$\varphi_{61} \to N \in \mathbb{N} \to \varphi_{50}$\lbl{A2}}
\AXC{$\mathrm{RFn}(f) \to c \in \mathbb{R} \to \alpha \in \mathbb{R} \to \varphi_{209}$\lbl{A1}}
\AXC{$\varphi_{282}$\lbl{A1}}
\BIC{$\mathrm{RFn}(f) \to c \in \mathbb{R} \to \alpha \in \mathbb{R} \to \mathrm{PNbhd}(f,c) \to \varphi_{202}$}
\AXC{$\varphi_{292}$\lbl{A1}}
\BIC{$(\varphi_{61} \to N \in \mathbb{N} \to \varphi_{50}) \to \mathrm{RFn}(f) \to c \in \mathbb{R} \to \alpha \in \mathbb{R} \to \varphi_{206}$}
\BIC{$\mathrm{RFn}(f) \to c \in \mathbb{R} \to \alpha \in \mathbb{R} \to \varphi_{206}$}
\end{ptree}
\noindent$\Pi_{4}$:
\begin{ptree}
\AXC{$\stackrel{\Pi_{3}}{\vdots}$}\noLine
\UIC{$\mathrm{RFn}(f) \to c \in \mathbb{R} \to \alpha \in \mathbb{R} \to \varphi_{206}$}
\AXC{$(\mathrm{RFn}(f) \to c \in \mathbb{R} \to \alpha \in \mathbb{R} \to \varphi_{206}) \to \varphi_{281}$\lbl{A1}}
\BIC{$\varphi_{281}$}
\end{ptree}
\noindent$\Pi_{5}$:
\begin{ptree}
\AXC{$\stackrel{\Pi_{2}}{\vdots}$}\noLine
\UIC{$\mathrm{RFn}(f) \to c \in \mathbb{R} \to \alpha \in \mathbb{R} \to \mathrm{PNbhd}(f,c) \to \varphi_{175}$}
\AXC{$\stackrel{\Pi_{4}}{\vdots}$}\noLine
\UIC{$\varphi_{281}$}
\BIC{$\mathrm{RFn}(f) \to c \in \mathbb{R} \to \alpha \in \mathbb{R} \to \mathrm{PNbhd}(f,c) \to \varphi_{210}$}
\end{ptree}
\noindent$\Pi_{6}$:
\begin{ptree}
\AXC{$z = w \to \lvert z - \alpha \rvert < e_{0} \to \lvert w - \alpha \rvert < e_{0}$\lbl{A6}}
\UIC{$\forall w\,(z = w \to \lvert z - \alpha \rvert < e_{0} \to \lvert w - \alpha \rvert < e_{0})$}
\UIC{$\forall z\,\forall w\,(z = w \to \lvert z - \alpha \rvert < e_{0} \to \lvert w - \alpha \rvert < e_{0})$}
\AXC{$\forall z\,\forall w\,(z = w \to \lvert z - \alpha \rvert < e_{0} \to \lvert w - \alpha \rvert < e_{0}) \to \varphi_{17}$\lbl{A2}}
\BIC{$\varphi_{17}$}
\AXC{$\varphi_{17} \to \varphi_{43}$\lbl{A2}}
\BIC{$\varphi_{43}$}
\end{ptree}
\noindent$\Pi_{7}$:
\begin{ptree}
\AXC{$\varphi_{48}$\lbl{A2}}
\AXC{$\mathrm{RFn}(f) \to c \in \mathbb{R} \to \alpha \in \mathbb{R} \to \mathrm{PNbhd}(f,c) \to \varphi_{205}$\lbl{A1}}
\AXC{$\varphi_{280}$\lbl{A1}}
\BIC{$\mathrm{RFn}(f) \to c \in \mathbb{R} \to \alpha \in \mathbb{R} \to \mathrm{PNbhd}(f,c) \to \varphi_{199}$}
\AXC{$\varphi_{289}$\lbl{A1}}
\BIC{$\varphi_{216}$}
\BIC{$\mathrm{RFn}(f) \to c \in \mathbb{R} \to \alpha \in \mathbb{R} \to \mathrm{PNbhd}(f,c) \to \varphi_{198}$}
\AXC{$\varphi_{321}$\lbl{A1}}
\BIC{$\varphi_{305}$}
\end{ptree}
\noindent$\Pi_{8}$:
\begin{ptree}
\AXC{$\stackrel{\Pi_{2}}{\vdots}$}\noLine
\UIC{$\mathrm{RFn}(f) \to c \in \mathbb{R} \to \alpha \in \mathbb{R} \to \mathrm{PNbhd}(f,c) \to \varphi_{175}$}
\AXC{$\stackrel{\Pi_{7}}{\vdots}$}\noLine
\UIC{$\varphi_{305}$}
\BIC{$\varphi_{260}$}
\end{ptree}
\noindent$\Pi_{9}$:
\begin{ptree}
\AXC{$x = x$\lbl{A5}}
\UIC{$\forall x\,(x = x)$}
\AXC{$\forall x\,(x = x) \to N = N$\lbl{A2}}
\BIC{$N = N$}
\AXC{$\varphi_{200}$\lbl{A1}}
\BIC{$\mathrm{RFn}(f) \to c \in \mathbb{R} \to \alpha \in \mathbb{R} \to \mathrm{PNbhd}(f,c) \to \varphi_{173}$}
\AXC{$\stackrel{\Pi_{8}}{\vdots}$}\noLine
\UIC{$\varphi_{260}$}
\BIC{$\mathrm{RFn}(f) \to c \in \mathbb{R} \to \alpha \in \mathbb{R} \to \mathrm{PNbhd}(f,c) \to \varphi_{189}$}
\end{ptree}
\noindent$\Pi_{10}$:
\begin{ptree}
\AXC{$\varphi_{61} \to n \in \mathbb{N} \to \varphi_{51}$\lbl{A2}}
\AXC{$\mathrm{RFn}(f) \to c \in \mathbb{R} \to \alpha \in \mathbb{R} \to \varphi_{209}$\lbl{A1}}
\AXC{$\varphi_{282}$\lbl{A1}}
\BIC{$\mathrm{RFn}(f) \to c \in \mathbb{R} \to \alpha \in \mathbb{R} \to \mathrm{PNbhd}(f,c) \to \varphi_{202}$}
\AXC{$\varphi_{293}$\lbl{A1}}
\BIC{$(\varphi_{61} \to n \in \mathbb{N} \to \varphi_{51}) \to \mathrm{RFn}(f) \to c \in \mathbb{R} \to \alpha \in \mathbb{R} \to \varphi_{207}$}
\BIC{$\mathrm{RFn}(f) \to c \in \mathbb{R} \to \alpha \in \mathbb{R} \to \varphi_{207}$}
\end{ptree}
\noindent$\Pi_{11}$:
\begin{ptree}
\AXC{$\forall a\,\forall n\,(\mathrm{Seq}(a) \to n \in \mathbb{N} \to a_{n} \in \mathbb{R})$\lbl{[S1]}}
\AXC{$\varphi_{21}$\lbl{A2}}
\BIC{$\forall n\,(\mathrm{Seq}(a) \to n \in \mathbb{N} \to a_{n} \in \mathbb{R})$}
\AXC{$\varphi_{13}$\lbl{A2}}
\BIC{$\mathrm{Seq}(a) \to n \in \mathbb{N} \to a_{n} \in \mathbb{R}$}
\AXC{$\varphi_{299}$\lbl{A1}}
\BIC{$(\mathrm{RFn}(f) \to c \in \mathbb{R} \to \alpha \in \mathbb{R} \to \varphi_{139}) \to \varphi_{242}$}
\end{ptree}
\noindent$\Pi_{12}$:
\begin{ptree}
\AXC{$\mathrm{RFn}(f) \to c \in \mathbb{R} \to \alpha \in \mathbb{R} \to \varphi_{209}$\lbl{A1}}
\AXC{$\varphi_{253}$\lbl{A1}}
\BIC{$\mathrm{RFn}(f) \to c \in \mathbb{R} \to \alpha \in \mathbb{R} \to \varphi_{139}$}
\AXC{$\stackrel{\Pi_{11}}{\vdots}$}\noLine
\UIC{$(\mathrm{RFn}(f) \to c \in \mathbb{R} \to \alpha \in \mathbb{R} \to \varphi_{139}) \to \varphi_{242}$}
\BIC{$\varphi_{242}$}
\end{ptree}
\noindent$\Pi_{13}$:
\begin{ptree}
\AXC{$\forall y\,\forall c\,(y \in \mathbb{R} \to c \in \mathbb{R} \to 0 < \lvert y - c \rvert \to y \neq c)$\lbl{[E37]}}
\AXC{$\varphi_{59}$\lbl{A2}}
\BIC{$\forall c\,(a_{n} \in \mathbb{R} \to c \in \mathbb{R} \to 0 < \lvert a_{n} - c \rvert \to a_{n} \neq c)$}
\AXC{$\varphi_{66}$\lbl{A2}}
\BIC{$a_{n} \in \mathbb{R} \to c \in \mathbb{R} \to 0 < \lvert a_{n} - c \rvert \to a_{n} \neq c$}
\end{ptree}
\noindent$\Pi_{14}$:
\begin{ptree}
\AXC{$\stackrel{\Pi_{13}}{\vdots}$}\noLine
\UIC{$a_{n} \in \mathbb{R} \to c \in \mathbb{R} \to 0 < \lvert a_{n} - c \rvert \to a_{n} \neq c$}
\AXC{$(a_{n} \in \mathbb{R} \to c \in \mathbb{R} \to 0 < \lvert a_{n} - c \rvert \to a_{n} \neq c) \to \varphi_{301}$\lbl{A1}}
\BIC{$\varphi_{301}$}
\end{ptree}
\noindent$\Pi_{15}$:
\begin{ptree}
\AXC{$\mathrm{RFn}(f) \to c \in \mathbb{R} \to c \in \mathbb{R}$\lbl{A1}}
\AXC{$\mathrm{RFn}(f) \to c \in \mathbb{R} \to \alpha \in \mathbb{R} \to \varphi_{147}$\lbl{A1}}
\AXC{$\stackrel{\Pi_{12}}{\vdots}$}\noLine
\UIC{$\varphi_{242}$}
\BIC{$\mathrm{RFn}(f) \to c \in \mathbb{R} \to \alpha \in \mathbb{R} \to \varphi_{153}$}
\AXC{$\stackrel{\Pi_{14}}{\vdots}$}\noLine
\UIC{$\varphi_{301}$}
\BIC{$(\mathrm{RFn}(f) \to c \in \mathbb{R} \to c \in \mathbb{R}) \to \varphi_{244}$}
\BIC{$\varphi_{244}$}
\end{ptree}
\noindent$\Pi_{16}$:
\begin{ptree}
\AXC{$\stackrel{\Pi_{10}}{\vdots}$}\noLine
\UIC{$\mathrm{RFn}(f) \to c \in \mathbb{R} \to \alpha \in \mathbb{R} \to \varphi_{207}$}
\AXC{$(\mathrm{RFn}(f) \to c \in \mathbb{R} \to \alpha \in \mathbb{R} \to \varphi_{207}) \to \varphi_{246}$\lbl{A1}}
\BIC{$\varphi_{246}$}
\end{ptree}
\noindent$\Pi_{17}$:
\begin{ptree}
\AXC{$\stackrel{\Pi_{10}}{\vdots}$}\noLine
\UIC{$\mathrm{RFn}(f) \to c \in \mathbb{R} \to \alpha \in \mathbb{R} \to \varphi_{207}$}
\AXC{$\varphi_{257}$\lbl{A1}}
\BIC{$\mathrm{RFn}(f) \to c \in \mathbb{R} \to \alpha \in \mathbb{R} \to \varphi_{159}$}
\AXC{$\stackrel{\Pi_{15}}{\vdots}$}\noLine
\UIC{$\varphi_{244}$}
\BIC{$\mathrm{RFn}(f) \to c \in \mathbb{R} \to \alpha \in \mathbb{R} \to \varphi_{149}$}
\AXC{$\stackrel{\Pi_{16}}{\vdots}$}\noLine
\UIC{$\varphi_{246}$}
\BIC{$\mathrm{RFn}(f) \to c \in \mathbb{R} \to \alpha \in \mathbb{R} \to \varphi_{163}$}
\AXC{$(\mathrm{RFn}(f) \to c \in \mathbb{R} \to \alpha \in \mathbb{R} \to \varphi_{163}) \to \varphi_{184}$\lbl{A1}}
\BIC{$\varphi_{184}$}
\UIC{$\forall n\,(\varphi_{184})$}
\end{ptree}
\noindent$\Pi_{18}$:
\begin{ptree}
\AXC{$\forall n\,(\varphi_{184}) \to \varphi_{185}$\lbl{A3}}
\AXC{$\stackrel{\Pi_{17}}{\vdots}$}\noLine
\UIC{$\forall n\,(\varphi_{184})$}
\AXC{$\forall n\,(\varphi_{184}) \to (\forall n\,(\varphi_{184}) \to \varphi_{185}) \to \varphi_{296}$\lbl{A1}}
\BIC{$(\forall n\,(\varphi_{184}) \to \varphi_{185}) \to \varphi_{296}$}
\BIC{$\varphi_{296}$}
\end{ptree}
\noindent$\Pi_{19}$:
\begin{ptree}
\AXC{$\varphi_{106} \land (\mathrm{Seq}(a) \land \varphi_{61}) \to \forall n\,(\varphi_{106} \land (\mathrm{Seq}(a) \land \varphi_{61}))$\lbl{A4}}
\AXC{$\stackrel{\Pi_{18}}{\vdots}$}\noLine
\UIC{$\varphi_{296}$}
\BIC{$\mathrm{RFn}(f) \to c \in \mathbb{R} \to \alpha \in \mathbb{R} \to \mathrm{PNbhd}(f,c) \to \varphi_{160}$}
\end{ptree}
\noindent$\Pi_{20}$:
\begin{ptree}
\AXC{$\varphi_{33}$\lbl{A2}}
\AXC{$\stackrel{\Pi_{19}}{\vdots}$}\noLine
\UIC{$\mathrm{RFn}(f) \to c \in \mathbb{R} \to \alpha \in \mathbb{R} \to \mathrm{PNbhd}(f,c) \to \varphi_{160}$}
\AXC{$\varphi_{265}$\lbl{A1}}
\BIC{$\varphi_{33} \to \mathrm{RFn}(f) \to c \in \mathbb{R} \to \alpha \in \mathbb{R} \to \varphi_{163}$}
\BIC{$\mathrm{RFn}(f) \to c \in \mathbb{R} \to \alpha \in \mathbb{R} \to \varphi_{163}$}
\AXC{$(\mathrm{RFn}(f) \to c \in \mathbb{R} \to \alpha \in \mathbb{R} \to \varphi_{163}) \to \varphi_{243}$\lbl{A1}}
\BIC{$\varphi_{243}$}
\end{ptree}
\noindent$\Pi_{21}$:
\begin{ptree}
\AXC{$\varphi_{178} \to \varphi_{179}$\lbl{A3}}
\AXC{$\mathrm{RFn}(f) \to c \in \mathbb{R} \to \alpha \in \mathbb{R} \to \varphi_{147}$\lbl{A1}}
\AXC{$\stackrel{\Pi_{20}}{\vdots}$}\noLine
\UIC{$\varphi_{243}$}
\BIC{$\mathrm{RFn}(f) \to c \in \mathbb{R} \to \alpha \in \mathbb{R} \to \varphi_{155}$}
\AXC{$\varphi_{245}$\lbl{A1}}
\BIC{$\varphi_{106} \land (\mathrm{Seq}(a) \land \varphi_{61}) \to n \in \mathbb{N} \to a_{n} \in \mathrm{dom}(f)$}
\UIC{$\varphi_{178}$}
\AXC{$\varphi_{178} \to (\varphi_{178} \to \varphi_{179}) \to \varphi_{295}$\lbl{A1}}
\BIC{$(\varphi_{178} \to \varphi_{179}) \to \varphi_{295}$}
\BIC{$\varphi_{295}$}
\end{ptree}
\noindent$\Pi_{22}$:
\begin{ptree}
\AXC{$\forall f\,\forall a\,(\mathrm{RFn}(f) \to \mathrm{Seq}(a) \to \varphi_{40})$\lbl{[S8]}}
\AXC{$\varphi_{82}$\lbl{A2}}
\BIC{$\forall a\,(\mathrm{RFn}(f) \to \mathrm{Seq}(a) \to \varphi_{40})$}
\AXC{$\varphi_{81}$\lbl{A2}}
\BIC{$\mathrm{RFn}(f) \to \mathrm{Seq}(a) \to \varphi_{40}$}
\AXC{$(\mathrm{RFn}(f) \to \mathrm{Seq}(a) \to \varphi_{40}) \to \varphi_{302}$\lbl{A1}}
\BIC{$\varphi_{302}$}
\end{ptree}
\noindent$\Pi_{23}$:
\begin{ptree}
\AXC{$\mathrm{RFn}(f) \to c \in \mathbb{R} \to \alpha \in \mathbb{R} \to \varphi_{209}$\lbl{A1}}
\AXC{$\varphi_{253}$\lbl{A1}}
\BIC{$\mathrm{RFn}(f) \to c \in \mathbb{R} \to \alpha \in \mathbb{R} \to \varphi_{139}$}
\AXC{$\mathrm{RFn}(f) \to \mathrm{RFn}(f)$\lbl{A1}}
\AXC{$\stackrel{\Pi_{22}}{\vdots}$}\noLine
\UIC{$\varphi_{302}$}
\BIC{$(\mathrm{RFn}(f) \to c \in \mathbb{R} \to \alpha \in \mathbb{R} \to \varphi_{139}) \to \varphi_{250}$}
\BIC{$\varphi_{250}$}
\end{ptree}
\noindent$\Pi_{24}$:
\begin{ptree}
\AXC{$\varphi_{35}$\lbl{A2}}
\AXC{$\varphi_{106} \land (\mathrm{Seq}(a) \land \varphi_{61}) \to \forall n\,(\varphi_{106} \land (\mathrm{Seq}(a) \land \varphi_{61}))$\lbl{A4}}
\AXC{$\stackrel{\Pi_{21}}{\vdots}$}\noLine
\UIC{$\varphi_{295}$}
\BIC{$\mathrm{RFn}(f) \to c \in \mathbb{R} \to \alpha \in \mathbb{R} \to \varphi_{157}$}
\AXC{$\stackrel{\Pi_{23}}{\vdots}$}\noLine
\UIC{$\varphi_{250}$}
\BIC{$\mathrm{RFn}(f) \to c \in \mathbb{R} \to \alpha \in \mathbb{R} \to \mathrm{PNbhd}(f,c) \to \varphi_{171}$}
\AXC{$\varphi_{266}$\lbl{A1}}
\BIC{$\mathrm{RFn}(f) \to c \in \mathbb{R} \to \alpha \in \mathbb{R} \to \mathrm{PNbhd}(f,c) \to \varphi_{195}$}
\AXC{$\varphi_{287}$\lbl{A1}}
\BIC{$\varphi_{215}$}
\BIC{$\mathrm{RFn}(f) \to c \in \mathbb{R} \to \alpha \in \mathbb{R} \to \mathrm{PNbhd}(f,c) \to \varphi_{193}$}
\end{ptree}
\noindent$\Pi_{25}$:
\begin{ptree}
\AXC{$\stackrel{\Pi_{2}}{\vdots}$}\noLine
\UIC{$\mathrm{RFn}(f) \to c \in \mathbb{R} \to \alpha \in \mathbb{R} \to \mathrm{PNbhd}(f,c) \to \varphi_{175}$}
\AXC{$\stackrel{\Pi_{24}}{\vdots}$}\noLine
\UIC{$\mathrm{RFn}(f) \to c \in \mathbb{R} \to \alpha \in \mathbb{R} \to \mathrm{PNbhd}(f,c) \to \varphi_{193}$}
\AXC{$\varphi_{309}$\lbl{A1}}
\BIC{$\varphi_{262}$}
\BIC{$\mathrm{RFn}(f) \to c \in \mathbb{R} \to \alpha \in \mathbb{R} \to \mathrm{PNbhd}(f,c) \to \varphi_{188}$}
\end{ptree}
\noindent$\Pi_{26}$:
\begin{ptree}
\AXC{$\stackrel{\Pi_{9}}{\vdots}$}\noLine
\UIC{$\mathrm{RFn}(f) \to c \in \mathbb{R} \to \alpha \in \mathbb{R} \to \mathrm{PNbhd}(f,c) \to \varphi_{189}$}
\AXC{$\stackrel{\Pi_{25}}{\vdots}$}\noLine
\UIC{$\mathrm{RFn}(f) \to c \in \mathbb{R} \to \alpha \in \mathbb{R} \to \mathrm{PNbhd}(f,c) \to \varphi_{188}$}
\AXC{$\varphi_{314}$\lbl{A1}}
\BIC{$\varphi_{286}$}
\BIC{$\varphi_{214}$}
\end{ptree}
\noindent$\Pi_{27}$:
\begin{ptree}
\AXC{$\stackrel{\Pi_{5}}{\vdots}$}\noLine
\UIC{$\mathrm{RFn}(f) \to c \in \mathbb{R} \to \alpha \in \mathbb{R} \to \mathrm{PNbhd}(f,c) \to \varphi_{210}$}
\AXC{$\stackrel{\Pi_{6}}{\vdots}$}\noLine
\UIC{$\varphi_{43}$}
\AXC{$\stackrel{\Pi_{26}}{\vdots}$}\noLine
\UIC{$\varphi_{214}$}
\BIC{$\mathrm{RFn}(f) \to c \in \mathbb{R} \to \alpha \in \mathbb{R} \to \mathrm{PNbhd}(f,c) \to \varphi_{186}$}
\AXC{$\varphi_{312}$\lbl{A1}}
\BIC{$\varphi_{279}$}
\BIC{$\mathrm{RFn}(f) \to c \in \mathbb{R} \to \alpha \in \mathbb{R} \to \mathrm{PNbhd}(f,c) \to \varphi_{174}$}
\end{ptree}
\noindent$\Pi_{28}$:
\begin{ptree}
\AXC{$\forall K\,\forall n\,(K \in \mathbb{N} \to n \in \mathbb{N} \to K \leq n \to \underline{n}^{-1} \leq \underline{K}^{-1})$\lbl{[C25]}}
\AXC{$\varphi_{53}$\lbl{A2}}
\BIC{$\forall n\,(K \in \mathbb{N} \to n \in \mathbb{N} \to K \leq n \to \underline{n}^{-1} \leq \underline{K}^{-1})$}
\AXC{$\varphi_{49}$\lbl{A2}}
\BIC{$K \in \mathbb{N} \to n \in \mathbb{N} \to K \leq n \to \underline{n}^{-1} \leq \underline{K}^{-1}$}
\end{ptree}
\noindent$\Pi_{29}$:
\begin{ptree}
\AXC{$\stackrel{\Pi_{28}}{\vdots}$}\noLine
\UIC{$K \in \mathbb{N} \to n \in \mathbb{N} \to K \leq n \to \underline{n}^{-1} \leq \underline{K}^{-1}$}
\AXC{$(K \in \mathbb{N} \to n \in \mathbb{N} \to K \leq n \to \underline{n}^{-1} \leq \underline{K}^{-1}) \to \varphi_{320}$\lbl{A1}}
\BIC{$\varphi_{320}$}
\end{ptree}
\noindent$\Pi_{30}$:
\begin{ptree}
\AXC{$\mathrm{RFn}(f) \to c \in \mathbb{R} \to \alpha \in \mathbb{R} \to \mathrm{PNbhd}(f,c) \to \varphi_{182}$\lbl{A1}}
\AXC{$\varphi_{256}$\lbl{A1}}
\BIC{$\mathrm{RFn}(f) \to c \in \mathbb{R} \to \alpha \in \mathbb{R} \to \mathrm{PNbhd}(f,c) \to \varphi_{170}$}
\AXC{$\stackrel{\Pi_{29}}{\vdots}$}\noLine
\UIC{$\varphi_{320}$}
\BIC{$\varphi_{306}$}
\end{ptree}
\noindent$\Pi_{31}$:
\begin{ptree}
\AXC{$\mathrm{RFn}(f) \to c \in \mathbb{R} \to \alpha \in \mathbb{R} \to \mathrm{PNbhd}(f,c) \to \varphi_{176}$\lbl{A1}}
\AXC{$\mathrm{RFn}(f) \to c \in \mathbb{R} \to \alpha \in \mathbb{R} \to \mathrm{PNbhd}(f,c) \to \varphi_{177}$\lbl{A1}}
\AXC{$\stackrel{\Pi_{30}}{\vdots}$}\noLine
\UIC{$\varphi_{306}$}
\BIC{$\varphi_{261}$}
\BIC{$\mathrm{RFn}(f) \to c \in \mathbb{R} \to \alpha \in \mathbb{R} \to \varphi_{191}$}
\end{ptree}
\noindent$\Pi_{32}$:
\begin{ptree}
\AXC{$\varphi_{61} \to n \in \mathbb{N} \to \varphi_{51}$\lbl{A2}}
\AXC{$\mathrm{RFn}(f) \to c \in \mathbb{R} \to \alpha \in \mathbb{R} \to \varphi_{209}$\lbl{A1}}
\AXC{$\varphi_{282}$\lbl{A1}}
\BIC{$\mathrm{RFn}(f) \to c \in \mathbb{R} \to \alpha \in \mathbb{R} \to \mathrm{PNbhd}(f,c) \to \varphi_{202}$}
\AXC{$\varphi_{293}$\lbl{A1}}
\BIC{$(\varphi_{61} \to n \in \mathbb{N} \to \varphi_{51}) \to \mathrm{RFn}(f) \to c \in \mathbb{R} \to \alpha \in \mathbb{R} \to \varphi_{207}$}
\BIC{$\mathrm{RFn}(f) \to c \in \mathbb{R} \to \alpha \in \mathbb{R} \to \varphi_{207}$}
\end{ptree}
\noindent$\Pi_{33}$:
\begin{ptree}
\AXC{$\stackrel{\Pi_{32}}{\vdots}$}\noLine
\UIC{$\mathrm{RFn}(f) \to c \in \mathbb{R} \to \alpha \in \mathbb{R} \to \varphi_{207}$}
\AXC{$(\mathrm{RFn}(f) \to c \in \mathbb{R} \to \alpha \in \mathbb{R} \to \varphi_{207}) \to \varphi_{283}$\lbl{A1}}
\BIC{$\varphi_{283}$}
\end{ptree}
\noindent$\Pi_{34}$:
\begin{ptree}
\AXC{$\mathrm{RFn}(f) \to c \in \mathbb{R} \to \alpha \in \mathbb{R} \to \mathrm{PNbhd}(f,c) \to \varphi_{177}$\lbl{A1}}
\AXC{$\stackrel{\Pi_{33}}{\vdots}$}\noLine
\UIC{$\varphi_{283}$}
\BIC{$\mathrm{RFn}(f) \to c \in \mathbb{R} \to \alpha \in \mathbb{R} \to \mathrm{PNbhd}(f,c) \to \varphi_{211}$}
\AXC{$\varphi_{284}$\lbl{A1}}
\BIC{$\mathrm{RFn}(f) \to c \in \mathbb{R} \to \alpha \in \mathbb{R} \to \varphi_{196}$}
\end{ptree}
\noindent$\Pi_{35}$:
\begin{ptree}
\AXC{$\forall x\,\forall y\,\forall z\,(x < y \to y \leq z \to x < z)$\lbl{[O3]}}
\AXC{$\varphi_{30}$\lbl{A2}}
\BIC{$\forall y\,\forall z\,(\lvert a_{n} - c \rvert < y \to y \leq z \to \lvert a_{n} - c \rvert < z)$}
\AXC{$\forall y\,\forall z\,(\lvert a_{n} - c \rvert < y \to y \leq z \to \lvert a_{n} - c \rvert < z) \to \varphi_{8}$\lbl{A2}}
\BIC{$\varphi_{8}$}
\AXC{$\varphi_{8} \to \varphi_{34}$\lbl{A2}}
\BIC{$\varphi_{34}$}
\end{ptree}
\noindent$\Pi_{36}$:
\begin{ptree}
\AXC{$\stackrel{\Pi_{34}}{\vdots}$}\noLine
\UIC{$\mathrm{RFn}(f) \to c \in \mathbb{R} \to \alpha \in \mathbb{R} \to \varphi_{196}$}
\AXC{$\stackrel{\Pi_{35}}{\vdots}$}\noLine
\UIC{$\varphi_{34}$}
\AXC{$\varphi_{34} \to (\mathrm{RFn}(f) \to c \in \mathbb{R} \to \alpha \in \mathbb{R} \to \varphi_{196}) \to \varphi_{271}$\lbl{A1}}
\BIC{$(\mathrm{RFn}(f) \to c \in \mathbb{R} \to \alpha \in \mathbb{R} \to \varphi_{196}) \to \varphi_{271}$}
\BIC{$\varphi_{271}$}
\end{ptree}
\noindent$\Pi_{37}$:
\begin{ptree}
\AXC{$\forall x\,\forall y\,\forall z\,(x < y \to y < z \to x < z)$\lbl{[O2]}}
\AXC{$\varphi_{36}$\lbl{A2}}
\BIC{$\forall y\,\forall z\,(\lvert a_{n} - c \rvert < y \to y < z \to \lvert a_{n} - c \rvert < z)$}
\AXC{$\forall y\,\forall z\,(\lvert a_{n} - c \rvert < y \to y < z \to \lvert a_{n} - c \rvert < z) \to \varphi_{16}$\lbl{A2}}
\BIC{$\varphi_{16}$}
\end{ptree}
\noindent$\Pi_{38}$:
\begin{ptree}
\AXC{$\stackrel{\Pi_{37}}{\vdots}$}\noLine
\UIC{$\varphi_{16}$}
\AXC{$\varphi_{16} \to \lvert a_{n} - c \rvert < \underline{K}^{-1} \to \underline{K}^{-1} < \varepsilon \to \lvert a_{n} - c \rvert < \varepsilon$\lbl{A2}}
\BIC{$\lvert a_{n} - c \rvert < \underline{K}^{-1} \to \underline{K}^{-1} < \varepsilon \to \lvert a_{n} - c \rvert < \varepsilon$}
\AXC{$\varphi_{313}$\lbl{A1}}
\BIC{$(\mathrm{RFn}(f) \to c \in \mathbb{R} \to \alpha \in \mathbb{R} \to \varphi_{197}) \to \varphi_{267}$}
\end{ptree}
\noindent$\Pi_{39}$:
\begin{ptree}
\AXC{$\stackrel{\Pi_{31}}{\vdots}$}\noLine
\UIC{$\mathrm{RFn}(f) \to c \in \mathbb{R} \to \alpha \in \mathbb{R} \to \varphi_{191}$}
\AXC{$\stackrel{\Pi_{36}}{\vdots}$}\noLine
\UIC{$\varphi_{271}$}
\BIC{$\mathrm{RFn}(f) \to c \in \mathbb{R} \to \alpha \in \mathbb{R} \to \varphi_{197}$}
\AXC{$\stackrel{\Pi_{38}}{\vdots}$}\noLine
\UIC{$(\mathrm{RFn}(f) \to c \in \mathbb{R} \to \alpha \in \mathbb{R} \to \varphi_{197}) \to \varphi_{267}$}
\BIC{$\varphi_{267}$}
\end{ptree}
\noindent$\Pi_{40}$:
\begin{ptree}
\AXC{$\varphi_{270}$\lbl{A3}}
\AXC{$\mathrm{RFn}(f) \to c \in \mathbb{R} \to \alpha \in \mathbb{R} \to \mathrm{PNbhd}(f,c) \to \varphi_{182}$\lbl{A1}}
\AXC{$\varphi_{263}$\lbl{A1}}
\BIC{$\mathrm{RFn}(f) \to c \in \mathbb{R} \to \alpha \in \mathbb{R} \to \varphi_{190}$}
\AXC{$\stackrel{\Pi_{39}}{\vdots}$}\noLine
\UIC{$\varphi_{267}$}
\BIC{$\mathrm{RFn}(f) \to c \in \mathbb{R} \to \alpha \in \mathbb{R} \to \varphi_{192}$}
\AXC{$\varphi_{268}$\lbl{A1}}
\BIC{$\varphi_{187} \to n \in \mathbb{N} \to K \leq n \to \lvert a_{n} - c \rvert < \varepsilon$}
\UIC{$\forall n\,(\varphi_{187} \to n \in \mathbb{N} \to K \leq n \to \lvert a_{n} - c \rvert < \varepsilon)$}
\AXC{$\varphi_{323}$\lbl{A1}}
\BIC{$\varphi_{270} \to (\varphi_{187} \to \forall n\,(\varphi_{187})) \to \mathrm{RFn}(f) \to c \in \mathbb{R} \to \alpha \in \mathbb{R} \to \varphi_{194}$}
\BIC{$(\varphi_{187} \to \forall n\,(\varphi_{187})) \to \mathrm{RFn}(f) \to c \in \mathbb{R} \to \alpha \in \mathbb{R} \to \varphi_{194}$}
\end{ptree}
\noindent$\Pi_{41}$:
\begin{ptree}
\AXC{$\mathrm{RFn}(f) \to c \in \mathbb{R} \to \alpha \in \mathbb{R} \to \mathrm{PNbhd}(f,c) \to \varphi_{182}$\lbl{A1}}
\AXC{$\varphi_{256}$\lbl{A1}}
\BIC{$\mathrm{RFn}(f) \to c \in \mathbb{R} \to \alpha \in \mathbb{R} \to \mathrm{PNbhd}(f,c) \to \varphi_{170}$}
\AXC{$\varphi_{308}$\lbl{A1}}
\BIC{$\varphi_{273}$}
\end{ptree}
\noindent$\Pi_{42}$:
\begin{ptree}
\AXC{$\varphi_{187} \to \forall n\,(\varphi_{187})$\lbl{A4}}
\AXC{$\stackrel{\Pi_{40}}{\vdots}$}\noLine
\UIC{$(\varphi_{187} \to \forall n\,(\varphi_{187})) \to \mathrm{RFn}(f) \to c \in \mathbb{R} \to \alpha \in \mathbb{R} \to \varphi_{194}$}
\BIC{$\mathrm{RFn}(f) \to c \in \mathbb{R} \to \alpha \in \mathbb{R} \to \varphi_{194}$}
\AXC{$\stackrel{\Pi_{41}}{\vdots}$}\noLine
\UIC{$\varphi_{273}$}
\BIC{$\mathrm{RFn}(f) \to c \in \mathbb{R} \to \alpha \in \mathbb{R} \to \varphi_{201}$}
\end{ptree}
\noindent$\Pi_{43}$:
\begin{ptree}
\AXC{$\stackrel{\Pi_{42}}{\vdots}$}\noLine
\UIC{$\mathrm{RFn}(f) \to c \in \mathbb{R} \to \alpha \in \mathbb{R} \to \varphi_{201}$}
\AXC{$(\mathrm{RFn}(f) \to c \in \mathbb{R} \to \alpha \in \mathbb{R} \to \varphi_{201}) \to \varphi_{218}$\lbl{A1}}
\BIC{$\varphi_{218}$}
\end{ptree}
\noindent$\Pi_{44}$:
\begin{ptree}
\AXC{$\varphi_{2} \to \neg (K \in \mathbb{N} \land \forall n \in \mathbb{N}\,(K \leq n \to \lvert a_{n} - c \rvert < \varepsilon))$\lbl{A2}}
\AXC{$\stackrel{\Pi_{43}}{\vdots}$}\noLine
\UIC{$\varphi_{218}$}
\BIC{$\mathrm{RFn}(f) \to c \in \mathbb{R} \to \alpha \in \mathbb{R} \to \varphi_{203}$}
\AXC{$\varphi_{276}$\lbl{A1}}
\BIC{$\varphi_{204} \to \neg (K \in \mathbb{N} \land \underline{K}^{-1} < \varepsilon)$}
\UIC{$\forall K\,(\varphi_{204} \to \neg (K \in \mathbb{N} \land \underline{K}^{-1} < \varepsilon))$}
\end{ptree}
\noindent$\Pi_{45}$:
\begin{ptree}
\AXC{$\mathrm{RFn}(f) \to c \in \mathbb{R} \to \alpha \in \mathbb{R} \to \varphi_{152}$\lbl{A1}}
\AXC{$\forall \varepsilon \in \mathbb{R}\,(0 < \varepsilon \to \exists K \in \mathbb{N}\,(\underline{K}^{-1} < \varepsilon))$\lbl{[C24]}}
\AXC{$\varphi_{56}$\lbl{A2}}
\BIC{$\varepsilon \in \mathbb{R} \to 0 < \varepsilon \to \exists K \in \mathbb{N}\,(\underline{K}^{-1} < \varepsilon)$}
\AXC{$\varphi_{304}$\lbl{A1}}
\BIC{$(\mathrm{RFn}(f) \to c \in \mathbb{R} \to \alpha \in \mathbb{R} \to \varphi_{152}) \to \varphi_{248}$}
\BIC{$\varphi_{248}$}
\end{ptree}
\noindent$\Pi_{46}$:
\begin{ptree}
\AXC{$\varphi_{277}$\lbl{A3}}
\AXC{$\stackrel{\Pi_{44}}{\vdots}$}\noLine
\UIC{$\forall K\,(\varphi_{204} \to \neg (K \in \mathbb{N} \land \underline{K}^{-1} < \varepsilon))$}
\AXC{$\mathrm{RFn}(f) \to c \in \mathbb{R} \to \alpha \in \mathbb{R} \to \varphi_{158}$\lbl{A1}}
\AXC{$\stackrel{\Pi_{45}}{\vdots}$}\noLine
\UIC{$\varphi_{248}$}
\BIC{$\mathrm{RFn}(f) \to c \in \mathbb{R} \to \alpha \in \mathbb{R} \to \mathrm{PNbhd}(f,c) \to \varphi_{164}$}
\AXC{$\varphi_{326}$\lbl{A1}}
\BIC{$\forall K\,(\varphi_{204} \to \neg (K \in \mathbb{N} \land \underline{K}^{-1} < \varepsilon)) \to \varphi_{277} \to \varphi_{307}$}
\BIC{$\varphi_{277} \to \varphi_{307}$}
\BIC{$\varphi_{307}$}
\end{ptree}
\noindent$\Pi_{47}$:
\begin{ptree}
\AXC{$\varphi_{264}$\lbl{A3}}
\AXC{$\varphi_{204} \to \forall K\,(\varphi_{204})$\lbl{A4}}
\AXC{$\stackrel{\Pi_{46}}{\vdots}$}\noLine
\UIC{$\varphi_{307}$}
\BIC{$\mathrm{RFn}(f) \to c \in \mathbb{R} \to \alpha \in \mathbb{R} \to \mathrm{PNbhd}(f,c) \to \varphi_{180}$}
\AXC{$\varphi_{259}$\lbl{A1}}
\BIC{$\varphi_{106} \land (\mathrm{Seq}(a) \land \varphi_{61}) \to \varphi_{20}$}
\UIC{$\forall \varepsilon\,(\varphi_{106} \land (\mathrm{Seq}(a) \land \varphi_{61}) \to \varphi_{20})$}
\AXC{$\forall \varepsilon\,(\varphi_{106} \land (\mathrm{Seq}(a) \land \varphi_{61}) \to \varphi_{20}) \to \varphi_{264} \to \varphi_{298}$\lbl{A1}}
\BIC{$\varphi_{264} \to \varphi_{298}$}
\BIC{$\varphi_{298}$}
\end{ptree}
\noindent$\Pi_{48}$:
\begin{ptree}
\AXC{$\forall a\,(\varphi_{60})$\lbl{DefConv}}
\AXC{$\forall a\,(\varphi_{60}) \to \varphi_{60}$\lbl{A2}}
\BIC{$\varphi_{60}$}
\AXC{$\varphi_{60} \to (\mathrm{Conv}(a,c) \leftrightarrow \mathrm{Seq}(a) \land c \in \mathbb{R} \land \varphi_{27})$\lbl{A2}}
\BIC{$\mathrm{Conv}(a,c) \leftrightarrow \mathrm{Seq}(a) \land c \in \mathbb{R} \land \varphi_{27}$}
\end{ptree}
\noindent$\Pi_{49}$:
\begin{ptree}
\AXC{$\stackrel{\Pi_{48}}{\vdots}$}\noLine
\UIC{$\mathrm{Conv}(a,c) \leftrightarrow \mathrm{Seq}(a) \land c \in \mathbb{R} \land \varphi_{27}$}
\AXC{$(\mathrm{Conv}(a,c) \leftrightarrow \mathrm{Seq}(a) \land c \in \mathbb{R} \land \varphi_{27}) \to \varphi_{303}$\lbl{A1}}
\BIC{$\varphi_{303}$}
\end{ptree}
\noindent$\Pi_{50}$:
\begin{ptree}
\AXC{$\mathrm{RFn}(f) \to c \in \mathbb{R} \to c \in \mathbb{R}$\lbl{A1}}
\AXC{$\mathrm{RFn}(f) \to c \in \mathbb{R} \to \alpha \in \mathbb{R} \to \varphi_{209}$\lbl{A1}}
\AXC{$\varphi_{253}$\lbl{A1}}
\BIC{$\mathrm{RFn}(f) \to c \in \mathbb{R} \to \alpha \in \mathbb{R} \to \varphi_{139}$}
\AXC{$\stackrel{\Pi_{49}}{\vdots}$}\noLine
\UIC{$\varphi_{303}$}
\BIC{$(\mathrm{RFn}(f) \to c \in \mathbb{R} \to c \in \mathbb{R}) \to \varphi_{249}$}
\BIC{$\varphi_{249}$}
\end{ptree}
\noindent$\Pi_{51}$:
\begin{ptree}
\AXC{$\varphi_{106} \land (\mathrm{Seq}(a) \land \varphi_{61}) \to \forall \varepsilon\,(\varphi_{106} \land (\mathrm{Seq}(a) \land \varphi_{61}))$\lbl{A4}}
\AXC{$\stackrel{\Pi_{47}}{\vdots}$}\noLine
\UIC{$\varphi_{298}$}
\BIC{$\mathrm{RFn}(f) \to c \in \mathbb{R} \to \alpha \in \mathbb{R} \to \mathrm{PNbhd}(f,c) \to \varphi_{183}$}
\AXC{$\stackrel{\Pi_{50}}{\vdots}$}\noLine
\UIC{$\varphi_{249}$}
\BIC{$\mathrm{RFn}(f) \to c \in \mathbb{R} \to \alpha \in \mathbb{R} \to \varphi_{142}$}
\end{ptree}
\noindent$\Pi_{52}$:
\begin{ptree}
\AXC{$\varphi_{69}$\lbl{A2}}
\AXC{$\mathrm{RFn}(f) \to c \in \mathbb{R} \to \alpha \in \mathbb{R} \to \mathrm{PNbhd}(f,c) \to \varphi_{70}$\lbl{A1}}
\AXC{$\varphi_{208}$\lbl{A1}}
\BIC{$\varphi_{108}$}
\BIC{$\mathrm{RFn}(f) \to c \in \mathbb{R} \to \alpha \in \mathbb{R} \to \mathrm{PNbhd}(f,c) \to \varphi_{69}$}
\AXC{$\varphi_{319}$\lbl{A1}}
\BIC{$\varphi_{316}$}
\end{ptree}
\noindent$\Pi_{53}$:
\begin{ptree}
\AXC{$\stackrel{\Pi_{51}}{\vdots}$}\noLine
\UIC{$\mathrm{RFn}(f) \to c \in \mathbb{R} \to \alpha \in \mathbb{R} \to \varphi_{142}$}
\AXC{$\mathrm{RFn}(f) \to c \in \mathbb{R} \to \alpha \in \mathbb{R} \to \varphi_{209}$\lbl{A1}}
\AXC{$\varphi_{253}$\lbl{A1}}
\BIC{$\mathrm{RFn}(f) \to c \in \mathbb{R} \to \alpha \in \mathbb{R} \to \varphi_{139}$}
\AXC{$\stackrel{\Pi_{52}}{\vdots}$}\noLine
\UIC{$\varphi_{316}$}
\BIC{$(\mathrm{RFn}(f) \to c \in \mathbb{R} \to \alpha \in \mathbb{R} \to \varphi_{142}) \to \varphi_{247}$}
\BIC{$\varphi_{247}$}
\end{ptree}
\noindent$\Pi_{54}$:
\begin{ptree}
\AXC{$\varphi_{47}$\lbl{[C3]}}
\AXC{$\varphi_{47} \to \forall \alpha\,(\varphi_{57})$\lbl{A2}}
\BIC{$\forall \alpha\,(\varphi_{57})$}
\AXC{$\forall \alpha\,(\varphi_{57}) \to \varphi_{57}$\lbl{A2}}
\BIC{$\varphi_{57}$}
\AXC{$\varphi_{57} \to \mathrm{Conv}(f \circ a,\alpha) \to e_{0} \in \mathbb{R} \to \varphi_{15}$\lbl{A2}}
\BIC{$\mathrm{Conv}(f \circ a,\alpha) \to e_{0} \in \mathbb{R} \to \varphi_{15}$}
\AXC{$\varphi_{315}$\lbl{A1}}
\BIC{$(\mathrm{RFn}(f) \to c \in \mathbb{R} \to \alpha \in \mathbb{R} \to \varphi_{148}) \to \varphi_{290}$}
\end{ptree}
\noindent$\Pi_{55}$:
\begin{ptree}
\AXC{$\stackrel{\Pi_{19}}{\vdots}$}\noLine
\UIC{$\mathrm{RFn}(f) \to c \in \mathbb{R} \to \alpha \in \mathbb{R} \to \mathrm{PNbhd}(f,c) \to \varphi_{160}$}
\AXC{$\stackrel{\Pi_{53}}{\vdots}$}\noLine
\UIC{$\varphi_{247}$}
\BIC{$\mathrm{RFn}(f) \to c \in \mathbb{R} \to \alpha \in \mathbb{R} \to \varphi_{148}$}
\AXC{$\stackrel{\Pi_{54}}{\vdots}$}\noLine
\UIC{$(\mathrm{RFn}(f) \to c \in \mathbb{R} \to \alpha \in \mathbb{R} \to \varphi_{148}) \to \varphi_{290}$}
\BIC{$\varphi_{290}$}
\end{ptree}
\noindent$\Pi_{56}$:
\begin{ptree}
\AXC{$\mathrm{RFn}(f) \to c \in \mathbb{R} \to \alpha \in \mathbb{R} \to \mathrm{PNbhd}(f,c) \to \varphi_{129}$\lbl{A1}}
\AXC{$\varphi_{230}$\lbl{A1}}
\BIC{$\mathrm{RFn}(f) \to c \in \mathbb{R} \to \alpha \in \mathbb{R} \to \mathrm{PNbhd}(f,c) \to \varphi_{95}$}
\AXC{$\stackrel{\Pi_{55}}{\vdots}$}\noLine
\UIC{$\varphi_{290}$}
\BIC{$\varphi_{231}$}
\end{ptree}
\noindent$\Pi_{57}$:
\begin{ptree}
\AXC{$\mathrm{RFn}(f) \to c \in \mathbb{R} \to \alpha \in \mathbb{R} \to \mathrm{PNbhd}(f,c) \to \varphi_{129}$\lbl{A1}}
\AXC{$\varphi_{229}$\lbl{A1}}
\BIC{$\mathrm{RFn}(f) \to c \in \mathbb{R} \to \alpha \in \mathbb{R} \to \mathrm{PNbhd}(f,c) \to \varphi_{94}$}
\AXC{$\stackrel{\Pi_{56}}{\vdots}$}\noLine
\UIC{$\varphi_{231}$}
\BIC{$\mathrm{RFn}(f) \to c \in \mathbb{R} \to \alpha \in \mathbb{R} \to \mathrm{PNbhd}(f,c) \to \varphi_{172}$}
\AXC{$\varphi_{325}$\lbl{A1}}
\BIC{$\forall N\,(\varphi_{106} \land (\mathrm{Seq}(a) \land \varphi_{61}) \land \neg 0 < 0 \to \varphi_{3}) \to \varphi_{258} \to \varphi_{294}$}
\end{ptree}
\noindent$\Pi_{58}$:
\begin{ptree}
\AXC{$\stackrel{\Pi_{27}}{\vdots}$}\noLine
\UIC{$\mathrm{RFn}(f) \to c \in \mathbb{R} \to \alpha \in \mathbb{R} \to \mathrm{PNbhd}(f,c) \to \varphi_{174}$}
\AXC{$\varphi_{255}$\lbl{A1}}
\BIC{$\varphi_{106} \land (\mathrm{Seq}(a) \land \varphi_{61}) \land \neg 0 < 0 \to \varphi_{3}$}
\UIC{$\forall N\,(\varphi_{106} \land (\mathrm{Seq}(a) \land \varphi_{61}) \land \neg 0 < 0 \to \varphi_{3})$}
\AXC{$\stackrel{\Pi_{57}}{\vdots}$}\noLine
\UIC{$\forall N\,(\varphi_{106} \land (\mathrm{Seq}(a) \land \varphi_{61}) \land \neg 0 < 0 \to \varphi_{3}) \to \varphi_{258} \to \varphi_{294}$}
\BIC{$\varphi_{258} \to \varphi_{294}$}
\end{ptree}
\noindent$\Pi_{59}$:
\begin{ptree}
\AXC{$\mathrm{RFn}(f) \to \mathrm{RFn}(f)$\lbl{A1}}
\AXC{$\forall f\,\forall x\,(\mathrm{RFn}(f) \to x \in \mathrm{dom}(f) \to x \in \mathbb{R})$\lbl{[S7]}}
\AXC{$\varphi_{9}$\lbl{A2}}
\BIC{$\forall x\,(\mathrm{RFn}(f) \to x \in \mathrm{dom}(f) \to x \in \mathbb{R})$}
\AXC{$\varphi_{5}$\lbl{A2}}
\BIC{$\mathrm{RFn}(f) \to x \in \mathrm{dom}(f) \to x \in \mathbb{R}$}
\AXC{$\varphi_{236}$\lbl{A1}}
\BIC{$(\mathrm{RFn}(f) \to \mathrm{RFn}(f)) \to \varphi_{232}$}
\BIC{$\varphi_{232}$}
\end{ptree}
\noindent$\Pi_{60}$:
\begin{ptree}
\AXC{$\mathrm{RFn}(f) \to c \in \mathbb{R} \to \alpha \in \mathbb{R} \to \mathrm{PNbhd}(f,c) \to \varphi_{161}$\lbl{A1}}
\AXC{$\varphi_{239}$\lbl{A1}}
\BIC{$\mathrm{RFn}(f) \to c \in \mathbb{R} \to \alpha \in \mathbb{R} \to \varphi_{128}$}
\AXC{$\stackrel{\Pi_{59}}{\vdots}$}\noLine
\UIC{$\varphi_{232}$}
\BIC{$\mathrm{RFn}(f) \to c \in \mathbb{R} \to \alpha \in \mathbb{R} \to \mathrm{PNbhd}(f,c) \to \varphi_{125}$}
\end{ptree}
\noindent$\Pi_{61}$:
\begin{ptree}
\AXC{$\stackrel{\Pi_{60}}{\vdots}$}\noLine
\UIC{$\mathrm{RFn}(f) \to c \in \mathbb{R} \to \alpha \in \mathbb{R} \to \mathrm{PNbhd}(f,c) \to \varphi_{125}$}
\AXC{$\mathrm{RFn}(f) \to c \in \mathbb{R} \to \alpha \in \mathbb{R} \to \mathrm{PNbhd}(f,c) \to \varphi_{161}$\lbl{A1}}
\AXC{$\varphi_{297}$\lbl{A1}}
\BIC{$\varphi_{241}$}
\BIC{$\mathrm{RFn}(f) \to c \in \mathbb{R} \to \alpha \in \mathbb{R} \to \mathrm{PNbhd}(f,c) \to \varphi_{166}$}
\end{ptree}
\noindent$\Pi_{62}$:
\begin{ptree}
\AXC{$\forall x\,(\neg \varphi_{37}) \to \neg \varphi_{37}$\lbl{A2}}
\AXC{$\stackrel{\Pi_{61}}{\vdots}$}\noLine
\UIC{$\mathrm{RFn}(f) \to c \in \mathbb{R} \to \alpha \in \mathbb{R} \to \mathrm{PNbhd}(f,c) \to \varphi_{166}$}
\AXC{$\varphi_{288}$\lbl{A1}}
\BIC{$\varphi_{219}$}
\BIC{$\mathrm{RFn}(f) \to c \in \mathbb{R} \to \alpha \in \mathbb{R} \to \mathrm{PNbhd}(f,c) \to \varphi_{168}$}
\AXC{$\varphi_{251}$\lbl{A1}}
\BIC{$\varphi_{106} \land k \in \mathbb{N} \land \neg \varphi_{39} \to \neg \neg \varphi_{24}$}
\UIC{$\forall x\,(\varphi_{106} \land k \in \mathbb{N} \land \neg \varphi_{39} \to \neg \neg \varphi_{24})$}
\end{ptree}
\noindent$\Pi_{63}$:
\begin{ptree}
\AXC{$\forall x\,(\neg \neg \varphi_{24}) \to \neg \neg \varphi_{24}$\lbl{A2}}
\AXC{$(\forall x\,(\neg \neg \varphi_{24}) \to \neg \neg \varphi_{24}) \to \forall x\,(\neg \neg \varphi_{24}) \to \varphi_{24}$\lbl{A1}}
\BIC{$\forall x\,(\neg \neg \varphi_{24}) \to \varphi_{24}$}
\UIC{$\forall x\,(\forall x\,(\neg \neg \varphi_{24}) \to \varphi_{24})$}
\end{ptree}
\noindent$\Pi_{64}$:
\begin{ptree}
\AXC{$\forall n \in \mathbb{N}\,(0 < \underline{n})$\lbl{[N15]}}
\AXC{$\forall n \in \mathbb{N}\,(0 < \underline{n}) \to k \in \mathbb{N} \to 0 < \underline{k}$\lbl{A2}}
\BIC{$k \in \mathbb{N} \to 0 < \underline{k}$}
\AXC{$(k \in \mathbb{N} \to 0 < \underline{k}) \to \varphi_{222}$\lbl{A1}}
\BIC{$\varphi_{222}$}
\end{ptree}
\noindent$\Pi_{65}$:
\begin{ptree}
\AXC{$\mathrm{RFn}(f) \to c \in \mathbb{R} \to \alpha \in \mathbb{R} \to \mathrm{PNbhd}(f,c) \to \varphi_{98}$\lbl{A1}}
\AXC{$\stackrel{\Pi_{64}}{\vdots}$}\noLine
\UIC{$\varphi_{222}$}
\BIC{$\mathrm{RFn}(f) \to c \in \mathbb{R} \to \alpha \in \mathbb{R} \to \mathrm{PNbhd}(f,c) \to \varphi_{101}$}
\end{ptree}
\noindent$\Pi_{66}$:
\begin{ptree}
\AXC{$\forall n \in \mathbb{N}\,(\underline{n} \in \mathbb{R})$\lbl{[N4]}}
\AXC{$\forall n \in \mathbb{N}\,(\underline{n} \in \mathbb{R}) \to k \in \mathbb{N} \to \underline{k} \in \mathbb{R}$\lbl{A2}}
\BIC{$k \in \mathbb{N} \to \underline{k} \in \mathbb{R}$}
\AXC{$(k \in \mathbb{N} \to \underline{k} \in \mathbb{R}) \to \varphi_{224}$\lbl{A1}}
\BIC{$\varphi_{224}$}
\end{ptree}
\noindent$\Pi_{67}$:
\begin{ptree}
\AXC{$\mathrm{RFn}(f) \to c \in \mathbb{R} \to \alpha \in \mathbb{R} \to \mathrm{PNbhd}(f,c) \to \varphi_{98}$\lbl{A1}}
\AXC{$\stackrel{\Pi_{66}}{\vdots}$}\noLine
\UIC{$\varphi_{224}$}
\BIC{$\mathrm{RFn}(f) \to c \in \mathbb{R} \to \alpha \in \mathbb{R} \to \mathrm{PNbhd}(f,c) \to \varphi_{102}$}
\end{ptree}
\noindent$\Pi_{68}$:
\begin{ptree}
\AXC{$\forall x \in \mathbb{R}\,(0 < x \to 0 < x^{-1})$\lbl{[F34]}}
\AXC{$\varphi_{25}$\lbl{A2}}
\BIC{$\underline{k} \in \mathbb{R} \to 0 < \underline{k} \to 0 < \underline{k}^{-1}$}
\AXC{$(\underline{k} \in \mathbb{R} \to 0 < \underline{k} \to 0 < \underline{k}^{-1}) \to \varphi_{274}$\lbl{A1}}
\BIC{$\varphi_{274}$}
\end{ptree}
\noindent$\Pi_{69}$:
\begin{ptree}
\AXC{$\stackrel{\Pi_{65}}{\vdots}$}\noLine
\UIC{$\mathrm{RFn}(f) \to c \in \mathbb{R} \to \alpha \in \mathbb{R} \to \mathrm{PNbhd}(f,c) \to \varphi_{101}$}
\AXC{$\stackrel{\Pi_{67}}{\vdots}$}\noLine
\UIC{$\mathrm{RFn}(f) \to c \in \mathbb{R} \to \alpha \in \mathbb{R} \to \mathrm{PNbhd}(f,c) \to \varphi_{102}$}
\AXC{$\stackrel{\Pi_{68}}{\vdots}$}\noLine
\UIC{$\varphi_{274}$}
\BIC{$\varphi_{225}$}
\BIC{$\mathrm{RFn}(f) \to c \in \mathbb{R} \to \alpha \in \mathbb{R} \to \mathrm{PNbhd}(f,c) \to \varphi_{103}$}
\end{ptree}
\noindent$\Pi_{70}$:
\begin{ptree}
\AXC{$x = x$\lbl{A5}}
\AXC{$x = y \to x = x \to y = x$\lbl{A6}}
\AXC{$(x = y \to x = x \to y = x) \to x = x \to x = y \to y = x$\lbl{A1}}
\BIC{$x = x \to x = y \to y = x$}
\BIC{$x = y \to y = x$}
\UIC{$\forall y\,(x = y \to y = x)$}
\UIC{$\forall x\,\forall y\,(x = y \to y = x)$}
\end{ptree}
\noindent$\Pi_{71}$:
\begin{ptree}
\AXC{$\stackrel{\Pi_{70}}{\vdots}$}\noLine
\UIC{$\forall x\,\forall y\,(x = y \to y = x)$}
\AXC{$\forall x\,\forall y\,(x = y \to y = x) \to \forall y\,(\underline{k} = y \to y = \underline{k})$\lbl{A2}}
\BIC{$\forall y\,(\underline{k} = y \to y = \underline{k})$}
\end{ptree}
\noindent$\Pi_{72}$:
\begin{ptree}
\AXC{$\forall x\,\forall y\,(x < y \to x \neq y)$\lbl{[F26]}}
\AXC{$\forall x\,\forall y\,(x < y \to x \neq y) \to \forall y\,(0 < y \to 0 \neq y)$\lbl{A2}}
\BIC{$\forall y\,(0 < y \to 0 \neq y)$}
\AXC{$\forall y\,(0 < y \to 0 \neq y) \to 0 < \underline{k} \to 0 \neq \underline{k}$\lbl{A2}}
\BIC{$0 < \underline{k} \to 0 \neq \underline{k}$}
\end{ptree}
\noindent$\Pi_{73}$:
\begin{ptree}
\AXC{$\stackrel{\Pi_{65}}{\vdots}$}\noLine
\UIC{$\mathrm{RFn}(f) \to c \in \mathbb{R} \to \alpha \in \mathbb{R} \to \mathrm{PNbhd}(f,c) \to \varphi_{101}$}
\AXC{$\stackrel{\Pi_{72}}{\vdots}$}\noLine
\UIC{$0 < \underline{k} \to 0 \neq \underline{k}$}
\AXC{$(0 < \underline{k} \to 0 \neq \underline{k}) \to \varphi_{223}$\lbl{A1}}
\BIC{$\varphi_{223}$}
\BIC{$\mathrm{RFn}(f) \to c \in \mathbb{R} \to \alpha \in \mathbb{R} \to \mathrm{PNbhd}(f,c) \to \varphi_{99}$}
\AXC{$\varphi_{227}$\lbl{A1}}
\BIC{$\varphi_{121}$}
\end{ptree}
\noindent$\Pi_{74}$:
\begin{ptree}
\AXC{$\stackrel{\Pi_{71}}{\vdots}$}\noLine
\UIC{$\forall y\,(\underline{k} = y \to y = \underline{k})$}
\AXC{$\forall y\,(\underline{k} = y \to y = \underline{k}) \to \underline{k} = 0 \to 0 = \underline{k}$\lbl{A2}}
\BIC{$\underline{k} = 0 \to 0 = \underline{k}$}
\AXC{$\stackrel{\Pi_{73}}{\vdots}$}\noLine
\UIC{$\varphi_{121}$}
\BIC{$\mathrm{RFn}(f) \to c \in \mathbb{R} \to \alpha \in \mathbb{R} \to \mathrm{PNbhd}(f,c) \to \varphi_{100}$}
\end{ptree}
\noindent$\Pi_{75}$:
\begin{ptree}
\AXC{$\forall x \in \mathbb{R}\,(x \neq 0 \to x^{-1} \in \mathbb{R})$\lbl{(R1)}}
\AXC{$\varphi_{31}$\lbl{A2}}
\BIC{$\underline{k} \in \mathbb{R} \to \underline{k} \neq 0 \to \underline{k}^{-1} \in \mathbb{R}$}
\AXC{$(\underline{k} \in \mathbb{R} \to \underline{k} \neq 0 \to \underline{k}^{-1} \in \mathbb{R}) \to \varphi_{275}$\lbl{A1}}
\BIC{$\varphi_{275}$}
\end{ptree}
\noindent$\Pi_{76}$:
\begin{ptree}
\AXC{$\stackrel{\Pi_{74}}{\vdots}$}\noLine
\UIC{$\mathrm{RFn}(f) \to c \in \mathbb{R} \to \alpha \in \mathbb{R} \to \mathrm{PNbhd}(f,c) \to \varphi_{100}$}
\AXC{$\stackrel{\Pi_{67}}{\vdots}$}\noLine
\UIC{$\mathrm{RFn}(f) \to c \in \mathbb{R} \to \alpha \in \mathbb{R} \to \mathrm{PNbhd}(f,c) \to \varphi_{102}$}
\AXC{$\stackrel{\Pi_{75}}{\vdots}$}\noLine
\UIC{$\varphi_{275}$}
\BIC{$\varphi_{226}$}
\BIC{$\mathrm{RFn}(f) \to c \in \mathbb{R} \to \alpha \in \mathbb{R} \to \mathrm{PNbhd}(f,c) \to \varphi_{104}$}
\end{ptree}
\noindent$\Pi_{77}$:
\begin{ptree}
\AXC{$\varphi_{44} \to \neg (\underline{k}^{-1} \in \mathbb{R} \land 0 < \underline{k}^{-1} \land \varphi_{29})$\lbl{A2}}
\AXC{$\mathrm{RFn}(f) \to c \in \mathbb{R} \to \alpha \in \mathbb{R} \to \mathrm{PNbhd}(f,c) \to \varphi_{129}$\lbl{A1}}
\AXC{$\varphi_{234}$\lbl{A1}}
\BIC{$\mathrm{RFn}(f) \to c \in \mathbb{R} \to \alpha \in \mathbb{R} \to \mathrm{PNbhd}(f,c) \to \varphi_{124}$}
\AXC{$\varphi_{269}$\lbl{A1}}
\BIC{$\varphi_{217}$}
\BIC{$\mathrm{RFn}(f) \to c \in \mathbb{R} \to \alpha \in \mathbb{R} \to \varphi_{134}$}
\end{ptree}
\noindent$\Pi_{78}$:
\begin{ptree}
\AXC{$\stackrel{\Pi_{77}}{\vdots}$}\noLine
\UIC{$\mathrm{RFn}(f) \to c \in \mathbb{R} \to \alpha \in \mathbb{R} \to \varphi_{134}$}
\AXC{$(\mathrm{RFn}(f) \to c \in \mathbb{R} \to \alpha \in \mathbb{R} \to \varphi_{134}) \to \varphi_{285}$\lbl{A1}}
\BIC{$\varphi_{285}$}
\end{ptree}
\noindent$\Pi_{79}$:
\begin{ptree}
\AXC{$\stackrel{\Pi_{69}}{\vdots}$}\noLine
\UIC{$\mathrm{RFn}(f) \to c \in \mathbb{R} \to \alpha \in \mathbb{R} \to \mathrm{PNbhd}(f,c) \to \varphi_{103}$}
\AXC{$\stackrel{\Pi_{76}}{\vdots}$}\noLine
\UIC{$\mathrm{RFn}(f) \to c \in \mathbb{R} \to \alpha \in \mathbb{R} \to \mathrm{PNbhd}(f,c) \to \varphi_{104}$}
\AXC{$\stackrel{\Pi_{78}}{\vdots}$}\noLine
\UIC{$\varphi_{285}$}
\BIC{$\varphi_{228}$}
\BIC{$\mathrm{RFn}(f) \to c \in \mathbb{R} \to \alpha \in \mathbb{R} \to \mathrm{PNbhd}(f,c) \to \varphi_{122}$}
\end{ptree}
\noindent$\Pi_{80}$:
\begin{ptree}
\AXC{$\stackrel{\Pi_{63}}{\vdots}$}\noLine
\UIC{$\forall x\,(\forall x\,(\neg \neg \varphi_{24}) \to \varphi_{24})$}
\AXC{$\stackrel{\Pi_{79}}{\vdots}$}\noLine
\UIC{$\mathrm{RFn}(f) \to c \in \mathbb{R} \to \alpha \in \mathbb{R} \to \mathrm{PNbhd}(f,c) \to \varphi_{122}$}
\AXC{$\varphi_{300}$\lbl{A1}}
\BIC{$\varphi_{254}$}
\BIC{$(\forall x\,(\forall x\,(\neg \neg \varphi_{24}) \to \varphi_{24}) \to \forall x\,\forall x\,(\neg \neg \varphi_{24}) \to \varphi_{29}) \to \varphi_{212}$}
\end{ptree}
\noindent$\Pi_{81}$:
\begin{ptree}
\AXC{$\forall x\,(\forall x\,(\neg \neg \varphi_{24}) \to \varphi_{24}) \to \forall x\,\forall x\,(\neg \neg \varphi_{24}) \to \varphi_{29}$\lbl{A3}}
\AXC{$\stackrel{\Pi_{80}}{\vdots}$}\noLine
\UIC{$(\forall x\,(\forall x\,(\neg \neg \varphi_{24}) \to \varphi_{24}) \to \forall x\,\forall x\,(\neg \neg \varphi_{24}) \to \varphi_{29}) \to \varphi_{212}$}
\BIC{$\varphi_{212}$}
\end{ptree}
\noindent$\Pi_{82}$:
\begin{ptree}
\AXC{$\varphi_{252}$\lbl{A3}}
\AXC{$\stackrel{\Pi_{62}}{\vdots}$}\noLine
\UIC{$\forall x\,(\varphi_{106} \land k \in \mathbb{N} \land \neg \varphi_{39} \to \neg \neg \varphi_{24})$}
\AXC{$\forall x\,(\neg \neg \varphi_{24}) \to \forall x\,\forall x\,(\neg \neg \varphi_{24})$\lbl{A4}}
\AXC{$\stackrel{\Pi_{81}}{\vdots}$}\noLine
\UIC{$\varphi_{212}$}
\BIC{$\mathrm{RFn}(f) \to c \in \mathbb{R} \to \alpha \in \mathbb{R} \to \mathrm{PNbhd}(f,c) \to \varphi_{123}$}
\AXC{$\varphi_{324}$\lbl{A1}}
\BIC{$\forall x\,(\varphi_{106} \land k \in \mathbb{N} \land \neg \varphi_{39} \to \neg \neg \varphi_{24}) \to \varphi_{252} \to \varphi_{291}$}
\BIC{$\varphi_{252} \to \varphi_{291}$}
\BIC{$\varphi_{291}$}
\end{ptree}
\noindent$\Pi_{83}$:
\begin{ptree}
\AXC{$\varphi_{106} \land k \in \mathbb{N} \land \neg \varphi_{39} \to \forall x\,(\varphi_{106} \land k \in \mathbb{N} \land \neg \varphi_{39})$\lbl{A4}}
\AXC{$\stackrel{\Pi_{82}}{\vdots}$}\noLine
\UIC{$\varphi_{291}$}
\BIC{$\mathrm{RFn}(f) \to c \in \mathbb{R} \to \alpha \in \mathbb{R} \to \mathrm{PNbhd}(f,c) \to \varphi_{126}$}
\AXC{$\varphi_{233}$\lbl{A1}}
\BIC{$\varphi_{106} \to k \in \mathbb{N} \to \varphi_{39}$}
\UIC{$\forall k\,(\varphi_{106} \to k \in \mathbb{N} \to \varphi_{39})$}
\AXC{$\forall k\,(\varphi_{106} \to k \in \mathbb{N} \to \varphi_{39}) \to \varphi_{318}$\lbl{A1}}
\BIC{$\varphi_{318}$}
\end{ptree}
\noindent$\Pi_{84}$:
\begin{ptree}
\AXC{$\forall f\,\forall c\,\forall \alpha\,\forall e_{0}\,(\forall k \in \mathbb{N}\,(\varphi_{39}) \to \varphi_{65})$\lbl{[CC1]}}
\AXC{$\varphi_{151}$\lbl{A2}}
\BIC{$\forall c\,\forall \alpha\,\forall e_{0}\,(\forall k \in \mathbb{N}\,(\varphi_{39}) \to \varphi_{65})$}
\AXC{$\varphi_{146}$\lbl{A2}}
\BIC{$\forall \alpha\,\forall e_{0}\,(\forall k \in \mathbb{N}\,(\varphi_{39}) \to \varphi_{65})$}
\AXC{$\varphi_{143}$\lbl{A2}}
\BIC{$\forall e_{0}\,(\forall k \in \mathbb{N}\,(\varphi_{39}) \to \varphi_{65})$}
\end{ptree}
\noindent$\Pi_{85}$:
\begin{ptree}
\AXC{$\stackrel{\Pi_{84}}{\vdots}$}\noLine
\UIC{$\forall e_{0}\,(\forall k \in \mathbb{N}\,(\varphi_{39}) \to \varphi_{65})$}
\AXC{$\forall e_{0}\,(\forall k \in \mathbb{N}\,(\varphi_{39}) \to \varphi_{65}) \to \forall k \in \mathbb{N}\,(\varphi_{39}) \to \varphi_{65}$\lbl{A2}}
\BIC{$\forall k \in \mathbb{N}\,(\varphi_{39}) \to \varphi_{65}$}
\AXC{$(\forall k \in \mathbb{N}\,(\varphi_{39}) \to \varphi_{65}) \to \varphi_{235}$\lbl{A1}}
\BIC{$\varphi_{235}$}
\end{ptree}
\noindent$\Pi_{86}$:
\begin{ptree}
\AXC{$\varphi_{106} \to \forall k\,(\varphi_{106})$\lbl{A4}}
\AXC{$\forall k\,(\varphi_{106} \to k \in \mathbb{N} \to \varphi_{39}) \to \forall k\,(\varphi_{106}) \to \forall k \in \mathbb{N}\,(\varphi_{39})$\lbl{A3}}
\AXC{$\stackrel{\Pi_{83}}{\vdots}$}\noLine
\UIC{$\varphi_{318}$}
\BIC{$\varphi_{278}$}
\BIC{$\mathrm{RFn}(f) \to c \in \mathbb{R} \to \alpha \in \mathbb{R} \to \mathrm{PNbhd}(f,c) \to \varphi_{127}$}
\AXC{$\stackrel{\Pi_{85}}{\vdots}$}\noLine
\UIC{$\varphi_{235}$}
\BIC{$\mathrm{RFn}(f) \to c \in \mathbb{R} \to \alpha \in \mathbb{R} \to \mathrm{PNbhd}(f,c) \to \varphi_{131}$}
\AXC{$\varphi_{327}$\lbl{A1}}
\BIC{$\forall a\,(\varphi_{106} \land \neg 0 < 0 \to \neg (\mathrm{Seq}(a) \land \varphi_{61})) \to \varphi_{240} \to \varphi_{317}$}
\end{ptree}
\noindent$\Pi_{87}$:
\begin{ptree}
\AXC{$\varphi_{240}$\lbl{A3}}
\AXC{$\varphi_{238}$\lbl{A4}}
\AXC{$\varphi_{258}$\lbl{A3}}
\AXC{$\stackrel{\Pi_{58}}{\vdots}$}\noLine
\UIC{$\varphi_{258} \to \varphi_{294}$}
\BIC{$\varphi_{294}$}
\BIC{$\mathrm{RFn}(f) \to c \in \mathbb{R} \to \alpha \in \mathbb{R} \to \mathrm{PNbhd}(f,c) \to \varphi_{135}$}
\AXC{$\varphi_{237}$\lbl{A1}}
\BIC{$\varphi_{106} \land \neg 0 < 0 \to \neg (\mathrm{Seq}(a) \land \varphi_{61})$}
\UIC{$\forall a\,(\varphi_{106} \land \neg 0 < 0 \to \neg (\mathrm{Seq}(a) \land \varphi_{61}))$}
\AXC{$\stackrel{\Pi_{86}}{\vdots}$}\noLine
\UIC{$\forall a\,(\varphi_{106} \land \neg 0 < 0 \to \neg (\mathrm{Seq}(a) \land \varphi_{61})) \to \varphi_{240} \to \varphi_{317}$}
\BIC{$\varphi_{240} \to \varphi_{317}$}
\BIC{$\varphi_{317}$}
\end{ptree}
\noindent$\Pi_{88}$:
\begin{ptree}
\AXC{$\varphi_{88}$\lbl{A4}}
\AXC{$\varphi_{169}$\lbl{A3}}
\AXC{$\forall a\,(\neg (\mathrm{Seq}(a) \land \varphi_{61})) \to \varphi_{64}$\lbl{A2}}
\UIC{$\forall q\,(\forall a\,(\neg (\mathrm{Seq}(a) \land \varphi_{61})) \to \varphi_{64})$}
\AXC{$\varphi_{106} \land \neg 0 < 0 \to \forall a\,(\varphi_{106} \land \neg 0 < 0)$\lbl{A4}}
\AXC{$\stackrel{\Pi_{87}}{\vdots}$}\noLine
\UIC{$\varphi_{317}$}
\BIC{$\forall q\,(\forall a\,(\neg (\mathrm{Seq}(a) \land \varphi_{61})) \to \varphi_{64}) \to \varphi_{169} \to \varphi_{213}$}
\BIC{$\varphi_{169} \to \varphi_{213}$}
\BIC{$\varphi_{213}$}
\BIC{$\mathrm{RFn}(f) \to c \in \mathbb{R} \to \alpha \in \mathbb{R} \to \mathrm{PNbhd}(f,c) \to \varphi_{93}$}
\end{ptree}
\noindent$\Pi_{89}$:
\begin{ptree}
\AXC{$\varphi_{83}$\lbl{A2}}
\AXC{$\varphi_{156}$\lbl{A1}}
\BIC{$\forall \varepsilon\,(\neg \neg (\varepsilon \in \mathbb{R} \to 0 < \varepsilon \to \varphi_{45})) \to \varepsilon \in \mathbb{R} \to 0 < \varepsilon \to \varphi_{45}$}
\UIC{$\varphi_{84}$}
\AXC{$\mathrm{RFn}(f) \to c \in \mathbb{R} \to \alpha \in \mathbb{R} \to \mathrm{PNbhd}(f,c) \to \varphi_{92}$\lbl{A1}}
\AXC{$\varphi_{310}$\lbl{A1}}
\BIC{$\varphi_{84} \to (\varphi_{84} \to \varphi_{86}) \to \varphi_{87} \to \mathrm{RFn}(f) \to c \in \mathbb{R} \to \alpha \in \mathbb{R} \to \varphi_{97}$}
\BIC{$(\varphi_{84} \to \varphi_{86}) \to \varphi_{87} \to \mathrm{RFn}(f) \to c \in \mathbb{R} \to \alpha \in \mathbb{R} \to \varphi_{97}$}
\end{ptree}
\noindent$\Pi_{90}$:
\begin{ptree}
\AXC{$\varphi_{87}$\lbl{A4}}
\AXC{$\varphi_{84} \to \varphi_{86}$\lbl{A3}}
\AXC{$\stackrel{\Pi_{89}}{\vdots}$}\noLine
\UIC{$(\varphi_{84} \to \varphi_{86}) \to \varphi_{87} \to \mathrm{RFn}(f) \to c \in \mathbb{R} \to \alpha \in \mathbb{R} \to \varphi_{97}$}
\BIC{$\varphi_{87} \to \mathrm{RFn}(f) \to c \in \mathbb{R} \to \alpha \in \mathbb{R} \to \varphi_{97}$}
\BIC{$\mathrm{RFn}(f) \to c \in \mathbb{R} \to \alpha \in \mathbb{R} \to \varphi_{97}$}
\AXC{$(\mathrm{RFn}(f) \to c \in \mathbb{R} \to \alpha \in \mathbb{R} \to \varphi_{97}) \to \varphi_{322}$\lbl{A1}}
\BIC{$\varphi_{322}$}
\end{ptree}
\noindent$\Pi_{91}$:
\begin{ptree}
\AXC{$\stackrel{\Pi_{88}}{\vdots}$}\noLine
\UIC{$\mathrm{RFn}(f) \to c \in \mathbb{R} \to \alpha \in \mathbb{R} \to \mathrm{PNbhd}(f,c) \to \varphi_{93}$}
\AXC{$\varphi_{220}$\lbl{A1}}
\BIC{$\varphi_{89} \to \neg \varphi_{62}$}
\UIC{$\forall e_{0}\,(\varphi_{89} \to \neg \varphi_{62})$}
\AXC{$\stackrel{\Pi_{90}}{\vdots}$}\noLine
\UIC{$\varphi_{322}$}
\BIC{$(\forall e_{0}\,(\varphi_{89} \to \neg \varphi_{62}) \to \forall e_{0}\,(\varphi_{89}) \to \forall e_{0}\,(\neg \varphi_{62})) \to \varphi_{311}$}
\end{ptree}
\noindent$\Pi_{92}$:
\begin{ptree}
\AXC{$\varphi_{89} \to \forall e_{0}\,(\varphi_{89})$\lbl{A4}}
\AXC{$\forall e_{0}\,(\varphi_{89} \to \neg \varphi_{62}) \to \forall e_{0}\,(\varphi_{89}) \to \forall e_{0}\,(\neg \varphi_{62})$\lbl{A3}}
\AXC{$\stackrel{\Pi_{91}}{\vdots}$}\noLine
\UIC{$(\forall e_{0}\,(\varphi_{89} \to \neg \varphi_{62}) \to \forall e_{0}\,(\varphi_{89}) \to \forall e_{0}\,(\neg \varphi_{62})) \to \varphi_{311}$}
\BIC{$\varphi_{311}$}
\BIC{$\forall \varepsilon\,(\forall e_{0}\,(\neg \varphi_{62}) \to \neg \neg (\varepsilon \in \mathbb{R} \to 0 < \varepsilon \to \varphi_{45})) \to \varphi_{162} \to \varphi_{165}$}
\end{ptree}
\noindent$\Pi_{93}$:
\begin{ptree}
\AXC{$\varphi_{162}$\lbl{A3}}
\AXC{$\forall e_{0}\,(\neg \varphi_{62}) \to \neg \neg (\varepsilon \in \mathbb{R} \to 0 < \varepsilon \to \varphi_{45})$\lbl{A2}}
\UIC{$\forall \varepsilon\,(\forall e_{0}\,(\neg \varphi_{62}) \to \neg \neg (\varepsilon \in \mathbb{R} \to 0 < \varepsilon \to \varphi_{45}))$}
\AXC{$\stackrel{\Pi_{92}}{\vdots}$}\noLine
\UIC{$\forall \varepsilon\,(\forall e_{0}\,(\neg \varphi_{62}) \to \neg \neg (\varepsilon \in \mathbb{R} \to 0 < \varepsilon \to \varphi_{45})) \to \varphi_{162} \to \varphi_{165}$}
\BIC{$\varphi_{162} \to \varphi_{165}$}
\BIC{$\varphi_{165}$}
\end{ptree}
\noindent$\Pi_{94}$:
\begin{ptree}
\AXC{$\stackrel{\Pi_{1}}{\vdots}$}\noLine
\UIC{$\mathrm{RFn}(f) \to c \in \mathbb{R} \to \alpha \in \mathbb{R} \to \mathrm{PNbhd}(f,c) \to \varphi_{80}$}
\AXC{$\forall e_{0}\,(\neg \varphi_{62}) \to \forall \varepsilon\,\forall e_{0}\,(\neg \varphi_{62})$\lbl{A4}}
\AXC{$\stackrel{\Pi_{93}}{\vdots}$}\noLine
\UIC{$\varphi_{165}$}
\BIC{$\mathrm{RFn}(f) \to c \in \mathbb{R} \to \alpha \in \mathbb{R} \to \mathrm{PNbhd}(f,c) \to \varphi_{79}$}
\AXC{$\varphi_{221}$\lbl{A1}}
\BIC{$\varphi_{167}$}
\BIC{$\mathrm{RFn}(f) \to c \in \mathbb{R} \to \alpha \in \mathbb{R} \to \mathrm{PNbhd}(f,c) \to \varphi_{78}$}
\end{ptree}
\noindent$\Pi_{95}$:
\begin{ptree}
\AXC{$\forall f\,\forall c\,\forall \alpha\,(\varphi_{73})$\lbl{DefFLim}}
\AXC{$\forall f\,\forall c\,\forall \alpha\,(\varphi_{73}) \to \forall c\,\forall \alpha\,(\varphi_{73})$\lbl{A2}}
\BIC{$\forall c\,\forall \alpha\,(\varphi_{73})$}
\AXC{$\forall c\,\forall \alpha\,(\varphi_{73}) \to \forall \alpha\,(\varphi_{73})$\lbl{A2}}
\BIC{$\forall \alpha\,(\varphi_{73})$}
\AXC{$\forall \alpha\,(\varphi_{73}) \to \varphi_{73}$\lbl{A2}}
\BIC{$\varphi_{73}$}
\end{ptree}
\noindent$\Pi_{96}$:
\begin{ptree}
\AXC{$\mathrm{RFn}(f) \to c \in \mathbb{R} \to c \in \mathbb{R}$\lbl{A1}}
\AXC{$\mathrm{RFn}(f) \to \mathrm{RFn}(f)$\lbl{A1}}
\AXC{$\stackrel{\Pi_{95}}{\vdots}$}\noLine
\UIC{$\varphi_{73}$}
\AXC{$\varphi_{73} \to \varphi_{181}$\lbl{A1}}
\BIC{$\varphi_{181}$}
\BIC{$(\mathrm{RFn}(f) \to c \in \mathbb{R} \to c \in \mathbb{R}) \to \varphi_{150}$}
\BIC{$\varphi_{150}$}
\end{ptree}
\noindent$\Pi_{97}$:
\begin{ptree}
\AXC{$\mathrm{RFn}(f) \to \varphi_{1}$\lbl{A1}}
\AXC{$\mathrm{RFn}(f) \to c \in \mathbb{R} \to \alpha \in \mathbb{R} \to \alpha \in \mathbb{R}$\lbl{A1}}
\AXC{$\stackrel{\Pi_{96}}{\vdots}$}\noLine
\UIC{$\varphi_{150}$}
\BIC{$(\mathrm{RFn}(f) \to \varphi_{1}) \to \varphi_{113}$}
\BIC{$\varphi_{113}$}
\end{ptree}
\noindent 主木:
\begin{ptree}
\AXC{$\stackrel{\Pi_{94}}{\vdots}$}\noLine
\UIC{$\mathrm{RFn}(f) \to c \in \mathbb{R} \to \alpha \in \mathbb{R} \to \mathrm{PNbhd}(f,c) \to \varphi_{78}$}
\AXC{$\stackrel{\Pi_{97}}{\vdots}$}\noLine
\UIC{$\varphi_{113}$}
\BIC{$\mathrm{RFn}(f) \to c \in \mathbb{R} \to \alpha \in \mathbb{R} \to \mathrm{PNbhd}(f,c) \to \varphi_{46}$}
\UIC{$\varphi_{67}$}
\UIC{$\forall c\,(\varphi_{67})$}
\UIC{$\forall f\,\forall c\,(\varphi_{67})$}
\end{ptree}
\begin{abbrevlist}
\item $\varphi_{1} :\equiv c \in \mathbb{R} \to \alpha \in \mathbb{R} \to \mathrm{PNbhd}(f,c) \to \mathrm{PNbhd}(f,c)$
\item $\varphi_{2} :\equiv \forall N\,(\neg (N \in \mathbb{N} \land \forall n \in \mathbb{N}\,(N \leq n \to \lvert a_{n} - c \rvert < \varepsilon)))$
\item $\varphi_{3} :\equiv \neg (N \in \mathbb{N} \land \forall n \in \mathbb{N}\,(N \leq n \to \lvert f \circ a(n) - \alpha \rvert < e_{0}))$
\item $\varphi_{4} :\equiv \forall n \in \mathbb{N}\,(a_{n} \in \mathrm{dom}(f) \land a_{n} \neq c) \to \mathrm{Conv}(f \circ a,\alpha)$
\item $\varphi_{5} :\equiv \forall x\,(\mathrm{RFn}(f) \to x \in \mathrm{dom}(f) \to x \in \mathbb{R}) \to \mathrm{RFn}(f) \to x \in \mathrm{dom}(f) \to x \in \mathbb{R}$
\item $\varphi_{6} :\equiv N \in \mathbb{N} \land \forall n \in \mathbb{N}\,(N \leq n \to \lvert f \circ a(n) - \alpha \rvert < e_{0}) \to N \leq N$
\item $\varphi_{7} :\equiv N \in \mathbb{N} \land \forall n \in \mathbb{N}\,(N \leq n \to \lvert f \circ a(n) - \alpha \rvert < e_{0}) \to 0 < 0$
\item $\varphi_{8} :\equiv \forall z\,(\lvert a_{n} - c \rvert < \underline{n}^{-1} \to \underline{n}^{-1} \leq z \to \lvert a_{n} - c \rvert < z)$
\item $\varphi_{9} :\equiv \forall f\,\forall x\,(\mathrm{RFn}(f) \to x \in \mathrm{dom}(f) \to x \in \mathbb{R}) \to \forall x\,(\mathrm{RFn}(f) \to x \in \mathrm{dom}(f) \to x \in \mathbb{R})$
\item $\varphi_{10} :\equiv K \in \mathbb{N} \land \underline{K}^{-1} < \varepsilon \to \forall n \in \mathbb{N}\,(K \leq n \to \lvert a_{n} - c \rvert < \varepsilon)$
\item $\varphi_{11} :\equiv N \in \mathbb{N} \land \forall n \in \mathbb{N}\,(N \leq n \to \lvert f \circ a(n) - \alpha \rvert < e_{0}) \to N \in \mathbb{N}$
\item $\varphi_{12} :\equiv \varepsilon \in \mathbb{R} \to 0 < \varepsilon \to K \in \mathbb{N} \land \underline{K}^{-1} < \varepsilon \to n \in \mathbb{N} \to K \leq n \to K \leq n$
\item $\varphi_{13} :\equiv \forall n\,(\mathrm{Seq}(a) \to n \in \mathbb{N} \to a_{n} \in \mathbb{R}) \to \mathrm{Seq}(a) \to n \in \mathbb{N} \to a_{n} \in \mathbb{R}$
\item $\varphi_{14} :\equiv 0 < \varepsilon \to K \in \mathbb{N} \land \underline{K}^{-1} < \varepsilon \to n \in \mathbb{N} \to K \leq n \to \underline{K}^{-1} < \varepsilon$
\item $\varphi_{15} :\equiv 0 < e_{0} \to \exists N \in \mathbb{N}\,\forall n \in \mathbb{N}\,(N \leq n \to \lvert f \circ a(n) - \alpha \rvert < e_{0})$
\item $\varphi_{16} :\equiv \forall z\,(\lvert a_{n} - c \rvert < \underline{K}^{-1} \to \underline{K}^{-1} < z \to \lvert a_{n} - c \rvert < z)$
\item $\varphi_{17} :\equiv \forall w\,(f \circ a(N) = w \to \lvert f \circ a(N) - \alpha \rvert < e_{0} \to \lvert w - \alpha \rvert < e_{0})$
\item $\varphi_{18} :\equiv K \in \mathbb{N} \land \underline{K}^{-1} < \varepsilon \to n \in \mathbb{N} \to K \leq n \to \lvert a_{n} - c \rvert < \underline{n}^{-1}$
\item $\varphi_{19} :\equiv K \in \mathbb{N} \land \underline{K}^{-1} < \varepsilon \to n \in \mathbb{N} \to K \leq n \to \lvert a_{n} - c \rvert < \underline{K}^{-1}$
\item $\varphi_{20} :\equiv \varepsilon \in \mathbb{R} \to 0 < \varepsilon \to \exists N \in \mathbb{N}\,\forall n \in \mathbb{N}\,(N \leq n \to \lvert a_{n} - c \rvert < \varepsilon)$
\item $\varphi_{21} :\equiv \forall a\,\forall n\,(\mathrm{Seq}(a) \to n \in \mathbb{N} \to a_{n} \in \mathbb{R}) \to \forall n\,(\mathrm{Seq}(a) \to n \in \mathbb{N} \to a_{n} \in \mathbb{R})$
\item $\varphi_{22} :\equiv 0 < \varepsilon \to K \in \mathbb{N} \land \underline{K}^{-1} < \varepsilon \to n \in \mathbb{N} \to K \leq n \to \underline{n}^{-1} \leq \underline{K}^{-1}$
\item $\varphi_{23} :\equiv 0 < \varepsilon \to K \in \mathbb{N} \land \underline{K}^{-1} < \varepsilon \to n \in \mathbb{N} \to K \leq n \to \lvert a_{n} - c \rvert < \varepsilon$
\item $\varphi_{24} :\equiv x \in \mathrm{dom}(f) \to 0 < \lvert x - c \rvert \to \lvert x - c \rvert < \underline{k}^{-1} \to \lvert f(x) - \alpha \rvert < e_{0}$
\item $\varphi_{25} :\equiv \forall x \in \mathbb{R}\,(0 < x \to 0 < x^{-1}) \to \underline{k} \in \mathbb{R} \to 0 < \underline{k} \to 0 < \underline{k}^{-1}$
\item $\varphi_{26} :\equiv \varepsilon \in \mathbb{R} \to 0 < \varepsilon \to K \in \mathbb{N} \land \underline{K}^{-1} < \varepsilon \to K \in \mathbb{N} \land \underline{K}^{-1} < \varepsilon$
\item $\varphi_{27} :\equiv \forall \varepsilon \in \mathbb{R}\,(0 < \varepsilon \to \exists N \in \mathbb{N}\,\forall n \in \mathbb{N}\,(N \leq n \to \lvert a_{n} - c \rvert < \varepsilon))$
\item $\varphi_{28} :\equiv K \in \mathbb{N} \land \underline{K}^{-1} < \varepsilon \to K \in \mathbb{N} \land \forall n \in \mathbb{N}\,(K \leq n \to \lvert a_{n} - c \rvert < \varepsilon)$
\item $\varphi_{29} :\equiv \forall x \in \mathrm{dom}(f)\,(0 < \lvert x - c \rvert \to \lvert x - c \rvert < \underline{k}^{-1} \to \lvert f(x) - \alpha \rvert < e_{0})$
\item $\varphi_{30} :\equiv \forall x\,\forall y\,\forall z\,(x < y \to y \leq z \to x < z) \to \forall y\,\forall z\,(\lvert a_{n} - c \rvert < y \to y \leq z \to \lvert a_{n} - c \rvert < z)$
\item $\varphi_{31} :\equiv \forall x \in \mathbb{R}\,(x \neq 0 \to x^{-1} \in \mathbb{R}) \to \underline{k} \in \mathbb{R} \to \underline{k} \neq 0 \to \underline{k}^{-1} \in \mathbb{R}$
\item $\varphi_{32} :\equiv K \in \mathbb{N} \land \underline{K}^{-1} < \varepsilon \to \exists N \in \mathbb{N}\,\forall n \in \mathbb{N}\,(N \leq n \to \lvert a_{n} - c \rvert < \varepsilon)$
\item $\varphi_{33} :\equiv \forall n \in \mathbb{N}\,(a_{n} \in \mathrm{dom}(f) \land a_{n} \neq c) \to n \in \mathbb{N} \to a_{n} \in \mathrm{dom}(f) \land a_{n} \neq c$
\item $\varphi_{34} :\equiv \lvert a_{n} - c \rvert < \underline{n}^{-1} \to \underline{n}^{-1} \leq \underline{K}^{-1} \to \lvert a_{n} - c \rvert < \underline{K}^{-1}$
\item $\varphi_{35} :\equiv \forall n \in \mathbb{N}\,(f \circ a(n) = f(a_{n})) \to N \in \mathbb{N} \to f \circ a(N) = f(a_{N})$
\item $\varphi_{36} :\equiv \forall x\,\forall y\,\forall z\,(x < y \to y < z \to x < z) \to \forall y\,\forall z\,(\lvert a_{n} - c \rvert < y \to y < z \to \lvert a_{n} - c \rvert < z)$
\item $\varphi_{37} :\equiv x \in \mathbb{R} \land (x \in \mathrm{dom}(f) \land 0 < \lvert x - c \rvert \land \lvert x - c \rvert < \underline{k}^{-1} \land \neg \lvert f(x) - \alpha \rvert < e_{0})$
\item $\varphi_{38} :\equiv N \in \mathbb{N} \land \forall n \in \mathbb{N}\,(N \leq n \to \lvert f \circ a(n) - \alpha \rvert < e_{0}) \to \lvert f(a_{N}) - \alpha \rvert < e_{0}$
\item $\varphi_{39} :\equiv \exists x \in \mathbb{R}\,(x \in \mathrm{dom}(f) \land 0 < \lvert x - c \rvert \land \lvert x - c \rvert < \underline{k}^{-1} \land \neg \lvert f(x) - \alpha \rvert < e_{0})$
\item $\varphi_{40} :\equiv \forall n \in \mathbb{N}\,(a_{n} \in \mathrm{dom}(f)) \to \mathrm{Seq}(f \circ a) \land \forall n \in \mathbb{N}\,(f \circ a(n) = f(a_{n}))$
\item $\varphi_{41} :\equiv N \in \mathbb{N} \land \forall n \in \mathbb{N}\,(N \leq n \to \lvert f \circ a(n) - \alpha \rvert < e_{0}) \to f \circ a(N) = f(a_{N})$
\item $\varphi_{42} :\equiv N \in \mathbb{N} \land \forall n \in \mathbb{N}\,(N \leq n \to \lvert f \circ a(n) - \alpha \rvert < e_{0}) \to \lvert f \circ a(N) - \alpha \rvert < e_{0}$
\item $\varphi_{43} :\equiv f \circ a(N) = f(a_{N}) \to \lvert f \circ a(N) - \alpha \rvert < e_{0} \to \lvert f(a_{N}) - \alpha \rvert < e_{0}$
\item $\varphi_{44} :\equiv \forall \delta\,(\neg (\delta \in \mathbb{R} \land 0 < \delta \land \forall x \in \mathrm{dom}(f)\,(0 < \lvert x - c \rvert \to \lvert x - c \rvert < \delta \to \lvert f(x) - \alpha \rvert < e_{0})))$
\item $\varphi_{45} :\equiv \exists \delta\,(\delta \in \mathbb{R} \land 0 < \delta \land \forall x \in \mathrm{dom}(f)\,(0 < \lvert x - c \rvert \to \lvert x - c \rvert < \delta \to \lvert f(x) - \alpha \rvert < \varepsilon))$
\item $\varphi_{46} :\equiv \forall a\,(\mathrm{Seq}(a) \to \mathrm{Conv}(a,c) \to \varphi_{4}) \to \mathrm{FLim}(f,c,\alpha)$
\item $\varphi_{47} :\equiv \forall a\,\forall \alpha\,\forall \varepsilon\,(\mathrm{Conv}(a,\alpha) \to \varepsilon \in \mathbb{R} \to 0 < \varepsilon \to \exists N \in \mathbb{N}\,\forall n \in \mathbb{N}\,(N \leq n \to \lvert a_{n} - \alpha \rvert < \varepsilon))$
\item $\varphi_{48} :\equiv \forall n \in \mathbb{N}\,(N \leq n \to \lvert f \circ a(n) - \alpha \rvert < e_{0}) \to N \in \mathbb{N} \to N \leq N \to \lvert f \circ a(N) - \alpha \rvert < e_{0}$
\item $\varphi_{49} :\equiv \forall n\,(K \in \mathbb{N} \to n \in \mathbb{N} \to K \leq n \to \underline{n}^{-1} \leq \underline{K}^{-1}) \to K \in \mathbb{N} \to n \in \mathbb{N} \to K \leq n \to \underline{n}^{-1} \leq \underline{K}^{-1}$
\item $\varphi_{50} :\equiv a_{N} \in \mathrm{dom}(f) \land 0 < \lvert a_{N} - c \rvert \land \lvert a_{N} - c \rvert < \underline{N}^{-1} \land \neg \lvert f(a_{N}) - \alpha \rvert < e_{0}$
\item $\varphi_{51} :\equiv a_{n} \in \mathrm{dom}(f) \land 0 < \lvert a_{n} - c \rvert \land \lvert a_{n} - c \rvert < \underline{n}^{-1} \land \neg \lvert f(a_{n}) - \alpha \rvert < e_{0}$
\item $\varphi_{52} :\equiv N \in \mathbb{N} \land \forall n \in \mathbb{N}\,(N \leq n \to \lvert f \circ a(n) - \alpha \rvert < e_{0}) \to N \in \mathbb{N} \to f \circ a(N) = f(a_{N})$
\item $\varphi_{53} :\equiv \forall K\,\forall n\,(K \in \mathbb{N} \to n \in \mathbb{N} \to K \leq n \to \underline{n}^{-1} \leq \underline{K}^{-1}) \to \forall n\,(K \in \mathbb{N} \to n \in \mathbb{N} \to K \leq n \to \underline{n}^{-1} \leq \underline{K}^{-1})$
\item $\varphi_{54} :\equiv N \in \mathbb{N} \land \forall n \in \mathbb{N}\,(N \leq n \to \lvert f \circ a(n) - \alpha \rvert < e_{0}) \to \forall n \in \mathbb{N}\,(f \circ a(n) = f(a_{n}))$
\item $\varphi_{55} :\equiv N \in \mathbb{N} \land \forall n \in \mathbb{N}\,(N \leq n \to \lvert f \circ a(n) - \alpha \rvert < e_{0}) \to N \in \mathbb{N} \to N \leq N \to \lvert f \circ a(N) - \alpha \rvert < e_{0}$
\item $\varphi_{56} :\equiv \forall \varepsilon \in \mathbb{R}\,(0 < \varepsilon \to \exists K \in \mathbb{N}\,(\underline{K}^{-1} < \varepsilon)) \to \varepsilon \in \mathbb{R} \to 0 < \varepsilon \to \exists K \in \mathbb{N}\,(\underline{K}^{-1} < \varepsilon)$
\item $\varphi_{57} :\equiv \forall \varepsilon\,(\mathrm{Conv}(f \circ a,\alpha) \to \varepsilon \in \mathbb{R} \to 0 < \varepsilon \to \exists N \in \mathbb{N}\,\forall n \in \mathbb{N}\,(N \leq n \to \lvert f \circ a(n) - \alpha \rvert < \varepsilon))$
\item $\varphi_{58} :\equiv N \in \mathbb{N} \land \forall n \in \mathbb{N}\,(N \leq n \to \lvert f \circ a(n) - \alpha \rvert < e_{0}) \to \forall n \in \mathbb{N}\,(N \leq n \to \lvert f \circ a(n) - \alpha \rvert < e_{0})$
\item $\varphi_{59} :\equiv \forall y\,\forall c\,(y \in \mathbb{R} \to c \in \mathbb{R} \to 0 < \lvert y - c \rvert \to y \neq c) \to \forall c\,(a_{n} \in \mathbb{R} \to c \in \mathbb{R} \to 0 < \lvert a_{n} - c \rvert \to a_{n} \neq c)$
\item $\varphi_{60} :\equiv \forall \alpha\,(\mathrm{Conv}(a,\alpha) \leftrightarrow \mathrm{Seq}(a) \land \alpha \in \mathbb{R} \land \forall \varepsilon \in \mathbb{R}\,(0 < \varepsilon \to \exists N \in \mathbb{N}\,\forall n \in \mathbb{N}\,(N \leq n \to \lvert a_{n} - \alpha \rvert < \varepsilon)))$
\item $\varphi_{61} :\equiv \forall k \in \mathbb{N}\,(a_{k} \in \mathrm{dom}(f) \land 0 < \lvert a_{k} - c \rvert \land \lvert a_{k} - c \rvert < \underline{k}^{-1} \land \neg \lvert f(a_{k}) - \alpha \rvert < e_{0})$
\item $\varphi_{62} :\equiv \neg (e_{0} \in \mathbb{R} \to 0 < e_{0} \to \exists \delta\,(\delta \in \mathbb{R} \land 0 < \delta \land \forall x \in \mathrm{dom}(f)\,(0 < \lvert x - c \rvert \to \lvert x - c \rvert < \delta \to \lvert f(x) - \alpha \rvert < e_{0})))$
\item $\varphi_{63} :\equiv N \in \mathbb{N} \land \forall n \in \mathbb{N}\,(N \leq n \to \lvert f \circ a(n) - \alpha \rvert < e_{0}) \to N \in \mathbb{N} \land \forall n \in \mathbb{N}\,(N \leq n \to \lvert f \circ a(n) - \alpha \rvert < e_{0})$
\item $\varphi_{64} :\equiv \neg (\mathrm{Seq}(q) \land \forall k \in \mathbb{N}\,(q(k) \in \mathrm{dom}(f) \land 0 < \lvert q(k) - c \rvert \land \lvert q(k) - c \rvert < \underline{k}^{-1} \land \neg \lvert f(q(k)) - \alpha \rvert < e_{0}))$
\item $\varphi_{65} :\equiv \exists q\,(\mathrm{Seq}(q) \land \forall k \in \mathbb{N}\,(q(k) \in \mathrm{dom}(f) \land 0 < \lvert q(k) - c \rvert \land \lvert q(k) - c \rvert < \underline{k}^{-1} \land \neg \lvert f(q(k)) - \alpha \rvert < e_{0}))$
\item $\varphi_{66} :\equiv \forall c\,(a_{n} \in \mathbb{R} \to c \in \mathbb{R} \to 0 < \lvert a_{n} - c \rvert \to a_{n} \neq c) \to a_{n} \in \mathbb{R} \to c \in \mathbb{R} \to 0 < \lvert a_{n} - c \rvert \to a_{n} \neq c$
\item $\varphi_{67} :\equiv \forall \alpha\,(\mathrm{RFn}(f) \to c \in \mathbb{R} \to \alpha \in \mathbb{R} \to \mathrm{PNbhd}(f,c) \to \varphi_{46})$
\item $\varphi_{68} :\equiv N \in \mathbb{N} \land \forall n \in \mathbb{N}\,(N \leq n \to \lvert f \circ a(n) - \alpha \rvert < e_{0}) \to \varphi_{50}$
\item $\varphi_{69} :\equiv \forall a\,(\mathrm{Seq}(a) \to \mathrm{Conv}(a,c) \to \varphi_{4}) \to \mathrm{Seq}(a) \to \mathrm{Conv}(a,c) \to \varphi_{4}$
\item $\varphi_{70} :\equiv \forall a\,(\mathrm{Seq}(a) \to \mathrm{Conv}(a,c) \to \varphi_{4}) \to \forall a\,(\mathrm{Seq}(a) \to \mathrm{Conv}(a,c) \to \varphi_{4})$
\item $\varphi_{71} :\equiv \varepsilon \in \mathbb{R} \to 0 < \varepsilon \to K \in \mathbb{N} \land \underline{K}^{-1} < \varepsilon \to n \in \mathbb{N} \to K \leq n \to \varphi_{51}$
\item $\varphi_{72} :\equiv \mathrm{Seq}(a) \land \varphi_{61} \to \varepsilon \in \mathbb{R} \to 0 < \varepsilon \to \exists K \in \mathbb{N}\,(\underline{K}^{-1} < \varepsilon)$
\item $\varphi_{73} :\equiv \mathrm{FLim}(f,c,\alpha) \leftrightarrow \mathrm{RFn}(f) \land c \in \mathbb{R} \land \alpha \in \mathbb{R} \land \mathrm{PNbhd}(f,c) \land \forall \varepsilon \in \mathbb{R}\,(0 < \varepsilon \to \varphi_{45})$
\item $\varphi_{74} :\equiv \mathrm{Seq}(a) \land \varphi_{61} \to \varepsilon \in \mathbb{R} \to 0 < \varepsilon \to K \in \mathbb{N} \land \underline{K}^{-1} < \varepsilon \to K \in \mathbb{N}$
\item $\varphi_{75} :\equiv \mathrm{Seq}(a) \land \varphi_{61} \to \mathrm{Seq}(f \circ a) \land \forall n \in \mathbb{N}\,(f \circ a(n) = f(a_{n}))$
\item $\varphi_{76} :\equiv \mathrm{Seq}(a) \land \varphi_{61} \to \exists N \in \mathbb{N}\,\forall n \in \mathbb{N}\,(N \leq n \to \lvert f \circ a(n) - \alpha \rvert < e_{0})$
\item $\varphi_{77} :\equiv \mathrm{Seq}(a) \land \varphi_{61} \to \varepsilon \in \mathbb{R} \to 0 < \varepsilon \to K \in \mathbb{N} \land \underline{K}^{-1} < \varepsilon \to n \in \mathbb{N} \to n \in \mathbb{N}$
\item $\varphi_{78} :\equiv \forall a\,(\mathrm{Seq}(a) \to \mathrm{Conv}(a,c) \to \varphi_{4}) \to \forall \varepsilon \in \mathbb{R}\,(0 < \varepsilon \to \varphi_{45})$
\item $\varphi_{79} :\equiv \forall a\,(\mathrm{Seq}(a) \to \mathrm{Conv}(a,c) \to \varphi_{4}) \to \neg \forall \varepsilon \in \mathbb{R}\,(0 < \varepsilon \to \varphi_{45}) \to 0 < 0$
\item $\varphi_{80} :\equiv \forall a\,(\mathrm{Seq}(a) \to \mathrm{Conv}(a,c) \to \varphi_{4}) \to \neg \forall \varepsilon \in \mathbb{R}\,(0 < \varepsilon \to \varphi_{45}) \to \neg 0 < 0$
\item $\varphi_{81} :\equiv \forall a\,(\mathrm{RFn}(f) \to \mathrm{Seq}(a) \to \varphi_{40}) \to \mathrm{RFn}(f) \to \mathrm{Seq}(a) \to \varphi_{40}$
\item $\varphi_{82} :\equiv \forall f\,\forall a\,(\mathrm{RFn}(f) \to \mathrm{Seq}(a) \to \varphi_{40}) \to \forall a\,(\mathrm{RFn}(f) \to \mathrm{Seq}(a) \to \varphi_{40})$
\item $\varphi_{83} :\equiv \forall \varepsilon\,(\neg \neg (\varepsilon \in \mathbb{R} \to 0 < \varepsilon \to \varphi_{45})) \to \neg \neg (\varepsilon \in \mathbb{R} \to 0 < \varepsilon \to \varphi_{45})$
\item $\varphi_{84} :\equiv \forall \varepsilon\,(\forall \varepsilon\,(\neg \neg (\varepsilon \in \mathbb{R} \to 0 < \varepsilon \to \varphi_{45})) \to \varepsilon \in \mathbb{R} \to 0 < \varepsilon \to \varphi_{45})$
\item $\varphi_{85} :\equiv \neg \forall \varepsilon \in \mathbb{R}\,(0 < \varepsilon \to \varphi_{45}) \to \exists \varepsilon\,(\neg (\varepsilon \in \mathbb{R} \to 0 < \varepsilon \to \varphi_{45}))$
\item $\varphi_{86} :\equiv \forall \varepsilon\,\forall \varepsilon\,(\neg \neg (\varepsilon \in \mathbb{R} \to 0 < \varepsilon \to \varphi_{45})) \to \forall \varepsilon \in \mathbb{R}\,(0 < \varepsilon \to \varphi_{45})$
\item $\varphi_{87} :\equiv \forall \varepsilon\,(\neg \neg (\varepsilon \in \mathbb{R} \to 0 < \varepsilon \to \varphi_{45})) \to \forall \varepsilon\,\forall \varepsilon\,(\neg \neg (\varepsilon \in \mathbb{R} \to 0 < \varepsilon \to \varphi_{45}))$
\item $\varphi_{88} :\equiv \forall a\,(\neg (\mathrm{Seq}(a) \land \varphi_{61})) \to \forall q\,\forall a\,(\neg (\mathrm{Seq}(a) \land \varphi_{61}))$
\item $\varphi_{89} :\equiv \mathrm{RFn}(f) \land c \in \mathbb{R} \land \alpha \in \mathbb{R} \land \mathrm{PNbhd}(f,c) \land \forall a\,(\mathrm{Seq}(a) \to \mathrm{Conv}(a,c) \to \varphi_{4}) \land \neg \forall \varepsilon \in \mathbb{R}\,(0 < \varepsilon \to \varphi_{45}) \land \neg 0 < 0$
\item $\varphi_{90} :\equiv \neg 0 < 0 \to \mathrm{RFn}(f) \to c \in \mathbb{R} \to \alpha \in \mathbb{R} \to \mathrm{PNbhd}(f,c) \to \varphi_{80}$
\item $\varphi_{91} :\equiv \varphi_{62} \to \mathrm{Seq}(a) \land \varphi_{61} \to \forall n \in \mathbb{N}\,(a_{n} \in \mathrm{dom}(f) \land a_{n} \neq c)$
\item $\varphi_{92} :\equiv \forall a\,(\mathrm{Seq}(a) \to \mathrm{Conv}(a,c) \to \varphi_{4}) \to \neg \forall \varepsilon \in \mathbb{R}\,(0 < \varepsilon \to \varphi_{45}) \to \neg \forall \varepsilon \in \mathbb{R}\,(0 < \varepsilon \to \varphi_{45})$
\item $\varphi_{93} :\equiv \forall a\,(\mathrm{Seq}(a) \to \mathrm{Conv}(a,c) \to \varphi_{4}) \to \neg \forall \varepsilon \in \mathbb{R}\,(0 < \varepsilon \to \varphi_{45}) \to \varphi_{62} \to 0 < 0$
\item $\varphi_{94} :\equiv \forall a\,(\mathrm{Seq}(a) \to \mathrm{Conv}(a,c) \to \varphi_{4}) \to \neg \forall \varepsilon \in \mathbb{R}\,(0 < \varepsilon \to \varphi_{45}) \to \varphi_{62} \to 0 < e_{0}$
\item $\varphi_{95} :\equiv \forall a\,(\mathrm{Seq}(a) \to \mathrm{Conv}(a,c) \to \varphi_{4}) \to \neg \forall \varepsilon \in \mathbb{R}\,(0 < \varepsilon \to \varphi_{45}) \to \varphi_{62} \to e_{0} \in \mathbb{R}$
\item $\varphi_{96} :\equiv \neg \forall \varepsilon \in \mathbb{R}\,(0 < \varepsilon \to \varphi_{45}) \to \varphi_{62} \to k \in \mathbb{N} \to \neg \varphi_{24} \to x \in \mathrm{dom}(f)$
\item $\varphi_{97} :\equiv \mathrm{PNbhd}(f,c) \to \forall a\,(\mathrm{Seq}(a) \to \mathrm{Conv}(a,c) \to \varphi_{4}) \to \varphi_{85}$
\item $\varphi_{98} :\equiv \forall a\,(\mathrm{Seq}(a) \to \mathrm{Conv}(a,c) \to \varphi_{4}) \to \neg \forall \varepsilon \in \mathbb{R}\,(0 < \varepsilon \to \varphi_{45}) \to \varphi_{62} \to k \in \mathbb{N} \to k \in \mathbb{N}$
\item $\varphi_{99} :\equiv \forall a\,(\mathrm{Seq}(a) \to \mathrm{Conv}(a,c) \to \varphi_{4}) \to \neg \forall \varepsilon \in \mathbb{R}\,(0 < \varepsilon \to \varphi_{45}) \to \varphi_{62} \to k \in \mathbb{N} \to 0 \neq \underline{k}$
\item $\varphi_{100} :\equiv \forall a\,(\mathrm{Seq}(a) \to \mathrm{Conv}(a,c) \to \varphi_{4}) \to \neg \forall \varepsilon \in \mathbb{R}\,(0 < \varepsilon \to \varphi_{45}) \to \varphi_{62} \to k \in \mathbb{N} \to \underline{k} \neq 0$
\item $\varphi_{101} :\equiv \forall a\,(\mathrm{Seq}(a) \to \mathrm{Conv}(a,c) \to \varphi_{4}) \to \neg \forall \varepsilon \in \mathbb{R}\,(0 < \varepsilon \to \varphi_{45}) \to \varphi_{62} \to k \in \mathbb{N} \to 0 < \underline{k}$
\item $\varphi_{102} :\equiv \forall a\,(\mathrm{Seq}(a) \to \mathrm{Conv}(a,c) \to \varphi_{4}) \to \neg \forall \varepsilon \in \mathbb{R}\,(0 < \varepsilon \to \varphi_{45}) \to \varphi_{62} \to k \in \mathbb{N} \to \underline{k} \in \mathbb{R}$
\item $\varphi_{103} :\equiv \forall a\,(\mathrm{Seq}(a) \to \mathrm{Conv}(a,c) \to \varphi_{4}) \to \neg \forall \varepsilon \in \mathbb{R}\,(0 < \varepsilon \to \varphi_{45}) \to \varphi_{62} \to k \in \mathbb{N} \to 0 < \underline{k}^{-1}$
\item $\varphi_{104} :\equiv \forall a\,(\mathrm{Seq}(a) \to \mathrm{Conv}(a,c) \to \varphi_{4}) \to \neg \forall \varepsilon \in \mathbb{R}\,(0 < \varepsilon \to \varphi_{45}) \to \varphi_{62} \to k \in \mathbb{N} \to \underline{k}^{-1} \in \mathbb{R}$
\item $\varphi_{105} :\equiv \neg \forall \varepsilon \in \mathbb{R}\,(0 < \varepsilon \to \varphi_{45}) \to \varphi_{62} \to \neg (\underline{k}^{-1} \in \mathbb{R} \land 0 < \underline{k}^{-1} \land \varphi_{29})$
\item $\varphi_{106} :\equiv \mathrm{RFn}(f) \land c \in \mathbb{R} \land \alpha \in \mathbb{R} \land \mathrm{PNbhd}(f,c) \land \forall a\,(\mathrm{Seq}(a) \to \mathrm{Conv}(a,c) \to \varphi_{4}) \land \neg \forall \varepsilon \in \mathbb{R}\,(0 < \varepsilon \to \varphi_{45}) \land \varphi_{62}$
\item $\varphi_{107} :\equiv \neg \forall \varepsilon \in \mathbb{R}\,(0 < \varepsilon \to \varphi_{45}) \to \varphi_{62} \to \mathrm{Seq}(a) \land \varphi_{61} \to \mathrm{Seq}(a)$
\item $\varphi_{108} :\equiv \varphi_{69} \to \mathrm{RFn}(f) \to c \in \mathbb{R} \to \alpha \in \mathbb{R} \to \mathrm{PNbhd}(f,c) \to \varphi_{69}$
\item $\varphi_{109} :\equiv \neg \forall \varepsilon \in \mathbb{R}\,(0 < \varepsilon \to \varphi_{45}) \to \varphi_{62} \to \mathrm{Seq}(a) \land \varphi_{61} \to \mathrm{Conv}(a,c)$
\item $\varphi_{110} :\equiv \neg \forall \varepsilon \in \mathbb{R}\,(0 < \varepsilon \to \varphi_{45}) \to \varphi_{62} \to \mathrm{Seq}(a) \land \varphi_{61} \to n \in \mathbb{N} \to n \in \mathbb{N}$
\item $\varphi_{111} :\equiv \neg \forall \varepsilon \in \mathbb{R}\,(0 < \varepsilon \to \varphi_{45}) \to \varphi_{62} \to \mathrm{Seq}(a) \land \varphi_{61} \to \mathrm{Conv}(f \circ a,\alpha)$
\item $\varphi_{112} :\equiv \neg \forall \varepsilon \in \mathbb{R}\,(0 < \varepsilon \to \varphi_{45}) \to \varphi_{62} \to \mathrm{Seq}(a) \land \varphi_{61} \to n \in \mathbb{N} \to a_{n} \neq c$
\item $\varphi_{113} :\equiv (\mathrm{RFn}(f) \to c \in \mathbb{R} \to \alpha \in \mathbb{R} \to \mathrm{PNbhd}(f,c) \to \varphi_{78}) \to \mathrm{RFn}(f) \to c \in \mathbb{R} \to \alpha \in \mathbb{R} \to \mathrm{PNbhd}(f,c) \to \varphi_{46}$
\item $\varphi_{114} :\equiv \neg \forall \varepsilon \in \mathbb{R}\,(0 < \varepsilon \to \varphi_{45}) \to \varphi_{62} \to \mathrm{Seq}(a) \land \varphi_{61} \to \varepsilon \in \mathbb{R} \to \varepsilon \in \mathbb{R}$
\item $\varphi_{115} :\equiv \neg \forall \varepsilon \in \mathbb{R}\,(0 < \varepsilon \to \varphi_{45}) \to \varphi_{62} \to \mathrm{Seq}(a) \land \varphi_{61} \to n \in \mathbb{N} \to a_{n} \in \mathbb{R}$
\item $\varphi_{116} :\equiv \neg \forall \varepsilon \in \mathbb{R}\,(0 < \varepsilon \to \varphi_{45}) \to \varphi_{62} \to \mathrm{Seq}(a) \land \varphi_{61} \to n \in \mathbb{N} \to a_{n} \in \mathrm{dom}(f)$
\item $\varphi_{117} :\equiv \neg \forall \varepsilon \in \mathbb{R}\,(0 < \varepsilon \to \varphi_{45}) \to \varphi_{62} \to \mathrm{Seq}(a) \land \varphi_{61} \to \forall n \in \mathbb{N}\,(a_{n} \in \mathrm{dom}(f))$
\item $\varphi_{118} :\equiv \neg \forall \varepsilon \in \mathbb{R}\,(0 < \varepsilon \to \varphi_{45}) \to \varphi_{62} \to \mathrm{Seq}(a) \land \varphi_{61} \to \varepsilon \in \mathbb{R} \to 0 < \varepsilon \to 0 < \varepsilon$
\item $\varphi_{119} :\equiv \neg \forall \varepsilon \in \mathbb{R}\,(0 < \varepsilon \to \varphi_{45}) \to \varphi_{62} \to \mathrm{Seq}(a) \land \varphi_{61} \to n \in \mathbb{N} \to 0 < \lvert a_{n} - c \rvert$
\item $\varphi_{120} :\equiv \neg \forall \varepsilon \in \mathbb{R}\,(0 < \varepsilon \to \varphi_{45}) \to \varphi_{62} \to \mathrm{Seq}(a) \land \varphi_{61} \to n \in \mathbb{N} \to a_{n} \in \mathrm{dom}(f) \land a_{n} \neq c$
\item $\varphi_{121} :\equiv (\underline{k} = 0 \to 0 = \underline{k}) \to \mathrm{RFn}(f) \to c \in \mathbb{R} \to \alpha \in \mathbb{R} \to \mathrm{PNbhd}(f,c) \to \varphi_{100}$
\item $\varphi_{122} :\equiv \forall a\,(\mathrm{Seq}(a) \to \mathrm{Conv}(a,c) \to \varphi_{4}) \to \neg \forall \varepsilon \in \mathbb{R}\,(0 < \varepsilon \to \varphi_{45}) \to \varphi_{62} \to k \in \mathbb{N} \to \neg \varphi_{29}$
\item $\varphi_{123} :\equiv \forall a\,(\mathrm{Seq}(a) \to \mathrm{Conv}(a,c) \to \varphi_{4}) \to \neg \forall \varepsilon \in \mathbb{R}\,(0 < \varepsilon \to \varphi_{45}) \to \varphi_{62} \to k \in \mathbb{N} \to \exists x\,(\neg \varphi_{24})$
\item $\varphi_{124} :\equiv \forall a\,(\mathrm{Seq}(a) \to \mathrm{Conv}(a,c) \to \varphi_{4}) \to \neg \forall \varepsilon \in \mathbb{R}\,(0 < \varepsilon \to \varphi_{45}) \to \varphi_{62} \to \varphi_{44}$
\item $\varphi_{125} :\equiv \forall a\,(\mathrm{Seq}(a) \to \mathrm{Conv}(a,c) \to \varphi_{4}) \to \neg \forall \varepsilon \in \mathbb{R}\,(0 < \varepsilon \to \varphi_{45}) \to \varphi_{62} \to k \in \mathbb{N} \to \neg \varphi_{24} \to x \in \mathbb{R}$
\item $\varphi_{126} :\equiv \forall a\,(\mathrm{Seq}(a) \to \mathrm{Conv}(a,c) \to \varphi_{4}) \to \neg \forall \varepsilon \in \mathbb{R}\,(0 < \varepsilon \to \varphi_{45}) \to \varphi_{62} \to k \in \mathbb{N} \to \varphi_{39}$
\item $\varphi_{127} :\equiv \forall a\,(\mathrm{Seq}(a) \to \mathrm{Conv}(a,c) \to \varphi_{4}) \to \neg \forall \varepsilon \in \mathbb{R}\,(0 < \varepsilon \to \varphi_{45}) \to \varphi_{62} \to \forall k \in \mathbb{N}\,(\varphi_{39})$
\item $\varphi_{128} :\equiv \mathrm{PNbhd}(f,c) \to \forall a\,(\mathrm{Seq}(a) \to \mathrm{Conv}(a,c) \to \varphi_{4}) \to \varphi_{96}$
\item $\varphi_{129} :\equiv \forall a\,(\mathrm{Seq}(a) \to \mathrm{Conv}(a,c) \to \varphi_{4}) \to \neg \forall \varepsilon \in \mathbb{R}\,(0 < \varepsilon \to \varphi_{45}) \to \varphi_{62} \to \varphi_{62}$
\item $\varphi_{130} :\equiv \neg \forall \varepsilon \in \mathbb{R}\,(0 < \varepsilon \to \varphi_{45}) \to \varphi_{62} \to \mathrm{Seq}(a) \land \varphi_{61} \to \varepsilon \in \mathbb{R} \to \varphi_{14}$
\item $\varphi_{131} :\equiv \forall a\,(\mathrm{Seq}(a) \to \mathrm{Conv}(a,c) \to \varphi_{4}) \to \neg \forall \varepsilon \in \mathbb{R}\,(0 < \varepsilon \to \varphi_{45}) \to \varphi_{62} \to \varphi_{65}$
\item $\varphi_{132} :\equiv \neg \forall \varepsilon \in \mathbb{R}\,(0 < \varepsilon \to \varphi_{45}) \to \varphi_{62} \to \mathrm{Seq}(a) \land \varphi_{61} \to \varepsilon \in \mathbb{R} \to \varphi_{22}$
\item $\varphi_{133} :\equiv \neg \forall \varepsilon \in \mathbb{R}\,(0 < \varepsilon \to \varphi_{45}) \to \varphi_{62} \to \mathrm{Seq}(a) \land \varphi_{61} \to \varepsilon \in \mathbb{R} \to \varphi_{23}$
\item $\varphi_{134} :\equiv \mathrm{PNbhd}(f,c) \to \forall a\,(\mathrm{Seq}(a) \to \mathrm{Conv}(a,c) \to \varphi_{4}) \to \varphi_{105}$
\item $\varphi_{135} :\equiv \forall a\,(\mathrm{Seq}(a) \to \mathrm{Conv}(a,c) \to \varphi_{4}) \to \neg \forall \varepsilon \in \mathbb{R}\,(0 < \varepsilon \to \varphi_{45}) \to \varphi_{62} \to \mathrm{Seq}(a) \land \varphi_{61} \to 0 < 0$
\item $\varphi_{136} :\equiv \neg \forall \varepsilon \in \mathbb{R}\,(0 < \varepsilon \to \varphi_{45}) \to \varphi_{62} \to \mathrm{Seq}(a) \land \varphi_{61} \to \varepsilon \in \mathbb{R} \to 0 < \varepsilon \to \varphi_{10}$
\item $\varphi_{137} :\equiv \neg \forall \varepsilon \in \mathbb{R}\,(0 < \varepsilon \to \varphi_{45}) \to \varphi_{62} \to \mathrm{Seq}(a) \land \varphi_{61} \to \varepsilon \in \mathbb{R} \to 0 < \varepsilon \to \varphi_{18}$
\item $\varphi_{138} :\equiv \neg \forall \varepsilon \in \mathbb{R}\,(0 < \varepsilon \to \varphi_{45}) \to \varphi_{62} \to \mathrm{Seq}(a) \land \varphi_{61} \to \varepsilon \in \mathbb{R} \to 0 < \varepsilon \to \varphi_{19}$
\item $\varphi_{139} :\equiv \mathrm{PNbhd}(f,c) \to \forall a\,(\mathrm{Seq}(a) \to \mathrm{Conv}(a,c) \to \varphi_{4}) \to \varphi_{107}$
\item $\varphi_{140} :\equiv \neg \forall \varepsilon \in \mathbb{R}\,(0 < \varepsilon \to \varphi_{45}) \to \varphi_{62} \to \mathrm{Seq}(a) \land \varphi_{61} \to \varepsilon \in \mathbb{R} \to 0 < \varepsilon \to \varphi_{28}$
\item $\varphi_{141} :\equiv \neg \forall \varepsilon \in \mathbb{R}\,(0 < \varepsilon \to \varphi_{45}) \to \varphi_{62} \to \mathrm{Seq}(a) \land \varphi_{61} \to \varepsilon \in \mathbb{R} \to 0 < \varepsilon \to \varphi_{32}$
\item $\varphi_{142} :\equiv \mathrm{PNbhd}(f,c) \to \forall a\,(\mathrm{Seq}(a) \to \mathrm{Conv}(a,c) \to \varphi_{4}) \to \varphi_{109}$
\item $\varphi_{143} :\equiv \forall \alpha\,\forall e_{0}\,(\forall k \in \mathbb{N}\,(\varphi_{39}) \to \varphi_{65}) \to \forall e_{0}\,(\forall k \in \mathbb{N}\,(\varphi_{39}) \to \varphi_{65})$
\item $\varphi_{144} :\equiv \neg \forall \varepsilon \in \mathbb{R}\,(0 < \varepsilon \to \varphi_{45}) \to \varphi_{62} \to \mathrm{Seq}(a) \land \varphi_{61} \to N \in \mathbb{N} \to \varphi_{50}$
\item $\varphi_{145} :\equiv \neg \forall \varepsilon \in \mathbb{R}\,(0 < \varepsilon \to \varphi_{45}) \to \varphi_{62} \to \mathrm{Seq}(a) \land \varphi_{61} \to n \in \mathbb{N} \to \varphi_{51}$
\item $\varphi_{146} :\equiv \forall c\,\forall \alpha\,\forall e_{0}\,(\forall k \in \mathbb{N}\,(\varphi_{39}) \to \varphi_{65}) \to \forall \alpha\,\forall e_{0}\,(\forall k \in \mathbb{N}\,(\varphi_{39}) \to \varphi_{65})$
\item $\varphi_{147} :\equiv \mathrm{PNbhd}(f,c) \to \forall a\,(\mathrm{Seq}(a) \to \mathrm{Conv}(a,c) \to \varphi_{4}) \to \varphi_{110}$
\item $\varphi_{148} :\equiv \mathrm{PNbhd}(f,c) \to \forall a\,(\mathrm{Seq}(a) \to \mathrm{Conv}(a,c) \to \varphi_{4}) \to \varphi_{111}$
\item $\varphi_{149} :\equiv \mathrm{PNbhd}(f,c) \to \forall a\,(\mathrm{Seq}(a) \to \mathrm{Conv}(a,c) \to \varphi_{4}) \to \varphi_{112}$
\item $\varphi_{150} :\equiv (\mathrm{RFn}(f) \to c \in \mathbb{R} \to \alpha \in \mathbb{R} \to \alpha \in \mathbb{R}) \to (\mathrm{RFn}(f) \to \varphi_{1}) \to \varphi_{113}$
\item $\varphi_{151} :\equiv \forall f\,\forall c\,\forall \alpha\,\forall e_{0}\,(\forall k \in \mathbb{N}\,(\varphi_{39}) \to \varphi_{65}) \to \forall c\,\forall \alpha\,\forall e_{0}\,(\forall k \in \mathbb{N}\,(\varphi_{39}) \to \varphi_{65})$
\item $\varphi_{152} :\equiv \mathrm{PNbhd}(f,c) \to \forall a\,(\mathrm{Seq}(a) \to \mathrm{Conv}(a,c) \to \varphi_{4}) \to \varphi_{114}$
\item $\varphi_{153} :\equiv \mathrm{PNbhd}(f,c) \to \forall a\,(\mathrm{Seq}(a) \to \mathrm{Conv}(a,c) \to \varphi_{4}) \to \varphi_{115}$
\item $\varphi_{154} :\equiv \neg \forall \varepsilon \in \mathbb{R}\,(0 < \varepsilon \to \varphi_{45}) \to \varphi_{62} \to \mathrm{Seq}(a) \land \varphi_{61} \to \mathrm{Seq}(a) \land \varphi_{61}$
\item $\varphi_{155} :\equiv \mathrm{PNbhd}(f,c) \to \forall a\,(\mathrm{Seq}(a) \to \mathrm{Conv}(a,c) \to \varphi_{4}) \to \varphi_{116}$
\item $\varphi_{156} :\equiv \varphi_{83} \to \forall \varepsilon\,(\neg \neg (\varepsilon \in \mathbb{R} \to 0 < \varepsilon \to \varphi_{45})) \to \varepsilon \in \mathbb{R} \to 0 < \varepsilon \to \varphi_{45}$
\item $\varphi_{157} :\equiv \mathrm{PNbhd}(f,c) \to \forall a\,(\mathrm{Seq}(a) \to \mathrm{Conv}(a,c) \to \varphi_{4}) \to \varphi_{117}$
\item $\varphi_{158} :\equiv \mathrm{PNbhd}(f,c) \to \forall a\,(\mathrm{Seq}(a) \to \mathrm{Conv}(a,c) \to \varphi_{4}) \to \varphi_{118}$
\item $\varphi_{159} :\equiv \mathrm{PNbhd}(f,c) \to \forall a\,(\mathrm{Seq}(a) \to \mathrm{Conv}(a,c) \to \varphi_{4}) \to \varphi_{119}$
\item $\varphi_{160} :\equiv \forall a\,(\mathrm{Seq}(a) \to \mathrm{Conv}(a,c) \to \varphi_{4}) \to \neg \forall \varepsilon \in \mathbb{R}\,(0 < \varepsilon \to \varphi_{45}) \to \varphi_{91}$
\item $\varphi_{161} :\equiv \forall a\,(\mathrm{Seq}(a) \to \mathrm{Conv}(a,c) \to \varphi_{4}) \to \neg \forall \varepsilon \in \mathbb{R}\,(0 < \varepsilon \to \varphi_{45}) \to \varphi_{62} \to k \in \mathbb{N} \to \neg \varphi_{24} \to \neg \varphi_{24}$
\item $\varphi_{162} :\equiv \forall \varepsilon\,(\forall e_{0}\,(\neg \varphi_{62}) \to \neg \neg (\varepsilon \in \mathbb{R} \to 0 < \varepsilon \to \varphi_{45})) \to \forall \varepsilon\,\forall e_{0}\,(\neg \varphi_{62}) \to \forall \varepsilon\,(\neg \neg (\varepsilon \in \mathbb{R} \to 0 < \varepsilon \to \varphi_{45}))$
\item $\varphi_{163} :\equiv \mathrm{PNbhd}(f,c) \to \forall a\,(\mathrm{Seq}(a) \to \mathrm{Conv}(a,c) \to \varphi_{4}) \to \varphi_{120}$
\item $\varphi_{164} :\equiv \forall a\,(\mathrm{Seq}(a) \to \mathrm{Conv}(a,c) \to \varphi_{4}) \to \neg \forall \varepsilon \in \mathbb{R}\,(0 < \varepsilon \to \varphi_{45}) \to \varphi_{62} \to \varphi_{72}$
\item $\varphi_{165} :\equiv (\forall e_{0}\,(\neg \varphi_{62}) \to \forall \varepsilon\,\forall e_{0}\,(\neg \varphi_{62})) \to \mathrm{RFn}(f) \to c \in \mathbb{R} \to \alpha \in \mathbb{R} \to \mathrm{PNbhd}(f,c) \to \varphi_{79}$
\item $\varphi_{166} :\equiv \forall a\,(\mathrm{Seq}(a) \to \mathrm{Conv}(a,c) \to \varphi_{4}) \to \neg \forall \varepsilon \in \mathbb{R}\,(0 < \varepsilon \to \varphi_{45}) \to \varphi_{62} \to k \in \mathbb{N} \to \neg \varphi_{24} \to \varphi_{37}$
\item $\varphi_{167} :\equiv (\mathrm{RFn}(f) \to c \in \mathbb{R} \to \alpha \in \mathbb{R} \to \mathrm{PNbhd}(f,c) \to \varphi_{80}) \to \mathrm{RFn}(f) \to c \in \mathbb{R} \to \alpha \in \mathbb{R} \to \mathrm{PNbhd}(f,c) \to \varphi_{78}$
\item $\varphi_{168} :\equiv \forall a\,(\mathrm{Seq}(a) \to \mathrm{Conv}(a,c) \to \varphi_{4}) \to \neg \forall \varepsilon \in \mathbb{R}\,(0 < \varepsilon \to \varphi_{45}) \to \varphi_{62} \to k \in \mathbb{N} \to \neg \varphi_{24} \to \varphi_{39}$
\item $\varphi_{169} :\equiv \forall q\,(\forall a\,(\neg (\mathrm{Seq}(a) \land \varphi_{61})) \to \varphi_{64}) \to \forall q\,\forall a\,(\neg (\mathrm{Seq}(a) \land \varphi_{61})) \to \forall q\,(\varphi_{64})$
\item $\varphi_{170} :\equiv \forall a\,(\mathrm{Seq}(a) \to \mathrm{Conv}(a,c) \to \varphi_{4}) \to \neg \forall \varepsilon \in \mathbb{R}\,(0 < \varepsilon \to \varphi_{45}) \to \varphi_{62} \to \varphi_{74}$
\item $\varphi_{171} :\equiv \forall a\,(\mathrm{Seq}(a) \to \mathrm{Conv}(a,c) \to \varphi_{4}) \to \neg \forall \varepsilon \in \mathbb{R}\,(0 < \varepsilon \to \varphi_{45}) \to \varphi_{62} \to \varphi_{75}$
\item $\varphi_{172} :\equiv \forall a\,(\mathrm{Seq}(a) \to \mathrm{Conv}(a,c) \to \varphi_{4}) \to \neg \forall \varepsilon \in \mathbb{R}\,(0 < \varepsilon \to \varphi_{45}) \to \varphi_{62} \to \varphi_{76}$
\item $\varphi_{173} :\equiv \forall a\,(\mathrm{Seq}(a) \to \mathrm{Conv}(a,c) \to \varphi_{4}) \to \neg \forall \varepsilon \in \mathbb{R}\,(0 < \varepsilon \to \varphi_{45}) \to \varphi_{62} \to \mathrm{Seq}(a) \land \varphi_{61} \to \varphi_{6}$
\item $\varphi_{174} :\equiv \forall a\,(\mathrm{Seq}(a) \to \mathrm{Conv}(a,c) \to \varphi_{4}) \to \neg \forall \varepsilon \in \mathbb{R}\,(0 < \varepsilon \to \varphi_{45}) \to \varphi_{62} \to \mathrm{Seq}(a) \land \varphi_{61} \to \varphi_{7}$
\item $\varphi_{175} :\equiv \forall a\,(\mathrm{Seq}(a) \to \mathrm{Conv}(a,c) \to \varphi_{4}) \to \neg \forall \varepsilon \in \mathbb{R}\,(0 < \varepsilon \to \varphi_{45}) \to \varphi_{62} \to \mathrm{Seq}(a) \land \varphi_{61} \to \varphi_{11}$
\item $\varphi_{176} :\equiv \forall a\,(\mathrm{Seq}(a) \to \mathrm{Conv}(a,c) \to \varphi_{4}) \to \neg \forall \varepsilon \in \mathbb{R}\,(0 < \varepsilon \to \varphi_{45}) \to \varphi_{62} \to \mathrm{Seq}(a) \land \varphi_{61} \to \varphi_{12}$
\item $\varphi_{177} :\equiv \forall a\,(\mathrm{Seq}(a) \to \mathrm{Conv}(a,c) \to \varphi_{4}) \to \neg \forall \varepsilon \in \mathbb{R}\,(0 < \varepsilon \to \varphi_{45}) \to \varphi_{62} \to \varphi_{77}$
\item $\varphi_{178} :\equiv \forall n\,(\varphi_{106} \land (\mathrm{Seq}(a) \land \varphi_{61}) \to n \in \mathbb{N} \to a_{n} \in \mathrm{dom}(f))$
\item $\varphi_{179} :\equiv \forall n\,(\varphi_{106} \land (\mathrm{Seq}(a) \land \varphi_{61})) \to \forall n \in \mathbb{N}\,(a_{n} \in \mathrm{dom}(f))$
\item $\varphi_{180} :\equiv \forall a\,(\mathrm{Seq}(a) \to \mathrm{Conv}(a,c) \to \varphi_{4}) \to \neg \forall \varepsilon \in \mathbb{R}\,(0 < \varepsilon \to \varphi_{45}) \to \varphi_{62} \to \mathrm{Seq}(a) \land \varphi_{61} \to \varphi_{20}$
\item $\varphi_{181} :\equiv (\mathrm{RFn}(f) \to \mathrm{RFn}(f)) \to (\mathrm{RFn}(f) \to c \in \mathbb{R} \to c \in \mathbb{R}) \to \varphi_{150}$
\item $\varphi_{182} :\equiv \forall a\,(\mathrm{Seq}(a) \to \mathrm{Conv}(a,c) \to \varphi_{4}) \to \neg \forall \varepsilon \in \mathbb{R}\,(0 < \varepsilon \to \varphi_{45}) \to \varphi_{62} \to \mathrm{Seq}(a) \land \varphi_{61} \to \varphi_{26}$
\item $\varphi_{183} :\equiv \forall a\,(\mathrm{Seq}(a) \to \mathrm{Conv}(a,c) \to \varphi_{4}) \to \neg \forall \varepsilon \in \mathbb{R}\,(0 < \varepsilon \to \varphi_{45}) \to \varphi_{62} \to \mathrm{Seq}(a) \land \varphi_{61} \to \varphi_{27}$
\item $\varphi_{184} :\equiv \varphi_{106} \land (\mathrm{Seq}(a) \land \varphi_{61}) \to n \in \mathbb{N} \to a_{n} \in \mathrm{dom}(f) \land a_{n} \neq c$
\item $\varphi_{185} :\equiv \forall n\,(\varphi_{106} \land (\mathrm{Seq}(a) \land \varphi_{61})) \to \forall n \in \mathbb{N}\,(a_{n} \in \mathrm{dom}(f) \land a_{n} \neq c)$
\item $\varphi_{186} :\equiv \forall a\,(\mathrm{Seq}(a) \to \mathrm{Conv}(a,c) \to \varphi_{4}) \to \neg \forall \varepsilon \in \mathbb{R}\,(0 < \varepsilon \to \varphi_{45}) \to \varphi_{62} \to \mathrm{Seq}(a) \land \varphi_{61} \to \varphi_{38}$
\item $\varphi_{187} :\equiv \varphi_{106} \land (\mathrm{Seq}(a) \land \varphi_{61}) \land \varepsilon \in \mathbb{R} \land 0 < \varepsilon \land (K \in \mathbb{N} \land \underline{K}^{-1} < \varepsilon)$
\item $\varphi_{188} :\equiv \forall a\,(\mathrm{Seq}(a) \to \mathrm{Conv}(a,c) \to \varphi_{4}) \to \neg \forall \varepsilon \in \mathbb{R}\,(0 < \varepsilon \to \varphi_{45}) \to \varphi_{62} \to \mathrm{Seq}(a) \land \varphi_{61} \to \varphi_{41}$
\item $\varphi_{189} :\equiv \forall a\,(\mathrm{Seq}(a) \to \mathrm{Conv}(a,c) \to \varphi_{4}) \to \neg \forall \varepsilon \in \mathbb{R}\,(0 < \varepsilon \to \varphi_{45}) \to \varphi_{62} \to \mathrm{Seq}(a) \land \varphi_{61} \to \varphi_{42}$
\item $\varphi_{190} :\equiv \mathrm{PNbhd}(f,c) \to \forall a\,(\mathrm{Seq}(a) \to \mathrm{Conv}(a,c) \to \varphi_{4}) \to \varphi_{130}$
\item $\varphi_{191} :\equiv \mathrm{PNbhd}(f,c) \to \forall a\,(\mathrm{Seq}(a) \to \mathrm{Conv}(a,c) \to \varphi_{4}) \to \varphi_{132}$
\item $\varphi_{192} :\equiv \mathrm{PNbhd}(f,c) \to \forall a\,(\mathrm{Seq}(a) \to \mathrm{Conv}(a,c) \to \varphi_{4}) \to \varphi_{133}$
\item $\varphi_{193} :\equiv \forall a\,(\mathrm{Seq}(a) \to \mathrm{Conv}(a,c) \to \varphi_{4}) \to \neg \forall \varepsilon \in \mathbb{R}\,(0 < \varepsilon \to \varphi_{45}) \to \varphi_{62} \to \mathrm{Seq}(a) \land \varphi_{61} \to \varphi_{52}$
\item $\varphi_{194} :\equiv \mathrm{PNbhd}(f,c) \to \forall a\,(\mathrm{Seq}(a) \to \mathrm{Conv}(a,c) \to \varphi_{4}) \to \varphi_{136}$
\item $\varphi_{195} :\equiv \forall a\,(\mathrm{Seq}(a) \to \mathrm{Conv}(a,c) \to \varphi_{4}) \to \neg \forall \varepsilon \in \mathbb{R}\,(0 < \varepsilon \to \varphi_{45}) \to \varphi_{62} \to \mathrm{Seq}(a) \land \varphi_{61} \to \varphi_{54}$
\item $\varphi_{196} :\equiv \mathrm{PNbhd}(f,c) \to \forall a\,(\mathrm{Seq}(a) \to \mathrm{Conv}(a,c) \to \varphi_{4}) \to \varphi_{137}$
\item $\varphi_{197} :\equiv \mathrm{PNbhd}(f,c) \to \forall a\,(\mathrm{Seq}(a) \to \mathrm{Conv}(a,c) \to \varphi_{4}) \to \varphi_{138}$
\item $\varphi_{198} :\equiv \forall a\,(\mathrm{Seq}(a) \to \mathrm{Conv}(a,c) \to \varphi_{4}) \to \neg \forall \varepsilon \in \mathbb{R}\,(0 < \varepsilon \to \varphi_{45}) \to \varphi_{62} \to \mathrm{Seq}(a) \land \varphi_{61} \to \varphi_{55}$
\item $\varphi_{199} :\equiv \forall a\,(\mathrm{Seq}(a) \to \mathrm{Conv}(a,c) \to \varphi_{4}) \to \neg \forall \varepsilon \in \mathbb{R}\,(0 < \varepsilon \to \varphi_{45}) \to \varphi_{62} \to \mathrm{Seq}(a) \land \varphi_{61} \to \varphi_{58}$
\item $\varphi_{200} :\equiv N = N \to \mathrm{RFn}(f) \to c \in \mathbb{R} \to \alpha \in \mathbb{R} \to \mathrm{PNbhd}(f,c) \to \varphi_{173}$
\item $\varphi_{201} :\equiv \mathrm{PNbhd}(f,c) \to \forall a\,(\mathrm{Seq}(a) \to \mathrm{Conv}(a,c) \to \varphi_{4}) \to \varphi_{140}$
\item $\varphi_{202} :\equiv \forall a\,(\mathrm{Seq}(a) \to \mathrm{Conv}(a,c) \to \varphi_{4}) \to \neg \forall \varepsilon \in \mathbb{R}\,(0 < \varepsilon \to \varphi_{45}) \to \varphi_{62} \to \mathrm{Seq}(a) \land \varphi_{61} \to \varphi_{61}$
\item $\varphi_{203} :\equiv \mathrm{PNbhd}(f,c) \to \forall a\,(\mathrm{Seq}(a) \to \mathrm{Conv}(a,c) \to \varphi_{4}) \to \varphi_{141}$
\item $\varphi_{204} :\equiv \varphi_{106} \land (\mathrm{Seq}(a) \land \varphi_{61}) \land \varepsilon \in \mathbb{R} \land 0 < \varepsilon \land \neg \exists N \in \mathbb{N}\,\forall n \in \mathbb{N}\,(N \leq n \to \lvert a_{n} - c \rvert < \varepsilon)$
\item $\varphi_{205} :\equiv \forall a\,(\mathrm{Seq}(a) \to \mathrm{Conv}(a,c) \to \varphi_{4}) \to \neg \forall \varepsilon \in \mathbb{R}\,(0 < \varepsilon \to \varphi_{45}) \to \varphi_{62} \to \mathrm{Seq}(a) \land \varphi_{61} \to \varphi_{63}$
\item $\varphi_{206} :\equiv \mathrm{PNbhd}(f,c) \to \forall a\,(\mathrm{Seq}(a) \to \mathrm{Conv}(a,c) \to \varphi_{4}) \to \varphi_{144}$
\item $\varphi_{207} :\equiv \mathrm{PNbhd}(f,c) \to \forall a\,(\mathrm{Seq}(a) \to \mathrm{Conv}(a,c) \to \varphi_{4}) \to \varphi_{145}$
\item $\varphi_{208} :\equiv (\mathrm{RFn}(f) \to c \in \mathbb{R} \to \alpha \in \mathbb{R} \to \mathrm{PNbhd}(f,c) \to \varphi_{70}) \to \varphi_{108}$
\item $\varphi_{209} :\equiv \mathrm{PNbhd}(f,c) \to \forall a\,(\mathrm{Seq}(a) \to \mathrm{Conv}(a,c) \to \varphi_{4}) \to \varphi_{154}$
\item $\varphi_{210} :\equiv \forall a\,(\mathrm{Seq}(a) \to \mathrm{Conv}(a,c) \to \varphi_{4}) \to \neg \forall \varepsilon \in \mathbb{R}\,(0 < \varepsilon \to \varphi_{45}) \to \varphi_{62} \to \mathrm{Seq}(a) \land \varphi_{61} \to \varphi_{68}$
\item $\varphi_{211} :\equiv \forall a\,(\mathrm{Seq}(a) \to \mathrm{Conv}(a,c) \to \varphi_{4}) \to \neg \forall \varepsilon \in \mathbb{R}\,(0 < \varepsilon \to \varphi_{45}) \to \varphi_{62} \to \mathrm{Seq}(a) \land \varphi_{61} \to \varphi_{71}$
\item $\varphi_{212} :\equiv (\forall x\,(\neg \neg \varphi_{24}) \to \forall x\,\forall x\,(\neg \neg \varphi_{24})) \to \mathrm{RFn}(f) \to c \in \mathbb{R} \to \alpha \in \mathbb{R} \to \mathrm{PNbhd}(f,c) \to \varphi_{123}$
\item $\varphi_{213} :\equiv \varphi_{88} \to \mathrm{RFn}(f) \to c \in \mathbb{R} \to \alpha \in \mathbb{R} \to \mathrm{PNbhd}(f,c) \to \varphi_{93}$
\item $\varphi_{214} :\equiv \varphi_{43} \to \mathrm{RFn}(f) \to c \in \mathbb{R} \to \alpha \in \mathbb{R} \to \mathrm{PNbhd}(f,c) \to \varphi_{186}$
\item $\varphi_{215} :\equiv \varphi_{35} \to \mathrm{RFn}(f) \to c \in \mathbb{R} \to \alpha \in \mathbb{R} \to \mathrm{PNbhd}(f,c) \to \varphi_{193}$
\item $\varphi_{216} :\equiv \varphi_{48} \to \mathrm{RFn}(f) \to c \in \mathbb{R} \to \alpha \in \mathbb{R} \to \mathrm{PNbhd}(f,c) \to \varphi_{198}$
\item $\varphi_{217} :\equiv (\varphi_{44} \to \neg (\underline{k}^{-1} \in \mathbb{R} \land 0 < \underline{k}^{-1} \land \varphi_{29})) \to \mathrm{RFn}(f) \to c \in \mathbb{R} \to \alpha \in \mathbb{R} \to \varphi_{134}$
\item $\varphi_{218} :\equiv (\varphi_{2} \to \neg (K \in \mathbb{N} \land \forall n \in \mathbb{N}\,(K \leq n \to \lvert a_{n} - c \rvert < \varepsilon))) \to \mathrm{RFn}(f) \to c \in \mathbb{R} \to \alpha \in \mathbb{R} \to \varphi_{203}$
\item $\varphi_{219} :\equiv (\forall x\,(\neg \varphi_{37}) \to \neg \varphi_{37}) \to \mathrm{RFn}(f) \to c \in \mathbb{R} \to \alpha \in \mathbb{R} \to \mathrm{PNbhd}(f,c) \to \varphi_{168}$
\item $\varphi_{220} :\equiv (\mathrm{RFn}(f) \to c \in \mathbb{R} \to \alpha \in \mathbb{R} \to \mathrm{PNbhd}(f,c) \to \varphi_{93}) \to \varphi_{89} \to \neg \varphi_{62}$
\item $\varphi_{221} :\equiv (\mathrm{RFn}(f) \to c \in \mathbb{R} \to \alpha \in \mathbb{R} \to \mathrm{PNbhd}(f,c) \to \varphi_{79}) \to \varphi_{167}$
\item $\varphi_{222} :\equiv (\mathrm{RFn}(f) \to c \in \mathbb{R} \to \alpha \in \mathbb{R} \to \mathrm{PNbhd}(f,c) \to \varphi_{98}) \to \mathrm{RFn}(f) \to c \in \mathbb{R} \to \alpha \in \mathbb{R} \to \mathrm{PNbhd}(f,c) \to \varphi_{101}$
\item $\varphi_{223} :\equiv (\mathrm{RFn}(f) \to c \in \mathbb{R} \to \alpha \in \mathbb{R} \to \mathrm{PNbhd}(f,c) \to \varphi_{101}) \to \mathrm{RFn}(f) \to c \in \mathbb{R} \to \alpha \in \mathbb{R} \to \mathrm{PNbhd}(f,c) \to \varphi_{99}$
\item $\varphi_{224} :\equiv (\mathrm{RFn}(f) \to c \in \mathbb{R} \to \alpha \in \mathbb{R} \to \mathrm{PNbhd}(f,c) \to \varphi_{98}) \to \mathrm{RFn}(f) \to c \in \mathbb{R} \to \alpha \in \mathbb{R} \to \mathrm{PNbhd}(f,c) \to \varphi_{102}$
\item $\varphi_{225} :\equiv (\mathrm{RFn}(f) \to c \in \mathbb{R} \to \alpha \in \mathbb{R} \to \mathrm{PNbhd}(f,c) \to \varphi_{101}) \to \mathrm{RFn}(f) \to c \in \mathbb{R} \to \alpha \in \mathbb{R} \to \mathrm{PNbhd}(f,c) \to \varphi_{103}$
\item $\varphi_{226} :\equiv (\mathrm{RFn}(f) \to c \in \mathbb{R} \to \alpha \in \mathbb{R} \to \mathrm{PNbhd}(f,c) \to \varphi_{100}) \to \mathrm{RFn}(f) \to c \in \mathbb{R} \to \alpha \in \mathbb{R} \to \mathrm{PNbhd}(f,c) \to \varphi_{104}$
\item $\varphi_{227} :\equiv (\mathrm{RFn}(f) \to c \in \mathbb{R} \to \alpha \in \mathbb{R} \to \mathrm{PNbhd}(f,c) \to \varphi_{99}) \to \varphi_{121}$
\item $\varphi_{228} :\equiv (\mathrm{RFn}(f) \to c \in \mathbb{R} \to \alpha \in \mathbb{R} \to \mathrm{PNbhd}(f,c) \to \varphi_{103}) \to \mathrm{RFn}(f) \to c \in \mathbb{R} \to \alpha \in \mathbb{R} \to \mathrm{PNbhd}(f,c) \to \varphi_{122}$
\item $\varphi_{229} :\equiv (\mathrm{RFn}(f) \to c \in \mathbb{R} \to \alpha \in \mathbb{R} \to \mathrm{PNbhd}(f,c) \to \varphi_{129}) \to \mathrm{RFn}(f) \to c \in \mathbb{R} \to \alpha \in \mathbb{R} \to \mathrm{PNbhd}(f,c) \to \varphi_{94}$
\item $\varphi_{230} :\equiv (\mathrm{RFn}(f) \to c \in \mathbb{R} \to \alpha \in \mathbb{R} \to \mathrm{PNbhd}(f,c) \to \varphi_{129}) \to \mathrm{RFn}(f) \to c \in \mathbb{R} \to \alpha \in \mathbb{R} \to \mathrm{PNbhd}(f,c) \to \varphi_{95}$
\item $\varphi_{231} :\equiv (\mathrm{RFn}(f) \to c \in \mathbb{R} \to \alpha \in \mathbb{R} \to \mathrm{PNbhd}(f,c) \to \varphi_{94}) \to \mathrm{RFn}(f) \to c \in \mathbb{R} \to \alpha \in \mathbb{R} \to \mathrm{PNbhd}(f,c) \to \varphi_{172}$
\item $\varphi_{232} :\equiv (\mathrm{RFn}(f) \to c \in \mathbb{R} \to \alpha \in \mathbb{R} \to \varphi_{128}) \to \mathrm{RFn}(f) \to c \in \mathbb{R} \to \alpha \in \mathbb{R} \to \mathrm{PNbhd}(f,c) \to \varphi_{125}$
\item $\varphi_{233} :\equiv (\mathrm{RFn}(f) \to c \in \mathbb{R} \to \alpha \in \mathbb{R} \to \mathrm{PNbhd}(f,c) \to \varphi_{126}) \to \varphi_{106} \to k \in \mathbb{N} \to \varphi_{39}$
\item $\varphi_{234} :\equiv (\mathrm{RFn}(f) \to c \in \mathbb{R} \to \alpha \in \mathbb{R} \to \mathrm{PNbhd}(f,c) \to \varphi_{129}) \to \mathrm{RFn}(f) \to c \in \mathbb{R} \to \alpha \in \mathbb{R} \to \mathrm{PNbhd}(f,c) \to \varphi_{124}$
\item $\varphi_{235} :\equiv (\mathrm{RFn}(f) \to c \in \mathbb{R} \to \alpha \in \mathbb{R} \to \mathrm{PNbhd}(f,c) \to \varphi_{127}) \to \mathrm{RFn}(f) \to c \in \mathbb{R} \to \alpha \in \mathbb{R} \to \mathrm{PNbhd}(f,c) \to \varphi_{131}$
\item $\varphi_{236} :\equiv (\mathrm{RFn}(f) \to x \in \mathrm{dom}(f) \to x \in \mathbb{R}) \to (\mathrm{RFn}(f) \to \mathrm{RFn}(f)) \to \varphi_{232}$
\item $\varphi_{237} :\equiv (\mathrm{RFn}(f) \to c \in \mathbb{R} \to \alpha \in \mathbb{R} \to \mathrm{PNbhd}(f,c) \to \varphi_{135}) \to \varphi_{106} \land \neg 0 < 0 \to \neg (\mathrm{Seq}(a) \land \varphi_{61})$
\item $\varphi_{238} :\equiv \varphi_{106} \land (\mathrm{Seq}(a) \land \varphi_{61}) \land \neg 0 < 0 \to \forall N\,(\varphi_{106} \land (\mathrm{Seq}(a) \land \varphi_{61}) \land \neg 0 < 0)$
\item $\varphi_{239} :\equiv (\mathrm{RFn}(f) \to c \in \mathbb{R} \to \alpha \in \mathbb{R} \to \mathrm{PNbhd}(f,c) \to \varphi_{161}) \to \mathrm{RFn}(f) \to c \in \mathbb{R} \to \alpha \in \mathbb{R} \to \varphi_{128}$
\item $\varphi_{240} :\equiv \forall a\,(\varphi_{106} \land \neg 0 < 0 \to \neg (\mathrm{Seq}(a) \land \varphi_{61})) \to \forall a\,(\varphi_{106} \land \neg 0 < 0) \to \forall a\,(\neg (\mathrm{Seq}(a) \land \varphi_{61}))$
\item $\varphi_{241} :\equiv (\mathrm{RFn}(f) \to c \in \mathbb{R} \to \alpha \in \mathbb{R} \to \mathrm{PNbhd}(f,c) \to \varphi_{125}) \to \mathrm{RFn}(f) \to c \in \mathbb{R} \to \alpha \in \mathbb{R} \to \mathrm{PNbhd}(f,c) \to \varphi_{166}$
\item $\varphi_{242} :\equiv (\mathrm{RFn}(f) \to c \in \mathbb{R} \to \alpha \in \mathbb{R} \to \varphi_{147}) \to \mathrm{RFn}(f) \to c \in \mathbb{R} \to \alpha \in \mathbb{R} \to \varphi_{153}$
\item $\varphi_{243} :\equiv (\mathrm{RFn}(f) \to c \in \mathbb{R} \to \alpha \in \mathbb{R} \to \varphi_{147}) \to \mathrm{RFn}(f) \to c \in \mathbb{R} \to \alpha \in \mathbb{R} \to \varphi_{155}$
\item $\varphi_{244} :\equiv (\mathrm{RFn}(f) \to c \in \mathbb{R} \to \alpha \in \mathbb{R} \to \varphi_{159}) \to \mathrm{RFn}(f) \to c \in \mathbb{R} \to \alpha \in \mathbb{R} \to \varphi_{149}$
\item $\varphi_{245} :\equiv (\mathrm{RFn}(f) \to c \in \mathbb{R} \to \alpha \in \mathbb{R} \to \varphi_{155}) \to \varphi_{106} \land (\mathrm{Seq}(a) \land \varphi_{61}) \to n \in \mathbb{N} \to a_{n} \in \mathrm{dom}(f)$
\item $\varphi_{246} :\equiv (\mathrm{RFn}(f) \to c \in \mathbb{R} \to \alpha \in \mathbb{R} \to \varphi_{149}) \to \mathrm{RFn}(f) \to c \in \mathbb{R} \to \alpha \in \mathbb{R} \to \varphi_{163}$
\item $\varphi_{247} :\equiv (\mathrm{RFn}(f) \to c \in \mathbb{R} \to \alpha \in \mathbb{R} \to \mathrm{PNbhd}(f,c) \to \varphi_{160}) \to \mathrm{RFn}(f) \to c \in \mathbb{R} \to \alpha \in \mathbb{R} \to \varphi_{148}$
\item $\varphi_{248} :\equiv (\mathrm{RFn}(f) \to c \in \mathbb{R} \to \alpha \in \mathbb{R} \to \varphi_{158}) \to \mathrm{RFn}(f) \to c \in \mathbb{R} \to \alpha \in \mathbb{R} \to \mathrm{PNbhd}(f,c) \to \varphi_{164}$
\item $\varphi_{249} :\equiv (\mathrm{RFn}(f) \to c \in \mathbb{R} \to \alpha \in \mathbb{R} \to \mathrm{PNbhd}(f,c) \to \varphi_{183}) \to \mathrm{RFn}(f) \to c \in \mathbb{R} \to \alpha \in \mathbb{R} \to \varphi_{142}$
\item $\varphi_{250} :\equiv (\mathrm{RFn}(f) \to c \in \mathbb{R} \to \alpha \in \mathbb{R} \to \varphi_{157}) \to \mathrm{RFn}(f) \to c \in \mathbb{R} \to \alpha \in \mathbb{R} \to \mathrm{PNbhd}(f,c) \to \varphi_{171}$
\item $\varphi_{251} :\equiv (\mathrm{RFn}(f) \to c \in \mathbb{R} \to \alpha \in \mathbb{R} \to \mathrm{PNbhd}(f,c) \to \varphi_{168}) \to \varphi_{106} \land k \in \mathbb{N} \land \neg \varphi_{39} \to \neg \neg \varphi_{24}$
\item $\varphi_{252} :\equiv \forall x\,(\varphi_{106} \land k \in \mathbb{N} \land \neg \varphi_{39} \to \neg \neg \varphi_{24}) \to \forall x\,(\varphi_{106} \land k \in \mathbb{N} \land \neg \varphi_{39}) \to \forall x\,(\neg \neg \varphi_{24})$
\item $\varphi_{253} :\equiv (\mathrm{RFn}(f) \to c \in \mathbb{R} \to \alpha \in \mathbb{R} \to \varphi_{209}) \to \mathrm{RFn}(f) \to c \in \mathbb{R} \to \alpha \in \mathbb{R} \to \varphi_{139}$
\item $\varphi_{254} :\equiv \forall x\,(\forall x\,(\neg \neg \varphi_{24}) \to \varphi_{24}) \to (\forall x\,(\forall x\,(\neg \neg \varphi_{24}) \to \varphi_{24}) \to \forall x\,\forall x\,(\neg \neg \varphi_{24}) \to \varphi_{29}) \to \varphi_{212}$
\item $\varphi_{255} :\equiv (\mathrm{RFn}(f) \to c \in \mathbb{R} \to \alpha \in \mathbb{R} \to \mathrm{PNbhd}(f,c) \to \varphi_{174}) \to \varphi_{106} \land (\mathrm{Seq}(a) \land \varphi_{61}) \land \neg 0 < 0 \to \varphi_{3}$
\item $\varphi_{256} :\equiv (\mathrm{RFn}(f) \to c \in \mathbb{R} \to \alpha \in \mathbb{R} \to \mathrm{PNbhd}(f,c) \to \varphi_{182}) \to \mathrm{RFn}(f) \to c \in \mathbb{R} \to \alpha \in \mathbb{R} \to \mathrm{PNbhd}(f,c) \to \varphi_{170}$
\item $\varphi_{257} :\equiv (\mathrm{RFn}(f) \to c \in \mathbb{R} \to \alpha \in \mathbb{R} \to \varphi_{207}) \to \mathrm{RFn}(f) \to c \in \mathbb{R} \to \alpha \in \mathbb{R} \to \varphi_{159}$
\item $\varphi_{258} :\equiv \forall N\,(\varphi_{106} \land (\mathrm{Seq}(a) \land \varphi_{61}) \land \neg 0 < 0 \to \varphi_{3}) \to \forall N\,(\varphi_{106} \land (\mathrm{Seq}(a) \land \varphi_{61}) \land \neg 0 < 0) \to \forall N\,(\varphi_{3})$
\item $\varphi_{259} :\equiv (\mathrm{RFn}(f) \to c \in \mathbb{R} \to \alpha \in \mathbb{R} \to \mathrm{PNbhd}(f,c) \to \varphi_{180}) \to \varphi_{106} \land (\mathrm{Seq}(a) \land \varphi_{61}) \to \varphi_{20}$
\item $\varphi_{260} :\equiv (\mathrm{RFn}(f) \to c \in \mathbb{R} \to \alpha \in \mathbb{R} \to \mathrm{PNbhd}(f,c) \to \varphi_{173}) \to \mathrm{RFn}(f) \to c \in \mathbb{R} \to \alpha \in \mathbb{R} \to \mathrm{PNbhd}(f,c) \to \varphi_{189}$
\item $\varphi_{261} :\equiv (\mathrm{RFn}(f) \to c \in \mathbb{R} \to \alpha \in \mathbb{R} \to \mathrm{PNbhd}(f,c) \to \varphi_{176}) \to \mathrm{RFn}(f) \to c \in \mathbb{R} \to \alpha \in \mathbb{R} \to \varphi_{191}$
\item $\varphi_{262} :\equiv (\mathrm{RFn}(f) \to c \in \mathbb{R} \to \alpha \in \mathbb{R} \to \mathrm{PNbhd}(f,c) \to \varphi_{175}) \to \mathrm{RFn}(f) \to c \in \mathbb{R} \to \alpha \in \mathbb{R} \to \mathrm{PNbhd}(f,c) \to \varphi_{188}$
\item $\varphi_{263} :\equiv (\mathrm{RFn}(f) \to c \in \mathbb{R} \to \alpha \in \mathbb{R} \to \mathrm{PNbhd}(f,c) \to \varphi_{182}) \to \mathrm{RFn}(f) \to c \in \mathbb{R} \to \alpha \in \mathbb{R} \to \varphi_{190}$
\item $\varphi_{264} :\equiv \forall \varepsilon\,(\varphi_{106} \land (\mathrm{Seq}(a) \land \varphi_{61}) \to \varphi_{20}) \to \forall \varepsilon\,(\varphi_{106} \land (\mathrm{Seq}(a) \land \varphi_{61})) \to \varphi_{27}$
\item $\varphi_{265} :\equiv (\mathrm{RFn}(f) \to c \in \mathbb{R} \to \alpha \in \mathbb{R} \to \mathrm{PNbhd}(f,c) \to \varphi_{160}) \to \varphi_{33} \to \mathrm{RFn}(f) \to c \in \mathbb{R} \to \alpha \in \mathbb{R} \to \varphi_{163}$
\item $\varphi_{266} :\equiv (\mathrm{RFn}(f) \to c \in \mathbb{R} \to \alpha \in \mathbb{R} \to \mathrm{PNbhd}(f,c) \to \varphi_{171}) \to \mathrm{RFn}(f) \to c \in \mathbb{R} \to \alpha \in \mathbb{R} \to \mathrm{PNbhd}(f,c) \to \varphi_{195}$
\item $\varphi_{267} :\equiv (\mathrm{RFn}(f) \to c \in \mathbb{R} \to \alpha \in \mathbb{R} \to \varphi_{190}) \to \mathrm{RFn}(f) \to c \in \mathbb{R} \to \alpha \in \mathbb{R} \to \varphi_{192}$
\item $\varphi_{268} :\equiv (\mathrm{RFn}(f) \to c \in \mathbb{R} \to \alpha \in \mathbb{R} \to \varphi_{192}) \to \varphi_{187} \to n \in \mathbb{N} \to K \leq n \to \lvert a_{n} - c \rvert < \varepsilon$
\item $\varphi_{269} :\equiv (\mathrm{RFn}(f) \to c \in \mathbb{R} \to \alpha \in \mathbb{R} \to \mathrm{PNbhd}(f,c) \to \varphi_{124}) \to \varphi_{217}$
\item $\varphi_{270} :\equiv \forall n\,(\varphi_{187} \to n \in \mathbb{N} \to K \leq n \to \lvert a_{n} - c \rvert < \varepsilon) \to \forall n\,(\varphi_{187}) \to \forall n \in \mathbb{N}\,(K \leq n \to \lvert a_{n} - c \rvert < \varepsilon)$
\item $\varphi_{271} :\equiv (\mathrm{RFn}(f) \to c \in \mathbb{R} \to \alpha \in \mathbb{R} \to \varphi_{191}) \to \mathrm{RFn}(f) \to c \in \mathbb{R} \to \alpha \in \mathbb{R} \to \varphi_{197}$
\item $\varphi_{272} :\equiv (\mathrm{RFn}(f) \to c \in \mathbb{R} \to \alpha \in \mathbb{R} \to \mathrm{PNbhd}(f,c) \to \varphi_{205}) \to \mathrm{RFn}(f) \to c \in \mathbb{R} \to \alpha \in \mathbb{R} \to \mathrm{PNbhd}(f,c) \to \varphi_{175}$
\item $\varphi_{273} :\equiv (\mathrm{RFn}(f) \to c \in \mathbb{R} \to \alpha \in \mathbb{R} \to \varphi_{194}) \to \mathrm{RFn}(f) \to c \in \mathbb{R} \to \alpha \in \mathbb{R} \to \varphi_{201}$
\item $\varphi_{274} :\equiv (\mathrm{RFn}(f) \to c \in \mathbb{R} \to \alpha \in \mathbb{R} \to \mathrm{PNbhd}(f,c) \to \varphi_{102}) \to \varphi_{225}$
\item $\varphi_{275} :\equiv (\mathrm{RFn}(f) \to c \in \mathbb{R} \to \alpha \in \mathbb{R} \to \mathrm{PNbhd}(f,c) \to \varphi_{102}) \to \varphi_{226}$
\item $\varphi_{276} :\equiv (\mathrm{RFn}(f) \to c \in \mathbb{R} \to \alpha \in \mathbb{R} \to \varphi_{203}) \to \varphi_{204} \to \neg (K \in \mathbb{N} \land \underline{K}^{-1} < \varepsilon)$
\item $\varphi_{277} :\equiv \forall K\,(\varphi_{204} \to \neg (K \in \mathbb{N} \land \underline{K}^{-1} < \varepsilon)) \to \forall K\,(\varphi_{204}) \to \forall K\,(\neg (K \in \mathbb{N} \land \underline{K}^{-1} < \varepsilon))$
\item $\varphi_{278} :\equiv (\varphi_{106} \to \forall k\,(\varphi_{106})) \to \mathrm{RFn}(f) \to c \in \mathbb{R} \to \alpha \in \mathbb{R} \to \mathrm{PNbhd}(f,c) \to \varphi_{127}$
\item $\varphi_{279} :\equiv (\mathrm{RFn}(f) \to c \in \mathbb{R} \to \alpha \in \mathbb{R} \to \mathrm{PNbhd}(f,c) \to \varphi_{210}) \to \mathrm{RFn}(f) \to c \in \mathbb{R} \to \alpha \in \mathbb{R} \to \mathrm{PNbhd}(f,c) \to \varphi_{174}$
\item $\varphi_{280} :\equiv (\mathrm{RFn}(f) \to c \in \mathbb{R} \to \alpha \in \mathbb{R} \to \mathrm{PNbhd}(f,c) \to \varphi_{205}) \to \mathrm{RFn}(f) \to c \in \mathbb{R} \to \alpha \in \mathbb{R} \to \mathrm{PNbhd}(f,c) \to \varphi_{199}$
\item $\varphi_{281} :\equiv (\mathrm{RFn}(f) \to c \in \mathbb{R} \to \alpha \in \mathbb{R} \to \mathrm{PNbhd}(f,c) \to \varphi_{175}) \to \mathrm{RFn}(f) \to c \in \mathbb{R} \to \alpha \in \mathbb{R} \to \mathrm{PNbhd}(f,c) \to \varphi_{210}$
\item $\varphi_{282} :\equiv (\mathrm{RFn}(f) \to c \in \mathbb{R} \to \alpha \in \mathbb{R} \to \varphi_{209}) \to \mathrm{RFn}(f) \to c \in \mathbb{R} \to \alpha \in \mathbb{R} \to \mathrm{PNbhd}(f,c) \to \varphi_{202}$
\item $\varphi_{283} :\equiv (\mathrm{RFn}(f) \to c \in \mathbb{R} \to \alpha \in \mathbb{R} \to \mathrm{PNbhd}(f,c) \to \varphi_{177}) \to \mathrm{RFn}(f) \to c \in \mathbb{R} \to \alpha \in \mathbb{R} \to \mathrm{PNbhd}(f,c) \to \varphi_{211}$
\item $\varphi_{284} :\equiv (\mathrm{RFn}(f) \to c \in \mathbb{R} \to \alpha \in \mathbb{R} \to \mathrm{PNbhd}(f,c) \to \varphi_{211}) \to \mathrm{RFn}(f) \to c \in \mathbb{R} \to \alpha \in \mathbb{R} \to \varphi_{196}$
\item $\varphi_{285} :\equiv (\mathrm{RFn}(f) \to c \in \mathbb{R} \to \alpha \in \mathbb{R} \to \mathrm{PNbhd}(f,c) \to \varphi_{104}) \to \varphi_{228}$
\item $\varphi_{286} :\equiv (\mathrm{RFn}(f) \to c \in \mathbb{R} \to \alpha \in \mathbb{R} \to \mathrm{PNbhd}(f,c) \to \varphi_{189}) \to \varphi_{214}$
\item $\varphi_{287} :\equiv (\mathrm{RFn}(f) \to c \in \mathbb{R} \to \alpha \in \mathbb{R} \to \mathrm{PNbhd}(f,c) \to \varphi_{195}) \to \varphi_{215}$
\item $\varphi_{288} :\equiv (\mathrm{RFn}(f) \to c \in \mathbb{R} \to \alpha \in \mathbb{R} \to \mathrm{PNbhd}(f,c) \to \varphi_{166}) \to \varphi_{219}$
\item $\varphi_{289} :\equiv (\mathrm{RFn}(f) \to c \in \mathbb{R} \to \alpha \in \mathbb{R} \to \mathrm{PNbhd}(f,c) \to \varphi_{199}) \to \varphi_{216}$
\item $\varphi_{290} :\equiv (\mathrm{RFn}(f) \to c \in \mathbb{R} \to \alpha \in \mathbb{R} \to \mathrm{PNbhd}(f,c) \to \varphi_{95}) \to \varphi_{231}$
\item $\varphi_{291} :\equiv (\varphi_{106} \land k \in \mathbb{N} \land \neg \varphi_{39} \to \forall x\,(\varphi_{106} \land k \in \mathbb{N} \land \neg \varphi_{39})) \to \mathrm{RFn}(f) \to c \in \mathbb{R} \to \alpha \in \mathbb{R} \to \mathrm{PNbhd}(f,c) \to \varphi_{126}$
\item $\varphi_{292} :\equiv (\mathrm{RFn}(f) \to c \in \mathbb{R} \to \alpha \in \mathbb{R} \to \mathrm{PNbhd}(f,c) \to \varphi_{202}) \to (\varphi_{61} \to N \in \mathbb{N} \to \varphi_{50}) \to \mathrm{RFn}(f) \to c \in \mathbb{R} \to \alpha \in \mathbb{R} \to \varphi_{206}$
\item $\varphi_{293} :\equiv (\mathrm{RFn}(f) \to c \in \mathbb{R} \to \alpha \in \mathbb{R} \to \mathrm{PNbhd}(f,c) \to \varphi_{202}) \to (\varphi_{61} \to n \in \mathbb{N} \to \varphi_{51}) \to \mathrm{RFn}(f) \to c \in \mathbb{R} \to \alpha \in \mathbb{R} \to \varphi_{207}$
\item $\varphi_{294} :\equiv \varphi_{238} \to \mathrm{RFn}(f) \to c \in \mathbb{R} \to \alpha \in \mathbb{R} \to \mathrm{PNbhd}(f,c) \to \varphi_{135}$
\item $\varphi_{295} :\equiv (\varphi_{106} \land (\mathrm{Seq}(a) \land \varphi_{61}) \to \forall n\,(\varphi_{106} \land (\mathrm{Seq}(a) \land \varphi_{61}))) \to \mathrm{RFn}(f) \to c \in \mathbb{R} \to \alpha \in \mathbb{R} \to \varphi_{157}$
\item $\varphi_{296} :\equiv (\varphi_{106} \land (\mathrm{Seq}(a) \land \varphi_{61}) \to \forall n\,(\varphi_{106} \land (\mathrm{Seq}(a) \land \varphi_{61}))) \to \mathrm{RFn}(f) \to c \in \mathbb{R} \to \alpha \in \mathbb{R} \to \mathrm{PNbhd}(f,c) \to \varphi_{160}$
\item $\varphi_{297} :\equiv (\mathrm{RFn}(f) \to c \in \mathbb{R} \to \alpha \in \mathbb{R} \to \mathrm{PNbhd}(f,c) \to \varphi_{161}) \to \varphi_{241}$
\item $\varphi_{298} :\equiv (\varphi_{106} \land (\mathrm{Seq}(a) \land \varphi_{61}) \to \forall \varepsilon\,(\varphi_{106} \land (\mathrm{Seq}(a) \land \varphi_{61}))) \to \mathrm{RFn}(f) \to c \in \mathbb{R} \to \alpha \in \mathbb{R} \to \mathrm{PNbhd}(f,c) \to \varphi_{183}$
\item $\varphi_{299} :\equiv (\mathrm{Seq}(a) \to n \in \mathbb{N} \to a_{n} \in \mathbb{R}) \to (\mathrm{RFn}(f) \to c \in \mathbb{R} \to \alpha \in \mathbb{R} \to \varphi_{139}) \to \varphi_{242}$
\item $\varphi_{300} :\equiv (\mathrm{RFn}(f) \to c \in \mathbb{R} \to \alpha \in \mathbb{R} \to \mathrm{PNbhd}(f,c) \to \varphi_{122}) \to \varphi_{254}$
\item $\varphi_{301} :\equiv (\mathrm{RFn}(f) \to c \in \mathbb{R} \to \alpha \in \mathbb{R} \to \varphi_{153}) \to (\mathrm{RFn}(f) \to c \in \mathbb{R} \to c \in \mathbb{R}) \to \varphi_{244}$
\item $\varphi_{302} :\equiv (\mathrm{RFn}(f) \to \mathrm{RFn}(f)) \to (\mathrm{RFn}(f) \to c \in \mathbb{R} \to \alpha \in \mathbb{R} \to \varphi_{139}) \to \varphi_{250}$
\item $\varphi_{303} :\equiv (\mathrm{RFn}(f) \to c \in \mathbb{R} \to \alpha \in \mathbb{R} \to \varphi_{139}) \to (\mathrm{RFn}(f) \to c \in \mathbb{R} \to c \in \mathbb{R}) \to \varphi_{249}$
\item $\varphi_{304} :\equiv (\varepsilon \in \mathbb{R} \to 0 < \varepsilon \to \exists K \in \mathbb{N}\,(\underline{K}^{-1} < \varepsilon)) \to (\mathrm{RFn}(f) \to c \in \mathbb{R} \to \alpha \in \mathbb{R} \to \varphi_{152}) \to \varphi_{248}$
\item $\varphi_{305} :\equiv (\mathrm{RFn}(f) \to c \in \mathbb{R} \to \alpha \in \mathbb{R} \to \mathrm{PNbhd}(f,c) \to \varphi_{175}) \to \varphi_{260}$
\item $\varphi_{306} :\equiv (\mathrm{RFn}(f) \to c \in \mathbb{R} \to \alpha \in \mathbb{R} \to \mathrm{PNbhd}(f,c) \to \varphi_{177}) \to \varphi_{261}$
\item $\varphi_{307} :\equiv (\varphi_{204} \to \forall K\,(\varphi_{204})) \to \mathrm{RFn}(f) \to c \in \mathbb{R} \to \alpha \in \mathbb{R} \to \mathrm{PNbhd}(f,c) \to \varphi_{180}$
\item $\varphi_{308} :\equiv (\mathrm{RFn}(f) \to c \in \mathbb{R} \to \alpha \in \mathbb{R} \to \mathrm{PNbhd}(f,c) \to \varphi_{170}) \to \varphi_{273}$
\item $\varphi_{309} :\equiv (\mathrm{RFn}(f) \to c \in \mathbb{R} \to \alpha \in \mathbb{R} \to \mathrm{PNbhd}(f,c) \to \varphi_{193}) \to \varphi_{262}$
\item $\varphi_{310} :\equiv (\mathrm{RFn}(f) \to c \in \mathbb{R} \to \alpha \in \mathbb{R} \to \mathrm{PNbhd}(f,c) \to \varphi_{92}) \to \varphi_{84} \to (\varphi_{84} \to \varphi_{86}) \to \varphi_{87} \to \mathrm{RFn}(f) \to c \in \mathbb{R} \to \alpha \in \mathbb{R} \to \varphi_{97}$
\item $\varphi_{311} :\equiv (\varphi_{89} \to \forall e_{0}\,(\varphi_{89})) \to \forall \varepsilon\,(\forall e_{0}\,(\neg \varphi_{62}) \to \neg \neg (\varepsilon \in \mathbb{R} \to 0 < \varepsilon \to \varphi_{45})) \to \varphi_{162} \to \varphi_{165}$
\item $\varphi_{312} :\equiv (\mathrm{RFn}(f) \to c \in \mathbb{R} \to \alpha \in \mathbb{R} \to \mathrm{PNbhd}(f,c) \to \varphi_{186}) \to \varphi_{279}$
\item $\varphi_{313} :\equiv (\lvert a_{n} - c \rvert < \underline{K}^{-1} \to \underline{K}^{-1} < \varepsilon \to \lvert a_{n} - c \rvert < \varepsilon) \to (\mathrm{RFn}(f) \to c \in \mathbb{R} \to \alpha \in \mathbb{R} \to \varphi_{197}) \to \varphi_{267}$
\item $\varphi_{314} :\equiv (\mathrm{RFn}(f) \to c \in \mathbb{R} \to \alpha \in \mathbb{R} \to \mathrm{PNbhd}(f,c) \to \varphi_{188}) \to \varphi_{286}$
\item $\varphi_{315} :\equiv (\mathrm{Conv}(f \circ a,\alpha) \to e_{0} \in \mathbb{R} \to \varphi_{15}) \to (\mathrm{RFn}(f) \to c \in \mathbb{R} \to \alpha \in \mathbb{R} \to \varphi_{148}) \to \varphi_{290}$
\item $\varphi_{316} :\equiv (\mathrm{RFn}(f) \to c \in \mathbb{R} \to \alpha \in \mathbb{R} \to \varphi_{139}) \to (\mathrm{RFn}(f) \to c \in \mathbb{R} \to \alpha \in \mathbb{R} \to \varphi_{142}) \to \varphi_{247}$
\item $\varphi_{317} :\equiv (\varphi_{106} \land \neg 0 < 0 \to \forall a\,(\varphi_{106} \land \neg 0 < 0)) \to \forall q\,(\forall a\,(\neg (\mathrm{Seq}(a) \land \varphi_{61})) \to \varphi_{64}) \to \varphi_{169} \to \varphi_{213}$
\item $\varphi_{318} :\equiv (\forall k\,(\varphi_{106} \to k \in \mathbb{N} \to \varphi_{39}) \to \forall k\,(\varphi_{106}) \to \forall k \in \mathbb{N}\,(\varphi_{39})) \to \varphi_{278}$
\item $\varphi_{319} :\equiv (\mathrm{RFn}(f) \to c \in \mathbb{R} \to \alpha \in \mathbb{R} \to \mathrm{PNbhd}(f,c) \to \varphi_{69}) \to \varphi_{316}$
\item $\varphi_{320} :\equiv (\mathrm{RFn}(f) \to c \in \mathbb{R} \to \alpha \in \mathbb{R} \to \mathrm{PNbhd}(f,c) \to \varphi_{170}) \to \varphi_{306}$
\item $\varphi_{321} :\equiv (\mathrm{RFn}(f) \to c \in \mathbb{R} \to \alpha \in \mathbb{R} \to \mathrm{PNbhd}(f,c) \to \varphi_{198}) \to \varphi_{305}$
\item $\varphi_{322} :\equiv \forall e_{0}\,(\varphi_{89} \to \neg \varphi_{62}) \to (\forall e_{0}\,(\varphi_{89} \to \neg \varphi_{62}) \to \forall e_{0}\,(\varphi_{89}) \to \forall e_{0}\,(\neg \varphi_{62})) \to \varphi_{311}$
\item $\varphi_{323} :\equiv \forall n\,(\varphi_{187} \to n \in \mathbb{N} \to K \leq n \to \lvert a_{n} - c \rvert < \varepsilon) \to \varphi_{270} \to (\varphi_{187} \to \forall n\,(\varphi_{187})) \to \mathrm{RFn}(f) \to c \in \mathbb{R} \to \alpha \in \mathbb{R} \to \varphi_{194}$
\item $\varphi_{324} :\equiv (\mathrm{RFn}(f) \to c \in \mathbb{R} \to \alpha \in \mathbb{R} \to \mathrm{PNbhd}(f,c) \to \varphi_{123}) \to \forall x\,(\varphi_{106} \land k \in \mathbb{N} \land \neg \varphi_{39} \to \neg \neg \varphi_{24}) \to \varphi_{252} \to \varphi_{291}$
\item $\varphi_{325} :\equiv (\mathrm{RFn}(f) \to c \in \mathbb{R} \to \alpha \in \mathbb{R} \to \mathrm{PNbhd}(f,c) \to \varphi_{172}) \to \forall N\,(\varphi_{106} \land (\mathrm{Seq}(a) \land \varphi_{61}) \land \neg 0 < 0 \to \varphi_{3}) \to \varphi_{258} \to \varphi_{294}$
\item $\varphi_{326} :\equiv (\mathrm{RFn}(f) \to c \in \mathbb{R} \to \alpha \in \mathbb{R} \to \mathrm{PNbhd}(f,c) \to \varphi_{164}) \to \forall K\,(\varphi_{204} \to \neg (K \in \mathbb{N} \land \underline{K}^{-1} < \varepsilon)) \to \varphi_{277} \to \varphi_{307}$
\item $\varphi_{327} :\equiv (\mathrm{RFn}(f) \to c \in \mathbb{R} \to \alpha \in \mathbb{R} \to \mathrm{PNbhd}(f,c) \to \varphi_{131}) \to \forall a\,(\varphi_{106} \land \neg 0 < 0 \to \neg (\mathrm{Seq}(a) \land \varphi_{61})) \to \varphi_{240} \to \varphi_{317}$
\end{abbrevlist}
\end{fullproof}

% 自動生成 (prover/build.js): FLimSeq  ノード数 66
\begin{fullproof}{定理1.7}{$\forall f\,\forall c\,\forall \alpha\,(\mathrm{FLim}(f,c,\alpha) \leftrightarrow \mathrm{RFn}(f) \land c \in \mathbb{R} \land \alpha \in \mathbb{R} \land \mathrm{PNbhd}(f,c) \land \forall a\,(\mathrm{Seq}(a) \to \mathrm{Conv}(a,c) \to \forall n \in \mathbb{N}\,(a_{n} \in \mathrm{dom}(f) \land a_{n} \neq c) \to \mathrm{Conv}(f \circ a,\alpha)))$}{66}\phantomsection\label{tree:FLimSeq}
\noindent$\Pi_{1}$:
\begin{ptree}
\AXC{$\forall f\,\forall c\,(\varphi_{11})$\lbl{定理1.7(⇐)}}
\AXC{$\forall f\,\forall c\,(\varphi_{11}) \to \forall c\,(\varphi_{11})$\lbl{A2}}
\BIC{$\forall c\,(\varphi_{11})$}
\AXC{$\forall c\,(\varphi_{11}) \to \varphi_{11}$\lbl{A2}}
\BIC{$\varphi_{11}$}
\AXC{$\varphi_{18}$\lbl{A2}}
\BIC{$\mathrm{RFn}(f) \to c \in \mathbb{R} \to \alpha \in \mathbb{R} \to \mathrm{PNbhd}(f,c) \to \varphi_{7}$}
\AXC{$\varphi_{23}$\lbl{A1}}
\BIC{$(\varphi_{10} \to \mathrm{RFn}(f)) \to \varphi_{22}$}
\end{ptree}
\noindent$\Pi_{2}$:
\begin{ptree}
\AXC{$\varphi_{10} \to \varphi_{10}$\lbl{A1}}
\AXC{$(\varphi_{10} \to \varphi_{10}) \to \varphi_{10} \to c \in \mathbb{R}$\lbl{A1}}
\BIC{$\varphi_{10} \to c \in \mathbb{R}$}
\AXC{$\varphi_{10} \to \varphi_{10}$\lbl{A1}}
\AXC{$(\varphi_{10} \to \varphi_{10}) \to \varphi_{10} \to \mathrm{RFn}(f)$\lbl{A1}}
\BIC{$\varphi_{10} \to \mathrm{RFn}(f)$}
\AXC{$\stackrel{\Pi_{1}}{\vdots}$}\noLine
\UIC{$(\varphi_{10} \to \mathrm{RFn}(f)) \to \varphi_{22}$}
\BIC{$\varphi_{22}$}
\BIC{$(\varphi_{10} \to \alpha \in \mathbb{R}) \to (\varphi_{10} \to \mathrm{PNbhd}(f,c)) \to \varphi_{19}$}
\end{ptree}
\noindent$\Pi_{3}$:
\begin{ptree}
\AXC{$\varphi_{10} \to \varphi_{10}$\lbl{A1}}
\AXC{$(\varphi_{10} \to \varphi_{10}) \to \varphi_{10} \to \alpha \in \mathbb{R}$\lbl{A1}}
\BIC{$\varphi_{10} \to \alpha \in \mathbb{R}$}
\AXC{$\stackrel{\Pi_{2}}{\vdots}$}\noLine
\UIC{$(\varphi_{10} \to \alpha \in \mathbb{R}) \to (\varphi_{10} \to \mathrm{PNbhd}(f,c)) \to \varphi_{19}$}
\BIC{$(\varphi_{10} \to \mathrm{PNbhd}(f,c)) \to \varphi_{19}$}
\end{ptree}
\noindent$\Pi_{4}$:
\begin{ptree}
\AXC{$\varphi_{10} \to \varphi_{10}$\lbl{A1}}
\AXC{$(\varphi_{10} \to \varphi_{10}) \to \varphi_{10} \to \mathrm{PNbhd}(f,c)$\lbl{A1}}
\BIC{$\varphi_{10} \to \mathrm{PNbhd}(f,c)$}
\AXC{$\stackrel{\Pi_{3}}{\vdots}$}\noLine
\UIC{$(\varphi_{10} \to \mathrm{PNbhd}(f,c)) \to \varphi_{19}$}
\BIC{$\varphi_{19}$}
\end{ptree}
\noindent$\Pi_{5}$:
\begin{ptree}
\AXC{$\varphi_{10} \to \varphi_{10}$\lbl{A1}}
\AXC{$(\varphi_{10} \to \varphi_{10}) \to \varphi_{10} \to \forall a\,(\mathrm{Seq}(a) \to \mathrm{Conv}(a,c) \to \varphi_{2})$\lbl{A1}}
\BIC{$\varphi_{10} \to \forall a\,(\mathrm{Seq}(a) \to \mathrm{Conv}(a,c) \to \varphi_{2})$}
\AXC{$\stackrel{\Pi_{4}}{\vdots}$}\noLine
\UIC{$\varphi_{19}$}
\BIC{$\varphi_{10} \to \mathrm{FLim}(f,c,\alpha)$}
\end{ptree}
\noindent$\Pi_{6}$:
\begin{ptree}
\AXC{$\forall f\,(\varphi_{4})$\lbl{定理1.7(⇒)}}
\AXC{$\forall f\,(\varphi_{4}) \to \varphi_{4}$\lbl{A2}}
\BIC{$\varphi_{4}$}
\AXC{$\varphi_{15}$\lbl{A2}}
\BIC{$\forall \alpha\,\forall a\,(\mathrm{FLim}(f,c,\alpha) \to \mathrm{Conv}(a,c) \to \varphi_{2})$}
\AXC{$\varphi_{14}$\lbl{A2}}
\BIC{$\forall a\,(\mathrm{FLim}(f,c,\alpha) \to \mathrm{Conv}(a,c) \to \varphi_{2})$}
\AXC{$\varphi_{12}$\lbl{A2}}
\BIC{$\mathrm{FLim}(f,c,\alpha) \to \mathrm{Conv}(a,c) \to \varphi_{2}$}
\AXC{$\varphi_{20}$\lbl{A1}}
\BIC{$(\mathrm{FLim}(f,c,\alpha) \to \mathrm{FLim}(f,c,\alpha)) \to \varphi_{1} \to \varphi_{16}$}
\end{ptree}
\noindent$\Pi_{7}$:
\begin{ptree}
\AXC{$\varphi_{1}$\lbl{A1}}
\AXC{$\mathrm{FLim}(f,c,\alpha) \to \mathrm{FLim}(f,c,\alpha)$\lbl{A1}}
\AXC{$\stackrel{\Pi_{6}}{\vdots}$}\noLine
\UIC{$(\mathrm{FLim}(f,c,\alpha) \to \mathrm{FLim}(f,c,\alpha)) \to \varphi_{1} \to \varphi_{16}$}
\BIC{$\varphi_{1} \to \varphi_{16}$}
\BIC{$\varphi_{16}$}
\end{ptree}
\noindent$\Pi_{8}$:
\begin{ptree}
\AXC{$\varphi_{5} \to \varphi_{8}$\lbl{A3}}
\AXC{$\mathrm{FLim}(f,c,\alpha) \to \mathrm{Seq}(a) \to \mathrm{Conv}(a,c) \to \varphi_{3}$\lbl{A1}}
\AXC{$\stackrel{\Pi_{7}}{\vdots}$}\noLine
\UIC{$\varphi_{16}$}
\BIC{$\mathrm{FLim}(f,c,\alpha) \to \mathrm{Seq}(a) \to \mathrm{Conv}(a,c) \to \varphi_{2}$}
\UIC{$\varphi_{5}$}
\AXC{$\varphi_{5} \to (\varphi_{5} \to \varphi_{8}) \to \varphi_{9}$\lbl{A1}}
\BIC{$(\varphi_{5} \to \varphi_{8}) \to \varphi_{9}$}
\BIC{$\varphi_{9}$}
\end{ptree}
\noindent$\Pi_{9}$:
\begin{ptree}
\AXC{$\forall f\,\forall c\,\forall \alpha\,(\varphi_{13})$\lbl{DefFLim}}
\AXC{$\forall f\,\forall c\,\forall \alpha\,(\varphi_{13}) \to \forall c\,\forall \alpha\,(\varphi_{13})$\lbl{A2}}
\BIC{$\forall c\,\forall \alpha\,(\varphi_{13})$}
\AXC{$\forall c\,\forall \alpha\,(\varphi_{13}) \to \forall \alpha\,(\varphi_{13})$\lbl{A2}}
\BIC{$\forall \alpha\,(\varphi_{13})$}
\AXC{$\forall \alpha\,(\varphi_{13}) \to \varphi_{13}$\lbl{A2}}
\BIC{$\varphi_{13}$}
\end{ptree}
\noindent$\Pi_{10}$:
\begin{ptree}
\AXC{$\stackrel{\Pi_{9}}{\vdots}$}\noLine
\UIC{$\varphi_{13}$}
\AXC{$\mathrm{FLim}(f,c,\alpha) \to \mathrm{FLim}(f,c,\alpha)$\lbl{A1}}
\AXC{$\varphi_{21}$\lbl{A1}}
\BIC{$\varphi_{13} \to \varphi_{6} \to \mathrm{FLim}(f,c,\alpha) \to \varphi_{10}$}
\BIC{$\varphi_{6} \to \mathrm{FLim}(f,c,\alpha) \to \varphi_{10}$}
\end{ptree}
\noindent$\Pi_{11}$:
\begin{ptree}
\AXC{$\mathrm{FLim}(f,c,\alpha) \to \forall a\,(\mathrm{FLim}(f,c,\alpha))$\lbl{A4}}
\AXC{$\stackrel{\Pi_{8}}{\vdots}$}\noLine
\UIC{$\varphi_{9}$}
\BIC{$\varphi_{6}$}
\AXC{$\stackrel{\Pi_{10}}{\vdots}$}\noLine
\UIC{$\varphi_{6} \to \mathrm{FLim}(f,c,\alpha) \to \varphi_{10}$}
\BIC{$\mathrm{FLim}(f,c,\alpha) \to \varphi_{10}$}
\end{ptree}
\noindent 主木:
\begin{ptree}
\AXC{$\stackrel{\Pi_{5}}{\vdots}$}\noLine
\UIC{$\varphi_{10} \to \mathrm{FLim}(f,c,\alpha)$}
\AXC{$\stackrel{\Pi_{11}}{\vdots}$}\noLine
\UIC{$\mathrm{FLim}(f,c,\alpha) \to \varphi_{10}$}
\AXC{$(\mathrm{FLim}(f,c,\alpha) \to \varphi_{10}) \to \varphi_{17}$\lbl{A1}}
\BIC{$\varphi_{17}$}
\BIC{$\mathrm{FLim}(f,c,\alpha) \leftrightarrow \varphi_{10}$}
\UIC{$\forall \alpha\,(\mathrm{FLim}(f,c,\alpha) \leftrightarrow \varphi_{10})$}
\UIC{$\forall c\,\forall \alpha\,(\mathrm{FLim}(f,c,\alpha) \leftrightarrow \varphi_{10})$}
\UIC{$\forall f\,\forall c\,\forall \alpha\,(\mathrm{FLim}(f,c,\alpha) \leftrightarrow \varphi_{10})$}
\end{ptree}
\begin{abbrevlist}
\item $\varphi_{1} :\equiv \mathrm{FLim}(f,c,\alpha) \to \mathrm{Seq}(a) \to \mathrm{Conv}(a,c) \to \mathrm{Conv}(a,c)$
\item $\varphi_{2} :\equiv \forall n \in \mathbb{N}\,(a_{n} \in \mathrm{dom}(f) \land a_{n} \neq c) \to \mathrm{Conv}(f \circ a,\alpha)$
\item $\varphi_{3} :\equiv \forall n \in \mathbb{N}\,(a_{n} \in \mathrm{dom}(f) \land a_{n} \neq c) \to \forall n \in \mathbb{N}\,(a_{n} \in \mathrm{dom}(f) \land a_{n} \neq c)$
\item $\varphi_{4} :\equiv \forall c\,\forall \alpha\,\forall a\,(\mathrm{FLim}(f,c,\alpha) \to \mathrm{Conv}(a,c) \to \varphi_{2})$
\item $\varphi_{5} :\equiv \forall a\,(\mathrm{FLim}(f,c,\alpha) \to \mathrm{Seq}(a) \to \mathrm{Conv}(a,c) \to \varphi_{2})$
\item $\varphi_{6} :\equiv \mathrm{FLim}(f,c,\alpha) \to \forall a\,(\mathrm{Seq}(a) \to \mathrm{Conv}(a,c) \to \varphi_{2})$
\item $\varphi_{7} :\equiv \forall a\,(\mathrm{Seq}(a) \to \mathrm{Conv}(a,c) \to \varphi_{2}) \to \mathrm{FLim}(f,c,\alpha)$
\item $\varphi_{8} :\equiv \forall a\,(\mathrm{FLim}(f,c,\alpha)) \to \forall a\,(\mathrm{Seq}(a) \to \mathrm{Conv}(a,c) \to \varphi_{2})$
\item $\varphi_{9} :\equiv (\mathrm{FLim}(f,c,\alpha) \to \forall a\,(\mathrm{FLim}(f,c,\alpha))) \to \varphi_{6}$
\item $\varphi_{10} :\equiv \mathrm{RFn}(f) \land c \in \mathbb{R} \land \alpha \in \mathbb{R} \land \mathrm{PNbhd}(f,c) \land \forall a\,(\mathrm{Seq}(a) \to \mathrm{Conv}(a,c) \to \varphi_{2})$
\item $\varphi_{11} :\equiv \forall \alpha\,(\mathrm{RFn}(f) \to c \in \mathbb{R} \to \alpha \in \mathbb{R} \to \mathrm{PNbhd}(f,c) \to \varphi_{7})$
\item $\varphi_{12} :\equiv \forall a\,(\mathrm{FLim}(f,c,\alpha) \to \mathrm{Conv}(a,c) \to \varphi_{2}) \to \mathrm{FLim}(f,c,\alpha) \to \mathrm{Conv}(a,c) \to \varphi_{2}$
\item $\varphi_{13} :\equiv \mathrm{FLim}(f,c,\alpha) \leftrightarrow \mathrm{RFn}(f) \land c \in \mathbb{R} \land \alpha \in \mathbb{R} \land \mathrm{PNbhd}(f,c) \land \forall \varepsilon \in \mathbb{R}\,(0 < \varepsilon \to \exists \delta\,(\delta \in \mathbb{R} \land 0 < \delta \land \forall x \in \mathrm{dom}(f)\,(0 < \lvert x - c \rvert \to \lvert x - c \rvert < \delta \to \lvert f(x) - \alpha \rvert < \varepsilon)))$
\item $\varphi_{14} :\equiv \forall \alpha\,\forall a\,(\mathrm{FLim}(f,c,\alpha) \to \mathrm{Conv}(a,c) \to \varphi_{2}) \to \forall a\,(\mathrm{FLim}(f,c,\alpha) \to \mathrm{Conv}(a,c) \to \varphi_{2})$
\item $\varphi_{15} :\equiv \varphi_{4} \to \forall \alpha\,\forall a\,(\mathrm{FLim}(f,c,\alpha) \to \mathrm{Conv}(a,c) \to \varphi_{2})$
\item $\varphi_{16} :\equiv (\mathrm{FLim}(f,c,\alpha) \to \mathrm{Seq}(a) \to \mathrm{Conv}(a,c) \to \varphi_{3}) \to \mathrm{FLim}(f,c,\alpha) \to \mathrm{Seq}(a) \to \mathrm{Conv}(a,c) \to \varphi_{2}$
\item $\varphi_{17} :\equiv (\varphi_{10} \to \mathrm{FLim}(f,c,\alpha)) \to (\mathrm{FLim}(f,c,\alpha) \leftrightarrow \varphi_{10})$
\item $\varphi_{18} :\equiv \varphi_{11} \to \mathrm{RFn}(f) \to c \in \mathbb{R} \to \alpha \in \mathbb{R} \to \mathrm{PNbhd}(f,c) \to \varphi_{7}$
\item $\varphi_{19} :\equiv (\varphi_{10} \to \forall a\,(\mathrm{Seq}(a) \to \mathrm{Conv}(a,c) \to \varphi_{2})) \to \varphi_{10} \to \mathrm{FLim}(f,c,\alpha)$
\item $\varphi_{20} :\equiv (\mathrm{FLim}(f,c,\alpha) \to \mathrm{Conv}(a,c) \to \varphi_{2}) \to (\mathrm{FLim}(f,c,\alpha) \to \mathrm{FLim}(f,c,\alpha)) \to \varphi_{1} \to \varphi_{16}$
\item $\varphi_{21} :\equiv (\mathrm{FLim}(f,c,\alpha) \to \mathrm{FLim}(f,c,\alpha)) \to \varphi_{13} \to \varphi_{6} \to \mathrm{FLim}(f,c,\alpha) \to \varphi_{10}$
\item $\varphi_{22} :\equiv (\varphi_{10} \to c \in \mathbb{R}) \to (\varphi_{10} \to \alpha \in \mathbb{R}) \to (\varphi_{10} \to \mathrm{PNbhd}(f,c)) \to \varphi_{19}$
\item $\varphi_{23} :\equiv (\mathrm{RFn}(f) \to c \in \mathbb{R} \to \alpha \in \mathbb{R} \to \mathrm{PNbhd}(f,c) \to \varphi_{7}) \to (\varphi_{10} \to \mathrm{RFn}(f)) \to \varphi_{22}$
\end{abbrevlist}
\end{fullproof}

% 自動生成 (prover/build.js): PNbhdInter  ノード数 157
\begin{fullproof}{[H2]}{$\forall f\,\forall g\,\forall c\,(\mathrm{PNbhd}(f,c) \to \mathrm{PNbhd}(g,c) \to \exists d\,(d \in \mathbb{R} \land 0 < d \land \forall x \in \mathbb{R}\,(0 < \lvert x - c \rvert \to \lvert x - c \rvert < d \to x \in \mathrm{dom}(f) \land x \in \mathrm{dom}(g))))$}{157}\phantomsection\label{tree:PNbhdInter}
\noindent$\Pi_{1}$:
\begin{ptree}
\AXC{$\mathrm{PNbhd}(f,c) \to \varphi_{60}$\lbl{A1}}
\AXC{$(\mathrm{PNbhd}(f,c) \to \varphi_{60}) \to \mathrm{PNbhd}(f,c) \to \varphi_{64}$\lbl{A1}}
\BIC{$\mathrm{PNbhd}(f,c) \to \varphi_{64}$}
\end{ptree}
\noindent$\Pi_{2}$:
\begin{ptree}
\AXC{$\forall x\,\forall y\,\forall z\,(x < y \to y \leq z \to x < z)$\lbl{[O3]}}
\AXC{$\varphi_{7}$\lbl{A2}}
\BIC{$\forall y\,\forall z\,(\lvert x - c \rvert < y \to y \leq z \to \lvert x - c \rvert < z)$}
\AXC{$\varphi_{21}$\lbl{A2}}
\BIC{$\forall z\,(\lvert x - c \rvert < d \to d \leq z \to \lvert x - c \rvert < z)$}
\AXC{$\varphi_{15}$\lbl{A2}}
\BIC{$\lvert x - c \rvert < d \to d \leq r_{2} \to \lvert x - c \rvert < r_{2}$}
\end{ptree}
\noindent$\Pi_{3}$:
\begin{ptree}
\AXC{$\mathrm{PNbhd}(f,c) \to \varphi_{68}$\lbl{A1}}
\AXC{$\stackrel{\Pi_{2}}{\vdots}$}\noLine
\UIC{$\lvert x - c \rvert < d \to d \leq r_{2} \to \lvert x - c \rvert < r_{2}$}
\AXC{$(\lvert x - c \rvert < d \to d \leq r_{2} \to \lvert x - c \rvert < r_{2}) \to \varphi_{96}$\lbl{A1}}
\BIC{$\varphi_{96}$}
\BIC{$(\mathrm{PNbhd}(f,c) \to \varphi_{64}) \to \mathrm{PNbhd}(f,c) \to \varphi_{69}$}
\end{ptree}
\noindent$\Pi_{4}$:
\begin{ptree}
\AXC{$\mathrm{PNbhd}(f,c) \to \mathrm{PNbhd}(g,c) \to \varphi_{62}$\lbl{A1}}
\AXC{$\varphi_{90}$\lbl{A1}}
\BIC{$\mathrm{PNbhd}(f,c) \to \varphi_{79}$}
\AXC{$(\mathrm{PNbhd}(f,c) \to \varphi_{79}) \to \varphi_{30} \to \mathrm{PNbhd}(f,c) \to \varphi_{77}$\lbl{A1}}
\BIC{$\varphi_{30} \to \mathrm{PNbhd}(f,c) \to \varphi_{77}$}
\end{ptree}
\noindent$\Pi_{5}$:
\begin{ptree}
\AXC{$\varphi_{61}$\lbl{A1}}
\AXC{$\varphi_{30}$\lbl{A2}}
\AXC{$\stackrel{\Pi_{4}}{\vdots}$}\noLine
\UIC{$\varphi_{30} \to \mathrm{PNbhd}(f,c) \to \varphi_{77}$}
\BIC{$\mathrm{PNbhd}(f,c) \to \varphi_{77}$}
\AXC{$(\mathrm{PNbhd}(f,c) \to \varphi_{77}) \to \varphi_{61} \to \varphi_{94}$\lbl{A1}}
\BIC{$\varphi_{61} \to \varphi_{94}$}
\BIC{$\varphi_{94}$}
\end{ptree}
\noindent$\Pi_{6}$:
\begin{ptree}
\AXC{$\stackrel{\Pi_{1}}{\vdots}$}\noLine
\UIC{$\mathrm{PNbhd}(f,c) \to \varphi_{64}$}
\AXC{$\stackrel{\Pi_{3}}{\vdots}$}\noLine
\UIC{$(\mathrm{PNbhd}(f,c) \to \varphi_{64}) \to \mathrm{PNbhd}(f,c) \to \varphi_{69}$}
\BIC{$\mathrm{PNbhd}(f,c) \to \varphi_{69}$}
\AXC{$\mathrm{PNbhd}(f,c) \to \varphi_{63}$\lbl{A1}}
\AXC{$\stackrel{\Pi_{5}}{\vdots}$}\noLine
\UIC{$\varphi_{94}$}
\BIC{$(\mathrm{PNbhd}(f,c) \to \varphi_{69}) \to \mathrm{PNbhd}(f,c) \to \varphi_{66}$}
\BIC{$\mathrm{PNbhd}(f,c) \to \varphi_{66}$}
\end{ptree}
\noindent$\Pi_{7}$:
\begin{ptree}
\AXC{$\mathrm{PNbhd}(f,c) \to \varphi_{60}$\lbl{A1}}
\AXC{$(\mathrm{PNbhd}(f,c) \to \varphi_{60}) \to \mathrm{PNbhd}(f,c) \to \varphi_{65}$\lbl{A1}}
\BIC{$\mathrm{PNbhd}(f,c) \to \varphi_{65}$}
\end{ptree}
\noindent$\Pi_{8}$:
\begin{ptree}
\AXC{$\forall x\,\forall y\,\forall z\,(x < y \to y \leq z \to x < z)$\lbl{[O3]}}
\AXC{$\varphi_{7}$\lbl{A2}}
\BIC{$\forall y\,\forall z\,(\lvert x - c \rvert < y \to y \leq z \to \lvert x - c \rvert < z)$}
\AXC{$\varphi_{21}$\lbl{A2}}
\BIC{$\forall z\,(\lvert x - c \rvert < d \to d \leq z \to \lvert x - c \rvert < z)$}
\AXC{$\varphi_{16}$\lbl{A2}}
\BIC{$\lvert x - c \rvert < d \to d \leq r_{1} \to \lvert x - c \rvert < r_{1}$}
\end{ptree}
\noindent$\Pi_{9}$:
\begin{ptree}
\AXC{$\mathrm{PNbhd}(f,c) \to \varphi_{68}$\lbl{A1}}
\AXC{$\stackrel{\Pi_{8}}{\vdots}$}\noLine
\UIC{$\lvert x - c \rvert < d \to d \leq r_{1} \to \lvert x - c \rvert < r_{1}$}
\AXC{$(\lvert x - c \rvert < d \to d \leq r_{1} \to \lvert x - c \rvert < r_{1}) \to \varphi_{97}$\lbl{A1}}
\BIC{$\varphi_{97}$}
\BIC{$(\mathrm{PNbhd}(f,c) \to \varphi_{65}) \to \mathrm{PNbhd}(f,c) \to \varphi_{70}$}
\end{ptree}
\noindent$\Pi_{10}$:
\begin{ptree}
\AXC{$\stackrel{\Pi_{7}}{\vdots}$}\noLine
\UIC{$\mathrm{PNbhd}(f,c) \to \varphi_{65}$}
\AXC{$\stackrel{\Pi_{9}}{\vdots}$}\noLine
\UIC{$(\mathrm{PNbhd}(f,c) \to \varphi_{65}) \to \mathrm{PNbhd}(f,c) \to \varphi_{70}$}
\BIC{$\mathrm{PNbhd}(f,c) \to \varphi_{70}$}
\end{ptree}
\noindent$\Pi_{11}$:
\begin{ptree}
\AXC{$\mathrm{PNbhd}(f,c) \to \varphi_{41}$\lbl{A1}}
\AXC{$(\mathrm{PNbhd}(f,c) \to \varphi_{41}) \to \mathrm{PNbhd}(f,c) \to \varphi_{80}$\lbl{A1}}
\BIC{$\mathrm{PNbhd}(f,c) \to \varphi_{80}$}
\end{ptree}
\noindent$\Pi_{12}$:
\begin{ptree}
\AXC{$\varphi_{31}$\lbl{A2}}
\AXC{$\stackrel{\Pi_{11}}{\vdots}$}\noLine
\UIC{$\mathrm{PNbhd}(f,c) \to \varphi_{80}$}
\AXC{$(\mathrm{PNbhd}(f,c) \to \varphi_{80}) \to \varphi_{31} \to \mathrm{PNbhd}(f,c) \to \varphi_{78}$\lbl{A1}}
\BIC{$\varphi_{31} \to \mathrm{PNbhd}(f,c) \to \varphi_{78}$}
\BIC{$\mathrm{PNbhd}(f,c) \to \varphi_{78}$}
\end{ptree}
\noindent$\Pi_{13}$:
\begin{ptree}
\AXC{$\mathrm{PNbhd}(f,c) \to \varphi_{63}$\lbl{A1}}
\AXC{$\varphi_{61}$\lbl{A1}}
\AXC{$\stackrel{\Pi_{12}}{\vdots}$}\noLine
\UIC{$\mathrm{PNbhd}(f,c) \to \varphi_{78}$}
\AXC{$(\mathrm{PNbhd}(f,c) \to \varphi_{78}) \to \varphi_{61} \to \varphi_{95}$\lbl{A1}}
\BIC{$\varphi_{61} \to \varphi_{95}$}
\BIC{$\varphi_{95}$}
\BIC{$(\mathrm{PNbhd}(f,c) \to \varphi_{70}) \to \mathrm{PNbhd}(f,c) \to \varphi_{67}$}
\end{ptree}
\noindent$\Pi_{14}$:
\begin{ptree}
\AXC{$\stackrel{\Pi_{6}}{\vdots}$}\noLine
\UIC{$\mathrm{PNbhd}(f,c) \to \varphi_{66}$}
\AXC{$\stackrel{\Pi_{10}}{\vdots}$}\noLine
\UIC{$\mathrm{PNbhd}(f,c) \to \varphi_{70}$}
\AXC{$\stackrel{\Pi_{13}}{\vdots}$}\noLine
\UIC{$(\mathrm{PNbhd}(f,c) \to \varphi_{70}) \to \mathrm{PNbhd}(f,c) \to \varphi_{67}$}
\BIC{$\mathrm{PNbhd}(f,c) \to \varphi_{67}$}
\AXC{$\varphi_{98}$\lbl{A1}}
\BIC{$(\mathrm{PNbhd}(f,c) \to \varphi_{66}) \to \varphi_{73}$}
\BIC{$\varphi_{73}$}
\end{ptree}
\noindent$\Pi_{15}$:
\begin{ptree}
\AXC{$\forall x\,(\varphi_{58} \to x \in \mathbb{R} \to \varphi_{3}) \to \forall x\,(\varphi_{58}) \to \forall x \in \mathbb{R}\,(\varphi_{3})$\lbl{A3}}
\AXC{$\stackrel{\Pi_{14}}{\vdots}$}\noLine
\UIC{$\varphi_{73}$}
\AXC{$\varphi_{73} \to \varphi_{58} \to x \in \mathbb{R} \to \varphi_{3}$\lbl{A1}}
\BIC{$\varphi_{58} \to x \in \mathbb{R} \to \varphi_{3}$}
\UIC{$\forall x\,(\varphi_{58} \to x \in \mathbb{R} \to \varphi_{3})$}
\AXC{$\forall x\,(\varphi_{58} \to x \in \mathbb{R} \to \varphi_{3}) \to \varphi_{101}$\lbl{A1}}
\BIC{$\varphi_{101}$}
\BIC{$(\varphi_{58} \to \forall x\,(\varphi_{58})) \to \varphi_{74}$}
\end{ptree}
\noindent$\Pi_{16}$:
\begin{ptree}
\AXC{$\mathrm{PNbhd}(f,c) \to \varphi_{60}$\lbl{A1}}
\AXC{$(\mathrm{PNbhd}(f,c) \to \varphi_{60}) \to \varphi_{74} \to \mathrm{PNbhd}(f,c) \to \varphi_{75}$\lbl{A1}}
\BIC{$\varphi_{74} \to \mathrm{PNbhd}(f,c) \to \varphi_{75}$}
\end{ptree}
\noindent$\Pi_{17}$:
\begin{ptree}
\AXC{$\varphi_{58} \to \forall x\,(\varphi_{58})$\lbl{A4}}
\AXC{$\stackrel{\Pi_{15}}{\vdots}$}\noLine
\UIC{$(\varphi_{58} \to \forall x\,(\varphi_{58})) \to \varphi_{74}$}
\BIC{$\varphi_{74}$}
\AXC{$\stackrel{\Pi_{16}}{\vdots}$}\noLine
\UIC{$\varphi_{74} \to \mathrm{PNbhd}(f,c) \to \varphi_{75}$}
\BIC{$\mathrm{PNbhd}(f,c) \to \varphi_{75}$}
\end{ptree}
\noindent$\Pi_{18}$:
\begin{ptree}
\AXC{$\varphi_{51}$\lbl{A2}}
\AXC{$\stackrel{\Pi_{17}}{\vdots}$}\noLine
\UIC{$\mathrm{PNbhd}(f,c) \to \varphi_{75}$}
\AXC{$(\mathrm{PNbhd}(f,c) \to \varphi_{75}) \to \varphi_{51} \to \mathrm{PNbhd}(f,c) \to \varphi_{76}$\lbl{A1}}
\BIC{$\varphi_{51} \to \mathrm{PNbhd}(f,c) \to \varphi_{76}$}
\BIC{$\mathrm{PNbhd}(f,c) \to \varphi_{76}$}
\AXC{$\varphi_{92}$\lbl{A1}}
\BIC{$\varphi_{72} \to \neg (d \in \mathbb{R} \land 0 < d \land d \leq r_{1} \land d \leq r_{2})$}
\UIC{$\forall d\,(\varphi_{72} \to \neg (d \in \mathbb{R} \land 0 < d \land d \leq r_{1} \land d \leq r_{2}))$}
\end{ptree}
\noindent$\Pi_{19}$:
\begin{ptree}
\AXC{$\varphi_{10}$\lbl{[D8]}}
\AXC{$\varphi_{10} \to \varphi_{11}$\lbl{A2}}
\BIC{$\varphi_{11}$}
\AXC{$\varphi_{11} \to r_{1} \in \mathbb{R} \to r_{2} \in \mathbb{R} \to \varphi_{2}$\lbl{A2}}
\BIC{$r_{1} \in \mathbb{R} \to r_{2} \in \mathbb{R} \to \varphi_{2}$}
\AXC{$(r_{1} \in \mathbb{R} \to r_{2} \in \mathbb{R} \to \varphi_{2}) \to \varphi_{100}$\lbl{A1}}
\BIC{$\varphi_{100}$}
\end{ptree}
\noindent$\Pi_{20}$:
\begin{ptree}
\AXC{$\mathrm{PNbhd}(f,c) \to \mathrm{PNbhd}(g,c) \to \varphi_{62}$\lbl{A1}}
\AXC{$\varphi_{88}$\lbl{A1}}
\BIC{$\mathrm{PNbhd}(f,c) \to \mathrm{PNbhd}(g,c) \to \varphi_{42}$}
\AXC{$\mathrm{PNbhd}(f,c) \to \varphi_{41}$\lbl{A1}}
\AXC{$\varphi_{84}$\lbl{A1}}
\BIC{$\mathrm{PNbhd}(f,c) \to \mathrm{PNbhd}(g,c) \to \varphi_{43}$}
\AXC{$\stackrel{\Pi_{19}}{\vdots}$}\noLine
\UIC{$\varphi_{100}$}
\BIC{$(\mathrm{PNbhd}(f,c) \to \mathrm{PNbhd}(g,c) \to \varphi_{42}) \to \varphi_{91}$}
\BIC{$\varphi_{91}$}
\end{ptree}
\noindent$\Pi_{21}$:
\begin{ptree}
\AXC{$\mathrm{PNbhd}(f,c) \to \varphi_{41}$\lbl{A1}}
\AXC{$(\mathrm{PNbhd}(f,c) \to \varphi_{41}) \to \mathrm{PNbhd}(f,c) \to \varphi_{48}$\lbl{A1}}
\BIC{$\mathrm{PNbhd}(f,c) \to \varphi_{48}$}
\AXC{$\stackrel{\Pi_{20}}{\vdots}$}\noLine
\UIC{$\varphi_{91}$}
\BIC{$(\mathrm{PNbhd}(f,c) \to \varphi_{47}) \to \varphi_{59}$}
\end{ptree}
\noindent$\Pi_{22}$:
\begin{ptree}
\AXC{$\mathrm{PNbhd}(f,c) \to \mathrm{PNbhd}(g,c) \to \varphi_{62}$\lbl{A1}}
\AXC{$\varphi_{87}$\lbl{A1}}
\BIC{$\mathrm{PNbhd}(f,c) \to \varphi_{47}$}
\AXC{$\stackrel{\Pi_{21}}{\vdots}$}\noLine
\UIC{$(\mathrm{PNbhd}(f,c) \to \varphi_{47}) \to \varphi_{59}$}
\BIC{$\varphi_{59}$}
\AXC{$\varphi_{104}$\lbl{A1}}
\BIC{$\forall d\,(\varphi_{72} \to \neg (d \in \mathbb{R} \land 0 < d \land d \leq r_{1} \land d \leq r_{2})) \to \varphi_{102}$}
\end{ptree}
\noindent$\Pi_{23}$:
\begin{ptree}
\AXC{$\stackrel{\Pi_{18}}{\vdots}$}\noLine
\UIC{$\forall d\,(\varphi_{72} \to \neg (d \in \mathbb{R} \land 0 < d \land d \leq r_{1} \land d \leq r_{2}))$}
\AXC{$\stackrel{\Pi_{22}}{\vdots}$}\noLine
\UIC{$\forall d\,(\varphi_{72} \to \neg (d \in \mathbb{R} \land 0 < d \land d \leq r_{1} \land d \leq r_{2})) \to \varphi_{102}$}
\BIC{$\varphi_{102}$}
\end{ptree}
\noindent$\Pi_{24}$:
\begin{ptree}
\AXC{$\varphi_{72} \to \forall d\,(\varphi_{72})$\lbl{A4}}
\AXC{$\forall d\,(\varphi_{72} \to \neg (d \in \mathbb{R} \land 0 < d \land d \leq r_{1} \land d \leq r_{2})) \to \varphi_{81}$\lbl{A3}}
\AXC{$\stackrel{\Pi_{23}}{\vdots}$}\noLine
\UIC{$\varphi_{102}$}
\BIC{$(\varphi_{72} \to \forall d\,(\varphi_{72})) \to \varphi_{71}$}
\BIC{$\varphi_{71}$}
\end{ptree}
\noindent$\Pi_{25}$:
\begin{ptree}
\AXC{$\forall f\,\forall c\,(\mathrm{PNbhd}(f,c) \leftrightarrow \varphi_{20})$\lbl{DefPNbhd}}
\AXC{$\varphi_{54}$\lbl{A2}}
\BIC{$\forall c\,(\mathrm{PNbhd}(g,c) \leftrightarrow \varphi_{19})$}
\AXC{$\varphi_{52}$\lbl{A2}}
\BIC{$\mathrm{PNbhd}(g,c) \leftrightarrow \varphi_{19}$}
\AXC{$\varphi_{1}$\lbl{A1}}
\AXC{$\varphi_{1} \to (\mathrm{PNbhd}(g,c) \leftrightarrow \varphi_{19}) \to \varphi_{49}$\lbl{A1}}
\BIC{$(\mathrm{PNbhd}(g,c) \leftrightarrow \varphi_{19}) \to \varphi_{49}$}
\BIC{$\varphi_{49}$}
\end{ptree}
\noindent$\Pi_{26}$:
\begin{ptree}
\AXC{$\varphi_{57} \to \forall r_{2}\,(\varphi_{57})$\lbl{A4}}
\AXC{$\varphi_{89}$\lbl{A3}}
\AXC{$\stackrel{\Pi_{24}}{\vdots}$}\noLine
\UIC{$\varphi_{71}$}
\AXC{$\varphi_{71} \to \varphi_{57} \to \neg (r_{2} \in \mathbb{R} \land 0 < r_{2} \land \varphi_{4})$\lbl{A1}}
\BIC{$\varphi_{57} \to \neg (r_{2} \in \mathbb{R} \land 0 < r_{2} \land \varphi_{4})$}
\UIC{$\forall r_{2}\,(\varphi_{57} \to \neg (r_{2} \in \mathbb{R} \land 0 < r_{2} \land \varphi_{4}))$}
\AXC{$\stackrel{\Pi_{25}}{\vdots}$}\noLine
\UIC{$\varphi_{49}$}
\AXC{$\varphi_{49} \to \varphi_{105}$\lbl{A1}}
\BIC{$\varphi_{105}$}
\BIC{$\varphi_{89} \to (\varphi_{57} \to \forall r_{2}\,(\varphi_{57})) \to \varphi_{99}$}
\BIC{$(\varphi_{57} \to \forall r_{2}\,(\varphi_{57})) \to \varphi_{99}$}
\BIC{$\varphi_{99}$}
\end{ptree}
\noindent$\Pi_{27}$:
\begin{ptree}
\AXC{$\varphi_{40}$\lbl{A4}}
\AXC{$\varphi_{83}$\lbl{A3}}
\AXC{$\forall r_{2}\,(\neg (r_{2} \in \mathbb{R} \land 0 < r_{2} \land \varphi_{4})) \to \varphi_{14}$\lbl{A2}}
\UIC{$\forall r\,(\forall r_{2}\,(\neg (r_{2} \in \mathbb{R} \land 0 < r_{2} \land \varphi_{4})) \to \varphi_{14})$}
\AXC{$\stackrel{\Pi_{26}}{\vdots}$}\noLine
\UIC{$\varphi_{99}$}
\BIC{$\varphi_{83} \to \varphi_{40} \to \mathrm{PNbhd}(f,c) \to \mathrm{PNbhd}(g,c) \to \varphi_{45}$}
\BIC{$\varphi_{40} \to \mathrm{PNbhd}(f,c) \to \mathrm{PNbhd}(g,c) \to \varphi_{45}$}
\BIC{$\mathrm{PNbhd}(f,c) \to \mathrm{PNbhd}(g,c) \to \varphi_{45}$}
\AXC{$\varphi_{85}$\lbl{A1}}
\BIC{$\varphi_{29} \to \neg (r_{1} \in \mathbb{R} \land 0 < r_{1} \land \varphi_{5})$}
\UIC{$\forall r_{1}\,(\varphi_{29} \to \neg (r_{1} \in \mathbb{R} \land 0 < r_{1} \land \varphi_{5}))$}
\end{ptree}
\noindent$\Pi_{28}$:
\begin{ptree}
\AXC{$\mathrm{PNbhd}(f,c) \to \mathrm{PNbhd}(f,c)$\lbl{A1}}
\AXC{$(\mathrm{PNbhd}(f,c) \to \mathrm{PNbhd}(f,c)) \to \varphi_{56}$\lbl{A1}}
\BIC{$\varphi_{56}$}
\end{ptree}
\noindent$\Pi_{29}$:
\begin{ptree}
\AXC{$\forall f\,\forall c\,(\mathrm{PNbhd}(f,c) \leftrightarrow \varphi_{20})$\lbl{DefPNbhd}}
\AXC{$\varphi_{55}$\lbl{A2}}
\BIC{$\forall c\,(\mathrm{PNbhd}(f,c) \leftrightarrow \varphi_{20})$}
\AXC{$\varphi_{53}$\lbl{A2}}
\BIC{$\mathrm{PNbhd}(f,c) \leftrightarrow \varphi_{20}$}
\AXC{$\stackrel{\Pi_{28}}{\vdots}$}\noLine
\UIC{$\varphi_{56}$}
\BIC{$\mathrm{PNbhd}(f,c) \to \mathrm{PNbhd}(g,c) \to \varphi_{20}$}
\end{ptree}
\noindent$\Pi_{30}$:
\begin{ptree}
\AXC{$\stackrel{\Pi_{29}}{\vdots}$}\noLine
\UIC{$\mathrm{PNbhd}(f,c) \to \mathrm{PNbhd}(g,c) \to \varphi_{20}$}
\AXC{$(\mathrm{PNbhd}(f,c) \to \mathrm{PNbhd}(g,c) \to \varphi_{20}) \to \varphi_{103}$\lbl{A1}}
\BIC{$\varphi_{103}$}
\end{ptree}
\noindent$\Pi_{31}$:
\begin{ptree}
\AXC{$\varphi_{29} \to \forall r_{1}\,(\varphi_{29})$\lbl{A4}}
\AXC{$\varphi_{86}$\lbl{A3}}
\AXC{$\stackrel{\Pi_{27}}{\vdots}$}\noLine
\UIC{$\forall r_{1}\,(\varphi_{29} \to \neg (r_{1} \in \mathbb{R} \land 0 < r_{1} \land \varphi_{5}))$}
\AXC{$\stackrel{\Pi_{30}}{\vdots}$}\noLine
\UIC{$\varphi_{103}$}
\BIC{$\varphi_{86} \to (\varphi_{29} \to \forall r_{1}\,(\varphi_{29})) \to \varphi_{93}$}
\BIC{$(\varphi_{29} \to \forall r_{1}\,(\varphi_{29})) \to \varphi_{93}$}
\BIC{$\varphi_{93}$}
\end{ptree}
\noindent 主木:
\begin{ptree}
\AXC{$\varphi_{39}$\lbl{A4}}
\AXC{$\varphi_{82}$\lbl{A3}}
\AXC{$\forall r_{1}\,(\neg (r_{1} \in \mathbb{R} \land 0 < r_{1} \land \varphi_{5})) \to \varphi_{13}$\lbl{A2}}
\UIC{$\forall r\,(\forall r_{1}\,(\neg (r_{1} \in \mathbb{R} \land 0 < r_{1} \land \varphi_{5})) \to \varphi_{13})$}
\AXC{$\stackrel{\Pi_{31}}{\vdots}$}\noLine
\UIC{$\varphi_{93}$}
\BIC{$\varphi_{82} \to \varphi_{39} \to \mathrm{PNbhd}(f,c) \to \varphi_{27}$}
\BIC{$\varphi_{39} \to \mathrm{PNbhd}(f,c) \to \varphi_{27}$}
\BIC{$\mathrm{PNbhd}(f,c) \to \varphi_{27}$}
\UIC{$\forall c\,(\mathrm{PNbhd}(f,c) \to \varphi_{27})$}
\UIC{$\forall g\,\forall c\,(\mathrm{PNbhd}(f,c) \to \varphi_{27})$}
\UIC{$\forall f\,\forall g\,\forall c\,(\mathrm{PNbhd}(f,c) \to \varphi_{27})$}
\end{ptree}
\begin{abbrevlist}
\item $\varphi_{1} :\equiv \mathrm{PNbhd}(f,c) \to \mathrm{PNbhd}(g,c) \to \mathrm{PNbhd}(g,c)$
\item $\varphi_{2} :\equiv 0 < r_{1} \to 0 < r_{2} \to \exists d\,(d \in \mathbb{R} \land 0 < d \land d \leq r_{1} \land d \leq r_{2})$
\item $\varphi_{3} :\equiv 0 < \lvert x - c \rvert \to \lvert x - c \rvert < d \to x \in \mathrm{dom}(f) \land x \in \mathrm{dom}(g)$
\item $\varphi_{4} :\equiv \forall x \in \mathbb{R}\,(0 < \lvert x - c \rvert \to \lvert x - c \rvert < r_{2} \to x \in \mathrm{dom}(g))$
\item $\varphi_{5} :\equiv \forall x \in \mathbb{R}\,(0 < \lvert x - c \rvert \to \lvert x - c \rvert < r_{1} \to x \in \mathrm{dom}(f))$
\item $\varphi_{6} :\equiv d \in \mathbb{R} \land 0 < d \land d \leq r_{1} \land d \leq r_{2} \to d \in \mathbb{R} \land 0 < d \land d \leq r_{1} \land d \leq r_{2}$
\item $\varphi_{7} :\equiv \forall x\,\forall y\,\forall z\,(x < y \to y \leq z \to x < z) \to \forall y\,\forall z\,(\lvert x - c \rvert < y \to y \leq z \to \lvert x - c \rvert < z)$
\item $\varphi_{8} :\equiv \lvert x - c \rvert < d \to x \in \mathbb{R} \to 0 < \lvert x - c \rvert \to \lvert x - c \rvert < r_{2} \to x \in \mathrm{dom}(g)$
\item $\varphi_{9} :\equiv \lvert x - c \rvert < d \to x \in \mathbb{R} \to 0 < \lvert x - c \rvert \to \lvert x - c \rvert < r_{1} \to x \in \mathrm{dom}(f)$
\item $\varphi_{10} :\equiv \forall p\,\forall q\,(p \in \mathbb{R} \to q \in \mathbb{R} \to 0 < p \to 0 < q \to \exists d\,(d \in \mathbb{R} \land 0 < d \land d \leq p \land d \leq q))$
\item $\varphi_{11} :\equiv \forall q\,(r_{1} \in \mathbb{R} \to q \in \mathbb{R} \to 0 < r_{1} \to 0 < q \to \exists d\,(d \in \mathbb{R} \land 0 < d \land d \leq r_{1} \land d \leq q))$
\item $\varphi_{12} :\equiv d \in \mathbb{R} \land 0 < d \land d \leq r_{1} \land d \leq r_{2} \to x \in \mathbb{R} \to 0 < \lvert x - c \rvert \to 0 < \lvert x - c \rvert$
\item $\varphi_{13} :\equiv \neg (r \in \mathbb{R} \land 0 < r \land \forall x \in \mathbb{R}\,(0 < \lvert x - c \rvert \to \lvert x - c \rvert < r \to x \in \mathrm{dom}(f)))$
\item $\varphi_{14} :\equiv \neg (r \in \mathbb{R} \land 0 < r \land \forall x \in \mathbb{R}\,(0 < \lvert x - c \rvert \to \lvert x - c \rvert < r \to x \in \mathrm{dom}(g)))$
\item $\varphi_{15} :\equiv \forall z\,(\lvert x - c \rvert < d \to d \leq z \to \lvert x - c \rvert < z) \to \lvert x - c \rvert < d \to d \leq r_{2} \to \lvert x - c \rvert < r_{2}$
\item $\varphi_{16} :\equiv \forall z\,(\lvert x - c \rvert < d \to d \leq z \to \lvert x - c \rvert < z) \to \lvert x - c \rvert < d \to d \leq r_{1} \to \lvert x - c \rvert < r_{1}$
\item $\varphi_{17} :\equiv d \in \mathbb{R} \land 0 < d \land d \leq r_{1} \land d \leq r_{2} \to x \in \mathbb{R} \to 0 < \lvert x - c \rvert \to \lvert x - c \rvert < d \to d \leq r_{2}$
\item $\varphi_{18} :\equiv d \in \mathbb{R} \land 0 < d \land d \leq r_{1} \land d \leq r_{2} \to x \in \mathbb{R} \to 0 < \lvert x - c \rvert \to \lvert x - c \rvert < d \to d \leq r_{1}$
\item $\varphi_{19} :\equiv \exists r\,(r \in \mathbb{R} \land 0 < r \land \forall x \in \mathbb{R}\,(0 < \lvert x - c \rvert \to \lvert x - c \rvert < r \to x \in \mathrm{dom}(g)))$
\item $\varphi_{20} :\equiv \exists r\,(r \in \mathbb{R} \land 0 < r \land \forall x \in \mathbb{R}\,(0 < \lvert x - c \rvert \to \lvert x - c \rvert < r \to x \in \mathrm{dom}(f)))$
\item $\varphi_{21} :\equiv \forall y\,\forall z\,(\lvert x - c \rvert < y \to y \leq z \to \lvert x - c \rvert < z) \to \forall z\,(\lvert x - c \rvert < d \to d \leq z \to \lvert x - c \rvert < z)$
\item $\varphi_{22} :\equiv d \in \mathbb{R} \land 0 < d \land d \leq r_{1} \land d \leq r_{2} \to x \in \mathbb{R} \to 0 < \lvert x - c \rvert \to \lvert x - c \rvert < d \to x \in \mathrm{dom}(g)$
\item $\varphi_{23} :\equiv d \in \mathbb{R} \land 0 < d \land d \leq r_{1} \land d \leq r_{2} \to x \in \mathbb{R} \to 0 < \lvert x - c \rvert \to \lvert x - c \rvert < d \to x \in \mathrm{dom}(f)$
\item $\varphi_{24} :\equiv d \in \mathbb{R} \land 0 < d \land d \leq r_{1} \land d \leq r_{2} \to x \in \mathbb{R} \to 0 < \lvert x - c \rvert \to \lvert x - c \rvert < d \to \lvert x - c \rvert < d$
\item $\varphi_{25} :\equiv d \in \mathbb{R} \land 0 < d \land d \leq r_{1} \land d \leq r_{2} \to x \in \mathbb{R} \to 0 < \lvert x - c \rvert \to \lvert x - c \rvert < d \to \lvert x - c \rvert < r_{2}$
\item $\varphi_{26} :\equiv d \in \mathbb{R} \land 0 < d \land d \leq r_{1} \land d \leq r_{2} \to x \in \mathbb{R} \to 0 < \lvert x - c \rvert \to \lvert x - c \rvert < d \to \lvert x - c \rvert < r_{1}$
\item $\varphi_{27} :\equiv \mathrm{PNbhd}(g,c) \to \exists d\,(d \in \mathbb{R} \land 0 < d \land \forall x \in \mathbb{R}\,(\varphi_{3}))$
\item $\varphi_{28} :\equiv r_{2} \in \mathbb{R} \land 0 < r_{2} \land \varphi_{4} \to \exists d\,(d \in \mathbb{R} \land 0 < d \land d \leq r_{1} \land d \leq r_{2})$
\item $\varphi_{29} :\equiv \mathrm{PNbhd}(f,c) \land \mathrm{PNbhd}(g,c) \land \neg \exists d\,(d \in \mathbb{R} \land 0 < d \land \forall x \in \mathbb{R}\,(\varphi_{3}))$
\item $\varphi_{30} :\equiv \varphi_{4} \to x \in \mathbb{R} \to 0 < \lvert x - c \rvert \to \lvert x - c \rvert < r_{2} \to x \in \mathrm{dom}(g)$
\item $\varphi_{31} :\equiv \varphi_{5} \to x \in \mathbb{R} \to 0 < \lvert x - c \rvert \to \lvert x - c \rvert < r_{1} \to x \in \mathrm{dom}(f)$
\item $\varphi_{32} :\equiv d \in \mathbb{R} \land 0 < d \land d \leq r_{1} \land d \leq r_{2} \to d \in \mathbb{R} \land 0 < d \land \forall x \in \mathbb{R}\,(\varphi_{3})$
\item $\varphi_{33} :\equiv d \in \mathbb{R} \land 0 < d \land d \leq r_{1} \land d \leq r_{2} \to \exists d\,(d \in \mathbb{R} \land 0 < d \land \forall x \in \mathbb{R}\,(\varphi_{3}))$
\item $\varphi_{34} :\equiv d \in \mathbb{R} \land 0 < d \land d \leq r_{1} \land d \leq r_{2} \to x \in \mathbb{R} \to 0 < \lvert x - c \rvert \to \varphi_{8}$
\item $\varphi_{35} :\equiv d \in \mathbb{R} \land 0 < d \land d \leq r_{1} \land d \leq r_{2} \to x \in \mathbb{R} \to 0 < \lvert x - c \rvert \to \varphi_{9}$
\item $\varphi_{36} :\equiv r_{2} \in \mathbb{R} \land 0 < r_{2} \land \varphi_{4} \to d \in \mathbb{R} \land 0 < d \land d \leq r_{1} \land d \leq r_{2} \to x \in \mathbb{R} \to x \in \mathbb{R}$
\item $\varphi_{37} :\equiv d \in \mathbb{R} \land 0 < d \land d \leq r_{1} \land d \leq r_{2} \to x \in \mathbb{R} \to 0 < \lvert x - c \rvert \to \lvert x - c \rvert < d \to \varphi_{4}$
\item $\varphi_{38} :\equiv d \in \mathbb{R} \land 0 < d \land d \leq r_{1} \land d \leq r_{2} \to x \in \mathbb{R} \to 0 < \lvert x - c \rvert \to \lvert x - c \rvert < d \to \varphi_{5}$
\item $\varphi_{39} :\equiv \forall r_{1}\,(\neg (r_{1} \in \mathbb{R} \land 0 < r_{1} \land \varphi_{5})) \to \forall r\,\forall r_{1}\,(\neg (r_{1} \in \mathbb{R} \land 0 < r_{1} \land \varphi_{5}))$
\item $\varphi_{40} :\equiv \forall r_{2}\,(\neg (r_{2} \in \mathbb{R} \land 0 < r_{2} \land \varphi_{4})) \to \forall r\,\forall r_{2}\,(\neg (r_{2} \in \mathbb{R} \land 0 < r_{2} \land \varphi_{4}))$
\item $\varphi_{41} :\equiv \mathrm{PNbhd}(g,c) \to r_{1} \in \mathbb{R} \land 0 < r_{1} \land \varphi_{5} \to r_{1} \in \mathbb{R} \land 0 < r_{1} \land \varphi_{5}$
\item $\varphi_{42} :\equiv r_{1} \in \mathbb{R} \land 0 < r_{1} \land \varphi_{5} \to r_{2} \in \mathbb{R} \land 0 < r_{2} \land \varphi_{4} \to r_{2} \in \mathbb{R}$
\item $\varphi_{43} :\equiv r_{1} \in \mathbb{R} \land 0 < r_{1} \land \varphi_{5} \to r_{2} \in \mathbb{R} \land 0 < r_{2} \land \varphi_{4} \to r_{1} \in \mathbb{R}$
\item $\varphi_{44} :\equiv r_{2} \in \mathbb{R} \land 0 < r_{2} \land \varphi_{4} \to \exists d\,(d \in \mathbb{R} \land 0 < d \land \forall x \in \mathbb{R}\,(\varphi_{3}))$
\item $\varphi_{45} :\equiv r_{1} \in \mathbb{R} \land 0 < r_{1} \land \varphi_{5} \to \exists d\,(d \in \mathbb{R} \land 0 < d \land \forall x \in \mathbb{R}\,(\varphi_{3}))$
\item $\varphi_{46} :\equiv r_{2} \in \mathbb{R} \land 0 < r_{2} \land \varphi_{4} \to d \in \mathbb{R} \land 0 < d \land d \leq r_{1} \land d \leq r_{2} \to x \in \mathbb{R} \to \varphi_{3}$
\item $\varphi_{47} :\equiv \mathrm{PNbhd}(g,c) \to r_{1} \in \mathbb{R} \land 0 < r_{1} \land \varphi_{5} \to r_{2} \in \mathbb{R} \land 0 < r_{2} \land \varphi_{4} \to 0 < r_{2}$
\item $\varphi_{48} :\equiv \mathrm{PNbhd}(g,c) \to r_{1} \in \mathbb{R} \land 0 < r_{1} \land \varphi_{5} \to r_{2} \in \mathbb{R} \land 0 < r_{2} \land \varphi_{4} \to 0 < r_{1}$
\item $\varphi_{49} :\equiv \mathrm{PNbhd}(f,c) \to \mathrm{PNbhd}(g,c) \to r_{1} \in \mathbb{R} \land 0 < r_{1} \land \varphi_{5} \to \varphi_{19}$
\item $\varphi_{50} :\equiv r_{2} \in \mathbb{R} \land 0 < r_{2} \land \varphi_{4} \to d \in \mathbb{R} \land 0 < d \land d \leq r_{1} \land d \leq r_{2} \to \forall x \in \mathbb{R}\,(\varphi_{3})$
\item $\varphi_{51} :\equiv \forall d\,(\neg (d \in \mathbb{R} \land 0 < d \land \forall x \in \mathbb{R}\,(\varphi_{3}))) \to \neg (d \in \mathbb{R} \land 0 < d \land \forall x \in \mathbb{R}\,(\varphi_{3}))$
\item $\varphi_{52} :\equiv \forall c\,(\mathrm{PNbhd}(g,c) \leftrightarrow \varphi_{19}) \to (\mathrm{PNbhd}(g,c) \leftrightarrow \varphi_{19})$
\item $\varphi_{53} :\equiv \forall c\,(\mathrm{PNbhd}(f,c) \leftrightarrow \varphi_{20}) \to (\mathrm{PNbhd}(f,c) \leftrightarrow \varphi_{20})$
\item $\varphi_{54} :\equiv \forall f\,\forall c\,(\mathrm{PNbhd}(f,c) \leftrightarrow \varphi_{20}) \to \forall c\,(\mathrm{PNbhd}(g,c) \leftrightarrow \varphi_{19})$
\item $\varphi_{55} :\equiv \forall f\,\forall c\,(\mathrm{PNbhd}(f,c) \leftrightarrow \varphi_{20}) \to \forall c\,(\mathrm{PNbhd}(f,c) \leftrightarrow \varphi_{20})$
\item $\varphi_{56} :\equiv (\mathrm{PNbhd}(f,c) \leftrightarrow \varphi_{20}) \to \mathrm{PNbhd}(f,c) \to \mathrm{PNbhd}(g,c) \to \varphi_{20}$
\item $\varphi_{57} :\equiv \mathrm{PNbhd}(f,c) \land \mathrm{PNbhd}(g,c) \land (r_{1} \in \mathbb{R} \land 0 < r_{1} \land \varphi_{5}) \land \neg \exists d\,(d \in \mathbb{R} \land 0 < d \land \forall x \in \mathbb{R}\,(\varphi_{3}))$
\item $\varphi_{58} :\equiv \mathrm{PNbhd}(f,c) \land \mathrm{PNbhd}(g,c) \land (r_{1} \in \mathbb{R} \land 0 < r_{1} \land \varphi_{5}) \land (r_{2} \in \mathbb{R} \land 0 < r_{2} \land \varphi_{4}) \land (d \in \mathbb{R} \land 0 < d \land d \leq r_{1} \land d \leq r_{2})$
\item $\varphi_{59} :\equiv \mathrm{PNbhd}(f,c) \to \mathrm{PNbhd}(g,c) \to r_{1} \in \mathbb{R} \land 0 < r_{1} \land \varphi_{5} \to \varphi_{28}$
\item $\varphi_{60} :\equiv \mathrm{PNbhd}(g,c) \to r_{1} \in \mathbb{R} \land 0 < r_{1} \land \varphi_{5} \to r_{2} \in \mathbb{R} \land 0 < r_{2} \land \varphi_{4} \to \varphi_{6}$
\item $\varphi_{61} :\equiv \mathrm{PNbhd}(f,c) \to \mathrm{PNbhd}(g,c) \to r_{1} \in \mathbb{R} \land 0 < r_{1} \land \varphi_{5} \to \varphi_{36}$
\item $\varphi_{62} :\equiv r_{1} \in \mathbb{R} \land 0 < r_{1} \land \varphi_{5} \to r_{2} \in \mathbb{R} \land 0 < r_{2} \land \varphi_{4} \to r_{2} \in \mathbb{R} \land 0 < r_{2} \land \varphi_{4}$
\item $\varphi_{63} :\equiv \mathrm{PNbhd}(g,c) \to r_{1} \in \mathbb{R} \land 0 < r_{1} \land \varphi_{5} \to r_{2} \in \mathbb{R} \land 0 < r_{2} \land \varphi_{4} \to \varphi_{12}$
\item $\varphi_{64} :\equiv \mathrm{PNbhd}(g,c) \to r_{1} \in \mathbb{R} \land 0 < r_{1} \land \varphi_{5} \to r_{2} \in \mathbb{R} \land 0 < r_{2} \land \varphi_{4} \to \varphi_{17}$
\item $\varphi_{65} :\equiv \mathrm{PNbhd}(g,c) \to r_{1} \in \mathbb{R} \land 0 < r_{1} \land \varphi_{5} \to r_{2} \in \mathbb{R} \land 0 < r_{2} \land \varphi_{4} \to \varphi_{18}$
\item $\varphi_{66} :\equiv \mathrm{PNbhd}(g,c) \to r_{1} \in \mathbb{R} \land 0 < r_{1} \land \varphi_{5} \to r_{2} \in \mathbb{R} \land 0 < r_{2} \land \varphi_{4} \to \varphi_{22}$
\item $\varphi_{67} :\equiv \mathrm{PNbhd}(g,c) \to r_{1} \in \mathbb{R} \land 0 < r_{1} \land \varphi_{5} \to r_{2} \in \mathbb{R} \land 0 < r_{2} \land \varphi_{4} \to \varphi_{23}$
\item $\varphi_{68} :\equiv \mathrm{PNbhd}(g,c) \to r_{1} \in \mathbb{R} \land 0 < r_{1} \land \varphi_{5} \to r_{2} \in \mathbb{R} \land 0 < r_{2} \land \varphi_{4} \to \varphi_{24}$
\item $\varphi_{69} :\equiv \mathrm{PNbhd}(g,c) \to r_{1} \in \mathbb{R} \land 0 < r_{1} \land \varphi_{5} \to r_{2} \in \mathbb{R} \land 0 < r_{2} \land \varphi_{4} \to \varphi_{25}$
\item $\varphi_{70} :\equiv \mathrm{PNbhd}(g,c) \to r_{1} \in \mathbb{R} \land 0 < r_{1} \land \varphi_{5} \to r_{2} \in \mathbb{R} \land 0 < r_{2} \land \varphi_{4} \to \varphi_{26}$
\item $\varphi_{71} :\equiv \mathrm{PNbhd}(f,c) \to \mathrm{PNbhd}(g,c) \to r_{1} \in \mathbb{R} \land 0 < r_{1} \land \varphi_{5} \to \varphi_{44}$
\item $\varphi_{72} :\equiv \mathrm{PNbhd}(f,c) \land \mathrm{PNbhd}(g,c) \land (r_{1} \in \mathbb{R} \land 0 < r_{1} \land \varphi_{5}) \land (r_{2} \in \mathbb{R} \land 0 < r_{2} \land \varphi_{4}) \land \neg \exists d\,(d \in \mathbb{R} \land 0 < d \land \forall x \in \mathbb{R}\,(\varphi_{3}))$
\item $\varphi_{73} :\equiv \mathrm{PNbhd}(f,c) \to \mathrm{PNbhd}(g,c) \to r_{1} \in \mathbb{R} \land 0 < r_{1} \land \varphi_{5} \to \varphi_{46}$
\item $\varphi_{74} :\equiv \mathrm{PNbhd}(f,c) \to \mathrm{PNbhd}(g,c) \to r_{1} \in \mathbb{R} \land 0 < r_{1} \land \varphi_{5} \to \varphi_{50}$
\item $\varphi_{75} :\equiv \mathrm{PNbhd}(g,c) \to r_{1} \in \mathbb{R} \land 0 < r_{1} \land \varphi_{5} \to r_{2} \in \mathbb{R} \land 0 < r_{2} \land \varphi_{4} \to \varphi_{32}$
\item $\varphi_{76} :\equiv \mathrm{PNbhd}(g,c) \to r_{1} \in \mathbb{R} \land 0 < r_{1} \land \varphi_{5} \to r_{2} \in \mathbb{R} \land 0 < r_{2} \land \varphi_{4} \to \varphi_{33}$
\item $\varphi_{77} :\equiv \mathrm{PNbhd}(g,c) \to r_{1} \in \mathbb{R} \land 0 < r_{1} \land \varphi_{5} \to r_{2} \in \mathbb{R} \land 0 < r_{2} \land \varphi_{4} \to \varphi_{34}$
\item $\varphi_{78} :\equiv \mathrm{PNbhd}(g,c) \to r_{1} \in \mathbb{R} \land 0 < r_{1} \land \varphi_{5} \to r_{2} \in \mathbb{R} \land 0 < r_{2} \land \varphi_{4} \to \varphi_{35}$
\item $\varphi_{79} :\equiv \mathrm{PNbhd}(g,c) \to r_{1} \in \mathbb{R} \land 0 < r_{1} \land \varphi_{5} \to r_{2} \in \mathbb{R} \land 0 < r_{2} \land \varphi_{4} \to \varphi_{37}$
\item $\varphi_{80} :\equiv \mathrm{PNbhd}(g,c) \to r_{1} \in \mathbb{R} \land 0 < r_{1} \land \varphi_{5} \to r_{2} \in \mathbb{R} \land 0 < r_{2} \land \varphi_{4} \to \varphi_{38}$
\item $\varphi_{81} :\equiv \forall d\,(\varphi_{72}) \to \forall d\,(\neg (d \in \mathbb{R} \land 0 < d \land d \leq r_{1} \land d \leq r_{2}))$
\item $\varphi_{82} :\equiv \forall r\,(\forall r_{1}\,(\neg (r_{1} \in \mathbb{R} \land 0 < r_{1} \land \varphi_{5})) \to \varphi_{13}) \to \forall r\,\forall r_{1}\,(\neg (r_{1} \in \mathbb{R} \land 0 < r_{1} \land \varphi_{5})) \to \forall r\,(\varphi_{13})$
\item $\varphi_{83} :\equiv \forall r\,(\forall r_{2}\,(\neg (r_{2} \in \mathbb{R} \land 0 < r_{2} \land \varphi_{4})) \to \varphi_{14}) \to \forall r\,\forall r_{2}\,(\neg (r_{2} \in \mathbb{R} \land 0 < r_{2} \land \varphi_{4})) \to \forall r\,(\varphi_{14})$
\item $\varphi_{84} :\equiv (\mathrm{PNbhd}(f,c) \to \varphi_{41}) \to \mathrm{PNbhd}(f,c) \to \mathrm{PNbhd}(g,c) \to \varphi_{43}$
\item $\varphi_{85} :\equiv (\mathrm{PNbhd}(f,c) \to \mathrm{PNbhd}(g,c) \to \varphi_{45}) \to \varphi_{29} \to \neg (r_{1} \in \mathbb{R} \land 0 < r_{1} \land \varphi_{5})$
\item $\varphi_{86} :\equiv \forall r_{1}\,(\varphi_{29} \to \neg (r_{1} \in \mathbb{R} \land 0 < r_{1} \land \varphi_{5})) \to \forall r_{1}\,(\varphi_{29}) \to \forall r_{1}\,(\neg (r_{1} \in \mathbb{R} \land 0 < r_{1} \land \varphi_{5}))$
\item $\varphi_{87} :\equiv (\mathrm{PNbhd}(f,c) \to \mathrm{PNbhd}(g,c) \to \varphi_{62}) \to \mathrm{PNbhd}(f,c) \to \varphi_{47}$
\item $\varphi_{88} :\equiv (\mathrm{PNbhd}(f,c) \to \mathrm{PNbhd}(g,c) \to \varphi_{62}) \to \mathrm{PNbhd}(f,c) \to \mathrm{PNbhd}(g,c) \to \varphi_{42}$
\item $\varphi_{89} :\equiv \forall r_{2}\,(\varphi_{57} \to \neg (r_{2} \in \mathbb{R} \land 0 < r_{2} \land \varphi_{4})) \to \forall r_{2}\,(\varphi_{57}) \to \forall r_{2}\,(\neg (r_{2} \in \mathbb{R} \land 0 < r_{2} \land \varphi_{4}))$
\item $\varphi_{90} :\equiv (\mathrm{PNbhd}(f,c) \to \mathrm{PNbhd}(g,c) \to \varphi_{62}) \to \mathrm{PNbhd}(f,c) \to \varphi_{79}$
\item $\varphi_{91} :\equiv (\mathrm{PNbhd}(f,c) \to \varphi_{48}) \to (\mathrm{PNbhd}(f,c) \to \varphi_{47}) \to \varphi_{59}$
\item $\varphi_{92} :\equiv (\mathrm{PNbhd}(f,c) \to \varphi_{76}) \to \varphi_{72} \to \neg (d \in \mathbb{R} \land 0 < d \land d \leq r_{1} \land d \leq r_{2})$
\item $\varphi_{93} :\equiv \forall r\,(\forall r_{1}\,(\neg (r_{1} \in \mathbb{R} \land 0 < r_{1} \land \varphi_{5})) \to \varphi_{13}) \to \varphi_{82} \to \varphi_{39} \to \mathrm{PNbhd}(f,c) \to \varphi_{27}$
\item $\varphi_{94} :\equiv (\mathrm{PNbhd}(f,c) \to \varphi_{63}) \to (\mathrm{PNbhd}(f,c) \to \varphi_{69}) \to \mathrm{PNbhd}(f,c) \to \varphi_{66}$
\item $\varphi_{95} :\equiv (\mathrm{PNbhd}(f,c) \to \varphi_{63}) \to (\mathrm{PNbhd}(f,c) \to \varphi_{70}) \to \mathrm{PNbhd}(f,c) \to \varphi_{67}$
\item $\varphi_{96} :\equiv (\mathrm{PNbhd}(f,c) \to \varphi_{68}) \to (\mathrm{PNbhd}(f,c) \to \varphi_{64}) \to \mathrm{PNbhd}(f,c) \to \varphi_{69}$
\item $\varphi_{97} :\equiv (\mathrm{PNbhd}(f,c) \to \varphi_{68}) \to (\mathrm{PNbhd}(f,c) \to \varphi_{65}) \to \mathrm{PNbhd}(f,c) \to \varphi_{70}$
\item $\varphi_{98} :\equiv (\mathrm{PNbhd}(f,c) \to \varphi_{67}) \to (\mathrm{PNbhd}(f,c) \to \varphi_{66}) \to \varphi_{73}$
\item $\varphi_{99} :\equiv \forall r\,(\forall r_{2}\,(\neg (r_{2} \in \mathbb{R} \land 0 < r_{2} \land \varphi_{4})) \to \varphi_{14}) \to \varphi_{83} \to \varphi_{40} \to \mathrm{PNbhd}(f,c) \to \mathrm{PNbhd}(g,c) \to \varphi_{45}$
\item $\varphi_{100} :\equiv (\mathrm{PNbhd}(f,c) \to \mathrm{PNbhd}(g,c) \to \varphi_{43}) \to (\mathrm{PNbhd}(f,c) \to \mathrm{PNbhd}(g,c) \to \varphi_{42}) \to \varphi_{91}$
\item $\varphi_{101} :\equiv (\forall x\,(\varphi_{58} \to x \in \mathbb{R} \to \varphi_{3}) \to \forall x\,(\varphi_{58}) \to \forall x \in \mathbb{R}\,(\varphi_{3})) \to (\varphi_{58} \to \forall x\,(\varphi_{58})) \to \varphi_{74}$
\item $\varphi_{102} :\equiv (\forall d\,(\varphi_{72} \to \neg (d \in \mathbb{R} \land 0 < d \land d \leq r_{1} \land d \leq r_{2})) \to \varphi_{81}) \to (\varphi_{72} \to \forall d\,(\varphi_{72})) \to \varphi_{71}$
\item $\varphi_{103} :\equiv \forall r_{1}\,(\varphi_{29} \to \neg (r_{1} \in \mathbb{R} \land 0 < r_{1} \land \varphi_{5})) \to \varphi_{86} \to (\varphi_{29} \to \forall r_{1}\,(\varphi_{29})) \to \varphi_{93}$
\item $\varphi_{104} :\equiv \varphi_{59} \to \forall d\,(\varphi_{72} \to \neg (d \in \mathbb{R} \land 0 < d \land d \leq r_{1} \land d \leq r_{2})) \to \varphi_{102}$
\item $\varphi_{105} :\equiv \forall r_{2}\,(\varphi_{57} \to \neg (r_{2} \in \mathbb{R} \land 0 < r_{2} \land \varphi_{4})) \to \varphi_{89} \to (\varphi_{57} \to \forall r_{2}\,(\varphi_{57})) \to \varphi_{99}$
\end{abbrevlist}
\end{fullproof}

% 自動生成 (prover/build.js): FLimAdd  ノード数 560
\begin{fullproof}{定理1.8(1)}{$\forall f\,\forall g\,\forall c\,\forall \alpha\,\forall \beta\,(\mathrm{FLim}(f,c,\alpha) \to \mathrm{FLim}(g,c,\beta) \to \mathrm{FLim}((f + g),c,\alpha + \beta))$}{560}\phantomsection\label{tree:FLimAdd}
\noindent$\Pi_{1}$:
\begin{ptree}
\AXC{$\forall f\,\forall c\,\forall \alpha\,(\mathrm{FLim}(f,c,\alpha) \to \mathrm{RFn}(f))$\lbl{[D2]}}
\AXC{$\varphi_{21}$\lbl{A2}}
\BIC{$\forall c\,\forall \alpha\,(\mathrm{FLim}(g,c,\alpha) \to \mathrm{RFn}(g))$}
\AXC{$\varphi_{12}$\lbl{A2}}
\BIC{$\forall \alpha\,(\mathrm{FLim}(g,c,\alpha) \to \mathrm{RFn}(g))$}
\AXC{$\varphi_{8}$\lbl{A2}}
\BIC{$\mathrm{FLim}(g,c,\beta) \to \mathrm{RFn}(g)$}
\end{ptree}
\noindent$\Pi_{2}$:
\begin{ptree}
\AXC{$\varphi_{3}$\lbl{A1}}
\AXC{$\stackrel{\Pi_{1}}{\vdots}$}\noLine
\UIC{$\mathrm{FLim}(g,c,\beta) \to \mathrm{RFn}(g)$}
\AXC{$(\mathrm{FLim}(g,c,\beta) \to \mathrm{RFn}(g)) \to \varphi_{57}$\lbl{A1}}
\BIC{$\varphi_{57}$}
\BIC{$\mathrm{FLim}(f,c,\alpha) \to \mathrm{FLim}(g,c,\beta) \to \mathrm{RFn}(g)$}
\end{ptree}
\noindent$\Pi_{3}$:
\begin{ptree}
\AXC{$\forall f\,\forall c\,\forall \alpha\,(\mathrm{FLim}(f,c,\alpha) \to \mathrm{RFn}(f))$\lbl{[D2]}}
\AXC{$\varphi_{22}$\lbl{A2}}
\BIC{$\forall c\,\forall \alpha\,(\mathrm{FLim}(f,c,\alpha) \to \mathrm{RFn}(f))$}
\AXC{$\varphi_{13}$\lbl{A2}}
\BIC{$\forall \alpha\,(\mathrm{FLim}(f,c,\alpha) \to \mathrm{RFn}(f))$}
\AXC{$\varphi_{9}$\lbl{A2}}
\BIC{$\mathrm{FLim}(f,c,\alpha) \to \mathrm{RFn}(f)$}
\end{ptree}
\noindent$\Pi_{4}$:
\begin{ptree}
\AXC{$\stackrel{\Pi_{3}}{\vdots}$}\noLine
\UIC{$\mathrm{FLim}(f,c,\alpha) \to \mathrm{RFn}(f)$}
\AXC{$(\mathrm{FLim}(f,c,\alpha) \to \mathrm{RFn}(f)) \to \varphi_{31}$\lbl{A1}}
\BIC{$\varphi_{31}$}
\end{ptree}
\noindent$\Pi_{5}$:
\begin{ptree}
\AXC{$\mathrm{FLim}(f,c,\alpha) \to \mathrm{FLim}(f,c,\alpha)$\lbl{A1}}
\AXC{$\stackrel{\Pi_{4}}{\vdots}$}\noLine
\UIC{$\varphi_{31}$}
\BIC{$\mathrm{FLim}(f,c,\alpha) \to \mathrm{FLim}(g,c,\beta) \to \mathrm{RFn}(f)$}
\end{ptree}
\noindent$\Pi_{6}$:
\begin{ptree}
\AXC{$\forall f\,(\varphi_{166})$\lbl{[S17]}}
\AXC{$\forall f\,(\varphi_{166}) \to \varphi_{166}$\lbl{A2}}
\BIC{$\varphi_{166}$}
\AXC{$\varphi_{166} \to \mathrm{RFn}(f) \to \mathrm{RFn}(g) \to \mathrm{RFn}((f + g)) \land \varphi_{11} \land \varphi_{23}$\lbl{A2}}
\BIC{$\mathrm{RFn}(f) \to \mathrm{RFn}(g) \to \mathrm{RFn}((f + g)) \land \varphi_{11} \land \varphi_{23}$}
\AXC{$\varphi_{262}$\lbl{A1}}
\BIC{$\varphi_{235}$}
\end{ptree}
\noindent$\Pi_{7}$:
\begin{ptree}
\AXC{$\stackrel{\Pi_{2}}{\vdots}$}\noLine
\UIC{$\mathrm{FLim}(f,c,\alpha) \to \mathrm{FLim}(g,c,\beta) \to \mathrm{RFn}(g)$}
\AXC{$\stackrel{\Pi_{5}}{\vdots}$}\noLine
\UIC{$\mathrm{FLim}(f,c,\alpha) \to \mathrm{FLim}(g,c,\beta) \to \mathrm{RFn}(f)$}
\AXC{$\stackrel{\Pi_{6}}{\vdots}$}\noLine
\UIC{$\varphi_{235}$}
\BIC{$\varphi_{222}$}
\BIC{$\varphi_{177}$}
\end{ptree}
\noindent$\Pi_{8}$:
\begin{ptree}
\AXC{$\stackrel{\Pi_{7}}{\vdots}$}\noLine
\UIC{$\varphi_{177}$}
\AXC{$\varphi_{177} \to \mathrm{FLim}(f,c,\alpha) \to \mathrm{FLim}(g,c,\beta) \to \varphi_{11}$\lbl{A1}}
\BIC{$\mathrm{FLim}(f,c,\alpha) \to \mathrm{FLim}(g,c,\beta) \to \varphi_{11}$}
\end{ptree}
\noindent$\Pi_{9}$:
\begin{ptree}
\AXC{$\varphi_{11} \to \varphi_{39}$\lbl{A2}}
\AXC{$\stackrel{\Pi_{8}}{\vdots}$}\noLine
\UIC{$\mathrm{FLim}(f,c,\alpha) \to \mathrm{FLim}(g,c,\beta) \to \varphi_{11}$}
\AXC{$(\mathrm{FLim}(f,c,\alpha) \to \mathrm{FLim}(g,c,\beta) \to \varphi_{11}) \to \varphi_{231}$\lbl{A1}}
\BIC{$\varphi_{231}$}
\BIC{$\mathrm{FLim}(f,c,\alpha) \to \mathrm{FLim}(g,c,\beta) \to \varphi_{39}$}
\end{ptree}
\noindent$\Pi_{10}$:
\begin{ptree}
\AXC{$\stackrel{\Pi_{9}}{\vdots}$}\noLine
\UIC{$\mathrm{FLim}(f,c,\alpha) \to \mathrm{FLim}(g,c,\beta) \to \varphi_{39}$}
\AXC{$\varphi_{122}$\lbl{A1}}
\AXC{$\varphi_{105}$\lbl{A2}}
\AXC{$\varphi_{178}$\lbl{A1}}
\AXC{$\varphi_{178} \to \varphi_{105} \to \varphi_{173}$\lbl{A1}}
\BIC{$\varphi_{105} \to \varphi_{173}$}
\BIC{$\varphi_{173}$}
\AXC{$\varphi_{173} \to \varphi_{260}$\lbl{A1}}
\BIC{$\varphi_{260}$}
\BIC{$(\mathrm{FLim}(f,c,\alpha) \to \mathrm{FLim}(g,c,\beta) \to \varphi_{39}) \to \varphi_{164}$}
\BIC{$\varphi_{164}$}
\end{ptree}
\noindent$\Pi_{11}$:
\begin{ptree}
\AXC{$\forall f\,(\varphi_{103})$\lbl{定理1.7(⇒)}}
\AXC{$\forall f\,(\varphi_{103}) \to \forall c\,(\varphi_{91})$\lbl{A2}}
\BIC{$\forall c\,(\varphi_{91})$}
\AXC{$\forall c\,(\varphi_{91}) \to \varphi_{91}$\lbl{A2}}
\BIC{$\varphi_{91}$}
\AXC{$\varphi_{91} \to \forall a\,(\mathrm{FLim}(g,c,\beta) \to \mathrm{Conv}(a,c) \to \varphi_{43})$\lbl{A2}}
\BIC{$\forall a\,(\mathrm{FLim}(g,c,\beta) \to \mathrm{Conv}(a,c) \to \varphi_{43})$}
\end{ptree}
\noindent$\Pi_{12}$:
\begin{ptree}
\AXC{$\mathrm{FLim}(f,c,\alpha) \to \varphi_{6}$\lbl{A1}}
\AXC{$\varphi_{3}$\lbl{A1}}
\AXC{$\stackrel{\Pi_{11}}{\vdots}$}\noLine
\UIC{$\forall a\,(\mathrm{FLim}(g,c,\beta) \to \mathrm{Conv}(a,c) \to \varphi_{43})$}
\AXC{$\varphi_{216}$\lbl{A2}}
\BIC{$\mathrm{FLim}(g,c,\beta) \to \mathrm{Conv}(a,c) \to \varphi_{43}$}
\AXC{$\varphi_{270}$\lbl{A1}}
\BIC{$\varphi_{3} \to (\mathrm{FLim}(f,c,\alpha) \to \varphi_{6}) \to \varphi_{170} \to \varphi_{123}$}
\BIC{$(\mathrm{FLim}(f,c,\alpha) \to \varphi_{6}) \to \varphi_{170} \to \varphi_{123}$}
\BIC{$\varphi_{170} \to \varphi_{123}$}
\end{ptree}
\noindent$\Pi_{13}$:
\begin{ptree}
\AXC{$\varphi_{98} \to \forall n\,(\varphi_{98})$\lbl{A4}}
\AXC{$\varphi_{251}$\lbl{A3}}
\AXC{$\stackrel{\Pi_{10}}{\vdots}$}\noLine
\UIC{$\varphi_{164}$}
\AXC{$\varphi_{164} \to \varphi_{98} \to n \in \mathbb{N} \to a_{n} \in \mathrm{dom}(g) \land a_{n} \neq c$\lbl{A1}}
\BIC{$\varphi_{98} \to n \in \mathbb{N} \to a_{n} \in \mathrm{dom}(g) \land a_{n} \neq c$}
\UIC{$\forall n\,(\varphi_{98} \to n \in \mathbb{N} \to a_{n} \in \mathrm{dom}(g) \land a_{n} \neq c)$}
\AXC{$\varphi_{283}$\lbl{A1}}
\BIC{$\varphi_{251} \to (\varphi_{98} \to \forall n\,(\varphi_{98})) \to \varphi_{170}$}
\BIC{$(\varphi_{98} \to \forall n\,(\varphi_{98})) \to \varphi_{170}$}
\BIC{$\varphi_{170}$}
\AXC{$\stackrel{\Pi_{12}}{\vdots}$}\noLine
\UIC{$\varphi_{170} \to \varphi_{123}$}
\BIC{$\varphi_{123}$}
\end{ptree}
\noindent$\Pi_{14}$:
\begin{ptree}
\AXC{$\varphi_{11} \to \varphi_{39}$\lbl{A2}}
\AXC{$\stackrel{\Pi_{8}}{\vdots}$}\noLine
\UIC{$\mathrm{FLim}(f,c,\alpha) \to \mathrm{FLim}(g,c,\beta) \to \varphi_{11}$}
\AXC{$(\mathrm{FLim}(f,c,\alpha) \to \mathrm{FLim}(g,c,\beta) \to \varphi_{11}) \to \varphi_{231}$\lbl{A1}}
\BIC{$\varphi_{231}$}
\BIC{$\mathrm{FLim}(f,c,\alpha) \to \mathrm{FLim}(g,c,\beta) \to \varphi_{39}$}
\end{ptree}
\noindent$\Pi_{15}$:
\begin{ptree}
\AXC{$\stackrel{\Pi_{14}}{\vdots}$}\noLine
\UIC{$\mathrm{FLim}(f,c,\alpha) \to \mathrm{FLim}(g,c,\beta) \to \varphi_{39}$}
\AXC{$\varphi_{122}$\lbl{A1}}
\AXC{$\varphi_{105}$\lbl{A2}}
\AXC{$\varphi_{178}$\lbl{A1}}
\AXC{$\varphi_{178} \to \varphi_{105} \to \varphi_{173}$\lbl{A1}}
\BIC{$\varphi_{105} \to \varphi_{173}$}
\BIC{$\varphi_{173}$}
\AXC{$\varphi_{173} \to \varphi_{261}$\lbl{A1}}
\BIC{$\varphi_{261}$}
\BIC{$(\mathrm{FLim}(f,c,\alpha) \to \mathrm{FLim}(g,c,\beta) \to \varphi_{39}) \to \varphi_{165}$}
\BIC{$\varphi_{165}$}
\end{ptree}
\noindent$\Pi_{16}$:
\begin{ptree}
\AXC{$\varphi_{98} \to \forall n\,(\varphi_{98})$\lbl{A4}}
\AXC{$\varphi_{252}$\lbl{A3}}
\AXC{$\stackrel{\Pi_{15}}{\vdots}$}\noLine
\UIC{$\varphi_{165}$}
\AXC{$\varphi_{165} \to \varphi_{98} \to n \in \mathbb{N} \to a_{n} \in \mathrm{dom}(f) \land a_{n} \neq c$\lbl{A1}}
\BIC{$\varphi_{98} \to n \in \mathbb{N} \to a_{n} \in \mathrm{dom}(f) \land a_{n} \neq c$}
\UIC{$\forall n\,(\varphi_{98} \to n \in \mathbb{N} \to a_{n} \in \mathrm{dom}(f) \land a_{n} \neq c)$}
\AXC{$\varphi_{284}$\lbl{A1}}
\BIC{$\varphi_{252} \to (\varphi_{98} \to \forall n\,(\varphi_{98})) \to \varphi_{171}$}
\BIC{$(\varphi_{98} \to \forall n\,(\varphi_{98})) \to \varphi_{171}$}
\BIC{$\varphi_{171}$}
\end{ptree}
\noindent$\Pi_{17}$:
\begin{ptree}
\AXC{$\forall f\,(\varphi_{103})$\lbl{定理1.7(⇒)}}
\AXC{$\forall f\,(\varphi_{103}) \to \varphi_{103}$\lbl{A2}}
\BIC{$\varphi_{103}$}
\AXC{$\varphi_{227}$\lbl{A2}}
\BIC{$\forall \alpha\,\forall a\,(\mathrm{FLim}(f,c,\alpha) \to \mathrm{Conv}(a,c) \to \varphi_{44})$}
\AXC{$\varphi_{224}$\lbl{A2}}
\BIC{$\forall a\,(\mathrm{FLim}(f,c,\alpha) \to \mathrm{Conv}(a,c) \to \varphi_{44})$}
\AXC{$\varphi_{217}$\lbl{A2}}
\BIC{$\mathrm{FLim}(f,c,\alpha) \to \mathrm{Conv}(a,c) \to \varphi_{44}$}
\end{ptree}
\noindent$\Pi_{18}$:
\begin{ptree}
\AXC{$\stackrel{\Pi_{17}}{\vdots}$}\noLine
\UIC{$\mathrm{FLim}(f,c,\alpha) \to \mathrm{Conv}(a,c) \to \varphi_{44}$}
\AXC{$(\mathrm{FLim}(f,c,\alpha) \to \mathrm{Conv}(a,c) \to \varphi_{44}) \to \varphi_{258}$\lbl{A1}}
\BIC{$\varphi_{258}$}
\end{ptree}
\noindent$\Pi_{19}$:
\begin{ptree}
\AXC{$\stackrel{\Pi_{16}}{\vdots}$}\noLine
\UIC{$\varphi_{171}$}
\AXC{$\mathrm{FLim}(f,c,\alpha) \to \varphi_{6}$\lbl{A1}}
\AXC{$\mathrm{FLim}(f,c,\alpha) \to \mathrm{FLim}(f,c,\alpha)$\lbl{A1}}
\AXC{$\stackrel{\Pi_{18}}{\vdots}$}\noLine
\UIC{$\varphi_{258}$}
\BIC{$(\mathrm{FLim}(f,c,\alpha) \to \varphi_{6}) \to \varphi_{171} \to \varphi_{124}$}
\BIC{$\varphi_{171} \to \varphi_{124}$}
\BIC{$\varphi_{124}$}
\end{ptree}
\noindent$\Pi_{20}$:
\begin{ptree}
\AXC{$\varphi_{27}$\lbl{定理2.1(1)}}
\AXC{$\varphi_{27} \to \forall \alpha\,(\varphi_{53})$\lbl{A2}}
\BIC{$\forall \alpha\,(\varphi_{53})$}
\AXC{$\forall \alpha\,(\varphi_{53}) \to \varphi_{53}$\lbl{A2}}
\BIC{$\varphi_{53}$}
\AXC{$\varphi_{53} \to \forall \beta\,(\varphi_{77})$\lbl{A2}}
\BIC{$\forall \beta\,(\varphi_{77})$}
\AXC{$\forall \beta\,(\varphi_{77}) \to \varphi_{77}$\lbl{A2}}
\BIC{$\varphi_{77}$}
\AXC{$\varphi_{77} \to \varphi_{124} \to \varphi_{123} \to \varphi_{167}$\lbl{A1}}
\BIC{$\varphi_{124} \to \varphi_{123} \to \varphi_{167}$}
\end{ptree}
\noindent$\Pi_{21}$:
\begin{ptree}
\AXC{$x = x$\lbl{A5}}
\AXC{$x = y \to x = x \to y = x$\lbl{A6}}
\AXC{$(x = y \to x = x \to y = x) \to x = x \to x = y \to y = x$\lbl{A1}}
\BIC{$x = x \to x = y \to y = x$}
\BIC{$x = y \to y = x$}
\UIC{$\forall y\,(x = y \to y = x)$}
\UIC{$\forall x\,\forall y\,(x = y \to y = x)$}
\end{ptree}
\noindent$\Pi_{22}$:
\begin{ptree}
\AXC{$y = z \to x = y \to x = z$\lbl{A6}}
\AXC{$(y = z \to x = y \to x = z) \to x = y \to y = z \to x = z$\lbl{A1}}
\BIC{$x = y \to y = z \to x = z$}
\UIC{$\forall z\,(x = y \to y = z \to x = z)$}
\UIC{$\forall y\,\forall z\,(x = y \to y = z \to x = z)$}
\UIC{$\forall x\,\forall y\,\forall z\,(x = y \to y = z \to x = z)$}
\AXC{$\forall x\,\forall y\,\forall z\,(x = y \to y = z \to x = z) \to \varphi_{45}$\lbl{A2}}
\BIC{$\varphi_{45}$}
\end{ptree}
\noindent$\Pi_{23}$:
\begin{ptree}
\AXC{$y = z \to x = y \to x = z$\lbl{A6}}
\AXC{$(y = z \to x = y \to x = z) \to x = y \to y = z \to x = z$\lbl{A1}}
\BIC{$x = y \to y = z \to x = z$}
\UIC{$\forall z\,(x = y \to y = z \to x = z)$}
\UIC{$\forall y\,\forall z\,(x = y \to y = z \to x = z)$}
\UIC{$\forall x\,\forall y\,\forall z\,(x = y \to y = z \to x = z)$}
\AXC{$\forall x\,\forall y\,\forall z\,(x = y \to y = z \to x = z) \to \varphi_{28}$\lbl{A2}}
\BIC{$\varphi_{28}$}
\end{ptree}
\noindent$\Pi_{24}$:
\begin{ptree}
\AXC{$\stackrel{\Pi_{23}}{\vdots}$}\noLine
\UIC{$\varphi_{28}$}
\AXC{$\varphi_{28} \to \varphi_{120}$\lbl{A2}}
\BIC{$\varphi_{120}$}
\AXC{$\varphi_{120} \to (f + g)(a_{n}) = f(a_{n}) + g(a_{n}) \to \varphi_{139}$\lbl{A2}}
\BIC{$(f + g)(a_{n}) = f(a_{n}) + g(a_{n}) \to \varphi_{139}$}
\end{ptree}
\noindent$\Pi_{25}$:
\begin{ptree}
\AXC{$y = z \to x = y \to x = z$\lbl{A6}}
\AXC{$(y = z \to x = y \to x = z) \to x = y \to y = z \to x = z$\lbl{A1}}
\BIC{$x = y \to y = z \to x = z$}
\UIC{$\forall z\,(x = y \to y = z \to x = z)$}
\UIC{$\forall y\,\forall z\,(x = y \to y = z \to x = z)$}
\UIC{$\forall x\,\forall y\,\forall z\,(x = y \to y = z \to x = z)$}
\AXC{$\forall x\,\forall y\,\forall z\,(x = y \to y = z \to x = z) \to \varphi_{80}$\lbl{A2}}
\BIC{$\varphi_{80}$}
\end{ptree}
\noindent$\Pi_{26}$:
\begin{ptree}
\AXC{$\stackrel{\Pi_{25}}{\vdots}$}\noLine
\UIC{$\varphi_{80}$}
\AXC{$\varphi_{80} \to \varphi_{180}$\lbl{A2}}
\BIC{$\varphi_{180}$}
\AXC{$\varphi_{180} \to f \circ a \oplus g \circ a(n) = f(a_{n}) + g \circ a(n) \to \varphi_{155}$\lbl{A2}}
\BIC{$f \circ a \oplus g \circ a(n) = f(a_{n}) + g \circ a(n) \to \varphi_{155}$}
\end{ptree}
\noindent$\Pi_{27}$:
\begin{ptree}
\AXC{$z = w \to \varphi_{112}$\lbl{A6}}
\UIC{$\forall w\,(z = w \to \varphi_{112})$}
\UIC{$\forall z\,\forall w\,(z = w \to \varphi_{112})$}
\AXC{$\forall z\,\forall w\,(z = w \to \varphi_{112}) \to \varphi_{179}$\lbl{A2}}
\BIC{$\varphi_{179}$}
\AXC{$\varphi_{179} \to g \circ a(n) = g(a_{n}) \to \varphi_{159}$\lbl{A2}}
\BIC{$g \circ a(n) = g(a_{n}) \to \varphi_{159}$}
\end{ptree}
\noindent$\Pi_{28}$:
\begin{ptree}
\AXC{$\varphi_{11} \to \varphi_{39}$\lbl{A2}}
\AXC{$\stackrel{\Pi_{8}}{\vdots}$}\noLine
\UIC{$\mathrm{FLim}(f,c,\alpha) \to \mathrm{FLim}(g,c,\beta) \to \varphi_{11}$}
\AXC{$(\mathrm{FLim}(f,c,\alpha) \to \mathrm{FLim}(g,c,\beta) \to \varphi_{11}) \to \varphi_{231}$\lbl{A1}}
\BIC{$\varphi_{231}$}
\BIC{$\mathrm{FLim}(f,c,\alpha) \to \mathrm{FLim}(g,c,\beta) \to \varphi_{39}$}
\end{ptree}
\noindent$\Pi_{29}$:
\begin{ptree}
\AXC{$\stackrel{\Pi_{28}}{\vdots}$}\noLine
\UIC{$\mathrm{FLim}(f,c,\alpha) \to \mathrm{FLim}(g,c,\beta) \to \varphi_{39}$}
\AXC{$\varphi_{122}$\lbl{A1}}
\AXC{$\varphi_{105}$\lbl{A2}}
\AXC{$\varphi_{178}$\lbl{A1}}
\AXC{$\varphi_{178} \to \varphi_{105} \to \varphi_{173}$\lbl{A1}}
\BIC{$\varphi_{105} \to \varphi_{173}$}
\BIC{$\varphi_{173}$}
\AXC{$\varphi_{173} \to \varphi_{256}$\lbl{A1}}
\BIC{$\varphi_{256}$}
\BIC{$(\mathrm{FLim}(f,c,\alpha) \to \mathrm{FLim}(g,c,\beta) \to \varphi_{39}) \to \varphi_{140}$}
\BIC{$\varphi_{140}$}
\end{ptree}
\noindent$\Pi_{30}$:
\begin{ptree}
\AXC{$\varphi_{98} \to \forall n\,(\varphi_{98})$\lbl{A4}}
\AXC{$\varphi_{246}$\lbl{A3}}
\AXC{$\stackrel{\Pi_{29}}{\vdots}$}\noLine
\UIC{$\varphi_{140}$}
\AXC{$\varphi_{140} \to \varphi_{98} \to n \in \mathbb{N} \to a_{n} \in \mathrm{dom}(g)$\lbl{A1}}
\BIC{$\varphi_{98} \to n \in \mathbb{N} \to a_{n} \in \mathrm{dom}(g)$}
\UIC{$\forall n\,(\varphi_{98} \to n \in \mathbb{N} \to a_{n} \in \mathrm{dom}(g))$}
\AXC{$\varphi_{279}$\lbl{A1}}
\BIC{$\varphi_{246} \to (\varphi_{98} \to \forall n\,(\varphi_{98})) \to \varphi_{143}$}
\BIC{$(\varphi_{98} \to \forall n\,(\varphi_{98})) \to \varphi_{143}$}
\BIC{$\varphi_{143}$}
\end{ptree}
\noindent$\Pi_{31}$:
\begin{ptree}
\AXC{$\forall f\,\forall a\,(\mathrm{RFn}(f) \to \mathrm{Seq}(a) \to \varphi_{102})$\lbl{[S8]}}
\AXC{$\varphi_{240}$\lbl{A2}}
\BIC{$\forall a\,(\mathrm{RFn}(g) \to \mathrm{Seq}(a) \to \varphi_{101})$}
\AXC{$\varphi_{238}$\lbl{A2}}
\BIC{$\mathrm{RFn}(g) \to \mathrm{Seq}(a) \to \varphi_{101}$}
\AXC{$(\mathrm{RFn}(g) \to \mathrm{Seq}(a) \to \varphi_{101}) \to \varphi_{263}$\lbl{A1}}
\BIC{$\varphi_{263}$}
\end{ptree}
\noindent$\Pi_{32}$:
\begin{ptree}
\AXC{$\stackrel{\Pi_{30}}{\vdots}$}\noLine
\UIC{$\varphi_{143}$}
\AXC{$\varphi_{5}$\lbl{A1}}
\AXC{$\stackrel{\Pi_{2}}{\vdots}$}\noLine
\UIC{$\mathrm{FLim}(f,c,\alpha) \to \mathrm{FLim}(g,c,\beta) \to \mathrm{RFn}(g)$}
\AXC{$\stackrel{\Pi_{31}}{\vdots}$}\noLine
\UIC{$\varphi_{263}$}
\BIC{$\varphi_{5} \to \varphi_{143} \to \varphi_{201}$}
\BIC{$\varphi_{143} \to \varphi_{201}$}
\BIC{$\varphi_{201}$}
\AXC{$\varphi_{201} \to \varphi_{186}$\lbl{A1}}
\BIC{$\varphi_{186}$}
\end{ptree}
\noindent$\Pi_{33}$:
\begin{ptree}
\AXC{$\varphi_{122}$\lbl{A1}}
\AXC{$\varphi_{85}$\lbl{A2}}
\AXC{$\stackrel{\Pi_{32}}{\vdots}$}\noLine
\UIC{$\varphi_{186}$}
\AXC{$\varphi_{186} \to \varphi_{85} \to \varphi_{182}$\lbl{A1}}
\BIC{$\varphi_{85} \to \varphi_{182}$}
\BIC{$\varphi_{182}$}
\AXC{$\varphi_{182} \to \varphi_{122} \to \varphi_{168}$\lbl{A1}}
\BIC{$\varphi_{122} \to \varphi_{168}$}
\BIC{$\varphi_{168}$}
\AXC{$\varphi_{168} \to \varphi_{267}$\lbl{A1}}
\BIC{$\varphi_{267}$}
\end{ptree}
\noindent$\Pi_{34}$:
\begin{ptree}
\AXC{$y = z \to x = y \to x = z$\lbl{A6}}
\AXC{$(y = z \to x = y \to x = z) \to x = y \to y = z \to x = z$\lbl{A1}}
\BIC{$x = y \to y = z \to x = z$}
\UIC{$\forall z\,(x = y \to y = z \to x = z)$}
\UIC{$\forall y\,\forall z\,(x = y \to y = z \to x = z)$}
\UIC{$\forall x\,\forall y\,\forall z\,(x = y \to y = z \to x = z)$}
\AXC{$\forall x\,\forall y\,\forall z\,(x = y \to y = z \to x = z) \to \varphi_{80}$\lbl{A2}}
\BIC{$\varphi_{80}$}
\end{ptree}
\noindent$\Pi_{35}$:
\begin{ptree}
\AXC{$\stackrel{\Pi_{34}}{\vdots}$}\noLine
\UIC{$\varphi_{80}$}
\AXC{$\varphi_{80} \to \varphi_{190}$\lbl{A2}}
\BIC{$\varphi_{190}$}
\AXC{$\varphi_{190} \to f \circ a \oplus g \circ a(n) = f \circ a(n) + g \circ a(n) \to \varphi_{172}$\lbl{A2}}
\BIC{$f \circ a \oplus g \circ a(n) = f \circ a(n) + g \circ a(n) \to \varphi_{172}$}
\end{ptree}
\noindent$\Pi_{36}$:
\begin{ptree}
\AXC{$z = w \to \varphi_{136}$\lbl{A6}}
\UIC{$\forall w\,(z = w \to \varphi_{136})$}
\UIC{$\forall z\,\forall w\,(z = w \to \varphi_{136})$}
\AXC{$\forall z\,\forall w\,(z = w \to \varphi_{136}) \to \varphi_{194}$\lbl{A2}}
\BIC{$\varphi_{194}$}
\AXC{$\varphi_{194} \to f \circ a(n) = f(a_{n}) \to \varphi_{184}$\lbl{A2}}
\BIC{$f \circ a(n) = f(a_{n}) \to \varphi_{184}$}
\end{ptree}
\noindent$\Pi_{37}$:
\begin{ptree}
\AXC{$\varphi_{11} \to \varphi_{39}$\lbl{A2}}
\AXC{$\stackrel{\Pi_{8}}{\vdots}$}\noLine
\UIC{$\mathrm{FLim}(f,c,\alpha) \to \mathrm{FLim}(g,c,\beta) \to \varphi_{11}$}
\AXC{$(\mathrm{FLim}(f,c,\alpha) \to \mathrm{FLim}(g,c,\beta) \to \varphi_{11}) \to \varphi_{231}$\lbl{A1}}
\BIC{$\varphi_{231}$}
\BIC{$\mathrm{FLim}(f,c,\alpha) \to \mathrm{FLim}(g,c,\beta) \to \varphi_{39}$}
\end{ptree}
\noindent$\Pi_{38}$:
\begin{ptree}
\AXC{$\stackrel{\Pi_{37}}{\vdots}$}\noLine
\UIC{$\mathrm{FLim}(f,c,\alpha) \to \mathrm{FLim}(g,c,\beta) \to \varphi_{39}$}
\AXC{$\varphi_{122}$\lbl{A1}}
\AXC{$\varphi_{105}$\lbl{A2}}
\AXC{$\varphi_{178}$\lbl{A1}}
\AXC{$\varphi_{178} \to \varphi_{105} \to \varphi_{173}$\lbl{A1}}
\BIC{$\varphi_{105} \to \varphi_{173}$}
\BIC{$\varphi_{173}$}
\AXC{$\varphi_{173} \to \varphi_{257}$\lbl{A1}}
\BIC{$\varphi_{257}$}
\BIC{$(\mathrm{FLim}(f,c,\alpha) \to \mathrm{FLim}(g,c,\beta) \to \varphi_{39}) \to \varphi_{141}$}
\BIC{$\varphi_{141}$}
\end{ptree}
\noindent$\Pi_{39}$:
\begin{ptree}
\AXC{$\varphi_{98} \to \forall n\,(\varphi_{98})$\lbl{A4}}
\AXC{$\varphi_{247}$\lbl{A3}}
\AXC{$\stackrel{\Pi_{38}}{\vdots}$}\noLine
\UIC{$\varphi_{141}$}
\AXC{$\varphi_{141} \to \varphi_{98} \to n \in \mathbb{N} \to a_{n} \in \mathrm{dom}(f)$\lbl{A1}}
\BIC{$\varphi_{98} \to n \in \mathbb{N} \to a_{n} \in \mathrm{dom}(f)$}
\UIC{$\forall n\,(\varphi_{98} \to n \in \mathbb{N} \to a_{n} \in \mathrm{dom}(f))$}
\AXC{$\varphi_{280}$\lbl{A1}}
\BIC{$\varphi_{247} \to (\varphi_{98} \to \forall n\,(\varphi_{98})) \to \varphi_{144}$}
\BIC{$(\varphi_{98} \to \forall n\,(\varphi_{98})) \to \varphi_{144}$}
\BIC{$\varphi_{144}$}
\end{ptree}
\noindent$\Pi_{40}$:
\begin{ptree}
\AXC{$\forall f\,\forall a\,(\mathrm{RFn}(f) \to \mathrm{Seq}(a) \to \varphi_{102})$\lbl{[S8]}}
\AXC{$\varphi_{241}$\lbl{A2}}
\BIC{$\forall a\,(\mathrm{RFn}(f) \to \mathrm{Seq}(a) \to \varphi_{102})$}
\AXC{$\varphi_{239}$\lbl{A2}}
\BIC{$\mathrm{RFn}(f) \to \mathrm{Seq}(a) \to \varphi_{102}$}
\AXC{$(\mathrm{RFn}(f) \to \mathrm{Seq}(a) \to \varphi_{102}) \to \varphi_{264}$\lbl{A1}}
\BIC{$\varphi_{264}$}
\end{ptree}
\noindent$\Pi_{41}$:
\begin{ptree}
\AXC{$\stackrel{\Pi_{39}}{\vdots}$}\noLine
\UIC{$\varphi_{144}$}
\AXC{$\varphi_{5}$\lbl{A1}}
\AXC{$\stackrel{\Pi_{5}}{\vdots}$}\noLine
\UIC{$\mathrm{FLim}(f,c,\alpha) \to \mathrm{FLim}(g,c,\beta) \to \mathrm{RFn}(f)$}
\AXC{$\stackrel{\Pi_{40}}{\vdots}$}\noLine
\UIC{$\varphi_{264}$}
\BIC{$\varphi_{5} \to \varphi_{144} \to \varphi_{202}$}
\BIC{$\varphi_{144} \to \varphi_{202}$}
\BIC{$\varphi_{202}$}
\AXC{$\varphi_{202} \to \varphi_{187}$\lbl{A1}}
\BIC{$\varphi_{187}$}
\end{ptree}
\noindent$\Pi_{42}$:
\begin{ptree}
\AXC{$\varphi_{122}$\lbl{A1}}
\AXC{$\varphi_{86}$\lbl{A2}}
\AXC{$\stackrel{\Pi_{41}}{\vdots}$}\noLine
\UIC{$\varphi_{187}$}
\AXC{$\varphi_{187} \to \varphi_{86} \to \varphi_{183}$\lbl{A1}}
\BIC{$\varphi_{86} \to \varphi_{183}$}
\BIC{$\varphi_{183}$}
\AXC{$\varphi_{183} \to \varphi_{122} \to \varphi_{169}$\lbl{A1}}
\BIC{$\varphi_{122} \to \varphi_{169}$}
\BIC{$\varphi_{169}$}
\AXC{$\varphi_{169} \to \varphi_{269}$\lbl{A1}}
\BIC{$\varphi_{269}$}
\end{ptree}
\noindent$\Pi_{43}$:
\begin{ptree}
\AXC{$\forall a\,\forall \alpha\,(\mathrm{Conv}(a,\alpha) \to \mathrm{Seq}(a))$\lbl{[C1]}}
\AXC{$\varphi_{46}$\lbl{A2}}
\BIC{$\forall \alpha\,(\mathrm{Conv}(g \circ a,\alpha) \to \mathrm{Seq}(g \circ a))$}
\AXC{$\varphi_{65}$\lbl{A2}}
\BIC{$\mathrm{Conv}(g \circ a,\beta) \to \mathrm{Seq}(g \circ a)$}
\AXC{$(\mathrm{Conv}(g \circ a,\beta) \to \mathrm{Seq}(g \circ a)) \to \varphi_{123} \to \varphi_{118}$\lbl{A1}}
\BIC{$\varphi_{123} \to \varphi_{118}$}
\end{ptree}
\noindent$\Pi_{44}$:
\begin{ptree}
\AXC{$\forall a\,\forall \alpha\,(\mathrm{Conv}(a,\alpha) \to \mathrm{Seq}(a))$\lbl{[C1]}}
\AXC{$\varphi_{47}$\lbl{A2}}
\BIC{$\forall \alpha\,(\mathrm{Conv}(f \circ a,\alpha) \to \mathrm{Seq}(f \circ a))$}
\AXC{$\varphi_{66}$\lbl{A2}}
\BIC{$\mathrm{Conv}(f \circ a,\alpha) \to \mathrm{Seq}(f \circ a)$}
\AXC{$(\mathrm{Conv}(f \circ a,\alpha) \to \mathrm{Seq}(f \circ a)) \to \varphi_{124} \to \varphi_{119}$\lbl{A1}}
\BIC{$\varphi_{124} \to \varphi_{119}$}
\end{ptree}
\noindent$\Pi_{45}$:
\begin{ptree}
\AXC{$\stackrel{\Pi_{13}}{\vdots}$}\noLine
\UIC{$\varphi_{123}$}
\AXC{$\stackrel{\Pi_{43}}{\vdots}$}\noLine
\UIC{$\varphi_{123} \to \varphi_{118}$}
\BIC{$\varphi_{118}$}
\AXC{$\stackrel{\Pi_{19}}{\vdots}$}\noLine
\UIC{$\varphi_{124}$}
\AXC{$\stackrel{\Pi_{44}}{\vdots}$}\noLine
\UIC{$\varphi_{124} \to \varphi_{119}$}
\BIC{$\varphi_{119}$}
\AXC{$\varphi_{78}$\lbl{[S2]}}
\AXC{$\varphi_{78} \to \varphi_{133}$\lbl{A2}}
\BIC{$\varphi_{133}$}
\AXC{$\varphi_{133} \to \varphi_{191}$\lbl{A2}}
\BIC{$\varphi_{191}$}
\AXC{$\varphi_{191} \to \varphi_{119} \to \varphi_{118} \to \varphi_{237}$\lbl{A1}}
\BIC{$\varphi_{119} \to \varphi_{118} \to \varphi_{237}$}
\BIC{$\varphi_{118} \to \varphi_{237}$}
\BIC{$\varphi_{237}$}
\end{ptree}
\noindent$\Pi_{46}$:
\begin{ptree}
\AXC{$\varphi_{122}$\lbl{A1}}
\AXC{$\varphi_{206}$\lbl{A2}}
\AXC{$\stackrel{\Pi_{45}}{\vdots}$}\noLine
\UIC{$\varphi_{237}$}
\AXC{$\varphi_{237} \to \varphi_{229}$\lbl{A1}}
\BIC{$\varphi_{229}$}
\AXC{$\varphi_{229} \to \varphi_{206} \to \varphi_{225}$\lbl{A1}}
\BIC{$\varphi_{206} \to \varphi_{225}$}
\BIC{$\varphi_{225}$}
\AXC{$\varphi_{225} \to \varphi_{122} \to \varphi_{219}$\lbl{A1}}
\BIC{$\varphi_{122} \to \varphi_{219}$}
\BIC{$\varphi_{219}$}
\AXC{$\varphi_{219} \to \varphi_{275}$\lbl{A1}}
\BIC{$\varphi_{275}$}
\end{ptree}
\noindent$\Pi_{47}$:
\begin{ptree}
\AXC{$\stackrel{\Pi_{36}}{\vdots}$}\noLine
\UIC{$f \circ a(n) = f(a_{n}) \to \varphi_{184}$}
\AXC{$x = x$\lbl{A5}}
\UIC{$\forall x\,(x = x)$}
\AXC{$\varphi_{73}$\lbl{A2}}
\BIC{$f \circ a(n) + g \circ a(n) = f \circ a(n) + g \circ a(n)$}
\AXC{$\stackrel{\Pi_{42}}{\vdots}$}\noLine
\UIC{$\varphi_{269}$}
\BIC{$(f \circ a(n) = f(a_{n}) \to \varphi_{184}) \to \varphi_{215}$}
\BIC{$\varphi_{215}$}
\AXC{$\stackrel{\Pi_{46}}{\vdots}$}\noLine
\UIC{$\varphi_{275}$}
\BIC{$(f \circ a \oplus g \circ a(n) = f \circ a(n) + g \circ a(n) \to \varphi_{172}) \to \varphi_{214}$}
\end{ptree}
\noindent$\Pi_{48}$:
\begin{ptree}
\AXC{$\stackrel{\Pi_{35}}{\vdots}$}\noLine
\UIC{$f \circ a \oplus g \circ a(n) = f \circ a(n) + g \circ a(n) \to \varphi_{172}$}
\AXC{$\stackrel{\Pi_{47}}{\vdots}$}\noLine
\UIC{$(f \circ a \oplus g \circ a(n) = f \circ a(n) + g \circ a(n) \to \varphi_{172}) \to \varphi_{214}$}
\BIC{$\varphi_{214}$}
\AXC{$\varphi_{214} \to \varphi_{273}$\lbl{A1}}
\BIC{$\varphi_{273}$}
\end{ptree}
\noindent$\Pi_{49}$:
\begin{ptree}
\AXC{$\stackrel{\Pi_{27}}{\vdots}$}\noLine
\UIC{$g \circ a(n) = g(a_{n}) \to \varphi_{159}$}
\AXC{$x = x$\lbl{A5}}
\UIC{$\forall x\,(x = x)$}
\AXC{$\varphi_{63}$\lbl{A2}}
\BIC{$f(a_{n}) + g \circ a(n) = f(a_{n}) + g \circ a(n)$}
\AXC{$\stackrel{\Pi_{33}}{\vdots}$}\noLine
\UIC{$\varphi_{267}$}
\BIC{$(g \circ a(n) = g(a_{n}) \to \varphi_{159}) \to \varphi_{207}$}
\BIC{$\varphi_{207}$}
\AXC{$\stackrel{\Pi_{48}}{\vdots}$}\noLine
\UIC{$\varphi_{273}$}
\BIC{$(f \circ a \oplus g \circ a(n) = f(a_{n}) + g \circ a(n) \to \varphi_{155}) \to \varphi_{210}$}
\end{ptree}
\noindent$\Pi_{50}$:
\begin{ptree}
\AXC{$\stackrel{\Pi_{26}}{\vdots}$}\noLine
\UIC{$f \circ a \oplus g \circ a(n) = f(a_{n}) + g \circ a(n) \to \varphi_{155}$}
\AXC{$\stackrel{\Pi_{49}}{\vdots}$}\noLine
\UIC{$(f \circ a \oplus g \circ a(n) = f(a_{n}) + g \circ a(n) \to \varphi_{155}) \to \varphi_{210}$}
\BIC{$\varphi_{210}$}
\AXC{$\varphi_{210} \to \varphi_{156} \to \varphi_{211}$\lbl{A1}}
\BIC{$\varphi_{156} \to \varphi_{211}$}
\end{ptree}
\noindent$\Pi_{51}$:
\begin{ptree}
\AXC{$\stackrel{\Pi_{21}}{\vdots}$}\noLine
\UIC{$\forall x\,\forall y\,(x = y \to y = x)$}
\AXC{$\varphi_{92}$\lbl{A2}}
\BIC{$\forall y\,(f \circ a \oplus g \circ a(n) = y \to y = f \circ a \oplus g \circ a(n))$}
\AXC{$\varphi_{234}$\lbl{A2}}
\BIC{$\varphi_{156}$}
\AXC{$\stackrel{\Pi_{50}}{\vdots}$}\noLine
\UIC{$\varphi_{156} \to \varphi_{211}$}
\BIC{$\varphi_{211}$}
\end{ptree}
\noindent$\Pi_{52}$:
\begin{ptree}
\AXC{$\stackrel{\Pi_{7}}{\vdots}$}\noLine
\UIC{$\varphi_{177}$}
\AXC{$\varphi_{177} \to \mathrm{FLim}(f,c,\alpha) \to \mathrm{FLim}(g,c,\beta) \to \varphi_{23}$\lbl{A1}}
\BIC{$\mathrm{FLim}(f,c,\alpha) \to \mathrm{FLim}(g,c,\beta) \to \varphi_{23}$}
\AXC{$(\mathrm{FLim}(f,c,\alpha) \to \mathrm{FLim}(g,c,\beta) \to \varphi_{23}) \to \varphi_{242}$\lbl{A1}}
\BIC{$\varphi_{242}$}
\end{ptree}
\noindent$\Pi_{53}$:
\begin{ptree}
\AXC{$\varphi_{122}$\lbl{A1}}
\AXC{$\varphi_{105}$\lbl{A2}}
\AXC{$\varphi_{178}$\lbl{A1}}
\AXC{$\varphi_{178} \to \varphi_{105} \to \varphi_{173}$\lbl{A1}}
\BIC{$\varphi_{105} \to \varphi_{173}$}
\BIC{$\varphi_{173}$}
\AXC{$\varphi_{23} \to \varphi_{64}$\lbl{A2}}
\AXC{$\stackrel{\Pi_{52}}{\vdots}$}\noLine
\UIC{$\varphi_{242}$}
\BIC{$\mathrm{FLim}(f,c,\alpha) \to \mathrm{FLim}(g,c,\beta) \to \varphi_{64}$}
\AXC{$\varphi_{272}$\lbl{A1}}
\BIC{$\varphi_{173} \to \varphi_{122} \to \varphi_{193}$}
\BIC{$\varphi_{122} \to \varphi_{193}$}
\BIC{$\varphi_{193}$}
\end{ptree}
\noindent$\Pi_{54}$:
\begin{ptree}
\AXC{$\stackrel{\Pi_{24}}{\vdots}$}\noLine
\UIC{$(f + g)(a_{n}) = f(a_{n}) + g(a_{n}) \to \varphi_{139}$}
\AXC{$\stackrel{\Pi_{51}}{\vdots}$}\noLine
\UIC{$\varphi_{211}$}
\AXC{$\stackrel{\Pi_{53}}{\vdots}$}\noLine
\UIC{$\varphi_{193}$}
\AXC{$\varphi_{278}$\lbl{A1}}
\BIC{$\varphi_{211} \to ((f + g)(a_{n}) = f(a_{n}) + g(a_{n}) \to \varphi_{139}) \to \varphi_{195}$}
\BIC{$((f + g)(a_{n}) = f(a_{n}) + g(a_{n}) \to \varphi_{139}) \to \varphi_{195}$}
\BIC{$\varphi_{195}$}
\end{ptree}
\noindent$\Pi_{55}$:
\begin{ptree}
\AXC{$\varphi_{122}$\lbl{A1}}
\AXC{$\varphi_{105}$\lbl{A2}}
\AXC{$\varphi_{178}$\lbl{A1}}
\AXC{$\varphi_{178} \to \varphi_{105} \to \varphi_{173}$\lbl{A1}}
\BIC{$\varphi_{105} \to \varphi_{173}$}
\BIC{$\varphi_{173}$}
\AXC{$\varphi_{173} \to \varphi_{122} \to \varphi_{152}$\lbl{A1}}
\BIC{$\varphi_{122} \to \varphi_{152}$}
\BIC{$\varphi_{152}$}
\AXC{$\varphi_{152} \to \varphi_{98} \to n \in \mathbb{N} \to a_{n} \in \mathrm{dom}((f + g))$\lbl{A1}}
\BIC{$\varphi_{98} \to n \in \mathbb{N} \to a_{n} \in \mathrm{dom}((f + g))$}
\UIC{$\forall n\,(\varphi_{98} \to n \in \mathbb{N} \to a_{n} \in \mathrm{dom}((f + g)))$}
\end{ptree}
\noindent$\Pi_{56}$:
\begin{ptree}
\AXC{$\varphi_{98} \to \forall n\,(\varphi_{98})$\lbl{A4}}
\AXC{$\varphi_{249}$\lbl{A3}}
\AXC{$\stackrel{\Pi_{55}}{\vdots}$}\noLine
\UIC{$\forall n\,(\varphi_{98} \to n \in \mathbb{N} \to a_{n} \in \mathrm{dom}((f + g)))$}
\AXC{$\varphi_{281}$\lbl{A1}}
\BIC{$\varphi_{249} \to (\varphi_{98} \to \forall n\,(\varphi_{98})) \to \varphi_{158}$}
\BIC{$(\varphi_{98} \to \forall n\,(\varphi_{98})) \to \varphi_{158}$}
\BIC{$\varphi_{158}$}
\end{ptree}
\noindent$\Pi_{57}$:
\begin{ptree}
\AXC{$\stackrel{\Pi_{7}}{\vdots}$}\noLine
\UIC{$\varphi_{177}$}
\AXC{$\varphi_{177} \to \varphi_{4}$\lbl{A1}}
\BIC{$\varphi_{4}$}
\end{ptree}
\noindent$\Pi_{58}$:
\begin{ptree}
\AXC{$\forall f\,\forall a\,(\mathrm{RFn}(f) \to \mathrm{Seq}(a) \to \varphi_{102})$\lbl{[S8]}}
\AXC{$\forall f\,\forall a\,(\mathrm{RFn}(f) \to \mathrm{Seq}(a) \to \varphi_{102}) \to \forall a\,(\varphi_{175})$\lbl{A2}}
\BIC{$\forall a\,(\varphi_{175})$}
\AXC{$\forall a\,(\varphi_{175}) \to \varphi_{175}$\lbl{A2}}
\BIC{$\varphi_{175}$}
\end{ptree}
\noindent$\Pi_{59}$:
\begin{ptree}
\AXC{$\stackrel{\Pi_{56}}{\vdots}$}\noLine
\UIC{$\varphi_{158}$}
\AXC{$\varphi_{5}$\lbl{A1}}
\AXC{$\stackrel{\Pi_{57}}{\vdots}$}\noLine
\UIC{$\varphi_{4}$}
\AXC{$\stackrel{\Pi_{58}}{\vdots}$}\noLine
\UIC{$\varphi_{175}$}
\AXC{$\varphi_{175} \to \varphi_{4} \to \varphi_{5} \to \varphi_{158} \to \varphi_{220}$\lbl{A1}}
\BIC{$\varphi_{4} \to \varphi_{5} \to \varphi_{158} \to \varphi_{220}$}
\BIC{$\varphi_{5} \to \varphi_{158} \to \varphi_{220}$}
\BIC{$\varphi_{158} \to \varphi_{220}$}
\BIC{$\varphi_{220}$}
\end{ptree}
\noindent$\Pi_{60}$:
\begin{ptree}
\AXC{$\varphi_{122}$\lbl{A1}}
\AXC{$\varphi_{125}$\lbl{A2}}
\AXC{$\stackrel{\Pi_{59}}{\vdots}$}\noLine
\UIC{$\varphi_{220}$}
\AXC{$\varphi_{220} \to \varphi_{197}$\lbl{A1}}
\BIC{$\varphi_{197}$}
\AXC{$\varphi_{197} \to \varphi_{125} \to \varphi_{196}$\lbl{A1}}
\BIC{$\varphi_{125} \to \varphi_{196}$}
\BIC{$\varphi_{196}$}
\AXC{$\varphi_{196} \to \varphi_{122} \to \varphi_{188}$\lbl{A1}}
\BIC{$\varphi_{122} \to \varphi_{188}$}
\BIC{$\varphi_{188}$}
\end{ptree}
\noindent$\Pi_{61}$:
\begin{ptree}
\AXC{$\stackrel{\Pi_{54}}{\vdots}$}\noLine
\UIC{$\varphi_{195}$}
\AXC{$\stackrel{\Pi_{60}}{\vdots}$}\noLine
\UIC{$\varphi_{188}$}
\AXC{$\varphi_{188} \to \varphi_{195} \to ((f + g) \circ a(n) = (f + g)(a_{n}) \to \varphi_{126}) \to \varphi_{198}$\lbl{A1}}
\BIC{$\varphi_{195} \to ((f + g) \circ a(n) = (f + g)(a_{n}) \to \varphi_{126}) \to \varphi_{198}$}
\BIC{$((f + g) \circ a(n) = (f + g)(a_{n}) \to \varphi_{126}) \to \varphi_{198}$}
\end{ptree}
\noindent$\Pi_{62}$:
\begin{ptree}
\AXC{$\stackrel{\Pi_{22}}{\vdots}$}\noLine
\UIC{$\varphi_{45}$}
\AXC{$\varphi_{45} \to \varphi_{106}$\lbl{A2}}
\BIC{$\varphi_{106}$}
\AXC{$\varphi_{106} \to (f + g) \circ a(n) = (f + g)(a_{n}) \to \varphi_{126}$\lbl{A2}}
\BIC{$(f + g) \circ a(n) = (f + g)(a_{n}) \to \varphi_{126}$}
\AXC{$\stackrel{\Pi_{61}}{\vdots}$}\noLine
\UIC{$((f + g) \circ a(n) = (f + g)(a_{n}) \to \varphi_{126}) \to \varphi_{198}$}
\BIC{$\varphi_{198}$}
\end{ptree}
\noindent$\Pi_{63}$:
\begin{ptree}
\AXC{$\stackrel{\Pi_{21}}{\vdots}$}\noLine
\UIC{$\forall x\,\forall y\,(x = y \to y = x)$}
\AXC{$\varphi_{58}$\lbl{A2}}
\BIC{$\forall y\,((f + g) \circ a(n) = y \to y = (f + g) \circ a(n))$}
\AXC{$\varphi_{218}$\lbl{A2}}
\BIC{$\varphi_{135}$}
\AXC{$\stackrel{\Pi_{62}}{\vdots}$}\noLine
\UIC{$\varphi_{198}$}
\AXC{$\varphi_{198} \to \varphi_{135} \to \varphi_{199}$\lbl{A1}}
\BIC{$\varphi_{135} \to \varphi_{199}$}
\BIC{$\varphi_{199}$}
\end{ptree}
\noindent$\Pi_{64}$:
\begin{ptree}
\AXC{$\varphi_{203} \to \varphi_{208}$\lbl{A3}}
\AXC{$\stackrel{\Pi_{63}}{\vdots}$}\noLine
\UIC{$\varphi_{199}$}
\AXC{$\varphi_{199} \to \varphi_{98} \to n \in \mathbb{N} \to f \circ a \oplus g \circ a(n) = (f + g) \circ a(n)$\lbl{A1}}
\BIC{$\varphi_{98} \to n \in \mathbb{N} \to f \circ a \oplus g \circ a(n) = (f + g) \circ a(n)$}
\UIC{$\varphi_{203}$}
\AXC{$\varphi_{203} \to (\varphi_{203} \to \varphi_{208}) \to (\varphi_{98} \to \forall n\,(\varphi_{98})) \to \varphi_{204}$\lbl{A1}}
\BIC{$(\varphi_{203} \to \varphi_{208}) \to (\varphi_{98} \to \forall n\,(\varphi_{98})) \to \varphi_{204}$}
\BIC{$(\varphi_{98} \to \forall n\,(\varphi_{98})) \to \varphi_{204}$}
\end{ptree}
\noindent$\Pi_{65}$:
\begin{ptree}
\AXC{$\stackrel{\Pi_{59}}{\vdots}$}\noLine
\UIC{$\varphi_{220}$}
\AXC{$\varphi_{220} \to \varphi_{131}$\lbl{A1}}
\BIC{$\varphi_{131}$}
\AXC{$\varphi_{51}$\lbl{[C13]}}
\AXC{$\varphi_{51} \to \varphi_{134}$\lbl{A2}}
\BIC{$\varphi_{134}$}
\AXC{$\varphi_{134} \to \varphi_{200}$\lbl{A2}}
\BIC{$\varphi_{200}$}
\AXC{$\varphi_{200} \to \varphi_{213}$\lbl{A2}}
\BIC{$\varphi_{213}$}
\AXC{$\varphi_{213} \to \varphi_{131} \to \varphi_{204} \to \varphi_{167} \to \varphi_{147}$\lbl{A1}}
\BIC{$\varphi_{131} \to \varphi_{204} \to \varphi_{167} \to \varphi_{147}$}
\BIC{$\varphi_{204} \to \varphi_{167} \to \varphi_{147}$}
\end{ptree}
\noindent$\Pi_{66}$:
\begin{ptree}
\AXC{$\stackrel{\Pi_{13}}{\vdots}$}\noLine
\UIC{$\varphi_{123}$}
\AXC{$\stackrel{\Pi_{19}}{\vdots}$}\noLine
\UIC{$\varphi_{124}$}
\AXC{$\stackrel{\Pi_{20}}{\vdots}$}\noLine
\UIC{$\varphi_{124} \to \varphi_{123} \to \varphi_{167}$}
\BIC{$\varphi_{123} \to \varphi_{167}$}
\BIC{$\varphi_{167}$}
\AXC{$\varphi_{98} \to \forall n\,(\varphi_{98})$\lbl{A4}}
\AXC{$\stackrel{\Pi_{64}}{\vdots}$}\noLine
\UIC{$(\varphi_{98} \to \forall n\,(\varphi_{98})) \to \varphi_{204}$}
\BIC{$\varphi_{204}$}
\AXC{$\stackrel{\Pi_{65}}{\vdots}$}\noLine
\UIC{$\varphi_{204} \to \varphi_{167} \to \varphi_{147}$}
\BIC{$\varphi_{167} \to \varphi_{147}$}
\BIC{$\varphi_{147}$}
\AXC{$\varphi_{147} \to \varphi_{148}$\lbl{A1}}
\BIC{$\varphi_{148}$}
\UIC{$\forall a\,(\varphi_{148})$}
\end{ptree}
\noindent$\Pi_{67}$:
\begin{ptree}
\AXC{$\varphi_{20}$\lbl{A4}}
\AXC{$\forall a\,(\varphi_{148}) \to \varphi_{157}$\lbl{A3}}
\AXC{$\stackrel{\Pi_{66}}{\vdots}$}\noLine
\UIC{$\forall a\,(\varphi_{148})$}
\AXC{$\varphi_{276}$\lbl{A1}}
\BIC{$(\forall a\,(\varphi_{148}) \to \varphi_{157}) \to \varphi_{20} \to \mathrm{FLim}(f,c,\alpha) \to \varphi_{127}$}
\BIC{$\varphi_{20} \to \mathrm{FLim}(f,c,\alpha) \to \varphi_{127}$}
\BIC{$\mathrm{FLim}(f,c,\alpha) \to \varphi_{127}$}
\end{ptree}
\noindent$\Pi_{68}$:
\begin{ptree}
\AXC{$\varphi_{93}$\lbl{A2}}
\AXC{$\stackrel{\Pi_{8}}{\vdots}$}\noLine
\UIC{$\mathrm{FLim}(f,c,\alpha) \to \mathrm{FLim}(g,c,\beta) \to \varphi_{11}$}
\AXC{$\varphi_{243}$\lbl{A1}}
\BIC{$\varphi_{93} \to \mathrm{FLim}(f,c,\alpha) \to \varphi_{29}$}
\BIC{$\mathrm{FLim}(f,c,\alpha) \to \varphi_{29}$}
\end{ptree}
\noindent$\Pi_{69}$:
\begin{ptree}
\AXC{$\mathrm{FLim}(f,c,\alpha) \to \mathrm{FLim}(g,c,\beta) \to \varphi_{223}$\lbl{A1}}
\AXC{$\varphi_{271}$\lbl{A1}}
\BIC{$\mathrm{FLim}(f,c,\alpha) \to \mathrm{FLim}(g,c,\beta) \to \varphi_{233}$}
\AXC{$(\mathrm{FLim}(f,c,\alpha) \to \mathrm{FLim}(g,c,\beta) \to \varphi_{233}) \to \varphi_{265}$\lbl{A1}}
\BIC{$\varphi_{265}$}
\end{ptree}
\noindent$\Pi_{70}$:
\begin{ptree}
\AXC{$\forall x \in \mathbb{R}\,(\varphi_{30}) \to x \in \mathbb{R} \to \varphi_{30}$\lbl{A2}}
\AXC{$\stackrel{\Pi_{69}}{\vdots}$}\noLine
\UIC{$\varphi_{265}$}
\BIC{$\mathrm{FLim}(f,c,\alpha) \to \mathrm{FLim}(g,c,\beta) \to \varphi_{232}$}
\AXC{$(\mathrm{FLim}(f,c,\alpha) \to \mathrm{FLim}(g,c,\beta) \to \varphi_{232}) \to \varphi_{282}$\lbl{A1}}
\BIC{$\varphi_{282}$}
\end{ptree}
\noindent$\Pi_{71}$:
\begin{ptree}
\AXC{$\mathrm{FLim}(f,c,\alpha) \to \mathrm{FLim}(g,c,\beta) \to \varphi_{163}$\lbl{A1}}
\AXC{$\mathrm{FLim}(f,c,\alpha) \to \varphi_{146}$\lbl{A1}}
\AXC{$\stackrel{\Pi_{70}}{\vdots}$}\noLine
\UIC{$\varphi_{282}$}
\BIC{$(\mathrm{FLim}(f,c,\alpha) \to \mathrm{FLim}(g,c,\beta) \to \varphi_{163}) \to \varphi_{268}$}
\BIC{$\varphi_{268}$}
\end{ptree}
\noindent$\Pi_{72}$:
\begin{ptree}
\AXC{$\stackrel{\Pi_{68}}{\vdots}$}\noLine
\UIC{$\mathrm{FLim}(f,c,\alpha) \to \varphi_{29}$}
\AXC{$\mathrm{FLim}(f,c,\alpha) \to \mathrm{FLim}(g,c,\beta) \to \varphi_{189}$\lbl{A1}}
\AXC{$\stackrel{\Pi_{71}}{\vdots}$}\noLine
\UIC{$\varphi_{268}$}
\BIC{$(\mathrm{FLim}(f,c,\alpha) \to \varphi_{29}) \to \mathrm{FLim}(f,c,\alpha) \to \varphi_{205}$}
\BIC{$\mathrm{FLim}(f,c,\alpha) \to \varphi_{205}$}
\end{ptree}
\noindent$\Pi_{73}$:
\begin{ptree}
\AXC{$\stackrel{\Pi_{72}}{\vdots}$}\noLine
\UIC{$\mathrm{FLim}(f,c,\alpha) \to \varphi_{205}$}
\AXC{$(\mathrm{FLim}(f,c,\alpha) \to \varphi_{205}) \to \varphi_{128} \to \varphi_{40}$\lbl{A1}}
\BIC{$\varphi_{128} \to \varphi_{40}$}
\UIC{$\forall x\,(\varphi_{128} \to \varphi_{40})$}
\AXC{$\forall x\,(\varphi_{128} \to \varphi_{40}) \to \varphi_{285}$\lbl{A1}}
\BIC{$\varphi_{285}$}
\end{ptree}
\noindent$\Pi_{74}$:
\begin{ptree}
\AXC{$\varphi_{128} \to \forall x\,(\varphi_{128})$\lbl{A4}}
\AXC{$\forall x\,(\varphi_{128} \to \varphi_{40}) \to \forall x\,(\varphi_{128}) \to \varphi_{48}$\lbl{A3}}
\AXC{$\stackrel{\Pi_{73}}{\vdots}$}\noLine
\UIC{$\varphi_{285}$}
\BIC{$(\varphi_{128} \to \forall x\,(\varphi_{128})) \to \mathrm{FLim}(f,c,\alpha) \to \varphi_{209}$}
\BIC{$\mathrm{FLim}(f,c,\alpha) \to \varphi_{209}$}
\end{ptree}
\noindent$\Pi_{75}$:
\begin{ptree}
\AXC{$\stackrel{\Pi_{74}}{\vdots}$}\noLine
\UIC{$\mathrm{FLim}(f,c,\alpha) \to \varphi_{209}$}
\AXC{$\mathrm{FLim}(f,c,\alpha) \to \mathrm{FLim}(g,c,\beta) \to \varphi_{223}$\lbl{A1}}
\AXC{$\varphi_{277}$\lbl{A1}}
\BIC{$(\mathrm{FLim}(f,c,\alpha) \to \varphi_{209}) \to \mathrm{FLim}(f,c,\alpha) \to \varphi_{226}$}
\BIC{$\mathrm{FLim}(f,c,\alpha) \to \varphi_{226}$}
\end{ptree}
\noindent$\Pi_{76}$:
\begin{ptree}
\AXC{$\varphi_{104}$\lbl{DefPNbhd}}
\AXC{$\varphi_{104} \to \forall c\,(\mathrm{PNbhd}((f + g),c) \leftrightarrow \varphi_{84})$\lbl{A2}}
\BIC{$\forall c\,(\mathrm{PNbhd}((f + g),c) \leftrightarrow \varphi_{84})$}
\AXC{$\varphi_{236}$\lbl{A2}}
\BIC{$\mathrm{PNbhd}((f + g),c) \leftrightarrow \varphi_{84}$}
\AXC{$(\mathrm{PNbhd}((f + g),c) \leftrightarrow \varphi_{84}) \to \varphi_{259}$\lbl{A1}}
\BIC{$\varphi_{259}$}
\end{ptree}
\noindent$\Pi_{77}$:
\begin{ptree}
\AXC{$\varphi_{87} \to \neg (d \in \mathbb{R} \land 0 < d \land \varphi_{48})$\lbl{A2}}
\AXC{$\stackrel{\Pi_{75}}{\vdots}$}\noLine
\UIC{$\mathrm{FLim}(f,c,\alpha) \to \varphi_{226}$}
\AXC{$(\mathrm{FLim}(f,c,\alpha) \to \varphi_{226}) \to \varphi_{266}$\lbl{A1}}
\BIC{$\varphi_{266}$}
\BIC{$\mathrm{FLim}(f,c,\alpha) \to \varphi_{230}$}
\AXC{$\stackrel{\Pi_{76}}{\vdots}$}\noLine
\UIC{$\varphi_{259}$}
\BIC{$\mathrm{FLim}(f,c,\alpha) \to \mathrm{FLim}(g,c,\beta) \to \varphi_{117}$}
\end{ptree}
\noindent$\Pi_{78}$:
\begin{ptree}
\AXC{$\forall f\,\forall c\,\forall \alpha\,(\mathrm{FLim}(f,c,\alpha) \to \mathrm{PNbhd}(f,c))$\lbl{[D5]}}
\AXC{$\varphi_{34}$\lbl{A2}}
\BIC{$\forall c\,\forall \alpha\,(\mathrm{FLim}(g,c,\alpha) \to \mathrm{PNbhd}(g,c))$}
\AXC{$\varphi_{24}$\lbl{A2}}
\BIC{$\forall \alpha\,(\mathrm{FLim}(g,c,\alpha) \to \mathrm{PNbhd}(g,c))$}
\AXC{$\varphi_{15}$\lbl{A2}}
\BIC{$\mathrm{FLim}(g,c,\beta) \to \mathrm{PNbhd}(g,c)$}
\end{ptree}
\noindent$\Pi_{79}$:
\begin{ptree}
\AXC{$\varphi_{3}$\lbl{A1}}
\AXC{$\stackrel{\Pi_{78}}{\vdots}$}\noLine
\UIC{$\mathrm{FLim}(g,c,\beta) \to \mathrm{PNbhd}(g,c)$}
\AXC{$(\mathrm{FLim}(g,c,\beta) \to \mathrm{PNbhd}(g,c)) \to \varphi_{3} \to \varphi_{1}$\lbl{A1}}
\BIC{$\varphi_{3} \to \varphi_{1}$}
\BIC{$\varphi_{1}$}
\end{ptree}
\noindent$\Pi_{80}$:
\begin{ptree}
\AXC{$\forall f\,\forall c\,\forall \alpha\,(\mathrm{FLim}(f,c,\alpha) \to \mathrm{PNbhd}(f,c))$\lbl{[D5]}}
\AXC{$\varphi_{35}$\lbl{A2}}
\BIC{$\forall c\,\forall \alpha\,(\mathrm{FLim}(f,c,\alpha) \to \mathrm{PNbhd}(f,c))$}
\AXC{$\varphi_{25}$\lbl{A2}}
\BIC{$\forall \alpha\,(\mathrm{FLim}(f,c,\alpha) \to \mathrm{PNbhd}(f,c))$}
\AXC{$\varphi_{16}$\lbl{A2}}
\BIC{$\mathrm{FLim}(f,c,\alpha) \to \mathrm{PNbhd}(f,c)$}
\AXC{$\varphi_{82}$\lbl{A1}}
\BIC{$(\mathrm{FLim}(f,c,\alpha) \to \mathrm{FLim}(f,c,\alpha)) \to \varphi_{2}$}
\end{ptree}
\noindent$\Pi_{81}$:
\begin{ptree}
\AXC{$\mathrm{FLim}(f,c,\alpha) \to \mathrm{FLim}(f,c,\alpha)$\lbl{A1}}
\AXC{$\stackrel{\Pi_{80}}{\vdots}$}\noLine
\UIC{$(\mathrm{FLim}(f,c,\alpha) \to \mathrm{FLim}(f,c,\alpha)) \to \varphi_{2}$}
\BIC{$\varphi_{2}$}
\end{ptree}
\noindent$\Pi_{82}$:
\begin{ptree}
\AXC{$\forall f\,\forall g\,\forall c\,(\mathrm{PNbhd}(f,c) \to \varphi_{111})$\lbl{[H2]}}
\AXC{$\varphi_{245}$\lbl{A2}}
\BIC{$\forall g\,\forall c\,(\mathrm{PNbhd}(f,c) \to \varphi_{111})$}
\AXC{$\varphi_{244}$\lbl{A2}}
\BIC{$\forall c\,(\mathrm{PNbhd}(f,c) \to \varphi_{111})$}
\AXC{$\forall c\,(\mathrm{PNbhd}(f,c) \to \varphi_{111}) \to \mathrm{PNbhd}(f,c) \to \varphi_{111}$\lbl{A2}}
\BIC{$\mathrm{PNbhd}(f,c) \to \varphi_{111}$}
\end{ptree}
\noindent$\Pi_{83}$:
\begin{ptree}
\AXC{$\stackrel{\Pi_{79}}{\vdots}$}\noLine
\UIC{$\varphi_{1}$}
\AXC{$\stackrel{\Pi_{81}}{\vdots}$}\noLine
\UIC{$\varphi_{2}$}
\AXC{$\stackrel{\Pi_{82}}{\vdots}$}\noLine
\UIC{$\mathrm{PNbhd}(f,c) \to \varphi_{111}$}
\AXC{$\varphi_{255}$\lbl{A1}}
\BIC{$\varphi_{2} \to \varphi_{1} \to \mathrm{FLim}(f,c,\alpha) \to \varphi_{115}$}
\BIC{$\varphi_{1} \to \mathrm{FLim}(f,c,\alpha) \to \varphi_{115}$}
\BIC{$\mathrm{FLim}(f,c,\alpha) \to \varphi_{115}$}
\AXC{$(\mathrm{FLim}(f,c,\alpha) \to \varphi_{115}) \to \varphi_{274}$\lbl{A1}}
\BIC{$\varphi_{274}$}
\end{ptree}
\noindent$\Pi_{84}$:
\begin{ptree}
\AXC{$\varphi_{10} \to \forall d\,(\varphi_{10})$\lbl{A4}}
\AXC{$\varphi_{253}$\lbl{A3}}
\AXC{$\stackrel{\Pi_{77}}{\vdots}$}\noLine
\UIC{$\mathrm{FLim}(f,c,\alpha) \to \mathrm{FLim}(g,c,\beta) \to \varphi_{117}$}
\AXC{$\varphi_{250}$\lbl{A1}}
\BIC{$\varphi_{10} \to \neg (d \in \mathbb{R} \land 0 < d \land \forall x \in \mathbb{R}\,(\varphi_{30}))$}
\UIC{$\forall d\,(\varphi_{10} \to \neg (d \in \mathbb{R} \land 0 < d \land \forall x \in \mathbb{R}\,(\varphi_{30})))$}
\AXC{$\stackrel{\Pi_{83}}{\vdots}$}\noLine
\UIC{$\varphi_{274}$}
\BIC{$\varphi_{253} \to (\varphi_{10} \to \forall d\,(\varphi_{10})) \to \varphi_{7}$}
\BIC{$(\varphi_{10} \to \forall d\,(\varphi_{10})) \to \varphi_{7}$}
\BIC{$\varphi_{7}$}
\end{ptree}
\noindent$\Pi_{85}$:
\begin{ptree}
\AXC{$\varphi_{3}$\lbl{A1}}
\AXC{$\forall f\,\forall c\,\forall \alpha\,(\mathrm{FLim}(f,c,\alpha) \to \alpha \in \mathbb{R})$\lbl{[D4]}}
\AXC{$\varphi_{41}$\lbl{A2}}
\BIC{$\forall c\,\forall \alpha\,(\mathrm{FLim}(g,c,\alpha) \to \alpha \in \mathbb{R})$}
\AXC{$\varphi_{32}$\lbl{A2}}
\BIC{$\forall \alpha\,(\mathrm{FLim}(g,c,\alpha) \to \alpha \in \mathbb{R})$}
\AXC{$\varphi_{18}$\lbl{A2}}
\BIC{$\mathrm{FLim}(g,c,\beta) \to \beta \in \mathbb{R}$}
\AXC{$\varphi_{107}$\lbl{A1}}
\BIC{$\varphi_{3} \to \mathrm{FLim}(f,c,\alpha) \to \mathrm{FLim}(g,c,\beta) \to \beta \in \mathbb{R}$}
\BIC{$\mathrm{FLim}(f,c,\alpha) \to \mathrm{FLim}(g,c,\beta) \to \beta \in \mathbb{R}$}
\end{ptree}
\noindent$\Pi_{86}$:
\begin{ptree}
\AXC{$\forall f\,\forall c\,\forall \alpha\,(\mathrm{FLim}(f,c,\alpha) \to \alpha \in \mathbb{R})$\lbl{[D4]}}
\AXC{$\varphi_{42}$\lbl{A2}}
\BIC{$\forall c\,\forall \alpha\,(\mathrm{FLim}(f,c,\alpha) \to \alpha \in \mathbb{R})$}
\AXC{$\varphi_{33}$\lbl{A2}}
\BIC{$\forall \alpha\,(\mathrm{FLim}(f,c,\alpha) \to \alpha \in \mathbb{R})$}
\AXC{$\varphi_{19}$\lbl{A2}}
\BIC{$\mathrm{FLim}(f,c,\alpha) \to \alpha \in \mathbb{R}$}
\end{ptree}
\noindent$\Pi_{87}$:
\begin{ptree}
\AXC{$\mathrm{FLim}(f,c,\alpha) \to \mathrm{FLim}(f,c,\alpha)$\lbl{A1}}
\AXC{$\stackrel{\Pi_{86}}{\vdots}$}\noLine
\UIC{$\mathrm{FLim}(f,c,\alpha) \to \alpha \in \mathbb{R}$}
\AXC{$(\mathrm{FLim}(f,c,\alpha) \to \alpha \in \mathbb{R}) \to \varphi_{38}$\lbl{A1}}
\BIC{$\varphi_{38}$}
\BIC{$\mathrm{FLim}(f,c,\alpha) \to \mathrm{FLim}(g,c,\beta) \to \alpha \in \mathbb{R}$}
\end{ptree}
\noindent$\Pi_{88}$:
\begin{ptree}
\AXC{$\forall x\,\forall y\,(x \in \mathbb{R} \to y \in \mathbb{R} \to x + y \in \mathbb{R})$\lbl{(R1)}}
\AXC{$\varphi_{79}$\lbl{A2}}
\BIC{$\forall y\,(\alpha \in \mathbb{R} \to y \in \mathbb{R} \to \alpha + y \in \mathbb{R})$}
\AXC{$\varphi_{74}$\lbl{A2}}
\BIC{$\alpha \in \mathbb{R} \to \beta \in \mathbb{R} \to \alpha + \beta \in \mathbb{R}$}
\AXC{$(\alpha \in \mathbb{R} \to \beta \in \mathbb{R} \to \alpha + \beta \in \mathbb{R}) \to \varphi_{132}$\lbl{A1}}
\BIC{$\varphi_{132}$}
\end{ptree}
\noindent$\Pi_{89}$:
\begin{ptree}
\AXC{$\stackrel{\Pi_{85}}{\vdots}$}\noLine
\UIC{$\mathrm{FLim}(f,c,\alpha) \to \mathrm{FLim}(g,c,\beta) \to \beta \in \mathbb{R}$}
\AXC{$\stackrel{\Pi_{87}}{\vdots}$}\noLine
\UIC{$\mathrm{FLim}(f,c,\alpha) \to \mathrm{FLim}(g,c,\beta) \to \alpha \in \mathbb{R}$}
\AXC{$\stackrel{\Pi_{88}}{\vdots}$}\noLine
\UIC{$\varphi_{132}$}
\BIC{$\varphi_{69}$}
\BIC{$\mathrm{FLim}(f,c,\alpha) \to \mathrm{FLim}(g,c,\beta) \to \alpha + \beta \in \mathbb{R}$}
\end{ptree}
\noindent$\Pi_{90}$:
\begin{ptree}
\AXC{$\forall f\,\forall c\,\forall \alpha\,(\mathrm{FLim}(f,c,\alpha) \to c \in \mathbb{R})$\lbl{[D3]}}
\AXC{$\varphi_{36}$\lbl{A2}}
\BIC{$\forall c\,\forall \alpha\,(\mathrm{FLim}(f,c,\alpha) \to c \in \mathbb{R})$}
\AXC{$\varphi_{26}$\lbl{A2}}
\BIC{$\forall \alpha\,(\mathrm{FLim}(f,c,\alpha) \to c \in \mathbb{R})$}
\AXC{$\varphi_{17}$\lbl{A2}}
\BIC{$\mathrm{FLim}(f,c,\alpha) \to c \in \mathbb{R}$}
\end{ptree}
\noindent$\Pi_{91}$:
\begin{ptree}
\AXC{$\mathrm{FLim}(f,c,\alpha) \to \mathrm{FLim}(f,c,\alpha)$\lbl{A1}}
\AXC{$\stackrel{\Pi_{90}}{\vdots}$}\noLine
\UIC{$\mathrm{FLim}(f,c,\alpha) \to c \in \mathbb{R}$}
\AXC{$(\mathrm{FLim}(f,c,\alpha) \to c \in \mathbb{R}) \to \varphi_{37}$\lbl{A1}}
\BIC{$\varphi_{37}$}
\BIC{$\mathrm{FLim}(f,c,\alpha) \to \mathrm{FLim}(g,c,\beta) \to c \in \mathbb{R}$}
\end{ptree}
\noindent$\Pi_{92}$:
\begin{ptree}
\AXC{$\stackrel{\Pi_{57}}{\vdots}$}\noLine
\UIC{$\varphi_{4}$}
\AXC{$\varphi_{185}$\lbl{定理1.7(⇐)}}
\AXC{$\varphi_{185} \to \forall c\,(\varphi_{221})$\lbl{A2}}
\BIC{$\forall c\,(\varphi_{221})$}
\AXC{$\forall c\,(\varphi_{221}) \to \varphi_{221}$\lbl{A2}}
\BIC{$\varphi_{221}$}
\AXC{$\varphi_{221} \to \varphi_{228}$\lbl{A2}}
\BIC{$\varphi_{228}$}
\AXC{$\varphi_{228} \to \varphi_{4} \to \varphi_{254}$\lbl{A1}}
\BIC{$\varphi_{4} \to \varphi_{254}$}
\BIC{$\varphi_{254}$}
\end{ptree}
\noindent$\Pi_{93}$:
\begin{ptree}
\AXC{$\stackrel{\Pi_{89}}{\vdots}$}\noLine
\UIC{$\mathrm{FLim}(f,c,\alpha) \to \mathrm{FLim}(g,c,\beta) \to \alpha + \beta \in \mathbb{R}$}
\AXC{$\stackrel{\Pi_{91}}{\vdots}$}\noLine
\UIC{$\mathrm{FLim}(f,c,\alpha) \to \mathrm{FLim}(g,c,\beta) \to c \in \mathbb{R}$}
\AXC{$\stackrel{\Pi_{92}}{\vdots}$}\noLine
\UIC{$\varphi_{254}$}
\BIC{$\varphi_{248}$}
\BIC{$\varphi_{7} \to (\mathrm{FLim}(f,c,\alpha) \to \varphi_{127}) \to \varphi_{14}$}
\end{ptree}
\noindent 主木:
\begin{ptree}
\AXC{$\stackrel{\Pi_{67}}{\vdots}$}\noLine
\UIC{$\mathrm{FLim}(f,c,\alpha) \to \varphi_{127}$}
\AXC{$\stackrel{\Pi_{84}}{\vdots}$}\noLine
\UIC{$\varphi_{7}$}
\AXC{$\stackrel{\Pi_{93}}{\vdots}$}\noLine
\UIC{$\varphi_{7} \to (\mathrm{FLim}(f,c,\alpha) \to \varphi_{127}) \to \varphi_{14}$}
\BIC{$(\mathrm{FLim}(f,c,\alpha) \to \varphi_{127}) \to \varphi_{14}$}
\BIC{$\varphi_{14}$}
\UIC{$\forall \beta\,(\varphi_{14})$}
\UIC{$\forall \alpha\,\forall \beta\,(\varphi_{14})$}
\UIC{$\forall c\,\forall \alpha\,\forall \beta\,(\varphi_{14})$}
\UIC{$\forall g\,\forall c\,\forall \alpha\,\forall \beta\,(\varphi_{14})$}
\UIC{$\forall f\,\forall g\,\forall c\,\forall \alpha\,\forall \beta\,(\varphi_{14})$}
\end{ptree}
\begin{abbrevlist}
\item $\varphi_{1} :\equiv \mathrm{FLim}(f,c,\alpha) \to \mathrm{FLim}(g,c,\beta) \to \mathrm{PNbhd}(g,c)$
\item $\varphi_{2} :\equiv \mathrm{FLim}(f,c,\alpha) \to \mathrm{FLim}(g,c,\beta) \to \mathrm{PNbhd}(f,c)$
\item $\varphi_{3} :\equiv \mathrm{FLim}(f,c,\alpha) \to \mathrm{FLim}(g,c,\beta) \to \mathrm{FLim}(g,c,\beta)$
\item $\varphi_{4} :\equiv \mathrm{FLim}(f,c,\alpha) \to \mathrm{FLim}(g,c,\beta) \to \mathrm{RFn}((f + g))$
\item $\varphi_{5} :\equiv \mathrm{FLim}(f,c,\alpha) \to \mathrm{FLim}(g,c,\beta) \to \mathrm{Seq}(a) \to \mathrm{Seq}(a)$
\item $\varphi_{6} :\equiv \mathrm{FLim}(g,c,\beta) \to \mathrm{Seq}(a) \to \mathrm{Conv}(a,c) \to \mathrm{Conv}(a,c)$
\item $\varphi_{7} :\equiv \mathrm{FLim}(f,c,\alpha) \to \mathrm{FLim}(g,c,\beta) \to \mathrm{PNbhd}((f + g),c)$
\item $\varphi_{8} :\equiv \forall \alpha\,(\mathrm{FLim}(g,c,\alpha) \to \mathrm{RFn}(g)) \to \mathrm{FLim}(g,c,\beta) \to \mathrm{RFn}(g)$
\item $\varphi_{9} :\equiv \forall \alpha\,(\mathrm{FLim}(f,c,\alpha) \to \mathrm{RFn}(f)) \to \mathrm{FLim}(f,c,\alpha) \to \mathrm{RFn}(f)$
\item $\varphi_{10} :\equiv \mathrm{FLim}(f,c,\alpha) \land \mathrm{FLim}(g,c,\beta) \land \neg \mathrm{PNbhd}((f + g),c)$
\item $\varphi_{11} :\equiv \forall x\,(x \in \mathrm{dom}((f + g)) \leftrightarrow x \in \mathrm{dom}(f) \land x \in \mathrm{dom}(g))$
\item $\varphi_{12} :\equiv \forall c\,\forall \alpha\,(\mathrm{FLim}(g,c,\alpha) \to \mathrm{RFn}(g)) \to \forall \alpha\,(\mathrm{FLim}(g,c,\alpha) \to \mathrm{RFn}(g))$
\item $\varphi_{13} :\equiv \forall c\,\forall \alpha\,(\mathrm{FLim}(f,c,\alpha) \to \mathrm{RFn}(f)) \to \forall \alpha\,(\mathrm{FLim}(f,c,\alpha) \to \mathrm{RFn}(f))$
\item $\varphi_{14} :\equiv \mathrm{FLim}(f,c,\alpha) \to \mathrm{FLim}(g,c,\beta) \to \mathrm{FLim}((f + g),c,\alpha + \beta)$
\item $\varphi_{15} :\equiv \forall \alpha\,(\mathrm{FLim}(g,c,\alpha) \to \mathrm{PNbhd}(g,c)) \to \mathrm{FLim}(g,c,\beta) \to \mathrm{PNbhd}(g,c)$
\item $\varphi_{16} :\equiv \forall \alpha\,(\mathrm{FLim}(f,c,\alpha) \to \mathrm{PNbhd}(f,c)) \to \mathrm{FLim}(f,c,\alpha) \to \mathrm{PNbhd}(f,c)$
\item $\varphi_{17} :\equiv \forall \alpha\,(\mathrm{FLim}(f,c,\alpha) \to c \in \mathbb{R}) \to \mathrm{FLim}(f,c,\alpha) \to c \in \mathbb{R}$
\item $\varphi_{18} :\equiv \forall \alpha\,(\mathrm{FLim}(g,c,\alpha) \to \alpha \in \mathbb{R}) \to \mathrm{FLim}(g,c,\beta) \to \beta \in \mathbb{R}$
\item $\varphi_{19} :\equiv \forall \alpha\,(\mathrm{FLim}(f,c,\alpha) \to \alpha \in \mathbb{R}) \to \mathrm{FLim}(f,c,\alpha) \to \alpha \in \mathbb{R}$
\item $\varphi_{20} :\equiv \mathrm{FLim}(f,c,\alpha) \land \mathrm{FLim}(g,c,\beta) \to \forall a\,(\mathrm{FLim}(f,c,\alpha) \land \mathrm{FLim}(g,c,\beta))$
\item $\varphi_{21} :\equiv \forall f\,\forall c\,\forall \alpha\,(\mathrm{FLim}(f,c,\alpha) \to \mathrm{RFn}(f)) \to \forall c\,\forall \alpha\,(\mathrm{FLim}(g,c,\alpha) \to \mathrm{RFn}(g))$
\item $\varphi_{22} :\equiv \forall f\,\forall c\,\forall \alpha\,(\mathrm{FLim}(f,c,\alpha) \to \mathrm{RFn}(f)) \to \forall c\,\forall \alpha\,(\mathrm{FLim}(f,c,\alpha) \to \mathrm{RFn}(f))$
\item $\varphi_{23} :\equiv \forall x \in \mathrm{dom}((f + g))\,((f + g)(x) = f(x) + g(x))$
\item $\varphi_{24} :\equiv \forall c\,\forall \alpha\,(\mathrm{FLim}(g,c,\alpha) \to \mathrm{PNbhd}(g,c)) \to \forall \alpha\,(\mathrm{FLim}(g,c,\alpha) \to \mathrm{PNbhd}(g,c))$
\item $\varphi_{25} :\equiv \forall c\,\forall \alpha\,(\mathrm{FLim}(f,c,\alpha) \to \mathrm{PNbhd}(f,c)) \to \forall \alpha\,(\mathrm{FLim}(f,c,\alpha) \to \mathrm{PNbhd}(f,c))$
\item $\varphi_{26} :\equiv \forall c\,\forall \alpha\,(\mathrm{FLim}(f,c,\alpha) \to c \in \mathbb{R}) \to \forall \alpha\,(\mathrm{FLim}(f,c,\alpha) \to c \in \mathbb{R})$
\item $\varphi_{27} :\equiv \forall a\,\forall \alpha\,\forall b\,\forall \beta\,(\mathrm{Conv}(a,\alpha) \to \mathrm{Conv}(b,\beta) \to \mathrm{Conv}(a \oplus b,\alpha + \beta))$
\item $\varphi_{28} :\equiv \forall y\,\forall z\,((f + g)(a_{n}) = y \to y = z \to (f + g)(a_{n}) = z)$
\item $\varphi_{29} :\equiv \mathrm{FLim}(g,c,\beta) \to (x \in \mathrm{dom}((f + g)) \leftrightarrow x \in \mathrm{dom}(f) \land x \in \mathrm{dom}(g))$
\item $\varphi_{30} :\equiv 0 < \lvert x - c \rvert \to \lvert x - c \rvert < d \to x \in \mathrm{dom}(f) \land x \in \mathrm{dom}(g)$
\item $\varphi_{31} :\equiv (\mathrm{FLim}(f,c,\alpha) \to \mathrm{FLim}(f,c,\alpha)) \to \mathrm{FLim}(f,c,\alpha) \to \mathrm{FLim}(g,c,\beta) \to \mathrm{RFn}(f)$
\item $\varphi_{32} :\equiv \forall c\,\forall \alpha\,(\mathrm{FLim}(g,c,\alpha) \to \alpha \in \mathbb{R}) \to \forall \alpha\,(\mathrm{FLim}(g,c,\alpha) \to \alpha \in \mathbb{R})$
\item $\varphi_{33} :\equiv \forall c\,\forall \alpha\,(\mathrm{FLim}(f,c,\alpha) \to \alpha \in \mathbb{R}) \to \forall \alpha\,(\mathrm{FLim}(f,c,\alpha) \to \alpha \in \mathbb{R})$
\item $\varphi_{34} :\equiv \forall f\,\forall c\,\forall \alpha\,(\mathrm{FLim}(f,c,\alpha) \to \mathrm{PNbhd}(f,c)) \to \forall c\,\forall \alpha\,(\mathrm{FLim}(g,c,\alpha) \to \mathrm{PNbhd}(g,c))$
\item $\varphi_{35} :\equiv \forall f\,\forall c\,\forall \alpha\,(\mathrm{FLim}(f,c,\alpha) \to \mathrm{PNbhd}(f,c)) \to \forall c\,\forall \alpha\,(\mathrm{FLim}(f,c,\alpha) \to \mathrm{PNbhd}(f,c))$
\item $\varphi_{36} :\equiv \forall f\,\forall c\,\forall \alpha\,(\mathrm{FLim}(f,c,\alpha) \to c \in \mathbb{R}) \to \forall c\,\forall \alpha\,(\mathrm{FLim}(f,c,\alpha) \to c \in \mathbb{R})$
\item $\varphi_{37} :\equiv (\mathrm{FLim}(f,c,\alpha) \to \mathrm{FLim}(f,c,\alpha)) \to \mathrm{FLim}(f,c,\alpha) \to \mathrm{FLim}(g,c,\beta) \to c \in \mathbb{R}$
\item $\varphi_{38} :\equiv (\mathrm{FLim}(f,c,\alpha) \to \mathrm{FLim}(f,c,\alpha)) \to \mathrm{FLim}(f,c,\alpha) \to \mathrm{FLim}(g,c,\beta) \to \alpha \in \mathbb{R}$
\item $\varphi_{39} :\equiv a_{n} \in \mathrm{dom}((f + g)) \leftrightarrow a_{n} \in \mathrm{dom}(f) \land a_{n} \in \mathrm{dom}(g)$
\item $\varphi_{40} :\equiv x \in \mathbb{R} \to 0 < \lvert x - c \rvert \to \lvert x - c \rvert < d \to x \in \mathrm{dom}((f + g))$
\item $\varphi_{41} :\equiv \forall f\,\forall c\,\forall \alpha\,(\mathrm{FLim}(f,c,\alpha) \to \alpha \in \mathbb{R}) \to \forall c\,\forall \alpha\,(\mathrm{FLim}(g,c,\alpha) \to \alpha \in \mathbb{R})$
\item $\varphi_{42} :\equiv \forall f\,\forall c\,\forall \alpha\,(\mathrm{FLim}(f,c,\alpha) \to \alpha \in \mathbb{R}) \to \forall c\,\forall \alpha\,(\mathrm{FLim}(f,c,\alpha) \to \alpha \in \mathbb{R})$
\item $\varphi_{43} :\equiv \forall n \in \mathbb{N}\,(a_{n} \in \mathrm{dom}(g) \land a_{n} \neq c) \to \mathrm{Conv}(g \circ a,\beta)$
\item $\varphi_{44} :\equiv \forall n \in \mathbb{N}\,(a_{n} \in \mathrm{dom}(f) \land a_{n} \neq c) \to \mathrm{Conv}(f \circ a,\alpha)$
\item $\varphi_{45} :\equiv \forall y\,\forall z\,((f + g) \circ a(n) = y \to y = z \to (f + g) \circ a(n) = z)$
\item $\varphi_{46} :\equiv \forall a\,\forall \alpha\,(\mathrm{Conv}(a,\alpha) \to \mathrm{Seq}(a)) \to \forall \alpha\,(\mathrm{Conv}(g \circ a,\alpha) \to \mathrm{Seq}(g \circ a))$
\item $\varphi_{47} :\equiv \forall a\,\forall \alpha\,(\mathrm{Conv}(a,\alpha) \to \mathrm{Seq}(a)) \to \forall \alpha\,(\mathrm{Conv}(f \circ a,\alpha) \to \mathrm{Seq}(f \circ a))$
\item $\varphi_{48} :\equiv \forall x \in \mathbb{R}\,(0 < \lvert x - c \rvert \to \lvert x - c \rvert < d \to x \in \mathrm{dom}((f + g)))$
\item $\varphi_{49} :\equiv \forall n \in \mathbb{N}\,(a_{n} \in \mathrm{dom}((f + g)) \land a_{n} \neq c) \to \mathrm{Seq}(g \circ a)$
\item $\varphi_{50} :\equiv \forall n \in \mathbb{N}\,(a_{n} \in \mathrm{dom}((f + g)) \land a_{n} \neq c) \to \mathrm{Seq}(f \circ a)$
\item $\varphi_{51} :\equiv \forall a\,\forall b\,\forall \alpha\,(\mathrm{Seq}(b) \to \forall n \in \mathbb{N}\,(a_{n} = b_{n}) \to \mathrm{Conv}(a,\alpha) \to \mathrm{Conv}(b,\alpha))$
\item $\varphi_{52} :\equiv n \in \mathbb{N} \to \forall n \in \mathbb{N}\,((f + g) \circ a(n) = (f + g)(a_{n}))$
\item $\varphi_{53} :\equiv \forall b\,\forall \beta\,(\mathrm{Conv}(f \circ a,\alpha) \to \mathrm{Conv}(b,\beta) \to \mathrm{Conv}(f \circ a \oplus b,\alpha + \beta))$
\item $\varphi_{54} :\equiv \forall n \in \mathbb{N}\,(a_{n} \in \mathrm{dom}((f + g)) \land a_{n} \neq c) \to n \in \mathbb{N} \to n \in \mathbb{N}$
\item $\varphi_{55} :\equiv \forall n \in \mathbb{N}\,(a_{n} \in \mathrm{dom}((f + g)) \land a_{n} \neq c) \to \mathrm{Conv}(g \circ a,\beta)$
\item $\varphi_{56} :\equiv \forall n \in \mathbb{N}\,(a_{n} \in \mathrm{dom}((f + g)) \land a_{n} \neq c) \to \mathrm{Conv}(f \circ a,\alpha)$
\item $\varphi_{57} :\equiv \varphi_{3} \to \mathrm{FLim}(f,c,\alpha) \to \mathrm{FLim}(g,c,\beta) \to \mathrm{RFn}(g)$
\item $\varphi_{58} :\equiv \forall x\,\forall y\,(x = y \to y = x) \to \forall y\,((f + g) \circ a(n) = y \to y = (f + g) \circ a(n))$
\item $\varphi_{59} :\equiv n \in \mathbb{N} \to \mathrm{Seq}(g \circ a) \land \forall n \in \mathbb{N}\,(g \circ a(n) = g(a_{n}))$
\item $\varphi_{60} :\equiv n \in \mathbb{N} \to \mathrm{Seq}(f \circ a) \land \forall n \in \mathbb{N}\,(f \circ a(n) = f(a_{n}))$
\item $\varphi_{61} :\equiv \forall n \in \mathbb{N}\,(a_{n} \in \mathrm{dom}((f + g)) \land a_{n} \neq c) \to \mathrm{Seq}((f + g) \circ a)$
\item $\varphi_{62} :\equiv \mathrm{Conv}(f \circ a \oplus g \circ a,\alpha + \beta) \to \mathrm{Conv}((f + g) \circ a,\alpha + \beta)$
\item $\varphi_{63} :\equiv \forall x\,(x = x) \to f(a_{n}) + g \circ a(n) = f(a_{n}) + g \circ a(n)$
\item $\varphi_{64} :\equiv a_{n} \in \mathrm{dom}((f + g)) \to (f + g)(a_{n}) = f(a_{n}) + g(a_{n})$
\item $\varphi_{65} :\equiv \forall \alpha\,(\mathrm{Conv}(g \circ a,\alpha) \to \mathrm{Seq}(g \circ a)) \to \mathrm{Conv}(g \circ a,\beta) \to \mathrm{Seq}(g \circ a)$
\item $\varphi_{66} :\equiv \forall \alpha\,(\mathrm{Conv}(f \circ a,\alpha) \to \mathrm{Seq}(f \circ a)) \to \mathrm{Conv}(f \circ a,\alpha) \to \mathrm{Seq}(f \circ a)$
\item $\varphi_{67} :\equiv \forall n \in \mathbb{N}\,(a_{n} \in \mathrm{dom}((f + g)) \land a_{n} \neq c) \to n \in \mathbb{N} \to a_{n} \in \mathrm{dom}(g)$
\item $\varphi_{68} :\equiv \forall n \in \mathbb{N}\,(a_{n} \in \mathrm{dom}((f + g)) \land a_{n} \neq c) \to n \in \mathbb{N} \to a_{n} \in \mathrm{dom}(f)$
\item $\varphi_{69} :\equiv (\mathrm{FLim}(f,c,\alpha) \to \mathrm{FLim}(g,c,\beta) \to \beta \in \mathbb{R}) \to \mathrm{FLim}(f,c,\alpha) \to \mathrm{FLim}(g,c,\beta) \to \alpha + \beta \in \mathbb{R}$
\item $\varphi_{70} :\equiv \forall n \in \mathbb{N}\,(a_{n} \in \mathrm{dom}((f + g)) \land a_{n} \neq c) \to \forall n \in \mathbb{N}\,(a_{n} \in \mathrm{dom}(g))$
\item $\varphi_{71} :\equiv \forall n \in \mathbb{N}\,(a_{n} \in \mathrm{dom}((f + g)) \land a_{n} \neq c) \to \forall n \in \mathbb{N}\,(a_{n} \in \mathrm{dom}(f))$
\item $\varphi_{72} :\equiv n \in \mathbb{N} \to f \circ a(n) + g \circ a(n) = f(a_{n}) + g \circ a(n)$
\item $\varphi_{73} :\equiv \forall x\,(x = x) \to f \circ a(n) + g \circ a(n) = f \circ a(n) + g \circ a(n)$
\item $\varphi_{74} :\equiv \forall y\,(\alpha \in \mathbb{R} \to y \in \mathbb{R} \to \alpha + y \in \mathbb{R}) \to \alpha \in \mathbb{R} \to \beta \in \mathbb{R} \to \alpha + \beta \in \mathbb{R}$
\item $\varphi_{75} :\equiv \mathrm{Seq}((f + g) \circ a) \land \forall n \in \mathbb{N}\,((f + g) \circ a(n) = (f + g)(a_{n}))$
\item $\varphi_{76} :\equiv \forall n \in \mathbb{N}\,(a_{n} \in \mathrm{dom}((f + g)) \land a_{n} \neq c) \to \mathrm{Conv}((f + g) \circ a,\alpha + \beta)$
\item $\varphi_{77} :\equiv \mathrm{Conv}(f \circ a,\alpha) \to \mathrm{Conv}(g \circ a,\beta) \to \mathrm{Conv}(f \circ a \oplus g \circ a,\alpha + \beta)$
\item $\varphi_{78} :\equiv \forall a\,\forall b\,(\mathrm{Seq}(a) \to \mathrm{Seq}(b) \to \mathrm{Seq}(a \oplus b) \land \forall n \in \mathbb{N}\,(a \oplus b(n) = a_{n} + b_{n}))$
\item $\varphi_{79} :\equiv \forall x\,\forall y\,(x \in \mathbb{R} \to y \in \mathbb{R} \to x + y \in \mathbb{R}) \to \forall y\,(\alpha \in \mathbb{R} \to y \in \mathbb{R} \to \alpha + y \in \mathbb{R})$
\item $\varphi_{80} :\equiv \forall y\,\forall z\,(f \circ a \oplus g \circ a(n) = y \to y = z \to f \circ a \oplus g \circ a(n) = z)$
\item $\varphi_{81} :\equiv \forall n \in \mathbb{N}\,(f \circ a \oplus g \circ a(n) = f \circ a(n) + g \circ a(n))$
\item $\varphi_{82} :\equiv (\mathrm{FLim}(f,c,\alpha) \to \mathrm{PNbhd}(f,c)) \to (\mathrm{FLim}(f,c,\alpha) \to \mathrm{FLim}(f,c,\alpha)) \to \varphi_{2}$
\item $\varphi_{83} :\equiv \forall n \in \mathbb{N}\,(a_{n} \in \mathrm{dom}((f + g)) \land a_{n} \neq c) \to n \in \mathbb{N} \to a_{n} \in \mathrm{dom}((f + g))$
\item $\varphi_{84} :\equiv \exists r\,(r \in \mathbb{R} \land 0 < r \land \forall x \in \mathbb{R}\,(0 < \lvert x - c \rvert \to \lvert x - c \rvert < r \to x \in \mathrm{dom}((f + g))))$
\item $\varphi_{85} :\equiv \forall n \in \mathbb{N}\,(g \circ a(n) = g(a_{n})) \to n \in \mathbb{N} \to g \circ a(n) = g(a_{n})$
\item $\varphi_{86} :\equiv \forall n \in \mathbb{N}\,(f \circ a(n) = f(a_{n})) \to n \in \mathbb{N} \to f \circ a(n) = f(a_{n})$
\item $\varphi_{87} :\equiv \forall r\,(\neg (r \in \mathbb{R} \land 0 < r \land \forall x \in \mathbb{R}\,(0 < \lvert x - c \rvert \to \lvert x - c \rvert < r \to x \in \mathrm{dom}((f + g)))))$
\item $\varphi_{88} :\equiv \forall n \in \mathbb{N}\,(a_{n} \in \mathrm{dom}((f + g)) \land a_{n} \neq c) \to \forall n \in \mathbb{N}\,(a_{n} \in \mathrm{dom}((f + g)))$
\item $\varphi_{89} :\equiv \forall n \in \mathbb{N}\,(a_{n} \in \mathrm{dom}((f + g)) \land a_{n} \neq c) \to n \in \mathbb{N} \to a_{n} \in \mathrm{dom}(g) \land a_{n} \neq c$
\item $\varphi_{90} :\equiv \forall n \in \mathbb{N}\,(a_{n} \in \mathrm{dom}((f + g)) \land a_{n} \neq c) \to n \in \mathbb{N} \to a_{n} \in \mathrm{dom}(f) \land a_{n} \neq c$
\item $\varphi_{91} :\equiv \forall \alpha\,\forall a\,(\mathrm{FLim}(g,c,\alpha) \to \mathrm{Conv}(a,c) \to \forall n \in \mathbb{N}\,(a_{n} \in \mathrm{dom}(g) \land a_{n} \neq c) \to \mathrm{Conv}(g \circ a,\alpha))$
\item $\varphi_{92} :\equiv \forall x\,\forall y\,(x = y \to y = x) \to \forall y\,(f \circ a \oplus g \circ a(n) = y \to y = f \circ a \oplus g \circ a(n))$
\item $\varphi_{93} :\equiv \varphi_{11} \to (x \in \mathrm{dom}((f + g)) \leftrightarrow x \in \mathrm{dom}(f) \land x \in \mathrm{dom}(g))$
\item $\varphi_{94} :\equiv \forall n \in \mathbb{N}\,(a_{n} \in \mathrm{dom}((f + g)) \land a_{n} \neq c) \to \mathrm{Conv}(f \circ a \oplus g \circ a,\alpha + \beta)$
\item $\varphi_{95} :\equiv \forall n \in \mathbb{N}\,(a_{n} \in \mathrm{dom}((f + g)) \land a_{n} \neq c) \to n \in \mathbb{N} \to g \circ a(n) = g(a_{n})$
\item $\varphi_{96} :\equiv \forall n \in \mathbb{N}\,(a_{n} \in \mathrm{dom}((f + g)) \land a_{n} \neq c) \to n \in \mathbb{N} \to f \circ a(n) = f(a_{n})$
\item $\varphi_{97} :\equiv n \in \mathbb{N} \to n \in \mathbb{N} \to f \circ a \oplus g \circ a(n) = f \circ a(n) + g \circ a(n)$
\item $\varphi_{98} :\equiv \mathrm{FLim}(f,c,\alpha) \land \mathrm{FLim}(g,c,\beta) \land \mathrm{Seq}(a) \land \mathrm{Conv}(a,c) \land \forall n \in \mathbb{N}\,(a_{n} \in \mathrm{dom}((f + g)) \land a_{n} \neq c)$
\item $\varphi_{99} :\equiv \forall n \in \mathbb{N}\,(a_{n} \in \mathrm{dom}((f + g)) \land a_{n} \neq c) \to \forall n \in \mathbb{N}\,(a_{n} \in \mathrm{dom}(g) \land a_{n} \neq c)$
\item $\varphi_{100} :\equiv \forall n \in \mathbb{N}\,(a_{n} \in \mathrm{dom}((f + g)) \land a_{n} \neq c) \to \forall n \in \mathbb{N}\,(a_{n} \in \mathrm{dom}(f) \land a_{n} \neq c)$
\item $\varphi_{101} :\equiv \forall n \in \mathbb{N}\,(a_{n} \in \mathrm{dom}(g)) \to \mathrm{Seq}(g \circ a) \land \forall n \in \mathbb{N}\,(g \circ a(n) = g(a_{n}))$
\item $\varphi_{102} :\equiv \forall n \in \mathbb{N}\,(a_{n} \in \mathrm{dom}(f)) \to \mathrm{Seq}(f \circ a) \land \forall n \in \mathbb{N}\,(f \circ a(n) = f(a_{n}))$
\item $\varphi_{103} :\equiv \forall c\,\forall \alpha\,\forall a\,(\mathrm{FLim}(f,c,\alpha) \to \mathrm{Conv}(a,c) \to \varphi_{44})$
\item $\varphi_{104} :\equiv \forall f\,\forall c\,(\mathrm{PNbhd}(f,c) \leftrightarrow \exists r\,(r \in \mathbb{R} \land 0 < r \land \forall x \in \mathbb{R}\,(0 < \lvert x - c \rvert \to \lvert x - c \rvert < r \to x \in \mathrm{dom}(f))))$
\item $\varphi_{105} :\equiv \forall n \in \mathbb{N}\,(a_{n} \in \mathrm{dom}((f + g)) \land a_{n} \neq c) \to n \in \mathbb{N} \to a_{n} \in \mathrm{dom}((f + g)) \land a_{n} \neq c$
\item $\varphi_{106} :\equiv \forall z\,((f + g) \circ a(n) = (f + g)(a_{n}) \to (f + g)(a_{n}) = z \to (f + g) \circ a(n) = z)$
\item $\varphi_{107} :\equiv (\mathrm{FLim}(g,c,\beta) \to \beta \in \mathbb{R}) \to \varphi_{3} \to \mathrm{FLim}(f,c,\alpha) \to \mathrm{FLim}(g,c,\beta) \to \beta \in \mathbb{R}$
\item $\varphi_{108} :\equiv \forall n \in \mathbb{N}\,(a_{n} \in \mathrm{dom}((f + g)) \land a_{n} \neq c) \to \forall n \in \mathbb{N}\,(a_{n} \in \mathrm{dom}((f + g)) \land a_{n} \neq c)$
\item $\varphi_{109} :\equiv \forall n \in \mathbb{N}\,(a_{n} \in \mathrm{dom}((f + g)) \land a_{n} \neq c) \to n \in \mathbb{N} \to n \in \mathbb{N} \to g \circ a(n) = g(a_{n})$
\item $\varphi_{110} :\equiv \forall n \in \mathbb{N}\,(a_{n} \in \mathrm{dom}((f + g)) \land a_{n} \neq c) \to n \in \mathbb{N} \to n \in \mathbb{N} \to f \circ a(n) = f(a_{n})$
\item $\varphi_{111} :\equiv \mathrm{PNbhd}(g,c) \to \exists d\,(d \in \mathbb{R} \land 0 < d \land \forall x \in \mathbb{R}\,(\varphi_{30}))$
\item $\varphi_{112} :\equiv f(a_{n}) + g \circ a(n) = f(a_{n}) + z \to f(a_{n}) + g \circ a(n) = f(a_{n}) + w$
\item $\varphi_{113} :\equiv \forall n \in \mathbb{N}\,(a_{n} \in \mathrm{dom}((f + g)) \land a_{n} \neq c) \to n \in \mathbb{N} \to \forall n \in \mathbb{N}\,(g \circ a(n) = g(a_{n}))$
\item $\varphi_{114} :\equiv \forall n \in \mathbb{N}\,(a_{n} \in \mathrm{dom}((f + g)) \land a_{n} \neq c) \to n \in \mathbb{N} \to \forall n \in \mathbb{N}\,(f \circ a(n) = f(a_{n}))$
\item $\varphi_{115} :\equiv \mathrm{FLim}(g,c,\beta) \to \exists d\,(d \in \mathbb{R} \land 0 < d \land \forall x \in \mathbb{R}\,(\varphi_{30}))$
\item $\varphi_{116} :\equiv \forall n \in \mathbb{N}\,(a_{n} \in \mathrm{dom}((f + g)) \land a_{n} \neq c) \to n \in \mathbb{N} \to (f + g) \circ a(n) = (f + g)(a_{n})$
\item $\varphi_{117} :\equiv d \in \mathbb{R} \land 0 < d \land \forall x \in \mathbb{R}\,(\varphi_{30}) \to \mathrm{PNbhd}((f + g),c)$
\item $\varphi_{118} :\equiv \mathrm{FLim}(f,c,\alpha) \to \mathrm{FLim}(g,c,\beta) \to \mathrm{Seq}(a) \to \mathrm{Conv}(a,c) \to \varphi_{49}$
\item $\varphi_{119} :\equiv \mathrm{FLim}(f,c,\alpha) \to \mathrm{FLim}(g,c,\beta) \to \mathrm{Seq}(a) \to \mathrm{Conv}(a,c) \to \varphi_{50}$
\item $\varphi_{120} :\equiv \forall z\,((f + g)(a_{n}) = f(a_{n}) + g(a_{n}) \to f(a_{n}) + g(a_{n}) = z \to (f + g)(a_{n}) = z)$
\item $\varphi_{121} :\equiv \forall n \in \mathbb{N}\,(a_{n} \in \mathrm{dom}((f + g)) \land a_{n} \neq c) \to n \in \mathbb{N} \to (f + g)(a_{n}) = f(a_{n}) + g(a_{n})$
\item $\varphi_{122} :\equiv \mathrm{FLim}(f,c,\alpha) \to \mathrm{FLim}(g,c,\beta) \to \mathrm{Seq}(a) \to \mathrm{Conv}(a,c) \to \varphi_{54}$
\item $\varphi_{123} :\equiv \mathrm{FLim}(f,c,\alpha) \to \mathrm{FLim}(g,c,\beta) \to \mathrm{Seq}(a) \to \mathrm{Conv}(a,c) \to \varphi_{55}$
\item $\varphi_{124} :\equiv \mathrm{FLim}(f,c,\alpha) \to \mathrm{FLim}(g,c,\beta) \to \mathrm{Seq}(a) \to \mathrm{Conv}(a,c) \to \varphi_{56}$
\item $\varphi_{125} :\equiv \forall n \in \mathbb{N}\,((f + g) \circ a(n) = (f + g)(a_{n})) \to n \in \mathbb{N} \to (f + g) \circ a(n) = (f + g)(a_{n})$
\item $\varphi_{126} :\equiv (f + g)(a_{n}) = f \circ a \oplus g \circ a(n) \to (f + g) \circ a(n) = f \circ a \oplus g \circ a(n)$
\item $\varphi_{127} :\equiv \mathrm{FLim}(g,c,\beta) \to \forall a\,(\mathrm{Seq}(a) \to \mathrm{Conv}(a,c) \to \varphi_{76})$
\item $\varphi_{128} :\equiv \mathrm{FLim}(f,c,\alpha) \land \mathrm{FLim}(g,c,\beta) \land (d \in \mathbb{R} \land 0 < d \land \forall x \in \mathbb{R}\,(\varphi_{30}))$
\item $\varphi_{129} :\equiv \forall n \in \mathbb{N}\,(a_{n} \in \mathrm{dom}((f + g)) \land a_{n} \neq c) \to n \in \mathbb{N} \to (f + g)(a_{n}) = f \circ a \oplus g \circ a(n)$
\item $\varphi_{130} :\equiv \forall n \in \mathbb{N}\,(a_{n} \in \mathrm{dom}((f + g)) \land a_{n} \neq c) \to n \in \mathbb{N} \to n \in \mathbb{N} \to (f + g) \circ a(n) = (f + g)(a_{n})$
\item $\varphi_{131} :\equiv \mathrm{FLim}(f,c,\alpha) \to \mathrm{FLim}(g,c,\beta) \to \mathrm{Seq}(a) \to \mathrm{Conv}(a,c) \to \varphi_{61}$
\item $\varphi_{132} :\equiv (\mathrm{FLim}(f,c,\alpha) \to \mathrm{FLim}(g,c,\beta) \to \alpha \in \mathbb{R}) \to \varphi_{69}$
\item $\varphi_{133} :\equiv \forall b\,(\mathrm{Seq}(f \circ a) \to \mathrm{Seq}(b) \to \mathrm{Seq}(f \circ a \oplus b) \land \forall n \in \mathbb{N}\,(f \circ a \oplus b(n) = f \circ a(n) + b_{n}))$
\item $\varphi_{134} :\equiv \forall b\,\forall \alpha\,(\mathrm{Seq}(b) \to \forall n \in \mathbb{N}\,(f \circ a \oplus g \circ a(n) = b_{n}) \to \mathrm{Conv}(f \circ a \oplus g \circ a,\alpha) \to \mathrm{Conv}(b,\alpha))$
\item $\varphi_{135} :\equiv (f + g) \circ a(n) = f \circ a \oplus g \circ a(n) \to f \circ a \oplus g \circ a(n) = (f + g) \circ a(n)$
\item $\varphi_{136} :\equiv f \circ a(n) + g \circ a(n) = z + g \circ a(n) \to f \circ a(n) + g \circ a(n) = w + g \circ a(n)$
\item $\varphi_{137} :\equiv \forall n \in \mathbb{N}\,(a_{n} \in \mathrm{dom}((f + g)) \land a_{n} \neq c) \to n \in \mathbb{N} \to (f + g) \circ a(n) = f \circ a \oplus g \circ a(n)$
\item $\varphi_{138} :\equiv \forall n \in \mathbb{N}\,(a_{n} \in \mathrm{dom}((f + g)) \land a_{n} \neq c) \to n \in \mathbb{N} \to f \circ a \oplus g \circ a(n) = (f + g) \circ a(n)$
\item $\varphi_{139} :\equiv f(a_{n}) + g(a_{n}) = f \circ a \oplus g \circ a(n) \to (f + g)(a_{n}) = f \circ a \oplus g \circ a(n)$
\item $\varphi_{140} :\equiv \mathrm{FLim}(f,c,\alpha) \to \mathrm{FLim}(g,c,\beta) \to \mathrm{Seq}(a) \to \mathrm{Conv}(a,c) \to \varphi_{67}$
\item $\varphi_{141} :\equiv \mathrm{FLim}(f,c,\alpha) \to \mathrm{FLim}(g,c,\beta) \to \mathrm{Seq}(a) \to \mathrm{Conv}(a,c) \to \varphi_{68}$
\item $\varphi_{142} :\equiv \forall n \in \mathbb{N}\,(a_{n} \in \mathrm{dom}((f + g)) \land a_{n} \neq c) \to \forall n \in \mathbb{N}\,(f \circ a \oplus g \circ a(n) = (f + g) \circ a(n))$
\item $\varphi_{143} :\equiv \mathrm{FLim}(f,c,\alpha) \to \mathrm{FLim}(g,c,\beta) \to \mathrm{Seq}(a) \to \mathrm{Conv}(a,c) \to \varphi_{70}$
\item $\varphi_{144} :\equiv \mathrm{FLim}(f,c,\alpha) \to \mathrm{FLim}(g,c,\beta) \to \mathrm{Seq}(a) \to \mathrm{Conv}(a,c) \to \varphi_{71}$
\item $\varphi_{145} :\equiv \forall n \in \mathbb{N}\,(a_{n} \in \mathrm{dom}((f + g)) \land a_{n} \neq c) \to n \in \mathbb{N} \to f(a_{n}) + g \circ a(n) = f(a_{n}) + g(a_{n})$
\item $\varphi_{146} :\equiv \mathrm{FLim}(g,c,\beta) \to d \in \mathbb{R} \land 0 < d \land \forall x \in \mathbb{R}\,(\varphi_{30}) \to x \in \mathbb{R} \to x \in \mathbb{R}$
\item $\varphi_{147} :\equiv \mathrm{FLim}(f,c,\alpha) \to \mathrm{FLim}(g,c,\beta) \to \mathrm{Seq}(a) \to \mathrm{Conv}(a,c) \to \varphi_{76}$
\item $\varphi_{148} :\equiv \mathrm{FLim}(f,c,\alpha) \land \mathrm{FLim}(g,c,\beta) \to \mathrm{Seq}(a) \to \mathrm{Conv}(a,c) \to \varphi_{76}$
\item $\varphi_{149} :\equiv \forall n \in \mathbb{N}\,(a_{n} \in \mathrm{dom}((f + g)) \land a_{n} \neq c) \to n \in \mathbb{N} \to f \circ a \oplus g \circ a(n) = f(a_{n}) + g(a_{n})$
\item $\varphi_{150} :\equiv \forall n \in \mathbb{N}\,(a_{n} \in \mathrm{dom}((f + g)) \land a_{n} \neq c) \to n \in \mathbb{N} \to f(a_{n}) + g(a_{n}) = f \circ a \oplus g \circ a(n)$
\item $\varphi_{151} :\equiv \forall a\,(\mathrm{Seq}(a) \to \mathrm{Conv}(a,c) \to \varphi_{76}) \to \mathrm{FLim}((f + g),c,\alpha + \beta)$
\item $\varphi_{152} :\equiv \mathrm{FLim}(f,c,\alpha) \to \mathrm{FLim}(g,c,\beta) \to \mathrm{Seq}(a) \to \mathrm{Conv}(a,c) \to \varphi_{83}$
\item $\varphi_{153} :\equiv \forall n \in \mathbb{N}\,(a_{n} \in \mathrm{dom}((f + g)) \land a_{n} \neq c) \to n \in \mathbb{N} \to f \circ a \oplus g \circ a(n) = f(a_{n}) + g \circ a(n)$
\item $\varphi_{154} :\equiv \mathrm{Conv}(a,c) \to \forall n \in \mathbb{N}\,(a_{n} \in \mathrm{dom}((f + g)) \land a_{n} \neq c) \to \varphi_{52}$
\item $\varphi_{155} :\equiv f(a_{n}) + g \circ a(n) = f(a_{n}) + g(a_{n}) \to f \circ a \oplus g \circ a(n) = f(a_{n}) + g(a_{n})$
\item $\varphi_{156} :\equiv f \circ a \oplus g \circ a(n) = f(a_{n}) + g(a_{n}) \to f(a_{n}) + g(a_{n}) = f \circ a \oplus g \circ a(n)$
\item $\varphi_{157} :\equiv \forall a\,(\mathrm{FLim}(f,c,\alpha) \land \mathrm{FLim}(g,c,\beta)) \to \forall a\,(\mathrm{Seq}(a) \to \mathrm{Conv}(a,c) \to \varphi_{76})$
\item $\varphi_{158} :\equiv \mathrm{FLim}(f,c,\alpha) \to \mathrm{FLim}(g,c,\beta) \to \mathrm{Seq}(a) \to \mathrm{Conv}(a,c) \to \varphi_{88}$
\item $\varphi_{159} :\equiv f(a_{n}) + g \circ a(n) = f(a_{n}) + g \circ a(n) \to f(a_{n}) + g \circ a(n) = f(a_{n}) + g(a_{n})$
\item $\varphi_{160} :\equiv \mathrm{Conv}(a,c) \to \forall n \in \mathbb{N}\,(a_{n} \in \mathrm{dom}((f + g)) \land a_{n} \neq c) \to \varphi_{59}$
\item $\varphi_{161} :\equiv \mathrm{Conv}(a,c) \to \forall n \in \mathbb{N}\,(a_{n} \in \mathrm{dom}((f + g)) \land a_{n} \neq c) \to \varphi_{60}$
\item $\varphi_{162} :\equiv \forall n \in \mathbb{N}\,(a_{n} \in \mathrm{dom}((f + g)) \land a_{n} \neq c) \to n \in \mathbb{N} \to f \circ a \oplus g \circ a(n) = f \circ a(n) + g \circ a(n)$
\item $\varphi_{163} :\equiv d \in \mathbb{R} \land 0 < d \land \forall x \in \mathbb{R}\,(\varphi_{30}) \to x \in \mathbb{R} \to 0 < \lvert x - c \rvert \to 0 < \lvert x - c \rvert$
\item $\varphi_{164} :\equiv \mathrm{FLim}(f,c,\alpha) \to \mathrm{FLim}(g,c,\beta) \to \mathrm{Seq}(a) \to \mathrm{Conv}(a,c) \to \varphi_{89}$
\item $\varphi_{165} :\equiv \mathrm{FLim}(f,c,\alpha) \to \mathrm{FLim}(g,c,\beta) \to \mathrm{Seq}(a) \to \mathrm{Conv}(a,c) \to \varphi_{90}$
\item $\varphi_{166} :\equiv \forall g\,(\mathrm{RFn}(f) \to \mathrm{RFn}(g) \to \mathrm{RFn}((f + g)) \land \varphi_{11} \land \varphi_{23})$
\item $\varphi_{167} :\equiv \mathrm{FLim}(f,c,\alpha) \to \mathrm{FLim}(g,c,\beta) \to \mathrm{Seq}(a) \to \mathrm{Conv}(a,c) \to \varphi_{94}$
\item $\varphi_{168} :\equiv \mathrm{FLim}(f,c,\alpha) \to \mathrm{FLim}(g,c,\beta) \to \mathrm{Seq}(a) \to \mathrm{Conv}(a,c) \to \varphi_{95}$
\item $\varphi_{169} :\equiv \mathrm{FLim}(f,c,\alpha) \to \mathrm{FLim}(g,c,\beta) \to \mathrm{Seq}(a) \to \mathrm{Conv}(a,c) \to \varphi_{96}$
\item $\varphi_{170} :\equiv \mathrm{FLim}(f,c,\alpha) \to \mathrm{FLim}(g,c,\beta) \to \mathrm{Seq}(a) \to \mathrm{Conv}(a,c) \to \varphi_{99}$
\item $\varphi_{171} :\equiv \mathrm{FLim}(f,c,\alpha) \to \mathrm{FLim}(g,c,\beta) \to \mathrm{Seq}(a) \to \mathrm{Conv}(a,c) \to \varphi_{100}$
\item $\varphi_{172} :\equiv f \circ a(n) + g \circ a(n) = f(a_{n}) + g \circ a(n) \to f \circ a \oplus g \circ a(n) = f(a_{n}) + g \circ a(n)$
\item $\varphi_{173} :\equiv \mathrm{FLim}(f,c,\alpha) \to \mathrm{FLim}(g,c,\beta) \to \mathrm{Seq}(a) \to \mathrm{Conv}(a,c) \to \varphi_{105}$
\item $\varphi_{174} :\equiv \mathrm{Conv}(a,c) \to \forall n \in \mathbb{N}\,(a_{n} \in \mathrm{dom}((f + g)) \land a_{n} \neq c) \to \varphi_{72}$
\item $\varphi_{175} :\equiv \mathrm{RFn}((f + g)) \to \mathrm{Seq}(a) \to \forall n \in \mathbb{N}\,(a_{n} \in \mathrm{dom}((f + g))) \to \varphi_{75}$
\item $\varphi_{176} :\equiv \mathrm{Conv}(a,c) \to \forall n \in \mathbb{N}\,(a_{n} \in \mathrm{dom}((f + g)) \land a_{n} \neq c) \to \varphi_{75}$
\item $\varphi_{177} :\equiv \mathrm{FLim}(f,c,\alpha) \to \mathrm{FLim}(g,c,\beta) \to \mathrm{RFn}((f + g)) \land \varphi_{11} \land \varphi_{23}$
\item $\varphi_{178} :\equiv \mathrm{FLim}(f,c,\alpha) \to \mathrm{FLim}(g,c,\beta) \to \mathrm{Seq}(a) \to \mathrm{Conv}(a,c) \to \varphi_{108}$
\item $\varphi_{179} :\equiv \forall w\,(g \circ a(n) = w \to f(a_{n}) + g \circ a(n) = f(a_{n}) + g \circ a(n) \to f(a_{n}) + g \circ a(n) = f(a_{n}) + w)$
\item $\varphi_{180} :\equiv \forall z\,(f \circ a \oplus g \circ a(n) = f(a_{n}) + g \circ a(n) \to f(a_{n}) + g \circ a(n) = z \to f \circ a \oplus g \circ a(n) = z)$
\item $\varphi_{181} :\equiv \forall n \in \mathbb{N}\,(a_{n} \in \mathrm{dom}((f + g)) \land a_{n} \neq c) \to n \in \mathbb{N} \to \varphi_{81}$
\item $\varphi_{182} :\equiv \mathrm{FLim}(f,c,\alpha) \to \mathrm{FLim}(g,c,\beta) \to \mathrm{Seq}(a) \to \mathrm{Conv}(a,c) \to \varphi_{109}$
\item $\varphi_{183} :\equiv \mathrm{FLim}(f,c,\alpha) \to \mathrm{FLim}(g,c,\beta) \to \mathrm{Seq}(a) \to \mathrm{Conv}(a,c) \to \varphi_{110}$
\item $\varphi_{184} :\equiv f \circ a(n) + g \circ a(n) = f \circ a(n) + g \circ a(n) \to f \circ a(n) + g \circ a(n) = f(a_{n}) + g \circ a(n)$
\item $\varphi_{185} :\equiv \forall f\,\forall c\,\forall \alpha\,(\mathrm{RFn}(f) \to c \in \mathbb{R} \to \alpha \in \mathbb{R} \to \mathrm{PNbhd}(f,c) \to \forall a\,(\mathrm{Seq}(a) \to \mathrm{Conv}(a,c) \to \varphi_{44}) \to \mathrm{FLim}(f,c,\alpha))$
\item $\varphi_{186} :\equiv \mathrm{FLim}(f,c,\alpha) \to \mathrm{FLim}(g,c,\beta) \to \mathrm{Seq}(a) \to \mathrm{Conv}(a,c) \to \varphi_{113}$
\item $\varphi_{187} :\equiv \mathrm{FLim}(f,c,\alpha) \to \mathrm{FLim}(g,c,\beta) \to \mathrm{Seq}(a) \to \mathrm{Conv}(a,c) \to \varphi_{114}$
\item $\varphi_{188} :\equiv \mathrm{FLim}(f,c,\alpha) \to \mathrm{FLim}(g,c,\beta) \to \mathrm{Seq}(a) \to \mathrm{Conv}(a,c) \to \varphi_{116}$
\item $\varphi_{189} :\equiv d \in \mathbb{R} \land 0 < d \land \forall x \in \mathbb{R}\,(\varphi_{30}) \to x \in \mathbb{R} \to 0 < \lvert x - c \rvert \to \lvert x - c \rvert < d \to \lvert x - c \rvert < d$
\item $\varphi_{190} :\equiv \forall z\,(f \circ a \oplus g \circ a(n) = f \circ a(n) + g \circ a(n) \to f \circ a(n) + g \circ a(n) = z \to f \circ a \oplus g \circ a(n) = z)$
\item $\varphi_{191} :\equiv \mathrm{Seq}(f \circ a) \to \mathrm{Seq}(g \circ a) \to \mathrm{Seq}(f \circ a \oplus g \circ a) \land \varphi_{81}$
\item $\varphi_{192} :\equiv \mathrm{Conv}(a,c) \to \forall n \in \mathbb{N}\,(a_{n} \in \mathrm{dom}((f + g)) \land a_{n} \neq c) \to \varphi_{97}$
\item $\varphi_{193} :\equiv \mathrm{FLim}(f,c,\alpha) \to \mathrm{FLim}(g,c,\beta) \to \mathrm{Seq}(a) \to \mathrm{Conv}(a,c) \to \varphi_{121}$
\item $\varphi_{194} :\equiv \forall w\,(f \circ a(n) = w \to f \circ a(n) + g \circ a(n) = f \circ a(n) + g \circ a(n) \to f \circ a(n) + g \circ a(n) = w + g \circ a(n))$
\item $\varphi_{195} :\equiv \mathrm{FLim}(f,c,\alpha) \to \mathrm{FLim}(g,c,\beta) \to \mathrm{Seq}(a) \to \mathrm{Conv}(a,c) \to \varphi_{129}$
\item $\varphi_{196} :\equiv \mathrm{FLim}(f,c,\alpha) \to \mathrm{FLim}(g,c,\beta) \to \mathrm{Seq}(a) \to \mathrm{Conv}(a,c) \to \varphi_{130}$
\item $\varphi_{197} :\equiv \mathrm{FLim}(f,c,\alpha) \to \mathrm{FLim}(g,c,\beta) \to \mathrm{Seq}(a) \to \varphi_{154}$
\item $\varphi_{198} :\equiv \mathrm{FLim}(f,c,\alpha) \to \mathrm{FLim}(g,c,\beta) \to \mathrm{Seq}(a) \to \mathrm{Conv}(a,c) \to \varphi_{137}$
\item $\varphi_{199} :\equiv \mathrm{FLim}(f,c,\alpha) \to \mathrm{FLim}(g,c,\beta) \to \mathrm{Seq}(a) \to \mathrm{Conv}(a,c) \to \varphi_{138}$
\item $\varphi_{200} :\equiv \forall \alpha\,(\mathrm{Seq}((f + g) \circ a) \to \forall n \in \mathbb{N}\,(f \circ a \oplus g \circ a(n) = (f + g) \circ a(n)) \to \mathrm{Conv}(f \circ a \oplus g \circ a,\alpha) \to \mathrm{Conv}((f + g) \circ a,\alpha))$
\item $\varphi_{201} :\equiv \mathrm{FLim}(f,c,\alpha) \to \mathrm{FLim}(g,c,\beta) \to \mathrm{Seq}(a) \to \varphi_{160}$
\item $\varphi_{202} :\equiv \mathrm{FLim}(f,c,\alpha) \to \mathrm{FLim}(g,c,\beta) \to \mathrm{Seq}(a) \to \varphi_{161}$
\item $\varphi_{203} :\equiv \forall n\,(\varphi_{98} \to n \in \mathbb{N} \to f \circ a \oplus g \circ a(n) = (f + g) \circ a(n))$
\item $\varphi_{204} :\equiv \mathrm{FLim}(f,c,\alpha) \to \mathrm{FLim}(g,c,\beta) \to \mathrm{Seq}(a) \to \mathrm{Conv}(a,c) \to \varphi_{142}$
\item $\varphi_{205} :\equiv \mathrm{FLim}(g,c,\beta) \to d \in \mathbb{R} \land 0 < d \land \forall x \in \mathbb{R}\,(\varphi_{30}) \to \varphi_{40}$
\item $\varphi_{206} :\equiv \varphi_{81} \to n \in \mathbb{N} \to f \circ a \oplus g \circ a(n) = f \circ a(n) + g \circ a(n)$
\item $\varphi_{207} :\equiv \mathrm{FLim}(f,c,\alpha) \to \mathrm{FLim}(g,c,\beta) \to \mathrm{Seq}(a) \to \mathrm{Conv}(a,c) \to \varphi_{145}$
\item $\varphi_{208} :\equiv \forall n\,(\varphi_{98}) \to \forall n \in \mathbb{N}\,(f \circ a \oplus g \circ a(n) = (f + g) \circ a(n))$
\item $\varphi_{209} :\equiv \mathrm{FLim}(g,c,\beta) \to d \in \mathbb{R} \land 0 < d \land \forall x \in \mathbb{R}\,(\varphi_{30}) \to \varphi_{48}$
\item $\varphi_{210} :\equiv \mathrm{FLim}(f,c,\alpha) \to \mathrm{FLim}(g,c,\beta) \to \mathrm{Seq}(a) \to \mathrm{Conv}(a,c) \to \varphi_{149}$
\item $\varphi_{211} :\equiv \mathrm{FLim}(f,c,\alpha) \to \mathrm{FLim}(g,c,\beta) \to \mathrm{Seq}(a) \to \mathrm{Conv}(a,c) \to \varphi_{150}$
\item $\varphi_{212} :\equiv \forall n \in \mathbb{N}\,(a_{n} \in \mathrm{dom}((f + g)) \land a_{n} \neq c) \to \mathrm{Seq}(f \circ a \oplus g \circ a) \land \varphi_{81}$
\item $\varphi_{213} :\equiv \mathrm{Seq}((f + g) \circ a) \to \forall n \in \mathbb{N}\,(f \circ a \oplus g \circ a(n) = (f + g) \circ a(n)) \to \varphi_{62}$
\item $\varphi_{214} :\equiv \mathrm{FLim}(f,c,\alpha) \to \mathrm{FLim}(g,c,\beta) \to \mathrm{Seq}(a) \to \mathrm{Conv}(a,c) \to \varphi_{153}$
\item $\varphi_{215} :\equiv \mathrm{FLim}(f,c,\alpha) \to \mathrm{FLim}(g,c,\beta) \to \mathrm{Seq}(a) \to \varphi_{174}$
\item $\varphi_{216} :\equiv \forall a\,(\mathrm{FLim}(g,c,\beta) \to \mathrm{Conv}(a,c) \to \varphi_{43}) \to \mathrm{FLim}(g,c,\beta) \to \mathrm{Conv}(a,c) \to \varphi_{43}$
\item $\varphi_{217} :\equiv \forall a\,(\mathrm{FLim}(f,c,\alpha) \to \mathrm{Conv}(a,c) \to \varphi_{44}) \to \mathrm{FLim}(f,c,\alpha) \to \mathrm{Conv}(a,c) \to \varphi_{44}$
\item $\varphi_{218} :\equiv \forall y\,((f + g) \circ a(n) = y \to y = (f + g) \circ a(n)) \to \varphi_{135}$
\item $\varphi_{219} :\equiv \mathrm{FLim}(f,c,\alpha) \to \mathrm{FLim}(g,c,\beta) \to \mathrm{Seq}(a) \to \mathrm{Conv}(a,c) \to \varphi_{162}$
\item $\varphi_{220} :\equiv \mathrm{FLim}(f,c,\alpha) \to \mathrm{FLim}(g,c,\beta) \to \mathrm{Seq}(a) \to \varphi_{176}$
\item $\varphi_{221} :\equiv \forall \alpha\,(\mathrm{RFn}((f + g)) \to c \in \mathbb{R} \to \alpha \in \mathbb{R} \to \mathrm{PNbhd}((f + g),c) \to \forall a\,(\mathrm{Seq}(a) \to \mathrm{Conv}(a,c) \to \forall n \in \mathbb{N}\,(a_{n} \in \mathrm{dom}((f + g)) \land a_{n} \neq c) \to \mathrm{Conv}((f + g) \circ a,\alpha)) \to \mathrm{FLim}((f + g),c,\alpha))$
\item $\varphi_{222} :\equiv (\mathrm{FLim}(f,c,\alpha) \to \mathrm{FLim}(g,c,\beta) \to \mathrm{RFn}(g)) \to \varphi_{177}$
\item $\varphi_{223} :\equiv d \in \mathbb{R} \land 0 < d \land \forall x \in \mathbb{R}\,(\varphi_{30}) \to d \in \mathbb{R} \land 0 < d \land \forall x \in \mathbb{R}\,(\varphi_{30})$
\item $\varphi_{224} :\equiv \forall \alpha\,\forall a\,(\mathrm{FLim}(f,c,\alpha) \to \mathrm{Conv}(a,c) \to \varphi_{44}) \to \forall a\,(\mathrm{FLim}(f,c,\alpha) \to \mathrm{Conv}(a,c) \to \varphi_{44})$
\item $\varphi_{225} :\equiv \mathrm{FLim}(f,c,\alpha) \to \mathrm{FLim}(g,c,\beta) \to \mathrm{Seq}(a) \to \varphi_{192}$
\item $\varphi_{226} :\equiv \mathrm{FLim}(g,c,\beta) \to d \in \mathbb{R} \land 0 < d \land \forall x \in \mathbb{R}\,(\varphi_{30}) \to d \in \mathbb{R} \land 0 < d \land \varphi_{48}$
\item $\varphi_{227} :\equiv \varphi_{103} \to \forall \alpha\,\forall a\,(\mathrm{FLim}(f,c,\alpha) \to \mathrm{Conv}(a,c) \to \varphi_{44})$
\item $\varphi_{228} :\equiv \mathrm{RFn}((f + g)) \to c \in \mathbb{R} \to \alpha + \beta \in \mathbb{R} \to \mathrm{PNbhd}((f + g),c) \to \varphi_{151}$
\item $\varphi_{229} :\equiv \mathrm{FLim}(f,c,\alpha) \to \mathrm{FLim}(g,c,\beta) \to \mathrm{Seq}(a) \to \mathrm{Conv}(a,c) \to \varphi_{181}$
\item $\varphi_{230} :\equiv \mathrm{FLim}(g,c,\beta) \to d \in \mathbb{R} \land 0 < d \land \forall x \in \mathbb{R}\,(\varphi_{30}) \to \varphi_{84}$
\item $\varphi_{231} :\equiv (\varphi_{11} \to \varphi_{39}) \to \mathrm{FLim}(f,c,\alpha) \to \mathrm{FLim}(g,c,\beta) \to \varphi_{39}$
\item $\varphi_{232} :\equiv d \in \mathbb{R} \land 0 < d \land \forall x \in \mathbb{R}\,(\varphi_{30}) \to x \in \mathbb{R} \to 0 < \lvert x - c \rvert \to \lvert x - c \rvert < d \to x \in \mathbb{R} \to \varphi_{30}$
\item $\varphi_{233} :\equiv d \in \mathbb{R} \land 0 < d \land \forall x \in \mathbb{R}\,(\varphi_{30}) \to x \in \mathbb{R} \to 0 < \lvert x - c \rvert \to \lvert x - c \rvert < d \to \forall x \in \mathbb{R}\,(\varphi_{30})$
\item $\varphi_{234} :\equiv \forall y\,(f \circ a \oplus g \circ a(n) = y \to y = f \circ a \oplus g \circ a(n)) \to \varphi_{156}$
\item $\varphi_{235} :\equiv (\mathrm{FLim}(f,c,\alpha) \to \mathrm{FLim}(g,c,\beta) \to \mathrm{RFn}(f)) \to \varphi_{222}$
\item $\varphi_{236} :\equiv \forall c\,(\mathrm{PNbhd}((f + g),c) \leftrightarrow \varphi_{84}) \to (\mathrm{PNbhd}((f + g),c) \leftrightarrow \varphi_{84})$
\item $\varphi_{237} :\equiv \mathrm{FLim}(f,c,\alpha) \to \mathrm{FLim}(g,c,\beta) \to \mathrm{Seq}(a) \to \mathrm{Conv}(a,c) \to \varphi_{212}$
\item $\varphi_{238} :\equiv \forall a\,(\mathrm{RFn}(g) \to \mathrm{Seq}(a) \to \varphi_{101}) \to \mathrm{RFn}(g) \to \mathrm{Seq}(a) \to \varphi_{101}$
\item $\varphi_{239} :\equiv \forall a\,(\mathrm{RFn}(f) \to \mathrm{Seq}(a) \to \varphi_{102}) \to \mathrm{RFn}(f) \to \mathrm{Seq}(a) \to \varphi_{102}$
\item $\varphi_{240} :\equiv \forall f\,\forall a\,(\mathrm{RFn}(f) \to \mathrm{Seq}(a) \to \varphi_{102}) \to \forall a\,(\mathrm{RFn}(g) \to \mathrm{Seq}(a) \to \varphi_{101})$
\item $\varphi_{241} :\equiv \forall f\,\forall a\,(\mathrm{RFn}(f) \to \mathrm{Seq}(a) \to \varphi_{102}) \to \forall a\,(\mathrm{RFn}(f) \to \mathrm{Seq}(a) \to \varphi_{102})$
\item $\varphi_{242} :\equiv (\varphi_{23} \to \varphi_{64}) \to \mathrm{FLim}(f,c,\alpha) \to \mathrm{FLim}(g,c,\beta) \to \varphi_{64}$
\item $\varphi_{243} :\equiv (\mathrm{FLim}(f,c,\alpha) \to \mathrm{FLim}(g,c,\beta) \to \varphi_{11}) \to \varphi_{93} \to \mathrm{FLim}(f,c,\alpha) \to \varphi_{29}$
\item $\varphi_{244} :\equiv \forall g\,\forall c\,(\mathrm{PNbhd}(f,c) \to \varphi_{111}) \to \forall c\,(\mathrm{PNbhd}(f,c) \to \varphi_{111})$
\item $\varphi_{245} :\equiv \forall f\,\forall g\,\forall c\,(\mathrm{PNbhd}(f,c) \to \varphi_{111}) \to \forall g\,\forall c\,(\mathrm{PNbhd}(f,c) \to \varphi_{111})$
\item $\varphi_{246} :\equiv \forall n\,(\varphi_{98} \to n \in \mathbb{N} \to a_{n} \in \mathrm{dom}(g)) \to \forall n\,(\varphi_{98}) \to \forall n \in \mathbb{N}\,(a_{n} \in \mathrm{dom}(g))$
\item $\varphi_{247} :\equiv \forall n\,(\varphi_{98} \to n \in \mathbb{N} \to a_{n} \in \mathrm{dom}(f)) \to \forall n\,(\varphi_{98}) \to \forall n \in \mathbb{N}\,(a_{n} \in \mathrm{dom}(f))$
\item $\varphi_{248} :\equiv (\mathrm{FLim}(f,c,\alpha) \to \mathrm{FLim}(g,c,\beta) \to \alpha + \beta \in \mathbb{R}) \to \varphi_{7} \to (\mathrm{FLim}(f,c,\alpha) \to \varphi_{127}) \to \varphi_{14}$
\item $\varphi_{249} :\equiv \forall n\,(\varphi_{98} \to n \in \mathbb{N} \to a_{n} \in \mathrm{dom}((f + g))) \to \forall n\,(\varphi_{98}) \to \forall n \in \mathbb{N}\,(a_{n} \in \mathrm{dom}((f + g)))$
\item $\varphi_{250} :\equiv (\mathrm{FLim}(f,c,\alpha) \to \mathrm{FLim}(g,c,\beta) \to \varphi_{117}) \to \varphi_{10} \to \neg (d \in \mathbb{R} \land 0 < d \land \forall x \in \mathbb{R}\,(\varphi_{30}))$
\item $\varphi_{251} :\equiv \forall n\,(\varphi_{98} \to n \in \mathbb{N} \to a_{n} \in \mathrm{dom}(g) \land a_{n} \neq c) \to \forall n\,(\varphi_{98}) \to \forall n \in \mathbb{N}\,(a_{n} \in \mathrm{dom}(g) \land a_{n} \neq c)$
\item $\varphi_{252} :\equiv \forall n\,(\varphi_{98} \to n \in \mathbb{N} \to a_{n} \in \mathrm{dom}(f) \land a_{n} \neq c) \to \forall n\,(\varphi_{98}) \to \forall n \in \mathbb{N}\,(a_{n} \in \mathrm{dom}(f) \land a_{n} \neq c)$
\item $\varphi_{253} :\equiv \forall d\,(\varphi_{10} \to \neg (d \in \mathbb{R} \land 0 < d \land \forall x \in \mathbb{R}\,(\varphi_{30}))) \to \forall d\,(\varphi_{10}) \to \forall d\,(\neg (d \in \mathbb{R} \land 0 < d \land \forall x \in \mathbb{R}\,(\varphi_{30})))$
\item $\varphi_{254} :\equiv (\mathrm{FLim}(f,c,\alpha) \to \mathrm{FLim}(g,c,\beta) \to c \in \mathbb{R}) \to \varphi_{248}$
\item $\varphi_{255} :\equiv (\mathrm{PNbhd}(f,c) \to \varphi_{111}) \to \varphi_{2} \to \varphi_{1} \to \mathrm{FLim}(f,c,\alpha) \to \varphi_{115}$
\item $\varphi_{256} :\equiv \varphi_{122} \to (\mathrm{FLim}(f,c,\alpha) \to \mathrm{FLim}(g,c,\beta) \to \varphi_{39}) \to \varphi_{140}$
\item $\varphi_{257} :\equiv \varphi_{122} \to (\mathrm{FLim}(f,c,\alpha) \to \mathrm{FLim}(g,c,\beta) \to \varphi_{39}) \to \varphi_{141}$
\item $\varphi_{258} :\equiv (\mathrm{FLim}(f,c,\alpha) \to \mathrm{FLim}(f,c,\alpha)) \to (\mathrm{FLim}(f,c,\alpha) \to \varphi_{6}) \to \varphi_{171} \to \varphi_{124}$
\item $\varphi_{259} :\equiv (\mathrm{FLim}(f,c,\alpha) \to \varphi_{230}) \to \mathrm{FLim}(f,c,\alpha) \to \mathrm{FLim}(g,c,\beta) \to \varphi_{117}$
\item $\varphi_{260} :\equiv \varphi_{122} \to (\mathrm{FLim}(f,c,\alpha) \to \mathrm{FLim}(g,c,\beta) \to \varphi_{39}) \to \varphi_{164}$
\item $\varphi_{261} :\equiv \varphi_{122} \to (\mathrm{FLim}(f,c,\alpha) \to \mathrm{FLim}(g,c,\beta) \to \varphi_{39}) \to \varphi_{165}$
\item $\varphi_{262} :\equiv (\mathrm{RFn}(f) \to \mathrm{RFn}(g) \to \mathrm{RFn}((f + g)) \land \varphi_{11} \land \varphi_{23}) \to \varphi_{235}$
\item $\varphi_{263} :\equiv (\mathrm{FLim}(f,c,\alpha) \to \mathrm{FLim}(g,c,\beta) \to \mathrm{RFn}(g)) \to \varphi_{5} \to \varphi_{143} \to \varphi_{201}$
\item $\varphi_{264} :\equiv (\mathrm{FLim}(f,c,\alpha) \to \mathrm{FLim}(g,c,\beta) \to \mathrm{RFn}(f)) \to \varphi_{5} \to \varphi_{144} \to \varphi_{202}$
\item $\varphi_{265} :\equiv (\forall x \in \mathbb{R}\,(\varphi_{30}) \to x \in \mathbb{R} \to \varphi_{30}) \to \mathrm{FLim}(f,c,\alpha) \to \mathrm{FLim}(g,c,\beta) \to \varphi_{232}$
\item $\varphi_{266} :\equiv (\varphi_{87} \to \neg (d \in \mathbb{R} \land 0 < d \land \varphi_{48})) \to \mathrm{FLim}(f,c,\alpha) \to \varphi_{230}$
\item $\varphi_{267} :\equiv f(a_{n}) + g \circ a(n) = f(a_{n}) + g \circ a(n) \to (g \circ a(n) = g(a_{n}) \to \varphi_{159}) \to \varphi_{207}$
\item $\varphi_{268} :\equiv (\mathrm{FLim}(f,c,\alpha) \to \mathrm{FLim}(g,c,\beta) \to \varphi_{189}) \to (\mathrm{FLim}(f,c,\alpha) \to \varphi_{29}) \to \mathrm{FLim}(f,c,\alpha) \to \varphi_{205}$
\item $\varphi_{269} :\equiv f \circ a(n) + g \circ a(n) = f \circ a(n) + g \circ a(n) \to (f \circ a(n) = f(a_{n}) \to \varphi_{184}) \to \varphi_{215}$
\item $\varphi_{270} :\equiv (\mathrm{FLim}(g,c,\beta) \to \mathrm{Conv}(a,c) \to \varphi_{43}) \to \varphi_{3} \to (\mathrm{FLim}(f,c,\alpha) \to \varphi_{6}) \to \varphi_{170} \to \varphi_{123}$
\item $\varphi_{271} :\equiv (\mathrm{FLim}(f,c,\alpha) \to \mathrm{FLim}(g,c,\beta) \to \varphi_{223}) \to \mathrm{FLim}(f,c,\alpha) \to \mathrm{FLim}(g,c,\beta) \to \varphi_{233}$
\item $\varphi_{272} :\equiv (\mathrm{FLim}(f,c,\alpha) \to \mathrm{FLim}(g,c,\beta) \to \varphi_{64}) \to \varphi_{173} \to \varphi_{122} \to \varphi_{193}$
\item $\varphi_{273} :\equiv \varphi_{207} \to (f \circ a \oplus g \circ a(n) = f(a_{n}) + g \circ a(n) \to \varphi_{155}) \to \varphi_{210}$
\item $\varphi_{274} :\equiv \forall d\,(\varphi_{10} \to \neg (d \in \mathbb{R} \land 0 < d \land \forall x \in \mathbb{R}\,(\varphi_{30}))) \to \varphi_{253} \to (\varphi_{10} \to \forall d\,(\varphi_{10})) \to \varphi_{7}$
\item $\varphi_{275} :\equiv \varphi_{215} \to (f \circ a \oplus g \circ a(n) = f \circ a(n) + g \circ a(n) \to \varphi_{172}) \to \varphi_{214}$
\item $\varphi_{276} :\equiv \forall a\,(\varphi_{148}) \to (\forall a\,(\varphi_{148}) \to \varphi_{157}) \to \varphi_{20} \to \mathrm{FLim}(f,c,\alpha) \to \varphi_{127}$
\item $\varphi_{277} :\equiv (\mathrm{FLim}(f,c,\alpha) \to \mathrm{FLim}(g,c,\beta) \to \varphi_{223}) \to (\mathrm{FLim}(f,c,\alpha) \to \varphi_{209}) \to \mathrm{FLim}(f,c,\alpha) \to \varphi_{226}$
\item $\varphi_{278} :\equiv \varphi_{193} \to \varphi_{211} \to ((f + g)(a_{n}) = f(a_{n}) + g(a_{n}) \to \varphi_{139}) \to \varphi_{195}$
\item $\varphi_{279} :\equiv \forall n\,(\varphi_{98} \to n \in \mathbb{N} \to a_{n} \in \mathrm{dom}(g)) \to \varphi_{246} \to (\varphi_{98} \to \forall n\,(\varphi_{98})) \to \varphi_{143}$
\item $\varphi_{280} :\equiv \forall n\,(\varphi_{98} \to n \in \mathbb{N} \to a_{n} \in \mathrm{dom}(f)) \to \varphi_{247} \to (\varphi_{98} \to \forall n\,(\varphi_{98})) \to \varphi_{144}$
\item $\varphi_{281} :\equiv \forall n\,(\varphi_{98} \to n \in \mathbb{N} \to a_{n} \in \mathrm{dom}((f + g))) \to \varphi_{249} \to (\varphi_{98} \to \forall n\,(\varphi_{98})) \to \varphi_{158}$
\item $\varphi_{282} :\equiv (\mathrm{FLim}(f,c,\alpha) \to \varphi_{146}) \to (\mathrm{FLim}(f,c,\alpha) \to \mathrm{FLim}(g,c,\beta) \to \varphi_{163}) \to \varphi_{268}$
\item $\varphi_{283} :\equiv \forall n\,(\varphi_{98} \to n \in \mathbb{N} \to a_{n} \in \mathrm{dom}(g) \land a_{n} \neq c) \to \varphi_{251} \to (\varphi_{98} \to \forall n\,(\varphi_{98})) \to \varphi_{170}$
\item $\varphi_{284} :\equiv \forall n\,(\varphi_{98} \to n \in \mathbb{N} \to a_{n} \in \mathrm{dom}(f) \land a_{n} \neq c) \to \varphi_{252} \to (\varphi_{98} \to \forall n\,(\varphi_{98})) \to \varphi_{171}$
\item $\varphi_{285} :\equiv (\forall x\,(\varphi_{128} \to \varphi_{40}) \to \forall x\,(\varphi_{128}) \to \varphi_{48}) \to (\varphi_{128} \to \forall x\,(\varphi_{128})) \to \mathrm{FLim}(f,c,\alpha) \to \varphi_{209}$
\end{abbrevlist}
\end{fullproof}

% 自動生成 (prover/build.js): FLimMul  ノード数 560
\begin{fullproof}{定理1.8(2)}{$\forall f\,\forall g\,\forall c\,\forall \alpha\,\forall \beta\,(\mathrm{FLim}(f,c,\alpha) \to \mathrm{FLim}(g,c,\beta) \to \mathrm{FLim}((f \cdot g),c,\alpha \cdot \beta))$}{560}\phantomsection\label{tree:FLimMul}
\noindent$\Pi_{1}$:
\begin{ptree}
\AXC{$\forall f\,\forall c\,\forall \alpha\,(\mathrm{FLim}(f,c,\alpha) \to \mathrm{RFn}(f))$\lbl{[D2]}}
\AXC{$\varphi_{21}$\lbl{A2}}
\BIC{$\forall c\,\forall \alpha\,(\mathrm{FLim}(g,c,\alpha) \to \mathrm{RFn}(g))$}
\AXC{$\varphi_{12}$\lbl{A2}}
\BIC{$\forall \alpha\,(\mathrm{FLim}(g,c,\alpha) \to \mathrm{RFn}(g))$}
\AXC{$\varphi_{8}$\lbl{A2}}
\BIC{$\mathrm{FLim}(g,c,\beta) \to \mathrm{RFn}(g)$}
\end{ptree}
\noindent$\Pi_{2}$:
\begin{ptree}
\AXC{$\varphi_{3}$\lbl{A1}}
\AXC{$\stackrel{\Pi_{1}}{\vdots}$}\noLine
\UIC{$\mathrm{FLim}(g,c,\beta) \to \mathrm{RFn}(g)$}
\AXC{$(\mathrm{FLim}(g,c,\beta) \to \mathrm{RFn}(g)) \to \varphi_{56}$\lbl{A1}}
\BIC{$\varphi_{56}$}
\BIC{$\mathrm{FLim}(f,c,\alpha) \to \mathrm{FLim}(g,c,\beta) \to \mathrm{RFn}(g)$}
\end{ptree}
\noindent$\Pi_{3}$:
\begin{ptree}
\AXC{$\forall f\,\forall c\,\forall \alpha\,(\mathrm{FLim}(f,c,\alpha) \to \mathrm{RFn}(f))$\lbl{[D2]}}
\AXC{$\varphi_{22}$\lbl{A2}}
\BIC{$\forall c\,\forall \alpha\,(\mathrm{FLim}(f,c,\alpha) \to \mathrm{RFn}(f))$}
\AXC{$\varphi_{13}$\lbl{A2}}
\BIC{$\forall \alpha\,(\mathrm{FLim}(f,c,\alpha) \to \mathrm{RFn}(f))$}
\AXC{$\varphi_{9}$\lbl{A2}}
\BIC{$\mathrm{FLim}(f,c,\alpha) \to \mathrm{RFn}(f)$}
\end{ptree}
\noindent$\Pi_{4}$:
\begin{ptree}
\AXC{$\stackrel{\Pi_{3}}{\vdots}$}\noLine
\UIC{$\mathrm{FLim}(f,c,\alpha) \to \mathrm{RFn}(f)$}
\AXC{$(\mathrm{FLim}(f,c,\alpha) \to \mathrm{RFn}(f)) \to \varphi_{30}$\lbl{A1}}
\BIC{$\varphi_{30}$}
\end{ptree}
\noindent$\Pi_{5}$:
\begin{ptree}
\AXC{$\mathrm{FLim}(f,c,\alpha) \to \mathrm{FLim}(f,c,\alpha)$\lbl{A1}}
\AXC{$\stackrel{\Pi_{4}}{\vdots}$}\noLine
\UIC{$\varphi_{30}$}
\BIC{$\mathrm{FLim}(f,c,\alpha) \to \mathrm{FLim}(g,c,\beta) \to \mathrm{RFn}(f)$}
\end{ptree}
\noindent$\Pi_{6}$:
\begin{ptree}
\AXC{$\forall f\,(\varphi_{166})$\lbl{[S18]}}
\AXC{$\forall f\,(\varphi_{166}) \to \varphi_{166}$\lbl{A2}}
\BIC{$\varphi_{166}$}
\AXC{$\varphi_{166} \to \mathrm{RFn}(f) \to \mathrm{RFn}(g) \to \mathrm{RFn}((f \cdot g)) \land \varphi_{11} \land \varphi_{23}$\lbl{A2}}
\BIC{$\mathrm{RFn}(f) \to \mathrm{RFn}(g) \to \mathrm{RFn}((f \cdot g)) \land \varphi_{11} \land \varphi_{23}$}
\AXC{$\varphi_{262}$\lbl{A1}}
\BIC{$\varphi_{234}$}
\end{ptree}
\noindent$\Pi_{7}$:
\begin{ptree}
\AXC{$\stackrel{\Pi_{2}}{\vdots}$}\noLine
\UIC{$\mathrm{FLim}(f,c,\alpha) \to \mathrm{FLim}(g,c,\beta) \to \mathrm{RFn}(g)$}
\AXC{$\stackrel{\Pi_{5}}{\vdots}$}\noLine
\UIC{$\mathrm{FLim}(f,c,\alpha) \to \mathrm{FLim}(g,c,\beta) \to \mathrm{RFn}(f)$}
\AXC{$\stackrel{\Pi_{6}}{\vdots}$}\noLine
\UIC{$\varphi_{234}$}
\BIC{$\varphi_{220}$}
\BIC{$\varphi_{177}$}
\end{ptree}
\noindent$\Pi_{8}$:
\begin{ptree}
\AXC{$\stackrel{\Pi_{7}}{\vdots}$}\noLine
\UIC{$\varphi_{177}$}
\AXC{$\varphi_{177} \to \mathrm{FLim}(f,c,\alpha) \to \mathrm{FLim}(g,c,\beta) \to \varphi_{11}$\lbl{A1}}
\BIC{$\mathrm{FLim}(f,c,\alpha) \to \mathrm{FLim}(g,c,\beta) \to \varphi_{11}$}
\end{ptree}
\noindent$\Pi_{9}$:
\begin{ptree}
\AXC{$\varphi_{11} \to \varphi_{39}$\lbl{A2}}
\AXC{$\stackrel{\Pi_{8}}{\vdots}$}\noLine
\UIC{$\mathrm{FLim}(f,c,\alpha) \to \mathrm{FLim}(g,c,\beta) \to \varphi_{11}$}
\AXC{$(\mathrm{FLim}(f,c,\alpha) \to \mathrm{FLim}(g,c,\beta) \to \varphi_{11}) \to \varphi_{231}$\lbl{A1}}
\BIC{$\varphi_{231}$}
\BIC{$\mathrm{FLim}(f,c,\alpha) \to \mathrm{FLim}(g,c,\beta) \to \varphi_{39}$}
\end{ptree}
\noindent$\Pi_{10}$:
\begin{ptree}
\AXC{$\stackrel{\Pi_{9}}{\vdots}$}\noLine
\UIC{$\mathrm{FLim}(f,c,\alpha) \to \mathrm{FLim}(g,c,\beta) \to \varphi_{39}$}
\AXC{$\varphi_{122}$\lbl{A1}}
\AXC{$\varphi_{105}$\lbl{A2}}
\AXC{$\varphi_{178}$\lbl{A1}}
\AXC{$\varphi_{178} \to \varphi_{105} \to \varphi_{172}$\lbl{A1}}
\BIC{$\varphi_{105} \to \varphi_{172}$}
\BIC{$\varphi_{172}$}
\AXC{$\varphi_{172} \to \varphi_{260}$\lbl{A1}}
\BIC{$\varphi_{260}$}
\BIC{$(\mathrm{FLim}(f,c,\alpha) \to \mathrm{FLim}(g,c,\beta) \to \varphi_{39}) \to \varphi_{164}$}
\BIC{$\varphi_{164}$}
\end{ptree}
\noindent$\Pi_{11}$:
\begin{ptree}
\AXC{$\forall f\,(\varphi_{103})$\lbl{定理1.7(⇒)}}
\AXC{$\forall f\,(\varphi_{103}) \to \forall c\,(\varphi_{91})$\lbl{A2}}
\BIC{$\forall c\,(\varphi_{91})$}
\AXC{$\forall c\,(\varphi_{91}) \to \varphi_{91}$\lbl{A2}}
\BIC{$\varphi_{91}$}
\AXC{$\varphi_{91} \to \forall a\,(\mathrm{FLim}(g,c,\beta) \to \mathrm{Conv}(a,c) \to \varphi_{43})$\lbl{A2}}
\BIC{$\forall a\,(\mathrm{FLim}(g,c,\beta) \to \mathrm{Conv}(a,c) \to \varphi_{43})$}
\end{ptree}
\noindent$\Pi_{12}$:
\begin{ptree}
\AXC{$\mathrm{FLim}(f,c,\alpha) \to \varphi_{6}$\lbl{A1}}
\AXC{$\varphi_{3}$\lbl{A1}}
\AXC{$\stackrel{\Pi_{11}}{\vdots}$}\noLine
\UIC{$\forall a\,(\mathrm{FLim}(g,c,\beta) \to \mathrm{Conv}(a,c) \to \varphi_{43})$}
\AXC{$\varphi_{216}$\lbl{A2}}
\BIC{$\mathrm{FLim}(g,c,\beta) \to \mathrm{Conv}(a,c) \to \varphi_{43}$}
\AXC{$\varphi_{270}$\lbl{A1}}
\BIC{$\varphi_{3} \to (\mathrm{FLim}(f,c,\alpha) \to \varphi_{6}) \to \varphi_{169} \to \varphi_{123}$}
\BIC{$(\mathrm{FLim}(f,c,\alpha) \to \varphi_{6}) \to \varphi_{169} \to \varphi_{123}$}
\BIC{$\varphi_{169} \to \varphi_{123}$}
\end{ptree}
\noindent$\Pi_{13}$:
\begin{ptree}
\AXC{$\varphi_{95} \to \forall n\,(\varphi_{95})$\lbl{A4}}
\AXC{$\varphi_{251}$\lbl{A3}}
\AXC{$\stackrel{\Pi_{10}}{\vdots}$}\noLine
\UIC{$\varphi_{164}$}
\AXC{$\varphi_{164} \to \varphi_{95} \to n \in \mathbb{N} \to a_{n} \in \mathrm{dom}(g) \land a_{n} \neq c$\lbl{A1}}
\BIC{$\varphi_{95} \to n \in \mathbb{N} \to a_{n} \in \mathrm{dom}(g) \land a_{n} \neq c$}
\UIC{$\forall n\,(\varphi_{95} \to n \in \mathbb{N} \to a_{n} \in \mathrm{dom}(g) \land a_{n} \neq c)$}
\AXC{$\varphi_{283}$\lbl{A1}}
\BIC{$\varphi_{251} \to (\varphi_{95} \to \forall n\,(\varphi_{95})) \to \varphi_{169}$}
\BIC{$(\varphi_{95} \to \forall n\,(\varphi_{95})) \to \varphi_{169}$}
\BIC{$\varphi_{169}$}
\AXC{$\stackrel{\Pi_{12}}{\vdots}$}\noLine
\UIC{$\varphi_{169} \to \varphi_{123}$}
\BIC{$\varphi_{123}$}
\end{ptree}
\noindent$\Pi_{14}$:
\begin{ptree}
\AXC{$\varphi_{11} \to \varphi_{39}$\lbl{A2}}
\AXC{$\stackrel{\Pi_{8}}{\vdots}$}\noLine
\UIC{$\mathrm{FLim}(f,c,\alpha) \to \mathrm{FLim}(g,c,\beta) \to \varphi_{11}$}
\AXC{$(\mathrm{FLim}(f,c,\alpha) \to \mathrm{FLim}(g,c,\beta) \to \varphi_{11}) \to \varphi_{231}$\lbl{A1}}
\BIC{$\varphi_{231}$}
\BIC{$\mathrm{FLim}(f,c,\alpha) \to \mathrm{FLim}(g,c,\beta) \to \varphi_{39}$}
\end{ptree}
\noindent$\Pi_{15}$:
\begin{ptree}
\AXC{$\stackrel{\Pi_{14}}{\vdots}$}\noLine
\UIC{$\mathrm{FLim}(f,c,\alpha) \to \mathrm{FLim}(g,c,\beta) \to \varphi_{39}$}
\AXC{$\varphi_{122}$\lbl{A1}}
\AXC{$\varphi_{105}$\lbl{A2}}
\AXC{$\varphi_{178}$\lbl{A1}}
\AXC{$\varphi_{178} \to \varphi_{105} \to \varphi_{172}$\lbl{A1}}
\BIC{$\varphi_{105} \to \varphi_{172}$}
\BIC{$\varphi_{172}$}
\AXC{$\varphi_{172} \to \varphi_{261}$\lbl{A1}}
\BIC{$\varphi_{261}$}
\BIC{$(\mathrm{FLim}(f,c,\alpha) \to \mathrm{FLim}(g,c,\beta) \to \varphi_{39}) \to \varphi_{165}$}
\BIC{$\varphi_{165}$}
\end{ptree}
\noindent$\Pi_{16}$:
\begin{ptree}
\AXC{$\varphi_{95} \to \forall n\,(\varphi_{95})$\lbl{A4}}
\AXC{$\varphi_{252}$\lbl{A3}}
\AXC{$\stackrel{\Pi_{15}}{\vdots}$}\noLine
\UIC{$\varphi_{165}$}
\AXC{$\varphi_{165} \to \varphi_{95} \to n \in \mathbb{N} \to a_{n} \in \mathrm{dom}(f) \land a_{n} \neq c$\lbl{A1}}
\BIC{$\varphi_{95} \to n \in \mathbb{N} \to a_{n} \in \mathrm{dom}(f) \land a_{n} \neq c$}
\UIC{$\forall n\,(\varphi_{95} \to n \in \mathbb{N} \to a_{n} \in \mathrm{dom}(f) \land a_{n} \neq c)$}
\AXC{$\varphi_{284}$\lbl{A1}}
\BIC{$\varphi_{252} \to (\varphi_{95} \to \forall n\,(\varphi_{95})) \to \varphi_{170}$}
\BIC{$(\varphi_{95} \to \forall n\,(\varphi_{95})) \to \varphi_{170}$}
\BIC{$\varphi_{170}$}
\end{ptree}
\noindent$\Pi_{17}$:
\begin{ptree}
\AXC{$\forall f\,(\varphi_{103})$\lbl{定理1.7(⇒)}}
\AXC{$\forall f\,(\varphi_{103}) \to \varphi_{103}$\lbl{A2}}
\BIC{$\varphi_{103}$}
\AXC{$\varphi_{227}$\lbl{A2}}
\BIC{$\forall \alpha\,\forall a\,(\mathrm{FLim}(f,c,\alpha) \to \mathrm{Conv}(a,c) \to \varphi_{44})$}
\AXC{$\varphi_{224}$\lbl{A2}}
\BIC{$\forall a\,(\mathrm{FLim}(f,c,\alpha) \to \mathrm{Conv}(a,c) \to \varphi_{44})$}
\AXC{$\varphi_{217}$\lbl{A2}}
\BIC{$\mathrm{FLim}(f,c,\alpha) \to \mathrm{Conv}(a,c) \to \varphi_{44}$}
\end{ptree}
\noindent$\Pi_{18}$:
\begin{ptree}
\AXC{$\stackrel{\Pi_{17}}{\vdots}$}\noLine
\UIC{$\mathrm{FLim}(f,c,\alpha) \to \mathrm{Conv}(a,c) \to \varphi_{44}$}
\AXC{$(\mathrm{FLim}(f,c,\alpha) \to \mathrm{Conv}(a,c) \to \varphi_{44}) \to \varphi_{258}$\lbl{A1}}
\BIC{$\varphi_{258}$}
\end{ptree}
\noindent$\Pi_{19}$:
\begin{ptree}
\AXC{$\stackrel{\Pi_{16}}{\vdots}$}\noLine
\UIC{$\varphi_{170}$}
\AXC{$\mathrm{FLim}(f,c,\alpha) \to \varphi_{6}$\lbl{A1}}
\AXC{$\mathrm{FLim}(f,c,\alpha) \to \mathrm{FLim}(f,c,\alpha)$\lbl{A1}}
\AXC{$\stackrel{\Pi_{18}}{\vdots}$}\noLine
\UIC{$\varphi_{258}$}
\BIC{$(\mathrm{FLim}(f,c,\alpha) \to \varphi_{6}) \to \varphi_{170} \to \varphi_{124}$}
\BIC{$\varphi_{170} \to \varphi_{124}$}
\BIC{$\varphi_{124}$}
\end{ptree}
\noindent$\Pi_{20}$:
\begin{ptree}
\AXC{$\varphi_{31}$\lbl{定理2.7}}
\AXC{$\varphi_{31} \to \forall \alpha\,(\varphi_{57})$\lbl{A2}}
\BIC{$\forall \alpha\,(\varphi_{57})$}
\AXC{$\forall \alpha\,(\varphi_{57}) \to \varphi_{57}$\lbl{A2}}
\BIC{$\varphi_{57}$}
\AXC{$\varphi_{57} \to \forall \beta\,(\varphi_{77})$\lbl{A2}}
\BIC{$\forall \beta\,(\varphi_{77})$}
\AXC{$\forall \beta\,(\varphi_{77}) \to \varphi_{77}$\lbl{A2}}
\BIC{$\varphi_{77}$}
\AXC{$\varphi_{77} \to \varphi_{124} \to \varphi_{123} \to \varphi_{171}$\lbl{A1}}
\BIC{$\varphi_{124} \to \varphi_{123} \to \varphi_{171}$}
\end{ptree}
\noindent$\Pi_{21}$:
\begin{ptree}
\AXC{$x = x$\lbl{A5}}
\AXC{$x = y \to x = x \to y = x$\lbl{A6}}
\AXC{$(x = y \to x = x \to y = x) \to x = x \to x = y \to y = x$\lbl{A1}}
\BIC{$x = x \to x = y \to y = x$}
\BIC{$x = y \to y = x$}
\UIC{$\forall y\,(x = y \to y = x)$}
\UIC{$\forall x\,\forall y\,(x = y \to y = x)$}
\end{ptree}
\noindent$\Pi_{22}$:
\begin{ptree}
\AXC{$y = z \to x = y \to x = z$\lbl{A6}}
\AXC{$(y = z \to x = y \to x = z) \to x = y \to y = z \to x = z$\lbl{A1}}
\BIC{$x = y \to y = z \to x = z$}
\UIC{$\forall z\,(x = y \to y = z \to x = z)$}
\UIC{$\forall y\,\forall z\,(x = y \to y = z \to x = z)$}
\UIC{$\forall x\,\forall y\,\forall z\,(x = y \to y = z \to x = z)$}
\AXC{$\forall x\,\forall y\,\forall z\,(x = y \to y = z \to x = z) \to \varphi_{45}$\lbl{A2}}
\BIC{$\varphi_{45}$}
\end{ptree}
\noindent$\Pi_{23}$:
\begin{ptree}
\AXC{$y = z \to x = y \to x = z$\lbl{A6}}
\AXC{$(y = z \to x = y \to x = z) \to x = y \to y = z \to x = z$\lbl{A1}}
\BIC{$x = y \to y = z \to x = z$}
\UIC{$\forall z\,(x = y \to y = z \to x = z)$}
\UIC{$\forall y\,\forall z\,(x = y \to y = z \to x = z)$}
\UIC{$\forall x\,\forall y\,\forall z\,(x = y \to y = z \to x = z)$}
\AXC{$\forall x\,\forall y\,\forall z\,(x = y \to y = z \to x = z) \to \varphi_{27}$\lbl{A2}}
\BIC{$\varphi_{27}$}
\end{ptree}
\noindent$\Pi_{24}$:
\begin{ptree}
\AXC{$\stackrel{\Pi_{23}}{\vdots}$}\noLine
\UIC{$\varphi_{27}$}
\AXC{$\varphi_{27} \to \varphi_{120}$\lbl{A2}}
\BIC{$\varphi_{120}$}
\AXC{$\varphi_{120} \to (f \cdot g)(a_{n}) = f(a_{n}) \cdot g(a_{n}) \to \varphi_{141}$\lbl{A2}}
\BIC{$(f \cdot g)(a_{n}) = f(a_{n}) \cdot g(a_{n}) \to \varphi_{141}$}
\end{ptree}
\noindent$\Pi_{25}$:
\begin{ptree}
\AXC{$y = z \to x = y \to x = z$\lbl{A6}}
\AXC{$(y = z \to x = y \to x = z) \to x = y \to y = z \to x = z$\lbl{A1}}
\BIC{$x = y \to y = z \to x = z$}
\UIC{$\forall z\,(x = y \to y = z \to x = z)$}
\UIC{$\forall y\,\forall z\,(x = y \to y = z \to x = z)$}
\UIC{$\forall x\,\forall y\,\forall z\,(x = y \to y = z \to x = z)$}
\AXC{$\forall x\,\forall y\,\forall z\,(x = y \to y = z \to x = z) \to \varphi_{83}$\lbl{A2}}
\BIC{$\varphi_{83}$}
\end{ptree}
\noindent$\Pi_{26}$:
\begin{ptree}
\AXC{$\stackrel{\Pi_{25}}{\vdots}$}\noLine
\UIC{$\varphi_{83}$}
\AXC{$\varphi_{83} \to \varphi_{184}$\lbl{A2}}
\BIC{$\varphi_{184}$}
\AXC{$\varphi_{184} \to f \circ a \otimes g \circ a(n) = f(a_{n}) \cdot g \circ a(n) \to \varphi_{156}$\lbl{A2}}
\BIC{$f \circ a \otimes g \circ a(n) = f(a_{n}) \cdot g \circ a(n) \to \varphi_{156}$}
\end{ptree}
\noindent$\Pi_{27}$:
\begin{ptree}
\AXC{$z = w \to \varphi_{112}$\lbl{A6}}
\UIC{$\forall w\,(z = w \to \varphi_{112})$}
\UIC{$\forall z\,\forall w\,(z = w \to \varphi_{112})$}
\AXC{$\forall z\,\forall w\,(z = w \to \varphi_{112}) \to \varphi_{179}$\lbl{A2}}
\BIC{$\varphi_{179}$}
\AXC{$\varphi_{179} \to g \circ a(n) = g(a_{n}) \to \varphi_{159}$\lbl{A2}}
\BIC{$g \circ a(n) = g(a_{n}) \to \varphi_{159}$}
\end{ptree}
\noindent$\Pi_{28}$:
\begin{ptree}
\AXC{$\varphi_{11} \to \varphi_{39}$\lbl{A2}}
\AXC{$\stackrel{\Pi_{8}}{\vdots}$}\noLine
\UIC{$\mathrm{FLim}(f,c,\alpha) \to \mathrm{FLim}(g,c,\beta) \to \varphi_{11}$}
\AXC{$(\mathrm{FLim}(f,c,\alpha) \to \mathrm{FLim}(g,c,\beta) \to \varphi_{11}) \to \varphi_{231}$\lbl{A1}}
\BIC{$\varphi_{231}$}
\BIC{$\mathrm{FLim}(f,c,\alpha) \to \mathrm{FLim}(g,c,\beta) \to \varphi_{39}$}
\end{ptree}
\noindent$\Pi_{29}$:
\begin{ptree}
\AXC{$\stackrel{\Pi_{28}}{\vdots}$}\noLine
\UIC{$\mathrm{FLim}(f,c,\alpha) \to \mathrm{FLim}(g,c,\beta) \to \varphi_{39}$}
\AXC{$\varphi_{122}$\lbl{A1}}
\AXC{$\varphi_{105}$\lbl{A2}}
\AXC{$\varphi_{178}$\lbl{A1}}
\AXC{$\varphi_{178} \to \varphi_{105} \to \varphi_{172}$\lbl{A1}}
\BIC{$\varphi_{105} \to \varphi_{172}$}
\BIC{$\varphi_{172}$}
\AXC{$\varphi_{172} \to \varphi_{256}$\lbl{A1}}
\BIC{$\varphi_{256}$}
\BIC{$(\mathrm{FLim}(f,c,\alpha) \to \mathrm{FLim}(g,c,\beta) \to \varphi_{39}) \to \varphi_{139}$}
\BIC{$\varphi_{139}$}
\end{ptree}
\noindent$\Pi_{30}$:
\begin{ptree}
\AXC{$\varphi_{95} \to \forall n\,(\varphi_{95})$\lbl{A4}}
\AXC{$\varphi_{246}$\lbl{A3}}
\AXC{$\stackrel{\Pi_{29}}{\vdots}$}\noLine
\UIC{$\varphi_{139}$}
\AXC{$\varphi_{139} \to \varphi_{95} \to n \in \mathbb{N} \to a_{n} \in \mathrm{dom}(g)$\lbl{A1}}
\BIC{$\varphi_{95} \to n \in \mathbb{N} \to a_{n} \in \mathrm{dom}(g)$}
\UIC{$\forall n\,(\varphi_{95} \to n \in \mathbb{N} \to a_{n} \in \mathrm{dom}(g))$}
\AXC{$\varphi_{279}$\lbl{A1}}
\BIC{$\varphi_{246} \to (\varphi_{95} \to \forall n\,(\varphi_{95})) \to \varphi_{143}$}
\BIC{$(\varphi_{95} \to \forall n\,(\varphi_{95})) \to \varphi_{143}$}
\BIC{$\varphi_{143}$}
\end{ptree}
\noindent$\Pi_{31}$:
\begin{ptree}
\AXC{$\forall f\,\forall a\,(\mathrm{RFn}(f) \to \mathrm{Seq}(a) \to \varphi_{100})$\lbl{[S8]}}
\AXC{$\varphi_{240}$\lbl{A2}}
\BIC{$\forall a\,(\mathrm{RFn}(g) \to \mathrm{Seq}(a) \to \varphi_{99})$}
\AXC{$\varphi_{238}$\lbl{A2}}
\BIC{$\mathrm{RFn}(g) \to \mathrm{Seq}(a) \to \varphi_{99}$}
\AXC{$(\mathrm{RFn}(g) \to \mathrm{Seq}(a) \to \varphi_{99}) \to \varphi_{263}$\lbl{A1}}
\BIC{$\varphi_{263}$}
\end{ptree}
\noindent$\Pi_{32}$:
\begin{ptree}
\AXC{$\stackrel{\Pi_{30}}{\vdots}$}\noLine
\UIC{$\varphi_{143}$}
\AXC{$\varphi_{5}$\lbl{A1}}
\AXC{$\stackrel{\Pi_{2}}{\vdots}$}\noLine
\UIC{$\mathrm{FLim}(f,c,\alpha) \to \mathrm{FLim}(g,c,\beta) \to \mathrm{RFn}(g)$}
\AXC{$\stackrel{\Pi_{31}}{\vdots}$}\noLine
\UIC{$\varphi_{263}$}
\BIC{$\varphi_{5} \to \varphi_{143} \to \varphi_{200}$}
\BIC{$\varphi_{143} \to \varphi_{200}$}
\BIC{$\varphi_{200}$}
\AXC{$\varphi_{200} \to \varphi_{186}$\lbl{A1}}
\BIC{$\varphi_{186}$}
\end{ptree}
\noindent$\Pi_{33}$:
\begin{ptree}
\AXC{$\varphi_{122}$\lbl{A1}}
\AXC{$\varphi_{85}$\lbl{A2}}
\AXC{$\stackrel{\Pi_{32}}{\vdots}$}\noLine
\UIC{$\varphi_{186}$}
\AXC{$\varphi_{186} \to \varphi_{85} \to \varphi_{180}$\lbl{A1}}
\BIC{$\varphi_{85} \to \varphi_{180}$}
\BIC{$\varphi_{180}$}
\AXC{$\varphi_{180} \to \varphi_{122} \to \varphi_{167}$\lbl{A1}}
\BIC{$\varphi_{122} \to \varphi_{167}$}
\BIC{$\varphi_{167}$}
\AXC{$\varphi_{167} \to \varphi_{267}$\lbl{A1}}
\BIC{$\varphi_{267}$}
\end{ptree}
\noindent$\Pi_{34}$:
\begin{ptree}
\AXC{$y = z \to x = y \to x = z$\lbl{A6}}
\AXC{$(y = z \to x = y \to x = z) \to x = y \to y = z \to x = z$\lbl{A1}}
\BIC{$x = y \to y = z \to x = z$}
\UIC{$\forall z\,(x = y \to y = z \to x = z)$}
\UIC{$\forall y\,\forall z\,(x = y \to y = z \to x = z)$}
\UIC{$\forall x\,\forall y\,\forall z\,(x = y \to y = z \to x = z)$}
\AXC{$\forall x\,\forall y\,\forall z\,(x = y \to y = z \to x = z) \to \varphi_{83}$\lbl{A2}}
\BIC{$\varphi_{83}$}
\end{ptree}
\noindent$\Pi_{35}$:
\begin{ptree}
\AXC{$\stackrel{\Pi_{34}}{\vdots}$}\noLine
\UIC{$\varphi_{83}$}
\AXC{$\varphi_{83} \to \varphi_{190}$\lbl{A2}}
\BIC{$\varphi_{190}$}
\AXC{$\varphi_{190} \to f \circ a \otimes g \circ a(n) = f \circ a(n) \cdot g \circ a(n) \to \varphi_{173}$\lbl{A2}}
\BIC{$f \circ a \otimes g \circ a(n) = f \circ a(n) \cdot g \circ a(n) \to \varphi_{173}$}
\end{ptree}
\noindent$\Pi_{36}$:
\begin{ptree}
\AXC{$z = w \to \varphi_{133}$\lbl{A6}}
\UIC{$\forall w\,(z = w \to \varphi_{133})$}
\UIC{$\forall z\,\forall w\,(z = w \to \varphi_{133})$}
\AXC{$\forall z\,\forall w\,(z = w \to \varphi_{133}) \to \varphi_{194}$\lbl{A2}}
\BIC{$\varphi_{194}$}
\AXC{$\varphi_{194} \to f \circ a(n) = f(a_{n}) \to \varphi_{182}$\lbl{A2}}
\BIC{$f \circ a(n) = f(a_{n}) \to \varphi_{182}$}
\end{ptree}
\noindent$\Pi_{37}$:
\begin{ptree}
\AXC{$\varphi_{11} \to \varphi_{39}$\lbl{A2}}
\AXC{$\stackrel{\Pi_{8}}{\vdots}$}\noLine
\UIC{$\mathrm{FLim}(f,c,\alpha) \to \mathrm{FLim}(g,c,\beta) \to \varphi_{11}$}
\AXC{$(\mathrm{FLim}(f,c,\alpha) \to \mathrm{FLim}(g,c,\beta) \to \varphi_{11}) \to \varphi_{231}$\lbl{A1}}
\BIC{$\varphi_{231}$}
\BIC{$\mathrm{FLim}(f,c,\alpha) \to \mathrm{FLim}(g,c,\beta) \to \varphi_{39}$}
\end{ptree}
\noindent$\Pi_{38}$:
\begin{ptree}
\AXC{$\stackrel{\Pi_{37}}{\vdots}$}\noLine
\UIC{$\mathrm{FLim}(f,c,\alpha) \to \mathrm{FLim}(g,c,\beta) \to \varphi_{39}$}
\AXC{$\varphi_{122}$\lbl{A1}}
\AXC{$\varphi_{105}$\lbl{A2}}
\AXC{$\varphi_{178}$\lbl{A1}}
\AXC{$\varphi_{178} \to \varphi_{105} \to \varphi_{172}$\lbl{A1}}
\BIC{$\varphi_{105} \to \varphi_{172}$}
\BIC{$\varphi_{172}$}
\AXC{$\varphi_{172} \to \varphi_{257}$\lbl{A1}}
\BIC{$\varphi_{257}$}
\BIC{$(\mathrm{FLim}(f,c,\alpha) \to \mathrm{FLim}(g,c,\beta) \to \varphi_{39}) \to \varphi_{140}$}
\BIC{$\varphi_{140}$}
\end{ptree}
\noindent$\Pi_{39}$:
\begin{ptree}
\AXC{$\varphi_{95} \to \forall n\,(\varphi_{95})$\lbl{A4}}
\AXC{$\varphi_{247}$\lbl{A3}}
\AXC{$\stackrel{\Pi_{38}}{\vdots}$}\noLine
\UIC{$\varphi_{140}$}
\AXC{$\varphi_{140} \to \varphi_{95} \to n \in \mathbb{N} \to a_{n} \in \mathrm{dom}(f)$\lbl{A1}}
\BIC{$\varphi_{95} \to n \in \mathbb{N} \to a_{n} \in \mathrm{dom}(f)$}
\UIC{$\forall n\,(\varphi_{95} \to n \in \mathbb{N} \to a_{n} \in \mathrm{dom}(f))$}
\AXC{$\varphi_{280}$\lbl{A1}}
\BIC{$\varphi_{247} \to (\varphi_{95} \to \forall n\,(\varphi_{95})) \to \varphi_{144}$}
\BIC{$(\varphi_{95} \to \forall n\,(\varphi_{95})) \to \varphi_{144}$}
\BIC{$\varphi_{144}$}
\end{ptree}
\noindent$\Pi_{40}$:
\begin{ptree}
\AXC{$\forall f\,\forall a\,(\mathrm{RFn}(f) \to \mathrm{Seq}(a) \to \varphi_{100})$\lbl{[S8]}}
\AXC{$\varphi_{241}$\lbl{A2}}
\BIC{$\forall a\,(\mathrm{RFn}(f) \to \mathrm{Seq}(a) \to \varphi_{100})$}
\AXC{$\varphi_{239}$\lbl{A2}}
\BIC{$\mathrm{RFn}(f) \to \mathrm{Seq}(a) \to \varphi_{100}$}
\AXC{$(\mathrm{RFn}(f) \to \mathrm{Seq}(a) \to \varphi_{100}) \to \varphi_{264}$\lbl{A1}}
\BIC{$\varphi_{264}$}
\end{ptree}
\noindent$\Pi_{41}$:
\begin{ptree}
\AXC{$\stackrel{\Pi_{39}}{\vdots}$}\noLine
\UIC{$\varphi_{144}$}
\AXC{$\varphi_{5}$\lbl{A1}}
\AXC{$\stackrel{\Pi_{5}}{\vdots}$}\noLine
\UIC{$\mathrm{FLim}(f,c,\alpha) \to \mathrm{FLim}(g,c,\beta) \to \mathrm{RFn}(f)$}
\AXC{$\stackrel{\Pi_{40}}{\vdots}$}\noLine
\UIC{$\varphi_{264}$}
\BIC{$\varphi_{5} \to \varphi_{144} \to \varphi_{201}$}
\BIC{$\varphi_{144} \to \varphi_{201}$}
\BIC{$\varphi_{201}$}
\AXC{$\varphi_{201} \to \varphi_{187}$\lbl{A1}}
\BIC{$\varphi_{187}$}
\end{ptree}
\noindent$\Pi_{42}$:
\begin{ptree}
\AXC{$\varphi_{122}$\lbl{A1}}
\AXC{$\varphi_{86}$\lbl{A2}}
\AXC{$\stackrel{\Pi_{41}}{\vdots}$}\noLine
\UIC{$\varphi_{187}$}
\AXC{$\varphi_{187} \to \varphi_{86} \to \varphi_{181}$\lbl{A1}}
\BIC{$\varphi_{86} \to \varphi_{181}$}
\BIC{$\varphi_{181}$}
\AXC{$\varphi_{181} \to \varphi_{122} \to \varphi_{168}$\lbl{A1}}
\BIC{$\varphi_{122} \to \varphi_{168}$}
\BIC{$\varphi_{168}$}
\AXC{$\varphi_{168} \to \varphi_{269}$\lbl{A1}}
\BIC{$\varphi_{269}$}
\end{ptree}
\noindent$\Pi_{43}$:
\begin{ptree}
\AXC{$\forall a\,\forall \alpha\,(\mathrm{Conv}(a,\alpha) \to \mathrm{Seq}(a))$\lbl{[C1]}}
\AXC{$\varphi_{46}$\lbl{A2}}
\BIC{$\forall \alpha\,(\mathrm{Conv}(g \circ a,\alpha) \to \mathrm{Seq}(g \circ a))$}
\AXC{$\varphi_{65}$\lbl{A2}}
\BIC{$\mathrm{Conv}(g \circ a,\beta) \to \mathrm{Seq}(g \circ a)$}
\AXC{$(\mathrm{Conv}(g \circ a,\beta) \to \mathrm{Seq}(g \circ a)) \to \varphi_{123} \to \varphi_{118}$\lbl{A1}}
\BIC{$\varphi_{123} \to \varphi_{118}$}
\end{ptree}
\noindent$\Pi_{44}$:
\begin{ptree}
\AXC{$\forall a\,\forall \alpha\,(\mathrm{Conv}(a,\alpha) \to \mathrm{Seq}(a))$\lbl{[C1]}}
\AXC{$\varphi_{47}$\lbl{A2}}
\BIC{$\forall \alpha\,(\mathrm{Conv}(f \circ a,\alpha) \to \mathrm{Seq}(f \circ a))$}
\AXC{$\varphi_{66}$\lbl{A2}}
\BIC{$\mathrm{Conv}(f \circ a,\alpha) \to \mathrm{Seq}(f \circ a)$}
\AXC{$(\mathrm{Conv}(f \circ a,\alpha) \to \mathrm{Seq}(f \circ a)) \to \varphi_{124} \to \varphi_{119}$\lbl{A1}}
\BIC{$\varphi_{124} \to \varphi_{119}$}
\end{ptree}
\noindent$\Pi_{45}$:
\begin{ptree}
\AXC{$\stackrel{\Pi_{13}}{\vdots}$}\noLine
\UIC{$\varphi_{123}$}
\AXC{$\stackrel{\Pi_{43}}{\vdots}$}\noLine
\UIC{$\varphi_{123} \to \varphi_{118}$}
\BIC{$\varphi_{118}$}
\AXC{$\stackrel{\Pi_{19}}{\vdots}$}\noLine
\UIC{$\varphi_{124}$}
\AXC{$\stackrel{\Pi_{44}}{\vdots}$}\noLine
\UIC{$\varphi_{124} \to \varphi_{119}$}
\BIC{$\varphi_{119}$}
\AXC{$\varphi_{79}$\lbl{[S4]}}
\AXC{$\varphi_{79} \to \varphi_{134}$\lbl{A2}}
\BIC{$\varphi_{134}$}
\AXC{$\varphi_{134} \to \varphi_{192}$\lbl{A2}}
\BIC{$\varphi_{192}$}
\AXC{$\varphi_{192} \to \varphi_{119} \to \varphi_{118} \to \varphi_{237}$\lbl{A1}}
\BIC{$\varphi_{119} \to \varphi_{118} \to \varphi_{237}$}
\BIC{$\varphi_{118} \to \varphi_{237}$}
\BIC{$\varphi_{237}$}
\end{ptree}
\noindent$\Pi_{46}$:
\begin{ptree}
\AXC{$\varphi_{122}$\lbl{A1}}
\AXC{$\varphi_{208}$\lbl{A2}}
\AXC{$\stackrel{\Pi_{45}}{\vdots}$}\noLine
\UIC{$\varphi_{237}$}
\AXC{$\varphi_{237} \to \varphi_{229}$\lbl{A1}}
\BIC{$\varphi_{229}$}
\AXC{$\varphi_{229} \to \varphi_{208} \to \varphi_{225}$\lbl{A1}}
\BIC{$\varphi_{208} \to \varphi_{225}$}
\BIC{$\varphi_{225}$}
\AXC{$\varphi_{225} \to \varphi_{122} \to \varphi_{221}$\lbl{A1}}
\BIC{$\varphi_{122} \to \varphi_{221}$}
\BIC{$\varphi_{221}$}
\AXC{$\varphi_{221} \to \varphi_{275}$\lbl{A1}}
\BIC{$\varphi_{275}$}
\end{ptree}
\noindent$\Pi_{47}$:
\begin{ptree}
\AXC{$\stackrel{\Pi_{36}}{\vdots}$}\noLine
\UIC{$f \circ a(n) = f(a_{n}) \to \varphi_{182}$}
\AXC{$x = x$\lbl{A5}}
\UIC{$\forall x\,(x = x)$}
\AXC{$\varphi_{73}$\lbl{A2}}
\BIC{$f \circ a(n) \cdot g \circ a(n) = f \circ a(n) \cdot g \circ a(n)$}
\AXC{$\stackrel{\Pi_{42}}{\vdots}$}\noLine
\UIC{$\varphi_{269}$}
\BIC{$(f \circ a(n) = f(a_{n}) \to \varphi_{182}) \to \varphi_{215}$}
\BIC{$\varphi_{215}$}
\AXC{$\stackrel{\Pi_{46}}{\vdots}$}\noLine
\UIC{$\varphi_{275}$}
\BIC{$(f \circ a \otimes g \circ a(n) = f \circ a(n) \cdot g \circ a(n) \to \varphi_{173}) \to \varphi_{214}$}
\end{ptree}
\noindent$\Pi_{48}$:
\begin{ptree}
\AXC{$\stackrel{\Pi_{35}}{\vdots}$}\noLine
\UIC{$f \circ a \otimes g \circ a(n) = f \circ a(n) \cdot g \circ a(n) \to \varphi_{173}$}
\AXC{$\stackrel{\Pi_{47}}{\vdots}$}\noLine
\UIC{$(f \circ a \otimes g \circ a(n) = f \circ a(n) \cdot g \circ a(n) \to \varphi_{173}) \to \varphi_{214}$}
\BIC{$\varphi_{214}$}
\AXC{$\varphi_{214} \to \varphi_{273}$\lbl{A1}}
\BIC{$\varphi_{273}$}
\end{ptree}
\noindent$\Pi_{49}$:
\begin{ptree}
\AXC{$\stackrel{\Pi_{27}}{\vdots}$}\noLine
\UIC{$g \circ a(n) = g(a_{n}) \to \varphi_{159}$}
\AXC{$x = x$\lbl{A5}}
\UIC{$\forall x\,(x = x)$}
\AXC{$\varphi_{62}$\lbl{A2}}
\BIC{$f(a_{n}) \cdot g \circ a(n) = f(a_{n}) \cdot g \circ a(n)$}
\AXC{$\stackrel{\Pi_{33}}{\vdots}$}\noLine
\UIC{$\varphi_{267}$}
\BIC{$(g \circ a(n) = g(a_{n}) \to \varphi_{159}) \to \varphi_{206}$}
\BIC{$\varphi_{206}$}
\AXC{$\stackrel{\Pi_{48}}{\vdots}$}\noLine
\UIC{$\varphi_{273}$}
\BIC{$(f \circ a \otimes g \circ a(n) = f(a_{n}) \cdot g \circ a(n) \to \varphi_{156}) \to \varphi_{210}$}
\end{ptree}
\noindent$\Pi_{50}$:
\begin{ptree}
\AXC{$\stackrel{\Pi_{26}}{\vdots}$}\noLine
\UIC{$f \circ a \otimes g \circ a(n) = f(a_{n}) \cdot g \circ a(n) \to \varphi_{156}$}
\AXC{$\stackrel{\Pi_{49}}{\vdots}$}\noLine
\UIC{$(f \circ a \otimes g \circ a(n) = f(a_{n}) \cdot g \circ a(n) \to \varphi_{156}) \to \varphi_{210}$}
\BIC{$\varphi_{210}$}
\AXC{$\varphi_{210} \to \varphi_{158} \to \varphi_{211}$\lbl{A1}}
\BIC{$\varphi_{158} \to \varphi_{211}$}
\end{ptree}
\noindent$\Pi_{51}$:
\begin{ptree}
\AXC{$\stackrel{\Pi_{21}}{\vdots}$}\noLine
\UIC{$\forall x\,\forall y\,(x = y \to y = x)$}
\AXC{$\varphi_{102}$\lbl{A2}}
\BIC{$\forall y\,(f \circ a \otimes g \circ a(n) = y \to y = f \circ a \otimes g \circ a(n))$}
\AXC{$\varphi_{235}$\lbl{A2}}
\BIC{$\varphi_{158}$}
\AXC{$\stackrel{\Pi_{50}}{\vdots}$}\noLine
\UIC{$\varphi_{158} \to \varphi_{211}$}
\BIC{$\varphi_{211}$}
\end{ptree}
\noindent$\Pi_{52}$:
\begin{ptree}
\AXC{$\stackrel{\Pi_{7}}{\vdots}$}\noLine
\UIC{$\varphi_{177}$}
\AXC{$\varphi_{177} \to \mathrm{FLim}(f,c,\alpha) \to \mathrm{FLim}(g,c,\beta) \to \varphi_{23}$\lbl{A1}}
\BIC{$\mathrm{FLim}(f,c,\alpha) \to \mathrm{FLim}(g,c,\beta) \to \varphi_{23}$}
\AXC{$(\mathrm{FLim}(f,c,\alpha) \to \mathrm{FLim}(g,c,\beta) \to \varphi_{23}) \to \varphi_{242}$\lbl{A1}}
\BIC{$\varphi_{242}$}
\end{ptree}
\noindent$\Pi_{53}$:
\begin{ptree}
\AXC{$\varphi_{122}$\lbl{A1}}
\AXC{$\varphi_{105}$\lbl{A2}}
\AXC{$\varphi_{178}$\lbl{A1}}
\AXC{$\varphi_{178} \to \varphi_{105} \to \varphi_{172}$\lbl{A1}}
\BIC{$\varphi_{105} \to \varphi_{172}$}
\BIC{$\varphi_{172}$}
\AXC{$\varphi_{23} \to \varphi_{64}$\lbl{A2}}
\AXC{$\stackrel{\Pi_{52}}{\vdots}$}\noLine
\UIC{$\varphi_{242}$}
\BIC{$\mathrm{FLim}(f,c,\alpha) \to \mathrm{FLim}(g,c,\beta) \to \varphi_{64}$}
\AXC{$\varphi_{272}$\lbl{A1}}
\BIC{$\varphi_{172} \to \varphi_{122} \to \varphi_{193}$}
\BIC{$\varphi_{122} \to \varphi_{193}$}
\BIC{$\varphi_{193}$}
\end{ptree}
\noindent$\Pi_{54}$:
\begin{ptree}
\AXC{$\stackrel{\Pi_{24}}{\vdots}$}\noLine
\UIC{$(f \cdot g)(a_{n}) = f(a_{n}) \cdot g(a_{n}) \to \varphi_{141}$}
\AXC{$\stackrel{\Pi_{51}}{\vdots}$}\noLine
\UIC{$\varphi_{211}$}
\AXC{$\stackrel{\Pi_{53}}{\vdots}$}\noLine
\UIC{$\varphi_{193}$}
\AXC{$\varphi_{278}$\lbl{A1}}
\BIC{$\varphi_{211} \to ((f \cdot g)(a_{n}) = f(a_{n}) \cdot g(a_{n}) \to \varphi_{141}) \to \varphi_{196}$}
\BIC{$((f \cdot g)(a_{n}) = f(a_{n}) \cdot g(a_{n}) \to \varphi_{141}) \to \varphi_{196}$}
\BIC{$\varphi_{196}$}
\end{ptree}
\noindent$\Pi_{55}$:
\begin{ptree}
\AXC{$\varphi_{122}$\lbl{A1}}
\AXC{$\varphi_{105}$\lbl{A2}}
\AXC{$\varphi_{178}$\lbl{A1}}
\AXC{$\varphi_{178} \to \varphi_{105} \to \varphi_{172}$\lbl{A1}}
\BIC{$\varphi_{105} \to \varphi_{172}$}
\BIC{$\varphi_{172}$}
\AXC{$\varphi_{172} \to \varphi_{122} \to \varphi_{152}$\lbl{A1}}
\BIC{$\varphi_{122} \to \varphi_{152}$}
\BIC{$\varphi_{152}$}
\AXC{$\varphi_{152} \to \varphi_{95} \to n \in \mathbb{N} \to a_{n} \in \mathrm{dom}((f \cdot g))$\lbl{A1}}
\BIC{$\varphi_{95} \to n \in \mathbb{N} \to a_{n} \in \mathrm{dom}((f \cdot g))$}
\UIC{$\forall n\,(\varphi_{95} \to n \in \mathbb{N} \to a_{n} \in \mathrm{dom}((f \cdot g)))$}
\end{ptree}
\noindent$\Pi_{56}$:
\begin{ptree}
\AXC{$\varphi_{95} \to \forall n\,(\varphi_{95})$\lbl{A4}}
\AXC{$\varphi_{249}$\lbl{A3}}
\AXC{$\stackrel{\Pi_{55}}{\vdots}$}\noLine
\UIC{$\forall n\,(\varphi_{95} \to n \in \mathbb{N} \to a_{n} \in \mathrm{dom}((f \cdot g)))$}
\AXC{$\varphi_{281}$\lbl{A1}}
\BIC{$\varphi_{249} \to (\varphi_{95} \to \forall n\,(\varphi_{95})) \to \varphi_{157}$}
\BIC{$(\varphi_{95} \to \forall n\,(\varphi_{95})) \to \varphi_{157}$}
\BIC{$\varphi_{157}$}
\end{ptree}
\noindent$\Pi_{57}$:
\begin{ptree}
\AXC{$\stackrel{\Pi_{7}}{\vdots}$}\noLine
\UIC{$\varphi_{177}$}
\AXC{$\varphi_{177} \to \varphi_{4}$\lbl{A1}}
\BIC{$\varphi_{4}$}
\end{ptree}
\noindent$\Pi_{58}$:
\begin{ptree}
\AXC{$\forall f\,\forall a\,(\mathrm{RFn}(f) \to \mathrm{Seq}(a) \to \varphi_{100})$\lbl{[S8]}}
\AXC{$\forall f\,\forall a\,(\mathrm{RFn}(f) \to \mathrm{Seq}(a) \to \varphi_{100}) \to \forall a\,(\varphi_{175})$\lbl{A2}}
\BIC{$\forall a\,(\varphi_{175})$}
\AXC{$\forall a\,(\varphi_{175}) \to \varphi_{175}$\lbl{A2}}
\BIC{$\varphi_{175}$}
\end{ptree}
\noindent$\Pi_{59}$:
\begin{ptree}
\AXC{$\stackrel{\Pi_{56}}{\vdots}$}\noLine
\UIC{$\varphi_{157}$}
\AXC{$\varphi_{5}$\lbl{A1}}
\AXC{$\stackrel{\Pi_{57}}{\vdots}$}\noLine
\UIC{$\varphi_{4}$}
\AXC{$\stackrel{\Pi_{58}}{\vdots}$}\noLine
\UIC{$\varphi_{175}$}
\AXC{$\varphi_{175} \to \varphi_{4} \to \varphi_{5} \to \varphi_{157} \to \varphi_{218}$\lbl{A1}}
\BIC{$\varphi_{4} \to \varphi_{5} \to \varphi_{157} \to \varphi_{218}$}
\BIC{$\varphi_{5} \to \varphi_{157} \to \varphi_{218}$}
\BIC{$\varphi_{157} \to \varphi_{218}$}
\BIC{$\varphi_{218}$}
\end{ptree}
\noindent$\Pi_{60}$:
\begin{ptree}
\AXC{$\varphi_{122}$\lbl{A1}}
\AXC{$\varphi_{125}$\lbl{A2}}
\AXC{$\stackrel{\Pi_{59}}{\vdots}$}\noLine
\UIC{$\varphi_{218}$}
\AXC{$\varphi_{218} \to \varphi_{197}$\lbl{A1}}
\BIC{$\varphi_{197}$}
\AXC{$\varphi_{197} \to \varphi_{125} \to \varphi_{195}$\lbl{A1}}
\BIC{$\varphi_{125} \to \varphi_{195}$}
\BIC{$\varphi_{195}$}
\AXC{$\varphi_{195} \to \varphi_{122} \to \varphi_{188}$\lbl{A1}}
\BIC{$\varphi_{122} \to \varphi_{188}$}
\BIC{$\varphi_{188}$}
\end{ptree}
\noindent$\Pi_{61}$:
\begin{ptree}
\AXC{$\stackrel{\Pi_{54}}{\vdots}$}\noLine
\UIC{$\varphi_{196}$}
\AXC{$\stackrel{\Pi_{60}}{\vdots}$}\noLine
\UIC{$\varphi_{188}$}
\AXC{$\varphi_{188} \to \varphi_{196} \to ((f \cdot g) \circ a(n) = (f \cdot g)(a_{n}) \to \varphi_{131}) \to \varphi_{198}$\lbl{A1}}
\BIC{$\varphi_{196} \to ((f \cdot g) \circ a(n) = (f \cdot g)(a_{n}) \to \varphi_{131}) \to \varphi_{198}$}
\BIC{$((f \cdot g) \circ a(n) = (f \cdot g)(a_{n}) \to \varphi_{131}) \to \varphi_{198}$}
\end{ptree}
\noindent$\Pi_{62}$:
\begin{ptree}
\AXC{$\stackrel{\Pi_{22}}{\vdots}$}\noLine
\UIC{$\varphi_{45}$}
\AXC{$\varphi_{45} \to \varphi_{106}$\lbl{A2}}
\BIC{$\varphi_{106}$}
\AXC{$\varphi_{106} \to (f \cdot g) \circ a(n) = (f \cdot g)(a_{n}) \to \varphi_{131}$\lbl{A2}}
\BIC{$(f \cdot g) \circ a(n) = (f \cdot g)(a_{n}) \to \varphi_{131}$}
\AXC{$\stackrel{\Pi_{61}}{\vdots}$}\noLine
\UIC{$((f \cdot g) \circ a(n) = (f \cdot g)(a_{n}) \to \varphi_{131}) \to \varphi_{198}$}
\BIC{$\varphi_{198}$}
\end{ptree}
\noindent$\Pi_{63}$:
\begin{ptree}
\AXC{$\stackrel{\Pi_{21}}{\vdots}$}\noLine
\UIC{$\forall x\,\forall y\,(x = y \to y = x)$}
\AXC{$\varphi_{58}$\lbl{A2}}
\BIC{$\forall y\,((f \cdot g) \circ a(n) = y \to y = (f \cdot g) \circ a(n))$}
\AXC{$\varphi_{222}$\lbl{A2}}
\BIC{$\varphi_{136}$}
\AXC{$\stackrel{\Pi_{62}}{\vdots}$}\noLine
\UIC{$\varphi_{198}$}
\AXC{$\varphi_{198} \to \varphi_{136} \to \varphi_{199}$\lbl{A1}}
\BIC{$\varphi_{136} \to \varphi_{199}$}
\BIC{$\varphi_{199}$}
\end{ptree}
\noindent$\Pi_{64}$:
\begin{ptree}
\AXC{$\varphi_{204} \to \varphi_{209}$\lbl{A3}}
\AXC{$\stackrel{\Pi_{63}}{\vdots}$}\noLine
\UIC{$\varphi_{199}$}
\AXC{$\varphi_{199} \to \varphi_{95} \to n \in \mathbb{N} \to f \circ a \otimes g \circ a(n) = (f \cdot g) \circ a(n)$\lbl{A1}}
\BIC{$\varphi_{95} \to n \in \mathbb{N} \to f \circ a \otimes g \circ a(n) = (f \cdot g) \circ a(n)$}
\UIC{$\varphi_{204}$}
\AXC{$\varphi_{204} \to (\varphi_{204} \to \varphi_{209}) \to (\varphi_{95} \to \forall n\,(\varphi_{95})) \to \varphi_{205}$\lbl{A1}}
\BIC{$(\varphi_{204} \to \varphi_{209}) \to (\varphi_{95} \to \forall n\,(\varphi_{95})) \to \varphi_{205}$}
\BIC{$(\varphi_{95} \to \forall n\,(\varphi_{95})) \to \varphi_{205}$}
\end{ptree}
\noindent$\Pi_{65}$:
\begin{ptree}
\AXC{$\stackrel{\Pi_{59}}{\vdots}$}\noLine
\UIC{$\varphi_{218}$}
\AXC{$\varphi_{218} \to \varphi_{129}$\lbl{A1}}
\BIC{$\varphi_{129}$}
\AXC{$\varphi_{51}$\lbl{[C13]}}
\AXC{$\varphi_{51} \to \varphi_{135}$\lbl{A2}}
\BIC{$\varphi_{135}$}
\AXC{$\varphi_{135} \to \varphi_{202}$\lbl{A2}}
\BIC{$\varphi_{202}$}
\AXC{$\varphi_{202} \to \varphi_{213}$\lbl{A2}}
\BIC{$\varphi_{213}$}
\AXC{$\varphi_{213} \to \varphi_{129} \to \varphi_{205} \to \varphi_{171} \to \varphi_{147}$\lbl{A1}}
\BIC{$\varphi_{129} \to \varphi_{205} \to \varphi_{171} \to \varphi_{147}$}
\BIC{$\varphi_{205} \to \varphi_{171} \to \varphi_{147}$}
\end{ptree}
\noindent$\Pi_{66}$:
\begin{ptree}
\AXC{$\stackrel{\Pi_{13}}{\vdots}$}\noLine
\UIC{$\varphi_{123}$}
\AXC{$\stackrel{\Pi_{19}}{\vdots}$}\noLine
\UIC{$\varphi_{124}$}
\AXC{$\stackrel{\Pi_{20}}{\vdots}$}\noLine
\UIC{$\varphi_{124} \to \varphi_{123} \to \varphi_{171}$}
\BIC{$\varphi_{123} \to \varphi_{171}$}
\BIC{$\varphi_{171}$}
\AXC{$\varphi_{95} \to \forall n\,(\varphi_{95})$\lbl{A4}}
\AXC{$\stackrel{\Pi_{64}}{\vdots}$}\noLine
\UIC{$(\varphi_{95} \to \forall n\,(\varphi_{95})) \to \varphi_{205}$}
\BIC{$\varphi_{205}$}
\AXC{$\stackrel{\Pi_{65}}{\vdots}$}\noLine
\UIC{$\varphi_{205} \to \varphi_{171} \to \varphi_{147}$}
\BIC{$\varphi_{171} \to \varphi_{147}$}
\BIC{$\varphi_{147}$}
\AXC{$\varphi_{147} \to \varphi_{148}$\lbl{A1}}
\BIC{$\varphi_{148}$}
\UIC{$\forall a\,(\varphi_{148})$}
\end{ptree}
\noindent$\Pi_{67}$:
\begin{ptree}
\AXC{$\varphi_{20}$\lbl{A4}}
\AXC{$\forall a\,(\varphi_{148}) \to \varphi_{155}$\lbl{A3}}
\AXC{$\stackrel{\Pi_{66}}{\vdots}$}\noLine
\UIC{$\forall a\,(\varphi_{148})$}
\AXC{$\varphi_{276}$\lbl{A1}}
\BIC{$(\forall a\,(\varphi_{148}) \to \varphi_{155}) \to \varphi_{20} \to \mathrm{FLim}(f,c,\alpha) \to \varphi_{126}$}
\BIC{$\varphi_{20} \to \mathrm{FLim}(f,c,\alpha) \to \varphi_{126}$}
\BIC{$\mathrm{FLim}(f,c,\alpha) \to \varphi_{126}$}
\end{ptree}
\noindent$\Pi_{68}$:
\begin{ptree}
\AXC{$\varphi_{92}$\lbl{A2}}
\AXC{$\stackrel{\Pi_{8}}{\vdots}$}\noLine
\UIC{$\mathrm{FLim}(f,c,\alpha) \to \mathrm{FLim}(g,c,\beta) \to \varphi_{11}$}
\AXC{$\varphi_{243}$\lbl{A1}}
\BIC{$\varphi_{92} \to \mathrm{FLim}(f,c,\alpha) \to \varphi_{28}$}
\BIC{$\mathrm{FLim}(f,c,\alpha) \to \varphi_{28}$}
\end{ptree}
\noindent$\Pi_{69}$:
\begin{ptree}
\AXC{$\mathrm{FLim}(f,c,\alpha) \to \mathrm{FLim}(g,c,\beta) \to \varphi_{223}$\lbl{A1}}
\AXC{$\varphi_{271}$\lbl{A1}}
\BIC{$\mathrm{FLim}(f,c,\alpha) \to \mathrm{FLim}(g,c,\beta) \to \varphi_{233}$}
\AXC{$(\mathrm{FLim}(f,c,\alpha) \to \mathrm{FLim}(g,c,\beta) \to \varphi_{233}) \to \varphi_{265}$\lbl{A1}}
\BIC{$\varphi_{265}$}
\end{ptree}
\noindent$\Pi_{70}$:
\begin{ptree}
\AXC{$\forall x \in \mathbb{R}\,(\varphi_{29}) \to x \in \mathbb{R} \to \varphi_{29}$\lbl{A2}}
\AXC{$\stackrel{\Pi_{69}}{\vdots}$}\noLine
\UIC{$\varphi_{265}$}
\BIC{$\mathrm{FLim}(f,c,\alpha) \to \mathrm{FLim}(g,c,\beta) \to \varphi_{232}$}
\AXC{$(\mathrm{FLim}(f,c,\alpha) \to \mathrm{FLim}(g,c,\beta) \to \varphi_{232}) \to \varphi_{282}$\lbl{A1}}
\BIC{$\varphi_{282}$}
\end{ptree}
\noindent$\Pi_{71}$:
\begin{ptree}
\AXC{$\mathrm{FLim}(f,c,\alpha) \to \mathrm{FLim}(g,c,\beta) \to \varphi_{163}$\lbl{A1}}
\AXC{$\mathrm{FLim}(f,c,\alpha) \to \varphi_{146}$\lbl{A1}}
\AXC{$\stackrel{\Pi_{70}}{\vdots}$}\noLine
\UIC{$\varphi_{282}$}
\BIC{$(\mathrm{FLim}(f,c,\alpha) \to \mathrm{FLim}(g,c,\beta) \to \varphi_{163}) \to \varphi_{268}$}
\BIC{$\varphi_{268}$}
\end{ptree}
\noindent$\Pi_{72}$:
\begin{ptree}
\AXC{$\stackrel{\Pi_{68}}{\vdots}$}\noLine
\UIC{$\mathrm{FLim}(f,c,\alpha) \to \varphi_{28}$}
\AXC{$\mathrm{FLim}(f,c,\alpha) \to \mathrm{FLim}(g,c,\beta) \to \varphi_{189}$\lbl{A1}}
\AXC{$\stackrel{\Pi_{71}}{\vdots}$}\noLine
\UIC{$\varphi_{268}$}
\BIC{$(\mathrm{FLim}(f,c,\alpha) \to \varphi_{28}) \to \mathrm{FLim}(f,c,\alpha) \to \varphi_{203}$}
\BIC{$\mathrm{FLim}(f,c,\alpha) \to \varphi_{203}$}
\end{ptree}
\noindent$\Pi_{73}$:
\begin{ptree}
\AXC{$\stackrel{\Pi_{72}}{\vdots}$}\noLine
\UIC{$\mathrm{FLim}(f,c,\alpha) \to \varphi_{203}$}
\AXC{$(\mathrm{FLim}(f,c,\alpha) \to \varphi_{203}) \to \varphi_{127} \to \varphi_{40}$\lbl{A1}}
\BIC{$\varphi_{127} \to \varphi_{40}$}
\UIC{$\forall x\,(\varphi_{127} \to \varphi_{40})$}
\AXC{$\forall x\,(\varphi_{127} \to \varphi_{40}) \to \varphi_{285}$\lbl{A1}}
\BIC{$\varphi_{285}$}
\end{ptree}
\noindent$\Pi_{74}$:
\begin{ptree}
\AXC{$\varphi_{127} \to \forall x\,(\varphi_{127})$\lbl{A4}}
\AXC{$\forall x\,(\varphi_{127} \to \varphi_{40}) \to \forall x\,(\varphi_{127}) \to \varphi_{48}$\lbl{A3}}
\AXC{$\stackrel{\Pi_{73}}{\vdots}$}\noLine
\UIC{$\varphi_{285}$}
\BIC{$(\varphi_{127} \to \forall x\,(\varphi_{127})) \to \mathrm{FLim}(f,c,\alpha) \to \varphi_{207}$}
\BIC{$\mathrm{FLim}(f,c,\alpha) \to \varphi_{207}$}
\end{ptree}
\noindent$\Pi_{75}$:
\begin{ptree}
\AXC{$\stackrel{\Pi_{74}}{\vdots}$}\noLine
\UIC{$\mathrm{FLim}(f,c,\alpha) \to \varphi_{207}$}
\AXC{$\mathrm{FLim}(f,c,\alpha) \to \mathrm{FLim}(g,c,\beta) \to \varphi_{223}$\lbl{A1}}
\AXC{$\varphi_{277}$\lbl{A1}}
\BIC{$(\mathrm{FLim}(f,c,\alpha) \to \varphi_{207}) \to \mathrm{FLim}(f,c,\alpha) \to \varphi_{226}$}
\BIC{$\mathrm{FLim}(f,c,\alpha) \to \varphi_{226}$}
\end{ptree}
\noindent$\Pi_{76}$:
\begin{ptree}
\AXC{$\varphi_{104}$\lbl{DefPNbhd}}
\AXC{$\varphi_{104} \to \forall c\,(\mathrm{PNbhd}((f \cdot g),c) \leftrightarrow \varphi_{84})$\lbl{A2}}
\BIC{$\forall c\,(\mathrm{PNbhd}((f \cdot g),c) \leftrightarrow \varphi_{84})$}
\AXC{$\varphi_{236}$\lbl{A2}}
\BIC{$\mathrm{PNbhd}((f \cdot g),c) \leftrightarrow \varphi_{84}$}
\AXC{$(\mathrm{PNbhd}((f \cdot g),c) \leftrightarrow \varphi_{84}) \to \varphi_{259}$\lbl{A1}}
\BIC{$\varphi_{259}$}
\end{ptree}
\noindent$\Pi_{77}$:
\begin{ptree}
\AXC{$\varphi_{87} \to \neg (d \in \mathbb{R} \land 0 < d \land \varphi_{48})$\lbl{A2}}
\AXC{$\stackrel{\Pi_{75}}{\vdots}$}\noLine
\UIC{$\mathrm{FLim}(f,c,\alpha) \to \varphi_{226}$}
\AXC{$(\mathrm{FLim}(f,c,\alpha) \to \varphi_{226}) \to \varphi_{266}$\lbl{A1}}
\BIC{$\varphi_{266}$}
\BIC{$\mathrm{FLim}(f,c,\alpha) \to \varphi_{230}$}
\AXC{$\stackrel{\Pi_{76}}{\vdots}$}\noLine
\UIC{$\varphi_{259}$}
\BIC{$\mathrm{FLim}(f,c,\alpha) \to \mathrm{FLim}(g,c,\beta) \to \varphi_{117}$}
\end{ptree}
\noindent$\Pi_{78}$:
\begin{ptree}
\AXC{$\forall f\,\forall c\,\forall \alpha\,(\mathrm{FLim}(f,c,\alpha) \to \mathrm{PNbhd}(f,c))$\lbl{[D5]}}
\AXC{$\varphi_{34}$\lbl{A2}}
\BIC{$\forall c\,\forall \alpha\,(\mathrm{FLim}(g,c,\alpha) \to \mathrm{PNbhd}(g,c))$}
\AXC{$\varphi_{24}$\lbl{A2}}
\BIC{$\forall \alpha\,(\mathrm{FLim}(g,c,\alpha) \to \mathrm{PNbhd}(g,c))$}
\AXC{$\varphi_{15}$\lbl{A2}}
\BIC{$\mathrm{FLim}(g,c,\beta) \to \mathrm{PNbhd}(g,c)$}
\end{ptree}
\noindent$\Pi_{79}$:
\begin{ptree}
\AXC{$\varphi_{3}$\lbl{A1}}
\AXC{$\stackrel{\Pi_{78}}{\vdots}$}\noLine
\UIC{$\mathrm{FLim}(g,c,\beta) \to \mathrm{PNbhd}(g,c)$}
\AXC{$(\mathrm{FLim}(g,c,\beta) \to \mathrm{PNbhd}(g,c)) \to \varphi_{3} \to \varphi_{1}$\lbl{A1}}
\BIC{$\varphi_{3} \to \varphi_{1}$}
\BIC{$\varphi_{1}$}
\end{ptree}
\noindent$\Pi_{80}$:
\begin{ptree}
\AXC{$\forall f\,\forall c\,\forall \alpha\,(\mathrm{FLim}(f,c,\alpha) \to \mathrm{PNbhd}(f,c))$\lbl{[D5]}}
\AXC{$\varphi_{35}$\lbl{A2}}
\BIC{$\forall c\,\forall \alpha\,(\mathrm{FLim}(f,c,\alpha) \to \mathrm{PNbhd}(f,c))$}
\AXC{$\varphi_{25}$\lbl{A2}}
\BIC{$\forall \alpha\,(\mathrm{FLim}(f,c,\alpha) \to \mathrm{PNbhd}(f,c))$}
\AXC{$\varphi_{16}$\lbl{A2}}
\BIC{$\mathrm{FLim}(f,c,\alpha) \to \mathrm{PNbhd}(f,c)$}
\AXC{$\varphi_{80}$\lbl{A1}}
\BIC{$(\mathrm{FLim}(f,c,\alpha) \to \mathrm{FLim}(f,c,\alpha)) \to \varphi_{2}$}
\end{ptree}
\noindent$\Pi_{81}$:
\begin{ptree}
\AXC{$\mathrm{FLim}(f,c,\alpha) \to \mathrm{FLim}(f,c,\alpha)$\lbl{A1}}
\AXC{$\stackrel{\Pi_{80}}{\vdots}$}\noLine
\UIC{$(\mathrm{FLim}(f,c,\alpha) \to \mathrm{FLim}(f,c,\alpha)) \to \varphi_{2}$}
\BIC{$\varphi_{2}$}
\end{ptree}
\noindent$\Pi_{82}$:
\begin{ptree}
\AXC{$\forall f\,\forall g\,\forall c\,(\mathrm{PNbhd}(f,c) \to \varphi_{111})$\lbl{[H2]}}
\AXC{$\varphi_{245}$\lbl{A2}}
\BIC{$\forall g\,\forall c\,(\mathrm{PNbhd}(f,c) \to \varphi_{111})$}
\AXC{$\varphi_{244}$\lbl{A2}}
\BIC{$\forall c\,(\mathrm{PNbhd}(f,c) \to \varphi_{111})$}
\AXC{$\forall c\,(\mathrm{PNbhd}(f,c) \to \varphi_{111}) \to \mathrm{PNbhd}(f,c) \to \varphi_{111}$\lbl{A2}}
\BIC{$\mathrm{PNbhd}(f,c) \to \varphi_{111}$}
\end{ptree}
\noindent$\Pi_{83}$:
\begin{ptree}
\AXC{$\stackrel{\Pi_{79}}{\vdots}$}\noLine
\UIC{$\varphi_{1}$}
\AXC{$\stackrel{\Pi_{81}}{\vdots}$}\noLine
\UIC{$\varphi_{2}$}
\AXC{$\stackrel{\Pi_{82}}{\vdots}$}\noLine
\UIC{$\mathrm{PNbhd}(f,c) \to \varphi_{111}$}
\AXC{$\varphi_{255}$\lbl{A1}}
\BIC{$\varphi_{2} \to \varphi_{1} \to \mathrm{FLim}(f,c,\alpha) \to \varphi_{115}$}
\BIC{$\varphi_{1} \to \mathrm{FLim}(f,c,\alpha) \to \varphi_{115}$}
\BIC{$\mathrm{FLim}(f,c,\alpha) \to \varphi_{115}$}
\AXC{$(\mathrm{FLim}(f,c,\alpha) \to \varphi_{115}) \to \varphi_{274}$\lbl{A1}}
\BIC{$\varphi_{274}$}
\end{ptree}
\noindent$\Pi_{84}$:
\begin{ptree}
\AXC{$\varphi_{10} \to \forall d\,(\varphi_{10})$\lbl{A4}}
\AXC{$\varphi_{253}$\lbl{A3}}
\AXC{$\stackrel{\Pi_{77}}{\vdots}$}\noLine
\UIC{$\mathrm{FLim}(f,c,\alpha) \to \mathrm{FLim}(g,c,\beta) \to \varphi_{117}$}
\AXC{$\varphi_{250}$\lbl{A1}}
\BIC{$\varphi_{10} \to \neg (d \in \mathbb{R} \land 0 < d \land \forall x \in \mathbb{R}\,(\varphi_{29}))$}
\UIC{$\forall d\,(\varphi_{10} \to \neg (d \in \mathbb{R} \land 0 < d \land \forall x \in \mathbb{R}\,(\varphi_{29})))$}
\AXC{$\stackrel{\Pi_{83}}{\vdots}$}\noLine
\UIC{$\varphi_{274}$}
\BIC{$\varphi_{253} \to (\varphi_{10} \to \forall d\,(\varphi_{10})) \to \varphi_{7}$}
\BIC{$(\varphi_{10} \to \forall d\,(\varphi_{10})) \to \varphi_{7}$}
\BIC{$\varphi_{7}$}
\end{ptree}
\noindent$\Pi_{85}$:
\begin{ptree}
\AXC{$\varphi_{3}$\lbl{A1}}
\AXC{$\forall f\,\forall c\,\forall \alpha\,(\mathrm{FLim}(f,c,\alpha) \to \alpha \in \mathbb{R})$\lbl{[D4]}}
\AXC{$\varphi_{41}$\lbl{A2}}
\BIC{$\forall c\,\forall \alpha\,(\mathrm{FLim}(g,c,\alpha) \to \alpha \in \mathbb{R})$}
\AXC{$\varphi_{32}$\lbl{A2}}
\BIC{$\forall \alpha\,(\mathrm{FLim}(g,c,\alpha) \to \alpha \in \mathbb{R})$}
\AXC{$\varphi_{18}$\lbl{A2}}
\BIC{$\mathrm{FLim}(g,c,\beta) \to \beta \in \mathbb{R}$}
\AXC{$\varphi_{107}$\lbl{A1}}
\BIC{$\varphi_{3} \to \mathrm{FLim}(f,c,\alpha) \to \mathrm{FLim}(g,c,\beta) \to \beta \in \mathbb{R}$}
\BIC{$\mathrm{FLim}(f,c,\alpha) \to \mathrm{FLim}(g,c,\beta) \to \beta \in \mathbb{R}$}
\end{ptree}
\noindent$\Pi_{86}$:
\begin{ptree}
\AXC{$\forall f\,\forall c\,\forall \alpha\,(\mathrm{FLim}(f,c,\alpha) \to \alpha \in \mathbb{R})$\lbl{[D4]}}
\AXC{$\varphi_{42}$\lbl{A2}}
\BIC{$\forall c\,\forall \alpha\,(\mathrm{FLim}(f,c,\alpha) \to \alpha \in \mathbb{R})$}
\AXC{$\varphi_{33}$\lbl{A2}}
\BIC{$\forall \alpha\,(\mathrm{FLim}(f,c,\alpha) \to \alpha \in \mathbb{R})$}
\AXC{$\varphi_{19}$\lbl{A2}}
\BIC{$\mathrm{FLim}(f,c,\alpha) \to \alpha \in \mathbb{R}$}
\end{ptree}
\noindent$\Pi_{87}$:
\begin{ptree}
\AXC{$\mathrm{FLim}(f,c,\alpha) \to \mathrm{FLim}(f,c,\alpha)$\lbl{A1}}
\AXC{$\stackrel{\Pi_{86}}{\vdots}$}\noLine
\UIC{$\mathrm{FLim}(f,c,\alpha) \to \alpha \in \mathbb{R}$}
\AXC{$(\mathrm{FLim}(f,c,\alpha) \to \alpha \in \mathbb{R}) \to \varphi_{38}$\lbl{A1}}
\BIC{$\varphi_{38}$}
\BIC{$\mathrm{FLim}(f,c,\alpha) \to \mathrm{FLim}(g,c,\beta) \to \alpha \in \mathbb{R}$}
\end{ptree}
\noindent$\Pi_{88}$:
\begin{ptree}
\AXC{$\forall x\,\forall y\,(x \in \mathbb{R} \to y \in \mathbb{R} \to x \cdot y \in \mathbb{R})$\lbl{(R1)}}
\AXC{$\varphi_{78}$\lbl{A2}}
\BIC{$\forall y\,(\alpha \in \mathbb{R} \to y \in \mathbb{R} \to \alpha \cdot y \in \mathbb{R})$}
\AXC{$\varphi_{74}$\lbl{A2}}
\BIC{$\alpha \in \mathbb{R} \to \beta \in \mathbb{R} \to \alpha \cdot \beta \in \mathbb{R}$}
\AXC{$(\alpha \in \mathbb{R} \to \beta \in \mathbb{R} \to \alpha \cdot \beta \in \mathbb{R}) \to \varphi_{130}$\lbl{A1}}
\BIC{$\varphi_{130}$}
\end{ptree}
\noindent$\Pi_{89}$:
\begin{ptree}
\AXC{$\stackrel{\Pi_{85}}{\vdots}$}\noLine
\UIC{$\mathrm{FLim}(f,c,\alpha) \to \mathrm{FLim}(g,c,\beta) \to \beta \in \mathbb{R}$}
\AXC{$\stackrel{\Pi_{87}}{\vdots}$}\noLine
\UIC{$\mathrm{FLim}(f,c,\alpha) \to \mathrm{FLim}(g,c,\beta) \to \alpha \in \mathbb{R}$}
\AXC{$\stackrel{\Pi_{88}}{\vdots}$}\noLine
\UIC{$\varphi_{130}$}
\BIC{$\varphi_{69}$}
\BIC{$\mathrm{FLim}(f,c,\alpha) \to \mathrm{FLim}(g,c,\beta) \to \alpha \cdot \beta \in \mathbb{R}$}
\end{ptree}
\noindent$\Pi_{90}$:
\begin{ptree}
\AXC{$\forall f\,\forall c\,\forall \alpha\,(\mathrm{FLim}(f,c,\alpha) \to c \in \mathbb{R})$\lbl{[D3]}}
\AXC{$\varphi_{36}$\lbl{A2}}
\BIC{$\forall c\,\forall \alpha\,(\mathrm{FLim}(f,c,\alpha) \to c \in \mathbb{R})$}
\AXC{$\varphi_{26}$\lbl{A2}}
\BIC{$\forall \alpha\,(\mathrm{FLim}(f,c,\alpha) \to c \in \mathbb{R})$}
\AXC{$\varphi_{17}$\lbl{A2}}
\BIC{$\mathrm{FLim}(f,c,\alpha) \to c \in \mathbb{R}$}
\end{ptree}
\noindent$\Pi_{91}$:
\begin{ptree}
\AXC{$\mathrm{FLim}(f,c,\alpha) \to \mathrm{FLim}(f,c,\alpha)$\lbl{A1}}
\AXC{$\stackrel{\Pi_{90}}{\vdots}$}\noLine
\UIC{$\mathrm{FLim}(f,c,\alpha) \to c \in \mathbb{R}$}
\AXC{$(\mathrm{FLim}(f,c,\alpha) \to c \in \mathbb{R}) \to \varphi_{37}$\lbl{A1}}
\BIC{$\varphi_{37}$}
\BIC{$\mathrm{FLim}(f,c,\alpha) \to \mathrm{FLim}(g,c,\beta) \to c \in \mathbb{R}$}
\end{ptree}
\noindent$\Pi_{92}$:
\begin{ptree}
\AXC{$\stackrel{\Pi_{57}}{\vdots}$}\noLine
\UIC{$\varphi_{4}$}
\AXC{$\varphi_{185}$\lbl{定理1.7(⇐)}}
\AXC{$\varphi_{185} \to \forall c\,(\varphi_{219})$\lbl{A2}}
\BIC{$\forall c\,(\varphi_{219})$}
\AXC{$\forall c\,(\varphi_{219}) \to \varphi_{219}$\lbl{A2}}
\BIC{$\varphi_{219}$}
\AXC{$\varphi_{219} \to \varphi_{228}$\lbl{A2}}
\BIC{$\varphi_{228}$}
\AXC{$\varphi_{228} \to \varphi_{4} \to \varphi_{254}$\lbl{A1}}
\BIC{$\varphi_{4} \to \varphi_{254}$}
\BIC{$\varphi_{254}$}
\end{ptree}
\noindent$\Pi_{93}$:
\begin{ptree}
\AXC{$\stackrel{\Pi_{89}}{\vdots}$}\noLine
\UIC{$\mathrm{FLim}(f,c,\alpha) \to \mathrm{FLim}(g,c,\beta) \to \alpha \cdot \beta \in \mathbb{R}$}
\AXC{$\stackrel{\Pi_{91}}{\vdots}$}\noLine
\UIC{$\mathrm{FLim}(f,c,\alpha) \to \mathrm{FLim}(g,c,\beta) \to c \in \mathbb{R}$}
\AXC{$\stackrel{\Pi_{92}}{\vdots}$}\noLine
\UIC{$\varphi_{254}$}
\BIC{$\varphi_{248}$}
\BIC{$\varphi_{7} \to (\mathrm{FLim}(f,c,\alpha) \to \varphi_{126}) \to \varphi_{14}$}
\end{ptree}
\noindent 主木:
\begin{ptree}
\AXC{$\stackrel{\Pi_{67}}{\vdots}$}\noLine
\UIC{$\mathrm{FLim}(f,c,\alpha) \to \varphi_{126}$}
\AXC{$\stackrel{\Pi_{84}}{\vdots}$}\noLine
\UIC{$\varphi_{7}$}
\AXC{$\stackrel{\Pi_{93}}{\vdots}$}\noLine
\UIC{$\varphi_{7} \to (\mathrm{FLim}(f,c,\alpha) \to \varphi_{126}) \to \varphi_{14}$}
\BIC{$(\mathrm{FLim}(f,c,\alpha) \to \varphi_{126}) \to \varphi_{14}$}
\BIC{$\varphi_{14}$}
\UIC{$\forall \beta\,(\varphi_{14})$}
\UIC{$\forall \alpha\,\forall \beta\,(\varphi_{14})$}
\UIC{$\forall c\,\forall \alpha\,\forall \beta\,(\varphi_{14})$}
\UIC{$\forall g\,\forall c\,\forall \alpha\,\forall \beta\,(\varphi_{14})$}
\UIC{$\forall f\,\forall g\,\forall c\,\forall \alpha\,\forall \beta\,(\varphi_{14})$}
\end{ptree}
\begin{abbrevlist}
\item $\varphi_{1} :\equiv \mathrm{FLim}(f,c,\alpha) \to \mathrm{FLim}(g,c,\beta) \to \mathrm{PNbhd}(g,c)$
\item $\varphi_{2} :\equiv \mathrm{FLim}(f,c,\alpha) \to \mathrm{FLim}(g,c,\beta) \to \mathrm{PNbhd}(f,c)$
\item $\varphi_{3} :\equiv \mathrm{FLim}(f,c,\alpha) \to \mathrm{FLim}(g,c,\beta) \to \mathrm{FLim}(g,c,\beta)$
\item $\varphi_{4} :\equiv \mathrm{FLim}(f,c,\alpha) \to \mathrm{FLim}(g,c,\beta) \to \mathrm{RFn}((f \cdot g))$
\item $\varphi_{5} :\equiv \mathrm{FLim}(f,c,\alpha) \to \mathrm{FLim}(g,c,\beta) \to \mathrm{Seq}(a) \to \mathrm{Seq}(a)$
\item $\varphi_{6} :\equiv \mathrm{FLim}(g,c,\beta) \to \mathrm{Seq}(a) \to \mathrm{Conv}(a,c) \to \mathrm{Conv}(a,c)$
\item $\varphi_{7} :\equiv \mathrm{FLim}(f,c,\alpha) \to \mathrm{FLim}(g,c,\beta) \to \mathrm{PNbhd}((f \cdot g),c)$
\item $\varphi_{8} :\equiv \forall \alpha\,(\mathrm{FLim}(g,c,\alpha) \to \mathrm{RFn}(g)) \to \mathrm{FLim}(g,c,\beta) \to \mathrm{RFn}(g)$
\item $\varphi_{9} :\equiv \forall \alpha\,(\mathrm{FLim}(f,c,\alpha) \to \mathrm{RFn}(f)) \to \mathrm{FLim}(f,c,\alpha) \to \mathrm{RFn}(f)$
\item $\varphi_{10} :\equiv \mathrm{FLim}(f,c,\alpha) \land \mathrm{FLim}(g,c,\beta) \land \neg \mathrm{PNbhd}((f \cdot g),c)$
\item $\varphi_{11} :\equiv \forall x\,(x \in \mathrm{dom}((f \cdot g)) \leftrightarrow x \in \mathrm{dom}(f) \land x \in \mathrm{dom}(g))$
\item $\varphi_{12} :\equiv \forall c\,\forall \alpha\,(\mathrm{FLim}(g,c,\alpha) \to \mathrm{RFn}(g)) \to \forall \alpha\,(\mathrm{FLim}(g,c,\alpha) \to \mathrm{RFn}(g))$
\item $\varphi_{13} :\equiv \forall c\,\forall \alpha\,(\mathrm{FLim}(f,c,\alpha) \to \mathrm{RFn}(f)) \to \forall \alpha\,(\mathrm{FLim}(f,c,\alpha) \to \mathrm{RFn}(f))$
\item $\varphi_{14} :\equiv \mathrm{FLim}(f,c,\alpha) \to \mathrm{FLim}(g,c,\beta) \to \mathrm{FLim}((f \cdot g),c,\alpha \cdot \beta)$
\item $\varphi_{15} :\equiv \forall \alpha\,(\mathrm{FLim}(g,c,\alpha) \to \mathrm{PNbhd}(g,c)) \to \mathrm{FLim}(g,c,\beta) \to \mathrm{PNbhd}(g,c)$
\item $\varphi_{16} :\equiv \forall \alpha\,(\mathrm{FLim}(f,c,\alpha) \to \mathrm{PNbhd}(f,c)) \to \mathrm{FLim}(f,c,\alpha) \to \mathrm{PNbhd}(f,c)$
\item $\varphi_{17} :\equiv \forall \alpha\,(\mathrm{FLim}(f,c,\alpha) \to c \in \mathbb{R}) \to \mathrm{FLim}(f,c,\alpha) \to c \in \mathbb{R}$
\item $\varphi_{18} :\equiv \forall \alpha\,(\mathrm{FLim}(g,c,\alpha) \to \alpha \in \mathbb{R}) \to \mathrm{FLim}(g,c,\beta) \to \beta \in \mathbb{R}$
\item $\varphi_{19} :\equiv \forall \alpha\,(\mathrm{FLim}(f,c,\alpha) \to \alpha \in \mathbb{R}) \to \mathrm{FLim}(f,c,\alpha) \to \alpha \in \mathbb{R}$
\item $\varphi_{20} :\equiv \mathrm{FLim}(f,c,\alpha) \land \mathrm{FLim}(g,c,\beta) \to \forall a\,(\mathrm{FLim}(f,c,\alpha) \land \mathrm{FLim}(g,c,\beta))$
\item $\varphi_{21} :\equiv \forall f\,\forall c\,\forall \alpha\,(\mathrm{FLim}(f,c,\alpha) \to \mathrm{RFn}(f)) \to \forall c\,\forall \alpha\,(\mathrm{FLim}(g,c,\alpha) \to \mathrm{RFn}(g))$
\item $\varphi_{22} :\equiv \forall f\,\forall c\,\forall \alpha\,(\mathrm{FLim}(f,c,\alpha) \to \mathrm{RFn}(f)) \to \forall c\,\forall \alpha\,(\mathrm{FLim}(f,c,\alpha) \to \mathrm{RFn}(f))$
\item $\varphi_{23} :\equiv \forall x \in \mathrm{dom}((f \cdot g))\,((f \cdot g)(x) = f(x) \cdot g(x))$
\item $\varphi_{24} :\equiv \forall c\,\forall \alpha\,(\mathrm{FLim}(g,c,\alpha) \to \mathrm{PNbhd}(g,c)) \to \forall \alpha\,(\mathrm{FLim}(g,c,\alpha) \to \mathrm{PNbhd}(g,c))$
\item $\varphi_{25} :\equiv \forall c\,\forall \alpha\,(\mathrm{FLim}(f,c,\alpha) \to \mathrm{PNbhd}(f,c)) \to \forall \alpha\,(\mathrm{FLim}(f,c,\alpha) \to \mathrm{PNbhd}(f,c))$
\item $\varphi_{26} :\equiv \forall c\,\forall \alpha\,(\mathrm{FLim}(f,c,\alpha) \to c \in \mathbb{R}) \to \forall \alpha\,(\mathrm{FLim}(f,c,\alpha) \to c \in \mathbb{R})$
\item $\varphi_{27} :\equiv \forall y\,\forall z\,((f \cdot g)(a_{n}) = y \to y = z \to (f \cdot g)(a_{n}) = z)$
\item $\varphi_{28} :\equiv \mathrm{FLim}(g,c,\beta) \to (x \in \mathrm{dom}((f \cdot g)) \leftrightarrow x \in \mathrm{dom}(f) \land x \in \mathrm{dom}(g))$
\item $\varphi_{29} :\equiv 0 < \lvert x - c \rvert \to \lvert x - c \rvert < d \to x \in \mathrm{dom}(f) \land x \in \mathrm{dom}(g)$
\item $\varphi_{30} :\equiv (\mathrm{FLim}(f,c,\alpha) \to \mathrm{FLim}(f,c,\alpha)) \to \mathrm{FLim}(f,c,\alpha) \to \mathrm{FLim}(g,c,\beta) \to \mathrm{RFn}(f)$
\item $\varphi_{31} :\equiv \forall a\,\forall \alpha\,\forall b\,\forall \beta\,(\mathrm{Conv}(a,\alpha) \to \mathrm{Conv}(b,\beta) \to \mathrm{Conv}(a \otimes b,\alpha \cdot \beta))$
\item $\varphi_{32} :\equiv \forall c\,\forall \alpha\,(\mathrm{FLim}(g,c,\alpha) \to \alpha \in \mathbb{R}) \to \forall \alpha\,(\mathrm{FLim}(g,c,\alpha) \to \alpha \in \mathbb{R})$
\item $\varphi_{33} :\equiv \forall c\,\forall \alpha\,(\mathrm{FLim}(f,c,\alpha) \to \alpha \in \mathbb{R}) \to \forall \alpha\,(\mathrm{FLim}(f,c,\alpha) \to \alpha \in \mathbb{R})$
\item $\varphi_{34} :\equiv \forall f\,\forall c\,\forall \alpha\,(\mathrm{FLim}(f,c,\alpha) \to \mathrm{PNbhd}(f,c)) \to \forall c\,\forall \alpha\,(\mathrm{FLim}(g,c,\alpha) \to \mathrm{PNbhd}(g,c))$
\item $\varphi_{35} :\equiv \forall f\,\forall c\,\forall \alpha\,(\mathrm{FLim}(f,c,\alpha) \to \mathrm{PNbhd}(f,c)) \to \forall c\,\forall \alpha\,(\mathrm{FLim}(f,c,\alpha) \to \mathrm{PNbhd}(f,c))$
\item $\varphi_{36} :\equiv \forall f\,\forall c\,\forall \alpha\,(\mathrm{FLim}(f,c,\alpha) \to c \in \mathbb{R}) \to \forall c\,\forall \alpha\,(\mathrm{FLim}(f,c,\alpha) \to c \in \mathbb{R})$
\item $\varphi_{37} :\equiv (\mathrm{FLim}(f,c,\alpha) \to \mathrm{FLim}(f,c,\alpha)) \to \mathrm{FLim}(f,c,\alpha) \to \mathrm{FLim}(g,c,\beta) \to c \in \mathbb{R}$
\item $\varphi_{38} :\equiv (\mathrm{FLim}(f,c,\alpha) \to \mathrm{FLim}(f,c,\alpha)) \to \mathrm{FLim}(f,c,\alpha) \to \mathrm{FLim}(g,c,\beta) \to \alpha \in \mathbb{R}$
\item $\varphi_{39} :\equiv a_{n} \in \mathrm{dom}((f \cdot g)) \leftrightarrow a_{n} \in \mathrm{dom}(f) \land a_{n} \in \mathrm{dom}(g)$
\item $\varphi_{40} :\equiv x \in \mathbb{R} \to 0 < \lvert x - c \rvert \to \lvert x - c \rvert < d \to x \in \mathrm{dom}((f \cdot g))$
\item $\varphi_{41} :\equiv \forall f\,\forall c\,\forall \alpha\,(\mathrm{FLim}(f,c,\alpha) \to \alpha \in \mathbb{R}) \to \forall c\,\forall \alpha\,(\mathrm{FLim}(g,c,\alpha) \to \alpha \in \mathbb{R})$
\item $\varphi_{42} :\equiv \forall f\,\forall c\,\forall \alpha\,(\mathrm{FLim}(f,c,\alpha) \to \alpha \in \mathbb{R}) \to \forall c\,\forall \alpha\,(\mathrm{FLim}(f,c,\alpha) \to \alpha \in \mathbb{R})$
\item $\varphi_{43} :\equiv \forall n \in \mathbb{N}\,(a_{n} \in \mathrm{dom}(g) \land a_{n} \neq c) \to \mathrm{Conv}(g \circ a,\beta)$
\item $\varphi_{44} :\equiv \forall n \in \mathbb{N}\,(a_{n} \in \mathrm{dom}(f) \land a_{n} \neq c) \to \mathrm{Conv}(f \circ a,\alpha)$
\item $\varphi_{45} :\equiv \forall y\,\forall z\,((f \cdot g) \circ a(n) = y \to y = z \to (f \cdot g) \circ a(n) = z)$
\item $\varphi_{46} :\equiv \forall a\,\forall \alpha\,(\mathrm{Conv}(a,\alpha) \to \mathrm{Seq}(a)) \to \forall \alpha\,(\mathrm{Conv}(g \circ a,\alpha) \to \mathrm{Seq}(g \circ a))$
\item $\varphi_{47} :\equiv \forall a\,\forall \alpha\,(\mathrm{Conv}(a,\alpha) \to \mathrm{Seq}(a)) \to \forall \alpha\,(\mathrm{Conv}(f \circ a,\alpha) \to \mathrm{Seq}(f \circ a))$
\item $\varphi_{48} :\equiv \forall x \in \mathbb{R}\,(0 < \lvert x - c \rvert \to \lvert x - c \rvert < d \to x \in \mathrm{dom}((f \cdot g)))$
\item $\varphi_{49} :\equiv \forall n \in \mathbb{N}\,(a_{n} \in \mathrm{dom}((f \cdot g)) \land a_{n} \neq c) \to \mathrm{Seq}(g \circ a)$
\item $\varphi_{50} :\equiv \forall n \in \mathbb{N}\,(a_{n} \in \mathrm{dom}((f \cdot g)) \land a_{n} \neq c) \to \mathrm{Seq}(f \circ a)$
\item $\varphi_{51} :\equiv \forall a\,\forall b\,\forall \alpha\,(\mathrm{Seq}(b) \to \forall n \in \mathbb{N}\,(a_{n} = b_{n}) \to \mathrm{Conv}(a,\alpha) \to \mathrm{Conv}(b,\alpha))$
\item $\varphi_{52} :\equiv n \in \mathbb{N} \to \forall n \in \mathbb{N}\,((f \cdot g) \circ a(n) = (f \cdot g)(a_{n}))$
\item $\varphi_{53} :\equiv \forall n \in \mathbb{N}\,(a_{n} \in \mathrm{dom}((f \cdot g)) \land a_{n} \neq c) \to n \in \mathbb{N} \to n \in \mathbb{N}$
\item $\varphi_{54} :\equiv \forall n \in \mathbb{N}\,(a_{n} \in \mathrm{dom}((f \cdot g)) \land a_{n} \neq c) \to \mathrm{Conv}(g \circ a,\beta)$
\item $\varphi_{55} :\equiv \forall n \in \mathbb{N}\,(a_{n} \in \mathrm{dom}((f \cdot g)) \land a_{n} \neq c) \to \mathrm{Conv}(f \circ a,\alpha)$
\item $\varphi_{56} :\equiv \varphi_{3} \to \mathrm{FLim}(f,c,\alpha) \to \mathrm{FLim}(g,c,\beta) \to \mathrm{RFn}(g)$
\item $\varphi_{57} :\equiv \forall b\,\forall \beta\,(\mathrm{Conv}(f \circ a,\alpha) \to \mathrm{Conv}(b,\beta) \to \mathrm{Conv}(f \circ a \otimes b,\alpha \cdot \beta))$
\item $\varphi_{58} :\equiv \forall x\,\forall y\,(x = y \to y = x) \to \forall y\,((f \cdot g) \circ a(n) = y \to y = (f \cdot g) \circ a(n))$
\item $\varphi_{59} :\equiv n \in \mathbb{N} \to \mathrm{Seq}(g \circ a) \land \forall n \in \mathbb{N}\,(g \circ a(n) = g(a_{n}))$
\item $\varphi_{60} :\equiv n \in \mathbb{N} \to \mathrm{Seq}(f \circ a) \land \forall n \in \mathbb{N}\,(f \circ a(n) = f(a_{n}))$
\item $\varphi_{61} :\equiv \forall n \in \mathbb{N}\,(a_{n} \in \mathrm{dom}((f \cdot g)) \land a_{n} \neq c) \to \mathrm{Seq}((f \cdot g) \circ a)$
\item $\varphi_{62} :\equiv \forall x\,(x = x) \to f(a_{n}) \cdot g \circ a(n) = f(a_{n}) \cdot g \circ a(n)$
\item $\varphi_{63} :\equiv \mathrm{Conv}(f \circ a \otimes g \circ a,\alpha \cdot \beta) \to \mathrm{Conv}((f \cdot g) \circ a,\alpha \cdot \beta)$
\item $\varphi_{64} :\equiv a_{n} \in \mathrm{dom}((f \cdot g)) \to (f \cdot g)(a_{n}) = f(a_{n}) \cdot g(a_{n})$
\item $\varphi_{65} :\equiv \forall \alpha\,(\mathrm{Conv}(g \circ a,\alpha) \to \mathrm{Seq}(g \circ a)) \to \mathrm{Conv}(g \circ a,\beta) \to \mathrm{Seq}(g \circ a)$
\item $\varphi_{66} :\equiv \forall \alpha\,(\mathrm{Conv}(f \circ a,\alpha) \to \mathrm{Seq}(f \circ a)) \to \mathrm{Conv}(f \circ a,\alpha) \to \mathrm{Seq}(f \circ a)$
\item $\varphi_{67} :\equiv \forall n \in \mathbb{N}\,(a_{n} \in \mathrm{dom}((f \cdot g)) \land a_{n} \neq c) \to n \in \mathbb{N} \to a_{n} \in \mathrm{dom}(g)$
\item $\varphi_{68} :\equiv \forall n \in \mathbb{N}\,(a_{n} \in \mathrm{dom}((f \cdot g)) \land a_{n} \neq c) \to n \in \mathbb{N} \to a_{n} \in \mathrm{dom}(f)$
\item $\varphi_{69} :\equiv (\mathrm{FLim}(f,c,\alpha) \to \mathrm{FLim}(g,c,\beta) \to \beta \in \mathbb{R}) \to \mathrm{FLim}(f,c,\alpha) \to \mathrm{FLim}(g,c,\beta) \to \alpha \cdot \beta \in \mathbb{R}$
\item $\varphi_{70} :\equiv \forall n \in \mathbb{N}\,(a_{n} \in \mathrm{dom}((f \cdot g)) \land a_{n} \neq c) \to \forall n \in \mathbb{N}\,(a_{n} \in \mathrm{dom}(g))$
\item $\varphi_{71} :\equiv \forall n \in \mathbb{N}\,(a_{n} \in \mathrm{dom}((f \cdot g)) \land a_{n} \neq c) \to \forall n \in \mathbb{N}\,(a_{n} \in \mathrm{dom}(f))$
\item $\varphi_{72} :\equiv n \in \mathbb{N} \to f \circ a(n) \cdot g \circ a(n) = f(a_{n}) \cdot g \circ a(n)$
\item $\varphi_{73} :\equiv \forall x\,(x = x) \to f \circ a(n) \cdot g \circ a(n) = f \circ a(n) \cdot g \circ a(n)$
\item $\varphi_{74} :\equiv \forall y\,(\alpha \in \mathbb{R} \to y \in \mathbb{R} \to \alpha \cdot y \in \mathbb{R}) \to \alpha \in \mathbb{R} \to \beta \in \mathbb{R} \to \alpha \cdot \beta \in \mathbb{R}$
\item $\varphi_{75} :\equiv \mathrm{Seq}((f \cdot g) \circ a) \land \forall n \in \mathbb{N}\,((f \cdot g) \circ a(n) = (f \cdot g)(a_{n}))$
\item $\varphi_{76} :\equiv \forall n \in \mathbb{N}\,(a_{n} \in \mathrm{dom}((f \cdot g)) \land a_{n} \neq c) \to \mathrm{Conv}((f \cdot g) \circ a,\alpha \cdot \beta)$
\item $\varphi_{77} :\equiv \mathrm{Conv}(f \circ a,\alpha) \to \mathrm{Conv}(g \circ a,\beta) \to \mathrm{Conv}(f \circ a \otimes g \circ a,\alpha \cdot \beta)$
\item $\varphi_{78} :\equiv \forall x\,\forall y\,(x \in \mathbb{R} \to y \in \mathbb{R} \to x \cdot y \in \mathbb{R}) \to \forall y\,(\alpha \in \mathbb{R} \to y \in \mathbb{R} \to \alpha \cdot y \in \mathbb{R})$
\item $\varphi_{79} :\equiv \forall a\,\forall b\,(\mathrm{Seq}(a) \to \mathrm{Seq}(b) \to \mathrm{Seq}(a \otimes b) \land \forall n \in \mathbb{N}\,(a \otimes b(n) = a_{n} \cdot b_{n}))$
\item $\varphi_{80} :\equiv (\mathrm{FLim}(f,c,\alpha) \to \mathrm{PNbhd}(f,c)) \to (\mathrm{FLim}(f,c,\alpha) \to \mathrm{FLim}(f,c,\alpha)) \to \varphi_{2}$
\item $\varphi_{81} :\equiv \forall n \in \mathbb{N}\,(f \circ a \otimes g \circ a(n) = f \circ a(n) \cdot g \circ a(n))$
\item $\varphi_{82} :\equiv \forall n \in \mathbb{N}\,(a_{n} \in \mathrm{dom}((f \cdot g)) \land a_{n} \neq c) \to n \in \mathbb{N} \to a_{n} \in \mathrm{dom}((f \cdot g))$
\item $\varphi_{83} :\equiv \forall y\,\forall z\,(f \circ a \otimes g \circ a(n) = y \to y = z \to f \circ a \otimes g \circ a(n) = z)$
\item $\varphi_{84} :\equiv \exists r\,(r \in \mathbb{R} \land 0 < r \land \forall x \in \mathbb{R}\,(0 < \lvert x - c \rvert \to \lvert x - c \rvert < r \to x \in \mathrm{dom}((f \cdot g))))$
\item $\varphi_{85} :\equiv \forall n \in \mathbb{N}\,(g \circ a(n) = g(a_{n})) \to n \in \mathbb{N} \to g \circ a(n) = g(a_{n})$
\item $\varphi_{86} :\equiv \forall n \in \mathbb{N}\,(f \circ a(n) = f(a_{n})) \to n \in \mathbb{N} \to f \circ a(n) = f(a_{n})$
\item $\varphi_{87} :\equiv \forall r\,(\neg (r \in \mathbb{R} \land 0 < r \land \forall x \in \mathbb{R}\,(0 < \lvert x - c \rvert \to \lvert x - c \rvert < r \to x \in \mathrm{dom}((f \cdot g)))))$
\item $\varphi_{88} :\equiv \forall n \in \mathbb{N}\,(a_{n} \in \mathrm{dom}((f \cdot g)) \land a_{n} \neq c) \to \forall n \in \mathbb{N}\,(a_{n} \in \mathrm{dom}((f \cdot g)))$
\item $\varphi_{89} :\equiv \forall n \in \mathbb{N}\,(a_{n} \in \mathrm{dom}((f \cdot g)) \land a_{n} \neq c) \to n \in \mathbb{N} \to a_{n} \in \mathrm{dom}(g) \land a_{n} \neq c$
\item $\varphi_{90} :\equiv \forall n \in \mathbb{N}\,(a_{n} \in \mathrm{dom}((f \cdot g)) \land a_{n} \neq c) \to n \in \mathbb{N} \to a_{n} \in \mathrm{dom}(f) \land a_{n} \neq c$
\item $\varphi_{91} :\equiv \forall \alpha\,\forall a\,(\mathrm{FLim}(g,c,\alpha) \to \mathrm{Conv}(a,c) \to \forall n \in \mathbb{N}\,(a_{n} \in \mathrm{dom}(g) \land a_{n} \neq c) \to \mathrm{Conv}(g \circ a,\alpha))$
\item $\varphi_{92} :\equiv \varphi_{11} \to (x \in \mathrm{dom}((f \cdot g)) \leftrightarrow x \in \mathrm{dom}(f) \land x \in \mathrm{dom}(g))$
\item $\varphi_{93} :\equiv \forall n \in \mathbb{N}\,(a_{n} \in \mathrm{dom}((f \cdot g)) \land a_{n} \neq c) \to n \in \mathbb{N} \to g \circ a(n) = g(a_{n})$
\item $\varphi_{94} :\equiv \forall n \in \mathbb{N}\,(a_{n} \in \mathrm{dom}((f \cdot g)) \land a_{n} \neq c) \to n \in \mathbb{N} \to f \circ a(n) = f(a_{n})$
\item $\varphi_{95} :\equiv \mathrm{FLim}(f,c,\alpha) \land \mathrm{FLim}(g,c,\beta) \land \mathrm{Seq}(a) \land \mathrm{Conv}(a,c) \land \forall n \in \mathbb{N}\,(a_{n} \in \mathrm{dom}((f \cdot g)) \land a_{n} \neq c)$
\item $\varphi_{96} :\equiv \forall n \in \mathbb{N}\,(a_{n} \in \mathrm{dom}((f \cdot g)) \land a_{n} \neq c) \to \forall n \in \mathbb{N}\,(a_{n} \in \mathrm{dom}(g) \land a_{n} \neq c)$
\item $\varphi_{97} :\equiv \forall n \in \mathbb{N}\,(a_{n} \in \mathrm{dom}((f \cdot g)) \land a_{n} \neq c) \to \forall n \in \mathbb{N}\,(a_{n} \in \mathrm{dom}(f) \land a_{n} \neq c)$
\item $\varphi_{98} :\equiv \forall n \in \mathbb{N}\,(a_{n} \in \mathrm{dom}((f \cdot g)) \land a_{n} \neq c) \to \mathrm{Conv}(f \circ a \otimes g \circ a,\alpha \cdot \beta)$
\item $\varphi_{99} :\equiv \forall n \in \mathbb{N}\,(a_{n} \in \mathrm{dom}(g)) \to \mathrm{Seq}(g \circ a) \land \forall n \in \mathbb{N}\,(g \circ a(n) = g(a_{n}))$
\item $\varphi_{100} :\equiv \forall n \in \mathbb{N}\,(a_{n} \in \mathrm{dom}(f)) \to \mathrm{Seq}(f \circ a) \land \forall n \in \mathbb{N}\,(f \circ a(n) = f(a_{n}))$
\item $\varphi_{101} :\equiv n \in \mathbb{N} \to n \in \mathbb{N} \to f \circ a \otimes g \circ a(n) = f \circ a(n) \cdot g \circ a(n)$
\item $\varphi_{102} :\equiv \forall x\,\forall y\,(x = y \to y = x) \to \forall y\,(f \circ a \otimes g \circ a(n) = y \to y = f \circ a \otimes g \circ a(n))$
\item $\varphi_{103} :\equiv \forall c\,\forall \alpha\,\forall a\,(\mathrm{FLim}(f,c,\alpha) \to \mathrm{Conv}(a,c) \to \varphi_{44})$
\item $\varphi_{104} :\equiv \forall f\,\forall c\,(\mathrm{PNbhd}(f,c) \leftrightarrow \exists r\,(r \in \mathbb{R} \land 0 < r \land \forall x \in \mathbb{R}\,(0 < \lvert x - c \rvert \to \lvert x - c \rvert < r \to x \in \mathrm{dom}(f))))$
\item $\varphi_{105} :\equiv \forall n \in \mathbb{N}\,(a_{n} \in \mathrm{dom}((f \cdot g)) \land a_{n} \neq c) \to n \in \mathbb{N} \to a_{n} \in \mathrm{dom}((f \cdot g)) \land a_{n} \neq c$
\item $\varphi_{106} :\equiv \forall z\,((f \cdot g) \circ a(n) = (f \cdot g)(a_{n}) \to (f \cdot g)(a_{n}) = z \to (f \cdot g) \circ a(n) = z)$
\item $\varphi_{107} :\equiv (\mathrm{FLim}(g,c,\beta) \to \beta \in \mathbb{R}) \to \varphi_{3} \to \mathrm{FLim}(f,c,\alpha) \to \mathrm{FLim}(g,c,\beta) \to \beta \in \mathbb{R}$
\item $\varphi_{108} :\equiv \forall n \in \mathbb{N}\,(a_{n} \in \mathrm{dom}((f \cdot g)) \land a_{n} \neq c) \to \forall n \in \mathbb{N}\,(a_{n} \in \mathrm{dom}((f \cdot g)) \land a_{n} \neq c)$
\item $\varphi_{109} :\equiv \forall n \in \mathbb{N}\,(a_{n} \in \mathrm{dom}((f \cdot g)) \land a_{n} \neq c) \to n \in \mathbb{N} \to n \in \mathbb{N} \to g \circ a(n) = g(a_{n})$
\item $\varphi_{110} :\equiv \forall n \in \mathbb{N}\,(a_{n} \in \mathrm{dom}((f \cdot g)) \land a_{n} \neq c) \to n \in \mathbb{N} \to n \in \mathbb{N} \to f \circ a(n) = f(a_{n})$
\item $\varphi_{111} :\equiv \mathrm{PNbhd}(g,c) \to \exists d\,(d \in \mathbb{R} \land 0 < d \land \forall x \in \mathbb{R}\,(\varphi_{29}))$
\item $\varphi_{112} :\equiv f(a_{n}) \cdot g \circ a(n) = f(a_{n}) \cdot z \to f(a_{n}) \cdot g \circ a(n) = f(a_{n}) \cdot w$
\item $\varphi_{113} :\equiv \forall n \in \mathbb{N}\,(a_{n} \in \mathrm{dom}((f \cdot g)) \land a_{n} \neq c) \to n \in \mathbb{N} \to \forall n \in \mathbb{N}\,(g \circ a(n) = g(a_{n}))$
\item $\varphi_{114} :\equiv \forall n \in \mathbb{N}\,(a_{n} \in \mathrm{dom}((f \cdot g)) \land a_{n} \neq c) \to n \in \mathbb{N} \to \forall n \in \mathbb{N}\,(f \circ a(n) = f(a_{n}))$
\item $\varphi_{115} :\equiv \mathrm{FLim}(g,c,\beta) \to \exists d\,(d \in \mathbb{R} \land 0 < d \land \forall x \in \mathbb{R}\,(\varphi_{29}))$
\item $\varphi_{116} :\equiv \forall n \in \mathbb{N}\,(a_{n} \in \mathrm{dom}((f \cdot g)) \land a_{n} \neq c) \to n \in \mathbb{N} \to (f \cdot g) \circ a(n) = (f \cdot g)(a_{n})$
\item $\varphi_{117} :\equiv d \in \mathbb{R} \land 0 < d \land \forall x \in \mathbb{R}\,(\varphi_{29}) \to \mathrm{PNbhd}((f \cdot g),c)$
\item $\varphi_{118} :\equiv \mathrm{FLim}(f,c,\alpha) \to \mathrm{FLim}(g,c,\beta) \to \mathrm{Seq}(a) \to \mathrm{Conv}(a,c) \to \varphi_{49}$
\item $\varphi_{119} :\equiv \mathrm{FLim}(f,c,\alpha) \to \mathrm{FLim}(g,c,\beta) \to \mathrm{Seq}(a) \to \mathrm{Conv}(a,c) \to \varphi_{50}$
\item $\varphi_{120} :\equiv \forall z\,((f \cdot g)(a_{n}) = f(a_{n}) \cdot g(a_{n}) \to f(a_{n}) \cdot g(a_{n}) = z \to (f \cdot g)(a_{n}) = z)$
\item $\varphi_{121} :\equiv \forall n \in \mathbb{N}\,(a_{n} \in \mathrm{dom}((f \cdot g)) \land a_{n} \neq c) \to n \in \mathbb{N} \to (f \cdot g)(a_{n}) = f(a_{n}) \cdot g(a_{n})$
\item $\varphi_{122} :\equiv \mathrm{FLim}(f,c,\alpha) \to \mathrm{FLim}(g,c,\beta) \to \mathrm{Seq}(a) \to \mathrm{Conv}(a,c) \to \varphi_{53}$
\item $\varphi_{123} :\equiv \mathrm{FLim}(f,c,\alpha) \to \mathrm{FLim}(g,c,\beta) \to \mathrm{Seq}(a) \to \mathrm{Conv}(a,c) \to \varphi_{54}$
\item $\varphi_{124} :\equiv \mathrm{FLim}(f,c,\alpha) \to \mathrm{FLim}(g,c,\beta) \to \mathrm{Seq}(a) \to \mathrm{Conv}(a,c) \to \varphi_{55}$
\item $\varphi_{125} :\equiv \forall n \in \mathbb{N}\,((f \cdot g) \circ a(n) = (f \cdot g)(a_{n})) \to n \in \mathbb{N} \to (f \cdot g) \circ a(n) = (f \cdot g)(a_{n})$
\item $\varphi_{126} :\equiv \mathrm{FLim}(g,c,\beta) \to \forall a\,(\mathrm{Seq}(a) \to \mathrm{Conv}(a,c) \to \varphi_{76})$
\item $\varphi_{127} :\equiv \mathrm{FLim}(f,c,\alpha) \land \mathrm{FLim}(g,c,\beta) \land (d \in \mathbb{R} \land 0 < d \land \forall x \in \mathbb{R}\,(\varphi_{29}))$
\item $\varphi_{128} :\equiv \forall n \in \mathbb{N}\,(a_{n} \in \mathrm{dom}((f \cdot g)) \land a_{n} \neq c) \to n \in \mathbb{N} \to n \in \mathbb{N} \to (f \cdot g) \circ a(n) = (f \cdot g)(a_{n})$
\item $\varphi_{129} :\equiv \mathrm{FLim}(f,c,\alpha) \to \mathrm{FLim}(g,c,\beta) \to \mathrm{Seq}(a) \to \mathrm{Conv}(a,c) \to \varphi_{61}$
\item $\varphi_{130} :\equiv (\mathrm{FLim}(f,c,\alpha) \to \mathrm{FLim}(g,c,\beta) \to \alpha \in \mathbb{R}) \to \varphi_{69}$
\item $\varphi_{131} :\equiv (f \cdot g)(a_{n}) = f \circ a \otimes g \circ a(n) \to (f \cdot g) \circ a(n) = f \circ a \otimes g \circ a(n)$
\item $\varphi_{132} :\equiv \forall n \in \mathbb{N}\,(a_{n} \in \mathrm{dom}((f \cdot g)) \land a_{n} \neq c) \to n \in \mathbb{N} \to (f \cdot g)(a_{n}) = f \circ a \otimes g \circ a(n)$
\item $\varphi_{133} :\equiv f \circ a(n) \cdot g \circ a(n) = z \cdot g \circ a(n) \to f \circ a(n) \cdot g \circ a(n) = w \cdot g \circ a(n)$
\item $\varphi_{134} :\equiv \forall b\,(\mathrm{Seq}(f \circ a) \to \mathrm{Seq}(b) \to \mathrm{Seq}(f \circ a \otimes b) \land \forall n \in \mathbb{N}\,(f \circ a \otimes b(n) = f \circ a(n) \cdot b_{n}))$
\item $\varphi_{135} :\equiv \forall b\,\forall \alpha\,(\mathrm{Seq}(b) \to \forall n \in \mathbb{N}\,(f \circ a \otimes g \circ a(n) = b_{n}) \to \mathrm{Conv}(f \circ a \otimes g \circ a,\alpha) \to \mathrm{Conv}(b,\alpha))$
\item $\varphi_{136} :\equiv (f \cdot g) \circ a(n) = f \circ a \otimes g \circ a(n) \to f \circ a \otimes g \circ a(n) = (f \cdot g) \circ a(n)$
\item $\varphi_{137} :\equiv \forall n \in \mathbb{N}\,(a_{n} \in \mathrm{dom}((f \cdot g)) \land a_{n} \neq c) \to n \in \mathbb{N} \to (f \cdot g) \circ a(n) = f \circ a \otimes g \circ a(n)$
\item $\varphi_{138} :\equiv \forall n \in \mathbb{N}\,(a_{n} \in \mathrm{dom}((f \cdot g)) \land a_{n} \neq c) \to n \in \mathbb{N} \to f \circ a \otimes g \circ a(n) = (f \cdot g) \circ a(n)$
\item $\varphi_{139} :\equiv \mathrm{FLim}(f,c,\alpha) \to \mathrm{FLim}(g,c,\beta) \to \mathrm{Seq}(a) \to \mathrm{Conv}(a,c) \to \varphi_{67}$
\item $\varphi_{140} :\equiv \mathrm{FLim}(f,c,\alpha) \to \mathrm{FLim}(g,c,\beta) \to \mathrm{Seq}(a) \to \mathrm{Conv}(a,c) \to \varphi_{68}$
\item $\varphi_{141} :\equiv f(a_{n}) \cdot g(a_{n}) = f \circ a \otimes g \circ a(n) \to (f \cdot g)(a_{n}) = f \circ a \otimes g \circ a(n)$
\item $\varphi_{142} :\equiv \forall n \in \mathbb{N}\,(a_{n} \in \mathrm{dom}((f \cdot g)) \land a_{n} \neq c) \to \forall n \in \mathbb{N}\,(f \circ a \otimes g \circ a(n) = (f \cdot g) \circ a(n))$
\item $\varphi_{143} :\equiv \mathrm{FLim}(f,c,\alpha) \to \mathrm{FLim}(g,c,\beta) \to \mathrm{Seq}(a) \to \mathrm{Conv}(a,c) \to \varphi_{70}$
\item $\varphi_{144} :\equiv \mathrm{FLim}(f,c,\alpha) \to \mathrm{FLim}(g,c,\beta) \to \mathrm{Seq}(a) \to \mathrm{Conv}(a,c) \to \varphi_{71}$
\item $\varphi_{145} :\equiv \forall n \in \mathbb{N}\,(a_{n} \in \mathrm{dom}((f \cdot g)) \land a_{n} \neq c) \to n \in \mathbb{N} \to f(a_{n}) \cdot g \circ a(n) = f(a_{n}) \cdot g(a_{n})$
\item $\varphi_{146} :\equiv \mathrm{FLim}(g,c,\beta) \to d \in \mathbb{R} \land 0 < d \land \forall x \in \mathbb{R}\,(\varphi_{29}) \to x \in \mathbb{R} \to x \in \mathbb{R}$
\item $\varphi_{147} :\equiv \mathrm{FLim}(f,c,\alpha) \to \mathrm{FLim}(g,c,\beta) \to \mathrm{Seq}(a) \to \mathrm{Conv}(a,c) \to \varphi_{76}$
\item $\varphi_{148} :\equiv \mathrm{FLim}(f,c,\alpha) \land \mathrm{FLim}(g,c,\beta) \to \mathrm{Seq}(a) \to \mathrm{Conv}(a,c) \to \varphi_{76}$
\item $\varphi_{149} :\equiv \forall n \in \mathbb{N}\,(a_{n} \in \mathrm{dom}((f \cdot g)) \land a_{n} \neq c) \to n \in \mathbb{N} \to f \circ a \otimes g \circ a(n) = f(a_{n}) \cdot g(a_{n})$
\item $\varphi_{150} :\equiv \forall n \in \mathbb{N}\,(a_{n} \in \mathrm{dom}((f \cdot g)) \land a_{n} \neq c) \to n \in \mathbb{N} \to f(a_{n}) \cdot g(a_{n}) = f \circ a \otimes g \circ a(n)$
\item $\varphi_{151} :\equiv \forall a\,(\mathrm{Seq}(a) \to \mathrm{Conv}(a,c) \to \varphi_{76}) \to \mathrm{FLim}((f \cdot g),c,\alpha \cdot \beta)$
\item $\varphi_{152} :\equiv \mathrm{FLim}(f,c,\alpha) \to \mathrm{FLim}(g,c,\beta) \to \mathrm{Seq}(a) \to \mathrm{Conv}(a,c) \to \varphi_{82}$
\item $\varphi_{153} :\equiv \mathrm{Conv}(a,c) \to \forall n \in \mathbb{N}\,(a_{n} \in \mathrm{dom}((f \cdot g)) \land a_{n} \neq c) \to \varphi_{52}$
\item $\varphi_{154} :\equiv \forall n \in \mathbb{N}\,(a_{n} \in \mathrm{dom}((f \cdot g)) \land a_{n} \neq c) \to n \in \mathbb{N} \to f \circ a \otimes g \circ a(n) = f(a_{n}) \cdot g \circ a(n)$
\item $\varphi_{155} :\equiv \forall a\,(\mathrm{FLim}(f,c,\alpha) \land \mathrm{FLim}(g,c,\beta)) \to \forall a\,(\mathrm{Seq}(a) \to \mathrm{Conv}(a,c) \to \varphi_{76})$
\item $\varphi_{156} :\equiv f(a_{n}) \cdot g \circ a(n) = f(a_{n}) \cdot g(a_{n}) \to f \circ a \otimes g \circ a(n) = f(a_{n}) \cdot g(a_{n})$
\item $\varphi_{157} :\equiv \mathrm{FLim}(f,c,\alpha) \to \mathrm{FLim}(g,c,\beta) \to \mathrm{Seq}(a) \to \mathrm{Conv}(a,c) \to \varphi_{88}$
\item $\varphi_{158} :\equiv f \circ a \otimes g \circ a(n) = f(a_{n}) \cdot g(a_{n}) \to f(a_{n}) \cdot g(a_{n}) = f \circ a \otimes g \circ a(n)$
\item $\varphi_{159} :\equiv f(a_{n}) \cdot g \circ a(n) = f(a_{n}) \cdot g \circ a(n) \to f(a_{n}) \cdot g \circ a(n) = f(a_{n}) \cdot g(a_{n})$
\item $\varphi_{160} :\equiv \mathrm{Conv}(a,c) \to \forall n \in \mathbb{N}\,(a_{n} \in \mathrm{dom}((f \cdot g)) \land a_{n} \neq c) \to \varphi_{59}$
\item $\varphi_{161} :\equiv \mathrm{Conv}(a,c) \to \forall n \in \mathbb{N}\,(a_{n} \in \mathrm{dom}((f \cdot g)) \land a_{n} \neq c) \to \varphi_{60}$
\item $\varphi_{162} :\equiv \forall n \in \mathbb{N}\,(a_{n} \in \mathrm{dom}((f \cdot g)) \land a_{n} \neq c) \to n \in \mathbb{N} \to f \circ a \otimes g \circ a(n) = f \circ a(n) \cdot g \circ a(n)$
\item $\varphi_{163} :\equiv d \in \mathbb{R} \land 0 < d \land \forall x \in \mathbb{R}\,(\varphi_{29}) \to x \in \mathbb{R} \to 0 < \lvert x - c \rvert \to 0 < \lvert x - c \rvert$
\item $\varphi_{164} :\equiv \mathrm{FLim}(f,c,\alpha) \to \mathrm{FLim}(g,c,\beta) \to \mathrm{Seq}(a) \to \mathrm{Conv}(a,c) \to \varphi_{89}$
\item $\varphi_{165} :\equiv \mathrm{FLim}(f,c,\alpha) \to \mathrm{FLim}(g,c,\beta) \to \mathrm{Seq}(a) \to \mathrm{Conv}(a,c) \to \varphi_{90}$
\item $\varphi_{166} :\equiv \forall g\,(\mathrm{RFn}(f) \to \mathrm{RFn}(g) \to \mathrm{RFn}((f \cdot g)) \land \varphi_{11} \land \varphi_{23})$
\item $\varphi_{167} :\equiv \mathrm{FLim}(f,c,\alpha) \to \mathrm{FLim}(g,c,\beta) \to \mathrm{Seq}(a) \to \mathrm{Conv}(a,c) \to \varphi_{93}$
\item $\varphi_{168} :\equiv \mathrm{FLim}(f,c,\alpha) \to \mathrm{FLim}(g,c,\beta) \to \mathrm{Seq}(a) \to \mathrm{Conv}(a,c) \to \varphi_{94}$
\item $\varphi_{169} :\equiv \mathrm{FLim}(f,c,\alpha) \to \mathrm{FLim}(g,c,\beta) \to \mathrm{Seq}(a) \to \mathrm{Conv}(a,c) \to \varphi_{96}$
\item $\varphi_{170} :\equiv \mathrm{FLim}(f,c,\alpha) \to \mathrm{FLim}(g,c,\beta) \to \mathrm{Seq}(a) \to \mathrm{Conv}(a,c) \to \varphi_{97}$
\item $\varphi_{171} :\equiv \mathrm{FLim}(f,c,\alpha) \to \mathrm{FLim}(g,c,\beta) \to \mathrm{Seq}(a) \to \mathrm{Conv}(a,c) \to \varphi_{98}$
\item $\varphi_{172} :\equiv \mathrm{FLim}(f,c,\alpha) \to \mathrm{FLim}(g,c,\beta) \to \mathrm{Seq}(a) \to \mathrm{Conv}(a,c) \to \varphi_{105}$
\item $\varphi_{173} :\equiv f \circ a(n) \cdot g \circ a(n) = f(a_{n}) \cdot g \circ a(n) \to f \circ a \otimes g \circ a(n) = f(a_{n}) \cdot g \circ a(n)$
\item $\varphi_{174} :\equiv \mathrm{Conv}(a,c) \to \forall n \in \mathbb{N}\,(a_{n} \in \mathrm{dom}((f \cdot g)) \land a_{n} \neq c) \to \varphi_{72}$
\item $\varphi_{175} :\equiv \mathrm{RFn}((f \cdot g)) \to \mathrm{Seq}(a) \to \forall n \in \mathbb{N}\,(a_{n} \in \mathrm{dom}((f \cdot g))) \to \varphi_{75}$
\item $\varphi_{176} :\equiv \mathrm{Conv}(a,c) \to \forall n \in \mathbb{N}\,(a_{n} \in \mathrm{dom}((f \cdot g)) \land a_{n} \neq c) \to \varphi_{75}$
\item $\varphi_{177} :\equiv \mathrm{FLim}(f,c,\alpha) \to \mathrm{FLim}(g,c,\beta) \to \mathrm{RFn}((f \cdot g)) \land \varphi_{11} \land \varphi_{23}$
\item $\varphi_{178} :\equiv \mathrm{FLim}(f,c,\alpha) \to \mathrm{FLim}(g,c,\beta) \to \mathrm{Seq}(a) \to \mathrm{Conv}(a,c) \to \varphi_{108}$
\item $\varphi_{179} :\equiv \forall w\,(g \circ a(n) = w \to f(a_{n}) \cdot g \circ a(n) = f(a_{n}) \cdot g \circ a(n) \to f(a_{n}) \cdot g \circ a(n) = f(a_{n}) \cdot w)$
\item $\varphi_{180} :\equiv \mathrm{FLim}(f,c,\alpha) \to \mathrm{FLim}(g,c,\beta) \to \mathrm{Seq}(a) \to \mathrm{Conv}(a,c) \to \varphi_{109}$
\item $\varphi_{181} :\equiv \mathrm{FLim}(f,c,\alpha) \to \mathrm{FLim}(g,c,\beta) \to \mathrm{Seq}(a) \to \mathrm{Conv}(a,c) \to \varphi_{110}$
\item $\varphi_{182} :\equiv f \circ a(n) \cdot g \circ a(n) = f \circ a(n) \cdot g \circ a(n) \to f \circ a(n) \cdot g \circ a(n) = f(a_{n}) \cdot g \circ a(n)$
\item $\varphi_{183} :\equiv \forall n \in \mathbb{N}\,(a_{n} \in \mathrm{dom}((f \cdot g)) \land a_{n} \neq c) \to n \in \mathbb{N} \to \varphi_{81}$
\item $\varphi_{184} :\equiv \forall z\,(f \circ a \otimes g \circ a(n) = f(a_{n}) \cdot g \circ a(n) \to f(a_{n}) \cdot g \circ a(n) = z \to f \circ a \otimes g \circ a(n) = z)$
\item $\varphi_{185} :\equiv \forall f\,\forall c\,\forall \alpha\,(\mathrm{RFn}(f) \to c \in \mathbb{R} \to \alpha \in \mathbb{R} \to \mathrm{PNbhd}(f,c) \to \forall a\,(\mathrm{Seq}(a) \to \mathrm{Conv}(a,c) \to \varphi_{44}) \to \mathrm{FLim}(f,c,\alpha))$
\item $\varphi_{186} :\equiv \mathrm{FLim}(f,c,\alpha) \to \mathrm{FLim}(g,c,\beta) \to \mathrm{Seq}(a) \to \mathrm{Conv}(a,c) \to \varphi_{113}$
\item $\varphi_{187} :\equiv \mathrm{FLim}(f,c,\alpha) \to \mathrm{FLim}(g,c,\beta) \to \mathrm{Seq}(a) \to \mathrm{Conv}(a,c) \to \varphi_{114}$
\item $\varphi_{188} :\equiv \mathrm{FLim}(f,c,\alpha) \to \mathrm{FLim}(g,c,\beta) \to \mathrm{Seq}(a) \to \mathrm{Conv}(a,c) \to \varphi_{116}$
\item $\varphi_{189} :\equiv d \in \mathbb{R} \land 0 < d \land \forall x \in \mathbb{R}\,(\varphi_{29}) \to x \in \mathbb{R} \to 0 < \lvert x - c \rvert \to \lvert x - c \rvert < d \to \lvert x - c \rvert < d$
\item $\varphi_{190} :\equiv \forall z\,(f \circ a \otimes g \circ a(n) = f \circ a(n) \cdot g \circ a(n) \to f \circ a(n) \cdot g \circ a(n) = z \to f \circ a \otimes g \circ a(n) = z)$
\item $\varphi_{191} :\equiv \mathrm{Conv}(a,c) \to \forall n \in \mathbb{N}\,(a_{n} \in \mathrm{dom}((f \cdot g)) \land a_{n} \neq c) \to \varphi_{101}$
\item $\varphi_{192} :\equiv \mathrm{Seq}(f \circ a) \to \mathrm{Seq}(g \circ a) \to \mathrm{Seq}(f \circ a \otimes g \circ a) \land \varphi_{81}$
\item $\varphi_{193} :\equiv \mathrm{FLim}(f,c,\alpha) \to \mathrm{FLim}(g,c,\beta) \to \mathrm{Seq}(a) \to \mathrm{Conv}(a,c) \to \varphi_{121}$
\item $\varphi_{194} :\equiv \forall w\,(f \circ a(n) = w \to f \circ a(n) \cdot g \circ a(n) = f \circ a(n) \cdot g \circ a(n) \to f \circ a(n) \cdot g \circ a(n) = w \cdot g \circ a(n))$
\item $\varphi_{195} :\equiv \mathrm{FLim}(f,c,\alpha) \to \mathrm{FLim}(g,c,\beta) \to \mathrm{Seq}(a) \to \mathrm{Conv}(a,c) \to \varphi_{128}$
\item $\varphi_{196} :\equiv \mathrm{FLim}(f,c,\alpha) \to \mathrm{FLim}(g,c,\beta) \to \mathrm{Seq}(a) \to \mathrm{Conv}(a,c) \to \varphi_{132}$
\item $\varphi_{197} :\equiv \mathrm{FLim}(f,c,\alpha) \to \mathrm{FLim}(g,c,\beta) \to \mathrm{Seq}(a) \to \varphi_{153}$
\item $\varphi_{198} :\equiv \mathrm{FLim}(f,c,\alpha) \to \mathrm{FLim}(g,c,\beta) \to \mathrm{Seq}(a) \to \mathrm{Conv}(a,c) \to \varphi_{137}$
\item $\varphi_{199} :\equiv \mathrm{FLim}(f,c,\alpha) \to \mathrm{FLim}(g,c,\beta) \to \mathrm{Seq}(a) \to \mathrm{Conv}(a,c) \to \varphi_{138}$
\item $\varphi_{200} :\equiv \mathrm{FLim}(f,c,\alpha) \to \mathrm{FLim}(g,c,\beta) \to \mathrm{Seq}(a) \to \varphi_{160}$
\item $\varphi_{201} :\equiv \mathrm{FLim}(f,c,\alpha) \to \mathrm{FLim}(g,c,\beta) \to \mathrm{Seq}(a) \to \varphi_{161}$
\item $\varphi_{202} :\equiv \forall \alpha\,(\mathrm{Seq}((f \cdot g) \circ a) \to \forall n \in \mathbb{N}\,(f \circ a \otimes g \circ a(n) = (f \cdot g) \circ a(n)) \to \mathrm{Conv}(f \circ a \otimes g \circ a,\alpha) \to \mathrm{Conv}((f \cdot g) \circ a,\alpha))$
\item $\varphi_{203} :\equiv \mathrm{FLim}(g,c,\beta) \to d \in \mathbb{R} \land 0 < d \land \forall x \in \mathbb{R}\,(\varphi_{29}) \to \varphi_{40}$
\item $\varphi_{204} :\equiv \forall n\,(\varphi_{95} \to n \in \mathbb{N} \to f \circ a \otimes g \circ a(n) = (f \cdot g) \circ a(n))$
\item $\varphi_{205} :\equiv \mathrm{FLim}(f,c,\alpha) \to \mathrm{FLim}(g,c,\beta) \to \mathrm{Seq}(a) \to \mathrm{Conv}(a,c) \to \varphi_{142}$
\item $\varphi_{206} :\equiv \mathrm{FLim}(f,c,\alpha) \to \mathrm{FLim}(g,c,\beta) \to \mathrm{Seq}(a) \to \mathrm{Conv}(a,c) \to \varphi_{145}$
\item $\varphi_{207} :\equiv \mathrm{FLim}(g,c,\beta) \to d \in \mathbb{R} \land 0 < d \land \forall x \in \mathbb{R}\,(\varphi_{29}) \to \varphi_{48}$
\item $\varphi_{208} :\equiv \varphi_{81} \to n \in \mathbb{N} \to f \circ a \otimes g \circ a(n) = f \circ a(n) \cdot g \circ a(n)$
\item $\varphi_{209} :\equiv \forall n\,(\varphi_{95}) \to \forall n \in \mathbb{N}\,(f \circ a \otimes g \circ a(n) = (f \cdot g) \circ a(n))$
\item $\varphi_{210} :\equiv \mathrm{FLim}(f,c,\alpha) \to \mathrm{FLim}(g,c,\beta) \to \mathrm{Seq}(a) \to \mathrm{Conv}(a,c) \to \varphi_{149}$
\item $\varphi_{211} :\equiv \mathrm{FLim}(f,c,\alpha) \to \mathrm{FLim}(g,c,\beta) \to \mathrm{Seq}(a) \to \mathrm{Conv}(a,c) \to \varphi_{150}$
\item $\varphi_{212} :\equiv \forall n \in \mathbb{N}\,(a_{n} \in \mathrm{dom}((f \cdot g)) \land a_{n} \neq c) \to \mathrm{Seq}(f \circ a \otimes g \circ a) \land \varphi_{81}$
\item $\varphi_{213} :\equiv \mathrm{Seq}((f \cdot g) \circ a) \to \forall n \in \mathbb{N}\,(f \circ a \otimes g \circ a(n) = (f \cdot g) \circ a(n)) \to \varphi_{63}$
\item $\varphi_{214} :\equiv \mathrm{FLim}(f,c,\alpha) \to \mathrm{FLim}(g,c,\beta) \to \mathrm{Seq}(a) \to \mathrm{Conv}(a,c) \to \varphi_{154}$
\item $\varphi_{215} :\equiv \mathrm{FLim}(f,c,\alpha) \to \mathrm{FLim}(g,c,\beta) \to \mathrm{Seq}(a) \to \varphi_{174}$
\item $\varphi_{216} :\equiv \forall a\,(\mathrm{FLim}(g,c,\beta) \to \mathrm{Conv}(a,c) \to \varphi_{43}) \to \mathrm{FLim}(g,c,\beta) \to \mathrm{Conv}(a,c) \to \varphi_{43}$
\item $\varphi_{217} :\equiv \forall a\,(\mathrm{FLim}(f,c,\alpha) \to \mathrm{Conv}(a,c) \to \varphi_{44}) \to \mathrm{FLim}(f,c,\alpha) \to \mathrm{Conv}(a,c) \to \varphi_{44}$
\item $\varphi_{218} :\equiv \mathrm{FLim}(f,c,\alpha) \to \mathrm{FLim}(g,c,\beta) \to \mathrm{Seq}(a) \to \varphi_{176}$
\item $\varphi_{219} :\equiv \forall \alpha\,(\mathrm{RFn}((f \cdot g)) \to c \in \mathbb{R} \to \alpha \in \mathbb{R} \to \mathrm{PNbhd}((f \cdot g),c) \to \forall a\,(\mathrm{Seq}(a) \to \mathrm{Conv}(a,c) \to \forall n \in \mathbb{N}\,(a_{n} \in \mathrm{dom}((f \cdot g)) \land a_{n} \neq c) \to \mathrm{Conv}((f \cdot g) \circ a,\alpha)) \to \mathrm{FLim}((f \cdot g),c,\alpha))$
\item $\varphi_{220} :\equiv (\mathrm{FLim}(f,c,\alpha) \to \mathrm{FLim}(g,c,\beta) \to \mathrm{RFn}(g)) \to \varphi_{177}$
\item $\varphi_{221} :\equiv \mathrm{FLim}(f,c,\alpha) \to \mathrm{FLim}(g,c,\beta) \to \mathrm{Seq}(a) \to \mathrm{Conv}(a,c) \to \varphi_{162}$
\item $\varphi_{222} :\equiv \forall y\,((f \cdot g) \circ a(n) = y \to y = (f \cdot g) \circ a(n)) \to \varphi_{136}$
\item $\varphi_{223} :\equiv d \in \mathbb{R} \land 0 < d \land \forall x \in \mathbb{R}\,(\varphi_{29}) \to d \in \mathbb{R} \land 0 < d \land \forall x \in \mathbb{R}\,(\varphi_{29})$
\item $\varphi_{224} :\equiv \forall \alpha\,\forall a\,(\mathrm{FLim}(f,c,\alpha) \to \mathrm{Conv}(a,c) \to \varphi_{44}) \to \forall a\,(\mathrm{FLim}(f,c,\alpha) \to \mathrm{Conv}(a,c) \to \varphi_{44})$
\item $\varphi_{225} :\equiv \mathrm{FLim}(f,c,\alpha) \to \mathrm{FLim}(g,c,\beta) \to \mathrm{Seq}(a) \to \varphi_{191}$
\item $\varphi_{226} :\equiv \mathrm{FLim}(g,c,\beta) \to d \in \mathbb{R} \land 0 < d \land \forall x \in \mathbb{R}\,(\varphi_{29}) \to d \in \mathbb{R} \land 0 < d \land \varphi_{48}$
\item $\varphi_{227} :\equiv \varphi_{103} \to \forall \alpha\,\forall a\,(\mathrm{FLim}(f,c,\alpha) \to \mathrm{Conv}(a,c) \to \varphi_{44})$
\item $\varphi_{228} :\equiv \mathrm{RFn}((f \cdot g)) \to c \in \mathbb{R} \to \alpha \cdot \beta \in \mathbb{R} \to \mathrm{PNbhd}((f \cdot g),c) \to \varphi_{151}$
\item $\varphi_{229} :\equiv \mathrm{FLim}(f,c,\alpha) \to \mathrm{FLim}(g,c,\beta) \to \mathrm{Seq}(a) \to \mathrm{Conv}(a,c) \to \varphi_{183}$
\item $\varphi_{230} :\equiv \mathrm{FLim}(g,c,\beta) \to d \in \mathbb{R} \land 0 < d \land \forall x \in \mathbb{R}\,(\varphi_{29}) \to \varphi_{84}$
\item $\varphi_{231} :\equiv (\varphi_{11} \to \varphi_{39}) \to \mathrm{FLim}(f,c,\alpha) \to \mathrm{FLim}(g,c,\beta) \to \varphi_{39}$
\item $\varphi_{232} :\equiv d \in \mathbb{R} \land 0 < d \land \forall x \in \mathbb{R}\,(\varphi_{29}) \to x \in \mathbb{R} \to 0 < \lvert x - c \rvert \to \lvert x - c \rvert < d \to x \in \mathbb{R} \to \varphi_{29}$
\item $\varphi_{233} :\equiv d \in \mathbb{R} \land 0 < d \land \forall x \in \mathbb{R}\,(\varphi_{29}) \to x \in \mathbb{R} \to 0 < \lvert x - c \rvert \to \lvert x - c \rvert < d \to \forall x \in \mathbb{R}\,(\varphi_{29})$
\item $\varphi_{234} :\equiv (\mathrm{FLim}(f,c,\alpha) \to \mathrm{FLim}(g,c,\beta) \to \mathrm{RFn}(f)) \to \varphi_{220}$
\item $\varphi_{235} :\equiv \forall y\,(f \circ a \otimes g \circ a(n) = y \to y = f \circ a \otimes g \circ a(n)) \to \varphi_{158}$
\item $\varphi_{236} :\equiv \forall c\,(\mathrm{PNbhd}((f \cdot g),c) \leftrightarrow \varphi_{84}) \to (\mathrm{PNbhd}((f \cdot g),c) \leftrightarrow \varphi_{84})$
\item $\varphi_{237} :\equiv \mathrm{FLim}(f,c,\alpha) \to \mathrm{FLim}(g,c,\beta) \to \mathrm{Seq}(a) \to \mathrm{Conv}(a,c) \to \varphi_{212}$
\item $\varphi_{238} :\equiv \forall a\,(\mathrm{RFn}(g) \to \mathrm{Seq}(a) \to \varphi_{99}) \to \mathrm{RFn}(g) \to \mathrm{Seq}(a) \to \varphi_{99}$
\item $\varphi_{239} :\equiv \forall a\,(\mathrm{RFn}(f) \to \mathrm{Seq}(a) \to \varphi_{100}) \to \mathrm{RFn}(f) \to \mathrm{Seq}(a) \to \varphi_{100}$
\item $\varphi_{240} :\equiv \forall f\,\forall a\,(\mathrm{RFn}(f) \to \mathrm{Seq}(a) \to \varphi_{100}) \to \forall a\,(\mathrm{RFn}(g) \to \mathrm{Seq}(a) \to \varphi_{99})$
\item $\varphi_{241} :\equiv \forall f\,\forall a\,(\mathrm{RFn}(f) \to \mathrm{Seq}(a) \to \varphi_{100}) \to \forall a\,(\mathrm{RFn}(f) \to \mathrm{Seq}(a) \to \varphi_{100})$
\item $\varphi_{242} :\equiv (\varphi_{23} \to \varphi_{64}) \to \mathrm{FLim}(f,c,\alpha) \to \mathrm{FLim}(g,c,\beta) \to \varphi_{64}$
\item $\varphi_{243} :\equiv (\mathrm{FLim}(f,c,\alpha) \to \mathrm{FLim}(g,c,\beta) \to \varphi_{11}) \to \varphi_{92} \to \mathrm{FLim}(f,c,\alpha) \to \varphi_{28}$
\item $\varphi_{244} :\equiv \forall g\,\forall c\,(\mathrm{PNbhd}(f,c) \to \varphi_{111}) \to \forall c\,(\mathrm{PNbhd}(f,c) \to \varphi_{111})$
\item $\varphi_{245} :\equiv \forall f\,\forall g\,\forall c\,(\mathrm{PNbhd}(f,c) \to \varphi_{111}) \to \forall g\,\forall c\,(\mathrm{PNbhd}(f,c) \to \varphi_{111})$
\item $\varphi_{246} :\equiv \forall n\,(\varphi_{95} \to n \in \mathbb{N} \to a_{n} \in \mathrm{dom}(g)) \to \forall n\,(\varphi_{95}) \to \forall n \in \mathbb{N}\,(a_{n} \in \mathrm{dom}(g))$
\item $\varphi_{247} :\equiv \forall n\,(\varphi_{95} \to n \in \mathbb{N} \to a_{n} \in \mathrm{dom}(f)) \to \forall n\,(\varphi_{95}) \to \forall n \in \mathbb{N}\,(a_{n} \in \mathrm{dom}(f))$
\item $\varphi_{248} :\equiv (\mathrm{FLim}(f,c,\alpha) \to \mathrm{FLim}(g,c,\beta) \to \alpha \cdot \beta \in \mathbb{R}) \to \varphi_{7} \to (\mathrm{FLim}(f,c,\alpha) \to \varphi_{126}) \to \varphi_{14}$
\item $\varphi_{249} :\equiv \forall n\,(\varphi_{95} \to n \in \mathbb{N} \to a_{n} \in \mathrm{dom}((f \cdot g))) \to \forall n\,(\varphi_{95}) \to \forall n \in \mathbb{N}\,(a_{n} \in \mathrm{dom}((f \cdot g)))$
\item $\varphi_{250} :\equiv (\mathrm{FLim}(f,c,\alpha) \to \mathrm{FLim}(g,c,\beta) \to \varphi_{117}) \to \varphi_{10} \to \neg (d \in \mathbb{R} \land 0 < d \land \forall x \in \mathbb{R}\,(\varphi_{29}))$
\item $\varphi_{251} :\equiv \forall n\,(\varphi_{95} \to n \in \mathbb{N} \to a_{n} \in \mathrm{dom}(g) \land a_{n} \neq c) \to \forall n\,(\varphi_{95}) \to \forall n \in \mathbb{N}\,(a_{n} \in \mathrm{dom}(g) \land a_{n} \neq c)$
\item $\varphi_{252} :\equiv \forall n\,(\varphi_{95} \to n \in \mathbb{N} \to a_{n} \in \mathrm{dom}(f) \land a_{n} \neq c) \to \forall n\,(\varphi_{95}) \to \forall n \in \mathbb{N}\,(a_{n} \in \mathrm{dom}(f) \land a_{n} \neq c)$
\item $\varphi_{253} :\equiv \forall d\,(\varphi_{10} \to \neg (d \in \mathbb{R} \land 0 < d \land \forall x \in \mathbb{R}\,(\varphi_{29}))) \to \forall d\,(\varphi_{10}) \to \forall d\,(\neg (d \in \mathbb{R} \land 0 < d \land \forall x \in \mathbb{R}\,(\varphi_{29})))$
\item $\varphi_{254} :\equiv (\mathrm{FLim}(f,c,\alpha) \to \mathrm{FLim}(g,c,\beta) \to c \in \mathbb{R}) \to \varphi_{248}$
\item $\varphi_{255} :\equiv (\mathrm{PNbhd}(f,c) \to \varphi_{111}) \to \varphi_{2} \to \varphi_{1} \to \mathrm{FLim}(f,c,\alpha) \to \varphi_{115}$
\item $\varphi_{256} :\equiv \varphi_{122} \to (\mathrm{FLim}(f,c,\alpha) \to \mathrm{FLim}(g,c,\beta) \to \varphi_{39}) \to \varphi_{139}$
\item $\varphi_{257} :\equiv \varphi_{122} \to (\mathrm{FLim}(f,c,\alpha) \to \mathrm{FLim}(g,c,\beta) \to \varphi_{39}) \to \varphi_{140}$
\item $\varphi_{258} :\equiv (\mathrm{FLim}(f,c,\alpha) \to \mathrm{FLim}(f,c,\alpha)) \to (\mathrm{FLim}(f,c,\alpha) \to \varphi_{6}) \to \varphi_{170} \to \varphi_{124}$
\item $\varphi_{259} :\equiv (\mathrm{FLim}(f,c,\alpha) \to \varphi_{230}) \to \mathrm{FLim}(f,c,\alpha) \to \mathrm{FLim}(g,c,\beta) \to \varphi_{117}$
\item $\varphi_{260} :\equiv \varphi_{122} \to (\mathrm{FLim}(f,c,\alpha) \to \mathrm{FLim}(g,c,\beta) \to \varphi_{39}) \to \varphi_{164}$
\item $\varphi_{261} :\equiv \varphi_{122} \to (\mathrm{FLim}(f,c,\alpha) \to \mathrm{FLim}(g,c,\beta) \to \varphi_{39}) \to \varphi_{165}$
\item $\varphi_{262} :\equiv (\mathrm{RFn}(f) \to \mathrm{RFn}(g) \to \mathrm{RFn}((f \cdot g)) \land \varphi_{11} \land \varphi_{23}) \to \varphi_{234}$
\item $\varphi_{263} :\equiv (\mathrm{FLim}(f,c,\alpha) \to \mathrm{FLim}(g,c,\beta) \to \mathrm{RFn}(g)) \to \varphi_{5} \to \varphi_{143} \to \varphi_{200}$
\item $\varphi_{264} :\equiv (\mathrm{FLim}(f,c,\alpha) \to \mathrm{FLim}(g,c,\beta) \to \mathrm{RFn}(f)) \to \varphi_{5} \to \varphi_{144} \to \varphi_{201}$
\item $\varphi_{265} :\equiv (\forall x \in \mathbb{R}\,(\varphi_{29}) \to x \in \mathbb{R} \to \varphi_{29}) \to \mathrm{FLim}(f,c,\alpha) \to \mathrm{FLim}(g,c,\beta) \to \varphi_{232}$
\item $\varphi_{266} :\equiv (\varphi_{87} \to \neg (d \in \mathbb{R} \land 0 < d \land \varphi_{48})) \to \mathrm{FLim}(f,c,\alpha) \to \varphi_{230}$
\item $\varphi_{267} :\equiv f(a_{n}) \cdot g \circ a(n) = f(a_{n}) \cdot g \circ a(n) \to (g \circ a(n) = g(a_{n}) \to \varphi_{159}) \to \varphi_{206}$
\item $\varphi_{268} :\equiv (\mathrm{FLim}(f,c,\alpha) \to \mathrm{FLim}(g,c,\beta) \to \varphi_{189}) \to (\mathrm{FLim}(f,c,\alpha) \to \varphi_{28}) \to \mathrm{FLim}(f,c,\alpha) \to \varphi_{203}$
\item $\varphi_{269} :\equiv f \circ a(n) \cdot g \circ a(n) = f \circ a(n) \cdot g \circ a(n) \to (f \circ a(n) = f(a_{n}) \to \varphi_{182}) \to \varphi_{215}$
\item $\varphi_{270} :\equiv (\mathrm{FLim}(g,c,\beta) \to \mathrm{Conv}(a,c) \to \varphi_{43}) \to \varphi_{3} \to (\mathrm{FLim}(f,c,\alpha) \to \varphi_{6}) \to \varphi_{169} \to \varphi_{123}$
\item $\varphi_{271} :\equiv (\mathrm{FLim}(f,c,\alpha) \to \mathrm{FLim}(g,c,\beta) \to \varphi_{223}) \to \mathrm{FLim}(f,c,\alpha) \to \mathrm{FLim}(g,c,\beta) \to \varphi_{233}$
\item $\varphi_{272} :\equiv (\mathrm{FLim}(f,c,\alpha) \to \mathrm{FLim}(g,c,\beta) \to \varphi_{64}) \to \varphi_{172} \to \varphi_{122} \to \varphi_{193}$
\item $\varphi_{273} :\equiv \varphi_{206} \to (f \circ a \otimes g \circ a(n) = f(a_{n}) \cdot g \circ a(n) \to \varphi_{156}) \to \varphi_{210}$
\item $\varphi_{274} :\equiv \forall d\,(\varphi_{10} \to \neg (d \in \mathbb{R} \land 0 < d \land \forall x \in \mathbb{R}\,(\varphi_{29}))) \to \varphi_{253} \to (\varphi_{10} \to \forall d\,(\varphi_{10})) \to \varphi_{7}$
\item $\varphi_{275} :\equiv \varphi_{215} \to (f \circ a \otimes g \circ a(n) = f \circ a(n) \cdot g \circ a(n) \to \varphi_{173}) \to \varphi_{214}$
\item $\varphi_{276} :\equiv \forall a\,(\varphi_{148}) \to (\forall a\,(\varphi_{148}) \to \varphi_{155}) \to \varphi_{20} \to \mathrm{FLim}(f,c,\alpha) \to \varphi_{126}$
\item $\varphi_{277} :\equiv (\mathrm{FLim}(f,c,\alpha) \to \mathrm{FLim}(g,c,\beta) \to \varphi_{223}) \to (\mathrm{FLim}(f,c,\alpha) \to \varphi_{207}) \to \mathrm{FLim}(f,c,\alpha) \to \varphi_{226}$
\item $\varphi_{278} :\equiv \varphi_{193} \to \varphi_{211} \to ((f \cdot g)(a_{n}) = f(a_{n}) \cdot g(a_{n}) \to \varphi_{141}) \to \varphi_{196}$
\item $\varphi_{279} :\equiv \forall n\,(\varphi_{95} \to n \in \mathbb{N} \to a_{n} \in \mathrm{dom}(g)) \to \varphi_{246} \to (\varphi_{95} \to \forall n\,(\varphi_{95})) \to \varphi_{143}$
\item $\varphi_{280} :\equiv \forall n\,(\varphi_{95} \to n \in \mathbb{N} \to a_{n} \in \mathrm{dom}(f)) \to \varphi_{247} \to (\varphi_{95} \to \forall n\,(\varphi_{95})) \to \varphi_{144}$
\item $\varphi_{281} :\equiv \forall n\,(\varphi_{95} \to n \in \mathbb{N} \to a_{n} \in \mathrm{dom}((f \cdot g))) \to \varphi_{249} \to (\varphi_{95} \to \forall n\,(\varphi_{95})) \to \varphi_{157}$
\item $\varphi_{282} :\equiv (\mathrm{FLim}(f,c,\alpha) \to \varphi_{146}) \to (\mathrm{FLim}(f,c,\alpha) \to \mathrm{FLim}(g,c,\beta) \to \varphi_{163}) \to \varphi_{268}$
\item $\varphi_{283} :\equiv \forall n\,(\varphi_{95} \to n \in \mathbb{N} \to a_{n} \in \mathrm{dom}(g) \land a_{n} \neq c) \to \varphi_{251} \to (\varphi_{95} \to \forall n\,(\varphi_{95})) \to \varphi_{169}$
\item $\varphi_{284} :\equiv \forall n\,(\varphi_{95} \to n \in \mathbb{N} \to a_{n} \in \mathrm{dom}(f) \land a_{n} \neq c) \to \varphi_{252} \to (\varphi_{95} \to \forall n\,(\varphi_{95})) \to \varphi_{170}$
\item $\varphi_{285} :\equiv (\forall x\,(\varphi_{127} \to \varphi_{40}) \to \forall x\,(\varphi_{127}) \to \varphi_{48}) \to (\varphi_{127} \to \forall x\,(\varphi_{127})) \to \mathrm{FLim}(f,c,\alpha) \to \varphi_{207}$
\end{abbrevlist}
\end{fullproof}

% 自動生成 (prover/build.js): FLimScal  ノード数 416
\begin{fullproof}{定理1.8(1')}{$\forall l\,\forall f\,\forall c\,\forall \alpha\,(l \in \mathbb{R} \to \mathrm{FLim}(f,c,\alpha) \to \mathrm{FLim}((l f),c,l \cdot \alpha))$}{416}\phantomsection\label{tree:FLimScal}
\noindent$\Pi_{1}$:
\begin{ptree}
\AXC{$\forall f\,\forall c\,\forall \alpha\,(\mathrm{FLim}(f,c,\alpha) \to \mathrm{RFn}(f))$\lbl{[D2]}}
\AXC{$\varphi_{14}$\lbl{A2}}
\BIC{$\forall c\,\forall \alpha\,(\mathrm{FLim}(f,c,\alpha) \to \mathrm{RFn}(f))$}
\AXC{$\varphi_{8}$\lbl{A2}}
\BIC{$\forall \alpha\,(\mathrm{FLim}(f,c,\alpha) \to \mathrm{RFn}(f))$}
\AXC{$\varphi_{3}$\lbl{A2}}
\BIC{$\mathrm{FLim}(f,c,\alpha) \to \mathrm{RFn}(f)$}
\end{ptree}
\noindent$\Pi_{2}$:
\begin{ptree}
\AXC{$\stackrel{\Pi_{1}}{\vdots}$}\noLine
\UIC{$\mathrm{FLim}(f,c,\alpha) \to \mathrm{RFn}(f)$}
\AXC{$(\mathrm{FLim}(f,c,\alpha) \to \mathrm{RFn}(f)) \to \varphi_{34}$\lbl{A1}}
\BIC{$\varphi_{34}$}
\end{ptree}
\noindent$\Pi_{3}$:
\begin{ptree}
\AXC{$l \in \mathbb{R} \to \mathrm{FLim}(f,c,\alpha) \to \mathrm{FLim}(f,c,\alpha)$\lbl{A1}}
\AXC{$\stackrel{\Pi_{2}}{\vdots}$}\noLine
\UIC{$\varphi_{34}$}
\BIC{$l \in \mathbb{R} \to \mathrm{FLim}(f,c,\alpha) \to \mathrm{RFn}(f)$}
\end{ptree}
\noindent$\Pi_{4}$:
\begin{ptree}
\AXC{$\forall l\,\forall f\,(l \in \mathbb{R} \to \mathrm{RFn}(f) \to \varphi_{87})$\lbl{[S19]}}
\AXC{$\varphi_{182}$\lbl{A2}}
\BIC{$\forall f\,(l \in \mathbb{R} \to \mathrm{RFn}(f) \to \varphi_{87})$}
\AXC{$\varphi_{181}$\lbl{A2}}
\BIC{$l \in \mathbb{R} \to \mathrm{RFn}(f) \to \varphi_{87}$}
\AXC{$(l \in \mathbb{R} \to \mathrm{RFn}(f) \to \varphi_{87}) \to (l \in \mathbb{R} \to l \in \mathbb{R}) \to \varphi_{153}$\lbl{A1}}
\BIC{$(l \in \mathbb{R} \to l \in \mathbb{R}) \to \varphi_{153}$}
\end{ptree}
\noindent$\Pi_{5}$:
\begin{ptree}
\AXC{$\stackrel{\Pi_{3}}{\vdots}$}\noLine
\UIC{$l \in \mathbb{R} \to \mathrm{FLim}(f,c,\alpha) \to \mathrm{RFn}(f)$}
\AXC{$l \in \mathbb{R} \to l \in \mathbb{R}$\lbl{A1}}
\AXC{$\stackrel{\Pi_{4}}{\vdots}$}\noLine
\UIC{$(l \in \mathbb{R} \to l \in \mathbb{R}) \to \varphi_{153}$}
\BIC{$\varphi_{153}$}
\BIC{$l \in \mathbb{R} \to \mathrm{FLim}(f,c,\alpha) \to \varphi_{87}$}
\end{ptree}
\noindent$\Pi_{6}$:
\begin{ptree}
\AXC{$\stackrel{\Pi_{5}}{\vdots}$}\noLine
\UIC{$l \in \mathbb{R} \to \mathrm{FLim}(f,c,\alpha) \to \varphi_{87}$}
\AXC{$(l \in \mathbb{R} \to \mathrm{FLim}(f,c,\alpha) \to \varphi_{87}) \to l \in \mathbb{R} \to \varphi_{6}$\lbl{A1}}
\BIC{$l \in \mathbb{R} \to \varphi_{6}$}
\end{ptree}
\noindent$\Pi_{7}$:
\begin{ptree}
\AXC{$\varphi_{103}$\lbl{A1}}
\AXC{$\varphi_{80}$\lbl{A2}}
\AXC{$\varphi_{134}$\lbl{A1}}
\AXC{$\varphi_{134} \to \varphi_{80} \to \varphi_{133}$\lbl{A1}}
\BIC{$\varphi_{80} \to \varphi_{133}$}
\BIC{$\varphi_{133}$}
\AXC{$\varphi_{133} \to \varphi_{103} \to (l \in \mathbb{R} \to \varphi_{17}) \to \varphi_{127}$\lbl{A1}}
\BIC{$\varphi_{103} \to (l \in \mathbb{R} \to \varphi_{17}) \to \varphi_{127}$}
\BIC{$(l \in \mathbb{R} \to \varphi_{17}) \to \varphi_{127}$}
\end{ptree}
\noindent$\Pi_{8}$:
\begin{ptree}
\AXC{$\varphi_{55}$\lbl{A2}}
\AXC{$\stackrel{\Pi_{6}}{\vdots}$}\noLine
\UIC{$l \in \mathbb{R} \to \varphi_{6}$}
\AXC{$(l \in \mathbb{R} \to \varphi_{6}) \to \varphi_{55} \to l \in \mathbb{R} \to \varphi_{17}$\lbl{A1}}
\BIC{$\varphi_{55} \to l \in \mathbb{R} \to \varphi_{17}$}
\BIC{$l \in \mathbb{R} \to \varphi_{17}$}
\AXC{$\stackrel{\Pi_{7}}{\vdots}$}\noLine
\UIC{$(l \in \mathbb{R} \to \varphi_{17}) \to \varphi_{127}$}
\BIC{$\varphi_{127}$}
\end{ptree}
\noindent$\Pi_{9}$:
\begin{ptree}
\AXC{$\varphi_{71} \to \forall n\,(\varphi_{71})$\lbl{A4}}
\AXC{$\varphi_{186}$\lbl{A3}}
\AXC{$\stackrel{\Pi_{8}}{\vdots}$}\noLine
\UIC{$\varphi_{127}$}
\AXC{$\varphi_{127} \to \varphi_{71} \to n \in \mathbb{N} \to a_{n} \in \mathrm{dom}(f) \land a_{n} \neq c$\lbl{A1}}
\BIC{$\varphi_{71} \to n \in \mathbb{N} \to a_{n} \in \mathrm{dom}(f) \land a_{n} \neq c$}
\UIC{$\forall n\,(\varphi_{71} \to n \in \mathbb{N} \to a_{n} \in \mathrm{dom}(f) \land a_{n} \neq c)$}
\AXC{$\varphi_{204}$\lbl{A1}}
\BIC{$\varphi_{186} \to (\varphi_{71} \to \forall n\,(\varphi_{71})) \to \varphi_{130}$}
\BIC{$(\varphi_{71} \to \forall n\,(\varphi_{71})) \to \varphi_{130}$}
\BIC{$\varphi_{130}$}
\end{ptree}
\noindent$\Pi_{10}$:
\begin{ptree}
\AXC{$\forall f\,(\varphi_{75})$\lbl{定理1.7(⇒)}}
\AXC{$\forall f\,(\varphi_{75}) \to \varphi_{75}$\lbl{A2}}
\BIC{$\varphi_{75}$}
\AXC{$\varphi_{170}$\lbl{A2}}
\BIC{$\forall \alpha\,\forall a\,(\mathrm{FLim}(f,c,\alpha) \to \mathrm{Conv}(a,c) \to \varphi_{25})$}
\AXC{$\varphi_{168}$\lbl{A2}}
\BIC{$\forall a\,(\mathrm{FLim}(f,c,\alpha) \to \mathrm{Conv}(a,c) \to \varphi_{25})$}
\AXC{$\varphi_{164}$\lbl{A2}}
\BIC{$\mathrm{FLim}(f,c,\alpha) \to \mathrm{Conv}(a,c) \to \varphi_{25}$}
\end{ptree}
\noindent$\Pi_{11}$:
\begin{ptree}
\AXC{$\stackrel{\Pi_{10}}{\vdots}$}\noLine
\UIC{$\mathrm{FLim}(f,c,\alpha) \to \mathrm{Conv}(a,c) \to \varphi_{25}$}
\AXC{$(\mathrm{FLim}(f,c,\alpha) \to \mathrm{Conv}(a,c) \to \varphi_{25}) \to \varphi_{193}$\lbl{A1}}
\BIC{$\varphi_{193}$}
\end{ptree}
\noindent$\Pi_{12}$:
\begin{ptree}
\AXC{$\stackrel{\Pi_{9}}{\vdots}$}\noLine
\UIC{$\varphi_{130}$}
\AXC{$l \in \mathbb{R} \to \varphi_{1}$\lbl{A1}}
\AXC{$l \in \mathbb{R} \to \mathrm{FLim}(f,c,\alpha) \to \mathrm{FLim}(f,c,\alpha)$\lbl{A1}}
\AXC{$\stackrel{\Pi_{11}}{\vdots}$}\noLine
\UIC{$\varphi_{193}$}
\BIC{$(l \in \mathbb{R} \to \varphi_{1}) \to \varphi_{130} \to \varphi_{104}$}
\BIC{$\varphi_{130} \to \varphi_{104}$}
\BIC{$\varphi_{104}$}
\end{ptree}
\noindent$\Pi_{13}$:
\begin{ptree}
\AXC{$\varphi_{12}$\lbl{定理2.1(2)}}
\AXC{$\varphi_{12} \to \varphi_{9}$\lbl{A2}}
\BIC{$\varphi_{9}$}
\AXC{$\varphi_{9} \to \forall \alpha\,(\varphi_{23})$\lbl{A2}}
\BIC{$\forall \alpha\,(\varphi_{23})$}
\AXC{$\forall \alpha\,(\varphi_{23}) \to \varphi_{23}$\lbl{A2}}
\BIC{$\varphi_{23}$}
\AXC{$\varphi_{23} \to (l \in \mathbb{R} \to l \in \mathbb{R}) \to \varphi_{104} \to \varphi_{118}$\lbl{A1}}
\BIC{$(l \in \mathbb{R} \to l \in \mathbb{R}) \to \varphi_{104} \to \varphi_{118}$}
\end{ptree}
\noindent$\Pi_{14}$:
\begin{ptree}
\AXC{$x = x$\lbl{A5}}
\AXC{$x = y \to x = x \to y = x$\lbl{A6}}
\AXC{$(x = y \to x = x \to y = x) \to x = x \to x = y \to y = x$\lbl{A1}}
\BIC{$x = x \to x = y \to y = x$}
\BIC{$x = y \to y = x$}
\UIC{$\forall y\,(x = y \to y = x)$}
\UIC{$\forall x\,\forall y\,(x = y \to y = x)$}
\end{ptree}
\noindent$\Pi_{15}$:
\begin{ptree}
\AXC{$y = z \to x = y \to x = z$\lbl{A6}}
\AXC{$(y = z \to x = y \to x = z) \to x = y \to y = z \to x = z$\lbl{A1}}
\BIC{$x = y \to y = z \to x = z$}
\UIC{$\forall z\,(x = y \to y = z \to x = z)$}
\UIC{$\forall y\,\forall z\,(x = y \to y = z \to x = z)$}
\UIC{$\forall x\,\forall y\,\forall z\,(x = y \to y = z \to x = z)$}
\AXC{$\forall x\,\forall y\,\forall z\,(x = y \to y = z \to x = z) \to \varphi_{32}$\lbl{A2}}
\BIC{$\varphi_{32}$}
\end{ptree}
\noindent$\Pi_{16}$:
\begin{ptree}
\AXC{$\stackrel{\Pi_{15}}{\vdots}$}\noLine
\UIC{$\varphi_{32}$}
\AXC{$\varphi_{32} \to \varphi_{86}$\lbl{A2}}
\BIC{$\varphi_{86}$}
\AXC{$\varphi_{86} \to (l f) \circ a(n) = (l f)(a_{n}) \to \varphi_{78}$\lbl{A2}}
\BIC{$(l f) \circ a(n) = (l f)(a_{n}) \to \varphi_{78}$}
\end{ptree}
\noindent$\Pi_{17}$:
\begin{ptree}
\AXC{$y = z \to x = y \to x = z$\lbl{A6}}
\AXC{$(y = z \to x = y \to x = z) \to x = y \to y = z \to x = z$\lbl{A1}}
\BIC{$x = y \to y = z \to x = z$}
\UIC{$\forall z\,(x = y \to y = z \to x = z)$}
\UIC{$\forall y\,\forall z\,(x = y \to y = z \to x = z)$}
\UIC{$\forall x\,\forall y\,\forall z\,(x = y \to y = z \to x = z)$}
\AXC{$\forall x\,\forall y\,\forall z\,(x = y \to y = z \to x = z) \to \varphi_{20}$\lbl{A2}}
\BIC{$\varphi_{20}$}
\end{ptree}
\noindent$\Pi_{18}$:
\begin{ptree}
\AXC{$y = z \to x = y \to x = z$\lbl{A6}}
\AXC{$(y = z \to x = y \to x = z) \to x = y \to y = z \to x = z$\lbl{A1}}
\BIC{$x = y \to y = z \to x = z$}
\UIC{$\forall z\,(x = y \to y = z \to x = z)$}
\UIC{$\forall y\,\forall z\,(x = y \to y = z \to x = z)$}
\UIC{$\forall x\,\forall y\,\forall z\,(x = y \to y = z \to x = z)$}
\AXC{$\forall x\,\forall y\,\forall z\,(x = y \to y = z \to x = z) \to \varphi_{28}$\lbl{A2}}
\BIC{$\varphi_{28}$}
\end{ptree}
\noindent$\Pi_{19}$:
\begin{ptree}
\AXC{$\stackrel{\Pi_{18}}{\vdots}$}\noLine
\UIC{$\varphi_{28}$}
\AXC{$\varphi_{28} \to \varphi_{82}$\lbl{A2}}
\BIC{$\varphi_{82}$}
\AXC{$\varphi_{82} \to l \odot f \circ a(n) = l \cdot f \circ a(n) \to \varphi_{61}$\lbl{A2}}
\BIC{$l \odot f \circ a(n) = l \cdot f \circ a(n) \to \varphi_{61}$}
\end{ptree}
\noindent$\Pi_{20}$:
\begin{ptree}
\AXC{$z = w \to l \cdot f \circ a(n) = l \cdot z \to l \cdot f \circ a(n) = l \cdot w$\lbl{A6}}
\UIC{$\varphi_{30}$}
\UIC{$\forall z\,(\varphi_{30})$}
\AXC{$\forall z\,(\varphi_{30}) \to \varphi_{76}$\lbl{A2}}
\BIC{$\varphi_{76}$}
\AXC{$\varphi_{76} \to f \circ a(n) = f(a_{n}) \to \varphi_{62}$\lbl{A2}}
\BIC{$f \circ a(n) = f(a_{n}) \to \varphi_{62}$}
\end{ptree}
\noindent$\Pi_{21}$:
\begin{ptree}
\AXC{$\varphi_{103}$\lbl{A1}}
\AXC{$\varphi_{80}$\lbl{A2}}
\AXC{$\varphi_{134}$\lbl{A1}}
\AXC{$\varphi_{134} \to \varphi_{80} \to \varphi_{133}$\lbl{A1}}
\BIC{$\varphi_{80} \to \varphi_{133}$}
\BIC{$\varphi_{133}$}
\AXC{$\varphi_{133} \to \varphi_{103} \to (l \in \mathbb{R} \to \varphi_{17}) \to \varphi_{116}$\lbl{A1}}
\BIC{$\varphi_{103} \to (l \in \mathbb{R} \to \varphi_{17}) \to \varphi_{116}$}
\BIC{$(l \in \mathbb{R} \to \varphi_{17}) \to \varphi_{116}$}
\end{ptree}
\noindent$\Pi_{22}$:
\begin{ptree}
\AXC{$\varphi_{55}$\lbl{A2}}
\AXC{$\stackrel{\Pi_{6}}{\vdots}$}\noLine
\UIC{$l \in \mathbb{R} \to \varphi_{6}$}
\AXC{$(l \in \mathbb{R} \to \varphi_{6}) \to \varphi_{55} \to l \in \mathbb{R} \to \varphi_{17}$\lbl{A1}}
\BIC{$\varphi_{55} \to l \in \mathbb{R} \to \varphi_{17}$}
\BIC{$l \in \mathbb{R} \to \varphi_{17}$}
\AXC{$\stackrel{\Pi_{21}}{\vdots}$}\noLine
\UIC{$(l \in \mathbb{R} \to \varphi_{17}) \to \varphi_{116}$}
\BIC{$\varphi_{116}$}
\end{ptree}
\noindent$\Pi_{23}$:
\begin{ptree}
\AXC{$\varphi_{71} \to \forall n\,(\varphi_{71})$\lbl{A4}}
\AXC{$\varphi_{183}$\lbl{A3}}
\AXC{$\stackrel{\Pi_{22}}{\vdots}$}\noLine
\UIC{$\varphi_{116}$}
\AXC{$\varphi_{116} \to \varphi_{71} \to n \in \mathbb{N} \to a_{n} \in \mathrm{dom}(f)$\lbl{A1}}
\BIC{$\varphi_{71} \to n \in \mathbb{N} \to a_{n} \in \mathrm{dom}(f)$}
\UIC{$\forall n\,(\varphi_{71} \to n \in \mathbb{N} \to a_{n} \in \mathrm{dom}(f))$}
\AXC{$\varphi_{202}$\lbl{A1}}
\BIC{$\varphi_{183} \to (\varphi_{71} \to \forall n\,(\varphi_{71})) \to \varphi_{119}$}
\BIC{$(\varphi_{71} \to \forall n\,(\varphi_{71})) \to \varphi_{119}$}
\BIC{$\varphi_{119}$}
\end{ptree}
\noindent$\Pi_{24}$:
\begin{ptree}
\AXC{$\forall f\,\forall a\,(\mathrm{RFn}(f) \to \mathrm{Seq}(a) \to \varphi_{72})$\lbl{[S8]}}
\AXC{$\varphi_{179}$\lbl{A2}}
\BIC{$\forall a\,(\mathrm{RFn}(f) \to \mathrm{Seq}(a) \to \varphi_{72})$}
\AXC{$\varphi_{177}$\lbl{A2}}
\BIC{$\mathrm{RFn}(f) \to \mathrm{Seq}(a) \to \varphi_{72}$}
\AXC{$\varphi_{198}$\lbl{A1}}
\BIC{$(l \in \mathbb{R} \to \mathrm{FLim}(f,c,\alpha) \to \mathrm{RFn}(f)) \to \varphi_{191}$}
\end{ptree}
\noindent$\Pi_{25}$:
\begin{ptree}
\AXC{$\stackrel{\Pi_{3}}{\vdots}$}\noLine
\UIC{$l \in \mathbb{R} \to \mathrm{FLim}(f,c,\alpha) \to \mathrm{RFn}(f)$}
\AXC{$\stackrel{\Pi_{24}}{\vdots}$}\noLine
\UIC{$(l \in \mathbb{R} \to \mathrm{FLim}(f,c,\alpha) \to \mathrm{RFn}(f)) \to \varphi_{191}$}
\BIC{$\varphi_{191}$}
\end{ptree}
\noindent$\Pi_{26}$:
\begin{ptree}
\AXC{$\stackrel{\Pi_{23}}{\vdots}$}\noLine
\UIC{$\varphi_{119}$}
\AXC{$l \in \mathbb{R} \to \mathrm{FLim}(f,c,\alpha) \to \mathrm{Seq}(a) \to \mathrm{Seq}(a)$\lbl{A1}}
\AXC{$\stackrel{\Pi_{25}}{\vdots}$}\noLine
\UIC{$\varphi_{191}$}
\BIC{$\varphi_{119} \to l \in \mathbb{R} \to \mathrm{FLim}(f,c,\alpha) \to \mathrm{Seq}(a) \to \varphi_{128}$}
\BIC{$l \in \mathbb{R} \to \mathrm{FLim}(f,c,\alpha) \to \mathrm{Seq}(a) \to \varphi_{128}$}
\end{ptree}
\noindent$\Pi_{27}$:
\begin{ptree}
\AXC{$\stackrel{\Pi_{26}}{\vdots}$}\noLine
\UIC{$l \in \mathbb{R} \to \mathrm{FLim}(f,c,\alpha) \to \mathrm{Seq}(a) \to \varphi_{128}$}
\AXC{$(l \in \mathbb{R} \to \mathrm{FLim}(f,c,\alpha) \to \mathrm{Seq}(a) \to \varphi_{128}) \to \varphi_{145}$\lbl{A1}}
\BIC{$\varphi_{145}$}
\AXC{$\varphi_{145} \to \varphi_{65} \to \varphi_{138}$\lbl{A1}}
\BIC{$\varphi_{65} \to \varphi_{138}$}
\end{ptree}
\noindent$\Pi_{28}$:
\begin{ptree}
\AXC{$\varphi_{103}$\lbl{A1}}
\AXC{$\varphi_{65}$\lbl{A2}}
\AXC{$\stackrel{\Pi_{27}}{\vdots}$}\noLine
\UIC{$\varphi_{65} \to \varphi_{138}$}
\BIC{$\varphi_{138}$}
\AXC{$\varphi_{138} \to \varphi_{103} \to \varphi_{129}$\lbl{A1}}
\BIC{$\varphi_{103} \to \varphi_{129}$}
\BIC{$\varphi_{129}$}
\AXC{$\varphi_{129} \to \varphi_{188}$\lbl{A1}}
\BIC{$\varphi_{188}$}
\end{ptree}
\noindent$\Pi_{29}$:
\begin{ptree}
\AXC{$\stackrel{\Pi_{20}}{\vdots}$}\noLine
\UIC{$f \circ a(n) = f(a_{n}) \to \varphi_{62}$}
\AXC{$x = x$\lbl{A5}}
\UIC{$\forall x\,(x = x)$}
\AXC{$\forall x\,(x = x) \to l \cdot f \circ a(n) = l \cdot f \circ a(n)$\lbl{A2}}
\BIC{$l \cdot f \circ a(n) = l \cdot f \circ a(n)$}
\AXC{$\stackrel{\Pi_{28}}{\vdots}$}\noLine
\UIC{$\varphi_{188}$}
\BIC{$(f \circ a(n) = f(a_{n}) \to \varphi_{62}) \to \varphi_{135}$}
\BIC{$\varphi_{135}$}
\end{ptree}
\noindent$\Pi_{30}$:
\begin{ptree}
\AXC{$\forall a\,\forall \alpha\,(\mathrm{Conv}(a,\alpha) \to \mathrm{Seq}(a))$\lbl{[C1]}}
\AXC{$\varphi_{31}$\lbl{A2}}
\BIC{$\forall \alpha\,(\mathrm{Conv}(f \circ a,\alpha) \to \mathrm{Seq}(f \circ a))$}
\AXC{$\varphi_{50}$\lbl{A2}}
\BIC{$\mathrm{Conv}(f \circ a,\alpha) \to \mathrm{Seq}(f \circ a)$}
\AXC{$(\mathrm{Conv}(f \circ a,\alpha) \to \mathrm{Seq}(f \circ a)) \to \varphi_{104} \to \varphi_{99}$\lbl{A1}}
\BIC{$\varphi_{104} \to \varphi_{99}$}
\end{ptree}
\noindent$\Pi_{31}$:
\begin{ptree}
\AXC{$\varphi_{51}$\lbl{[S3]}}
\AXC{$\varphi_{51} \to \varphi_{49}$\lbl{A2}}
\BIC{$\varphi_{49}$}
\AXC{$\varphi_{49} \to l \in \mathbb{R} \to \mathrm{Seq}(f \circ a) \to \varphi_{59}$\lbl{A2}}
\BIC{$l \in \mathbb{R} \to \mathrm{Seq}(f \circ a) \to \varphi_{59}$}
\AXC{$(l \in \mathbb{R} \to \mathrm{Seq}(f \circ a) \to \varphi_{59}) \to \varphi_{189}$\lbl{A1}}
\BIC{$\varphi_{189}$}
\end{ptree}
\noindent$\Pi_{32}$:
\begin{ptree}
\AXC{$\stackrel{\Pi_{12}}{\vdots}$}\noLine
\UIC{$\varphi_{104}$}
\AXC{$\stackrel{\Pi_{30}}{\vdots}$}\noLine
\UIC{$\varphi_{104} \to \varphi_{99}$}
\BIC{$\varphi_{99}$}
\AXC{$l \in \mathbb{R} \to l \in \mathbb{R}$\lbl{A1}}
\AXC{$\stackrel{\Pi_{31}}{\vdots}$}\noLine
\UIC{$\varphi_{189}$}
\BIC{$\varphi_{99} \to l \in \mathbb{R} \to \mathrm{FLim}(f,c,\alpha) \to \mathrm{Seq}(a) \to \varphi_{137}$}
\BIC{$l \in \mathbb{R} \to \mathrm{FLim}(f,c,\alpha) \to \mathrm{Seq}(a) \to \varphi_{137}$}
\end{ptree}
\noindent$\Pi_{33}$:
\begin{ptree}
\AXC{$\stackrel{\Pi_{32}}{\vdots}$}\noLine
\UIC{$l \in \mathbb{R} \to \mathrm{FLim}(f,c,\alpha) \to \mathrm{Seq}(a) \to \varphi_{137}$}
\AXC{$(l \in \mathbb{R} \to \mathrm{FLim}(f,c,\alpha) \to \mathrm{Seq}(a) \to \varphi_{137}) \to \varphi_{157}$\lbl{A1}}
\BIC{$\varphi_{157}$}
\AXC{$\varphi_{157} \to \varphi_{107} \to \varphi_{155}$\lbl{A1}}
\BIC{$\varphi_{107} \to \varphi_{155}$}
\end{ptree}
\noindent$\Pi_{34}$:
\begin{ptree}
\AXC{$\varphi_{103}$\lbl{A1}}
\AXC{$\varphi_{107}$\lbl{A2}}
\AXC{$\stackrel{\Pi_{33}}{\vdots}$}\noLine
\UIC{$\varphi_{107} \to \varphi_{155}$}
\BIC{$\varphi_{155}$}
\AXC{$\varphi_{155} \to \varphi_{103} \to \varphi_{146}$\lbl{A1}}
\BIC{$\varphi_{103} \to \varphi_{146}$}
\BIC{$\varphi_{146}$}
\AXC{$\varphi_{146} \to \varphi_{135} \to (l \odot f \circ a(n) = l \cdot f \circ a(n) \to \varphi_{61}) \to \varphi_{139}$\lbl{A1}}
\BIC{$\varphi_{135} \to (l \odot f \circ a(n) = l \cdot f \circ a(n) \to \varphi_{61}) \to \varphi_{139}$}
\end{ptree}
\noindent$\Pi_{35}$:
\begin{ptree}
\AXC{$\stackrel{\Pi_{19}}{\vdots}$}\noLine
\UIC{$l \odot f \circ a(n) = l \cdot f \circ a(n) \to \varphi_{61}$}
\AXC{$\stackrel{\Pi_{29}}{\vdots}$}\noLine
\UIC{$\varphi_{135}$}
\AXC{$\stackrel{\Pi_{34}}{\vdots}$}\noLine
\UIC{$\varphi_{135} \to (l \odot f \circ a(n) = l \cdot f \circ a(n) \to \varphi_{61}) \to \varphi_{139}$}
\BIC{$(l \odot f \circ a(n) = l \cdot f \circ a(n) \to \varphi_{61}) \to \varphi_{139}$}
\BIC{$\varphi_{139}$}
\AXC{$\varphi_{139} \to \varphi_{64} \to \varphi_{140}$\lbl{A1}}
\BIC{$\varphi_{64} \to \varphi_{140}$}
\end{ptree}
\noindent$\Pi_{36}$:
\begin{ptree}
\AXC{$\stackrel{\Pi_{14}}{\vdots}$}\noLine
\UIC{$\forall x\,\forall y\,(x = y \to y = x)$}
\AXC{$\varphi_{42}$\lbl{A2}}
\BIC{$\forall y\,(l \odot f \circ a(n) = y \to y = l \odot f \circ a(n))$}
\AXC{$\forall y\,(l \odot f \circ a(n) = y \to y = l \odot f \circ a(n)) \to \varphi_{64}$\lbl{A2}}
\BIC{$\varphi_{64}$}
\AXC{$\stackrel{\Pi_{35}}{\vdots}$}\noLine
\UIC{$\varphi_{64} \to \varphi_{140}$}
\BIC{$\varphi_{140}$}
\end{ptree}
\noindent$\Pi_{37}$:
\begin{ptree}
\AXC{$\stackrel{\Pi_{5}}{\vdots}$}\noLine
\UIC{$l \in \mathbb{R} \to \mathrm{FLim}(f,c,\alpha) \to \varphi_{87}$}
\AXC{$(l \in \mathbb{R} \to \mathrm{FLim}(f,c,\alpha) \to \varphi_{87}) \to l \in \mathbb{R} \to \varphi_{29}$\lbl{A1}}
\BIC{$l \in \mathbb{R} \to \varphi_{29}$}
\end{ptree}
\noindent$\Pi_{38}$:
\begin{ptree}
\AXC{$\varphi_{112}$\lbl{A2}}
\AXC{$\stackrel{\Pi_{37}}{\vdots}$}\noLine
\UIC{$l \in \mathbb{R} \to \varphi_{29}$}
\AXC{$(l \in \mathbb{R} \to \varphi_{29}) \to \varphi_{112} \to l \in \mathbb{R} \to \varphi_{54}$\lbl{A1}}
\BIC{$\varphi_{112} \to l \in \mathbb{R} \to \varphi_{54}$}
\BIC{$l \in \mathbb{R} \to \varphi_{54}$}
\AXC{$(l \in \mathbb{R} \to \varphi_{54}) \to \varphi_{133} \to \varphi_{103} \to \varphi_{136}$\lbl{A1}}
\BIC{$\varphi_{133} \to \varphi_{103} \to \varphi_{136}$}
\end{ptree}
\noindent$\Pi_{39}$:
\begin{ptree}
\AXC{$\varphi_{103}$\lbl{A1}}
\AXC{$\varphi_{80}$\lbl{A2}}
\AXC{$\varphi_{134}$\lbl{A1}}
\AXC{$\varphi_{134} \to \varphi_{80} \to \varphi_{133}$\lbl{A1}}
\BIC{$\varphi_{80} \to \varphi_{133}$}
\BIC{$\varphi_{133}$}
\AXC{$\stackrel{\Pi_{38}}{\vdots}$}\noLine
\UIC{$\varphi_{133} \to \varphi_{103} \to \varphi_{136}$}
\BIC{$\varphi_{103} \to \varphi_{136}$}
\BIC{$\varphi_{136}$}
\end{ptree}
\noindent$\Pi_{40}$:
\begin{ptree}
\AXC{$\stackrel{\Pi_{36}}{\vdots}$}\noLine
\UIC{$\varphi_{140}$}
\AXC{$\stackrel{\Pi_{39}}{\vdots}$}\noLine
\UIC{$\varphi_{136}$}
\AXC{$\varphi_{136} \to \varphi_{140} \to ((l f)(a_{n}) = l \cdot f(a_{n}) \to \varphi_{68}) \to \varphi_{147}$\lbl{A1}}
\BIC{$\varphi_{140} \to ((l f)(a_{n}) = l \cdot f(a_{n}) \to \varphi_{68}) \to \varphi_{147}$}
\BIC{$((l f)(a_{n}) = l \cdot f(a_{n}) \to \varphi_{68}) \to \varphi_{147}$}
\end{ptree}
\noindent$\Pi_{41}$:
\begin{ptree}
\AXC{$\stackrel{\Pi_{17}}{\vdots}$}\noLine
\UIC{$\varphi_{20}$}
\AXC{$\varphi_{20} \to \varphi_{67}$\lbl{A2}}
\BIC{$\varphi_{67}$}
\AXC{$\varphi_{67} \to (l f)(a_{n}) = l \cdot f(a_{n}) \to \varphi_{68}$\lbl{A2}}
\BIC{$(l f)(a_{n}) = l \cdot f(a_{n}) \to \varphi_{68}$}
\AXC{$\stackrel{\Pi_{40}}{\vdots}$}\noLine
\UIC{$((l f)(a_{n}) = l \cdot f(a_{n}) \to \varphi_{68}) \to \varphi_{147}$}
\BIC{$\varphi_{147}$}
\end{ptree}
\noindent$\Pi_{42}$:
\begin{ptree}
\AXC{$\varphi_{103}$\lbl{A1}}
\AXC{$\varphi_{80}$\lbl{A2}}
\AXC{$\varphi_{134}$\lbl{A1}}
\AXC{$\varphi_{134} \to \varphi_{80} \to \varphi_{133}$\lbl{A1}}
\BIC{$\varphi_{80} \to \varphi_{133}$}
\BIC{$\varphi_{133}$}
\AXC{$\varphi_{133} \to \varphi_{103} \to \varphi_{123}$\lbl{A1}}
\BIC{$\varphi_{103} \to \varphi_{123}$}
\BIC{$\varphi_{123}$}
\AXC{$\varphi_{123} \to \varphi_{71} \to n \in \mathbb{N} \to a_{n} \in \mathrm{dom}((l f))$\lbl{A1}}
\BIC{$\varphi_{71} \to n \in \mathbb{N} \to a_{n} \in \mathrm{dom}((l f))$}
\UIC{$\forall n\,(\varphi_{71} \to n \in \mathbb{N} \to a_{n} \in \mathrm{dom}((l f)))$}
\end{ptree}
\noindent$\Pi_{43}$:
\begin{ptree}
\AXC{$\varphi_{71} \to \forall n\,(\varphi_{71})$\lbl{A4}}
\AXC{$\varphi_{185}$\lbl{A3}}
\AXC{$\stackrel{\Pi_{42}}{\vdots}$}\noLine
\UIC{$\forall n\,(\varphi_{71} \to n \in \mathbb{N} \to a_{n} \in \mathrm{dom}((l f)))$}
\AXC{$\varphi_{203}$\lbl{A1}}
\BIC{$\varphi_{185} \to (\varphi_{71} \to \forall n\,(\varphi_{71})) \to \varphi_{125}$}
\BIC{$(\varphi_{71} \to \forall n\,(\varphi_{71})) \to \varphi_{125}$}
\BIC{$\varphi_{125}$}
\end{ptree}
\noindent$\Pi_{44}$:
\begin{ptree}
\AXC{$\stackrel{\Pi_{5}}{\vdots}$}\noLine
\UIC{$l \in \mathbb{R} \to \mathrm{FLim}(f,c,\alpha) \to \varphi_{87}$}
\AXC{$\varphi_{160}$\lbl{A1}}
\BIC{$l \in \mathbb{R} \to \mathrm{FLim}(f,c,\alpha) \to \mathrm{RFn}((l f))$}
\end{ptree}
\noindent$\Pi_{45}$:
\begin{ptree}
\AXC{$\forall f\,\forall a\,(\mathrm{RFn}(f) \to \mathrm{Seq}(a) \to \varphi_{72})$\lbl{[S8]}}
\AXC{$\forall f\,\forall a\,(\mathrm{RFn}(f) \to \mathrm{Seq}(a) \to \varphi_{72}) \to \forall a\,(\varphi_{143})$\lbl{A2}}
\BIC{$\forall a\,(\varphi_{143})$}
\AXC{$\forall a\,(\varphi_{143}) \to \varphi_{143}$\lbl{A2}}
\BIC{$\varphi_{143}$}
\end{ptree}
\noindent$\Pi_{46}$:
\begin{ptree}
\AXC{$\stackrel{\Pi_{44}}{\vdots}$}\noLine
\UIC{$l \in \mathbb{R} \to \mathrm{FLim}(f,c,\alpha) \to \mathrm{RFn}((l f))$}
\AXC{$\stackrel{\Pi_{45}}{\vdots}$}\noLine
\UIC{$\varphi_{143}$}
\AXC{$\varphi_{201}$\lbl{A1}}
\BIC{$(l \in \mathbb{R} \to \mathrm{FLim}(f,c,\alpha) \to \mathrm{RFn}((l f))) \to \varphi_{194}$}
\BIC{$\varphi_{194}$}
\end{ptree}
\noindent$\Pi_{47}$:
\begin{ptree}
\AXC{$\stackrel{\Pi_{43}}{\vdots}$}\noLine
\UIC{$\varphi_{125}$}
\AXC{$l \in \mathbb{R} \to \mathrm{FLim}(f,c,\alpha) \to \mathrm{Seq}(a) \to \mathrm{Seq}(a)$\lbl{A1}}
\AXC{$\stackrel{\Pi_{46}}{\vdots}$}\noLine
\UIC{$\varphi_{194}$}
\BIC{$\varphi_{125} \to l \in \mathbb{R} \to \mathrm{FLim}(f,c,\alpha) \to \mathrm{Seq}(a) \to \varphi_{142}$}
\BIC{$l \in \mathbb{R} \to \mathrm{FLim}(f,c,\alpha) \to \mathrm{Seq}(a) \to \varphi_{142}$}
\end{ptree}
\noindent$\Pi_{48}$:
\begin{ptree}
\AXC{$\stackrel{\Pi_{47}}{\vdots}$}\noLine
\UIC{$l \in \mathbb{R} \to \mathrm{FLim}(f,c,\alpha) \to \mathrm{Seq}(a) \to \varphi_{142}$}
\AXC{$(l \in \mathbb{R} \to \mathrm{FLim}(f,c,\alpha) \to \mathrm{Seq}(a) \to \varphi_{142}) \to \varphi_{158}$\lbl{A1}}
\BIC{$\varphi_{158}$}
\AXC{$\varphi_{158} \to \varphi_{110} \to \varphi_{156}$\lbl{A1}}
\BIC{$\varphi_{110} \to \varphi_{156}$}
\end{ptree}
\noindent$\Pi_{49}$:
\begin{ptree}
\AXC{$\varphi_{103}$\lbl{A1}}
\AXC{$\varphi_{110}$\lbl{A2}}
\AXC{$\stackrel{\Pi_{48}}{\vdots}$}\noLine
\UIC{$\varphi_{110} \to \varphi_{156}$}
\BIC{$\varphi_{156}$}
\AXC{$\varphi_{156} \to \varphi_{103} \to \varphi_{148}$\lbl{A1}}
\BIC{$\varphi_{103} \to \varphi_{148}$}
\BIC{$\varphi_{148}$}
\AXC{$\varphi_{148} \to \varphi_{147} \to ((l f) \circ a(n) = (l f)(a_{n}) \to \varphi_{78}) \to \varphi_{150}$\lbl{A1}}
\BIC{$\varphi_{147} \to ((l f) \circ a(n) = (l f)(a_{n}) \to \varphi_{78}) \to \varphi_{150}$}
\end{ptree}
\noindent$\Pi_{50}$:
\begin{ptree}
\AXC{$\stackrel{\Pi_{16}}{\vdots}$}\noLine
\UIC{$(l f) \circ a(n) = (l f)(a_{n}) \to \varphi_{78}$}
\AXC{$\stackrel{\Pi_{41}}{\vdots}$}\noLine
\UIC{$\varphi_{147}$}
\AXC{$\stackrel{\Pi_{49}}{\vdots}$}\noLine
\UIC{$\varphi_{147} \to ((l f) \circ a(n) = (l f)(a_{n}) \to \varphi_{78}) \to \varphi_{150}$}
\BIC{$((l f) \circ a(n) = (l f)(a_{n}) \to \varphi_{78}) \to \varphi_{150}$}
\BIC{$\varphi_{150}$}
\AXC{$\varphi_{150} \to \varphi_{81} \to \varphi_{151}$\lbl{A1}}
\BIC{$\varphi_{81} \to \varphi_{151}$}
\end{ptree}
\noindent$\Pi_{51}$:
\begin{ptree}
\AXC{$\stackrel{\Pi_{14}}{\vdots}$}\noLine
\UIC{$\forall x\,\forall y\,(x = y \to y = x)$}
\AXC{$\varphi_{46}$\lbl{A2}}
\BIC{$\forall y\,((l f) \circ a(n) = y \to y = (l f) \circ a(n))$}
\AXC{$\forall y\,((l f) \circ a(n) = y \to y = (l f) \circ a(n)) \to \varphi_{81}$\lbl{A2}}
\BIC{$\varphi_{81}$}
\AXC{$\stackrel{\Pi_{50}}{\vdots}$}\noLine
\UIC{$\varphi_{81} \to \varphi_{151}$}
\BIC{$\varphi_{151}$}
\end{ptree}
\noindent$\Pi_{52}$:
\begin{ptree}
\AXC{$\varphi_{71} \to \forall n\,(\varphi_{71})$\lbl{A4}}
\AXC{$\varphi_{190}$\lbl{A3}}
\AXC{$\stackrel{\Pi_{51}}{\vdots}$}\noLine
\UIC{$\varphi_{151}$}
\AXC{$\varphi_{151} \to \varphi_{71} \to n \in \mathbb{N} \to l \odot f \circ a(n) = (l f) \circ a(n)$\lbl{A1}}
\BIC{$\varphi_{71} \to n \in \mathbb{N} \to l \odot f \circ a(n) = (l f) \circ a(n)$}
\UIC{$\forall n\,(\varphi_{71} \to n \in \mathbb{N} \to l \odot f \circ a(n) = (l f) \circ a(n))$}
\AXC{$\varphi_{205}$\lbl{A1}}
\BIC{$\varphi_{190} \to (\varphi_{71} \to \forall n\,(\varphi_{71})) \to \varphi_{152}$}
\BIC{$(\varphi_{71} \to \forall n\,(\varphi_{71})) \to \varphi_{152}$}
\BIC{$\varphi_{152}$}
\end{ptree}
\noindent$\Pi_{53}$:
\begin{ptree}
\AXC{$\stackrel{\Pi_{47}}{\vdots}$}\noLine
\UIC{$l \in \mathbb{R} \to \mathrm{FLim}(f,c,\alpha) \to \mathrm{Seq}(a) \to \varphi_{142}$}
\AXC{$(l \in \mathbb{R} \to \mathrm{FLim}(f,c,\alpha) \to \mathrm{Seq}(a) \to \varphi_{142}) \to \varphi_{109}$\lbl{A1}}
\BIC{$\varphi_{109}$}
\end{ptree}
\noindent$\Pi_{54}$:
\begin{ptree}
\AXC{$\stackrel{\Pi_{52}}{\vdots}$}\noLine
\UIC{$\varphi_{152}$}
\AXC{$\stackrel{\Pi_{53}}{\vdots}$}\noLine
\UIC{$\varphi_{109}$}
\AXC{$\varphi_{35}$\lbl{[C13]}}
\AXC{$\varphi_{35} \to \varphi_{79}$\lbl{A2}}
\BIC{$\varphi_{79}$}
\AXC{$\varphi_{79} \to \varphi_{141}$\lbl{A2}}
\BIC{$\varphi_{141}$}
\AXC{$\varphi_{141} \to \varphi_{149}$\lbl{A2}}
\BIC{$\varphi_{149}$}
\AXC{$\varphi_{149} \to \varphi_{109} \to \varphi_{152} \to \varphi_{118} \to \varphi_{120}$\lbl{A1}}
\BIC{$\varphi_{109} \to \varphi_{152} \to \varphi_{118} \to \varphi_{120}$}
\BIC{$\varphi_{152} \to \varphi_{118} \to \varphi_{120}$}
\BIC{$\varphi_{118} \to \varphi_{120}$}
\end{ptree}
\noindent$\Pi_{55}$:
\begin{ptree}
\AXC{$\stackrel{\Pi_{12}}{\vdots}$}\noLine
\UIC{$\varphi_{104}$}
\AXC{$l \in \mathbb{R} \to l \in \mathbb{R}$\lbl{A1}}
\AXC{$\stackrel{\Pi_{13}}{\vdots}$}\noLine
\UIC{$(l \in \mathbb{R} \to l \in \mathbb{R}) \to \varphi_{104} \to \varphi_{118}$}
\BIC{$\varphi_{104} \to \varphi_{118}$}
\BIC{$\varphi_{118}$}
\AXC{$\stackrel{\Pi_{54}}{\vdots}$}\noLine
\UIC{$\varphi_{118} \to \varphi_{120}$}
\BIC{$\varphi_{120}$}
\AXC{$\varphi_{120} \to \varphi_{121}$\lbl{A1}}
\BIC{$\varphi_{121}$}
\UIC{$\forall a\,(\varphi_{121})$}
\end{ptree}
\noindent$\Pi_{56}$:
\begin{ptree}
\AXC{$\varphi_{7}$\lbl{A4}}
\AXC{$\forall a\,(\varphi_{121}) \to \varphi_{124}$\lbl{A3}}
\AXC{$\stackrel{\Pi_{55}}{\vdots}$}\noLine
\UIC{$\forall a\,(\varphi_{121})$}
\AXC{$\forall a\,(\varphi_{121}) \to (\forall a\,(\varphi_{121}) \to \varphi_{124}) \to \varphi_{7} \to l \in \mathbb{R} \to \varphi_{108}$\lbl{A1}}
\BIC{$(\forall a\,(\varphi_{121}) \to \varphi_{124}) \to \varphi_{7} \to l \in \mathbb{R} \to \varphi_{108}$}
\BIC{$\varphi_{7} \to l \in \mathbb{R} \to \varphi_{108}$}
\BIC{$l \in \mathbb{R} \to \varphi_{108}$}
\end{ptree}
\noindent$\Pi_{57}$:
\begin{ptree}
\AXC{$l \in \mathbb{R} \to \varphi_{161}$\lbl{A1}}
\AXC{$(l \in \mathbb{R} \to \varphi_{161}) \to l \in \mathbb{R} \to \mathrm{FLim}(f,c,\alpha) \to \varphi_{166}$\lbl{A1}}
\BIC{$l \in \mathbb{R} \to \mathrm{FLim}(f,c,\alpha) \to \varphi_{166}$}
\end{ptree}
\noindent$\Pi_{58}$:
\begin{ptree}
\AXC{$\varphi_{106}$\lbl{A2}}
\AXC{$\stackrel{\Pi_{57}}{\vdots}$}\noLine
\UIC{$l \in \mathbb{R} \to \mathrm{FLim}(f,c,\alpha) \to \varphi_{166}$}
\AXC{$(l \in \mathbb{R} \to \mathrm{FLim}(f,c,\alpha) \to \varphi_{166}) \to \varphi_{106} \to l \in \mathbb{R} \to \varphi_{173}$\lbl{A1}}
\BIC{$\varphi_{106} \to l \in \mathbb{R} \to \varphi_{173}$}
\BIC{$l \in \mathbb{R} \to \varphi_{173}$}
\end{ptree}
\noindent$\Pi_{59}$:
\begin{ptree}
\AXC{$l \in \mathbb{R} \to \mathrm{FLim}(f,c,\alpha) \to \varphi_{115}$\lbl{A1}}
\AXC{$l \in \mathbb{R} \to \varphi_{102}$\lbl{A1}}
\AXC{$\stackrel{\Pi_{58}}{\vdots}$}\noLine
\UIC{$l \in \mathbb{R} \to \varphi_{173}$}
\AXC{$(l \in \mathbb{R} \to \varphi_{173}) \to (l \in \mathbb{R} \to \varphi_{102}) \to \varphi_{200}$\lbl{A1}}
\BIC{$(l \in \mathbb{R} \to \varphi_{102}) \to \varphi_{200}$}
\BIC{$\varphi_{200}$}
\BIC{$(l \in \mathbb{R} \to \mathrm{FLim}(f,c,\alpha) \to \varphi_{132}) \to \varphi_{178}$}
\end{ptree}
\noindent$\Pi_{60}$:
\begin{ptree}
\AXC{$l \in \mathbb{R} \to \mathrm{FLim}(f,c,\alpha) \to \varphi_{132}$\lbl{A1}}
\AXC{$\stackrel{\Pi_{59}}{\vdots}$}\noLine
\UIC{$(l \in \mathbb{R} \to \mathrm{FLim}(f,c,\alpha) \to \varphi_{132}) \to \varphi_{178}$}
\BIC{$\varphi_{178}$}
\end{ptree}
\noindent$\Pi_{61}$:
\begin{ptree}
\AXC{$\varphi_{37}$\lbl{A2}}
\AXC{$\stackrel{\Pi_{6}}{\vdots}$}\noLine
\UIC{$l \in \mathbb{R} \to \varphi_{6}$}
\AXC{$(l \in \mathbb{R} \to \varphi_{6}) \to \varphi_{37} \to l \in \mathbb{R} \to \varphi_{4}$\lbl{A1}}
\BIC{$\varphi_{37} \to l \in \mathbb{R} \to \varphi_{4}$}
\BIC{$l \in \mathbb{R} \to \varphi_{4}$}
\AXC{$\stackrel{\Pi_{60}}{\vdots}$}\noLine
\UIC{$\varphi_{178}$}
\BIC{$l \in \mathbb{R} \to \mathrm{FLim}(f,c,\alpha) \to r \in \mathbb{R} \land 0 < r \land \varphi_{19} \to \varphi_{27}$}
\AXC{$\varphi_{192}$\lbl{A1}}
\BIC{$\varphi_{159}$}
\UIC{$\forall x\,(\varphi_{159})$}
\end{ptree}
\noindent$\Pi_{62}$:
\begin{ptree}
\AXC{$\varphi_{174}$\lbl{A4}}
\AXC{$\forall x\,(\varphi_{159}) \to \varphi_{90} \to \varphi_{33}$\lbl{A3}}
\AXC{$\stackrel{\Pi_{61}}{\vdots}$}\noLine
\UIC{$\forall x\,(\varphi_{159})$}
\AXC{$\forall x\,(\varphi_{159}) \to (\forall x\,(\varphi_{159}) \to \varphi_{90} \to \varphi_{33}) \to \varphi_{196}$\lbl{A1}}
\BIC{$(\forall x\,(\varphi_{159}) \to \varphi_{90} \to \varphi_{33}) \to \varphi_{196}$}
\BIC{$\varphi_{196}$}
\BIC{$l \in \mathbb{R} \to \mathrm{FLim}(f,c,\alpha) \to r \in \mathbb{R} \land 0 < r \land \varphi_{19} \to \varphi_{33}$}
\end{ptree}
\noindent$\Pi_{63}$:
\begin{ptree}
\AXC{$\stackrel{\Pi_{62}}{\vdots}$}\noLine
\UIC{$l \in \mathbb{R} \to \mathrm{FLim}(f,c,\alpha) \to r \in \mathbb{R} \land 0 < r \land \varphi_{19} \to \varphi_{33}$}
\AXC{$l \in \mathbb{R} \to \varphi_{161}$\lbl{A1}}
\AXC{$(l \in \mathbb{R} \to \varphi_{161}) \to \varphi_{195}$\lbl{A1}}
\BIC{$\varphi_{195}$}
\BIC{$l \in \mathbb{R} \to \varphi_{163}$}
\end{ptree}
\noindent$\Pi_{64}$:
\begin{ptree}
\AXC{$\varphi_{162}$\lbl{A2}}
\AXC{$\stackrel{\Pi_{63}}{\vdots}$}\noLine
\UIC{$l \in \mathbb{R} \to \varphi_{163}$}
\AXC{$(l \in \mathbb{R} \to \varphi_{163}) \to \varphi_{162} \to l \in \mathbb{R} \to \varphi_{165}$\lbl{A1}}
\BIC{$\varphi_{162} \to l \in \mathbb{R} \to \varphi_{165}$}
\BIC{$l \in \mathbb{R} \to \varphi_{165}$}
\end{ptree}
\noindent$\Pi_{65}$:
\begin{ptree}
\AXC{$\stackrel{\Pi_{64}}{\vdots}$}\noLine
\UIC{$l \in \mathbb{R} \to \varphi_{165}$}
\AXC{$\varphi_{77}$\lbl{DefPNbhd}}
\AXC{$\varphi_{77} \to \varphi_{96}$\lbl{A2}}
\BIC{$\varphi_{96}$}
\AXC{$\varphi_{175}$\lbl{A2}}
\BIC{$\mathrm{PNbhd}((l f),c) \leftrightarrow \exists r\,(r \in \mathbb{R} \land 0 < r \land \varphi_{33})$}
\AXC{$\varphi_{197}$\lbl{A1}}
\BIC{$(l \in \mathbb{R} \to \varphi_{165}) \to l \in \mathbb{R} \to \varphi_{98}$}
\BIC{$l \in \mathbb{R} \to \varphi_{98}$}
\end{ptree}
\noindent$\Pi_{66}$:
\begin{ptree}
\AXC{$\stackrel{\Pi_{65}}{\vdots}$}\noLine
\UIC{$l \in \mathbb{R} \to \varphi_{98}$}
\AXC{$(l \in \mathbb{R} \to \varphi_{98}) \to \varphi_{2} \to \neg (r \in \mathbb{R} \land 0 < r \land \varphi_{19})$\lbl{A1}}
\BIC{$\varphi_{2} \to \neg (r \in \mathbb{R} \land 0 < r \land \varphi_{19})$}
\UIC{$\forall r\,(\varphi_{2} \to \neg (r \in \mathbb{R} \land 0 < r \land \varphi_{19}))$}
\end{ptree}
\noindent$\Pi_{67}$:
\begin{ptree}
\AXC{$\forall f\,\forall c\,\forall \alpha\,(\mathrm{FLim}(f,c,\alpha) \to \mathrm{PNbhd}(f,c))$\lbl{[D5]}}
\AXC{$\varphi_{21}$\lbl{A2}}
\BIC{$\forall c\,\forall \alpha\,(\mathrm{FLim}(f,c,\alpha) \to \mathrm{PNbhd}(f,c))$}
\AXC{$\varphi_{15}$\lbl{A2}}
\BIC{$\forall \alpha\,(\mathrm{FLim}(f,c,\alpha) \to \mathrm{PNbhd}(f,c))$}
\AXC{$\varphi_{10}$\lbl{A2}}
\BIC{$\mathrm{FLim}(f,c,\alpha) \to \mathrm{PNbhd}(f,c)$}
\end{ptree}
\noindent$\Pi_{68}$:
\begin{ptree}
\AXC{$\stackrel{\Pi_{67}}{\vdots}$}\noLine
\UIC{$\mathrm{FLim}(f,c,\alpha) \to \mathrm{PNbhd}(f,c)$}
\AXC{$(\mathrm{FLim}(f,c,\alpha) \to \mathrm{PNbhd}(f,c)) \to \varphi_{38}$\lbl{A1}}
\BIC{$\varphi_{38}$}
\end{ptree}
\noindent$\Pi_{69}$:
\begin{ptree}
\AXC{$l \in \mathbb{R} \to \mathrm{FLim}(f,c,\alpha) \to \mathrm{FLim}(f,c,\alpha)$\lbl{A1}}
\AXC{$\stackrel{\Pi_{68}}{\vdots}$}\noLine
\UIC{$\varphi_{38}$}
\BIC{$l \in \mathbb{R} \to \mathrm{FLim}(f,c,\alpha) \to \mathrm{PNbhd}(f,c)$}
\AXC{$\varphi_{180}$\lbl{A1}}
\BIC{$(\mathrm{PNbhd}(f,c) \leftrightarrow \exists r\,(r \in \mathbb{R} \land 0 < r \land \varphi_{19})) \to \varphi_{88}$}
\end{ptree}
\noindent$\Pi_{70}$:
\begin{ptree}
\AXC{$\varphi_{77}$\lbl{DefPNbhd}}
\AXC{$\varphi_{171}$\lbl{A2}}
\BIC{$\forall c\,(\mathrm{PNbhd}(f,c) \leftrightarrow \exists r\,(r \in \mathbb{R} \land 0 < r \land \varphi_{19}))$}
\AXC{$\varphi_{169}$\lbl{A2}}
\BIC{$\mathrm{PNbhd}(f,c) \leftrightarrow \exists r\,(r \in \mathbb{R} \land 0 < r \land \varphi_{19})$}
\AXC{$\stackrel{\Pi_{69}}{\vdots}$}\noLine
\UIC{$(\mathrm{PNbhd}(f,c) \leftrightarrow \exists r\,(r \in \mathbb{R} \land 0 < r \land \varphi_{19})) \to \varphi_{88}$}
\BIC{$\varphi_{88}$}
\end{ptree}
\noindent$\Pi_{71}$:
\begin{ptree}
\AXC{$\varphi_{184}$\lbl{A3}}
\AXC{$\stackrel{\Pi_{66}}{\vdots}$}\noLine
\UIC{$\forall r\,(\varphi_{2} \to \neg (r \in \mathbb{R} \land 0 < r \land \varphi_{19}))$}
\AXC{$\stackrel{\Pi_{70}}{\vdots}$}\noLine
\UIC{$\varphi_{88}$}
\AXC{$\varphi_{88} \to \forall r\,(\varphi_{2} \to \neg (r \in \mathbb{R} \land 0 < r \land \varphi_{19})) \to \varphi_{184} \to \varphi_{126}$\lbl{A1}}
\BIC{$\forall r\,(\varphi_{2} \to \neg (r \in \mathbb{R} \land 0 < r \land \varphi_{19})) \to \varphi_{184} \to \varphi_{126}$}
\BIC{$\varphi_{184} \to \varphi_{126}$}
\BIC{$\varphi_{126}$}
\end{ptree}
\noindent$\Pi_{72}$:
\begin{ptree}
\AXC{$\forall f\,\forall c\,\forall \alpha\,(\mathrm{FLim}(f,c,\alpha) \to \alpha \in \mathbb{R})$\lbl{[D4]}}
\AXC{$\varphi_{24}$\lbl{A2}}
\BIC{$\forall c\,\forall \alpha\,(\mathrm{FLim}(f,c,\alpha) \to \alpha \in \mathbb{R})$}
\AXC{$\varphi_{18}$\lbl{A2}}
\BIC{$\forall \alpha\,(\mathrm{FLim}(f,c,\alpha) \to \alpha \in \mathbb{R})$}
\AXC{$\varphi_{13}$\lbl{A2}}
\BIC{$\mathrm{FLim}(f,c,\alpha) \to \alpha \in \mathbb{R}$}
\end{ptree}
\noindent$\Pi_{73}$:
\begin{ptree}
\AXC{$\stackrel{\Pi_{72}}{\vdots}$}\noLine
\UIC{$\mathrm{FLim}(f,c,\alpha) \to \alpha \in \mathbb{R}$}
\AXC{$(\mathrm{FLim}(f,c,\alpha) \to \alpha \in \mathbb{R}) \to \varphi_{43}$\lbl{A1}}
\BIC{$\varphi_{43}$}
\end{ptree}
\noindent$\Pi_{74}$:
\begin{ptree}
\AXC{$l \in \mathbb{R} \to l \in \mathbb{R}$\lbl{A1}}
\AXC{$\forall x\,\forall y\,(x \in \mathbb{R} \to y \in \mathbb{R} \to x \cdot y \in \mathbb{R})$\lbl{(R1)}}
\AXC{$\varphi_{56}$\lbl{A2}}
\BIC{$\forall y\,(l \in \mathbb{R} \to y \in \mathbb{R} \to l \cdot y \in \mathbb{R})$}
\AXC{$\varphi_{52}$\lbl{A2}}
\BIC{$l \in \mathbb{R} \to \alpha \in \mathbb{R} \to l \cdot \alpha \in \mathbb{R}$}
\AXC{$\varphi_{131}$\lbl{A1}}
\BIC{$(l \in \mathbb{R} \to l \in \mathbb{R}) \to \varphi_{47}$}
\BIC{$\varphi_{47}$}
\end{ptree}
\noindent$\Pi_{75}$:
\begin{ptree}
\AXC{$l \in \mathbb{R} \to \mathrm{FLim}(f,c,\alpha) \to \mathrm{FLim}(f,c,\alpha)$\lbl{A1}}
\AXC{$\stackrel{\Pi_{73}}{\vdots}$}\noLine
\UIC{$\varphi_{43}$}
\BIC{$l \in \mathbb{R} \to \mathrm{FLim}(f,c,\alpha) \to \alpha \in \mathbb{R}$}
\AXC{$\stackrel{\Pi_{74}}{\vdots}$}\noLine
\UIC{$\varphi_{47}$}
\BIC{$l \in \mathbb{R} \to \mathrm{FLim}(f,c,\alpha) \to l \cdot \alpha \in \mathbb{R}$}
\end{ptree}
\noindent$\Pi_{76}$:
\begin{ptree}
\AXC{$\forall f\,\forall c\,\forall \alpha\,(\mathrm{FLim}(f,c,\alpha) \to c \in \mathbb{R})$\lbl{[D3]}}
\AXC{$\varphi_{22}$\lbl{A2}}
\BIC{$\forall c\,\forall \alpha\,(\mathrm{FLim}(f,c,\alpha) \to c \in \mathbb{R})$}
\AXC{$\varphi_{16}$\lbl{A2}}
\BIC{$\forall \alpha\,(\mathrm{FLim}(f,c,\alpha) \to c \in \mathbb{R})$}
\AXC{$\varphi_{11}$\lbl{A2}}
\BIC{$\mathrm{FLim}(f,c,\alpha) \to c \in \mathbb{R}$}
\end{ptree}
\noindent$\Pi_{77}$:
\begin{ptree}
\AXC{$\stackrel{\Pi_{76}}{\vdots}$}\noLine
\UIC{$\mathrm{FLim}(f,c,\alpha) \to c \in \mathbb{R}$}
\AXC{$(\mathrm{FLim}(f,c,\alpha) \to c \in \mathbb{R}) \to \varphi_{39}$\lbl{A1}}
\BIC{$\varphi_{39}$}
\end{ptree}
\noindent$\Pi_{78}$:
\begin{ptree}
\AXC{$\varphi_{144}$\lbl{定理1.7(⇐)}}
\AXC{$\varphi_{144} \to \forall c\,(\varphi_{167})$\lbl{A2}}
\BIC{$\forall c\,(\varphi_{167})$}
\AXC{$\forall c\,(\varphi_{167}) \to \varphi_{167}$\lbl{A2}}
\BIC{$\varphi_{167}$}
\AXC{$\varphi_{167} \to \varphi_{172}$\lbl{A2}}
\BIC{$\varphi_{172}$}
\AXC{$\varphi_{199}$\lbl{A1}}
\BIC{$(l \in \mathbb{R} \to \mathrm{FLim}(f,c,\alpha) \to \mathrm{RFn}((l f))) \to \varphi_{187}$}
\end{ptree}
\noindent$\Pi_{79}$:
\begin{ptree}
\AXC{$\stackrel{\Pi_{44}}{\vdots}$}\noLine
\UIC{$l \in \mathbb{R} \to \mathrm{FLim}(f,c,\alpha) \to \mathrm{RFn}((l f))$}
\AXC{$\stackrel{\Pi_{78}}{\vdots}$}\noLine
\UIC{$(l \in \mathbb{R} \to \mathrm{FLim}(f,c,\alpha) \to \mathrm{RFn}((l f))) \to \varphi_{187}$}
\BIC{$\varphi_{187}$}
\end{ptree}
\noindent$\Pi_{80}$:
\begin{ptree}
\AXC{$\stackrel{\Pi_{75}}{\vdots}$}\noLine
\UIC{$l \in \mathbb{R} \to \mathrm{FLim}(f,c,\alpha) \to l \cdot \alpha \in \mathbb{R}$}
\AXC{$l \in \mathbb{R} \to \mathrm{FLim}(f,c,\alpha) \to \mathrm{FLim}(f,c,\alpha)$\lbl{A1}}
\AXC{$\stackrel{\Pi_{77}}{\vdots}$}\noLine
\UIC{$\varphi_{39}$}
\BIC{$l \in \mathbb{R} \to \mathrm{FLim}(f,c,\alpha) \to c \in \mathbb{R}$}
\AXC{$\stackrel{\Pi_{79}}{\vdots}$}\noLine
\UIC{$\varphi_{187}$}
\BIC{$(l \in \mathbb{R} \to \mathrm{FLim}(f,c,\alpha) \to l \cdot \alpha \in \mathbb{R}) \to \varphi_{176}$}
\BIC{$\varphi_{176}$}
\end{ptree}
\noindent 主木:
\begin{ptree}
\AXC{$\stackrel{\Pi_{56}}{\vdots}$}\noLine
\UIC{$l \in \mathbb{R} \to \varphi_{108}$}
\AXC{$\varphi_{2} \to \forall r\,(\varphi_{2})$\lbl{A4}}
\AXC{$\stackrel{\Pi_{71}}{\vdots}$}\noLine
\UIC{$\varphi_{126}$}
\BIC{$l \in \mathbb{R} \to \mathrm{FLim}(f,c,\alpha) \to \mathrm{PNbhd}((l f),c)$}
\AXC{$\stackrel{\Pi_{80}}{\vdots}$}\noLine
\UIC{$\varphi_{176}$}
\BIC{$(l \in \mathbb{R} \to \varphi_{108}) \to \varphi_{5}$}
\BIC{$\varphi_{5}$}
\UIC{$\forall \alpha\,(\varphi_{5})$}
\UIC{$\forall c\,\forall \alpha\,(\varphi_{5})$}
\UIC{$\forall f\,\forall c\,\forall \alpha\,(\varphi_{5})$}
\UIC{$\forall l\,\forall f\,\forall c\,\forall \alpha\,(\varphi_{5})$}
\end{ptree}
\begin{abbrevlist}
\item $\varphi_{1} :\equiv \mathrm{FLim}(f,c,\alpha) \to \mathrm{Seq}(a) \to \mathrm{Conv}(a,c) \to \mathrm{Conv}(a,c)$
\item $\varphi_{2} :\equiv l \in \mathbb{R} \land \mathrm{FLim}(f,c,\alpha) \land \neg \mathrm{PNbhd}((l f),c)$
\item $\varphi_{3} :\equiv \forall \alpha\,(\mathrm{FLim}(f,c,\alpha) \to \mathrm{RFn}(f)) \to \mathrm{FLim}(f,c,\alpha) \to \mathrm{RFn}(f)$
\item $\varphi_{4} :\equiv \mathrm{FLim}(f,c,\alpha) \to (x \in \mathrm{dom}((l f)) \leftrightarrow x \in \mathrm{dom}(f))$
\item $\varphi_{5} :\equiv l \in \mathbb{R} \to \mathrm{FLim}(f,c,\alpha) \to \mathrm{FLim}((l f),c,l \cdot \alpha)$
\item $\varphi_{6} :\equiv \mathrm{FLim}(f,c,\alpha) \to \forall x\,(x \in \mathrm{dom}((l f)) \leftrightarrow x \in \mathrm{dom}(f))$
\item $\varphi_{7} :\equiv l \in \mathbb{R} \land \mathrm{FLim}(f,c,\alpha) \to \forall a\,(l \in \mathbb{R} \land \mathrm{FLim}(f,c,\alpha))$
\item $\varphi_{8} :\equiv \forall c\,\forall \alpha\,(\mathrm{FLim}(f,c,\alpha) \to \mathrm{RFn}(f)) \to \forall \alpha\,(\mathrm{FLim}(f,c,\alpha) \to \mathrm{RFn}(f))$
\item $\varphi_{9} :\equiv \forall a\,\forall \alpha\,(l \in \mathbb{R} \to \mathrm{Conv}(a,\alpha) \to \mathrm{Conv}(l \odot a,l \cdot \alpha))$
\item $\varphi_{10} :\equiv \forall \alpha\,(\mathrm{FLim}(f,c,\alpha) \to \mathrm{PNbhd}(f,c)) \to \mathrm{FLim}(f,c,\alpha) \to \mathrm{PNbhd}(f,c)$
\item $\varphi_{11} :\equiv \forall \alpha\,(\mathrm{FLim}(f,c,\alpha) \to c \in \mathbb{R}) \to \mathrm{FLim}(f,c,\alpha) \to c \in \mathbb{R}$
\item $\varphi_{12} :\equiv \forall c\,\forall a\,\forall \alpha\,(c \in \mathbb{R} \to \mathrm{Conv}(a,\alpha) \to \mathrm{Conv}(c \odot a,c \cdot \alpha))$
\item $\varphi_{13} :\equiv \forall \alpha\,(\mathrm{FLim}(f,c,\alpha) \to \alpha \in \mathbb{R}) \to \mathrm{FLim}(f,c,\alpha) \to \alpha \in \mathbb{R}$
\item $\varphi_{14} :\equiv \forall f\,\forall c\,\forall \alpha\,(\mathrm{FLim}(f,c,\alpha) \to \mathrm{RFn}(f)) \to \forall c\,\forall \alpha\,(\mathrm{FLim}(f,c,\alpha) \to \mathrm{RFn}(f))$
\item $\varphi_{15} :\equiv \forall c\,\forall \alpha\,(\mathrm{FLim}(f,c,\alpha) \to \mathrm{PNbhd}(f,c)) \to \forall \alpha\,(\mathrm{FLim}(f,c,\alpha) \to \mathrm{PNbhd}(f,c))$
\item $\varphi_{16} :\equiv \forall c\,\forall \alpha\,(\mathrm{FLim}(f,c,\alpha) \to c \in \mathbb{R}) \to \forall \alpha\,(\mathrm{FLim}(f,c,\alpha) \to c \in \mathbb{R})$
\item $\varphi_{17} :\equiv \mathrm{FLim}(f,c,\alpha) \to (a_{n} \in \mathrm{dom}((l f)) \leftrightarrow a_{n} \in \mathrm{dom}(f))$
\item $\varphi_{18} :\equiv \forall c\,\forall \alpha\,(\mathrm{FLim}(f,c,\alpha) \to \alpha \in \mathbb{R}) \to \forall \alpha\,(\mathrm{FLim}(f,c,\alpha) \to \alpha \in \mathbb{R})$
\item $\varphi_{19} :\equiv \forall x \in \mathbb{R}\,(0 < \lvert x - c \rvert \to \lvert x - c \rvert < r \to x \in \mathrm{dom}(f))$
\item $\varphi_{20} :\equiv \forall y\,\forall z\,((l f)(a_{n}) = y \to y = z \to (l f)(a_{n}) = z)$
\item $\varphi_{21} :\equiv \forall f\,\forall c\,\forall \alpha\,(\mathrm{FLim}(f,c,\alpha) \to \mathrm{PNbhd}(f,c)) \to \forall c\,\forall \alpha\,(\mathrm{FLim}(f,c,\alpha) \to \mathrm{PNbhd}(f,c))$
\item $\varphi_{22} :\equiv \forall f\,\forall c\,\forall \alpha\,(\mathrm{FLim}(f,c,\alpha) \to c \in \mathbb{R}) \to \forall c\,\forall \alpha\,(\mathrm{FLim}(f,c,\alpha) \to c \in \mathbb{R})$
\item $\varphi_{23} :\equiv l \in \mathbb{R} \to \mathrm{Conv}(f \circ a,\alpha) \to \mathrm{Conv}(l \odot f \circ a,l \cdot \alpha)$
\item $\varphi_{24} :\equiv \forall f\,\forall c\,\forall \alpha\,(\mathrm{FLim}(f,c,\alpha) \to \alpha \in \mathbb{R}) \to \forall c\,\forall \alpha\,(\mathrm{FLim}(f,c,\alpha) \to \alpha \in \mathbb{R})$
\item $\varphi_{25} :\equiv \forall n \in \mathbb{N}\,(a_{n} \in \mathrm{dom}(f) \land a_{n} \neq c) \to \mathrm{Conv}(f \circ a,\alpha)$
\item $\varphi_{26} :\equiv \mathrm{Conv}(l \odot f \circ a,l \cdot \alpha) \to \mathrm{Conv}((l f) \circ a,l \cdot \alpha)$
\item $\varphi_{27} :\equiv x \in \mathbb{R} \to 0 < \lvert x - c \rvert \to \lvert x - c \rvert < r \to x \in \mathrm{dom}((l f))$
\item $\varphi_{28} :\equiv \forall y\,\forall z\,(l \odot f \circ a(n) = y \to y = z \to l \odot f \circ a(n) = z)$
\item $\varphi_{29} :\equiv \mathrm{FLim}(f,c,\alpha) \to \forall x \in \mathrm{dom}((l f))\,((l f)(x) = l \cdot f(x))$
\item $\varphi_{30} :\equiv \forall w\,(z = w \to l \cdot f \circ a(n) = l \cdot z \to l \cdot f \circ a(n) = l \cdot w)$
\item $\varphi_{31} :\equiv \forall a\,\forall \alpha\,(\mathrm{Conv}(a,\alpha) \to \mathrm{Seq}(a)) \to \forall \alpha\,(\mathrm{Conv}(f \circ a,\alpha) \to \mathrm{Seq}(f \circ a))$
\item $\varphi_{32} :\equiv \forall y\,\forall z\,((l f) \circ a(n) = y \to y = z \to (l f) \circ a(n) = z)$
\item $\varphi_{33} :\equiv \forall x \in \mathbb{R}\,(0 < \lvert x - c \rvert \to \lvert x - c \rvert < r \to x \in \mathrm{dom}((l f)))$
\item $\varphi_{34} :\equiv (l \in \mathbb{R} \to \mathrm{FLim}(f,c,\alpha) \to \mathrm{FLim}(f,c,\alpha)) \to l \in \mathbb{R} \to \mathrm{FLim}(f,c,\alpha) \to \mathrm{RFn}(f)$
\item $\varphi_{35} :\equiv \forall a\,\forall b\,\forall \alpha\,(\mathrm{Seq}(b) \to \forall n \in \mathbb{N}\,(a_{n} = b_{n}) \to \mathrm{Conv}(a,\alpha) \to \mathrm{Conv}(b,\alpha))$
\item $\varphi_{36} :\equiv \forall n \in \mathbb{N}\,(a_{n} \in \mathrm{dom}((l f)) \land a_{n} \neq c) \to \mathrm{Seq}(f \circ a)$
\item $\varphi_{37} :\equiv \forall x\,(x \in \mathrm{dom}((l f)) \leftrightarrow x \in \mathrm{dom}(f)) \to (x \in \mathrm{dom}((l f)) \leftrightarrow x \in \mathrm{dom}(f))$
\item $\varphi_{38} :\equiv (l \in \mathbb{R} \to \mathrm{FLim}(f,c,\alpha) \to \mathrm{FLim}(f,c,\alpha)) \to l \in \mathbb{R} \to \mathrm{FLim}(f,c,\alpha) \to \mathrm{PNbhd}(f,c)$
\item $\varphi_{39} :\equiv (l \in \mathbb{R} \to \mathrm{FLim}(f,c,\alpha) \to \mathrm{FLim}(f,c,\alpha)) \to l \in \mathbb{R} \to \mathrm{FLim}(f,c,\alpha) \to c \in \mathbb{R}$
\item $\varphi_{40} :\equiv \forall n \in \mathbb{N}\,(a_{n} \in \mathrm{dom}((l f)) \land a_{n} \neq c) \to n \in \mathbb{N} \to n \in \mathbb{N}$
\item $\varphi_{41} :\equiv \forall n \in \mathbb{N}\,(a_{n} \in \mathrm{dom}((l f)) \land a_{n} \neq c) \to \mathrm{Conv}(f \circ a,\alpha)$
\item $\varphi_{42} :\equiv \forall x\,\forall y\,(x = y \to y = x) \to \forall y\,(l \odot f \circ a(n) = y \to y = l \odot f \circ a(n))$
\item $\varphi_{43} :\equiv (l \in \mathbb{R} \to \mathrm{FLim}(f,c,\alpha) \to \mathrm{FLim}(f,c,\alpha)) \to l \in \mathbb{R} \to \mathrm{FLim}(f,c,\alpha) \to \alpha \in \mathbb{R}$
\item $\varphi_{44} :\equiv n \in \mathbb{N} \to \mathrm{Seq}(f \circ a) \land \forall n \in \mathbb{N}\,(f \circ a(n) = f(a_{n}))$
\item $\varphi_{45} :\equiv \lvert x - c \rvert < r \to x \in \mathbb{R} \to 0 < \lvert x - c \rvert \to \lvert x - c \rvert < r \to x \in \mathrm{dom}(f)$
\item $\varphi_{46} :\equiv \forall x\,\forall y\,(x = y \to y = x) \to \forall y\,((l f) \circ a(n) = y \to y = (l f) \circ a(n))$
\item $\varphi_{47} :\equiv (l \in \mathbb{R} \to \mathrm{FLim}(f,c,\alpha) \to \alpha \in \mathbb{R}) \to l \in \mathbb{R} \to \mathrm{FLim}(f,c,\alpha) \to l \cdot \alpha \in \mathbb{R}$
\item $\varphi_{48} :\equiv \forall n \in \mathbb{N}\,(a_{n} \in \mathrm{dom}((l f)) \land a_{n} \neq c) \to \mathrm{Seq}((l f) \circ a)$
\item $\varphi_{49} :\equiv \forall a\,(l \in \mathbb{R} \to \mathrm{Seq}(a) \to \mathrm{Seq}(l \odot a) \land \forall n \in \mathbb{N}\,(l \odot a(n) = l \cdot a_{n}))$
\item $\varphi_{50} :\equiv \forall \alpha\,(\mathrm{Conv}(f \circ a,\alpha) \to \mathrm{Seq}(f \circ a)) \to \mathrm{Conv}(f \circ a,\alpha) \to \mathrm{Seq}(f \circ a)$
\item $\varphi_{51} :\equiv \forall c\,\forall a\,(c \in \mathbb{R} \to \mathrm{Seq}(a) \to \mathrm{Seq}(c \odot a) \land \forall n \in \mathbb{N}\,(c \odot a(n) = c \cdot a_{n}))$
\item $\varphi_{52} :\equiv \forall y\,(l \in \mathbb{R} \to y \in \mathbb{R} \to l \cdot y \in \mathbb{R}) \to l \in \mathbb{R} \to \alpha \in \mathbb{R} \to l \cdot \alpha \in \mathbb{R}$
\item $\varphi_{53} :\equiv \forall n \in \mathbb{N}\,(a_{n} \in \mathrm{dom}((l f)) \land a_{n} \neq c) \to n \in \mathbb{N} \to a_{n} \in \mathrm{dom}(f)$
\item $\varphi_{54} :\equiv \mathrm{FLim}(f,c,\alpha) \to a_{n} \in \mathrm{dom}((l f)) \to (l f)(a_{n}) = l \cdot f(a_{n})$
\item $\varphi_{55} :\equiv \forall x\,(x \in \mathrm{dom}((l f)) \leftrightarrow x \in \mathrm{dom}(f)) \to (a_{n} \in \mathrm{dom}((l f)) \leftrightarrow a_{n} \in \mathrm{dom}(f))$
\item $\varphi_{56} :\equiv \forall x\,\forall y\,(x \in \mathbb{R} \to y \in \mathbb{R} \to x \cdot y \in \mathbb{R}) \to \forall y\,(l \in \mathbb{R} \to y \in \mathbb{R} \to l \cdot y \in \mathbb{R})$
\item $\varphi_{57} :\equiv \forall n \in \mathbb{N}\,(a_{n} \in \mathrm{dom}((l f)) \land a_{n} \neq c) \to \mathrm{Conv}(l \odot f \circ a,l \cdot \alpha)$
\item $\varphi_{58} :\equiv \forall n \in \mathbb{N}\,(a_{n} \in \mathrm{dom}((l f)) \land a_{n} \neq c) \to \forall n \in \mathbb{N}\,(a_{n} \in \mathrm{dom}(f))$
\item $\varphi_{59} :\equiv \mathrm{Seq}(l \odot f \circ a) \land \forall n \in \mathbb{N}\,(l \odot f \circ a(n) = l \cdot f \circ a(n))$
\item $\varphi_{60} :\equiv \forall n \in \mathbb{N}\,(a_{n} \in \mathrm{dom}((l f)) \land a_{n} \neq c) \to \mathrm{Conv}((l f) \circ a,l \cdot \alpha)$
\item $\varphi_{61} :\equiv l \cdot f \circ a(n) = l \cdot f(a_{n}) \to l \odot f \circ a(n) = l \cdot f(a_{n})$
\item $\varphi_{62} :\equiv l \cdot f \circ a(n) = l \cdot f \circ a(n) \to l \cdot f \circ a(n) = l \cdot f(a_{n})$
\item $\varphi_{63} :\equiv \mathrm{Seq}((l f) \circ a) \land \forall n \in \mathbb{N}\,((l f) \circ a(n) = (l f)(a_{n}))$
\item $\varphi_{64} :\equiv l \odot f \circ a(n) = l \cdot f(a_{n}) \to l \cdot f(a_{n}) = l \odot f \circ a(n)$
\item $\varphi_{65} :\equiv \forall n \in \mathbb{N}\,(f \circ a(n) = f(a_{n})) \to n \in \mathbb{N} \to f \circ a(n) = f(a_{n})$
\item $\varphi_{66} :\equiv \forall n \in \mathbb{N}\,(a_{n} \in \mathrm{dom}((l f)) \land a_{n} \neq c) \to n \in \mathbb{N} \to a_{n} \in \mathrm{dom}((l f))$
\item $\varphi_{67} :\equiv \forall z\,((l f)(a_{n}) = l \cdot f(a_{n}) \to l \cdot f(a_{n}) = z \to (l f)(a_{n}) = z)$
\item $\varphi_{68} :\equiv l \cdot f(a_{n}) = l \odot f \circ a(n) \to (l f)(a_{n}) = l \odot f \circ a(n)$
\item $\varphi_{69} :\equiv \forall n \in \mathbb{N}\,(a_{n} \in \mathrm{dom}((l f)) \land a_{n} \neq c) \to \forall n \in \mathbb{N}\,(a_{n} \in \mathrm{dom}((l f)))$
\item $\varphi_{70} :\equiv \forall n \in \mathbb{N}\,(a_{n} \in \mathrm{dom}((l f)) \land a_{n} \neq c) \to n \in \mathbb{N} \to a_{n} \in \mathrm{dom}(f) \land a_{n} \neq c$
\item $\varphi_{71} :\equiv l \in \mathbb{R} \land \mathrm{FLim}(f,c,\alpha) \land \mathrm{Seq}(a) \land \mathrm{Conv}(a,c) \land \forall n \in \mathbb{N}\,(a_{n} \in \mathrm{dom}((l f)) \land a_{n} \neq c)$
\item $\varphi_{72} :\equiv \forall n \in \mathbb{N}\,(a_{n} \in \mathrm{dom}(f)) \to \mathrm{Seq}(f \circ a) \land \forall n \in \mathbb{N}\,(f \circ a(n) = f(a_{n}))$
\item $\varphi_{73} :\equiv \forall n \in \mathbb{N}\,(a_{n} \in \mathrm{dom}((l f)) \land a_{n} \neq c) \to n \in \mathbb{N} \to f \circ a(n) = f(a_{n})$
\item $\varphi_{74} :\equiv \forall n \in \mathbb{N}\,(a_{n} \in \mathrm{dom}((l f)) \land a_{n} \neq c) \to \forall n \in \mathbb{N}\,(a_{n} \in \mathrm{dom}(f) \land a_{n} \neq c)$
\item $\varphi_{75} :\equiv \forall c\,\forall \alpha\,\forall a\,(\mathrm{FLim}(f,c,\alpha) \to \mathrm{Conv}(a,c) \to \varphi_{25})$
\item $\varphi_{76} :\equiv \forall w\,(f \circ a(n) = w \to l \cdot f \circ a(n) = l \cdot f \circ a(n) \to l \cdot f \circ a(n) = l \cdot w)$
\item $\varphi_{77} :\equiv \forall f\,\forall c\,(\mathrm{PNbhd}(f,c) \leftrightarrow \exists r\,(r \in \mathbb{R} \land 0 < r \land \varphi_{19}))$
\item $\varphi_{78} :\equiv (l f)(a_{n}) = l \odot f \circ a(n) \to (l f) \circ a(n) = l \odot f \circ a(n)$
\item $\varphi_{79} :\equiv \forall b\,\forall \alpha\,(\mathrm{Seq}(b) \to \forall n \in \mathbb{N}\,(l \odot f \circ a(n) = b_{n}) \to \mathrm{Conv}(l \odot f \circ a,\alpha) \to \mathrm{Conv}(b,\alpha))$
\item $\varphi_{80} :\equiv \forall n \in \mathbb{N}\,(a_{n} \in \mathrm{dom}((l f)) \land a_{n} \neq c) \to n \in \mathbb{N} \to a_{n} \in \mathrm{dom}((l f)) \land a_{n} \neq c$
\item $\varphi_{81} :\equiv (l f) \circ a(n) = l \odot f \circ a(n) \to l \odot f \circ a(n) = (l f) \circ a(n)$
\item $\varphi_{82} :\equiv \forall z\,(l \odot f \circ a(n) = l \cdot f \circ a(n) \to l \cdot f \circ a(n) = z \to l \odot f \circ a(n) = z)$
\item $\varphi_{83} :\equiv \forall n \in \mathbb{N}\,(a_{n} \in \mathrm{dom}((l f)) \land a_{n} \neq c) \to \forall n \in \mathbb{N}\,(a_{n} \in \mathrm{dom}((l f)) \land a_{n} \neq c)$
\item $\varphi_{84} :\equiv \forall n \in \mathbb{N}\,(a_{n} \in \mathrm{dom}((l f)) \land a_{n} \neq c) \to n \in \mathbb{N} \to l \cdot f \circ a(n) = l \cdot f(a_{n})$
\item $\varphi_{85} :\equiv \forall n \in \mathbb{N}\,(a_{n} \in \mathrm{dom}((l f)) \land a_{n} \neq c) \to n \in \mathbb{N} \to (l f)(a_{n}) = l \cdot f(a_{n})$
\item $\varphi_{86} :\equiv \forall z\,((l f) \circ a(n) = (l f)(a_{n}) \to (l f)(a_{n}) = z \to (l f) \circ a(n) = z)$
\item $\varphi_{87} :\equiv \mathrm{RFn}((l f)) \land \forall x\,(x \in \mathrm{dom}((l f)) \leftrightarrow x \in \mathrm{dom}(f)) \land \forall x \in \mathrm{dom}((l f))\,((l f)(x) = l \cdot f(x))$
\item $\varphi_{88} :\equiv l \in \mathbb{R} \to \mathrm{FLim}(f,c,\alpha) \to \exists r\,(r \in \mathbb{R} \land 0 < r \land \varphi_{19})$
\item $\varphi_{89} :\equiv \forall n \in \mathbb{N}\,(a_{n} \in \mathrm{dom}((l f)) \land a_{n} \neq c) \to n \in \mathbb{N} \to n \in \mathbb{N} \to f \circ a(n) = f(a_{n})$
\item $\varphi_{90} :\equiv \forall x\,(l \in \mathbb{R} \land \mathrm{FLim}(f,c,\alpha) \land (r \in \mathbb{R} \land 0 < r \land \varphi_{19}))$
\item $\varphi_{91} :\equiv \forall n \in \mathbb{N}\,(a_{n} \in \mathrm{dom}((l f)) \land a_{n} \neq c) \to n \in \mathbb{N} \to l \odot f \circ a(n) = l \cdot f(a_{n})$
\item $\varphi_{92} :\equiv \forall n \in \mathbb{N}\,(a_{n} \in \mathrm{dom}((l f)) \land a_{n} \neq c) \to n \in \mathbb{N} \to l \cdot f(a_{n}) = l \odot f \circ a(n)$
\item $\varphi_{93} :\equiv \forall n \in \mathbb{N}\,(a_{n} \in \mathrm{dom}((l f)) \land a_{n} \neq c) \to n \in \mathbb{N} \to \forall n \in \mathbb{N}\,(f \circ a(n) = f(a_{n}))$
\item $\varphi_{94} :\equiv \forall n \in \mathbb{N}\,(a_{n} \in \mathrm{dom}((l f)) \land a_{n} \neq c) \to n \in \mathbb{N} \to l \odot f \circ a(n) = l \cdot f \circ a(n)$
\item $\varphi_{95} :\equiv \forall n \in \mathbb{N}\,(a_{n} \in \mathrm{dom}((l f)) \land a_{n} \neq c) \to n \in \mathbb{N} \to (l f)(a_{n}) = l \odot f \circ a(n)$
\item $\varphi_{96} :\equiv \forall c\,(\mathrm{PNbhd}((l f),c) \leftrightarrow \exists r\,(r \in \mathbb{R} \land 0 < r \land \varphi_{33}))$
\item $\varphi_{97} :\equiv \forall n \in \mathbb{N}\,(a_{n} \in \mathrm{dom}((l f)) \land a_{n} \neq c) \to n \in \mathbb{N} \to (l f) \circ a(n) = (l f)(a_{n})$
\item $\varphi_{98} :\equiv \mathrm{FLim}(f,c,\alpha) \to r \in \mathbb{R} \land 0 < r \land \varphi_{19} \to \mathrm{PNbhd}((l f),c)$
\item $\varphi_{99} :\equiv l \in \mathbb{R} \to \mathrm{FLim}(f,c,\alpha) \to \mathrm{Seq}(a) \to \mathrm{Conv}(a,c) \to \varphi_{36}$
\item $\varphi_{100} :\equiv \forall n \in \mathbb{N}\,(a_{n} \in \mathrm{dom}((l f)) \land a_{n} \neq c) \to n \in \mathbb{N} \to (l f) \circ a(n) = l \odot f \circ a(n)$
\item $\varphi_{101} :\equiv \forall n \in \mathbb{N}\,(a_{n} \in \mathrm{dom}((l f)) \land a_{n} \neq c) \to n \in \mathbb{N} \to l \odot f \circ a(n) = (l f) \circ a(n)$
\item $\varphi_{102} :\equiv \mathrm{FLim}(f,c,\alpha) \to r \in \mathbb{R} \land 0 < r \land \varphi_{19} \to x \in \mathbb{R} \to x \in \mathbb{R}$
\item $\varphi_{103} :\equiv l \in \mathbb{R} \to \mathrm{FLim}(f,c,\alpha) \to \mathrm{Seq}(a) \to \mathrm{Conv}(a,c) \to \varphi_{40}$
\item $\varphi_{104} :\equiv l \in \mathbb{R} \to \mathrm{FLim}(f,c,\alpha) \to \mathrm{Seq}(a) \to \mathrm{Conv}(a,c) \to \varphi_{41}$
\item $\varphi_{105} :\equiv \forall n \in \mathbb{N}\,(a_{n} \in \mathrm{dom}((l f)) \land a_{n} \neq c) \to \forall n \in \mathbb{N}\,(l \odot f \circ a(n) = (l f) \circ a(n))$
\item $\varphi_{106} :\equiv \varphi_{19} \to x \in \mathbb{R} \to 0 < \lvert x - c \rvert \to \lvert x - c \rvert < r \to x \in \mathrm{dom}(f)$
\item $\varphi_{107} :\equiv \forall n \in \mathbb{N}\,(l \odot f \circ a(n) = l \cdot f \circ a(n)) \to n \in \mathbb{N} \to l \odot f \circ a(n) = l \cdot f \circ a(n)$
\item $\varphi_{108} :\equiv \mathrm{FLim}(f,c,\alpha) \to \forall a\,(\mathrm{Seq}(a) \to \mathrm{Conv}(a,c) \to \varphi_{60})$
\item $\varphi_{109} :\equiv l \in \mathbb{R} \to \mathrm{FLim}(f,c,\alpha) \to \mathrm{Seq}(a) \to \mathrm{Conv}(a,c) \to \varphi_{48}$
\item $\varphi_{110} :\equiv \forall n \in \mathbb{N}\,((l f) \circ a(n) = (l f)(a_{n})) \to n \in \mathbb{N} \to (l f) \circ a(n) = (l f)(a_{n})$
\item $\varphi_{111} :\equiv \forall n \in \mathbb{N}\,(a_{n} \in \mathrm{dom}((l f)) \land a_{n} \neq c) \to n \in \mathbb{N} \to n \in \mathbb{N} \to l \odot f \circ a(n) = l \cdot f \circ a(n)$
\item $\varphi_{112} :\equiv \forall x \in \mathrm{dom}((l f))\,((l f)(x) = l \cdot f(x)) \to a_{n} \in \mathrm{dom}((l f)) \to (l f)(a_{n}) = l \cdot f(a_{n})$
\item $\varphi_{113} :\equiv \forall n \in \mathbb{N}\,(a_{n} \in \mathrm{dom}((l f)) \land a_{n} \neq c) \to n \in \mathbb{N} \to n \in \mathbb{N} \to (l f) \circ a(n) = (l f)(a_{n})$
\item $\varphi_{114} :\equiv \forall n \in \mathbb{N}\,(a_{n} \in \mathrm{dom}((l f)) \land a_{n} \neq c) \to n \in \mathbb{N} \to \forall n \in \mathbb{N}\,(l \odot f \circ a(n) = l \cdot f \circ a(n))$
\item $\varphi_{115} :\equiv r \in \mathbb{R} \land 0 < r \land \varphi_{19} \to x \in \mathbb{R} \to 0 < \lvert x - c \rvert \to 0 < \lvert x - c \rvert$
\item $\varphi_{116} :\equiv l \in \mathbb{R} \to \mathrm{FLim}(f,c,\alpha) \to \mathrm{Seq}(a) \to \mathrm{Conv}(a,c) \to \varphi_{53}$
\item $\varphi_{117} :\equiv \forall n \in \mathbb{N}\,(a_{n} \in \mathrm{dom}((l f)) \land a_{n} \neq c) \to n \in \mathbb{N} \to \forall n \in \mathbb{N}\,((l f) \circ a(n) = (l f)(a_{n}))$
\item $\varphi_{118} :\equiv l \in \mathbb{R} \to \mathrm{FLim}(f,c,\alpha) \to \mathrm{Seq}(a) \to \mathrm{Conv}(a,c) \to \varphi_{57}$
\item $\varphi_{119} :\equiv l \in \mathbb{R} \to \mathrm{FLim}(f,c,\alpha) \to \mathrm{Seq}(a) \to \mathrm{Conv}(a,c) \to \varphi_{58}$
\item $\varphi_{120} :\equiv l \in \mathbb{R} \to \mathrm{FLim}(f,c,\alpha) \to \mathrm{Seq}(a) \to \mathrm{Conv}(a,c) \to \varphi_{60}$
\item $\varphi_{121} :\equiv l \in \mathbb{R} \land \mathrm{FLim}(f,c,\alpha) \to \mathrm{Seq}(a) \to \mathrm{Conv}(a,c) \to \varphi_{60}$
\item $\varphi_{122} :\equiv \forall a\,(\mathrm{Seq}(a) \to \mathrm{Conv}(a,c) \to \varphi_{60}) \to \mathrm{FLim}((l f),c,l \cdot \alpha)$
\item $\varphi_{123} :\equiv l \in \mathbb{R} \to \mathrm{FLim}(f,c,\alpha) \to \mathrm{Seq}(a) \to \mathrm{Conv}(a,c) \to \varphi_{66}$
\item $\varphi_{124} :\equiv \forall a\,(l \in \mathbb{R} \land \mathrm{FLim}(f,c,\alpha)) \to \forall a\,(\mathrm{Seq}(a) \to \mathrm{Conv}(a,c) \to \varphi_{60})$
\item $\varphi_{125} :\equiv l \in \mathbb{R} \to \mathrm{FLim}(f,c,\alpha) \to \mathrm{Seq}(a) \to \mathrm{Conv}(a,c) \to \varphi_{69}$
\item $\varphi_{126} :\equiv (\varphi_{2} \to \forall r\,(\varphi_{2})) \to l \in \mathbb{R} \to \mathrm{FLim}(f,c,\alpha) \to \mathrm{PNbhd}((l f),c)$
\item $\varphi_{127} :\equiv l \in \mathbb{R} \to \mathrm{FLim}(f,c,\alpha) \to \mathrm{Seq}(a) \to \mathrm{Conv}(a,c) \to \varphi_{70}$
\item $\varphi_{128} :\equiv \mathrm{Conv}(a,c) \to \forall n \in \mathbb{N}\,(a_{n} \in \mathrm{dom}((l f)) \land a_{n} \neq c) \to \varphi_{44}$
\item $\varphi_{129} :\equiv l \in \mathbb{R} \to \mathrm{FLim}(f,c,\alpha) \to \mathrm{Seq}(a) \to \mathrm{Conv}(a,c) \to \varphi_{73}$
\item $\varphi_{130} :\equiv l \in \mathbb{R} \to \mathrm{FLim}(f,c,\alpha) \to \mathrm{Seq}(a) \to \mathrm{Conv}(a,c) \to \varphi_{74}$
\item $\varphi_{131} :\equiv (l \in \mathbb{R} \to \alpha \in \mathbb{R} \to l \cdot \alpha \in \mathbb{R}) \to (l \in \mathbb{R} \to l \in \mathbb{R}) \to \varphi_{47}$
\item $\varphi_{132} :\equiv r \in \mathbb{R} \land 0 < r \land \varphi_{19} \to x \in \mathbb{R} \to 0 < \lvert x - c \rvert \to \lvert x - c \rvert < r \to \lvert x - c \rvert < r$
\item $\varphi_{133} :\equiv l \in \mathbb{R} \to \mathrm{FLim}(f,c,\alpha) \to \mathrm{Seq}(a) \to \mathrm{Conv}(a,c) \to \varphi_{80}$
\item $\varphi_{134} :\equiv l \in \mathbb{R} \to \mathrm{FLim}(f,c,\alpha) \to \mathrm{Seq}(a) \to \mathrm{Conv}(a,c) \to \varphi_{83}$
\item $\varphi_{135} :\equiv l \in \mathbb{R} \to \mathrm{FLim}(f,c,\alpha) \to \mathrm{Seq}(a) \to \mathrm{Conv}(a,c) \to \varphi_{84}$
\item $\varphi_{136} :\equiv l \in \mathbb{R} \to \mathrm{FLim}(f,c,\alpha) \to \mathrm{Seq}(a) \to \mathrm{Conv}(a,c) \to \varphi_{85}$
\item $\varphi_{137} :\equiv \mathrm{Conv}(a,c) \to \forall n \in \mathbb{N}\,(a_{n} \in \mathrm{dom}((l f)) \land a_{n} \neq c) \to \varphi_{59}$
\item $\varphi_{138} :\equiv l \in \mathbb{R} \to \mathrm{FLim}(f,c,\alpha) \to \mathrm{Seq}(a) \to \mathrm{Conv}(a,c) \to \varphi_{89}$
\item $\varphi_{139} :\equiv l \in \mathbb{R} \to \mathrm{FLim}(f,c,\alpha) \to \mathrm{Seq}(a) \to \mathrm{Conv}(a,c) \to \varphi_{91}$
\item $\varphi_{140} :\equiv l \in \mathbb{R} \to \mathrm{FLim}(f,c,\alpha) \to \mathrm{Seq}(a) \to \mathrm{Conv}(a,c) \to \varphi_{92}$
\item $\varphi_{141} :\equiv \forall \alpha\,(\mathrm{Seq}((l f) \circ a) \to \forall n \in \mathbb{N}\,(l \odot f \circ a(n) = (l f) \circ a(n)) \to \mathrm{Conv}(l \odot f \circ a,\alpha) \to \mathrm{Conv}((l f) \circ a,\alpha))$
\item $\varphi_{142} :\equiv \mathrm{Conv}(a,c) \to \forall n \in \mathbb{N}\,(a_{n} \in \mathrm{dom}((l f)) \land a_{n} \neq c) \to \varphi_{63}$
\item $\varphi_{143} :\equiv \mathrm{RFn}((l f)) \to \mathrm{Seq}(a) \to \forall n \in \mathbb{N}\,(a_{n} \in \mathrm{dom}((l f))) \to \varphi_{63}$
\item $\varphi_{144} :\equiv \forall f\,\forall c\,\forall \alpha\,(\mathrm{RFn}(f) \to c \in \mathbb{R} \to \alpha \in \mathbb{R} \to \mathrm{PNbhd}(f,c) \to \forall a\,(\mathrm{Seq}(a) \to \mathrm{Conv}(a,c) \to \varphi_{25}) \to \mathrm{FLim}(f,c,\alpha))$
\item $\varphi_{145} :\equiv l \in \mathbb{R} \to \mathrm{FLim}(f,c,\alpha) \to \mathrm{Seq}(a) \to \mathrm{Conv}(a,c) \to \varphi_{93}$
\item $\varphi_{146} :\equiv l \in \mathbb{R} \to \mathrm{FLim}(f,c,\alpha) \to \mathrm{Seq}(a) \to \mathrm{Conv}(a,c) \to \varphi_{94}$
\item $\varphi_{147} :\equiv l \in \mathbb{R} \to \mathrm{FLim}(f,c,\alpha) \to \mathrm{Seq}(a) \to \mathrm{Conv}(a,c) \to \varphi_{95}$
\item $\varphi_{148} :\equiv l \in \mathbb{R} \to \mathrm{FLim}(f,c,\alpha) \to \mathrm{Seq}(a) \to \mathrm{Conv}(a,c) \to \varphi_{97}$
\item $\varphi_{149} :\equiv \mathrm{Seq}((l f) \circ a) \to \forall n \in \mathbb{N}\,(l \odot f \circ a(n) = (l f) \circ a(n)) \to \varphi_{26}$
\item $\varphi_{150} :\equiv l \in \mathbb{R} \to \mathrm{FLim}(f,c,\alpha) \to \mathrm{Seq}(a) \to \mathrm{Conv}(a,c) \to \varphi_{100}$
\item $\varphi_{151} :\equiv l \in \mathbb{R} \to \mathrm{FLim}(f,c,\alpha) \to \mathrm{Seq}(a) \to \mathrm{Conv}(a,c) \to \varphi_{101}$
\item $\varphi_{152} :\equiv l \in \mathbb{R} \to \mathrm{FLim}(f,c,\alpha) \to \mathrm{Seq}(a) \to \mathrm{Conv}(a,c) \to \varphi_{105}$
\item $\varphi_{153} :\equiv (l \in \mathbb{R} \to \mathrm{FLim}(f,c,\alpha) \to \mathrm{RFn}(f)) \to l \in \mathbb{R} \to \mathrm{FLim}(f,c,\alpha) \to \varphi_{87}$
\item $\varphi_{154} :\equiv \forall n\,(\varphi_{71}) \to \forall n \in \mathbb{N}\,(l \odot f \circ a(n) = (l f) \circ a(n))$
\item $\varphi_{155} :\equiv l \in \mathbb{R} \to \mathrm{FLim}(f,c,\alpha) \to \mathrm{Seq}(a) \to \mathrm{Conv}(a,c) \to \varphi_{111}$
\item $\varphi_{156} :\equiv l \in \mathbb{R} \to \mathrm{FLim}(f,c,\alpha) \to \mathrm{Seq}(a) \to \mathrm{Conv}(a,c) \to \varphi_{113}$
\item $\varphi_{157} :\equiv l \in \mathbb{R} \to \mathrm{FLim}(f,c,\alpha) \to \mathrm{Seq}(a) \to \mathrm{Conv}(a,c) \to \varphi_{114}$
\item $\varphi_{158} :\equiv l \in \mathbb{R} \to \mathrm{FLim}(f,c,\alpha) \to \mathrm{Seq}(a) \to \mathrm{Conv}(a,c) \to \varphi_{117}$
\item $\varphi_{159} :\equiv l \in \mathbb{R} \land \mathrm{FLim}(f,c,\alpha) \land (r \in \mathbb{R} \land 0 < r \land \varphi_{19}) \to \varphi_{27}$
\item $\varphi_{160} :\equiv (l \in \mathbb{R} \to \mathrm{FLim}(f,c,\alpha) \to \varphi_{87}) \to l \in \mathbb{R} \to \mathrm{FLim}(f,c,\alpha) \to \mathrm{RFn}((l f))$
\item $\varphi_{161} :\equiv \mathrm{FLim}(f,c,\alpha) \to r \in \mathbb{R} \land 0 < r \land \varphi_{19} \to r \in \mathbb{R} \land 0 < r \land \varphi_{19}$
\item $\varphi_{162} :\equiv \forall r\,(\neg (r \in \mathbb{R} \land 0 < r \land \varphi_{33})) \to \neg (r \in \mathbb{R} \land 0 < r \land \varphi_{33})$
\item $\varphi_{163} :\equiv \mathrm{FLim}(f,c,\alpha) \to r \in \mathbb{R} \land 0 < r \land \varphi_{19} \to r \in \mathbb{R} \land 0 < r \land \varphi_{33}$
\item $\varphi_{164} :\equiv \forall a\,(\mathrm{FLim}(f,c,\alpha) \to \mathrm{Conv}(a,c) \to \varphi_{25}) \to \mathrm{FLim}(f,c,\alpha) \to \mathrm{Conv}(a,c) \to \varphi_{25}$
\item $\varphi_{165} :\equiv \mathrm{FLim}(f,c,\alpha) \to r \in \mathbb{R} \land 0 < r \land \varphi_{19} \to \exists r\,(r \in \mathbb{R} \land 0 < r \land \varphi_{33})$
\item $\varphi_{166} :\equiv r \in \mathbb{R} \land 0 < r \land \varphi_{19} \to x \in \mathbb{R} \to 0 < \lvert x - c \rvert \to \lvert x - c \rvert < r \to \varphi_{19}$
\item $\varphi_{167} :\equiv \forall \alpha\,(\mathrm{RFn}((l f)) \to c \in \mathbb{R} \to \alpha \in \mathbb{R} \to \mathrm{PNbhd}((l f),c) \to \forall a\,(\mathrm{Seq}(a) \to \mathrm{Conv}(a,c) \to \forall n \in \mathbb{N}\,(a_{n} \in \mathrm{dom}((l f)) \land a_{n} \neq c) \to \mathrm{Conv}((l f) \circ a,\alpha)) \to \mathrm{FLim}((l f),c,\alpha))$
\item $\varphi_{168} :\equiv \forall \alpha\,\forall a\,(\mathrm{FLim}(f,c,\alpha) \to \mathrm{Conv}(a,c) \to \varphi_{25}) \to \forall a\,(\mathrm{FLim}(f,c,\alpha) \to \mathrm{Conv}(a,c) \to \varphi_{25})$
\item $\varphi_{169} :\equiv \forall c\,(\mathrm{PNbhd}(f,c) \leftrightarrow \exists r\,(r \in \mathbb{R} \land 0 < r \land \varphi_{19})) \to (\mathrm{PNbhd}(f,c) \leftrightarrow \exists r\,(r \in \mathbb{R} \land 0 < r \land \varphi_{19}))$
\item $\varphi_{170} :\equiv \varphi_{75} \to \forall \alpha\,\forall a\,(\mathrm{FLim}(f,c,\alpha) \to \mathrm{Conv}(a,c) \to \varphi_{25})$
\item $\varphi_{171} :\equiv \varphi_{77} \to \forall c\,(\mathrm{PNbhd}(f,c) \leftrightarrow \exists r\,(r \in \mathbb{R} \land 0 < r \land \varphi_{19}))$
\item $\varphi_{172} :\equiv \mathrm{RFn}((l f)) \to c \in \mathbb{R} \to l \cdot \alpha \in \mathbb{R} \to \mathrm{PNbhd}((l f),c) \to \varphi_{122}$
\item $\varphi_{173} :\equiv \mathrm{FLim}(f,c,\alpha) \to r \in \mathbb{R} \land 0 < r \land \varphi_{19} \to x \in \mathbb{R} \to 0 < \lvert x - c \rvert \to \varphi_{45}$
\item $\varphi_{174} :\equiv l \in \mathbb{R} \land \mathrm{FLim}(f,c,\alpha) \land (r \in \mathbb{R} \land 0 < r \land \varphi_{19}) \to \varphi_{90}$
\item $\varphi_{175} :\equiv \varphi_{96} \to (\mathrm{PNbhd}((l f),c) \leftrightarrow \exists r\,(r \in \mathbb{R} \land 0 < r \land \varphi_{33}))$
\item $\varphi_{176} :\equiv (l \in \mathbb{R} \to \mathrm{FLim}(f,c,\alpha) \to \mathrm{PNbhd}((l f),c)) \to (l \in \mathbb{R} \to \varphi_{108}) \to \varphi_{5}$
\item $\varphi_{177} :\equiv \forall a\,(\mathrm{RFn}(f) \to \mathrm{Seq}(a) \to \varphi_{72}) \to \mathrm{RFn}(f) \to \mathrm{Seq}(a) \to \varphi_{72}$
\item $\varphi_{178} :\equiv (l \in \mathbb{R} \to \varphi_{4}) \to l \in \mathbb{R} \to \mathrm{FLim}(f,c,\alpha) \to r \in \mathbb{R} \land 0 < r \land \varphi_{19} \to \varphi_{27}$
\item $\varphi_{179} :\equiv \forall f\,\forall a\,(\mathrm{RFn}(f) \to \mathrm{Seq}(a) \to \varphi_{72}) \to \forall a\,(\mathrm{RFn}(f) \to \mathrm{Seq}(a) \to \varphi_{72})$
\item $\varphi_{180} :\equiv (l \in \mathbb{R} \to \mathrm{FLim}(f,c,\alpha) \to \mathrm{PNbhd}(f,c)) \to (\mathrm{PNbhd}(f,c) \leftrightarrow \exists r\,(r \in \mathbb{R} \land 0 < r \land \varphi_{19})) \to \varphi_{88}$
\item $\varphi_{181} :\equiv \forall f\,(l \in \mathbb{R} \to \mathrm{RFn}(f) \to \varphi_{87}) \to l \in \mathbb{R} \to \mathrm{RFn}(f) \to \varphi_{87}$
\item $\varphi_{182} :\equiv \forall l\,\forall f\,(l \in \mathbb{R} \to \mathrm{RFn}(f) \to \varphi_{87}) \to \forall f\,(l \in \mathbb{R} \to \mathrm{RFn}(f) \to \varphi_{87})$
\item $\varphi_{183} :\equiv \forall n\,(\varphi_{71} \to n \in \mathbb{N} \to a_{n} \in \mathrm{dom}(f)) \to \forall n\,(\varphi_{71}) \to \forall n \in \mathbb{N}\,(a_{n} \in \mathrm{dom}(f))$
\item $\varphi_{184} :\equiv \forall r\,(\varphi_{2} \to \neg (r \in \mathbb{R} \land 0 < r \land \varphi_{19})) \to \forall r\,(\varphi_{2}) \to \forall r\,(\neg (r \in \mathbb{R} \land 0 < r \land \varphi_{19}))$
\item $\varphi_{185} :\equiv \forall n\,(\varphi_{71} \to n \in \mathbb{N} \to a_{n} \in \mathrm{dom}((l f))) \to \forall n\,(\varphi_{71}) \to \forall n \in \mathbb{N}\,(a_{n} \in \mathrm{dom}((l f)))$
\item $\varphi_{186} :\equiv \forall n\,(\varphi_{71} \to n \in \mathbb{N} \to a_{n} \in \mathrm{dom}(f) \land a_{n} \neq c) \to \forall n\,(\varphi_{71}) \to \forall n \in \mathbb{N}\,(a_{n} \in \mathrm{dom}(f) \land a_{n} \neq c)$
\item $\varphi_{187} :\equiv (l \in \mathbb{R} \to \mathrm{FLim}(f,c,\alpha) \to c \in \mathbb{R}) \to (l \in \mathbb{R} \to \mathrm{FLim}(f,c,\alpha) \to l \cdot \alpha \in \mathbb{R}) \to \varphi_{176}$
\item $\varphi_{188} :\equiv l \cdot f \circ a(n) = l \cdot f \circ a(n) \to (f \circ a(n) = f(a_{n}) \to \varphi_{62}) \to \varphi_{135}$
\item $\varphi_{189} :\equiv (l \in \mathbb{R} \to l \in \mathbb{R}) \to \varphi_{99} \to l \in \mathbb{R} \to \mathrm{FLim}(f,c,\alpha) \to \mathrm{Seq}(a) \to \varphi_{137}$
\item $\varphi_{190} :\equiv \forall n\,(\varphi_{71} \to n \in \mathbb{N} \to l \odot f \circ a(n) = (l f) \circ a(n)) \to \varphi_{154}$
\item $\varphi_{191} :\equiv (l \in \mathbb{R} \to \mathrm{FLim}(f,c,\alpha) \to \mathrm{Seq}(a) \to \mathrm{Seq}(a)) \to \varphi_{119} \to l \in \mathbb{R} \to \mathrm{FLim}(f,c,\alpha) \to \mathrm{Seq}(a) \to \varphi_{128}$
\item $\varphi_{192} :\equiv (l \in \mathbb{R} \to \mathrm{FLim}(f,c,\alpha) \to r \in \mathbb{R} \land 0 < r \land \varphi_{19} \to \varphi_{27}) \to \varphi_{159}$
\item $\varphi_{193} :\equiv (l \in \mathbb{R} \to \mathrm{FLim}(f,c,\alpha) \to \mathrm{FLim}(f,c,\alpha)) \to (l \in \mathbb{R} \to \varphi_{1}) \to \varphi_{130} \to \varphi_{104}$
\item $\varphi_{194} :\equiv (l \in \mathbb{R} \to \mathrm{FLim}(f,c,\alpha) \to \mathrm{Seq}(a) \to \mathrm{Seq}(a)) \to \varphi_{125} \to l \in \mathbb{R} \to \mathrm{FLim}(f,c,\alpha) \to \mathrm{Seq}(a) \to \varphi_{142}$
\item $\varphi_{195} :\equiv (l \in \mathbb{R} \to \mathrm{FLim}(f,c,\alpha) \to r \in \mathbb{R} \land 0 < r \land \varphi_{19} \to \varphi_{33}) \to l \in \mathbb{R} \to \varphi_{163}$
\item $\varphi_{196} :\equiv \varphi_{174} \to l \in \mathbb{R} \to \mathrm{FLim}(f,c,\alpha) \to r \in \mathbb{R} \land 0 < r \land \varphi_{19} \to \varphi_{33}$
\item $\varphi_{197} :\equiv (\mathrm{PNbhd}((l f),c) \leftrightarrow \exists r\,(r \in \mathbb{R} \land 0 < r \land \varphi_{33})) \to (l \in \mathbb{R} \to \varphi_{165}) \to l \in \mathbb{R} \to \varphi_{98}$
\item $\varphi_{198} :\equiv (\mathrm{RFn}(f) \to \mathrm{Seq}(a) \to \varphi_{72}) \to (l \in \mathbb{R} \to \mathrm{FLim}(f,c,\alpha) \to \mathrm{RFn}(f)) \to \varphi_{191}$
\item $\varphi_{199} :\equiv \varphi_{172} \to (l \in \mathbb{R} \to \mathrm{FLim}(f,c,\alpha) \to \mathrm{RFn}((l f))) \to \varphi_{187}$
\item $\varphi_{200} :\equiv (l \in \mathbb{R} \to \mathrm{FLim}(f,c,\alpha) \to \varphi_{115}) \to (l \in \mathbb{R} \to \mathrm{FLim}(f,c,\alpha) \to \varphi_{132}) \to \varphi_{178}$
\item $\varphi_{201} :\equiv \varphi_{143} \to (l \in \mathbb{R} \to \mathrm{FLim}(f,c,\alpha) \to \mathrm{RFn}((l f))) \to \varphi_{194}$
\item $\varphi_{202} :\equiv \forall n\,(\varphi_{71} \to n \in \mathbb{N} \to a_{n} \in \mathrm{dom}(f)) \to \varphi_{183} \to (\varphi_{71} \to \forall n\,(\varphi_{71})) \to \varphi_{119}$
\item $\varphi_{203} :\equiv \forall n\,(\varphi_{71} \to n \in \mathbb{N} \to a_{n} \in \mathrm{dom}((l f))) \to \varphi_{185} \to (\varphi_{71} \to \forall n\,(\varphi_{71})) \to \varphi_{125}$
\item $\varphi_{204} :\equiv \forall n\,(\varphi_{71} \to n \in \mathbb{N} \to a_{n} \in \mathrm{dom}(f) \land a_{n} \neq c) \to \varphi_{186} \to (\varphi_{71} \to \forall n\,(\varphi_{71})) \to \varphi_{130}$
\item $\varphi_{205} :\equiv \forall n\,(\varphi_{71} \to n \in \mathbb{N} \to l \odot f \circ a(n) = (l f) \circ a(n)) \to \varphi_{190} \to (\varphi_{71} \to \forall n\,(\varphi_{71})) \to \varphi_{152}$
\end{abbrevlist}
\end{fullproof}

% 自動生成 (prover/build.js): FLimDiv  ノード数 1002
\begin{fullproof}{定理1.8(3)}{$\forall f\,\forall g\,\forall c\,\forall \alpha\,\forall \beta\,(\mathrm{FLim}(f,c,\alpha) \to \mathrm{FLim}(g,c,\beta) \to \beta \neq 0 \to \mathrm{FLim}((f / g),c,\alpha \cdot \beta^{-1}))$}{1002}\phantomsection\label{tree:FLimDiv}
\noindent$\Pi_{1}$:
\begin{ptree}
\AXC{$x = x$\lbl{A5}}
\AXC{$x = y \to x = x \to y = x$\lbl{A6}}
\AXC{$(x = y \to x = x \to y = x) \to x = x \to x = y \to y = x$\lbl{A1}}
\BIC{$x = x \to x = y \to y = x$}
\BIC{$x = y \to y = x$}
\UIC{$\forall y\,(x = y \to y = x)$}
\UIC{$\forall x\,\forall y\,(x = y \to y = x)$}
\end{ptree}
\noindent$\Pi_{2}$:
\begin{ptree}
\AXC{$\forall f\,\forall c\,\forall \alpha\,(\mathrm{FLim}(f,c,\alpha) \to \mathrm{RFn}(f))$\lbl{[D2]}}
\AXC{$\varphi_{35}$\lbl{A2}}
\BIC{$\forall c\,\forall \alpha\,(\mathrm{FLim}(g,c,\alpha) \to \mathrm{RFn}(g))$}
\AXC{$\varphi_{22}$\lbl{A2}}
\BIC{$\forall \alpha\,(\mathrm{FLim}(g,c,\alpha) \to \mathrm{RFn}(g))$}
\AXC{$\varphi_{7}$\lbl{A2}}
\BIC{$\mathrm{FLim}(g,c,\beta) \to \mathrm{RFn}(g)$}
\end{ptree}
\noindent$\Pi_{3}$:
\begin{ptree}
\AXC{$\varphi_{2}$\lbl{A1}}
\AXC{$\stackrel{\Pi_{2}}{\vdots}$}\noLine
\UIC{$\mathrm{FLim}(g,c,\beta) \to \mathrm{RFn}(g)$}
\AXC{$(\mathrm{FLim}(g,c,\beta) \to \mathrm{RFn}(g)) \to \varphi_{2} \to \varphi_{4}$\lbl{A1}}
\BIC{$\varphi_{2} \to \varphi_{4}$}
\BIC{$\varphi_{4}$}
\end{ptree}
\noindent$\Pi_{4}$:
\begin{ptree}
\AXC{$\forall f\,\forall c\,\forall \alpha\,(\mathrm{FLim}(f,c,\alpha) \to \mathrm{RFn}(f))$\lbl{[D2]}}
\AXC{$\varphi_{36}$\lbl{A2}}
\BIC{$\forall c\,\forall \alpha\,(\mathrm{FLim}(f,c,\alpha) \to \mathrm{RFn}(f))$}
\AXC{$\varphi_{23}$\lbl{A2}}
\BIC{$\forall \alpha\,(\mathrm{FLim}(f,c,\alpha) \to \mathrm{RFn}(f))$}
\AXC{$\varphi_{8}$\lbl{A2}}
\BIC{$\mathrm{FLim}(f,c,\alpha) \to \mathrm{RFn}(f)$}
\AXC{$\varphi_{141}$\lbl{A1}}
\BIC{$(\mathrm{FLim}(f,c,\alpha) \to \mathrm{FLim}(f,c,\alpha)) \to \varphi_{5}$}
\end{ptree}
\noindent$\Pi_{5}$:
\begin{ptree}
\AXC{$\mathrm{FLim}(f,c,\alpha) \to \mathrm{FLim}(f,c,\alpha)$\lbl{A1}}
\AXC{$\stackrel{\Pi_{4}}{\vdots}$}\noLine
\UIC{$(\mathrm{FLim}(f,c,\alpha) \to \mathrm{FLim}(f,c,\alpha)) \to \varphi_{5}$}
\BIC{$\varphi_{5}$}
\end{ptree}
\noindent$\Pi_{6}$:
\begin{ptree}
\AXC{$\forall f\,\forall g\,(\varphi_{293})$\lbl{[S20]}}
\AXC{$\forall f\,\forall g\,(\varphi_{293}) \to \forall g\,(\varphi_{293})$\lbl{A2}}
\BIC{$\forall g\,(\varphi_{293})$}
\AXC{$\forall g\,(\varphi_{293}) \to \varphi_{293}$\lbl{A2}}
\BIC{$\varphi_{293}$}
\AXC{$\varphi_{293} \to \varphi_{5} \to \varphi_{4} \to \mathrm{FLim}(f,c,\alpha) \to \varphi_{299}$\lbl{A1}}
\BIC{$\varphi_{5} \to \varphi_{4} \to \mathrm{FLim}(f,c,\alpha) \to \varphi_{299}$}
\end{ptree}
\noindent$\Pi_{7}$:
\begin{ptree}
\AXC{$\stackrel{\Pi_{3}}{\vdots}$}\noLine
\UIC{$\varphi_{4}$}
\AXC{$\stackrel{\Pi_{5}}{\vdots}$}\noLine
\UIC{$\varphi_{5}$}
\AXC{$\stackrel{\Pi_{6}}{\vdots}$}\noLine
\UIC{$\varphi_{5} \to \varphi_{4} \to \mathrm{FLim}(f,c,\alpha) \to \varphi_{299}$}
\BIC{$\varphi_{4} \to \mathrm{FLim}(f,c,\alpha) \to \varphi_{299}$}
\BIC{$\mathrm{FLim}(f,c,\alpha) \to \varphi_{299}$}
\end{ptree}
\noindent$\Pi_{8}$:
\begin{ptree}
\AXC{$\stackrel{\Pi_{7}}{\vdots}$}\noLine
\UIC{$\mathrm{FLim}(f,c,\alpha) \to \varphi_{299}$}
\AXC{$(\mathrm{FLim}(f,c,\alpha) \to \varphi_{299}) \to \varphi_{168}$\lbl{A1}}
\BIC{$\varphi_{168}$}
\end{ptree}
\noindent$\Pi_{9}$:
\begin{ptree}
\AXC{$\mathrm{FLim}(f,c,\alpha) \to \varphi_{206}$\lbl{A1}}
\AXC{$\varphi_{184}$\lbl{A2}}
\AXC{$\mathrm{FLim}(f,c,\alpha) \to \varphi_{269}$\lbl{A1}}
\AXC{$\varphi_{472}$\lbl{A1}}
\BIC{$\varphi_{184} \to \mathrm{FLim}(f,c,\alpha) \to \varphi_{263}$}
\BIC{$\mathrm{FLim}(f,c,\alpha) \to \varphi_{263}$}
\AXC{$\varphi_{526}$\lbl{A1}}
\BIC{$(\mathrm{FLim}(f,c,\alpha) \to \varphi_{206}) \to \varphi_{432}$}
\BIC{$\varphi_{432}$}
\end{ptree}
\noindent$\Pi_{10}$:
\begin{ptree}
\AXC{$\forall x\,(\varphi_{42}) \to \varphi_{122}$\lbl{A2}}
\AXC{$\stackrel{\Pi_{8}}{\vdots}$}\noLine
\UIC{$\varphi_{168}$}
\AXC{$\varphi_{168} \to \varphi_{400}$\lbl{A1}}
\BIC{$\varphi_{400}$}
\BIC{$\mathrm{FLim}(f,c,\alpha) \to \mathrm{FLim}(g,c,\beta) \to \beta \neq 0 \to \varphi_{122}$}
\AXC{$\stackrel{\Pi_{9}}{\vdots}$}\noLine
\UIC{$\varphi_{432}$}
\BIC{$\mathrm{FLim}(f,c,\alpha) \to \varphi_{316}$}
\AXC{$(\mathrm{FLim}(f,c,\alpha) \to \varphi_{316}) \to \varphi_{190} \to n \in \mathbb{N} \to \varphi_{81}$\lbl{A1}}
\BIC{$\varphi_{190} \to n \in \mathbb{N} \to \varphi_{81}$}
\UIC{$\forall n\,(\varphi_{190} \to n \in \mathbb{N} \to \varphi_{81})$}
\end{ptree}
\noindent$\Pi_{11}$:
\begin{ptree}
\AXC{$\forall n\,(\varphi_{190} \to n \in \mathbb{N} \to \varphi_{81}) \to \forall n\,(\varphi_{190}) \to \forall n \in \mathbb{N}\,(\varphi_{81})$\lbl{A3}}
\AXC{$\stackrel{\Pi_{10}}{\vdots}$}\noLine
\UIC{$\forall n\,(\varphi_{190} \to n \in \mathbb{N} \to \varphi_{81})$}
\AXC{$\forall n\,(\varphi_{190} \to n \in \mathbb{N} \to \varphi_{81}) \to \varphi_{559}$\lbl{A1}}
\BIC{$\varphi_{559}$}
\BIC{$(\varphi_{190} \to \forall n\,(\varphi_{190})) \to \mathrm{FLim}(f,c,\alpha) \to \varphi_{318}$}
\end{ptree}
\noindent$\Pi_{12}$:
\begin{ptree}
\AXC{$\varphi_{190} \to \forall n\,(\varphi_{190})$\lbl{A4}}
\AXC{$\stackrel{\Pi_{11}}{\vdots}$}\noLine
\UIC{$(\varphi_{190} \to \forall n\,(\varphi_{190})) \to \mathrm{FLim}(f,c,\alpha) \to \varphi_{318}$}
\BIC{$\mathrm{FLim}(f,c,\alpha) \to \varphi_{318}$}
\end{ptree}
\noindent$\Pi_{13}$:
\begin{ptree}
\AXC{$\forall n \in \mathbb{N}\,(\varphi_{81}) \to n \in \mathbb{N} \to \varphi_{81}$\lbl{A2}}
\AXC{$\stackrel{\Pi_{12}}{\vdots}$}\noLine
\UIC{$\mathrm{FLim}(f,c,\alpha) \to \varphi_{318}$}
\AXC{$(\mathrm{FLim}(f,c,\alpha) \to \varphi_{318}) \to \varphi_{456}$\lbl{A1}}
\BIC{$\varphi_{456}$}
\BIC{$\mathrm{FLim}(f,c,\alpha) \to \varphi_{316}$}
\AXC{$\varphi_{486}$\lbl{A1}}
\BIC{$(\mathrm{FLim}(f,c,\alpha) \to \varphi_{206}) \to \mathrm{FLim}(f,c,\alpha) \to \varphi_{217}$}
\end{ptree}
\noindent$\Pi_{14}$:
\begin{ptree}
\AXC{$\varphi_{409}$\lbl{A3}}
\AXC{$\mathrm{FLim}(f,c,\alpha) \to \varphi_{206}$\lbl{A1}}
\AXC{$\stackrel{\Pi_{13}}{\vdots}$}\noLine
\UIC{$(\mathrm{FLim}(f,c,\alpha) \to \varphi_{206}) \to \mathrm{FLim}(f,c,\alpha) \to \varphi_{217}$}
\BIC{$\mathrm{FLim}(f,c,\alpha) \to \varphi_{217}$}
\AXC{$\varphi_{405}$\lbl{A1}}
\BIC{$\varphi_{190} \to n \in \mathbb{N} \to a_{n} \in \mathrm{dom}(g)$}
\UIC{$\forall n\,(\varphi_{190} \to n \in \mathbb{N} \to a_{n} \in \mathrm{dom}(g))$}
\AXC{$\varphi_{555}$\lbl{A1}}
\BIC{$\varphi_{409} \to (\varphi_{190} \to \forall n\,(\varphi_{190})) \to \mathrm{FLim}(f,c,\alpha) \to \varphi_{226}$}
\BIC{$(\varphi_{190} \to \forall n\,(\varphi_{190})) \to \mathrm{FLim}(f,c,\alpha) \to \varphi_{226}$}
\end{ptree}
\noindent$\Pi_{15}$:
\begin{ptree}
\AXC{$\varphi_{190} \to \forall n\,(\varphi_{190})$\lbl{A4}}
\AXC{$\stackrel{\Pi_{14}}{\vdots}$}\noLine
\UIC{$(\varphi_{190} \to \forall n\,(\varphi_{190})) \to \mathrm{FLim}(f,c,\alpha) \to \varphi_{226}$}
\BIC{$\mathrm{FLim}(f,c,\alpha) \to \varphi_{226}$}
\end{ptree}
\noindent$\Pi_{16}$:
\begin{ptree}
\AXC{$\varphi_{13}$\lbl{A1}}
\AXC{$\stackrel{\Pi_{3}}{\vdots}$}\noLine
\UIC{$\varphi_{4}$}
\AXC{$\forall f\,\forall a\,(\mathrm{RFn}(f) \to \mathrm{Seq}(a) \to \varphi_{178})$\lbl{[S8]}}
\AXC{$\varphi_{386}$\lbl{A2}}
\BIC{$\forall a\,(\mathrm{RFn}(g) \to \mathrm{Seq}(a) \to \varphi_{175})$}
\AXC{$\varphi_{384}$\lbl{A2}}
\BIC{$\mathrm{RFn}(g) \to \mathrm{Seq}(a) \to \varphi_{175}$}
\AXC{$\varphi_{511}$\lbl{A1}}
\BIC{$\varphi_{4} \to \varphi_{13} \to (\mathrm{FLim}(f,c,\alpha) \to \varphi_{226}) \to \varphi_{321}$}
\BIC{$\varphi_{13} \to (\mathrm{FLim}(f,c,\alpha) \to \varphi_{226}) \to \varphi_{321}$}
\BIC{$(\mathrm{FLim}(f,c,\alpha) \to \varphi_{226}) \to \varphi_{321}$}
\end{ptree}
\noindent$\Pi_{17}$:
\begin{ptree}
\AXC{$\stackrel{\Pi_{15}}{\vdots}$}\noLine
\UIC{$\mathrm{FLim}(f,c,\alpha) \to \varphi_{226}$}
\AXC{$\stackrel{\Pi_{16}}{\vdots}$}\noLine
\UIC{$(\mathrm{FLim}(f,c,\alpha) \to \varphi_{226}) \to \varphi_{321}$}
\BIC{$\varphi_{321}$}
\AXC{$\varphi_{321} \to \mathrm{FLim}(f,c,\alpha) \to \varphi_{281}$\lbl{A1}}
\BIC{$\mathrm{FLim}(f,c,\alpha) \to \varphi_{281}$}
\AXC{$\varphi_{473}$\lbl{A1}}
\BIC{$\varphi_{149} \to \mathrm{FLim}(f,c,\alpha) \to \varphi_{274}$}
\end{ptree}
\noindent$\Pi_{18}$:
\begin{ptree}
\AXC{$\mathrm{FLim}(f,c,\alpha) \to \varphi_{206}$\lbl{A1}}
\AXC{$\varphi_{149}$\lbl{A2}}
\AXC{$\stackrel{\Pi_{17}}{\vdots}$}\noLine
\UIC{$\varphi_{149} \to \mathrm{FLim}(f,c,\alpha) \to \varphi_{274}$}
\BIC{$\mathrm{FLim}(f,c,\alpha) \to \varphi_{274}$}
\AXC{$\varphi_{482}$\lbl{A1}}
\BIC{$(\mathrm{FLim}(f,c,\alpha) \to \varphi_{206}) \to \mathrm{FLim}(f,c,\alpha) \to \varphi_{251}$}
\BIC{$\mathrm{FLim}(f,c,\alpha) \to \varphi_{251}$}
\end{ptree}
\noindent$\Pi_{19}$:
\begin{ptree}
\AXC{$\stackrel{\Pi_{1}}{\vdots}$}\noLine
\UIC{$\forall x\,\forall y\,(x = y \to y = x)$}
\AXC{$\varphi_{48}$\lbl{A2}}
\BIC{$\forall y\,(g \circ a(n) = y \to y = g \circ a(n))$}
\AXC{$\varphi_{205}$\lbl{A2}}
\BIC{$g \circ a(n) = g(a_{n}) \to g(a_{n}) = g \circ a(n)$}
\AXC{$\stackrel{\Pi_{18}}{\vdots}$}\noLine
\UIC{$\mathrm{FLim}(f,c,\alpha) \to \varphi_{251}$}
\AXC{$(\mathrm{FLim}(f,c,\alpha) \to \varphi_{251}) \to \varphi_{380}$\lbl{A1}}
\BIC{$\varphi_{380}$}
\BIC{$\mathrm{FLim}(f,c,\alpha) \to \varphi_{252}$}
\end{ptree}
\noindent$\Pi_{20}$:
\begin{ptree}
\AXC{$\stackrel{\Pi_{1}}{\vdots}$}\noLine
\UIC{$\forall x\,\forall y\,(x = y \to y = x)$}
\AXC{$\varphi_{29}$\lbl{A2}}
\BIC{$\forall y\,(g(a_{n}) = y \to y = g(a_{n}))$}
\AXC{$\varphi_{192}$\lbl{A2}}
\BIC{$g(a_{n}) = g \circ a(n) \to g \circ a(n) = g(a_{n})$}
\AXC{$\stackrel{\Pi_{19}}{\vdots}$}\noLine
\UIC{$\mathrm{FLim}(f,c,\alpha) \to \varphi_{252}$}
\AXC{$(\mathrm{FLim}(f,c,\alpha) \to \varphi_{252}) \to \varphi_{381}$\lbl{A1}}
\BIC{$\varphi_{381}$}
\BIC{$\mathrm{FLim}(f,c,\alpha) \to \varphi_{251}$}
\end{ptree}
\noindent$\Pi_{21}$:
\begin{ptree}
\AXC{$\forall n \in \mathbb{N}\,(\varphi_{81}) \to n \in \mathbb{N} \to \varphi_{81}$\lbl{A2}}
\AXC{$\stackrel{\Pi_{12}}{\vdots}$}\noLine
\UIC{$\mathrm{FLim}(f,c,\alpha) \to \varphi_{318}$}
\AXC{$(\mathrm{FLim}(f,c,\alpha) \to \varphi_{318}) \to \varphi_{456}$\lbl{A1}}
\BIC{$\varphi_{456}$}
\BIC{$\mathrm{FLim}(f,c,\alpha) \to \varphi_{316}$}
\AXC{$\varphi_{485}$\lbl{A1}}
\BIC{$(\mathrm{FLim}(f,c,\alpha) \to \varphi_{206}) \to \mathrm{FLim}(f,c,\alpha) \to \varphi_{216}$}
\end{ptree}
\noindent$\Pi_{22}$:
\begin{ptree}
\AXC{$\mathrm{FLim}(f,c,\alpha) \to \varphi_{206}$\lbl{A1}}
\AXC{$\stackrel{\Pi_{21}}{\vdots}$}\noLine
\UIC{$(\mathrm{FLim}(f,c,\alpha) \to \varphi_{206}) \to \mathrm{FLim}(f,c,\alpha) \to \varphi_{216}$}
\BIC{$\mathrm{FLim}(f,c,\alpha) \to \varphi_{216}$}
\end{ptree}
\noindent$\Pi_{23}$:
\begin{ptree}
\AXC{$\stackrel{\Pi_{22}}{\vdots}$}\noLine
\UIC{$\mathrm{FLim}(f,c,\alpha) \to \varphi_{216}$}
\AXC{$\stackrel{\Pi_{19}}{\vdots}$}\noLine
\UIC{$\mathrm{FLim}(f,c,\alpha) \to \varphi_{252}$}
\AXC{$\varphi_{546}$\lbl{A1}}
\BIC{$(\mathrm{FLim}(f,c,\alpha) \to \varphi_{216}) \to \varphi_{481}$}
\BIC{$\varphi_{481}$}
\end{ptree}
\noindent$\Pi_{24}$:
\begin{ptree}
\AXC{$z = w \to z = 0 \to w = 0$\lbl{A6}}
\UIC{$\forall w\,(z = w \to z = 0 \to w = 0)$}
\UIC{$\forall z\,\forall w\,(z = w \to z = 0 \to w = 0)$}
\AXC{$\varphi_{58}$\lbl{A2}}
\BIC{$\forall w\,(g(a_{n}) = w \to g(a_{n}) = 0 \to w = 0)$}
\AXC{$\varphi_{211}$\lbl{A2}}
\BIC{$g(a_{n}) = g \circ a(n) \to g(a_{n}) = 0 \to g \circ a(n) = 0$}
\AXC{$\stackrel{\Pi_{23}}{\vdots}$}\noLine
\UIC{$\varphi_{481}$}
\BIC{$\varphi_{460}$}
\end{ptree}
\noindent$\Pi_{25}$:
\begin{ptree}
\AXC{$\stackrel{\Pi_{20}}{\vdots}$}\noLine
\UIC{$\mathrm{FLim}(f,c,\alpha) \to \varphi_{251}$}
\AXC{$z = w \to z = 0 \to w = 0$\lbl{A6}}
\UIC{$\forall w\,(z = w \to z = 0 \to w = 0)$}
\UIC{$\forall z\,\forall w\,(z = w \to z = 0 \to w = 0)$}
\AXC{$\varphi_{88}$\lbl{A2}}
\BIC{$\forall w\,(g \circ a(n) = w \to g \circ a(n) = 0 \to w = 0)$}
\AXC{$\varphi_{225}$\lbl{A2}}
\BIC{$g \circ a(n) = g(a_{n}) \to g \circ a(n) = 0 \to g(a_{n}) = 0$}
\AXC{$\stackrel{\Pi_{24}}{\vdots}$}\noLine
\UIC{$\varphi_{460}$}
\BIC{$(\mathrm{FLim}(f,c,\alpha) \to \varphi_{251}) \to \mathrm{FLim}(f,c,\alpha) \to \varphi_{227}$}
\BIC{$\mathrm{FLim}(f,c,\alpha) \to \varphi_{227}$}
\AXC{$\varphi_{408}$\lbl{A1}}
\BIC{$\varphi_{190} \to n \in \mathbb{N} \to g \circ a(n) \neq 0$}
\UIC{$\forall n\,(\varphi_{190} \to n \in \mathbb{N} \to g \circ a(n) \neq 0)$}
\end{ptree}
\noindent$\Pi_{26}$:
\begin{ptree}
\AXC{$\varphi_{190} \to \forall n\,(\varphi_{190})$\lbl{A4}}
\AXC{$\varphi_{413}$\lbl{A3}}
\AXC{$\stackrel{\Pi_{25}}{\vdots}$}\noLine
\UIC{$\forall n\,(\varphi_{190} \to n \in \mathbb{N} \to g \circ a(n) \neq 0)$}
\AXC{$\varphi_{558}$\lbl{A1}}
\BIC{$\varphi_{413} \to (\varphi_{190} \to \forall n\,(\varphi_{190})) \to \mathrm{FLim}(f,c,\alpha) \to \varphi_{233}$}
\BIC{$(\varphi_{190} \to \forall n\,(\varphi_{190})) \to \mathrm{FLim}(f,c,\alpha) \to \varphi_{233}$}
\BIC{$\mathrm{FLim}(f,c,\alpha) \to \varphi_{233}$}
\end{ptree}
\noindent$\Pi_{27}$:
\begin{ptree}
\AXC{$\forall n \in \mathbb{N}\,(\varphi_{81}) \to n \in \mathbb{N} \to \varphi_{81}$\lbl{A2}}
\AXC{$\stackrel{\Pi_{12}}{\vdots}$}\noLine
\UIC{$\mathrm{FLim}(f,c,\alpha) \to \varphi_{318}$}
\AXC{$(\mathrm{FLim}(f,c,\alpha) \to \varphi_{318}) \to \varphi_{456}$\lbl{A1}}
\BIC{$\varphi_{456}$}
\BIC{$\mathrm{FLim}(f,c,\alpha) \to \varphi_{316}$}
\AXC{$\varphi_{490}$\lbl{A1}}
\BIC{$(\mathrm{FLim}(f,c,\alpha) \to \varphi_{206}) \to \mathrm{FLim}(f,c,\alpha) \to \varphi_{245}$}
\end{ptree}
\noindent$\Pi_{28}$:
\begin{ptree}
\AXC{$\varphi_{430}$\lbl{A3}}
\AXC{$\mathrm{FLim}(f,c,\alpha) \to \varphi_{206}$\lbl{A1}}
\AXC{$\stackrel{\Pi_{27}}{\vdots}$}\noLine
\UIC{$(\mathrm{FLim}(f,c,\alpha) \to \varphi_{206}) \to \mathrm{FLim}(f,c,\alpha) \to \varphi_{245}$}
\BIC{$\mathrm{FLim}(f,c,\alpha) \to \varphi_{245}$}
\AXC{$\varphi_{424}$\lbl{A1}}
\BIC{$\varphi_{190} \to n \in \mathbb{N} \to a_{n} \in \mathrm{dom}(g) \land a_{n} \neq c$}
\UIC{$\forall n\,(\varphi_{190} \to n \in \mathbb{N} \to a_{n} \in \mathrm{dom}(g) \land a_{n} \neq c)$}
\AXC{$\varphi_{562}$\lbl{A1}}
\BIC{$\varphi_{430} \to (\varphi_{190} \to \forall n\,(\varphi_{190})) \to \mathrm{FLim}(f,c,\alpha) \to \varphi_{255}$}
\BIC{$(\varphi_{190} \to \forall n\,(\varphi_{190})) \to \mathrm{FLim}(f,c,\alpha) \to \varphi_{255}$}
\end{ptree}
\noindent$\Pi_{29}$:
\begin{ptree}
\AXC{$\varphi_{190} \to \forall n\,(\varphi_{190})$\lbl{A4}}
\AXC{$\stackrel{\Pi_{28}}{\vdots}$}\noLine
\UIC{$(\varphi_{190} \to \forall n\,(\varphi_{190})) \to \mathrm{FLim}(f,c,\alpha) \to \varphi_{255}$}
\BIC{$\mathrm{FLim}(f,c,\alpha) \to \varphi_{255}$}
\end{ptree}
\noindent$\Pi_{30}$:
\begin{ptree}
\AXC{$\forall f\,(\varphi_{181})$\lbl{定理1.7(⇒)}}
\AXC{$\forall f\,(\varphi_{181}) \to \forall c\,(\varphi_{169})$\lbl{A2}}
\BIC{$\forall c\,(\varphi_{169})$}
\AXC{$\forall c\,(\varphi_{169}) \to \varphi_{169}$\lbl{A2}}
\BIC{$\varphi_{169}$}
\AXC{$\varphi_{169} \to \forall a\,(\mathrm{FLim}(g,c,\beta) \to \mathrm{Conv}(a,c) \to \varphi_{71})$\lbl{A2}}
\BIC{$\forall a\,(\mathrm{FLim}(g,c,\beta) \to \mathrm{Conv}(a,c) \to \varphi_{71})$}
\end{ptree}
\noindent$\Pi_{31}$:
\begin{ptree}
\AXC{$\stackrel{\Pi_{30}}{\vdots}$}\noLine
\UIC{$\forall a\,(\mathrm{FLim}(g,c,\beta) \to \mathrm{Conv}(a,c) \to \varphi_{71})$}
\AXC{$\varphi_{327}$\lbl{A2}}
\BIC{$\mathrm{FLim}(g,c,\beta) \to \mathrm{Conv}(a,c) \to \varphi_{71}$}
\AXC{$(\mathrm{FLim}(g,c,\beta) \to \mathrm{Conv}(a,c) \to \varphi_{71}) \to \varphi_{2} \to \varphi_{450}$\lbl{A1}}
\BIC{$\varphi_{2} \to \varphi_{450}$}
\end{ptree}
\noindent$\Pi_{32}$:
\begin{ptree}
\AXC{$\stackrel{\Pi_{29}}{\vdots}$}\noLine
\UIC{$\mathrm{FLim}(f,c,\alpha) \to \varphi_{255}$}
\AXC{$\mathrm{FLim}(f,c,\alpha) \to \varphi_{14}$\lbl{A1}}
\AXC{$\varphi_{2}$\lbl{A1}}
\AXC{$\stackrel{\Pi_{31}}{\vdots}$}\noLine
\UIC{$\varphi_{2} \to \varphi_{450}$}
\BIC{$\varphi_{450}$}
\BIC{$(\mathrm{FLim}(f,c,\alpha) \to \varphi_{255}) \to \mathrm{FLim}(f,c,\alpha) \to \varphi_{207}$}
\BIC{$\mathrm{FLim}(f,c,\alpha) \to \varphi_{207}$}
\end{ptree}
\noindent$\Pi_{33}$:
\begin{ptree}
\AXC{$\forall n \in \mathbb{N}\,(\varphi_{81}) \to n \in \mathbb{N} \to \varphi_{81}$\lbl{A2}}
\AXC{$\stackrel{\Pi_{12}}{\vdots}$}\noLine
\UIC{$\mathrm{FLim}(f,c,\alpha) \to \varphi_{318}$}
\AXC{$(\mathrm{FLim}(f,c,\alpha) \to \varphi_{318}) \to \varphi_{456}$\lbl{A1}}
\BIC{$\varphi_{456}$}
\BIC{$\mathrm{FLim}(f,c,\alpha) \to \varphi_{316}$}
\AXC{$\varphi_{491}$\lbl{A1}}
\BIC{$(\mathrm{FLim}(f,c,\alpha) \to \varphi_{206}) \to \mathrm{FLim}(f,c,\alpha) \to \varphi_{246}$}
\end{ptree}
\noindent$\Pi_{34}$:
\begin{ptree}
\AXC{$\varphi_{431}$\lbl{A3}}
\AXC{$\mathrm{FLim}(f,c,\alpha) \to \varphi_{206}$\lbl{A1}}
\AXC{$\stackrel{\Pi_{33}}{\vdots}$}\noLine
\UIC{$(\mathrm{FLim}(f,c,\alpha) \to \varphi_{206}) \to \mathrm{FLim}(f,c,\alpha) \to \varphi_{246}$}
\BIC{$\mathrm{FLim}(f,c,\alpha) \to \varphi_{246}$}
\AXC{$\varphi_{425}$\lbl{A1}}
\BIC{$\varphi_{190} \to n \in \mathbb{N} \to a_{n} \in \mathrm{dom}(f) \land a_{n} \neq c$}
\UIC{$\forall n\,(\varphi_{190} \to n \in \mathbb{N} \to a_{n} \in \mathrm{dom}(f) \land a_{n} \neq c)$}
\AXC{$\varphi_{563}$\lbl{A1}}
\BIC{$\varphi_{431} \to (\varphi_{190} \to \forall n\,(\varphi_{190})) \to \mathrm{FLim}(f,c,\alpha) \to \varphi_{256}$}
\BIC{$(\varphi_{190} \to \forall n\,(\varphi_{190})) \to \mathrm{FLim}(f,c,\alpha) \to \varphi_{256}$}
\end{ptree}
\noindent$\Pi_{35}$:
\begin{ptree}
\AXC{$\forall f\,(\varphi_{181})$\lbl{定理1.7(⇒)}}
\AXC{$\forall f\,(\varphi_{181}) \to \varphi_{181}$\lbl{A2}}
\BIC{$\varphi_{181}$}
\AXC{$\varphi_{351}$\lbl{A2}}
\BIC{$\forall \alpha\,\forall a\,(\mathrm{FLim}(f,c,\alpha) \to \mathrm{Conv}(a,c) \to \varphi_{72})$}
\AXC{$\varphi_{339}$\lbl{A2}}
\BIC{$\forall a\,(\mathrm{FLim}(f,c,\alpha) \to \mathrm{Conv}(a,c) \to \varphi_{72})$}
\AXC{$\varphi_{328}$\lbl{A2}}
\BIC{$\mathrm{FLim}(f,c,\alpha) \to \mathrm{Conv}(a,c) \to \varphi_{72}$}
\AXC{$\varphi_{489}$\lbl{A1}}
\BIC{$(\mathrm{FLim}(f,c,\alpha) \to \mathrm{FLim}(f,c,\alpha)) \to \varphi_{451}$}
\end{ptree}
\noindent$\Pi_{36}$:
\begin{ptree}
\AXC{$\mathrm{FLim}(f,c,\alpha) \to \mathrm{FLim}(f,c,\alpha)$\lbl{A1}}
\AXC{$\stackrel{\Pi_{35}}{\vdots}$}\noLine
\UIC{$(\mathrm{FLim}(f,c,\alpha) \to \mathrm{FLim}(f,c,\alpha)) \to \varphi_{451}$}
\BIC{$\varphi_{451}$}
\end{ptree}
\noindent$\Pi_{37}$:
\begin{ptree}
\AXC{$\varphi_{190} \to \forall n\,(\varphi_{190})$\lbl{A4}}
\AXC{$\stackrel{\Pi_{34}}{\vdots}$}\noLine
\UIC{$(\varphi_{190} \to \forall n\,(\varphi_{190})) \to \mathrm{FLim}(f,c,\alpha) \to \varphi_{256}$}
\BIC{$\mathrm{FLim}(f,c,\alpha) \to \varphi_{256}$}
\AXC{$\mathrm{FLim}(f,c,\alpha) \to \varphi_{14}$\lbl{A1}}
\AXC{$\stackrel{\Pi_{36}}{\vdots}$}\noLine
\UIC{$\varphi_{451}$}
\BIC{$(\mathrm{FLim}(f,c,\alpha) \to \varphi_{256}) \to \mathrm{FLim}(f,c,\alpha) \to \varphi_{208}$}
\BIC{$\mathrm{FLim}(f,c,\alpha) \to \varphi_{208}$}
\end{ptree}
\noindent$\Pi_{38}$:
\begin{ptree}
\AXC{$\varphi_{189}$\lbl{定理2.9}}
\AXC{$\varphi_{189} \to \forall \alpha\,(\varphi_{213})$\lbl{A2}}
\BIC{$\forall \alpha\,(\varphi_{213})$}
\AXC{$\forall \alpha\,(\varphi_{213}) \to \varphi_{213}$\lbl{A2}}
\BIC{$\varphi_{213}$}
\AXC{$\varphi_{213} \to \varphi_{270}$\lbl{A2}}
\BIC{$\varphi_{270}$}
\AXC{$\varphi_{416}$\lbl{A2}}
\BIC{$\mathrm{Conv}(f \circ a,\alpha) \to \mathrm{Conv}(g \circ a,\beta) \to \beta \neq 0 \to \varphi_{153}$}
\AXC{$\varphi_{549}$\lbl{A1}}
\BIC{$\varphi_{528}$}
\end{ptree}
\noindent$\Pi_{39}$:
\begin{ptree}
\AXC{$\stackrel{\Pi_{26}}{\vdots}$}\noLine
\UIC{$\mathrm{FLim}(f,c,\alpha) \to \varphi_{233}$}
\AXC{$\varphi_{3}$\lbl{A1}}
\AXC{$\stackrel{\Pi_{32}}{\vdots}$}\noLine
\UIC{$\mathrm{FLim}(f,c,\alpha) \to \varphi_{207}$}
\AXC{$\stackrel{\Pi_{37}}{\vdots}$}\noLine
\UIC{$\mathrm{FLim}(f,c,\alpha) \to \varphi_{208}$}
\AXC{$\stackrel{\Pi_{38}}{\vdots}$}\noLine
\UIC{$\varphi_{528}$}
\BIC{$(\mathrm{FLim}(f,c,\alpha) \to \varphi_{207}) \to \varphi_{454}$}
\BIC{$\varphi_{454}$}
\BIC{$(\mathrm{FLim}(f,c,\alpha) \to \varphi_{233}) \to \mathrm{FLim}(f,c,\alpha) \to \varphi_{276}$}
\BIC{$\mathrm{FLim}(f,c,\alpha) \to \varphi_{276}$}
\end{ptree}
\noindent$\Pi_{40}$:
\begin{ptree}
\AXC{$y = z \to x = y \to x = z$\lbl{A6}}
\AXC{$(y = z \to x = y \to x = z) \to x = y \to y = z \to x = z$\lbl{A1}}
\BIC{$x = y \to y = z \to x = z$}
\UIC{$\forall z\,(x = y \to y = z \to x = z)$}
\UIC{$\forall y\,\forall z\,(x = y \to y = z \to x = z)$}
\UIC{$\forall x\,\forall y\,\forall z\,(x = y \to y = z \to x = z)$}
\AXC{$\forall x\,\forall y\,\forall z\,(x = y \to y = z \to x = z) \to \varphi_{74}$\lbl{A2}}
\BIC{$\varphi_{74}$}
\end{ptree}
\noindent$\Pi_{41}$:
\begin{ptree}
\AXC{$\stackrel{\Pi_{40}}{\vdots}$}\noLine
\UIC{$\varphi_{74}$}
\AXC{$\varphi_{74} \to \varphi_{188}$\lbl{A2}}
\BIC{$\varphi_{188}$}
\AXC{$\varphi_{188} \to (f / g) \circ a(n) = (f / g)(a_{n}) \to \varphi_{248}$\lbl{A2}}
\BIC{$(f / g) \circ a(n) = (f / g)(a_{n}) \to \varphi_{248}$}
\end{ptree}
\noindent$\Pi_{42}$:
\begin{ptree}
\AXC{$y = z \to x = y \to x = z$\lbl{A6}}
\AXC{$(y = z \to x = y \to x = z) \to x = y \to y = z \to x = z$\lbl{A1}}
\BIC{$x = y \to y = z \to x = z$}
\UIC{$\forall z\,(x = y \to y = z \to x = z)$}
\UIC{$\forall y\,\forall z\,(x = y \to y = z \to x = z)$}
\UIC{$\forall x\,\forall y\,\forall z\,(x = y \to y = z \to x = z)$}
\AXC{$\forall x\,\forall y\,\forall z\,(x = y \to y = z \to x = z) \to \varphi_{44}$\lbl{A2}}
\BIC{$\varphi_{44}$}
\end{ptree}
\noindent$\Pi_{43}$:
\begin{ptree}
\AXC{$y = z \to x = y \to x = z$\lbl{A6}}
\AXC{$(y = z \to x = y \to x = z) \to x = y \to y = z \to x = z$\lbl{A1}}
\BIC{$x = y \to y = z \to x = z$}
\UIC{$\forall z\,(x = y \to y = z \to x = z)$}
\UIC{$\forall y\,\forall z\,(x = y \to y = z \to x = z)$}
\UIC{$\forall x\,\forall y\,\forall z\,(x = y \to y = z \to x = z)$}
\AXC{$\forall x\,\forall y\,\forall z\,(x = y \to y = z \to x = z) \to \varphi_{186}$\lbl{A2}}
\BIC{$\varphi_{186}$}
\end{ptree}
\noindent$\Pi_{44}$:
\begin{ptree}
\AXC{$z = w \to \varphi_{234}$\lbl{A6}}
\UIC{$\forall w\,(z = w \to \varphi_{234})$}
\UIC{$\forall z\,\forall w\,(z = w \to \varphi_{234})$}
\AXC{$\forall z\,\forall w\,(z = w \to \varphi_{234}) \to \varphi_{310}$\lbl{A2}}
\BIC{$\varphi_{310}$}
\AXC{$\varphi_{310} \to g \circ a(n) = g(a_{n}) \to \varphi_{295}$\lbl{A2}}
\BIC{$g \circ a(n) = g(a_{n}) \to \varphi_{295}$}
\end{ptree}
\noindent$\Pi_{45}$:
\begin{ptree}
\AXC{$\varphi_{13}$\lbl{A1}}
\AXC{$\stackrel{\Pi_{3}}{\vdots}$}\noLine
\UIC{$\varphi_{4}$}
\AXC{$\forall f\,\forall a\,(\mathrm{RFn}(f) \to \mathrm{Seq}(a) \to \varphi_{178})$\lbl{[S8]}}
\AXC{$\varphi_{386}$\lbl{A2}}
\BIC{$\forall a\,(\mathrm{RFn}(g) \to \mathrm{Seq}(a) \to \varphi_{175})$}
\AXC{$\varphi_{384}$\lbl{A2}}
\BIC{$\mathrm{RFn}(g) \to \mathrm{Seq}(a) \to \varphi_{175}$}
\AXC{$\varphi_{511}$\lbl{A1}}
\BIC{$\varphi_{4} \to \varphi_{13} \to (\mathrm{FLim}(f,c,\alpha) \to \varphi_{226}) \to \varphi_{321}$}
\BIC{$\varphi_{13} \to (\mathrm{FLim}(f,c,\alpha) \to \varphi_{226}) \to \varphi_{321}$}
\BIC{$(\mathrm{FLim}(f,c,\alpha) \to \varphi_{226}) \to \varphi_{321}$}
\end{ptree}
\noindent$\Pi_{46}$:
\begin{ptree}
\AXC{$\stackrel{\Pi_{15}}{\vdots}$}\noLine
\UIC{$\mathrm{FLim}(f,c,\alpha) \to \varphi_{226}$}
\AXC{$\stackrel{\Pi_{45}}{\vdots}$}\noLine
\UIC{$(\mathrm{FLim}(f,c,\alpha) \to \varphi_{226}) \to \varphi_{321}$}
\BIC{$\varphi_{321}$}
\AXC{$\varphi_{321} \to \mathrm{FLim}(f,c,\alpha) \to \varphi_{281}$\lbl{A1}}
\BIC{$\mathrm{FLim}(f,c,\alpha) \to \varphi_{281}$}
\AXC{$\varphi_{473}$\lbl{A1}}
\BIC{$\varphi_{149} \to \mathrm{FLim}(f,c,\alpha) \to \varphi_{274}$}
\end{ptree}
\noindent$\Pi_{47}$:
\begin{ptree}
\AXC{$\mathrm{FLim}(f,c,\alpha) \to \varphi_{206}$\lbl{A1}}
\AXC{$\varphi_{149}$\lbl{A2}}
\AXC{$\stackrel{\Pi_{46}}{\vdots}$}\noLine
\UIC{$\varphi_{149} \to \mathrm{FLim}(f,c,\alpha) \to \varphi_{274}$}
\BIC{$\mathrm{FLim}(f,c,\alpha) \to \varphi_{274}$}
\AXC{$\varphi_{482}$\lbl{A1}}
\BIC{$(\mathrm{FLim}(f,c,\alpha) \to \varphi_{206}) \to \mathrm{FLim}(f,c,\alpha) \to \varphi_{251}$}
\BIC{$\mathrm{FLim}(f,c,\alpha) \to \varphi_{251}$}
\end{ptree}
\noindent$\Pi_{48}$:
\begin{ptree}
\AXC{$x = x$\lbl{A5}}
\UIC{$\forall x\,(x = x)$}
\AXC{$\varphi_{133}$\lbl{A2}}
\BIC{$f(a_{n}) \cdot g \circ a(n)^{-1} = f(a_{n}) \cdot g \circ a(n)^{-1}$}
\AXC{$\stackrel{\Pi_{47}}{\vdots}$}\noLine
\UIC{$\mathrm{FLim}(f,c,\alpha) \to \varphi_{251}$}
\AXC{$(\mathrm{FLim}(f,c,\alpha) \to \varphi_{251}) \to \varphi_{492}$\lbl{A1}}
\BIC{$\varphi_{492}$}
\BIC{$(g \circ a(n) = g(a_{n}) \to \varphi_{295}) \to \varphi_{344}$}
\end{ptree}
\noindent$\Pi_{49}$:
\begin{ptree}
\AXC{$y = z \to x = y \to x = z$\lbl{A6}}
\AXC{$(y = z \to x = y \to x = z) \to x = y \to y = z \to x = z$\lbl{A1}}
\BIC{$x = y \to y = z \to x = z$}
\UIC{$\forall z\,(x = y \to y = z \to x = z)$}
\UIC{$\forall y\,\forall z\,(x = y \to y = z \to x = z)$}
\UIC{$\forall x\,\forall y\,\forall z\,(x = y \to y = z \to x = z)$}
\AXC{$\forall x\,\forall y\,\forall z\,(x = y \to y = z \to x = z) \to \varphi_{186}$\lbl{A2}}
\BIC{$\varphi_{186}$}
\end{ptree}
\noindent$\Pi_{50}$:
\begin{ptree}
\AXC{$z = w \to \varphi_{265}$\lbl{A6}}
\UIC{$\forall w\,(z = w \to \varphi_{265})$}
\UIC{$\forall z\,\forall w\,(z = w \to \varphi_{265})$}
\AXC{$\forall z\,\forall w\,(z = w \to \varphi_{265}) \to \varphi_{325}$\lbl{A2}}
\BIC{$\varphi_{325}$}
\AXC{$\varphi_{325} \to f \circ a(n) = f(a_{n}) \to \varphi_{312}$\lbl{A2}}
\BIC{$f \circ a(n) = f(a_{n}) \to \varphi_{312}$}
\end{ptree}
\noindent$\Pi_{51}$:
\begin{ptree}
\AXC{$\mathrm{FLim}(f,c,\alpha) \to \varphi_{206}$\lbl{A1}}
\AXC{$\varphi_{184}$\lbl{A2}}
\AXC{$\mathrm{FLim}(f,c,\alpha) \to \varphi_{269}$\lbl{A1}}
\AXC{$\varphi_{472}$\lbl{A1}}
\BIC{$\varphi_{184} \to \mathrm{FLim}(f,c,\alpha) \to \varphi_{263}$}
\BIC{$\mathrm{FLim}(f,c,\alpha) \to \varphi_{263}$}
\AXC{$\varphi_{515}$\lbl{A1}}
\BIC{$(\mathrm{FLim}(f,c,\alpha) \to \varphi_{206}) \to \varphi_{391}$}
\BIC{$\varphi_{391}$}
\end{ptree}
\noindent$\Pi_{52}$:
\begin{ptree}
\AXC{$\varphi_{190} \to \forall n\,(\varphi_{190})$\lbl{A4}}
\AXC{$\varphi_{410}$\lbl{A3}}
\AXC{$\forall x\,(\varphi_{42}) \to \varphi_{122}$\lbl{A2}}
\AXC{$\stackrel{\Pi_{8}}{\vdots}$}\noLine
\UIC{$\varphi_{168}$}
\AXC{$\varphi_{168} \to \varphi_{400}$\lbl{A1}}
\BIC{$\varphi_{400}$}
\BIC{$\mathrm{FLim}(f,c,\alpha) \to \mathrm{FLim}(g,c,\beta) \to \beta \neq 0 \to \varphi_{122}$}
\AXC{$\stackrel{\Pi_{51}}{\vdots}$}\noLine
\UIC{$\varphi_{391}$}
\BIC{$\mathrm{FLim}(f,c,\alpha) \to \varphi_{218}$}
\AXC{$\varphi_{406}$\lbl{A1}}
\BIC{$\varphi_{190} \to n \in \mathbb{N} \to a_{n} \in \mathrm{dom}(f)$}
\UIC{$\forall n\,(\varphi_{190} \to n \in \mathbb{N} \to a_{n} \in \mathrm{dom}(f))$}
\AXC{$\varphi_{556}$\lbl{A1}}
\BIC{$\varphi_{410} \to (\varphi_{190} \to \forall n\,(\varphi_{190})) \to \mathrm{FLim}(f,c,\alpha) \to \varphi_{228}$}
\BIC{$(\varphi_{190} \to \forall n\,(\varphi_{190})) \to \mathrm{FLim}(f,c,\alpha) \to \varphi_{228}$}
\BIC{$\mathrm{FLim}(f,c,\alpha) \to \varphi_{228}$}
\end{ptree}
\noindent$\Pi_{53}$:
\begin{ptree}
\AXC{$\varphi_{13}$\lbl{A1}}
\AXC{$\stackrel{\Pi_{5}}{\vdots}$}\noLine
\UIC{$\varphi_{5}$}
\AXC{$\forall f\,\forall a\,(\mathrm{RFn}(f) \to \mathrm{Seq}(a) \to \varphi_{178})$\lbl{[S8]}}
\AXC{$\varphi_{387}$\lbl{A2}}
\BIC{$\forall a\,(\mathrm{RFn}(f) \to \mathrm{Seq}(a) \to \varphi_{178})$}
\AXC{$\varphi_{385}$\lbl{A2}}
\BIC{$\mathrm{RFn}(f) \to \mathrm{Seq}(a) \to \varphi_{178}$}
\AXC{$\varphi_{512}$\lbl{A1}}
\BIC{$\varphi_{5} \to \varphi_{13} \to (\mathrm{FLim}(f,c,\alpha) \to \varphi_{228}) \to \varphi_{322}$}
\BIC{$\varphi_{13} \to (\mathrm{FLim}(f,c,\alpha) \to \varphi_{228}) \to \varphi_{322}$}
\BIC{$(\mathrm{FLim}(f,c,\alpha) \to \varphi_{228}) \to \varphi_{322}$}
\end{ptree}
\noindent$\Pi_{54}$:
\begin{ptree}
\AXC{$\stackrel{\Pi_{52}}{\vdots}$}\noLine
\UIC{$\mathrm{FLim}(f,c,\alpha) \to \varphi_{228}$}
\AXC{$\stackrel{\Pi_{53}}{\vdots}$}\noLine
\UIC{$(\mathrm{FLim}(f,c,\alpha) \to \varphi_{228}) \to \varphi_{322}$}
\BIC{$\varphi_{322}$}
\AXC{$\varphi_{322} \to \mathrm{FLim}(f,c,\alpha) \to \varphi_{282}$\lbl{A1}}
\BIC{$\mathrm{FLim}(f,c,\alpha) \to \varphi_{282}$}
\AXC{$\varphi_{474}$\lbl{A1}}
\BIC{$\varphi_{150} \to \mathrm{FLim}(f,c,\alpha) \to \varphi_{275}$}
\end{ptree}
\noindent$\Pi_{55}$:
\begin{ptree}
\AXC{$\mathrm{FLim}(f,c,\alpha) \to \varphi_{206}$\lbl{A1}}
\AXC{$\varphi_{150}$\lbl{A2}}
\AXC{$\stackrel{\Pi_{54}}{\vdots}$}\noLine
\UIC{$\varphi_{150} \to \mathrm{FLim}(f,c,\alpha) \to \varphi_{275}$}
\BIC{$\mathrm{FLim}(f,c,\alpha) \to \varphi_{275}$}
\AXC{$\varphi_{483}$\lbl{A1}}
\BIC{$(\mathrm{FLim}(f,c,\alpha) \to \varphi_{206}) \to \mathrm{FLim}(f,c,\alpha) \to \varphi_{253}$}
\BIC{$\mathrm{FLim}(f,c,\alpha) \to \varphi_{253}$}
\AXC{$\varphi_{533}$\lbl{A1}}
\BIC{$\varphi_{131} \to (f \circ a(n) = f(a_{n}) \to \varphi_{312}) \to \varphi_{357}$}
\end{ptree}
\noindent$\Pi_{56}$:
\begin{ptree}
\AXC{$\stackrel{\Pi_{50}}{\vdots}$}\noLine
\UIC{$f \circ a(n) = f(a_{n}) \to \varphi_{312}$}
\AXC{$x = x$\lbl{A5}}
\UIC{$\forall x\,(x = x)$}
\AXC{$\forall x\,(x = x) \to \varphi_{131}$\lbl{A2}}
\BIC{$\varphi_{131}$}
\AXC{$\stackrel{\Pi_{55}}{\vdots}$}\noLine
\UIC{$\varphi_{131} \to (f \circ a(n) = f(a_{n}) \to \varphi_{312}) \to \varphi_{357}$}
\BIC{$(f \circ a(n) = f(a_{n}) \to \varphi_{312}) \to \varphi_{357}$}
\BIC{$\varphi_{357}$}
\end{ptree}
\noindent$\Pi_{57}$:
\begin{ptree}
\AXC{$y = z \to x = y \to x = z$\lbl{A6}}
\AXC{$(y = z \to x = y \to x = z) \to x = y \to y = z \to x = z$\lbl{A1}}
\BIC{$x = y \to y = z \to x = z$}
\UIC{$\forall z\,(x = y \to y = z \to x = z)$}
\UIC{$\forall y\,\forall z\,(x = y \to y = z \to x = z)$}
\UIC{$\forall x\,\forall y\,\forall z\,(x = y \to y = z \to x = z)$}
\AXC{$\forall x\,\forall y\,\forall z\,(x = y \to y = z \to x = z) \to \varphi_{186}$\lbl{A2}}
\BIC{$\varphi_{186}$}
\end{ptree}
\noindent$\Pi_{58}$:
\begin{ptree}
\AXC{$\stackrel{\Pi_{32}}{\vdots}$}\noLine
\UIC{$\mathrm{FLim}(f,c,\alpha) \to \varphi_{207}$}
\AXC{$\forall a\,\forall \alpha\,(\mathrm{Conv}(a,\alpha) \to \mathrm{Seq}(a))$\lbl{[C1]}}
\AXC{$\varphi_{76}$\lbl{A2}}
\BIC{$\forall \alpha\,(\mathrm{Conv}(g \circ a,\alpha) \to \mathrm{Seq}(g \circ a))$}
\AXC{$\varphi_{111}$\lbl{A2}}
\BIC{$\mathrm{Conv}(g \circ a,\beta) \to \mathrm{Seq}(g \circ a)$}
\AXC{$\varphi_{419}$\lbl{A1}}
\BIC{$(\mathrm{FLim}(f,c,\alpha) \to \varphi_{207}) \to \mathrm{FLim}(f,c,\alpha) \to \varphi_{197}$}
\BIC{$\mathrm{FLim}(f,c,\alpha) \to \varphi_{197}$}
\end{ptree}
\noindent$\Pi_{59}$:
\begin{ptree}
\AXC{$\varphi_{154}$\lbl{[S5]}}
\AXC{$\varphi_{154} \to \mathrm{Seq}(g \circ a) \to \forall n \in \mathbb{N}\,(g \circ a(n) \neq 0) \to \varphi_{126}$\lbl{A2}}
\BIC{$\mathrm{Seq}(g \circ a) \to \forall n \in \mathbb{N}\,(g \circ a(n) \neq 0) \to \varphi_{126}$}
\AXC{$(\mathrm{Seq}(g \circ a) \to \forall n \in \mathbb{N}\,(g \circ a(n) \neq 0) \to \varphi_{126}) \to \varphi_{487}$\lbl{A1}}
\BIC{$\varphi_{487}$}
\end{ptree}
\noindent$\Pi_{60}$:
\begin{ptree}
\AXC{$\stackrel{\Pi_{26}}{\vdots}$}\noLine
\UIC{$\mathrm{FLim}(f,c,\alpha) \to \varphi_{233}$}
\AXC{$\stackrel{\Pi_{58}}{\vdots}$}\noLine
\UIC{$\mathrm{FLim}(f,c,\alpha) \to \varphi_{197}$}
\AXC{$\stackrel{\Pi_{59}}{\vdots}$}\noLine
\UIC{$\varphi_{487}$}
\BIC{$(\mathrm{FLim}(f,c,\alpha) \to \varphi_{233}) \to \varphi_{340}$}
\BIC{$\varphi_{340}$}
\end{ptree}
\noindent$\Pi_{61}$:
\begin{ptree}
\AXC{$\mathrm{FLim}(f,c,\alpha) \to \varphi_{206}$\lbl{A1}}
\AXC{$\varphi_{214}$\lbl{A2}}
\AXC{$\stackrel{\Pi_{60}}{\vdots}$}\noLine
\UIC{$\varphi_{340}$}
\AXC{$\varphi_{340} \to \mathrm{FLim}(f,c,\alpha) \to \varphi_{306}$\lbl{A1}}
\BIC{$\mathrm{FLim}(f,c,\alpha) \to \varphi_{306}$}
\AXC{$\varphi_{495}$\lbl{A1}}
\BIC{$\varphi_{214} \to \mathrm{FLim}(f,c,\alpha) \to \varphi_{300}$}
\BIC{$\mathrm{FLim}(f,c,\alpha) \to \varphi_{300}$}
\AXC{$\varphi_{493}$\lbl{A1}}
\BIC{$(\mathrm{FLim}(f,c,\alpha) \to \varphi_{206}) \to \mathrm{FLim}(f,c,\alpha) \to \varphi_{285}$}
\BIC{$\mathrm{FLim}(f,c,\alpha) \to \varphi_{285}$}
\end{ptree}
\noindent$\Pi_{62}$:
\begin{ptree}
\AXC{$x = x$\lbl{A5}}
\UIC{$\forall x\,(x = x)$}
\AXC{$\forall x\,(x = x) \to \varphi_{145}$\lbl{A2}}
\BIC{$\varphi_{145}$}
\AXC{$\stackrel{\Pi_{61}}{\vdots}$}\noLine
\UIC{$\mathrm{FLim}(f,c,\alpha) \to \varphi_{285}$}
\AXC{$(\mathrm{FLim}(f,c,\alpha) \to \varphi_{285}) \to \varphi_{145} \to \varphi_{480}$\lbl{A1}}
\BIC{$\varphi_{145} \to \varphi_{480}$}
\BIC{$\varphi_{480}$}
\end{ptree}
\noindent$\Pi_{63}$:
\begin{ptree}
\AXC{$z = w \to \varphi_{257}$\lbl{A6}}
\UIC{$\forall w\,(z = w \to \varphi_{257})$}
\UIC{$\forall z\,\forall w\,(z = w \to \varphi_{257})$}
\AXC{$\forall z\,\forall w\,(z = w \to \varphi_{257}) \to \varphi_{337}$\lbl{A2}}
\BIC{$\varphi_{337}$}
\AXC{$\varphi_{337} \to g \circ a^{\ast}(n) = g \circ a(n)^{-1} \to \varphi_{145} \to \varphi_{139}$\lbl{A2}}
\BIC{$g \circ a^{\ast}(n) = g \circ a(n)^{-1} \to \varphi_{145} \to \varphi_{139}$}
\AXC{$\stackrel{\Pi_{62}}{\vdots}$}\noLine
\UIC{$\varphi_{480}$}
\BIC{$\mathrm{FLim}(f,c,\alpha) \to \varphi_{342}$}
\end{ptree}
\noindent$\Pi_{64}$:
\begin{ptree}
\AXC{$\stackrel{\Pi_{37}}{\vdots}$}\noLine
\UIC{$\mathrm{FLim}(f,c,\alpha) \to \varphi_{208}$}
\AXC{$\forall a\,\forall \alpha\,(\mathrm{Conv}(a,\alpha) \to \mathrm{Seq}(a))$\lbl{[C1]}}
\AXC{$\varphi_{77}$\lbl{A2}}
\BIC{$\forall \alpha\,(\mathrm{Conv}(f \circ a,\alpha) \to \mathrm{Seq}(f \circ a))$}
\AXC{$\varphi_{112}$\lbl{A2}}
\BIC{$\mathrm{Conv}(f \circ a,\alpha) \to \mathrm{Seq}(f \circ a)$}
\AXC{$\varphi_{420}$\lbl{A1}}
\BIC{$(\mathrm{FLim}(f,c,\alpha) \to \varphi_{208}) \to \mathrm{FLim}(f,c,\alpha) \to \varphi_{198}$}
\BIC{$\mathrm{FLim}(f,c,\alpha) \to \varphi_{198}$}
\end{ptree}
\noindent$\Pi_{65}$:
\begin{ptree}
\AXC{$\stackrel{\Pi_{64}}{\vdots}$}\noLine
\UIC{$\mathrm{FLim}(f,c,\alpha) \to \varphi_{198}$}
\AXC{$\varphi_{137}$\lbl{[S4]}}
\AXC{$\varphi_{137} \to \varphi_{232}$\lbl{A2}}
\BIC{$\varphi_{232}$}
\AXC{$\varphi_{232} \to \varphi_{334}$\lbl{A2}}
\BIC{$\varphi_{334}$}
\AXC{$\varphi_{334} \to \varphi_{503}$\lbl{A1}}
\BIC{$\varphi_{503}$}
\BIC{$(\mathrm{FLim}(f,c,\alpha) \to \varphi_{210}) \to \mathrm{FLim}(f,c,\alpha) \to \varphi_{388}$}
\end{ptree}
\noindent$\Pi_{66}$:
\begin{ptree}
\AXC{$\stackrel{\Pi_{60}}{\vdots}$}\noLine
\UIC{$\varphi_{340}$}
\AXC{$\varphi_{340} \to \mathrm{FLim}(f,c,\alpha) \to \varphi_{210}$\lbl{A1}}
\BIC{$\mathrm{FLim}(f,c,\alpha) \to \varphi_{210}$}
\AXC{$\stackrel{\Pi_{65}}{\vdots}$}\noLine
\UIC{$(\mathrm{FLim}(f,c,\alpha) \to \varphi_{210}) \to \mathrm{FLim}(f,c,\alpha) \to \varphi_{388}$}
\BIC{$\mathrm{FLim}(f,c,\alpha) \to \varphi_{388}$}
\end{ptree}
\noindent$\Pi_{67}$:
\begin{ptree}
\AXC{$\varphi_{185} \to \varphi_{182}$\lbl{A2}}
\AXC{$\stackrel{\Pi_{66}}{\vdots}$}\noLine
\UIC{$\mathrm{FLim}(f,c,\alpha) \to \varphi_{388}$}
\AXC{$(\mathrm{FLim}(f,c,\alpha) \to \varphi_{388}) \to \mathrm{FLim}(f,c,\alpha) \to \varphi_{368}$\lbl{A1}}
\BIC{$\mathrm{FLim}(f,c,\alpha) \to \varphi_{368}$}
\AXC{$\varphi_{531}$\lbl{A1}}
\BIC{$(\varphi_{185} \to \varphi_{182}) \to \mathrm{FLim}(f,c,\alpha) \to \varphi_{362}$}
\BIC{$\mathrm{FLim}(f,c,\alpha) \to \varphi_{362}$}
\end{ptree}
\noindent$\Pi_{68}$:
\begin{ptree}
\AXC{$\mathrm{FLim}(f,c,\alpha) \to \varphi_{206}$\lbl{A1}}
\AXC{$\stackrel{\Pi_{67}}{\vdots}$}\noLine
\UIC{$\mathrm{FLim}(f,c,\alpha) \to \varphi_{362}$}
\AXC{$\varphi_{517}$\lbl{A1}}
\BIC{$(\mathrm{FLim}(f,c,\alpha) \to \varphi_{206}) \to \varphi_{367}$}
\BIC{$\varphi_{367}$}
\AXC{$\varphi_{367} \to (\mathrm{FLim}(f,c,\alpha) \to \varphi_{342}) \to \varphi_{401} \to \varphi_{363}$\lbl{A1}}
\BIC{$(\mathrm{FLim}(f,c,\alpha) \to \varphi_{342}) \to \varphi_{401} \to \varphi_{363}$}
\end{ptree}
\noindent$\Pi_{69}$:
\begin{ptree}
\AXC{$\stackrel{\Pi_{57}}{\vdots}$}\noLine
\UIC{$\varphi_{186}$}
\AXC{$\varphi_{186} \to \varphi_{332}$\lbl{A2}}
\BIC{$\varphi_{332}$}
\AXC{$\varphi_{332} \to \varphi_{401}$\lbl{A2}}
\BIC{$\varphi_{401}$}
\AXC{$\stackrel{\Pi_{63}}{\vdots}$}\noLine
\UIC{$\mathrm{FLim}(f,c,\alpha) \to \varphi_{342}$}
\AXC{$\stackrel{\Pi_{68}}{\vdots}$}\noLine
\UIC{$(\mathrm{FLim}(f,c,\alpha) \to \varphi_{342}) \to \varphi_{401} \to \varphi_{363}$}
\BIC{$\varphi_{401} \to \varphi_{363}$}
\BIC{$\varphi_{363}$}
\end{ptree}
\noindent$\Pi_{70}$:
\begin{ptree}
\AXC{$\stackrel{\Pi_{49}}{\vdots}$}\noLine
\UIC{$\varphi_{186}$}
\AXC{$\varphi_{186} \to \varphi_{324}$\lbl{A2}}
\BIC{$\varphi_{324}$}
\AXC{$\varphi_{324} \to \varphi_{393}$\lbl{A2}}
\BIC{$\varphi_{393}$}
\AXC{$\stackrel{\Pi_{56}}{\vdots}$}\noLine
\UIC{$\varphi_{357}$}
\AXC{$\stackrel{\Pi_{69}}{\vdots}$}\noLine
\UIC{$\varphi_{363}$}
\AXC{$\varphi_{363} \to \varphi_{357} \to \varphi_{393} \to \varphi_{358}$\lbl{A1}}
\BIC{$\varphi_{357} \to \varphi_{393} \to \varphi_{358}$}
\BIC{$\varphi_{393} \to \varphi_{358}$}
\BIC{$\varphi_{358}$}
\AXC{$\varphi_{358} \to \varphi_{344} \to \varphi_{484}$\lbl{A1}}
\BIC{$\varphi_{344} \to \varphi_{484}$}
\end{ptree}
\noindent$\Pi_{71}$:
\begin{ptree}
\AXC{$\stackrel{\Pi_{44}}{\vdots}$}\noLine
\UIC{$g \circ a(n) = g(a_{n}) \to \varphi_{295}$}
\AXC{$\stackrel{\Pi_{48}}{\vdots}$}\noLine
\UIC{$(g \circ a(n) = g(a_{n}) \to \varphi_{295}) \to \varphi_{344}$}
\BIC{$\varphi_{344}$}
\AXC{$\stackrel{\Pi_{70}}{\vdots}$}\noLine
\UIC{$\varphi_{344} \to \varphi_{484}$}
\BIC{$\varphi_{484}$}
\end{ptree}
\noindent$\Pi_{72}$:
\begin{ptree}
\AXC{$\stackrel{\Pi_{43}}{\vdots}$}\noLine
\UIC{$\varphi_{186}$}
\AXC{$\varphi_{186} \to \varphi_{317}$\lbl{A2}}
\BIC{$\varphi_{317}$}
\AXC{$\varphi_{475}$\lbl{A2}}
\BIC{$f \circ a \otimes g \circ a^{\ast}(n) = f(a_{n}) \cdot g \circ a(n)^{-1} \to \varphi_{294}$}
\AXC{$\stackrel{\Pi_{71}}{\vdots}$}\noLine
\UIC{$\varphi_{484}$}
\BIC{$\varphi_{347}$}
\AXC{$\varphi_{347} \to \varphi_{298} \to \varphi_{348}$\lbl{A1}}
\BIC{$\varphi_{298} \to \varphi_{348}$}
\end{ptree}
\noindent$\Pi_{73}$:
\begin{ptree}
\AXC{$\stackrel{\Pi_{1}}{\vdots}$}\noLine
\UIC{$\forall x\,\forall y\,(x = y \to y = x)$}
\AXC{$\forall x\,\forall y\,(x = y \to y = x) \to \varphi_{163}$\lbl{A2}}
\BIC{$\varphi_{163}$}
\AXC{$\varphi_{163} \to \varphi_{298}$\lbl{A2}}
\BIC{$\varphi_{298}$}
\AXC{$\stackrel{\Pi_{72}}{\vdots}$}\noLine
\UIC{$\varphi_{298} \to \varphi_{348}$}
\BIC{$\varphi_{348}$}
\end{ptree}
\noindent$\Pi_{74}$:
\begin{ptree}
\AXC{$\varphi_{54} \to \varphi_{127}$\lbl{A2}}
\AXC{$\stackrel{\Pi_{7}}{\vdots}$}\noLine
\UIC{$\mathrm{FLim}(f,c,\alpha) \to \varphi_{299}$}
\AXC{$\varphi_{407}$\lbl{A1}}
\BIC{$\mathrm{FLim}(f,c,\alpha) \to \mathrm{FLim}(g,c,\beta) \to \beta \neq 0 \to \varphi_{54}$}
\AXC{$\varphi_{461}$\lbl{A1}}
\BIC{$\varphi_{403}$}
\BIC{$\mathrm{FLim}(f,c,\alpha) \to \mathrm{FLim}(g,c,\beta) \to \beta \neq 0 \to \varphi_{127}$}
\AXC{$\varphi_{521}$\lbl{A1}}
\BIC{$\varphi_{488}$}
\end{ptree}
\noindent$\Pi_{75}$:
\begin{ptree}
\AXC{$\mathrm{FLim}(f,c,\alpha) \to \varphi_{206}$\lbl{A1}}
\AXC{$\varphi_{184}$\lbl{A2}}
\AXC{$\mathrm{FLim}(f,c,\alpha) \to \varphi_{269}$\lbl{A1}}
\AXC{$\varphi_{472}$\lbl{A1}}
\BIC{$\varphi_{184} \to \mathrm{FLim}(f,c,\alpha) \to \varphi_{263}$}
\BIC{$\mathrm{FLim}(f,c,\alpha) \to \varphi_{263}$}
\AXC{$\stackrel{\Pi_{74}}{\vdots}$}\noLine
\UIC{$\varphi_{488}$}
\BIC{$(\mathrm{FLim}(f,c,\alpha) \to \varphi_{206}) \to \mathrm{FLim}(f,c,\alpha) \to \varphi_{302}$}
\BIC{$\mathrm{FLim}(f,c,\alpha) \to \varphi_{302}$}
\end{ptree}
\noindent$\Pi_{76}$:
\begin{ptree}
\AXC{$\stackrel{\Pi_{73}}{\vdots}$}\noLine
\UIC{$\varphi_{348}$}
\AXC{$\stackrel{\Pi_{75}}{\vdots}$}\noLine
\UIC{$\mathrm{FLim}(f,c,\alpha) \to \varphi_{302}$}
\AXC{$(\mathrm{FLim}(f,c,\alpha) \to \varphi_{302}) \to \varphi_{348} \to \varphi_{466}$\lbl{A1}}
\BIC{$\varphi_{348} \to \varphi_{466}$}
\BIC{$\varphi_{466}$}
\end{ptree}
\noindent$\Pi_{77}$:
\begin{ptree}
\AXC{$\stackrel{\Pi_{42}}{\vdots}$}\noLine
\UIC{$\varphi_{44}$}
\AXC{$\varphi_{44} \to \varphi_{231}$\lbl{A2}}
\BIC{$\varphi_{231}$}
\AXC{$\varphi_{231} \to (f / g)(a_{n}) = f(a_{n}) \cdot g(a_{n})^{-1} \to \varphi_{272}$\lbl{A2}}
\BIC{$(f / g)(a_{n}) = f(a_{n}) \cdot g(a_{n})^{-1} \to \varphi_{272}$}
\AXC{$\stackrel{\Pi_{76}}{\vdots}$}\noLine
\UIC{$\varphi_{466}$}
\BIC{$\mathrm{FLim}(f,c,\alpha) \to \varphi_{311}$}
\end{ptree}
\noindent$\Pi_{78}$:
\begin{ptree}
\AXC{$\mathrm{FLim}(f,c,\alpha) \to \varphi_{206}$\lbl{A1}}
\AXC{$\varphi_{184}$\lbl{A2}}
\AXC{$\mathrm{FLim}(f,c,\alpha) \to \varphi_{269}$\lbl{A1}}
\AXC{$\varphi_{472}$\lbl{A1}}
\BIC{$\varphi_{184} \to \mathrm{FLim}(f,c,\alpha) \to \varphi_{263}$}
\BIC{$\mathrm{FLim}(f,c,\alpha) \to \varphi_{263}$}
\AXC{$\varphi_{476}$\lbl{A1}}
\BIC{$(\mathrm{FLim}(f,c,\alpha) \to \varphi_{206}) \to \mathrm{FLim}(f,c,\alpha) \to \varphi_{236}$}
\BIC{$\mathrm{FLim}(f,c,\alpha) \to \varphi_{236}$}
\AXC{$\varphi_{414}$\lbl{A1}}
\BIC{$\varphi_{190} \to n \in \mathbb{N} \to a_{n} \in \mathrm{dom}((f / g))$}
\UIC{$\forall n\,(\varphi_{190} \to n \in \mathbb{N} \to a_{n} \in \mathrm{dom}((f / g)))$}
\end{ptree}
\noindent$\Pi_{79}$:
\begin{ptree}
\AXC{$\varphi_{190} \to \forall n\,(\varphi_{190})$\lbl{A4}}
\AXC{$\varphi_{421}$\lbl{A3}}
\AXC{$\stackrel{\Pi_{78}}{\vdots}$}\noLine
\UIC{$\forall n\,(\varphi_{190} \to n \in \mathbb{N} \to a_{n} \in \mathrm{dom}((f / g)))$}
\AXC{$\varphi_{561}$\lbl{A1}}
\BIC{$\varphi_{421} \to (\varphi_{190} \to \forall n\,(\varphi_{190})) \to \mathrm{FLim}(f,c,\alpha) \to \varphi_{239}$}
\BIC{$(\varphi_{190} \to \forall n\,(\varphi_{190})) \to \mathrm{FLim}(f,c,\alpha) \to \varphi_{239}$}
\BIC{$\mathrm{FLim}(f,c,\alpha) \to \varphi_{239}$}
\end{ptree}
\noindent$\Pi_{80}$:
\begin{ptree}
\AXC{$\stackrel{\Pi_{7}}{\vdots}$}\noLine
\UIC{$\mathrm{FLim}(f,c,\alpha) \to \varphi_{299}$}
\AXC{$(\mathrm{FLim}(f,c,\alpha) \to \varphi_{299}) \to \varphi_{12}$\lbl{A1}}
\BIC{$\varphi_{12}$}
\end{ptree}
\noindent$\Pi_{81}$:
\begin{ptree}
\AXC{$\forall f\,\forall a\,(\mathrm{RFn}(f) \to \mathrm{Seq}(a) \to \varphi_{178})$\lbl{[S8]}}
\AXC{$\forall f\,\forall a\,(\mathrm{RFn}(f) \to \mathrm{Seq}(a) \to \varphi_{178}) \to \forall a\,(\varphi_{280})$\lbl{A2}}
\BIC{$\forall a\,(\varphi_{280})$}
\AXC{$\forall a\,(\varphi_{280}) \to \varphi_{280}$\lbl{A2}}
\BIC{$\varphi_{280}$}
\end{ptree}
\noindent$\Pi_{82}$:
\begin{ptree}
\AXC{$\stackrel{\Pi_{79}}{\vdots}$}\noLine
\UIC{$\mathrm{FLim}(f,c,\alpha) \to \varphi_{239}$}
\AXC{$\varphi_{13}$\lbl{A1}}
\AXC{$\stackrel{\Pi_{80}}{\vdots}$}\noLine
\UIC{$\varphi_{12}$}
\AXC{$\stackrel{\Pi_{81}}{\vdots}$}\noLine
\UIC{$\varphi_{280}$}
\AXC{$\varphi_{280} \to \varphi_{12} \to \varphi_{13} \to (\mathrm{FLim}(f,c,\alpha) \to \varphi_{239}) \to \varphi_{341}$\lbl{A1}}
\BIC{$\varphi_{12} \to \varphi_{13} \to (\mathrm{FLim}(f,c,\alpha) \to \varphi_{239}) \to \varphi_{341}$}
\BIC{$\varphi_{13} \to (\mathrm{FLim}(f,c,\alpha) \to \varphi_{239}) \to \varphi_{341}$}
\BIC{$(\mathrm{FLim}(f,c,\alpha) \to \varphi_{239}) \to \varphi_{341}$}
\BIC{$\varphi_{341}$}
\end{ptree}
\noindent$\Pi_{83}$:
\begin{ptree}
\AXC{$\mathrm{FLim}(f,c,\alpha) \to \varphi_{206}$\lbl{A1}}
\AXC{$\varphi_{215}$\lbl{A2}}
\AXC{$\stackrel{\Pi_{82}}{\vdots}$}\noLine
\UIC{$\varphi_{341}$}
\AXC{$\varphi_{341} \to \varphi_{320}$\lbl{A1}}
\BIC{$\varphi_{320}$}
\AXC{$\varphi_{320} \to \varphi_{215} \to \mathrm{FLim}(f,c,\alpha) \to \varphi_{303}$\lbl{A1}}
\BIC{$\varphi_{215} \to \mathrm{FLim}(f,c,\alpha) \to \varphi_{303}$}
\BIC{$\mathrm{FLim}(f,c,\alpha) \to \varphi_{303}$}
\AXC{$\varphi_{494}$\lbl{A1}}
\BIC{$(\mathrm{FLim}(f,c,\alpha) \to \varphi_{206}) \to \mathrm{FLim}(f,c,\alpha) \to \varphi_{286}$}
\BIC{$\mathrm{FLim}(f,c,\alpha) \to \varphi_{286}$}
\end{ptree}
\noindent$\Pi_{84}$:
\begin{ptree}
\AXC{$\stackrel{\Pi_{41}}{\vdots}$}\noLine
\UIC{$(f / g) \circ a(n) = (f / g)(a_{n}) \to \varphi_{248}$}
\AXC{$\stackrel{\Pi_{77}}{\vdots}$}\noLine
\UIC{$\mathrm{FLim}(f,c,\alpha) \to \varphi_{311}$}
\AXC{$\stackrel{\Pi_{83}}{\vdots}$}\noLine
\UIC{$\mathrm{FLim}(f,c,\alpha) \to \varphi_{286}$}
\AXC{$\varphi_{545}$\lbl{A1}}
\BIC{$(\mathrm{FLim}(f,c,\alpha) \to \varphi_{311}) \to \varphi_{463}$}
\BIC{$\varphi_{463}$}
\BIC{$\mathrm{FLim}(f,c,\alpha) \to \varphi_{313}$}
\end{ptree}
\noindent$\Pi_{85}$:
\begin{ptree}
\AXC{$\stackrel{\Pi_{1}}{\vdots}$}\noLine
\UIC{$\forall x\,\forall y\,(x = y \to y = x)$}
\AXC{$\varphi_{100}$\lbl{A2}}
\BIC{$\forall y\,((f / g) \circ a(n) = y \to y = (f / g) \circ a(n))$}
\AXC{$\varphi_{356}$\lbl{A2}}
\BIC{$\varphi_{259}$}
\AXC{$\stackrel{\Pi_{84}}{\vdots}$}\noLine
\UIC{$\mathrm{FLim}(f,c,\alpha) \to \varphi_{313}$}
\AXC{$\varphi_{508}$\lbl{A1}}
\BIC{$\varphi_{259} \to \mathrm{FLim}(f,c,\alpha) \to \varphi_{314}$}
\BIC{$\mathrm{FLim}(f,c,\alpha) \to \varphi_{314}$}
\end{ptree}
\noindent$\Pi_{86}$:
\begin{ptree}
\AXC{$\forall n\,(\varphi_{330}) \to \forall n\,(\varphi_{190}) \to \varphi_{118}$\lbl{A3}}
\AXC{$\stackrel{\Pi_{85}}{\vdots}$}\noLine
\UIC{$\mathrm{FLim}(f,c,\alpha) \to \varphi_{314}$}
\AXC{$(\mathrm{FLim}(f,c,\alpha) \to \varphi_{314}) \to \varphi_{330}$\lbl{A1}}
\BIC{$\varphi_{330}$}
\UIC{$\forall n\,(\varphi_{330})$}
\AXC{$\varphi_{566}$\lbl{A1}}
\BIC{$(\forall n\,(\varphi_{330}) \to \forall n\,(\varphi_{190}) \to \varphi_{118}) \to (\varphi_{190} \to \forall n\,(\varphi_{190})) \to \varphi_{335}$}
\BIC{$(\varphi_{190} \to \forall n\,(\varphi_{190})) \to \varphi_{335}$}
\end{ptree}
\noindent$\Pi_{87}$:
\begin{ptree}
\AXC{$\varphi_{86}$\lbl{[C13]}}
\AXC{$\varphi_{86} \to \varphi_{258}$\lbl{A2}}
\BIC{$\varphi_{258}$}
\AXC{$\varphi_{258} \to \varphi_{333}$\lbl{A2}}
\BIC{$\varphi_{333}$}
\AXC{$\varphi_{333} \to \mathrm{Seq}((f / g) \circ a) \to \varphi_{118} \to \varphi_{160}$\lbl{A2}}
\BIC{$\mathrm{Seq}((f / g) \circ a) \to \varphi_{118} \to \varphi_{160}$}
\AXC{$\varphi_{564}$\lbl{A1}}
\BIC{$(\mathrm{FLim}(f,c,\alpha) \to \varphi_{212}) \to \varphi_{501}$}
\end{ptree}
\noindent$\Pi_{88}$:
\begin{ptree}
\AXC{$\stackrel{\Pi_{82}}{\vdots}$}\noLine
\UIC{$\varphi_{341}$}
\AXC{$\varphi_{341} \to \mathrm{FLim}(f,c,\alpha) \to \varphi_{212}$\lbl{A1}}
\BIC{$\mathrm{FLim}(f,c,\alpha) \to \varphi_{212}$}
\AXC{$\stackrel{\Pi_{87}}{\vdots}$}\noLine
\UIC{$(\mathrm{FLim}(f,c,\alpha) \to \varphi_{212}) \to \varphi_{501}$}
\BIC{$\varphi_{501}$}
\end{ptree}
\noindent$\Pi_{89}$:
\begin{ptree}
\AXC{$\stackrel{\Pi_{39}}{\vdots}$}\noLine
\UIC{$\mathrm{FLim}(f,c,\alpha) \to \varphi_{276}$}
\AXC{$\varphi_{190} \to \forall n\,(\varphi_{190})$\lbl{A4}}
\AXC{$\stackrel{\Pi_{86}}{\vdots}$}\noLine
\UIC{$(\varphi_{190} \to \forall n\,(\varphi_{190})) \to \varphi_{335}$}
\BIC{$\varphi_{335}$}
\AXC{$\stackrel{\Pi_{88}}{\vdots}$}\noLine
\UIC{$\varphi_{501}$}
\BIC{$(\mathrm{FLim}(f,c,\alpha) \to \varphi_{276}) \to \mathrm{FLim}(f,c,\alpha) \to \varphi_{238}$}
\BIC{$\mathrm{FLim}(f,c,\alpha) \to \varphi_{238}$}
\end{ptree}
\noindent$\Pi_{90}$:
\begin{ptree}
\AXC{$\stackrel{\Pi_{89}}{\vdots}$}\noLine
\UIC{$\mathrm{FLim}(f,c,\alpha) \to \varphi_{238}$}
\AXC{$(\mathrm{FLim}(f,c,\alpha) \to \varphi_{238}) \to \varphi_{267}$\lbl{A1}}
\BIC{$\varphi_{267}$}
\UIC{$\forall a\,(\varphi_{267})$}
\AXC{$\forall a\,(\varphi_{267}) \to \varphi_{509}$\lbl{A1}}
\BIC{$\varphi_{509}$}
\end{ptree}
\noindent$\Pi_{91}$:
\begin{ptree}
\AXC{$\mathrm{FLim}(f,c,\alpha) \land \mathrm{FLim}(g,c,\beta) \land \beta \neq 0 \to \varphi_{1}$\lbl{A4}}
\AXC{$\forall a\,(\varphi_{267}) \to \varphi_{1} \to \forall a\,(\mathrm{Seq}(a) \to \mathrm{Conv}(a,c) \to \varphi_{148})$\lbl{A3}}
\AXC{$\stackrel{\Pi_{90}}{\vdots}$}\noLine
\UIC{$\varphi_{509}$}
\BIC{$\varphi_{378}$}
\BIC{$\mathrm{FLim}(f,c,\alpha) \to \varphi_{242}$}
\end{ptree}
\noindent$\Pi_{92}$:
\begin{ptree}
\AXC{$\stackrel{\Pi_{1}}{\vdots}$}\noLine
\UIC{$\forall x\,\forall y\,(x = y \to y = x)$}
\AXC{$\forall x\,\forall y\,(x = y \to y = x) \to \forall y\,(0 = y \to y = 0)$\lbl{A2}}
\BIC{$\forall y\,(0 = y \to y = 0)$}
\end{ptree}
\noindent$\Pi_{93}$:
\begin{ptree}
\AXC{$\stackrel{\Pi_{92}}{\vdots}$}\noLine
\UIC{$\forall y\,(0 = y \to y = 0)$}
\AXC{$\forall y\,(0 = y \to y = 0) \to 0 = \lvert g(x) \rvert \to \lvert g(x) \rvert = 0$\lbl{A2}}
\BIC{$0 = \lvert g(x) \rvert \to \lvert g(x) \rvert = 0$}
\end{ptree}
\noindent$\Pi_{94}$:
\begin{ptree}
\AXC{$y = z \to x = y \to x = z$\lbl{A6}}
\AXC{$(y = z \to x = y \to x = z) \to x = y \to y = z \to x = z$\lbl{A1}}
\BIC{$x = y \to y = z \to x = z$}
\UIC{$\forall z\,(x = y \to y = z \to x = z)$}
\UIC{$\forall y\,\forall z\,(x = y \to y = z \to x = z)$}
\UIC{$\forall x\,\forall y\,\forall z\,(x = y \to y = z \to x = z)$}
\AXC{$\varphi_{67}$\lbl{A2}}
\BIC{$\forall y\,\forall z\,(\lvert g(x) \rvert = y \to y = z \to \lvert g(x) \rvert = z)$}
\AXC{$\varphi_{151}$\lbl{A2}}
\BIC{$\forall z\,(\lvert g(x) \rvert = \lvert 0 \rvert \to \lvert 0 \rvert = z \to \lvert g(x) \rvert = z)$}
\AXC{$\varphi_{164}$\lbl{A2}}
\BIC{$\lvert g(x) \rvert = \lvert 0 \rvert \to \lvert 0 \rvert = 0 \to \lvert g(x) \rvert = 0$}
\end{ptree}
\noindent$\Pi_{95}$:
\begin{ptree}
\AXC{$x = x$\lbl{A5}}
\UIC{$\forall x\,(x = x)$}
\AXC{$\forall x\,(x = x) \to 0 = 0$\lbl{A2}}
\BIC{$0 = 0$}
\AXC{$0 = 0 \to \mathrm{FLim}(f,c,\alpha) \to \varphi_{435}$\lbl{A1}}
\BIC{$\mathrm{FLim}(f,c,\alpha) \to \varphi_{435}$}
\end{ptree}
\noindent$\Pi_{96}$:
\begin{ptree}
\AXC{$\forall x\,(0 \leq x \to \lvert x \rvert = x)$\lbl{[A1]}}
\AXC{$\forall x\,(0 \leq x \to \lvert x \rvert = x) \to 0 \leq 0 \to \lvert 0 \rvert = 0$\lbl{A2}}
\BIC{$0 \leq 0 \to \lvert 0 \rvert = 0$}
\AXC{$\varphi_{544}$\lbl{A1}}
\BIC{$(\mathrm{FLim}(f,c,\alpha) \to \varphi_{435}) \to \mathrm{FLim}(f,c,\alpha) \to \varphi_{437}$}
\end{ptree}
\noindent$\Pi_{97}$:
\begin{ptree}
\AXC{$\stackrel{\Pi_{95}}{\vdots}$}\noLine
\UIC{$\mathrm{FLim}(f,c,\alpha) \to \varphi_{435}$}
\AXC{$\stackrel{\Pi_{96}}{\vdots}$}\noLine
\UIC{$(\mathrm{FLim}(f,c,\alpha) \to \varphi_{435}) \to \mathrm{FLim}(f,c,\alpha) \to \varphi_{437}$}
\BIC{$\mathrm{FLim}(f,c,\alpha) \to \varphi_{437}$}
\end{ptree}
\noindent$\Pi_{98}$:
\begin{ptree}
\AXC{$\mathrm{FLim}(f,c,\alpha) \to \varphi_{442}$\lbl{A1}}
\AXC{$(\mathrm{FLim}(f,c,\alpha) \to \varphi_{442}) \to \mathrm{FLim}(f,c,\alpha) \to \varphi_{439}$\lbl{A1}}
\BIC{$\mathrm{FLim}(f,c,\alpha) \to \varphi_{439}$}
\AXC{$\varphi_{557}$\lbl{A1}}
\BIC{$\lvert g(x) \rvert = \lvert g(x) \rvert \to \varphi_{37} \to \mathrm{FLim}(f,c,\alpha) \to \varphi_{449}$}
\end{ptree}
\noindent$\Pi_{99}$:
\begin{ptree}
\AXC{$x = x$\lbl{A5}}
\UIC{$\forall x\,(x = x)$}
\AXC{$\forall x\,(x = x) \to \lvert g(x) \rvert = \lvert g(x) \rvert$\lbl{A2}}
\BIC{$\lvert g(x) \rvert = \lvert g(x) \rvert$}
\AXC{$\stackrel{\Pi_{98}}{\vdots}$}\noLine
\UIC{$\lvert g(x) \rvert = \lvert g(x) \rvert \to \varphi_{37} \to \mathrm{FLim}(f,c,\alpha) \to \varphi_{449}$}
\BIC{$\varphi_{37} \to \mathrm{FLim}(f,c,\alpha) \to \varphi_{449}$}
\end{ptree}
\noindent$\Pi_{100}$:
\begin{ptree}
\AXC{$z = w \to \lvert g(x) \rvert = \lvert z \rvert \to \lvert g(x) \rvert = \lvert w \rvert$\lbl{A6}}
\UIC{$\forall w\,(z = w \to \lvert g(x) \rvert = \lvert z \rvert \to \lvert g(x) \rvert = \lvert w \rvert)$}
\UIC{$\varphi_{21}$}
\AXC{$\varphi_{21} \to \varphi_{49}$\lbl{A2}}
\BIC{$\varphi_{49}$}
\AXC{$\varphi_{49} \to \varphi_{37}$\lbl{A2}}
\BIC{$\varphi_{37}$}
\AXC{$\stackrel{\Pi_{99}}{\vdots}$}\noLine
\UIC{$\varphi_{37} \to \mathrm{FLim}(f,c,\alpha) \to \varphi_{449}$}
\BIC{$\mathrm{FLim}(f,c,\alpha) \to \varphi_{449}$}
\end{ptree}
\noindent$\Pi_{101}$:
\begin{ptree}
\AXC{$\stackrel{\Pi_{94}}{\vdots}$}\noLine
\UIC{$\lvert g(x) \rvert = \lvert 0 \rvert \to \lvert 0 \rvert = 0 \to \lvert g(x) \rvert = 0$}
\AXC{$\stackrel{\Pi_{97}}{\vdots}$}\noLine
\UIC{$\mathrm{FLim}(f,c,\alpha) \to \varphi_{437}$}
\AXC{$\stackrel{\Pi_{100}}{\vdots}$}\noLine
\UIC{$\mathrm{FLim}(f,c,\alpha) \to \varphi_{449}$}
\AXC{$\varphi_{572}$\lbl{A1}}
\BIC{$(\mathrm{FLim}(f,c,\alpha) \to \varphi_{437}) \to \varphi_{467}$}
\BIC{$\varphi_{467}$}
\BIC{$\mathrm{FLim}(f,c,\alpha) \to \varphi_{444}$}
\end{ptree}
\noindent$\Pi_{102}$:
\begin{ptree}
\AXC{$\stackrel{\Pi_{1}}{\vdots}$}\noLine
\UIC{$\forall x\,\forall y\,(x = y \to y = x)$}
\AXC{$\varphi_{16}$\lbl{A2}}
\BIC{$\forall y\,(\lvert g(x) \rvert = y \to y = \lvert g(x) \rvert)$}
\AXC{$\varphi_{66}$\lbl{A2}}
\BIC{$\lvert g(x) \rvert = 0 \to 0 = \lvert g(x) \rvert$}
\AXC{$\stackrel{\Pi_{101}}{\vdots}$}\noLine
\UIC{$\mathrm{FLim}(f,c,\alpha) \to \varphi_{444}$}
\AXC{$(\mathrm{FLim}(f,c,\alpha) \to \varphi_{444}) \to \varphi_{464}$\lbl{A1}}
\BIC{$\varphi_{464}$}
\BIC{$\mathrm{FLim}(f,c,\alpha) \to \varphi_{445}$}
\end{ptree}
\noindent$\Pi_{103}$:
\begin{ptree}
\AXC{$\stackrel{\Pi_{93}}{\vdots}$}\noLine
\UIC{$0 = \lvert g(x) \rvert \to \lvert g(x) \rvert = 0$}
\AXC{$\stackrel{\Pi_{102}}{\vdots}$}\noLine
\UIC{$\mathrm{FLim}(f,c,\alpha) \to \varphi_{445}$}
\AXC{$(\mathrm{FLim}(f,c,\alpha) \to \varphi_{445}) \to \varphi_{465}$\lbl{A1}}
\BIC{$\varphi_{465}$}
\BIC{$\mathrm{FLim}(f,c,\alpha) \to \varphi_{444}$}
\end{ptree}
\noindent$\Pi_{104}$:
\begin{ptree}
\AXC{$\forall x\,(\neg x < x)$\lbl{[O1]}}
\AXC{$\forall x\,(\neg x < x) \to \neg 0 < 0$\lbl{A2}}
\BIC{$\neg 0 < 0$}
\AXC{$\stackrel{\Pi_{102}}{\vdots}$}\noLine
\UIC{$\mathrm{FLim}(f,c,\alpha) \to \varphi_{445}$}
\AXC{$\varphi_{575}$\lbl{A1}}
\BIC{$\neg 0 < 0 \to (0 = \lvert g(x) \rvert \to 0 < 0 \to 0 < \lvert g(x) \rvert) \to \varphi_{548}$}
\BIC{$(0 = \lvert g(x) \rvert \to 0 < 0 \to 0 < \lvert g(x) \rvert) \to \varphi_{548}$}
\end{ptree}
\noindent$\Pi_{105}$:
\begin{ptree}
\AXC{$z = w \to 0 < z \to 0 < w$\lbl{A6}}
\UIC{$\forall w\,(z = w \to 0 < z \to 0 < w)$}
\UIC{$\forall z\,\forall w\,(z = w \to 0 < z \to 0 < w)$}
\AXC{$\varphi_{17}$\lbl{A2}}
\BIC{$\forall w\,(0 = w \to 0 < 0 \to 0 < w)$}
\AXC{$\varphi_{68}$\lbl{A2}}
\BIC{$0 = \lvert g(x) \rvert \to 0 < 0 \to 0 < \lvert g(x) \rvert$}
\AXC{$\stackrel{\Pi_{104}}{\vdots}$}\noLine
\UIC{$(0 = \lvert g(x) \rvert \to 0 < 0 \to 0 < \lvert g(x) \rvert) \to \varphi_{548}$}
\BIC{$\varphi_{548}$}
\end{ptree}
\noindent$\Pi_{106}$:
\begin{ptree}
\AXC{$\mathrm{FLim}(f,c,\alpha) \to \mathrm{FLim}(g,c,\beta) \to \beta \neq 0 \to \varphi_{402}$\lbl{A1}}
\AXC{$\varphi_{530}$\lbl{A1}}
\BIC{$\mathrm{FLim}(f,c,\alpha) \to \mathrm{FLim}(g,c,\beta) \to \beta \neq 0 \to \varphi_{418}$}
\end{ptree}
\noindent$\Pi_{107}$:
\begin{ptree}
\AXC{$\forall x\,\forall y\,\forall z\,(x < y \to y \leq z \to x < z)$\lbl{[O3]}}
\AXC{$\varphi_{95}$\lbl{A2}}
\BIC{$\forall y\,\forall z\,(\lvert x - c \rvert < y \to y \leq z \to \lvert x - c \rvert < z)$}
\AXC{$\varphi_{128}$\lbl{A2}}
\BIC{$\forall z\,(\lvert x - c \rvert < r \to r \leq z \to \lvert x - c \rvert < z)$}
\AXC{$\varphi_{115}$\lbl{A2}}
\BIC{$\lvert x - c \rvert < r \to r \leq \delta \to \lvert x - c \rvert < \delta$}
\end{ptree}
\noindent$\Pi_{108}$:
\begin{ptree}
\AXC{$\stackrel{\Pi_{107}}{\vdots}$}\noLine
\UIC{$\lvert x - c \rvert < r \to r \leq \delta \to \lvert x - c \rvert < \delta$}
\AXC{$(\lvert x - c \rvert < r \to r \leq \delta \to \lvert x - c \rvert < \delta) \to \varphi_{568}$\lbl{A1}}
\BIC{$\varphi_{568}$}
\end{ptree}
\noindent$\Pi_{109}$:
\begin{ptree}
\AXC{$\stackrel{\Pi_{106}}{\vdots}$}\noLine
\UIC{$\mathrm{FLim}(f,c,\alpha) \to \mathrm{FLim}(g,c,\beta) \to \beta \neq 0 \to \varphi_{418}$}
\AXC{$\mathrm{FLim}(f,c,\alpha) \to \mathrm{FLim}(g,c,\beta) \to \beta \neq 0 \to \varphi_{426}$\lbl{A1}}
\AXC{$\stackrel{\Pi_{108}}{\vdots}$}\noLine
\UIC{$\varphi_{568}$}
\BIC{$\varphi_{536}$}
\BIC{$\mathrm{FLim}(f,c,\alpha) \to \mathrm{FLim}(g,c,\beta) \to \beta \neq 0 \to \varphi_{428}$}
\end{ptree}
\noindent$\Pi_{110}$:
\begin{ptree}
\AXC{$\mathrm{FLim}(f,c,\alpha) \to \mathrm{FLim}(g,c,\beta) \to \beta \neq 0 \to \varphi_{402}$\lbl{A1}}
\AXC{$\varphi_{529}$\lbl{A1}}
\BIC{$\mathrm{FLim}(f,c,\alpha) \to \mathrm{FLim}(g,c,\beta) \to \beta \neq 0 \to \varphi_{417}$}
\end{ptree}
\noindent$\Pi_{111}$:
\begin{ptree}
\AXC{$\forall x\,\forall y\,\forall z\,(x < y \to y \leq z \to x < z)$\lbl{[O3]}}
\AXC{$\varphi_{95}$\lbl{A2}}
\BIC{$\forall y\,\forall z\,(\lvert x - c \rvert < y \to y \leq z \to \lvert x - c \rvert < z)$}
\AXC{$\varphi_{128}$\lbl{A2}}
\BIC{$\forall z\,(\lvert x - c \rvert < r \to r \leq z \to \lvert x - c \rvert < z)$}
\AXC{$\varphi_{109}$\lbl{A2}}
\BIC{$\lvert x - c \rvert < r \to r \leq d \to \lvert x - c \rvert < d$}
\end{ptree}
\noindent$\Pi_{112}$:
\begin{ptree}
\AXC{$\stackrel{\Pi_{111}}{\vdots}$}\noLine
\UIC{$\lvert x - c \rvert < r \to r \leq d \to \lvert x - c \rvert < d$}
\AXC{$(\lvert x - c \rvert < r \to r \leq d \to \lvert x - c \rvert < d) \to \varphi_{567}$\lbl{A1}}
\BIC{$\varphi_{567}$}
\end{ptree}
\noindent$\Pi_{113}$:
\begin{ptree}
\AXC{$\stackrel{\Pi_{110}}{\vdots}$}\noLine
\UIC{$\mathrm{FLim}(f,c,\alpha) \to \mathrm{FLim}(g,c,\beta) \to \beta \neq 0 \to \varphi_{417}$}
\AXC{$\mathrm{FLim}(f,c,\alpha) \to \mathrm{FLim}(g,c,\beta) \to \beta \neq 0 \to \varphi_{426}$\lbl{A1}}
\AXC{$\stackrel{\Pi_{112}}{\vdots}$}\noLine
\UIC{$\varphi_{567}$}
\BIC{$\varphi_{535}$}
\BIC{$\mathrm{FLim}(f,c,\alpha) \to \mathrm{FLim}(g,c,\beta) \to \beta \neq 0 \to \varphi_{427}$}
\end{ptree}
\noindent$\Pi_{114}$:
\begin{ptree}
\AXC{$\forall x \in \mathbb{R}\,(\varphi_{46}) \to x \in \mathbb{R} \to \varphi_{46}$\lbl{A2}}
\AXC{$\mathrm{FLim}(f,c,\alpha) \to \varphi_{434}$\lbl{A1}}
\AXC{$\varphi_{547}$\lbl{A1}}
\BIC{$\mathrm{FLim}(f,c,\alpha) \to \mathrm{FLim}(g,c,\beta) \to \beta \neq 0 \to \varphi_{459}$}
\AXC{$\varphi_{565}$\lbl{A1}}
\BIC{$\varphi_{514}$}
\BIC{$\mathrm{FLim}(f,c,\alpha) \to \mathrm{FLim}(g,c,\beta) \to \beta \neq 0 \to \varphi_{458}$}
\AXC{$\varphi_{580}$\lbl{A1}}
\BIC{$(\mathrm{FLim}(f,c,\alpha) \to \varphi_{411}) \to \varphi_{569}$}
\end{ptree}
\noindent$\Pi_{115}$:
\begin{ptree}
\AXC{$\mathrm{FLim}(f,c,\alpha) \to \mathrm{FLim}(g,c,\beta) \to \beta \neq 0 \to \varphi_{412}$\lbl{A1}}
\AXC{$\mathrm{FLim}(f,c,\alpha) \to \varphi_{411}$\lbl{A1}}
\AXC{$\stackrel{\Pi_{114}}{\vdots}$}\noLine
\UIC{$(\mathrm{FLim}(f,c,\alpha) \to \varphi_{411}) \to \varphi_{569}$}
\BIC{$\varphi_{569}$}
\BIC{$\varphi_{540}$}
\end{ptree}
\noindent$\Pi_{116}$:
\begin{ptree}
\AXC{$\stackrel{\Pi_{113}}{\vdots}$}\noLine
\UIC{$\mathrm{FLim}(f,c,\alpha) \to \mathrm{FLim}(g,c,\beta) \to \beta \neq 0 \to \varphi_{427}$}
\AXC{$\stackrel{\Pi_{115}}{\vdots}$}\noLine
\UIC{$\varphi_{540}$}
\BIC{$\mathrm{FLim}(f,c,\alpha) \to \varphi_{443}$}
\end{ptree}
\noindent$\Pi_{117}$:
\begin{ptree}
\AXC{$\mathrm{FLim}(f,c,\alpha) \to \mathrm{FLim}(g,c,\beta) \to \beta \neq 0 \to \varphi_{382}$\lbl{A1}}
\AXC{$\varphi_{500}$\lbl{A1}}
\BIC{$\mathrm{FLim}(f,c,\alpha) \to \mathrm{FLim}(g,c,\beta) \to \beta \neq 0 \to \varphi_{359}$}
\AXC{$\varphi_{532}$\lbl{A1}}
\BIC{$\varphi_{469}$}
\end{ptree}
\noindent$\Pi_{118}$:
\begin{ptree}
\AXC{$\stackrel{\Pi_{116}}{\vdots}$}\noLine
\UIC{$\mathrm{FLim}(f,c,\alpha) \to \varphi_{443}$}
\AXC{$\forall x \in \mathrm{dom}(g)\,(\varphi_{94}) \to x \in \mathrm{dom}(g) \to \varphi_{94}$\lbl{A2}}
\AXC{$\stackrel{\Pi_{117}}{\vdots}$}\noLine
\UIC{$\varphi_{469}$}
\BIC{$\mathrm{FLim}(f,c,\alpha) \to \mathrm{FLim}(g,c,\beta) \to \beta \neq 0 \to \varphi_{352}$}
\AXC{$\varphi_{579}$\lbl{A1}}
\BIC{$(\mathrm{FLim}(f,c,\alpha) \to \varphi_{443}) \to \varphi_{571}$}
\BIC{$\varphi_{571}$}
\end{ptree}
\noindent$\Pi_{119}$:
\begin{ptree}
\AXC{$\stackrel{\Pi_{109}}{\vdots}$}\noLine
\UIC{$\mathrm{FLim}(f,c,\alpha) \to \mathrm{FLim}(g,c,\beta) \to \beta \neq 0 \to \varphi_{428}$}
\AXC{$\mathrm{FLim}(f,c,\alpha) \to \mathrm{FLim}(g,c,\beta) \to \beta \neq 0 \to \varphi_{412}$\lbl{A1}}
\AXC{$\stackrel{\Pi_{118}}{\vdots}$}\noLine
\UIC{$\varphi_{571}$}
\BIC{$\varphi_{543}$}
\BIC{$\mathrm{FLim}(f,c,\alpha) \to \varphi_{453}$}
\end{ptree}
\noindent$\Pi_{120}$:
\begin{ptree}
\AXC{$\forall f\,\forall c\,\forall \alpha\,(\mathrm{FLim}(f,c,\alpha) \to \alpha \in \mathbb{R})$\lbl{[D4]}}
\AXC{$\varphi_{69}$\lbl{A2}}
\BIC{$\forall c\,\forall \alpha\,(\mathrm{FLim}(g,c,\alpha) \to \alpha \in \mathbb{R})$}
\AXC{$\varphi_{50}$\lbl{A2}}
\BIC{$\forall \alpha\,(\mathrm{FLim}(g,c,\alpha) \to \alpha \in \mathbb{R})$}
\AXC{$\varphi_{30}$\lbl{A2}}
\BIC{$\mathrm{FLim}(g,c,\beta) \to \beta \in \mathbb{R}$}
\end{ptree}
\noindent$\Pi_{121}$:
\begin{ptree}
\AXC{$\varphi_{2}$\lbl{A1}}
\AXC{$\stackrel{\Pi_{120}}{\vdots}$}\noLine
\UIC{$\mathrm{FLim}(g,c,\beta) \to \beta \in \mathbb{R}$}
\AXC{$(\mathrm{FLim}(g,c,\beta) \to \beta \in \mathbb{R}) \to \varphi_{2} \to \varphi_{9}$\lbl{A1}}
\BIC{$\varphi_{2} \to \varphi_{9}$}
\BIC{$\varphi_{9}$}
\end{ptree}
\noindent$\Pi_{122}$:
\begin{ptree}
\AXC{$\forall f\,\forall x\,(\mathrm{RFn}(f) \to x \in \mathrm{dom}(f) \to f(x) \in \mathbb{R})$\lbl{[S6]}}
\AXC{$\varphi_{129}$\lbl{A2}}
\BIC{$\forall x\,(\mathrm{RFn}(g) \to x \in \mathrm{dom}(g) \to g(x) \in \mathbb{R})$}
\AXC{$\varphi_{110}$\lbl{A2}}
\BIC{$\mathrm{RFn}(g) \to x \in \mathrm{dom}(g) \to g(x) \in \mathbb{R}$}
\end{ptree}
\noindent$\Pi_{123}$:
\begin{ptree}
\AXC{$\stackrel{\Pi_{3}}{\vdots}$}\noLine
\UIC{$\varphi_{4}$}
\AXC{$\stackrel{\Pi_{122}}{\vdots}$}\noLine
\UIC{$\mathrm{RFn}(g) \to x \in \mathrm{dom}(g) \to g(x) \in \mathbb{R}$}
\AXC{$(\mathrm{RFn}(g) \to x \in \mathrm{dom}(g) \to g(x) \in \mathbb{R}) \to \varphi_{4} \to \varphi_{537}$\lbl{A1}}
\BIC{$\varphi_{4} \to \varphi_{537}$}
\BIC{$\varphi_{537}$}
\end{ptree}
\noindent$\Pi_{124}$:
\begin{ptree}
\AXC{$\varphi_{138}$\lbl{[E26]}}
\AXC{$\varphi_{138} \to \forall \beta\,(g(x) \in \mathbb{R} \to \beta \in \mathbb{R} \to \varphi_{87})$\lbl{A2}}
\BIC{$\forall \beta\,(g(x) \in \mathbb{R} \to \beta \in \mathbb{R} \to \varphi_{87})$}
\AXC{$\varphi_{370}$\lbl{A2}}
\BIC{$g(x) \in \mathbb{R} \to \beta \in \mathbb{R} \to \varphi_{87}$}
\AXC{$(g(x) \in \mathbb{R} \to \beta \in \mathbb{R} \to \varphi_{87}) \to \varphi_{573}$\lbl{A1}}
\BIC{$\varphi_{573}$}
\end{ptree}
\noindent$\Pi_{125}$:
\begin{ptree}
\AXC{$\stackrel{\Pi_{119}}{\vdots}$}\noLine
\UIC{$\mathrm{FLim}(f,c,\alpha) \to \varphi_{453}$}
\AXC{$\stackrel{\Pi_{121}}{\vdots}$}\noLine
\UIC{$\varphi_{9}$}
\AXC{$\stackrel{\Pi_{116}}{\vdots}$}\noLine
\UIC{$\mathrm{FLim}(f,c,\alpha) \to \varphi_{443}$}
\AXC{$\varphi_{538}$\lbl{A1}}
\BIC{$\mathrm{FLim}(f,c,\alpha) \to \mathrm{FLim}(g,c,\beta) \to \beta \neq 0 \to \varphi_{423}$}
\AXC{$\stackrel{\Pi_{123}}{\vdots}$}\noLine
\UIC{$\varphi_{537}$}
\BIC{$\mathrm{FLim}(f,c,\alpha) \to \mathrm{FLim}(g,c,\beta) \to \beta \neq 0 \to \varphi_{429}$}
\AXC{$\stackrel{\Pi_{124}}{\vdots}$}\noLine
\UIC{$\varphi_{573}$}
\BIC{$\varphi_{550}$}
\BIC{$(\mathrm{FLim}(f,c,\alpha) \to \varphi_{453}) \to \mathrm{FLim}(f,c,\alpha) \to \varphi_{448}$}
\BIC{$\mathrm{FLim}(f,c,\alpha) \to \varphi_{448}$}
\end{ptree}
\noindent$\Pi_{126}$:
\begin{ptree}
\AXC{$\varphi_{3}$\lbl{A1}}
\AXC{$\stackrel{\Pi_{121}}{\vdots}$}\noLine
\UIC{$\varphi_{9}$}
\AXC{$\forall x \in \mathbb{R}\,(x \neq 0 \to 0 < \lvert x \rvert)$\lbl{[A14]}}
\AXC{$\varphi_{63}$\lbl{A2}}
\BIC{$\beta \in \mathbb{R} \to \beta \neq 0 \to 0 < \lvert \beta \rvert$}
\AXC{$(\beta \in \mathbb{R} \to \beta \neq 0 \to 0 < \lvert \beta \rvert) \to \varphi_{9} \to \varphi_{3} \to \varphi_{11}$\lbl{A1}}
\BIC{$\varphi_{9} \to \varphi_{3} \to \varphi_{11}$}
\BIC{$\varphi_{3} \to \varphi_{11}$}
\BIC{$\varphi_{11}$}
\end{ptree}
\noindent$\Pi_{127}$:
\begin{ptree}
\AXC{$\forall x \in \mathbb{R}\,(\lvert x \rvert \in \mathbb{R})$\lbl{[A5]}}
\AXC{$\forall x \in \mathbb{R}\,(\lvert x \rvert \in \mathbb{R}) \to \beta \in \mathbb{R} \to \lvert \beta \rvert \in \mathbb{R}$\lbl{A2}}
\BIC{$\beta \in \mathbb{R} \to \lvert \beta \rvert \in \mathbb{R}$}
\AXC{$(\beta \in \mathbb{R} \to \lvert \beta \rvert \in \mathbb{R}) \to \varphi_{9} \to \varphi_{18}$\lbl{A1}}
\BIC{$\varphi_{9} \to \varphi_{18}$}
\end{ptree}
\noindent$\Pi_{128}$:
\begin{ptree}
\AXC{$\stackrel{\Pi_{121}}{\vdots}$}\noLine
\UIC{$\varphi_{9}$}
\AXC{$\stackrel{\Pi_{127}}{\vdots}$}\noLine
\UIC{$\varphi_{9} \to \varphi_{18}$}
\BIC{$\varphi_{18}$}
\end{ptree}
\noindent$\Pi_{129}$:
\begin{ptree}
\AXC{$\stackrel{\Pi_{128}}{\vdots}$}\noLine
\UIC{$\varphi_{18}$}
\AXC{$\forall e \in \mathbb{R}\,(0 < e \to 0 < e \cdot 2^{-1})$\lbl{[F35]}}
\AXC{$\varphi_{152}$\lbl{A2}}
\BIC{$\lvert \beta \rvert \in \mathbb{R} \to 0 < \lvert \beta \rvert \to 0 < \lvert \beta \rvert \cdot 2^{-1}$}
\AXC{$(\lvert \beta \rvert \in \mathbb{R} \to 0 < \lvert \beta \rvert \to 0 < \lvert \beta \rvert \cdot 2^{-1}) \to \varphi_{18} \to \varphi_{11} \to \varphi_{33}$\lbl{A1}}
\BIC{$\varphi_{18} \to \varphi_{11} \to \varphi_{33}$}
\BIC{$\varphi_{11} \to \varphi_{33}$}
\end{ptree}
\noindent$\Pi_{130}$:
\begin{ptree}
\AXC{$\stackrel{\Pi_{126}}{\vdots}$}\noLine
\UIC{$\varphi_{11}$}
\AXC{$\stackrel{\Pi_{129}}{\vdots}$}\noLine
\UIC{$\varphi_{11} \to \varphi_{33}$}
\BIC{$\varphi_{33}$}
\end{ptree}
\noindent$\Pi_{131}$:
\begin{ptree}
\AXC{$\forall x\,\forall y\,\forall z\,(x < y \to y < z \to x < z)$\lbl{[O2]}}
\AXC{$\varphi_{51}$\lbl{A2}}
\BIC{$\forall y\,\forall z\,(0 < y \to y < z \to 0 < z)$}
\AXC{$\varphi_{121}$\lbl{A2}}
\BIC{$\forall z\,(0 < \lvert \beta \rvert \cdot 2^{-1} \to \lvert \beta \rvert \cdot 2^{-1} < z \to 0 < z)$}
\AXC{$\forall z\,(0 < \lvert \beta \rvert \cdot 2^{-1} \to \lvert \beta \rvert \cdot 2^{-1} < z \to 0 < z) \to \varphi_{96}$\lbl{A2}}
\BIC{$\varphi_{96}$}
\end{ptree}
\noindent$\Pi_{132}$:
\begin{ptree}
\AXC{$\stackrel{\Pi_{125}}{\vdots}$}\noLine
\UIC{$\mathrm{FLim}(f,c,\alpha) \to \varphi_{448}$}
\AXC{$\stackrel{\Pi_{130}}{\vdots}$}\noLine
\UIC{$\varphi_{33}$}
\AXC{$\stackrel{\Pi_{131}}{\vdots}$}\noLine
\UIC{$\varphi_{96}$}
\AXC{$\varphi_{96} \to \varphi_{551}$\lbl{A1}}
\BIC{$\varphi_{551}$}
\BIC{$(\mathrm{FLim}(f,c,\alpha) \to \varphi_{448}) \to \mathrm{FLim}(f,c,\alpha) \to \varphi_{446}$}
\BIC{$\mathrm{FLim}(f,c,\alpha) \to \varphi_{446}$}
\AXC{$(\mathrm{FLim}(f,c,\alpha) \to \varphi_{446}) \to \varphi_{541}$\lbl{A1}}
\BIC{$\varphi_{541}$}
\end{ptree}
\noindent$\Pi_{133}$:
\begin{ptree}
\AXC{$\stackrel{\Pi_{103}}{\vdots}$}\noLine
\UIC{$\mathrm{FLim}(f,c,\alpha) \to \varphi_{444}$}
\AXC{$z = w \to 0 < z \to 0 < w$\lbl{A6}}
\UIC{$\forall w\,(z = w \to 0 < z \to 0 < w)$}
\UIC{$\forall z\,\forall w\,(z = w \to 0 < z \to 0 < w)$}
\AXC{$\varphi_{89}$\lbl{A2}}
\BIC{$\forall w\,(\lvert g(x) \rvert = w \to 0 < \lvert g(x) \rvert \to 0 < w)$}
\AXC{$\varphi_{134}$\lbl{A2}}
\BIC{$\lvert g(x) \rvert = 0 \to 0 < \lvert g(x) \rvert \to 0 < 0$}
\AXC{$\stackrel{\Pi_{105}}{\vdots}$}\noLine
\UIC{$\varphi_{548}$}
\BIC{$(\mathrm{FLim}(f,c,\alpha) \to \varphi_{444}) \to \mathrm{FLim}(f,c,\alpha) \to \varphi_{447}$}
\BIC{$\mathrm{FLim}(f,c,\alpha) \to \varphi_{447}$}
\AXC{$\stackrel{\Pi_{132}}{\vdots}$}\noLine
\UIC{$\varphi_{541}$}
\BIC{$\mathrm{FLim}(f,c,\alpha) \to \mathrm{FLim}(g,c,\beta) \to \beta \neq 0 \to \varphi_{422}$}
\end{ptree}
\noindent$\Pi_{134}$:
\begin{ptree}
\AXC{$\stackrel{\Pi_{133}}{\vdots}$}\noLine
\UIC{$\mathrm{FLim}(f,c,\alpha) \to \mathrm{FLim}(g,c,\beta) \to \beta \neq 0 \to \varphi_{422}$}
\AXC{$\stackrel{\Pi_{116}}{\vdots}$}\noLine
\UIC{$\mathrm{FLim}(f,c,\alpha) \to \varphi_{443}$}
\AXC{$(\mathrm{FLim}(f,c,\alpha) \to \varphi_{443}) \to \varphi_{553}$\lbl{A1}}
\BIC{$\varphi_{553}$}
\BIC{$\varphi_{478}$}
\end{ptree}
\noindent$\Pi_{135}$:
\begin{ptree}
\AXC{$\forall x\,(\varphi_{42}) \to \varphi_{42}$\lbl{A2}}
\AXC{$\stackrel{\Pi_{8}}{\vdots}$}\noLine
\UIC{$\varphi_{168}$}
\AXC{$\varphi_{168} \to \varphi_{377}$\lbl{A1}}
\BIC{$\varphi_{377}$}
\BIC{$\mathrm{FLim}(f,c,\alpha) \to \mathrm{FLim}(g,c,\beta) \to \beta \neq 0 \to \varphi_{42}$}
\AXC{$\stackrel{\Pi_{134}}{\vdots}$}\noLine
\UIC{$\varphi_{478}$}
\BIC{$\mathrm{FLim}(f,c,\alpha) \to \varphi_{436}$}
\AXC{$(\mathrm{FLim}(f,c,\alpha) \to \varphi_{436}) \to \varphi_{404} \to \varphi_{65}$\lbl{A1}}
\BIC{$\varphi_{404} \to \varphi_{65}$}
\UIC{$\forall x\,(\varphi_{404} \to \varphi_{65})$}
\end{ptree}
\noindent$\Pi_{136}$:
\begin{ptree}
\AXC{$\varphi_{404} \to \forall x\,(\varphi_{404})$\lbl{A4}}
\AXC{$\forall x\,(\varphi_{404} \to \varphi_{65}) \to \forall x\,(\varphi_{404}) \to \varphi_{78}$\lbl{A3}}
\AXC{$\stackrel{\Pi_{135}}{\vdots}$}\noLine
\UIC{$\forall x\,(\varphi_{404} \to \varphi_{65})$}
\AXC{$\forall x\,(\varphi_{404} \to \varphi_{65}) \to \varphi_{578}$\lbl{A1}}
\BIC{$\varphi_{578}$}
\BIC{$(\varphi_{404} \to \forall x\,(\varphi_{404})) \to \mathrm{FLim}(f,c,\alpha) \to \varphi_{438}$}
\BIC{$\mathrm{FLim}(f,c,\alpha) \to \varphi_{438}$}
\end{ptree}
\noindent$\Pi_{137}$:
\begin{ptree}
\AXC{$\stackrel{\Pi_{136}}{\vdots}$}\noLine
\UIC{$\mathrm{FLim}(f,c,\alpha) \to \varphi_{438}$}
\AXC{$\mathrm{FLim}(f,c,\alpha) \to \mathrm{FLim}(g,c,\beta) \to \beta \neq 0 \to \varphi_{402}$\lbl{A1}}
\AXC{$\varphi_{570}$\lbl{A1}}
\BIC{$(\mathrm{FLim}(f,c,\alpha) \to \varphi_{438}) \to \mathrm{FLim}(f,c,\alpha) \to \varphi_{455}$}
\BIC{$\mathrm{FLim}(f,c,\alpha) \to \varphi_{455}$}
\end{ptree}
\noindent$\Pi_{138}$:
\begin{ptree}
\AXC{$\varphi_{323}$\lbl{A2}}
\AXC{$\stackrel{\Pi_{137}}{\vdots}$}\noLine
\UIC{$\mathrm{FLim}(f,c,\alpha) \to \varphi_{455}$}
\AXC{$(\mathrm{FLim}(f,c,\alpha) \to \varphi_{455}) \to \varphi_{516}$\lbl{A1}}
\BIC{$\varphi_{516}$}
\BIC{$\mathrm{FLim}(f,c,\alpha) \to \mathrm{FLim}(g,c,\beta) \to \beta \neq 0 \to \varphi_{440}$}
\end{ptree}
\noindent$\Pi_{139}$:
\begin{ptree}
\AXC{$\stackrel{\Pi_{138}}{\vdots}$}\noLine
\UIC{$\mathrm{FLim}(f,c,\alpha) \to \mathrm{FLim}(g,c,\beta) \to \beta \neq 0 \to \varphi_{440}$}
\AXC{$\varphi_{183}$\lbl{DefPNbhd}}
\AXC{$\varphi_{183} \to \varphi_{204}$\lbl{A2}}
\BIC{$\varphi_{204}$}
\AXC{$\varphi_{383}$\lbl{A2}}
\BIC{$\mathrm{PNbhd}((f / g),c) \leftrightarrow \exists r\,(r \in \mathbb{R} \land 0 < r \land \varphi_{78})$}
\AXC{$\varphi_{554}$\lbl{A1}}
\BIC{$\varphi_{534}$}
\BIC{$\mathrm{FLim}(f,c,\alpha) \to \mathrm{FLim}(g,c,\beta) \to \beta \neq 0 \to \varphi_{394}$}
\end{ptree}
\noindent$\Pi_{140}$:
\begin{ptree}
\AXC{$\stackrel{\Pi_{139}}{\vdots}$}\noLine
\UIC{$\mathrm{FLim}(f,c,\alpha) \to \mathrm{FLim}(g,c,\beta) \to \beta \neq 0 \to \varphi_{394}$}
\AXC{$\varphi_{523}$\lbl{A1}}
\BIC{$\varphi_{395} \to \neg (r \in \mathbb{R} \land 0 < r \land r \leq d \land r \leq \delta)$}
\UIC{$\forall r\,(\varphi_{395} \to \neg (r \in \mathbb{R} \land 0 < r \land r \leq d \land r \leq \delta))$}
\end{ptree}
\noindent$\Pi_{141}$:
\begin{ptree}
\AXC{$\mathrm{FLim}(f,c,\alpha) \to \mathrm{FLim}(g,c,\beta) \to \beta \neq 0 \to \varphi_{382}$\lbl{A1}}
\AXC{$\varphi_{470}$\lbl{A1}}
\BIC{$\mathrm{FLim}(f,c,\alpha) \to \varphi_{262}$}
\end{ptree}
\noindent$\Pi_{142}$:
\begin{ptree}
\AXC{$\forall d\,(\forall u_{4}\,(\varphi_{85}) \to \varphi_{62}) \to \forall d\,\forall u_{4}\,(\varphi_{85}) \to \forall d\,(\varphi_{62})$\lbl{A3}}
\AXC{$\forall u_{4}\,(\varphi_{85}) \to \varphi_{62}$\lbl{A2}}
\UIC{$\forall d\,(\forall u_{4}\,(\varphi_{85}) \to \varphi_{62})$}
\AXC{$\forall d\,(\forall u_{4}\,(\varphi_{85}) \to \varphi_{62}) \to \varphi_{522}$\lbl{A1}}
\BIC{$\varphi_{522}$}
\BIC{$(\forall u_{4}\,(\varphi_{85}) \to \forall d\,\forall u_{4}\,(\varphi_{85})) \to \forall u_{4}\,(\varphi_{85}) \to \forall d\,(\varphi_{62})$}
\end{ptree}
\noindent$\Pi_{143}$:
\begin{ptree}
\AXC{$\forall u_{4}\,(\varphi_{85}) \to \forall d\,\forall u_{4}\,(\varphi_{85})$\lbl{A4}}
\AXC{$\stackrel{\Pi_{142}}{\vdots}$}\noLine
\UIC{$(\forall u_{4}\,(\varphi_{85}) \to \forall d\,\forall u_{4}\,(\varphi_{85})) \to \forall u_{4}\,(\varphi_{85}) \to \forall d\,(\varphi_{62})$}
\BIC{$\forall u_{4}\,(\varphi_{85}) \to \forall d\,(\varphi_{62})$}
\end{ptree}
\noindent$\Pi_{144}$:
\begin{ptree}
\AXC{$\forall u_{4}\,(\forall d\,(\varphi_{62}) \to \varphi_{85}) \to \forall u_{4}\,\forall d\,(\varphi_{62}) \to \forall u_{4}\,(\varphi_{85})$\lbl{A3}}
\AXC{$\forall d\,(\varphi_{62}) \to \varphi_{85}$\lbl{A2}}
\UIC{$\forall u_{4}\,(\forall d\,(\varphi_{62}) \to \varphi_{85})$}
\AXC{$\forall u_{4}\,(\forall d\,(\varphi_{62}) \to \varphi_{85}) \to \varphi_{520}$\lbl{A1}}
\BIC{$\varphi_{520}$}
\BIC{$(\forall d\,(\varphi_{62}) \to \forall u_{4}\,\forall d\,(\varphi_{62})) \to \forall d\,(\varphi_{62}) \to \forall u_{4}\,(\varphi_{85})$}
\end{ptree}
\noindent$\Pi_{145}$:
\begin{ptree}
\AXC{$\forall d\,(\varphi_{62}) \to \forall u_{4}\,\forall d\,(\varphi_{62})$\lbl{A4}}
\AXC{$\stackrel{\Pi_{144}}{\vdots}$}\noLine
\UIC{$(\forall d\,(\varphi_{62}) \to \forall u_{4}\,\forall d\,(\varphi_{62})) \to \forall d\,(\varphi_{62}) \to \forall u_{4}\,(\varphi_{85})$}
\BIC{$\forall d\,(\varphi_{62}) \to \forall u_{4}\,(\varphi_{85})$}
\end{ptree}
\noindent$\Pi_{146}$:
\begin{ptree}
\AXC{$\stackrel{\Pi_{145}}{\vdots}$}\noLine
\UIC{$\forall d\,(\varphi_{62}) \to \forall u_{4}\,(\varphi_{85})$}
\AXC{$\forall d\,(\varphi_{62}) \leftrightarrow \forall d\,(\varphi_{62})$\lbl{A1}}
\AXC{$(\forall d\,(\varphi_{62}) \leftrightarrow \forall d\,(\varphi_{62})) \to \varphi_{477}$\lbl{A1}}
\BIC{$\varphi_{477}$}
\BIC{$(\forall u_{4}\,(\varphi_{85}) \to \forall d\,(\varphi_{62})) \to (\forall d\,(\varphi_{62}) \leftrightarrow \forall u_{4}\,(\varphi_{85}))$}
\end{ptree}
\noindent$\Pi_{147}$:
\begin{ptree}
\AXC{$\stackrel{\Pi_{143}}{\vdots}$}\noLine
\UIC{$\forall u_{4}\,(\varphi_{85}) \to \forall d\,(\varphi_{62})$}
\AXC{$\stackrel{\Pi_{146}}{\vdots}$}\noLine
\UIC{$(\forall u_{4}\,(\varphi_{85}) \to \forall d\,(\varphi_{62})) \to (\forall d\,(\varphi_{62}) \leftrightarrow \forall u_{4}\,(\varphi_{85}))$}
\BIC{$\forall d\,(\varphi_{62}) \leftrightarrow \forall u_{4}\,(\varphi_{85})$}
\AXC{$\varphi_{415}$\lbl{A1}}
\BIC{$\neg \forall d\,(\varphi_{62}) \leftrightarrow \neg \forall u_{4}\,(\varphi_{85})$}
\end{ptree}
\noindent$\Pi_{148}$:
\begin{ptree}
\AXC{$\stackrel{\Pi_{147}}{\vdots}$}\noLine
\UIC{$\neg \forall d\,(\varphi_{62}) \leftrightarrow \neg \forall u_{4}\,(\varphi_{85})$}
\AXC{$0 < v_{767} \leftrightarrow 0 < v_{767}$\lbl{A1}}
\AXC{$(0 < v_{767} \leftrightarrow 0 < v_{767}) \to \varphi_{433}$\lbl{A1}}
\BIC{$\varphi_{433}$}
\BIC{$0 < v_{767} \to \neg \forall d\,(\varphi_{62}) \leftrightarrow 0 < v_{767} \to \neg \forall u_{4}\,(\varphi_{85})$}
\end{ptree}
\noindent$\Pi_{149}$:
\begin{ptree}
\AXC{$\stackrel{\Pi_{148}}{\vdots}$}\noLine
\UIC{$0 < v_{767} \to \neg \forall d\,(\varphi_{62}) \leftrightarrow 0 < v_{767} \to \neg \forall u_{4}\,(\varphi_{85})$}
\AXC{$0 < v_{766} \leftrightarrow 0 < v_{766}$\lbl{A1}}
\AXC{$(0 < v_{766} \leftrightarrow 0 < v_{766}) \to \varphi_{457}$\lbl{A1}}
\BIC{$\varphi_{457}$}
\BIC{$\varphi_{43} \leftrightarrow 0 < v_{766} \to 0 < v_{767} \to \neg \forall u_{4}\,(\varphi_{85})$}
\end{ptree}
\noindent$\Pi_{150}$:
\begin{ptree}
\AXC{$\stackrel{\Pi_{149}}{\vdots}$}\noLine
\UIC{$\varphi_{43} \leftrightarrow 0 < v_{766} \to 0 < v_{767} \to \neg \forall u_{4}\,(\varphi_{85})$}
\AXC{$v_{767} \in \mathbb{R} \leftrightarrow v_{767} \in \mathbb{R}$\lbl{A1}}
\AXC{$\varphi_{479}$\lbl{A1}}
\BIC{$(\varphi_{43} \leftrightarrow 0 < v_{766} \to 0 < v_{767} \to \neg \forall u_{4}\,(\varphi_{85})) \to \varphi_{376}$}
\BIC{$\varphi_{376}$}
\end{ptree}
\noindent$\Pi_{151}$:
\begin{ptree}
\AXC{$\stackrel{\Pi_{150}}{\vdots}$}\noLine
\UIC{$\varphi_{376}$}
\AXC{$v_{766} \in \mathbb{R} \leftrightarrow v_{766} \in \mathbb{R}$\lbl{A1}}
\AXC{$\varphi_{497}$\lbl{A1}}
\BIC{$\varphi_{376} \to (v_{766} \in \mathbb{R} \to v_{767} \in \mathbb{R} \to \varphi_{43} \leftrightarrow \varphi_{230})$}
\BIC{$v_{766} \in \mathbb{R} \to v_{767} \in \mathbb{R} \to \varphi_{43} \leftrightarrow \varphi_{230}$}
\end{ptree}
\noindent$\Pi_{152}$:
\begin{ptree}
\AXC{$\varphi_{505}$\lbl{A3}}
\AXC{$\stackrel{\Pi_{151}}{\vdots}$}\noLine
\UIC{$v_{766} \in \mathbb{R} \to v_{767} \in \mathbb{R} \to \varphi_{43} \leftrightarrow \varphi_{230}$}
\AXC{$\varphi_{498}$\lbl{A1}}
\BIC{$\varphi_{230} \to v_{766} \in \mathbb{R} \to v_{767} \in \mathbb{R} \to \varphi_{43}$}
\UIC{$\forall v_{767}\,(\varphi_{230} \to v_{766} \in \mathbb{R} \to v_{767} \in \mathbb{R} \to \varphi_{43})$}
\AXC{$\stackrel{\Pi_{151}}{\vdots}$}\noLine
\UIC{$v_{766} \in \mathbb{R} \to v_{767} \in \mathbb{R} \to \varphi_{43} \leftrightarrow \varphi_{230}$}
\AXC{$\varphi_{499}$\lbl{A1}}
\BIC{$(v_{766} \in \mathbb{R} \to v_{767} \in \mathbb{R} \to \varphi_{43}) \to \varphi_{230}$}
\UIC{$\forall v_{767}\,((v_{766} \in \mathbb{R} \to v_{767} \in \mathbb{R} \to \varphi_{43}) \to \varphi_{230})$}
\AXC{$\varphi_{581}$\lbl{A1}}
\BIC{$\forall v_{767}\,(\varphi_{230} \to v_{766} \in \mathbb{R} \to v_{767} \in \mathbb{R} \to \varphi_{43}) \to \varphi_{505} \to \varphi_{552}$}
\BIC{$\varphi_{505} \to \varphi_{552}$}
\BIC{$\varphi_{552}$}
\end{ptree}
\noindent$\Pi_{153}$:
\begin{ptree}
\AXC{$\varphi_{504}$\lbl{A3}}
\AXC{$\stackrel{\Pi_{152}}{\vdots}$}\noLine
\UIC{$\varphi_{552}$}
\BIC{$\forall v_{767}\,(v_{766} \in \mathbb{R} \to v_{767} \in \mathbb{R} \to \varphi_{43}) \leftrightarrow \forall v_{767}\,(\varphi_{230})$}
\end{ptree}
\noindent$\Pi_{154}$:
\begin{ptree}
\AXC{$\stackrel{\Pi_{153}}{\vdots}$}\noLine
\UIC{$\forall v_{767}\,(v_{766} \in \mathbb{R} \to v_{767} \in \mathbb{R} \to \varphi_{43}) \leftrightarrow \forall v_{767}\,(\varphi_{230})$}
\AXC{$\varphi_{507}$\lbl{A1}}
\BIC{$\forall v_{767}\,(v_{766} \in \mathbb{R} \to v_{767} \in \mathbb{R} \to \varphi_{43}) \to \forall v_{767}\,(\varphi_{230})$}
\UIC{$\varphi_{397}$}
\AXC{$\varphi_{397} \to \varphi_{396} \to \varphi_{576}$\lbl{A1}}
\BIC{$\varphi_{396} \to \varphi_{576}$}
\end{ptree}
\noindent$\Pi_{155}$:
\begin{ptree}
\AXC{$\stackrel{\Pi_{153}}{\vdots}$}\noLine
\UIC{$\forall v_{767}\,(v_{766} \in \mathbb{R} \to v_{767} \in \mathbb{R} \to \varphi_{43}) \leftrightarrow \forall v_{767}\,(\varphi_{230})$}
\AXC{$\varphi_{506}$\lbl{A1}}
\BIC{$\forall v_{767}\,(\varphi_{230}) \to \forall v_{767}\,(v_{766} \in \mathbb{R} \to v_{767} \in \mathbb{R} \to \varphi_{43})$}
\UIC{$\varphi_{396}$}
\AXC{$\stackrel{\Pi_{154}}{\vdots}$}\noLine
\UIC{$\varphi_{396} \to \varphi_{576}$}
\BIC{$\varphi_{576}$}
\end{ptree}
\noindent$\Pi_{156}$:
\begin{ptree}
\AXC{$\varphi_{396} \to \varphi_{187} \to \forall v_{766}\,\forall v_{767}\,(v_{766} \in \mathbb{R} \to v_{767} \in \mathbb{R} \to \varphi_{43})$\lbl{A3}}
\AXC{$\varphi_{397} \to \forall v_{766}\,\forall v_{767}\,(v_{766} \in \mathbb{R} \to v_{767} \in \mathbb{R} \to \varphi_{43}) \to \varphi_{187}$\lbl{A3}}
\AXC{$\stackrel{\Pi_{155}}{\vdots}$}\noLine
\UIC{$\varphi_{576}$}
\BIC{$\varphi_{560}$}
\BIC{$\forall v_{766}\,\forall v_{767}\,(v_{766} \in \mathbb{R} \to v_{767} \in \mathbb{R} \to \varphi_{43}) \leftrightarrow \varphi_{187}$}
\end{ptree}
\noindent$\Pi_{157}$:
\begin{ptree}
\AXC{$\varphi_{105}$\lbl{[D8]}}
\AXC{$\varphi_{105} \to \varphi_{116}$\lbl{A2}}
\BIC{$\varphi_{116}$}
\AXC{$\varphi_{116} \to v_{766} \in \mathbb{R} \to v_{767} \in \mathbb{R} \to \varphi_{43}$\lbl{A2}}
\BIC{$v_{766} \in \mathbb{R} \to v_{767} \in \mathbb{R} \to \varphi_{43}$}
\UIC{$\forall v_{767}\,(v_{766} \in \mathbb{R} \to v_{767} \in \mathbb{R} \to \varphi_{43})$}
\UIC{$\forall v_{766}\,\forall v_{767}\,(v_{766} \in \mathbb{R} \to v_{767} \in \mathbb{R} \to \varphi_{43})$}
\AXC{$\forall v_{766}\,\forall v_{767}\,(v_{766} \in \mathbb{R} \to v_{767} \in \mathbb{R} \to \varphi_{43}) \to \varphi_{462}$\lbl{A1}}
\BIC{$\varphi_{462}$}
\end{ptree}
\noindent$\Pi_{158}$:
\begin{ptree}
\AXC{$\stackrel{\Pi_{156}}{\vdots}$}\noLine
\UIC{$\forall v_{766}\,\forall v_{767}\,(v_{766} \in \mathbb{R} \to v_{767} \in \mathbb{R} \to \varphi_{43}) \leftrightarrow \varphi_{187}$}
\AXC{$\stackrel{\Pi_{157}}{\vdots}$}\noLine
\UIC{$\varphi_{462}$}
\BIC{$\varphi_{187}$}
\AXC{$\varphi_{187} \to \varphi_{142}$\lbl{A2}}
\BIC{$\varphi_{142}$}
\AXC{$\varphi_{142} \to d \in \mathbb{R} \to \delta \in \mathbb{R} \to \varphi_{32}$\lbl{A2}}
\BIC{$d \in \mathbb{R} \to \delta \in \mathbb{R} \to \varphi_{32}$}
\end{ptree}
\noindent$\Pi_{159}$:
\begin{ptree}
\AXC{$\mathrm{FLim}(f,c,\alpha) \to \varphi_{434}$\lbl{A1}}
\AXC{$\varphi_{525}$\lbl{A1}}
\BIC{$\mathrm{FLim}(f,c,\alpha) \to \mathrm{FLim}(g,c,\beta) \to \beta \neq 0 \to \varphi_{379}$}
\AXC{$\stackrel{\Pi_{158}}{\vdots}$}\noLine
\UIC{$d \in \mathbb{R} \to \delta \in \mathbb{R} \to \varphi_{32}$}
\AXC{$(d \in \mathbb{R} \to \delta \in \mathbb{R} \to \varphi_{32}) \to \varphi_{574}$\lbl{A1}}
\BIC{$\varphi_{574}$}
\BIC{$(\mathrm{FLim}(f,c,\alpha) \to \varphi_{268}) \to \varphi_{542}$}
\end{ptree}
\noindent$\Pi_{160}$:
\begin{ptree}
\AXC{$\mathrm{FLim}(f,c,\alpha) \to \mathrm{FLim}(g,c,\beta) \to \beta \neq 0 \to \varphi_{382}$\lbl{A1}}
\AXC{$\varphi_{471}$\lbl{A1}}
\BIC{$\mathrm{FLim}(f,c,\alpha) \to \varphi_{268}$}
\AXC{$\stackrel{\Pi_{159}}{\vdots}$}\noLine
\UIC{$(\mathrm{FLim}(f,c,\alpha) \to \varphi_{268}) \to \varphi_{542}$}
\BIC{$\varphi_{542}$}
\end{ptree}
\noindent$\Pi_{161}$:
\begin{ptree}
\AXC{$\stackrel{\Pi_{141}}{\vdots}$}\noLine
\UIC{$\mathrm{FLim}(f,c,\alpha) \to \varphi_{262}$}
\AXC{$\mathrm{FLim}(f,c,\alpha) \to \varphi_{434}$\lbl{A1}}
\AXC{$\varphi_{524}$\lbl{A1}}
\BIC{$\mathrm{FLim}(f,c,\alpha) \to \mathrm{FLim}(g,c,\beta) \to \beta \neq 0 \to \varphi_{374}$}
\AXC{$\stackrel{\Pi_{160}}{\vdots}$}\noLine
\UIC{$\varphi_{542}$}
\BIC{$(\mathrm{FLim}(f,c,\alpha) \to \varphi_{262}) \to \mathrm{FLim}(f,c,\alpha) \to \varphi_{399}$}
\BIC{$\mathrm{FLim}(f,c,\alpha) \to \varphi_{399}$}
\end{ptree}
\noindent$\Pi_{162}$:
\begin{ptree}
\AXC{$\stackrel{\Pi_{140}}{\vdots}$}\noLine
\UIC{$\forall r\,(\varphi_{395} \to \neg (r \in \mathbb{R} \land 0 < r \land r \leq d \land r \leq \delta))$}
\AXC{$\stackrel{\Pi_{161}}{\vdots}$}\noLine
\UIC{$\mathrm{FLim}(f,c,\alpha) \to \varphi_{399}$}
\AXC{$(\mathrm{FLim}(f,c,\alpha) \to \varphi_{399}) \to \varphi_{582}$\lbl{A1}}
\BIC{$\varphi_{582}$}
\BIC{$\varphi_{527} \to (\varphi_{395} \to \forall r\,(\varphi_{395})) \to \varphi_{519}$}
\end{ptree}
\noindent$\Pi_{163}$:
\begin{ptree}
\AXC{$\varphi_{92}$\lbl{A4}}
\AXC{$\forall u_{4}\,(\varphi_{80}) \to \varphi_{107}$\lbl{A3}}
\AXC{$\varphi_{80}$\lbl{A2}}
\UIC{$\forall u_{4}\,(\varphi_{80})$}
\AXC{$\varphi_{395} \to \forall r\,(\varphi_{395})$\lbl{A4}}
\AXC{$\varphi_{527}$\lbl{A3}}
\AXC{$\stackrel{\Pi_{162}}{\vdots}$}\noLine
\UIC{$\varphi_{527} \to (\varphi_{395} \to \forall r\,(\varphi_{395})) \to \varphi_{519}$}
\BIC{$(\varphi_{395} \to \forall r\,(\varphi_{395})) \to \varphi_{519}$}
\BIC{$\varphi_{519}$}
\BIC{$(\forall u_{4}\,(\varphi_{80}) \to \varphi_{107}) \to \varphi_{92} \to \mathrm{FLim}(f,c,\alpha) \to \varphi_{389}$}
\BIC{$\varphi_{92} \to \mathrm{FLim}(f,c,\alpha) \to \varphi_{389}$}
\BIC{$\mathrm{FLim}(f,c,\alpha) \to \varphi_{389}$}
\end{ptree}
\noindent$\Pi_{164}$:
\begin{ptree}
\AXC{$\forall f\,\forall c\,\forall \alpha\,(\mathrm{FLim}(f,c,\alpha) \to \mathrm{PNbhd}(f,c))$\lbl{[D5]}}
\AXC{$\varphi_{59}$\lbl{A2}}
\BIC{$\forall c\,\forall \alpha\,(\mathrm{FLim}(g,c,\alpha) \to \mathrm{PNbhd}(g,c))$}
\AXC{$\varphi_{38}$\lbl{A2}}
\BIC{$\forall \alpha\,(\mathrm{FLim}(g,c,\alpha) \to \mathrm{PNbhd}(g,c))$}
\AXC{$\varphi_{25}$\lbl{A2}}
\BIC{$\mathrm{FLim}(g,c,\beta) \to \mathrm{PNbhd}(g,c)$}
\end{ptree}
\noindent$\Pi_{165}$:
\begin{ptree}
\AXC{$\forall f\,\forall c\,\forall \alpha\,(\mathrm{FLim}(f,c,\alpha) \to \mathrm{PNbhd}(f,c))$\lbl{[D5]}}
\AXC{$\varphi_{60}$\lbl{A2}}
\BIC{$\forall c\,\forall \alpha\,(\mathrm{FLim}(f,c,\alpha) \to \mathrm{PNbhd}(f,c))$}
\AXC{$\varphi_{39}$\lbl{A2}}
\BIC{$\forall \alpha\,(\mathrm{FLim}(f,c,\alpha) \to \mathrm{PNbhd}(f,c))$}
\AXC{$\varphi_{26}$\lbl{A2}}
\BIC{$\mathrm{FLim}(f,c,\alpha) \to \mathrm{PNbhd}(f,c)$}
\end{ptree}
\noindent$\Pi_{166}$:
\begin{ptree}
\AXC{$\stackrel{\Pi_{165}}{\vdots}$}\noLine
\UIC{$\mathrm{FLim}(f,c,\alpha) \to \mathrm{PNbhd}(f,c)$}
\AXC{$(\mathrm{FLim}(f,c,\alpha) \to \mathrm{PNbhd}(f,c)) \to \varphi_{326}$\lbl{A1}}
\BIC{$\varphi_{326}$}
\end{ptree}
\noindent$\Pi_{167}$:
\begin{ptree}
\AXC{$\forall f\,\forall g\,\forall c\,(\mathrm{PNbhd}(f,c) \to \varphi_{196})$\lbl{[H2]}}
\AXC{$\varphi_{392}$\lbl{A2}}
\BIC{$\forall g\,\forall c\,(\mathrm{PNbhd}(f,c) \to \varphi_{196})$}
\AXC{$\varphi_{390}$\lbl{A2}}
\BIC{$\forall c\,(\mathrm{PNbhd}(f,c) \to \varphi_{196})$}
\AXC{$\forall c\,(\mathrm{PNbhd}(f,c) \to \varphi_{196}) \to \mathrm{PNbhd}(f,c) \to \varphi_{196}$\lbl{A2}}
\BIC{$\mathrm{PNbhd}(f,c) \to \varphi_{196}$}
\end{ptree}
\noindent$\Pi_{168}$:
\begin{ptree}
\AXC{$\mathrm{FLim}(f,c,\alpha) \to \mathrm{FLim}(f,c,\alpha)$\lbl{A1}}
\AXC{$\stackrel{\Pi_{166}}{\vdots}$}\noLine
\UIC{$\varphi_{326}$}
\BIC{$\mathrm{FLim}(f,c,\alpha) \to \mathrm{FLim}(g,c,\beta) \to \beta \neq 0 \to \varphi_{224}$}
\AXC{$\stackrel{\Pi_{167}}{\vdots}$}\noLine
\UIC{$\mathrm{PNbhd}(f,c) \to \varphi_{196}$}
\AXC{$(\mathrm{PNbhd}(f,c) \to \varphi_{196}) \to \varphi_{518}$\lbl{A1}}
\BIC{$\varphi_{518}$}
\BIC{$\varphi_{468}$}
\end{ptree}
\noindent$\Pi_{169}$:
\begin{ptree}
\AXC{$\varphi_{2}$\lbl{A1}}
\AXC{$\stackrel{\Pi_{164}}{\vdots}$}\noLine
\UIC{$\mathrm{FLim}(g,c,\beta) \to \mathrm{PNbhd}(g,c)$}
\AXC{$(\mathrm{FLim}(g,c,\beta) \to \mathrm{PNbhd}(g,c)) \to \varphi_{349}$\lbl{A1}}
\BIC{$\varphi_{349}$}
\BIC{$\mathrm{FLim}(f,c,\alpha) \to \mathrm{FLim}(g,c,\beta) \to \beta \neq 0 \to \varphi_{223}$}
\AXC{$\stackrel{\Pi_{168}}{\vdots}$}\noLine
\UIC{$\varphi_{468}$}
\BIC{$\mathrm{FLim}(f,c,\alpha) \to \mathrm{FLim}(g,c,\beta) \to \beta \neq 0 \to \varphi_{371}$}
\end{ptree}
\noindent$\Pi_{170}$:
\begin{ptree}
\AXC{$\varphi_{513}$\lbl{A3}}
\AXC{$\stackrel{\Pi_{163}}{\vdots}$}\noLine
\UIC{$\mathrm{FLim}(f,c,\alpha) \to \varphi_{389}$}
\AXC{$\varphi_{510}$\lbl{A1}}
\BIC{$\varphi_{307} \to \neg (d \in \mathbb{R} \land 0 < d \land \forall x \in \mathbb{R}\,(\varphi_{46}))$}
\UIC{$\forall d\,(\varphi_{307} \to \neg (d \in \mathbb{R} \land 0 < d \land \forall x \in \mathbb{R}\,(\varphi_{46})))$}
\AXC{$\stackrel{\Pi_{169}}{\vdots}$}\noLine
\UIC{$\mathrm{FLim}(f,c,\alpha) \to \mathrm{FLim}(g,c,\beta) \to \beta \neq 0 \to \varphi_{371}$}
\AXC{$\varphi_{577}$\lbl{A1}}
\BIC{$\forall d\,(\varphi_{307} \to \neg (d \in \mathbb{R} \land 0 < d \land \forall x \in \mathbb{R}\,(\varphi_{46}))) \to \varphi_{513} \to \varphi_{502}$}
\BIC{$\varphi_{513} \to \varphi_{502}$}
\BIC{$\varphi_{502}$}
\end{ptree}
\noindent$\Pi_{171}$:
\begin{ptree}
\AXC{$2^{-1} \in \mathbb{R}$\lbl{[F31]}}
\AXC{$\stackrel{\Pi_{128}}{\vdots}$}\noLine
\UIC{$\varphi_{18}$}
\AXC{$\forall x\,\forall y\,(x \in \mathbb{R} \to y \in \mathbb{R} \to x \cdot y \in \mathbb{R})$\lbl{(R1)}}
\AXC{$\varphi_{170}$\lbl{A2}}
\BIC{$\forall y\,(\lvert \beta \rvert \in \mathbb{R} \to y \in \mathbb{R} \to \lvert \beta \rvert \cdot y \in \mathbb{R})$}
\AXC{$\varphi_{201}$\lbl{A2}}
\BIC{$\lvert \beta \rvert \in \mathbb{R} \to 2^{-1} \in \mathbb{R} \to \lvert \beta \rvert \cdot 2^{-1} \in \mathbb{R}$}
\AXC{$\varphi_{331}$\lbl{A1}}
\BIC{$\varphi_{18} \to 2^{-1} \in \mathbb{R} \to \varphi_{40}$}
\BIC{$2^{-1} \in \mathbb{R} \to \varphi_{40}$}
\BIC{$\varphi_{40}$}
\end{ptree}
\noindent$\Pi_{172}$:
\begin{ptree}
\AXC{$\varphi_{287}$\lbl{[D1]}}
\AXC{$\varphi_{287} \to \forall c\,(\varphi_{271})$\lbl{A2}}
\BIC{$\forall c\,(\varphi_{271})$}
\AXC{$\forall c\,(\varphi_{271}) \to \varphi_{271}$\lbl{A2}}
\BIC{$\varphi_{271}$}
\AXC{$\varphi_{271} \to \varphi_{264}$\lbl{A2}}
\BIC{$\varphi_{264}$}
\AXC{$\varphi_{264} \to \mathrm{FLim}(g,c,\beta) \to \lvert \beta \rvert \cdot 2^{-1} \in \mathbb{R} \to \varphi_{254}$\lbl{A2}}
\BIC{$\mathrm{FLim}(g,c,\beta) \to \lvert \beta \rvert \cdot 2^{-1} \in \mathbb{R} \to \varphi_{254}$}
\end{ptree}
\noindent$\Pi_{173}$:
\begin{ptree}
\AXC{$\stackrel{\Pi_{130}}{\vdots}$}\noLine
\UIC{$\varphi_{33}$}
\AXC{$\stackrel{\Pi_{171}}{\vdots}$}\noLine
\UIC{$\varphi_{40}$}
\AXC{$\varphi_{2}$\lbl{A1}}
\AXC{$\stackrel{\Pi_{172}}{\vdots}$}\noLine
\UIC{$\mathrm{FLim}(g,c,\beta) \to \lvert \beta \rvert \cdot 2^{-1} \in \mathbb{R} \to \varphi_{254}$}
\AXC{$\varphi_{496}$\lbl{A1}}
\BIC{$\varphi_{2} \to \varphi_{40} \to \varphi_{33} \to \mathrm{FLim}(f,c,\alpha) \to \varphi_{247}$}
\BIC{$\varphi_{40} \to \varphi_{33} \to \mathrm{FLim}(f,c,\alpha) \to \varphi_{247}$}
\BIC{$\varphi_{33} \to \mathrm{FLim}(f,c,\alpha) \to \varphi_{247}$}
\BIC{$\mathrm{FLim}(f,c,\alpha) \to \varphi_{247}$}
\end{ptree}
\noindent$\Pi_{174}$:
\begin{ptree}
\AXC{$\stackrel{\Pi_{173}}{\vdots}$}\noLine
\UIC{$\mathrm{FLim}(f,c,\alpha) \to \varphi_{247}$}
\AXC{$(\mathrm{FLim}(f,c,\alpha) \to \varphi_{247}) \to \varphi_{539}$\lbl{A1}}
\BIC{$\varphi_{539}$}
\end{ptree}
\noindent$\Pi_{175}$:
\begin{ptree}
\AXC{$\varphi_{28} \to \forall \delta\,(\varphi_{28})$\lbl{A4}}
\AXC{$\varphi_{452}$\lbl{A3}}
\AXC{$\varphi_{307} \to \forall d\,(\varphi_{307})$\lbl{A4}}
\AXC{$\stackrel{\Pi_{170}}{\vdots}$}\noLine
\UIC{$\varphi_{502}$}
\BIC{$\mathrm{FLim}(f,c,\alpha) \to \mathrm{FLim}(g,c,\beta) \to \beta \neq 0 \to \varphi_{237}$}
\AXC{$\varphi_{441}$\lbl{A1}}
\BIC{$\varphi_{28} \to \neg (\delta \in \mathbb{R} \land 0 < \delta \land \forall x \in \mathrm{dom}(g)\,(\varphi_{94}))$}
\UIC{$\forall \delta\,(\varphi_{28} \to \neg (\delta \in \mathbb{R} \land 0 < \delta \land \forall x \in \mathrm{dom}(g)\,(\varphi_{94})))$}
\AXC{$\stackrel{\Pi_{174}}{\vdots}$}\noLine
\UIC{$\varphi_{539}$}
\BIC{$\varphi_{452} \to (\varphi_{28} \to \forall \delta\,(\varphi_{28})) \to \varphi_{24}$}
\BIC{$(\varphi_{28} \to \forall \delta\,(\varphi_{28})) \to \varphi_{24}$}
\BIC{$\varphi_{24}$}
\end{ptree}
\noindent$\Pi_{176}$:
\begin{ptree}
\AXC{$\varphi_{3}$\lbl{A1}}
\AXC{$\stackrel{\Pi_{121}}{\vdots}$}\noLine
\UIC{$\varphi_{9}$}
\AXC{$\forall x \in \mathbb{R}\,(x \neq 0 \to x^{-1} \in \mathbb{R})$\lbl{(R1)}}
\AXC{$\varphi_{93}$\lbl{A2}}
\BIC{$\beta \in \mathbb{R} \to \beta \neq 0 \to \beta^{-1} \in \mathbb{R}$}
\AXC{$(\beta \in \mathbb{R} \to \beta \neq 0 \to \beta^{-1} \in \mathbb{R}) \to \varphi_{9} \to \varphi_{3} \to \varphi_{19}$\lbl{A1}}
\BIC{$\varphi_{9} \to \varphi_{3} \to \varphi_{19}$}
\BIC{$\varphi_{3} \to \varphi_{19}$}
\BIC{$\varphi_{19}$}
\end{ptree}
\noindent$\Pi_{177}$:
\begin{ptree}
\AXC{$\forall f\,\forall c\,\forall \alpha\,(\mathrm{FLim}(f,c,\alpha) \to \alpha \in \mathbb{R})$\lbl{[D4]}}
\AXC{$\varphi_{70}$\lbl{A2}}
\BIC{$\forall c\,\forall \alpha\,(\mathrm{FLim}(f,c,\alpha) \to \alpha \in \mathbb{R})$}
\AXC{$\varphi_{52}$\lbl{A2}}
\BIC{$\forall \alpha\,(\mathrm{FLim}(f,c,\alpha) \to \alpha \in \mathbb{R})$}
\AXC{$\varphi_{31}$\lbl{A2}}
\BIC{$\mathrm{FLim}(f,c,\alpha) \to \alpha \in \mathbb{R}$}
\AXC{$\varphi_{180}$\lbl{A1}}
\BIC{$(\mathrm{FLim}(f,c,\alpha) \to \mathrm{FLim}(f,c,\alpha)) \to \varphi_{10}$}
\end{ptree}
\noindent$\Pi_{178}$:
\begin{ptree}
\AXC{$\forall x\,\forall y\,(x \in \mathbb{R} \to y \in \mathbb{R} \to x \cdot y \in \mathbb{R})$\lbl{(R1)}}
\AXC{$\varphi_{135}$\lbl{A2}}
\BIC{$\forall y\,(\alpha \in \mathbb{R} \to y \in \mathbb{R} \to \alpha \cdot y \in \mathbb{R})$}
\AXC{$\varphi_{165}$\lbl{A2}}
\BIC{$\alpha \in \mathbb{R} \to \beta^{-1} \in \mathbb{R} \to \alpha \cdot \beta^{-1} \in \mathbb{R}$}
\AXC{$(\alpha \in \mathbb{R} \to \beta^{-1} \in \mathbb{R} \to \alpha \cdot \beta^{-1} \in \mathbb{R}) \to \varphi_{10} \to \varphi_{19} \to \varphi_{34}$\lbl{A1}}
\BIC{$\varphi_{10} \to \varphi_{19} \to \varphi_{34}$}
\end{ptree}
\noindent$\Pi_{179}$:
\begin{ptree}
\AXC{$\mathrm{FLim}(f,c,\alpha) \to \mathrm{FLim}(f,c,\alpha)$\lbl{A1}}
\AXC{$\stackrel{\Pi_{177}}{\vdots}$}\noLine
\UIC{$(\mathrm{FLim}(f,c,\alpha) \to \mathrm{FLim}(f,c,\alpha)) \to \varphi_{10}$}
\BIC{$\varphi_{10}$}
\AXC{$\stackrel{\Pi_{178}}{\vdots}$}\noLine
\UIC{$\varphi_{10} \to \varphi_{19} \to \varphi_{34}$}
\BIC{$\varphi_{19} \to \varphi_{34}$}
\end{ptree}
\noindent$\Pi_{180}$:
\begin{ptree}
\AXC{$\forall f\,\forall c\,\forall \alpha\,(\mathrm{FLim}(f,c,\alpha) \to c \in \mathbb{R})$\lbl{[D3]}}
\AXC{$\varphi_{61}$\lbl{A2}}
\BIC{$\forall c\,\forall \alpha\,(\mathrm{FLim}(f,c,\alpha) \to c \in \mathbb{R})$}
\AXC{$\varphi_{41}$\lbl{A2}}
\BIC{$\forall \alpha\,(\mathrm{FLim}(f,c,\alpha) \to c \in \mathbb{R})$}
\AXC{$\varphi_{27}$\lbl{A2}}
\BIC{$\mathrm{FLim}(f,c,\alpha) \to c \in \mathbb{R}$}
\AXC{$\varphi_{171}$\lbl{A1}}
\BIC{$(\mathrm{FLim}(f,c,\alpha) \to \mathrm{FLim}(f,c,\alpha)) \to \varphi_{6}$}
\end{ptree}
\noindent$\Pi_{181}$:
\begin{ptree}
\AXC{$\stackrel{\Pi_{80}}{\vdots}$}\noLine
\UIC{$\varphi_{12}$}
\AXC{$\varphi_{290}$\lbl{定理1.7(⇐)}}
\AXC{$\varphi_{290} \to \forall c\,(\varphi_{329})$\lbl{A2}}
\BIC{$\forall c\,(\varphi_{329})$}
\AXC{$\forall c\,(\varphi_{329}) \to \varphi_{329}$\lbl{A2}}
\BIC{$\varphi_{329}$}
\AXC{$\varphi_{329} \to \varphi_{373}$\lbl{A2}}
\BIC{$\varphi_{373}$}
\AXC{$\varphi_{373} \to \varphi_{12} \to \varphi_{6} \to \varphi_{34} \to \varphi_{398}$\lbl{A1}}
\BIC{$\varphi_{12} \to \varphi_{6} \to \varphi_{34} \to \varphi_{398}$}
\BIC{$\varphi_{6} \to \varphi_{34} \to \varphi_{398}$}
\end{ptree}
\noindent$\Pi_{182}$:
\begin{ptree}
\AXC{$\mathrm{FLim}(f,c,\alpha) \to \mathrm{FLim}(f,c,\alpha)$\lbl{A1}}
\AXC{$\stackrel{\Pi_{180}}{\vdots}$}\noLine
\UIC{$(\mathrm{FLim}(f,c,\alpha) \to \mathrm{FLim}(f,c,\alpha)) \to \varphi_{6}$}
\BIC{$\varphi_{6}$}
\AXC{$\stackrel{\Pi_{181}}{\vdots}$}\noLine
\UIC{$\varphi_{6} \to \varphi_{34} \to \varphi_{398}$}
\BIC{$\varphi_{34} \to \varphi_{398}$}
\end{ptree}
\noindent 主木:
\begin{ptree}
\AXC{$\stackrel{\Pi_{91}}{\vdots}$}\noLine
\UIC{$\mathrm{FLim}(f,c,\alpha) \to \varphi_{242}$}
\AXC{$\stackrel{\Pi_{175}}{\vdots}$}\noLine
\UIC{$\varphi_{24}$}
\AXC{$\stackrel{\Pi_{176}}{\vdots}$}\noLine
\UIC{$\varphi_{19}$}
\AXC{$\stackrel{\Pi_{179}}{\vdots}$}\noLine
\UIC{$\varphi_{19} \to \varphi_{34}$}
\BIC{$\varphi_{34}$}
\AXC{$\stackrel{\Pi_{182}}{\vdots}$}\noLine
\UIC{$\varphi_{34} \to \varphi_{398}$}
\BIC{$\varphi_{398}$}
\BIC{$(\mathrm{FLim}(f,c,\alpha) \to \varphi_{242}) \to \mathrm{FLim}(f,c,\alpha) \to \varphi_{20}$}
\BIC{$\mathrm{FLim}(f,c,\alpha) \to \varphi_{20}$}
\UIC{$\forall \beta\,(\mathrm{FLim}(f,c,\alpha) \to \varphi_{20})$}
\UIC{$\forall \alpha\,\forall \beta\,(\mathrm{FLim}(f,c,\alpha) \to \varphi_{20})$}
\UIC{$\forall c\,\forall \alpha\,\forall \beta\,(\mathrm{FLim}(f,c,\alpha) \to \varphi_{20})$}
\UIC{$\forall g\,\forall c\,\forall \alpha\,\forall \beta\,(\mathrm{FLim}(f,c,\alpha) \to \varphi_{20})$}
\UIC{$\forall f\,\forall g\,\forall c\,\forall \alpha\,\forall \beta\,(\mathrm{FLim}(f,c,\alpha) \to \varphi_{20})$}
\end{ptree}
\begin{abbrevlist}
\item $\varphi_{1} :\equiv \forall a\,(\mathrm{FLim}(f,c,\alpha) \land \mathrm{FLim}(g,c,\beta) \land \beta \neq 0)$
\item $\varphi_{2} :\equiv \mathrm{FLim}(f,c,\alpha) \to \mathrm{FLim}(g,c,\beta) \to \mathrm{FLim}(g,c,\beta)$
\item $\varphi_{3} :\equiv \mathrm{FLim}(f,c,\alpha) \to \mathrm{FLim}(g,c,\beta) \to \beta \neq 0 \to \beta \neq 0$
\item $\varphi_{4} :\equiv \mathrm{FLim}(f,c,\alpha) \to \mathrm{FLim}(g,c,\beta) \to \beta \neq 0 \to \mathrm{RFn}(g)$
\item $\varphi_{5} :\equiv \mathrm{FLim}(f,c,\alpha) \to \mathrm{FLim}(g,c,\beta) \to \beta \neq 0 \to \mathrm{RFn}(f)$
\item $\varphi_{6} :\equiv \mathrm{FLim}(f,c,\alpha) \to \mathrm{FLim}(g,c,\beta) \to \beta \neq 0 \to c \in \mathbb{R}$
\item $\varphi_{7} :\equiv \forall \alpha\,(\mathrm{FLim}(g,c,\alpha) \to \mathrm{RFn}(g)) \to \mathrm{FLim}(g,c,\beta) \to \mathrm{RFn}(g)$
\item $\varphi_{8} :\equiv \forall \alpha\,(\mathrm{FLim}(f,c,\alpha) \to \mathrm{RFn}(f)) \to \mathrm{FLim}(f,c,\alpha) \to \mathrm{RFn}(f)$
\item $\varphi_{9} :\equiv \mathrm{FLim}(f,c,\alpha) \to \mathrm{FLim}(g,c,\beta) \to \beta \neq 0 \to \beta \in \mathbb{R}$
\item $\varphi_{10} :\equiv \mathrm{FLim}(f,c,\alpha) \to \mathrm{FLim}(g,c,\beta) \to \beta \neq 0 \to \alpha \in \mathbb{R}$
\item $\varphi_{11} :\equiv \mathrm{FLim}(f,c,\alpha) \to \mathrm{FLim}(g,c,\beta) \to \beta \neq 0 \to 0 < \lvert \beta \rvert$
\item $\varphi_{12} :\equiv \mathrm{FLim}(f,c,\alpha) \to \mathrm{FLim}(g,c,\beta) \to \beta \neq 0 \to \mathrm{RFn}((f / g))$
\item $\varphi_{13} :\equiv \mathrm{FLim}(f,c,\alpha) \to \mathrm{FLim}(g,c,\beta) \to \beta \neq 0 \to \mathrm{Seq}(a) \to \mathrm{Seq}(a)$
\item $\varphi_{14} :\equiv \mathrm{FLim}(g,c,\beta) \to \beta \neq 0 \to \mathrm{Seq}(a) \to \mathrm{Conv}(a,c) \to \mathrm{Conv}(a,c)$
\item $\varphi_{15} :\equiv r \in \mathbb{R} \land 0 < r \land r \leq d \land r \leq \delta \to \mathrm{PNbhd}((f / g),c)$
\item $\varphi_{16} :\equiv \forall x\,\forall y\,(x = y \to y = x) \to \forall y\,(\lvert g(x) \rvert = y \to y = \lvert g(x) \rvert)$
\item $\varphi_{17} :\equiv \forall z\,\forall w\,(z = w \to 0 < z \to 0 < w) \to \forall w\,(0 = w \to 0 < 0 \to 0 < w)$
\item $\varphi_{18} :\equiv \mathrm{FLim}(f,c,\alpha) \to \mathrm{FLim}(g,c,\beta) \to \beta \neq 0 \to \lvert \beta \rvert \in \mathbb{R}$
\item $\varphi_{19} :\equiv \mathrm{FLim}(f,c,\alpha) \to \mathrm{FLim}(g,c,\beta) \to \beta \neq 0 \to \beta^{-1} \in \mathbb{R}$
\item $\varphi_{20} :\equiv \mathrm{FLim}(g,c,\beta) \to \beta \neq 0 \to \mathrm{FLim}((f / g),c,\alpha \cdot \beta^{-1})$
\item $\varphi_{21} :\equiv \forall z\,\forall w\,(z = w \to \lvert g(x) \rvert = \lvert z \rvert \to \lvert g(x) \rvert = \lvert w \rvert)$
\item $\varphi_{22} :\equiv \forall c\,\forall \alpha\,(\mathrm{FLim}(g,c,\alpha) \to \mathrm{RFn}(g)) \to \forall \alpha\,(\mathrm{FLim}(g,c,\alpha) \to \mathrm{RFn}(g))$
\item $\varphi_{23} :\equiv \forall c\,\forall \alpha\,(\mathrm{FLim}(f,c,\alpha) \to \mathrm{RFn}(f)) \to \forall \alpha\,(\mathrm{FLim}(f,c,\alpha) \to \mathrm{RFn}(f))$
\item $\varphi_{24} :\equiv \mathrm{FLim}(f,c,\alpha) \to \mathrm{FLim}(g,c,\beta) \to \beta \neq 0 \to \mathrm{PNbhd}((f / g),c)$
\item $\varphi_{25} :\equiv \forall \alpha\,(\mathrm{FLim}(g,c,\alpha) \to \mathrm{PNbhd}(g,c)) \to \mathrm{FLim}(g,c,\beta) \to \mathrm{PNbhd}(g,c)$
\item $\varphi_{26} :\equiv \forall \alpha\,(\mathrm{FLim}(f,c,\alpha) \to \mathrm{PNbhd}(f,c)) \to \mathrm{FLim}(f,c,\alpha) \to \mathrm{PNbhd}(f,c)$
\item $\varphi_{27} :\equiv \forall \alpha\,(\mathrm{FLim}(f,c,\alpha) \to c \in \mathbb{R}) \to \mathrm{FLim}(f,c,\alpha) \to c \in \mathbb{R}$
\item $\varphi_{28} :\equiv \mathrm{FLim}(f,c,\alpha) \land \mathrm{FLim}(g,c,\beta) \land \beta \neq 0 \land \neg \mathrm{PNbhd}((f / g),c)$
\item $\varphi_{29} :\equiv \forall x\,\forall y\,(x = y \to y = x) \to \forall y\,(g(a_{n}) = y \to y = g(a_{n}))$
\item $\varphi_{30} :\equiv \forall \alpha\,(\mathrm{FLim}(g,c,\alpha) \to \alpha \in \mathbb{R}) \to \mathrm{FLim}(g,c,\beta) \to \beta \in \mathbb{R}$
\item $\varphi_{31} :\equiv \forall \alpha\,(\mathrm{FLim}(f,c,\alpha) \to \alpha \in \mathbb{R}) \to \mathrm{FLim}(f,c,\alpha) \to \alpha \in \mathbb{R}$
\item $\varphi_{32} :\equiv 0 < d \to 0 < \delta \to \exists u_{4}\,(u_{4} \in \mathbb{R} \land 0 < u_{4} \land u_{4} \leq d \land u_{4} \leq \delta)$
\item $\varphi_{33} :\equiv \mathrm{FLim}(f,c,\alpha) \to \mathrm{FLim}(g,c,\beta) \to \beta \neq 0 \to 0 < \lvert \beta \rvert \cdot 2^{-1}$
\item $\varphi_{34} :\equiv \mathrm{FLim}(f,c,\alpha) \to \mathrm{FLim}(g,c,\beta) \to \beta \neq 0 \to \alpha \cdot \beta^{-1} \in \mathbb{R}$
\item $\varphi_{35} :\equiv \forall f\,\forall c\,\forall \alpha\,(\mathrm{FLim}(f,c,\alpha) \to \mathrm{RFn}(f)) \to \forall c\,\forall \alpha\,(\mathrm{FLim}(g,c,\alpha) \to \mathrm{RFn}(g))$
\item $\varphi_{36} :\equiv \forall f\,\forall c\,\forall \alpha\,(\mathrm{FLim}(f,c,\alpha) \to \mathrm{RFn}(f)) \to \forall c\,\forall \alpha\,(\mathrm{FLim}(f,c,\alpha) \to \mathrm{RFn}(f))$
\item $\varphi_{37} :\equiv g(x) = 0 \to \lvert g(x) \rvert = \lvert g(x) \rvert \to \lvert g(x) \rvert = \lvert 0 \rvert$
\item $\varphi_{38} :\equiv \forall c\,\forall \alpha\,(\mathrm{FLim}(g,c,\alpha) \to \mathrm{PNbhd}(g,c)) \to \forall \alpha\,(\mathrm{FLim}(g,c,\alpha) \to \mathrm{PNbhd}(g,c))$
\item $\varphi_{39} :\equiv \forall c\,\forall \alpha\,(\mathrm{FLim}(f,c,\alpha) \to \mathrm{PNbhd}(f,c)) \to \forall \alpha\,(\mathrm{FLim}(f,c,\alpha) \to \mathrm{PNbhd}(f,c))$
\item $\varphi_{40} :\equiv \mathrm{FLim}(f,c,\alpha) \to \mathrm{FLim}(g,c,\beta) \to \beta \neq 0 \to \lvert \beta \rvert \cdot 2^{-1} \in \mathbb{R}$
\item $\varphi_{41} :\equiv \forall c\,\forall \alpha\,(\mathrm{FLim}(f,c,\alpha) \to c \in \mathbb{R}) \to \forall \alpha\,(\mathrm{FLim}(f,c,\alpha) \to c \in \mathbb{R})$
\item $\varphi_{42} :\equiv x \in \mathrm{dom}((f / g)) \leftrightarrow x \in \mathrm{dom}(f) \land x \in \mathrm{dom}(g) \land g(x) \neq 0$
\item $\varphi_{43} :\equiv 0 < v_{766} \to 0 < v_{767} \to \exists d\,(d \in \mathbb{R} \land 0 < d \land d \leq v_{766} \land d \leq v_{767})$
\item $\varphi_{44} :\equiv \forall y\,\forall z\,((f / g)(a_{n}) = y \to y = z \to (f / g)(a_{n}) = z)$
\item $\varphi_{45} :\equiv 0 < \lvert x - c \rvert \to \lvert x - c \rvert < r \to \neg g(x) \neq 0 \to \neg g(x) \neq 0$
\item $\varphi_{46} :\equiv 0 < \lvert x - c \rvert \to \lvert x - c \rvert < d \to x \in \mathrm{dom}(f) \land x \in \mathrm{dom}(g)$
\item $\varphi_{47} :\equiv 0 < \lvert x - c \rvert \to \lvert x - c \rvert < r \to x \in \mathrm{dom}(f) \land x \in \mathrm{dom}(g)$
\item $\varphi_{48} :\equiv \forall x\,\forall y\,(x = y \to y = x) \to \forall y\,(g \circ a(n) = y \to y = g \circ a(n))$
\item $\varphi_{49} :\equiv \forall w\,(g(x) = w \to \lvert g(x) \rvert = \lvert g(x) \rvert \to \lvert g(x) \rvert = \lvert w \rvert)$
\item $\varphi_{50} :\equiv \forall c\,\forall \alpha\,(\mathrm{FLim}(g,c,\alpha) \to \alpha \in \mathbb{R}) \to \forall \alpha\,(\mathrm{FLim}(g,c,\alpha) \to \alpha \in \mathbb{R})$
\item $\varphi_{51} :\equiv \forall x\,\forall y\,\forall z\,(x < y \to y < z \to x < z) \to \forall y\,\forall z\,(0 < y \to y < z \to 0 < z)$
\item $\varphi_{52} :\equiv \forall c\,\forall \alpha\,(\mathrm{FLim}(f,c,\alpha) \to \alpha \in \mathbb{R}) \to \forall \alpha\,(\mathrm{FLim}(f,c,\alpha) \to \alpha \in \mathbb{R})$
\item $\varphi_{53} :\equiv r \in \mathbb{R} \land 0 < r \land r \leq d \land r \leq \delta \to r \in \mathbb{R} \land 0 < r \land r \leq d \land r \leq \delta$
\item $\varphi_{54} :\equiv \forall x \in \mathrm{dom}((f / g))\,((f / g)(x) = f(x) \cdot g(x)^{-1})$
\item $\varphi_{55} :\equiv 0 < \lvert x - c \rvert \to \lvert x - c \rvert < r \to \neg g(x) \neq 0 \to \lvert g(x) \rvert = 0$
\item $\varphi_{56} :\equiv 0 < \lvert x - c \rvert \to \lvert x - c \rvert < r \to \neg g(x) \neq 0 \to 0 = \lvert g(x) \rvert$
\item $\varphi_{57} :\equiv x \in \mathbb{R} \to 0 < \lvert x - c \rvert \to \lvert x - c \rvert < r \to \neg g(x) \neq 0 \to 0 \leq 0$
\item $\varphi_{58} :\equiv \forall z\,\forall w\,(z = w \to z = 0 \to w = 0) \to \forall w\,(g(a_{n}) = w \to g(a_{n}) = 0 \to w = 0)$
\item $\varphi_{59} :\equiv \forall f\,\forall c\,\forall \alpha\,(\mathrm{FLim}(f,c,\alpha) \to \mathrm{PNbhd}(f,c)) \to \forall c\,\forall \alpha\,(\mathrm{FLim}(g,c,\alpha) \to \mathrm{PNbhd}(g,c))$
\item $\varphi_{60} :\equiv \forall f\,\forall c\,\forall \alpha\,(\mathrm{FLim}(f,c,\alpha) \to \mathrm{PNbhd}(f,c)) \to \forall c\,\forall \alpha\,(\mathrm{FLim}(f,c,\alpha) \to \mathrm{PNbhd}(f,c))$
\item $\varphi_{61} :\equiv \forall f\,\forall c\,\forall \alpha\,(\mathrm{FLim}(f,c,\alpha) \to c \in \mathbb{R}) \to \forall c\,\forall \alpha\,(\mathrm{FLim}(f,c,\alpha) \to c \in \mathbb{R})$
\item $\varphi_{62} :\equiv \neg \neg (\neg (\neg (d \in \mathbb{R} \to \neg 0 < d) \to \neg (\neg d < v_{766} \to d = v_{766})) \to \neg (\neg d < v_{767} \to d = v_{767}))$
\item $\varphi_{63} :\equiv \forall x \in \mathbb{R}\,(x \neq 0 \to 0 < \lvert x \rvert) \to \beta \in \mathbb{R} \to \beta \neq 0 \to 0 < \lvert \beta \rvert$
\item $\varphi_{64} :\equiv 0 < \lvert x - c \rvert \to \lvert x - c \rvert < r \to \neg g(x) \neq 0 \to 0 < \lvert g(x) \rvert$
\item $\varphi_{65} :\equiv x \in \mathbb{R} \to 0 < \lvert x - c \rvert \to \lvert x - c \rvert < r \to x \in \mathrm{dom}((f / g))$
\item $\varphi_{66} :\equiv \forall y\,(\lvert g(x) \rvert = y \to y = \lvert g(x) \rvert) \to \lvert g(x) \rvert = 0 \to 0 = \lvert g(x) \rvert$
\item $\varphi_{67} :\equiv \forall x\,\forall y\,\forall z\,(x = y \to y = z \to x = z) \to \forall y\,\forall z\,(\lvert g(x) \rvert = y \to y = z \to \lvert g(x) \rvert = z)$
\item $\varphi_{68} :\equiv \forall w\,(0 = w \to 0 < 0 \to 0 < w) \to 0 = \lvert g(x) \rvert \to 0 < 0 \to 0 < \lvert g(x) \rvert$
\item $\varphi_{69} :\equiv \forall f\,\forall c\,\forall \alpha\,(\mathrm{FLim}(f,c,\alpha) \to \alpha \in \mathbb{R}) \to \forall c\,\forall \alpha\,(\mathrm{FLim}(g,c,\alpha) \to \alpha \in \mathbb{R})$
\item $\varphi_{70} :\equiv \forall f\,\forall c\,\forall \alpha\,(\mathrm{FLim}(f,c,\alpha) \to \alpha \in \mathbb{R}) \to \forall c\,\forall \alpha\,(\mathrm{FLim}(f,c,\alpha) \to \alpha \in \mathbb{R})$
\item $\varphi_{71} :\equiv \forall n \in \mathbb{N}\,(a_{n} \in \mathrm{dom}(g) \land a_{n} \neq c) \to \mathrm{Conv}(g \circ a,\beta)$
\item $\varphi_{72} :\equiv \forall n \in \mathbb{N}\,(a_{n} \in \mathrm{dom}(f) \land a_{n} \neq c) \to \mathrm{Conv}(f \circ a,\alpha)$
\item $\varphi_{73} :\equiv 0 < \lvert x - c \rvert \to \lvert x - c \rvert < r \to \neg g(x) \neq 0 \to \neg 0 < \lvert g(x) \rvert$
\item $\varphi_{74} :\equiv \forall y\,\forall z\,((f / g) \circ a(n) = y \to y = z \to (f / g) \circ a(n) = z)$
\item $\varphi_{75} :\equiv 0 < \lvert x - c \rvert \to \lvert x - c \rvert < r \to \neg g(x) \neq 0 \to \lvert g(x) \rvert = \lvert 0 \rvert$
\item $\varphi_{76} :\equiv \forall a\,\forall \alpha\,(\mathrm{Conv}(a,\alpha) \to \mathrm{Seq}(a)) \to \forall \alpha\,(\mathrm{Conv}(g \circ a,\alpha) \to \mathrm{Seq}(g \circ a))$
\item $\varphi_{77} :\equiv \forall a\,\forall \alpha\,(\mathrm{Conv}(a,\alpha) \to \mathrm{Seq}(a)) \to \forall \alpha\,(\mathrm{Conv}(f \circ a,\alpha) \to \mathrm{Seq}(f \circ a))$
\item $\varphi_{78} :\equiv \forall x \in \mathbb{R}\,(0 < \lvert x - c \rvert \to \lvert x - c \rvert < r \to x \in \mathrm{dom}((f / g)))$
\item $\varphi_{79} :\equiv x \in \mathbb{R} \to 0 < \lvert x - c \rvert \to \lvert x - c \rvert < r \to \neg g(x) \neq 0 \to \lvert 0 \rvert = 0$
\item $\varphi_{80} :\equiv \forall r\,(\neg (r \in \mathbb{R} \land 0 < r \land r \leq d \land r \leq \delta)) \to \neg (u_{4} \in \mathbb{R} \land 0 < u_{4} \land u_{4} \leq d \land u_{4} \leq \delta)$
\item $\varphi_{81} :\equiv a_{n} \in \mathrm{dom}(f) \land a_{n} \in \mathrm{dom}(g) \land g(a_{n}) \neq 0 \land a_{n} \neq c$
\item $\varphi_{82} :\equiv \forall n \in \mathbb{N}\,(a_{n} \in \mathrm{dom}((f / g)) \land a_{n} \neq c) \to \mathrm{Seq}(g \circ a)$
\item $\varphi_{83} :\equiv \forall n \in \mathbb{N}\,(a_{n} \in \mathrm{dom}((f / g)) \land a_{n} \neq c) \to \mathrm{Seq}(f \circ a)$
\item $\varphi_{84} :\equiv x \in \mathbb{R} \to 0 < \lvert x - c \rvert \to \lvert x - c \rvert < r \to \neg g(x) \neq 0 \to g(x) = 0$
\item $\varphi_{85} :\equiv \neg \neg (\neg (\neg (u_{4} \in \mathbb{R} \to \neg 0 < u_{4}) \to \neg (\neg u_{4} < v_{766} \to u_{4} = v_{766})) \to \neg (\neg u_{4} < v_{767} \to u_{4} = v_{767}))$
\item $\varphi_{86} :\equiv \forall a\,\forall b\,\forall \alpha\,(\mathrm{Seq}(b) \to \forall n \in \mathbb{N}\,(a_{n} = b_{n}) \to \mathrm{Conv}(a,\alpha) \to \mathrm{Conv}(b,\alpha))$
\item $\varphi_{87} :\equiv \lvert g(x) - \beta \rvert < \lvert \beta \rvert \cdot 2^{-1} \to \lvert \beta \rvert \cdot 2^{-1} < \lvert g(x) \rvert$
\item $\varphi_{88} :\equiv \forall z\,\forall w\,(z = w \to z = 0 \to w = 0) \to \forall w\,(g \circ a(n) = w \to g \circ a(n) = 0 \to w = 0)$
\item $\varphi_{89} :\equiv \forall z\,\forall w\,(z = w \to 0 < z \to 0 < w) \to \forall w\,(\lvert g(x) \rvert = w \to 0 < \lvert g(x) \rvert \to 0 < w)$
\item $\varphi_{90} :\equiv n \in \mathbb{N} \to \forall n \in \mathbb{N}\,((f / g) \circ a(n) = (f / g)(a_{n}))$
\item $\varphi_{91} :\equiv 0 < \lvert x - c \rvert \to \lvert x - c \rvert < r \to \lvert g(x) - \beta \rvert < \lvert \beta \rvert \cdot 2^{-1}$
\item $\varphi_{92} :\equiv \forall r\,(\neg (r \in \mathbb{R} \land 0 < r \land r \leq d \land r \leq \delta)) \to \forall u_{4}\,\forall r\,(\neg (r \in \mathbb{R} \land 0 < r \land r \leq d \land r \leq \delta))$
\item $\varphi_{93} :\equiv \forall x \in \mathbb{R}\,(x \neq 0 \to x^{-1} \in \mathbb{R}) \to \beta \in \mathbb{R} \to \beta \neq 0 \to \beta^{-1} \in \mathbb{R}$
\item $\varphi_{94} :\equiv 0 < \lvert x - c \rvert \to \lvert x - c \rvert < \delta \to \lvert g(x) - \beta \rvert < \lvert \beta \rvert \cdot 2^{-1}$
\item $\varphi_{95} :\equiv \forall x\,\forall y\,\forall z\,(x < y \to y \leq z \to x < z) \to \forall y\,\forall z\,(\lvert x - c \rvert < y \to y \leq z \to \lvert x - c \rvert < z)$
\item $\varphi_{96} :\equiv 0 < \lvert \beta \rvert \cdot 2^{-1} \to \lvert \beta \rvert \cdot 2^{-1} < \lvert g(x) \rvert \to 0 < \lvert g(x) \rvert$
\item $\varphi_{97} :\equiv \forall n \in \mathbb{N}\,(a_{n} \in \mathrm{dom}((f / g)) \land a_{n} \neq c) \to n \in \mathbb{N} \to n \in \mathbb{N}$
\item $\varphi_{98} :\equiv \forall n \in \mathbb{N}\,(a_{n} \in \mathrm{dom}((f / g)) \land a_{n} \neq c) \to \mathrm{Conv}(g \circ a,\beta)$
\item $\varphi_{99} :\equiv \forall n \in \mathbb{N}\,(a_{n} \in \mathrm{dom}((f / g)) \land a_{n} \neq c) \to \mathrm{Conv}(f \circ a,\alpha)$
\item $\varphi_{100} :\equiv \forall x\,\forall y\,(x = y \to y = x) \to \forall y\,((f / g) \circ a(n) = y \to y = (f / g) \circ a(n))$
\item $\varphi_{101} :\equiv n \in \mathbb{N} \to \mathrm{Seq}(g \circ a) \land \forall n \in \mathbb{N}\,(g \circ a(n) = g(a_{n}))$
\item $\varphi_{102} :\equiv n \in \mathbb{N} \to \mathrm{Seq}(f \circ a) \land \forall n \in \mathbb{N}\,(f \circ a(n) = f(a_{n}))$
\item $\varphi_{103} :\equiv \forall n \in \mathbb{N}\,(a_{n} \in \mathrm{dom}((f / g)) \land a_{n} \neq c) \to \mathrm{Seq}(g \circ a^{\ast})$
\item $\varphi_{104} :\equiv r \in \mathbb{R} \land 0 < r \land r \leq d \land r \leq \delta \to x \in \mathbb{R} \to 0 < \lvert x - c \rvert \to 0 < \lvert x - c \rvert$
\item $\varphi_{105} :\equiv \forall p\,\forall q\,(p \in \mathbb{R} \to q \in \mathbb{R} \to 0 < p \to 0 < q \to \exists d\,(d \in \mathbb{R} \land 0 < d \land d \leq p \land d \leq q))$
\item $\varphi_{106} :\equiv \forall n \in \mathbb{N}\,(a_{n} \in \mathrm{dom}((f / g)) \land a_{n} \neq c) \to \mathrm{Seq}((f / g) \circ a)$
\item $\varphi_{107} :\equiv \forall u_{4}\,\forall r\,(\neg (r \in \mathbb{R} \land 0 < r \land r \leq d \land r \leq \delta)) \to \forall u_{4}\,(\neg (u_{4} \in \mathbb{R} \land 0 < u_{4} \land u_{4} \leq d \land u_{4} \leq \delta))$
\item $\varphi_{108} :\equiv x \in \mathbb{R} \to 0 < \lvert x - c \rvert \to \lvert x - c \rvert < r \to \lvert \beta \rvert \cdot 2^{-1} < \lvert g(x) \rvert$
\item $\varphi_{109} :\equiv \forall z\,(\lvert x - c \rvert < r \to r \leq z \to \lvert x - c \rvert < z) \to \lvert x - c \rvert < r \to r \leq d \to \lvert x - c \rvert < d$
\item $\varphi_{110} :\equiv \forall x\,(\mathrm{RFn}(g) \to x \in \mathrm{dom}(g) \to g(x) \in \mathbb{R}) \to \mathrm{RFn}(g) \to x \in \mathrm{dom}(g) \to g(x) \in \mathbb{R}$
\item $\varphi_{111} :\equiv \forall \alpha\,(\mathrm{Conv}(g \circ a,\alpha) \to \mathrm{Seq}(g \circ a)) \to \mathrm{Conv}(g \circ a,\beta) \to \mathrm{Seq}(g \circ a)$
\item $\varphi_{112} :\equiv \forall \alpha\,(\mathrm{Conv}(f \circ a,\alpha) \to \mathrm{Seq}(f \circ a)) \to \mathrm{Conv}(f \circ a,\alpha) \to \mathrm{Seq}(f \circ a)$
\item $\varphi_{113} :\equiv \forall n \in \mathbb{N}\,(a_{n} \in \mathrm{dom}((f / g)) \land a_{n} \neq c) \to n \in \mathbb{N} \to g(a_{n}) \neq 0$
\item $\varphi_{114} :\equiv r \in \mathbb{R} \land 0 < r \land r \leq d \land r \leq \delta \to x \in \mathbb{R} \to 0 < \lvert x - c \rvert \to \lvert x - c \rvert < r \to r \leq d$
\item $\varphi_{115} :\equiv \forall z\,(\lvert x - c \rvert < r \to r \leq z \to \lvert x - c \rvert < z) \to \lvert x - c \rvert < r \to r \leq \delta \to \lvert x - c \rvert < \delta$
\item $\varphi_{116} :\equiv \forall q\,(v_{766} \in \mathbb{R} \to q \in \mathbb{R} \to 0 < v_{766} \to 0 < q \to \exists d\,(d \in \mathbb{R} \land 0 < d \land d \leq v_{766} \land d \leq q))$
\item $\varphi_{117} :\equiv \forall n \in \mathbb{N}\,(a_{n} \in \mathrm{dom}((f / g)) \land a_{n} \neq c) \to n \in \mathbb{N} \to a_{n} \in \mathrm{dom}(g)$
\item $\varphi_{118} :\equiv \forall n \in \mathbb{N}\,(f \circ a \otimes g \circ a^{\ast}(n) = (f / g) \circ a(n))$
\item $\varphi_{119} :\equiv \forall n \in \mathbb{N}\,(a_{n} \in \mathrm{dom}((f / g)) \land a_{n} \neq c) \to n \in \mathbb{N} \to a_{n} \in \mathrm{dom}(f)$
\item $\varphi_{120} :\equiv r \in \mathbb{R} \land 0 < r \land r \leq d \land r \leq \delta \to x \in \mathbb{R} \to 0 < \lvert x - c \rvert \to \lvert x - c \rvert < r \to r \leq \delta$
\item $\varphi_{121} :\equiv \forall y\,\forall z\,(0 < y \to y < z \to 0 < z) \to \forall z\,(0 < \lvert \beta \rvert \cdot 2^{-1} \to \lvert \beta \rvert \cdot 2^{-1} < z \to 0 < z)$
\item $\varphi_{122} :\equiv a_{n} \in \mathrm{dom}((f / g)) \leftrightarrow a_{n} \in \mathrm{dom}(f) \land a_{n} \in \mathrm{dom}(g) \land g(a_{n}) \neq 0$
\item $\varphi_{123} :\equiv \forall n \in \mathbb{N}\,(a_{n} \in \mathrm{dom}((f / g)) \land a_{n} \neq c) \to \forall n \in \mathbb{N}\,(a_{n} \in \mathrm{dom}(g))$
\item $\varphi_{124} :\equiv \forall n \in \mathbb{N}\,(a_{n} \in \mathrm{dom}((f / g)) \land a_{n} \neq c) \to n \in \mathbb{N} \to g \circ a(n) \neq 0$
\item $\varphi_{125} :\equiv \forall n \in \mathbb{N}\,(a_{n} \in \mathrm{dom}((f / g)) \land a_{n} \neq c) \to \forall n \in \mathbb{N}\,(a_{n} \in \mathrm{dom}(f))$
\item $\varphi_{126} :\equiv \mathrm{Seq}(g \circ a^{\ast}) \land \forall n \in \mathbb{N}\,(g \circ a^{\ast}(n) = g \circ a(n)^{-1})$
\item $\varphi_{127} :\equiv a_{n} \in \mathrm{dom}((f / g)) \to (f / g)(a_{n}) = f(a_{n}) \cdot g(a_{n})^{-1}$
\item $\varphi_{128} :\equiv \forall y\,\forall z\,(\lvert x - c \rvert < y \to y \leq z \to \lvert x - c \rvert < z) \to \forall z\,(\lvert x - c \rvert < r \to r \leq z \to \lvert x - c \rvert < z)$
\item $\varphi_{129} :\equiv \forall f\,\forall x\,(\mathrm{RFn}(f) \to x \in \mathrm{dom}(f) \to f(x) \in \mathbb{R}) \to \forall x\,(\mathrm{RFn}(g) \to x \in \mathrm{dom}(g) \to g(x) \in \mathbb{R})$
\item $\varphi_{130} :\equiv \mathrm{Seq}((f / g) \circ a) \land \forall n \in \mathbb{N}\,((f / g) \circ a(n) = (f / g)(a_{n}))$
\item $\varphi_{131} :\equiv f \circ a(n) \cdot g \circ a(n)^{-1} = f \circ a(n) \cdot g \circ a(n)^{-1}$
\item $\varphi_{132} :\equiv n \in \mathbb{N} \to f(a_{n}) \cdot g \circ a(n)^{-1} = f(a_{n}) \cdot g(a_{n})^{-1}$
\item $\varphi_{133} :\equiv \forall x\,(x = x) \to f(a_{n}) \cdot g \circ a(n)^{-1} = f(a_{n}) \cdot g \circ a(n)^{-1}$
\item $\varphi_{134} :\equiv \forall w\,(\lvert g(x) \rvert = w \to 0 < \lvert g(x) \rvert \to 0 < w) \to \lvert g(x) \rvert = 0 \to 0 < \lvert g(x) \rvert \to 0 < 0$
\item $\varphi_{135} :\equiv \forall x\,\forall y\,(x \in \mathbb{R} \to y \in \mathbb{R} \to x \cdot y \in \mathbb{R}) \to \forall y\,(\alpha \in \mathbb{R} \to y \in \mathbb{R} \to \alpha \cdot y \in \mathbb{R})$
\item $\varphi_{136} :\equiv \forall n \in \mathbb{N}\,(a_{n} \in \mathrm{dom}((f / g)) \land a_{n} \neq c) \to \forall n \in \mathbb{N}\,(g \circ a(n) \neq 0)$
\item $\varphi_{137} :\equiv \forall a\,\forall b\,(\mathrm{Seq}(a) \to \mathrm{Seq}(b) \to \mathrm{Seq}(a \otimes b) \land \forall n \in \mathbb{N}\,(a \otimes b(n) = a_{n} \cdot b_{n}))$
\item $\varphi_{138} :\equiv \forall x\,\forall \beta\,(x \in \mathbb{R} \to \beta \in \mathbb{R} \to \lvert x - \beta \rvert < \lvert \beta \rvert \cdot 2^{-1} \to \lvert \beta \rvert \cdot 2^{-1} < \lvert x \rvert)$
\item $\varphi_{139} :\equiv f \circ a(n) \cdot g \circ a^{\ast}(n) = f \circ a(n) \cdot g \circ a(n)^{-1}$
\item $\varphi_{140} :\equiv r \in \mathbb{R} \land 0 < r \land r \leq d \land r \leq \delta \to x \in \mathbb{R} \to 0 < \lvert x - c \rvert \to \lvert x - c \rvert < r \to g(x) \neq 0$
\item $\varphi_{141} :\equiv (\mathrm{FLim}(f,c,\alpha) \to \mathrm{RFn}(f)) \to (\mathrm{FLim}(f,c,\alpha) \to \mathrm{FLim}(f,c,\alpha)) \to \varphi_{5}$
\item $\varphi_{142} :\equiv \forall v_{767}\,(d \in \mathbb{R} \to v_{767} \in \mathbb{R} \to 0 < d \to 0 < v_{767} \to \exists u_{4}\,(u_{4} \in \mathbb{R} \land 0 < u_{4} \land u_{4} \leq d \land u_{4} \leq v_{767}))$
\item $\varphi_{143} :\equiv \forall n \in \mathbb{N}\,(a_{n} \in \mathrm{dom}((f / g)) \land a_{n} \neq c) \to n \in \mathbb{N} \to a_{n} \in \mathrm{dom}((f / g))$
\item $\varphi_{144} :\equiv r \in \mathbb{R} \land 0 < r \land r \leq d \land r \leq \delta \to x \in \mathbb{R} \to 0 < \lvert x - c \rvert \to \lvert x - c \rvert < r \to x \in \mathrm{dom}(g)$
\item $\varphi_{145} :\equiv f \circ a(n) \cdot g \circ a^{\ast}(n) = f \circ a(n) \cdot g \circ a^{\ast}(n)$
\item $\varphi_{146} :\equiv n \in \mathbb{N} \to f \circ a \otimes g \circ a^{\ast}(n) = f(a_{n}) \cdot g(a_{n})^{-1}$
\item $\varphi_{147} :\equiv n \in \mathbb{N} \to f(a_{n}) \cdot g(a_{n})^{-1} = f \circ a \otimes g \circ a^{\ast}(n)$
\item $\varphi_{148} :\equiv \forall n \in \mathbb{N}\,(a_{n} \in \mathrm{dom}((f / g)) \land a_{n} \neq c) \to \mathrm{Conv}((f / g) \circ a,\alpha \cdot \beta^{-1})$
\item $\varphi_{149} :\equiv \forall n \in \mathbb{N}\,(g \circ a(n) = g(a_{n})) \to n \in \mathbb{N} \to g \circ a(n) = g(a_{n})$
\item $\varphi_{150} :\equiv \forall n \in \mathbb{N}\,(f \circ a(n) = f(a_{n})) \to n \in \mathbb{N} \to f \circ a(n) = f(a_{n})$
\item $\varphi_{151} :\equiv \forall y\,\forall z\,(\lvert g(x) \rvert = y \to y = z \to \lvert g(x) \rvert = z) \to \forall z\,(\lvert g(x) \rvert = \lvert 0 \rvert \to \lvert 0 \rvert = z \to \lvert g(x) \rvert = z)$
\item $\varphi_{152} :\equiv \forall e \in \mathbb{R}\,(0 < e \to 0 < e \cdot 2^{-1}) \to \lvert \beta \rvert \in \mathbb{R} \to 0 < \lvert \beta \rvert \to 0 < \lvert \beta \rvert \cdot 2^{-1}$
\item $\varphi_{153} :\equiv \forall n \in \mathbb{N}\,(g \circ a(n) \neq 0) \to \mathrm{Conv}(f \circ a \otimes g \circ a^{\ast},\alpha \cdot \beta^{-1})$
\item $\varphi_{154} :\equiv \forall a\,(\mathrm{Seq}(a) \to \forall n \in \mathbb{N}\,(a_{n} \neq 0) \to \mathrm{Seq}(a^{\ast}) \land \forall n \in \mathbb{N}\,(a^{\ast}(n) = a_{n}^{-1}))$
\item $\varphi_{155} :\equiv \forall n \in \mathbb{N}\,(a_{n} \in \mathrm{dom}((f / g)) \land a_{n} \neq c) \to \forall n \in \mathbb{N}\,(a_{n} \in \mathrm{dom}((f / g)))$
\item $\varphi_{156} :\equiv r \in \mathbb{R} \land 0 < r \land r \leq d \land r \leq \delta \to x \in \mathbb{R} \to 0 < \lvert x - c \rvert \to \lvert x - c \rvert < r \to \lvert x - c \rvert < r$
\item $\varphi_{157} :\equiv r \in \mathbb{R} \land 0 < r \land r \leq d \land r \leq \delta \to x \in \mathbb{R} \to 0 < \lvert x - c \rvert \to \lvert x - c \rvert < r \to \lvert x - c \rvert < d$
\item $\varphi_{158} :\equiv n \in \mathbb{N} \to f \circ a(n) \cdot g \circ a(n)^{-1} = f(a_{n}) \cdot g \circ a(n)^{-1}$
\item $\varphi_{159} :\equiv n \in \mathbb{N} \to f \circ a \otimes g \circ a^{\ast}(n) = f(a_{n}) \cdot g \circ a(n)^{-1}$
\item $\varphi_{160} :\equiv \mathrm{Conv}(f \circ a \otimes g \circ a^{\ast},\alpha \cdot \beta^{-1}) \to \mathrm{Conv}((f / g) \circ a,\alpha \cdot \beta^{-1})$
\item $\varphi_{161} :\equiv r \in \mathbb{R} \land 0 < r \land r \leq d \land r \leq \delta \to x \in \mathbb{R} \to 0 < \lvert x - c \rvert \to \lvert x - c \rvert < r \to \lvert x - c \rvert < \delta$
\item $\varphi_{162} :\equiv r \in \mathbb{R} \land 0 < r \land r \leq d \land r \leq \delta \to x \in \mathbb{R} \to 0 < \lvert x - c \rvert \to \lvert x - c \rvert < r \to g(x) \in \mathbb{R}$
\item $\varphi_{163} :\equiv \forall y\,(f \circ a \otimes g \circ a^{\ast}(n) = y \to y = f \circ a \otimes g \circ a^{\ast}(n))$
\item $\varphi_{164} :\equiv \forall z\,(\lvert g(x) \rvert = \lvert 0 \rvert \to \lvert 0 \rvert = z \to \lvert g(x) \rvert = z) \to \lvert g(x) \rvert = \lvert 0 \rvert \to \lvert 0 \rvert = 0 \to \lvert g(x) \rvert = 0$
\item $\varphi_{165} :\equiv \forall y\,(\alpha \in \mathbb{R} \to y \in \mathbb{R} \to \alpha \cdot y \in \mathbb{R}) \to \alpha \in \mathbb{R} \to \beta^{-1} \in \mathbb{R} \to \alpha \cdot \beta^{-1} \in \mathbb{R}$
\item $\varphi_{166} :\equiv \forall n \in \mathbb{N}\,(a_{n} \in \mathrm{dom}((f / g)) \land a_{n} \neq c) \to n \in \mathbb{N} \to a_{n} \in \mathrm{dom}(g) \land a_{n} \neq c$
\item $\varphi_{167} :\equiv \forall n \in \mathbb{N}\,(a_{n} \in \mathrm{dom}((f / g)) \land a_{n} \neq c) \to n \in \mathbb{N} \to a_{n} \in \mathrm{dom}(f) \land a_{n} \neq c$
\item $\varphi_{168} :\equiv \mathrm{FLim}(f,c,\alpha) \to \mathrm{FLim}(g,c,\beta) \to \beta \neq 0 \to \forall x\,(\varphi_{42})$
\item $\varphi_{169} :\equiv \forall \alpha\,\forall a\,(\mathrm{FLim}(g,c,\alpha) \to \mathrm{Conv}(a,c) \to \forall n \in \mathbb{N}\,(a_{n} \in \mathrm{dom}(g) \land a_{n} \neq c) \to \mathrm{Conv}(g \circ a,\alpha))$
\item $\varphi_{170} :\equiv \forall x\,\forall y\,(x \in \mathbb{R} \to y \in \mathbb{R} \to x \cdot y \in \mathbb{R}) \to \forall y\,(\lvert \beta \rvert \in \mathbb{R} \to y \in \mathbb{R} \to \lvert \beta \rvert \cdot y \in \mathbb{R})$
\item $\varphi_{171} :\equiv (\mathrm{FLim}(f,c,\alpha) \to c \in \mathbb{R}) \to (\mathrm{FLim}(f,c,\alpha) \to \mathrm{FLim}(f,c,\alpha)) \to \varphi_{6}$
\item $\varphi_{172} :\equiv \forall n \in \mathbb{N}\,(a_{n} \in \mathrm{dom}((f / g)) \land a_{n} \neq c) \to n \in \mathbb{N} \to g \circ a(n) = g(a_{n})$
\item $\varphi_{173} :\equiv \forall n \in \mathbb{N}\,(a_{n} \in \mathrm{dom}((f / g)) \land a_{n} \neq c) \to n \in \mathbb{N} \to g(a_{n}) = g \circ a(n)$
\item $\varphi_{174} :\equiv \forall n \in \mathbb{N}\,(a_{n} \in \mathrm{dom}((f / g)) \land a_{n} \neq c) \to n \in \mathbb{N} \to f \circ a(n) = f(a_{n})$
\item $\varphi_{175} :\equiv \forall n \in \mathbb{N}\,(a_{n} \in \mathrm{dom}(g)) \to \mathrm{Seq}(g \circ a) \land \forall n \in \mathbb{N}\,(g \circ a(n) = g(a_{n}))$
\item $\varphi_{176} :\equiv \forall n \in \mathbb{N}\,(a_{n} \in \mathrm{dom}((f / g)) \land a_{n} \neq c) \to \forall n \in \mathbb{N}\,(a_{n} \in \mathrm{dom}(g) \land a_{n} \neq c)$
\item $\varphi_{177} :\equiv \forall n \in \mathbb{N}\,(a_{n} \in \mathrm{dom}((f / g)) \land a_{n} \neq c) \to \forall n \in \mathbb{N}\,(a_{n} \in \mathrm{dom}(f) \land a_{n} \neq c)$
\item $\varphi_{178} :\equiv \forall n \in \mathbb{N}\,(a_{n} \in \mathrm{dom}(f)) \to \mathrm{Seq}(f \circ a) \land \forall n \in \mathbb{N}\,(f \circ a(n) = f(a_{n}))$
\item $\varphi_{179} :\equiv n \in \mathbb{N} \to f \circ a \otimes g \circ a^{\ast}(n) = f \circ a(n) \cdot g \circ a(n)^{-1}$
\item $\varphi_{180} :\equiv (\mathrm{FLim}(f,c,\alpha) \to \alpha \in \mathbb{R}) \to (\mathrm{FLim}(f,c,\alpha) \to \mathrm{FLim}(f,c,\alpha)) \to \varphi_{10}$
\item $\varphi_{181} :\equiv \forall c\,\forall \alpha\,\forall a\,(\mathrm{FLim}(f,c,\alpha) \to \mathrm{Conv}(a,c) \to \varphi_{72})$
\item $\varphi_{182} :\equiv n \in \mathbb{N} \to f \circ a \otimes g \circ a^{\ast}(n) = f \circ a(n) \cdot g \circ a^{\ast}(n)$
\item $\varphi_{183} :\equiv \forall f\,\forall c\,(\mathrm{PNbhd}(f,c) \leftrightarrow \exists r\,(r \in \mathbb{R} \land 0 < r \land \forall x \in \mathbb{R}\,(0 < \lvert x - c \rvert \to \lvert x - c \rvert < r \to x \in \mathrm{dom}(f))))$
\item $\varphi_{184} :\equiv \forall n \in \mathbb{N}\,(a_{n} \in \mathrm{dom}((f / g)) \land a_{n} \neq c) \to n \in \mathbb{N} \to a_{n} \in \mathrm{dom}((f / g)) \land a_{n} \neq c$
\item $\varphi_{185} :\equiv \forall n \in \mathbb{N}\,(f \circ a \otimes g \circ a^{\ast}(n) = f \circ a(n) \cdot g \circ a^{\ast}(n))$
\item $\varphi_{186} :\equiv \forall y\,\forall z\,(f \circ a \otimes g \circ a^{\ast}(n) = y \to y = z \to f \circ a \otimes g \circ a^{\ast}(n) = z)$
\item $\varphi_{187} :\equiv \forall v_{766}\,\forall v_{767}\,(\varphi_{230})$
\item $\varphi_{188} :\equiv \forall z\,((f / g) \circ a(n) = (f / g)(a_{n}) \to (f / g)(a_{n}) = z \to (f / g) \circ a(n) = z)$
\item $\varphi_{189} :\equiv \forall a\,\forall \alpha\,\forall b\,\forall \beta\,(\mathrm{Conv}(a,\alpha) \to \mathrm{Conv}(b,\beta) \to \beta \neq 0 \to \forall n \in \mathbb{N}\,(b_{n} \neq 0) \to \mathrm{Conv}(a \otimes b^{\ast},\alpha \cdot \beta^{-1}))$
\item $\varphi_{190} :\equiv \mathrm{FLim}(f,c,\alpha) \land \mathrm{FLim}(g,c,\beta) \land \beta \neq 0 \land \mathrm{Seq}(a) \land \mathrm{Conv}(a,c) \land \forall n \in \mathbb{N}\,(a_{n} \in \mathrm{dom}((f / g)) \land a_{n} \neq c)$
\item $\varphi_{191} :\equiv \forall n \in \mathbb{N}\,(a_{n} \in \mathrm{dom}((f / g)) \land a_{n} \neq c) \to \forall n \in \mathbb{N}\,(a_{n} \in \mathrm{dom}((f / g)) \land a_{n} \neq c)$
\item $\varphi_{192} :\equiv \forall y\,(g(a_{n}) = y \to y = g(a_{n})) \to g(a_{n}) = g \circ a(n) \to g \circ a(n) = g(a_{n})$
\item $\varphi_{193} :\equiv \forall n \in \mathbb{N}\,(a_{n} \in \mathrm{dom}((f / g)) \land a_{n} \neq c) \to n \in \mathbb{N} \to n \in \mathbb{N} \to g \circ a(n) = g(a_{n})$
\item $\varphi_{194} :\equiv \forall n \in \mathbb{N}\,(a_{n} \in \mathrm{dom}((f / g)) \land a_{n} \neq c) \to n \in \mathbb{N} \to n \in \mathbb{N} \to f \circ a(n) = f(a_{n})$
\item $\varphi_{195} :\equiv \forall n \in \mathbb{N}\,(a_{n} \in \mathrm{dom}((f / g)) \land a_{n} \neq c) \to \mathrm{Conv}(f \circ a \otimes g \circ a^{\ast},\alpha \cdot \beta^{-1})$
\item $\varphi_{196} :\equiv \mathrm{PNbhd}(g,c) \to \exists d\,(d \in \mathbb{R} \land 0 < d \land \forall x \in \mathbb{R}\,(\varphi_{46}))$
\item $\varphi_{197} :\equiv \mathrm{FLim}(g,c,\beta) \to \beta \neq 0 \to \mathrm{Seq}(a) \to \mathrm{Conv}(a,c) \to \varphi_{82}$
\item $\varphi_{198} :\equiv \mathrm{FLim}(g,c,\beta) \to \beta \neq 0 \to \mathrm{Seq}(a) \to \mathrm{Conv}(a,c) \to \varphi_{83}$
\item $\varphi_{199} :\equiv \forall n \in \mathbb{N}\,(a_{n} \in \mathrm{dom}((f / g)) \land a_{n} \neq c) \to n \in \mathbb{N} \to \forall n \in \mathbb{N}\,(g \circ a(n) = g(a_{n}))$
\item $\varphi_{200} :\equiv \forall n \in \mathbb{N}\,(a_{n} \in \mathrm{dom}((f / g)) \land a_{n} \neq c) \to n \in \mathbb{N} \to \forall n \in \mathbb{N}\,(f \circ a(n) = f(a_{n}))$
\item $\varphi_{201} :\equiv \forall y\,(\lvert \beta \rvert \in \mathbb{R} \to y \in \mathbb{R} \to \lvert \beta \rvert \cdot y \in \mathbb{R}) \to \lvert \beta \rvert \in \mathbb{R} \to 2^{-1} \in \mathbb{R} \to \lvert \beta \rvert \cdot 2^{-1} \in \mathbb{R}$
\item $\varphi_{202} :\equiv \forall n \in \mathbb{N}\,(a_{n} \in \mathrm{dom}((f / g)) \land a_{n} \neq c) \to n \in \mathbb{N} \to g \circ a^{\ast}(n) = g \circ a(n)^{-1}$
\item $\varphi_{203} :\equiv \forall n \in \mathbb{N}\,(a_{n} \in \mathrm{dom}((f / g)) \land a_{n} \neq c) \to n \in \mathbb{N} \to (f / g) \circ a(n) = (f / g)(a_{n})$
\item $\varphi_{204} :\equiv \forall c\,(\mathrm{PNbhd}((f / g),c) \leftrightarrow \exists r\,(r \in \mathbb{R} \land 0 < r \land \varphi_{78}))$
\item $\varphi_{205} :\equiv \forall y\,(g \circ a(n) = y \to y = g \circ a(n)) \to g \circ a(n) = g(a_{n}) \to g(a_{n}) = g \circ a(n)$
\item $\varphi_{206} :\equiv \mathrm{FLim}(g,c,\beta) \to \beta \neq 0 \to \mathrm{Seq}(a) \to \mathrm{Conv}(a,c) \to \varphi_{97}$
\item $\varphi_{207} :\equiv \mathrm{FLim}(g,c,\beta) \to \beta \neq 0 \to \mathrm{Seq}(a) \to \mathrm{Conv}(a,c) \to \varphi_{98}$
\item $\varphi_{208} :\equiv \mathrm{FLim}(g,c,\beta) \to \beta \neq 0 \to \mathrm{Seq}(a) \to \mathrm{Conv}(a,c) \to \varphi_{99}$
\item $\varphi_{209} :\equiv d \in \mathbb{R} \land 0 < d \land \forall x \in \mathbb{R}\,(\varphi_{46}) \to \mathrm{PNbhd}((f / g),c)$
\item $\varphi_{210} :\equiv \mathrm{FLim}(g,c,\beta) \to \beta \neq 0 \to \mathrm{Seq}(a) \to \mathrm{Conv}(a,c) \to \varphi_{103}$
\item $\varphi_{211} :\equiv \forall w\,(g(a_{n}) = w \to g(a_{n}) = 0 \to w = 0) \to g(a_{n}) = g \circ a(n) \to g(a_{n}) = 0 \to g \circ a(n) = 0$
\item $\varphi_{212} :\equiv \mathrm{FLim}(g,c,\beta) \to \beta \neq 0 \to \mathrm{Seq}(a) \to \mathrm{Conv}(a,c) \to \varphi_{106}$
\item $\varphi_{213} :\equiv \forall b\,\forall \beta\,(\mathrm{Conv}(f \circ a,\alpha) \to \mathrm{Conv}(b,\beta) \to \beta \neq 0 \to \forall n \in \mathbb{N}\,(b_{n} \neq 0) \to \mathrm{Conv}(f \circ a \otimes b^{\ast},\alpha \cdot \beta^{-1}))$
\item $\varphi_{214} :\equiv \forall n \in \mathbb{N}\,(g \circ a^{\ast}(n) = g \circ a(n)^{-1}) \to n \in \mathbb{N} \to g \circ a^{\ast}(n) = g \circ a(n)^{-1}$
\item $\varphi_{215} :\equiv \forall n \in \mathbb{N}\,((f / g) \circ a(n) = (f / g)(a_{n})) \to n \in \mathbb{N} \to (f / g) \circ a(n) = (f / g)(a_{n})$
\item $\varphi_{216} :\equiv \mathrm{FLim}(g,c,\beta) \to \beta \neq 0 \to \mathrm{Seq}(a) \to \mathrm{Conv}(a,c) \to \varphi_{113}$
\item $\varphi_{217} :\equiv \mathrm{FLim}(g,c,\beta) \to \beta \neq 0 \to \mathrm{Seq}(a) \to \mathrm{Conv}(a,c) \to \varphi_{117}$
\item $\varphi_{218} :\equiv \mathrm{FLim}(g,c,\beta) \to \beta \neq 0 \to \mathrm{Seq}(a) \to \mathrm{Conv}(a,c) \to \varphi_{119}$
\item $\varphi_{219} :\equiv \forall n \in \mathbb{N}\,(a_{n} \in \mathrm{dom}((f / g)) \land a_{n} \neq c) \to n \in \mathbb{N} \to n \in \mathbb{N} \to g \circ a^{\ast}(n) = g \circ a(n)^{-1}$
\item $\varphi_{220} :\equiv \forall n \in \mathbb{N}\,(a_{n} \in \mathrm{dom}((f / g)) \land a_{n} \neq c) \to n \in \mathbb{N} \to (f / g)(a_{n}) = f(a_{n}) \cdot g(a_{n})^{-1}$
\item $\varphi_{221} :\equiv \forall n \in \mathbb{N}\,(a_{n} \in \mathrm{dom}((f / g)) \land a_{n} \neq c) \to n \in \mathbb{N} \to n \in \mathbb{N} \to (f / g) \circ a(n) = (f / g)(a_{n})$
\item $\varphi_{222} :\equiv r \in \mathbb{R} \land 0 < r \land r \leq d \land r \leq \delta \to \exists r\,(r \in \mathbb{R} \land 0 < r \land \varphi_{78})$
\item $\varphi_{223} :\equiv \delta \in \mathbb{R} \land 0 < \delta \land \forall x \in \mathrm{dom}(g)\,(\varphi_{94}) \to \mathrm{PNbhd}(g,c)$
\item $\varphi_{224} :\equiv \delta \in \mathbb{R} \land 0 < \delta \land \forall x \in \mathrm{dom}(g)\,(\varphi_{94}) \to \mathrm{PNbhd}(f,c)$
\item $\varphi_{225} :\equiv \forall w\,(g \circ a(n) = w \to g \circ a(n) = 0 \to w = 0) \to g \circ a(n) = g(a_{n}) \to g \circ a(n) = 0 \to g(a_{n}) = 0$
\item $\varphi_{226} :\equiv \mathrm{FLim}(g,c,\beta) \to \beta \neq 0 \to \mathrm{Seq}(a) \to \mathrm{Conv}(a,c) \to \varphi_{123}$
\item $\varphi_{227} :\equiv \mathrm{FLim}(g,c,\beta) \to \beta \neq 0 \to \mathrm{Seq}(a) \to \mathrm{Conv}(a,c) \to \varphi_{124}$
\item $\varphi_{228} :\equiv \mathrm{FLim}(g,c,\beta) \to \beta \neq 0 \to \mathrm{Seq}(a) \to \mathrm{Conv}(a,c) \to \varphi_{125}$
\item $\varphi_{229} :\equiv \forall n \in \mathbb{N}\,(a_{n} \in \mathrm{dom}((f / g)) \land a_{n} \neq c) \to n \in \mathbb{N} \to \forall n \in \mathbb{N}\,(g \circ a^{\ast}(n) = g \circ a(n)^{-1})$
\item $\varphi_{230} :\equiv v_{766} \in \mathbb{R} \to v_{767} \in \mathbb{R} \to 0 < v_{766} \to 0 < v_{767} \to \neg \forall u_{4}\,(\varphi_{85})$
\item $\varphi_{231} :\equiv \forall z\,((f / g)(a_{n}) = f(a_{n}) \cdot g(a_{n})^{-1} \to f(a_{n}) \cdot g(a_{n})^{-1} = z \to (f / g)(a_{n}) = z)$
\item $\varphi_{232} :\equiv \forall b\,(\mathrm{Seq}(f \circ a) \to \mathrm{Seq}(b) \to \mathrm{Seq}(f \circ a \otimes b) \land \forall n \in \mathbb{N}\,(f \circ a \otimes b(n) = f \circ a(n) \cdot b_{n}))$
\item $\varphi_{233} :\equiv \mathrm{FLim}(g,c,\beta) \to \beta \neq 0 \to \mathrm{Seq}(a) \to \mathrm{Conv}(a,c) \to \varphi_{136}$
\item $\varphi_{234} :\equiv f(a_{n}) \cdot g \circ a(n)^{-1} = f(a_{n}) \cdot z^{-1} \to f(a_{n}) \cdot g \circ a(n)^{-1} = f(a_{n}) \cdot w^{-1}$
\item $\varphi_{235} :\equiv \forall n \in \mathbb{N}\,(a_{n} \in \mathrm{dom}((f / g)) \land a_{n} \neq c) \to n \in \mathbb{N} \to (f / g)(a_{n}) = f \circ a \otimes g \circ a^{\ast}(n)$
\item $\varphi_{236} :\equiv \mathrm{FLim}(g,c,\beta) \to \beta \neq 0 \to \mathrm{Seq}(a) \to \mathrm{Conv}(a,c) \to \varphi_{143}$
\item $\varphi_{237} :\equiv \delta \in \mathbb{R} \land 0 < \delta \land \forall x \in \mathrm{dom}(g)\,(\varphi_{94}) \to \mathrm{PNbhd}((f / g),c)$
\item $\varphi_{238} :\equiv \mathrm{FLim}(g,c,\beta) \to \beta \neq 0 \to \mathrm{Seq}(a) \to \mathrm{Conv}(a,c) \to \varphi_{148}$
\item $\varphi_{239} :\equiv \mathrm{FLim}(g,c,\beta) \to \beta \neq 0 \to \mathrm{Seq}(a) \to \mathrm{Conv}(a,c) \to \varphi_{155}$
\item $\varphi_{240} :\equiv \forall n \in \mathbb{N}\,(a_{n} \in \mathrm{dom}((f / g)) \land a_{n} \neq c) \to n \in \mathbb{N} \to (f / g) \circ a(n) = f \circ a \otimes g \circ a^{\ast}(n)$
\item $\varphi_{241} :\equiv \forall n \in \mathbb{N}\,(a_{n} \in \mathrm{dom}((f / g)) \land a_{n} \neq c) \to n \in \mathbb{N} \to f \circ a \otimes g \circ a^{\ast}(n) = (f / g) \circ a(n)$
\item $\varphi_{242} :\equiv \mathrm{FLim}(g,c,\beta) \to \beta \neq 0 \to \forall a\,(\mathrm{Seq}(a) \to \mathrm{Conv}(a,c) \to \varphi_{148})$
\item $\varphi_{243} :\equiv d \in \mathbb{R} \land 0 < d \land \forall x \in \mathbb{R}\,(\varphi_{46}) \to \exists u_{4}\,(u_{4} \in \mathbb{R} \land 0 < u_{4} \land u_{4} \leq d \land u_{4} \leq \delta)$
\item $\varphi_{244} :\equiv \forall n \in \mathbb{N}\,(a_{n} \in \mathrm{dom}((f / g)) \land a_{n} \neq c) \to n \in \mathbb{N} \to \varphi_{81}$
\item $\varphi_{245} :\equiv \mathrm{FLim}(g,c,\beta) \to \beta \neq 0 \to \mathrm{Seq}(a) \to \mathrm{Conv}(a,c) \to \varphi_{166}$
\item $\varphi_{246} :\equiv \mathrm{FLim}(g,c,\beta) \to \beta \neq 0 \to \mathrm{Seq}(a) \to \mathrm{Conv}(a,c) \to \varphi_{167}$
\item $\varphi_{247} :\equiv \mathrm{FLim}(g,c,\beta) \to \beta \neq 0 \to \exists \delta\,(\delta \in \mathbb{R} \land 0 < \delta \land \forall x \in \mathrm{dom}(g)\,(\varphi_{94}))$
\item $\varphi_{248} :\equiv (f / g)(a_{n}) = f \circ a \otimes g \circ a^{\ast}(n) \to (f / g) \circ a(n) = f \circ a \otimes g \circ a^{\ast}(n)$
\item $\varphi_{249} :\equiv \mathrm{Conv}(a,c) \to \forall n \in \mathbb{N}\,(a_{n} \in \mathrm{dom}((f / g)) \land a_{n} \neq c) \to \varphi_{90}$
\item $\varphi_{250} :\equiv \forall n \in \mathbb{N}\,(a_{n} \in \mathrm{dom}((f / g)) \land a_{n} \neq c) \to \forall n \in \mathbb{N}\,(\varphi_{81})$
\item $\varphi_{251} :\equiv \mathrm{FLim}(g,c,\beta) \to \beta \neq 0 \to \mathrm{Seq}(a) \to \mathrm{Conv}(a,c) \to \varphi_{172}$
\item $\varphi_{252} :\equiv \mathrm{FLim}(g,c,\beta) \to \beta \neq 0 \to \mathrm{Seq}(a) \to \mathrm{Conv}(a,c) \to \varphi_{173}$
\item $\varphi_{253} :\equiv \mathrm{FLim}(g,c,\beta) \to \beta \neq 0 \to \mathrm{Seq}(a) \to \mathrm{Conv}(a,c) \to \varphi_{174}$
\item $\varphi_{254} :\equiv 0 < \lvert \beta \rvert \cdot 2^{-1} \to \exists \delta\,(\delta \in \mathbb{R} \land 0 < \delta \land \forall x \in \mathrm{dom}(g)\,(\varphi_{94}))$
\item $\varphi_{255} :\equiv \mathrm{FLim}(g,c,\beta) \to \beta \neq 0 \to \mathrm{Seq}(a) \to \mathrm{Conv}(a,c) \to \varphi_{176}$
\item $\varphi_{256} :\equiv \mathrm{FLim}(g,c,\beta) \to \beta \neq 0 \to \mathrm{Seq}(a) \to \mathrm{Conv}(a,c) \to \varphi_{177}$
\item $\varphi_{257} :\equiv f \circ a(n) \cdot g \circ a^{\ast}(n) = f \circ a(n) \cdot z \to f \circ a(n) \cdot g \circ a^{\ast}(n) = f \circ a(n) \cdot w$
\item $\varphi_{258} :\equiv \forall b\,\forall \alpha\,(\mathrm{Seq}(b) \to \forall n \in \mathbb{N}\,(f \circ a \otimes g \circ a^{\ast}(n) = b_{n}) \to \mathrm{Conv}(f \circ a \otimes g \circ a^{\ast},\alpha) \to \mathrm{Conv}(b,\alpha))$
\item $\varphi_{259} :\equiv (f / g) \circ a(n) = f \circ a \otimes g \circ a^{\ast}(n) \to f \circ a \otimes g \circ a^{\ast}(n) = (f / g) \circ a(n)$
\item $\varphi_{260} :\equiv \mathrm{Conv}(a,c) \to \forall n \in \mathbb{N}\,(a_{n} \in \mathrm{dom}((f / g)) \land a_{n} \neq c) \to \varphi_{101}$
\item $\varphi_{261} :\equiv \mathrm{Conv}(a,c) \to \forall n \in \mathbb{N}\,(a_{n} \in \mathrm{dom}((f / g)) \land a_{n} \neq c) \to \varphi_{102}$
\item $\varphi_{262} :\equiv \mathrm{FLim}(g,c,\beta) \to \beta \neq 0 \to \delta \in \mathbb{R} \land 0 < \delta \land \forall x \in \mathrm{dom}(g)\,(\varphi_{94}) \to 0 < \delta$
\item $\varphi_{263} :\equiv \mathrm{FLim}(g,c,\beta) \to \beta \neq 0 \to \mathrm{Seq}(a) \to \mathrm{Conv}(a,c) \to \varphi_{184}$
\item $\varphi_{264} :\equiv \forall \varepsilon\,(\mathrm{FLim}(g,c,\beta) \to \varepsilon \in \mathbb{R} \to 0 < \varepsilon \to \exists \delta\,(\delta \in \mathbb{R} \land 0 < \delta \land \forall x \in \mathrm{dom}(g)\,(0 < \lvert x - c \rvert \to \lvert x - c \rvert < \delta \to \lvert g(x) - \beta \rvert < \varepsilon)))$
\item $\varphi_{265} :\equiv f \circ a(n) \cdot g \circ a(n)^{-1} = z \cdot g \circ a(n)^{-1} \to f \circ a(n) \cdot g \circ a(n)^{-1} = w \cdot g \circ a(n)^{-1}$
\item $\varphi_{266} :\equiv \forall a\,(\mathrm{Seq}(a) \to \mathrm{Conv}(a,c) \to \varphi_{148}) \to \mathrm{FLim}((f / g),c,\alpha \cdot \beta^{-1})$
\item $\varphi_{267} :\equiv \mathrm{FLim}(f,c,\alpha) \land \mathrm{FLim}(g,c,\beta) \land \beta \neq 0 \to \mathrm{Seq}(a) \to \mathrm{Conv}(a,c) \to \varphi_{148}$
\item $\varphi_{268} :\equiv \mathrm{FLim}(g,c,\beta) \to \beta \neq 0 \to \delta \in \mathbb{R} \land 0 < \delta \land \forall x \in \mathrm{dom}(g)\,(\varphi_{94}) \to \delta \in \mathbb{R}$
\item $\varphi_{269} :\equiv \mathrm{FLim}(g,c,\beta) \to \beta \neq 0 \to \mathrm{Seq}(a) \to \mathrm{Conv}(a,c) \to \varphi_{191}$
\item $\varphi_{270} :\equiv \forall \beta\,(\mathrm{Conv}(f \circ a,\alpha) \to \mathrm{Conv}(g \circ a,\beta) \to \beta \neq 0 \to \varphi_{153})$
\item $\varphi_{271} :\equiv \forall \alpha\,\forall \varepsilon\,(\mathrm{FLim}(g,c,\alpha) \to \varepsilon \in \mathbb{R} \to 0 < \varepsilon \to \exists \delta\,(\delta \in \mathbb{R} \land 0 < \delta \land \forall x \in \mathrm{dom}(g)\,(0 < \lvert x - c \rvert \to \lvert x - c \rvert < \delta \to \lvert g(x) - \alpha \rvert < \varepsilon)))$
\item $\varphi_{272} :\equiv f(a_{n}) \cdot g(a_{n})^{-1} = f \circ a \otimes g \circ a^{\ast}(n) \to (f / g)(a_{n}) = f \circ a \otimes g \circ a^{\ast}(n)$
\item $\varphi_{273} :\equiv \mathrm{Conv}(a,c) \to \forall n \in \mathbb{N}\,(a_{n} \in \mathrm{dom}((f / g)) \land a_{n} \neq c) \to \varphi_{118}$
\item $\varphi_{274} :\equiv \mathrm{FLim}(g,c,\beta) \to \beta \neq 0 \to \mathrm{Seq}(a) \to \mathrm{Conv}(a,c) \to \varphi_{193}$
\item $\varphi_{275} :\equiv \mathrm{FLim}(g,c,\beta) \to \beta \neq 0 \to \mathrm{Seq}(a) \to \mathrm{Conv}(a,c) \to \varphi_{194}$
\item $\varphi_{276} :\equiv \mathrm{FLim}(g,c,\beta) \to \beta \neq 0 \to \mathrm{Seq}(a) \to \mathrm{Conv}(a,c) \to \varphi_{195}$
\item $\varphi_{277} :\equiv \mathrm{Conv}(a,c) \to \forall n \in \mathbb{N}\,(a_{n} \in \mathrm{dom}((f / g)) \land a_{n} \neq c) \to \varphi_{126}$
\item $\varphi_{278} :\equiv d \in \mathbb{R} \land 0 < d \land \forall x \in \mathbb{R}\,(\varphi_{46}) \to r \in \mathbb{R} \land 0 < r \land r \leq d \land r \leq \delta \to x \in \mathbb{R} \to x \in \mathbb{R}$
\item $\varphi_{279} :\equiv r \in \mathbb{R} \land 0 < r \land r \leq d \land r \leq \delta \to x \in \mathbb{R} \to 0 < \lvert x - c \rvert \to \lvert x - c \rvert < r \to x \in \mathbb{R} \to \varphi_{46}$
\item $\varphi_{280} :\equiv \mathrm{RFn}((f / g)) \to \mathrm{Seq}(a) \to \forall n \in \mathbb{N}\,(a_{n} \in \mathrm{dom}((f / g))) \to \varphi_{130}$
\item $\varphi_{281} :\equiv \mathrm{FLim}(g,c,\beta) \to \beta \neq 0 \to \mathrm{Seq}(a) \to \mathrm{Conv}(a,c) \to \varphi_{199}$
\item $\varphi_{282} :\equiv \mathrm{FLim}(g,c,\beta) \to \beta \neq 0 \to \mathrm{Seq}(a) \to \mathrm{Conv}(a,c) \to \varphi_{200}$
\item $\varphi_{283} :\equiv \mathrm{Conv}(a,c) \to \forall n \in \mathbb{N}\,(a_{n} \in \mathrm{dom}((f / g)) \land a_{n} \neq c) \to \varphi_{130}$
\item $\varphi_{284} :\equiv \mathrm{Conv}(a,c) \to \forall n \in \mathbb{N}\,(a_{n} \in \mathrm{dom}((f / g)) \land a_{n} \neq c) \to \varphi_{132}$
\item $\varphi_{285} :\equiv \mathrm{FLim}(g,c,\beta) \to \beta \neq 0 \to \mathrm{Seq}(a) \to \mathrm{Conv}(a,c) \to \varphi_{202}$
\item $\varphi_{286} :\equiv \mathrm{FLim}(g,c,\beta) \to \beta \neq 0 \to \mathrm{Seq}(a) \to \mathrm{Conv}(a,c) \to \varphi_{203}$
\item $\varphi_{287} :\equiv \forall f\,\forall c\,\forall \alpha\,\forall \varepsilon\,(\mathrm{FLim}(f,c,\alpha) \to \varepsilon \in \mathbb{R} \to 0 < \varepsilon \to \exists \delta\,(\delta \in \mathbb{R} \land 0 < \delta \land \forall x \in \mathrm{dom}(f)\,(0 < \lvert x - c \rvert \to \lvert x - c \rvert < \delta \to \lvert f(x) - \alpha \rvert < \varepsilon)))$
\item $\varphi_{288} :\equiv r \in \mathbb{R} \land 0 < r \land r \leq d \land r \leq \delta \to x \in \mathbb{R} \to 0 < \lvert x - c \rvert \to \lvert x - c \rvert < r \to \forall x \in \mathbb{R}\,(\varphi_{46})$
\item $\varphi_{289} :\equiv \forall n \in \mathbb{N}\,(a_{n} \in \mathrm{dom}((f / g)) \land a_{n} \neq c) \to n \in \mathbb{N} \to \varphi_{139}$
\item $\varphi_{290} :\equiv \forall f\,\forall c\,\forall \alpha\,(\mathrm{RFn}(f) \to c \in \mathbb{R} \to \alpha \in \mathbb{R} \to \mathrm{PNbhd}(f,c) \to \forall a\,(\mathrm{Seq}(a) \to \mathrm{Conv}(a,c) \to \varphi_{72}) \to \mathrm{FLim}(f,c,\alpha))$
\item $\varphi_{291} :\equiv \mathrm{Conv}(a,c) \to \forall n \in \mathbb{N}\,(a_{n} \in \mathrm{dom}((f / g)) \land a_{n} \neq c) \to \varphi_{146}$
\item $\varphi_{292} :\equiv \mathrm{Conv}(a,c) \to \forall n \in \mathbb{N}\,(a_{n} \in \mathrm{dom}((f / g)) \land a_{n} \neq c) \to \varphi_{147}$
\item $\varphi_{293} :\equiv \mathrm{RFn}(f) \to \mathrm{RFn}(g) \to \mathrm{RFn}((f / g)) \land \forall x\,(\varphi_{42}) \land \varphi_{54}$
\item $\varphi_{294} :\equiv f(a_{n}) \cdot g \circ a(n)^{-1} = f(a_{n}) \cdot g(a_{n})^{-1} \to f \circ a \otimes g \circ a^{\ast}(n) = f(a_{n}) \cdot g(a_{n})^{-1}$
\item $\varphi_{295} :\equiv f(a_{n}) \cdot g \circ a(n)^{-1} = f(a_{n}) \cdot g \circ a(n)^{-1} \to f(a_{n}) \cdot g \circ a(n)^{-1} = f(a_{n}) \cdot g(a_{n})^{-1}$
\item $\varphi_{296} :\equiv \mathrm{Conv}(a,c) \to \forall n \in \mathbb{N}\,(a_{n} \in \mathrm{dom}((f / g)) \land a_{n} \neq c) \to \varphi_{158}$
\item $\varphi_{297} :\equiv \mathrm{Conv}(a,c) \to \forall n \in \mathbb{N}\,(a_{n} \in \mathrm{dom}((f / g)) \land a_{n} \neq c) \to \varphi_{159}$
\item $\varphi_{298} :\equiv f \circ a \otimes g \circ a^{\ast}(n) = f(a_{n}) \cdot g(a_{n})^{-1} \to f(a_{n}) \cdot g(a_{n})^{-1} = f \circ a \otimes g \circ a^{\ast}(n)$
\item $\varphi_{299} :\equiv \mathrm{FLim}(g,c,\beta) \to \beta \neq 0 \to \mathrm{RFn}((f / g)) \land \forall x\,(\varphi_{42}) \land \varphi_{54}$
\item $\varphi_{300} :\equiv \mathrm{FLim}(g,c,\beta) \to \beta \neq 0 \to \mathrm{Seq}(a) \to \mathrm{Conv}(a,c) \to \varphi_{219}$
\item $\varphi_{301} :\equiv \mathrm{Conv}(a,c) \to \forall n \in \mathbb{N}\,(a_{n} \in \mathrm{dom}((f / g)) \land a_{n} \neq c) \to \varphi_{179}$
\item $\varphi_{302} :\equiv \mathrm{FLim}(g,c,\beta) \to \beta \neq 0 \to \mathrm{Seq}(a) \to \mathrm{Conv}(a,c) \to \varphi_{220}$
\item $\varphi_{303} :\equiv \mathrm{FLim}(g,c,\beta) \to \beta \neq 0 \to \mathrm{Seq}(a) \to \mathrm{Conv}(a,c) \to \varphi_{221}$
\item $\varphi_{304} :\equiv \forall n \in \mathbb{N}\,(a_{n} \in \mathrm{dom}((f / g)) \land a_{n} \neq c) \to n \in \mathbb{N} \to \varphi_{182}$
\item $\varphi_{305} :\equiv \mathrm{Conv}(a,c) \to \forall n \in \mathbb{N}\,(a_{n} \in \mathrm{dom}((f / g)) \land a_{n} \neq c) \to \varphi_{182}$
\item $\varphi_{306} :\equiv \mathrm{FLim}(g,c,\beta) \to \beta \neq 0 \to \mathrm{Seq}(a) \to \mathrm{Conv}(a,c) \to \varphi_{229}$
\item $\varphi_{307} :\equiv \mathrm{FLim}(f,c,\alpha) \land \mathrm{FLim}(g,c,\beta) \land \beta \neq 0 \land (\delta \in \mathbb{R} \land 0 < \delta \land \forall x \in \mathrm{dom}(g)\,(\varphi_{94})) \land \neg \mathrm{PNbhd}((f / g),c)$
\item $\varphi_{308} :\equiv \forall n \in \mathbb{N}\,(a_{n} \in \mathrm{dom}((f / g)) \land a_{n} \neq c) \to n \in \mathbb{N} \to \varphi_{185}$
\item $\varphi_{309} :\equiv f \circ a(n) \cdot g \circ a(n)^{-1} = f(a_{n}) \cdot g \circ a(n)^{-1} \to f \circ a \otimes g \circ a^{\ast}(n) = f(a_{n}) \cdot g \circ a(n)^{-1}$
\item $\varphi_{310} :\equiv \forall w\,(g \circ a(n) = w \to f(a_{n}) \cdot g \circ a(n)^{-1} = f(a_{n}) \cdot g \circ a(n)^{-1} \to f(a_{n}) \cdot g \circ a(n)^{-1} = f(a_{n}) \cdot w^{-1})$
\item $\varphi_{311} :\equiv \mathrm{FLim}(g,c,\beta) \to \beta \neq 0 \to \mathrm{Seq}(a) \to \mathrm{Conv}(a,c) \to \varphi_{235}$
\item $\varphi_{312} :\equiv \varphi_{131} \to f \circ a(n) \cdot g \circ a(n)^{-1} = f(a_{n}) \cdot g \circ a(n)^{-1}$
\item $\varphi_{313} :\equiv \mathrm{FLim}(g,c,\beta) \to \beta \neq 0 \to \mathrm{Seq}(a) \to \mathrm{Conv}(a,c) \to \varphi_{240}$
\item $\varphi_{314} :\equiv \mathrm{FLim}(g,c,\beta) \to \beta \neq 0 \to \mathrm{Seq}(a) \to \mathrm{Conv}(a,c) \to \varphi_{241}$
\item $\varphi_{315} :\equiv \forall \delta\,(\varphi_{28}) \to \forall \delta\,(\neg (\delta \in \mathbb{R} \land 0 < \delta \land \forall x \in \mathrm{dom}(g)\,(\varphi_{94})))$
\item $\varphi_{316} :\equiv \mathrm{FLim}(g,c,\beta) \to \beta \neq 0 \to \mathrm{Seq}(a) \to \mathrm{Conv}(a,c) \to \varphi_{244}$
\item $\varphi_{317} :\equiv \forall z\,(f \circ a \otimes g \circ a^{\ast}(n) = f(a_{n}) \cdot g \circ a(n)^{-1} \to f(a_{n}) \cdot g \circ a(n)^{-1} = z \to f \circ a \otimes g \circ a^{\ast}(n) = z)$
\item $\varphi_{318} :\equiv \mathrm{FLim}(g,c,\beta) \to \beta \neq 0 \to \mathrm{Seq}(a) \to \mathrm{Conv}(a,c) \to \varphi_{250}$
\item $\varphi_{319} :\equiv \varphi_{139} \to f \circ a \otimes g \circ a^{\ast}(n) = f \circ a(n) \cdot g \circ a(n)^{-1}$
\item $\varphi_{320} :\equiv \mathrm{FLim}(f,c,\alpha) \to \mathrm{FLim}(g,c,\beta) \to \beta \neq 0 \to \mathrm{Seq}(a) \to \varphi_{249}$
\item $\varphi_{321} :\equiv \mathrm{FLim}(f,c,\alpha) \to \mathrm{FLim}(g,c,\beta) \to \beta \neq 0 \to \mathrm{Seq}(a) \to \varphi_{260}$
\item $\varphi_{322} :\equiv \mathrm{FLim}(f,c,\alpha) \to \mathrm{FLim}(g,c,\beta) \to \beta \neq 0 \to \mathrm{Seq}(a) \to \varphi_{261}$
\item $\varphi_{323} :\equiv \forall r\,(\neg (r \in \mathbb{R} \land 0 < r \land \varphi_{78})) \to \neg (r \in \mathbb{R} \land 0 < r \land \varphi_{78})$
\item $\varphi_{324} :\equiv \forall z\,(f \circ a \otimes g \circ a^{\ast}(n) = f \circ a(n) \cdot g \circ a(n)^{-1} \to f \circ a(n) \cdot g \circ a(n)^{-1} = z \to f \circ a \otimes g \circ a^{\ast}(n) = z)$
\item $\varphi_{325} :\equiv \forall w\,(f \circ a(n) = w \to \varphi_{131} \to f \circ a(n) \cdot g \circ a(n)^{-1} = w \cdot g \circ a(n)^{-1})$
\item $\varphi_{326} :\equiv (\mathrm{FLim}(f,c,\alpha) \to \mathrm{FLim}(f,c,\alpha)) \to \mathrm{FLim}(f,c,\alpha) \to \mathrm{FLim}(g,c,\beta) \to \beta \neq 0 \to \varphi_{224}$
\item $\varphi_{327} :\equiv \forall a\,(\mathrm{FLim}(g,c,\beta) \to \mathrm{Conv}(a,c) \to \varphi_{71}) \to \mathrm{FLim}(g,c,\beta) \to \mathrm{Conv}(a,c) \to \varphi_{71}$
\item $\varphi_{328} :\equiv \forall a\,(\mathrm{FLim}(f,c,\alpha) \to \mathrm{Conv}(a,c) \to \varphi_{72}) \to \mathrm{FLim}(f,c,\alpha) \to \mathrm{Conv}(a,c) \to \varphi_{72}$
\item $\varphi_{329} :\equiv \forall \alpha\,(\mathrm{RFn}((f / g)) \to c \in \mathbb{R} \to \alpha \in \mathbb{R} \to \mathrm{PNbhd}((f / g),c) \to \forall a\,(\mathrm{Seq}(a) \to \mathrm{Conv}(a,c) \to \forall n \in \mathbb{N}\,(a_{n} \in \mathrm{dom}((f / g)) \land a_{n} \neq c) \to \mathrm{Conv}((f / g) \circ a,\alpha)) \to \mathrm{FLim}((f / g),c,\alpha))$
\item $\varphi_{330} :\equiv \varphi_{190} \to n \in \mathbb{N} \to f \circ a \otimes g \circ a^{\ast}(n) = (f / g) \circ a(n)$
\item $\varphi_{331} :\equiv (\lvert \beta \rvert \in \mathbb{R} \to 2^{-1} \in \mathbb{R} \to \lvert \beta \rvert \cdot 2^{-1} \in \mathbb{R}) \to \varphi_{18} \to 2^{-1} \in \mathbb{R} \to \varphi_{40}$
\item $\varphi_{332} :\equiv \forall z\,(f \circ a \otimes g \circ a^{\ast}(n) = f \circ a(n) \cdot g \circ a^{\ast}(n) \to f \circ a(n) \cdot g \circ a^{\ast}(n) = z \to f \circ a \otimes g \circ a^{\ast}(n) = z)$
\item $\varphi_{333} :\equiv \forall \alpha\,(\mathrm{Seq}((f / g) \circ a) \to \varphi_{118} \to \mathrm{Conv}(f \circ a \otimes g \circ a^{\ast},\alpha) \to \mathrm{Conv}((f / g) \circ a,\alpha))$
\item $\varphi_{334} :\equiv \mathrm{Seq}(f \circ a) \to \mathrm{Seq}(g \circ a^{\ast}) \to \mathrm{Seq}(f \circ a \otimes g \circ a^{\ast}) \land \varphi_{185}$
\item $\varphi_{335} :\equiv \mathrm{FLim}(f,c,\alpha) \to \mathrm{FLim}(g,c,\beta) \to \beta \neq 0 \to \mathrm{Seq}(a) \to \varphi_{273}$
\item $\varphi_{336} :\equiv d \in \mathbb{R} \land 0 < d \land \forall x \in \mathbb{R}\,(\varphi_{46}) \to d \in \mathbb{R} \land 0 < d \land \forall x \in \mathbb{R}\,(\varphi_{46})$
\item $\varphi_{337} :\equiv \forall w\,(g \circ a^{\ast}(n) = w \to \varphi_{145} \to f \circ a(n) \cdot g \circ a^{\ast}(n) = f \circ a(n) \cdot w)$
\item $\varphi_{338} :\equiv d \in \mathbb{R} \land 0 < d \land \forall x \in \mathbb{R}\,(\varphi_{46}) \to r \in \mathbb{R} \land 0 < r \land r \leq d \land r \leq \delta \to \varphi_{57}$
\item $\varphi_{339} :\equiv \forall \alpha\,\forall a\,(\mathrm{FLim}(f,c,\alpha) \to \mathrm{Conv}(a,c) \to \varphi_{72}) \to \forall a\,(\mathrm{FLim}(f,c,\alpha) \to \mathrm{Conv}(a,c) \to \varphi_{72})$
\item $\varphi_{340} :\equiv \mathrm{FLim}(f,c,\alpha) \to \mathrm{FLim}(g,c,\beta) \to \beta \neq 0 \to \mathrm{Seq}(a) \to \varphi_{277}$
\item $\varphi_{341} :\equiv \mathrm{FLim}(f,c,\alpha) \to \mathrm{FLim}(g,c,\beta) \to \beta \neq 0 \to \mathrm{Seq}(a) \to \varphi_{283}$
\item $\varphi_{342} :\equiv \mathrm{FLim}(g,c,\beta) \to \beta \neq 0 \to \mathrm{Seq}(a) \to \mathrm{Conv}(a,c) \to \varphi_{289}$
\item $\varphi_{343} :\equiv d \in \mathbb{R} \land 0 < d \land \forall x \in \mathbb{R}\,(\varphi_{46}) \to r \in \mathbb{R} \land 0 < r \land r \leq d \land r \leq \delta \to \varphi_{65}$
\item $\varphi_{344} :\equiv \mathrm{FLim}(f,c,\alpha) \to \mathrm{FLim}(g,c,\beta) \to \beta \neq 0 \to \mathrm{Seq}(a) \to \varphi_{284}$
\item $\varphi_{345} :\equiv d \in \mathbb{R} \land 0 < d \land \forall x \in \mathbb{R}\,(\varphi_{46}) \to r \in \mathbb{R} \land 0 < r \land r \leq d \land r \leq \delta \to \varphi_{79}$
\item $\varphi_{346} :\equiv d \in \mathbb{R} \land 0 < d \land \forall x \in \mathbb{R}\,(\varphi_{46}) \to r \in \mathbb{R} \land 0 < r \land r \leq d \land r \leq \delta \to \varphi_{78}$
\item $\varphi_{347} :\equiv \mathrm{FLim}(f,c,\alpha) \to \mathrm{FLim}(g,c,\beta) \to \beta \neq 0 \to \mathrm{Seq}(a) \to \varphi_{291}$
\item $\varphi_{348} :\equiv \mathrm{FLim}(f,c,\alpha) \to \mathrm{FLim}(g,c,\beta) \to \beta \neq 0 \to \mathrm{Seq}(a) \to \varphi_{292}$
\item $\varphi_{349} :\equiv \varphi_{2} \to \mathrm{FLim}(f,c,\alpha) \to \mathrm{FLim}(g,c,\beta) \to \beta \neq 0 \to \varphi_{223}$
\item $\varphi_{350} :\equiv d \in \mathbb{R} \land 0 < d \land \forall x \in \mathbb{R}\,(\varphi_{46}) \to r \in \mathbb{R} \land 0 < r \land r \leq d \land r \leq \delta \to \varphi_{84}$
\item $\varphi_{351} :\equiv \varphi_{181} \to \forall \alpha\,\forall a\,(\mathrm{FLim}(f,c,\alpha) \to \mathrm{Conv}(a,c) \to \varphi_{72})$
\item $\varphi_{352} :\equiv \delta \in \mathbb{R} \land 0 < \delta \land \forall x \in \mathrm{dom}(g)\,(\varphi_{94}) \to x \in \mathrm{dom}(g) \to \varphi_{94}$
\item $\varphi_{353} :\equiv \forall n \in \mathbb{N}\,(a_{n} \in \mathrm{dom}((f / g)) \land a_{n} \neq c) \to \mathrm{Seq}(f \circ a \otimes g \circ a^{\ast}) \land \varphi_{185}$
\item $\varphi_{354} :\equiv d \in \mathbb{R} \land 0 < d \land \forall x \in \mathbb{R}\,(\varphi_{46}) \to r \in \mathbb{R} \land 0 < r \land r \leq d \land r \leq \delta \to x \in \mathbb{R} \to \varphi_{45}$
\item $\varphi_{355} :\equiv d \in \mathbb{R} \land 0 < d \land \forall x \in \mathbb{R}\,(\varphi_{46}) \to r \in \mathbb{R} \land 0 < r \land r \leq d \land r \leq \delta \to x \in \mathbb{R} \to \varphi_{47}$
\item $\varphi_{356} :\equiv \forall y\,((f / g) \circ a(n) = y \to y = (f / g) \circ a(n)) \to \varphi_{259}$
\item $\varphi_{357} :\equiv \mathrm{FLim}(f,c,\alpha) \to \mathrm{FLim}(g,c,\beta) \to \beta \neq 0 \to \mathrm{Seq}(a) \to \varphi_{296}$
\item $\varphi_{358} :\equiv \mathrm{FLim}(f,c,\alpha) \to \mathrm{FLim}(g,c,\beta) \to \beta \neq 0 \to \mathrm{Seq}(a) \to \varphi_{297}$
\item $\varphi_{359} :\equiv \delta \in \mathbb{R} \land 0 < \delta \land \forall x \in \mathrm{dom}(g)\,(\varphi_{94}) \to \forall x \in \mathrm{dom}(g)\,(\varphi_{94})$
\item $\varphi_{360} :\equiv d \in \mathbb{R} \land 0 < d \land \forall x \in \mathbb{R}\,(\varphi_{46}) \to r \in \mathbb{R} \land 0 < r \land r \leq d \land r \leq \delta \to x \in \mathbb{R} \to \varphi_{55}$
\item $\varphi_{361} :\equiv d \in \mathbb{R} \land 0 < d \land \forall x \in \mathbb{R}\,(\varphi_{46}) \to r \in \mathbb{R} \land 0 < r \land r \leq d \land r \leq \delta \to x \in \mathbb{R} \to \varphi_{56}$
\item $\varphi_{362} :\equiv \mathrm{FLim}(g,c,\beta) \to \beta \neq 0 \to \mathrm{Seq}(a) \to \mathrm{Conv}(a,c) \to \varphi_{304}$
\item $\varphi_{363} :\equiv \mathrm{FLim}(f,c,\alpha) \to \mathrm{FLim}(g,c,\beta) \to \beta \neq 0 \to \mathrm{Seq}(a) \to \varphi_{301}$
\item $\varphi_{364} :\equiv d \in \mathbb{R} \land 0 < d \land \forall x \in \mathbb{R}\,(\varphi_{46}) \to r \in \mathbb{R} \land 0 < r \land r \leq d \land r \leq \delta \to x \in \mathbb{R} \to \varphi_{64}$
\item $\varphi_{365} :\equiv d \in \mathbb{R} \land 0 < d \land \forall x \in \mathbb{R}\,(\varphi_{46}) \to r \in \mathbb{R} \land 0 < r \land r \leq d \land r \leq \delta \to x \in \mathbb{R} \to \varphi_{73}$
\item $\varphi_{366} :\equiv d \in \mathbb{R} \land 0 < d \land \forall x \in \mathbb{R}\,(\varphi_{46}) \to r \in \mathbb{R} \land 0 < r \land r \leq d \land r \leq \delta \to \varphi_{108}$
\item $\varphi_{367} :\equiv \mathrm{FLim}(f,c,\alpha) \to \mathrm{FLim}(g,c,\beta) \to \beta \neq 0 \to \mathrm{Seq}(a) \to \varphi_{305}$
\item $\varphi_{368} :\equiv \mathrm{FLim}(g,c,\beta) \to \beta \neq 0 \to \mathrm{Seq}(a) \to \mathrm{Conv}(a,c) \to \varphi_{308}$
\item $\varphi_{369} :\equiv d \in \mathbb{R} \land 0 < d \land \forall x \in \mathbb{R}\,(\varphi_{46}) \to r \in \mathbb{R} \land 0 < r \land r \leq d \land r \leq \delta \to x \in \mathbb{R} \to \varphi_{75}$
\item $\varphi_{370} :\equiv \forall \beta\,(g(x) \in \mathbb{R} \to \beta \in \mathbb{R} \to \varphi_{87}) \to g(x) \in \mathbb{R} \to \beta \in \mathbb{R} \to \varphi_{87}$
\item $\varphi_{371} :\equiv \delta \in \mathbb{R} \land 0 < \delta \land \forall x \in \mathrm{dom}(g)\,(\varphi_{94}) \to \exists d\,(d \in \mathbb{R} \land 0 < d \land \forall x \in \mathbb{R}\,(\varphi_{46}))$
\item $\varphi_{372} :\equiv d \in \mathbb{R} \land 0 < d \land \forall x \in \mathbb{R}\,(\varphi_{46}) \to r \in \mathbb{R} \land 0 < r \land r \leq d \land r \leq \delta \to x \in \mathbb{R} \to \varphi_{91}$
\item $\varphi_{373} :\equiv \mathrm{RFn}((f / g)) \to c \in \mathbb{R} \to \alpha \cdot \beta^{-1} \in \mathbb{R} \to \mathrm{PNbhd}((f / g),c) \to \varphi_{266}$
\item $\varphi_{374} :\equiv \delta \in \mathbb{R} \land 0 < \delta \land \forall x \in \mathrm{dom}(g)\,(\varphi_{94}) \to d \in \mathbb{R} \land 0 < d \land \forall x \in \mathbb{R}\,(\varphi_{46}) \to 0 < d$
\item $\varphi_{375} :\equiv d \in \mathbb{R} \land 0 < d \land \forall x \in \mathbb{R}\,(\varphi_{46}) \to r \in \mathbb{R} \land 0 < r \land r \leq d \land r \leq \delta \to r \in \mathbb{R} \land 0 < r \land \varphi_{78}$
\item $\varphi_{376} :\equiv v_{767} \in \mathbb{R} \to \varphi_{43} \leftrightarrow v_{767} \in \mathbb{R} \to 0 < v_{766} \to 0 < v_{767} \to \neg \forall u_{4}\,(\varphi_{85})$
\item $\varphi_{377} :\equiv (\forall x\,(\varphi_{42}) \to \varphi_{42}) \to \mathrm{FLim}(f,c,\alpha) \to \mathrm{FLim}(g,c,\beta) \to \beta \neq 0 \to \varphi_{42}$
\item $\varphi_{378} :\equiv (\mathrm{FLim}(f,c,\alpha) \land \mathrm{FLim}(g,c,\beta) \land \beta \neq 0 \to \varphi_{1}) \to \mathrm{FLim}(f,c,\alpha) \to \varphi_{242}$
\item $\varphi_{379} :\equiv \delta \in \mathbb{R} \land 0 < \delta \land \forall x \in \mathrm{dom}(g)\,(\varphi_{94}) \to d \in \mathbb{R} \land 0 < d \land \forall x \in \mathbb{R}\,(\varphi_{46}) \to d \in \mathbb{R}$
\item $\varphi_{380} :\equiv (g \circ a(n) = g(a_{n}) \to g(a_{n}) = g \circ a(n)) \to \mathrm{FLim}(f,c,\alpha) \to \varphi_{252}$
\item $\varphi_{381} :\equiv (g(a_{n}) = g \circ a(n) \to g \circ a(n) = g(a_{n})) \to \mathrm{FLim}(f,c,\alpha) \to \varphi_{251}$
\item $\varphi_{382} :\equiv \delta \in \mathbb{R} \land 0 < \delta \land \forall x \in \mathrm{dom}(g)\,(\varphi_{94}) \to \delta \in \mathbb{R} \land 0 < \delta \land \forall x \in \mathrm{dom}(g)\,(\varphi_{94})$
\item $\varphi_{383} :\equiv \varphi_{204} \to (\mathrm{PNbhd}((f / g),c) \leftrightarrow \exists r\,(r \in \mathbb{R} \land 0 < r \land \varphi_{78}))$
\item $\varphi_{384} :\equiv \forall a\,(\mathrm{RFn}(g) \to \mathrm{Seq}(a) \to \varphi_{175}) \to \mathrm{RFn}(g) \to \mathrm{Seq}(a) \to \varphi_{175}$
\item $\varphi_{385} :\equiv \forall a\,(\mathrm{RFn}(f) \to \mathrm{Seq}(a) \to \varphi_{178}) \to \mathrm{RFn}(f) \to \mathrm{Seq}(a) \to \varphi_{178}$
\item $\varphi_{386} :\equiv \forall f\,\forall a\,(\mathrm{RFn}(f) \to \mathrm{Seq}(a) \to \varphi_{178}) \to \forall a\,(\mathrm{RFn}(g) \to \mathrm{Seq}(a) \to \varphi_{175})$
\item $\varphi_{387} :\equiv \forall f\,\forall a\,(\mathrm{RFn}(f) \to \mathrm{Seq}(a) \to \varphi_{178}) \to \forall a\,(\mathrm{RFn}(f) \to \mathrm{Seq}(a) \to \varphi_{178})$
\item $\varphi_{388} :\equiv \mathrm{FLim}(g,c,\beta) \to \beta \neq 0 \to \mathrm{Seq}(a) \to \mathrm{Conv}(a,c) \to \varphi_{353}$
\item $\varphi_{389} :\equiv \mathrm{FLim}(g,c,\beta) \to \beta \neq 0 \to \delta \in \mathbb{R} \land 0 < \delta \land \forall x \in \mathrm{dom}(g)\,(\varphi_{94}) \to \varphi_{209}$
\item $\varphi_{390} :\equiv \forall g\,\forall c\,(\mathrm{PNbhd}(f,c) \to \varphi_{196}) \to \forall c\,(\mathrm{PNbhd}(f,c) \to \varphi_{196})$
\item $\varphi_{391} :\equiv (\mathrm{FLim}(f,c,\alpha) \to \mathrm{FLim}(g,c,\beta) \to \beta \neq 0 \to \varphi_{122}) \to \mathrm{FLim}(f,c,\alpha) \to \varphi_{218}$
\item $\varphi_{392} :\equiv \forall f\,\forall g\,\forall c\,(\mathrm{PNbhd}(f,c) \to \varphi_{196}) \to \forall g\,\forall c\,(\mathrm{PNbhd}(f,c) \to \varphi_{196})$
\item $\varphi_{393} :\equiv f \circ a \otimes g \circ a^{\ast}(n) = f \circ a(n) \cdot g \circ a(n)^{-1} \to \varphi_{309}$
\item $\varphi_{394} :\equiv \delta \in \mathbb{R} \land 0 < \delta \land \forall x \in \mathrm{dom}(g)\,(\varphi_{94}) \to d \in \mathbb{R} \land 0 < d \land \forall x \in \mathbb{R}\,(\varphi_{46}) \to \varphi_{15}$
\item $\varphi_{395} :\equiv \mathrm{FLim}(f,c,\alpha) \land \mathrm{FLim}(g,c,\beta) \land \beta \neq 0 \land (\delta \in \mathbb{R} \land 0 < \delta \land \forall x \in \mathrm{dom}(g)\,(\varphi_{94})) \land (d \in \mathbb{R} \land 0 < d \land \forall x \in \mathbb{R}\,(\varphi_{46})) \land \neg \mathrm{PNbhd}((f / g),c)$
\item $\varphi_{396} :\equiv \forall v_{766}\,(\forall v_{767}\,(\varphi_{230}) \to \forall v_{767}\,(v_{766} \in \mathbb{R} \to v_{767} \in \mathbb{R} \to \varphi_{43}))$
\item $\varphi_{397} :\equiv \forall v_{766}\,(\forall v_{767}\,(v_{766} \in \mathbb{R} \to v_{767} \in \mathbb{R} \to \varphi_{43}) \to \forall v_{767}\,(\varphi_{230}))$
\item $\varphi_{398} :\equiv \varphi_{24} \to (\mathrm{FLim}(f,c,\alpha) \to \varphi_{242}) \to \mathrm{FLim}(f,c,\alpha) \to \varphi_{20}$
\item $\varphi_{399} :\equiv \mathrm{FLim}(g,c,\beta) \to \beta \neq 0 \to \delta \in \mathbb{R} \land 0 < \delta \land \forall x \in \mathrm{dom}(g)\,(\varphi_{94}) \to \varphi_{243}$
\item $\varphi_{400} :\equiv (\forall x\,(\varphi_{42}) \to \varphi_{122}) \to \mathrm{FLim}(f,c,\alpha) \to \mathrm{FLim}(g,c,\beta) \to \beta \neq 0 \to \varphi_{122}$
\item $\varphi_{401} :\equiv f \circ a \otimes g \circ a^{\ast}(n) = f \circ a(n) \cdot g \circ a^{\ast}(n) \to \varphi_{319}$
\item $\varphi_{402} :\equiv \delta \in \mathbb{R} \land 0 < \delta \land \forall x \in \mathrm{dom}(g)\,(\varphi_{94}) \to d \in \mathbb{R} \land 0 < d \land \forall x \in \mathbb{R}\,(\varphi_{46}) \to \varphi_{53}$
\item $\varphi_{403} :\equiv (\varphi_{54} \to \varphi_{127}) \to \mathrm{FLim}(f,c,\alpha) \to \mathrm{FLim}(g,c,\beta) \to \beta \neq 0 \to \varphi_{127}$
\item $\varphi_{404} :\equiv \mathrm{FLim}(f,c,\alpha) \land \mathrm{FLim}(g,c,\beta) \land \beta \neq 0 \land (\delta \in \mathbb{R} \land 0 < \delta \land \forall x \in \mathrm{dom}(g)\,(\varphi_{94})) \land (d \in \mathbb{R} \land 0 < d \land \forall x \in \mathbb{R}\,(\varphi_{46})) \land (r \in \mathbb{R} \land 0 < r \land r \leq d \land r \leq \delta)$
\item $\varphi_{405} :\equiv (\mathrm{FLim}(f,c,\alpha) \to \varphi_{217}) \to \varphi_{190} \to n \in \mathbb{N} \to a_{n} \in \mathrm{dom}(g)$
\item $\varphi_{406} :\equiv (\mathrm{FLim}(f,c,\alpha) \to \varphi_{218}) \to \varphi_{190} \to n \in \mathbb{N} \to a_{n} \in \mathrm{dom}(f)$
\item $\varphi_{407} :\equiv (\mathrm{FLim}(f,c,\alpha) \to \varphi_{299}) \to \mathrm{FLim}(f,c,\alpha) \to \mathrm{FLim}(g,c,\beta) \to \beta \neq 0 \to \varphi_{54}$
\item $\varphi_{408} :\equiv (\mathrm{FLim}(f,c,\alpha) \to \varphi_{227}) \to \varphi_{190} \to n \in \mathbb{N} \to g \circ a(n) \neq 0$
\item $\varphi_{409} :\equiv \forall n\,(\varphi_{190} \to n \in \mathbb{N} \to a_{n} \in \mathrm{dom}(g)) \to \forall n\,(\varphi_{190}) \to \forall n \in \mathbb{N}\,(a_{n} \in \mathrm{dom}(g))$
\item $\varphi_{410} :\equiv \forall n\,(\varphi_{190} \to n \in \mathbb{N} \to a_{n} \in \mathrm{dom}(f)) \to \forall n\,(\varphi_{190}) \to \forall n \in \mathbb{N}\,(a_{n} \in \mathrm{dom}(f))$
\item $\varphi_{411} :\equiv \mathrm{FLim}(g,c,\beta) \to \beta \neq 0 \to \delta \in \mathbb{R} \land 0 < \delta \land \forall x \in \mathrm{dom}(g)\,(\varphi_{94}) \to \varphi_{278}$
\item $\varphi_{412} :\equiv \delta \in \mathbb{R} \land 0 < \delta \land \forall x \in \mathrm{dom}(g)\,(\varphi_{94}) \to d \in \mathbb{R} \land 0 < d \land \forall x \in \mathbb{R}\,(\varphi_{46}) \to \varphi_{104}$
\item $\varphi_{413} :\equiv \forall n\,(\varphi_{190} \to n \in \mathbb{N} \to g \circ a(n) \neq 0) \to \forall n\,(\varphi_{190}) \to \forall n \in \mathbb{N}\,(g \circ a(n) \neq 0)$
\item $\varphi_{414} :\equiv (\mathrm{FLim}(f,c,\alpha) \to \varphi_{236}) \to \varphi_{190} \to n \in \mathbb{N} \to a_{n} \in \mathrm{dom}((f / g))$
\item $\varphi_{415} :\equiv (\forall d\,(\varphi_{62}) \leftrightarrow \forall u_{4}\,(\varphi_{85})) \to (\neg \forall d\,(\varphi_{62}) \leftrightarrow \neg \forall u_{4}\,(\varphi_{85}))$
\item $\varphi_{416} :\equiv \varphi_{270} \to \mathrm{Conv}(f \circ a,\alpha) \to \mathrm{Conv}(g \circ a,\beta) \to \beta \neq 0 \to \varphi_{153}$
\item $\varphi_{417} :\equiv \delta \in \mathbb{R} \land 0 < \delta \land \forall x \in \mathrm{dom}(g)\,(\varphi_{94}) \to d \in \mathbb{R} \land 0 < d \land \forall x \in \mathbb{R}\,(\varphi_{46}) \to \varphi_{114}$
\item $\varphi_{418} :\equiv \delta \in \mathbb{R} \land 0 < \delta \land \forall x \in \mathrm{dom}(g)\,(\varphi_{94}) \to d \in \mathbb{R} \land 0 < d \land \forall x \in \mathbb{R}\,(\varphi_{46}) \to \varphi_{120}$
\item $\varphi_{419} :\equiv (\mathrm{Conv}(g \circ a,\beta) \to \mathrm{Seq}(g \circ a)) \to (\mathrm{FLim}(f,c,\alpha) \to \varphi_{207}) \to \mathrm{FLim}(f,c,\alpha) \to \varphi_{197}$
\item $\varphi_{420} :\equiv (\mathrm{Conv}(f \circ a,\alpha) \to \mathrm{Seq}(f \circ a)) \to (\mathrm{FLim}(f,c,\alpha) \to \varphi_{208}) \to \mathrm{FLim}(f,c,\alpha) \to \varphi_{198}$
\item $\varphi_{421} :\equiv \forall n\,(\varphi_{190} \to n \in \mathbb{N} \to a_{n} \in \mathrm{dom}((f / g))) \to \forall n\,(\varphi_{190}) \to \forall n \in \mathbb{N}\,(a_{n} \in \mathrm{dom}((f / g)))$
\item $\varphi_{422} :\equiv \delta \in \mathbb{R} \land 0 < \delta \land \forall x \in \mathrm{dom}(g)\,(\varphi_{94}) \to d \in \mathbb{R} \land 0 < d \land \forall x \in \mathbb{R}\,(\varphi_{46}) \to \varphi_{140}$
\item $\varphi_{423} :\equiv \delta \in \mathbb{R} \land 0 < \delta \land \forall x \in \mathrm{dom}(g)\,(\varphi_{94}) \to d \in \mathbb{R} \land 0 < d \land \forall x \in \mathbb{R}\,(\varphi_{46}) \to \varphi_{144}$
\item $\varphi_{424} :\equiv (\mathrm{FLim}(f,c,\alpha) \to \varphi_{245}) \to \varphi_{190} \to n \in \mathbb{N} \to a_{n} \in \mathrm{dom}(g) \land a_{n} \neq c$
\item $\varphi_{425} :\equiv (\mathrm{FLim}(f,c,\alpha) \to \varphi_{246}) \to \varphi_{190} \to n \in \mathbb{N} \to a_{n} \in \mathrm{dom}(f) \land a_{n} \neq c$
\item $\varphi_{426} :\equiv \delta \in \mathbb{R} \land 0 < \delta \land \forall x \in \mathrm{dom}(g)\,(\varphi_{94}) \to d \in \mathbb{R} \land 0 < d \land \forall x \in \mathbb{R}\,(\varphi_{46}) \to \varphi_{156}$
\item $\varphi_{427} :\equiv \delta \in \mathbb{R} \land 0 < \delta \land \forall x \in \mathrm{dom}(g)\,(\varphi_{94}) \to d \in \mathbb{R} \land 0 < d \land \forall x \in \mathbb{R}\,(\varphi_{46}) \to \varphi_{157}$
\item $\varphi_{428} :\equiv \delta \in \mathbb{R} \land 0 < \delta \land \forall x \in \mathrm{dom}(g)\,(\varphi_{94}) \to d \in \mathbb{R} \land 0 < d \land \forall x \in \mathbb{R}\,(\varphi_{46}) \to \varphi_{161}$
\item $\varphi_{429} :\equiv \delta \in \mathbb{R} \land 0 < \delta \land \forall x \in \mathrm{dom}(g)\,(\varphi_{94}) \to d \in \mathbb{R} \land 0 < d \land \forall x \in \mathbb{R}\,(\varphi_{46}) \to \varphi_{162}$
\item $\varphi_{430} :\equiv \forall n\,(\varphi_{190} \to n \in \mathbb{N} \to a_{n} \in \mathrm{dom}(g) \land a_{n} \neq c) \to \forall n\,(\varphi_{190}) \to \forall n \in \mathbb{N}\,(a_{n} \in \mathrm{dom}(g) \land a_{n} \neq c)$
\item $\varphi_{431} :\equiv \forall n\,(\varphi_{190} \to n \in \mathbb{N} \to a_{n} \in \mathrm{dom}(f) \land a_{n} \neq c) \to \forall n\,(\varphi_{190}) \to \forall n \in \mathbb{N}\,(a_{n} \in \mathrm{dom}(f) \land a_{n} \neq c)$
\item $\varphi_{432} :\equiv (\mathrm{FLim}(f,c,\alpha) \to \mathrm{FLim}(g,c,\beta) \to \beta \neq 0 \to \varphi_{122}) \to \mathrm{FLim}(f,c,\alpha) \to \varphi_{316}$
\item $\varphi_{433} :\equiv (\neg \forall d\,(\varphi_{62}) \leftrightarrow \neg \forall u_{4}\,(\varphi_{85})) \to (0 < v_{767} \to \neg \forall d\,(\varphi_{62}) \leftrightarrow 0 < v_{767} \to \neg \forall u_{4}\,(\varphi_{85}))$
\item $\varphi_{434} :\equiv \mathrm{FLim}(g,c,\beta) \to \beta \neq 0 \to \delta \in \mathbb{R} \land 0 < \delta \land \forall x \in \mathrm{dom}(g)\,(\varphi_{94}) \to \varphi_{336}$
\item $\varphi_{435} :\equiv \mathrm{FLim}(g,c,\beta) \to \beta \neq 0 \to \delta \in \mathbb{R} \land 0 < \delta \land \forall x \in \mathrm{dom}(g)\,(\varphi_{94}) \to \varphi_{338}$
\item $\varphi_{436} :\equiv \mathrm{FLim}(g,c,\beta) \to \beta \neq 0 \to \delta \in \mathbb{R} \land 0 < \delta \land \forall x \in \mathrm{dom}(g)\,(\varphi_{94}) \to \varphi_{343}$
\item $\varphi_{437} :\equiv \mathrm{FLim}(g,c,\beta) \to \beta \neq 0 \to \delta \in \mathbb{R} \land 0 < \delta \land \forall x \in \mathrm{dom}(g)\,(\varphi_{94}) \to \varphi_{345}$
\item $\varphi_{438} :\equiv \mathrm{FLim}(g,c,\beta) \to \beta \neq 0 \to \delta \in \mathbb{R} \land 0 < \delta \land \forall x \in \mathrm{dom}(g)\,(\varphi_{94}) \to \varphi_{346}$
\item $\varphi_{439} :\equiv \mathrm{FLim}(g,c,\beta) \to \beta \neq 0 \to \delta \in \mathbb{R} \land 0 < \delta \land \forall x \in \mathrm{dom}(g)\,(\varphi_{94}) \to \varphi_{350}$
\item $\varphi_{440} :\equiv \delta \in \mathbb{R} \land 0 < \delta \land \forall x \in \mathrm{dom}(g)\,(\varphi_{94}) \to d \in \mathbb{R} \land 0 < d \land \forall x \in \mathbb{R}\,(\varphi_{46}) \to \varphi_{222}$
\item $\varphi_{441} :\equiv (\mathrm{FLim}(f,c,\alpha) \to \mathrm{FLim}(g,c,\beta) \to \beta \neq 0 \to \varphi_{237}) \to \varphi_{28} \to \neg (\delta \in \mathbb{R} \land 0 < \delta \land \forall x \in \mathrm{dom}(g)\,(\varphi_{94}))$
\item $\varphi_{442} :\equiv \mathrm{FLim}(g,c,\beta) \to \beta \neq 0 \to \delta \in \mathbb{R} \land 0 < \delta \land \forall x \in \mathrm{dom}(g)\,(\varphi_{94}) \to \varphi_{354}$
\item $\varphi_{443} :\equiv \mathrm{FLim}(g,c,\beta) \to \beta \neq 0 \to \delta \in \mathbb{R} \land 0 < \delta \land \forall x \in \mathrm{dom}(g)\,(\varphi_{94}) \to \varphi_{355}$
\item $\varphi_{444} :\equiv \mathrm{FLim}(g,c,\beta) \to \beta \neq 0 \to \delta \in \mathbb{R} \land 0 < \delta \land \forall x \in \mathrm{dom}(g)\,(\varphi_{94}) \to \varphi_{360}$
\item $\varphi_{445} :\equiv \mathrm{FLim}(g,c,\beta) \to \beta \neq 0 \to \delta \in \mathbb{R} \land 0 < \delta \land \forall x \in \mathrm{dom}(g)\,(\varphi_{94}) \to \varphi_{361}$
\item $\varphi_{446} :\equiv \mathrm{FLim}(g,c,\beta) \to \beta \neq 0 \to \delta \in \mathbb{R} \land 0 < \delta \land \forall x \in \mathrm{dom}(g)\,(\varphi_{94}) \to \varphi_{364}$
\item $\varphi_{447} :\equiv \mathrm{FLim}(g,c,\beta) \to \beta \neq 0 \to \delta \in \mathbb{R} \land 0 < \delta \land \forall x \in \mathrm{dom}(g)\,(\varphi_{94}) \to \varphi_{365}$
\item $\varphi_{448} :\equiv \mathrm{FLim}(g,c,\beta) \to \beta \neq 0 \to \delta \in \mathbb{R} \land 0 < \delta \land \forall x \in \mathrm{dom}(g)\,(\varphi_{94}) \to \varphi_{366}$
\item $\varphi_{449} :\equiv \mathrm{FLim}(g,c,\beta) \to \beta \neq 0 \to \delta \in \mathbb{R} \land 0 < \delta \land \forall x \in \mathrm{dom}(g)\,(\varphi_{94}) \to \varphi_{369}$
\item $\varphi_{450} :\equiv (\mathrm{FLim}(f,c,\alpha) \to \varphi_{14}) \to (\mathrm{FLim}(f,c,\alpha) \to \varphi_{255}) \to \mathrm{FLim}(f,c,\alpha) \to \varphi_{207}$
\item $\varphi_{451} :\equiv (\mathrm{FLim}(f,c,\alpha) \to \varphi_{14}) \to (\mathrm{FLim}(f,c,\alpha) \to \varphi_{256}) \to \mathrm{FLim}(f,c,\alpha) \to \varphi_{208}$
\item $\varphi_{452} :\equiv \forall \delta\,(\varphi_{28} \to \neg (\delta \in \mathbb{R} \land 0 < \delta \land \forall x \in \mathrm{dom}(g)\,(\varphi_{94}))) \to \varphi_{315}$
\item $\varphi_{453} :\equiv \mathrm{FLim}(g,c,\beta) \to \beta \neq 0 \to \delta \in \mathbb{R} \land 0 < \delta \land \forall x \in \mathrm{dom}(g)\,(\varphi_{94}) \to \varphi_{372}$
\item $\varphi_{454} :\equiv \varphi_{3} \to (\mathrm{FLim}(f,c,\alpha) \to \varphi_{233}) \to \mathrm{FLim}(f,c,\alpha) \to \varphi_{276}$
\item $\varphi_{455} :\equiv \mathrm{FLim}(g,c,\beta) \to \beta \neq 0 \to \delta \in \mathbb{R} \land 0 < \delta \land \forall x \in \mathrm{dom}(g)\,(\varphi_{94}) \to \varphi_{375}$
\item $\varphi_{456} :\equiv (\forall n \in \mathbb{N}\,(\varphi_{81}) \to n \in \mathbb{N} \to \varphi_{81}) \to \mathrm{FLim}(f,c,\alpha) \to \varphi_{316}$
\item $\varphi_{457} :\equiv (0 < v_{767} \to \neg \forall d\,(\varphi_{62}) \leftrightarrow 0 < v_{767} \to \neg \forall u_{4}\,(\varphi_{85})) \to (\varphi_{43} \leftrightarrow 0 < v_{766} \to 0 < v_{767} \to \neg \forall u_{4}\,(\varphi_{85}))$
\item $\varphi_{458} :\equiv \delta \in \mathbb{R} \land 0 < \delta \land \forall x \in \mathrm{dom}(g)\,(\varphi_{94}) \to d \in \mathbb{R} \land 0 < d \land \forall x \in \mathbb{R}\,(\varphi_{46}) \to \varphi_{279}$
\item $\varphi_{459} :\equiv \delta \in \mathbb{R} \land 0 < \delta \land \forall x \in \mathrm{dom}(g)\,(\varphi_{94}) \to d \in \mathbb{R} \land 0 < d \land \forall x \in \mathbb{R}\,(\varphi_{46}) \to \varphi_{288}$
\item $\varphi_{460} :\equiv (g \circ a(n) = g(a_{n}) \to g \circ a(n) = 0 \to g(a_{n}) = 0) \to (\mathrm{FLim}(f,c,\alpha) \to \varphi_{251}) \to \mathrm{FLim}(f,c,\alpha) \to \varphi_{227}$
\item $\varphi_{461} :\equiv (\mathrm{FLim}(f,c,\alpha) \to \mathrm{FLim}(g,c,\beta) \to \beta \neq 0 \to \varphi_{54}) \to \varphi_{403}$
\item $\varphi_{462} :\equiv (\forall v_{766}\,\forall v_{767}\,(v_{766} \in \mathbb{R} \to v_{767} \in \mathbb{R} \to \varphi_{43}) \leftrightarrow \varphi_{187}) \to \varphi_{187}$
\item $\varphi_{463} :\equiv ((f / g) \circ a(n) = (f / g)(a_{n}) \to \varphi_{248}) \to \mathrm{FLim}(f,c,\alpha) \to \varphi_{313}$
\item $\varphi_{464} :\equiv (\lvert g(x) \rvert = 0 \to 0 = \lvert g(x) \rvert) \to \mathrm{FLim}(f,c,\alpha) \to \varphi_{445}$
\item $\varphi_{465} :\equiv (0 = \lvert g(x) \rvert \to \lvert g(x) \rvert = 0) \to \mathrm{FLim}(f,c,\alpha) \to \varphi_{444}$
\item $\varphi_{466} :\equiv ((f / g)(a_{n}) = f(a_{n}) \cdot g(a_{n})^{-1} \to \varphi_{272}) \to \mathrm{FLim}(f,c,\alpha) \to \varphi_{311}$
\item $\varphi_{467} :\equiv (\lvert g(x) \rvert = \lvert 0 \rvert \to \lvert 0 \rvert = 0 \to \lvert g(x) \rvert = 0) \to \mathrm{FLim}(f,c,\alpha) \to \varphi_{444}$
\item $\varphi_{468} :\equiv (\mathrm{FLim}(f,c,\alpha) \to \mathrm{FLim}(g,c,\beta) \to \beta \neq 0 \to \varphi_{223}) \to \mathrm{FLim}(f,c,\alpha) \to \mathrm{FLim}(g,c,\beta) \to \beta \neq 0 \to \varphi_{371}$
\item $\varphi_{469} :\equiv (\forall x \in \mathrm{dom}(g)\,(\varphi_{94}) \to x \in \mathrm{dom}(g) \to \varphi_{94}) \to \mathrm{FLim}(f,c,\alpha) \to \mathrm{FLim}(g,c,\beta) \to \beta \neq 0 \to \varphi_{352}$
\item $\varphi_{470} :\equiv (\mathrm{FLim}(f,c,\alpha) \to \mathrm{FLim}(g,c,\beta) \to \beta \neq 0 \to \varphi_{382}) \to \mathrm{FLim}(f,c,\alpha) \to \varphi_{262}$
\item $\varphi_{471} :\equiv (\mathrm{FLim}(f,c,\alpha) \to \mathrm{FLim}(g,c,\beta) \to \beta \neq 0 \to \varphi_{382}) \to \mathrm{FLim}(f,c,\alpha) \to \varphi_{268}$
\item $\varphi_{472} :\equiv (\mathrm{FLim}(f,c,\alpha) \to \varphi_{269}) \to \varphi_{184} \to \mathrm{FLim}(f,c,\alpha) \to \varphi_{263}$
\item $\varphi_{473} :\equiv (\mathrm{FLim}(f,c,\alpha) \to \varphi_{281}) \to \varphi_{149} \to \mathrm{FLim}(f,c,\alpha) \to \varphi_{274}$
\item $\varphi_{474} :\equiv (\mathrm{FLim}(f,c,\alpha) \to \varphi_{282}) \to \varphi_{150} \to \mathrm{FLim}(f,c,\alpha) \to \varphi_{275}$
\item $\varphi_{475} :\equiv \varphi_{317} \to f \circ a \otimes g \circ a^{\ast}(n) = f(a_{n}) \cdot g \circ a(n)^{-1} \to \varphi_{294}$
\item $\varphi_{476} :\equiv (\mathrm{FLim}(f,c,\alpha) \to \varphi_{263}) \to (\mathrm{FLim}(f,c,\alpha) \to \varphi_{206}) \to \mathrm{FLim}(f,c,\alpha) \to \varphi_{236}$
\item $\varphi_{477} :\equiv (\forall d\,(\varphi_{62}) \to \forall u_{4}\,(\varphi_{85})) \to (\forall u_{4}\,(\varphi_{85}) \to \forall d\,(\varphi_{62})) \to (\forall d\,(\varphi_{62}) \leftrightarrow \forall u_{4}\,(\varphi_{85}))$
\item $\varphi_{478} :\equiv (\mathrm{FLim}(f,c,\alpha) \to \mathrm{FLim}(g,c,\beta) \to \beta \neq 0 \to \varphi_{42}) \to \mathrm{FLim}(f,c,\alpha) \to \varphi_{436}$
\item $\varphi_{479} :\equiv (v_{767} \in \mathbb{R} \leftrightarrow v_{767} \in \mathbb{R}) \to (\varphi_{43} \leftrightarrow 0 < v_{766} \to 0 < v_{767} \to \neg \forall u_{4}\,(\varphi_{85})) \to \varphi_{376}$
\item $\varphi_{480} :\equiv (g \circ a^{\ast}(n) = g \circ a(n)^{-1} \to \varphi_{145} \to \varphi_{139}) \to \mathrm{FLim}(f,c,\alpha) \to \varphi_{342}$
\item $\varphi_{481} :\equiv (g(a_{n}) = g \circ a(n) \to g(a_{n}) = 0 \to g \circ a(n) = 0) \to \varphi_{460}$
\item $\varphi_{482} :\equiv (\mathrm{FLim}(f,c,\alpha) \to \varphi_{274}) \to (\mathrm{FLim}(f,c,\alpha) \to \varphi_{206}) \to \mathrm{FLim}(f,c,\alpha) \to \varphi_{251}$
\item $\varphi_{483} :\equiv (\mathrm{FLim}(f,c,\alpha) \to \varphi_{275}) \to (\mathrm{FLim}(f,c,\alpha) \to \varphi_{206}) \to \mathrm{FLim}(f,c,\alpha) \to \varphi_{253}$
\item $\varphi_{484} :\equiv (f \circ a \otimes g \circ a^{\ast}(n) = f(a_{n}) \cdot g \circ a(n)^{-1} \to \varphi_{294}) \to \varphi_{347}$
\item $\varphi_{485} :\equiv (\mathrm{FLim}(f,c,\alpha) \to \varphi_{316}) \to (\mathrm{FLim}(f,c,\alpha) \to \varphi_{206}) \to \mathrm{FLim}(f,c,\alpha) \to \varphi_{216}$
\item $\varphi_{486} :\equiv (\mathrm{FLim}(f,c,\alpha) \to \varphi_{316}) \to (\mathrm{FLim}(f,c,\alpha) \to \varphi_{206}) \to \mathrm{FLim}(f,c,\alpha) \to \varphi_{217}$
\item $\varphi_{487} :\equiv (\mathrm{FLim}(f,c,\alpha) \to \varphi_{197}) \to (\mathrm{FLim}(f,c,\alpha) \to \varphi_{233}) \to \varphi_{340}$
\item $\varphi_{488} :\equiv (\mathrm{FLim}(f,c,\alpha) \to \varphi_{263}) \to (\mathrm{FLim}(f,c,\alpha) \to \varphi_{206}) \to \mathrm{FLim}(f,c,\alpha) \to \varphi_{302}$
\item $\varphi_{489} :\equiv (\mathrm{FLim}(f,c,\alpha) \to \mathrm{Conv}(a,c) \to \varphi_{72}) \to (\mathrm{FLim}(f,c,\alpha) \to \mathrm{FLim}(f,c,\alpha)) \to \varphi_{451}$
\item $\varphi_{490} :\equiv (\mathrm{FLim}(f,c,\alpha) \to \varphi_{316}) \to (\mathrm{FLim}(f,c,\alpha) \to \varphi_{206}) \to \mathrm{FLim}(f,c,\alpha) \to \varphi_{245}$
\item $\varphi_{491} :\equiv (\mathrm{FLim}(f,c,\alpha) \to \varphi_{316}) \to (\mathrm{FLim}(f,c,\alpha) \to \varphi_{206}) \to \mathrm{FLim}(f,c,\alpha) \to \varphi_{246}$
\item $\varphi_{492} :\equiv f(a_{n}) \cdot g \circ a(n)^{-1} = f(a_{n}) \cdot g \circ a(n)^{-1} \to (g \circ a(n) = g(a_{n}) \to \varphi_{295}) \to \varphi_{344}$
\item $\varphi_{493} :\equiv (\mathrm{FLim}(f,c,\alpha) \to \varphi_{300}) \to (\mathrm{FLim}(f,c,\alpha) \to \varphi_{206}) \to \mathrm{FLim}(f,c,\alpha) \to \varphi_{285}$
\item $\varphi_{494} :\equiv (\mathrm{FLim}(f,c,\alpha) \to \varphi_{303}) \to (\mathrm{FLim}(f,c,\alpha) \to \varphi_{206}) \to \mathrm{FLim}(f,c,\alpha) \to \varphi_{286}$
\item $\varphi_{495} :\equiv (\mathrm{FLim}(f,c,\alpha) \to \varphi_{306}) \to \varphi_{214} \to \mathrm{FLim}(f,c,\alpha) \to \varphi_{300}$
\item $\varphi_{496} :\equiv (\mathrm{FLim}(g,c,\beta) \to \lvert \beta \rvert \cdot 2^{-1} \in \mathbb{R} \to \varphi_{254}) \to \varphi_{2} \to \varphi_{40} \to \varphi_{33} \to \mathrm{FLim}(f,c,\alpha) \to \varphi_{247}$
\item $\varphi_{497} :\equiv (v_{766} \in \mathbb{R} \leftrightarrow v_{766} \in \mathbb{R}) \to \varphi_{376} \to (v_{766} \in \mathbb{R} \to v_{767} \in \mathbb{R} \to \varphi_{43} \leftrightarrow \varphi_{230})$
\item $\varphi_{498} :\equiv (v_{766} \in \mathbb{R} \to v_{767} \in \mathbb{R} \to \varphi_{43} \leftrightarrow \varphi_{230}) \to \varphi_{230} \to v_{766} \in \mathbb{R} \to v_{767} \in \mathbb{R} \to \varphi_{43}$
\item $\varphi_{499} :\equiv (v_{766} \in \mathbb{R} \to v_{767} \in \mathbb{R} \to \varphi_{43} \leftrightarrow \varphi_{230}) \to (v_{766} \in \mathbb{R} \to v_{767} \in \mathbb{R} \to \varphi_{43}) \to \varphi_{230}$
\item $\varphi_{500} :\equiv (\mathrm{FLim}(f,c,\alpha) \to \mathrm{FLim}(g,c,\beta) \to \beta \neq 0 \to \varphi_{382}) \to \mathrm{FLim}(f,c,\alpha) \to \mathrm{FLim}(g,c,\beta) \to \beta \neq 0 \to \varphi_{359}$
\item $\varphi_{501} :\equiv \varphi_{335} \to (\mathrm{FLim}(f,c,\alpha) \to \varphi_{276}) \to \mathrm{FLim}(f,c,\alpha) \to \varphi_{238}$
\item $\varphi_{502} :\equiv (\varphi_{307} \to \forall d\,(\varphi_{307})) \to \mathrm{FLim}(f,c,\alpha) \to \mathrm{FLim}(g,c,\beta) \to \beta \neq 0 \to \varphi_{237}$
\item $\varphi_{503} :\equiv (\mathrm{FLim}(f,c,\alpha) \to \varphi_{198}) \to (\mathrm{FLim}(f,c,\alpha) \to \varphi_{210}) \to \mathrm{FLim}(f,c,\alpha) \to \varphi_{388}$
\item $\varphi_{504} :\equiv \forall v_{767}\,(\varphi_{230} \to v_{766} \in \mathbb{R} \to v_{767} \in \mathbb{R} \to \varphi_{43}) \to \forall v_{767}\,(\varphi_{230}) \to \forall v_{767}\,(v_{766} \in \mathbb{R} \to v_{767} \in \mathbb{R} \to \varphi_{43})$
\item $\varphi_{505} :\equiv \forall v_{767}\,((v_{766} \in \mathbb{R} \to v_{767} \in \mathbb{R} \to \varphi_{43}) \to \varphi_{230}) \to \forall v_{767}\,(v_{766} \in \mathbb{R} \to v_{767} \in \mathbb{R} \to \varphi_{43}) \to \forall v_{767}\,(\varphi_{230})$
\item $\varphi_{506} :\equiv (\forall v_{767}\,(v_{766} \in \mathbb{R} \to v_{767} \in \mathbb{R} \to \varphi_{43}) \leftrightarrow \forall v_{767}\,(\varphi_{230})) \to \forall v_{767}\,(\varphi_{230}) \to \forall v_{767}\,(v_{766} \in \mathbb{R} \to v_{767} \in \mathbb{R} \to \varphi_{43})$
\item $\varphi_{507} :\equiv (\forall v_{767}\,(v_{766} \in \mathbb{R} \to v_{767} \in \mathbb{R} \to \varphi_{43}) \leftrightarrow \forall v_{767}\,(\varphi_{230})) \to \forall v_{767}\,(v_{766} \in \mathbb{R} \to v_{767} \in \mathbb{R} \to \varphi_{43}) \to \forall v_{767}\,(\varphi_{230})$
\item $\varphi_{508} :\equiv (\mathrm{FLim}(f,c,\alpha) \to \varphi_{313}) \to \varphi_{259} \to \mathrm{FLim}(f,c,\alpha) \to \varphi_{314}$
\item $\varphi_{509} :\equiv (\forall a\,(\varphi_{267}) \to \varphi_{1} \to \forall a\,(\mathrm{Seq}(a) \to \mathrm{Conv}(a,c) \to \varphi_{148})) \to \varphi_{378}$
\item $\varphi_{510} :\equiv (\mathrm{FLim}(f,c,\alpha) \to \varphi_{389}) \to \varphi_{307} \to \neg (d \in \mathbb{R} \land 0 < d \land \forall x \in \mathbb{R}\,(\varphi_{46}))$
\item $\varphi_{511} :\equiv (\mathrm{RFn}(g) \to \mathrm{Seq}(a) \to \varphi_{175}) \to \varphi_{4} \to \varphi_{13} \to (\mathrm{FLim}(f,c,\alpha) \to \varphi_{226}) \to \varphi_{321}$
\item $\varphi_{512} :\equiv (\mathrm{RFn}(f) \to \mathrm{Seq}(a) \to \varphi_{178}) \to \varphi_{5} \to \varphi_{13} \to (\mathrm{FLim}(f,c,\alpha) \to \varphi_{228}) \to \varphi_{322}$
\item $\varphi_{513} :\equiv \forall d\,(\varphi_{307} \to \neg (d \in \mathbb{R} \land 0 < d \land \forall x \in \mathbb{R}\,(\varphi_{46}))) \to \forall d\,(\varphi_{307}) \to \forall d\,(\neg (d \in \mathbb{R} \land 0 < d \land \forall x \in \mathbb{R}\,(\varphi_{46})))$
\item $\varphi_{514} :\equiv (\forall x \in \mathbb{R}\,(\varphi_{46}) \to x \in \mathbb{R} \to \varphi_{46}) \to \mathrm{FLim}(f,c,\alpha) \to \mathrm{FLim}(g,c,\beta) \to \beta \neq 0 \to \varphi_{458}$
\item $\varphi_{515} :\equiv (\mathrm{FLim}(f,c,\alpha) \to \varphi_{263}) \to (\mathrm{FLim}(f,c,\alpha) \to \varphi_{206}) \to \varphi_{391}$
\item $\varphi_{516} :\equiv \varphi_{323} \to \mathrm{FLim}(f,c,\alpha) \to \mathrm{FLim}(g,c,\beta) \to \beta \neq 0 \to \varphi_{440}$
\item $\varphi_{517} :\equiv (\mathrm{FLim}(f,c,\alpha) \to \varphi_{362}) \to (\mathrm{FLim}(f,c,\alpha) \to \varphi_{206}) \to \varphi_{367}$
\item $\varphi_{518} :\equiv (\mathrm{FLim}(f,c,\alpha) \to \mathrm{FLim}(g,c,\beta) \to \beta \neq 0 \to \varphi_{224}) \to \varphi_{468}$
\item $\varphi_{519} :\equiv \forall u_{4}\,(\varphi_{80}) \to (\forall u_{4}\,(\varphi_{80}) \to \varphi_{107}) \to \varphi_{92} \to \mathrm{FLim}(f,c,\alpha) \to \varphi_{389}$
\item $\varphi_{520} :\equiv (\forall u_{4}\,(\forall d\,(\varphi_{62}) \to \varphi_{85}) \to \forall u_{4}\,\forall d\,(\varphi_{62}) \to \forall u_{4}\,(\varphi_{85})) \to (\forall d\,(\varphi_{62}) \to \forall u_{4}\,\forall d\,(\varphi_{62})) \to \forall d\,(\varphi_{62}) \to \forall u_{4}\,(\varphi_{85})$
\item $\varphi_{521} :\equiv (\mathrm{FLim}(f,c,\alpha) \to \mathrm{FLim}(g,c,\beta) \to \beta \neq 0 \to \varphi_{127}) \to \varphi_{488}$
\item $\varphi_{522} :\equiv (\forall d\,(\forall u_{4}\,(\varphi_{85}) \to \varphi_{62}) \to \forall d\,\forall u_{4}\,(\varphi_{85}) \to \forall d\,(\varphi_{62})) \to (\forall u_{4}\,(\varphi_{85}) \to \forall d\,\forall u_{4}\,(\varphi_{85})) \to \forall u_{4}\,(\varphi_{85}) \to \forall d\,(\varphi_{62})$
\item $\varphi_{523} :\equiv (\mathrm{FLim}(f,c,\alpha) \to \mathrm{FLim}(g,c,\beta) \to \beta \neq 0 \to \varphi_{394}) \to \varphi_{395} \to \neg (r \in \mathbb{R} \land 0 < r \land r \leq d \land r \leq \delta)$
\item $\varphi_{524} :\equiv (\mathrm{FLim}(f,c,\alpha) \to \varphi_{434}) \to \mathrm{FLim}(f,c,\alpha) \to \mathrm{FLim}(g,c,\beta) \to \beta \neq 0 \to \varphi_{374}$
\item $\varphi_{525} :\equiv (\mathrm{FLim}(f,c,\alpha) \to \varphi_{434}) \to \mathrm{FLim}(f,c,\alpha) \to \mathrm{FLim}(g,c,\beta) \to \beta \neq 0 \to \varphi_{379}$
\item $\varphi_{526} :\equiv (\mathrm{FLim}(f,c,\alpha) \to \varphi_{263}) \to (\mathrm{FLim}(f,c,\alpha) \to \varphi_{206}) \to \varphi_{432}$
\item $\varphi_{527} :\equiv \forall r\,(\varphi_{395} \to \neg (r \in \mathbb{R} \land 0 < r \land r \leq d \land r \leq \delta)) \to \forall r\,(\varphi_{395}) \to \forall r\,(\neg (r \in \mathbb{R} \land 0 < r \land r \leq d \land r \leq \delta))$
\item $\varphi_{528} :\equiv (\mathrm{FLim}(f,c,\alpha) \to \varphi_{208}) \to (\mathrm{FLim}(f,c,\alpha) \to \varphi_{207}) \to \varphi_{454}$
\item $\varphi_{529} :\equiv (\mathrm{FLim}(f,c,\alpha) \to \mathrm{FLim}(g,c,\beta) \to \beta \neq 0 \to \varphi_{402}) \to \mathrm{FLim}(f,c,\alpha) \to \mathrm{FLim}(g,c,\beta) \to \beta \neq 0 \to \varphi_{417}$
\item $\varphi_{530} :\equiv (\mathrm{FLim}(f,c,\alpha) \to \mathrm{FLim}(g,c,\beta) \to \beta \neq 0 \to \varphi_{402}) \to \mathrm{FLim}(f,c,\alpha) \to \mathrm{FLim}(g,c,\beta) \to \beta \neq 0 \to \varphi_{418}$
\item $\varphi_{531} :\equiv (\mathrm{FLim}(f,c,\alpha) \to \varphi_{368}) \to (\varphi_{185} \to \varphi_{182}) \to \mathrm{FLim}(f,c,\alpha) \to \varphi_{362}$
\item $\varphi_{532} :\equiv (\mathrm{FLim}(f,c,\alpha) \to \mathrm{FLim}(g,c,\beta) \to \beta \neq 0 \to \varphi_{359}) \to \varphi_{469}$
\item $\varphi_{533} :\equiv (\mathrm{FLim}(f,c,\alpha) \to \varphi_{253}) \to \varphi_{131} \to (f \circ a(n) = f(a_{n}) \to \varphi_{312}) \to \varphi_{357}$
\item $\varphi_{534} :\equiv (\mathrm{FLim}(f,c,\alpha) \to \mathrm{FLim}(g,c,\beta) \to \beta \neq 0 \to \varphi_{440}) \to \mathrm{FLim}(f,c,\alpha) \to \mathrm{FLim}(g,c,\beta) \to \beta \neq 0 \to \varphi_{394}$
\item $\varphi_{535} :\equiv (\mathrm{FLim}(f,c,\alpha) \to \mathrm{FLim}(g,c,\beta) \to \beta \neq 0 \to \varphi_{417}) \to \mathrm{FLim}(f,c,\alpha) \to \mathrm{FLim}(g,c,\beta) \to \beta \neq 0 \to \varphi_{427}$
\item $\varphi_{536} :\equiv (\mathrm{FLim}(f,c,\alpha) \to \mathrm{FLim}(g,c,\beta) \to \beta \neq 0 \to \varphi_{418}) \to \mathrm{FLim}(f,c,\alpha) \to \mathrm{FLim}(g,c,\beta) \to \beta \neq 0 \to \varphi_{428}$
\item $\varphi_{537} :\equiv (\mathrm{FLim}(f,c,\alpha) \to \mathrm{FLim}(g,c,\beta) \to \beta \neq 0 \to \varphi_{423}) \to \mathrm{FLim}(f,c,\alpha) \to \mathrm{FLim}(g,c,\beta) \to \beta \neq 0 \to \varphi_{429}$
\item $\varphi_{538} :\equiv (\mathrm{FLim}(f,c,\alpha) \to \varphi_{443}) \to \mathrm{FLim}(f,c,\alpha) \to \mathrm{FLim}(g,c,\beta) \to \beta \neq 0 \to \varphi_{423}$
\item $\varphi_{539} :\equiv \forall \delta\,(\varphi_{28} \to \neg (\delta \in \mathbb{R} \land 0 < \delta \land \forall x \in \mathrm{dom}(g)\,(\varphi_{94}))) \to \varphi_{452} \to (\varphi_{28} \to \forall \delta\,(\varphi_{28})) \to \varphi_{24}$
\item $\varphi_{540} :\equiv (\mathrm{FLim}(f,c,\alpha) \to \mathrm{FLim}(g,c,\beta) \to \beta \neq 0 \to \varphi_{427}) \to \mathrm{FLim}(f,c,\alpha) \to \varphi_{443}$
\item $\varphi_{541} :\equiv (\mathrm{FLim}(f,c,\alpha) \to \varphi_{447}) \to \mathrm{FLim}(f,c,\alpha) \to \mathrm{FLim}(g,c,\beta) \to \beta \neq 0 \to \varphi_{422}$
\item $\varphi_{542} :\equiv (\mathrm{FLim}(f,c,\alpha) \to \mathrm{FLim}(g,c,\beta) \to \beta \neq 0 \to \varphi_{374}) \to (\mathrm{FLim}(f,c,\alpha) \to \varphi_{262}) \to \mathrm{FLim}(f,c,\alpha) \to \varphi_{399}$
\item $\varphi_{543} :\equiv (\mathrm{FLim}(f,c,\alpha) \to \mathrm{FLim}(g,c,\beta) \to \beta \neq 0 \to \varphi_{428}) \to \mathrm{FLim}(f,c,\alpha) \to \varphi_{453}$
\item $\varphi_{544} :\equiv (0 \leq 0 \to \lvert 0 \rvert = 0) \to (\mathrm{FLim}(f,c,\alpha) \to \varphi_{435}) \to \mathrm{FLim}(f,c,\alpha) \to \varphi_{437}$
\item $\varphi_{545} :\equiv (\mathrm{FLim}(f,c,\alpha) \to \varphi_{286}) \to (\mathrm{FLim}(f,c,\alpha) \to \varphi_{311}) \to \varphi_{463}$
\item $\varphi_{546} :\equiv (\mathrm{FLim}(f,c,\alpha) \to \varphi_{252}) \to (\mathrm{FLim}(f,c,\alpha) \to \varphi_{216}) \to \varphi_{481}$
\item $\varphi_{547} :\equiv (\mathrm{FLim}(f,c,\alpha) \to \varphi_{434}) \to \mathrm{FLim}(f,c,\alpha) \to \mathrm{FLim}(g,c,\beta) \to \beta \neq 0 \to \varphi_{459}$
\item $\varphi_{548} :\equiv (\lvert g(x) \rvert = 0 \to 0 < \lvert g(x) \rvert \to 0 < 0) \to (\mathrm{FLim}(f,c,\alpha) \to \varphi_{444}) \to \mathrm{FLim}(f,c,\alpha) \to \varphi_{447}$
\item $\varphi_{549} :\equiv (\mathrm{Conv}(f \circ a,\alpha) \to \mathrm{Conv}(g \circ a,\beta) \to \beta \neq 0 \to \varphi_{153}) \to \varphi_{528}$
\item $\varphi_{550} :\equiv \varphi_{9} \to (\mathrm{FLim}(f,c,\alpha) \to \varphi_{453}) \to \mathrm{FLim}(f,c,\alpha) \to \varphi_{448}$
\item $\varphi_{551} :\equiv \varphi_{33} \to (\mathrm{FLim}(f,c,\alpha) \to \varphi_{448}) \to \mathrm{FLim}(f,c,\alpha) \to \varphi_{446}$
\item $\varphi_{552} :\equiv \varphi_{504} \to (\forall v_{767}\,(v_{766} \in \mathbb{R} \to v_{767} \in \mathbb{R} \to \varphi_{43}) \leftrightarrow \forall v_{767}\,(\varphi_{230}))$
\item $\varphi_{553} :\equiv (\mathrm{FLim}(f,c,\alpha) \to \mathrm{FLim}(g,c,\beta) \to \beta \neq 0 \to \varphi_{422}) \to \varphi_{478}$
\item $\varphi_{554} :\equiv (\mathrm{PNbhd}((f / g),c) \leftrightarrow \exists r\,(r \in \mathbb{R} \land 0 < r \land \varphi_{78})) \to \varphi_{534}$
\item $\varphi_{555} :\equiv \forall n\,(\varphi_{190} \to n \in \mathbb{N} \to a_{n} \in \mathrm{dom}(g)) \to \varphi_{409} \to (\varphi_{190} \to \forall n\,(\varphi_{190})) \to \mathrm{FLim}(f,c,\alpha) \to \varphi_{226}$
\item $\varphi_{556} :\equiv \forall n\,(\varphi_{190} \to n \in \mathbb{N} \to a_{n} \in \mathrm{dom}(f)) \to \varphi_{410} \to (\varphi_{190} \to \forall n\,(\varphi_{190})) \to \mathrm{FLim}(f,c,\alpha) \to \varphi_{228}$
\item $\varphi_{557} :\equiv (\mathrm{FLim}(f,c,\alpha) \to \varphi_{439}) \to \lvert g(x) \rvert = \lvert g(x) \rvert \to \varphi_{37} \to \mathrm{FLim}(f,c,\alpha) \to \varphi_{449}$
\item $\varphi_{558} :\equiv \forall n\,(\varphi_{190} \to n \in \mathbb{N} \to g \circ a(n) \neq 0) \to \varphi_{413} \to (\varphi_{190} \to \forall n\,(\varphi_{190})) \to \mathrm{FLim}(f,c,\alpha) \to \varphi_{233}$
\item $\varphi_{559} :\equiv (\forall n\,(\varphi_{190} \to n \in \mathbb{N} \to \varphi_{81}) \to \forall n\,(\varphi_{190}) \to \forall n \in \mathbb{N}\,(\varphi_{81})) \to (\varphi_{190} \to \forall n\,(\varphi_{190})) \to \mathrm{FLim}(f,c,\alpha) \to \varphi_{318}$
\item $\varphi_{560} :\equiv (\varphi_{396} \to \varphi_{187} \to \forall v_{766}\,\forall v_{767}\,(v_{766} \in \mathbb{R} \to v_{767} \in \mathbb{R} \to \varphi_{43})) \to (\forall v_{766}\,\forall v_{767}\,(v_{766} \in \mathbb{R} \to v_{767} \in \mathbb{R} \to \varphi_{43}) \leftrightarrow \varphi_{187})$
\item $\varphi_{561} :\equiv \forall n\,(\varphi_{190} \to n \in \mathbb{N} \to a_{n} \in \mathrm{dom}((f / g))) \to \varphi_{421} \to (\varphi_{190} \to \forall n\,(\varphi_{190})) \to \mathrm{FLim}(f,c,\alpha) \to \varphi_{239}$
\item $\varphi_{562} :\equiv \forall n\,(\varphi_{190} \to n \in \mathbb{N} \to a_{n} \in \mathrm{dom}(g) \land a_{n} \neq c) \to \varphi_{430} \to (\varphi_{190} \to \forall n\,(\varphi_{190})) \to \mathrm{FLim}(f,c,\alpha) \to \varphi_{255}$
\item $\varphi_{563} :\equiv \forall n\,(\varphi_{190} \to n \in \mathbb{N} \to a_{n} \in \mathrm{dom}(f) \land a_{n} \neq c) \to \varphi_{431} \to (\varphi_{190} \to \forall n\,(\varphi_{190})) \to \mathrm{FLim}(f,c,\alpha) \to \varphi_{256}$
\item $\varphi_{564} :\equiv (\mathrm{Seq}((f / g) \circ a) \to \varphi_{118} \to \varphi_{160}) \to (\mathrm{FLim}(f,c,\alpha) \to \varphi_{212}) \to \varphi_{501}$
\item $\varphi_{565} :\equiv (\mathrm{FLim}(f,c,\alpha) \to \mathrm{FLim}(g,c,\beta) \to \beta \neq 0 \to \varphi_{459}) \to \varphi_{514}$
\item $\varphi_{566} :\equiv \forall n\,(\varphi_{330}) \to (\forall n\,(\varphi_{330}) \to \forall n\,(\varphi_{190}) \to \varphi_{118}) \to (\varphi_{190} \to \forall n\,(\varphi_{190})) \to \varphi_{335}$
\item $\varphi_{567} :\equiv (\mathrm{FLim}(f,c,\alpha) \to \mathrm{FLim}(g,c,\beta) \to \beta \neq 0 \to \varphi_{426}) \to \varphi_{535}$
\item $\varphi_{568} :\equiv (\mathrm{FLim}(f,c,\alpha) \to \mathrm{FLim}(g,c,\beta) \to \beta \neq 0 \to \varphi_{426}) \to \varphi_{536}$
\item $\varphi_{569} :\equiv (\mathrm{FLim}(f,c,\alpha) \to \mathrm{FLim}(g,c,\beta) \to \beta \neq 0 \to \varphi_{412}) \to \varphi_{540}$
\item $\varphi_{570} :\equiv (\mathrm{FLim}(f,c,\alpha) \to \mathrm{FLim}(g,c,\beta) \to \beta \neq 0 \to \varphi_{402}) \to (\mathrm{FLim}(f,c,\alpha) \to \varphi_{438}) \to \mathrm{FLim}(f,c,\alpha) \to \varphi_{455}$
\item $\varphi_{571} :\equiv (\mathrm{FLim}(f,c,\alpha) \to \mathrm{FLim}(g,c,\beta) \to \beta \neq 0 \to \varphi_{412}) \to \varphi_{543}$
\item $\varphi_{572} :\equiv (\mathrm{FLim}(f,c,\alpha) \to \varphi_{449}) \to (\mathrm{FLim}(f,c,\alpha) \to \varphi_{437}) \to \varphi_{467}$
\item $\varphi_{573} :\equiv (\mathrm{FLim}(f,c,\alpha) \to \mathrm{FLim}(g,c,\beta) \to \beta \neq 0 \to \varphi_{429}) \to \varphi_{550}$
\item $\varphi_{574} :\equiv (\mathrm{FLim}(f,c,\alpha) \to \mathrm{FLim}(g,c,\beta) \to \beta \neq 0 \to \varphi_{379}) \to (\mathrm{FLim}(f,c,\alpha) \to \varphi_{268}) \to \varphi_{542}$
\item $\varphi_{575} :\equiv (\mathrm{FLim}(f,c,\alpha) \to \varphi_{445}) \to \neg 0 < 0 \to (0 = \lvert g(x) \rvert \to 0 < 0 \to 0 < \lvert g(x) \rvert) \to \varphi_{548}$
\item $\varphi_{576} :\equiv (\varphi_{397} \to \forall v_{766}\,\forall v_{767}\,(v_{766} \in \mathbb{R} \to v_{767} \in \mathbb{R} \to \varphi_{43}) \to \varphi_{187}) \to \varphi_{560}$
\item $\varphi_{577} :\equiv (\mathrm{FLim}(f,c,\alpha) \to \mathrm{FLim}(g,c,\beta) \to \beta \neq 0 \to \varphi_{371}) \to \forall d\,(\varphi_{307} \to \neg (d \in \mathbb{R} \land 0 < d \land \forall x \in \mathbb{R}\,(\varphi_{46}))) \to \varphi_{513} \to \varphi_{502}$
\item $\varphi_{578} :\equiv (\forall x\,(\varphi_{404} \to \varphi_{65}) \to \forall x\,(\varphi_{404}) \to \varphi_{78}) \to (\varphi_{404} \to \forall x\,(\varphi_{404})) \to \mathrm{FLim}(f,c,\alpha) \to \varphi_{438}$
\item $\varphi_{579} :\equiv (\mathrm{FLim}(f,c,\alpha) \to \mathrm{FLim}(g,c,\beta) \to \beta \neq 0 \to \varphi_{352}) \to (\mathrm{FLim}(f,c,\alpha) \to \varphi_{443}) \to \varphi_{571}$
\item $\varphi_{580} :\equiv (\mathrm{FLim}(f,c,\alpha) \to \mathrm{FLim}(g,c,\beta) \to \beta \neq 0 \to \varphi_{458}) \to (\mathrm{FLim}(f,c,\alpha) \to \varphi_{411}) \to \varphi_{569}$
\item $\varphi_{581} :\equiv \forall v_{767}\,((v_{766} \in \mathbb{R} \to v_{767} \in \mathbb{R} \to \varphi_{43}) \to \varphi_{230}) \to \forall v_{767}\,(\varphi_{230} \to v_{766} \in \mathbb{R} \to v_{767} \in \mathbb{R} \to \varphi_{43}) \to \varphi_{505} \to \varphi_{552}$
\item $\varphi_{582} :\equiv \forall r\,(\varphi_{395} \to \neg (r \in \mathbb{R} \land 0 < r \land r \leq d \land r \leq \delta)) \to \varphi_{527} \to (\varphi_{395} \to \forall r\,(\varphi_{395})) \to \varphi_{519}$
\end{abbrevlist}
\end{fullproof}

% 自動生成 (prover/build.js): FLimLe  ノード数 574
\begin{fullproof}{定理1.9}{$\forall f\,\forall g\,\forall c\,\forall \alpha\,\forall \beta\,\forall r\,(\mathrm{FLim}(f,c,\alpha) \to \mathrm{FLim}(g,c,\beta) \to r \in \mathbb{R} \to 0 < r \to \forall x \in \mathrm{dom}(f)\,(x \in \mathrm{dom}(g) \to 0 < \lvert x - c \rvert \to \lvert x - c \rvert < r \to f(x) \leq g(x)) \to \alpha \leq \beta)$}{574}\phantomsection\label{tree:FLimLe}
\noindent$\Pi_{1}$:
\begin{ptree}
\AXC{$z = w \to f \circ a(n) = z \to f \circ a(n) = w$\lbl{A6}}
\UIC{$\forall w\,(z = w \to f \circ a(n) = z \to f \circ a(n) = w)$}
\UIC{$\forall z\,\forall w\,(z = w \to f \circ a(n) = z \to f \circ a(n) = w)$}
\AXC{$\forall z\,\forall w\,(z = w \to f \circ a(n) = z \to f \circ a(n) = w) \to \varphi_{32}$\lbl{A2}}
\BIC{$\varphi_{32}$}
\AXC{$\varphi_{32} \to \varphi_{62}$\lbl{A2}}
\BIC{$\varphi_{62}$}
\end{ptree}
\noindent$\Pi_{2}$:
\begin{ptree}
\AXC{$z = w \to f \circ a(n) < z \to f \circ a(n) < w$\lbl{A6}}
\UIC{$\forall w\,(z = w \to f \circ a(n) < z \to f \circ a(n) < w)$}
\UIC{$\forall z\,\forall w\,(z = w \to f \circ a(n) < z \to f \circ a(n) < w)$}
\AXC{$\forall z\,\forall w\,(z = w \to f \circ a(n) < z \to f \circ a(n) < w) \to \varphi_{40}$\lbl{A2}}
\BIC{$\varphi_{40}$}
\AXC{$\varphi_{40} \to \varphi_{63}$\lbl{A2}}
\BIC{$\varphi_{63}$}
\end{ptree}
\noindent$\Pi_{3}$:
\begin{ptree}
\AXC{$z = w \to z = g(a_{n}) \to w = g(a_{n})$\lbl{A6}}
\UIC{$\forall w\,(z = w \to z = g(a_{n}) \to w = g(a_{n}))$}
\UIC{$\forall z\,\forall w\,(z = w \to z = g(a_{n}) \to w = g(a_{n}))$}
\AXC{$\forall z\,\forall w\,(z = w \to z = g(a_{n}) \to w = g(a_{n})) \to \varphi_{26}$\lbl{A2}}
\BIC{$\varphi_{26}$}
\AXC{$\varphi_{26} \to \varphi_{56}$\lbl{A2}}
\BIC{$\varphi_{56}$}
\end{ptree}
\noindent$\Pi_{4}$:
\begin{ptree}
\AXC{$z = w \to z < g(a_{n}) \to w < g(a_{n})$\lbl{A6}}
\UIC{$\forall w\,(z = w \to z < g(a_{n}) \to w < g(a_{n}))$}
\UIC{$\forall z\,\forall w\,(z = w \to z < g(a_{n}) \to w < g(a_{n}))$}
\AXC{$\forall z\,\forall w\,(z = w \to z < g(a_{n}) \to w < g(a_{n})) \to \varphi_{28}$\lbl{A2}}
\BIC{$\varphi_{28}$}
\AXC{$\varphi_{28} \to \varphi_{60}$\lbl{A2}}
\BIC{$\varphi_{60}$}
\end{ptree}
\noindent$\Pi_{5}$:
\begin{ptree}
\AXC{$\varphi_{71}$\lbl{A2}}
\AXC{$\varphi_{205}$\lbl{A1}}
\AXC{$\varphi_{205} \to \mathrm{FLim}(f,c,\alpha) \to \varphi_{202}$\lbl{A1}}
\BIC{$\mathrm{FLim}(f,c,\alpha) \to \varphi_{202}$}
\AXC{$\varphi_{246}$\lbl{A1}}
\BIC{$\varphi_{71} \to \mathrm{FLim}(f,c,\alpha) \to \varphi_{201}$}
\BIC{$\mathrm{FLim}(f,c,\alpha) \to \varphi_{201}$}
\end{ptree}
\noindent$\Pi_{6}$:
\begin{ptree}
\AXC{$\stackrel{\Pi_{5}}{\vdots}$}\noLine
\UIC{$\mathrm{FLim}(f,c,\alpha) \to \varphi_{201}$}
\AXC{$(\mathrm{FLim}(f,c,\alpha) \to \varphi_{201}) \to \mathrm{FLim}(f,c,\alpha) \to \varphi_{180}$\lbl{A1}}
\BIC{$\mathrm{FLim}(f,c,\alpha) \to \varphi_{180}$}
\end{ptree}
\noindent$\Pi_{7}$:
\begin{ptree}
\AXC{$\forall x\,\forall y\,\forall z\,(x < y \to y \leq z \to x < z)$\lbl{[O3]}}
\AXC{$\varphi_{46}$\lbl{A2}}
\BIC{$\forall y\,\forall z\,(\lvert a_{n} - c \rvert < y \to y \leq z \to \lvert a_{n} - c \rvert < z)$}
\AXC{$\varphi_{68}$\lbl{A2}}
\BIC{$\forall z\,(\lvert a_{n} - c \rvert < \rho \to \rho \leq z \to \lvert a_{n} - c \rvert < z)$}
\AXC{$\varphi_{64}$\lbl{A2}}
\BIC{$\lvert a_{n} - c \rvert < \rho \to \rho \leq r \to \lvert a_{n} - c \rvert < r$}
\end{ptree}
\noindent$\Pi_{8}$:
\begin{ptree}
\AXC{$\stackrel{\Pi_{6}}{\vdots}$}\noLine
\UIC{$\mathrm{FLim}(f,c,\alpha) \to \varphi_{180}$}
\AXC{$\stackrel{\Pi_{7}}{\vdots}$}\noLine
\UIC{$\lvert a_{n} - c \rvert < \rho \to \rho \leq r \to \lvert a_{n} - c \rvert < r$}
\AXC{$(\lvert a_{n} - c \rvert < \rho \to \rho \leq r \to \lvert a_{n} - c \rvert < r) \to \varphi_{249}$\lbl{A1}}
\BIC{$\varphi_{249}$}
\BIC{$\varphi_{170} \to \mathrm{FLim}(f,c,\alpha) \to \varphi_{174}$}
\end{ptree}
\noindent$\Pi_{9}$:
\begin{ptree}
\AXC{$\mathrm{FLim}(f,c,\alpha) \to \varphi_{154}$\lbl{A1}}
\AXC{$(\mathrm{FLim}(f,c,\alpha) \to \varphi_{154}) \to \varphi_{170}$\lbl{A1}}
\BIC{$\varphi_{170}$}
\AXC{$\stackrel{\Pi_{8}}{\vdots}$}\noLine
\UIC{$\varphi_{170} \to \mathrm{FLim}(f,c,\alpha) \to \varphi_{174}$}
\BIC{$\mathrm{FLim}(f,c,\alpha) \to \varphi_{174}$}
\end{ptree}
\noindent$\Pi_{10}$:
\begin{ptree}
\AXC{$\forall x\,\forall y\,\forall z\,(x < y \to y \leq z \to x < z)$\lbl{[O3]}}
\AXC{$\varphi_{46}$\lbl{A2}}
\BIC{$\forall y\,\forall z\,(\lvert a_{n} - c \rvert < y \to y \leq z \to \lvert a_{n} - c \rvert < z)$}
\AXC{$\varphi_{68}$\lbl{A2}}
\BIC{$\forall z\,(\lvert a_{n} - c \rvert < \rho \to \rho \leq z \to \lvert a_{n} - c \rvert < z)$}
\AXC{$\varphi_{65}$\lbl{A2}}
\BIC{$\lvert a_{n} - c \rvert < \rho \to \rho \leq d \to \lvert a_{n} - c \rvert < d$}
\end{ptree}
\noindent$\Pi_{11}$:
\begin{ptree}
\AXC{$\stackrel{\Pi_{6}}{\vdots}$}\noLine
\UIC{$\mathrm{FLim}(f,c,\alpha) \to \varphi_{180}$}
\AXC{$\stackrel{\Pi_{10}}{\vdots}$}\noLine
\UIC{$\lvert a_{n} - c \rvert < \rho \to \rho \leq d \to \lvert a_{n} - c \rvert < d$}
\AXC{$(\lvert a_{n} - c \rvert < \rho \to \rho \leq d \to \lvert a_{n} - c \rvert < d) \to \varphi_{250}$\lbl{A1}}
\BIC{$\varphi_{250}$}
\BIC{$\varphi_{171} \to \mathrm{FLim}(f,c,\alpha) \to \varphi_{175}$}
\end{ptree}
\noindent$\Pi_{12}$:
\begin{ptree}
\AXC{$\mathrm{FLim}(f,c,\alpha) \to \varphi_{154}$\lbl{A1}}
\AXC{$(\mathrm{FLim}(f,c,\alpha) \to \varphi_{154}) \to \varphi_{171}$\lbl{A1}}
\BIC{$\varphi_{171}$}
\AXC{$\stackrel{\Pi_{11}}{\vdots}$}\noLine
\UIC{$\varphi_{171} \to \mathrm{FLim}(f,c,\alpha) \to \varphi_{175}$}
\BIC{$\mathrm{FLim}(f,c,\alpha) \to \varphi_{175}$}
\end{ptree}
\noindent$\Pi_{13}$:
\begin{ptree}
\AXC{$\forall a\,\forall n\,(\mathrm{Seq}(a) \to n \in \mathbb{N} \to a_{n} \in \mathbb{R})$\lbl{[S1]}}
\AXC{$\varphi_{41}$\lbl{A2}}
\BIC{$\forall n\,(\mathrm{Seq}(a) \to n \in \mathbb{N} \to a_{n} \in \mathbb{R})$}
\AXC{$\varphi_{33}$\lbl{A2}}
\BIC{$\mathrm{Seq}(a) \to n \in \mathbb{N} \to a_{n} \in \mathbb{R}$}
\AXC{$(\mathrm{Seq}(a) \to n \in \mathbb{N} \to a_{n} \in \mathbb{R}) \to \varphi_{165} \to \varphi_{177} \to \varphi_{183}$\lbl{A1}}
\BIC{$\varphi_{165} \to \varphi_{177} \to \varphi_{183}$}
\end{ptree}
\noindent$\Pi_{14}$:
\begin{ptree}
\AXC{$\varphi_{177}$\lbl{A1}}
\AXC{$\varphi_{205}$\lbl{A1}}
\AXC{$\varphi_{205} \to \varphi_{165}$\lbl{A1}}
\BIC{$\varphi_{165}$}
\AXC{$\stackrel{\Pi_{13}}{\vdots}$}\noLine
\UIC{$\varphi_{165} \to \varphi_{177} \to \varphi_{183}$}
\BIC{$\varphi_{177} \to \varphi_{183}$}
\BIC{$\varphi_{183}$}
\end{ptree}
\noindent$\Pi_{15}$:
\begin{ptree}
\AXC{$\mathrm{FLim}(f,c,\alpha) \to \mathrm{FLim}(g,c,\beta) \to \varphi_{158}$\lbl{A1}}
\AXC{$(\mathrm{FLim}(f,c,\alpha) \to \mathrm{FLim}(g,c,\beta) \to \varphi_{158}) \to \varphi_{206}$\lbl{A1}}
\BIC{$\varphi_{206}$}
\end{ptree}
\noindent$\Pi_{16}$:
\begin{ptree}
\AXC{$\varphi_{31} \to a_{n} \in \mathbb{R} \to \varphi_{45}$\lbl{A2}}
\AXC{$\stackrel{\Pi_{15}}{\vdots}$}\noLine
\UIC{$\varphi_{206}$}
\AXC{$\varphi_{206} \to (\varphi_{31} \to a_{n} \in \mathbb{R} \to \varphi_{45}) \to \mathrm{FLim}(f,c,\alpha) \to \varphi_{209}$\lbl{A1}}
\BIC{$(\varphi_{31} \to a_{n} \in \mathbb{R} \to \varphi_{45}) \to \mathrm{FLim}(f,c,\alpha) \to \varphi_{209}$}
\BIC{$\mathrm{FLim}(f,c,\alpha) \to \varphi_{209}$}
\end{ptree}
\noindent$\Pi_{17}$:
\begin{ptree}
\AXC{$\stackrel{\Pi_{5}}{\vdots}$}\noLine
\UIC{$\mathrm{FLim}(f,c,\alpha) \to \varphi_{201}$}
\AXC{$\stackrel{\Pi_{14}}{\vdots}$}\noLine
\UIC{$\varphi_{183}$}
\AXC{$\stackrel{\Pi_{16}}{\vdots}$}\noLine
\UIC{$\mathrm{FLim}(f,c,\alpha) \to \varphi_{209}$}
\AXC{$(\mathrm{FLim}(f,c,\alpha) \to \varphi_{209}) \to \varphi_{183} \to \varphi_{260}$\lbl{A1}}
\BIC{$\varphi_{183} \to \varphi_{260}$}
\BIC{$\varphi_{260}$}
\BIC{$(\mathrm{FLim}(f,c,\alpha) \to \varphi_{175}) \to \mathrm{FLim}(f,c,\alpha) \to \varphi_{195}$}
\end{ptree}
\noindent$\Pi_{18}$:
\begin{ptree}
\AXC{$\varphi_{42} \to a_{n} \in \mathrm{dom}(f) \to \varphi_{54}$\lbl{A2}}
\AXC{$\varphi_{122}$\lbl{A1}}
\AXC{$\varphi_{228}$\lbl{A1}}
\BIC{$(\varphi_{42} \to a_{n} \in \mathrm{dom}(f) \to \varphi_{54}) \to \mathrm{FLim}(f,c,\alpha) \to \varphi_{136}$}
\BIC{$\mathrm{FLim}(f,c,\alpha) \to \varphi_{136}$}
\AXC{$\varphi_{269}$\lbl{A1}}
\BIC{$(\mathrm{FLim}(f,c,\alpha) \to \varphi_{195}) \to \varphi_{259}$}
\end{ptree}
\noindent$\Pi_{19}$:
\begin{ptree}
\AXC{$\stackrel{\Pi_{12}}{\vdots}$}\noLine
\UIC{$\mathrm{FLim}(f,c,\alpha) \to \varphi_{175}$}
\AXC{$\stackrel{\Pi_{17}}{\vdots}$}\noLine
\UIC{$(\mathrm{FLim}(f,c,\alpha) \to \varphi_{175}) \to \mathrm{FLim}(f,c,\alpha) \to \varphi_{195}$}
\BIC{$\mathrm{FLim}(f,c,\alpha) \to \varphi_{195}$}
\AXC{$\stackrel{\Pi_{18}}{\vdots}$}\noLine
\UIC{$(\mathrm{FLim}(f,c,\alpha) \to \varphi_{195}) \to \varphi_{259}$}
\BIC{$\varphi_{259}$}
\end{ptree}
\noindent$\Pi_{20}$:
\begin{ptree}
\AXC{$\stackrel{\Pi_{9}}{\vdots}$}\noLine
\UIC{$\mathrm{FLim}(f,c,\alpha) \to \varphi_{174}$}
\AXC{$\stackrel{\Pi_{5}}{\vdots}$}\noLine
\UIC{$\mathrm{FLim}(f,c,\alpha) \to \varphi_{201}$}
\AXC{$\stackrel{\Pi_{19}}{\vdots}$}\noLine
\UIC{$\varphi_{259}$}
\BIC{$(\mathrm{FLim}(f,c,\alpha) \to \varphi_{174}) \to \mathrm{FLim}(f,c,\alpha) \to \varphi_{182}$}
\BIC{$\mathrm{FLim}(f,c,\alpha) \to \varphi_{182}$}
\end{ptree}
\noindent$\Pi_{21}$:
\begin{ptree}
\AXC{$x = x$\lbl{A5}}
\AXC{$x = y \to x = x \to y = x$\lbl{A6}}
\AXC{$(x = y \to x = x \to y = x) \to x = x \to x = y \to y = x$\lbl{A1}}
\BIC{$x = x \to x = y \to y = x$}
\BIC{$x = y \to y = x$}
\UIC{$\forall y\,(x = y \to y = x)$}
\UIC{$\forall x\,\forall y\,(x = y \to y = x)$}
\end{ptree}
\noindent$\Pi_{22}$:
\begin{ptree}
\AXC{$\varphi_{71}$\lbl{A2}}
\AXC{$\varphi_{205}$\lbl{A1}}
\AXC{$\varphi_{205} \to \mathrm{FLim}(f,c,\alpha) \to \varphi_{202}$\lbl{A1}}
\BIC{$\mathrm{FLim}(f,c,\alpha) \to \varphi_{202}$}
\AXC{$\varphi_{246}$\lbl{A1}}
\BIC{$\varphi_{71} \to \mathrm{FLim}(f,c,\alpha) \to \varphi_{201}$}
\BIC{$\mathrm{FLim}(f,c,\alpha) \to \varphi_{201}$}
\end{ptree}
\noindent$\Pi_{23}$:
\begin{ptree}
\AXC{$\forall f\,\forall c\,\forall \alpha\,(\mathrm{FLim}(f,c,\alpha) \to c \in \mathbb{R})$\lbl{[D3]}}
\AXC{$\varphi_{20}$\lbl{A2}}
\BIC{$\forall c\,\forall \alpha\,(\mathrm{FLim}(f,c,\alpha) \to c \in \mathbb{R})$}
\AXC{$\varphi_{14}$\lbl{A2}}
\BIC{$\forall \alpha\,(\mathrm{FLim}(f,c,\alpha) \to c \in \mathbb{R})$}
\AXC{$\varphi_{7}$\lbl{A2}}
\BIC{$\mathrm{FLim}(f,c,\alpha) \to c \in \mathbb{R}$}
\AXC{$\varphi_{117}$\lbl{A1}}
\BIC{$(\mathrm{FLim}(f,c,\alpha) \to \mathrm{FLim}(f,c,\alpha)) \to \varphi_{82}$}
\end{ptree}
\noindent$\Pi_{24}$:
\begin{ptree}
\AXC{$\mathrm{FLim}(f,c,\alpha) \to \mathrm{FLim}(f,c,\alpha)$\lbl{A1}}
\AXC{$\stackrel{\Pi_{23}}{\vdots}$}\noLine
\UIC{$(\mathrm{FLim}(f,c,\alpha) \to \mathrm{FLim}(f,c,\alpha)) \to \varphi_{82}$}
\BIC{$\varphi_{82}$}
\end{ptree}
\noindent$\Pi_{25}$:
\begin{ptree}
\AXC{$\stackrel{\Pi_{24}}{\vdots}$}\noLine
\UIC{$\varphi_{82}$}
\AXC{$\stackrel{\Pi_{14}}{\vdots}$}\noLine
\UIC{$\varphi_{183}$}
\AXC{$\forall y\,\forall c\,(y \in \mathbb{R} \to c \in \mathbb{R} \to 0 < \lvert y - c \rvert \to y \neq c)$\lbl{[E37]}}
\AXC{$\varphi_{70}$\lbl{A2}}
\BIC{$\forall c\,(a_{n} \in \mathbb{R} \to c \in \mathbb{R} \to 0 < \lvert a_{n} - c \rvert \to a_{n} \neq c)$}
\AXC{$\varphi_{77}$\lbl{A2}}
\BIC{$a_{n} \in \mathbb{R} \to c \in \mathbb{R} \to 0 < \lvert a_{n} - c \rvert \to a_{n} \neq c$}
\AXC{$\varphi_{262}$\lbl{A1}}
\BIC{$\varphi_{183} \to \varphi_{82} \to (\mathrm{FLim}(f,c,\alpha) \to \varphi_{176}) \to \varphi_{181}$}
\BIC{$\varphi_{82} \to (\mathrm{FLim}(f,c,\alpha) \to \varphi_{176}) \to \varphi_{181}$}
\BIC{$(\mathrm{FLim}(f,c,\alpha) \to \varphi_{176}) \to \varphi_{181}$}
\end{ptree}
\noindent$\Pi_{26}$:
\begin{ptree}
\AXC{$\stackrel{\Pi_{22}}{\vdots}$}\noLine
\UIC{$\mathrm{FLim}(f,c,\alpha) \to \varphi_{201}$}
\AXC{$(\mathrm{FLim}(f,c,\alpha) \to \varphi_{201}) \to \mathrm{FLim}(f,c,\alpha) \to \varphi_{176}$\lbl{A1}}
\BIC{$\mathrm{FLim}(f,c,\alpha) \to \varphi_{176}$}
\AXC{$\stackrel{\Pi_{25}}{\vdots}$}\noLine
\UIC{$(\mathrm{FLim}(f,c,\alpha) \to \varphi_{176}) \to \varphi_{181}$}
\BIC{$\varphi_{181}$}
\end{ptree}
\noindent$\Pi_{27}$:
\begin{ptree}
\AXC{$\forall x\,\forall y\,\forall z\,(x < y \to y \leq z \to x < z)$\lbl{[O3]}}
\AXC{$\varphi_{46}$\lbl{A2}}
\BIC{$\forall y\,\forall z\,(\lvert a_{n} - c \rvert < y \to y \leq z \to \lvert a_{n} - c \rvert < z)$}
\AXC{$\varphi_{68}$\lbl{A2}}
\BIC{$\forall z\,(\lvert a_{n} - c \rvert < \rho \to \rho \leq z \to \lvert a_{n} - c \rvert < z)$}
\AXC{$\varphi_{65}$\lbl{A2}}
\BIC{$\lvert a_{n} - c \rvert < \rho \to \rho \leq d \to \lvert a_{n} - c \rvert < d$}
\end{ptree}
\noindent$\Pi_{28}$:
\begin{ptree}
\AXC{$\stackrel{\Pi_{27}}{\vdots}$}\noLine
\UIC{$\lvert a_{n} - c \rvert < \rho \to \rho \leq d \to \lvert a_{n} - c \rvert < d$}
\AXC{$(\lvert a_{n} - c \rvert < \rho \to \rho \leq d \to \lvert a_{n} - c \rvert < d) \to \varphi_{250}$\lbl{A1}}
\BIC{$\varphi_{250}$}
\end{ptree}
\noindent$\Pi_{29}$:
\begin{ptree}
\AXC{$\stackrel{\Pi_{22}}{\vdots}$}\noLine
\UIC{$\mathrm{FLim}(f,c,\alpha) \to \varphi_{201}$}
\AXC{$(\mathrm{FLim}(f,c,\alpha) \to \varphi_{201}) \to \mathrm{FLim}(f,c,\alpha) \to \varphi_{180}$\lbl{A1}}
\BIC{$\mathrm{FLim}(f,c,\alpha) \to \varphi_{180}$}
\AXC{$\stackrel{\Pi_{28}}{\vdots}$}\noLine
\UIC{$\varphi_{250}$}
\BIC{$\varphi_{171} \to \mathrm{FLim}(f,c,\alpha) \to \varphi_{175}$}
\end{ptree}
\noindent$\Pi_{30}$:
\begin{ptree}
\AXC{$\mathrm{FLim}(f,c,\alpha) \to \varphi_{154}$\lbl{A1}}
\AXC{$(\mathrm{FLim}(f,c,\alpha) \to \varphi_{154}) \to \varphi_{171}$\lbl{A1}}
\BIC{$\varphi_{171}$}
\AXC{$\stackrel{\Pi_{29}}{\vdots}$}\noLine
\UIC{$\varphi_{171} \to \mathrm{FLim}(f,c,\alpha) \to \varphi_{175}$}
\BIC{$\mathrm{FLim}(f,c,\alpha) \to \varphi_{175}$}
\end{ptree}
\noindent$\Pi_{31}$:
\begin{ptree}
\AXC{$\mathrm{FLim}(f,c,\alpha) \to \mathrm{FLim}(g,c,\beta) \to \varphi_{158}$\lbl{A1}}
\AXC{$(\mathrm{FLim}(f,c,\alpha) \to \mathrm{FLim}(g,c,\beta) \to \varphi_{158}) \to \varphi_{206}$\lbl{A1}}
\BIC{$\varphi_{206}$}
\end{ptree}
\noindent$\Pi_{32}$:
\begin{ptree}
\AXC{$\varphi_{31} \to a_{n} \in \mathbb{R} \to \varphi_{45}$\lbl{A2}}
\AXC{$\stackrel{\Pi_{31}}{\vdots}$}\noLine
\UIC{$\varphi_{206}$}
\AXC{$\varphi_{206} \to (\varphi_{31} \to a_{n} \in \mathbb{R} \to \varphi_{45}) \to \mathrm{FLim}(f,c,\alpha) \to \varphi_{209}$\lbl{A1}}
\BIC{$(\varphi_{31} \to a_{n} \in \mathbb{R} \to \varphi_{45}) \to \mathrm{FLim}(f,c,\alpha) \to \varphi_{209}$}
\BIC{$\mathrm{FLim}(f,c,\alpha) \to \varphi_{209}$}
\end{ptree}
\noindent$\Pi_{33}$:
\begin{ptree}
\AXC{$\stackrel{\Pi_{22}}{\vdots}$}\noLine
\UIC{$\mathrm{FLim}(f,c,\alpha) \to \varphi_{201}$}
\AXC{$\stackrel{\Pi_{14}}{\vdots}$}\noLine
\UIC{$\varphi_{183}$}
\AXC{$\stackrel{\Pi_{32}}{\vdots}$}\noLine
\UIC{$\mathrm{FLim}(f,c,\alpha) \to \varphi_{209}$}
\AXC{$(\mathrm{FLim}(f,c,\alpha) \to \varphi_{209}) \to \varphi_{183} \to \varphi_{260}$\lbl{A1}}
\BIC{$\varphi_{183} \to \varphi_{260}$}
\BIC{$\varphi_{260}$}
\BIC{$(\mathrm{FLim}(f,c,\alpha) \to \varphi_{175}) \to \mathrm{FLim}(f,c,\alpha) \to \varphi_{195}$}
\end{ptree}
\noindent$\Pi_{34}$:
\begin{ptree}
\AXC{$\stackrel{\Pi_{26}}{\vdots}$}\noLine
\UIC{$\varphi_{181}$}
\AXC{$\stackrel{\Pi_{30}}{\vdots}$}\noLine
\UIC{$\mathrm{FLim}(f,c,\alpha) \to \varphi_{175}$}
\AXC{$\stackrel{\Pi_{33}}{\vdots}$}\noLine
\UIC{$(\mathrm{FLim}(f,c,\alpha) \to \varphi_{175}) \to \mathrm{FLim}(f,c,\alpha) \to \varphi_{195}$}
\BIC{$\mathrm{FLim}(f,c,\alpha) \to \varphi_{195}$}
\AXC{$\varphi_{255}$\lbl{A1}}
\BIC{$\varphi_{181} \to \mathrm{FLim}(f,c,\alpha) \to \varphi_{185}$}
\BIC{$\mathrm{FLim}(f,c,\alpha) \to \varphi_{185}$}
\AXC{$\varphi_{234}$\lbl{A1}}
\BIC{$\varphi_{162} \to n \in \mathbb{N} \to a_{n} \in \mathrm{dom}(f) \land a_{n} \neq c$}
\UIC{$\forall n\,(\varphi_{162} \to n \in \mathbb{N} \to a_{n} \in \mathrm{dom}(f) \land a_{n} \neq c)$}
\end{ptree}
\noindent$\Pi_{35}$:
\begin{ptree}
\AXC{$\varphi_{162} \to \forall n\,(\varphi_{162})$\lbl{A4}}
\AXC{$\varphi_{236}$\lbl{A3}}
\AXC{$\stackrel{\Pi_{34}}{\vdots}$}\noLine
\UIC{$\forall n\,(\varphi_{162} \to n \in \mathbb{N} \to a_{n} \in \mathrm{dom}(f) \land a_{n} \neq c)$}
\AXC{$\varphi_{273}$\lbl{A1}}
\BIC{$\varphi_{236} \to (\varphi_{162} \to \forall n\,(\varphi_{162})) \to \mathrm{FLim}(f,c,\alpha) \to \varphi_{192}$}
\BIC{$(\varphi_{162} \to \forall n\,(\varphi_{162})) \to \mathrm{FLim}(f,c,\alpha) \to \varphi_{192}$}
\BIC{$\mathrm{FLim}(f,c,\alpha) \to \varphi_{192}$}
\end{ptree}
\noindent$\Pi_{36}$:
\begin{ptree}
\AXC{$\varphi_{231}$\lbl{A3}}
\AXC{$\varphi_{177}$\lbl{A1}}
\AXC{$\varphi_{47}$\lbl{A2}}
\AXC{$\stackrel{\Pi_{35}}{\vdots}$}\noLine
\UIC{$\mathrm{FLim}(f,c,\alpha) \to \varphi_{192}$}
\AXC{$\varphi_{242}$\lbl{A1}}
\BIC{$\varphi_{47} \to \mathrm{FLim}(f,c,\alpha) \to \varphi_{185}$}
\BIC{$\mathrm{FLim}(f,c,\alpha) \to \varphi_{185}$}
\AXC{$\varphi_{252}$\lbl{A1}}
\BIC{$\varphi_{177} \to \mathrm{FLim}(f,c,\alpha) \to \varphi_{168}$}
\BIC{$\mathrm{FLim}(f,c,\alpha) \to \varphi_{168}$}
\AXC{$\varphi_{229}$\lbl{A1}}
\BIC{$\varphi_{162} \to n \in \mathbb{N} \to a_{n} \in \mathrm{dom}(f)$}
\UIC{$\forall n\,(\varphi_{162} \to n \in \mathbb{N} \to a_{n} \in \mathrm{dom}(f))$}
\AXC{$\varphi_{271}$\lbl{A1}}
\BIC{$\varphi_{231} \to (\varphi_{162} \to \forall n\,(\varphi_{162})) \to \mathrm{FLim}(f,c,\alpha) \to \varphi_{172}$}
\BIC{$(\varphi_{162} \to \forall n\,(\varphi_{162})) \to \mathrm{FLim}(f,c,\alpha) \to \varphi_{172}$}
\end{ptree}
\noindent$\Pi_{37}$:
\begin{ptree}
\AXC{$\forall f\,\forall c\,\forall \alpha\,(\mathrm{FLim}(f,c,\alpha) \to \mathrm{RFn}(f))$\lbl{[D2]}}
\AXC{$\varphi_{12}$\lbl{A2}}
\BIC{$\forall c\,\forall \alpha\,(\mathrm{FLim}(f,c,\alpha) \to \mathrm{RFn}(f))$}
\AXC{$\varphi_{5}$\lbl{A2}}
\BIC{$\forall \alpha\,(\mathrm{FLim}(f,c,\alpha) \to \mathrm{RFn}(f))$}
\AXC{$\varphi_{2}$\lbl{A2}}
\BIC{$\mathrm{FLim}(f,c,\alpha) \to \mathrm{RFn}(f)$}
\AXC{$\varphi_{114}$\lbl{A1}}
\BIC{$(\mathrm{FLim}(f,c,\alpha) \to \mathrm{FLim}(f,c,\alpha)) \to \varphi_{79}$}
\end{ptree}
\noindent$\Pi_{38}$:
\begin{ptree}
\AXC{$\mathrm{FLim}(f,c,\alpha) \to \mathrm{FLim}(f,c,\alpha)$\lbl{A1}}
\AXC{$\stackrel{\Pi_{37}}{\vdots}$}\noLine
\UIC{$(\mathrm{FLim}(f,c,\alpha) \to \mathrm{FLim}(f,c,\alpha)) \to \varphi_{79}$}
\BIC{$\varphi_{79}$}
\end{ptree}
\noindent$\Pi_{39}$:
\begin{ptree}
\AXC{$\varphi_{205}$\lbl{A1}}
\AXC{$\varphi_{205} \to \varphi_{165}$\lbl{A1}}
\BIC{$\varphi_{165}$}
\AXC{$\stackrel{\Pi_{38}}{\vdots}$}\noLine
\UIC{$\varphi_{79}$}
\AXC{$\forall f\,\forall a\,(\mathrm{RFn}(f) \to \mathrm{Seq}(a) \to \varphi_{57})$\lbl{[S8]}}
\AXC{$\varphi_{144}$\lbl{A2}}
\BIC{$\forall a\,(\mathrm{RFn}(f) \to \mathrm{Seq}(a) \to \varphi_{57})$}
\AXC{$\varphi_{137}$\lbl{A2}}
\BIC{$\mathrm{RFn}(f) \to \mathrm{Seq}(a) \to \varphi_{57}$}
\AXC{$\varphi_{264}$\lbl{A1}}
\BIC{$\varphi_{79} \to \varphi_{165} \to (\mathrm{FLim}(f,c,\alpha) \to \varphi_{172}) \to \varphi_{203}$}
\BIC{$\varphi_{165} \to (\mathrm{FLim}(f,c,\alpha) \to \varphi_{172}) \to \varphi_{203}$}
\BIC{$(\mathrm{FLim}(f,c,\alpha) \to \varphi_{172}) \to \varphi_{203}$}
\end{ptree}
\noindent$\Pi_{40}$:
\begin{ptree}
\AXC{$\varphi_{162} \to \forall n\,(\varphi_{162})$\lbl{A4}}
\AXC{$\stackrel{\Pi_{36}}{\vdots}$}\noLine
\UIC{$(\varphi_{162} \to \forall n\,(\varphi_{162})) \to \mathrm{FLim}(f,c,\alpha) \to \varphi_{172}$}
\BIC{$\mathrm{FLim}(f,c,\alpha) \to \varphi_{172}$}
\AXC{$\stackrel{\Pi_{39}}{\vdots}$}\noLine
\UIC{$(\mathrm{FLim}(f,c,\alpha) \to \varphi_{172}) \to \varphi_{203}$}
\BIC{$\varphi_{203}$}
\end{ptree}
\noindent$\Pi_{41}$:
\begin{ptree}
\AXC{$\varphi_{177}$\lbl{A1}}
\AXC{$\varphi_{52}$\lbl{A2}}
\AXC{$\stackrel{\Pi_{40}}{\vdots}$}\noLine
\UIC{$\varphi_{203}$}
\AXC{$\varphi_{203} \to \mathrm{FLim}(f,c,\alpha) \to \varphi_{199}$\lbl{A1}}
\BIC{$\mathrm{FLim}(f,c,\alpha) \to \varphi_{199}$}
\AXC{$\varphi_{244}$\lbl{A1}}
\BIC{$\varphi_{52} \to \mathrm{FLim}(f,c,\alpha) \to \varphi_{197}$}
\BIC{$\mathrm{FLim}(f,c,\alpha) \to \varphi_{197}$}
\AXC{$\varphi_{257}$\lbl{A1}}
\BIC{$\varphi_{177} \to \mathrm{FLim}(f,c,\alpha) \to \varphi_{187}$}
\BIC{$\mathrm{FLim}(f,c,\alpha) \to \varphi_{187}$}
\end{ptree}
\noindent$\Pi_{42}$:
\begin{ptree}
\AXC{$\stackrel{\Pi_{21}}{\vdots}$}\noLine
\UIC{$\forall x\,\forall y\,(x = y \to y = x)$}
\AXC{$\varphi_{18}$\lbl{A2}}
\BIC{$\forall y\,(f \circ a(n) = y \to y = f \circ a(n))$}
\AXC{$\varphi_{66}$\lbl{A2}}
\BIC{$f \circ a(n) = f(a_{n}) \to f(a_{n}) = f \circ a(n)$}
\AXC{$\stackrel{\Pi_{41}}{\vdots}$}\noLine
\UIC{$\mathrm{FLim}(f,c,\alpha) \to \varphi_{187}$}
\AXC{$(\mathrm{FLim}(f,c,\alpha) \to \varphi_{187}) \to \varphi_{210}$\lbl{A1}}
\BIC{$\varphi_{210}$}
\BIC{$\mathrm{FLim}(f,c,\alpha) \to \varphi_{188}$}
\end{ptree}
\noindent$\Pi_{43}$:
\begin{ptree}
\AXC{$\stackrel{\Pi_{3}}{\vdots}$}\noLine
\UIC{$\varphi_{56}$}
\AXC{$\stackrel{\Pi_{4}}{\vdots}$}\noLine
\UIC{$\varphi_{60}$}
\AXC{$\stackrel{\Pi_{20}}{\vdots}$}\noLine
\UIC{$\mathrm{FLim}(f,c,\alpha) \to \varphi_{182}$}
\AXC{$\stackrel{\Pi_{42}}{\vdots}$}\noLine
\UIC{$\mathrm{FLim}(f,c,\alpha) \to \varphi_{188}$}
\AXC{$(\mathrm{FLim}(f,c,\alpha) \to \varphi_{188}) \to \varphi_{247}$\lbl{A1}}
\BIC{$\varphi_{247}$}
\BIC{$\varphi_{60} \to \varphi_{56} \to \mathrm{FLim}(f,c,\alpha) \to \varphi_{189}$}
\BIC{$\varphi_{56} \to \mathrm{FLim}(f,c,\alpha) \to \varphi_{189}$}
\BIC{$\mathrm{FLim}(f,c,\alpha) \to \varphi_{189}$}
\end{ptree}
\noindent$\Pi_{44}$:
\begin{ptree}
\AXC{$\varphi_{71}$\lbl{A2}}
\AXC{$\varphi_{205}$\lbl{A1}}
\AXC{$\varphi_{205} \to \mathrm{FLim}(f,c,\alpha) \to \varphi_{202}$\lbl{A1}}
\BIC{$\mathrm{FLim}(f,c,\alpha) \to \varphi_{202}$}
\AXC{$\varphi_{246}$\lbl{A1}}
\BIC{$\varphi_{71} \to \mathrm{FLim}(f,c,\alpha) \to \varphi_{201}$}
\BIC{$\mathrm{FLim}(f,c,\alpha) \to \varphi_{201}$}
\end{ptree}
\noindent$\Pi_{45}$:
\begin{ptree}
\AXC{$\stackrel{\Pi_{24}}{\vdots}$}\noLine
\UIC{$\varphi_{82}$}
\AXC{$\stackrel{\Pi_{14}}{\vdots}$}\noLine
\UIC{$\varphi_{183}$}
\AXC{$\forall y\,\forall c\,(y \in \mathbb{R} \to c \in \mathbb{R} \to 0 < \lvert y - c \rvert \to y \neq c)$\lbl{[E37]}}
\AXC{$\varphi_{70}$\lbl{A2}}
\BIC{$\forall c\,(a_{n} \in \mathbb{R} \to c \in \mathbb{R} \to 0 < \lvert a_{n} - c \rvert \to a_{n} \neq c)$}
\AXC{$\varphi_{77}$\lbl{A2}}
\BIC{$a_{n} \in \mathbb{R} \to c \in \mathbb{R} \to 0 < \lvert a_{n} - c \rvert \to a_{n} \neq c$}
\AXC{$\varphi_{262}$\lbl{A1}}
\BIC{$\varphi_{183} \to \varphi_{82} \to (\mathrm{FLim}(f,c,\alpha) \to \varphi_{176}) \to \varphi_{181}$}
\BIC{$\varphi_{82} \to (\mathrm{FLim}(f,c,\alpha) \to \varphi_{176}) \to \varphi_{181}$}
\BIC{$(\mathrm{FLim}(f,c,\alpha) \to \varphi_{176}) \to \varphi_{181}$}
\end{ptree}
\noindent$\Pi_{46}$:
\begin{ptree}
\AXC{$\stackrel{\Pi_{44}}{\vdots}$}\noLine
\UIC{$\mathrm{FLim}(f,c,\alpha) \to \varphi_{201}$}
\AXC{$(\mathrm{FLim}(f,c,\alpha) \to \varphi_{201}) \to \mathrm{FLim}(f,c,\alpha) \to \varphi_{176}$\lbl{A1}}
\BIC{$\mathrm{FLim}(f,c,\alpha) \to \varphi_{176}$}
\AXC{$\stackrel{\Pi_{45}}{\vdots}$}\noLine
\UIC{$(\mathrm{FLim}(f,c,\alpha) \to \varphi_{176}) \to \varphi_{181}$}
\BIC{$\varphi_{181}$}
\end{ptree}
\noindent$\Pi_{47}$:
\begin{ptree}
\AXC{$\forall x\,\forall y\,\forall z\,(x < y \to y \leq z \to x < z)$\lbl{[O3]}}
\AXC{$\varphi_{46}$\lbl{A2}}
\BIC{$\forall y\,\forall z\,(\lvert a_{n} - c \rvert < y \to y \leq z \to \lvert a_{n} - c \rvert < z)$}
\AXC{$\varphi_{68}$\lbl{A2}}
\BIC{$\forall z\,(\lvert a_{n} - c \rvert < \rho \to \rho \leq z \to \lvert a_{n} - c \rvert < z)$}
\AXC{$\varphi_{65}$\lbl{A2}}
\BIC{$\lvert a_{n} - c \rvert < \rho \to \rho \leq d \to \lvert a_{n} - c \rvert < d$}
\end{ptree}
\noindent$\Pi_{48}$:
\begin{ptree}
\AXC{$\stackrel{\Pi_{47}}{\vdots}$}\noLine
\UIC{$\lvert a_{n} - c \rvert < \rho \to \rho \leq d \to \lvert a_{n} - c \rvert < d$}
\AXC{$(\lvert a_{n} - c \rvert < \rho \to \rho \leq d \to \lvert a_{n} - c \rvert < d) \to \varphi_{250}$\lbl{A1}}
\BIC{$\varphi_{250}$}
\end{ptree}
\noindent$\Pi_{49}$:
\begin{ptree}
\AXC{$\stackrel{\Pi_{44}}{\vdots}$}\noLine
\UIC{$\mathrm{FLim}(f,c,\alpha) \to \varphi_{201}$}
\AXC{$(\mathrm{FLim}(f,c,\alpha) \to \varphi_{201}) \to \mathrm{FLim}(f,c,\alpha) \to \varphi_{180}$\lbl{A1}}
\BIC{$\mathrm{FLim}(f,c,\alpha) \to \varphi_{180}$}
\AXC{$\stackrel{\Pi_{48}}{\vdots}$}\noLine
\UIC{$\varphi_{250}$}
\BIC{$\varphi_{171} \to \mathrm{FLim}(f,c,\alpha) \to \varphi_{175}$}
\end{ptree}
\noindent$\Pi_{50}$:
\begin{ptree}
\AXC{$\mathrm{FLim}(f,c,\alpha) \to \varphi_{154}$\lbl{A1}}
\AXC{$(\mathrm{FLim}(f,c,\alpha) \to \varphi_{154}) \to \varphi_{171}$\lbl{A1}}
\BIC{$\varphi_{171}$}
\AXC{$\stackrel{\Pi_{49}}{\vdots}$}\noLine
\UIC{$\varphi_{171} \to \mathrm{FLim}(f,c,\alpha) \to \varphi_{175}$}
\BIC{$\mathrm{FLim}(f,c,\alpha) \to \varphi_{175}$}
\end{ptree}
\noindent$\Pi_{51}$:
\begin{ptree}
\AXC{$\mathrm{FLim}(f,c,\alpha) \to \mathrm{FLim}(g,c,\beta) \to \varphi_{158}$\lbl{A1}}
\AXC{$(\mathrm{FLim}(f,c,\alpha) \to \mathrm{FLim}(g,c,\beta) \to \varphi_{158}) \to \varphi_{206}$\lbl{A1}}
\BIC{$\varphi_{206}$}
\end{ptree}
\noindent$\Pi_{52}$:
\begin{ptree}
\AXC{$\varphi_{31} \to a_{n} \in \mathbb{R} \to \varphi_{45}$\lbl{A2}}
\AXC{$\stackrel{\Pi_{51}}{\vdots}$}\noLine
\UIC{$\varphi_{206}$}
\AXC{$\varphi_{206} \to (\varphi_{31} \to a_{n} \in \mathbb{R} \to \varphi_{45}) \to \mathrm{FLim}(f,c,\alpha) \to \varphi_{209}$\lbl{A1}}
\BIC{$(\varphi_{31} \to a_{n} \in \mathbb{R} \to \varphi_{45}) \to \mathrm{FLim}(f,c,\alpha) \to \varphi_{209}$}
\BIC{$\mathrm{FLim}(f,c,\alpha) \to \varphi_{209}$}
\end{ptree}
\noindent$\Pi_{53}$:
\begin{ptree}
\AXC{$\stackrel{\Pi_{44}}{\vdots}$}\noLine
\UIC{$\mathrm{FLim}(f,c,\alpha) \to \varphi_{201}$}
\AXC{$\stackrel{\Pi_{14}}{\vdots}$}\noLine
\UIC{$\varphi_{183}$}
\AXC{$\stackrel{\Pi_{52}}{\vdots}$}\noLine
\UIC{$\mathrm{FLim}(f,c,\alpha) \to \varphi_{209}$}
\AXC{$(\mathrm{FLim}(f,c,\alpha) \to \varphi_{209}) \to \varphi_{183} \to \varphi_{260}$\lbl{A1}}
\BIC{$\varphi_{183} \to \varphi_{260}$}
\BIC{$\varphi_{260}$}
\BIC{$(\mathrm{FLim}(f,c,\alpha) \to \varphi_{175}) \to \mathrm{FLim}(f,c,\alpha) \to \varphi_{195}$}
\end{ptree}
\noindent$\Pi_{54}$:
\begin{ptree}
\AXC{$\stackrel{\Pi_{46}}{\vdots}$}\noLine
\UIC{$\varphi_{181}$}
\AXC{$\stackrel{\Pi_{50}}{\vdots}$}\noLine
\UIC{$\mathrm{FLim}(f,c,\alpha) \to \varphi_{175}$}
\AXC{$\stackrel{\Pi_{53}}{\vdots}$}\noLine
\UIC{$(\mathrm{FLim}(f,c,\alpha) \to \varphi_{175}) \to \mathrm{FLim}(f,c,\alpha) \to \varphi_{195}$}
\BIC{$\mathrm{FLim}(f,c,\alpha) \to \varphi_{195}$}
\AXC{$\varphi_{256}$\lbl{A1}}
\BIC{$\varphi_{181} \to \mathrm{FLim}(f,c,\alpha) \to \varphi_{186}$}
\BIC{$\mathrm{FLim}(f,c,\alpha) \to \varphi_{186}$}
\AXC{$\varphi_{235}$\lbl{A1}}
\BIC{$\varphi_{162} \to n \in \mathbb{N} \to a_{n} \in \mathrm{dom}(g) \land a_{n} \neq c$}
\UIC{$\forall n\,(\varphi_{162} \to n \in \mathbb{N} \to a_{n} \in \mathrm{dom}(g) \land a_{n} \neq c)$}
\end{ptree}
\noindent$\Pi_{55}$:
\begin{ptree}
\AXC{$\varphi_{162} \to \forall n\,(\varphi_{162})$\lbl{A4}}
\AXC{$\varphi_{237}$\lbl{A3}}
\AXC{$\stackrel{\Pi_{54}}{\vdots}$}\noLine
\UIC{$\forall n\,(\varphi_{162} \to n \in \mathbb{N} \to a_{n} \in \mathrm{dom}(g) \land a_{n} \neq c)$}
\AXC{$\varphi_{274}$\lbl{A1}}
\BIC{$\varphi_{237} \to (\varphi_{162} \to \forall n\,(\varphi_{162})) \to \mathrm{FLim}(f,c,\alpha) \to \varphi_{193}$}
\BIC{$(\varphi_{162} \to \forall n\,(\varphi_{162})) \to \mathrm{FLim}(f,c,\alpha) \to \varphi_{193}$}
\BIC{$\mathrm{FLim}(f,c,\alpha) \to \varphi_{193}$}
\end{ptree}
\noindent$\Pi_{56}$:
\begin{ptree}
\AXC{$\varphi_{232}$\lbl{A3}}
\AXC{$\varphi_{177}$\lbl{A1}}
\AXC{$\varphi_{48}$\lbl{A2}}
\AXC{$\stackrel{\Pi_{55}}{\vdots}$}\noLine
\UIC{$\mathrm{FLim}(f,c,\alpha) \to \varphi_{193}$}
\AXC{$\varphi_{243}$\lbl{A1}}
\BIC{$\varphi_{48} \to \mathrm{FLim}(f,c,\alpha) \to \varphi_{186}$}
\BIC{$\mathrm{FLim}(f,c,\alpha) \to \varphi_{186}$}
\AXC{$\varphi_{253}$\lbl{A1}}
\BIC{$\varphi_{177} \to \mathrm{FLim}(f,c,\alpha) \to \varphi_{169}$}
\BIC{$\mathrm{FLim}(f,c,\alpha) \to \varphi_{169}$}
\AXC{$\varphi_{230}$\lbl{A1}}
\BIC{$\varphi_{162} \to n \in \mathbb{N} \to a_{n} \in \mathrm{dom}(g)$}
\UIC{$\forall n\,(\varphi_{162} \to n \in \mathbb{N} \to a_{n} \in \mathrm{dom}(g))$}
\AXC{$\varphi_{272}$\lbl{A1}}
\BIC{$\varphi_{232} \to (\varphi_{162} \to \forall n\,(\varphi_{162})) \to \mathrm{FLim}(f,c,\alpha) \to \varphi_{173}$}
\BIC{$(\varphi_{162} \to \forall n\,(\varphi_{162})) \to \mathrm{FLim}(f,c,\alpha) \to \varphi_{173}$}
\end{ptree}
\noindent$\Pi_{57}$:
\begin{ptree}
\AXC{$\forall f\,\forall c\,\forall \alpha\,(\mathrm{FLim}(f,c,\alpha) \to \mathrm{RFn}(f))$\lbl{[D2]}}
\AXC{$\varphi_{13}$\lbl{A2}}
\BIC{$\forall c\,\forall \alpha\,(\mathrm{FLim}(g,c,\alpha) \to \mathrm{RFn}(g))$}
\AXC{$\varphi_{6}$\lbl{A2}}
\BIC{$\forall \alpha\,(\mathrm{FLim}(g,c,\alpha) \to \mathrm{RFn}(g))$}
\AXC{$\varphi_{3}$\lbl{A2}}
\BIC{$\mathrm{FLim}(g,c,\beta) \to \mathrm{RFn}(g)$}
\end{ptree}
\noindent$\Pi_{58}$:
\begin{ptree}
\AXC{$\varphi_{1}$\lbl{A1}}
\AXC{$\stackrel{\Pi_{57}}{\vdots}$}\noLine
\UIC{$\mathrm{FLim}(g,c,\beta) \to \mathrm{RFn}(g)$}
\AXC{$(\mathrm{FLim}(g,c,\beta) \to \mathrm{RFn}(g)) \to \varphi_{1} \to \varphi_{80}$\lbl{A1}}
\BIC{$\varphi_{1} \to \varphi_{80}$}
\BIC{$\varphi_{80}$}
\end{ptree}
\noindent$\Pi_{59}$:
\begin{ptree}
\AXC{$\varphi_{205}$\lbl{A1}}
\AXC{$\varphi_{205} \to \varphi_{165}$\lbl{A1}}
\BIC{$\varphi_{165}$}
\AXC{$\stackrel{\Pi_{58}}{\vdots}$}\noLine
\UIC{$\varphi_{80}$}
\AXC{$\forall f\,\forall a\,(\mathrm{RFn}(f) \to \mathrm{Seq}(a) \to \varphi_{57})$\lbl{[S8]}}
\AXC{$\varphi_{145}$\lbl{A2}}
\BIC{$\forall a\,(\mathrm{RFn}(g) \to \mathrm{Seq}(a) \to \varphi_{58})$}
\AXC{$\varphi_{138}$\lbl{A2}}
\BIC{$\mathrm{RFn}(g) \to \mathrm{Seq}(a) \to \varphi_{58}$}
\AXC{$\varphi_{265}$\lbl{A1}}
\BIC{$\varphi_{80} \to \varphi_{165} \to (\mathrm{FLim}(f,c,\alpha) \to \varphi_{173}) \to \varphi_{204}$}
\BIC{$\varphi_{165} \to (\mathrm{FLim}(f,c,\alpha) \to \varphi_{173}) \to \varphi_{204}$}
\BIC{$(\mathrm{FLim}(f,c,\alpha) \to \varphi_{173}) \to \varphi_{204}$}
\end{ptree}
\noindent$\Pi_{60}$:
\begin{ptree}
\AXC{$\varphi_{162} \to \forall n\,(\varphi_{162})$\lbl{A4}}
\AXC{$\stackrel{\Pi_{56}}{\vdots}$}\noLine
\UIC{$(\varphi_{162} \to \forall n\,(\varphi_{162})) \to \mathrm{FLim}(f,c,\alpha) \to \varphi_{173}$}
\BIC{$\mathrm{FLim}(f,c,\alpha) \to \varphi_{173}$}
\AXC{$\stackrel{\Pi_{59}}{\vdots}$}\noLine
\UIC{$(\mathrm{FLim}(f,c,\alpha) \to \varphi_{173}) \to \varphi_{204}$}
\BIC{$\varphi_{204}$}
\end{ptree}
\noindent$\Pi_{61}$:
\begin{ptree}
\AXC{$\varphi_{177}$\lbl{A1}}
\AXC{$\varphi_{53}$\lbl{A2}}
\AXC{$\stackrel{\Pi_{60}}{\vdots}$}\noLine
\UIC{$\varphi_{204}$}
\AXC{$\varphi_{204} \to \mathrm{FLim}(f,c,\alpha) \to \varphi_{200}$\lbl{A1}}
\BIC{$\mathrm{FLim}(f,c,\alpha) \to \varphi_{200}$}
\AXC{$\varphi_{245}$\lbl{A1}}
\BIC{$\varphi_{53} \to \mathrm{FLim}(f,c,\alpha) \to \varphi_{198}$}
\BIC{$\mathrm{FLim}(f,c,\alpha) \to \varphi_{198}$}
\AXC{$\varphi_{258}$\lbl{A1}}
\BIC{$\varphi_{177} \to \mathrm{FLim}(f,c,\alpha) \to \varphi_{190}$}
\BIC{$\mathrm{FLim}(f,c,\alpha) \to \varphi_{190}$}
\end{ptree}
\noindent$\Pi_{62}$:
\begin{ptree}
\AXC{$\stackrel{\Pi_{21}}{\vdots}$}\noLine
\UIC{$\forall x\,\forall y\,(x = y \to y = x)$}
\AXC{$\varphi_{19}$\lbl{A2}}
\BIC{$\forall y\,(g \circ a(n) = y \to y = g \circ a(n))$}
\AXC{$\varphi_{67}$\lbl{A2}}
\BIC{$g \circ a(n) = g(a_{n}) \to g(a_{n}) = g \circ a(n)$}
\AXC{$\stackrel{\Pi_{61}}{\vdots}$}\noLine
\UIC{$\mathrm{FLim}(f,c,\alpha) \to \varphi_{190}$}
\AXC{$(\mathrm{FLim}(f,c,\alpha) \to \varphi_{190}) \to \varphi_{211}$\lbl{A1}}
\BIC{$\varphi_{211}$}
\BIC{$\mathrm{FLim}(f,c,\alpha) \to \varphi_{191}$}
\end{ptree}
\noindent$\Pi_{63}$:
\begin{ptree}
\AXC{$\stackrel{\Pi_{1}}{\vdots}$}\noLine
\UIC{$\varphi_{62}$}
\AXC{$\stackrel{\Pi_{2}}{\vdots}$}\noLine
\UIC{$\varphi_{63}$}
\AXC{$\stackrel{\Pi_{43}}{\vdots}$}\noLine
\UIC{$\mathrm{FLim}(f,c,\alpha) \to \varphi_{189}$}
\AXC{$\stackrel{\Pi_{62}}{\vdots}$}\noLine
\UIC{$\mathrm{FLim}(f,c,\alpha) \to \varphi_{191}$}
\AXC{$(\mathrm{FLim}(f,c,\alpha) \to \varphi_{191}) \to \varphi_{248}$\lbl{A1}}
\BIC{$\varphi_{248}$}
\BIC{$\varphi_{63} \to \varphi_{62} \to \mathrm{FLim}(f,c,\alpha) \to \varphi_{194}$}
\BIC{$\varphi_{62} \to \mathrm{FLim}(f,c,\alpha) \to \varphi_{194}$}
\BIC{$\mathrm{FLim}(f,c,\alpha) \to \varphi_{194}$}
\end{ptree}
\noindent$\Pi_{64}$:
\begin{ptree}
\AXC{$\varphi_{162} \to \forall n\,(\varphi_{162})$\lbl{A4}}
\AXC{$\varphi_{239}$\lbl{A3}}
\AXC{$\stackrel{\Pi_{63}}{\vdots}$}\noLine
\UIC{$\mathrm{FLim}(f,c,\alpha) \to \varphi_{194}$}
\AXC{$\varphi_{238}$\lbl{A1}}
\BIC{$\varphi_{162} \to n \in \mathbb{N} \to f \circ a(n) \leq g \circ a(n)$}
\UIC{$\forall n\,(\varphi_{162} \to n \in \mathbb{N} \to f \circ a(n) \leq g \circ a(n))$}
\AXC{$\varphi_{275}$\lbl{A1}}
\BIC{$\varphi_{239} \to (\varphi_{162} \to \forall n\,(\varphi_{162})) \to \mathrm{FLim}(f,c,\alpha) \to \varphi_{196}$}
\BIC{$(\varphi_{162} \to \forall n\,(\varphi_{162})) \to \mathrm{FLim}(f,c,\alpha) \to \varphi_{196}$}
\BIC{$\mathrm{FLim}(f,c,\alpha) \to \varphi_{196}$}
\end{ptree}
\noindent$\Pi_{65}$:
\begin{ptree}
\AXC{$\forall f\,(\varphi_{59})$\lbl{定理1.7(⇒)}}
\AXC{$\forall f\,(\varphi_{59}) \to \forall c\,(\varphi_{55})$\lbl{A2}}
\BIC{$\forall c\,(\varphi_{55})$}
\AXC{$\forall c\,(\varphi_{55}) \to \varphi_{55}$\lbl{A2}}
\BIC{$\varphi_{55}$}
\AXC{$\varphi_{55} \to \forall a\,(\mathrm{FLim}(g,c,\beta) \to \mathrm{Conv}(a,c) \to \varphi_{24})$\lbl{A2}}
\BIC{$\forall a\,(\mathrm{FLim}(g,c,\beta) \to \mathrm{Conv}(a,c) \to \varphi_{24})$}
\end{ptree}
\noindent$\Pi_{66}$:
\begin{ptree}
\AXC{$\varphi_{205}$\lbl{A1}}
\AXC{$\varphi_{205} \to \varphi_{166}$\lbl{A1}}
\BIC{$\varphi_{166}$}
\AXC{$\varphi_{1}$\lbl{A1}}
\AXC{$\stackrel{\Pi_{65}}{\vdots}$}\noLine
\UIC{$\forall a\,(\mathrm{FLim}(g,c,\beta) \to \mathrm{Conv}(a,c) \to \varphi_{24})$}
\AXC{$\varphi_{115}$\lbl{A2}}
\BIC{$\mathrm{FLim}(g,c,\beta) \to \mathrm{Conv}(a,c) \to \varphi_{24}$}
\AXC{$\varphi_{261}$\lbl{A1}}
\BIC{$\varphi_{1} \to \varphi_{166} \to (\mathrm{FLim}(f,c,\alpha) \to \varphi_{193}) \to \varphi_{178}$}
\BIC{$\varphi_{166} \to (\mathrm{FLim}(f,c,\alpha) \to \varphi_{193}) \to \varphi_{178}$}
\BIC{$(\mathrm{FLim}(f,c,\alpha) \to \varphi_{193}) \to \varphi_{178}$}
\end{ptree}
\noindent$\Pi_{67}$:
\begin{ptree}
\AXC{$\forall f\,(\varphi_{59})$\lbl{定理1.7(⇒)}}
\AXC{$\forall f\,(\varphi_{59}) \to \varphi_{59}$\lbl{A2}}
\BIC{$\varphi_{59}$}
\AXC{$\varphi_{119}$\lbl{A2}}
\BIC{$\forall \alpha\,\forall a\,(\mathrm{FLim}(f,c,\alpha) \to \mathrm{Conv}(a,c) \to \varphi_{25})$}
\AXC{$\varphi_{118}$\lbl{A2}}
\BIC{$\forall a\,(\mathrm{FLim}(f,c,\alpha) \to \mathrm{Conv}(a,c) \to \varphi_{25})$}
\AXC{$\varphi_{116}$\lbl{A2}}
\BIC{$\mathrm{FLim}(f,c,\alpha) \to \mathrm{Conv}(a,c) \to \varphi_{25}$}
\end{ptree}
\noindent$\Pi_{68}$:
\begin{ptree}
\AXC{$\stackrel{\Pi_{67}}{\vdots}$}\noLine
\UIC{$\mathrm{FLim}(f,c,\alpha) \to \mathrm{Conv}(a,c) \to \varphi_{25}$}
\AXC{$(\mathrm{FLim}(f,c,\alpha) \to \mathrm{Conv}(a,c) \to \varphi_{25}) \to \varphi_{254}$\lbl{A1}}
\BIC{$\varphi_{254}$}
\end{ptree}
\noindent$\Pi_{69}$:
\begin{ptree}
\AXC{$\stackrel{\Pi_{35}}{\vdots}$}\noLine
\UIC{$\mathrm{FLim}(f,c,\alpha) \to \varphi_{192}$}
\AXC{$\varphi_{205}$\lbl{A1}}
\AXC{$\varphi_{205} \to \varphi_{166}$\lbl{A1}}
\BIC{$\varphi_{166}$}
\AXC{$\mathrm{FLim}(f,c,\alpha) \to \mathrm{FLim}(f,c,\alpha)$\lbl{A1}}
\AXC{$\stackrel{\Pi_{68}}{\vdots}$}\noLine
\UIC{$\varphi_{254}$}
\BIC{$\varphi_{166} \to (\mathrm{FLim}(f,c,\alpha) \to \varphi_{192}) \to \varphi_{179}$}
\BIC{$(\mathrm{FLim}(f,c,\alpha) \to \varphi_{192}) \to \varphi_{179}$}
\BIC{$\varphi_{179}$}
\end{ptree}
\noindent$\Pi_{70}$:
\begin{ptree}
\AXC{$\stackrel{\Pi_{69}}{\vdots}$}\noLine
\UIC{$\varphi_{179}$}
\AXC{$\varphi_{29}$\lbl{定理2.5}}
\AXC{$\varphi_{29} \to \varphi_{51}$\lbl{A2}}
\BIC{$\varphi_{51}$}
\AXC{$\varphi_{51} \to \forall \alpha\,\forall \beta\,(\mathrm{Conv}(f \circ a,\alpha) \to \varphi_{38})$\lbl{A2}}
\BIC{$\forall \alpha\,\forall \beta\,(\mathrm{Conv}(f \circ a,\alpha) \to \varphi_{38})$}
\AXC{$\varphi_{124}$\lbl{A2}}
\BIC{$\forall \beta\,(\mathrm{Conv}(f \circ a,\alpha) \to \varphi_{38})$}
\AXC{$\varphi_{121}$\lbl{A2}}
\BIC{$\mathrm{Conv}(f \circ a,\alpha) \to \varphi_{38}$}
\AXC{$\varphi_{267}$\lbl{A1}}
\BIC{$\varphi_{179} \to \varphi_{178} \to (\mathrm{FLim}(f,c,\alpha) \to \varphi_{196}) \to \varphi_{164}$}
\BIC{$\varphi_{178} \to (\mathrm{FLim}(f,c,\alpha) \to \varphi_{196}) \to \varphi_{164}$}
\end{ptree}
\noindent$\Pi_{71}$:
\begin{ptree}
\AXC{$\stackrel{\Pi_{55}}{\vdots}$}\noLine
\UIC{$\mathrm{FLim}(f,c,\alpha) \to \varphi_{193}$}
\AXC{$\stackrel{\Pi_{66}}{\vdots}$}\noLine
\UIC{$(\mathrm{FLim}(f,c,\alpha) \to \varphi_{193}) \to \varphi_{178}$}
\BIC{$\varphi_{178}$}
\AXC{$\stackrel{\Pi_{70}}{\vdots}$}\noLine
\UIC{$\varphi_{178} \to (\mathrm{FLim}(f,c,\alpha) \to \varphi_{196}) \to \varphi_{164}$}
\BIC{$(\mathrm{FLim}(f,c,\alpha) \to \varphi_{196}) \to \varphi_{164}$}
\end{ptree}
\noindent$\Pi_{72}$:
\begin{ptree}
\AXC{$\forall c\,(\varphi_{69})$\lbl{[H1]}}
\AXC{$\forall c\,(\varphi_{69}) \to \varphi_{69}$\lbl{A2}}
\BIC{$\varphi_{69}$}
\AXC{$\varphi_{69} \to c \in \mathbb{R} \to \rho \in \mathbb{R} \to 0 < \rho \to \exists a\,(\varphi_{39})$\lbl{A2}}
\BIC{$c \in \mathbb{R} \to \rho \in \mathbb{R} \to 0 < \rho \to \exists a\,(\varphi_{39})$}
\end{ptree}
\noindent$\Pi_{73}$:
\begin{ptree}
\AXC{$\stackrel{\Pi_{24}}{\vdots}$}\noLine
\UIC{$\varphi_{82}$}
\AXC{$\stackrel{\Pi_{72}}{\vdots}$}\noLine
\UIC{$c \in \mathbb{R} \to \rho \in \mathbb{R} \to 0 < \rho \to \exists a\,(\varphi_{39})$}
\AXC{$(c \in \mathbb{R} \to \rho \in \mathbb{R} \to 0 < \rho \to \exists a\,(\varphi_{39})) \to \varphi_{82} \to \varphi_{153} \to \varphi_{152} \to \varphi_{163}$\lbl{A1}}
\BIC{$\varphi_{82} \to \varphi_{153} \to \varphi_{152} \to \varphi_{163}$}
\BIC{$\varphi_{153} \to \varphi_{152} \to \varphi_{163}$}
\end{ptree}
\noindent$\Pi_{74}$:
\begin{ptree}
\AXC{$\mathrm{FLim}(f,c,\alpha) \to \varphi_{154}$\lbl{A1}}
\AXC{$(\mathrm{FLim}(f,c,\alpha) \to \varphi_{154}) \to \varphi_{153}$\lbl{A1}}
\BIC{$\varphi_{153}$}
\AXC{$\stackrel{\Pi_{73}}{\vdots}$}\noLine
\UIC{$\varphi_{153} \to \varphi_{152} \to \varphi_{163}$}
\BIC{$\varphi_{152} \to \varphi_{163}$}
\end{ptree}
\noindent$\Pi_{75}$:
\begin{ptree}
\AXC{$\mathrm{FLim}(f,c,\alpha) \to \varphi_{154}$\lbl{A1}}
\AXC{$(\mathrm{FLim}(f,c,\alpha) \to \varphi_{154}) \to \varphi_{152}$\lbl{A1}}
\BIC{$\varphi_{152}$}
\AXC{$\stackrel{\Pi_{74}}{\vdots}$}\noLine
\UIC{$\varphi_{152} \to \varphi_{163}$}
\BIC{$\varphi_{163}$}
\AXC{$\varphi_{163} \to \forall a\,(\varphi_{151} \to \neg \varphi_{39}) \to \varphi_{266}$\lbl{A1}}
\BIC{$\forall a\,(\varphi_{151} \to \neg \varphi_{39}) \to \varphi_{266}$}
\end{ptree}
\noindent$\Pi_{76}$:
\begin{ptree}
\AXC{$\stackrel{\Pi_{64}}{\vdots}$}\noLine
\UIC{$\mathrm{FLim}(f,c,\alpha) \to \varphi_{196}$}
\AXC{$\stackrel{\Pi_{71}}{\vdots}$}\noLine
\UIC{$(\mathrm{FLim}(f,c,\alpha) \to \varphi_{196}) \to \varphi_{164}$}
\BIC{$\varphi_{164}$}
\AXC{$\varphi_{164} \to \varphi_{151} \to \neg \varphi_{39}$\lbl{A1}}
\BIC{$\varphi_{151} \to \neg \varphi_{39}$}
\UIC{$\forall a\,(\varphi_{151} \to \neg \varphi_{39})$}
\AXC{$\stackrel{\Pi_{75}}{\vdots}$}\noLine
\UIC{$\forall a\,(\varphi_{151} \to \neg \varphi_{39}) \to \varphi_{266}$}
\BIC{$\varphi_{266}$}
\end{ptree}
\noindent$\Pi_{77}$:
\begin{ptree}
\AXC{$\varphi_{151} \to \forall a\,(\varphi_{151})$\lbl{A4}}
\AXC{$\forall a\,(\varphi_{151} \to \neg \varphi_{39}) \to \forall a\,(\varphi_{151}) \to \forall a\,(\neg \varphi_{39})$\lbl{A3}}
\AXC{$\stackrel{\Pi_{76}}{\vdots}$}\noLine
\UIC{$\varphi_{266}$}
\BIC{$(\varphi_{151} \to \forall a\,(\varphi_{151})) \to \varphi_{150}$}
\BIC{$\varphi_{150}$}
\AXC{$\varphi_{150} \to \varphi_{133} \to \neg (\rho \in \mathbb{R} \land 0 < \rho \land \rho \leq d \land \rho \leq r)$\lbl{A1}}
\BIC{$\varphi_{133} \to \neg (\rho \in \mathbb{R} \land 0 < \rho \land \rho \leq d \land \rho \leq r)$}
\UIC{$\forall \rho\,(\varphi_{133} \to \neg (\rho \in \mathbb{R} \land 0 < \rho \land \rho \leq d \land \rho \leq r))$}
\end{ptree}
\noindent$\Pi_{78}$:
\begin{ptree}
\AXC{$\forall d\,(\forall u_{5}\,(\varphi_{27}) \to \varphi_{23}) \to \forall d\,\forall u_{5}\,(\varphi_{27}) \to \forall d\,(\varphi_{23})$\lbl{A3}}
\AXC{$\forall u_{5}\,(\varphi_{27}) \to \varphi_{23}$\lbl{A2}}
\UIC{$\forall d\,(\forall u_{5}\,(\varphi_{27}) \to \varphi_{23})$}
\AXC{$\forall d\,(\forall u_{5}\,(\varphi_{27}) \to \varphi_{23}) \to \varphi_{227}$\lbl{A1}}
\BIC{$\varphi_{227}$}
\BIC{$(\forall u_{5}\,(\varphi_{27}) \to \forall d\,\forall u_{5}\,(\varphi_{27})) \to \forall u_{5}\,(\varphi_{27}) \to \forall d\,(\varphi_{23})$}
\end{ptree}
\noindent$\Pi_{79}$:
\begin{ptree}
\AXC{$\forall u_{5}\,(\varphi_{27}) \to \forall d\,\forall u_{5}\,(\varphi_{27})$\lbl{A4}}
\AXC{$\stackrel{\Pi_{78}}{\vdots}$}\noLine
\UIC{$(\forall u_{5}\,(\varphi_{27}) \to \forall d\,\forall u_{5}\,(\varphi_{27})) \to \forall u_{5}\,(\varphi_{27}) \to \forall d\,(\varphi_{23})$}
\BIC{$\forall u_{5}\,(\varphi_{27}) \to \forall d\,(\varphi_{23})$}
\end{ptree}
\noindent$\Pi_{80}$:
\begin{ptree}
\AXC{$\forall u_{5}\,(\forall d\,(\varphi_{23}) \to \varphi_{27}) \to \forall u_{5}\,\forall d\,(\varphi_{23}) \to \forall u_{5}\,(\varphi_{27})$\lbl{A3}}
\AXC{$\forall d\,(\varphi_{23}) \to \varphi_{27}$\lbl{A2}}
\UIC{$\forall u_{5}\,(\forall d\,(\varphi_{23}) \to \varphi_{27})$}
\AXC{$\forall u_{5}\,(\forall d\,(\varphi_{23}) \to \varphi_{27}) \to \varphi_{226}$\lbl{A1}}
\BIC{$\varphi_{226}$}
\BIC{$(\forall d\,(\varphi_{23}) \to \forall u_{5}\,\forall d\,(\varphi_{23})) \to \forall d\,(\varphi_{23}) \to \forall u_{5}\,(\varphi_{27})$}
\end{ptree}
\noindent$\Pi_{81}$:
\begin{ptree}
\AXC{$\forall d\,(\varphi_{23}) \to \forall u_{5}\,\forall d\,(\varphi_{23})$\lbl{A4}}
\AXC{$\stackrel{\Pi_{80}}{\vdots}$}\noLine
\UIC{$(\forall d\,(\varphi_{23}) \to \forall u_{5}\,\forall d\,(\varphi_{23})) \to \forall d\,(\varphi_{23}) \to \forall u_{5}\,(\varphi_{27})$}
\BIC{$\forall d\,(\varphi_{23}) \to \forall u_{5}\,(\varphi_{27})$}
\end{ptree}
\noindent$\Pi_{82}$:
\begin{ptree}
\AXC{$\stackrel{\Pi_{81}}{\vdots}$}\noLine
\UIC{$\forall d\,(\varphi_{23}) \to \forall u_{5}\,(\varphi_{27})$}
\AXC{$\forall d\,(\varphi_{23}) \leftrightarrow \forall d\,(\varphi_{23})$\lbl{A1}}
\AXC{$(\forall d\,(\varphi_{23}) \leftrightarrow \forall d\,(\varphi_{23})) \to \varphi_{207}$\lbl{A1}}
\BIC{$\varphi_{207}$}
\BIC{$(\forall u_{5}\,(\varphi_{27}) \to \forall d\,(\varphi_{23})) \to (\forall d\,(\varphi_{23}) \leftrightarrow \forall u_{5}\,(\varphi_{27}))$}
\end{ptree}
\noindent$\Pi_{83}$:
\begin{ptree}
\AXC{$\stackrel{\Pi_{79}}{\vdots}$}\noLine
\UIC{$\forall u_{5}\,(\varphi_{27}) \to \forall d\,(\varphi_{23})$}
\AXC{$\stackrel{\Pi_{82}}{\vdots}$}\noLine
\UIC{$(\forall u_{5}\,(\varphi_{27}) \to \forall d\,(\varphi_{23})) \to (\forall d\,(\varphi_{23}) \leftrightarrow \forall u_{5}\,(\varphi_{27}))$}
\BIC{$\forall d\,(\varphi_{23}) \leftrightarrow \forall u_{5}\,(\varphi_{27})$}
\AXC{$\varphi_{155}$\lbl{A1}}
\BIC{$\neg \forall d\,(\varphi_{23}) \leftrightarrow \neg \forall u_{5}\,(\varphi_{27})$}
\end{ptree}
\noindent$\Pi_{84}$:
\begin{ptree}
\AXC{$\stackrel{\Pi_{83}}{\vdots}$}\noLine
\UIC{$\neg \forall d\,(\varphi_{23}) \leftrightarrow \neg \forall u_{5}\,(\varphi_{27})$}
\AXC{$0 < v_{769} \leftrightarrow 0 < v_{769}$\lbl{A1}}
\AXC{$(0 < v_{769} \leftrightarrow 0 < v_{769}) \to \varphi_{161}$\lbl{A1}}
\BIC{$\varphi_{161}$}
\BIC{$0 < v_{769} \to \neg \forall d\,(\varphi_{23}) \leftrightarrow 0 < v_{769} \to \neg \forall u_{5}\,(\varphi_{27})$}
\end{ptree}
\noindent$\Pi_{85}$:
\begin{ptree}
\AXC{$\stackrel{\Pi_{84}}{\vdots}$}\noLine
\UIC{$0 < v_{769} \to \neg \forall d\,(\varphi_{23}) \leftrightarrow 0 < v_{769} \to \neg \forall u_{5}\,(\varphi_{27})$}
\AXC{$0 < v_{768} \leftrightarrow 0 < v_{768}$\lbl{A1}}
\AXC{$(0 < v_{768} \leftrightarrow 0 < v_{768}) \to \varphi_{167}$\lbl{A1}}
\BIC{$\varphi_{167}$}
\BIC{$\varphi_{17} \leftrightarrow 0 < v_{768} \to 0 < v_{769} \to \neg \forall u_{5}\,(\varphi_{27})$}
\end{ptree}
\noindent$\Pi_{86}$:
\begin{ptree}
\AXC{$\stackrel{\Pi_{85}}{\vdots}$}\noLine
\UIC{$\varphi_{17} \leftrightarrow 0 < v_{768} \to 0 < v_{769} \to \neg \forall u_{5}\,(\varphi_{27})$}
\AXC{$v_{769} \in \mathbb{R} \leftrightarrow v_{769} \in \mathbb{R}$\lbl{A1}}
\AXC{$\varphi_{208}$\lbl{A1}}
\BIC{$(\varphi_{17} \leftrightarrow 0 < v_{768} \to 0 < v_{769} \to \neg \forall u_{5}\,(\varphi_{27})) \to \varphi_{125}$}
\BIC{$\varphi_{125}$}
\end{ptree}
\noindent$\Pi_{87}$:
\begin{ptree}
\AXC{$\stackrel{\Pi_{86}}{\vdots}$}\noLine
\UIC{$\varphi_{125}$}
\AXC{$v_{768} \in \mathbb{R} \leftrightarrow v_{768} \in \mathbb{R}$\lbl{A1}}
\AXC{$\varphi_{214}$\lbl{A1}}
\BIC{$\varphi_{125} \to (v_{768} \in \mathbb{R} \to v_{769} \in \mathbb{R} \to \varphi_{17} \leftrightarrow \varphi_{73})$}
\BIC{$v_{768} \in \mathbb{R} \to v_{769} \in \mathbb{R} \to \varphi_{17} \leftrightarrow \varphi_{73}$}
\end{ptree}
\noindent$\Pi_{88}$:
\begin{ptree}
\AXC{$\varphi_{218}$\lbl{A3}}
\AXC{$\stackrel{\Pi_{87}}{\vdots}$}\noLine
\UIC{$v_{768} \in \mathbb{R} \to v_{769} \in \mathbb{R} \to \varphi_{17} \leftrightarrow \varphi_{73}$}
\AXC{$\varphi_{215}$\lbl{A1}}
\BIC{$\varphi_{73} \to v_{768} \in \mathbb{R} \to v_{769} \in \mathbb{R} \to \varphi_{17}$}
\UIC{$\forall v_{769}\,(\varphi_{73} \to v_{768} \in \mathbb{R} \to v_{769} \in \mathbb{R} \to \varphi_{17})$}
\AXC{$\stackrel{\Pi_{87}}{\vdots}$}\noLine
\UIC{$v_{768} \in \mathbb{R} \to v_{769} \in \mathbb{R} \to \varphi_{17} \leftrightarrow \varphi_{73}$}
\AXC{$\varphi_{216}$\lbl{A1}}
\BIC{$(v_{768} \in \mathbb{R} \to v_{769} \in \mathbb{R} \to \varphi_{17}) \to \varphi_{73}$}
\UIC{$\forall v_{769}\,((v_{768} \in \mathbb{R} \to v_{769} \in \mathbb{R} \to \varphi_{17}) \to \varphi_{73})$}
\AXC{$\varphi_{268}$\lbl{A1}}
\BIC{$\forall v_{769}\,(\varphi_{73} \to v_{768} \in \mathbb{R} \to v_{769} \in \mathbb{R} \to \varphi_{17}) \to \varphi_{218} \to \varphi_{233}$}
\BIC{$\varphi_{218} \to \varphi_{233}$}
\BIC{$\varphi_{233}$}
\end{ptree}
\noindent$\Pi_{89}$:
\begin{ptree}
\AXC{$\varphi_{217}$\lbl{A3}}
\AXC{$\stackrel{\Pi_{88}}{\vdots}$}\noLine
\UIC{$\varphi_{233}$}
\BIC{$\forall v_{769}\,(v_{768} \in \mathbb{R} \to v_{769} \in \mathbb{R} \to \varphi_{17}) \leftrightarrow \forall v_{769}\,(\varphi_{73})$}
\end{ptree}
\noindent$\Pi_{90}$:
\begin{ptree}
\AXC{$\stackrel{\Pi_{89}}{\vdots}$}\noLine
\UIC{$\forall v_{769}\,(v_{768} \in \mathbb{R} \to v_{769} \in \mathbb{R} \to \varphi_{17}) \leftrightarrow \forall v_{769}\,(\varphi_{73})$}
\AXC{$\varphi_{221}$\lbl{A1}}
\BIC{$\forall v_{769}\,(v_{768} \in \mathbb{R} \to v_{769} \in \mathbb{R} \to \varphi_{17}) \to \forall v_{769}\,(\varphi_{73})$}
\UIC{$\varphi_{148}$}
\AXC{$\varphi_{148} \to \varphi_{147} \to \varphi_{263}$\lbl{A1}}
\BIC{$\varphi_{147} \to \varphi_{263}$}
\end{ptree}
\noindent$\Pi_{91}$:
\begin{ptree}
\AXC{$\stackrel{\Pi_{89}}{\vdots}$}\noLine
\UIC{$\forall v_{769}\,(v_{768} \in \mathbb{R} \to v_{769} \in \mathbb{R} \to \varphi_{17}) \leftrightarrow \forall v_{769}\,(\varphi_{73})$}
\AXC{$\varphi_{220}$\lbl{A1}}
\BIC{$\forall v_{769}\,(\varphi_{73}) \to \forall v_{769}\,(v_{768} \in \mathbb{R} \to v_{769} \in \mathbb{R} \to \varphi_{17})$}
\UIC{$\varphi_{147}$}
\AXC{$\stackrel{\Pi_{90}}{\vdots}$}\noLine
\UIC{$\varphi_{147} \to \varphi_{263}$}
\BIC{$\varphi_{263}$}
\end{ptree}
\noindent$\Pi_{92}$:
\begin{ptree}
\AXC{$\varphi_{147} \to \varphi_{61} \to \forall v_{768}\,\forall v_{769}\,(v_{768} \in \mathbb{R} \to v_{769} \in \mathbb{R} \to \varphi_{17})$\lbl{A3}}
\AXC{$\varphi_{148} \to \forall v_{768}\,\forall v_{769}\,(v_{768} \in \mathbb{R} \to v_{769} \in \mathbb{R} \to \varphi_{17}) \to \varphi_{61}$\lbl{A3}}
\AXC{$\stackrel{\Pi_{91}}{\vdots}$}\noLine
\UIC{$\varphi_{263}$}
\BIC{$\varphi_{240}$}
\BIC{$\forall v_{768}\,\forall v_{769}\,(v_{768} \in \mathbb{R} \to v_{769} \in \mathbb{R} \to \varphi_{17}) \leftrightarrow \varphi_{61}$}
\end{ptree}
\noindent$\Pi_{93}$:
\begin{ptree}
\AXC{$\varphi_{37}$\lbl{[D8]}}
\AXC{$\varphi_{37} \to \varphi_{43}$\lbl{A2}}
\BIC{$\varphi_{43}$}
\AXC{$\varphi_{43} \to v_{768} \in \mathbb{R} \to v_{769} \in \mathbb{R} \to \varphi_{17}$\lbl{A2}}
\BIC{$v_{768} \in \mathbb{R} \to v_{769} \in \mathbb{R} \to \varphi_{17}$}
\UIC{$\forall v_{769}\,(v_{768} \in \mathbb{R} \to v_{769} \in \mathbb{R} \to \varphi_{17})$}
\UIC{$\forall v_{768}\,\forall v_{769}\,(v_{768} \in \mathbb{R} \to v_{769} \in \mathbb{R} \to \varphi_{17})$}
\AXC{$\forall v_{768}\,\forall v_{769}\,(v_{768} \in \mathbb{R} \to v_{769} \in \mathbb{R} \to \varphi_{17}) \to \varphi_{184}$\lbl{A1}}
\BIC{$\varphi_{184}$}
\end{ptree}
\noindent$\Pi_{94}$:
\begin{ptree}
\AXC{$\stackrel{\Pi_{92}}{\vdots}$}\noLine
\UIC{$\forall v_{768}\,\forall v_{769}\,(v_{768} \in \mathbb{R} \to v_{769} \in \mathbb{R} \to \varphi_{17}) \leftrightarrow \varphi_{61}$}
\AXC{$\stackrel{\Pi_{93}}{\vdots}$}\noLine
\UIC{$\varphi_{184}$}
\BIC{$\varphi_{61}$}
\AXC{$\varphi_{61} \to \varphi_{50}$\lbl{A2}}
\BIC{$\varphi_{50}$}
\AXC{$\varphi_{50} \to d \in \mathbb{R} \to r \in \mathbb{R} \to \varphi_{10}$\lbl{A2}}
\BIC{$d \in \mathbb{R} \to r \in \mathbb{R} \to \varphi_{10}$}
\end{ptree}
\noindent$\Pi_{95}$:
\begin{ptree}
\AXC{$\mathrm{FLim}(f,c,\alpha) \to \mathrm{FLim}(g,c,\beta) \to \varphi_{158}$\lbl{A1}}
\AXC{$\varphi_{224}$\lbl{A1}}
\BIC{$\mathrm{FLim}(f,c,\alpha) \to \varphi_{131}$}
\AXC{$\stackrel{\Pi_{94}}{\vdots}$}\noLine
\UIC{$d \in \mathbb{R} \to r \in \mathbb{R} \to \varphi_{10}$}
\AXC{$(d \in \mathbb{R} \to r \in \mathbb{R} \to \varphi_{10}) \to \varphi_{241}$\lbl{A1}}
\BIC{$\varphi_{241}$}
\BIC{$\varphi_{4} \to (\mathrm{FLim}(f,c,\alpha) \to \varphi_{128}) \to \varphi_{11} \to \varphi_{149}$}
\end{ptree}
\noindent$\Pi_{96}$:
\begin{ptree}
\AXC{$\mathrm{FLim}(f,c,\alpha) \to \mathrm{FLim}(g,c,\beta) \to \varphi_{158}$\lbl{A1}}
\AXC{$\varphi_{222}$\lbl{A1}}
\BIC{$\mathrm{FLim}(f,c,\alpha) \to \varphi_{128}$}
\AXC{$\varphi_{4}$\lbl{A1}}
\AXC{$\stackrel{\Pi_{95}}{\vdots}$}\noLine
\UIC{$\varphi_{4} \to (\mathrm{FLim}(f,c,\alpha) \to \varphi_{128}) \to \varphi_{11} \to \varphi_{149}$}
\BIC{$(\mathrm{FLim}(f,c,\alpha) \to \varphi_{128}) \to \varphi_{11} \to \varphi_{149}$}
\BIC{$\varphi_{11} \to \varphi_{149}$}
\end{ptree}
\noindent$\Pi_{97}$:
\begin{ptree}
\AXC{$\varphi_{223}$\lbl{A3}}
\AXC{$\stackrel{\Pi_{77}}{\vdots}$}\noLine
\UIC{$\forall \rho\,(\varphi_{133} \to \neg (\rho \in \mathbb{R} \land 0 < \rho \land \rho \leq d \land \rho \leq r))$}
\AXC{$\varphi_{11}$\lbl{A1}}
\AXC{$\stackrel{\Pi_{96}}{\vdots}$}\noLine
\UIC{$\varphi_{11} \to \varphi_{149}$}
\BIC{$\varphi_{149}$}
\AXC{$\varphi_{149} \to \varphi_{270}$\lbl{A1}}
\BIC{$\varphi_{270}$}
\BIC{$\varphi_{223} \to (\varphi_{133} \to \forall \rho\,(\varphi_{133})) \to \varphi_{225}$}
\BIC{$(\varphi_{133} \to \forall \rho\,(\varphi_{133})) \to \varphi_{225}$}
\end{ptree}
\noindent$\Pi_{98}$:
\begin{ptree}
\AXC{$\varphi_{49}$\lbl{A4}}
\AXC{$\forall u_{5}\,(\varphi_{36}) \to \varphi_{44}$\lbl{A3}}
\AXC{$\varphi_{36}$\lbl{A2}}
\UIC{$\forall u_{5}\,(\varphi_{36})$}
\AXC{$\varphi_{133} \to \forall \rho\,(\varphi_{133})$\lbl{A4}}
\AXC{$\stackrel{\Pi_{97}}{\vdots}$}\noLine
\UIC{$(\varphi_{133} \to \forall \rho\,(\varphi_{133})) \to \varphi_{225}$}
\BIC{$\varphi_{225}$}
\BIC{$(\forall u_{5}\,(\varphi_{36}) \to \varphi_{44}) \to \varphi_{49} \to \mathrm{FLim}(f,c,\alpha) \to \varphi_{127}$}
\BIC{$\varphi_{49} \to \mathrm{FLim}(f,c,\alpha) \to \varphi_{127}$}
\BIC{$\mathrm{FLim}(f,c,\alpha) \to \varphi_{127}$}
\AXC{$\varphi_{212}$\lbl{A1}}
\BIC{$\varphi_{81} \to \neg (d \in \mathbb{R} \land 0 < d \land \varphi_{31})$}
\UIC{$\forall d\,(\varphi_{81} \to \neg (d \in \mathbb{R} \land 0 < d \land \varphi_{31}))$}
\end{ptree}
\noindent$\Pi_{99}$:
\begin{ptree}
\AXC{$\varphi_{1}$\lbl{A1}}
\AXC{$\forall f\,\forall c\,\forall \alpha\,(\mathrm{FLim}(f,c,\alpha) \to \mathrm{PNbhd}(f,c))$\lbl{[D5]}}
\AXC{$\varphi_{21}$\lbl{A2}}
\BIC{$\forall c\,\forall \alpha\,(\mathrm{FLim}(g,c,\alpha) \to \mathrm{PNbhd}(g,c))$}
\AXC{$\varphi_{15}$\lbl{A2}}
\BIC{$\forall \alpha\,(\mathrm{FLim}(g,c,\alpha) \to \mathrm{PNbhd}(g,c))$}
\AXC{$\varphi_{8}$\lbl{A2}}
\BIC{$\mathrm{FLim}(g,c,\beta) \to \mathrm{PNbhd}(g,c)$}
\AXC{$\varphi_{120}$\lbl{A1}}
\BIC{$\varphi_{1} \to \mathrm{FLim}(f,c,\alpha) \to \varphi_{74}$}
\BIC{$\mathrm{FLim}(f,c,\alpha) \to \varphi_{74}$}
\end{ptree}
\noindent$\Pi_{100}$:
\begin{ptree}
\AXC{$\forall f\,\forall c\,\forall \alpha\,(\mathrm{FLim}(f,c,\alpha) \to \mathrm{PNbhd}(f,c))$\lbl{[D5]}}
\AXC{$\varphi_{22}$\lbl{A2}}
\BIC{$\forall c\,\forall \alpha\,(\mathrm{FLim}(f,c,\alpha) \to \mathrm{PNbhd}(f,c))$}
\AXC{$\varphi_{16}$\lbl{A2}}
\BIC{$\forall \alpha\,(\mathrm{FLim}(f,c,\alpha) \to \mathrm{PNbhd}(f,c))$}
\AXC{$\varphi_{9}$\lbl{A2}}
\BIC{$\mathrm{FLim}(f,c,\alpha) \to \mathrm{PNbhd}(f,c)$}
\end{ptree}
\noindent$\Pi_{101}$:
\begin{ptree}
\AXC{$\stackrel{\Pi_{100}}{\vdots}$}\noLine
\UIC{$\mathrm{FLim}(f,c,\alpha) \to \mathrm{PNbhd}(f,c)$}
\AXC{$(\mathrm{FLim}(f,c,\alpha) \to \mathrm{PNbhd}(f,c)) \to \varphi_{104}$\lbl{A1}}
\BIC{$\varphi_{104}$}
\end{ptree}
\noindent$\Pi_{102}$:
\begin{ptree}
\AXC{$\forall f\,\forall g\,\forall c\,(\varphi_{72})$\lbl{[H2]}}
\AXC{$\forall f\,\forall g\,\forall c\,(\varphi_{72}) \to \forall g\,\forall c\,(\varphi_{72})$\lbl{A2}}
\BIC{$\forall g\,\forall c\,(\varphi_{72})$}
\AXC{$\forall g\,\forall c\,(\varphi_{72}) \to \forall c\,(\varphi_{72})$\lbl{A2}}
\BIC{$\forall c\,(\varphi_{72})$}
\AXC{$\forall c\,(\varphi_{72}) \to \varphi_{72}$\lbl{A2}}
\BIC{$\varphi_{72}$}
\end{ptree}
\noindent$\Pi_{103}$:
\begin{ptree}
\AXC{$\stackrel{\Pi_{99}}{\vdots}$}\noLine
\UIC{$\mathrm{FLim}(f,c,\alpha) \to \varphi_{74}$}
\AXC{$\mathrm{FLim}(f,c,\alpha) \to \mathrm{FLim}(f,c,\alpha)$\lbl{A1}}
\AXC{$\stackrel{\Pi_{101}}{\vdots}$}\noLine
\UIC{$\varphi_{104}$}
\BIC{$\mathrm{FLim}(f,c,\alpha) \to \varphi_{75}$}
\AXC{$\stackrel{\Pi_{102}}{\vdots}$}\noLine
\UIC{$\varphi_{72}$}
\AXC{$\varphi_{72} \to \varphi_{219}$\lbl{A1}}
\BIC{$\varphi_{219}$}
\BIC{$(\mathrm{FLim}(f,c,\alpha) \to \varphi_{74}) \to \mathrm{FLim}(f,c,\alpha) \to \varphi_{123}$}
\BIC{$\mathrm{FLim}(f,c,\alpha) \to \varphi_{123}$}
\end{ptree}
\noindent$\Pi_{104}$:
\begin{ptree}
\AXC{$\varphi_{213}$\lbl{A3}}
\AXC{$\stackrel{\Pi_{98}}{\vdots}$}\noLine
\UIC{$\forall d\,(\varphi_{81} \to \neg (d \in \mathbb{R} \land 0 < d \land \varphi_{31}))$}
\AXC{$\stackrel{\Pi_{103}}{\vdots}$}\noLine
\UIC{$\mathrm{FLim}(f,c,\alpha) \to \varphi_{123}$}
\AXC{$(\mathrm{FLim}(f,c,\alpha) \to \varphi_{123}) \to \varphi_{251}$\lbl{A1}}
\BIC{$\varphi_{251}$}
\BIC{$\varphi_{213} \to (\varphi_{81} \to \forall d\,(\varphi_{81})) \to \varphi_{78}$}
\BIC{$(\varphi_{81} \to \forall d\,(\varphi_{81})) \to \varphi_{78}$}
\end{ptree}
\noindent 主木:
\begin{ptree}
\AXC{$\varphi_{81} \to \forall d\,(\varphi_{81})$\lbl{A4}}
\AXC{$\stackrel{\Pi_{104}}{\vdots}$}\noLine
\UIC{$(\varphi_{81} \to \forall d\,(\varphi_{81})) \to \varphi_{78}$}
\BIC{$\varphi_{78}$}
\UIC{$\forall r\,(\varphi_{78})$}
\UIC{$\forall \beta\,\forall r\,(\varphi_{78})$}
\UIC{$\forall \alpha\,\forall \beta\,\forall r\,(\varphi_{78})$}
\UIC{$\forall c\,\forall \alpha\,\forall \beta\,\forall r\,(\varphi_{78})$}
\UIC{$\forall g\,\forall c\,\forall \alpha\,\forall \beta\,\forall r\,(\varphi_{78})$}
\UIC{$\forall f\,\forall g\,\forall c\,\forall \alpha\,\forall \beta\,\forall r\,(\varphi_{78})$}
\end{ptree}
\begin{abbrevlist}
\item $\varphi_{1} :\equiv \mathrm{FLim}(f,c,\alpha) \to \mathrm{FLim}(g,c,\beta) \to \mathrm{FLim}(g,c,\beta)$
\item $\varphi_{2} :\equiv \forall \alpha\,(\mathrm{FLim}(f,c,\alpha) \to \mathrm{RFn}(f)) \to \mathrm{FLim}(f,c,\alpha) \to \mathrm{RFn}(f)$
\item $\varphi_{3} :\equiv \forall \alpha\,(\mathrm{FLim}(g,c,\alpha) \to \mathrm{RFn}(g)) \to \mathrm{FLim}(g,c,\beta) \to \mathrm{RFn}(g)$
\item $\varphi_{4} :\equiv \mathrm{FLim}(f,c,\alpha) \to \mathrm{FLim}(g,c,\beta) \to r \in \mathbb{R} \to r \in \mathbb{R}$
\item $\varphi_{5} :\equiv \forall c\,\forall \alpha\,(\mathrm{FLim}(f,c,\alpha) \to \mathrm{RFn}(f)) \to \forall \alpha\,(\mathrm{FLim}(f,c,\alpha) \to \mathrm{RFn}(f))$
\item $\varphi_{6} :\equiv \forall c\,\forall \alpha\,(\mathrm{FLim}(g,c,\alpha) \to \mathrm{RFn}(g)) \to \forall \alpha\,(\mathrm{FLim}(g,c,\alpha) \to \mathrm{RFn}(g))$
\item $\varphi_{7} :\equiv \forall \alpha\,(\mathrm{FLim}(f,c,\alpha) \to c \in \mathbb{R}) \to \mathrm{FLim}(f,c,\alpha) \to c \in \mathbb{R}$
\item $\varphi_{8} :\equiv \forall \alpha\,(\mathrm{FLim}(g,c,\alpha) \to \mathrm{PNbhd}(g,c)) \to \mathrm{FLim}(g,c,\beta) \to \mathrm{PNbhd}(g,c)$
\item $\varphi_{9} :\equiv \forall \alpha\,(\mathrm{FLim}(f,c,\alpha) \to \mathrm{PNbhd}(f,c)) \to \mathrm{FLim}(f,c,\alpha) \to \mathrm{PNbhd}(f,c)$
\item $\varphi_{10} :\equiv 0 < d \to 0 < r \to \exists u_{5}\,(u_{5} \in \mathbb{R} \land 0 < u_{5} \land u_{5} \leq d \land u_{5} \leq r)$
\item $\varphi_{11} :\equiv \mathrm{FLim}(f,c,\alpha) \to \mathrm{FLim}(g,c,\beta) \to r \in \mathbb{R} \to 0 < r \to 0 < r$
\item $\varphi_{12} :\equiv \forall f\,\forall c\,\forall \alpha\,(\mathrm{FLim}(f,c,\alpha) \to \mathrm{RFn}(f)) \to \forall c\,\forall \alpha\,(\mathrm{FLim}(f,c,\alpha) \to \mathrm{RFn}(f))$
\item $\varphi_{13} :\equiv \forall f\,\forall c\,\forall \alpha\,(\mathrm{FLim}(f,c,\alpha) \to \mathrm{RFn}(f)) \to \forall c\,\forall \alpha\,(\mathrm{FLim}(g,c,\alpha) \to \mathrm{RFn}(g))$
\item $\varphi_{14} :\equiv \forall c\,\forall \alpha\,(\mathrm{FLim}(f,c,\alpha) \to c \in \mathbb{R}) \to \forall \alpha\,(\mathrm{FLim}(f,c,\alpha) \to c \in \mathbb{R})$
\item $\varphi_{15} :\equiv \forall c\,\forall \alpha\,(\mathrm{FLim}(g,c,\alpha) \to \mathrm{PNbhd}(g,c)) \to \forall \alpha\,(\mathrm{FLim}(g,c,\alpha) \to \mathrm{PNbhd}(g,c))$
\item $\varphi_{16} :\equiv \forall c\,\forall \alpha\,(\mathrm{FLim}(f,c,\alpha) \to \mathrm{PNbhd}(f,c)) \to \forall \alpha\,(\mathrm{FLim}(f,c,\alpha) \to \mathrm{PNbhd}(f,c))$
\item $\varphi_{17} :\equiv 0 < v_{768} \to 0 < v_{769} \to \exists d\,(d \in \mathbb{R} \land 0 < d \land d \leq v_{768} \land d \leq v_{769})$
\item $\varphi_{18} :\equiv \forall x\,\forall y\,(x = y \to y = x) \to \forall y\,(f \circ a(n) = y \to y = f \circ a(n))$
\item $\varphi_{19} :\equiv \forall x\,\forall y\,(x = y \to y = x) \to \forall y\,(g \circ a(n) = y \to y = g \circ a(n))$
\item $\varphi_{20} :\equiv \forall f\,\forall c\,\forall \alpha\,(\mathrm{FLim}(f,c,\alpha) \to c \in \mathbb{R}) \to \forall c\,\forall \alpha\,(\mathrm{FLim}(f,c,\alpha) \to c \in \mathbb{R})$
\item $\varphi_{21} :\equiv \forall f\,\forall c\,\forall \alpha\,(\mathrm{FLim}(f,c,\alpha) \to \mathrm{PNbhd}(f,c)) \to \forall c\,\forall \alpha\,(\mathrm{FLim}(g,c,\alpha) \to \mathrm{PNbhd}(g,c))$
\item $\varphi_{22} :\equiv \forall f\,\forall c\,\forall \alpha\,(\mathrm{FLim}(f,c,\alpha) \to \mathrm{PNbhd}(f,c)) \to \forall c\,\forall \alpha\,(\mathrm{FLim}(f,c,\alpha) \to \mathrm{PNbhd}(f,c))$
\item $\varphi_{23} :\equiv \neg \neg (\neg (\neg (d \in \mathbb{R} \to \neg 0 < d) \to \neg (\neg d < v_{768} \to d = v_{768})) \to \neg (\neg d < v_{769} \to d = v_{769}))$
\item $\varphi_{24} :\equiv \forall n \in \mathbb{N}\,(a_{n} \in \mathrm{dom}(g) \land a_{n} \neq c) \to \mathrm{Conv}(g \circ a,\beta)$
\item $\varphi_{25} :\equiv \forall n \in \mathbb{N}\,(a_{n} \in \mathrm{dom}(f) \land a_{n} \neq c) \to \mathrm{Conv}(f \circ a,\alpha)$
\item $\varphi_{26} :\equiv \forall w\,(f(a_{n}) = w \to f(a_{n}) = g(a_{n}) \to w = g(a_{n}))$
\item $\varphi_{27} :\equiv \neg \neg (\neg (\neg (u_{5} \in \mathbb{R} \to \neg 0 < u_{5}) \to \neg (\neg u_{5} < v_{768} \to u_{5} = v_{768})) \to \neg (\neg u_{5} < v_{769} \to u_{5} = v_{769}))$
\item $\varphi_{28} :\equiv \forall w\,(f(a_{n}) = w \to f(a_{n}) < g(a_{n}) \to w < g(a_{n}))$
\item $\varphi_{29} :\equiv \forall a\,\forall b\,\forall \alpha\,\forall \beta\,(\mathrm{Conv}(a,\alpha) \to \mathrm{Conv}(b,\beta) \to \forall n \in \mathbb{N}\,(a_{n} \leq b_{n}) \to \alpha \leq \beta)$
\item $\varphi_{30} :\equiv \rho \in \mathbb{R} \land 0 < \rho \land \rho \leq d \land \rho \leq r \to \rho \in \mathbb{R} \land 0 < \rho \land \rho \leq d \land \rho \leq r$
\item $\varphi_{31} :\equiv \forall x \in \mathbb{R}\,(0 < \lvert x - c \rvert \to \lvert x - c \rvert < d \to x \in \mathrm{dom}(f) \land x \in \mathrm{dom}(g))$
\item $\varphi_{32} :\equiv \forall w\,(g(a_{n}) = w \to f \circ a(n) = g(a_{n}) \to f \circ a(n) = w)$
\item $\varphi_{33} :\equiv \forall n\,(\mathrm{Seq}(a) \to n \in \mathbb{N} \to a_{n} \in \mathbb{R}) \to \mathrm{Seq}(a) \to n \in \mathbb{N} \to a_{n} \in \mathbb{R}$
\item $\varphi_{34} :\equiv n \in \mathbb{N} \to \mathrm{Seq}(f \circ a) \land \forall n \in \mathbb{N}\,(f \circ a(n) = f(a_{n}))$
\item $\varphi_{35} :\equiv n \in \mathbb{N} \to \mathrm{Seq}(g \circ a) \land \forall n \in \mathbb{N}\,(g \circ a(n) = g(a_{n}))$
\item $\varphi_{36} :\equiv \forall \rho\,(\neg (\rho \in \mathbb{R} \land 0 < \rho \land \rho \leq d \land \rho \leq r)) \to \neg (u_{5} \in \mathbb{R} \land 0 < u_{5} \land u_{5} \leq d \land u_{5} \leq r)$
\item $\varphi_{37} :\equiv \forall p\,\forall q\,(p \in \mathbb{R} \to q \in \mathbb{R} \to 0 < p \to 0 < q \to \exists d\,(d \in \mathbb{R} \land 0 < d \land d \leq p \land d \leq q))$
\item $\varphi_{38} :\equiv \mathrm{Conv}(g \circ a,\beta) \to \forall n \in \mathbb{N}\,(f \circ a(n) \leq g \circ a(n)) \to \alpha \leq \beta$
\item $\varphi_{39} :\equiv \mathrm{Seq}(a) \land \mathrm{Conv}(a,c) \land \forall n \in \mathbb{N}\,(0 < \lvert a_{n} - c \rvert \land \lvert a_{n} - c \rvert < \rho)$
\item $\varphi_{40} :\equiv \forall w\,(g(a_{n}) = w \to f \circ a(n) < g(a_{n}) \to f \circ a(n) < w)$
\item $\varphi_{41} :\equiv \forall a\,\forall n\,(\mathrm{Seq}(a) \to n \in \mathbb{N} \to a_{n} \in \mathbb{R}) \to \forall n\,(\mathrm{Seq}(a) \to n \in \mathbb{N} \to a_{n} \in \mathbb{R})$
\item $\varphi_{42} :\equiv \forall x \in \mathrm{dom}(f)\,(x \in \mathrm{dom}(g) \to 0 < \lvert x - c \rvert \to \lvert x - c \rvert < r \to f(x) \leq g(x))$
\item $\varphi_{43} :\equiv \forall q\,(v_{768} \in \mathbb{R} \to q \in \mathbb{R} \to 0 < v_{768} \to 0 < q \to \exists d\,(d \in \mathbb{R} \land 0 < d \land d \leq v_{768} \land d \leq q))$
\item $\varphi_{44} :\equiv \forall u_{5}\,\forall \rho\,(\neg (\rho \in \mathbb{R} \land 0 < \rho \land \rho \leq d \land \rho \leq r)) \to \forall u_{5}\,(\neg (u_{5} \in \mathbb{R} \land 0 < u_{5} \land u_{5} \leq d \land u_{5} \leq r))$
\item $\varphi_{45} :\equiv 0 < \lvert a_{n} - c \rvert \to \lvert a_{n} - c \rvert < d \to a_{n} \in \mathrm{dom}(f) \land a_{n} \in \mathrm{dom}(g)$
\item $\varphi_{46} :\equiv \forall x\,\forall y\,\forall z\,(x < y \to y \leq z \to x < z) \to \forall y\,\forall z\,(\lvert a_{n} - c \rvert < y \to y \leq z \to \lvert a_{n} - c \rvert < z)$
\item $\varphi_{47} :\equiv \forall n \in \mathbb{N}\,(a_{n} \in \mathrm{dom}(f) \land a_{n} \neq c) \to n \in \mathbb{N} \to a_{n} \in \mathrm{dom}(f) \land a_{n} \neq c$
\item $\varphi_{48} :\equiv \forall n \in \mathbb{N}\,(a_{n} \in \mathrm{dom}(g) \land a_{n} \neq c) \to n \in \mathbb{N} \to a_{n} \in \mathrm{dom}(g) \land a_{n} \neq c$
\item $\varphi_{49} :\equiv \forall \rho\,(\neg (\rho \in \mathbb{R} \land 0 < \rho \land \rho \leq d \land \rho \leq r)) \to \forall u_{5}\,\forall \rho\,(\neg (\rho \in \mathbb{R} \land 0 < \rho \land \rho \leq d \land \rho \leq r))$
\item $\varphi_{50} :\equiv \forall v_{769}\,(d \in \mathbb{R} \to v_{769} \in \mathbb{R} \to 0 < d \to 0 < v_{769} \to \exists u_{5}\,(u_{5} \in \mathbb{R} \land 0 < u_{5} \land u_{5} \leq d \land u_{5} \leq v_{769}))$
\item $\varphi_{51} :\equiv \forall b\,\forall \alpha\,\forall \beta\,(\mathrm{Conv}(f \circ a,\alpha) \to \mathrm{Conv}(b,\beta) \to \forall n \in \mathbb{N}\,(f \circ a(n) \leq b_{n}) \to \alpha \leq \beta)$
\item $\varphi_{52} :\equiv \forall n \in \mathbb{N}\,(f \circ a(n) = f(a_{n})) \to n \in \mathbb{N} \to f \circ a(n) = f(a_{n})$
\item $\varphi_{53} :\equiv \forall n \in \mathbb{N}\,(g \circ a(n) = g(a_{n})) \to n \in \mathbb{N} \to g \circ a(n) = g(a_{n})$
\item $\varphi_{54} :\equiv a_{n} \in \mathrm{dom}(g) \to 0 < \lvert a_{n} - c \rvert \to \lvert a_{n} - c \rvert < r \to f(a_{n}) \leq g(a_{n})$
\item $\varphi_{55} :\equiv \forall \alpha\,\forall a\,(\mathrm{FLim}(g,c,\alpha) \to \mathrm{Conv}(a,c) \to \forall n \in \mathbb{N}\,(a_{n} \in \mathrm{dom}(g) \land a_{n} \neq c) \to \mathrm{Conv}(g \circ a,\alpha))$
\item $\varphi_{56} :\equiv f(a_{n}) = f \circ a(n) \to f(a_{n}) = g(a_{n}) \to f \circ a(n) = g(a_{n})$
\item $\varphi_{57} :\equiv \forall n \in \mathbb{N}\,(a_{n} \in \mathrm{dom}(f)) \to \mathrm{Seq}(f \circ a) \land \forall n \in \mathbb{N}\,(f \circ a(n) = f(a_{n}))$
\item $\varphi_{58} :\equiv \forall n \in \mathbb{N}\,(a_{n} \in \mathrm{dom}(g)) \to \mathrm{Seq}(g \circ a) \land \forall n \in \mathbb{N}\,(g \circ a(n) = g(a_{n}))$
\item $\varphi_{59} :\equiv \forall c\,\forall \alpha\,\forall a\,(\mathrm{FLim}(f,c,\alpha) \to \mathrm{Conv}(a,c) \to \varphi_{25})$
\item $\varphi_{60} :\equiv f(a_{n}) = f \circ a(n) \to f(a_{n}) < g(a_{n}) \to f \circ a(n) < g(a_{n})$
\item $\varphi_{61} :\equiv \forall v_{768}\,\forall v_{769}\,(\varphi_{73})$
\item $\varphi_{62} :\equiv g(a_{n}) = g \circ a(n) \to f \circ a(n) = g(a_{n}) \to f \circ a(n) = g \circ a(n)$
\item $\varphi_{63} :\equiv g(a_{n}) = g \circ a(n) \to f \circ a(n) < g(a_{n}) \to f \circ a(n) < g \circ a(n)$
\item $\varphi_{64} :\equiv \forall z\,(\lvert a_{n} - c \rvert < \rho \to \rho \leq z \to \lvert a_{n} - c \rvert < z) \to \lvert a_{n} - c \rvert < \rho \to \rho \leq r \to \lvert a_{n} - c \rvert < r$
\item $\varphi_{65} :\equiv \forall z\,(\lvert a_{n} - c \rvert < \rho \to \rho \leq z \to \lvert a_{n} - c \rvert < z) \to \lvert a_{n} - c \rvert < \rho \to \rho \leq d \to \lvert a_{n} - c \rvert < d$
\item $\varphi_{66} :\equiv \forall y\,(f \circ a(n) = y \to y = f \circ a(n)) \to f \circ a(n) = f(a_{n}) \to f(a_{n}) = f \circ a(n)$
\item $\varphi_{67} :\equiv \forall y\,(g \circ a(n) = y \to y = g \circ a(n)) \to g \circ a(n) = g(a_{n}) \to g(a_{n}) = g \circ a(n)$
\item $\varphi_{68} :\equiv \forall y\,\forall z\,(\lvert a_{n} - c \rvert < y \to y \leq z \to \lvert a_{n} - c \rvert < z) \to \forall z\,(\lvert a_{n} - c \rvert < \rho \to \rho \leq z \to \lvert a_{n} - c \rvert < z)$
\item $\varphi_{69} :\equiv \forall r\,(c \in \mathbb{R} \to r \in \mathbb{R} \to 0 < r \to \exists a\,(\mathrm{Seq}(a) \land \mathrm{Conv}(a,c) \land \forall n \in \mathbb{N}\,(0 < \lvert a_{n} - c \rvert \land \lvert a_{n} - c \rvert < r)))$
\item $\varphi_{70} :\equiv \forall y\,\forall c\,(y \in \mathbb{R} \to c \in \mathbb{R} \to 0 < \lvert y - c \rvert \to y \neq c) \to \forall c\,(a_{n} \in \mathbb{R} \to c \in \mathbb{R} \to 0 < \lvert a_{n} - c \rvert \to a_{n} \neq c)$
\item $\varphi_{71} :\equiv \forall n \in \mathbb{N}\,(0 < \lvert a_{n} - c \rvert \land \lvert a_{n} - c \rvert < \rho) \to n \in \mathbb{N} \to 0 < \lvert a_{n} - c \rvert \land \lvert a_{n} - c \rvert < \rho$
\item $\varphi_{72} :\equiv \mathrm{PNbhd}(f,c) \to \mathrm{PNbhd}(g,c) \to \exists d\,(d \in \mathbb{R} \land 0 < d \land \varphi_{31})$
\item $\varphi_{73} :\equiv v_{768} \in \mathbb{R} \to v_{769} \in \mathbb{R} \to 0 < v_{768} \to 0 < v_{769} \to \neg \forall u_{5}\,(\varphi_{27})$
\item $\varphi_{74} :\equiv \mathrm{FLim}(g,c,\beta) \to r \in \mathbb{R} \to 0 < r \to \varphi_{42} \to \mathrm{PNbhd}(g,c)$
\item $\varphi_{75} :\equiv \mathrm{FLim}(g,c,\beta) \to r \in \mathbb{R} \to 0 < r \to \varphi_{42} \to \mathrm{PNbhd}(f,c)$
\item $\varphi_{76} :\equiv d \in \mathbb{R} \land 0 < d \land \varphi_{31} \to \exists u_{5}\,(u_{5} \in \mathbb{R} \land 0 < u_{5} \land u_{5} \leq d \land u_{5} \leq r)$
\item $\varphi_{77} :\equiv \forall c\,(a_{n} \in \mathbb{R} \to c \in \mathbb{R} \to 0 < \lvert a_{n} - c \rvert \to a_{n} \neq c) \to a_{n} \in \mathbb{R} \to c \in \mathbb{R} \to 0 < \lvert a_{n} - c \rvert \to a_{n} \neq c$
\item $\varphi_{78} :\equiv \mathrm{FLim}(f,c,\alpha) \to \mathrm{FLim}(g,c,\beta) \to r \in \mathbb{R} \to 0 < r \to \varphi_{42} \to \alpha \leq \beta$
\item $\varphi_{79} :\equiv \mathrm{FLim}(f,c,\alpha) \to \mathrm{FLim}(g,c,\beta) \to r \in \mathbb{R} \to 0 < r \to \varphi_{42} \to \mathrm{RFn}(f)$
\item $\varphi_{80} :\equiv \mathrm{FLim}(f,c,\alpha) \to \mathrm{FLim}(g,c,\beta) \to r \in \mathbb{R} \to 0 < r \to \varphi_{42} \to \mathrm{RFn}(g)$
\item $\varphi_{81} :\equiv \mathrm{FLim}(f,c,\alpha) \land \mathrm{FLim}(g,c,\beta) \land r \in \mathbb{R} \land 0 < r \land \varphi_{42} \land \neg \alpha \leq \beta$
\item $\varphi_{82} :\equiv \mathrm{FLim}(f,c,\alpha) \to \mathrm{FLim}(g,c,\beta) \to r \in \mathbb{R} \to 0 < r \to \varphi_{42} \to c \in \mathbb{R}$
\item $\varphi_{83} :\equiv d \in \mathbb{R} \land 0 < d \land \varphi_{31} \to \rho \in \mathbb{R} \land 0 < \rho \land \rho \leq d \land \rho \leq r \to \alpha \leq \beta$
\item $\varphi_{84} :\equiv d \in \mathbb{R} \land 0 < d \land \varphi_{31} \to \rho \in \mathbb{R} \land 0 < \rho \land \rho \leq d \land \rho \leq r \to 0 < \rho$
\item $\varphi_{85} :\equiv d \in \mathbb{R} \land 0 < d \land \varphi_{31} \to \rho \in \mathbb{R} \land 0 < \rho \land \rho \leq d \land \rho \leq r \to \rho \in \mathbb{R}$
\item $\varphi_{86} :\equiv \rho \in \mathbb{R} \land 0 < \rho \land \rho \leq d \land \rho \leq r \to \varphi_{39} \to n \in \mathbb{N} \to a_{n} \in \mathrm{dom}(f)$
\item $\varphi_{87} :\equiv \rho \in \mathbb{R} \land 0 < \rho \land \rho \leq d \land \rho \leq r \to \varphi_{39} \to n \in \mathbb{N} \to a_{n} \in \mathrm{dom}(g)$
\item $\varphi_{88} :\equiv \rho \in \mathbb{R} \land 0 < \rho \land \rho \leq d \land \rho \leq r \to \varphi_{39} \to \forall n \in \mathbb{N}\,(a_{n} \in \mathrm{dom}(f))$
\item $\varphi_{89} :\equiv \rho \in \mathbb{R} \land 0 < \rho \land \rho \leq d \land \rho \leq r \to \varphi_{39} \to \forall n \in \mathbb{N}\,(a_{n} \in \mathrm{dom}(g))$
\item $\varphi_{90} :\equiv \rho \in \mathbb{R} \land 0 < \rho \land \rho \leq d \land \rho \leq r \to \varphi_{39} \to n \in \mathbb{N} \to \lvert a_{n} - c \rvert < r$
\item $\varphi_{91} :\equiv \rho \in \mathbb{R} \land 0 < \rho \land \rho \leq d \land \rho \leq r \to \varphi_{39} \to n \in \mathbb{N} \to \lvert a_{n} - c \rvert < d$
\item $\varphi_{92} :\equiv \rho \in \mathbb{R} \land 0 < \rho \land \rho \leq d \land \rho \leq r \to \varphi_{39} \to n \in \mathbb{N} \to 0 < \lvert a_{n} - c \rvert$
\item $\varphi_{93} :\equiv \rho \in \mathbb{R} \land 0 < \rho \land \rho \leq d \land \rho \leq r \to \varphi_{39} \to n \in \mathbb{N} \to \lvert a_{n} - c \rvert < \rho$
\item $\varphi_{94} :\equiv \rho \in \mathbb{R} \land 0 < \rho \land \rho \leq d \land \rho \leq r \to \varphi_{39} \to n \in \mathbb{N} \to f(a_{n}) \leq g(a_{n})$
\item $\varphi_{95} :\equiv \rho \in \mathbb{R} \land 0 < \rho \land \rho \leq d \land \rho \leq r \to \varphi_{39} \to n \in \mathbb{N} \to a_{n} \in \mathrm{dom}(f) \land a_{n} \neq c$
\item $\varphi_{96} :\equiv \rho \in \mathbb{R} \land 0 < \rho \land \rho \leq d \land \rho \leq r \to \varphi_{39} \to n \in \mathbb{N} \to a_{n} \in \mathrm{dom}(g) \land a_{n} \neq c$
\item $\varphi_{97} :\equiv \rho \in \mathbb{R} \land 0 < \rho \land \rho \leq d \land \rho \leq r \to \varphi_{39} \to n \in \mathbb{N} \to f \circ a(n) = f(a_{n})$
\item $\varphi_{98} :\equiv \rho \in \mathbb{R} \land 0 < \rho \land \rho \leq d \land \rho \leq r \to \varphi_{39} \to n \in \mathbb{N} \to f(a_{n}) = f \circ a(n)$
\item $\varphi_{99} :\equiv \rho \in \mathbb{R} \land 0 < \rho \land \rho \leq d \land \rho \leq r \to \varphi_{39} \to n \in \mathbb{N} \to f \circ a(n) \leq g(a_{n})$
\item $\varphi_{100} :\equiv \rho \in \mathbb{R} \land 0 < \rho \land \rho \leq d \land \rho \leq r \to \varphi_{39} \to n \in \mathbb{N} \to g \circ a(n) = g(a_{n})$
\item $\varphi_{101} :\equiv \rho \in \mathbb{R} \land 0 < \rho \land \rho \leq d \land \rho \leq r \to \varphi_{39} \to n \in \mathbb{N} \to g(a_{n}) = g \circ a(n)$
\item $\varphi_{102} :\equiv \rho \in \mathbb{R} \land 0 < \rho \land \rho \leq d \land \rho \leq r \to \varphi_{39} \to \forall n \in \mathbb{N}\,(a_{n} \in \mathrm{dom}(f) \land a_{n} \neq c)$
\item $\varphi_{103} :\equiv \rho \in \mathbb{R} \land 0 < \rho \land \rho \leq d \land \rho \leq r \to \varphi_{39} \to \forall n \in \mathbb{N}\,(a_{n} \in \mathrm{dom}(g) \land a_{n} \neq c)$
\item $\varphi_{104} :\equiv (\mathrm{FLim}(f,c,\alpha) \to \mathrm{FLim}(f,c,\alpha)) \to \mathrm{FLim}(f,c,\alpha) \to \varphi_{75}$
\item $\varphi_{105} :\equiv \rho \in \mathbb{R} \land 0 < \rho \land \rho \leq d \land \rho \leq r \to \varphi_{39} \to n \in \mathbb{N} \to f \circ a(n) \leq g \circ a(n)$
\item $\varphi_{106} :\equiv \rho \in \mathbb{R} \land 0 < \rho \land \rho \leq d \land \rho \leq r \to \varphi_{39} \to n \in \mathbb{N} \to a_{n} \in \mathrm{dom}(f) \land a_{n} \in \mathrm{dom}(g)$
\item $\varphi_{107} :\equiv \rho \in \mathbb{R} \land 0 < \rho \land \rho \leq d \land \rho \leq r \to \varphi_{39} \to \forall n \in \mathbb{N}\,(f \circ a(n) \leq g \circ a(n))$
\item $\varphi_{108} :\equiv \rho \in \mathbb{R} \land 0 < \rho \land \rho \leq d \land \rho \leq r \to \varphi_{39} \to n \in \mathbb{N} \to n \in \mathbb{N} \to f \circ a(n) = f(a_{n})$
\item $\varphi_{109} :\equiv \rho \in \mathbb{R} \land 0 < \rho \land \rho \leq d \land \rho \leq r \to \varphi_{39} \to n \in \mathbb{N} \to n \in \mathbb{N} \to g \circ a(n) = g(a_{n})$
\item $\varphi_{110} :\equiv \rho \in \mathbb{R} \land 0 < \rho \land \rho \leq d \land \rho \leq r \to \varphi_{39} \to n \in \mathbb{N} \to \forall n \in \mathbb{N}\,(f \circ a(n) = f(a_{n}))$
\item $\varphi_{111} :\equiv \rho \in \mathbb{R} \land 0 < \rho \land \rho \leq d \land \rho \leq r \to \varphi_{39} \to n \in \mathbb{N} \to \forall n \in \mathbb{N}\,(g \circ a(n) = g(a_{n}))$
\item $\varphi_{112} :\equiv \rho \in \mathbb{R} \land 0 < \rho \land \rho \leq d \land \rho \leq r \to \varphi_{39} \to n \in \mathbb{N} \to 0 < \lvert a_{n} - c \rvert \land \lvert a_{n} - c \rvert < \rho$
\item $\varphi_{113} :\equiv \rho \in \mathbb{R} \land 0 < \rho \land \rho \leq d \land \rho \leq r \to \varphi_{39} \to \forall n \in \mathbb{N}\,(0 < \lvert a_{n} - c \rvert \land \lvert a_{n} - c \rvert < \rho)$
\item $\varphi_{114} :\equiv (\mathrm{FLim}(f,c,\alpha) \to \mathrm{RFn}(f)) \to (\mathrm{FLim}(f,c,\alpha) \to \mathrm{FLim}(f,c,\alpha)) \to \varphi_{79}$
\item $\varphi_{115} :\equiv \forall a\,(\mathrm{FLim}(g,c,\beta) \to \mathrm{Conv}(a,c) \to \varphi_{24}) \to \mathrm{FLim}(g,c,\beta) \to \mathrm{Conv}(a,c) \to \varphi_{24}$
\item $\varphi_{116} :\equiv \forall a\,(\mathrm{FLim}(f,c,\alpha) \to \mathrm{Conv}(a,c) \to \varphi_{25}) \to \mathrm{FLim}(f,c,\alpha) \to \mathrm{Conv}(a,c) \to \varphi_{25}$
\item $\varphi_{117} :\equiv (\mathrm{FLim}(f,c,\alpha) \to c \in \mathbb{R}) \to (\mathrm{FLim}(f,c,\alpha) \to \mathrm{FLim}(f,c,\alpha)) \to \varphi_{82}$
\item $\varphi_{118} :\equiv \forall \alpha\,\forall a\,(\mathrm{FLim}(f,c,\alpha) \to \mathrm{Conv}(a,c) \to \varphi_{25}) \to \forall a\,(\mathrm{FLim}(f,c,\alpha) \to \mathrm{Conv}(a,c) \to \varphi_{25})$
\item $\varphi_{119} :\equiv \varphi_{59} \to \forall \alpha\,\forall a\,(\mathrm{FLim}(f,c,\alpha) \to \mathrm{Conv}(a,c) \to \varphi_{25})$
\item $\varphi_{120} :\equiv (\mathrm{FLim}(g,c,\beta) \to \mathrm{PNbhd}(g,c)) \to \varphi_{1} \to \mathrm{FLim}(f,c,\alpha) \to \varphi_{74}$
\item $\varphi_{121} :\equiv \forall \beta\,(\mathrm{Conv}(f \circ a,\alpha) \to \varphi_{38}) \to \mathrm{Conv}(f \circ a,\alpha) \to \varphi_{38}$
\item $\varphi_{122} :\equiv \mathrm{FLim}(f,c,\alpha) \to \mathrm{FLim}(g,c,\beta) \to r \in \mathbb{R} \to 0 < r \to \varphi_{42} \to \varphi_{42}$
\item $\varphi_{123} :\equiv \mathrm{FLim}(g,c,\beta) \to r \in \mathbb{R} \to 0 < r \to \varphi_{42} \to \exists d\,(d \in \mathbb{R} \land 0 < d \land \varphi_{31})$
\item $\varphi_{124} :\equiv \forall \alpha\,\forall \beta\,(\mathrm{Conv}(f \circ a,\alpha) \to \varphi_{38}) \to \forall \beta\,(\mathrm{Conv}(f \circ a,\alpha) \to \varphi_{38})$
\item $\varphi_{125} :\equiv v_{769} \in \mathbb{R} \to \varphi_{17} \leftrightarrow v_{769} \in \mathbb{R} \to 0 < v_{768} \to 0 < v_{769} \to \neg \forall u_{5}\,(\varphi_{27})$
\item $\varphi_{126} :\equiv d \in \mathbb{R} \land 0 < d \land \varphi_{31} \to \rho \in \mathbb{R} \land 0 < \rho \land \rho \leq d \land \rho \leq r \to \exists a\,(\varphi_{39})$
\item $\varphi_{127} :\equiv \mathrm{FLim}(g,c,\beta) \to r \in \mathbb{R} \to 0 < r \to \varphi_{42} \to d \in \mathbb{R} \land 0 < d \land \varphi_{31} \to \alpha \leq \beta$
\item $\varphi_{128} :\equiv \mathrm{FLim}(g,c,\beta) \to r \in \mathbb{R} \to 0 < r \to \varphi_{42} \to d \in \mathbb{R} \land 0 < d \land \varphi_{31} \to 0 < d$
\item $\varphi_{129} :\equiv d \in \mathbb{R} \land 0 < d \land \varphi_{31} \to \rho \in \mathbb{R} \land 0 < \rho \land \rho \leq d \land \rho \leq r \to \varphi_{39} \to \alpha \leq \beta$
\item $\varphi_{130} :\equiv d \in \mathbb{R} \land 0 < d \land \varphi_{31} \to \rho \in \mathbb{R} \land 0 < \rho \land \rho \leq d \land \rho \leq r \to \varphi_{39} \to \mathrm{Seq}(a)$
\item $\varphi_{131} :\equiv \mathrm{FLim}(g,c,\beta) \to r \in \mathbb{R} \to 0 < r \to \varphi_{42} \to d \in \mathbb{R} \land 0 < d \land \varphi_{31} \to d \in \mathbb{R}$
\item $\varphi_{132} :\equiv d \in \mathbb{R} \land 0 < d \land \varphi_{31} \to \rho \in \mathbb{R} \land 0 < \rho \land \rho \leq d \land \rho \leq r \to \varphi_{39} \to \mathrm{Conv}(a,c)$
\item $\varphi_{133} :\equiv \mathrm{FLim}(f,c,\alpha) \land \mathrm{FLim}(g,c,\beta) \land r \in \mathbb{R} \land 0 < r \land \varphi_{42} \land (d \in \mathbb{R} \land 0 < d \land \varphi_{31}) \land \neg \alpha \leq \beta$
\item $\varphi_{134} :\equiv d \in \mathbb{R} \land 0 < d \land \varphi_{31} \to \rho \in \mathbb{R} \land 0 < \rho \land \rho \leq d \land \rho \leq r \to \varphi_{39} \to n \in \mathbb{N} \to \rho \leq r$
\item $\varphi_{135} :\equiv d \in \mathbb{R} \land 0 < d \land \varphi_{31} \to \rho \in \mathbb{R} \land 0 < \rho \land \rho \leq d \land \rho \leq r \to \varphi_{39} \to n \in \mathbb{N} \to \rho \leq d$
\item $\varphi_{136} :\equiv \mathrm{FLim}(g,c,\beta) \to r \in \mathbb{R} \to 0 < r \to \varphi_{42} \to a_{n} \in \mathrm{dom}(f) \to \varphi_{54}$
\item $\varphi_{137} :\equiv \forall a\,(\mathrm{RFn}(f) \to \mathrm{Seq}(a) \to \varphi_{57}) \to \mathrm{RFn}(f) \to \mathrm{Seq}(a) \to \varphi_{57}$
\item $\varphi_{138} :\equiv \forall a\,(\mathrm{RFn}(g) \to \mathrm{Seq}(a) \to \varphi_{58}) \to \mathrm{RFn}(g) \to \mathrm{Seq}(a) \to \varphi_{58}$
\item $\varphi_{139} :\equiv \rho \in \mathbb{R} \land 0 < \rho \land \rho \leq d \land \rho \leq r \to \varphi_{39} \to n \in \mathbb{N} \to a_{n} \in \mathbb{R} \to \varphi_{45}$
\item $\varphi_{140} :\equiv d \in \mathbb{R} \land 0 < d \land \varphi_{31} \to \rho \in \mathbb{R} \land 0 < \rho \land \rho \leq d \land \rho \leq r \to \varphi_{39} \to n \in \mathbb{N} \to n \in \mathbb{N}$
\item $\varphi_{141} :\equiv d \in \mathbb{R} \land 0 < d \land \varphi_{31} \to \rho \in \mathbb{R} \land 0 < \rho \land \rho \leq d \land \rho \leq r \to \varphi_{39} \to \mathrm{Conv}(g \circ a,\beta)$
\item $\varphi_{142} :\equiv d \in \mathbb{R} \land 0 < d \land \varphi_{31} \to \rho \in \mathbb{R} \land 0 < \rho \land \rho \leq d \land \rho \leq r \to \varphi_{39} \to \mathrm{Conv}(f \circ a,\alpha)$
\item $\varphi_{143} :\equiv d \in \mathbb{R} \land 0 < d \land \varphi_{31} \to \rho \in \mathbb{R} \land 0 < \rho \land \rho \leq d \land \rho \leq r \to \varphi_{39} \to n \in \mathbb{N} \to a_{n} \neq c$
\item $\varphi_{144} :\equiv \forall f\,\forall a\,(\mathrm{RFn}(f) \to \mathrm{Seq}(a) \to \varphi_{57}) \to \forall a\,(\mathrm{RFn}(f) \to \mathrm{Seq}(a) \to \varphi_{57})$
\item $\varphi_{145} :\equiv \forall f\,\forall a\,(\mathrm{RFn}(f) \to \mathrm{Seq}(a) \to \varphi_{57}) \to \forall a\,(\mathrm{RFn}(g) \to \mathrm{Seq}(a) \to \varphi_{58})$
\item $\varphi_{146} :\equiv d \in \mathbb{R} \land 0 < d \land \varphi_{31} \to \rho \in \mathbb{R} \land 0 < \rho \land \rho \leq d \land \rho \leq r \to \varphi_{39} \to n \in \mathbb{N} \to a_{n} \in \mathbb{R}$
\item $\varphi_{147} :\equiv \forall v_{768}\,(\forall v_{769}\,(\varphi_{73}) \to \forall v_{769}\,(v_{768} \in \mathbb{R} \to v_{769} \in \mathbb{R} \to \varphi_{17}))$
\item $\varphi_{148} :\equiv \forall v_{768}\,(\forall v_{769}\,(v_{768} \in \mathbb{R} \to v_{769} \in \mathbb{R} \to \varphi_{17}) \to \forall v_{769}\,(\varphi_{73}))$
\item $\varphi_{149} :\equiv \mathrm{FLim}(f,c,\alpha) \to \mathrm{FLim}(g,c,\beta) \to r \in \mathbb{R} \to 0 < r \to \varphi_{42} \to \varphi_{76}$
\item $\varphi_{150} :\equiv \mathrm{FLim}(f,c,\alpha) \to \mathrm{FLim}(g,c,\beta) \to r \in \mathbb{R} \to 0 < r \to \varphi_{42} \to \varphi_{83}$
\item $\varphi_{151} :\equiv \mathrm{FLim}(f,c,\alpha) \land \mathrm{FLim}(g,c,\beta) \land r \in \mathbb{R} \land 0 < r \land \varphi_{42} \land (d \in \mathbb{R} \land 0 < d \land \varphi_{31}) \land (\rho \in \mathbb{R} \land 0 < \rho \land \rho \leq d \land \rho \leq r) \land \neg \alpha \leq \beta$
\item $\varphi_{152} :\equiv \mathrm{FLim}(f,c,\alpha) \to \mathrm{FLim}(g,c,\beta) \to r \in \mathbb{R} \to 0 < r \to \varphi_{42} \to \varphi_{84}$
\item $\varphi_{153} :\equiv \mathrm{FLim}(f,c,\alpha) \to \mathrm{FLim}(g,c,\beta) \to r \in \mathbb{R} \to 0 < r \to \varphi_{42} \to \varphi_{85}$
\item $\varphi_{154} :\equiv \mathrm{FLim}(g,c,\beta) \to r \in \mathbb{R} \to 0 < r \to \varphi_{42} \to d \in \mathbb{R} \land 0 < d \land \varphi_{31} \to \varphi_{30}$
\item $\varphi_{155} :\equiv (\forall d\,(\varphi_{23}) \leftrightarrow \forall u_{5}\,(\varphi_{27})) \to (\neg \forall d\,(\varphi_{23}) \leftrightarrow \neg \forall u_{5}\,(\varphi_{27}))$
\item $\varphi_{156} :\equiv d \in \mathbb{R} \land 0 < d \land \varphi_{31} \to \rho \in \mathbb{R} \land 0 < \rho \land \rho \leq d \land \rho \leq r \to \varphi_{39} \to \varphi_{34}$
\item $\varphi_{157} :\equiv d \in \mathbb{R} \land 0 < d \land \varphi_{31} \to \rho \in \mathbb{R} \land 0 < \rho \land \rho \leq d \land \rho \leq r \to \varphi_{39} \to \varphi_{35}$
\item $\varphi_{158} :\equiv r \in \mathbb{R} \to 0 < r \to \varphi_{42} \to d \in \mathbb{R} \land 0 < d \land \varphi_{31} \to d \in \mathbb{R} \land 0 < d \land \varphi_{31}$
\item $\varphi_{159} :\equiv d \in \mathbb{R} \land 0 < d \land \varphi_{31} \to \rho \in \mathbb{R} \land 0 < \rho \land \rho \leq d \land \rho \leq r \to \varphi_{39} \to \varphi_{39}$
\item $\varphi_{160} :\equiv d \in \mathbb{R} \land 0 < d \land \varphi_{31} \to \rho \in \mathbb{R} \land 0 < \rho \land \rho \leq d \land \rho \leq r \to \varphi_{39} \to n \in \mathbb{N} \to \varphi_{31}$
\item $\varphi_{161} :\equiv (\neg \forall d\,(\varphi_{23}) \leftrightarrow \neg \forall u_{5}\,(\varphi_{27})) \to (0 < v_{769} \to \neg \forall d\,(\varphi_{23}) \leftrightarrow 0 < v_{769} \to \neg \forall u_{5}\,(\varphi_{27}))$
\item $\varphi_{162} :\equiv \mathrm{FLim}(f,c,\alpha) \land \mathrm{FLim}(g,c,\beta) \land r \in \mathbb{R} \land 0 < r \land \varphi_{42} \land (d \in \mathbb{R} \land 0 < d \land \varphi_{31}) \land (\rho \in \mathbb{R} \land 0 < \rho \land \rho \leq d \land \rho \leq r) \land \varphi_{39}$
\item $\varphi_{163} :\equiv \mathrm{FLim}(f,c,\alpha) \to \mathrm{FLim}(g,c,\beta) \to r \in \mathbb{R} \to 0 < r \to \varphi_{42} \to \varphi_{126}$
\item $\varphi_{164} :\equiv \mathrm{FLim}(f,c,\alpha) \to \mathrm{FLim}(g,c,\beta) \to r \in \mathbb{R} \to 0 < r \to \varphi_{42} \to \varphi_{129}$
\item $\varphi_{165} :\equiv \mathrm{FLim}(f,c,\alpha) \to \mathrm{FLim}(g,c,\beta) \to r \in \mathbb{R} \to 0 < r \to \varphi_{42} \to \varphi_{130}$
\item $\varphi_{166} :\equiv \mathrm{FLim}(f,c,\alpha) \to \mathrm{FLim}(g,c,\beta) \to r \in \mathbb{R} \to 0 < r \to \varphi_{42} \to \varphi_{132}$
\item $\varphi_{167} :\equiv (0 < v_{769} \to \neg \forall d\,(\varphi_{23}) \leftrightarrow 0 < v_{769} \to \neg \forall u_{5}\,(\varphi_{27})) \to (\varphi_{17} \leftrightarrow 0 < v_{768} \to 0 < v_{769} \to \neg \forall u_{5}\,(\varphi_{27}))$
\item $\varphi_{168} :\equiv \mathrm{FLim}(g,c,\beta) \to r \in \mathbb{R} \to 0 < r \to \varphi_{42} \to d \in \mathbb{R} \land 0 < d \land \varphi_{31} \to \varphi_{86}$
\item $\varphi_{169} :\equiv \mathrm{FLim}(g,c,\beta) \to r \in \mathbb{R} \to 0 < r \to \varphi_{42} \to d \in \mathbb{R} \land 0 < d \land \varphi_{31} \to \varphi_{87}$
\item $\varphi_{170} :\equiv \mathrm{FLim}(f,c,\alpha) \to \mathrm{FLim}(g,c,\beta) \to r \in \mathbb{R} \to 0 < r \to \varphi_{42} \to \varphi_{134}$
\item $\varphi_{171} :\equiv \mathrm{FLim}(f,c,\alpha) \to \mathrm{FLim}(g,c,\beta) \to r \in \mathbb{R} \to 0 < r \to \varphi_{42} \to \varphi_{135}$
\item $\varphi_{172} :\equiv \mathrm{FLim}(g,c,\beta) \to r \in \mathbb{R} \to 0 < r \to \varphi_{42} \to d \in \mathbb{R} \land 0 < d \land \varphi_{31} \to \varphi_{88}$
\item $\varphi_{173} :\equiv \mathrm{FLim}(g,c,\beta) \to r \in \mathbb{R} \to 0 < r \to \varphi_{42} \to d \in \mathbb{R} \land 0 < d \land \varphi_{31} \to \varphi_{89}$
\item $\varphi_{174} :\equiv \mathrm{FLim}(g,c,\beta) \to r \in \mathbb{R} \to 0 < r \to \varphi_{42} \to d \in \mathbb{R} \land 0 < d \land \varphi_{31} \to \varphi_{90}$
\item $\varphi_{175} :\equiv \mathrm{FLim}(g,c,\beta) \to r \in \mathbb{R} \to 0 < r \to \varphi_{42} \to d \in \mathbb{R} \land 0 < d \land \varphi_{31} \to \varphi_{91}$
\item $\varphi_{176} :\equiv \mathrm{FLim}(g,c,\beta) \to r \in \mathbb{R} \to 0 < r \to \varphi_{42} \to d \in \mathbb{R} \land 0 < d \land \varphi_{31} \to \varphi_{92}$
\item $\varphi_{177} :\equiv \mathrm{FLim}(f,c,\alpha) \to \mathrm{FLim}(g,c,\beta) \to r \in \mathbb{R} \to 0 < r \to \varphi_{42} \to \varphi_{140}$
\item $\varphi_{178} :\equiv \mathrm{FLim}(f,c,\alpha) \to \mathrm{FLim}(g,c,\beta) \to r \in \mathbb{R} \to 0 < r \to \varphi_{42} \to \varphi_{141}$
\item $\varphi_{179} :\equiv \mathrm{FLim}(f,c,\alpha) \to \mathrm{FLim}(g,c,\beta) \to r \in \mathbb{R} \to 0 < r \to \varphi_{42} \to \varphi_{142}$
\item $\varphi_{180} :\equiv \mathrm{FLim}(g,c,\beta) \to r \in \mathbb{R} \to 0 < r \to \varphi_{42} \to d \in \mathbb{R} \land 0 < d \land \varphi_{31} \to \varphi_{93}$
\item $\varphi_{181} :\equiv \mathrm{FLim}(f,c,\alpha) \to \mathrm{FLim}(g,c,\beta) \to r \in \mathbb{R} \to 0 < r \to \varphi_{42} \to \varphi_{143}$
\item $\varphi_{182} :\equiv \mathrm{FLim}(g,c,\beta) \to r \in \mathbb{R} \to 0 < r \to \varphi_{42} \to d \in \mathbb{R} \land 0 < d \land \varphi_{31} \to \varphi_{94}$
\item $\varphi_{183} :\equiv \mathrm{FLim}(f,c,\alpha) \to \mathrm{FLim}(g,c,\beta) \to r \in \mathbb{R} \to 0 < r \to \varphi_{42} \to \varphi_{146}$
\item $\varphi_{184} :\equiv (\forall v_{768}\,\forall v_{769}\,(v_{768} \in \mathbb{R} \to v_{769} \in \mathbb{R} \to \varphi_{17}) \leftrightarrow \varphi_{61}) \to \varphi_{61}$
\item $\varphi_{185} :\equiv \mathrm{FLim}(g,c,\beta) \to r \in \mathbb{R} \to 0 < r \to \varphi_{42} \to d \in \mathbb{R} \land 0 < d \land \varphi_{31} \to \varphi_{95}$
\item $\varphi_{186} :\equiv \mathrm{FLim}(g,c,\beta) \to r \in \mathbb{R} \to 0 < r \to \varphi_{42} \to d \in \mathbb{R} \land 0 < d \land \varphi_{31} \to \varphi_{96}$
\item $\varphi_{187} :\equiv \mathrm{FLim}(g,c,\beta) \to r \in \mathbb{R} \to 0 < r \to \varphi_{42} \to d \in \mathbb{R} \land 0 < d \land \varphi_{31} \to \varphi_{97}$
\item $\varphi_{188} :\equiv \mathrm{FLim}(g,c,\beta) \to r \in \mathbb{R} \to 0 < r \to \varphi_{42} \to d \in \mathbb{R} \land 0 < d \land \varphi_{31} \to \varphi_{98}$
\item $\varphi_{189} :\equiv \mathrm{FLim}(g,c,\beta) \to r \in \mathbb{R} \to 0 < r \to \varphi_{42} \to d \in \mathbb{R} \land 0 < d \land \varphi_{31} \to \varphi_{99}$
\item $\varphi_{190} :\equiv \mathrm{FLim}(g,c,\beta) \to r \in \mathbb{R} \to 0 < r \to \varphi_{42} \to d \in \mathbb{R} \land 0 < d \land \varphi_{31} \to \varphi_{100}$
\item $\varphi_{191} :\equiv \mathrm{FLim}(g,c,\beta) \to r \in \mathbb{R} \to 0 < r \to \varphi_{42} \to d \in \mathbb{R} \land 0 < d \land \varphi_{31} \to \varphi_{101}$
\item $\varphi_{192} :\equiv \mathrm{FLim}(g,c,\beta) \to r \in \mathbb{R} \to 0 < r \to \varphi_{42} \to d \in \mathbb{R} \land 0 < d \land \varphi_{31} \to \varphi_{102}$
\item $\varphi_{193} :\equiv \mathrm{FLim}(g,c,\beta) \to r \in \mathbb{R} \to 0 < r \to \varphi_{42} \to d \in \mathbb{R} \land 0 < d \land \varphi_{31} \to \varphi_{103}$
\item $\varphi_{194} :\equiv \mathrm{FLim}(g,c,\beta) \to r \in \mathbb{R} \to 0 < r \to \varphi_{42} \to d \in \mathbb{R} \land 0 < d \land \varphi_{31} \to \varphi_{105}$
\item $\varphi_{195} :\equiv \mathrm{FLim}(g,c,\beta) \to r \in \mathbb{R} \to 0 < r \to \varphi_{42} \to d \in \mathbb{R} \land 0 < d \land \varphi_{31} \to \varphi_{106}$
\item $\varphi_{196} :\equiv \mathrm{FLim}(g,c,\beta) \to r \in \mathbb{R} \to 0 < r \to \varphi_{42} \to d \in \mathbb{R} \land 0 < d \land \varphi_{31} \to \varphi_{107}$
\item $\varphi_{197} :\equiv \mathrm{FLim}(g,c,\beta) \to r \in \mathbb{R} \to 0 < r \to \varphi_{42} \to d \in \mathbb{R} \land 0 < d \land \varphi_{31} \to \varphi_{108}$
\item $\varphi_{198} :\equiv \mathrm{FLim}(g,c,\beta) \to r \in \mathbb{R} \to 0 < r \to \varphi_{42} \to d \in \mathbb{R} \land 0 < d \land \varphi_{31} \to \varphi_{109}$
\item $\varphi_{199} :\equiv \mathrm{FLim}(g,c,\beta) \to r \in \mathbb{R} \to 0 < r \to \varphi_{42} \to d \in \mathbb{R} \land 0 < d \land \varphi_{31} \to \varphi_{110}$
\item $\varphi_{200} :\equiv \mathrm{FLim}(g,c,\beta) \to r \in \mathbb{R} \to 0 < r \to \varphi_{42} \to d \in \mathbb{R} \land 0 < d \land \varphi_{31} \to \varphi_{111}$
\item $\varphi_{201} :\equiv \mathrm{FLim}(g,c,\beta) \to r \in \mathbb{R} \to 0 < r \to \varphi_{42} \to d \in \mathbb{R} \land 0 < d \land \varphi_{31} \to \varphi_{112}$
\item $\varphi_{202} :\equiv \mathrm{FLim}(g,c,\beta) \to r \in \mathbb{R} \to 0 < r \to \varphi_{42} \to d \in \mathbb{R} \land 0 < d \land \varphi_{31} \to \varphi_{113}$
\item $\varphi_{203} :\equiv \mathrm{FLim}(f,c,\alpha) \to \mathrm{FLim}(g,c,\beta) \to r \in \mathbb{R} \to 0 < r \to \varphi_{42} \to \varphi_{156}$
\item $\varphi_{204} :\equiv \mathrm{FLim}(f,c,\alpha) \to \mathrm{FLim}(g,c,\beta) \to r \in \mathbb{R} \to 0 < r \to \varphi_{42} \to \varphi_{157}$
\item $\varphi_{205} :\equiv \mathrm{FLim}(f,c,\alpha) \to \mathrm{FLim}(g,c,\beta) \to r \in \mathbb{R} \to 0 < r \to \varphi_{42} \to \varphi_{159}$
\item $\varphi_{206} :\equiv \mathrm{FLim}(f,c,\alpha) \to \mathrm{FLim}(g,c,\beta) \to r \in \mathbb{R} \to 0 < r \to \varphi_{42} \to \varphi_{160}$
\item $\varphi_{207} :\equiv (\forall d\,(\varphi_{23}) \to \forall u_{5}\,(\varphi_{27})) \to (\forall u_{5}\,(\varphi_{27}) \to \forall d\,(\varphi_{23})) \to (\forall d\,(\varphi_{23}) \leftrightarrow \forall u_{5}\,(\varphi_{27}))$
\item $\varphi_{208} :\equiv (v_{769} \in \mathbb{R} \leftrightarrow v_{769} \in \mathbb{R}) \to (\varphi_{17} \leftrightarrow 0 < v_{768} \to 0 < v_{769} \to \neg \forall u_{5}\,(\varphi_{27})) \to \varphi_{125}$
\item $\varphi_{209} :\equiv \mathrm{FLim}(g,c,\beta) \to r \in \mathbb{R} \to 0 < r \to \varphi_{42} \to d \in \mathbb{R} \land 0 < d \land \varphi_{31} \to \varphi_{139}$
\item $\varphi_{210} :\equiv (f \circ a(n) = f(a_{n}) \to f(a_{n}) = f \circ a(n)) \to \mathrm{FLim}(f,c,\alpha) \to \varphi_{188}$
\item $\varphi_{211} :\equiv (g \circ a(n) = g(a_{n}) \to g(a_{n}) = g \circ a(n)) \to \mathrm{FLim}(f,c,\alpha) \to \varphi_{191}$
\item $\varphi_{212} :\equiv (\mathrm{FLim}(f,c,\alpha) \to \varphi_{127}) \to \varphi_{81} \to \neg (d \in \mathbb{R} \land 0 < d \land \varphi_{31})$
\item $\varphi_{213} :\equiv \forall d\,(\varphi_{81} \to \neg (d \in \mathbb{R} \land 0 < d \land \varphi_{31})) \to \forall d\,(\varphi_{81}) \to \forall d\,(\neg (d \in \mathbb{R} \land 0 < d \land \varphi_{31}))$
\item $\varphi_{214} :\equiv (v_{768} \in \mathbb{R} \leftrightarrow v_{768} \in \mathbb{R}) \to \varphi_{125} \to (v_{768} \in \mathbb{R} \to v_{769} \in \mathbb{R} \to \varphi_{17} \leftrightarrow \varphi_{73})$
\item $\varphi_{215} :\equiv (v_{768} \in \mathbb{R} \to v_{769} \in \mathbb{R} \to \varphi_{17} \leftrightarrow \varphi_{73}) \to \varphi_{73} \to v_{768} \in \mathbb{R} \to v_{769} \in \mathbb{R} \to \varphi_{17}$
\item $\varphi_{216} :\equiv (v_{768} \in \mathbb{R} \to v_{769} \in \mathbb{R} \to \varphi_{17} \leftrightarrow \varphi_{73}) \to (v_{768} \in \mathbb{R} \to v_{769} \in \mathbb{R} \to \varphi_{17}) \to \varphi_{73}$
\item $\varphi_{217} :\equiv \forall v_{769}\,(\varphi_{73} \to v_{768} \in \mathbb{R} \to v_{769} \in \mathbb{R} \to \varphi_{17}) \to \forall v_{769}\,(\varphi_{73}) \to \forall v_{769}\,(v_{768} \in \mathbb{R} \to v_{769} \in \mathbb{R} \to \varphi_{17})$
\item $\varphi_{218} :\equiv \forall v_{769}\,((v_{768} \in \mathbb{R} \to v_{769} \in \mathbb{R} \to \varphi_{17}) \to \varphi_{73}) \to \forall v_{769}\,(v_{768} \in \mathbb{R} \to v_{769} \in \mathbb{R} \to \varphi_{17}) \to \forall v_{769}\,(\varphi_{73})$
\item $\varphi_{219} :\equiv (\mathrm{FLim}(f,c,\alpha) \to \varphi_{75}) \to (\mathrm{FLim}(f,c,\alpha) \to \varphi_{74}) \to \mathrm{FLim}(f,c,\alpha) \to \varphi_{123}$
\item $\varphi_{220} :\equiv (\forall v_{769}\,(v_{768} \in \mathbb{R} \to v_{769} \in \mathbb{R} \to \varphi_{17}) \leftrightarrow \forall v_{769}\,(\varphi_{73})) \to \forall v_{769}\,(\varphi_{73}) \to \forall v_{769}\,(v_{768} \in \mathbb{R} \to v_{769} \in \mathbb{R} \to \varphi_{17})$
\item $\varphi_{221} :\equiv (\forall v_{769}\,(v_{768} \in \mathbb{R} \to v_{769} \in \mathbb{R} \to \varphi_{17}) \leftrightarrow \forall v_{769}\,(\varphi_{73})) \to \forall v_{769}\,(v_{768} \in \mathbb{R} \to v_{769} \in \mathbb{R} \to \varphi_{17}) \to \forall v_{769}\,(\varphi_{73})$
\item $\varphi_{222} :\equiv (\mathrm{FLim}(f,c,\alpha) \to \mathrm{FLim}(g,c,\beta) \to \varphi_{158}) \to \mathrm{FLim}(f,c,\alpha) \to \varphi_{128}$
\item $\varphi_{223} :\equiv \forall \rho\,(\varphi_{133} \to \neg (\rho \in \mathbb{R} \land 0 < \rho \land \rho \leq d \land \rho \leq r)) \to \forall \rho\,(\varphi_{133}) \to \forall \rho\,(\neg (\rho \in \mathbb{R} \land 0 < \rho \land \rho \leq d \land \rho \leq r))$
\item $\varphi_{224} :\equiv (\mathrm{FLim}(f,c,\alpha) \to \mathrm{FLim}(g,c,\beta) \to \varphi_{158}) \to \mathrm{FLim}(f,c,\alpha) \to \varphi_{131}$
\item $\varphi_{225} :\equiv \forall u_{5}\,(\varphi_{36}) \to (\forall u_{5}\,(\varphi_{36}) \to \varphi_{44}) \to \varphi_{49} \to \mathrm{FLim}(f,c,\alpha) \to \varphi_{127}$
\item $\varphi_{226} :\equiv (\forall u_{5}\,(\forall d\,(\varphi_{23}) \to \varphi_{27}) \to \forall u_{5}\,\forall d\,(\varphi_{23}) \to \forall u_{5}\,(\varphi_{27})) \to (\forall d\,(\varphi_{23}) \to \forall u_{5}\,\forall d\,(\varphi_{23})) \to \forall d\,(\varphi_{23}) \to \forall u_{5}\,(\varphi_{27})$
\item $\varphi_{227} :\equiv (\forall d\,(\forall u_{5}\,(\varphi_{27}) \to \varphi_{23}) \to \forall d\,\forall u_{5}\,(\varphi_{27}) \to \forall d\,(\varphi_{23})) \to (\forall u_{5}\,(\varphi_{27}) \to \forall d\,\forall u_{5}\,(\varphi_{27})) \to \forall u_{5}\,(\varphi_{27}) \to \forall d\,(\varphi_{23})$
\item $\varphi_{228} :\equiv \varphi_{122} \to (\varphi_{42} \to a_{n} \in \mathrm{dom}(f) \to \varphi_{54}) \to \mathrm{FLim}(f,c,\alpha) \to \varphi_{136}$
\item $\varphi_{229} :\equiv (\mathrm{FLim}(f,c,\alpha) \to \varphi_{168}) \to \varphi_{162} \to n \in \mathbb{N} \to a_{n} \in \mathrm{dom}(f)$
\item $\varphi_{230} :\equiv (\mathrm{FLim}(f,c,\alpha) \to \varphi_{169}) \to \varphi_{162} \to n \in \mathbb{N} \to a_{n} \in \mathrm{dom}(g)$
\item $\varphi_{231} :\equiv \forall n\,(\varphi_{162} \to n \in \mathbb{N} \to a_{n} \in \mathrm{dom}(f)) \to \forall n\,(\varphi_{162}) \to \forall n \in \mathbb{N}\,(a_{n} \in \mathrm{dom}(f))$
\item $\varphi_{232} :\equiv \forall n\,(\varphi_{162} \to n \in \mathbb{N} \to a_{n} \in \mathrm{dom}(g)) \to \forall n\,(\varphi_{162}) \to \forall n \in \mathbb{N}\,(a_{n} \in \mathrm{dom}(g))$
\item $\varphi_{233} :\equiv \varphi_{217} \to (\forall v_{769}\,(v_{768} \in \mathbb{R} \to v_{769} \in \mathbb{R} \to \varphi_{17}) \leftrightarrow \forall v_{769}\,(\varphi_{73}))$
\item $\varphi_{234} :\equiv (\mathrm{FLim}(f,c,\alpha) \to \varphi_{185}) \to \varphi_{162} \to n \in \mathbb{N} \to a_{n} \in \mathrm{dom}(f) \land a_{n} \neq c$
\item $\varphi_{235} :\equiv (\mathrm{FLim}(f,c,\alpha) \to \varphi_{186}) \to \varphi_{162} \to n \in \mathbb{N} \to a_{n} \in \mathrm{dom}(g) \land a_{n} \neq c$
\item $\varphi_{236} :\equiv \forall n\,(\varphi_{162} \to n \in \mathbb{N} \to a_{n} \in \mathrm{dom}(f) \land a_{n} \neq c) \to \forall n\,(\varphi_{162}) \to \forall n \in \mathbb{N}\,(a_{n} \in \mathrm{dom}(f) \land a_{n} \neq c)$
\item $\varphi_{237} :\equiv \forall n\,(\varphi_{162} \to n \in \mathbb{N} \to a_{n} \in \mathrm{dom}(g) \land a_{n} \neq c) \to \forall n\,(\varphi_{162}) \to \forall n \in \mathbb{N}\,(a_{n} \in \mathrm{dom}(g) \land a_{n} \neq c)$
\item $\varphi_{238} :\equiv (\mathrm{FLim}(f,c,\alpha) \to \varphi_{194}) \to \varphi_{162} \to n \in \mathbb{N} \to f \circ a(n) \leq g \circ a(n)$
\item $\varphi_{239} :\equiv \forall n\,(\varphi_{162} \to n \in \mathbb{N} \to f \circ a(n) \leq g \circ a(n)) \to \forall n\,(\varphi_{162}) \to \forall n \in \mathbb{N}\,(f \circ a(n) \leq g \circ a(n))$
\item $\varphi_{240} :\equiv (\varphi_{147} \to \varphi_{61} \to \forall v_{768}\,\forall v_{769}\,(v_{768} \in \mathbb{R} \to v_{769} \in \mathbb{R} \to \varphi_{17})) \to (\forall v_{768}\,\forall v_{769}\,(v_{768} \in \mathbb{R} \to v_{769} \in \mathbb{R} \to \varphi_{17}) \leftrightarrow \varphi_{61})$
\item $\varphi_{241} :\equiv (\mathrm{FLim}(f,c,\alpha) \to \varphi_{131}) \to \varphi_{4} \to (\mathrm{FLim}(f,c,\alpha) \to \varphi_{128}) \to \varphi_{11} \to \varphi_{149}$
\item $\varphi_{242} :\equiv (\mathrm{FLim}(f,c,\alpha) \to \varphi_{192}) \to \varphi_{47} \to \mathrm{FLim}(f,c,\alpha) \to \varphi_{185}$
\item $\varphi_{243} :\equiv (\mathrm{FLim}(f,c,\alpha) \to \varphi_{193}) \to \varphi_{48} \to \mathrm{FLim}(f,c,\alpha) \to \varphi_{186}$
\item $\varphi_{244} :\equiv (\mathrm{FLim}(f,c,\alpha) \to \varphi_{199}) \to \varphi_{52} \to \mathrm{FLim}(f,c,\alpha) \to \varphi_{197}$
\item $\varphi_{245} :\equiv (\mathrm{FLim}(f,c,\alpha) \to \varphi_{200}) \to \varphi_{53} \to \mathrm{FLim}(f,c,\alpha) \to \varphi_{198}$
\item $\varphi_{246} :\equiv (\mathrm{FLim}(f,c,\alpha) \to \varphi_{202}) \to \varphi_{71} \to \mathrm{FLim}(f,c,\alpha) \to \varphi_{201}$
\item $\varphi_{247} :\equiv (\mathrm{FLim}(f,c,\alpha) \to \varphi_{182}) \to \varphi_{60} \to \varphi_{56} \to \mathrm{FLim}(f,c,\alpha) \to \varphi_{189}$
\item $\varphi_{248} :\equiv (\mathrm{FLim}(f,c,\alpha) \to \varphi_{189}) \to \varphi_{63} \to \varphi_{62} \to \mathrm{FLim}(f,c,\alpha) \to \varphi_{194}$
\item $\varphi_{249} :\equiv (\mathrm{FLim}(f,c,\alpha) \to \varphi_{180}) \to \varphi_{170} \to \mathrm{FLim}(f,c,\alpha) \to \varphi_{174}$
\item $\varphi_{250} :\equiv (\mathrm{FLim}(f,c,\alpha) \to \varphi_{180}) \to \varphi_{171} \to \mathrm{FLim}(f,c,\alpha) \to \varphi_{175}$
\item $\varphi_{251} :\equiv \forall d\,(\varphi_{81} \to \neg (d \in \mathbb{R} \land 0 < d \land \varphi_{31})) \to \varphi_{213} \to (\varphi_{81} \to \forall d\,(\varphi_{81})) \to \varphi_{78}$
\item $\varphi_{252} :\equiv (\mathrm{FLim}(f,c,\alpha) \to \varphi_{185}) \to \varphi_{177} \to \mathrm{FLim}(f,c,\alpha) \to \varphi_{168}$
\item $\varphi_{253} :\equiv (\mathrm{FLim}(f,c,\alpha) \to \varphi_{186}) \to \varphi_{177} \to \mathrm{FLim}(f,c,\alpha) \to \varphi_{169}$
\item $\varphi_{254} :\equiv (\mathrm{FLim}(f,c,\alpha) \to \mathrm{FLim}(f,c,\alpha)) \to \varphi_{166} \to (\mathrm{FLim}(f,c,\alpha) \to \varphi_{192}) \to \varphi_{179}$
\item $\varphi_{255} :\equiv (\mathrm{FLim}(f,c,\alpha) \to \varphi_{195}) \to \varphi_{181} \to \mathrm{FLim}(f,c,\alpha) \to \varphi_{185}$
\item $\varphi_{256} :\equiv (\mathrm{FLim}(f,c,\alpha) \to \varphi_{195}) \to \varphi_{181} \to \mathrm{FLim}(f,c,\alpha) \to \varphi_{186}$
\item $\varphi_{257} :\equiv (\mathrm{FLim}(f,c,\alpha) \to \varphi_{197}) \to \varphi_{177} \to \mathrm{FLim}(f,c,\alpha) \to \varphi_{187}$
\item $\varphi_{258} :\equiv (\mathrm{FLim}(f,c,\alpha) \to \varphi_{198}) \to \varphi_{177} \to \mathrm{FLim}(f,c,\alpha) \to \varphi_{190}$
\item $\varphi_{259} :\equiv (\mathrm{FLim}(f,c,\alpha) \to \varphi_{201}) \to (\mathrm{FLim}(f,c,\alpha) \to \varphi_{174}) \to \mathrm{FLim}(f,c,\alpha) \to \varphi_{182}$
\item $\varphi_{260} :\equiv (\mathrm{FLim}(f,c,\alpha) \to \varphi_{201}) \to (\mathrm{FLim}(f,c,\alpha) \to \varphi_{175}) \to \mathrm{FLim}(f,c,\alpha) \to \varphi_{195}$
\item $\varphi_{261} :\equiv (\mathrm{FLim}(g,c,\beta) \to \mathrm{Conv}(a,c) \to \varphi_{24}) \to \varphi_{1} \to \varphi_{166} \to (\mathrm{FLim}(f,c,\alpha) \to \varphi_{193}) \to \varphi_{178}$
\item $\varphi_{262} :\equiv (a_{n} \in \mathbb{R} \to c \in \mathbb{R} \to 0 < \lvert a_{n} - c \rvert \to a_{n} \neq c) \to \varphi_{183} \to \varphi_{82} \to (\mathrm{FLim}(f,c,\alpha) \to \varphi_{176}) \to \varphi_{181}$
\item $\varphi_{263} :\equiv (\varphi_{148} \to \forall v_{768}\,\forall v_{769}\,(v_{768} \in \mathbb{R} \to v_{769} \in \mathbb{R} \to \varphi_{17}) \to \varphi_{61}) \to \varphi_{240}$
\item $\varphi_{264} :\equiv (\mathrm{RFn}(f) \to \mathrm{Seq}(a) \to \varphi_{57}) \to \varphi_{79} \to \varphi_{165} \to (\mathrm{FLim}(f,c,\alpha) \to \varphi_{172}) \to \varphi_{203}$
\item $\varphi_{265} :\equiv (\mathrm{RFn}(g) \to \mathrm{Seq}(a) \to \varphi_{58}) \to \varphi_{80} \to \varphi_{165} \to (\mathrm{FLim}(f,c,\alpha) \to \varphi_{173}) \to \varphi_{204}$
\item $\varphi_{266} :\equiv (\forall a\,(\varphi_{151} \to \neg \varphi_{39}) \to \forall a\,(\varphi_{151}) \to \forall a\,(\neg \varphi_{39})) \to (\varphi_{151} \to \forall a\,(\varphi_{151})) \to \varphi_{150}$
\item $\varphi_{267} :\equiv (\mathrm{Conv}(f \circ a,\alpha) \to \varphi_{38}) \to \varphi_{179} \to \varphi_{178} \to (\mathrm{FLim}(f,c,\alpha) \to \varphi_{196}) \to \varphi_{164}$
\item $\varphi_{268} :\equiv \forall v_{769}\,((v_{768} \in \mathbb{R} \to v_{769} \in \mathbb{R} \to \varphi_{17}) \to \varphi_{73}) \to \forall v_{769}\,(\varphi_{73} \to v_{768} \in \mathbb{R} \to v_{769} \in \mathbb{R} \to \varphi_{17}) \to \varphi_{218} \to \varphi_{233}$
\item $\varphi_{269} :\equiv (\mathrm{FLim}(f,c,\alpha) \to \varphi_{136}) \to (\mathrm{FLim}(f,c,\alpha) \to \varphi_{195}) \to \varphi_{259}$
\item $\varphi_{270} :\equiv \forall \rho\,(\varphi_{133} \to \neg (\rho \in \mathbb{R} \land 0 < \rho \land \rho \leq d \land \rho \leq r)) \to \varphi_{223} \to (\varphi_{133} \to \forall \rho\,(\varphi_{133})) \to \varphi_{225}$
\item $\varphi_{271} :\equiv \forall n\,(\varphi_{162} \to n \in \mathbb{N} \to a_{n} \in \mathrm{dom}(f)) \to \varphi_{231} \to (\varphi_{162} \to \forall n\,(\varphi_{162})) \to \mathrm{FLim}(f,c,\alpha) \to \varphi_{172}$
\item $\varphi_{272} :\equiv \forall n\,(\varphi_{162} \to n \in \mathbb{N} \to a_{n} \in \mathrm{dom}(g)) \to \varphi_{232} \to (\varphi_{162} \to \forall n\,(\varphi_{162})) \to \mathrm{FLim}(f,c,\alpha) \to \varphi_{173}$
\item $\varphi_{273} :\equiv \forall n\,(\varphi_{162} \to n \in \mathbb{N} \to a_{n} \in \mathrm{dom}(f) \land a_{n} \neq c) \to \varphi_{236} \to (\varphi_{162} \to \forall n\,(\varphi_{162})) \to \mathrm{FLim}(f,c,\alpha) \to \varphi_{192}$
\item $\varphi_{274} :\equiv \forall n\,(\varphi_{162} \to n \in \mathbb{N} \to a_{n} \in \mathrm{dom}(g) \land a_{n} \neq c) \to \varphi_{237} \to (\varphi_{162} \to \forall n\,(\varphi_{162})) \to \mathrm{FLim}(f,c,\alpha) \to \varphi_{193}$
\item $\varphi_{275} :\equiv \forall n\,(\varphi_{162} \to n \in \mathbb{N} \to f \circ a(n) \leq g \circ a(n)) \to \varphi_{239} \to (\varphi_{162} \to \forall n\,(\varphi_{162})) \to \mathrm{FLim}(f,c,\alpha) \to \varphi_{196}$
\end{abbrevlist}
\end{fullproof}

% 自動生成 (prover/build.js): FCauchy  ノード数 946
\begin{fullproof}{定理1.13}{$\forall f\,\forall c\,(\mathrm{RFn}(f) \to c \in \mathbb{R} \to \mathrm{PNbhd}(f,c) \to (\exists \alpha\,(\mathrm{FLim}(f,c,\alpha)) \leftrightarrow \forall \varepsilon \in \mathbb{R}\,(0 < \varepsilon \to \exists \delta\,(\delta \in \mathbb{R} \land 0 < \delta \land \forall x\,\forall y\,(x \in \mathrm{dom}(f) \to y \in \mathrm{dom}(f) \to 0 < \lvert x - c \rvert \to \lvert x - c \rvert < \delta \to 0 < \lvert y - c \rvert \to \lvert y - c \rvert < \delta \to \lvert f(x) - f(y) \rvert < \varepsilon)))))$}{946}\phantomsection\label{tree:FCauchy}
\noindent$\Pi_{1}$:
\begin{ptree}
\AXC{$z = w \to \lvert z - \alpha \rvert < \varepsilon \cdot 2^{-1} \to \lvert w - \alpha \rvert < \varepsilon \cdot 2^{-1}$\lbl{A6}}
\UIC{$\varphi_{25}$}
\UIC{$\forall z\,(\varphi_{25})$}
\AXC{$\forall z\,(\varphi_{25}) \to \varphi_{143}$\lbl{A2}}
\BIC{$\varphi_{143}$}
\AXC{$\varphi_{143} \to f \circ p(N_{0} \cup N_{1}) = f(p(N_{0} \cup N_{1})) \to \varphi_{129}$\lbl{A2}}
\BIC{$f \circ p(N_{0} \cup N_{1}) = f(p(N_{0} \cup N_{1})) \to \varphi_{129}$}
\end{ptree}
\noindent$\Pi_{2}$:
\begin{ptree}
\AXC{$\forall m\,\forall n\,(n \leq m \cup n)$\lbl{[C8]}}
\AXC{$\forall m\,\forall n\,(n \leq m \cup n) \to \forall n\,(n \leq N_{0} \cup n)$\lbl{A2}}
\BIC{$\forall n\,(n \leq N_{0} \cup n)$}
\AXC{$\forall n\,(n \leq N_{0} \cup n) \to N_{1} \leq N_{0} \cup N_{1}$\lbl{A2}}
\BIC{$N_{1} \leq N_{0} \cup N_{1}$}
\end{ptree}
\noindent$\Pi_{3}$:
\begin{ptree}
\AXC{$\forall m\,\forall n\,(m \in \mathbb{N} \to n \in \mathbb{N} \to m \cup n \in \mathbb{N})$\lbl{[N3]}}
\AXC{$\varphi_{101}$\lbl{A2}}
\BIC{$\forall n\,(N_{0} \in \mathbb{N} \to n \in \mathbb{N} \to N_{0} \cup n \in \mathbb{N})$}
\AXC{$\varphi_{94}$\lbl{A2}}
\BIC{$N_{0} \in \mathbb{N} \to N_{1} \in \mathbb{N} \to N_{0} \cup N_{1} \in \mathbb{N}$}
\AXC{$\varphi_{749}$\lbl{A1}}
\BIC{$(\mathrm{RFn}(f) \to c \in \mathbb{R} \to \mathrm{PNbhd}(f,c) \to \varphi_{560}) \to \varphi_{705}$}
\end{ptree}
\noindent$\Pi_{4}$:
\begin{ptree}
\AXC{$\mathrm{RFn}(f) \to c \in \mathbb{R} \to \mathrm{PNbhd}(f,c) \to \varphi_{561}$\lbl{A1}}
\AXC{$\varphi_{704}$\lbl{A1}}
\BIC{$\mathrm{RFn}(f) \to c \in \mathbb{R} \to \mathrm{PNbhd}(f,c) \to \varphi_{560}$}
\AXC{$\stackrel{\Pi_{3}}{\vdots}$}\noLine
\UIC{$(\mathrm{RFn}(f) \to c \in \mathbb{R} \to \mathrm{PNbhd}(f,c) \to \varphi_{560}) \to \varphi_{705}$}
\BIC{$\varphi_{705}$}
\end{ptree}
\noindent$\Pi_{5}$:
\begin{ptree}
\AXC{$\mathrm{RFn}(f) \to c \in \mathbb{R} \to \mathrm{PNbhd}(f,c) \to \varphi_{594}$\lbl{A1}}
\AXC{$\varphi_{713}$\lbl{A1}}
\BIC{$\mathrm{RFn}(f) \to c \in \mathbb{R} \to \mathrm{PNbhd}(f,c) \to \varphi_{559}$}
\AXC{$\stackrel{\Pi_{4}}{\vdots}$}\noLine
\UIC{$\varphi_{705}$}
\BIC{$\mathrm{RFn}(f) \to c \in \mathbb{R} \to \mathrm{PNbhd}(f,c) \to \varphi_{562}$}
\end{ptree}
\noindent$\Pi_{6}$:
\begin{ptree}
\AXC{$\mathrm{RFn}(f) \to c \in \mathbb{R} \to \mathrm{PNbhd}(f,c) \to \varphi_{594}$\lbl{A1}}
\AXC{$\varphi_{719}$\lbl{A1}}
\BIC{$\mathrm{RFn}(f) \to c \in \mathbb{R} \to \mathrm{PNbhd}(f,c) \to \varphi_{585}$}
\AXC{$(\mathrm{RFn}(f) \to c \in \mathbb{R} \to \mathrm{PNbhd}(f,c) \to \varphi_{585}) \to \varphi_{640}$\lbl{A1}}
\BIC{$\varphi_{640}$}
\end{ptree}
\noindent$\Pi_{7}$:
\begin{ptree}
\AXC{$\forall n \in \mathbb{N}\,(N_{1} \leq n \to \lvert f \circ p(n) - \alpha \rvert < \varepsilon \cdot 2^{-1}) \to \varphi_{96}$\lbl{A2}}
\AXC{$\stackrel{\Pi_{6}}{\vdots}$}\noLine
\UIC{$\varphi_{640}$}
\BIC{$\mathrm{RFn}(f) \to c \in \mathbb{R} \to \mathrm{PNbhd}(f,c) \to \varphi_{602}$}
\AXC{$\varphi_{752}$\lbl{A1}}
\BIC{$(\mathrm{RFn}(f) \to c \in \mathbb{R} \to \mathrm{PNbhd}(f,c) \to \varphi_{562}) \to \varphi_{606}$}
\end{ptree}
\noindent$\Pi_{8}$:
\begin{ptree}
\AXC{$\stackrel{\Pi_{2}}{\vdots}$}\noLine
\UIC{$N_{1} \leq N_{0} \cup N_{1}$}
\AXC{$\stackrel{\Pi_{5}}{\vdots}$}\noLine
\UIC{$\mathrm{RFn}(f) \to c \in \mathbb{R} \to \mathrm{PNbhd}(f,c) \to \varphi_{562}$}
\AXC{$\stackrel{\Pi_{7}}{\vdots}$}\noLine
\UIC{$(\mathrm{RFn}(f) \to c \in \mathbb{R} \to \mathrm{PNbhd}(f,c) \to \varphi_{562}) \to \varphi_{606}$}
\BIC{$\varphi_{606}$}
\BIC{$\mathrm{RFn}(f) \to c \in \mathbb{R} \to \mathrm{PNbhd}(f,c) \to \varphi_{578}$}
\end{ptree}
\noindent$\Pi_{9}$:
\begin{ptree}
\AXC{$\forall a\,\forall n\,(\mathrm{Seq}(a) \to n \in \mathbb{N} \to a_{n} \in \mathbb{R})$\lbl{[S1]}}
\AXC{$\varphi_{81}$\lbl{A2}}
\BIC{$\forall n\,(\mathrm{Seq}(p) \to n \in \mathbb{N} \to p(n) \in \mathbb{R})$}
\AXC{$\varphi_{64}$\lbl{A2}}
\BIC{$\mathrm{Seq}(p) \to n \in \mathbb{N} \to p(n) \in \mathbb{R}$}
\AXC{$\varphi_{665}$\lbl{A1}}
\BIC{$(\mathrm{RFn}(f) \to c \in \mathbb{R} \to \mathrm{PNbhd}(f,c) \to \varphi_{343}) \to \varphi_{601}$}
\end{ptree}
\noindent$\Pi_{10}$:
\begin{ptree}
\AXC{$\mathrm{RFn}(f) \to c \in \mathbb{R} \to \mathrm{PNbhd}(f,c) \to \varphi_{401}$\lbl{A1}}
\AXC{$\varphi_{618}$\lbl{A1}}
\BIC{$\mathrm{RFn}(f) \to c \in \mathbb{R} \to \mathrm{PNbhd}(f,c) \to \varphi_{343}$}
\AXC{$\stackrel{\Pi_{9}}{\vdots}$}\noLine
\UIC{$(\mathrm{RFn}(f) \to c \in \mathbb{R} \to \mathrm{PNbhd}(f,c) \to \varphi_{343}) \to \varphi_{601}$}
\BIC{$\varphi_{601}$}
\end{ptree}
\noindent$\Pi_{11}$:
\begin{ptree}
\AXC{$\mathrm{RFn}(f) \to c \in \mathbb{R} \to \mathrm{PNbhd}(f,c) \to \varphi_{336}$\lbl{A1}}
\AXC{$\varphi_{614}$\lbl{A1}}
\BIC{$\mathrm{RFn}(f) \to c \in \mathbb{R} \to \mathrm{PNbhd}(f,c) \to \varphi_{390}$}
\AXC{$(\mathrm{RFn}(f) \to c \in \mathbb{R} \to \mathrm{PNbhd}(f,c) \to \varphi_{390}) \to \varphi_{500}$\lbl{A1}}
\BIC{$\varphi_{500}$}
\end{ptree}
\noindent$\Pi_{12}$:
\begin{ptree}
\AXC{$\mathrm{RFn}(f) \to c \in \mathbb{R} \to \mathrm{PNbhd}(f,c) \to \varphi_{355}$\lbl{A1}}
\AXC{$\stackrel{\Pi_{10}}{\vdots}$}\noLine
\UIC{$\varphi_{601}$}
\BIC{$\mathrm{RFn}(f) \to c \in \mathbb{R} \to \mathrm{PNbhd}(f,c) \to \varphi_{359}$}
\AXC{$\varphi_{16} \to p(n) \in \mathbb{R} \to \varphi_{27}$\lbl{A2}}
\AXC{$\stackrel{\Pi_{11}}{\vdots}$}\noLine
\UIC{$\varphi_{500}$}
\BIC{$\mathrm{RFn}(f) \to c \in \mathbb{R} \to \mathrm{PNbhd}(f,c) \to \varphi_{422}$}
\AXC{$\varphi_{729}$\lbl{A1}}
\BIC{$(\mathrm{RFn}(f) \to c \in \mathbb{R} \to \mathrm{PNbhd}(f,c) \to \varphi_{359}) \to \varphi_{620}$}
\BIC{$\varphi_{620}$}
\end{ptree}
\noindent$\Pi_{13}$:
\begin{ptree}
\AXC{$\varphi_{146}$\lbl{A2}}
\AXC{$\mathrm{RFn}(f) \to c \in \mathbb{R} \to \mathrm{PNbhd}(f,c) \to \varphi_{401}$\lbl{A1}}
\AXC{$\varphi_{627}$\lbl{A1}}
\BIC{$\mathrm{RFn}(f) \to c \in \mathbb{R} \to \mathrm{PNbhd}(f,c) \to \varphi_{381}$}
\AXC{$\varphi_{641}$\lbl{A1}}
\BIC{$\varphi_{146} \to \mathrm{RFn}(f) \to c \in \mathbb{R} \to \mathrm{PNbhd}(f,c) \to \varphi_{379}$}
\BIC{$\mathrm{RFn}(f) \to c \in \mathbb{R} \to \mathrm{PNbhd}(f,c) \to \varphi_{379}$}
\AXC{$\stackrel{\Pi_{12}}{\vdots}$}\noLine
\UIC{$\varphi_{620}$}
\BIC{$\mathrm{RFn}(f) \to c \in \mathbb{R} \to \mathrm{PNbhd}(f,c) \to \varphi_{361}$}
\AXC{$\varphi_{605}$\lbl{A1}}
\BIC{$\varphi_{362} \to n \in \mathbb{N} \to p(n) \in \mathrm{dom}(f)$}
\UIC{$\forall n\,(\varphi_{362} \to n \in \mathbb{N} \to p(n) \in \mathrm{dom}(f))$}
\end{ptree}
\noindent$\Pi_{14}$:
\begin{ptree}
\AXC{$\varphi_{610}$\lbl{A3}}
\AXC{$\stackrel{\Pi_{13}}{\vdots}$}\noLine
\UIC{$\forall n\,(\varphi_{362} \to n \in \mathbb{N} \to p(n) \in \mathrm{dom}(f))$}
\AXC{$\forall n\,(\varphi_{362} \to n \in \mathbb{N} \to p(n) \in \mathrm{dom}(f)) \to \varphi_{610} \to \varphi_{658}$\lbl{A1}}
\BIC{$\varphi_{610} \to \varphi_{658}$}
\BIC{$\varphi_{658}$}
\end{ptree}
\noindent$\Pi_{15}$:
\begin{ptree}
\AXC{$\forall f\,(\varphi_{147})$\lbl{[S8]}}
\AXC{$\forall f\,(\varphi_{147}) \to \varphi_{147}$\lbl{A2}}
\BIC{$\varphi_{147}$}
\AXC{$\varphi_{147} \to \mathrm{RFn}(f) \to \mathrm{Seq}(p) \to \varphi_{124}$\lbl{A2}}
\BIC{$\mathrm{RFn}(f) \to \mathrm{Seq}(p) \to \varphi_{124}$}
\AXC{$(\mathrm{RFn}(f) \to \mathrm{Seq}(p) \to \varphi_{124}) \to \varphi_{667}$\lbl{A1}}
\BIC{$\varphi_{667}$}
\end{ptree}
\noindent$\Pi_{16}$:
\begin{ptree}
\AXC{$\mathrm{RFn}(f) \to c \in \mathbb{R} \to \mathrm{PNbhd}(f,c) \to \varphi_{401}$\lbl{A1}}
\AXC{$\varphi_{618}$\lbl{A1}}
\BIC{$\mathrm{RFn}(f) \to c \in \mathbb{R} \to \mathrm{PNbhd}(f,c) \to \varphi_{343}$}
\AXC{$\mathrm{RFn}(f) \to \mathrm{RFn}(f)$\lbl{A1}}
\AXC{$\stackrel{\Pi_{15}}{\vdots}$}\noLine
\UIC{$\varphi_{667}$}
\BIC{$(\mathrm{RFn}(f) \to c \in \mathbb{R} \to \mathrm{PNbhd}(f,c) \to \varphi_{343}) \to \varphi_{622}$}
\BIC{$\varphi_{622}$}
\end{ptree}
\noindent$\Pi_{17}$:
\begin{ptree}
\AXC{$\varphi_{362} \to \forall n\,(\varphi_{362})$\lbl{A4}}
\AXC{$\stackrel{\Pi_{14}}{\vdots}$}\noLine
\UIC{$\varphi_{658}$}
\BIC{$\mathrm{RFn}(f) \to c \in \mathbb{R} \to \mathrm{PNbhd}(f,c) \to \varphi_{363}$}
\AXC{$\stackrel{\Pi_{16}}{\vdots}$}\noLine
\UIC{$\varphi_{622}$}
\BIC{$\mathrm{RFn}(f) \to c \in \mathbb{R} \to \mathrm{PNbhd}(f,c) \to \varphi_{384}$}
\end{ptree}
\noindent$\Pi_{18}$:
\begin{ptree}
\AXC{$\stackrel{\Pi_{17}}{\vdots}$}\noLine
\UIC{$\mathrm{RFn}(f) \to c \in \mathbb{R} \to \mathrm{PNbhd}(f,c) \to \varphi_{384}$}
\AXC{$\varphi_{666}$\lbl{A1}}
\BIC{$\mathrm{RFn}(f) \to c \in \mathbb{R} \to \mathrm{PNbhd}(f,c) \to \varphi_{574}$}
\AXC{$(\mathrm{RFn}(f) \to c \in \mathbb{R} \to \mathrm{PNbhd}(f,c) \to \varphi_{574}) \to \varphi_{632}$\lbl{A1}}
\BIC{$\varphi_{632}$}
\end{ptree}
\noindent$\Pi_{19}$:
\begin{ptree}
\AXC{$\forall n \in \mathbb{N}\,(f \circ p(n) = f(p(n))) \to \varphi_{46}$\lbl{A2}}
\AXC{$\stackrel{\Pi_{18}}{\vdots}$}\noLine
\UIC{$\varphi_{632}$}
\BIC{$\mathrm{RFn}(f) \to c \in \mathbb{R} \to \mathrm{PNbhd}(f,c) \to \varphi_{589}$}
\AXC{$(\mathrm{RFn}(f) \to c \in \mathbb{R} \to \mathrm{PNbhd}(f,c) \to \varphi_{589}) \to \varphi_{710}$\lbl{A1}}
\BIC{$\varphi_{710}$}
\end{ptree}
\noindent$\Pi_{20}$:
\begin{ptree}
\AXC{$\stackrel{\Pi_{8}}{\vdots}$}\noLine
\UIC{$\mathrm{RFn}(f) \to c \in \mathbb{R} \to \mathrm{PNbhd}(f,c) \to \varphi_{578}$}
\AXC{$\stackrel{\Pi_{5}}{\vdots}$}\noLine
\UIC{$\mathrm{RFn}(f) \to c \in \mathbb{R} \to \mathrm{PNbhd}(f,c) \to \varphi_{562}$}
\AXC{$\stackrel{\Pi_{19}}{\vdots}$}\noLine
\UIC{$\varphi_{710}$}
\BIC{$\mathrm{RFn}(f) \to c \in \mathbb{R} \to \mathrm{PNbhd}(f,c) \to \varphi_{577}$}
\AXC{$\varphi_{754}$\lbl{A1}}
\BIC{$(\mathrm{RFn}(f) \to c \in \mathbb{R} \to \mathrm{PNbhd}(f,c) \to \varphi_{578}) \to \varphi_{636}$}
\BIC{$\varphi_{636}$}
\end{ptree}
\noindent$\Pi_{21}$:
\begin{ptree}
\AXC{$\forall m\,\forall n\,(m \leq m \cup n)$\lbl{[C7]}}
\AXC{$\forall m\,\forall n\,(m \leq m \cup n) \to \forall n\,(N_{0} \leq N_{0} \cup n)$\lbl{A2}}
\BIC{$\forall n\,(N_{0} \leq N_{0} \cup n)$}
\AXC{$\forall n\,(N_{0} \leq N_{0} \cup n) \to N_{0} \leq N_{0} \cup N_{1}$\lbl{A2}}
\BIC{$N_{0} \leq N_{0} \cup N_{1}$}
\end{ptree}
\noindent$\Pi_{22}$:
\begin{ptree}
\AXC{$\mathrm{RFn}(f) \to c \in \mathbb{R} \to \mathrm{PNbhd}(f,c) \to \varphi_{561}$\lbl{A1}}
\AXC{$\varphi_{715}$\lbl{A1}}
\BIC{$\mathrm{RFn}(f) \to c \in \mathbb{R} \to \mathrm{PNbhd}(f,c) \to \varphi_{596}$}
\AXC{$(\mathrm{RFn}(f) \to c \in \mathbb{R} \to \mathrm{PNbhd}(f,c) \to \varphi_{596}) \to \varphi_{650}$\lbl{A1}}
\BIC{$\varphi_{650}$}
\end{ptree}
\noindent$\Pi_{23}$:
\begin{ptree}
\AXC{$\stackrel{\Pi_{5}}{\vdots}$}\noLine
\UIC{$\mathrm{RFn}(f) \to c \in \mathbb{R} \to \mathrm{PNbhd}(f,c) \to \varphi_{562}$}
\AXC{$\varphi_{77} \to N_{0} \cup N_{1} \in \mathbb{N} \to N_{0} \leq N_{0} \cup N_{1} \to \varphi_{61}$\lbl{A2}}
\AXC{$\stackrel{\Pi_{22}}{\vdots}$}\noLine
\UIC{$\varphi_{650}$}
\BIC{$\mathrm{RFn}(f) \to c \in \mathbb{R} \to \mathrm{PNbhd}(f,c) \to \varphi_{613}$}
\AXC{$\varphi_{753}$\lbl{A1}}
\BIC{$(\mathrm{RFn}(f) \to c \in \mathbb{R} \to \mathrm{PNbhd}(f,c) \to \varphi_{562}) \to \varphi_{619}$}
\BIC{$\varphi_{619}$}
\end{ptree}
\noindent$\Pi_{24}$:
\begin{ptree}
\AXC{$\mathrm{RFn}(f) \to c \in \mathbb{R} \to \mathrm{PNbhd}(f,c) \to \varphi_{401}$\lbl{A1}}
\AXC{$\varphi_{627}$\lbl{A1}}
\BIC{$\mathrm{RFn}(f) \to c \in \mathbb{R} \to \mathrm{PNbhd}(f,c) \to \varphi_{381}$}
\AXC{$(\mathrm{RFn}(f) \to c \in \mathbb{R} \to \mathrm{PNbhd}(f,c) \to \varphi_{381}) \to \varphi_{495}$\lbl{A1}}
\BIC{$\varphi_{495}$}
\end{ptree}
\noindent$\Pi_{25}$:
\begin{ptree}
\AXC{$\forall n \in \mathbb{N}\,(0 < \lvert p(n) - c \rvert \land \lvert p(n) - c \rvert < r) \to \varphi_{85}$\lbl{A2}}
\AXC{$\stackrel{\Pi_{24}}{\vdots}$}\noLine
\UIC{$\varphi_{495}$}
\BIC{$\mathrm{RFn}(f) \to c \in \mathbb{R} \to \mathrm{PNbhd}(f,c) \to \varphi_{404}$}
\AXC{$(\mathrm{RFn}(f) \to c \in \mathbb{R} \to \mathrm{PNbhd}(f,c) \to \varphi_{404}) \to \varphi_{712}$\lbl{A1}}
\BIC{$\varphi_{712}$}
\end{ptree}
\noindent$\Pi_{26}$:
\begin{ptree}
\AXC{$\stackrel{\Pi_{5}}{\vdots}$}\noLine
\UIC{$\mathrm{RFn}(f) \to c \in \mathbb{R} \to \mathrm{PNbhd}(f,c) \to \varphi_{562}$}
\AXC{$\stackrel{\Pi_{25}}{\vdots}$}\noLine
\UIC{$\varphi_{712}$}
\BIC{$\mathrm{RFn}(f) \to c \in \mathbb{R} \to \mathrm{PNbhd}(f,c) \to \varphi_{584}$}
\end{ptree}
\noindent$\Pi_{27}$:
\begin{ptree}
\AXC{$\forall a\,\forall n\,(\mathrm{Seq}(a) \to n \in \mathbb{N} \to a_{n} \in \mathbb{R})$\lbl{[S1]}}
\AXC{$\varphi_{81}$\lbl{A2}}
\BIC{$\forall n\,(\mathrm{Seq}(p) \to n \in \mathbb{N} \to p(n) \in \mathbb{R})$}
\AXC{$\varphi_{118}$\lbl{A2}}
\BIC{$\mathrm{Seq}(p) \to N_{0} \cup N_{1} \in \mathbb{N} \to p(N_{0} \cup N_{1}) \in \mathbb{R}$}
\AXC{$\varphi_{739}$\lbl{A1}}
\BIC{$(\mathrm{RFn}(f) \to c \in \mathbb{R} \to \mathrm{PNbhd}(f,c) \to \varphi_{343}) \to \varphi_{707}$}
\end{ptree}
\noindent$\Pi_{28}$:
\begin{ptree}
\AXC{$\mathrm{RFn}(f) \to c \in \mathbb{R} \to \mathrm{PNbhd}(f,c) \to \varphi_{401}$\lbl{A1}}
\AXC{$\varphi_{618}$\lbl{A1}}
\BIC{$\mathrm{RFn}(f) \to c \in \mathbb{R} \to \mathrm{PNbhd}(f,c) \to \varphi_{343}$}
\AXC{$\stackrel{\Pi_{27}}{\vdots}$}\noLine
\UIC{$(\mathrm{RFn}(f) \to c \in \mathbb{R} \to \mathrm{PNbhd}(f,c) \to \varphi_{343}) \to \varphi_{707}$}
\BIC{$\varphi_{707}$}
\end{ptree}
\noindent$\Pi_{29}$:
\begin{ptree}
\AXC{$\mathrm{RFn}(f) \to c \in \mathbb{R} \to \mathrm{PNbhd}(f,c) \to \varphi_{336}$\lbl{A1}}
\AXC{$\varphi_{661}$\lbl{A1}}
\BIC{$\mathrm{RFn}(f) \to c \in \mathbb{R} \to \mathrm{PNbhd}(f,c) \to \varphi_{583}$}
\AXC{$(\mathrm{RFn}(f) \to c \in \mathbb{R} \to \mathrm{PNbhd}(f,c) \to \varphi_{583}) \to \varphi_{649}$\lbl{A1}}
\BIC{$\varphi_{649}$}
\end{ptree}
\noindent$\Pi_{30}$:
\begin{ptree}
\AXC{$\stackrel{\Pi_{5}}{\vdots}$}\noLine
\UIC{$\mathrm{RFn}(f) \to c \in \mathbb{R} \to \mathrm{PNbhd}(f,c) \to \varphi_{562}$}
\AXC{$\stackrel{\Pi_{28}}{\vdots}$}\noLine
\UIC{$\varphi_{707}$}
\BIC{$\mathrm{RFn}(f) \to c \in \mathbb{R} \to \mathrm{PNbhd}(f,c) \to \varphi_{566}$}
\AXC{$\varphi_{16} \to p(N_{0} \cup N_{1}) \in \mathbb{R} \to \varphi_{112}$\lbl{A2}}
\AXC{$\stackrel{\Pi_{29}}{\vdots}$}\noLine
\UIC{$\varphi_{649}$}
\BIC{$\mathrm{RFn}(f) \to c \in \mathbb{R} \to \mathrm{PNbhd}(f,c) \to \varphi_{616}$}
\AXC{$\varphi_{766}$\lbl{A1}}
\BIC{$(\mathrm{RFn}(f) \to c \in \mathbb{R} \to \mathrm{PNbhd}(f,c) \to \varphi_{566}) \to \varphi_{714}$}
\BIC{$\varphi_{714}$}
\end{ptree}
\noindent$\Pi_{31}$:
\begin{ptree}
\AXC{$\stackrel{\Pi_{26}}{\vdots}$}\noLine
\UIC{$\mathrm{RFn}(f) \to c \in \mathbb{R} \to \mathrm{PNbhd}(f,c) \to \varphi_{584}$}
\AXC{$\stackrel{\Pi_{30}}{\vdots}$}\noLine
\UIC{$\varphi_{714}$}
\BIC{$\mathrm{RFn}(f) \to c \in \mathbb{R} \to \mathrm{PNbhd}(f,c) \to \varphi_{569}$}
\end{ptree}
\noindent$\Pi_{32}$:
\begin{ptree}
\AXC{$\mathrm{RFn}(f) \to c \in \mathbb{R} \to \mathrm{PNbhd}(f,c) \to \varphi_{557}$\lbl{A1}}
\AXC{$\varphi_{725}$\lbl{A1}}
\BIC{$\mathrm{RFn}(f) \to c \in \mathbb{R} \to \mathrm{PNbhd}(f,c) \to \varphi_{630}$}
\AXC{$(\mathrm{RFn}(f) \to c \in \mathbb{R} \to \mathrm{PNbhd}(f,c) \to \varphi_{630}) \to \varphi_{663}$\lbl{A1}}
\BIC{$\varphi_{663}$}
\end{ptree}
\noindent$\Pi_{33}$:
\begin{ptree}
\AXC{$\forall x\,(\varphi_{176}) \to \varphi_{176}$\lbl{A2}}
\AXC{$\stackrel{\Pi_{32}}{\vdots}$}\noLine
\UIC{$\varphi_{663}$}
\BIC{$\mathrm{RFn}(f) \to c \in \mathbb{R} \to \mathrm{PNbhd}(f,c) \to \varphi_{629}$}
\AXC{$(\mathrm{RFn}(f) \to c \in \mathbb{R} \to \mathrm{PNbhd}(f,c) \to \varphi_{629}) \to \varphi_{683}$\lbl{A1}}
\BIC{$\varphi_{683}$}
\end{ptree}
\noindent$\Pi_{34}$:
\begin{ptree}
\AXC{$\varphi_{176} \to x \in \mathrm{dom}(f) \to p(N_{0} \cup N_{1}) \in \mathrm{dom}(f) \to \varphi_{196}$\lbl{A2}}
\AXC{$\stackrel{\Pi_{33}}{\vdots}$}\noLine
\UIC{$\varphi_{683}$}
\BIC{$\mathrm{RFn}(f) \to c \in \mathbb{R} \to \mathrm{PNbhd}(f,c) \to \varphi_{637}$}
\AXC{$\varphi_{781}$\lbl{A1}}
\BIC{$(\mathrm{RFn}(f) \to c \in \mathbb{R} \to \mathrm{PNbhd}(f,c) \to \varphi_{568}) \to \varphi_{776}$}
\end{ptree}
\noindent$\Pi_{35}$:
\begin{ptree}
\AXC{$\mathrm{RFn}(f) \to c \in \mathbb{R} \to \mathrm{PNbhd}(f,c) \to \varphi_{568}$\lbl{A1}}
\AXC{$\stackrel{\Pi_{34}}{\vdots}$}\noLine
\UIC{$(\mathrm{RFn}(f) \to c \in \mathbb{R} \to \mathrm{PNbhd}(f,c) \to \varphi_{568}) \to \varphi_{776}$}
\BIC{$\varphi_{776}$}
\end{ptree}
\noindent$\Pi_{36}$:
\begin{ptree}
\AXC{$\mathrm{RFn}(f) \to c \in \mathbb{R} \to \mathrm{PNbhd}(f,c) \to \varphi_{575}$\lbl{A1}}
\AXC{$\stackrel{\Pi_{31}}{\vdots}$}\noLine
\UIC{$\mathrm{RFn}(f) \to c \in \mathbb{R} \to \mathrm{PNbhd}(f,c) \to \varphi_{569}$}
\AXC{$\stackrel{\Pi_{35}}{\vdots}$}\noLine
\UIC{$\varphi_{776}$}
\BIC{$(\mathrm{RFn}(f) \to c \in \mathbb{R} \to \mathrm{PNbhd}(f,c) \to \varphi_{575}) \to \varphi_{767}$}
\BIC{$\varphi_{767}$}
\end{ptree}
\noindent$\Pi_{37}$:
\begin{ptree}
\AXC{$\stackrel{\Pi_{26}}{\vdots}$}\noLine
\UIC{$\mathrm{RFn}(f) \to c \in \mathbb{R} \to \mathrm{PNbhd}(f,c) \to \varphi_{584}$}
\AXC{$\mathrm{RFn}(f) \to c \in \mathbb{R} \to \mathrm{PNbhd}(f,c) \to \varphi_{586}$\lbl{A1}}
\AXC{$\stackrel{\Pi_{36}}{\vdots}$}\noLine
\UIC{$\varphi_{767}$}
\BIC{$(\mathrm{RFn}(f) \to c \in \mathbb{R} \to \mathrm{PNbhd}(f,c) \to \varphi_{584}) \to \varphi_{724}$}
\BIC{$\varphi_{724}$}
\end{ptree}
\noindent$\Pi_{38}$:
\begin{ptree}
\AXC{$\forall x\,\forall y\,(x \in \mathbb{R} \to y \in \mathbb{R} \to x \cdot y \in \mathbb{R})$\lbl{(R1)}}
\AXC{$\varphi_{111}$\lbl{A2}}
\BIC{$\forall y\,(\varepsilon \in \mathbb{R} \to y \in \mathbb{R} \to \varepsilon \cdot y \in \mathbb{R})$}
\AXC{$\varphi_{123}$\lbl{A2}}
\BIC{$\varepsilon \in \mathbb{R} \to 2^{-1} \in \mathbb{R} \to \varepsilon \cdot 2^{-1} \in \mathbb{R}$}
\AXC{$\varphi_{633}$\lbl{A1}}
\BIC{$(\mathrm{RFn}(f) \to c \in \mathbb{R} \to \mathrm{PNbhd}(f,c) \to \varphi_{373}) \to \varphi_{433}$}
\end{ptree}
\noindent$\Pi_{39}$:
\begin{ptree}
\AXC{$2^{-1} \in \mathbb{R}$\lbl{[F31]}}
\AXC{$\mathrm{RFn}(f) \to c \in \mathbb{R} \to \mathrm{PNbhd}(f,c) \to \varphi_{373}$\lbl{A1}}
\AXC{$\stackrel{\Pi_{38}}{\vdots}$}\noLine
\UIC{$(\mathrm{RFn}(f) \to c \in \mathbb{R} \to \mathrm{PNbhd}(f,c) \to \varphi_{373}) \to \varphi_{433}$}
\BIC{$\varphi_{433}$}
\BIC{$\mathrm{RFn}(f) \to c \in \mathbb{R} \to \mathrm{PNbhd}(f,c) \to \varphi_{385}$}
\end{ptree}
\noindent$\Pi_{40}$:
\begin{ptree}
\AXC{$\forall a\,\forall \alpha\,(\mathrm{Conv}(a,\alpha) \to \alpha \in \mathbb{R})$\lbl{[C2]}}
\AXC{$\varphi_{39}$\lbl{A2}}
\BIC{$\forall \alpha\,(\mathrm{Conv}(f \circ p,\alpha) \to \alpha \in \mathbb{R})$}
\AXC{$\varphi_{53}$\lbl{A2}}
\BIC{$\mathrm{Conv}(f \circ p,\alpha) \to \alpha \in \mathbb{R}$}
\AXC{$(\mathrm{Conv}(f \circ p,\alpha) \to \alpha \in \mathbb{R}) \to \varphi_{615}$\lbl{A1}}
\BIC{$\varphi_{615}$}
\end{ptree}
\noindent$\Pi_{41}$:
\begin{ptree}
\AXC{$\mathrm{RFn}(f) \to c \in \mathbb{R} \to \mathrm{PNbhd}(f,c) \to \varphi_{371}$\lbl{A1}}
\AXC{$\stackrel{\Pi_{40}}{\vdots}$}\noLine
\UIC{$\varphi_{615}$}
\BIC{$\mathrm{RFn}(f) \to c \in \mathbb{R} \to \mathrm{PNbhd}(f,c) \to \varphi_{365}$}
\end{ptree}
\noindent$\Pi_{42}$:
\begin{ptree}
\AXC{$\forall f\,\forall x\,(\mathrm{RFn}(f) \to x \in \mathrm{dom}(f) \to f(x) \in \mathbb{R})$\lbl{[S6]}}
\AXC{$\varphi_{102}$\lbl{A2}}
\BIC{$\forall x\,(\mathrm{RFn}(f) \to x \in \mathrm{dom}(f) \to f(x) \in \mathbb{R})$}
\AXC{$\forall x\,(\mathrm{RFn}(f) \to x \in \mathrm{dom}(f) \to f(x) \in \mathbb{R}) \to \varphi_{40}$\lbl{A2}}
\BIC{$\varphi_{40}$}
\end{ptree}
\noindent$\Pi_{43}$:
\begin{ptree}
\AXC{$\mathrm{RFn}(f) \to \mathrm{RFn}(f)$\lbl{A1}}
\AXC{$\stackrel{\Pi_{42}}{\vdots}$}\noLine
\UIC{$\varphi_{40}$}
\AXC{$\varphi_{40} \to (\mathrm{RFn}(f) \to \mathrm{RFn}(f)) \to \varphi_{717}$\lbl{A1}}
\BIC{$(\mathrm{RFn}(f) \to \mathrm{RFn}(f)) \to \varphi_{717}$}
\BIC{$\varphi_{717}$}
\end{ptree}
\noindent$\Pi_{44}$:
\begin{ptree}
\AXC{$\forall f\,\forall x\,(\mathrm{RFn}(f) \to x \in \mathrm{dom}(f) \to f(x) \in \mathbb{R})$\lbl{[S6]}}
\AXC{$\varphi_{102}$\lbl{A2}}
\BIC{$\forall x\,(\mathrm{RFn}(f) \to x \in \mathrm{dom}(f) \to f(x) \in \mathbb{R})$}
\AXC{$\varphi_{82}$\lbl{A2}}
\BIC{$\mathrm{RFn}(f) \to x \in \mathrm{dom}(f) \to f(x) \in \mathbb{R}$}
\AXC{$\varphi_{721}$\lbl{A1}}
\BIC{$(\mathrm{RFn}(f) \to \mathrm{RFn}(f)) \to \varphi_{716}$}
\end{ptree}
\noindent$\Pi_{45}$:
\begin{ptree}
\AXC{$\mathrm{RFn}(f) \to c \in \mathbb{R} \to \mathrm{PNbhd}(f,c) \to \varphi_{568}$\lbl{A1}}
\AXC{$\mathrm{RFn}(f) \to \mathrm{RFn}(f)$\lbl{A1}}
\AXC{$\stackrel{\Pi_{44}}{\vdots}$}\noLine
\UIC{$(\mathrm{RFn}(f) \to \mathrm{RFn}(f)) \to \varphi_{716}$}
\BIC{$\varphi_{716}$}
\BIC{$\mathrm{RFn}(f) \to c \in \mathbb{R} \to \mathrm{PNbhd}(f,c) \to \varphi_{587}$}
\end{ptree}
\noindent$\Pi_{46}$:
\begin{ptree}
\AXC{$\varphi_{142}$\lbl{[C27]}}
\AXC{$\varphi_{142} \to \varphi_{155}$\lbl{A2}}
\BIC{$\varphi_{155}$}
\AXC{$\varphi_{155} \to \varphi_{210}$\lbl{A2}}
\BIC{$\varphi_{210}$}
\AXC{$\varphi_{210} \to \varphi_{211}$\lbl{A2}}
\BIC{$\varphi_{211}$}
\AXC{$\varphi_{403}$\lbl{A2}}
\BIC{$f(x) \in \mathbb{R} \to f(p(N_{0} \cup N_{1})) \in \mathbb{R} \to \alpha \in \mathbb{R} \to \varphi_{208}$}
\AXC{$\varphi_{778}$\lbl{A1}}
\BIC{$(\mathrm{RFn}(f) \to c \in \mathbb{R} \to \mathrm{PNbhd}(f,c) \to \varphi_{587}) \to \varphi_{772}$}
\end{ptree}
\noindent$\Pi_{47}$:
\begin{ptree}
\AXC{$\stackrel{\Pi_{45}}{\vdots}$}\noLine
\UIC{$\mathrm{RFn}(f) \to c \in \mathbb{R} \to \mathrm{PNbhd}(f,c) \to \varphi_{587}$}
\AXC{$\stackrel{\Pi_{46}}{\vdots}$}\noLine
\UIC{$(\mathrm{RFn}(f) \to c \in \mathbb{R} \to \mathrm{PNbhd}(f,c) \to \varphi_{587}) \to \varphi_{772}$}
\BIC{$\varphi_{772}$}
\end{ptree}
\noindent$\Pi_{48}$:
\begin{ptree}
\AXC{$\stackrel{\Pi_{41}}{\vdots}$}\noLine
\UIC{$\mathrm{RFn}(f) \to c \in \mathbb{R} \to \mathrm{PNbhd}(f,c) \to \varphi_{365}$}
\AXC{$\stackrel{\Pi_{31}}{\vdots}$}\noLine
\UIC{$\mathrm{RFn}(f) \to c \in \mathbb{R} \to \mathrm{PNbhd}(f,c) \to \varphi_{569}$}
\AXC{$\stackrel{\Pi_{43}}{\vdots}$}\noLine
\UIC{$\varphi_{717}$}
\BIC{$\mathrm{RFn}(f) \to c \in \mathbb{R} \to \mathrm{PNbhd}(f,c) \to \varphi_{598}$}
\AXC{$\stackrel{\Pi_{47}}{\vdots}$}\noLine
\UIC{$\varphi_{772}$}
\BIC{$(\mathrm{RFn}(f) \to c \in \mathbb{R} \to \mathrm{PNbhd}(f,c) \to \varphi_{365}) \to \varphi_{764}$}
\BIC{$\varphi_{764}$}
\end{ptree}
\noindent$\Pi_{49}$:
\begin{ptree}
\AXC{$\stackrel{\Pi_{21}}{\vdots}$}\noLine
\UIC{$N_{0} \leq N_{0} \cup N_{1}$}
\AXC{$\stackrel{\Pi_{23}}{\vdots}$}\noLine
\UIC{$\varphi_{619}$}
\BIC{$\mathrm{RFn}(f) \to c \in \mathbb{R} \to \mathrm{PNbhd}(f,c) \to \varphi_{593}$}
\AXC{$\varphi_{722}$\lbl{A1}}
\BIC{$\mathrm{RFn}(f) \to c \in \mathbb{R} \to \mathrm{PNbhd}(f,c) \to \varphi_{597}$}
\AXC{$\stackrel{\Pi_{37}}{\vdots}$}\noLine
\UIC{$\varphi_{724}$}
\BIC{$\mathrm{RFn}(f) \to c \in \mathbb{R} \to \mathrm{PNbhd}(f,c) \to \varphi_{611}$}
\AXC{$\stackrel{\Pi_{39}}{\vdots}$}\noLine
\UIC{$\mathrm{RFn}(f) \to c \in \mathbb{R} \to \mathrm{PNbhd}(f,c) \to \varphi_{385}$}
\AXC{$\stackrel{\Pi_{48}}{\vdots}$}\noLine
\UIC{$\varphi_{764}$}
\BIC{$(\mathrm{RFn}(f) \to c \in \mathbb{R} \to \mathrm{PNbhd}(f,c) \to \varphi_{611}) \to \varphi_{720}$}
\BIC{$\varphi_{720}$}
\end{ptree}
\noindent$\Pi_{50}$:
\begin{ptree}
\AXC{$\forall e \in \mathbb{R}\,(e \cdot 2^{-1} + e \cdot 2^{-1} = e)$\lbl{[F36]}}
\AXC{$\varphi_{109}$\lbl{A2}}
\BIC{$\varepsilon \in \mathbb{R} \to \varepsilon \cdot 2^{-1} + \varepsilon \cdot 2^{-1} = \varepsilon$}
\AXC{$(\varepsilon \in \mathbb{R} \to \varepsilon \cdot 2^{-1} + \varepsilon \cdot 2^{-1} = \varepsilon) \to \varphi_{668}$\lbl{A1}}
\BIC{$\varphi_{668}$}
\end{ptree}
\noindent$\Pi_{51}$:
\begin{ptree}
\AXC{$\mathrm{RFn}(f) \to c \in \mathbb{R} \to \mathrm{PNbhd}(f,c) \to \varphi_{373}$\lbl{A1}}
\AXC{$\stackrel{\Pi_{50}}{\vdots}$}\noLine
\UIC{$\varphi_{668}$}
\BIC{$\mathrm{RFn}(f) \to c \in \mathbb{R} \to \mathrm{PNbhd}(f,c) \to \varphi_{600}$}
\AXC{$\varphi_{756}$\lbl{A1}}
\BIC{$(\mathrm{RFn}(f) \to c \in \mathbb{R} \to \mathrm{PNbhd}(f,c) \to \varphi_{607}) \to \varphi_{631}$}
\end{ptree}
\noindent$\Pi_{52}$:
\begin{ptree}
\AXC{$\stackrel{\Pi_{1}}{\vdots}$}\noLine
\UIC{$f \circ p(N_{0} \cup N_{1}) = f(p(N_{0} \cup N_{1})) \to \varphi_{129}$}
\AXC{$\stackrel{\Pi_{20}}{\vdots}$}\noLine
\UIC{$\varphi_{636}$}
\BIC{$\mathrm{RFn}(f) \to c \in \mathbb{R} \to \mathrm{PNbhd}(f,c) \to \varphi_{576}$}
\AXC{$\stackrel{\Pi_{49}}{\vdots}$}\noLine
\UIC{$\varphi_{720}$}
\BIC{$\mathrm{RFn}(f) \to c \in \mathbb{R} \to \mathrm{PNbhd}(f,c) \to \varphi_{607}$}
\AXC{$\stackrel{\Pi_{51}}{\vdots}$}\noLine
\UIC{$(\mathrm{RFn}(f) \to c \in \mathbb{R} \to \mathrm{PNbhd}(f,c) \to \varphi_{607}) \to \varphi_{631}$}
\BIC{$\varphi_{631}$}
\end{ptree}
\noindent$\Pi_{53}$:
\begin{ptree}
\AXC{$z = w \to \lvert f(x) - \alpha \rvert < z \to \lvert f(x) - \alpha \rvert < w$\lbl{A6}}
\UIC{$\forall w\,(z = w \to \lvert f(x) - \alpha \rvert < z \to \lvert f(x) - \alpha \rvert < w)$}
\UIC{$\varphi_{10}$}
\AXC{$\varphi_{10} \to \varphi_{139}$\lbl{A2}}
\BIC{$\varphi_{139}$}
\AXC{$\varphi_{139} \to \varepsilon \cdot 2^{-1} + \varepsilon \cdot 2^{-1} = \varepsilon \to \varphi_{66}$\lbl{A2}}
\BIC{$\varepsilon \cdot 2^{-1} + \varepsilon \cdot 2^{-1} = \varepsilon \to \varphi_{66}$}
\AXC{$\stackrel{\Pi_{52}}{\vdots}$}\noLine
\UIC{$\varphi_{631}$}
\BIC{$\mathrm{RFn}(f) \to c \in \mathbb{R} \to \mathrm{PNbhd}(f,c) \to \varphi_{592}$}
\end{ptree}
\noindent$\Pi_{54}$:
\begin{ptree}
\AXC{$\stackrel{\Pi_{53}}{\vdots}$}\noLine
\UIC{$\mathrm{RFn}(f) \to c \in \mathbb{R} \to \mathrm{PNbhd}(f,c) \to \varphi_{592}$}
\AXC{$(\mathrm{RFn}(f) \to c \in \mathbb{R} \to \mathrm{PNbhd}(f,c) \to \varphi_{592}) \to \varphi_{570} \to \varphi_{59}$\lbl{A1}}
\BIC{$\varphi_{570} \to \varphi_{59}$}
\UIC{$\forall x\,(\varphi_{570} \to \varphi_{59})$}
\end{ptree}
\noindent$\Pi_{55}$:
\begin{ptree}
\AXC{$\forall x\,(\varphi_{570} \to \varphi_{59}) \to \forall x\,(\varphi_{570}) \to \varphi_{69}$\lbl{A3}}
\AXC{$\stackrel{\Pi_{54}}{\vdots}$}\noLine
\UIC{$\forall x\,(\varphi_{570} \to \varphi_{59})$}
\AXC{$\forall x\,(\varphi_{570} \to \varphi_{59}) \to (\forall x\,(\varphi_{570} \to \varphi_{59}) \to \forall x\,(\varphi_{570}) \to \varphi_{69}) \to \varphi_{748}$\lbl{A1}}
\BIC{$(\forall x\,(\varphi_{570} \to \varphi_{59}) \to \forall x\,(\varphi_{570}) \to \varphi_{69}) \to \varphi_{748}$}
\BIC{$\varphi_{748}$}
\end{ptree}
\noindent$\Pi_{56}$:
\begin{ptree}
\AXC{$\mathrm{RFn}(f) \to c \in \mathbb{R} \to \mathrm{PNbhd}(f,c) \to \varphi_{557}$\lbl{A1}}
\AXC{$(\mathrm{RFn}(f) \to c \in \mathbb{R} \to \mathrm{PNbhd}(f,c) \to \varphi_{557}) \to \varphi_{723}$\lbl{A1}}
\BIC{$\varphi_{723}$}
\end{ptree}
\noindent$\Pi_{57}$:
\begin{ptree}
\AXC{$\varphi_{224}$\lbl{A2}}
\AXC{$\varphi_{570} \to \forall x\,(\varphi_{570})$\lbl{A4}}
\AXC{$\stackrel{\Pi_{55}}{\vdots}$}\noLine
\UIC{$\varphi_{748}$}
\BIC{$\mathrm{RFn}(f) \to c \in \mathbb{R} \to \mathrm{PNbhd}(f,c) \to \varphi_{595}$}
\AXC{$\stackrel{\Pi_{56}}{\vdots}$}\noLine
\UIC{$\varphi_{723}$}
\BIC{$\mathrm{RFn}(f) \to c \in \mathbb{R} \to \mathrm{PNbhd}(f,c) \to \varphi_{608}$}
\AXC{$\varphi_{733}$\lbl{A1}}
\BIC{$\varphi_{224} \to \mathrm{RFn}(f) \to c \in \mathbb{R} \to \mathrm{PNbhd}(f,c) \to \varphi_{609}$}
\BIC{$\mathrm{RFn}(f) \to c \in \mathbb{R} \to \mathrm{PNbhd}(f,c) \to \varphi_{609}$}
\AXC{$\varphi_{727}$\lbl{A1}}
\BIC{$\varphi_{582} \to \neg \varphi_{67}$}
\UIC{$\forall N_{1}\,(\varphi_{582} \to \neg \varphi_{67})$}
\end{ptree}
\noindent$\Pi_{58}$:
\begin{ptree}
\AXC{$\mathrm{RFn}(f) \to c \in \mathbb{R} \to \mathrm{PNbhd}(f,c) \to \varphi_{373}$\lbl{A1}}
\AXC{$\forall e \in \mathbb{R}\,(0 < e \to 0 < e \cdot 2^{-1})$\lbl{[F35]}}
\AXC{$\varphi_{83}$\lbl{A2}}
\BIC{$\varepsilon \in \mathbb{R} \to 0 < \varepsilon \to 0 < \varepsilon \cdot 2^{-1}$}
\AXC{$\varphi_{685}$\lbl{A1}}
\BIC{$(\mathrm{RFn}(f) \to c \in \mathbb{R} \to \mathrm{PNbhd}(f,c) \to \varphi_{373}) \to \varphi_{625}$}
\BIC{$\varphi_{625}$}
\end{ptree}
\noindent$\Pi_{59}$:
\begin{ptree}
\AXC{$\mathrm{RFn}(f) \to c \in \mathbb{R} \to \mathrm{PNbhd}(f,c) \to \varphi_{377}$\lbl{A1}}
\AXC{$\stackrel{\Pi_{58}}{\vdots}$}\noLine
\UIC{$\varphi_{625}$}
\BIC{$\mathrm{RFn}(f) \to c \in \mathbb{R} \to \mathrm{PNbhd}(f,c) \to \varphi_{382}$}
\end{ptree}
\noindent$\Pi_{60}$:
\begin{ptree}
\AXC{$\varphi_{131}$\lbl{[C3]}}
\AXC{$\varphi_{131} \to \forall \alpha\,(\varphi_{148})$\lbl{A2}}
\BIC{$\forall \alpha\,(\varphi_{148})$}
\AXC{$\forall \alpha\,(\varphi_{148}) \to \varphi_{148}$\lbl{A2}}
\BIC{$\varphi_{148}$}
\AXC{$\varphi_{277}$\lbl{A2}}
\BIC{$\mathrm{Conv}(f \circ p,\alpha) \to \varepsilon \cdot 2^{-1} \in \mathbb{R} \to 0 < \varepsilon \cdot 2^{-1} \to \varphi_{70}$}
\AXC{$\varphi_{746}$\lbl{A1}}
\BIC{$\varphi_{741}$}
\end{ptree}
\noindent$\Pi_{61}$:
\begin{ptree}
\AXC{$\stackrel{\Pi_{39}}{\vdots}$}\noLine
\UIC{$\mathrm{RFn}(f) \to c \in \mathbb{R} \to \mathrm{PNbhd}(f,c) \to \varphi_{385}$}
\AXC{$\mathrm{RFn}(f) \to c \in \mathbb{R} \to \mathrm{PNbhd}(f,c) \to \varphi_{371}$\lbl{A1}}
\AXC{$\stackrel{\Pi_{60}}{\vdots}$}\noLine
\UIC{$\varphi_{741}$}
\BIC{$(\mathrm{RFn}(f) \to c \in \mathbb{R} \to \mathrm{PNbhd}(f,c) \to \varphi_{385}) \to \varphi_{662}$}
\BIC{$\varphi_{662}$}
\end{ptree}
\noindent$\Pi_{62}$:
\begin{ptree}
\AXC{$\stackrel{\Pi_{57}}{\vdots}$}\noLine
\UIC{$\forall N_{1}\,(\varphi_{582} \to \neg \varphi_{67})$}
\AXC{$\stackrel{\Pi_{59}}{\vdots}$}\noLine
\UIC{$\mathrm{RFn}(f) \to c \in \mathbb{R} \to \mathrm{PNbhd}(f,c) \to \varphi_{382}$}
\AXC{$\stackrel{\Pi_{61}}{\vdots}$}\noLine
\UIC{$\varphi_{662}$}
\BIC{$\mathrm{RFn}(f) \to c \in \mathbb{R} \to \mathrm{PNbhd}(f,c) \to \varphi_{558}$}
\AXC{$\varphi_{780}$\lbl{A1}}
\BIC{$\forall N_{1}\,(\varphi_{582} \to \neg \varphi_{67}) \to \varphi_{774}$}
\BIC{$\varphi_{774}$}
\end{ptree}
\noindent$\Pi_{63}$:
\begin{ptree}
\AXC{$\forall N_{1}\,(\neg \varphi_{67}) \to \varphi_{60}$\lbl{A2}}
\UIC{$\forall N\,(\forall N_{1}\,(\neg \varphi_{67}) \to \varphi_{60})$}
\AXC{$\varphi_{582} \to \forall N_{1}\,(\varphi_{582})$\lbl{A4}}
\AXC{$\forall N_{1}\,(\varphi_{582} \to \neg \varphi_{67}) \to \forall N_{1}\,(\varphi_{582}) \to \forall N_{1}\,(\neg \varphi_{67})$\lbl{A3}}
\AXC{$\stackrel{\Pi_{62}}{\vdots}$}\noLine
\UIC{$\varphi_{774}$}
\BIC{$(\varphi_{582} \to \forall N_{1}\,(\varphi_{582})) \to \forall N\,(\forall N_{1}\,(\neg \varphi_{67}) \to \varphi_{60}) \to \varphi_{669}$}
\BIC{$\forall N\,(\forall N_{1}\,(\neg \varphi_{67}) \to \varphi_{60}) \to \varphi_{669}$}
\BIC{$\varphi_{669}$}
\end{ptree}
\noindent$\Pi_{64}$:
\begin{ptree}
\AXC{$\forall N_{1}\,(\neg \varphi_{67}) \to \forall N\,\forall N_{1}\,(\neg \varphi_{67})$\lbl{A4}}
\AXC{$\forall N\,(\forall N_{1}\,(\neg \varphi_{67}) \to \varphi_{60}) \to \forall N\,\forall N_{1}\,(\neg \varphi_{67}) \to \forall N\,(\varphi_{60})$\lbl{A3}}
\AXC{$\stackrel{\Pi_{63}}{\vdots}$}\noLine
\UIC{$\varphi_{669}$}
\BIC{$\varphi_{635}$}
\BIC{$\mathrm{RFn}(f) \to c \in \mathbb{R} \to \mathrm{PNbhd}(f,c) \to \varphi_{567}$}
\AXC{$\varphi_{708}$\lbl{A1}}
\BIC{$\varphi_{550} \to \neg (N_{0} \in \mathbb{N} \land \varphi_{77})$}
\UIC{$\forall N_{0}\,(\varphi_{550} \to \neg (N_{0} \in \mathbb{N} \land \varphi_{77}))$}
\end{ptree}
\noindent$\Pi_{65}$:
\begin{ptree}
\AXC{$\varphi_{187}$\lbl{[C9]}}
\AXC{$\varphi_{187} \to \varphi_{182}$\lbl{A2}}
\BIC{$\varphi_{182}$}
\AXC{$\varphi_{182} \to \varphi_{174}$\lbl{A2}}
\BIC{$\varphi_{174}$}
\AXC{$\varphi_{174} \to \varphi_{201}$\lbl{A2}}
\BIC{$\varphi_{201}$}
\AXC{$\varphi_{201} \to \varphi_{195}$\lbl{A2}}
\BIC{$\varphi_{195}$}
\AXC{$\varphi_{195} \to \varphi_{183}$\lbl{A2}}
\BIC{$\varphi_{183}$}
\end{ptree}
\noindent$\Pi_{66}$:
\begin{ptree}
\AXC{$\stackrel{\Pi_{65}}{\vdots}$}\noLine
\UIC{$\varphi_{183}$}
\AXC{$\varphi_{183} \to (\mathrm{RFn}(f) \to c \in \mathbb{R} \to \mathrm{PNbhd}(f,c) \to \varphi_{349}) \to \varphi_{747}$\lbl{A1}}
\BIC{$(\mathrm{RFn}(f) \to c \in \mathbb{R} \to \mathrm{PNbhd}(f,c) \to \varphi_{349}) \to \varphi_{747}$}
\end{ptree}
\noindent$\Pi_{67}$:
\begin{ptree}
\AXC{$\mathrm{RFn}(f) \to c \in \mathbb{R} \to \mathrm{PNbhd}(f,c) \to \varphi_{401}$\lbl{A1}}
\AXC{$\varphi_{621}$\lbl{A1}}
\BIC{$\mathrm{RFn}(f) \to c \in \mathbb{R} \to \mathrm{PNbhd}(f,c) \to \varphi_{349}$}
\AXC{$\stackrel{\Pi_{66}}{\vdots}$}\noLine
\UIC{$(\mathrm{RFn}(f) \to c \in \mathbb{R} \to \mathrm{PNbhd}(f,c) \to \varphi_{349}) \to \varphi_{747}$}
\BIC{$\varphi_{747}$}
\end{ptree}
\noindent$\Pi_{68}$:
\begin{ptree}
\AXC{$\mathrm{RFn}(f) \to c \in \mathbb{R} \to \mathrm{PNbhd}(f,c) \to \varphi_{557}$\lbl{A1}}
\AXC{$\varphi_{688}$\lbl{A1}}
\BIC{$\mathrm{RFn}(f) \to c \in \mathbb{R} \to \mathrm{PNbhd}(f,c) \to \varphi_{485}$}
\AXC{$\mathrm{RFn}(f) \to c \in \mathbb{R} \to \mathrm{PNbhd}(f,c) \to \varphi_{371}$\lbl{A1}}
\AXC{$\stackrel{\Pi_{67}}{\vdots}$}\noLine
\UIC{$\varphi_{747}$}
\BIC{$(\mathrm{RFn}(f) \to c \in \mathbb{R} \to \mathrm{PNbhd}(f,c) \to \varphi_{485}) \to \varphi_{671}$}
\BIC{$\varphi_{671}$}
\end{ptree}
\noindent$\Pi_{69}$:
\begin{ptree}
\AXC{$\mathrm{RFn}(f) \to c \in \mathbb{R} \to \mathrm{PNbhd}(f,c) \to \varphi_{557}$\lbl{A1}}
\AXC{$\varphi_{686}$\lbl{A1}}
\BIC{$\mathrm{RFn}(f) \to c \in \mathbb{R} \to \mathrm{PNbhd}(f,c) \to \varphi_{484}$}
\AXC{$\stackrel{\Pi_{68}}{\vdots}$}\noLine
\UIC{$\varphi_{671}$}
\BIC{$\mathrm{RFn}(f) \to c \in \mathbb{R} \to \mathrm{PNbhd}(f,c) \to \varphi_{522}$}
\end{ptree}
\noindent$\Pi_{70}$:
\begin{ptree}
\AXC{$\stackrel{\Pi_{69}}{\vdots}$}\noLine
\UIC{$\mathrm{RFn}(f) \to c \in \mathbb{R} \to \mathrm{PNbhd}(f,c) \to \varphi_{522}$}
\AXC{$(\mathrm{RFn}(f) \to c \in \mathbb{R} \to \mathrm{PNbhd}(f,c) \to \varphi_{522}) \to \varphi_{777}$\lbl{A1}}
\BIC{$\varphi_{777}$}
\end{ptree}
\noindent$\Pi_{71}$:
\begin{ptree}
\AXC{$\forall N_{0}\,(\neg (N_{0} \in \mathbb{N} \land \varphi_{77})) \to \varphi_{110}$\lbl{A2}}
\UIC{$\forall N\,(\forall N_{0}\,(\neg (N_{0} \in \mathbb{N} \land \varphi_{77})) \to \varphi_{110})$}
\AXC{$\varphi_{550} \to \forall N_{0}\,(\varphi_{550})$\lbl{A4}}
\AXC{$\varphi_{709}$\lbl{A3}}
\AXC{$\stackrel{\Pi_{64}}{\vdots}$}\noLine
\UIC{$\forall N_{0}\,(\varphi_{550} \to \neg (N_{0} \in \mathbb{N} \land \varphi_{77}))$}
\AXC{$\stackrel{\Pi_{70}}{\vdots}$}\noLine
\UIC{$\varphi_{777}$}
\BIC{$\varphi_{709} \to (\varphi_{550} \to \forall N_{0}\,(\varphi_{550})) \to \varphi_{697}$}
\BIC{$(\varphi_{550} \to \forall N_{0}\,(\varphi_{550})) \to \varphi_{697}$}
\BIC{$\varphi_{697}$}
\BIC{$\varphi_{360} \to \varphi_{214} \to \mathrm{RFn}(f) \to c \in \mathbb{R} \to \mathrm{PNbhd}(f,c) \to \varphi_{532}$}
\end{ptree}
\noindent$\Pi_{72}$:
\begin{ptree}
\AXC{$\varphi_{214}$\lbl{A4}}
\AXC{$\varphi_{360}$\lbl{A3}}
\AXC{$\stackrel{\Pi_{71}}{\vdots}$}\noLine
\UIC{$\varphi_{360} \to \varphi_{214} \to \mathrm{RFn}(f) \to c \in \mathbb{R} \to \mathrm{PNbhd}(f,c) \to \varphi_{532}$}
\BIC{$\varphi_{214} \to \mathrm{RFn}(f) \to c \in \mathbb{R} \to \mathrm{PNbhd}(f,c) \to \varphi_{532}$}
\BIC{$\mathrm{RFn}(f) \to c \in \mathbb{R} \to \mathrm{PNbhd}(f,c) \to \varphi_{532}$}
\AXC{$\varphi_{692}$\lbl{A1}}
\BIC{$\varphi_{464} \to \neg (\delta \in \mathbb{R} \land 0 < \delta \land \forall x\,(\varphi_{176}))$}
\UIC{$\forall \delta\,(\varphi_{464} \to \neg (\delta \in \mathbb{R} \land 0 < \delta \land \forall x\,(\varphi_{176})))$}
\end{ptree}
\noindent$\Pi_{73}$:
\begin{ptree}
\AXC{$\varphi_{372}$\lbl{A2}}
\AXC{$\mathrm{RFn}(f) \to c \in \mathbb{R} \to \mathrm{PNbhd}(f,c) \to \varphi_{357}$\lbl{A1}}
\AXC{$\varphi_{664}$\lbl{A1}}
\BIC{$\varphi_{372} \to \mathrm{RFn}(f) \to c \in \mathbb{R} \to \mathrm{PNbhd}(f,c) \to \varphi_{372}$}
\BIC{$\mathrm{RFn}(f) \to c \in \mathbb{R} \to \mathrm{PNbhd}(f,c) \to \varphi_{372}$}
\AXC{$\varphi_{732}$\lbl{A1}}
\BIC{$(\mathrm{RFn}(f) \to c \in \mathbb{R} \to \mathrm{PNbhd}(f,c) \to \varphi_{385}) \to \varphi_{648}$}
\end{ptree}
\noindent$\Pi_{74}$:
\begin{ptree}
\AXC{$\stackrel{\Pi_{39}}{\vdots}$}\noLine
\UIC{$\mathrm{RFn}(f) \to c \in \mathbb{R} \to \mathrm{PNbhd}(f,c) \to \varphi_{385}$}
\AXC{$\stackrel{\Pi_{73}}{\vdots}$}\noLine
\UIC{$(\mathrm{RFn}(f) \to c \in \mathbb{R} \to \mathrm{PNbhd}(f,c) \to \varphi_{385}) \to \varphi_{648}$}
\BIC{$\varphi_{648}$}
\end{ptree}
\noindent$\Pi_{75}$:
\begin{ptree}
\AXC{$\stackrel{\Pi_{72}}{\vdots}$}\noLine
\UIC{$\forall \delta\,(\varphi_{464} \to \neg (\delta \in \mathbb{R} \land 0 < \delta \land \forall x\,(\varphi_{176})))$}
\AXC{$\stackrel{\Pi_{59}}{\vdots}$}\noLine
\UIC{$\mathrm{RFn}(f) \to c \in \mathbb{R} \to \mathrm{PNbhd}(f,c) \to \varphi_{382}$}
\AXC{$\stackrel{\Pi_{74}}{\vdots}$}\noLine
\UIC{$\varphi_{648}$}
\BIC{$\mathrm{RFn}(f) \to c \in \mathbb{R} \to \mathrm{PNbhd}(f,c) \to \varphi_{478}$}
\AXC{$\varphi_{773}$\lbl{A1}}
\BIC{$\forall \delta\,(\varphi_{464} \to \neg (\delta \in \mathbb{R} \land 0 < \delta \land \forall x\,(\varphi_{176}))) \to \varphi_{693} \to \varphi_{711}$}
\BIC{$\varphi_{693} \to \varphi_{711}$}
\end{ptree}
\noindent$\Pi_{76}$:
\begin{ptree}
\AXC{$\varphi_{464} \to \forall \delta\,(\varphi_{464})$\lbl{A4}}
\AXC{$\varphi_{693}$\lbl{A3}}
\AXC{$\stackrel{\Pi_{75}}{\vdots}$}\noLine
\UIC{$\varphi_{693} \to \varphi_{711}$}
\BIC{$\varphi_{711}$}
\BIC{$\mathrm{RFn}(f) \to c \in \mathbb{R} \to \mathrm{PNbhd}(f,c) \to \varphi_{457}$}
\AXC{$(\mathrm{RFn}(f) \to c \in \mathbb{R} \to \mathrm{PNbhd}(f,c) \to \varphi_{457}) \to \varphi_{463}$\lbl{A1}}
\BIC{$\varphi_{463}$}
\UIC{$\forall \varepsilon\,(\varphi_{463})$}
\end{ptree}
\noindent$\Pi_{77}$:
\begin{ptree}
\AXC{$\forall f\,\forall c\,\forall \alpha\,(\varphi_{209})$\lbl{DefFLim}}
\AXC{$\forall f\,\forall c\,\forall \alpha\,(\varphi_{209}) \to \forall c\,\forall \alpha\,(\varphi_{209})$\lbl{A2}}
\BIC{$\forall c\,\forall \alpha\,(\varphi_{209})$}
\AXC{$\forall c\,\forall \alpha\,(\varphi_{209}) \to \forall \alpha\,(\varphi_{209})$\lbl{A2}}
\BIC{$\forall \alpha\,(\varphi_{209})$}
\AXC{$\forall \alpha\,(\varphi_{209}) \to \varphi_{209}$\lbl{A2}}
\BIC{$\varphi_{209}$}
\end{ptree}
\noindent$\Pi_{78}$:
\begin{ptree}
\AXC{$\mathrm{RFn}(f) \to c \in \mathbb{R} \to c \in \mathbb{R}$\lbl{A1}}
\AXC{$\mathrm{RFn}(f) \to \mathrm{RFn}(f)$\lbl{A1}}
\AXC{$\stackrel{\Pi_{77}}{\vdots}$}\noLine
\UIC{$\varphi_{209}$}
\AXC{$\varphi_{209} \to (\mathrm{RFn}(f) \to \mathrm{RFn}(f)) \to \varphi_{696}$\lbl{A1}}
\BIC{$(\mathrm{RFn}(f) \to \mathrm{RFn}(f)) \to \varphi_{696}$}
\BIC{$\varphi_{696}$}
\BIC{$(\mathrm{RFn}(f) \to c \in \mathbb{R} \to \mathrm{PNbhd}(f,c) \to \varphi_{365}) \to \varphi_{2} \to \varphi_{634}$}
\end{ptree}
\noindent$\Pi_{79}$:
\begin{ptree}
\AXC{$\varphi_{2}$\lbl{A1}}
\AXC{$\stackrel{\Pi_{41}}{\vdots}$}\noLine
\UIC{$\mathrm{RFn}(f) \to c \in \mathbb{R} \to \mathrm{PNbhd}(f,c) \to \varphi_{365}$}
\AXC{$\stackrel{\Pi_{78}}{\vdots}$}\noLine
\UIC{$(\mathrm{RFn}(f) \to c \in \mathbb{R} \to \mathrm{PNbhd}(f,c) \to \varphi_{365}) \to \varphi_{2} \to \varphi_{634}$}
\BIC{$\varphi_{2} \to \varphi_{634}$}
\BIC{$\varphi_{634}$}
\end{ptree}
\noindent$\Pi_{80}$:
\begin{ptree}
\AXC{$\varphi_{599}$\lbl{A4}}
\AXC{$\forall \varepsilon\,(\varphi_{463}) \to \varphi_{468}$\lbl{A3}}
\AXC{$\stackrel{\Pi_{76}}{\vdots}$}\noLine
\UIC{$\forall \varepsilon\,(\varphi_{463})$}
\AXC{$\forall \varepsilon\,(\varphi_{463}) \to \varphi_{750}$\lbl{A1}}
\BIC{$\varphi_{750}$}
\BIC{$\varphi_{599} \to \mathrm{RFn}(f) \to c \in \mathbb{R} \to \mathrm{PNbhd}(f,c) \to \varphi_{459}$}
\BIC{$\mathrm{RFn}(f) \to c \in \mathbb{R} \to \mathrm{PNbhd}(f,c) \to \varphi_{459}$}
\AXC{$\stackrel{\Pi_{79}}{\vdots}$}\noLine
\UIC{$\varphi_{634}$}
\BIC{$\mathrm{RFn}(f) \to c \in \mathbb{R} \to \mathrm{PNbhd}(f,c) \to \varphi_{366}$}
\end{ptree}
\noindent$\Pi_{81}$:
\begin{ptree}
\AXC{$\stackrel{\Pi_{80}}{\vdots}$}\noLine
\UIC{$\mathrm{RFn}(f) \to c \in \mathbb{R} \to \mathrm{PNbhd}(f,c) \to \varphi_{366}$}
\AXC{$(\mathrm{RFn}(f) \to c \in \mathbb{R} \to \mathrm{PNbhd}(f,c) \to \varphi_{366}) \to \varphi_{430}$\lbl{A1}}
\BIC{$\varphi_{430}$}
\end{ptree}
\noindent$\Pi_{82}$:
\begin{ptree}
\AXC{$\forall \alpha\,(\neg \mathrm{FLim}(f,c,\alpha)) \to \neg \mathrm{FLim}(f,c,\alpha)$\lbl{A2}}
\AXC{$\stackrel{\Pi_{81}}{\vdots}$}\noLine
\UIC{$\varphi_{430}$}
\BIC{$\mathrm{RFn}(f) \to c \in \mathbb{R} \to \mathrm{PNbhd}(f,c) \to \varphi_{368}$}
\AXC{$(\mathrm{RFn}(f) \to c \in \mathbb{R} \to \mathrm{PNbhd}(f,c) \to \varphi_{368}) \to \varphi_{389}$\lbl{A1}}
\BIC{$\varphi_{389}$}
\UIC{$\forall \alpha\,(\varphi_{389})$}
\end{ptree}
\noindent$\Pi_{83}$:
\begin{ptree}
\AXC{$z = w \to \lvert f \circ p(m) - z \rvert < \varepsilon \to \lvert f \circ p(m) - w \rvert < \varepsilon$\lbl{A6}}
\UIC{$\varphi_{71}$}
\UIC{$\forall z\,(\varphi_{71})$}
\AXC{$\forall z\,(\varphi_{71}) \to \varphi_{130}$\lbl{A2}}
\BIC{$\varphi_{130}$}
\AXC{$\varphi_{130} \to f(p(n)) = f \circ p(n) \to \varphi_{125}$\lbl{A2}}
\BIC{$f(p(n)) = f \circ p(n) \to \varphi_{125}$}
\end{ptree}
\noindent$\Pi_{84}$:
\begin{ptree}
\AXC{$\varphi_{88}$\lbl{A2}}
\AXC{$\mathrm{RFn}(f) \to c \in \mathbb{R} \to \mathrm{PNbhd}(f,c) \to \varphi_{507}$\lbl{A1}}
\AXC{$\varphi_{687}$\lbl{A1}}
\BIC{$\mathrm{RFn}(f) \to c \in \mathbb{R} \to \mathrm{PNbhd}(f,c) \to \varphi_{527}$}
\AXC{$\varphi_{698}$\lbl{A1}}
\BIC{$\varphi_{88} \to \mathrm{RFn}(f) \to c \in \mathbb{R} \to \mathrm{PNbhd}(f,c) \to \varphi_{524}$}
\BIC{$\mathrm{RFn}(f) \to c \in \mathbb{R} \to \mathrm{PNbhd}(f,c) \to \varphi_{524}$}
\AXC{$\varphi_{757}$\lbl{A1}}
\BIC{$(\mathrm{RFn}(f) \to c \in \mathbb{R} \to \mathrm{PNbhd}(f,c) \to \varphi_{497}) \to \varphi_{680}$}
\end{ptree}
\noindent$\Pi_{85}$:
\begin{ptree}
\AXC{$\mathrm{RFn}(f) \to c \in \mathbb{R} \to \mathrm{PNbhd}(f,c) \to \varphi_{497}$\lbl{A1}}
\AXC{$\stackrel{\Pi_{84}}{\vdots}$}\noLine
\UIC{$(\mathrm{RFn}(f) \to c \in \mathbb{R} \to \mathrm{PNbhd}(f,c) \to \varphi_{497}) \to \varphi_{680}$}
\BIC{$\varphi_{680}$}
\end{ptree}
\noindent$\Pi_{86}$:
\begin{ptree}
\AXC{$\varphi_{146}$\lbl{A2}}
\AXC{$\mathrm{RFn}(f) \to c \in \mathbb{R} \to \mathrm{PNbhd}(f,c) \to \varphi_{401}$\lbl{A1}}
\AXC{$\varphi_{627}$\lbl{A1}}
\BIC{$\mathrm{RFn}(f) \to c \in \mathbb{R} \to \mathrm{PNbhd}(f,c) \to \varphi_{381}$}
\AXC{$\varphi_{641}$\lbl{A1}}
\BIC{$\varphi_{146} \to \mathrm{RFn}(f) \to c \in \mathbb{R} \to \mathrm{PNbhd}(f,c) \to \varphi_{379}$}
\BIC{$\mathrm{RFn}(f) \to c \in \mathbb{R} \to \mathrm{PNbhd}(f,c) \to \varphi_{379}$}
\end{ptree}
\noindent$\Pi_{87}$:
\begin{ptree}
\AXC{$\stackrel{\Pi_{86}}{\vdots}$}\noLine
\UIC{$\mathrm{RFn}(f) \to c \in \mathbb{R} \to \mathrm{PNbhd}(f,c) \to \varphi_{379}$}
\AXC{$(\mathrm{RFn}(f) \to c \in \mathbb{R} \to \mathrm{PNbhd}(f,c) \to \varphi_{379}) \to \varphi_{684}$\lbl{A1}}
\BIC{$\varphi_{684}$}
\end{ptree}
\noindent$\Pi_{88}$:
\begin{ptree}
\AXC{$\mathrm{RFn}(f) \to c \in \mathbb{R} \to \mathrm{PNbhd}(f,c) \to \varphi_{497}$\lbl{A1}}
\AXC{$\stackrel{\Pi_{87}}{\vdots}$}\noLine
\UIC{$\varphi_{684}$}
\BIC{$\mathrm{RFn}(f) \to c \in \mathbb{R} \to \mathrm{PNbhd}(f,c) \to \varphi_{533}$}
\end{ptree}
\noindent$\Pi_{89}$:
\begin{ptree}
\AXC{$\varphi_{89}$\lbl{A2}}
\AXC{$\mathrm{RFn}(f) \to c \in \mathbb{R} \to \mathrm{PNbhd}(f,c) \to \varphi_{507}$\lbl{A1}}
\AXC{$\varphi_{687}$\lbl{A1}}
\BIC{$\mathrm{RFn}(f) \to c \in \mathbb{R} \to \mathrm{PNbhd}(f,c) \to \varphi_{527}$}
\AXC{$\varphi_{699}$\lbl{A1}}
\BIC{$\varphi_{89} \to \mathrm{RFn}(f) \to c \in \mathbb{R} \to \mathrm{PNbhd}(f,c) \to \varphi_{525}$}
\BIC{$\mathrm{RFn}(f) \to c \in \mathbb{R} \to \mathrm{PNbhd}(f,c) \to \varphi_{525}$}
\AXC{$\varphi_{755}$\lbl{A1}}
\BIC{$(\mathrm{RFn}(f) \to c \in \mathbb{R} \to \mathrm{PNbhd}(f,c) \to \varphi_{494}) \to \varphi_{679}$}
\end{ptree}
\noindent$\Pi_{90}$:
\begin{ptree}
\AXC{$\mathrm{RFn}(f) \to c \in \mathbb{R} \to \mathrm{PNbhd}(f,c) \to \varphi_{494}$\lbl{A1}}
\AXC{$\stackrel{\Pi_{89}}{\vdots}$}\noLine
\UIC{$(\mathrm{RFn}(f) \to c \in \mathbb{R} \to \mathrm{PNbhd}(f,c) \to \varphi_{494}) \to \varphi_{679}$}
\BIC{$\varphi_{679}$}
\end{ptree}
\noindent$\Pi_{91}$:
\begin{ptree}
\AXC{$\varphi_{149}$\lbl{A2}}
\AXC{$\mathrm{RFn}(f) \to c \in \mathbb{R} \to \mathrm{PNbhd}(f,c) \to \varphi_{401}$\lbl{A1}}
\AXC{$\varphi_{627}$\lbl{A1}}
\BIC{$\mathrm{RFn}(f) \to c \in \mathbb{R} \to \mathrm{PNbhd}(f,c) \to \varphi_{381}$}
\AXC{$\varphi_{642}$\lbl{A1}}
\BIC{$\varphi_{149} \to \mathrm{RFn}(f) \to c \in \mathbb{R} \to \mathrm{PNbhd}(f,c) \to \varphi_{380}$}
\BIC{$\mathrm{RFn}(f) \to c \in \mathbb{R} \to \mathrm{PNbhd}(f,c) \to \varphi_{380}$}
\end{ptree}
\noindent$\Pi_{92}$:
\begin{ptree}
\AXC{$\stackrel{\Pi_{91}}{\vdots}$}\noLine
\UIC{$\mathrm{RFn}(f) \to c \in \mathbb{R} \to \mathrm{PNbhd}(f,c) \to \varphi_{380}$}
\AXC{$(\mathrm{RFn}(f) \to c \in \mathbb{R} \to \mathrm{PNbhd}(f,c) \to \varphi_{380}) \to \varphi_{682}$\lbl{A1}}
\BIC{$\varphi_{682}$}
\end{ptree}
\noindent$\Pi_{93}$:
\begin{ptree}
\AXC{$\mathrm{RFn}(f) \to c \in \mathbb{R} \to \mathrm{PNbhd}(f,c) \to \varphi_{494}$\lbl{A1}}
\AXC{$\stackrel{\Pi_{92}}{\vdots}$}\noLine
\UIC{$\varphi_{682}$}
\BIC{$\mathrm{RFn}(f) \to c \in \mathbb{R} \to \mathrm{PNbhd}(f,c) \to \varphi_{534}$}
\end{ptree}
\noindent$\Pi_{94}$:
\begin{ptree}
\AXC{$\forall a\,\forall n\,(\mathrm{Seq}(a) \to n \in \mathbb{N} \to a_{n} \in \mathbb{R})$\lbl{[S1]}}
\AXC{$\varphi_{81}$\lbl{A2}}
\BIC{$\forall n\,(\mathrm{Seq}(p) \to n \in \mathbb{N} \to p(n) \in \mathbb{R})$}
\AXC{$\varphi_{64}$\lbl{A2}}
\BIC{$\mathrm{Seq}(p) \to n \in \mathbb{N} \to p(n) \in \mathbb{R}$}
\AXC{$\varphi_{730}$\lbl{A1}}
\BIC{$(\mathrm{RFn}(f) \to c \in \mathbb{R} \to \mathrm{PNbhd}(f,c) \to \varphi_{343}) \to \varphi_{676}$}
\end{ptree}
\noindent$\Pi_{95}$:
\begin{ptree}
\AXC{$\mathrm{RFn}(f) \to c \in \mathbb{R} \to \mathrm{PNbhd}(f,c) \to \varphi_{401}$\lbl{A1}}
\AXC{$\varphi_{618}$\lbl{A1}}
\BIC{$\mathrm{RFn}(f) \to c \in \mathbb{R} \to \mathrm{PNbhd}(f,c) \to \varphi_{343}$}
\AXC{$\stackrel{\Pi_{94}}{\vdots}$}\noLine
\UIC{$(\mathrm{RFn}(f) \to c \in \mathbb{R} \to \mathrm{PNbhd}(f,c) \to \varphi_{343}) \to \varphi_{676}$}
\BIC{$\varphi_{676}$}
\end{ptree}
\noindent$\Pi_{96}$:
\begin{ptree}
\AXC{$\mathrm{RFn}(f) \to c \in \mathbb{R} \to \mathrm{PNbhd}(f,c) \to \varphi_{336}$\lbl{A1}}
\AXC{$\varphi_{652}$\lbl{A1}}
\BIC{$\mathrm{RFn}(f) \to c \in \mathbb{R} \to \mathrm{PNbhd}(f,c) \to \varphi_{544}$}
\AXC{$(\mathrm{RFn}(f) \to c \in \mathbb{R} \to \mathrm{PNbhd}(f,c) \to \varphi_{544}) \to \varphi_{623}$\lbl{A1}}
\BIC{$\varphi_{623}$}
\end{ptree}
\noindent$\Pi_{97}$:
\begin{ptree}
\AXC{$\mathrm{RFn}(f) \to c \in \mathbb{R} \to \mathrm{PNbhd}(f,c) \to \varphi_{497}$\lbl{A1}}
\AXC{$\stackrel{\Pi_{95}}{\vdots}$}\noLine
\UIC{$\varphi_{676}$}
\BIC{$\mathrm{RFn}(f) \to c \in \mathbb{R} \to \mathrm{PNbhd}(f,c) \to \varphi_{502}$}
\AXC{$\varphi_{16} \to p(n) \in \mathbb{R} \to \varphi_{27}$\lbl{A2}}
\AXC{$\stackrel{\Pi_{96}}{\vdots}$}\noLine
\UIC{$\varphi_{623}$}
\BIC{$\mathrm{RFn}(f) \to c \in \mathbb{R} \to \mathrm{PNbhd}(f,c) \to \varphi_{552}$}
\AXC{$\varphi_{759}$\lbl{A1}}
\BIC{$(\mathrm{RFn}(f) \to c \in \mathbb{R} \to \mathrm{PNbhd}(f,c) \to \varphi_{502}) \to \varphi_{689}$}
\BIC{$\varphi_{689}$}
\end{ptree}
\noindent$\Pi_{98}$:
\begin{ptree}
\AXC{$\forall a\,\forall n\,(\mathrm{Seq}(a) \to n \in \mathbb{N} \to a_{n} \in \mathbb{R})$\lbl{[S1]}}
\AXC{$\varphi_{81}$\lbl{A2}}
\BIC{$\forall n\,(\mathrm{Seq}(p) \to n \in \mathbb{N} \to p(n) \in \mathbb{R})$}
\AXC{$\varphi_{65}$\lbl{A2}}
\BIC{$\mathrm{Seq}(p) \to m \in \mathbb{N} \to p(m) \in \mathbb{R}$}
\AXC{$\varphi_{728}$\lbl{A1}}
\BIC{$(\mathrm{RFn}(f) \to c \in \mathbb{R} \to \mathrm{PNbhd}(f,c) \to \varphi_{343}) \to \varphi_{675}$}
\end{ptree}
\noindent$\Pi_{99}$:
\begin{ptree}
\AXC{$\mathrm{RFn}(f) \to c \in \mathbb{R} \to \mathrm{PNbhd}(f,c) \to \varphi_{401}$\lbl{A1}}
\AXC{$\varphi_{618}$\lbl{A1}}
\BIC{$\mathrm{RFn}(f) \to c \in \mathbb{R} \to \mathrm{PNbhd}(f,c) \to \varphi_{343}$}
\AXC{$\stackrel{\Pi_{98}}{\vdots}$}\noLine
\UIC{$(\mathrm{RFn}(f) \to c \in \mathbb{R} \to \mathrm{PNbhd}(f,c) \to \varphi_{343}) \to \varphi_{675}$}
\BIC{$\varphi_{675}$}
\end{ptree}
\noindent$\Pi_{100}$:
\begin{ptree}
\AXC{$\mathrm{RFn}(f) \to c \in \mathbb{R} \to \mathrm{PNbhd}(f,c) \to \varphi_{336}$\lbl{A1}}
\AXC{$\varphi_{652}$\lbl{A1}}
\BIC{$\mathrm{RFn}(f) \to c \in \mathbb{R} \to \mathrm{PNbhd}(f,c) \to \varphi_{544}$}
\AXC{$(\mathrm{RFn}(f) \to c \in \mathbb{R} \to \mathrm{PNbhd}(f,c) \to \varphi_{544}) \to \varphi_{624}$\lbl{A1}}
\BIC{$\varphi_{624}$}
\end{ptree}
\noindent$\Pi_{101}$:
\begin{ptree}
\AXC{$\mathrm{RFn}(f) \to c \in \mathbb{R} \to \mathrm{PNbhd}(f,c) \to \varphi_{494}$\lbl{A1}}
\AXC{$\stackrel{\Pi_{99}}{\vdots}$}\noLine
\UIC{$\varphi_{675}$}
\BIC{$\mathrm{RFn}(f) \to c \in \mathbb{R} \to \mathrm{PNbhd}(f,c) \to \varphi_{503}$}
\AXC{$\varphi_{16} \to p(m) \in \mathbb{R} \to \varphi_{30}$\lbl{A2}}
\AXC{$\stackrel{\Pi_{100}}{\vdots}$}\noLine
\UIC{$\varphi_{624}$}
\BIC{$\mathrm{RFn}(f) \to c \in \mathbb{R} \to \mathrm{PNbhd}(f,c) \to \varphi_{553}$}
\AXC{$\varphi_{760}$\lbl{A1}}
\BIC{$(\mathrm{RFn}(f) \to c \in \mathbb{R} \to \mathrm{PNbhd}(f,c) \to \varphi_{503}) \to \varphi_{690}$}
\BIC{$\varphi_{690}$}
\end{ptree}
\noindent$\Pi_{102}$:
\begin{ptree}
\AXC{$\mathrm{RFn}(f) \to c \in \mathbb{R} \to \mathrm{PNbhd}(f,c) \to \varphi_{546}$\lbl{A1}}
\AXC{$\varphi_{703}$\lbl{A1}}
\BIC{$\mathrm{RFn}(f) \to c \in \mathbb{R} \to \mathrm{PNbhd}(f,c) \to \varphi_{565}$}
\AXC{$(\mathrm{RFn}(f) \to c \in \mathbb{R} \to \mathrm{PNbhd}(f,c) \to \varphi_{565}) \to \varphi_{653}$\lbl{A1}}
\BIC{$\varphi_{653}$}
\end{ptree}
\noindent$\Pi_{103}$:
\begin{ptree}
\AXC{$\forall x\,(\varphi_{161}) \to \varphi_{199}$\lbl{A2}}
\AXC{$\stackrel{\Pi_{102}}{\vdots}$}\noLine
\UIC{$\varphi_{653}$}
\BIC{$\mathrm{RFn}(f) \to c \in \mathbb{R} \to \mathrm{PNbhd}(f,c) \to \varphi_{573}$}
\AXC{$(\mathrm{RFn}(f) \to c \in \mathbb{R} \to \mathrm{PNbhd}(f,c) \to \varphi_{573}) \to \varphi_{657}$\lbl{A1}}
\BIC{$\varphi_{657}$}
\end{ptree}
\noindent$\Pi_{104}$:
\begin{ptree}
\AXC{$\stackrel{\Pi_{93}}{\vdots}$}\noLine
\UIC{$\mathrm{RFn}(f) \to c \in \mathbb{R} \to \mathrm{PNbhd}(f,c) \to \varphi_{534}$}
\AXC{$\stackrel{\Pi_{101}}{\vdots}$}\noLine
\UIC{$\varphi_{690}$}
\BIC{$\mathrm{RFn}(f) \to c \in \mathbb{R} \to \mathrm{PNbhd}(f,c) \to \varphi_{505}$}
\AXC{$\varphi_{199} \to p(m) \in \mathrm{dom}(f) \to \varphi_{202}$\lbl{A2}}
\AXC{$\stackrel{\Pi_{103}}{\vdots}$}\noLine
\UIC{$\varphi_{657}$}
\BIC{$\mathrm{RFn}(f) \to c \in \mathbb{R} \to \mathrm{PNbhd}(f,c) \to \varphi_{588}$}
\AXC{$\varphi_{779}$\lbl{A1}}
\BIC{$(\mathrm{RFn}(f) \to c \in \mathbb{R} \to \mathrm{PNbhd}(f,c) \to \varphi_{505}) \to \varphi_{770}$}
\BIC{$\varphi_{770}$}
\end{ptree}
\noindent$\Pi_{105}$:
\begin{ptree}
\AXC{$\stackrel{\Pi_{93}}{\vdots}$}\noLine
\UIC{$\mathrm{RFn}(f) \to c \in \mathbb{R} \to \mathrm{PNbhd}(f,c) \to \varphi_{534}$}
\AXC{$\stackrel{\Pi_{88}}{\vdots}$}\noLine
\UIC{$\mathrm{RFn}(f) \to c \in \mathbb{R} \to \mathrm{PNbhd}(f,c) \to \varphi_{533}$}
\AXC{$\stackrel{\Pi_{97}}{\vdots}$}\noLine
\UIC{$\varphi_{689}$}
\BIC{$\mathrm{RFn}(f) \to c \in \mathbb{R} \to \mathrm{PNbhd}(f,c) \to \varphi_{504}$}
\AXC{$\stackrel{\Pi_{104}}{\vdots}$}\noLine
\UIC{$\varphi_{770}$}
\BIC{$(\mathrm{RFn}(f) \to c \in \mathbb{R} \to \mathrm{PNbhd}(f,c) \to \varphi_{534}) \to \varphi_{758}$}
\BIC{$\varphi_{758}$}
\end{ptree}
\noindent$\Pi_{106}$:
\begin{ptree}
\AXC{$\stackrel{\Pi_{88}}{\vdots}$}\noLine
\UIC{$\mathrm{RFn}(f) \to c \in \mathbb{R} \to \mathrm{PNbhd}(f,c) \to \varphi_{533}$}
\AXC{$\mathrm{RFn}(f) \to c \in \mathbb{R} \to \mathrm{PNbhd}(f,c) \to \varphi_{498}$\lbl{A1}}
\AXC{$\stackrel{\Pi_{90}}{\vdots}$}\noLine
\UIC{$\varphi_{679}$}
\BIC{$\mathrm{RFn}(f) \to c \in \mathbb{R} \to \mathrm{PNbhd}(f,c) \to \varphi_{512}$}
\AXC{$\stackrel{\Pi_{105}}{\vdots}$}\noLine
\UIC{$\varphi_{758}$}
\BIC{$(\mathrm{RFn}(f) \to c \in \mathbb{R} \to \mathrm{PNbhd}(f,c) \to \varphi_{533}) \to \varphi_{691}$}
\BIC{$\varphi_{691}$}
\end{ptree}
\noindent$\Pi_{107}$:
\begin{ptree}
\AXC{$x = x$\lbl{A5}}
\AXC{$x = y \to x = x \to y = x$\lbl{A6}}
\AXC{$(x = y \to x = x \to y = x) \to x = x \to x = y \to y = x$\lbl{A1}}
\BIC{$x = x \to x = y \to y = x$}
\BIC{$x = y \to y = x$}
\UIC{$\forall y\,(x = y \to y = x)$}
\UIC{$\forall x\,\forall y\,(x = y \to y = x)$}
\end{ptree}
\noindent$\Pi_{108}$:
\begin{ptree}
\AXC{$\varphi_{119}$\lbl{A2}}
\AXC{$\stackrel{\Pi_{17}}{\vdots}$}\noLine
\UIC{$\mathrm{RFn}(f) \to c \in \mathbb{R} \to \mathrm{PNbhd}(f,c) \to \varphi_{384}$}
\AXC{$\varphi_{655}$\lbl{A1}}
\BIC{$\mathrm{RFn}(f) \to c \in \mathbb{R} \to \mathrm{PNbhd}(f,c) \to \varphi_{535}$}
\AXC{$\varphi_{701}$\lbl{A1}}
\BIC{$\varphi_{119} \to \mathrm{RFn}(f) \to c \in \mathbb{R} \to \mathrm{PNbhd}(f,c) \to \varphi_{530}$}
\BIC{$\mathrm{RFn}(f) \to c \in \mathbb{R} \to \mathrm{PNbhd}(f,c) \to \varphi_{530}$}
\end{ptree}
\noindent$\Pi_{109}$:
\begin{ptree}
\AXC{$\stackrel{\Pi_{108}}{\vdots}$}\noLine
\UIC{$\mathrm{RFn}(f) \to c \in \mathbb{R} \to \mathrm{PNbhd}(f,c) \to \varphi_{530}$}
\AXC{$(\mathrm{RFn}(f) \to c \in \mathbb{R} \to \mathrm{PNbhd}(f,c) \to \varphi_{530}) \to \varphi_{678}$\lbl{A1}}
\BIC{$\varphi_{678}$}
\end{ptree}
\noindent$\Pi_{110}$:
\begin{ptree}
\AXC{$\mathrm{RFn}(f) \to c \in \mathbb{R} \to \mathrm{PNbhd}(f,c) \to \varphi_{494}$\lbl{A1}}
\AXC{$\stackrel{\Pi_{109}}{\vdots}$}\noLine
\UIC{$\varphi_{678}$}
\BIC{$\mathrm{RFn}(f) \to c \in \mathbb{R} \to \mathrm{PNbhd}(f,c) \to \varphi_{517}$}
\AXC{$(\mathrm{RFn}(f) \to c \in \mathbb{R} \to \mathrm{PNbhd}(f,c) \to \varphi_{517}) \to \varphi_{563}$\lbl{A1}}
\BIC{$\varphi_{563}$}
\end{ptree}
\noindent$\Pi_{111}$:
\begin{ptree}
\AXC{$\stackrel{\Pi_{107}}{\vdots}$}\noLine
\UIC{$\forall x\,\forall y\,(x = y \to y = x)$}
\AXC{$\varphi_{13}$\lbl{A2}}
\BIC{$\forall y\,(f \circ p(m) = y \to y = f \circ p(m))$}
\AXC{$\varphi_{144}$\lbl{A2}}
\BIC{$f \circ p(m) = f(p(m)) \to f(p(m)) = f \circ p(m)$}
\AXC{$\stackrel{\Pi_{110}}{\vdots}$}\noLine
\UIC{$\varphi_{563}$}
\BIC{$\mathrm{RFn}(f) \to c \in \mathbb{R} \to \mathrm{PNbhd}(f,c) \to \varphi_{518}$}
\AXC{$\varphi_{743}$\lbl{A1}}
\BIC{$(\mathrm{RFn}(f) \to c \in \mathbb{R} \to \mathrm{PNbhd}(f,c) \to \varphi_{529}) \to \varphi_{603}$}
\end{ptree}
\noindent$\Pi_{112}$:
\begin{ptree}
\AXC{$\mathrm{RFn}(f) \to c \in \mathbb{R} \to \mathrm{PNbhd}(f,c) \to \varphi_{499}$\lbl{A1}}
\AXC{$\stackrel{\Pi_{85}}{\vdots}$}\noLine
\UIC{$\varphi_{680}$}
\BIC{$\mathrm{RFn}(f) \to c \in \mathbb{R} \to \mathrm{PNbhd}(f,c) \to \varphi_{511}$}
\AXC{$\stackrel{\Pi_{106}}{\vdots}$}\noLine
\UIC{$\varphi_{691}$}
\BIC{$\mathrm{RFn}(f) \to c \in \mathbb{R} \to \mathrm{PNbhd}(f,c) \to \varphi_{529}$}
\AXC{$\stackrel{\Pi_{111}}{\vdots}$}\noLine
\UIC{$(\mathrm{RFn}(f) \to c \in \mathbb{R} \to \mathrm{PNbhd}(f,c) \to \varphi_{529}) \to \varphi_{603}$}
\BIC{$\varphi_{603}$}
\end{ptree}
\noindent$\Pi_{113}$:
\begin{ptree}
\AXC{$\varphi_{41}$\lbl{A6}}
\UIC{$\forall w\,(\varphi_{41})$}
\UIC{$\forall z\,\forall w\,(\varphi_{41})$}
\AXC{$\forall z\,\forall w\,(\varphi_{41}) \to \varphi_{121}$\lbl{A2}}
\BIC{$\varphi_{121}$}
\AXC{$\varphi_{121} \to f(p(m)) = f \circ p(m) \to \varphi_{115}$\lbl{A2}}
\BIC{$f(p(m)) = f \circ p(m) \to \varphi_{115}$}
\AXC{$\stackrel{\Pi_{112}}{\vdots}$}\noLine
\UIC{$\varphi_{603}$}
\BIC{$\mathrm{RFn}(f) \to c \in \mathbb{R} \to \mathrm{PNbhd}(f,c) \to \varphi_{536}$}
\end{ptree}
\noindent$\Pi_{114}$:
\begin{ptree}
\AXC{$\varphi_{120}$\lbl{A2}}
\AXC{$\stackrel{\Pi_{17}}{\vdots}$}\noLine
\UIC{$\mathrm{RFn}(f) \to c \in \mathbb{R} \to \mathrm{PNbhd}(f,c) \to \varphi_{384}$}
\AXC{$\varphi_{655}$\lbl{A1}}
\BIC{$\mathrm{RFn}(f) \to c \in \mathbb{R} \to \mathrm{PNbhd}(f,c) \to \varphi_{535}$}
\AXC{$\varphi_{702}$\lbl{A1}}
\BIC{$\varphi_{120} \to \mathrm{RFn}(f) \to c \in \mathbb{R} \to \mathrm{PNbhd}(f,c) \to \varphi_{531}$}
\BIC{$\mathrm{RFn}(f) \to c \in \mathbb{R} \to \mathrm{PNbhd}(f,c) \to \varphi_{531}$}
\end{ptree}
\noindent$\Pi_{115}$:
\begin{ptree}
\AXC{$\stackrel{\Pi_{114}}{\vdots}$}\noLine
\UIC{$\mathrm{RFn}(f) \to c \in \mathbb{R} \to \mathrm{PNbhd}(f,c) \to \varphi_{531}$}
\AXC{$(\mathrm{RFn}(f) \to c \in \mathbb{R} \to \mathrm{PNbhd}(f,c) \to \varphi_{531}) \to \varphi_{681}$\lbl{A1}}
\BIC{$\varphi_{681}$}
\end{ptree}
\noindent$\Pi_{116}$:
\begin{ptree}
\AXC{$\mathrm{RFn}(f) \to c \in \mathbb{R} \to \mathrm{PNbhd}(f,c) \to \varphi_{497}$\lbl{A1}}
\AXC{$\stackrel{\Pi_{115}}{\vdots}$}\noLine
\UIC{$\varphi_{681}$}
\BIC{$\mathrm{RFn}(f) \to c \in \mathbb{R} \to \mathrm{PNbhd}(f,c) \to \varphi_{519}$}
\AXC{$(\mathrm{RFn}(f) \to c \in \mathbb{R} \to \mathrm{PNbhd}(f,c) \to \varphi_{519}) \to \varphi_{564}$\lbl{A1}}
\BIC{$\varphi_{564}$}
\end{ptree}
\noindent$\Pi_{117}$:
\begin{ptree}
\AXC{$\stackrel{\Pi_{107}}{\vdots}$}\noLine
\UIC{$\forall x\,\forall y\,(x = y \to y = x)$}
\AXC{$\varphi_{14}$\lbl{A2}}
\BIC{$\forall y\,(f \circ p(n) = y \to y = f \circ p(n))$}
\AXC{$\varphi_{145}$\lbl{A2}}
\BIC{$f \circ p(n) = f(p(n)) \to f(p(n)) = f \circ p(n)$}
\AXC{$\stackrel{\Pi_{116}}{\vdots}$}\noLine
\UIC{$\varphi_{564}$}
\BIC{$\mathrm{RFn}(f) \to c \in \mathbb{R} \to \mathrm{PNbhd}(f,c) \to \varphi_{520}$}
\AXC{$\varphi_{745}$\lbl{A1}}
\BIC{$(\mathrm{RFn}(f) \to c \in \mathbb{R} \to \mathrm{PNbhd}(f,c) \to \varphi_{536}) \to \varphi_{612}$}
\end{ptree}
\noindent$\Pi_{118}$:
\begin{ptree}
\AXC{$\stackrel{\Pi_{113}}{\vdots}$}\noLine
\UIC{$\mathrm{RFn}(f) \to c \in \mathbb{R} \to \mathrm{PNbhd}(f,c) \to \varphi_{536}$}
\AXC{$\stackrel{\Pi_{117}}{\vdots}$}\noLine
\UIC{$(\mathrm{RFn}(f) \to c \in \mathbb{R} \to \mathrm{PNbhd}(f,c) \to \varphi_{536}) \to \varphi_{612}$}
\BIC{$\varphi_{612}$}
\end{ptree}
\noindent$\Pi_{119}$:
\begin{ptree}
\AXC{$\stackrel{\Pi_{83}}{\vdots}$}\noLine
\UIC{$f(p(n)) = f \circ p(n) \to \varphi_{125}$}
\AXC{$\stackrel{\Pi_{118}}{\vdots}$}\noLine
\UIC{$\varphi_{612}$}
\BIC{$\mathrm{RFn}(f) \to c \in \mathbb{R} \to \mathrm{PNbhd}(f,c) \to \varphi_{539}$}
\AXC{$(\mathrm{RFn}(f) \to c \in \mathbb{R} \to \mathrm{PNbhd}(f,c) \to \varphi_{539}) \to \varphi_{496} \to \varphi_{90}$\lbl{A1}}
\BIC{$\varphi_{496} \to \varphi_{90}$}
\UIC{$\forall n\,(\varphi_{496} \to \varphi_{90})$}
\AXC{$\varphi_{768}$\lbl{A1}}
\BIC{$(\forall n\,(\varphi_{496} \to \varphi_{90}) \to \forall n\,(\varphi_{496}) \to \forall n\,(\varphi_{90})) \to \varphi_{734}$}
\end{ptree}
\noindent$\Pi_{120}$:
\begin{ptree}
\AXC{$\varphi_{496} \to \forall n\,(\varphi_{496})$\lbl{A4}}
\AXC{$\forall n\,(\varphi_{496} \to \varphi_{90}) \to \forall n\,(\varphi_{496}) \to \forall n\,(\varphi_{90})$\lbl{A3}}
\AXC{$\stackrel{\Pi_{119}}{\vdots}$}\noLine
\UIC{$(\forall n\,(\varphi_{496} \to \varphi_{90}) \to \forall n\,(\varphi_{496}) \to \forall n\,(\varphi_{90})) \to \varphi_{734}$}
\BIC{$\varphi_{734}$}
\BIC{$\mathrm{RFn}(f) \to c \in \mathbb{R} \to \mathrm{PNbhd}(f,c) \to \varphi_{540}$}
\AXC{$\varphi_{694}$\lbl{A1}}
\BIC{$\varphi_{496} \to \forall n\,(\varphi_{90})$}
\UIC{$\forall m\,(\varphi_{496} \to \forall n\,(\varphi_{90}))$}
\end{ptree}
\noindent$\Pi_{121}$:
\begin{ptree}
\AXC{$\varphi_{496} \to \forall m\,(\varphi_{496})$\lbl{A4}}
\AXC{$\forall m\,(\varphi_{496} \to \forall n\,(\varphi_{90})) \to \forall m\,(\varphi_{496}) \to \forall m\,\forall n\,(\varphi_{90})$\lbl{A3}}
\AXC{$\stackrel{\Pi_{120}}{\vdots}$}\noLine
\UIC{$\forall m\,(\varphi_{496} \to \forall n\,(\varphi_{90}))$}
\AXC{$\varphi_{769}$\lbl{A1}}
\BIC{$(\forall m\,(\varphi_{496} \to \forall n\,(\varphi_{90})) \to \forall m\,(\varphi_{496}) \to \forall m\,\forall n\,(\varphi_{90})) \to \varphi_{736}$}
\BIC{$\varphi_{736}$}
\BIC{$\mathrm{RFn}(f) \to c \in \mathbb{R} \to \mathrm{PNbhd}(f,c) \to \varphi_{542}$}
\end{ptree}
\noindent$\Pi_{122}$:
\begin{ptree}
\AXC{$\mathrm{RFn}(f) \to c \in \mathbb{R} \to \mathrm{PNbhd}(f,c) \to \varphi_{507}$\lbl{A1}}
\AXC{$\varphi_{674}$\lbl{A1}}
\BIC{$\mathrm{RFn}(f) \to c \in \mathbb{R} \to \mathrm{PNbhd}(f,c) \to \varphi_{491}$}
\AXC{$(\mathrm{RFn}(f) \to c \in \mathbb{R} \to \mathrm{PNbhd}(f,c) \to \varphi_{491}) \to \varphi_{695}$\lbl{A1}}
\BIC{$\varphi_{695}$}
\end{ptree}
\noindent$\Pi_{123}$:
\begin{ptree}
\AXC{$\varphi_{225}$\lbl{A2}}
\AXC{$\stackrel{\Pi_{121}}{\vdots}$}\noLine
\UIC{$\mathrm{RFn}(f) \to c \in \mathbb{R} \to \mathrm{PNbhd}(f,c) \to \varphi_{542}$}
\AXC{$\stackrel{\Pi_{122}}{\vdots}$}\noLine
\UIC{$\varphi_{695}$}
\BIC{$\mathrm{RFn}(f) \to c \in \mathbb{R} \to \mathrm{PNbhd}(f,c) \to \varphi_{545}$}
\AXC{$\varphi_{718}$\lbl{A1}}
\BIC{$\varphi_{225} \to \mathrm{RFn}(f) \to c \in \mathbb{R} \to \mathrm{PNbhd}(f,c) \to \varphi_{547}$}
\BIC{$\mathrm{RFn}(f) \to c \in \mathbb{R} \to \mathrm{PNbhd}(f,c) \to \varphi_{547}$}
\end{ptree}
\noindent$\Pi_{124}$:
\begin{ptree}
\AXC{$\stackrel{\Pi_{123}}{\vdots}$}\noLine
\UIC{$\mathrm{RFn}(f) \to c \in \mathbb{R} \to \mathrm{PNbhd}(f,c) \to \varphi_{547}$}
\AXC{$(\mathrm{RFn}(f) \to c \in \mathbb{R} \to \mathrm{PNbhd}(f,c) \to \varphi_{547}) \to \varphi_{556}$\lbl{A1}}
\BIC{$\varphi_{556}$}
\UIC{$\forall N\,(\varphi_{556})$}
\end{ptree}
\noindent$\Pi_{125}$:
\begin{ptree}
\AXC{$\varphi_{131}$\lbl{[C3]}}
\AXC{$\varphi_{131} \to \varphi_{127}$\lbl{A2}}
\BIC{$\varphi_{127}$}
\AXC{$\varphi_{127} \to \varphi_{122}$\lbl{A2}}
\BIC{$\varphi_{122}$}
\AXC{$\varphi_{122} \to \mathrm{Conv}(p,c) \to \varphi_{74}$\lbl{A2}}
\BIC{$\mathrm{Conv}(p,c) \to \varphi_{74}$}
\AXC{$(\mathrm{Conv}(p,c) \to \varphi_{74}) \to \varphi_{742}$\lbl{A1}}
\BIC{$\varphi_{742}$}
\end{ptree}
\noindent$\Pi_{126}$:
\begin{ptree}
\AXC{$\mathrm{RFn}(f) \to c \in \mathbb{R} \to \mathrm{PNbhd}(f,c) \to \varphi_{401}$\lbl{A1}}
\AXC{$\varphi_{621}$\lbl{A1}}
\BIC{$\mathrm{RFn}(f) \to c \in \mathbb{R} \to \mathrm{PNbhd}(f,c) \to \varphi_{349}$}
\AXC{$\stackrel{\Pi_{125}}{\vdots}$}\noLine
\UIC{$\varphi_{742}$}
\BIC{$(\mathrm{RFn}(f) \to c \in \mathbb{R} \to \mathrm{PNbhd}(f,c) \to \varphi_{470}) \to \varphi_{660}$}
\end{ptree}
\noindent$\Pi_{127}$:
\begin{ptree}
\AXC{$\mathrm{RFn}(f) \to c \in \mathbb{R} \to \mathrm{PNbhd}(f,c) \to \varphi_{546}$\lbl{A1}}
\AXC{$\varphi_{673}$\lbl{A1}}
\BIC{$\mathrm{RFn}(f) \to c \in \mathbb{R} \to \mathrm{PNbhd}(f,c) \to \varphi_{470}$}
\AXC{$\stackrel{\Pi_{126}}{\vdots}$}\noLine
\UIC{$(\mathrm{RFn}(f) \to c \in \mathbb{R} \to \mathrm{PNbhd}(f,c) \to \varphi_{470}) \to \varphi_{660}$}
\BIC{$\varphi_{660}$}
\end{ptree}
\noindent$\Pi_{128}$:
\begin{ptree}
\AXC{$\stackrel{\Pi_{124}}{\vdots}$}\noLine
\UIC{$\forall N\,(\varphi_{556})$}
\AXC{$\mathrm{RFn}(f) \to c \in \mathbb{R} \to \mathrm{PNbhd}(f,c) \to \varphi_{546}$\lbl{A1}}
\AXC{$\varphi_{672}$\lbl{A1}}
\BIC{$\mathrm{RFn}(f) \to c \in \mathbb{R} \to \mathrm{PNbhd}(f,c) \to \varphi_{469}$}
\AXC{$\stackrel{\Pi_{127}}{\vdots}$}\noLine
\UIC{$\varphi_{660}$}
\BIC{$\mathrm{RFn}(f) \to c \in \mathbb{R} \to \mathrm{PNbhd}(f,c) \to \varphi_{488}$}
\AXC{$\varphi_{775}$\lbl{A1}}
\BIC{$\forall N\,(\varphi_{556}) \to (\forall N\,(\varphi_{556}) \to \forall N\,(\varphi_{537}) \to \varphi_{7}) \to \varphi_{738}$}
\BIC{$(\forall N\,(\varphi_{556}) \to \forall N\,(\varphi_{537}) \to \varphi_{7}) \to \varphi_{738}$}
\end{ptree}
\noindent$\Pi_{129}$:
\begin{ptree}
\AXC{$\varphi_{537} \to \forall N\,(\varphi_{537})$\lbl{A4}}
\AXC{$\forall N\,(\varphi_{556}) \to \forall N\,(\varphi_{537}) \to \varphi_{7}$\lbl{A3}}
\AXC{$\stackrel{\Pi_{128}}{\vdots}$}\noLine
\UIC{$(\forall N\,(\varphi_{556}) \to \forall N\,(\varphi_{537}) \to \varphi_{7}) \to \varphi_{738}$}
\BIC{$\varphi_{738}$}
\BIC{$\mathrm{RFn}(f) \to c \in \mathbb{R} \to \mathrm{PNbhd}(f,c) \to \varphi_{501}$}
\end{ptree}
\noindent$\Pi_{130}$:
\begin{ptree}
\AXC{$\stackrel{\Pi_{129}}{\vdots}$}\noLine
\UIC{$\mathrm{RFn}(f) \to c \in \mathbb{R} \to \mathrm{PNbhd}(f,c) \to \varphi_{501}$}
\AXC{$(\mathrm{RFn}(f) \to c \in \mathbb{R} \to \mathrm{PNbhd}(f,c) \to \varphi_{501}) \to \varphi_{538}$\lbl{A1}}
\BIC{$\varphi_{538}$}
\UIC{$\forall \delta\,(\varphi_{538})$}
\end{ptree}
\noindent$\Pi_{131}$:
\begin{ptree}
\AXC{$\varphi_{354}$\lbl{A2}}
\AXC{$\mathrm{RFn}(f) \to c \in \mathbb{R} \to \mathrm{PNbhd}(f,c) \to \varphi_{357}$\lbl{A1}}
\AXC{$\varphi_{659}$\lbl{A1}}
\BIC{$\varphi_{354} \to \mathrm{RFn}(f) \to c \in \mathbb{R} \to \mathrm{PNbhd}(f,c) \to \varphi_{354}$}
\BIC{$\mathrm{RFn}(f) \to c \in \mathbb{R} \to \mathrm{PNbhd}(f,c) \to \varphi_{354}$}
\AXC{$\varphi_{731}$\lbl{A1}}
\BIC{$(\mathrm{RFn}(f) \to c \in \mathbb{R} \to \mathrm{PNbhd}(f,c) \to \varphi_{358}) \to \varphi_{638}$}
\end{ptree}
\noindent$\Pi_{132}$:
\begin{ptree}
\AXC{$\mathrm{RFn}(f) \to c \in \mathbb{R} \to \mathrm{PNbhd}(f,c) \to \varphi_{358}$\lbl{A1}}
\AXC{$\stackrel{\Pi_{131}}{\vdots}$}\noLine
\UIC{$(\mathrm{RFn}(f) \to c \in \mathbb{R} \to \mathrm{PNbhd}(f,c) \to \varphi_{358}) \to \varphi_{638}$}
\BIC{$\varphi_{638}$}
\end{ptree}
\noindent$\Pi_{133}$:
\begin{ptree}
\AXC{$\forall \delta\,(\varphi_{538}) \to \varphi_{541}$\lbl{A3}}
\AXC{$\stackrel{\Pi_{130}}{\vdots}$}\noLine
\UIC{$\forall \delta\,(\varphi_{538})$}
\AXC{$\mathrm{RFn}(f) \to c \in \mathbb{R} \to \mathrm{PNbhd}(f,c) \to \varphi_{364}$\lbl{A1}}
\AXC{$\stackrel{\Pi_{132}}{\vdots}$}\noLine
\UIC{$\varphi_{638}$}
\BIC{$\mathrm{RFn}(f) \to c \in \mathbb{R} \to \mathrm{PNbhd}(f,c) \to \varphi_{466}$}
\AXC{$\varphi_{771}$\lbl{A1}}
\BIC{$\forall \delta\,(\varphi_{538}) \to \varphi_{762}$}
\BIC{$\varphi_{762}$}
\BIC{$\varphi_{644} \to \mathrm{RFn}(f) \to c \in \mathbb{R} \to \mathrm{PNbhd}(f,c) \to \varphi_{447}$}
\end{ptree}
\noindent$\Pi_{134}$:
\begin{ptree}
\AXC{$\varphi_{362} \to \forall \varepsilon\,(\varphi_{362})$\lbl{A4}}
\AXC{$\varphi_{646}$\lbl{A3}}
\AXC{$\varphi_{644}$\lbl{A4}}
\AXC{$\stackrel{\Pi_{133}}{\vdots}$}\noLine
\UIC{$\varphi_{644} \to \mathrm{RFn}(f) \to c \in \mathbb{R} \to \mathrm{PNbhd}(f,c) \to \varphi_{447}$}
\BIC{$\mathrm{RFn}(f) \to c \in \mathbb{R} \to \mathrm{PNbhd}(f,c) \to \varphi_{447}$}
\AXC{$\varphi_{643}$\lbl{A1}}
\BIC{$\varphi_{362} \to \varepsilon \in \mathbb{R} \to 0 < \varepsilon \to \exists N \in \mathbb{N}\,\forall m\,\forall n\,(\varphi_{90})$}
\UIC{$\forall \varepsilon\,(\varphi_{362} \to \varepsilon \in \mathbb{R} \to 0 < \varepsilon \to \exists N \in \mathbb{N}\,\forall m\,\forall n\,(\varphi_{90}))$}
\AXC{$\varphi_{761}$\lbl{A1}}
\BIC{$\varphi_{646} \to \varphi_{670}$}
\BIC{$\varphi_{670}$}
\BIC{$\mathrm{RFn}(f) \to c \in \mathbb{R} \to \mathrm{PNbhd}(f,c) \to \varphi_{448}$}
\end{ptree}
\noindent$\Pi_{135}$:
\begin{ptree}
\AXC{$\varphi_{162}$\lbl{DefCauchy}}
\AXC{$\varphi_{162} \to \varphi_{205}$\lbl{A2}}
\BIC{$\varphi_{205}$}
\AXC{$\varphi_{205} \to (\mathrm{RFn}(f) \to c \in \mathbb{R} \to \mathrm{PNbhd}(f,c) \to \varphi_{352}) \to \varphi_{628}$\lbl{A1}}
\BIC{$(\mathrm{RFn}(f) \to c \in \mathbb{R} \to \mathrm{PNbhd}(f,c) \to \varphi_{352}) \to \varphi_{628}$}
\end{ptree}
\noindent$\Pi_{136}$:
\begin{ptree}
\AXC{$\stackrel{\Pi_{17}}{\vdots}$}\noLine
\UIC{$\mathrm{RFn}(f) \to c \in \mathbb{R} \to \mathrm{PNbhd}(f,c) \to \varphi_{384}$}
\AXC{$\varphi_{617}$\lbl{A1}}
\BIC{$\mathrm{RFn}(f) \to c \in \mathbb{R} \to \mathrm{PNbhd}(f,c) \to \varphi_{352}$}
\AXC{$\stackrel{\Pi_{135}}{\vdots}$}\noLine
\UIC{$(\mathrm{RFn}(f) \to c \in \mathbb{R} \to \mathrm{PNbhd}(f,c) \to \varphi_{352}) \to \varphi_{628}$}
\BIC{$\varphi_{628}$}
\end{ptree}
\noindent$\Pi_{137}$:
\begin{ptree}
\AXC{$\forall a\,(\mathrm{Cgt}(a) \leftrightarrow \mathrm{Cauchy}(a))$\lbl{定理6.11}}
\AXC{$\varphi_{11}$\lbl{A2}}
\BIC{$\mathrm{Cgt}(f \circ p) \leftrightarrow \mathrm{Cauchy}(f \circ p)$}
\AXC{$\varphi_{626}$\lbl{A1}}
\BIC{$(\mathrm{RFn}(f) \to c \in \mathbb{R} \to \mathrm{PNbhd}(f,c) \to \varphi_{353}) \to \varphi_{427}$}
\end{ptree}
\noindent$\Pi_{138}$:
\begin{ptree}
\AXC{$\stackrel{\Pi_{134}}{\vdots}$}\noLine
\UIC{$\mathrm{RFn}(f) \to c \in \mathbb{R} \to \mathrm{PNbhd}(f,c) \to \varphi_{448}$}
\AXC{$\stackrel{\Pi_{136}}{\vdots}$}\noLine
\UIC{$\varphi_{628}$}
\BIC{$\mathrm{RFn}(f) \to c \in \mathbb{R} \to \mathrm{PNbhd}(f,c) \to \varphi_{353}$}
\AXC{$\stackrel{\Pi_{137}}{\vdots}$}\noLine
\UIC{$(\mathrm{RFn}(f) \to c \in \mathbb{R} \to \mathrm{PNbhd}(f,c) \to \varphi_{353}) \to \varphi_{427}$}
\BIC{$\varphi_{427}$}
\end{ptree}
\noindent$\Pi_{139}$:
\begin{ptree}
\AXC{$\stackrel{\Pi_{82}}{\vdots}$}\noLine
\UIC{$\forall \alpha\,(\varphi_{389})$}
\AXC{$\forall a\,(\mathrm{Cgt}(a) \leftrightarrow \exists \alpha\,(\mathrm{Conv}(a,\alpha)))$\lbl{DefCgt}}
\AXC{$\varphi_{45}$\lbl{A2}}
\BIC{$\mathrm{Cgt}(f \circ p) \leftrightarrow \exists \alpha\,(\mathrm{Conv}(f \circ p,\alpha))$}
\AXC{$\stackrel{\Pi_{138}}{\vdots}$}\noLine
\UIC{$\varphi_{427}$}
\BIC{$\mathrm{RFn}(f) \to c \in \mathbb{R} \to \mathrm{PNbhd}(f,c) \to \varphi_{356}$}
\AXC{$\varphi_{763}$\lbl{A1}}
\BIC{$\forall \alpha\,(\varphi_{389}) \to \varphi_{740}$}
\BIC{$\varphi_{740}$}
\end{ptree}
\noindent$\Pi_{140}$:
\begin{ptree}
\AXC{$\forall c\,\forall r\,(c \in \mathbb{R} \to r \in \mathbb{R} \to 0 < r \to \varphi_{87})$\lbl{[H1]}}
\AXC{$\varphi_{260}$\lbl{A2}}
\BIC{$\forall r\,(c \in \mathbb{R} \to r \in \mathbb{R} \to 0 < r \to \varphi_{87})$}
\AXC{$\varphi_{257}$\lbl{A2}}
\BIC{$c \in \mathbb{R} \to r \in \mathbb{R} \to 0 < r \to \varphi_{87}$}
\AXC{$(c \in \mathbb{R} \to r \in \mathbb{R} \to 0 < r \to \varphi_{87}) \to \varphi_{647}$\lbl{A1}}
\BIC{$\varphi_{647}$}
\end{ptree}
\noindent$\Pi_{141}$:
\begin{ptree}
\AXC{$\mathrm{RFn}(f) \to c \in \mathbb{R} \to \mathrm{PNbhd}(f,c) \to \varphi_{336}$\lbl{A1}}
\AXC{$\varphi_{551}$\lbl{A1}}
\BIC{$\mathrm{RFn}(f) \to c \in \mathbb{R} \to \mathrm{PNbhd}(f,c) \to \varphi_{301}$}
\AXC{$\mathrm{RFn}(f) \to c \in \mathbb{R} \to c \in \mathbb{R}$\lbl{A1}}
\AXC{$\stackrel{\Pi_{140}}{\vdots}$}\noLine
\UIC{$\varphi_{647}$}
\BIC{$(\mathrm{RFn}(f) \to c \in \mathbb{R} \to \mathrm{PNbhd}(f,c) \to \varphi_{301}) \to \varphi_{549}$}
\BIC{$\varphi_{549}$}
\end{ptree}
\noindent$\Pi_{142}$:
\begin{ptree}
\AXC{$\mathrm{RFn}(f) \to c \in \mathbb{R} \to \mathrm{PNbhd}(f,c) \to \varphi_{336}$\lbl{A1}}
\AXC{$\varphi_{548}$\lbl{A1}}
\BIC{$\mathrm{RFn}(f) \to c \in \mathbb{R} \to \mathrm{PNbhd}(f,c) \to \varphi_{293}$}
\AXC{$\stackrel{\Pi_{141}}{\vdots}$}\noLine
\UIC{$\varphi_{549}$}
\BIC{$\mathrm{RFn}(f) \to c \in \mathbb{R} \to \mathrm{PNbhd}(f,c) \to \varphi_{337}$}
\AXC{$\varphi_{765}$\lbl{A1}}
\BIC{$\forall p\,(\varphi_{320} \to \neg \varphi_{76}) \to \varphi_{751}$}
\end{ptree}
\noindent$\Pi_{143}$:
\begin{ptree}
\AXC{$\varphi_{591}$\lbl{A4}}
\AXC{$\forall \alpha\,(\varphi_{389}) \to \varphi_{400}$\lbl{A3}}
\AXC{$\stackrel{\Pi_{139}}{\vdots}$}\noLine
\UIC{$\varphi_{740}$}
\BIC{$\varphi_{591} \to \mathrm{RFn}(f) \to c \in \mathbb{R} \to \mathrm{PNbhd}(f,c) \to \varphi_{351}$}
\BIC{$\mathrm{RFn}(f) \to c \in \mathbb{R} \to \mathrm{PNbhd}(f,c) \to \varphi_{351}$}
\AXC{$\varphi_{590}$\lbl{A1}}
\BIC{$\varphi_{320} \to \neg \varphi_{76}$}
\UIC{$\forall p\,(\varphi_{320} \to \neg \varphi_{76})$}
\AXC{$\stackrel{\Pi_{142}}{\vdots}$}\noLine
\UIC{$\forall p\,(\varphi_{320} \to \neg \varphi_{76}) \to \varphi_{751}$}
\BIC{$\varphi_{751}$}
\end{ptree}
\noindent$\Pi_{144}$:
\begin{ptree}
\AXC{$\forall p\,(\neg \varphi_{76}) \to \varphi_{78}$\lbl{A2}}
\UIC{$\forall a\,(\forall p\,(\neg \varphi_{76}) \to \varphi_{78})$}
\AXC{$\varphi_{320} \to \forall p\,(\varphi_{320})$\lbl{A4}}
\AXC{$\forall p\,(\varphi_{320} \to \neg \varphi_{76}) \to \forall p\,(\varphi_{320}) \to \forall p\,(\neg \varphi_{76})$\lbl{A3}}
\AXC{$\stackrel{\Pi_{143}}{\vdots}$}\noLine
\UIC{$\varphi_{751}$}
\BIC{$(\varphi_{320} \to \forall p\,(\varphi_{320})) \to \forall a\,(\forall p\,(\neg \varphi_{76}) \to \varphi_{78}) \to \varphi_{604}$}
\BIC{$\forall a\,(\forall p\,(\neg \varphi_{76}) \to \varphi_{78}) \to \varphi_{604}$}
\BIC{$\varphi_{604}$}
\end{ptree}
\noindent$\Pi_{145}$:
\begin{ptree}
\AXC{$\forall p\,(\neg \varphi_{76}) \to \forall a\,\forall p\,(\neg \varphi_{76})$\lbl{A4}}
\AXC{$\forall a\,(\forall p\,(\neg \varphi_{76}) \to \varphi_{78}) \to \forall a\,\forall p\,(\neg \varphi_{76}) \to \forall a\,(\varphi_{78})$\lbl{A3}}
\AXC{$\stackrel{\Pi_{144}}{\vdots}$}\noLine
\UIC{$\varphi_{604}$}
\BIC{$\varphi_{452}$}
\BIC{$\mathrm{RFn}(f) \to c \in \mathbb{R} \to \mathrm{PNbhd}(f,c) \to \varphi_{303}$}
\AXC{$\varphi_{508}$\lbl{A1}}
\BIC{$\varphi_{261} \to \neg (r \in \mathbb{R} \land 0 < r \land \varphi_{16})$}
\UIC{$\forall r\,(\varphi_{261} \to \neg (r \in \mathbb{R} \land 0 < r \land \varphi_{16}))$}
\end{ptree}
\noindent$\Pi_{146}$:
\begin{ptree}
\AXC{$\varphi_{128}$\lbl{DefPNbhd}}
\AXC{$\varphi_{221}$\lbl{A2}}
\BIC{$\forall c\,(\mathrm{PNbhd}(f,c) \leftrightarrow \exists r\,(r \in \mathbb{R} \land 0 < r \land \varphi_{16}))$}
\AXC{$\varphi_{216}$\lbl{A2}}
\BIC{$\mathrm{PNbhd}(f,c) \leftrightarrow \exists r\,(r \in \mathbb{R} \land 0 < r \land \varphi_{16})$}
\AXC{$\varphi_{2}$\lbl{A1}}
\AXC{$\varphi_{2} \to \varphi_{374}$\lbl{A1}}
\BIC{$\varphi_{374}$}
\BIC{$\mathrm{RFn}(f) \to c \in \mathbb{R} \to \mathrm{PNbhd}(f,c) \to \varphi_{285}$}
\AXC{$\varphi_{737}$\lbl{A1}}
\BIC{$\forall r\,(\varphi_{261} \to \neg (r \in \mathbb{R} \land 0 < r \land \varphi_{16})) \to \varphi_{515} \to \varphi_{554}$}
\end{ptree}
\noindent$\Pi_{147}$:
\begin{ptree}
\AXC{$\varphi_{261} \to \forall r\,(\varphi_{261})$\lbl{A4}}
\AXC{$\varphi_{515}$\lbl{A3}}
\AXC{$\stackrel{\Pi_{145}}{\vdots}$}\noLine
\UIC{$\forall r\,(\varphi_{261} \to \neg (r \in \mathbb{R} \land 0 < r \land \varphi_{16}))$}
\AXC{$\stackrel{\Pi_{146}}{\vdots}$}\noLine
\UIC{$\forall r\,(\varphi_{261} \to \neg (r \in \mathbb{R} \land 0 < r \land \varphi_{16})) \to \varphi_{515} \to \varphi_{554}$}
\BIC{$\varphi_{515} \to \varphi_{554}$}
\BIC{$\varphi_{554}$}
\BIC{$\mathrm{RFn}(f) \to c \in \mathbb{R} \to \mathrm{PNbhd}(f,c) \to \varphi_{227}$}
\end{ptree}
\noindent$\Pi_{148}$:
\begin{ptree}
\AXC{$z = w \to z < \varepsilon \cdot 2^{-1} \to w < \varepsilon \cdot 2^{-1}$\lbl{A6}}
\UIC{$\forall w\,(z = w \to z < \varepsilon \cdot 2^{-1} \to w < \varepsilon \cdot 2^{-1})$}
\UIC{$\forall z\,\forall w\,(z = w \to z < \varepsilon \cdot 2^{-1} \to w < \varepsilon \cdot 2^{-1})$}
\AXC{$\forall z\,\forall w\,(z = w \to z < \varepsilon \cdot 2^{-1} \to w < \varepsilon \cdot 2^{-1}) \to \varphi_{72}$\lbl{A2}}
\BIC{$\varphi_{72}$}
\AXC{$\varphi_{72} \to \varphi_{136}$\lbl{A2}}
\BIC{$\varphi_{136}$}
\end{ptree}
\noindent$\Pi_{149}$:
\begin{ptree}
\AXC{$\mathrm{RFn}(f) \to \varphi_{310}$\lbl{A1}}
\AXC{$(\mathrm{RFn}(f) \to \varphi_{310}) \to \mathrm{RFn}(f) \to \varphi_{286}$\lbl{A1}}
\BIC{$\mathrm{RFn}(f) \to \varphi_{286}$}
\AXC{$(\mathrm{RFn}(f) \to \varphi_{286}) \to (\varphi_{99} \to \varphi_{91}) \to \mathrm{RFn}(f) \to \varphi_{283}$\lbl{A1}}
\BIC{$(\varphi_{99} \to \varphi_{91}) \to \mathrm{RFn}(f) \to \varphi_{283}$}
\end{ptree}
\noindent$\Pi_{150}$:
\begin{ptree}
\AXC{$\mathrm{RFn}(f) \to c \in \mathbb{R} \to \mathrm{PNbhd}(f,c) \to \varphi_{223}$\lbl{A1}}
\AXC{$\varphi_{99} \to \varphi_{91}$\lbl{A2}}
\AXC{$\stackrel{\Pi_{149}}{\vdots}$}\noLine
\UIC{$(\varphi_{99} \to \varphi_{91}) \to \mathrm{RFn}(f) \to \varphi_{283}$}
\BIC{$\mathrm{RFn}(f) \to \varphi_{283}$}
\AXC{$(\mathrm{RFn}(f) \to \varphi_{283}) \to \varphi_{656}$\lbl{A1}}
\BIC{$\varphi_{656}$}
\BIC{$(\mathrm{RFn}(f) \to c \in \mathbb{R} \to \mathrm{PNbhd}(f,c) \to \varphi_{278}) \to \varphi_{506}$}
\end{ptree}
\noindent$\Pi_{151}$:
\begin{ptree}
\AXC{$\mathrm{RFn}(f) \to c \in \mathbb{R} \to \mathrm{PNbhd}(f,c) \to \varphi_{278}$\lbl{A1}}
\AXC{$\stackrel{\Pi_{150}}{\vdots}$}\noLine
\UIC{$(\mathrm{RFn}(f) \to c \in \mathbb{R} \to \mathrm{PNbhd}(f,c) \to \varphi_{278}) \to \varphi_{506}$}
\BIC{$\varphi_{506}$}
\end{ptree}
\noindent$\Pi_{152}$:
\begin{ptree}
\AXC{$\forall f\,\forall c\,\forall \alpha\,(\mathrm{FLim}(f,c,\alpha) \to \alpha \in \mathbb{R})$\lbl{[D4]}}
\AXC{$\varphi_{38}$\lbl{A2}}
\BIC{$\forall c\,\forall \alpha\,(\mathrm{FLim}(f,c,\alpha) \to \alpha \in \mathbb{R})$}
\AXC{$\varphi_{15}$\lbl{A2}}
\BIC{$\forall \alpha\,(\mathrm{FLim}(f,c,\alpha) \to \alpha \in \mathbb{R})$}
\AXC{$\varphi_{5}$\lbl{A2}}
\BIC{$\mathrm{FLim}(f,c,\alpha) \to \alpha \in \mathbb{R}$}
\end{ptree}
\noindent$\Pi_{153}$:
\begin{ptree}
\AXC{$\mathrm{RFn}(f) \to c \in \mathbb{R} \to \mathrm{PNbhd}(f,c) \to \varphi_{4}$\lbl{A1}}
\AXC{$\stackrel{\Pi_{152}}{\vdots}$}\noLine
\UIC{$\mathrm{FLim}(f,c,\alpha) \to \alpha \in \mathbb{R}$}
\AXC{$(\mathrm{FLim}(f,c,\alpha) \to \alpha \in \mathbb{R}) \to \varphi_{191}$\lbl{A1}}
\BIC{$\varphi_{191}$}
\BIC{$\mathrm{RFn}(f) \to c \in \mathbb{R} \to \mathrm{PNbhd}(f,c) \to \varphi_{3}$}
\end{ptree}
\noindent$\Pi_{154}$:
\begin{ptree}
\AXC{$\forall f\,\forall x\,(\mathrm{RFn}(f) \to x \in \mathrm{dom}(f) \to f(x) \in \mathbb{R})$\lbl{[S6]}}
\AXC{$\varphi_{102}$\lbl{A2}}
\BIC{$\forall x\,(\mathrm{RFn}(f) \to x \in \mathrm{dom}(f) \to f(x) \in \mathbb{R})$}
\AXC{$\varphi_{84}$\lbl{A2}}
\BIC{$\mathrm{RFn}(f) \to y \in \mathrm{dom}(f) \to f(y) \in \mathbb{R}$}
\AXC{$\varphi_{492}$\lbl{A1}}
\BIC{$(\mathrm{RFn}(f) \to \mathrm{RFn}(f)) \to \varphi_{471}$}
\end{ptree}
\noindent$\Pi_{155}$:
\begin{ptree}
\AXC{$\mathrm{RFn}(f) \to c \in \mathbb{R} \to \mathrm{PNbhd}(f,c) \to \varphi_{223}$\lbl{A1}}
\AXC{$\mathrm{RFn}(f) \to \mathrm{RFn}(f)$\lbl{A1}}
\AXC{$\stackrel{\Pi_{154}}{\vdots}$}\noLine
\UIC{$(\mathrm{RFn}(f) \to \mathrm{RFn}(f)) \to \varphi_{471}$}
\BIC{$\varphi_{471}$}
\BIC{$\mathrm{RFn}(f) \to c \in \mathbb{R} \to \mathrm{PNbhd}(f,c) \to \varphi_{295}$}
\end{ptree}
\noindent$\Pi_{156}$:
\begin{ptree}
\AXC{$\forall x\,\forall y\,(x \in \mathbb{R} \to y \in \mathbb{R} \to \lvert x - y \rvert = \lvert y - x \rvert)$\lbl{[A13]}}
\AXC{$\forall x\,\forall y\,(x \in \mathbb{R} \to y \in \mathbb{R} \to \lvert x - y \rvert = \lvert y - x \rvert) \to \varphi_{6}$\lbl{A2}}
\BIC{$\varphi_{6}$}
\AXC{$\varphi_{6} \to \varphi_{20}$\lbl{A2}}
\BIC{$\varphi_{20}$}
\UIC{$\forall v_{771}\,(\varphi_{20})$}
\UIC{$\forall v_{770}\,\forall v_{771}\,(\varphi_{20})$}
\end{ptree}
\noindent$\Pi_{157}$:
\begin{ptree}
\AXC{$\stackrel{\Pi_{156}}{\vdots}$}\noLine
\UIC{$\forall v_{770}\,\forall v_{771}\,(\varphi_{20})$}
\AXC{$\forall v_{770}\,\forall v_{771}\,(\varphi_{20}) \to \varphi_{75}$\lbl{A2}}
\BIC{$\varphi_{75}$}
\AXC{$\varphi_{185}$\lbl{A2}}
\BIC{$f(y) \in \mathbb{R} \to \alpha \in \mathbb{R} \to \lvert f(y) - \alpha \rvert = \lvert \alpha - f(y) \rvert$}
\AXC{$\varphi_{579}$\lbl{A1}}
\BIC{$(\mathrm{RFn}(f) \to c \in \mathbb{R} \to \mathrm{PNbhd}(f,c) \to \varphi_{295}) \to \varphi_{375}$}
\end{ptree}
\noindent$\Pi_{158}$:
\begin{ptree}
\AXC{$\stackrel{\Pi_{155}}{\vdots}$}\noLine
\UIC{$\mathrm{RFn}(f) \to c \in \mathbb{R} \to \mathrm{PNbhd}(f,c) \to \varphi_{295}$}
\AXC{$\stackrel{\Pi_{157}}{\vdots}$}\noLine
\UIC{$(\mathrm{RFn}(f) \to c \in \mathbb{R} \to \mathrm{PNbhd}(f,c) \to \varphi_{295}) \to \varphi_{375}$}
\BIC{$\varphi_{375}$}
\end{ptree}
\noindent$\Pi_{159}$:
\begin{ptree}
\AXC{$\stackrel{\Pi_{153}}{\vdots}$}\noLine
\UIC{$\mathrm{RFn}(f) \to c \in \mathbb{R} \to \mathrm{PNbhd}(f,c) \to \varphi_{3}$}
\AXC{$\stackrel{\Pi_{158}}{\vdots}$}\noLine
\UIC{$\varphi_{375}$}
\BIC{$\mathrm{RFn}(f) \to \varphi_{322}$}
\AXC{$(\mathrm{RFn}(f) \to \varphi_{322}) \to \varphi_{571}$\lbl{A1}}
\BIC{$\varphi_{571}$}
\end{ptree}
\noindent$\Pi_{160}$:
\begin{ptree}
\AXC{$z = w \to z < \varepsilon \cdot 2^{-1} \to w < \varepsilon \cdot 2^{-1}$\lbl{A6}}
\UIC{$\forall w\,(z = w \to z < \varepsilon \cdot 2^{-1} \to w < \varepsilon \cdot 2^{-1})$}
\UIC{$\forall z\,\forall w\,(z = w \to z < \varepsilon \cdot 2^{-1} \to w < \varepsilon \cdot 2^{-1})$}
\AXC{$\forall z\,\forall w\,(z = w \to z < \varepsilon \cdot 2^{-1} \to w < \varepsilon \cdot 2^{-1}) \to \varphi_{73}$\lbl{A2}}
\BIC{$\varphi_{73}$}
\AXC{$\varphi_{73} \to \varphi_{137}$\lbl{A2}}
\BIC{$\varphi_{137}$}
\end{ptree}
\noindent$\Pi_{161}$:
\begin{ptree}
\AXC{$\mathrm{RFn}(f) \to \varphi_{310}$\lbl{A1}}
\AXC{$(\mathrm{RFn}(f) \to \varphi_{310}) \to \mathrm{RFn}(f) \to \varphi_{286}$\lbl{A1}}
\BIC{$\mathrm{RFn}(f) \to \varphi_{286}$}
\AXC{$(\mathrm{RFn}(f) \to \varphi_{286}) \to (\varphi_{99} \to \varphi_{92}) \to \mathrm{RFn}(f) \to \varphi_{284}$\lbl{A1}}
\BIC{$(\varphi_{99} \to \varphi_{92}) \to \mathrm{RFn}(f) \to \varphi_{284}$}
\end{ptree}
\noindent$\Pi_{162}$:
\begin{ptree}
\AXC{$\mathrm{RFn}(f) \to c \in \mathbb{R} \to \mathrm{PNbhd}(f,c) \to \varphi_{212}$\lbl{A1}}
\AXC{$\varphi_{99} \to \varphi_{92}$\lbl{A2}}
\AXC{$\stackrel{\Pi_{161}}{\vdots}$}\noLine
\UIC{$(\varphi_{99} \to \varphi_{92}) \to \mathrm{RFn}(f) \to \varphi_{284}$}
\BIC{$\mathrm{RFn}(f) \to \varphi_{284}$}
\AXC{$(\mathrm{RFn}(f) \to \varphi_{284}) \to \varphi_{651}$\lbl{A1}}
\BIC{$\varphi_{651}$}
\BIC{$(\mathrm{RFn}(f) \to \varphi_{272}) \to \varphi_{493}$}
\end{ptree}
\noindent$\Pi_{163}$:
\begin{ptree}
\AXC{$\forall f\,\forall x\,(\mathrm{RFn}(f) \to x \in \mathrm{dom}(f) \to f(x) \in \mathbb{R})$\lbl{[S6]}}
\AXC{$\varphi_{102}$\lbl{A2}}
\BIC{$\forall x\,(\mathrm{RFn}(f) \to x \in \mathrm{dom}(f) \to f(x) \in \mathbb{R})$}
\AXC{$\varphi_{82}$\lbl{A2}}
\BIC{$\mathrm{RFn}(f) \to x \in \mathrm{dom}(f) \to f(x) \in \mathbb{R}$}
\AXC{$\varphi_{489}$\lbl{A1}}
\BIC{$(\mathrm{RFn}(f) \to \mathrm{RFn}(f)) \to \varphi_{465}$}
\end{ptree}
\noindent$\Pi_{164}$:
\begin{ptree}
\AXC{$\mathrm{RFn}(f) \to c \in \mathbb{R} \to \mathrm{PNbhd}(f,c) \to \varphi_{212}$\lbl{A1}}
\AXC{$\mathrm{RFn}(f) \to \mathrm{RFn}(f)$\lbl{A1}}
\AXC{$\stackrel{\Pi_{163}}{\vdots}$}\noLine
\UIC{$(\mathrm{RFn}(f) \to \mathrm{RFn}(f)) \to \varphi_{465}$}
\BIC{$\varphi_{465}$}
\BIC{$\mathrm{RFn}(f) \to c \in \mathbb{R} \to \mathrm{PNbhd}(f,c) \to \varphi_{296}$}
\end{ptree}
\noindent$\Pi_{165}$:
\begin{ptree}
\AXC{$\forall x\,\forall y\,(x \in \mathbb{R} \to y \in \mathbb{R} \to \lvert x - y \rvert = \lvert y - x \rvert)$\lbl{[A13]}}
\AXC{$\forall x\,\forall y\,(x \in \mathbb{R} \to y \in \mathbb{R} \to \lvert x - y \rvert = \lvert y - x \rvert) \to \varphi_{42}$\lbl{A2}}
\BIC{$\varphi_{42}$}
\AXC{$\varphi_{169}$\lbl{A2}}
\BIC{$f(x) \in \mathbb{R} \to \alpha \in \mathbb{R} \to \lvert f(x) - \alpha \rvert = \lvert \alpha - f(x) \rvert$}
\end{ptree}
\noindent$\Pi_{166}$:
\begin{ptree}
\AXC{$\stackrel{\Pi_{164}}{\vdots}$}\noLine
\UIC{$\mathrm{RFn}(f) \to c \in \mathbb{R} \to \mathrm{PNbhd}(f,c) \to \varphi_{296}$}
\AXC{$\stackrel{\Pi_{165}}{\vdots}$}\noLine
\UIC{$f(x) \in \mathbb{R} \to \alpha \in \mathbb{R} \to \lvert f(x) - \alpha \rvert = \lvert \alpha - f(x) \rvert$}
\AXC{$\varphi_{580}$\lbl{A1}}
\BIC{$(\mathrm{RFn}(f) \to c \in \mathbb{R} \to \mathrm{PNbhd}(f,c) \to \varphi_{296}) \to \varphi_{376}$}
\BIC{$\varphi_{376}$}
\end{ptree}
\noindent$\Pi_{167}$:
\begin{ptree}
\AXC{$\stackrel{\Pi_{153}}{\vdots}$}\noLine
\UIC{$\mathrm{RFn}(f) \to c \in \mathbb{R} \to \mathrm{PNbhd}(f,c) \to \varphi_{3}$}
\AXC{$\stackrel{\Pi_{166}}{\vdots}$}\noLine
\UIC{$\varphi_{376}$}
\BIC{$\mathrm{RFn}(f) \to \varphi_{323}$}
\AXC{$(\mathrm{RFn}(f) \to \varphi_{323}) \to \varphi_{572}$\lbl{A1}}
\BIC{$\varphi_{572}$}
\end{ptree}
\noindent$\Pi_{168}$:
\begin{ptree}
\AXC{$\stackrel{\Pi_{160}}{\vdots}$}\noLine
\UIC{$\varphi_{137}$}
\AXC{$\mathrm{RFn}(f) \to \varphi_{280}$\lbl{A1}}
\AXC{$\mathrm{RFn}(f) \to \varphi_{272}$\lbl{A1}}
\AXC{$\stackrel{\Pi_{162}}{\vdots}$}\noLine
\UIC{$(\mathrm{RFn}(f) \to \varphi_{272}) \to \varphi_{493}$}
\BIC{$\varphi_{493}$}
\BIC{$\mathrm{RFn}(f) \to c \in \mathbb{R} \to \mathrm{PNbhd}(f,c) \to \varphi_{313}$}
\AXC{$\stackrel{\Pi_{167}}{\vdots}$}\noLine
\UIC{$\varphi_{572}$}
\BIC{$\varphi_{137} \to \mathrm{RFn}(f) \to c \in \mathbb{R} \to \mathrm{PNbhd}(f,c) \to \varphi_{314}$}
\BIC{$\mathrm{RFn}(f) \to c \in \mathbb{R} \to \mathrm{PNbhd}(f,c) \to \varphi_{314}$}
\end{ptree}
\noindent$\Pi_{169}$:
\begin{ptree}
\AXC{$\forall x\,\forall y\,(x \in \mathbb{R} \to y \in \mathbb{R} \to x \cdot y \in \mathbb{R})$\lbl{(R1)}}
\AXC{$\varphi_{111}$\lbl{A2}}
\BIC{$\forall y\,(\varepsilon \in \mathbb{R} \to y \in \mathbb{R} \to \varepsilon \cdot y \in \mathbb{R})$}
\AXC{$\varphi_{123}$\lbl{A2}}
\BIC{$\varepsilon \in \mathbb{R} \to 2^{-1} \in \mathbb{R} \to \varepsilon \cdot 2^{-1} \in \mathbb{R}$}
\AXC{$\varphi_{437}$\lbl{A1}}
\BIC{$(\mathrm{RFn}(f) \to c \in \mathbb{R} \to \mathrm{PNbhd}(f,c) \to \varphi_{8}) \to \varphi_{327}$}
\end{ptree}
\noindent$\Pi_{170}$:
\begin{ptree}
\AXC{$2^{-1} \in \mathbb{R}$\lbl{[F31]}}
\AXC{$\mathrm{RFn}(f) \to c \in \mathbb{R} \to \mathrm{PNbhd}(f,c) \to \varphi_{8}$\lbl{A1}}
\AXC{$\stackrel{\Pi_{169}}{\vdots}$}\noLine
\UIC{$(\mathrm{RFn}(f) \to c \in \mathbb{R} \to \mathrm{PNbhd}(f,c) \to \varphi_{8}) \to \varphi_{327}$}
\BIC{$\varphi_{327}$}
\BIC{$\mathrm{RFn}(f) \to c \in \mathbb{R} \to \mathrm{PNbhd}(f,c) \to \varphi_{302}$}
\end{ptree}
\noindent$\Pi_{171}$:
\begin{ptree}
\AXC{$\varphi_{158}$\lbl{[C10]}}
\AXC{$\varphi_{158} \to \varphi_{170}$\lbl{A2}}
\BIC{$\varphi_{170}$}
\AXC{$\varphi_{170} \to \varphi_{171}$\lbl{A2}}
\BIC{$\varphi_{171}$}
\AXC{$\varphi_{171} \to \varphi_{175}$\lbl{A2}}
\BIC{$\varphi_{175}$}
\AXC{$\varphi_{175} \to v_{772} \in \mathbb{R} \to \varphi_{152}$\lbl{A2}}
\BIC{$v_{772} \in \mathbb{R} \to \varphi_{152}$}
\UIC{$\forall v_{775}\,(v_{772} \in \mathbb{R} \to \varphi_{152})$}
\UIC{$\forall v_{774}\,\forall v_{775}\,(v_{772} \in \mathbb{R} \to \varphi_{152})$}
\UIC{$\forall v_{773}\,\forall v_{774}\,\forall v_{775}\,(v_{772} \in \mathbb{R} \to \varphi_{152})$}
\UIC{$\forall v_{772}\,\forall v_{773}\,\forall v_{774}\,\forall v_{775}\,(v_{772} \in \mathbb{R} \to \varphi_{152})$}
\end{ptree}
\noindent$\Pi_{172}$:
\begin{ptree}
\AXC{$\stackrel{\Pi_{171}}{\vdots}$}\noLine
\UIC{$\forall v_{772}\,\forall v_{773}\,\forall v_{774}\,\forall v_{775}\,(v_{772} \in \mathbb{R} \to \varphi_{152})$}
\AXC{$\forall v_{772}\,\forall v_{773}\,\forall v_{774}\,\forall v_{775}\,(v_{772} \in \mathbb{R} \to \varphi_{152}) \to \varphi_{192}$\lbl{A2}}
\BIC{$\varphi_{192}$}
\AXC{$\varphi_{192} \to \varphi_{198}$\lbl{A2}}
\BIC{$\varphi_{198}$}
\end{ptree}
\noindent$\Pi_{173}$:
\begin{ptree}
\AXC{$\stackrel{\Pi_{153}}{\vdots}$}\noLine
\UIC{$\mathrm{RFn}(f) \to c \in \mathbb{R} \to \mathrm{PNbhd}(f,c) \to \varphi_{3}$}
\AXC{$\stackrel{\Pi_{172}}{\vdots}$}\noLine
\UIC{$\varphi_{198}$}
\AXC{$\varphi_{198} \to \varphi_{203}$\lbl{A2}}
\BIC{$\varphi_{203}$}
\AXC{$\varphi_{203} \to \alpha \in \mathbb{R} \to f(x) \in \mathbb{R} \to \varphi_{204}$\lbl{A2}}
\BIC{$\alpha \in \mathbb{R} \to f(x) \in \mathbb{R} \to \varphi_{204}$}
\AXC{$\varphi_{744}$\lbl{A1}}
\BIC{$(\mathrm{RFn}(f) \to c \in \mathbb{R} \to \mathrm{PNbhd}(f,c) \to \varphi_{3}) \to \varphi_{735}$}
\BIC{$\varphi_{735}$}
\end{ptree}
\noindent$\Pi_{174}$:
\begin{ptree}
\AXC{$\stackrel{\Pi_{155}}{\vdots}$}\noLine
\UIC{$\mathrm{RFn}(f) \to c \in \mathbb{R} \to \mathrm{PNbhd}(f,c) \to \varphi_{295}$}
\AXC{$\stackrel{\Pi_{164}}{\vdots}$}\noLine
\UIC{$\mathrm{RFn}(f) \to c \in \mathbb{R} \to \mathrm{PNbhd}(f,c) \to \varphi_{296}$}
\AXC{$\stackrel{\Pi_{173}}{\vdots}$}\noLine
\UIC{$\varphi_{735}$}
\BIC{$(\mathrm{RFn}(f) \to c \in \mathbb{R} \to \mathrm{PNbhd}(f,c) \to \varphi_{295}) \to \varphi_{677}$}
\BIC{$\varphi_{677}$}
\end{ptree}
\noindent$\Pi_{175}$:
\begin{ptree}
\AXC{$\stackrel{\Pi_{168}}{\vdots}$}\noLine
\UIC{$\mathrm{RFn}(f) \to c \in \mathbb{R} \to \mathrm{PNbhd}(f,c) \to \varphi_{314}$}
\AXC{$\stackrel{\Pi_{170}}{\vdots}$}\noLine
\UIC{$\mathrm{RFn}(f) \to c \in \mathbb{R} \to \mathrm{PNbhd}(f,c) \to \varphi_{302}$}
\AXC{$\stackrel{\Pi_{174}}{\vdots}$}\noLine
\UIC{$\varphi_{677}$}
\BIC{$(\mathrm{RFn}(f) \to c \in \mathbb{R} \to \mathrm{PNbhd}(f,c) \to \varphi_{314}) \to \varphi_{543}$}
\BIC{$\varphi_{543}$}
\end{ptree}
\noindent$\Pi_{176}$:
\begin{ptree}
\AXC{$\stackrel{\Pi_{148}}{\vdots}$}\noLine
\UIC{$\varphi_{136}$}
\AXC{$\mathrm{RFn}(f) \to c \in \mathbb{R} \to \mathrm{PNbhd}(f,c) \to \varphi_{294}$\lbl{A1}}
\AXC{$\stackrel{\Pi_{151}}{\vdots}$}\noLine
\UIC{$\varphi_{506}$}
\BIC{$\mathrm{RFn}(f) \to c \in \mathbb{R} \to \mathrm{PNbhd}(f,c) \to \varphi_{311}$}
\AXC{$\stackrel{\Pi_{159}}{\vdots}$}\noLine
\UIC{$\varphi_{571}$}
\BIC{$\varphi_{136} \to \mathrm{RFn}(f) \to c \in \mathbb{R} \to \mathrm{PNbhd}(f,c) \to \varphi_{312}$}
\BIC{$\mathrm{RFn}(f) \to c \in \mathbb{R} \to \mathrm{PNbhd}(f,c) \to \varphi_{312}$}
\AXC{$\stackrel{\Pi_{175}}{\vdots}$}\noLine
\UIC{$\varphi_{543}$}
\BIC{$\mathrm{RFn}(f) \to \varphi_{331}$}
\end{ptree}
\noindent$\Pi_{177}$:
\begin{ptree}
\AXC{$\forall e \in \mathbb{R}\,(e \cdot 2^{-1} + e \cdot 2^{-1} = e)$\lbl{[F36]}}
\AXC{$\varphi_{109}$\lbl{A2}}
\BIC{$\varepsilon \in \mathbb{R} \to \varepsilon \cdot 2^{-1} + \varepsilon \cdot 2^{-1} = \varepsilon$}
\AXC{$(\varepsilon \in \mathbb{R} \to \varepsilon \cdot 2^{-1} + \varepsilon \cdot 2^{-1} = \varepsilon) \to \varphi_{383}$\lbl{A1}}
\BIC{$\varphi_{383}$}
\end{ptree}
\noindent$\Pi_{178}$:
\begin{ptree}
\AXC{$\stackrel{\Pi_{176}}{\vdots}$}\noLine
\UIC{$\mathrm{RFn}(f) \to \varphi_{331}$}
\AXC{$\mathrm{RFn}(f) \to c \in \mathbb{R} \to \mathrm{PNbhd}(f,c) \to \varphi_{8}$\lbl{A1}}
\AXC{$\stackrel{\Pi_{177}}{\vdots}$}\noLine
\UIC{$\varphi_{383}$}
\BIC{$\mathrm{RFn}(f) \to c \in \mathbb{R} \to \mathrm{PNbhd}(f,c) \to \varphi_{308}$}
\AXC{$\varphi_{654}$\lbl{A1}}
\BIC{$(\mathrm{RFn}(f) \to \varphi_{331}) \to \varphi_{421}$}
\BIC{$\varphi_{421}$}
\end{ptree}
\noindent$\Pi_{179}$:
\begin{ptree}
\AXC{$\varphi_{22}$\lbl{A6}}
\UIC{$\forall w\,(\varphi_{22})$}
\UIC{$\forall z\,\forall w\,(\varphi_{22})$}
\AXC{$\forall z\,\forall w\,(\varphi_{22}) \to \varphi_{150}$\lbl{A2}}
\BIC{$\varphi_{150}$}
\AXC{$\varphi_{150} \to \varepsilon \cdot 2^{-1} + \varepsilon \cdot 2^{-1} = \varepsilon \to \varphi_{104}$\lbl{A2}}
\BIC{$\varepsilon \cdot 2^{-1} + \varepsilon \cdot 2^{-1} = \varepsilon \to \varphi_{104}$}
\AXC{$\stackrel{\Pi_{178}}{\vdots}$}\noLine
\UIC{$\varphi_{421}$}
\BIC{$\mathrm{RFn}(f) \to c \in \mathbb{R} \to \mathrm{PNbhd}(f,c) \to \varphi_{306}$}
\end{ptree}
\noindent$\Pi_{180}$:
\begin{ptree}
\AXC{$\stackrel{\Pi_{179}}{\vdots}$}\noLine
\UIC{$\mathrm{RFn}(f) \to c \in \mathbb{R} \to \mathrm{PNbhd}(f,c) \to \varphi_{306}$}
\AXC{$(\mathrm{RFn}(f) \to c \in \mathbb{R} \to \mathrm{PNbhd}(f,c) \to \varphi_{306}) \to \varphi_{326}$\lbl{A1}}
\BIC{$\varphi_{326}$}
\UIC{$\forall y\,(\varphi_{326})$}
\end{ptree}
\noindent$\Pi_{181}$:
\begin{ptree}
\AXC{$\varphi_{528}$\lbl{A3}}
\AXC{$\varphi_{370}$\lbl{A4}}
\AXC{$\varphi_{523}$\lbl{A3}}
\AXC{$\stackrel{\Pi_{180}}{\vdots}$}\noLine
\UIC{$\forall y\,(\varphi_{326})$}
\AXC{$\forall y\,(\varphi_{326}) \to \varphi_{523} \to \varphi_{370} \to \mathrm{RFn}(f) \to \varphi_{319}$\lbl{A1}}
\BIC{$\varphi_{523} \to \varphi_{370} \to \mathrm{RFn}(f) \to \varphi_{319}$}
\BIC{$\varphi_{370} \to \mathrm{RFn}(f) \to \varphi_{319}$}
\BIC{$\mathrm{RFn}(f) \to \varphi_{319}$}
\AXC{$\varphi_{521}$\lbl{A1}}
\BIC{$\varphi_{19} \land \varepsilon \in \mathbb{R} \land 0 < \varepsilon \land (\delta \in \mathbb{R} \land 0 < \delta \land \varphi_{99}) \to \varphi_{161}$}
\UIC{$\forall x\,(\varphi_{19} \land \varepsilon \in \mathbb{R} \land 0 < \varepsilon \land (\delta \in \mathbb{R} \land 0 < \delta \land \varphi_{99}) \to \varphi_{161})$}
\AXC{$\varphi_{726}$\lbl{A1}}
\BIC{$\varphi_{528} \to \varphi_{369} \to \mathrm{RFn}(f) \to \varphi_{321}$}
\BIC{$\varphi_{369} \to \mathrm{RFn}(f) \to \varphi_{321}$}
\end{ptree}
\noindent$\Pi_{182}$:
\begin{ptree}
\AXC{$\varphi_{318}$\lbl{A2}}
\AXC{$\varphi_{369}$\lbl{A4}}
\AXC{$\stackrel{\Pi_{181}}{\vdots}$}\noLine
\UIC{$\varphi_{369} \to \mathrm{RFn}(f) \to \varphi_{321}$}
\BIC{$\mathrm{RFn}(f) \to \varphi_{321}$}
\AXC{$\mathrm{RFn}(f) \to \varphi_{310}$\lbl{A1}}
\AXC{$\varphi_{639}$\lbl{A1}}
\BIC{$(\mathrm{RFn}(f) \to \varphi_{321}) \to \mathrm{RFn}(f) \to c \in \mathbb{R} \to \varphi_{328}$}
\BIC{$\mathrm{RFn}(f) \to c \in \mathbb{R} \to \varphi_{328}$}
\AXC{$\varphi_{645}$\lbl{A1}}
\BIC{$\varphi_{318} \to \mathrm{RFn}(f) \to c \in \mathbb{R} \to \mathrm{PNbhd}(f,c) \to \varphi_{325}$}
\BIC{$\mathrm{RFn}(f) \to c \in \mathbb{R} \to \mathrm{PNbhd}(f,c) \to \varphi_{325}$}
\end{ptree}
\noindent$\Pi_{183}$:
\begin{ptree}
\AXC{$\mathrm{RFn}(f) \to c \in \mathbb{R} \to \mathrm{PNbhd}(f,c) \to \varphi_{8}$\lbl{A1}}
\AXC{$\forall e \in \mathbb{R}\,(0 < e \to 0 < e \cdot 2^{-1})$\lbl{[F35]}}
\AXC{$\varphi_{83}$\lbl{A2}}
\BIC{$\varepsilon \in \mathbb{R} \to 0 < \varepsilon \to 0 < \varepsilon \cdot 2^{-1}$}
\AXC{$\varphi_{330}$\lbl{A1}}
\BIC{$(\mathrm{RFn}(f) \to c \in \mathbb{R} \to \mathrm{PNbhd}(f,c) \to \varphi_{8}) \to \varphi_{219}$}
\BIC{$\varphi_{219}$}
\end{ptree}
\noindent$\Pi_{184}$:
\begin{ptree}
\AXC{$\forall x\,\forall y\,(x \in \mathbb{R} \to y \in \mathbb{R} \to x \cdot y \in \mathbb{R})$\lbl{(R1)}}
\AXC{$\varphi_{111}$\lbl{A2}}
\BIC{$\forall y\,(\varepsilon \in \mathbb{R} \to y \in \mathbb{R} \to \varepsilon \cdot y \in \mathbb{R})$}
\AXC{$\varphi_{123}$\lbl{A2}}
\BIC{$\varepsilon \in \mathbb{R} \to 2^{-1} \in \mathbb{R} \to \varepsilon \cdot 2^{-1} \in \mathbb{R}$}
\AXC{$\varphi_{279}$\lbl{A1}}
\BIC{$(\mathrm{RFn}(f) \to c \in \mathbb{R} \to \mathrm{PNbhd}(f,c) \to \varphi_{8}) \to \varphi_{151}$}
\end{ptree}
\noindent$\Pi_{185}$:
\begin{ptree}
\AXC{$2^{-1} \in \mathbb{R}$\lbl{[F31]}}
\AXC{$\mathrm{RFn}(f) \to c \in \mathbb{R} \to \mathrm{PNbhd}(f,c) \to \varphi_{8}$\lbl{A1}}
\AXC{$\stackrel{\Pi_{184}}{\vdots}$}\noLine
\UIC{$(\mathrm{RFn}(f) \to c \in \mathbb{R} \to \mathrm{PNbhd}(f,c) \to \varphi_{8}) \to \varphi_{151}$}
\BIC{$\varphi_{151}$}
\BIC{$\mathrm{RFn}(f) \to c \in \mathbb{R} \to \mathrm{PNbhd}(f,c) \to \varphi_{68}$}
\end{ptree}
\noindent$\Pi_{186}$:
\begin{ptree}
\AXC{$\forall f\,\forall c\,\forall \alpha\,(\varphi_{179})$\lbl{[D1]}}
\AXC{$\forall f\,\forall c\,\forall \alpha\,(\varphi_{179}) \to \forall c\,\forall \alpha\,(\varphi_{179})$\lbl{A2}}
\BIC{$\forall c\,\forall \alpha\,(\varphi_{179})$}
\AXC{$\forall c\,\forall \alpha\,(\varphi_{179}) \to \forall \alpha\,(\varphi_{179})$\lbl{A2}}
\BIC{$\forall \alpha\,(\varphi_{179})$}
\AXC{$\forall \alpha\,(\varphi_{179}) \to \varphi_{179}$\lbl{A2}}
\BIC{$\varphi_{179}$}
\end{ptree}
\noindent$\Pi_{187}$:
\begin{ptree}
\AXC{$\stackrel{\Pi_{186}}{\vdots}$}\noLine
\UIC{$\varphi_{179}$}
\AXC{$\varphi_{179} \to \mathrm{FLim}(f,c,\alpha) \to \varphi_{194}$\lbl{A2}}
\BIC{$\mathrm{FLim}(f,c,\alpha) \to \varphi_{194}$}
\AXC{$(\mathrm{FLim}(f,c,\alpha) \to \varphi_{194}) \to \varphi_{451}$\lbl{A1}}
\BIC{$\varphi_{451}$}
\end{ptree}
\noindent$\Pi_{188}$:
\begin{ptree}
\AXC{$\stackrel{\Pi_{185}}{\vdots}$}\noLine
\UIC{$\mathrm{RFn}(f) \to c \in \mathbb{R} \to \mathrm{PNbhd}(f,c) \to \varphi_{68}$}
\AXC{$\mathrm{RFn}(f) \to c \in \mathbb{R} \to \mathrm{PNbhd}(f,c) \to \varphi_{4}$\lbl{A1}}
\AXC{$\stackrel{\Pi_{187}}{\vdots}$}\noLine
\UIC{$\varphi_{451}$}
\BIC{$(\mathrm{RFn}(f) \to c \in \mathbb{R} \to \mathrm{PNbhd}(f,c) \to \varphi_{68}) \to \varphi_{316}$}
\BIC{$\varphi_{316}$}
\end{ptree}
\noindent$\Pi_{189}$:
\begin{ptree}
\AXC{$\mathrm{RFn}(f) \to c \in \mathbb{R} \to \mathrm{PNbhd}(f,c) \to \varphi_{21}$\lbl{A1}}
\AXC{$\stackrel{\Pi_{183}}{\vdots}$}\noLine
\UIC{$\varphi_{219}$}
\BIC{$\mathrm{RFn}(f) \to c \in \mathbb{R} \to \mathrm{PNbhd}(f,c) \to \varphi_{57}$}
\AXC{$\stackrel{\Pi_{188}}{\vdots}$}\noLine
\UIC{$\varphi_{316}$}
\BIC{$\mathrm{RFn}(f) \to \varphi_{213}$}
\AXC{$(\mathrm{RFn}(f) \to \varphi_{213}) \to \forall \delta\,(\varphi_{341}) \to \varphi_{700}$\lbl{A1}}
\BIC{$\forall \delta\,(\varphi_{341}) \to \varphi_{700}$}
\end{ptree}
\noindent$\Pi_{190}$:
\begin{ptree}
\AXC{$\stackrel{\Pi_{182}}{\vdots}$}\noLine
\UIC{$\mathrm{RFn}(f) \to c \in \mathbb{R} \to \mathrm{PNbhd}(f,c) \to \varphi_{325}$}
\AXC{$(\mathrm{RFn}(f) \to c \in \mathbb{R} \to \mathrm{PNbhd}(f,c) \to \varphi_{325}) \to \varphi_{341}$\lbl{A1}}
\BIC{$\varphi_{341}$}
\UIC{$\forall \delta\,(\varphi_{341})$}
\AXC{$\stackrel{\Pi_{189}}{\vdots}$}\noLine
\UIC{$\forall \delta\,(\varphi_{341}) \to \varphi_{700}$}
\BIC{$\varphi_{700}$}
\end{ptree}
\noindent$\Pi_{191}$:
\begin{ptree}
\AXC{$\varphi_{456}$\lbl{A4}}
\AXC{$\forall \delta\,(\varphi_{341}) \to \varphi_{274} \to \forall \delta\,(\neg (\delta \in \mathbb{R} \land 0 < \delta \land \varphi_{99}))$\lbl{A3}}
\AXC{$\stackrel{\Pi_{190}}{\vdots}$}\noLine
\UIC{$\varphi_{700}$}
\BIC{$\varphi_{456} \to \mathrm{RFn}(f) \to \varphi_{264}$}
\BIC{$\mathrm{RFn}(f) \to \varphi_{264}$}
\AXC{$\varphi_{454}$\lbl{A1}}
\BIC{$\varphi_{19} \to \varepsilon \in \mathbb{R} \to 0 < \varepsilon \to \exists \delta\,(\delta \in \mathbb{R} \land 0 < \delta \land \forall x\,(\varphi_{161}))$}
\UIC{$\varphi_{273}$}
\end{ptree}
\noindent$\Pi_{192}$:
\begin{ptree}
\AXC{$\varphi_{19} \to \forall \varepsilon\,(\varphi_{19})$\lbl{A4}}
\AXC{$\varphi_{273} \to \varphi_{276}$\lbl{A3}}
\AXC{$\stackrel{\Pi_{191}}{\vdots}$}\noLine
\UIC{$\varphi_{273}$}
\AXC{$\varphi_{273} \to (\varphi_{273} \to \varphi_{276}) \to (\varphi_{19} \to \forall \varepsilon\,(\varphi_{19})) \to \mathrm{RFn}(f) \to \varphi_{270}$\lbl{A1}}
\BIC{$(\varphi_{273} \to \varphi_{276}) \to (\varphi_{19} \to \forall \varepsilon\,(\varphi_{19})) \to \mathrm{RFn}(f) \to \varphi_{270}$}
\BIC{$(\varphi_{19} \to \forall \varepsilon\,(\varphi_{19})) \to \mathrm{RFn}(f) \to \varphi_{270}$}
\BIC{$\mathrm{RFn}(f) \to \varphi_{270}$}
\end{ptree}
\noindent$\Pi_{193}$:
\begin{ptree}
\AXC{$\stackrel{\Pi_{192}}{\vdots}$}\noLine
\UIC{$\mathrm{RFn}(f) \to \varphi_{270}$}
\AXC{$(\mathrm{RFn}(f) \to \varphi_{270}) \to \varphi_{262} \to \neg \mathrm{FLim}(f,c,\alpha)$\lbl{A1}}
\BIC{$\varphi_{262} \to \neg \mathrm{FLim}(f,c,\alpha)$}
\UIC{$\forall \alpha\,(\varphi_{262} \to \neg \mathrm{FLim}(f,c,\alpha))$}
\end{ptree}
\noindent$\Pi_{194}$:
\begin{ptree}
\AXC{$\varphi_{462}$\lbl{A3}}
\AXC{$\stackrel{\Pi_{193}}{\vdots}$}\noLine
\UIC{$\forall \alpha\,(\varphi_{262} \to \neg \mathrm{FLim}(f,c,\alpha))$}
\AXC{$\mathrm{RFn}(f) \to c \in \mathbb{R} \to \mathrm{PNbhd}(f,c) \to \varphi_{1}$\lbl{A1}}
\AXC{$\varphi_{706}$\lbl{A1}}
\BIC{$\forall \alpha\,(\varphi_{262} \to \neg \mathrm{FLim}(f,c,\alpha)) \to \varphi_{462} \to \varphi_{555}$}
\BIC{$\varphi_{462} \to \varphi_{555}$}
\BIC{$\varphi_{555}$}
\end{ptree}
\noindent$\Pi_{195}$:
\begin{ptree}
\AXC{$\varphi_{262} \to \forall \alpha\,(\varphi_{262})$\lbl{A4}}
\AXC{$\stackrel{\Pi_{194}}{\vdots}$}\noLine
\UIC{$\varphi_{555}$}
\BIC{$\mathrm{RFn}(f) \to c \in \mathbb{R} \to \mathrm{PNbhd}(f,c) \to \varphi_{228}$}
\AXC{$(\mathrm{RFn}(f) \to c \in \mathbb{R} \to \mathrm{PNbhd}(f,c) \to \varphi_{228}) \to \varphi_{439}$\lbl{A1}}
\BIC{$\varphi_{439}$}
\end{ptree}
\noindent 主木:
\begin{ptree}
\AXC{$\stackrel{\Pi_{147}}{\vdots}$}\noLine
\UIC{$\mathrm{RFn}(f) \to c \in \mathbb{R} \to \mathrm{PNbhd}(f,c) \to \varphi_{227}$}
\AXC{$\stackrel{\Pi_{195}}{\vdots}$}\noLine
\UIC{$\varphi_{439}$}
\BIC{$\mathrm{RFn}(f) \to c \in \mathbb{R} \to \mathrm{PNbhd}(f,c) \to \varphi_{229}$}
\UIC{$\forall c\,(\mathrm{RFn}(f) \to c \in \mathbb{R} \to \mathrm{PNbhd}(f,c) \to \varphi_{229})$}
\UIC{$\forall f\,\forall c\,(\mathrm{RFn}(f) \to c \in \mathbb{R} \to \mathrm{PNbhd}(f,c) \to \varphi_{229})$}
\end{ptree}
\begin{abbrevlist}
\item $\varphi_{1} :\equiv \exists \alpha\,(\mathrm{FLim}(f,c,\alpha)) \to \exists \alpha\,(\mathrm{FLim}(f,c,\alpha))$
\item $\varphi_{2} :\equiv \mathrm{RFn}(f) \to c \in \mathbb{R} \to \mathrm{PNbhd}(f,c) \to \mathrm{PNbhd}(f,c)$
\item $\varphi_{3} :\equiv \exists \alpha\,(\mathrm{FLim}(f,c,\alpha)) \to \mathrm{FLim}(f,c,\alpha) \to \alpha \in \mathbb{R}$
\item $\varphi_{4} :\equiv \exists \alpha\,(\mathrm{FLim}(f,c,\alpha)) \to \mathrm{FLim}(f,c,\alpha) \to \mathrm{FLim}(f,c,\alpha)$
\item $\varphi_{5} :\equiv \forall \alpha\,(\mathrm{FLim}(f,c,\alpha) \to \alpha \in \mathbb{R}) \to \mathrm{FLim}(f,c,\alpha) \to \alpha \in \mathbb{R}$
\item $\varphi_{6} :\equiv \forall y\,(v_{770} \in \mathbb{R} \to y \in \mathbb{R} \to \lvert v_{770} - y \rvert = \lvert y - v_{770} \rvert)$
\item $\varphi_{7} :\equiv \forall N\,(\neg (N \in \mathbb{N} \land \forall n \in \mathbb{N}\,(N \leq n \to \lvert p(n) - c \rvert < \delta)))$
\item $\varphi_{8} :\equiv \exists \alpha\,(\mathrm{FLim}(f,c,\alpha)) \to \mathrm{FLim}(f,c,\alpha) \to \varepsilon \in \mathbb{R} \to \varepsilon \in \mathbb{R}$
\item $\varphi_{9} :\equiv \lvert v_{772} - v_{774} \rvert < v_{775} \to \lvert v_{773} - v_{774} \rvert < v_{775} + v_{775}$
\item $\varphi_{10} :\equiv \forall z\,\forall w\,(z = w \to \lvert f(x) - \alpha \rvert < z \to \lvert f(x) - \alpha \rvert < w)$
\item $\varphi_{11} :\equiv \forall a\,(\mathrm{Cgt}(a) \leftrightarrow \mathrm{Cauchy}(a)) \to (\mathrm{Cgt}(f \circ p) \leftrightarrow \mathrm{Cauchy}(f \circ p))$
\item $\varphi_{12} :\equiv x \in \mathrm{dom}(f) \to y \in \mathrm{dom}(f) \to 0 < \lvert x - c \rvert \to 0 < \lvert x - c \rvert$
\item $\varphi_{13} :\equiv \forall x\,\forall y\,(x = y \to y = x) \to \forall y\,(f \circ p(m) = y \to y = f \circ p(m))$
\item $\varphi_{14} :\equiv \forall x\,\forall y\,(x = y \to y = x) \to \forall y\,(f \circ p(n) = y \to y = f \circ p(n))$
\item $\varphi_{15} :\equiv \forall c\,\forall \alpha\,(\mathrm{FLim}(f,c,\alpha) \to \alpha \in \mathbb{R}) \to \forall \alpha\,(\mathrm{FLim}(f,c,\alpha) \to \alpha \in \mathbb{R})$
\item $\varphi_{16} :\equiv \forall x \in \mathbb{R}\,(0 < \lvert x - c \rvert \to \lvert x - c \rvert < r \to x \in \mathrm{dom}(f))$
\item $\varphi_{17} :\equiv n \in \mathbb{N} \to N \leq m \to N \leq n \to n \in \mathbb{N} \to N \leq n \to \lvert p(n) - c \rvert < \delta$
\item $\varphi_{18} :\equiv n \in \mathbb{N} \to N \leq m \to N \leq n \to m \in \mathbb{N} \to N \leq m \to \lvert p(m) - c \rvert < \delta$
\item $\varphi_{19} :\equiv \mathrm{RFn}(f) \land c \in \mathbb{R} \land \mathrm{PNbhd}(f,c) \land \exists \alpha\,(\mathrm{FLim}(f,c,\alpha)) \land \mathrm{FLim}(f,c,\alpha)$
\item $\varphi_{20} :\equiv v_{770} \in \mathbb{R} \to v_{771} \in \mathbb{R} \to \lvert v_{770} - v_{771} \rvert = \lvert v_{771} - v_{770} \rvert$
\item $\varphi_{21} :\equiv \exists \alpha\,(\mathrm{FLim}(f,c,\alpha)) \to \mathrm{FLim}(f,c,\alpha) \to \varepsilon \in \mathbb{R} \to 0 < \varepsilon \to 0 < \varepsilon$
\item $\varphi_{22} :\equiv z = w \to \lvert f(x) - f(y) \rvert < z \to \lvert f(x) - f(y) \rvert < w$
\item $\varphi_{23} :\equiv n \in \mathbb{N} \to N \leq m \to N \leq n \to \forall n \in \mathbb{N}\,(N \leq n \to \lvert p(n) - c \rvert < \delta)$
\item $\varphi_{24} :\equiv N \in \mathbb{N} \land \forall n \in \mathbb{N}\,(N \leq n \to \lvert p(n) - c \rvert < \delta) \to N \in \mathbb{N}$
\item $\varphi_{25} :\equiv \forall w\,(z = w \to \lvert z - \alpha \rvert < \varepsilon \cdot 2^{-1} \to \lvert w - \alpha \rvert < \varepsilon \cdot 2^{-1})$
\item $\varphi_{26} :\equiv n \in \mathbb{N} \to N \leq m \to N \leq n \to \lvert f(p(m)) - f(p(n)) \rvert < \varepsilon$
\item $\varphi_{27} :\equiv 0 < \lvert p(n) - c \rvert \to \lvert p(n) - c \rvert < r \to p(n) \in \mathrm{dom}(f)$
\item $\varphi_{28} :\equiv x \in \mathrm{dom}(f) \to 0 < \lvert x - c \rvert \to \lvert x - c \rvert < \delta \to \lvert x - c \rvert < \delta$
\item $\varphi_{29} :\equiv x \in \mathrm{dom}(f) \to 0 < \lvert x - c \rvert \to \lvert x - c \rvert < \delta \to f(x) \in \mathbb{R}$
\item $\varphi_{30} :\equiv 0 < \lvert p(m) - c \rvert \to \lvert p(m) - c \rvert < r \to p(m) \in \mathrm{dom}(f)$
\item $\varphi_{31} :\equiv n \in \mathbb{N} \to N \leq m \to N \leq n \to m \in \mathbb{N} \to f \circ p(m) = f(p(m))$
\item $\varphi_{32} :\equiv m \in \mathbb{N} \to n \in \mathbb{N} \to N \leq m \to N \leq n \to f \circ p(m) = f(p(m))$
\item $\varphi_{33} :\equiv m \in \mathbb{N} \to n \in \mathbb{N} \to N \leq m \to N \leq n \to f(p(m)) = f \circ p(m)$
\item $\varphi_{34} :\equiv n \in \mathbb{N} \to N \leq m \to N \leq n \to n \in \mathbb{N} \to f \circ p(n) = f(p(n))$
\item $\varphi_{35} :\equiv m \in \mathbb{N} \to n \in \mathbb{N} \to N \leq m \to N \leq n \to f \circ p(n) = f(p(n))$
\item $\varphi_{36} :\equiv m \in \mathbb{N} \to n \in \mathbb{N} \to N \leq m \to N \leq n \to f(p(n)) = f \circ p(n)$
\item $\varphi_{37} :\equiv y \in \mathrm{dom}(f) \to 0 < \lvert x - c \rvert \to \lvert x - c \rvert < \delta \to \lvert x - c \rvert < \delta$
\item $\varphi_{38} :\equiv \forall f\,\forall c\,\forall \alpha\,(\mathrm{FLim}(f,c,\alpha) \to \alpha \in \mathbb{R}) \to \forall c\,\forall \alpha\,(\mathrm{FLim}(f,c,\alpha) \to \alpha \in \mathbb{R})$
\item $\varphi_{39} :\equiv \forall a\,\forall \alpha\,(\mathrm{Conv}(a,\alpha) \to \alpha \in \mathbb{R}) \to \forall \alpha\,(\mathrm{Conv}(f \circ p,\alpha) \to \alpha \in \mathbb{R})$
\item $\varphi_{40} :\equiv \mathrm{RFn}(f) \to p(N_{0} \cup N_{1}) \in \mathrm{dom}(f) \to f(p(N_{0} \cup N_{1})) \in \mathbb{R}$
\item $\varphi_{41} :\equiv z = w \to \lvert z - f(p(n)) \rvert < \varepsilon \to \lvert w - f(p(n)) \rvert < \varepsilon$
\item $\varphi_{42} :\equiv \forall y\,(f(x) \in \mathbb{R} \to y \in \mathbb{R} \to \lvert f(x) - y \rvert = \lvert y - f(x) \rvert)$
\item $\varphi_{43} :\equiv n \in \mathbb{N} \to N \leq m \to N \leq n \to 0 < \lvert p(n) - c \rvert \land \lvert p(n) - c \rvert < r$
\item $\varphi_{44} :\equiv n \in \mathbb{N} \to N \leq m \to N \leq n \to 0 < \lvert p(m) - c \rvert \land \lvert p(m) - c \rvert < r$
\item $\varphi_{45} :\equiv \forall a\,(\mathrm{Cgt}(a) \leftrightarrow \exists \alpha\,(\mathrm{Conv}(a,\alpha))) \to (\mathrm{Cgt}(f \circ p) \leftrightarrow \exists \alpha\,(\mathrm{Conv}(f \circ p,\alpha)))$
\item $\varphi_{46} :\equiv N_{0} \cup N_{1} \in \mathbb{N} \to f \circ p(N_{0} \cup N_{1}) = f(p(N_{0} \cup N_{1}))$
\item $\varphi_{47} :\equiv n \in \mathbb{N} \to N \leq m \to N \leq n \to \forall n \in \mathbb{N}\,(f \circ p(n) = f(p(n)))$
\item $\varphi_{48} :\equiv n \in \mathbb{N} \to N \leq m \to N \leq n \to \lvert f \circ p(m) - f(p(n)) \rvert < \varepsilon$
\item $\varphi_{49} :\equiv 0 < \lvert x - c \rvert \to \lvert x - c \rvert < \delta \to 0 < \lvert y - c \rvert \to 0 < \lvert y - c \rvert$
\item $\varphi_{50} :\equiv \lvert x - c \rvert < \delta \to 0 < \lvert y - c \rvert \to \lvert y - c \rvert < \delta \to \lvert y - c \rvert < \delta$
\item $\varphi_{51} :\equiv \lvert x - c \rvert < \delta \to 0 < \lvert y - c \rvert \to \lvert y - c \rvert < \delta \to f(y) \in \mathbb{R}$
\item $\varphi_{52} :\equiv \lvert x - c \rvert < \delta \to 0 < \lvert y - c \rvert \to \lvert y - c \rvert < \delta \to f(x) \in \mathbb{R}$
\item $\varphi_{53} :\equiv \forall \alpha\,(\mathrm{Conv}(f \circ p,\alpha) \to \alpha \in \mathbb{R}) \to \mathrm{Conv}(f \circ p,\alpha) \to \alpha \in \mathbb{R}$
\item $\varphi_{54} :\equiv \lvert x - c \rvert < \delta \to \lvert f(x) - f(p(N_{0} \cup N_{1})) \rvert < \varepsilon \cdot 2^{-1}$
\item $\varphi_{55} :\equiv 0 < \lvert y - c \rvert \to \lvert y - c \rvert < \delta \to \lvert f(y) - \alpha \rvert = \lvert \alpha - f(y) \rvert$
\item $\varphi_{56} :\equiv 0 < \lvert y - c \rvert \to \lvert y - c \rvert < \delta \to \lvert f(x) - \alpha \rvert = \lvert \alpha - f(x) \rvert$
\item $\varphi_{57} :\equiv \exists \alpha\,(\mathrm{FLim}(f,c,\alpha)) \to \mathrm{FLim}(f,c,\alpha) \to \varepsilon \in \mathbb{R} \to 0 < \varepsilon \to 0 < \varepsilon \cdot 2^{-1}$
\item $\varphi_{58} :\equiv \lvert y - c \rvert < \delta \to \lvert f(x) - f(y) \rvert < \varepsilon \cdot 2^{-1} + \varepsilon \cdot 2^{-1}$
\item $\varphi_{59} :\equiv x \in \mathrm{dom}(f) \to 0 < \lvert x - c \rvert \to \lvert x - c \rvert < \delta \to \lvert f(x) - \alpha \rvert < \varepsilon$
\item $\varphi_{60} :\equiv \neg (N \in \mathbb{N} \land \forall n \in \mathbb{N}\,(N \leq n \to \lvert f \circ p(n) - \alpha \rvert < \varepsilon \cdot 2^{-1}))$
\item $\varphi_{61} :\equiv \lvert p(N_{0} \cup N_{1}) - c \rvert < \delta \land \lvert f \circ p(N_{0} \cup N_{1}) - \alpha \rvert < \delta$
\item $\varphi_{62} :\equiv N \in \mathbb{N} \land \forall n \in \mathbb{N}\,(N \leq n \to \lvert p(n) - c \rvert < \delta) \to m \in \mathbb{N} \to m \in \mathbb{N}$
\item $\varphi_{63} :\equiv \lvert x - c \rvert < \delta \to 0 < \lvert y - c \rvert \to \lvert y - c \rvert < \delta \to \varepsilon \cdot 2^{-1} \in \mathbb{R}$
\item $\varphi_{64} :\equiv \forall n\,(\mathrm{Seq}(p) \to n \in \mathbb{N} \to p(n) \in \mathbb{R}) \to \mathrm{Seq}(p) \to n \in \mathbb{N} \to p(n) \in \mathbb{R}$
\item $\varphi_{65} :\equiv \forall n\,(\mathrm{Seq}(p) \to n \in \mathbb{N} \to p(n) \in \mathbb{R}) \to \mathrm{Seq}(p) \to m \in \mathbb{N} \to p(m) \in \mathbb{R}$
\item $\varphi_{66} :\equiv \lvert f(x) - \alpha \rvert < \varepsilon \cdot 2^{-1} + \varepsilon \cdot 2^{-1} \to \lvert f(x) - \alpha \rvert < \varepsilon$
\item $\varphi_{67} :\equiv N_{1} \in \mathbb{N} \land \forall n \in \mathbb{N}\,(N_{1} \leq n \to \lvert f \circ p(n) - \alpha \rvert < \varepsilon \cdot 2^{-1})$
\item $\varphi_{68} :\equiv \exists \alpha\,(\mathrm{FLim}(f,c,\alpha)) \to \mathrm{FLim}(f,c,\alpha) \to \varepsilon \in \mathbb{R} \to 0 < \varepsilon \to \varepsilon \cdot 2^{-1} \in \mathbb{R}$
\item $\varphi_{69} :\equiv \forall x \in \mathrm{dom}(f)\,(0 < \lvert x - c \rvert \to \lvert x - c \rvert < \delta \to \lvert f(x) - \alpha \rvert < \varepsilon)$
\item $\varphi_{70} :\equiv \exists N \in \mathbb{N}\,\forall n \in \mathbb{N}\,(N \leq n \to \lvert f \circ p(n) - \alpha \rvert < \varepsilon \cdot 2^{-1})$
\item $\varphi_{71} :\equiv \forall w\,(z = w \to \lvert f \circ p(m) - z \rvert < \varepsilon \to \lvert f \circ p(m) - w \rvert < \varepsilon)$
\item $\varphi_{72} :\equiv \forall w\,(\lvert f(y) - \alpha \rvert = w \to \lvert f(y) - \alpha \rvert < \varepsilon \cdot 2^{-1} \to w < \varepsilon \cdot 2^{-1})$
\item $\varphi_{73} :\equiv \forall w\,(\lvert f(x) - \alpha \rvert = w \to \lvert f(x) - \alpha \rvert < \varepsilon \cdot 2^{-1} \to w < \varepsilon \cdot 2^{-1})$
\item $\varphi_{74} :\equiv \delta \in \mathbb{R} \to 0 < \delta \to \exists N \in \mathbb{N}\,\forall n \in \mathbb{N}\,(N \leq n \to \lvert p(n) - c \rvert < \delta)$
\item $\varphi_{75} :\equiv \forall v_{771}\,(f(y) \in \mathbb{R} \to v_{771} \in \mathbb{R} \to \lvert f(y) - v_{771} \rvert = \lvert v_{771} - f(y) \rvert)$
\item $\varphi_{76} :\equiv \mathrm{Seq}(p) \land \mathrm{Conv}(p,c) \land \forall n \in \mathbb{N}\,(0 < \lvert p(n) - c \rvert \land \lvert p(n) - c \rvert < r)$
\item $\varphi_{77} :\equiv \forall n \in \mathbb{N}\,(N_{0} \leq n \to \lvert p(n) - c \rvert < \delta \land \lvert f \circ p(n) - \alpha \rvert < \delta)$
\item $\varphi_{78} :\equiv \neg (\mathrm{Seq}(a) \land \mathrm{Conv}(a,c) \land \forall n \in \mathbb{N}\,(0 < \lvert a_{n} - c \rvert \land \lvert a_{n} - c \rvert < r))$
\item $\varphi_{79} :\equiv x \in \mathrm{dom}(f) \to 0 < \lvert x - c \rvert \to \lvert x - c \rvert < \delta \to \lvert p(N_{0} \cup N_{1}) - c \rvert < \delta$
\item $\varphi_{80} :\equiv x \in \mathrm{dom}(f) \to 0 < \lvert x - c \rvert \to \lvert x - c \rvert < \delta \to f(p(N_{0} \cup N_{1})) \in \mathbb{R}$
\item $\varphi_{81} :\equiv \forall a\,\forall n\,(\mathrm{Seq}(a) \to n \in \mathbb{N} \to a_{n} \in \mathbb{R}) \to \forall n\,(\mathrm{Seq}(p) \to n \in \mathbb{N} \to p(n) \in \mathbb{R})$
\item $\varphi_{82} :\equiv \forall x\,(\mathrm{RFn}(f) \to x \in \mathrm{dom}(f) \to f(x) \in \mathbb{R}) \to \mathrm{RFn}(f) \to x \in \mathrm{dom}(f) \to f(x) \in \mathbb{R}$
\item $\varphi_{83} :\equiv \forall e \in \mathbb{R}\,(0 < e \to 0 < e \cdot 2^{-1}) \to \varepsilon \in \mathbb{R} \to 0 < \varepsilon \to 0 < \varepsilon \cdot 2^{-1}$
\item $\varphi_{84} :\equiv \forall x\,(\mathrm{RFn}(f) \to x \in \mathrm{dom}(f) \to f(x) \in \mathbb{R}) \to \mathrm{RFn}(f) \to y \in \mathrm{dom}(f) \to f(y) \in \mathbb{R}$
\item $\varphi_{85} :\equiv N_{0} \cup N_{1} \in \mathbb{N} \to 0 < \lvert p(N_{0} \cup N_{1}) - c \rvert \land \lvert p(N_{0} \cup N_{1}) - c \rvert < r$
\item $\varphi_{86} :\equiv x \in \mathrm{dom}(f) \to 0 < \lvert x - c \rvert \to \lvert x - c \rvert < \delta \to \varepsilon \cdot 2^{-1} + \varepsilon \cdot 2^{-1} = \varepsilon$
\item $\varphi_{87} :\equiv \exists a\,(\mathrm{Seq}(a) \land \mathrm{Conv}(a,c) \land \forall n \in \mathbb{N}\,(0 < \lvert a_{n} - c \rvert \land \lvert a_{n} - c \rvert < r))$
\item $\varphi_{88} :\equiv \forall n \in \mathbb{N}\,(N \leq n \to \lvert p(n) - c \rvert < \delta) \to n \in \mathbb{N} \to N \leq n \to \lvert p(n) - c \rvert < \delta$
\item $\varphi_{89} :\equiv \forall n \in \mathbb{N}\,(N \leq n \to \lvert p(n) - c \rvert < \delta) \to m \in \mathbb{N} \to N \leq m \to \lvert p(m) - c \rvert < \delta$
\item $\varphi_{90} :\equiv m \in \mathbb{N} \to n \in \mathbb{N} \to N \leq m \to N \leq n \to \lvert f \circ p(m) - f \circ p(n) \rvert < \varepsilon$
\item $\varphi_{91} :\equiv y \in \mathrm{dom}(f) \to 0 < \lvert y - c \rvert \to \lvert y - c \rvert < \delta \to \lvert f(y) - \alpha \rvert < \varepsilon \cdot 2^{-1}$
\item $\varphi_{92} :\equiv x \in \mathrm{dom}(f) \to 0 < \lvert x - c \rvert \to \lvert x - c \rvert < \delta \to \lvert f(x) - \alpha \rvert < \varepsilon \cdot 2^{-1}$
\item $\varphi_{93} :\equiv \lvert x - c \rvert < \delta \to 0 < \lvert y - c \rvert \to \lvert y - c \rvert < \delta \to \lvert f(x) - f(y) \rvert < \varepsilon$
\item $\varphi_{94} :\equiv \forall n\,(N_{0} \in \mathbb{N} \to n \in \mathbb{N} \to N_{0} \cup n \in \mathbb{N}) \to N_{0} \in \mathbb{N} \to N_{1} \in \mathbb{N} \to N_{0} \cup N_{1} \in \mathbb{N}$
\item $\varphi_{95} :\equiv 0 < \lvert x - c \rvert \to \lvert x - c \rvert < \delta \to \lvert f(x) - \alpha \rvert < \varepsilon \cdot 2^{-1} + \varepsilon \cdot 2^{-1}$
\item $\varphi_{96} :\equiv N_{0} \cup N_{1} \in \mathbb{N} \to N_{1} \leq N_{0} \cup N_{1} \to \lvert f \circ p(N_{0} \cup N_{1}) - \alpha \rvert < \varepsilon \cdot 2^{-1}$
\item $\varphi_{97} :\equiv N \in \mathbb{N} \land \forall n \in \mathbb{N}\,(N \leq n \to \lvert p(n) - c \rvert < \delta) \to m \in \mathbb{N} \to n \in \mathbb{N} \to n \in \mathbb{N}$
\item $\varphi_{98} :\equiv N \in \mathbb{N} \land \forall n \in \mathbb{N}\,(N \leq n \to \lvert p(n) - c \rvert < \delta) \to m \in \mathbb{N} \to n \in \mathbb{N} \to N \leq m \to N \leq m$
\item $\varphi_{99} :\equiv \forall x \in \mathrm{dom}(f)\,(0 < \lvert x - c \rvert \to \lvert x - c \rvert < \delta \to \lvert f(x) - \alpha \rvert < \varepsilon \cdot 2^{-1})$
\item $\varphi_{100} :\equiv \lvert x - c \rvert < \delta \to 0 < \lvert y - c \rvert \to \lvert y - c \rvert < \delta \to \varepsilon \cdot 2^{-1} + \varepsilon \cdot 2^{-1} = \varepsilon$
\item $\varphi_{101} :\equiv \forall m\,\forall n\,(m \in \mathbb{N} \to n \in \mathbb{N} \to m \cup n \in \mathbb{N}) \to \forall n\,(N_{0} \in \mathbb{N} \to n \in \mathbb{N} \to N_{0} \cup n \in \mathbb{N})$
\item $\varphi_{102} :\equiv \forall f\,\forall x\,(\mathrm{RFn}(f) \to x \in \mathrm{dom}(f) \to f(x) \in \mathbb{R}) \to \forall x\,(\mathrm{RFn}(f) \to x \in \mathrm{dom}(f) \to f(x) \in \mathbb{R})$
\item $\varphi_{103} :\equiv \lvert p(N_{0} \cup N_{1}) - c \rvert < \delta \to \lvert f(x) - f(p(N_{0} \cup N_{1})) \rvert < \varepsilon \cdot 2^{-1}$
\item $\varphi_{104} :\equiv \lvert f(x) - f(y) \rvert < \varepsilon \cdot 2^{-1} + \varepsilon \cdot 2^{-1} \to \lvert f(x) - f(y) \rvert < \varepsilon$
\item $\varphi_{105} :\equiv \lvert x - c \rvert < \delta \to 0 < \lvert y - c \rvert \to \lvert y - c \rvert < \delta \to \lvert f(y) - \alpha \rvert < \varepsilon \cdot 2^{-1}$
\item $\varphi_{106} :\equiv \lvert x - c \rvert < \delta \to 0 < \lvert y - c \rvert \to \lvert y - c \rvert < \delta \to \lvert \alpha - f(y) \rvert < \varepsilon \cdot 2^{-1}$
\item $\varphi_{107} :\equiv \lvert x - c \rvert < \delta \to 0 < \lvert y - c \rvert \to \lvert y - c \rvert < \delta \to \lvert f(x) - \alpha \rvert < \varepsilon \cdot 2^{-1}$
\item $\varphi_{108} :\equiv \lvert x - c \rvert < \delta \to 0 < \lvert y - c \rvert \to \lvert y - c \rvert < \delta \to \lvert \alpha - f(x) \rvert < \varepsilon \cdot 2^{-1}$
\item $\varphi_{109} :\equiv \forall e \in \mathbb{R}\,(e \cdot 2^{-1} + e \cdot 2^{-1} = e) \to \varepsilon \in \mathbb{R} \to \varepsilon \cdot 2^{-1} + \varepsilon \cdot 2^{-1} = \varepsilon$
\item $\varphi_{110} :\equiv \neg (N \in \mathbb{N} \land \forall n \in \mathbb{N}\,(N \leq n \to \lvert p(n) - c \rvert < \delta \land \lvert f \circ p(n) - \alpha \rvert < \delta))$
\item $\varphi_{111} :\equiv \forall x\,\forall y\,(x \in \mathbb{R} \to y \in \mathbb{R} \to x \cdot y \in \mathbb{R}) \to \forall y\,(\varepsilon \in \mathbb{R} \to y \in \mathbb{R} \to \varepsilon \cdot y \in \mathbb{R})$
\item $\varphi_{112} :\equiv 0 < \lvert p(N_{0} \cup N_{1}) - c \rvert \to \lvert p(N_{0} \cup N_{1}) - c \rvert < r \to p(N_{0} \cup N_{1}) \in \mathrm{dom}(f)$
\item $\varphi_{113} :\equiv \lvert \alpha - f(y) \rvert < \varepsilon \cdot 2^{-1} \to \lvert f(x) - f(y) \rvert < \varepsilon \cdot 2^{-1} + \varepsilon \cdot 2^{-1}$
\item $\varphi_{114} :\equiv \exists N \in \mathbb{N}\,\forall n \in \mathbb{N}\,(N \leq n \to \lvert p(n) - c \rvert < \delta \land \lvert f \circ p(n) - \alpha \rvert < \delta)$
\item $\varphi_{115} :\equiv \lvert f(p(m)) - f(p(n)) \rvert < \varepsilon \to \lvert f \circ p(m) - f(p(n)) \rvert < \varepsilon$
\item $\varphi_{116} :\equiv N \in \mathbb{N} \land \forall n \in \mathbb{N}\,(N \leq n \to \lvert p(n) - c \rvert < \delta) \to m \in \mathbb{N} \to n \in \mathbb{N} \to N \leq m \to N \leq n \to N \leq n$
\item $\varphi_{117} :\equiv 0 < \lvert p(n) - c \rvert \to \lvert p(n) - c \rvert < \delta \to \lvert f(p(m)) - f(p(n)) \rvert < \varepsilon$
\item $\varphi_{118} :\equiv \forall n\,(\mathrm{Seq}(p) \to n \in \mathbb{N} \to p(n) \in \mathbb{R}) \to \mathrm{Seq}(p) \to N_{0} \cup N_{1} \in \mathbb{N} \to p(N_{0} \cup N_{1}) \in \mathbb{R}$
\item $\varphi_{119} :\equiv \forall n \in \mathbb{N}\,(f \circ p(n) = f(p(n))) \to m \in \mathbb{N} \to f \circ p(m) = f(p(m))$
\item $\varphi_{120} :\equiv \forall n \in \mathbb{N}\,(f \circ p(n) = f(p(n))) \to n \in \mathbb{N} \to f \circ p(n) = f(p(n))$
\item $\varphi_{121} :\equiv \forall w\,(f(p(m)) = w \to \lvert f(p(m)) - f(p(n)) \rvert < \varepsilon \to \lvert w - f(p(n)) \rvert < \varepsilon)$
\item $\varphi_{122} :\equiv \forall \varepsilon\,(\mathrm{Conv}(p,c) \to \varepsilon \in \mathbb{R} \to 0 < \varepsilon \to \exists N \in \mathbb{N}\,\forall n \in \mathbb{N}\,(N \leq n \to \lvert p(n) - c \rvert < \varepsilon))$
\item $\varphi_{123} :\equiv \forall y\,(\varepsilon \in \mathbb{R} \to y \in \mathbb{R} \to \varepsilon \cdot y \in \mathbb{R}) \to \varepsilon \in \mathbb{R} \to 2^{-1} \in \mathbb{R} \to \varepsilon \cdot 2^{-1} \in \mathbb{R}$
\item $\varphi_{124} :\equiv \forall n \in \mathbb{N}\,(p(n) \in \mathrm{dom}(f)) \to \mathrm{Seq}(f \circ p) \land \forall n \in \mathbb{N}\,(f \circ p(n) = f(p(n)))$
\item $\varphi_{125} :\equiv \lvert f \circ p(m) - f(p(n)) \rvert < \varepsilon \to \lvert f \circ p(m) - f \circ p(n) \rvert < \varepsilon$
\item $\varphi_{126} :\equiv \lvert f(p(N_{0} \cup N_{1})) - \alpha \rvert < \varepsilon \cdot 2^{-1} \to \lvert f(x) - \alpha \rvert < \varepsilon \cdot 2^{-1} + \varepsilon \cdot 2^{-1}$
\item $\varphi_{127} :\equiv \forall \alpha\,\forall \varepsilon\,(\mathrm{Conv}(p,\alpha) \to \varepsilon \in \mathbb{R} \to 0 < \varepsilon \to \exists N \in \mathbb{N}\,\forall n \in \mathbb{N}\,(N \leq n \to \lvert p(n) - \alpha \rvert < \varepsilon))$
\item $\varphi_{128} :\equiv \forall f\,\forall c\,(\mathrm{PNbhd}(f,c) \leftrightarrow \exists r\,(r \in \mathbb{R} \land 0 < r \land \varphi_{16}))$
\item $\varphi_{129} :\equiv \lvert f \circ p(N_{0} \cup N_{1}) - \alpha \rvert < \varepsilon \cdot 2^{-1} \to \lvert f(p(N_{0} \cup N_{1})) - \alpha \rvert < \varepsilon \cdot 2^{-1}$
\item $\varphi_{130} :\equiv \forall w\,(f(p(n)) = w \to \lvert f \circ p(m) - f(p(n)) \rvert < \varepsilon \to \lvert f \circ p(m) - w \rvert < \varepsilon)$
\item $\varphi_{131} :\equiv \forall a\,\forall \alpha\,\forall \varepsilon\,(\mathrm{Conv}(a,\alpha) \to \varepsilon \in \mathbb{R} \to 0 < \varepsilon \to \exists N \in \mathbb{N}\,\forall n \in \mathbb{N}\,(N \leq n \to \lvert a_{n} - \alpha \rvert < \varepsilon))$
\item $\varphi_{132} :\equiv N \in \mathbb{N} \land \forall n \in \mathbb{N}\,(N \leq n \to \lvert p(n) - c \rvert < \delta) \to m \in \mathbb{N} \to n \in \mathbb{N} \to N \leq m \to N \leq n \to p(n) \in \mathbb{R}$
\item $\varphi_{133} :\equiv N \in \mathbb{N} \land \forall n \in \mathbb{N}\,(N \leq n \to \lvert p(n) - c \rvert < \delta) \to m \in \mathbb{N} \to n \in \mathbb{N} \to N \leq m \to N \leq n \to p(m) \in \mathbb{R}$
\item $\varphi_{134} :\equiv N \in \mathbb{N} \land \forall n \in \mathbb{N}\,(N \leq n \to \lvert p(n) - c \rvert < \delta) \to m \in \mathbb{N} \to n \in \mathbb{N} \to N \leq m \to N \leq n \to p(n) \in \mathrm{dom}(f)$
\item $\varphi_{135} :\equiv N \in \mathbb{N} \land \forall n \in \mathbb{N}\,(N \leq n \to \lvert p(n) - c \rvert < \delta) \to m \in \mathbb{N} \to n \in \mathbb{N} \to N \leq m \to N \leq n \to p(m) \in \mathrm{dom}(f)$
\item $\varphi_{136} :\equiv \lvert f(y) - \alpha \rvert = \lvert \alpha - f(y) \rvert \to \lvert f(y) - \alpha \rvert < \varepsilon \cdot 2^{-1} \to \lvert \alpha - f(y) \rvert < \varepsilon \cdot 2^{-1}$
\item $\varphi_{137} :\equiv \lvert f(x) - \alpha \rvert = \lvert \alpha - f(x) \rvert \to \lvert f(x) - \alpha \rvert < \varepsilon \cdot 2^{-1} \to \lvert \alpha - f(x) \rvert < \varepsilon \cdot 2^{-1}$
\item $\varphi_{138} :\equiv N \in \mathbb{N} \land \forall n \in \mathbb{N}\,(N \leq n \to \lvert p(n) - c \rvert < \delta) \to N \in \mathbb{N} \land \forall n \in \mathbb{N}\,(N \leq n \to \lvert p(n) - c \rvert < \delta)$
\item $\varphi_{139} :\equiv \forall w\,(\varepsilon \cdot 2^{-1} + \varepsilon \cdot 2^{-1} = w \to \lvert f(x) - \alpha \rvert < \varepsilon \cdot 2^{-1} + \varepsilon \cdot 2^{-1} \to \lvert f(x) - \alpha \rvert < w)$
\item $\varphi_{140} :\equiv N \in \mathbb{N} \land \forall n \in \mathbb{N}\,(N \leq n \to \lvert p(n) - c \rvert < \delta) \to m \in \mathbb{N} \to n \in \mathbb{N} \to N \leq m \to N \leq n \to \lvert p(n) - c \rvert < \delta$
\item $\varphi_{141} :\equiv N \in \mathbb{N} \land \forall n \in \mathbb{N}\,(N \leq n \to \lvert p(n) - c \rvert < \delta) \to m \in \mathbb{N} \to n \in \mathbb{N} \to N \leq m \to N \leq n \to \lvert p(m) - c \rvert < \delta$
\item $\varphi_{142} :\equiv \forall x\,\forall y\,\forall z\,\forall e\,(x \in \mathbb{R} \to y \in \mathbb{R} \to z \in \mathbb{R} \to e \in \mathbb{R} \to \lvert x - y \rvert < e \to \lvert y - z \rvert < e \to \lvert x - z \rvert < e + e)$
\item $\varphi_{143} :\equiv \forall w\,(f \circ p(N_{0} \cup N_{1}) = w \to \lvert f \circ p(N_{0} \cup N_{1}) - \alpha \rvert < \varepsilon \cdot 2^{-1} \to \lvert w - \alpha \rvert < \varepsilon \cdot 2^{-1})$
\item $\varphi_{144} :\equiv \forall y\,(f \circ p(m) = y \to y = f \circ p(m)) \to f \circ p(m) = f(p(m)) \to f(p(m)) = f \circ p(m)$
\item $\varphi_{145} :\equiv \forall y\,(f \circ p(n) = y \to y = f \circ p(n)) \to f \circ p(n) = f(p(n)) \to f(p(n)) = f \circ p(n)$
\item $\varphi_{146} :\equiv \forall n \in \mathbb{N}\,(0 < \lvert p(n) - c \rvert \land \lvert p(n) - c \rvert < r) \to n \in \mathbb{N} \to 0 < \lvert p(n) - c \rvert \land \lvert p(n) - c \rvert < r$
\item $\varphi_{147} :\equiv \forall a\,(\mathrm{RFn}(f) \to \mathrm{Seq}(a) \to \forall n \in \mathbb{N}\,(a_{n} \in \mathrm{dom}(f)) \to \mathrm{Seq}(f \circ a) \land \forall n \in \mathbb{N}\,(f \circ a(n) = f(a_{n})))$
\item $\varphi_{148} :\equiv \forall \varepsilon\,(\mathrm{Conv}(f \circ p,\alpha) \to \varepsilon \in \mathbb{R} \to 0 < \varepsilon \to \exists N \in \mathbb{N}\,\forall n \in \mathbb{N}\,(N \leq n \to \lvert f \circ p(n) - \alpha \rvert < \varepsilon))$
\item $\varphi_{149} :\equiv \forall n \in \mathbb{N}\,(0 < \lvert p(n) - c \rvert \land \lvert p(n) - c \rvert < r) \to m \in \mathbb{N} \to 0 < \lvert p(m) - c \rvert \land \lvert p(m) - c \rvert < r$
\item $\varphi_{150} :\equiv \forall w\,(\varepsilon \cdot 2^{-1} + \varepsilon \cdot 2^{-1} = w \to \lvert f(x) - f(y) \rvert < \varepsilon \cdot 2^{-1} + \varepsilon \cdot 2^{-1} \to \lvert f(x) - f(y) \rvert < w)$
\item $\varphi_{151} :\equiv 2^{-1} \in \mathbb{R} \to \mathrm{RFn}(f) \to c \in \mathbb{R} \to \mathrm{PNbhd}(f,c) \to \varphi_{68}$
\item $\varphi_{152} :\equiv v_{773} \in \mathbb{R} \to v_{774} \in \mathbb{R} \to v_{775} \in \mathbb{R} \to \lvert v_{772} - v_{773} \rvert < v_{775} \to \varphi_{9}$
\item $\varphi_{153} :\equiv N \in \mathbb{N} \land \forall n \in \mathbb{N}\,(N \leq n \to \lvert p(n) - c \rvert < \delta) \to m \in \mathbb{N} \to \varphi_{17}$
\item $\varphi_{154} :\equiv N \in \mathbb{N} \land \forall n \in \mathbb{N}\,(N \leq n \to \lvert p(n) - c \rvert < \delta) \to m \in \mathbb{N} \to \varphi_{18}$
\item $\varphi_{155} :\equiv \forall y\,\forall z\,\forall e\,(f(x) \in \mathbb{R} \to y \in \mathbb{R} \to z \in \mathbb{R} \to e \in \mathbb{R} \to \lvert f(x) - y \rvert < e \to \lvert y - z \rvert < e \to \lvert f(x) - z \rvert < e + e)$
\item $\varphi_{156} :\equiv N \in \mathbb{N} \land \forall n \in \mathbb{N}\,(N \leq n \to \lvert p(n) - c \rvert < \delta) \to m \in \mathbb{N} \to \varphi_{23}$
\item $\varphi_{157} :\equiv N \in \mathbb{N} \land \forall n \in \mathbb{N}\,(N \leq n \to \lvert p(n) - c \rvert < \delta) \to m \in \mathbb{N} \to \varphi_{26}$
\item $\varphi_{158} :\equiv \forall x\,\forall \alpha\,\forall \beta\,\forall \varepsilon\,(x \in \mathbb{R} \to \alpha \in \mathbb{R} \to \beta \in \mathbb{R} \to \varepsilon \in \mathbb{R} \to \lvert x - \alpha \rvert < \varepsilon \to \lvert x - \beta \rvert < \varepsilon \to \lvert \alpha - \beta \rvert < \varepsilon + \varepsilon)$
\item $\varphi_{159} :\equiv N \in \mathbb{N} \land \forall n \in \mathbb{N}\,(N \leq n \to \lvert p(n) - c \rvert < \delta) \to m \in \mathbb{N} \to \varphi_{31}$
\item $\varphi_{160} :\equiv N \in \mathbb{N} \land \forall n \in \mathbb{N}\,(N \leq n \to \lvert p(n) - c \rvert < \delta) \to m \in \mathbb{N} \to \varphi_{34}$
\item $\varphi_{161} :\equiv \forall y\,(x \in \mathrm{dom}(f) \to y \in \mathrm{dom}(f) \to 0 < \lvert x - c \rvert \to \varphi_{93})$
\item $\varphi_{162} :\equiv \forall a\,(\mathrm{Cauchy}(a) \leftrightarrow \mathrm{Seq}(a) \land \forall \varepsilon \in \mathbb{R}\,(0 < \varepsilon \to \exists N \in \mathbb{N}\,\forall m\,\forall n\,(m \in \mathbb{N} \to n \in \mathbb{N} \to N \leq m \to N \leq n \to \lvert a_{m} - a_{n} \rvert < \varepsilon)))$
\item $\varphi_{163} :\equiv N \in \mathbb{N} \land \forall n \in \mathbb{N}\,(N \leq n \to \lvert p(n) - c \rvert < \delta) \to m \in \mathbb{N} \to \varphi_{43}$
\item $\varphi_{164} :\equiv N \in \mathbb{N} \land \forall n \in \mathbb{N}\,(N \leq n \to \lvert p(n) - c \rvert < \delta) \to m \in \mathbb{N} \to \varphi_{44}$
\item $\varphi_{165} :\equiv N \in \mathbb{N} \land \forall n \in \mathbb{N}\,(N \leq n \to \lvert p(n) - c \rvert < \delta) \to m \in \mathbb{N} \to \varphi_{47}$
\item $\varphi_{166} :\equiv N \in \mathbb{N} \land \forall n \in \mathbb{N}\,(N \leq n \to \lvert p(n) - c \rvert < \delta) \to m \in \mathbb{N} \to \varphi_{48}$
\item $\varphi_{167} :\equiv x \in \mathrm{dom}(f) \to y \in \mathrm{dom}(f) \to 0 < \lvert x - c \rvert \to \lvert x - c \rvert < \delta \to \varphi_{55}$
\item $\varphi_{168} :\equiv x \in \mathrm{dom}(f) \to y \in \mathrm{dom}(f) \to 0 < \lvert x - c \rvert \to \lvert x - c \rvert < \delta \to \varphi_{56}$
\item $\varphi_{169} :\equiv \varphi_{42} \to f(x) \in \mathbb{R} \to \alpha \in \mathbb{R} \to \lvert f(x) - \alpha \rvert = \lvert \alpha - f(x) \rvert$
\item $\varphi_{170} :\equiv \forall \alpha\,\forall \beta\,\forall \varepsilon\,(v_{772} \in \mathbb{R} \to \alpha \in \mathbb{R} \to \beta \in \mathbb{R} \to \varepsilon \in \mathbb{R} \to \lvert v_{772} - \alpha \rvert < \varepsilon \to \lvert v_{772} - \beta \rvert < \varepsilon \to \lvert \alpha - \beta \rvert < \varepsilon + \varepsilon)$
\item $\varphi_{171} :\equiv \forall \beta\,\forall \varepsilon\,(v_{772} \in \mathbb{R} \to v_{773} \in \mathbb{R} \to \beta \in \mathbb{R} \to \varepsilon \in \mathbb{R} \to \lvert v_{772} - v_{773} \rvert < \varepsilon \to \lvert v_{772} - \beta \rvert < \varepsilon \to \lvert v_{773} - \beta \rvert < \varepsilon + \varepsilon)$
\item $\varphi_{172} :\equiv \varphi_{67} \to 0 < \lvert p(N_{0} \cup N_{1}) - c \rvert \land \lvert p(N_{0} \cup N_{1}) - c \rvert < r$
\item $\varphi_{173} :\equiv y \in \mathrm{dom}(f) \to 0 < \lvert x - c \rvert \to \lvert x - c \rvert < \delta \to 0 < \lvert y - c \rvert \to \varphi_{58}$
\item $\varphi_{174} :\equiv \forall b\,\forall \beta\,\forall \varepsilon\,(\mathrm{Conv}(p,c) \to \mathrm{Conv}(b,\beta) \to \varepsilon \in \mathbb{R} \to 0 < \varepsilon \to \exists N \in \mathbb{N}\,\forall n \in \mathbb{N}\,(N \leq n \to \lvert p(n) - c \rvert < \varepsilon \land \lvert b_{n} - \beta \rvert < \varepsilon))$
\item $\varphi_{175} :\equiv \forall \varepsilon\,(v_{772} \in \mathbb{R} \to v_{773} \in \mathbb{R} \to v_{774} \in \mathbb{R} \to \varepsilon \in \mathbb{R} \to \lvert v_{772} - v_{773} \rvert < \varepsilon \to \lvert v_{772} - v_{774} \rvert < \varepsilon \to \lvert v_{773} - v_{774} \rvert < \varepsilon + \varepsilon)$
\item $\varphi_{176} :\equiv \forall y\,(x \in \mathrm{dom}(f) \to y \in \mathrm{dom}(f) \to 0 < \lvert x - c \rvert \to \lvert x - c \rvert < \delta \to 0 < \lvert y - c \rvert \to \lvert y - c \rvert < \delta \to \lvert f(x) - f(y) \rvert < \varepsilon \cdot 2^{-1})$
\item $\varphi_{177} :\equiv N \in \mathbb{N} \land \forall n \in \mathbb{N}\,(N \leq n \to \lvert p(n) - c \rvert < \delta) \to \forall n\,(\varphi_{90})$
\item $\varphi_{178} :\equiv N \in \mathbb{N} \land \forall n \in \mathbb{N}\,(N \leq n \to \lvert p(n) - c \rvert < \delta) \to \forall m\,\forall n\,(\varphi_{90})$
\item $\varphi_{179} :\equiv \forall \varepsilon\,(\mathrm{FLim}(f,c,\alpha) \to \varepsilon \in \mathbb{R} \to 0 < \varepsilon \to \exists \delta\,(\delta \in \mathbb{R} \land 0 < \delta \land \varphi_{69}))$
\item $\varphi_{180} :\equiv \varphi_{76} \to \mathrm{Seq}(f \circ p) \land \forall n \in \mathbb{N}\,(f \circ p(n) = f(p(n)))$
\item $\varphi_{181} :\equiv 0 < \varepsilon \to \delta \in \mathbb{R} \land 0 < \delta \land \varphi_{99} \to x \in \mathrm{dom}(f) \to x \in \mathrm{dom}(f)$
\item $\varphi_{182} :\equiv \forall \alpha\,\forall b\,\forall \beta\,\forall \varepsilon\,(\mathrm{Conv}(p,\alpha) \to \mathrm{Conv}(b,\beta) \to \varepsilon \in \mathbb{R} \to 0 < \varepsilon \to \exists N \in \mathbb{N}\,\forall n \in \mathbb{N}\,(N \leq n \to \lvert p(n) - \alpha \rvert < \varepsilon \land \lvert b_{n} - \beta \rvert < \varepsilon))$
\item $\varphi_{183} :\equiv \mathrm{Conv}(p,c) \to \mathrm{Conv}(f \circ p,\alpha) \to \delta \in \mathbb{R} \to 0 < \delta \to \varphi_{114}$
\item $\varphi_{184} :\equiv \varphi_{76} \to \mathrm{Conv}(f \circ p,\alpha) \to \varepsilon \in \mathbb{R} \to 0 < \varepsilon \to \varepsilon \cdot 2^{-1} \in \mathbb{R}$
\item $\varphi_{185} :\equiv \varphi_{75} \to f(y) \in \mathbb{R} \to \alpha \in \mathbb{R} \to \lvert f(y) - \alpha \rvert = \lvert \alpha - f(y) \rvert$
\item $\varphi_{186} :\equiv \mathrm{Conv}(f \circ p,\alpha) \to \varepsilon \in \mathbb{R} \to 0 < \varepsilon \to \exists \delta\,(\delta \in \mathbb{R} \land 0 < \delta \land \varphi_{69})$
\item $\varphi_{187} :\equiv \forall a\,\forall \alpha\,\forall b\,\forall \beta\,\forall \varepsilon\,(\mathrm{Conv}(a,\alpha) \to \mathrm{Conv}(b,\beta) \to \varepsilon \in \mathbb{R} \to 0 < \varepsilon \to \exists N \in \mathbb{N}\,\forall n \in \mathbb{N}\,(N \leq n \to \lvert a_{n} - \alpha \rvert < \varepsilon \land \lvert b_{n} - \beta \rvert < \varepsilon))$
\item $\varphi_{188} :\equiv N \in \mathbb{N} \land \forall n \in \mathbb{N}\,(N \leq n \to \lvert p(n) - c \rvert < \delta) \to m \in \mathbb{N} \to n \in \mathbb{N} \to N \leq m \to N \leq n \to \varphi_{16}$
\item $\varphi_{189} :\equiv \delta \in \mathbb{R} \land 0 < \delta \land \varphi_{99} \to x \in \mathrm{dom}(f) \to y \in \mathrm{dom}(f) \to y \in \mathrm{dom}(f)$
\item $\varphi_{190} :\equiv \mathrm{FLim}(f,c,\alpha) \to \varepsilon \in \mathbb{R} \to 0 < \varepsilon \to \exists \delta\,(\delta \in \mathbb{R} \land 0 < \delta \land \varphi_{99})$
\item $\varphi_{191} :\equiv (\mathrm{RFn}(f) \to c \in \mathbb{R} \to \mathrm{PNbhd}(f,c) \to \varphi_{4}) \to \mathrm{RFn}(f) \to c \in \mathbb{R} \to \mathrm{PNbhd}(f,c) \to \varphi_{3}$
\item $\varphi_{192} :\equiv \forall v_{773}\,\forall v_{774}\,\forall v_{775}\,(\alpha \in \mathbb{R} \to v_{773} \in \mathbb{R} \to v_{774} \in \mathbb{R} \to v_{775} \in \mathbb{R} \to \lvert \alpha - v_{773} \rvert < v_{775} \to \lvert \alpha - v_{774} \rvert < v_{775} \to \lvert v_{773} - v_{774} \rvert < v_{775} + v_{775})$
\item $\varphi_{193} :\equiv \mathrm{Conv}(f \circ p,\alpha) \to \forall \varepsilon \in \mathbb{R}\,(0 < \varepsilon \to \exists \delta\,(\delta \in \mathbb{R} \land 0 < \delta \land \varphi_{69}))$
\item $\varphi_{194} :\equiv \varepsilon \cdot 2^{-1} \in \mathbb{R} \to 0 < \varepsilon \cdot 2^{-1} \to \exists \delta\,(\delta \in \mathbb{R} \land 0 < \delta \land \varphi_{99})$
\item $\varphi_{195} :\equiv \forall \varepsilon\,(\mathrm{Conv}(p,c) \to \mathrm{Conv}(f \circ p,\alpha) \to \varepsilon \in \mathbb{R} \to 0 < \varepsilon \to \exists N \in \mathbb{N}\,\forall n \in \mathbb{N}\,(N \leq n \to \lvert p(n) - c \rvert < \varepsilon \land \lvert f \circ p(n) - \alpha \rvert < \varepsilon))$
\item $\varphi_{196} :\equiv 0 < \lvert x - c \rvert \to \lvert x - c \rvert < \delta \to 0 < \lvert p(N_{0} \cup N_{1}) - c \rvert \to \varphi_{103}$
\item $\varphi_{197} :\equiv N \in \mathbb{N} \land \forall n \in \mathbb{N}\,(N \leq n \to \lvert p(n) - c \rvert < \delta) \to N \in \mathbb{N} \land \forall m\,\forall n\,(\varphi_{90})$
\item $\varphi_{198} :\equiv \forall v_{774}\,\forall v_{775}\,(\alpha \in \mathbb{R} \to f(x) \in \mathbb{R} \to v_{774} \in \mathbb{R} \to v_{775} \in \mathbb{R} \to \lvert \alpha - f(x) \rvert < v_{775} \to \lvert \alpha - v_{774} \rvert < v_{775} \to \lvert f(x) - v_{774} \rvert < v_{775} + v_{775})$
\item $\varphi_{199} :\equiv \forall y\,(p(m) \in \mathrm{dom}(f) \to y \in \mathrm{dom}(f) \to 0 < \lvert p(m) - c \rvert \to \lvert p(m) - c \rvert < \delta \to 0 < \lvert y - c \rvert \to \lvert y - c \rvert < \delta \to \lvert f(p(m)) - f(y) \rvert < \varepsilon)$
\item $\varphi_{200} :\equiv N \in \mathbb{N} \land \forall n \in \mathbb{N}\,(N \leq n \to \lvert p(n) - c \rvert < \delta) \to \exists N \in \mathbb{N}\,\forall m\,\forall n\,(\varphi_{90})$
\item $\varphi_{201} :\equiv \forall \beta\,\forall \varepsilon\,(\mathrm{Conv}(p,c) \to \mathrm{Conv}(f \circ p,\beta) \to \varepsilon \in \mathbb{R} \to 0 < \varepsilon \to \exists N \in \mathbb{N}\,\forall n \in \mathbb{N}\,(N \leq n \to \lvert p(n) - c \rvert < \varepsilon \land \lvert f \circ p(n) - \beta \rvert < \varepsilon))$
\item $\varphi_{202} :\equiv p(n) \in \mathrm{dom}(f) \to 0 < \lvert p(m) - c \rvert \to \lvert p(m) - c \rvert < \delta \to \varphi_{117}$
\item $\varphi_{203} :\equiv \forall v_{775}\,(\alpha \in \mathbb{R} \to f(x) \in \mathbb{R} \to f(y) \in \mathbb{R} \to v_{775} \in \mathbb{R} \to \lvert \alpha - f(x) \rvert < v_{775} \to \lvert \alpha - f(y) \rvert < v_{775} \to \lvert f(x) - f(y) \rvert < v_{775} + v_{775})$
\item $\varphi_{204} :\equiv f(y) \in \mathbb{R} \to \varepsilon \cdot 2^{-1} \in \mathbb{R} \to \lvert \alpha - f(x) \rvert < \varepsilon \cdot 2^{-1} \to \varphi_{113}$
\item $\varphi_{205} :\equiv \mathrm{Cauchy}(f \circ p) \leftrightarrow \mathrm{Seq}(f \circ p) \land \forall \varepsilon \in \mathbb{R}\,(0 < \varepsilon \to \exists N \in \mathbb{N}\,\forall m\,\forall n\,(\varphi_{90}))$
\item $\varphi_{206} :\equiv N \in \mathbb{N} \land \forall n \in \mathbb{N}\,(N \leq n \to \lvert p(n) - c \rvert < \delta) \to m \in \mathbb{N} \to n \in \mathbb{N} \to N \leq m \to N \leq n \to p(n) \in \mathbb{R} \to \varphi_{27}$
\item $\varphi_{207} :\equiv N \in \mathbb{N} \land \forall n \in \mathbb{N}\,(N \leq n \to \lvert p(n) - c \rvert < \delta) \to m \in \mathbb{N} \to n \in \mathbb{N} \to N \leq m \to N \leq n \to p(m) \in \mathbb{R} \to \varphi_{30}$
\item $\varphi_{208} :\equiv \varepsilon \cdot 2^{-1} \in \mathbb{R} \to \lvert f(x) - f(p(N_{0} \cup N_{1})) \rvert < \varepsilon \cdot 2^{-1} \to \varphi_{126}$
\item $\varphi_{209} :\equiv \mathrm{FLim}(f,c,\alpha) \leftrightarrow \mathrm{RFn}(f) \land c \in \mathbb{R} \land \alpha \in \mathbb{R} \land \mathrm{PNbhd}(f,c) \land \forall \varepsilon \in \mathbb{R}\,(0 < \varepsilon \to \exists \delta\,(\delta \in \mathbb{R} \land 0 < \delta \land \varphi_{69}))$
\item $\varphi_{210} :\equiv \forall z\,\forall e\,(f(x) \in \mathbb{R} \to f(p(N_{0} \cup N_{1})) \in \mathbb{R} \to z \in \mathbb{R} \to e \in \mathbb{R} \to \lvert f(x) - f(p(N_{0} \cup N_{1})) \rvert < e \to \lvert f(p(N_{0} \cup N_{1})) - z \rvert < e \to \lvert f(x) - z \rvert < e + e)$
\item $\varphi_{211} :\equiv \forall e\,(f(x) \in \mathbb{R} \to f(p(N_{0} \cup N_{1})) \in \mathbb{R} \to \alpha \in \mathbb{R} \to e \in \mathbb{R} \to \lvert f(x) - f(p(N_{0} \cup N_{1})) \rvert < e \to \lvert f(p(N_{0} \cup N_{1})) - \alpha \rvert < e \to \lvert f(x) - \alpha \rvert < e + e)$
\item $\varphi_{212} :\equiv \exists \alpha\,(\mathrm{FLim}(f,c,\alpha)) \to \mathrm{FLim}(f,c,\alpha) \to \varepsilon \in \mathbb{R} \to \varphi_{181}$
\item $\varphi_{213} :\equiv c \in \mathbb{R} \to \mathrm{PNbhd}(f,c) \to \exists \alpha\,(\mathrm{FLim}(f,c,\alpha)) \to \varphi_{190}$
\item $\varphi_{214} :\equiv \forall N_{0}\,(\neg (N_{0} \in \mathbb{N} \land \varphi_{77})) \to \forall N\,\forall N_{0}\,(\neg (N_{0} \in \mathbb{N} \land \varphi_{77}))$
\item $\varphi_{215} :\equiv r \in \mathbb{R} \land 0 < r \land \varphi_{16} \to \varphi_{76} \to \forall n \in \mathbb{N}\,(p(n) \in \mathrm{dom}(f))$
\item $\varphi_{216} :\equiv \forall c\,(\mathrm{PNbhd}(f,c) \leftrightarrow \exists r\,(r \in \mathbb{R} \land 0 < r \land \varphi_{16})) \to (\mathrm{PNbhd}(f,c) \leftrightarrow \exists r\,(r \in \mathbb{R} \land 0 < r \land \varphi_{16}))$
\item $\varphi_{217} :\equiv r \in \mathbb{R} \land 0 < r \land \varphi_{16} \to \varphi_{76} \to \mathrm{Conv}(f \circ p,\alpha) \to \mathrm{FLim}(f,c,\alpha)$
\item $\varphi_{218} :\equiv \mathrm{FLim}(f,c,\alpha) \to \varepsilon \in \mathbb{R} \to 0 < \varepsilon \to \exists \delta\,(\delta \in \mathbb{R} \land 0 < \delta \land \forall x\,(\varphi_{161}))$
\item $\varphi_{219} :\equiv (\mathrm{RFn}(f) \to c \in \mathbb{R} \to \mathrm{PNbhd}(f,c) \to \varphi_{21}) \to \mathrm{RFn}(f) \to c \in \mathbb{R} \to \mathrm{PNbhd}(f,c) \to \varphi_{57}$
\item $\varphi_{220} :\equiv r \in \mathbb{R} \land 0 < r \land \varphi_{16} \to \varphi_{76} \to \mathrm{Conv}(f \circ p,\alpha) \to \exists \alpha\,(\mathrm{FLim}(f,c,\alpha))$
\item $\varphi_{221} :\equiv \varphi_{128} \to \forall c\,(\mathrm{PNbhd}(f,c) \leftrightarrow \exists r\,(r \in \mathbb{R} \land 0 < r \land \varphi_{16}))$
\item $\varphi_{222} :\equiv \mathrm{FLim}(f,c,\alpha) \to \forall \varepsilon \in \mathbb{R}\,(0 < \varepsilon \to \exists \delta\,(\delta \in \mathbb{R} \land 0 < \delta \land \forall x\,(\varphi_{161})))$
\item $\varphi_{223} :\equiv \exists \alpha\,(\mathrm{FLim}(f,c,\alpha)) \to \mathrm{FLim}(f,c,\alpha) \to \varepsilon \in \mathbb{R} \to 0 < \varepsilon \to \varphi_{189}$
\item $\varphi_{224} :\equiv \forall \delta\,(\neg (\delta \in \mathbb{R} \land 0 < \delta \land \varphi_{69})) \to \neg (\delta \in \mathbb{R} \land 0 < \delta \land \varphi_{69})$
\item $\varphi_{225} :\equiv \forall N\,(\neg (N \in \mathbb{N} \land \forall m\,\forall n\,(\varphi_{90}))) \to \neg (N \in \mathbb{N} \land \forall m\,\forall n\,(\varphi_{90}))$
\item $\varphi_{226} :\equiv r \in \mathbb{R} \land 0 < r \land \varphi_{16} \to \varphi_{76} \to \mathrm{Conv}(f \circ p,\alpha) \to \mathrm{Conv}(f \circ p,\alpha)$
\item $\varphi_{227} :\equiv \forall \varepsilon \in \mathbb{R}\,(0 < \varepsilon \to \exists \delta\,(\delta \in \mathbb{R} \land 0 < \delta \land \forall x\,(\varphi_{161}))) \to \exists \alpha\,(\mathrm{FLim}(f,c,\alpha))$
\item $\varphi_{228} :\equiv \exists \alpha\,(\mathrm{FLim}(f,c,\alpha)) \to \forall \varepsilon \in \mathbb{R}\,(0 < \varepsilon \to \exists \delta\,(\delta \in \mathbb{R} \land 0 < \delta \land \forall x\,(\varphi_{161})))$
\item $\varphi_{229} :\equiv \exists \alpha\,(\mathrm{FLim}(f,c,\alpha)) \leftrightarrow \forall \varepsilon \in \mathbb{R}\,(0 < \varepsilon \to \exists \delta\,(\delta \in \mathbb{R} \land 0 < \delta \land \forall x\,(\varphi_{161})))$
\item $\varphi_{230} :\equiv \mathrm{FLim}(f,c,\alpha) \to \varepsilon \in \mathbb{R} \to 0 < \varepsilon \to \delta \in \mathbb{R} \land 0 < \delta \land \varphi_{99} \to \varphi_{12}$
\item $\varphi_{231} :\equiv r \in \mathbb{R} \land 0 < r \land \varphi_{16} \to \varphi_{76} \to \mathrm{Conv}(f \circ p,\alpha) \to \varepsilon \in \mathbb{R} \to \varepsilon \in \mathbb{R}$
\item $\varphi_{232} :\equiv \varepsilon \cdot 2^{-1} \in \mathbb{R} \to 0 < \varepsilon \cdot 2^{-1} \to \exists \delta\,(\delta \in \mathbb{R} \land 0 < \delta \land \forall x\,(\varphi_{176}))$
\item $\varphi_{233} :\equiv m \in \mathbb{N} \to n \in \mathbb{N} \to N \leq m \to N \leq n \to p(m) \in \mathrm{dom}(f) \to \varphi_{202}$
\item $\varphi_{234} :\equiv \delta \in \mathbb{R} \land 0 < \delta \land \forall x\,(\varphi_{161}) \to \exists N \in \mathbb{N}\,\forall n \in \mathbb{N}\,(N \leq n \to \lvert p(n) - c \rvert < \delta)$
\item $\varphi_{235} :\equiv \lvert x - c \rvert < \delta \to x \in \mathrm{dom}(f) \to p(N_{0} \cup N_{1}) \in \mathrm{dom}(f) \to \varphi_{196}$
\item $\varphi_{236} :\equiv \mathrm{Conv}(f \circ p,\alpha) \to \varepsilon \in \mathbb{R} \to 0 < \varepsilon \to \exists \delta\,(\delta \in \mathbb{R} \land 0 < \delta \land \forall x\,(\varphi_{176}))$
\item $\varphi_{237} :\equiv 0 < \varepsilon \to \delta \in \mathbb{R} \land 0 < \delta \land \varphi_{99} \to x \in \mathrm{dom}(f) \to y \in \mathrm{dom}(f) \to \varphi_{49}$
\item $\varphi_{238} :\equiv N_{0} \in \mathbb{N} \land \varphi_{77} \to \varphi_{67} \to \forall n \in \mathbb{N}\,(f \circ p(n) = f(p(n)))$
\item $\varphi_{239} :\equiv N_{0} \in \mathbb{N} \land \varphi_{77} \to \varphi_{67} \to x \in \mathrm{dom}(f) \to 0 < \lvert x - c \rvert \to 0 < \lvert x - c \rvert$
\item $\varphi_{240} :\equiv r \in \mathbb{R} \land 0 < r \land \varphi_{16} \to \varphi_{76} \to \mathrm{Conv}(f \circ p,\alpha) \to \varepsilon \in \mathbb{R} \to 0 < \varepsilon \to 0 < \varepsilon$
\item $\varphi_{241} :\equiv N_{0} \in \mathbb{N} \land \varphi_{77} \to \varphi_{67} \to \lvert f(p(N_{0} \cup N_{1})) - \alpha \rvert < \varepsilon \cdot 2^{-1}$
\item $\varphi_{242} :\equiv N_{0} \in \mathbb{N} \land \varphi_{77} \to \varphi_{67} \to f \circ p(N_{0} \cup N_{1}) = f(p(N_{0} \cup N_{1}))$
\item $\varphi_{243} :\equiv r \in \mathbb{R} \land 0 < r \land \varphi_{16} \to \varphi_{76} \to n \in \mathbb{N} \to 0 < \lvert p(n) - c \rvert \land \lvert p(n) - c \rvert < r$
\item $\varphi_{244} :\equiv \mathrm{Conv}(f \circ p,\alpha) \to \varepsilon \in \mathbb{R} \to 0 < \varepsilon \to \delta \in \mathbb{R} \land 0 < \delta \land \forall x\,(\varphi_{176}) \to 0 < \delta$
\item $\varphi_{245} :\equiv r \in \mathbb{R} \land 0 < r \land \varphi_{16} \to \varphi_{76} \to m \in \mathbb{N} \to 0 < \lvert p(m) - c \rvert \land \lvert p(m) - c \rvert < r$
\item $\varphi_{246} :\equiv N \in \mathbb{N} \land \forall n \in \mathbb{N}\,(N \leq n \to \lvert p(n) - c \rvert < \delta) \to m \in \mathbb{N} \to n \in \mathbb{N} \to N \leq m \to N \leq n \to \forall x\,(\varphi_{161})$
\item $\varphi_{247} :\equiv r \in \mathbb{R} \land 0 < r \land \varphi_{16} \to \varphi_{76} \to \forall n \in \mathbb{N}\,(0 < \lvert p(n) - c \rvert \land \lvert p(n) - c \rvert < r)$
\item $\varphi_{248} :\equiv \mathrm{Conv}(f \circ p,\alpha) \to \varepsilon \in \mathbb{R} \to 0 < \varepsilon \to \delta \in \mathbb{R} \land 0 < \delta \land \forall x\,(\varphi_{176}) \to \delta \in \mathbb{R}$
\item $\varphi_{249} :\equiv \delta \in \mathbb{R} \land 0 < \delta \land \varphi_{99} \to x \in \mathrm{dom}(f) \to y \in \mathrm{dom}(f) \to 0 < \lvert x - c \rvert \to \varphi_{50}$
\item $\varphi_{250} :\equiv \delta \in \mathbb{R} \land 0 < \delta \land \varphi_{99} \to x \in \mathrm{dom}(f) \to y \in \mathrm{dom}(f) \to 0 < \lvert x - c \rvert \to \varphi_{51}$
\item $\varphi_{251} :\equiv \delta \in \mathbb{R} \land 0 < \delta \land \varphi_{99} \to x \in \mathrm{dom}(f) \to y \in \mathrm{dom}(f) \to 0 < \lvert x - c \rvert \to \varphi_{52}$
\item $\varphi_{252} :\equiv r \in \mathbb{R} \land 0 < r \land \varphi_{16} \to \varphi_{76} \to \mathrm{Conv}(f \circ p,\alpha) \to \varepsilon \in \mathbb{R} \to 0 < \varepsilon \to 0 < \varepsilon \cdot 2^{-1}$
\item $\varphi_{253} :\equiv \mathrm{FLim}(f,c,\alpha) \to \varepsilon \in \mathbb{R} \to 0 < \varepsilon \to \delta \in \mathbb{R} \land 0 < \delta \land \varphi_{99} \to x \in \mathrm{dom}(f) \to \varphi_{37}$
\item $\varphi_{254} :\equiv \delta \in \mathbb{R} \land 0 < \delta \land \varphi_{99} \to x \in \mathrm{dom}(f) \to y \in \mathrm{dom}(f) \to 0 < \lvert x - c \rvert \to \varphi_{63}$
\item $\varphi_{255} :\equiv \mathrm{FLim}(f,c,\alpha) \to \varepsilon \in \mathbb{R} \to 0 < \varepsilon \to \delta \in \mathbb{R} \land 0 < \delta \land \varphi_{99} \to \varphi_{91}$
\item $\varphi_{256} :\equiv \mathrm{FLim}(f,c,\alpha) \to \varepsilon \in \mathbb{R} \to 0 < \varepsilon \to \delta \in \mathbb{R} \land 0 < \delta \land \varphi_{99} \to \varphi_{92}$
\item $\varphi_{257} :\equiv \forall r\,(c \in \mathbb{R} \to r \in \mathbb{R} \to 0 < r \to \varphi_{87}) \to c \in \mathbb{R} \to r \in \mathbb{R} \to 0 < r \to \varphi_{87}$
\item $\varphi_{258} :\equiv \mathrm{FLim}(f,c,\alpha) \to \varepsilon \in \mathbb{R} \to 0 < \varepsilon \to \delta \in \mathbb{R} \land 0 < \delta \land \varphi_{99} \to \varphi_{99}$
\item $\varphi_{259} :\equiv N_{0} \in \mathbb{N} \land \varphi_{77} \to \varphi_{67} \to \forall n \in \mathbb{N}\,(N_{1} \leq n \to \lvert f \circ p(n) - \alpha \rvert < \varepsilon \cdot 2^{-1})$
\item $\varphi_{260} :\equiv \forall c\,\forall r\,(c \in \mathbb{R} \to r \in \mathbb{R} \to 0 < r \to \varphi_{87}) \to \forall r\,(c \in \mathbb{R} \to r \in \mathbb{R} \to 0 < r \to \varphi_{87})$
\item $\varphi_{261} :\equiv \mathrm{RFn}(f) \land c \in \mathbb{R} \land \mathrm{PNbhd}(f,c) \land \forall \varepsilon \in \mathbb{R}\,(0 < \varepsilon \to \exists \delta\,(\delta \in \mathbb{R} \land 0 < \delta \land \forall x\,(\varphi_{161}))) \land \neg \exists \alpha\,(\mathrm{FLim}(f,c,\alpha))$
\item $\varphi_{262} :\equiv \mathrm{RFn}(f) \land c \in \mathbb{R} \land \mathrm{PNbhd}(f,c) \land \exists \alpha\,(\mathrm{FLim}(f,c,\alpha)) \land \neg \forall \varepsilon \in \mathbb{R}\,(0 < \varepsilon \to \exists \delta\,(\delta \in \mathbb{R} \land 0 < \delta \land \forall x\,(\varphi_{161})))$
\item $\varphi_{263} :\equiv \delta \in \mathbb{R} \land 0 < \delta \land \varphi_{99} \to x \in \mathrm{dom}(f) \to y \in \mathrm{dom}(f) \to 0 < \lvert x - c \rvert \to \varphi_{93}$
\item $\varphi_{264} :\equiv c \in \mathbb{R} \to \mathrm{PNbhd}(f,c) \to \exists \alpha\,(\mathrm{FLim}(f,c,\alpha)) \to \varphi_{218}$
\item $\varphi_{265} :\equiv \delta \in \mathbb{R} \land 0 < \delta \land \varphi_{99} \to x \in \mathrm{dom}(f) \to y \in \mathrm{dom}(f) \to 0 < \lvert x - c \rvert \to \varphi_{100}$
\item $\varphi_{266} :\equiv \delta \in \mathbb{R} \land 0 < \delta \land \varphi_{99} \to x \in \mathrm{dom}(f) \to y \in \mathrm{dom}(f) \to 0 < \lvert x - c \rvert \to \varphi_{105}$
\item $\varphi_{267} :\equiv \delta \in \mathbb{R} \land 0 < \delta \land \varphi_{99} \to x \in \mathrm{dom}(f) \to y \in \mathrm{dom}(f) \to 0 < \lvert x - c \rvert \to \varphi_{106}$
\item $\varphi_{268} :\equiv \delta \in \mathbb{R} \land 0 < \delta \land \varphi_{99} \to x \in \mathrm{dom}(f) \to y \in \mathrm{dom}(f) \to 0 < \lvert x - c \rvert \to \varphi_{107}$
\item $\varphi_{269} :\equiv \delta \in \mathbb{R} \land 0 < \delta \land \varphi_{99} \to x \in \mathrm{dom}(f) \to y \in \mathrm{dom}(f) \to 0 < \lvert x - c \rvert \to \varphi_{108}$
\item $\varphi_{270} :\equiv c \in \mathbb{R} \to \mathrm{PNbhd}(f,c) \to \exists \alpha\,(\mathrm{FLim}(f,c,\alpha)) \to \varphi_{222}$
\item $\varphi_{271} :\equiv N \in \mathbb{N} \land \forall n \in \mathbb{N}\,(N \leq n \to \lvert p(n) - c \rvert < \delta) \to m \in \mathbb{N} \to n \in \mathbb{N} \to N \leq m \to N \leq n \to \varphi_{199}$
\item $\varphi_{272} :\equiv c \in \mathbb{R} \to \mathrm{PNbhd}(f,c) \to \exists \alpha\,(\mathrm{FLim}(f,c,\alpha)) \to \varphi_{230}$
\item $\varphi_{273} :\equiv \forall \varepsilon\,(\varphi_{19} \to \varepsilon \in \mathbb{R} \to 0 < \varepsilon \to \exists \delta\,(\delta \in \mathbb{R} \land 0 < \delta \land \forall x\,(\varphi_{161})))$
\item $\varphi_{274} :\equiv \forall \delta\,(\varphi_{19} \land \varepsilon \in \mathbb{R} \land 0 < \varepsilon \land \neg \exists \delta\,(\delta \in \mathbb{R} \land 0 < \delta \land \forall x\,(\varphi_{161})))$
\item $\varphi_{275} :\equiv \mathrm{FLim}(f,c,\alpha) \to \varepsilon \in \mathbb{R} \to 0 < \varepsilon \to \delta \in \mathbb{R} \land 0 < \delta \land \varphi_{99} \to \delta \in \mathbb{R} \land 0 < \delta \land \varphi_{99}$
\item $\varphi_{276} :\equiv \forall \varepsilon\,(\varphi_{19}) \to \forall \varepsilon \in \mathbb{R}\,(0 < \varepsilon \to \exists \delta\,(\delta \in \mathbb{R} \land 0 < \delta \land \forall x\,(\varphi_{161})))$
\item $\varphi_{277} :\equiv \varphi_{148} \to \mathrm{Conv}(f \circ p,\alpha) \to \varepsilon \cdot 2^{-1} \in \mathbb{R} \to 0 < \varepsilon \cdot 2^{-1} \to \varphi_{70}$
\item $\varphi_{278} :\equiv \exists \alpha\,(\mathrm{FLim}(f,c,\alpha)) \to \mathrm{FLim}(f,c,\alpha) \to \varepsilon \in \mathbb{R} \to \varphi_{237}$
\item $\varphi_{279} :\equiv (\varepsilon \in \mathbb{R} \to 2^{-1} \in \mathbb{R} \to \varepsilon \cdot 2^{-1} \in \mathbb{R}) \to (\mathrm{RFn}(f) \to c \in \mathbb{R} \to \mathrm{PNbhd}(f,c) \to \varphi_{8}) \to \varphi_{151}$
\item $\varphi_{280} :\equiv c \in \mathbb{R} \to \mathrm{PNbhd}(f,c) \to \exists \alpha\,(\mathrm{FLim}(f,c,\alpha)) \to \varphi_{253}$
\item $\varphi_{281} :\equiv \varphi_{67} \to x \in \mathrm{dom}(f) \to 0 < \lvert x - c \rvert \to \lvert x - c \rvert < \delta \to \forall x\,(\varphi_{176})$
\item $\varphi_{282} :\equiv 0 < \varepsilon \to \delta \in \mathbb{R} \land 0 < \delta \land \forall x\,(\varphi_{161}) \to \exists N \in \mathbb{N}\,\forall m\,\forall n\,(\varphi_{90})$
\item $\varphi_{283} :\equiv c \in \mathbb{R} \to \mathrm{PNbhd}(f,c) \to \exists \alpha\,(\mathrm{FLim}(f,c,\alpha)) \to \varphi_{255}$
\item $\varphi_{284} :\equiv c \in \mathbb{R} \to \mathrm{PNbhd}(f,c) \to \exists \alpha\,(\mathrm{FLim}(f,c,\alpha)) \to \varphi_{256}$
\item $\varphi_{285} :\equiv \forall \varepsilon \in \mathbb{R}\,(0 < \varepsilon \to \exists \delta\,(\delta \in \mathbb{R} \land 0 < \delta \land \forall x\,(\varphi_{161}))) \to \exists r\,(r \in \mathbb{R} \land 0 < r \land \varphi_{16})$
\item $\varphi_{286} :\equiv c \in \mathbb{R} \to \mathrm{PNbhd}(f,c) \to \exists \alpha\,(\mathrm{FLim}(f,c,\alpha)) \to \varphi_{258}$
\item $\varphi_{287} :\equiv \delta \in \mathbb{R} \land 0 < \delta \land \forall x\,(\varphi_{161}) \to N \in \mathbb{N} \land \forall n \in \mathbb{N}\,(N \leq n \to \lvert p(n) - c \rvert < \delta) \to \varphi_{32}$
\item $\varphi_{288} :\equiv \delta \in \mathbb{R} \land 0 < \delta \land \forall x\,(\varphi_{161}) \to N \in \mathbb{N} \land \forall n \in \mathbb{N}\,(N \leq n \to \lvert p(n) - c \rvert < \delta) \to \varphi_{33}$
\item $\varphi_{289} :\equiv \delta \in \mathbb{R} \land 0 < \delta \land \forall x\,(\varphi_{161}) \to N \in \mathbb{N} \land \forall n \in \mathbb{N}\,(N \leq n \to \lvert p(n) - c \rvert < \delta) \to \varphi_{35}$
\item $\varphi_{290} :\equiv \delta \in \mathbb{R} \land 0 < \delta \land \forall x\,(\varphi_{161}) \to N \in \mathbb{N} \land \forall n \in \mathbb{N}\,(N \leq n \to \lvert p(n) - c \rvert < \delta) \to \varphi_{36}$
\item $\varphi_{291} :\equiv \mathrm{FLim}(f,c,\alpha) \to \varepsilon \in \mathbb{R} \to 0 < \varepsilon \to \delta \in \mathbb{R} \land 0 < \delta \land \varphi_{99} \to \varphi_{161}$
\item $\varphi_{292} :\equiv 0 < \varepsilon \to \delta \in \mathbb{R} \land 0 < \delta \land \forall x\,(\varphi_{176}) \to \exists \delta\,(\delta \in \mathbb{R} \land 0 < \delta \land \varphi_{69})$
\item $\varphi_{293} :\equiv \forall \varepsilon \in \mathbb{R}\,(0 < \varepsilon \to \exists \delta\,(\delta \in \mathbb{R} \land 0 < \delta \land \forall x\,(\varphi_{161}))) \to r \in \mathbb{R} \land 0 < r \land \varphi_{16} \to 0 < r$
\item $\varphi_{294} :\equiv \exists \alpha\,(\mathrm{FLim}(f,c,\alpha)) \to \mathrm{FLim}(f,c,\alpha) \to \varepsilon \in \mathbb{R} \to 0 < \varepsilon \to \varphi_{249}$
\item $\varphi_{295} :\equiv \exists \alpha\,(\mathrm{FLim}(f,c,\alpha)) \to \mathrm{FLim}(f,c,\alpha) \to \varepsilon \in \mathbb{R} \to 0 < \varepsilon \to \varphi_{250}$
\item $\varphi_{296} :\equiv \exists \alpha\,(\mathrm{FLim}(f,c,\alpha)) \to \mathrm{FLim}(f,c,\alpha) \to \varepsilon \in \mathbb{R} \to 0 < \varepsilon \to \varphi_{251}$
\item $\varphi_{297} :\equiv 0 < \varepsilon \to \delta \in \mathbb{R} \land 0 < \delta \land \varphi_{99} \to \exists \delta\,(\delta \in \mathbb{R} \land 0 < \delta \land \forall x\,(\varphi_{161}))$
\item $\varphi_{298} :\equiv \mathrm{FLim}(f,c,\alpha) \to \varepsilon \in \mathbb{R} \to 0 < \varepsilon \to \delta \in \mathbb{R} \land 0 < \delta \land \varphi_{99} \to \forall x\,(\varphi_{161})$
\item $\varphi_{299} :\equiv \mathrm{FLim}(f,c,\alpha) \to \varepsilon \in \mathbb{R} \to 0 < \varepsilon \to \delta \in \mathbb{R} \land 0 < \delta \land \varphi_{99} \to \varphi_{167}$
\item $\varphi_{300} :\equiv \mathrm{FLim}(f,c,\alpha) \to \varepsilon \in \mathbb{R} \to 0 < \varepsilon \to \delta \in \mathbb{R} \land 0 < \delta \land \varphi_{99} \to \varphi_{168}$
\item $\varphi_{301} :\equiv \forall \varepsilon \in \mathbb{R}\,(0 < \varepsilon \to \exists \delta\,(\delta \in \mathbb{R} \land 0 < \delta \land \forall x\,(\varphi_{161}))) \to r \in \mathbb{R} \land 0 < r \land \varphi_{16} \to r \in \mathbb{R}$
\item $\varphi_{302} :\equiv \exists \alpha\,(\mathrm{FLim}(f,c,\alpha)) \to \mathrm{FLim}(f,c,\alpha) \to \varepsilon \in \mathbb{R} \to 0 < \varepsilon \to \varphi_{254}$
\item $\varphi_{303} :\equiv \forall \varepsilon \in \mathbb{R}\,(0 < \varepsilon \to \exists \delta\,(\delta \in \mathbb{R} \land 0 < \delta \land \forall x\,(\varphi_{161}))) \to r \in \mathbb{R} \land 0 < r \land \varphi_{16} \to \exists \alpha\,(\mathrm{FLim}(f,c,\alpha))$
\item $\varphi_{304} :\equiv r \in \mathbb{R} \land 0 < r \land \varphi_{16} \to \varphi_{76} \to \varepsilon \in \mathbb{R} \to 0 < \varepsilon \to \exists N \in \mathbb{N}\,\forall m\,\forall n\,(\varphi_{90})$
\item $\varphi_{305} :\equiv \varepsilon \in \mathbb{R} \to 0 < \varepsilon \to \delta \in \mathbb{R} \land 0 < \delta \land \varphi_{99} \to \delta \in \mathbb{R} \land 0 < \delta \land \forall x\,(\varphi_{161})$
\item $\varphi_{306} :\equiv \exists \alpha\,(\mathrm{FLim}(f,c,\alpha)) \to \mathrm{FLim}(f,c,\alpha) \to \varepsilon \in \mathbb{R} \to 0 < \varepsilon \to \varphi_{263}$
\item $\varphi_{307} :\equiv r \in \mathbb{R} \land 0 < r \land \varphi_{16} \to \varphi_{76} \to \forall \varepsilon \in \mathbb{R}\,(0 < \varepsilon \to \exists N \in \mathbb{N}\,\forall m\,\forall n\,(\varphi_{90}))$
\item $\varphi_{308} :\equiv \exists \alpha\,(\mathrm{FLim}(f,c,\alpha)) \to \mathrm{FLim}(f,c,\alpha) \to \varepsilon \in \mathbb{R} \to 0 < \varepsilon \to \varphi_{265}$
\item $\varphi_{309} :\equiv \delta \in \mathbb{R} \land 0 < \delta \land \forall x\,(\varphi_{161}) \to N \in \mathbb{N} \land \forall n \in \mathbb{N}\,(N \leq n \to \lvert p(n) - c \rvert < \delta) \to \varphi_{90}$
\item $\varphi_{310} :\equiv c \in \mathbb{R} \to \mathrm{PNbhd}(f,c) \to \exists \alpha\,(\mathrm{FLim}(f,c,\alpha)) \to \varphi_{275}$
\item $\varphi_{311} :\equiv \exists \alpha\,(\mathrm{FLim}(f,c,\alpha)) \to \mathrm{FLim}(f,c,\alpha) \to \varepsilon \in \mathbb{R} \to 0 < \varepsilon \to \varphi_{266}$
\item $\varphi_{312} :\equiv \exists \alpha\,(\mathrm{FLim}(f,c,\alpha)) \to \mathrm{FLim}(f,c,\alpha) \to \varepsilon \in \mathbb{R} \to 0 < \varepsilon \to \varphi_{267}$
\item $\varphi_{313} :\equiv \exists \alpha\,(\mathrm{FLim}(f,c,\alpha)) \to \mathrm{FLim}(f,c,\alpha) \to \varepsilon \in \mathbb{R} \to 0 < \varepsilon \to \varphi_{268}$
\item $\varphi_{314} :\equiv \exists \alpha\,(\mathrm{FLim}(f,c,\alpha)) \to \mathrm{FLim}(f,c,\alpha) \to \varepsilon \in \mathbb{R} \to 0 < \varepsilon \to \varphi_{269}$
\item $\varphi_{315} :\equiv \mathrm{FLim}(f,c,\alpha) \to \varepsilon \in \mathbb{R} \to 0 < \varepsilon \to \delta \in \mathbb{R} \land 0 < \delta \land \varphi_{99} \to x \in \mathrm{dom}(f) \to \varphi_{173}$
\item $\varphi_{316} :\equiv (\mathrm{RFn}(f) \to c \in \mathbb{R} \to \mathrm{PNbhd}(f,c) \to \varphi_{57}) \to \mathrm{RFn}(f) \to \varphi_{213}$
\item $\varphi_{317} :\equiv \mathrm{Conv}(f \circ p,\alpha) \to \varepsilon \in \mathbb{R} \to 0 < \varepsilon \to \delta \in \mathbb{R} \land 0 < \delta \land \forall x\,(\varphi_{176}) \to \varphi_{114}$
\item $\varphi_{318} :\equiv \forall \delta\,(\neg (\delta \in \mathbb{R} \land 0 < \delta \land \forall x\,(\varphi_{161}))) \to \neg (\delta \in \mathbb{R} \land 0 < \delta \land \forall x\,(\varphi_{161}))$
\item $\varphi_{319} :\equiv c \in \mathbb{R} \to \mathrm{PNbhd}(f,c) \to \exists \alpha\,(\mathrm{FLim}(f,c,\alpha)) \to \varphi_{291}$
\item $\varphi_{320} :\equiv \mathrm{RFn}(f) \land c \in \mathbb{R} \land \mathrm{PNbhd}(f,c) \land \forall \varepsilon \in \mathbb{R}\,(0 < \varepsilon \to \exists \delta\,(\delta \in \mathbb{R} \land 0 < \delta \land \forall x\,(\varphi_{161}))) \land (r \in \mathbb{R} \land 0 < r \land \varphi_{16}) \land \neg \exists \alpha\,(\mathrm{FLim}(f,c,\alpha))$
\item $\varphi_{321} :\equiv c \in \mathbb{R} \to \mathrm{PNbhd}(f,c) \to \exists \alpha\,(\mathrm{FLim}(f,c,\alpha)) \to \varphi_{298}$
\item $\varphi_{322} :\equiv c \in \mathbb{R} \to \mathrm{PNbhd}(f,c) \to \exists \alpha\,(\mathrm{FLim}(f,c,\alpha)) \to \varphi_{299}$
\item $\varphi_{323} :\equiv c \in \mathbb{R} \to \mathrm{PNbhd}(f,c) \to \exists \alpha\,(\mathrm{FLim}(f,c,\alpha)) \to \varphi_{300}$
\item $\varphi_{324} :\equiv 0 < \varepsilon \to \delta \in \mathbb{R} \land 0 < \delta \land \forall x\,(\varphi_{161}) \to \delta \in \mathbb{R} \land 0 < \delta \land \forall x\,(\varphi_{161})$
\item $\varphi_{325} :\equiv \exists \alpha\,(\mathrm{FLim}(f,c,\alpha)) \to \mathrm{FLim}(f,c,\alpha) \to \varepsilon \in \mathbb{R} \to \varphi_{297}$
\item $\varphi_{326} :\equiv \varphi_{19} \land \varepsilon \in \mathbb{R} \land 0 < \varepsilon \land (\delta \in \mathbb{R} \land 0 < \delta \land \varphi_{99}) \to x \in \mathrm{dom}(f) \to y \in \mathrm{dom}(f) \to 0 < \lvert x - c \rvert \to \varphi_{93}$
\item $\varphi_{327} :\equiv 2^{-1} \in \mathbb{R} \to \mathrm{RFn}(f) \to c \in \mathbb{R} \to \mathrm{PNbhd}(f,c) \to \varphi_{302}$
\item $\varphi_{328} :\equiv \mathrm{PNbhd}(f,c) \to \exists \alpha\,(\mathrm{FLim}(f,c,\alpha)) \to \mathrm{FLim}(f,c,\alpha) \to \varphi_{305}$
\item $\varphi_{329} :\equiv \forall x\,(\varphi_{19} \land \varepsilon \in \mathbb{R} \land 0 < \varepsilon \land (\delta \in \mathbb{R} \land 0 < \delta \land \varphi_{99})) \to \forall x\,(\varphi_{161})$
\item $\varphi_{330} :\equiv (\varepsilon \in \mathbb{R} \to 0 < \varepsilon \to 0 < \varepsilon \cdot 2^{-1}) \to (\mathrm{RFn}(f) \to c \in \mathbb{R} \to \mathrm{PNbhd}(f,c) \to \varphi_{8}) \to \varphi_{219}$
\item $\varphi_{331} :\equiv c \in \mathbb{R} \to \mathrm{PNbhd}(f,c) \to \exists \alpha\,(\mathrm{FLim}(f,c,\alpha)) \to \varphi_{315}$
\item $\varphi_{332} :\equiv 0 < \varepsilon \to \delta \in \mathbb{R} \land 0 < \delta \land \forall x\,(\varphi_{176}) \to \delta \in \mathbb{R} \land 0 < \delta \land \forall x\,(\varphi_{176})$
\item $\varphi_{333} :\equiv r \in \mathbb{R} \land 0 < r \land \varphi_{16} \to \varphi_{76} \to \varepsilon \in \mathbb{R} \to 0 < \varepsilon \to \exists \delta\,(\delta \in \mathbb{R} \land 0 < \delta \land \forall x\,(\varphi_{161}))$
\item $\varphi_{334} :\equiv r \in \mathbb{R} \land 0 < r \land \varphi_{16} \to \varphi_{76} \to \varepsilon \in \mathbb{R} \to 0 < \varepsilon \to \delta \in \mathbb{R} \land 0 < \delta \land \forall x\,(\varphi_{161}) \to 0 < \delta$
\item $\varphi_{335} :\equiv \delta \in \mathbb{R} \land 0 < \delta \land \forall x\,(\varphi_{176}) \to N_{0} \in \mathbb{N} \land \varphi_{77} \to \varphi_{67} \to N_{0} \cup N_{1} \in \mathbb{N}$
\item $\varphi_{336} :\equiv \forall \varepsilon \in \mathbb{R}\,(0 < \varepsilon \to \exists \delta\,(\delta \in \mathbb{R} \land 0 < \delta \land \forall x\,(\varphi_{161}))) \to r \in \mathbb{R} \land 0 < r \land \varphi_{16} \to r \in \mathbb{R} \land 0 < r \land \varphi_{16}$
\item $\varphi_{337} :\equiv \forall \varepsilon \in \mathbb{R}\,(0 < \varepsilon \to \exists \delta\,(\delta \in \mathbb{R} \land 0 < \delta \land \forall x\,(\varphi_{161}))) \to r \in \mathbb{R} \land 0 < r \land \varphi_{16} \to \varphi_{87}$
\item $\varphi_{338} :\equiv r \in \mathbb{R} \land 0 < r \land \varphi_{16} \to \varphi_{76} \to \varepsilon \in \mathbb{R} \to 0 < \varepsilon \to \delta \in \mathbb{R} \land 0 < \delta \land \forall x\,(\varphi_{161}) \to \delta \in \mathbb{R}$
\item $\varphi_{339} :\equiv 0 < \varepsilon \to \delta \in \mathbb{R} \land 0 < \delta \land \forall x\,(\varphi_{176}) \to N_{0} \in \mathbb{N} \land \varphi_{77} \to \varphi_{67} \to N_{1} \in \mathbb{N}$
\item $\varphi_{340} :\equiv 0 < \varepsilon \to \delta \in \mathbb{R} \land 0 < \delta \land \forall x\,(\varphi_{176}) \to N_{0} \in \mathbb{N} \land \varphi_{77} \to \varphi_{67} \to N_{0} \in \mathbb{N}$
\item $\varphi_{341} :\equiv \varphi_{19} \land \varepsilon \in \mathbb{R} \land 0 < \varepsilon \land \neg \exists \delta\,(\delta \in \mathbb{R} \land 0 < \delta \land \forall x\,(\varphi_{161})) \to \neg (\delta \in \mathbb{R} \land 0 < \delta \land \varphi_{99})$
\item $\varphi_{342} :\equiv 0 < \varepsilon \to \delta \in \mathbb{R} \land 0 < \delta \land \forall x\,(\varphi_{176}) \to N_{0} \in \mathbb{N} \land \varphi_{77} \to N_{0} \in \mathbb{N} \land \varphi_{77}$
\item $\varphi_{343} :\equiv \forall \varepsilon \in \mathbb{R}\,(0 < \varepsilon \to \exists \delta\,(\delta \in \mathbb{R} \land 0 < \delta \land \forall x\,(\varphi_{161}))) \to r \in \mathbb{R} \land 0 < r \land \varphi_{16} \to \varphi_{76} \to \mathrm{Seq}(p)$
\item $\varphi_{344} :\equiv N_{0} \in \mathbb{N} \land \varphi_{77} \to \varphi_{67} \to x \in \mathrm{dom}(f) \to 0 < \lvert x - c \rvert \to \lvert x - c \rvert < \delta \to \varphi_{176}$
\item $\varphi_{345} :\equiv \varepsilon \in \mathbb{R} \to 0 < \varepsilon \to \delta \in \mathbb{R} \land 0 < \delta \land \forall x\,(\varphi_{176}) \to N_{0} \in \mathbb{N} \land \varphi_{77} \to \varphi_{70}$
\item $\varphi_{346} :\equiv \delta \in \mathbb{R} \land 0 < \delta \land \forall x\,(\varphi_{176}) \to N_{0} \in \mathbb{N} \land \varphi_{77} \to \varphi_{67} \to p(N_{0} \cup N_{1}) \in \mathbb{R}$
\item $\varphi_{347} :\equiv \delta \in \mathbb{R} \land 0 < \delta \land \forall x\,(\varphi_{176}) \to N_{0} \in \mathbb{N} \land \varphi_{77} \to \exists \delta\,(\delta \in \mathbb{R} \land 0 < \delta \land \varphi_{69})$
\item $\varphi_{348} :\equiv \delta \in \mathbb{R} \land 0 < \delta \land \forall x\,(\varphi_{176}) \to N_{0} \in \mathbb{N} \land \varphi_{77} \to \varphi_{67} \to x \in \mathrm{dom}(f) \to x \in \mathrm{dom}(f)$
\item $\varphi_{349} :\equiv \forall \varepsilon \in \mathbb{R}\,(0 < \varepsilon \to \exists \delta\,(\delta \in \mathbb{R} \land 0 < \delta \land \forall x\,(\varphi_{161}))) \to r \in \mathbb{R} \land 0 < r \land \varphi_{16} \to \varphi_{76} \to \mathrm{Conv}(p,c)$
\item $\varphi_{350} :\equiv \delta \in \mathbb{R} \land 0 < \delta \land \forall x\,(\varphi_{176}) \to N_{0} \in \mathbb{N} \land \varphi_{77} \to \varphi_{67} \to p(N_{0} \cup N_{1}) \in \mathrm{dom}(f)$
\item $\varphi_{351} :\equiv \forall \varepsilon \in \mathbb{R}\,(0 < \varepsilon \to \exists \delta\,(\delta \in \mathbb{R} \land 0 < \delta \land \forall x\,(\varphi_{161}))) \to r \in \mathbb{R} \land 0 < r \land \varphi_{16} \to \varphi_{76} \to \exists \alpha\,(\mathrm{FLim}(f,c,\alpha))$
\item $\varphi_{352} :\equiv \forall \varepsilon \in \mathbb{R}\,(0 < \varepsilon \to \exists \delta\,(\delta \in \mathbb{R} \land 0 < \delta \land \forall x\,(\varphi_{161}))) \to r \in \mathbb{R} \land 0 < r \land \varphi_{16} \to \varphi_{76} \to \mathrm{Seq}(f \circ p)$
\item $\varphi_{353} :\equiv \forall \varepsilon \in \mathbb{R}\,(0 < \varepsilon \to \exists \delta\,(\delta \in \mathbb{R} \land 0 < \delta \land \forall x\,(\varphi_{161}))) \to r \in \mathbb{R} \land 0 < r \land \varphi_{16} \to \varphi_{76} \to \mathrm{Cauchy}(f \circ p)$
\item $\varphi_{354} :\equiv \forall \varepsilon \in \mathbb{R}\,(0 < \varepsilon \to \exists \delta\,(\delta \in \mathbb{R} \land 0 < \delta \land \forall x\,(\varphi_{161}))) \to \varepsilon \in \mathbb{R} \to 0 < \varepsilon \to \exists \delta\,(\delta \in \mathbb{R} \land 0 < \delta \land \forall x\,(\varphi_{161}))$
\item $\varphi_{355} :\equiv \forall \varepsilon \in \mathbb{R}\,(0 < \varepsilon \to \exists \delta\,(\delta \in \mathbb{R} \land 0 < \delta \land \forall x\,(\varphi_{161}))) \to r \in \mathbb{R} \land 0 < r \land \varphi_{16} \to \varphi_{76} \to n \in \mathbb{N} \to n \in \mathbb{N}$
\item $\varphi_{356} :\equiv \forall \varepsilon \in \mathbb{R}\,(0 < \varepsilon \to \exists \delta\,(\delta \in \mathbb{R} \land 0 < \delta \land \forall x\,(\varphi_{161}))) \to r \in \mathbb{R} \land 0 < r \land \varphi_{16} \to \varphi_{76} \to \exists \alpha\,(\mathrm{Conv}(f \circ p,\alpha))$
\item $\varphi_{357} :\equiv \forall \varepsilon \in \mathbb{R}\,(0 < \varepsilon \to \exists \delta\,(\delta \in \mathbb{R} \land 0 < \delta \land \forall x\,(\varphi_{161}))) \to \forall \varepsilon \in \mathbb{R}\,(0 < \varepsilon \to \exists \delta\,(\delta \in \mathbb{R} \land 0 < \delta \land \forall x\,(\varphi_{161})))$
\item $\varphi_{358} :\equiv \forall \varepsilon \in \mathbb{R}\,(0 < \varepsilon \to \exists \delta\,(\delta \in \mathbb{R} \land 0 < \delta \land \forall x\,(\varphi_{161}))) \to r \in \mathbb{R} \land 0 < r \land \varphi_{16} \to \varphi_{76} \to \varepsilon \in \mathbb{R} \to \varepsilon \in \mathbb{R}$
\item $\varphi_{359} :\equiv \forall \varepsilon \in \mathbb{R}\,(0 < \varepsilon \to \exists \delta\,(\delta \in \mathbb{R} \land 0 < \delta \land \forall x\,(\varphi_{161}))) \to r \in \mathbb{R} \land 0 < r \land \varphi_{16} \to \varphi_{76} \to n \in \mathbb{N} \to p(n) \in \mathbb{R}$
\item $\varphi_{360} :\equiv \forall N\,(\forall N_{0}\,(\neg (N_{0} \in \mathbb{N} \land \varphi_{77})) \to \varphi_{110}) \to \forall N\,\forall N_{0}\,(\neg (N_{0} \in \mathbb{N} \land \varphi_{77})) \to \forall N\,(\varphi_{110})$
\item $\varphi_{361} :\equiv \forall \varepsilon \in \mathbb{R}\,(0 < \varepsilon \to \exists \delta\,(\delta \in \mathbb{R} \land 0 < \delta \land \forall x\,(\varphi_{161}))) \to r \in \mathbb{R} \land 0 < r \land \varphi_{16} \to \varphi_{76} \to n \in \mathbb{N} \to p(n) \in \mathrm{dom}(f)$
\item $\varphi_{362} :\equiv \mathrm{RFn}(f) \land c \in \mathbb{R} \land \mathrm{PNbhd}(f,c) \land \forall \varepsilon \in \mathbb{R}\,(0 < \varepsilon \to \exists \delta\,(\delta \in \mathbb{R} \land 0 < \delta \land \forall x\,(\varphi_{161}))) \land (r \in \mathbb{R} \land 0 < r \land \varphi_{16}) \land \varphi_{76}$
\item $\varphi_{363} :\equiv \forall \varepsilon \in \mathbb{R}\,(0 < \varepsilon \to \exists \delta\,(\delta \in \mathbb{R} \land 0 < \delta \land \forall x\,(\varphi_{161}))) \to \varphi_{215}$
\item $\varphi_{364} :\equiv \forall \varepsilon \in \mathbb{R}\,(0 < \varepsilon \to \exists \delta\,(\delta \in \mathbb{R} \land 0 < \delta \land \forall x\,(\varphi_{161}))) \to r \in \mathbb{R} \land 0 < r \land \varphi_{16} \to \varphi_{76} \to \varepsilon \in \mathbb{R} \to 0 < \varepsilon \to 0 < \varepsilon$
\item $\varphi_{365} :\equiv \forall \varepsilon \in \mathbb{R}\,(0 < \varepsilon \to \exists \delta\,(\delta \in \mathbb{R} \land 0 < \delta \land \forall x\,(\varphi_{161}))) \to r \in \mathbb{R} \land 0 < r \land \varphi_{16} \to \varphi_{76} \to \mathrm{Conv}(f \circ p,\alpha) \to \alpha \in \mathbb{R}$
\item $\varphi_{366} :\equiv \forall \varepsilon \in \mathbb{R}\,(0 < \varepsilon \to \exists \delta\,(\delta \in \mathbb{R} \land 0 < \delta \land \forall x\,(\varphi_{161}))) \to \varphi_{217}$
\item $\varphi_{367} :\equiv \delta \in \mathbb{R} \land 0 < \delta \land \forall x\,(\varphi_{176}) \to N_{0} \in \mathbb{N} \land \varphi_{77} \to \varphi_{67} \to \lvert f \circ p(N_{0} \cup N_{1}) - \alpha \rvert < \varepsilon \cdot 2^{-1}$
\item $\varphi_{368} :\equiv \forall \varepsilon \in \mathbb{R}\,(0 < \varepsilon \to \exists \delta\,(\delta \in \mathbb{R} \land 0 < \delta \land \forall x\,(\varphi_{161}))) \to \varphi_{220}$
\item $\varphi_{369} :\equiv \varphi_{19} \land \varepsilon \in \mathbb{R} \land 0 < \varepsilon \land (\delta \in \mathbb{R} \land 0 < \delta \land \varphi_{99}) \to \forall x\,(\varphi_{19} \land \varepsilon \in \mathbb{R} \land 0 < \varepsilon \land (\delta \in \mathbb{R} \land 0 < \delta \land \varphi_{99}))$
\item $\varphi_{370} :\equiv \varphi_{19} \land \varepsilon \in \mathbb{R} \land 0 < \varepsilon \land (\delta \in \mathbb{R} \land 0 < \delta \land \varphi_{99}) \to \forall y\,(\varphi_{19} \land \varepsilon \in \mathbb{R} \land 0 < \varepsilon \land (\delta \in \mathbb{R} \land 0 < \delta \land \varphi_{99}))$
\item $\varphi_{371} :\equiv \forall \varepsilon \in \mathbb{R}\,(0 < \varepsilon \to \exists \delta\,(\delta \in \mathbb{R} \land 0 < \delta \land \forall x\,(\varphi_{161}))) \to \varphi_{226}$
\item $\varphi_{372} :\equiv \forall \varepsilon \in \mathbb{R}\,(0 < \varepsilon \to \exists \delta\,(\delta \in \mathbb{R} \land 0 < \delta \land \forall x\,(\varphi_{161}))) \to \varphi_{232}$
\item $\varphi_{373} :\equiv \forall \varepsilon \in \mathbb{R}\,(0 < \varepsilon \to \exists \delta\,(\delta \in \mathbb{R} \land 0 < \delta \land \forall x\,(\varphi_{161}))) \to \varphi_{231}$
\item $\varphi_{374} :\equiv (\mathrm{PNbhd}(f,c) \leftrightarrow \exists r\,(r \in \mathbb{R} \land 0 < r \land \varphi_{16})) \to \mathrm{RFn}(f) \to c \in \mathbb{R} \to \mathrm{PNbhd}(f,c) \to \varphi_{285}$
\item $\varphi_{375} :\equiv (\mathrm{RFn}(f) \to c \in \mathbb{R} \to \mathrm{PNbhd}(f,c) \to \varphi_{3}) \to \mathrm{RFn}(f) \to \varphi_{322}$
\item $\varphi_{376} :\equiv (\mathrm{RFn}(f) \to c \in \mathbb{R} \to \mathrm{PNbhd}(f,c) \to \varphi_{3}) \to \mathrm{RFn}(f) \to \varphi_{323}$
\item $\varphi_{377} :\equiv \forall \varepsilon \in \mathbb{R}\,(0 < \varepsilon \to \exists \delta\,(\delta \in \mathbb{R} \land 0 < \delta \land \forall x\,(\varphi_{161}))) \to \varphi_{240}$
\item $\varphi_{378} :\equiv r \in \mathbb{R} \land 0 < r \land \varphi_{16} \to \varphi_{76} \to \varepsilon \in \mathbb{R} \to 0 < \varepsilon \to \delta \in \mathbb{R} \land 0 < \delta \land \forall x\,(\varphi_{161}) \to \varphi_{24}$
\item $\varphi_{379} :\equiv \forall \varepsilon \in \mathbb{R}\,(0 < \varepsilon \to \exists \delta\,(\delta \in \mathbb{R} \land 0 < \delta \land \forall x\,(\varphi_{161}))) \to \varphi_{243}$
\item $\varphi_{380} :\equiv \forall \varepsilon \in \mathbb{R}\,(0 < \varepsilon \to \exists \delta\,(\delta \in \mathbb{R} \land 0 < \delta \land \forall x\,(\varphi_{161}))) \to \varphi_{245}$
\item $\varphi_{381} :\equiv \forall \varepsilon \in \mathbb{R}\,(0 < \varepsilon \to \exists \delta\,(\delta \in \mathbb{R} \land 0 < \delta \land \forall x\,(\varphi_{161}))) \to \varphi_{247}$
\item $\varphi_{382} :\equiv \forall \varepsilon \in \mathbb{R}\,(0 < \varepsilon \to \exists \delta\,(\delta \in \mathbb{R} \land 0 < \delta \land \forall x\,(\varphi_{161}))) \to \varphi_{252}$
\item $\varphi_{383} :\equiv (\mathrm{RFn}(f) \to c \in \mathbb{R} \to \mathrm{PNbhd}(f,c) \to \varphi_{8}) \to \mathrm{RFn}(f) \to c \in \mathbb{R} \to \mathrm{PNbhd}(f,c) \to \varphi_{308}$
\item $\varphi_{384} :\equiv \forall \varepsilon \in \mathbb{R}\,(0 < \varepsilon \to \exists \delta\,(\delta \in \mathbb{R} \land 0 < \delta \land \forall x\,(\varphi_{161}))) \to r \in \mathbb{R} \land 0 < r \land \varphi_{16} \to \varphi_{180}$
\item $\varphi_{385} :\equiv \forall \varepsilon \in \mathbb{R}\,(0 < \varepsilon \to \exists \delta\,(\delta \in \mathbb{R} \land 0 < \delta \land \forall x\,(\varphi_{161}))) \to r \in \mathbb{R} \land 0 < r \land \varphi_{16} \to \varphi_{184}$
\item $\varphi_{386} :\equiv r \in \mathbb{R} \land 0 < r \land \varphi_{16} \to \varphi_{76} \to \varepsilon \in \mathbb{R} \to 0 < \varepsilon \to \delta \in \mathbb{R} \land 0 < \delta \land \forall x\,(\varphi_{161}) \to \varphi_{62}$
\item $\varphi_{387} :\equiv \varepsilon \in \mathbb{R} \to 0 < \varepsilon \to \delta \in \mathbb{R} \land 0 < \delta \land \forall x\,(\varphi_{176}) \to N_{0} \in \mathbb{N} \land \varphi_{77} \to \varphi_{67} \to \varphi_{16}$
\item $\varphi_{388} :\equiv \varepsilon \in \mathbb{R} \to 0 < \varepsilon \to \delta \in \mathbb{R} \land 0 < \delta \land \forall x\,(\varphi_{176}) \to N_{0} \in \mathbb{N} \land \varphi_{77} \to \varphi_{172}$
\item $\varphi_{389} :\equiv \varphi_{362} \land \neg \exists \alpha\,(\mathrm{FLim}(f,c,\alpha)) \to \neg \mathrm{Conv}(f \circ p,\alpha)$
\item $\varphi_{390} :\equiv \forall \varepsilon \in \mathbb{R}\,(0 < \varepsilon \to \exists \delta\,(\delta \in \mathbb{R} \land 0 < \delta \land \forall x\,(\varphi_{161}))) \to r \in \mathbb{R} \land 0 < r \land \varphi_{16} \to \varphi_{76} \to n \in \mathbb{N} \to \varphi_{16}$
\item $\varphi_{391} :\equiv \delta \in \mathbb{R} \land 0 < \delta \land \forall x\,(\varphi_{161}) \to N \in \mathbb{N} \land \forall n \in \mathbb{N}\,(N \leq n \to \lvert p(n) - c \rvert < \delta) \to \varphi_{233}$
\item $\varphi_{392} :\equiv \varepsilon \in \mathbb{R} \to 0 < \varepsilon \to \delta \in \mathbb{R} \land 0 < \delta \land \forall x\,(\varphi_{176}) \to N_{0} \in \mathbb{N} \land \varphi_{77} \to \varphi_{67} \to \varphi_{28}$
\item $\varphi_{393} :\equiv \varepsilon \in \mathbb{R} \to 0 < \varepsilon \to \delta \in \mathbb{R} \land 0 < \delta \land \forall x\,(\varphi_{176}) \to N_{0} \in \mathbb{N} \land \varphi_{77} \to \varphi_{67} \to \varphi_{29}$
\item $\varphi_{394} :\equiv \mathrm{Conv}(f \circ p,\alpha) \to \varepsilon \in \mathbb{R} \to 0 < \varepsilon \to \delta \in \mathbb{R} \land 0 < \delta \land \forall x\,(\varphi_{176}) \to \varphi_{238}$
\item $\varphi_{395} :\equiv \varepsilon \in \mathbb{R} \to 0 < \varepsilon \to \delta \in \mathbb{R} \land 0 < \delta \land \forall x\,(\varphi_{176}) \to N_{0} \in \mathbb{N} \land \varphi_{77} \to \varphi_{67} \to \varphi_{46}$
\item $\varphi_{396} :\equiv \mathrm{Conv}(f \circ p,\alpha) \to \varepsilon \in \mathbb{R} \to 0 < \varepsilon \to \delta \in \mathbb{R} \land 0 < \delta \land \forall x\,(\varphi_{176}) \to \varphi_{239}$
\item $\varphi_{397} :\equiv \mathrm{Conv}(f \circ p,\alpha) \to \varepsilon \in \mathbb{R} \to 0 < \varepsilon \to \delta \in \mathbb{R} \land 0 < \delta \land \forall x\,(\varphi_{176}) \to \varphi_{241}$
\item $\varphi_{398} :\equiv r \in \mathbb{R} \land 0 < r \land \varphi_{16} \to \varphi_{76} \to \varepsilon \in \mathbb{R} \to 0 < \varepsilon \to \delta \in \mathbb{R} \land 0 < \delta \land \forall x\,(\varphi_{161}) \to \varphi_{97}$
\item $\varphi_{399} :\equiv r \in \mathbb{R} \land 0 < r \land \varphi_{16} \to \varphi_{76} \to \varepsilon \in \mathbb{R} \to 0 < \varepsilon \to \delta \in \mathbb{R} \land 0 < \delta \land \forall x\,(\varphi_{161}) \to \varphi_{98}$
\item $\varphi_{400} :\equiv \forall \alpha\,(\varphi_{362} \land \neg \exists \alpha\,(\mathrm{FLim}(f,c,\alpha))) \to \forall \alpha\,(\neg \mathrm{Conv}(f \circ p,\alpha))$
\item $\varphi_{401} :\equiv \forall \varepsilon \in \mathbb{R}\,(0 < \varepsilon \to \exists \delta\,(\delta \in \mathbb{R} \land 0 < \delta \land \forall x\,(\varphi_{161}))) \to r \in \mathbb{R} \land 0 < r \land \varphi_{16} \to \varphi_{76} \to \varphi_{76}$
\item $\varphi_{402} :\equiv \mathrm{Conv}(f \circ p,\alpha) \to \varepsilon \in \mathbb{R} \to 0 < \varepsilon \to \delta \in \mathbb{R} \land 0 < \delta \land \forall x\,(\varphi_{176}) \to \varphi_{242}$
\item $\varphi_{403} :\equiv \varphi_{211} \to f(x) \in \mathbb{R} \to f(p(N_{0} \cup N_{1})) \in \mathbb{R} \to \alpha \in \mathbb{R} \to \varphi_{208}$
\item $\varphi_{404} :\equiv \forall \varepsilon \in \mathbb{R}\,(0 < \varepsilon \to \exists \delta\,(\delta \in \mathbb{R} \land 0 < \delta \land \forall x\,(\varphi_{161}))) \to r \in \mathbb{R} \land 0 < r \land \varphi_{16} \to \varphi_{76} \to \varphi_{85}$
\item $\varphi_{405} :\equiv \varepsilon \in \mathbb{R} \to 0 < \varepsilon \to \delta \in \mathbb{R} \land 0 < \delta \land \forall x\,(\varphi_{176}) \to N_{0} \in \mathbb{N} \land \varphi_{77} \to \varphi_{67} \to \varphi_{59}$
\item $\varphi_{406} :\equiv \varepsilon \in \mathbb{R} \to 0 < \varepsilon \to \delta \in \mathbb{R} \land 0 < \delta \land \forall x\,(\varphi_{176}) \to N_{0} \in \mathbb{N} \land \varphi_{77} \to \varphi_{67} \to \varphi_{61}$
\item $\varphi_{407} :\equiv \delta \in \mathbb{R} \land 0 < \delta \land \forall x\,(\varphi_{176}) \to N_{0} \in \mathbb{N} \land \varphi_{77} \to \varphi_{67} \to x \in \mathrm{dom}(f) \to \varphi_{95}$
\item $\varphi_{408} :\equiv \delta \in \mathbb{R} \land 0 < \delta \land \forall x\,(\varphi_{176}) \to N_{0} \in \mathbb{N} \land \varphi_{77} \to \varphi_{67} \to \delta \in \mathbb{R} \land 0 < \delta \land \varphi_{69}$
\item $\varphi_{409} :\equiv \varepsilon \in \mathbb{R} \to 0 < \varepsilon \to \delta \in \mathbb{R} \land 0 < \delta \land \forall x\,(\varphi_{176}) \to N_{0} \in \mathbb{N} \land \varphi_{77} \to \varphi_{67} \to \varphi_{67}$
\item $\varphi_{410} :\equiv r \in \mathbb{R} \land 0 < r \land \varphi_{16} \to \varphi_{76} \to \varepsilon \in \mathbb{R} \to 0 < \varepsilon \to \delta \in \mathbb{R} \land 0 < \delta \land \forall x\,(\varphi_{161}) \to \varphi_{116}$
\item $\varphi_{411} :\equiv \varepsilon \in \mathbb{R} \to 0 < \varepsilon \to \delta \in \mathbb{R} \land 0 < \delta \land \forall x\,(\varphi_{176}) \to N_{0} \in \mathbb{N} \land \varphi_{77} \to \varphi_{67} \to \varphi_{69}$
\item $\varphi_{412} :\equiv \varepsilon \in \mathbb{R} \to 0 < \varepsilon \to \delta \in \mathbb{R} \land 0 < \delta \land \forall x\,(\varphi_{176}) \to N_{0} \in \mathbb{N} \land \varphi_{77} \to \varphi_{67} \to \varphi_{77}$
\item $\varphi_{413} :\equiv \delta \in \mathbb{R} \land 0 < \delta \land \forall x\,(\varphi_{176}) \to N_{0} \in \mathbb{N} \land \varphi_{77} \to \varphi_{67} \to \exists \delta\,(\delta \in \mathbb{R} \land 0 < \delta \land \varphi_{69})$
\item $\varphi_{414} :\equiv \delta \in \mathbb{R} \land 0 < \delta \land \forall x\,(\varphi_{176}) \to N_{0} \in \mathbb{N} \land \varphi_{77} \to \varphi_{67} \to x \in \mathrm{dom}(f) \to 0 < \lvert x - c \rvert \to \varphi_{54}$
\item $\varphi_{415} :\equiv \varepsilon \in \mathbb{R} \to 0 < \varepsilon \to \delta \in \mathbb{R} \land 0 < \delta \land \forall x\,(\varphi_{176}) \to N_{0} \in \mathbb{N} \land \varphi_{77} \to \varphi_{67} \to \varphi_{79}$
\item $\varphi_{416} :\equiv \varepsilon \in \mathbb{R} \to 0 < \varepsilon \to \delta \in \mathbb{R} \land 0 < \delta \land \forall x\,(\varphi_{176}) \to N_{0} \in \mathbb{N} \land \varphi_{77} \to \varphi_{67} \to \varphi_{80}$
\item $\varphi_{417} :\equiv \varepsilon \in \mathbb{R} \to 0 < \varepsilon \to \delta \in \mathbb{R} \land 0 < \delta \land \forall x\,(\varphi_{176}) \to N_{0} \in \mathbb{N} \land \varphi_{77} \to \varphi_{67} \to \varphi_{86}$
\item $\varphi_{418} :\equiv \delta \in \mathbb{R} \land 0 < \delta \land \forall x\,(\varphi_{176}) \to N_{0} \in \mathbb{N} \land \varphi_{77} \to \varphi_{67} \to N_{0} \cup N_{1} \in \mathbb{N} \to N_{0} \leq N_{0} \cup N_{1} \to \varphi_{61}$
\item $\varphi_{419} :\equiv \mathrm{Conv}(f \circ p,\alpha) \to \varepsilon \in \mathbb{R} \to 0 < \varepsilon \to \delta \in \mathbb{R} \land 0 < \delta \land \forall x\,(\varphi_{176}) \to \varphi_{259}$
\item $\varphi_{420} :\equiv \varepsilon \in \mathbb{R} \to 0 < \varepsilon \to \delta \in \mathbb{R} \land 0 < \delta \land \forall x\,(\varphi_{176}) \to N_{0} \in \mathbb{N} \land \varphi_{77} \to \varphi_{67} \to \varphi_{96}$
\item $\varphi_{421} :\equiv (\varepsilon \cdot 2^{-1} + \varepsilon \cdot 2^{-1} = \varepsilon \to \varphi_{104}) \to \mathrm{RFn}(f) \to c \in \mathbb{R} \to \mathrm{PNbhd}(f,c) \to \varphi_{306}$
\item $\varphi_{422} :\equiv \forall \varepsilon \in \mathbb{R}\,(0 < \varepsilon \to \exists \delta\,(\delta \in \mathbb{R} \land 0 < \delta \land \forall x\,(\varphi_{161}))) \to r \in \mathbb{R} \land 0 < r \land \varphi_{16} \to \varphi_{76} \to n \in \mathbb{N} \to p(n) \in \mathbb{R} \to \varphi_{27}$
\item $\varphi_{423} :\equiv r \in \mathbb{R} \land 0 < r \land \varphi_{16} \to \varphi_{76} \to \varepsilon \in \mathbb{R} \to 0 < \varepsilon \to \delta \in \mathbb{R} \land 0 < \delta \land \forall x\,(\varphi_{161}) \to \varphi_{132}$
\item $\varphi_{424} :\equiv r \in \mathbb{R} \land 0 < r \land \varphi_{16} \to \varphi_{76} \to \varepsilon \in \mathbb{R} \to 0 < \varepsilon \to \delta \in \mathbb{R} \land 0 < \delta \land \forall x\,(\varphi_{161}) \to \varphi_{133}$
\item $\varphi_{425} :\equiv r \in \mathbb{R} \land 0 < r \land \varphi_{16} \to \varphi_{76} \to \varepsilon \in \mathbb{R} \to 0 < \varepsilon \to \delta \in \mathbb{R} \land 0 < \delta \land \forall x\,(\varphi_{161}) \to \varphi_{134}$
\item $\varphi_{426} :\equiv r \in \mathbb{R} \land 0 < r \land \varphi_{16} \to \varphi_{76} \to \varepsilon \in \mathbb{R} \to 0 < \varepsilon \to \delta \in \mathbb{R} \land 0 < \delta \land \forall x\,(\varphi_{161}) \to \varphi_{135}$
\item $\varphi_{427} :\equiv (\mathrm{Cgt}(f \circ p) \leftrightarrow \exists \alpha\,(\mathrm{Conv}(f \circ p,\alpha))) \to \mathrm{RFn}(f) \to c \in \mathbb{R} \to \mathrm{PNbhd}(f,c) \to \varphi_{356}$
\item $\varphi_{428} :\equiv r \in \mathbb{R} \land 0 < r \land \varphi_{16} \to \varphi_{76} \to \varepsilon \in \mathbb{R} \to 0 < \varepsilon \to \delta \in \mathbb{R} \land 0 < \delta \land \forall x\,(\varphi_{161}) \to \varphi_{138}$
\item $\varphi_{429} :\equiv \delta \in \mathbb{R} \land 0 < \delta \land \forall x\,(\varphi_{176}) \to N_{0} \in \mathbb{N} \land \varphi_{77} \to \varphi_{67} \to p(N_{0} \cup N_{1}) \in \mathbb{R} \to \varphi_{112}$
\item $\varphi_{430} :\equiv (\forall \alpha\,(\neg \mathrm{FLim}(f,c,\alpha)) \to \neg \mathrm{FLim}(f,c,\alpha)) \to \mathrm{RFn}(f) \to c \in \mathbb{R} \to \mathrm{PNbhd}(f,c) \to \varphi_{368}$
\item $\varphi_{431} :\equiv r \in \mathbb{R} \land 0 < r \land \varphi_{16} \to \varphi_{76} \to \varepsilon \in \mathbb{R} \to 0 < \varepsilon \to \delta \in \mathbb{R} \land 0 < \delta \land \forall x\,(\varphi_{161}) \to \varphi_{140}$
\item $\varphi_{432} :\equiv r \in \mathbb{R} \land 0 < r \land \varphi_{16} \to \varphi_{76} \to \varepsilon \in \mathbb{R} \to 0 < \varepsilon \to \delta \in \mathbb{R} \land 0 < \delta \land \forall x\,(\varphi_{161}) \to \varphi_{141}$
\item $\varphi_{433} :\equiv 2^{-1} \in \mathbb{R} \to \mathrm{RFn}(f) \to c \in \mathbb{R} \to \mathrm{PNbhd}(f,c) \to \varphi_{385}$
\item $\varphi_{434} :\equiv r \in \mathbb{R} \land 0 < r \land \varphi_{16} \to \varphi_{76} \to \varepsilon \in \mathbb{R} \to 0 < \varepsilon \to \delta \in \mathbb{R} \land 0 < \delta \land \forall x\,(\varphi_{161}) \to \varphi_{153}$
\item $\varphi_{435} :\equiv r \in \mathbb{R} \land 0 < r \land \varphi_{16} \to \varphi_{76} \to \varepsilon \in \mathbb{R} \to 0 < \varepsilon \to \delta \in \mathbb{R} \land 0 < \delta \land \forall x\,(\varphi_{161}) \to \varphi_{154}$
\item $\varphi_{436} :\equiv r \in \mathbb{R} \land 0 < r \land \varphi_{16} \to \varphi_{76} \to \varepsilon \in \mathbb{R} \to 0 < \varepsilon \to \delta \in \mathbb{R} \land 0 < \delta \land \forall x\,(\varphi_{161}) \to \varphi_{156}$
\item $\varphi_{437} :\equiv (\varepsilon \in \mathbb{R} \to 2^{-1} \in \mathbb{R} \to \varepsilon \cdot 2^{-1} \in \mathbb{R}) \to (\mathrm{RFn}(f) \to c \in \mathbb{R} \to \mathrm{PNbhd}(f,c) \to \varphi_{8}) \to \varphi_{327}$
\item $\varphi_{438} :\equiv r \in \mathbb{R} \land 0 < r \land \varphi_{16} \to \varphi_{76} \to \varepsilon \in \mathbb{R} \to 0 < \varepsilon \to \delta \in \mathbb{R} \land 0 < \delta \land \forall x\,(\varphi_{161}) \to \varphi_{157}$
\item $\varphi_{439} :\equiv (\mathrm{RFn}(f) \to c \in \mathbb{R} \to \mathrm{PNbhd}(f,c) \to \varphi_{227}) \to \mathrm{RFn}(f) \to c \in \mathbb{R} \to \mathrm{PNbhd}(f,c) \to \varphi_{229}$
\item $\varphi_{440} :\equiv r \in \mathbb{R} \land 0 < r \land \varphi_{16} \to \varphi_{76} \to \varepsilon \in \mathbb{R} \to 0 < \varepsilon \to \delta \in \mathbb{R} \land 0 < \delta \land \forall x\,(\varphi_{161}) \to \varphi_{159}$
\item $\varphi_{441} :\equiv r \in \mathbb{R} \land 0 < r \land \varphi_{16} \to \varphi_{76} \to \varepsilon \in \mathbb{R} \to 0 < \varepsilon \to \delta \in \mathbb{R} \land 0 < \delta \land \forall x\,(\varphi_{161}) \to \varphi_{160}$
\item $\varphi_{442} :\equiv r \in \mathbb{R} \land 0 < r \land \varphi_{16} \to \varphi_{76} \to \mathrm{Conv}(f \circ p,\alpha) \to \varepsilon \in \mathbb{R} \to \varphi_{292}$
\item $\varphi_{443} :\equiv r \in \mathbb{R} \land 0 < r \land \varphi_{16} \to \varphi_{76} \to \varepsilon \in \mathbb{R} \to 0 < \varepsilon \to \delta \in \mathbb{R} \land 0 < \delta \land \forall x\,(\varphi_{161}) \to \varphi_{163}$
\item $\varphi_{444} :\equiv r \in \mathbb{R} \land 0 < r \land \varphi_{16} \to \varphi_{76} \to \varepsilon \in \mathbb{R} \to 0 < \varepsilon \to \delta \in \mathbb{R} \land 0 < \delta \land \forall x\,(\varphi_{161}) \to \varphi_{164}$
\item $\varphi_{445} :\equiv r \in \mathbb{R} \land 0 < r \land \varphi_{16} \to \varphi_{76} \to \varepsilon \in \mathbb{R} \to 0 < \varepsilon \to \delta \in \mathbb{R} \land 0 < \delta \land \forall x\,(\varphi_{161}) \to \varphi_{165}$
\item $\varphi_{446} :\equiv r \in \mathbb{R} \land 0 < r \land \varphi_{16} \to \varphi_{76} \to \varepsilon \in \mathbb{R} \to 0 < \varepsilon \to \delta \in \mathbb{R} \land 0 < \delta \land \forall x\,(\varphi_{161}) \to \varphi_{166}$
\item $\varphi_{447} :\equiv \forall \varepsilon \in \mathbb{R}\,(0 < \varepsilon \to \exists \delta\,(\delta \in \mathbb{R} \land 0 < \delta \land \forall x\,(\varphi_{161}))) \to \varphi_{304}$
\item $\varphi_{448} :\equiv \forall \varepsilon \in \mathbb{R}\,(0 < \varepsilon \to \exists \delta\,(\delta \in \mathbb{R} \land 0 < \delta \land \forall x\,(\varphi_{161}))) \to \varphi_{307}$
\item $\varphi_{449} :\equiv r \in \mathbb{R} \land 0 < r \land \varphi_{16} \to \varphi_{76} \to \varepsilon \in \mathbb{R} \to 0 < \varepsilon \to \delta \in \mathbb{R} \land 0 < \delta \land \forall x\,(\varphi_{161}) \to \varphi_{177}$
\item $\varphi_{450} :\equiv r \in \mathbb{R} \land 0 < r \land \varphi_{16} \to \varphi_{76} \to \varepsilon \in \mathbb{R} \to 0 < \varepsilon \to \delta \in \mathbb{R} \land 0 < \delta \land \forall x\,(\varphi_{161}) \to \varphi_{178}$
\item $\varphi_{451} :\equiv (\mathrm{RFn}(f) \to c \in \mathbb{R} \to \mathrm{PNbhd}(f,c) \to \varphi_{4}) \to (\mathrm{RFn}(f) \to c \in \mathbb{R} \to \mathrm{PNbhd}(f,c) \to \varphi_{68}) \to \varphi_{316}$
\item $\varphi_{452} :\equiv (\forall p\,(\neg \varphi_{76}) \to \forall a\,\forall p\,(\neg \varphi_{76})) \to \mathrm{RFn}(f) \to c \in \mathbb{R} \to \mathrm{PNbhd}(f,c) \to \varphi_{303}$
\item $\varphi_{453} :\equiv r \in \mathbb{R} \land 0 < r \land \varphi_{16} \to \varphi_{76} \to \varepsilon \in \mathbb{R} \to 0 < \varepsilon \to \delta \in \mathbb{R} \land 0 < \delta \land \forall x\,(\varphi_{161}) \to \varphi_{188}$
\item $\varphi_{454} :\equiv (\mathrm{RFn}(f) \to \varphi_{264}) \to \varphi_{19} \to \varepsilon \in \mathbb{R} \to 0 < \varepsilon \to \exists \delta\,(\delta \in \mathbb{R} \land 0 < \delta \land \forall x\,(\varphi_{161}))$
\item $\varphi_{455} :\equiv r \in \mathbb{R} \land 0 < r \land \varphi_{16} \to \varphi_{76} \to \varepsilon \in \mathbb{R} \to 0 < \varepsilon \to \delta \in \mathbb{R} \land 0 < \delta \land \forall x\,(\varphi_{161}) \to \varphi_{197}$
\item $\varphi_{456} :\equiv \varphi_{19} \land \varepsilon \in \mathbb{R} \land 0 < \varepsilon \land \neg \exists \delta\,(\delta \in \mathbb{R} \land 0 < \delta \land \forall x\,(\varphi_{161})) \to \varphi_{274}$
\item $\varphi_{457} :\equiv \forall \varepsilon \in \mathbb{R}\,(0 < \varepsilon \to \exists \delta\,(\delta \in \mathbb{R} \land 0 < \delta \land \forall x\,(\varphi_{161}))) \to r \in \mathbb{R} \land 0 < r \land \varphi_{16} \to \varphi_{76} \to \varphi_{186}$
\item $\varphi_{458} :\equiv r \in \mathbb{R} \land 0 < r \land \varphi_{16} \to \varphi_{76} \to \varepsilon \in \mathbb{R} \to 0 < \varepsilon \to \delta \in \mathbb{R} \land 0 < \delta \land \forall x\,(\varphi_{161}) \to \varphi_{200}$
\item $\varphi_{459} :\equiv \forall \varepsilon \in \mathbb{R}\,(0 < \varepsilon \to \exists \delta\,(\delta \in \mathbb{R} \land 0 < \delta \land \forall x\,(\varphi_{161}))) \to r \in \mathbb{R} \land 0 < r \land \varphi_{16} \to \varphi_{76} \to \varphi_{193}$
\item $\varphi_{460} :\equiv r \in \mathbb{R} \land 0 < r \land \varphi_{16} \to \varphi_{76} \to \varepsilon \in \mathbb{R} \to 0 < \varepsilon \to \delta \in \mathbb{R} \land 0 < \delta \land \forall x\,(\varphi_{161}) \to \varphi_{206}$
\item $\varphi_{461} :\equiv r \in \mathbb{R} \land 0 < r \land \varphi_{16} \to \varphi_{76} \to \varepsilon \in \mathbb{R} \to 0 < \varepsilon \to \delta \in \mathbb{R} \land 0 < \delta \land \forall x\,(\varphi_{161}) \to \varphi_{207}$
\item $\varphi_{462} :\equiv \forall \alpha\,(\varphi_{262} \to \neg \mathrm{FLim}(f,c,\alpha)) \to \forall \alpha\,(\varphi_{262}) \to \forall \alpha\,(\neg \mathrm{FLim}(f,c,\alpha))$
\item $\varphi_{463} :\equiv \varphi_{362} \land \mathrm{Conv}(f \circ p,\alpha) \to \varepsilon \in \mathbb{R} \to 0 < \varepsilon \to \exists \delta\,(\delta \in \mathbb{R} \land 0 < \delta \land \varphi_{69})$
\item $\varphi_{464} :\equiv \varphi_{362} \land \mathrm{Conv}(f \circ p,\alpha) \land \varepsilon \in \mathbb{R} \land 0 < \varepsilon \land \neg \exists \delta\,(\delta \in \mathbb{R} \land 0 < \delta \land \varphi_{69})$
\item $\varphi_{465} :\equiv (\mathrm{RFn}(f) \to c \in \mathbb{R} \to \mathrm{PNbhd}(f,c) \to \varphi_{212}) \to \mathrm{RFn}(f) \to c \in \mathbb{R} \to \mathrm{PNbhd}(f,c) \to \varphi_{296}$
\item $\varphi_{466} :\equiv \forall \varepsilon \in \mathbb{R}\,(0 < \varepsilon \to \exists \delta\,(\delta \in \mathbb{R} \land 0 < \delta \land \forall x\,(\varphi_{161}))) \to \varphi_{333}$
\item $\varphi_{467} :\equiv r \in \mathbb{R} \land 0 < r \land \varphi_{16} \to \varphi_{76} \to \mathrm{Conv}(f \circ p,\alpha) \to \varepsilon \in \mathbb{R} \to \varphi_{332}$
\item $\varphi_{468} :\equiv \forall \varepsilon\,(\varphi_{362} \land \mathrm{Conv}(f \circ p,\alpha)) \to \forall \varepsilon \in \mathbb{R}\,(0 < \varepsilon \to \exists \delta\,(\delta \in \mathbb{R} \land 0 < \delta \land \varphi_{69}))$
\item $\varphi_{469} :\equiv \forall \varepsilon \in \mathbb{R}\,(0 < \varepsilon \to \exists \delta\,(\delta \in \mathbb{R} \land 0 < \delta \land \forall x\,(\varphi_{161}))) \to \varphi_{334}$
\item $\varphi_{470} :\equiv \forall \varepsilon \in \mathbb{R}\,(0 < \varepsilon \to \exists \delta\,(\delta \in \mathbb{R} \land 0 < \delta \land \forall x\,(\varphi_{161}))) \to \varphi_{338}$
\item $\varphi_{471} :\equiv (\mathrm{RFn}(f) \to c \in \mathbb{R} \to \mathrm{PNbhd}(f,c) \to \varphi_{223}) \to \mathrm{RFn}(f) \to c \in \mathbb{R} \to \mathrm{PNbhd}(f,c) \to \varphi_{295}$
\item $\varphi_{472} :\equiv \varepsilon \in \mathbb{R} \to 0 < \varepsilon \to \delta \in \mathbb{R} \land 0 < \delta \land \forall x\,(\varphi_{176}) \to N_{0} \in \mathbb{N} \land \varphi_{77} \to \varphi_{281}$
\item $\varphi_{473} :\equiv r \in \mathbb{R} \land 0 < r \land \varphi_{16} \to \varphi_{76} \to \mathrm{Conv}(f \circ p,\alpha) \to \varepsilon \in \mathbb{R} \to \varphi_{339}$
\item $\varphi_{474} :\equiv r \in \mathbb{R} \land 0 < r \land \varphi_{16} \to \varphi_{76} \to \mathrm{Conv}(f \circ p,\alpha) \to \varepsilon \in \mathbb{R} \to \varphi_{340}$
\item $\varphi_{475} :\equiv r \in \mathbb{R} \land 0 < r \land \varphi_{16} \to \varphi_{76} \to \mathrm{Conv}(f \circ p,\alpha) \to \varepsilon \in \mathbb{R} \to \varphi_{342}$
\item $\varphi_{476} :\equiv r \in \mathbb{R} \land 0 < r \land \varphi_{16} \to \varphi_{76} \to \mathrm{Conv}(f \circ p,\alpha) \to \varepsilon \in \mathbb{R} \to 0 < \varepsilon \to \varphi_{335}$
\item $\varphi_{477} :\equiv r \in \mathbb{R} \land 0 < r \land \varphi_{16} \to \varphi_{76} \to \varepsilon \in \mathbb{R} \to 0 < \varepsilon \to \delta \in \mathbb{R} \land 0 < \delta \land \forall x\,(\varphi_{161}) \to \varphi_{246}$
\item $\varphi_{478} :\equiv \forall \varepsilon \in \mathbb{R}\,(0 < \varepsilon \to \exists \delta\,(\delta \in \mathbb{R} \land 0 < \delta \land \forall x\,(\varphi_{161}))) \to r \in \mathbb{R} \land 0 < r \land \varphi_{16} \to \varphi_{76} \to \varphi_{236}$
\item $\varphi_{479} :\equiv r \in \mathbb{R} \land 0 < r \land \varphi_{16} \to \varphi_{76} \to \mathrm{Conv}(f \circ p,\alpha) \to \varepsilon \in \mathbb{R} \to 0 < \varepsilon \to \varphi_{346}$
\item $\varphi_{480} :\equiv r \in \mathbb{R} \land 0 < r \land \varphi_{16} \to \varphi_{76} \to \mathrm{Conv}(f \circ p,\alpha) \to \varepsilon \in \mathbb{R} \to 0 < \varepsilon \to \varphi_{347}$
\item $\varphi_{481} :\equiv r \in \mathbb{R} \land 0 < r \land \varphi_{16} \to \varphi_{76} \to \mathrm{Conv}(f \circ p,\alpha) \to \varepsilon \in \mathbb{R} \to 0 < \varepsilon \to \varphi_{348}$
\item $\varphi_{482} :\equiv r \in \mathbb{R} \land 0 < r \land \varphi_{16} \to \varphi_{76} \to \mathrm{Conv}(f \circ p,\alpha) \to \varepsilon \in \mathbb{R} \to 0 < \varepsilon \to \varphi_{350}$
\item $\varphi_{483} :\equiv \mathrm{Conv}(f \circ p,\alpha) \to \varepsilon \in \mathbb{R} \to 0 < \varepsilon \to \delta \in \mathbb{R} \land 0 < \delta \land \forall x\,(\varphi_{176}) \to \varphi_{344}$
\item $\varphi_{484} :\equiv \forall \varepsilon \in \mathbb{R}\,(0 < \varepsilon \to \exists \delta\,(\delta \in \mathbb{R} \land 0 < \delta \land \forall x\,(\varphi_{161}))) \to r \in \mathbb{R} \land 0 < r \land \varphi_{16} \to \varphi_{76} \to \varphi_{244}$
\item $\varphi_{485} :\equiv \forall \varepsilon \in \mathbb{R}\,(0 < \varepsilon \to \exists \delta\,(\delta \in \mathbb{R} \land 0 < \delta \land \forall x\,(\varphi_{161}))) \to r \in \mathbb{R} \land 0 < r \land \varphi_{16} \to \varphi_{76} \to \varphi_{248}$
\item $\varphi_{486} :\equiv \delta \in \mathbb{R} \land 0 < \delta \land \forall x\,(\varphi_{176}) \to N_{0} \in \mathbb{N} \land \varphi_{77} \to \varphi_{67} \to x \in \mathrm{dom}(f) \to 0 < \lvert x - c \rvert \to \varphi_{235}$
\item $\varphi_{487} :\equiv r \in \mathbb{R} \land 0 < r \land \varphi_{16} \to \varphi_{76} \to \varepsilon \in \mathbb{R} \to 0 < \varepsilon \to \delta \in \mathbb{R} \land 0 < \delta \land \forall x\,(\varphi_{161}) \to \varphi_{271}$
\item $\varphi_{488} :\equiv \forall \varepsilon \in \mathbb{R}\,(0 < \varepsilon \to \exists \delta\,(\delta \in \mathbb{R} \land 0 < \delta \land \forall x\,(\varphi_{161}))) \to r \in \mathbb{R} \land 0 < r \land \varphi_{16} \to \varphi_{76} \to \varepsilon \in \mathbb{R} \to 0 < \varepsilon \to \varphi_{234}$
\item $\varphi_{489} :\equiv (\mathrm{RFn}(f) \to x \in \mathrm{dom}(f) \to f(x) \in \mathbb{R}) \to (\mathrm{RFn}(f) \to \mathrm{RFn}(f)) \to \varphi_{465}$
\item $\varphi_{490} :\equiv r \in \mathbb{R} \land 0 < r \land \varphi_{16} \to \varphi_{76} \to \mathrm{Conv}(f \circ p,\alpha) \to \varepsilon \in \mathbb{R} \to 0 < \varepsilon \to \varphi_{367}$
\item $\varphi_{491} :\equiv \forall \varepsilon \in \mathbb{R}\,(0 < \varepsilon \to \exists \delta\,(\delta \in \mathbb{R} \land 0 < \delta \land \forall x\,(\varphi_{161}))) \to \varphi_{378}$
\item $\varphi_{492} :\equiv (\mathrm{RFn}(f) \to y \in \mathrm{dom}(f) \to f(y) \in \mathbb{R}) \to (\mathrm{RFn}(f) \to \mathrm{RFn}(f)) \to \varphi_{471}$
\item $\varphi_{493} :\equiv (\mathrm{RFn}(f) \to \varphi_{280}) \to \mathrm{RFn}(f) \to c \in \mathbb{R} \to \mathrm{PNbhd}(f,c) \to \varphi_{313}$
\item $\varphi_{494} :\equiv \forall \varepsilon \in \mathbb{R}\,(0 < \varepsilon \to \exists \delta\,(\delta \in \mathbb{R} \land 0 < \delta \land \forall x\,(\varphi_{161}))) \to \varphi_{386}$
\item $\varphi_{495} :\equiv (\forall n \in \mathbb{N}\,(0 < \lvert p(n) - c \rvert \land \lvert p(n) - c \rvert < r) \to \varphi_{85}) \to \mathrm{RFn}(f) \to c \in \mathbb{R} \to \mathrm{PNbhd}(f,c) \to \varphi_{404}$
\item $\varphi_{496} :\equiv \varphi_{362} \land \varepsilon \in \mathbb{R} \land 0 < \varepsilon \land (\delta \in \mathbb{R} \land 0 < \delta \land \forall x\,(\varphi_{161})) \land (N \in \mathbb{N} \land \forall n \in \mathbb{N}\,(N \leq n \to \lvert p(n) - c \rvert < \delta))$
\item $\varphi_{497} :\equiv \forall \varepsilon \in \mathbb{R}\,(0 < \varepsilon \to \exists \delta\,(\delta \in \mathbb{R} \land 0 < \delta \land \forall x\,(\varphi_{161}))) \to \varphi_{398}$
\item $\varphi_{498} :\equiv \forall \varepsilon \in \mathbb{R}\,(0 < \varepsilon \to \exists \delta\,(\delta \in \mathbb{R} \land 0 < \delta \land \forall x\,(\varphi_{161}))) \to \varphi_{399}$
\item $\varphi_{499} :\equiv \forall \varepsilon \in \mathbb{R}\,(0 < \varepsilon \to \exists \delta\,(\delta \in \mathbb{R} \land 0 < \delta \land \forall x\,(\varphi_{161}))) \to \varphi_{410}$
\item $\varphi_{500} :\equiv (\varphi_{16} \to p(n) \in \mathbb{R} \to \varphi_{27}) \to \mathrm{RFn}(f) \to c \in \mathbb{R} \to \mathrm{PNbhd}(f,c) \to \varphi_{422}$
\item $\varphi_{501} :\equiv \forall \varepsilon \in \mathbb{R}\,(0 < \varepsilon \to \exists \delta\,(\delta \in \mathbb{R} \land 0 < \delta \land \forall x\,(\varphi_{161}))) \to r \in \mathbb{R} \land 0 < r \land \varphi_{16} \to \varphi_{76} \to \varepsilon \in \mathbb{R} \to \varphi_{282}$
\item $\varphi_{502} :\equiv \forall \varepsilon \in \mathbb{R}\,(0 < \varepsilon \to \exists \delta\,(\delta \in \mathbb{R} \land 0 < \delta \land \forall x\,(\varphi_{161}))) \to \varphi_{423}$
\item $\varphi_{503} :\equiv \forall \varepsilon \in \mathbb{R}\,(0 < \varepsilon \to \exists \delta\,(\delta \in \mathbb{R} \land 0 < \delta \land \forall x\,(\varphi_{161}))) \to \varphi_{424}$
\item $\varphi_{504} :\equiv \forall \varepsilon \in \mathbb{R}\,(0 < \varepsilon \to \exists \delta\,(\delta \in \mathbb{R} \land 0 < \delta \land \forall x\,(\varphi_{161}))) \to \varphi_{425}$
\item $\varphi_{505} :\equiv \forall \varepsilon \in \mathbb{R}\,(0 < \varepsilon \to \exists \delta\,(\delta \in \mathbb{R} \land 0 < \delta \land \forall x\,(\varphi_{161}))) \to \varphi_{426}$
\item $\varphi_{506} :\equiv (\mathrm{RFn}(f) \to c \in \mathbb{R} \to \mathrm{PNbhd}(f,c) \to \varphi_{294}) \to \mathrm{RFn}(f) \to c \in \mathbb{R} \to \mathrm{PNbhd}(f,c) \to \varphi_{311}$
\item $\varphi_{507} :\equiv \forall \varepsilon \in \mathbb{R}\,(0 < \varepsilon \to \exists \delta\,(\delta \in \mathbb{R} \land 0 < \delta \land \forall x\,(\varphi_{161}))) \to \varphi_{428}$
\item $\varphi_{508} :\equiv (\mathrm{RFn}(f) \to c \in \mathbb{R} \to \mathrm{PNbhd}(f,c) \to \varphi_{303}) \to \varphi_{261} \to \neg (r \in \mathbb{R} \land 0 < r \land \varphi_{16})$
\item $\varphi_{509} :\equiv r \in \mathbb{R} \land 0 < r \land \varphi_{16} \to \varphi_{76} \to \mathrm{Conv}(f \circ p,\alpha) \to \varepsilon \in \mathbb{R} \to 0 < \varepsilon \to \varphi_{407}$
\item $\varphi_{510} :\equiv r \in \mathbb{R} \land 0 < r \land \varphi_{16} \to \varphi_{76} \to \mathrm{Conv}(f \circ p,\alpha) \to \varepsilon \in \mathbb{R} \to 0 < \varepsilon \to \varphi_{408}$
\item $\varphi_{511} :\equiv \forall \varepsilon \in \mathbb{R}\,(0 < \varepsilon \to \exists \delta\,(\delta \in \mathbb{R} \land 0 < \delta \land \forall x\,(\varphi_{161}))) \to \varphi_{431}$
\item $\varphi_{512} :\equiv \forall \varepsilon \in \mathbb{R}\,(0 < \varepsilon \to \exists \delta\,(\delta \in \mathbb{R} \land 0 < \delta \land \forall x\,(\varphi_{161}))) \to \varphi_{432}$
\item $\varphi_{513} :\equiv r \in \mathbb{R} \land 0 < r \land \varphi_{16} \to \varphi_{76} \to \mathrm{Conv}(f \circ p,\alpha) \to \varepsilon \in \mathbb{R} \to 0 < \varepsilon \to \varphi_{413}$
\item $\varphi_{514} :\equiv r \in \mathbb{R} \land 0 < r \land \varphi_{16} \to \varphi_{76} \to \mathrm{Conv}(f \circ p,\alpha) \to \varepsilon \in \mathbb{R} \to 0 < \varepsilon \to \varphi_{414}$
\item $\varphi_{515} :\equiv \forall r\,(\varphi_{261} \to \neg (r \in \mathbb{R} \land 0 < r \land \varphi_{16})) \to \forall r\,(\varphi_{261}) \to \forall r\,(\neg (r \in \mathbb{R} \land 0 < r \land \varphi_{16}))$
\item $\varphi_{516} :\equiv r \in \mathbb{R} \land 0 < r \land \varphi_{16} \to \varphi_{76} \to \mathrm{Conv}(f \circ p,\alpha) \to \varepsilon \in \mathbb{R} \to 0 < \varepsilon \to \varphi_{418}$
\item $\varphi_{517} :\equiv \forall \varepsilon \in \mathbb{R}\,(0 < \varepsilon \to \exists \delta\,(\delta \in \mathbb{R} \land 0 < \delta \land \forall x\,(\varphi_{161}))) \to r \in \mathbb{R} \land 0 < r \land \varphi_{16} \to \varphi_{76} \to \varepsilon \in \mathbb{R} \to 0 < \varepsilon \to \varphi_{287}$
\item $\varphi_{518} :\equiv \forall \varepsilon \in \mathbb{R}\,(0 < \varepsilon \to \exists \delta\,(\delta \in \mathbb{R} \land 0 < \delta \land \forall x\,(\varphi_{161}))) \to r \in \mathbb{R} \land 0 < r \land \varphi_{16} \to \varphi_{76} \to \varepsilon \in \mathbb{R} \to 0 < \varepsilon \to \varphi_{288}$
\item $\varphi_{519} :\equiv \forall \varepsilon \in \mathbb{R}\,(0 < \varepsilon \to \exists \delta\,(\delta \in \mathbb{R} \land 0 < \delta \land \forall x\,(\varphi_{161}))) \to r \in \mathbb{R} \land 0 < r \land \varphi_{16} \to \varphi_{76} \to \varepsilon \in \mathbb{R} \to 0 < \varepsilon \to \varphi_{289}$
\item $\varphi_{520} :\equiv \forall \varepsilon \in \mathbb{R}\,(0 < \varepsilon \to \exists \delta\,(\delta \in \mathbb{R} \land 0 < \delta \land \forall x\,(\varphi_{161}))) \to r \in \mathbb{R} \land 0 < r \land \varphi_{16} \to \varphi_{76} \to \varepsilon \in \mathbb{R} \to 0 < \varepsilon \to \varphi_{290}$
\item $\varphi_{521} :\equiv (\mathrm{RFn}(f) \to \varphi_{319}) \to \varphi_{19} \land \varepsilon \in \mathbb{R} \land 0 < \varepsilon \land (\delta \in \mathbb{R} \land 0 < \delta \land \varphi_{99}) \to \varphi_{161}$
\item $\varphi_{522} :\equiv \forall \varepsilon \in \mathbb{R}\,(0 < \varepsilon \to \exists \delta\,(\delta \in \mathbb{R} \land 0 < \delta \land \forall x\,(\varphi_{161}))) \to r \in \mathbb{R} \land 0 < r \land \varphi_{16} \to \varphi_{76} \to \varphi_{317}$
\item $\varphi_{523} :\equiv \forall y\,(\varphi_{326}) \to \forall y\,(\varphi_{19} \land \varepsilon \in \mathbb{R} \land 0 < \varepsilon \land (\delta \in \mathbb{R} \land 0 < \delta \land \varphi_{99})) \to \varphi_{161}$
\item $\varphi_{524} :\equiv \forall \varepsilon \in \mathbb{R}\,(0 < \varepsilon \to \exists \delta\,(\delta \in \mathbb{R} \land 0 < \delta \land \forall x\,(\varphi_{161}))) \to \varphi_{434}$
\item $\varphi_{525} :\equiv \forall \varepsilon \in \mathbb{R}\,(0 < \varepsilon \to \exists \delta\,(\delta \in \mathbb{R} \land 0 < \delta \land \forall x\,(\varphi_{161}))) \to \varphi_{435}$
\item $\varphi_{526} :\equiv r \in \mathbb{R} \land 0 < r \land \varphi_{16} \to \varphi_{76} \to \mathrm{Conv}(f \circ p,\alpha) \to \varepsilon \in \mathbb{R} \to 0 < \varepsilon \to \varphi_{429}$
\item $\varphi_{527} :\equiv \forall \varepsilon \in \mathbb{R}\,(0 < \varepsilon \to \exists \delta\,(\delta \in \mathbb{R} \land 0 < \delta \land \forall x\,(\varphi_{161}))) \to \varphi_{436}$
\item $\varphi_{528} :\equiv \forall x\,(\varphi_{19} \land \varepsilon \in \mathbb{R} \land 0 < \varepsilon \land (\delta \in \mathbb{R} \land 0 < \delta \land \varphi_{99}) \to \varphi_{161}) \to \varphi_{329}$
\item $\varphi_{529} :\equiv \forall \varepsilon \in \mathbb{R}\,(0 < \varepsilon \to \exists \delta\,(\delta \in \mathbb{R} \land 0 < \delta \land \forall x\,(\varphi_{161}))) \to \varphi_{438}$
\item $\varphi_{530} :\equiv \forall \varepsilon \in \mathbb{R}\,(0 < \varepsilon \to \exists \delta\,(\delta \in \mathbb{R} \land 0 < \delta \land \forall x\,(\varphi_{161}))) \to \varphi_{440}$
\item $\varphi_{531} :\equiv \forall \varepsilon \in \mathbb{R}\,(0 < \varepsilon \to \exists \delta\,(\delta \in \mathbb{R} \land 0 < \delta \land \forall x\,(\varphi_{161}))) \to \varphi_{441}$
\item $\varphi_{532} :\equiv \forall \varepsilon \in \mathbb{R}\,(0 < \varepsilon \to \exists \delta\,(\delta \in \mathbb{R} \land 0 < \delta \land \forall x\,(\varphi_{161}))) \to \varphi_{442}$
\item $\varphi_{533} :\equiv \forall \varepsilon \in \mathbb{R}\,(0 < \varepsilon \to \exists \delta\,(\delta \in \mathbb{R} \land 0 < \delta \land \forall x\,(\varphi_{161}))) \to \varphi_{443}$
\item $\varphi_{534} :\equiv \forall \varepsilon \in \mathbb{R}\,(0 < \varepsilon \to \exists \delta\,(\delta \in \mathbb{R} \land 0 < \delta \land \forall x\,(\varphi_{161}))) \to \varphi_{444}$
\item $\varphi_{535} :\equiv \forall \varepsilon \in \mathbb{R}\,(0 < \varepsilon \to \exists \delta\,(\delta \in \mathbb{R} \land 0 < \delta \land \forall x\,(\varphi_{161}))) \to \varphi_{445}$
\item $\varphi_{536} :\equiv \forall \varepsilon \in \mathbb{R}\,(0 < \varepsilon \to \exists \delta\,(\delta \in \mathbb{R} \land 0 < \delta \land \forall x\,(\varphi_{161}))) \to \varphi_{446}$
\item $\varphi_{537} :\equiv \varphi_{362} \land \varepsilon \in \mathbb{R} \land 0 < \varepsilon \land (\delta \in \mathbb{R} \land 0 < \delta \land \forall x\,(\varphi_{161})) \land \neg \exists N \in \mathbb{N}\,\forall m\,\forall n\,(\varphi_{90})$
\item $\varphi_{538} :\equiv \varphi_{362} \land \varepsilon \in \mathbb{R} \land 0 < \varepsilon \land \neg \exists N \in \mathbb{N}\,\forall m\,\forall n\,(\varphi_{90}) \to \neg (\delta \in \mathbb{R} \land 0 < \delta \land \forall x\,(\varphi_{161}))$
\item $\varphi_{539} :\equiv \forall \varepsilon \in \mathbb{R}\,(0 < \varepsilon \to \exists \delta\,(\delta \in \mathbb{R} \land 0 < \delta \land \forall x\,(\varphi_{161}))) \to r \in \mathbb{R} \land 0 < r \land \varphi_{16} \to \varphi_{76} \to \varepsilon \in \mathbb{R} \to 0 < \varepsilon \to \varphi_{309}$
\item $\varphi_{540} :\equiv \forall \varepsilon \in \mathbb{R}\,(0 < \varepsilon \to \exists \delta\,(\delta \in \mathbb{R} \land 0 < \delta \land \forall x\,(\varphi_{161}))) \to \varphi_{449}$
\item $\varphi_{541} :\equiv \forall \delta\,(\varphi_{362} \land \varepsilon \in \mathbb{R} \land 0 < \varepsilon \land \neg \exists N \in \mathbb{N}\,\forall m\,\forall n\,(\varphi_{90})) \to \forall \delta\,(\neg (\delta \in \mathbb{R} \land 0 < \delta \land \forall x\,(\varphi_{161})))$
\item $\varphi_{542} :\equiv \forall \varepsilon \in \mathbb{R}\,(0 < \varepsilon \to \exists \delta\,(\delta \in \mathbb{R} \land 0 < \delta \land \forall x\,(\varphi_{161}))) \to \varphi_{450}$
\item $\varphi_{543} :\equiv (\mathrm{RFn}(f) \to c \in \mathbb{R} \to \mathrm{PNbhd}(f,c) \to \varphi_{312}) \to \mathrm{RFn}(f) \to \varphi_{331}$
\item $\varphi_{544} :\equiv \forall \varepsilon \in \mathbb{R}\,(0 < \varepsilon \to \exists \delta\,(\delta \in \mathbb{R} \land 0 < \delta \land \forall x\,(\varphi_{161}))) \to \varphi_{453}$
\item $\varphi_{545} :\equiv \forall \varepsilon \in \mathbb{R}\,(0 < \varepsilon \to \exists \delta\,(\delta \in \mathbb{R} \land 0 < \delta \land \forall x\,(\varphi_{161}))) \to \varphi_{455}$
\item $\varphi_{546} :\equiv \forall \varepsilon \in \mathbb{R}\,(0 < \varepsilon \to \exists \delta\,(\delta \in \mathbb{R} \land 0 < \delta \land \forall x\,(\varphi_{161}))) \to r \in \mathbb{R} \land 0 < r \land \varphi_{16} \to \varphi_{76} \to \varepsilon \in \mathbb{R} \to \varphi_{324}$
\item $\varphi_{547} :\equiv \forall \varepsilon \in \mathbb{R}\,(0 < \varepsilon \to \exists \delta\,(\delta \in \mathbb{R} \land 0 < \delta \land \forall x\,(\varphi_{161}))) \to \varphi_{458}$
\item $\varphi_{548} :\equiv (\mathrm{RFn}(f) \to c \in \mathbb{R} \to \mathrm{PNbhd}(f,c) \to \varphi_{336}) \to \mathrm{RFn}(f) \to c \in \mathbb{R} \to \mathrm{PNbhd}(f,c) \to \varphi_{293}$
\item $\varphi_{549} :\equiv (\mathrm{RFn}(f) \to c \in \mathbb{R} \to \mathrm{PNbhd}(f,c) \to \varphi_{293}) \to \mathrm{RFn}(f) \to c \in \mathbb{R} \to \mathrm{PNbhd}(f,c) \to \varphi_{337}$
\item $\varphi_{550} :\equiv \varphi_{362} \land \mathrm{Conv}(f \circ p,\alpha) \land \varepsilon \in \mathbb{R} \land 0 < \varepsilon \land (\delta \in \mathbb{R} \land 0 < \delta \land \forall x\,(\varphi_{176})) \land \neg \exists \delta\,(\delta \in \mathbb{R} \land 0 < \delta \land \varphi_{69})$
\item $\varphi_{551} :\equiv (\mathrm{RFn}(f) \to c \in \mathbb{R} \to \mathrm{PNbhd}(f,c) \to \varphi_{336}) \to \mathrm{RFn}(f) \to c \in \mathbb{R} \to \mathrm{PNbhd}(f,c) \to \varphi_{301}$
\item $\varphi_{552} :\equiv \forall \varepsilon \in \mathbb{R}\,(0 < \varepsilon \to \exists \delta\,(\delta \in \mathbb{R} \land 0 < \delta \land \forall x\,(\varphi_{161}))) \to \varphi_{460}$
\item $\varphi_{553} :\equiv \forall \varepsilon \in \mathbb{R}\,(0 < \varepsilon \to \exists \delta\,(\delta \in \mathbb{R} \land 0 < \delta \land \forall x\,(\varphi_{161}))) \to \varphi_{461}$
\item $\varphi_{554} :\equiv (\varphi_{261} \to \forall r\,(\varphi_{261})) \to \mathrm{RFn}(f) \to c \in \mathbb{R} \to \mathrm{PNbhd}(f,c) \to \varphi_{227}$
\item $\varphi_{555} :\equiv (\varphi_{262} \to \forall \alpha\,(\varphi_{262})) \to \mathrm{RFn}(f) \to c \in \mathbb{R} \to \mathrm{PNbhd}(f,c) \to \varphi_{228}$
\item $\varphi_{556} :\equiv \varphi_{537} \to \neg (N \in \mathbb{N} \land \forall n \in \mathbb{N}\,(N \leq n \to \lvert p(n) - c \rvert < \delta))$
\item $\varphi_{557} :\equiv \forall \varepsilon \in \mathbb{R}\,(0 < \varepsilon \to \exists \delta\,(\delta \in \mathbb{R} \land 0 < \delta \land \forall x\,(\varphi_{161}))) \to \varphi_{467}$
\item $\varphi_{558} :\equiv \forall \varepsilon \in \mathbb{R}\,(0 < \varepsilon \to \exists \delta\,(\delta \in \mathbb{R} \land 0 < \delta \land \forall x\,(\varphi_{161}))) \to r \in \mathbb{R} \land 0 < r \land \varphi_{16} \to \varphi_{76} \to \mathrm{Conv}(f \circ p,\alpha) \to \varphi_{345}$
\item $\varphi_{559} :\equiv \forall \varepsilon \in \mathbb{R}\,(0 < \varepsilon \to \exists \delta\,(\delta \in \mathbb{R} \land 0 < \delta \land \forall x\,(\varphi_{161}))) \to \varphi_{473}$
\item $\varphi_{560} :\equiv \forall \varepsilon \in \mathbb{R}\,(0 < \varepsilon \to \exists \delta\,(\delta \in \mathbb{R} \land 0 < \delta \land \forall x\,(\varphi_{161}))) \to \varphi_{474}$
\item $\varphi_{561} :\equiv \forall \varepsilon \in \mathbb{R}\,(0 < \varepsilon \to \exists \delta\,(\delta \in \mathbb{R} \land 0 < \delta \land \forall x\,(\varphi_{161}))) \to \varphi_{475}$
\item $\varphi_{562} :\equiv \forall \varepsilon \in \mathbb{R}\,(0 < \varepsilon \to \exists \delta\,(\delta \in \mathbb{R} \land 0 < \delta \land \forall x\,(\varphi_{161}))) \to \varphi_{476}$
\item $\varphi_{563} :\equiv (f \circ p(m) = f(p(m)) \to f(p(m)) = f \circ p(m)) \to \mathrm{RFn}(f) \to c \in \mathbb{R} \to \mathrm{PNbhd}(f,c) \to \varphi_{518}$
\item $\varphi_{564} :\equiv (f \circ p(n) = f(p(n)) \to f(p(n)) = f \circ p(n)) \to \mathrm{RFn}(f) \to c \in \mathbb{R} \to \mathrm{PNbhd}(f,c) \to \varphi_{520}$
\item $\varphi_{565} :\equiv \forall \varepsilon \in \mathbb{R}\,(0 < \varepsilon \to \exists \delta\,(\delta \in \mathbb{R} \land 0 < \delta \land \forall x\,(\varphi_{161}))) \to \varphi_{477}$
\item $\varphi_{566} :\equiv \forall \varepsilon \in \mathbb{R}\,(0 < \varepsilon \to \exists \delta\,(\delta \in \mathbb{R} \land 0 < \delta \land \forall x\,(\varphi_{161}))) \to \varphi_{479}$
\item $\varphi_{567} :\equiv \forall \varepsilon \in \mathbb{R}\,(0 < \varepsilon \to \exists \delta\,(\delta \in \mathbb{R} \land 0 < \delta \land \forall x\,(\varphi_{161}))) \to \varphi_{480}$
\item $\varphi_{568} :\equiv \forall \varepsilon \in \mathbb{R}\,(0 < \varepsilon \to \exists \delta\,(\delta \in \mathbb{R} \land 0 < \delta \land \forall x\,(\varphi_{161}))) \to \varphi_{481}$
\item $\varphi_{569} :\equiv \forall \varepsilon \in \mathbb{R}\,(0 < \varepsilon \to \exists \delta\,(\delta \in \mathbb{R} \land 0 < \delta \land \forall x\,(\varphi_{161}))) \to \varphi_{482}$
\item $\varphi_{570} :\equiv \varphi_{362} \land \mathrm{Conv}(f \circ p,\alpha) \land \varepsilon \in \mathbb{R} \land 0 < \varepsilon \land (\delta \in \mathbb{R} \land 0 < \delta \land \forall x\,(\varphi_{176})) \land (N_{0} \in \mathbb{N} \land \varphi_{77}) \land \varphi_{67}$
\item $\varphi_{571} :\equiv (\mathrm{RFn}(f) \to c \in \mathbb{R} \to \mathrm{PNbhd}(f,c) \to \varphi_{311}) \to \varphi_{136} \to \mathrm{RFn}(f) \to c \in \mathbb{R} \to \mathrm{PNbhd}(f,c) \to \varphi_{312}$
\item $\varphi_{572} :\equiv (\mathrm{RFn}(f) \to c \in \mathbb{R} \to \mathrm{PNbhd}(f,c) \to \varphi_{313}) \to \varphi_{137} \to \mathrm{RFn}(f) \to c \in \mathbb{R} \to \mathrm{PNbhd}(f,c) \to \varphi_{314}$
\item $\varphi_{573} :\equiv \forall \varepsilon \in \mathbb{R}\,(0 < \varepsilon \to \exists \delta\,(\delta \in \mathbb{R} \land 0 < \delta \land \forall x\,(\varphi_{161}))) \to \varphi_{487}$
\item $\varphi_{574} :\equiv \forall \varepsilon \in \mathbb{R}\,(0 < \varepsilon \to \exists \delta\,(\delta \in \mathbb{R} \land 0 < \delta \land \forall x\,(\varphi_{161}))) \to r \in \mathbb{R} \land 0 < r \land \varphi_{16} \to \varphi_{76} \to \varphi_{394}$
\item $\varphi_{575} :\equiv \forall \varepsilon \in \mathbb{R}\,(0 < \varepsilon \to \exists \delta\,(\delta \in \mathbb{R} \land 0 < \delta \land \forall x\,(\varphi_{161}))) \to r \in \mathbb{R} \land 0 < r \land \varphi_{16} \to \varphi_{76} \to \varphi_{396}$
\item $\varphi_{576} :\equiv \forall \varepsilon \in \mathbb{R}\,(0 < \varepsilon \to \exists \delta\,(\delta \in \mathbb{R} \land 0 < \delta \land \forall x\,(\varphi_{161}))) \to r \in \mathbb{R} \land 0 < r \land \varphi_{16} \to \varphi_{76} \to \varphi_{397}$
\item $\varphi_{577} :\equiv \forall \varepsilon \in \mathbb{R}\,(0 < \varepsilon \to \exists \delta\,(\delta \in \mathbb{R} \land 0 < \delta \land \forall x\,(\varphi_{161}))) \to r \in \mathbb{R} \land 0 < r \land \varphi_{16} \to \varphi_{76} \to \varphi_{402}$
\item $\varphi_{578} :\equiv \forall \varepsilon \in \mathbb{R}\,(0 < \varepsilon \to \exists \delta\,(\delta \in \mathbb{R} \land 0 < \delta \land \forall x\,(\varphi_{161}))) \to \varphi_{490}$
\item $\varphi_{579} :\equiv (f(y) \in \mathbb{R} \to \alpha \in \mathbb{R} \to \lvert f(y) - \alpha \rvert = \lvert \alpha - f(y) \rvert) \to (\mathrm{RFn}(f) \to c \in \mathbb{R} \to \mathrm{PNbhd}(f,c) \to \varphi_{295}) \to \varphi_{375}$
\item $\varphi_{580} :\equiv (f(x) \in \mathbb{R} \to \alpha \in \mathbb{R} \to \lvert f(x) - \alpha \rvert = \lvert \alpha - f(x) \rvert) \to (\mathrm{RFn}(f) \to c \in \mathbb{R} \to \mathrm{PNbhd}(f,c) \to \varphi_{296}) \to \varphi_{376}$
\item $\varphi_{581} :\equiv r \in \mathbb{R} \land 0 < r \land \varphi_{16} \to \varphi_{76} \to \mathrm{Conv}(f \circ p,\alpha) \to \varepsilon \in \mathbb{R} \to 0 < \varepsilon \to \varphi_{486}$
\item $\varphi_{582} :\equiv \varphi_{362} \land \mathrm{Conv}(f \circ p,\alpha) \land \varepsilon \in \mathbb{R} \land 0 < \varepsilon \land (\delta \in \mathbb{R} \land 0 < \delta \land \forall x\,(\varphi_{176})) \land (N_{0} \in \mathbb{N} \land \varphi_{77}) \land \neg \exists \delta\,(\delta \in \mathbb{R} \land 0 < \delta \land \varphi_{69})$
\item $\varphi_{583} :\equiv \forall \varepsilon \in \mathbb{R}\,(0 < \varepsilon \to \exists \delta\,(\delta \in \mathbb{R} \land 0 < \delta \land \forall x\,(\varphi_{161}))) \to r \in \mathbb{R} \land 0 < r \land \varphi_{16} \to \varphi_{76} \to \mathrm{Conv}(f \circ p,\alpha) \to \varphi_{387}$
\item $\varphi_{584} :\equiv \forall \varepsilon \in \mathbb{R}\,(0 < \varepsilon \to \exists \delta\,(\delta \in \mathbb{R} \land 0 < \delta \land \forall x\,(\varphi_{161}))) \to r \in \mathbb{R} \land 0 < r \land \varphi_{16} \to \varphi_{76} \to \mathrm{Conv}(f \circ p,\alpha) \to \varphi_{388}$
\item $\varphi_{585} :\equiv \forall \varepsilon \in \mathbb{R}\,(0 < \varepsilon \to \exists \delta\,(\delta \in \mathbb{R} \land 0 < \delta \land \forall x\,(\varphi_{161}))) \to r \in \mathbb{R} \land 0 < r \land \varphi_{16} \to \varphi_{76} \to \varphi_{419}$
\item $\varphi_{586} :\equiv \forall \varepsilon \in \mathbb{R}\,(0 < \varepsilon \to \exists \delta\,(\delta \in \mathbb{R} \land 0 < \delta \land \forall x\,(\varphi_{161}))) \to r \in \mathbb{R} \land 0 < r \land \varphi_{16} \to \varphi_{76} \to \mathrm{Conv}(f \circ p,\alpha) \to \varphi_{392}$
\item $\varphi_{587} :\equiv \forall \varepsilon \in \mathbb{R}\,(0 < \varepsilon \to \exists \delta\,(\delta \in \mathbb{R} \land 0 < \delta \land \forall x\,(\varphi_{161}))) \to r \in \mathbb{R} \land 0 < r \land \varphi_{16} \to \varphi_{76} \to \mathrm{Conv}(f \circ p,\alpha) \to \varphi_{393}$
\item $\varphi_{588} :\equiv \forall \varepsilon \in \mathbb{R}\,(0 < \varepsilon \to \exists \delta\,(\delta \in \mathbb{R} \land 0 < \delta \land \forall x\,(\varphi_{161}))) \to r \in \mathbb{R} \land 0 < r \land \varphi_{16} \to \varphi_{76} \to \varepsilon \in \mathbb{R} \to 0 < \varepsilon \to \varphi_{391}$
\item $\varphi_{589} :\equiv \forall \varepsilon \in \mathbb{R}\,(0 < \varepsilon \to \exists \delta\,(\delta \in \mathbb{R} \land 0 < \delta \land \forall x\,(\varphi_{161}))) \to r \in \mathbb{R} \land 0 < r \land \varphi_{16} \to \varphi_{76} \to \mathrm{Conv}(f \circ p,\alpha) \to \varphi_{395}$
\item $\varphi_{590} :\equiv (\mathrm{RFn}(f) \to c \in \mathbb{R} \to \mathrm{PNbhd}(f,c) \to \varphi_{351}) \to \varphi_{320} \to \neg \varphi_{76}$
\item $\varphi_{591} :\equiv \varphi_{362} \land \neg \exists \alpha\,(\mathrm{FLim}(f,c,\alpha)) \to \forall \alpha\,(\varphi_{362} \land \neg \exists \alpha\,(\mathrm{FLim}(f,c,\alpha)))$
\item $\varphi_{592} :\equiv \forall \varepsilon \in \mathbb{R}\,(0 < \varepsilon \to \exists \delta\,(\delta \in \mathbb{R} \land 0 < \delta \land \forall x\,(\varphi_{161}))) \to r \in \mathbb{R} \land 0 < r \land \varphi_{16} \to \varphi_{76} \to \mathrm{Conv}(f \circ p,\alpha) \to \varphi_{405}$
\item $\varphi_{593} :\equiv \forall \varepsilon \in \mathbb{R}\,(0 < \varepsilon \to \exists \delta\,(\delta \in \mathbb{R} \land 0 < \delta \land \forall x\,(\varphi_{161}))) \to r \in \mathbb{R} \land 0 < r \land \varphi_{16} \to \varphi_{76} \to \mathrm{Conv}(f \circ p,\alpha) \to \varphi_{406}$
\item $\varphi_{594} :\equiv \forall \varepsilon \in \mathbb{R}\,(0 < \varepsilon \to \exists \delta\,(\delta \in \mathbb{R} \land 0 < \delta \land \forall x\,(\varphi_{161}))) \to r \in \mathbb{R} \land 0 < r \land \varphi_{16} \to \varphi_{76} \to \mathrm{Conv}(f \circ p,\alpha) \to \varphi_{409}$
\item $\varphi_{595} :\equiv \forall \varepsilon \in \mathbb{R}\,(0 < \varepsilon \to \exists \delta\,(\delta \in \mathbb{R} \land 0 < \delta \land \forall x\,(\varphi_{161}))) \to r \in \mathbb{R} \land 0 < r \land \varphi_{16} \to \varphi_{76} \to \mathrm{Conv}(f \circ p,\alpha) \to \varphi_{411}$
\item $\varphi_{596} :\equiv \forall \varepsilon \in \mathbb{R}\,(0 < \varepsilon \to \exists \delta\,(\delta \in \mathbb{R} \land 0 < \delta \land \forall x\,(\varphi_{161}))) \to r \in \mathbb{R} \land 0 < r \land \varphi_{16} \to \varphi_{76} \to \mathrm{Conv}(f \circ p,\alpha) \to \varphi_{412}$
\item $\varphi_{597} :\equiv \forall \varepsilon \in \mathbb{R}\,(0 < \varepsilon \to \exists \delta\,(\delta \in \mathbb{R} \land 0 < \delta \land \forall x\,(\varphi_{161}))) \to r \in \mathbb{R} \land 0 < r \land \varphi_{16} \to \varphi_{76} \to \mathrm{Conv}(f \circ p,\alpha) \to \varphi_{415}$
\item $\varphi_{598} :\equiv \forall \varepsilon \in \mathbb{R}\,(0 < \varepsilon \to \exists \delta\,(\delta \in \mathbb{R} \land 0 < \delta \land \forall x\,(\varphi_{161}))) \to r \in \mathbb{R} \land 0 < r \land \varphi_{16} \to \varphi_{76} \to \mathrm{Conv}(f \circ p,\alpha) \to \varphi_{416}$
\item $\varphi_{599} :\equiv \varphi_{362} \land \mathrm{Conv}(f \circ p,\alpha) \to \forall \varepsilon\,(\varphi_{362} \land \mathrm{Conv}(f \circ p,\alpha))$
\item $\varphi_{600} :\equiv \forall \varepsilon \in \mathbb{R}\,(0 < \varepsilon \to \exists \delta\,(\delta \in \mathbb{R} \land 0 < \delta \land \forall x\,(\varphi_{161}))) \to r \in \mathbb{R} \land 0 < r \land \varphi_{16} \to \varphi_{76} \to \mathrm{Conv}(f \circ p,\alpha) \to \varphi_{417}$
\item $\varphi_{601} :\equiv (\mathrm{RFn}(f) \to c \in \mathbb{R} \to \mathrm{PNbhd}(f,c) \to \varphi_{355}) \to \mathrm{RFn}(f) \to c \in \mathbb{R} \to \mathrm{PNbhd}(f,c) \to \varphi_{359}$
\item $\varphi_{602} :\equiv \forall \varepsilon \in \mathbb{R}\,(0 < \varepsilon \to \exists \delta\,(\delta \in \mathbb{R} \land 0 < \delta \land \forall x\,(\varphi_{161}))) \to r \in \mathbb{R} \land 0 < r \land \varphi_{16} \to \varphi_{76} \to \mathrm{Conv}(f \circ p,\alpha) \to \varphi_{420}$
\item $\varphi_{603} :\equiv (f(p(m)) = f \circ p(m) \to \varphi_{115}) \to \mathrm{RFn}(f) \to c \in \mathbb{R} \to \mathrm{PNbhd}(f,c) \to \varphi_{536}$
\item $\varphi_{604} :\equiv (\forall a\,(\forall p\,(\neg \varphi_{76}) \to \varphi_{78}) \to \forall a\,\forall p\,(\neg \varphi_{76}) \to \forall a\,(\varphi_{78})) \to \varphi_{452}$
\item $\varphi_{605} :\equiv (\mathrm{RFn}(f) \to c \in \mathbb{R} \to \mathrm{PNbhd}(f,c) \to \varphi_{361}) \to \varphi_{362} \to n \in \mathbb{N} \to p(n) \in \mathrm{dom}(f)$
\item $\varphi_{606} :\equiv N_{1} \leq N_{0} \cup N_{1} \to \mathrm{RFn}(f) \to c \in \mathbb{R} \to \mathrm{PNbhd}(f,c) \to \varphi_{578}$
\item $\varphi_{607} :\equiv \forall \varepsilon \in \mathbb{R}\,(0 < \varepsilon \to \exists \delta\,(\delta \in \mathbb{R} \land 0 < \delta \land \forall x\,(\varphi_{161}))) \to \varphi_{509}$
\item $\varphi_{608} :\equiv \forall \varepsilon \in \mathbb{R}\,(0 < \varepsilon \to \exists \delta\,(\delta \in \mathbb{R} \land 0 < \delta \land \forall x\,(\varphi_{161}))) \to \varphi_{510}$
\item $\varphi_{609} :\equiv \forall \varepsilon \in \mathbb{R}\,(0 < \varepsilon \to \exists \delta\,(\delta \in \mathbb{R} \land 0 < \delta \land \forall x\,(\varphi_{161}))) \to \varphi_{513}$
\item $\varphi_{610} :\equiv \forall n\,(\varphi_{362} \to n \in \mathbb{N} \to p(n) \in \mathrm{dom}(f)) \to \forall n\,(\varphi_{362}) \to \forall n \in \mathbb{N}\,(p(n) \in \mathrm{dom}(f))$
\item $\varphi_{611} :\equiv \forall \varepsilon \in \mathbb{R}\,(0 < \varepsilon \to \exists \delta\,(\delta \in \mathbb{R} \land 0 < \delta \land \forall x\,(\varphi_{161}))) \to \varphi_{514}$
\item $\varphi_{612} :\equiv (f(p(n)) = f \circ p(n) \to \varphi_{125}) \to \mathrm{RFn}(f) \to c \in \mathbb{R} \to \mathrm{PNbhd}(f,c) \to \varphi_{539}$
\item $\varphi_{613} :\equiv \forall \varepsilon \in \mathbb{R}\,(0 < \varepsilon \to \exists \delta\,(\delta \in \mathbb{R} \land 0 < \delta \land \forall x\,(\varphi_{161}))) \to \varphi_{516}$
\item $\varphi_{614} :\equiv (\mathrm{RFn}(f) \to c \in \mathbb{R} \to \mathrm{PNbhd}(f,c) \to \varphi_{336}) \to \mathrm{RFn}(f) \to c \in \mathbb{R} \to \mathrm{PNbhd}(f,c) \to \varphi_{390}$
\item $\varphi_{615} :\equiv (\mathrm{RFn}(f) \to c \in \mathbb{R} \to \mathrm{PNbhd}(f,c) \to \varphi_{371}) \to \mathrm{RFn}(f) \to c \in \mathbb{R} \to \mathrm{PNbhd}(f,c) \to \varphi_{365}$
\item $\varphi_{616} :\equiv \forall \varepsilon \in \mathbb{R}\,(0 < \varepsilon \to \exists \delta\,(\delta \in \mathbb{R} \land 0 < \delta \land \forall x\,(\varphi_{161}))) \to \varphi_{526}$
\item $\varphi_{617} :\equiv (\mathrm{RFn}(f) \to c \in \mathbb{R} \to \mathrm{PNbhd}(f,c) \to \varphi_{384}) \to \mathrm{RFn}(f) \to c \in \mathbb{R} \to \mathrm{PNbhd}(f,c) \to \varphi_{352}$
\item $\varphi_{618} :\equiv (\mathrm{RFn}(f) \to c \in \mathbb{R} \to \mathrm{PNbhd}(f,c) \to \varphi_{401}) \to \mathrm{RFn}(f) \to c \in \mathbb{R} \to \mathrm{PNbhd}(f,c) \to \varphi_{343}$
\item $\varphi_{619} :\equiv N_{0} \leq N_{0} \cup N_{1} \to \mathrm{RFn}(f) \to c \in \mathbb{R} \to \mathrm{PNbhd}(f,c) \to \varphi_{593}$
\item $\varphi_{620} :\equiv (\mathrm{RFn}(f) \to c \in \mathbb{R} \to \mathrm{PNbhd}(f,c) \to \varphi_{379}) \to \mathrm{RFn}(f) \to c \in \mathbb{R} \to \mathrm{PNbhd}(f,c) \to \varphi_{361}$
\item $\varphi_{621} :\equiv (\mathrm{RFn}(f) \to c \in \mathbb{R} \to \mathrm{PNbhd}(f,c) \to \varphi_{401}) \to \mathrm{RFn}(f) \to c \in \mathbb{R} \to \mathrm{PNbhd}(f,c) \to \varphi_{349}$
\item $\varphi_{622} :\equiv (\mathrm{RFn}(f) \to c \in \mathbb{R} \to \mathrm{PNbhd}(f,c) \to \varphi_{363}) \to \mathrm{RFn}(f) \to c \in \mathbb{R} \to \mathrm{PNbhd}(f,c) \to \varphi_{384}$
\item $\varphi_{623} :\equiv (\varphi_{16} \to p(n) \in \mathbb{R} \to \varphi_{27}) \to \mathrm{RFn}(f) \to c \in \mathbb{R} \to \mathrm{PNbhd}(f,c) \to \varphi_{552}$
\item $\varphi_{624} :\equiv (\varphi_{16} \to p(m) \in \mathbb{R} \to \varphi_{30}) \to \mathrm{RFn}(f) \to c \in \mathbb{R} \to \mathrm{PNbhd}(f,c) \to \varphi_{553}$
\item $\varphi_{625} :\equiv (\mathrm{RFn}(f) \to c \in \mathbb{R} \to \mathrm{PNbhd}(f,c) \to \varphi_{377}) \to \mathrm{RFn}(f) \to c \in \mathbb{R} \to \mathrm{PNbhd}(f,c) \to \varphi_{382}$
\item $\varphi_{626} :\equiv (\mathrm{Cgt}(f \circ p) \leftrightarrow \mathrm{Cauchy}(f \circ p)) \to (\mathrm{RFn}(f) \to c \in \mathbb{R} \to \mathrm{PNbhd}(f,c) \to \varphi_{353}) \to \varphi_{427}$
\item $\varphi_{627} :\equiv (\mathrm{RFn}(f) \to c \in \mathbb{R} \to \mathrm{PNbhd}(f,c) \to \varphi_{401}) \to \mathrm{RFn}(f) \to c \in \mathbb{R} \to \mathrm{PNbhd}(f,c) \to \varphi_{381}$
\item $\varphi_{628} :\equiv (\mathrm{RFn}(f) \to c \in \mathbb{R} \to \mathrm{PNbhd}(f,c) \to \varphi_{448}) \to \mathrm{RFn}(f) \to c \in \mathbb{R} \to \mathrm{PNbhd}(f,c) \to \varphi_{353}$
\item $\varphi_{629} :\equiv \forall \varepsilon \in \mathbb{R}\,(0 < \varepsilon \to \exists \delta\,(\delta \in \mathbb{R} \land 0 < \delta \land \forall x\,(\varphi_{161}))) \to r \in \mathbb{R} \land 0 < r \land \varphi_{16} \to \varphi_{76} \to \varphi_{483}$
\item $\varphi_{630} :\equiv \forall \varepsilon \in \mathbb{R}\,(0 < \varepsilon \to \exists \delta\,(\delta \in \mathbb{R} \land 0 < \delta \land \forall x\,(\varphi_{161}))) \to r \in \mathbb{R} \land 0 < r \land \varphi_{16} \to \varphi_{76} \to \mathrm{Conv}(f \circ p,\alpha) \to \varphi_{472}$
\item $\varphi_{631} :\equiv (\varepsilon \cdot 2^{-1} + \varepsilon \cdot 2^{-1} = \varepsilon \to \varphi_{66}) \to \mathrm{RFn}(f) \to c \in \mathbb{R} \to \mathrm{PNbhd}(f,c) \to \varphi_{592}$
\item $\varphi_{632} :\equiv (\forall n \in \mathbb{N}\,(f \circ p(n) = f(p(n))) \to \varphi_{46}) \to \mathrm{RFn}(f) \to c \in \mathbb{R} \to \mathrm{PNbhd}(f,c) \to \varphi_{589}$
\item $\varphi_{633} :\equiv (\varepsilon \in \mathbb{R} \to 2^{-1} \in \mathbb{R} \to \varepsilon \cdot 2^{-1} \in \mathbb{R}) \to (\mathrm{RFn}(f) \to c \in \mathbb{R} \to \mathrm{PNbhd}(f,c) \to \varphi_{373}) \to \varphi_{433}$
\item $\varphi_{634} :\equiv (\mathrm{RFn}(f) \to c \in \mathbb{R} \to \mathrm{PNbhd}(f,c) \to \varphi_{459}) \to \mathrm{RFn}(f) \to c \in \mathbb{R} \to \mathrm{PNbhd}(f,c) \to \varphi_{366}$
\item $\varphi_{635} :\equiv (\forall N_{1}\,(\neg \varphi_{67}) \to \forall N\,\forall N_{1}\,(\neg \varphi_{67})) \to \mathrm{RFn}(f) \to c \in \mathbb{R} \to \mathrm{PNbhd}(f,c) \to \varphi_{567}$
\item $\varphi_{636} :\equiv (f \circ p(N_{0} \cup N_{1}) = f(p(N_{0} \cup N_{1})) \to \varphi_{129}) \to \mathrm{RFn}(f) \to c \in \mathbb{R} \to \mathrm{PNbhd}(f,c) \to \varphi_{576}$
\item $\varphi_{637} :\equiv \forall \varepsilon \in \mathbb{R}\,(0 < \varepsilon \to \exists \delta\,(\delta \in \mathbb{R} \land 0 < \delta \land \forall x\,(\varphi_{161}))) \to \varphi_{581}$
\item $\varphi_{638} :\equiv (\mathrm{RFn}(f) \to c \in \mathbb{R} \to \mathrm{PNbhd}(f,c) \to \varphi_{364}) \to \mathrm{RFn}(f) \to c \in \mathbb{R} \to \mathrm{PNbhd}(f,c) \to \varphi_{466}$
\item $\varphi_{639} :\equiv (\mathrm{RFn}(f) \to \varphi_{310}) \to (\mathrm{RFn}(f) \to \varphi_{321}) \to \mathrm{RFn}(f) \to c \in \mathbb{R} \to \varphi_{328}$
\item $\varphi_{640} :\equiv (\forall n \in \mathbb{N}\,(N_{1} \leq n \to \lvert f \circ p(n) - \alpha \rvert < \varepsilon \cdot 2^{-1}) \to \varphi_{96}) \to \mathrm{RFn}(f) \to c \in \mathbb{R} \to \mathrm{PNbhd}(f,c) \to \varphi_{602}$
\item $\varphi_{641} :\equiv (\mathrm{RFn}(f) \to c \in \mathbb{R} \to \mathrm{PNbhd}(f,c) \to \varphi_{381}) \to \varphi_{146} \to \mathrm{RFn}(f) \to c \in \mathbb{R} \to \mathrm{PNbhd}(f,c) \to \varphi_{379}$
\item $\varphi_{642} :\equiv (\mathrm{RFn}(f) \to c \in \mathbb{R} \to \mathrm{PNbhd}(f,c) \to \varphi_{381}) \to \varphi_{149} \to \mathrm{RFn}(f) \to c \in \mathbb{R} \to \mathrm{PNbhd}(f,c) \to \varphi_{380}$
\item $\varphi_{643} :\equiv (\mathrm{RFn}(f) \to c \in \mathbb{R} \to \mathrm{PNbhd}(f,c) \to \varphi_{447}) \to \varphi_{362} \to \varepsilon \in \mathbb{R} \to 0 < \varepsilon \to \exists N \in \mathbb{N}\,\forall m\,\forall n\,(\varphi_{90})$
\item $\varphi_{644} :\equiv \varphi_{362} \land \varepsilon \in \mathbb{R} \land 0 < \varepsilon \land \neg \exists N \in \mathbb{N}\,\forall m\,\forall n\,(\varphi_{90}) \to \forall \delta\,(\varphi_{362} \land \varepsilon \in \mathbb{R} \land 0 < \varepsilon \land \neg \exists N \in \mathbb{N}\,\forall m\,\forall n\,(\varphi_{90}))$
\item $\varphi_{645} :\equiv (\mathrm{RFn}(f) \to c \in \mathbb{R} \to \varphi_{328}) \to \varphi_{318} \to \mathrm{RFn}(f) \to c \in \mathbb{R} \to \mathrm{PNbhd}(f,c) \to \varphi_{325}$
\item $\varphi_{646} :\equiv \forall \varepsilon\,(\varphi_{362} \to \varepsilon \in \mathbb{R} \to 0 < \varepsilon \to \exists N \in \mathbb{N}\,\forall m\,\forall n\,(\varphi_{90})) \to \forall \varepsilon\,(\varphi_{362}) \to \forall \varepsilon \in \mathbb{R}\,(0 < \varepsilon \to \exists N \in \mathbb{N}\,\forall m\,\forall n\,(\varphi_{90}))$
\item $\varphi_{647} :\equiv (\mathrm{RFn}(f) \to c \in \mathbb{R} \to c \in \mathbb{R}) \to (\mathrm{RFn}(f) \to c \in \mathbb{R} \to \mathrm{PNbhd}(f,c) \to \varphi_{301}) \to \varphi_{549}$
\item $\varphi_{648} :\equiv (\mathrm{RFn}(f) \to c \in \mathbb{R} \to \mathrm{PNbhd}(f,c) \to \varphi_{382}) \to \mathrm{RFn}(f) \to c \in \mathbb{R} \to \mathrm{PNbhd}(f,c) \to \varphi_{478}$
\item $\varphi_{649} :\equiv (\varphi_{16} \to p(N_{0} \cup N_{1}) \in \mathbb{R} \to \varphi_{112}) \to \mathrm{RFn}(f) \to c \in \mathbb{R} \to \mathrm{PNbhd}(f,c) \to \varphi_{616}$
\item $\varphi_{650} :\equiv (\varphi_{77} \to N_{0} \cup N_{1} \in \mathbb{N} \to N_{0} \leq N_{0} \cup N_{1} \to \varphi_{61}) \to \mathrm{RFn}(f) \to c \in \mathbb{R} \to \mathrm{PNbhd}(f,c) \to \varphi_{613}$
\item $\varphi_{651} :\equiv (\mathrm{RFn}(f) \to c \in \mathbb{R} \to \mathrm{PNbhd}(f,c) \to \varphi_{212}) \to (\mathrm{RFn}(f) \to \varphi_{272}) \to \varphi_{493}$
\item $\varphi_{652} :\equiv (\mathrm{RFn}(f) \to c \in \mathbb{R} \to \mathrm{PNbhd}(f,c) \to \varphi_{336}) \to \mathrm{RFn}(f) \to c \in \mathbb{R} \to \mathrm{PNbhd}(f,c) \to \varphi_{544}$
\item $\varphi_{653} :\equiv (\forall x\,(\varphi_{161}) \to \varphi_{199}) \to \mathrm{RFn}(f) \to c \in \mathbb{R} \to \mathrm{PNbhd}(f,c) \to \varphi_{573}$
\item $\varphi_{654} :\equiv (\mathrm{RFn}(f) \to c \in \mathbb{R} \to \mathrm{PNbhd}(f,c) \to \varphi_{308}) \to (\mathrm{RFn}(f) \to \varphi_{331}) \to \varphi_{421}$
\item $\varphi_{655} :\equiv (\mathrm{RFn}(f) \to c \in \mathbb{R} \to \mathrm{PNbhd}(f,c) \to \varphi_{384}) \to \mathrm{RFn}(f) \to c \in \mathbb{R} \to \mathrm{PNbhd}(f,c) \to \varphi_{535}$
\item $\varphi_{656} :\equiv (\mathrm{RFn}(f) \to c \in \mathbb{R} \to \mathrm{PNbhd}(f,c) \to \varphi_{223}) \to (\mathrm{RFn}(f) \to c \in \mathbb{R} \to \mathrm{PNbhd}(f,c) \to \varphi_{278}) \to \varphi_{506}$
\item $\varphi_{657} :\equiv (\varphi_{199} \to p(m) \in \mathrm{dom}(f) \to \varphi_{202}) \to \mathrm{RFn}(f) \to c \in \mathbb{R} \to \mathrm{PNbhd}(f,c) \to \varphi_{588}$
\item $\varphi_{658} :\equiv (\varphi_{362} \to \forall n\,(\varphi_{362})) \to \mathrm{RFn}(f) \to c \in \mathbb{R} \to \mathrm{PNbhd}(f,c) \to \varphi_{363}$
\item $\varphi_{659} :\equiv (\mathrm{RFn}(f) \to c \in \mathbb{R} \to \mathrm{PNbhd}(f,c) \to \varphi_{357}) \to \varphi_{354} \to \mathrm{RFn}(f) \to c \in \mathbb{R} \to \mathrm{PNbhd}(f,c) \to \varphi_{354}$
\item $\varphi_{660} :\equiv (\mathrm{RFn}(f) \to c \in \mathbb{R} \to \mathrm{PNbhd}(f,c) \to \varphi_{469}) \to \mathrm{RFn}(f) \to c \in \mathbb{R} \to \mathrm{PNbhd}(f,c) \to \varphi_{488}$
\item $\varphi_{661} :\equiv (\mathrm{RFn}(f) \to c \in \mathbb{R} \to \mathrm{PNbhd}(f,c) \to \varphi_{336}) \to \mathrm{RFn}(f) \to c \in \mathbb{R} \to \mathrm{PNbhd}(f,c) \to \varphi_{583}$
\item $\varphi_{662} :\equiv (\mathrm{RFn}(f) \to c \in \mathbb{R} \to \mathrm{PNbhd}(f,c) \to \varphi_{382}) \to \mathrm{RFn}(f) \to c \in \mathbb{R} \to \mathrm{PNbhd}(f,c) \to \varphi_{558}$
\item $\varphi_{663} :\equiv (\forall x\,(\varphi_{176}) \to \varphi_{176}) \to \mathrm{RFn}(f) \to c \in \mathbb{R} \to \mathrm{PNbhd}(f,c) \to \varphi_{629}$
\item $\varphi_{664} :\equiv (\mathrm{RFn}(f) \to c \in \mathbb{R} \to \mathrm{PNbhd}(f,c) \to \varphi_{357}) \to \varphi_{372} \to \mathrm{RFn}(f) \to c \in \mathbb{R} \to \mathrm{PNbhd}(f,c) \to \varphi_{372}$
\item $\varphi_{665} :\equiv (\mathrm{Seq}(p) \to n \in \mathbb{N} \to p(n) \in \mathbb{R}) \to (\mathrm{RFn}(f) \to c \in \mathbb{R} \to \mathrm{PNbhd}(f,c) \to \varphi_{343}) \to \varphi_{601}$
\item $\varphi_{666} :\equiv (\mathrm{RFn}(f) \to c \in \mathbb{R} \to \mathrm{PNbhd}(f,c) \to \varphi_{384}) \to \mathrm{RFn}(f) \to c \in \mathbb{R} \to \mathrm{PNbhd}(f,c) \to \varphi_{574}$
\item $\varphi_{667} :\equiv (\mathrm{RFn}(f) \to \mathrm{RFn}(f)) \to (\mathrm{RFn}(f) \to c \in \mathbb{R} \to \mathrm{PNbhd}(f,c) \to \varphi_{343}) \to \varphi_{622}$
\item $\varphi_{668} :\equiv (\mathrm{RFn}(f) \to c \in \mathbb{R} \to \mathrm{PNbhd}(f,c) \to \varphi_{373}) \to \mathrm{RFn}(f) \to c \in \mathbb{R} \to \mathrm{PNbhd}(f,c) \to \varphi_{600}$
\item $\varphi_{669} :\equiv (\forall N\,(\forall N_{1}\,(\neg \varphi_{67}) \to \varphi_{60}) \to \forall N\,\forall N_{1}\,(\neg \varphi_{67}) \to \forall N\,(\varphi_{60})) \to \varphi_{635}$
\item $\varphi_{670} :\equiv (\varphi_{362} \to \forall \varepsilon\,(\varphi_{362})) \to \mathrm{RFn}(f) \to c \in \mathbb{R} \to \mathrm{PNbhd}(f,c) \to \varphi_{448}$
\item $\varphi_{671} :\equiv (\mathrm{RFn}(f) \to c \in \mathbb{R} \to \mathrm{PNbhd}(f,c) \to \varphi_{484}) \to \mathrm{RFn}(f) \to c \in \mathbb{R} \to \mathrm{PNbhd}(f,c) \to \varphi_{522}$
\item $\varphi_{672} :\equiv (\mathrm{RFn}(f) \to c \in \mathbb{R} \to \mathrm{PNbhd}(f,c) \to \varphi_{546}) \to \mathrm{RFn}(f) \to c \in \mathbb{R} \to \mathrm{PNbhd}(f,c) \to \varphi_{469}$
\item $\varphi_{673} :\equiv (\mathrm{RFn}(f) \to c \in \mathbb{R} \to \mathrm{PNbhd}(f,c) \to \varphi_{546}) \to \mathrm{RFn}(f) \to c \in \mathbb{R} \to \mathrm{PNbhd}(f,c) \to \varphi_{470}$
\item $\varphi_{674} :\equiv (\mathrm{RFn}(f) \to c \in \mathbb{R} \to \mathrm{PNbhd}(f,c) \to \varphi_{507}) \to \mathrm{RFn}(f) \to c \in \mathbb{R} \to \mathrm{PNbhd}(f,c) \to \varphi_{491}$
\item $\varphi_{675} :\equiv (\mathrm{RFn}(f) \to c \in \mathbb{R} \to \mathrm{PNbhd}(f,c) \to \varphi_{494}) \to \mathrm{RFn}(f) \to c \in \mathbb{R} \to \mathrm{PNbhd}(f,c) \to \varphi_{503}$
\item $\varphi_{676} :\equiv (\mathrm{RFn}(f) \to c \in \mathbb{R} \to \mathrm{PNbhd}(f,c) \to \varphi_{497}) \to \mathrm{RFn}(f) \to c \in \mathbb{R} \to \mathrm{PNbhd}(f,c) \to \varphi_{502}$
\item $\varphi_{677} :\equiv (\mathrm{RFn}(f) \to c \in \mathbb{R} \to \mathrm{PNbhd}(f,c) \to \varphi_{302}) \to (\mathrm{RFn}(f) \to c \in \mathbb{R} \to \mathrm{PNbhd}(f,c) \to \varphi_{314}) \to \varphi_{543}$
\item $\varphi_{678} :\equiv (\mathrm{RFn}(f) \to c \in \mathbb{R} \to \mathrm{PNbhd}(f,c) \to \varphi_{494}) \to \mathrm{RFn}(f) \to c \in \mathbb{R} \to \mathrm{PNbhd}(f,c) \to \varphi_{517}$
\item $\varphi_{679} :\equiv (\mathrm{RFn}(f) \to c \in \mathbb{R} \to \mathrm{PNbhd}(f,c) \to \varphi_{498}) \to \mathrm{RFn}(f) \to c \in \mathbb{R} \to \mathrm{PNbhd}(f,c) \to \varphi_{512}$
\item $\varphi_{680} :\equiv (\mathrm{RFn}(f) \to c \in \mathbb{R} \to \mathrm{PNbhd}(f,c) \to \varphi_{499}) \to \mathrm{RFn}(f) \to c \in \mathbb{R} \to \mathrm{PNbhd}(f,c) \to \varphi_{511}$
\item $\varphi_{681} :\equiv (\mathrm{RFn}(f) \to c \in \mathbb{R} \to \mathrm{PNbhd}(f,c) \to \varphi_{497}) \to \mathrm{RFn}(f) \to c \in \mathbb{R} \to \mathrm{PNbhd}(f,c) \to \varphi_{519}$
\item $\varphi_{682} :\equiv (\mathrm{RFn}(f) \to c \in \mathbb{R} \to \mathrm{PNbhd}(f,c) \to \varphi_{494}) \to \mathrm{RFn}(f) \to c \in \mathbb{R} \to \mathrm{PNbhd}(f,c) \to \varphi_{534}$
\item $\varphi_{683} :\equiv (\varphi_{176} \to x \in \mathrm{dom}(f) \to p(N_{0} \cup N_{1}) \in \mathrm{dom}(f) \to \varphi_{196}) \to \mathrm{RFn}(f) \to c \in \mathbb{R} \to \mathrm{PNbhd}(f,c) \to \varphi_{637}$
\item $\varphi_{684} :\equiv (\mathrm{RFn}(f) \to c \in \mathbb{R} \to \mathrm{PNbhd}(f,c) \to \varphi_{497}) \to \mathrm{RFn}(f) \to c \in \mathbb{R} \to \mathrm{PNbhd}(f,c) \to \varphi_{533}$
\item $\varphi_{685} :\equiv (\varepsilon \in \mathbb{R} \to 0 < \varepsilon \to 0 < \varepsilon \cdot 2^{-1}) \to (\mathrm{RFn}(f) \to c \in \mathbb{R} \to \mathrm{PNbhd}(f,c) \to \varphi_{373}) \to \varphi_{625}$
\item $\varphi_{686} :\equiv (\mathrm{RFn}(f) \to c \in \mathbb{R} \to \mathrm{PNbhd}(f,c) \to \varphi_{557}) \to \mathrm{RFn}(f) \to c \in \mathbb{R} \to \mathrm{PNbhd}(f,c) \to \varphi_{484}$
\item $\varphi_{687} :\equiv (\mathrm{RFn}(f) \to c \in \mathbb{R} \to \mathrm{PNbhd}(f,c) \to \varphi_{507}) \to \mathrm{RFn}(f) \to c \in \mathbb{R} \to \mathrm{PNbhd}(f,c) \to \varphi_{527}$
\item $\varphi_{688} :\equiv (\mathrm{RFn}(f) \to c \in \mathbb{R} \to \mathrm{PNbhd}(f,c) \to \varphi_{557}) \to \mathrm{RFn}(f) \to c \in \mathbb{R} \to \mathrm{PNbhd}(f,c) \to \varphi_{485}$
\item $\varphi_{689} :\equiv (\mathrm{RFn}(f) \to c \in \mathbb{R} \to \mathrm{PNbhd}(f,c) \to \varphi_{533}) \to \mathrm{RFn}(f) \to c \in \mathbb{R} \to \mathrm{PNbhd}(f,c) \to \varphi_{504}$
\item $\varphi_{690} :\equiv (\mathrm{RFn}(f) \to c \in \mathbb{R} \to \mathrm{PNbhd}(f,c) \to \varphi_{534}) \to \mathrm{RFn}(f) \to c \in \mathbb{R} \to \mathrm{PNbhd}(f,c) \to \varphi_{505}$
\item $\varphi_{691} :\equiv (\mathrm{RFn}(f) \to c \in \mathbb{R} \to \mathrm{PNbhd}(f,c) \to \varphi_{511}) \to \mathrm{RFn}(f) \to c \in \mathbb{R} \to \mathrm{PNbhd}(f,c) \to \varphi_{529}$
\item $\varphi_{692} :\equiv (\mathrm{RFn}(f) \to c \in \mathbb{R} \to \mathrm{PNbhd}(f,c) \to \varphi_{532}) \to \varphi_{464} \to \neg (\delta \in \mathbb{R} \land 0 < \delta \land \forall x\,(\varphi_{176}))$
\item $\varphi_{693} :\equiv \forall \delta\,(\varphi_{464} \to \neg (\delta \in \mathbb{R} \land 0 < \delta \land \forall x\,(\varphi_{176}))) \to \forall \delta\,(\varphi_{464}) \to \forall \delta\,(\neg (\delta \in \mathbb{R} \land 0 < \delta \land \forall x\,(\varphi_{176})))$
\item $\varphi_{694} :\equiv (\mathrm{RFn}(f) \to c \in \mathbb{R} \to \mathrm{PNbhd}(f,c) \to \varphi_{540}) \to \varphi_{496} \to \forall n\,(\varphi_{90})$
\item $\varphi_{695} :\equiv (\mathrm{RFn}(f) \to c \in \mathbb{R} \to \mathrm{PNbhd}(f,c) \to \varphi_{542}) \to \mathrm{RFn}(f) \to c \in \mathbb{R} \to \mathrm{PNbhd}(f,c) \to \varphi_{545}$
\item $\varphi_{696} :\equiv (\mathrm{RFn}(f) \to c \in \mathbb{R} \to c \in \mathbb{R}) \to (\mathrm{RFn}(f) \to c \in \mathbb{R} \to \mathrm{PNbhd}(f,c) \to \varphi_{365}) \to \varphi_{2} \to \varphi_{634}$
\item $\varphi_{697} :\equiv \forall N\,(\forall N_{0}\,(\neg (N_{0} \in \mathbb{N} \land \varphi_{77})) \to \varphi_{110}) \to \varphi_{360} \to \varphi_{214} \to \mathrm{RFn}(f) \to c \in \mathbb{R} \to \mathrm{PNbhd}(f,c) \to \varphi_{532}$
\item $\varphi_{698} :\equiv (\mathrm{RFn}(f) \to c \in \mathbb{R} \to \mathrm{PNbhd}(f,c) \to \varphi_{527}) \to \varphi_{88} \to \mathrm{RFn}(f) \to c \in \mathbb{R} \to \mathrm{PNbhd}(f,c) \to \varphi_{524}$
\item $\varphi_{699} :\equiv (\mathrm{RFn}(f) \to c \in \mathbb{R} \to \mathrm{PNbhd}(f,c) \to \varphi_{527}) \to \varphi_{89} \to \mathrm{RFn}(f) \to c \in \mathbb{R} \to \mathrm{PNbhd}(f,c) \to \varphi_{525}$
\item $\varphi_{700} :\equiv (\forall \delta\,(\varphi_{341}) \to \varphi_{274} \to \forall \delta\,(\neg (\delta \in \mathbb{R} \land 0 < \delta \land \varphi_{99}))) \to \varphi_{456} \to \mathrm{RFn}(f) \to \varphi_{264}$
\item $\varphi_{701} :\equiv (\mathrm{RFn}(f) \to c \in \mathbb{R} \to \mathrm{PNbhd}(f,c) \to \varphi_{535}) \to \varphi_{119} \to \mathrm{RFn}(f) \to c \in \mathbb{R} \to \mathrm{PNbhd}(f,c) \to \varphi_{530}$
\item $\varphi_{702} :\equiv (\mathrm{RFn}(f) \to c \in \mathbb{R} \to \mathrm{PNbhd}(f,c) \to \varphi_{535}) \to \varphi_{120} \to \mathrm{RFn}(f) \to c \in \mathbb{R} \to \mathrm{PNbhd}(f,c) \to \varphi_{531}$
\item $\varphi_{703} :\equiv (\mathrm{RFn}(f) \to c \in \mathbb{R} \to \mathrm{PNbhd}(f,c) \to \varphi_{546}) \to \mathrm{RFn}(f) \to c \in \mathbb{R} \to \mathrm{PNbhd}(f,c) \to \varphi_{565}$
\item $\varphi_{704} :\equiv (\mathrm{RFn}(f) \to c \in \mathbb{R} \to \mathrm{PNbhd}(f,c) \to \varphi_{561}) \to \mathrm{RFn}(f) \to c \in \mathbb{R} \to \mathrm{PNbhd}(f,c) \to \varphi_{560}$
\item $\varphi_{705} :\equiv (\mathrm{RFn}(f) \to c \in \mathbb{R} \to \mathrm{PNbhd}(f,c) \to \varphi_{559}) \to \mathrm{RFn}(f) \to c \in \mathbb{R} \to \mathrm{PNbhd}(f,c) \to \varphi_{562}$
\item $\varphi_{706} :\equiv (\mathrm{RFn}(f) \to c \in \mathbb{R} \to \mathrm{PNbhd}(f,c) \to \varphi_{1}) \to \forall \alpha\,(\varphi_{262} \to \neg \mathrm{FLim}(f,c,\alpha)) \to \varphi_{462} \to \varphi_{555}$
\item $\varphi_{707} :\equiv (\mathrm{RFn}(f) \to c \in \mathbb{R} \to \mathrm{PNbhd}(f,c) \to \varphi_{562}) \to \mathrm{RFn}(f) \to c \in \mathbb{R} \to \mathrm{PNbhd}(f,c) \to \varphi_{566}$
\item $\varphi_{708} :\equiv (\mathrm{RFn}(f) \to c \in \mathbb{R} \to \mathrm{PNbhd}(f,c) \to \varphi_{567}) \to \varphi_{550} \to \neg (N_{0} \in \mathbb{N} \land \varphi_{77})$
\item $\varphi_{709} :\equiv \forall N_{0}\,(\varphi_{550} \to \neg (N_{0} \in \mathbb{N} \land \varphi_{77})) \to \forall N_{0}\,(\varphi_{550}) \to \forall N_{0}\,(\neg (N_{0} \in \mathbb{N} \land \varphi_{77}))$
\item $\varphi_{710} :\equiv (\mathrm{RFn}(f) \to c \in \mathbb{R} \to \mathrm{PNbhd}(f,c) \to \varphi_{562}) \to \mathrm{RFn}(f) \to c \in \mathbb{R} \to \mathrm{PNbhd}(f,c) \to \varphi_{577}$
\item $\varphi_{711} :\equiv (\varphi_{464} \to \forall \delta\,(\varphi_{464})) \to \mathrm{RFn}(f) \to c \in \mathbb{R} \to \mathrm{PNbhd}(f,c) \to \varphi_{457}$
\item $\varphi_{712} :\equiv (\mathrm{RFn}(f) \to c \in \mathbb{R} \to \mathrm{PNbhd}(f,c) \to \varphi_{562}) \to \mathrm{RFn}(f) \to c \in \mathbb{R} \to \mathrm{PNbhd}(f,c) \to \varphi_{584}$
\item $\varphi_{713} :\equiv (\mathrm{RFn}(f) \to c \in \mathbb{R} \to \mathrm{PNbhd}(f,c) \to \varphi_{594}) \to \mathrm{RFn}(f) \to c \in \mathbb{R} \to \mathrm{PNbhd}(f,c) \to \varphi_{559}$
\item $\varphi_{714} :\equiv (\mathrm{RFn}(f) \to c \in \mathbb{R} \to \mathrm{PNbhd}(f,c) \to \varphi_{584}) \to \mathrm{RFn}(f) \to c \in \mathbb{R} \to \mathrm{PNbhd}(f,c) \to \varphi_{569}$
\item $\varphi_{715} :\equiv (\mathrm{RFn}(f) \to c \in \mathbb{R} \to \mathrm{PNbhd}(f,c) \to \varphi_{561}) \to \mathrm{RFn}(f) \to c \in \mathbb{R} \to \mathrm{PNbhd}(f,c) \to \varphi_{596}$
\item $\varphi_{716} :\equiv (\mathrm{RFn}(f) \to c \in \mathbb{R} \to \mathrm{PNbhd}(f,c) \to \varphi_{568}) \to \mathrm{RFn}(f) \to c \in \mathbb{R} \to \mathrm{PNbhd}(f,c) \to \varphi_{587}$
\item $\varphi_{717} :\equiv (\mathrm{RFn}(f) \to c \in \mathbb{R} \to \mathrm{PNbhd}(f,c) \to \varphi_{569}) \to \mathrm{RFn}(f) \to c \in \mathbb{R} \to \mathrm{PNbhd}(f,c) \to \varphi_{598}$
\item $\varphi_{718} :\equiv (\mathrm{RFn}(f) \to c \in \mathbb{R} \to \mathrm{PNbhd}(f,c) \to \varphi_{545}) \to \varphi_{225} \to \mathrm{RFn}(f) \to c \in \mathbb{R} \to \mathrm{PNbhd}(f,c) \to \varphi_{547}$
\item $\varphi_{719} :\equiv (\mathrm{RFn}(f) \to c \in \mathbb{R} \to \mathrm{PNbhd}(f,c) \to \varphi_{594}) \to \mathrm{RFn}(f) \to c \in \mathbb{R} \to \mathrm{PNbhd}(f,c) \to \varphi_{585}$
\item $\varphi_{720} :\equiv (\mathrm{RFn}(f) \to c \in \mathbb{R} \to \mathrm{PNbhd}(f,c) \to \varphi_{576}) \to \mathrm{RFn}(f) \to c \in \mathbb{R} \to \mathrm{PNbhd}(f,c) \to \varphi_{607}$
\item $\varphi_{721} :\equiv (\mathrm{RFn}(f) \to x \in \mathrm{dom}(f) \to f(x) \in \mathbb{R}) \to (\mathrm{RFn}(f) \to \mathrm{RFn}(f)) \to \varphi_{716}$
\item $\varphi_{722} :\equiv (\mathrm{RFn}(f) \to c \in \mathbb{R} \to \mathrm{PNbhd}(f,c) \to \varphi_{593}) \to \mathrm{RFn}(f) \to c \in \mathbb{R} \to \mathrm{PNbhd}(f,c) \to \varphi_{597}$
\item $\varphi_{723} :\equiv (\mathrm{RFn}(f) \to c \in \mathbb{R} \to \mathrm{PNbhd}(f,c) \to \varphi_{595}) \to \mathrm{RFn}(f) \to c \in \mathbb{R} \to \mathrm{PNbhd}(f,c) \to \varphi_{608}$
\item $\varphi_{724} :\equiv (\mathrm{RFn}(f) \to c \in \mathbb{R} \to \mathrm{PNbhd}(f,c) \to \varphi_{597}) \to \mathrm{RFn}(f) \to c \in \mathbb{R} \to \mathrm{PNbhd}(f,c) \to \varphi_{611}$
\item $\varphi_{725} :\equiv (\mathrm{RFn}(f) \to c \in \mathbb{R} \to \mathrm{PNbhd}(f,c) \to \varphi_{557}) \to \mathrm{RFn}(f) \to c \in \mathbb{R} \to \mathrm{PNbhd}(f,c) \to \varphi_{630}$
\item $\varphi_{726} :\equiv \forall x\,(\varphi_{19} \land \varepsilon \in \mathbb{R} \land 0 < \varepsilon \land (\delta \in \mathbb{R} \land 0 < \delta \land \varphi_{99}) \to \varphi_{161}) \to \varphi_{528} \to \varphi_{369} \to \mathrm{RFn}(f) \to \varphi_{321}$
\item $\varphi_{727} :\equiv (\mathrm{RFn}(f) \to c \in \mathbb{R} \to \mathrm{PNbhd}(f,c) \to \varphi_{609}) \to \varphi_{582} \to \neg \varphi_{67}$
\item $\varphi_{728} :\equiv (\mathrm{Seq}(p) \to m \in \mathbb{N} \to p(m) \in \mathbb{R}) \to (\mathrm{RFn}(f) \to c \in \mathbb{R} \to \mathrm{PNbhd}(f,c) \to \varphi_{343}) \to \varphi_{675}$
\item $\varphi_{729} :\equiv (\mathrm{RFn}(f) \to c \in \mathbb{R} \to \mathrm{PNbhd}(f,c) \to \varphi_{422}) \to (\mathrm{RFn}(f) \to c \in \mathbb{R} \to \mathrm{PNbhd}(f,c) \to \varphi_{359}) \to \varphi_{620}$
\item $\varphi_{730} :\equiv (\mathrm{Seq}(p) \to n \in \mathbb{N} \to p(n) \in \mathbb{R}) \to (\mathrm{RFn}(f) \to c \in \mathbb{R} \to \mathrm{PNbhd}(f,c) \to \varphi_{343}) \to \varphi_{676}$
\item $\varphi_{731} :\equiv (\mathrm{RFn}(f) \to c \in \mathbb{R} \to \mathrm{PNbhd}(f,c) \to \varphi_{354}) \to (\mathrm{RFn}(f) \to c \in \mathbb{R} \to \mathrm{PNbhd}(f,c) \to \varphi_{358}) \to \varphi_{638}$
\item $\varphi_{732} :\equiv (\mathrm{RFn}(f) \to c \in \mathbb{R} \to \mathrm{PNbhd}(f,c) \to \varphi_{372}) \to (\mathrm{RFn}(f) \to c \in \mathbb{R} \to \mathrm{PNbhd}(f,c) \to \varphi_{385}) \to \varphi_{648}$
\item $\varphi_{733} :\equiv (\mathrm{RFn}(f) \to c \in \mathbb{R} \to \mathrm{PNbhd}(f,c) \to \varphi_{608}) \to \varphi_{224} \to \mathrm{RFn}(f) \to c \in \mathbb{R} \to \mathrm{PNbhd}(f,c) \to \varphi_{609}$
\item $\varphi_{734} :\equiv (\varphi_{496} \to \forall n\,(\varphi_{496})) \to \mathrm{RFn}(f) \to c \in \mathbb{R} \to \mathrm{PNbhd}(f,c) \to \varphi_{540}$
\item $\varphi_{735} :\equiv (\mathrm{RFn}(f) \to c \in \mathbb{R} \to \mathrm{PNbhd}(f,c) \to \varphi_{296}) \to (\mathrm{RFn}(f) \to c \in \mathbb{R} \to \mathrm{PNbhd}(f,c) \to \varphi_{295}) \to \varphi_{677}$
\item $\varphi_{736} :\equiv (\varphi_{496} \to \forall m\,(\varphi_{496})) \to \mathrm{RFn}(f) \to c \in \mathbb{R} \to \mathrm{PNbhd}(f,c) \to \varphi_{542}$
\item $\varphi_{737} :\equiv (\mathrm{RFn}(f) \to c \in \mathbb{R} \to \mathrm{PNbhd}(f,c) \to \varphi_{285}) \to \forall r\,(\varphi_{261} \to \neg (r \in \mathbb{R} \land 0 < r \land \varphi_{16})) \to \varphi_{515} \to \varphi_{554}$
\item $\varphi_{738} :\equiv (\varphi_{537} \to \forall N\,(\varphi_{537})) \to \mathrm{RFn}(f) \to c \in \mathbb{R} \to \mathrm{PNbhd}(f,c) \to \varphi_{501}$
\item $\varphi_{739} :\equiv (\mathrm{Seq}(p) \to N_{0} \cup N_{1} \in \mathbb{N} \to p(N_{0} \cup N_{1}) \in \mathbb{R}) \to (\mathrm{RFn}(f) \to c \in \mathbb{R} \to \mathrm{PNbhd}(f,c) \to \varphi_{343}) \to \varphi_{707}$
\item $\varphi_{740} :\equiv (\forall \alpha\,(\varphi_{389}) \to \varphi_{400}) \to \varphi_{591} \to \mathrm{RFn}(f) \to c \in \mathbb{R} \to \mathrm{PNbhd}(f,c) \to \varphi_{351}$
\item $\varphi_{741} :\equiv (\mathrm{RFn}(f) \to c \in \mathbb{R} \to \mathrm{PNbhd}(f,c) \to \varphi_{371}) \to (\mathrm{RFn}(f) \to c \in \mathbb{R} \to \mathrm{PNbhd}(f,c) \to \varphi_{385}) \to \varphi_{662}$
\item $\varphi_{742} :\equiv (\mathrm{RFn}(f) \to c \in \mathbb{R} \to \mathrm{PNbhd}(f,c) \to \varphi_{349}) \to (\mathrm{RFn}(f) \to c \in \mathbb{R} \to \mathrm{PNbhd}(f,c) \to \varphi_{470}) \to \varphi_{660}$
\item $\varphi_{743} :\equiv (\mathrm{RFn}(f) \to c \in \mathbb{R} \to \mathrm{PNbhd}(f,c) \to \varphi_{518}) \to (\mathrm{RFn}(f) \to c \in \mathbb{R} \to \mathrm{PNbhd}(f,c) \to \varphi_{529}) \to \varphi_{603}$
\item $\varphi_{744} :\equiv (\alpha \in \mathbb{R} \to f(x) \in \mathbb{R} \to \varphi_{204}) \to (\mathrm{RFn}(f) \to c \in \mathbb{R} \to \mathrm{PNbhd}(f,c) \to \varphi_{3}) \to \varphi_{735}$
\item $\varphi_{745} :\equiv (\mathrm{RFn}(f) \to c \in \mathbb{R} \to \mathrm{PNbhd}(f,c) \to \varphi_{520}) \to (\mathrm{RFn}(f) \to c \in \mathbb{R} \to \mathrm{PNbhd}(f,c) \to \varphi_{536}) \to \varphi_{612}$
\item $\varphi_{746} :\equiv (\mathrm{Conv}(f \circ p,\alpha) \to \varepsilon \cdot 2^{-1} \in \mathbb{R} \to 0 < \varepsilon \cdot 2^{-1} \to \varphi_{70}) \to \varphi_{741}$
\item $\varphi_{747} :\equiv (\mathrm{RFn}(f) \to c \in \mathbb{R} \to \mathrm{PNbhd}(f,c) \to \varphi_{371}) \to (\mathrm{RFn}(f) \to c \in \mathbb{R} \to \mathrm{PNbhd}(f,c) \to \varphi_{485}) \to \varphi_{671}$
\item $\varphi_{748} :\equiv (\varphi_{570} \to \forall x\,(\varphi_{570})) \to \mathrm{RFn}(f) \to c \in \mathbb{R} \to \mathrm{PNbhd}(f,c) \to \varphi_{595}$
\item $\varphi_{749} :\equiv (N_{0} \in \mathbb{N} \to N_{1} \in \mathbb{N} \to N_{0} \cup N_{1} \in \mathbb{N}) \to (\mathrm{RFn}(f) \to c \in \mathbb{R} \to \mathrm{PNbhd}(f,c) \to \varphi_{560}) \to \varphi_{705}$
\item $\varphi_{750} :\equiv (\forall \varepsilon\,(\varphi_{463}) \to \varphi_{468}) \to \varphi_{599} \to \mathrm{RFn}(f) \to c \in \mathbb{R} \to \mathrm{PNbhd}(f,c) \to \varphi_{459}$
\item $\varphi_{751} :\equiv (\forall p\,(\varphi_{320} \to \neg \varphi_{76}) \to \forall p\,(\varphi_{320}) \to \forall p\,(\neg \varphi_{76})) \to (\varphi_{320} \to \forall p\,(\varphi_{320})) \to \forall a\,(\forall p\,(\neg \varphi_{76}) \to \varphi_{78}) \to \varphi_{604}$
\item $\varphi_{752} :\equiv (\mathrm{RFn}(f) \to c \in \mathbb{R} \to \mathrm{PNbhd}(f,c) \to \varphi_{602}) \to (\mathrm{RFn}(f) \to c \in \mathbb{R} \to \mathrm{PNbhd}(f,c) \to \varphi_{562}) \to \varphi_{606}$
\item $\varphi_{753} :\equiv (\mathrm{RFn}(f) \to c \in \mathbb{R} \to \mathrm{PNbhd}(f,c) \to \varphi_{613}) \to (\mathrm{RFn}(f) \to c \in \mathbb{R} \to \mathrm{PNbhd}(f,c) \to \varphi_{562}) \to \varphi_{619}$
\item $\varphi_{754} :\equiv (\mathrm{RFn}(f) \to c \in \mathbb{R} \to \mathrm{PNbhd}(f,c) \to \varphi_{577}) \to (\mathrm{RFn}(f) \to c \in \mathbb{R} \to \mathrm{PNbhd}(f,c) \to \varphi_{578}) \to \varphi_{636}$
\item $\varphi_{755} :\equiv (\mathrm{RFn}(f) \to c \in \mathbb{R} \to \mathrm{PNbhd}(f,c) \to \varphi_{525}) \to (\mathrm{RFn}(f) \to c \in \mathbb{R} \to \mathrm{PNbhd}(f,c) \to \varphi_{494}) \to \varphi_{679}$
\item $\varphi_{756} :\equiv (\mathrm{RFn}(f) \to c \in \mathbb{R} \to \mathrm{PNbhd}(f,c) \to \varphi_{600}) \to (\mathrm{RFn}(f) \to c \in \mathbb{R} \to \mathrm{PNbhd}(f,c) \to \varphi_{607}) \to \varphi_{631}$
\item $\varphi_{757} :\equiv (\mathrm{RFn}(f) \to c \in \mathbb{R} \to \mathrm{PNbhd}(f,c) \to \varphi_{524}) \to (\mathrm{RFn}(f) \to c \in \mathbb{R} \to \mathrm{PNbhd}(f,c) \to \varphi_{497}) \to \varphi_{680}$
\item $\varphi_{758} :\equiv (\mathrm{RFn}(f) \to c \in \mathbb{R} \to \mathrm{PNbhd}(f,c) \to \varphi_{512}) \to (\mathrm{RFn}(f) \to c \in \mathbb{R} \to \mathrm{PNbhd}(f,c) \to \varphi_{533}) \to \varphi_{691}$
\item $\varphi_{759} :\equiv (\mathrm{RFn}(f) \to c \in \mathbb{R} \to \mathrm{PNbhd}(f,c) \to \varphi_{552}) \to (\mathrm{RFn}(f) \to c \in \mathbb{R} \to \mathrm{PNbhd}(f,c) \to \varphi_{502}) \to \varphi_{689}$
\item $\varphi_{760} :\equiv (\mathrm{RFn}(f) \to c \in \mathbb{R} \to \mathrm{PNbhd}(f,c) \to \varphi_{553}) \to (\mathrm{RFn}(f) \to c \in \mathbb{R} \to \mathrm{PNbhd}(f,c) \to \varphi_{503}) \to \varphi_{690}$
\item $\varphi_{761} :\equiv \forall \varepsilon\,(\varphi_{362} \to \varepsilon \in \mathbb{R} \to 0 < \varepsilon \to \exists N \in \mathbb{N}\,\forall m\,\forall n\,(\varphi_{90})) \to \varphi_{646} \to \varphi_{670}$
\item $\varphi_{762} :\equiv (\forall \delta\,(\varphi_{538}) \to \varphi_{541}) \to \varphi_{644} \to \mathrm{RFn}(f) \to c \in \mathbb{R} \to \mathrm{PNbhd}(f,c) \to \varphi_{447}$
\item $\varphi_{763} :\equiv (\mathrm{RFn}(f) \to c \in \mathbb{R} \to \mathrm{PNbhd}(f,c) \to \varphi_{356}) \to \forall \alpha\,(\varphi_{389}) \to \varphi_{740}$
\item $\varphi_{764} :\equiv (\mathrm{RFn}(f) \to c \in \mathbb{R} \to \mathrm{PNbhd}(f,c) \to \varphi_{385}) \to (\mathrm{RFn}(f) \to c \in \mathbb{R} \to \mathrm{PNbhd}(f,c) \to \varphi_{611}) \to \varphi_{720}$
\item $\varphi_{765} :\equiv (\mathrm{RFn}(f) \to c \in \mathbb{R} \to \mathrm{PNbhd}(f,c) \to \varphi_{337}) \to \forall p\,(\varphi_{320} \to \neg \varphi_{76}) \to \varphi_{751}$
\item $\varphi_{766} :\equiv (\mathrm{RFn}(f) \to c \in \mathbb{R} \to \mathrm{PNbhd}(f,c) \to \varphi_{616}) \to (\mathrm{RFn}(f) \to c \in \mathbb{R} \to \mathrm{PNbhd}(f,c) \to \varphi_{566}) \to \varphi_{714}$
\item $\varphi_{767} :\equiv (\mathrm{RFn}(f) \to c \in \mathbb{R} \to \mathrm{PNbhd}(f,c) \to \varphi_{586}) \to (\mathrm{RFn}(f) \to c \in \mathbb{R} \to \mathrm{PNbhd}(f,c) \to \varphi_{584}) \to \varphi_{724}$
\item $\varphi_{768} :\equiv \forall n\,(\varphi_{496} \to \varphi_{90}) \to (\forall n\,(\varphi_{496} \to \varphi_{90}) \to \forall n\,(\varphi_{496}) \to \forall n\,(\varphi_{90})) \to \varphi_{734}$
\item $\varphi_{769} :\equiv \forall m\,(\varphi_{496} \to \forall n\,(\varphi_{90})) \to (\forall m\,(\varphi_{496} \to \forall n\,(\varphi_{90})) \to \forall m\,(\varphi_{496}) \to \forall m\,\forall n\,(\varphi_{90})) \to \varphi_{736}$
\item $\varphi_{770} :\equiv (\mathrm{RFn}(f) \to c \in \mathbb{R} \to \mathrm{PNbhd}(f,c) \to \varphi_{504}) \to (\mathrm{RFn}(f) \to c \in \mathbb{R} \to \mathrm{PNbhd}(f,c) \to \varphi_{534}) \to \varphi_{758}$
\item $\varphi_{771} :\equiv (\mathrm{RFn}(f) \to c \in \mathbb{R} \to \mathrm{PNbhd}(f,c) \to \varphi_{466}) \to \forall \delta\,(\varphi_{538}) \to \varphi_{762}$
\item $\varphi_{772} :\equiv (\mathrm{RFn}(f) \to c \in \mathbb{R} \to \mathrm{PNbhd}(f,c) \to \varphi_{598}) \to (\mathrm{RFn}(f) \to c \in \mathbb{R} \to \mathrm{PNbhd}(f,c) \to \varphi_{365}) \to \varphi_{764}$
\item $\varphi_{773} :\equiv (\mathrm{RFn}(f) \to c \in \mathbb{R} \to \mathrm{PNbhd}(f,c) \to \varphi_{478}) \to \forall \delta\,(\varphi_{464} \to \neg (\delta \in \mathbb{R} \land 0 < \delta \land \forall x\,(\varphi_{176}))) \to \varphi_{693} \to \varphi_{711}$
\item $\varphi_{774} :\equiv (\forall N_{1}\,(\varphi_{582} \to \neg \varphi_{67}) \to \forall N_{1}\,(\varphi_{582}) \to \forall N_{1}\,(\neg \varphi_{67})) \to (\varphi_{582} \to \forall N_{1}\,(\varphi_{582})) \to \forall N\,(\forall N_{1}\,(\neg \varphi_{67}) \to \varphi_{60}) \to \varphi_{669}$
\item $\varphi_{775} :\equiv (\mathrm{RFn}(f) \to c \in \mathbb{R} \to \mathrm{PNbhd}(f,c) \to \varphi_{488}) \to \forall N\,(\varphi_{556}) \to (\forall N\,(\varphi_{556}) \to \forall N\,(\varphi_{537}) \to \varphi_{7}) \to \varphi_{738}$
\item $\varphi_{776} :\equiv (\mathrm{RFn}(f) \to c \in \mathbb{R} \to \mathrm{PNbhd}(f,c) \to \varphi_{569}) \to (\mathrm{RFn}(f) \to c \in \mathbb{R} \to \mathrm{PNbhd}(f,c) \to \varphi_{575}) \to \varphi_{767}$
\item $\varphi_{777} :\equiv \forall N_{0}\,(\varphi_{550} \to \neg (N_{0} \in \mathbb{N} \land \varphi_{77})) \to \varphi_{709} \to (\varphi_{550} \to \forall N_{0}\,(\varphi_{550})) \to \varphi_{697}$
\item $\varphi_{778} :\equiv (f(x) \in \mathbb{R} \to f(p(N_{0} \cup N_{1})) \in \mathbb{R} \to \alpha \in \mathbb{R} \to \varphi_{208}) \to (\mathrm{RFn}(f) \to c \in \mathbb{R} \to \mathrm{PNbhd}(f,c) \to \varphi_{587}) \to \varphi_{772}$
\item $\varphi_{779} :\equiv (\mathrm{RFn}(f) \to c \in \mathbb{R} \to \mathrm{PNbhd}(f,c) \to \varphi_{588}) \to (\mathrm{RFn}(f) \to c \in \mathbb{R} \to \mathrm{PNbhd}(f,c) \to \varphi_{505}) \to \varphi_{770}$
\item $\varphi_{780} :\equiv (\mathrm{RFn}(f) \to c \in \mathbb{R} \to \mathrm{PNbhd}(f,c) \to \varphi_{558}) \to \forall N_{1}\,(\varphi_{582} \to \neg \varphi_{67}) \to \varphi_{774}$
\item $\varphi_{781} :\equiv (\mathrm{RFn}(f) \to c \in \mathbb{R} \to \mathrm{PNbhd}(f,c) \to \varphi_{637}) \to (\mathrm{RFn}(f) \to c \in \mathbb{R} \to \mathrm{PNbhd}(f,c) \to \varphi_{568}) \to \varphi_{776}$
\end{abbrevlist}
\end{fullproof}

% 自動生成 (prover/build.js): RightAbs  ノード数 87
\begin{fullproof}{[H3]}{$\forall x\,\forall c\,\forall d\,(x \in \mathbb{R} \to c \in \mathbb{R} \to d \in \mathbb{R} \to c < x \to x < c + d \to 0 < \lvert x - c \rvert \land \lvert x - c \rvert < d)$}{87}\phantomsection\label{tree:RightAbs}
\noindent$\Pi_{1}$:
\begin{ptree}
\AXC{$\varphi_{5}$\lbl{[E32]}}
\AXC{$\varphi_{5} \to \varphi_{4}$\lbl{A2}}
\BIC{$\varphi_{4}$}
\AXC{$\varphi_{17}$\lbl{A2}}
\BIC{$\forall y\,(c \in \mathbb{R} \to d \in \mathbb{R} \to y \in \mathbb{R} \to y < c + d \to y - c < d)$}
\AXC{$\varphi_{16}$\lbl{A2}}
\BIC{$c \in \mathbb{R} \to d \in \mathbb{R} \to x \in \mathbb{R} \to x < c + d \to x - c < d$}
\AXC{$(c \in \mathbb{R} \to d \in \mathbb{R} \to x \in \mathbb{R} \to x < c + d \to x - c < d) \to \varphi_{28}$\lbl{A1}}
\BIC{$\varphi_{28}$}
\end{ptree}
\noindent$\Pi_{2}$:
\begin{ptree}
\AXC{$x \in \mathbb{R} \to c \in \mathbb{R} \to d \in \mathbb{R} \to c < x \to x < c + d \to d \in \mathbb{R}$\lbl{A1}}
\AXC{$x \in \mathbb{R} \to c \in \mathbb{R} \to d \in \mathbb{R} \to c < x \to x < c + d \to c \in \mathbb{R}$\lbl{A1}}
\AXC{$\stackrel{\Pi_{1}}{\vdots}$}\noLine
\UIC{$\varphi_{28}$}
\BIC{$\varphi_{27}$}
\BIC{$\varphi_{25}$}
\end{ptree}
\noindent$\Pi_{3}$:
\begin{ptree}
\AXC{$x \in \mathbb{R} \to c \in \mathbb{R} \to d \in \mathbb{R} \to c < x \to x < c + d \to x < c + d$\lbl{A1}}
\AXC{$x \in \mathbb{R} \to c \in \mathbb{R} \to d \in \mathbb{R} \to c < x \to x < c + d \to x \in \mathbb{R}$\lbl{A1}}
\AXC{$\stackrel{\Pi_{2}}{\vdots}$}\noLine
\UIC{$\varphi_{25}$}
\BIC{$\varphi_{20}$}
\BIC{$x \in \mathbb{R} \to c \in \mathbb{R} \to d \in \mathbb{R} \to c < x \to x < c + d \to x - c < d$}
\end{ptree}
\noindent$\Pi_{4}$:
\begin{ptree}
\AXC{$x = x$\lbl{A5}}
\AXC{$x = y \to x = x \to y = x$\lbl{A6}}
\AXC{$(x = y \to x = x \to y = x) \to x = x \to x = y \to y = x$\lbl{A1}}
\BIC{$x = x \to x = y \to y = x$}
\BIC{$x = y \to y = x$}
\UIC{$\forall y\,(x = y \to y = x)$}
\UIC{$\forall x\,\forall y\,(x = y \to y = x)$}
\end{ptree}
\noindent$\Pi_{5}$:
\begin{ptree}
\AXC{$\forall x\,\forall y\,(x \in \mathbb{R} \to y \in \mathbb{R} \to x < y \to 0 < y - x)$\lbl{[F24]}}
\AXC{$\varphi_{15}$\lbl{A2}}
\BIC{$\forall y\,(c \in \mathbb{R} \to y \in \mathbb{R} \to c < y \to 0 < y - c)$}
\AXC{$\varphi_{14}$\lbl{A2}}
\BIC{$c \in \mathbb{R} \to x \in \mathbb{R} \to c < x \to 0 < x - c$}
\AXC{$(c \in \mathbb{R} \to x \in \mathbb{R} \to c < x \to 0 < x - c) \to \varphi_{26}$\lbl{A1}}
\BIC{$\varphi_{26}$}
\end{ptree}
\noindent$\Pi_{6}$:
\begin{ptree}
\AXC{$x \in \mathbb{R} \to c \in \mathbb{R} \to d \in \mathbb{R} \to c < x \to x < c + d \to x \in \mathbb{R}$\lbl{A1}}
\AXC{$x \in \mathbb{R} \to c \in \mathbb{R} \to d \in \mathbb{R} \to c < x \to x < c + d \to c \in \mathbb{R}$\lbl{A1}}
\AXC{$\stackrel{\Pi_{5}}{\vdots}$}\noLine
\UIC{$\varphi_{26}$}
\BIC{$\varphi_{24}$}
\BIC{$\varphi_{18}$}
\end{ptree}
\noindent$\Pi_{7}$:
\begin{ptree}
\AXC{$x \in \mathbb{R} \to c \in \mathbb{R} \to d \in \mathbb{R} \to c < x \to x < c + d \to c < x$\lbl{A1}}
\AXC{$\stackrel{\Pi_{6}}{\vdots}$}\noLine
\UIC{$\varphi_{18}$}
\BIC{$x \in \mathbb{R} \to c \in \mathbb{R} \to d \in \mathbb{R} \to c < x \to x < c + d \to 0 < x - c$}
\end{ptree}
\noindent$\Pi_{8}$:
\begin{ptree}
\AXC{$\forall x\,(0 \leq x \to \lvert x \rvert = x)$\lbl{[A1]}}
\AXC{$\forall x\,(0 \leq x \to \lvert x \rvert = x) \to 0 \leq x - c \to \lvert x - c \rvert = x - c$\lbl{A2}}
\BIC{$0 \leq x - c \to \lvert x - c \rvert = x - c$}
\end{ptree}
\noindent$\Pi_{9}$:
\begin{ptree}
\AXC{$\stackrel{\Pi_{8}}{\vdots}$}\noLine
\UIC{$0 \leq x - c \to \lvert x - c \rvert = x - c$}
\AXC{$(0 \leq x - c \to \lvert x - c \rvert = x - c) \to \varphi_{21}$\lbl{A1}}
\BIC{$\varphi_{21}$}
\end{ptree}
\noindent$\Pi_{10}$:
\begin{ptree}
\AXC{$\stackrel{\Pi_{7}}{\vdots}$}\noLine
\UIC{$x \in \mathbb{R} \to c \in \mathbb{R} \to d \in \mathbb{R} \to c < x \to x < c + d \to 0 < x - c$}
\AXC{$\varphi_{19}$\lbl{A1}}
\BIC{$x \in \mathbb{R} \to c \in \mathbb{R} \to d \in \mathbb{R} \to c < x \to x < c + d \to 0 \leq x - c$}
\AXC{$\stackrel{\Pi_{9}}{\vdots}$}\noLine
\UIC{$\varphi_{21}$}
\BIC{$\varphi_{9}$}
\end{ptree}
\noindent$\Pi_{11}$:
\begin{ptree}
\AXC{$\stackrel{\Pi_{4}}{\vdots}$}\noLine
\UIC{$\forall x\,\forall y\,(x = y \to y = x)$}
\AXC{$\varphi_{1}$\lbl{A2}}
\BIC{$\forall y\,(\lvert x - c \rvert = y \to y = \lvert x - c \rvert)$}
\AXC{$\varphi_{6}$\lbl{A2}}
\BIC{$\lvert x - c \rvert = x - c \to x - c = \lvert x - c \rvert$}
\AXC{$\stackrel{\Pi_{10}}{\vdots}$}\noLine
\UIC{$\varphi_{9}$}
\AXC{$\varphi_{9} \to (\lvert x - c \rvert = x - c \to x - c = \lvert x - c \rvert) \to \varphi_{10}$\lbl{A1}}
\BIC{$(\lvert x - c \rvert = x - c \to x - c = \lvert x - c \rvert) \to \varphi_{10}$}
\BIC{$\varphi_{10}$}
\end{ptree}
\noindent$\Pi_{12}$:
\begin{ptree}
\AXC{$\stackrel{\Pi_{3}}{\vdots}$}\noLine
\UIC{$x \in \mathbb{R} \to c \in \mathbb{R} \to d \in \mathbb{R} \to c < x \to x < c + d \to x - c < d$}
\AXC{$\stackrel{\Pi_{11}}{\vdots}$}\noLine
\UIC{$\varphi_{10}$}
\AXC{$\varphi_{10} \to \varphi_{22}$\lbl{A1}}
\BIC{$\varphi_{22}$}
\BIC{$(x - c = \lvert x - c \rvert \to x - c < d \to \lvert x - c \rvert < d) \to \varphi_{7}$}
\end{ptree}
\noindent$\Pi_{13}$:
\begin{ptree}
\AXC{$z = w \to z < d \to w < d$\lbl{A6}}
\UIC{$\forall w\,(z = w \to z < d \to w < d)$}
\UIC{$\forall z\,\forall w\,(z = w \to z < d \to w < d)$}
\AXC{$\varphi_{2}$\lbl{A2}}
\BIC{$\forall w\,(x - c = w \to x - c < d \to w < d)$}
\AXC{$\varphi_{12}$\lbl{A2}}
\BIC{$x - c = \lvert x - c \rvert \to x - c < d \to \lvert x - c \rvert < d$}
\AXC{$\stackrel{\Pi_{12}}{\vdots}$}\noLine
\UIC{$(x - c = \lvert x - c \rvert \to x - c < d \to \lvert x - c \rvert < d) \to \varphi_{7}$}
\BIC{$\varphi_{7}$}
\end{ptree}
\noindent$\Pi_{14}$:
\begin{ptree}
\AXC{$\stackrel{\Pi_{4}}{\vdots}$}\noLine
\UIC{$\forall x\,\forall y\,(x = y \to y = x)$}
\AXC{$\varphi_{1}$\lbl{A2}}
\BIC{$\forall y\,(\lvert x - c \rvert = y \to y = \lvert x - c \rvert)$}
\AXC{$\varphi_{6}$\lbl{A2}}
\BIC{$\lvert x - c \rvert = x - c \to x - c = \lvert x - c \rvert$}
\AXC{$\stackrel{\Pi_{10}}{\vdots}$}\noLine
\UIC{$\varphi_{9}$}
\AXC{$\varphi_{9} \to (\lvert x - c \rvert = x - c \to x - c = \lvert x - c \rvert) \to \varphi_{10}$\lbl{A1}}
\BIC{$(\lvert x - c \rvert = x - c \to x - c = \lvert x - c \rvert) \to \varphi_{10}$}
\BIC{$\varphi_{10}$}
\end{ptree}
\noindent$\Pi_{15}$:
\begin{ptree}
\AXC{$\stackrel{\Pi_{7}}{\vdots}$}\noLine
\UIC{$x \in \mathbb{R} \to c \in \mathbb{R} \to d \in \mathbb{R} \to c < x \to x < c + d \to 0 < x - c$}
\AXC{$\stackrel{\Pi_{14}}{\vdots}$}\noLine
\UIC{$\varphi_{10}$}
\AXC{$\varphi_{10} \to \varphi_{23}$\lbl{A1}}
\BIC{$\varphi_{23}$}
\BIC{$(x - c = \lvert x - c \rvert \to 0 < x - c \to 0 < \lvert x - c \rvert) \to \varphi_{8}$}
\end{ptree}
\noindent$\Pi_{16}$:
\begin{ptree}
\AXC{$z = w \to 0 < z \to 0 < w$\lbl{A6}}
\UIC{$\forall w\,(z = w \to 0 < z \to 0 < w)$}
\UIC{$\forall z\,\forall w\,(z = w \to 0 < z \to 0 < w)$}
\AXC{$\varphi_{3}$\lbl{A2}}
\BIC{$\forall w\,(x - c = w \to 0 < x - c \to 0 < w)$}
\AXC{$\varphi_{13}$\lbl{A2}}
\BIC{$x - c = \lvert x - c \rvert \to 0 < x - c \to 0 < \lvert x - c \rvert$}
\AXC{$\stackrel{\Pi_{15}}{\vdots}$}\noLine
\UIC{$(x - c = \lvert x - c \rvert \to 0 < x - c \to 0 < \lvert x - c \rvert) \to \varphi_{8}$}
\BIC{$\varphi_{8}$}
\end{ptree}
\noindent 主木:
\begin{ptree}
\AXC{$\stackrel{\Pi_{13}}{\vdots}$}\noLine
\UIC{$\varphi_{7}$}
\AXC{$\stackrel{\Pi_{16}}{\vdots}$}\noLine
\UIC{$\varphi_{8}$}
\AXC{$\varphi_{8} \to \varphi_{7} \to x \in \mathbb{R} \to \varphi_{11}$\lbl{A1}}
\BIC{$\varphi_{7} \to x \in \mathbb{R} \to \varphi_{11}$}
\BIC{$x \in \mathbb{R} \to \varphi_{11}$}
\UIC{$\forall d\,(x \in \mathbb{R} \to \varphi_{11})$}
\UIC{$\forall c\,\forall d\,(x \in \mathbb{R} \to \varphi_{11})$}
\UIC{$\forall x\,\forall c\,\forall d\,(x \in \mathbb{R} \to \varphi_{11})$}
\end{ptree}
\begin{abbrevlist}
\item $\varphi_{1} :\equiv \forall x\,\forall y\,(x = y \to y = x) \to \forall y\,(\lvert x - c \rvert = y \to y = \lvert x - c \rvert)$
\item $\varphi_{2} :\equiv \forall z\,\forall w\,(z = w \to z < d \to w < d) \to \forall w\,(x - c = w \to x - c < d \to w < d)$
\item $\varphi_{3} :\equiv \forall z\,\forall w\,(z = w \to 0 < z \to 0 < w) \to \forall w\,(x - c = w \to 0 < x - c \to 0 < w)$
\item $\varphi_{4} :\equiv \forall e\,\forall y\,(c \in \mathbb{R} \to e \in \mathbb{R} \to y \in \mathbb{R} \to y < c + e \to y - c < e)$
\item $\varphi_{5} :\equiv \forall x\,\forall e\,\forall y\,(x \in \mathbb{R} \to e \in \mathbb{R} \to y \in \mathbb{R} \to y < x + e \to y - x < e)$
\item $\varphi_{6} :\equiv \forall y\,(\lvert x - c \rvert = y \to y = \lvert x - c \rvert) \to \lvert x - c \rvert = x - c \to x - c = \lvert x - c \rvert$
\item $\varphi_{7} :\equiv x \in \mathbb{R} \to c \in \mathbb{R} \to d \in \mathbb{R} \to c < x \to x < c + d \to \lvert x - c \rvert < d$
\item $\varphi_{8} :\equiv x \in \mathbb{R} \to c \in \mathbb{R} \to d \in \mathbb{R} \to c < x \to x < c + d \to 0 < \lvert x - c \rvert$
\item $\varphi_{9} :\equiv x \in \mathbb{R} \to c \in \mathbb{R} \to d \in \mathbb{R} \to c < x \to x < c + d \to \lvert x - c \rvert = x - c$
\item $\varphi_{10} :\equiv x \in \mathbb{R} \to c \in \mathbb{R} \to d \in \mathbb{R} \to c < x \to x < c + d \to x - c = \lvert x - c \rvert$
\item $\varphi_{11} :\equiv c \in \mathbb{R} \to d \in \mathbb{R} \to c < x \to x < c + d \to 0 < \lvert x - c \rvert \land \lvert x - c \rvert < d$
\item $\varphi_{12} :\equiv \forall w\,(x - c = w \to x - c < d \to w < d) \to x - c = \lvert x - c \rvert \to x - c < d \to \lvert x - c \rvert < d$
\item $\varphi_{13} :\equiv \forall w\,(x - c = w \to 0 < x - c \to 0 < w) \to x - c = \lvert x - c \rvert \to 0 < x - c \to 0 < \lvert x - c \rvert$
\item $\varphi_{14} :\equiv \forall y\,(c \in \mathbb{R} \to y \in \mathbb{R} \to c < y \to 0 < y - c) \to c \in \mathbb{R} \to x \in \mathbb{R} \to c < x \to 0 < x - c$
\item $\varphi_{15} :\equiv \forall x\,\forall y\,(x \in \mathbb{R} \to y \in \mathbb{R} \to x < y \to 0 < y - x) \to \forall y\,(c \in \mathbb{R} \to y \in \mathbb{R} \to c < y \to 0 < y - c)$
\item $\varphi_{16} :\equiv \forall y\,(c \in \mathbb{R} \to d \in \mathbb{R} \to y \in \mathbb{R} \to y < c + d \to y - c < d) \to c \in \mathbb{R} \to d \in \mathbb{R} \to x \in \mathbb{R} \to x < c + d \to x - c < d$
\item $\varphi_{17} :\equiv \varphi_{4} \to \forall y\,(c \in \mathbb{R} \to d \in \mathbb{R} \to y \in \mathbb{R} \to y < c + d \to y - c < d)$
\item $\varphi_{18} :\equiv (x \in \mathbb{R} \to c \in \mathbb{R} \to d \in \mathbb{R} \to c < x \to x < c + d \to c < x) \to x \in \mathbb{R} \to c \in \mathbb{R} \to d \in \mathbb{R} \to c < x \to x < c + d \to 0 < x - c$
\item $\varphi_{19} :\equiv (x \in \mathbb{R} \to c \in \mathbb{R} \to d \in \mathbb{R} \to c < x \to x < c + d \to 0 < x - c) \to x \in \mathbb{R} \to c \in \mathbb{R} \to d \in \mathbb{R} \to c < x \to x < c + d \to 0 \leq x - c$
\item $\varphi_{20} :\equiv (x \in \mathbb{R} \to c \in \mathbb{R} \to d \in \mathbb{R} \to c < x \to x < c + d \to x < c + d) \to x \in \mathbb{R} \to c \in \mathbb{R} \to d \in \mathbb{R} \to c < x \to x < c + d \to x - c < d$
\item $\varphi_{21} :\equiv (x \in \mathbb{R} \to c \in \mathbb{R} \to d \in \mathbb{R} \to c < x \to x < c + d \to 0 \leq x - c) \to \varphi_{9}$
\item $\varphi_{22} :\equiv (x \in \mathbb{R} \to c \in \mathbb{R} \to d \in \mathbb{R} \to c < x \to x < c + d \to x - c < d) \to (x - c = \lvert x - c \rvert \to x - c < d \to \lvert x - c \rvert < d) \to \varphi_{7}$
\item $\varphi_{23} :\equiv (x \in \mathbb{R} \to c \in \mathbb{R} \to d \in \mathbb{R} \to c < x \to x < c + d \to 0 < x - c) \to (x - c = \lvert x - c \rvert \to 0 < x - c \to 0 < \lvert x - c \rvert) \to \varphi_{8}$
\item $\varphi_{24} :\equiv (x \in \mathbb{R} \to c \in \mathbb{R} \to d \in \mathbb{R} \to c < x \to x < c + d \to x \in \mathbb{R}) \to \varphi_{18}$
\item $\varphi_{25} :\equiv (x \in \mathbb{R} \to c \in \mathbb{R} \to d \in \mathbb{R} \to c < x \to x < c + d \to x \in \mathbb{R}) \to \varphi_{20}$
\item $\varphi_{26} :\equiv (x \in \mathbb{R} \to c \in \mathbb{R} \to d \in \mathbb{R} \to c < x \to x < c + d \to c \in \mathbb{R}) \to \varphi_{24}$
\item $\varphi_{27} :\equiv (x \in \mathbb{R} \to c \in \mathbb{R} \to d \in \mathbb{R} \to c < x \to x < c + d \to d \in \mathbb{R}) \to \varphi_{25}$
\item $\varphi_{28} :\equiv (x \in \mathbb{R} \to c \in \mathbb{R} \to d \in \mathbb{R} \to c < x \to x < c + d \to c \in \mathbb{R}) \to \varphi_{27}$
\end{abbrevlist}
\end{fullproof}

% 自動生成 (prover/build.js): LtAddSwap2  ノード数 180
\begin{fullproof}{[H4]}{$\forall c\,\forall d\,\forall x\,(c \in \mathbb{R} \to d \in \mathbb{R} \to x \in \mathbb{R} \to c < d + x \to c - d < x)$}{180}\phantomsection\label{tree:LtAddSwap2}
\noindent$\Pi_{1}$:
\begin{ptree}
\AXC{$z = w \to c + -d < z \to c + -d < w$\lbl{A6}}
\UIC{$\forall w\,(z = w \to c + -d < z \to c + -d < w)$}
\UIC{$\forall z\,\forall w\,(z = w \to c + -d < z \to c + -d < w)$}
\AXC{$\forall z\,\forall w\,(z = w \to c + -d < z \to c + -d < w) \to \varphi_{13}$\lbl{A2}}
\BIC{$\varphi_{13}$}
\end{ptree}
\noindent$\Pi_{2}$:
\begin{ptree}
\AXC{$\forall x \in \mathbb{R}\,(-x \in \mathbb{R})$\lbl{(R1)}}
\AXC{$\forall x \in \mathbb{R}\,(-x \in \mathbb{R}) \to d \in \mathbb{R} \to -d \in \mathbb{R}$\lbl{A2}}
\BIC{$d \in \mathbb{R} \to -d \in \mathbb{R}$}
\AXC{$(d \in \mathbb{R} \to -d \in \mathbb{R}) \to \varphi_{42}$\lbl{A1}}
\BIC{$\varphi_{42}$}
\end{ptree}
\noindent$\Pi_{3}$:
\begin{ptree}
\AXC{$c \in \mathbb{R} \to d \in \mathbb{R} \to x \in \mathbb{R} \to c < d + x \to d \in \mathbb{R}$\lbl{A1}}
\AXC{$\stackrel{\Pi_{2}}{\vdots}$}\noLine
\UIC{$\varphi_{42}$}
\BIC{$c \in \mathbb{R} \to d \in \mathbb{R} \to x \in \mathbb{R} \to c < d + x \to -d \in \mathbb{R}$}
\end{ptree}
\noindent$\Pi_{4}$:
\begin{ptree}
\AXC{$\forall x\,\forall y\,(x \in \mathbb{R} \to y \in \mathbb{R} \to x + y \in \mathbb{R})$\lbl{(R1)}}
\AXC{$\varphi_{32}$\lbl{A2}}
\BIC{$\forall y\,(d \in \mathbb{R} \to y \in \mathbb{R} \to d + y \in \mathbb{R})$}
\AXC{$\varphi_{30}$\lbl{A2}}
\BIC{$d \in \mathbb{R} \to x \in \mathbb{R} \to d + x \in \mathbb{R}$}
\AXC{$\varphi_{49}$\lbl{A1}}
\BIC{$(c \in \mathbb{R} \to d \in \mathbb{R} \to x \in \mathbb{R} \to c < d + x \to d \in \mathbb{R}) \to \varphi_{43}$}
\end{ptree}
\noindent$\Pi_{5}$:
\begin{ptree}
\AXC{$c \in \mathbb{R} \to d \in \mathbb{R} \to x \in \mathbb{R} \to c < d + x \to d \in \mathbb{R}$\lbl{A1}}
\AXC{$\stackrel{\Pi_{4}}{\vdots}$}\noLine
\UIC{$(c \in \mathbb{R} \to d \in \mathbb{R} \to x \in \mathbb{R} \to c < d + x \to d \in \mathbb{R}) \to \varphi_{43}$}
\BIC{$\varphi_{43}$}
\end{ptree}
\noindent$\Pi_{6}$:
\begin{ptree}
\AXC{$c \in \mathbb{R} \to d \in \mathbb{R} \to x \in \mathbb{R} \to c < d + x \to c \in \mathbb{R}$\lbl{A1}}
\AXC{$\varphi_{9}$\lbl{(R14)}}
\AXC{$\varphi_{9} \to \varphi_{7}$\lbl{A2}}
\BIC{$\varphi_{7}$}
\AXC{$\varphi_{7} \to \varphi_{24}$\lbl{A2}}
\BIC{$\varphi_{24}$}
\AXC{$\varphi_{24} \to \varphi_{33}$\lbl{A2}}
\BIC{$\varphi_{33}$}
\AXC{$\varphi_{57}$\lbl{A1}}
\BIC{$(c \in \mathbb{R} \to d \in \mathbb{R} \to x \in \mathbb{R} \to c < d + x \to c \in \mathbb{R}) \to \varphi_{55}$}
\BIC{$\varphi_{55}$}
\end{ptree}
\noindent$\Pi_{7}$:
\begin{ptree}
\AXC{$\stackrel{\Pi_{3}}{\vdots}$}\noLine
\UIC{$c \in \mathbb{R} \to d \in \mathbb{R} \to x \in \mathbb{R} \to c < d + x \to -d \in \mathbb{R}$}
\AXC{$c \in \mathbb{R} \to d \in \mathbb{R} \to x \in \mathbb{R} \to c < d + x \to x \in \mathbb{R}$\lbl{A1}}
\AXC{$\stackrel{\Pi_{5}}{\vdots}$}\noLine
\UIC{$\varphi_{43}$}
\BIC{$c \in \mathbb{R} \to d \in \mathbb{R} \to x \in \mathbb{R} \to c < d + x \to d + x \in \mathbb{R}$}
\AXC{$\stackrel{\Pi_{6}}{\vdots}$}\noLine
\UIC{$\varphi_{55}$}
\BIC{$\varphi_{47}$}
\BIC{$(c \in \mathbb{R} \to d \in \mathbb{R} \to x \in \mathbb{R} \to c < d + x \to c < d + x) \to \varphi_{25}$}
\end{ptree}
\noindent$\Pi_{8}$:
\begin{ptree}
\AXC{$c \in \mathbb{R} \to d \in \mathbb{R} \to x \in \mathbb{R} \to c < d + x \to c < d + x$\lbl{A1}}
\AXC{$\stackrel{\Pi_{7}}{\vdots}$}\noLine
\UIC{$(c \in \mathbb{R} \to d \in \mathbb{R} \to x \in \mathbb{R} \to c < d + x \to c < d + x) \to \varphi_{25}$}
\BIC{$\varphi_{25}$}
\end{ptree}
\noindent$\Pi_{9}$:
\begin{ptree}
\AXC{$y = z \to x = y \to x = z$\lbl{A6}}
\AXC{$(y = z \to x = y \to x = z) \to x = y \to y = z \to x = z$\lbl{A1}}
\BIC{$x = y \to y = z \to x = z$}
\UIC{$\forall z\,(x = y \to y = z \to x = z)$}
\UIC{$\forall y\,\forall z\,(x = y \to y = z \to x = z)$}
\UIC{$\forall x\,\forall y\,\forall z\,(x = y \to y = z \to x = z)$}
\AXC{$\varphi_{18}$\lbl{A2}}
\BIC{$\forall y\,\forall z\,(d + x + -d = y \to y = z \to d + x + -d = z)$}
\AXC{$\varphi_{35}$\lbl{A2}}
\BIC{$\forall z\,(d + x + -d = x + 0 \to x + 0 = z \to d + x + -d = z)$}
\AXC{$\varphi_{36}$\lbl{A2}}
\BIC{$d + x + -d = x + 0 \to x + 0 = x \to d + x + -d = x$}
\end{ptree}
\noindent$\Pi_{10}$:
\begin{ptree}
\AXC{$\forall x \in \mathbb{R}\,(x + 0 = x)$\lbl{(R4)}}
\AXC{$\forall x \in \mathbb{R}\,(x + 0 = x) \to x \in \mathbb{R} \to x + 0 = x$\lbl{A2}}
\BIC{$x \in \mathbb{R} \to x + 0 = x$}
\AXC{$(x \in \mathbb{R} \to x + 0 = x) \to \varphi_{38}$\lbl{A1}}
\BIC{$\varphi_{38}$}
\end{ptree}
\noindent$\Pi_{11}$:
\begin{ptree}
\AXC{$y = z \to x = y \to x = z$\lbl{A6}}
\AXC{$(y = z \to x = y \to x = z) \to x = y \to y = z \to x = z$\lbl{A1}}
\BIC{$x = y \to y = z \to x = z$}
\UIC{$\forall z\,(x = y \to y = z \to x = z)$}
\UIC{$\forall y\,\forall z\,(x = y \to y = z \to x = z)$}
\UIC{$\forall x\,\forall y\,\forall z\,(x = y \to y = z \to x = z)$}
\AXC{$\varphi_{18}$\lbl{A2}}
\BIC{$\forall y\,\forall z\,(d + x + -d = y \to y = z \to d + x + -d = z)$}
\end{ptree}
\noindent$\Pi_{12}$:
\begin{ptree}
\AXC{$\stackrel{\Pi_{11}}{\vdots}$}\noLine
\UIC{$\forall y\,\forall z\,(d + x + -d = y \to y = z \to d + x + -d = z)$}
\AXC{$\forall y\,\forall z\,(d + x + -d = y \to y = z \to d + x + -d = z) \to \varphi_{19}$\lbl{A2}}
\BIC{$\varphi_{19}$}
\AXC{$\varphi_{19} \to \varphi_{27}$\lbl{A2}}
\BIC{$\varphi_{27}$}
\end{ptree}
\noindent$\Pi_{13}$:
\begin{ptree}
\AXC{$\forall x \in \mathbb{R}\,(x + -x = 0)$\lbl{(R5)}}
\AXC{$\forall x \in \mathbb{R}\,(x + -x = 0) \to d \in \mathbb{R} \to d + -d = 0$\lbl{A2}}
\BIC{$d \in \mathbb{R} \to d + -d = 0$}
\AXC{$(d \in \mathbb{R} \to d + -d = 0) \to \varphi_{40}$\lbl{A1}}
\BIC{$\varphi_{40}$}
\end{ptree}
\noindent$\Pi_{14}$:
\begin{ptree}
\AXC{$c \in \mathbb{R} \to d \in \mathbb{R} \to x \in \mathbb{R} \to c < d + x \to d \in \mathbb{R}$\lbl{A1}}
\AXC{$\stackrel{\Pi_{13}}{\vdots}$}\noLine
\UIC{$\varphi_{40}$}
\BIC{$c \in \mathbb{R} \to d \in \mathbb{R} \to x \in \mathbb{R} \to c < d + x \to d + -d = 0$}
\AXC{$\varphi_{50}$\lbl{A1}}
\BIC{$x + (d + -d) = x + (d + -d) \to \varphi_{14} \to \varphi_{15}$}
\end{ptree}
\noindent$\Pi_{15}$:
\begin{ptree}
\AXC{$x = x$\lbl{A5}}
\UIC{$\forall x\,(x = x)$}
\AXC{$\forall x\,(x = x) \to x + (d + -d) = x + (d + -d)$\lbl{A2}}
\BIC{$x + (d + -d) = x + (d + -d)$}
\AXC{$\stackrel{\Pi_{14}}{\vdots}$}\noLine
\UIC{$x + (d + -d) = x + (d + -d) \to \varphi_{14} \to \varphi_{15}$}
\BIC{$\varphi_{14} \to \varphi_{15}$}
\end{ptree}
\noindent$\Pi_{16}$:
\begin{ptree}
\AXC{$z = w \to x + (d + -d) = x + z \to x + (d + -d) = x + w$\lbl{A6}}
\UIC{$\varphi_{1}$}
\UIC{$\forall z\,(\varphi_{1})$}
\AXC{$\forall z\,(\varphi_{1}) \to \varphi_{20}$\lbl{A2}}
\BIC{$\varphi_{20}$}
\AXC{$\varphi_{20} \to \varphi_{14}$\lbl{A2}}
\BIC{$\varphi_{14}$}
\AXC{$\stackrel{\Pi_{15}}{\vdots}$}\noLine
\UIC{$\varphi_{14} \to \varphi_{15}$}
\BIC{$\varphi_{15}$}
\end{ptree}
\noindent$\Pi_{17}$:
\begin{ptree}
\AXC{$y = z \to x = y \to x = z$\lbl{A6}}
\AXC{$(y = z \to x = y \to x = z) \to x = y \to y = z \to x = z$\lbl{A1}}
\BIC{$x = y \to y = z \to x = z$}
\UIC{$\forall z\,(x = y \to y = z \to x = z)$}
\UIC{$\forall y\,\forall z\,(x = y \to y = z \to x = z)$}
\UIC{$\forall x\,\forall y\,\forall z\,(x = y \to y = z \to x = z)$}
\AXC{$\varphi_{18}$\lbl{A2}}
\BIC{$\forall y\,\forall z\,(d + x + -d = y \to y = z \to d + x + -d = z)$}
\end{ptree}
\noindent$\Pi_{18}$:
\begin{ptree}
\AXC{$\stackrel{\Pi_{17}}{\vdots}$}\noLine
\UIC{$\forall y\,\forall z\,(d + x + -d = y \to y = z \to d + x + -d = z)$}
\AXC{$\forall y\,\forall z\,(d + x + -d = y \to y = z \to d + x + -d = z) \to \varphi_{21}$\lbl{A2}}
\BIC{$\varphi_{21}$}
\end{ptree}
\noindent$\Pi_{19}$:
\begin{ptree}
\AXC{$\forall x\,(\varphi_{4})$\lbl{(R2)}}
\AXC{$\forall x\,(\varphi_{4}) \to \varphi_{4}$\lbl{A2}}
\BIC{$\varphi_{4}$}
\AXC{$\varphi_{4} \to \varphi_{3}$\lbl{A2}}
\BIC{$\varphi_{3}$}
\AXC{$\varphi_{3} \to \varphi_{17}$\lbl{A2}}
\BIC{$\varphi_{17}$}
\AXC{$\varphi_{17} \to \varphi_{54}$\lbl{A1}}
\BIC{$\varphi_{54}$}
\end{ptree}
\noindent$\Pi_{20}$:
\begin{ptree}
\AXC{$c \in \mathbb{R} \to d \in \mathbb{R} \to x \in \mathbb{R} \to c < d + x \to d \in \mathbb{R}$\lbl{A1}}
\AXC{$c \in \mathbb{R} \to d \in \mathbb{R} \to x \in \mathbb{R} \to c < d + x \to x \in \mathbb{R}$\lbl{A1}}
\AXC{$\stackrel{\Pi_{19}}{\vdots}$}\noLine
\UIC{$\varphi_{54}$}
\BIC{$(c \in \mathbb{R} \to d \in \mathbb{R} \to x \in \mathbb{R} \to c < d + x \to d \in \mathbb{R}) \to \varphi_{45}$}
\BIC{$\varphi_{45}$}
\end{ptree}
\noindent$\Pi_{21}$:
\begin{ptree}
\AXC{$\forall x\,\forall y\,(x \in \mathbb{R} \to y \in \mathbb{R} \to x + y = y + x)$\lbl{(R3)}}
\AXC{$\varphi_{31}$\lbl{A2}}
\BIC{$\forall y\,(d \in \mathbb{R} \to y \in \mathbb{R} \to d + y = y + d)$}
\AXC{$\varphi_{26}$\lbl{A2}}
\BIC{$d \in \mathbb{R} \to x \in \mathbb{R} \to d + x = x + d$}
\AXC{$\varphi_{48}$\lbl{A1}}
\BIC{$(c \in \mathbb{R} \to d \in \mathbb{R} \to x \in \mathbb{R} \to c < d + x \to d \in \mathbb{R}) \to \varphi_{41}$}
\end{ptree}
\noindent$\Pi_{22}$:
\begin{ptree}
\AXC{$c \in \mathbb{R} \to d \in \mathbb{R} \to x \in \mathbb{R} \to c < d + x \to d \in \mathbb{R}$\lbl{A1}}
\AXC{$\stackrel{\Pi_{21}}{\vdots}$}\noLine
\UIC{$(c \in \mathbb{R} \to d \in \mathbb{R} \to x \in \mathbb{R} \to c < d + x \to d \in \mathbb{R}) \to \varphi_{41}$}
\BIC{$\varphi_{41}$}
\end{ptree}
\noindent$\Pi_{23}$:
\begin{ptree}
\AXC{$c \in \mathbb{R} \to d \in \mathbb{R} \to x \in \mathbb{R} \to c < d + x \to x \in \mathbb{R}$\lbl{A1}}
\AXC{$\stackrel{\Pi_{22}}{\vdots}$}\noLine
\UIC{$\varphi_{41}$}
\BIC{$c \in \mathbb{R} \to d \in \mathbb{R} \to x \in \mathbb{R} \to c < d + x \to d + x = x + d$}
\AXC{$\varphi_{52}$\lbl{A1}}
\BIC{$d + x + -d = d + x + -d \to \varphi_{28} \to \varphi_{29}$}
\end{ptree}
\noindent$\Pi_{24}$:
\begin{ptree}
\AXC{$x = x$\lbl{A5}}
\UIC{$\forall x\,(x = x)$}
\AXC{$\forall x\,(x = x) \to d + x + -d = d + x + -d$\lbl{A2}}
\BIC{$d + x + -d = d + x + -d$}
\AXC{$\stackrel{\Pi_{23}}{\vdots}$}\noLine
\UIC{$d + x + -d = d + x + -d \to \varphi_{28} \to \varphi_{29}$}
\BIC{$\varphi_{28} \to \varphi_{29}$}
\end{ptree}
\noindent$\Pi_{25}$:
\begin{ptree}
\AXC{$z = w \to d + x + -d = z + -d \to d + x + -d = w + -d$\lbl{A6}}
\UIC{$\varphi_{5}$}
\UIC{$\forall z\,(\varphi_{5})$}
\AXC{$\forall z\,(\varphi_{5}) \to \varphi_{22}$\lbl{A2}}
\BIC{$\varphi_{22}$}
\AXC{$\varphi_{22} \to \varphi_{28}$\lbl{A2}}
\BIC{$\varphi_{28}$}
\AXC{$\stackrel{\Pi_{24}}{\vdots}$}\noLine
\UIC{$\varphi_{28} \to \varphi_{29}$}
\BIC{$\varphi_{29}$}
\AXC{$\varphi_{29} \to \varphi_{53}$\lbl{A1}}
\BIC{$\varphi_{53}$}
\end{ptree}
\noindent$\Pi_{26}$:
\begin{ptree}
\AXC{$\stackrel{\Pi_{18}}{\vdots}$}\noLine
\UIC{$\varphi_{21}$}
\AXC{$\varphi_{21} \to d + x + -d = x + d + -d \to \varphi_{6}$\lbl{A2}}
\BIC{$d + x + -d = x + d + -d \to \varphi_{6}$}
\AXC{$\stackrel{\Pi_{3}}{\vdots}$}\noLine
\UIC{$c \in \mathbb{R} \to d \in \mathbb{R} \to x \in \mathbb{R} \to c < d + x \to -d \in \mathbb{R}$}
\AXC{$\stackrel{\Pi_{20}}{\vdots}$}\noLine
\UIC{$\varphi_{45}$}
\BIC{$c \in \mathbb{R} \to \varphi_{10}$}
\AXC{$\stackrel{\Pi_{25}}{\vdots}$}\noLine
\UIC{$\varphi_{53}$}
\BIC{$(d + x + -d = x + d + -d \to \varphi_{6}) \to c \in \mathbb{R} \to \varphi_{11}$}
\BIC{$c \in \mathbb{R} \to \varphi_{11}$}
\end{ptree}
\noindent$\Pi_{27}$:
\begin{ptree}
\AXC{$\stackrel{\Pi_{12}}{\vdots}$}\noLine
\UIC{$\varphi_{27}$}
\AXC{$\stackrel{\Pi_{16}}{\vdots}$}\noLine
\UIC{$\varphi_{15}$}
\AXC{$\stackrel{\Pi_{26}}{\vdots}$}\noLine
\UIC{$c \in \mathbb{R} \to \varphi_{11}$}
\AXC{$(c \in \mathbb{R} \to \varphi_{11}) \to \varphi_{15} \to \varphi_{27} \to \varphi_{16}$\lbl{A1}}
\BIC{$\varphi_{15} \to \varphi_{27} \to \varphi_{16}$}
\BIC{$\varphi_{27} \to \varphi_{16}$}
\BIC{$\varphi_{16}$}
\AXC{$\varphi_{51}$\lbl{A1}}
\BIC{$(c \in \mathbb{R} \to d \in \mathbb{R} \to x \in \mathbb{R} \to c < d + x \to x + 0 = x) \to \varphi_{39}$}
\end{ptree}
\noindent$\Pi_{28}$:
\begin{ptree}
\AXC{$c \in \mathbb{R} \to d \in \mathbb{R} \to x \in \mathbb{R} \to c < d + x \to x \in \mathbb{R}$\lbl{A1}}
\AXC{$\stackrel{\Pi_{10}}{\vdots}$}\noLine
\UIC{$\varphi_{38}$}
\BIC{$c \in \mathbb{R} \to d \in \mathbb{R} \to x \in \mathbb{R} \to c < d + x \to x + 0 = x$}
\AXC{$\stackrel{\Pi_{27}}{\vdots}$}\noLine
\UIC{$(c \in \mathbb{R} \to d \in \mathbb{R} \to x \in \mathbb{R} \to c < d + x \to x + 0 = x) \to \varphi_{39}$}
\BIC{$\varphi_{39}$}
\end{ptree}
\noindent$\Pi_{29}$:
\begin{ptree}
\AXC{$\stackrel{\Pi_{8}}{\vdots}$}\noLine
\UIC{$\varphi_{25}$}
\AXC{$\stackrel{\Pi_{9}}{\vdots}$}\noLine
\UIC{$d + x + -d = x + 0 \to x + 0 = x \to d + x + -d = x$}
\AXC{$\stackrel{\Pi_{28}}{\vdots}$}\noLine
\UIC{$\varphi_{39}$}
\BIC{$c \in \mathbb{R} \to d \in \mathbb{R} \to x \in \mathbb{R} \to c < d + x \to d + x + -d = x$}
\AXC{$\varphi_{56}$\lbl{A1}}
\BIC{$\varphi_{25} \to \varphi_{44}$}
\BIC{$\varphi_{44}$}
\end{ptree}
\noindent$\Pi_{30}$:
\begin{ptree}
\AXC{$\stackrel{\Pi_{1}}{\vdots}$}\noLine
\UIC{$\varphi_{13}$}
\AXC{$\varphi_{13} \to d + x + -d = x \to c + -d < d + x + -d \to c + -d < x$\lbl{A2}}
\BIC{$d + x + -d = x \to c + -d < d + x + -d \to c + -d < x$}
\AXC{$\stackrel{\Pi_{29}}{\vdots}$}\noLine
\UIC{$\varphi_{44}$}
\BIC{$c \in \mathbb{R} \to d \in \mathbb{R} \to x \in \mathbb{R} \to c < d + x \to c + -d < x$}
\end{ptree}
\noindent$\Pi_{31}$:
\begin{ptree}
\AXC{$x = x$\lbl{A5}}
\AXC{$x = y \to x = x \to y = x$\lbl{A6}}
\AXC{$(x = y \to x = x \to y = x) \to x = x \to x = y \to y = x$\lbl{A1}}
\BIC{$x = x \to x = y \to y = x$}
\BIC{$x = y \to y = x$}
\UIC{$\forall y\,(x = y \to y = x)$}
\UIC{$\forall x\,\forall y\,(x = y \to y = x)$}
\end{ptree}
\noindent$\Pi_{32}$:
\begin{ptree}
\AXC{$\stackrel{\Pi_{31}}{\vdots}$}\noLine
\UIC{$\forall x\,\forall y\,(x = y \to y = x)$}
\AXC{$\forall x\,\forall y\,(x = y \to y = x) \to \forall y\,(c - d = y \to y = c - d)$\lbl{A2}}
\BIC{$\forall y\,(c - d = y \to y = c - d)$}
\AXC{$\varphi_{2}$\lbl{A2}}
\BIC{$c - d = c + -d \to c + -d = c - d$}
\end{ptree}
\noindent$\Pi_{33}$:
\begin{ptree}
\AXC{$\forall x\,\forall y\,(x - y = x + -y)$\lbl{Def$-$}}
\AXC{$\forall x\,\forall y\,(x - y = x + -y) \to \forall y\,(c - y = c + -y)$\lbl{A2}}
\BIC{$\forall y\,(c - y = c + -y)$}
\end{ptree}
\noindent$\Pi_{34}$:
\begin{ptree}
\AXC{$\stackrel{\Pi_{32}}{\vdots}$}\noLine
\UIC{$c - d = c + -d \to c + -d = c - d$}
\AXC{$\stackrel{\Pi_{33}}{\vdots}$}\noLine
\UIC{$\forall y\,(c - y = c + -y)$}
\AXC{$\forall y\,(c - y = c + -y) \to c - d = c + -d$\lbl{A2}}
\BIC{$c - d = c + -d$}
\AXC{$\varphi_{23}$\lbl{A1}}
\BIC{$(c - d = c + -d \to c + -d = c - d) \to c + -d = c - d$}
\BIC{$c + -d = c - d$}
\end{ptree}
\noindent$\Pi_{35}$:
\begin{ptree}
\AXC{$\stackrel{\Pi_{30}}{\vdots}$}\noLine
\UIC{$c \in \mathbb{R} \to d \in \mathbb{R} \to x \in \mathbb{R} \to c < d + x \to c + -d < x$}
\AXC{$\stackrel{\Pi_{34}}{\vdots}$}\noLine
\UIC{$c + -d = c - d$}
\AXC{$c + -d = c - d \to \varphi_{46}$\lbl{A1}}
\BIC{$\varphi_{46}$}
\BIC{$\varphi_{37}$}
\end{ptree}
\noindent 主木:
\begin{ptree}
\AXC{$z = w \to z < x \to w < x$\lbl{A6}}
\UIC{$\forall w\,(z = w \to z < x \to w < x)$}
\UIC{$\forall z\,\forall w\,(z = w \to z < x \to w < x)$}
\AXC{$\varphi_{12}$\lbl{A2}}
\BIC{$\forall w\,(c + -d = w \to c + -d < x \to w < x)$}
\AXC{$\varphi_{34}$\lbl{A2}}
\BIC{$c + -d = c - d \to c + -d < x \to c - d < x$}
\AXC{$\stackrel{\Pi_{35}}{\vdots}$}\noLine
\UIC{$\varphi_{37}$}
\BIC{$c \in \mathbb{R} \to d \in \mathbb{R} \to x \in \mathbb{R} \to c < d + x \to c - d < x$}
\UIC{$\forall x\,(c \in \mathbb{R} \to d \in \mathbb{R} \to x \in \mathbb{R} \to c < d + x \to c - d < x)$}
\UIC{$\varphi_{8}$}
\UIC{$\forall c\,(\varphi_{8})$}
\end{ptree}
\begin{abbrevlist}
\item $\varphi_{1} :\equiv \forall w\,(z = w \to x + (d + -d) = x + z \to x + (d + -d) = x + w)$
\item $\varphi_{2} :\equiv \forall y\,(c - d = y \to y = c - d) \to c - d = c + -d \to c + -d = c - d$
\item $\varphi_{3} :\equiv \forall z\,(x \in \mathbb{R} \to d \in \mathbb{R} \to z \in \mathbb{R} \to x + d + z = x + (d + z))$
\item $\varphi_{4} :\equiv \forall y\,\forall z\,(x \in \mathbb{R} \to y \in \mathbb{R} \to z \in \mathbb{R} \to x + y + z = x + (y + z))$
\item $\varphi_{5} :\equiv \forall w\,(z = w \to d + x + -d = z + -d \to d + x + -d = w + -d)$
\item $\varphi_{6} :\equiv x + d + -d = x + (d + -d) \to d + x + -d = x + (d + -d)$
\item $\varphi_{7} :\equiv \forall y\,\forall z\,(c \in \mathbb{R} \to y \in \mathbb{R} \to z \in \mathbb{R} \to c < y \to c + z < y + z)$
\item $\varphi_{8} :\equiv \forall d\,\forall x\,(c \in \mathbb{R} \to d \in \mathbb{R} \to x \in \mathbb{R} \to c < d + x \to c - d < x)$
\item $\varphi_{9} :\equiv \forall x\,\forall y\,\forall z\,(x \in \mathbb{R} \to y \in \mathbb{R} \to z \in \mathbb{R} \to x < y \to x + z < y + z)$
\item $\varphi_{10} :\equiv d \in \mathbb{R} \to x \in \mathbb{R} \to c < d + x \to x + d + -d = x + (d + -d)$
\item $\varphi_{11} :\equiv d \in \mathbb{R} \to x \in \mathbb{R} \to c < d + x \to d + x + -d = x + (d + -d)$
\item $\varphi_{12} :\equiv \forall z\,\forall w\,(z = w \to z < x \to w < x) \to \forall w\,(c + -d = w \to c + -d < x \to w < x)$
\item $\varphi_{13} :\equiv \forall w\,(d + x + -d = w \to c + -d < d + x + -d \to c + -d < w)$
\item $\varphi_{14} :\equiv d + -d = 0 \to x + (d + -d) = x + (d + -d) \to x + (d + -d) = x + 0$
\item $\varphi_{15} :\equiv c \in \mathbb{R} \to d \in \mathbb{R} \to x \in \mathbb{R} \to c < d + x \to x + (d + -d) = x + 0$
\item $\varphi_{16} :\equiv c \in \mathbb{R} \to d \in \mathbb{R} \to x \in \mathbb{R} \to c < d + x \to d + x + -d = x + 0$
\item $\varphi_{17} :\equiv x \in \mathbb{R} \to d \in \mathbb{R} \to -d \in \mathbb{R} \to x + d + -d = x + (d + -d)$
\item $\varphi_{18} :\equiv \forall x\,\forall y\,\forall z\,(x = y \to y = z \to x = z) \to \forall y\,\forall z\,(d + x + -d = y \to y = z \to d + x + -d = z)$
\item $\varphi_{19} :\equiv \forall z\,(d + x + -d = x + (d + -d) \to x + (d + -d) = z \to d + x + -d = z)$
\item $\varphi_{20} :\equiv \forall w\,(d + -d = w \to x + (d + -d) = x + (d + -d) \to x + (d + -d) = x + w)$
\item $\varphi_{21} :\equiv \forall z\,(d + x + -d = x + d + -d \to x + d + -d = z \to d + x + -d = z)$
\item $\varphi_{22} :\equiv \forall w\,(d + x = w \to d + x + -d = d + x + -d \to d + x + -d = w + -d)$
\item $\varphi_{23} :\equiv c - d = c + -d \to (c - d = c + -d \to c + -d = c - d) \to c + -d = c - d$
\item $\varphi_{24} :\equiv \forall z\,(c \in \mathbb{R} \to d + x \in \mathbb{R} \to z \in \mathbb{R} \to c < d + x \to c + z < d + x + z)$
\item $\varphi_{25} :\equiv c \in \mathbb{R} \to d \in \mathbb{R} \to x \in \mathbb{R} \to c < d + x \to c + -d < d + x + -d$
\item $\varphi_{26} :\equiv \forall y\,(d \in \mathbb{R} \to y \in \mathbb{R} \to d + y = y + d) \to d \in \mathbb{R} \to x \in \mathbb{R} \to d + x = x + d$
\item $\varphi_{27} :\equiv d + x + -d = x + (d + -d) \to x + (d + -d) = x + 0 \to d + x + -d = x + 0$
\item $\varphi_{28} :\equiv d + x = x + d \to d + x + -d = d + x + -d \to d + x + -d = x + d + -d$
\item $\varphi_{29} :\equiv c \in \mathbb{R} \to d \in \mathbb{R} \to x \in \mathbb{R} \to c < d + x \to d + x + -d = x + d + -d$
\item $\varphi_{30} :\equiv \forall y\,(d \in \mathbb{R} \to y \in \mathbb{R} \to d + y \in \mathbb{R}) \to d \in \mathbb{R} \to x \in \mathbb{R} \to d + x \in \mathbb{R}$
\item $\varphi_{31} :\equiv \forall x\,\forall y\,(x \in \mathbb{R} \to y \in \mathbb{R} \to x + y = y + x) \to \forall y\,(d \in \mathbb{R} \to y \in \mathbb{R} \to d + y = y + d)$
\item $\varphi_{32} :\equiv \forall x\,\forall y\,(x \in \mathbb{R} \to y \in \mathbb{R} \to x + y \in \mathbb{R}) \to \forall y\,(d \in \mathbb{R} \to y \in \mathbb{R} \to d + y \in \mathbb{R})$
\item $\varphi_{33} :\equiv c \in \mathbb{R} \to d + x \in \mathbb{R} \to -d \in \mathbb{R} \to c < d + x \to c + -d < d + x + -d$
\item $\varphi_{34} :\equiv \forall w\,(c + -d = w \to c + -d < x \to w < x) \to c + -d = c - d \to c + -d < x \to c - d < x$
\item $\varphi_{35} :\equiv \forall y\,\forall z\,(d + x + -d = y \to y = z \to d + x + -d = z) \to \forall z\,(d + x + -d = x + 0 \to x + 0 = z \to d + x + -d = z)$
\item $\varphi_{36} :\equiv \forall z\,(d + x + -d = x + 0 \to x + 0 = z \to d + x + -d = z) \to d + x + -d = x + 0 \to x + 0 = x \to d + x + -d = x$
\item $\varphi_{37} :\equiv (c + -d = c - d \to c + -d < x \to c - d < x) \to c \in \mathbb{R} \to d \in \mathbb{R} \to x \in \mathbb{R} \to c < d + x \to c - d < x$
\item $\varphi_{38} :\equiv (c \in \mathbb{R} \to d \in \mathbb{R} \to x \in \mathbb{R} \to c < d + x \to x \in \mathbb{R}) \to c \in \mathbb{R} \to d \in \mathbb{R} \to x \in \mathbb{R} \to c < d + x \to x + 0 = x$
\item $\varphi_{39} :\equiv (d + x + -d = x + 0 \to x + 0 = x \to d + x + -d = x) \to c \in \mathbb{R} \to d \in \mathbb{R} \to x \in \mathbb{R} \to c < d + x \to d + x + -d = x$
\item $\varphi_{40} :\equiv (c \in \mathbb{R} \to d \in \mathbb{R} \to x \in \mathbb{R} \to c < d + x \to d \in \mathbb{R}) \to c \in \mathbb{R} \to d \in \mathbb{R} \to x \in \mathbb{R} \to c < d + x \to d + -d = 0$
\item $\varphi_{41} :\equiv (c \in \mathbb{R} \to d \in \mathbb{R} \to x \in \mathbb{R} \to c < d + x \to x \in \mathbb{R}) \to c \in \mathbb{R} \to d \in \mathbb{R} \to x \in \mathbb{R} \to c < d + x \to d + x = x + d$
\item $\varphi_{42} :\equiv (c \in \mathbb{R} \to d \in \mathbb{R} \to x \in \mathbb{R} \to c < d + x \to d \in \mathbb{R}) \to c \in \mathbb{R} \to d \in \mathbb{R} \to x \in \mathbb{R} \to c < d + x \to -d \in \mathbb{R}$
\item $\varphi_{43} :\equiv (c \in \mathbb{R} \to d \in \mathbb{R} \to x \in \mathbb{R} \to c < d + x \to x \in \mathbb{R}) \to c \in \mathbb{R} \to d \in \mathbb{R} \to x \in \mathbb{R} \to c < d + x \to d + x \in \mathbb{R}$
\item $\varphi_{44} :\equiv (d + x + -d = x \to c + -d < d + x + -d \to c + -d < x) \to c \in \mathbb{R} \to d \in \mathbb{R} \to x \in \mathbb{R} \to c < d + x \to c + -d < x$
\item $\varphi_{45} :\equiv (c \in \mathbb{R} \to d \in \mathbb{R} \to x \in \mathbb{R} \to c < d + x \to -d \in \mathbb{R}) \to c \in \mathbb{R} \to \varphi_{10}$
\item $\varphi_{46} :\equiv (c \in \mathbb{R} \to d \in \mathbb{R} \to x \in \mathbb{R} \to c < d + x \to c + -d < x) \to \varphi_{37}$
\item $\varphi_{47} :\equiv (c \in \mathbb{R} \to d \in \mathbb{R} \to x \in \mathbb{R} \to c < d + x \to -d \in \mathbb{R}) \to (c \in \mathbb{R} \to d \in \mathbb{R} \to x \in \mathbb{R} \to c < d + x \to c < d + x) \to \varphi_{25}$
\item $\varphi_{48} :\equiv (d \in \mathbb{R} \to x \in \mathbb{R} \to d + x = x + d) \to (c \in \mathbb{R} \to d \in \mathbb{R} \to x \in \mathbb{R} \to c < d + x \to d \in \mathbb{R}) \to \varphi_{41}$
\item $\varphi_{49} :\equiv (d \in \mathbb{R} \to x \in \mathbb{R} \to d + x \in \mathbb{R}) \to (c \in \mathbb{R} \to d \in \mathbb{R} \to x \in \mathbb{R} \to c < d + x \to d \in \mathbb{R}) \to \varphi_{43}$
\item $\varphi_{50} :\equiv (c \in \mathbb{R} \to d \in \mathbb{R} \to x \in \mathbb{R} \to c < d + x \to d + -d = 0) \to x + (d + -d) = x + (d + -d) \to \varphi_{14} \to \varphi_{15}$
\item $\varphi_{51} :\equiv \varphi_{16} \to (c \in \mathbb{R} \to d \in \mathbb{R} \to x \in \mathbb{R} \to c < d + x \to x + 0 = x) \to \varphi_{39}$
\item $\varphi_{52} :\equiv (c \in \mathbb{R} \to d \in \mathbb{R} \to x \in \mathbb{R} \to c < d + x \to d + x = x + d) \to d + x + -d = d + x + -d \to \varphi_{28} \to \varphi_{29}$
\item $\varphi_{53} :\equiv (c \in \mathbb{R} \to \varphi_{10}) \to (d + x + -d = x + d + -d \to \varphi_{6}) \to c \in \mathbb{R} \to \varphi_{11}$
\item $\varphi_{54} :\equiv (c \in \mathbb{R} \to d \in \mathbb{R} \to x \in \mathbb{R} \to c < d + x \to x \in \mathbb{R}) \to (c \in \mathbb{R} \to d \in \mathbb{R} \to x \in \mathbb{R} \to c < d + x \to d \in \mathbb{R}) \to \varphi_{45}$
\item $\varphi_{55} :\equiv (c \in \mathbb{R} \to d \in \mathbb{R} \to x \in \mathbb{R} \to c < d + x \to d + x \in \mathbb{R}) \to \varphi_{47}$
\item $\varphi_{56} :\equiv (c \in \mathbb{R} \to d \in \mathbb{R} \to x \in \mathbb{R} \to c < d + x \to d + x + -d = x) \to \varphi_{25} \to \varphi_{44}$
\item $\varphi_{57} :\equiv \varphi_{33} \to (c \in \mathbb{R} \to d \in \mathbb{R} \to x \in \mathbb{R} \to c < d + x \to c \in \mathbb{R}) \to \varphi_{55}$
\end{abbrevlist}
\end{fullproof}

% 自動生成 (prover/build.js): LeftAbs  ノード数 182
\begin{fullproof}{[H5]}{$\forall x\,\forall c\,\forall d\,(x \in \mathbb{R} \to c \in \mathbb{R} \to d \in \mathbb{R} \to c - d < x \to x < c \to 0 < \lvert x - c \rvert \land \lvert x - c \rvert < d)$}{182}\phantomsection\label{tree:LeftAbs}
\noindent$\Pi_{1}$:
\begin{ptree}
\AXC{$z = w \to -(x - c) < z \to -(x - c) < w$\lbl{A6}}
\UIC{$\forall w\,(z = w \to -(x - c) < z \to -(x - c) < w)$}
\UIC{$\forall z\,\forall w\,(z = w \to -(x - c) < z \to -(x - c) < w)$}
\AXC{$\forall z\,\forall w\,(z = w \to -(x - c) < z \to -(x - c) < w) \to \varphi_{4}$\lbl{A2}}
\BIC{$\varphi_{4}$}
\AXC{$\varphi_{39}$\lbl{A2}}
\BIC{$-(-d) = d \to -(x - c) < -(-d) \to -(x - c) < d$}
\end{ptree}
\noindent$\Pi_{2}$:
\begin{ptree}
\AXC{$\varphi_{16}$\lbl{[E31]}}
\AXC{$\varphi_{16} \to \varphi_{12}$\lbl{A2}}
\BIC{$\varphi_{12}$}
\AXC{$\varphi_{12} \to \varphi_{8}$\lbl{A2}}
\BIC{$\varphi_{8}$}
\AXC{$\varphi_{8} \to c \in \mathbb{R} \to d \in \mathbb{R} \to x \in \mathbb{R} \to c - d < x \to -d < x - c$\lbl{A2}}
\BIC{$c \in \mathbb{R} \to d \in \mathbb{R} \to x \in \mathbb{R} \to c - d < x \to -d < x - c$}
\AXC{$\varphi_{62}$\lbl{A1}}
\BIC{$\varphi_{61}$}
\end{ptree}
\noindent$\Pi_{3}$:
\begin{ptree}
\AXC{$x \in \mathbb{R} \to c \in \mathbb{R} \to d \in \mathbb{R} \to c - d < x \to x < c \to d \in \mathbb{R}$\lbl{A1}}
\AXC{$x \in \mathbb{R} \to c \in \mathbb{R} \to d \in \mathbb{R} \to c - d < x \to x < c \to c \in \mathbb{R}$\lbl{A1}}
\AXC{$\stackrel{\Pi_{2}}{\vdots}$}\noLine
\UIC{$\varphi_{61}$}
\BIC{$\varphi_{59}$}
\BIC{$\varphi_{55}$}
\end{ptree}
\noindent$\Pi_{4}$:
\begin{ptree}
\AXC{$x \in \mathbb{R} \to c \in \mathbb{R} \to d \in \mathbb{R} \to c - d < x \to x < c \to c - d < x$\lbl{A1}}
\AXC{$x \in \mathbb{R} \to c \in \mathbb{R} \to d \in \mathbb{R} \to c - d < x \to x < c \to x \in \mathbb{R}$\lbl{A1}}
\AXC{$\stackrel{\Pi_{3}}{\vdots}$}\noLine
\UIC{$\varphi_{55}$}
\BIC{$\varphi_{48}$}
\BIC{$\varphi_{17}$}
\end{ptree}
\noindent$\Pi_{5}$:
\begin{ptree}
\AXC{$\forall x\,\forall y\,(x \in \mathbb{R} \to y \in \mathbb{R} \to x - y \in \mathbb{R})$\lbl{[F1]}}
\AXC{$\varphi_{30}$\lbl{A2}}
\BIC{$\forall y\,(x \in \mathbb{R} \to y \in \mathbb{R} \to x - y \in \mathbb{R})$}
\AXC{$\varphi_{28}$\lbl{A2}}
\BIC{$x \in \mathbb{R} \to c \in \mathbb{R} \to x - c \in \mathbb{R}$}
\AXC{$(x \in \mathbb{R} \to c \in \mathbb{R} \to x - c \in \mathbb{R}) \to \varphi_{54}$\lbl{A1}}
\BIC{$\varphi_{54}$}
\end{ptree}
\noindent$\Pi_{6}$:
\begin{ptree}
\AXC{$x \in \mathbb{R} \to c \in \mathbb{R} \to d \in \mathbb{R} \to c - d < x \to x < c \to c \in \mathbb{R}$\lbl{A1}}
\AXC{$x \in \mathbb{R} \to c \in \mathbb{R} \to d \in \mathbb{R} \to c - d < x \to x < c \to x \in \mathbb{R}$\lbl{A1}}
\AXC{$\stackrel{\Pi_{5}}{\vdots}$}\noLine
\UIC{$\varphi_{54}$}
\BIC{$\varphi_{45}$}
\BIC{$\varphi_{15}$}
\end{ptree}
\noindent$\Pi_{7}$:
\begin{ptree}
\AXC{$\forall x \in \mathbb{R}\,(-x \in \mathbb{R})$\lbl{(R1)}}
\AXC{$\forall x \in \mathbb{R}\,(-x \in \mathbb{R}) \to d \in \mathbb{R} \to -d \in \mathbb{R}$\lbl{A2}}
\BIC{$d \in \mathbb{R} \to -d \in \mathbb{R}$}
\AXC{$(d \in \mathbb{R} \to -d \in \mathbb{R}) \to \varphi_{46}$\lbl{A1}}
\BIC{$\varphi_{46}$}
\end{ptree}
\noindent$\Pi_{8}$:
\begin{ptree}
\AXC{$\forall x\,\forall y\,(x \in \mathbb{R} \to y \in \mathbb{R} \to x < y \to -y < -x)$\lbl{[F22]}}
\AXC{$\varphi_{40}$\lbl{A2}}
\BIC{$\forall y\,(-d \in \mathbb{R} \to y \in \mathbb{R} \to -d < y \to -y < -(-d))$}
\AXC{$\forall y\,(-d \in \mathbb{R} \to y \in \mathbb{R} \to -d < y \to -y < -(-d)) \to \varphi_{25}$\lbl{A2}}
\BIC{$\varphi_{25}$}
\end{ptree}
\noindent$\Pi_{9}$:
\begin{ptree}
\AXC{$\stackrel{\Pi_{4}}{\vdots}$}\noLine
\UIC{$\varphi_{17}$}
\AXC{$\stackrel{\Pi_{6}}{\vdots}$}\noLine
\UIC{$\varphi_{15}$}
\AXC{$x \in \mathbb{R} \to c \in \mathbb{R} \to d \in \mathbb{R} \to c - d < x \to x < c \to d \in \mathbb{R}$\lbl{A1}}
\AXC{$\stackrel{\Pi_{7}}{\vdots}$}\noLine
\UIC{$\varphi_{46}$}
\BIC{$x \in \mathbb{R} \to c \in \mathbb{R} \to d \in \mathbb{R} \to c - d < x \to x < c \to -d \in \mathbb{R}$}
\AXC{$\stackrel{\Pi_{8}}{\vdots}$}\noLine
\UIC{$\varphi_{25}$}
\AXC{$\varphi_{25} \to \varphi_{60}$\lbl{A1}}
\BIC{$\varphi_{60}$}
\BIC{$\varphi_{15} \to \varphi_{17} \to x \in \mathbb{R} \to \varphi_{14}$}
\BIC{$\varphi_{17} \to x \in \mathbb{R} \to \varphi_{14}$}
\BIC{$x \in \mathbb{R} \to \varphi_{14}$}
\end{ptree}
\noindent$\Pi_{10}$:
\begin{ptree}
\AXC{$\forall x \in \mathbb{R}\,(-(-x) = x)$\lbl{[F5]}}
\AXC{$\forall x \in \mathbb{R}\,(-(-x) = x) \to d \in \mathbb{R} \to -(-d) = d$\lbl{A2}}
\BIC{$d \in \mathbb{R} \to -(-d) = d$}
\AXC{$(d \in \mathbb{R} \to -(-d) = d) \to \varphi_{44}$\lbl{A1}}
\BIC{$\varphi_{44}$}
\end{ptree}
\noindent$\Pi_{11}$:
\begin{ptree}
\AXC{$\stackrel{\Pi_{9}}{\vdots}$}\noLine
\UIC{$x \in \mathbb{R} \to \varphi_{14}$}
\AXC{$x \in \mathbb{R} \to c \in \mathbb{R} \to d \in \mathbb{R} \to c - d < x \to x < c \to d \in \mathbb{R}$\lbl{A1}}
\AXC{$\stackrel{\Pi_{10}}{\vdots}$}\noLine
\UIC{$\varphi_{44}$}
\BIC{$\varphi_{13}$}
\AXC{$\varphi_{13} \to (x \in \mathbb{R} \to \varphi_{14}) \to \varphi_{41}$\lbl{A1}}
\BIC{$(x \in \mathbb{R} \to \varphi_{14}) \to \varphi_{41}$}
\BIC{$\varphi_{41}$}
\end{ptree}
\noindent$\Pi_{12}$:
\begin{ptree}
\AXC{$\forall x\,\forall y\,(x \in \mathbb{R} \to y \in \mathbb{R} \to -(y - x) = x - y)$\lbl{[F14]}}
\AXC{$\varphi_{34}$\lbl{A2}}
\BIC{$\forall y\,(c \in \mathbb{R} \to y \in \mathbb{R} \to -(y - c) = c - y)$}
\AXC{$\varphi_{31}$\lbl{A2}}
\BIC{$c \in \mathbb{R} \to x \in \mathbb{R} \to -(x - c) = c - x$}
\AXC{$(c \in \mathbb{R} \to x \in \mathbb{R} \to -(x - c) = c - x) \to \varphi_{56}$\lbl{A1}}
\BIC{$\varphi_{56}$}
\end{ptree}
\noindent$\Pi_{13}$:
\begin{ptree}
\AXC{$x \in \mathbb{R} \to c \in \mathbb{R} \to d \in \mathbb{R} \to c - d < x \to x < c \to x \in \mathbb{R}$\lbl{A1}}
\AXC{$x \in \mathbb{R} \to c \in \mathbb{R} \to d \in \mathbb{R} \to c - d < x \to x < c \to c \in \mathbb{R}$\lbl{A1}}
\AXC{$\stackrel{\Pi_{12}}{\vdots}$}\noLine
\UIC{$\varphi_{56}$}
\BIC{$\varphi_{49}$}
\BIC{$\varphi_{21}$}
\end{ptree}
\noindent$\Pi_{14}$:
\begin{ptree}
\AXC{$\stackrel{\Pi_{1}}{\vdots}$}\noLine
\UIC{$-(-d) = d \to -(x - c) < -(-d) \to -(x - c) < d$}
\AXC{$\stackrel{\Pi_{11}}{\vdots}$}\noLine
\UIC{$\varphi_{41}$}
\BIC{$\varphi_{18}$}
\AXC{$\stackrel{\Pi_{13}}{\vdots}$}\noLine
\UIC{$\varphi_{21}$}
\AXC{$\varphi_{21} \to \varphi_{18} \to \varphi_{38}$\lbl{A1}}
\BIC{$\varphi_{18} \to \varphi_{38}$}
\BIC{$\varphi_{38}$}
\end{ptree}
\noindent$\Pi_{15}$:
\begin{ptree}
\AXC{$x = x$\lbl{A5}}
\AXC{$x = y \to x = x \to y = x$\lbl{A6}}
\AXC{$(x = y \to x = x \to y = x) \to x = x \to x = y \to y = x$\lbl{A1}}
\BIC{$x = x \to x = y \to y = x$}
\BIC{$x = y \to y = x$}
\UIC{$\forall y\,(x = y \to y = x)$}
\UIC{$\forall x\,\forall y\,(x = y \to y = x)$}
\end{ptree}
\noindent$\Pi_{16}$:
\begin{ptree}
\AXC{$y = z \to x = y \to x = z$\lbl{A6}}
\AXC{$(y = z \to x = y \to x = z) \to x = y \to y = z \to x = z$\lbl{A1}}
\BIC{$x = y \to y = z \to x = z$}
\UIC{$\forall z\,(x = y \to y = z \to x = z)$}
\UIC{$\forall y\,\forall z\,(x = y \to y = z \to x = z)$}
\UIC{$\forall x\,\forall y\,\forall z\,(x = y \to y = z \to x = z)$}
\AXC{$\varphi_{7}$\lbl{A2}}
\BIC{$\forall y\,\forall z\,(\lvert x - c \rvert = y \to y = z \to \lvert x - c \rvert = z)$}
\end{ptree}
\noindent$\Pi_{17}$:
\begin{ptree}
\AXC{$\stackrel{\Pi_{16}}{\vdots}$}\noLine
\UIC{$\forall y\,\forall z\,(\lvert x - c \rvert = y \to y = z \to \lvert x - c \rvert = z)$}
\AXC{$\forall y\,\forall z\,(\lvert x - c \rvert = y \to y = z \to \lvert x - c \rvert = z) \to \varphi_{2}$\lbl{A2}}
\BIC{$\varphi_{2}$}
\AXC{$\varphi_{2} \to \varphi_{6}$\lbl{A2}}
\BIC{$\varphi_{6}$}
\end{ptree}
\noindent$\Pi_{18}$:
\begin{ptree}
\AXC{$\forall x\,\forall y\,(x \in \mathbb{R} \to y \in \mathbb{R} \to x < y \to 0 < y - x)$\lbl{[F24]}}
\AXC{$\varphi_{35}$\lbl{A2}}
\BIC{$\forall y\,(x \in \mathbb{R} \to y \in \mathbb{R} \to x < y \to 0 < y - x)$}
\AXC{$\varphi_{33}$\lbl{A2}}
\BIC{$x \in \mathbb{R} \to c \in \mathbb{R} \to x < c \to 0 < c - x$}
\AXC{$(x \in \mathbb{R} \to c \in \mathbb{R} \to x < c \to 0 < c - x) \to \varphi_{58}$\lbl{A1}}
\BIC{$\varphi_{58}$}
\end{ptree}
\noindent$\Pi_{19}$:
\begin{ptree}
\AXC{$x \in \mathbb{R} \to c \in \mathbb{R} \to d \in \mathbb{R} \to c - d < x \to x < c \to c \in \mathbb{R}$\lbl{A1}}
\AXC{$x \in \mathbb{R} \to c \in \mathbb{R} \to d \in \mathbb{R} \to c - d < x \to x < c \to x \in \mathbb{R}$\lbl{A1}}
\AXC{$\stackrel{\Pi_{18}}{\vdots}$}\noLine
\UIC{$\varphi_{58}$}
\BIC{$\varphi_{53}$}
\BIC{$\varphi_{42}$}
\end{ptree}
\noindent$\Pi_{20}$:
\begin{ptree}
\AXC{$x \in \mathbb{R} \to c \in \mathbb{R} \to d \in \mathbb{R} \to c - d < x \to x < c \to x < c$\lbl{A1}}
\AXC{$\stackrel{\Pi_{19}}{\vdots}$}\noLine
\UIC{$\varphi_{42}$}
\BIC{$x \in \mathbb{R} \to c \in \mathbb{R} \to d \in \mathbb{R} \to c - d < x \to x < c \to 0 < c - x$}
\end{ptree}
\noindent$\Pi_{21}$:
\begin{ptree}
\AXC{$\forall x\,(0 \leq x \to \lvert x \rvert = x)$\lbl{[A1]}}
\AXC{$\forall x\,(0 \leq x \to \lvert x \rvert = x) \to 0 \leq c - x \to \lvert c - x \rvert = c - x$\lbl{A2}}
\BIC{$0 \leq c - x \to \lvert c - x \rvert = c - x$}
\end{ptree}
\noindent$\Pi_{22}$:
\begin{ptree}
\AXC{$\stackrel{\Pi_{21}}{\vdots}$}\noLine
\UIC{$0 \leq c - x \to \lvert c - x \rvert = c - x$}
\AXC{$(0 \leq c - x \to \lvert c - x \rvert = c - x) \to \varphi_{47}$\lbl{A1}}
\BIC{$\varphi_{47}$}
\end{ptree}
\noindent$\Pi_{23}$:
\begin{ptree}
\AXC{$\stackrel{\Pi_{20}}{\vdots}$}\noLine
\UIC{$x \in \mathbb{R} \to c \in \mathbb{R} \to d \in \mathbb{R} \to c - d < x \to x < c \to 0 < c - x$}
\AXC{$\varphi_{43}$\lbl{A1}}
\BIC{$x \in \mathbb{R} \to c \in \mathbb{R} \to d \in \mathbb{R} \to c - d < x \to x < c \to 0 \leq c - x$}
\AXC{$\stackrel{\Pi_{22}}{\vdots}$}\noLine
\UIC{$\varphi_{47}$}
\BIC{$\varphi_{22}$}
\end{ptree}
\noindent$\Pi_{24}$:
\begin{ptree}
\AXC{$\forall x\,\forall y\,(x \in \mathbb{R} \to y \in \mathbb{R} \to \lvert x - y \rvert = \lvert y - x \rvert)$\lbl{[A13]}}
\AXC{$\varphi_{37}$\lbl{A2}}
\BIC{$\forall y\,(x \in \mathbb{R} \to y \in \mathbb{R} \to \lvert x - y \rvert = \lvert y - x \rvert)$}
\AXC{$\varphi_{36}$\lbl{A2}}
\BIC{$x \in \mathbb{R} \to c \in \mathbb{R} \to \lvert x - c \rvert = \lvert c - x \rvert$}
\end{ptree}
\noindent$\Pi_{25}$:
\begin{ptree}
\AXC{$\stackrel{\Pi_{24}}{\vdots}$}\noLine
\UIC{$x \in \mathbb{R} \to c \in \mathbb{R} \to \lvert x - c \rvert = \lvert c - x \rvert$}
\AXC{$(x \in \mathbb{R} \to c \in \mathbb{R} \to \lvert x - c \rvert = \lvert c - x \rvert) \to \varphi_{57}$\lbl{A1}}
\BIC{$\varphi_{57}$}
\end{ptree}
\noindent$\Pi_{26}$:
\begin{ptree}
\AXC{$x \in \mathbb{R} \to c \in \mathbb{R} \to d \in \mathbb{R} \to c - d < x \to x < c \to c \in \mathbb{R}$\lbl{A1}}
\AXC{$x \in \mathbb{R} \to c \in \mathbb{R} \to d \in \mathbb{R} \to c - d < x \to x < c \to x \in \mathbb{R}$\lbl{A1}}
\AXC{$\stackrel{\Pi_{25}}{\vdots}$}\noLine
\UIC{$\varphi_{57}$}
\BIC{$\varphi_{50}$}
\BIC{$x \in \mathbb{R} \to \varphi_{9}$}
\end{ptree}
\noindent$\Pi_{27}$:
\begin{ptree}
\AXC{$\stackrel{\Pi_{17}}{\vdots}$}\noLine
\UIC{$\varphi_{6}$}
\AXC{$\stackrel{\Pi_{23}}{\vdots}$}\noLine
\UIC{$\varphi_{22}$}
\AXC{$\stackrel{\Pi_{26}}{\vdots}$}\noLine
\UIC{$x \in \mathbb{R} \to \varphi_{9}$}
\AXC{$(x \in \mathbb{R} \to \varphi_{9}) \to \varphi_{22} \to \varphi_{6} \to \varphi_{23}$\lbl{A1}}
\BIC{$\varphi_{22} \to \varphi_{6} \to \varphi_{23}$}
\BIC{$\varphi_{6} \to \varphi_{23}$}
\BIC{$\varphi_{23}$}
\end{ptree}
\noindent$\Pi_{28}$:
\begin{ptree}
\AXC{$\stackrel{\Pi_{15}}{\vdots}$}\noLine
\UIC{$\forall x\,\forall y\,(x = y \to y = x)$}
\AXC{$\varphi_{1}$\lbl{A2}}
\BIC{$\forall y\,(\lvert x - c \rvert = y \to y = \lvert x - c \rvert)$}
\AXC{$\varphi_{11}$\lbl{A2}}
\BIC{$\lvert x - c \rvert = c - x \to c - x = \lvert x - c \rvert$}
\AXC{$\stackrel{\Pi_{27}}{\vdots}$}\noLine
\UIC{$\varphi_{23}$}
\AXC{$\varphi_{23} \to (\lvert x - c \rvert = c - x \to c - x = \lvert x - c \rvert) \to \varphi_{24}$\lbl{A1}}
\BIC{$(\lvert x - c \rvert = c - x \to c - x = \lvert x - c \rvert) \to \varphi_{24}$}
\BIC{$\varphi_{24}$}
\end{ptree}
\noindent$\Pi_{29}$:
\begin{ptree}
\AXC{$z = w \to z < d \to w < d$\lbl{A6}}
\UIC{$\forall w\,(z = w \to z < d \to w < d)$}
\UIC{$\forall z\,\forall w\,(z = w \to z < d \to w < d)$}
\AXC{$\varphi_{10}$\lbl{A2}}
\BIC{$\forall w\,(-(x - c) = w \to -(x - c) < d \to w < d)$}
\AXC{$\varphi_{32}$\lbl{A2}}
\BIC{$-(x - c) = c - x \to -(x - c) < d \to c - x < d$}
\AXC{$\stackrel{\Pi_{14}}{\vdots}$}\noLine
\UIC{$\varphi_{38}$}
\BIC{$x \in \mathbb{R} \to c \in \mathbb{R} \to d \in \mathbb{R} \to c - d < x \to x < c \to c - x < d$}
\AXC{$\stackrel{\Pi_{28}}{\vdots}$}\noLine
\UIC{$\varphi_{24}$}
\AXC{$\varphi_{24} \to \varphi_{51}$\lbl{A1}}
\BIC{$\varphi_{51}$}
\BIC{$(c - x = \lvert x - c \rvert \to c - x < d \to \lvert x - c \rvert < d) \to \varphi_{19}$}
\end{ptree}
\noindent$\Pi_{30}$:
\begin{ptree}
\AXC{$z = w \to z < d \to w < d$\lbl{A6}}
\UIC{$\forall w\,(z = w \to z < d \to w < d)$}
\UIC{$\forall z\,\forall w\,(z = w \to z < d \to w < d)$}
\AXC{$\varphi_{3}$\lbl{A2}}
\BIC{$\forall w\,(c - x = w \to c - x < d \to w < d)$}
\AXC{$\varphi_{27}$\lbl{A2}}
\BIC{$c - x = \lvert x - c \rvert \to c - x < d \to \lvert x - c \rvert < d$}
\AXC{$\stackrel{\Pi_{29}}{\vdots}$}\noLine
\UIC{$(c - x = \lvert x - c \rvert \to c - x < d \to \lvert x - c \rvert < d) \to \varphi_{19}$}
\BIC{$\varphi_{19}$}
\end{ptree}
\noindent$\Pi_{31}$:
\begin{ptree}
\AXC{$\stackrel{\Pi_{15}}{\vdots}$}\noLine
\UIC{$\forall x\,\forall y\,(x = y \to y = x)$}
\AXC{$\varphi_{1}$\lbl{A2}}
\BIC{$\forall y\,(\lvert x - c \rvert = y \to y = \lvert x - c \rvert)$}
\AXC{$\varphi_{11}$\lbl{A2}}
\BIC{$\lvert x - c \rvert = c - x \to c - x = \lvert x - c \rvert$}
\AXC{$\stackrel{\Pi_{27}}{\vdots}$}\noLine
\UIC{$\varphi_{23}$}
\AXC{$\varphi_{23} \to (\lvert x - c \rvert = c - x \to c - x = \lvert x - c \rvert) \to \varphi_{24}$\lbl{A1}}
\BIC{$(\lvert x - c \rvert = c - x \to c - x = \lvert x - c \rvert) \to \varphi_{24}$}
\BIC{$\varphi_{24}$}
\end{ptree}
\noindent$\Pi_{32}$:
\begin{ptree}
\AXC{$\stackrel{\Pi_{20}}{\vdots}$}\noLine
\UIC{$x \in \mathbb{R} \to c \in \mathbb{R} \to d \in \mathbb{R} \to c - d < x \to x < c \to 0 < c - x$}
\AXC{$\stackrel{\Pi_{31}}{\vdots}$}\noLine
\UIC{$\varphi_{24}$}
\AXC{$\varphi_{24} \to \varphi_{52}$\lbl{A1}}
\BIC{$\varphi_{52}$}
\BIC{$(c - x = \lvert x - c \rvert \to 0 < c - x \to 0 < \lvert x - c \rvert) \to \varphi_{20}$}
\end{ptree}
\noindent$\Pi_{33}$:
\begin{ptree}
\AXC{$z = w \to 0 < z \to 0 < w$\lbl{A6}}
\UIC{$\forall w\,(z = w \to 0 < z \to 0 < w)$}
\UIC{$\forall z\,\forall w\,(z = w \to 0 < z \to 0 < w)$}
\AXC{$\varphi_{5}$\lbl{A2}}
\BIC{$\forall w\,(c - x = w \to 0 < c - x \to 0 < w)$}
\AXC{$\varphi_{29}$\lbl{A2}}
\BIC{$c - x = \lvert x - c \rvert \to 0 < c - x \to 0 < \lvert x - c \rvert$}
\AXC{$\stackrel{\Pi_{32}}{\vdots}$}\noLine
\UIC{$(c - x = \lvert x - c \rvert \to 0 < c - x \to 0 < \lvert x - c \rvert) \to \varphi_{20}$}
\BIC{$\varphi_{20}$}
\end{ptree}
\noindent 主木:
\begin{ptree}
\AXC{$\stackrel{\Pi_{30}}{\vdots}$}\noLine
\UIC{$\varphi_{19}$}
\AXC{$\stackrel{\Pi_{33}}{\vdots}$}\noLine
\UIC{$\varphi_{20}$}
\AXC{$\varphi_{20} \to \varphi_{19} \to x \in \mathbb{R} \to \varphi_{26}$\lbl{A1}}
\BIC{$\varphi_{19} \to x \in \mathbb{R} \to \varphi_{26}$}
\BIC{$x \in \mathbb{R} \to \varphi_{26}$}
\UIC{$\forall d\,(x \in \mathbb{R} \to \varphi_{26})$}
\UIC{$\forall c\,\forall d\,(x \in \mathbb{R} \to \varphi_{26})$}
\UIC{$\forall x\,\forall c\,\forall d\,(x \in \mathbb{R} \to \varphi_{26})$}
\end{ptree}
\begin{abbrevlist}
\item $\varphi_{1} :\equiv \forall x\,\forall y\,(x = y \to y = x) \to \forall y\,(\lvert x - c \rvert = y \to y = \lvert x - c \rvert)$
\item $\varphi_{2} :\equiv \forall z\,(\lvert x - c \rvert = \lvert c - x \rvert \to \lvert c - x \rvert = z \to \lvert x - c \rvert = z)$
\item $\varphi_{3} :\equiv \forall z\,\forall w\,(z = w \to z < d \to w < d) \to \forall w\,(c - x = w \to c - x < d \to w < d)$
\item $\varphi_{4} :\equiv \forall w\,(-(-d) = w \to -(x - c) < -(-d) \to -(x - c) < w)$
\item $\varphi_{5} :\equiv \forall z\,\forall w\,(z = w \to 0 < z \to 0 < w) \to \forall w\,(c - x = w \to 0 < c - x \to 0 < w)$
\item $\varphi_{6} :\equiv \lvert x - c \rvert = \lvert c - x \rvert \to \lvert c - x \rvert = c - x \to \lvert x - c \rvert = c - x$
\item $\varphi_{7} :\equiv \forall x\,\forall y\,\forall z\,(x = y \to y = z \to x = z) \to \forall y\,\forall z\,(\lvert x - c \rvert = y \to y = z \to \lvert x - c \rvert = z)$
\item $\varphi_{8} :\equiv \forall y\,(c \in \mathbb{R} \to d \in \mathbb{R} \to y \in \mathbb{R} \to c - d < y \to -d < y - c)$
\item $\varphi_{9} :\equiv c \in \mathbb{R} \to d \in \mathbb{R} \to c - d < x \to x < c \to \lvert x - c \rvert = \lvert c - x \rvert$
\item $\varphi_{10} :\equiv \forall z\,\forall w\,(z = w \to z < d \to w < d) \to \forall w\,(-(x - c) = w \to -(x - c) < d \to w < d)$
\item $\varphi_{11} :\equiv \forall y\,(\lvert x - c \rvert = y \to y = \lvert x - c \rvert) \to \lvert x - c \rvert = c - x \to c - x = \lvert x - c \rvert$
\item $\varphi_{12} :\equiv \forall e\,\forall y\,(c \in \mathbb{R} \to e \in \mathbb{R} \to y \in \mathbb{R} \to c - e < y \to -e < y - c)$
\item $\varphi_{13} :\equiv x \in \mathbb{R} \to c \in \mathbb{R} \to d \in \mathbb{R} \to c - d < x \to x < c \to -(-d) = d$
\item $\varphi_{14} :\equiv c \in \mathbb{R} \to d \in \mathbb{R} \to c - d < x \to x < c \to -(x - c) < -(-d)$
\item $\varphi_{15} :\equiv x \in \mathbb{R} \to c \in \mathbb{R} \to d \in \mathbb{R} \to c - d < x \to x < c \to x - c \in \mathbb{R}$
\item $\varphi_{16} :\equiv \forall x\,\forall e\,\forall y\,(x \in \mathbb{R} \to e \in \mathbb{R} \to y \in \mathbb{R} \to x - e < y \to -e < y - x)$
\item $\varphi_{17} :\equiv x \in \mathbb{R} \to c \in \mathbb{R} \to d \in \mathbb{R} \to c - d < x \to x < c \to -d < x - c$
\item $\varphi_{18} :\equiv x \in \mathbb{R} \to c \in \mathbb{R} \to d \in \mathbb{R} \to c - d < x \to x < c \to -(x - c) < d$
\item $\varphi_{19} :\equiv x \in \mathbb{R} \to c \in \mathbb{R} \to d \in \mathbb{R} \to c - d < x \to x < c \to \lvert x - c \rvert < d$
\item $\varphi_{20} :\equiv x \in \mathbb{R} \to c \in \mathbb{R} \to d \in \mathbb{R} \to c - d < x \to x < c \to 0 < \lvert x - c \rvert$
\item $\varphi_{21} :\equiv x \in \mathbb{R} \to c \in \mathbb{R} \to d \in \mathbb{R} \to c - d < x \to x < c \to -(x - c) = c - x$
\item $\varphi_{22} :\equiv x \in \mathbb{R} \to c \in \mathbb{R} \to d \in \mathbb{R} \to c - d < x \to x < c \to \lvert c - x \rvert = c - x$
\item $\varphi_{23} :\equiv x \in \mathbb{R} \to c \in \mathbb{R} \to d \in \mathbb{R} \to c - d < x \to x < c \to \lvert x - c \rvert = c - x$
\item $\varphi_{24} :\equiv x \in \mathbb{R} \to c \in \mathbb{R} \to d \in \mathbb{R} \to c - d < x \to x < c \to c - x = \lvert x - c \rvert$
\item $\varphi_{25} :\equiv -d \in \mathbb{R} \to x - c \in \mathbb{R} \to -d < x - c \to -(x - c) < -(-d)$
\item $\varphi_{26} :\equiv c \in \mathbb{R} \to d \in \mathbb{R} \to c - d < x \to x < c \to 0 < \lvert x - c \rvert \land \lvert x - c \rvert < d$
\item $\varphi_{27} :\equiv \forall w\,(c - x = w \to c - x < d \to w < d) \to c - x = \lvert x - c \rvert \to c - x < d \to \lvert x - c \rvert < d$
\item $\varphi_{28} :\equiv \forall y\,(x \in \mathbb{R} \to y \in \mathbb{R} \to x - y \in \mathbb{R}) \to x \in \mathbb{R} \to c \in \mathbb{R} \to x - c \in \mathbb{R}$
\item $\varphi_{29} :\equiv \forall w\,(c - x = w \to 0 < c - x \to 0 < w) \to c - x = \lvert x - c \rvert \to 0 < c - x \to 0 < \lvert x - c \rvert$
\item $\varphi_{30} :\equiv \forall x\,\forall y\,(x \in \mathbb{R} \to y \in \mathbb{R} \to x - y \in \mathbb{R}) \to \forall y\,(x \in \mathbb{R} \to y \in \mathbb{R} \to x - y \in \mathbb{R})$
\item $\varphi_{31} :\equiv \forall y\,(c \in \mathbb{R} \to y \in \mathbb{R} \to -(y - c) = c - y) \to c \in \mathbb{R} \to x \in \mathbb{R} \to -(x - c) = c - x$
\item $\varphi_{32} :\equiv \forall w\,(-(x - c) = w \to -(x - c) < d \to w < d) \to -(x - c) = c - x \to -(x - c) < d \to c - x < d$
\item $\varphi_{33} :\equiv \forall y\,(x \in \mathbb{R} \to y \in \mathbb{R} \to x < y \to 0 < y - x) \to x \in \mathbb{R} \to c \in \mathbb{R} \to x < c \to 0 < c - x$
\item $\varphi_{34} :\equiv \forall x\,\forall y\,(x \in \mathbb{R} \to y \in \mathbb{R} \to -(y - x) = x - y) \to \forall y\,(c \in \mathbb{R} \to y \in \mathbb{R} \to -(y - c) = c - y)$
\item $\varphi_{35} :\equiv \forall x\,\forall y\,(x \in \mathbb{R} \to y \in \mathbb{R} \to x < y \to 0 < y - x) \to \forall y\,(x \in \mathbb{R} \to y \in \mathbb{R} \to x < y \to 0 < y - x)$
\item $\varphi_{36} :\equiv \forall y\,(x \in \mathbb{R} \to y \in \mathbb{R} \to \lvert x - y \rvert = \lvert y - x \rvert) \to x \in \mathbb{R} \to c \in \mathbb{R} \to \lvert x - c \rvert = \lvert c - x \rvert$
\item $\varphi_{37} :\equiv \forall x\,\forall y\,(x \in \mathbb{R} \to y \in \mathbb{R} \to \lvert x - y \rvert = \lvert y - x \rvert) \to \forall y\,(x \in \mathbb{R} \to y \in \mathbb{R} \to \lvert x - y \rvert = \lvert y - x \rvert)$
\item $\varphi_{38} :\equiv (-(x - c) = c - x \to -(x - c) < d \to c - x < d) \to x \in \mathbb{R} \to c \in \mathbb{R} \to d \in \mathbb{R} \to c - d < x \to x < c \to c - x < d$
\item $\varphi_{39} :\equiv \varphi_{4} \to -(-d) = d \to -(x - c) < -(-d) \to -(x - c) < d$
\item $\varphi_{40} :\equiv \forall x\,\forall y\,(x \in \mathbb{R} \to y \in \mathbb{R} \to x < y \to -y < -x) \to \forall y\,(-d \in \mathbb{R} \to y \in \mathbb{R} \to -d < y \to -y < -(-d))$
\item $\varphi_{41} :\equiv (-(-d) = d \to -(x - c) < -(-d) \to -(x - c) < d) \to \varphi_{18}$
\item $\varphi_{42} :\equiv (x \in \mathbb{R} \to c \in \mathbb{R} \to d \in \mathbb{R} \to c - d < x \to x < c \to x < c) \to x \in \mathbb{R} \to c \in \mathbb{R} \to d \in \mathbb{R} \to c - d < x \to x < c \to 0 < c - x$
\item $\varphi_{43} :\equiv (x \in \mathbb{R} \to c \in \mathbb{R} \to d \in \mathbb{R} \to c - d < x \to x < c \to 0 < c - x) \to x \in \mathbb{R} \to c \in \mathbb{R} \to d \in \mathbb{R} \to c - d < x \to x < c \to 0 \leq c - x$
\item $\varphi_{44} :\equiv (x \in \mathbb{R} \to c \in \mathbb{R} \to d \in \mathbb{R} \to c - d < x \to x < c \to d \in \mathbb{R}) \to \varphi_{13}$
\item $\varphi_{45} :\equiv (x \in \mathbb{R} \to c \in \mathbb{R} \to d \in \mathbb{R} \to c - d < x \to x < c \to c \in \mathbb{R}) \to \varphi_{15}$
\item $\varphi_{46} :\equiv (x \in \mathbb{R} \to c \in \mathbb{R} \to d \in \mathbb{R} \to c - d < x \to x < c \to d \in \mathbb{R}) \to x \in \mathbb{R} \to c \in \mathbb{R} \to d \in \mathbb{R} \to c - d < x \to x < c \to -d \in \mathbb{R}$
\item $\varphi_{47} :\equiv (x \in \mathbb{R} \to c \in \mathbb{R} \to d \in \mathbb{R} \to c - d < x \to x < c \to 0 \leq c - x) \to \varphi_{22}$
\item $\varphi_{48} :\equiv (x \in \mathbb{R} \to c \in \mathbb{R} \to d \in \mathbb{R} \to c - d < x \to x < c \to c - d < x) \to \varphi_{17}$
\item $\varphi_{49} :\equiv (x \in \mathbb{R} \to c \in \mathbb{R} \to d \in \mathbb{R} \to c - d < x \to x < c \to x \in \mathbb{R}) \to \varphi_{21}$
\item $\varphi_{50} :\equiv (x \in \mathbb{R} \to c \in \mathbb{R} \to d \in \mathbb{R} \to c - d < x \to x < c \to c \in \mathbb{R}) \to x \in \mathbb{R} \to \varphi_{9}$
\item $\varphi_{51} :\equiv (x \in \mathbb{R} \to c \in \mathbb{R} \to d \in \mathbb{R} \to c - d < x \to x < c \to c - x < d) \to (c - x = \lvert x - c \rvert \to c - x < d \to \lvert x - c \rvert < d) \to \varphi_{19}$
\item $\varphi_{52} :\equiv (x \in \mathbb{R} \to c \in \mathbb{R} \to d \in \mathbb{R} \to c - d < x \to x < c \to 0 < c - x) \to (c - x = \lvert x - c \rvert \to 0 < c - x \to 0 < \lvert x - c \rvert) \to \varphi_{20}$
\item $\varphi_{53} :\equiv (x \in \mathbb{R} \to c \in \mathbb{R} \to d \in \mathbb{R} \to c - d < x \to x < c \to c \in \mathbb{R}) \to \varphi_{42}$
\item $\varphi_{54} :\equiv (x \in \mathbb{R} \to c \in \mathbb{R} \to d \in \mathbb{R} \to c - d < x \to x < c \to x \in \mathbb{R}) \to \varphi_{45}$
\item $\varphi_{55} :\equiv (x \in \mathbb{R} \to c \in \mathbb{R} \to d \in \mathbb{R} \to c - d < x \to x < c \to x \in \mathbb{R}) \to \varphi_{48}$
\item $\varphi_{56} :\equiv (x \in \mathbb{R} \to c \in \mathbb{R} \to d \in \mathbb{R} \to c - d < x \to x < c \to c \in \mathbb{R}) \to \varphi_{49}$
\item $\varphi_{57} :\equiv (x \in \mathbb{R} \to c \in \mathbb{R} \to d \in \mathbb{R} \to c - d < x \to x < c \to x \in \mathbb{R}) \to \varphi_{50}$
\item $\varphi_{58} :\equiv (x \in \mathbb{R} \to c \in \mathbb{R} \to d \in \mathbb{R} \to c - d < x \to x < c \to x \in \mathbb{R}) \to \varphi_{53}$
\item $\varphi_{59} :\equiv (x \in \mathbb{R} \to c \in \mathbb{R} \to d \in \mathbb{R} \to c - d < x \to x < c \to d \in \mathbb{R}) \to \varphi_{55}$
\item $\varphi_{60} :\equiv (x \in \mathbb{R} \to c \in \mathbb{R} \to d \in \mathbb{R} \to c - d < x \to x < c \to -d \in \mathbb{R}) \to \varphi_{15} \to \varphi_{17} \to x \in \mathbb{R} \to \varphi_{14}$
\item $\varphi_{61} :\equiv (x \in \mathbb{R} \to c \in \mathbb{R} \to d \in \mathbb{R} \to c - d < x \to x < c \to c \in \mathbb{R}) \to \varphi_{59}$
\item $\varphi_{62} :\equiv (c \in \mathbb{R} \to d \in \mathbb{R} \to x \in \mathbb{R} \to c - d < x \to -d < x - c) \to \varphi_{61}$
\end{abbrevlist}
\end{fullproof}

% 自動生成 (prover/build.js): NeSides  ノード数 21
\begin{fullproof}{[H6]}{$\forall x\,\forall c\,(x \in \mathbb{R} \to c \in \mathbb{R} \to 0 < \lvert x - c \rvert \to x < c \lor c < x)$}{21}\phantomsection\label{tree:NeSides}
\noindent$\Pi_{1}$:
\begin{ptree}
\AXC{$\forall y\,\forall c\,(y \in \mathbb{R} \to c \in \mathbb{R} \to 0 < \lvert y - c \rvert \to y \neq c)$\lbl{[E37]}}
\AXC{$\varphi_{5}$\lbl{A2}}
\BIC{$\forall c\,(x \in \mathbb{R} \to c \in \mathbb{R} \to 0 < \lvert x - c \rvert \to x \neq c)$}
\AXC{$\varphi_{3}$\lbl{A2}}
\BIC{$x \in \mathbb{R} \to c \in \mathbb{R} \to 0 < \lvert x - c \rvert \to x \neq c$}
\end{ptree}
\noindent$\Pi_{2}$:
\begin{ptree}
\AXC{$\forall x\,\forall y\,(x \in \mathbb{R} \to y \in \mathbb{R} \to x < y \lor x = y \lor y < x)$\lbl{(R13)}}
\AXC{$\varphi_{4}$\lbl{A2}}
\BIC{$\forall y\,(x \in \mathbb{R} \to y \in \mathbb{R} \to x < y \lor x = y \lor y < x)$}
\AXC{$\varphi_{2}$\lbl{A2}}
\BIC{$x \in \mathbb{R} \to c \in \mathbb{R} \to x < c \lor x = c \lor c < x$}
\AXC{$\varphi_{9}$\lbl{A1}}
\BIC{$(x \in \mathbb{R} \to c \in \mathbb{R} \to 0 < \lvert x - c \rvert \to x \in \mathbb{R}) \to \varphi_{7}$}
\end{ptree}
\noindent$\Pi_{3}$:
\begin{ptree}
\AXC{$x \in \mathbb{R} \to c \in \mathbb{R} \to 0 < \lvert x - c \rvert \to x \in \mathbb{R}$\lbl{A1}}
\AXC{$\stackrel{\Pi_{2}}{\vdots}$}\noLine
\UIC{$(x \in \mathbb{R} \to c \in \mathbb{R} \to 0 < \lvert x - c \rvert \to x \in \mathbb{R}) \to \varphi_{7}$}
\BIC{$\varphi_{7}$}
\end{ptree}
\noindent 主木:
\begin{ptree}
\AXC{$\stackrel{\Pi_{1}}{\vdots}$}\noLine
\UIC{$x \in \mathbb{R} \to c \in \mathbb{R} \to 0 < \lvert x - c \rvert \to x \neq c$}
\AXC{$x \in \mathbb{R} \to c \in \mathbb{R} \to 0 < \lvert x - c \rvert \to c \in \mathbb{R}$\lbl{A1}}
\AXC{$\stackrel{\Pi_{3}}{\vdots}$}\noLine
\UIC{$\varphi_{7}$}
\BIC{$x \in \mathbb{R} \to c \in \mathbb{R} \to 0 < \lvert x - c \rvert \to x < c \lor x = c \lor c < x$}
\AXC{$\varphi_{8}$\lbl{A1}}
\BIC{$\varphi_{6}$}
\BIC{$x \in \mathbb{R} \to c \in \mathbb{R} \to 0 < \lvert x - c \rvert \to x < c \lor c < x$}
\UIC{$\forall c\,(x \in \mathbb{R} \to c \in \mathbb{R} \to 0 < \lvert x - c \rvert \to x < c \lor c < x)$}
\UIC{$\varphi_{1}$}
\end{ptree}
\begin{abbrevlist}
\item $\varphi_{1} :\equiv \forall x\,\forall c\,(x \in \mathbb{R} \to c \in \mathbb{R} \to 0 < \lvert x - c \rvert \to x < c \lor c < x)$
\item $\varphi_{2} :\equiv \forall y\,(x \in \mathbb{R} \to y \in \mathbb{R} \to x < y \lor x = y \lor y < x) \to x \in \mathbb{R} \to c \in \mathbb{R} \to x < c \lor x = c \lor c < x$
\item $\varphi_{3} :\equiv \forall c\,(x \in \mathbb{R} \to c \in \mathbb{R} \to 0 < \lvert x - c \rvert \to x \neq c) \to x \in \mathbb{R} \to c \in \mathbb{R} \to 0 < \lvert x - c \rvert \to x \neq c$
\item $\varphi_{4} :\equiv \forall x\,\forall y\,(x \in \mathbb{R} \to y \in \mathbb{R} \to x < y \lor x = y \lor y < x) \to \forall y\,(x \in \mathbb{R} \to y \in \mathbb{R} \to x < y \lor x = y \lor y < x)$
\item $\varphi_{5} :\equiv \forall y\,\forall c\,(y \in \mathbb{R} \to c \in \mathbb{R} \to 0 < \lvert y - c \rvert \to y \neq c) \to \forall c\,(x \in \mathbb{R} \to c \in \mathbb{R} \to 0 < \lvert x - c \rvert \to x \neq c)$
\item $\varphi_{6} :\equiv (x \in \mathbb{R} \to c \in \mathbb{R} \to 0 < \lvert x - c \rvert \to x \neq c) \to x \in \mathbb{R} \to c \in \mathbb{R} \to 0 < \lvert x - c \rvert \to x < c \lor c < x$
\item $\varphi_{7} :\equiv (x \in \mathbb{R} \to c \in \mathbb{R} \to 0 < \lvert x - c \rvert \to c \in \mathbb{R}) \to x \in \mathbb{R} \to c \in \mathbb{R} \to 0 < \lvert x - c \rvert \to x < c \lor x = c \lor c < x$
\item $\varphi_{8} :\equiv (x \in \mathbb{R} \to c \in \mathbb{R} \to 0 < \lvert x - c \rvert \to x < c \lor x = c \lor c < x) \to \varphi_{6}$
\item $\varphi_{9} :\equiv (x \in \mathbb{R} \to c \in \mathbb{R} \to x < c \lor x = c \lor c < x) \to (x \in \mathbb{R} \to c \in \mathbb{R} \to 0 < \lvert x - c \rvert \to x \in \mathbb{R}) \to \varphi_{7}$
\end{abbrevlist}
\end{fullproof}

% 自動生成 (prover/build.js): SideR  ノード数 74
\begin{fullproof}{[H7]}{$\forall x\,\forall c\,\forall d\,(x \in \mathbb{R} \to c \in \mathbb{R} \to d \in \mathbb{R} \to c < x \to \lvert x - c \rvert < d \to x < c + d)$}{74}\phantomsection\label{tree:SideR}
\noindent$\Pi_{1}$:
\begin{ptree}
\AXC{$\forall x\,\forall y\,(x \in \mathbb{R} \to y \in \mathbb{R} \to x < y \to 0 < y - x)$\lbl{[F24]}}
\AXC{$\varphi_{20}$\lbl{A2}}
\BIC{$\forall y\,(c \in \mathbb{R} \to y \in \mathbb{R} \to c < y \to 0 < y - c)$}
\AXC{$\varphi_{19}$\lbl{A2}}
\BIC{$c \in \mathbb{R} \to x \in \mathbb{R} \to c < x \to 0 < x - c$}
\end{ptree}
\noindent$\Pi_{2}$:
\begin{ptree}
\AXC{$\varphi_{6}$\lbl{A1}}
\AXC{$\varphi_{7}$\lbl{A1}}
\AXC{$\stackrel{\Pi_{1}}{\vdots}$}\noLine
\UIC{$c \in \mathbb{R} \to x \in \mathbb{R} \to c < x \to 0 < x - c$}
\AXC{$(c \in \mathbb{R} \to x \in \mathbb{R} \to c < x \to 0 < x - c) \to \varphi_{7} \to \varphi_{6} \to \varphi_{23}$\lbl{A1}}
\BIC{$\varphi_{7} \to \varphi_{6} \to \varphi_{23}$}
\BIC{$\varphi_{6} \to \varphi_{23}$}
\BIC{$\varphi_{23}$}
\end{ptree}
\noindent$\Pi_{3}$:
\begin{ptree}
\AXC{$x \in \mathbb{R} \to c \in \mathbb{R} \to d \in \mathbb{R} \to c < x \to \lvert x - c \rvert < d \to c < x$\lbl{A1}}
\AXC{$\stackrel{\Pi_{2}}{\vdots}$}\noLine
\UIC{$\varphi_{23}$}
\BIC{$\varphi_{10}$}
\AXC{$\varphi_{10} \to \varphi_{5}$\lbl{A1}}
\BIC{$\varphi_{5}$}
\end{ptree}
\noindent$\Pi_{4}$:
\begin{ptree}
\AXC{$\forall x\,(0 \leq x \to \lvert x \rvert = x)$\lbl{[A1]}}
\AXC{$\forall x\,(0 \leq x \to \lvert x \rvert = x) \to 0 \leq x - c \to \lvert x - c \rvert = x - c$\lbl{A2}}
\BIC{$0 \leq x - c \to \lvert x - c \rvert = x - c$}
\end{ptree}
\noindent$\Pi_{5}$:
\begin{ptree}
\AXC{$\varphi_{15}$\lbl{A1}}
\AXC{$\stackrel{\Pi_{3}}{\vdots}$}\noLine
\UIC{$\varphi_{5}$}
\AXC{$\stackrel{\Pi_{4}}{\vdots}$}\noLine
\UIC{$0 \leq x - c \to \lvert x - c \rvert = x - c$}
\AXC{$(0 \leq x - c \to \lvert x - c \rvert = x - c) \to \varphi_{5} \to x \in \mathbb{R} \to \varphi_{3}$\lbl{A1}}
\BIC{$\varphi_{5} \to x \in \mathbb{R} \to \varphi_{3}$}
\BIC{$x \in \mathbb{R} \to \varphi_{3}$}
\AXC{$\varphi_{24}$\lbl{A1}}
\BIC{$\varphi_{15} \to (\lvert x - c \rvert = x - c \to \lvert x - c \rvert < d \to x - c < d) \to \varphi_{11}$}
\BIC{$(\lvert x - c \rvert = x - c \to \lvert x - c \rvert < d \to x - c < d) \to \varphi_{11}$}
\end{ptree}
\noindent$\Pi_{6}$:
\begin{ptree}
\AXC{$\varphi_{8}$\lbl{A1}}
\AXC{$\varphi_{7}$\lbl{A1}}
\AXC{$\varphi_{6}$\lbl{A1}}
\AXC{$\forall x\,(\varphi_{2})$\lbl{[B7]}}
\AXC{$\forall x\,(\varphi_{2}) \to \varphi_{2}$\lbl{A2}}
\BIC{$\varphi_{2}$}
\AXC{$\varphi_{22}$\lbl{A2}}
\BIC{$\forall y\,(x \in \mathbb{R} \to c \in \mathbb{R} \to y \in \mathbb{R} \to x - c < y \to x < y + c)$}
\AXC{$\varphi_{21}$\lbl{A2}}
\BIC{$x \in \mathbb{R} \to c \in \mathbb{R} \to d \in \mathbb{R} \to x - c < d \to x < d + c$}
\AXC{$\varphi_{25}$\lbl{A1}}
\BIC{$\varphi_{6} \to \varphi_{7} \to \varphi_{8} \to \varphi_{11} \to \varphi_{12}$}
\BIC{$\varphi_{7} \to \varphi_{8} \to \varphi_{11} \to \varphi_{12}$}
\BIC{$\varphi_{8} \to \varphi_{11} \to \varphi_{12}$}
\BIC{$\varphi_{11} \to \varphi_{12}$}
\end{ptree}
\noindent$\Pi_{7}$:
\begin{ptree}
\AXC{$z = w \to z < d \to w < d$\lbl{A6}}
\UIC{$\forall w\,(z = w \to z < d \to w < d)$}
\UIC{$\forall z\,\forall w\,(z = w \to z < d \to w < d)$}
\AXC{$\varphi_{4}$\lbl{A2}}
\BIC{$\forall w\,(\lvert x - c \rvert = w \to \lvert x - c \rvert < d \to w < d)$}
\AXC{$\varphi_{18}$\lbl{A2}}
\BIC{$\lvert x - c \rvert = x - c \to \lvert x - c \rvert < d \to x - c < d$}
\AXC{$\stackrel{\Pi_{5}}{\vdots}$}\noLine
\UIC{$(\lvert x - c \rvert = x - c \to \lvert x - c \rvert < d \to x - c < d) \to \varphi_{11}$}
\BIC{$\varphi_{11}$}
\AXC{$\stackrel{\Pi_{6}}{\vdots}$}\noLine
\UIC{$\varphi_{11} \to \varphi_{12}$}
\BIC{$\varphi_{12}$}
\end{ptree}
\noindent$\Pi_{8}$:
\begin{ptree}
\AXC{$\forall x\,\forall y\,(x \in \mathbb{R} \to y \in \mathbb{R} \to x + y = y + x)$\lbl{(R3)}}
\AXC{$\varphi_{17}$\lbl{A2}}
\BIC{$\forall y\,(d \in \mathbb{R} \to y \in \mathbb{R} \to d + y = y + d)$}
\AXC{$\varphi_{16}$\lbl{A2}}
\BIC{$d \in \mathbb{R} \to c \in \mathbb{R} \to d + c = c + d$}
\AXC{$(d \in \mathbb{R} \to c \in \mathbb{R} \to d + c = c + d) \to \varphi_{8} \to \varphi_{7} \to \varphi_{14}$\lbl{A1}}
\BIC{$\varphi_{8} \to \varphi_{7} \to \varphi_{14}$}
\end{ptree}
\noindent$\Pi_{9}$:
\begin{ptree}
\AXC{$\stackrel{\Pi_{7}}{\vdots}$}\noLine
\UIC{$\varphi_{12}$}
\AXC{$\varphi_{7}$\lbl{A1}}
\AXC{$\varphi_{8}$\lbl{A1}}
\AXC{$\stackrel{\Pi_{8}}{\vdots}$}\noLine
\UIC{$\varphi_{8} \to \varphi_{7} \to \varphi_{14}$}
\BIC{$\varphi_{7} \to \varphi_{14}$}
\BIC{$\varphi_{14}$}
\AXC{$\varphi_{14} \to \varphi_{12} \to (d + c = c + d \to x < d + c \to x < c + d) \to \varphi_{13}$\lbl{A1}}
\BIC{$\varphi_{12} \to (d + c = c + d \to x < d + c \to x < c + d) \to \varphi_{13}$}
\BIC{$(d + c = c + d \to x < d + c \to x < c + d) \to \varphi_{13}$}
\end{ptree}
\noindent 主木:
\begin{ptree}
\AXC{$z = w \to x < z \to x < w$\lbl{A6}}
\UIC{$\forall w\,(z = w \to x < z \to x < w)$}
\UIC{$\forall z\,\forall w\,(z = w \to x < z \to x < w)$}
\AXC{$\varphi_{1}$\lbl{A2}}
\BIC{$\forall w\,(d + c = w \to x < d + c \to x < w)$}
\AXC{$\varphi_{9}$\lbl{A2}}
\BIC{$d + c = c + d \to x < d + c \to x < c + d$}
\AXC{$\stackrel{\Pi_{9}}{\vdots}$}\noLine
\UIC{$(d + c = c + d \to x < d + c \to x < c + d) \to \varphi_{13}$}
\BIC{$\varphi_{13}$}
\UIC{$\forall d\,(\varphi_{13})$}
\UIC{$\forall c\,\forall d\,(\varphi_{13})$}
\UIC{$\forall x\,\forall c\,\forall d\,(\varphi_{13})$}
\end{ptree}
\begin{abbrevlist}
\item $\varphi_{1} :\equiv \forall z\,\forall w\,(z = w \to x < z \to x < w) \to \forall w\,(d + c = w \to x < d + c \to x < w)$
\item $\varphi_{2} :\equiv \forall a\,\forall y\,(x \in \mathbb{R} \to a \in \mathbb{R} \to y \in \mathbb{R} \to x - a < y \to x < y + a)$
\item $\varphi_{3} :\equiv c \in \mathbb{R} \to d \in \mathbb{R} \to c < x \to \lvert x - c \rvert < d \to \lvert x - c \rvert = x - c$
\item $\varphi_{4} :\equiv \forall z\,\forall w\,(z = w \to z < d \to w < d) \to \forall w\,(\lvert x - c \rvert = w \to \lvert x - c \rvert < d \to w < d)$
\item $\varphi_{5} :\equiv x \in \mathbb{R} \to c \in \mathbb{R} \to d \in \mathbb{R} \to c < x \to \lvert x - c \rvert < d \to 0 \leq x - c$
\item $\varphi_{6} :\equiv x \in \mathbb{R} \to c \in \mathbb{R} \to d \in \mathbb{R} \to c < x \to \lvert x - c \rvert < d \to x \in \mathbb{R}$
\item $\varphi_{7} :\equiv x \in \mathbb{R} \to c \in \mathbb{R} \to d \in \mathbb{R} \to c < x \to \lvert x - c \rvert < d \to c \in \mathbb{R}$
\item $\varphi_{8} :\equiv x \in \mathbb{R} \to c \in \mathbb{R} \to d \in \mathbb{R} \to c < x \to \lvert x - c \rvert < d \to d \in \mathbb{R}$
\item $\varphi_{9} :\equiv \forall w\,(d + c = w \to x < d + c \to x < w) \to d + c = c + d \to x < d + c \to x < c + d$
\item $\varphi_{10} :\equiv x \in \mathbb{R} \to c \in \mathbb{R} \to d \in \mathbb{R} \to c < x \to \lvert x - c \rvert < d \to 0 < x - c$
\item $\varphi_{11} :\equiv x \in \mathbb{R} \to c \in \mathbb{R} \to d \in \mathbb{R} \to c < x \to \lvert x - c \rvert < d \to x - c < d$
\item $\varphi_{12} :\equiv x \in \mathbb{R} \to c \in \mathbb{R} \to d \in \mathbb{R} \to c < x \to \lvert x - c \rvert < d \to x < d + c$
\item $\varphi_{13} :\equiv x \in \mathbb{R} \to c \in \mathbb{R} \to d \in \mathbb{R} \to c < x \to \lvert x - c \rvert < d \to x < c + d$
\item $\varphi_{14} :\equiv x \in \mathbb{R} \to c \in \mathbb{R} \to d \in \mathbb{R} \to c < x \to \lvert x - c \rvert < d \to d + c = c + d$
\item $\varphi_{15} :\equiv x \in \mathbb{R} \to c \in \mathbb{R} \to d \in \mathbb{R} \to c < x \to \lvert x - c \rvert < d \to \lvert x - c \rvert < d$
\item $\varphi_{16} :\equiv \forall y\,(d \in \mathbb{R} \to y \in \mathbb{R} \to d + y = y + d) \to d \in \mathbb{R} \to c \in \mathbb{R} \to d + c = c + d$
\item $\varphi_{17} :\equiv \forall x\,\forall y\,(x \in \mathbb{R} \to y \in \mathbb{R} \to x + y = y + x) \to \forall y\,(d \in \mathbb{R} \to y \in \mathbb{R} \to d + y = y + d)$
\item $\varphi_{18} :\equiv \forall w\,(\lvert x - c \rvert = w \to \lvert x - c \rvert < d \to w < d) \to \lvert x - c \rvert = x - c \to \lvert x - c \rvert < d \to x - c < d$
\item $\varphi_{19} :\equiv \forall y\,(c \in \mathbb{R} \to y \in \mathbb{R} \to c < y \to 0 < y - c) \to c \in \mathbb{R} \to x \in \mathbb{R} \to c < x \to 0 < x - c$
\item $\varphi_{20} :\equiv \forall x\,\forall y\,(x \in \mathbb{R} \to y \in \mathbb{R} \to x < y \to 0 < y - x) \to \forall y\,(c \in \mathbb{R} \to y \in \mathbb{R} \to c < y \to 0 < y - c)$
\item $\varphi_{21} :\equiv \forall y\,(x \in \mathbb{R} \to c \in \mathbb{R} \to y \in \mathbb{R} \to x - c < y \to x < y + c) \to x \in \mathbb{R} \to c \in \mathbb{R} \to d \in \mathbb{R} \to x - c < d \to x < d + c$
\item $\varphi_{22} :\equiv \varphi_{2} \to \forall y\,(x \in \mathbb{R} \to c \in \mathbb{R} \to y \in \mathbb{R} \to x - c < y \to x < y + c)$
\item $\varphi_{23} :\equiv (x \in \mathbb{R} \to c \in \mathbb{R} \to d \in \mathbb{R} \to c < x \to \lvert x - c \rvert < d \to c < x) \to \varphi_{10}$
\item $\varphi_{24} :\equiv (x \in \mathbb{R} \to \varphi_{3}) \to \varphi_{15} \to (\lvert x - c \rvert = x - c \to \lvert x - c \rvert < d \to x - c < d) \to \varphi_{11}$
\item $\varphi_{25} :\equiv (x \in \mathbb{R} \to c \in \mathbb{R} \to d \in \mathbb{R} \to x - c < d \to x < d + c) \to \varphi_{6} \to \varphi_{7} \to \varphi_{8} \to \varphi_{11} \to \varphi_{12}$
\end{abbrevlist}
\end{fullproof}

% 自動生成 (prover/build.js): SideL  ノード数 95
\begin{fullproof}{[H8]}{$\forall x\,\forall c\,\forall d\,(x \in \mathbb{R} \to c \in \mathbb{R} \to d \in \mathbb{R} \to x < c \to \lvert x - c \rvert < d \to c - d < x)$}{95}\phantomsection\label{tree:SideL}
\noindent$\Pi_{1}$:
\begin{ptree}
\AXC{$y = z \to x = y \to x = z$\lbl{A6}}
\AXC{$(y = z \to x = y \to x = z) \to x = y \to y = z \to x = z$\lbl{A1}}
\BIC{$x = y \to y = z \to x = z$}
\UIC{$\forall z\,(x = y \to y = z \to x = z)$}
\UIC{$\forall y\,\forall z\,(x = y \to y = z \to x = z)$}
\UIC{$\forall x\,\forall y\,\forall z\,(x = y \to y = z \to x = z)$}
\AXC{$\varphi_{3}$\lbl{A2}}
\BIC{$\forall y\,\forall z\,(\lvert x - c \rvert = y \to y = z \to \lvert x - c \rvert = z)$}
\end{ptree}
\noindent$\Pi_{2}$:
\begin{ptree}
\AXC{$\stackrel{\Pi_{1}}{\vdots}$}\noLine
\UIC{$\forall y\,\forall z\,(\lvert x - c \rvert = y \to y = z \to \lvert x - c \rvert = z)$}
\AXC{$\forall y\,\forall z\,(\lvert x - c \rvert = y \to y = z \to \lvert x - c \rvert = z) \to \varphi_{1}$\lbl{A2}}
\BIC{$\varphi_{1}$}
\AXC{$\varphi_{1} \to \varphi_{2}$\lbl{A2}}
\BIC{$\varphi_{2}$}
\end{ptree}
\noindent$\Pi_{3}$:
\begin{ptree}
\AXC{$\forall x\,\forall y\,(x \in \mathbb{R} \to y \in \mathbb{R} \to x < y \to 0 < y - x)$\lbl{[F24]}}
\AXC{$\varphi_{22}$\lbl{A2}}
\BIC{$\forall y\,(x \in \mathbb{R} \to y \in \mathbb{R} \to x < y \to 0 < y - x)$}
\AXC{$\varphi_{21}$\lbl{A2}}
\BIC{$x \in \mathbb{R} \to c \in \mathbb{R} \to x < c \to 0 < c - x$}
\end{ptree}
\noindent$\Pi_{4}$:
\begin{ptree}
\AXC{$\varphi_{11}$\lbl{A1}}
\AXC{$\varphi_{12}$\lbl{A1}}
\AXC{$\stackrel{\Pi_{3}}{\vdots}$}\noLine
\UIC{$x \in \mathbb{R} \to c \in \mathbb{R} \to x < c \to 0 < c - x$}
\AXC{$(x \in \mathbb{R} \to c \in \mathbb{R} \to x < c \to 0 < c - x) \to \varphi_{12} \to \varphi_{11} \to \varphi_{29}$\lbl{A1}}
\BIC{$\varphi_{12} \to \varphi_{11} \to \varphi_{29}$}
\BIC{$\varphi_{11} \to \varphi_{29}$}
\BIC{$\varphi_{29}$}
\end{ptree}
\noindent$\Pi_{5}$:
\begin{ptree}
\AXC{$x \in \mathbb{R} \to c \in \mathbb{R} \to d \in \mathbb{R} \to x < c \to \lvert x - c \rvert < d \to x < c$\lbl{A1}}
\AXC{$\stackrel{\Pi_{4}}{\vdots}$}\noLine
\UIC{$\varphi_{29}$}
\BIC{$\varphi_{15}$}
\AXC{$\varphi_{15} \to \varphi_{10}$\lbl{A1}}
\BIC{$\varphi_{10}$}
\end{ptree}
\noindent$\Pi_{6}$:
\begin{ptree}
\AXC{$\forall x\,(0 \leq x \to \lvert x \rvert = x)$\lbl{[A1]}}
\AXC{$\forall x\,(0 \leq x \to \lvert x \rvert = x) \to 0 \leq c - x \to \lvert c - x \rvert = c - x$\lbl{A2}}
\BIC{$0 \leq c - x \to \lvert c - x \rvert = c - x$}
\end{ptree}
\noindent$\Pi_{7}$:
\begin{ptree}
\AXC{$\stackrel{\Pi_{5}}{\vdots}$}\noLine
\UIC{$\varphi_{10}$}
\AXC{$\stackrel{\Pi_{6}}{\vdots}$}\noLine
\UIC{$0 \leq c - x \to \lvert c - x \rvert = c - x$}
\AXC{$(0 \leq c - x \to \lvert c - x \rvert = c - x) \to \varphi_{10} \to x \in \mathbb{R} \to \varphi_{6}$\lbl{A1}}
\BIC{$\varphi_{10} \to x \in \mathbb{R} \to \varphi_{6}$}
\BIC{$x \in \mathbb{R} \to \varphi_{6}$}
\end{ptree}
\noindent$\Pi_{8}$:
\begin{ptree}
\AXC{$\varphi_{11}$\lbl{A1}}
\AXC{$\varphi_{12}$\lbl{A1}}
\AXC{$\forall x\,\forall y\,(x \in \mathbb{R} \to y \in \mathbb{R} \to \lvert x - y \rvert = \lvert y - x \rvert)$\lbl{[A13]}}
\AXC{$\varphi_{24}$\lbl{A2}}
\BIC{$\forall y\,(x \in \mathbb{R} \to y \in \mathbb{R} \to \lvert x - y \rvert = \lvert y - x \rvert)$}
\AXC{$\varphi_{23}$\lbl{A2}}
\BIC{$x \in \mathbb{R} \to c \in \mathbb{R} \to \lvert x - c \rvert = \lvert c - x \rvert$}
\AXC{$\varphi_{30}$\lbl{A1}}
\BIC{$\varphi_{12} \to \varphi_{11} \to x \in \mathbb{R} \to \varphi_{14}$}
\BIC{$\varphi_{11} \to x \in \mathbb{R} \to \varphi_{14}$}
\BIC{$x \in \mathbb{R} \to \varphi_{14}$}
\end{ptree}
\noindent$\Pi_{9}$:
\begin{ptree}
\AXC{$\stackrel{\Pi_{2}}{\vdots}$}\noLine
\UIC{$\varphi_{2}$}
\AXC{$\stackrel{\Pi_{7}}{\vdots}$}\noLine
\UIC{$x \in \mathbb{R} \to \varphi_{6}$}
\AXC{$\stackrel{\Pi_{8}}{\vdots}$}\noLine
\UIC{$x \in \mathbb{R} \to \varphi_{14}$}
\AXC{$(x \in \mathbb{R} \to \varphi_{14}) \to (x \in \mathbb{R} \to \varphi_{6}) \to \varphi_{2} \to x \in \mathbb{R} \to \varphi_{7}$\lbl{A1}}
\BIC{$(x \in \mathbb{R} \to \varphi_{6}) \to \varphi_{2} \to x \in \mathbb{R} \to \varphi_{7}$}
\BIC{$\varphi_{2} \to x \in \mathbb{R} \to \varphi_{7}$}
\BIC{$x \in \mathbb{R} \to \varphi_{7}$}
\AXC{$\varphi_{31}$\lbl{A1}}
\BIC{$\varphi_{19} \to (\lvert x - c \rvert = c - x \to \lvert x - c \rvert < d \to c - x < d) \to \varphi_{16}$}
\end{ptree}
\noindent$\Pi_{10}$:
\begin{ptree}
\AXC{$z = w \to z < d \to w < d$\lbl{A6}}
\UIC{$\forall w\,(z = w \to z < d \to w < d)$}
\UIC{$\forall z\,\forall w\,(z = w \to z < d \to w < d)$}
\AXC{$\varphi_{9}$\lbl{A2}}
\BIC{$\forall w\,(\lvert x - c \rvert = w \to \lvert x - c \rvert < d \to w < d)$}
\AXC{$\varphi_{20}$\lbl{A2}}
\BIC{$\lvert x - c \rvert = c - x \to \lvert x - c \rvert < d \to c - x < d$}
\AXC{$\varphi_{19}$\lbl{A1}}
\AXC{$\stackrel{\Pi_{9}}{\vdots}$}\noLine
\UIC{$\varphi_{19} \to (\lvert x - c \rvert = c - x \to \lvert x - c \rvert < d \to c - x < d) \to \varphi_{16}$}
\BIC{$(\lvert x - c \rvert = c - x \to \lvert x - c \rvert < d \to c - x < d) \to \varphi_{16}$}
\BIC{$\varphi_{16}$}
\end{ptree}
\noindent$\Pi_{11}$:
\begin{ptree}
\AXC{$\varphi_{13}$\lbl{A1}}
\AXC{$\varphi_{12}$\lbl{A1}}
\AXC{$\varphi_{11}$\lbl{A1}}
\AXC{$\varphi_{8}$\lbl{[B7]}}
\AXC{$\varphi_{8} \to \varphi_{4}$\lbl{A2}}
\BIC{$\varphi_{4}$}
\AXC{$\varphi_{27}$\lbl{A2}}
\BIC{$\forall y\,(c \in \mathbb{R} \to x \in \mathbb{R} \to y \in \mathbb{R} \to c - x < y \to c < y + x)$}
\AXC{$\varphi_{25}$\lbl{A2}}
\BIC{$c \in \mathbb{R} \to x \in \mathbb{R} \to d \in \mathbb{R} \to c - x < d \to c < d + x$}
\AXC{$\varphi_{32}$\lbl{A1}}
\BIC{$\varphi_{11} \to \varphi_{12} \to \varphi_{13} \to \varphi_{16} \to \varphi_{17}$}
\BIC{$\varphi_{12} \to \varphi_{13} \to \varphi_{16} \to \varphi_{17}$}
\BIC{$\varphi_{13} \to \varphi_{16} \to \varphi_{17}$}
\BIC{$\varphi_{16} \to \varphi_{17}$}
\end{ptree}
\noindent$\Pi_{12}$:
\begin{ptree}
\AXC{$\varphi_{12}$\lbl{A1}}
\AXC{$\varphi_{13}$\lbl{A1}}
\AXC{$\varphi_{11}$\lbl{A1}}
\AXC{$\forall c\,(\varphi_{5})$\lbl{[H4]}}
\AXC{$\forall c\,(\varphi_{5}) \to \varphi_{5}$\lbl{A2}}
\BIC{$\varphi_{5}$}
\AXC{$\varphi_{28}$\lbl{A2}}
\BIC{$\forall x\,(c \in \mathbb{R} \to d \in \mathbb{R} \to x \in \mathbb{R} \to c < d + x \to c - d < x)$}
\AXC{$\varphi_{26}$\lbl{A2}}
\BIC{$c \in \mathbb{R} \to d \in \mathbb{R} \to x \in \mathbb{R} \to c < d + x \to c - d < x$}
\AXC{$\varphi_{33}$\lbl{A1}}
\BIC{$\varphi_{11} \to \varphi_{13} \to \varphi_{12} \to \varphi_{17} \to \varphi_{18}$}
\BIC{$\varphi_{13} \to \varphi_{12} \to \varphi_{17} \to \varphi_{18}$}
\BIC{$\varphi_{12} \to \varphi_{17} \to \varphi_{18}$}
\BIC{$\varphi_{17} \to \varphi_{18}$}
\end{ptree}
\noindent 主木:
\begin{ptree}
\AXC{$\stackrel{\Pi_{10}}{\vdots}$}\noLine
\UIC{$\varphi_{16}$}
\AXC{$\stackrel{\Pi_{11}}{\vdots}$}\noLine
\UIC{$\varphi_{16} \to \varphi_{17}$}
\BIC{$\varphi_{17}$}
\AXC{$\stackrel{\Pi_{12}}{\vdots}$}\noLine
\UIC{$\varphi_{17} \to \varphi_{18}$}
\BIC{$\varphi_{18}$}
\UIC{$\forall d\,(\varphi_{18})$}
\UIC{$\forall c\,\forall d\,(\varphi_{18})$}
\UIC{$\forall x\,\forall c\,\forall d\,(\varphi_{18})$}
\end{ptree}
\begin{abbrevlist}
\item $\varphi_{1} :\equiv \forall z\,(\lvert x - c \rvert = \lvert c - x \rvert \to \lvert c - x \rvert = z \to \lvert x - c \rvert = z)$
\item $\varphi_{2} :\equiv \lvert x - c \rvert = \lvert c - x \rvert \to \lvert c - x \rvert = c - x \to \lvert x - c \rvert = c - x$
\item $\varphi_{3} :\equiv \forall x\,\forall y\,\forall z\,(x = y \to y = z \to x = z) \to \forall y\,\forall z\,(\lvert x - c \rvert = y \to y = z \to \lvert x - c \rvert = z)$
\item $\varphi_{4} :\equiv \forall a\,\forall y\,(c \in \mathbb{R} \to a \in \mathbb{R} \to y \in \mathbb{R} \to c - a < y \to c < y + a)$
\item $\varphi_{5} :\equiv \forall d\,\forall x\,(c \in \mathbb{R} \to d \in \mathbb{R} \to x \in \mathbb{R} \to c < d + x \to c - d < x)$
\item $\varphi_{6} :\equiv c \in \mathbb{R} \to d \in \mathbb{R} \to x < c \to \lvert x - c \rvert < d \to \lvert c - x \rvert = c - x$
\item $\varphi_{7} :\equiv c \in \mathbb{R} \to d \in \mathbb{R} \to x < c \to \lvert x - c \rvert < d \to \lvert x - c \rvert = c - x$
\item $\varphi_{8} :\equiv \forall x\,\forall a\,\forall y\,(x \in \mathbb{R} \to a \in \mathbb{R} \to y \in \mathbb{R} \to x - a < y \to x < y + a)$
\item $\varphi_{9} :\equiv \forall z\,\forall w\,(z = w \to z < d \to w < d) \to \forall w\,(\lvert x - c \rvert = w \to \lvert x - c \rvert < d \to w < d)$
\item $\varphi_{10} :\equiv x \in \mathbb{R} \to c \in \mathbb{R} \to d \in \mathbb{R} \to x < c \to \lvert x - c \rvert < d \to 0 \leq c - x$
\item $\varphi_{11} :\equiv x \in \mathbb{R} \to c \in \mathbb{R} \to d \in \mathbb{R} \to x < c \to \lvert x - c \rvert < d \to c \in \mathbb{R}$
\item $\varphi_{12} :\equiv x \in \mathbb{R} \to c \in \mathbb{R} \to d \in \mathbb{R} \to x < c \to \lvert x - c \rvert < d \to x \in \mathbb{R}$
\item $\varphi_{13} :\equiv x \in \mathbb{R} \to c \in \mathbb{R} \to d \in \mathbb{R} \to x < c \to \lvert x - c \rvert < d \to d \in \mathbb{R}$
\item $\varphi_{14} :\equiv c \in \mathbb{R} \to d \in \mathbb{R} \to x < c \to \lvert x - c \rvert < d \to \lvert x - c \rvert = \lvert c - x \rvert$
\item $\varphi_{15} :\equiv x \in \mathbb{R} \to c \in \mathbb{R} \to d \in \mathbb{R} \to x < c \to \lvert x - c \rvert < d \to 0 < c - x$
\item $\varphi_{16} :\equiv x \in \mathbb{R} \to c \in \mathbb{R} \to d \in \mathbb{R} \to x < c \to \lvert x - c \rvert < d \to c - x < d$
\item $\varphi_{17} :\equiv x \in \mathbb{R} \to c \in \mathbb{R} \to d \in \mathbb{R} \to x < c \to \lvert x - c \rvert < d \to c < d + x$
\item $\varphi_{18} :\equiv x \in \mathbb{R} \to c \in \mathbb{R} \to d \in \mathbb{R} \to x < c \to \lvert x - c \rvert < d \to c - d < x$
\item $\varphi_{19} :\equiv x \in \mathbb{R} \to c \in \mathbb{R} \to d \in \mathbb{R} \to x < c \to \lvert x - c \rvert < d \to \lvert x - c \rvert < d$
\item $\varphi_{20} :\equiv \forall w\,(\lvert x - c \rvert = w \to \lvert x - c \rvert < d \to w < d) \to \lvert x - c \rvert = c - x \to \lvert x - c \rvert < d \to c - x < d$
\item $\varphi_{21} :\equiv \forall y\,(x \in \mathbb{R} \to y \in \mathbb{R} \to x < y \to 0 < y - x) \to x \in \mathbb{R} \to c \in \mathbb{R} \to x < c \to 0 < c - x$
\item $\varphi_{22} :\equiv \forall x\,\forall y\,(x \in \mathbb{R} \to y \in \mathbb{R} \to x < y \to 0 < y - x) \to \forall y\,(x \in \mathbb{R} \to y \in \mathbb{R} \to x < y \to 0 < y - x)$
\item $\varphi_{23} :\equiv \forall y\,(x \in \mathbb{R} \to y \in \mathbb{R} \to \lvert x - y \rvert = \lvert y - x \rvert) \to x \in \mathbb{R} \to c \in \mathbb{R} \to \lvert x - c \rvert = \lvert c - x \rvert$
\item $\varphi_{24} :\equiv \forall x\,\forall y\,(x \in \mathbb{R} \to y \in \mathbb{R} \to \lvert x - y \rvert = \lvert y - x \rvert) \to \forall y\,(x \in \mathbb{R} \to y \in \mathbb{R} \to \lvert x - y \rvert = \lvert y - x \rvert)$
\item $\varphi_{25} :\equiv \forall y\,(c \in \mathbb{R} \to x \in \mathbb{R} \to y \in \mathbb{R} \to c - x < y \to c < y + x) \to c \in \mathbb{R} \to x \in \mathbb{R} \to d \in \mathbb{R} \to c - x < d \to c < d + x$
\item $\varphi_{26} :\equiv \forall x\,(c \in \mathbb{R} \to d \in \mathbb{R} \to x \in \mathbb{R} \to c < d + x \to c - d < x) \to c \in \mathbb{R} \to d \in \mathbb{R} \to x \in \mathbb{R} \to c < d + x \to c - d < x$
\item $\varphi_{27} :\equiv \varphi_{4} \to \forall y\,(c \in \mathbb{R} \to x \in \mathbb{R} \to y \in \mathbb{R} \to c - x < y \to c < y + x)$
\item $\varphi_{28} :\equiv \varphi_{5} \to \forall x\,(c \in \mathbb{R} \to d \in \mathbb{R} \to x \in \mathbb{R} \to c < d + x \to c - d < x)$
\item $\varphi_{29} :\equiv (x \in \mathbb{R} \to c \in \mathbb{R} \to d \in \mathbb{R} \to x < c \to \lvert x - c \rvert < d \to x < c) \to \varphi_{15}$
\item $\varphi_{30} :\equiv (x \in \mathbb{R} \to c \in \mathbb{R} \to \lvert x - c \rvert = \lvert c - x \rvert) \to \varphi_{12} \to \varphi_{11} \to x \in \mathbb{R} \to \varphi_{14}$
\item $\varphi_{31} :\equiv (x \in \mathbb{R} \to \varphi_{7}) \to \varphi_{19} \to (\lvert x - c \rvert = c - x \to \lvert x - c \rvert < d \to c - x < d) \to \varphi_{16}$
\item $\varphi_{32} :\equiv (c \in \mathbb{R} \to x \in \mathbb{R} \to d \in \mathbb{R} \to c - x < d \to c < d + x) \to \varphi_{11} \to \varphi_{12} \to \varphi_{13} \to \varphi_{16} \to \varphi_{17}$
\item $\varphi_{33} :\equiv (c \in \mathbb{R} \to d \in \mathbb{R} \to x \in \mathbb{R} \to c < d + x \to c - d < x) \to \varphi_{11} \to \varphi_{13} \to \varphi_{12} \to \varphi_{17} \to \varphi_{18}$
\end{abbrevlist}
\end{fullproof}

% 自動生成 (prover/build.js): FLimLR  ノード数 805
\begin{fullproof}{定理1.19}{$\forall f\,\forall c\,\forall \alpha\,(\mathrm{FLim}(f,c,\alpha) \leftrightarrow \mathrm{FLimR}(f,c,\alpha) \land \mathrm{FLimL}(f,c,\alpha))$}{805}\phantomsection\label{tree:FLimLR}
\noindent$\Pi_{1}$:
\begin{ptree}
\AXC{$\forall x\,\forall y\,\forall z\,(x < y \to y \leq z \to x < z)$\lbl{[O3]}}
\AXC{$\varphi_{65}$\lbl{A2}}
\BIC{$\forall y\,\forall z\,(\lvert x - c \rvert < y \to y \leq z \to \lvert x - c \rvert < z)$}
\AXC{$\varphi_{97}$\lbl{A2}}
\BIC{$\forall z\,(\lvert x - c \rvert < \delta \to \delta \leq z \to \lvert x - c \rvert < z)$}
\AXC{$\varphi_{94}$\lbl{A2}}
\BIC{$\lvert x - c \rvert < \delta \to \delta \leq d_{1} \to \lvert x - c \rvert < d_{1}$}
\end{ptree}
\noindent$\Pi_{2}$:
\begin{ptree}
\AXC{$\varphi_{321}$\lbl{A1}}
\AXC{$\stackrel{\Pi_{1}}{\vdots}$}\noLine
\UIC{$\lvert x - c \rvert < \delta \to \delta \leq d_{1} \to \lvert x - c \rvert < d_{1}$}
\AXC{$(\lvert x - c \rvert < \delta \to \delta \leq d_{1} \to \lvert x - c \rvert < d_{1}) \to \varphi_{321} \to \varphi_{315} \to \varphi_{326}$\lbl{A1}}
\BIC{$\varphi_{321} \to \varphi_{315} \to \varphi_{326}$}
\BIC{$\varphi_{315} \to \varphi_{326}$}
\end{ptree}
\noindent$\Pi_{3}$:
\begin{ptree}
\AXC{$\forall f\,\forall c\,\forall \alpha\,(\varphi_{189})$\lbl{DefFLimR}}
\AXC{$\forall f\,\forall c\,\forall \alpha\,(\varphi_{189}) \to \forall c\,\forall \alpha\,(\varphi_{189})$\lbl{A2}}
\BIC{$\forall c\,\forall \alpha\,(\varphi_{189})$}
\AXC{$\forall c\,\forall \alpha\,(\varphi_{189}) \to \forall \alpha\,(\varphi_{189})$\lbl{A2}}
\BIC{$\forall \alpha\,(\varphi_{189})$}
\AXC{$\forall \alpha\,(\varphi_{189}) \to \varphi_{189}$\lbl{A2}}
\BIC{$\varphi_{189}$}
\end{ptree}
\noindent$\Pi_{4}$:
\begin{ptree}
\AXC{$\stackrel{\Pi_{3}}{\vdots}$}\noLine
\UIC{$\varphi_{189}$}
\AXC{$\varphi_{17}$\lbl{A1}}
\AXC{$\varphi_{17} \to \varphi_{251}$\lbl{A1}}
\BIC{$\varphi_{251}$}
\BIC{$\mathrm{FLimR}(f,c,\alpha) \land \mathrm{FLimL}(f,c,\alpha) \to c \in \mathbb{R}$}
\end{ptree}
\noindent$\Pi_{5}$:
\begin{ptree}
\AXC{$\stackrel{\Pi_{3}}{\vdots}$}\noLine
\UIC{$\varphi_{189}$}
\AXC{$\varphi_{17}$\lbl{A1}}
\AXC{$\varphi_{17} \to \varphi_{189} \to \varphi_{1}$\lbl{A1}}
\BIC{$\varphi_{189} \to \varphi_{1}$}
\BIC{$\varphi_{1}$}
\end{ptree}
\noindent$\Pi_{6}$:
\begin{ptree}
\AXC{$\forall f\,\forall x\,(\mathrm{RFn}(f) \to x \in \mathrm{dom}(f) \to x \in \mathbb{R})$\lbl{[S7]}}
\AXC{$\varphi_{66}$\lbl{A2}}
\BIC{$\forall x\,(\mathrm{RFn}(f) \to x \in \mathrm{dom}(f) \to x \in \mathbb{R})$}
\AXC{$\varphi_{49}$\lbl{A2}}
\BIC{$\mathrm{RFn}(f) \to x \in \mathrm{dom}(f) \to x \in \mathbb{R}$}
\end{ptree}
\noindent$\Pi_{7}$:
\begin{ptree}
\AXC{$\varphi_{296}$\lbl{A1}}
\AXC{$\stackrel{\Pi_{5}}{\vdots}$}\noLine
\UIC{$\varphi_{1}$}
\AXC{$\stackrel{\Pi_{6}}{\vdots}$}\noLine
\UIC{$\mathrm{RFn}(f) \to x \in \mathrm{dom}(f) \to x \in \mathbb{R}$}
\AXC{$(\mathrm{RFn}(f) \to x \in \mathrm{dom}(f) \to x \in \mathbb{R}) \to \varphi_{1} \to \varphi_{296} \to \varphi_{313}$\lbl{A1}}
\BIC{$\varphi_{1} \to \varphi_{296} \to \varphi_{313}$}
\BIC{$\varphi_{296} \to \varphi_{313}$}
\BIC{$\varphi_{313}$}
\end{ptree}
\noindent$\Pi_{8}$:
\begin{ptree}
\AXC{$\stackrel{\Pi_{7}}{\vdots}$}\noLine
\UIC{$\varphi_{313}$}
\AXC{$\forall x\,\forall c\,(\varphi_{72})$\lbl{[H7]}}
\AXC{$\forall x\,\forall c\,(\varphi_{72}) \to \forall c\,(\varphi_{72})$\lbl{A2}}
\BIC{$\forall c\,(\varphi_{72})$}
\AXC{$\forall c\,(\varphi_{72}) \to \varphi_{72}$\lbl{A2}}
\BIC{$\varphi_{72}$}
\AXC{$\varphi_{72} \to \varphi_{73}$\lbl{A2}}
\BIC{$\varphi_{73}$}
\AXC{$\varphi_{73} \to \varphi_{313} \to \varphi_{400}$\lbl{A1}}
\BIC{$\varphi_{313} \to \varphi_{400}$}
\BIC{$\varphi_{400}$}
\end{ptree}
\noindent$\Pi_{9}$:
\begin{ptree}
\AXC{$\stackrel{\Pi_{4}}{\vdots}$}\noLine
\UIC{$\mathrm{FLimR}(f,c,\alpha) \land \mathrm{FLimL}(f,c,\alpha) \to c \in \mathbb{R}$}
\AXC{$\stackrel{\Pi_{8}}{\vdots}$}\noLine
\UIC{$\varphi_{400}$}
\BIC{$\varphi_{263} \to \varphi_{317} \to \varphi_{326} \to \varphi_{324}$}
\end{ptree}
\noindent$\Pi_{10}$:
\begin{ptree}
\AXC{$\varphi_{317}$\lbl{A1}}
\AXC{$\mathrm{FLimR}(f,c,\alpha) \land \mathrm{FLimL}(f,c,\alpha) \to \varphi_{227}$\lbl{A1}}
\AXC{$\varphi_{358}$\lbl{A1}}
\BIC{$\varphi_{263}$}
\AXC{$\stackrel{\Pi_{9}}{\vdots}$}\noLine
\UIC{$\varphi_{263} \to \varphi_{317} \to \varphi_{326} \to \varphi_{324}$}
\BIC{$\varphi_{317} \to \varphi_{326} \to \varphi_{324}$}
\BIC{$\varphi_{326} \to \varphi_{324}$}
\end{ptree}
\noindent$\Pi_{11}$:
\begin{ptree}
\AXC{$\varphi_{129}$\lbl{A2}}
\AXC{$\mathrm{FLimR}(f,c,\alpha) \land \mathrm{FLimL}(f,c,\alpha) \to \varphi_{227}$\lbl{A1}}
\AXC{$\varphi_{374}$\lbl{A1}}
\BIC{$\varphi_{346}$}
\AXC{$\varphi_{346} \to \varphi_{129} \to \varphi_{344}$\lbl{A1}}
\BIC{$\varphi_{129} \to \varphi_{344}$}
\BIC{$\varphi_{344}$}
\AXC{$\varphi_{344} \to \varphi_{296} \to \varphi_{317} \to \varphi_{324} \to \varphi_{334}$\lbl{A1}}
\BIC{$\varphi_{296} \to \varphi_{317} \to \varphi_{324} \to \varphi_{334}$}
\end{ptree}
\noindent$\Pi_{12}$:
\begin{ptree}
\AXC{$\forall x\,\forall y\,\forall z\,(x < y \to y \leq z \to x < z)$\lbl{[O3]}}
\AXC{$\varphi_{65}$\lbl{A2}}
\BIC{$\forall y\,\forall z\,(\lvert x - c \rvert < y \to y \leq z \to \lvert x - c \rvert < z)$}
\AXC{$\varphi_{97}$\lbl{A2}}
\BIC{$\forall z\,(\lvert x - c \rvert < \delta \to \delta \leq z \to \lvert x - c \rvert < z)$}
\AXC{$\varphi_{95}$\lbl{A2}}
\BIC{$\lvert x - c \rvert < \delta \to \delta \leq d_{2} \to \lvert x - c \rvert < d_{2}$}
\end{ptree}
\noindent$\Pi_{13}$:
\begin{ptree}
\AXC{$\varphi_{321}$\lbl{A1}}
\AXC{$\stackrel{\Pi_{12}}{\vdots}$}\noLine
\UIC{$\lvert x - c \rvert < \delta \to \delta \leq d_{2} \to \lvert x - c \rvert < d_{2}$}
\AXC{$(\lvert x - c \rvert < \delta \to \delta \leq d_{2} \to \lvert x - c \rvert < d_{2}) \to \varphi_{321} \to \varphi_{316} \to \varphi_{327}$\lbl{A1}}
\BIC{$\varphi_{321} \to \varphi_{316} \to \varphi_{327}$}
\BIC{$\varphi_{316} \to \varphi_{327}$}
\end{ptree}
\noindent$\Pi_{14}$:
\begin{ptree}
\AXC{$\stackrel{\Pi_{7}}{\vdots}$}\noLine
\UIC{$\varphi_{313}$}
\AXC{$\forall x\,\forall c\,(\varphi_{74})$\lbl{[H8]}}
\AXC{$\forall x\,\forall c\,(\varphi_{74}) \to \forall c\,(\varphi_{74})$\lbl{A2}}
\BIC{$\forall c\,(\varphi_{74})$}
\AXC{$\forall c\,(\varphi_{74}) \to \varphi_{74}$\lbl{A2}}
\BIC{$\varphi_{74}$}
\AXC{$\varphi_{74} \to \varphi_{75}$\lbl{A2}}
\BIC{$\varphi_{75}$}
\AXC{$\varphi_{75} \to \varphi_{313} \to \varphi_{401}$\lbl{A1}}
\BIC{$\varphi_{313} \to \varphi_{401}$}
\BIC{$\varphi_{401}$}
\end{ptree}
\noindent$\Pi_{15}$:
\begin{ptree}
\AXC{$\varphi_{318}$\lbl{A1}}
\AXC{$\varphi_{312}$\lbl{A1}}
\AXC{$\varphi_{312} \to \varphi_{264}$\lbl{A1}}
\BIC{$\varphi_{264}$}
\AXC{$\stackrel{\Pi_{4}}{\vdots}$}\noLine
\UIC{$\mathrm{FLimR}(f,c,\alpha) \land \mathrm{FLimL}(f,c,\alpha) \to c \in \mathbb{R}$}
\AXC{$\stackrel{\Pi_{14}}{\vdots}$}\noLine
\UIC{$\varphi_{401}$}
\BIC{$\varphi_{264} \to \varphi_{318} \to \varphi_{327} \to \varphi_{325}$}
\BIC{$\varphi_{318} \to \varphi_{327} \to \varphi_{325}$}
\BIC{$\varphi_{327} \to \varphi_{325}$}
\end{ptree}
\noindent$\Pi_{16}$:
\begin{ptree}
\AXC{$\varphi_{296}$\lbl{A1}}
\AXC{$\varphi_{130}$\lbl{A2}}
\AXC{$\varphi_{312}$\lbl{A1}}
\AXC{$\varphi_{312} \to \varphi_{347}$\lbl{A1}}
\BIC{$\varphi_{347}$}
\AXC{$\varphi_{347} \to \varphi_{130} \to \varphi_{345}$\lbl{A1}}
\BIC{$\varphi_{130} \to \varphi_{345}$}
\BIC{$\varphi_{345}$}
\AXC{$\varphi_{345} \to \varphi_{296} \to \varphi_{325} \to \varphi_{318} \to \varphi_{335}$\lbl{A1}}
\BIC{$\varphi_{296} \to \varphi_{325} \to \varphi_{318} \to \varphi_{335}$}
\BIC{$\varphi_{325} \to \varphi_{318} \to \varphi_{335}$}
\end{ptree}
\noindent$\Pi_{17}$:
\begin{ptree}
\AXC{$\stackrel{\Pi_{7}}{\vdots}$}\noLine
\UIC{$\varphi_{313}$}
\AXC{$\varphi_{22}$\lbl{[H6]}}
\AXC{$\varphi_{126}$\lbl{A2}}
\BIC{$\forall c\,(x \in \mathbb{R} \to c \in \mathbb{R} \to 0 < \lvert x - c \rvert \to x < c \lor c < x)$}
\AXC{$\varphi_{122}$\lbl{A2}}
\BIC{$x \in \mathbb{R} \to c \in \mathbb{R} \to 0 < \lvert x - c \rvert \to x < c \lor c < x$}
\AXC{$\varphi_{393}$\lbl{A1}}
\BIC{$\varphi_{313} \to \varphi_{377}$}
\BIC{$\varphi_{377}$}
\end{ptree}
\noindent$\Pi_{18}$:
\begin{ptree}
\AXC{$\varphi_{305}$\lbl{A1}}
\AXC{$\stackrel{\Pi_{4}}{\vdots}$}\noLine
\UIC{$\mathrm{FLimR}(f,c,\alpha) \land \mathrm{FLimL}(f,c,\alpha) \to c \in \mathbb{R}$}
\AXC{$\stackrel{\Pi_{17}}{\vdots}$}\noLine
\UIC{$\varphi_{377}$}
\BIC{$\varphi_{305} \to \varphi_{319}$}
\BIC{$\varphi_{319}$}
\AXC{$\varphi_{319} \to \varphi_{335} \to \varphi_{334} \to \varphi_{328}$\lbl{A1}}
\BIC{$\varphi_{335} \to \varphi_{334} \to \varphi_{328}$}
\end{ptree}
\noindent$\Pi_{19}$:
\begin{ptree}
\AXC{$\varphi_{318}$\lbl{A1}}
\AXC{$\varphi_{298}$\lbl{A1}}
\AXC{$\varphi_{298} \to \varphi_{316}$\lbl{A1}}
\BIC{$\varphi_{316}$}
\AXC{$\stackrel{\Pi_{13}}{\vdots}$}\noLine
\UIC{$\varphi_{316} \to \varphi_{327}$}
\BIC{$\varphi_{327}$}
\AXC{$\stackrel{\Pi_{15}}{\vdots}$}\noLine
\UIC{$\varphi_{327} \to \varphi_{325}$}
\BIC{$\varphi_{325}$}
\AXC{$\stackrel{\Pi_{16}}{\vdots}$}\noLine
\UIC{$\varphi_{325} \to \varphi_{318} \to \varphi_{335}$}
\BIC{$\varphi_{318} \to \varphi_{335}$}
\BIC{$\varphi_{335}$}
\AXC{$\stackrel{\Pi_{18}}{\vdots}$}\noLine
\UIC{$\varphi_{335} \to \varphi_{334} \to \varphi_{328}$}
\BIC{$\varphi_{334} \to \varphi_{328}$}
\end{ptree}
\noindent$\Pi_{20}$:
\begin{ptree}
\AXC{$\varphi_{298}$\lbl{A1}}
\AXC{$\varphi_{298} \to \varphi_{315}$\lbl{A1}}
\BIC{$\varphi_{315}$}
\AXC{$\stackrel{\Pi_{2}}{\vdots}$}\noLine
\UIC{$\varphi_{315} \to \varphi_{326}$}
\BIC{$\varphi_{326}$}
\AXC{$\stackrel{\Pi_{10}}{\vdots}$}\noLine
\UIC{$\varphi_{326} \to \varphi_{324}$}
\BIC{$\varphi_{324}$}
\AXC{$\varphi_{317}$\lbl{A1}}
\AXC{$\varphi_{296}$\lbl{A1}}
\AXC{$\stackrel{\Pi_{11}}{\vdots}$}\noLine
\UIC{$\varphi_{296} \to \varphi_{317} \to \varphi_{324} \to \varphi_{334}$}
\BIC{$\varphi_{317} \to \varphi_{324} \to \varphi_{334}$}
\BIC{$\varphi_{324} \to \varphi_{334}$}
\BIC{$\varphi_{334}$}
\AXC{$\stackrel{\Pi_{19}}{\vdots}$}\noLine
\UIC{$\varphi_{334} \to \varphi_{328}$}
\BIC{$\varphi_{328}$}
\end{ptree}
\noindent$\Pi_{21}$:
\begin{ptree}
\AXC{$\varphi_{274} \to \forall x\,(\varphi_{274})$\lbl{A4}}
\AXC{$\forall x\,(\varphi_{274} \to \varphi_{69}) \to \forall x\,(\varphi_{274}) \to \varphi_{78}$\lbl{A3}}
\AXC{$\stackrel{\Pi_{20}}{\vdots}$}\noLine
\UIC{$\varphi_{328}$}
\AXC{$\varphi_{328} \to \varphi_{274} \to \varphi_{69}$\lbl{A1}}
\BIC{$\varphi_{274} \to \varphi_{69}$}
\UIC{$\forall x\,(\varphi_{274} \to \varphi_{69})$}
\AXC{$\varphi_{406}$\lbl{A1}}
\BIC{$(\forall x\,(\varphi_{274} \to \varphi_{69}) \to \forall x\,(\varphi_{274}) \to \varphi_{78}) \to (\varphi_{274} \to \forall x\,(\varphi_{274})) \to \varphi_{331}$}
\BIC{$(\varphi_{274} \to \forall x\,(\varphi_{274})) \to \varphi_{331}$}
\BIC{$\varphi_{331}$}
\end{ptree}
\noindent$\Pi_{22}$:
\begin{ptree}
\AXC{$\varphi_{240}$\lbl{A2}}
\AXC{$\stackrel{\Pi_{21}}{\vdots}$}\noLine
\UIC{$\varphi_{331}$}
\AXC{$\varphi_{298}$\lbl{A1}}
\AXC{$\varphi_{298} \to \varphi_{331} \to \varphi_{341}$\lbl{A1}}
\BIC{$\varphi_{331} \to \varphi_{341}$}
\BIC{$\varphi_{341}$}
\AXC{$\varphi_{341} \to \varphi_{240} \to \varphi_{342}$\lbl{A1}}
\BIC{$\varphi_{240} \to \varphi_{342}$}
\BIC{$\varphi_{342}$}
\end{ptree}
\noindent$\Pi_{23}$:
\begin{ptree}
\AXC{$\varphi_{79}$\lbl{[D8]}}
\AXC{$\varphi_{79} \to \varphi_{80}$\lbl{A2}}
\BIC{$\varphi_{80}$}
\AXC{$\varphi_{80} \to d_{1} \in \mathbb{R} \to d_{2} \in \mathbb{R} \to \varphi_{9}$\lbl{A2}}
\BIC{$d_{1} \in \mathbb{R} \to d_{2} \in \mathbb{R} \to \varphi_{9}$}
\AXC{$(d_{1} \in \mathbb{R} \to d_{2} \in \mathbb{R} \to \varphi_{9}) \to \varphi_{263} \to \varphi_{264} \to \varphi_{261} \to \varphi_{260} \to \varphi_{273}$\lbl{A1}}
\BIC{$\varphi_{263} \to \varphi_{264} \to \varphi_{261} \to \varphi_{260} \to \varphi_{273}$}
\end{ptree}
\noindent$\Pi_{24}$:
\begin{ptree}
\AXC{$\varphi_{312}$\lbl{A1}}
\AXC{$\varphi_{312} \to \varphi_{264}$\lbl{A1}}
\BIC{$\varphi_{264}$}
\AXC{$\mathrm{FLimR}(f,c,\alpha) \land \mathrm{FLimL}(f,c,\alpha) \to \varphi_{227}$\lbl{A1}}
\AXC{$\varphi_{358}$\lbl{A1}}
\BIC{$\varphi_{263}$}
\AXC{$\stackrel{\Pi_{23}}{\vdots}$}\noLine
\UIC{$\varphi_{263} \to \varphi_{264} \to \varphi_{261} \to \varphi_{260} \to \varphi_{273}$}
\BIC{$\varphi_{264} \to \varphi_{261} \to \varphi_{260} \to \varphi_{273}$}
\BIC{$\varphi_{261} \to \varphi_{260} \to \varphi_{273}$}
\end{ptree}
\noindent$\Pi_{25}$:
\begin{ptree}
\AXC{$\varphi_{312}$\lbl{A1}}
\AXC{$\varphi_{312} \to \varphi_{260}$\lbl{A1}}
\BIC{$\varphi_{260}$}
\AXC{$\mathrm{FLimR}(f,c,\alpha) \land \mathrm{FLimL}(f,c,\alpha) \to \varphi_{227}$\lbl{A1}}
\AXC{$\varphi_{357}$\lbl{A1}}
\BIC{$\varphi_{261}$}
\AXC{$\stackrel{\Pi_{24}}{\vdots}$}\noLine
\UIC{$\varphi_{261} \to \varphi_{260} \to \varphi_{273}$}
\BIC{$\varphi_{260} \to \varphi_{273}$}
\BIC{$\varphi_{273}$}
\end{ptree}
\noindent$\Pi_{26}$:
\begin{ptree}
\AXC{$\stackrel{\Pi_{22}}{\vdots}$}\noLine
\UIC{$\varphi_{342}$}
\AXC{$\varphi_{342} \to \varphi_{323} \to \neg (\delta \in \mathbb{R} \land 0 < \delta \land \delta \leq d_{1} \land \delta \leq d_{2})$\lbl{A1}}
\BIC{$\varphi_{323} \to \neg (\delta \in \mathbb{R} \land 0 < \delta \land \delta \leq d_{1} \land \delta \leq d_{2})$}
\UIC{$\forall \delta\,(\varphi_{323} \to \neg (\delta \in \mathbb{R} \land 0 < \delta \land \delta \leq d_{1} \land \delta \leq d_{2}))$}
\AXC{$\stackrel{\Pi_{25}}{\vdots}$}\noLine
\UIC{$\varphi_{273}$}
\AXC{$\varphi_{409}$\lbl{A1}}
\BIC{$\forall \delta\,(\varphi_{323} \to \neg (\delta \in \mathbb{R} \land 0 < \delta \land \delta \leq d_{1} \land \delta \leq d_{2})) \to \varphi_{407}$}
\BIC{$\varphi_{407}$}
\end{ptree}
\noindent$\Pi_{27}$:
\begin{ptree}
\AXC{$\forall d\,(\varphi_{64}) \to \varphi_{85}$\lbl{A3}}
\AXC{$\varphi_{64}$\lbl{A2}}
\UIC{$\forall d\,(\varphi_{64})$}
\AXC{$\varphi_{323} \to \forall \delta\,(\varphi_{323})$\lbl{A4}}
\AXC{$\forall \delta\,(\varphi_{323} \to \neg (\delta \in \mathbb{R} \land 0 < \delta \land \delta \leq d_{1} \land \delta \leq d_{2})) \to \varphi_{343}$\lbl{A3}}
\AXC{$\stackrel{\Pi_{26}}{\vdots}$}\noLine
\UIC{$\varphi_{407}$}
\BIC{$(\varphi_{323} \to \forall \delta\,(\varphi_{323})) \to \forall d\,(\varphi_{64}) \to (\forall d\,(\varphi_{64}) \to \varphi_{85}) \to \varphi_{90} \to \varphi_{322}$}
\BIC{$\forall d\,(\varphi_{64}) \to (\forall d\,(\varphi_{64}) \to \varphi_{85}) \to \varphi_{90} \to \varphi_{322}$}
\BIC{$(\forall d\,(\varphi_{64}) \to \varphi_{85}) \to \varphi_{90} \to \varphi_{322}$}
\BIC{$\varphi_{90} \to \varphi_{322}$}
\end{ptree}
\noindent$\Pi_{28}$:
\begin{ptree}
\AXC{$\forall f\,\forall c\,\forall \alpha\,(\varphi_{190})$\lbl{DefFLimL}}
\AXC{$\forall f\,\forall c\,\forall \alpha\,(\varphi_{190}) \to \forall c\,\forall \alpha\,(\varphi_{190})$\lbl{A2}}
\BIC{$\forall c\,\forall \alpha\,(\varphi_{190})$}
\AXC{$\forall c\,\forall \alpha\,(\varphi_{190}) \to \forall \alpha\,(\varphi_{190})$\lbl{A2}}
\BIC{$\forall \alpha\,(\varphi_{190})$}
\AXC{$\forall \alpha\,(\varphi_{190}) \to \varphi_{190}$\lbl{A2}}
\BIC{$\varphi_{190}$}
\end{ptree}
\noindent$\Pi_{29}$:
\begin{ptree}
\AXC{$\varphi_{15}$\lbl{A1}}
\AXC{$\varphi_{258}$\lbl{A2}}
\AXC{$\stackrel{\Pi_{28}}{\vdots}$}\noLine
\UIC{$\varphi_{190}$}
\AXC{$\varphi_{17}$\lbl{A1}}
\AXC{$\varphi_{17} \to \varphi_{190} \to \varphi_{161}$\lbl{A1}}
\BIC{$\varphi_{190} \to \varphi_{161}$}
\BIC{$\varphi_{161}$}
\AXC{$\varphi_{161} \to \varphi_{258} \to \varphi_{156}$\lbl{A1}}
\BIC{$\varphi_{258} \to \varphi_{156}$}
\BIC{$\varphi_{156}$}
\AXC{$\varphi_{156} \to \varphi_{15} \to \varphi_{23} \to \varphi_{257}$\lbl{A1}}
\BIC{$\varphi_{15} \to \varphi_{23} \to \varphi_{257}$}
\BIC{$\varphi_{23} \to \varphi_{257}$}
\end{ptree}
\noindent$\Pi_{30}$:
\begin{ptree}
\AXC{$\varphi_{23}$\lbl{A1}}
\AXC{$\stackrel{\Pi_{29}}{\vdots}$}\noLine
\UIC{$\varphi_{23} \to \varphi_{257}$}
\BIC{$\varphi_{257}$}
\AXC{$\varphi_{257} \to \forall d_{2}\,(\varphi_{267} \to \neg (d_{2} \in \mathbb{R} \land 0 < d_{2} \land \varphi_{31})) \to \varphi_{408}$\lbl{A1}}
\BIC{$\forall d_{2}\,(\varphi_{267} \to \neg (d_{2} \in \mathbb{R} \land 0 < d_{2} \land \varphi_{31})) \to \varphi_{408}$}
\end{ptree}
\noindent$\Pi_{31}$:
\begin{ptree}
\AXC{$\varphi_{90}$\lbl{A4}}
\AXC{$\stackrel{\Pi_{27}}{\vdots}$}\noLine
\UIC{$\varphi_{90} \to \varphi_{322}$}
\BIC{$\varphi_{322}$}
\AXC{$\varphi_{322} \to \varphi_{267} \to \neg (d_{2} \in \mathbb{R} \land 0 < d_{2} \land \varphi_{31})$\lbl{A1}}
\BIC{$\varphi_{267} \to \neg (d_{2} \in \mathbb{R} \land 0 < d_{2} \land \varphi_{31})$}
\UIC{$\forall d_{2}\,(\varphi_{267} \to \neg (d_{2} \in \mathbb{R} \land 0 < d_{2} \land \varphi_{31}))$}
\AXC{$\stackrel{\Pi_{30}}{\vdots}$}\noLine
\UIC{$\forall d_{2}\,(\varphi_{267} \to \neg (d_{2} \in \mathbb{R} \land 0 < d_{2} \land \varphi_{31})) \to \varphi_{408}$}
\BIC{$\varphi_{408}$}
\end{ptree}
\noindent$\Pi_{32}$:
\begin{ptree}
\AXC{$\varphi_{200}$\lbl{A4}}
\AXC{$\forall \delta\,(\varphi_{192}) \to \varphi_{204}$\lbl{A3}}
\AXC{$\varphi_{192}$\lbl{A2}}
\UIC{$\forall \delta\,(\varphi_{192})$}
\AXC{$\varphi_{267} \to \forall d_{2}\,(\varphi_{267})$\lbl{A4}}
\AXC{$\varphi_{376}$\lbl{A3}}
\AXC{$\stackrel{\Pi_{31}}{\vdots}$}\noLine
\UIC{$\varphi_{408}$}
\BIC{$(\varphi_{267} \to \forall d_{2}\,(\varphi_{267})) \to \forall \delta\,(\varphi_{192}) \to (\forall \delta\,(\varphi_{192}) \to \varphi_{204}) \to \varphi_{200} \to \varphi_{266}$}
\BIC{$\forall \delta\,(\varphi_{192}) \to (\forall \delta\,(\varphi_{192}) \to \varphi_{204}) \to \varphi_{200} \to \varphi_{266}$}
\BIC{$(\forall \delta\,(\varphi_{192}) \to \varphi_{204}) \to \varphi_{200} \to \varphi_{266}$}
\BIC{$\varphi_{200} \to \varphi_{266}$}
\BIC{$\varphi_{266}$}
\end{ptree}
\noindent$\Pi_{33}$:
\begin{ptree}
\AXC{$\varphi_{259}$\lbl{A2}}
\AXC{$\stackrel{\Pi_{3}}{\vdots}$}\noLine
\UIC{$\varphi_{189}$}
\AXC{$\varphi_{17}$\lbl{A1}}
\AXC{$\varphi_{17} \to \varphi_{189} \to \varphi_{162}$\lbl{A1}}
\BIC{$\varphi_{189} \to \varphi_{162}$}
\BIC{$\varphi_{162}$}
\AXC{$\varphi_{162} \to \varphi_{259} \to \varphi_{157}$\lbl{A1}}
\BIC{$\varphi_{259} \to \varphi_{157}$}
\BIC{$\varphi_{157}$}
\end{ptree}
\noindent$\Pi_{34}$:
\begin{ptree}
\AXC{$\stackrel{\Pi_{32}}{\vdots}$}\noLine
\UIC{$\varphi_{266}$}
\AXC{$\varphi_{266} \to \varphi_{174} \to \neg (d_{1} \in \mathbb{R} \land 0 < d_{1} \land \varphi_{29})$\lbl{A1}}
\BIC{$\varphi_{174} \to \neg (d_{1} \in \mathbb{R} \land 0 < d_{1} \land \varphi_{29})$}
\UIC{$\forall d_{1}\,(\varphi_{174} \to \neg (d_{1} \in \mathbb{R} \land 0 < d_{1} \land \varphi_{29}))$}
\AXC{$\stackrel{\Pi_{33}}{\vdots}$}\noLine
\UIC{$\varphi_{157}$}
\AXC{$\varphi_{157} \to \forall d_{1}\,(\varphi_{174} \to \neg (d_{1} \in \mathbb{R} \land 0 < d_{1} \land \varphi_{29})) \to \varphi_{405}$\lbl{A1}}
\BIC{$\forall d_{1}\,(\varphi_{174} \to \neg (d_{1} \in \mathbb{R} \land 0 < d_{1} \land \varphi_{29})) \to \varphi_{405}$}
\BIC{$\varphi_{405}$}
\end{ptree}
\noindent$\Pi_{35}$:
\begin{ptree}
\AXC{$\varphi_{199}$\lbl{A4}}
\AXC{$\forall \delta\,(\varphi_{191}) \to \varphi_{203}$\lbl{A3}}
\AXC{$\varphi_{191}$\lbl{A2}}
\UIC{$\forall \delta\,(\varphi_{191})$}
\AXC{$\varphi_{174} \to \forall d_{1}\,(\varphi_{174})$\lbl{A4}}
\AXC{$\varphi_{362}$\lbl{A3}}
\AXC{$\stackrel{\Pi_{34}}{\vdots}$}\noLine
\UIC{$\varphi_{405}$}
\BIC{$(\varphi_{174} \to \forall d_{1}\,(\varphi_{174})) \to \forall \delta\,(\varphi_{191}) \to (\forall \delta\,(\varphi_{191}) \to \varphi_{203}) \to \varphi_{199} \to \varphi_{173}$}
\BIC{$\forall \delta\,(\varphi_{191}) \to (\forall \delta\,(\varphi_{191}) \to \varphi_{203}) \to \varphi_{199} \to \varphi_{173}$}
\BIC{$(\forall \delta\,(\varphi_{191}) \to \varphi_{203}) \to \varphi_{199} \to \varphi_{173}$}
\BIC{$\varphi_{199} \to \varphi_{173}$}
\BIC{$\varphi_{173}$}
\UIC{$\forall \varepsilon\,(\varphi_{173})$}
\end{ptree}
\noindent$\Pi_{36}$:
\begin{ptree}
\AXC{$\varphi_{21}$\lbl{A4}}
\AXC{$\forall \varepsilon\,(\varphi_{173}) \to \varphi_{182}$\lbl{A3}}
\AXC{$\stackrel{\Pi_{35}}{\vdots}$}\noLine
\UIC{$\forall \varepsilon\,(\varphi_{173})$}
\AXC{$\forall \varepsilon\,(\varphi_{173}) \to (\forall \varepsilon\,(\varphi_{173}) \to \varphi_{182}) \to \varphi_{21} \to \varphi_{177}$\lbl{A1}}
\BIC{$(\forall \varepsilon\,(\varphi_{173}) \to \varphi_{182}) \to \varphi_{21} \to \varphi_{177}$}
\BIC{$\varphi_{21} \to \varphi_{177}$}
\BIC{$\varphi_{177}$}
\end{ptree}
\noindent$\Pi_{37}$:
\begin{ptree}
\AXC{$\forall x\,\forall y\,\forall z\,(x < y \to y \leq z \to x < z)$\lbl{[O3]}}
\AXC{$\varphi_{65}$\lbl{A2}}
\BIC{$\forall y\,\forall z\,(\lvert x - c \rvert < y \to y \leq z \to \lvert x - c \rvert < z)$}
\AXC{$\varphi_{96}$\lbl{A2}}
\BIC{$\forall z\,(\lvert x - c \rvert < d \to d \leq z \to \lvert x - c \rvert < z)$}
\AXC{$\varphi_{91}$\lbl{A2}}
\BIC{$\lvert x - c \rvert < d \to d \leq r_{1} \to \lvert x - c \rvert < r_{1}$}
\end{ptree}
\noindent$\Pi_{38}$:
\begin{ptree}
\AXC{$\varphi_{284}$\lbl{A1}}
\AXC{$\stackrel{\Pi_{37}}{\vdots}$}\noLine
\UIC{$\lvert x - c \rvert < d \to d \leq r_{1} \to \lvert x - c \rvert < r_{1}$}
\AXC{$(\lvert x - c \rvert < d \to d \leq r_{1} \to \lvert x - c \rvert < r_{1}) \to \varphi_{284} \to \varphi_{277} \to \varphi_{292}$\lbl{A1}}
\BIC{$\varphi_{284} \to \varphi_{277} \to \varphi_{292}$}
\BIC{$\varphi_{277} \to \varphi_{292}$}
\end{ptree}
\noindent$\Pi_{39}$:
\begin{ptree}
\AXC{$\forall x\,\forall c\,(\varphi_{72})$\lbl{[H7]}}
\AXC{$\forall x\,\forall c\,(\varphi_{72}) \to \forall c\,(\varphi_{72})$\lbl{A2}}
\BIC{$\forall c\,(\varphi_{72})$}
\AXC{$\forall c\,(\varphi_{72}) \to \varphi_{72}$\lbl{A2}}
\BIC{$\varphi_{72}$}
\AXC{$\varphi_{72} \to \varphi_{76}$\lbl{A2}}
\BIC{$\varphi_{76}$}
\AXC{$\varphi_{76} \to \varphi_{398}$\lbl{A1}}
\BIC{$\varphi_{398}$}
\end{ptree}
\noindent$\Pi_{40}$:
\begin{ptree}
\AXC{$\stackrel{\Pi_{4}}{\vdots}$}\noLine
\UIC{$\mathrm{FLimR}(f,c,\alpha) \land \mathrm{FLimL}(f,c,\alpha) \to c \in \mathbb{R}$}
\AXC{$\mathrm{FLimR}(f,c,\alpha) \land \mathrm{FLimL}(f,c,\alpha) \to \varphi_{235}$\lbl{A1}}
\AXC{$\stackrel{\Pi_{39}}{\vdots}$}\noLine
\UIC{$\varphi_{398}$}
\BIC{$\varphi_{391}$}
\BIC{$\varphi_{201} \to \varphi_{280} \to \varphi_{292} \to \varphi_{285}$}
\end{ptree}
\noindent$\Pi_{41}$:
\begin{ptree}
\AXC{$\varphi_{280}$\lbl{A1}}
\AXC{$\mathrm{FLimR}(f,c,\alpha) \land \mathrm{FLimL}(f,c,\alpha) \to \varphi_{56} \to \varphi_{56}$\lbl{A1}}
\AXC{$\varphi_{340}$\lbl{A1}}
\BIC{$\varphi_{201}$}
\AXC{$\stackrel{\Pi_{40}}{\vdots}$}\noLine
\UIC{$\varphi_{201} \to \varphi_{280} \to \varphi_{292} \to \varphi_{285}$}
\BIC{$\varphi_{280} \to \varphi_{292} \to \varphi_{285}$}
\BIC{$\varphi_{292} \to \varphi_{285}$}
\end{ptree}
\noindent$\Pi_{42}$:
\begin{ptree}
\AXC{$\varphi_{107}$\lbl{A2}}
\AXC{$\mathrm{FLimR}(f,c,\alpha) \land \mathrm{FLimL}(f,c,\alpha) \to \varphi_{56} \to \varphi_{56}$\lbl{A1}}
\AXC{$\varphi_{361}$\lbl{A1}}
\BIC{$\varphi_{306}$}
\AXC{$\varphi_{306} \to \varphi_{107} \to \varphi_{301}$\lbl{A1}}
\BIC{$\varphi_{107} \to \varphi_{301}$}
\BIC{$\varphi_{301}$}
\AXC{$\varphi_{301} \to \varphi_{389}$\lbl{A1}}
\BIC{$\varphi_{389}$}
\end{ptree}
\noindent$\Pi_{43}$:
\begin{ptree}
\AXC{$\varphi_{280}$\lbl{A1}}
\AXC{$\mathrm{FLimR}(f,c,\alpha) \land \mathrm{FLimL}(f,c,\alpha) \to \varphi_{235}$\lbl{A1}}
\AXC{$\stackrel{\Pi_{42}}{\vdots}$}\noLine
\UIC{$\varphi_{389}$}
\BIC{$\varphi_{280} \to \varphi_{285} \to \varphi_{286}$}
\BIC{$\varphi_{285} \to \varphi_{286}$}
\end{ptree}
\noindent$\Pi_{44}$:
\begin{ptree}
\AXC{$\varphi_{265}$\lbl{A1}}
\AXC{$\varphi_{265} \to \varphi_{277}$\lbl{A1}}
\BIC{$\varphi_{277}$}
\AXC{$\stackrel{\Pi_{38}}{\vdots}$}\noLine
\UIC{$\varphi_{277} \to \varphi_{292}$}
\BIC{$\varphi_{292}$}
\AXC{$\stackrel{\Pi_{41}}{\vdots}$}\noLine
\UIC{$\varphi_{292} \to \varphi_{285}$}
\BIC{$\varphi_{285}$}
\AXC{$\stackrel{\Pi_{43}}{\vdots}$}\noLine
\UIC{$\varphi_{285} \to \varphi_{286}$}
\BIC{$\varphi_{286}$}
\end{ptree}
\noindent$\Pi_{45}$:
\begin{ptree}
\AXC{$\forall x\,\forall y\,\forall z\,(x < y \to y \leq z \to x < z)$\lbl{[O3]}}
\AXC{$\varphi_{65}$\lbl{A2}}
\BIC{$\forall y\,\forall z\,(\lvert x - c \rvert < y \to y \leq z \to \lvert x - c \rvert < z)$}
\AXC{$\varphi_{96}$\lbl{A2}}
\BIC{$\forall z\,(\lvert x - c \rvert < d \to d \leq z \to \lvert x - c \rvert < z)$}
\AXC{$\varphi_{92}$\lbl{A2}}
\BIC{$\lvert x - c \rvert < d \to d \leq r_{2} \to \lvert x - c \rvert < r_{2}$}
\end{ptree}
\noindent$\Pi_{46}$:
\begin{ptree}
\AXC{$\varphi_{284}$\lbl{A1}}
\AXC{$\stackrel{\Pi_{45}}{\vdots}$}\noLine
\UIC{$\lvert x - c \rvert < d \to d \leq r_{2} \to \lvert x - c \rvert < r_{2}$}
\AXC{$(\lvert x - c \rvert < d \to d \leq r_{2} \to \lvert x - c \rvert < r_{2}) \to \varphi_{284} \to \varphi_{278} \to \varphi_{293}$\lbl{A1}}
\BIC{$\varphi_{284} \to \varphi_{278} \to \varphi_{293}$}
\BIC{$\varphi_{278} \to \varphi_{293}$}
\end{ptree}
\noindent$\Pi_{47}$:
\begin{ptree}
\AXC{$\forall x\,\forall c\,(\varphi_{74})$\lbl{[H8]}}
\AXC{$\forall x\,\forall c\,(\varphi_{74}) \to \forall c\,(\varphi_{74})$\lbl{A2}}
\BIC{$\forall c\,(\varphi_{74})$}
\AXC{$\forall c\,(\varphi_{74}) \to \varphi_{74}$\lbl{A2}}
\BIC{$\varphi_{74}$}
\AXC{$\varphi_{74} \to \varphi_{77}$\lbl{A2}}
\BIC{$\varphi_{77}$}
\AXC{$\varphi_{77} \to \varphi_{399}$\lbl{A1}}
\BIC{$\varphi_{399}$}
\end{ptree}
\noindent$\Pi_{48}$:
\begin{ptree}
\AXC{$\stackrel{\Pi_{4}}{\vdots}$}\noLine
\UIC{$\mathrm{FLimR}(f,c,\alpha) \land \mathrm{FLimL}(f,c,\alpha) \to c \in \mathbb{R}$}
\AXC{$\mathrm{FLimR}(f,c,\alpha) \land \mathrm{FLimL}(f,c,\alpha) \to \varphi_{235}$\lbl{A1}}
\AXC{$\stackrel{\Pi_{47}}{\vdots}$}\noLine
\UIC{$\varphi_{399}$}
\BIC{$\varphi_{392}$}
\BIC{$\varphi_{202} \to \varphi_{281} \to \varphi_{293} \to \varphi_{287}$}
\end{ptree}
\noindent$\Pi_{49}$:
\begin{ptree}
\AXC{$\varphi_{265}$\lbl{A1}}
\AXC{$\varphi_{265} \to \varphi_{278}$\lbl{A1}}
\BIC{$\varphi_{278}$}
\AXC{$\stackrel{\Pi_{46}}{\vdots}$}\noLine
\UIC{$\varphi_{278} \to \varphi_{293}$}
\BIC{$\varphi_{293}$}
\AXC{$\varphi_{281}$\lbl{A1}}
\AXC{$\varphi_{268}$\lbl{A1}}
\AXC{$\varphi_{268} \to \varphi_{202}$\lbl{A1}}
\BIC{$\varphi_{202}$}
\AXC{$\stackrel{\Pi_{48}}{\vdots}$}\noLine
\UIC{$\varphi_{202} \to \varphi_{281} \to \varphi_{293} \to \varphi_{287}$}
\BIC{$\varphi_{281} \to \varphi_{293} \to \varphi_{287}$}
\BIC{$\varphi_{293} \to \varphi_{287}$}
\BIC{$\varphi_{287}$}
\end{ptree}
\noindent$\Pi_{50}$:
\begin{ptree}
\AXC{$\varphi_{108}$\lbl{A2}}
\AXC{$\varphi_{268}$\lbl{A1}}
\AXC{$\varphi_{268} \to \varphi_{307}$\lbl{A1}}
\BIC{$\varphi_{307}$}
\AXC{$\varphi_{307} \to \varphi_{108} \to \varphi_{302}$\lbl{A1}}
\BIC{$\varphi_{108} \to \varphi_{302}$}
\BIC{$\varphi_{302}$}
\AXC{$\varphi_{302} \to \varphi_{390}$\lbl{A1}}
\BIC{$\varphi_{390}$}
\end{ptree}
\noindent$\Pi_{51}$:
\begin{ptree}
\AXC{$\varphi_{281}$\lbl{A1}}
\AXC{$\stackrel{\Pi_{49}}{\vdots}$}\noLine
\UIC{$\varphi_{287}$}
\AXC{$\mathrm{FLimR}(f,c,\alpha) \land \mathrm{FLimL}(f,c,\alpha) \to \varphi_{235}$\lbl{A1}}
\AXC{$\stackrel{\Pi_{50}}{\vdots}$}\noLine
\UIC{$\varphi_{390}$}
\BIC{$\varphi_{287} \to \varphi_{281} \to \varphi_{288}$}
\BIC{$\varphi_{281} \to \varphi_{288}$}
\BIC{$\varphi_{288}$}
\end{ptree}
\noindent$\Pi_{52}$:
\begin{ptree}
\AXC{$\varphi_{22}$\lbl{[H6]}}
\AXC{$\varphi_{126}$\lbl{A2}}
\BIC{$\forall c\,(x \in \mathbb{R} \to c \in \mathbb{R} \to 0 < \lvert x - c \rvert \to x < c \lor c < x)$}
\AXC{$\varphi_{122}$\lbl{A2}}
\BIC{$x \in \mathbb{R} \to c \in \mathbb{R} \to 0 < \lvert x - c \rvert \to x < c \lor c < x$}
\AXC{$(x \in \mathbb{R} \to c \in \mathbb{R} \to 0 < \lvert x - c \rvert \to x < c \lor c < x) \to \varphi_{384}$\lbl{A1}}
\BIC{$\varphi_{384}$}
\end{ptree}
\noindent$\Pi_{53}$:
\begin{ptree}
\AXC{$\stackrel{\Pi_{4}}{\vdots}$}\noLine
\UIC{$\mathrm{FLimR}(f,c,\alpha) \land \mathrm{FLimL}(f,c,\alpha) \to c \in \mathbb{R}$}
\AXC{$\mathrm{FLimR}(f,c,\alpha) \land \mathrm{FLimL}(f,c,\alpha) \to \varphi_{235}$\lbl{A1}}
\AXC{$\stackrel{\Pi_{52}}{\vdots}$}\noLine
\UIC{$\varphi_{384}$}
\BIC{$\varphi_{369}$}
\BIC{$\varphi_{270} \to \varphi_{282}$}
\end{ptree}
\noindent$\Pi_{54}$:
\begin{ptree}
\AXC{$\stackrel{\Pi_{44}}{\vdots}$}\noLine
\UIC{$\varphi_{286}$}
\AXC{$\stackrel{\Pi_{51}}{\vdots}$}\noLine
\UIC{$\varphi_{288}$}
\AXC{$\varphi_{270}$\lbl{A1}}
\AXC{$\stackrel{\Pi_{53}}{\vdots}$}\noLine
\UIC{$\varphi_{270} \to \varphi_{282}$}
\BIC{$\varphi_{282}$}
\AXC{$\varphi_{282} \to \varphi_{288} \to \varphi_{286} \to \varphi_{275}$\lbl{A1}}
\BIC{$\varphi_{288} \to \varphi_{286} \to \varphi_{275}$}
\BIC{$\varphi_{286} \to \varphi_{275}$}
\BIC{$\varphi_{275}$}
\AXC{$\varphi_{275} \to \varphi_{276}$\lbl{A1}}
\BIC{$\varphi_{276}$}
\UIC{$\forall x\,(\varphi_{276})$}
\AXC{$\varphi_{404}$\lbl{A1}}
\BIC{$(\forall x\,(\varphi_{276}) \to \forall x\,(\varphi_{241}) \to \varphi_{25}) \to (\varphi_{241} \to \forall x\,(\varphi_{241})) \to \varphi_{283}$}
\end{ptree}
\noindent$\Pi_{55}$:
\begin{ptree}
\AXC{$\varphi_{241} \to \forall x\,(\varphi_{241})$\lbl{A4}}
\AXC{$\forall x\,(\varphi_{276}) \to \forall x\,(\varphi_{241}) \to \varphi_{25}$\lbl{A3}}
\AXC{$\stackrel{\Pi_{54}}{\vdots}$}\noLine
\UIC{$(\forall x\,(\varphi_{276}) \to \forall x\,(\varphi_{241}) \to \varphi_{25}) \to (\varphi_{241} \to \forall x\,(\varphi_{241})) \to \varphi_{283}$}
\BIC{$(\varphi_{241} \to \forall x\,(\varphi_{241})) \to \varphi_{283}$}
\BIC{$\varphi_{283}$}
\end{ptree}
\noindent$\Pi_{56}$:
\begin{ptree}
\AXC{$\varphi_{112}$\lbl{DefPNbhd}}
\AXC{$\varphi_{239}$\lbl{A2}}
\BIC{$\forall c\,(\mathrm{PNbhd}(f,c) \leftrightarrow \exists r\,(r \in \mathbb{R} \land 0 < r \land \varphi_{26}))$}
\AXC{$\varphi_{228}$\lbl{A2}}
\BIC{$\mathrm{PNbhd}(f,c) \leftrightarrow \exists r\,(r \in \mathbb{R} \land 0 < r \land \varphi_{26})$}
\AXC{$(\mathrm{PNbhd}(f,c) \leftrightarrow \exists r\,(r \in \mathbb{R} \land 0 < r \land \varphi_{26})) \to \varphi_{363}$\lbl{A1}}
\BIC{$\varphi_{363}$}
\end{ptree}
\noindent$\Pi_{57}$:
\begin{ptree}
\AXC{$\varphi_{183}$\lbl{A2}}
\AXC{$\stackrel{\Pi_{55}}{\vdots}$}\noLine
\UIC{$\varphi_{283}$}
\AXC{$\varphi_{265}$\lbl{A1}}
\AXC{$\varphi_{265} \to \varphi_{368}$\lbl{A1}}
\BIC{$\varphi_{368}$}
\BIC{$\mathrm{FLimR}(f,c,\alpha) \land \mathrm{FLimL}(f,c,\alpha) \to \varphi_{56} \to \varphi_{211}$}
\AXC{$\varphi_{379}$\lbl{A1}}
\BIC{$\varphi_{183} \to \varphi_{297}$}
\BIC{$\varphi_{297}$}
\AXC{$\stackrel{\Pi_{56}}{\vdots}$}\noLine
\UIC{$\varphi_{363}$}
\BIC{$\mathrm{FLimR}(f,c,\alpha) \land \mathrm{FLimL}(f,c,\alpha) \to \varphi_{56} \to \varphi_{121}$}
\AXC{$\varphi_{354}$\lbl{A1}}
\BIC{$\varphi_{207} \to \neg (d \in \mathbb{R} \land 0 < d \land d \leq r_{1} \land d \leq r_{2})$}
\UIC{$\forall d\,(\varphi_{207} \to \neg (d \in \mathbb{R} \land 0 < d \land d \leq r_{1} \land d \leq r_{2}))$}
\end{ptree}
\noindent$\Pi_{58}$:
\begin{ptree}
\AXC{$\varphi_{79}$\lbl{[D8]}}
\AXC{$\varphi_{79} \to \varphi_{81}$\lbl{A2}}
\BIC{$\varphi_{81}$}
\AXC{$\varphi_{81} \to r_{1} \in \mathbb{R} \to r_{2} \in \mathbb{R} \to \varphi_{10}$\lbl{A2}}
\BIC{$r_{1} \in \mathbb{R} \to r_{2} \in \mathbb{R} \to \varphi_{10}$}
\AXC{$(r_{1} \in \mathbb{R} \to r_{2} \in \mathbb{R} \to \varphi_{10}) \to \varphi_{201} \to \varphi_{202} \to \varphi_{197} \to \varphi_{196} \to \varphi_{244}$\lbl{A1}}
\BIC{$\varphi_{201} \to \varphi_{202} \to \varphi_{197} \to \varphi_{196} \to \varphi_{244}$}
\end{ptree}
\noindent$\Pi_{59}$:
\begin{ptree}
\AXC{$\varphi_{268}$\lbl{A1}}
\AXC{$\varphi_{268} \to \varphi_{202}$\lbl{A1}}
\BIC{$\varphi_{202}$}
\AXC{$\mathrm{FLimR}(f,c,\alpha) \land \mathrm{FLimL}(f,c,\alpha) \to \varphi_{56} \to \varphi_{56}$\lbl{A1}}
\AXC{$\varphi_{340}$\lbl{A1}}
\BIC{$\varphi_{201}$}
\AXC{$\stackrel{\Pi_{58}}{\vdots}$}\noLine
\UIC{$\varphi_{201} \to \varphi_{202} \to \varphi_{197} \to \varphi_{196} \to \varphi_{244}$}
\BIC{$\varphi_{202} \to \varphi_{197} \to \varphi_{196} \to \varphi_{244}$}
\BIC{$\varphi_{197} \to \varphi_{196} \to \varphi_{244}$}
\end{ptree}
\noindent$\Pi_{60}$:
\begin{ptree}
\AXC{$\varphi_{268}$\lbl{A1}}
\AXC{$\varphi_{268} \to \varphi_{196}$\lbl{A1}}
\BIC{$\varphi_{196}$}
\AXC{$\mathrm{FLimR}(f,c,\alpha) \land \mathrm{FLimL}(f,c,\alpha) \to \varphi_{56} \to \varphi_{56}$\lbl{A1}}
\AXC{$\varphi_{338}$\lbl{A1}}
\BIC{$\varphi_{197}$}
\AXC{$\stackrel{\Pi_{59}}{\vdots}$}\noLine
\UIC{$\varphi_{197} \to \varphi_{196} \to \varphi_{244}$}
\BIC{$\varphi_{196} \to \varphi_{244}$}
\BIC{$\varphi_{244}$}
\end{ptree}
\noindent$\Pi_{61}$:
\begin{ptree}
\AXC{$\stackrel{\Pi_{57}}{\vdots}$}\noLine
\UIC{$\forall d\,(\varphi_{207} \to \neg (d \in \mathbb{R} \land 0 < d \land d \leq r_{1} \land d \leq r_{2}))$}
\AXC{$\stackrel{\Pi_{60}}{\vdots}$}\noLine
\UIC{$\varphi_{244}$}
\AXC{$\varphi_{403}$\lbl{A1}}
\BIC{$\forall d\,(\varphi_{207} \to \neg (d \in \mathbb{R} \land 0 < d \land d \leq r_{1} \land d \leq r_{2})) \to \varphi_{388}$}
\BIC{$\varphi_{388}$}
\end{ptree}
\noindent$\Pi_{62}$:
\begin{ptree}
\AXC{$\varphi_{207} \to \forall d\,(\varphi_{207})$\lbl{A4}}
\AXC{$\forall d\,(\varphi_{207} \to \neg (d \in \mathbb{R} \land 0 < d \land d \leq r_{1} \land d \leq r_{2})) \to \varphi_{256}$\lbl{A3}}
\AXC{$\stackrel{\Pi_{61}}{\vdots}$}\noLine
\UIC{$\varphi_{388}$}
\BIC{$(\varphi_{207} \to \forall d\,(\varphi_{207})) \to \varphi_{198}$}
\BIC{$\varphi_{198}$}
\AXC{$\varphi_{198} \to \varphi_{118} \to \neg \varphi_{57}$\lbl{A1}}
\BIC{$\varphi_{118} \to \neg \varphi_{57}$}
\UIC{$\forall r_{2}\,(\varphi_{118} \to \neg \varphi_{57})$}
\end{ptree}
\noindent$\Pi_{63}$:
\begin{ptree}
\AXC{$\stackrel{\Pi_{28}}{\vdots}$}\noLine
\UIC{$\varphi_{190}$}
\AXC{$\varphi_{17}$\lbl{A1}}
\AXC{$\varphi_{17} \to \varphi_{190} \to (\mathrm{LNbhd}(f,c) \leftrightarrow \exists r\,(\varphi_{47})) \to \varphi_{110}$\lbl{A1}}
\BIC{$\varphi_{190} \to (\mathrm{LNbhd}(f,c) \leftrightarrow \exists r\,(\varphi_{47})) \to \varphi_{110}$}
\BIC{$(\mathrm{LNbhd}(f,c) \leftrightarrow \exists r\,(\varphi_{47})) \to \varphi_{110}$}
\end{ptree}
\noindent$\Pi_{64}$:
\begin{ptree}
\AXC{$\forall f\,\forall c\,(\mathrm{LNbhd}(f,c) \leftrightarrow \exists r\,(\varphi_{47}))$\lbl{DefLNbhd}}
\AXC{$\varphi_{187}$\lbl{A2}}
\BIC{$\forall c\,(\mathrm{LNbhd}(f,c) \leftrightarrow \exists r\,(\varphi_{47}))$}
\AXC{$\varphi_{184}$\lbl{A2}}
\BIC{$\mathrm{LNbhd}(f,c) \leftrightarrow \exists r\,(\varphi_{47})$}
\AXC{$\stackrel{\Pi_{63}}{\vdots}$}\noLine
\UIC{$(\mathrm{LNbhd}(f,c) \leftrightarrow \exists r\,(\varphi_{47})) \to \varphi_{110}$}
\BIC{$\varphi_{110}$}
\end{ptree}
\noindent$\Pi_{65}$:
\begin{ptree}
\AXC{$\stackrel{\Pi_{62}}{\vdots}$}\noLine
\UIC{$\forall r_{2}\,(\varphi_{118} \to \neg \varphi_{57})$}
\AXC{$\stackrel{\Pi_{64}}{\vdots}$}\noLine
\UIC{$\varphi_{110}$}
\AXC{$\varphi_{110} \to \varphi_{186}$\lbl{A1}}
\BIC{$\varphi_{186}$}
\AXC{$\varphi_{186} \to \forall r_{2}\,(\varphi_{118} \to \neg \varphi_{57}) \to \varphi_{402}$\lbl{A1}}
\BIC{$\forall r_{2}\,(\varphi_{118} \to \neg \varphi_{57}) \to \varphi_{402}$}
\BIC{$\varphi_{402}$}
\end{ptree}
\noindent$\Pi_{66}$:
\begin{ptree}
\AXC{$\varphi_{300}$\lbl{A3}}
\AXC{$\forall r_{2}\,(\neg \varphi_{57}) \to \neg \varphi_{47}$\lbl{A2}}
\UIC{$\forall r\,(\forall r_{2}\,(\neg \varphi_{57}) \to \neg \varphi_{47})$}
\AXC{$\varphi_{118} \to \forall r_{2}\,(\varphi_{118})$\lbl{A4}}
\AXC{$\forall r_{2}\,(\varphi_{118} \to \neg \varphi_{57}) \to \forall r_{2}\,(\varphi_{118}) \to \forall r_{2}\,(\neg \varphi_{57})$\lbl{A3}}
\AXC{$\stackrel{\Pi_{65}}{\vdots}$}\noLine
\UIC{$\varphi_{402}$}
\BIC{$(\varphi_{118} \to \forall r_{2}\,(\varphi_{118})) \to \varphi_{378}$}
\BIC{$\varphi_{378}$}
\BIC{$\varphi_{300} \to (\forall r_{2}\,(\neg \varphi_{57}) \to \forall r\,\forall r_{2}\,(\neg \varphi_{57})) \to \varphi_{117}$}
\BIC{$(\forall r_{2}\,(\neg \varphi_{57}) \to \forall r\,\forall r_{2}\,(\neg \varphi_{57})) \to \varphi_{117}$}
\end{ptree}
\noindent$\Pi_{67}$:
\begin{ptree}
\AXC{$\forall r_{2}\,(\neg \varphi_{57}) \to \forall r\,\forall r_{2}\,(\neg \varphi_{57})$\lbl{A4}}
\AXC{$\stackrel{\Pi_{66}}{\vdots}$}\noLine
\UIC{$(\forall r_{2}\,(\neg \varphi_{57}) \to \forall r\,\forall r_{2}\,(\neg \varphi_{57})) \to \varphi_{117}$}
\BIC{$\varphi_{117}$}
\AXC{$\varphi_{117} \to \varphi_{3} \to \neg \varphi_{56}$\lbl{A1}}
\BIC{$\varphi_{3} \to \neg \varphi_{56}$}
\UIC{$\forall r_{1}\,(\varphi_{3} \to \neg \varphi_{56})$}
\end{ptree}
\noindent$\Pi_{68}$:
\begin{ptree}
\AXC{$\stackrel{\Pi_{3}}{\vdots}$}\noLine
\UIC{$\varphi_{189}$}
\AXC{$\varphi_{17}$\lbl{A1}}
\AXC{$\varphi_{17} \to \varphi_{189} \to (\mathrm{RNbhd}(f,c) \leftrightarrow \exists r\,(\varphi_{48})) \to \varphi_{111}$\lbl{A1}}
\BIC{$\varphi_{189} \to (\mathrm{RNbhd}(f,c) \leftrightarrow \exists r\,(\varphi_{48})) \to \varphi_{111}$}
\BIC{$(\mathrm{RNbhd}(f,c) \leftrightarrow \exists r\,(\varphi_{48})) \to \varphi_{111}$}
\end{ptree}
\noindent$\Pi_{69}$:
\begin{ptree}
\AXC{$\forall f\,\forall c\,(\mathrm{RNbhd}(f,c) \leftrightarrow \exists r\,(\varphi_{48}))$\lbl{DefRNbhd}}
\AXC{$\varphi_{188}$\lbl{A2}}
\BIC{$\forall c\,(\mathrm{RNbhd}(f,c) \leftrightarrow \exists r\,(\varphi_{48}))$}
\AXC{$\varphi_{185}$\lbl{A2}}
\BIC{$\mathrm{RNbhd}(f,c) \leftrightarrow \exists r\,(\varphi_{48})$}
\AXC{$\stackrel{\Pi_{68}}{\vdots}$}\noLine
\UIC{$(\mathrm{RNbhd}(f,c) \leftrightarrow \exists r\,(\varphi_{48})) \to \varphi_{111}$}
\BIC{$\varphi_{111}$}
\end{ptree}
\noindent$\Pi_{70}$:
\begin{ptree}
\AXC{$\forall r_{1}\,(\varphi_{3} \to \neg \varphi_{56}) \to \forall r_{1}\,(\varphi_{3}) \to \forall r_{1}\,(\neg \varphi_{56})$\lbl{A3}}
\AXC{$\stackrel{\Pi_{67}}{\vdots}$}\noLine
\UIC{$\forall r_{1}\,(\varphi_{3} \to \neg \varphi_{56})$}
\AXC{$\stackrel{\Pi_{69}}{\vdots}$}\noLine
\UIC{$\varphi_{111}$}
\AXC{$\varphi_{111} \to \forall r_{1}\,(\varphi_{3} \to \neg \varphi_{56}) \to \varphi_{387}$\lbl{A1}}
\BIC{$\forall r_{1}\,(\varphi_{3} \to \neg \varphi_{56}) \to \varphi_{387}$}
\BIC{$\varphi_{387}$}
\BIC{$(\varphi_{3} \to \forall r_{1}\,(\varphi_{3})) \to \varphi_{375}$}
\end{ptree}
\noindent$\Pi_{71}$:
\begin{ptree}
\AXC{$\varphi_{299}$\lbl{A3}}
\AXC{$\forall r_{1}\,(\neg \varphi_{56}) \to \neg \varphi_{48}$\lbl{A2}}
\UIC{$\forall r\,(\forall r_{1}\,(\neg \varphi_{56}) \to \neg \varphi_{48})$}
\AXC{$\varphi_{3} \to \forall r_{1}\,(\varphi_{3})$\lbl{A4}}
\AXC{$\stackrel{\Pi_{70}}{\vdots}$}\noLine
\UIC{$(\varphi_{3} \to \forall r_{1}\,(\varphi_{3})) \to \varphi_{375}$}
\BIC{$\varphi_{375}$}
\BIC{$\varphi_{299} \to (\forall r_{1}\,(\neg \varphi_{56}) \to \forall r\,\forall r_{1}\,(\neg \varphi_{56})) \to \varphi_{2}$}
\BIC{$(\forall r_{1}\,(\neg \varphi_{56}) \to \forall r\,\forall r_{1}\,(\neg \varphi_{56})) \to \varphi_{2}$}
\end{ptree}
\noindent$\Pi_{72}$:
\begin{ptree}
\AXC{$\forall r_{1}\,(\neg \varphi_{56}) \to \forall r\,\forall r_{1}\,(\neg \varphi_{56})$\lbl{A4}}
\AXC{$\stackrel{\Pi_{71}}{\vdots}$}\noLine
\UIC{$(\forall r_{1}\,(\neg \varphi_{56}) \to \forall r\,\forall r_{1}\,(\neg \varphi_{56})) \to \varphi_{2}$}
\BIC{$\varphi_{2}$}
\end{ptree}
\noindent$\Pi_{73}$:
\begin{ptree}
\AXC{$\forall f\,\forall c\,\forall \alpha\,(\varphi_{217})$\lbl{DefFLim}}
\AXC{$\forall f\,\forall c\,\forall \alpha\,(\varphi_{217}) \to \forall c\,\forall \alpha\,(\varphi_{217})$\lbl{A2}}
\BIC{$\forall c\,\forall \alpha\,(\varphi_{217})$}
\AXC{$\forall c\,\forall \alpha\,(\varphi_{217}) \to \forall \alpha\,(\varphi_{217})$\lbl{A2}}
\BIC{$\forall \alpha\,(\varphi_{217})$}
\AXC{$\forall \alpha\,(\varphi_{217}) \to \varphi_{217}$\lbl{A2}}
\BIC{$\varphi_{217}$}
\end{ptree}
\noindent$\Pi_{74}$:
\begin{ptree}
\AXC{$\stackrel{\Pi_{4}}{\vdots}$}\noLine
\UIC{$\mathrm{FLimR}(f,c,\alpha) \land \mathrm{FLimL}(f,c,\alpha) \to c \in \mathbb{R}$}
\AXC{$\stackrel{\Pi_{5}}{\vdots}$}\noLine
\UIC{$\varphi_{1}$}
\AXC{$\stackrel{\Pi_{73}}{\vdots}$}\noLine
\UIC{$\varphi_{217}$}
\AXC{$\varphi_{217} \to \varphi_{1} \to \varphi_{320}$\lbl{A1}}
\BIC{$\varphi_{1} \to \varphi_{320}$}
\BIC{$\varphi_{320}$}
\BIC{$\varphi_{295}$}
\end{ptree}
\noindent$\Pi_{75}$:
\begin{ptree}
\AXC{$\forall f\,\forall c\,\forall \alpha\,(\mathrm{FLim}(f,c,\alpha) \to c \in \mathbb{R})$\lbl{[D3]}}
\AXC{$\varphi_{38}$\lbl{A2}}
\BIC{$\forall c\,\forall \alpha\,(\mathrm{FLim}(f,c,\alpha) \to c \in \mathbb{R})$}
\AXC{$\varphi_{19}$\lbl{A2}}
\BIC{$\forall \alpha\,(\mathrm{FLim}(f,c,\alpha) \to c \in \mathbb{R})$}
\AXC{$\varphi_{11}$\lbl{A2}}
\BIC{$\mathrm{FLim}(f,c,\alpha) \to c \in \mathbb{R}$}
\end{ptree}
\noindent$\Pi_{76}$:
\begin{ptree}
\AXC{$\forall f\,\forall c\,\forall \alpha\,(\mathrm{FLim}(f,c,\alpha) \to \mathrm{RFn}(f))$\lbl{[D2]}}
\AXC{$\varphi_{18}$\lbl{A2}}
\BIC{$\forall c\,\forall \alpha\,(\mathrm{FLim}(f,c,\alpha) \to \mathrm{RFn}(f))$}
\AXC{$\varphi_{8}$\lbl{A2}}
\BIC{$\forall \alpha\,(\mathrm{FLim}(f,c,\alpha) \to \mathrm{RFn}(f))$}
\AXC{$\varphi_{7}$\lbl{A2}}
\BIC{$\mathrm{FLim}(f,c,\alpha) \to \mathrm{RFn}(f)$}
\end{ptree}
\noindent$\Pi_{77}$:
\begin{ptree}
\AXC{$\forall f\,\forall x\,(\mathrm{RFn}(f) \to x \in \mathrm{dom}(f) \to x \in \mathbb{R})$\lbl{[S7]}}
\AXC{$\varphi_{66}$\lbl{A2}}
\BIC{$\forall x\,(\mathrm{RFn}(f) \to x \in \mathrm{dom}(f) \to x \in \mathbb{R})$}
\AXC{$\varphi_{49}$\lbl{A2}}
\BIC{$\mathrm{RFn}(f) \to x \in \mathrm{dom}(f) \to x \in \mathbb{R}$}
\AXC{$\varphi_{348}$\lbl{A1}}
\BIC{$(\mathrm{FLim}(f,c,\alpha) \to \mathrm{RFn}(f)) \to \varphi_{336}$}
\end{ptree}
\noindent$\Pi_{78}$:
\begin{ptree}
\AXC{$\mathrm{FLim}(f,c,\alpha) \to \varepsilon \in \mathbb{R} \to \varphi_{151}$\lbl{A1}}
\AXC{$\stackrel{\Pi_{76}}{\vdots}$}\noLine
\UIC{$\mathrm{FLim}(f,c,\alpha) \to \mathrm{RFn}(f)$}
\AXC{$\stackrel{\Pi_{77}}{\vdots}$}\noLine
\UIC{$(\mathrm{FLim}(f,c,\alpha) \to \mathrm{RFn}(f)) \to \varphi_{336}$}
\BIC{$\varphi_{336}$}
\BIC{$\mathrm{FLim}(f,c,\alpha) \to \varepsilon \in \mathbb{R} \to 0 < \varepsilon \to \varphi_{163}$}
\end{ptree}
\noindent$\Pi_{79}$:
\begin{ptree}
\AXC{$\forall x\,\forall c\,(\varphi_{99})$\lbl{[H5]}}
\AXC{$\forall x\,\forall c\,(\varphi_{99}) \to \forall c\,(\varphi_{99})$\lbl{A2}}
\BIC{$\forall c\,(\varphi_{99})$}
\AXC{$\forall c\,(\varphi_{99}) \to \varphi_{99}$\lbl{A2}}
\BIC{$\varphi_{99}$}
\AXC{$\varphi_{99} \to x \in \mathbb{R} \to \varphi_{86}$\lbl{A2}}
\BIC{$x \in \mathbb{R} \to \varphi_{86}$}
\end{ptree}
\noindent$\Pi_{80}$:
\begin{ptree}
\AXC{$\stackrel{\Pi_{75}}{\vdots}$}\noLine
\UIC{$\mathrm{FLim}(f,c,\alpha) \to c \in \mathbb{R}$}
\AXC{$\stackrel{\Pi_{78}}{\vdots}$}\noLine
\UIC{$\mathrm{FLim}(f,c,\alpha) \to \varepsilon \in \mathbb{R} \to 0 < \varepsilon \to \varphi_{163}$}
\AXC{$\stackrel{\Pi_{79}}{\vdots}$}\noLine
\UIC{$x \in \mathbb{R} \to \varphi_{86}$}
\AXC{$(x \in \mathbb{R} \to \varphi_{86}) \to \varphi_{386}$\lbl{A1}}
\BIC{$\varphi_{386}$}
\BIC{$(\mathrm{FLim}(f,c,\alpha) \to c \in \mathbb{R}) \to \varphi_{168} \to \varphi_{371}$}
\BIC{$\varphi_{168} \to \varphi_{371}$}
\end{ptree}
\noindent$\Pi_{81}$:
\begin{ptree}
\AXC{$\mathrm{FLim}(f,c,\alpha) \to \varepsilon \in \mathbb{R} \to 0 < \varepsilon \to \varphi_{155}$\lbl{A1}}
\AXC{$\varphi_{262}$\lbl{A1}}
\AXC{$\varphi_{262} \to \varphi_{168}$\lbl{A1}}
\BIC{$\varphi_{168}$}
\AXC{$\stackrel{\Pi_{80}}{\vdots}$}\noLine
\UIC{$\varphi_{168} \to \varphi_{371}$}
\BIC{$\varphi_{371}$}
\BIC{$(\mathrm{FLim}(f,c,\alpha) \to \varepsilon \in \mathbb{R} \to 0 < \varepsilon \to \varphi_{160}) \to \varphi_{245}$}
\end{ptree}
\noindent$\Pi_{82}$:
\begin{ptree}
\AXC{$\mathrm{FLim}(f,c,\alpha) \to \varepsilon \in \mathbb{R} \to 0 < \varepsilon \to \varphi_{160}$\lbl{A1}}
\AXC{$\stackrel{\Pi_{81}}{\vdots}$}\noLine
\UIC{$(\mathrm{FLim}(f,c,\alpha) \to \varepsilon \in \mathbb{R} \to 0 < \varepsilon \to \varphi_{160}) \to \varphi_{245}$}
\BIC{$\varphi_{245}$}
\end{ptree}
\noindent$\Pi_{83}$:
\begin{ptree}
\AXC{$\varphi_{78} \to \varphi_{69}$\lbl{A2}}
\AXC{$\varphi_{262}$\lbl{A1}}
\AXC{$\varphi_{262} \to \varphi_{250}$\lbl{A1}}
\BIC{$\varphi_{250}$}
\AXC{$\varphi_{250} \to (\varphi_{78} \to \varphi_{69}) \to \varphi_{247}$\lbl{A1}}
\BIC{$(\varphi_{78} \to \varphi_{69}) \to \varphi_{247}$}
\BIC{$\varphi_{247}$}
\AXC{$\varphi_{247} \to \varphi_{372}$\lbl{A1}}
\BIC{$\varphi_{372}$}
\end{ptree}
\noindent$\Pi_{84}$:
\begin{ptree}
\AXC{$\stackrel{\Pi_{82}}{\vdots}$}\noLine
\UIC{$\varphi_{245}$}
\AXC{$\mathrm{FLim}(f,c,\alpha) \to \varepsilon \in \mathbb{R} \to \varphi_{151}$\lbl{A1}}
\AXC{$\stackrel{\Pi_{83}}{\vdots}$}\noLine
\UIC{$\varphi_{372}$}
\BIC{$\varphi_{245} \to \mathrm{FLim}(f,c,\alpha) \to \varepsilon \in \mathbb{R} \to 0 < \varepsilon \to \varphi_{178}$}
\BIC{$\mathrm{FLim}(f,c,\alpha) \to \varepsilon \in \mathbb{R} \to 0 < \varepsilon \to \varphi_{178}$}
\AXC{$(\mathrm{FLim}(f,c,\alpha) \to \varepsilon \in \mathbb{R} \to 0 < \varepsilon \to \varphi_{178}) \to \varphi_{232}$\lbl{A1}}
\BIC{$\varphi_{232}$}
\UIC{$\forall x\,(\varphi_{232})$}
\end{ptree}
\noindent$\Pi_{85}$:
\begin{ptree}
\AXC{$\varphi_{150} \to \forall x\,(\varphi_{150})$\lbl{A4}}
\AXC{$\forall x\,(\varphi_{232}) \to \forall x\,(\varphi_{150}) \to \varphi_{32}$\lbl{A3}}
\AXC{$\stackrel{\Pi_{84}}{\vdots}$}\noLine
\UIC{$\forall x\,(\varphi_{232})$}
\AXC{$\varphi_{394}$\lbl{A1}}
\BIC{$(\forall x\,(\varphi_{232}) \to \forall x\,(\varphi_{150}) \to \varphi_{32}) \to (\varphi_{150} \to \forall x\,(\varphi_{150})) \to \varphi_{237}$}
\BIC{$(\varphi_{150} \to \forall x\,(\varphi_{150})) \to \varphi_{237}$}
\BIC{$\varphi_{237}$}
\end{ptree}
\noindent$\Pi_{86}$:
\begin{ptree}
\AXC{$\varphi_{193}$\lbl{A2}}
\AXC{$\stackrel{\Pi_{85}}{\vdots}$}\noLine
\UIC{$\varphi_{237}$}
\AXC{$\varphi_{262}$\lbl{A1}}
\AXC{$\varphi_{262} \to \varphi_{237} \to \varphi_{254}$\lbl{A1}}
\BIC{$\varphi_{237} \to \varphi_{254}$}
\BIC{$\varphi_{254}$}
\AXC{$\varphi_{254} \to \varphi_{193} \to \mathrm{FLim}(f,c,\alpha) \to \varphi_{248}$\lbl{A1}}
\BIC{$\varphi_{193} \to \mathrm{FLim}(f,c,\alpha) \to \varphi_{248}$}
\BIC{$\mathrm{FLim}(f,c,\alpha) \to \varphi_{248}$}
\AXC{$\varphi_{355}$\lbl{A1}}
\BIC{$\varphi_{133} \to \neg (\delta \in \mathbb{R} \land 0 < \delta \land \varphi_{78})$}
\UIC{$\forall \delta\,(\varphi_{133} \to \neg (\delta \in \mathbb{R} \land 0 < \delta \land \varphi_{78}))$}
\end{ptree}
\noindent$\Pi_{87}$:
\begin{ptree}
\AXC{$\forall f\,\forall c\,\forall \alpha\,\forall \varepsilon\,(\varphi_{154})$\lbl{[D1]}}
\AXC{$\forall f\,\forall c\,\forall \alpha\,\forall \varepsilon\,(\varphi_{154}) \to \forall c\,\forall \alpha\,\forall \varepsilon\,(\varphi_{154})$\lbl{A2}}
\BIC{$\forall c\,\forall \alpha\,\forall \varepsilon\,(\varphi_{154})$}
\AXC{$\forall c\,\forall \alpha\,\forall \varepsilon\,(\varphi_{154}) \to \forall \alpha\,\forall \varepsilon\,(\varphi_{154})$\lbl{A2}}
\BIC{$\forall \alpha\,\forall \varepsilon\,(\varphi_{154})$}
\end{ptree}
\noindent$\Pi_{88}$:
\begin{ptree}
\AXC{$\stackrel{\Pi_{87}}{\vdots}$}\noLine
\UIC{$\forall \alpha\,\forall \varepsilon\,(\varphi_{154})$}
\AXC{$\forall \alpha\,\forall \varepsilon\,(\varphi_{154}) \to \forall \varepsilon\,(\varphi_{154})$\lbl{A2}}
\BIC{$\forall \varepsilon\,(\varphi_{154})$}
\AXC{$\forall \varepsilon\,(\varphi_{154}) \to \varphi_{154}$\lbl{A2}}
\BIC{$\varphi_{154}$}
\AXC{$\varphi_{154} \to \varphi_{396}$\lbl{A1}}
\BIC{$\varphi_{396}$}
\end{ptree}
\noindent$\Pi_{89}$:
\begin{ptree}
\AXC{$\varphi_{133} \to \forall \delta\,(\varphi_{133})$\lbl{A4}}
\AXC{$\varphi_{359}$\lbl{A3}}
\AXC{$\stackrel{\Pi_{86}}{\vdots}$}\noLine
\UIC{$\forall \delta\,(\varphi_{133} \to \neg (\delta \in \mathbb{R} \land 0 < \delta \land \varphi_{78}))$}
\AXC{$\stackrel{\Pi_{88}}{\vdots}$}\noLine
\UIC{$\varphi_{396}$}
\BIC{$\varphi_{359} \to (\varphi_{133} \to \forall \delta\,(\varphi_{133})) \to \varphi_{131}$}
\BIC{$(\varphi_{133} \to \forall \delta\,(\varphi_{133})) \to \varphi_{131}$}
\BIC{$\varphi_{131}$}
\UIC{$\forall \varepsilon\,(\varphi_{131})$}
\AXC{$\forall \varepsilon\,(\varphi_{131}) \to (\forall \varepsilon\,(\varphi_{131}) \to \varphi_{145}) \to \varphi_{180}$\lbl{A1}}
\BIC{$(\forall \varepsilon\,(\varphi_{131}) \to \varphi_{145}) \to \varphi_{180}$}
\end{ptree}
\noindent$\Pi_{90}$:
\begin{ptree}
\AXC{$\forall f\,\forall x\,(\mathrm{RFn}(f) \to x \in \mathrm{dom}(f) \to x \in \mathbb{R})$\lbl{[S7]}}
\AXC{$\varphi_{66}$\lbl{A2}}
\BIC{$\forall x\,(\mathrm{RFn}(f) \to x \in \mathrm{dom}(f) \to x \in \mathbb{R})$}
\AXC{$\varphi_{49}$\lbl{A2}}
\BIC{$\mathrm{RFn}(f) \to x \in \mathrm{dom}(f) \to x \in \mathbb{R}$}
\AXC{$\varphi_{349}$\lbl{A1}}
\BIC{$(\mathrm{FLim}(f,c,\alpha) \to \mathrm{RFn}(f)) \to \varphi_{337}$}
\end{ptree}
\noindent$\Pi_{91}$:
\begin{ptree}
\AXC{$\mathrm{FLim}(f,c,\alpha) \to \varepsilon \in \mathbb{R} \to \varphi_{151}$\lbl{A1}}
\AXC{$\stackrel{\Pi_{76}}{\vdots}$}\noLine
\UIC{$\mathrm{FLim}(f,c,\alpha) \to \mathrm{RFn}(f)$}
\AXC{$\stackrel{\Pi_{90}}{\vdots}$}\noLine
\UIC{$(\mathrm{FLim}(f,c,\alpha) \to \mathrm{RFn}(f)) \to \varphi_{337}$}
\BIC{$\varphi_{337}$}
\BIC{$\mathrm{FLim}(f,c,\alpha) \to \varepsilon \in \mathbb{R} \to 0 < \varepsilon \to \varphi_{164}$}
\end{ptree}
\noindent$\Pi_{92}$:
\begin{ptree}
\AXC{$\forall x\,\forall c\,(\varphi_{100})$\lbl{[H3]}}
\AXC{$\forall x\,\forall c\,(\varphi_{100}) \to \forall c\,(\varphi_{100})$\lbl{A2}}
\BIC{$\forall c\,(\varphi_{100})$}
\AXC{$\forall c\,(\varphi_{100}) \to \varphi_{100}$\lbl{A2}}
\BIC{$\varphi_{100}$}
\AXC{$\varphi_{100} \to x \in \mathbb{R} \to \varphi_{87}$\lbl{A2}}
\BIC{$x \in \mathbb{R} \to \varphi_{87}$}
\end{ptree}
\noindent$\Pi_{93}$:
\begin{ptree}
\AXC{$\stackrel{\Pi_{75}}{\vdots}$}\noLine
\UIC{$\mathrm{FLim}(f,c,\alpha) \to c \in \mathbb{R}$}
\AXC{$\stackrel{\Pi_{91}}{\vdots}$}\noLine
\UIC{$\mathrm{FLim}(f,c,\alpha) \to \varepsilon \in \mathbb{R} \to 0 < \varepsilon \to \varphi_{164}$}
\AXC{$\stackrel{\Pi_{92}}{\vdots}$}\noLine
\UIC{$x \in \mathbb{R} \to \varphi_{87}$}
\AXC{$(x \in \mathbb{R} \to \varphi_{87}) \to \varphi_{385}$\lbl{A1}}
\BIC{$\varphi_{385}$}
\BIC{$(\mathrm{FLim}(f,c,\alpha) \to c \in \mathbb{R}) \to \varphi_{168} \to \varphi_{370}$}
\BIC{$\varphi_{168} \to \varphi_{370}$}
\end{ptree}
\noindent$\Pi_{94}$:
\begin{ptree}
\AXC{$\mathrm{FLim}(f,c,\alpha) \to \varepsilon \in \mathbb{R} \to \varphi_{153}$\lbl{A1}}
\AXC{$\varphi_{262}$\lbl{A1}}
\AXC{$\varphi_{262} \to \varphi_{168}$\lbl{A1}}
\BIC{$\varphi_{168}$}
\AXC{$\stackrel{\Pi_{93}}{\vdots}$}\noLine
\UIC{$\varphi_{168} \to \varphi_{370}$}
\BIC{$\varphi_{370}$}
\BIC{$(\mathrm{FLim}(f,c,\alpha) \to \varepsilon \in \mathbb{R} \to 0 < \varepsilon \to \varphi_{167}) \to \varphi_{246}$}
\end{ptree}
\noindent$\Pi_{95}$:
\begin{ptree}
\AXC{$\mathrm{FLim}(f,c,\alpha) \to \varepsilon \in \mathbb{R} \to 0 < \varepsilon \to \varphi_{167}$\lbl{A1}}
\AXC{$\stackrel{\Pi_{94}}{\vdots}$}\noLine
\UIC{$(\mathrm{FLim}(f,c,\alpha) \to \varepsilon \in \mathbb{R} \to 0 < \varepsilon \to \varphi_{167}) \to \varphi_{246}$}
\BIC{$\varphi_{246}$}
\end{ptree}
\noindent$\Pi_{96}$:
\begin{ptree}
\AXC{$\varphi_{78} \to \varphi_{69}$\lbl{A2}}
\AXC{$\varphi_{262}$\lbl{A1}}
\AXC{$\varphi_{262} \to \varphi_{250}$\lbl{A1}}
\BIC{$\varphi_{250}$}
\AXC{$\varphi_{250} \to (\varphi_{78} \to \varphi_{69}) \to \varphi_{247}$\lbl{A1}}
\BIC{$(\varphi_{78} \to \varphi_{69}) \to \varphi_{247}$}
\BIC{$\varphi_{247}$}
\AXC{$\varphi_{247} \to \varphi_{373}$\lbl{A1}}
\BIC{$\varphi_{373}$}
\end{ptree}
\noindent$\Pi_{97}$:
\begin{ptree}
\AXC{$\stackrel{\Pi_{95}}{\vdots}$}\noLine
\UIC{$\varphi_{246}$}
\AXC{$\mathrm{FLim}(f,c,\alpha) \to \varepsilon \in \mathbb{R} \to \varphi_{151}$\lbl{A1}}
\AXC{$\stackrel{\Pi_{96}}{\vdots}$}\noLine
\UIC{$\varphi_{373}$}
\BIC{$\varphi_{246} \to \mathrm{FLim}(f,c,\alpha) \to \varepsilon \in \mathbb{R} \to 0 < \varepsilon \to \varphi_{179}$}
\BIC{$\mathrm{FLim}(f,c,\alpha) \to \varepsilon \in \mathbb{R} \to 0 < \varepsilon \to \varphi_{179}$}
\AXC{$(\mathrm{FLim}(f,c,\alpha) \to \varepsilon \in \mathbb{R} \to 0 < \varepsilon \to \varphi_{179}) \to \varphi_{233}$\lbl{A1}}
\BIC{$\varphi_{233}$}
\UIC{$\forall x\,(\varphi_{233})$}
\end{ptree}
\noindent$\Pi_{98}$:
\begin{ptree}
\AXC{$\varphi_{150} \to \forall x\,(\varphi_{150})$\lbl{A4}}
\AXC{$\forall x\,(\varphi_{233}) \to \forall x\,(\varphi_{150}) \to \varphi_{33}$\lbl{A3}}
\AXC{$\stackrel{\Pi_{97}}{\vdots}$}\noLine
\UIC{$\forall x\,(\varphi_{233})$}
\AXC{$\varphi_{395}$\lbl{A1}}
\BIC{$(\forall x\,(\varphi_{233}) \to \forall x\,(\varphi_{150}) \to \varphi_{33}) \to (\varphi_{150} \to \forall x\,(\varphi_{150})) \to \varphi_{238}$}
\BIC{$(\varphi_{150} \to \forall x\,(\varphi_{150})) \to \varphi_{238}$}
\BIC{$\varphi_{238}$}
\end{ptree}
\noindent$\Pi_{99}$:
\begin{ptree}
\AXC{$\varphi_{194}$\lbl{A2}}
\AXC{$\stackrel{\Pi_{98}}{\vdots}$}\noLine
\UIC{$\varphi_{238}$}
\AXC{$\varphi_{262}$\lbl{A1}}
\AXC{$\varphi_{262} \to \varphi_{238} \to \varphi_{255}$\lbl{A1}}
\BIC{$\varphi_{238} \to \varphi_{255}$}
\BIC{$\varphi_{255}$}
\AXC{$\varphi_{255} \to \varphi_{194} \to \mathrm{FLim}(f,c,\alpha) \to \varphi_{249}$\lbl{A1}}
\BIC{$\varphi_{194} \to \mathrm{FLim}(f,c,\alpha) \to \varphi_{249}$}
\BIC{$\mathrm{FLim}(f,c,\alpha) \to \varphi_{249}$}
\AXC{$\varphi_{356}$\lbl{A1}}
\BIC{$\varphi_{134} \to \neg (\delta \in \mathbb{R} \land 0 < \delta \land \varphi_{78})$}
\UIC{$\forall \delta\,(\varphi_{134} \to \neg (\delta \in \mathbb{R} \land 0 < \delta \land \varphi_{78}))$}
\end{ptree}
\noindent$\Pi_{100}$:
\begin{ptree}
\AXC{$\forall f\,\forall c\,\forall \alpha\,\forall \varepsilon\,(\varphi_{154})$\lbl{[D1]}}
\AXC{$\forall f\,\forall c\,\forall \alpha\,\forall \varepsilon\,(\varphi_{154}) \to \forall c\,\forall \alpha\,\forall \varepsilon\,(\varphi_{154})$\lbl{A2}}
\BIC{$\forall c\,\forall \alpha\,\forall \varepsilon\,(\varphi_{154})$}
\AXC{$\forall c\,\forall \alpha\,\forall \varepsilon\,(\varphi_{154}) \to \forall \alpha\,\forall \varepsilon\,(\varphi_{154})$\lbl{A2}}
\BIC{$\forall \alpha\,\forall \varepsilon\,(\varphi_{154})$}
\end{ptree}
\noindent$\Pi_{101}$:
\begin{ptree}
\AXC{$\stackrel{\Pi_{100}}{\vdots}$}\noLine
\UIC{$\forall \alpha\,\forall \varepsilon\,(\varphi_{154})$}
\AXC{$\forall \alpha\,\forall \varepsilon\,(\varphi_{154}) \to \forall \varepsilon\,(\varphi_{154})$\lbl{A2}}
\BIC{$\forall \varepsilon\,(\varphi_{154})$}
\AXC{$\forall \varepsilon\,(\varphi_{154}) \to \varphi_{154}$\lbl{A2}}
\BIC{$\varphi_{154}$}
\AXC{$\varphi_{154} \to \varphi_{397}$\lbl{A1}}
\BIC{$\varphi_{397}$}
\end{ptree}
\noindent$\Pi_{102}$:
\begin{ptree}
\AXC{$\varphi_{134} \to \forall \delta\,(\varphi_{134})$\lbl{A4}}
\AXC{$\varphi_{360}$\lbl{A3}}
\AXC{$\stackrel{\Pi_{99}}{\vdots}$}\noLine
\UIC{$\forall \delta\,(\varphi_{134} \to \neg (\delta \in \mathbb{R} \land 0 < \delta \land \varphi_{78}))$}
\AXC{$\stackrel{\Pi_{101}}{\vdots}$}\noLine
\UIC{$\varphi_{397}$}
\BIC{$\varphi_{360} \to (\varphi_{134} \to \forall \delta\,(\varphi_{134})) \to \varphi_{132}$}
\BIC{$(\varphi_{134} \to \forall \delta\,(\varphi_{134})) \to \varphi_{132}$}
\BIC{$\varphi_{132}$}
\UIC{$\forall \varepsilon\,(\varphi_{132})$}
\AXC{$\forall \varepsilon\,(\varphi_{132}) \to (\forall \varepsilon\,(\varphi_{132}) \to \varphi_{146}) \to \varphi_{181}$\lbl{A1}}
\BIC{$(\forall \varepsilon\,(\varphi_{132}) \to \varphi_{146}) \to \varphi_{181}$}
\end{ptree}
\noindent$\Pi_{103}$:
\begin{ptree}
\AXC{$\mathrm{FLim}(f,c,\alpha) \to \forall \varepsilon\,(\mathrm{FLim}(f,c,\alpha))$\lbl{A4}}
\AXC{$\forall \varepsilon\,(\varphi_{132}) \to \varphi_{146}$\lbl{A3}}
\AXC{$\stackrel{\Pi_{102}}{\vdots}$}\noLine
\UIC{$(\forall \varepsilon\,(\varphi_{132}) \to \varphi_{146}) \to \varphi_{181}$}
\BIC{$\varphi_{181}$}
\BIC{$\varphi_{139}$}
\end{ptree}
\noindent$\Pi_{104}$:
\begin{ptree}
\AXC{$\forall x\,\forall c\,(\varphi_{99})$\lbl{[H5]}}
\AXC{$\forall x\,\forall c\,(\varphi_{99}) \to \forall c\,(\varphi_{99})$\lbl{A2}}
\BIC{$\forall c\,(\varphi_{99})$}
\AXC{$\forall c\,(\varphi_{99}) \to \varphi_{99}$\lbl{A2}}
\BIC{$\varphi_{99}$}
\AXC{$\varphi_{99} \to x \in \mathbb{R} \to \varphi_{83}$\lbl{A2}}
\BIC{$x \in \mathbb{R} \to \varphi_{83}$}
\AXC{$\varphi_{381}$\lbl{A1}}
\BIC{$\varphi_{124} \to (\mathrm{FLim}(f,c,\alpha) \to c \in \mathbb{R}) \to \varphi_{367}$}
\end{ptree}
\noindent$\Pi_{105}$:
\begin{ptree}
\AXC{$\stackrel{\Pi_{75}}{\vdots}$}\noLine
\UIC{$\mathrm{FLim}(f,c,\alpha) \to c \in \mathbb{R}$}
\AXC{$\varphi_{124}$\lbl{A1}}
\AXC{$\stackrel{\Pi_{104}}{\vdots}$}\noLine
\UIC{$\varphi_{124} \to (\mathrm{FLim}(f,c,\alpha) \to c \in \mathbb{R}) \to \varphi_{367}$}
\BIC{$(\mathrm{FLim}(f,c,\alpha) \to c \in \mathbb{R}) \to \varphi_{367}$}
\BIC{$\varphi_{367}$}
\end{ptree}
\noindent$\Pi_{106}$:
\begin{ptree}
\AXC{$\varphi_{147}$\lbl{A1}}
\AXC{$\varphi_{142}$\lbl{A1}}
\AXC{$\varphi_{195}$\lbl{A1}}
\AXC{$\varphi_{195} \to \mathrm{FLim}(f,c,\alpha) \to r \in \mathbb{R} \land 0 < r \land \varphi_{26} \to r \in \mathbb{R}$\lbl{A1}}
\BIC{$\mathrm{FLim}(f,c,\alpha) \to r \in \mathbb{R} \land 0 < r \land \varphi_{26} \to r \in \mathbb{R}$}
\AXC{$\stackrel{\Pi_{105}}{\vdots}$}\noLine
\UIC{$\varphi_{367}$}
\BIC{$\varphi_{142} \to \varphi_{147} \to \mathrm{FLim}(f,c,\alpha) \to \varphi_{169}$}
\BIC{$\varphi_{147} \to \mathrm{FLim}(f,c,\alpha) \to \varphi_{169}$}
\BIC{$\mathrm{FLim}(f,c,\alpha) \to \varphi_{169}$}
\end{ptree}
\noindent$\Pi_{107}$:
\begin{ptree}
\AXC{$\stackrel{\Pi_{106}}{\vdots}$}\noLine
\UIC{$\mathrm{FLim}(f,c,\alpha) \to \varphi_{169}$}
\AXC{$\varphi_{124}$\lbl{A1}}
\AXC{$\varphi_{127}$\lbl{A2}}
\AXC{$\varphi_{195}$\lbl{A1}}
\AXC{$\varphi_{195} \to \varphi_{219}$\lbl{A1}}
\BIC{$\varphi_{219}$}
\AXC{$\varphi_{219} \to \varphi_{127} \to \varphi_{215}$\lbl{A1}}
\BIC{$\varphi_{127} \to \varphi_{215}$}
\BIC{$\varphi_{215}$}
\AXC{$\varphi_{215} \to \varphi_{351}$\lbl{A1}}
\BIC{$\varphi_{351}$}
\BIC{$(\mathrm{FLim}(f,c,\alpha) \to \varphi_{169}) \to \mathrm{FLim}(f,c,\alpha) \to \varphi_{136}$}
\BIC{$\mathrm{FLim}(f,c,\alpha) \to \varphi_{136}$}
\end{ptree}
\noindent$\Pi_{108}$:
\begin{ptree}
\AXC{$\varphi_{225}$\lbl{A4}}
\AXC{$\forall x\,(\varphi_{158}) \to \varphi_{165}$\lbl{A3}}
\AXC{$\stackrel{\Pi_{107}}{\vdots}$}\noLine
\UIC{$\mathrm{FLim}(f,c,\alpha) \to \varphi_{136}$}
\AXC{$(\mathrm{FLim}(f,c,\alpha) \to \varphi_{136}) \to \varphi_{158}$\lbl{A1}}
\BIC{$\varphi_{158}$}
\UIC{$\forall x\,(\varphi_{158})$}
\AXC{$\varphi_{382}$\lbl{A1}}
\BIC{$(\forall x\,(\varphi_{158}) \to \varphi_{165}) \to \varphi_{225} \to \mathrm{FLim}(f,c,\alpha) \to \varphi_{140}$}
\BIC{$\varphi_{225} \to \mathrm{FLim}(f,c,\alpha) \to \varphi_{140}$}
\BIC{$\mathrm{FLim}(f,c,\alpha) \to \varphi_{140}$}
\end{ptree}
\noindent$\Pi_{109}$:
\begin{ptree}
\AXC{$\stackrel{\Pi_{108}}{\vdots}$}\noLine
\UIC{$\mathrm{FLim}(f,c,\alpha) \to \varphi_{140}$}
\AXC{$\varphi_{195}$\lbl{A1}}
\AXC{$\varphi_{195} \to \varphi_{308}$\lbl{A1}}
\BIC{$\varphi_{308}$}
\BIC{$\mathrm{FLim}(f,c,\alpha) \to r \in \mathbb{R} \land 0 < r \land \varphi_{26} \to \varphi_{47}$}
\AXC{$(\mathrm{FLim}(f,c,\alpha) \to r \in \mathbb{R} \land 0 < r \land \varphi_{26} \to \varphi_{47}) \to \varphi_{309}$\lbl{A1}}
\BIC{$\varphi_{309}$}
\end{ptree}
\noindent$\Pi_{110}$:
\begin{ptree}
\AXC{$\forall f\,\forall c\,(\mathrm{LNbhd}(f,c) \leftrightarrow \exists r\,(\varphi_{47}))$\lbl{DefLNbhd}}
\AXC{$\varphi_{187}$\lbl{A2}}
\BIC{$\forall c\,(\mathrm{LNbhd}(f,c) \leftrightarrow \exists r\,(\varphi_{47}))$}
\AXC{$\varphi_{184}$\lbl{A2}}
\BIC{$\mathrm{LNbhd}(f,c) \leftrightarrow \exists r\,(\varphi_{47})$}
\AXC{$(\mathrm{LNbhd}(f,c) \leftrightarrow \exists r\,(\varphi_{47})) \to \varphi_{289}$\lbl{A1}}
\BIC{$\varphi_{289}$}
\end{ptree}
\noindent$\Pi_{111}$:
\begin{ptree}
\AXC{$\forall r\,(\neg \varphi_{47}) \to \neg \varphi_{47}$\lbl{A2}}
\AXC{$\stackrel{\Pi_{109}}{\vdots}$}\noLine
\UIC{$\varphi_{309}$}
\BIC{$\mathrm{FLim}(f,c,\alpha) \to r \in \mathbb{R} \land 0 < r \land \varphi_{26} \to \exists r\,(\varphi_{47})$}
\AXC{$\stackrel{\Pi_{110}}{\vdots}$}\noLine
\UIC{$\varphi_{289}$}
\BIC{$\varphi_{113}$}
\AXC{$\varphi_{113} \to \varphi_{115}$\lbl{A1}}
\BIC{$\varphi_{115}$}
\UIC{$\forall r\,(\varphi_{115})$}
\end{ptree}
\noindent$\Pi_{112}$:
\begin{ptree}
\AXC{$\forall f\,\forall c\,\forall \alpha\,(\mathrm{FLim}(f,c,\alpha) \to \mathrm{PNbhd}(f,c))$\lbl{[D5]}}
\AXC{$\varphi_{39}$\lbl{A2}}
\BIC{$\forall c\,\forall \alpha\,(\mathrm{FLim}(f,c,\alpha) \to \mathrm{PNbhd}(f,c))$}
\AXC{$\varphi_{20}$\lbl{A2}}
\BIC{$\forall \alpha\,(\mathrm{FLim}(f,c,\alpha) \to \mathrm{PNbhd}(f,c))$}
\AXC{$\varphi_{12}$\lbl{A2}}
\BIC{$\mathrm{FLim}(f,c,\alpha) \to \mathrm{PNbhd}(f,c)$}
\end{ptree}
\noindent$\Pi_{113}$:
\begin{ptree}
\AXC{$\stackrel{\Pi_{112}}{\vdots}$}\noLine
\UIC{$\mathrm{FLim}(f,c,\alpha) \to \mathrm{PNbhd}(f,c)$}
\AXC{$(\mathrm{FLim}(f,c,\alpha) \to \mathrm{PNbhd}(f,c)) \to \varphi_{226}$\lbl{A1}}
\BIC{$\varphi_{226}$}
\end{ptree}
\noindent$\Pi_{114}$:
\begin{ptree}
\AXC{$\varphi_{112}$\lbl{DefPNbhd}}
\AXC{$\varphi_{239}$\lbl{A2}}
\BIC{$\forall c\,(\mathrm{PNbhd}(f,c) \leftrightarrow \exists r\,(r \in \mathbb{R} \land 0 < r \land \varphi_{26}))$}
\AXC{$\varphi_{228}$\lbl{A2}}
\BIC{$\mathrm{PNbhd}(f,c) \leftrightarrow \exists r\,(r \in \mathbb{R} \land 0 < r \land \varphi_{26})$}
\AXC{$\stackrel{\Pi_{113}}{\vdots}$}\noLine
\UIC{$\varphi_{226}$}
\BIC{$\mathrm{FLim}(f,c,\alpha) \to \exists r\,(r \in \mathbb{R} \land 0 < r \land \varphi_{26})$}
\end{ptree}
\noindent$\Pi_{115}$:
\begin{ptree}
\AXC{$\varphi_{13}$\lbl{A4}}
\AXC{$\forall r\,(\varphi_{115}) \to \varphi_{119}$\lbl{A3}}
\AXC{$\stackrel{\Pi_{111}}{\vdots}$}\noLine
\UIC{$\forall r\,(\varphi_{115})$}
\AXC{$\stackrel{\Pi_{114}}{\vdots}$}\noLine
\UIC{$\mathrm{FLim}(f,c,\alpha) \to \exists r\,(r \in \mathbb{R} \land 0 < r \land \varphi_{26})$}
\AXC{$\varphi_{364}$\lbl{A1}}
\BIC{$\forall r\,(\varphi_{115}) \to \varphi_{303}$}
\BIC{$\varphi_{303}$}
\BIC{$\varphi_{13} \to \mathrm{FLim}(f,c,\alpha) \to \mathrm{LNbhd}(f,c)$}
\BIC{$\mathrm{FLim}(f,c,\alpha) \to \mathrm{LNbhd}(f,c)$}
\end{ptree}
\noindent$\Pi_{116}$:
\begin{ptree}
\AXC{$\forall x\,\forall c\,(\varphi_{100})$\lbl{[H3]}}
\AXC{$\forall x\,\forall c\,(\varphi_{100}) \to \forall c\,(\varphi_{100})$\lbl{A2}}
\BIC{$\forall c\,(\varphi_{100})$}
\AXC{$\forall c\,(\varphi_{100}) \to \varphi_{100}$\lbl{A2}}
\BIC{$\varphi_{100}$}
\AXC{$\varphi_{100} \to x \in \mathbb{R} \to \varphi_{84}$\lbl{A2}}
\BIC{$x \in \mathbb{R} \to \varphi_{84}$}
\AXC{$\varphi_{380}$\lbl{A1}}
\BIC{$\varphi_{124} \to (\mathrm{FLim}(f,c,\alpha) \to c \in \mathbb{R}) \to \varphi_{366}$}
\end{ptree}
\noindent$\Pi_{117}$:
\begin{ptree}
\AXC{$\stackrel{\Pi_{75}}{\vdots}$}\noLine
\UIC{$\mathrm{FLim}(f,c,\alpha) \to c \in \mathbb{R}$}
\AXC{$\varphi_{124}$\lbl{A1}}
\AXC{$\stackrel{\Pi_{116}}{\vdots}$}\noLine
\UIC{$\varphi_{124} \to (\mathrm{FLim}(f,c,\alpha) \to c \in \mathbb{R}) \to \varphi_{366}$}
\BIC{$(\mathrm{FLim}(f,c,\alpha) \to c \in \mathbb{R}) \to \varphi_{366}$}
\BIC{$\varphi_{366}$}
\end{ptree}
\noindent$\Pi_{118}$:
\begin{ptree}
\AXC{$\mathrm{FLim}(f,c,\alpha) \to \varphi_{135}$\lbl{A1}}
\AXC{$\varphi_{128}$\lbl{A1}}
\AXC{$\varphi_{195}$\lbl{A1}}
\AXC{$\varphi_{195} \to \mathrm{FLim}(f,c,\alpha) \to r \in \mathbb{R} \land 0 < r \land \varphi_{26} \to r \in \mathbb{R}$\lbl{A1}}
\BIC{$\mathrm{FLim}(f,c,\alpha) \to r \in \mathbb{R} \land 0 < r \land \varphi_{26} \to r \in \mathbb{R}$}
\AXC{$\stackrel{\Pi_{117}}{\vdots}$}\noLine
\UIC{$\varphi_{366}$}
\BIC{$\varphi_{353}$}
\BIC{$(\mathrm{FLim}(f,c,\alpha) \to \varphi_{135}) \to \mathrm{FLim}(f,c,\alpha) \to \varphi_{170}$}
\BIC{$\mathrm{FLim}(f,c,\alpha) \to \varphi_{170}$}
\end{ptree}
\noindent$\Pi_{119}$:
\begin{ptree}
\AXC{$\stackrel{\Pi_{118}}{\vdots}$}\noLine
\UIC{$\mathrm{FLim}(f,c,\alpha) \to \varphi_{170}$}
\AXC{$\varphi_{124}$\lbl{A1}}
\AXC{$\varphi_{127}$\lbl{A2}}
\AXC{$\varphi_{195}$\lbl{A1}}
\AXC{$\varphi_{195} \to \varphi_{220}$\lbl{A1}}
\BIC{$\varphi_{220}$}
\AXC{$\varphi_{220} \to \varphi_{127} \to \varphi_{216}$\lbl{A1}}
\BIC{$\varphi_{127} \to \varphi_{216}$}
\BIC{$\varphi_{216}$}
\AXC{$\varphi_{216} \to \varphi_{352}$\lbl{A1}}
\BIC{$\varphi_{352}$}
\BIC{$(\mathrm{FLim}(f,c,\alpha) \to \varphi_{170}) \to \mathrm{FLim}(f,c,\alpha) \to \varphi_{137}$}
\BIC{$\mathrm{FLim}(f,c,\alpha) \to \varphi_{137}$}
\end{ptree}
\noindent$\Pi_{120}$:
\begin{ptree}
\AXC{$\varphi_{225}$\lbl{A4}}
\AXC{$\forall x\,(\varphi_{159}) \to \varphi_{166}$\lbl{A3}}
\AXC{$\stackrel{\Pi_{119}}{\vdots}$}\noLine
\UIC{$\mathrm{FLim}(f,c,\alpha) \to \varphi_{137}$}
\AXC{$(\mathrm{FLim}(f,c,\alpha) \to \varphi_{137}) \to \varphi_{159}$\lbl{A1}}
\BIC{$\varphi_{159}$}
\UIC{$\forall x\,(\varphi_{159})$}
\AXC{$\varphi_{383}$\lbl{A1}}
\BIC{$(\forall x\,(\varphi_{159}) \to \varphi_{166}) \to \varphi_{225} \to \mathrm{FLim}(f,c,\alpha) \to \varphi_{141}$}
\BIC{$\varphi_{225} \to \mathrm{FLim}(f,c,\alpha) \to \varphi_{141}$}
\BIC{$\mathrm{FLim}(f,c,\alpha) \to \varphi_{141}$}
\end{ptree}
\noindent$\Pi_{121}$:
\begin{ptree}
\AXC{$\stackrel{\Pi_{120}}{\vdots}$}\noLine
\UIC{$\mathrm{FLim}(f,c,\alpha) \to \varphi_{141}$}
\AXC{$\varphi_{195}$\lbl{A1}}
\AXC{$\varphi_{195} \to \varphi_{310}$\lbl{A1}}
\BIC{$\varphi_{310}$}
\BIC{$\mathrm{FLim}(f,c,\alpha) \to r \in \mathbb{R} \land 0 < r \land \varphi_{26} \to \varphi_{48}$}
\AXC{$(\mathrm{FLim}(f,c,\alpha) \to r \in \mathbb{R} \land 0 < r \land \varphi_{26} \to \varphi_{48}) \to \varphi_{311}$\lbl{A1}}
\BIC{$\varphi_{311}$}
\end{ptree}
\noindent$\Pi_{122}$:
\begin{ptree}
\AXC{$\forall f\,\forall c\,(\mathrm{RNbhd}(f,c) \leftrightarrow \exists r\,(\varphi_{48}))$\lbl{DefRNbhd}}
\AXC{$\varphi_{188}$\lbl{A2}}
\BIC{$\forall c\,(\mathrm{RNbhd}(f,c) \leftrightarrow \exists r\,(\varphi_{48}))$}
\AXC{$\varphi_{185}$\lbl{A2}}
\BIC{$\mathrm{RNbhd}(f,c) \leftrightarrow \exists r\,(\varphi_{48})$}
\AXC{$(\mathrm{RNbhd}(f,c) \leftrightarrow \exists r\,(\varphi_{48})) \to \varphi_{290}$\lbl{A1}}
\BIC{$\varphi_{290}$}
\end{ptree}
\noindent$\Pi_{123}$:
\begin{ptree}
\AXC{$\forall r\,(\neg \varphi_{48}) \to \neg \varphi_{48}$\lbl{A2}}
\AXC{$\stackrel{\Pi_{121}}{\vdots}$}\noLine
\UIC{$\varphi_{311}$}
\BIC{$\mathrm{FLim}(f,c,\alpha) \to r \in \mathbb{R} \land 0 < r \land \varphi_{26} \to \exists r\,(\varphi_{48})$}
\AXC{$\stackrel{\Pi_{122}}{\vdots}$}\noLine
\UIC{$\varphi_{290}$}
\BIC{$\varphi_{114}$}
\AXC{$\varphi_{114} \to \varphi_{116}$\lbl{A1}}
\BIC{$\varphi_{116}$}
\UIC{$\forall r\,(\varphi_{116})$}
\end{ptree}
\noindent$\Pi_{124}$:
\begin{ptree}
\AXC{$\varphi_{14}$\lbl{A4}}
\AXC{$\forall r\,(\varphi_{116}) \to \varphi_{120}$\lbl{A3}}
\AXC{$\stackrel{\Pi_{123}}{\vdots}$}\noLine
\UIC{$\forall r\,(\varphi_{116})$}
\AXC{$\stackrel{\Pi_{114}}{\vdots}$}\noLine
\UIC{$\mathrm{FLim}(f,c,\alpha) \to \exists r\,(r \in \mathbb{R} \land 0 < r \land \varphi_{26})$}
\AXC{$\varphi_{365}$\lbl{A1}}
\BIC{$\forall r\,(\varphi_{116}) \to \varphi_{304}$}
\BIC{$\varphi_{304}$}
\BIC{$\varphi_{14} \to \mathrm{FLim}(f,c,\alpha) \to \mathrm{RNbhd}(f,c)$}
\BIC{$\mathrm{FLim}(f,c,\alpha) \to \mathrm{RNbhd}(f,c)$}
\end{ptree}
\noindent$\Pi_{125}$:
\begin{ptree}
\AXC{$\forall f\,\forall c\,\forall \alpha\,(\mathrm{FLim}(f,c,\alpha) \to \alpha \in \mathbb{R})$\lbl{[D4]}}
\AXC{$\varphi_{41}$\lbl{A2}}
\BIC{$\forall c\,\forall \alpha\,(\mathrm{FLim}(f,c,\alpha) \to \alpha \in \mathbb{R})$}
\AXC{$\varphi_{24}$\lbl{A2}}
\BIC{$\forall \alpha\,(\mathrm{FLim}(f,c,\alpha) \to \alpha \in \mathbb{R})$}
\AXC{$\varphi_{16}$\lbl{A2}}
\BIC{$\mathrm{FLim}(f,c,\alpha) \to \alpha \in \mathbb{R}$}
\end{ptree}
\noindent$\Pi_{126}$:
\begin{ptree}
\AXC{$\stackrel{\Pi_{76}}{\vdots}$}\noLine
\UIC{$\mathrm{FLim}(f,c,\alpha) \to \mathrm{RFn}(f)$}
\AXC{$\stackrel{\Pi_{28}}{\vdots}$}\noLine
\UIC{$\varphi_{190}$}
\AXC{$\stackrel{\Pi_{3}}{\vdots}$}\noLine
\UIC{$\varphi_{189}$}
\AXC{$\varphi_{189} \to \varphi_{190} \to (\mathrm{FLim}(f,c,\alpha) \to \mathrm{RFn}(f)) \to \varphi_{350}$\lbl{A1}}
\BIC{$\varphi_{190} \to (\mathrm{FLim}(f,c,\alpha) \to \mathrm{RFn}(f)) \to \varphi_{350}$}
\BIC{$(\mathrm{FLim}(f,c,\alpha) \to \mathrm{RFn}(f)) \to \varphi_{350}$}
\BIC{$\varphi_{350}$}
\end{ptree}
\noindent$\Pi_{127}$:
\begin{ptree}
\AXC{$\stackrel{\Pi_{124}}{\vdots}$}\noLine
\UIC{$\mathrm{FLim}(f,c,\alpha) \to \mathrm{RNbhd}(f,c)$}
\AXC{$\stackrel{\Pi_{125}}{\vdots}$}\noLine
\UIC{$\mathrm{FLim}(f,c,\alpha) \to \alpha \in \mathbb{R}$}
\AXC{$\stackrel{\Pi_{75}}{\vdots}$}\noLine
\UIC{$\mathrm{FLim}(f,c,\alpha) \to c \in \mathbb{R}$}
\AXC{$\stackrel{\Pi_{126}}{\vdots}$}\noLine
\UIC{$\varphi_{350}$}
\BIC{$(\mathrm{FLim}(f,c,\alpha) \to \alpha \in \mathbb{R}) \to \varphi_{339}$}
\BIC{$\varphi_{339}$}
\BIC{$(\mathrm{FLim}(f,c,\alpha) \to \mathrm{LNbhd}(f,c)) \to \varphi_{139} \to \varphi_{138} \to \varphi_{5}$}
\end{ptree}
\noindent$\Pi_{128}$:
\begin{ptree}
\AXC{$\stackrel{\Pi_{103}}{\vdots}$}\noLine
\UIC{$\varphi_{139}$}
\AXC{$\stackrel{\Pi_{115}}{\vdots}$}\noLine
\UIC{$\mathrm{FLim}(f,c,\alpha) \to \mathrm{LNbhd}(f,c)$}
\AXC{$\stackrel{\Pi_{127}}{\vdots}$}\noLine
\UIC{$(\mathrm{FLim}(f,c,\alpha) \to \mathrm{LNbhd}(f,c)) \to \varphi_{139} \to \varphi_{138} \to \varphi_{5}$}
\BIC{$\varphi_{139} \to \varphi_{138} \to \varphi_{5}$}
\BIC{$\varphi_{138} \to \varphi_{5}$}
\end{ptree}
\noindent$\Pi_{129}$:
\begin{ptree}
\AXC{$\mathrm{FLim}(f,c,\alpha) \to \forall \varepsilon\,(\mathrm{FLim}(f,c,\alpha))$\lbl{A4}}
\AXC{$\forall \varepsilon\,(\varphi_{131}) \to \varphi_{145}$\lbl{A3}}
\AXC{$\stackrel{\Pi_{89}}{\vdots}$}\noLine
\UIC{$(\forall \varepsilon\,(\varphi_{131}) \to \varphi_{145}) \to \varphi_{180}$}
\BIC{$\varphi_{180}$}
\BIC{$\varphi_{138}$}
\AXC{$\stackrel{\Pi_{128}}{\vdots}$}\noLine
\UIC{$\varphi_{138} \to \varphi_{5}$}
\BIC{$\varphi_{5}$}
\end{ptree}
\noindent 主木:
\begin{ptree}
\AXC{$\stackrel{\Pi_{36}}{\vdots}$}\noLine
\UIC{$\varphi_{177}$}
\AXC{$\stackrel{\Pi_{72}}{\vdots}$}\noLine
\UIC{$\varphi_{2}$}
\AXC{$\stackrel{\Pi_{3}}{\vdots}$}\noLine
\UIC{$\varphi_{189}$}
\AXC{$\varphi_{17}$\lbl{A1}}
\AXC{$\varphi_{17} \to \varphi_{252}$\lbl{A1}}
\BIC{$\varphi_{252}$}
\BIC{$\mathrm{FLimR}(f,c,\alpha) \land \mathrm{FLimL}(f,c,\alpha) \to \alpha \in \mathbb{R}$}
\AXC{$\stackrel{\Pi_{74}}{\vdots}$}\noLine
\UIC{$\varphi_{295}$}
\BIC{$\varphi_{2} \to \varphi_{177} \to \varphi_{4}$}
\BIC{$\varphi_{177} \to \varphi_{4}$}
\BIC{$\varphi_{4}$}
\AXC{$\stackrel{\Pi_{129}}{\vdots}$}\noLine
\UIC{$\varphi_{5}$}
\AXC{$\varphi_{5} \to \varphi_{4} \to \varphi_{6}$\lbl{A1}}
\BIC{$\varphi_{4} \to \varphi_{6}$}
\BIC{$\varphi_{6}$}
\UIC{$\forall \alpha\,(\varphi_{6})$}
\UIC{$\forall c\,\forall \alpha\,(\varphi_{6})$}
\UIC{$\forall f\,\forall c\,\forall \alpha\,(\varphi_{6})$}
\end{ptree}
\begin{abbrevlist}
\item $\varphi_{1} :\equiv \mathrm{FLimR}(f,c,\alpha) \land \mathrm{FLimL}(f,c,\alpha) \to \mathrm{RFn}(f)$
\item $\varphi_{2} :\equiv \mathrm{FLimR}(f,c,\alpha) \land \mathrm{FLimL}(f,c,\alpha) \to \mathrm{PNbhd}(f,c)$
\item $\varphi_{3} :\equiv \mathrm{FLimR}(f,c,\alpha) \land \mathrm{FLimL}(f,c,\alpha) \land \neg \mathrm{PNbhd}(f,c)$
\item $\varphi_{4} :\equiv \mathrm{FLimR}(f,c,\alpha) \land \mathrm{FLimL}(f,c,\alpha) \to \mathrm{FLim}(f,c,\alpha)$
\item $\varphi_{5} :\equiv \mathrm{FLim}(f,c,\alpha) \to \mathrm{FLimR}(f,c,\alpha) \land \mathrm{FLimL}(f,c,\alpha)$
\item $\varphi_{6} :\equiv \mathrm{FLim}(f,c,\alpha) \leftrightarrow \mathrm{FLimR}(f,c,\alpha) \land \mathrm{FLimL}(f,c,\alpha)$
\item $\varphi_{7} :\equiv \forall \alpha\,(\mathrm{FLim}(f,c,\alpha) \to \mathrm{RFn}(f)) \to \mathrm{FLim}(f,c,\alpha) \to \mathrm{RFn}(f)$
\item $\varphi_{8} :\equiv \forall c\,\forall \alpha\,(\mathrm{FLim}(f,c,\alpha) \to \mathrm{RFn}(f)) \to \forall \alpha\,(\mathrm{FLim}(f,c,\alpha) \to \mathrm{RFn}(f))$
\item $\varphi_{9} :\equiv 0 < d_{1} \to 0 < d_{2} \to \exists d\,(d \in \mathbb{R} \land 0 < d \land d \leq d_{1} \land d \leq d_{2})$
\item $\varphi_{10} :\equiv 0 < r_{1} \to 0 < r_{2} \to \exists d\,(d \in \mathbb{R} \land 0 < d \land d \leq r_{1} \land d \leq r_{2})$
\item $\varphi_{11} :\equiv \forall \alpha\,(\mathrm{FLim}(f,c,\alpha) \to c \in \mathbb{R}) \to \mathrm{FLim}(f,c,\alpha) \to c \in \mathbb{R}$
\item $\varphi_{12} :\equiv \forall \alpha\,(\mathrm{FLim}(f,c,\alpha) \to \mathrm{PNbhd}(f,c)) \to \mathrm{FLim}(f,c,\alpha) \to \mathrm{PNbhd}(f,c)$
\item $\varphi_{13} :\equiv \mathrm{FLim}(f,c,\alpha) \land \neg \mathrm{LNbhd}(f,c) \to \forall r\,(\mathrm{FLim}(f,c,\alpha) \land \neg \mathrm{LNbhd}(f,c))$
\item $\varphi_{14} :\equiv \mathrm{FLim}(f,c,\alpha) \land \neg \mathrm{RNbhd}(f,c) \to \forall r\,(\mathrm{FLim}(f,c,\alpha) \land \neg \mathrm{RNbhd}(f,c))$
\item $\varphi_{15} :\equiv \mathrm{FLimR}(f,c,\alpha) \land \mathrm{FLimL}(f,c,\alpha) \to \varepsilon \in \mathbb{R} \to \varepsilon \in \mathbb{R}$
\item $\varphi_{16} :\equiv \forall \alpha\,(\mathrm{FLim}(f,c,\alpha) \to \alpha \in \mathbb{R}) \to \mathrm{FLim}(f,c,\alpha) \to \alpha \in \mathbb{R}$
\item $\varphi_{17} :\equiv \mathrm{FLimR}(f,c,\alpha) \land \mathrm{FLimL}(f,c,\alpha) \to \mathrm{FLimR}(f,c,\alpha) \land \mathrm{FLimL}(f,c,\alpha)$
\item $\varphi_{18} :\equiv \forall f\,\forall c\,\forall \alpha\,(\mathrm{FLim}(f,c,\alpha) \to \mathrm{RFn}(f)) \to \forall c\,\forall \alpha\,(\mathrm{FLim}(f,c,\alpha) \to \mathrm{RFn}(f))$
\item $\varphi_{19} :\equiv \forall c\,\forall \alpha\,(\mathrm{FLim}(f,c,\alpha) \to c \in \mathbb{R}) \to \forall \alpha\,(\mathrm{FLim}(f,c,\alpha) \to c \in \mathbb{R})$
\item $\varphi_{20} :\equiv \forall c\,\forall \alpha\,(\mathrm{FLim}(f,c,\alpha) \to \mathrm{PNbhd}(f,c)) \to \forall \alpha\,(\mathrm{FLim}(f,c,\alpha) \to \mathrm{PNbhd}(f,c))$
\item $\varphi_{21} :\equiv \mathrm{FLimR}(f,c,\alpha) \land \mathrm{FLimL}(f,c,\alpha) \to \forall \varepsilon\,(\mathrm{FLimR}(f,c,\alpha) \land \mathrm{FLimL}(f,c,\alpha))$
\item $\varphi_{22} :\equiv \forall x\,\forall c\,(x \in \mathbb{R} \to c \in \mathbb{R} \to 0 < \lvert x - c \rvert \to x < c \lor c < x)$
\item $\varphi_{23} :\equiv \mathrm{FLimR}(f,c,\alpha) \land \mathrm{FLimL}(f,c,\alpha) \to \varepsilon \in \mathbb{R} \to 0 < \varepsilon \to 0 < \varepsilon$
\item $\varphi_{24} :\equiv \forall c\,\forall \alpha\,(\mathrm{FLim}(f,c,\alpha) \to \alpha \in \mathbb{R}) \to \forall \alpha\,(\mathrm{FLim}(f,c,\alpha) \to \alpha \in \mathbb{R})$
\item $\varphi_{25} :\equiv \forall x \in \mathbb{R}\,(0 < \lvert x - c \rvert \to \lvert x - c \rvert < d \to x \in \mathrm{dom}(f))$
\item $\varphi_{26} :\equiv \forall x \in \mathbb{R}\,(0 < \lvert x - c \rvert \to \lvert x - c \rvert < r \to x \in \mathrm{dom}(f))$
\item $\varphi_{27} :\equiv \delta \in \mathbb{R} \land 0 < \delta \land \delta \leq d_{1} \land \delta \leq d_{2} \to x \in \mathrm{dom}(f) \to x \in \mathrm{dom}(f)$
\item $\varphi_{28} :\equiv x \in \mathrm{dom}(f) \to 0 < \lvert x - c \rvert \to \lvert x - c \rvert < \delta \to c < x \to \delta \leq d_{1}$
\item $\varphi_{29} :\equiv \forall x \in \mathrm{dom}(f)\,(c < x \to x < c + d_{1} \to \lvert f(x) - \alpha \rvert < \varepsilon)$
\item $\varphi_{30} :\equiv x \in \mathrm{dom}(f) \to 0 < \lvert x - c \rvert \to \lvert x - c \rvert < \delta \to x < c \to \delta \leq d_{2}$
\item $\varphi_{31} :\equiv \forall x \in \mathrm{dom}(f)\,(c - d_{2} < x \to x < c \to \lvert f(x) - \alpha \rvert < \varepsilon)$
\item $\varphi_{32} :\equiv \forall x \in \mathrm{dom}(f)\,(c - \delta < x \to x < c \to \lvert f(x) - \alpha \rvert < \varepsilon)$
\item $\varphi_{33} :\equiv \forall x \in \mathrm{dom}(f)\,(c < x \to x < c + \delta \to \lvert f(x) - \alpha \rvert < \varepsilon)$
\item $\varphi_{34} :\equiv x \in \mathrm{dom}(f) \to 0 < \lvert x - c \rvert \to \lvert x - c \rvert < \delta \to c < x \to c < x$
\item $\varphi_{35} :\equiv x \in \mathrm{dom}(f) \to 0 < \lvert x - c \rvert \to \lvert x - c \rvert < \delta \to x < c \to x < c$
\item $\varphi_{36} :\equiv x \in \mathrm{dom}(f) \to 0 < \lvert x - c \rvert \to \lvert x - c \rvert < \delta \to x < c \lor c < x$
\item $\varphi_{37} :\equiv d \in \mathbb{R} \land 0 < d \land d \leq r_{1} \land d \leq r_{2} \to d \in \mathbb{R} \land 0 < d \land d \leq r_{1} \land d \leq r_{2}$
\item $\varphi_{38} :\equiv \forall f\,\forall c\,\forall \alpha\,(\mathrm{FLim}(f,c,\alpha) \to c \in \mathbb{R}) \to \forall c\,\forall \alpha\,(\mathrm{FLim}(f,c,\alpha) \to c \in \mathbb{R})$
\item $\varphi_{39} :\equiv \forall f\,\forall c\,\forall \alpha\,(\mathrm{FLim}(f,c,\alpha) \to \mathrm{PNbhd}(f,c)) \to \forall c\,\forall \alpha\,(\mathrm{FLim}(f,c,\alpha) \to \mathrm{PNbhd}(f,c))$
\item $\varphi_{40} :\equiv x \in \mathrm{dom}(f) \to 0 < \lvert x - c \rvert \to \lvert x - c \rvert < \delta \to \lvert x - c \rvert < \delta$
\item $\varphi_{41} :\equiv \forall f\,\forall c\,\forall \alpha\,(\mathrm{FLim}(f,c,\alpha) \to \alpha \in \mathbb{R}) \to \forall c\,\forall \alpha\,(\mathrm{FLim}(f,c,\alpha) \to \alpha \in \mathbb{R})$
\item $\varphi_{42} :\equiv x \in \mathbb{R} \to 0 < \lvert x - c \rvert \to \lvert x - c \rvert < d \to c < x \to x < c + r_{1}$
\item $\varphi_{43} :\equiv x \in \mathbb{R} \to 0 < \lvert x - c \rvert \to \lvert x - c \rvert < d \to c < x \to x \in \mathrm{dom}(f)$
\item $\varphi_{44} :\equiv x \in \mathbb{R} \to 0 < \lvert x - c \rvert \to \lvert x - c \rvert < d \to x < c \to c - r_{2} < x$
\item $\varphi_{45} :\equiv x \in \mathbb{R} \to 0 < \lvert x - c \rvert \to \lvert x - c \rvert < d \to x < c \to x \in \mathrm{dom}(f)$
\item $\varphi_{46} :\equiv x < c \to x \in \mathbb{R} \to 0 < \lvert x - c \rvert \to \lvert x - c \rvert < r \to x \in \mathrm{dom}(f)$
\item $\varphi_{47} :\equiv r \in \mathbb{R} \land 0 < r \land \forall x \in \mathbb{R}\,(c - r < x \to x < c \to x \in \mathrm{dom}(f))$
\item $\varphi_{48} :\equiv r \in \mathbb{R} \land 0 < r \land \forall x \in \mathbb{R}\,(c < x \to x < c + r \to x \in \mathrm{dom}(f))$
\item $\varphi_{49} :\equiv \forall x\,(\mathrm{RFn}(f) \to x \in \mathrm{dom}(f) \to x \in \mathbb{R}) \to \mathrm{RFn}(f) \to x \in \mathrm{dom}(f) \to x \in \mathbb{R}$
\item $\varphi_{50} :\equiv c < x \to x \in \mathrm{dom}(f) \to c < x \to x < c + d_{1} \to \lvert f(x) - \alpha \rvert < \varepsilon$
\item $\varphi_{51} :\equiv x < c \to x \in \mathrm{dom}(f) \to c - d_{2} < x \to x < c \to \lvert f(x) - \alpha \rvert < \varepsilon$
\item $\varphi_{52} :\equiv x \in \mathrm{dom}(f) \to 0 < \lvert x - c \rvert \to \lvert x - c \rvert < \delta \to c < x \to x < c + d_{1}$
\item $\varphi_{53} :\equiv x \in \mathrm{dom}(f) \to 0 < \lvert x - c \rvert \to \lvert x - c \rvert < \delta \to x < c \to c - d_{2} < x$
\item $\varphi_{54} :\equiv x \in \mathrm{dom}(f) \to c - \delta < x \to x < c \to 0 < \lvert x - c \rvert \land \lvert x - c \rvert < \delta$
\item $\varphi_{55} :\equiv x \in \mathrm{dom}(f) \to c < x \to x < c + \delta \to 0 < \lvert x - c \rvert \land \lvert x - c \rvert < \delta$
\item $\varphi_{56} :\equiv r_{1} \in \mathbb{R} \land 0 < r_{1} \land \forall x \in \mathbb{R}\,(c < x \to x < c + r_{1} \to x \in \mathrm{dom}(f))$
\item $\varphi_{57} :\equiv r_{2} \in \mathbb{R} \land 0 < r_{2} \land \forall x \in \mathbb{R}\,(c - r_{2} < x \to x < c \to x \in \mathrm{dom}(f))$
\item $\varphi_{58} :\equiv \delta \in \mathbb{R} \land 0 < \delta \land \delta \leq d_{1} \land \delta \leq d_{2} \to \delta \in \mathbb{R} \land 0 < \delta \land \delta \leq d_{1} \land \delta \leq d_{2}$
\item $\varphi_{59} :\equiv x \in \mathbb{R} \to 0 < \lvert x - c \rvert \to \lvert x - c \rvert < d \to c < x \to \lvert x - c \rvert < r_{1}$
\item $\varphi_{60} :\equiv \lvert x - c \rvert < d \to c < x \to x \in \mathbb{R} \to c < x \to x < c + r_{1} \to x \in \mathrm{dom}(f)$
\item $\varphi_{61} :\equiv x \in \mathbb{R} \to 0 < \lvert x - c \rvert \to \lvert x - c \rvert < d \to x < c \to \lvert x - c \rvert < r_{2}$
\item $\varphi_{62} :\equiv \lvert x - c \rvert < d \to x < c \to x \in \mathbb{R} \to c - r_{2} < x \to x < c \to x \in \mathrm{dom}(f)$
\item $\varphi_{63} :\equiv x < c + r \to x \in \mathbb{R} \to 0 < \lvert x - c \rvert \to \lvert x - c \rvert < r \to x \in \mathrm{dom}(f)$
\item $\varphi_{64} :\equiv \forall \delta\,(\neg (\delta \in \mathbb{R} \land 0 < \delta \land \delta \leq d_{1} \land \delta \leq d_{2})) \to \neg (d \in \mathbb{R} \land 0 < d \land d \leq d_{1} \land d \leq d_{2})$
\item $\varphi_{65} :\equiv \forall x\,\forall y\,\forall z\,(x < y \to y \leq z \to x < z) \to \forall y\,\forall z\,(\lvert x - c \rvert < y \to y \leq z \to \lvert x - c \rvert < z)$
\item $\varphi_{66} :\equiv \forall f\,\forall x\,(\mathrm{RFn}(f) \to x \in \mathrm{dom}(f) \to x \in \mathbb{R}) \to \forall x\,(\mathrm{RFn}(f) \to x \in \mathrm{dom}(f) \to x \in \mathbb{R})$
\item $\varphi_{67} :\equiv x \in \mathrm{dom}(f) \to 0 < \lvert x - c \rvert \to \lvert x - c \rvert < \delta \to c < x \to \lvert x - c \rvert < d_{1}$
\item $\varphi_{68} :\equiv x \in \mathrm{dom}(f) \to 0 < \lvert x - c \rvert \to \lvert x - c \rvert < \delta \to x < c \to \lvert x - c \rvert < d_{2}$
\item $\varphi_{69} :\equiv x \in \mathrm{dom}(f) \to 0 < \lvert x - c \rvert \to \lvert x - c \rvert < \delta \to \lvert f(x) - \alpha \rvert < \varepsilon$
\item $\varphi_{70} :\equiv \lvert x - c \rvert < d \to c < x \to \forall x \in \mathbb{R}\,(c < x \to x < c + r_{1} \to x \in \mathrm{dom}(f))$
\item $\varphi_{71} :\equiv \lvert x - c \rvert < d \to x < c \to \forall x \in \mathbb{R}\,(c - r_{2} < x \to x < c \to x \in \mathrm{dom}(f))$
\item $\varphi_{72} :\equiv \forall d\,(x \in \mathbb{R} \to c \in \mathbb{R} \to d \in \mathbb{R} \to c < x \to \lvert x - c \rvert < d \to x < c + d)$
\item $\varphi_{73} :\equiv x \in \mathbb{R} \to c \in \mathbb{R} \to d_{1} \in \mathbb{R} \to c < x \to \lvert x - c \rvert < d_{1} \to x < c + d_{1}$
\item $\varphi_{74} :\equiv \forall d\,(x \in \mathbb{R} \to c \in \mathbb{R} \to d \in \mathbb{R} \to x < c \to \lvert x - c \rvert < d \to c - d < x)$
\item $\varphi_{75} :\equiv x \in \mathbb{R} \to c \in \mathbb{R} \to d_{2} \in \mathbb{R} \to x < c \to \lvert x - c \rvert < d_{2} \to c - d_{2} < x$
\item $\varphi_{76} :\equiv x \in \mathbb{R} \to c \in \mathbb{R} \to r_{1} \in \mathbb{R} \to c < x \to \lvert x - c \rvert < r_{1} \to x < c + r_{1}$
\item $\varphi_{77} :\equiv x \in \mathbb{R} \to c \in \mathbb{R} \to r_{2} \in \mathbb{R} \to x < c \to \lvert x - c \rvert < r_{2} \to c - r_{2} < x$
\item $\varphi_{78} :\equiv \forall x \in \mathrm{dom}(f)\,(0 < \lvert x - c \rvert \to \lvert x - c \rvert < \delta \to \lvert f(x) - \alpha \rvert < \varepsilon)$
\item $\varphi_{79} :\equiv \forall p\,\forall q\,(p \in \mathbb{R} \to q \in \mathbb{R} \to 0 < p \to 0 < q \to \exists d\,(d \in \mathbb{R} \land 0 < d \land d \leq p \land d \leq q))$
\item $\varphi_{80} :\equiv \forall q\,(d_{1} \in \mathbb{R} \to q \in \mathbb{R} \to 0 < d_{1} \to 0 < q \to \exists d\,(d \in \mathbb{R} \land 0 < d \land d \leq d_{1} \land d \leq q))$
\item $\varphi_{81} :\equiv \forall q\,(r_{1} \in \mathbb{R} \to q \in \mathbb{R} \to 0 < r_{1} \to 0 < q \to \exists d\,(d \in \mathbb{R} \land 0 < d \land d \leq r_{1} \land d \leq q))$
\item $\varphi_{82} :\equiv d \in \mathbb{R} \land 0 < d \land d \leq r_{1} \land d \leq r_{2} \to x \in \mathbb{R} \to 0 < \lvert x - c \rvert \to 0 < \lvert x - c \rvert$
\item $\varphi_{83} :\equiv c \in \mathbb{R} \to r \in \mathbb{R} \to c - r < x \to x < c \to 0 < \lvert x - c \rvert \land \lvert x - c \rvert < r$
\item $\varphi_{84} :\equiv c \in \mathbb{R} \to r \in \mathbb{R} \to c < x \to x < c + r \to 0 < \lvert x - c \rvert \land \lvert x - c \rvert < r$
\item $\varphi_{85} :\equiv \forall d\,\forall \delta\,(\neg (\delta \in \mathbb{R} \land 0 < \delta \land \delta \leq d_{1} \land \delta \leq d_{2})) \to \forall d\,(\neg (d \in \mathbb{R} \land 0 < d \land d \leq d_{1} \land d \leq d_{2}))$
\item $\varphi_{86} :\equiv c \in \mathbb{R} \to \delta \in \mathbb{R} \to c - \delta < x \to x < c \to 0 < \lvert x - c \rvert \land \lvert x - c \rvert < \delta$
\item $\varphi_{87} :\equiv c \in \mathbb{R} \to \delta \in \mathbb{R} \to c < x \to x < c + \delta \to 0 < \lvert x - c \rvert \land \lvert x - c \rvert < \delta$
\item $\varphi_{88} :\equiv x \in \mathrm{dom}(f) \to 0 < \lvert x - c \rvert \to \lvert x - c \rvert < \delta \to c < x \to \lvert f(x) - \alpha \rvert < \varepsilon$
\item $\varphi_{89} :\equiv x \in \mathrm{dom}(f) \to 0 < \lvert x - c \rvert \to \lvert x - c \rvert < \delta \to x < c \to \lvert f(x) - \alpha \rvert < \varepsilon$
\item $\varphi_{90} :\equiv \forall \delta\,(\neg (\delta \in \mathbb{R} \land 0 < \delta \land \delta \leq d_{1} \land \delta \leq d_{2})) \to \forall d\,\forall \delta\,(\neg (\delta \in \mathbb{R} \land 0 < \delta \land \delta \leq d_{1} \land \delta \leq d_{2}))$
\item $\varphi_{91} :\equiv \forall z\,(\lvert x - c \rvert < d \to d \leq z \to \lvert x - c \rvert < z) \to \lvert x - c \rvert < d \to d \leq r_{1} \to \lvert x - c \rvert < r_{1}$
\item $\varphi_{92} :\equiv \forall z\,(\lvert x - c \rvert < d \to d \leq z \to \lvert x - c \rvert < z) \to \lvert x - c \rvert < d \to d \leq r_{2} \to \lvert x - c \rvert < r_{2}$
\item $\varphi_{93} :\equiv \delta \in \mathbb{R} \land 0 < \delta \land \delta \leq d_{1} \land \delta \leq d_{2} \to x \in \mathrm{dom}(f) \to 0 < \lvert x - c \rvert \to 0 < \lvert x - c \rvert$
\item $\varphi_{94} :\equiv \forall z\,(\lvert x - c \rvert < \delta \to \delta \leq z \to \lvert x - c \rvert < z) \to \lvert x - c \rvert < \delta \to \delta \leq d_{1} \to \lvert x - c \rvert < d_{1}$
\item $\varphi_{95} :\equiv \forall z\,(\lvert x - c \rvert < \delta \to \delta \leq z \to \lvert x - c \rvert < z) \to \lvert x - c \rvert < \delta \to \delta \leq d_{2} \to \lvert x - c \rvert < d_{2}$
\item $\varphi_{96} :\equiv \forall y\,\forall z\,(\lvert x - c \rvert < y \to y \leq z \to \lvert x - c \rvert < z) \to \forall z\,(\lvert x - c \rvert < d \to d \leq z \to \lvert x - c \rvert < z)$
\item $\varphi_{97} :\equiv \forall y\,\forall z\,(\lvert x - c \rvert < y \to y \leq z \to \lvert x - c \rvert < z) \to \forall z\,(\lvert x - c \rvert < \delta \to \delta \leq z \to \lvert x - c \rvert < z)$
\item $\varphi_{98} :\equiv d \in \mathbb{R} \land 0 < d \land d \leq r_{1} \land d \leq r_{2} \to x \in \mathbb{R} \to 0 < \lvert x - c \rvert \to \lvert x - c \rvert < d \to x \in \mathrm{dom}(f)$
\item $\varphi_{99} :\equiv \forall d\,(x \in \mathbb{R} \to c \in \mathbb{R} \to d \in \mathbb{R} \to c - d < x \to x < c \to 0 < \lvert x - c \rvert \land \lvert x - c \rvert < d)$
\item $\varphi_{100} :\equiv \forall d\,(x \in \mathbb{R} \to c \in \mathbb{R} \to d \in \mathbb{R} \to c < x \to x < c + d \to 0 < \lvert x - c \rvert \land \lvert x - c \rvert < d)$
\item $\varphi_{101} :\equiv d \in \mathbb{R} \land 0 < d \land d \leq r_{1} \land d \leq r_{2} \to x \in \mathbb{R} \to 0 < \lvert x - c \rvert \to \lvert x - c \rvert < d \to c < x \to d \leq r_{1}$
\item $\varphi_{102} :\equiv d \in \mathbb{R} \land 0 < d \land d \leq r_{1} \land d \leq r_{2} \to x \in \mathbb{R} \to 0 < \lvert x - c \rvert \to \lvert x - c \rvert < d \to x < c \to d \leq r_{2}$
\item $\varphi_{103} :\equiv d \in \mathbb{R} \land 0 < d \land d \leq r_{1} \land d \leq r_{2} \to x \in \mathbb{R} \to 0 < \lvert x - c \rvert \to \lvert x - c \rvert < d \to c < x \to c < x$
\item $\varphi_{104} :\equiv d \in \mathbb{R} \land 0 < d \land d \leq r_{1} \land d \leq r_{2} \to x \in \mathbb{R} \to 0 < \lvert x - c \rvert \to \lvert x - c \rvert < d \to x < c \to x < c$
\item $\varphi_{105} :\equiv d \in \mathbb{R} \land 0 < d \land d \leq r_{1} \land d \leq r_{2} \to x \in \mathbb{R} \to 0 < \lvert x - c \rvert \to \lvert x - c \rvert < d \to x < c \lor c < x$
\item $\varphi_{106} :\equiv d \in \mathbb{R} \land 0 < d \land d \leq r_{1} \land d \leq r_{2} \to x \in \mathbb{R} \to 0 < \lvert x - c \rvert \to \lvert x - c \rvert < d \to \lvert x - c \rvert < d$
\item $\varphi_{107} :\equiv \forall x \in \mathbb{R}\,(c < x \to x < c + r_{1} \to x \in \mathrm{dom}(f)) \to x \in \mathbb{R} \to c < x \to x < c + r_{1} \to x \in \mathrm{dom}(f)$
\item $\varphi_{108} :\equiv \forall x \in \mathbb{R}\,(c - r_{2} < x \to x < c \to x \in \mathrm{dom}(f)) \to x \in \mathbb{R} \to c - r_{2} < x \to x < c \to x \in \mathrm{dom}(f)$
\item $\varphi_{109} :\equiv \delta \in \mathbb{R} \land 0 < \delta \land \delta \leq d_{1} \land \delta \leq d_{2} \to x \in \mathrm{dom}(f) \to 0 < \lvert x - c \rvert \to \lvert x - c \rvert < \delta \to x \in \mathbb{R}$
\item $\varphi_{110} :\equiv \mathrm{FLimR}(f,c,\alpha) \land \mathrm{FLimL}(f,c,\alpha) \to \exists r\,(\varphi_{47})$
\item $\varphi_{111} :\equiv \mathrm{FLimR}(f,c,\alpha) \land \mathrm{FLimL}(f,c,\alpha) \to \exists r\,(\varphi_{48})$
\item $\varphi_{112} :\equiv \forall f\,\forall c\,(\mathrm{PNbhd}(f,c) \leftrightarrow \exists r\,(r \in \mathbb{R} \land 0 < r \land \varphi_{26}))$
\item $\varphi_{113} :\equiv \mathrm{FLim}(f,c,\alpha) \to r \in \mathbb{R} \land 0 < r \land \varphi_{26} \to \mathrm{LNbhd}(f,c)$
\item $\varphi_{114} :\equiv \mathrm{FLim}(f,c,\alpha) \to r \in \mathbb{R} \land 0 < r \land \varphi_{26} \to \mathrm{RNbhd}(f,c)$
\item $\varphi_{115} :\equiv \mathrm{FLim}(f,c,\alpha) \land \neg \mathrm{LNbhd}(f,c) \to \neg (r \in \mathbb{R} \land 0 < r \land \varphi_{26})$
\item $\varphi_{116} :\equiv \mathrm{FLim}(f,c,\alpha) \land \neg \mathrm{RNbhd}(f,c) \to \neg (r \in \mathbb{R} \land 0 < r \land \varphi_{26})$
\item $\varphi_{117} :\equiv \mathrm{FLimR}(f,c,\alpha) \land \mathrm{FLimL}(f,c,\alpha) \to \varphi_{56} \to \mathrm{PNbhd}(f,c)$
\item $\varphi_{118} :\equiv \mathrm{FLimR}(f,c,\alpha) \land \mathrm{FLimL}(f,c,\alpha) \land \varphi_{56} \land \neg \mathrm{PNbhd}(f,c)$
\item $\varphi_{119} :\equiv \forall r\,(\mathrm{FLim}(f,c,\alpha) \land \neg \mathrm{LNbhd}(f,c)) \to \forall r\,(\neg (r \in \mathbb{R} \land 0 < r \land \varphi_{26}))$
\item $\varphi_{120} :\equiv \forall r\,(\mathrm{FLim}(f,c,\alpha) \land \neg \mathrm{RNbhd}(f,c)) \to \forall r\,(\neg (r \in \mathbb{R} \land 0 < r \land \varphi_{26}))$
\item $\varphi_{121} :\equiv \varphi_{57} \to d \in \mathbb{R} \land 0 < d \land d \leq r_{1} \land d \leq r_{2} \to \mathrm{PNbhd}(f,c)$
\item $\varphi_{122} :\equiv \forall c\,(x \in \mathbb{R} \to c \in \mathbb{R} \to 0 < \lvert x - c \rvert \to x < c \lor c < x) \to x \in \mathbb{R} \to c \in \mathbb{R} \to 0 < \lvert x - c \rvert \to x < c \lor c < x$
\item $\varphi_{123} :\equiv d \in \mathbb{R} \land 0 < d \land d \leq r_{1} \land d \leq r_{2} \to \exists r\,(r \in \mathbb{R} \land 0 < r \land \varphi_{26})$
\item $\varphi_{124} :\equiv \mathrm{FLim}(f,c,\alpha) \to r \in \mathbb{R} \land 0 < r \land \varphi_{26} \to x \in \mathbb{R} \to x \in \mathbb{R}$
\item $\varphi_{125} :\equiv d_{2} \in \mathbb{R} \land 0 < d_{2} \land \varphi_{31} \to \exists d\,(d \in \mathbb{R} \land 0 < d \land d \leq d_{1} \land d \leq d_{2})$
\item $\varphi_{126} :\equiv \varphi_{22} \to \forall c\,(x \in \mathbb{R} \to c \in \mathbb{R} \to 0 < \lvert x - c \rvert \to x < c \lor c < x)$
\item $\varphi_{127} :\equiv \varphi_{26} \to x \in \mathbb{R} \to 0 < \lvert x - c \rvert \to \lvert x - c \rvert < r \to x \in \mathrm{dom}(f)$
\item $\varphi_{128} :\equiv \mathrm{FLim}(f,c,\alpha) \to r \in \mathbb{R} \land 0 < r \land \varphi_{26} \to x \in \mathbb{R} \to c < x \to c < x$
\item $\varphi_{129} :\equiv \varphi_{29} \to x \in \mathrm{dom}(f) \to c < x \to x < c + d_{1} \to \lvert f(x) - \alpha \rvert < \varepsilon$
\item $\varphi_{130} :\equiv \varphi_{31} \to x \in \mathrm{dom}(f) \to c - d_{2} < x \to x < c \to \lvert f(x) - \alpha \rvert < \varepsilon$
\item $\varphi_{131} :\equiv \mathrm{FLim}(f,c,\alpha) \to \varepsilon \in \mathbb{R} \to 0 < \varepsilon \to \exists \delta\,(\delta \in \mathbb{R} \land 0 < \delta \land \varphi_{32})$
\item $\varphi_{132} :\equiv \mathrm{FLim}(f,c,\alpha) \to \varepsilon \in \mathbb{R} \to 0 < \varepsilon \to \exists \delta\,(\delta \in \mathbb{R} \land 0 < \delta \land \varphi_{33})$
\item $\varphi_{133} :\equiv \mathrm{FLim}(f,c,\alpha) \land \varepsilon \in \mathbb{R} \land 0 < \varepsilon \land \neg \exists \delta\,(\delta \in \mathbb{R} \land 0 < \delta \land \varphi_{32})$
\item $\varphi_{134} :\equiv \mathrm{FLim}(f,c,\alpha) \land \varepsilon \in \mathbb{R} \land 0 < \varepsilon \land \neg \exists \delta\,(\delta \in \mathbb{R} \land 0 < \delta \land \varphi_{33})$
\item $\varphi_{135} :\equiv r \in \mathbb{R} \land 0 < r \land \varphi_{26} \to x \in \mathbb{R} \to c < x \to x < c + r \to x < c + r$
\item $\varphi_{136} :\equiv r \in \mathbb{R} \land 0 < r \land \varphi_{26} \to x \in \mathbb{R} \to c - r < x \to x < c \to x \in \mathrm{dom}(f)$
\item $\varphi_{137} :\equiv r \in \mathbb{R} \land 0 < r \land \varphi_{26} \to x \in \mathbb{R} \to c < x \to x < c + r \to x \in \mathrm{dom}(f)$
\item $\varphi_{138} :\equiv \mathrm{FLim}(f,c,\alpha) \to \forall \varepsilon \in \mathbb{R}\,(0 < \varepsilon \to \exists \delta\,(\delta \in \mathbb{R} \land 0 < \delta \land \varphi_{32}))$
\item $\varphi_{139} :\equiv \mathrm{FLim}(f,c,\alpha) \to \forall \varepsilon \in \mathbb{R}\,(0 < \varepsilon \to \exists \delta\,(\delta \in \mathbb{R} \land 0 < \delta \land \varphi_{33}))$
\item $\varphi_{140} :\equiv r \in \mathbb{R} \land 0 < r \land \varphi_{26} \to \forall x \in \mathbb{R}\,(c - r < x \to x < c \to x \in \mathrm{dom}(f))$
\item $\varphi_{141} :\equiv r \in \mathbb{R} \land 0 < r \land \varphi_{26} \to \forall x \in \mathbb{R}\,(c < x \to x < c + r \to x \in \mathrm{dom}(f))$
\item $\varphi_{142} :\equiv \mathrm{FLim}(f,c,\alpha) \to r \in \mathbb{R} \land 0 < r \land \varphi_{26} \to x \in \mathbb{R} \to c - r < x \to c - r < x$
\item $\varphi_{143} :\equiv d \in \mathbb{R} \land 0 < d \land d \leq r_{1} \land d \leq r_{2} \to x \in \mathbb{R} \to 0 < \lvert x - c \rvert \to \varphi_{60}$
\item $\varphi_{144} :\equiv d \in \mathbb{R} \land 0 < d \land d \leq r_{1} \land d \leq r_{2} \to x \in \mathbb{R} \to 0 < \lvert x - c \rvert \to \varphi_{62}$
\item $\varphi_{145} :\equiv \forall \varepsilon\,(\mathrm{FLim}(f,c,\alpha)) \to \forall \varepsilon \in \mathbb{R}\,(0 < \varepsilon \to \exists \delta\,(\delta \in \mathbb{R} \land 0 < \delta \land \varphi_{32}))$
\item $\varphi_{146} :\equiv \forall \varepsilon\,(\mathrm{FLim}(f,c,\alpha)) \to \forall \varepsilon \in \mathbb{R}\,(0 < \varepsilon \to \exists \delta\,(\delta \in \mathbb{R} \land 0 < \delta \land \varphi_{33}))$
\item $\varphi_{147} :\equiv \mathrm{FLim}(f,c,\alpha) \to r \in \mathbb{R} \land 0 < r \land \varphi_{26} \to x \in \mathbb{R} \to c - r < x \to x < c \to x < c$
\item $\varphi_{148} :\equiv d \in \mathbb{R} \land 0 < d \land d \leq r_{1} \land d \leq r_{2} \to x \in \mathbb{R} \to 0 < \lvert x - c \rvert \to \varphi_{70}$
\item $\varphi_{149} :\equiv d \in \mathbb{R} \land 0 < d \land d \leq r_{1} \land d \leq r_{2} \to x \in \mathbb{R} \to 0 < \lvert x - c \rvert \to \varphi_{71}$
\item $\varphi_{150} :\equiv \mathrm{FLim}(f,c,\alpha) \land \varepsilon \in \mathbb{R} \land 0 < \varepsilon \land (\delta \in \mathbb{R} \land 0 < \delta \land \varphi_{78})$
\item $\varphi_{151} :\equiv 0 < \varepsilon \to \delta \in \mathbb{R} \land 0 < \delta \land \varphi_{78} \to x \in \mathrm{dom}(f) \to x \in \mathrm{dom}(f)$
\item $\varphi_{152} :\equiv \delta \in \mathbb{R} \land 0 < \delta \land \delta \leq d_{1} \land \delta \leq d_{2} \to \exists \delta\,(\delta \in \mathbb{R} \land 0 < \delta \land \varphi_{78})$
\item $\varphi_{153} :\equiv 0 < \varepsilon \to \delta \in \mathbb{R} \land 0 < \delta \land \varphi_{78} \to x \in \mathrm{dom}(f) \to c < x \to c < x$
\item $\varphi_{154} :\equiv \mathrm{FLim}(f,c,\alpha) \to \varepsilon \in \mathbb{R} \to 0 < \varepsilon \to \exists \delta\,(\delta \in \mathbb{R} \land 0 < \delta \land \varphi_{78})$
\item $\varphi_{155} :\equiv \delta \in \mathbb{R} \land 0 < \delta \land \varphi_{78} \to x \in \mathrm{dom}(f) \to c - \delta < x \to c - \delta < x$
\item $\varphi_{156} :\equiv \mathrm{FLimR}(f,c,\alpha) \land \mathrm{FLimL}(f,c,\alpha) \to \varepsilon \in \mathbb{R} \to 0 < \varepsilon \to \exists \delta\,(\delta \in \mathbb{R} \land 0 < \delta \land \varphi_{32})$
\item $\varphi_{157} :\equiv \mathrm{FLimR}(f,c,\alpha) \land \mathrm{FLimL}(f,c,\alpha) \to \varepsilon \in \mathbb{R} \to 0 < \varepsilon \to \exists \delta\,(\delta \in \mathbb{R} \land 0 < \delta \land \varphi_{33})$
\item $\varphi_{158} :\equiv \mathrm{FLim}(f,c,\alpha) \land (r \in \mathbb{R} \land 0 < r \land \varphi_{26}) \to x \in \mathbb{R} \to c - r < x \to x < c \to x \in \mathrm{dom}(f)$
\item $\varphi_{159} :\equiv \mathrm{FLim}(f,c,\alpha) \land (r \in \mathbb{R} \land 0 < r \land \varphi_{26}) \to x \in \mathbb{R} \to c < x \to x < c + r \to x \in \mathrm{dom}(f)$
\item $\varphi_{160} :\equiv \delta \in \mathbb{R} \land 0 < \delta \land \varphi_{78} \to x \in \mathrm{dom}(f) \to c - \delta < x \to x < c \to x < c$
\item $\varphi_{161} :\equiv \mathrm{FLimR}(f,c,\alpha) \land \mathrm{FLimL}(f,c,\alpha) \to \forall \varepsilon \in \mathbb{R}\,(0 < \varepsilon \to \exists \delta\,(\delta \in \mathbb{R} \land 0 < \delta \land \varphi_{32}))$
\item $\varphi_{162} :\equiv \mathrm{FLimR}(f,c,\alpha) \land \mathrm{FLimL}(f,c,\alpha) \to \forall \varepsilon \in \mathbb{R}\,(0 < \varepsilon \to \exists \delta\,(\delta \in \mathbb{R} \land 0 < \delta \land \varphi_{33}))$
\item $\varphi_{163} :\equiv \delta \in \mathbb{R} \land 0 < \delta \land \varphi_{78} \to x \in \mathrm{dom}(f) \to c - \delta < x \to x < c \to x \in \mathbb{R}$
\item $\varphi_{164} :\equiv \delta \in \mathbb{R} \land 0 < \delta \land \varphi_{78} \to x \in \mathrm{dom}(f) \to c < x \to x < c + \delta \to x \in \mathbb{R}$
\item $\varphi_{165} :\equiv \forall x\,(\mathrm{FLim}(f,c,\alpha) \land (r \in \mathbb{R} \land 0 < r \land \varphi_{26})) \to \forall x \in \mathbb{R}\,(c - r < x \to x < c \to x \in \mathrm{dom}(f))$
\item $\varphi_{166} :\equiv \forall x\,(\mathrm{FLim}(f,c,\alpha) \land (r \in \mathbb{R} \land 0 < r \land \varphi_{26})) \to \forall x \in \mathbb{R}\,(c < x \to x < c + r \to x \in \mathrm{dom}(f))$
\item $\varphi_{167} :\equiv \delta \in \mathbb{R} \land 0 < \delta \land \varphi_{78} \to x \in \mathrm{dom}(f) \to c < x \to x < c + \delta \to x < c + \delta$
\item $\varphi_{168} :\equiv \mathrm{FLim}(f,c,\alpha) \to \varepsilon \in \mathbb{R} \to 0 < \varepsilon \to \delta \in \mathbb{R} \land 0 < \delta \land \varphi_{78} \to \delta \in \mathbb{R}$
\item $\varphi_{169} :\equiv r \in \mathbb{R} \land 0 < r \land \varphi_{26} \to x \in \mathbb{R} \to c - r < x \to x < c \to 0 < \lvert x - c \rvert \land \lvert x - c \rvert < r$
\item $\varphi_{170} :\equiv r \in \mathbb{R} \land 0 < r \land \varphi_{26} \to x \in \mathbb{R} \to c < x \to x < c + r \to 0 < \lvert x - c \rvert \land \lvert x - c \rvert < r$
\item $\varphi_{171} :\equiv \delta \in \mathbb{R} \land 0 < \delta \land \delta \leq d_{1} \land \delta \leq d_{2} \to x \in \mathrm{dom}(f) \to 0 < \lvert x - c \rvert \to \lvert x - c \rvert < \delta \to \varphi_{50}$
\item $\varphi_{172} :\equiv \delta \in \mathbb{R} \land 0 < \delta \land \delta \leq d_{1} \land \delta \leq d_{2} \to x \in \mathrm{dom}(f) \to 0 < \lvert x - c \rvert \to \lvert x - c \rvert < \delta \to \varphi_{51}$
\item $\varphi_{173} :\equiv \mathrm{FLimR}(f,c,\alpha) \land \mathrm{FLimL}(f,c,\alpha) \to \varepsilon \in \mathbb{R} \to 0 < \varepsilon \to \exists \delta\,(\delta \in \mathbb{R} \land 0 < \delta \land \varphi_{78})$
\item $\varphi_{174} :\equiv \mathrm{FLimR}(f,c,\alpha) \land \mathrm{FLimL}(f,c,\alpha) \land \varepsilon \in \mathbb{R} \land 0 < \varepsilon \land \neg \exists \delta\,(\delta \in \mathbb{R} \land 0 < \delta \land \varphi_{78})$
\item $\varphi_{175} :\equiv \delta \in \mathbb{R} \land 0 < \delta \land \delta \leq d_{1} \land \delta \leq d_{2} \to x \in \mathrm{dom}(f) \to 0 < \lvert x - c \rvert \to \lvert x - c \rvert < \delta \to c < x \to \varphi_{29}$
\item $\varphi_{176} :\equiv \delta \in \mathbb{R} \land 0 < \delta \land \delta \leq d_{1} \land \delta \leq d_{2} \to x \in \mathrm{dom}(f) \to 0 < \lvert x - c \rvert \to \lvert x - c \rvert < \delta \to x < c \to \varphi_{31}$
\item $\varphi_{177} :\equiv \mathrm{FLimR}(f,c,\alpha) \land \mathrm{FLimL}(f,c,\alpha) \to \forall \varepsilon \in \mathbb{R}\,(0 < \varepsilon \to \exists \delta\,(\delta \in \mathbb{R} \land 0 < \delta \land \varphi_{78}))$
\item $\varphi_{178} :\equiv \delta \in \mathbb{R} \land 0 < \delta \land \varphi_{78} \to x \in \mathrm{dom}(f) \to c - \delta < x \to x < c \to \lvert f(x) - \alpha \rvert < \varepsilon$
\item $\varphi_{179} :\equiv \delta \in \mathbb{R} \land 0 < \delta \land \varphi_{78} \to x \in \mathrm{dom}(f) \to c < x \to x < c + \delta \to \lvert f(x) - \alpha \rvert < \varepsilon$
\item $\varphi_{180} :\equiv (\mathrm{FLim}(f,c,\alpha) \to \forall \varepsilon\,(\mathrm{FLim}(f,c,\alpha))) \to \varphi_{138}$
\item $\varphi_{181} :\equiv (\mathrm{FLim}(f,c,\alpha) \to \forall \varepsilon\,(\mathrm{FLim}(f,c,\alpha))) \to \varphi_{139}$
\item $\varphi_{182} :\equiv \forall \varepsilon\,(\mathrm{FLimR}(f,c,\alpha) \land \mathrm{FLimL}(f,c,\alpha)) \to \forall \varepsilon \in \mathbb{R}\,(0 < \varepsilon \to \exists \delta\,(\delta \in \mathbb{R} \land 0 < \delta \land \varphi_{78}))$
\item $\varphi_{183} :\equiv \forall r\,(\neg (r \in \mathbb{R} \land 0 < r \land \varphi_{26})) \to \neg (d \in \mathbb{R} \land 0 < d \land \varphi_{25})$
\item $\varphi_{184} :\equiv \forall c\,(\mathrm{LNbhd}(f,c) \leftrightarrow \exists r\,(\varphi_{47})) \to (\mathrm{LNbhd}(f,c) \leftrightarrow \exists r\,(\varphi_{47}))$
\item $\varphi_{185} :\equiv \forall c\,(\mathrm{RNbhd}(f,c) \leftrightarrow \exists r\,(\varphi_{48})) \to (\mathrm{RNbhd}(f,c) \leftrightarrow \exists r\,(\varphi_{48}))$
\item $\varphi_{186} :\equiv \mathrm{FLimR}(f,c,\alpha) \land \mathrm{FLimL}(f,c,\alpha) \to \varphi_{56} \to \exists r\,(\varphi_{47})$
\item $\varphi_{187} :\equiv \forall f\,\forall c\,(\mathrm{LNbhd}(f,c) \leftrightarrow \exists r\,(\varphi_{47})) \to \forall c\,(\mathrm{LNbhd}(f,c) \leftrightarrow \exists r\,(\varphi_{47}))$
\item $\varphi_{188} :\equiv \forall f\,\forall c\,(\mathrm{RNbhd}(f,c) \leftrightarrow \exists r\,(\varphi_{48})) \to \forall c\,(\mathrm{RNbhd}(f,c) \leftrightarrow \exists r\,(\varphi_{48}))$
\item $\varphi_{189} :\equiv \mathrm{FLimR}(f,c,\alpha) \leftrightarrow \mathrm{RFn}(f) \land c \in \mathbb{R} \land \alpha \in \mathbb{R} \land \mathrm{RNbhd}(f,c) \land \forall \varepsilon \in \mathbb{R}\,(0 < \varepsilon \to \exists \delta\,(\delta \in \mathbb{R} \land 0 < \delta \land \varphi_{33}))$
\item $\varphi_{190} :\equiv \mathrm{FLimL}(f,c,\alpha) \leftrightarrow \mathrm{RFn}(f) \land c \in \mathbb{R} \land \alpha \in \mathbb{R} \land \mathrm{LNbhd}(f,c) \land \forall \varepsilon \in \mathbb{R}\,(0 < \varepsilon \to \exists \delta\,(\delta \in \mathbb{R} \land 0 < \delta \land \varphi_{32}))$
\item $\varphi_{191} :\equiv \forall d_{1}\,(\neg (d_{1} \in \mathbb{R} \land 0 < d_{1} \land \varphi_{29})) \to \neg (\delta \in \mathbb{R} \land 0 < \delta \land \varphi_{33})$
\item $\varphi_{192} :\equiv \forall d_{2}\,(\neg (d_{2} \in \mathbb{R} \land 0 < d_{2} \land \varphi_{31})) \to \neg (\delta \in \mathbb{R} \land 0 < \delta \land \varphi_{32})$
\item $\varphi_{193} :\equiv \forall \delta\,(\neg (\delta \in \mathbb{R} \land 0 < \delta \land \varphi_{32})) \to \neg (\delta \in \mathbb{R} \land 0 < \delta \land \varphi_{32})$
\item $\varphi_{194} :\equiv \forall \delta\,(\neg (\delta \in \mathbb{R} \land 0 < \delta \land \varphi_{33})) \to \neg (\delta \in \mathbb{R} \land 0 < \delta \land \varphi_{33})$
\item $\varphi_{195} :\equiv \mathrm{FLim}(f,c,\alpha) \to r \in \mathbb{R} \land 0 < r \land \varphi_{26} \to r \in \mathbb{R} \land 0 < r \land \varphi_{26}$
\item $\varphi_{196} :\equiv \mathrm{FLimR}(f,c,\alpha) \land \mathrm{FLimL}(f,c,\alpha) \to \varphi_{56} \to \varphi_{57} \to 0 < r_{2}$
\item $\varphi_{197} :\equiv \mathrm{FLimR}(f,c,\alpha) \land \mathrm{FLimL}(f,c,\alpha) \to \varphi_{56} \to \varphi_{57} \to 0 < r_{1}$
\item $\varphi_{198} :\equiv \mathrm{FLimR}(f,c,\alpha) \land \mathrm{FLimL}(f,c,\alpha) \to \varphi_{56} \to \varphi_{57} \to \mathrm{PNbhd}(f,c)$
\item $\varphi_{199} :\equiv \forall d_{1}\,(\neg (d_{1} \in \mathbb{R} \land 0 < d_{1} \land \varphi_{29})) \to \forall \delta\,\forall d_{1}\,(\neg (d_{1} \in \mathbb{R} \land 0 < d_{1} \land \varphi_{29}))$
\item $\varphi_{200} :\equiv \forall d_{2}\,(\neg (d_{2} \in \mathbb{R} \land 0 < d_{2} \land \varphi_{31})) \to \forall \delta\,\forall d_{2}\,(\neg (d_{2} \in \mathbb{R} \land 0 < d_{2} \land \varphi_{31}))$
\item $\varphi_{201} :\equiv \mathrm{FLimR}(f,c,\alpha) \land \mathrm{FLimL}(f,c,\alpha) \to \varphi_{56} \to \varphi_{57} \to r_{1} \in \mathbb{R}$
\item $\varphi_{202} :\equiv \mathrm{FLimR}(f,c,\alpha) \land \mathrm{FLimL}(f,c,\alpha) \to \varphi_{56} \to \varphi_{57} \to r_{2} \in \mathbb{R}$
\item $\varphi_{203} :\equiv \forall \delta\,\forall d_{1}\,(\neg (d_{1} \in \mathbb{R} \land 0 < d_{1} \land \varphi_{29})) \to \forall \delta\,(\neg (\delta \in \mathbb{R} \land 0 < \delta \land \varphi_{33}))$
\item $\varphi_{204} :\equiv \forall \delta\,\forall d_{2}\,(\neg (d_{2} \in \mathbb{R} \land 0 < d_{2} \land \varphi_{31})) \to \forall \delta\,(\neg (\delta \in \mathbb{R} \land 0 < \delta \land \varphi_{32}))$
\item $\varphi_{205} :\equiv d_{1} \in \mathbb{R} \land 0 < d_{1} \land \varphi_{29} \to d_{2} \in \mathbb{R} \land 0 < d_{2} \land \varphi_{31} \to d_{1} \in \mathbb{R}$
\item $\varphi_{206} :\equiv d_{1} \in \mathbb{R} \land 0 < d_{1} \land \varphi_{29} \to d_{2} \in \mathbb{R} \land 0 < d_{2} \land \varphi_{31} \to d_{2} \in \mathbb{R}$
\item $\varphi_{207} :\equiv \mathrm{FLimR}(f,c,\alpha) \land \mathrm{FLimL}(f,c,\alpha) \land \varphi_{56} \land \varphi_{57} \land \neg \mathrm{PNbhd}(f,c)$
\item $\varphi_{208} :\equiv 0 < \varepsilon \to d_{1} \in \mathbb{R} \land 0 < d_{1} \land \varphi_{29} \to \exists \delta\,(\delta \in \mathbb{R} \land 0 < \delta \land \varphi_{32})$
\item $\varphi_{209} :\equiv d_{2} \in \mathbb{R} \land 0 < d_{2} \land \varphi_{31} \to \delta \in \mathbb{R} \land 0 < \delta \land \delta \leq d_{1} \land \delta \leq d_{2} \to \varphi_{28}$
\item $\varphi_{210} :\equiv d_{2} \in \mathbb{R} \land 0 < d_{2} \land \varphi_{31} \to \delta \in \mathbb{R} \land 0 < \delta \land \delta \leq d_{1} \land \delta \leq d_{2} \to \varphi_{30}$
\item $\varphi_{211} :\equiv \varphi_{57} \to d \in \mathbb{R} \land 0 < d \land d \leq r_{1} \land d \leq r_{2} \to d \in \mathbb{R} \land 0 < d \land \varphi_{25}$
\item $\varphi_{212} :\equiv d_{2} \in \mathbb{R} \land 0 < d_{2} \land \varphi_{31} \to \delta \in \mathbb{R} \land 0 < \delta \land \delta \leq d_{1} \land \delta \leq d_{2} \to \varphi_{34}$
\item $\varphi_{213} :\equiv d_{2} \in \mathbb{R} \land 0 < d_{2} \land \varphi_{31} \to \delta \in \mathbb{R} \land 0 < \delta \land \delta \leq d_{1} \land \delta \leq d_{2} \to \varphi_{35}$
\item $\varphi_{214} :\equiv d_{2} \in \mathbb{R} \land 0 < d_{2} \land \varphi_{31} \to \delta \in \mathbb{R} \land 0 < \delta \land \delta \leq d_{1} \land \delta \leq d_{2} \to \varphi_{36}$
\item $\varphi_{215} :\equiv \mathrm{FLim}(f,c,\alpha) \to r \in \mathbb{R} \land 0 < r \land \varphi_{26} \to x \in \mathbb{R} \to c - r < x \to \varphi_{46}$
\item $\varphi_{216} :\equiv \mathrm{FLim}(f,c,\alpha) \to r \in \mathbb{R} \land 0 < r \land \varphi_{26} \to x \in \mathbb{R} \to c < x \to \varphi_{63}$
\item $\varphi_{217} :\equiv \mathrm{FLim}(f,c,\alpha) \leftrightarrow \mathrm{RFn}(f) \land c \in \mathbb{R} \land \alpha \in \mathbb{R} \land \mathrm{PNbhd}(f,c) \land \forall \varepsilon \in \mathbb{R}\,(0 < \varepsilon \to \exists \delta\,(\delta \in \mathbb{R} \land 0 < \delta \land \varphi_{78}))$
\item $\varphi_{218} :\equiv d_{2} \in \mathbb{R} \land 0 < d_{2} \land \varphi_{31} \to \delta \in \mathbb{R} \land 0 < \delta \land \delta \leq d_{1} \land \delta \leq d_{2} \to \varphi_{40}$
\item $\varphi_{219} :\equiv \mathrm{FLim}(f,c,\alpha) \to r \in \mathbb{R} \land 0 < r \land \varphi_{26} \to x \in \mathbb{R} \to c - r < x \to x < c \to \varphi_{26}$
\item $\varphi_{220} :\equiv \mathrm{FLim}(f,c,\alpha) \to r \in \mathbb{R} \land 0 < r \land \varphi_{26} \to x \in \mathbb{R} \to c < x \to x < c + r \to \varphi_{26}$
\item $\varphi_{221} :\equiv 0 < \varepsilon \to d_{1} \in \mathbb{R} \land 0 < d_{1} \land \varphi_{29} \to d_{2} \in \mathbb{R} \land 0 < d_{2} \land \varphi_{31} \to 0 < d_{2}$
\item $\varphi_{222} :\equiv 0 < \varepsilon \to d_{1} \in \mathbb{R} \land 0 < d_{1} \land \varphi_{29} \to d_{2} \in \mathbb{R} \land 0 < d_{2} \land \varphi_{31} \to 0 < d_{1}$
\item $\varphi_{223} :\equiv d_{2} \in \mathbb{R} \land 0 < d_{2} \land \varphi_{31} \to \delta \in \mathbb{R} \land 0 < \delta \land \delta \leq d_{1} \land \delta \leq d_{2} \to \varphi_{52}$
\item $\varphi_{224} :\equiv d_{2} \in \mathbb{R} \land 0 < d_{2} \land \varphi_{31} \to \delta \in \mathbb{R} \land 0 < \delta \land \delta \leq d_{1} \land \delta \leq d_{2} \to \varphi_{53}$
\item $\varphi_{225} :\equiv \mathrm{FLim}(f,c,\alpha) \land (r \in \mathbb{R} \land 0 < r \land \varphi_{26}) \to \forall x\,(\mathrm{FLim}(f,c,\alpha) \land (r \in \mathbb{R} \land 0 < r \land \varphi_{26}))$
\item $\varphi_{226} :\equiv (\mathrm{PNbhd}(f,c) \leftrightarrow \exists r\,(r \in \mathbb{R} \land 0 < r \land \varphi_{26})) \to \mathrm{FLim}(f,c,\alpha) \to \exists r\,(r \in \mathbb{R} \land 0 < r \land \varphi_{26})$
\item $\varphi_{227} :\equiv \varepsilon \in \mathbb{R} \to 0 < \varepsilon \to d_{1} \in \mathbb{R} \land 0 < d_{1} \land \varphi_{29} \to d_{1} \in \mathbb{R} \land 0 < d_{1} \land \varphi_{29}$
\item $\varphi_{228} :\equiv \forall c\,(\mathrm{PNbhd}(f,c) \leftrightarrow \exists r\,(r \in \mathbb{R} \land 0 < r \land \varphi_{26})) \to (\mathrm{PNbhd}(f,c) \leftrightarrow \exists r\,(r \in \mathbb{R} \land 0 < r \land \varphi_{26}))$
\item $\varphi_{229} :\equiv d_{2} \in \mathbb{R} \land 0 < d_{2} \land \varphi_{31} \to \delta \in \mathbb{R} \land 0 < \delta \land \delta \leq d_{1} \land \delta \leq d_{2} \to \varphi_{67}$
\item $\varphi_{230} :\equiv d_{2} \in \mathbb{R} \land 0 < d_{2} \land \varphi_{31} \to \delta \in \mathbb{R} \land 0 < \delta \land \delta \leq d_{1} \land \delta \leq d_{2} \to \varphi_{68}$
\item $\varphi_{231} :\equiv d_{2} \in \mathbb{R} \land 0 < d_{2} \land \varphi_{31} \to \delta \in \mathbb{R} \land 0 < \delta \land \delta \leq d_{1} \land \delta \leq d_{2} \to \varphi_{69}$
\item $\varphi_{232} :\equiv \varphi_{150} \to x \in \mathrm{dom}(f) \to c - \delta < x \to x < c \to \lvert f(x) - \alpha \rvert < \varepsilon$
\item $\varphi_{233} :\equiv \varphi_{150} \to x \in \mathrm{dom}(f) \to c < x \to x < c + \delta \to \lvert f(x) - \alpha \rvert < \varepsilon$
\item $\varphi_{234} :\equiv 0 < \varepsilon \to d_{1} \in \mathbb{R} \land 0 < d_{1} \land \varphi_{29} \to \exists \delta\,(\delta \in \mathbb{R} \land 0 < \delta \land \varphi_{78})$
\item $\varphi_{235} :\equiv \varphi_{56} \to \varphi_{57} \to d \in \mathbb{R} \land 0 < d \land d \leq r_{1} \land d \leq r_{2} \to x \in \mathbb{R} \to x \in \mathbb{R}$
\item $\varphi_{236} :\equiv d_{2} \in \mathbb{R} \land 0 < d_{2} \land \varphi_{31} \to \delta \in \mathbb{R} \land 0 < \delta \land \delta \leq d_{1} \land \delta \leq d_{2} \to \varphi_{78}$
\item $\varphi_{237} :\equiv \mathrm{FLim}(f,c,\alpha) \to \varepsilon \in \mathbb{R} \to 0 < \varepsilon \to \delta \in \mathbb{R} \land 0 < \delta \land \varphi_{78} \to \varphi_{32}$
\item $\varphi_{238} :\equiv \mathrm{FLim}(f,c,\alpha) \to \varepsilon \in \mathbb{R} \to 0 < \varepsilon \to \delta \in \mathbb{R} \land 0 < \delta \land \varphi_{78} \to \varphi_{33}$
\item $\varphi_{239} :\equiv \varphi_{112} \to \forall c\,(\mathrm{PNbhd}(f,c) \leftrightarrow \exists r\,(r \in \mathbb{R} \land 0 < r \land \varphi_{26}))$
\item $\varphi_{240} :\equiv \forall \delta\,(\neg (\delta \in \mathbb{R} \land 0 < \delta \land \varphi_{78})) \to \neg (\delta \in \mathbb{R} \land 0 < \delta \land \varphi_{78})$
\item $\varphi_{241} :\equiv \mathrm{FLimR}(f,c,\alpha) \land \mathrm{FLimL}(f,c,\alpha) \land \varphi_{56} \land \varphi_{57} \land (d \in \mathbb{R} \land 0 < d \land d \leq r_{1} \land d \leq r_{2})$
\item $\varphi_{242} :\equiv d_{2} \in \mathbb{R} \land 0 < d_{2} \land \varphi_{31} \to \delta \in \mathbb{R} \land 0 < \delta \land \delta \leq d_{1} \land \delta \leq d_{2} \to \varphi_{88}$
\item $\varphi_{243} :\equiv d_{2} \in \mathbb{R} \land 0 < d_{2} \land \varphi_{31} \to \delta \in \mathbb{R} \land 0 < \delta \land \delta \leq d_{1} \land \delta \leq d_{2} \to \varphi_{89}$
\item $\varphi_{244} :\equiv \mathrm{FLimR}(f,c,\alpha) \land \mathrm{FLimL}(f,c,\alpha) \to \varphi_{56} \to \varphi_{57} \to \exists d\,(d \in \mathbb{R} \land 0 < d \land d \leq r_{1} \land d \leq r_{2})$
\item $\varphi_{245} :\equiv \mathrm{FLim}(f,c,\alpha) \to \varepsilon \in \mathbb{R} \to 0 < \varepsilon \to \delta \in \mathbb{R} \land 0 < \delta \land \varphi_{78} \to \varphi_{54}$
\item $\varphi_{246} :\equiv \mathrm{FLim}(f,c,\alpha) \to \varepsilon \in \mathbb{R} \to 0 < \varepsilon \to \delta \in \mathbb{R} \land 0 < \delta \land \varphi_{78} \to \varphi_{55}$
\item $\varphi_{247} :\equiv \mathrm{FLim}(f,c,\alpha) \to \varepsilon \in \mathbb{R} \to 0 < \varepsilon \to \delta \in \mathbb{R} \land 0 < \delta \land \varphi_{78} \to \varphi_{69}$
\item $\varphi_{248} :\equiv \varepsilon \in \mathbb{R} \to 0 < \varepsilon \to \delta \in \mathbb{R} \land 0 < \delta \land \varphi_{78} \to \exists \delta\,(\delta \in \mathbb{R} \land 0 < \delta \land \varphi_{32})$
\item $\varphi_{249} :\equiv \varepsilon \in \mathbb{R} \to 0 < \varepsilon \to \delta \in \mathbb{R} \land 0 < \delta \land \varphi_{78} \to \exists \delta\,(\delta \in \mathbb{R} \land 0 < \delta \land \varphi_{33})$
\item $\varphi_{250} :\equiv \mathrm{FLim}(f,c,\alpha) \to \varepsilon \in \mathbb{R} \to 0 < \varepsilon \to \delta \in \mathbb{R} \land 0 < \delta \land \varphi_{78} \to \varphi_{78}$
\item $\varphi_{251} :\equiv \varphi_{189} \to \mathrm{FLimR}(f,c,\alpha) \land \mathrm{FLimL}(f,c,\alpha) \to c \in \mathbb{R}$
\item $\varphi_{252} :\equiv \varphi_{189} \to \mathrm{FLimR}(f,c,\alpha) \land \mathrm{FLimL}(f,c,\alpha) \to \alpha \in \mathbb{R}$
\item $\varphi_{253} :\equiv d_{2} \in \mathbb{R} \land 0 < d_{2} \land \varphi_{31} \to \delta \in \mathbb{R} \land 0 < \delta \land \delta \leq d_{1} \land \delta \leq d_{2} \to \delta \in \mathbb{R} \land 0 < \delta \land \varphi_{78}$
\item $\varphi_{254} :\equiv \mathrm{FLim}(f,c,\alpha) \to \varepsilon \in \mathbb{R} \to 0 < \varepsilon \to \delta \in \mathbb{R} \land 0 < \delta \land \varphi_{78} \to \delta \in \mathbb{R} \land 0 < \delta \land \varphi_{32}$
\item $\varphi_{255} :\equiv \mathrm{FLim}(f,c,\alpha) \to \varepsilon \in \mathbb{R} \to 0 < \varepsilon \to \delta \in \mathbb{R} \land 0 < \delta \land \varphi_{78} \to \delta \in \mathbb{R} \land 0 < \delta \land \varphi_{33}$
\item $\varphi_{256} :\equiv \forall d\,(\varphi_{207}) \to \forall d\,(\neg (d \in \mathbb{R} \land 0 < d \land d \leq r_{1} \land d \leq r_{2}))$
\item $\varphi_{257} :\equiv \mathrm{FLimR}(f,c,\alpha) \land \mathrm{FLimL}(f,c,\alpha) \to \varepsilon \in \mathbb{R} \to \varphi_{208}$
\item $\varphi_{258} :\equiv \forall \varepsilon \in \mathbb{R}\,(0 < \varepsilon \to \exists \delta\,(\delta \in \mathbb{R} \land 0 < \delta \land \varphi_{32})) \to \varepsilon \in \mathbb{R} \to 0 < \varepsilon \to \exists \delta\,(\delta \in \mathbb{R} \land 0 < \delta \land \varphi_{32})$
\item $\varphi_{259} :\equiv \forall \varepsilon \in \mathbb{R}\,(0 < \varepsilon \to \exists \delta\,(\delta \in \mathbb{R} \land 0 < \delta \land \varphi_{33})) \to \varepsilon \in \mathbb{R} \to 0 < \varepsilon \to \exists \delta\,(\delta \in \mathbb{R} \land 0 < \delta \land \varphi_{33})$
\item $\varphi_{260} :\equiv \mathrm{FLimR}(f,c,\alpha) \land \mathrm{FLimL}(f,c,\alpha) \to \varepsilon \in \mathbb{R} \to \varphi_{221}$
\item $\varphi_{261} :\equiv \mathrm{FLimR}(f,c,\alpha) \land \mathrm{FLimL}(f,c,\alpha) \to \varepsilon \in \mathbb{R} \to \varphi_{222}$
\item $\varphi_{262} :\equiv \mathrm{FLim}(f,c,\alpha) \to \varepsilon \in \mathbb{R} \to 0 < \varepsilon \to \delta \in \mathbb{R} \land 0 < \delta \land \varphi_{78} \to \delta \in \mathbb{R} \land 0 < \delta \land \varphi_{78}$
\item $\varphi_{263} :\equiv \mathrm{FLimR}(f,c,\alpha) \land \mathrm{FLimL}(f,c,\alpha) \to \varepsilon \in \mathbb{R} \to 0 < \varepsilon \to \varphi_{205}$
\item $\varphi_{264} :\equiv \mathrm{FLimR}(f,c,\alpha) \land \mathrm{FLimL}(f,c,\alpha) \to \varepsilon \in \mathbb{R} \to 0 < \varepsilon \to \varphi_{206}$
\item $\varphi_{265} :\equiv \mathrm{FLimR}(f,c,\alpha) \land \mathrm{FLimL}(f,c,\alpha) \to \varphi_{56} \to \varphi_{57} \to \varphi_{37}$
\item $\varphi_{266} :\equiv \mathrm{FLimR}(f,c,\alpha) \land \mathrm{FLimL}(f,c,\alpha) \to \varepsilon \in \mathbb{R} \to \varphi_{234}$
\item $\varphi_{267} :\equiv \mathrm{FLimR}(f,c,\alpha) \land \mathrm{FLimL}(f,c,\alpha) \land \varepsilon \in \mathbb{R} \land 0 < \varepsilon \land (d_{1} \in \mathbb{R} \land 0 < d_{1} \land \varphi_{29}) \land \neg \exists \delta\,(\delta \in \mathbb{R} \land 0 < \delta \land \varphi_{78})$
\item $\varphi_{268} :\equiv \mathrm{FLimR}(f,c,\alpha) \land \mathrm{FLimL}(f,c,\alpha) \to \varphi_{56} \to \varphi_{57} \to \varphi_{57}$
\item $\varphi_{269} :\equiv 0 < \varepsilon \to d_{1} \in \mathbb{R} \land 0 < d_{1} \land \varphi_{29} \to d_{2} \in \mathbb{R} \land 0 < d_{2} \land \varphi_{31} \to \varphi_{27}$
\item $\varphi_{270} :\equiv \mathrm{FLimR}(f,c,\alpha) \land \mathrm{FLimL}(f,c,\alpha) \to \varphi_{56} \to \varphi_{57} \to \varphi_{82}$
\item $\varphi_{271} :\equiv 0 < \varepsilon \to d_{1} \in \mathbb{R} \land 0 < d_{1} \land \varphi_{29} \to d_{2} \in \mathbb{R} \land 0 < d_{2} \land \varphi_{31} \to \varphi_{58}$
\item $\varphi_{272} :\equiv d_{1} \in \mathbb{R} \land 0 < d_{1} \land \varphi_{29} \to d_{2} \in \mathbb{R} \land 0 < d_{2} \land \varphi_{31} \to d_{2} \in \mathbb{R} \land 0 < d_{2} \land \varphi_{31}$
\item $\varphi_{273} :\equiv \mathrm{FLimR}(f,c,\alpha) \land \mathrm{FLimL}(f,c,\alpha) \to \varepsilon \in \mathbb{R} \to 0 < \varepsilon \to d_{1} \in \mathbb{R} \land 0 < d_{1} \land \varphi_{29} \to \varphi_{125}$
\item $\varphi_{274} :\equiv \mathrm{FLimR}(f,c,\alpha) \land \mathrm{FLimL}(f,c,\alpha) \land \varepsilon \in \mathbb{R} \land 0 < \varepsilon \land (d_{1} \in \mathbb{R} \land 0 < d_{1} \land \varphi_{29}) \land (d_{2} \in \mathbb{R} \land 0 < d_{2} \land \varphi_{31}) \land (\delta \in \mathbb{R} \land 0 < \delta \land \delta \leq d_{1} \land \delta \leq d_{2})$
\item $\varphi_{275} :\equiv \mathrm{FLimR}(f,c,\alpha) \land \mathrm{FLimL}(f,c,\alpha) \to \varphi_{56} \to \varphi_{57} \to \varphi_{98}$
\item $\varphi_{276} :\equiv \varphi_{241} \to x \in \mathbb{R} \to 0 < \lvert x - c \rvert \to \lvert x - c \rvert < d \to x \in \mathrm{dom}(f)$
\item $\varphi_{277} :\equiv \mathrm{FLimR}(f,c,\alpha) \land \mathrm{FLimL}(f,c,\alpha) \to \varphi_{56} \to \varphi_{57} \to \varphi_{101}$
\item $\varphi_{278} :\equiv \mathrm{FLimR}(f,c,\alpha) \land \mathrm{FLimL}(f,c,\alpha) \to \varphi_{56} \to \varphi_{57} \to \varphi_{102}$
\item $\varphi_{279} :\equiv 0 < \varepsilon \to d_{1} \in \mathbb{R} \land 0 < d_{1} \land \varphi_{29} \to d_{2} \in \mathbb{R} \land 0 < d_{2} \land \varphi_{31} \to \varphi_{93}$
\item $\varphi_{280} :\equiv \mathrm{FLimR}(f,c,\alpha) \land \mathrm{FLimL}(f,c,\alpha) \to \varphi_{56} \to \varphi_{57} \to \varphi_{103}$
\item $\varphi_{281} :\equiv \mathrm{FLimR}(f,c,\alpha) \land \mathrm{FLimL}(f,c,\alpha) \to \varphi_{56} \to \varphi_{57} \to \varphi_{104}$
\item $\varphi_{282} :\equiv \mathrm{FLimR}(f,c,\alpha) \land \mathrm{FLimL}(f,c,\alpha) \to \varphi_{56} \to \varphi_{57} \to \varphi_{105}$
\item $\varphi_{283} :\equiv \mathrm{FLimR}(f,c,\alpha) \land \mathrm{FLimL}(f,c,\alpha) \to \varphi_{56} \to \varphi_{57} \to d \in \mathbb{R} \land 0 < d \land d \leq r_{1} \land d \leq r_{2} \to \varphi_{25}$
\item $\varphi_{284} :\equiv \mathrm{FLimR}(f,c,\alpha) \land \mathrm{FLimL}(f,c,\alpha) \to \varphi_{56} \to \varphi_{57} \to \varphi_{106}$
\item $\varphi_{285} :\equiv \mathrm{FLimR}(f,c,\alpha) \land \mathrm{FLimL}(f,c,\alpha) \to \varphi_{56} \to \varphi_{57} \to d \in \mathbb{R} \land 0 < d \land d \leq r_{1} \land d \leq r_{2} \to \varphi_{42}$
\item $\varphi_{286} :\equiv \mathrm{FLimR}(f,c,\alpha) \land \mathrm{FLimL}(f,c,\alpha) \to \varphi_{56} \to \varphi_{57} \to d \in \mathbb{R} \land 0 < d \land d \leq r_{1} \land d \leq r_{2} \to \varphi_{43}$
\item $\varphi_{287} :\equiv \mathrm{FLimR}(f,c,\alpha) \land \mathrm{FLimL}(f,c,\alpha) \to \varphi_{56} \to \varphi_{57} \to d \in \mathbb{R} \land 0 < d \land d \leq r_{1} \land d \leq r_{2} \to \varphi_{44}$
\item $\varphi_{288} :\equiv \mathrm{FLimR}(f,c,\alpha) \land \mathrm{FLimL}(f,c,\alpha) \to \varphi_{56} \to \varphi_{57} \to d \in \mathbb{R} \land 0 < d \land d \leq r_{1} \land d \leq r_{2} \to \varphi_{45}$
\item $\varphi_{289} :\equiv (\mathrm{FLim}(f,c,\alpha) \to r \in \mathbb{R} \land 0 < r \land \varphi_{26} \to \exists r\,(\varphi_{47})) \to \varphi_{113}$
\item $\varphi_{290} :\equiv (\mathrm{FLim}(f,c,\alpha) \to r \in \mathbb{R} \land 0 < r \land \varphi_{26} \to \exists r\,(\varphi_{48})) \to \varphi_{114}$
\item $\varphi_{291} :\equiv d_{1} \in \mathbb{R} \land 0 < d_{1} \land \varphi_{29} \to d_{2} \in \mathbb{R} \land 0 < d_{2} \land \varphi_{31} \to \exists \delta\,(\delta \in \mathbb{R} \land 0 < \delta \land \varphi_{78})$
\item $\varphi_{292} :\equiv \mathrm{FLimR}(f,c,\alpha) \land \mathrm{FLimL}(f,c,\alpha) \to \varphi_{56} \to \varphi_{57} \to d \in \mathbb{R} \land 0 < d \land d \leq r_{1} \land d \leq r_{2} \to \varphi_{59}$
\item $\varphi_{293} :\equiv \mathrm{FLimR}(f,c,\alpha) \land \mathrm{FLimL}(f,c,\alpha) \to \varphi_{56} \to \varphi_{57} \to d \in \mathbb{R} \land 0 < d \land d \leq r_{1} \land d \leq r_{2} \to \varphi_{61}$
\item $\varphi_{294} :\equiv 0 < \varepsilon \to d_{1} \in \mathbb{R} \land 0 < d_{1} \land \varphi_{29} \to d_{2} \in \mathbb{R} \land 0 < d_{2} \land \varphi_{31} \to \varphi_{109}$
\item $\varphi_{295} :\equiv (\mathrm{FLimR}(f,c,\alpha) \land \mathrm{FLimL}(f,c,\alpha) \to \alpha \in \mathbb{R}) \to \varphi_{2} \to \varphi_{177} \to \varphi_{4}$
\item $\varphi_{296} :\equiv \mathrm{FLimR}(f,c,\alpha) \land \mathrm{FLimL}(f,c,\alpha) \to \varepsilon \in \mathbb{R} \to \varphi_{269}$
\item $\varphi_{297} :\equiv \mathrm{FLimR}(f,c,\alpha) \land \mathrm{FLimL}(f,c,\alpha) \to \varphi_{56} \to \varphi_{57} \to \varphi_{123}$
\item $\varphi_{298} :\equiv \mathrm{FLimR}(f,c,\alpha) \land \mathrm{FLimL}(f,c,\alpha) \to \varepsilon \in \mathbb{R} \to \varphi_{271}$
\item $\varphi_{299} :\equiv \forall r\,(\forall r_{1}\,(\neg \varphi_{56}) \to \neg \varphi_{48}) \to \forall r\,\forall r_{1}\,(\neg \varphi_{56}) \to \forall r\,(\neg \varphi_{48})$
\item $\varphi_{300} :\equiv \forall r\,(\forall r_{2}\,(\neg \varphi_{57}) \to \neg \varphi_{47}) \to \forall r\,\forall r_{2}\,(\neg \varphi_{57}) \to \forall r\,(\neg \varphi_{47})$
\item $\varphi_{301} :\equiv \mathrm{FLimR}(f,c,\alpha) \land \mathrm{FLimL}(f,c,\alpha) \to \varphi_{56} \to \varphi_{57} \to \varphi_{143}$
\item $\varphi_{302} :\equiv \mathrm{FLimR}(f,c,\alpha) \land \mathrm{FLimL}(f,c,\alpha) \to \varphi_{56} \to \varphi_{57} \to \varphi_{144}$
\item $\varphi_{303} :\equiv (\forall r\,(\varphi_{115}) \to \varphi_{119}) \to \varphi_{13} \to \mathrm{FLim}(f,c,\alpha) \to \mathrm{LNbhd}(f,c)$
\item $\varphi_{304} :\equiv (\forall r\,(\varphi_{116}) \to \varphi_{120}) \to \varphi_{14} \to \mathrm{FLim}(f,c,\alpha) \to \mathrm{RNbhd}(f,c)$
\item $\varphi_{305} :\equiv \mathrm{FLimR}(f,c,\alpha) \land \mathrm{FLimL}(f,c,\alpha) \to \varepsilon \in \mathbb{R} \to \varphi_{279}$
\item $\varphi_{306} :\equiv \mathrm{FLimR}(f,c,\alpha) \land \mathrm{FLimL}(f,c,\alpha) \to \varphi_{56} \to \varphi_{57} \to \varphi_{148}$
\item $\varphi_{307} :\equiv \mathrm{FLimR}(f,c,\alpha) \land \mathrm{FLimL}(f,c,\alpha) \to \varphi_{56} \to \varphi_{57} \to \varphi_{149}$
\item $\varphi_{308} :\equiv (\mathrm{FLim}(f,c,\alpha) \to \varphi_{140}) \to \mathrm{FLim}(f,c,\alpha) \to r \in \mathbb{R} \land 0 < r \land \varphi_{26} \to \varphi_{47}$
\item $\varphi_{309} :\equiv (\forall r\,(\neg \varphi_{47}) \to \neg \varphi_{47}) \to \mathrm{FLim}(f,c,\alpha) \to r \in \mathbb{R} \land 0 < r \land \varphi_{26} \to \exists r\,(\varphi_{47})$
\item $\varphi_{310} :\equiv (\mathrm{FLim}(f,c,\alpha) \to \varphi_{141}) \to \mathrm{FLim}(f,c,\alpha) \to r \in \mathbb{R} \land 0 < r \land \varphi_{26} \to \varphi_{48}$
\item $\varphi_{311} :\equiv (\forall r\,(\neg \varphi_{48}) \to \neg \varphi_{48}) \to \mathrm{FLim}(f,c,\alpha) \to r \in \mathbb{R} \land 0 < r \land \varphi_{26} \to \exists r\,(\varphi_{48})$
\item $\varphi_{312} :\equiv \mathrm{FLimR}(f,c,\alpha) \land \mathrm{FLimL}(f,c,\alpha) \to \varepsilon \in \mathbb{R} \to 0 < \varepsilon \to \varphi_{272}$
\item $\varphi_{313} :\equiv \mathrm{FLimR}(f,c,\alpha) \land \mathrm{FLimL}(f,c,\alpha) \to \varepsilon \in \mathbb{R} \to \varphi_{294}$
\item $\varphi_{314} :\equiv 0 < \varepsilon \to d_{1} \in \mathbb{R} \land 0 < d_{1} \land \varphi_{29} \to d_{2} \in \mathbb{R} \land 0 < d_{2} \land \varphi_{31} \to \varphi_{152}$
\item $\varphi_{315} :\equiv \mathrm{FLimR}(f,c,\alpha) \land \mathrm{FLimL}(f,c,\alpha) \to \varepsilon \in \mathbb{R} \to 0 < \varepsilon \to d_{1} \in \mathbb{R} \land 0 < d_{1} \land \varphi_{29} \to \varphi_{209}$
\item $\varphi_{316} :\equiv \mathrm{FLimR}(f,c,\alpha) \land \mathrm{FLimL}(f,c,\alpha) \to \varepsilon \in \mathbb{R} \to 0 < \varepsilon \to d_{1} \in \mathbb{R} \land 0 < d_{1} \land \varphi_{29} \to \varphi_{210}$
\item $\varphi_{317} :\equiv \mathrm{FLimR}(f,c,\alpha) \land \mathrm{FLimL}(f,c,\alpha) \to \varepsilon \in \mathbb{R} \to 0 < \varepsilon \to d_{1} \in \mathbb{R} \land 0 < d_{1} \land \varphi_{29} \to \varphi_{212}$
\item $\varphi_{318} :\equiv \mathrm{FLimR}(f,c,\alpha) \land \mathrm{FLimL}(f,c,\alpha) \to \varepsilon \in \mathbb{R} \to 0 < \varepsilon \to d_{1} \in \mathbb{R} \land 0 < d_{1} \land \varphi_{29} \to \varphi_{213}$
\item $\varphi_{319} :\equiv \mathrm{FLimR}(f,c,\alpha) \land \mathrm{FLimL}(f,c,\alpha) \to \varepsilon \in \mathbb{R} \to 0 < \varepsilon \to d_{1} \in \mathbb{R} \land 0 < d_{1} \land \varphi_{29} \to \varphi_{214}$
\item $\varphi_{320} :\equiv (\mathrm{FLimR}(f,c,\alpha) \land \mathrm{FLimL}(f,c,\alpha) \to c \in \mathbb{R}) \to \varphi_{295}$
\item $\varphi_{321} :\equiv \mathrm{FLimR}(f,c,\alpha) \land \mathrm{FLimL}(f,c,\alpha) \to \varepsilon \in \mathbb{R} \to 0 < \varepsilon \to d_{1} \in \mathbb{R} \land 0 < d_{1} \land \varphi_{29} \to \varphi_{218}$
\item $\varphi_{322} :\equiv \mathrm{FLimR}(f,c,\alpha) \land \mathrm{FLimL}(f,c,\alpha) \to \varepsilon \in \mathbb{R} \to 0 < \varepsilon \to \varphi_{291}$
\item $\varphi_{323} :\equiv \mathrm{FLimR}(f,c,\alpha) \land \mathrm{FLimL}(f,c,\alpha) \land \varepsilon \in \mathbb{R} \land 0 < \varepsilon \land (d_{1} \in \mathbb{R} \land 0 < d_{1} \land \varphi_{29}) \land (d_{2} \in \mathbb{R} \land 0 < d_{2} \land \varphi_{31}) \land \neg \exists \delta\,(\delta \in \mathbb{R} \land 0 < \delta \land \varphi_{78})$
\item $\varphi_{324} :\equiv \mathrm{FLimR}(f,c,\alpha) \land \mathrm{FLimL}(f,c,\alpha) \to \varepsilon \in \mathbb{R} \to 0 < \varepsilon \to d_{1} \in \mathbb{R} \land 0 < d_{1} \land \varphi_{29} \to \varphi_{223}$
\item $\varphi_{325} :\equiv \mathrm{FLimR}(f,c,\alpha) \land \mathrm{FLimL}(f,c,\alpha) \to \varepsilon \in \mathbb{R} \to 0 < \varepsilon \to d_{1} \in \mathbb{R} \land 0 < d_{1} \land \varphi_{29} \to \varphi_{224}$
\item $\varphi_{326} :\equiv \mathrm{FLimR}(f,c,\alpha) \land \mathrm{FLimL}(f,c,\alpha) \to \varepsilon \in \mathbb{R} \to 0 < \varepsilon \to d_{1} \in \mathbb{R} \land 0 < d_{1} \land \varphi_{29} \to \varphi_{229}$
\item $\varphi_{327} :\equiv \mathrm{FLimR}(f,c,\alpha) \land \mathrm{FLimL}(f,c,\alpha) \to \varepsilon \in \mathbb{R} \to 0 < \varepsilon \to d_{1} \in \mathbb{R} \land 0 < d_{1} \land \varphi_{29} \to \varphi_{230}$
\item $\varphi_{328} :\equiv \mathrm{FLimR}(f,c,\alpha) \land \mathrm{FLimL}(f,c,\alpha) \to \varepsilon \in \mathbb{R} \to 0 < \varepsilon \to d_{1} \in \mathbb{R} \land 0 < d_{1} \land \varphi_{29} \to \varphi_{231}$
\item $\varphi_{329} :\equiv 0 < \varepsilon \to d_{1} \in \mathbb{R} \land 0 < d_{1} \land \varphi_{29} \to d_{2} \in \mathbb{R} \land 0 < d_{2} \land \varphi_{31} \to \varphi_{171}$
\item $\varphi_{330} :\equiv 0 < \varepsilon \to d_{1} \in \mathbb{R} \land 0 < d_{1} \land \varphi_{29} \to d_{2} \in \mathbb{R} \land 0 < d_{2} \land \varphi_{31} \to \varphi_{172}$
\item $\varphi_{331} :\equiv \mathrm{FLimR}(f,c,\alpha) \land \mathrm{FLimL}(f,c,\alpha) \to \varepsilon \in \mathbb{R} \to 0 < \varepsilon \to d_{1} \in \mathbb{R} \land 0 < d_{1} \land \varphi_{29} \to \varphi_{236}$
\item $\varphi_{332} :\equiv 0 < \varepsilon \to d_{1} \in \mathbb{R} \land 0 < d_{1} \land \varphi_{29} \to d_{2} \in \mathbb{R} \land 0 < d_{2} \land \varphi_{31} \to \varphi_{175}$
\item $\varphi_{333} :\equiv 0 < \varepsilon \to d_{1} \in \mathbb{R} \land 0 < d_{1} \land \varphi_{29} \to d_{2} \in \mathbb{R} \land 0 < d_{2} \land \varphi_{31} \to \varphi_{176}$
\item $\varphi_{334} :\equiv \mathrm{FLimR}(f,c,\alpha) \land \mathrm{FLimL}(f,c,\alpha) \to \varepsilon \in \mathbb{R} \to 0 < \varepsilon \to d_{1} \in \mathbb{R} \land 0 < d_{1} \land \varphi_{29} \to \varphi_{242}$
\item $\varphi_{335} :\equiv \mathrm{FLimR}(f,c,\alpha) \land \mathrm{FLimL}(f,c,\alpha) \to \varepsilon \in \mathbb{R} \to 0 < \varepsilon \to d_{1} \in \mathbb{R} \land 0 < d_{1} \land \varphi_{29} \to \varphi_{243}$
\item $\varphi_{336} :\equiv (\mathrm{FLim}(f,c,\alpha) \to \varepsilon \in \mathbb{R} \to \varphi_{151}) \to \mathrm{FLim}(f,c,\alpha) \to \varepsilon \in \mathbb{R} \to 0 < \varepsilon \to \varphi_{163}$
\item $\varphi_{337} :\equiv (\mathrm{FLim}(f,c,\alpha) \to \varepsilon \in \mathbb{R} \to \varphi_{151}) \to \mathrm{FLim}(f,c,\alpha) \to \varepsilon \in \mathbb{R} \to 0 < \varepsilon \to \varphi_{164}$
\item $\varphi_{338} :\equiv (\mathrm{FLimR}(f,c,\alpha) \land \mathrm{FLimL}(f,c,\alpha) \to \varphi_{56} \to \varphi_{56}) \to \varphi_{197}$
\item $\varphi_{339} :\equiv (\mathrm{FLim}(f,c,\alpha) \to \mathrm{RNbhd}(f,c)) \to (\mathrm{FLim}(f,c,\alpha) \to \mathrm{LNbhd}(f,c)) \to \varphi_{139} \to \varphi_{138} \to \varphi_{5}$
\item $\varphi_{340} :\equiv (\mathrm{FLimR}(f,c,\alpha) \land \mathrm{FLimL}(f,c,\alpha) \to \varphi_{56} \to \varphi_{56}) \to \varphi_{201}$
\item $\varphi_{341} :\equiv \mathrm{FLimR}(f,c,\alpha) \land \mathrm{FLimL}(f,c,\alpha) \to \varepsilon \in \mathbb{R} \to 0 < \varepsilon \to d_{1} \in \mathbb{R} \land 0 < d_{1} \land \varphi_{29} \to \varphi_{253}$
\item $\varphi_{342} :\equiv \mathrm{FLimR}(f,c,\alpha) \land \mathrm{FLimL}(f,c,\alpha) \to \varepsilon \in \mathbb{R} \to \varphi_{314}$
\item $\varphi_{343} :\equiv \forall \delta\,(\varphi_{323}) \to \forall \delta\,(\neg (\delta \in \mathbb{R} \land 0 < \delta \land \delta \leq d_{1} \land \delta \leq d_{2}))$
\item $\varphi_{344} :\equiv \mathrm{FLimR}(f,c,\alpha) \land \mathrm{FLimL}(f,c,\alpha) \to \varepsilon \in \mathbb{R} \to \varphi_{329}$
\item $\varphi_{345} :\equiv \mathrm{FLimR}(f,c,\alpha) \land \mathrm{FLimL}(f,c,\alpha) \to \varepsilon \in \mathbb{R} \to \varphi_{330}$
\item $\varphi_{346} :\equiv \mathrm{FLimR}(f,c,\alpha) \land \mathrm{FLimL}(f,c,\alpha) \to \varepsilon \in \mathbb{R} \to \varphi_{332}$
\item $\varphi_{347} :\equiv \mathrm{FLimR}(f,c,\alpha) \land \mathrm{FLimL}(f,c,\alpha) \to \varepsilon \in \mathbb{R} \to \varphi_{333}$
\item $\varphi_{348} :\equiv (\mathrm{RFn}(f) \to x \in \mathrm{dom}(f) \to x \in \mathbb{R}) \to (\mathrm{FLim}(f,c,\alpha) \to \mathrm{RFn}(f)) \to \varphi_{336}$
\item $\varphi_{349} :\equiv (\mathrm{RFn}(f) \to x \in \mathrm{dom}(f) \to x \in \mathbb{R}) \to (\mathrm{FLim}(f,c,\alpha) \to \mathrm{RFn}(f)) \to \varphi_{337}$
\item $\varphi_{350} :\equiv (\mathrm{FLim}(f,c,\alpha) \to c \in \mathbb{R}) \to (\mathrm{FLim}(f,c,\alpha) \to \alpha \in \mathbb{R}) \to \varphi_{339}$
\item $\varphi_{351} :\equiv \varphi_{124} \to (\mathrm{FLim}(f,c,\alpha) \to \varphi_{169}) \to \mathrm{FLim}(f,c,\alpha) \to \varphi_{136}$
\item $\varphi_{352} :\equiv \varphi_{124} \to (\mathrm{FLim}(f,c,\alpha) \to \varphi_{170}) \to \mathrm{FLim}(f,c,\alpha) \to \varphi_{137}$
\item $\varphi_{353} :\equiv \varphi_{128} \to (\mathrm{FLim}(f,c,\alpha) \to \varphi_{135}) \to \mathrm{FLim}(f,c,\alpha) \to \varphi_{170}$
\item $\varphi_{354} :\equiv (\mathrm{FLimR}(f,c,\alpha) \land \mathrm{FLimL}(f,c,\alpha) \to \varphi_{56} \to \varphi_{121}) \to \varphi_{207} \to \neg (d \in \mathbb{R} \land 0 < d \land d \leq r_{1} \land d \leq r_{2})$
\item $\varphi_{355} :\equiv (\mathrm{FLim}(f,c,\alpha) \to \varphi_{248}) \to \varphi_{133} \to \neg (\delta \in \mathbb{R} \land 0 < \delta \land \varphi_{78})$
\item $\varphi_{356} :\equiv (\mathrm{FLim}(f,c,\alpha) \to \varphi_{249}) \to \varphi_{134} \to \neg (\delta \in \mathbb{R} \land 0 < \delta \land \varphi_{78})$
\item $\varphi_{357} :\equiv (\mathrm{FLimR}(f,c,\alpha) \land \mathrm{FLimL}(f,c,\alpha) \to \varphi_{227}) \to \varphi_{261}$
\item $\varphi_{358} :\equiv (\mathrm{FLimR}(f,c,\alpha) \land \mathrm{FLimL}(f,c,\alpha) \to \varphi_{227}) \to \varphi_{263}$
\item $\varphi_{359} :\equiv \forall \delta\,(\varphi_{133} \to \neg (\delta \in \mathbb{R} \land 0 < \delta \land \varphi_{78})) \to \forall \delta\,(\varphi_{133}) \to \forall \delta\,(\neg (\delta \in \mathbb{R} \land 0 < \delta \land \varphi_{78}))$
\item $\varphi_{360} :\equiv \forall \delta\,(\varphi_{134} \to \neg (\delta \in \mathbb{R} \land 0 < \delta \land \varphi_{78})) \to \forall \delta\,(\varphi_{134}) \to \forall \delta\,(\neg (\delta \in \mathbb{R} \land 0 < \delta \land \varphi_{78}))$
\item $\varphi_{361} :\equiv (\mathrm{FLimR}(f,c,\alpha) \land \mathrm{FLimL}(f,c,\alpha) \to \varphi_{56} \to \varphi_{56}) \to \varphi_{306}$
\item $\varphi_{362} :\equiv \forall d_{1}\,(\varphi_{174} \to \neg (d_{1} \in \mathbb{R} \land 0 < d_{1} \land \varphi_{29})) \to \forall d_{1}\,(\varphi_{174}) \to \forall d_{1}\,(\neg (d_{1} \in \mathbb{R} \land 0 < d_{1} \land \varphi_{29}))$
\item $\varphi_{363} :\equiv \varphi_{297} \to \mathrm{FLimR}(f,c,\alpha) \land \mathrm{FLimL}(f,c,\alpha) \to \varphi_{56} \to \varphi_{121}$
\item $\varphi_{364} :\equiv (\mathrm{FLim}(f,c,\alpha) \to \exists r\,(r \in \mathbb{R} \land 0 < r \land \varphi_{26})) \to \forall r\,(\varphi_{115}) \to \varphi_{303}$
\item $\varphi_{365} :\equiv (\mathrm{FLim}(f,c,\alpha) \to \exists r\,(r \in \mathbb{R} \land 0 < r \land \varphi_{26})) \to \forall r\,(\varphi_{116}) \to \varphi_{304}$
\item $\varphi_{366} :\equiv (\mathrm{FLim}(f,c,\alpha) \to r \in \mathbb{R} \land 0 < r \land \varphi_{26} \to r \in \mathbb{R}) \to \varphi_{353}$
\item $\varphi_{367} :\equiv (\mathrm{FLim}(f,c,\alpha) \to r \in \mathbb{R} \land 0 < r \land \varphi_{26} \to r \in \mathbb{R}) \to \varphi_{142} \to \varphi_{147} \to \mathrm{FLim}(f,c,\alpha) \to \varphi_{169}$
\item $\varphi_{368} :\equiv \varphi_{283} \to \mathrm{FLimR}(f,c,\alpha) \land \mathrm{FLimL}(f,c,\alpha) \to \varphi_{56} \to \varphi_{211}$
\item $\varphi_{369} :\equiv (\mathrm{FLimR}(f,c,\alpha) \land \mathrm{FLimL}(f,c,\alpha) \to c \in \mathbb{R}) \to \varphi_{270} \to \varphi_{282}$
\item $\varphi_{370} :\equiv (\mathrm{FLim}(f,c,\alpha) \to \varepsilon \in \mathbb{R} \to \varphi_{153}) \to (\mathrm{FLim}(f,c,\alpha) \to \varepsilon \in \mathbb{R} \to 0 < \varepsilon \to \varphi_{167}) \to \varphi_{246}$
\item $\varphi_{371} :\equiv (\mathrm{FLim}(f,c,\alpha) \to \varepsilon \in \mathbb{R} \to 0 < \varepsilon \to \varphi_{155}) \to (\mathrm{FLim}(f,c,\alpha) \to \varepsilon \in \mathbb{R} \to 0 < \varepsilon \to \varphi_{160}) \to \varphi_{245}$
\item $\varphi_{372} :\equiv (\mathrm{FLim}(f,c,\alpha) \to \varepsilon \in \mathbb{R} \to \varphi_{151}) \to \varphi_{245} \to \mathrm{FLim}(f,c,\alpha) \to \varepsilon \in \mathbb{R} \to 0 < \varepsilon \to \varphi_{178}$
\item $\varphi_{373} :\equiv (\mathrm{FLim}(f,c,\alpha) \to \varepsilon \in \mathbb{R} \to \varphi_{151}) \to \varphi_{246} \to \mathrm{FLim}(f,c,\alpha) \to \varepsilon \in \mathbb{R} \to 0 < \varepsilon \to \varphi_{179}$
\item $\varphi_{374} :\equiv (\mathrm{FLimR}(f,c,\alpha) \land \mathrm{FLimL}(f,c,\alpha) \to \varphi_{227}) \to \varphi_{346}$
\item $\varphi_{375} :\equiv \forall r\,(\forall r_{1}\,(\neg \varphi_{56}) \to \neg \varphi_{48}) \to \varphi_{299} \to (\forall r_{1}\,(\neg \varphi_{56}) \to \forall r\,\forall r_{1}\,(\neg \varphi_{56})) \to \varphi_{2}$
\item $\varphi_{376} :\equiv \forall d_{2}\,(\varphi_{267} \to \neg (d_{2} \in \mathbb{R} \land 0 < d_{2} \land \varphi_{31})) \to \forall d_{2}\,(\varphi_{267}) \to \forall d_{2}\,(\neg (d_{2} \in \mathbb{R} \land 0 < d_{2} \land \varphi_{31}))$
\item $\varphi_{377} :\equiv (\mathrm{FLimR}(f,c,\alpha) \land \mathrm{FLimL}(f,c,\alpha) \to c \in \mathbb{R}) \to \varphi_{305} \to \varphi_{319}$
\item $\varphi_{378} :\equiv \forall r\,(\forall r_{2}\,(\neg \varphi_{57}) \to \neg \varphi_{47}) \to \varphi_{300} \to (\forall r_{2}\,(\neg \varphi_{57}) \to \forall r\,\forall r_{2}\,(\neg \varphi_{57})) \to \varphi_{117}$
\item $\varphi_{379} :\equiv (\mathrm{FLimR}(f,c,\alpha) \land \mathrm{FLimL}(f,c,\alpha) \to \varphi_{56} \to \varphi_{211}) \to \varphi_{183} \to \varphi_{297}$
\item $\varphi_{380} :\equiv (x \in \mathbb{R} \to \varphi_{84}) \to \varphi_{124} \to (\mathrm{FLim}(f,c,\alpha) \to c \in \mathbb{R}) \to \varphi_{366}$
\item $\varphi_{381} :\equiv (x \in \mathbb{R} \to \varphi_{83}) \to \varphi_{124} \to (\mathrm{FLim}(f,c,\alpha) \to c \in \mathbb{R}) \to \varphi_{367}$
\item $\varphi_{382} :\equiv \forall x\,(\varphi_{158}) \to (\forall x\,(\varphi_{158}) \to \varphi_{165}) \to \varphi_{225} \to \mathrm{FLim}(f,c,\alpha) \to \varphi_{140}$
\item $\varphi_{383} :\equiv \forall x\,(\varphi_{159}) \to (\forall x\,(\varphi_{159}) \to \varphi_{166}) \to \varphi_{225} \to \mathrm{FLim}(f,c,\alpha) \to \varphi_{141}$
\item $\varphi_{384} :\equiv (\mathrm{FLimR}(f,c,\alpha) \land \mathrm{FLimL}(f,c,\alpha) \to \varphi_{235}) \to \varphi_{369}$
\item $\varphi_{385} :\equiv (\mathrm{FLim}(f,c,\alpha) \to \varepsilon \in \mathbb{R} \to 0 < \varepsilon \to \varphi_{164}) \to (\mathrm{FLim}(f,c,\alpha) \to c \in \mathbb{R}) \to \varphi_{168} \to \varphi_{370}$
\item $\varphi_{386} :\equiv (\mathrm{FLim}(f,c,\alpha) \to \varepsilon \in \mathbb{R} \to 0 < \varepsilon \to \varphi_{163}) \to (\mathrm{FLim}(f,c,\alpha) \to c \in \mathbb{R}) \to \varphi_{168} \to \varphi_{371}$
\item $\varphi_{387} :\equiv (\forall r_{1}\,(\varphi_{3} \to \neg \varphi_{56}) \to \forall r_{1}\,(\varphi_{3}) \to \forall r_{1}\,(\neg \varphi_{56})) \to (\varphi_{3} \to \forall r_{1}\,(\varphi_{3})) \to \varphi_{375}$
\item $\varphi_{388} :\equiv (\forall d\,(\varphi_{207} \to \neg (d \in \mathbb{R} \land 0 < d \land d \leq r_{1} \land d \leq r_{2})) \to \varphi_{256}) \to (\varphi_{207} \to \forall d\,(\varphi_{207})) \to \varphi_{198}$
\item $\varphi_{389} :\equiv (\mathrm{FLimR}(f,c,\alpha) \land \mathrm{FLimL}(f,c,\alpha) \to \varphi_{235}) \to \varphi_{280} \to \varphi_{285} \to \varphi_{286}$
\item $\varphi_{390} :\equiv (\mathrm{FLimR}(f,c,\alpha) \land \mathrm{FLimL}(f,c,\alpha) \to \varphi_{235}) \to \varphi_{287} \to \varphi_{281} \to \varphi_{288}$
\item $\varphi_{391} :\equiv (\mathrm{FLimR}(f,c,\alpha) \land \mathrm{FLimL}(f,c,\alpha) \to c \in \mathbb{R}) \to \varphi_{201} \to \varphi_{280} \to \varphi_{292} \to \varphi_{285}$
\item $\varphi_{392} :\equiv (\mathrm{FLimR}(f,c,\alpha) \land \mathrm{FLimL}(f,c,\alpha) \to c \in \mathbb{R}) \to \varphi_{202} \to \varphi_{281} \to \varphi_{293} \to \varphi_{287}$
\item $\varphi_{393} :\equiv (x \in \mathbb{R} \to c \in \mathbb{R} \to 0 < \lvert x - c \rvert \to x < c \lor c < x) \to \varphi_{313} \to \varphi_{377}$
\item $\varphi_{394} :\equiv \forall x\,(\varphi_{232}) \to (\forall x\,(\varphi_{232}) \to \forall x\,(\varphi_{150}) \to \varphi_{32}) \to (\varphi_{150} \to \forall x\,(\varphi_{150})) \to \varphi_{237}$
\item $\varphi_{395} :\equiv \forall x\,(\varphi_{233}) \to (\forall x\,(\varphi_{233}) \to \forall x\,(\varphi_{150}) \to \varphi_{33}) \to (\varphi_{150} \to \forall x\,(\varphi_{150})) \to \varphi_{238}$
\item $\varphi_{396} :\equiv \forall \delta\,(\varphi_{133} \to \neg (\delta \in \mathbb{R} \land 0 < \delta \land \varphi_{78})) \to \varphi_{359} \to (\varphi_{133} \to \forall \delta\,(\varphi_{133})) \to \varphi_{131}$
\item $\varphi_{397} :\equiv \forall \delta\,(\varphi_{134} \to \neg (\delta \in \mathbb{R} \land 0 < \delta \land \varphi_{78})) \to \varphi_{360} \to (\varphi_{134} \to \forall \delta\,(\varphi_{134})) \to \varphi_{132}$
\item $\varphi_{398} :\equiv (\mathrm{FLimR}(f,c,\alpha) \land \mathrm{FLimL}(f,c,\alpha) \to \varphi_{235}) \to \varphi_{391}$
\item $\varphi_{399} :\equiv (\mathrm{FLimR}(f,c,\alpha) \land \mathrm{FLimL}(f,c,\alpha) \to \varphi_{235}) \to \varphi_{392}$
\item $\varphi_{400} :\equiv (\mathrm{FLimR}(f,c,\alpha) \land \mathrm{FLimL}(f,c,\alpha) \to c \in \mathbb{R}) \to \varphi_{263} \to \varphi_{317} \to \varphi_{326} \to \varphi_{324}$
\item $\varphi_{401} :\equiv (\mathrm{FLimR}(f,c,\alpha) \land \mathrm{FLimL}(f,c,\alpha) \to c \in \mathbb{R}) \to \varphi_{264} \to \varphi_{318} \to \varphi_{327} \to \varphi_{325}$
\item $\varphi_{402} :\equiv (\forall r_{2}\,(\varphi_{118} \to \neg \varphi_{57}) \to \forall r_{2}\,(\varphi_{118}) \to \forall r_{2}\,(\neg \varphi_{57})) \to (\varphi_{118} \to \forall r_{2}\,(\varphi_{118})) \to \varphi_{378}$
\item $\varphi_{403} :\equiv \varphi_{244} \to \forall d\,(\varphi_{207} \to \neg (d \in \mathbb{R} \land 0 < d \land d \leq r_{1} \land d \leq r_{2})) \to \varphi_{388}$
\item $\varphi_{404} :\equiv \forall x\,(\varphi_{276}) \to (\forall x\,(\varphi_{276}) \to \forall x\,(\varphi_{241}) \to \varphi_{25}) \to (\varphi_{241} \to \forall x\,(\varphi_{241})) \to \varphi_{283}$
\item $\varphi_{405} :\equiv \varphi_{362} \to (\varphi_{174} \to \forall d_{1}\,(\varphi_{174})) \to \forall \delta\,(\varphi_{191}) \to (\forall \delta\,(\varphi_{191}) \to \varphi_{203}) \to \varphi_{199} \to \varphi_{173}$
\item $\varphi_{406} :\equiv \forall x\,(\varphi_{274} \to \varphi_{69}) \to (\forall x\,(\varphi_{274} \to \varphi_{69}) \to \forall x\,(\varphi_{274}) \to \varphi_{78}) \to (\varphi_{274} \to \forall x\,(\varphi_{274})) \to \varphi_{331}$
\item $\varphi_{407} :\equiv (\forall \delta\,(\varphi_{323} \to \neg (\delta \in \mathbb{R} \land 0 < \delta \land \delta \leq d_{1} \land \delta \leq d_{2})) \to \varphi_{343}) \to (\varphi_{323} \to \forall \delta\,(\varphi_{323})) \to \forall d\,(\varphi_{64}) \to (\forall d\,(\varphi_{64}) \to \varphi_{85}) \to \varphi_{90} \to \varphi_{322}$
\item $\varphi_{408} :\equiv \varphi_{376} \to (\varphi_{267} \to \forall d_{2}\,(\varphi_{267})) \to \forall \delta\,(\varphi_{192}) \to (\forall \delta\,(\varphi_{192}) \to \varphi_{204}) \to \varphi_{200} \to \varphi_{266}$
\item $\varphi_{409} :\equiv \varphi_{273} \to \forall \delta\,(\varphi_{323} \to \neg (\delta \in \mathbb{R} \land 0 < \delta \land \delta \leq d_{1} \land \delta \leq d_{2})) \to \varphi_{407}$
\end{abbrevlist}
\end{fullproof}

\section{第4章 2節: 連続性 (定義2.1, 2.2, 2.8, 定理2.3, 2.4, 2.6, 2.7, 2.9, 2.10)}
% 自動生成 (prover/build.js): ContInSeq  ノード数 248
\begin{fullproof}{[G1]}{$\forall f\,\forall A\,\forall c\,\forall a\,(\mathrm{ContIn}(f,A,c) \to \mathrm{Seq}(a) \to \forall n \in \mathbb{N}\,(a_{n} \in A) \to \mathrm{Conv}(a,c) \to \mathrm{Conv}(f \circ a,f(c)))$}{248}\phantomsection\label{tree:ContInSeq}
\noindent$\Pi_{1}$:
\begin{ptree}
\AXC{$z = w \to \lvert z - f(c) \rvert < \varepsilon \to \lvert w - f(c) \rvert < \varepsilon$\lbl{A6}}
\UIC{$\forall w\,(z = w \to \lvert z - f(c) \rvert < \varepsilon \to \lvert w - f(c) \rvert < \varepsilon)$}
\UIC{$\varphi_{5}$}
\AXC{$\varphi_{5} \to \varphi_{32}$\lbl{A2}}
\BIC{$\varphi_{32}$}
\AXC{$\varphi_{32} \to \varphi_{48}$\lbl{A2}}
\BIC{$\varphi_{48}$}
\end{ptree}
\noindent$\Pi_{2}$:
\begin{ptree}
\AXC{$\varphi_{33}$\lbl{A2}}
\AXC{$\mathrm{ContIn}(f,A,c) \to \mathrm{Seq}(a) \to \varphi_{94}$\lbl{A1}}
\AXC{$\varphi_{132}$\lbl{A1}}
\BIC{$\mathrm{ContIn}(f,A,c) \to \mathrm{Seq}(a) \to \varphi_{99}$}
\AXC{$\varphi_{138}$\lbl{A1}}
\BIC{$\varphi_{33} \to \mathrm{ContIn}(f,A,c) \to \mathrm{Seq}(a) \to \varphi_{96}$}
\BIC{$\mathrm{ContIn}(f,A,c) \to \mathrm{Seq}(a) \to \varphi_{96}$}
\AXC{$\varphi_{145}$\lbl{A1}}
\BIC{$(\mathrm{ContIn}(f,A,c) \to \mathrm{Seq}(a) \to \varphi_{81}) \to \varphi_{127}$}
\end{ptree}
\noindent$\Pi_{3}$:
\begin{ptree}
\AXC{$\mathrm{ContIn}(f,A,c) \to \mathrm{Seq}(a) \to \varphi_{81}$\lbl{A1}}
\AXC{$\stackrel{\Pi_{2}}{\vdots}$}\noLine
\UIC{$(\mathrm{ContIn}(f,A,c) \to \mathrm{Seq}(a) \to \varphi_{81}) \to \varphi_{127}$}
\BIC{$\varphi_{127}$}
\end{ptree}
\noindent$\Pi_{4}$:
\begin{ptree}
\AXC{$\forall n \in \mathbb{N}\,(a_{n} \in A) \to n \in \mathbb{N} \to a_{n} \in A$\lbl{A2}}
\AXC{$\varphi_{28}$\lbl{A1}}
\AXC{$\varphi_{28} \to (\forall n \in \mathbb{N}\,(a_{n} \in A) \to n \in \mathbb{N} \to a_{n} \in A) \to \varphi_{26}$\lbl{A1}}
\BIC{$(\forall n \in \mathbb{N}\,(a_{n} \in A) \to n \in \mathbb{N} \to a_{n} \in A) \to \varphi_{26}$}
\BIC{$\varphi_{26}$}
\end{ptree}
\noindent$\Pi_{5}$:
\begin{ptree}
\AXC{$\stackrel{\Pi_{4}}{\vdots}$}\noLine
\UIC{$\varphi_{26}$}
\AXC{$\varphi_{58}$\lbl{A2}}
\AXC{$\mathrm{ContIn}(f,A,c) \to \mathrm{Seq}(a) \to \varphi_{86}$\lbl{A1}}
\AXC{$\varphi_{131}$\lbl{A1}}
\BIC{$\mathrm{ContIn}(f,A,c) \to \mathrm{Seq}(a) \to \varphi_{107}$}
\AXC{$\varphi_{144}$\lbl{A1}}
\BIC{$\varphi_{58} \to \mathrm{ContIn}(f,A,c) \to \mathrm{Seq}(a) \to \varphi_{112}$}
\BIC{$\mathrm{ContIn}(f,A,c) \to \mathrm{Seq}(a) \to \varphi_{112}$}
\AXC{$\varphi_{146}$\lbl{A1}}
\BIC{$\varphi_{26} \to (\mathrm{ContIn}(f,A,c) \to \mathrm{Seq}(a) \to \varphi_{81}) \to \varphi_{130}$}
\BIC{$(\mathrm{ContIn}(f,A,c) \to \mathrm{Seq}(a) \to \varphi_{81}) \to \varphi_{130}$}
\end{ptree}
\noindent$\Pi_{6}$:
\begin{ptree}
\AXC{$\mathrm{ContIn}(f,A,c) \to \mathrm{Seq}(a) \to \varphi_{81}$\lbl{A1}}
\AXC{$\stackrel{\Pi_{5}}{\vdots}$}\noLine
\UIC{$(\mathrm{ContIn}(f,A,c) \to \mathrm{Seq}(a) \to \varphi_{81}) \to \varphi_{130}$}
\BIC{$\varphi_{130}$}
\end{ptree}
\noindent$\Pi_{7}$:
\begin{ptree}
\AXC{$x = x$\lbl{A5}}
\AXC{$x = y \to x = x \to y = x$\lbl{A6}}
\AXC{$(x = y \to x = x \to y = x) \to x = x \to x = y \to y = x$\lbl{A1}}
\BIC{$x = x \to x = y \to y = x$}
\BIC{$x = y \to y = x$}
\UIC{$\forall y\,(x = y \to y = x)$}
\UIC{$\forall x\,\forall y\,(x = y \to y = x)$}
\AXC{$\varphi_{8}$\lbl{A2}}
\BIC{$\forall y\,(f \circ a(n) = y \to y = f \circ a(n))$}
\AXC{$\varphi_{50}$\lbl{A2}}
\BIC{$f \circ a(n) = f(a_{n}) \to f(a_{n}) = f \circ a(n)$}
\end{ptree}
\noindent$\Pi_{8}$:
\begin{ptree}
\AXC{$\forall n \in \mathbb{N}\,(a_{n} \in A) \to n \in \mathbb{N} \to a_{n} \in A$\lbl{A2}}
\AXC{$\varphi_{28}$\lbl{A1}}
\AXC{$\varphi_{28} \to (\forall n \in \mathbb{N}\,(a_{n} \in A) \to n \in \mathbb{N} \to a_{n} \in A) \to \varphi_{26}$\lbl{A1}}
\BIC{$(\forall n \in \mathbb{N}\,(a_{n} \in A) \to n \in \mathbb{N} \to a_{n} \in A) \to \varphi_{26}$}
\BIC{$\varphi_{26}$}
\AXC{$\varphi_{26} \to \varphi_{72}$\lbl{A1}}
\BIC{$\varphi_{72}$}
\end{ptree}
\noindent$\Pi_{9}$:
\begin{ptree}
\AXC{$\forall f\,\forall A\,\forall c\,(\varphi_{69})$\lbl{DefContIn}}
\AXC{$\forall f\,\forall A\,\forall c\,(\varphi_{69}) \to \forall A\,\forall c\,(\varphi_{69})$\lbl{A2}}
\BIC{$\forall A\,\forall c\,(\varphi_{69})$}
\AXC{$\forall A\,\forall c\,(\varphi_{69}) \to \forall c\,(\varphi_{69})$\lbl{A2}}
\BIC{$\forall c\,(\varphi_{69})$}
\AXC{$\forall c\,(\varphi_{69}) \to \varphi_{69}$\lbl{A2}}
\BIC{$\varphi_{69}$}
\end{ptree}
\noindent$\Pi_{10}$:
\begin{ptree}
\AXC{$\stackrel{\Pi_{9}}{\vdots}$}\noLine
\UIC{$\varphi_{69}$}
\AXC{$\mathrm{ContIn}(f,A,c) \to \mathrm{ContIn}(f,A,c)$\lbl{A1}}
\AXC{$\varphi_{105}$\lbl{A1}}
\BIC{$\varphi_{69} \to \mathrm{ContIn}(f,A,c) \to \varphi_{10}$}
\BIC{$\mathrm{ContIn}(f,A,c) \to \varphi_{10}$}
\end{ptree}
\noindent$\Pi_{11}$:
\begin{ptree}
\AXC{$\forall X\,\forall Y\,\forall x\,(X \subseteq Y \to x \in X \to x \in Y)$\lbl{[B5]}}
\AXC{$\varphi_{16}$\lbl{A2}}
\BIC{$\forall Y\,\forall x\,(A \subseteq Y \to x \in A \to x \in Y)$}
\AXC{$\varphi_{19}$\lbl{A2}}
\BIC{$\forall x\,(A \subseteq \mathrm{dom}(f) \to x \in A \to x \in \mathrm{dom}(f))$}
\AXC{$\varphi_{35}$\lbl{A2}}
\BIC{$A \subseteq \mathrm{dom}(f) \to a_{n} \in A \to a_{n} \in \mathrm{dom}(f)$}
\AXC{$\varphi_{114}$\lbl{A1}}
\BIC{$(\mathrm{ContIn}(f,A,c) \to \varphi_{10}) \to \varphi_{73}$}
\end{ptree}
\noindent$\Pi_{12}$:
\begin{ptree}
\AXC{$\stackrel{\Pi_{10}}{\vdots}$}\noLine
\UIC{$\mathrm{ContIn}(f,A,c) \to \varphi_{10}$}
\AXC{$\stackrel{\Pi_{11}}{\vdots}$}\noLine
\UIC{$(\mathrm{ContIn}(f,A,c) \to \varphi_{10}) \to \varphi_{73}$}
\BIC{$\varphi_{73}$}
\end{ptree}
\noindent$\Pi_{13}$:
\begin{ptree}
\AXC{$\mathrm{ContIn}(f,A,c) \to \mathrm{Seq}(a) \to \varphi_{6}$\lbl{A1}}
\AXC{$\stackrel{\Pi_{8}}{\vdots}$}\noLine
\UIC{$\varphi_{72}$}
\BIC{$\mathrm{ContIn}(f,A,c) \to \mathrm{Seq}(a) \to \varphi_{11}$}
\AXC{$\stackrel{\Pi_{12}}{\vdots}$}\noLine
\UIC{$\varphi_{73}$}
\BIC{$\mathrm{ContIn}(f,A,c) \to \mathrm{Seq}(a) \to \varphi_{17}$}
\AXC{$\varphi_{74}$\lbl{A1}}
\BIC{$\varphi_{9} \to n \in \mathbb{N} \to a_{n} \in \mathrm{dom}(f)$}
\UIC{$\forall n\,(\varphi_{9} \to n \in \mathbb{N} \to a_{n} \in \mathrm{dom}(f))$}
\end{ptree}
\noindent$\Pi_{14}$:
\begin{ptree}
\AXC{$\varphi_{9} \to \forall n\,(\varphi_{9})$\lbl{A4}}
\AXC{$\varphi_{75}$\lbl{A3}}
\AXC{$\stackrel{\Pi_{13}}{\vdots}$}\noLine
\UIC{$\forall n\,(\varphi_{9} \to n \in \mathbb{N} \to a_{n} \in \mathrm{dom}(f))$}
\AXC{$\forall n\,(\varphi_{9} \to n \in \mathbb{N} \to a_{n} \in \mathrm{dom}(f)) \to \varphi_{75} \to \varphi_{82}$\lbl{A1}}
\BIC{$\varphi_{75} \to \varphi_{82}$}
\BIC{$\varphi_{82}$}
\BIC{$\mathrm{ContIn}(f,A,c) \to \mathrm{Seq}(a) \to \varphi_{20}$}
\end{ptree}
\noindent$\Pi_{15}$:
\begin{ptree}
\AXC{$\forall f\,\forall A\,\forall c\,(\varphi_{69})$\lbl{DefContIn}}
\AXC{$\forall f\,\forall A\,\forall c\,(\varphi_{69}) \to \forall A\,\forall c\,(\varphi_{69})$\lbl{A2}}
\BIC{$\forall A\,\forall c\,(\varphi_{69})$}
\AXC{$\forall A\,\forall c\,(\varphi_{69}) \to \forall c\,(\varphi_{69})$\lbl{A2}}
\BIC{$\forall c\,(\varphi_{69})$}
\AXC{$\forall c\,(\varphi_{69}) \to \varphi_{69}$\lbl{A2}}
\BIC{$\varphi_{69}$}
\end{ptree}
\noindent$\Pi_{16}$:
\begin{ptree}
\AXC{$\stackrel{\Pi_{15}}{\vdots}$}\noLine
\UIC{$\varphi_{69}$}
\AXC{$\mathrm{ContIn}(f,A,c) \to \mathrm{ContIn}(f,A,c)$\lbl{A1}}
\AXC{$\varphi_{98}$\lbl{A1}}
\BIC{$\varphi_{69} \to \mathrm{ContIn}(f,A,c) \to \varphi_{2}$}
\BIC{$\mathrm{ContIn}(f,A,c) \to \varphi_{2}$}
\end{ptree}
\noindent$\Pi_{17}$:
\begin{ptree}
\AXC{$\stackrel{\Pi_{16}}{\vdots}$}\noLine
\UIC{$\mathrm{ContIn}(f,A,c) \to \varphi_{2}$}
\AXC{$\forall f\,\forall a\,(\mathrm{RFn}(f) \to \mathrm{Seq}(a) \to \varphi_{39})$\lbl{[S8]}}
\AXC{$\varphi_{88}$\lbl{A2}}
\BIC{$\forall a\,(\mathrm{RFn}(f) \to \mathrm{Seq}(a) \to \varphi_{39})$}
\AXC{$\varphi_{87}$\lbl{A2}}
\BIC{$\mathrm{RFn}(f) \to \mathrm{Seq}(a) \to \varphi_{39}$}
\AXC{$\varphi_{123}$\lbl{A1}}
\BIC{$(\mathrm{ContIn}(f,A,c) \to \varphi_{2}) \to \varphi_{93}$}
\BIC{$\varphi_{93}$}
\end{ptree}
\noindent$\Pi_{18}$:
\begin{ptree}
\AXC{$\stackrel{\Pi_{14}}{\vdots}$}\noLine
\UIC{$\mathrm{ContIn}(f,A,c) \to \mathrm{Seq}(a) \to \varphi_{20}$}
\AXC{$\mathrm{ContIn}(f,A,c) \to \mathrm{Seq}(a) \to \mathrm{Seq}(a)$\lbl{A1}}
\AXC{$\stackrel{\Pi_{17}}{\vdots}$}\noLine
\UIC{$\varphi_{93}$}
\BIC{$(\mathrm{ContIn}(f,A,c) \to \mathrm{Seq}(a) \to \varphi_{20}) \to \varphi_{55}$}
\BIC{$\varphi_{55}$}
\end{ptree}
\noindent$\Pi_{19}$:
\begin{ptree}
\AXC{$\varphi_{36}$\lbl{A2}}
\AXC{$\stackrel{\Pi_{18}}{\vdots}$}\noLine
\UIC{$\varphi_{55}$}
\AXC{$\varphi_{55} \to \mathrm{ContIn}(f,A,c) \to \mathrm{Seq}(a) \to \varphi_{104}$\lbl{A1}}
\BIC{$\mathrm{ContIn}(f,A,c) \to \mathrm{Seq}(a) \to \varphi_{104}$}
\AXC{$\varphi_{139}$\lbl{A1}}
\BIC{$\varphi_{36} \to \mathrm{ContIn}(f,A,c) \to \mathrm{Seq}(a) \to \varphi_{103}$}
\BIC{$\mathrm{ContIn}(f,A,c) \to \mathrm{Seq}(a) \to \varphi_{103}$}
\end{ptree}
\noindent$\Pi_{20}$:
\begin{ptree}
\AXC{$\stackrel{\Pi_{19}}{\vdots}$}\noLine
\UIC{$\mathrm{ContIn}(f,A,c) \to \mathrm{Seq}(a) \to \varphi_{103}$}
\AXC{$(\mathrm{ContIn}(f,A,c) \to \mathrm{Seq}(a) \to \varphi_{103}) \to \varphi_{129}$\lbl{A1}}
\BIC{$\varphi_{129}$}
\end{ptree}
\noindent$\Pi_{21}$:
\begin{ptree}
\AXC{$\mathrm{ContIn}(f,A,c) \to \mathrm{Seq}(a) \to \varphi_{81}$\lbl{A1}}
\AXC{$\stackrel{\Pi_{20}}{\vdots}$}\noLine
\UIC{$\varphi_{129}$}
\BIC{$\mathrm{ContIn}(f,A,c) \to \mathrm{Seq}(a) \to \varphi_{90}$}
\AXC{$(\mathrm{ContIn}(f,A,c) \to \mathrm{Seq}(a) \to \varphi_{90}) \to \varphi_{118}$\lbl{A1}}
\BIC{$\varphi_{118}$}
\end{ptree}
\noindent$\Pi_{22}$:
\begin{ptree}
\AXC{$\stackrel{\Pi_{7}}{\vdots}$}\noLine
\UIC{$f \circ a(n) = f(a_{n}) \to f(a_{n}) = f \circ a(n)$}
\AXC{$\stackrel{\Pi_{21}}{\vdots}$}\noLine
\UIC{$\varphi_{118}$}
\BIC{$\mathrm{ContIn}(f,A,c) \to \mathrm{Seq}(a) \to \varphi_{91}$}
\AXC{$(\mathrm{ContIn}(f,A,c) \to \mathrm{Seq}(a) \to \varphi_{91}) \to \varphi_{140}$\lbl{A1}}
\BIC{$\varphi_{140}$}
\end{ptree}
\noindent$\Pi_{23}$:
\begin{ptree}
\AXC{$\stackrel{\Pi_{1}}{\vdots}$}\noLine
\UIC{$\varphi_{48}$}
\AXC{$\mathrm{ContIn}(f,A,c) \to \mathrm{Seq}(a) \to \varphi_{80}$\lbl{A1}}
\AXC{$\stackrel{\Pi_{3}}{\vdots}$}\noLine
\UIC{$\varphi_{127}$}
\BIC{$\mathrm{ContIn}(f,A,c) \to \mathrm{Seq}(a) \to \varphi_{89}$}
\AXC{$\stackrel{\Pi_{6}}{\vdots}$}\noLine
\UIC{$\varphi_{130}$}
\BIC{$\mathrm{ContIn}(f,A,c) \to \mathrm{Seq}(a) \to \varphi_{95}$}
\AXC{$\stackrel{\Pi_{22}}{\vdots}$}\noLine
\UIC{$\varphi_{140}$}
\BIC{$\varphi_{48} \to \mathrm{ContIn}(f,A,c) \to \mathrm{Seq}(a) \to \varphi_{97}$}
\BIC{$\mathrm{ContIn}(f,A,c) \to \mathrm{Seq}(a) \to \varphi_{97}$}
\AXC{$\varphi_{133}$\lbl{A1}}
\BIC{$\varphi_{83} \to n \in \mathbb{N} \to N \leq n \to \lvert f \circ a(n) - f(c) \rvert < \varepsilon$}
\UIC{$\forall n\,(\varphi_{83} \to n \in \mathbb{N} \to N \leq n \to \lvert f \circ a(n) - f(c) \rvert < \varepsilon)$}
\AXC{$\varphi_{149}$\lbl{A1}}
\BIC{$\varphi_{134} \to \varphi_{141}$}
\end{ptree}
\noindent$\Pi_{24}$:
\begin{ptree}
\AXC{$\mathrm{ContIn}(f,A,c) \to \mathrm{Seq}(a) \to \varphi_{94}$\lbl{A1}}
\AXC{$\varphi_{128}$\lbl{A1}}
\BIC{$\mathrm{ContIn}(f,A,c) \to \mathrm{Seq}(a) \to \varphi_{76}$}
\AXC{$(\mathrm{ContIn}(f,A,c) \to \mathrm{Seq}(a) \to \varphi_{76}) \to \varphi_{136}$\lbl{A1}}
\BIC{$\varphi_{136}$}
\end{ptree}
\noindent$\Pi_{25}$:
\begin{ptree}
\AXC{$\forall N\,(\neg \varphi_{18}) \to \neg \varphi_{18}$\lbl{A2}}
\AXC{$\varphi_{83} \to \forall n\,(\varphi_{83})$\lbl{A4}}
\AXC{$\varphi_{134}$\lbl{A3}}
\AXC{$\stackrel{\Pi_{23}}{\vdots}$}\noLine
\UIC{$\varphi_{134} \to \varphi_{141}$}
\BIC{$\varphi_{141}$}
\BIC{$\mathrm{ContIn}(f,A,c) \to \mathrm{Seq}(a) \to \varphi_{102}$}
\AXC{$\stackrel{\Pi_{24}}{\vdots}$}\noLine
\UIC{$\varphi_{136}$}
\BIC{$\mathrm{ContIn}(f,A,c) \to \varphi_{108}$}
\AXC{$\varphi_{142}$\lbl{A1}}
\BIC{$(\forall N\,(\neg \varphi_{18}) \to \neg \varphi_{18}) \to \mathrm{ContIn}(f,A,c) \to \varphi_{109}$}
\BIC{$\mathrm{ContIn}(f,A,c) \to \varphi_{109}$}
\end{ptree}
\noindent$\Pi_{26}$:
\begin{ptree}
\AXC{$\forall a\,(\varphi_{41})$\lbl{[C3]}}
\AXC{$\forall a\,(\varphi_{41}) \to \varphi_{41}$\lbl{A2}}
\BIC{$\varphi_{41}$}
\AXC{$\varphi_{41} \to \varphi_{38}$\lbl{A2}}
\BIC{$\varphi_{38}$}
\AXC{$\varphi_{38} \to \mathrm{Conv}(a,c) \to \varphi_{27}$\lbl{A2}}
\BIC{$\mathrm{Conv}(a,c) \to \varphi_{27}$}
\AXC{$\varphi_{143}$\lbl{A1}}
\BIC{$(\mathrm{ContIn}(f,A,c) \to \mathrm{Seq}(a) \to \varphi_{1}) \to \varphi_{137}$}
\end{ptree}
\noindent$\Pi_{27}$:
\begin{ptree}
\AXC{$\mathrm{ContIn}(f,A,c) \to \mathrm{Seq}(a) \to \varphi_{1}$\lbl{A1}}
\AXC{$\stackrel{\Pi_{26}}{\vdots}$}\noLine
\UIC{$(\mathrm{ContIn}(f,A,c) \to \mathrm{Seq}(a) \to \varphi_{1}) \to \varphi_{137}$}
\BIC{$\varphi_{137}$}
\end{ptree}
\noindent$\Pi_{28}$:
\begin{ptree}
\AXC{$\mathrm{ContIn}(f,A,c) \to \mathrm{Seq}(a) \to \varphi_{86}$\lbl{A1}}
\AXC{$\varphi_{122}$\lbl{A1}}
\BIC{$\mathrm{ContIn}(f,A,c) \to \mathrm{Seq}(a) \to \varphi_{68}$}
\AXC{$\stackrel{\Pi_{27}}{\vdots}$}\noLine
\UIC{$\varphi_{137}$}
\BIC{$(\mathrm{ContIn}(f,A,c) \to \mathrm{Seq}(a) \to \varphi_{66}) \to \varphi_{85}$}
\end{ptree}
\noindent$\Pi_{29}$:
\begin{ptree}
\AXC{$\mathrm{ContIn}(f,A,c) \to \mathrm{Seq}(a) \to \varphi_{86}$\lbl{A1}}
\AXC{$\varphi_{121}$\lbl{A1}}
\BIC{$\mathrm{ContIn}(f,A,c) \to \mathrm{Seq}(a) \to \varphi_{66}$}
\AXC{$\stackrel{\Pi_{28}}{\vdots}$}\noLine
\UIC{$(\mathrm{ContIn}(f,A,c) \to \mathrm{Seq}(a) \to \varphi_{66}) \to \varphi_{85}$}
\BIC{$\varphi_{85}$}
\end{ptree}
\noindent$\Pi_{30}$:
\begin{ptree}
\AXC{$\stackrel{\Pi_{25}}{\vdots}$}\noLine
\UIC{$\mathrm{ContIn}(f,A,c) \to \varphi_{109}$}
\AXC{$(\mathrm{ContIn}(f,A,c) \to \varphi_{109}) \to \varphi_{113}$\lbl{A1}}
\BIC{$\varphi_{113}$}
\UIC{$\forall N\,(\varphi_{113})$}
\AXC{$\stackrel{\Pi_{29}}{\vdots}$}\noLine
\UIC{$\varphi_{85}$}
\AXC{$\varphi_{85} \to \forall N\,(\varphi_{113}) \to \varphi_{148}$\lbl{A1}}
\BIC{$\forall N\,(\varphi_{113}) \to \varphi_{148}$}
\BIC{$\varphi_{148}$}
\end{ptree}
\noindent$\Pi_{31}$:
\begin{ptree}
\AXC{$\forall f\,\forall A\,\forall c\,(\varphi_{69})$\lbl{DefContIn}}
\AXC{$\forall f\,\forall A\,\forall c\,(\varphi_{69}) \to \forall A\,\forall c\,(\varphi_{69})$\lbl{A2}}
\BIC{$\forall A\,\forall c\,(\varphi_{69})$}
\AXC{$\forall A\,\forall c\,(\varphi_{69}) \to \forall c\,(\varphi_{69})$\lbl{A2}}
\BIC{$\forall c\,(\varphi_{69})$}
\AXC{$\forall c\,(\varphi_{69}) \to \varphi_{69}$\lbl{A2}}
\BIC{$\varphi_{69}$}
\end{ptree}
\noindent$\Pi_{32}$:
\begin{ptree}
\AXC{$\varphi_{84}$\lbl{A2}}
\AXC{$\stackrel{\Pi_{31}}{\vdots}$}\noLine
\UIC{$\varphi_{69}$}
\AXC{$\mathrm{ContIn}(f,A,c) \to \mathrm{ContIn}(f,A,c)$\lbl{A1}}
\AXC{$\varphi_{119}$\lbl{A1}}
\BIC{$\varphi_{69} \to \mathrm{ContIn}(f,A,c) \to \mathrm{Seq}(a) \to \varphi_{65}$}
\BIC{$\mathrm{ContIn}(f,A,c) \to \mathrm{Seq}(a) \to \varphi_{65}$}
\AXC{$\varphi_{135}$\lbl{A1}}
\BIC{$\varphi_{84} \to \mathrm{ContIn}(f,A,c) \to \mathrm{Seq}(a) \to \varphi_{63}$}
\BIC{$\mathrm{ContIn}(f,A,c) \to \mathrm{Seq}(a) \to \varphi_{63}$}
\end{ptree}
\noindent$\Pi_{33}$:
\begin{ptree}
\AXC{$\stackrel{\Pi_{32}}{\vdots}$}\noLine
\UIC{$\mathrm{ContIn}(f,A,c) \to \mathrm{Seq}(a) \to \varphi_{63}$}
\AXC{$(\mathrm{ContIn}(f,A,c) \to \mathrm{Seq}(a) \to \varphi_{63}) \to \varphi_{147}$\lbl{A1}}
\BIC{$\varphi_{147}$}
\end{ptree}
\noindent$\Pi_{34}$:
\begin{ptree}
\AXC{$\varphi_{126}$\lbl{A3}}
\AXC{$\varphi_{125}$\lbl{A4}}
\AXC{$\forall N\,(\varphi_{113}) \to \forall N\,(\varphi_{9} \land \varepsilon \in \mathbb{R} \land 0 < \varepsilon \land \varphi_{30} \land \neg \varphi_{21}) \to \varphi_{4}$\lbl{A3}}
\AXC{$\stackrel{\Pi_{30}}{\vdots}$}\noLine
\UIC{$\varphi_{148}$}
\BIC{$\varphi_{125} \to \mathrm{ContIn}(f,A,c) \to \mathrm{Seq}(a) \to \varphi_{78}$}
\BIC{$\mathrm{ContIn}(f,A,c) \to \mathrm{Seq}(a) \to \varphi_{78}$}
\AXC{$\varphi_{124}$\lbl{A1}}
\BIC{$\varphi_{9} \land \varepsilon \in \mathbb{R} \land 0 < \varepsilon \land \neg \varphi_{21} \to \neg \varphi_{30}$}
\UIC{$\forall \delta\,(\varphi_{9} \land \varepsilon \in \mathbb{R} \land 0 < \varepsilon \land \neg \varphi_{21} \to \neg \varphi_{30})$}
\AXC{$\stackrel{\Pi_{33}}{\vdots}$}\noLine
\UIC{$\varphi_{147}$}
\BIC{$\varphi_{126} \to \varphi_{116} \to \mathrm{ContIn}(f,A,c) \to \mathrm{Seq}(a) \to \varphi_{57}$}
\BIC{$\varphi_{116} \to \mathrm{ContIn}(f,A,c) \to \mathrm{Seq}(a) \to \varphi_{57}$}
\end{ptree}
\noindent$\Pi_{35}$:
\begin{ptree}
\AXC{$\varphi_{116}$\lbl{A4}}
\AXC{$\stackrel{\Pi_{34}}{\vdots}$}\noLine
\UIC{$\varphi_{116} \to \mathrm{ContIn}(f,A,c) \to \mathrm{Seq}(a) \to \varphi_{57}$}
\BIC{$\mathrm{ContIn}(f,A,c) \to \mathrm{Seq}(a) \to \varphi_{57}$}
\AXC{$\varphi_{115}$\lbl{A1}}
\BIC{$\varphi_{9} \to \varepsilon \in \mathbb{R} \to 0 < \varepsilon \to \varphi_{21}$}
\UIC{$\forall \varepsilon\,(\varphi_{9} \to \varepsilon \in \mathbb{R} \to 0 < \varepsilon \to \varphi_{21})$}
\AXC{$\forall \varepsilon\,(\varphi_{9} \to \varepsilon \in \mathbb{R} \to 0 < \varepsilon \to \varphi_{21}) \to \varphi_{117} \to \varphi_{111}$\lbl{A1}}
\BIC{$\varphi_{117} \to \varphi_{111}$}
\end{ptree}
\noindent$\Pi_{36}$:
\begin{ptree}
\AXC{$\forall f\,\forall A\,\forall c\,(\varphi_{69})$\lbl{DefContIn}}
\AXC{$\forall f\,\forall A\,\forall c\,(\varphi_{69}) \to \forall A\,\forall c\,(\varphi_{69})$\lbl{A2}}
\BIC{$\forall A\,\forall c\,(\varphi_{69})$}
\AXC{$\forall A\,\forall c\,(\varphi_{69}) \to \forall c\,(\varphi_{69})$\lbl{A2}}
\BIC{$\forall c\,(\varphi_{69})$}
\AXC{$\forall c\,(\varphi_{69}) \to \varphi_{69}$\lbl{A2}}
\BIC{$\varphi_{69}$}
\end{ptree}
\noindent$\Pi_{37}$:
\begin{ptree}
\AXC{$\stackrel{\Pi_{36}}{\vdots}$}\noLine
\UIC{$\varphi_{69}$}
\AXC{$\mathrm{ContIn}(f,A,c) \to \mathrm{ContIn}(f,A,c)$\lbl{A1}}
\AXC{$\varphi_{100}$\lbl{A1}}
\BIC{$\varphi_{69} \to \mathrm{ContIn}(f,A,c) \to \varphi_{3}$}
\BIC{$\mathrm{ContIn}(f,A,c) \to \varphi_{3}$}
\end{ptree}
\noindent$\Pi_{38}$:
\begin{ptree}
\AXC{$\forall X\,\forall Y\,\forall x\,(X \subseteq Y \to x \in X \to x \in Y)$\lbl{[B5]}}
\AXC{$\varphi_{16}$\lbl{A2}}
\BIC{$\forall Y\,\forall x\,(A \subseteq Y \to x \in A \to x \in Y)$}
\AXC{$\varphi_{19}$\lbl{A2}}
\BIC{$\forall x\,(A \subseteq \mathrm{dom}(f) \to x \in A \to x \in \mathrm{dom}(f))$}
\AXC{$\varphi_{24}$\lbl{A2}}
\BIC{$A \subseteq \mathrm{dom}(f) \to c \in A \to c \in \mathrm{dom}(f)$}
\AXC{$\varphi_{101}$\lbl{A1}}
\BIC{$(\mathrm{ContIn}(f,A,c) \to \varphi_{10}) \to \varphi_{60}$}
\end{ptree}
\noindent$\Pi_{39}$:
\begin{ptree}
\AXC{$\stackrel{\Pi_{37}}{\vdots}$}\noLine
\UIC{$\mathrm{ContIn}(f,A,c) \to \varphi_{3}$}
\AXC{$\stackrel{\Pi_{10}}{\vdots}$}\noLine
\UIC{$\mathrm{ContIn}(f,A,c) \to \varphi_{10}$}
\AXC{$\stackrel{\Pi_{38}}{\vdots}$}\noLine
\UIC{$(\mathrm{ContIn}(f,A,c) \to \varphi_{10}) \to \varphi_{60}$}
\BIC{$\varphi_{60}$}
\BIC{$\mathrm{ContIn}(f,A,c) \to \varphi_{7}$}
\end{ptree}
\noindent$\Pi_{40}$:
\begin{ptree}
\AXC{$\forall f\,\forall x\,(\mathrm{RFn}(f) \to x \in \mathrm{dom}(f) \to f(x) \in \mathbb{R})$\lbl{[S6]}}
\AXC{$\varphi_{34}$\lbl{A2}}
\BIC{$\forall x\,(\mathrm{RFn}(f) \to x \in \mathrm{dom}(f) \to f(x) \in \mathbb{R})$}
\AXC{$\varphi_{29}$\lbl{A2}}
\BIC{$\mathrm{RFn}(f) \to c \in \mathrm{dom}(f) \to f(c) \in \mathbb{R}$}
\AXC{$\varphi_{106}$\lbl{A1}}
\BIC{$(\mathrm{ContIn}(f,A,c) \to \varphi_{2}) \to \varphi_{67}$}
\end{ptree}
\noindent$\Pi_{41}$:
\begin{ptree}
\AXC{$\stackrel{\Pi_{39}}{\vdots}$}\noLine
\UIC{$\mathrm{ContIn}(f,A,c) \to \varphi_{7}$}
\AXC{$\stackrel{\Pi_{16}}{\vdots}$}\noLine
\UIC{$\mathrm{ContIn}(f,A,c) \to \varphi_{2}$}
\AXC{$\stackrel{\Pi_{40}}{\vdots}$}\noLine
\UIC{$(\mathrm{ContIn}(f,A,c) \to \varphi_{2}) \to \varphi_{67}$}
\BIC{$\varphi_{67}$}
\BIC{$\mathrm{ContIn}(f,A,c) \to \varphi_{12}$}
\end{ptree}
\noindent$\Pi_{42}$:
\begin{ptree}
\AXC{$\stackrel{\Pi_{18}}{\vdots}$}\noLine
\UIC{$\varphi_{55}$}
\AXC{$\varphi_{55} \to \mathrm{ContIn}(f,A,c) \to \varphi_{13}$\lbl{A1}}
\BIC{$\mathrm{ContIn}(f,A,c) \to \varphi_{13}$}
\AXC{$\varphi_{54}$\lbl{DefConv}}
\AXC{$\varphi_{54} \to \varphi_{64}$\lbl{A2}}
\BIC{$\varphi_{64}$}
\AXC{$\varphi_{64} \to \varphi_{71}$\lbl{A2}}
\BIC{$\varphi_{71}$}
\AXC{$\varphi_{71} \to \varphi_{120}$\lbl{A1}}
\BIC{$\varphi_{120}$}
\BIC{$(\mathrm{ContIn}(f,A,c) \to \varphi_{12}) \to \varphi_{92}$}
\end{ptree}
\noindent 主木:
\begin{ptree}
\AXC{$\varphi_{9} \to \forall \varepsilon\,(\varphi_{9})$\lbl{A4}}
\AXC{$\varphi_{117}$\lbl{A3}}
\AXC{$\stackrel{\Pi_{35}}{\vdots}$}\noLine
\UIC{$\varphi_{117} \to \varphi_{111}$}
\BIC{$\varphi_{111}$}
\BIC{$\mathrm{ContIn}(f,A,c) \to \mathrm{Seq}(a) \to \varphi_{59}$}
\AXC{$\stackrel{\Pi_{41}}{\vdots}$}\noLine
\UIC{$\mathrm{ContIn}(f,A,c) \to \varphi_{12}$}
\AXC{$\stackrel{\Pi_{42}}{\vdots}$}\noLine
\UIC{$(\mathrm{ContIn}(f,A,c) \to \varphi_{12}) \to \varphi_{92}$}
\BIC{$\varphi_{92}$}
\BIC{$\mathrm{ContIn}(f,A,c) \to \mathrm{Seq}(a) \to \varphi_{15}$}
\UIC{$\forall a\,(\mathrm{ContIn}(f,A,c) \to \mathrm{Seq}(a) \to \varphi_{15})$}
\UIC{$\forall c\,\forall a\,(\mathrm{ContIn}(f,A,c) \to \mathrm{Seq}(a) \to \varphi_{15})$}
\UIC{$\forall A\,\forall c\,\forall a\,(\mathrm{ContIn}(f,A,c) \to \mathrm{Seq}(a) \to \varphi_{15})$}
\UIC{$\varphi_{40}$}
\end{ptree}
\begin{abbrevlist}
\item $\varphi_{1} :\equiv \forall n \in \mathbb{N}\,(a_{n} \in A) \to \mathrm{Conv}(a,c) \to \mathrm{Conv}(a,c)$
\item $\varphi_{2} :\equiv \mathrm{Seq}(a) \to \forall n \in \mathbb{N}\,(a_{n} \in A) \to \mathrm{Conv}(a,c) \to \mathrm{RFn}(f)$
\item $\varphi_{3} :\equiv \mathrm{Seq}(a) \to \forall n \in \mathbb{N}\,(a_{n} \in A) \to \mathrm{Conv}(a,c) \to c \in A$
\item $\varphi_{4} :\equiv \forall N\,(\neg (N \in \mathbb{N} \land \forall n \in \mathbb{N}\,(N \leq n \to \lvert a_{n} - c \rvert < \delta)))$
\item $\varphi_{5} :\equiv \forall z\,\forall w\,(z = w \to \lvert z - f(c) \rvert < \varepsilon \to \lvert w - f(c) \rvert < \varepsilon)$
\item $\varphi_{6} :\equiv \forall n \in \mathbb{N}\,(a_{n} \in A) \to \mathrm{Conv}(a,c) \to n \in \mathbb{N} \to n \in \mathbb{N}$
\item $\varphi_{7} :\equiv \mathrm{Seq}(a) \to \forall n \in \mathbb{N}\,(a_{n} \in A) \to \mathrm{Conv}(a,c) \to c \in \mathrm{dom}(f)$
\item $\varphi_{8} :\equiv \forall x\,\forall y\,(x = y \to y = x) \to \forall y\,(f \circ a(n) = y \to y = f \circ a(n))$
\item $\varphi_{9} :\equiv \mathrm{ContIn}(f,A,c) \land \mathrm{Seq}(a) \land \forall n \in \mathbb{N}\,(a_{n} \in A) \land \mathrm{Conv}(a,c)$
\item $\varphi_{10} :\equiv \mathrm{Seq}(a) \to \forall n \in \mathbb{N}\,(a_{n} \in A) \to \mathrm{Conv}(a,c) \to A \subseteq \mathrm{dom}(f)$
\item $\varphi_{11} :\equiv \forall n \in \mathbb{N}\,(a_{n} \in A) \to \mathrm{Conv}(a,c) \to n \in \mathbb{N} \to a_{n} \in A$
\item $\varphi_{12} :\equiv \mathrm{Seq}(a) \to \forall n \in \mathbb{N}\,(a_{n} \in A) \to \mathrm{Conv}(a,c) \to f(c) \in \mathbb{R}$
\item $\varphi_{13} :\equiv \mathrm{Seq}(a) \to \forall n \in \mathbb{N}\,(a_{n} \in A) \to \mathrm{Conv}(a,c) \to \mathrm{Seq}(f \circ a)$
\item $\varphi_{14} :\equiv N \leq n \to \forall x \in A\,(\lvert x - c \rvert < \delta \to \lvert f(x) - f(c) \rvert < \varepsilon)$
\item $\varphi_{15} :\equiv \forall n \in \mathbb{N}\,(a_{n} \in A) \to \mathrm{Conv}(a,c) \to \mathrm{Conv}(f \circ a,f(c))$
\item $\varphi_{16} :\equiv \forall X\,\forall Y\,\forall x\,(X \subseteq Y \to x \in X \to x \in Y) \to \forall Y\,\forall x\,(A \subseteq Y \to x \in A \to x \in Y)$
\item $\varphi_{17} :\equiv \forall n \in \mathbb{N}\,(a_{n} \in A) \to \mathrm{Conv}(a,c) \to n \in \mathbb{N} \to a_{n} \in \mathrm{dom}(f)$
\item $\varphi_{18} :\equiv N \in \mathbb{N} \land \forall n \in \mathbb{N}\,(N \leq n \to \lvert f \circ a(n) - f(c) \rvert < \varepsilon)$
\item $\varphi_{19} :\equiv \forall Y\,\forall x\,(A \subseteq Y \to x \in A \to x \in Y) \to \forall x\,(A \subseteq \mathrm{dom}(f) \to x \in A \to x \in \mathrm{dom}(f))$
\item $\varphi_{20} :\equiv \forall n \in \mathbb{N}\,(a_{n} \in A) \to \mathrm{Conv}(a,c) \to \forall n \in \mathbb{N}\,(a_{n} \in \mathrm{dom}(f))$
\item $\varphi_{21} :\equiv \exists N \in \mathbb{N}\,\forall n \in \mathbb{N}\,(N \leq n \to \lvert f \circ a(n) - f(c) \rvert < \varepsilon)$
\item $\varphi_{22} :\equiv N \in \mathbb{N} \land \forall n \in \mathbb{N}\,(N \leq n \to \lvert a_{n} - c \rvert < \delta) \to n \in \mathbb{N} \to N \leq n \to N \leq n$
\item $\varphi_{23} :\equiv N \in \mathbb{N} \land \forall n \in \mathbb{N}\,(N \leq n \to \lvert a_{n} - c \rvert < \delta) \to n \in \mathbb{N} \to n \in \mathbb{N}$
\item $\varphi_{24} :\equiv \forall x\,(A \subseteq \mathrm{dom}(f) \to x \in A \to x \in \mathrm{dom}(f)) \to A \subseteq \mathrm{dom}(f) \to c \in A \to c \in \mathrm{dom}(f)$
\item $\varphi_{25} :\equiv \mathrm{Conv}(a,c) \to \mathrm{Seq}(f \circ a) \land \forall n \in \mathbb{N}\,(f \circ a(n) = f(a_{n}))$
\item $\varphi_{26} :\equiv \mathrm{ContIn}(f,A,c) \to \mathrm{Seq}(a) \to \forall n \in \mathbb{N}\,(a_{n} \in A) \to n \in \mathbb{N} \to a_{n} \in A$
\item $\varphi_{27} :\equiv \delta \in \mathbb{R} \to 0 < \delta \to \exists N \in \mathbb{N}\,\forall n \in \mathbb{N}\,(N \leq n \to \lvert a_{n} - c \rvert < \delta)$
\item $\varphi_{28} :\equiv \mathrm{ContIn}(f,A,c) \to \mathrm{Seq}(a) \to \forall n \in \mathbb{N}\,(a_{n} \in A) \to \forall n \in \mathbb{N}\,(a_{n} \in A)$
\item $\varphi_{29} :\equiv \forall x\,(\mathrm{RFn}(f) \to x \in \mathrm{dom}(f) \to f(x) \in \mathbb{R}) \to \mathrm{RFn}(f) \to c \in \mathrm{dom}(f) \to f(c) \in \mathbb{R}$
\item $\varphi_{30} :\equiv \delta \in \mathbb{R} \land 0 < \delta \land \forall x \in A\,(\lvert x - c \rvert < \delta \to \lvert f(x) - f(c) \rvert < \varepsilon)$
\item $\varphi_{31} :\equiv N \leq n \to a_{n} \in A \to \lvert a_{n} - c \rvert < \delta \to \lvert f(a_{n}) - f(c) \rvert < \varepsilon$
\item $\varphi_{32} :\equiv \forall w\,(f(a_{n}) = w \to \lvert f(a_{n}) - f(c) \rvert < \varepsilon \to \lvert w - f(c) \rvert < \varepsilon)$
\item $\varphi_{33} :\equiv \forall n \in \mathbb{N}\,(N \leq n \to \lvert a_{n} - c \rvert < \delta) \to n \in \mathbb{N} \to N \leq n \to \lvert a_{n} - c \rvert < \delta$
\item $\varphi_{34} :\equiv \forall f\,\forall x\,(\mathrm{RFn}(f) \to x \in \mathrm{dom}(f) \to f(x) \in \mathbb{R}) \to \forall x\,(\mathrm{RFn}(f) \to x \in \mathrm{dom}(f) \to f(x) \in \mathbb{R})$
\item $\varphi_{35} :\equiv \forall x\,(A \subseteq \mathrm{dom}(f) \to x \in A \to x \in \mathrm{dom}(f)) \to A \subseteq \mathrm{dom}(f) \to a_{n} \in A \to a_{n} \in \mathrm{dom}(f)$
\item $\varphi_{36} :\equiv \forall n \in \mathbb{N}\,(f \circ a(n) = f(a_{n})) \to n \in \mathbb{N} \to f \circ a(n) = f(a_{n})$
\item $\varphi_{37} :\equiv N \in \mathbb{N} \land \forall n \in \mathbb{N}\,(N \leq n \to \lvert a_{n} - c \rvert < \delta) \to n \in \mathbb{N} \to N \leq n \to \lvert a_{n} - c \rvert < \delta$
\item $\varphi_{38} :\equiv \forall \varepsilon\,(\mathrm{Conv}(a,c) \to \varepsilon \in \mathbb{R} \to 0 < \varepsilon \to \exists N \in \mathbb{N}\,\forall n \in \mathbb{N}\,(N \leq n \to \lvert a_{n} - c \rvert < \varepsilon))$
\item $\varphi_{39} :\equiv \forall n \in \mathbb{N}\,(a_{n} \in \mathrm{dom}(f)) \to \mathrm{Seq}(f \circ a) \land \forall n \in \mathbb{N}\,(f \circ a(n) = f(a_{n}))$
\item $\varphi_{40} :\equiv \forall f\,\forall A\,\forall c\,\forall a\,(\mathrm{ContIn}(f,A,c) \to \mathrm{Seq}(a) \to \varphi_{15})$
\item $\varphi_{41} :\equiv \forall \alpha\,\forall \varepsilon\,(\mathrm{Conv}(a,\alpha) \to \varepsilon \in \mathbb{R} \to 0 < \varepsilon \to \exists N \in \mathbb{N}\,\forall n \in \mathbb{N}\,(N \leq n \to \lvert a_{n} - \alpha \rvert < \varepsilon))$
\item $\varphi_{42} :\equiv N \in \mathbb{N} \land \forall n \in \mathbb{N}\,(N \leq n \to \lvert a_{n} - c \rvert < \delta) \to n \in \mathbb{N} \to N \leq n \to f \circ a(n) = f(a_{n})$
\item $\varphi_{43} :\equiv N \in \mathbb{N} \land \forall n \in \mathbb{N}\,(N \leq n \to \lvert a_{n} - c \rvert < \delta) \to n \in \mathbb{N} \to N \leq n \to f(a_{n}) = f \circ a(n)$
\item $\varphi_{44} :\equiv N \in \mathbb{N} \land \forall n \in \mathbb{N}\,(N \leq n \to \lvert a_{n} - c \rvert < \delta) \to N \in \mathbb{N} \land \forall n \in \mathbb{N}\,(N \leq n \to \lvert a_{n} - c \rvert < \delta)$
\item $\varphi_{45} :\equiv N \in \mathbb{N} \land \forall n \in \mathbb{N}\,(N \leq n \to \lvert a_{n} - c \rvert < \delta) \to n \in \mathbb{N} \to N \leq n \to \lvert f(a_{n}) - f(c) \rvert < \varepsilon$
\item $\varphi_{46} :\equiv N \in \mathbb{N} \land \forall n \in \mathbb{N}\,(N \leq n \to \lvert a_{n} - c \rvert < \delta) \to n \in \mathbb{N} \to N \leq n \to n \in \mathbb{N} \to N \leq n \to \lvert a_{n} - c \rvert < \delta$
\item $\varphi_{47} :\equiv N \in \mathbb{N} \land \forall n \in \mathbb{N}\,(N \leq n \to \lvert a_{n} - c \rvert < \delta) \to n \in \mathbb{N} \to N \leq n \to \lvert f \circ a(n) - f(c) \rvert < \varepsilon$
\item $\varphi_{48} :\equiv f(a_{n}) = f \circ a(n) \to \lvert f(a_{n}) - f(c) \rvert < \varepsilon \to \lvert f \circ a(n) - f(c) \rvert < \varepsilon$
\item $\varphi_{49} :\equiv N \in \mathbb{N} \land \forall n \in \mathbb{N}\,(N \leq n \to \lvert a_{n} - c \rvert < \delta) \to n \in \mathbb{N} \to N \leq n \to \forall n \in \mathbb{N}\,(N \leq n \to \lvert a_{n} - c \rvert < \delta)$
\item $\varphi_{50} :\equiv \forall y\,(f \circ a(n) = y \to y = f \circ a(n)) \to f \circ a(n) = f(a_{n}) \to f(a_{n}) = f \circ a(n)$
\item $\varphi_{51} :\equiv N \in \mathbb{N} \land \forall n \in \mathbb{N}\,(N \leq n \to \lvert a_{n} - c \rvert < \delta) \to \forall n \in \mathbb{N}\,(N \leq n \to \lvert f \circ a(n) - f(c) \rvert < \varepsilon)$
\item $\varphi_{52} :\equiv N \in \mathbb{N} \land \forall n \in \mathbb{N}\,(N \leq n \to \lvert a_{n} - c \rvert < \delta) \to n \in \mathbb{N} \to N \leq n \to n \in \mathbb{N} \to f \circ a(n) = f(a_{n})$
\item $\varphi_{53} :\equiv N \in \mathbb{N} \land \forall n \in \mathbb{N}\,(N \leq n \to \lvert a_{n} - c \rvert < \delta) \to n \in \mathbb{N} \to N \leq n \to \forall n \in \mathbb{N}\,(f \circ a(n) = f(a_{n}))$
\item $\varphi_{54} :\equiv \forall a\,\forall \alpha\,(\mathrm{Conv}(a,\alpha) \leftrightarrow \mathrm{Seq}(a) \land \alpha \in \mathbb{R} \land \forall \varepsilon \in \mathbb{R}\,(0 < \varepsilon \to \exists N \in \mathbb{N}\,\forall n \in \mathbb{N}\,(N \leq n \to \lvert a_{n} - \alpha \rvert < \varepsilon)))$
\item $\varphi_{55} :\equiv \mathrm{ContIn}(f,A,c) \to \mathrm{Seq}(a) \to \forall n \in \mathbb{N}\,(a_{n} \in A) \to \varphi_{25}$
\item $\varphi_{56} :\equiv N \in \mathbb{N} \land \forall n \in \mathbb{N}\,(N \leq n \to \lvert a_{n} - c \rvert < \delta) \to n \in \mathbb{N} \to \varphi_{14}$
\item $\varphi_{57} :\equiv \forall n \in \mathbb{N}\,(a_{n} \in A) \to \mathrm{Conv}(a,c) \to \varepsilon \in \mathbb{R} \to 0 < \varepsilon \to \varphi_{21}$
\item $\varphi_{58} :\equiv \forall x \in A\,(\lvert x - c \rvert < \delta \to \lvert f(x) - f(c) \rvert < \varepsilon) \to a_{n} \in A \to \lvert a_{n} - c \rvert < \delta \to \lvert f(a_{n}) - f(c) \rvert < \varepsilon$
\item $\varphi_{59} :\equiv \forall n \in \mathbb{N}\,(a_{n} \in A) \to \mathrm{Conv}(a,c) \to \forall \varepsilon \in \mathbb{R}\,(0 < \varepsilon \to \varphi_{21})$
\item $\varphi_{60} :\equiv (\mathrm{ContIn}(f,A,c) \to \varphi_{3}) \to \mathrm{ContIn}(f,A,c) \to \varphi_{7}$
\item $\varphi_{61} :\equiv N \in \mathbb{N} \land \forall n \in \mathbb{N}\,(N \leq n \to \lvert a_{n} - c \rvert < \delta) \to n \in \mathbb{N} \to \varphi_{31}$
\item $\varphi_{62} :\equiv \varphi_{30} \to N \in \mathbb{N} \land \forall n \in \mathbb{N}\,(N \leq n \to \lvert a_{n} - c \rvert < \delta) \to N \in \mathbb{N}$
\item $\varphi_{63} :\equiv \forall n \in \mathbb{N}\,(a_{n} \in A) \to \mathrm{Conv}(a,c) \to \varepsilon \in \mathbb{R} \to 0 < \varepsilon \to \exists \delta\,(\varphi_{30})$
\item $\varphi_{64} :\equiv \forall \alpha\,(\mathrm{Conv}(f \circ a,\alpha) \leftrightarrow \mathrm{Seq}(f \circ a) \land \alpha \in \mathbb{R} \land \forall \varepsilon \in \mathbb{R}\,(0 < \varepsilon \to \exists N \in \mathbb{N}\,\forall n \in \mathbb{N}\,(N \leq n \to \lvert f \circ a(n) - \alpha \rvert < \varepsilon)))$
\item $\varphi_{65} :\equiv \forall n \in \mathbb{N}\,(a_{n} \in A) \to \mathrm{Conv}(a,c) \to \forall \varepsilon \in \mathbb{R}\,(0 < \varepsilon \to \exists \delta\,(\varphi_{30}))$
\item $\varphi_{66} :\equiv \forall n \in \mathbb{N}\,(a_{n} \in A) \to \mathrm{Conv}(a,c) \to \varepsilon \in \mathbb{R} \to 0 < \varepsilon \to \varphi_{30} \to 0 < \delta$
\item $\varphi_{67} :\equiv (\mathrm{ContIn}(f,A,c) \to \varphi_{7}) \to \mathrm{ContIn}(f,A,c) \to \varphi_{12}$
\item $\varphi_{68} :\equiv \forall n \in \mathbb{N}\,(a_{n} \in A) \to \mathrm{Conv}(a,c) \to \varepsilon \in \mathbb{R} \to 0 < \varepsilon \to \varphi_{30} \to \delta \in \mathbb{R}$
\item $\varphi_{69} :\equiv \mathrm{ContIn}(f,A,c) \leftrightarrow \mathrm{RFn}(f) \land A \subseteq \mathrm{dom}(f) \land c \in A \land \forall \varepsilon \in \mathbb{R}\,(0 < \varepsilon \to \exists \delta\,(\varphi_{30}))$
\item $\varphi_{70} :\equiv \varepsilon \in \mathbb{R} \to 0 < \varepsilon \to \varphi_{30} \to \exists N \in \mathbb{N}\,\forall n \in \mathbb{N}\,(N \leq n \to \lvert a_{n} - c \rvert < \delta)$
\item $\varphi_{71} :\equiv \mathrm{Conv}(f \circ a,f(c)) \leftrightarrow \mathrm{Seq}(f \circ a) \land f(c) \in \mathbb{R} \land \forall \varepsilon \in \mathbb{R}\,(0 < \varepsilon \to \varphi_{21})$
\item $\varphi_{72} :\equiv (\mathrm{ContIn}(f,A,c) \to \mathrm{Seq}(a) \to \varphi_{6}) \to \mathrm{ContIn}(f,A,c) \to \mathrm{Seq}(a) \to \varphi_{11}$
\item $\varphi_{73} :\equiv (\mathrm{ContIn}(f,A,c) \to \mathrm{Seq}(a) \to \varphi_{11}) \to \mathrm{ContIn}(f,A,c) \to \mathrm{Seq}(a) \to \varphi_{17}$
\item $\varphi_{74} :\equiv (\mathrm{ContIn}(f,A,c) \to \mathrm{Seq}(a) \to \varphi_{17}) \to \varphi_{9} \to n \in \mathbb{N} \to a_{n} \in \mathrm{dom}(f)$
\item $\varphi_{75} :\equiv \forall n\,(\varphi_{9} \to n \in \mathbb{N} \to a_{n} \in \mathrm{dom}(f)) \to \forall n\,(\varphi_{9}) \to \forall n \in \mathbb{N}\,(a_{n} \in \mathrm{dom}(f))$
\item $\varphi_{76} :\equiv \forall n \in \mathbb{N}\,(a_{n} \in A) \to \mathrm{Conv}(a,c) \to \varepsilon \in \mathbb{R} \to 0 < \varepsilon \to \varphi_{62}$
\item $\varphi_{77} :\equiv 0 < \varepsilon \to \varphi_{30} \to N \in \mathbb{N} \land \forall n \in \mathbb{N}\,(N \leq n \to \lvert a_{n} - c \rvert < \delta) \to \varphi_{18}$
\item $\varphi_{78} :\equiv \forall n \in \mathbb{N}\,(a_{n} \in A) \to \mathrm{Conv}(a,c) \to \varepsilon \in \mathbb{R} \to 0 < \varepsilon \to \varphi_{30} \to \varphi_{21}$
\item $\varphi_{79} :\equiv 0 < \varepsilon \to \varphi_{30} \to N \in \mathbb{N} \land \forall n \in \mathbb{N}\,(N \leq n \to \lvert a_{n} - c \rvert < \delta) \to \varphi_{21}$
\item $\varphi_{80} :\equiv \forall n \in \mathbb{N}\,(a_{n} \in A) \to \mathrm{Conv}(a,c) \to \varepsilon \in \mathbb{R} \to 0 < \varepsilon \to \varphi_{30} \to \varphi_{22}$
\item $\varphi_{81} :\equiv \forall n \in \mathbb{N}\,(a_{n} \in A) \to \mathrm{Conv}(a,c) \to \varepsilon \in \mathbb{R} \to 0 < \varepsilon \to \varphi_{30} \to \varphi_{23}$
\item $\varphi_{82} :\equiv (\varphi_{9} \to \forall n\,(\varphi_{9})) \to \mathrm{ContIn}(f,A,c) \to \mathrm{Seq}(a) \to \varphi_{20}$
\item $\varphi_{83} :\equiv \varphi_{9} \land \varepsilon \in \mathbb{R} \land 0 < \varepsilon \land \varphi_{30} \land (N \in \mathbb{N} \land \forall n \in \mathbb{N}\,(N \leq n \to \lvert a_{n} - c \rvert < \delta))$
\item $\varphi_{84} :\equiv \forall \varepsilon \in \mathbb{R}\,(0 < \varepsilon \to \exists \delta\,(\varphi_{30})) \to \varepsilon \in \mathbb{R} \to 0 < \varepsilon \to \exists \delta\,(\varphi_{30})$
\item $\varphi_{85} :\equiv \mathrm{ContIn}(f,A,c) \to \mathrm{Seq}(a) \to \forall n \in \mathbb{N}\,(a_{n} \in A) \to \mathrm{Conv}(a,c) \to \varphi_{70}$
\item $\varphi_{86} :\equiv \forall n \in \mathbb{N}\,(a_{n} \in A) \to \mathrm{Conv}(a,c) \to \varepsilon \in \mathbb{R} \to 0 < \varepsilon \to \varphi_{30} \to \varphi_{30}$
\item $\varphi_{87} :\equiv \forall a\,(\mathrm{RFn}(f) \to \mathrm{Seq}(a) \to \varphi_{39}) \to \mathrm{RFn}(f) \to \mathrm{Seq}(a) \to \varphi_{39}$
\item $\varphi_{88} :\equiv \forall f\,\forall a\,(\mathrm{RFn}(f) \to \mathrm{Seq}(a) \to \varphi_{39}) \to \forall a\,(\mathrm{RFn}(f) \to \mathrm{Seq}(a) \to \varphi_{39})$
\item $\varphi_{89} :\equiv \forall n \in \mathbb{N}\,(a_{n} \in A) \to \mathrm{Conv}(a,c) \to \varepsilon \in \mathbb{R} \to 0 < \varepsilon \to \varphi_{30} \to \varphi_{37}$
\item $\varphi_{90} :\equiv \forall n \in \mathbb{N}\,(a_{n} \in A) \to \mathrm{Conv}(a,c) \to \varepsilon \in \mathbb{R} \to 0 < \varepsilon \to \varphi_{30} \to \varphi_{42}$
\item $\varphi_{91} :\equiv \forall n \in \mathbb{N}\,(a_{n} \in A) \to \mathrm{Conv}(a,c) \to \varepsilon \in \mathbb{R} \to 0 < \varepsilon \to \varphi_{30} \to \varphi_{43}$
\item $\varphi_{92} :\equiv (\mathrm{ContIn}(f,A,c) \to \mathrm{Seq}(a) \to \varphi_{59}) \to \mathrm{ContIn}(f,A,c) \to \mathrm{Seq}(a) \to \varphi_{15}$
\item $\varphi_{93} :\equiv (\mathrm{ContIn}(f,A,c) \to \mathrm{Seq}(a) \to \mathrm{Seq}(a)) \to (\mathrm{ContIn}(f,A,c) \to \mathrm{Seq}(a) \to \varphi_{20}) \to \varphi_{55}$
\item $\varphi_{94} :\equiv \forall n \in \mathbb{N}\,(a_{n} \in A) \to \mathrm{Conv}(a,c) \to \varepsilon \in \mathbb{R} \to 0 < \varepsilon \to \varphi_{30} \to \varphi_{44}$
\item $\varphi_{95} :\equiv \forall n \in \mathbb{N}\,(a_{n} \in A) \to \mathrm{Conv}(a,c) \to \varepsilon \in \mathbb{R} \to 0 < \varepsilon \to \varphi_{30} \to \varphi_{45}$
\item $\varphi_{96} :\equiv \forall n \in \mathbb{N}\,(a_{n} \in A) \to \mathrm{Conv}(a,c) \to \varepsilon \in \mathbb{R} \to 0 < \varepsilon \to \varphi_{30} \to \varphi_{46}$
\item $\varphi_{97} :\equiv \forall n \in \mathbb{N}\,(a_{n} \in A) \to \mathrm{Conv}(a,c) \to \varepsilon \in \mathbb{R} \to 0 < \varepsilon \to \varphi_{30} \to \varphi_{47}$
\item $\varphi_{98} :\equiv (\mathrm{ContIn}(f,A,c) \to \mathrm{ContIn}(f,A,c)) \to \varphi_{69} \to \mathrm{ContIn}(f,A,c) \to \varphi_{2}$
\item $\varphi_{99} :\equiv \forall n \in \mathbb{N}\,(a_{n} \in A) \to \mathrm{Conv}(a,c) \to \varepsilon \in \mathbb{R} \to 0 < \varepsilon \to \varphi_{30} \to \varphi_{49}$
\item $\varphi_{100} :\equiv (\mathrm{ContIn}(f,A,c) \to \mathrm{ContIn}(f,A,c)) \to \varphi_{69} \to \mathrm{ContIn}(f,A,c) \to \varphi_{3}$
\item $\varphi_{101} :\equiv (A \subseteq \mathrm{dom}(f) \to c \in A \to c \in \mathrm{dom}(f)) \to (\mathrm{ContIn}(f,A,c) \to \varphi_{10}) \to \varphi_{60}$
\item $\varphi_{102} :\equiv \forall n \in \mathbb{N}\,(a_{n} \in A) \to \mathrm{Conv}(a,c) \to \varepsilon \in \mathbb{R} \to 0 < \varepsilon \to \varphi_{30} \to \varphi_{51}$
\item $\varphi_{103} :\equiv \forall n \in \mathbb{N}\,(a_{n} \in A) \to \mathrm{Conv}(a,c) \to \varepsilon \in \mathbb{R} \to 0 < \varepsilon \to \varphi_{30} \to \varphi_{52}$
\item $\varphi_{104} :\equiv \forall n \in \mathbb{N}\,(a_{n} \in A) \to \mathrm{Conv}(a,c) \to \varepsilon \in \mathbb{R} \to 0 < \varepsilon \to \varphi_{30} \to \varphi_{53}$
\item $\varphi_{105} :\equiv (\mathrm{ContIn}(f,A,c) \to \mathrm{ContIn}(f,A,c)) \to \varphi_{69} \to \mathrm{ContIn}(f,A,c) \to \varphi_{10}$
\item $\varphi_{106} :\equiv (\mathrm{RFn}(f) \to c \in \mathrm{dom}(f) \to f(c) \in \mathbb{R}) \to (\mathrm{ContIn}(f,A,c) \to \varphi_{2}) \to \varphi_{67}$
\item $\varphi_{107} :\equiv \forall n \in \mathbb{N}\,(a_{n} \in A) \to \mathrm{Conv}(a,c) \to \varepsilon \in \mathbb{R} \to 0 < \varepsilon \to \varphi_{30} \to \varphi_{56}$
\item $\varphi_{108} :\equiv \mathrm{Seq}(a) \to \forall n \in \mathbb{N}\,(a_{n} \in A) \to \mathrm{Conv}(a,c) \to \varepsilon \in \mathbb{R} \to \varphi_{77}$
\item $\varphi_{109} :\equiv \mathrm{Seq}(a) \to \forall n \in \mathbb{N}\,(a_{n} \in A) \to \mathrm{Conv}(a,c) \to \varepsilon \in \mathbb{R} \to \varphi_{79}$
\item $\varphi_{110} :\equiv \forall n\,(\varphi_{83}) \to \forall n \in \mathbb{N}\,(N \leq n \to \lvert f \circ a(n) - f(c) \rvert < \varepsilon)$
\item $\varphi_{111} :\equiv (\varphi_{9} \to \forall \varepsilon\,(\varphi_{9})) \to \mathrm{ContIn}(f,A,c) \to \mathrm{Seq}(a) \to \varphi_{59}$
\item $\varphi_{112} :\equiv \forall n \in \mathbb{N}\,(a_{n} \in A) \to \mathrm{Conv}(a,c) \to \varepsilon \in \mathbb{R} \to 0 < \varepsilon \to \varphi_{30} \to \varphi_{61}$
\item $\varphi_{113} :\equiv \varphi_{9} \land \varepsilon \in \mathbb{R} \land 0 < \varepsilon \land \varphi_{30} \land \neg \varphi_{21} \to \neg (N \in \mathbb{N} \land \forall n \in \mathbb{N}\,(N \leq n \to \lvert a_{n} - c \rvert < \delta))$
\item $\varphi_{114} :\equiv (A \subseteq \mathrm{dom}(f) \to a_{n} \in A \to a_{n} \in \mathrm{dom}(f)) \to (\mathrm{ContIn}(f,A,c) \to \varphi_{10}) \to \varphi_{73}$
\item $\varphi_{115} :\equiv (\mathrm{ContIn}(f,A,c) \to \mathrm{Seq}(a) \to \varphi_{57}) \to \varphi_{9} \to \varepsilon \in \mathbb{R} \to 0 < \varepsilon \to \varphi_{21}$
\item $\varphi_{116} :\equiv \varphi_{9} \land \varepsilon \in \mathbb{R} \land 0 < \varepsilon \land \neg \varphi_{21} \to \forall \delta\,(\varphi_{9} \land \varepsilon \in \mathbb{R} \land 0 < \varepsilon \land \neg \varphi_{21})$
\item $\varphi_{117} :\equiv \forall \varepsilon\,(\varphi_{9} \to \varepsilon \in \mathbb{R} \to 0 < \varepsilon \to \varphi_{21}) \to \forall \varepsilon\,(\varphi_{9}) \to \forall \varepsilon \in \mathbb{R}\,(0 < \varepsilon \to \varphi_{21})$
\item $\varphi_{118} :\equiv (f \circ a(n) = f(a_{n}) \to f(a_{n}) = f \circ a(n)) \to \mathrm{ContIn}(f,A,c) \to \mathrm{Seq}(a) \to \varphi_{91}$
\item $\varphi_{119} :\equiv (\mathrm{ContIn}(f,A,c) \to \mathrm{ContIn}(f,A,c)) \to \varphi_{69} \to \mathrm{ContIn}(f,A,c) \to \mathrm{Seq}(a) \to \varphi_{65}$
\item $\varphi_{120} :\equiv (\mathrm{ContIn}(f,A,c) \to \varphi_{13}) \to (\mathrm{ContIn}(f,A,c) \to \varphi_{12}) \to \varphi_{92}$
\item $\varphi_{121} :\equiv (\mathrm{ContIn}(f,A,c) \to \mathrm{Seq}(a) \to \varphi_{86}) \to \mathrm{ContIn}(f,A,c) \to \mathrm{Seq}(a) \to \varphi_{66}$
\item $\varphi_{122} :\equiv (\mathrm{ContIn}(f,A,c) \to \mathrm{Seq}(a) \to \varphi_{86}) \to \mathrm{ContIn}(f,A,c) \to \mathrm{Seq}(a) \to \varphi_{68}$
\item $\varphi_{123} :\equiv (\mathrm{RFn}(f) \to \mathrm{Seq}(a) \to \varphi_{39}) \to (\mathrm{ContIn}(f,A,c) \to \varphi_{2}) \to \varphi_{93}$
\item $\varphi_{124} :\equiv (\mathrm{ContIn}(f,A,c) \to \mathrm{Seq}(a) \to \varphi_{78}) \to \varphi_{9} \land \varepsilon \in \mathbb{R} \land 0 < \varepsilon \land \neg \varphi_{21} \to \neg \varphi_{30}$
\item $\varphi_{125} :\equiv \varphi_{9} \land \varepsilon \in \mathbb{R} \land 0 < \varepsilon \land \varphi_{30} \land \neg \varphi_{21} \to \forall N\,(\varphi_{9} \land \varepsilon \in \mathbb{R} \land 0 < \varepsilon \land \varphi_{30} \land \neg \varphi_{21})$
\item $\varphi_{126} :\equiv \forall \delta\,(\varphi_{9} \land \varepsilon \in \mathbb{R} \land 0 < \varepsilon \land \neg \varphi_{21} \to \neg \varphi_{30}) \to \forall \delta\,(\varphi_{9} \land \varepsilon \in \mathbb{R} \land 0 < \varepsilon \land \neg \varphi_{21}) \to \forall \delta\,(\neg \varphi_{30})$
\item $\varphi_{127} :\equiv (\mathrm{ContIn}(f,A,c) \to \mathrm{Seq}(a) \to \varphi_{80}) \to \mathrm{ContIn}(f,A,c) \to \mathrm{Seq}(a) \to \varphi_{89}$
\item $\varphi_{128} :\equiv (\mathrm{ContIn}(f,A,c) \to \mathrm{Seq}(a) \to \varphi_{94}) \to \mathrm{ContIn}(f,A,c) \to \mathrm{Seq}(a) \to \varphi_{76}$
\item $\varphi_{129} :\equiv (\mathrm{ContIn}(f,A,c) \to \mathrm{Seq}(a) \to \varphi_{81}) \to \mathrm{ContIn}(f,A,c) \to \mathrm{Seq}(a) \to \varphi_{90}$
\item $\varphi_{130} :\equiv (\mathrm{ContIn}(f,A,c) \to \mathrm{Seq}(a) \to \varphi_{89}) \to \mathrm{ContIn}(f,A,c) \to \mathrm{Seq}(a) \to \varphi_{95}$
\item $\varphi_{131} :\equiv (\mathrm{ContIn}(f,A,c) \to \mathrm{Seq}(a) \to \varphi_{86}) \to \mathrm{ContIn}(f,A,c) \to \mathrm{Seq}(a) \to \varphi_{107}$
\item $\varphi_{132} :\equiv (\mathrm{ContIn}(f,A,c) \to \mathrm{Seq}(a) \to \varphi_{94}) \to \mathrm{ContIn}(f,A,c) \to \mathrm{Seq}(a) \to \varphi_{99}$
\item $\varphi_{133} :\equiv (\mathrm{ContIn}(f,A,c) \to \mathrm{Seq}(a) \to \varphi_{97}) \to \varphi_{83} \to n \in \mathbb{N} \to N \leq n \to \lvert f \circ a(n) - f(c) \rvert < \varepsilon$
\item $\varphi_{134} :\equiv \forall n\,(\varphi_{83} \to n \in \mathbb{N} \to N \leq n \to \lvert f \circ a(n) - f(c) \rvert < \varepsilon) \to \varphi_{110}$
\item $\varphi_{135} :\equiv (\mathrm{ContIn}(f,A,c) \to \mathrm{Seq}(a) \to \varphi_{65}) \to \varphi_{84} \to \mathrm{ContIn}(f,A,c) \to \mathrm{Seq}(a) \to \varphi_{63}$
\item $\varphi_{136} :\equiv (\mathrm{ContIn}(f,A,c) \to \mathrm{Seq}(a) \to \varphi_{102}) \to \mathrm{ContIn}(f,A,c) \to \varphi_{108}$
\item $\varphi_{137} :\equiv (\mathrm{ContIn}(f,A,c) \to \mathrm{Seq}(a) \to \varphi_{68}) \to (\mathrm{ContIn}(f,A,c) \to \mathrm{Seq}(a) \to \varphi_{66}) \to \varphi_{85}$
\item $\varphi_{138} :\equiv (\mathrm{ContIn}(f,A,c) \to \mathrm{Seq}(a) \to \varphi_{99}) \to \varphi_{33} \to \mathrm{ContIn}(f,A,c) \to \mathrm{Seq}(a) \to \varphi_{96}$
\item $\varphi_{139} :\equiv (\mathrm{ContIn}(f,A,c) \to \mathrm{Seq}(a) \to \varphi_{104}) \to \varphi_{36} \to \mathrm{ContIn}(f,A,c) \to \mathrm{Seq}(a) \to \varphi_{103}$
\item $\varphi_{140} :\equiv (\mathrm{ContIn}(f,A,c) \to \mathrm{Seq}(a) \to \varphi_{95}) \to \varphi_{48} \to \mathrm{ContIn}(f,A,c) \to \mathrm{Seq}(a) \to \varphi_{97}$
\item $\varphi_{141} :\equiv (\varphi_{83} \to \forall n\,(\varphi_{83})) \to \mathrm{ContIn}(f,A,c) \to \mathrm{Seq}(a) \to \varphi_{102}$
\item $\varphi_{142} :\equiv (\mathrm{ContIn}(f,A,c) \to \varphi_{108}) \to (\forall N\,(\neg \varphi_{18}) \to \neg \varphi_{18}) \to \mathrm{ContIn}(f,A,c) \to \varphi_{109}$
\item $\varphi_{143} :\equiv (\mathrm{Conv}(a,c) \to \varphi_{27}) \to (\mathrm{ContIn}(f,A,c) \to \mathrm{Seq}(a) \to \varphi_{1}) \to \varphi_{137}$
\item $\varphi_{144} :\equiv (\mathrm{ContIn}(f,A,c) \to \mathrm{Seq}(a) \to \varphi_{107}) \to \varphi_{58} \to \mathrm{ContIn}(f,A,c) \to \mathrm{Seq}(a) \to \varphi_{112}$
\item $\varphi_{145} :\equiv (\mathrm{ContIn}(f,A,c) \to \mathrm{Seq}(a) \to \varphi_{96}) \to (\mathrm{ContIn}(f,A,c) \to \mathrm{Seq}(a) \to \varphi_{81}) \to \varphi_{127}$
\item $\varphi_{146} :\equiv (\mathrm{ContIn}(f,A,c) \to \mathrm{Seq}(a) \to \varphi_{112}) \to \varphi_{26} \to (\mathrm{ContIn}(f,A,c) \to \mathrm{Seq}(a) \to \varphi_{81}) \to \varphi_{130}$
\item $\varphi_{147} :\equiv \forall \delta\,(\varphi_{9} \land \varepsilon \in \mathbb{R} \land 0 < \varepsilon \land \neg \varphi_{21} \to \neg \varphi_{30}) \to \varphi_{126} \to \varphi_{116} \to \mathrm{ContIn}(f,A,c) \to \mathrm{Seq}(a) \to \varphi_{57}$
\item $\varphi_{148} :\equiv (\forall N\,(\varphi_{113}) \to \forall N\,(\varphi_{9} \land \varepsilon \in \mathbb{R} \land 0 < \varepsilon \land \varphi_{30} \land \neg \varphi_{21}) \to \varphi_{4}) \to \varphi_{125} \to \mathrm{ContIn}(f,A,c) \to \mathrm{Seq}(a) \to \varphi_{78}$
\item $\varphi_{149} :\equiv \forall n\,(\varphi_{83} \to n \in \mathbb{N} \to N \leq n \to \lvert f \circ a(n) - f(c) \rvert < \varepsilon) \to \varphi_{134} \to \varphi_{141}$
\end{abbrevlist}
\end{fullproof}

% 自動生成 (prover/build.js): ContSeqInv  ノード数 466
\begin{fullproof}{[G2]}{$\forall f\,\forall A\,\forall c\,(\mathrm{RFn}(f) \to A \subseteq \mathrm{dom}(f) \to c \in A \to \forall a\,(\mathrm{Seq}(a) \to \forall n \in \mathbb{N}\,(a_{n} \in A) \to \mathrm{Conv}(a,c) \to \mathrm{Conv}(f \circ a,f(c))) \to \mathrm{ContIn}(f,A,c))$}{466}\phantomsection\label{tree:ContSeqInv}
\noindent$\Pi_{1}$:
\begin{ptree}
\AXC{$\forall x\,(\neg x < x)$\lbl{[O1]}}
\AXC{$\forall x\,(\neg x < x) \to \neg 0 < 0$\lbl{A2}}
\BIC{$\neg 0 < 0$}
\AXC{$\neg 0 < 0 \to \mathrm{RFn}(f) \to A \subseteq \mathrm{dom}(f) \to c \in A \to \varphi_{70}$\lbl{A1}}
\BIC{$\mathrm{RFn}(f) \to A \subseteq \mathrm{dom}(f) \to c \in A \to \varphi_{70}$}
\end{ptree}
\noindent$\Pi_{2}$:
\begin{ptree}
\AXC{$\mathrm{RFn}(f) \to A \subseteq \mathrm{dom}(f) \to c \in A \to \varphi_{195}$\lbl{A1}}
\AXC{$(\mathrm{RFn}(f) \to A \subseteq \mathrm{dom}(f) \to c \in A \to \varphi_{195}) \to \varphi_{178}$\lbl{A1}}
\BIC{$\varphi_{178}$}
\end{ptree}
\noindent$\Pi_{3}$:
\begin{ptree}
\AXC{$\varphi_{45} \to N \in \mathbb{N} \to \varphi_{32}$\lbl{A2}}
\AXC{$\varphi_{193}$\lbl{A1}}
\AXC{$\varphi_{193} \to \mathrm{RFn}(f) \to A \subseteq \mathrm{dom}(f) \to c \in A \to \varphi_{177}$\lbl{A1}}
\BIC{$\mathrm{RFn}(f) \to A \subseteq \mathrm{dom}(f) \to c \in A \to \varphi_{177}$}
\AXC{$\varphi_{244}$\lbl{A1}}
\BIC{$(\varphi_{45} \to N \in \mathbb{N} \to \varphi_{32}) \to \varphi_{184}$}
\BIC{$\varphi_{184}$}
\end{ptree}
\noindent$\Pi_{4}$:
\begin{ptree}
\AXC{$z = w \to \lvert z - f(c) \rvert < e_{0} \to \lvert w - f(c) \rvert < e_{0}$\lbl{A6}}
\UIC{$\varphi_{1}$}
\UIC{$\forall z\,(\varphi_{1})$}
\AXC{$\forall z\,(\varphi_{1}) \to \varphi_{39}$\lbl{A2}}
\BIC{$\varphi_{39}$}
\AXC{$\varphi_{39} \to f \circ a(N) = f(a_{N}) \to \varphi_{20}$\lbl{A2}}
\BIC{$f \circ a(N) = f(a_{N}) \to \varphi_{20}$}
\end{ptree}
\noindent$\Pi_{5}$:
\begin{ptree}
\AXC{$x = x$\lbl{A5}}
\UIC{$\forall x\,(x = x)$}
\AXC{$\forall x\,(x = x) \to N = N$\lbl{A2}}
\BIC{$N = N$}
\AXC{$N = N \to \mathrm{RFn}(f) \to A \subseteq \mathrm{dom}(f) \to c \in A \to \varphi_{169}$\lbl{A1}}
\BIC{$\mathrm{RFn}(f) \to A \subseteq \mathrm{dom}(f) \to c \in A \to \varphi_{169}$}
\end{ptree}
\noindent$\Pi_{6}$:
\begin{ptree}
\AXC{$\stackrel{\Pi_{2}}{\vdots}$}\noLine
\UIC{$\varphi_{178}$}
\AXC{$\varphi_{54}$\lbl{A2}}
\AXC{$\mathrm{RFn}(f) \to A \subseteq \mathrm{dom}(f) \to c \in A \to \varphi_{195}$\lbl{A1}}
\AXC{$\varphi_{241}$\lbl{A1}}
\BIC{$\mathrm{RFn}(f) \to A \subseteq \mathrm{dom}(f) \to \varphi_{192}$}
\AXC{$\varphi_{243}$\lbl{A1}}
\BIC{$\varphi_{54} \to \mathrm{RFn}(f) \to A \subseteq \mathrm{dom}(f) \to \varphi_{190}$}
\BIC{$\mathrm{RFn}(f) \to A \subseteq \mathrm{dom}(f) \to \varphi_{190}$}
\AXC{$\varphi_{262}$\lbl{A1}}
\BIC{$\varphi_{178} \to (\mathrm{RFn}(f) \to A \subseteq \mathrm{dom}(f) \to c \in A \to \varphi_{169}) \to \varphi_{196}$}
\BIC{$(\mathrm{RFn}(f) \to A \subseteq \mathrm{dom}(f) \to c \in A \to \varphi_{169}) \to \varphi_{196}$}
\end{ptree}
\noindent$\Pi_{7}$:
\begin{ptree}
\AXC{$\stackrel{\Pi_{5}}{\vdots}$}\noLine
\UIC{$\mathrm{RFn}(f) \to A \subseteq \mathrm{dom}(f) \to c \in A \to \varphi_{169}$}
\AXC{$\stackrel{\Pi_{6}}{\vdots}$}\noLine
\UIC{$(\mathrm{RFn}(f) \to A \subseteq \mathrm{dom}(f) \to c \in A \to \varphi_{169}) \to \varphi_{196}$}
\BIC{$\varphi_{196}$}
\end{ptree}
\noindent$\Pi_{8}$:
\begin{ptree}
\AXC{$\varphi_{45} \to n \in \mathbb{N} \to \varphi_{33}$\lbl{A2}}
\AXC{$\varphi_{193}$\lbl{A1}}
\AXC{$\varphi_{193} \to \mathrm{RFn}(f) \to A \subseteq \mathrm{dom}(f) \to c \in A \to \varphi_{177}$\lbl{A1}}
\BIC{$\mathrm{RFn}(f) \to A \subseteq \mathrm{dom}(f) \to c \in A \to \varphi_{177}$}
\AXC{$\varphi_{245}$\lbl{A1}}
\BIC{$(\varphi_{45} \to n \in \mathbb{N} \to \varphi_{33}) \to \varphi_{185}$}
\BIC{$\varphi_{185}$}
\AXC{$\varphi_{185} \to \varphi_{156}$\lbl{A1}}
\BIC{$\varphi_{156}$}
\end{ptree}
\noindent$\Pi_{9}$:
\begin{ptree}
\AXC{$\stackrel{\Pi_{8}}{\vdots}$}\noLine
\UIC{$\varphi_{156}$}
\AXC{$\varphi_{156} \to \varphi_{97} \land (\mathrm{Seq}(a) \land \varphi_{45}) \to n \in \mathbb{N} \to a_{n} \in A$\lbl{A1}}
\BIC{$\varphi_{97} \land (\mathrm{Seq}(a) \land \varphi_{45}) \to n \in \mathbb{N} \to a_{n} \in A$}
\UIC{$\forall n\,(\varphi_{97} \land (\mathrm{Seq}(a) \land \varphi_{45}) \to n \in \mathbb{N} \to a_{n} \in A)$}
\AXC{$\forall n\,(\varphi_{97} \land (\mathrm{Seq}(a) \land \varphi_{45}) \to n \in \mathbb{N} \to a_{n} \in A) \to \varphi_{223} \to \varphi_{247}$\lbl{A1}}
\BIC{$\varphi_{223} \to \varphi_{247}$}
\end{ptree}
\noindent$\Pi_{10}$:
\begin{ptree}
\AXC{$\varphi_{97} \land (\mathrm{Seq}(a) \land \varphi_{45}) \to \forall n\,(\varphi_{97} \land (\mathrm{Seq}(a) \land \varphi_{45}))$\lbl{A4}}
\AXC{$\varphi_{223}$\lbl{A3}}
\AXC{$\stackrel{\Pi_{9}}{\vdots}$}\noLine
\UIC{$\varphi_{223} \to \varphi_{247}$}
\BIC{$\varphi_{247}$}
\BIC{$\varphi_{159}$}
\end{ptree}
\noindent$\Pi_{11}$:
\begin{ptree}
\AXC{$\forall n \in \mathbb{N}\,(a_{n} \in A) \to n \in \mathbb{N} \to a_{n} \in A$\lbl{A2}}
\AXC{$\stackrel{\Pi_{10}}{\vdots}$}\noLine
\UIC{$\varphi_{159}$}
\AXC{$\varphi_{159} \to (\forall n \in \mathbb{N}\,(a_{n} \in A) \to n \in \mathbb{N} \to a_{n} \in A) \to \varphi_{156}$\lbl{A1}}
\BIC{$(\forall n \in \mathbb{N}\,(a_{n} \in A) \to n \in \mathbb{N} \to a_{n} \in A) \to \varphi_{156}$}
\BIC{$\varphi_{156}$}
\end{ptree}
\noindent$\Pi_{12}$:
\begin{ptree}
\AXC{$\forall X\,\forall Y\,\forall x\,(X \subseteq Y \to x \in X \to x \in Y)$\lbl{[B5]}}
\AXC{$\varphi_{5}$\lbl{A2}}
\BIC{$\forall Y\,\forall x\,(A \subseteq Y \to x \in A \to x \in Y)$}
\AXC{$\varphi_{10}$\lbl{A2}}
\BIC{$\forall x\,(A \subseteq \mathrm{dom}(f) \to x \in A \to x \in \mathrm{dom}(f))$}
\AXC{$\varphi_{40}$\lbl{A2}}
\BIC{$A \subseteq \mathrm{dom}(f) \to a_{n} \in A \to a_{n} \in \mathrm{dom}(f)$}
\AXC{$\varphi_{226}$\lbl{A1}}
\BIC{$(\mathrm{RFn}(f) \to A \subseteq \mathrm{dom}(f) \to A \subseteq \mathrm{dom}(f)) \to \varphi_{156} \to \varphi_{160}$}
\end{ptree}
\noindent$\Pi_{13}$:
\begin{ptree}
\AXC{$\stackrel{\Pi_{11}}{\vdots}$}\noLine
\UIC{$\varphi_{156}$}
\AXC{$\mathrm{RFn}(f) \to A \subseteq \mathrm{dom}(f) \to A \subseteq \mathrm{dom}(f)$\lbl{A1}}
\AXC{$\stackrel{\Pi_{12}}{\vdots}$}\noLine
\UIC{$(\mathrm{RFn}(f) \to A \subseteq \mathrm{dom}(f) \to A \subseteq \mathrm{dom}(f)) \to \varphi_{156} \to \varphi_{160}$}
\BIC{$\varphi_{156} \to \varphi_{160}$}
\BIC{$\varphi_{160}$}
\end{ptree}
\noindent$\Pi_{14}$:
\begin{ptree}
\AXC{$\varphi_{161} \to \varphi_{165}$\lbl{A3}}
\AXC{$\stackrel{\Pi_{13}}{\vdots}$}\noLine
\UIC{$\varphi_{160}$}
\AXC{$\varphi_{160} \to \varphi_{97} \land (\mathrm{Seq}(a) \land \varphi_{45}) \to n \in \mathbb{N} \to a_{n} \in \mathrm{dom}(f)$\lbl{A1}}
\BIC{$\varphi_{97} \land (\mathrm{Seq}(a) \land \varphi_{45}) \to n \in \mathbb{N} \to a_{n} \in \mathrm{dom}(f)$}
\UIC{$\varphi_{161}$}
\AXC{$\varphi_{161} \to (\varphi_{161} \to \varphi_{165}) \to \varphi_{248}$\lbl{A1}}
\BIC{$(\varphi_{161} \to \varphi_{165}) \to \varphi_{248}$}
\BIC{$\varphi_{248}$}
\end{ptree}
\noindent$\Pi_{15}$:
\begin{ptree}
\AXC{$\forall f\,\forall a\,(\mathrm{RFn}(f) \to \mathrm{Seq}(a) \to \varphi_{44})$\lbl{[S8]}}
\AXC{$\varphi_{84}$\lbl{A2}}
\BIC{$\forall a\,(\mathrm{RFn}(f) \to \mathrm{Seq}(a) \to \varphi_{44})$}
\AXC{$\varphi_{77}$\lbl{A2}}
\BIC{$\mathrm{RFn}(f) \to \mathrm{Seq}(a) \to \varphi_{44}$}
\AXC{$(\mathrm{RFn}(f) \to \mathrm{Seq}(a) \to \varphi_{44}) \to \varphi_{251}$\lbl{A1}}
\BIC{$\varphi_{251}$}
\end{ptree}
\noindent$\Pi_{16}$:
\begin{ptree}
\AXC{$\varphi_{193}$\lbl{A1}}
\AXC{$\varphi_{193} \to \varphi_{147}$\lbl{A1}}
\BIC{$\varphi_{147}$}
\AXC{$\mathrm{RFn}(f) \to \mathrm{RFn}(f)$\lbl{A1}}
\AXC{$\stackrel{\Pi_{15}}{\vdots}$}\noLine
\UIC{$\varphi_{251}$}
\BIC{$\varphi_{147} \to \varphi_{162} \to \mathrm{RFn}(f) \to A \subseteq \mathrm{dom}(f) \to \varphi_{168}$}
\BIC{$\varphi_{162} \to \mathrm{RFn}(f) \to A \subseteq \mathrm{dom}(f) \to \varphi_{168}$}
\end{ptree}
\noindent$\Pi_{17}$:
\begin{ptree}
\AXC{$\varphi_{97} \land (\mathrm{Seq}(a) \land \varphi_{45}) \to \forall n\,(\varphi_{97} \land (\mathrm{Seq}(a) \land \varphi_{45}))$\lbl{A4}}
\AXC{$\stackrel{\Pi_{14}}{\vdots}$}\noLine
\UIC{$\varphi_{248}$}
\BIC{$\varphi_{162}$}
\AXC{$\stackrel{\Pi_{16}}{\vdots}$}\noLine
\UIC{$\varphi_{162} \to \mathrm{RFn}(f) \to A \subseteq \mathrm{dom}(f) \to \varphi_{168}$}
\BIC{$\mathrm{RFn}(f) \to A \subseteq \mathrm{dom}(f) \to \varphi_{168}$}
\AXC{$\varphi_{233}$\lbl{A1}}
\BIC{$\mathrm{RFn}(f) \to A \subseteq \mathrm{dom}(f) \to \varphi_{183}$}
\end{ptree}
\noindent$\Pi_{18}$:
\begin{ptree}
\AXC{$\varphi_{42}$\lbl{A2}}
\AXC{$\stackrel{\Pi_{17}}{\vdots}$}\noLine
\UIC{$\mathrm{RFn}(f) \to A \subseteq \mathrm{dom}(f) \to \varphi_{183}$}
\AXC{$(\mathrm{RFn}(f) \to A \subseteq \mathrm{dom}(f) \to \varphi_{183}) \to \varphi_{42} \to \varphi_{198}$\lbl{A1}}
\BIC{$\varphi_{42} \to \varphi_{198}$}
\BIC{$\varphi_{198}$}
\AXC{$\varphi_{198} \to \varphi_{178} \to \varphi_{189}$\lbl{A1}}
\BIC{$\varphi_{178} \to \varphi_{189}$}
\end{ptree}
\noindent$\Pi_{19}$:
\begin{ptree}
\AXC{$\stackrel{\Pi_{4}}{\vdots}$}\noLine
\UIC{$f \circ a(N) = f(a_{N}) \to \varphi_{20}$}
\AXC{$\stackrel{\Pi_{7}}{\vdots}$}\noLine
\UIC{$\varphi_{196}$}
\AXC{$\stackrel{\Pi_{2}}{\vdots}$}\noLine
\UIC{$\varphi_{178}$}
\AXC{$\stackrel{\Pi_{18}}{\vdots}$}\noLine
\UIC{$\varphi_{178} \to \varphi_{189}$}
\BIC{$\varphi_{189}$}
\AXC{$\varphi_{189} \to \varphi_{196} \to (f \circ a(N) = f(a_{N}) \to \varphi_{20}) \to \varphi_{191}$\lbl{A1}}
\BIC{$\varphi_{196} \to (f \circ a(N) = f(a_{N}) \to \varphi_{20}) \to \varphi_{191}$}
\BIC{$(f \circ a(N) = f(a_{N}) \to \varphi_{20}) \to \varphi_{191}$}
\BIC{$\varphi_{191}$}
\end{ptree}
\noindent$\Pi_{20}$:
\begin{ptree}
\AXC{$\stackrel{\Pi_{2}}{\vdots}$}\noLine
\UIC{$\varphi_{178}$}
\AXC{$\stackrel{\Pi_{3}}{\vdots}$}\noLine
\UIC{$\varphi_{184}$}
\AXC{$\varphi_{184} \to \varphi_{178} \to \mathrm{RFn}(f) \to A \subseteq \mathrm{dom}(f) \to c \in A \to \varphi_{200}$\lbl{A1}}
\BIC{$\varphi_{178} \to \mathrm{RFn}(f) \to A \subseteq \mathrm{dom}(f) \to c \in A \to \varphi_{200}$}
\BIC{$\mathrm{RFn}(f) \to A \subseteq \mathrm{dom}(f) \to c \in A \to \varphi_{200}$}
\AXC{$\stackrel{\Pi_{19}}{\vdots}$}\noLine
\UIC{$\varphi_{191}$}
\AXC{$\varphi_{191} \to \varphi_{237}$\lbl{A1}}
\BIC{$\varphi_{237}$}
\BIC{$\mathrm{RFn}(f) \to A \subseteq \mathrm{dom}(f) \to c \in A \to \varphi_{171}$}
\AXC{$\varphi_{228}$\lbl{A1}}
\BIC{$\varphi_{97} \land (\mathrm{Seq}(a) \land \varphi_{45}) \land \neg 0 < 0 \to \neg \varphi_{9}$}
\UIC{$\forall N\,(\varphi_{97} \land (\mathrm{Seq}(a) \land \varphi_{45}) \land \neg 0 < 0 \to \neg \varphi_{9})$}
\end{ptree}
\noindent$\Pi_{21}$:
\begin{ptree}
\AXC{$\mathrm{RFn}(f) \to A \subseteq \mathrm{dom}(f) \to c \in A \to \varphi_{173}$\lbl{A1}}
\AXC{$(\mathrm{RFn}(f) \to A \subseteq \mathrm{dom}(f) \to c \in A \to \varphi_{173}) \to \varphi_{180}$\lbl{A1}}
\BIC{$\varphi_{180}$}
\end{ptree}
\noindent$\Pi_{22}$:
\begin{ptree}
\AXC{$\forall K\,\forall n\,(K \in \mathbb{N} \to n \in \mathbb{N} \to K \leq n \to \underline{n}^{-1} \leq \underline{K}^{-1})$\lbl{[C25]}}
\AXC{$\varphi_{51}$\lbl{A2}}
\BIC{$\forall n\,(K \in \mathbb{N} \to n \in \mathbb{N} \to K \leq n \to \underline{n}^{-1} \leq \underline{K}^{-1})$}
\AXC{$\varphi_{47}$\lbl{A2}}
\BIC{$K \in \mathbb{N} \to n \in \mathbb{N} \to K \leq n \to \underline{n}^{-1} \leq \underline{K}^{-1}$}
\AXC{$\varphi_{261}$\lbl{A1}}
\BIC{$(\mathrm{RFn}(f) \to A \subseteq \mathrm{dom}(f) \to \varphi_{167}) \to \varphi_{253}$}
\end{ptree}
\noindent$\Pi_{23}$:
\begin{ptree}
\AXC{$\mathrm{RFn}(f) \to A \subseteq \mathrm{dom}(f) \to c \in A \to \varphi_{173}$\lbl{A1}}
\AXC{$\varphi_{227}$\lbl{A1}}
\BIC{$\mathrm{RFn}(f) \to A \subseteq \mathrm{dom}(f) \to \varphi_{167}$}
\AXC{$\stackrel{\Pi_{22}}{\vdots}$}\noLine
\UIC{$(\mathrm{RFn}(f) \to A \subseteq \mathrm{dom}(f) \to \varphi_{167}) \to \varphi_{253}$}
\BIC{$\varphi_{253}$}
\end{ptree}
\noindent$\Pi_{24}$:
\begin{ptree}
\AXC{$\mathrm{RFn}(f) \to A \subseteq \mathrm{dom}(f) \to c \in A \to \varphi_{170}$\lbl{A1}}
\AXC{$\mathrm{RFn}(f) \to A \subseteq \mathrm{dom}(f) \to \varphi_{174}$\lbl{A1}}
\AXC{$\stackrel{\Pi_{23}}{\vdots}$}\noLine
\UIC{$\varphi_{253}$}
\BIC{$(\mathrm{RFn}(f) \to A \subseteq \mathrm{dom}(f) \to c \in A \to \varphi_{170}) \to \varphi_{181}$}
\BIC{$\varphi_{181}$}
\end{ptree}
\noindent$\Pi_{25}$:
\begin{ptree}
\AXC{$\varphi_{45} \to n \in \mathbb{N} \to \varphi_{33}$\lbl{A2}}
\AXC{$\varphi_{193}$\lbl{A1}}
\AXC{$\varphi_{193} \to \mathrm{RFn}(f) \to A \subseteq \mathrm{dom}(f) \to c \in A \to \varphi_{177}$\lbl{A1}}
\BIC{$\mathrm{RFn}(f) \to A \subseteq \mathrm{dom}(f) \to c \in A \to \varphi_{177}$}
\AXC{$\varphi_{245}$\lbl{A1}}
\BIC{$(\varphi_{45} \to n \in \mathbb{N} \to \varphi_{33}) \to \varphi_{185}$}
\BIC{$\varphi_{185}$}
\AXC{$\varphi_{185} \to \varphi_{236}$\lbl{A1}}
\BIC{$\varphi_{236}$}
\end{ptree}
\noindent$\Pi_{26}$:
\begin{ptree}
\AXC{$\mathrm{RFn}(f) \to A \subseteq \mathrm{dom}(f) \to \varphi_{174}$\lbl{A1}}
\AXC{$\stackrel{\Pi_{25}}{\vdots}$}\noLine
\UIC{$\varphi_{236}$}
\BIC{$\mathrm{RFn}(f) \to A \subseteq \mathrm{dom}(f) \to c \in A \to \varphi_{201}$}
\AXC{$(\mathrm{RFn}(f) \to A \subseteq \mathrm{dom}(f) \to c \in A \to \varphi_{201}) \to \varphi_{187}$\lbl{A1}}
\BIC{$\varphi_{187}$}
\end{ptree}
\noindent$\Pi_{27}$:
\begin{ptree}
\AXC{$\forall x\,\forall y\,\forall z\,(x < y \to y \leq z \to x < z)$\lbl{[O3]}}
\AXC{$\varphi_{36}$\lbl{A2}}
\BIC{$\forall y\,\forall z\,(\lvert a_{n} - c \rvert < y \to y \leq z \to \lvert a_{n} - c \rvert < z)$}
\AXC{$\forall y\,\forall z\,(\lvert a_{n} - c \rvert < y \to y \leq z \to \lvert a_{n} - c \rvert < z) \to \varphi_{12}$\lbl{A2}}
\BIC{$\varphi_{12}$}
\AXC{$\varphi_{12} \to \varphi_{41}$\lbl{A2}}
\BIC{$\varphi_{41}$}
\end{ptree}
\noindent$\Pi_{28}$:
\begin{ptree}
\AXC{$\forall x\,\forall y\,\forall z\,(x < y \to y < z \to x < z)$\lbl{[O2]}}
\AXC{$\varphi_{43}$\lbl{A2}}
\BIC{$\forall y\,\forall z\,(\lvert a_{n} - c \rvert < y \to y < z \to \lvert a_{n} - c \rvert < z)$}
\AXC{$\forall y\,\forall z\,(\lvert a_{n} - c \rvert < y \to y < z \to \lvert a_{n} - c \rvert < z) \to \varphi_{19}$\lbl{A2}}
\BIC{$\varphi_{19}$}
\end{ptree}
\noindent$\Pi_{29}$:
\begin{ptree}
\AXC{$\stackrel{\Pi_{28}}{\vdots}$}\noLine
\UIC{$\varphi_{19}$}
\AXC{$\varphi_{19} \to \lvert a_{n} - c \rvert < \underline{K}^{-1} \to \underline{K}^{-1} < \varepsilon \to \lvert a_{n} - c \rvert < \varepsilon$\lbl{A2}}
\BIC{$\lvert a_{n} - c \rvert < \underline{K}^{-1} \to \underline{K}^{-1} < \varepsilon \to \lvert a_{n} - c \rvert < \varepsilon$}
\AXC{$\varphi_{258}$\lbl{A1}}
\BIC{$\varphi_{188} \to \varphi_{180} \to \varphi_{182}$}
\end{ptree}
\noindent$\Pi_{30}$:
\begin{ptree}
\AXC{$\stackrel{\Pi_{21}}{\vdots}$}\noLine
\UIC{$\varphi_{180}$}
\AXC{$\stackrel{\Pi_{24}}{\vdots}$}\noLine
\UIC{$\varphi_{181}$}
\AXC{$\stackrel{\Pi_{26}}{\vdots}$}\noLine
\UIC{$\varphi_{187}$}
\AXC{$\stackrel{\Pi_{27}}{\vdots}$}\noLine
\UIC{$\varphi_{41}$}
\AXC{$\varphi_{41} \to \varphi_{187} \to \varphi_{181} \to \varphi_{188}$\lbl{A1}}
\BIC{$\varphi_{187} \to \varphi_{181} \to \varphi_{188}$}
\BIC{$\varphi_{181} \to \varphi_{188}$}
\BIC{$\varphi_{188}$}
\AXC{$\stackrel{\Pi_{29}}{\vdots}$}\noLine
\UIC{$\varphi_{188} \to \varphi_{180} \to \varphi_{182}$}
\BIC{$\varphi_{180} \to \varphi_{182}$}
\BIC{$\varphi_{182}$}
\end{ptree}
\noindent$\Pi_{31}$:
\begin{ptree}
\AXC{$\varphi_{175} \to \forall n\,(\varphi_{175})$\lbl{A4}}
\AXC{$\varphi_{234}$\lbl{A3}}
\AXC{$\stackrel{\Pi_{30}}{\vdots}$}\noLine
\UIC{$\varphi_{182}$}
\AXC{$\varphi_{182} \to \varphi_{175} \to n \in \mathbb{N} \to K \leq n \to \lvert a_{n} - c \rvert < \varepsilon$\lbl{A1}}
\BIC{$\varphi_{175} \to n \in \mathbb{N} \to K \leq n \to \lvert a_{n} - c \rvert < \varepsilon$}
\UIC{$\forall n\,(\varphi_{175} \to n \in \mathbb{N} \to K \leq n \to \lvert a_{n} - c \rvert < \varepsilon)$}
\AXC{$\varphi_{265}$\lbl{A1}}
\BIC{$\varphi_{234} \to (\varphi_{175} \to \forall n\,(\varphi_{175})) \to \varphi_{186}$}
\BIC{$(\varphi_{175} \to \forall n\,(\varphi_{175})) \to \varphi_{186}$}
\BIC{$\varphi_{186}$}
\end{ptree}
\noindent$\Pi_{32}$:
\begin{ptree}
\AXC{$\stackrel{\Pi_{31}}{\vdots}$}\noLine
\UIC{$\varphi_{186}$}
\AXC{$\mathrm{RFn}(f) \to A \subseteq \mathrm{dom}(f) \to c \in A \to \varphi_{173}$\lbl{A1}}
\AXC{$\varphi_{227}$\lbl{A1}}
\BIC{$\mathrm{RFn}(f) \to A \subseteq \mathrm{dom}(f) \to \varphi_{167}$}
\AXC{$(\mathrm{RFn}(f) \to A \subseteq \mathrm{dom}(f) \to \varphi_{167}) \to \varphi_{186} \to \varphi_{194}$\lbl{A1}}
\BIC{$\varphi_{186} \to \varphi_{194}$}
\BIC{$\varphi_{194}$}
\end{ptree}
\noindent$\Pi_{33}$:
\begin{ptree}
\AXC{$\varphi_{2} \to \neg (K \in \mathbb{N} \land \forall n \in \mathbb{N}\,(K \leq n \to \lvert a_{n} - c \rvert < \varepsilon))$\lbl{A2}}
\AXC{$\stackrel{\Pi_{32}}{\vdots}$}\noLine
\UIC{$\varphi_{194}$}
\AXC{$\varphi_{194} \to \varphi_{203}$\lbl{A1}}
\BIC{$\varphi_{203}$}
\BIC{$\varphi_{197}$}
\AXC{$\varphi_{197} \to \varphi_{179} \to \neg (K \in \mathbb{N} \land \underline{K}^{-1} < \varepsilon)$\lbl{A1}}
\BIC{$\varphi_{179} \to \neg (K \in \mathbb{N} \land \underline{K}^{-1} < \varepsilon)$}
\UIC{$\forall K\,(\varphi_{179} \to \neg (K \in \mathbb{N} \land \underline{K}^{-1} < \varepsilon))$}
\end{ptree}
\noindent$\Pi_{34}$:
\begin{ptree}
\AXC{$\varphi_{163}$\lbl{A1}}
\AXC{$\varphi_{157}$\lbl{A1}}
\AXC{$\forall \varepsilon \in \mathbb{R}\,(0 < \varepsilon \to \exists K \in \mathbb{N}\,(\underline{K}^{-1} < \varepsilon))$\lbl{[C24]}}
\AXC{$\varphi_{55}$\lbl{A2}}
\BIC{$\varepsilon \in \mathbb{R} \to 0 < \varepsilon \to \exists K \in \mathbb{N}\,(\underline{K}^{-1} < \varepsilon)$}
\AXC{$\varphi_{252}$\lbl{A1}}
\BIC{$\varphi_{157} \to \varphi_{163} \to \mathrm{RFn}(f) \to A \subseteq \mathrm{dom}(f) \to \varphi_{164}$}
\BIC{$\varphi_{163} \to \mathrm{RFn}(f) \to A \subseteq \mathrm{dom}(f) \to \varphi_{164}$}
\BIC{$\mathrm{RFn}(f) \to A \subseteq \mathrm{dom}(f) \to \varphi_{164}$}
\AXC{$\varphi_{268}$\lbl{A1}}
\BIC{$\forall K\,(\varphi_{179} \to \neg (K \in \mathbb{N} \land \underline{K}^{-1} < \varepsilon)) \to \varphi_{238} \to \varphi_{254}$}
\end{ptree}
\noindent$\Pi_{35}$:
\begin{ptree}
\AXC{$\varphi_{179} \to \forall K\,(\varphi_{179})$\lbl{A4}}
\AXC{$\varphi_{238}$\lbl{A3}}
\AXC{$\stackrel{\Pi_{33}}{\vdots}$}\noLine
\UIC{$\forall K\,(\varphi_{179} \to \neg (K \in \mathbb{N} \land \underline{K}^{-1} < \varepsilon))$}
\AXC{$\stackrel{\Pi_{34}}{\vdots}$}\noLine
\UIC{$\forall K\,(\varphi_{179} \to \neg (K \in \mathbb{N} \land \underline{K}^{-1} < \varepsilon)) \to \varphi_{238} \to \varphi_{254}$}
\BIC{$\varphi_{238} \to \varphi_{254}$}
\BIC{$\varphi_{254}$}
\BIC{$\mathrm{RFn}(f) \to A \subseteq \mathrm{dom}(f) \to c \in A \to \varphi_{172}$}
\end{ptree}
\noindent$\Pi_{36}$:
\begin{ptree}
\AXC{$\varphi_{231}$\lbl{A3}}
\AXC{$\stackrel{\Pi_{35}}{\vdots}$}\noLine
\UIC{$\mathrm{RFn}(f) \to A \subseteq \mathrm{dom}(f) \to c \in A \to \varphi_{172}$}
\AXC{$\varphi_{229}$\lbl{A1}}
\BIC{$\varphi_{97} \land (\mathrm{Seq}(a) \land \varphi_{45}) \to \varphi_{24}$}
\UIC{$\forall \varepsilon\,(\varphi_{97} \land (\mathrm{Seq}(a) \land \varphi_{45}) \to \varphi_{24})$}
\AXC{$\forall \varepsilon\,(\varphi_{97} \land (\mathrm{Seq}(a) \land \varphi_{45}) \to \varphi_{24}) \to \varphi_{231} \to \varphi_{250}$\lbl{A1}}
\BIC{$\varphi_{231} \to \varphi_{250}$}
\BIC{$\varphi_{250}$}
\end{ptree}
\noindent$\Pi_{37}$:
\begin{ptree}
\AXC{$\forall X\,\forall Y\,\forall x\,(X \subseteq Y \to x \in X \to x \in Y)$\lbl{[B5]}}
\AXC{$\varphi_{5}$\lbl{A2}}
\BIC{$\forall Y\,\forall x\,(A \subseteq Y \to x \in A \to x \in Y)$}
\AXC{$\varphi_{10}$\lbl{A2}}
\BIC{$\forall x\,(A \subseteq \mathrm{dom}(f) \to x \in A \to x \in \mathrm{dom}(f))$}
\AXC{$\varphi_{16}$\lbl{A2}}
\BIC{$A \subseteq \mathrm{dom}(f) \to c \in A \to c \in \mathrm{dom}(f)$}
\AXC{$\varphi_{81}$\lbl{A1}}
\BIC{$(\mathrm{RFn}(f) \to A \subseteq \mathrm{dom}(f) \to A \subseteq \mathrm{dom}(f)) \to \varphi_{63}$}
\end{ptree}
\noindent$\Pi_{38}$:
\begin{ptree}
\AXC{$\mathrm{RFn}(f) \to A \subseteq \mathrm{dom}(f) \to A \subseteq \mathrm{dom}(f)$\lbl{A1}}
\AXC{$\stackrel{\Pi_{37}}{\vdots}$}\noLine
\UIC{$(\mathrm{RFn}(f) \to A \subseteq \mathrm{dom}(f) \to A \subseteq \mathrm{dom}(f)) \to \varphi_{63}$}
\BIC{$\varphi_{63}$}
\end{ptree}
\noindent$\Pi_{39}$:
\begin{ptree}
\AXC{$\mathrm{RFn}(f) \to \mathrm{RFn}(f)$\lbl{A1}}
\AXC{$\forall f\,\forall x\,(\mathrm{RFn}(f) \to x \in \mathrm{dom}(f) \to x \in \mathbb{R})$\lbl{[S7]}}
\AXC{$\varphi_{13}$\lbl{A2}}
\BIC{$\forall x\,(\mathrm{RFn}(f) \to x \in \mathrm{dom}(f) \to x \in \mathbb{R})$}
\AXC{$\varphi_{6}$\lbl{A2}}
\BIC{$\mathrm{RFn}(f) \to c \in \mathrm{dom}(f) \to c \in \mathbb{R}$}
\AXC{$\varphi_{85}$\lbl{A1}}
\BIC{$(\mathrm{RFn}(f) \to \mathrm{RFn}(f)) \to (\mathrm{RFn}(f) \to \varphi_{52}) \to \varphi_{56}$}
\BIC{$(\mathrm{RFn}(f) \to \varphi_{52}) \to \varphi_{56}$}
\end{ptree}
\noindent$\Pi_{40}$:
\begin{ptree}
\AXC{$\mathrm{RFn}(f) \to A \subseteq \mathrm{dom}(f) \to c \in A \to c \in A$\lbl{A1}}
\AXC{$\stackrel{\Pi_{38}}{\vdots}$}\noLine
\UIC{$\varphi_{63}$}
\BIC{$\mathrm{RFn}(f) \to \varphi_{52}$}
\AXC{$\stackrel{\Pi_{39}}{\vdots}$}\noLine
\UIC{$(\mathrm{RFn}(f) \to \varphi_{52}) \to \varphi_{56}$}
\BIC{$\varphi_{56}$}
\end{ptree}
\noindent$\Pi_{41}$:
\begin{ptree}
\AXC{$\forall a\,(\varphi_{58})$\lbl{DefConv}}
\AXC{$\forall a\,(\varphi_{58}) \to \varphi_{58}$\lbl{A2}}
\BIC{$\varphi_{58}$}
\AXC{$\varphi_{58} \to (\mathrm{Conv}(a,c) \leftrightarrow \mathrm{Seq}(a) \land c \in \mathbb{R} \land \varphi_{30})$\lbl{A2}}
\BIC{$\mathrm{Conv}(a,c) \leftrightarrow \mathrm{Seq}(a) \land c \in \mathbb{R} \land \varphi_{30}$}
\AXC{$\varphi_{257}$\lbl{A1}}
\BIC{$\varphi_{147} \to \varphi_{56} \to (\mathrm{RFn}(f) \to A \subseteq \mathrm{dom}(f) \to c \in A \to \varphi_{176}) \to \varphi_{149}$}
\end{ptree}
\noindent$\Pi_{42}$:
\begin{ptree}
\AXC{$\stackrel{\Pi_{40}}{\vdots}$}\noLine
\UIC{$\varphi_{56}$}
\AXC{$\varphi_{193}$\lbl{A1}}
\AXC{$\varphi_{193} \to \varphi_{147}$\lbl{A1}}
\BIC{$\varphi_{147}$}
\AXC{$\stackrel{\Pi_{41}}{\vdots}$}\noLine
\UIC{$\varphi_{147} \to \varphi_{56} \to (\mathrm{RFn}(f) \to A \subseteq \mathrm{dom}(f) \to c \in A \to \varphi_{176}) \to \varphi_{149}$}
\BIC{$\varphi_{56} \to (\mathrm{RFn}(f) \to A \subseteq \mathrm{dom}(f) \to c \in A \to \varphi_{176}) \to \varphi_{149}$}
\BIC{$(\mathrm{RFn}(f) \to A \subseteq \mathrm{dom}(f) \to c \in A \to \varphi_{176}) \to \varphi_{149}$}
\end{ptree}
\noindent$\Pi_{43}$:
\begin{ptree}
\AXC{$\varphi_{97} \land (\mathrm{Seq}(a) \land \varphi_{45}) \to \forall \varepsilon\,(\varphi_{97} \land (\mathrm{Seq}(a) \land \varphi_{45}))$\lbl{A4}}
\AXC{$\stackrel{\Pi_{36}}{\vdots}$}\noLine
\UIC{$\varphi_{250}$}
\BIC{$\mathrm{RFn}(f) \to A \subseteq \mathrm{dom}(f) \to c \in A \to \varphi_{176}$}
\AXC{$\stackrel{\Pi_{42}}{\vdots}$}\noLine
\UIC{$(\mathrm{RFn}(f) \to A \subseteq \mathrm{dom}(f) \to c \in A \to \varphi_{176}) \to \varphi_{149}$}
\BIC{$\varphi_{149}$}
\end{ptree}
\noindent$\Pi_{44}$:
\begin{ptree}
\AXC{$\forall a\,(\mathrm{Seq}(a) \to \varphi_{3}) \to \mathrm{Seq}(a) \to \varphi_{3}$\lbl{A2}}
\AXC{$\mathrm{RFn}(f) \to \varphi_{68}$\lbl{A1}}
\AXC{$(\mathrm{RFn}(f) \to \varphi_{68}) \to \varphi_{101}$\lbl{A1}}
\BIC{$\varphi_{101}$}
\BIC{$\mathrm{RFn}(f) \to \varphi_{67}$}
\end{ptree}
\noindent$\Pi_{45}$:
\begin{ptree}
\AXC{$\stackrel{\Pi_{43}}{\vdots}$}\noLine
\UIC{$\varphi_{149}$}
\AXC{$\stackrel{\Pi_{10}}{\vdots}$}\noLine
\UIC{$\varphi_{159}$}
\AXC{$\varphi_{193}$\lbl{A1}}
\AXC{$\varphi_{193} \to \varphi_{147}$\lbl{A1}}
\BIC{$\varphi_{147}$}
\AXC{$\stackrel{\Pi_{44}}{\vdots}$}\noLine
\UIC{$\mathrm{RFn}(f) \to \varphi_{67}$}
\AXC{$(\mathrm{RFn}(f) \to \varphi_{67}) \to \varphi_{147} \to \varphi_{159} \to \varphi_{149} \to \varphi_{158}$\lbl{A1}}
\BIC{$\varphi_{147} \to \varphi_{159} \to \varphi_{149} \to \varphi_{158}$}
\BIC{$\varphi_{159} \to \varphi_{149} \to \varphi_{158}$}
\BIC{$\varphi_{149} \to \varphi_{158}$}
\BIC{$\varphi_{158}$}
\end{ptree}
\noindent$\Pi_{46}$:
\begin{ptree}
\AXC{$\varphi_{46}$\lbl{[C3]}}
\AXC{$\varphi_{46} \to \varphi_{57}$\lbl{A2}}
\BIC{$\varphi_{57}$}
\AXC{$\varphi_{57} \to \varphi_{59}$\lbl{A2}}
\BIC{$\varphi_{59}$}
\AXC{$\varphi_{59} \to \mathrm{Conv}(f \circ a,f(c)) \to e_{0} \in \mathbb{R} \to 0 < e_{0} \to \varphi_{11}$\lbl{A2}}
\BIC{$\mathrm{Conv}(f \circ a,f(c)) \to e_{0} \in \mathbb{R} \to 0 < e_{0} \to \varphi_{11}$}
\AXC{$\varphi_{259}$\lbl{A1}}
\BIC{$\varphi_{158} \to (\mathrm{RFn}(f) \to A \subseteq \mathrm{dom}(f) \to c \in A \to \varphi_{89}) \to \varphi_{213}$}
\end{ptree}
\noindent$\Pi_{47}$:
\begin{ptree}
\AXC{$\mathrm{RFn}(f) \to A \subseteq \mathrm{dom}(f) \to \varphi_{123}$\lbl{A1}}
\AXC{$\varphi_{211}$\lbl{A1}}
\BIC{$\mathrm{RFn}(f) \to A \subseteq \mathrm{dom}(f) \to c \in A \to \varphi_{89}$}
\AXC{$\stackrel{\Pi_{45}}{\vdots}$}\noLine
\UIC{$\varphi_{158}$}
\AXC{$\stackrel{\Pi_{46}}{\vdots}$}\noLine
\UIC{$\varphi_{158} \to (\mathrm{RFn}(f) \to A \subseteq \mathrm{dom}(f) \to c \in A \to \varphi_{89}) \to \varphi_{213}$}
\BIC{$(\mathrm{RFn}(f) \to A \subseteq \mathrm{dom}(f) \to c \in A \to \varphi_{89}) \to \varphi_{213}$}
\BIC{$\varphi_{213}$}
\end{ptree}
\noindent$\Pi_{48}$:
\begin{ptree}
\AXC{$\mathrm{RFn}(f) \to A \subseteq \mathrm{dom}(f) \to \varphi_{123}$\lbl{A1}}
\AXC{$\varphi_{210}$\lbl{A1}}
\BIC{$\mathrm{RFn}(f) \to A \subseteq \mathrm{dom}(f) \to c \in A \to \varphi_{88}$}
\AXC{$\stackrel{\Pi_{47}}{\vdots}$}\noLine
\UIC{$\varphi_{213}$}
\BIC{$\mathrm{RFn}(f) \to A \subseteq \mathrm{dom}(f) \to c \in A \to \varphi_{166}$}
\AXC{$(\mathrm{RFn}(f) \to A \subseteq \mathrm{dom}(f) \to c \in A \to \varphi_{166}) \to \varphi_{264}$\lbl{A1}}
\BIC{$\varphi_{264}$}
\end{ptree}
\noindent$\Pi_{49}$:
\begin{ptree}
\AXC{$\varphi_{220}$\lbl{A4}}
\AXC{$\varphi_{230}$\lbl{A3}}
\AXC{$\stackrel{\Pi_{20}}{\vdots}$}\noLine
\UIC{$\forall N\,(\varphi_{97} \land (\mathrm{Seq}(a) \land \varphi_{45}) \land \neg 0 < 0 \to \neg \varphi_{9})$}
\AXC{$\stackrel{\Pi_{48}}{\vdots}$}\noLine
\UIC{$\varphi_{264}$}
\BIC{$\varphi_{230} \to \varphi_{220} \to \mathrm{RFn}(f) \to A \subseteq \mathrm{dom}(f) \to c \in A \to \varphi_{122}$}
\BIC{$\varphi_{220} \to \mathrm{RFn}(f) \to A \subseteq \mathrm{dom}(f) \to c \in A \to \varphi_{122}$}
\BIC{$\mathrm{RFn}(f) \to A \subseteq \mathrm{dom}(f) \to c \in A \to \varphi_{122}$}
\AXC{$\varphi_{219}$\lbl{A1}}
\BIC{$\varphi_{97} \land \neg 0 < 0 \to \neg (\mathrm{Seq}(a) \land \varphi_{45})$}
\UIC{$\forall a\,(\varphi_{97} \land \neg 0 < 0 \to \neg (\mathrm{Seq}(a) \land \varphi_{45}))$}
\end{ptree}
\noindent$\Pi_{50}$:
\begin{ptree}
\AXC{$\forall X\,\forall Y\,\forall x\,(X \subseteq Y \to x \in X \to x \in Y)$\lbl{[B5]}}
\AXC{$\varphi_{5}$\lbl{A2}}
\BIC{$\forall Y\,\forall x\,(A \subseteq Y \to x \in A \to x \in Y)$}
\AXC{$\varphi_{10}$\lbl{A2}}
\BIC{$\forall x\,(A \subseteq \mathrm{dom}(f) \to x \in A \to x \in \mathrm{dom}(f))$}
\AXC{$\varphi_{17}$\lbl{A2}}
\BIC{$A \subseteq \mathrm{dom}(f) \to x \in A \to x \in \mathrm{dom}(f)$}
\end{ptree}
\noindent$\Pi_{51}$:
\begin{ptree}
\AXC{$\stackrel{\Pi_{50}}{\vdots}$}\noLine
\UIC{$A \subseteq \mathrm{dom}(f) \to x \in A \to x \in \mathrm{dom}(f)$}
\AXC{$(A \subseteq \mathrm{dom}(f) \to x \in A \to x \in \mathrm{dom}(f)) \to \varphi_{217}$\lbl{A1}}
\BIC{$\varphi_{217}$}
\end{ptree}
\noindent$\Pi_{52}$:
\begin{ptree}
\AXC{$\mathrm{RFn}(f) \to A \subseteq \mathrm{dom}(f) \to c \in A \to \varphi_{146}$\lbl{A1}}
\AXC{$\varphi_{218}$\lbl{A1}}
\BIC{$\mathrm{RFn}(f) \to A \subseteq \mathrm{dom}(f) \to c \in A \to \varphi_{116}$}
\AXC{$\mathrm{RFn}(f) \to A \subseteq \mathrm{dom}(f) \to A \subseteq \mathrm{dom}(f)$\lbl{A1}}
\AXC{$\stackrel{\Pi_{51}}{\vdots}$}\noLine
\UIC{$\varphi_{217}$}
\BIC{$(\mathrm{RFn}(f) \to A \subseteq \mathrm{dom}(f) \to c \in A \to \varphi_{116}) \to \varphi_{127}$}
\BIC{$\varphi_{127}$}
\end{ptree}
\noindent$\Pi_{53}$:
\begin{ptree}
\AXC{$\forall f\,\forall x\,(\mathrm{RFn}(f) \to x \in \mathrm{dom}(f) \to x \in \mathbb{R})$\lbl{[S7]}}
\AXC{$\varphi_{13}$\lbl{A2}}
\BIC{$\forall x\,(\mathrm{RFn}(f) \to x \in \mathrm{dom}(f) \to x \in \mathbb{R})$}
\AXC{$\varphi_{7}$\lbl{A2}}
\BIC{$\mathrm{RFn}(f) \to x \in \mathrm{dom}(f) \to x \in \mathbb{R}$}
\AXC{$(\mathrm{RFn}(f) \to x \in \mathrm{dom}(f) \to x \in \mathbb{R}) \to \varphi_{215}$\lbl{A1}}
\BIC{$\varphi_{215}$}
\end{ptree}
\noindent$\Pi_{54}$:
\begin{ptree}
\AXC{$\mathrm{RFn}(f) \to A \subseteq \mathrm{dom}(f) \to c \in A \to \varphi_{146}$\lbl{A1}}
\AXC{$(\mathrm{RFn}(f) \to A \subseteq \mathrm{dom}(f) \to c \in A \to \varphi_{146}) \to \varphi_{222}$\lbl{A1}}
\BIC{$\varphi_{222}$}
\end{ptree}
\noindent$\Pi_{55}$:
\begin{ptree}
\AXC{$\stackrel{\Pi_{52}}{\vdots}$}\noLine
\UIC{$\varphi_{127}$}
\AXC{$\mathrm{RFn}(f) \to \mathrm{RFn}(f)$\lbl{A1}}
\AXC{$\stackrel{\Pi_{53}}{\vdots}$}\noLine
\UIC{$\varphi_{215}$}
\BIC{$\varphi_{127} \to \mathrm{RFn}(f) \to A \subseteq \mathrm{dom}(f) \to c \in A \to \varphi_{117}$}
\BIC{$\mathrm{RFn}(f) \to A \subseteq \mathrm{dom}(f) \to c \in A \to \varphi_{117}$}
\AXC{$\stackrel{\Pi_{54}}{\vdots}$}\noLine
\UIC{$\varphi_{222}$}
\BIC{$\mathrm{RFn}(f) \to A \subseteq \mathrm{dom}(f) \to c \in A \to \varphi_{152}$}
\AXC{$(\mathrm{RFn}(f) \to A \subseteq \mathrm{dom}(f) \to c \in A \to \varphi_{152}) \to \varphi_{202}$\lbl{A1}}
\BIC{$\varphi_{202}$}
\end{ptree}
\noindent$\Pi_{56}$:
\begin{ptree}
\AXC{$\forall n \in \mathbb{N}\,(0 < \underline{n})$\lbl{[N15]}}
\AXC{$\forall n \in \mathbb{N}\,(0 < \underline{n}) \to k \in \mathbb{N} \to 0 < \underline{k}$\lbl{A2}}
\BIC{$k \in \mathbb{N} \to 0 < \underline{k}$}
\AXC{$(k \in \mathbb{N} \to 0 < \underline{k}) \to \varphi_{204}$\lbl{A1}}
\BIC{$\varphi_{204}$}
\end{ptree}
\noindent$\Pi_{57}$:
\begin{ptree}
\AXC{$\mathrm{RFn}(f) \to A \subseteq \mathrm{dom}(f) \to c \in A \to \varphi_{91}$\lbl{A1}}
\AXC{$\stackrel{\Pi_{56}}{\vdots}$}\noLine
\UIC{$\varphi_{204}$}
\BIC{$\mathrm{RFn}(f) \to A \subseteq \mathrm{dom}(f) \to c \in A \to \varphi_{94}$}
\end{ptree}
\noindent$\Pi_{58}$:
\begin{ptree}
\AXC{$\forall n \in \mathbb{N}\,(\underline{n} \in \mathbb{R})$\lbl{[N4]}}
\AXC{$\forall n \in \mathbb{N}\,(\underline{n} \in \mathbb{R}) \to k \in \mathbb{N} \to \underline{k} \in \mathbb{R}$\lbl{A2}}
\BIC{$k \in \mathbb{N} \to \underline{k} \in \mathbb{R}$}
\AXC{$(k \in \mathbb{N} \to \underline{k} \in \mathbb{R}) \to \varphi_{206}$\lbl{A1}}
\BIC{$\varphi_{206}$}
\end{ptree}
\noindent$\Pi_{59}$:
\begin{ptree}
\AXC{$\mathrm{RFn}(f) \to A \subseteq \mathrm{dom}(f) \to c \in A \to \varphi_{91}$\lbl{A1}}
\AXC{$\stackrel{\Pi_{58}}{\vdots}$}\noLine
\UIC{$\varphi_{206}$}
\BIC{$\mathrm{RFn}(f) \to A \subseteq \mathrm{dom}(f) \to c \in A \to \varphi_{95}$}
\end{ptree}
\noindent$\Pi_{60}$:
\begin{ptree}
\AXC{$\stackrel{\Pi_{59}}{\vdots}$}\noLine
\UIC{$\mathrm{RFn}(f) \to A \subseteq \mathrm{dom}(f) \to c \in A \to \varphi_{95}$}
\AXC{$\forall x \in \mathbb{R}\,(0 < x \to 0 < x^{-1})$\lbl{[F34]}}
\AXC{$\varphi_{28}$\lbl{A2}}
\BIC{$\underline{k} \in \mathbb{R} \to 0 < \underline{k} \to 0 < \underline{k}^{-1}$}
\AXC{$\varphi_{239}$\lbl{A1}}
\BIC{$(\mathrm{RFn}(f) \to A \subseteq \mathrm{dom}(f) \to c \in A \to \varphi_{95}) \to \varphi_{207}$}
\BIC{$\varphi_{207}$}
\end{ptree}
\noindent$\Pi_{61}$:
\begin{ptree}
\AXC{$x = x$\lbl{A5}}
\AXC{$x = y \to x = x \to y = x$\lbl{A6}}
\AXC{$(x = y \to x = x \to y = x) \to x = x \to x = y \to y = x$\lbl{A1}}
\BIC{$x = x \to x = y \to y = x$}
\BIC{$x = y \to y = x$}
\UIC{$\forall y\,(x = y \to y = x)$}
\UIC{$\forall x\,\forall y\,(x = y \to y = x)$}
\end{ptree}
\noindent$\Pi_{62}$:
\begin{ptree}
\AXC{$\stackrel{\Pi_{61}}{\vdots}$}\noLine
\UIC{$\forall x\,\forall y\,(x = y \to y = x)$}
\AXC{$\forall x\,\forall y\,(x = y \to y = x) \to \forall y\,(\underline{k} = y \to y = \underline{k})$\lbl{A2}}
\BIC{$\forall y\,(\underline{k} = y \to y = \underline{k})$}
\end{ptree}
\noindent$\Pi_{63}$:
\begin{ptree}
\AXC{$\forall x\,\forall y\,(x < y \to x \neq y)$\lbl{[F26]}}
\AXC{$\forall x\,\forall y\,(x < y \to x \neq y) \to \forall y\,(0 < y \to 0 \neq y)$\lbl{A2}}
\BIC{$\forall y\,(0 < y \to 0 \neq y)$}
\AXC{$\forall y\,(0 < y \to 0 \neq y) \to 0 < \underline{k} \to 0 \neq \underline{k}$\lbl{A2}}
\BIC{$0 < \underline{k} \to 0 \neq \underline{k}$}
\end{ptree}
\noindent$\Pi_{64}$:
\begin{ptree}
\AXC{$\stackrel{\Pi_{57}}{\vdots}$}\noLine
\UIC{$\mathrm{RFn}(f) \to A \subseteq \mathrm{dom}(f) \to c \in A \to \varphi_{94}$}
\AXC{$\stackrel{\Pi_{63}}{\vdots}$}\noLine
\UIC{$0 < \underline{k} \to 0 \neq \underline{k}$}
\AXC{$(0 < \underline{k} \to 0 \neq \underline{k}) \to \varphi_{205}$\lbl{A1}}
\BIC{$\varphi_{205}$}
\BIC{$\mathrm{RFn}(f) \to A \subseteq \mathrm{dom}(f) \to c \in A \to \varphi_{92}$}
\end{ptree}
\noindent$\Pi_{65}$:
\begin{ptree}
\AXC{$\stackrel{\Pi_{64}}{\vdots}$}\noLine
\UIC{$\mathrm{RFn}(f) \to A \subseteq \mathrm{dom}(f) \to c \in A \to \varphi_{92}$}
\AXC{$(\mathrm{RFn}(f) \to A \subseteq \mathrm{dom}(f) \to c \in A \to \varphi_{92}) \to \varphi_{112}$\lbl{A1}}
\BIC{$\varphi_{112}$}
\end{ptree}
\noindent$\Pi_{66}$:
\begin{ptree}
\AXC{$\stackrel{\Pi_{59}}{\vdots}$}\noLine
\UIC{$\mathrm{RFn}(f) \to A \subseteq \mathrm{dom}(f) \to c \in A \to \varphi_{95}$}
\AXC{$\forall x \in \mathbb{R}\,(x \neq 0 \to x^{-1} \in \mathbb{R})$\lbl{(R1)}}
\AXC{$\varphi_{37}$\lbl{A2}}
\BIC{$\underline{k} \in \mathbb{R} \to \underline{k} \neq 0 \to \underline{k}^{-1} \in \mathbb{R}$}
\AXC{$\varphi_{240}$\lbl{A1}}
\BIC{$(\mathrm{RFn}(f) \to A \subseteq \mathrm{dom}(f) \to c \in A \to \varphi_{95}) \to \varphi_{208}$}
\BIC{$\varphi_{208}$}
\end{ptree}
\noindent$\Pi_{67}$:
\begin{ptree}
\AXC{$\stackrel{\Pi_{62}}{\vdots}$}\noLine
\UIC{$\forall y\,(\underline{k} = y \to y = \underline{k})$}
\AXC{$\forall y\,(\underline{k} = y \to y = \underline{k}) \to \underline{k} = 0 \to 0 = \underline{k}$\lbl{A2}}
\BIC{$\underline{k} = 0 \to 0 = \underline{k}$}
\AXC{$\stackrel{\Pi_{65}}{\vdots}$}\noLine
\UIC{$\varphi_{112}$}
\BIC{$\mathrm{RFn}(f) \to A \subseteq \mathrm{dom}(f) \to c \in A \to \varphi_{93}$}
\AXC{$\stackrel{\Pi_{66}}{\vdots}$}\noLine
\UIC{$\varphi_{208}$}
\BIC{$\mathrm{RFn}(f) \to A \subseteq \mathrm{dom}(f) \to c \in A \to \varphi_{98}$}
\end{ptree}
\noindent$\Pi_{68}$:
\begin{ptree}
\AXC{$\varphi_{34} \to \neg (\underline{k}^{-1} \in \mathbb{R} \land 0 < \underline{k}^{-1} \land \varphi_{8})$\lbl{A2}}
\AXC{$\mathrm{RFn}(f) \to A \subseteq \mathrm{dom}(f) \to \varphi_{123}$\lbl{A1}}
\AXC{$\varphi_{214}$\lbl{A1}}
\BIC{$\mathrm{RFn}(f) \to A \subseteq \mathrm{dom}(f) \to \varphi_{119}$}
\AXC{$\varphi_{232}$\lbl{A1}}
\BIC{$(\varphi_{34} \to \neg (\underline{k}^{-1} \in \mathbb{R} \land 0 < \underline{k}^{-1} \land \varphi_{8})) \to \varphi_{142}$}
\BIC{$\varphi_{142}$}
\end{ptree}
\noindent$\Pi_{69}$:
\begin{ptree}
\AXC{$\stackrel{\Pi_{67}}{\vdots}$}\noLine
\UIC{$\mathrm{RFn}(f) \to A \subseteq \mathrm{dom}(f) \to c \in A \to \varphi_{98}$}
\AXC{$\stackrel{\Pi_{68}}{\vdots}$}\noLine
\UIC{$\varphi_{142}$}
\AXC{$\varphi_{142} \to (\mathrm{RFn}(f) \to A \subseteq \mathrm{dom}(f) \to c \in A \to \varphi_{98}) \to \varphi_{209}$\lbl{A1}}
\BIC{$(\mathrm{RFn}(f) \to A \subseteq \mathrm{dom}(f) \to c \in A \to \varphi_{98}) \to \varphi_{209}$}
\BIC{$\varphi_{209}$}
\end{ptree}
\noindent$\Pi_{70}$:
\begin{ptree}
\AXC{$\stackrel{\Pi_{57}}{\vdots}$}\noLine
\UIC{$\mathrm{RFn}(f) \to A \subseteq \mathrm{dom}(f) \to c \in A \to \varphi_{94}$}
\AXC{$\stackrel{\Pi_{60}}{\vdots}$}\noLine
\UIC{$\varphi_{207}$}
\BIC{$\mathrm{RFn}(f) \to A \subseteq \mathrm{dom}(f) \to c \in A \to \varphi_{96}$}
\AXC{$\stackrel{\Pi_{69}}{\vdots}$}\noLine
\UIC{$\varphi_{209}$}
\BIC{$\mathrm{RFn}(f) \to A \subseteq \mathrm{dom}(f) \to c \in A \to \varphi_{113}$}
\AXC{$\varphi_{249}$\lbl{A1}}
\BIC{$\forall x\,(\varphi_{60}) \to (\forall x\,(\varphi_{60}) \to \forall x\,\forall x\,(\neg \varphi_{4}) \to \varphi_{8}) \to \varphi_{199}$}
\end{ptree}
\noindent$\Pi_{71}$:
\begin{ptree}
\AXC{$\forall x\,(\neg \varphi_{4}) \to \neg \varphi_{4}$\lbl{A2}}
\AXC{$(\forall x\,(\neg \varphi_{4}) \to \neg \varphi_{4}) \to \varphi_{60}$\lbl{A1}}
\BIC{$\varphi_{60}$}
\UIC{$\forall x\,(\varphi_{60})$}
\AXC{$\stackrel{\Pi_{70}}{\vdots}$}\noLine
\UIC{$\forall x\,(\varphi_{60}) \to (\forall x\,(\varphi_{60}) \to \forall x\,\forall x\,(\neg \varphi_{4}) \to \varphi_{8}) \to \varphi_{199}$}
\BIC{$(\forall x\,(\varphi_{60}) \to \forall x\,\forall x\,(\neg \varphi_{4}) \to \varphi_{8}) \to \varphi_{199}$}
\end{ptree}
\noindent$\Pi_{72}$:
\begin{ptree}
\AXC{$\forall x\,(\neg \varphi_{4}) \to \forall x\,\forall x\,(\neg \varphi_{4})$\lbl{A4}}
\AXC{$\forall x\,(\varphi_{60}) \to \forall x\,\forall x\,(\neg \varphi_{4}) \to \varphi_{8}$\lbl{A3}}
\AXC{$\stackrel{\Pi_{71}}{\vdots}$}\noLine
\UIC{$(\forall x\,(\varphi_{60}) \to \forall x\,\forall x\,(\neg \varphi_{4}) \to \varphi_{8}) \to \varphi_{199}$}
\BIC{$\varphi_{199}$}
\BIC{$\mathrm{RFn}(f) \to A \subseteq \mathrm{dom}(f) \to c \in A \to \varphi_{114}$}
\AXC{$\varphi_{266}$\lbl{A1}}
\BIC{$\forall x\,(\varphi_{97} \land k \in \mathbb{N} \land \neg \varphi_{27} \to \neg \varphi_{4}) \to \varphi_{225} \to \varphi_{246}$}
\end{ptree}
\noindent$\Pi_{73}$:
\begin{ptree}
\AXC{$\varphi_{225}$\lbl{A3}}
\AXC{$\forall x\,(\neg \varphi_{21}) \to \neg \varphi_{21}$\lbl{A2}}
\AXC{$\stackrel{\Pi_{55}}{\vdots}$}\noLine
\UIC{$\varphi_{202}$}
\BIC{$\mathrm{RFn}(f) \to A \subseteq \mathrm{dom}(f) \to c \in A \to \varphi_{153}$}
\AXC{$\varphi_{224}$\lbl{A1}}
\BIC{$\varphi_{97} \land k \in \mathbb{N} \land \neg \varphi_{27} \to \neg \varphi_{4}$}
\UIC{$\forall x\,(\varphi_{97} \land k \in \mathbb{N} \land \neg \varphi_{27} \to \neg \varphi_{4})$}
\AXC{$\stackrel{\Pi_{72}}{\vdots}$}\noLine
\UIC{$\forall x\,(\varphi_{97} \land k \in \mathbb{N} \land \neg \varphi_{27} \to \neg \varphi_{4}) \to \varphi_{225} \to \varphi_{246}$}
\BIC{$\varphi_{225} \to \varphi_{246}$}
\BIC{$\varphi_{246}$}
\end{ptree}
\noindent$\Pi_{74}$:
\begin{ptree}
\AXC{$\varphi_{97} \land k \in \mathbb{N} \land \neg \varphi_{27} \to \forall x\,(\varphi_{97} \land k \in \mathbb{N} \land \neg \varphi_{27})$\lbl{A4}}
\AXC{$\stackrel{\Pi_{73}}{\vdots}$}\noLine
\UIC{$\varphi_{246}$}
\BIC{$\mathrm{RFn}(f) \to A \subseteq \mathrm{dom}(f) \to c \in A \to \varphi_{118}$}
\AXC{$\varphi_{212}$\lbl{A1}}
\BIC{$\varphi_{97} \to k \in \mathbb{N} \to \varphi_{27}$}
\UIC{$\forall k\,(\varphi_{97} \to k \in \mathbb{N} \to \varphi_{27})$}
\AXC{$\forall k\,(\varphi_{97} \to k \in \mathbb{N} \to \varphi_{27}) \to \varphi_{260}$\lbl{A1}}
\BIC{$\varphi_{260}$}
\end{ptree}
\noindent$\Pi_{75}$:
\begin{ptree}
\AXC{$\forall A\,\forall c\,\forall f\,\forall e_{0}\,(\forall k \in \mathbb{N}\,(\varphi_{27}) \to \varphi_{50})$\lbl{[CC2]}}
\AXC{$\varphi_{134}$\lbl{A2}}
\BIC{$\forall c\,\forall f\,\forall e_{0}\,(\forall k \in \mathbb{N}\,(\varphi_{27}) \to \varphi_{50})$}
\AXC{$\varphi_{130}$\lbl{A2}}
\BIC{$\forall f\,\forall e_{0}\,(\forall k \in \mathbb{N}\,(\varphi_{27}) \to \varphi_{50})$}
\AXC{$\varphi_{126}$\lbl{A2}}
\BIC{$\forall e_{0}\,(\forall k \in \mathbb{N}\,(\varphi_{27}) \to \varphi_{50})$}
\end{ptree}
\noindent$\Pi_{76}$:
\begin{ptree}
\AXC{$\stackrel{\Pi_{75}}{\vdots}$}\noLine
\UIC{$\forall e_{0}\,(\forall k \in \mathbb{N}\,(\varphi_{27}) \to \varphi_{50})$}
\AXC{$\forall e_{0}\,(\forall k \in \mathbb{N}\,(\varphi_{27}) \to \varphi_{50}) \to \forall k \in \mathbb{N}\,(\varphi_{27}) \to \varphi_{50}$\lbl{A2}}
\BIC{$\forall k \in \mathbb{N}\,(\varphi_{27}) \to \varphi_{50}$}
\AXC{$(\forall k \in \mathbb{N}\,(\varphi_{27}) \to \varphi_{50}) \to \varphi_{216}$\lbl{A1}}
\BIC{$\varphi_{216}$}
\end{ptree}
\noindent$\Pi_{77}$:
\begin{ptree}
\AXC{$\varphi_{97} \to \forall k\,(\varphi_{97})$\lbl{A4}}
\AXC{$\forall k\,(\varphi_{97} \to k \in \mathbb{N} \to \varphi_{27}) \to \forall k\,(\varphi_{97}) \to \forall k \in \mathbb{N}\,(\varphi_{27})$\lbl{A3}}
\AXC{$\stackrel{\Pi_{74}}{\vdots}$}\noLine
\UIC{$\varphi_{260}$}
\BIC{$\varphi_{235}$}
\BIC{$\mathrm{RFn}(f) \to A \subseteq \mathrm{dom}(f) \to c \in A \to \varphi_{120}$}
\AXC{$\stackrel{\Pi_{76}}{\vdots}$}\noLine
\UIC{$\varphi_{216}$}
\BIC{$\mathrm{RFn}(f) \to A \subseteq \mathrm{dom}(f) \to \varphi_{124}$}
\end{ptree}
\noindent$\Pi_{78}$:
\begin{ptree}
\AXC{$\stackrel{\Pi_{49}}{\vdots}$}\noLine
\UIC{$\forall a\,(\varphi_{97} \land \neg 0 < 0 \to \neg (\mathrm{Seq}(a) \land \varphi_{45}))$}
\AXC{$\stackrel{\Pi_{77}}{\vdots}$}\noLine
\UIC{$\mathrm{RFn}(f) \to A \subseteq \mathrm{dom}(f) \to \varphi_{124}$}
\AXC{$(\mathrm{RFn}(f) \to A \subseteq \mathrm{dom}(f) \to \varphi_{124}) \to \varphi_{267}$\lbl{A1}}
\BIC{$\varphi_{267}$}
\BIC{$\varphi_{221} \to (\varphi_{97} \land \neg 0 < 0 \to \forall a\,(\varphi_{97} \land \neg 0 < 0)) \to \varphi_{242}$}
\end{ptree}
\noindent$\Pi_{79}$:
\begin{ptree}
\AXC{$\varphi_{97} \land \neg 0 < 0 \to \forall a\,(\varphi_{97} \land \neg 0 < 0)$\lbl{A4}}
\AXC{$\varphi_{221}$\lbl{A3}}
\AXC{$\stackrel{\Pi_{78}}{\vdots}$}\noLine
\UIC{$\varphi_{221} \to (\varphi_{97} \land \neg 0 < 0 \to \forall a\,(\varphi_{97} \land \neg 0 < 0)) \to \varphi_{242}$}
\BIC{$(\varphi_{97} \land \neg 0 < 0 \to \forall a\,(\varphi_{97} \land \neg 0 < 0)) \to \varphi_{242}$}
\BIC{$\varphi_{242}$}
\end{ptree}
\noindent$\Pi_{80}$:
\begin{ptree}
\AXC{$\varphi_{73}$\lbl{A4}}
\AXC{$\varphi_{148}$\lbl{A3}}
\AXC{$\forall a\,(\neg (\mathrm{Seq}(a) \land \varphi_{45})) \to \varphi_{49}$\lbl{A2}}
\UIC{$\forall q\,(\forall a\,(\neg (\mathrm{Seq}(a) \land \varphi_{45})) \to \varphi_{49})$}
\AXC{$\stackrel{\Pi_{79}}{\vdots}$}\noLine
\UIC{$\varphi_{242}$}
\BIC{$\varphi_{148} \to \varphi_{73} \to \mathrm{RFn}(f) \to A \subseteq \mathrm{dom}(f) \to c \in A \to \varphi_{87}$}
\BIC{$\varphi_{73} \to \mathrm{RFn}(f) \to A \subseteq \mathrm{dom}(f) \to c \in A \to \varphi_{87}$}
\BIC{$\mathrm{RFn}(f) \to A \subseteq \mathrm{dom}(f) \to c \in A \to \varphi_{87}$}
\end{ptree}
\noindent$\Pi_{81}$:
\begin{ptree}
\AXC{$\stackrel{\Pi_{80}}{\vdots}$}\noLine
\UIC{$\mathrm{RFn}(f) \to A \subseteq \mathrm{dom}(f) \to c \in A \to \varphi_{87}$}
\AXC{$(\mathrm{RFn}(f) \to A \subseteq \mathrm{dom}(f) \to c \in A \to \varphi_{87}) \to \varphi_{78} \to \neg \varphi_{48}$\lbl{A1}}
\BIC{$\varphi_{78} \to \neg \varphi_{48}$}
\UIC{$\forall e_{0}\,(\varphi_{78} \to \neg \varphi_{48})$}
\end{ptree}
\noindent$\Pi_{82}$:
\begin{ptree}
\AXC{$\varphi_{75} \to \varphi_{80}$\lbl{A3}}
\AXC{$\varphi_{74}$\lbl{A2}}
\AXC{$\varphi_{150}$\lbl{A1}}
\BIC{$\forall \varepsilon\,(\neg \neg (\varepsilon \in \mathbb{R} \to 0 < \varepsilon \to \varphi_{35})) \to \varepsilon \in \mathbb{R} \to 0 < \varepsilon \to \varphi_{35}$}
\UIC{$\varphi_{75}$}
\AXC{$\mathrm{RFn}(f) \to A \subseteq \mathrm{dom}(f) \to c \in A \to \varphi_{86}$\lbl{A1}}
\AXC{$\varphi_{255}$\lbl{A1}}
\BIC{$\varphi_{75} \to (\varphi_{75} \to \varphi_{80}) \to \varphi_{82} \to \varphi_{102}$}
\BIC{$(\varphi_{75} \to \varphi_{80}) \to \varphi_{82} \to \varphi_{102}$}
\BIC{$\varphi_{82} \to \varphi_{102}$}
\end{ptree}
\noindent$\Pi_{83}$:
\begin{ptree}
\AXC{$\forall e_{0}\,(\varphi_{78} \to \neg \varphi_{48}) \to \forall e_{0}\,(\varphi_{78}) \to \forall e_{0}\,(\neg \varphi_{48})$\lbl{A3}}
\AXC{$\stackrel{\Pi_{81}}{\vdots}$}\noLine
\UIC{$\forall e_{0}\,(\varphi_{78} \to \neg \varphi_{48})$}
\AXC{$\varphi_{82}$\lbl{A4}}
\AXC{$\stackrel{\Pi_{82}}{\vdots}$}\noLine
\UIC{$\varphi_{82} \to \varphi_{102}$}
\BIC{$\varphi_{102}$}
\AXC{$\varphi_{102} \to \varphi_{263}$\lbl{A1}}
\BIC{$\varphi_{263}$}
\BIC{$(\forall e_{0}\,(\varphi_{78} \to \neg \varphi_{48}) \to \forall e_{0}\,(\varphi_{78}) \to \forall e_{0}\,(\neg \varphi_{48})) \to \varphi_{256}$}
\BIC{$\varphi_{256}$}
\end{ptree}
\noindent$\Pi_{84}$:
\begin{ptree}
\AXC{$\varphi_{154}$\lbl{A3}}
\AXC{$\forall e_{0}\,(\neg \varphi_{48}) \to \neg \neg (\varepsilon \in \mathbb{R} \to 0 < \varepsilon \to \varphi_{35})$\lbl{A2}}
\UIC{$\forall \varepsilon\,(\forall e_{0}\,(\neg \varphi_{48}) \to \neg \neg (\varepsilon \in \mathbb{R} \to 0 < \varepsilon \to \varphi_{35}))$}
\AXC{$\varphi_{78} \to \forall e_{0}\,(\varphi_{78})$\lbl{A4}}
\AXC{$\stackrel{\Pi_{83}}{\vdots}$}\noLine
\UIC{$\varphi_{256}$}
\BIC{$\forall \varepsilon\,(\forall e_{0}\,(\neg \varphi_{48}) \to \neg \neg (\varepsilon \in \mathbb{R} \to 0 < \varepsilon \to \varphi_{35})) \to \varphi_{154} \to \varphi_{155}$}
\BIC{$\varphi_{154} \to \varphi_{155}$}
\BIC{$\varphi_{155}$}
\end{ptree}
\noindent$\Pi_{85}$:
\begin{ptree}
\AXC{$\forall e_{0}\,(\neg \varphi_{48}) \to \forall \varepsilon\,\forall e_{0}\,(\neg \varphi_{48})$\lbl{A4}}
\AXC{$\stackrel{\Pi_{84}}{\vdots}$}\noLine
\UIC{$\varphi_{155}$}
\BIC{$\mathrm{RFn}(f) \to A \subseteq \mathrm{dom}(f) \to \varphi_{71}$}
\AXC{$(\mathrm{RFn}(f) \to A \subseteq \mathrm{dom}(f) \to \varphi_{71}) \to \varphi_{151}$\lbl{A1}}
\BIC{$\varphi_{151}$}
\end{ptree}
\noindent$\Pi_{86}$:
\begin{ptree}
\AXC{$\forall f\,\forall A\,\forall c\,(\varphi_{62})$\lbl{DefContIn}}
\AXC{$\forall f\,\forall A\,\forall c\,(\varphi_{62}) \to \forall A\,\forall c\,(\varphi_{62})$\lbl{A2}}
\BIC{$\forall A\,\forall c\,(\varphi_{62})$}
\AXC{$\forall A\,\forall c\,(\varphi_{62}) \to \forall c\,(\varphi_{62})$\lbl{A2}}
\BIC{$\forall c\,(\varphi_{62})$}
\AXC{$\forall c\,(\varphi_{62}) \to \varphi_{62}$\lbl{A2}}
\BIC{$\varphi_{62}$}
\end{ptree}
\noindent$\Pi_{87}$:
\begin{ptree}
\AXC{$\mathrm{RFn}(f) \to A \subseteq \mathrm{dom}(f) \to A \subseteq \mathrm{dom}(f)$\lbl{A1}}
\AXC{$\mathrm{RFn}(f) \to \mathrm{RFn}(f)$\lbl{A1}}
\AXC{$\stackrel{\Pi_{86}}{\vdots}$}\noLine
\UIC{$\varphi_{62}$}
\AXC{$\varphi_{62} \to \varphi_{140}$\lbl{A1}}
\BIC{$\varphi_{140}$}
\BIC{$(\mathrm{RFn}(f) \to A \subseteq \mathrm{dom}(f) \to A \subseteq \mathrm{dom}(f)) \to \varphi_{115}$}
\BIC{$\varphi_{115}$}
\end{ptree}
\noindent 主木:
\begin{ptree}
\AXC{$\stackrel{\Pi_{1}}{\vdots}$}\noLine
\UIC{$\mathrm{RFn}(f) \to A \subseteq \mathrm{dom}(f) \to c \in A \to \varphi_{70}$}
\AXC{$\stackrel{\Pi_{85}}{\vdots}$}\noLine
\UIC{$\varphi_{151}$}
\BIC{$\mathrm{RFn}(f) \to \varphi_{72}$}
\AXC{$\mathrm{RFn}(f) \to A \subseteq \mathrm{dom}(f) \to c \in A \to c \in A$\lbl{A1}}
\AXC{$\stackrel{\Pi_{87}}{\vdots}$}\noLine
\UIC{$\varphi_{115}$}
\BIC{$(\mathrm{RFn}(f) \to \varphi_{72}) \to \mathrm{RFn}(f) \to \varphi_{53}$}
\BIC{$\mathrm{RFn}(f) \to \varphi_{53}$}
\UIC{$\forall c\,(\mathrm{RFn}(f) \to \varphi_{53})$}
\UIC{$\forall A\,\forall c\,(\mathrm{RFn}(f) \to \varphi_{53})$}
\UIC{$\forall f\,\forall A\,\forall c\,(\mathrm{RFn}(f) \to \varphi_{53})$}
\end{ptree}
\begin{abbrevlist}
\item $\varphi_{1} :\equiv \forall w\,(z = w \to \lvert z - f(c) \rvert < e_{0} \to \lvert w - f(c) \rvert < e_{0})$
\item $\varphi_{2} :\equiv \forall N\,(\neg (N \in \mathbb{N} \land \forall n \in \mathbb{N}\,(N \leq n \to \lvert a_{n} - c \rvert < \varepsilon)))$
\item $\varphi_{3} :\equiv \forall n \in \mathbb{N}\,(a_{n} \in A) \to \mathrm{Conv}(a,c) \to \mathrm{Conv}(f \circ a,f(c))$
\item $\varphi_{4} :\equiv \neg (x \in A \to \lvert x - c \rvert < \underline{k}^{-1} \to \lvert f(x) - f(c) \rvert < e_{0})$
\item $\varphi_{5} :\equiv \forall X\,\forall Y\,\forall x\,(X \subseteq Y \to x \in X \to x \in Y) \to \forall Y\,\forall x\,(A \subseteq Y \to x \in A \to x \in Y)$
\item $\varphi_{6} :\equiv \forall x\,(\mathrm{RFn}(f) \to x \in \mathrm{dom}(f) \to x \in \mathbb{R}) \to \mathrm{RFn}(f) \to c \in \mathrm{dom}(f) \to c \in \mathbb{R}$
\item $\varphi_{7} :\equiv \forall x\,(\mathrm{RFn}(f) \to x \in \mathrm{dom}(f) \to x \in \mathbb{R}) \to \mathrm{RFn}(f) \to x \in \mathrm{dom}(f) \to x \in \mathbb{R}$
\item $\varphi_{8} :\equiv \forall x \in A\,(\lvert x - c \rvert < \underline{k}^{-1} \to \lvert f(x) - f(c) \rvert < e_{0})$
\item $\varphi_{9} :\equiv N \in \mathbb{N} \land \forall n \in \mathbb{N}\,(N \leq n \to \lvert f \circ a(n) - f(c) \rvert < e_{0})$
\item $\varphi_{10} :\equiv \forall Y\,\forall x\,(A \subseteq Y \to x \in A \to x \in Y) \to \forall x\,(A \subseteq \mathrm{dom}(f) \to x \in A \to x \in \mathrm{dom}(f))$
\item $\varphi_{11} :\equiv \exists N \in \mathbb{N}\,\forall n \in \mathbb{N}\,(N \leq n \to \lvert f \circ a(n) - f(c) \rvert < e_{0})$
\item $\varphi_{12} :\equiv \forall z\,(\lvert a_{n} - c \rvert < \underline{n}^{-1} \to \underline{n}^{-1} \leq z \to \lvert a_{n} - c \rvert < z)$
\item $\varphi_{13} :\equiv \forall f\,\forall x\,(\mathrm{RFn}(f) \to x \in \mathrm{dom}(f) \to x \in \mathbb{R}) \to \forall x\,(\mathrm{RFn}(f) \to x \in \mathrm{dom}(f) \to x \in \mathbb{R})$
\item $\varphi_{14} :\equiv K \in \mathbb{N} \land \underline{K}^{-1} < \varepsilon \to \forall n \in \mathbb{N}\,(K \leq n \to \lvert a_{n} - c \rvert < \varepsilon)$
\item $\varphi_{15} :\equiv \varepsilon \in \mathbb{R} \to 0 < \varepsilon \to K \in \mathbb{N} \land \underline{K}^{-1} < \varepsilon \to n \in \mathbb{N} \to K \leq n \to K \leq n$
\item $\varphi_{16} :\equiv \forall x\,(A \subseteq \mathrm{dom}(f) \to x \in A \to x \in \mathrm{dom}(f)) \to A \subseteq \mathrm{dom}(f) \to c \in A \to c \in \mathrm{dom}(f)$
\item $\varphi_{17} :\equiv \forall x\,(A \subseteq \mathrm{dom}(f) \to x \in A \to x \in \mathrm{dom}(f)) \to A \subseteq \mathrm{dom}(f) \to x \in A \to x \in \mathrm{dom}(f)$
\item $\varphi_{18} :\equiv 0 < \varepsilon \to K \in \mathbb{N} \land \underline{K}^{-1} < \varepsilon \to n \in \mathbb{N} \to K \leq n \to \underline{K}^{-1} < \varepsilon$
\item $\varphi_{19} :\equiv \forall z\,(\lvert a_{n} - c \rvert < \underline{K}^{-1} \to \underline{K}^{-1} < z \to \lvert a_{n} - c \rvert < z)$
\item $\varphi_{20} :\equiv \lvert f \circ a(N) - f(c) \rvert < e_{0} \to \lvert f(a_{N}) - f(c) \rvert < e_{0}$
\item $\varphi_{21} :\equiv x \in \mathbb{R} \land (x \in A \land \lvert x - c \rvert < \underline{k}^{-1} \land \neg \lvert f(x) - f(c) \rvert < e_{0})$
\item $\varphi_{22} :\equiv K \in \mathbb{N} \land \underline{K}^{-1} < \varepsilon \to n \in \mathbb{N} \to K \leq n \to \lvert a_{n} - c \rvert < \underline{n}^{-1}$
\item $\varphi_{23} :\equiv K \in \mathbb{N} \land \underline{K}^{-1} < \varepsilon \to n \in \mathbb{N} \to K \leq n \to \lvert a_{n} - c \rvert < \underline{K}^{-1}$
\item $\varphi_{24} :\equiv \varepsilon \in \mathbb{R} \to 0 < \varepsilon \to \exists N \in \mathbb{N}\,\forall n \in \mathbb{N}\,(N \leq n \to \lvert a_{n} - c \rvert < \varepsilon)$
\item $\varphi_{25} :\equiv 0 < \varepsilon \to K \in \mathbb{N} \land \underline{K}^{-1} < \varepsilon \to n \in \mathbb{N} \to K \leq n \to \underline{n}^{-1} \leq \underline{K}^{-1}$
\item $\varphi_{26} :\equiv 0 < \varepsilon \to K \in \mathbb{N} \land \underline{K}^{-1} < \varepsilon \to n \in \mathbb{N} \to K \leq n \to \lvert a_{n} - c \rvert < \varepsilon$
\item $\varphi_{27} :\equiv \exists x \in \mathbb{R}\,(x \in A \land \lvert x - c \rvert < \underline{k}^{-1} \land \neg \lvert f(x) - f(c) \rvert < e_{0})$
\item $\varphi_{28} :\equiv \forall x \in \mathbb{R}\,(0 < x \to 0 < x^{-1}) \to \underline{k} \in \mathbb{R} \to 0 < \underline{k} \to 0 < \underline{k}^{-1}$
\item $\varphi_{29} :\equiv \varepsilon \in \mathbb{R} \to 0 < \varepsilon \to K \in \mathbb{N} \land \underline{K}^{-1} < \varepsilon \to K \in \mathbb{N} \land \underline{K}^{-1} < \varepsilon$
\item $\varphi_{30} :\equiv \forall \varepsilon \in \mathbb{R}\,(0 < \varepsilon \to \exists N \in \mathbb{N}\,\forall n \in \mathbb{N}\,(N \leq n \to \lvert a_{n} - c \rvert < \varepsilon))$
\item $\varphi_{31} :\equiv K \in \mathbb{N} \land \underline{K}^{-1} < \varepsilon \to K \in \mathbb{N} \land \forall n \in \mathbb{N}\,(K \leq n \to \lvert a_{n} - c \rvert < \varepsilon)$
\item $\varphi_{32} :\equiv a_{N} \in A \land \lvert a_{N} - c \rvert < \underline{N}^{-1} \land \neg \lvert f(a_{N}) - f(c) \rvert < e_{0}$
\item $\varphi_{33} :\equiv a_{n} \in A \land \lvert a_{n} - c \rvert < \underline{n}^{-1} \land \neg \lvert f(a_{n}) - f(c) \rvert < e_{0}$
\item $\varphi_{34} :\equiv \forall \delta\,(\neg (\delta \in \mathbb{R} \land 0 < \delta \land \forall x \in A\,(\lvert x - c \rvert < \delta \to \lvert f(x) - f(c) \rvert < e_{0})))$
\item $\varphi_{35} :\equiv \exists \delta\,(\delta \in \mathbb{R} \land 0 < \delta \land \forall x \in A\,(\lvert x - c \rvert < \delta \to \lvert f(x) - f(c) \rvert < \varepsilon))$
\item $\varphi_{36} :\equiv \forall x\,\forall y\,\forall z\,(x < y \to y \leq z \to x < z) \to \forall y\,\forall z\,(\lvert a_{n} - c \rvert < y \to y \leq z \to \lvert a_{n} - c \rvert < z)$
\item $\varphi_{37} :\equiv \forall x \in \mathbb{R}\,(x \neq 0 \to x^{-1} \in \mathbb{R}) \to \underline{k} \in \mathbb{R} \to \underline{k} \neq 0 \to \underline{k}^{-1} \in \mathbb{R}$
\item $\varphi_{38} :\equiv K \in \mathbb{N} \land \underline{K}^{-1} < \varepsilon \to \exists N \in \mathbb{N}\,\forall n \in \mathbb{N}\,(N \leq n \to \lvert a_{n} - c \rvert < \varepsilon)$
\item $\varphi_{39} :\equiv \forall w\,(f \circ a(N) = w \to \lvert f \circ a(N) - f(c) \rvert < e_{0} \to \lvert w - f(c) \rvert < e_{0})$
\item $\varphi_{40} :\equiv \forall x\,(A \subseteq \mathrm{dom}(f) \to x \in A \to x \in \mathrm{dom}(f)) \to A \subseteq \mathrm{dom}(f) \to a_{n} \in A \to a_{n} \in \mathrm{dom}(f)$
\item $\varphi_{41} :\equiv \lvert a_{n} - c \rvert < \underline{n}^{-1} \to \underline{n}^{-1} \leq \underline{K}^{-1} \to \lvert a_{n} - c \rvert < \underline{K}^{-1}$
\item $\varphi_{42} :\equiv \forall n \in \mathbb{N}\,(f \circ a(n) = f(a_{n})) \to N \in \mathbb{N} \to f \circ a(N) = f(a_{N})$
\item $\varphi_{43} :\equiv \forall x\,\forall y\,\forall z\,(x < y \to y < z \to x < z) \to \forall y\,\forall z\,(\lvert a_{n} - c \rvert < y \to y < z \to \lvert a_{n} - c \rvert < z)$
\item $\varphi_{44} :\equiv \forall n \in \mathbb{N}\,(a_{n} \in \mathrm{dom}(f)) \to \mathrm{Seq}(f \circ a) \land \forall n \in \mathbb{N}\,(f \circ a(n) = f(a_{n}))$
\item $\varphi_{45} :\equiv \forall k \in \mathbb{N}\,(a_{k} \in A \land \lvert a_{k} - c \rvert < \underline{k}^{-1} \land \neg \lvert f(a_{k}) - f(c) \rvert < e_{0})$
\item $\varphi_{46} :\equiv \forall a\,\forall \alpha\,\forall \varepsilon\,(\mathrm{Conv}(a,\alpha) \to \varepsilon \in \mathbb{R} \to 0 < \varepsilon \to \exists N \in \mathbb{N}\,\forall n \in \mathbb{N}\,(N \leq n \to \lvert a_{n} - \alpha \rvert < \varepsilon))$
\item $\varphi_{47} :\equiv \forall n\,(K \in \mathbb{N} \to n \in \mathbb{N} \to K \leq n \to \underline{n}^{-1} \leq \underline{K}^{-1}) \to K \in \mathbb{N} \to n \in \mathbb{N} \to K \leq n \to \underline{n}^{-1} \leq \underline{K}^{-1}$
\item $\varphi_{48} :\equiv \neg (e_{0} \in \mathbb{R} \to 0 < e_{0} \to \exists \delta\,(\delta \in \mathbb{R} \land 0 < \delta \land \forall x \in A\,(\lvert x - c \rvert < \delta \to \lvert f(x) - f(c) \rvert < e_{0})))$
\item $\varphi_{49} :\equiv \neg (\mathrm{Seq}(q) \land \forall k \in \mathbb{N}\,(q(k) \in A \land \lvert q(k) - c \rvert < \underline{k}^{-1} \land \neg \lvert f(q(k)) - f(c) \rvert < e_{0}))$
\item $\varphi_{50} :\equiv \exists q\,(\mathrm{Seq}(q) \land \forall k \in \mathbb{N}\,(q(k) \in A \land \lvert q(k) - c \rvert < \underline{k}^{-1} \land \neg \lvert f(q(k)) - f(c) \rvert < e_{0}))$
\item $\varphi_{51} :\equiv \forall K\,\forall n\,(K \in \mathbb{N} \to n \in \mathbb{N} \to K \leq n \to \underline{n}^{-1} \leq \underline{K}^{-1}) \to \forall n\,(K \in \mathbb{N} \to n \in \mathbb{N} \to K \leq n \to \underline{n}^{-1} \leq \underline{K}^{-1})$
\item $\varphi_{52} :\equiv A \subseteq \mathrm{dom}(f) \to c \in A \to \forall a\,(\mathrm{Seq}(a) \to \varphi_{3}) \to c \in \mathrm{dom}(f)$
\item $\varphi_{53} :\equiv A \subseteq \mathrm{dom}(f) \to c \in A \to \forall a\,(\mathrm{Seq}(a) \to \varphi_{3}) \to \mathrm{ContIn}(f,A,c)$
\item $\varphi_{54} :\equiv \forall n \in \mathbb{N}\,(N \leq n \to \lvert f \circ a(n) - f(c) \rvert < e_{0}) \to N \in \mathbb{N} \to N \leq N \to \lvert f \circ a(N) - f(c) \rvert < e_{0}$
\item $\varphi_{55} :\equiv \forall \varepsilon \in \mathbb{R}\,(0 < \varepsilon \to \exists K \in \mathbb{N}\,(\underline{K}^{-1} < \varepsilon)) \to \varepsilon \in \mathbb{R} \to 0 < \varepsilon \to \exists K \in \mathbb{N}\,(\underline{K}^{-1} < \varepsilon)$
\item $\varphi_{56} :\equiv \mathrm{RFn}(f) \to A \subseteq \mathrm{dom}(f) \to c \in A \to \forall a\,(\mathrm{Seq}(a) \to \varphi_{3}) \to c \in \mathbb{R}$
\item $\varphi_{57} :\equiv \forall \alpha\,\forall \varepsilon\,(\mathrm{Conv}(f \circ a,\alpha) \to \varepsilon \in \mathbb{R} \to 0 < \varepsilon \to \exists N \in \mathbb{N}\,\forall n \in \mathbb{N}\,(N \leq n \to \lvert f \circ a(n) - \alpha \rvert < \varepsilon))$
\item $\varphi_{58} :\equiv \forall \alpha\,(\mathrm{Conv}(a,\alpha) \leftrightarrow \mathrm{Seq}(a) \land \alpha \in \mathbb{R} \land \forall \varepsilon \in \mathbb{R}\,(0 < \varepsilon \to \exists N \in \mathbb{N}\,\forall n \in \mathbb{N}\,(N \leq n \to \lvert a_{n} - \alpha \rvert < \varepsilon)))$
\item $\varphi_{59} :\equiv \forall \varepsilon\,(\mathrm{Conv}(f \circ a,f(c)) \to \varepsilon \in \mathbb{R} \to 0 < \varepsilon \to \exists N \in \mathbb{N}\,\forall n \in \mathbb{N}\,(N \leq n \to \lvert f \circ a(n) - f(c) \rvert < \varepsilon))$
\item $\varphi_{60} :\equiv \forall x\,(\neg \varphi_{4}) \to x \in A \to \lvert x - c \rvert < \underline{k}^{-1} \to \lvert f(x) - f(c) \rvert < e_{0}$
\item $\varphi_{61} :\equiv \varepsilon \in \mathbb{R} \to 0 < \varepsilon \to K \in \mathbb{N} \land \underline{K}^{-1} < \varepsilon \to n \in \mathbb{N} \to K \leq n \to \varphi_{33}$
\item $\varphi_{62} :\equiv \mathrm{ContIn}(f,A,c) \leftrightarrow \mathrm{RFn}(f) \land A \subseteq \mathrm{dom}(f) \land c \in A \land \forall \varepsilon \in \mathbb{R}\,(0 < \varepsilon \to \varphi_{35})$
\item $\varphi_{63} :\equiv (\mathrm{RFn}(f) \to A \subseteq \mathrm{dom}(f) \to c \in A \to c \in A) \to \mathrm{RFn}(f) \to \varphi_{52}$
\item $\varphi_{64} :\equiv \mathrm{Seq}(a) \land \varphi_{45} \to \varepsilon \in \mathbb{R} \to 0 < \varepsilon \to \exists K \in \mathbb{N}\,(\underline{K}^{-1} < \varepsilon)$
\item $\varphi_{65} :\equiv \mathrm{Seq}(a) \land \varphi_{45} \to \varepsilon \in \mathbb{R} \to 0 < \varepsilon \to K \in \mathbb{N} \land \underline{K}^{-1} < \varepsilon \to K \in \mathbb{N}$
\item $\varphi_{66} :\equiv \mathrm{Seq}(a) \land \varphi_{45} \to \mathrm{Seq}(f \circ a) \land \forall n \in \mathbb{N}\,(f \circ a(n) = f(a_{n}))$
\item $\varphi_{67} :\equiv A \subseteq \mathrm{dom}(f) \to c \in A \to \forall a\,(\mathrm{Seq}(a) \to \varphi_{3}) \to \mathrm{Seq}(a) \to \varphi_{3}$
\item $\varphi_{68} :\equiv A \subseteq \mathrm{dom}(f) \to c \in A \to \forall a\,(\mathrm{Seq}(a) \to \varphi_{3}) \to \forall a\,(\mathrm{Seq}(a) \to \varphi_{3})$
\item $\varphi_{69} :\equiv \mathrm{Seq}(a) \land \varphi_{45} \to \varepsilon \in \mathbb{R} \to 0 < \varepsilon \to K \in \mathbb{N} \land \underline{K}^{-1} < \varepsilon \to n \in \mathbb{N} \to n \in \mathbb{N}$
\item $\varphi_{70} :\equiv \forall a\,(\mathrm{Seq}(a) \to \varphi_{3}) \to \neg \forall \varepsilon \in \mathbb{R}\,(0 < \varepsilon \to \varphi_{35}) \to \neg 0 < 0$
\item $\varphi_{71} :\equiv c \in A \to \forall a\,(\mathrm{Seq}(a) \to \varphi_{3}) \to \neg \forall \varepsilon \in \mathbb{R}\,(0 < \varepsilon \to \varphi_{35}) \to 0 < 0$
\item $\varphi_{72} :\equiv A \subseteq \mathrm{dom}(f) \to c \in A \to \forall a\,(\mathrm{Seq}(a) \to \varphi_{3}) \to \forall \varepsilon \in \mathbb{R}\,(0 < \varepsilon \to \varphi_{35})$
\item $\varphi_{73} :\equiv \forall a\,(\neg (\mathrm{Seq}(a) \land \varphi_{45})) \to \forall q\,\forall a\,(\neg (\mathrm{Seq}(a) \land \varphi_{45}))$
\item $\varphi_{74} :\equiv \forall \varepsilon\,(\neg \neg (\varepsilon \in \mathbb{R} \to 0 < \varepsilon \to \varphi_{35})) \to \neg \neg (\varepsilon \in \mathbb{R} \to 0 < \varepsilon \to \varphi_{35})$
\item $\varphi_{75} :\equiv \forall \varepsilon\,(\forall \varepsilon\,(\neg \neg (\varepsilon \in \mathbb{R} \to 0 < \varepsilon \to \varphi_{35})) \to \varepsilon \in \mathbb{R} \to 0 < \varepsilon \to \varphi_{35})$
\item $\varphi_{76} :\equiv \neg \forall \varepsilon \in \mathbb{R}\,(0 < \varepsilon \to \varphi_{35}) \to \exists \varepsilon\,(\neg (\varepsilon \in \mathbb{R} \to 0 < \varepsilon \to \varphi_{35}))$
\item $\varphi_{77} :\equiv \forall a\,(\mathrm{RFn}(f) \to \mathrm{Seq}(a) \to \varphi_{44}) \to \mathrm{RFn}(f) \to \mathrm{Seq}(a) \to \varphi_{44}$
\item $\varphi_{78} :\equiv \mathrm{RFn}(f) \land A \subseteq \mathrm{dom}(f) \land c \in A \land \forall a\,(\mathrm{Seq}(a) \to \varphi_{3}) \land \neg \forall \varepsilon \in \mathbb{R}\,(0 < \varepsilon \to \varphi_{35}) \land \neg 0 < 0$
\item $\varphi_{79} :\equiv \mathrm{Seq}(a) \land \varphi_{45} \to \varphi_{9} \to N \in \mathbb{N} \to N \leq N \to \lvert f \circ a(N) - f(c) \rvert < e_{0}$
\item $\varphi_{80} :\equiv \forall \varepsilon\,\forall \varepsilon\,(\neg \neg (\varepsilon \in \mathbb{R} \to 0 < \varepsilon \to \varphi_{35})) \to \forall \varepsilon \in \mathbb{R}\,(0 < \varepsilon \to \varphi_{35})$
\item $\varphi_{81} :\equiv (A \subseteq \mathrm{dom}(f) \to c \in A \to c \in \mathrm{dom}(f)) \to (\mathrm{RFn}(f) \to A \subseteq \mathrm{dom}(f) \to A \subseteq \mathrm{dom}(f)) \to \varphi_{63}$
\item $\varphi_{82} :\equiv \forall \varepsilon\,(\neg \neg (\varepsilon \in \mathbb{R} \to 0 < \varepsilon \to \varphi_{35})) \to \forall \varepsilon\,\forall \varepsilon\,(\neg \neg (\varepsilon \in \mathbb{R} \to 0 < \varepsilon \to \varphi_{35}))$
\item $\varphi_{83} :\equiv \mathrm{Seq}(a) \land \varphi_{45} \to \varphi_{9} \to \forall n \in \mathbb{N}\,(N \leq n \to \lvert f \circ a(n) - f(c) \rvert < e_{0})$
\item $\varphi_{84} :\equiv \forall f\,\forall a\,(\mathrm{RFn}(f) \to \mathrm{Seq}(a) \to \varphi_{44}) \to \forall a\,(\mathrm{RFn}(f) \to \mathrm{Seq}(a) \to \varphi_{44})$
\item $\varphi_{85} :\equiv (\mathrm{RFn}(f) \to c \in \mathrm{dom}(f) \to c \in \mathbb{R}) \to (\mathrm{RFn}(f) \to \mathrm{RFn}(f)) \to (\mathrm{RFn}(f) \to \varphi_{52}) \to \varphi_{56}$
\item $\varphi_{86} :\equiv \forall a\,(\mathrm{Seq}(a) \to \varphi_{3}) \to \neg \forall \varepsilon \in \mathbb{R}\,(0 < \varepsilon \to \varphi_{35}) \to \neg \forall \varepsilon \in \mathbb{R}\,(0 < \varepsilon \to \varphi_{35})$
\item $\varphi_{87} :\equiv \forall a\,(\mathrm{Seq}(a) \to \varphi_{3}) \to \neg \forall \varepsilon \in \mathbb{R}\,(0 < \varepsilon \to \varphi_{35}) \to \varphi_{48} \to 0 < 0$
\item $\varphi_{88} :\equiv \forall a\,(\mathrm{Seq}(a) \to \varphi_{3}) \to \neg \forall \varepsilon \in \mathbb{R}\,(0 < \varepsilon \to \varphi_{35}) \to \varphi_{48} \to 0 < e_{0}$
\item $\varphi_{89} :\equiv \forall a\,(\mathrm{Seq}(a) \to \varphi_{3}) \to \neg \forall \varepsilon \in \mathbb{R}\,(0 < \varepsilon \to \varphi_{35}) \to \varphi_{48} \to e_{0} \in \mathbb{R}$
\item $\varphi_{90} :\equiv \neg \forall \varepsilon \in \mathbb{R}\,(0 < \varepsilon \to \varphi_{35}) \to \varphi_{48} \to k \in \mathbb{N} \to \varphi_{4} \to x \in \mathrm{dom}(f)$
\item $\varphi_{91} :\equiv \forall a\,(\mathrm{Seq}(a) \to \varphi_{3}) \to \neg \forall \varepsilon \in \mathbb{R}\,(0 < \varepsilon \to \varphi_{35}) \to \varphi_{48} \to k \in \mathbb{N} \to k \in \mathbb{N}$
\item $\varphi_{92} :\equiv \forall a\,(\mathrm{Seq}(a) \to \varphi_{3}) \to \neg \forall \varepsilon \in \mathbb{R}\,(0 < \varepsilon \to \varphi_{35}) \to \varphi_{48} \to k \in \mathbb{N} \to 0 \neq \underline{k}$
\item $\varphi_{93} :\equiv \forall a\,(\mathrm{Seq}(a) \to \varphi_{3}) \to \neg \forall \varepsilon \in \mathbb{R}\,(0 < \varepsilon \to \varphi_{35}) \to \varphi_{48} \to k \in \mathbb{N} \to \underline{k} \neq 0$
\item $\varphi_{94} :\equiv \forall a\,(\mathrm{Seq}(a) \to \varphi_{3}) \to \neg \forall \varepsilon \in \mathbb{R}\,(0 < \varepsilon \to \varphi_{35}) \to \varphi_{48} \to k \in \mathbb{N} \to 0 < \underline{k}$
\item $\varphi_{95} :\equiv \forall a\,(\mathrm{Seq}(a) \to \varphi_{3}) \to \neg \forall \varepsilon \in \mathbb{R}\,(0 < \varepsilon \to \varphi_{35}) \to \varphi_{48} \to k \in \mathbb{N} \to \underline{k} \in \mathbb{R}$
\item $\varphi_{96} :\equiv \forall a\,(\mathrm{Seq}(a) \to \varphi_{3}) \to \neg \forall \varepsilon \in \mathbb{R}\,(0 < \varepsilon \to \varphi_{35}) \to \varphi_{48} \to k \in \mathbb{N} \to 0 < \underline{k}^{-1}$
\item $\varphi_{97} :\equiv \mathrm{RFn}(f) \land A \subseteq \mathrm{dom}(f) \land c \in A \land \forall a\,(\mathrm{Seq}(a) \to \varphi_{3}) \land \neg \forall \varepsilon \in \mathbb{R}\,(0 < \varepsilon \to \varphi_{35}) \land \varphi_{48}$
\item $\varphi_{98} :\equiv \forall a\,(\mathrm{Seq}(a) \to \varphi_{3}) \to \neg \forall \varepsilon \in \mathbb{R}\,(0 < \varepsilon \to \varphi_{35}) \to \varphi_{48} \to k \in \mathbb{N} \to \underline{k}^{-1} \in \mathbb{R}$
\item $\varphi_{99} :\equiv \neg \forall \varepsilon \in \mathbb{R}\,(0 < \varepsilon \to \varphi_{35}) \to \varphi_{48} \to \neg (\underline{k}^{-1} \in \mathbb{R} \land 0 < \underline{k}^{-1} \land \varphi_{8})$
\item $\varphi_{100} :\equiv \varphi_{48} \to \mathrm{Seq}(a) \land \varphi_{45} \to \varphi_{9} \to \forall n \in \mathbb{N}\,(f \circ a(n) = f(a_{n}))$
\item $\varphi_{101} :\equiv (\forall a\,(\mathrm{Seq}(a) \to \varphi_{3}) \to \mathrm{Seq}(a) \to \varphi_{3}) \to \mathrm{RFn}(f) \to \varphi_{67}$
\item $\varphi_{102} :\equiv \mathrm{RFn}(f) \to A \subseteq \mathrm{dom}(f) \to c \in A \to \forall a\,(\mathrm{Seq}(a) \to \varphi_{3}) \to \varphi_{76}$
\item $\varphi_{103} :\equiv \neg \forall \varepsilon \in \mathbb{R}\,(0 < \varepsilon \to \varphi_{35}) \to \varphi_{48} \to \mathrm{Seq}(a) \land \varphi_{45} \to \mathrm{Seq}(a)$
\item $\varphi_{104} :\equiv \neg \forall \varepsilon \in \mathbb{R}\,(0 < \varepsilon \to \varphi_{35}) \to \varphi_{48} \to \mathrm{Seq}(a) \land \varphi_{45} \to \mathrm{Conv}(a,c)$
\item $\varphi_{105} :\equiv \neg \forall \varepsilon \in \mathbb{R}\,(0 < \varepsilon \to \varphi_{35}) \to \varphi_{48} \to \mathrm{Seq}(a) \land \varphi_{45} \to n \in \mathbb{N} \to a_{n} \in A$
\item $\varphi_{106} :\equiv \neg \forall \varepsilon \in \mathbb{R}\,(0 < \varepsilon \to \varphi_{35}) \to \varphi_{48} \to \mathrm{Seq}(a) \land \varphi_{45} \to \varepsilon \in \mathbb{R} \to \varepsilon \in \mathbb{R}$
\item $\varphi_{107} :\equiv \neg \forall \varepsilon \in \mathbb{R}\,(0 < \varepsilon \to \varphi_{35}) \to \varphi_{48} \to \mathrm{Seq}(a) \land \varphi_{45} \to \mathrm{Conv}(f \circ a,f(c))$
\item $\varphi_{108} :\equiv \neg \forall \varepsilon \in \mathbb{R}\,(0 < \varepsilon \to \varphi_{35}) \to \varphi_{48} \to \mathrm{Seq}(a) \land \varphi_{45} \to \forall n \in \mathbb{N}\,(a_{n} \in A)$
\item $\varphi_{109} :\equiv \neg \forall \varepsilon \in \mathbb{R}\,(0 < \varepsilon \to \varphi_{35}) \to \varphi_{48} \to \mathrm{Seq}(a) \land \varphi_{45} \to n \in \mathbb{N} \to a_{n} \in \mathrm{dom}(f)$
\item $\varphi_{110} :\equiv \neg \forall \varepsilon \in \mathbb{R}\,(0 < \varepsilon \to \varphi_{35}) \to \varphi_{48} \to \mathrm{Seq}(a) \land \varphi_{45} \to \forall n \in \mathbb{N}\,(a_{n} \in \mathrm{dom}(f))$
\item $\varphi_{111} :\equiv \neg \forall \varepsilon \in \mathbb{R}\,(0 < \varepsilon \to \varphi_{35}) \to \varphi_{48} \to \mathrm{Seq}(a) \land \varphi_{45} \to \varepsilon \in \mathbb{R} \to 0 < \varepsilon \to 0 < \varepsilon$
\item $\varphi_{112} :\equiv (\underline{k} = 0 \to 0 = \underline{k}) \to \mathrm{RFn}(f) \to A \subseteq \mathrm{dom}(f) \to c \in A \to \varphi_{93}$
\item $\varphi_{113} :\equiv \forall a\,(\mathrm{Seq}(a) \to \varphi_{3}) \to \neg \forall \varepsilon \in \mathbb{R}\,(0 < \varepsilon \to \varphi_{35}) \to \varphi_{48} \to k \in \mathbb{N} \to \neg \varphi_{8}$
\item $\varphi_{114} :\equiv \forall a\,(\mathrm{Seq}(a) \to \varphi_{3}) \to \neg \forall \varepsilon \in \mathbb{R}\,(0 < \varepsilon \to \varphi_{35}) \to \varphi_{48} \to k \in \mathbb{N} \to \exists x\,(\varphi_{4})$
\item $\varphi_{115} :\equiv (\mathrm{RFn}(f) \to A \subseteq \mathrm{dom}(f) \to c \in A \to c \in A) \to (\mathrm{RFn}(f) \to \varphi_{72}) \to \mathrm{RFn}(f) \to \varphi_{53}$
\item $\varphi_{116} :\equiv \forall a\,(\mathrm{Seq}(a) \to \varphi_{3}) \to \neg \forall \varepsilon \in \mathbb{R}\,(0 < \varepsilon \to \varphi_{35}) \to \varphi_{48} \to k \in \mathbb{N} \to \varphi_{4} \to x \in A$
\item $\varphi_{117} :\equiv \forall a\,(\mathrm{Seq}(a) \to \varphi_{3}) \to \neg \forall \varepsilon \in \mathbb{R}\,(0 < \varepsilon \to \varphi_{35}) \to \varphi_{48} \to k \in \mathbb{N} \to \varphi_{4} \to x \in \mathbb{R}$
\item $\varphi_{118} :\equiv \forall a\,(\mathrm{Seq}(a) \to \varphi_{3}) \to \neg \forall \varepsilon \in \mathbb{R}\,(0 < \varepsilon \to \varphi_{35}) \to \varphi_{48} \to k \in \mathbb{N} \to \varphi_{27}$
\item $\varphi_{119} :\equiv c \in A \to \forall a\,(\mathrm{Seq}(a) \to \varphi_{3}) \to \neg \forall \varepsilon \in \mathbb{R}\,(0 < \varepsilon \to \varphi_{35}) \to \varphi_{48} \to \varphi_{34}$
\item $\varphi_{120} :\equiv \forall a\,(\mathrm{Seq}(a) \to \varphi_{3}) \to \neg \forall \varepsilon \in \mathbb{R}\,(0 < \varepsilon \to \varphi_{35}) \to \varphi_{48} \to \forall k \in \mathbb{N}\,(\varphi_{27})$
\item $\varphi_{121} :\equiv \neg \forall \varepsilon \in \mathbb{R}\,(0 < \varepsilon \to \varphi_{35}) \to \varphi_{48} \to \mathrm{Seq}(a) \land \varphi_{45} \to \varphi_{9} \to N \in \mathbb{N}$
\item $\varphi_{122} :\equiv \forall a\,(\mathrm{Seq}(a) \to \varphi_{3}) \to \neg \forall \varepsilon \in \mathbb{R}\,(0 < \varepsilon \to \varphi_{35}) \to \varphi_{48} \to \mathrm{Seq}(a) \land \varphi_{45} \to 0 < 0$
\item $\varphi_{123} :\equiv c \in A \to \forall a\,(\mathrm{Seq}(a) \to \varphi_{3}) \to \neg \forall \varepsilon \in \mathbb{R}\,(0 < \varepsilon \to \varphi_{35}) \to \varphi_{48} \to \varphi_{48}$
\item $\varphi_{124} :\equiv c \in A \to \forall a\,(\mathrm{Seq}(a) \to \varphi_{3}) \to \neg \forall \varepsilon \in \mathbb{R}\,(0 < \varepsilon \to \varphi_{35}) \to \varphi_{48} \to \varphi_{50}$
\item $\varphi_{125} :\equiv \neg \forall \varepsilon \in \mathbb{R}\,(0 < \varepsilon \to \varphi_{35}) \to \varphi_{48} \to \mathrm{Seq}(a) \land \varphi_{45} \to \varepsilon \in \mathbb{R} \to \varphi_{18}$
\item $\varphi_{126} :\equiv \forall f\,\forall e_{0}\,(\forall k \in \mathbb{N}\,(\varphi_{27}) \to \varphi_{50}) \to \forall e_{0}\,(\forall k \in \mathbb{N}\,(\varphi_{27}) \to \varphi_{50})$
\item $\varphi_{127} :\equiv \mathrm{RFn}(f) \to A \subseteq \mathrm{dom}(f) \to c \in A \to \forall a\,(\mathrm{Seq}(a) \to \varphi_{3}) \to \varphi_{90}$
\item $\varphi_{128} :\equiv \neg \forall \varepsilon \in \mathbb{R}\,(0 < \varepsilon \to \varphi_{35}) \to \varphi_{48} \to \mathrm{Seq}(a) \land \varphi_{45} \to \varepsilon \in \mathbb{R} \to \varphi_{25}$
\item $\varphi_{129} :\equiv \neg \forall \varepsilon \in \mathbb{R}\,(0 < \varepsilon \to \varphi_{35}) \to \varphi_{48} \to \mathrm{Seq}(a) \land \varphi_{45} \to \varepsilon \in \mathbb{R} \to \varphi_{26}$
\item $\varphi_{130} :\equiv \forall c\,\forall f\,\forall e_{0}\,(\forall k \in \mathbb{N}\,(\varphi_{27}) \to \varphi_{50}) \to \forall f\,\forall e_{0}\,(\forall k \in \mathbb{N}\,(\varphi_{27}) \to \varphi_{50})$
\item $\varphi_{131} :\equiv \neg \forall \varepsilon \in \mathbb{R}\,(0 < \varepsilon \to \varphi_{35}) \to \varphi_{48} \to \mathrm{Seq}(a) \land \varphi_{45} \to N \in \mathbb{N} \to \varphi_{32}$
\item $\varphi_{132} :\equiv \neg \forall \varepsilon \in \mathbb{R}\,(0 < \varepsilon \to \varphi_{35}) \to \varphi_{48} \to \mathrm{Seq}(a) \land \varphi_{45} \to n \in \mathbb{N} \to \varphi_{33}$
\item $\varphi_{133} :\equiv \neg \forall \varepsilon \in \mathbb{R}\,(0 < \varepsilon \to \varphi_{35}) \to \varphi_{48} \to \mathrm{Seq}(a) \land \varphi_{45} \to \varepsilon \in \mathbb{R} \to 0 < \varepsilon \to \varphi_{14}$
\item $\varphi_{134} :\equiv \forall A\,\forall c\,\forall f\,\forall e_{0}\,(\forall k \in \mathbb{N}\,(\varphi_{27}) \to \varphi_{50}) \to \forall c\,\forall f\,\forall e_{0}\,(\forall k \in \mathbb{N}\,(\varphi_{27}) \to \varphi_{50})$
\item $\varphi_{135} :\equiv \neg \forall \varepsilon \in \mathbb{R}\,(0 < \varepsilon \to \varphi_{35}) \to \varphi_{48} \to \mathrm{Seq}(a) \land \varphi_{45} \to \varepsilon \in \mathbb{R} \to 0 < \varepsilon \to \varphi_{22}$
\item $\varphi_{136} :\equiv \neg \forall \varepsilon \in \mathbb{R}\,(0 < \varepsilon \to \varphi_{35}) \to \varphi_{48} \to \mathrm{Seq}(a) \land \varphi_{45} \to \varepsilon \in \mathbb{R} \to 0 < \varepsilon \to \varphi_{23}$
\item $\varphi_{137} :\equiv \neg \forall \varepsilon \in \mathbb{R}\,(0 < \varepsilon \to \varphi_{35}) \to \varphi_{48} \to \mathrm{Seq}(a) \land \varphi_{45} \to \varphi_{9} \to f \circ a(N) = f(a_{N})$
\item $\varphi_{138} :\equiv \neg \forall \varepsilon \in \mathbb{R}\,(0 < \varepsilon \to \varphi_{35}) \to \varphi_{48} \to \mathrm{Seq}(a) \land \varphi_{45} \to \varphi_{9} \to \lvert f(a_{N}) - f(c) \rvert < e_{0}$
\item $\varphi_{139} :\equiv \neg \forall \varepsilon \in \mathbb{R}\,(0 < \varepsilon \to \varphi_{35}) \to \varphi_{48} \to \mathrm{Seq}(a) \land \varphi_{45} \to \mathrm{Seq}(a) \land \varphi_{45}$
\item $\varphi_{140} :\equiv (\mathrm{RFn}(f) \to \mathrm{RFn}(f)) \to (\mathrm{RFn}(f) \to A \subseteq \mathrm{dom}(f) \to A \subseteq \mathrm{dom}(f)) \to \varphi_{115}$
\item $\varphi_{141} :\equiv \neg \forall \varepsilon \in \mathbb{R}\,(0 < \varepsilon \to \varphi_{35}) \to \varphi_{48} \to \mathrm{Seq}(a) \land \varphi_{45} \to \varepsilon \in \mathbb{R} \to 0 < \varepsilon \to \varphi_{31}$
\item $\varphi_{142} :\equiv \mathrm{RFn}(f) \to A \subseteq \mathrm{dom}(f) \to c \in A \to \forall a\,(\mathrm{Seq}(a) \to \varphi_{3}) \to \varphi_{99}$
\item $\varphi_{143} :\equiv \neg \forall \varepsilon \in \mathbb{R}\,(0 < \varepsilon \to \varphi_{35}) \to \varphi_{48} \to \mathrm{Seq}(a) \land \varphi_{45} \to \varphi_{9} \to \lvert f \circ a(N) - f(c) \rvert < e_{0}$
\item $\varphi_{144} :\equiv \neg \forall \varepsilon \in \mathbb{R}\,(0 < \varepsilon \to \varphi_{35}) \to \varphi_{48} \to \mathrm{Seq}(a) \land \varphi_{45} \to \varepsilon \in \mathbb{R} \to 0 < \varepsilon \to \varphi_{38}$
\item $\varphi_{145} :\equiv \neg \forall \varepsilon \in \mathbb{R}\,(0 < \varepsilon \to \varphi_{35}) \to \varphi_{48} \to \mathrm{Seq}(a) \land \varphi_{45} \to \varphi_{9} \to N \in \mathbb{N} \to f \circ a(N) = f(a_{N})$
\item $\varphi_{146} :\equiv \forall a\,(\mathrm{Seq}(a) \to \varphi_{3}) \to \neg \forall \varepsilon \in \mathbb{R}\,(0 < \varepsilon \to \varphi_{35}) \to \varphi_{48} \to k \in \mathbb{N} \to \varphi_{4} \to \varphi_{4}$
\item $\varphi_{147} :\equiv \mathrm{RFn}(f) \to A \subseteq \mathrm{dom}(f) \to c \in A \to \forall a\,(\mathrm{Seq}(a) \to \varphi_{3}) \to \varphi_{103}$
\item $\varphi_{148} :\equiv \forall q\,(\forall a\,(\neg (\mathrm{Seq}(a) \land \varphi_{45})) \to \varphi_{49}) \to \forall q\,\forall a\,(\neg (\mathrm{Seq}(a) \land \varphi_{45})) \to \forall q\,(\varphi_{49})$
\item $\varphi_{149} :\equiv \mathrm{RFn}(f) \to A \subseteq \mathrm{dom}(f) \to c \in A \to \forall a\,(\mathrm{Seq}(a) \to \varphi_{3}) \to \varphi_{104}$
\item $\varphi_{150} :\equiv \varphi_{74} \to \forall \varepsilon\,(\neg \neg (\varepsilon \in \mathbb{R} \to 0 < \varepsilon \to \varphi_{35})) \to \varepsilon \in \mathbb{R} \to 0 < \varepsilon \to \varphi_{35}$
\item $\varphi_{151} :\equiv (\mathrm{RFn}(f) \to A \subseteq \mathrm{dom}(f) \to c \in A \to \varphi_{70}) \to \mathrm{RFn}(f) \to \varphi_{72}$
\item $\varphi_{152} :\equiv \forall a\,(\mathrm{Seq}(a) \to \varphi_{3}) \to \neg \forall \varepsilon \in \mathbb{R}\,(0 < \varepsilon \to \varphi_{35}) \to \varphi_{48} \to k \in \mathbb{N} \to \varphi_{4} \to \varphi_{21}$
\item $\varphi_{153} :\equiv \forall a\,(\mathrm{Seq}(a) \to \varphi_{3}) \to \neg \forall \varepsilon \in \mathbb{R}\,(0 < \varepsilon \to \varphi_{35}) \to \varphi_{48} \to k \in \mathbb{N} \to \varphi_{4} \to \varphi_{27}$
\item $\varphi_{154} :\equiv \forall \varepsilon\,(\forall e_{0}\,(\neg \varphi_{48}) \to \neg \neg (\varepsilon \in \mathbb{R} \to 0 < \varepsilon \to \varphi_{35})) \to \forall \varepsilon\,\forall e_{0}\,(\neg \varphi_{48}) \to \forall \varepsilon\,(\neg \neg (\varepsilon \in \mathbb{R} \to 0 < \varepsilon \to \varphi_{35}))$
\item $\varphi_{155} :\equiv (\forall e_{0}\,(\neg \varphi_{48}) \to \forall \varepsilon\,\forall e_{0}\,(\neg \varphi_{48})) \to \mathrm{RFn}(f) \to A \subseteq \mathrm{dom}(f) \to \varphi_{71}$
\item $\varphi_{156} :\equiv \mathrm{RFn}(f) \to A \subseteq \mathrm{dom}(f) \to c \in A \to \forall a\,(\mathrm{Seq}(a) \to \varphi_{3}) \to \varphi_{105}$
\item $\varphi_{157} :\equiv \mathrm{RFn}(f) \to A \subseteq \mathrm{dom}(f) \to c \in A \to \forall a\,(\mathrm{Seq}(a) \to \varphi_{3}) \to \varphi_{106}$
\item $\varphi_{158} :\equiv \mathrm{RFn}(f) \to A \subseteq \mathrm{dom}(f) \to c \in A \to \forall a\,(\mathrm{Seq}(a) \to \varphi_{3}) \to \varphi_{107}$
\item $\varphi_{159} :\equiv \mathrm{RFn}(f) \to A \subseteq \mathrm{dom}(f) \to c \in A \to \forall a\,(\mathrm{Seq}(a) \to \varphi_{3}) \to \varphi_{108}$
\item $\varphi_{160} :\equiv \mathrm{RFn}(f) \to A \subseteq \mathrm{dom}(f) \to c \in A \to \forall a\,(\mathrm{Seq}(a) \to \varphi_{3}) \to \varphi_{109}$
\item $\varphi_{161} :\equiv \forall n\,(\varphi_{97} \land (\mathrm{Seq}(a) \land \varphi_{45}) \to n \in \mathbb{N} \to a_{n} \in \mathrm{dom}(f))$
\item $\varphi_{162} :\equiv \mathrm{RFn}(f) \to A \subseteq \mathrm{dom}(f) \to c \in A \to \forall a\,(\mathrm{Seq}(a) \to \varphi_{3}) \to \varphi_{110}$
\item $\varphi_{163} :\equiv \mathrm{RFn}(f) \to A \subseteq \mathrm{dom}(f) \to c \in A \to \forall a\,(\mathrm{Seq}(a) \to \varphi_{3}) \to \varphi_{111}$
\item $\varphi_{164} :\equiv c \in A \to \forall a\,(\mathrm{Seq}(a) \to \varphi_{3}) \to \neg \forall \varepsilon \in \mathbb{R}\,(0 < \varepsilon \to \varphi_{35}) \to \varphi_{48} \to \varphi_{64}$
\item $\varphi_{165} :\equiv \forall n\,(\varphi_{97} \land (\mathrm{Seq}(a) \land \varphi_{45})) \to \forall n \in \mathbb{N}\,(a_{n} \in \mathrm{dom}(f))$
\item $\varphi_{166} :\equiv \forall a\,(\mathrm{Seq}(a) \to \varphi_{3}) \to \neg \forall \varepsilon \in \mathbb{R}\,(0 < \varepsilon \to \varphi_{35}) \to \varphi_{48} \to \mathrm{Seq}(a) \land \varphi_{45} \to \varphi_{11}$
\item $\varphi_{167} :\equiv c \in A \to \forall a\,(\mathrm{Seq}(a) \to \varphi_{3}) \to \neg \forall \varepsilon \in \mathbb{R}\,(0 < \varepsilon \to \varphi_{35}) \to \varphi_{48} \to \varphi_{65}$
\item $\varphi_{168} :\equiv c \in A \to \forall a\,(\mathrm{Seq}(a) \to \varphi_{3}) \to \neg \forall \varepsilon \in \mathbb{R}\,(0 < \varepsilon \to \varphi_{35}) \to \varphi_{48} \to \varphi_{66}$
\item $\varphi_{169} :\equiv \forall a\,(\mathrm{Seq}(a) \to \varphi_{3}) \to \neg \forall \varepsilon \in \mathbb{R}\,(0 < \varepsilon \to \varphi_{35}) \to \varphi_{48} \to \mathrm{Seq}(a) \land \varphi_{45} \to \varphi_{9} \to N \leq N$
\item $\varphi_{170} :\equiv \forall a\,(\mathrm{Seq}(a) \to \varphi_{3}) \to \neg \forall \varepsilon \in \mathbb{R}\,(0 < \varepsilon \to \varphi_{35}) \to \varphi_{48} \to \mathrm{Seq}(a) \land \varphi_{45} \to \varphi_{15}$
\item $\varphi_{171} :\equiv \forall a\,(\mathrm{Seq}(a) \to \varphi_{3}) \to \neg \forall \varepsilon \in \mathbb{R}\,(0 < \varepsilon \to \varphi_{35}) \to \varphi_{48} \to \mathrm{Seq}(a) \land \varphi_{45} \to \varphi_{9} \to 0 < 0$
\item $\varphi_{172} :\equiv \forall a\,(\mathrm{Seq}(a) \to \varphi_{3}) \to \neg \forall \varepsilon \in \mathbb{R}\,(0 < \varepsilon \to \varphi_{35}) \to \varphi_{48} \to \mathrm{Seq}(a) \land \varphi_{45} \to \varphi_{24}$
\item $\varphi_{173} :\equiv \forall a\,(\mathrm{Seq}(a) \to \varphi_{3}) \to \neg \forall \varepsilon \in \mathbb{R}\,(0 < \varepsilon \to \varphi_{35}) \to \varphi_{48} \to \mathrm{Seq}(a) \land \varphi_{45} \to \varphi_{29}$
\item $\varphi_{174} :\equiv c \in A \to \forall a\,(\mathrm{Seq}(a) \to \varphi_{3}) \to \neg \forall \varepsilon \in \mathbb{R}\,(0 < \varepsilon \to \varphi_{35}) \to \varphi_{48} \to \varphi_{69}$
\item $\varphi_{175} :\equiv \varphi_{97} \land (\mathrm{Seq}(a) \land \varphi_{45}) \land \varepsilon \in \mathbb{R} \land 0 < \varepsilon \land (K \in \mathbb{N} \land \underline{K}^{-1} < \varepsilon)$
\item $\varphi_{176} :\equiv \forall a\,(\mathrm{Seq}(a) \to \varphi_{3}) \to \neg \forall \varepsilon \in \mathbb{R}\,(0 < \varepsilon \to \varphi_{35}) \to \varphi_{48} \to \mathrm{Seq}(a) \land \varphi_{45} \to \varphi_{30}$
\item $\varphi_{177} :\equiv \forall a\,(\mathrm{Seq}(a) \to \varphi_{3}) \to \neg \forall \varepsilon \in \mathbb{R}\,(0 < \varepsilon \to \varphi_{35}) \to \varphi_{48} \to \mathrm{Seq}(a) \land \varphi_{45} \to \varphi_{45}$
\item $\varphi_{178} :\equiv \mathrm{RFn}(f) \to A \subseteq \mathrm{dom}(f) \to c \in A \to \forall a\,(\mathrm{Seq}(a) \to \varphi_{3}) \to \varphi_{121}$
\item $\varphi_{179} :\equiv \varphi_{97} \land (\mathrm{Seq}(a) \land \varphi_{45}) \land \varepsilon \in \mathbb{R} \land 0 < \varepsilon \land \neg \exists N \in \mathbb{N}\,\forall n \in \mathbb{N}\,(N \leq n \to \lvert a_{n} - c \rvert < \varepsilon)$
\item $\varphi_{180} :\equiv \mathrm{RFn}(f) \to A \subseteq \mathrm{dom}(f) \to c \in A \to \forall a\,(\mathrm{Seq}(a) \to \varphi_{3}) \to \varphi_{125}$
\item $\varphi_{181} :\equiv \mathrm{RFn}(f) \to A \subseteq \mathrm{dom}(f) \to c \in A \to \forall a\,(\mathrm{Seq}(a) \to \varphi_{3}) \to \varphi_{128}$
\item $\varphi_{182} :\equiv \mathrm{RFn}(f) \to A \subseteq \mathrm{dom}(f) \to c \in A \to \forall a\,(\mathrm{Seq}(a) \to \varphi_{3}) \to \varphi_{129}$
\item $\varphi_{183} :\equiv c \in A \to \forall a\,(\mathrm{Seq}(a) \to \varphi_{3}) \to \neg \forall \varepsilon \in \mathbb{R}\,(0 < \varepsilon \to \varphi_{35}) \to \varphi_{100}$
\item $\varphi_{184} :\equiv \mathrm{RFn}(f) \to A \subseteq \mathrm{dom}(f) \to c \in A \to \forall a\,(\mathrm{Seq}(a) \to \varphi_{3}) \to \varphi_{131}$
\item $\varphi_{185} :\equiv \mathrm{RFn}(f) \to A \subseteq \mathrm{dom}(f) \to c \in A \to \forall a\,(\mathrm{Seq}(a) \to \varphi_{3}) \to \varphi_{132}$
\item $\varphi_{186} :\equiv \mathrm{RFn}(f) \to A \subseteq \mathrm{dom}(f) \to c \in A \to \forall a\,(\mathrm{Seq}(a) \to \varphi_{3}) \to \varphi_{133}$
\item $\varphi_{187} :\equiv \mathrm{RFn}(f) \to A \subseteq \mathrm{dom}(f) \to c \in A \to \forall a\,(\mathrm{Seq}(a) \to \varphi_{3}) \to \varphi_{135}$
\item $\varphi_{188} :\equiv \mathrm{RFn}(f) \to A \subseteq \mathrm{dom}(f) \to c \in A \to \forall a\,(\mathrm{Seq}(a) \to \varphi_{3}) \to \varphi_{136}$
\item $\varphi_{189} :\equiv \mathrm{RFn}(f) \to A \subseteq \mathrm{dom}(f) \to c \in A \to \forall a\,(\mathrm{Seq}(a) \to \varphi_{3}) \to \varphi_{137}$
\item $\varphi_{190} :\equiv c \in A \to \forall a\,(\mathrm{Seq}(a) \to \varphi_{3}) \to \neg \forall \varepsilon \in \mathbb{R}\,(0 < \varepsilon \to \varphi_{35}) \to \varphi_{48} \to \varphi_{79}$
\item $\varphi_{191} :\equiv \mathrm{RFn}(f) \to A \subseteq \mathrm{dom}(f) \to c \in A \to \forall a\,(\mathrm{Seq}(a) \to \varphi_{3}) \to \varphi_{138}$
\item $\varphi_{192} :\equiv c \in A \to \forall a\,(\mathrm{Seq}(a) \to \varphi_{3}) \to \neg \forall \varepsilon \in \mathbb{R}\,(0 < \varepsilon \to \varphi_{35}) \to \varphi_{48} \to \varphi_{83}$
\item $\varphi_{193} :\equiv \mathrm{RFn}(f) \to A \subseteq \mathrm{dom}(f) \to c \in A \to \forall a\,(\mathrm{Seq}(a) \to \varphi_{3}) \to \varphi_{139}$
\item $\varphi_{194} :\equiv \mathrm{RFn}(f) \to A \subseteq \mathrm{dom}(f) \to c \in A \to \forall a\,(\mathrm{Seq}(a) \to \varphi_{3}) \to \varphi_{141}$
\item $\varphi_{195} :\equiv \forall a\,(\mathrm{Seq}(a) \to \varphi_{3}) \to \neg \forall \varepsilon \in \mathbb{R}\,(0 < \varepsilon \to \varphi_{35}) \to \varphi_{48} \to \mathrm{Seq}(a) \land \varphi_{45} \to \varphi_{9} \to \varphi_{9}$
\item $\varphi_{196} :\equiv \mathrm{RFn}(f) \to A \subseteq \mathrm{dom}(f) \to c \in A \to \forall a\,(\mathrm{Seq}(a) \to \varphi_{3}) \to \varphi_{143}$
\item $\varphi_{197} :\equiv \mathrm{RFn}(f) \to A \subseteq \mathrm{dom}(f) \to c \in A \to \forall a\,(\mathrm{Seq}(a) \to \varphi_{3}) \to \varphi_{144}$
\item $\varphi_{198} :\equiv \mathrm{RFn}(f) \to A \subseteq \mathrm{dom}(f) \to c \in A \to \forall a\,(\mathrm{Seq}(a) \to \varphi_{3}) \to \varphi_{145}$
\item $\varphi_{199} :\equiv (\forall x\,(\neg \varphi_{4}) \to \forall x\,\forall x\,(\neg \varphi_{4})) \to \mathrm{RFn}(f) \to A \subseteq \mathrm{dom}(f) \to c \in A \to \varphi_{114}$
\item $\varphi_{200} :\equiv \forall a\,(\mathrm{Seq}(a) \to \varphi_{3}) \to \neg \forall \varepsilon \in \mathbb{R}\,(0 < \varepsilon \to \varphi_{35}) \to \varphi_{48} \to \mathrm{Seq}(a) \land \varphi_{45} \to \varphi_{9} \to \varphi_{32}$
\item $\varphi_{201} :\equiv \forall a\,(\mathrm{Seq}(a) \to \varphi_{3}) \to \neg \forall \varepsilon \in \mathbb{R}\,(0 < \varepsilon \to \varphi_{35}) \to \varphi_{48} \to \mathrm{Seq}(a) \land \varphi_{45} \to \varphi_{61}$
\item $\varphi_{202} :\equiv (\forall x\,(\neg \varphi_{21}) \to \neg \varphi_{21}) \to \mathrm{RFn}(f) \to A \subseteq \mathrm{dom}(f) \to c \in A \to \varphi_{153}$
\item $\varphi_{203} :\equiv (\varphi_{2} \to \neg (K \in \mathbb{N} \land \forall n \in \mathbb{N}\,(K \leq n \to \lvert a_{n} - c \rvert < \varepsilon))) \to \varphi_{197}$
\item $\varphi_{204} :\equiv (\mathrm{RFn}(f) \to A \subseteq \mathrm{dom}(f) \to c \in A \to \varphi_{91}) \to \mathrm{RFn}(f) \to A \subseteq \mathrm{dom}(f) \to c \in A \to \varphi_{94}$
\item $\varphi_{205} :\equiv (\mathrm{RFn}(f) \to A \subseteq \mathrm{dom}(f) \to c \in A \to \varphi_{94}) \to \mathrm{RFn}(f) \to A \subseteq \mathrm{dom}(f) \to c \in A \to \varphi_{92}$
\item $\varphi_{206} :\equiv (\mathrm{RFn}(f) \to A \subseteq \mathrm{dom}(f) \to c \in A \to \varphi_{91}) \to \mathrm{RFn}(f) \to A \subseteq \mathrm{dom}(f) \to c \in A \to \varphi_{95}$
\item $\varphi_{207} :\equiv (\mathrm{RFn}(f) \to A \subseteq \mathrm{dom}(f) \to c \in A \to \varphi_{94}) \to \mathrm{RFn}(f) \to A \subseteq \mathrm{dom}(f) \to c \in A \to \varphi_{96}$
\item $\varphi_{208} :\equiv (\mathrm{RFn}(f) \to A \subseteq \mathrm{dom}(f) \to c \in A \to \varphi_{93}) \to \mathrm{RFn}(f) \to A \subseteq \mathrm{dom}(f) \to c \in A \to \varphi_{98}$
\item $\varphi_{209} :\equiv (\mathrm{RFn}(f) \to A \subseteq \mathrm{dom}(f) \to c \in A \to \varphi_{96}) \to \mathrm{RFn}(f) \to A \subseteq \mathrm{dom}(f) \to c \in A \to \varphi_{113}$
\item $\varphi_{210} :\equiv (\mathrm{RFn}(f) \to A \subseteq \mathrm{dom}(f) \to \varphi_{123}) \to \mathrm{RFn}(f) \to A \subseteq \mathrm{dom}(f) \to c \in A \to \varphi_{88}$
\item $\varphi_{211} :\equiv (\mathrm{RFn}(f) \to A \subseteq \mathrm{dom}(f) \to \varphi_{123}) \to \mathrm{RFn}(f) \to A \subseteq \mathrm{dom}(f) \to c \in A \to \varphi_{89}$
\item $\varphi_{212} :\equiv (\mathrm{RFn}(f) \to A \subseteq \mathrm{dom}(f) \to c \in A \to \varphi_{118}) \to \varphi_{97} \to k \in \mathbb{N} \to \varphi_{27}$
\item $\varphi_{213} :\equiv (\mathrm{RFn}(f) \to A \subseteq \mathrm{dom}(f) \to c \in A \to \varphi_{88}) \to \mathrm{RFn}(f) \to A \subseteq \mathrm{dom}(f) \to c \in A \to \varphi_{166}$
\item $\varphi_{214} :\equiv (\mathrm{RFn}(f) \to A \subseteq \mathrm{dom}(f) \to \varphi_{123}) \to \mathrm{RFn}(f) \to A \subseteq \mathrm{dom}(f) \to \varphi_{119}$
\item $\varphi_{215} :\equiv (\mathrm{RFn}(f) \to \mathrm{RFn}(f)) \to \varphi_{127} \to \mathrm{RFn}(f) \to A \subseteq \mathrm{dom}(f) \to c \in A \to \varphi_{117}$
\item $\varphi_{216} :\equiv (\mathrm{RFn}(f) \to A \subseteq \mathrm{dom}(f) \to c \in A \to \varphi_{120}) \to \mathrm{RFn}(f) \to A \subseteq \mathrm{dom}(f) \to \varphi_{124}$
\item $\varphi_{217} :\equiv (\mathrm{RFn}(f) \to A \subseteq \mathrm{dom}(f) \to A \subseteq \mathrm{dom}(f)) \to (\mathrm{RFn}(f) \to A \subseteq \mathrm{dom}(f) \to c \in A \to \varphi_{116}) \to \varphi_{127}$
\item $\varphi_{218} :\equiv (\mathrm{RFn}(f) \to A \subseteq \mathrm{dom}(f) \to c \in A \to \varphi_{146}) \to \mathrm{RFn}(f) \to A \subseteq \mathrm{dom}(f) \to c \in A \to \varphi_{116}$
\item $\varphi_{219} :\equiv (\mathrm{RFn}(f) \to A \subseteq \mathrm{dom}(f) \to c \in A \to \varphi_{122}) \to \varphi_{97} \land \neg 0 < 0 \to \neg (\mathrm{Seq}(a) \land \varphi_{45})$
\item $\varphi_{220} :\equiv \varphi_{97} \land (\mathrm{Seq}(a) \land \varphi_{45}) \land \neg 0 < 0 \to \forall N\,(\varphi_{97} \land (\mathrm{Seq}(a) \land \varphi_{45}) \land \neg 0 < 0)$
\item $\varphi_{221} :\equiv \forall a\,(\varphi_{97} \land \neg 0 < 0 \to \neg (\mathrm{Seq}(a) \land \varphi_{45})) \to \forall a\,(\varphi_{97} \land \neg 0 < 0) \to \forall a\,(\neg (\mathrm{Seq}(a) \land \varphi_{45}))$
\item $\varphi_{222} :\equiv (\mathrm{RFn}(f) \to A \subseteq \mathrm{dom}(f) \to c \in A \to \varphi_{117}) \to \mathrm{RFn}(f) \to A \subseteq \mathrm{dom}(f) \to c \in A \to \varphi_{152}$
\item $\varphi_{223} :\equiv \forall n\,(\varphi_{97} \land (\mathrm{Seq}(a) \land \varphi_{45}) \to n \in \mathbb{N} \to a_{n} \in A) \to \forall n\,(\varphi_{97} \land (\mathrm{Seq}(a) \land \varphi_{45})) \to \forall n \in \mathbb{N}\,(a_{n} \in A)$
\item $\varphi_{224} :\equiv (\mathrm{RFn}(f) \to A \subseteq \mathrm{dom}(f) \to c \in A \to \varphi_{153}) \to \varphi_{97} \land k \in \mathbb{N} \land \neg \varphi_{27} \to \neg \varphi_{4}$
\item $\varphi_{225} :\equiv \forall x\,(\varphi_{97} \land k \in \mathbb{N} \land \neg \varphi_{27} \to \neg \varphi_{4}) \to \forall x\,(\varphi_{97} \land k \in \mathbb{N} \land \neg \varphi_{27}) \to \forall x\,(\neg \varphi_{4})$
\item $\varphi_{226} :\equiv (A \subseteq \mathrm{dom}(f) \to a_{n} \in A \to a_{n} \in \mathrm{dom}(f)) \to (\mathrm{RFn}(f) \to A \subseteq \mathrm{dom}(f) \to A \subseteq \mathrm{dom}(f)) \to \varphi_{156} \to \varphi_{160}$
\item $\varphi_{227} :\equiv (\mathrm{RFn}(f) \to A \subseteq \mathrm{dom}(f) \to c \in A \to \varphi_{173}) \to \mathrm{RFn}(f) \to A \subseteq \mathrm{dom}(f) \to \varphi_{167}$
\item $\varphi_{228} :\equiv (\mathrm{RFn}(f) \to A \subseteq \mathrm{dom}(f) \to c \in A \to \varphi_{171}) \to \varphi_{97} \land (\mathrm{Seq}(a) \land \varphi_{45}) \land \neg 0 < 0 \to \neg \varphi_{9}$
\item $\varphi_{229} :\equiv (\mathrm{RFn}(f) \to A \subseteq \mathrm{dom}(f) \to c \in A \to \varphi_{172}) \to \varphi_{97} \land (\mathrm{Seq}(a) \land \varphi_{45}) \to \varphi_{24}$
\item $\varphi_{230} :\equiv \forall N\,(\varphi_{97} \land (\mathrm{Seq}(a) \land \varphi_{45}) \land \neg 0 < 0 \to \neg \varphi_{9}) \to \forall N\,(\varphi_{97} \land (\mathrm{Seq}(a) \land \varphi_{45}) \land \neg 0 < 0) \to \forall N\,(\neg \varphi_{9})$
\item $\varphi_{231} :\equiv \forall \varepsilon\,(\varphi_{97} \land (\mathrm{Seq}(a) \land \varphi_{45}) \to \varphi_{24}) \to \forall \varepsilon\,(\varphi_{97} \land (\mathrm{Seq}(a) \land \varphi_{45})) \to \varphi_{30}$
\item $\varphi_{232} :\equiv (\mathrm{RFn}(f) \to A \subseteq \mathrm{dom}(f) \to \varphi_{119}) \to (\varphi_{34} \to \neg (\underline{k}^{-1} \in \mathbb{R} \land 0 < \underline{k}^{-1} \land \varphi_{8})) \to \varphi_{142}$
\item $\varphi_{233} :\equiv (\mathrm{RFn}(f) \to A \subseteq \mathrm{dom}(f) \to \varphi_{168}) \to \mathrm{RFn}(f) \to A \subseteq \mathrm{dom}(f) \to \varphi_{183}$
\item $\varphi_{234} :\equiv \forall n\,(\varphi_{175} \to n \in \mathbb{N} \to K \leq n \to \lvert a_{n} - c \rvert < \varepsilon) \to \forall n\,(\varphi_{175}) \to \forall n \in \mathbb{N}\,(K \leq n \to \lvert a_{n} - c \rvert < \varepsilon)$
\item $\varphi_{235} :\equiv (\varphi_{97} \to \forall k\,(\varphi_{97})) \to \mathrm{RFn}(f) \to A \subseteq \mathrm{dom}(f) \to c \in A \to \varphi_{120}$
\item $\varphi_{236} :\equiv (\mathrm{RFn}(f) \to A \subseteq \mathrm{dom}(f) \to \varphi_{174}) \to \mathrm{RFn}(f) \to A \subseteq \mathrm{dom}(f) \to c \in A \to \varphi_{201}$
\item $\varphi_{237} :\equiv (\mathrm{RFn}(f) \to A \subseteq \mathrm{dom}(f) \to c \in A \to \varphi_{200}) \to \mathrm{RFn}(f) \to A \subseteq \mathrm{dom}(f) \to c \in A \to \varphi_{171}$
\item $\varphi_{238} :\equiv \forall K\,(\varphi_{179} \to \neg (K \in \mathbb{N} \land \underline{K}^{-1} < \varepsilon)) \to \forall K\,(\varphi_{179}) \to \forall K\,(\neg (K \in \mathbb{N} \land \underline{K}^{-1} < \varepsilon))$
\item $\varphi_{239} :\equiv (\underline{k} \in \mathbb{R} \to 0 < \underline{k} \to 0 < \underline{k}^{-1}) \to (\mathrm{RFn}(f) \to A \subseteq \mathrm{dom}(f) \to c \in A \to \varphi_{95}) \to \varphi_{207}$
\item $\varphi_{240} :\equiv (\underline{k} \in \mathbb{R} \to \underline{k} \neq 0 \to \underline{k}^{-1} \in \mathbb{R}) \to (\mathrm{RFn}(f) \to A \subseteq \mathrm{dom}(f) \to c \in A \to \varphi_{95}) \to \varphi_{208}$
\item $\varphi_{241} :\equiv (\mathrm{RFn}(f) \to A \subseteq \mathrm{dom}(f) \to c \in A \to \varphi_{195}) \to \mathrm{RFn}(f) \to A \subseteq \mathrm{dom}(f) \to \varphi_{192}$
\item $\varphi_{242} :\equiv \forall q\,(\forall a\,(\neg (\mathrm{Seq}(a) \land \varphi_{45})) \to \varphi_{49}) \to \varphi_{148} \to \varphi_{73} \to \mathrm{RFn}(f) \to A \subseteq \mathrm{dom}(f) \to c \in A \to \varphi_{87}$
\item $\varphi_{243} :\equiv (\mathrm{RFn}(f) \to A \subseteq \mathrm{dom}(f) \to \varphi_{192}) \to \varphi_{54} \to \mathrm{RFn}(f) \to A \subseteq \mathrm{dom}(f) \to \varphi_{190}$
\item $\varphi_{244} :\equiv (\mathrm{RFn}(f) \to A \subseteq \mathrm{dom}(f) \to c \in A \to \varphi_{177}) \to (\varphi_{45} \to N \in \mathbb{N} \to \varphi_{32}) \to \varphi_{184}$
\item $\varphi_{245} :\equiv (\mathrm{RFn}(f) \to A \subseteq \mathrm{dom}(f) \to c \in A \to \varphi_{177}) \to (\varphi_{45} \to n \in \mathbb{N} \to \varphi_{33}) \to \varphi_{185}$
\item $\varphi_{246} :\equiv (\varphi_{97} \land k \in \mathbb{N} \land \neg \varphi_{27} \to \forall x\,(\varphi_{97} \land k \in \mathbb{N} \land \neg \varphi_{27})) \to \mathrm{RFn}(f) \to A \subseteq \mathrm{dom}(f) \to c \in A \to \varphi_{118}$
\item $\varphi_{247} :\equiv (\varphi_{97} \land (\mathrm{Seq}(a) \land \varphi_{45}) \to \forall n\,(\varphi_{97} \land (\mathrm{Seq}(a) \land \varphi_{45}))) \to \varphi_{159}$
\item $\varphi_{248} :\equiv (\varphi_{97} \land (\mathrm{Seq}(a) \land \varphi_{45}) \to \forall n\,(\varphi_{97} \land (\mathrm{Seq}(a) \land \varphi_{45}))) \to \varphi_{162}$
\item $\varphi_{249} :\equiv (\mathrm{RFn}(f) \to A \subseteq \mathrm{dom}(f) \to c \in A \to \varphi_{113}) \to \forall x\,(\varphi_{60}) \to (\forall x\,(\varphi_{60}) \to \forall x\,\forall x\,(\neg \varphi_{4}) \to \varphi_{8}) \to \varphi_{199}$
\item $\varphi_{250} :\equiv (\varphi_{97} \land (\mathrm{Seq}(a) \land \varphi_{45}) \to \forall \varepsilon\,(\varphi_{97} \land (\mathrm{Seq}(a) \land \varphi_{45}))) \to \mathrm{RFn}(f) \to A \subseteq \mathrm{dom}(f) \to c \in A \to \varphi_{176}$
\item $\varphi_{251} :\equiv (\mathrm{RFn}(f) \to \mathrm{RFn}(f)) \to \varphi_{147} \to \varphi_{162} \to \mathrm{RFn}(f) \to A \subseteq \mathrm{dom}(f) \to \varphi_{168}$
\item $\varphi_{252} :\equiv (\varepsilon \in \mathbb{R} \to 0 < \varepsilon \to \exists K \in \mathbb{N}\,(\underline{K}^{-1} < \varepsilon)) \to \varphi_{157} \to \varphi_{163} \to \mathrm{RFn}(f) \to A \subseteq \mathrm{dom}(f) \to \varphi_{164}$
\item $\varphi_{253} :\equiv (\mathrm{RFn}(f) \to A \subseteq \mathrm{dom}(f) \to \varphi_{174}) \to (\mathrm{RFn}(f) \to A \subseteq \mathrm{dom}(f) \to c \in A \to \varphi_{170}) \to \varphi_{181}$
\item $\varphi_{254} :\equiv (\varphi_{179} \to \forall K\,(\varphi_{179})) \to \mathrm{RFn}(f) \to A \subseteq \mathrm{dom}(f) \to c \in A \to \varphi_{172}$
\item $\varphi_{255} :\equiv (\mathrm{RFn}(f) \to A \subseteq \mathrm{dom}(f) \to c \in A \to \varphi_{86}) \to \varphi_{75} \to (\varphi_{75} \to \varphi_{80}) \to \varphi_{82} \to \varphi_{102}$
\item $\varphi_{256} :\equiv (\varphi_{78} \to \forall e_{0}\,(\varphi_{78})) \to \forall \varepsilon\,(\forall e_{0}\,(\neg \varphi_{48}) \to \neg \neg (\varepsilon \in \mathbb{R} \to 0 < \varepsilon \to \varphi_{35})) \to \varphi_{154} \to \varphi_{155}$
\item $\varphi_{257} :\equiv (\mathrm{Conv}(a,c) \leftrightarrow \mathrm{Seq}(a) \land c \in \mathbb{R} \land \varphi_{30}) \to \varphi_{147} \to \varphi_{56} \to (\mathrm{RFn}(f) \to A \subseteq \mathrm{dom}(f) \to c \in A \to \varphi_{176}) \to \varphi_{149}$
\item $\varphi_{258} :\equiv (\lvert a_{n} - c \rvert < \underline{K}^{-1} \to \underline{K}^{-1} < \varepsilon \to \lvert a_{n} - c \rvert < \varepsilon) \to \varphi_{188} \to \varphi_{180} \to \varphi_{182}$
\item $\varphi_{259} :\equiv (\mathrm{Conv}(f \circ a,f(c)) \to e_{0} \in \mathbb{R} \to 0 < e_{0} \to \varphi_{11}) \to \varphi_{158} \to (\mathrm{RFn}(f) \to A \subseteq \mathrm{dom}(f) \to c \in A \to \varphi_{89}) \to \varphi_{213}$
\item $\varphi_{260} :\equiv (\forall k\,(\varphi_{97} \to k \in \mathbb{N} \to \varphi_{27}) \to \forall k\,(\varphi_{97}) \to \forall k \in \mathbb{N}\,(\varphi_{27})) \to \varphi_{235}$
\item $\varphi_{261} :\equiv (K \in \mathbb{N} \to n \in \mathbb{N} \to K \leq n \to \underline{n}^{-1} \leq \underline{K}^{-1}) \to (\mathrm{RFn}(f) \to A \subseteq \mathrm{dom}(f) \to \varphi_{167}) \to \varphi_{253}$
\item $\varphi_{262} :\equiv (\mathrm{RFn}(f) \to A \subseteq \mathrm{dom}(f) \to \varphi_{190}) \to \varphi_{178} \to (\mathrm{RFn}(f) \to A \subseteq \mathrm{dom}(f) \to c \in A \to \varphi_{169}) \to \varphi_{196}$
\item $\varphi_{263} :\equiv \forall e_{0}\,(\varphi_{78} \to \neg \varphi_{48}) \to (\forall e_{0}\,(\varphi_{78} \to \neg \varphi_{48}) \to \forall e_{0}\,(\varphi_{78}) \to \forall e_{0}\,(\neg \varphi_{48})) \to \varphi_{256}$
\item $\varphi_{264} :\equiv \forall N\,(\varphi_{97} \land (\mathrm{Seq}(a) \land \varphi_{45}) \land \neg 0 < 0 \to \neg \varphi_{9}) \to \varphi_{230} \to \varphi_{220} \to \mathrm{RFn}(f) \to A \subseteq \mathrm{dom}(f) \to c \in A \to \varphi_{122}$
\item $\varphi_{265} :\equiv \forall n\,(\varphi_{175} \to n \in \mathbb{N} \to K \leq n \to \lvert a_{n} - c \rvert < \varepsilon) \to \varphi_{234} \to (\varphi_{175} \to \forall n\,(\varphi_{175})) \to \varphi_{186}$
\item $\varphi_{266} :\equiv (\mathrm{RFn}(f) \to A \subseteq \mathrm{dom}(f) \to c \in A \to \varphi_{114}) \to \forall x\,(\varphi_{97} \land k \in \mathbb{N} \land \neg \varphi_{27} \to \neg \varphi_{4}) \to \varphi_{225} \to \varphi_{246}$
\item $\varphi_{267} :\equiv \forall a\,(\varphi_{97} \land \neg 0 < 0 \to \neg (\mathrm{Seq}(a) \land \varphi_{45})) \to \varphi_{221} \to (\varphi_{97} \land \neg 0 < 0 \to \forall a\,(\varphi_{97} \land \neg 0 < 0)) \to \varphi_{242}$
\item $\varphi_{268} :\equiv (\mathrm{RFn}(f) \to A \subseteq \mathrm{dom}(f) \to \varphi_{164}) \to \forall K\,(\varphi_{179} \to \neg (K \in \mathbb{N} \land \underline{K}^{-1} < \varepsilon)) \to \varphi_{238} \to \varphi_{254}$
\end{abbrevlist}
\end{fullproof}

% 自動生成 (prover/build.js): ContInSub  ノード数 159
\begin{fullproof}{[G3]}{$\forall f\,\forall A\,\forall B\,\forall c\,(\mathrm{ContIn}(f,A,c) \to B \subseteq A \to c \in B \to \mathrm{ContIn}(f,B,c))$}{159}\phantomsection\label{tree:ContInSub}
\noindent$\Pi_{1}$:
\begin{ptree}
\AXC{$\forall X\,\forall Y\,\forall x\,(X \subseteq Y \to x \in X \to x \in Y)$\lbl{[B5]}}
\AXC{$\varphi_{9}$\lbl{A2}}
\BIC{$\forall Y\,\forall x\,(B \subseteq Y \to x \in B \to x \in Y)$}
\AXC{$\varphi_{7}$\lbl{A2}}
\BIC{$\forall x\,(B \subseteq A \to x \in B \to x \in A)$}
\AXC{$\forall x\,(B \subseteq A \to x \in B \to x \in A) \to B \subseteq A \to x \in B \to x \in A$\lbl{A2}}
\BIC{$B \subseteq A \to x \in B \to x \in A$}
\end{ptree}
\noindent$\Pi_{2}$:
\begin{ptree}
\AXC{$\varphi_{37}$\lbl{A1}}
\AXC{$\mathrm{ContIn}(f,A,c) \to B \subseteq A \to B \subseteq A$\lbl{A1}}
\AXC{$\stackrel{\Pi_{1}}{\vdots}$}\noLine
\UIC{$B \subseteq A \to x \in B \to x \in A$}
\AXC{$(B \subseteq A \to x \in B \to x \in A) \to \varphi_{58}$\lbl{A1}}
\BIC{$\varphi_{58}$}
\BIC{$\varphi_{37} \to \mathrm{ContIn}(f,A,c) \to \varphi_{41}$}
\BIC{$\mathrm{ContIn}(f,A,c) \to \varphi_{41}$}
\end{ptree}
\noindent$\Pi_{3}$:
\begin{ptree}
\AXC{$\varphi_{24}$\lbl{A2}}
\AXC{$\varphi_{46}$\lbl{A1}}
\AXC{$\varphi_{46} \to \varphi_{51}$\lbl{A1}}
\BIC{$\varphi_{51}$}
\AXC{$\varphi_{51} \to \varphi_{24} \to \varphi_{49}$\lbl{A1}}
\BIC{$\varphi_{24} \to \varphi_{49}$}
\BIC{$\varphi_{49}$}
\AXC{$\varphi_{49} \to (\mathrm{ContIn}(f,A,c) \to \varphi_{41}) \to \varphi_{59}$\lbl{A1}}
\BIC{$(\mathrm{ContIn}(f,A,c) \to \varphi_{41}) \to \varphi_{59}$}
\end{ptree}
\noindent$\Pi_{4}$:
\begin{ptree}
\AXC{$\mathrm{ContIn}(f,A,c) \to B \subseteq A \to \varphi_{40}$\lbl{A1}}
\AXC{$\stackrel{\Pi_{2}}{\vdots}$}\noLine
\UIC{$\mathrm{ContIn}(f,A,c) \to \varphi_{41}$}
\AXC{$\stackrel{\Pi_{3}}{\vdots}$}\noLine
\UIC{$(\mathrm{ContIn}(f,A,c) \to \varphi_{41}) \to \varphi_{59}$}
\BIC{$\varphi_{59}$}
\BIC{$\mathrm{ContIn}(f,A,c) \to B \subseteq A \to c \in B \to \varepsilon \in \mathbb{R} \to \varphi_{31}$}
\AXC{$\varphi_{60}$\lbl{A1}}
\BIC{$\varphi_{25} \to x \in B \to \lvert x - c \rvert < \delta \to \lvert f(x) - f(c) \rvert < \varepsilon$}
\UIC{$\varphi_{43}$}
\end{ptree}
\noindent$\Pi_{5}$:
\begin{ptree}
\AXC{$\varphi_{25} \to \forall x\,(\varphi_{25})$\lbl{A4}}
\AXC{$\varphi_{43} \to \varphi_{44}$\lbl{A3}}
\AXC{$\stackrel{\Pi_{4}}{\vdots}$}\noLine
\UIC{$\varphi_{43}$}
\AXC{$\varphi_{43} \to (\varphi_{43} \to \varphi_{44}) \to \varphi_{62}$\lbl{A1}}
\BIC{$(\varphi_{43} \to \varphi_{44}) \to \varphi_{62}$}
\BIC{$\varphi_{62}$}
\BIC{$\mathrm{ContIn}(f,A,c) \to B \subseteq A \to c \in B \to \varepsilon \in \mathbb{R} \to \varphi_{34}$}
\AXC{$\varphi_{46}$\lbl{A1}}
\AXC{$\varphi_{46} \to \varphi_{61}$\lbl{A1}}
\BIC{$\varphi_{61}$}
\BIC{$\varphi_{47}$}
\end{ptree}
\noindent$\Pi_{6}$:
\begin{ptree}
\AXC{$\forall \delta\,(\neg \varphi_{20}) \to \neg \varphi_{20}$\lbl{A2}}
\AXC{$\stackrel{\Pi_{5}}{\vdots}$}\noLine
\UIC{$\varphi_{47}$}
\AXC{$\varphi_{47} \to (\forall \delta\,(\neg \varphi_{20}) \to \neg \varphi_{20}) \to \varphi_{50}$\lbl{A1}}
\BIC{$(\forall \delta\,(\neg \varphi_{20}) \to \neg \varphi_{20}) \to \varphi_{50}$}
\BIC{$\varphi_{50}$}
\AXC{$\varphi_{50} \to \varphi_{29} \to \neg \varphi_{19}$\lbl{A1}}
\BIC{$\varphi_{29} \to \neg \varphi_{19}$}
\UIC{$\forall \delta\,(\varphi_{29} \to \neg \varphi_{19})$}
\end{ptree}
\noindent$\Pi_{7}$:
\begin{ptree}
\AXC{$\forall f\,\forall A\,\forall c\,(\varphi_{38})$\lbl{DefContIn}}
\AXC{$\forall f\,\forall A\,\forall c\,(\varphi_{38}) \to \forall A\,\forall c\,(\varphi_{38})$\lbl{A2}}
\BIC{$\forall A\,\forall c\,(\varphi_{38})$}
\AXC{$\forall A\,\forall c\,(\varphi_{38}) \to \forall c\,(\varphi_{38})$\lbl{A2}}
\BIC{$\forall c\,(\varphi_{38})$}
\AXC{$\forall c\,(\varphi_{38}) \to \varphi_{38}$\lbl{A2}}
\BIC{$\varphi_{38}$}
\end{ptree}
\noindent$\Pi_{8}$:
\begin{ptree}
\AXC{$\varphi_{48}$\lbl{A2}}
\AXC{$\stackrel{\Pi_{7}}{\vdots}$}\noLine
\UIC{$\varphi_{38}$}
\AXC{$\mathrm{ContIn}(f,A,c) \to \mathrm{ContIn}(f,A,c)$\lbl{A1}}
\AXC{$\varphi_{57}$\lbl{A1}}
\BIC{$\varphi_{38} \to \mathrm{ContIn}(f,A,c) \to \varphi_{36}$}
\BIC{$\mathrm{ContIn}(f,A,c) \to \varphi_{36}$}
\AXC{$\varphi_{63}$\lbl{A1}}
\BIC{$\varphi_{48} \to \mathrm{ContIn}(f,A,c) \to \varphi_{33}$}
\BIC{$\mathrm{ContIn}(f,A,c) \to \varphi_{33}$}
\end{ptree}
\noindent$\Pi_{9}$:
\begin{ptree}
\AXC{$\varphi_{8}$\lbl{A1}}
\AXC{$\varphi_{5}$\lbl{A1}}
\AXC{$\stackrel{\Pi_{8}}{\vdots}$}\noLine
\UIC{$\mathrm{ContIn}(f,A,c) \to \varphi_{33}$}
\AXC{$(\mathrm{ContIn}(f,A,c) \to \varphi_{33}) \to \varphi_{5} \to \varphi_{8} \to \varphi_{26}$\lbl{A1}}
\BIC{$\varphi_{5} \to \varphi_{8} \to \varphi_{26}$}
\BIC{$\varphi_{8} \to \varphi_{26}$}
\BIC{$\varphi_{26}$}
\end{ptree}
\noindent$\Pi_{10}$:
\begin{ptree}
\AXC{$\forall \delta\,(\varphi_{29} \to \neg \varphi_{19}) \to \forall \delta\,(\varphi_{29}) \to \forall \delta\,(\neg \varphi_{19})$\lbl{A3}}
\AXC{$\stackrel{\Pi_{6}}{\vdots}$}\noLine
\UIC{$\forall \delta\,(\varphi_{29} \to \neg \varphi_{19})$}
\AXC{$\stackrel{\Pi_{9}}{\vdots}$}\noLine
\UIC{$\varphi_{26}$}
\AXC{$\varphi_{26} \to \forall \delta\,(\varphi_{29} \to \neg \varphi_{19}) \to \varphi_{64}$\lbl{A1}}
\BIC{$\forall \delta\,(\varphi_{29} \to \neg \varphi_{19}) \to \varphi_{64}$}
\BIC{$\varphi_{64}$}
\BIC{$(\varphi_{29} \to \forall \delta\,(\varphi_{29})) \to \varphi_{27}$}
\end{ptree}
\noindent$\Pi_{11}$:
\begin{ptree}
\AXC{$\varphi_{29} \to \forall \delta\,(\varphi_{29})$\lbl{A4}}
\AXC{$\stackrel{\Pi_{10}}{\vdots}$}\noLine
\UIC{$(\varphi_{29} \to \forall \delta\,(\varphi_{29})) \to \varphi_{27}$}
\BIC{$\varphi_{27}$}
\AXC{$\varphi_{27} \to \varphi_{28}$\lbl{A1}}
\BIC{$\varphi_{28}$}
\UIC{$\forall \varepsilon\,(\varphi_{28})$}
\AXC{$\forall \varepsilon\,(\varphi_{28}) \to (\forall \varepsilon\,(\varphi_{28}) \to \varphi_{32}) \to \varphi_{13} \to \varphi_{30}$\lbl{A1}}
\BIC{$(\forall \varepsilon\,(\varphi_{28}) \to \varphi_{32}) \to \varphi_{13} \to \varphi_{30}$}
\end{ptree}
\noindent$\Pi_{12}$:
\begin{ptree}
\AXC{$\forall X\,\forall Y\,\forall x\,(X \subseteq Y \to x \in X \to x \in Y)$\lbl{[B5]}}
\AXC{$\varphi_{9}$\lbl{A2}}
\BIC{$\forall Y\,\forall x\,(B \subseteq Y \to x \in B \to x \in Y)$}
\AXC{$\varphi_{7}$\lbl{A2}}
\BIC{$\forall x\,(B \subseteq A \to x \in B \to x \in A)$}
\AXC{$\forall x\,(B \subseteq A \to x \in B \to x \in A) \to B \subseteq A \to u \in B \to u \in A$\lbl{A2}}
\BIC{$B \subseteq A \to u \in B \to u \in A$}
\end{ptree}
\noindent$\Pi_{13}$:
\begin{ptree}
\AXC{$\mathrm{ContIn}(f,A,c) \to B \subseteq A \to B \subseteq A$\lbl{A1}}
\AXC{$\stackrel{\Pi_{12}}{\vdots}$}\noLine
\UIC{$B \subseteq A \to u \in B \to u \in A$}
\AXC{$\varphi_{42}$\lbl{A1}}
\BIC{$(\mathrm{ContIn}(f,A,c) \to B \subseteq A \to B \subseteq A) \to \varphi_{21}$}
\BIC{$\varphi_{21}$}
\end{ptree}
\noindent$\Pi_{14}$:
\begin{ptree}
\AXC{$\forall f\,\forall A\,\forall c\,(\varphi_{38})$\lbl{DefContIn}}
\AXC{$\forall f\,\forall A\,\forall c\,(\varphi_{38}) \to \forall A\,\forall c\,(\varphi_{38})$\lbl{A2}}
\BIC{$\forall A\,\forall c\,(\varphi_{38})$}
\AXC{$\forall A\,\forall c\,(\varphi_{38}) \to \forall c\,(\varphi_{38})$\lbl{A2}}
\BIC{$\forall c\,(\varphi_{38})$}
\AXC{$\forall c\,(\varphi_{38}) \to \varphi_{38}$\lbl{A2}}
\BIC{$\varphi_{38}$}
\end{ptree}
\noindent$\Pi_{15}$:
\begin{ptree}
\AXC{$\forall X\,\forall Y\,\forall x\,(X \subseteq Y \to x \in X \to x \in Y)$\lbl{[B5]}}
\AXC{$\varphi_{11}$\lbl{A2}}
\BIC{$\forall Y\,\forall x\,(A \subseteq Y \to x \in A \to x \in Y)$}
\AXC{$\varphi_{14}$\lbl{A2}}
\BIC{$\forall x\,(A \subseteq \mathrm{dom}(f) \to x \in A \to x \in \mathrm{dom}(f))$}
\AXC{$\varphi_{16}$\lbl{A2}}
\BIC{$A \subseteq \mathrm{dom}(f) \to u \in A \to u \in \mathrm{dom}(f)$}
\end{ptree}
\noindent$\Pi_{16}$:
\begin{ptree}
\AXC{$\stackrel{\Pi_{15}}{\vdots}$}\noLine
\UIC{$A \subseteq \mathrm{dom}(f) \to u \in A \to u \in \mathrm{dom}(f)$}
\AXC{$(A \subseteq \mathrm{dom}(f) \to u \in A \to u \in \mathrm{dom}(f)) \to \varphi_{35}$\lbl{A1}}
\BIC{$\varphi_{35}$}
\end{ptree}
\noindent$\Pi_{17}$:
\begin{ptree}
\AXC{$\stackrel{\Pi_{14}}{\vdots}$}\noLine
\UIC{$\varphi_{38}$}
\AXC{$\mathrm{ContIn}(f,A,c) \to \mathrm{ContIn}(f,A,c)$\lbl{A1}}
\AXC{$\varphi_{53}$\lbl{A1}}
\BIC{$\varphi_{45}$}
\BIC{$\mathrm{ContIn}(f,A,c) \to B \subseteq A \to c \in B \to A \subseteq \mathrm{dom}(f)$}
\AXC{$\stackrel{\Pi_{16}}{\vdots}$}\noLine
\UIC{$\varphi_{35}$}
\BIC{$\varphi_{23}$}
\end{ptree}
\noindent$\Pi_{18}$:
\begin{ptree}
\AXC{$\mathrm{ContIn}(f,A,c) \to B \subseteq A \to c \in B \to u \in B \to u \in B$\lbl{A1}}
\AXC{$\stackrel{\Pi_{13}}{\vdots}$}\noLine
\UIC{$\varphi_{21}$}
\BIC{$\mathrm{ContIn}(f,A,c) \to B \subseteq A \to c \in B \to u \in B \to u \in A$}
\AXC{$\stackrel{\Pi_{17}}{\vdots}$}\noLine
\UIC{$\varphi_{23}$}
\BIC{$\varphi_{2}$}
\AXC{$\varphi_{2} \to \varphi_{3}$\lbl{A1}}
\BIC{$\varphi_{3}$}
\UIC{$\forall u\,(\varphi_{3})$}
\end{ptree}
\noindent$\Pi_{19}$:
\begin{ptree}
\AXC{$\forall x\,\forall y\,(x \subseteq y \leftrightarrow \forall u \in x\,(u \in y))$\lbl{Def$\subseteq$}}
\AXC{$\varphi_{12}$\lbl{A2}}
\BIC{$\forall y\,(B \subseteq y \leftrightarrow \forall u \in B\,(u \in y))$}
\AXC{$\varphi_{15}$\lbl{A2}}
\BIC{$B \subseteq \mathrm{dom}(f) \leftrightarrow \forall u \in B\,(u \in \mathrm{dom}(f))$}
\AXC{$(B \subseteq \mathrm{dom}(f) \leftrightarrow \forall u \in B\,(u \in \mathrm{dom}(f))) \to \varphi_{22}$\lbl{A1}}
\BIC{$\varphi_{22}$}
\end{ptree}
\noindent$\Pi_{20}$:
\begin{ptree}
\AXC{$\varphi_{10}$\lbl{A4}}
\AXC{$\forall u\,(\varphi_{3}) \to \varphi_{6}$\lbl{A3}}
\AXC{$\stackrel{\Pi_{18}}{\vdots}$}\noLine
\UIC{$\forall u\,(\varphi_{3})$}
\AXC{$\forall u\,(\varphi_{3}) \to (\forall u\,(\varphi_{3}) \to \varphi_{6}) \to \varphi_{10} \to \varphi_{4}$\lbl{A1}}
\BIC{$(\forall u\,(\varphi_{3}) \to \varphi_{6}) \to \varphi_{10} \to \varphi_{4}$}
\BIC{$\varphi_{10} \to \varphi_{4}$}
\BIC{$\varphi_{4}$}
\AXC{$\stackrel{\Pi_{19}}{\vdots}$}\noLine
\UIC{$\varphi_{22}$}
\BIC{$\mathrm{ContIn}(f,A,c) \to B \subseteq A \to c \in B \to B \subseteq \mathrm{dom}(f)$}
\end{ptree}
\noindent$\Pi_{21}$:
\begin{ptree}
\AXC{$\forall f\,\forall A\,\forall c\,(\varphi_{38})$\lbl{DefContIn}}
\AXC{$\forall f\,\forall A\,\forall c\,(\varphi_{38}) \to \forall A\,\forall c\,(\varphi_{38})$\lbl{A2}}
\BIC{$\forall A\,\forall c\,(\varphi_{38})$}
\AXC{$\forall A\,\forall c\,(\varphi_{38}) \to \forall c\,(\varphi_{38})$\lbl{A2}}
\BIC{$\forall c\,(\varphi_{38})$}
\AXC{$\forall c\,(\varphi_{38}) \to \varphi_{38}$\lbl{A2}}
\BIC{$\varphi_{38}$}
\end{ptree}
\noindent$\Pi_{22}$:
\begin{ptree}
\AXC{$\forall f\,\forall A\,\forall c\,(\varphi_{38})$\lbl{DefContIn}}
\AXC{$\forall f\,\forall A\,\forall c\,(\varphi_{38}) \to \forall A\,\forall c\,(\varphi_{38})$\lbl{A2}}
\BIC{$\forall A\,\forall c\,(\varphi_{38})$}
\AXC{$\forall A\,\forall c\,(\varphi_{38}) \to \forall c\,(\varphi_{39})$\lbl{A2}}
\BIC{$\forall c\,(\varphi_{39})$}
\AXC{$\forall c\,(\varphi_{39}) \to \varphi_{39}$\lbl{A2}}
\BIC{$\varphi_{39}$}
\end{ptree}
\noindent$\Pi_{23}$:
\begin{ptree}
\AXC{$\stackrel{\Pi_{21}}{\vdots}$}\noLine
\UIC{$\varphi_{38}$}
\AXC{$\mathrm{ContIn}(f,A,c) \to \mathrm{ContIn}(f,A,c)$\lbl{A1}}
\AXC{$\varphi_{52}$\lbl{A1}}
\BIC{$\varphi_{38} \to \mathrm{ContIn}(f,A,c) \to B \subseteq A \to c \in B \to \mathrm{RFn}(f)$}
\BIC{$\mathrm{ContIn}(f,A,c) \to B \subseteq A \to c \in B \to \mathrm{RFn}(f)$}
\AXC{$\stackrel{\Pi_{22}}{\vdots}$}\noLine
\UIC{$\varphi_{39}$}
\AXC{$\varphi_{39} \to \varphi_{56}$\lbl{A1}}
\BIC{$\varphi_{56}$}
\BIC{$\varphi_{55}$}
\end{ptree}
\noindent$\Pi_{24}$:
\begin{ptree}
\AXC{$\mathrm{ContIn}(f,A,c) \to B \subseteq A \to c \in B \to c \in B$\lbl{A1}}
\AXC{$\stackrel{\Pi_{20}}{\vdots}$}\noLine
\UIC{$\mathrm{ContIn}(f,A,c) \to B \subseteq A \to c \in B \to B \subseteq \mathrm{dom}(f)$}
\AXC{$\stackrel{\Pi_{23}}{\vdots}$}\noLine
\UIC{$\varphi_{55}$}
\BIC{$\varphi_{54}$}
\BIC{$\varphi_{30} \to \varphi_{1}$}
\end{ptree}
\noindent 主木:
\begin{ptree}
\AXC{$\varphi_{13}$\lbl{A4}}
\AXC{$\forall \varepsilon\,(\varphi_{28}) \to \varphi_{32}$\lbl{A3}}
\AXC{$\stackrel{\Pi_{11}}{\vdots}$}\noLine
\UIC{$(\forall \varepsilon\,(\varphi_{28}) \to \varphi_{32}) \to \varphi_{13} \to \varphi_{30}$}
\BIC{$\varphi_{13} \to \varphi_{30}$}
\BIC{$\varphi_{30}$}
\AXC{$\stackrel{\Pi_{24}}{\vdots}$}\noLine
\UIC{$\varphi_{30} \to \varphi_{1}$}
\BIC{$\varphi_{1}$}
\UIC{$\forall c\,(\varphi_{1})$}
\UIC{$\forall B\,\forall c\,(\varphi_{1})$}
\UIC{$\forall A\,\forall B\,\forall c\,(\varphi_{1})$}
\UIC{$\forall f\,\forall A\,\forall B\,\forall c\,(\varphi_{1})$}
\end{ptree}
\begin{abbrevlist}
\item $\varphi_{1} :\equiv \mathrm{ContIn}(f,A,c) \to B \subseteq A \to c \in B \to \mathrm{ContIn}(f,B,c)$
\item $\varphi_{2} :\equiv \mathrm{ContIn}(f,A,c) \to B \subseteq A \to c \in B \to u \in B \to u \in \mathrm{dom}(f)$
\item $\varphi_{3} :\equiv \mathrm{ContIn}(f,A,c) \land B \subseteq A \land c \in B \to u \in B \to u \in \mathrm{dom}(f)$
\item $\varphi_{4} :\equiv \mathrm{ContIn}(f,A,c) \to B \subseteq A \to c \in B \to \forall u \in B\,(u \in \mathrm{dom}(f))$
\item $\varphi_{5} :\equiv \mathrm{ContIn}(f,A,c) \to B \subseteq A \to c \in B \to \varepsilon \in \mathbb{R} \to \varepsilon \in \mathbb{R}$
\item $\varphi_{6} :\equiv \forall u\,(\mathrm{ContIn}(f,A,c) \land B \subseteq A \land c \in B) \to \forall u \in B\,(u \in \mathrm{dom}(f))$
\item $\varphi_{7} :\equiv \forall Y\,\forall x\,(B \subseteq Y \to x \in B \to x \in Y) \to \forall x\,(B \subseteq A \to x \in B \to x \in A)$
\item $\varphi_{8} :\equiv \mathrm{ContIn}(f,A,c) \to B \subseteq A \to c \in B \to \varepsilon \in \mathbb{R} \to 0 < \varepsilon \to 0 < \varepsilon$
\item $\varphi_{9} :\equiv \forall X\,\forall Y\,\forall x\,(X \subseteq Y \to x \in X \to x \in Y) \to \forall Y\,\forall x\,(B \subseteq Y \to x \in B \to x \in Y)$
\item $\varphi_{10} :\equiv \mathrm{ContIn}(f,A,c) \land B \subseteq A \land c \in B \to \forall u\,(\mathrm{ContIn}(f,A,c) \land B \subseteq A \land c \in B)$
\item $\varphi_{11} :\equiv \forall X\,\forall Y\,\forall x\,(X \subseteq Y \to x \in X \to x \in Y) \to \forall Y\,\forall x\,(A \subseteq Y \to x \in A \to x \in Y)$
\item $\varphi_{12} :\equiv \forall x\,\forall y\,(x \subseteq y \leftrightarrow \forall u \in x\,(u \in y)) \to \forall y\,(B \subseteq y \leftrightarrow \forall u \in B\,(u \in y))$
\item $\varphi_{13} :\equiv \mathrm{ContIn}(f,A,c) \land B \subseteq A \land c \in B \to \forall \varepsilon\,(\mathrm{ContIn}(f,A,c) \land B \subseteq A \land c \in B)$
\item $\varphi_{14} :\equiv \forall Y\,\forall x\,(A \subseteq Y \to x \in A \to x \in Y) \to \forall x\,(A \subseteq \mathrm{dom}(f) \to x \in A \to x \in \mathrm{dom}(f))$
\item $\varphi_{15} :\equiv \forall y\,(B \subseteq y \leftrightarrow \forall u \in B\,(u \in y)) \to (B \subseteq \mathrm{dom}(f) \leftrightarrow \forall u \in B\,(u \in \mathrm{dom}(f)))$
\item $\varphi_{16} :\equiv \forall x\,(A \subseteq \mathrm{dom}(f) \to x \in A \to x \in \mathrm{dom}(f)) \to A \subseteq \mathrm{dom}(f) \to u \in A \to u \in \mathrm{dom}(f)$
\item $\varphi_{17} :\equiv \lvert x - c \rvert < \delta \to x \in A \to \lvert x - c \rvert < \delta \to \lvert f(x) - f(c) \rvert < \varepsilon$
\item $\varphi_{18} :\equiv \lvert x - c \rvert < \delta \to \forall x \in A\,(\lvert x - c \rvert < \delta \to \lvert f(x) - f(c) \rvert < \varepsilon)$
\item $\varphi_{19} :\equiv \delta \in \mathbb{R} \land 0 < \delta \land \forall x \in A\,(\lvert x - c \rvert < \delta \to \lvert f(x) - f(c) \rvert < \varepsilon)$
\item $\varphi_{20} :\equiv \delta \in \mathbb{R} \land 0 < \delta \land \forall x \in B\,(\lvert x - c \rvert < \delta \to \lvert f(x) - f(c) \rvert < \varepsilon)$
\item $\varphi_{21} :\equiv (\mathrm{ContIn}(f,A,c) \to B \subseteq A \to c \in B \to u \in B \to u \in B) \to \mathrm{ContIn}(f,A,c) \to B \subseteq A \to c \in B \to u \in B \to u \in A$
\item $\varphi_{22} :\equiv \varphi_{4} \to \mathrm{ContIn}(f,A,c) \to B \subseteq A \to c \in B \to B \subseteq \mathrm{dom}(f)$
\item $\varphi_{23} :\equiv (\mathrm{ContIn}(f,A,c) \to B \subseteq A \to c \in B \to u \in B \to u \in A) \to \varphi_{2}$
\item $\varphi_{24} :\equiv \forall x \in A\,(\lvert x - c \rvert < \delta \to \lvert f(x) - f(c) \rvert < \varepsilon) \to x \in A \to \lvert x - c \rvert < \delta \to \lvert f(x) - f(c) \rvert < \varepsilon$
\item $\varphi_{25} :\equiv \mathrm{ContIn}(f,A,c) \land B \subseteq A \land c \in B \land \varepsilon \in \mathbb{R} \land 0 < \varepsilon \land \varphi_{19}$
\item $\varphi_{26} :\equiv \mathrm{ContIn}(f,A,c) \to B \subseteq A \to c \in B \to \varepsilon \in \mathbb{R} \to 0 < \varepsilon \to \exists \delta\,(\varphi_{19})$
\item $\varphi_{27} :\equiv \mathrm{ContIn}(f,A,c) \to B \subseteq A \to c \in B \to \varepsilon \in \mathbb{R} \to 0 < \varepsilon \to \exists \delta\,(\varphi_{20})$
\item $\varphi_{28} :\equiv \mathrm{ContIn}(f,A,c) \land B \subseteq A \land c \in B \to \varepsilon \in \mathbb{R} \to 0 < \varepsilon \to \exists \delta\,(\varphi_{20})$
\item $\varphi_{29} :\equiv \mathrm{ContIn}(f,A,c) \land B \subseteq A \land c \in B \land \varepsilon \in \mathbb{R} \land 0 < \varepsilon \land \neg \exists \delta\,(\varphi_{20})$
\item $\varphi_{30} :\equiv \mathrm{ContIn}(f,A,c) \to B \subseteq A \to c \in B \to \forall \varepsilon \in \mathbb{R}\,(0 < \varepsilon \to \exists \delta\,(\varphi_{20}))$
\item $\varphi_{31} :\equiv 0 < \varepsilon \to \varphi_{19} \to x \in B \to \lvert x - c \rvert < \delta \to \lvert f(x) - f(c) \rvert < \varepsilon$
\item $\varphi_{32} :\equiv \forall \varepsilon\,(\mathrm{ContIn}(f,A,c) \land B \subseteq A \land c \in B) \to \forall \varepsilon \in \mathbb{R}\,(0 < \varepsilon \to \exists \delta\,(\varphi_{20}))$
\item $\varphi_{33} :\equiv B \subseteq A \to c \in B \to \varepsilon \in \mathbb{R} \to 0 < \varepsilon \to \varepsilon \in \mathbb{R} \to 0 < \varepsilon \to \exists \delta\,(\varphi_{19})$
\item $\varphi_{34} :\equiv 0 < \varepsilon \to \varphi_{19} \to \forall x \in B\,(\lvert x - c \rvert < \delta \to \lvert f(x) - f(c) \rvert < \varepsilon)$
\item $\varphi_{35} :\equiv (\mathrm{ContIn}(f,A,c) \to B \subseteq A \to c \in B \to A \subseteq \mathrm{dom}(f)) \to \varphi_{23}$
\item $\varphi_{36} :\equiv B \subseteq A \to c \in B \to \varepsilon \in \mathbb{R} \to 0 < \varepsilon \to \forall \varepsilon \in \mathbb{R}\,(0 < \varepsilon \to \exists \delta\,(\varphi_{19}))$
\item $\varphi_{37} :\equiv \mathrm{ContIn}(f,A,c) \to B \subseteq A \to c \in B \to \varepsilon \in \mathbb{R} \to 0 < \varepsilon \to \varphi_{19} \to x \in B \to x \in B$
\item $\varphi_{38} :\equiv \mathrm{ContIn}(f,A,c) \leftrightarrow \mathrm{RFn}(f) \land A \subseteq \mathrm{dom}(f) \land c \in A \land \forall \varepsilon \in \mathbb{R}\,(0 < \varepsilon \to \exists \delta\,(\varphi_{19}))$
\item $\varphi_{39} :\equiv \mathrm{ContIn}(f,B,c) \leftrightarrow \mathrm{RFn}(f) \land B \subseteq \mathrm{dom}(f) \land c \in B \land \forall \varepsilon \in \mathbb{R}\,(0 < \varepsilon \to \exists \delta\,(\varphi_{20}))$
\item $\varphi_{40} :\equiv c \in B \to \varepsilon \in \mathbb{R} \to 0 < \varepsilon \to \varphi_{19} \to x \in B \to \lvert x - c \rvert < \delta \to \lvert x - c \rvert < \delta$
\item $\varphi_{41} :\equiv B \subseteq A \to c \in B \to \varepsilon \in \mathbb{R} \to 0 < \varepsilon \to \varphi_{19} \to x \in B \to \lvert x - c \rvert < \delta \to x \in A$
\item $\varphi_{42} :\equiv (B \subseteq A \to u \in B \to u \in A) \to (\mathrm{ContIn}(f,A,c) \to B \subseteq A \to B \subseteq A) \to \varphi_{21}$
\item $\varphi_{43} :\equiv \forall x\,(\varphi_{25} \to x \in B \to \lvert x - c \rvert < \delta \to \lvert f(x) - f(c) \rvert < \varepsilon)$
\item $\varphi_{44} :\equiv \forall x\,(\varphi_{25}) \to \forall x \in B\,(\lvert x - c \rvert < \delta \to \lvert f(x) - f(c) \rvert < \varepsilon)$
\item $\varphi_{45} :\equiv \varphi_{38} \to \mathrm{ContIn}(f,A,c) \to B \subseteq A \to c \in B \to A \subseteq \mathrm{dom}(f)$
\item $\varphi_{46} :\equiv \mathrm{ContIn}(f,A,c) \to B \subseteq A \to c \in B \to \varepsilon \in \mathbb{R} \to 0 < \varepsilon \to \varphi_{19} \to \varphi_{19}$
\item $\varphi_{47} :\equiv \mathrm{ContIn}(f,A,c) \to B \subseteq A \to c \in B \to \varepsilon \in \mathbb{R} \to 0 < \varepsilon \to \varphi_{19} \to \varphi_{20}$
\item $\varphi_{48} :\equiv \forall \varepsilon \in \mathbb{R}\,(0 < \varepsilon \to \exists \delta\,(\varphi_{19})) \to \varepsilon \in \mathbb{R} \to 0 < \varepsilon \to \exists \delta\,(\varphi_{19})$
\item $\varphi_{49} :\equiv \mathrm{ContIn}(f,A,c) \to B \subseteq A \to c \in B \to \varepsilon \in \mathbb{R} \to 0 < \varepsilon \to \varphi_{19} \to x \in B \to \varphi_{17}$
\item $\varphi_{50} :\equiv \mathrm{ContIn}(f,A,c) \to B \subseteq A \to c \in B \to \varepsilon \in \mathbb{R} \to 0 < \varepsilon \to \varphi_{19} \to \exists \delta\,(\varphi_{20})$
\item $\varphi_{51} :\equiv \mathrm{ContIn}(f,A,c) \to B \subseteq A \to c \in B \to \varepsilon \in \mathbb{R} \to 0 < \varepsilon \to \varphi_{19} \to x \in B \to \varphi_{18}$
\item $\varphi_{52} :\equiv (\mathrm{ContIn}(f,A,c) \to \mathrm{ContIn}(f,A,c)) \to \varphi_{38} \to \mathrm{ContIn}(f,A,c) \to B \subseteq A \to c \in B \to \mathrm{RFn}(f)$
\item $\varphi_{53} :\equiv (\mathrm{ContIn}(f,A,c) \to \mathrm{ContIn}(f,A,c)) \to \varphi_{45}$
\item $\varphi_{54} :\equiv (\mathrm{ContIn}(f,A,c) \to B \subseteq A \to c \in B \to c \in B) \to \varphi_{30} \to \varphi_{1}$
\item $\varphi_{55} :\equiv (\mathrm{ContIn}(f,A,c) \to B \subseteq A \to c \in B \to B \subseteq \mathrm{dom}(f)) \to \varphi_{54}$
\item $\varphi_{56} :\equiv (\mathrm{ContIn}(f,A,c) \to B \subseteq A \to c \in B \to \mathrm{RFn}(f)) \to \varphi_{55}$
\item $\varphi_{57} :\equiv (\mathrm{ContIn}(f,A,c) \to \mathrm{ContIn}(f,A,c)) \to \varphi_{38} \to \mathrm{ContIn}(f,A,c) \to \varphi_{36}$
\item $\varphi_{58} :\equiv (\mathrm{ContIn}(f,A,c) \to B \subseteq A \to B \subseteq A) \to \varphi_{37} \to \mathrm{ContIn}(f,A,c) \to \varphi_{41}$
\item $\varphi_{59} :\equiv (\mathrm{ContIn}(f,A,c) \to B \subseteq A \to \varphi_{40}) \to \mathrm{ContIn}(f,A,c) \to B \subseteq A \to c \in B \to \varepsilon \in \mathbb{R} \to \varphi_{31}$
\item $\varphi_{60} :\equiv (\mathrm{ContIn}(f,A,c) \to B \subseteq A \to c \in B \to \varepsilon \in \mathbb{R} \to \varphi_{31}) \to \varphi_{25} \to x \in B \to \lvert x - c \rvert < \delta \to \lvert f(x) - f(c) \rvert < \varepsilon$
\item $\varphi_{61} :\equiv (\mathrm{ContIn}(f,A,c) \to B \subseteq A \to c \in B \to \varepsilon \in \mathbb{R} \to \varphi_{34}) \to \varphi_{47}$
\item $\varphi_{62} :\equiv (\varphi_{25} \to \forall x\,(\varphi_{25})) \to \mathrm{ContIn}(f,A,c) \to B \subseteq A \to c \in B \to \varepsilon \in \mathbb{R} \to \varphi_{34}$
\item $\varphi_{63} :\equiv (\mathrm{ContIn}(f,A,c) \to \varphi_{36}) \to \varphi_{48} \to \mathrm{ContIn}(f,A,c) \to \varphi_{33}$
\item $\varphi_{64} :\equiv (\forall \delta\,(\varphi_{29} \to \neg \varphi_{19}) \to \forall \delta\,(\varphi_{29}) \to \forall \delta\,(\neg \varphi_{19})) \to (\varphi_{29} \to \forall \delta\,(\varphi_{29})) \to \varphi_{27}$
\end{abbrevlist}
\end{fullproof}

% 自動生成 (prover/build.js): ContInAdd  ノード数 567
\begin{fullproof}{定理2.3(和)}{$\forall f\,\forall g\,\forall A\,\forall c\,(\mathrm{ContIn}(f,A,c) \to \mathrm{ContIn}(g,A,c) \to \mathrm{ContIn}((f + g),A,c))$}{567}\phantomsection\label{tree:ContInAdd}
\noindent$\Pi_{1}$:
\begin{ptree}
\AXC{$\varphi_{7}$\lbl{A1}}
\AXC{$\varphi_{118}$\lbl{[G1]}}
\AXC{$\varphi_{209}$\lbl{A2}}
\BIC{$\forall A\,\forall c\,\forall a\,(\mathrm{ContIn}(g,A,c) \to \mathrm{Seq}(a) \to \varphi_{39})$}
\AXC{$\varphi_{207}$\lbl{A2}}
\BIC{$\forall c\,\forall a\,(\mathrm{ContIn}(g,A,c) \to \mathrm{Seq}(a) \to \varphi_{39})$}
\AXC{$\varphi_{204}$\lbl{A2}}
\BIC{$\forall a\,(\mathrm{ContIn}(g,A,c) \to \mathrm{Seq}(a) \to \varphi_{39})$}
\AXC{$\varphi_{201}$\lbl{A2}}
\BIC{$\mathrm{ContIn}(g,A,c) \to \mathrm{Seq}(a) \to \varphi_{39}$}
\AXC{$\varphi_{272}$\lbl{A1}}
\BIC{$\varphi_{7} \to \varphi_{11} \to (\mathrm{ContIn}(f,A,c) \to \varphi_{78}) \to \varphi_{91} \to \varphi_{123}$}
\BIC{$\varphi_{11} \to (\mathrm{ContIn}(f,A,c) \to \varphi_{78}) \to \varphi_{91} \to \varphi_{123}$}
\end{ptree}
\noindent$\Pi_{2}$:
\begin{ptree}
\AXC{$\varphi_{91}$\lbl{A1}}
\AXC{$\mathrm{ContIn}(f,A,c) \to \varphi_{78}$\lbl{A1}}
\AXC{$\varphi_{11}$\lbl{A1}}
\AXC{$\stackrel{\Pi_{1}}{\vdots}$}\noLine
\UIC{$\varphi_{11} \to (\mathrm{ContIn}(f,A,c) \to \varphi_{78}) \to \varphi_{91} \to \varphi_{123}$}
\BIC{$(\mathrm{ContIn}(f,A,c) \to \varphi_{78}) \to \varphi_{91} \to \varphi_{123}$}
\BIC{$\varphi_{91} \to \varphi_{123}$}
\BIC{$\varphi_{123}$}
\end{ptree}
\noindent$\Pi_{3}$:
\begin{ptree}
\AXC{$\varphi_{118}$\lbl{[G1]}}
\AXC{$\varphi_{210}$\lbl{A2}}
\BIC{$\forall A\,\forall c\,\forall a\,(\mathrm{ContIn}(f,A,c) \to \mathrm{Seq}(a) \to \varphi_{40})$}
\AXC{$\varphi_{208}$\lbl{A2}}
\BIC{$\forall c\,\forall a\,(\mathrm{ContIn}(f,A,c) \to \mathrm{Seq}(a) \to \varphi_{40})$}
\AXC{$\varphi_{205}$\lbl{A2}}
\BIC{$\forall a\,(\mathrm{ContIn}(f,A,c) \to \mathrm{Seq}(a) \to \varphi_{40})$}
\AXC{$\varphi_{202}$\lbl{A2}}
\BIC{$\mathrm{ContIn}(f,A,c) \to \mathrm{Seq}(a) \to \varphi_{40}$}
\end{ptree}
\noindent$\Pi_{4}$:
\begin{ptree}
\AXC{$\stackrel{\Pi_{3}}{\vdots}$}\noLine
\UIC{$\mathrm{ContIn}(f,A,c) \to \mathrm{Seq}(a) \to \varphi_{40}$}
\AXC{$(\mathrm{ContIn}(f,A,c) \to \mathrm{Seq}(a) \to \varphi_{40}) \to \varphi_{256}$\lbl{A1}}
\BIC{$\varphi_{256}$}
\end{ptree}
\noindent$\Pi_{5}$:
\begin{ptree}
\AXC{$\varphi_{91}$\lbl{A1}}
\AXC{$\mathrm{ContIn}(f,A,c) \to \varphi_{78}$\lbl{A1}}
\AXC{$\varphi_{11}$\lbl{A1}}
\AXC{$\mathrm{ContIn}(f,A,c) \to \mathrm{ContIn}(f,A,c)$\lbl{A1}}
\AXC{$\stackrel{\Pi_{4}}{\vdots}$}\noLine
\UIC{$\varphi_{256}$}
\BIC{$\varphi_{11} \to (\mathrm{ContIn}(f,A,c) \to \varphi_{78}) \to \varphi_{91} \to \varphi_{124}$}
\BIC{$(\mathrm{ContIn}(f,A,c) \to \varphi_{78}) \to \varphi_{91} \to \varphi_{124}$}
\BIC{$\varphi_{91} \to \varphi_{124}$}
\BIC{$\varphi_{124}$}
\end{ptree}
\noindent$\Pi_{6}$:
\begin{ptree}
\AXC{$\varphi_{29}$\lbl{定理2.1(1)}}
\AXC{$\varphi_{29} \to \varphi_{75}$\lbl{A2}}
\BIC{$\varphi_{75}$}
\AXC{$\varphi_{75} \to \varphi_{88}$\lbl{A2}}
\BIC{$\varphi_{88}$}
\AXC{$\varphi_{88} \to \varphi_{121}$\lbl{A2}}
\BIC{$\varphi_{121}$}
\AXC{$\varphi_{121} \to \mathrm{Conv}(f \circ a,f(c)) \to \varphi_{83}$\lbl{A2}}
\BIC{$\mathrm{Conv}(f \circ a,f(c)) \to \varphi_{83}$}
\AXC{$\varphi_{268}$\lbl{A1}}
\BIC{$\varphi_{124} \to \varphi_{123} \to \mathrm{ContIn}(f,A,c) \to \varphi_{141}$}
\end{ptree}
\noindent$\Pi_{7}$:
\begin{ptree}
\AXC{$x = x$\lbl{A5}}
\AXC{$x = y \to x = x \to y = x$\lbl{A6}}
\AXC{$(x = y \to x = x \to y = x) \to x = x \to x = y \to y = x$\lbl{A1}}
\BIC{$x = x \to x = y \to y = x$}
\BIC{$x = y \to y = x$}
\UIC{$\forall y\,(x = y \to y = x)$}
\UIC{$\forall x\,\forall y\,(x = y \to y = x)$}
\end{ptree}
\noindent$\Pi_{8}$:
\begin{ptree}
\AXC{$y = z \to x = y \to x = z$\lbl{A6}}
\AXC{$(y = z \to x = y \to x = z) \to x = y \to y = z \to x = z$\lbl{A1}}
\BIC{$x = y \to y = z \to x = z$}
\UIC{$\forall z\,(x = y \to y = z \to x = z)$}
\UIC{$\forall y\,\forall z\,(x = y \to y = z \to x = z)$}
\UIC{$\forall x\,\forall y\,\forall z\,(x = y \to y = z \to x = z)$}
\AXC{$\forall x\,\forall y\,\forall z\,(x = y \to y = z \to x = z) \to \varphi_{45}$\lbl{A2}}
\BIC{$\varphi_{45}$}
\end{ptree}
\noindent$\Pi_{9}$:
\begin{ptree}
\AXC{$\stackrel{\Pi_{8}}{\vdots}$}\noLine
\UIC{$\varphi_{45}$}
\AXC{$\varphi_{45} \to \varphi_{151}$\lbl{A2}}
\BIC{$\varphi_{151}$}
\AXC{$\varphi_{151} \to (f + g) \circ a(n) = f(a_{n}) + g(a_{n}) \to \varphi_{163}$\lbl{A2}}
\BIC{$(f + g) \circ a(n) = f(a_{n}) + g(a_{n}) \to \varphi_{163}$}
\end{ptree}
\noindent$\Pi_{10}$:
\begin{ptree}
\AXC{$y = z \to x = y \to x = z$\lbl{A6}}
\AXC{$(y = z \to x = y \to x = z) \to x = y \to y = z \to x = z$\lbl{A1}}
\BIC{$x = y \to y = z \to x = z$}
\UIC{$\forall z\,(x = y \to y = z \to x = z)$}
\UIC{$\forall y\,\forall z\,(x = y \to y = z \to x = z)$}
\UIC{$\forall x\,\forall y\,\forall z\,(x = y \to y = z \to x = z)$}
\AXC{$\forall x\,\forall y\,\forall z\,(x = y \to y = z \to x = z) \to \varphi_{100}$\lbl{A2}}
\BIC{$\varphi_{100}$}
\end{ptree}
\noindent$\Pi_{11}$:
\begin{ptree}
\AXC{$\stackrel{\Pi_{10}}{\vdots}$}\noLine
\UIC{$\varphi_{100}$}
\AXC{$\varphi_{100} \to \varphi_{186}$\lbl{A2}}
\BIC{$\varphi_{186}$}
\AXC{$\varphi_{186} \to f \circ a \oplus g \circ a(n) = f(a_{n}) + g \circ a(n) \to \varphi_{168}$\lbl{A2}}
\BIC{$f \circ a \oplus g \circ a(n) = f(a_{n}) + g \circ a(n) \to \varphi_{168}$}
\end{ptree}
\noindent$\Pi_{12}$:
\begin{ptree}
\AXC{$\forall n \in \mathbb{N}\,(a_{n} \in A) \to n \in \mathbb{N} \to a_{n} \in A$\lbl{A2}}
\AXC{$\mathrm{ContIn}(f,A,c) \to \varphi_{78}$\lbl{A1}}
\AXC{$(\mathrm{ContIn}(f,A,c) \to \varphi_{78}) \to \varphi_{177}$\lbl{A1}}
\BIC{$\varphi_{177}$}
\BIC{$\mathrm{ContIn}(f,A,c) \to \varphi_{73}$}
\end{ptree}
\noindent$\Pi_{13}$:
\begin{ptree}
\AXC{$\forall f\,\forall A\,\forall c\,(\varphi_{191})$\lbl{DefContIn}}
\AXC{$\forall f\,\forall A\,\forall c\,(\varphi_{191}) \to \forall A\,\forall c\,(\varphi_{190})$\lbl{A2}}
\BIC{$\forall A\,\forall c\,(\varphi_{190})$}
\AXC{$\forall A\,\forall c\,(\varphi_{190}) \to \forall c\,(\varphi_{190})$\lbl{A2}}
\BIC{$\forall c\,(\varphi_{190})$}
\AXC{$\forall c\,(\varphi_{190}) \to \varphi_{190}$\lbl{A2}}
\BIC{$\varphi_{190}$}
\end{ptree}
\noindent$\Pi_{14}$:
\begin{ptree}
\AXC{$\stackrel{\Pi_{13}}{\vdots}$}\noLine
\UIC{$\varphi_{190}$}
\AXC{$\varphi_{7}$\lbl{A1}}
\AXC{$\varphi_{7} \to \varphi_{190} \to \varphi_{8}$\lbl{A1}}
\BIC{$\varphi_{190} \to \varphi_{8}$}
\BIC{$\varphi_{8}$}
\end{ptree}
\noindent$\Pi_{15}$:
\begin{ptree}
\AXC{$\stackrel{\Pi_{14}}{\vdots}$}\noLine
\UIC{$\varphi_{8}$}
\AXC{$\forall X\,\forall Y\,\forall x\,(X \subseteq Y \to x \in X \to x \in Y)$\lbl{[B5]}}
\AXC{$\varphi_{46}$\lbl{A2}}
\BIC{$\forall Y\,\forall x\,(A \subseteq Y \to x \in A \to x \in Y)$}
\AXC{$\varphi_{57}$\lbl{A2}}
\BIC{$\forall x\,(A \subseteq \mathrm{dom}(g) \to x \in A \to x \in \mathrm{dom}(g))$}
\AXC{$\varphi_{98}$\lbl{A2}}
\BIC{$A \subseteq \mathrm{dom}(g) \to a_{n} \in A \to a_{n} \in \mathrm{dom}(g)$}
\AXC{$\varphi_{245}$\lbl{A1}}
\BIC{$\varphi_{8} \to \varphi_{122} \to \varphi_{128}$}
\BIC{$\varphi_{122} \to \varphi_{128}$}
\end{ptree}
\noindent$\Pi_{16}$:
\begin{ptree}
\AXC{$\varphi_{111}$\lbl{A1}}
\AXC{$\stackrel{\Pi_{12}}{\vdots}$}\noLine
\UIC{$\mathrm{ContIn}(f,A,c) \to \varphi_{73}$}
\AXC{$(\mathrm{ContIn}(f,A,c) \to \varphi_{73}) \to \varphi_{111} \to \varphi_{122}$\lbl{A1}}
\BIC{$\varphi_{111} \to \varphi_{122}$}
\BIC{$\varphi_{122}$}
\AXC{$\stackrel{\Pi_{15}}{\vdots}$}\noLine
\UIC{$\varphi_{122} \to \varphi_{128}$}
\BIC{$\varphi_{128}$}
\end{ptree}
\noindent$\Pi_{17}$:
\begin{ptree}
\AXC{$\varphi_{69} \to \forall n\,(\varphi_{69})$\lbl{A4}}
\AXC{$\varphi_{219}$\lbl{A3}}
\AXC{$\stackrel{\Pi_{16}}{\vdots}$}\noLine
\UIC{$\varphi_{128}$}
\AXC{$\varphi_{128} \to \varphi_{69} \to n \in \mathbb{N} \to a_{n} \in \mathrm{dom}(g)$\lbl{A1}}
\BIC{$\varphi_{69} \to n \in \mathbb{N} \to a_{n} \in \mathrm{dom}(g)$}
\UIC{$\forall n\,(\varphi_{69} \to n \in \mathbb{N} \to a_{n} \in \mathrm{dom}(g))$}
\AXC{$\varphi_{278}$\lbl{A1}}
\BIC{$\varphi_{219} \to (\varphi_{69} \to \forall n\,(\varphi_{69})) \to \varphi_{131}$}
\BIC{$(\varphi_{69} \to \forall n\,(\varphi_{69})) \to \varphi_{131}$}
\BIC{$\varphi_{131}$}
\end{ptree}
\noindent$\Pi_{18}$:
\begin{ptree}
\AXC{$\forall f\,\forall A\,\forall c\,(\varphi_{191})$\lbl{DefContIn}}
\AXC{$\forall f\,\forall A\,\forall c\,(\varphi_{191}) \to \forall A\,\forall c\,(\varphi_{190})$\lbl{A2}}
\BIC{$\forall A\,\forall c\,(\varphi_{190})$}
\AXC{$\forall A\,\forall c\,(\varphi_{190}) \to \forall c\,(\varphi_{190})$\lbl{A2}}
\BIC{$\forall c\,(\varphi_{190})$}
\AXC{$\forall c\,(\varphi_{190}) \to \varphi_{190}$\lbl{A2}}
\BIC{$\varphi_{190}$}
\end{ptree}
\noindent$\Pi_{19}$:
\begin{ptree}
\AXC{$\stackrel{\Pi_{18}}{\vdots}$}\noLine
\UIC{$\varphi_{190}$}
\AXC{$\varphi_{7}$\lbl{A1}}
\AXC{$\varphi_{7} \to \varphi_{190} \to \varphi_{4}$\lbl{A1}}
\BIC{$\varphi_{190} \to \varphi_{4}$}
\BIC{$\varphi_{4}$}
\end{ptree}
\noindent$\Pi_{20}$:
\begin{ptree}
\AXC{$\varphi_{11}$\lbl{A1}}
\AXC{$\stackrel{\Pi_{19}}{\vdots}$}\noLine
\UIC{$\varphi_{4}$}
\AXC{$\forall f\,\forall a\,(\mathrm{RFn}(f) \to \mathrm{Seq}(a) \to \varphi_{117})$\lbl{[S8]}}
\AXC{$\varphi_{228}$\lbl{A2}}
\BIC{$\forall a\,(\mathrm{RFn}(g) \to \mathrm{Seq}(a) \to \varphi_{116})$}
\AXC{$\varphi_{224}$\lbl{A2}}
\BIC{$\mathrm{RFn}(g) \to \mathrm{Seq}(a) \to \varphi_{116}$}
\AXC{$\varphi_{269}$\lbl{A1}}
\BIC{$\varphi_{4} \to \varphi_{11} \to \varphi_{131} \to \mathrm{ContIn}(f,A,c) \to \varphi_{158}$}
\BIC{$\varphi_{11} \to \varphi_{131} \to \mathrm{ContIn}(f,A,c) \to \varphi_{158}$}
\BIC{$\varphi_{131} \to \mathrm{ContIn}(f,A,c) \to \varphi_{158}$}
\end{ptree}
\noindent$\Pi_{21}$:
\begin{ptree}
\AXC{$\varphi_{105}$\lbl{A2}}
\AXC{$\stackrel{\Pi_{17}}{\vdots}$}\noLine
\UIC{$\varphi_{131}$}
\AXC{$\stackrel{\Pi_{20}}{\vdots}$}\noLine
\UIC{$\varphi_{131} \to \mathrm{ContIn}(f,A,c) \to \varphi_{158}$}
\BIC{$\mathrm{ContIn}(f,A,c) \to \varphi_{158}$}
\AXC{$\varphi_{238}$\lbl{A1}}
\BIC{$\mathrm{ContIn}(f,A,c) \to \varphi_{144}$}
\AXC{$\varphi_{257}$\lbl{A1}}
\BIC{$\varphi_{105} \to \mathrm{ContIn}(f,A,c) \to \varphi_{142}$}
\BIC{$\mathrm{ContIn}(f,A,c) \to \varphi_{142}$}
\end{ptree}
\noindent$\Pi_{22}$:
\begin{ptree}
\AXC{$\varphi_{111}$\lbl{A1}}
\AXC{$\stackrel{\Pi_{21}}{\vdots}$}\noLine
\UIC{$\mathrm{ContIn}(f,A,c) \to \varphi_{142}$}
\AXC{$(\mathrm{ContIn}(f,A,c) \to \varphi_{142}) \to \varphi_{111} \to \varphi_{146}$\lbl{A1}}
\BIC{$\varphi_{111} \to \varphi_{146}$}
\BIC{$\varphi_{146}$}
\end{ptree}
\noindent$\Pi_{23}$:
\begin{ptree}
\AXC{$x = x$\lbl{A5}}
\UIC{$\forall x\,(x = x)$}
\AXC{$\varphi_{77}$\lbl{A2}}
\BIC{$f(a_{n}) + g \circ a(n) = f(a_{n}) + g \circ a(n)$}
\AXC{$\stackrel{\Pi_{22}}{\vdots}$}\noLine
\UIC{$\varphi_{146}$}
\AXC{$\varphi_{146} \to f(a_{n}) + g \circ a(n) = f(a_{n}) + g \circ a(n) \to \varphi_{251}$\lbl{A1}}
\BIC{$f(a_{n}) + g \circ a(n) = f(a_{n}) + g \circ a(n) \to \varphi_{251}$}
\BIC{$\varphi_{251}$}
\end{ptree}
\noindent$\Pi_{24}$:
\begin{ptree}
\AXC{$z = w \to \varphi_{133}$\lbl{A6}}
\UIC{$\forall w\,(z = w \to \varphi_{133})$}
\UIC{$\forall z\,\forall w\,(z = w \to \varphi_{133})$}
\AXC{$\forall z\,\forall w\,(z = w \to \varphi_{133}) \to \varphi_{184}$\lbl{A2}}
\BIC{$\varphi_{184}$}
\AXC{$\varphi_{184} \to g \circ a(n) = g(a_{n}) \to \varphi_{171}$\lbl{A2}}
\BIC{$g \circ a(n) = g(a_{n}) \to \varphi_{171}$}
\AXC{$\stackrel{\Pi_{23}}{\vdots}$}\noLine
\UIC{$\varphi_{251}$}
\BIC{$\mathrm{ContIn}(f,A,c) \to \varphi_{172}$}
\end{ptree}
\noindent$\Pi_{25}$:
\begin{ptree}
\AXC{$y = z \to x = y \to x = z$\lbl{A6}}
\AXC{$(y = z \to x = y \to x = z) \to x = y \to y = z \to x = z$\lbl{A1}}
\BIC{$x = y \to y = z \to x = z$}
\UIC{$\forall z\,(x = y \to y = z \to x = z)$}
\UIC{$\forall y\,\forall z\,(x = y \to y = z \to x = z)$}
\UIC{$\forall x\,\forall y\,\forall z\,(x = y \to y = z \to x = z)$}
\AXC{$\forall x\,\forall y\,\forall z\,(x = y \to y = z \to x = z) \to \varphi_{100}$\lbl{A2}}
\BIC{$\varphi_{100}$}
\end{ptree}
\noindent$\Pi_{26}$:
\begin{ptree}
\AXC{$\forall n \in \mathbb{N}\,(a_{n} \in A) \to n \in \mathbb{N} \to a_{n} \in A$\lbl{A2}}
\AXC{$\mathrm{ContIn}(f,A,c) \to \varphi_{78}$\lbl{A1}}
\AXC{$(\mathrm{ContIn}(f,A,c) \to \varphi_{78}) \to \varphi_{177}$\lbl{A1}}
\BIC{$\varphi_{177}$}
\BIC{$\mathrm{ContIn}(f,A,c) \to \varphi_{73}$}
\end{ptree}
\noindent$\Pi_{27}$:
\begin{ptree}
\AXC{$\forall f\,\forall A\,\forall c\,(\varphi_{191})$\lbl{DefContIn}}
\AXC{$\forall f\,\forall A\,\forall c\,(\varphi_{191}) \to \forall A\,\forall c\,(\varphi_{191})$\lbl{A2}}
\BIC{$\forall A\,\forall c\,(\varphi_{191})$}
\AXC{$\forall A\,\forall c\,(\varphi_{191}) \to \forall c\,(\varphi_{191})$\lbl{A2}}
\BIC{$\forall c\,(\varphi_{191})$}
\AXC{$\forall c\,(\varphi_{191}) \to \varphi_{191}$\lbl{A2}}
\BIC{$\varphi_{191}$}
\end{ptree}
\noindent$\Pi_{28}$:
\begin{ptree}
\AXC{$\stackrel{\Pi_{27}}{\vdots}$}\noLine
\UIC{$\varphi_{191}$}
\AXC{$\mathrm{ContIn}(f,A,c) \to \mathrm{ContIn}(f,A,c)$\lbl{A1}}
\AXC{$\varphi_{223}$\lbl{A1}}
\BIC{$\varphi_{191} \to \varphi_{9}$}
\BIC{$\varphi_{9}$}
\end{ptree}
\noindent$\Pi_{29}$:
\begin{ptree}
\AXC{$\stackrel{\Pi_{28}}{\vdots}$}\noLine
\UIC{$\varphi_{9}$}
\AXC{$\forall X\,\forall Y\,\forall x\,(X \subseteq Y \to x \in X \to x \in Y)$\lbl{[B5]}}
\AXC{$\varphi_{46}$\lbl{A2}}
\BIC{$\forall Y\,\forall x\,(A \subseteq Y \to x \in A \to x \in Y)$}
\AXC{$\varphi_{58}$\lbl{A2}}
\BIC{$\forall x\,(A \subseteq \mathrm{dom}(f) \to x \in A \to x \in \mathrm{dom}(f))$}
\AXC{$\varphi_{99}$\lbl{A2}}
\BIC{$A \subseteq \mathrm{dom}(f) \to a_{n} \in A \to a_{n} \in \mathrm{dom}(f)$}
\AXC{$\varphi_{246}$\lbl{A1}}
\BIC{$\varphi_{9} \to \varphi_{122} \to \varphi_{129}$}
\BIC{$\varphi_{122} \to \varphi_{129}$}
\end{ptree}
\noindent$\Pi_{30}$:
\begin{ptree}
\AXC{$\varphi_{111}$\lbl{A1}}
\AXC{$\stackrel{\Pi_{26}}{\vdots}$}\noLine
\UIC{$\mathrm{ContIn}(f,A,c) \to \varphi_{73}$}
\AXC{$(\mathrm{ContIn}(f,A,c) \to \varphi_{73}) \to \varphi_{111} \to \varphi_{122}$\lbl{A1}}
\BIC{$\varphi_{111} \to \varphi_{122}$}
\BIC{$\varphi_{122}$}
\AXC{$\stackrel{\Pi_{29}}{\vdots}$}\noLine
\UIC{$\varphi_{122} \to \varphi_{129}$}
\BIC{$\varphi_{129}$}
\end{ptree}
\noindent$\Pi_{31}$:
\begin{ptree}
\AXC{$\varphi_{69} \to \forall n\,(\varphi_{69})$\lbl{A4}}
\AXC{$\varphi_{220}$\lbl{A3}}
\AXC{$\stackrel{\Pi_{30}}{\vdots}$}\noLine
\UIC{$\varphi_{129}$}
\AXC{$\varphi_{129} \to \varphi_{69} \to n \in \mathbb{N} \to a_{n} \in \mathrm{dom}(f)$\lbl{A1}}
\BIC{$\varphi_{69} \to n \in \mathbb{N} \to a_{n} \in \mathrm{dom}(f)$}
\UIC{$\forall n\,(\varphi_{69} \to n \in \mathbb{N} \to a_{n} \in \mathrm{dom}(f))$}
\AXC{$\varphi_{279}$\lbl{A1}}
\BIC{$\varphi_{220} \to (\varphi_{69} \to \forall n\,(\varphi_{69})) \to \varphi_{132}$}
\BIC{$(\varphi_{69} \to \forall n\,(\varphi_{69})) \to \varphi_{132}$}
\BIC{$\varphi_{132}$}
\end{ptree}
\noindent$\Pi_{32}$:
\begin{ptree}
\AXC{$\forall f\,\forall A\,\forall c\,(\varphi_{191})$\lbl{DefContIn}}
\AXC{$\forall f\,\forall A\,\forall c\,(\varphi_{191}) \to \forall A\,\forall c\,(\varphi_{191})$\lbl{A2}}
\BIC{$\forall A\,\forall c\,(\varphi_{191})$}
\AXC{$\forall A\,\forall c\,(\varphi_{191}) \to \forall c\,(\varphi_{191})$\lbl{A2}}
\BIC{$\forall c\,(\varphi_{191})$}
\AXC{$\forall c\,(\varphi_{191}) \to \varphi_{191}$\lbl{A2}}
\BIC{$\varphi_{191}$}
\end{ptree}
\noindent$\Pi_{33}$:
\begin{ptree}
\AXC{$\stackrel{\Pi_{32}}{\vdots}$}\noLine
\UIC{$\varphi_{191}$}
\AXC{$\mathrm{ContIn}(f,A,c) \to \mathrm{ContIn}(f,A,c)$\lbl{A1}}
\AXC{$\varphi_{221}$\lbl{A1}}
\BIC{$\varphi_{191} \to \varphi_{5}$}
\BIC{$\varphi_{5}$}
\end{ptree}
\noindent$\Pi_{34}$:
\begin{ptree}
\AXC{$\varphi_{11}$\lbl{A1}}
\AXC{$\stackrel{\Pi_{33}}{\vdots}$}\noLine
\UIC{$\varphi_{5}$}
\AXC{$\forall f\,\forall a\,(\mathrm{RFn}(f) \to \mathrm{Seq}(a) \to \varphi_{117})$\lbl{[S8]}}
\AXC{$\varphi_{229}$\lbl{A2}}
\BIC{$\forall a\,(\mathrm{RFn}(f) \to \mathrm{Seq}(a) \to \varphi_{117})$}
\AXC{$\varphi_{225}$\lbl{A2}}
\BIC{$\mathrm{RFn}(f) \to \mathrm{Seq}(a) \to \varphi_{117}$}
\AXC{$\varphi_{270}$\lbl{A1}}
\BIC{$\varphi_{5} \to \varphi_{11} \to \varphi_{132} \to \mathrm{ContIn}(f,A,c) \to \varphi_{159}$}
\BIC{$\varphi_{11} \to \varphi_{132} \to \mathrm{ContIn}(f,A,c) \to \varphi_{159}$}
\BIC{$\varphi_{132} \to \mathrm{ContIn}(f,A,c) \to \varphi_{159}$}
\end{ptree}
\noindent$\Pi_{35}$:
\begin{ptree}
\AXC{$\varphi_{106}$\lbl{A2}}
\AXC{$\stackrel{\Pi_{31}}{\vdots}$}\noLine
\UIC{$\varphi_{132}$}
\AXC{$\stackrel{\Pi_{34}}{\vdots}$}\noLine
\UIC{$\varphi_{132} \to \mathrm{ContIn}(f,A,c) \to \varphi_{159}$}
\BIC{$\mathrm{ContIn}(f,A,c) \to \varphi_{159}$}
\AXC{$\varphi_{239}$\lbl{A1}}
\BIC{$\mathrm{ContIn}(f,A,c) \to \varphi_{145}$}
\AXC{$\varphi_{258}$\lbl{A1}}
\BIC{$\varphi_{106} \to \mathrm{ContIn}(f,A,c) \to \varphi_{143}$}
\BIC{$\mathrm{ContIn}(f,A,c) \to \varphi_{143}$}
\end{ptree}
\noindent$\Pi_{36}$:
\begin{ptree}
\AXC{$\varphi_{111}$\lbl{A1}}
\AXC{$\stackrel{\Pi_{35}}{\vdots}$}\noLine
\UIC{$\mathrm{ContIn}(f,A,c) \to \varphi_{143}$}
\AXC{$(\mathrm{ContIn}(f,A,c) \to \varphi_{143}) \to \varphi_{111} \to \varphi_{147}$\lbl{A1}}
\BIC{$\varphi_{111} \to \varphi_{147}$}
\BIC{$\varphi_{147}$}
\AXC{$\varphi_{276}$\lbl{A1}}
\BIC{$f \circ a(n) + g \circ a(n) = f \circ a(n) + g \circ a(n) \to \varphi_{254}$}
\end{ptree}
\noindent$\Pi_{37}$:
\begin{ptree}
\AXC{$x = x$\lbl{A5}}
\UIC{$\forall x\,(x = x)$}
\AXC{$\varphi_{92}$\lbl{A2}}
\BIC{$f \circ a(n) + g \circ a(n) = f \circ a(n) + g \circ a(n)$}
\AXC{$\stackrel{\Pi_{36}}{\vdots}$}\noLine
\UIC{$f \circ a(n) + g \circ a(n) = f \circ a(n) + g \circ a(n) \to \varphi_{254}$}
\BIC{$\varphi_{254}$}
\end{ptree}
\noindent$\Pi_{38}$:
\begin{ptree}
\AXC{$z = w \to \varphi_{155}$\lbl{A6}}
\UIC{$\forall w\,(z = w \to \varphi_{155})$}
\UIC{$\forall z\,\forall w\,(z = w \to \varphi_{155})$}
\AXC{$\forall z\,\forall w\,(z = w \to \varphi_{155}) \to \varphi_{197}$\lbl{A2}}
\BIC{$\varphi_{197}$}
\AXC{$\varphi_{197} \to f \circ a(n) = f(a_{n}) \to \varphi_{188}$\lbl{A2}}
\BIC{$f \circ a(n) = f(a_{n}) \to \varphi_{188}$}
\AXC{$\stackrel{\Pi_{37}}{\vdots}$}\noLine
\UIC{$\varphi_{254}$}
\BIC{$\mathrm{ContIn}(f,A,c) \to \varphi_{179}$}
\end{ptree}
\noindent$\Pi_{39}$:
\begin{ptree}
\AXC{$\forall a\,\forall \alpha\,(\mathrm{Conv}(a,\alpha) \to \mathrm{Seq}(a))$\lbl{[C1]}}
\AXC{$\varphi_{47}$\lbl{A2}}
\BIC{$\forall \alpha\,(\mathrm{Conv}(g \circ a,\alpha) \to \mathrm{Seq}(g \circ a))$}
\AXC{$\varphi_{95}$\lbl{A2}}
\BIC{$\mathrm{Conv}(g \circ a,g(c)) \to \mathrm{Seq}(g \circ a)$}
\AXC{$(\mathrm{Conv}(g \circ a,g(c)) \to \mathrm{Seq}(g \circ a)) \to \varphi_{211}$\lbl{A1}}
\BIC{$\varphi_{211}$}
\end{ptree}
\noindent$\Pi_{40}$:
\begin{ptree}
\AXC{$\forall a\,\forall \alpha\,(\mathrm{Conv}(a,\alpha) \to \mathrm{Seq}(a))$\lbl{[C1]}}
\AXC{$\varphi_{48}$\lbl{A2}}
\BIC{$\forall \alpha\,(\mathrm{Conv}(f \circ a,\alpha) \to \mathrm{Seq}(f \circ a))$}
\AXC{$\varphi_{96}$\lbl{A2}}
\BIC{$\mathrm{Conv}(f \circ a,f(c)) \to \mathrm{Seq}(f \circ a)$}
\AXC{$(\mathrm{Conv}(f \circ a,f(c)) \to \mathrm{Seq}(f \circ a)) \to \varphi_{212}$\lbl{A1}}
\BIC{$\varphi_{212}$}
\end{ptree}
\noindent$\Pi_{41}$:
\begin{ptree}
\AXC{$\stackrel{\Pi_{5}}{\vdots}$}\noLine
\UIC{$\varphi_{124}$}
\AXC{$\stackrel{\Pi_{40}}{\vdots}$}\noLine
\UIC{$\varphi_{212}$}
\BIC{$\mathrm{ContIn}(f,A,c) \to \mathrm{ContIn}(g,A,c) \to \varphi_{38}$}
\AXC{$\varphi_{97}$\lbl{[S2]}}
\AXC{$\varphi_{97} \to \varphi_{152}$\lbl{A2}}
\BIC{$\varphi_{152}$}
\AXC{$\varphi_{152} \to \varphi_{194}$\lbl{A2}}
\BIC{$\varphi_{194}$}
\AXC{$\varphi_{194} \to \varphi_{264}$\lbl{A1}}
\BIC{$\varphi_{264}$}
\BIC{$\varphi_{243}$}
\end{ptree}
\noindent$\Pi_{42}$:
\begin{ptree}
\AXC{$\varphi_{200}$\lbl{A2}}
\AXC{$\stackrel{\Pi_{2}}{\vdots}$}\noLine
\UIC{$\varphi_{123}$}
\AXC{$\stackrel{\Pi_{39}}{\vdots}$}\noLine
\UIC{$\varphi_{211}$}
\BIC{$\mathrm{ContIn}(f,A,c) \to \mathrm{ContIn}(g,A,c) \to \varphi_{37}$}
\AXC{$\stackrel{\Pi_{41}}{\vdots}$}\noLine
\UIC{$\varphi_{243}$}
\BIC{$\varphi_{214}$}
\AXC{$\varphi_{261}$\lbl{A1}}
\BIC{$\mathrm{ContIn}(f,A,c) \to \mathrm{ContIn}(g,A,c) \to \varphi_{180}$}
\AXC{$\varphi_{275}$\lbl{A1}}
\BIC{$\varphi_{200} \to \mathrm{ContIn}(f,A,c) \to \varphi_{192}$}
\BIC{$\mathrm{ContIn}(f,A,c) \to \varphi_{192}$}
\AXC{$\varphi_{267}$\lbl{A1}}
\BIC{$\varphi_{111} \to \mathrm{ContIn}(f,A,c) \to \varphi_{181}$}
\end{ptree}
\noindent$\Pi_{43}$:
\begin{ptree}
\AXC{$\stackrel{\Pi_{38}}{\vdots}$}\noLine
\UIC{$\mathrm{ContIn}(f,A,c) \to \varphi_{179}$}
\AXC{$\varphi_{111}$\lbl{A1}}
\AXC{$\stackrel{\Pi_{42}}{\vdots}$}\noLine
\UIC{$\varphi_{111} \to \mathrm{ContIn}(f,A,c) \to \varphi_{181}$}
\BIC{$\mathrm{ContIn}(f,A,c) \to \varphi_{181}$}
\AXC{$\varphi_{284}$\lbl{A1}}
\BIC{$(\mathrm{ContIn}(f,A,c) \to \varphi_{179}) \to \varphi_{263}$}
\BIC{$\varphi_{263}$}
\end{ptree}
\noindent$\Pi_{44}$:
\begin{ptree}
\AXC{$\stackrel{\Pi_{25}}{\vdots}$}\noLine
\UIC{$\varphi_{100}$}
\AXC{$\varphi_{100} \to \varphi_{193}$\lbl{A2}}
\BIC{$\varphi_{193}$}
\AXC{$\varphi_{193} \to f \circ a \oplus g \circ a(n) = f \circ a(n) + g \circ a(n) \to \varphi_{178}$\lbl{A2}}
\BIC{$f \circ a \oplus g \circ a(n) = f \circ a(n) + g \circ a(n) \to \varphi_{178}$}
\AXC{$\stackrel{\Pi_{43}}{\vdots}$}\noLine
\UIC{$\varphi_{263}$}
\BIC{$\mathrm{ContIn}(f,A,c) \to \varphi_{176}$}
\AXC{$\varphi_{283}$\lbl{A1}}
\BIC{$(\mathrm{ContIn}(f,A,c) \to \varphi_{172}) \to \varphi_{259}$}
\end{ptree}
\noindent$\Pi_{45}$:
\begin{ptree}
\AXC{$\stackrel{\Pi_{11}}{\vdots}$}\noLine
\UIC{$f \circ a \oplus g \circ a(n) = f(a_{n}) + g \circ a(n) \to \varphi_{168}$}
\AXC{$\stackrel{\Pi_{24}}{\vdots}$}\noLine
\UIC{$\mathrm{ContIn}(f,A,c) \to \varphi_{172}$}
\AXC{$\stackrel{\Pi_{44}}{\vdots}$}\noLine
\UIC{$(\mathrm{ContIn}(f,A,c) \to \varphi_{172}) \to \varphi_{259}$}
\BIC{$\varphi_{259}$}
\BIC{$\mathrm{ContIn}(f,A,c) \to \varphi_{173}$}
\end{ptree}
\noindent$\Pi_{46}$:
\begin{ptree}
\AXC{$\stackrel{\Pi_{7}}{\vdots}$}\noLine
\UIC{$\forall x\,\forall y\,(x = y \to y = x)$}
\AXC{$\varphi_{110}$\lbl{A2}}
\BIC{$\forall y\,(f \circ a \oplus g \circ a(n) = y \to y = f \circ a \oplus g \circ a(n))$}
\AXC{$\varphi_{218}$\lbl{A2}}
\BIC{$\varphi_{169}$}
\AXC{$\stackrel{\Pi_{45}}{\vdots}$}\noLine
\UIC{$\mathrm{ContIn}(f,A,c) \to \varphi_{173}$}
\AXC{$\varphi_{271}$\lbl{A1}}
\BIC{$\varphi_{169} \to \mathrm{ContIn}(f,A,c) \to \varphi_{174}$}
\BIC{$\mathrm{ContIn}(f,A,c) \to \varphi_{174}$}
\end{ptree}
\noindent$\Pi_{47}$:
\begin{ptree}
\AXC{$y = z \to x = y \to x = z$\lbl{A6}}
\AXC{$(y = z \to x = y \to x = z) \to x = y \to y = z \to x = z$\lbl{A1}}
\BIC{$x = y \to y = z \to x = z$}
\UIC{$\forall z\,(x = y \to y = z \to x = z)$}
\UIC{$\forall y\,\forall z\,(x = y \to y = z \to x = z)$}
\UIC{$\forall x\,\forall y\,\forall z\,(x = y \to y = z \to x = z)$}
\AXC{$\forall x\,\forall y\,\forall z\,(x = y \to y = z \to x = z) \to \varphi_{45}$\lbl{A2}}
\BIC{$\varphi_{45}$}
\end{ptree}
\noindent$\Pi_{48}$:
\begin{ptree}
\AXC{$\stackrel{\Pi_{47}}{\vdots}$}\noLine
\UIC{$\varphi_{45}$}
\AXC{$\varphi_{45} \to \varphi_{126}$\lbl{A2}}
\BIC{$\varphi_{126}$}
\AXC{$\varphi_{126} \to (f + g) \circ a(n) = (f + g)(a_{n}) \to \varphi_{136}$\lbl{A2}}
\BIC{$(f + g) \circ a(n) = (f + g)(a_{n}) \to \varphi_{136}$}
\end{ptree}
\noindent$\Pi_{49}$:
\begin{ptree}
\AXC{$\stackrel{\Pi_{33}}{\vdots}$}\noLine
\UIC{$\varphi_{5}$}
\AXC{$\forall f\,(\varphi_{175})$\lbl{[S17]}}
\AXC{$\forall f\,(\varphi_{175}) \to \varphi_{175}$\lbl{A2}}
\BIC{$\varphi_{175}$}
\AXC{$\varphi_{175} \to \mathrm{RFn}(f) \to \mathrm{RFn}(g) \to \mathrm{RFn}((f + g)) \land \varphi_{16} \land \varphi_{27}$\lbl{A2}}
\BIC{$\mathrm{RFn}(f) \to \mathrm{RFn}(g) \to \mathrm{RFn}((f + g)) \land \varphi_{16} \land \varphi_{27}$}
\AXC{$\varphi_{260}$\lbl{A1}}
\BIC{$\varphi_{5} \to \varphi_{4} \to \varphi_{187}$}
\BIC{$\varphi_{4} \to \varphi_{187}$}
\end{ptree}
\noindent$\Pi_{50}$:
\begin{ptree}
\AXC{$\stackrel{\Pi_{19}}{\vdots}$}\noLine
\UIC{$\varphi_{4}$}
\AXC{$\stackrel{\Pi_{49}}{\vdots}$}\noLine
\UIC{$\varphi_{4} \to \varphi_{187}$}
\BIC{$\varphi_{187}$}
\end{ptree}
\noindent$\Pi_{51}$:
\begin{ptree}
\AXC{$\stackrel{\Pi_{50}}{\vdots}$}\noLine
\UIC{$\varphi_{187}$}
\AXC{$\varphi_{226}$\lbl{A1}}
\BIC{$\mathrm{ContIn}(f,A,c) \to \mathrm{ContIn}(g,A,c) \to \varphi_{16}$}
\end{ptree}
\noindent$\Pi_{52}$:
\begin{ptree}
\AXC{$\varphi_{16} \to \varphi_{44}$\lbl{A2}}
\AXC{$\stackrel{\Pi_{51}}{\vdots}$}\noLine
\UIC{$\mathrm{ContIn}(f,A,c) \to \mathrm{ContIn}(g,A,c) \to \varphi_{16}$}
\AXC{$\varphi_{241}$\lbl{A1}}
\BIC{$\varphi_{217}$}
\BIC{$\mathrm{ContIn}(f,A,c) \to \mathrm{ContIn}(g,A,c) \to \varphi_{44}$}
\end{ptree}
\noindent$\Pi_{53}$:
\begin{ptree}
\AXC{$\forall n \in \mathbb{N}\,(a_{n} \in A) \to n \in \mathbb{N} \to a_{n} \in A$\lbl{A2}}
\AXC{$\mathrm{ContIn}(f,A,c) \to \varphi_{78}$\lbl{A1}}
\AXC{$(\mathrm{ContIn}(f,A,c) \to \varphi_{78}) \to \varphi_{177}$\lbl{A1}}
\BIC{$\varphi_{177}$}
\BIC{$\mathrm{ContIn}(f,A,c) \to \varphi_{73}$}
\end{ptree}
\noindent$\Pi_{54}$:
\begin{ptree}
\AXC{$\varphi_{111}$\lbl{A1}}
\AXC{$\stackrel{\Pi_{53}}{\vdots}$}\noLine
\UIC{$\mathrm{ContIn}(f,A,c) \to \varphi_{73}$}
\AXC{$(\mathrm{ContIn}(f,A,c) \to \varphi_{73}) \to \varphi_{111} \to \varphi_{122}$\lbl{A1}}
\BIC{$\varphi_{111} \to \varphi_{122}$}
\BIC{$\varphi_{122}$}
\end{ptree}
\noindent$\Pi_{55}$:
\begin{ptree}
\AXC{$\stackrel{\Pi_{54}}{\vdots}$}\noLine
\UIC{$\varphi_{122}$}
\AXC{$\stackrel{\Pi_{14}}{\vdots}$}\noLine
\UIC{$\varphi_{8}$}
\AXC{$\forall X\,\forall Y\,\forall x\,(X \subseteq Y \to x \in X \to x \in Y)$\lbl{[B5]}}
\AXC{$\varphi_{46}$\lbl{A2}}
\BIC{$\forall Y\,\forall x\,(A \subseteq Y \to x \in A \to x \in Y)$}
\AXC{$\varphi_{57}$\lbl{A2}}
\BIC{$\forall x\,(A \subseteq \mathrm{dom}(g) \to x \in A \to x \in \mathrm{dom}(g))$}
\AXC{$\varphi_{98}$\lbl{A2}}
\BIC{$A \subseteq \mathrm{dom}(g) \to a_{n} \in A \to a_{n} \in \mathrm{dom}(g)$}
\AXC{$\varphi_{245}$\lbl{A1}}
\BIC{$\varphi_{8} \to \varphi_{122} \to \varphi_{128}$}
\BIC{$\varphi_{122} \to \varphi_{128}$}
\BIC{$\varphi_{128}$}
\end{ptree}
\noindent$\Pi_{56}$:
\begin{ptree}
\AXC{$\stackrel{\Pi_{54}}{\vdots}$}\noLine
\UIC{$\varphi_{122}$}
\AXC{$\stackrel{\Pi_{28}}{\vdots}$}\noLine
\UIC{$\varphi_{9}$}
\AXC{$\forall X\,\forall Y\,\forall x\,(X \subseteq Y \to x \in X \to x \in Y)$\lbl{[B5]}}
\AXC{$\varphi_{46}$\lbl{A2}}
\BIC{$\forall Y\,\forall x\,(A \subseteq Y \to x \in A \to x \in Y)$}
\AXC{$\varphi_{58}$\lbl{A2}}
\BIC{$\forall x\,(A \subseteq \mathrm{dom}(f) \to x \in A \to x \in \mathrm{dom}(f))$}
\AXC{$\varphi_{99}$\lbl{A2}}
\BIC{$A \subseteq \mathrm{dom}(f) \to a_{n} \in A \to a_{n} \in \mathrm{dom}(f)$}
\AXC{$\varphi_{246}$\lbl{A1}}
\BIC{$\varphi_{9} \to \varphi_{122} \to \varphi_{129}$}
\BIC{$\varphi_{122} \to \varphi_{129}$}
\BIC{$\varphi_{129}$}
\end{ptree}
\noindent$\Pi_{57}$:
\begin{ptree}
\AXC{$\stackrel{\Pi_{52}}{\vdots}$}\noLine
\UIC{$\mathrm{ContIn}(f,A,c) \to \mathrm{ContIn}(g,A,c) \to \varphi_{44}$}
\AXC{$\stackrel{\Pi_{55}}{\vdots}$}\noLine
\UIC{$\varphi_{128}$}
\AXC{$\stackrel{\Pi_{56}}{\vdots}$}\noLine
\UIC{$\varphi_{129}$}
\AXC{$\varphi_{129} \to \varphi_{128} \to \varphi_{216}$\lbl{A1}}
\BIC{$\varphi_{128} \to \varphi_{216}$}
\BIC{$\varphi_{216}$}
\BIC{$\varphi_{138}$}
\end{ptree}
\noindent$\Pi_{58}$:
\begin{ptree}
\AXC{$\varphi_{69} \to \forall n\,(\varphi_{69})$\lbl{A4}}
\AXC{$\varphi_{227}$\lbl{A3}}
\AXC{$\stackrel{\Pi_{57}}{\vdots}$}\noLine
\UIC{$\varphi_{138}$}
\AXC{$\varphi_{138} \to \varphi_{69} \to n \in \mathbb{N} \to a_{n} \in \mathrm{dom}((f + g))$\lbl{A1}}
\BIC{$\varphi_{69} \to n \in \mathbb{N} \to a_{n} \in \mathrm{dom}((f + g))$}
\UIC{$\forall n\,(\varphi_{69} \to n \in \mathbb{N} \to a_{n} \in \mathrm{dom}((f + g)))$}
\AXC{$\varphi_{281}$\lbl{A1}}
\BIC{$\varphi_{227} \to (\varphi_{69} \to \forall n\,(\varphi_{69})) \to \varphi_{139}$}
\BIC{$(\varphi_{69} \to \forall n\,(\varphi_{69})) \to \varphi_{139}$}
\BIC{$\varphi_{139}$}
\end{ptree}
\noindent$\Pi_{59}$:
\begin{ptree}
\AXC{$\stackrel{\Pi_{50}}{\vdots}$}\noLine
\UIC{$\varphi_{187}$}
\AXC{$\varphi_{231}$\lbl{A1}}
\BIC{$\mathrm{ContIn}(f,A,c) \to \mathrm{ContIn}(g,A,c) \to \varphi_{27}$}
\end{ptree}
\noindent$\Pi_{60}$:
\begin{ptree}
\AXC{$\varphi_{27} \to \varphi_{82}$\lbl{A2}}
\AXC{$\stackrel{\Pi_{59}}{\vdots}$}\noLine
\UIC{$\mathrm{ContIn}(f,A,c) \to \mathrm{ContIn}(g,A,c) \to \varphi_{27}$}
\AXC{$\varphi_{252}$\lbl{A1}}
\BIC{$\varphi_{233}$}
\BIC{$\mathrm{ContIn}(f,A,c) \to \mathrm{ContIn}(g,A,c) \to \varphi_{82}$}
\AXC{$\varphi_{274}$\lbl{A1}}
\BIC{$\varphi_{138} \to \varphi_{111} \to \mathrm{ContIn}(f,A,c) \to \varphi_{160}$}
\end{ptree}
\noindent$\Pi_{61}$:
\begin{ptree}
\AXC{$\varphi_{111}$\lbl{A1}}
\AXC{$\varphi_{76}$\lbl{A2}}
\AXC{$\stackrel{\Pi_{58}}{\vdots}$}\noLine
\UIC{$\varphi_{139}$}
\AXC{$\varphi_{139} \to \varphi_{76} \to \varphi_{138}$\lbl{A1}}
\BIC{$\varphi_{76} \to \varphi_{138}$}
\BIC{$\varphi_{138}$}
\AXC{$\stackrel{\Pi_{60}}{\vdots}$}\noLine
\UIC{$\varphi_{138} \to \varphi_{111} \to \mathrm{ContIn}(f,A,c) \to \varphi_{160}$}
\BIC{$\varphi_{111} \to \mathrm{ContIn}(f,A,c) \to \varphi_{160}$}
\BIC{$\mathrm{ContIn}(f,A,c) \to \varphi_{160}$}
\end{ptree}
\noindent$\Pi_{62}$:
\begin{ptree}
\AXC{$\stackrel{\Pi_{50}}{\vdots}$}\noLine
\UIC{$\varphi_{187}$}
\AXC{$\varphi_{187} \to \varphi_{10}$\lbl{A1}}
\BIC{$\varphi_{10}$}
\end{ptree}
\noindent$\Pi_{63}$:
\begin{ptree}
\AXC{$\forall f\,\forall a\,(\mathrm{RFn}(f) \to \mathrm{Seq}(a) \to \varphi_{117})$\lbl{[S8]}}
\AXC{$\forall f\,\forall a\,(\mathrm{RFn}(f) \to \mathrm{Seq}(a) \to \varphi_{117}) \to \forall a\,(\varphi_{183})$\lbl{A2}}
\BIC{$\forall a\,(\varphi_{183})$}
\AXC{$\forall a\,(\varphi_{183}) \to \varphi_{183}$\lbl{A2}}
\BIC{$\varphi_{183}$}
\end{ptree}
\noindent$\Pi_{64}$:
\begin{ptree}
\AXC{$\stackrel{\Pi_{58}}{\vdots}$}\noLine
\UIC{$\varphi_{139}$}
\AXC{$\varphi_{11}$\lbl{A1}}
\AXC{$\stackrel{\Pi_{62}}{\vdots}$}\noLine
\UIC{$\varphi_{10}$}
\AXC{$\stackrel{\Pi_{63}}{\vdots}$}\noLine
\UIC{$\varphi_{183}$}
\AXC{$\varphi_{183} \to \varphi_{10} \to \varphi_{11} \to \varphi_{139} \to \mathrm{ContIn}(f,A,c) \to \varphi_{182}$\lbl{A1}}
\BIC{$\varphi_{10} \to \varphi_{11} \to \varphi_{139} \to \mathrm{ContIn}(f,A,c) \to \varphi_{182}$}
\BIC{$\varphi_{11} \to \varphi_{139} \to \mathrm{ContIn}(f,A,c) \to \varphi_{182}$}
\BIC{$\varphi_{139} \to \mathrm{ContIn}(f,A,c) \to \varphi_{182}$}
\BIC{$\mathrm{ContIn}(f,A,c) \to \varphi_{182}$}
\end{ptree}
\noindent$\Pi_{65}$:
\begin{ptree}
\AXC{$\varphi_{111}$\lbl{A1}}
\AXC{$\varphi_{148}$\lbl{A2}}
\AXC{$\stackrel{\Pi_{64}}{\vdots}$}\noLine
\UIC{$\mathrm{ContIn}(f,A,c) \to \varphi_{182}$}
\AXC{$\varphi_{249}$\lbl{A1}}
\BIC{$\mathrm{ContIn}(f,A,c) \to \varphi_{165}$}
\AXC{$\varphi_{265}$\lbl{A1}}
\BIC{$\varphi_{148} \to \mathrm{ContIn}(f,A,c) \to \varphi_{164}$}
\BIC{$\mathrm{ContIn}(f,A,c) \to \varphi_{164}$}
\AXC{$\varphi_{262}$\lbl{A1}}
\BIC{$\varphi_{111} \to \mathrm{ContIn}(f,A,c) \to \varphi_{150}$}
\BIC{$\mathrm{ContIn}(f,A,c) \to \varphi_{150}$}
\AXC{$\varphi_{277}$\lbl{A1}}
\BIC{$(\mathrm{ContIn}(f,A,c) \to \varphi_{160}) \to \varphi_{247}$}
\end{ptree}
\noindent$\Pi_{66}$:
\begin{ptree}
\AXC{$\stackrel{\Pi_{48}}{\vdots}$}\noLine
\UIC{$(f + g) \circ a(n) = (f + g)(a_{n}) \to \varphi_{136}$}
\AXC{$\stackrel{\Pi_{61}}{\vdots}$}\noLine
\UIC{$\mathrm{ContIn}(f,A,c) \to \varphi_{160}$}
\AXC{$\stackrel{\Pi_{65}}{\vdots}$}\noLine
\UIC{$(\mathrm{ContIn}(f,A,c) \to \varphi_{160}) \to \varphi_{247}$}
\BIC{$\varphi_{247}$}
\BIC{$\mathrm{ContIn}(f,A,c) \to \varphi_{162}$}
\end{ptree}
\noindent$\Pi_{67}$:
\begin{ptree}
\AXC{$\stackrel{\Pi_{9}}{\vdots}$}\noLine
\UIC{$(f + g) \circ a(n) = f(a_{n}) + g(a_{n}) \to \varphi_{163}$}
\AXC{$\stackrel{\Pi_{46}}{\vdots}$}\noLine
\UIC{$\mathrm{ContIn}(f,A,c) \to \varphi_{174}$}
\AXC{$\stackrel{\Pi_{66}}{\vdots}$}\noLine
\UIC{$\mathrm{ContIn}(f,A,c) \to \varphi_{162}$}
\AXC{$\varphi_{282}$\lbl{A1}}
\BIC{$(\mathrm{ContIn}(f,A,c) \to \varphi_{174}) \to \varphi_{253}$}
\BIC{$\varphi_{253}$}
\BIC{$\mathrm{ContIn}(f,A,c) \to \varphi_{166}$}
\end{ptree}
\noindent$\Pi_{68}$:
\begin{ptree}
\AXC{$\stackrel{\Pi_{7}}{\vdots}$}\noLine
\UIC{$\forall x\,\forall y\,(x = y \to y = x)$}
\AXC{$\varphi_{64}$\lbl{A2}}
\BIC{$\forall y\,((f + g) \circ a(n) = y \to y = (f + g) \circ a(n))$}
\AXC{$\varphi_{206}$\lbl{A2}}
\BIC{$\varphi_{154}$}
\AXC{$\stackrel{\Pi_{67}}{\vdots}$}\noLine
\UIC{$\mathrm{ContIn}(f,A,c) \to \varphi_{166}$}
\AXC{$\varphi_{266}$\lbl{A1}}
\BIC{$\varphi_{154} \to \mathrm{ContIn}(f,A,c) \to \varphi_{167}$}
\BIC{$\mathrm{ContIn}(f,A,c) \to \varphi_{167}$}
\AXC{$\varphi_{244}$\lbl{A1}}
\BIC{$\varphi_{69} \to n \in \mathbb{N} \to f \circ a \oplus g \circ a(n) = (f + g) \circ a(n)$}
\UIC{$\varphi_{185}$}
\end{ptree}
\noindent$\Pi_{69}$:
\begin{ptree}
\AXC{$\varphi_{55}$\lbl{[C13]}}
\AXC{$\varphi_{55} \to \varphi_{153}$\lbl{A2}}
\BIC{$\varphi_{153}$}
\AXC{$\varphi_{153} \to \varphi_{199}$\lbl{A2}}
\BIC{$\varphi_{199}$}
\AXC{$\varphi_{199} \to \varphi_{215}$\lbl{A2}}
\BIC{$\varphi_{215}$}
\AXC{$\varphi_{215} \to \varphi_{120} \to (\mathrm{ContIn}(f,A,c) \to \varphi_{170}) \to \varphi_{234}$\lbl{A1}}
\BIC{$\varphi_{120} \to (\mathrm{ContIn}(f,A,c) \to \varphi_{170}) \to \varphi_{234}$}
\end{ptree}
\noindent$\Pi_{70}$:
\begin{ptree}
\AXC{$\stackrel{\Pi_{64}}{\vdots}$}\noLine
\UIC{$\mathrm{ContIn}(f,A,c) \to \varphi_{182}$}
\AXC{$(\mathrm{ContIn}(f,A,c) \to \varphi_{182}) \to \varphi_{120}$\lbl{A1}}
\BIC{$\varphi_{120}$}
\AXC{$\stackrel{\Pi_{69}}{\vdots}$}\noLine
\UIC{$\varphi_{120} \to (\mathrm{ContIn}(f,A,c) \to \varphi_{170}) \to \varphi_{234}$}
\BIC{$(\mathrm{ContIn}(f,A,c) \to \varphi_{170}) \to \varphi_{234}$}
\end{ptree}
\noindent$\Pi_{71}$:
\begin{ptree}
\AXC{$\varphi_{69} \to \forall n\,(\varphi_{69})$\lbl{A4}}
\AXC{$\varphi_{185} \to \varphi_{189}$\lbl{A3}}
\AXC{$\stackrel{\Pi_{68}}{\vdots}$}\noLine
\UIC{$\varphi_{185}$}
\AXC{$\varphi_{185} \to \varphi_{280}$\lbl{A1}}
\BIC{$\varphi_{280}$}
\BIC{$(\varphi_{69} \to \forall n\,(\varphi_{69})) \to \mathrm{ContIn}(f,A,c) \to \varphi_{170}$}
\BIC{$\mathrm{ContIn}(f,A,c) \to \varphi_{170}$}
\AXC{$\stackrel{\Pi_{70}}{\vdots}$}\noLine
\UIC{$(\mathrm{ContIn}(f,A,c) \to \varphi_{170}) \to \varphi_{234}$}
\BIC{$\varphi_{234}$}
\end{ptree}
\noindent$\Pi_{72}$:
\begin{ptree}
\AXC{$\stackrel{\Pi_{7}}{\vdots}$}\noLine
\UIC{$\forall x\,\forall y\,(x = y \to y = x)$}
\AXC{$\varphi_{23}$\lbl{A2}}
\BIC{$\forall y\,((f + g)(c) = y \to y = (f + g)(c))$}
\AXC{$\forall y\,((f + g)(c) = y \to y = (f + g)(c)) \to \varphi_{56}$\lbl{A2}}
\BIC{$\varphi_{56}$}
\end{ptree}
\noindent$\Pi_{73}$:
\begin{ptree}
\AXC{$\forall f\,\forall A\,\forall c\,(\varphi_{191})$\lbl{DefContIn}}
\AXC{$\forall f\,\forall A\,\forall c\,(\varphi_{191}) \to \forall A\,\forall c\,(\varphi_{191})$\lbl{A2}}
\BIC{$\forall A\,\forall c\,(\varphi_{191})$}
\AXC{$\forall A\,\forall c\,(\varphi_{191}) \to \forall c\,(\varphi_{191})$\lbl{A2}}
\BIC{$\forall c\,(\varphi_{191})$}
\AXC{$\forall c\,(\varphi_{191}) \to \varphi_{191}$\lbl{A2}}
\BIC{$\varphi_{191}$}
\end{ptree}
\noindent$\Pi_{74}$:
\begin{ptree}
\AXC{$\stackrel{\Pi_{73}}{\vdots}$}\noLine
\UIC{$\varphi_{191}$}
\AXC{$\mathrm{ContIn}(f,A,c) \to \mathrm{ContIn}(f,A,c)$\lbl{A1}}
\AXC{$\varphi_{222}$\lbl{A1}}
\BIC{$\varphi_{191} \to \varphi_{6}$}
\BIC{$\varphi_{6}$}
\end{ptree}
\noindent$\Pi_{75}$:
\begin{ptree}
\AXC{$\forall X\,\forall Y\,\forall x\,(X \subseteq Y \to x \in X \to x \in Y)$\lbl{[B5]}}
\AXC{$\varphi_{46}$\lbl{A2}}
\BIC{$\forall Y\,\forall x\,(A \subseteq Y \to x \in A \to x \in Y)$}
\AXC{$\varphi_{57}$\lbl{A2}}
\BIC{$\forall x\,(A \subseteq \mathrm{dom}(g) \to x \in A \to x \in \mathrm{dom}(g))$}
\AXC{$\varphi_{65}$\lbl{A2}}
\BIC{$A \subseteq \mathrm{dom}(g) \to c \in A \to c \in \mathrm{dom}(g)$}
\end{ptree}
\noindent$\Pi_{76}$:
\begin{ptree}
\AXC{$\stackrel{\Pi_{74}}{\vdots}$}\noLine
\UIC{$\varphi_{6}$}
\AXC{$\stackrel{\Pi_{14}}{\vdots}$}\noLine
\UIC{$\varphi_{8}$}
\AXC{$\stackrel{\Pi_{75}}{\vdots}$}\noLine
\UIC{$A \subseteq \mathrm{dom}(g) \to c \in A \to c \in \mathrm{dom}(g)$}
\AXC{$(A \subseteq \mathrm{dom}(g) \to c \in A \to c \in \mathrm{dom}(g)) \to \varphi_{8} \to \varphi_{156}$\lbl{A1}}
\BIC{$\varphi_{8} \to \varphi_{156}$}
\BIC{$\varphi_{156}$}
\BIC{$\mathrm{ContIn}(f,A,c) \to \mathrm{ContIn}(g,A,c) \to \varphi_{31}$}
\end{ptree}
\noindent$\Pi_{77}$:
\begin{ptree}
\AXC{$\forall X\,\forall Y\,\forall x\,(X \subseteq Y \to x \in X \to x \in Y)$\lbl{[B5]}}
\AXC{$\varphi_{46}$\lbl{A2}}
\BIC{$\forall Y\,\forall x\,(A \subseteq Y \to x \in A \to x \in Y)$}
\AXC{$\varphi_{58}$\lbl{A2}}
\BIC{$\forall x\,(A \subseteq \mathrm{dom}(f) \to x \in A \to x \in \mathrm{dom}(f))$}
\AXC{$\varphi_{66}$\lbl{A2}}
\BIC{$A \subseteq \mathrm{dom}(f) \to c \in A \to c \in \mathrm{dom}(f)$}
\end{ptree}
\noindent$\Pi_{78}$:
\begin{ptree}
\AXC{$\stackrel{\Pi_{74}}{\vdots}$}\noLine
\UIC{$\varphi_{6}$}
\AXC{$\stackrel{\Pi_{28}}{\vdots}$}\noLine
\UIC{$\varphi_{9}$}
\AXC{$\stackrel{\Pi_{77}}{\vdots}$}\noLine
\UIC{$A \subseteq \mathrm{dom}(f) \to c \in A \to c \in \mathrm{dom}(f)$}
\AXC{$(A \subseteq \mathrm{dom}(f) \to c \in A \to c \in \mathrm{dom}(f)) \to \varphi_{9} \to \varphi_{157}$\lbl{A1}}
\BIC{$\varphi_{9} \to \varphi_{157}$}
\BIC{$\varphi_{157}$}
\BIC{$\mathrm{ContIn}(f,A,c) \to \mathrm{ContIn}(g,A,c) \to \varphi_{32}$}
\end{ptree}
\noindent$\Pi_{79}$:
\begin{ptree}
\AXC{$\stackrel{\Pi_{76}}{\vdots}$}\noLine
\UIC{$\mathrm{ContIn}(f,A,c) \to \mathrm{ContIn}(g,A,c) \to \varphi_{31}$}
\AXC{$\stackrel{\Pi_{78}}{\vdots}$}\noLine
\UIC{$\mathrm{ContIn}(f,A,c) \to \mathrm{ContIn}(g,A,c) \to \varphi_{32}$}
\AXC{$\varphi_{255}$\lbl{A1}}
\BIC{$\varphi_{237}$}
\BIC{$(\mathrm{ContIn}(f,A,c) \to \varphi_{33}) \to \varphi_{108}$}
\end{ptree}
\noindent$\Pi_{80}$:
\begin{ptree}
\AXC{$\varphi_{112}$\lbl{A2}}
\AXC{$\stackrel{\Pi_{51}}{\vdots}$}\noLine
\UIC{$\mathrm{ContIn}(f,A,c) \to \mathrm{ContIn}(g,A,c) \to \varphi_{16}$}
\AXC{$\varphi_{235}$\lbl{A1}}
\BIC{$\varphi_{112} \to \mathrm{ContIn}(f,A,c) \to \varphi_{33}$}
\BIC{$\mathrm{ContIn}(f,A,c) \to \varphi_{33}$}
\AXC{$\stackrel{\Pi_{79}}{\vdots}$}\noLine
\UIC{$(\mathrm{ContIn}(f,A,c) \to \varphi_{33}) \to \varphi_{108}$}
\BIC{$\varphi_{108}$}
\end{ptree}
\noindent$\Pi_{81}$:
\begin{ptree}
\AXC{$\stackrel{\Pi_{80}}{\vdots}$}\noLine
\UIC{$\varphi_{108}$}
\AXC{$\varphi_{27} \to \varphi_{24}$\lbl{A2}}
\AXC{$\stackrel{\Pi_{59}}{\vdots}$}\noLine
\UIC{$\mathrm{ContIn}(f,A,c) \to \mathrm{ContIn}(g,A,c) \to \varphi_{27}$}
\AXC{$\varphi_{240}$\lbl{A1}}
\BIC{$\varphi_{213}$}
\BIC{$\mathrm{ContIn}(f,A,c) \to \mathrm{ContIn}(g,A,c) \to \varphi_{24}$}
\AXC{$\varphi_{242}$\lbl{A1}}
\BIC{$\varphi_{108} \to \varphi_{134}$}
\BIC{$\varphi_{134}$}
\end{ptree}
\noindent$\Pi_{82}$:
\begin{ptree}
\AXC{$\stackrel{\Pi_{72}}{\vdots}$}\noLine
\UIC{$\varphi_{56}$}
\AXC{$\stackrel{\Pi_{81}}{\vdots}$}\noLine
\UIC{$\varphi_{134}$}
\AXC{$\varphi_{134} \to \varphi_{56} \to \varphi_{135}$\lbl{A1}}
\BIC{$\varphi_{56} \to \varphi_{135}$}
\BIC{$\varphi_{135}$}
\AXC{$\varphi_{135} \to (\mathrm{ContIn}(f,A,c) \to \varphi_{130}) \to \varphi_{230}$\lbl{A1}}
\BIC{$(\mathrm{ContIn}(f,A,c) \to \varphi_{130}) \to \varphi_{230}$}
\end{ptree}
\noindent$\Pi_{83}$:
\begin{ptree}
\AXC{$\stackrel{\Pi_{2}}{\vdots}$}\noLine
\UIC{$\varphi_{123}$}
\AXC{$\stackrel{\Pi_{5}}{\vdots}$}\noLine
\UIC{$\varphi_{124}$}
\AXC{$\stackrel{\Pi_{6}}{\vdots}$}\noLine
\UIC{$\varphi_{124} \to \varphi_{123} \to \mathrm{ContIn}(f,A,c) \to \varphi_{141}$}
\BIC{$\varphi_{123} \to \mathrm{ContIn}(f,A,c) \to \varphi_{141}$}
\BIC{$\mathrm{ContIn}(f,A,c) \to \varphi_{141}$}
\AXC{$\stackrel{\Pi_{71}}{\vdots}$}\noLine
\UIC{$\varphi_{234}$}
\BIC{$\mathrm{ContIn}(f,A,c) \to \varphi_{130}$}
\AXC{$\stackrel{\Pi_{82}}{\vdots}$}\noLine
\UIC{$(\mathrm{ContIn}(f,A,c) \to \varphi_{130}) \to \varphi_{230}$}
\BIC{$\varphi_{230}$}
\end{ptree}
\noindent$\Pi_{84}$:
\begin{ptree}
\AXC{$z = w \to \varphi_{21}$\lbl{A6}}
\UIC{$\forall w\,(z = w \to \varphi_{21})$}
\UIC{$\forall z\,\forall w\,(z = w \to \varphi_{21})$}
\AXC{$\forall z\,\forall w\,(z = w \to \varphi_{21}) \to \varphi_{109}$\lbl{A2}}
\BIC{$\varphi_{109}$}
\AXC{$\varphi_{109} \to f(c) + g(c) = (f + g)(c) \to \varphi_{89}$\lbl{A2}}
\BIC{$f(c) + g(c) = (f + g)(c) \to \varphi_{89}$}
\AXC{$\stackrel{\Pi_{83}}{\vdots}$}\noLine
\UIC{$\varphi_{230}$}
\BIC{$\mathrm{ContIn}(f,A,c) \to \varphi_{125}$}
\end{ptree}
\noindent$\Pi_{85}$:
\begin{ptree}
\AXC{$\varphi_{232}$\lbl{A3}}
\AXC{$\stackrel{\Pi_{84}}{\vdots}$}\noLine
\UIC{$\mathrm{ContIn}(f,A,c) \to \varphi_{125}$}
\AXC{$(\mathrm{ContIn}(f,A,c) \to \varphi_{125}) \to \varphi_{140}$\lbl{A1}}
\BIC{$\varphi_{140}$}
\UIC{$\forall a\,(\varphi_{140})$}
\AXC{$\forall a\,(\varphi_{140}) \to \varphi_{232} \to \varphi_{203}$\lbl{A1}}
\BIC{$\varphi_{232} \to \varphi_{203}$}
\BIC{$\varphi_{203}$}
\end{ptree}
\noindent$\Pi_{86}$:
\begin{ptree}
\AXC{$\varphi_{113}$\lbl{A2}}
\AXC{$\stackrel{\Pi_{51}}{\vdots}$}\noLine
\UIC{$\mathrm{ContIn}(f,A,c) \to \mathrm{ContIn}(g,A,c) \to \varphi_{16}$}
\AXC{$\varphi_{236}$\lbl{A1}}
\BIC{$\varphi_{113} \to \mathrm{ContIn}(f,A,c) \to \varphi_{34}$}
\BIC{$\mathrm{ContIn}(f,A,c) \to \varphi_{34}$}
\end{ptree}
\noindent$\Pi_{87}$:
\begin{ptree}
\AXC{$\forall X\,\forall Y\,\forall x\,(X \subseteq Y \to x \in X \to x \in Y)$\lbl{[B5]}}
\AXC{$\varphi_{46}$\lbl{A2}}
\BIC{$\forall Y\,\forall x\,(A \subseteq Y \to x \in A \to x \in Y)$}
\AXC{$\varphi_{57}$\lbl{A2}}
\BIC{$\forall x\,(A \subseteq \mathrm{dom}(g) \to x \in A \to x \in \mathrm{dom}(g))$}
\AXC{$\varphi_{67}$\lbl{A2}}
\BIC{$A \subseteq \mathrm{dom}(g) \to u \in A \to u \in \mathrm{dom}(g)$}
\end{ptree}
\noindent$\Pi_{88}$:
\begin{ptree}
\AXC{$\varphi_{12}$\lbl{A1}}
\AXC{$\stackrel{\Pi_{14}}{\vdots}$}\noLine
\UIC{$\varphi_{8}$}
\AXC{$\stackrel{\Pi_{87}}{\vdots}$}\noLine
\UIC{$A \subseteq \mathrm{dom}(g) \to u \in A \to u \in \mathrm{dom}(g)$}
\AXC{$(A \subseteq \mathrm{dom}(g) \to u \in A \to u \in \mathrm{dom}(g)) \to \varphi_{8} \to \varphi_{12} \to \varphi_{17}$\lbl{A1}}
\BIC{$\varphi_{8} \to \varphi_{12} \to \varphi_{17}$}
\BIC{$\varphi_{12} \to \varphi_{17}$}
\BIC{$\varphi_{17}$}
\end{ptree}
\noindent$\Pi_{89}$:
\begin{ptree}
\AXC{$\forall X\,\forall Y\,\forall x\,(X \subseteq Y \to x \in X \to x \in Y)$\lbl{[B5]}}
\AXC{$\varphi_{46}$\lbl{A2}}
\BIC{$\forall Y\,\forall x\,(A \subseteq Y \to x \in A \to x \in Y)$}
\AXC{$\varphi_{58}$\lbl{A2}}
\BIC{$\forall x\,(A \subseteq \mathrm{dom}(f) \to x \in A \to x \in \mathrm{dom}(f))$}
\AXC{$\varphi_{68}$\lbl{A2}}
\BIC{$A \subseteq \mathrm{dom}(f) \to u \in A \to u \in \mathrm{dom}(f)$}
\end{ptree}
\noindent$\Pi_{90}$:
\begin{ptree}
\AXC{$\varphi_{12}$\lbl{A1}}
\AXC{$\stackrel{\Pi_{28}}{\vdots}$}\noLine
\UIC{$\varphi_{9}$}
\AXC{$\stackrel{\Pi_{89}}{\vdots}$}\noLine
\UIC{$A \subseteq \mathrm{dom}(f) \to u \in A \to u \in \mathrm{dom}(f)$}
\AXC{$(A \subseteq \mathrm{dom}(f) \to u \in A \to u \in \mathrm{dom}(f)) \to \varphi_{9} \to \varphi_{12} \to \varphi_{18}$\lbl{A1}}
\BIC{$\varphi_{9} \to \varphi_{12} \to \varphi_{18}$}
\BIC{$\varphi_{12} \to \varphi_{18}$}
\BIC{$\varphi_{18}$}
\end{ptree}
\noindent$\Pi_{91}$:
\begin{ptree}
\AXC{$\stackrel{\Pi_{86}}{\vdots}$}\noLine
\UIC{$\mathrm{ContIn}(f,A,c) \to \varphi_{34}$}
\AXC{$\stackrel{\Pi_{88}}{\vdots}$}\noLine
\UIC{$\varphi_{17}$}
\AXC{$\stackrel{\Pi_{90}}{\vdots}$}\noLine
\UIC{$\varphi_{18}$}
\AXC{$\varphi_{18} \to \varphi_{17} \to (\mathrm{ContIn}(f,A,c) \to \varphi_{34}) \to \varphi_{25}$\lbl{A1}}
\BIC{$\varphi_{17} \to (\mathrm{ContIn}(f,A,c) \to \varphi_{34}) \to \varphi_{25}$}
\BIC{$(\mathrm{ContIn}(f,A,c) \to \varphi_{34}) \to \varphi_{25}$}
\BIC{$\varphi_{25}$}
\AXC{$\varphi_{25} \to \varphi_{26}$\lbl{A1}}
\BIC{$\varphi_{26}$}
\UIC{$\forall u\,(\varphi_{26})$}
\end{ptree}
\noindent$\Pi_{92}$:
\begin{ptree}
\AXC{$\forall u\,(\varphi_{26}) \to \varphi_{2} \to \forall u \in A\,(u \in \mathrm{dom}((f + g)))$\lbl{A3}}
\AXC{$\stackrel{\Pi_{91}}{\vdots}$}\noLine
\UIC{$\forall u\,(\varphi_{26})$}
\AXC{$\varphi_{248}$\lbl{A1}}
\BIC{$(\forall u\,(\varphi_{26}) \to \varphi_{2} \to \forall u \in A\,(u \in \mathrm{dom}((f + g)))) \to \varphi_{149}$}
\BIC{$\varphi_{149}$}
\end{ptree}
\noindent$\Pi_{93}$:
\begin{ptree}
\AXC{$\forall x\,\forall y\,(x \subseteq y \leftrightarrow \forall u \in x\,(u \in y))$\lbl{Def$\subseteq$}}
\AXC{$\varphi_{49}$\lbl{A2}}
\BIC{$\forall y\,(A \subseteq y \leftrightarrow \forall u \in A\,(u \in y))$}
\AXC{$\forall y\,(A \subseteq y \leftrightarrow \forall u \in A\,(u \in y)) \to \varphi_{20}$\lbl{A2}}
\BIC{$\varphi_{20}$}
\end{ptree}
\noindent$\Pi_{94}$:
\begin{ptree}
\AXC{$\mathrm{ContIn}(f,A,c) \land \mathrm{ContIn}(g,A,c) \to \varphi_{2}$\lbl{A4}}
\AXC{$\stackrel{\Pi_{92}}{\vdots}$}\noLine
\UIC{$\varphi_{149}$}
\BIC{$\mathrm{ContIn}(f,A,c) \to \varphi_{13}$}
\AXC{$\stackrel{\Pi_{93}}{\vdots}$}\noLine
\UIC{$\varphi_{20}$}
\AXC{$\varphi_{20} \to (\mathrm{ContIn}(f,A,c) \to \varphi_{13}) \to \varphi_{14}$\lbl{A1}}
\BIC{$(\mathrm{ContIn}(f,A,c) \to \varphi_{13}) \to \varphi_{14}$}
\BIC{$\varphi_{14}$}
\end{ptree}
\noindent$\Pi_{95}$:
\begin{ptree}
\AXC{$\stackrel{\Pi_{74}}{\vdots}$}\noLine
\UIC{$\varphi_{6}$}
\AXC{$\stackrel{\Pi_{94}}{\vdots}$}\noLine
\UIC{$\varphi_{14}$}
\AXC{$\stackrel{\Pi_{62}}{\vdots}$}\noLine
\UIC{$\varphi_{10}$}
\AXC{$\varphi_{161}$\lbl{[G2]}}
\AXC{$\varphi_{161} \to \forall A\,(\varphi_{195})$\lbl{A2}}
\BIC{$\forall A\,(\varphi_{195})$}
\AXC{$\forall A\,(\varphi_{195}) \to \varphi_{195}$\lbl{A2}}
\BIC{$\varphi_{195}$}
\AXC{$\varphi_{250}$\lbl{A2}}
\BIC{$\mathrm{RFn}((f + g)) \to A \subseteq \mathrm{dom}((f + g)) \to c \in A \to \varphi_{137}$}
\AXC{$\varphi_{273}$\lbl{A1}}
\BIC{$\varphi_{10} \to \varphi_{14} \to \varphi_{6} \to \varphi_{198}$}
\BIC{$\varphi_{14} \to \varphi_{6} \to \varphi_{198}$}
\BIC{$\varphi_{6} \to \varphi_{198}$}
\BIC{$\varphi_{198}$}
\end{ptree}
\noindent 主木:
\begin{ptree}
\AXC{$\mathrm{ContIn}(f,A,c) \land \mathrm{ContIn}(g,A,c) \to \varphi_{1}$\lbl{A4}}
\AXC{$\stackrel{\Pi_{85}}{\vdots}$}\noLine
\UIC{$\varphi_{203}$}
\BIC{$\mathrm{ContIn}(f,A,c) \to \varphi_{127}$}
\AXC{$\stackrel{\Pi_{95}}{\vdots}$}\noLine
\UIC{$\varphi_{198}$}
\BIC{$\mathrm{ContIn}(f,A,c) \to \varphi_{3}$}
\UIC{$\forall c\,(\mathrm{ContIn}(f,A,c) \to \varphi_{3})$}
\UIC{$\forall A\,\forall c\,(\mathrm{ContIn}(f,A,c) \to \varphi_{3})$}
\UIC{$\forall g\,\forall A\,\forall c\,(\mathrm{ContIn}(f,A,c) \to \varphi_{3})$}
\UIC{$\forall f\,\forall g\,\forall A\,\forall c\,(\mathrm{ContIn}(f,A,c) \to \varphi_{3})$}
\end{ptree}
\begin{abbrevlist}
\item $\varphi_{1} :\equiv \forall a\,(\mathrm{ContIn}(f,A,c) \land \mathrm{ContIn}(g,A,c))$
\item $\varphi_{2} :\equiv \forall u\,(\mathrm{ContIn}(f,A,c) \land \mathrm{ContIn}(g,A,c))$
\item $\varphi_{3} :\equiv \mathrm{ContIn}(g,A,c) \to \mathrm{ContIn}((f + g),A,c)$
\item $\varphi_{4} :\equiv \mathrm{ContIn}(f,A,c) \to \mathrm{ContIn}(g,A,c) \to \mathrm{RFn}(g)$
\item $\varphi_{5} :\equiv \mathrm{ContIn}(f,A,c) \to \mathrm{ContIn}(g,A,c) \to \mathrm{RFn}(f)$
\item $\varphi_{6} :\equiv \mathrm{ContIn}(f,A,c) \to \mathrm{ContIn}(g,A,c) \to c \in A$
\item $\varphi_{7} :\equiv \mathrm{ContIn}(f,A,c) \to \mathrm{ContIn}(g,A,c) \to \mathrm{ContIn}(g,A,c)$
\item $\varphi_{8} :\equiv \mathrm{ContIn}(f,A,c) \to \mathrm{ContIn}(g,A,c) \to A \subseteq \mathrm{dom}(g)$
\item $\varphi_{9} :\equiv \mathrm{ContIn}(f,A,c) \to \mathrm{ContIn}(g,A,c) \to A \subseteq \mathrm{dom}(f)$
\item $\varphi_{10} :\equiv \mathrm{ContIn}(f,A,c) \to \mathrm{ContIn}(g,A,c) \to \mathrm{RFn}((f + g))$
\item $\varphi_{11} :\equiv \mathrm{ContIn}(f,A,c) \to \mathrm{ContIn}(g,A,c) \to \mathrm{Seq}(a) \to \mathrm{Seq}(a)$
\item $\varphi_{12} :\equiv \mathrm{ContIn}(f,A,c) \to \mathrm{ContIn}(g,A,c) \to u \in A \to u \in A$
\item $\varphi_{13} :\equiv \mathrm{ContIn}(g,A,c) \to \forall u \in A\,(u \in \mathrm{dom}((f + g)))$
\item $\varphi_{14} :\equiv \mathrm{ContIn}(f,A,c) \to \mathrm{ContIn}(g,A,c) \to A \subseteq \mathrm{dom}((f + g))$
\item $\varphi_{15} :\equiv \forall n \in \mathbb{N}\,(a_{n} \in A) \to \mathrm{Conv}(a,c) \to \mathrm{Conv}(a,c)$
\item $\varphi_{16} :\equiv \forall x\,(x \in \mathrm{dom}((f + g)) \leftrightarrow x \in \mathrm{dom}(f) \land x \in \mathrm{dom}(g))$
\item $\varphi_{17} :\equiv \mathrm{ContIn}(f,A,c) \to \mathrm{ContIn}(g,A,c) \to u \in A \to u \in \mathrm{dom}(g)$
\item $\varphi_{18} :\equiv \mathrm{ContIn}(f,A,c) \to \mathrm{ContIn}(g,A,c) \to u \in A \to u \in \mathrm{dom}(f)$
\item $\varphi_{19} :\equiv \mathrm{Conv}(a,c) \to \mathrm{Conv}((f + g) \circ a,(f + g)(c))$
\item $\varphi_{20} :\equiv A \subseteq \mathrm{dom}((f + g)) \leftrightarrow \forall u \in A\,(u \in \mathrm{dom}((f + g)))$
\item $\varphi_{21} :\equiv \mathrm{Conv}((f + g) \circ a,z) \to \mathrm{Conv}((f + g) \circ a,w)$
\item $\varphi_{22} :\equiv \mathrm{Conv}(a,c) \to \mathrm{Conv}((f + g) \circ a,f(c) + g(c))$
\item $\varphi_{23} :\equiv \forall x\,\forall y\,(x = y \to y = x) \to \forall y\,((f + g)(c) = y \to y = (f + g)(c))$
\item $\varphi_{24} :\equiv c \in \mathrm{dom}((f + g)) \to (f + g)(c) = f(c) + g(c)$
\item $\varphi_{25} :\equiv \mathrm{ContIn}(f,A,c) \to \mathrm{ContIn}(g,A,c) \to u \in A \to u \in \mathrm{dom}((f + g))$
\item $\varphi_{26} :\equiv \mathrm{ContIn}(f,A,c) \land \mathrm{ContIn}(g,A,c) \to u \in A \to u \in \mathrm{dom}((f + g))$
\item $\varphi_{27} :\equiv \forall x \in \mathrm{dom}((f + g))\,((f + g)(x) = f(x) + g(x))$
\item $\varphi_{28} :\equiv \forall n \in \mathbb{N}\,(a_{n} \in A) \to \mathrm{Conv}(a,c) \to c \in \mathrm{dom}((f + g))$
\item $\varphi_{29} :\equiv \forall a\,\forall \alpha\,\forall b\,\forall \beta\,(\mathrm{Conv}(a,\alpha) \to \mathrm{Conv}(b,\beta) \to \mathrm{Conv}(a \oplus b,\alpha + \beta))$
\item $\varphi_{30} :\equiv \forall n \in \mathbb{N}\,(a_{n} \in A) \to \mathrm{Conv}(a,c) \to n \in \mathbb{N} \to n \in \mathbb{N}$
\item $\varphi_{31} :\equiv \mathrm{Seq}(a) \to \forall n \in \mathbb{N}\,(a_{n} \in A) \to \mathrm{Conv}(a,c) \to c \in \mathrm{dom}(g)$
\item $\varphi_{32} :\equiv \mathrm{Seq}(a) \to \forall n \in \mathbb{N}\,(a_{n} \in A) \to \mathrm{Conv}(a,c) \to c \in \mathrm{dom}(f)$
\item $\varphi_{33} :\equiv \mathrm{ContIn}(g,A,c) \to (c \in \mathrm{dom}((f + g)) \leftrightarrow c \in \mathrm{dom}(f) \land c \in \mathrm{dom}(g))$
\item $\varphi_{34} :\equiv \mathrm{ContIn}(g,A,c) \to (u \in \mathrm{dom}((f + g)) \leftrightarrow u \in \mathrm{dom}(f) \land u \in \mathrm{dom}(g))$
\item $\varphi_{35} :\equiv \forall n \in \mathbb{N}\,(a_{n} \in A) \to \mathrm{Conv}(a,c) \to \mathrm{Seq}((f + g) \circ a)$
\item $\varphi_{36} :\equiv \forall n \in \mathbb{N}\,(a_{n} \in A) \to \mathrm{Conv}(a,c) \to n \in \mathbb{N} \to a_{n} \in A$
\item $\varphi_{37} :\equiv \mathrm{Seq}(a) \to \forall n \in \mathbb{N}\,(a_{n} \in A) \to \mathrm{Conv}(a,c) \to \mathrm{Seq}(g \circ a)$
\item $\varphi_{38} :\equiv \mathrm{Seq}(a) \to \forall n \in \mathbb{N}\,(a_{n} \in A) \to \mathrm{Conv}(a,c) \to \mathrm{Seq}(f \circ a)$
\item $\varphi_{39} :\equiv \forall n \in \mathbb{N}\,(a_{n} \in A) \to \mathrm{Conv}(a,c) \to \mathrm{Conv}(g \circ a,g(c))$
\item $\varphi_{40} :\equiv \forall n \in \mathbb{N}\,(a_{n} \in A) \to \mathrm{Conv}(a,c) \to \mathrm{Conv}(f \circ a,f(c))$
\item $\varphi_{41} :\equiv \mathrm{Conv}(a,c) \to \mathrm{Conv}(f \circ a \oplus g \circ a,f(c) + g(c))$
\item $\varphi_{42} :\equiv \mathrm{Conv}(a,c) \to n \in \mathbb{N} \to n \in \mathbb{N} \to g \circ a(n) = g(a_{n})$
\item $\varphi_{43} :\equiv \mathrm{Conv}(a,c) \to n \in \mathbb{N} \to n \in \mathbb{N} \to f \circ a(n) = f(a_{n})$
\item $\varphi_{44} :\equiv a_{n} \in \mathrm{dom}((f + g)) \leftrightarrow a_{n} \in \mathrm{dom}(f) \land a_{n} \in \mathrm{dom}(g)$
\item $\varphi_{45} :\equiv \forall y\,\forall z\,((f + g) \circ a(n) = y \to y = z \to (f + g) \circ a(n) = z)$
\item $\varphi_{46} :\equiv \forall X\,\forall Y\,\forall x\,(X \subseteq Y \to x \in X \to x \in Y) \to \forall Y\,\forall x\,(A \subseteq Y \to x \in A \to x \in Y)$
\item $\varphi_{47} :\equiv \forall a\,\forall \alpha\,(\mathrm{Conv}(a,\alpha) \to \mathrm{Seq}(a)) \to \forall \alpha\,(\mathrm{Conv}(g \circ a,\alpha) \to \mathrm{Seq}(g \circ a))$
\item $\varphi_{48} :\equiv \forall a\,\forall \alpha\,(\mathrm{Conv}(a,\alpha) \to \mathrm{Seq}(a)) \to \forall \alpha\,(\mathrm{Conv}(f \circ a,\alpha) \to \mathrm{Seq}(f \circ a))$
\item $\varphi_{49} :\equiv \forall x\,\forall y\,(x \subseteq y \leftrightarrow \forall u \in x\,(u \in y)) \to \forall y\,(A \subseteq y \leftrightarrow \forall u \in A\,(u \in y))$
\item $\varphi_{50} :\equiv \mathrm{Conv}(a,c) \to n \in \mathbb{N} \to \forall n \in \mathbb{N}\,(g \circ a(n) = g(a_{n}))$
\item $\varphi_{51} :\equiv \mathrm{Conv}(a,c) \to n \in \mathbb{N} \to \forall n \in \mathbb{N}\,(f \circ a(n) = f(a_{n}))$
\item $\varphi_{52} :\equiv \forall n \in \mathbb{N}\,(a_{n} \in A) \to \mathrm{Conv}(a,c) \to n \in \mathbb{N} \to a_{n} \in \mathrm{dom}(g)$
\item $\varphi_{53} :\equiv \forall n \in \mathbb{N}\,(a_{n} \in A) \to \mathrm{Conv}(a,c) \to n \in \mathbb{N} \to a_{n} \in \mathrm{dom}(f)$
\item $\varphi_{54} :\equiv \mathrm{Conv}(a,c) \to n \in \mathbb{N} \to (f + g) \circ a(n) = (f + g)(a_{n})$
\item $\varphi_{55} :\equiv \forall a\,\forall b\,\forall \alpha\,(\mathrm{Seq}(b) \to \forall n \in \mathbb{N}\,(a_{n} = b_{n}) \to \mathrm{Conv}(a,\alpha) \to \mathrm{Conv}(b,\alpha))$
\item $\varphi_{56} :\equiv (f + g)(c) = f(c) + g(c) \to f(c) + g(c) = (f + g)(c)$
\item $\varphi_{57} :\equiv \forall Y\,\forall x\,(A \subseteq Y \to x \in A \to x \in Y) \to \forall x\,(A \subseteq \mathrm{dom}(g) \to x \in A \to x \in \mathrm{dom}(g))$
\item $\varphi_{58} :\equiv \forall Y\,\forall x\,(A \subseteq Y \to x \in A \to x \in Y) \to \forall x\,(A \subseteq \mathrm{dom}(f) \to x \in A \to x \in \mathrm{dom}(f))$
\item $\varphi_{59} :\equiv \forall n \in \mathbb{N}\,(a_{n} \in A) \to \mathrm{Conv}(a,c) \to \forall n \in \mathbb{N}\,(a_{n} \in \mathrm{dom}(g))$
\item $\varphi_{60} :\equiv \forall n \in \mathbb{N}\,(a_{n} \in A) \to \mathrm{Conv}(a,c) \to \forall n \in \mathbb{N}\,(a_{n} \in \mathrm{dom}(f))$
\item $\varphi_{61} :\equiv n \in \mathbb{N} \to \forall n \in \mathbb{N}\,((f + g) \circ a(n) = (f + g)(a_{n}))$
\item $\varphi_{62} :\equiv \forall n \in \mathbb{N}\,(a_{n} \in A) \to \mathrm{Conv}(a,c) \to (f + g)(c) = f(c) + g(c)$
\item $\varphi_{63} :\equiv \forall n \in \mathbb{N}\,(a_{n} \in A) \to \mathrm{Conv}(a,c) \to f(c) + g(c) = (f + g)(c)$
\item $\varphi_{64} :\equiv \forall x\,\forall y\,(x = y \to y = x) \to \forall y\,((f + g) \circ a(n) = y \to y = (f + g) \circ a(n))$
\item $\varphi_{65} :\equiv \forall x\,(A \subseteq \mathrm{dom}(g) \to x \in A \to x \in \mathrm{dom}(g)) \to A \subseteq \mathrm{dom}(g) \to c \in A \to c \in \mathrm{dom}(g)$
\item $\varphi_{66} :\equiv \forall x\,(A \subseteq \mathrm{dom}(f) \to x \in A \to x \in \mathrm{dom}(f)) \to A \subseteq \mathrm{dom}(f) \to c \in A \to c \in \mathrm{dom}(f)$
\item $\varphi_{67} :\equiv \forall x\,(A \subseteq \mathrm{dom}(g) \to x \in A \to x \in \mathrm{dom}(g)) \to A \subseteq \mathrm{dom}(g) \to u \in A \to u \in \mathrm{dom}(g)$
\item $\varphi_{68} :\equiv \forall x\,(A \subseteq \mathrm{dom}(f) \to x \in A \to x \in \mathrm{dom}(f)) \to A \subseteq \mathrm{dom}(f) \to u \in A \to u \in \mathrm{dom}(f)$
\item $\varphi_{69} :\equiv \mathrm{ContIn}(f,A,c) \land \mathrm{ContIn}(g,A,c) \land \mathrm{Seq}(a) \land \forall n \in \mathbb{N}\,(a_{n} \in A) \land \mathrm{Conv}(a,c)$
\item $\varphi_{70} :\equiv \mathrm{Conv}(a,c) \to \mathrm{Seq}(g \circ a) \land \forall n \in \mathbb{N}\,(g \circ a(n) = g(a_{n}))$
\item $\varphi_{71} :\equiv \mathrm{Conv}(a,c) \to \mathrm{Seq}(f \circ a) \land \forall n \in \mathbb{N}\,(f \circ a(n) = f(a_{n}))$
\item $\varphi_{72} :\equiv \mathrm{Conv}(a,c) \to n \in \mathbb{N} \to (f + g)(a_{n}) = f(a_{n}) + g(a_{n})$
\item $\varphi_{73} :\equiv \mathrm{ContIn}(g,A,c) \to \mathrm{Seq}(a) \to \forall n \in \mathbb{N}\,(a_{n} \in A) \to n \in \mathbb{N} \to a_{n} \in A$
\item $\varphi_{74} :\equiv \forall n \in \mathbb{N}\,(a_{n} \in A) \to \mathrm{Conv}(a,c) \to n \in \mathbb{N} \to a_{n} \in \mathrm{dom}((f + g))$
\item $\varphi_{75} :\equiv \forall \alpha\,\forall b\,\forall \beta\,(\mathrm{Conv}(f \circ a,\alpha) \to \mathrm{Conv}(b,\beta) \to \mathrm{Conv}(f \circ a \oplus b,\alpha + \beta))$
\item $\varphi_{76} :\equiv \forall n \in \mathbb{N}\,(a_{n} \in \mathrm{dom}((f + g))) \to n \in \mathbb{N} \to a_{n} \in \mathrm{dom}((f + g))$
\item $\varphi_{77} :\equiv \forall x\,(x = x) \to f(a_{n}) + g \circ a(n) = f(a_{n}) + g \circ a(n)$
\item $\varphi_{78} :\equiv \mathrm{ContIn}(g,A,c) \to \mathrm{Seq}(a) \to \forall n \in \mathbb{N}\,(a_{n} \in A) \to \forall n \in \mathbb{N}\,(a_{n} \in A)$
\item $\varphi_{79} :\equiv \forall n \in \mathbb{N}\,(a_{n} \in A) \to \mathrm{Conv}(a,c) \to \forall n \in \mathbb{N}\,(a_{n} \in \mathrm{dom}((f + g)))$
\item $\varphi_{80} :\equiv \mathrm{Conv}(a,c) \to n \in \mathbb{N} \to (f + g) \circ a(n) = f(a_{n}) + g(a_{n})$
\item $\varphi_{81} :\equiv \mathrm{Conv}(a,c) \to n \in \mathbb{N} \to n \in \mathbb{N} \to (f + g) \circ a(n) = (f + g)(a_{n})$
\item $\varphi_{82} :\equiv a_{n} \in \mathrm{dom}((f + g)) \to (f + g)(a_{n}) = f(a_{n}) + g(a_{n})$
\item $\varphi_{83} :\equiv \mathrm{Conv}(g \circ a,g(c)) \to \mathrm{Conv}(f \circ a \oplus g \circ a,f(c) + g(c))$
\item $\varphi_{84} :\equiv \forall n \in \mathbb{N}\,(a_{n} \in A) \to \mathrm{Conv}(a,c) \to n \in \mathbb{N} \to g \circ a(n) = g(a_{n})$
\item $\varphi_{85} :\equiv \forall n \in \mathbb{N}\,(a_{n} \in A) \to \mathrm{Conv}(a,c) \to n \in \mathbb{N} \to f \circ a(n) = f(a_{n})$
\item $\varphi_{86} :\equiv \mathrm{Conv}(a,c) \to n \in \mathbb{N} \to (f + g) \circ a(n) = f \circ a \oplus g \circ a(n)$
\item $\varphi_{87} :\equiv \mathrm{Conv}(a,c) \to n \in \mathbb{N} \to f \circ a \oplus g \circ a(n) = (f + g) \circ a(n)$
\item $\varphi_{88} :\equiv \forall b\,\forall \beta\,(\mathrm{Conv}(f \circ a,f(c)) \to \mathrm{Conv}(b,\beta) \to \mathrm{Conv}(f \circ a \oplus b,f(c) + \beta))$
\item $\varphi_{89} :\equiv \mathrm{Conv}((f + g) \circ a,f(c) + g(c)) \to \mathrm{Conv}((f + g) \circ a,(f + g)(c))$
\item $\varphi_{90} :\equiv n \in \mathbb{N} \to f \circ a(n) + g \circ a(n) = f(a_{n}) + g \circ a(n)$
\item $\varphi_{91} :\equiv \mathrm{ContIn}(f,A,c) \to \mathrm{ContIn}(g,A,c) \to \mathrm{Seq}(a) \to \varphi_{15}$
\item $\varphi_{92} :\equiv \forall x\,(x = x) \to f \circ a(n) + g \circ a(n) = f \circ a(n) + g \circ a(n)$
\item $\varphi_{93} :\equiv \mathrm{Seq}((f + g) \circ a) \land \forall n \in \mathbb{N}\,((f + g) \circ a(n) = (f + g)(a_{n}))$
\item $\varphi_{94} :\equiv \mathrm{Conv}(a,c) \to \forall n \in \mathbb{N}\,(f \circ a \oplus g \circ a(n) = (f + g) \circ a(n))$
\item $\varphi_{95} :\equiv \forall \alpha\,(\mathrm{Conv}(g \circ a,\alpha) \to \mathrm{Seq}(g \circ a)) \to \mathrm{Conv}(g \circ a,g(c)) \to \mathrm{Seq}(g \circ a)$
\item $\varphi_{96} :\equiv \forall \alpha\,(\mathrm{Conv}(f \circ a,\alpha) \to \mathrm{Seq}(f \circ a)) \to \mathrm{Conv}(f \circ a,f(c)) \to \mathrm{Seq}(f \circ a)$
\item $\varphi_{97} :\equiv \forall a\,\forall b\,(\mathrm{Seq}(a) \to \mathrm{Seq}(b) \to \mathrm{Seq}(a \oplus b) \land \forall n \in \mathbb{N}\,(a \oplus b(n) = a_{n} + b_{n}))$
\item $\varphi_{98} :\equiv \forall x\,(A \subseteq \mathrm{dom}(g) \to x \in A \to x \in \mathrm{dom}(g)) \to A \subseteq \mathrm{dom}(g) \to a_{n} \in A \to a_{n} \in \mathrm{dom}(g)$
\item $\varphi_{99} :\equiv \forall x\,(A \subseteq \mathrm{dom}(f) \to x \in A \to x \in \mathrm{dom}(f)) \to A \subseteq \mathrm{dom}(f) \to a_{n} \in A \to a_{n} \in \mathrm{dom}(f)$
\item $\varphi_{100} :\equiv \forall y\,\forall z\,(f \circ a \oplus g \circ a(n) = y \to y = z \to f \circ a \oplus g \circ a(n) = z)$
\item $\varphi_{101} :\equiv \mathrm{Conv}(a,c) \to n \in \mathbb{N} \to f(a_{n}) + g \circ a(n) = f(a_{n}) + g(a_{n})$
\item $\varphi_{102} :\equiv \forall n \in \mathbb{N}\,(f \circ a \oplus g \circ a(n) = f \circ a(n) + g \circ a(n))$
\item $\varphi_{103} :\equiv \mathrm{Conv}(a,c) \to n \in \mathbb{N} \to f \circ a \oplus g \circ a(n) = f(a_{n}) + g(a_{n})$
\item $\varphi_{104} :\equiv \mathrm{Conv}(a,c) \to n \in \mathbb{N} \to f(a_{n}) + g(a_{n}) = f \circ a \oplus g \circ a(n)$
\item $\varphi_{105} :\equiv \forall n \in \mathbb{N}\,(g \circ a(n) = g(a_{n})) \to n \in \mathbb{N} \to g \circ a(n) = g(a_{n})$
\item $\varphi_{106} :\equiv \forall n \in \mathbb{N}\,(f \circ a(n) = f(a_{n})) \to n \in \mathbb{N} \to f \circ a(n) = f(a_{n})$
\item $\varphi_{107} :\equiv \mathrm{Conv}(a,c) \to n \in \mathbb{N} \to f \circ a \oplus g \circ a(n) = f(a_{n}) + g \circ a(n)$
\item $\varphi_{108} :\equiv \mathrm{ContIn}(f,A,c) \to \mathrm{ContIn}(g,A,c) \to \mathrm{Seq}(a) \to \varphi_{28}$
\item $\varphi_{109} :\equiv \forall w\,(f(c) + g(c) = w \to \mathrm{Conv}((f + g) \circ a,f(c) + g(c)) \to \mathrm{Conv}((f + g) \circ a,w))$
\item $\varphi_{110} :\equiv \forall x\,\forall y\,(x = y \to y = x) \to \forall y\,(f \circ a \oplus g \circ a(n) = y \to y = f \circ a \oplus g \circ a(n))$
\item $\varphi_{111} :\equiv \mathrm{ContIn}(f,A,c) \to \mathrm{ContIn}(g,A,c) \to \mathrm{Seq}(a) \to \varphi_{30}$
\item $\varphi_{112} :\equiv \varphi_{16} \to (c \in \mathrm{dom}((f + g)) \leftrightarrow c \in \mathrm{dom}(f) \land c \in \mathrm{dom}(g))$
\item $\varphi_{113} :\equiv \varphi_{16} \to (u \in \mathrm{dom}((f + g)) \leftrightarrow u \in \mathrm{dom}(f) \land u \in \mathrm{dom}(g))$
\item $\varphi_{114} :\equiv n \in \mathbb{N} \to n \in \mathbb{N} \to f \circ a \oplus g \circ a(n) = f \circ a(n) + g \circ a(n)$
\item $\varphi_{115} :\equiv \mathrm{Conv}(a,c) \to n \in \mathbb{N} \to f \circ a \oplus g \circ a(n) = f \circ a(n) + g \circ a(n)$
\item $\varphi_{116} :\equiv \forall n \in \mathbb{N}\,(a_{n} \in \mathrm{dom}(g)) \to \mathrm{Seq}(g \circ a) \land \forall n \in \mathbb{N}\,(g \circ a(n) = g(a_{n}))$
\item $\varphi_{117} :\equiv \forall n \in \mathbb{N}\,(a_{n} \in \mathrm{dom}(f)) \to \mathrm{Seq}(f \circ a) \land \forall n \in \mathbb{N}\,(f \circ a(n) = f(a_{n}))$
\item $\varphi_{118} :\equiv \forall f\,\forall A\,\forall c\,\forall a\,(\mathrm{ContIn}(f,A,c) \to \mathrm{Seq}(a) \to \varphi_{40})$
\item $\varphi_{119} :\equiv \mathrm{Conv}(f \circ a \oplus g \circ a,f(c) + g(c)) \to \mathrm{Conv}((f + g) \circ a,f(c) + g(c))$
\item $\varphi_{120} :\equiv \mathrm{ContIn}(f,A,c) \to \mathrm{ContIn}(g,A,c) \to \mathrm{Seq}(a) \to \varphi_{35}$
\item $\varphi_{121} :\equiv \forall \beta\,(\mathrm{Conv}(f \circ a,f(c)) \to \mathrm{Conv}(g \circ a,\beta) \to \mathrm{Conv}(f \circ a \oplus g \circ a,f(c) + \beta))$
\item $\varphi_{122} :\equiv \mathrm{ContIn}(f,A,c) \to \mathrm{ContIn}(g,A,c) \to \mathrm{Seq}(a) \to \varphi_{36}$
\item $\varphi_{123} :\equiv \mathrm{ContIn}(f,A,c) \to \mathrm{ContIn}(g,A,c) \to \mathrm{Seq}(a) \to \varphi_{39}$
\item $\varphi_{124} :\equiv \mathrm{ContIn}(f,A,c) \to \mathrm{ContIn}(g,A,c) \to \mathrm{Seq}(a) \to \varphi_{40}$
\item $\varphi_{125} :\equiv \mathrm{ContIn}(g,A,c) \to \mathrm{Seq}(a) \to \forall n \in \mathbb{N}\,(a_{n} \in A) \to \varphi_{19}$
\item $\varphi_{126} :\equiv \forall z\,((f + g) \circ a(n) = (f + g)(a_{n}) \to (f + g)(a_{n}) = z \to (f + g) \circ a(n) = z)$
\item $\varphi_{127} :\equiv \mathrm{ContIn}(g,A,c) \to \forall a\,(\mathrm{Seq}(a) \to \forall n \in \mathbb{N}\,(a_{n} \in A) \to \varphi_{19})$
\item $\varphi_{128} :\equiv \mathrm{ContIn}(f,A,c) \to \mathrm{ContIn}(g,A,c) \to \mathrm{Seq}(a) \to \varphi_{52}$
\item $\varphi_{129} :\equiv \mathrm{ContIn}(f,A,c) \to \mathrm{ContIn}(g,A,c) \to \mathrm{Seq}(a) \to \varphi_{53}$
\item $\varphi_{130} :\equiv \mathrm{ContIn}(g,A,c) \to \mathrm{Seq}(a) \to \forall n \in \mathbb{N}\,(a_{n} \in A) \to \varphi_{22}$
\item $\varphi_{131} :\equiv \mathrm{ContIn}(f,A,c) \to \mathrm{ContIn}(g,A,c) \to \mathrm{Seq}(a) \to \varphi_{59}$
\item $\varphi_{132} :\equiv \mathrm{ContIn}(f,A,c) \to \mathrm{ContIn}(g,A,c) \to \mathrm{Seq}(a) \to \varphi_{60}$
\item $\varphi_{133} :\equiv f(a_{n}) + g \circ a(n) = f(a_{n}) + z \to f(a_{n}) + g \circ a(n) = f(a_{n}) + w$
\item $\varphi_{134} :\equiv \mathrm{ContIn}(f,A,c) \to \mathrm{ContIn}(g,A,c) \to \mathrm{Seq}(a) \to \varphi_{62}$
\item $\varphi_{135} :\equiv \mathrm{ContIn}(f,A,c) \to \mathrm{ContIn}(g,A,c) \to \mathrm{Seq}(a) \to \varphi_{63}$
\item $\varphi_{136} :\equiv (f + g)(a_{n}) = f(a_{n}) + g(a_{n}) \to (f + g) \circ a(n) = f(a_{n}) + g(a_{n})$
\item $\varphi_{137} :\equiv \forall a\,(\mathrm{Seq}(a) \to \forall n \in \mathbb{N}\,(a_{n} \in A) \to \varphi_{19}) \to \mathrm{ContIn}((f + g),A,c)$
\item $\varphi_{138} :\equiv \mathrm{ContIn}(f,A,c) \to \mathrm{ContIn}(g,A,c) \to \mathrm{Seq}(a) \to \varphi_{74}$
\item $\varphi_{139} :\equiv \mathrm{ContIn}(f,A,c) \to \mathrm{ContIn}(g,A,c) \to \mathrm{Seq}(a) \to \varphi_{79}$
\item $\varphi_{140} :\equiv \mathrm{ContIn}(f,A,c) \land \mathrm{ContIn}(g,A,c) \to \mathrm{Seq}(a) \to \forall n \in \mathbb{N}\,(a_{n} \in A) \to \varphi_{19}$
\item $\varphi_{141} :\equiv \mathrm{ContIn}(g,A,c) \to \mathrm{Seq}(a) \to \forall n \in \mathbb{N}\,(a_{n} \in A) \to \varphi_{41}$
\item $\varphi_{142} :\equiv \mathrm{ContIn}(g,A,c) \to \mathrm{Seq}(a) \to \forall n \in \mathbb{N}\,(a_{n} \in A) \to \varphi_{42}$
\item $\varphi_{143} :\equiv \mathrm{ContIn}(g,A,c) \to \mathrm{Seq}(a) \to \forall n \in \mathbb{N}\,(a_{n} \in A) \to \varphi_{43}$
\item $\varphi_{144} :\equiv \mathrm{ContIn}(g,A,c) \to \mathrm{Seq}(a) \to \forall n \in \mathbb{N}\,(a_{n} \in A) \to \varphi_{50}$
\item $\varphi_{145} :\equiv \mathrm{ContIn}(g,A,c) \to \mathrm{Seq}(a) \to \forall n \in \mathbb{N}\,(a_{n} \in A) \to \varphi_{51}$
\item $\varphi_{146} :\equiv \mathrm{ContIn}(f,A,c) \to \mathrm{ContIn}(g,A,c) \to \mathrm{Seq}(a) \to \varphi_{84}$
\item $\varphi_{147} :\equiv \mathrm{ContIn}(f,A,c) \to \mathrm{ContIn}(g,A,c) \to \mathrm{Seq}(a) \to \varphi_{85}$
\item $\varphi_{148} :\equiv \forall n \in \mathbb{N}\,((f + g) \circ a(n) = (f + g)(a_{n})) \to n \in \mathbb{N} \to (f + g) \circ a(n) = (f + g)(a_{n})$
\item $\varphi_{149} :\equiv (\mathrm{ContIn}(f,A,c) \land \mathrm{ContIn}(g,A,c) \to \varphi_{2}) \to \mathrm{ContIn}(f,A,c) \to \varphi_{13}$
\item $\varphi_{150} :\equiv \mathrm{ContIn}(g,A,c) \to \mathrm{Seq}(a) \to \forall n \in \mathbb{N}\,(a_{n} \in A) \to \varphi_{54}$
\item $\varphi_{151} :\equiv \forall z\,((f + g) \circ a(n) = f(a_{n}) + g(a_{n}) \to f(a_{n}) + g(a_{n}) = z \to (f + g) \circ a(n) = z)$
\item $\varphi_{152} :\equiv \forall b\,(\mathrm{Seq}(f \circ a) \to \mathrm{Seq}(b) \to \mathrm{Seq}(f \circ a \oplus b) \land \forall n \in \mathbb{N}\,(f \circ a \oplus b(n) = f \circ a(n) + b_{n}))$
\item $\varphi_{153} :\equiv \forall b\,\forall \alpha\,(\mathrm{Seq}(b) \to \forall n \in \mathbb{N}\,(f \circ a \oplus g \circ a(n) = b_{n}) \to \mathrm{Conv}(f \circ a \oplus g \circ a,\alpha) \to \mathrm{Conv}(b,\alpha))$
\item $\varphi_{154} :\equiv (f + g) \circ a(n) = f \circ a \oplus g \circ a(n) \to f \circ a \oplus g \circ a(n) = (f + g) \circ a(n)$
\item $\varphi_{155} :\equiv f \circ a(n) + g \circ a(n) = z + g \circ a(n) \to f \circ a(n) + g \circ a(n) = w + g \circ a(n)$
\item $\varphi_{156} :\equiv \varphi_{6} \to \mathrm{ContIn}(f,A,c) \to \mathrm{ContIn}(g,A,c) \to \varphi_{31}$
\item $\varphi_{157} :\equiv \varphi_{6} \to \mathrm{ContIn}(f,A,c) \to \mathrm{ContIn}(g,A,c) \to \varphi_{32}$
\item $\varphi_{158} :\equiv \mathrm{ContIn}(g,A,c) \to \mathrm{Seq}(a) \to \forall n \in \mathbb{N}\,(a_{n} \in A) \to \varphi_{70}$
\item $\varphi_{159} :\equiv \mathrm{ContIn}(g,A,c) \to \mathrm{Seq}(a) \to \forall n \in \mathbb{N}\,(a_{n} \in A) \to \varphi_{71}$
\item $\varphi_{160} :\equiv \mathrm{ContIn}(g,A,c) \to \mathrm{Seq}(a) \to \forall n \in \mathbb{N}\,(a_{n} \in A) \to \varphi_{72}$
\item $\varphi_{161} :\equiv \forall f\,\forall A\,\forall c\,(\mathrm{RFn}(f) \to A \subseteq \mathrm{dom}(f) \to c \in A \to \forall a\,(\mathrm{Seq}(a) \to \varphi_{40}) \to \mathrm{ContIn}(f,A,c))$
\item $\varphi_{162} :\equiv \mathrm{ContIn}(g,A,c) \to \mathrm{Seq}(a) \to \forall n \in \mathbb{N}\,(a_{n} \in A) \to \varphi_{80}$
\item $\varphi_{163} :\equiv f(a_{n}) + g(a_{n}) = f \circ a \oplus g \circ a(n) \to (f + g) \circ a(n) = f \circ a \oplus g \circ a(n)$
\item $\varphi_{164} :\equiv \mathrm{ContIn}(g,A,c) \to \mathrm{Seq}(a) \to \forall n \in \mathbb{N}\,(a_{n} \in A) \to \varphi_{81}$
\item $\varphi_{165} :\equiv \mathrm{ContIn}(g,A,c) \to \mathrm{Seq}(a) \to \forall n \in \mathbb{N}\,(a_{n} \in A) \to \mathrm{Conv}(a,c) \to \varphi_{61}$
\item $\varphi_{166} :\equiv \mathrm{ContIn}(g,A,c) \to \mathrm{Seq}(a) \to \forall n \in \mathbb{N}\,(a_{n} \in A) \to \varphi_{86}$
\item $\varphi_{167} :\equiv \mathrm{ContIn}(g,A,c) \to \mathrm{Seq}(a) \to \forall n \in \mathbb{N}\,(a_{n} \in A) \to \varphi_{87}$
\item $\varphi_{168} :\equiv f(a_{n}) + g \circ a(n) = f(a_{n}) + g(a_{n}) \to f \circ a \oplus g \circ a(n) = f(a_{n}) + g(a_{n})$
\item $\varphi_{169} :\equiv f \circ a \oplus g \circ a(n) = f(a_{n}) + g(a_{n}) \to f(a_{n}) + g(a_{n}) = f \circ a \oplus g \circ a(n)$
\item $\varphi_{170} :\equiv \mathrm{ContIn}(g,A,c) \to \mathrm{Seq}(a) \to \forall n \in \mathbb{N}\,(a_{n} \in A) \to \varphi_{94}$
\item $\varphi_{171} :\equiv f(a_{n}) + g \circ a(n) = f(a_{n}) + g \circ a(n) \to f(a_{n}) + g \circ a(n) = f(a_{n}) + g(a_{n})$
\item $\varphi_{172} :\equiv \mathrm{ContIn}(g,A,c) \to \mathrm{Seq}(a) \to \forall n \in \mathbb{N}\,(a_{n} \in A) \to \varphi_{101}$
\item $\varphi_{173} :\equiv \mathrm{ContIn}(g,A,c) \to \mathrm{Seq}(a) \to \forall n \in \mathbb{N}\,(a_{n} \in A) \to \varphi_{103}$
\item $\varphi_{174} :\equiv \mathrm{ContIn}(g,A,c) \to \mathrm{Seq}(a) \to \forall n \in \mathbb{N}\,(a_{n} \in A) \to \varphi_{104}$
\item $\varphi_{175} :\equiv \forall g\,(\mathrm{RFn}(f) \to \mathrm{RFn}(g) \to \mathrm{RFn}((f + g)) \land \varphi_{16} \land \varphi_{27})$
\item $\varphi_{176} :\equiv \mathrm{ContIn}(g,A,c) \to \mathrm{Seq}(a) \to \forall n \in \mathbb{N}\,(a_{n} \in A) \to \varphi_{107}$
\item $\varphi_{177} :\equiv (\forall n \in \mathbb{N}\,(a_{n} \in A) \to n \in \mathbb{N} \to a_{n} \in A) \to \mathrm{ContIn}(f,A,c) \to \varphi_{73}$
\item $\varphi_{178} :\equiv f \circ a(n) + g \circ a(n) = f(a_{n}) + g \circ a(n) \to f \circ a \oplus g \circ a(n) = f(a_{n}) + g \circ a(n)$
\item $\varphi_{179} :\equiv \mathrm{ContIn}(g,A,c) \to \mathrm{Seq}(a) \to \forall n \in \mathbb{N}\,(a_{n} \in A) \to \mathrm{Conv}(a,c) \to \varphi_{90}$
\item $\varphi_{180} :\equiv \mathrm{Seq}(a) \to \forall n \in \mathbb{N}\,(a_{n} \in A) \to \mathrm{Conv}(a,c) \to n \in \mathbb{N} \to \varphi_{102}$
\item $\varphi_{181} :\equiv \mathrm{ContIn}(g,A,c) \to \mathrm{Seq}(a) \to \forall n \in \mathbb{N}\,(a_{n} \in A) \to \varphi_{115}$
\item $\varphi_{182} :\equiv \mathrm{ContIn}(g,A,c) \to \mathrm{Seq}(a) \to \forall n \in \mathbb{N}\,(a_{n} \in A) \to \mathrm{Conv}(a,c) \to \varphi_{93}$
\item $\varphi_{183} :\equiv \mathrm{RFn}((f + g)) \to \mathrm{Seq}(a) \to \forall n \in \mathbb{N}\,(a_{n} \in \mathrm{dom}((f + g))) \to \varphi_{93}$
\item $\varphi_{184} :\equiv \forall w\,(g \circ a(n) = w \to f(a_{n}) + g \circ a(n) = f(a_{n}) + g \circ a(n) \to f(a_{n}) + g \circ a(n) = f(a_{n}) + w)$
\item $\varphi_{185} :\equiv \forall n\,(\varphi_{69} \to n \in \mathbb{N} \to f \circ a \oplus g \circ a(n) = (f + g) \circ a(n))$
\item $\varphi_{186} :\equiv \forall z\,(f \circ a \oplus g \circ a(n) = f(a_{n}) + g \circ a(n) \to f(a_{n}) + g \circ a(n) = z \to f \circ a \oplus g \circ a(n) = z)$
\item $\varphi_{187} :\equiv \mathrm{ContIn}(f,A,c) \to \mathrm{ContIn}(g,A,c) \to \mathrm{RFn}((f + g)) \land \varphi_{16} \land \varphi_{27}$
\item $\varphi_{188} :\equiv f \circ a(n) + g \circ a(n) = f \circ a(n) + g \circ a(n) \to f \circ a(n) + g \circ a(n) = f(a_{n}) + g \circ a(n)$
\item $\varphi_{189} :\equiv \forall n\,(\varphi_{69}) \to \forall n \in \mathbb{N}\,(f \circ a \oplus g \circ a(n) = (f + g) \circ a(n))$
\item $\varphi_{190} :\equiv \mathrm{ContIn}(g,A,c) \leftrightarrow \mathrm{RFn}(g) \land A \subseteq \mathrm{dom}(g) \land c \in A \land \forall \varepsilon \in \mathbb{R}\,(0 < \varepsilon \to \exists \delta\,(\delta \in \mathbb{R} \land 0 < \delta \land \forall x \in A\,(\lvert x - c \rvert < \delta \to \lvert g(x) - g(c) \rvert < \varepsilon)))$
\item $\varphi_{191} :\equiv \mathrm{ContIn}(f,A,c) \leftrightarrow \mathrm{RFn}(f) \land A \subseteq \mathrm{dom}(f) \land c \in A \land \forall \varepsilon \in \mathbb{R}\,(0 < \varepsilon \to \exists \delta\,(\delta \in \mathbb{R} \land 0 < \delta \land \forall x \in A\,(\lvert x - c \rvert < \delta \to \lvert f(x) - f(c) \rvert < \varepsilon)))$
\item $\varphi_{192} :\equiv \mathrm{ContIn}(g,A,c) \to \mathrm{Seq}(a) \to \forall n \in \mathbb{N}\,(a_{n} \in A) \to \mathrm{Conv}(a,c) \to \varphi_{114}$
\item $\varphi_{193} :\equiv \forall z\,(f \circ a \oplus g \circ a(n) = f \circ a(n) + g \circ a(n) \to f \circ a(n) + g \circ a(n) = z \to f \circ a \oplus g \circ a(n) = z)$
\item $\varphi_{194} :\equiv \mathrm{Seq}(f \circ a) \to \mathrm{Seq}(g \circ a) \to \mathrm{Seq}(f \circ a \oplus g \circ a) \land \varphi_{102}$
\item $\varphi_{195} :\equiv \forall c\,(\mathrm{RFn}((f + g)) \to A \subseteq \mathrm{dom}((f + g)) \to c \in A \to \varphi_{137})$
\item $\varphi_{196} :\equiv \forall n \in \mathbb{N}\,(a_{n} \in A) \to \mathrm{Conv}(a,c) \to \mathrm{Seq}(f \circ a \oplus g \circ a) \land \varphi_{102}$
\item $\varphi_{197} :\equiv \forall w\,(f \circ a(n) = w \to f \circ a(n) + g \circ a(n) = f \circ a(n) + g \circ a(n) \to f \circ a(n) + g \circ a(n) = w + g \circ a(n))$
\item $\varphi_{198} :\equiv (\mathrm{ContIn}(f,A,c) \to \varphi_{127}) \to \mathrm{ContIn}(f,A,c) \to \varphi_{3}$
\item $\varphi_{199} :\equiv \forall \alpha\,(\mathrm{Seq}((f + g) \circ a) \to \forall n \in \mathbb{N}\,(f \circ a \oplus g \circ a(n) = (f + g) \circ a(n)) \to \mathrm{Conv}(f \circ a \oplus g \circ a,\alpha) \to \mathrm{Conv}((f + g) \circ a,\alpha))$
\item $\varphi_{200} :\equiv \varphi_{102} \to n \in \mathbb{N} \to f \circ a \oplus g \circ a(n) = f \circ a(n) + g \circ a(n)$
\item $\varphi_{201} :\equiv \forall a\,(\mathrm{ContIn}(g,A,c) \to \mathrm{Seq}(a) \to \varphi_{39}) \to \mathrm{ContIn}(g,A,c) \to \mathrm{Seq}(a) \to \varphi_{39}$
\item $\varphi_{202} :\equiv \forall a\,(\mathrm{ContIn}(f,A,c) \to \mathrm{Seq}(a) \to \varphi_{40}) \to \mathrm{ContIn}(f,A,c) \to \mathrm{Seq}(a) \to \varphi_{40}$
\item $\varphi_{203} :\equiv (\mathrm{ContIn}(f,A,c) \land \mathrm{ContIn}(g,A,c) \to \varphi_{1}) \to \mathrm{ContIn}(f,A,c) \to \varphi_{127}$
\item $\varphi_{204} :\equiv \forall c\,\forall a\,(\mathrm{ContIn}(g,A,c) \to \mathrm{Seq}(a) \to \varphi_{39}) \to \forall a\,(\mathrm{ContIn}(g,A,c) \to \mathrm{Seq}(a) \to \varphi_{39})$
\item $\varphi_{205} :\equiv \forall c\,\forall a\,(\mathrm{ContIn}(f,A,c) \to \mathrm{Seq}(a) \to \varphi_{40}) \to \forall a\,(\mathrm{ContIn}(f,A,c) \to \mathrm{Seq}(a) \to \varphi_{40})$
\item $\varphi_{206} :\equiv \forall y\,((f + g) \circ a(n) = y \to y = (f + g) \circ a(n)) \to \varphi_{154}$
\item $\varphi_{207} :\equiv \forall A\,\forall c\,\forall a\,(\mathrm{ContIn}(g,A,c) \to \mathrm{Seq}(a) \to \varphi_{39}) \to \forall c\,\forall a\,(\mathrm{ContIn}(g,A,c) \to \mathrm{Seq}(a) \to \varphi_{39})$
\item $\varphi_{208} :\equiv \forall A\,\forall c\,\forall a\,(\mathrm{ContIn}(f,A,c) \to \mathrm{Seq}(a) \to \varphi_{40}) \to \forall c\,\forall a\,(\mathrm{ContIn}(f,A,c) \to \mathrm{Seq}(a) \to \varphi_{40})$
\item $\varphi_{209} :\equiv \varphi_{118} \to \forall A\,\forall c\,\forall a\,(\mathrm{ContIn}(g,A,c) \to \mathrm{Seq}(a) \to \varphi_{39})$
\item $\varphi_{210} :\equiv \varphi_{118} \to \forall A\,\forall c\,\forall a\,(\mathrm{ContIn}(f,A,c) \to \mathrm{Seq}(a) \to \varphi_{40})$
\item $\varphi_{211} :\equiv \varphi_{123} \to \mathrm{ContIn}(f,A,c) \to \mathrm{ContIn}(g,A,c) \to \varphi_{37}$
\item $\varphi_{212} :\equiv \varphi_{124} \to \mathrm{ContIn}(f,A,c) \to \mathrm{ContIn}(g,A,c) \to \varphi_{38}$
\item $\varphi_{213} :\equiv (\varphi_{27} \to \varphi_{24}) \to \mathrm{ContIn}(f,A,c) \to \mathrm{ContIn}(g,A,c) \to \varphi_{24}$
\item $\varphi_{214} :\equiv \mathrm{ContIn}(f,A,c) \to \mathrm{ContIn}(g,A,c) \to \mathrm{Seq}(a) \to \varphi_{196}$
\item $\varphi_{215} :\equiv \mathrm{Seq}((f + g) \circ a) \to \forall n \in \mathbb{N}\,(f \circ a \oplus g \circ a(n) = (f + g) \circ a(n)) \to \varphi_{119}$
\item $\varphi_{216} :\equiv (\mathrm{ContIn}(f,A,c) \to \mathrm{ContIn}(g,A,c) \to \varphi_{44}) \to \varphi_{138}$
\item $\varphi_{217} :\equiv (\varphi_{16} \to \varphi_{44}) \to \mathrm{ContIn}(f,A,c) \to \mathrm{ContIn}(g,A,c) \to \varphi_{44}$
\item $\varphi_{218} :\equiv \forall y\,(f \circ a \oplus g \circ a(n) = y \to y = f \circ a \oplus g \circ a(n)) \to \varphi_{169}$
\item $\varphi_{219} :\equiv \forall n\,(\varphi_{69} \to n \in \mathbb{N} \to a_{n} \in \mathrm{dom}(g)) \to \forall n\,(\varphi_{69}) \to \forall n \in \mathbb{N}\,(a_{n} \in \mathrm{dom}(g))$
\item $\varphi_{220} :\equiv \forall n\,(\varphi_{69} \to n \in \mathbb{N} \to a_{n} \in \mathrm{dom}(f)) \to \forall n\,(\varphi_{69}) \to \forall n \in \mathbb{N}\,(a_{n} \in \mathrm{dom}(f))$
\item $\varphi_{221} :\equiv (\mathrm{ContIn}(f,A,c) \to \mathrm{ContIn}(f,A,c)) \to \varphi_{191} \to \varphi_{5}$
\item $\varphi_{222} :\equiv (\mathrm{ContIn}(f,A,c) \to \mathrm{ContIn}(f,A,c)) \to \varphi_{191} \to \varphi_{6}$
\item $\varphi_{223} :\equiv (\mathrm{ContIn}(f,A,c) \to \mathrm{ContIn}(f,A,c)) \to \varphi_{191} \to \varphi_{9}$
\item $\varphi_{224} :\equiv \forall a\,(\mathrm{RFn}(g) \to \mathrm{Seq}(a) \to \varphi_{116}) \to \mathrm{RFn}(g) \to \mathrm{Seq}(a) \to \varphi_{116}$
\item $\varphi_{225} :\equiv \forall a\,(\mathrm{RFn}(f) \to \mathrm{Seq}(a) \to \varphi_{117}) \to \mathrm{RFn}(f) \to \mathrm{Seq}(a) \to \varphi_{117}$
\item $\varphi_{226} :\equiv \varphi_{187} \to \mathrm{ContIn}(f,A,c) \to \mathrm{ContIn}(g,A,c) \to \varphi_{16}$
\item $\varphi_{227} :\equiv \forall n\,(\varphi_{69} \to n \in \mathbb{N} \to a_{n} \in \mathrm{dom}((f + g))) \to \forall n\,(\varphi_{69}) \to \forall n \in \mathbb{N}\,(a_{n} \in \mathrm{dom}((f + g)))$
\item $\varphi_{228} :\equiv \forall f\,\forall a\,(\mathrm{RFn}(f) \to \mathrm{Seq}(a) \to \varphi_{117}) \to \forall a\,(\mathrm{RFn}(g) \to \mathrm{Seq}(a) \to \varphi_{116})$
\item $\varphi_{229} :\equiv \forall f\,\forall a\,(\mathrm{RFn}(f) \to \mathrm{Seq}(a) \to \varphi_{117}) \to \forall a\,(\mathrm{RFn}(f) \to \mathrm{Seq}(a) \to \varphi_{117})$
\item $\varphi_{230} :\equiv (f(c) + g(c) = (f + g)(c) \to \varphi_{89}) \to \mathrm{ContIn}(f,A,c) \to \varphi_{125}$
\item $\varphi_{231} :\equiv \varphi_{187} \to \mathrm{ContIn}(f,A,c) \to \mathrm{ContIn}(g,A,c) \to \varphi_{27}$
\item $\varphi_{232} :\equiv \forall a\,(\varphi_{140}) \to \varphi_{1} \to \forall a\,(\mathrm{Seq}(a) \to \forall n \in \mathbb{N}\,(a_{n} \in A) \to \varphi_{19})$
\item $\varphi_{233} :\equiv (\varphi_{27} \to \varphi_{82}) \to \mathrm{ContIn}(f,A,c) \to \mathrm{ContIn}(g,A,c) \to \varphi_{82}$
\item $\varphi_{234} :\equiv (\mathrm{ContIn}(f,A,c) \to \varphi_{141}) \to \mathrm{ContIn}(f,A,c) \to \varphi_{130}$
\item $\varphi_{235} :\equiv (\mathrm{ContIn}(f,A,c) \to \mathrm{ContIn}(g,A,c) \to \varphi_{16}) \to \varphi_{112} \to \mathrm{ContIn}(f,A,c) \to \varphi_{33}$
\item $\varphi_{236} :\equiv (\mathrm{ContIn}(f,A,c) \to \mathrm{ContIn}(g,A,c) \to \varphi_{16}) \to \varphi_{113} \to \mathrm{ContIn}(f,A,c) \to \varphi_{34}$
\item $\varphi_{237} :\equiv (\mathrm{ContIn}(f,A,c) \to \mathrm{ContIn}(g,A,c) \to \varphi_{31}) \to (\mathrm{ContIn}(f,A,c) \to \varphi_{33}) \to \varphi_{108}$
\item $\varphi_{238} :\equiv (\mathrm{ContIn}(f,A,c) \to \varphi_{158}) \to \mathrm{ContIn}(f,A,c) \to \varphi_{144}$
\item $\varphi_{239} :\equiv (\mathrm{ContIn}(f,A,c) \to \varphi_{159}) \to \mathrm{ContIn}(f,A,c) \to \varphi_{145}$
\item $\varphi_{240} :\equiv (\mathrm{ContIn}(f,A,c) \to \mathrm{ContIn}(g,A,c) \to \varphi_{27}) \to \varphi_{213}$
\item $\varphi_{241} :\equiv (\mathrm{ContIn}(f,A,c) \to \mathrm{ContIn}(g,A,c) \to \varphi_{16}) \to \varphi_{217}$
\item $\varphi_{242} :\equiv (\mathrm{ContIn}(f,A,c) \to \mathrm{ContIn}(g,A,c) \to \varphi_{24}) \to \varphi_{108} \to \varphi_{134}$
\item $\varphi_{243} :\equiv (\mathrm{ContIn}(f,A,c) \to \mathrm{ContIn}(g,A,c) \to \varphi_{37}) \to \varphi_{214}$
\item $\varphi_{244} :\equiv (\mathrm{ContIn}(f,A,c) \to \varphi_{167}) \to \varphi_{69} \to n \in \mathbb{N} \to f \circ a \oplus g \circ a(n) = (f + g) \circ a(n)$
\item $\varphi_{245} :\equiv (A \subseteq \mathrm{dom}(g) \to a_{n} \in A \to a_{n} \in \mathrm{dom}(g)) \to \varphi_{8} \to \varphi_{122} \to \varphi_{128}$
\item $\varphi_{246} :\equiv (A \subseteq \mathrm{dom}(f) \to a_{n} \in A \to a_{n} \in \mathrm{dom}(f)) \to \varphi_{9} \to \varphi_{122} \to \varphi_{129}$
\item $\varphi_{247} :\equiv ((f + g) \circ a(n) = (f + g)(a_{n}) \to \varphi_{136}) \to \mathrm{ContIn}(f,A,c) \to \varphi_{162}$
\item $\varphi_{248} :\equiv \forall u\,(\varphi_{26}) \to (\forall u\,(\varphi_{26}) \to \varphi_{2} \to \forall u \in A\,(u \in \mathrm{dom}((f + g)))) \to \varphi_{149}$
\item $\varphi_{249} :\equiv (\mathrm{ContIn}(f,A,c) \to \varphi_{182}) \to \mathrm{ContIn}(f,A,c) \to \varphi_{165}$
\item $\varphi_{250} :\equiv \varphi_{195} \to \mathrm{RFn}((f + g)) \to A \subseteq \mathrm{dom}((f + g)) \to c \in A \to \varphi_{137}$
\item $\varphi_{251} :\equiv (g \circ a(n) = g(a_{n}) \to \varphi_{171}) \to \mathrm{ContIn}(f,A,c) \to \varphi_{172}$
\item $\varphi_{252} :\equiv (\mathrm{ContIn}(f,A,c) \to \mathrm{ContIn}(g,A,c) \to \varphi_{27}) \to \varphi_{233}$
\item $\varphi_{253} :\equiv ((f + g) \circ a(n) = f(a_{n}) + g(a_{n}) \to \varphi_{163}) \to \mathrm{ContIn}(f,A,c) \to \varphi_{166}$
\item $\varphi_{254} :\equiv (f \circ a(n) = f(a_{n}) \to \varphi_{188}) \to \mathrm{ContIn}(f,A,c) \to \varphi_{179}$
\item $\varphi_{255} :\equiv (\mathrm{ContIn}(f,A,c) \to \mathrm{ContIn}(g,A,c) \to \varphi_{32}) \to \varphi_{237}$
\item $\varphi_{256} :\equiv (\mathrm{ContIn}(f,A,c) \to \mathrm{ContIn}(f,A,c)) \to \varphi_{11} \to (\mathrm{ContIn}(f,A,c) \to \varphi_{78}) \to \varphi_{91} \to \varphi_{124}$
\item $\varphi_{257} :\equiv (\mathrm{ContIn}(f,A,c) \to \varphi_{144}) \to \varphi_{105} \to \mathrm{ContIn}(f,A,c) \to \varphi_{142}$
\item $\varphi_{258} :\equiv (\mathrm{ContIn}(f,A,c) \to \varphi_{145}) \to \varphi_{106} \to \mathrm{ContIn}(f,A,c) \to \varphi_{143}$
\item $\varphi_{259} :\equiv (f \circ a \oplus g \circ a(n) = f(a_{n}) + g \circ a(n) \to \varphi_{168}) \to \mathrm{ContIn}(f,A,c) \to \varphi_{173}$
\item $\varphi_{260} :\equiv (\mathrm{RFn}(f) \to \mathrm{RFn}(g) \to \mathrm{RFn}((f + g)) \land \varphi_{16} \land \varphi_{27}) \to \varphi_{5} \to \varphi_{4} \to \varphi_{187}$
\item $\varphi_{261} :\equiv \varphi_{214} \to \mathrm{ContIn}(f,A,c) \to \mathrm{ContIn}(g,A,c) \to \varphi_{180}$
\item $\varphi_{262} :\equiv (\mathrm{ContIn}(f,A,c) \to \varphi_{164}) \to \varphi_{111} \to \mathrm{ContIn}(f,A,c) \to \varphi_{150}$
\item $\varphi_{263} :\equiv (f \circ a \oplus g \circ a(n) = f \circ a(n) + g \circ a(n) \to \varphi_{178}) \to \mathrm{ContIn}(f,A,c) \to \varphi_{176}$
\item $\varphi_{264} :\equiv (\mathrm{ContIn}(f,A,c) \to \mathrm{ContIn}(g,A,c) \to \varphi_{38}) \to \varphi_{243}$
\item $\varphi_{265} :\equiv (\mathrm{ContIn}(f,A,c) \to \varphi_{165}) \to \varphi_{148} \to \mathrm{ContIn}(f,A,c) \to \varphi_{164}$
\item $\varphi_{266} :\equiv (\mathrm{ContIn}(f,A,c) \to \varphi_{166}) \to \varphi_{154} \to \mathrm{ContIn}(f,A,c) \to \varphi_{167}$
\item $\varphi_{267} :\equiv (\mathrm{ContIn}(f,A,c) \to \varphi_{192}) \to \varphi_{111} \to \mathrm{ContIn}(f,A,c) \to \varphi_{181}$
\item $\varphi_{268} :\equiv (\mathrm{Conv}(f \circ a,f(c)) \to \varphi_{83}) \to \varphi_{124} \to \varphi_{123} \to \mathrm{ContIn}(f,A,c) \to \varphi_{141}$
\item $\varphi_{269} :\equiv (\mathrm{RFn}(g) \to \mathrm{Seq}(a) \to \varphi_{116}) \to \varphi_{4} \to \varphi_{11} \to \varphi_{131} \to \mathrm{ContIn}(f,A,c) \to \varphi_{158}$
\item $\varphi_{270} :\equiv (\mathrm{RFn}(f) \to \mathrm{Seq}(a) \to \varphi_{117}) \to \varphi_{5} \to \varphi_{11} \to \varphi_{132} \to \mathrm{ContIn}(f,A,c) \to \varphi_{159}$
\item $\varphi_{271} :\equiv (\mathrm{ContIn}(f,A,c) \to \varphi_{173}) \to \varphi_{169} \to \mathrm{ContIn}(f,A,c) \to \varphi_{174}$
\item $\varphi_{272} :\equiv (\mathrm{ContIn}(g,A,c) \to \mathrm{Seq}(a) \to \varphi_{39}) \to \varphi_{7} \to \varphi_{11} \to (\mathrm{ContIn}(f,A,c) \to \varphi_{78}) \to \varphi_{91} \to \varphi_{123}$
\item $\varphi_{273} :\equiv (\mathrm{RFn}((f + g)) \to A \subseteq \mathrm{dom}((f + g)) \to c \in A \to \varphi_{137}) \to \varphi_{10} \to \varphi_{14} \to \varphi_{6} \to \varphi_{198}$
\item $\varphi_{274} :\equiv (\mathrm{ContIn}(f,A,c) \to \mathrm{ContIn}(g,A,c) \to \varphi_{82}) \to \varphi_{138} \to \varphi_{111} \to \mathrm{ContIn}(f,A,c) \to \varphi_{160}$
\item $\varphi_{275} :\equiv (\mathrm{ContIn}(f,A,c) \to \mathrm{ContIn}(g,A,c) \to \varphi_{180}) \to \varphi_{200} \to \mathrm{ContIn}(f,A,c) \to \varphi_{192}$
\item $\varphi_{276} :\equiv \varphi_{147} \to f \circ a(n) + g \circ a(n) = f \circ a(n) + g \circ a(n) \to \varphi_{254}$
\item $\varphi_{277} :\equiv (\mathrm{ContIn}(f,A,c) \to \varphi_{150}) \to (\mathrm{ContIn}(f,A,c) \to \varphi_{160}) \to \varphi_{247}$
\item $\varphi_{278} :\equiv \forall n\,(\varphi_{69} \to n \in \mathbb{N} \to a_{n} \in \mathrm{dom}(g)) \to \varphi_{219} \to (\varphi_{69} \to \forall n\,(\varphi_{69})) \to \varphi_{131}$
\item $\varphi_{279} :\equiv \forall n\,(\varphi_{69} \to n \in \mathbb{N} \to a_{n} \in \mathrm{dom}(f)) \to \varphi_{220} \to (\varphi_{69} \to \forall n\,(\varphi_{69})) \to \varphi_{132}$
\item $\varphi_{280} :\equiv (\varphi_{185} \to \varphi_{189}) \to (\varphi_{69} \to \forall n\,(\varphi_{69})) \to \mathrm{ContIn}(f,A,c) \to \varphi_{170}$
\item $\varphi_{281} :\equiv \forall n\,(\varphi_{69} \to n \in \mathbb{N} \to a_{n} \in \mathrm{dom}((f + g))) \to \varphi_{227} \to (\varphi_{69} \to \forall n\,(\varphi_{69})) \to \varphi_{139}$
\item $\varphi_{282} :\equiv (\mathrm{ContIn}(f,A,c) \to \varphi_{162}) \to (\mathrm{ContIn}(f,A,c) \to \varphi_{174}) \to \varphi_{253}$
\item $\varphi_{283} :\equiv (\mathrm{ContIn}(f,A,c) \to \varphi_{176}) \to (\mathrm{ContIn}(f,A,c) \to \varphi_{172}) \to \varphi_{259}$
\item $\varphi_{284} :\equiv (\mathrm{ContIn}(f,A,c) \to \varphi_{181}) \to (\mathrm{ContIn}(f,A,c) \to \varphi_{179}) \to \varphi_{263}$
\end{abbrevlist}
\end{fullproof}

% 自動生成 (prover/build.js): ContInMul  ノード数 567
\begin{fullproof}{定理2.3(積)}{$\forall f\,\forall g\,\forall A\,\forall c\,(\mathrm{ContIn}(f,A,c) \to \mathrm{ContIn}(g,A,c) \to \mathrm{ContIn}((f \cdot g),A,c))$}{567}\phantomsection\label{tree:ContInMul}
\noindent$\Pi_{1}$:
\begin{ptree}
\AXC{$\varphi_{7}$\lbl{A1}}
\AXC{$\varphi_{117}$\lbl{[G1]}}
\AXC{$\varphi_{209}$\lbl{A2}}
\BIC{$\forall A\,\forall c\,\forall a\,(\mathrm{ContIn}(g,A,c) \to \mathrm{Seq}(a) \to \varphi_{39})$}
\AXC{$\varphi_{207}$\lbl{A2}}
\BIC{$\forall c\,\forall a\,(\mathrm{ContIn}(g,A,c) \to \mathrm{Seq}(a) \to \varphi_{39})$}
\AXC{$\varphi_{204}$\lbl{A2}}
\BIC{$\forall a\,(\mathrm{ContIn}(g,A,c) \to \mathrm{Seq}(a) \to \varphi_{39})$}
\AXC{$\varphi_{201}$\lbl{A2}}
\BIC{$\mathrm{ContIn}(g,A,c) \to \mathrm{Seq}(a) \to \varphi_{39}$}
\AXC{$\varphi_{272}$\lbl{A1}}
\BIC{$\varphi_{7} \to \varphi_{11} \to (\mathrm{ContIn}(f,A,c) \to \varphi_{78}) \to \varphi_{91} \to \varphi_{122}$}
\BIC{$\varphi_{11} \to (\mathrm{ContIn}(f,A,c) \to \varphi_{78}) \to \varphi_{91} \to \varphi_{122}$}
\end{ptree}
\noindent$\Pi_{2}$:
\begin{ptree}
\AXC{$\varphi_{91}$\lbl{A1}}
\AXC{$\mathrm{ContIn}(f,A,c) \to \varphi_{78}$\lbl{A1}}
\AXC{$\varphi_{11}$\lbl{A1}}
\AXC{$\stackrel{\Pi_{1}}{\vdots}$}\noLine
\UIC{$\varphi_{11} \to (\mathrm{ContIn}(f,A,c) \to \varphi_{78}) \to \varphi_{91} \to \varphi_{122}$}
\BIC{$(\mathrm{ContIn}(f,A,c) \to \varphi_{78}) \to \varphi_{91} \to \varphi_{122}$}
\BIC{$\varphi_{91} \to \varphi_{122}$}
\BIC{$\varphi_{122}$}
\end{ptree}
\noindent$\Pi_{3}$:
\begin{ptree}
\AXC{$\varphi_{117}$\lbl{[G1]}}
\AXC{$\varphi_{210}$\lbl{A2}}
\BIC{$\forall A\,\forall c\,\forall a\,(\mathrm{ContIn}(f,A,c) \to \mathrm{Seq}(a) \to \varphi_{40})$}
\AXC{$\varphi_{208}$\lbl{A2}}
\BIC{$\forall c\,\forall a\,(\mathrm{ContIn}(f,A,c) \to \mathrm{Seq}(a) \to \varphi_{40})$}
\AXC{$\varphi_{205}$\lbl{A2}}
\BIC{$\forall a\,(\mathrm{ContIn}(f,A,c) \to \mathrm{Seq}(a) \to \varphi_{40})$}
\AXC{$\varphi_{202}$\lbl{A2}}
\BIC{$\mathrm{ContIn}(f,A,c) \to \mathrm{Seq}(a) \to \varphi_{40}$}
\end{ptree}
\noindent$\Pi_{4}$:
\begin{ptree}
\AXC{$\stackrel{\Pi_{3}}{\vdots}$}\noLine
\UIC{$\mathrm{ContIn}(f,A,c) \to \mathrm{Seq}(a) \to \varphi_{40}$}
\AXC{$(\mathrm{ContIn}(f,A,c) \to \mathrm{Seq}(a) \to \varphi_{40}) \to \varphi_{256}$\lbl{A1}}
\BIC{$\varphi_{256}$}
\end{ptree}
\noindent$\Pi_{5}$:
\begin{ptree}
\AXC{$\varphi_{91}$\lbl{A1}}
\AXC{$\mathrm{ContIn}(f,A,c) \to \varphi_{78}$\lbl{A1}}
\AXC{$\varphi_{11}$\lbl{A1}}
\AXC{$\mathrm{ContIn}(f,A,c) \to \mathrm{ContIn}(f,A,c)$\lbl{A1}}
\AXC{$\stackrel{\Pi_{4}}{\vdots}$}\noLine
\UIC{$\varphi_{256}$}
\BIC{$\varphi_{11} \to (\mathrm{ContIn}(f,A,c) \to \varphi_{78}) \to \varphi_{91} \to \varphi_{123}$}
\BIC{$(\mathrm{ContIn}(f,A,c) \to \varphi_{78}) \to \varphi_{91} \to \varphi_{123}$}
\BIC{$\varphi_{91} \to \varphi_{123}$}
\BIC{$\varphi_{123}$}
\end{ptree}
\noindent$\Pi_{6}$:
\begin{ptree}
\AXC{$\varphi_{32}$\lbl{定理2.7}}
\AXC{$\varphi_{32} \to \varphi_{76}$\lbl{A2}}
\BIC{$\varphi_{76}$}
\AXC{$\varphi_{76} \to \varphi_{88}$\lbl{A2}}
\BIC{$\varphi_{88}$}
\AXC{$\varphi_{88} \to \varphi_{124}$\lbl{A2}}
\BIC{$\varphi_{124}$}
\AXC{$\varphi_{124} \to \mathrm{Conv}(f \circ a,f(c)) \to \varphi_{83}$\lbl{A2}}
\BIC{$\mathrm{Conv}(f \circ a,f(c)) \to \varphi_{83}$}
\AXC{$\varphi_{270}$\lbl{A1}}
\BIC{$\varphi_{123} \to \varphi_{122} \to \mathrm{ContIn}(f,A,c) \to \varphi_{143}$}
\end{ptree}
\noindent$\Pi_{7}$:
\begin{ptree}
\AXC{$x = x$\lbl{A5}}
\AXC{$x = y \to x = x \to y = x$\lbl{A6}}
\AXC{$(x = y \to x = x \to y = x) \to x = x \to x = y \to y = x$\lbl{A1}}
\BIC{$x = x \to x = y \to y = x$}
\BIC{$x = y \to y = x$}
\UIC{$\forall y\,(x = y \to y = x)$}
\UIC{$\forall x\,\forall y\,(x = y \to y = x)$}
\end{ptree}
\noindent$\Pi_{8}$:
\begin{ptree}
\AXC{$y = z \to x = y \to x = z$\lbl{A6}}
\AXC{$(y = z \to x = y \to x = z) \to x = y \to y = z \to x = z$\lbl{A1}}
\BIC{$x = y \to y = z \to x = z$}
\UIC{$\forall z\,(x = y \to y = z \to x = z)$}
\UIC{$\forall y\,\forall z\,(x = y \to y = z \to x = z)$}
\UIC{$\forall x\,\forall y\,\forall z\,(x = y \to y = z \to x = z)$}
\AXC{$\forall x\,\forall y\,\forall z\,(x = y \to y = z \to x = z) \to \varphi_{45}$\lbl{A2}}
\BIC{$\varphi_{45}$}
\end{ptree}
\noindent$\Pi_{9}$:
\begin{ptree}
\AXC{$\stackrel{\Pi_{8}}{\vdots}$}\noLine
\UIC{$\varphi_{45}$}
\AXC{$\varphi_{45} \to \varphi_{151}$\lbl{A2}}
\BIC{$\varphi_{151}$}
\AXC{$\varphi_{151} \to (f \cdot g) \circ a(n) = f(a_{n}) \cdot g(a_{n}) \to \varphi_{164}$\lbl{A2}}
\BIC{$(f \cdot g) \circ a(n) = f(a_{n}) \cdot g(a_{n}) \to \varphi_{164}$}
\end{ptree}
\noindent$\Pi_{10}$:
\begin{ptree}
\AXC{$y = z \to x = y \to x = z$\lbl{A6}}
\AXC{$(y = z \to x = y \to x = z) \to x = y \to y = z \to x = z$\lbl{A1}}
\BIC{$x = y \to y = z \to x = z$}
\UIC{$\forall z\,(x = y \to y = z \to x = z)$}
\UIC{$\forall y\,\forall z\,(x = y \to y = z \to x = z)$}
\UIC{$\forall x\,\forall y\,\forall z\,(x = y \to y = z \to x = z)$}
\AXC{$\forall x\,\forall y\,\forall z\,(x = y \to y = z \to x = z) \to \varphi_{102}$\lbl{A2}}
\BIC{$\varphi_{102}$}
\end{ptree}
\noindent$\Pi_{11}$:
\begin{ptree}
\AXC{$\stackrel{\Pi_{10}}{\vdots}$}\noLine
\UIC{$\varphi_{102}$}
\AXC{$\varphi_{102} \to \varphi_{188}$\lbl{A2}}
\BIC{$\varphi_{188}$}
\AXC{$\varphi_{188} \to f \circ a \otimes g \circ a(n) = f(a_{n}) \cdot g \circ a(n) \to \varphi_{168}$\lbl{A2}}
\BIC{$f \circ a \otimes g \circ a(n) = f(a_{n}) \cdot g \circ a(n) \to \varphi_{168}$}
\end{ptree}
\noindent$\Pi_{12}$:
\begin{ptree}
\AXC{$\forall n \in \mathbb{N}\,(a_{n} \in A) \to n \in \mathbb{N} \to a_{n} \in A$\lbl{A2}}
\AXC{$\mathrm{ContIn}(f,A,c) \to \varphi_{78}$\lbl{A1}}
\AXC{$(\mathrm{ContIn}(f,A,c) \to \varphi_{78}) \to \varphi_{176}$\lbl{A1}}
\BIC{$\varphi_{176}$}
\BIC{$\mathrm{ContIn}(f,A,c) \to \varphi_{73}$}
\end{ptree}
\noindent$\Pi_{13}$:
\begin{ptree}
\AXC{$\forall f\,\forall A\,\forall c\,(\varphi_{191})$\lbl{DefContIn}}
\AXC{$\forall f\,\forall A\,\forall c\,(\varphi_{191}) \to \forall A\,\forall c\,(\varphi_{190})$\lbl{A2}}
\BIC{$\forall A\,\forall c\,(\varphi_{190})$}
\AXC{$\forall A\,\forall c\,(\varphi_{190}) \to \forall c\,(\varphi_{190})$\lbl{A2}}
\BIC{$\forall c\,(\varphi_{190})$}
\AXC{$\forall c\,(\varphi_{190}) \to \varphi_{190}$\lbl{A2}}
\BIC{$\varphi_{190}$}
\end{ptree}
\noindent$\Pi_{14}$:
\begin{ptree}
\AXC{$\stackrel{\Pi_{13}}{\vdots}$}\noLine
\UIC{$\varphi_{190}$}
\AXC{$\varphi_{7}$\lbl{A1}}
\AXC{$\varphi_{7} \to \varphi_{190} \to \varphi_{8}$\lbl{A1}}
\BIC{$\varphi_{190} \to \varphi_{8}$}
\BIC{$\varphi_{8}$}
\end{ptree}
\noindent$\Pi_{15}$:
\begin{ptree}
\AXC{$\stackrel{\Pi_{14}}{\vdots}$}\noLine
\UIC{$\varphi_{8}$}
\AXC{$\forall X\,\forall Y\,\forall x\,(X \subseteq Y \to x \in X \to x \in Y)$\lbl{[B5]}}
\AXC{$\varphi_{46}$\lbl{A2}}
\BIC{$\forall Y\,\forall x\,(A \subseteq Y \to x \in A \to x \in Y)$}
\AXC{$\varphi_{57}$\lbl{A2}}
\BIC{$\forall x\,(A \subseteq \mathrm{dom}(g) \to x \in A \to x \in \mathrm{dom}(g))$}
\AXC{$\varphi_{97}$\lbl{A2}}
\BIC{$A \subseteq \mathrm{dom}(g) \to a_{n} \in A \to a_{n} \in \mathrm{dom}(g)$}
\AXC{$\varphi_{244}$\lbl{A1}}
\BIC{$\varphi_{8} \to \varphi_{121} \to \varphi_{128}$}
\BIC{$\varphi_{121} \to \varphi_{128}$}
\end{ptree}
\noindent$\Pi_{16}$:
\begin{ptree}
\AXC{$\varphi_{110}$\lbl{A1}}
\AXC{$\stackrel{\Pi_{12}}{\vdots}$}\noLine
\UIC{$\mathrm{ContIn}(f,A,c) \to \varphi_{73}$}
\AXC{$(\mathrm{ContIn}(f,A,c) \to \varphi_{73}) \to \varphi_{110} \to \varphi_{121}$\lbl{A1}}
\BIC{$\varphi_{110} \to \varphi_{121}$}
\BIC{$\varphi_{121}$}
\AXC{$\stackrel{\Pi_{15}}{\vdots}$}\noLine
\UIC{$\varphi_{121} \to \varphi_{128}$}
\BIC{$\varphi_{128}$}
\end{ptree}
\noindent$\Pi_{17}$:
\begin{ptree}
\AXC{$\varphi_{69} \to \forall n\,(\varphi_{69})$\lbl{A4}}
\AXC{$\varphi_{218}$\lbl{A3}}
\AXC{$\stackrel{\Pi_{16}}{\vdots}$}\noLine
\UIC{$\varphi_{128}$}
\AXC{$\varphi_{128} \to \varphi_{69} \to n \in \mathbb{N} \to a_{n} \in \mathrm{dom}(g)$\lbl{A1}}
\BIC{$\varphi_{69} \to n \in \mathbb{N} \to a_{n} \in \mathrm{dom}(g)$}
\UIC{$\forall n\,(\varphi_{69} \to n \in \mathbb{N} \to a_{n} \in \mathrm{dom}(g))$}
\AXC{$\varphi_{278}$\lbl{A1}}
\BIC{$\varphi_{218} \to (\varphi_{69} \to \forall n\,(\varphi_{69})) \to \varphi_{131}$}
\BIC{$(\varphi_{69} \to \forall n\,(\varphi_{69})) \to \varphi_{131}$}
\BIC{$\varphi_{131}$}
\end{ptree}
\noindent$\Pi_{18}$:
\begin{ptree}
\AXC{$\forall f\,\forall A\,\forall c\,(\varphi_{191})$\lbl{DefContIn}}
\AXC{$\forall f\,\forall A\,\forall c\,(\varphi_{191}) \to \forall A\,\forall c\,(\varphi_{190})$\lbl{A2}}
\BIC{$\forall A\,\forall c\,(\varphi_{190})$}
\AXC{$\forall A\,\forall c\,(\varphi_{190}) \to \forall c\,(\varphi_{190})$\lbl{A2}}
\BIC{$\forall c\,(\varphi_{190})$}
\AXC{$\forall c\,(\varphi_{190}) \to \varphi_{190}$\lbl{A2}}
\BIC{$\varphi_{190}$}
\end{ptree}
\noindent$\Pi_{19}$:
\begin{ptree}
\AXC{$\stackrel{\Pi_{18}}{\vdots}$}\noLine
\UIC{$\varphi_{190}$}
\AXC{$\varphi_{7}$\lbl{A1}}
\AXC{$\varphi_{7} \to \varphi_{190} \to \varphi_{4}$\lbl{A1}}
\BIC{$\varphi_{190} \to \varphi_{4}$}
\BIC{$\varphi_{4}$}
\end{ptree}
\noindent$\Pi_{20}$:
\begin{ptree}
\AXC{$\varphi_{11}$\lbl{A1}}
\AXC{$\stackrel{\Pi_{19}}{\vdots}$}\noLine
\UIC{$\varphi_{4}$}
\AXC{$\forall f\,\forall a\,(\mathrm{RFn}(f) \to \mathrm{Seq}(a) \to \varphi_{114})$\lbl{[S8]}}
\AXC{$\varphi_{228}$\lbl{A2}}
\BIC{$\forall a\,(\mathrm{RFn}(g) \to \mathrm{Seq}(a) \to \varphi_{113})$}
\AXC{$\varphi_{224}$\lbl{A2}}
\BIC{$\mathrm{RFn}(g) \to \mathrm{Seq}(a) \to \varphi_{113}$}
\AXC{$\varphi_{268}$\lbl{A1}}
\BIC{$\varphi_{4} \to \varphi_{11} \to \varphi_{131} \to \mathrm{ContIn}(f,A,c) \to \varphi_{158}$}
\BIC{$\varphi_{11} \to \varphi_{131} \to \mathrm{ContIn}(f,A,c) \to \varphi_{158}$}
\BIC{$\varphi_{131} \to \mathrm{ContIn}(f,A,c) \to \varphi_{158}$}
\end{ptree}
\noindent$\Pi_{21}$:
\begin{ptree}
\AXC{$\varphi_{105}$\lbl{A2}}
\AXC{$\stackrel{\Pi_{17}}{\vdots}$}\noLine
\UIC{$\varphi_{131}$}
\AXC{$\stackrel{\Pi_{20}}{\vdots}$}\noLine
\UIC{$\varphi_{131} \to \mathrm{ContIn}(f,A,c) \to \varphi_{158}$}
\BIC{$\mathrm{ContIn}(f,A,c) \to \varphi_{158}$}
\AXC{$\varphi_{238}$\lbl{A1}}
\BIC{$\mathrm{ContIn}(f,A,c) \to \varphi_{144}$}
\AXC{$\varphi_{257}$\lbl{A1}}
\BIC{$\varphi_{105} \to \mathrm{ContIn}(f,A,c) \to \varphi_{141}$}
\BIC{$\mathrm{ContIn}(f,A,c) \to \varphi_{141}$}
\end{ptree}
\noindent$\Pi_{22}$:
\begin{ptree}
\AXC{$\varphi_{110}$\lbl{A1}}
\AXC{$\stackrel{\Pi_{21}}{\vdots}$}\noLine
\UIC{$\mathrm{ContIn}(f,A,c) \to \varphi_{141}$}
\AXC{$(\mathrm{ContIn}(f,A,c) \to \varphi_{141}) \to \varphi_{110} \to \varphi_{146}$\lbl{A1}}
\BIC{$\varphi_{110} \to \varphi_{146}$}
\BIC{$\varphi_{146}$}
\end{ptree}
\noindent$\Pi_{23}$:
\begin{ptree}
\AXC{$x = x$\lbl{A5}}
\UIC{$\forall x\,(x = x)$}
\AXC{$\varphi_{77}$\lbl{A2}}
\BIC{$f(a_{n}) \cdot g \circ a(n) = f(a_{n}) \cdot g \circ a(n)$}
\AXC{$\stackrel{\Pi_{22}}{\vdots}$}\noLine
\UIC{$\varphi_{146}$}
\AXC{$\varphi_{146} \to f(a_{n}) \cdot g \circ a(n) = f(a_{n}) \cdot g \circ a(n) \to \varphi_{251}$\lbl{A1}}
\BIC{$f(a_{n}) \cdot g \circ a(n) = f(a_{n}) \cdot g \circ a(n) \to \varphi_{251}$}
\BIC{$\varphi_{251}$}
\end{ptree}
\noindent$\Pi_{24}$:
\begin{ptree}
\AXC{$z = w \to \varphi_{133}$\lbl{A6}}
\UIC{$\forall w\,(z = w \to \varphi_{133})$}
\UIC{$\forall z\,\forall w\,(z = w \to \varphi_{133})$}
\AXC{$\forall z\,\forall w\,(z = w \to \varphi_{133}) \to \varphi_{184}$\lbl{A2}}
\BIC{$\varphi_{184}$}
\AXC{$\varphi_{184} \to g \circ a(n) = g(a_{n}) \to \varphi_{171}$\lbl{A2}}
\BIC{$g \circ a(n) = g(a_{n}) \to \varphi_{171}$}
\AXC{$\stackrel{\Pi_{23}}{\vdots}$}\noLine
\UIC{$\varphi_{251}$}
\BIC{$\mathrm{ContIn}(f,A,c) \to \varphi_{172}$}
\end{ptree}
\noindent$\Pi_{25}$:
\begin{ptree}
\AXC{$y = z \to x = y \to x = z$\lbl{A6}}
\AXC{$(y = z \to x = y \to x = z) \to x = y \to y = z \to x = z$\lbl{A1}}
\BIC{$x = y \to y = z \to x = z$}
\UIC{$\forall z\,(x = y \to y = z \to x = z)$}
\UIC{$\forall y\,\forall z\,(x = y \to y = z \to x = z)$}
\UIC{$\forall x\,\forall y\,\forall z\,(x = y \to y = z \to x = z)$}
\AXC{$\forall x\,\forall y\,\forall z\,(x = y \to y = z \to x = z) \to \varphi_{102}$\lbl{A2}}
\BIC{$\varphi_{102}$}
\end{ptree}
\noindent$\Pi_{26}$:
\begin{ptree}
\AXC{$\forall n \in \mathbb{N}\,(a_{n} \in A) \to n \in \mathbb{N} \to a_{n} \in A$\lbl{A2}}
\AXC{$\mathrm{ContIn}(f,A,c) \to \varphi_{78}$\lbl{A1}}
\AXC{$(\mathrm{ContIn}(f,A,c) \to \varphi_{78}) \to \varphi_{176}$\lbl{A1}}
\BIC{$\varphi_{176}$}
\BIC{$\mathrm{ContIn}(f,A,c) \to \varphi_{73}$}
\end{ptree}
\noindent$\Pi_{27}$:
\begin{ptree}
\AXC{$\forall f\,\forall A\,\forall c\,(\varphi_{191})$\lbl{DefContIn}}
\AXC{$\forall f\,\forall A\,\forall c\,(\varphi_{191}) \to \forall A\,\forall c\,(\varphi_{191})$\lbl{A2}}
\BIC{$\forall A\,\forall c\,(\varphi_{191})$}
\AXC{$\forall A\,\forall c\,(\varphi_{191}) \to \forall c\,(\varphi_{191})$\lbl{A2}}
\BIC{$\forall c\,(\varphi_{191})$}
\AXC{$\forall c\,(\varphi_{191}) \to \varphi_{191}$\lbl{A2}}
\BIC{$\varphi_{191}$}
\end{ptree}
\noindent$\Pi_{28}$:
\begin{ptree}
\AXC{$\stackrel{\Pi_{27}}{\vdots}$}\noLine
\UIC{$\varphi_{191}$}
\AXC{$\mathrm{ContIn}(f,A,c) \to \mathrm{ContIn}(f,A,c)$\lbl{A1}}
\AXC{$\varphi_{223}$\lbl{A1}}
\BIC{$\varphi_{191} \to \varphi_{9}$}
\BIC{$\varphi_{9}$}
\end{ptree}
\noindent$\Pi_{29}$:
\begin{ptree}
\AXC{$\stackrel{\Pi_{28}}{\vdots}$}\noLine
\UIC{$\varphi_{9}$}
\AXC{$\forall X\,\forall Y\,\forall x\,(X \subseteq Y \to x \in X \to x \in Y)$\lbl{[B5]}}
\AXC{$\varphi_{46}$\lbl{A2}}
\BIC{$\forall Y\,\forall x\,(A \subseteq Y \to x \in A \to x \in Y)$}
\AXC{$\varphi_{58}$\lbl{A2}}
\BIC{$\forall x\,(A \subseteq \mathrm{dom}(f) \to x \in A \to x \in \mathrm{dom}(f))$}
\AXC{$\varphi_{98}$\lbl{A2}}
\BIC{$A \subseteq \mathrm{dom}(f) \to a_{n} \in A \to a_{n} \in \mathrm{dom}(f)$}
\AXC{$\varphi_{245}$\lbl{A1}}
\BIC{$\varphi_{9} \to \varphi_{121} \to \varphi_{129}$}
\BIC{$\varphi_{121} \to \varphi_{129}$}
\end{ptree}
\noindent$\Pi_{30}$:
\begin{ptree}
\AXC{$\varphi_{110}$\lbl{A1}}
\AXC{$\stackrel{\Pi_{26}}{\vdots}$}\noLine
\UIC{$\mathrm{ContIn}(f,A,c) \to \varphi_{73}$}
\AXC{$(\mathrm{ContIn}(f,A,c) \to \varphi_{73}) \to \varphi_{110} \to \varphi_{121}$\lbl{A1}}
\BIC{$\varphi_{110} \to \varphi_{121}$}
\BIC{$\varphi_{121}$}
\AXC{$\stackrel{\Pi_{29}}{\vdots}$}\noLine
\UIC{$\varphi_{121} \to \varphi_{129}$}
\BIC{$\varphi_{129}$}
\end{ptree}
\noindent$\Pi_{31}$:
\begin{ptree}
\AXC{$\varphi_{69} \to \forall n\,(\varphi_{69})$\lbl{A4}}
\AXC{$\varphi_{219}$\lbl{A3}}
\AXC{$\stackrel{\Pi_{30}}{\vdots}$}\noLine
\UIC{$\varphi_{129}$}
\AXC{$\varphi_{129} \to \varphi_{69} \to n \in \mathbb{N} \to a_{n} \in \mathrm{dom}(f)$\lbl{A1}}
\BIC{$\varphi_{69} \to n \in \mathbb{N} \to a_{n} \in \mathrm{dom}(f)$}
\UIC{$\forall n\,(\varphi_{69} \to n \in \mathbb{N} \to a_{n} \in \mathrm{dom}(f))$}
\AXC{$\varphi_{279}$\lbl{A1}}
\BIC{$\varphi_{219} \to (\varphi_{69} \to \forall n\,(\varphi_{69})) \to \varphi_{132}$}
\BIC{$(\varphi_{69} \to \forall n\,(\varphi_{69})) \to \varphi_{132}$}
\BIC{$\varphi_{132}$}
\end{ptree}
\noindent$\Pi_{32}$:
\begin{ptree}
\AXC{$\forall f\,\forall A\,\forall c\,(\varphi_{191})$\lbl{DefContIn}}
\AXC{$\forall f\,\forall A\,\forall c\,(\varphi_{191}) \to \forall A\,\forall c\,(\varphi_{191})$\lbl{A2}}
\BIC{$\forall A\,\forall c\,(\varphi_{191})$}
\AXC{$\forall A\,\forall c\,(\varphi_{191}) \to \forall c\,(\varphi_{191})$\lbl{A2}}
\BIC{$\forall c\,(\varphi_{191})$}
\AXC{$\forall c\,(\varphi_{191}) \to \varphi_{191}$\lbl{A2}}
\BIC{$\varphi_{191}$}
\end{ptree}
\noindent$\Pi_{33}$:
\begin{ptree}
\AXC{$\stackrel{\Pi_{32}}{\vdots}$}\noLine
\UIC{$\varphi_{191}$}
\AXC{$\mathrm{ContIn}(f,A,c) \to \mathrm{ContIn}(f,A,c)$\lbl{A1}}
\AXC{$\varphi_{221}$\lbl{A1}}
\BIC{$\varphi_{191} \to \varphi_{5}$}
\BIC{$\varphi_{5}$}
\end{ptree}
\noindent$\Pi_{34}$:
\begin{ptree}
\AXC{$\varphi_{11}$\lbl{A1}}
\AXC{$\stackrel{\Pi_{33}}{\vdots}$}\noLine
\UIC{$\varphi_{5}$}
\AXC{$\forall f\,\forall a\,(\mathrm{RFn}(f) \to \mathrm{Seq}(a) \to \varphi_{114})$\lbl{[S8]}}
\AXC{$\varphi_{229}$\lbl{A2}}
\BIC{$\forall a\,(\mathrm{RFn}(f) \to \mathrm{Seq}(a) \to \varphi_{114})$}
\AXC{$\varphi_{225}$\lbl{A2}}
\BIC{$\mathrm{RFn}(f) \to \mathrm{Seq}(a) \to \varphi_{114}$}
\AXC{$\varphi_{269}$\lbl{A1}}
\BIC{$\varphi_{5} \to \varphi_{11} \to \varphi_{132} \to \mathrm{ContIn}(f,A,c) \to \varphi_{159}$}
\BIC{$\varphi_{11} \to \varphi_{132} \to \mathrm{ContIn}(f,A,c) \to \varphi_{159}$}
\BIC{$\varphi_{132} \to \mathrm{ContIn}(f,A,c) \to \varphi_{159}$}
\end{ptree}
\noindent$\Pi_{35}$:
\begin{ptree}
\AXC{$\varphi_{106}$\lbl{A2}}
\AXC{$\stackrel{\Pi_{31}}{\vdots}$}\noLine
\UIC{$\varphi_{132}$}
\AXC{$\stackrel{\Pi_{34}}{\vdots}$}\noLine
\UIC{$\varphi_{132} \to \mathrm{ContIn}(f,A,c) \to \varphi_{159}$}
\BIC{$\mathrm{ContIn}(f,A,c) \to \varphi_{159}$}
\AXC{$\varphi_{239}$\lbl{A1}}
\BIC{$\mathrm{ContIn}(f,A,c) \to \varphi_{145}$}
\AXC{$\varphi_{258}$\lbl{A1}}
\BIC{$\varphi_{106} \to \mathrm{ContIn}(f,A,c) \to \varphi_{142}$}
\BIC{$\mathrm{ContIn}(f,A,c) \to \varphi_{142}$}
\end{ptree}
\noindent$\Pi_{36}$:
\begin{ptree}
\AXC{$\varphi_{110}$\lbl{A1}}
\AXC{$\stackrel{\Pi_{35}}{\vdots}$}\noLine
\UIC{$\mathrm{ContIn}(f,A,c) \to \varphi_{142}$}
\AXC{$(\mathrm{ContIn}(f,A,c) \to \varphi_{142}) \to \varphi_{110} \to \varphi_{147}$\lbl{A1}}
\BIC{$\varphi_{110} \to \varphi_{147}$}
\BIC{$\varphi_{147}$}
\AXC{$\varphi_{276}$\lbl{A1}}
\BIC{$f \circ a(n) \cdot g \circ a(n) = f \circ a(n) \cdot g \circ a(n) \to \varphi_{254}$}
\end{ptree}
\noindent$\Pi_{37}$:
\begin{ptree}
\AXC{$x = x$\lbl{A5}}
\UIC{$\forall x\,(x = x)$}
\AXC{$\varphi_{92}$\lbl{A2}}
\BIC{$f \circ a(n) \cdot g \circ a(n) = f \circ a(n) \cdot g \circ a(n)$}
\AXC{$\stackrel{\Pi_{36}}{\vdots}$}\noLine
\UIC{$f \circ a(n) \cdot g \circ a(n) = f \circ a(n) \cdot g \circ a(n) \to \varphi_{254}$}
\BIC{$\varphi_{254}$}
\end{ptree}
\noindent$\Pi_{38}$:
\begin{ptree}
\AXC{$z = w \to \varphi_{152}$\lbl{A6}}
\UIC{$\forall w\,(z = w \to \varphi_{152})$}
\UIC{$\forall z\,\forall w\,(z = w \to \varphi_{152})$}
\AXC{$\forall z\,\forall w\,(z = w \to \varphi_{152}) \to \varphi_{196}$\lbl{A2}}
\BIC{$\varphi_{196}$}
\AXC{$\varphi_{196} \to f \circ a(n) = f(a_{n}) \to \varphi_{187}$\lbl{A2}}
\BIC{$f \circ a(n) = f(a_{n}) \to \varphi_{187}$}
\AXC{$\stackrel{\Pi_{37}}{\vdots}$}\noLine
\UIC{$\varphi_{254}$}
\BIC{$\mathrm{ContIn}(f,A,c) \to \varphi_{178}$}
\end{ptree}
\noindent$\Pi_{39}$:
\begin{ptree}
\AXC{$\forall a\,\forall \alpha\,(\mathrm{Conv}(a,\alpha) \to \mathrm{Seq}(a))$\lbl{[C1]}}
\AXC{$\varphi_{47}$\lbl{A2}}
\BIC{$\forall \alpha\,(\mathrm{Conv}(g \circ a,\alpha) \to \mathrm{Seq}(g \circ a))$}
\AXC{$\varphi_{94}$\lbl{A2}}
\BIC{$\mathrm{Conv}(g \circ a,g(c)) \to \mathrm{Seq}(g \circ a)$}
\AXC{$(\mathrm{Conv}(g \circ a,g(c)) \to \mathrm{Seq}(g \circ a)) \to \varphi_{211}$\lbl{A1}}
\BIC{$\varphi_{211}$}
\end{ptree}
\noindent$\Pi_{40}$:
\begin{ptree}
\AXC{$\forall a\,\forall \alpha\,(\mathrm{Conv}(a,\alpha) \to \mathrm{Seq}(a))$\lbl{[C1]}}
\AXC{$\varphi_{48}$\lbl{A2}}
\BIC{$\forall \alpha\,(\mathrm{Conv}(f \circ a,\alpha) \to \mathrm{Seq}(f \circ a))$}
\AXC{$\varphi_{95}$\lbl{A2}}
\BIC{$\mathrm{Conv}(f \circ a,f(c)) \to \mathrm{Seq}(f \circ a)$}
\AXC{$(\mathrm{Conv}(f \circ a,f(c)) \to \mathrm{Seq}(f \circ a)) \to \varphi_{212}$\lbl{A1}}
\BIC{$\varphi_{212}$}
\end{ptree}
\noindent$\Pi_{41}$:
\begin{ptree}
\AXC{$\stackrel{\Pi_{5}}{\vdots}$}\noLine
\UIC{$\varphi_{123}$}
\AXC{$\stackrel{\Pi_{40}}{\vdots}$}\noLine
\UIC{$\varphi_{212}$}
\BIC{$\mathrm{ContIn}(f,A,c) \to \mathrm{ContIn}(g,A,c) \to \varphi_{38}$}
\AXC{$\varphi_{99}$\lbl{[S4]}}
\AXC{$\varphi_{99} \to \varphi_{153}$\lbl{A2}}
\BIC{$\varphi_{153}$}
\AXC{$\varphi_{153} \to \varphi_{195}$\lbl{A2}}
\BIC{$\varphi_{195}$}
\AXC{$\varphi_{195} \to \varphi_{264}$\lbl{A1}}
\BIC{$\varphi_{264}$}
\BIC{$\varphi_{243}$}
\end{ptree}
\noindent$\Pi_{42}$:
\begin{ptree}
\AXC{$\varphi_{200}$\lbl{A2}}
\AXC{$\stackrel{\Pi_{2}}{\vdots}$}\noLine
\UIC{$\varphi_{122}$}
\AXC{$\stackrel{\Pi_{39}}{\vdots}$}\noLine
\UIC{$\varphi_{211}$}
\BIC{$\mathrm{ContIn}(f,A,c) \to \mathrm{ContIn}(g,A,c) \to \varphi_{37}$}
\AXC{$\stackrel{\Pi_{41}}{\vdots}$}\noLine
\UIC{$\varphi_{243}$}
\BIC{$\varphi_{214}$}
\AXC{$\varphi_{262}$\lbl{A1}}
\BIC{$\mathrm{ContIn}(f,A,c) \to \mathrm{ContIn}(g,A,c) \to \varphi_{180}$}
\AXC{$\varphi_{275}$\lbl{A1}}
\BIC{$\varphi_{200} \to \mathrm{ContIn}(f,A,c) \to \varphi_{192}$}
\BIC{$\mathrm{ContIn}(f,A,c) \to \varphi_{192}$}
\AXC{$\varphi_{267}$\lbl{A1}}
\BIC{$\varphi_{110} \to \mathrm{ContIn}(f,A,c) \to \varphi_{182}$}
\end{ptree}
\noindent$\Pi_{43}$:
\begin{ptree}
\AXC{$\stackrel{\Pi_{38}}{\vdots}$}\noLine
\UIC{$\mathrm{ContIn}(f,A,c) \to \varphi_{178}$}
\AXC{$\varphi_{110}$\lbl{A1}}
\AXC{$\stackrel{\Pi_{42}}{\vdots}$}\noLine
\UIC{$\varphi_{110} \to \mathrm{ContIn}(f,A,c) \to \varphi_{182}$}
\BIC{$\mathrm{ContIn}(f,A,c) \to \varphi_{182}$}
\AXC{$\varphi_{284}$\lbl{A1}}
\BIC{$(\mathrm{ContIn}(f,A,c) \to \varphi_{178}) \to \varphi_{263}$}
\BIC{$\varphi_{263}$}
\end{ptree}
\noindent$\Pi_{44}$:
\begin{ptree}
\AXC{$\stackrel{\Pi_{25}}{\vdots}$}\noLine
\UIC{$\varphi_{102}$}
\AXC{$\varphi_{102} \to \varphi_{194}$\lbl{A2}}
\BIC{$\varphi_{194}$}
\AXC{$\varphi_{194} \to f \circ a \otimes g \circ a(n) = f \circ a(n) \cdot g \circ a(n) \to \varphi_{179}$\lbl{A2}}
\BIC{$f \circ a \otimes g \circ a(n) = f \circ a(n) \cdot g \circ a(n) \to \varphi_{179}$}
\AXC{$\stackrel{\Pi_{43}}{\vdots}$}\noLine
\UIC{$\varphi_{263}$}
\BIC{$\mathrm{ContIn}(f,A,c) \to \varphi_{177}$}
\AXC{$\varphi_{283}$\lbl{A1}}
\BIC{$(\mathrm{ContIn}(f,A,c) \to \varphi_{172}) \to \varphi_{260}$}
\end{ptree}
\noindent$\Pi_{45}$:
\begin{ptree}
\AXC{$\stackrel{\Pi_{11}}{\vdots}$}\noLine
\UIC{$f \circ a \otimes g \circ a(n) = f(a_{n}) \cdot g \circ a(n) \to \varphi_{168}$}
\AXC{$\stackrel{\Pi_{24}}{\vdots}$}\noLine
\UIC{$\mathrm{ContIn}(f,A,c) \to \varphi_{172}$}
\AXC{$\stackrel{\Pi_{44}}{\vdots}$}\noLine
\UIC{$(\mathrm{ContIn}(f,A,c) \to \varphi_{172}) \to \varphi_{260}$}
\BIC{$\varphi_{260}$}
\BIC{$\mathrm{ContIn}(f,A,c) \to \varphi_{173}$}
\end{ptree}
\noindent$\Pi_{46}$:
\begin{ptree}
\AXC{$\stackrel{\Pi_{7}}{\vdots}$}\noLine
\UIC{$\forall x\,\forall y\,(x = y \to y = x)$}
\AXC{$\varphi_{118}$\lbl{A2}}
\BIC{$\forall y\,(f \circ a \otimes g \circ a(n) = y \to y = f \circ a \otimes g \circ a(n))$}
\AXC{$\varphi_{220}$\lbl{A2}}
\BIC{$\varphi_{170}$}
\AXC{$\stackrel{\Pi_{45}}{\vdots}$}\noLine
\UIC{$\mathrm{ContIn}(f,A,c) \to \varphi_{173}$}
\AXC{$\varphi_{271}$\lbl{A1}}
\BIC{$\varphi_{170} \to \mathrm{ContIn}(f,A,c) \to \varphi_{174}$}
\BIC{$\mathrm{ContIn}(f,A,c) \to \varphi_{174}$}
\end{ptree}
\noindent$\Pi_{47}$:
\begin{ptree}
\AXC{$y = z \to x = y \to x = z$\lbl{A6}}
\AXC{$(y = z \to x = y \to x = z) \to x = y \to y = z \to x = z$\lbl{A1}}
\BIC{$x = y \to y = z \to x = z$}
\UIC{$\forall z\,(x = y \to y = z \to x = z)$}
\UIC{$\forall y\,\forall z\,(x = y \to y = z \to x = z)$}
\UIC{$\forall x\,\forall y\,\forall z\,(x = y \to y = z \to x = z)$}
\AXC{$\forall x\,\forall y\,\forall z\,(x = y \to y = z \to x = z) \to \varphi_{45}$\lbl{A2}}
\BIC{$\varphi_{45}$}
\end{ptree}
\noindent$\Pi_{48}$:
\begin{ptree}
\AXC{$\stackrel{\Pi_{47}}{\vdots}$}\noLine
\UIC{$\varphi_{45}$}
\AXC{$\varphi_{45} \to \varphi_{126}$\lbl{A2}}
\BIC{$\varphi_{126}$}
\AXC{$\varphi_{126} \to (f \cdot g) \circ a(n) = (f \cdot g)(a_{n}) \to \varphi_{136}$\lbl{A2}}
\BIC{$(f \cdot g) \circ a(n) = (f \cdot g)(a_{n}) \to \varphi_{136}$}
\end{ptree}
\noindent$\Pi_{49}$:
\begin{ptree}
\AXC{$\stackrel{\Pi_{33}}{\vdots}$}\noLine
\UIC{$\varphi_{5}$}
\AXC{$\forall f\,(\varphi_{175})$\lbl{[S18]}}
\AXC{$\forall f\,(\varphi_{175}) \to \varphi_{175}$\lbl{A2}}
\BIC{$\varphi_{175}$}
\AXC{$\varphi_{175} \to \mathrm{RFn}(f) \to \mathrm{RFn}(g) \to \mathrm{RFn}((f \cdot g)) \land \varphi_{16} \land \varphi_{27}$\lbl{A2}}
\BIC{$\mathrm{RFn}(f) \to \mathrm{RFn}(g) \to \mathrm{RFn}((f \cdot g)) \land \varphi_{16} \land \varphi_{27}$}
\AXC{$\varphi_{259}$\lbl{A1}}
\BIC{$\varphi_{5} \to \varphi_{4} \to \varphi_{185}$}
\BIC{$\varphi_{4} \to \varphi_{185}$}
\end{ptree}
\noindent$\Pi_{50}$:
\begin{ptree}
\AXC{$\stackrel{\Pi_{19}}{\vdots}$}\noLine
\UIC{$\varphi_{4}$}
\AXC{$\stackrel{\Pi_{49}}{\vdots}$}\noLine
\UIC{$\varphi_{4} \to \varphi_{185}$}
\BIC{$\varphi_{185}$}
\end{ptree}
\noindent$\Pi_{51}$:
\begin{ptree}
\AXC{$\stackrel{\Pi_{50}}{\vdots}$}\noLine
\UIC{$\varphi_{185}$}
\AXC{$\varphi_{226}$\lbl{A1}}
\BIC{$\mathrm{ContIn}(f,A,c) \to \mathrm{ContIn}(g,A,c) \to \varphi_{16}$}
\end{ptree}
\noindent$\Pi_{52}$:
\begin{ptree}
\AXC{$\varphi_{16} \to \varphi_{43}$\lbl{A2}}
\AXC{$\stackrel{\Pi_{51}}{\vdots}$}\noLine
\UIC{$\mathrm{ContIn}(f,A,c) \to \mathrm{ContIn}(g,A,c) \to \varphi_{16}$}
\AXC{$\varphi_{241}$\lbl{A1}}
\BIC{$\varphi_{217}$}
\BIC{$\mathrm{ContIn}(f,A,c) \to \mathrm{ContIn}(g,A,c) \to \varphi_{43}$}
\end{ptree}
\noindent$\Pi_{53}$:
\begin{ptree}
\AXC{$\forall n \in \mathbb{N}\,(a_{n} \in A) \to n \in \mathbb{N} \to a_{n} \in A$\lbl{A2}}
\AXC{$\mathrm{ContIn}(f,A,c) \to \varphi_{78}$\lbl{A1}}
\AXC{$(\mathrm{ContIn}(f,A,c) \to \varphi_{78}) \to \varphi_{176}$\lbl{A1}}
\BIC{$\varphi_{176}$}
\BIC{$\mathrm{ContIn}(f,A,c) \to \varphi_{73}$}
\end{ptree}
\noindent$\Pi_{54}$:
\begin{ptree}
\AXC{$\varphi_{110}$\lbl{A1}}
\AXC{$\stackrel{\Pi_{53}}{\vdots}$}\noLine
\UIC{$\mathrm{ContIn}(f,A,c) \to \varphi_{73}$}
\AXC{$(\mathrm{ContIn}(f,A,c) \to \varphi_{73}) \to \varphi_{110} \to \varphi_{121}$\lbl{A1}}
\BIC{$\varphi_{110} \to \varphi_{121}$}
\BIC{$\varphi_{121}$}
\end{ptree}
\noindent$\Pi_{55}$:
\begin{ptree}
\AXC{$\stackrel{\Pi_{54}}{\vdots}$}\noLine
\UIC{$\varphi_{121}$}
\AXC{$\stackrel{\Pi_{14}}{\vdots}$}\noLine
\UIC{$\varphi_{8}$}
\AXC{$\forall X\,\forall Y\,\forall x\,(X \subseteq Y \to x \in X \to x \in Y)$\lbl{[B5]}}
\AXC{$\varphi_{46}$\lbl{A2}}
\BIC{$\forall Y\,\forall x\,(A \subseteq Y \to x \in A \to x \in Y)$}
\AXC{$\varphi_{57}$\lbl{A2}}
\BIC{$\forall x\,(A \subseteq \mathrm{dom}(g) \to x \in A \to x \in \mathrm{dom}(g))$}
\AXC{$\varphi_{97}$\lbl{A2}}
\BIC{$A \subseteq \mathrm{dom}(g) \to a_{n} \in A \to a_{n} \in \mathrm{dom}(g)$}
\AXC{$\varphi_{244}$\lbl{A1}}
\BIC{$\varphi_{8} \to \varphi_{121} \to \varphi_{128}$}
\BIC{$\varphi_{121} \to \varphi_{128}$}
\BIC{$\varphi_{128}$}
\end{ptree}
\noindent$\Pi_{56}$:
\begin{ptree}
\AXC{$\stackrel{\Pi_{54}}{\vdots}$}\noLine
\UIC{$\varphi_{121}$}
\AXC{$\stackrel{\Pi_{28}}{\vdots}$}\noLine
\UIC{$\varphi_{9}$}
\AXC{$\forall X\,\forall Y\,\forall x\,(X \subseteq Y \to x \in X \to x \in Y)$\lbl{[B5]}}
\AXC{$\varphi_{46}$\lbl{A2}}
\BIC{$\forall Y\,\forall x\,(A \subseteq Y \to x \in A \to x \in Y)$}
\AXC{$\varphi_{58}$\lbl{A2}}
\BIC{$\forall x\,(A \subseteq \mathrm{dom}(f) \to x \in A \to x \in \mathrm{dom}(f))$}
\AXC{$\varphi_{98}$\lbl{A2}}
\BIC{$A \subseteq \mathrm{dom}(f) \to a_{n} \in A \to a_{n} \in \mathrm{dom}(f)$}
\AXC{$\varphi_{245}$\lbl{A1}}
\BIC{$\varphi_{9} \to \varphi_{121} \to \varphi_{129}$}
\BIC{$\varphi_{121} \to \varphi_{129}$}
\BIC{$\varphi_{129}$}
\end{ptree}
\noindent$\Pi_{57}$:
\begin{ptree}
\AXC{$\stackrel{\Pi_{52}}{\vdots}$}\noLine
\UIC{$\mathrm{ContIn}(f,A,c) \to \mathrm{ContIn}(g,A,c) \to \varphi_{43}$}
\AXC{$\stackrel{\Pi_{55}}{\vdots}$}\noLine
\UIC{$\varphi_{128}$}
\AXC{$\stackrel{\Pi_{56}}{\vdots}$}\noLine
\UIC{$\varphi_{129}$}
\AXC{$\varphi_{129} \to \varphi_{128} \to \varphi_{216}$\lbl{A1}}
\BIC{$\varphi_{128} \to \varphi_{216}$}
\BIC{$\varphi_{216}$}
\BIC{$\varphi_{138}$}
\end{ptree}
\noindent$\Pi_{58}$:
\begin{ptree}
\AXC{$\varphi_{69} \to \forall n\,(\varphi_{69})$\lbl{A4}}
\AXC{$\varphi_{227}$\lbl{A3}}
\AXC{$\stackrel{\Pi_{57}}{\vdots}$}\noLine
\UIC{$\varphi_{138}$}
\AXC{$\varphi_{138} \to \varphi_{69} \to n \in \mathbb{N} \to a_{n} \in \mathrm{dom}((f \cdot g))$\lbl{A1}}
\BIC{$\varphi_{69} \to n \in \mathbb{N} \to a_{n} \in \mathrm{dom}((f \cdot g))$}
\UIC{$\forall n\,(\varphi_{69} \to n \in \mathbb{N} \to a_{n} \in \mathrm{dom}((f \cdot g)))$}
\AXC{$\varphi_{281}$\lbl{A1}}
\BIC{$\varphi_{227} \to (\varphi_{69} \to \forall n\,(\varphi_{69})) \to \varphi_{139}$}
\BIC{$(\varphi_{69} \to \forall n\,(\varphi_{69})) \to \varphi_{139}$}
\BIC{$\varphi_{139}$}
\end{ptree}
\noindent$\Pi_{59}$:
\begin{ptree}
\AXC{$\stackrel{\Pi_{50}}{\vdots}$}\noLine
\UIC{$\varphi_{185}$}
\AXC{$\varphi_{231}$\lbl{A1}}
\BIC{$\mathrm{ContIn}(f,A,c) \to \mathrm{ContIn}(g,A,c) \to \varphi_{27}$}
\end{ptree}
\noindent$\Pi_{60}$:
\begin{ptree}
\AXC{$\varphi_{27} \to \varphi_{82}$\lbl{A2}}
\AXC{$\stackrel{\Pi_{59}}{\vdots}$}\noLine
\UIC{$\mathrm{ContIn}(f,A,c) \to \mathrm{ContIn}(g,A,c) \to \varphi_{27}$}
\AXC{$\varphi_{252}$\lbl{A1}}
\BIC{$\varphi_{233}$}
\BIC{$\mathrm{ContIn}(f,A,c) \to \mathrm{ContIn}(g,A,c) \to \varphi_{82}$}
\AXC{$\varphi_{274}$\lbl{A1}}
\BIC{$\varphi_{138} \to \varphi_{110} \to \mathrm{ContIn}(f,A,c) \to \varphi_{160}$}
\end{ptree}
\noindent$\Pi_{61}$:
\begin{ptree}
\AXC{$\varphi_{110}$\lbl{A1}}
\AXC{$\varphi_{75}$\lbl{A2}}
\AXC{$\stackrel{\Pi_{58}}{\vdots}$}\noLine
\UIC{$\varphi_{139}$}
\AXC{$\varphi_{139} \to \varphi_{75} \to \varphi_{138}$\lbl{A1}}
\BIC{$\varphi_{75} \to \varphi_{138}$}
\BIC{$\varphi_{138}$}
\AXC{$\stackrel{\Pi_{60}}{\vdots}$}\noLine
\UIC{$\varphi_{138} \to \varphi_{110} \to \mathrm{ContIn}(f,A,c) \to \varphi_{160}$}
\BIC{$\varphi_{110} \to \mathrm{ContIn}(f,A,c) \to \varphi_{160}$}
\BIC{$\mathrm{ContIn}(f,A,c) \to \varphi_{160}$}
\end{ptree}
\noindent$\Pi_{62}$:
\begin{ptree}
\AXC{$\stackrel{\Pi_{50}}{\vdots}$}\noLine
\UIC{$\varphi_{185}$}
\AXC{$\varphi_{185} \to \varphi_{10}$\lbl{A1}}
\BIC{$\varphi_{10}$}
\end{ptree}
\noindent$\Pi_{63}$:
\begin{ptree}
\AXC{$\forall f\,\forall a\,(\mathrm{RFn}(f) \to \mathrm{Seq}(a) \to \varphi_{114})$\lbl{[S8]}}
\AXC{$\forall f\,\forall a\,(\mathrm{RFn}(f) \to \mathrm{Seq}(a) \to \varphi_{114}) \to \forall a\,(\varphi_{183})$\lbl{A2}}
\BIC{$\forall a\,(\varphi_{183})$}
\AXC{$\forall a\,(\varphi_{183}) \to \varphi_{183}$\lbl{A2}}
\BIC{$\varphi_{183}$}
\end{ptree}
\noindent$\Pi_{64}$:
\begin{ptree}
\AXC{$\stackrel{\Pi_{58}}{\vdots}$}\noLine
\UIC{$\varphi_{139}$}
\AXC{$\varphi_{11}$\lbl{A1}}
\AXC{$\stackrel{\Pi_{62}}{\vdots}$}\noLine
\UIC{$\varphi_{10}$}
\AXC{$\stackrel{\Pi_{63}}{\vdots}$}\noLine
\UIC{$\varphi_{183}$}
\AXC{$\varphi_{183} \to \varphi_{10} \to \varphi_{11} \to \varphi_{139} \to \mathrm{ContIn}(f,A,c) \to \varphi_{181}$\lbl{A1}}
\BIC{$\varphi_{10} \to \varphi_{11} \to \varphi_{139} \to \mathrm{ContIn}(f,A,c) \to \varphi_{181}$}
\BIC{$\varphi_{11} \to \varphi_{139} \to \mathrm{ContIn}(f,A,c) \to \varphi_{181}$}
\BIC{$\varphi_{139} \to \mathrm{ContIn}(f,A,c) \to \varphi_{181}$}
\BIC{$\mathrm{ContIn}(f,A,c) \to \varphi_{181}$}
\end{ptree}
\noindent$\Pi_{65}$:
\begin{ptree}
\AXC{$\varphi_{110}$\lbl{A1}}
\AXC{$\varphi_{148}$\lbl{A2}}
\AXC{$\stackrel{\Pi_{64}}{\vdots}$}\noLine
\UIC{$\mathrm{ContIn}(f,A,c) \to \varphi_{181}$}
\AXC{$\varphi_{249}$\lbl{A1}}
\BIC{$\mathrm{ContIn}(f,A,c) \to \varphi_{165}$}
\AXC{$\varphi_{265}$\lbl{A1}}
\BIC{$\varphi_{148} \to \mathrm{ContIn}(f,A,c) \to \varphi_{163}$}
\BIC{$\mathrm{ContIn}(f,A,c) \to \varphi_{163}$}
\AXC{$\varphi_{261}$\lbl{A1}}
\BIC{$\varphi_{110} \to \mathrm{ContIn}(f,A,c) \to \varphi_{150}$}
\BIC{$\mathrm{ContIn}(f,A,c) \to \varphi_{150}$}
\AXC{$\varphi_{277}$\lbl{A1}}
\BIC{$(\mathrm{ContIn}(f,A,c) \to \varphi_{160}) \to \varphi_{247}$}
\end{ptree}
\noindent$\Pi_{66}$:
\begin{ptree}
\AXC{$\stackrel{\Pi_{48}}{\vdots}$}\noLine
\UIC{$(f \cdot g) \circ a(n) = (f \cdot g)(a_{n}) \to \varphi_{136}$}
\AXC{$\stackrel{\Pi_{61}}{\vdots}$}\noLine
\UIC{$\mathrm{ContIn}(f,A,c) \to \varphi_{160}$}
\AXC{$\stackrel{\Pi_{65}}{\vdots}$}\noLine
\UIC{$(\mathrm{ContIn}(f,A,c) \to \varphi_{160}) \to \varphi_{247}$}
\BIC{$\varphi_{247}$}
\BIC{$\mathrm{ContIn}(f,A,c) \to \varphi_{162}$}
\end{ptree}
\noindent$\Pi_{67}$:
\begin{ptree}
\AXC{$\stackrel{\Pi_{9}}{\vdots}$}\noLine
\UIC{$(f \cdot g) \circ a(n) = f(a_{n}) \cdot g(a_{n}) \to \varphi_{164}$}
\AXC{$\stackrel{\Pi_{46}}{\vdots}$}\noLine
\UIC{$\mathrm{ContIn}(f,A,c) \to \varphi_{174}$}
\AXC{$\stackrel{\Pi_{66}}{\vdots}$}\noLine
\UIC{$\mathrm{ContIn}(f,A,c) \to \varphi_{162}$}
\AXC{$\varphi_{282}$\lbl{A1}}
\BIC{$(\mathrm{ContIn}(f,A,c) \to \varphi_{174}) \to \varphi_{253}$}
\BIC{$\varphi_{253}$}
\BIC{$\mathrm{ContIn}(f,A,c) \to \varphi_{166}$}
\end{ptree}
\noindent$\Pi_{68}$:
\begin{ptree}
\AXC{$\stackrel{\Pi_{7}}{\vdots}$}\noLine
\UIC{$\forall x\,\forall y\,(x = y \to y = x)$}
\AXC{$\varphi_{64}$\lbl{A2}}
\BIC{$\forall y\,((f \cdot g) \circ a(n) = y \to y = (f \cdot g) \circ a(n))$}
\AXC{$\varphi_{206}$\lbl{A2}}
\BIC{$\varphi_{155}$}
\AXC{$\stackrel{\Pi_{67}}{\vdots}$}\noLine
\UIC{$\mathrm{ContIn}(f,A,c) \to \varphi_{166}$}
\AXC{$\varphi_{266}$\lbl{A1}}
\BIC{$\varphi_{155} \to \mathrm{ContIn}(f,A,c) \to \varphi_{167}$}
\BIC{$\mathrm{ContIn}(f,A,c) \to \varphi_{167}$}
\AXC{$\varphi_{246}$\lbl{A1}}
\BIC{$\varphi_{69} \to n \in \mathbb{N} \to f \circ a \otimes g \circ a(n) = (f \cdot g) \circ a(n)$}
\UIC{$\varphi_{186}$}
\end{ptree}
\noindent$\Pi_{69}$:
\begin{ptree}
\AXC{$\varphi_{55}$\lbl{[C13]}}
\AXC{$\varphi_{55} \to \varphi_{154}$\lbl{A2}}
\BIC{$\varphi_{154}$}
\AXC{$\varphi_{154} \to \varphi_{199}$\lbl{A2}}
\BIC{$\varphi_{199}$}
\AXC{$\varphi_{199} \to \varphi_{215}$\lbl{A2}}
\BIC{$\varphi_{215}$}
\AXC{$\varphi_{215} \to \varphi_{119} \to (\mathrm{ContIn}(f,A,c) \to \varphi_{169}) \to \varphi_{236}$\lbl{A1}}
\BIC{$\varphi_{119} \to (\mathrm{ContIn}(f,A,c) \to \varphi_{169}) \to \varphi_{236}$}
\end{ptree}
\noindent$\Pi_{70}$:
\begin{ptree}
\AXC{$\stackrel{\Pi_{64}}{\vdots}$}\noLine
\UIC{$\mathrm{ContIn}(f,A,c) \to \varphi_{181}$}
\AXC{$(\mathrm{ContIn}(f,A,c) \to \varphi_{181}) \to \varphi_{119}$\lbl{A1}}
\BIC{$\varphi_{119}$}
\AXC{$\stackrel{\Pi_{69}}{\vdots}$}\noLine
\UIC{$\varphi_{119} \to (\mathrm{ContIn}(f,A,c) \to \varphi_{169}) \to \varphi_{236}$}
\BIC{$(\mathrm{ContIn}(f,A,c) \to \varphi_{169}) \to \varphi_{236}$}
\end{ptree}
\noindent$\Pi_{71}$:
\begin{ptree}
\AXC{$\varphi_{69} \to \forall n\,(\varphi_{69})$\lbl{A4}}
\AXC{$\varphi_{186} \to \varphi_{189}$\lbl{A3}}
\AXC{$\stackrel{\Pi_{68}}{\vdots}$}\noLine
\UIC{$\varphi_{186}$}
\AXC{$\varphi_{186} \to \varphi_{280}$\lbl{A1}}
\BIC{$\varphi_{280}$}
\BIC{$(\varphi_{69} \to \forall n\,(\varphi_{69})) \to \mathrm{ContIn}(f,A,c) \to \varphi_{169}$}
\BIC{$\mathrm{ContIn}(f,A,c) \to \varphi_{169}$}
\AXC{$\stackrel{\Pi_{70}}{\vdots}$}\noLine
\UIC{$(\mathrm{ContIn}(f,A,c) \to \varphi_{169}) \to \varphi_{236}$}
\BIC{$\varphi_{236}$}
\end{ptree}
\noindent$\Pi_{72}$:
\begin{ptree}
\AXC{$\stackrel{\Pi_{7}}{\vdots}$}\noLine
\UIC{$\forall x\,\forall y\,(x = y \to y = x)$}
\AXC{$\varphi_{23}$\lbl{A2}}
\BIC{$\forall y\,((f \cdot g)(c) = y \to y = (f \cdot g)(c))$}
\AXC{$\forall y\,((f \cdot g)(c) = y \to y = (f \cdot g)(c)) \to \varphi_{56}$\lbl{A2}}
\BIC{$\varphi_{56}$}
\end{ptree}
\noindent$\Pi_{73}$:
\begin{ptree}
\AXC{$\forall f\,\forall A\,\forall c\,(\varphi_{191})$\lbl{DefContIn}}
\AXC{$\forall f\,\forall A\,\forall c\,(\varphi_{191}) \to \forall A\,\forall c\,(\varphi_{191})$\lbl{A2}}
\BIC{$\forall A\,\forall c\,(\varphi_{191})$}
\AXC{$\forall A\,\forall c\,(\varphi_{191}) \to \forall c\,(\varphi_{191})$\lbl{A2}}
\BIC{$\forall c\,(\varphi_{191})$}
\AXC{$\forall c\,(\varphi_{191}) \to \varphi_{191}$\lbl{A2}}
\BIC{$\varphi_{191}$}
\end{ptree}
\noindent$\Pi_{74}$:
\begin{ptree}
\AXC{$\stackrel{\Pi_{73}}{\vdots}$}\noLine
\UIC{$\varphi_{191}$}
\AXC{$\mathrm{ContIn}(f,A,c) \to \mathrm{ContIn}(f,A,c)$\lbl{A1}}
\AXC{$\varphi_{222}$\lbl{A1}}
\BIC{$\varphi_{191} \to \varphi_{6}$}
\BIC{$\varphi_{6}$}
\end{ptree}
\noindent$\Pi_{75}$:
\begin{ptree}
\AXC{$\forall X\,\forall Y\,\forall x\,(X \subseteq Y \to x \in X \to x \in Y)$\lbl{[B5]}}
\AXC{$\varphi_{46}$\lbl{A2}}
\BIC{$\forall Y\,\forall x\,(A \subseteq Y \to x \in A \to x \in Y)$}
\AXC{$\varphi_{57}$\lbl{A2}}
\BIC{$\forall x\,(A \subseteq \mathrm{dom}(g) \to x \in A \to x \in \mathrm{dom}(g))$}
\AXC{$\varphi_{65}$\lbl{A2}}
\BIC{$A \subseteq \mathrm{dom}(g) \to c \in A \to c \in \mathrm{dom}(g)$}
\end{ptree}
\noindent$\Pi_{76}$:
\begin{ptree}
\AXC{$\stackrel{\Pi_{74}}{\vdots}$}\noLine
\UIC{$\varphi_{6}$}
\AXC{$\stackrel{\Pi_{14}}{\vdots}$}\noLine
\UIC{$\varphi_{8}$}
\AXC{$\stackrel{\Pi_{75}}{\vdots}$}\noLine
\UIC{$A \subseteq \mathrm{dom}(g) \to c \in A \to c \in \mathrm{dom}(g)$}
\AXC{$(A \subseteq \mathrm{dom}(g) \to c \in A \to c \in \mathrm{dom}(g)) \to \varphi_{8} \to \varphi_{156}$\lbl{A1}}
\BIC{$\varphi_{8} \to \varphi_{156}$}
\BIC{$\varphi_{156}$}
\BIC{$\mathrm{ContIn}(f,A,c) \to \mathrm{ContIn}(g,A,c) \to \varphi_{30}$}
\end{ptree}
\noindent$\Pi_{77}$:
\begin{ptree}
\AXC{$\forall X\,\forall Y\,\forall x\,(X \subseteq Y \to x \in X \to x \in Y)$\lbl{[B5]}}
\AXC{$\varphi_{46}$\lbl{A2}}
\BIC{$\forall Y\,\forall x\,(A \subseteq Y \to x \in A \to x \in Y)$}
\AXC{$\varphi_{58}$\lbl{A2}}
\BIC{$\forall x\,(A \subseteq \mathrm{dom}(f) \to x \in A \to x \in \mathrm{dom}(f))$}
\AXC{$\varphi_{66}$\lbl{A2}}
\BIC{$A \subseteq \mathrm{dom}(f) \to c \in A \to c \in \mathrm{dom}(f)$}
\end{ptree}
\noindent$\Pi_{78}$:
\begin{ptree}
\AXC{$\stackrel{\Pi_{74}}{\vdots}$}\noLine
\UIC{$\varphi_{6}$}
\AXC{$\stackrel{\Pi_{28}}{\vdots}$}\noLine
\UIC{$\varphi_{9}$}
\AXC{$\stackrel{\Pi_{77}}{\vdots}$}\noLine
\UIC{$A \subseteq \mathrm{dom}(f) \to c \in A \to c \in \mathrm{dom}(f)$}
\AXC{$(A \subseteq \mathrm{dom}(f) \to c \in A \to c \in \mathrm{dom}(f)) \to \varphi_{9} \to \varphi_{157}$\lbl{A1}}
\BIC{$\varphi_{9} \to \varphi_{157}$}
\BIC{$\varphi_{157}$}
\BIC{$\mathrm{ContIn}(f,A,c) \to \mathrm{ContIn}(g,A,c) \to \varphi_{31}$}
\end{ptree}
\noindent$\Pi_{79}$:
\begin{ptree}
\AXC{$\stackrel{\Pi_{76}}{\vdots}$}\noLine
\UIC{$\mathrm{ContIn}(f,A,c) \to \mathrm{ContIn}(g,A,c) \to \varphi_{30}$}
\AXC{$\stackrel{\Pi_{78}}{\vdots}$}\noLine
\UIC{$\mathrm{ContIn}(f,A,c) \to \mathrm{ContIn}(g,A,c) \to \varphi_{31}$}
\AXC{$\varphi_{255}$\lbl{A1}}
\BIC{$\varphi_{237}$}
\BIC{$(\mathrm{ContIn}(f,A,c) \to \varphi_{33}) \to \varphi_{108}$}
\end{ptree}
\noindent$\Pi_{80}$:
\begin{ptree}
\AXC{$\varphi_{111}$\lbl{A2}}
\AXC{$\stackrel{\Pi_{51}}{\vdots}$}\noLine
\UIC{$\mathrm{ContIn}(f,A,c) \to \mathrm{ContIn}(g,A,c) \to \varphi_{16}$}
\AXC{$\varphi_{234}$\lbl{A1}}
\BIC{$\varphi_{111} \to \mathrm{ContIn}(f,A,c) \to \varphi_{33}$}
\BIC{$\mathrm{ContIn}(f,A,c) \to \varphi_{33}$}
\AXC{$\stackrel{\Pi_{79}}{\vdots}$}\noLine
\UIC{$(\mathrm{ContIn}(f,A,c) \to \varphi_{33}) \to \varphi_{108}$}
\BIC{$\varphi_{108}$}
\end{ptree}
\noindent$\Pi_{81}$:
\begin{ptree}
\AXC{$\stackrel{\Pi_{80}}{\vdots}$}\noLine
\UIC{$\varphi_{108}$}
\AXC{$\varphi_{27} \to \varphi_{24}$\lbl{A2}}
\AXC{$\stackrel{\Pi_{59}}{\vdots}$}\noLine
\UIC{$\mathrm{ContIn}(f,A,c) \to \mathrm{ContIn}(g,A,c) \to \varphi_{27}$}
\AXC{$\varphi_{240}$\lbl{A1}}
\BIC{$\varphi_{213}$}
\BIC{$\mathrm{ContIn}(f,A,c) \to \mathrm{ContIn}(g,A,c) \to \varphi_{24}$}
\AXC{$\varphi_{242}$\lbl{A1}}
\BIC{$\varphi_{108} \to \varphi_{134}$}
\BIC{$\varphi_{134}$}
\end{ptree}
\noindent$\Pi_{82}$:
\begin{ptree}
\AXC{$\stackrel{\Pi_{72}}{\vdots}$}\noLine
\UIC{$\varphi_{56}$}
\AXC{$\stackrel{\Pi_{81}}{\vdots}$}\noLine
\UIC{$\varphi_{134}$}
\AXC{$\varphi_{134} \to \varphi_{56} \to \varphi_{135}$\lbl{A1}}
\BIC{$\varphi_{56} \to \varphi_{135}$}
\BIC{$\varphi_{135}$}
\AXC{$\varphi_{135} \to (\mathrm{ContIn}(f,A,c) \to \varphi_{130}) \to \varphi_{230}$\lbl{A1}}
\BIC{$(\mathrm{ContIn}(f,A,c) \to \varphi_{130}) \to \varphi_{230}$}
\end{ptree}
\noindent$\Pi_{83}$:
\begin{ptree}
\AXC{$\stackrel{\Pi_{2}}{\vdots}$}\noLine
\UIC{$\varphi_{122}$}
\AXC{$\stackrel{\Pi_{5}}{\vdots}$}\noLine
\UIC{$\varphi_{123}$}
\AXC{$\stackrel{\Pi_{6}}{\vdots}$}\noLine
\UIC{$\varphi_{123} \to \varphi_{122} \to \mathrm{ContIn}(f,A,c) \to \varphi_{143}$}
\BIC{$\varphi_{122} \to \mathrm{ContIn}(f,A,c) \to \varphi_{143}$}
\BIC{$\mathrm{ContIn}(f,A,c) \to \varphi_{143}$}
\AXC{$\stackrel{\Pi_{71}}{\vdots}$}\noLine
\UIC{$\varphi_{236}$}
\BIC{$\mathrm{ContIn}(f,A,c) \to \varphi_{130}$}
\AXC{$\stackrel{\Pi_{82}}{\vdots}$}\noLine
\UIC{$(\mathrm{ContIn}(f,A,c) \to \varphi_{130}) \to \varphi_{230}$}
\BIC{$\varphi_{230}$}
\end{ptree}
\noindent$\Pi_{84}$:
\begin{ptree}
\AXC{$z = w \to \varphi_{21}$\lbl{A6}}
\UIC{$\forall w\,(z = w \to \varphi_{21})$}
\UIC{$\forall z\,\forall w\,(z = w \to \varphi_{21})$}
\AXC{$\forall z\,\forall w\,(z = w \to \varphi_{21}) \to \varphi_{109}$\lbl{A2}}
\BIC{$\varphi_{109}$}
\AXC{$\varphi_{109} \to f(c) \cdot g(c) = (f \cdot g)(c) \to \varphi_{89}$\lbl{A2}}
\BIC{$f(c) \cdot g(c) = (f \cdot g)(c) \to \varphi_{89}$}
\AXC{$\stackrel{\Pi_{83}}{\vdots}$}\noLine
\UIC{$\varphi_{230}$}
\BIC{$\mathrm{ContIn}(f,A,c) \to \varphi_{125}$}
\end{ptree}
\noindent$\Pi_{85}$:
\begin{ptree}
\AXC{$\varphi_{232}$\lbl{A3}}
\AXC{$\stackrel{\Pi_{84}}{\vdots}$}\noLine
\UIC{$\mathrm{ContIn}(f,A,c) \to \varphi_{125}$}
\AXC{$(\mathrm{ContIn}(f,A,c) \to \varphi_{125}) \to \varphi_{140}$\lbl{A1}}
\BIC{$\varphi_{140}$}
\UIC{$\forall a\,(\varphi_{140})$}
\AXC{$\forall a\,(\varphi_{140}) \to \varphi_{232} \to \varphi_{203}$\lbl{A1}}
\BIC{$\varphi_{232} \to \varphi_{203}$}
\BIC{$\varphi_{203}$}
\end{ptree}
\noindent$\Pi_{86}$:
\begin{ptree}
\AXC{$\varphi_{112}$\lbl{A2}}
\AXC{$\stackrel{\Pi_{51}}{\vdots}$}\noLine
\UIC{$\mathrm{ContIn}(f,A,c) \to \mathrm{ContIn}(g,A,c) \to \varphi_{16}$}
\AXC{$\varphi_{235}$\lbl{A1}}
\BIC{$\varphi_{112} \to \mathrm{ContIn}(f,A,c) \to \varphi_{34}$}
\BIC{$\mathrm{ContIn}(f,A,c) \to \varphi_{34}$}
\end{ptree}
\noindent$\Pi_{87}$:
\begin{ptree}
\AXC{$\forall X\,\forall Y\,\forall x\,(X \subseteq Y \to x \in X \to x \in Y)$\lbl{[B5]}}
\AXC{$\varphi_{46}$\lbl{A2}}
\BIC{$\forall Y\,\forall x\,(A \subseteq Y \to x \in A \to x \in Y)$}
\AXC{$\varphi_{57}$\lbl{A2}}
\BIC{$\forall x\,(A \subseteq \mathrm{dom}(g) \to x \in A \to x \in \mathrm{dom}(g))$}
\AXC{$\varphi_{67}$\lbl{A2}}
\BIC{$A \subseteq \mathrm{dom}(g) \to u \in A \to u \in \mathrm{dom}(g)$}
\end{ptree}
\noindent$\Pi_{88}$:
\begin{ptree}
\AXC{$\varphi_{12}$\lbl{A1}}
\AXC{$\stackrel{\Pi_{14}}{\vdots}$}\noLine
\UIC{$\varphi_{8}$}
\AXC{$\stackrel{\Pi_{87}}{\vdots}$}\noLine
\UIC{$A \subseteq \mathrm{dom}(g) \to u \in A \to u \in \mathrm{dom}(g)$}
\AXC{$(A \subseteq \mathrm{dom}(g) \to u \in A \to u \in \mathrm{dom}(g)) \to \varphi_{8} \to \varphi_{12} \to \varphi_{17}$\lbl{A1}}
\BIC{$\varphi_{8} \to \varphi_{12} \to \varphi_{17}$}
\BIC{$\varphi_{12} \to \varphi_{17}$}
\BIC{$\varphi_{17}$}
\end{ptree}
\noindent$\Pi_{89}$:
\begin{ptree}
\AXC{$\forall X\,\forall Y\,\forall x\,(X \subseteq Y \to x \in X \to x \in Y)$\lbl{[B5]}}
\AXC{$\varphi_{46}$\lbl{A2}}
\BIC{$\forall Y\,\forall x\,(A \subseteq Y \to x \in A \to x \in Y)$}
\AXC{$\varphi_{58}$\lbl{A2}}
\BIC{$\forall x\,(A \subseteq \mathrm{dom}(f) \to x \in A \to x \in \mathrm{dom}(f))$}
\AXC{$\varphi_{68}$\lbl{A2}}
\BIC{$A \subseteq \mathrm{dom}(f) \to u \in A \to u \in \mathrm{dom}(f)$}
\end{ptree}
\noindent$\Pi_{90}$:
\begin{ptree}
\AXC{$\varphi_{12}$\lbl{A1}}
\AXC{$\stackrel{\Pi_{28}}{\vdots}$}\noLine
\UIC{$\varphi_{9}$}
\AXC{$\stackrel{\Pi_{89}}{\vdots}$}\noLine
\UIC{$A \subseteq \mathrm{dom}(f) \to u \in A \to u \in \mathrm{dom}(f)$}
\AXC{$(A \subseteq \mathrm{dom}(f) \to u \in A \to u \in \mathrm{dom}(f)) \to \varphi_{9} \to \varphi_{12} \to \varphi_{18}$\lbl{A1}}
\BIC{$\varphi_{9} \to \varphi_{12} \to \varphi_{18}$}
\BIC{$\varphi_{12} \to \varphi_{18}$}
\BIC{$\varphi_{18}$}
\end{ptree}
\noindent$\Pi_{91}$:
\begin{ptree}
\AXC{$\stackrel{\Pi_{86}}{\vdots}$}\noLine
\UIC{$\mathrm{ContIn}(f,A,c) \to \varphi_{34}$}
\AXC{$\stackrel{\Pi_{88}}{\vdots}$}\noLine
\UIC{$\varphi_{17}$}
\AXC{$\stackrel{\Pi_{90}}{\vdots}$}\noLine
\UIC{$\varphi_{18}$}
\AXC{$\varphi_{18} \to \varphi_{17} \to (\mathrm{ContIn}(f,A,c) \to \varphi_{34}) \to \varphi_{25}$\lbl{A1}}
\BIC{$\varphi_{17} \to (\mathrm{ContIn}(f,A,c) \to \varphi_{34}) \to \varphi_{25}$}
\BIC{$(\mathrm{ContIn}(f,A,c) \to \varphi_{34}) \to \varphi_{25}$}
\BIC{$\varphi_{25}$}
\AXC{$\varphi_{25} \to \varphi_{26}$\lbl{A1}}
\BIC{$\varphi_{26}$}
\UIC{$\forall u\,(\varphi_{26})$}
\end{ptree}
\noindent$\Pi_{92}$:
\begin{ptree}
\AXC{$\forall u\,(\varphi_{26}) \to \varphi_{2} \to \forall u \in A\,(u \in \mathrm{dom}((f \cdot g)))$\lbl{A3}}
\AXC{$\stackrel{\Pi_{91}}{\vdots}$}\noLine
\UIC{$\forall u\,(\varphi_{26})$}
\AXC{$\varphi_{248}$\lbl{A1}}
\BIC{$(\forall u\,(\varphi_{26}) \to \varphi_{2} \to \forall u \in A\,(u \in \mathrm{dom}((f \cdot g)))) \to \varphi_{149}$}
\BIC{$\varphi_{149}$}
\end{ptree}
\noindent$\Pi_{93}$:
\begin{ptree}
\AXC{$\forall x\,\forall y\,(x \subseteq y \leftrightarrow \forall u \in x\,(u \in y))$\lbl{Def$\subseteq$}}
\AXC{$\varphi_{49}$\lbl{A2}}
\BIC{$\forall y\,(A \subseteq y \leftrightarrow \forall u \in A\,(u \in y))$}
\AXC{$\forall y\,(A \subseteq y \leftrightarrow \forall u \in A\,(u \in y)) \to \varphi_{20}$\lbl{A2}}
\BIC{$\varphi_{20}$}
\end{ptree}
\noindent$\Pi_{94}$:
\begin{ptree}
\AXC{$\mathrm{ContIn}(f,A,c) \land \mathrm{ContIn}(g,A,c) \to \varphi_{2}$\lbl{A4}}
\AXC{$\stackrel{\Pi_{92}}{\vdots}$}\noLine
\UIC{$\varphi_{149}$}
\BIC{$\mathrm{ContIn}(f,A,c) \to \varphi_{13}$}
\AXC{$\stackrel{\Pi_{93}}{\vdots}$}\noLine
\UIC{$\varphi_{20}$}
\AXC{$\varphi_{20} \to (\mathrm{ContIn}(f,A,c) \to \varphi_{13}) \to \varphi_{14}$\lbl{A1}}
\BIC{$(\mathrm{ContIn}(f,A,c) \to \varphi_{13}) \to \varphi_{14}$}
\BIC{$\varphi_{14}$}
\end{ptree}
\noindent$\Pi_{95}$:
\begin{ptree}
\AXC{$\stackrel{\Pi_{74}}{\vdots}$}\noLine
\UIC{$\varphi_{6}$}
\AXC{$\stackrel{\Pi_{94}}{\vdots}$}\noLine
\UIC{$\varphi_{14}$}
\AXC{$\stackrel{\Pi_{62}}{\vdots}$}\noLine
\UIC{$\varphi_{10}$}
\AXC{$\varphi_{161}$\lbl{[G2]}}
\AXC{$\varphi_{161} \to \forall A\,(\varphi_{193})$\lbl{A2}}
\BIC{$\forall A\,(\varphi_{193})$}
\AXC{$\forall A\,(\varphi_{193}) \to \varphi_{193}$\lbl{A2}}
\BIC{$\varphi_{193}$}
\AXC{$\varphi_{250}$\lbl{A2}}
\BIC{$\mathrm{RFn}((f \cdot g)) \to A \subseteq \mathrm{dom}((f \cdot g)) \to c \in A \to \varphi_{137}$}
\AXC{$\varphi_{273}$\lbl{A1}}
\BIC{$\varphi_{10} \to \varphi_{14} \to \varphi_{6} \to \varphi_{198}$}
\BIC{$\varphi_{14} \to \varphi_{6} \to \varphi_{198}$}
\BIC{$\varphi_{6} \to \varphi_{198}$}
\BIC{$\varphi_{198}$}
\end{ptree}
\noindent 主木:
\begin{ptree}
\AXC{$\mathrm{ContIn}(f,A,c) \land \mathrm{ContIn}(g,A,c) \to \varphi_{1}$\lbl{A4}}
\AXC{$\stackrel{\Pi_{85}}{\vdots}$}\noLine
\UIC{$\varphi_{203}$}
\BIC{$\mathrm{ContIn}(f,A,c) \to \varphi_{127}$}
\AXC{$\stackrel{\Pi_{95}}{\vdots}$}\noLine
\UIC{$\varphi_{198}$}
\BIC{$\mathrm{ContIn}(f,A,c) \to \varphi_{3}$}
\UIC{$\forall c\,(\mathrm{ContIn}(f,A,c) \to \varphi_{3})$}
\UIC{$\forall A\,\forall c\,(\mathrm{ContIn}(f,A,c) \to \varphi_{3})$}
\UIC{$\forall g\,\forall A\,\forall c\,(\mathrm{ContIn}(f,A,c) \to \varphi_{3})$}
\UIC{$\forall f\,\forall g\,\forall A\,\forall c\,(\mathrm{ContIn}(f,A,c) \to \varphi_{3})$}
\end{ptree}
\begin{abbrevlist}
\item $\varphi_{1} :\equiv \forall a\,(\mathrm{ContIn}(f,A,c) \land \mathrm{ContIn}(g,A,c))$
\item $\varphi_{2} :\equiv \forall u\,(\mathrm{ContIn}(f,A,c) \land \mathrm{ContIn}(g,A,c))$
\item $\varphi_{3} :\equiv \mathrm{ContIn}(g,A,c) \to \mathrm{ContIn}((f \cdot g),A,c)$
\item $\varphi_{4} :\equiv \mathrm{ContIn}(f,A,c) \to \mathrm{ContIn}(g,A,c) \to \mathrm{RFn}(g)$
\item $\varphi_{5} :\equiv \mathrm{ContIn}(f,A,c) \to \mathrm{ContIn}(g,A,c) \to \mathrm{RFn}(f)$
\item $\varphi_{6} :\equiv \mathrm{ContIn}(f,A,c) \to \mathrm{ContIn}(g,A,c) \to c \in A$
\item $\varphi_{7} :\equiv \mathrm{ContIn}(f,A,c) \to \mathrm{ContIn}(g,A,c) \to \mathrm{ContIn}(g,A,c)$
\item $\varphi_{8} :\equiv \mathrm{ContIn}(f,A,c) \to \mathrm{ContIn}(g,A,c) \to A \subseteq \mathrm{dom}(g)$
\item $\varphi_{9} :\equiv \mathrm{ContIn}(f,A,c) \to \mathrm{ContIn}(g,A,c) \to A \subseteq \mathrm{dom}(f)$
\item $\varphi_{10} :\equiv \mathrm{ContIn}(f,A,c) \to \mathrm{ContIn}(g,A,c) \to \mathrm{RFn}((f \cdot g))$
\item $\varphi_{11} :\equiv \mathrm{ContIn}(f,A,c) \to \mathrm{ContIn}(g,A,c) \to \mathrm{Seq}(a) \to \mathrm{Seq}(a)$
\item $\varphi_{12} :\equiv \mathrm{ContIn}(f,A,c) \to \mathrm{ContIn}(g,A,c) \to u \in A \to u \in A$
\item $\varphi_{13} :\equiv \mathrm{ContIn}(g,A,c) \to \forall u \in A\,(u \in \mathrm{dom}((f \cdot g)))$
\item $\varphi_{14} :\equiv \mathrm{ContIn}(f,A,c) \to \mathrm{ContIn}(g,A,c) \to A \subseteq \mathrm{dom}((f \cdot g))$
\item $\varphi_{15} :\equiv \forall n \in \mathbb{N}\,(a_{n} \in A) \to \mathrm{Conv}(a,c) \to \mathrm{Conv}(a,c)$
\item $\varphi_{16} :\equiv \forall x\,(x \in \mathrm{dom}((f \cdot g)) \leftrightarrow x \in \mathrm{dom}(f) \land x \in \mathrm{dom}(g))$
\item $\varphi_{17} :\equiv \mathrm{ContIn}(f,A,c) \to \mathrm{ContIn}(g,A,c) \to u \in A \to u \in \mathrm{dom}(g)$
\item $\varphi_{18} :\equiv \mathrm{ContIn}(f,A,c) \to \mathrm{ContIn}(g,A,c) \to u \in A \to u \in \mathrm{dom}(f)$
\item $\varphi_{19} :\equiv \mathrm{Conv}(a,c) \to \mathrm{Conv}((f \cdot g) \circ a,(f \cdot g)(c))$
\item $\varphi_{20} :\equiv A \subseteq \mathrm{dom}((f \cdot g)) \leftrightarrow \forall u \in A\,(u \in \mathrm{dom}((f \cdot g)))$
\item $\varphi_{21} :\equiv \mathrm{Conv}((f \cdot g) \circ a,z) \to \mathrm{Conv}((f \cdot g) \circ a,w)$
\item $\varphi_{22} :\equiv \mathrm{Conv}(a,c) \to \mathrm{Conv}((f \cdot g) \circ a,f(c) \cdot g(c))$
\item $\varphi_{23} :\equiv \forall x\,\forall y\,(x = y \to y = x) \to \forall y\,((f \cdot g)(c) = y \to y = (f \cdot g)(c))$
\item $\varphi_{24} :\equiv c \in \mathrm{dom}((f \cdot g)) \to (f \cdot g)(c) = f(c) \cdot g(c)$
\item $\varphi_{25} :\equiv \mathrm{ContIn}(f,A,c) \to \mathrm{ContIn}(g,A,c) \to u \in A \to u \in \mathrm{dom}((f \cdot g))$
\item $\varphi_{26} :\equiv \mathrm{ContIn}(f,A,c) \land \mathrm{ContIn}(g,A,c) \to u \in A \to u \in \mathrm{dom}((f \cdot g))$
\item $\varphi_{27} :\equiv \forall x \in \mathrm{dom}((f \cdot g))\,((f \cdot g)(x) = f(x) \cdot g(x))$
\item $\varphi_{28} :\equiv \forall n \in \mathbb{N}\,(a_{n} \in A) \to \mathrm{Conv}(a,c) \to c \in \mathrm{dom}((f \cdot g))$
\item $\varphi_{29} :\equiv \forall n \in \mathbb{N}\,(a_{n} \in A) \to \mathrm{Conv}(a,c) \to n \in \mathbb{N} \to n \in \mathbb{N}$
\item $\varphi_{30} :\equiv \mathrm{Seq}(a) \to \forall n \in \mathbb{N}\,(a_{n} \in A) \to \mathrm{Conv}(a,c) \to c \in \mathrm{dom}(g)$
\item $\varphi_{31} :\equiv \mathrm{Seq}(a) \to \forall n \in \mathbb{N}\,(a_{n} \in A) \to \mathrm{Conv}(a,c) \to c \in \mathrm{dom}(f)$
\item $\varphi_{32} :\equiv \forall a\,\forall \alpha\,\forall b\,\forall \beta\,(\mathrm{Conv}(a,\alpha) \to \mathrm{Conv}(b,\beta) \to \mathrm{Conv}(a \otimes b,\alpha \cdot \beta))$
\item $\varphi_{33} :\equiv \mathrm{ContIn}(g,A,c) \to (c \in \mathrm{dom}((f \cdot g)) \leftrightarrow c \in \mathrm{dom}(f) \land c \in \mathrm{dom}(g))$
\item $\varphi_{34} :\equiv \mathrm{ContIn}(g,A,c) \to (u \in \mathrm{dom}((f \cdot g)) \leftrightarrow u \in \mathrm{dom}(f) \land u \in \mathrm{dom}(g))$
\item $\varphi_{35} :\equiv \forall n \in \mathbb{N}\,(a_{n} \in A) \to \mathrm{Conv}(a,c) \to \mathrm{Seq}((f \cdot g) \circ a)$
\item $\varphi_{36} :\equiv \forall n \in \mathbb{N}\,(a_{n} \in A) \to \mathrm{Conv}(a,c) \to n \in \mathbb{N} \to a_{n} \in A$
\item $\varphi_{37} :\equiv \mathrm{Seq}(a) \to \forall n \in \mathbb{N}\,(a_{n} \in A) \to \mathrm{Conv}(a,c) \to \mathrm{Seq}(g \circ a)$
\item $\varphi_{38} :\equiv \mathrm{Seq}(a) \to \forall n \in \mathbb{N}\,(a_{n} \in A) \to \mathrm{Conv}(a,c) \to \mathrm{Seq}(f \circ a)$
\item $\varphi_{39} :\equiv \forall n \in \mathbb{N}\,(a_{n} \in A) \to \mathrm{Conv}(a,c) \to \mathrm{Conv}(g \circ a,g(c))$
\item $\varphi_{40} :\equiv \forall n \in \mathbb{N}\,(a_{n} \in A) \to \mathrm{Conv}(a,c) \to \mathrm{Conv}(f \circ a,f(c))$
\item $\varphi_{41} :\equiv \mathrm{Conv}(a,c) \to n \in \mathbb{N} \to n \in \mathbb{N} \to g \circ a(n) = g(a_{n})$
\item $\varphi_{42} :\equiv \mathrm{Conv}(a,c) \to n \in \mathbb{N} \to n \in \mathbb{N} \to f \circ a(n) = f(a_{n})$
\item $\varphi_{43} :\equiv a_{n} \in \mathrm{dom}((f \cdot g)) \leftrightarrow a_{n} \in \mathrm{dom}(f) \land a_{n} \in \mathrm{dom}(g)$
\item $\varphi_{44} :\equiv \mathrm{Conv}(a,c) \to \mathrm{Conv}(f \circ a \otimes g \circ a,f(c) \cdot g(c))$
\item $\varphi_{45} :\equiv \forall y\,\forall z\,((f \cdot g) \circ a(n) = y \to y = z \to (f \cdot g) \circ a(n) = z)$
\item $\varphi_{46} :\equiv \forall X\,\forall Y\,\forall x\,(X \subseteq Y \to x \in X \to x \in Y) \to \forall Y\,\forall x\,(A \subseteq Y \to x \in A \to x \in Y)$
\item $\varphi_{47} :\equiv \forall a\,\forall \alpha\,(\mathrm{Conv}(a,\alpha) \to \mathrm{Seq}(a)) \to \forall \alpha\,(\mathrm{Conv}(g \circ a,\alpha) \to \mathrm{Seq}(g \circ a))$
\item $\varphi_{48} :\equiv \forall a\,\forall \alpha\,(\mathrm{Conv}(a,\alpha) \to \mathrm{Seq}(a)) \to \forall \alpha\,(\mathrm{Conv}(f \circ a,\alpha) \to \mathrm{Seq}(f \circ a))$
\item $\varphi_{49} :\equiv \forall x\,\forall y\,(x \subseteq y \leftrightarrow \forall u \in x\,(u \in y)) \to \forall y\,(A \subseteq y \leftrightarrow \forall u \in A\,(u \in y))$
\item $\varphi_{50} :\equiv \mathrm{Conv}(a,c) \to n \in \mathbb{N} \to \forall n \in \mathbb{N}\,(g \circ a(n) = g(a_{n}))$
\item $\varphi_{51} :\equiv \mathrm{Conv}(a,c) \to n \in \mathbb{N} \to \forall n \in \mathbb{N}\,(f \circ a(n) = f(a_{n}))$
\item $\varphi_{52} :\equiv \forall n \in \mathbb{N}\,(a_{n} \in A) \to \mathrm{Conv}(a,c) \to n \in \mathbb{N} \to a_{n} \in \mathrm{dom}(g)$
\item $\varphi_{53} :\equiv \forall n \in \mathbb{N}\,(a_{n} \in A) \to \mathrm{Conv}(a,c) \to n \in \mathbb{N} \to a_{n} \in \mathrm{dom}(f)$
\item $\varphi_{54} :\equiv \mathrm{Conv}(a,c) \to n \in \mathbb{N} \to (f \cdot g) \circ a(n) = (f \cdot g)(a_{n})$
\item $\varphi_{55} :\equiv \forall a\,\forall b\,\forall \alpha\,(\mathrm{Seq}(b) \to \forall n \in \mathbb{N}\,(a_{n} = b_{n}) \to \mathrm{Conv}(a,\alpha) \to \mathrm{Conv}(b,\alpha))$
\item $\varphi_{56} :\equiv (f \cdot g)(c) = f(c) \cdot g(c) \to f(c) \cdot g(c) = (f \cdot g)(c)$
\item $\varphi_{57} :\equiv \forall Y\,\forall x\,(A \subseteq Y \to x \in A \to x \in Y) \to \forall x\,(A \subseteq \mathrm{dom}(g) \to x \in A \to x \in \mathrm{dom}(g))$
\item $\varphi_{58} :\equiv \forall Y\,\forall x\,(A \subseteq Y \to x \in A \to x \in Y) \to \forall x\,(A \subseteq \mathrm{dom}(f) \to x \in A \to x \in \mathrm{dom}(f))$
\item $\varphi_{59} :\equiv \forall n \in \mathbb{N}\,(a_{n} \in A) \to \mathrm{Conv}(a,c) \to \forall n \in \mathbb{N}\,(a_{n} \in \mathrm{dom}(g))$
\item $\varphi_{60} :\equiv \forall n \in \mathbb{N}\,(a_{n} \in A) \to \mathrm{Conv}(a,c) \to \forall n \in \mathbb{N}\,(a_{n} \in \mathrm{dom}(f))$
\item $\varphi_{61} :\equiv n \in \mathbb{N} \to \forall n \in \mathbb{N}\,((f \cdot g) \circ a(n) = (f \cdot g)(a_{n}))$
\item $\varphi_{62} :\equiv \forall n \in \mathbb{N}\,(a_{n} \in A) \to \mathrm{Conv}(a,c) \to (f \cdot g)(c) = f(c) \cdot g(c)$
\item $\varphi_{63} :\equiv \forall n \in \mathbb{N}\,(a_{n} \in A) \to \mathrm{Conv}(a,c) \to f(c) \cdot g(c) = (f \cdot g)(c)$
\item $\varphi_{64} :\equiv \forall x\,\forall y\,(x = y \to y = x) \to \forall y\,((f \cdot g) \circ a(n) = y \to y = (f \cdot g) \circ a(n))$
\item $\varphi_{65} :\equiv \forall x\,(A \subseteq \mathrm{dom}(g) \to x \in A \to x \in \mathrm{dom}(g)) \to A \subseteq \mathrm{dom}(g) \to c \in A \to c \in \mathrm{dom}(g)$
\item $\varphi_{66} :\equiv \forall x\,(A \subseteq \mathrm{dom}(f) \to x \in A \to x \in \mathrm{dom}(f)) \to A \subseteq \mathrm{dom}(f) \to c \in A \to c \in \mathrm{dom}(f)$
\item $\varphi_{67} :\equiv \forall x\,(A \subseteq \mathrm{dom}(g) \to x \in A \to x \in \mathrm{dom}(g)) \to A \subseteq \mathrm{dom}(g) \to u \in A \to u \in \mathrm{dom}(g)$
\item $\varphi_{68} :\equiv \forall x\,(A \subseteq \mathrm{dom}(f) \to x \in A \to x \in \mathrm{dom}(f)) \to A \subseteq \mathrm{dom}(f) \to u \in A \to u \in \mathrm{dom}(f)$
\item $\varphi_{69} :\equiv \mathrm{ContIn}(f,A,c) \land \mathrm{ContIn}(g,A,c) \land \mathrm{Seq}(a) \land \forall n \in \mathbb{N}\,(a_{n} \in A) \land \mathrm{Conv}(a,c)$
\item $\varphi_{70} :\equiv \mathrm{Conv}(a,c) \to \mathrm{Seq}(g \circ a) \land \forall n \in \mathbb{N}\,(g \circ a(n) = g(a_{n}))$
\item $\varphi_{71} :\equiv \mathrm{Conv}(a,c) \to \mathrm{Seq}(f \circ a) \land \forall n \in \mathbb{N}\,(f \circ a(n) = f(a_{n}))$
\item $\varphi_{72} :\equiv \mathrm{Conv}(a,c) \to n \in \mathbb{N} \to (f \cdot g)(a_{n}) = f(a_{n}) \cdot g(a_{n})$
\item $\varphi_{73} :\equiv \mathrm{ContIn}(g,A,c) \to \mathrm{Seq}(a) \to \forall n \in \mathbb{N}\,(a_{n} \in A) \to n \in \mathbb{N} \to a_{n} \in A$
\item $\varphi_{74} :\equiv \forall n \in \mathbb{N}\,(a_{n} \in A) \to \mathrm{Conv}(a,c) \to n \in \mathbb{N} \to a_{n} \in \mathrm{dom}((f \cdot g))$
\item $\varphi_{75} :\equiv \forall n \in \mathbb{N}\,(a_{n} \in \mathrm{dom}((f \cdot g))) \to n \in \mathbb{N} \to a_{n} \in \mathrm{dom}((f \cdot g))$
\item $\varphi_{76} :\equiv \forall \alpha\,\forall b\,\forall \beta\,(\mathrm{Conv}(f \circ a,\alpha) \to \mathrm{Conv}(b,\beta) \to \mathrm{Conv}(f \circ a \otimes b,\alpha \cdot \beta))$
\item $\varphi_{77} :\equiv \forall x\,(x = x) \to f(a_{n}) \cdot g \circ a(n) = f(a_{n}) \cdot g \circ a(n)$
\item $\varphi_{78} :\equiv \mathrm{ContIn}(g,A,c) \to \mathrm{Seq}(a) \to \forall n \in \mathbb{N}\,(a_{n} \in A) \to \forall n \in \mathbb{N}\,(a_{n} \in A)$
\item $\varphi_{79} :\equiv \forall n \in \mathbb{N}\,(a_{n} \in A) \to \mathrm{Conv}(a,c) \to \forall n \in \mathbb{N}\,(a_{n} \in \mathrm{dom}((f \cdot g)))$
\item $\varphi_{80} :\equiv \mathrm{Conv}(a,c) \to n \in \mathbb{N} \to (f \cdot g) \circ a(n) = f(a_{n}) \cdot g(a_{n})$
\item $\varphi_{81} :\equiv \mathrm{Conv}(a,c) \to n \in \mathbb{N} \to n \in \mathbb{N} \to (f \cdot g) \circ a(n) = (f \cdot g)(a_{n})$
\item $\varphi_{82} :\equiv a_{n} \in \mathrm{dom}((f \cdot g)) \to (f \cdot g)(a_{n}) = f(a_{n}) \cdot g(a_{n})$
\item $\varphi_{83} :\equiv \mathrm{Conv}(g \circ a,g(c)) \to \mathrm{Conv}(f \circ a \otimes g \circ a,f(c) \cdot g(c))$
\item $\varphi_{84} :\equiv \forall n \in \mathbb{N}\,(a_{n} \in A) \to \mathrm{Conv}(a,c) \to n \in \mathbb{N} \to g \circ a(n) = g(a_{n})$
\item $\varphi_{85} :\equiv \forall n \in \mathbb{N}\,(a_{n} \in A) \to \mathrm{Conv}(a,c) \to n \in \mathbb{N} \to f \circ a(n) = f(a_{n})$
\item $\varphi_{86} :\equiv \mathrm{Conv}(a,c) \to n \in \mathbb{N} \to (f \cdot g) \circ a(n) = f \circ a \otimes g \circ a(n)$
\item $\varphi_{87} :\equiv \mathrm{Conv}(a,c) \to n \in \mathbb{N} \to f \circ a \otimes g \circ a(n) = (f \cdot g) \circ a(n)$
\item $\varphi_{88} :\equiv \forall b\,\forall \beta\,(\mathrm{Conv}(f \circ a,f(c)) \to \mathrm{Conv}(b,\beta) \to \mathrm{Conv}(f \circ a \otimes b,f(c) \cdot \beta))$
\item $\varphi_{89} :\equiv \mathrm{Conv}((f \cdot g) \circ a,f(c) \cdot g(c)) \to \mathrm{Conv}((f \cdot g) \circ a,(f \cdot g)(c))$
\item $\varphi_{90} :\equiv n \in \mathbb{N} \to f \circ a(n) \cdot g \circ a(n) = f(a_{n}) \cdot g \circ a(n)$
\item $\varphi_{91} :\equiv \mathrm{ContIn}(f,A,c) \to \mathrm{ContIn}(g,A,c) \to \mathrm{Seq}(a) \to \varphi_{15}$
\item $\varphi_{92} :\equiv \forall x\,(x = x) \to f \circ a(n) \cdot g \circ a(n) = f \circ a(n) \cdot g \circ a(n)$
\item $\varphi_{93} :\equiv \mathrm{Seq}((f \cdot g) \circ a) \land \forall n \in \mathbb{N}\,((f \cdot g) \circ a(n) = (f \cdot g)(a_{n}))$
\item $\varphi_{94} :\equiv \forall \alpha\,(\mathrm{Conv}(g \circ a,\alpha) \to \mathrm{Seq}(g \circ a)) \to \mathrm{Conv}(g \circ a,g(c)) \to \mathrm{Seq}(g \circ a)$
\item $\varphi_{95} :\equiv \forall \alpha\,(\mathrm{Conv}(f \circ a,\alpha) \to \mathrm{Seq}(f \circ a)) \to \mathrm{Conv}(f \circ a,f(c)) \to \mathrm{Seq}(f \circ a)$
\item $\varphi_{96} :\equiv \mathrm{Conv}(a,c) \to \forall n \in \mathbb{N}\,(f \circ a \otimes g \circ a(n) = (f \cdot g) \circ a(n))$
\item $\varphi_{97} :\equiv \forall x\,(A \subseteq \mathrm{dom}(g) \to x \in A \to x \in \mathrm{dom}(g)) \to A \subseteq \mathrm{dom}(g) \to a_{n} \in A \to a_{n} \in \mathrm{dom}(g)$
\item $\varphi_{98} :\equiv \forall x\,(A \subseteq \mathrm{dom}(f) \to x \in A \to x \in \mathrm{dom}(f)) \to A \subseteq \mathrm{dom}(f) \to a_{n} \in A \to a_{n} \in \mathrm{dom}(f)$
\item $\varphi_{99} :\equiv \forall a\,\forall b\,(\mathrm{Seq}(a) \to \mathrm{Seq}(b) \to \mathrm{Seq}(a \otimes b) \land \forall n \in \mathbb{N}\,(a \otimes b(n) = a_{n} \cdot b_{n}))$
\item $\varphi_{100} :\equiv \mathrm{Conv}(a,c) \to n \in \mathbb{N} \to f(a_{n}) \cdot g \circ a(n) = f(a_{n}) \cdot g(a_{n})$
\item $\varphi_{101} :\equiv \forall n \in \mathbb{N}\,(f \circ a \otimes g \circ a(n) = f \circ a(n) \cdot g \circ a(n))$
\item $\varphi_{102} :\equiv \forall y\,\forall z\,(f \circ a \otimes g \circ a(n) = y \to y = z \to f \circ a \otimes g \circ a(n) = z)$
\item $\varphi_{103} :\equiv \mathrm{Conv}(a,c) \to n \in \mathbb{N} \to f \circ a \otimes g \circ a(n) = f(a_{n}) \cdot g(a_{n})$
\item $\varphi_{104} :\equiv \mathrm{Conv}(a,c) \to n \in \mathbb{N} \to f(a_{n}) \cdot g(a_{n}) = f \circ a \otimes g \circ a(n)$
\item $\varphi_{105} :\equiv \forall n \in \mathbb{N}\,(g \circ a(n) = g(a_{n})) \to n \in \mathbb{N} \to g \circ a(n) = g(a_{n})$
\item $\varphi_{106} :\equiv \forall n \in \mathbb{N}\,(f \circ a(n) = f(a_{n})) \to n \in \mathbb{N} \to f \circ a(n) = f(a_{n})$
\item $\varphi_{107} :\equiv \mathrm{Conv}(a,c) \to n \in \mathbb{N} \to f \circ a \otimes g \circ a(n) = f(a_{n}) \cdot g \circ a(n)$
\item $\varphi_{108} :\equiv \mathrm{ContIn}(f,A,c) \to \mathrm{ContIn}(g,A,c) \to \mathrm{Seq}(a) \to \varphi_{28}$
\item $\varphi_{109} :\equiv \forall w\,(f(c) \cdot g(c) = w \to \mathrm{Conv}((f \cdot g) \circ a,f(c) \cdot g(c)) \to \mathrm{Conv}((f \cdot g) \circ a,w))$
\item $\varphi_{110} :\equiv \mathrm{ContIn}(f,A,c) \to \mathrm{ContIn}(g,A,c) \to \mathrm{Seq}(a) \to \varphi_{29}$
\item $\varphi_{111} :\equiv \varphi_{16} \to (c \in \mathrm{dom}((f \cdot g)) \leftrightarrow c \in \mathrm{dom}(f) \land c \in \mathrm{dom}(g))$
\item $\varphi_{112} :\equiv \varphi_{16} \to (u \in \mathrm{dom}((f \cdot g)) \leftrightarrow u \in \mathrm{dom}(f) \land u \in \mathrm{dom}(g))$
\item $\varphi_{113} :\equiv \forall n \in \mathbb{N}\,(a_{n} \in \mathrm{dom}(g)) \to \mathrm{Seq}(g \circ a) \land \forall n \in \mathbb{N}\,(g \circ a(n) = g(a_{n}))$
\item $\varphi_{114} :\equiv \forall n \in \mathbb{N}\,(a_{n} \in \mathrm{dom}(f)) \to \mathrm{Seq}(f \circ a) \land \forall n \in \mathbb{N}\,(f \circ a(n) = f(a_{n}))$
\item $\varphi_{115} :\equiv n \in \mathbb{N} \to n \in \mathbb{N} \to f \circ a \otimes g \circ a(n) = f \circ a(n) \cdot g \circ a(n)$
\item $\varphi_{116} :\equiv \mathrm{Conv}(a,c) \to n \in \mathbb{N} \to f \circ a \otimes g \circ a(n) = f \circ a(n) \cdot g \circ a(n)$
\item $\varphi_{117} :\equiv \forall f\,\forall A\,\forall c\,\forall a\,(\mathrm{ContIn}(f,A,c) \to \mathrm{Seq}(a) \to \varphi_{40})$
\item $\varphi_{118} :\equiv \forall x\,\forall y\,(x = y \to y = x) \to \forall y\,(f \circ a \otimes g \circ a(n) = y \to y = f \circ a \otimes g \circ a(n))$
\item $\varphi_{119} :\equiv \mathrm{ContIn}(f,A,c) \to \mathrm{ContIn}(g,A,c) \to \mathrm{Seq}(a) \to \varphi_{35}$
\item $\varphi_{120} :\equiv \mathrm{Conv}(f \circ a \otimes g \circ a,f(c) \cdot g(c)) \to \mathrm{Conv}((f \cdot g) \circ a,f(c) \cdot g(c))$
\item $\varphi_{121} :\equiv \mathrm{ContIn}(f,A,c) \to \mathrm{ContIn}(g,A,c) \to \mathrm{Seq}(a) \to \varphi_{36}$
\item $\varphi_{122} :\equiv \mathrm{ContIn}(f,A,c) \to \mathrm{ContIn}(g,A,c) \to \mathrm{Seq}(a) \to \varphi_{39}$
\item $\varphi_{123} :\equiv \mathrm{ContIn}(f,A,c) \to \mathrm{ContIn}(g,A,c) \to \mathrm{Seq}(a) \to \varphi_{40}$
\item $\varphi_{124} :\equiv \forall \beta\,(\mathrm{Conv}(f \circ a,f(c)) \to \mathrm{Conv}(g \circ a,\beta) \to \mathrm{Conv}(f \circ a \otimes g \circ a,f(c) \cdot \beta))$
\item $\varphi_{125} :\equiv \mathrm{ContIn}(g,A,c) \to \mathrm{Seq}(a) \to \forall n \in \mathbb{N}\,(a_{n} \in A) \to \varphi_{19}$
\item $\varphi_{126} :\equiv \forall z\,((f \cdot g) \circ a(n) = (f \cdot g)(a_{n}) \to (f \cdot g)(a_{n}) = z \to (f \cdot g) \circ a(n) = z)$
\item $\varphi_{127} :\equiv \mathrm{ContIn}(g,A,c) \to \forall a\,(\mathrm{Seq}(a) \to \forall n \in \mathbb{N}\,(a_{n} \in A) \to \varphi_{19})$
\item $\varphi_{128} :\equiv \mathrm{ContIn}(f,A,c) \to \mathrm{ContIn}(g,A,c) \to \mathrm{Seq}(a) \to \varphi_{52}$
\item $\varphi_{129} :\equiv \mathrm{ContIn}(f,A,c) \to \mathrm{ContIn}(g,A,c) \to \mathrm{Seq}(a) \to \varphi_{53}$
\item $\varphi_{130} :\equiv \mathrm{ContIn}(g,A,c) \to \mathrm{Seq}(a) \to \forall n \in \mathbb{N}\,(a_{n} \in A) \to \varphi_{22}$
\item $\varphi_{131} :\equiv \mathrm{ContIn}(f,A,c) \to \mathrm{ContIn}(g,A,c) \to \mathrm{Seq}(a) \to \varphi_{59}$
\item $\varphi_{132} :\equiv \mathrm{ContIn}(f,A,c) \to \mathrm{ContIn}(g,A,c) \to \mathrm{Seq}(a) \to \varphi_{60}$
\item $\varphi_{133} :\equiv f(a_{n}) \cdot g \circ a(n) = f(a_{n}) \cdot z \to f(a_{n}) \cdot g \circ a(n) = f(a_{n}) \cdot w$
\item $\varphi_{134} :\equiv \mathrm{ContIn}(f,A,c) \to \mathrm{ContIn}(g,A,c) \to \mathrm{Seq}(a) \to \varphi_{62}$
\item $\varphi_{135} :\equiv \mathrm{ContIn}(f,A,c) \to \mathrm{ContIn}(g,A,c) \to \mathrm{Seq}(a) \to \varphi_{63}$
\item $\varphi_{136} :\equiv (f \cdot g)(a_{n}) = f(a_{n}) \cdot g(a_{n}) \to (f \cdot g) \circ a(n) = f(a_{n}) \cdot g(a_{n})$
\item $\varphi_{137} :\equiv \forall a\,(\mathrm{Seq}(a) \to \forall n \in \mathbb{N}\,(a_{n} \in A) \to \varphi_{19}) \to \mathrm{ContIn}((f \cdot g),A,c)$
\item $\varphi_{138} :\equiv \mathrm{ContIn}(f,A,c) \to \mathrm{ContIn}(g,A,c) \to \mathrm{Seq}(a) \to \varphi_{74}$
\item $\varphi_{139} :\equiv \mathrm{ContIn}(f,A,c) \to \mathrm{ContIn}(g,A,c) \to \mathrm{Seq}(a) \to \varphi_{79}$
\item $\varphi_{140} :\equiv \mathrm{ContIn}(f,A,c) \land \mathrm{ContIn}(g,A,c) \to \mathrm{Seq}(a) \to \forall n \in \mathbb{N}\,(a_{n} \in A) \to \varphi_{19}$
\item $\varphi_{141} :\equiv \mathrm{ContIn}(g,A,c) \to \mathrm{Seq}(a) \to \forall n \in \mathbb{N}\,(a_{n} \in A) \to \varphi_{41}$
\item $\varphi_{142} :\equiv \mathrm{ContIn}(g,A,c) \to \mathrm{Seq}(a) \to \forall n \in \mathbb{N}\,(a_{n} \in A) \to \varphi_{42}$
\item $\varphi_{143} :\equiv \mathrm{ContIn}(g,A,c) \to \mathrm{Seq}(a) \to \forall n \in \mathbb{N}\,(a_{n} \in A) \to \varphi_{44}$
\item $\varphi_{144} :\equiv \mathrm{ContIn}(g,A,c) \to \mathrm{Seq}(a) \to \forall n \in \mathbb{N}\,(a_{n} \in A) \to \varphi_{50}$
\item $\varphi_{145} :\equiv \mathrm{ContIn}(g,A,c) \to \mathrm{Seq}(a) \to \forall n \in \mathbb{N}\,(a_{n} \in A) \to \varphi_{51}$
\item $\varphi_{146} :\equiv \mathrm{ContIn}(f,A,c) \to \mathrm{ContIn}(g,A,c) \to \mathrm{Seq}(a) \to \varphi_{84}$
\item $\varphi_{147} :\equiv \mathrm{ContIn}(f,A,c) \to \mathrm{ContIn}(g,A,c) \to \mathrm{Seq}(a) \to \varphi_{85}$
\item $\varphi_{148} :\equiv \forall n \in \mathbb{N}\,((f \cdot g) \circ a(n) = (f \cdot g)(a_{n})) \to n \in \mathbb{N} \to (f \cdot g) \circ a(n) = (f \cdot g)(a_{n})$
\item $\varphi_{149} :\equiv (\mathrm{ContIn}(f,A,c) \land \mathrm{ContIn}(g,A,c) \to \varphi_{2}) \to \mathrm{ContIn}(f,A,c) \to \varphi_{13}$
\item $\varphi_{150} :\equiv \mathrm{ContIn}(g,A,c) \to \mathrm{Seq}(a) \to \forall n \in \mathbb{N}\,(a_{n} \in A) \to \varphi_{54}$
\item $\varphi_{151} :\equiv \forall z\,((f \cdot g) \circ a(n) = f(a_{n}) \cdot g(a_{n}) \to f(a_{n}) \cdot g(a_{n}) = z \to (f \cdot g) \circ a(n) = z)$
\item $\varphi_{152} :\equiv f \circ a(n) \cdot g \circ a(n) = z \cdot g \circ a(n) \to f \circ a(n) \cdot g \circ a(n) = w \cdot g \circ a(n)$
\item $\varphi_{153} :\equiv \forall b\,(\mathrm{Seq}(f \circ a) \to \mathrm{Seq}(b) \to \mathrm{Seq}(f \circ a \otimes b) \land \forall n \in \mathbb{N}\,(f \circ a \otimes b(n) = f \circ a(n) \cdot b_{n}))$
\item $\varphi_{154} :\equiv \forall b\,\forall \alpha\,(\mathrm{Seq}(b) \to \forall n \in \mathbb{N}\,(f \circ a \otimes g \circ a(n) = b_{n}) \to \mathrm{Conv}(f \circ a \otimes g \circ a,\alpha) \to \mathrm{Conv}(b,\alpha))$
\item $\varphi_{155} :\equiv (f \cdot g) \circ a(n) = f \circ a \otimes g \circ a(n) \to f \circ a \otimes g \circ a(n) = (f \cdot g) \circ a(n)$
\item $\varphi_{156} :\equiv \varphi_{6} \to \mathrm{ContIn}(f,A,c) \to \mathrm{ContIn}(g,A,c) \to \varphi_{30}$
\item $\varphi_{157} :\equiv \varphi_{6} \to \mathrm{ContIn}(f,A,c) \to \mathrm{ContIn}(g,A,c) \to \varphi_{31}$
\item $\varphi_{158} :\equiv \mathrm{ContIn}(g,A,c) \to \mathrm{Seq}(a) \to \forall n \in \mathbb{N}\,(a_{n} \in A) \to \varphi_{70}$
\item $\varphi_{159} :\equiv \mathrm{ContIn}(g,A,c) \to \mathrm{Seq}(a) \to \forall n \in \mathbb{N}\,(a_{n} \in A) \to \varphi_{71}$
\item $\varphi_{160} :\equiv \mathrm{ContIn}(g,A,c) \to \mathrm{Seq}(a) \to \forall n \in \mathbb{N}\,(a_{n} \in A) \to \varphi_{72}$
\item $\varphi_{161} :\equiv \forall f\,\forall A\,\forall c\,(\mathrm{RFn}(f) \to A \subseteq \mathrm{dom}(f) \to c \in A \to \forall a\,(\mathrm{Seq}(a) \to \varphi_{40}) \to \mathrm{ContIn}(f,A,c))$
\item $\varphi_{162} :\equiv \mathrm{ContIn}(g,A,c) \to \mathrm{Seq}(a) \to \forall n \in \mathbb{N}\,(a_{n} \in A) \to \varphi_{80}$
\item $\varphi_{163} :\equiv \mathrm{ContIn}(g,A,c) \to \mathrm{Seq}(a) \to \forall n \in \mathbb{N}\,(a_{n} \in A) \to \varphi_{81}$
\item $\varphi_{164} :\equiv f(a_{n}) \cdot g(a_{n}) = f \circ a \otimes g \circ a(n) \to (f \cdot g) \circ a(n) = f \circ a \otimes g \circ a(n)$
\item $\varphi_{165} :\equiv \mathrm{ContIn}(g,A,c) \to \mathrm{Seq}(a) \to \forall n \in \mathbb{N}\,(a_{n} \in A) \to \mathrm{Conv}(a,c) \to \varphi_{61}$
\item $\varphi_{166} :\equiv \mathrm{ContIn}(g,A,c) \to \mathrm{Seq}(a) \to \forall n \in \mathbb{N}\,(a_{n} \in A) \to \varphi_{86}$
\item $\varphi_{167} :\equiv \mathrm{ContIn}(g,A,c) \to \mathrm{Seq}(a) \to \forall n \in \mathbb{N}\,(a_{n} \in A) \to \varphi_{87}$
\item $\varphi_{168} :\equiv f(a_{n}) \cdot g \circ a(n) = f(a_{n}) \cdot g(a_{n}) \to f \circ a \otimes g \circ a(n) = f(a_{n}) \cdot g(a_{n})$
\item $\varphi_{169} :\equiv \mathrm{ContIn}(g,A,c) \to \mathrm{Seq}(a) \to \forall n \in \mathbb{N}\,(a_{n} \in A) \to \varphi_{96}$
\item $\varphi_{170} :\equiv f \circ a \otimes g \circ a(n) = f(a_{n}) \cdot g(a_{n}) \to f(a_{n}) \cdot g(a_{n}) = f \circ a \otimes g \circ a(n)$
\item $\varphi_{171} :\equiv f(a_{n}) \cdot g \circ a(n) = f(a_{n}) \cdot g \circ a(n) \to f(a_{n}) \cdot g \circ a(n) = f(a_{n}) \cdot g(a_{n})$
\item $\varphi_{172} :\equiv \mathrm{ContIn}(g,A,c) \to \mathrm{Seq}(a) \to \forall n \in \mathbb{N}\,(a_{n} \in A) \to \varphi_{100}$
\item $\varphi_{173} :\equiv \mathrm{ContIn}(g,A,c) \to \mathrm{Seq}(a) \to \forall n \in \mathbb{N}\,(a_{n} \in A) \to \varphi_{103}$
\item $\varphi_{174} :\equiv \mathrm{ContIn}(g,A,c) \to \mathrm{Seq}(a) \to \forall n \in \mathbb{N}\,(a_{n} \in A) \to \varphi_{104}$
\item $\varphi_{175} :\equiv \forall g\,(\mathrm{RFn}(f) \to \mathrm{RFn}(g) \to \mathrm{RFn}((f \cdot g)) \land \varphi_{16} \land \varphi_{27})$
\item $\varphi_{176} :\equiv (\forall n \in \mathbb{N}\,(a_{n} \in A) \to n \in \mathbb{N} \to a_{n} \in A) \to \mathrm{ContIn}(f,A,c) \to \varphi_{73}$
\item $\varphi_{177} :\equiv \mathrm{ContIn}(g,A,c) \to \mathrm{Seq}(a) \to \forall n \in \mathbb{N}\,(a_{n} \in A) \to \varphi_{107}$
\item $\varphi_{178} :\equiv \mathrm{ContIn}(g,A,c) \to \mathrm{Seq}(a) \to \forall n \in \mathbb{N}\,(a_{n} \in A) \to \mathrm{Conv}(a,c) \to \varphi_{90}$
\item $\varphi_{179} :\equiv f \circ a(n) \cdot g \circ a(n) = f(a_{n}) \cdot g \circ a(n) \to f \circ a \otimes g \circ a(n) = f(a_{n}) \cdot g \circ a(n)$
\item $\varphi_{180} :\equiv \mathrm{Seq}(a) \to \forall n \in \mathbb{N}\,(a_{n} \in A) \to \mathrm{Conv}(a,c) \to n \in \mathbb{N} \to \varphi_{101}$
\item $\varphi_{181} :\equiv \mathrm{ContIn}(g,A,c) \to \mathrm{Seq}(a) \to \forall n \in \mathbb{N}\,(a_{n} \in A) \to \mathrm{Conv}(a,c) \to \varphi_{93}$
\item $\varphi_{182} :\equiv \mathrm{ContIn}(g,A,c) \to \mathrm{Seq}(a) \to \forall n \in \mathbb{N}\,(a_{n} \in A) \to \varphi_{116}$
\item $\varphi_{183} :\equiv \mathrm{RFn}((f \cdot g)) \to \mathrm{Seq}(a) \to \forall n \in \mathbb{N}\,(a_{n} \in \mathrm{dom}((f \cdot g))) \to \varphi_{93}$
\item $\varphi_{184} :\equiv \forall w\,(g \circ a(n) = w \to f(a_{n}) \cdot g \circ a(n) = f(a_{n}) \cdot g \circ a(n) \to f(a_{n}) \cdot g \circ a(n) = f(a_{n}) \cdot w)$
\item $\varphi_{185} :\equiv \mathrm{ContIn}(f,A,c) \to \mathrm{ContIn}(g,A,c) \to \mathrm{RFn}((f \cdot g)) \land \varphi_{16} \land \varphi_{27}$
\item $\varphi_{186} :\equiv \forall n\,(\varphi_{69} \to n \in \mathbb{N} \to f \circ a \otimes g \circ a(n) = (f \cdot g) \circ a(n))$
\item $\varphi_{187} :\equiv f \circ a(n) \cdot g \circ a(n) = f \circ a(n) \cdot g \circ a(n) \to f \circ a(n) \cdot g \circ a(n) = f(a_{n}) \cdot g \circ a(n)$
\item $\varphi_{188} :\equiv \forall z\,(f \circ a \otimes g \circ a(n) = f(a_{n}) \cdot g \circ a(n) \to f(a_{n}) \cdot g \circ a(n) = z \to f \circ a \otimes g \circ a(n) = z)$
\item $\varphi_{189} :\equiv \forall n\,(\varphi_{69}) \to \forall n \in \mathbb{N}\,(f \circ a \otimes g \circ a(n) = (f \cdot g) \circ a(n))$
\item $\varphi_{190} :\equiv \mathrm{ContIn}(g,A,c) \leftrightarrow \mathrm{RFn}(g) \land A \subseteq \mathrm{dom}(g) \land c \in A \land \forall \varepsilon \in \mathbb{R}\,(0 < \varepsilon \to \exists \delta\,(\delta \in \mathbb{R} \land 0 < \delta \land \forall x \in A\,(\lvert x - c \rvert < \delta \to \lvert g(x) - g(c) \rvert < \varepsilon)))$
\item $\varphi_{191} :\equiv \mathrm{ContIn}(f,A,c) \leftrightarrow \mathrm{RFn}(f) \land A \subseteq \mathrm{dom}(f) \land c \in A \land \forall \varepsilon \in \mathbb{R}\,(0 < \varepsilon \to \exists \delta\,(\delta \in \mathbb{R} \land 0 < \delta \land \forall x \in A\,(\lvert x - c \rvert < \delta \to \lvert f(x) - f(c) \rvert < \varepsilon)))$
\item $\varphi_{192} :\equiv \mathrm{ContIn}(g,A,c) \to \mathrm{Seq}(a) \to \forall n \in \mathbb{N}\,(a_{n} \in A) \to \mathrm{Conv}(a,c) \to \varphi_{115}$
\item $\varphi_{193} :\equiv \forall c\,(\mathrm{RFn}((f \cdot g)) \to A \subseteq \mathrm{dom}((f \cdot g)) \to c \in A \to \varphi_{137})$
\item $\varphi_{194} :\equiv \forall z\,(f \circ a \otimes g \circ a(n) = f \circ a(n) \cdot g \circ a(n) \to f \circ a(n) \cdot g \circ a(n) = z \to f \circ a \otimes g \circ a(n) = z)$
\item $\varphi_{195} :\equiv \mathrm{Seq}(f \circ a) \to \mathrm{Seq}(g \circ a) \to \mathrm{Seq}(f \circ a \otimes g \circ a) \land \varphi_{101}$
\item $\varphi_{196} :\equiv \forall w\,(f \circ a(n) = w \to f \circ a(n) \cdot g \circ a(n) = f \circ a(n) \cdot g \circ a(n) \to f \circ a(n) \cdot g \circ a(n) = w \cdot g \circ a(n))$
\item $\varphi_{197} :\equiv \forall n \in \mathbb{N}\,(a_{n} \in A) \to \mathrm{Conv}(a,c) \to \mathrm{Seq}(f \circ a \otimes g \circ a) \land \varphi_{101}$
\item $\varphi_{198} :\equiv (\mathrm{ContIn}(f,A,c) \to \varphi_{127}) \to \mathrm{ContIn}(f,A,c) \to \varphi_{3}$
\item $\varphi_{199} :\equiv \forall \alpha\,(\mathrm{Seq}((f \cdot g) \circ a) \to \forall n \in \mathbb{N}\,(f \circ a \otimes g \circ a(n) = (f \cdot g) \circ a(n)) \to \mathrm{Conv}(f \circ a \otimes g \circ a,\alpha) \to \mathrm{Conv}((f \cdot g) \circ a,\alpha))$
\item $\varphi_{200} :\equiv \varphi_{101} \to n \in \mathbb{N} \to f \circ a \otimes g \circ a(n) = f \circ a(n) \cdot g \circ a(n)$
\item $\varphi_{201} :\equiv \forall a\,(\mathrm{ContIn}(g,A,c) \to \mathrm{Seq}(a) \to \varphi_{39}) \to \mathrm{ContIn}(g,A,c) \to \mathrm{Seq}(a) \to \varphi_{39}$
\item $\varphi_{202} :\equiv \forall a\,(\mathrm{ContIn}(f,A,c) \to \mathrm{Seq}(a) \to \varphi_{40}) \to \mathrm{ContIn}(f,A,c) \to \mathrm{Seq}(a) \to \varphi_{40}$
\item $\varphi_{203} :\equiv (\mathrm{ContIn}(f,A,c) \land \mathrm{ContIn}(g,A,c) \to \varphi_{1}) \to \mathrm{ContIn}(f,A,c) \to \varphi_{127}$
\item $\varphi_{204} :\equiv \forall c\,\forall a\,(\mathrm{ContIn}(g,A,c) \to \mathrm{Seq}(a) \to \varphi_{39}) \to \forall a\,(\mathrm{ContIn}(g,A,c) \to \mathrm{Seq}(a) \to \varphi_{39})$
\item $\varphi_{205} :\equiv \forall c\,\forall a\,(\mathrm{ContIn}(f,A,c) \to \mathrm{Seq}(a) \to \varphi_{40}) \to \forall a\,(\mathrm{ContIn}(f,A,c) \to \mathrm{Seq}(a) \to \varphi_{40})$
\item $\varphi_{206} :\equiv \forall y\,((f \cdot g) \circ a(n) = y \to y = (f \cdot g) \circ a(n)) \to \varphi_{155}$
\item $\varphi_{207} :\equiv \forall A\,\forall c\,\forall a\,(\mathrm{ContIn}(g,A,c) \to \mathrm{Seq}(a) \to \varphi_{39}) \to \forall c\,\forall a\,(\mathrm{ContIn}(g,A,c) \to \mathrm{Seq}(a) \to \varphi_{39})$
\item $\varphi_{208} :\equiv \forall A\,\forall c\,\forall a\,(\mathrm{ContIn}(f,A,c) \to \mathrm{Seq}(a) \to \varphi_{40}) \to \forall c\,\forall a\,(\mathrm{ContIn}(f,A,c) \to \mathrm{Seq}(a) \to \varphi_{40})$
\item $\varphi_{209} :\equiv \varphi_{117} \to \forall A\,\forall c\,\forall a\,(\mathrm{ContIn}(g,A,c) \to \mathrm{Seq}(a) \to \varphi_{39})$
\item $\varphi_{210} :\equiv \varphi_{117} \to \forall A\,\forall c\,\forall a\,(\mathrm{ContIn}(f,A,c) \to \mathrm{Seq}(a) \to \varphi_{40})$
\item $\varphi_{211} :\equiv \varphi_{122} \to \mathrm{ContIn}(f,A,c) \to \mathrm{ContIn}(g,A,c) \to \varphi_{37}$
\item $\varphi_{212} :\equiv \varphi_{123} \to \mathrm{ContIn}(f,A,c) \to \mathrm{ContIn}(g,A,c) \to \varphi_{38}$
\item $\varphi_{213} :\equiv (\varphi_{27} \to \varphi_{24}) \to \mathrm{ContIn}(f,A,c) \to \mathrm{ContIn}(g,A,c) \to \varphi_{24}$
\item $\varphi_{214} :\equiv \mathrm{ContIn}(f,A,c) \to \mathrm{ContIn}(g,A,c) \to \mathrm{Seq}(a) \to \varphi_{197}$
\item $\varphi_{215} :\equiv \mathrm{Seq}((f \cdot g) \circ a) \to \forall n \in \mathbb{N}\,(f \circ a \otimes g \circ a(n) = (f \cdot g) \circ a(n)) \to \varphi_{120}$
\item $\varphi_{216} :\equiv (\mathrm{ContIn}(f,A,c) \to \mathrm{ContIn}(g,A,c) \to \varphi_{43}) \to \varphi_{138}$
\item $\varphi_{217} :\equiv (\varphi_{16} \to \varphi_{43}) \to \mathrm{ContIn}(f,A,c) \to \mathrm{ContIn}(g,A,c) \to \varphi_{43}$
\item $\varphi_{218} :\equiv \forall n\,(\varphi_{69} \to n \in \mathbb{N} \to a_{n} \in \mathrm{dom}(g)) \to \forall n\,(\varphi_{69}) \to \forall n \in \mathbb{N}\,(a_{n} \in \mathrm{dom}(g))$
\item $\varphi_{219} :\equiv \forall n\,(\varphi_{69} \to n \in \mathbb{N} \to a_{n} \in \mathrm{dom}(f)) \to \forall n\,(\varphi_{69}) \to \forall n \in \mathbb{N}\,(a_{n} \in \mathrm{dom}(f))$
\item $\varphi_{220} :\equiv \forall y\,(f \circ a \otimes g \circ a(n) = y \to y = f \circ a \otimes g \circ a(n)) \to \varphi_{170}$
\item $\varphi_{221} :\equiv (\mathrm{ContIn}(f,A,c) \to \mathrm{ContIn}(f,A,c)) \to \varphi_{191} \to \varphi_{5}$
\item $\varphi_{222} :\equiv (\mathrm{ContIn}(f,A,c) \to \mathrm{ContIn}(f,A,c)) \to \varphi_{191} \to \varphi_{6}$
\item $\varphi_{223} :\equiv (\mathrm{ContIn}(f,A,c) \to \mathrm{ContIn}(f,A,c)) \to \varphi_{191} \to \varphi_{9}$
\item $\varphi_{224} :\equiv \forall a\,(\mathrm{RFn}(g) \to \mathrm{Seq}(a) \to \varphi_{113}) \to \mathrm{RFn}(g) \to \mathrm{Seq}(a) \to \varphi_{113}$
\item $\varphi_{225} :\equiv \forall a\,(\mathrm{RFn}(f) \to \mathrm{Seq}(a) \to \varphi_{114}) \to \mathrm{RFn}(f) \to \mathrm{Seq}(a) \to \varphi_{114}$
\item $\varphi_{226} :\equiv \varphi_{185} \to \mathrm{ContIn}(f,A,c) \to \mathrm{ContIn}(g,A,c) \to \varphi_{16}$
\item $\varphi_{227} :\equiv \forall n\,(\varphi_{69} \to n \in \mathbb{N} \to a_{n} \in \mathrm{dom}((f \cdot g))) \to \forall n\,(\varphi_{69}) \to \forall n \in \mathbb{N}\,(a_{n} \in \mathrm{dom}((f \cdot g)))$
\item $\varphi_{228} :\equiv \forall f\,\forall a\,(\mathrm{RFn}(f) \to \mathrm{Seq}(a) \to \varphi_{114}) \to \forall a\,(\mathrm{RFn}(g) \to \mathrm{Seq}(a) \to \varphi_{113})$
\item $\varphi_{229} :\equiv \forall f\,\forall a\,(\mathrm{RFn}(f) \to \mathrm{Seq}(a) \to \varphi_{114}) \to \forall a\,(\mathrm{RFn}(f) \to \mathrm{Seq}(a) \to \varphi_{114})$
\item $\varphi_{230} :\equiv (f(c) \cdot g(c) = (f \cdot g)(c) \to \varphi_{89}) \to \mathrm{ContIn}(f,A,c) \to \varphi_{125}$
\item $\varphi_{231} :\equiv \varphi_{185} \to \mathrm{ContIn}(f,A,c) \to \mathrm{ContIn}(g,A,c) \to \varphi_{27}$
\item $\varphi_{232} :\equiv \forall a\,(\varphi_{140}) \to \varphi_{1} \to \forall a\,(\mathrm{Seq}(a) \to \forall n \in \mathbb{N}\,(a_{n} \in A) \to \varphi_{19})$
\item $\varphi_{233} :\equiv (\varphi_{27} \to \varphi_{82}) \to \mathrm{ContIn}(f,A,c) \to \mathrm{ContIn}(g,A,c) \to \varphi_{82}$
\item $\varphi_{234} :\equiv (\mathrm{ContIn}(f,A,c) \to \mathrm{ContIn}(g,A,c) \to \varphi_{16}) \to \varphi_{111} \to \mathrm{ContIn}(f,A,c) \to \varphi_{33}$
\item $\varphi_{235} :\equiv (\mathrm{ContIn}(f,A,c) \to \mathrm{ContIn}(g,A,c) \to \varphi_{16}) \to \varphi_{112} \to \mathrm{ContIn}(f,A,c) \to \varphi_{34}$
\item $\varphi_{236} :\equiv (\mathrm{ContIn}(f,A,c) \to \varphi_{143}) \to \mathrm{ContIn}(f,A,c) \to \varphi_{130}$
\item $\varphi_{237} :\equiv (\mathrm{ContIn}(f,A,c) \to \mathrm{ContIn}(g,A,c) \to \varphi_{30}) \to (\mathrm{ContIn}(f,A,c) \to \varphi_{33}) \to \varphi_{108}$
\item $\varphi_{238} :\equiv (\mathrm{ContIn}(f,A,c) \to \varphi_{158}) \to \mathrm{ContIn}(f,A,c) \to \varphi_{144}$
\item $\varphi_{239} :\equiv (\mathrm{ContIn}(f,A,c) \to \varphi_{159}) \to \mathrm{ContIn}(f,A,c) \to \varphi_{145}$
\item $\varphi_{240} :\equiv (\mathrm{ContIn}(f,A,c) \to \mathrm{ContIn}(g,A,c) \to \varphi_{27}) \to \varphi_{213}$
\item $\varphi_{241} :\equiv (\mathrm{ContIn}(f,A,c) \to \mathrm{ContIn}(g,A,c) \to \varphi_{16}) \to \varphi_{217}$
\item $\varphi_{242} :\equiv (\mathrm{ContIn}(f,A,c) \to \mathrm{ContIn}(g,A,c) \to \varphi_{24}) \to \varphi_{108} \to \varphi_{134}$
\item $\varphi_{243} :\equiv (\mathrm{ContIn}(f,A,c) \to \mathrm{ContIn}(g,A,c) \to \varphi_{37}) \to \varphi_{214}$
\item $\varphi_{244} :\equiv (A \subseteq \mathrm{dom}(g) \to a_{n} \in A \to a_{n} \in \mathrm{dom}(g)) \to \varphi_{8} \to \varphi_{121} \to \varphi_{128}$
\item $\varphi_{245} :\equiv (A \subseteq \mathrm{dom}(f) \to a_{n} \in A \to a_{n} \in \mathrm{dom}(f)) \to \varphi_{9} \to \varphi_{121} \to \varphi_{129}$
\item $\varphi_{246} :\equiv (\mathrm{ContIn}(f,A,c) \to \varphi_{167}) \to \varphi_{69} \to n \in \mathbb{N} \to f \circ a \otimes g \circ a(n) = (f \cdot g) \circ a(n)$
\item $\varphi_{247} :\equiv ((f \cdot g) \circ a(n) = (f \cdot g)(a_{n}) \to \varphi_{136}) \to \mathrm{ContIn}(f,A,c) \to \varphi_{162}$
\item $\varphi_{248} :\equiv \forall u\,(\varphi_{26}) \to (\forall u\,(\varphi_{26}) \to \varphi_{2} \to \forall u \in A\,(u \in \mathrm{dom}((f \cdot g)))) \to \varphi_{149}$
\item $\varphi_{249} :\equiv (\mathrm{ContIn}(f,A,c) \to \varphi_{181}) \to \mathrm{ContIn}(f,A,c) \to \varphi_{165}$
\item $\varphi_{250} :\equiv \varphi_{193} \to \mathrm{RFn}((f \cdot g)) \to A \subseteq \mathrm{dom}((f \cdot g)) \to c \in A \to \varphi_{137}$
\item $\varphi_{251} :\equiv (g \circ a(n) = g(a_{n}) \to \varphi_{171}) \to \mathrm{ContIn}(f,A,c) \to \varphi_{172}$
\item $\varphi_{252} :\equiv (\mathrm{ContIn}(f,A,c) \to \mathrm{ContIn}(g,A,c) \to \varphi_{27}) \to \varphi_{233}$
\item $\varphi_{253} :\equiv ((f \cdot g) \circ a(n) = f(a_{n}) \cdot g(a_{n}) \to \varphi_{164}) \to \mathrm{ContIn}(f,A,c) \to \varphi_{166}$
\item $\varphi_{254} :\equiv (f \circ a(n) = f(a_{n}) \to \varphi_{187}) \to \mathrm{ContIn}(f,A,c) \to \varphi_{178}$
\item $\varphi_{255} :\equiv (\mathrm{ContIn}(f,A,c) \to \mathrm{ContIn}(g,A,c) \to \varphi_{31}) \to \varphi_{237}$
\item $\varphi_{256} :\equiv (\mathrm{ContIn}(f,A,c) \to \mathrm{ContIn}(f,A,c)) \to \varphi_{11} \to (\mathrm{ContIn}(f,A,c) \to \varphi_{78}) \to \varphi_{91} \to \varphi_{123}$
\item $\varphi_{257} :\equiv (\mathrm{ContIn}(f,A,c) \to \varphi_{144}) \to \varphi_{105} \to \mathrm{ContIn}(f,A,c) \to \varphi_{141}$
\item $\varphi_{258} :\equiv (\mathrm{ContIn}(f,A,c) \to \varphi_{145}) \to \varphi_{106} \to \mathrm{ContIn}(f,A,c) \to \varphi_{142}$
\item $\varphi_{259} :\equiv (\mathrm{RFn}(f) \to \mathrm{RFn}(g) \to \mathrm{RFn}((f \cdot g)) \land \varphi_{16} \land \varphi_{27}) \to \varphi_{5} \to \varphi_{4} \to \varphi_{185}$
\item $\varphi_{260} :\equiv (f \circ a \otimes g \circ a(n) = f(a_{n}) \cdot g \circ a(n) \to \varphi_{168}) \to \mathrm{ContIn}(f,A,c) \to \varphi_{173}$
\item $\varphi_{261} :\equiv (\mathrm{ContIn}(f,A,c) \to \varphi_{163}) \to \varphi_{110} \to \mathrm{ContIn}(f,A,c) \to \varphi_{150}$
\item $\varphi_{262} :\equiv \varphi_{214} \to \mathrm{ContIn}(f,A,c) \to \mathrm{ContIn}(g,A,c) \to \varphi_{180}$
\item $\varphi_{263} :\equiv (f \circ a \otimes g \circ a(n) = f \circ a(n) \cdot g \circ a(n) \to \varphi_{179}) \to \mathrm{ContIn}(f,A,c) \to \varphi_{177}$
\item $\varphi_{264} :\equiv (\mathrm{ContIn}(f,A,c) \to \mathrm{ContIn}(g,A,c) \to \varphi_{38}) \to \varphi_{243}$
\item $\varphi_{265} :\equiv (\mathrm{ContIn}(f,A,c) \to \varphi_{165}) \to \varphi_{148} \to \mathrm{ContIn}(f,A,c) \to \varphi_{163}$
\item $\varphi_{266} :\equiv (\mathrm{ContIn}(f,A,c) \to \varphi_{166}) \to \varphi_{155} \to \mathrm{ContIn}(f,A,c) \to \varphi_{167}$
\item $\varphi_{267} :\equiv (\mathrm{ContIn}(f,A,c) \to \varphi_{192}) \to \varphi_{110} \to \mathrm{ContIn}(f,A,c) \to \varphi_{182}$
\item $\varphi_{268} :\equiv (\mathrm{RFn}(g) \to \mathrm{Seq}(a) \to \varphi_{113}) \to \varphi_{4} \to \varphi_{11} \to \varphi_{131} \to \mathrm{ContIn}(f,A,c) \to \varphi_{158}$
\item $\varphi_{269} :\equiv (\mathrm{RFn}(f) \to \mathrm{Seq}(a) \to \varphi_{114}) \to \varphi_{5} \to \varphi_{11} \to \varphi_{132} \to \mathrm{ContIn}(f,A,c) \to \varphi_{159}$
\item $\varphi_{270} :\equiv (\mathrm{Conv}(f \circ a,f(c)) \to \varphi_{83}) \to \varphi_{123} \to \varphi_{122} \to \mathrm{ContIn}(f,A,c) \to \varphi_{143}$
\item $\varphi_{271} :\equiv (\mathrm{ContIn}(f,A,c) \to \varphi_{173}) \to \varphi_{170} \to \mathrm{ContIn}(f,A,c) \to \varphi_{174}$
\item $\varphi_{272} :\equiv (\mathrm{ContIn}(g,A,c) \to \mathrm{Seq}(a) \to \varphi_{39}) \to \varphi_{7} \to \varphi_{11} \to (\mathrm{ContIn}(f,A,c) \to \varphi_{78}) \to \varphi_{91} \to \varphi_{122}$
\item $\varphi_{273} :\equiv (\mathrm{RFn}((f \cdot g)) \to A \subseteq \mathrm{dom}((f \cdot g)) \to c \in A \to \varphi_{137}) \to \varphi_{10} \to \varphi_{14} \to \varphi_{6} \to \varphi_{198}$
\item $\varphi_{274} :\equiv (\mathrm{ContIn}(f,A,c) \to \mathrm{ContIn}(g,A,c) \to \varphi_{82}) \to \varphi_{138} \to \varphi_{110} \to \mathrm{ContIn}(f,A,c) \to \varphi_{160}$
\item $\varphi_{275} :\equiv (\mathrm{ContIn}(f,A,c) \to \mathrm{ContIn}(g,A,c) \to \varphi_{180}) \to \varphi_{200} \to \mathrm{ContIn}(f,A,c) \to \varphi_{192}$
\item $\varphi_{276} :\equiv \varphi_{147} \to f \circ a(n) \cdot g \circ a(n) = f \circ a(n) \cdot g \circ a(n) \to \varphi_{254}$
\item $\varphi_{277} :\equiv (\mathrm{ContIn}(f,A,c) \to \varphi_{150}) \to (\mathrm{ContIn}(f,A,c) \to \varphi_{160}) \to \varphi_{247}$
\item $\varphi_{278} :\equiv \forall n\,(\varphi_{69} \to n \in \mathbb{N} \to a_{n} \in \mathrm{dom}(g)) \to \varphi_{218} \to (\varphi_{69} \to \forall n\,(\varphi_{69})) \to \varphi_{131}$
\item $\varphi_{279} :\equiv \forall n\,(\varphi_{69} \to n \in \mathbb{N} \to a_{n} \in \mathrm{dom}(f)) \to \varphi_{219} \to (\varphi_{69} \to \forall n\,(\varphi_{69})) \to \varphi_{132}$
\item $\varphi_{280} :\equiv (\varphi_{186} \to \varphi_{189}) \to (\varphi_{69} \to \forall n\,(\varphi_{69})) \to \mathrm{ContIn}(f,A,c) \to \varphi_{169}$
\item $\varphi_{281} :\equiv \forall n\,(\varphi_{69} \to n \in \mathbb{N} \to a_{n} \in \mathrm{dom}((f \cdot g))) \to \varphi_{227} \to (\varphi_{69} \to \forall n\,(\varphi_{69})) \to \varphi_{139}$
\item $\varphi_{282} :\equiv (\mathrm{ContIn}(f,A,c) \to \varphi_{162}) \to (\mathrm{ContIn}(f,A,c) \to \varphi_{174}) \to \varphi_{253}$
\item $\varphi_{283} :\equiv (\mathrm{ContIn}(f,A,c) \to \varphi_{177}) \to (\mathrm{ContIn}(f,A,c) \to \varphi_{172}) \to \varphi_{260}$
\item $\varphi_{284} :\equiv (\mathrm{ContIn}(f,A,c) \to \varphi_{182}) \to (\mathrm{ContIn}(f,A,c) \to \varphi_{178}) \to \varphi_{263}$
\end{abbrevlist}
\end{fullproof}

% 自動生成 (prover/build.js): ContInScal  ノード数 416
\begin{fullproof}{定理2.3(定数倍)}{$\forall l\,\forall f\,\forall A\,\forall c\,(l \in \mathbb{R} \to \mathrm{ContIn}(f,A,c) \to \mathrm{ContIn}((l f),A,c))$}{416}\phantomsection\label{tree:ContInScal}
\noindent$\Pi_{1}$:
\begin{ptree}
\AXC{$\varphi_{22}$\lbl{A6}}
\UIC{$\forall w\,(\varphi_{22})$}
\UIC{$\forall z\,\forall w\,(\varphi_{22})$}
\AXC{$\forall z\,\forall w\,(\varphi_{22}) \to \varphi_{82}$\lbl{A2}}
\BIC{$\varphi_{82}$}
\AXC{$\varphi_{82} \to l \cdot f(c) = (l f)(c) \to \varphi_{73}$\lbl{A2}}
\BIC{$l \cdot f(c) = (l f)(c) \to \varphi_{73}$}
\end{ptree}
\noindent$\Pi_{2}$:
\begin{ptree}
\AXC{$\varphi_{1}$\lbl{A1}}
\AXC{$\varphi_{95}$\lbl{[G1]}}
\AXC{$\varphi_{153}$\lbl{A2}}
\BIC{$\forall A\,\forall c\,\forall a\,(\mathrm{ContIn}(f,A,c) \to \mathrm{Seq}(a) \to \varphi_{29})$}
\AXC{$\varphi_{150}$\lbl{A2}}
\BIC{$\forall c\,\forall a\,(\mathrm{ContIn}(f,A,c) \to \mathrm{Seq}(a) \to \varphi_{29})$}
\AXC{$\varphi_{148}$\lbl{A2}}
\BIC{$\forall a\,(\mathrm{ContIn}(f,A,c) \to \mathrm{Seq}(a) \to \varphi_{29})$}
\AXC{$\varphi_{146}$\lbl{A2}}
\BIC{$\mathrm{ContIn}(f,A,c) \to \mathrm{Seq}(a) \to \varphi_{29}$}
\AXC{$\varphi_{201}$\lbl{A1}}
\BIC{$\varphi_{1} \to \varphi_{2} \to (l \in \mathbb{R} \to \varphi_{69}) \to \varphi_{147}$}
\BIC{$\varphi_{2} \to (l \in \mathbb{R} \to \varphi_{69}) \to \varphi_{147}$}
\end{ptree}
\noindent$\Pi_{3}$:
\begin{ptree}
\AXC{$l \in \mathbb{R} \to \mathrm{ContIn}(f,A,c) \to \mathrm{Seq}(a) \to \varphi_{8}$\lbl{A1}}
\AXC{$l \in \mathbb{R} \to \varphi_{69}$\lbl{A1}}
\AXC{$\varphi_{2}$\lbl{A1}}
\AXC{$\stackrel{\Pi_{2}}{\vdots}$}\noLine
\UIC{$\varphi_{2} \to (l \in \mathbb{R} \to \varphi_{69}) \to \varphi_{147}$}
\BIC{$(l \in \mathbb{R} \to \varphi_{69}) \to \varphi_{147}$}
\BIC{$\varphi_{147}$}
\BIC{$l \in \mathbb{R} \to \mathrm{ContIn}(f,A,c) \to \mathrm{Seq}(a) \to \varphi_{29}$}
\end{ptree}
\noindent$\Pi_{4}$:
\begin{ptree}
\AXC{$\varphi_{13}$\lbl{定理2.1(2)}}
\AXC{$\varphi_{13} \to \varphi_{10}$\lbl{A2}}
\BIC{$\varphi_{10}$}
\AXC{$\varphi_{10} \to \varphi_{42}$\lbl{A2}}
\BIC{$\varphi_{42}$}
\AXC{$\varphi_{42} \to l \in \mathbb{R} \to \varphi_{25}$\lbl{A2}}
\BIC{$l \in \mathbb{R} \to \varphi_{25}$}
\AXC{$(l \in \mathbb{R} \to \varphi_{25}) \to (l \in \mathbb{R} \to l \in \mathbb{R}) \to \varphi_{157}$\lbl{A1}}
\BIC{$(l \in \mathbb{R} \to l \in \mathbb{R}) \to \varphi_{157}$}
\end{ptree}
\noindent$\Pi_{5}$:
\begin{ptree}
\AXC{$\stackrel{\Pi_{3}}{\vdots}$}\noLine
\UIC{$l \in \mathbb{R} \to \mathrm{ContIn}(f,A,c) \to \mathrm{Seq}(a) \to \varphi_{29}$}
\AXC{$l \in \mathbb{R} \to l \in \mathbb{R}$\lbl{A1}}
\AXC{$\stackrel{\Pi_{4}}{\vdots}$}\noLine
\UIC{$(l \in \mathbb{R} \to l \in \mathbb{R}) \to \varphi_{157}$}
\BIC{$\varphi_{157}$}
\BIC{$l \in \mathbb{R} \to \mathrm{ContIn}(f,A,c) \to \mathrm{Seq}(a) \to \varphi_{64}$}
\end{ptree}
\noindent$\Pi_{6}$:
\begin{ptree}
\AXC{$x = x$\lbl{A5}}
\AXC{$x = y \to x = x \to y = x$\lbl{A6}}
\AXC{$(x = y \to x = x \to y = x) \to x = x \to x = y \to y = x$\lbl{A1}}
\BIC{$x = x \to x = y \to y = x$}
\BIC{$x = y \to y = x$}
\UIC{$\forall y\,(x = y \to y = x)$}
\UIC{$\forall x\,\forall y\,(x = y \to y = x)$}
\end{ptree}
\noindent$\Pi_{7}$:
\begin{ptree}
\AXC{$\stackrel{\Pi_{6}}{\vdots}$}\noLine
\UIC{$\forall x\,\forall y\,(x = y \to y = x)$}
\AXC{$\varphi_{63}$\lbl{A2}}
\BIC{$\forall y\,((l f) \circ a(n) = y \to y = (l f) \circ a(n))$}
\AXC{$\forall y\,((l f) \circ a(n) = y \to y = (l f) \circ a(n)) \to \varphi_{100}$\lbl{A2}}
\BIC{$\varphi_{100}$}
\end{ptree}
\noindent$\Pi_{8}$:
\begin{ptree}
\AXC{$y = z \to x = y \to x = z$\lbl{A6}}
\AXC{$(y = z \to x = y \to x = z) \to x = y \to y = z \to x = z$\lbl{A1}}
\BIC{$x = y \to y = z \to x = z$}
\UIC{$\forall z\,(x = y \to y = z \to x = z)$}
\UIC{$\forall y\,\forall z\,(x = y \to y = z \to x = z)$}
\UIC{$\forall x\,\forall y\,\forall z\,(x = y \to y = z \to x = z)$}
\AXC{$\forall x\,\forall y\,\forall z\,(x = y \to y = z \to x = z) \to \varphi_{43}$\lbl{A2}}
\BIC{$\varphi_{43}$}
\end{ptree}
\noindent$\Pi_{9}$:
\begin{ptree}
\AXC{$\stackrel{\Pi_{6}}{\vdots}$}\noLine
\UIC{$\forall x\,\forall y\,(x = y \to y = x)$}
\AXC{$\varphi_{57}$\lbl{A2}}
\BIC{$\forall y\,(l \odot f \circ a(n) = y \to y = l \odot f \circ a(n))$}
\AXC{$\forall y\,(l \odot f \circ a(n) = y \to y = l \odot f \circ a(n)) \to \varphi_{91}$\lbl{A2}}
\BIC{$\varphi_{91}$}
\end{ptree}
\noindent$\Pi_{10}$:
\begin{ptree}
\AXC{$y = z \to x = y \to x = z$\lbl{A6}}
\AXC{$(y = z \to x = y \to x = z) \to x = y \to y = z \to x = z$\lbl{A1}}
\BIC{$x = y \to y = z \to x = z$}
\UIC{$\forall z\,(x = y \to y = z \to x = z)$}
\UIC{$\forall y\,\forall z\,(x = y \to y = z \to x = z)$}
\UIC{$\forall x\,\forall y\,\forall z\,(x = y \to y = z \to x = z)$}
\AXC{$\forall x\,\forall y\,\forall z\,(x = y \to y = z \to x = z) \to \varphi_{34}$\lbl{A2}}
\BIC{$\varphi_{34}$}
\end{ptree}
\noindent$\Pi_{11}$:
\begin{ptree}
\AXC{$z = w \to l \cdot f \circ a(n) = l \cdot z \to l \cdot f \circ a(n) = l \cdot w$\lbl{A6}}
\UIC{$\varphi_{36}$}
\UIC{$\forall z\,(\varphi_{36})$}
\AXC{$\forall z\,(\varphi_{36}) \to \varphi_{97}$\lbl{A2}}
\BIC{$\varphi_{97}$}
\AXC{$\varphi_{97} \to f \circ a(n) = f(a_{n}) \to \varphi_{88}$\lbl{A2}}
\BIC{$f \circ a(n) = f(a_{n}) \to \varphi_{88}$}
\end{ptree}
\noindent$\Pi_{12}$:
\begin{ptree}
\AXC{$\forall n \in \mathbb{N}\,(a_{n} \in A) \to n \in \mathbb{N} \to a_{n} \in A$\lbl{A2}}
\AXC{$l \in \mathbb{R} \to \varphi_{69}$\lbl{A1}}
\AXC{$\varphi_{169}$\lbl{A1}}
\BIC{$(\forall n \in \mathbb{N}\,(a_{n} \in A) \to n \in \mathbb{N} \to a_{n} \in A) \to l \in \mathbb{R} \to \varphi_{61}$}
\BIC{$l \in \mathbb{R} \to \varphi_{61}$}
\AXC{$(l \in \mathbb{R} \to \varphi_{61}) \to \varphi_{152}$\lbl{A1}}
\BIC{$\varphi_{152}$}
\end{ptree}
\noindent$\Pi_{13}$:
\begin{ptree}
\AXC{$l \in \mathbb{R} \to \mathrm{ContIn}(f,A,c) \to \mathrm{Seq}(a) \to \varphi_{17}$\lbl{A1}}
\AXC{$\stackrel{\Pi_{12}}{\vdots}$}\noLine
\UIC{$\varphi_{152}$}
\BIC{$l \in \mathbb{R} \to \mathrm{ContIn}(f,A,c) \to \mathrm{Seq}(a) \to \varphi_{26}$}
\end{ptree}
\noindent$\Pi_{14}$:
\begin{ptree}
\AXC{$\forall f\,\forall A\,\forall c\,(\varphi_{139})$\lbl{DefContIn}}
\AXC{$\forall f\,\forall A\,\forall c\,(\varphi_{139}) \to \forall A\,\forall c\,(\varphi_{139})$\lbl{A2}}
\BIC{$\forall A\,\forall c\,(\varphi_{139})$}
\AXC{$\forall A\,\forall c\,(\varphi_{139}) \to \forall c\,(\varphi_{139})$\lbl{A2}}
\BIC{$\forall c\,(\varphi_{139})$}
\AXC{$\forall c\,(\varphi_{139}) \to \varphi_{139}$\lbl{A2}}
\BIC{$\varphi_{139}$}
\end{ptree}
\noindent$\Pi_{15}$:
\begin{ptree}
\AXC{$\stackrel{\Pi_{14}}{\vdots}$}\noLine
\UIC{$\varphi_{139}$}
\AXC{$\varphi_{1}$\lbl{A1}}
\AXC{$\varphi_{171}$\lbl{A1}}
\BIC{$\varphi_{139} \to l \in \mathbb{R} \to \mathrm{ContIn}(f,A,c) \to A \subseteq \mathrm{dom}(f)$}
\BIC{$l \in \mathbb{R} \to \mathrm{ContIn}(f,A,c) \to A \subseteq \mathrm{dom}(f)$}
\end{ptree}
\noindent$\Pi_{16}$:
\begin{ptree}
\AXC{$\forall X\,\forall Y\,\forall x\,(X \subseteq Y \to x \in X \to x \in Y)$\lbl{[B5]}}
\AXC{$\varphi_{37}$\lbl{A2}}
\BIC{$\forall Y\,\forall x\,(A \subseteq Y \to x \in A \to x \in Y)$}
\AXC{$\varphi_{50}$\lbl{A2}}
\BIC{$\forall x\,(A \subseteq \mathrm{dom}(f) \to x \in A \to x \in \mathrm{dom}(f))$}
\AXC{$\varphi_{85}$\lbl{A2}}
\BIC{$A \subseteq \mathrm{dom}(f) \to a_{n} \in A \to a_{n} \in \mathrm{dom}(f)$}
\end{ptree}
\noindent$\Pi_{17}$:
\begin{ptree}
\AXC{$\stackrel{\Pi_{16}}{\vdots}$}\noLine
\UIC{$A \subseteq \mathrm{dom}(f) \to a_{n} \in A \to a_{n} \in \mathrm{dom}(f)$}
\AXC{$(A \subseteq \mathrm{dom}(f) \to a_{n} \in A \to a_{n} \in \mathrm{dom}(f)) \to \varphi_{172}$\lbl{A1}}
\BIC{$\varphi_{172}$}
\end{ptree}
\noindent$\Pi_{18}$:
\begin{ptree}
\AXC{$\stackrel{\Pi_{13}}{\vdots}$}\noLine
\UIC{$l \in \mathbb{R} \to \mathrm{ContIn}(f,A,c) \to \mathrm{Seq}(a) \to \varphi_{26}$}
\AXC{$\stackrel{\Pi_{15}}{\vdots}$}\noLine
\UIC{$l \in \mathbb{R} \to \mathrm{ContIn}(f,A,c) \to A \subseteq \mathrm{dom}(f)$}
\AXC{$\stackrel{\Pi_{17}}{\vdots}$}\noLine
\UIC{$\varphi_{172}$}
\BIC{$\varphi_{155}$}
\BIC{$l \in \mathbb{R} \to \mathrm{ContIn}(f,A,c) \to \mathrm{Seq}(a) \to \varphi_{44}$}
\end{ptree}
\noindent$\Pi_{19}$:
\begin{ptree}
\AXC{$\varphi_{159}$\lbl{A3}}
\AXC{$\stackrel{\Pi_{18}}{\vdots}$}\noLine
\UIC{$l \in \mathbb{R} \to \mathrm{ContIn}(f,A,c) \to \mathrm{Seq}(a) \to \varphi_{44}$}
\AXC{$\varphi_{156}$\lbl{A1}}
\BIC{$\varphi_{54} \to n \in \mathbb{N} \to a_{n} \in \mathrm{dom}(f)$}
\UIC{$\forall n\,(\varphi_{54} \to n \in \mathbb{N} \to a_{n} \in \mathrm{dom}(f))$}
\AXC{$\forall n\,(\varphi_{54} \to n \in \mathbb{N} \to a_{n} \in \mathrm{dom}(f)) \to \varphi_{159} \to \varphi_{175}$\lbl{A1}}
\BIC{$\varphi_{159} \to \varphi_{175}$}
\BIC{$\varphi_{175}$}
\end{ptree}
\noindent$\Pi_{20}$:
\begin{ptree}
\AXC{$\forall f\,\forall A\,\forall c\,(\varphi_{139})$\lbl{DefContIn}}
\AXC{$\forall f\,\forall A\,\forall c\,(\varphi_{139}) \to \forall A\,\forall c\,(\varphi_{139})$\lbl{A2}}
\BIC{$\forall A\,\forall c\,(\varphi_{139})$}
\AXC{$\forall A\,\forall c\,(\varphi_{139}) \to \forall c\,(\varphi_{139})$\lbl{A2}}
\BIC{$\forall c\,(\varphi_{139})$}
\AXC{$\forall c\,(\varphi_{139}) \to \varphi_{139}$\lbl{A2}}
\BIC{$\varphi_{139}$}
\end{ptree}
\noindent$\Pi_{21}$:
\begin{ptree}
\AXC{$\stackrel{\Pi_{20}}{\vdots}$}\noLine
\UIC{$\varphi_{139}$}
\AXC{$\varphi_{1}$\lbl{A1}}
\AXC{$\varphi_{1} \to \varphi_{139} \to l \in \mathbb{R} \to \mathrm{ContIn}(f,A,c) \to \mathrm{RFn}(f)$\lbl{A1}}
\BIC{$\varphi_{139} \to l \in \mathbb{R} \to \mathrm{ContIn}(f,A,c) \to \mathrm{RFn}(f)$}
\BIC{$l \in \mathbb{R} \to \mathrm{ContIn}(f,A,c) \to \mathrm{RFn}(f)$}
\end{ptree}
\noindent$\Pi_{22}$:
\begin{ptree}
\AXC{$\forall f\,\forall a\,(\mathrm{RFn}(f) \to \mathrm{Seq}(a) \to \varphi_{94})$\lbl{[S8]}}
\AXC{$\varphi_{170}$\lbl{A2}}
\BIC{$\forall a\,(\mathrm{RFn}(f) \to \mathrm{Seq}(a) \to \varphi_{94})$}
\AXC{$\varphi_{166}$\lbl{A2}}
\BIC{$\mathrm{RFn}(f) \to \mathrm{Seq}(a) \to \varphi_{94}$}
\AXC{$(\mathrm{RFn}(f) \to \mathrm{Seq}(a) \to \varphi_{94}) \to \varphi_{190}$\lbl{A1}}
\BIC{$\varphi_{190}$}
\end{ptree}
\noindent$\Pi_{23}$:
\begin{ptree}
\AXC{$\varphi_{54} \to \forall n\,(\varphi_{54})$\lbl{A4}}
\AXC{$\stackrel{\Pi_{19}}{\vdots}$}\noLine
\UIC{$\varphi_{175}$}
\BIC{$l \in \mathbb{R} \to \mathrm{ContIn}(f,A,c) \to \mathrm{Seq}(a) \to \varphi_{53}$}
\AXC{$\varphi_{2}$\lbl{A1}}
\AXC{$\stackrel{\Pi_{21}}{\vdots}$}\noLine
\UIC{$l \in \mathbb{R} \to \mathrm{ContIn}(f,A,c) \to \mathrm{RFn}(f)$}
\AXC{$\stackrel{\Pi_{22}}{\vdots}$}\noLine
\UIC{$\varphi_{190}$}
\BIC{$\varphi_{2} \to \varphi_{173}$}
\BIC{$\varphi_{173}$}
\BIC{$l \in \mathbb{R} \to \varphi_{125}$}
\end{ptree}
\noindent$\Pi_{24}$:
\begin{ptree}
\AXC{$\varphi_{92}$\lbl{A2}}
\AXC{$\stackrel{\Pi_{23}}{\vdots}$}\noLine
\UIC{$l \in \mathbb{R} \to \varphi_{125}$}
\AXC{$(l \in \mathbb{R} \to \varphi_{125}) \to l \in \mathbb{R} \to \varphi_{114}$\lbl{A1}}
\BIC{$l \in \mathbb{R} \to \varphi_{114}$}
\AXC{$(l \in \mathbb{R} \to \varphi_{114}) \to \varphi_{92} \to l \in \mathbb{R} \to \varphi_{108}$\lbl{A1}}
\BIC{$\varphi_{92} \to l \in \mathbb{R} \to \varphi_{108}$}
\BIC{$l \in \mathbb{R} \to \varphi_{108}$}
\end{ptree}
\noindent$\Pi_{25}$:
\begin{ptree}
\AXC{$l \in \mathbb{R} \to \mathrm{ContIn}(f,A,c) \to \mathrm{Seq}(a) \to \varphi_{17}$\lbl{A1}}
\AXC{$\stackrel{\Pi_{24}}{\vdots}$}\noLine
\UIC{$l \in \mathbb{R} \to \varphi_{108}$}
\AXC{$(l \in \mathbb{R} \to \varphi_{108}) \to \varphi_{158}$\lbl{A1}}
\BIC{$\varphi_{158}$}
\BIC{$l \in \mathbb{R} \to \mathrm{ContIn}(f,A,c) \to \mathrm{Seq}(a) \to \varphi_{77}$}
\AXC{$\varphi_{199}$\lbl{A1}}
\BIC{$\varphi_{187}$}
\end{ptree}
\noindent$\Pi_{26}$:
\begin{ptree}
\AXC{$\forall a\,\forall \alpha\,(\mathrm{Conv}(a,\alpha) \to \mathrm{Seq}(a))$\lbl{[C1]}}
\AXC{$\varphi_{38}$\lbl{A2}}
\BIC{$\forall \alpha\,(\mathrm{Conv}(f \circ a,\alpha) \to \mathrm{Seq}(f \circ a))$}
\AXC{$\varphi_{83}$\lbl{A2}}
\BIC{$\mathrm{Conv}(f \circ a,f(c)) \to \mathrm{Seq}(f \circ a)$}
\AXC{$(\mathrm{Conv}(f \circ a,f(c)) \to \mathrm{Seq}(f \circ a)) \to \varphi_{151}$\lbl{A1}}
\BIC{$\varphi_{151}$}
\end{ptree}
\noindent$\Pi_{27}$:
\begin{ptree}
\AXC{$\varphi_{76}$\lbl{[S3]}}
\AXC{$\varphi_{76} \to \varphi_{70}$\lbl{A2}}
\BIC{$\varphi_{70}$}
\AXC{$\varphi_{70} \to l \in \mathbb{R} \to \mathrm{Seq}(f \circ a) \to \varphi_{84}$\lbl{A2}}
\BIC{$l \in \mathbb{R} \to \mathrm{Seq}(f \circ a) \to \varphi_{84}$}
\AXC{$(l \in \mathbb{R} \to \mathrm{Seq}(f \circ a) \to \varphi_{84}) \to \varphi_{182}$\lbl{A1}}
\BIC{$\varphi_{182}$}
\end{ptree}
\noindent$\Pi_{28}$:
\begin{ptree}
\AXC{$\stackrel{\Pi_{3}}{\vdots}$}\noLine
\UIC{$l \in \mathbb{R} \to \mathrm{ContIn}(f,A,c) \to \mathrm{Seq}(a) \to \varphi_{29}$}
\AXC{$\stackrel{\Pi_{26}}{\vdots}$}\noLine
\UIC{$\varphi_{151}$}
\BIC{$l \in \mathbb{R} \to \mathrm{ContIn}(f,A,c) \to \varphi_{27}$}
\AXC{$l \in \mathbb{R} \to l \in \mathbb{R}$\lbl{A1}}
\AXC{$\stackrel{\Pi_{27}}{\vdots}$}\noLine
\UIC{$\varphi_{182}$}
\BIC{$(l \in \mathbb{R} \to \mathrm{ContIn}(f,A,c) \to \varphi_{27}) \to l \in \mathbb{R} \to \varphi_{134}$}
\BIC{$l \in \mathbb{R} \to \varphi_{134}$}
\end{ptree}
\noindent$\Pi_{29}$:
\begin{ptree}
\AXC{$\varphi_{118}$\lbl{A2}}
\AXC{$\stackrel{\Pi_{28}}{\vdots}$}\noLine
\UIC{$l \in \mathbb{R} \to \varphi_{134}$}
\AXC{$(l \in \mathbb{R} \to \varphi_{134}) \to l \in \mathbb{R} \to \varphi_{131}$\lbl{A1}}
\BIC{$l \in \mathbb{R} \to \varphi_{131}$}
\AXC{$(l \in \mathbb{R} \to \varphi_{131}) \to \varphi_{118} \to l \in \mathbb{R} \to \varphi_{129}$\lbl{A1}}
\BIC{$\varphi_{118} \to l \in \mathbb{R} \to \varphi_{129}$}
\BIC{$l \in \mathbb{R} \to \varphi_{129}$}
\end{ptree}
\noindent$\Pi_{30}$:
\begin{ptree}
\AXC{$l \in \mathbb{R} \to \mathrm{ContIn}(f,A,c) \to \mathrm{Seq}(a) \to \varphi_{17}$\lbl{A1}}
\AXC{$\stackrel{\Pi_{29}}{\vdots}$}\noLine
\UIC{$l \in \mathbb{R} \to \varphi_{129}$}
\AXC{$(l \in \mathbb{R} \to \varphi_{129}) \to \varphi_{163}$\lbl{A1}}
\BIC{$\varphi_{163}$}
\BIC{$l \in \mathbb{R} \to \varphi_{117}$}
\AXC{$(l \in \mathbb{R} \to \varphi_{117}) \to \varphi_{196}$\lbl{A1}}
\BIC{$\varphi_{196}$}
\end{ptree}
\noindent$\Pi_{31}$:
\begin{ptree}
\AXC{$\stackrel{\Pi_{11}}{\vdots}$}\noLine
\UIC{$f \circ a(n) = f(a_{n}) \to \varphi_{88}$}
\AXC{$x = x$\lbl{A5}}
\UIC{$\forall x\,(x = x)$}
\AXC{$\forall x\,(x = x) \to l \cdot f \circ a(n) = l \cdot f \circ a(n)$\lbl{A2}}
\BIC{$l \cdot f \circ a(n) = l \cdot f \circ a(n)$}
\AXC{$\stackrel{\Pi_{25}}{\vdots}$}\noLine
\UIC{$\varphi_{187}$}
\BIC{$(f \circ a(n) = f(a_{n}) \to \varphi_{88}) \to l \in \mathbb{R} \to \varphi_{106}$}
\BIC{$l \in \mathbb{R} \to \varphi_{106}$}
\AXC{$\stackrel{\Pi_{30}}{\vdots}$}\noLine
\UIC{$\varphi_{196}$}
\BIC{$(l \odot f \circ a(n) = l \cdot f \circ a(n) \to \varphi_{87}) \to l \in \mathbb{R} \to \varphi_{109}$}
\end{ptree}
\noindent$\Pi_{32}$:
\begin{ptree}
\AXC{$\stackrel{\Pi_{10}}{\vdots}$}\noLine
\UIC{$\varphi_{34}$}
\AXC{$\varphi_{34} \to \varphi_{101}$\lbl{A2}}
\BIC{$\varphi_{101}$}
\AXC{$\varphi_{101} \to l \odot f \circ a(n) = l \cdot f \circ a(n) \to \varphi_{87}$\lbl{A2}}
\BIC{$l \odot f \circ a(n) = l \cdot f \circ a(n) \to \varphi_{87}$}
\AXC{$\stackrel{\Pi_{31}}{\vdots}$}\noLine
\UIC{$(l \odot f \circ a(n) = l \cdot f \circ a(n) \to \varphi_{87}) \to l \in \mathbb{R} \to \varphi_{109}$}
\BIC{$l \in \mathbb{R} \to \varphi_{109}$}
\end{ptree}
\noindent$\Pi_{33}$:
\begin{ptree}
\AXC{$y = z \to x = y \to x = z$\lbl{A6}}
\AXC{$(y = z \to x = y \to x = z) \to x = y \to y = z \to x = z$\lbl{A1}}
\BIC{$x = y \to y = z \to x = z$}
\UIC{$\forall z\,(x = y \to y = z \to x = z)$}
\UIC{$\forall y\,\forall z\,(x = y \to y = z \to x = z)$}
\UIC{$\forall x\,\forall y\,\forall z\,(x = y \to y = z \to x = z)$}
\AXC{$\forall x\,\forall y\,\forall z\,(x = y \to y = z \to x = z) \to \varphi_{43}$\lbl{A2}}
\BIC{$\varphi_{43}$}
\end{ptree}
\noindent$\Pi_{34}$:
\begin{ptree}
\AXC{$\stackrel{\Pi_{33}}{\vdots}$}\noLine
\UIC{$\varphi_{43}$}
\AXC{$\varphi_{43} \to \varphi_{104}$\lbl{A2}}
\BIC{$\varphi_{104}$}
\AXC{$\varphi_{104} \to (l f) \circ a(n) = (l f)(a_{n}) \to \varphi_{89}$\lbl{A2}}
\BIC{$(l f) \circ a(n) = (l f)(a_{n}) \to \varphi_{89}$}
\end{ptree}
\noindent$\Pi_{35}$:
\begin{ptree}
\AXC{$\forall l\,\forall f\,(l \in \mathbb{R} \to \mathrm{RFn}(f) \to \varphi_{105})$\lbl{[S19]}}
\AXC{$\varphi_{178}$\lbl{A2}}
\BIC{$\forall f\,(l \in \mathbb{R} \to \mathrm{RFn}(f) \to \varphi_{105})$}
\AXC{$\varphi_{176}$\lbl{A2}}
\BIC{$l \in \mathbb{R} \to \mathrm{RFn}(f) \to \varphi_{105}$}
\AXC{$(l \in \mathbb{R} \to \mathrm{RFn}(f) \to \varphi_{105}) \to (l \in \mathbb{R} \to l \in \mathbb{R}) \to \varphi_{143}$\lbl{A1}}
\BIC{$(l \in \mathbb{R} \to l \in \mathbb{R}) \to \varphi_{143}$}
\end{ptree}
\noindent$\Pi_{36}$:
\begin{ptree}
\AXC{$\stackrel{\Pi_{21}}{\vdots}$}\noLine
\UIC{$l \in \mathbb{R} \to \mathrm{ContIn}(f,A,c) \to \mathrm{RFn}(f)$}
\AXC{$l \in \mathbb{R} \to l \in \mathbb{R}$\lbl{A1}}
\AXC{$\stackrel{\Pi_{35}}{\vdots}$}\noLine
\UIC{$(l \in \mathbb{R} \to l \in \mathbb{R}) \to \varphi_{143}$}
\BIC{$\varphi_{143}$}
\BIC{$l \in \mathbb{R} \to \mathrm{ContIn}(f,A,c) \to \varphi_{105}$}
\end{ptree}
\noindent$\Pi_{37}$:
\begin{ptree}
\AXC{$\stackrel{\Pi_{36}}{\vdots}$}\noLine
\UIC{$l \in \mathbb{R} \to \mathrm{ContIn}(f,A,c) \to \varphi_{105}$}
\AXC{$(l \in \mathbb{R} \to \mathrm{ContIn}(f,A,c) \to \varphi_{105}) \to l \in \mathbb{R} \to \varphi_{9}$\lbl{A1}}
\BIC{$l \in \mathbb{R} \to \varphi_{9}$}
\end{ptree}
\noindent$\Pi_{38}$:
\begin{ptree}
\AXC{$\varphi_{79}$\lbl{A2}}
\AXC{$\stackrel{\Pi_{37}}{\vdots}$}\noLine
\UIC{$l \in \mathbb{R} \to \varphi_{9}$}
\AXC{$(l \in \mathbb{R} \to \varphi_{9}) \to \varphi_{79} \to l \in \mathbb{R} \to \varphi_{20}$\lbl{A1}}
\BIC{$\varphi_{79} \to l \in \mathbb{R} \to \varphi_{20}$}
\BIC{$l \in \mathbb{R} \to \varphi_{20}$}
\end{ptree}
\noindent$\Pi_{39}$:
\begin{ptree}
\AXC{$\forall n \in \mathbb{N}\,(a_{n} \in A) \to n \in \mathbb{N} \to a_{n} \in A$\lbl{A2}}
\AXC{$l \in \mathbb{R} \to \varphi_{69}$\lbl{A1}}
\AXC{$\varphi_{169}$\lbl{A1}}
\BIC{$(\forall n \in \mathbb{N}\,(a_{n} \in A) \to n \in \mathbb{N} \to a_{n} \in A) \to l \in \mathbb{R} \to \varphi_{61}$}
\BIC{$l \in \mathbb{R} \to \varphi_{61}$}
\AXC{$(l \in \mathbb{R} \to \varphi_{61}) \to \varphi_{152}$\lbl{A1}}
\BIC{$\varphi_{152}$}
\end{ptree}
\noindent$\Pi_{40}$:
\begin{ptree}
\AXC{$l \in \mathbb{R} \to \mathrm{ContIn}(f,A,c) \to \mathrm{Seq}(a) \to \varphi_{17}$\lbl{A1}}
\AXC{$\stackrel{\Pi_{39}}{\vdots}$}\noLine
\UIC{$\varphi_{152}$}
\BIC{$l \in \mathbb{R} \to \mathrm{ContIn}(f,A,c) \to \mathrm{Seq}(a) \to \varphi_{26}$}
\end{ptree}
\noindent$\Pi_{41}$:
\begin{ptree}
\AXC{$\forall X\,\forall Y\,\forall x\,(X \subseteq Y \to x \in X \to x \in Y)$\lbl{[B5]}}
\AXC{$\varphi_{37}$\lbl{A2}}
\BIC{$\forall Y\,\forall x\,(A \subseteq Y \to x \in A \to x \in Y)$}
\AXC{$\varphi_{50}$\lbl{A2}}
\BIC{$\forall x\,(A \subseteq \mathrm{dom}(f) \to x \in A \to x \in \mathrm{dom}(f))$}
\AXC{$\varphi_{85}$\lbl{A2}}
\BIC{$A \subseteq \mathrm{dom}(f) \to a_{n} \in A \to a_{n} \in \mathrm{dom}(f)$}
\end{ptree}
\noindent$\Pi_{42}$:
\begin{ptree}
\AXC{$\stackrel{\Pi_{41}}{\vdots}$}\noLine
\UIC{$A \subseteq \mathrm{dom}(f) \to a_{n} \in A \to a_{n} \in \mathrm{dom}(f)$}
\AXC{$(A \subseteq \mathrm{dom}(f) \to a_{n} \in A \to a_{n} \in \mathrm{dom}(f)) \to \varphi_{172}$\lbl{A1}}
\BIC{$\varphi_{172}$}
\end{ptree}
\noindent$\Pi_{43}$:
\begin{ptree}
\AXC{$\stackrel{\Pi_{40}}{\vdots}$}\noLine
\UIC{$l \in \mathbb{R} \to \mathrm{ContIn}(f,A,c) \to \mathrm{Seq}(a) \to \varphi_{26}$}
\AXC{$\stackrel{\Pi_{15}}{\vdots}$}\noLine
\UIC{$l \in \mathbb{R} \to \mathrm{ContIn}(f,A,c) \to A \subseteq \mathrm{dom}(f)$}
\AXC{$\stackrel{\Pi_{42}}{\vdots}$}\noLine
\UIC{$\varphi_{172}$}
\BIC{$\varphi_{155}$}
\BIC{$l \in \mathbb{R} \to \mathrm{ContIn}(f,A,c) \to \mathrm{Seq}(a) \to \varphi_{44}$}
\end{ptree}
\noindent$\Pi_{44}$:
\begin{ptree}
\AXC{$\stackrel{\Pi_{38}}{\vdots}$}\noLine
\UIC{$l \in \mathbb{R} \to \varphi_{20}$}
\AXC{$\stackrel{\Pi_{43}}{\vdots}$}\noLine
\UIC{$l \in \mathbb{R} \to \mathrm{ContIn}(f,A,c) \to \mathrm{Seq}(a) \to \varphi_{44}$}
\AXC{$\varphi_{185}$\lbl{A1}}
\BIC{$\varphi_{149}$}
\BIC{$l \in \mathbb{R} \to \mathrm{ContIn}(f,A,c) \to \mathrm{Seq}(a) \to \varphi_{65}$}
\AXC{$\varphi_{161}$\lbl{A1}}
\BIC{$\varphi_{54} \to n \in \mathbb{N} \to a_{n} \in \mathrm{dom}((l f))$}
\UIC{$\forall n\,(\varphi_{54} \to n \in \mathbb{N} \to a_{n} \in \mathrm{dom}((l f)))$}
\end{ptree}
\noindent$\Pi_{45}$:
\begin{ptree}
\AXC{$\varphi_{167}$\lbl{A3}}
\AXC{$\stackrel{\Pi_{44}}{\vdots}$}\noLine
\UIC{$\forall n\,(\varphi_{54} \to n \in \mathbb{N} \to a_{n} \in \mathrm{dom}((l f)))$}
\AXC{$\forall n\,(\varphi_{54} \to n \in \mathbb{N} \to a_{n} \in \mathrm{dom}((l f))) \to \varphi_{167} \to \varphi_{179}$\lbl{A1}}
\BIC{$\varphi_{167} \to \varphi_{179}$}
\BIC{$\varphi_{179}$}
\end{ptree}
\noindent$\Pi_{46}$:
\begin{ptree}
\AXC{$\varphi_{54} \to \forall n\,(\varphi_{54})$\lbl{A4}}
\AXC{$\stackrel{\Pi_{45}}{\vdots}$}\noLine
\UIC{$\varphi_{179}$}
\BIC{$l \in \mathbb{R} \to \mathrm{ContIn}(f,A,c) \to \mathrm{Seq}(a) \to \varphi_{71}$}
\end{ptree}
\noindent$\Pi_{47}$:
\begin{ptree}
\AXC{$\stackrel{\Pi_{36}}{\vdots}$}\noLine
\UIC{$l \in \mathbb{R} \to \mathrm{ContIn}(f,A,c) \to \varphi_{105}$}
\AXC{$(l \in \mathbb{R} \to \mathrm{ContIn}(f,A,c) \to \varphi_{105}) \to l \in \mathbb{R} \to \varphi_{41}$\lbl{A1}}
\BIC{$l \in \mathbb{R} \to \varphi_{41}$}
\end{ptree}
\noindent$\Pi_{48}$:
\begin{ptree}
\AXC{$\varphi_{122}$\lbl{A2}}
\AXC{$\stackrel{\Pi_{47}}{\vdots}$}\noLine
\UIC{$l \in \mathbb{R} \to \varphi_{41}$}
\AXC{$(l \in \mathbb{R} \to \varphi_{41}) \to \varphi_{122} \to l \in \mathbb{R} \to \varphi_{80}$\lbl{A1}}
\BIC{$\varphi_{122} \to l \in \mathbb{R} \to \varphi_{80}$}
\BIC{$l \in \mathbb{R} \to \varphi_{80}$}
\AXC{$(l \in \mathbb{R} \to \varphi_{80}) \to \varphi_{192}$\lbl{A1}}
\BIC{$\varphi_{192}$}
\end{ptree}
\noindent$\Pi_{49}$:
\begin{ptree}
\AXC{$\varphi_{68}$\lbl{A2}}
\AXC{$\stackrel{\Pi_{46}}{\vdots}$}\noLine
\UIC{$l \in \mathbb{R} \to \mathrm{ContIn}(f,A,c) \to \mathrm{Seq}(a) \to \varphi_{71}$}
\AXC{$\varphi_{189}$\lbl{A1}}
\BIC{$\varphi_{68} \to l \in \mathbb{R} \to \mathrm{ContIn}(f,A,c) \to \mathrm{Seq}(a) \to \varphi_{65}$}
\BIC{$l \in \mathbb{R} \to \mathrm{ContIn}(f,A,c) \to \mathrm{Seq}(a) \to \varphi_{65}$}
\AXC{$\stackrel{\Pi_{48}}{\vdots}$}\noLine
\UIC{$\varphi_{192}$}
\BIC{$\varphi_{160}$}
\end{ptree}
\noindent$\Pi_{50}$:
\begin{ptree}
\AXC{$\stackrel{\Pi_{36}}{\vdots}$}\noLine
\UIC{$l \in \mathbb{R} \to \mathrm{ContIn}(f,A,c) \to \varphi_{105}$}
\AXC{$\varphi_{145}$\lbl{A1}}
\BIC{$l \in \mathbb{R} \to \mathrm{ContIn}(f,A,c) \to \mathrm{RFn}((l f))$}
\end{ptree}
\noindent$\Pi_{51}$:
\begin{ptree}
\AXC{$\forall f\,\forall a\,(\mathrm{RFn}(f) \to \mathrm{Seq}(a) \to \varphi_{94})$\lbl{[S8]}}
\AXC{$\forall f\,\forall a\,(\mathrm{RFn}(f) \to \mathrm{Seq}(a) \to \varphi_{94}) \to \forall a\,(\varphi_{137})$\lbl{A2}}
\BIC{$\forall a\,(\varphi_{137})$}
\AXC{$\forall a\,(\varphi_{137}) \to \varphi_{137}$\lbl{A2}}
\BIC{$\varphi_{137}$}
\end{ptree}
\noindent$\Pi_{52}$:
\begin{ptree}
\AXC{$\varphi_{2}$\lbl{A1}}
\AXC{$\stackrel{\Pi_{50}}{\vdots}$}\noLine
\UIC{$l \in \mathbb{R} \to \mathrm{ContIn}(f,A,c) \to \mathrm{RFn}((l f))$}
\AXC{$\stackrel{\Pi_{51}}{\vdots}$}\noLine
\UIC{$\varphi_{137}$}
\AXC{$\varphi_{137} \to \varphi_{195}$\lbl{A1}}
\BIC{$\varphi_{195}$}
\BIC{$\varphi_{2} \to \varphi_{183}$}
\BIC{$\varphi_{183}$}
\end{ptree}
\noindent$\Pi_{53}$:
\begin{ptree}
\AXC{$\stackrel{\Pi_{46}}{\vdots}$}\noLine
\UIC{$l \in \mathbb{R} \to \mathrm{ContIn}(f,A,c) \to \mathrm{Seq}(a) \to \varphi_{71}$}
\AXC{$\stackrel{\Pi_{52}}{\vdots}$}\noLine
\UIC{$\varphi_{183}$}
\BIC{$l \in \mathbb{R} \to \varphi_{135}$}
\end{ptree}
\noindent$\Pi_{54}$:
\begin{ptree}
\AXC{$\varphi_{120}$\lbl{A2}}
\AXC{$\stackrel{\Pi_{53}}{\vdots}$}\noLine
\UIC{$l \in \mathbb{R} \to \varphi_{135}$}
\AXC{$(l \in \mathbb{R} \to \varphi_{135}) \to l \in \mathbb{R} \to \varphi_{132}$\lbl{A1}}
\BIC{$l \in \mathbb{R} \to \varphi_{132}$}
\AXC{$(l \in \mathbb{R} \to \varphi_{132}) \to \varphi_{120} \to l \in \mathbb{R} \to \varphi_{130}$\lbl{A1}}
\BIC{$\varphi_{120} \to l \in \mathbb{R} \to \varphi_{130}$}
\BIC{$l \in \mathbb{R} \to \varphi_{130}$}
\end{ptree}
\noindent$\Pi_{55}$:
\begin{ptree}
\AXC{$l \in \mathbb{R} \to \mathrm{ContIn}(f,A,c) \to \mathrm{Seq}(a) \to \varphi_{17}$\lbl{A1}}
\AXC{$\stackrel{\Pi_{54}}{\vdots}$}\noLine
\UIC{$l \in \mathbb{R} \to \varphi_{130}$}
\AXC{$(l \in \mathbb{R} \to \varphi_{130}) \to \varphi_{164}$\lbl{A1}}
\BIC{$\varphi_{164}$}
\BIC{$l \in \mathbb{R} \to \varphi_{119}$}
\AXC{$(l \in \mathbb{R} \to \varphi_{119}) \to \varphi_{197}$\lbl{A1}}
\BIC{$\varphi_{197}$}
\end{ptree}
\noindent$\Pi_{56}$:
\begin{ptree}
\AXC{$l \in \mathbb{R} \to \mathrm{ContIn}(f,A,c) \to \mathrm{Seq}(a) \to \varphi_{17}$\lbl{A1}}
\AXC{$\stackrel{\Pi_{49}}{\vdots}$}\noLine
\UIC{$\varphi_{160}$}
\BIC{$l \in \mathbb{R} \to \varphi_{107}$}
\AXC{$\stackrel{\Pi_{55}}{\vdots}$}\noLine
\UIC{$\varphi_{197}$}
\BIC{$((l f) \circ a(n) = (l f)(a_{n}) \to \varphi_{89}) \to l \in \mathbb{R} \to \varphi_{113}$}
\end{ptree}
\noindent$\Pi_{57}$:
\begin{ptree}
\AXC{$\stackrel{\Pi_{34}}{\vdots}$}\noLine
\UIC{$(l f) \circ a(n) = (l f)(a_{n}) \to \varphi_{89}$}
\AXC{$\stackrel{\Pi_{56}}{\vdots}$}\noLine
\UIC{$((l f) \circ a(n) = (l f)(a_{n}) \to \varphi_{89}) \to l \in \mathbb{R} \to \varphi_{113}$}
\BIC{$l \in \mathbb{R} \to \varphi_{113}$}
\AXC{$(l \in \mathbb{R} \to \varphi_{113}) \to \varphi_{198}$\lbl{A1}}
\BIC{$\varphi_{198}$}
\end{ptree}
\noindent$\Pi_{58}$:
\begin{ptree}
\AXC{$\stackrel{\Pi_{9}}{\vdots}$}\noLine
\UIC{$\varphi_{91}$}
\AXC{$\stackrel{\Pi_{32}}{\vdots}$}\noLine
\UIC{$l \in \mathbb{R} \to \varphi_{109}$}
\AXC{$(l \in \mathbb{R} \to \varphi_{109}) \to \varphi_{91} \to l \in \mathbb{R} \to \varphi_{110}$\lbl{A1}}
\BIC{$\varphi_{91} \to l \in \mathbb{R} \to \varphi_{110}$}
\BIC{$l \in \mathbb{R} \to \varphi_{110}$}
\AXC{$\stackrel{\Pi_{57}}{\vdots}$}\noLine
\UIC{$\varphi_{198}$}
\BIC{$((l f) \circ a(n) = l \cdot f(a_{n}) \to \varphi_{93}) \to l \in \mathbb{R} \to \varphi_{123}$}
\end{ptree}
\noindent$\Pi_{59}$:
\begin{ptree}
\AXC{$\stackrel{\Pi_{8}}{\vdots}$}\noLine
\UIC{$\varphi_{43}$}
\AXC{$\varphi_{43} \to \varphi_{96}$\lbl{A2}}
\BIC{$\varphi_{96}$}
\AXC{$\varphi_{96} \to (l f) \circ a(n) = l \cdot f(a_{n}) \to \varphi_{93}$\lbl{A2}}
\BIC{$(l f) \circ a(n) = l \cdot f(a_{n}) \to \varphi_{93}$}
\AXC{$\stackrel{\Pi_{58}}{\vdots}$}\noLine
\UIC{$((l f) \circ a(n) = l \cdot f(a_{n}) \to \varphi_{93}) \to l \in \mathbb{R} \to \varphi_{123}$}
\BIC{$l \in \mathbb{R} \to \varphi_{123}$}
\end{ptree}
\noindent$\Pi_{60}$:
\begin{ptree}
\AXC{$\varphi_{184}$\lbl{A3}}
\AXC{$\stackrel{\Pi_{7}}{\vdots}$}\noLine
\UIC{$\varphi_{100}$}
\AXC{$\stackrel{\Pi_{59}}{\vdots}$}\noLine
\UIC{$l \in \mathbb{R} \to \varphi_{123}$}
\AXC{$(l \in \mathbb{R} \to \varphi_{123}) \to \varphi_{100} \to l \in \mathbb{R} \to \varphi_{124}$\lbl{A1}}
\BIC{$\varphi_{100} \to l \in \mathbb{R} \to \varphi_{124}$}
\BIC{$l \in \mathbb{R} \to \varphi_{124}$}
\AXC{$\varphi_{181}$\lbl{A1}}
\BIC{$\varphi_{54} \to n \in \mathbb{N} \to l \odot f \circ a(n) = (l f) \circ a(n)$}
\UIC{$\forall n\,(\varphi_{54} \to n \in \mathbb{N} \to l \odot f \circ a(n) = (l f) \circ a(n))$}
\AXC{$\varphi_{203}$\lbl{A1}}
\BIC{$\varphi_{184} \to (\varphi_{54} \to \forall n\,(\varphi_{54})) \to l \in \mathbb{R} \to \varphi_{126}$}
\BIC{$(\varphi_{54} \to \forall n\,(\varphi_{54})) \to l \in \mathbb{R} \to \varphi_{126}$}
\end{ptree}
\noindent$\Pi_{61}$:
\begin{ptree}
\AXC{$\stackrel{\Pi_{53}}{\vdots}$}\noLine
\UIC{$l \in \mathbb{R} \to \varphi_{135}$}
\AXC{$\varphi_{177}$\lbl{A1}}
\BIC{$l \in \mathbb{R} \to \mathrm{ContIn}(f,A,c) \to \mathrm{Seq}(a) \to \varphi_{28}$}
\AXC{$\varphi_{48}$\lbl{[C13]}}
\AXC{$\varphi_{48} \to \varphi_{98}$\lbl{A2}}
\BIC{$\varphi_{98}$}
\AXC{$\varphi_{98} \to \varphi_{136}$\lbl{A2}}
\BIC{$\varphi_{136}$}
\AXC{$\varphi_{136} \to \varphi_{144}$\lbl{A2}}
\BIC{$\varphi_{144}$}
\AXC{$\varphi_{144} \to \varphi_{202}$\lbl{A1}}
\BIC{$\varphi_{202}$}
\BIC{$(l \in \mathbb{R} \to \varphi_{126}) \to \varphi_{162}$}
\end{ptree}
\noindent$\Pi_{62}$:
\begin{ptree}
\AXC{$\varphi_{54} \to \forall n\,(\varphi_{54})$\lbl{A4}}
\AXC{$\stackrel{\Pi_{60}}{\vdots}$}\noLine
\UIC{$(\varphi_{54} \to \forall n\,(\varphi_{54})) \to l \in \mathbb{R} \to \varphi_{126}$}
\BIC{$l \in \mathbb{R} \to \varphi_{126}$}
\AXC{$\stackrel{\Pi_{61}}{\vdots}$}\noLine
\UIC{$(l \in \mathbb{R} \to \varphi_{126}) \to \varphi_{162}$}
\BIC{$\varphi_{162}$}
\end{ptree}
\noindent$\Pi_{63}$:
\begin{ptree}
\AXC{$\forall f\,\forall A\,\forall c\,(\varphi_{139})$\lbl{DefContIn}}
\AXC{$\forall f\,\forall A\,\forall c\,(\varphi_{139}) \to \forall A\,\forall c\,(\varphi_{139})$\lbl{A2}}
\BIC{$\forall A\,\forall c\,(\varphi_{139})$}
\AXC{$\forall A\,\forall c\,(\varphi_{139}) \to \forall c\,(\varphi_{139})$\lbl{A2}}
\BIC{$\forall c\,(\varphi_{139})$}
\AXC{$\forall c\,(\varphi_{139}) \to \varphi_{139}$\lbl{A2}}
\BIC{$\varphi_{139}$}
\end{ptree}
\noindent$\Pi_{64}$:
\begin{ptree}
\AXC{$\stackrel{\Pi_{63}}{\vdots}$}\noLine
\UIC{$\varphi_{139}$}
\AXC{$\varphi_{1}$\lbl{A1}}
\AXC{$\varphi_{1} \to \varphi_{139} \to l \in \mathbb{R} \to \mathrm{ContIn}(f,A,c) \to c \in A$\lbl{A1}}
\BIC{$\varphi_{139} \to l \in \mathbb{R} \to \mathrm{ContIn}(f,A,c) \to c \in A$}
\BIC{$l \in \mathbb{R} \to \mathrm{ContIn}(f,A,c) \to c \in A$}
\end{ptree}
\noindent$\Pi_{65}$:
\begin{ptree}
\AXC{$\forall X\,\forall Y\,\forall x\,(X \subseteq Y \to x \in X \to x \in Y)$\lbl{[B5]}}
\AXC{$\varphi_{37}$\lbl{A2}}
\BIC{$\forall Y\,\forall x\,(A \subseteq Y \to x \in A \to x \in Y)$}
\AXC{$\varphi_{50}$\lbl{A2}}
\BIC{$\forall x\,(A \subseteq \mathrm{dom}(f) \to x \in A \to x \in \mathrm{dom}(f))$}
\AXC{$\varphi_{58}$\lbl{A2}}
\BIC{$A \subseteq \mathrm{dom}(f) \to c \in A \to c \in \mathrm{dom}(f)$}
\end{ptree}
\noindent$\Pi_{66}$:
\begin{ptree}
\AXC{$\stackrel{\Pi_{65}}{\vdots}$}\noLine
\UIC{$A \subseteq \mathrm{dom}(f) \to c \in A \to c \in \mathrm{dom}(f)$}
\AXC{$(A \subseteq \mathrm{dom}(f) \to c \in A \to c \in \mathrm{dom}(f)) \to \varphi_{140}$\lbl{A1}}
\BIC{$\varphi_{140}$}
\end{ptree}
\noindent$\Pi_{67}$:
\begin{ptree}
\AXC{$\stackrel{\Pi_{64}}{\vdots}$}\noLine
\UIC{$l \in \mathbb{R} \to \mathrm{ContIn}(f,A,c) \to c \in A$}
\AXC{$\stackrel{\Pi_{15}}{\vdots}$}\noLine
\UIC{$l \in \mathbb{R} \to \mathrm{ContIn}(f,A,c) \to A \subseteq \mathrm{dom}(f)$}
\AXC{$\stackrel{\Pi_{66}}{\vdots}$}\noLine
\UIC{$\varphi_{140}$}
\BIC{$\varphi_{115}$}
\BIC{$l \in \mathbb{R} \to \mathrm{ContIn}(f,A,c) \to \varphi_{18}$}
\end{ptree}
\noindent$\Pi_{68}$:
\begin{ptree}
\AXC{$\stackrel{\Pi_{67}}{\vdots}$}\noLine
\UIC{$l \in \mathbb{R} \to \mathrm{ContIn}(f,A,c) \to \varphi_{18}$}
\AXC{$(l \in \mathbb{R} \to \mathrm{ContIn}(f,A,c) \to \varphi_{18}) \to \varphi_{138}$\lbl{A1}}
\BIC{$\varphi_{138}$}
\end{ptree}
\noindent$\Pi_{69}$:
\begin{ptree}
\AXC{$\varphi_{102}$\lbl{A2}}
\AXC{$\stackrel{\Pi_{47}}{\vdots}$}\noLine
\UIC{$l \in \mathbb{R} \to \varphi_{41}$}
\AXC{$(l \in \mathbb{R} \to \varphi_{41}) \to \varphi_{102} \to l \in \mathbb{R} \to \varphi_{31}$\lbl{A1}}
\BIC{$\varphi_{102} \to l \in \mathbb{R} \to \varphi_{31}$}
\BIC{$l \in \mathbb{R} \to \varphi_{31}$}
\AXC{$(l \in \mathbb{R} \to \varphi_{31}) \to \varphi_{154}$\lbl{A1}}
\BIC{$\varphi_{154}$}
\end{ptree}
\noindent$\Pi_{70}$:
\begin{ptree}
\AXC{$\varphi_{51}$\lbl{A2}}
\AXC{$\stackrel{\Pi_{37}}{\vdots}$}\noLine
\UIC{$l \in \mathbb{R} \to \varphi_{9}$}
\AXC{$(l \in \mathbb{R} \to \varphi_{9}) \to \varphi_{51} \to l \in \mathbb{R} \to \varphi_{6}$\lbl{A1}}
\BIC{$\varphi_{51} \to l \in \mathbb{R} \to \varphi_{6}$}
\BIC{$l \in \mathbb{R} \to \varphi_{6}$}
\AXC{$\stackrel{\Pi_{68}}{\vdots}$}\noLine
\UIC{$\varphi_{138}$}
\BIC{$l \in \mathbb{R} \to \mathrm{ContIn}(f,A,c) \to \mathrm{Seq}(a) \to \varphi_{19}$}
\AXC{$\stackrel{\Pi_{69}}{\vdots}$}\noLine
\UIC{$\varphi_{154}$}
\BIC{$l \in \mathbb{R} \to \mathrm{ContIn}(f,A,c) \to \mathrm{Seq}(a) \to \varphi_{45}$}
\AXC{$\varphi_{180}$\lbl{A1}}
\BIC{$\varphi_{142}$}
\end{ptree}
\noindent$\Pi_{71}$:
\begin{ptree}
\AXC{$\stackrel{\Pi_{6}}{\vdots}$}\noLine
\UIC{$\forall x\,\forall y\,(x = y \to y = x)$}
\AXC{$\varphi_{16}$\lbl{A2}}
\BIC{$\forall y\,((l f)(c) = y \to y = (l f)(c))$}
\AXC{$\varphi_{99}$\lbl{A2}}
\BIC{$(l f)(c) = l \cdot f(c) \to l \cdot f(c) = (l f)(c)$}
\AXC{$\stackrel{\Pi_{70}}{\vdots}$}\noLine
\UIC{$\varphi_{142}$}
\BIC{$l \in \mathbb{R} \to \mathrm{ContIn}(f,A,c) \to \mathrm{Seq}(a) \to \varphi_{46}$}
\AXC{$\varphi_{200}$\lbl{A1}}
\BIC{$\varphi_{194}$}
\end{ptree}
\noindent$\Pi_{72}$:
\begin{ptree}
\AXC{$\stackrel{\Pi_{1}}{\vdots}$}\noLine
\UIC{$l \cdot f(c) = (l f)(c) \to \varphi_{73}$}
\AXC{$\stackrel{\Pi_{5}}{\vdots}$}\noLine
\UIC{$l \in \mathbb{R} \to \mathrm{ContIn}(f,A,c) \to \mathrm{Seq}(a) \to \varphi_{64}$}
\AXC{$\stackrel{\Pi_{62}}{\vdots}$}\noLine
\UIC{$\varphi_{162}$}
\BIC{$l \in \mathbb{R} \to \mathrm{ContIn}(f,A,c) \to \mathrm{Seq}(a) \to \varphi_{67}$}
\AXC{$\stackrel{\Pi_{71}}{\vdots}$}\noLine
\UIC{$\varphi_{194}$}
\BIC{$\varphi_{165}$}
\BIC{$l \in \mathbb{R} \to \mathrm{ContIn}(f,A,c) \to \mathrm{Seq}(a) \to \varphi_{75}$}
\AXC{$\varphi_{168}$\lbl{A1}}
\BIC{$l \in \mathbb{R} \land \mathrm{ContIn}(f,A,c) \to \mathrm{Seq}(a) \to \varphi_{75}$}
\UIC{$\varphi_{111}$}
\end{ptree}
\noindent$\Pi_{73}$:
\begin{ptree}
\AXC{$\forall X\,\forall Y\,\forall x\,(X \subseteq Y \to x \in X \to x \in Y)$\lbl{[B5]}}
\AXC{$\varphi_{37}$\lbl{A2}}
\BIC{$\forall Y\,\forall x\,(A \subseteq Y \to x \in A \to x \in Y)$}
\AXC{$\varphi_{50}$\lbl{A2}}
\BIC{$\forall x\,(A \subseteq \mathrm{dom}(f) \to x \in A \to x \in \mathrm{dom}(f))$}
\AXC{$\varphi_{59}$\lbl{A2}}
\BIC{$A \subseteq \mathrm{dom}(f) \to u \in A \to u \in \mathrm{dom}(f)$}
\end{ptree}
\noindent$\Pi_{74}$:
\begin{ptree}
\AXC{$\stackrel{\Pi_{73}}{\vdots}$}\noLine
\UIC{$A \subseteq \mathrm{dom}(f) \to u \in A \to u \in \mathrm{dom}(f)$}
\AXC{$(A \subseteq \mathrm{dom}(f) \to u \in A \to u \in \mathrm{dom}(f)) \to \varphi_{128}$\lbl{A1}}
\BIC{$\varphi_{128}$}
\end{ptree}
\noindent$\Pi_{75}$:
\begin{ptree}
\AXC{$l \in \mathbb{R} \to \mathrm{ContIn}(f,A,c) \to u \in A \to u \in A$\lbl{A1}}
\AXC{$\stackrel{\Pi_{15}}{\vdots}$}\noLine
\UIC{$l \in \mathbb{R} \to \mathrm{ContIn}(f,A,c) \to A \subseteq \mathrm{dom}(f)$}
\AXC{$\stackrel{\Pi_{74}}{\vdots}$}\noLine
\UIC{$\varphi_{128}$}
\BIC{$(l \in \mathbb{R} \to \mathrm{ContIn}(f,A,c) \to u \in A \to u \in A) \to \varphi_{4}$}
\BIC{$\varphi_{4}$}
\end{ptree}
\noindent$\Pi_{76}$:
\begin{ptree}
\AXC{$\varphi_{52}$\lbl{A2}}
\AXC{$\stackrel{\Pi_{37}}{\vdots}$}\noLine
\UIC{$l \in \mathbb{R} \to \varphi_{9}$}
\AXC{$(l \in \mathbb{R} \to \varphi_{9}) \to \varphi_{52} \to l \in \mathbb{R} \to \varphi_{7}$\lbl{A1}}
\BIC{$\varphi_{52} \to l \in \mathbb{R} \to \varphi_{7}$}
\BIC{$l \in \mathbb{R} \to \varphi_{7}$}
\AXC{$\stackrel{\Pi_{75}}{\vdots}$}\noLine
\UIC{$\varphi_{4}$}
\AXC{$\varphi_{4} \to (l \in \mathbb{R} \to \varphi_{7}) \to \varphi_{14}$\lbl{A1}}
\BIC{$(l \in \mathbb{R} \to \varphi_{7}) \to \varphi_{14}$}
\BIC{$\varphi_{14}$}
\end{ptree}
\noindent$\Pi_{77}$:
\begin{ptree}
\AXC{$\varphi_{12}$\lbl{A4}}
\AXC{$\forall u\,(\varphi_{15}) \to \varphi_{21}$\lbl{A3}}
\AXC{$\stackrel{\Pi_{76}}{\vdots}$}\noLine
\UIC{$\varphi_{14}$}
\AXC{$\varphi_{14} \to \varphi_{15}$\lbl{A1}}
\BIC{$\varphi_{15}$}
\UIC{$\forall u\,(\varphi_{15})$}
\AXC{$\forall u\,(\varphi_{15}) \to (\forall u\,(\varphi_{15}) \to \varphi_{21}) \to \varphi_{12} \to l \in \mathbb{R} \to \varphi_{3}$\lbl{A1}}
\BIC{$(\forall u\,(\varphi_{15}) \to \varphi_{21}) \to \varphi_{12} \to l \in \mathbb{R} \to \varphi_{3}$}
\BIC{$\varphi_{12} \to l \in \mathbb{R} \to \varphi_{3}$}
\BIC{$l \in \mathbb{R} \to \varphi_{3}$}
\end{ptree}
\noindent$\Pi_{78}$:
\begin{ptree}
\AXC{$\stackrel{\Pi_{77}}{\vdots}$}\noLine
\UIC{$l \in \mathbb{R} \to \varphi_{3}$}
\AXC{$\forall x\,\forall y\,(x \subseteq y \leftrightarrow \forall u \in x\,(u \in y))$\lbl{Def$\subseteq$}}
\AXC{$\varphi_{39}$\lbl{A2}}
\BIC{$\forall y\,(A \subseteq y \leftrightarrow \forall u \in A\,(u \in y))$}
\AXC{$\varphi_{86}$\lbl{A2}}
\BIC{$A \subseteq \mathrm{dom}((l f)) \leftrightarrow \forall u \in A\,(u \in \mathrm{dom}((l f)))$}
\AXC{$\varphi_{141}$\lbl{A1}}
\BIC{$\varphi_{103}$}
\BIC{$l \in \mathbb{R} \to \mathrm{ContIn}(f,A,c) \to A \subseteq \mathrm{dom}((l f))$}
\end{ptree}
\noindent$\Pi_{79}$:
\begin{ptree}
\AXC{$\varphi_{127}$\lbl{[G2]}}
\AXC{$\varphi_{188}$\lbl{A2}}
\BIC{$\forall A\,\forall c\,(\mathrm{RFn}((l f)) \to A \subseteq \mathrm{dom}((l f)) \to \varphi_{121})$}
\AXC{$\varphi_{193}$\lbl{A2}}
\BIC{$\forall c\,(\mathrm{RFn}((l f)) \to A \subseteq \mathrm{dom}((l f)) \to \varphi_{121})$}
\AXC{$\varphi_{191}$\lbl{A2}}
\BIC{$\mathrm{RFn}((l f)) \to A \subseteq \mathrm{dom}((l f)) \to \varphi_{121}$}
\end{ptree}
\noindent$\Pi_{80}$:
\begin{ptree}
\AXC{$\stackrel{\Pi_{79}}{\vdots}$}\noLine
\UIC{$\mathrm{RFn}((l f)) \to A \subseteq \mathrm{dom}((l f)) \to \varphi_{121}$}
\AXC{$(\mathrm{RFn}((l f)) \to A \subseteq \mathrm{dom}((l f)) \to \varphi_{121}) \to \varphi_{186}$\lbl{A1}}
\BIC{$\varphi_{186}$}
\end{ptree}
\noindent$\Pi_{81}$:
\begin{ptree}
\AXC{$\stackrel{\Pi_{78}}{\vdots}$}\noLine
\UIC{$l \in \mathbb{R} \to \mathrm{ContIn}(f,A,c) \to A \subseteq \mathrm{dom}((l f))$}
\AXC{$\stackrel{\Pi_{50}}{\vdots}$}\noLine
\UIC{$l \in \mathbb{R} \to \mathrm{ContIn}(f,A,c) \to \mathrm{RFn}((l f))$}
\AXC{$\stackrel{\Pi_{80}}{\vdots}$}\noLine
\UIC{$\varphi_{186}$}
\BIC{$\varphi_{174}$}
\BIC{$(l \in \mathbb{R} \to \mathrm{ContIn}(f,A,c) \to c \in A) \to \varphi_{112} \to \varphi_{5}$}
\end{ptree}
\noindent$\Pi_{82}$:
\begin{ptree}
\AXC{$\stackrel{\Pi_{64}}{\vdots}$}\noLine
\UIC{$l \in \mathbb{R} \to \mathrm{ContIn}(f,A,c) \to c \in A$}
\AXC{$\stackrel{\Pi_{81}}{\vdots}$}\noLine
\UIC{$(l \in \mathbb{R} \to \mathrm{ContIn}(f,A,c) \to c \in A) \to \varphi_{112} \to \varphi_{5}$}
\BIC{$\varphi_{112} \to \varphi_{5}$}
\end{ptree}
\noindent 主木:
\begin{ptree}
\AXC{$\varphi_{11}$\lbl{A4}}
\AXC{$\varphi_{111} \to \varphi_{116}$\lbl{A3}}
\AXC{$\stackrel{\Pi_{72}}{\vdots}$}\noLine
\UIC{$\varphi_{111}$}
\AXC{$\varphi_{111} \to (\varphi_{111} \to \varphi_{116}) \to \varphi_{11} \to \varphi_{112}$\lbl{A1}}
\BIC{$(\varphi_{111} \to \varphi_{116}) \to \varphi_{11} \to \varphi_{112}$}
\BIC{$\varphi_{11} \to \varphi_{112}$}
\BIC{$\varphi_{112}$}
\AXC{$\stackrel{\Pi_{82}}{\vdots}$}\noLine
\UIC{$\varphi_{112} \to \varphi_{5}$}
\BIC{$\varphi_{5}$}
\UIC{$\forall c\,(\varphi_{5})$}
\UIC{$\forall A\,\forall c\,(\varphi_{5})$}
\UIC{$\forall f\,\forall A\,\forall c\,(\varphi_{5})$}
\UIC{$\forall l\,\forall f\,\forall A\,\forall c\,(\varphi_{5})$}
\end{ptree}
\begin{abbrevlist}
\item $\varphi_{1} :\equiv l \in \mathbb{R} \to \mathrm{ContIn}(f,A,c) \to \mathrm{ContIn}(f,A,c)$
\item $\varphi_{2} :\equiv l \in \mathbb{R} \to \mathrm{ContIn}(f,A,c) \to \mathrm{Seq}(a) \to \mathrm{Seq}(a)$
\item $\varphi_{3} :\equiv \mathrm{ContIn}(f,A,c) \to \forall u \in A\,(u \in \mathrm{dom}((l f)))$
\item $\varphi_{4} :\equiv l \in \mathbb{R} \to \mathrm{ContIn}(f,A,c) \to u \in A \to u \in \mathrm{dom}(f)$
\item $\varphi_{5} :\equiv l \in \mathbb{R} \to \mathrm{ContIn}(f,A,c) \to \mathrm{ContIn}((l f),A,c)$
\item $\varphi_{6} :\equiv \mathrm{ContIn}(f,A,c) \to (c \in \mathrm{dom}((l f)) \leftrightarrow c \in \mathrm{dom}(f))$
\item $\varphi_{7} :\equiv \mathrm{ContIn}(f,A,c) \to (u \in \mathrm{dom}((l f)) \leftrightarrow u \in \mathrm{dom}(f))$
\item $\varphi_{8} :\equiv \forall n \in \mathbb{N}\,(a_{n} \in A) \to \mathrm{Conv}(a,c) \to \mathrm{Conv}(a,c)$
\item $\varphi_{9} :\equiv \mathrm{ContIn}(f,A,c) \to \forall x\,(x \in \mathrm{dom}((l f)) \leftrightarrow x \in \mathrm{dom}(f))$
\item $\varphi_{10} :\equiv \forall a\,\forall \alpha\,(l \in \mathbb{R} \to \mathrm{Conv}(a,\alpha) \to \mathrm{Conv}(l \odot a,l \cdot \alpha))$
\item $\varphi_{11} :\equiv l \in \mathbb{R} \land \mathrm{ContIn}(f,A,c) \to \forall a\,(l \in \mathbb{R} \land \mathrm{ContIn}(f,A,c))$
\item $\varphi_{12} :\equiv l \in \mathbb{R} \land \mathrm{ContIn}(f,A,c) \to \forall u\,(l \in \mathbb{R} \land \mathrm{ContIn}(f,A,c))$
\item $\varphi_{13} :\equiv \forall c\,\forall a\,\forall \alpha\,(c \in \mathbb{R} \to \mathrm{Conv}(a,\alpha) \to \mathrm{Conv}(c \odot a,c \cdot \alpha))$
\item $\varphi_{14} :\equiv l \in \mathbb{R} \to \mathrm{ContIn}(f,A,c) \to u \in A \to u \in \mathrm{dom}((l f))$
\item $\varphi_{15} :\equiv l \in \mathbb{R} \land \mathrm{ContIn}(f,A,c) \to u \in A \to u \in \mathrm{dom}((l f))$
\item $\varphi_{16} :\equiv \forall x\,\forall y\,(x = y \to y = x) \to \forall y\,((l f)(c) = y \to y = (l f)(c))$
\item $\varphi_{17} :\equiv \forall n \in \mathbb{N}\,(a_{n} \in A) \to \mathrm{Conv}(a,c) \to n \in \mathbb{N} \to n \in \mathbb{N}$
\item $\varphi_{18} :\equiv \mathrm{Seq}(a) \to \forall n \in \mathbb{N}\,(a_{n} \in A) \to \mathrm{Conv}(a,c) \to c \in \mathrm{dom}(f)$
\item $\varphi_{19} :\equiv \forall n \in \mathbb{N}\,(a_{n} \in A) \to \mathrm{Conv}(a,c) \to c \in \mathrm{dom}((l f))$
\item $\varphi_{20} :\equiv \mathrm{ContIn}(f,A,c) \to (a_{n} \in \mathrm{dom}((l f)) \leftrightarrow a_{n} \in \mathrm{dom}(f))$
\item $\varphi_{21} :\equiv \forall u\,(l \in \mathbb{R} \land \mathrm{ContIn}(f,A,c)) \to \forall u \in A\,(u \in \mathrm{dom}((l f)))$
\item $\varphi_{22} :\equiv z = w \to \mathrm{Conv}((l f) \circ a,z) \to \mathrm{Conv}((l f) \circ a,w)$
\item $\varphi_{23} :\equiv \mathrm{Conv}(a,c) \to n \in \mathbb{N} \to l \cdot f \circ a(n) = l \cdot f(a_{n})$
\item $\varphi_{24} :\equiv \mathrm{Conv}(a,c) \to n \in \mathbb{N} \to (l f)(a_{n}) = l \cdot f(a_{n})$
\item $\varphi_{25} :\equiv \mathrm{Conv}(f \circ a,f(c)) \to \mathrm{Conv}(l \odot f \circ a,l \cdot f(c))$
\item $\varphi_{26} :\equiv \forall n \in \mathbb{N}\,(a_{n} \in A) \to \mathrm{Conv}(a,c) \to n \in \mathbb{N} \to a_{n} \in A$
\item $\varphi_{27} :\equiv \mathrm{Seq}(a) \to \forall n \in \mathbb{N}\,(a_{n} \in A) \to \mathrm{Conv}(a,c) \to \mathrm{Seq}(f \circ a)$
\item $\varphi_{28} :\equiv \forall n \in \mathbb{N}\,(a_{n} \in A) \to \mathrm{Conv}(a,c) \to \mathrm{Seq}((l f) \circ a)$
\item $\varphi_{29} :\equiv \forall n \in \mathbb{N}\,(a_{n} \in A) \to \mathrm{Conv}(a,c) \to \mathrm{Conv}(f \circ a,f(c))$
\item $\varphi_{30} :\equiv \mathrm{Conv}(a,c) \to n \in \mathbb{N} \to n \in \mathbb{N} \to f \circ a(n) = f(a_{n})$
\item $\varphi_{31} :\equiv \mathrm{ContIn}(f,A,c) \to c \in \mathrm{dom}((l f)) \to (l f)(c) = l \cdot f(c)$
\item $\varphi_{32} :\equiv \mathrm{Conv}(a,c) \to n \in \mathbb{N} \to l \odot f \circ a(n) = l \cdot f(a_{n})$
\item $\varphi_{33} :\equiv \mathrm{Conv}(a,c) \to n \in \mathbb{N} \to l \cdot f(a_{n}) = l \odot f \circ a(n)$
\item $\varphi_{34} :\equiv \forall y\,\forall z\,(l \odot f \circ a(n) = y \to y = z \to l \odot f \circ a(n) = z)$
\item $\varphi_{35} :\equiv \mathrm{Conv}(a,c) \to n \in \mathbb{N} \to (l f) \circ a(n) = l \cdot f(a_{n})$
\item $\varphi_{36} :\equiv \forall w\,(z = w \to l \cdot f \circ a(n) = l \cdot z \to l \cdot f \circ a(n) = l \cdot w)$
\item $\varphi_{37} :\equiv \forall X\,\forall Y\,\forall x\,(X \subseteq Y \to x \in X \to x \in Y) \to \forall Y\,\forall x\,(A \subseteq Y \to x \in A \to x \in Y)$
\item $\varphi_{38} :\equiv \forall a\,\forall \alpha\,(\mathrm{Conv}(a,\alpha) \to \mathrm{Seq}(a)) \to \forall \alpha\,(\mathrm{Conv}(f \circ a,\alpha) \to \mathrm{Seq}(f \circ a))$
\item $\varphi_{39} :\equiv \forall x\,\forall y\,(x \subseteq y \leftrightarrow \forall u \in x\,(u \in y)) \to \forall y\,(A \subseteq y \leftrightarrow \forall u \in A\,(u \in y))$
\item $\varphi_{40} :\equiv \mathrm{Conv}(a,c) \to n \in \mathbb{N} \to \forall n \in \mathbb{N}\,(f \circ a(n) = f(a_{n}))$
\item $\varphi_{41} :\equiv \mathrm{ContIn}(f,A,c) \to \forall x \in \mathrm{dom}((l f))\,((l f)(x) = l \cdot f(x))$
\item $\varphi_{42} :\equiv \forall \alpha\,(l \in \mathbb{R} \to \mathrm{Conv}(f \circ a,\alpha) \to \mathrm{Conv}(l \odot f \circ a,l \cdot \alpha))$
\item $\varphi_{43} :\equiv \forall y\,\forall z\,((l f) \circ a(n) = y \to y = z \to (l f) \circ a(n) = z)$
\item $\varphi_{44} :\equiv \forall n \in \mathbb{N}\,(a_{n} \in A) \to \mathrm{Conv}(a,c) \to n \in \mathbb{N} \to a_{n} \in \mathrm{dom}(f)$
\item $\varphi_{45} :\equiv \forall n \in \mathbb{N}\,(a_{n} \in A) \to \mathrm{Conv}(a,c) \to (l f)(c) = l \cdot f(c)$
\item $\varphi_{46} :\equiv \forall n \in \mathbb{N}\,(a_{n} \in A) \to \mathrm{Conv}(a,c) \to l \cdot f(c) = (l f)(c)$
\item $\varphi_{47} :\equiv \mathrm{Conv}(a,c) \to n \in \mathbb{N} \to l \odot f \circ a(n) = l \cdot f \circ a(n)$
\item $\varphi_{48} :\equiv \forall a\,\forall b\,\forall \alpha\,(\mathrm{Seq}(b) \to \forall n \in \mathbb{N}\,(a_{n} = b_{n}) \to \mathrm{Conv}(a,\alpha) \to \mathrm{Conv}(b,\alpha))$
\item $\varphi_{49} :\equiv \mathrm{Conv}(a,c) \to n \in \mathbb{N} \to (l f) \circ a(n) = (l f)(a_{n})$
\item $\varphi_{50} :\equiv \forall Y\,\forall x\,(A \subseteq Y \to x \in A \to x \in Y) \to \forall x\,(A \subseteq \mathrm{dom}(f) \to x \in A \to x \in \mathrm{dom}(f))$
\item $\varphi_{51} :\equiv \forall x\,(x \in \mathrm{dom}((l f)) \leftrightarrow x \in \mathrm{dom}(f)) \to (c \in \mathrm{dom}((l f)) \leftrightarrow c \in \mathrm{dom}(f))$
\item $\varphi_{52} :\equiv \forall x\,(x \in \mathrm{dom}((l f)) \leftrightarrow x \in \mathrm{dom}(f)) \to (u \in \mathrm{dom}((l f)) \leftrightarrow u \in \mathrm{dom}(f))$
\item $\varphi_{53} :\equiv \forall n \in \mathbb{N}\,(a_{n} \in A) \to \mathrm{Conv}(a,c) \to \forall n \in \mathbb{N}\,(a_{n} \in \mathrm{dom}(f))$
\item $\varphi_{54} :\equiv l \in \mathbb{R} \land \mathrm{ContIn}(f,A,c) \land \mathrm{Seq}(a) \land \forall n \in \mathbb{N}\,(a_{n} \in A) \land \mathrm{Conv}(a,c)$
\item $\varphi_{55} :\equiv \mathrm{Conv}(a,c) \to n \in \mathbb{N} \to (l f) \circ a(n) = l \odot f \circ a(n)$
\item $\varphi_{56} :\equiv \mathrm{Conv}(a,c) \to n \in \mathbb{N} \to l \odot f \circ a(n) = (l f) \circ a(n)$
\item $\varphi_{57} :\equiv \forall x\,\forall y\,(x = y \to y = x) \to \forall y\,(l \odot f \circ a(n) = y \to y = l \odot f \circ a(n))$
\item $\varphi_{58} :\equiv \forall x\,(A \subseteq \mathrm{dom}(f) \to x \in A \to x \in \mathrm{dom}(f)) \to A \subseteq \mathrm{dom}(f) \to c \in A \to c \in \mathrm{dom}(f)$
\item $\varphi_{59} :\equiv \forall x\,(A \subseteq \mathrm{dom}(f) \to x \in A \to x \in \mathrm{dom}(f)) \to A \subseteq \mathrm{dom}(f) \to u \in A \to u \in \mathrm{dom}(f)$
\item $\varphi_{60} :\equiv \mathrm{Conv}(a,c) \to \mathrm{Seq}(f \circ a) \land \forall n \in \mathbb{N}\,(f \circ a(n) = f(a_{n}))$
\item $\varphi_{61} :\equiv \mathrm{ContIn}(f,A,c) \to \mathrm{Seq}(a) \to \forall n \in \mathbb{N}\,(a_{n} \in A) \to n \in \mathbb{N} \to a_{n} \in A$
\item $\varphi_{62} :\equiv \mathrm{Conv}(a,c) \to \forall n \in \mathbb{N}\,(l \odot f \circ a(n) = (l f) \circ a(n))$
\item $\varphi_{63} :\equiv \forall x\,\forall y\,(x = y \to y = x) \to \forall y\,((l f) \circ a(n) = y \to y = (l f) \circ a(n))$
\item $\varphi_{64} :\equiv \forall n \in \mathbb{N}\,(a_{n} \in A) \to \mathrm{Conv}(a,c) \to \mathrm{Conv}(l \odot f \circ a,l \cdot f(c))$
\item $\varphi_{65} :\equiv \forall n \in \mathbb{N}\,(a_{n} \in A) \to \mathrm{Conv}(a,c) \to n \in \mathbb{N} \to a_{n} \in \mathrm{dom}((l f))$
\item $\varphi_{66} :\equiv \mathrm{Conv}(l \odot f \circ a,l \cdot f(c)) \to \mathrm{Conv}((l f) \circ a,l \cdot f(c))$
\item $\varphi_{67} :\equiv \forall n \in \mathbb{N}\,(a_{n} \in A) \to \mathrm{Conv}(a,c) \to \mathrm{Conv}((l f) \circ a,l \cdot f(c))$
\item $\varphi_{68} :\equiv \forall n \in \mathbb{N}\,(a_{n} \in \mathrm{dom}((l f))) \to n \in \mathbb{N} \to a_{n} \in \mathrm{dom}((l f))$
\item $\varphi_{69} :\equiv \mathrm{ContIn}(f,A,c) \to \mathrm{Seq}(a) \to \forall n \in \mathbb{N}\,(a_{n} \in A) \to \forall n \in \mathbb{N}\,(a_{n} \in A)$
\item $\varphi_{70} :\equiv \forall a\,(l \in \mathbb{R} \to \mathrm{Seq}(a) \to \mathrm{Seq}(l \odot a) \land \forall n \in \mathbb{N}\,(l \odot a(n) = l \cdot a_{n}))$
\item $\varphi_{71} :\equiv \forall n \in \mathbb{N}\,(a_{n} \in A) \to \mathrm{Conv}(a,c) \to \forall n \in \mathbb{N}\,(a_{n} \in \mathrm{dom}((l f)))$
\item $\varphi_{72} :\equiv \mathrm{Conv}(a,c) \to n \in \mathbb{N} \to n \in \mathbb{N} \to l \odot f \circ a(n) = l \cdot f \circ a(n)$
\item $\varphi_{73} :\equiv \mathrm{Conv}((l f) \circ a,l \cdot f(c)) \to \mathrm{Conv}((l f) \circ a,(l f)(c))$
\item $\varphi_{74} :\equiv \mathrm{Conv}(a,c) \to n \in \mathbb{N} \to n \in \mathbb{N} \to (l f) \circ a(n) = (l f)(a_{n})$
\item $\varphi_{75} :\equiv \forall n \in \mathbb{N}\,(a_{n} \in A) \to \mathrm{Conv}(a,c) \to \mathrm{Conv}((l f) \circ a,(l f)(c))$
\item $\varphi_{76} :\equiv \forall c\,\forall a\,(c \in \mathbb{R} \to \mathrm{Seq}(a) \to \mathrm{Seq}(c \odot a) \land \forall n \in \mathbb{N}\,(c \odot a(n) = c \cdot a_{n}))$
\item $\varphi_{77} :\equiv \forall n \in \mathbb{N}\,(a_{n} \in A) \to \mathrm{Conv}(a,c) \to n \in \mathbb{N} \to f \circ a(n) = f(a_{n})$
\item $\varphi_{78} :\equiv \mathrm{Conv}(a,c) \to n \in \mathbb{N} \to \forall n \in \mathbb{N}\,(l \odot f \circ a(n) = l \cdot f \circ a(n))$
\item $\varphi_{79} :\equiv \forall x\,(x \in \mathrm{dom}((l f)) \leftrightarrow x \in \mathrm{dom}(f)) \to (a_{n} \in \mathrm{dom}((l f)) \leftrightarrow a_{n} \in \mathrm{dom}(f))$
\item $\varphi_{80} :\equiv \mathrm{ContIn}(f,A,c) \to a_{n} \in \mathrm{dom}((l f)) \to (l f)(a_{n}) = l \cdot f(a_{n})$
\item $\varphi_{81} :\equiv \mathrm{Conv}(a,c) \to n \in \mathbb{N} \to \forall n \in \mathbb{N}\,((l f) \circ a(n) = (l f)(a_{n}))$
\item $\varphi_{82} :\equiv \forall w\,(l \cdot f(c) = w \to \mathrm{Conv}((l f) \circ a,l \cdot f(c)) \to \mathrm{Conv}((l f) \circ a,w))$
\item $\varphi_{83} :\equiv \forall \alpha\,(\mathrm{Conv}(f \circ a,\alpha) \to \mathrm{Seq}(f \circ a)) \to \mathrm{Conv}(f \circ a,f(c)) \to \mathrm{Seq}(f \circ a)$
\item $\varphi_{84} :\equiv \mathrm{Seq}(l \odot f \circ a) \land \forall n \in \mathbb{N}\,(l \odot f \circ a(n) = l \cdot f \circ a(n))$
\item $\varphi_{85} :\equiv \forall x\,(A \subseteq \mathrm{dom}(f) \to x \in A \to x \in \mathrm{dom}(f)) \to A \subseteq \mathrm{dom}(f) \to a_{n} \in A \to a_{n} \in \mathrm{dom}(f)$
\item $\varphi_{86} :\equiv \forall y\,(A \subseteq y \leftrightarrow \forall u \in A\,(u \in y)) \to (A \subseteq \mathrm{dom}((l f)) \leftrightarrow \forall u \in A\,(u \in \mathrm{dom}((l f))))$
\item $\varphi_{87} :\equiv l \cdot f \circ a(n) = l \cdot f(a_{n}) \to l \odot f \circ a(n) = l \cdot f(a_{n})$
\item $\varphi_{88} :\equiv l \cdot f \circ a(n) = l \cdot f \circ a(n) \to l \cdot f \circ a(n) = l \cdot f(a_{n})$
\item $\varphi_{89} :\equiv (l f)(a_{n}) = l \cdot f(a_{n}) \to (l f) \circ a(n) = l \cdot f(a_{n})$
\item $\varphi_{90} :\equiv \mathrm{Seq}((l f) \circ a) \land \forall n \in \mathbb{N}\,((l f) \circ a(n) = (l f)(a_{n}))$
\item $\varphi_{91} :\equiv l \odot f \circ a(n) = l \cdot f(a_{n}) \to l \cdot f(a_{n}) = l \odot f \circ a(n)$
\item $\varphi_{92} :\equiv \forall n \in \mathbb{N}\,(f \circ a(n) = f(a_{n})) \to n \in \mathbb{N} \to f \circ a(n) = f(a_{n})$
\item $\varphi_{93} :\equiv l \cdot f(a_{n}) = l \odot f \circ a(n) \to (l f) \circ a(n) = l \odot f \circ a(n)$
\item $\varphi_{94} :\equiv \forall n \in \mathbb{N}\,(a_{n} \in \mathrm{dom}(f)) \to \mathrm{Seq}(f \circ a) \land \forall n \in \mathbb{N}\,(f \circ a(n) = f(a_{n}))$
\item $\varphi_{95} :\equiv \forall f\,\forall A\,\forall c\,\forall a\,(\mathrm{ContIn}(f,A,c) \to \mathrm{Seq}(a) \to \varphi_{29})$
\item $\varphi_{96} :\equiv \forall z\,((l f) \circ a(n) = l \cdot f(a_{n}) \to l \cdot f(a_{n}) = z \to (l f) \circ a(n) = z)$
\item $\varphi_{97} :\equiv \forall w\,(f \circ a(n) = w \to l \cdot f \circ a(n) = l \cdot f \circ a(n) \to l \cdot f \circ a(n) = l \cdot w)$
\item $\varphi_{98} :\equiv \forall b\,\forall \alpha\,(\mathrm{Seq}(b) \to \forall n \in \mathbb{N}\,(l \odot f \circ a(n) = b_{n}) \to \mathrm{Conv}(l \odot f \circ a,\alpha) \to \mathrm{Conv}(b,\alpha))$
\item $\varphi_{99} :\equiv \forall y\,((l f)(c) = y \to y = (l f)(c)) \to (l f)(c) = l \cdot f(c) \to l \cdot f(c) = (l f)(c)$
\item $\varphi_{100} :\equiv (l f) \circ a(n) = l \odot f \circ a(n) \to l \odot f \circ a(n) = (l f) \circ a(n)$
\item $\varphi_{101} :\equiv \forall z\,(l \odot f \circ a(n) = l \cdot f \circ a(n) \to l \cdot f \circ a(n) = z \to l \odot f \circ a(n) = z)$
\item $\varphi_{102} :\equiv \forall x \in \mathrm{dom}((l f))\,((l f)(x) = l \cdot f(x)) \to c \in \mathrm{dom}((l f)) \to (l f)(c) = l \cdot f(c)$
\item $\varphi_{103} :\equiv (l \in \mathbb{R} \to \varphi_{3}) \to l \in \mathbb{R} \to \mathrm{ContIn}(f,A,c) \to A \subseteq \mathrm{dom}((l f))$
\item $\varphi_{104} :\equiv \forall z\,((l f) \circ a(n) = (l f)(a_{n}) \to (l f)(a_{n}) = z \to (l f) \circ a(n) = z)$
\item $\varphi_{105} :\equiv \mathrm{RFn}((l f)) \land \forall x\,(x \in \mathrm{dom}((l f)) \leftrightarrow x \in \mathrm{dom}(f)) \land \forall x \in \mathrm{dom}((l f))\,((l f)(x) = l \cdot f(x))$
\item $\varphi_{106} :\equiv \mathrm{ContIn}(f,A,c) \to \mathrm{Seq}(a) \to \forall n \in \mathbb{N}\,(a_{n} \in A) \to \varphi_{23}$
\item $\varphi_{107} :\equiv \mathrm{ContIn}(f,A,c) \to \mathrm{Seq}(a) \to \forall n \in \mathbb{N}\,(a_{n} \in A) \to \varphi_{24}$
\item $\varphi_{108} :\equiv \mathrm{ContIn}(f,A,c) \to \mathrm{Seq}(a) \to \forall n \in \mathbb{N}\,(a_{n} \in A) \to \varphi_{30}$
\item $\varphi_{109} :\equiv \mathrm{ContIn}(f,A,c) \to \mathrm{Seq}(a) \to \forall n \in \mathbb{N}\,(a_{n} \in A) \to \varphi_{32}$
\item $\varphi_{110} :\equiv \mathrm{ContIn}(f,A,c) \to \mathrm{Seq}(a) \to \forall n \in \mathbb{N}\,(a_{n} \in A) \to \varphi_{33}$
\item $\varphi_{111} :\equiv \forall a\,(l \in \mathbb{R} \land \mathrm{ContIn}(f,A,c) \to \mathrm{Seq}(a) \to \varphi_{75})$
\item $\varphi_{112} :\equiv l \in \mathbb{R} \to \mathrm{ContIn}(f,A,c) \to \forall a\,(\mathrm{Seq}(a) \to \varphi_{75})$
\item $\varphi_{113} :\equiv \mathrm{ContIn}(f,A,c) \to \mathrm{Seq}(a) \to \forall n \in \mathbb{N}\,(a_{n} \in A) \to \varphi_{35}$
\item $\varphi_{114} :\equiv \mathrm{ContIn}(f,A,c) \to \mathrm{Seq}(a) \to \forall n \in \mathbb{N}\,(a_{n} \in A) \to \varphi_{40}$
\item $\varphi_{115} :\equiv (l \in \mathbb{R} \to \mathrm{ContIn}(f,A,c) \to c \in A) \to l \in \mathbb{R} \to \mathrm{ContIn}(f,A,c) \to \varphi_{18}$
\item $\varphi_{116} :\equiv \forall a\,(l \in \mathbb{R} \land \mathrm{ContIn}(f,A,c)) \to \forall a\,(\mathrm{Seq}(a) \to \varphi_{75})$
\item $\varphi_{117} :\equiv \mathrm{ContIn}(f,A,c) \to \mathrm{Seq}(a) \to \forall n \in \mathbb{N}\,(a_{n} \in A) \to \varphi_{47}$
\item $\varphi_{118} :\equiv \forall n \in \mathbb{N}\,(l \odot f \circ a(n) = l \cdot f \circ a(n)) \to n \in \mathbb{N} \to l \odot f \circ a(n) = l \cdot f \circ a(n)$
\item $\varphi_{119} :\equiv \mathrm{ContIn}(f,A,c) \to \mathrm{Seq}(a) \to \forall n \in \mathbb{N}\,(a_{n} \in A) \to \varphi_{49}$
\item $\varphi_{120} :\equiv \forall n \in \mathbb{N}\,((l f) \circ a(n) = (l f)(a_{n})) \to n \in \mathbb{N} \to (l f) \circ a(n) = (l f)(a_{n})$
\item $\varphi_{121} :\equiv c \in A \to \forall a\,(\mathrm{Seq}(a) \to \varphi_{75}) \to \mathrm{ContIn}((l f),A,c)$
\item $\varphi_{122} :\equiv \forall x \in \mathrm{dom}((l f))\,((l f)(x) = l \cdot f(x)) \to a_{n} \in \mathrm{dom}((l f)) \to (l f)(a_{n}) = l \cdot f(a_{n})$
\item $\varphi_{123} :\equiv \mathrm{ContIn}(f,A,c) \to \mathrm{Seq}(a) \to \forall n \in \mathbb{N}\,(a_{n} \in A) \to \varphi_{55}$
\item $\varphi_{124} :\equiv \mathrm{ContIn}(f,A,c) \to \mathrm{Seq}(a) \to \forall n \in \mathbb{N}\,(a_{n} \in A) \to \varphi_{56}$
\item $\varphi_{125} :\equiv \mathrm{ContIn}(f,A,c) \to \mathrm{Seq}(a) \to \forall n \in \mathbb{N}\,(a_{n} \in A) \to \varphi_{60}$
\item $\varphi_{126} :\equiv \mathrm{ContIn}(f,A,c) \to \mathrm{Seq}(a) \to \forall n \in \mathbb{N}\,(a_{n} \in A) \to \varphi_{62}$
\item $\varphi_{127} :\equiv \forall f\,\forall A\,\forall c\,(\mathrm{RFn}(f) \to A \subseteq \mathrm{dom}(f) \to c \in A \to \forall a\,(\mathrm{Seq}(a) \to \varphi_{29}) \to \mathrm{ContIn}(f,A,c))$
\item $\varphi_{128} :\equiv (l \in \mathbb{R} \to \mathrm{ContIn}(f,A,c) \to A \subseteq \mathrm{dom}(f)) \to (l \in \mathbb{R} \to \mathrm{ContIn}(f,A,c) \to u \in A \to u \in A) \to \varphi_{4}$
\item $\varphi_{129} :\equiv \mathrm{ContIn}(f,A,c) \to \mathrm{Seq}(a) \to \forall n \in \mathbb{N}\,(a_{n} \in A) \to \varphi_{72}$
\item $\varphi_{130} :\equiv \mathrm{ContIn}(f,A,c) \to \mathrm{Seq}(a) \to \forall n \in \mathbb{N}\,(a_{n} \in A) \to \varphi_{74}$
\item $\varphi_{131} :\equiv \mathrm{ContIn}(f,A,c) \to \mathrm{Seq}(a) \to \forall n \in \mathbb{N}\,(a_{n} \in A) \to \varphi_{78}$
\item $\varphi_{132} :\equiv \mathrm{ContIn}(f,A,c) \to \mathrm{Seq}(a) \to \forall n \in \mathbb{N}\,(a_{n} \in A) \to \varphi_{81}$
\item $\varphi_{133} :\equiv \forall n\,(\varphi_{54}) \to \forall n \in \mathbb{N}\,(l \odot f \circ a(n) = (l f) \circ a(n))$
\item $\varphi_{134} :\equiv \mathrm{ContIn}(f,A,c) \to \mathrm{Seq}(a) \to \forall n \in \mathbb{N}\,(a_{n} \in A) \to \mathrm{Conv}(a,c) \to \varphi_{84}$
\item $\varphi_{135} :\equiv \mathrm{ContIn}(f,A,c) \to \mathrm{Seq}(a) \to \forall n \in \mathbb{N}\,(a_{n} \in A) \to \mathrm{Conv}(a,c) \to \varphi_{90}$
\item $\varphi_{136} :\equiv \forall \alpha\,(\mathrm{Seq}((l f) \circ a) \to \forall n \in \mathbb{N}\,(l \odot f \circ a(n) = (l f) \circ a(n)) \to \mathrm{Conv}(l \odot f \circ a,\alpha) \to \mathrm{Conv}((l f) \circ a,\alpha))$
\item $\varphi_{137} :\equiv \mathrm{RFn}((l f)) \to \mathrm{Seq}(a) \to \forall n \in \mathbb{N}\,(a_{n} \in \mathrm{dom}((l f))) \to \varphi_{90}$
\item $\varphi_{138} :\equiv (l \in \mathbb{R} \to \varphi_{6}) \to l \in \mathbb{R} \to \mathrm{ContIn}(f,A,c) \to \mathrm{Seq}(a) \to \varphi_{19}$
\item $\varphi_{139} :\equiv \mathrm{ContIn}(f,A,c) \leftrightarrow \mathrm{RFn}(f) \land A \subseteq \mathrm{dom}(f) \land c \in A \land \forall \varepsilon \in \mathbb{R}\,(0 < \varepsilon \to \exists \delta\,(\delta \in \mathbb{R} \land 0 < \delta \land \forall x \in A\,(\lvert x - c \rvert < \delta \to \lvert f(x) - f(c) \rvert < \varepsilon)))$
\item $\varphi_{140} :\equiv (l \in \mathbb{R} \to \mathrm{ContIn}(f,A,c) \to A \subseteq \mathrm{dom}(f)) \to \varphi_{115}$
\item $\varphi_{141} :\equiv (A \subseteq \mathrm{dom}((l f)) \leftrightarrow \forall u \in A\,(u \in \mathrm{dom}((l f)))) \to \varphi_{103}$
\item $\varphi_{142} :\equiv ((l f)(c) = l \cdot f(c) \to l \cdot f(c) = (l f)(c)) \to l \in \mathbb{R} \to \mathrm{ContIn}(f,A,c) \to \mathrm{Seq}(a) \to \varphi_{46}$
\item $\varphi_{143} :\equiv (l \in \mathbb{R} \to \mathrm{ContIn}(f,A,c) \to \mathrm{RFn}(f)) \to l \in \mathbb{R} \to \mathrm{ContIn}(f,A,c) \to \varphi_{105}$
\item $\varphi_{144} :\equiv \mathrm{Seq}((l f) \circ a) \to \forall n \in \mathbb{N}\,(l \odot f \circ a(n) = (l f) \circ a(n)) \to \varphi_{66}$
\item $\varphi_{145} :\equiv (l \in \mathbb{R} \to \mathrm{ContIn}(f,A,c) \to \varphi_{105}) \to l \in \mathbb{R} \to \mathrm{ContIn}(f,A,c) \to \mathrm{RFn}((l f))$
\item $\varphi_{146} :\equiv \forall a\,(\mathrm{ContIn}(f,A,c) \to \mathrm{Seq}(a) \to \varphi_{29}) \to \mathrm{ContIn}(f,A,c) \to \mathrm{Seq}(a) \to \varphi_{29}$
\item $\varphi_{147} :\equiv (l \in \mathbb{R} \to \mathrm{ContIn}(f,A,c) \to \mathrm{Seq}(a) \to \varphi_{8}) \to l \in \mathbb{R} \to \mathrm{ContIn}(f,A,c) \to \mathrm{Seq}(a) \to \varphi_{29}$
\item $\varphi_{148} :\equiv \forall c\,\forall a\,(\mathrm{ContIn}(f,A,c) \to \mathrm{Seq}(a) \to \varphi_{29}) \to \forall a\,(\mathrm{ContIn}(f,A,c) \to \mathrm{Seq}(a) \to \varphi_{29})$
\item $\varphi_{149} :\equiv (l \in \mathbb{R} \to \varphi_{20}) \to l \in \mathbb{R} \to \mathrm{ContIn}(f,A,c) \to \mathrm{Seq}(a) \to \varphi_{65}$
\item $\varphi_{150} :\equiv \forall A\,\forall c\,\forall a\,(\mathrm{ContIn}(f,A,c) \to \mathrm{Seq}(a) \to \varphi_{29}) \to \forall c\,\forall a\,(\mathrm{ContIn}(f,A,c) \to \mathrm{Seq}(a) \to \varphi_{29})$
\item $\varphi_{151} :\equiv (l \in \mathbb{R} \to \mathrm{ContIn}(f,A,c) \to \mathrm{Seq}(a) \to \varphi_{29}) \to l \in \mathbb{R} \to \mathrm{ContIn}(f,A,c) \to \varphi_{27}$
\item $\varphi_{152} :\equiv (l \in \mathbb{R} \to \mathrm{ContIn}(f,A,c) \to \mathrm{Seq}(a) \to \varphi_{17}) \to l \in \mathbb{R} \to \mathrm{ContIn}(f,A,c) \to \mathrm{Seq}(a) \to \varphi_{26}$
\item $\varphi_{153} :\equiv \varphi_{95} \to \forall A\,\forall c\,\forall a\,(\mathrm{ContIn}(f,A,c) \to \mathrm{Seq}(a) \to \varphi_{29})$
\item $\varphi_{154} :\equiv (l \in \mathbb{R} \to \mathrm{ContIn}(f,A,c) \to \mathrm{Seq}(a) \to \varphi_{19}) \to l \in \mathbb{R} \to \mathrm{ContIn}(f,A,c) \to \mathrm{Seq}(a) \to \varphi_{45}$
\item $\varphi_{155} :\equiv (l \in \mathbb{R} \to \mathrm{ContIn}(f,A,c) \to \mathrm{Seq}(a) \to \varphi_{26}) \to l \in \mathbb{R} \to \mathrm{ContIn}(f,A,c) \to \mathrm{Seq}(a) \to \varphi_{44}$
\item $\varphi_{156} :\equiv (l \in \mathbb{R} \to \mathrm{ContIn}(f,A,c) \to \mathrm{Seq}(a) \to \varphi_{44}) \to \varphi_{54} \to n \in \mathbb{N} \to a_{n} \in \mathrm{dom}(f)$
\item $\varphi_{157} :\equiv (l \in \mathbb{R} \to \mathrm{ContIn}(f,A,c) \to \mathrm{Seq}(a) \to \varphi_{29}) \to l \in \mathbb{R} \to \mathrm{ContIn}(f,A,c) \to \mathrm{Seq}(a) \to \varphi_{64}$
\item $\varphi_{158} :\equiv (l \in \mathbb{R} \to \mathrm{ContIn}(f,A,c) \to \mathrm{Seq}(a) \to \varphi_{17}) \to l \in \mathbb{R} \to \mathrm{ContIn}(f,A,c) \to \mathrm{Seq}(a) \to \varphi_{77}$
\item $\varphi_{159} :\equiv \forall n\,(\varphi_{54} \to n \in \mathbb{N} \to a_{n} \in \mathrm{dom}(f)) \to \forall n\,(\varphi_{54}) \to \forall n \in \mathbb{N}\,(a_{n} \in \mathrm{dom}(f))$
\item $\varphi_{160} :\equiv (l \in \mathbb{R} \to \mathrm{ContIn}(f,A,c) \to \mathrm{Seq}(a) \to \varphi_{17}) \to l \in \mathbb{R} \to \varphi_{107}$
\item $\varphi_{161} :\equiv (l \in \mathbb{R} \to \mathrm{ContIn}(f,A,c) \to \mathrm{Seq}(a) \to \varphi_{65}) \to \varphi_{54} \to n \in \mathbb{N} \to a_{n} \in \mathrm{dom}((l f))$
\item $\varphi_{162} :\equiv (l \in \mathbb{R} \to \mathrm{ContIn}(f,A,c) \to \mathrm{Seq}(a) \to \varphi_{64}) \to l \in \mathbb{R} \to \mathrm{ContIn}(f,A,c) \to \mathrm{Seq}(a) \to \varphi_{67}$
\item $\varphi_{163} :\equiv (l \in \mathbb{R} \to \mathrm{ContIn}(f,A,c) \to \mathrm{Seq}(a) \to \varphi_{17}) \to l \in \mathbb{R} \to \varphi_{117}$
\item $\varphi_{164} :\equiv (l \in \mathbb{R} \to \mathrm{ContIn}(f,A,c) \to \mathrm{Seq}(a) \to \varphi_{17}) \to l \in \mathbb{R} \to \varphi_{119}$
\item $\varphi_{165} :\equiv (l \cdot f(c) = (l f)(c) \to \varphi_{73}) \to l \in \mathbb{R} \to \mathrm{ContIn}(f,A,c) \to \mathrm{Seq}(a) \to \varphi_{75}$
\item $\varphi_{166} :\equiv \forall a\,(\mathrm{RFn}(f) \to \mathrm{Seq}(a) \to \varphi_{94}) \to \mathrm{RFn}(f) \to \mathrm{Seq}(a) \to \varphi_{94}$
\item $\varphi_{167} :\equiv \forall n\,(\varphi_{54} \to n \in \mathbb{N} \to a_{n} \in \mathrm{dom}((l f))) \to \forall n\,(\varphi_{54}) \to \forall n \in \mathbb{N}\,(a_{n} \in \mathrm{dom}((l f)))$
\item $\varphi_{168} :\equiv (l \in \mathbb{R} \to \mathrm{ContIn}(f,A,c) \to \mathrm{Seq}(a) \to \varphi_{75}) \to l \in \mathbb{R} \land \mathrm{ContIn}(f,A,c) \to \mathrm{Seq}(a) \to \varphi_{75}$
\item $\varphi_{169} :\equiv (l \in \mathbb{R} \to \varphi_{69}) \to (\forall n \in \mathbb{N}\,(a_{n} \in A) \to n \in \mathbb{N} \to a_{n} \in A) \to l \in \mathbb{R} \to \varphi_{61}$
\item $\varphi_{170} :\equiv \forall f\,\forall a\,(\mathrm{RFn}(f) \to \mathrm{Seq}(a) \to \varphi_{94}) \to \forall a\,(\mathrm{RFn}(f) \to \mathrm{Seq}(a) \to \varphi_{94})$
\item $\varphi_{171} :\equiv \varphi_{1} \to \varphi_{139} \to l \in \mathbb{R} \to \mathrm{ContIn}(f,A,c) \to A \subseteq \mathrm{dom}(f)$
\item $\varphi_{172} :\equiv (l \in \mathbb{R} \to \mathrm{ContIn}(f,A,c) \to A \subseteq \mathrm{dom}(f)) \to \varphi_{155}$
\item $\varphi_{173} :\equiv (l \in \mathbb{R} \to \mathrm{ContIn}(f,A,c) \to \mathrm{Seq}(a) \to \varphi_{53}) \to l \in \mathbb{R} \to \varphi_{125}$
\item $\varphi_{174} :\equiv (l \in \mathbb{R} \to \mathrm{ContIn}(f,A,c) \to A \subseteq \mathrm{dom}((l f))) \to (l \in \mathbb{R} \to \mathrm{ContIn}(f,A,c) \to c \in A) \to \varphi_{112} \to \varphi_{5}$
\item $\varphi_{175} :\equiv (\varphi_{54} \to \forall n\,(\varphi_{54})) \to l \in \mathbb{R} \to \mathrm{ContIn}(f,A,c) \to \mathrm{Seq}(a) \to \varphi_{53}$
\item $\varphi_{176} :\equiv \forall f\,(l \in \mathbb{R} \to \mathrm{RFn}(f) \to \varphi_{105}) \to l \in \mathbb{R} \to \mathrm{RFn}(f) \to \varphi_{105}$
\item $\varphi_{177} :\equiv (l \in \mathbb{R} \to \varphi_{135}) \to l \in \mathbb{R} \to \mathrm{ContIn}(f,A,c) \to \mathrm{Seq}(a) \to \varphi_{28}$
\item $\varphi_{178} :\equiv \forall l\,\forall f\,(l \in \mathbb{R} \to \mathrm{RFn}(f) \to \varphi_{105}) \to \forall f\,(l \in \mathbb{R} \to \mathrm{RFn}(f) \to \varphi_{105})$
\item $\varphi_{179} :\equiv (\varphi_{54} \to \forall n\,(\varphi_{54})) \to l \in \mathbb{R} \to \mathrm{ContIn}(f,A,c) \to \mathrm{Seq}(a) \to \varphi_{71}$
\item $\varphi_{180} :\equiv (l \in \mathbb{R} \to \mathrm{ContIn}(f,A,c) \to \mathrm{Seq}(a) \to \varphi_{45}) \to \varphi_{142}$
\item $\varphi_{181} :\equiv (l \in \mathbb{R} \to \varphi_{124}) \to \varphi_{54} \to n \in \mathbb{N} \to l \odot f \circ a(n) = (l f) \circ a(n)$
\item $\varphi_{182} :\equiv (l \in \mathbb{R} \to l \in \mathbb{R}) \to (l \in \mathbb{R} \to \mathrm{ContIn}(f,A,c) \to \varphi_{27}) \to l \in \mathbb{R} \to \varphi_{134}$
\item $\varphi_{183} :\equiv (l \in \mathbb{R} \to \mathrm{ContIn}(f,A,c) \to \mathrm{Seq}(a) \to \varphi_{71}) \to l \in \mathbb{R} \to \varphi_{135}$
\item $\varphi_{184} :\equiv \forall n\,(\varphi_{54} \to n \in \mathbb{N} \to l \odot f \circ a(n) = (l f) \circ a(n)) \to \varphi_{133}$
\item $\varphi_{185} :\equiv (l \in \mathbb{R} \to \mathrm{ContIn}(f,A,c) \to \mathrm{Seq}(a) \to \varphi_{44}) \to \varphi_{149}$
\item $\varphi_{186} :\equiv (l \in \mathbb{R} \to \mathrm{ContIn}(f,A,c) \to \mathrm{RFn}((l f))) \to \varphi_{174}$
\item $\varphi_{187} :\equiv l \cdot f \circ a(n) = l \cdot f \circ a(n) \to (f \circ a(n) = f(a_{n}) \to \varphi_{88}) \to l \in \mathbb{R} \to \varphi_{106}$
\item $\varphi_{188} :\equiv \varphi_{127} \to \forall A\,\forall c\,(\mathrm{RFn}((l f)) \to A \subseteq \mathrm{dom}((l f)) \to \varphi_{121})$
\item $\varphi_{189} :\equiv (l \in \mathbb{R} \to \mathrm{ContIn}(f,A,c) \to \mathrm{Seq}(a) \to \varphi_{71}) \to \varphi_{68} \to l \in \mathbb{R} \to \mathrm{ContIn}(f,A,c) \to \mathrm{Seq}(a) \to \varphi_{65}$
\item $\varphi_{190} :\equiv (l \in \mathbb{R} \to \mathrm{ContIn}(f,A,c) \to \mathrm{RFn}(f)) \to \varphi_{2} \to \varphi_{173}$
\item $\varphi_{191} :\equiv \forall c\,(\mathrm{RFn}((l f)) \to A \subseteq \mathrm{dom}((l f)) \to \varphi_{121}) \to \mathrm{RFn}((l f)) \to A \subseteq \mathrm{dom}((l f)) \to \varphi_{121}$
\item $\varphi_{192} :\equiv (l \in \mathbb{R} \to \mathrm{ContIn}(f,A,c) \to \mathrm{Seq}(a) \to \varphi_{65}) \to \varphi_{160}$
\item $\varphi_{193} :\equiv \forall A\,\forall c\,(\mathrm{RFn}((l f)) \to A \subseteq \mathrm{dom}((l f)) \to \varphi_{121}) \to \forall c\,(\mathrm{RFn}((l f)) \to A \subseteq \mathrm{dom}((l f)) \to \varphi_{121})$
\item $\varphi_{194} :\equiv (l \in \mathbb{R} \to \mathrm{ContIn}(f,A,c) \to \mathrm{Seq}(a) \to \varphi_{67}) \to \varphi_{165}$
\item $\varphi_{195} :\equiv (l \in \mathbb{R} \to \mathrm{ContIn}(f,A,c) \to \mathrm{RFn}((l f))) \to \varphi_{2} \to \varphi_{183}$
\item $\varphi_{196} :\equiv (l \in \mathbb{R} \to \varphi_{106}) \to (l \odot f \circ a(n) = l \cdot f \circ a(n) \to \varphi_{87}) \to l \in \mathbb{R} \to \varphi_{109}$
\item $\varphi_{197} :\equiv (l \in \mathbb{R} \to \varphi_{107}) \to ((l f) \circ a(n) = (l f)(a_{n}) \to \varphi_{89}) \to l \in \mathbb{R} \to \varphi_{113}$
\item $\varphi_{198} :\equiv (l \in \mathbb{R} \to \varphi_{110}) \to ((l f) \circ a(n) = l \cdot f(a_{n}) \to \varphi_{93}) \to l \in \mathbb{R} \to \varphi_{123}$
\item $\varphi_{199} :\equiv (l \in \mathbb{R} \to \mathrm{ContIn}(f,A,c) \to \mathrm{Seq}(a) \to \varphi_{77}) \to \varphi_{187}$
\item $\varphi_{200} :\equiv (l \in \mathbb{R} \to \mathrm{ContIn}(f,A,c) \to \mathrm{Seq}(a) \to \varphi_{46}) \to \varphi_{194}$
\item $\varphi_{201} :\equiv (\mathrm{ContIn}(f,A,c) \to \mathrm{Seq}(a) \to \varphi_{29}) \to \varphi_{1} \to \varphi_{2} \to (l \in \mathbb{R} \to \varphi_{69}) \to \varphi_{147}$
\item $\varphi_{202} :\equiv (l \in \mathbb{R} \to \mathrm{ContIn}(f,A,c) \to \mathrm{Seq}(a) \to \varphi_{28}) \to (l \in \mathbb{R} \to \varphi_{126}) \to \varphi_{162}$
\item $\varphi_{203} :\equiv \forall n\,(\varphi_{54} \to n \in \mathbb{N} \to l \odot f \circ a(n) = (l f) \circ a(n)) \to \varphi_{184} \to (\varphi_{54} \to \forall n\,(\varphi_{54})) \to l \in \mathbb{R} \to \varphi_{126}$
\end{abbrevlist}
\end{fullproof}

% 自動生成 (prover/build.js): ContInAbs  ノード数 394
\begin{fullproof}{定理2.3(絶対値)}{$\forall f\,\forall A\,\forall c\,(\mathrm{ContIn}(f,A,c) \to \mathrm{ContIn}(\lvert f \rvert,A,c))$}{394}\phantomsection\label{tree:ContInAbs}
\noindent$\Pi_{1}$:
\begin{ptree}
\AXC{$\varphi_{97}$\lbl{[G1]}}
\AXC{$\varphi_{149}$\lbl{A2}}
\BIC{$\forall A\,\forall c\,\forall a\,(\mathrm{ContIn}(f,A,c) \to \mathrm{Seq}(a) \to \varphi_{30})$}
\AXC{$\varphi_{145}$\lbl{A2}}
\BIC{$\forall c\,\forall a\,(\mathrm{ContIn}(f,A,c) \to \mathrm{Seq}(a) \to \varphi_{30})$}
\AXC{$\varphi_{142}$\lbl{A2}}
\BIC{$\forall a\,(\mathrm{ContIn}(f,A,c) \to \mathrm{Seq}(a) \to \varphi_{30})$}
\AXC{$\varphi_{139}$\lbl{A2}}
\BIC{$\mathrm{ContIn}(f,A,c) \to \mathrm{Seq}(a) \to \varphi_{30}$}
\end{ptree}
\noindent$\Pi_{2}$:
\begin{ptree}
\AXC{$\forall a\,\forall \alpha\,(\mathrm{Conv}(a,\alpha) \to \mathrm{Conv}(\lvert a \rvert,\lvert \alpha \rvert))$\lbl{定理2.15(1)}}
\AXC{$\varphi_{89}$\lbl{A2}}
\BIC{$\varphi_{10}$}
\AXC{$\varphi_{10} \to \varphi_{19}$\lbl{A2}}
\BIC{$\varphi_{19}$}
\AXC{$\varphi_{19} \to \varphi_{143}$\lbl{A1}}
\BIC{$\varphi_{143}$}
\end{ptree}
\noindent$\Pi_{3}$:
\begin{ptree}
\AXC{$x = x$\lbl{A5}}
\AXC{$x = y \to x = x \to y = x$\lbl{A6}}
\AXC{$(x = y \to x = x \to y = x) \to x = x \to x = y \to y = x$\lbl{A1}}
\BIC{$x = x \to x = y \to y = x$}
\BIC{$x = y \to y = x$}
\UIC{$\forall y\,(x = y \to y = x)$}
\UIC{$\forall x\,\forall y\,(x = y \to y = x)$}
\end{ptree}
\noindent$\Pi_{4}$:
\begin{ptree}
\AXC{$\stackrel{\Pi_{3}}{\vdots}$}\noLine
\UIC{$\forall x\,\forall y\,(x = y \to y = x)$}
\AXC{$\varphi_{48}$\lbl{A2}}
\BIC{$\forall y\,(\lvert f \rvert \circ a(n) = y \to y = \lvert f \rvert \circ a(n))$}
\AXC{$\forall y\,(\lvert f \rvert \circ a(n) = y \to y = \lvert f \rvert \circ a(n)) \to \varphi_{92}$\lbl{A2}}
\BIC{$\varphi_{92}$}
\end{ptree}
\noindent$\Pi_{5}$:
\begin{ptree}
\AXC{$y = z \to x = y \to x = z$\lbl{A6}}
\AXC{$(y = z \to x = y \to x = z) \to x = y \to y = z \to x = z$\lbl{A1}}
\BIC{$x = y \to y = z \to x = z$}
\UIC{$\forall z\,(x = y \to y = z \to x = z)$}
\UIC{$\forall y\,\forall z\,(x = y \to y = z \to x = z)$}
\UIC{$\forall x\,\forall y\,\forall z\,(x = y \to y = z \to x = z)$}
\AXC{$\forall x\,\forall y\,\forall z\,(x = y \to y = z \to x = z) \to \varphi_{21}$\lbl{A2}}
\BIC{$\varphi_{21}$}
\end{ptree}
\noindent$\Pi_{6}$:
\begin{ptree}
\AXC{$\stackrel{\Pi_{3}}{\vdots}$}\noLine
\UIC{$\forall x\,\forall y\,(x = y \to y = x)$}
\AXC{$\varphi_{49}$\lbl{A2}}
\BIC{$\forall y\,(\lvert f \circ a \rvert(n) = y \to y = \lvert f \circ a \rvert(n))$}
\AXC{$\forall y\,(\lvert f \circ a \rvert(n) = y \to y = \lvert f \circ a \rvert(n)) \to \varphi_{81}$\lbl{A2}}
\BIC{$\varphi_{81}$}
\end{ptree}
\noindent$\Pi_{7}$:
\begin{ptree}
\AXC{$y = z \to x = y \to x = z$\lbl{A6}}
\AXC{$(y = z \to x = y \to x = z) \to x = y \to y = z \to x = z$\lbl{A1}}
\BIC{$x = y \to y = z \to x = z$}
\UIC{$\forall z\,(x = y \to y = z \to x = z)$}
\UIC{$\forall y\,\forall z\,(x = y \to y = z \to x = z)$}
\UIC{$\forall x\,\forall y\,\forall z\,(x = y \to y = z \to x = z)$}
\AXC{$\forall x\,\forall y\,\forall z\,(x = y \to y = z \to x = z) \to \varphi_{22}$\lbl{A2}}
\BIC{$\varphi_{22}$}
\end{ptree}
\noindent$\Pi_{8}$:
\begin{ptree}
\AXC{$z = w \to \lvert f \circ a(n) \rvert = \lvert z \rvert \to \lvert f \circ a(n) \rvert = \lvert w \rvert$\lbl{A6}}
\UIC{$\varphi_{37}$}
\UIC{$\forall z\,(\varphi_{37})$}
\AXC{$\forall z\,(\varphi_{37}) \to \varphi_{99}$\lbl{A2}}
\BIC{$\varphi_{99}$}
\AXC{$\varphi_{99} \to f \circ a(n) = f(a_{n}) \to \varphi_{84}$\lbl{A2}}
\BIC{$f \circ a(n) = f(a_{n}) \to \varphi_{84}$}
\end{ptree}
\noindent$\Pi_{9}$:
\begin{ptree}
\AXC{$\forall n \in \mathbb{N}\,(a_{n} \in A) \to n \in \mathbb{N} \to a_{n} \in A$\lbl{A2}}
\AXC{$\varphi_{68}$\lbl{A1}}
\AXC{$\varphi_{68} \to (\forall n \in \mathbb{N}\,(a_{n} \in A) \to n \in \mathbb{N} \to a_{n} \in A) \to \varphi_{63}$\lbl{A1}}
\BIC{$(\forall n \in \mathbb{N}\,(a_{n} \in A) \to n \in \mathbb{N} \to a_{n} \in A) \to \varphi_{63}$}
\BIC{$\varphi_{63}$}
\AXC{$\varphi_{63} \to \varphi_{138}$\lbl{A1}}
\BIC{$\varphi_{138}$}
\end{ptree}
\noindent$\Pi_{10}$:
\begin{ptree}
\AXC{$\forall f\,\forall A\,\forall c\,(\varphi_{134})$\lbl{DefContIn}}
\AXC{$\forall f\,\forall A\,\forall c\,(\varphi_{134}) \to \forall A\,\forall c\,(\varphi_{134})$\lbl{A2}}
\BIC{$\forall A\,\forall c\,(\varphi_{134})$}
\AXC{$\forall A\,\forall c\,(\varphi_{134}) \to \forall c\,(\varphi_{134})$\lbl{A2}}
\BIC{$\forall c\,(\varphi_{134})$}
\AXC{$\forall c\,(\varphi_{134}) \to \varphi_{134}$\lbl{A2}}
\BIC{$\varphi_{134}$}
\end{ptree}
\noindent$\Pi_{11}$:
\begin{ptree}
\AXC{$\stackrel{\Pi_{10}}{\vdots}$}\noLine
\UIC{$\varphi_{134}$}
\AXC{$\mathrm{ContIn}(f,A,c) \to \mathrm{ContIn}(f,A,c)$\lbl{A1}}
\AXC{$\varphi_{158}$\lbl{A1}}
\BIC{$\varphi_{134} \to \mathrm{ContIn}(f,A,c) \to A \subseteq \mathrm{dom}(f)$}
\BIC{$\mathrm{ContIn}(f,A,c) \to A \subseteq \mathrm{dom}(f)$}
\end{ptree}
\noindent$\Pi_{12}$:
\begin{ptree}
\AXC{$\forall X\,\forall Y\,\forall x\,(X \subseteq Y \to x \in X \to x \in Y)$\lbl{[B5]}}
\AXC{$\varphi_{38}$\lbl{A2}}
\BIC{$\forall Y\,\forall x\,(A \subseteq Y \to x \in A \to x \in Y)$}
\AXC{$\varphi_{50}$\lbl{A2}}
\BIC{$\forall x\,(A \subseteq \mathrm{dom}(f) \to x \in A \to x \in \mathrm{dom}(f))$}
\AXC{$\varphi_{83}$\lbl{A2}}
\BIC{$A \subseteq \mathrm{dom}(f) \to a_{n} \in A \to a_{n} \in \mathrm{dom}(f)$}
\AXC{$\varphi_{166}$\lbl{A1}}
\BIC{$(\mathrm{ContIn}(f,A,c) \to A \subseteq \mathrm{dom}(f)) \to \varphi_{140}$}
\end{ptree}
\noindent$\Pi_{13}$:
\begin{ptree}
\AXC{$\stackrel{\Pi_{11}}{\vdots}$}\noLine
\UIC{$\mathrm{ContIn}(f,A,c) \to A \subseteq \mathrm{dom}(f)$}
\AXC{$\stackrel{\Pi_{12}}{\vdots}$}\noLine
\UIC{$(\mathrm{ContIn}(f,A,c) \to A \subseteq \mathrm{dom}(f)) \to \varphi_{140}$}
\BIC{$\varphi_{140}$}
\end{ptree}
\noindent$\Pi_{14}$:
\begin{ptree}
\AXC{$\mathrm{ContIn}(f,A,c) \to \mathrm{Seq}(a) \to \varphi_{15}$\lbl{A1}}
\AXC{$\stackrel{\Pi_{9}}{\vdots}$}\noLine
\UIC{$\varphi_{138}$}
\BIC{$\mathrm{ContIn}(f,A,c) \to \mathrm{Seq}(a) \to \varphi_{25}$}
\AXC{$\stackrel{\Pi_{13}}{\vdots}$}\noLine
\UIC{$\varphi_{140}$}
\BIC{$\mathrm{ContIn}(f,A,c) \to \mathrm{Seq}(a) \to \varphi_{44}$}
\AXC{$\varphi_{141}$\lbl{A1}}
\BIC{$\varphi_{17} \to n \in \mathbb{N} \to a_{n} \in \mathrm{dom}(f)$}
\UIC{$\forall n\,(\varphi_{17} \to n \in \mathbb{N} \to a_{n} \in \mathrm{dom}(f))$}
\end{ptree}
\noindent$\Pi_{15}$:
\begin{ptree}
\AXC{$\varphi_{17} \to \forall n\,(\varphi_{17})$\lbl{A4}}
\AXC{$\varphi_{146}$\lbl{A3}}
\AXC{$\stackrel{\Pi_{14}}{\vdots}$}\noLine
\UIC{$\forall n\,(\varphi_{17} \to n \in \mathbb{N} \to a_{n} \in \mathrm{dom}(f))$}
\AXC{$\forall n\,(\varphi_{17} \to n \in \mathbb{N} \to a_{n} \in \mathrm{dom}(f)) \to \varphi_{146} \to \varphi_{159}$\lbl{A1}}
\BIC{$\varphi_{146} \to \varphi_{159}$}
\BIC{$\varphi_{159}$}
\BIC{$\mathrm{ContIn}(f,A,c) \to \mathrm{Seq}(a) \to \varphi_{51}$}
\end{ptree}
\noindent$\Pi_{16}$:
\begin{ptree}
\AXC{$\forall f\,\forall A\,\forall c\,(\varphi_{134})$\lbl{DefContIn}}
\AXC{$\forall f\,\forall A\,\forall c\,(\varphi_{134}) \to \forall A\,\forall c\,(\varphi_{134})$\lbl{A2}}
\BIC{$\forall A\,\forall c\,(\varphi_{134})$}
\AXC{$\forall A\,\forall c\,(\varphi_{134}) \to \forall c\,(\varphi_{134})$\lbl{A2}}
\BIC{$\forall c\,(\varphi_{134})$}
\AXC{$\forall c\,(\varphi_{134}) \to \varphi_{134}$\lbl{A2}}
\BIC{$\varphi_{134}$}
\end{ptree}
\noindent$\Pi_{17}$:
\begin{ptree}
\AXC{$\stackrel{\Pi_{16}}{\vdots}$}\noLine
\UIC{$\varphi_{134}$}
\AXC{$\mathrm{ContIn}(f,A,c) \to \mathrm{ContIn}(f,A,c)$\lbl{A1}}
\AXC{$\varphi_{156}$\lbl{A1}}
\BIC{$\varphi_{134} \to \mathrm{ContIn}(f,A,c) \to \mathrm{RFn}(f)$}
\BIC{$\mathrm{ContIn}(f,A,c) \to \mathrm{RFn}(f)$}
\end{ptree}
\noindent$\Pi_{18}$:
\begin{ptree}
\AXC{$\stackrel{\Pi_{17}}{\vdots}$}\noLine
\UIC{$\mathrm{ContIn}(f,A,c) \to \mathrm{RFn}(f)$}
\AXC{$\forall f\,\forall a\,(\mathrm{RFn}(f) \to \mathrm{Seq}(a) \to \varphi_{96})$\lbl{[S8]}}
\AXC{$\varphi_{162}$\lbl{A2}}
\BIC{$\forall a\,(\mathrm{RFn}(f) \to \mathrm{Seq}(a) \to \varphi_{96})$}
\AXC{$\varphi_{161}$\lbl{A2}}
\BIC{$\mathrm{RFn}(f) \to \mathrm{Seq}(a) \to \varphi_{96}$}
\AXC{$\varphi_{175}$\lbl{A1}}
\BIC{$(\mathrm{ContIn}(f,A,c) \to \mathrm{RFn}(f)) \to \varphi_{165}$}
\BIC{$\varphi_{165}$}
\end{ptree}
\noindent$\Pi_{19}$:
\begin{ptree}
\AXC{$\stackrel{\Pi_{15}}{\vdots}$}\noLine
\UIC{$\mathrm{ContIn}(f,A,c) \to \mathrm{Seq}(a) \to \varphi_{51}$}
\AXC{$\mathrm{ContIn}(f,A,c) \to \mathrm{Seq}(a) \to \mathrm{Seq}(a)$\lbl{A1}}
\AXC{$\stackrel{\Pi_{18}}{\vdots}$}\noLine
\UIC{$\varphi_{165}$}
\BIC{$(\mathrm{ContIn}(f,A,c) \to \mathrm{Seq}(a) \to \varphi_{51}) \to \varphi_{117}$}
\BIC{$\varphi_{117}$}
\AXC{$\varphi_{117} \to \varphi_{111}$\lbl{A1}}
\BIC{$\varphi_{111}$}
\end{ptree}
\noindent$\Pi_{20}$:
\begin{ptree}
\AXC{$\mathrm{ContIn}(f,A,c) \to \mathrm{Seq}(a) \to \varphi_{15}$\lbl{A1}}
\AXC{$\varphi_{86}$\lbl{A2}}
\AXC{$\stackrel{\Pi_{19}}{\vdots}$}\noLine
\UIC{$\varphi_{111}$}
\AXC{$\varphi_{111} \to \varphi_{86} \to \varphi_{108}$\lbl{A1}}
\BIC{$\varphi_{86} \to \varphi_{108}$}
\BIC{$\varphi_{108}$}
\AXC{$\varphi_{108} \to \varphi_{144}$\lbl{A1}}
\BIC{$\varphi_{144}$}
\BIC{$\mathrm{ContIn}(f,A,c) \to \mathrm{Seq}(a) \to \varphi_{77}$}
\end{ptree}
\noindent$\Pi_{21}$:
\begin{ptree}
\AXC{$\stackrel{\Pi_{20}}{\vdots}$}\noLine
\UIC{$\mathrm{ContIn}(f,A,c) \to \mathrm{Seq}(a) \to \varphi_{77}$}
\AXC{$(\mathrm{ContIn}(f,A,c) \to \mathrm{Seq}(a) \to \varphi_{77}) \to \varphi_{169}$\lbl{A1}}
\BIC{$\varphi_{169}$}
\end{ptree}
\noindent$\Pi_{22}$:
\begin{ptree}
\AXC{$\stackrel{\Pi_{8}}{\vdots}$}\noLine
\UIC{$f \circ a(n) = f(a_{n}) \to \varphi_{84}$}
\AXC{$x = x$\lbl{A5}}
\UIC{$\forall x\,(x = x)$}
\AXC{$\forall x\,(x = x) \to \lvert f \circ a(n) \rvert = \lvert f \circ a(n) \rvert$\lbl{A2}}
\BIC{$\lvert f \circ a(n) \rvert = \lvert f \circ a(n) \rvert$}
\AXC{$\stackrel{\Pi_{21}}{\vdots}$}\noLine
\UIC{$\varphi_{169}$}
\BIC{$(f \circ a(n) = f(a_{n}) \to \varphi_{84}) \to \varphi_{103}$}
\BIC{$\varphi_{103}$}
\end{ptree}
\noindent$\Pi_{23}$:
\begin{ptree}
\AXC{$\forall a\,\forall \alpha\,(\mathrm{Conv}(a,\alpha) \to \mathrm{Seq}(a))$\lbl{[C1]}}
\AXC{$\varphi_{39}$\lbl{A2}}
\BIC{$\forall \alpha\,(\mathrm{Conv}(f \circ a,\alpha) \to \mathrm{Seq}(f \circ a))$}
\AXC{$\varphi_{82}$\lbl{A2}}
\BIC{$\mathrm{Conv}(f \circ a,f(c)) \to \mathrm{Seq}(f \circ a)$}
\AXC{$(\mathrm{Conv}(f \circ a,f(c)) \to \mathrm{Seq}(f \circ a)) \to \varphi_{136}$\lbl{A1}}
\BIC{$\varphi_{136}$}
\end{ptree}
\noindent$\Pi_{24}$:
\begin{ptree}
\AXC{$\stackrel{\Pi_{1}}{\vdots}$}\noLine
\UIC{$\mathrm{ContIn}(f,A,c) \to \mathrm{Seq}(a) \to \varphi_{30}$}
\AXC{$\stackrel{\Pi_{23}}{\vdots}$}\noLine
\UIC{$\varphi_{136}$}
\BIC{$\mathrm{ContIn}(f,A,c) \to \varphi_{26}$}
\AXC{$\varphi_{24}$\lbl{[S11]}}
\AXC{$\varphi_{24} \to \mathrm{Seq}(f \circ a) \to \varphi_{78}$\lbl{A2}}
\BIC{$\mathrm{Seq}(f \circ a) \to \varphi_{78}$}
\AXC{$\varphi_{172}$\lbl{A1}}
\BIC{$(\mathrm{ContIn}(f,A,c) \to \varphi_{26}) \to \varphi_{132}$}
\BIC{$\varphi_{132}$}
\end{ptree}
\noindent$\Pi_{25}$:
\begin{ptree}
\AXC{$\varphi_{110}$\lbl{A2}}
\AXC{$\stackrel{\Pi_{24}}{\vdots}$}\noLine
\UIC{$\varphi_{132}$}
\AXC{$\varphi_{132} \to \varphi_{123}$\lbl{A1}}
\BIC{$\varphi_{123}$}
\AXC{$\varphi_{123} \to \varphi_{110} \to \varphi_{120}$\lbl{A1}}
\BIC{$\varphi_{110} \to \varphi_{120}$}
\BIC{$\varphi_{120}$}
\AXC{$\varphi_{120} \to (\mathrm{ContIn}(f,A,c) \to \mathrm{Seq}(a) \to \varphi_{15}) \to \varphi_{112}$\lbl{A1}}
\BIC{$(\mathrm{ContIn}(f,A,c) \to \mathrm{Seq}(a) \to \varphi_{15}) \to \varphi_{112}$}
\end{ptree}
\noindent$\Pi_{26}$:
\begin{ptree}
\AXC{$\mathrm{ContIn}(f,A,c) \to \mathrm{Seq}(a) \to \varphi_{15}$\lbl{A1}}
\AXC{$\stackrel{\Pi_{25}}{\vdots}$}\noLine
\UIC{$(\mathrm{ContIn}(f,A,c) \to \mathrm{Seq}(a) \to \varphi_{15}) \to \varphi_{112}$}
\BIC{$\varphi_{112}$}
\end{ptree}
\noindent$\Pi_{27}$:
\begin{ptree}
\AXC{$\stackrel{\Pi_{22}}{\vdots}$}\noLine
\UIC{$\varphi_{103}$}
\AXC{$\stackrel{\Pi_{26}}{\vdots}$}\noLine
\UIC{$\varphi_{112}$}
\AXC{$\varphi_{112} \to \varphi_{103} \to (\lvert f \circ a \rvert(n) = \lvert f \circ a(n) \rvert \to \varphi_{80}) \to \varphi_{104}$\lbl{A1}}
\BIC{$\varphi_{103} \to (\lvert f \circ a \rvert(n) = \lvert f \circ a(n) \rvert \to \varphi_{80}) \to \varphi_{104}$}
\BIC{$(\lvert f \circ a \rvert(n) = \lvert f \circ a(n) \rvert \to \varphi_{80}) \to \varphi_{104}$}
\end{ptree}
\noindent$\Pi_{28}$:
\begin{ptree}
\AXC{$\stackrel{\Pi_{7}}{\vdots}$}\noLine
\UIC{$\varphi_{22}$}
\AXC{$\varphi_{22} \to \varphi_{100}$\lbl{A2}}
\BIC{$\varphi_{100}$}
\AXC{$\varphi_{100} \to \lvert f \circ a \rvert(n) = \lvert f \circ a(n) \rvert \to \varphi_{80}$\lbl{A2}}
\BIC{$\lvert f \circ a \rvert(n) = \lvert f \circ a(n) \rvert \to \varphi_{80}$}
\AXC{$\stackrel{\Pi_{27}}{\vdots}$}\noLine
\UIC{$(\lvert f \circ a \rvert(n) = \lvert f \circ a(n) \rvert \to \varphi_{80}) \to \varphi_{104}$}
\BIC{$\varphi_{104}$}
\end{ptree}
\noindent$\Pi_{29}$:
\begin{ptree}
\AXC{$y = z \to x = y \to x = z$\lbl{A6}}
\AXC{$(y = z \to x = y \to x = z) \to x = y \to y = z \to x = z$\lbl{A1}}
\BIC{$x = y \to y = z \to x = z$}
\UIC{$\forall z\,(x = y \to y = z \to x = z)$}
\UIC{$\forall y\,\forall z\,(x = y \to y = z \to x = z)$}
\UIC{$\forall x\,\forall y\,\forall z\,(x = y \to y = z \to x = z)$}
\AXC{$\forall x\,\forall y\,\forall z\,(x = y \to y = z \to x = z) \to \varphi_{21}$\lbl{A2}}
\BIC{$\varphi_{21}$}
\end{ptree}
\noindent$\Pi_{30}$:
\begin{ptree}
\AXC{$\stackrel{\Pi_{29}}{\vdots}$}\noLine
\UIC{$\varphi_{21}$}
\AXC{$\varphi_{21} \to \varphi_{93}$\lbl{A2}}
\BIC{$\varphi_{93}$}
\AXC{$\varphi_{93} \to \lvert f \rvert \circ a(n) = \lvert f \rvert(a_{n}) \to \varphi_{79}$\lbl{A2}}
\BIC{$\lvert f \rvert \circ a(n) = \lvert f \rvert(a_{n}) \to \varphi_{79}$}
\end{ptree}
\noindent$\Pi_{31}$:
\begin{ptree}
\AXC{$\forall f\,(\mathrm{RFn}(f) \to \varphi_{94})$\lbl{[S21]}}
\AXC{$\forall f\,(\mathrm{RFn}(f) \to \varphi_{94}) \to \mathrm{RFn}(f) \to \varphi_{94}$\lbl{A2}}
\BIC{$\mathrm{RFn}(f) \to \varphi_{94}$}
\AXC{$(\mathrm{RFn}(f) \to \varphi_{94}) \to \varphi_{124}$\lbl{A1}}
\BIC{$\varphi_{124}$}
\end{ptree}
\noindent$\Pi_{32}$:
\begin{ptree}
\AXC{$\stackrel{\Pi_{17}}{\vdots}$}\noLine
\UIC{$\mathrm{ContIn}(f,A,c) \to \mathrm{RFn}(f)$}
\AXC{$\stackrel{\Pi_{31}}{\vdots}$}\noLine
\UIC{$\varphi_{124}$}
\BIC{$\mathrm{ContIn}(f,A,c) \to \varphi_{94}$}
\end{ptree}
\noindent$\Pi_{33}$:
\begin{ptree}
\AXC{$\stackrel{\Pi_{32}}{\vdots}$}\noLine
\UIC{$\mathrm{ContIn}(f,A,c) \to \varphi_{94}$}
\AXC{$(\mathrm{ContIn}(f,A,c) \to \varphi_{94}) \to \varphi_{8}$\lbl{A1}}
\BIC{$\varphi_{8}$}
\end{ptree}
\noindent$\Pi_{34}$:
\begin{ptree}
\AXC{$\forall n \in \mathbb{N}\,(a_{n} \in A) \to n \in \mathbb{N} \to a_{n} \in A$\lbl{A2}}
\AXC{$\varphi_{68}$\lbl{A1}}
\AXC{$\varphi_{68} \to (\forall n \in \mathbb{N}\,(a_{n} \in A) \to n \in \mathbb{N} \to a_{n} \in A) \to \varphi_{63}$\lbl{A1}}
\BIC{$(\forall n \in \mathbb{N}\,(a_{n} \in A) \to n \in \mathbb{N} \to a_{n} \in A) \to \varphi_{63}$}
\BIC{$\varphi_{63}$}
\AXC{$\varphi_{63} \to \varphi_{138}$\lbl{A1}}
\BIC{$\varphi_{138}$}
\end{ptree}
\noindent$\Pi_{35}$:
\begin{ptree}
\AXC{$\forall X\,\forall Y\,\forall x\,(X \subseteq Y \to x \in X \to x \in Y)$\lbl{[B5]}}
\AXC{$\varphi_{38}$\lbl{A2}}
\BIC{$\forall Y\,\forall x\,(A \subseteq Y \to x \in A \to x \in Y)$}
\AXC{$\varphi_{50}$\lbl{A2}}
\BIC{$\forall x\,(A \subseteq \mathrm{dom}(f) \to x \in A \to x \in \mathrm{dom}(f))$}
\AXC{$\varphi_{83}$\lbl{A2}}
\BIC{$A \subseteq \mathrm{dom}(f) \to a_{n} \in A \to a_{n} \in \mathrm{dom}(f)$}
\AXC{$\varphi_{166}$\lbl{A1}}
\BIC{$(\mathrm{ContIn}(f,A,c) \to A \subseteq \mathrm{dom}(f)) \to \varphi_{140}$}
\end{ptree}
\noindent$\Pi_{36}$:
\begin{ptree}
\AXC{$\stackrel{\Pi_{11}}{\vdots}$}\noLine
\UIC{$\mathrm{ContIn}(f,A,c) \to A \subseteq \mathrm{dom}(f)$}
\AXC{$\stackrel{\Pi_{35}}{\vdots}$}\noLine
\UIC{$(\mathrm{ContIn}(f,A,c) \to A \subseteq \mathrm{dom}(f)) \to \varphi_{140}$}
\BIC{$\varphi_{140}$}
\end{ptree}
\noindent$\Pi_{37}$:
\begin{ptree}
\AXC{$\mathrm{ContIn}(f,A,c) \to \mathrm{Seq}(a) \to \varphi_{15}$\lbl{A1}}
\AXC{$\stackrel{\Pi_{34}}{\vdots}$}\noLine
\UIC{$\varphi_{138}$}
\BIC{$\mathrm{ContIn}(f,A,c) \to \mathrm{Seq}(a) \to \varphi_{25}$}
\AXC{$\stackrel{\Pi_{36}}{\vdots}$}\noLine
\UIC{$\varphi_{140}$}
\BIC{$\mathrm{ContIn}(f,A,c) \to \mathrm{Seq}(a) \to \varphi_{44}$}
\AXC{$\varphi_{163}$\lbl{A1}}
\BIC{$\varphi_{13} \to \mathrm{ContIn}(f,A,c) \to \mathrm{Seq}(a) \to \varphi_{56}$}
\end{ptree}
\noindent$\Pi_{38}$:
\begin{ptree}
\AXC{$\varphi_{69}$\lbl{A2}}
\AXC{$\stackrel{\Pi_{33}}{\vdots}$}\noLine
\UIC{$\varphi_{8}$}
\AXC{$\varphi_{8} \to \varphi_{69} \to \varphi_{13}$\lbl{A1}}
\BIC{$\varphi_{69} \to \varphi_{13}$}
\BIC{$\varphi_{13}$}
\AXC{$\stackrel{\Pi_{37}}{\vdots}$}\noLine
\UIC{$\varphi_{13} \to \mathrm{ContIn}(f,A,c) \to \mathrm{Seq}(a) \to \varphi_{56}$}
\BIC{$\mathrm{ContIn}(f,A,c) \to \mathrm{Seq}(a) \to \varphi_{56}$}
\AXC{$\varphi_{148}$\lbl{A1}}
\BIC{$\varphi_{17} \to n \in \mathbb{N} \to a_{n} \in \mathrm{dom}(\lvert f \rvert)$}
\UIC{$\forall n\,(\varphi_{17} \to n \in \mathbb{N} \to a_{n} \in \mathrm{dom}(\lvert f \rvert))$}
\end{ptree}
\noindent$\Pi_{39}$:
\begin{ptree}
\AXC{$\varphi_{154}$\lbl{A3}}
\AXC{$\stackrel{\Pi_{38}}{\vdots}$}\noLine
\UIC{$\forall n\,(\varphi_{17} \to n \in \mathbb{N} \to a_{n} \in \mathrm{dom}(\lvert f \rvert))$}
\AXC{$\forall n\,(\varphi_{17} \to n \in \mathbb{N} \to a_{n} \in \mathrm{dom}(\lvert f \rvert)) \to \varphi_{154} \to \varphi_{160}$\lbl{A1}}
\BIC{$\varphi_{154} \to \varphi_{160}$}
\BIC{$\varphi_{160}$}
\end{ptree}
\noindent$\Pi_{40}$:
\begin{ptree}
\AXC{$\varphi_{17} \to \forall n\,(\varphi_{17})$\lbl{A4}}
\AXC{$\stackrel{\Pi_{39}}{\vdots}$}\noLine
\UIC{$\varphi_{160}$}
\BIC{$\mathrm{ContIn}(f,A,c) \to \mathrm{Seq}(a) \to \varphi_{66}$}
\end{ptree}
\noindent$\Pi_{41}$:
\begin{ptree}
\AXC{$\stackrel{\Pi_{32}}{\vdots}$}\noLine
\UIC{$\mathrm{ContIn}(f,A,c) \to \varphi_{94}$}
\AXC{$(\mathrm{ContIn}(f,A,c) \to \varphi_{94}) \to \varphi_{20}$\lbl{A1}}
\BIC{$\varphi_{20}$}
\end{ptree}
\noindent$\Pi_{42}$:
\begin{ptree}
\AXC{$\varphi_{102}$\lbl{A2}}
\AXC{$\stackrel{\Pi_{41}}{\vdots}$}\noLine
\UIC{$\varphi_{20}$}
\AXC{$\varphi_{20} \to \varphi_{102} \to \varphi_{70}$\lbl{A1}}
\BIC{$\varphi_{102} \to \varphi_{70}$}
\BIC{$\varphi_{70}$}
\AXC{$\varphi_{70} \to (\mathrm{ContIn}(f,A,c) \to \mathrm{Seq}(a) \to \varphi_{56}) \to \varphi_{150}$\lbl{A1}}
\BIC{$(\mathrm{ContIn}(f,A,c) \to \mathrm{Seq}(a) \to \varphi_{56}) \to \varphi_{150}$}
\end{ptree}
\noindent$\Pi_{43}$:
\begin{ptree}
\AXC{$\varphi_{54}$\lbl{A2}}
\AXC{$\stackrel{\Pi_{40}}{\vdots}$}\noLine
\UIC{$\mathrm{ContIn}(f,A,c) \to \mathrm{Seq}(a) \to \varphi_{66}$}
\AXC{$\varphi_{167}$\lbl{A1}}
\BIC{$\varphi_{54} \to \mathrm{ContIn}(f,A,c) \to \mathrm{Seq}(a) \to \varphi_{56}$}
\BIC{$\mathrm{ContIn}(f,A,c) \to \mathrm{Seq}(a) \to \varphi_{56}$}
\AXC{$\stackrel{\Pi_{42}}{\vdots}$}\noLine
\UIC{$(\mathrm{ContIn}(f,A,c) \to \mathrm{Seq}(a) \to \varphi_{56}) \to \varphi_{150}$}
\BIC{$\varphi_{150}$}
\end{ptree}
\noindent$\Pi_{44}$:
\begin{ptree}
\AXC{$\stackrel{\Pi_{32}}{\vdots}$}\noLine
\UIC{$\mathrm{ContIn}(f,A,c) \to \varphi_{94}$}
\AXC{$\varphi_{127}$\lbl{A1}}
\BIC{$\mathrm{ContIn}(f,A,c) \to \mathrm{RFn}(\lvert f \rvert)$}
\end{ptree}
\noindent$\Pi_{45}$:
\begin{ptree}
\AXC{$\forall f\,\forall a\,(\mathrm{RFn}(f) \to \mathrm{Seq}(a) \to \varphi_{96})$\lbl{[S8]}}
\AXC{$\forall f\,\forall a\,(\mathrm{RFn}(f) \to \mathrm{Seq}(a) \to \varphi_{96}) \to \forall a\,(\varphi_{128})$\lbl{A2}}
\BIC{$\forall a\,(\varphi_{128})$}
\AXC{$\forall a\,(\varphi_{128}) \to \varphi_{128}$\lbl{A2}}
\BIC{$\varphi_{128}$}
\end{ptree}
\noindent$\Pi_{46}$:
\begin{ptree}
\AXC{$\stackrel{\Pi_{44}}{\vdots}$}\noLine
\UIC{$\mathrm{ContIn}(f,A,c) \to \mathrm{RFn}(\lvert f \rvert)$}
\AXC{$\stackrel{\Pi_{45}}{\vdots}$}\noLine
\UIC{$\varphi_{128}$}
\AXC{$\varphi_{128} \to (\mathrm{ContIn}(f,A,c) \to \mathrm{RFn}(\lvert f \rvert)) \to \varphi_{168}$\lbl{A1}}
\BIC{$(\mathrm{ContIn}(f,A,c) \to \mathrm{RFn}(\lvert f \rvert)) \to \varphi_{168}$}
\BIC{$\varphi_{168}$}
\end{ptree}
\noindent$\Pi_{47}$:
\begin{ptree}
\AXC{$\stackrel{\Pi_{40}}{\vdots}$}\noLine
\UIC{$\mathrm{ContIn}(f,A,c) \to \mathrm{Seq}(a) \to \varphi_{66}$}
\AXC{$\mathrm{ContIn}(f,A,c) \to \mathrm{Seq}(a) \to \mathrm{Seq}(a)$\lbl{A1}}
\AXC{$\stackrel{\Pi_{46}}{\vdots}$}\noLine
\UIC{$\varphi_{168}$}
\BIC{$(\mathrm{ContIn}(f,A,c) \to \mathrm{Seq}(a) \to \varphi_{66}) \to \varphi_{130}$}
\BIC{$\varphi_{130}$}
\end{ptree}
\noindent$\Pi_{48}$:
\begin{ptree}
\AXC{$\varphi_{101}$\lbl{A2}}
\AXC{$\stackrel{\Pi_{47}}{\vdots}$}\noLine
\UIC{$\varphi_{130}$}
\AXC{$\varphi_{130} \to \varphi_{121}$\lbl{A1}}
\BIC{$\varphi_{121}$}
\AXC{$\varphi_{121} \to \varphi_{101} \to \varphi_{118}$\lbl{A1}}
\BIC{$\varphi_{101} \to \varphi_{118}$}
\BIC{$\varphi_{118}$}
\AXC{$\varphi_{118} \to (\mathrm{ContIn}(f,A,c) \to \mathrm{Seq}(a) \to \varphi_{15}) \to \varphi_{109}$\lbl{A1}}
\BIC{$(\mathrm{ContIn}(f,A,c) \to \mathrm{Seq}(a) \to \varphi_{15}) \to \varphi_{109}$}
\end{ptree}
\noindent$\Pi_{49}$:
\begin{ptree}
\AXC{$\mathrm{ContIn}(f,A,c) \to \mathrm{Seq}(a) \to \varphi_{15}$\lbl{A1}}
\AXC{$\stackrel{\Pi_{48}}{\vdots}$}\noLine
\UIC{$(\mathrm{ContIn}(f,A,c) \to \mathrm{Seq}(a) \to \varphi_{15}) \to \varphi_{109}$}
\BIC{$\varphi_{109}$}
\AXC{$\varphi_{109} \to \varphi_{173}$\lbl{A1}}
\BIC{$\varphi_{173}$}
\end{ptree}
\noindent$\Pi_{50}$:
\begin{ptree}
\AXC{$\stackrel{\Pi_{30}}{\vdots}$}\noLine
\UIC{$\lvert f \rvert \circ a(n) = \lvert f \rvert(a_{n}) \to \varphi_{79}$}
\AXC{$\mathrm{ContIn}(f,A,c) \to \mathrm{Seq}(a) \to \varphi_{15}$\lbl{A1}}
\AXC{$\stackrel{\Pi_{43}}{\vdots}$}\noLine
\UIC{$\varphi_{150}$}
\BIC{$\mathrm{ContIn}(f,A,c) \to \mathrm{Seq}(a) \to \varphi_{85}$}
\AXC{$\stackrel{\Pi_{49}}{\vdots}$}\noLine
\UIC{$\varphi_{173}$}
\BIC{$(\lvert f \rvert \circ a(n) = \lvert f \rvert(a_{n}) \to \varphi_{79}) \to \varphi_{106}$}
\BIC{$\varphi_{106}$}
\end{ptree}
\noindent$\Pi_{51}$:
\begin{ptree}
\AXC{$\stackrel{\Pi_{6}}{\vdots}$}\noLine
\UIC{$\varphi_{81}$}
\AXC{$\stackrel{\Pi_{28}}{\vdots}$}\noLine
\UIC{$\varphi_{104}$}
\AXC{$\varphi_{104} \to \varphi_{81} \to \varphi_{105}$\lbl{A1}}
\BIC{$\varphi_{81} \to \varphi_{105}$}
\BIC{$\varphi_{105}$}
\AXC{$\stackrel{\Pi_{50}}{\vdots}$}\noLine
\UIC{$\varphi_{106}$}
\AXC{$\varphi_{106} \to \varphi_{105} \to (\lvert f \rvert \circ a(n) = \lvert f(a_{n}) \rvert \to \varphi_{88}) \to \varphi_{113}$\lbl{A1}}
\BIC{$\varphi_{105} \to (\lvert f \rvert \circ a(n) = \lvert f(a_{n}) \rvert \to \varphi_{88}) \to \varphi_{113}$}
\BIC{$(\lvert f \rvert \circ a(n) = \lvert f(a_{n}) \rvert \to \varphi_{88}) \to \varphi_{113}$}
\end{ptree}
\noindent$\Pi_{52}$:
\begin{ptree}
\AXC{$\stackrel{\Pi_{5}}{\vdots}$}\noLine
\UIC{$\varphi_{21}$}
\AXC{$\varphi_{21} \to \varphi_{90}$\lbl{A2}}
\BIC{$\varphi_{90}$}
\AXC{$\varphi_{90} \to \lvert f \rvert \circ a(n) = \lvert f(a_{n}) \rvert \to \varphi_{88}$\lbl{A2}}
\BIC{$\lvert f \rvert \circ a(n) = \lvert f(a_{n}) \rvert \to \varphi_{88}$}
\AXC{$\stackrel{\Pi_{51}}{\vdots}$}\noLine
\UIC{$(\lvert f \rvert \circ a(n) = \lvert f(a_{n}) \rvert \to \varphi_{88}) \to \varphi_{113}$}
\BIC{$\varphi_{113}$}
\end{ptree}
\noindent$\Pi_{53}$:
\begin{ptree}
\AXC{$\varphi_{164}$\lbl{A3}}
\AXC{$\stackrel{\Pi_{4}}{\vdots}$}\noLine
\UIC{$\varphi_{92}$}
\AXC{$\stackrel{\Pi_{52}}{\vdots}$}\noLine
\UIC{$\varphi_{113}$}
\AXC{$\varphi_{113} \to \varphi_{92} \to \varphi_{114}$\lbl{A1}}
\BIC{$\varphi_{92} \to \varphi_{114}$}
\BIC{$\varphi_{114}$}
\AXC{$\varphi_{114} \to \varphi_{17} \to n \in \mathbb{N} \to \lvert f \circ a \rvert(n) = \lvert f \rvert \circ a(n)$\lbl{A1}}
\BIC{$\varphi_{17} \to n \in \mathbb{N} \to \lvert f \circ a \rvert(n) = \lvert f \rvert \circ a(n)$}
\UIC{$\forall n\,(\varphi_{17} \to n \in \mathbb{N} \to \lvert f \circ a \rvert(n) = \lvert f \rvert \circ a(n))$}
\AXC{$\varphi_{178}$\lbl{A1}}
\BIC{$\varphi_{164} \to (\varphi_{17} \to \forall n\,(\varphi_{17})) \to \varphi_{115}$}
\BIC{$(\varphi_{17} \to \forall n\,(\varphi_{17})) \to \varphi_{115}$}
\end{ptree}
\noindent$\Pi_{54}$:
\begin{ptree}
\AXC{$\varphi_{47}$\lbl{[C13]}}
\AXC{$\varphi_{47} \to \varphi_{98}$\lbl{A2}}
\BIC{$\varphi_{98}$}
\AXC{$\varphi_{98} \to \varphi_{129}$\lbl{A2}}
\BIC{$\varphi_{129}$}
\AXC{$\varphi_{129} \to \varphi_{133}$\lbl{A2}}
\BIC{$\varphi_{133}$}
\AXC{$\varphi_{133} \to (\mathrm{ContIn}(f,A,c) \to \mathrm{Seq}(a) \to \varphi_{18}) \to \varphi_{115} \to \varphi_{151}$\lbl{A1}}
\BIC{$(\mathrm{ContIn}(f,A,c) \to \mathrm{Seq}(a) \to \varphi_{18}) \to \varphi_{115} \to \varphi_{151}$}
\end{ptree}
\noindent$\Pi_{55}$:
\begin{ptree}
\AXC{$\stackrel{\Pi_{47}}{\vdots}$}\noLine
\UIC{$\varphi_{130}$}
\AXC{$\varphi_{130} \to \mathrm{ContIn}(f,A,c) \to \mathrm{Seq}(a) \to \varphi_{18}$\lbl{A1}}
\BIC{$\mathrm{ContIn}(f,A,c) \to \mathrm{Seq}(a) \to \varphi_{18}$}
\AXC{$\stackrel{\Pi_{54}}{\vdots}$}\noLine
\UIC{$(\mathrm{ContIn}(f,A,c) \to \mathrm{Seq}(a) \to \varphi_{18}) \to \varphi_{115} \to \varphi_{151}$}
\BIC{$\varphi_{115} \to \varphi_{151}$}
\end{ptree}
\noindent$\Pi_{56}$:
\begin{ptree}
\AXC{$\stackrel{\Pi_{1}}{\vdots}$}\noLine
\UIC{$\mathrm{ContIn}(f,A,c) \to \mathrm{Seq}(a) \to \varphi_{30}$}
\AXC{$\stackrel{\Pi_{2}}{\vdots}$}\noLine
\UIC{$\varphi_{143}$}
\BIC{$\mathrm{ContIn}(f,A,c) \to \mathrm{Seq}(a) \to \varphi_{60}$}
\AXC{$\varphi_{17} \to \forall n\,(\varphi_{17})$\lbl{A4}}
\AXC{$\stackrel{\Pi_{53}}{\vdots}$}\noLine
\UIC{$(\varphi_{17} \to \forall n\,(\varphi_{17})) \to \varphi_{115}$}
\BIC{$\varphi_{115}$}
\AXC{$\stackrel{\Pi_{55}}{\vdots}$}\noLine
\UIC{$\varphi_{115} \to \varphi_{151}$}
\BIC{$\varphi_{151}$}
\BIC{$\mathrm{ContIn}(f,A,c) \to \mathrm{Seq}(a) \to \varphi_{62}$}
\end{ptree}
\noindent$\Pi_{57}$:
\begin{ptree}
\AXC{$\forall f\,\forall A\,\forall c\,(\varphi_{134})$\lbl{DefContIn}}
\AXC{$\forall f\,\forall A\,\forall c\,(\varphi_{134}) \to \forall A\,\forall c\,(\varphi_{134})$\lbl{A2}}
\BIC{$\forall A\,\forall c\,(\varphi_{134})$}
\AXC{$\forall A\,\forall c\,(\varphi_{134}) \to \forall c\,(\varphi_{134})$\lbl{A2}}
\BIC{$\forall c\,(\varphi_{134})$}
\AXC{$\forall c\,(\varphi_{134}) \to \varphi_{134}$\lbl{A2}}
\BIC{$\varphi_{134}$}
\end{ptree}
\noindent$\Pi_{58}$:
\begin{ptree}
\AXC{$\stackrel{\Pi_{57}}{\vdots}$}\noLine
\UIC{$\varphi_{134}$}
\AXC{$\mathrm{ContIn}(f,A,c) \to \mathrm{ContIn}(f,A,c)$\lbl{A1}}
\AXC{$\varphi_{157}$\lbl{A1}}
\BIC{$\varphi_{134} \to \mathrm{ContIn}(f,A,c) \to c \in A$}
\BIC{$\mathrm{ContIn}(f,A,c) \to c \in A$}
\end{ptree}
\noindent$\Pi_{59}$:
\begin{ptree}
\AXC{$\forall X\,\forall Y\,\forall x\,(X \subseteq Y \to x \in X \to x \in Y)$\lbl{[B5]}}
\AXC{$\varphi_{38}$\lbl{A2}}
\BIC{$\forall Y\,\forall x\,(A \subseteq Y \to x \in A \to x \in Y)$}
\AXC{$\varphi_{50}$\lbl{A2}}
\BIC{$\forall x\,(A \subseteq \mathrm{dom}(f) \to x \in A \to x \in \mathrm{dom}(f))$}
\AXC{$\varphi_{57}$\lbl{A2}}
\BIC{$A \subseteq \mathrm{dom}(f) \to c \in A \to c \in \mathrm{dom}(f)$}
\AXC{$\varphi_{135}$\lbl{A1}}
\BIC{$(\mathrm{ContIn}(f,A,c) \to A \subseteq \mathrm{dom}(f)) \to \varphi_{95}$}
\end{ptree}
\noindent$\Pi_{60}$:
\begin{ptree}
\AXC{$\stackrel{\Pi_{11}}{\vdots}$}\noLine
\UIC{$\mathrm{ContIn}(f,A,c) \to A \subseteq \mathrm{dom}(f)$}
\AXC{$\stackrel{\Pi_{59}}{\vdots}$}\noLine
\UIC{$(\mathrm{ContIn}(f,A,c) \to A \subseteq \mathrm{dom}(f)) \to \varphi_{95}$}
\BIC{$\varphi_{95}$}
\end{ptree}
\noindent$\Pi_{61}$:
\begin{ptree}
\AXC{$\varphi_{33}$\lbl{A2}}
\AXC{$\stackrel{\Pi_{33}}{\vdots}$}\noLine
\UIC{$\varphi_{8}$}
\AXC{$\varphi_{8} \to \varphi_{33} \to \varphi_{5}$\lbl{A1}}
\BIC{$\varphi_{33} \to \varphi_{5}$}
\BIC{$\varphi_{5}$}
\AXC{$\stackrel{\Pi_{58}}{\vdots}$}\noLine
\UIC{$\mathrm{ContIn}(f,A,c) \to c \in A$}
\AXC{$\stackrel{\Pi_{60}}{\vdots}$}\noLine
\UIC{$\varphi_{95}$}
\BIC{$\mathrm{ContIn}(f,A,c) \to \varphi_{16}$}
\AXC{$\varphi_{153}$\lbl{A1}}
\BIC{$\varphi_{5} \to \mathrm{ContIn}(f,A,c) \to \mathrm{Seq}(a) \to \varphi_{12}$}
\BIC{$\mathrm{ContIn}(f,A,c) \to \mathrm{Seq}(a) \to \varphi_{12}$}
\end{ptree}
\noindent$\Pi_{62}$:
\begin{ptree}
\AXC{$\stackrel{\Pi_{61}}{\vdots}$}\noLine
\UIC{$\mathrm{ContIn}(f,A,c) \to \mathrm{Seq}(a) \to \varphi_{12}$}
\AXC{$\varphi_{91}$\lbl{A2}}
\AXC{$\stackrel{\Pi_{41}}{\vdots}$}\noLine
\UIC{$\varphi_{20}$}
\AXC{$\varphi_{20} \to \varphi_{91} \to \varphi_{14}$\lbl{A1}}
\BIC{$\varphi_{91} \to \varphi_{14}$}
\BIC{$\varphi_{14}$}
\AXC{$\varphi_{14} \to \varphi_{137}$\lbl{A1}}
\BIC{$\varphi_{137}$}
\BIC{$\mathrm{ContIn}(f,A,c) \to \mathrm{Seq}(a) \to \varphi_{35}$}
\end{ptree}
\noindent$\Pi_{63}$:
\begin{ptree}
\AXC{$\stackrel{\Pi_{62}}{\vdots}$}\noLine
\UIC{$\mathrm{ContIn}(f,A,c) \to \mathrm{Seq}(a) \to \varphi_{35}$}
\AXC{$(\mathrm{ContIn}(f,A,c) \to \mathrm{Seq}(a) \to \varphi_{35}) \to \varphi_{131}$\lbl{A1}}
\BIC{$\varphi_{131}$}
\end{ptree}
\noindent$\Pi_{64}$:
\begin{ptree}
\AXC{$\stackrel{\Pi_{56}}{\vdots}$}\noLine
\UIC{$\mathrm{ContIn}(f,A,c) \to \mathrm{Seq}(a) \to \varphi_{62}$}
\AXC{$\stackrel{\Pi_{3}}{\vdots}$}\noLine
\UIC{$\forall x\,\forall y\,(x = y \to y = x)$}
\AXC{$\varphi_{9}$\lbl{A2}}
\BIC{$\forall y\,(\lvert f \rvert(c) = y \to y = \lvert f \rvert(c))$}
\AXC{$\varphi_{87}$\lbl{A2}}
\BIC{$\lvert f \rvert(c) = \lvert f(c) \rvert \to \lvert f(c) \rvert = \lvert f \rvert(c)$}
\AXC{$\stackrel{\Pi_{63}}{\vdots}$}\noLine
\UIC{$\varphi_{131}$}
\BIC{$\mathrm{ContIn}(f,A,c) \to \mathrm{Seq}(a) \to \varphi_{36}$}
\AXC{$\varphi_{176}$\lbl{A1}}
\BIC{$(\mathrm{ContIn}(f,A,c) \to \mathrm{Seq}(a) \to \varphi_{62}) \to \varphi_{152}$}
\BIC{$\varphi_{152}$}
\end{ptree}
\noindent$\Pi_{65}$:
\begin{ptree}
\AXC{$\varphi_{155}$\lbl{A3}}
\AXC{$\varphi_{11}$\lbl{A6}}
\UIC{$\forall w\,(\varphi_{11})$}
\UIC{$\forall z\,\forall w\,(\varphi_{11})$}
\AXC{$\forall z\,\forall w\,(\varphi_{11}) \to \varphi_{73}$\lbl{A2}}
\BIC{$\varphi_{73}$}
\AXC{$\varphi_{73} \to \lvert f(c) \rvert = \lvert f \rvert(c) \to \varphi_{55}$\lbl{A2}}
\BIC{$\lvert f(c) \rvert = \lvert f \rvert(c) \to \varphi_{55}$}
\AXC{$\stackrel{\Pi_{64}}{\vdots}$}\noLine
\UIC{$\varphi_{152}$}
\BIC{$\mathrm{ContIn}(f,A,c) \to \mathrm{Seq}(a) \to \varphi_{65}$}
\UIC{$\forall a\,(\mathrm{ContIn}(f,A,c) \to \mathrm{Seq}(a) \to \varphi_{65})$}
\AXC{$\varphi_{177}$\lbl{A1}}
\BIC{$\varphi_{155} \to \varphi_{1} \to \mathrm{ContIn}(f,A,c) \to \forall a\,(\mathrm{Seq}(a) \to \varphi_{65})$}
\BIC{$\varphi_{1} \to \mathrm{ContIn}(f,A,c) \to \forall a\,(\mathrm{Seq}(a) \to \varphi_{65})$}
\end{ptree}
\noindent$\Pi_{66}$:
\begin{ptree}
\AXC{$\forall X\,\forall Y\,\forall x\,(X \subseteq Y \to x \in X \to x \in Y)$\lbl{[B5]}}
\AXC{$\varphi_{38}$\lbl{A2}}
\BIC{$\forall Y\,\forall x\,(A \subseteq Y \to x \in A \to x \in Y)$}
\AXC{$\varphi_{50}$\lbl{A2}}
\BIC{$\forall x\,(A \subseteq \mathrm{dom}(f) \to x \in A \to x \in \mathrm{dom}(f))$}
\AXC{$\varphi_{58}$\lbl{A2}}
\BIC{$A \subseteq \mathrm{dom}(f) \to u \in A \to u \in \mathrm{dom}(f)$}
\AXC{$\varphi_{125}$\lbl{A1}}
\BIC{$(\mathrm{ContIn}(f,A,c) \to A \subseteq \mathrm{dom}(f)) \to \varphi_{40}$}
\end{ptree}
\noindent$\Pi_{67}$:
\begin{ptree}
\AXC{$\stackrel{\Pi_{11}}{\vdots}$}\noLine
\UIC{$\mathrm{ContIn}(f,A,c) \to A \subseteq \mathrm{dom}(f)$}
\AXC{$\stackrel{\Pi_{66}}{\vdots}$}\noLine
\UIC{$(\mathrm{ContIn}(f,A,c) \to A \subseteq \mathrm{dom}(f)) \to \varphi_{40}$}
\BIC{$\varphi_{40}$}
\end{ptree}
\noindent$\Pi_{68}$:
\begin{ptree}
\AXC{$\varphi_{34}$\lbl{A2}}
\AXC{$\stackrel{\Pi_{33}}{\vdots}$}\noLine
\UIC{$\varphi_{8}$}
\AXC{$\varphi_{8} \to \varphi_{34} \to \varphi_{6}$\lbl{A1}}
\BIC{$\varphi_{34} \to \varphi_{6}$}
\BIC{$\varphi_{6}$}
\AXC{$\mathrm{ContIn}(f,A,c) \to u \in A \to u \in A$\lbl{A1}}
\AXC{$\stackrel{\Pi_{67}}{\vdots}$}\noLine
\UIC{$\varphi_{40}$}
\BIC{$\mathrm{ContIn}(f,A,c) \to u \in A \to u \in \mathrm{dom}(f)$}
\AXC{$\varphi_{116}$\lbl{A1}}
\BIC{$\varphi_{6} \to \mathrm{ContIn}(f,A,c) \to u \in A \to u \in \mathrm{dom}(\lvert f \rvert)$}
\BIC{$\mathrm{ContIn}(f,A,c) \to u \in A \to u \in \mathrm{dom}(\lvert f \rvert)$}
\UIC{$\varphi_{4}$}
\end{ptree}
\noindent$\Pi_{69}$:
\begin{ptree}
\AXC{$\forall x\,\forall y\,(x \subseteq y \leftrightarrow \forall u \in x\,(u \in y))$\lbl{Def$\subseteq$}}
\AXC{$\varphi_{41}$\lbl{A2}}
\BIC{$\forall y\,(A \subseteq y \leftrightarrow \forall u \in A\,(u \in y))$}
\AXC{$\varphi_{76}$\lbl{A2}}
\BIC{$A \subseteq \mathrm{dom}(\lvert f \rvert) \leftrightarrow \forall u \in A\,(u \in \mathrm{dom}(\lvert f \rvert))$}
\AXC{$\varphi_{119}$\lbl{A1}}
\BIC{$\varphi_{67}$}
\end{ptree}
\noindent$\Pi_{70}$:
\begin{ptree}
\AXC{$\varphi_{122}$\lbl{[G2]}}
\AXC{$\varphi_{122} \to \forall A\,\forall c\,(\mathrm{RFn}(\lvert f \rvert) \to A \subseteq \mathrm{dom}(\lvert f \rvert) \to \varphi_{107})$\lbl{A2}}
\BIC{$\forall A\,\forall c\,(\mathrm{RFn}(\lvert f \rvert) \to A \subseteq \mathrm{dom}(\lvert f \rvert) \to \varphi_{107})$}
\AXC{$\varphi_{171}$\lbl{A2}}
\BIC{$\forall c\,(\mathrm{RFn}(\lvert f \rvert) \to A \subseteq \mathrm{dom}(\lvert f \rvert) \to \varphi_{107})$}
\AXC{$\varphi_{170}$\lbl{A2}}
\BIC{$\mathrm{RFn}(\lvert f \rvert) \to A \subseteq \mathrm{dom}(\lvert f \rvert) \to \varphi_{107}$}
\AXC{$\varphi_{174}$\lbl{A1}}
\BIC{$(\mathrm{ContIn}(f,A,c) \to \mathrm{RFn}(\lvert f \rvert)) \to \varphi_{147}$}
\end{ptree}
\noindent$\Pi_{71}$:
\begin{ptree}
\AXC{$\stackrel{\Pi_{44}}{\vdots}$}\noLine
\UIC{$\mathrm{ContIn}(f,A,c) \to \mathrm{RFn}(\lvert f \rvert)$}
\AXC{$\stackrel{\Pi_{70}}{\vdots}$}\noLine
\UIC{$(\mathrm{ContIn}(f,A,c) \to \mathrm{RFn}(\lvert f \rvert)) \to \varphi_{147}$}
\BIC{$\varphi_{147}$}
\end{ptree}
\noindent$\Pi_{72}$:
\begin{ptree}
\AXC{$\varphi_{2}$\lbl{A4}}
\AXC{$\varphi_{4} \to \varphi_{7}$\lbl{A3}}
\AXC{$\stackrel{\Pi_{68}}{\vdots}$}\noLine
\UIC{$\varphi_{4}$}
\AXC{$\varphi_{4} \to (\varphi_{4} \to \varphi_{7}) \to \varphi_{59}$\lbl{A1}}
\BIC{$(\varphi_{4} \to \varphi_{7}) \to \varphi_{59}$}
\BIC{$\varphi_{59}$}
\BIC{$\mathrm{ContIn}(f,A,c) \to \forall u \in A\,(u \in \mathrm{dom}(\lvert f \rvert))$}
\AXC{$\stackrel{\Pi_{69}}{\vdots}$}\noLine
\UIC{$\varphi_{67}$}
\BIC{$\mathrm{ContIn}(f,A,c) \to A \subseteq \mathrm{dom}(\lvert f \rvert)$}
\AXC{$\stackrel{\Pi_{71}}{\vdots}$}\noLine
\UIC{$\varphi_{147}$}
\BIC{$(\mathrm{ContIn}(f,A,c) \to c \in A) \to \varphi_{126}$}
\end{ptree}
\noindent$\Pi_{73}$:
\begin{ptree}
\AXC{$\stackrel{\Pi_{58}}{\vdots}$}\noLine
\UIC{$\mathrm{ContIn}(f,A,c) \to c \in A$}
\AXC{$\stackrel{\Pi_{72}}{\vdots}$}\noLine
\UIC{$(\mathrm{ContIn}(f,A,c) \to c \in A) \to \varphi_{126}$}
\BIC{$\varphi_{126}$}
\end{ptree}
\noindent 主木:
\begin{ptree}
\AXC{$\varphi_{1}$\lbl{A4}}
\AXC{$\stackrel{\Pi_{65}}{\vdots}$}\noLine
\UIC{$\varphi_{1} \to \mathrm{ContIn}(f,A,c) \to \forall a\,(\mathrm{Seq}(a) \to \varphi_{65})$}
\BIC{$\mathrm{ContIn}(f,A,c) \to \forall a\,(\mathrm{Seq}(a) \to \varphi_{65})$}
\AXC{$\stackrel{\Pi_{73}}{\vdots}$}\noLine
\UIC{$\varphi_{126}$}
\BIC{$\mathrm{ContIn}(f,A,c) \to \mathrm{ContIn}(\lvert f \rvert,A,c)$}
\UIC{$\varphi_{3}$}
\UIC{$\forall A\,(\varphi_{3})$}
\UIC{$\forall f\,\forall A\,(\varphi_{3})$}
\end{ptree}
\begin{abbrevlist}
\item $\varphi_{1} :\equiv \mathrm{ContIn}(f,A,c) \to \forall a\,(\mathrm{ContIn}(f,A,c))$
\item $\varphi_{2} :\equiv \mathrm{ContIn}(f,A,c) \to \forall u\,(\mathrm{ContIn}(f,A,c))$
\item $\varphi_{3} :\equiv \forall c\,(\mathrm{ContIn}(f,A,c) \to \mathrm{ContIn}(\lvert f \rvert,A,c))$
\item $\varphi_{4} :\equiv \forall u\,(\mathrm{ContIn}(f,A,c) \to u \in A \to u \in \mathrm{dom}(\lvert f \rvert))$
\item $\varphi_{5} :\equiv \mathrm{ContIn}(f,A,c) \to (c \in \mathrm{dom}(\lvert f \rvert) \leftrightarrow c \in \mathrm{dom}(f))$
\item $\varphi_{6} :\equiv \mathrm{ContIn}(f,A,c) \to (u \in \mathrm{dom}(\lvert f \rvert) \leftrightarrow u \in \mathrm{dom}(f))$
\item $\varphi_{7} :\equiv \forall u\,(\mathrm{ContIn}(f,A,c)) \to \forall u \in A\,(u \in \mathrm{dom}(\lvert f \rvert))$
\item $\varphi_{8} :\equiv \mathrm{ContIn}(f,A,c) \to \forall x\,(x \in \mathrm{dom}(\lvert f \rvert) \leftrightarrow x \in \mathrm{dom}(f))$
\item $\varphi_{9} :\equiv \forall x\,\forall y\,(x = y \to y = x) \to \forall y\,(\lvert f \rvert(c) = y \to y = \lvert f \rvert(c))$
\item $\varphi_{10} :\equiv \forall \alpha\,(\mathrm{Conv}(f \circ a,\alpha) \to \mathrm{Conv}(\lvert f \circ a \rvert,\lvert \alpha \rvert))$
\item $\varphi_{11} :\equiv z = w \to \mathrm{Conv}(\lvert f \rvert \circ a,z) \to \mathrm{Conv}(\lvert f \rvert \circ a,w)$
\item $\varphi_{12} :\equiv \forall n \in \mathbb{N}\,(a_{n} \in A) \to \mathrm{Conv}(a,c) \to c \in \mathrm{dom}(\lvert f \rvert)$
\item $\varphi_{13} :\equiv \mathrm{ContIn}(f,A,c) \to (a_{n} \in \mathrm{dom}(\lvert f \rvert) \leftrightarrow a_{n} \in \mathrm{dom}(f))$
\item $\varphi_{14} :\equiv \mathrm{ContIn}(f,A,c) \to c \in \mathrm{dom}(\lvert f \rvert) \to \lvert f \rvert(c) = \lvert f(c) \rvert$
\item $\varphi_{15} :\equiv \forall n \in \mathbb{N}\,(a_{n} \in A) \to \mathrm{Conv}(a,c) \to n \in \mathbb{N} \to n \in \mathbb{N}$
\item $\varphi_{16} :\equiv \mathrm{Seq}(a) \to \forall n \in \mathbb{N}\,(a_{n} \in A) \to \mathrm{Conv}(a,c) \to c \in \mathrm{dom}(f)$
\item $\varphi_{17} :\equiv \mathrm{ContIn}(f,A,c) \land \mathrm{Seq}(a) \land \forall n \in \mathbb{N}\,(a_{n} \in A) \land \mathrm{Conv}(a,c)$
\item $\varphi_{18} :\equiv \forall n \in \mathbb{N}\,(a_{n} \in A) \to \mathrm{Conv}(a,c) \to \mathrm{Seq}(\lvert f \rvert \circ a)$
\item $\varphi_{19} :\equiv \mathrm{Conv}(f \circ a,f(c)) \to \mathrm{Conv}(\lvert f \circ a \rvert,\lvert f(c) \rvert)$
\item $\varphi_{20} :\equiv \mathrm{ContIn}(f,A,c) \to \forall x \in \mathrm{dom}(\lvert f \rvert)\,(\lvert f \rvert(x) = \lvert f(x) \rvert)$
\item $\varphi_{21} :\equiv \forall y\,\forall z\,(\lvert f \rvert \circ a(n) = y \to y = z \to \lvert f \rvert \circ a(n) = z)$
\item $\varphi_{22} :\equiv \forall y\,\forall z\,(\lvert f \circ a \rvert(n) = y \to y = z \to \lvert f \circ a \rvert(n) = z)$
\item $\varphi_{23} :\equiv \mathrm{Conv}(a,c) \to n \in \mathbb{N} \to \lvert f \circ a(n) \rvert = \lvert f(a_{n}) \rvert$
\item $\varphi_{24} :\equiv \forall a\,(\mathrm{Seq}(a) \to \mathrm{Seq}(\lvert a \rvert) \land \forall n \in \mathbb{N}\,(\lvert a \rvert(n) = \lvert a_{n} \rvert))$
\item $\varphi_{25} :\equiv \forall n \in \mathbb{N}\,(a_{n} \in A) \to \mathrm{Conv}(a,c) \to n \in \mathbb{N} \to a_{n} \in A$
\item $\varphi_{26} :\equiv \mathrm{Seq}(a) \to \forall n \in \mathbb{N}\,(a_{n} \in A) \to \mathrm{Conv}(a,c) \to \mathrm{Seq}(f \circ a)$
\item $\varphi_{27} :\equiv \mathrm{Conv}(a,c) \to n \in \mathbb{N} \to \lvert f \circ a \rvert(n) = \lvert f(a_{n}) \rvert$
\item $\varphi_{28} :\equiv \mathrm{Conv}(a,c) \to n \in \mathbb{N} \to \lvert f(a_{n}) \rvert = \lvert f \circ a \rvert(n)$
\item $\varphi_{29} :\equiv \mathrm{Conv}(a,c) \to n \in \mathbb{N} \to \lvert f \rvert \circ a(n) = \lvert f(a_{n}) \rvert$
\item $\varphi_{30} :\equiv \forall n \in \mathbb{N}\,(a_{n} \in A) \to \mathrm{Conv}(a,c) \to \mathrm{Conv}(f \circ a,f(c))$
\item $\varphi_{31} :\equiv \mathrm{Conv}(a,c) \to n \in \mathbb{N} \to n \in \mathbb{N} \to f \circ a(n) = f(a_{n})$
\item $\varphi_{32} :\equiv \mathrm{Conv}(a,c) \to n \in \mathbb{N} \to \lvert f \rvert \circ a(n) = \lvert f \rvert(a_{n})$
\item $\varphi_{33} :\equiv \forall x\,(x \in \mathrm{dom}(\lvert f \rvert) \leftrightarrow x \in \mathrm{dom}(f)) \to (c \in \mathrm{dom}(\lvert f \rvert) \leftrightarrow c \in \mathrm{dom}(f))$
\item $\varphi_{34} :\equiv \forall x\,(x \in \mathrm{dom}(\lvert f \rvert) \leftrightarrow x \in \mathrm{dom}(f)) \to (u \in \mathrm{dom}(\lvert f \rvert) \leftrightarrow u \in \mathrm{dom}(f))$
\item $\varphi_{35} :\equiv \forall n \in \mathbb{N}\,(a_{n} \in A) \to \mathrm{Conv}(a,c) \to \lvert f \rvert(c) = \lvert f(c) \rvert$
\item $\varphi_{36} :\equiv \forall n \in \mathbb{N}\,(a_{n} \in A) \to \mathrm{Conv}(a,c) \to \lvert f(c) \rvert = \lvert f \rvert(c)$
\item $\varphi_{37} :\equiv \forall w\,(z = w \to \lvert f \circ a(n) \rvert = \lvert z \rvert \to \lvert f \circ a(n) \rvert = \lvert w \rvert)$
\item $\varphi_{38} :\equiv \forall X\,\forall Y\,\forall x\,(X \subseteq Y \to x \in X \to x \in Y) \to \forall Y\,\forall x\,(A \subseteq Y \to x \in A \to x \in Y)$
\item $\varphi_{39} :\equiv \forall a\,\forall \alpha\,(\mathrm{Conv}(a,\alpha) \to \mathrm{Seq}(a)) \to \forall \alpha\,(\mathrm{Conv}(f \circ a,\alpha) \to \mathrm{Seq}(f \circ a))$
\item $\varphi_{40} :\equiv (\mathrm{ContIn}(f,A,c) \to u \in A \to u \in A) \to \mathrm{ContIn}(f,A,c) \to u \in A \to u \in \mathrm{dom}(f)$
\item $\varphi_{41} :\equiv \forall x\,\forall y\,(x \subseteq y \leftrightarrow \forall u \in x\,(u \in y)) \to \forall y\,(A \subseteq y \leftrightarrow \forall u \in A\,(u \in y))$
\item $\varphi_{42} :\equiv \mathrm{Conv}(a,c) \to n \in \mathbb{N} \to \forall n \in \mathbb{N}\,(f \circ a(n) = f(a_{n}))$
\item $\varphi_{43} :\equiv \mathrm{Conv}(a,c) \to n \in \mathbb{N} \to \lvert f \circ a \rvert(n) = \lvert f \circ a(n) \rvert$
\item $\varphi_{44} :\equiv \forall n \in \mathbb{N}\,(a_{n} \in A) \to \mathrm{Conv}(a,c) \to n \in \mathbb{N} \to a_{n} \in \mathrm{dom}(f)$
\item $\varphi_{45} :\equiv \mathrm{Conv}(a,c) \to n \in \mathbb{N} \to \lvert f \rvert \circ a(n) = \lvert f \circ a \rvert(n)$
\item $\varphi_{46} :\equiv \mathrm{Conv}(a,c) \to n \in \mathbb{N} \to \lvert f \circ a \rvert(n) = \lvert f \rvert \circ a(n)$
\item $\varphi_{47} :\equiv \forall a\,\forall b\,\forall \alpha\,(\mathrm{Seq}(b) \to \forall n \in \mathbb{N}\,(a_{n} = b_{n}) \to \mathrm{Conv}(a,\alpha) \to \mathrm{Conv}(b,\alpha))$
\item $\varphi_{48} :\equiv \forall x\,\forall y\,(x = y \to y = x) \to \forall y\,(\lvert f \rvert \circ a(n) = y \to y = \lvert f \rvert \circ a(n))$
\item $\varphi_{49} :\equiv \forall x\,\forall y\,(x = y \to y = x) \to \forall y\,(\lvert f \circ a \rvert(n) = y \to y = \lvert f \circ a \rvert(n))$
\item $\varphi_{50} :\equiv \forall Y\,\forall x\,(A \subseteq Y \to x \in A \to x \in Y) \to \forall x\,(A \subseteq \mathrm{dom}(f) \to x \in A \to x \in \mathrm{dom}(f))$
\item $\varphi_{51} :\equiv \forall n \in \mathbb{N}\,(a_{n} \in A) \to \mathrm{Conv}(a,c) \to \forall n \in \mathbb{N}\,(a_{n} \in \mathrm{dom}(f))$
\item $\varphi_{52} :\equiv \mathrm{Conv}(a,c) \to \forall n \in \mathbb{N}\,(\lvert f \circ a \rvert(n) = \lvert f \rvert \circ a(n))$
\item $\varphi_{53} :\equiv \mathrm{Conv}(\lvert f \circ a \rvert,\lvert f(c) \rvert) \to \mathrm{Conv}(\lvert f \rvert \circ a,\lvert f(c) \rvert)$
\item $\varphi_{54} :\equiv \forall n \in \mathbb{N}\,(a_{n} \in \mathrm{dom}(\lvert f \rvert)) \to n \in \mathbb{N} \to a_{n} \in \mathrm{dom}(\lvert f \rvert)$
\item $\varphi_{55} :\equiv \mathrm{Conv}(\lvert f \rvert \circ a,\lvert f(c) \rvert) \to \mathrm{Conv}(\lvert f \rvert \circ a,\lvert f \rvert(c))$
\item $\varphi_{56} :\equiv \forall n \in \mathbb{N}\,(a_{n} \in A) \to \mathrm{Conv}(a,c) \to n \in \mathbb{N} \to a_{n} \in \mathrm{dom}(\lvert f \rvert)$
\item $\varphi_{57} :\equiv \forall x\,(A \subseteq \mathrm{dom}(f) \to x \in A \to x \in \mathrm{dom}(f)) \to A \subseteq \mathrm{dom}(f) \to c \in A \to c \in \mathrm{dom}(f)$
\item $\varphi_{58} :\equiv \forall x\,(A \subseteq \mathrm{dom}(f) \to x \in A \to x \in \mathrm{dom}(f)) \to A \subseteq \mathrm{dom}(f) \to u \in A \to u \in \mathrm{dom}(f)$
\item $\varphi_{59} :\equiv \varphi_{2} \to \mathrm{ContIn}(f,A,c) \to \forall u \in A\,(u \in \mathrm{dom}(\lvert f \rvert))$
\item $\varphi_{60} :\equiv \forall n \in \mathbb{N}\,(a_{n} \in A) \to \mathrm{Conv}(a,c) \to \mathrm{Conv}(\lvert f \circ a \rvert,\lvert f(c) \rvert)$
\item $\varphi_{61} :\equiv \mathrm{Conv}(a,c) \to \mathrm{Seq}(f \circ a) \land \forall n \in \mathbb{N}\,(f \circ a(n) = f(a_{n}))$
\item $\varphi_{62} :\equiv \forall n \in \mathbb{N}\,(a_{n} \in A) \to \mathrm{Conv}(a,c) \to \mathrm{Conv}(\lvert f \rvert \circ a,\lvert f(c) \rvert)$
\item $\varphi_{63} :\equiv \mathrm{ContIn}(f,A,c) \to \mathrm{Seq}(a) \to \forall n \in \mathbb{N}\,(a_{n} \in A) \to n \in \mathbb{N} \to a_{n} \in A$
\item $\varphi_{64} :\equiv \mathrm{Conv}(a,c) \to n \in \mathbb{N} \to n \in \mathbb{N} \to \lvert f \rvert \circ a(n) = \lvert f \rvert(a_{n})$
\item $\varphi_{65} :\equiv \forall n \in \mathbb{N}\,(a_{n} \in A) \to \mathrm{Conv}(a,c) \to \mathrm{Conv}(\lvert f \rvert \circ a,\lvert f \rvert(c))$
\item $\varphi_{66} :\equiv \forall n \in \mathbb{N}\,(a_{n} \in A) \to \mathrm{Conv}(a,c) \to \forall n \in \mathbb{N}\,(a_{n} \in \mathrm{dom}(\lvert f \rvert))$
\item $\varphi_{67} :\equiv (\mathrm{ContIn}(f,A,c) \to \forall u \in A\,(u \in \mathrm{dom}(\lvert f \rvert))) \to \mathrm{ContIn}(f,A,c) \to A \subseteq \mathrm{dom}(\lvert f \rvert)$
\item $\varphi_{68} :\equiv \mathrm{ContIn}(f,A,c) \to \mathrm{Seq}(a) \to \forall n \in \mathbb{N}\,(a_{n} \in A) \to \forall n \in \mathbb{N}\,(a_{n} \in A)$
\item $\varphi_{69} :\equiv \forall x\,(x \in \mathrm{dom}(\lvert f \rvert) \leftrightarrow x \in \mathrm{dom}(f)) \to (a_{n} \in \mathrm{dom}(\lvert f \rvert) \leftrightarrow a_{n} \in \mathrm{dom}(f))$
\item $\varphi_{70} :\equiv \mathrm{ContIn}(f,A,c) \to a_{n} \in \mathrm{dom}(\lvert f \rvert) \to \lvert f \rvert(a_{n}) = \lvert f(a_{n}) \rvert$
\item $\varphi_{71} :\equiv \mathrm{Conv}(a,c) \to n \in \mathbb{N} \to n \in \mathbb{N} \to \lvert f \circ a \rvert(n) = \lvert f \circ a(n) \rvert$
\item $\varphi_{72} :\equiv \mathrm{Conv}(a,c) \to n \in \mathbb{N} \to \forall n \in \mathbb{N}\,(\lvert f \rvert \circ a(n) = \lvert f \rvert(a_{n}))$
\item $\varphi_{73} :\equiv \forall w\,(\lvert f(c) \rvert = w \to \mathrm{Conv}(\lvert f \rvert \circ a,\lvert f(c) \rvert) \to \mathrm{Conv}(\lvert f \rvert \circ a,w))$
\item $\varphi_{74} :\equiv \mathrm{Seq}(\lvert f \rvert \circ a) \land \forall n \in \mathbb{N}\,(\lvert f \rvert \circ a(n) = \lvert f \rvert(a_{n}))$
\item $\varphi_{75} :\equiv \mathrm{Conv}(a,c) \to n \in \mathbb{N} \to \forall n \in \mathbb{N}\,(\lvert f \circ a \rvert(n) = \lvert f \circ a(n) \rvert)$
\item $\varphi_{76} :\equiv \forall y\,(A \subseteq y \leftrightarrow \forall u \in A\,(u \in y)) \to (A \subseteq \mathrm{dom}(\lvert f \rvert) \leftrightarrow \forall u \in A\,(u \in \mathrm{dom}(\lvert f \rvert)))$
\item $\varphi_{77} :\equiv \forall n \in \mathbb{N}\,(a_{n} \in A) \to \mathrm{Conv}(a,c) \to n \in \mathbb{N} \to f \circ a(n) = f(a_{n})$
\item $\varphi_{78} :\equiv \mathrm{Seq}(\lvert f \circ a \rvert) \land \forall n \in \mathbb{N}\,(\lvert f \circ a \rvert(n) = \lvert f \circ a(n) \rvert)$
\item $\varphi_{79} :\equiv \lvert f \rvert(a_{n}) = \lvert f(a_{n}) \rvert \to \lvert f \rvert \circ a(n) = \lvert f(a_{n}) \rvert$
\item $\varphi_{80} :\equiv \lvert f \circ a(n) \rvert = \lvert f(a_{n}) \rvert \to \lvert f \circ a \rvert(n) = \lvert f(a_{n}) \rvert$
\item $\varphi_{81} :\equiv \lvert f \circ a \rvert(n) = \lvert f(a_{n}) \rvert \to \lvert f(a_{n}) \rvert = \lvert f \circ a \rvert(n)$
\item $\varphi_{82} :\equiv \forall \alpha\,(\mathrm{Conv}(f \circ a,\alpha) \to \mathrm{Seq}(f \circ a)) \to \mathrm{Conv}(f \circ a,f(c)) \to \mathrm{Seq}(f \circ a)$
\item $\varphi_{83} :\equiv \forall x\,(A \subseteq \mathrm{dom}(f) \to x \in A \to x \in \mathrm{dom}(f)) \to A \subseteq \mathrm{dom}(f) \to a_{n} \in A \to a_{n} \in \mathrm{dom}(f)$
\item $\varphi_{84} :\equiv \lvert f \circ a(n) \rvert = \lvert f \circ a(n) \rvert \to \lvert f \circ a(n) \rvert = \lvert f(a_{n}) \rvert$
\item $\varphi_{85} :\equiv \forall n \in \mathbb{N}\,(a_{n} \in A) \to \mathrm{Conv}(a,c) \to n \in \mathbb{N} \to \lvert f \rvert(a_{n}) = \lvert f(a_{n}) \rvert$
\item $\varphi_{86} :\equiv \forall n \in \mathbb{N}\,(f \circ a(n) = f(a_{n})) \to n \in \mathbb{N} \to f \circ a(n) = f(a_{n})$
\item $\varphi_{87} :\equiv \forall y\,(\lvert f \rvert(c) = y \to y = \lvert f \rvert(c)) \to \lvert f \rvert(c) = \lvert f(c) \rvert \to \lvert f(c) \rvert = \lvert f \rvert(c)$
\item $\varphi_{88} :\equiv \lvert f(a_{n}) \rvert = \lvert f \circ a \rvert(n) \to \lvert f \rvert \circ a(n) = \lvert f \circ a \rvert(n)$
\item $\varphi_{89} :\equiv \forall a\,\forall \alpha\,(\mathrm{Conv}(a,\alpha) \to \mathrm{Conv}(\lvert a \rvert,\lvert \alpha \rvert)) \to \varphi_{10}$
\item $\varphi_{90} :\equiv \forall z\,(\lvert f \rvert \circ a(n) = \lvert f(a_{n}) \rvert \to \lvert f(a_{n}) \rvert = z \to \lvert f \rvert \circ a(n) = z)$
\item $\varphi_{91} :\equiv \forall x \in \mathrm{dom}(\lvert f \rvert)\,(\lvert f \rvert(x) = \lvert f(x) \rvert) \to c \in \mathrm{dom}(\lvert f \rvert) \to \lvert f \rvert(c) = \lvert f(c) \rvert$
\item $\varphi_{92} :\equiv \lvert f \rvert \circ a(n) = \lvert f \circ a \rvert(n) \to \lvert f \circ a \rvert(n) = \lvert f \rvert \circ a(n)$
\item $\varphi_{93} :\equiv \forall z\,(\lvert f \rvert \circ a(n) = \lvert f \rvert(a_{n}) \to \lvert f \rvert(a_{n}) = z \to \lvert f \rvert \circ a(n) = z)$
\item $\varphi_{94} :\equiv \mathrm{RFn}(\lvert f \rvert) \land \forall x\,(x \in \mathrm{dom}(\lvert f \rvert) \leftrightarrow x \in \mathrm{dom}(f)) \land \forall x \in \mathrm{dom}(\lvert f \rvert)\,(\lvert f \rvert(x) = \lvert f(x) \rvert)$
\item $\varphi_{95} :\equiv (\mathrm{ContIn}(f,A,c) \to c \in A) \to \mathrm{ContIn}(f,A,c) \to \varphi_{16}$
\item $\varphi_{96} :\equiv \forall n \in \mathbb{N}\,(a_{n} \in \mathrm{dom}(f)) \to \mathrm{Seq}(f \circ a) \land \forall n \in \mathbb{N}\,(f \circ a(n) = f(a_{n}))$
\item $\varphi_{97} :\equiv \forall f\,\forall A\,\forall c\,\forall a\,(\mathrm{ContIn}(f,A,c) \to \mathrm{Seq}(a) \to \varphi_{30})$
\item $\varphi_{98} :\equiv \forall b\,\forall \alpha\,(\mathrm{Seq}(b) \to \forall n \in \mathbb{N}\,(\lvert f \circ a \rvert(n) = b_{n}) \to \mathrm{Conv}(\lvert f \circ a \rvert,\alpha) \to \mathrm{Conv}(b,\alpha))$
\item $\varphi_{99} :\equiv \forall w\,(f \circ a(n) = w \to \lvert f \circ a(n) \rvert = \lvert f \circ a(n) \rvert \to \lvert f \circ a(n) \rvert = \lvert w \rvert)$
\item $\varphi_{100} :\equiv \forall z\,(\lvert f \circ a \rvert(n) = \lvert f \circ a(n) \rvert \to \lvert f \circ a(n) \rvert = z \to \lvert f \circ a \rvert(n) = z)$
\item $\varphi_{101} :\equiv \forall n \in \mathbb{N}\,(\lvert f \rvert \circ a(n) = \lvert f \rvert(a_{n})) \to n \in \mathbb{N} \to \lvert f \rvert \circ a(n) = \lvert f \rvert(a_{n})$
\item $\varphi_{102} :\equiv \forall x \in \mathrm{dom}(\lvert f \rvert)\,(\lvert f \rvert(x) = \lvert f(x) \rvert) \to a_{n} \in \mathrm{dom}(\lvert f \rvert) \to \lvert f \rvert(a_{n}) = \lvert f(a_{n}) \rvert$
\item $\varphi_{103} :\equiv \mathrm{ContIn}(f,A,c) \to \mathrm{Seq}(a) \to \forall n \in \mathbb{N}\,(a_{n} \in A) \to \varphi_{23}$
\item $\varphi_{104} :\equiv \mathrm{ContIn}(f,A,c) \to \mathrm{Seq}(a) \to \forall n \in \mathbb{N}\,(a_{n} \in A) \to \varphi_{27}$
\item $\varphi_{105} :\equiv \mathrm{ContIn}(f,A,c) \to \mathrm{Seq}(a) \to \forall n \in \mathbb{N}\,(a_{n} \in A) \to \varphi_{28}$
\item $\varphi_{106} :\equiv \mathrm{ContIn}(f,A,c) \to \mathrm{Seq}(a) \to \forall n \in \mathbb{N}\,(a_{n} \in A) \to \varphi_{29}$
\item $\varphi_{107} :\equiv c \in A \to \forall a\,(\mathrm{Seq}(a) \to \varphi_{65}) \to \mathrm{ContIn}(\lvert f \rvert,A,c)$
\item $\varphi_{108} :\equiv \mathrm{ContIn}(f,A,c) \to \mathrm{Seq}(a) \to \forall n \in \mathbb{N}\,(a_{n} \in A) \to \varphi_{31}$
\item $\varphi_{109} :\equiv \mathrm{ContIn}(f,A,c) \to \mathrm{Seq}(a) \to \forall n \in \mathbb{N}\,(a_{n} \in A) \to \varphi_{32}$
\item $\varphi_{110} :\equiv \forall n \in \mathbb{N}\,(\lvert f \circ a \rvert(n) = \lvert f \circ a(n) \rvert) \to n \in \mathbb{N} \to \lvert f \circ a \rvert(n) = \lvert f \circ a(n) \rvert$
\item $\varphi_{111} :\equiv \mathrm{ContIn}(f,A,c) \to \mathrm{Seq}(a) \to \forall n \in \mathbb{N}\,(a_{n} \in A) \to \varphi_{42}$
\item $\varphi_{112} :\equiv \mathrm{ContIn}(f,A,c) \to \mathrm{Seq}(a) \to \forall n \in \mathbb{N}\,(a_{n} \in A) \to \varphi_{43}$
\item $\varphi_{113} :\equiv \mathrm{ContIn}(f,A,c) \to \mathrm{Seq}(a) \to \forall n \in \mathbb{N}\,(a_{n} \in A) \to \varphi_{45}$
\item $\varphi_{114} :\equiv \mathrm{ContIn}(f,A,c) \to \mathrm{Seq}(a) \to \forall n \in \mathbb{N}\,(a_{n} \in A) \to \varphi_{46}$
\item $\varphi_{115} :\equiv \mathrm{ContIn}(f,A,c) \to \mathrm{Seq}(a) \to \forall n \in \mathbb{N}\,(a_{n} \in A) \to \varphi_{52}$
\item $\varphi_{116} :\equiv (\mathrm{ContIn}(f,A,c) \to u \in A \to u \in \mathrm{dom}(f)) \to \varphi_{6} \to \mathrm{ContIn}(f,A,c) \to u \in A \to u \in \mathrm{dom}(\lvert f \rvert)$
\item $\varphi_{117} :\equiv \mathrm{ContIn}(f,A,c) \to \mathrm{Seq}(a) \to \forall n \in \mathbb{N}\,(a_{n} \in A) \to \varphi_{61}$
\item $\varphi_{118} :\equiv \mathrm{ContIn}(f,A,c) \to \mathrm{Seq}(a) \to \forall n \in \mathbb{N}\,(a_{n} \in A) \to \varphi_{64}$
\item $\varphi_{119} :\equiv (A \subseteq \mathrm{dom}(\lvert f \rvert) \leftrightarrow \forall u \in A\,(u \in \mathrm{dom}(\lvert f \rvert))) \to \varphi_{67}$
\item $\varphi_{120} :\equiv \mathrm{ContIn}(f,A,c) \to \mathrm{Seq}(a) \to \forall n \in \mathbb{N}\,(a_{n} \in A) \to \varphi_{71}$
\item $\varphi_{121} :\equiv \mathrm{ContIn}(f,A,c) \to \mathrm{Seq}(a) \to \forall n \in \mathbb{N}\,(a_{n} \in A) \to \varphi_{72}$
\item $\varphi_{122} :\equiv \forall f\,\forall A\,\forall c\,(\mathrm{RFn}(f) \to A \subseteq \mathrm{dom}(f) \to c \in A \to \forall a\,(\mathrm{Seq}(a) \to \varphi_{30}) \to \mathrm{ContIn}(f,A,c))$
\item $\varphi_{123} :\equiv \mathrm{ContIn}(f,A,c) \to \mathrm{Seq}(a) \to \forall n \in \mathbb{N}\,(a_{n} \in A) \to \varphi_{75}$
\item $\varphi_{124} :\equiv (\mathrm{ContIn}(f,A,c) \to \mathrm{RFn}(f)) \to \mathrm{ContIn}(f,A,c) \to \varphi_{94}$
\item $\varphi_{125} :\equiv (A \subseteq \mathrm{dom}(f) \to u \in A \to u \in \mathrm{dom}(f)) \to (\mathrm{ContIn}(f,A,c) \to A \subseteq \mathrm{dom}(f)) \to \varphi_{40}$
\item $\varphi_{126} :\equiv (\mathrm{ContIn}(f,A,c) \to \forall a\,(\mathrm{Seq}(a) \to \varphi_{65})) \to \mathrm{ContIn}(f,A,c) \to \mathrm{ContIn}(\lvert f \rvert,A,c)$
\item $\varphi_{127} :\equiv (\mathrm{ContIn}(f,A,c) \to \varphi_{94}) \to \mathrm{ContIn}(f,A,c) \to \mathrm{RFn}(\lvert f \rvert)$
\item $\varphi_{128} :\equiv \mathrm{RFn}(\lvert f \rvert) \to \mathrm{Seq}(a) \to \forall n \in \mathbb{N}\,(a_{n} \in \mathrm{dom}(\lvert f \rvert)) \to \varphi_{74}$
\item $\varphi_{129} :\equiv \forall \alpha\,(\mathrm{Seq}(\lvert f \rvert \circ a) \to \forall n \in \mathbb{N}\,(\lvert f \circ a \rvert(n) = \lvert f \rvert \circ a(n)) \to \mathrm{Conv}(\lvert f \circ a \rvert,\alpha) \to \mathrm{Conv}(\lvert f \rvert \circ a,\alpha))$
\item $\varphi_{130} :\equiv \mathrm{ContIn}(f,A,c) \to \mathrm{Seq}(a) \to \forall n \in \mathbb{N}\,(a_{n} \in A) \to \mathrm{Conv}(a,c) \to \varphi_{74}$
\item $\varphi_{131} :\equiv (\lvert f \rvert(c) = \lvert f(c) \rvert \to \lvert f(c) \rvert = \lvert f \rvert(c)) \to \mathrm{ContIn}(f,A,c) \to \mathrm{Seq}(a) \to \varphi_{36}$
\item $\varphi_{132} :\equiv \mathrm{ContIn}(f,A,c) \to \mathrm{Seq}(a) \to \forall n \in \mathbb{N}\,(a_{n} \in A) \to \mathrm{Conv}(a,c) \to \varphi_{78}$
\item $\varphi_{133} :\equiv \mathrm{Seq}(\lvert f \rvert \circ a) \to \forall n \in \mathbb{N}\,(\lvert f \circ a \rvert(n) = \lvert f \rvert \circ a(n)) \to \varphi_{53}$
\item $\varphi_{134} :\equiv \mathrm{ContIn}(f,A,c) \leftrightarrow \mathrm{RFn}(f) \land A \subseteq \mathrm{dom}(f) \land c \in A \land \forall \varepsilon \in \mathbb{R}\,(0 < \varepsilon \to \exists \delta\,(\delta \in \mathbb{R} \land 0 < \delta \land \forall x \in A\,(\lvert x - c \rvert < \delta \to \lvert f(x) - f(c) \rvert < \varepsilon)))$
\item $\varphi_{135} :\equiv (A \subseteq \mathrm{dom}(f) \to c \in A \to c \in \mathrm{dom}(f)) \to (\mathrm{ContIn}(f,A,c) \to A \subseteq \mathrm{dom}(f)) \to \varphi_{95}$
\item $\varphi_{136} :\equiv (\mathrm{ContIn}(f,A,c) \to \mathrm{Seq}(a) \to \varphi_{30}) \to \mathrm{ContIn}(f,A,c) \to \varphi_{26}$
\item $\varphi_{137} :\equiv (\mathrm{ContIn}(f,A,c) \to \mathrm{Seq}(a) \to \varphi_{12}) \to \mathrm{ContIn}(f,A,c) \to \mathrm{Seq}(a) \to \varphi_{35}$
\item $\varphi_{138} :\equiv (\mathrm{ContIn}(f,A,c) \to \mathrm{Seq}(a) \to \varphi_{15}) \to \mathrm{ContIn}(f,A,c) \to \mathrm{Seq}(a) \to \varphi_{25}$
\item $\varphi_{139} :\equiv \forall a\,(\mathrm{ContIn}(f,A,c) \to \mathrm{Seq}(a) \to \varphi_{30}) \to \mathrm{ContIn}(f,A,c) \to \mathrm{Seq}(a) \to \varphi_{30}$
\item $\varphi_{140} :\equiv (\mathrm{ContIn}(f,A,c) \to \mathrm{Seq}(a) \to \varphi_{25}) \to \mathrm{ContIn}(f,A,c) \to \mathrm{Seq}(a) \to \varphi_{44}$
\item $\varphi_{141} :\equiv (\mathrm{ContIn}(f,A,c) \to \mathrm{Seq}(a) \to \varphi_{44}) \to \varphi_{17} \to n \in \mathbb{N} \to a_{n} \in \mathrm{dom}(f)$
\item $\varphi_{142} :\equiv \forall c\,\forall a\,(\mathrm{ContIn}(f,A,c) \to \mathrm{Seq}(a) \to \varphi_{30}) \to \forall a\,(\mathrm{ContIn}(f,A,c) \to \mathrm{Seq}(a) \to \varphi_{30})$
\item $\varphi_{143} :\equiv (\mathrm{ContIn}(f,A,c) \to \mathrm{Seq}(a) \to \varphi_{30}) \to \mathrm{ContIn}(f,A,c) \to \mathrm{Seq}(a) \to \varphi_{60}$
\item $\varphi_{144} :\equiv (\mathrm{ContIn}(f,A,c) \to \mathrm{Seq}(a) \to \varphi_{15}) \to \mathrm{ContIn}(f,A,c) \to \mathrm{Seq}(a) \to \varphi_{77}$
\item $\varphi_{145} :\equiv \forall A\,\forall c\,\forall a\,(\mathrm{ContIn}(f,A,c) \to \mathrm{Seq}(a) \to \varphi_{30}) \to \forall c\,\forall a\,(\mathrm{ContIn}(f,A,c) \to \mathrm{Seq}(a) \to \varphi_{30})$
\item $\varphi_{146} :\equiv \forall n\,(\varphi_{17} \to n \in \mathbb{N} \to a_{n} \in \mathrm{dom}(f)) \to \forall n\,(\varphi_{17}) \to \forall n \in \mathbb{N}\,(a_{n} \in \mathrm{dom}(f))$
\item $\varphi_{147} :\equiv (\mathrm{ContIn}(f,A,c) \to A \subseteq \mathrm{dom}(\lvert f \rvert)) \to (\mathrm{ContIn}(f,A,c) \to c \in A) \to \varphi_{126}$
\item $\varphi_{148} :\equiv (\mathrm{ContIn}(f,A,c) \to \mathrm{Seq}(a) \to \varphi_{56}) \to \varphi_{17} \to n \in \mathbb{N} \to a_{n} \in \mathrm{dom}(\lvert f \rvert)$
\item $\varphi_{149} :\equiv \varphi_{97} \to \forall A\,\forall c\,\forall a\,(\mathrm{ContIn}(f,A,c) \to \mathrm{Seq}(a) \to \varphi_{30})$
\item $\varphi_{150} :\equiv (\mathrm{ContIn}(f,A,c) \to \mathrm{Seq}(a) \to \varphi_{15}) \to \mathrm{ContIn}(f,A,c) \to \mathrm{Seq}(a) \to \varphi_{85}$
\item $\varphi_{151} :\equiv (\mathrm{ContIn}(f,A,c) \to \mathrm{Seq}(a) \to \varphi_{60}) \to \mathrm{ContIn}(f,A,c) \to \mathrm{Seq}(a) \to \varphi_{62}$
\item $\varphi_{152} :\equiv (\lvert f(c) \rvert = \lvert f \rvert(c) \to \varphi_{55}) \to \mathrm{ContIn}(f,A,c) \to \mathrm{Seq}(a) \to \varphi_{65}$
\item $\varphi_{153} :\equiv (\mathrm{ContIn}(f,A,c) \to \varphi_{16}) \to \varphi_{5} \to \mathrm{ContIn}(f,A,c) \to \mathrm{Seq}(a) \to \varphi_{12}$
\item $\varphi_{154} :\equiv \forall n\,(\varphi_{17} \to n \in \mathbb{N} \to a_{n} \in \mathrm{dom}(\lvert f \rvert)) \to \forall n\,(\varphi_{17}) \to \forall n \in \mathbb{N}\,(a_{n} \in \mathrm{dom}(\lvert f \rvert))$
\item $\varphi_{155} :\equiv \forall a\,(\mathrm{ContIn}(f,A,c) \to \mathrm{Seq}(a) \to \varphi_{65}) \to \forall a\,(\mathrm{ContIn}(f,A,c)) \to \forall a\,(\mathrm{Seq}(a) \to \varphi_{65})$
\item $\varphi_{156} :\equiv (\mathrm{ContIn}(f,A,c) \to \mathrm{ContIn}(f,A,c)) \to \varphi_{134} \to \mathrm{ContIn}(f,A,c) \to \mathrm{RFn}(f)$
\item $\varphi_{157} :\equiv (\mathrm{ContIn}(f,A,c) \to \mathrm{ContIn}(f,A,c)) \to \varphi_{134} \to \mathrm{ContIn}(f,A,c) \to c \in A$
\item $\varphi_{158} :\equiv (\mathrm{ContIn}(f,A,c) \to \mathrm{ContIn}(f,A,c)) \to \varphi_{134} \to \mathrm{ContIn}(f,A,c) \to A \subseteq \mathrm{dom}(f)$
\item $\varphi_{159} :\equiv (\varphi_{17} \to \forall n\,(\varphi_{17})) \to \mathrm{ContIn}(f,A,c) \to \mathrm{Seq}(a) \to \varphi_{51}$
\item $\varphi_{160} :\equiv (\varphi_{17} \to \forall n\,(\varphi_{17})) \to \mathrm{ContIn}(f,A,c) \to \mathrm{Seq}(a) \to \varphi_{66}$
\item $\varphi_{161} :\equiv \forall a\,(\mathrm{RFn}(f) \to \mathrm{Seq}(a) \to \varphi_{96}) \to \mathrm{RFn}(f) \to \mathrm{Seq}(a) \to \varphi_{96}$
\item $\varphi_{162} :\equiv \forall f\,\forall a\,(\mathrm{RFn}(f) \to \mathrm{Seq}(a) \to \varphi_{96}) \to \forall a\,(\mathrm{RFn}(f) \to \mathrm{Seq}(a) \to \varphi_{96})$
\item $\varphi_{163} :\equiv (\mathrm{ContIn}(f,A,c) \to \mathrm{Seq}(a) \to \varphi_{44}) \to \varphi_{13} \to \mathrm{ContIn}(f,A,c) \to \mathrm{Seq}(a) \to \varphi_{56}$
\item $\varphi_{164} :\equiv \forall n\,(\varphi_{17} \to n \in \mathbb{N} \to \lvert f \circ a \rvert(n) = \lvert f \rvert \circ a(n)) \to \forall n\,(\varphi_{17}) \to \forall n \in \mathbb{N}\,(\lvert f \circ a \rvert(n) = \lvert f \rvert \circ a(n))$
\item $\varphi_{165} :\equiv (\mathrm{ContIn}(f,A,c) \to \mathrm{Seq}(a) \to \mathrm{Seq}(a)) \to (\mathrm{ContIn}(f,A,c) \to \mathrm{Seq}(a) \to \varphi_{51}) \to \varphi_{117}$
\item $\varphi_{166} :\equiv (A \subseteq \mathrm{dom}(f) \to a_{n} \in A \to a_{n} \in \mathrm{dom}(f)) \to (\mathrm{ContIn}(f,A,c) \to A \subseteq \mathrm{dom}(f)) \to \varphi_{140}$
\item $\varphi_{167} :\equiv (\mathrm{ContIn}(f,A,c) \to \mathrm{Seq}(a) \to \varphi_{66}) \to \varphi_{54} \to \mathrm{ContIn}(f,A,c) \to \mathrm{Seq}(a) \to \varphi_{56}$
\item $\varphi_{168} :\equiv (\mathrm{ContIn}(f,A,c) \to \mathrm{Seq}(a) \to \mathrm{Seq}(a)) \to (\mathrm{ContIn}(f,A,c) \to \mathrm{Seq}(a) \to \varphi_{66}) \to \varphi_{130}$
\item $\varphi_{169} :\equiv \lvert f \circ a(n) \rvert = \lvert f \circ a(n) \rvert \to (f \circ a(n) = f(a_{n}) \to \varphi_{84}) \to \varphi_{103}$
\item $\varphi_{170} :\equiv \forall c\,(\mathrm{RFn}(\lvert f \rvert) \to A \subseteq \mathrm{dom}(\lvert f \rvert) \to \varphi_{107}) \to \mathrm{RFn}(\lvert f \rvert) \to A \subseteq \mathrm{dom}(\lvert f \rvert) \to \varphi_{107}$
\item $\varphi_{171} :\equiv \forall A\,\forall c\,(\mathrm{RFn}(\lvert f \rvert) \to A \subseteq \mathrm{dom}(\lvert f \rvert) \to \varphi_{107}) \to \forall c\,(\mathrm{RFn}(\lvert f \rvert) \to A \subseteq \mathrm{dom}(\lvert f \rvert) \to \varphi_{107})$
\item $\varphi_{172} :\equiv (\mathrm{Seq}(f \circ a) \to \varphi_{78}) \to (\mathrm{ContIn}(f,A,c) \to \varphi_{26}) \to \varphi_{132}$
\item $\varphi_{173} :\equiv (\mathrm{ContIn}(f,A,c) \to \mathrm{Seq}(a) \to \varphi_{85}) \to (\lvert f \rvert \circ a(n) = \lvert f \rvert(a_{n}) \to \varphi_{79}) \to \varphi_{106}$
\item $\varphi_{174} :\equiv (\mathrm{RFn}(\lvert f \rvert) \to A \subseteq \mathrm{dom}(\lvert f \rvert) \to \varphi_{107}) \to (\mathrm{ContIn}(f,A,c) \to \mathrm{RFn}(\lvert f \rvert)) \to \varphi_{147}$
\item $\varphi_{175} :\equiv (\mathrm{RFn}(f) \to \mathrm{Seq}(a) \to \varphi_{96}) \to (\mathrm{ContIn}(f,A,c) \to \mathrm{RFn}(f)) \to \varphi_{165}$
\item $\varphi_{176} :\equiv (\mathrm{ContIn}(f,A,c) \to \mathrm{Seq}(a) \to \varphi_{36}) \to (\mathrm{ContIn}(f,A,c) \to \mathrm{Seq}(a) \to \varphi_{62}) \to \varphi_{152}$
\item $\varphi_{177} :\equiv \forall a\,(\mathrm{ContIn}(f,A,c) \to \mathrm{Seq}(a) \to \varphi_{65}) \to \varphi_{155} \to \varphi_{1} \to \mathrm{ContIn}(f,A,c) \to \forall a\,(\mathrm{Seq}(a) \to \varphi_{65})$
\item $\varphi_{178} :\equiv \forall n\,(\varphi_{17} \to n \in \mathbb{N} \to \lvert f \circ a \rvert(n) = \lvert f \rvert \circ a(n)) \to \varphi_{164} \to (\varphi_{17} \to \forall n\,(\varphi_{17})) \to \varphi_{115}$
\end{abbrevlist}
\end{fullproof}

% 自動生成 (prover/build.js): ContInDiv  ノード数 770
\begin{fullproof}{定理2.3(商)}{$\forall f\,\forall g\,\forall A\,\forall c\,(\mathrm{ContIn}(f,A,c) \to \mathrm{ContIn}(g,A,c) \to \forall x \in A\,(g(x) \neq 0) \to \mathrm{ContIn}((f / g),A,c))$}{770}\phantomsection\label{tree:ContInDiv}
\noindent$\Pi_{1}$:
\begin{ptree}
\AXC{$z = w \to \varphi_{13}$\lbl{A6}}
\UIC{$\forall w\,(z = w \to \varphi_{13})$}
\UIC{$\forall z\,\forall w\,(z = w \to \varphi_{13})$}
\AXC{$\forall z\,\forall w\,(z = w \to \varphi_{13}) \to \varphi_{137}$\lbl{A2}}
\BIC{$\varphi_{137}$}
\AXC{$\varphi_{137} \to f(c) \cdot g(c)^{-1} = (f / g)(c) \to \varphi_{113}$\lbl{A2}}
\BIC{$f(c) \cdot g(c)^{-1} = (f / g)(c) \to \varphi_{113}$}
\end{ptree}
\noindent$\Pi_{2}$:
\begin{ptree}
\AXC{$x = x$\lbl{A5}}
\AXC{$x = y \to x = x \to y = x$\lbl{A6}}
\AXC{$(x = y \to x = x \to y = x) \to x = x \to x = y \to y = x$\lbl{A1}}
\BIC{$x = x \to x = y \to y = x$}
\BIC{$x = y \to y = x$}
\UIC{$\forall y\,(x = y \to y = x)$}
\UIC{$\forall x\,\forall y\,(x = y \to y = x)$}
\end{ptree}
\noindent$\Pi_{3}$:
\begin{ptree}
\AXC{$\forall n \in \mathbb{N}\,(a_{n} \in A) \to n \in \mathbb{N} \to a_{n} \in A$\lbl{A2}}
\AXC{$\mathrm{ContIn}(f,A,c) \to \varphi_{159}$\lbl{A1}}
\AXC{$(\mathrm{ContIn}(f,A,c) \to \varphi_{159}) \to \varphi_{260}$\lbl{A1}}
\BIC{$\varphi_{260}$}
\BIC{$\mathrm{ContIn}(f,A,c) \to \varphi_{146}$}
\end{ptree}
\noindent$\Pi_{4}$:
\begin{ptree}
\AXC{$\forall f\,\forall A\,\forall c\,(\varphi_{235})$\lbl{DefContIn}}
\AXC{$\forall f\,\forall A\,\forall c\,(\varphi_{235}) \to \forall A\,\forall c\,(\varphi_{234})$\lbl{A2}}
\BIC{$\forall A\,\forall c\,(\varphi_{234})$}
\AXC{$\forall A\,\forall c\,(\varphi_{234}) \to \forall c\,(\varphi_{234})$\lbl{A2}}
\BIC{$\forall c\,(\varphi_{234})$}
\AXC{$\forall c\,(\varphi_{234}) \to \varphi_{234}$\lbl{A2}}
\BIC{$\varphi_{234}$}
\end{ptree}
\noindent$\Pi_{5}$:
\begin{ptree}
\AXC{$\forall X\,\forall Y\,\forall x\,(X \subseteq Y \to x \in X \to x \in Y)$\lbl{[B5]}}
\AXC{$\varphi_{54}$\lbl{A2}}
\BIC{$\forall Y\,\forall x\,(A \subseteq Y \to x \in A \to x \in Y)$}
\AXC{$\varphi_{66}$\lbl{A2}}
\BIC{$\forall x\,(A \subseteq \mathrm{dom}(g) \to x \in A \to x \in \mathrm{dom}(g))$}
\AXC{$\varphi_{107}$\lbl{A2}}
\BIC{$A \subseteq \mathrm{dom}(g) \to a_{n} \in A \to a_{n} \in \mathrm{dom}(g)$}
\end{ptree}
\noindent$\Pi_{6}$:
\begin{ptree}
\AXC{$\stackrel{\Pi_{5}}{\vdots}$}\noLine
\UIC{$A \subseteq \mathrm{dom}(g) \to a_{n} \in A \to a_{n} \in \mathrm{dom}(g)$}
\AXC{$(A \subseteq \mathrm{dom}(g) \to a_{n} \in A \to a_{n} \in \mathrm{dom}(g)) \to \varphi_{158} \to \varphi_{304}$\lbl{A1}}
\BIC{$\varphi_{158} \to \varphi_{304}$}
\end{ptree}
\noindent$\Pi_{7}$:
\begin{ptree}
\AXC{$\stackrel{\Pi_{3}}{\vdots}$}\noLine
\UIC{$\mathrm{ContIn}(f,A,c) \to \varphi_{146}$}
\AXC{$\stackrel{\Pi_{4}}{\vdots}$}\noLine
\UIC{$\varphi_{234}$}
\AXC{$\varphi_{1}$\lbl{A1}}
\AXC{$\varphi_{1} \to \varphi_{234} \to \varphi_{158}$\lbl{A1}}
\BIC{$\varphi_{234} \to \varphi_{158}$}
\BIC{$\varphi_{158}$}
\AXC{$\stackrel{\Pi_{6}}{\vdots}$}\noLine
\UIC{$\varphi_{158} \to \varphi_{304}$}
\BIC{$\varphi_{304}$}
\BIC{$\mathrm{ContIn}(f,A,c) \to \varphi_{164}$}
\end{ptree}
\noindent$\Pi_{8}$:
\begin{ptree}
\AXC{$\stackrel{\Pi_{7}}{\vdots}$}\noLine
\UIC{$\mathrm{ContIn}(f,A,c) \to \varphi_{164}$}
\AXC{$(\mathrm{ContIn}(f,A,c) \to \varphi_{164}) \to \varphi_{193}$\lbl{A1}}
\BIC{$\varphi_{193}$}
\UIC{$\forall n\,(\varphi_{193})$}
\AXC{$\forall n\,(\varphi_{193}) \to \varphi_{389}$\lbl{A1}}
\BIC{$\varphi_{389}$}
\end{ptree}
\noindent$\Pi_{9}$:
\begin{ptree}
\AXC{$\forall f\,\forall A\,\forall c\,(\varphi_{235})$\lbl{DefContIn}}
\AXC{$\forall f\,\forall A\,\forall c\,(\varphi_{235}) \to \forall A\,\forall c\,(\varphi_{234})$\lbl{A2}}
\BIC{$\forall A\,\forall c\,(\varphi_{234})$}
\AXC{$\forall A\,\forall c\,(\varphi_{234}) \to \forall c\,(\varphi_{234})$\lbl{A2}}
\BIC{$\forall c\,(\varphi_{234})$}
\AXC{$\forall c\,(\varphi_{234}) \to \varphi_{234}$\lbl{A2}}
\BIC{$\varphi_{234}$}
\end{ptree}
\noindent$\Pi_{10}$:
\begin{ptree}
\AXC{$\forall f\,\forall a\,(\mathrm{RFn}(f) \to \mathrm{Seq}(a) \to \varphi_{127})$\lbl{[S8]}}
\AXC{$\varphi_{287}$\lbl{A2}}
\BIC{$\forall a\,(\mathrm{RFn}(g) \to \mathrm{Seq}(a) \to \varphi_{126})$}
\AXC{$\varphi_{284}$\lbl{A2}}
\BIC{$\mathrm{RFn}(g) \to \mathrm{Seq}(a) \to \varphi_{126}$}
\AXC{$(\mathrm{RFn}(g) \to \mathrm{Seq}(a) \to \varphi_{126}) \to \varphi_{142} \to \varphi_{141} \to \varphi_{326}$\lbl{A1}}
\BIC{$\varphi_{142} \to \varphi_{141} \to \varphi_{326}$}
\end{ptree}
\noindent$\Pi_{11}$:
\begin{ptree}
\AXC{$\varphi_{141}$\lbl{A1}}
\AXC{$\stackrel{\Pi_{9}}{\vdots}$}\noLine
\UIC{$\varphi_{234}$}
\AXC{$\varphi_{1}$\lbl{A1}}
\AXC{$\varphi_{1} \to \varphi_{234} \to \varphi_{142}$\lbl{A1}}
\BIC{$\varphi_{234} \to \varphi_{142}$}
\BIC{$\varphi_{142}$}
\AXC{$\stackrel{\Pi_{10}}{\vdots}$}\noLine
\UIC{$\varphi_{142} \to \varphi_{141} \to \varphi_{326}$}
\BIC{$\varphi_{141} \to \varphi_{326}$}
\BIC{$\varphi_{326}$}
\end{ptree}
\noindent$\Pi_{12}$:
\begin{ptree}
\AXC{$\varphi_{119}$\lbl{A2}}
\AXC{$\varphi_{278}$\lbl{A4}}
\AXC{$\forall n\,(\varphi_{193}) \to \varphi_{131} \to \forall n \in \mathbb{N}\,(a_{n} \in \mathrm{dom}(g))$\lbl{A3}}
\AXC{$\stackrel{\Pi_{8}}{\vdots}$}\noLine
\UIC{$\varphi_{389}$}
\BIC{$\varphi_{278} \to \mathrm{ContIn}(f,A,c) \to \varphi_{170}$}
\BIC{$\mathrm{ContIn}(f,A,c) \to \varphi_{170}$}
\AXC{$\stackrel{\Pi_{11}}{\vdots}$}\noLine
\UIC{$\varphi_{326}$}
\BIC{$\mathrm{ContIn}(f,A,c) \to \mathrm{ContIn}(g,A,c) \to \varphi_{203}$}
\AXC{$\varphi_{336}$\lbl{A1}}
\BIC{$\mathrm{ContIn}(f,A,c) \to \mathrm{ContIn}(g,A,c) \to \varphi_{191}$}
\AXC{$\varphi_{358}$\lbl{A1}}
\BIC{$\varphi_{291}$}
\BIC{$\mathrm{ContIn}(f,A,c) \to \mathrm{ContIn}(g,A,c) \to \varphi_{187}$}
\end{ptree}
\noindent$\Pi_{13}$:
\begin{ptree}
\AXC{$\mathrm{ContIn}(f,A,c) \to \varphi_{140}$\lbl{A1}}
\AXC{$\stackrel{\Pi_{12}}{\vdots}$}\noLine
\UIC{$\mathrm{ContIn}(f,A,c) \to \mathrm{ContIn}(g,A,c) \to \varphi_{187}$}
\AXC{$\varphi_{363}$\lbl{A1}}
\BIC{$\varphi_{308}$}
\BIC{$\mathrm{ContIn}(f,A,c) \to \varphi_{196}$}
\end{ptree}
\noindent$\Pi_{14}$:
\begin{ptree}
\AXC{$\stackrel{\Pi_{13}}{\vdots}$}\noLine
\UIC{$\mathrm{ContIn}(f,A,c) \to \varphi_{196}$}
\AXC{$(\mathrm{ContIn}(f,A,c) \to \varphi_{196}) \to \varphi_{280}$\lbl{A1}}
\BIC{$\varphi_{280}$}
\end{ptree}
\noindent$\Pi_{15}$:
\begin{ptree}
\AXC{$\stackrel{\Pi_{2}}{\vdots}$}\noLine
\UIC{$\forall x\,\forall y\,(x = y \to y = x)$}
\AXC{$\varphi_{35}$\lbl{A2}}
\BIC{$\forall y\,(g \circ a(n) = y \to y = g \circ a(n))$}
\AXC{$\varphi_{154}$\lbl{A2}}
\BIC{$g \circ a(n) = g(a_{n}) \to g(a_{n}) = g \circ a(n)$}
\AXC{$\stackrel{\Pi_{14}}{\vdots}$}\noLine
\UIC{$\varphi_{280}$}
\BIC{$\mathrm{ContIn}(f,A,c) \to \varphi_{197}$}
\end{ptree}
\noindent$\Pi_{16}$:
\begin{ptree}
\AXC{$\stackrel{\Pi_{15}}{\vdots}$}\noLine
\UIC{$\mathrm{ContIn}(f,A,c) \to \varphi_{197}$}
\AXC{$(\mathrm{ContIn}(f,A,c) \to \varphi_{197}) \to \varphi_{281}$\lbl{A1}}
\BIC{$\varphi_{281}$}
\end{ptree}
\noindent$\Pi_{17}$:
\begin{ptree}
\AXC{$\forall n \in \mathbb{N}\,(a_{n} \in A) \to n \in \mathbb{N} \to a_{n} \in A$\lbl{A2}}
\AXC{$\mathrm{ContIn}(f,A,c) \to \varphi_{159}$\lbl{A1}}
\AXC{$(\mathrm{ContIn}(f,A,c) \to \varphi_{159}) \to \varphi_{260}$\lbl{A1}}
\BIC{$\varphi_{260}$}
\BIC{$\mathrm{ContIn}(f,A,c) \to \varphi_{146}$}
\end{ptree}
\noindent$\Pi_{18}$:
\begin{ptree}
\AXC{$\forall x \in A\,(g(x) \neq 0) \to a_{n} \in A \to g(a_{n}) \neq 0$\lbl{A2}}
\AXC{$\mathrm{ContIn}(f,A,c) \to \varphi_{38}$\lbl{A1}}
\AXC{$(\mathrm{ContIn}(f,A,c) \to \varphi_{38}) \to \varphi_{218}$\lbl{A1}}
\BIC{$\varphi_{218}$}
\BIC{$\mathrm{ContIn}(f,A,c) \to \varphi_{70}$}
\end{ptree}
\noindent$\Pi_{19}$:
\begin{ptree}
\AXC{$\mathrm{ContIn}(f,A,c) \to \varphi_{140}$\lbl{A1}}
\AXC{$\stackrel{\Pi_{17}}{\vdots}$}\noLine
\UIC{$\mathrm{ContIn}(f,A,c) \to \varphi_{146}$}
\AXC{$\stackrel{\Pi_{18}}{\vdots}$}\noLine
\UIC{$\mathrm{ContIn}(f,A,c) \to \varphi_{70}$}
\AXC{$\varphi_{377}$\lbl{A1}}
\BIC{$(\mathrm{ContIn}(f,A,c) \to \varphi_{146}) \to \varphi_{298}$}
\BIC{$\varphi_{298}$}
\BIC{$\mathrm{ContIn}(f,A,c) \to \varphi_{161}$}
\end{ptree}
\noindent$\Pi_{20}$:
\begin{ptree}
\AXC{$\stackrel{\Pi_{19}}{\vdots}$}\noLine
\UIC{$\mathrm{ContIn}(f,A,c) \to \varphi_{161}$}
\AXC{$\stackrel{\Pi_{15}}{\vdots}$}\noLine
\UIC{$\mathrm{ContIn}(f,A,c) \to \varphi_{197}$}
\AXC{$\varphi_{399}$\lbl{A1}}
\BIC{$(\mathrm{ContIn}(f,A,c) \to \varphi_{161}) \to \varphi_{366}$}
\BIC{$\varphi_{366}$}
\end{ptree}
\noindent$\Pi_{21}$:
\begin{ptree}
\AXC{$z = w \to z = 0 \to w = 0$\lbl{A6}}
\UIC{$\forall w\,(z = w \to z = 0 \to w = 0)$}
\UIC{$\forall z\,\forall w\,(z = w \to z = 0 \to w = 0)$}
\AXC{$\varphi_{43}$\lbl{A2}}
\BIC{$\forall w\,(g(a_{n}) = w \to g(a_{n}) = 0 \to w = 0)$}
\AXC{$\varphi_{167}$\lbl{A2}}
\BIC{$g(a_{n}) = g \circ a(n) \to g(a_{n}) = 0 \to g \circ a(n) = 0$}
\AXC{$\stackrel{\Pi_{20}}{\vdots}$}\noLine
\UIC{$\varphi_{366}$}
\BIC{$\varphi_{340}$}
\end{ptree}
\noindent$\Pi_{22}$:
\begin{ptree}
\AXC{$z = w \to z = 0 \to w = 0$\lbl{A6}}
\UIC{$\forall w\,(z = w \to z = 0 \to w = 0)$}
\UIC{$\forall z\,\forall w\,(z = w \to z = 0 \to w = 0)$}
\AXC{$\varphi_{67}$\lbl{A2}}
\BIC{$\forall w\,(g \circ a(n) = w \to g \circ a(n) = 0 \to w = 0)$}
\AXC{$\varphi_{182}$\lbl{A2}}
\BIC{$g \circ a(n) = g(a_{n}) \to g \circ a(n) = 0 \to g(a_{n}) = 0$}
\AXC{$\stackrel{\Pi_{21}}{\vdots}$}\noLine
\UIC{$\varphi_{340}$}
\BIC{$\varphi_{313}$}
\end{ptree}
\noindent$\Pi_{23}$:
\begin{ptree}
\AXC{$\stackrel{\Pi_{2}}{\vdots}$}\noLine
\UIC{$\forall x\,\forall y\,(x = y \to y = x)$}
\AXC{$\varphi_{11}$\lbl{A2}}
\BIC{$\forall y\,(g(a_{n}) = y \to y = g(a_{n}))$}
\AXC{$\varphi_{138}$\lbl{A2}}
\BIC{$g(a_{n}) = g \circ a(n) \to g \circ a(n) = g(a_{n})$}
\AXC{$\stackrel{\Pi_{16}}{\vdots}$}\noLine
\UIC{$\varphi_{281}$}
\BIC{$\mathrm{ContIn}(f,A,c) \to \varphi_{196}$}
\AXC{$\stackrel{\Pi_{22}}{\vdots}$}\noLine
\UIC{$\varphi_{313}$}
\BIC{$\mathrm{ContIn}(f,A,c) \to \varphi_{171}$}
\AXC{$(\mathrm{ContIn}(f,A,c) \to \varphi_{171}) \to \varphi_{201}$\lbl{A1}}
\BIC{$\varphi_{201}$}
\UIC{$\forall n\,(\varphi_{201})$}
\end{ptree}
\noindent$\Pi_{24}$:
\begin{ptree}
\AXC{$\varphi_{278}$\lbl{A4}}
\AXC{$\forall n\,(\varphi_{201}) \to \varphi_{131} \to \forall n \in \mathbb{N}\,(g \circ a(n) \neq 0)$\lbl{A3}}
\AXC{$\stackrel{\Pi_{23}}{\vdots}$}\noLine
\UIC{$\forall n\,(\varphi_{201})$}
\AXC{$\forall n\,(\varphi_{201}) \to \varphi_{391}$\lbl{A1}}
\BIC{$\varphi_{391}$}
\BIC{$\varphi_{278} \to \mathrm{ContIn}(f,A,c) \to \varphi_{178}$}
\BIC{$\mathrm{ContIn}(f,A,c) \to \varphi_{178}$}
\end{ptree}
\noindent$\Pi_{25}$:
\begin{ptree}
\AXC{$\forall f\,\forall A\,\forall c\,(\varphi_{235})$\lbl{DefContIn}}
\AXC{$\forall f\,\forall A\,\forall c\,(\varphi_{235}) \to \forall A\,\forall c\,(\varphi_{234})$\lbl{A2}}
\BIC{$\forall A\,\forall c\,(\varphi_{234})$}
\AXC{$\forall A\,\forall c\,(\varphi_{234}) \to \forall c\,(\varphi_{234})$\lbl{A2}}
\BIC{$\forall c\,(\varphi_{234})$}
\AXC{$\forall c\,(\varphi_{234}) \to \varphi_{234}$\lbl{A2}}
\BIC{$\varphi_{234}$}
\end{ptree}
\noindent$\Pi_{26}$:
\begin{ptree}
\AXC{$\forall x \in A\,(g(x) \neq 0) \to c \in A \to g(c) \neq 0$\lbl{A2}}
\AXC{$\mathrm{ContIn}(f,A,c) \to \varphi_{38}$\lbl{A1}}
\AXC{$(\mathrm{ContIn}(f,A,c) \to \varphi_{38}) \to \varphi_{173}$\lbl{A1}}
\BIC{$\varphi_{173}$}
\BIC{$\mathrm{ContIn}(f,A,c) \to \varphi_{20}$}
\end{ptree}
\noindent$\Pi_{27}$:
\begin{ptree}
\AXC{$\stackrel{\Pi_{25}}{\vdots}$}\noLine
\UIC{$\varphi_{234}$}
\AXC{$\varphi_{1}$\lbl{A1}}
\AXC{$\varphi_{1} \to \varphi_{234} \to \varphi_{143}$\lbl{A1}}
\BIC{$\varphi_{234} \to \varphi_{143}$}
\BIC{$\varphi_{143}$}
\AXC{$\stackrel{\Pi_{26}}{\vdots}$}\noLine
\UIC{$\mathrm{ContIn}(f,A,c) \to \varphi_{20}$}
\AXC{$(\mathrm{ContIn}(f,A,c) \to \varphi_{20}) \to \varphi_{143} \to \varphi_{147}$\lbl{A1}}
\BIC{$\varphi_{143} \to \varphi_{147}$}
\BIC{$\varphi_{147}$}
\end{ptree}
\noindent$\Pi_{28}$:
\begin{ptree}
\AXC{$\varphi_{1}$\lbl{A1}}
\AXC{$\varphi_{129}$\lbl{[G1]}}
\AXC{$\varphi_{274}$\lbl{A2}}
\BIC{$\forall A\,\forall c\,\forall a\,(\mathrm{ContIn}(g,A,c) \to \mathrm{Seq}(a) \to \varphi_{48})$}
\AXC{$\varphi_{270}$\lbl{A2}}
\BIC{$\forall c\,\forall a\,(\mathrm{ContIn}(g,A,c) \to \mathrm{Seq}(a) \to \varphi_{48})$}
\AXC{$\varphi_{264}$\lbl{A2}}
\BIC{$\forall a\,(\mathrm{ContIn}(g,A,c) \to \mathrm{Seq}(a) \to \varphi_{48})$}
\AXC{$\varphi_{257}$\lbl{A2}}
\BIC{$\mathrm{ContIn}(g,A,c) \to \mathrm{Seq}(a) \to \varphi_{48}$}
\AXC{$\varphi_{385}$\lbl{A1}}
\BIC{$\varphi_{1} \to (\mathrm{ContIn}(f,A,c) \to \varphi_{12}) \to \varphi_{342}$}
\BIC{$(\mathrm{ContIn}(f,A,c) \to \varphi_{12}) \to \varphi_{342}$}
\end{ptree}
\noindent$\Pi_{29}$:
\begin{ptree}
\AXC{$\mathrm{ContIn}(f,A,c) \to \varphi_{12}$\lbl{A1}}
\AXC{$\stackrel{\Pi_{28}}{\vdots}$}\noLine
\UIC{$(\mathrm{ContIn}(f,A,c) \to \varphi_{12}) \to \varphi_{342}$}
\BIC{$\varphi_{342}$}
\end{ptree}
\noindent$\Pi_{30}$:
\begin{ptree}
\AXC{$\mathrm{ContIn}(f,A,c) \to \varphi_{125}$\lbl{A1}}
\AXC{$\mathrm{ContIn}(f,A,c) \to \mathrm{ContIn}(g,A,c) \to \varphi_{111}$\lbl{A1}}
\AXC{$\stackrel{\Pi_{29}}{\vdots}$}\noLine
\UIC{$\varphi_{342}$}
\BIC{$\varphi_{289}$}
\BIC{$\mathrm{ContIn}(f,A,c) \to \varphi_{150}$}
\end{ptree}
\noindent$\Pi_{31}$:
\begin{ptree}
\AXC{$\varphi_{129}$\lbl{[G1]}}
\AXC{$\varphi_{275}$\lbl{A2}}
\BIC{$\forall A\,\forall c\,\forall a\,(\mathrm{ContIn}(f,A,c) \to \mathrm{Seq}(a) \to \varphi_{49})$}
\AXC{$\varphi_{271}$\lbl{A2}}
\BIC{$\forall c\,\forall a\,(\mathrm{ContIn}(f,A,c) \to \mathrm{Seq}(a) \to \varphi_{49})$}
\AXC{$\varphi_{265}$\lbl{A2}}
\BIC{$\forall a\,(\mathrm{ContIn}(f,A,c) \to \mathrm{Seq}(a) \to \varphi_{49})$}
\AXC{$\varphi_{258}$\lbl{A2}}
\BIC{$\mathrm{ContIn}(f,A,c) \to \mathrm{Seq}(a) \to \varphi_{49}$}
\end{ptree}
\noindent$\Pi_{32}$:
\begin{ptree}
\AXC{$\stackrel{\Pi_{31}}{\vdots}$}\noLine
\UIC{$\mathrm{ContIn}(f,A,c) \to \mathrm{Seq}(a) \to \varphi_{49}$}
\AXC{$(\mathrm{ContIn}(f,A,c) \to \mathrm{Seq}(a) \to \varphi_{49}) \to \varphi_{373}$\lbl{A1}}
\BIC{$\varphi_{373}$}
\end{ptree}
\noindent$\Pi_{33}$:
\begin{ptree}
\AXC{$\mathrm{ContIn}(f,A,c) \to \varphi_{12}$\lbl{A1}}
\AXC{$\mathrm{ContIn}(f,A,c) \to \mathrm{ContIn}(f,A,c)$\lbl{A1}}
\AXC{$\stackrel{\Pi_{32}}{\vdots}$}\noLine
\UIC{$\varphi_{373}$}
\BIC{$(\mathrm{ContIn}(f,A,c) \to \varphi_{12}) \to \varphi_{343}$}
\BIC{$\varphi_{343}$}
\end{ptree}
\noindent$\Pi_{34}$:
\begin{ptree}
\AXC{$\mathrm{ContIn}(f,A,c) \to \varphi_{125}$\lbl{A1}}
\AXC{$\mathrm{ContIn}(f,A,c) \to \mathrm{ContIn}(g,A,c) \to \varphi_{111}$\lbl{A1}}
\AXC{$\stackrel{\Pi_{33}}{\vdots}$}\noLine
\UIC{$\varphi_{343}$}
\BIC{$\varphi_{290}$}
\BIC{$\mathrm{ContIn}(f,A,c) \to \varphi_{151}$}
\end{ptree}
\noindent$\Pi_{35}$:
\begin{ptree}
\AXC{$\varphi_{136}$\lbl{定理2.9}}
\AXC{$\varphi_{136} \to \varphi_{180}$\lbl{A2}}
\BIC{$\varphi_{180}$}
\AXC{$\varphi_{180} \to \varphi_{190}$\lbl{A2}}
\BIC{$\varphi_{190}$}
\AXC{$\varphi_{190} \to \varphi_{228}$\lbl{A2}}
\BIC{$\varphi_{228}$}
\AXC{$\varphi_{228} \to \varphi_{248}$\lbl{A2}}
\BIC{$\varphi_{248}$}
\AXC{$\varphi_{248} \to \varphi_{396}$\lbl{A1}}
\BIC{$\varphi_{396}$}
\end{ptree}
\noindent$\Pi_{36}$:
\begin{ptree}
\AXC{$\stackrel{\Pi_{24}}{\vdots}$}\noLine
\UIC{$\mathrm{ContIn}(f,A,c) \to \varphi_{178}$}
\AXC{$\stackrel{\Pi_{27}}{\vdots}$}\noLine
\UIC{$\varphi_{147}$}
\AXC{$\stackrel{\Pi_{30}}{\vdots}$}\noLine
\UIC{$\mathrm{ContIn}(f,A,c) \to \varphi_{150}$}
\AXC{$\stackrel{\Pi_{34}}{\vdots}$}\noLine
\UIC{$\mathrm{ContIn}(f,A,c) \to \varphi_{151}$}
\AXC{$\stackrel{\Pi_{35}}{\vdots}$}\noLine
\UIC{$\varphi_{396}$}
\BIC{$(\mathrm{ContIn}(f,A,c) \to \varphi_{150}) \to \varphi_{147} \to \varphi_{328}$}
\BIC{$\varphi_{147} \to \varphi_{328}$}
\BIC{$\varphi_{328}$}
\BIC{$\mathrm{ContIn}(f,A,c) \to \mathrm{ContIn}(g,A,c) \to \varphi_{210}$}
\end{ptree}
\noindent$\Pi_{37}$:
\begin{ptree}
\AXC{$y = z \to x = y \to x = z$\lbl{A6}}
\AXC{$(y = z \to x = y \to x = z) \to x = y \to y = z \to x = z$\lbl{A1}}
\BIC{$x = y \to y = z \to x = z$}
\UIC{$\forall z\,(x = y \to y = z \to x = z)$}
\UIC{$\forall y\,\forall z\,(x = y \to y = z \to x = z)$}
\UIC{$\forall x\,\forall y\,\forall z\,(x = y \to y = z \to x = z)$}
\AXC{$\forall x\,\forall y\,\forall z\,(x = y \to y = z \to x = z) \to \varphi_{51}$\lbl{A2}}
\BIC{$\varphi_{51}$}
\end{ptree}
\noindent$\Pi_{38}$:
\begin{ptree}
\AXC{$\stackrel{\Pi_{2}}{\vdots}$}\noLine
\UIC{$\forall x\,\forall y\,(x = y \to y = x)$}
\AXC{$\forall x\,\forall y\,(x = y \to y = x) \to \varphi_{124}$\lbl{A2}}
\BIC{$\varphi_{124}$}
\AXC{$\varphi_{124} \to \varphi_{239}$\lbl{A2}}
\BIC{$\varphi_{239}$}
\end{ptree}
\noindent$\Pi_{39}$:
\begin{ptree}
\AXC{$y = z \to x = y \to x = z$\lbl{A6}}
\AXC{$(y = z \to x = y \to x = z) \to x = y \to y = z \to x = z$\lbl{A1}}
\BIC{$x = y \to y = z \to x = z$}
\UIC{$\forall z\,(x = y \to y = z \to x = z)$}
\UIC{$\forall y\,\forall z\,(x = y \to y = z \to x = z)$}
\UIC{$\forall x\,\forall y\,\forall z\,(x = y \to y = z \to x = z)$}
\AXC{$\forall x\,\forall y\,\forall z\,(x = y \to y = z \to x = z) \to \varphi_{134}$\lbl{A2}}
\BIC{$\varphi_{134}$}
\end{ptree}
\noindent$\Pi_{40}$:
\begin{ptree}
\AXC{$\forall n \in \mathbb{N}\,(a_{n} \in A) \to n \in \mathbb{N} \to a_{n} \in A$\lbl{A2}}
\AXC{$\mathrm{ContIn}(f,A,c) \to \mathrm{ContIn}(g,A,c) \to \varphi_{111}$\lbl{A1}}
\AXC{$\varphi_{319}$\lbl{A1}}
\BIC{$\varphi_{250}$}
\BIC{$\mathrm{ContIn}(f,A,c) \to \mathrm{ContIn}(g,A,c) \to \varphi_{101}$}
\AXC{$\varphi_{349}$\lbl{A1}}
\BIC{$\varphi_{292}$}
\end{ptree}
\noindent$\Pi_{41}$:
\begin{ptree}
\AXC{$\forall f\,\forall A\,\forall c\,(\varphi_{235})$\lbl{DefContIn}}
\AXC{$\forall f\,\forall A\,\forall c\,(\varphi_{235}) \to \forall A\,\forall c\,(\varphi_{234})$\lbl{A2}}
\BIC{$\forall A\,\forall c\,(\varphi_{234})$}
\AXC{$\forall A\,\forall c\,(\varphi_{234}) \to \forall c\,(\varphi_{234})$\lbl{A2}}
\BIC{$\forall c\,(\varphi_{234})$}
\AXC{$\forall c\,(\varphi_{234}) \to \varphi_{234}$\lbl{A2}}
\BIC{$\varphi_{234}$}
\end{ptree}
\noindent$\Pi_{42}$:
\begin{ptree}
\AXC{$\stackrel{\Pi_{41}}{\vdots}$}\noLine
\UIC{$\varphi_{234}$}
\AXC{$\varphi_{1}$\lbl{A1}}
\AXC{$\varphi_{1} \to \varphi_{234} \to \mathrm{ContIn}(f,A,c) \to \varphi_{7}$\lbl{A1}}
\BIC{$\varphi_{234} \to \mathrm{ContIn}(f,A,c) \to \varphi_{7}$}
\BIC{$\mathrm{ContIn}(f,A,c) \to \varphi_{7}$}
\end{ptree}
\noindent$\Pi_{43}$:
\begin{ptree}
\AXC{$\forall X\,\forall Y\,\forall x\,(X \subseteq Y \to x \in X \to x \in Y)$\lbl{[B5]}}
\AXC{$\varphi_{54}$\lbl{A2}}
\BIC{$\forall Y\,\forall x\,(A \subseteq Y \to x \in A \to x \in Y)$}
\AXC{$\varphi_{66}$\lbl{A2}}
\BIC{$\forall x\,(A \subseteq \mathrm{dom}(g) \to x \in A \to x \in \mathrm{dom}(g))$}
\AXC{$\varphi_{107}$\lbl{A2}}
\BIC{$A \subseteq \mathrm{dom}(g) \to a_{n} \in A \to a_{n} \in \mathrm{dom}(g)$}
\AXC{$\varphi_{351}$\lbl{A1}}
\BIC{$(\mathrm{ContIn}(f,A,c) \to \varphi_{7}) \to \varphi_{304}$}
\end{ptree}
\noindent$\Pi_{44}$:
\begin{ptree}
\AXC{$\mathrm{ContIn}(f,A,c) \to \varphi_{140}$\lbl{A1}}
\AXC{$\stackrel{\Pi_{40}}{\vdots}$}\noLine
\UIC{$\varphi_{292}$}
\BIC{$\mathrm{ContIn}(f,A,c) \to \varphi_{146}$}
\AXC{$\stackrel{\Pi_{42}}{\vdots}$}\noLine
\UIC{$\mathrm{ContIn}(f,A,c) \to \varphi_{7}$}
\AXC{$\stackrel{\Pi_{43}}{\vdots}$}\noLine
\UIC{$(\mathrm{ContIn}(f,A,c) \to \varphi_{7}) \to \varphi_{304}$}
\BIC{$\varphi_{304}$}
\BIC{$\mathrm{ContIn}(f,A,c) \to \varphi_{164}$}
\end{ptree}
\noindent$\Pi_{45}$:
\begin{ptree}
\AXC{$\stackrel{\Pi_{44}}{\vdots}$}\noLine
\UIC{$\mathrm{ContIn}(f,A,c) \to \varphi_{164}$}
\AXC{$(\mathrm{ContIn}(f,A,c) \to \varphi_{164}) \to \varphi_{193}$\lbl{A1}}
\BIC{$\varphi_{193}$}
\UIC{$\forall n\,(\varphi_{193})$}
\AXC{$\forall n\,(\varphi_{193}) \to \varphi_{389}$\lbl{A1}}
\BIC{$\varphi_{389}$}
\end{ptree}
\noindent$\Pi_{46}$:
\begin{ptree}
\AXC{$\forall f\,\forall A\,\forall c\,(\varphi_{235})$\lbl{DefContIn}}
\AXC{$\forall f\,\forall A\,\forall c\,(\varphi_{235}) \to \forall A\,\forall c\,(\varphi_{234})$\lbl{A2}}
\BIC{$\forall A\,\forall c\,(\varphi_{234})$}
\AXC{$\forall A\,\forall c\,(\varphi_{234}) \to \forall c\,(\varphi_{234})$\lbl{A2}}
\BIC{$\forall c\,(\varphi_{234})$}
\AXC{$\forall c\,(\varphi_{234}) \to \varphi_{234}$\lbl{A2}}
\BIC{$\varphi_{234}$}
\end{ptree}
\noindent$\Pi_{47}$:
\begin{ptree}
\AXC{$\stackrel{\Pi_{46}}{\vdots}$}\noLine
\UIC{$\varphi_{234}$}
\AXC{$\varphi_{1}$\lbl{A1}}
\AXC{$\varphi_{1} \to \varphi_{234} \to \mathrm{ContIn}(f,A,c) \to \varphi_{2}$\lbl{A1}}
\BIC{$\varphi_{234} \to \mathrm{ContIn}(f,A,c) \to \varphi_{2}$}
\BIC{$\mathrm{ContIn}(f,A,c) \to \varphi_{2}$}
\end{ptree}
\noindent$\Pi_{48}$:
\begin{ptree}
\AXC{$\forall f\,\forall a\,(\mathrm{RFn}(f) \to \mathrm{Seq}(a) \to \varphi_{127})$\lbl{[S8]}}
\AXC{$\varphi_{287}$\lbl{A2}}
\BIC{$\forall a\,(\mathrm{RFn}(g) \to \mathrm{Seq}(a) \to \varphi_{126})$}
\AXC{$\varphi_{284}$\lbl{A2}}
\BIC{$\mathrm{RFn}(g) \to \mathrm{Seq}(a) \to \varphi_{126}$}
\AXC{$(\mathrm{RFn}(g) \to \mathrm{Seq}(a) \to \varphi_{126}) \to \varphi_{367}$\lbl{A1}}
\BIC{$\varphi_{367}$}
\end{ptree}
\noindent$\Pi_{49}$:
\begin{ptree}
\AXC{$\mathrm{ContIn}(f,A,c) \to \varphi_{12}$\lbl{A1}}
\AXC{$\stackrel{\Pi_{47}}{\vdots}$}\noLine
\UIC{$\mathrm{ContIn}(f,A,c) \to \varphi_{2}$}
\AXC{$\stackrel{\Pi_{48}}{\vdots}$}\noLine
\UIC{$\varphi_{367}$}
\BIC{$(\mathrm{ContIn}(f,A,c) \to \varphi_{12}) \to \varphi_{326}$}
\BIC{$\varphi_{326}$}
\end{ptree}
\noindent$\Pi_{50}$:
\begin{ptree}
\AXC{$\varphi_{119}$\lbl{A2}}
\AXC{$\varphi_{278}$\lbl{A4}}
\AXC{$\forall n\,(\varphi_{193}) \to \varphi_{131} \to \forall n \in \mathbb{N}\,(a_{n} \in \mathrm{dom}(g))$\lbl{A3}}
\AXC{$\stackrel{\Pi_{45}}{\vdots}$}\noLine
\UIC{$\varphi_{389}$}
\BIC{$\varphi_{278} \to \mathrm{ContIn}(f,A,c) \to \varphi_{170}$}
\BIC{$\mathrm{ContIn}(f,A,c) \to \varphi_{170}$}
\AXC{$\stackrel{\Pi_{49}}{\vdots}$}\noLine
\UIC{$\varphi_{326}$}
\BIC{$\mathrm{ContIn}(f,A,c) \to \mathrm{ContIn}(g,A,c) \to \varphi_{203}$}
\AXC{$\varphi_{336}$\lbl{A1}}
\BIC{$\mathrm{ContIn}(f,A,c) \to \mathrm{ContIn}(g,A,c) \to \varphi_{191}$}
\AXC{$\varphi_{358}$\lbl{A1}}
\BIC{$\varphi_{291}$}
\BIC{$\mathrm{ContIn}(f,A,c) \to \mathrm{ContIn}(g,A,c) \to \varphi_{187}$}
\end{ptree}
\noindent$\Pi_{51}$:
\begin{ptree}
\AXC{$\mathrm{ContIn}(f,A,c) \to \varphi_{140}$\lbl{A1}}
\AXC{$\stackrel{\Pi_{50}}{\vdots}$}\noLine
\UIC{$\mathrm{ContIn}(f,A,c) \to \mathrm{ContIn}(g,A,c) \to \varphi_{187}$}
\AXC{$\varphi_{363}$\lbl{A1}}
\BIC{$\varphi_{308}$}
\BIC{$\mathrm{ContIn}(f,A,c) \to \varphi_{196}$}
\end{ptree}
\noindent$\Pi_{52}$:
\begin{ptree}
\AXC{$x = x$\lbl{A5}}
\UIC{$\forall x\,(x = x)$}
\AXC{$\varphi_{108}$\lbl{A2}}
\BIC{$f(a_{n}) \cdot g \circ a(n)^{-1} = f(a_{n}) \cdot g \circ a(n)^{-1}$}
\AXC{$\stackrel{\Pi_{51}}{\vdots}$}\noLine
\UIC{$\mathrm{ContIn}(f,A,c) \to \varphi_{196}$}
\AXC{$(\mathrm{ContIn}(f,A,c) \to \varphi_{196}) \to \varphi_{374}$\lbl{A1}}
\BIC{$\varphi_{374}$}
\BIC{$\varphi_{353}$}
\end{ptree}
\noindent$\Pi_{53}$:
\begin{ptree}
\AXC{$z = w \to \varphi_{189}$\lbl{A6}}
\UIC{$\forall w\,(z = w \to \varphi_{189})$}
\UIC{$\forall z\,\forall w\,(z = w \to \varphi_{189})$}
\AXC{$\forall z\,\forall w\,(z = w \to \varphi_{189}) \to \varphi_{249}$\lbl{A2}}
\BIC{$\varphi_{249}$}
\AXC{$\varphi_{249} \to g \circ a(n) = g(a_{n}) \to \varphi_{230}$\lbl{A2}}
\BIC{$g \circ a(n) = g(a_{n}) \to \varphi_{230}$}
\AXC{$\stackrel{\Pi_{52}}{\vdots}$}\noLine
\UIC{$\varphi_{353}$}
\BIC{$\mathrm{ContIn}(f,A,c) \to \mathrm{ContIn}(g,A,c) \to \varphi_{237}$}
\end{ptree}
\noindent$\Pi_{54}$:
\begin{ptree}
\AXC{$y = z \to x = y \to x = z$\lbl{A6}}
\AXC{$(y = z \to x = y \to x = z) \to x = y \to y = z \to x = z$\lbl{A1}}
\BIC{$x = y \to y = z \to x = z$}
\UIC{$\forall z\,(x = y \to y = z \to x = z)$}
\UIC{$\forall y\,\forall z\,(x = y \to y = z \to x = z)$}
\UIC{$\forall x\,\forall y\,\forall z\,(x = y \to y = z \to x = z)$}
\AXC{$\forall x\,\forall y\,\forall z\,(x = y \to y = z \to x = z) \to \varphi_{134}$\lbl{A2}}
\BIC{$\varphi_{134}$}
\end{ptree}
\noindent$\Pi_{55}$:
\begin{ptree}
\AXC{$\forall n \in \mathbb{N}\,(a_{n} \in A) \to n \in \mathbb{N} \to a_{n} \in A$\lbl{A2}}
\AXC{$\mathrm{ContIn}(f,A,c) \to \mathrm{ContIn}(g,A,c) \to \varphi_{111}$\lbl{A1}}
\AXC{$\varphi_{319}$\lbl{A1}}
\BIC{$\varphi_{250}$}
\BIC{$\mathrm{ContIn}(f,A,c) \to \mathrm{ContIn}(g,A,c) \to \varphi_{101}$}
\AXC{$\varphi_{349}$\lbl{A1}}
\BIC{$\varphi_{292}$}
\end{ptree}
\noindent$\Pi_{56}$:
\begin{ptree}
\AXC{$\forall f\,\forall A\,\forall c\,(\varphi_{235})$\lbl{DefContIn}}
\AXC{$\forall f\,\forall A\,\forall c\,(\varphi_{235}) \to \forall A\,\forall c\,(\varphi_{235})$\lbl{A2}}
\BIC{$\forall A\,\forall c\,(\varphi_{235})$}
\AXC{$\forall A\,\forall c\,(\varphi_{235}) \to \forall c\,(\varphi_{235})$\lbl{A2}}
\BIC{$\forall c\,(\varphi_{235})$}
\AXC{$\forall c\,(\varphi_{235}) \to \varphi_{235}$\lbl{A2}}
\BIC{$\varphi_{235}$}
\end{ptree}
\noindent$\Pi_{57}$:
\begin{ptree}
\AXC{$\stackrel{\Pi_{56}}{\vdots}$}\noLine
\UIC{$\varphi_{235}$}
\AXC{$\mathrm{ContIn}(f,A,c) \to \mathrm{ContIn}(f,A,c)$\lbl{A1}}
\AXC{$\varphi_{300}$\lbl{A1}}
\BIC{$\varphi_{235} \to \mathrm{ContIn}(f,A,c) \to \varphi_{8}$}
\BIC{$\mathrm{ContIn}(f,A,c) \to \varphi_{8}$}
\end{ptree}
\noindent$\Pi_{58}$:
\begin{ptree}
\AXC{$\forall X\,\forall Y\,\forall x\,(X \subseteq Y \to x \in X \to x \in Y)$\lbl{[B5]}}
\AXC{$\varphi_{54}$\lbl{A2}}
\BIC{$\forall Y\,\forall x\,(A \subseteq Y \to x \in A \to x \in Y)$}
\AXC{$\varphi_{68}$\lbl{A2}}
\BIC{$\forall x\,(A \subseteq \mathrm{dom}(f) \to x \in A \to x \in \mathrm{dom}(f))$}
\AXC{$\varphi_{109}$\lbl{A2}}
\BIC{$A \subseteq \mathrm{dom}(f) \to a_{n} \in A \to a_{n} \in \mathrm{dom}(f)$}
\AXC{$\varphi_{352}$\lbl{A1}}
\BIC{$(\mathrm{ContIn}(f,A,c) \to \varphi_{8}) \to \varphi_{305}$}
\end{ptree}
\noindent$\Pi_{59}$:
\begin{ptree}
\AXC{$\mathrm{ContIn}(f,A,c) \to \varphi_{140}$\lbl{A1}}
\AXC{$\stackrel{\Pi_{55}}{\vdots}$}\noLine
\UIC{$\varphi_{292}$}
\BIC{$\mathrm{ContIn}(f,A,c) \to \varphi_{146}$}
\AXC{$\stackrel{\Pi_{57}}{\vdots}$}\noLine
\UIC{$\mathrm{ContIn}(f,A,c) \to \varphi_{8}$}
\AXC{$\stackrel{\Pi_{58}}{\vdots}$}\noLine
\UIC{$(\mathrm{ContIn}(f,A,c) \to \varphi_{8}) \to \varphi_{305}$}
\BIC{$\varphi_{305}$}
\BIC{$\mathrm{ContIn}(f,A,c) \to \varphi_{165}$}
\end{ptree}
\noindent$\Pi_{60}$:
\begin{ptree}
\AXC{$\stackrel{\Pi_{59}}{\vdots}$}\noLine
\UIC{$\mathrm{ContIn}(f,A,c) \to \varphi_{165}$}
\AXC{$(\mathrm{ContIn}(f,A,c) \to \varphi_{165}) \to \varphi_{194}$\lbl{A1}}
\BIC{$\varphi_{194}$}
\UIC{$\forall n\,(\varphi_{194})$}
\AXC{$\forall n\,(\varphi_{194}) \to \varphi_{390}$\lbl{A1}}
\BIC{$\varphi_{390}$}
\end{ptree}
\noindent$\Pi_{61}$:
\begin{ptree}
\AXC{$\forall f\,\forall A\,\forall c\,(\varphi_{235})$\lbl{DefContIn}}
\AXC{$\forall f\,\forall A\,\forall c\,(\varphi_{235}) \to \forall A\,\forall c\,(\varphi_{235})$\lbl{A2}}
\BIC{$\forall A\,\forall c\,(\varphi_{235})$}
\AXC{$\forall A\,\forall c\,(\varphi_{235}) \to \forall c\,(\varphi_{235})$\lbl{A2}}
\BIC{$\forall c\,(\varphi_{235})$}
\AXC{$\forall c\,(\varphi_{235}) \to \varphi_{235}$\lbl{A2}}
\BIC{$\varphi_{235}$}
\end{ptree}
\noindent$\Pi_{62}$:
\begin{ptree}
\AXC{$\stackrel{\Pi_{61}}{\vdots}$}\noLine
\UIC{$\varphi_{235}$}
\AXC{$\mathrm{ContIn}(f,A,c) \to \mathrm{ContIn}(f,A,c)$\lbl{A1}}
\AXC{$\varphi_{294}$\lbl{A1}}
\BIC{$\varphi_{235} \to \mathrm{ContIn}(f,A,c) \to \varphi_{3}$}
\BIC{$\mathrm{ContIn}(f,A,c) \to \varphi_{3}$}
\end{ptree}
\noindent$\Pi_{63}$:
\begin{ptree}
\AXC{$\forall f\,\forall a\,(\mathrm{RFn}(f) \to \mathrm{Seq}(a) \to \varphi_{127})$\lbl{[S8]}}
\AXC{$\varphi_{288}$\lbl{A2}}
\BIC{$\forall a\,(\mathrm{RFn}(f) \to \mathrm{Seq}(a) \to \varphi_{127})$}
\AXC{$\varphi_{285}$\lbl{A2}}
\BIC{$\mathrm{RFn}(f) \to \mathrm{Seq}(a) \to \varphi_{127}$}
\AXC{$(\mathrm{RFn}(f) \to \mathrm{Seq}(a) \to \varphi_{127}) \to \varphi_{368}$\lbl{A1}}
\BIC{$\varphi_{368}$}
\end{ptree}
\noindent$\Pi_{64}$:
\begin{ptree}
\AXC{$\mathrm{ContIn}(f,A,c) \to \varphi_{12}$\lbl{A1}}
\AXC{$\stackrel{\Pi_{62}}{\vdots}$}\noLine
\UIC{$\mathrm{ContIn}(f,A,c) \to \varphi_{3}$}
\AXC{$\stackrel{\Pi_{63}}{\vdots}$}\noLine
\UIC{$\varphi_{368}$}
\BIC{$(\mathrm{ContIn}(f,A,c) \to \varphi_{12}) \to \varphi_{327}$}
\BIC{$\varphi_{327}$}
\end{ptree}
\noindent$\Pi_{65}$:
\begin{ptree}
\AXC{$\varphi_{120}$\lbl{A2}}
\AXC{$\varphi_{278}$\lbl{A4}}
\AXC{$\forall n\,(\varphi_{194}) \to \varphi_{131} \to \forall n \in \mathbb{N}\,(a_{n} \in \mathrm{dom}(f))$\lbl{A3}}
\AXC{$\stackrel{\Pi_{60}}{\vdots}$}\noLine
\UIC{$\varphi_{390}$}
\BIC{$\varphi_{278} \to \mathrm{ContIn}(f,A,c) \to \varphi_{172}$}
\BIC{$\mathrm{ContIn}(f,A,c) \to \varphi_{172}$}
\AXC{$\stackrel{\Pi_{64}}{\vdots}$}\noLine
\UIC{$\varphi_{327}$}
\BIC{$\mathrm{ContIn}(f,A,c) \to \mathrm{ContIn}(g,A,c) \to \varphi_{204}$}
\AXC{$\varphi_{337}$\lbl{A1}}
\BIC{$\mathrm{ContIn}(f,A,c) \to \mathrm{ContIn}(g,A,c) \to \varphi_{192}$}
\AXC{$\varphi_{359}$\lbl{A1}}
\BIC{$\varphi_{293}$}
\BIC{$\mathrm{ContIn}(f,A,c) \to \mathrm{ContIn}(g,A,c) \to \varphi_{188}$}
\end{ptree}
\noindent$\Pi_{66}$:
\begin{ptree}
\AXC{$\mathrm{ContIn}(f,A,c) \to \varphi_{140}$\lbl{A1}}
\AXC{$\stackrel{\Pi_{65}}{\vdots}$}\noLine
\UIC{$\mathrm{ContIn}(f,A,c) \to \mathrm{ContIn}(g,A,c) \to \varphi_{188}$}
\AXC{$\varphi_{364}$\lbl{A1}}
\BIC{$\varphi_{309}$}
\BIC{$\mathrm{ContIn}(f,A,c) \to \varphi_{198}$}
\end{ptree}
\noindent$\Pi_{67}$:
\begin{ptree}
\AXC{$x = x$\lbl{A5}}
\UIC{$\forall x\,(x = x)$}
\AXC{$\forall x\,(x = x) \to \varphi_{103}$\lbl{A2}}
\BIC{$\varphi_{103}$}
\AXC{$\stackrel{\Pi_{66}}{\vdots}$}\noLine
\UIC{$\mathrm{ContIn}(f,A,c) \to \varphi_{198}$}
\AXC{$(\mathrm{ContIn}(f,A,c) \to \varphi_{198}) \to \varphi_{103} \to \varphi_{356}$\lbl{A1}}
\BIC{$\varphi_{103} \to \varphi_{356}$}
\BIC{$\varphi_{356}$}
\end{ptree}
\noindent$\Pi_{68}$:
\begin{ptree}
\AXC{$z = w \to \varphi_{215}$\lbl{A6}}
\UIC{$\forall w\,(z = w \to \varphi_{215})$}
\UIC{$\forall z\,\forall w\,(z = w \to \varphi_{215})$}
\AXC{$\forall z\,\forall w\,(z = w \to \varphi_{215}) \to \varphi_{263}$\lbl{A2}}
\BIC{$\varphi_{263}$}
\AXC{$\varphi_{263} \to f \circ a(n) = f(a_{n}) \to \varphi_{252}$\lbl{A2}}
\BIC{$f \circ a(n) = f(a_{n}) \to \varphi_{252}$}
\AXC{$\stackrel{\Pi_{67}}{\vdots}$}\noLine
\UIC{$\varphi_{356}$}
\BIC{$\mathrm{ContIn}(f,A,c) \to \mathrm{ContIn}(g,A,c) \to \varphi_{245}$}
\end{ptree}
\noindent$\Pi_{69}$:
\begin{ptree}
\AXC{$y = z \to x = y \to x = z$\lbl{A6}}
\AXC{$(y = z \to x = y \to x = z) \to x = y \to y = z \to x = z$\lbl{A1}}
\BIC{$x = y \to y = z \to x = z$}
\UIC{$\forall z\,(x = y \to y = z \to x = z)$}
\UIC{$\forall y\,\forall z\,(x = y \to y = z \to x = z)$}
\UIC{$\forall x\,\forall y\,\forall z\,(x = y \to y = z \to x = z)$}
\AXC{$\forall x\,\forall y\,\forall z\,(x = y \to y = z \to x = z) \to \varphi_{134}$\lbl{A2}}
\BIC{$\varphi_{134}$}
\end{ptree}
\noindent$\Pi_{70}$:
\begin{ptree}
\AXC{$z = w \to \varphi_{207}$\lbl{A6}}
\UIC{$\forall w\,(z = w \to \varphi_{207})$}
\UIC{$\forall z\,\forall w\,(z = w \to \varphi_{207})$}
\AXC{$\forall z\,\forall w\,(z = w \to \varphi_{207}) \to \varphi_{272}$\lbl{A2}}
\BIC{$\varphi_{272}$}
\AXC{$\varphi_{272} \to g \circ a^{\ast}(n) = g \circ a(n)^{-1} \to \varphi_{114} \to \varphi_{112}$\lbl{A2}}
\BIC{$g \circ a^{\ast}(n) = g \circ a(n)^{-1} \to \varphi_{114} \to \varphi_{112}$}
\end{ptree}
\noindent$\Pi_{71}$:
\begin{ptree}
\AXC{$\stackrel{\Pi_{30}}{\vdots}$}\noLine
\UIC{$\mathrm{ContIn}(f,A,c) \to \varphi_{150}$}
\AXC{$\forall a\,\forall \alpha\,(\mathrm{Conv}(a,\alpha) \to \mathrm{Seq}(a))$\lbl{[C1]}}
\AXC{$\varphi_{55}$\lbl{A2}}
\BIC{$\forall \alpha\,(\mathrm{Conv}(g \circ a,\alpha) \to \mathrm{Seq}(g \circ a))$}
\AXC{$\varphi_{104}$\lbl{A2}}
\BIC{$\mathrm{Conv}(g \circ a,g(c)) \to \mathrm{Seq}(g \circ a)$}
\AXC{$\varphi_{320}$\lbl{A1}}
\BIC{$(\mathrm{ContIn}(f,A,c) \to \varphi_{150}) \to \varphi_{168}$}
\BIC{$\varphi_{168}$}
\end{ptree}
\noindent$\Pi_{72}$:
\begin{ptree}
\AXC{$\stackrel{\Pi_{24}}{\vdots}$}\noLine
\UIC{$\mathrm{ContIn}(f,A,c) \to \varphi_{178}$}
\AXC{$\stackrel{\Pi_{71}}{\vdots}$}\noLine
\UIC{$\varphi_{168}$}
\AXC{$\varphi_{121}$\lbl{[S5]}}
\AXC{$\varphi_{121} \to \mathrm{Seq}(g \circ a) \to \forall n \in \mathbb{N}\,(g \circ a(n) \neq 0) \to \varphi_{99}$\lbl{A2}}
\BIC{$\mathrm{Seq}(g \circ a) \to \forall n \in \mathbb{N}\,(g \circ a(n) \neq 0) \to \varphi_{99}$}
\AXC{$\varphi_{386}$\lbl{A1}}
\BIC{$\varphi_{168} \to \varphi_{334}$}
\BIC{$\varphi_{334}$}
\BIC{$\mathrm{ContIn}(f,A,c) \to \mathrm{ContIn}(g,A,c) \to \varphi_{231}$}
\end{ptree}
\noindent$\Pi_{73}$:
\begin{ptree}
\AXC{$\varphi_{175}$\lbl{A2}}
\AXC{$\stackrel{\Pi_{72}}{\vdots}$}\noLine
\UIC{$\mathrm{ContIn}(f,A,c) \to \mathrm{ContIn}(g,A,c) \to \varphi_{231}$}
\AXC{$\varphi_{345}$\lbl{A1}}
\BIC{$\mathrm{ContIn}(f,A,c) \to \mathrm{ContIn}(g,A,c) \to \varphi_{216}$}
\AXC{$\varphi_{375}$\lbl{A1}}
\BIC{$\varphi_{323}$}
\BIC{$\mathrm{ContIn}(f,A,c) \to \mathrm{ContIn}(g,A,c) \to \varphi_{212}$}
\AXC{$\varphi_{370}$\lbl{A1}}
\BIC{$\varphi_{316}$}
\end{ptree}
\noindent$\Pi_{74}$:
\begin{ptree}
\AXC{$x = x$\lbl{A5}}
\UIC{$\forall x\,(x = x)$}
\AXC{$\forall x\,(x = x) \to \varphi_{114}$\lbl{A2}}
\BIC{$\varphi_{114}$}
\AXC{$\mathrm{ContIn}(f,A,c) \to \varphi_{140}$\lbl{A1}}
\AXC{$\stackrel{\Pi_{73}}{\vdots}$}\noLine
\UIC{$\varphi_{316}$}
\BIC{$\mathrm{ContIn}(f,A,c) \to \mathrm{ContIn}(g,A,c) \to \varphi_{199}$}
\AXC{$\varphi_{398}$\lbl{A1}}
\BIC{$\varphi_{114} \to (g \circ a^{\ast}(n) = g \circ a(n)^{-1} \to \varphi_{114} \to \varphi_{112}) \to \varphi_{277}$}
\BIC{$(g \circ a^{\ast}(n) = g \circ a(n)^{-1} \to \varphi_{114} \to \varphi_{112}) \to \varphi_{277}$}
\end{ptree}
\noindent$\Pi_{75}$:
\begin{ptree}
\AXC{$\stackrel{\Pi_{34}}{\vdots}$}\noLine
\UIC{$\mathrm{ContIn}(f,A,c) \to \varphi_{151}$}
\AXC{$\forall a\,\forall \alpha\,(\mathrm{Conv}(a,\alpha) \to \mathrm{Seq}(a))$\lbl{[C1]}}
\AXC{$\varphi_{56}$\lbl{A2}}
\BIC{$\forall \alpha\,(\mathrm{Conv}(f \circ a,\alpha) \to \mathrm{Seq}(f \circ a))$}
\AXC{$\varphi_{105}$\lbl{A2}}
\BIC{$\mathrm{Conv}(f \circ a,f(c)) \to \mathrm{Seq}(f \circ a)$}
\AXC{$\varphi_{321}$\lbl{A1}}
\BIC{$(\mathrm{ContIn}(f,A,c) \to \varphi_{151}) \to \varphi_{169}$}
\BIC{$\varphi_{169}$}
\end{ptree}
\noindent$\Pi_{76}$:
\begin{ptree}
\AXC{$\stackrel{\Pi_{75}}{\vdots}$}\noLine
\UIC{$\varphi_{169}$}
\AXC{$\varphi_{110}$\lbl{[S4]}}
\AXC{$\varphi_{110} \to \varphi_{183}$\lbl{A2}}
\BIC{$\varphi_{183}$}
\AXC{$\varphi_{183} \to \varphi_{269}$\lbl{A2}}
\BIC{$\varphi_{269}$}
\AXC{$\varphi_{269} \to \varphi_{169} \to \varphi_{350}$\lbl{A1}}
\BIC{$\varphi_{169} \to \varphi_{350}$}
\BIC{$\varphi_{350}$}
\end{ptree}
\noindent$\Pi_{77}$:
\begin{ptree}
\AXC{$\stackrel{\Pi_{72}}{\vdots}$}\noLine
\UIC{$\mathrm{ContIn}(f,A,c) \to \mathrm{ContIn}(g,A,c) \to \varphi_{231}$}
\AXC{$\varphi_{329}$\lbl{A1}}
\BIC{$\mathrm{ContIn}(f,A,c) \to \varphi_{144}$}
\AXC{$\stackrel{\Pi_{76}}{\vdots}$}\noLine
\UIC{$\varphi_{350}$}
\BIC{$\mathrm{ContIn}(f,A,c) \to \varphi_{286}$}
\AXC{$(\mathrm{ContIn}(f,A,c) \to \varphi_{286}) \to \varphi_{282}$\lbl{A1}}
\BIC{$\varphi_{282}$}
\end{ptree}
\noindent$\Pi_{78}$:
\begin{ptree}
\AXC{$\mathrm{ContIn}(f,A,c) \to \varphi_{140}$\lbl{A1}}
\AXC{$\varphi_{133} \to \varphi_{130}$\lbl{A2}}
\AXC{$\stackrel{\Pi_{77}}{\vdots}$}\noLine
\UIC{$\varphi_{282}$}
\AXC{$\varphi_{282} \to (\varphi_{133} \to \varphi_{130}) \to \varphi_{279}$\lbl{A1}}
\BIC{$(\varphi_{133} \to \varphi_{130}) \to \varphi_{279}$}
\BIC{$\varphi_{279}$}
\AXC{$\varphi_{279} \to \varphi_{335}$\lbl{A1}}
\BIC{$\varphi_{335}$}
\BIC{$\mathrm{ContIn}(f,A,c) \to \mathrm{ContIn}(g,A,c) \to \varphi_{253}$}
\AXC{$\varphi_{404}$\lbl{A1}}
\BIC{$\varphi_{277} \to \varphi_{379}$}
\end{ptree}
\noindent$\Pi_{79}$:
\begin{ptree}
\AXC{$\stackrel{\Pi_{70}}{\vdots}$}\noLine
\UIC{$g \circ a^{\ast}(n) = g \circ a(n)^{-1} \to \varphi_{114} \to \varphi_{112}$}
\AXC{$\stackrel{\Pi_{74}}{\vdots}$}\noLine
\UIC{$(g \circ a^{\ast}(n) = g \circ a(n)^{-1} \to \varphi_{114} \to \varphi_{112}) \to \varphi_{277}$}
\BIC{$\varphi_{277}$}
\AXC{$\stackrel{\Pi_{78}}{\vdots}$}\noLine
\UIC{$\varphi_{277} \to \varphi_{379}$}
\BIC{$\varphi_{379}$}
\end{ptree}
\noindent$\Pi_{80}$:
\begin{ptree}
\AXC{$\stackrel{\Pi_{68}}{\vdots}$}\noLine
\UIC{$\mathrm{ContIn}(f,A,c) \to \mathrm{ContIn}(g,A,c) \to \varphi_{245}$}
\AXC{$\stackrel{\Pi_{69}}{\vdots}$}\noLine
\UIC{$\varphi_{134}$}
\AXC{$\varphi_{134} \to \varphi_{267}$\lbl{A2}}
\BIC{$\varphi_{267}$}
\AXC{$\varphi_{267} \to \varphi_{310}$\lbl{A2}}
\BIC{$\varphi_{310}$}
\AXC{$\stackrel{\Pi_{79}}{\vdots}$}\noLine
\UIC{$\varphi_{379}$}
\BIC{$\mathrm{ContIn}(f,A,c) \to \mathrm{ContIn}(g,A,c) \to \varphi_{251}$}
\AXC{$\varphi_{402}$\lbl{A1}}
\BIC{$\varphi_{394}$}
\BIC{$\varphi_{371}$}
\end{ptree}
\noindent$\Pi_{81}$:
\begin{ptree}
\AXC{$\stackrel{\Pi_{53}}{\vdots}$}\noLine
\UIC{$\mathrm{ContIn}(f,A,c) \to \mathrm{ContIn}(g,A,c) \to \varphi_{237}$}
\AXC{$\stackrel{\Pi_{54}}{\vdots}$}\noLine
\UIC{$\varphi_{134}$}
\AXC{$\varphi_{134} \to \varphi_{262}$\lbl{A2}}
\BIC{$\varphi_{262}$}
\AXC{$\varphi_{262} \to \varphi_{302}$\lbl{A2}}
\BIC{$\varphi_{302}$}
\AXC{$\stackrel{\Pi_{80}}{\vdots}$}\noLine
\UIC{$\varphi_{371}$}
\BIC{$\mathrm{ContIn}(f,A,c) \to \mathrm{ContIn}(g,A,c) \to \varphi_{246}$}
\AXC{$\varphi_{401}$\lbl{A1}}
\BIC{$\varphi_{392}$}
\BIC{$\varphi_{369}$}
\end{ptree}
\noindent$\Pi_{82}$:
\begin{ptree}
\AXC{$y = z \to x = y \to x = z$\lbl{A6}}
\AXC{$(y = z \to x = y \to x = z) \to x = y \to y = z \to x = z$\lbl{A1}}
\BIC{$x = y \to y = z \to x = z$}
\UIC{$\forall z\,(x = y \to y = z \to x = z)$}
\UIC{$\forall y\,\forall z\,(x = y \to y = z \to x = z)$}
\UIC{$\forall x\,\forall y\,\forall z\,(x = y \to y = z \to x = z)$}
\AXC{$\forall x\,\forall y\,\forall z\,(x = y \to y = z \to x = z) \to \varphi_{51}$\lbl{A2}}
\BIC{$\varphi_{51}$}
\end{ptree}
\noindent$\Pi_{83}$:
\begin{ptree}
\AXC{$\forall f\,(\varphi_{238})$\lbl{[S20]}}
\AXC{$\forall f\,(\varphi_{238}) \to \varphi_{238}$\lbl{A2}}
\BIC{$\varphi_{238}$}
\AXC{$\varphi_{238} \to \mathrm{RFn}(f) \to \mathrm{RFn}(g) \to \mathrm{RFn}((f / g)) \land \varphi_{39} \land \varphi_{41}$\lbl{A2}}
\BIC{$\mathrm{RFn}(f) \to \mathrm{RFn}(g) \to \mathrm{RFn}((f / g)) \land \varphi_{39} \land \varphi_{41}$}
\AXC{$\varphi_{380}$\lbl{A1}}
\BIC{$(\mathrm{ContIn}(f,A,c) \to \varphi_{3}) \to \varphi_{301}$}
\end{ptree}
\noindent$\Pi_{84}$:
\begin{ptree}
\AXC{$\stackrel{\Pi_{47}}{\vdots}$}\noLine
\UIC{$\mathrm{ContIn}(f,A,c) \to \varphi_{2}$}
\AXC{$\stackrel{\Pi_{62}}{\vdots}$}\noLine
\UIC{$\mathrm{ContIn}(f,A,c) \to \varphi_{3}$}
\AXC{$\stackrel{\Pi_{83}}{\vdots}$}\noLine
\UIC{$(\mathrm{ContIn}(f,A,c) \to \varphi_{3}) \to \varphi_{301}$}
\BIC{$\varphi_{301}$}
\BIC{$\mathrm{ContIn}(f,A,c) \to \varphi_{259}$}
\end{ptree}
\noindent$\Pi_{85}$:
\begin{ptree}
\AXC{$\stackrel{\Pi_{84}}{\vdots}$}\noLine
\UIC{$\mathrm{ContIn}(f,A,c) \to \varphi_{259}$}
\AXC{$(\mathrm{ContIn}(f,A,c) \to \varphi_{259}) \to \varphi_{160}$\lbl{A1}}
\BIC{$\varphi_{160}$}
\end{ptree}
\noindent$\Pi_{86}$:
\begin{ptree}
\AXC{$\forall n \in \mathbb{N}\,(a_{n} \in A) \to n \in \mathbb{N} \to a_{n} \in A$\lbl{A2}}
\AXC{$\mathrm{ContIn}(f,A,c) \to \mathrm{ContIn}(g,A,c) \to \varphi_{111}$\lbl{A1}}
\AXC{$\varphi_{319}$\lbl{A1}}
\BIC{$\varphi_{250}$}
\BIC{$\mathrm{ContIn}(f,A,c) \to \mathrm{ContIn}(g,A,c) \to \varphi_{101}$}
\AXC{$\varphi_{349}$\lbl{A1}}
\BIC{$\varphi_{292}$}
\end{ptree}
\noindent$\Pi_{87}$:
\begin{ptree}
\AXC{$\mathrm{ContIn}(f,A,c) \to \varphi_{140}$\lbl{A1}}
\AXC{$\stackrel{\Pi_{86}}{\vdots}$}\noLine
\UIC{$\varphi_{292}$}
\BIC{$\mathrm{ContIn}(f,A,c) \to \varphi_{146}$}
\end{ptree}
\noindent$\Pi_{88}$:
\begin{ptree}
\AXC{$\forall x \in A\,(g(x) \neq 0) \to a_{n} \in A \to g(a_{n}) \neq 0$\lbl{A2}}
\AXC{$\mathrm{ContIn}(f,A,c) \to \varphi_{38}$\lbl{A1}}
\AXC{$(\mathrm{ContIn}(f,A,c) \to \varphi_{38}) \to \varphi_{218}$\lbl{A1}}
\BIC{$\varphi_{218}$}
\BIC{$\mathrm{ContIn}(f,A,c) \to \varphi_{70}$}
\end{ptree}
\noindent$\Pi_{89}$:
\begin{ptree}
\AXC{$\stackrel{\Pi_{87}}{\vdots}$}\noLine
\UIC{$\mathrm{ContIn}(f,A,c) \to \varphi_{146}$}
\AXC{$\stackrel{\Pi_{88}}{\vdots}$}\noLine
\UIC{$\mathrm{ContIn}(f,A,c) \to \varphi_{70}$}
\AXC{$(\mathrm{ContIn}(f,A,c) \to \varphi_{70}) \to \varphi_{303}$\lbl{A1}}
\BIC{$\varphi_{303}$}
\BIC{$\mathrm{ContIn}(f,A,c) \to \varphi_{161}$}
\end{ptree}
\noindent$\Pi_{90}$:
\begin{ptree}
\AXC{$\forall X\,\forall Y\,\forall x\,(X \subseteq Y \to x \in X \to x \in Y)$\lbl{[B5]}}
\AXC{$\varphi_{54}$\lbl{A2}}
\BIC{$\forall Y\,\forall x\,(A \subseteq Y \to x \in A \to x \in Y)$}
\AXC{$\varphi_{66}$\lbl{A2}}
\BIC{$\forall x\,(A \subseteq \mathrm{dom}(g) \to x \in A \to x \in \mathrm{dom}(g))$}
\AXC{$\varphi_{107}$\lbl{A2}}
\BIC{$A \subseteq \mathrm{dom}(g) \to a_{n} \in A \to a_{n} \in \mathrm{dom}(g)$}
\AXC{$\varphi_{351}$\lbl{A1}}
\BIC{$(\mathrm{ContIn}(f,A,c) \to \varphi_{7}) \to \varphi_{304}$}
\end{ptree}
\noindent$\Pi_{91}$:
\begin{ptree}
\AXC{$\stackrel{\Pi_{87}}{\vdots}$}\noLine
\UIC{$\mathrm{ContIn}(f,A,c) \to \varphi_{146}$}
\AXC{$\stackrel{\Pi_{42}}{\vdots}$}\noLine
\UIC{$\mathrm{ContIn}(f,A,c) \to \varphi_{7}$}
\AXC{$\stackrel{\Pi_{90}}{\vdots}$}\noLine
\UIC{$(\mathrm{ContIn}(f,A,c) \to \varphi_{7}) \to \varphi_{304}$}
\BIC{$\varphi_{304}$}
\BIC{$\mathrm{ContIn}(f,A,c) \to \varphi_{164}$}
\end{ptree}
\noindent$\Pi_{92}$:
\begin{ptree}
\AXC{$\forall X\,\forall Y\,\forall x\,(X \subseteq Y \to x \in X \to x \in Y)$\lbl{[B5]}}
\AXC{$\varphi_{54}$\lbl{A2}}
\BIC{$\forall Y\,\forall x\,(A \subseteq Y \to x \in A \to x \in Y)$}
\AXC{$\varphi_{68}$\lbl{A2}}
\BIC{$\forall x\,(A \subseteq \mathrm{dom}(f) \to x \in A \to x \in \mathrm{dom}(f))$}
\AXC{$\varphi_{109}$\lbl{A2}}
\BIC{$A \subseteq \mathrm{dom}(f) \to a_{n} \in A \to a_{n} \in \mathrm{dom}(f)$}
\AXC{$\varphi_{352}$\lbl{A1}}
\BIC{$(\mathrm{ContIn}(f,A,c) \to \varphi_{8}) \to \varphi_{305}$}
\end{ptree}
\noindent$\Pi_{93}$:
\begin{ptree}
\AXC{$\stackrel{\Pi_{87}}{\vdots}$}\noLine
\UIC{$\mathrm{ContIn}(f,A,c) \to \varphi_{146}$}
\AXC{$\stackrel{\Pi_{57}}{\vdots}$}\noLine
\UIC{$\mathrm{ContIn}(f,A,c) \to \varphi_{8}$}
\AXC{$\stackrel{\Pi_{92}}{\vdots}$}\noLine
\UIC{$(\mathrm{ContIn}(f,A,c) \to \varphi_{8}) \to \varphi_{305}$}
\BIC{$\varphi_{305}$}
\BIC{$\mathrm{ContIn}(f,A,c) \to \varphi_{165}$}
\AXC{$\varphi_{395}$\lbl{A1}}
\BIC{$(\mathrm{ContIn}(f,A,c) \to \varphi_{164}) \to \varphi_{357}$}
\end{ptree}
\noindent$\Pi_{94}$:
\begin{ptree}
\AXC{$\stackrel{\Pi_{89}}{\vdots}$}\noLine
\UIC{$\mathrm{ContIn}(f,A,c) \to \varphi_{161}$}
\AXC{$\stackrel{\Pi_{91}}{\vdots}$}\noLine
\UIC{$\mathrm{ContIn}(f,A,c) \to \varphi_{164}$}
\AXC{$\stackrel{\Pi_{93}}{\vdots}$}\noLine
\UIC{$(\mathrm{ContIn}(f,A,c) \to \varphi_{164}) \to \varphi_{357}$}
\BIC{$\varphi_{357}$}
\BIC{$\varphi_{205} \to \mathrm{ContIn}(f,A,c) \to \varphi_{181}$}
\end{ptree}
\noindent$\Pi_{95}$:
\begin{ptree}
\AXC{$\varphi_{39} \to \varphi_{98}$\lbl{A2}}
\AXC{$\stackrel{\Pi_{85}}{\vdots}$}\noLine
\UIC{$\varphi_{160}$}
\AXC{$\varphi_{160} \to (\varphi_{39} \to \varphi_{98}) \to \varphi_{205}$\lbl{A1}}
\BIC{$(\varphi_{39} \to \varphi_{98}) \to \varphi_{205}$}
\BIC{$\varphi_{205}$}
\AXC{$\stackrel{\Pi_{94}}{\vdots}$}\noLine
\UIC{$\varphi_{205} \to \mathrm{ContIn}(f,A,c) \to \varphi_{181}$}
\BIC{$\mathrm{ContIn}(f,A,c) \to \varphi_{181}$}
\end{ptree}
\noindent$\Pi_{96}$:
\begin{ptree}
\AXC{$\stackrel{\Pi_{95}}{\vdots}$}\noLine
\UIC{$\mathrm{ContIn}(f,A,c) \to \varphi_{181}$}
\AXC{$(\mathrm{ContIn}(f,A,c) \to \varphi_{181}) \to \varphi_{206}$\lbl{A1}}
\BIC{$\varphi_{206}$}
\UIC{$\forall n\,(\varphi_{206})$}
\AXC{$\forall n\,(\varphi_{206}) \to \varphi_{393}$\lbl{A1}}
\BIC{$\varphi_{393}$}
\end{ptree}
\noindent$\Pi_{97}$:
\begin{ptree}
\AXC{$\varphi_{278}$\lbl{A4}}
\AXC{$\forall n\,(\varphi_{206}) \to \varphi_{131} \to \forall n \in \mathbb{N}\,(a_{n} \in \mathrm{dom}((f / g)))$\lbl{A3}}
\AXC{$\stackrel{\Pi_{96}}{\vdots}$}\noLine
\UIC{$\varphi_{393}$}
\BIC{$\varphi_{278} \to \mathrm{ContIn}(f,A,c) \to \varphi_{184}$}
\BIC{$\mathrm{ContIn}(f,A,c) \to \varphi_{184}$}
\end{ptree}
\noindent$\Pi_{98}$:
\begin{ptree}
\AXC{$\stackrel{\Pi_{84}}{\vdots}$}\noLine
\UIC{$\mathrm{ContIn}(f,A,c) \to \varphi_{259}$}
\AXC{$(\mathrm{ContIn}(f,A,c) \to \varphi_{259}) \to \varphi_{163}$\lbl{A1}}
\BIC{$\varphi_{163}$}
\end{ptree}
\noindent$\Pi_{99}$:
\begin{ptree}
\AXC{$\varphi_{41} \to \varphi_{100}$\lbl{A2}}
\AXC{$\stackrel{\Pi_{98}}{\vdots}$}\noLine
\UIC{$\varphi_{163}$}
\AXC{$\varphi_{163} \to (\varphi_{41} \to \varphi_{100}) \to \varphi_{208}$\lbl{A1}}
\BIC{$(\varphi_{41} \to \varphi_{100}) \to \varphi_{208}$}
\BIC{$\varphi_{208}$}
\AXC{$\varphi_{208} \to (\mathrm{ContIn}(f,A,c) \to \varphi_{181}) \to \varphi_{325}$\lbl{A1}}
\BIC{$(\mathrm{ContIn}(f,A,c) \to \varphi_{181}) \to \varphi_{325}$}
\end{ptree}
\noindent$\Pi_{100}$:
\begin{ptree}
\AXC{$\varphi_{83}$\lbl{A2}}
\AXC{$\stackrel{\Pi_{97}}{\vdots}$}\noLine
\UIC{$\mathrm{ContIn}(f,A,c) \to \varphi_{184}$}
\AXC{$\varphi_{344}$\lbl{A1}}
\BIC{$\varphi_{83} \to \mathrm{ContIn}(f,A,c) \to \varphi_{181}$}
\BIC{$\mathrm{ContIn}(f,A,c) \to \varphi_{181}$}
\AXC{$\stackrel{\Pi_{99}}{\vdots}$}\noLine
\UIC{$(\mathrm{ContIn}(f,A,c) \to \varphi_{181}) \to \varphi_{325}$}
\BIC{$\varphi_{325}$}
\end{ptree}
\noindent$\Pi_{101}$:
\begin{ptree}
\AXC{$\stackrel{\Pi_{84}}{\vdots}$}\noLine
\UIC{$\mathrm{ContIn}(f,A,c) \to \varphi_{259}$}
\AXC{$\varphi_{307}$\lbl{A1}}
\BIC{$\mathrm{ContIn}(f,A,c) \to \varphi_{10}$}
\end{ptree}
\noindent$\Pi_{102}$:
\begin{ptree}
\AXC{$\forall f\,\forall a\,(\mathrm{RFn}(f) \to \mathrm{Seq}(a) \to \varphi_{127})$\lbl{[S8]}}
\AXC{$\forall f\,\forall a\,(\mathrm{RFn}(f) \to \mathrm{Seq}(a) \to \varphi_{127}) \to \forall a\,(\varphi_{223})$\lbl{A2}}
\BIC{$\forall a\,(\varphi_{223})$}
\AXC{$\forall a\,(\varphi_{223}) \to \varphi_{223}$\lbl{A2}}
\BIC{$\varphi_{223}$}
\end{ptree}
\noindent$\Pi_{103}$:
\begin{ptree}
\AXC{$\mathrm{ContIn}(f,A,c) \to \varphi_{12}$\lbl{A1}}
\AXC{$\stackrel{\Pi_{101}}{\vdots}$}\noLine
\UIC{$\mathrm{ContIn}(f,A,c) \to \varphi_{10}$}
\AXC{$\stackrel{\Pi_{102}}{\vdots}$}\noLine
\UIC{$\varphi_{223}$}
\AXC{$\varphi_{223} \to \varphi_{378}$\lbl{A1}}
\BIC{$\varphi_{378}$}
\BIC{$(\mathrm{ContIn}(f,A,c) \to \varphi_{12}) \to \varphi_{338}$}
\BIC{$\varphi_{338}$}
\end{ptree}
\noindent$\Pi_{104}$:
\begin{ptree}
\AXC{$\stackrel{\Pi_{97}}{\vdots}$}\noLine
\UIC{$\mathrm{ContIn}(f,A,c) \to \varphi_{184}$}
\AXC{$\stackrel{\Pi_{103}}{\vdots}$}\noLine
\UIC{$\varphi_{338}$}
\BIC{$\mathrm{ContIn}(f,A,c) \to \mathrm{ContIn}(g,A,c) \to \varphi_{233}$}
\end{ptree}
\noindent$\Pi_{105}$:
\begin{ptree}
\AXC{$\varphi_{176}$\lbl{A2}}
\AXC{$\stackrel{\Pi_{104}}{\vdots}$}\noLine
\UIC{$\mathrm{ContIn}(f,A,c) \to \mathrm{ContIn}(g,A,c) \to \varphi_{233}$}
\AXC{$\varphi_{346}$\lbl{A1}}
\BIC{$\mathrm{ContIn}(f,A,c) \to \mathrm{ContIn}(g,A,c) \to \varphi_{217}$}
\AXC{$\varphi_{376}$\lbl{A1}}
\BIC{$\varphi_{324}$}
\BIC{$\mathrm{ContIn}(f,A,c) \to \mathrm{ContIn}(g,A,c) \to \varphi_{214}$}
\AXC{$\varphi_{372}$\lbl{A1}}
\BIC{$\varphi_{317}$}
\end{ptree}
\noindent$\Pi_{106}$:
\begin{ptree}
\AXC{$\mathrm{ContIn}(f,A,c) \to \varphi_{140}$\lbl{A1}}
\AXC{$\stackrel{\Pi_{100}}{\vdots}$}\noLine
\UIC{$\varphi_{325}$}
\BIC{$\mathrm{ContIn}(f,A,c) \to \mathrm{ContIn}(g,A,c) \to \varphi_{213}$}
\AXC{$\mathrm{ContIn}(f,A,c) \to \varphi_{140}$\lbl{A1}}
\AXC{$\stackrel{\Pi_{105}}{\vdots}$}\noLine
\UIC{$\varphi_{317}$}
\BIC{$\mathrm{ContIn}(f,A,c) \to \mathrm{ContIn}(g,A,c) \to \varphi_{200}$}
\AXC{$\varphi_{397}$\lbl{A1}}
\BIC{$\varphi_{383}$}
\BIC{$\varphi_{339}$}
\end{ptree}
\noindent$\Pi_{107}$:
\begin{ptree}
\AXC{$\stackrel{\Pi_{82}}{\vdots}$}\noLine
\UIC{$\varphi_{51}$}
\AXC{$\varphi_{51} \to \varphi_{135}$\lbl{A2}}
\BIC{$\varphi_{135}$}
\AXC{$\varphi_{135} \to (f / g) \circ a(n) = (f / g)(a_{n}) \to \varphi_{179}$\lbl{A2}}
\BIC{$(f / g) \circ a(n) = (f / g)(a_{n}) \to \varphi_{179}$}
\AXC{$\stackrel{\Pi_{106}}{\vdots}$}\noLine
\UIC{$\varphi_{339}$}
\BIC{$\mathrm{ContIn}(f,A,c) \to \mathrm{ContIn}(g,A,c) \to \varphi_{219}$}
\AXC{$\varphi_{400}$\lbl{A1}}
\BIC{$\varphi_{388}$}
\end{ptree}
\noindent$\Pi_{108}$:
\begin{ptree}
\AXC{$\stackrel{\Pi_{38}}{\vdots}$}\noLine
\UIC{$\varphi_{239}$}
\AXC{$\stackrel{\Pi_{39}}{\vdots}$}\noLine
\UIC{$\varphi_{134}$}
\AXC{$\varphi_{134} \to \varphi_{255}$\lbl{A2}}
\BIC{$\varphi_{255}$}
\AXC{$\varphi_{360}$\lbl{A2}}
\BIC{$f \circ a \otimes g \circ a^{\ast}(n) = f(a_{n}) \cdot g \circ a(n)^{-1} \to \varphi_{229}$}
\AXC{$\stackrel{\Pi_{81}}{\vdots}$}\noLine
\UIC{$\varphi_{369}$}
\BIC{$\mathrm{ContIn}(f,A,c) \to \mathrm{ContIn}(g,A,c) \to \varphi_{242}$}
\AXC{$\varphi_{384}$\lbl{A1}}
\BIC{$\varphi_{341}$}
\BIC{$\mathrm{ContIn}(f,A,c) \to \mathrm{ContIn}(g,A,c) \to \varphi_{243}$}
\AXC{$\stackrel{\Pi_{107}}{\vdots}$}\noLine
\UIC{$\varphi_{388}$}
\BIC{$\varphi_{355}$}
\end{ptree}
\noindent$\Pi_{109}$:
\begin{ptree}
\AXC{$\stackrel{\Pi_{37}}{\vdots}$}\noLine
\UIC{$\varphi_{51}$}
\AXC{$\varphi_{51} \to \varphi_{202}$\lbl{A2}}
\BIC{$\varphi_{202}$}
\AXC{$\varphi_{202} \to (f / g) \circ a(n) = f(a_{n}) \cdot g(a_{n})^{-1} \to \varphi_{224}$\lbl{A2}}
\BIC{$(f / g) \circ a(n) = f(a_{n}) \cdot g(a_{n})^{-1} \to \varphi_{224}$}
\AXC{$\stackrel{\Pi_{108}}{\vdots}$}\noLine
\UIC{$\varphi_{355}$}
\BIC{$\mathrm{ContIn}(f,A,c) \to \mathrm{ContIn}(g,A,c) \to \varphi_{225}$}
\AXC{$\varphi_{382}$\lbl{A1}}
\BIC{$\varphi_{333}$}
\end{ptree}
\noindent$\Pi_{110}$:
\begin{ptree}
\AXC{$\stackrel{\Pi_{2}}{\vdots}$}\noLine
\UIC{$\forall x\,\forall y\,(x = y \to y = x)$}
\AXC{$\varphi_{75}$\lbl{A2}}
\BIC{$\forall y\,((f / g) \circ a(n) = y \to y = (f / g) \circ a(n))$}
\AXC{$\varphi_{276}$\lbl{A2}}
\BIC{$\varphi_{211}$}
\AXC{$\stackrel{\Pi_{109}}{\vdots}$}\noLine
\UIC{$\varphi_{333}$}
\BIC{$\mathrm{ContIn}(f,A,c) \to \mathrm{ContIn}(g,A,c) \to \varphi_{226}$}
\AXC{$\varphi_{347}$\lbl{A1}}
\BIC{$\varphi_{261}$}
\UIC{$\forall n\,(\varphi_{261})$}
\end{ptree}
\noindent$\Pi_{111}$:
\begin{ptree}
\AXC{$\varphi_{65}$\lbl{[C13]}}
\AXC{$\varphi_{65} \to \varphi_{209}$\lbl{A2}}
\BIC{$\varphi_{209}$}
\AXC{$\varphi_{209} \to \varphi_{268}$\lbl{A2}}
\BIC{$\varphi_{268}$}
\AXC{$\varphi_{268} \to \mathrm{Seq}((f / g) \circ a) \to \varphi_{92} \to \varphi_{166}$\lbl{A2}}
\BIC{$\mathrm{Seq}((f / g) \circ a) \to \varphi_{92} \to \varphi_{166}$}
\AXC{$\varphi_{403}$\lbl{A1}}
\BIC{$(\mathrm{ContIn}(f,A,c) \to \varphi_{145}) \to \varphi_{381}$}
\end{ptree}
\noindent$\Pi_{112}$:
\begin{ptree}
\AXC{$\stackrel{\Pi_{104}}{\vdots}$}\noLine
\UIC{$\mathrm{ContIn}(f,A,c) \to \mathrm{ContIn}(g,A,c) \to \varphi_{233}$}
\AXC{$\varphi_{330}$\lbl{A1}}
\BIC{$\mathrm{ContIn}(f,A,c) \to \varphi_{145}$}
\AXC{$\stackrel{\Pi_{111}}{\vdots}$}\noLine
\UIC{$(\mathrm{ContIn}(f,A,c) \to \varphi_{145}) \to \varphi_{381}$}
\BIC{$\varphi_{381}$}
\end{ptree}
\noindent$\Pi_{113}$:
\begin{ptree}
\AXC{$\varphi_{278}$\lbl{A4}}
\AXC{$\forall n\,(\varphi_{261}) \to \varphi_{131} \to \varphi_{92}$\lbl{A3}}
\AXC{$\stackrel{\Pi_{110}}{\vdots}$}\noLine
\UIC{$\forall n\,(\varphi_{261})$}
\AXC{$\forall n\,(\varphi_{261}) \to (\forall n\,(\varphi_{261}) \to \varphi_{131} \to \varphi_{92}) \to \varphi_{354}$\lbl{A1}}
\BIC{$(\forall n\,(\varphi_{261}) \to \varphi_{131} \to \varphi_{92}) \to \varphi_{354}$}
\BIC{$\varphi_{354}$}
\BIC{$\mathrm{ContIn}(f,A,c) \to \mathrm{ContIn}(g,A,c) \to \varphi_{227}$}
\AXC{$\stackrel{\Pi_{112}}{\vdots}$}\noLine
\UIC{$\varphi_{381}$}
\BIC{$\varphi_{331}$}
\end{ptree}
\noindent$\Pi_{114}$:
\begin{ptree}
\AXC{$\stackrel{\Pi_{2}}{\vdots}$}\noLine
\UIC{$\forall x\,\forall y\,(x = y \to y = x)$}
\AXC{$\varphi_{16}$\lbl{A2}}
\BIC{$\forall y\,((f / g)(c) = y \to y = (f / g)(c))$}
\AXC{$\forall y\,((f / g)(c) = y \to y = (f / g)(c)) \to \varphi_{88}$\lbl{A2}}
\BIC{$\varphi_{88}$}
\end{ptree}
\noindent$\Pi_{115}$:
\begin{ptree}
\AXC{$\forall f\,\forall A\,\forall c\,(\varphi_{235})$\lbl{DefContIn}}
\AXC{$\forall f\,\forall A\,\forall c\,(\varphi_{235}) \to \forall A\,\forall c\,(\varphi_{235})$\lbl{A2}}
\BIC{$\forall A\,\forall c\,(\varphi_{235})$}
\AXC{$\forall A\,\forall c\,(\varphi_{235}) \to \forall c\,(\varphi_{235})$\lbl{A2}}
\BIC{$\forall c\,(\varphi_{235})$}
\AXC{$\forall c\,(\varphi_{235}) \to \varphi_{235}$\lbl{A2}}
\BIC{$\varphi_{235}$}
\end{ptree}
\noindent$\Pi_{116}$:
\begin{ptree}
\AXC{$\stackrel{\Pi_{115}}{\vdots}$}\noLine
\UIC{$\varphi_{235}$}
\AXC{$\mathrm{ContIn}(f,A,c) \to \mathrm{ContIn}(f,A,c)$\lbl{A1}}
\AXC{$\varphi_{295}$\lbl{A1}}
\BIC{$\varphi_{235} \to \varphi_{36}$}
\BIC{$\varphi_{36}$}
\end{ptree}
\noindent$\Pi_{117}$:
\begin{ptree}
\AXC{$\forall x \in A\,(g(x) \neq 0) \to c \in A \to g(c) \neq 0$\lbl{A2}}
\AXC{$\mathrm{ContIn}(f,A,c) \to \varphi_{38}$\lbl{A1}}
\AXC{$(\mathrm{ContIn}(f,A,c) \to \varphi_{38}) \to \varphi_{173}$\lbl{A1}}
\BIC{$\varphi_{173}$}
\BIC{$\mathrm{ContIn}(f,A,c) \to \varphi_{20}$}
\end{ptree}
\noindent$\Pi_{118}$:
\begin{ptree}
\AXC{$\stackrel{\Pi_{116}}{\vdots}$}\noLine
\UIC{$\varphi_{36}$}
\AXC{$\stackrel{\Pi_{117}}{\vdots}$}\noLine
\UIC{$\mathrm{ContIn}(f,A,c) \to \varphi_{20}$}
\AXC{$(\mathrm{ContIn}(f,A,c) \to \varphi_{20}) \to \varphi_{36} \to \varphi_{147}$\lbl{A1}}
\BIC{$\varphi_{36} \to \varphi_{147}$}
\BIC{$\varphi_{147}$}
\end{ptree}
\noindent$\Pi_{119}$:
\begin{ptree}
\AXC{$\forall X\,\forall Y\,\forall x\,(X \subseteq Y \to x \in X \to x \in Y)$\lbl{[B5]}}
\AXC{$\varphi_{54}$\lbl{A2}}
\BIC{$\forall Y\,\forall x\,(A \subseteq Y \to x \in A \to x \in Y)$}
\AXC{$\varphi_{66}$\lbl{A2}}
\BIC{$\forall x\,(A \subseteq \mathrm{dom}(g) \to x \in A \to x \in \mathrm{dom}(g))$}
\AXC{$\varphi_{76}$\lbl{A2}}
\BIC{$A \subseteq \mathrm{dom}(g) \to c \in A \to c \in \mathrm{dom}(g)$}
\AXC{$\varphi_{314}$\lbl{A1}}
\BIC{$(\mathrm{ContIn}(f,A,c) \to \varphi_{7}) \to \varphi_{36} \to \varphi_{155}$}
\end{ptree}
\noindent$\Pi_{120}$:
\begin{ptree}
\AXC{$\stackrel{\Pi_{116}}{\vdots}$}\noLine
\UIC{$\varphi_{36}$}
\AXC{$\stackrel{\Pi_{42}}{\vdots}$}\noLine
\UIC{$\mathrm{ContIn}(f,A,c) \to \varphi_{7}$}
\AXC{$\stackrel{\Pi_{119}}{\vdots}$}\noLine
\UIC{$(\mathrm{ContIn}(f,A,c) \to \varphi_{7}) \to \varphi_{36} \to \varphi_{155}$}
\BIC{$\varphi_{36} \to \varphi_{155}$}
\BIC{$\varphi_{155}$}
\end{ptree}
\noindent$\Pi_{121}$:
\begin{ptree}
\AXC{$\forall X\,\forall Y\,\forall x\,(X \subseteq Y \to x \in X \to x \in Y)$\lbl{[B5]}}
\AXC{$\varphi_{54}$\lbl{A2}}
\BIC{$\forall Y\,\forall x\,(A \subseteq Y \to x \in A \to x \in Y)$}
\AXC{$\varphi_{68}$\lbl{A2}}
\BIC{$\forall x\,(A \subseteq \mathrm{dom}(f) \to x \in A \to x \in \mathrm{dom}(f))$}
\AXC{$\varphi_{77}$\lbl{A2}}
\BIC{$A \subseteq \mathrm{dom}(f) \to c \in A \to c \in \mathrm{dom}(f)$}
\AXC{$\varphi_{315}$\lbl{A1}}
\BIC{$(\mathrm{ContIn}(f,A,c) \to \varphi_{8}) \to \varphi_{36} \to \varphi_{156}$}
\end{ptree}
\noindent$\Pi_{122}$:
\begin{ptree}
\AXC{$\stackrel{\Pi_{116}}{\vdots}$}\noLine
\UIC{$\varphi_{36}$}
\AXC{$\stackrel{\Pi_{57}}{\vdots}$}\noLine
\UIC{$\mathrm{ContIn}(f,A,c) \to \varphi_{8}$}
\AXC{$\stackrel{\Pi_{121}}{\vdots}$}\noLine
\UIC{$(\mathrm{ContIn}(f,A,c) \to \varphi_{8}) \to \varphi_{36} \to \varphi_{156}$}
\BIC{$\varphi_{36} \to \varphi_{156}$}
\BIC{$\varphi_{156}$}
\end{ptree}
\noindent$\Pi_{123}$:
\begin{ptree}
\AXC{$\stackrel{\Pi_{118}}{\vdots}$}\noLine
\UIC{$\varphi_{147}$}
\AXC{$\stackrel{\Pi_{120}}{\vdots}$}\noLine
\UIC{$\varphi_{155}$}
\AXC{$\stackrel{\Pi_{122}}{\vdots}$}\noLine
\UIC{$\varphi_{156}$}
\AXC{$\varphi_{156} \to \varphi_{155} \to \varphi_{147} \to \varphi_{148} \to \mathrm{ContIn}(f,A,c) \to \varphi_{139}$\lbl{A1}}
\BIC{$\varphi_{155} \to \varphi_{147} \to \varphi_{148} \to \mathrm{ContIn}(f,A,c) \to \varphi_{139}$}
\BIC{$\varphi_{147} \to \varphi_{148} \to \mathrm{ContIn}(f,A,c) \to \varphi_{139}$}
\BIC{$\varphi_{148} \to \mathrm{ContIn}(f,A,c) \to \varphi_{139}$}
\end{ptree}
\noindent$\Pi_{124}$:
\begin{ptree}
\AXC{$\varphi_{39} \to \varphi_{23}$\lbl{A2}}
\AXC{$\stackrel{\Pi_{85}}{\vdots}$}\noLine
\UIC{$\varphi_{160}$}
\AXC{$\varphi_{160} \to (\varphi_{39} \to \varphi_{23}) \to \varphi_{148}$\lbl{A1}}
\BIC{$(\varphi_{39} \to \varphi_{23}) \to \varphi_{148}$}
\BIC{$\varphi_{148}$}
\AXC{$\stackrel{\Pi_{123}}{\vdots}$}\noLine
\UIC{$\varphi_{148} \to \mathrm{ContIn}(f,A,c) \to \varphi_{139}$}
\BIC{$\mathrm{ContIn}(f,A,c) \to \varphi_{139}$}
\end{ptree}
\noindent$\Pi_{125}$:
\begin{ptree}
\AXC{$\stackrel{\Pi_{124}}{\vdots}$}\noLine
\UIC{$\mathrm{ContIn}(f,A,c) \to \varphi_{139}$}
\AXC{$\varphi_{41} \to \varphi_{32}$\lbl{A2}}
\AXC{$\stackrel{\Pi_{98}}{\vdots}$}\noLine
\UIC{$\varphi_{163}$}
\AXC{$\varphi_{163} \to (\varphi_{41} \to \varphi_{32}) \to \varphi_{157}$\lbl{A1}}
\BIC{$(\varphi_{41} \to \varphi_{32}) \to \varphi_{157}$}
\BIC{$\varphi_{157}$}
\AXC{$\varphi_{157} \to \varphi_{306}$\lbl{A1}}
\BIC{$\varphi_{306}$}
\BIC{$\mathrm{ContIn}(f,A,c) \to \varphi_{185}$}
\AXC{$\varphi_{348}$\lbl{A1}}
\BIC{$\varphi_{88} \to \mathrm{ContIn}(f,A,c) \to \varphi_{186}$}
\end{ptree}
\noindent$\Pi_{126}$:
\begin{ptree}
\AXC{$\stackrel{\Pi_{114}}{\vdots}$}\noLine
\UIC{$\varphi_{88}$}
\AXC{$\stackrel{\Pi_{125}}{\vdots}$}\noLine
\UIC{$\varphi_{88} \to \mathrm{ContIn}(f,A,c) \to \varphi_{186}$}
\BIC{$\mathrm{ContIn}(f,A,c) \to \varphi_{186}$}
\AXC{$(\mathrm{ContIn}(f,A,c) \to \varphi_{186}) \to \varphi_{362}$\lbl{A1}}
\BIC{$\varphi_{362}$}
\end{ptree}
\noindent$\Pi_{127}$:
\begin{ptree}
\AXC{$\stackrel{\Pi_{1}}{\vdots}$}\noLine
\UIC{$f(c) \cdot g(c)^{-1} = (f / g)(c) \to \varphi_{113}$}
\AXC{$\stackrel{\Pi_{36}}{\vdots}$}\noLine
\UIC{$\mathrm{ContIn}(f,A,c) \to \mathrm{ContIn}(g,A,c) \to \varphi_{210}$}
\AXC{$\stackrel{\Pi_{113}}{\vdots}$}\noLine
\UIC{$\varphi_{331}$}
\BIC{$\mathrm{ContIn}(f,A,c) \to \mathrm{ContIn}(g,A,c) \to \varphi_{177}$}
\AXC{$\stackrel{\Pi_{126}}{\vdots}$}\noLine
\UIC{$\varphi_{362}$}
\BIC{$\varphi_{312}$}
\BIC{$\mathrm{ContIn}(f,A,c) \to \mathrm{ContIn}(g,A,c) \to \varphi_{152}$}
\end{ptree}
\noindent$\Pi_{128}$:
\begin{ptree}
\AXC{$\stackrel{\Pi_{127}}{\vdots}$}\noLine
\UIC{$\mathrm{ContIn}(f,A,c) \to \mathrm{ContIn}(g,A,c) \to \varphi_{152}$}
\AXC{$\varphi_{318}$\lbl{A1}}
\BIC{$\varphi_{6} \to \mathrm{Seq}(a) \to \forall n \in \mathbb{N}\,(a_{n} \in A) \to \varphi_{5}$}
\UIC{$\forall a\,(\varphi_{6} \to \mathrm{Seq}(a) \to \forall n \in \mathbb{N}\,(a_{n} \in A) \to \varphi_{5})$}
\AXC{$\forall a\,(\varphi_{6} \to \mathrm{Seq}(a) \to \forall n \in \mathbb{N}\,(a_{n} \in A) \to \varphi_{5}) \to \varphi_{322} \to \varphi_{299}$\lbl{A1}}
\BIC{$\varphi_{322} \to \varphi_{299}$}
\end{ptree}
\noindent$\Pi_{129}$:
\begin{ptree}
\AXC{$\forall x \in A\,(g(x) \neq 0) \to u \in A \to g(u) \neq 0$\lbl{A2}}
\AXC{$\mathrm{ContIn}(f,A,c) \to \varphi_{38}$\lbl{A1}}
\AXC{$(\mathrm{ContIn}(f,A,c) \to \varphi_{38}) \to \varphi_{174}$\lbl{A1}}
\BIC{$\varphi_{174}$}
\BIC{$\mathrm{ContIn}(f,A,c) \to \varphi_{26}$}
\end{ptree}
\noindent$\Pi_{130}$:
\begin{ptree}
\AXC{$\forall X\,\forall Y\,\forall x\,(X \subseteq Y \to x \in X \to x \in Y)$\lbl{[B5]}}
\AXC{$\varphi_{54}$\lbl{A2}}
\BIC{$\forall Y\,\forall x\,(A \subseteq Y \to x \in A \to x \in Y)$}
\AXC{$\varphi_{66}$\lbl{A2}}
\BIC{$\forall x\,(A \subseteq \mathrm{dom}(g) \to x \in A \to x \in \mathrm{dom}(g))$}
\AXC{$\varphi_{78}$\lbl{A2}}
\BIC{$A \subseteq \mathrm{dom}(g) \to u \in A \to u \in \mathrm{dom}(g)$}
\AXC{$\varphi_{296}$\lbl{A1}}
\BIC{$(\mathrm{ContIn}(f,A,c) \to \varphi_{7}) \to \varphi_{220}$}
\end{ptree}
\noindent$\Pi_{131}$:
\begin{ptree}
\AXC{$\mathrm{ContIn}(f,A,c) \to \varphi_{18}$\lbl{A1}}
\AXC{$\stackrel{\Pi_{42}}{\vdots}$}\noLine
\UIC{$\mathrm{ContIn}(f,A,c) \to \varphi_{7}$}
\AXC{$\stackrel{\Pi_{130}}{\vdots}$}\noLine
\UIC{$(\mathrm{ContIn}(f,A,c) \to \varphi_{7}) \to \varphi_{220}$}
\BIC{$\varphi_{220}$}
\BIC{$\mathrm{ContIn}(f,A,c) \to \varphi_{33}$}
\end{ptree}
\noindent$\Pi_{132}$:
\begin{ptree}
\AXC{$\forall X\,\forall Y\,\forall x\,(X \subseteq Y \to x \in X \to x \in Y)$\lbl{[B5]}}
\AXC{$\varphi_{54}$\lbl{A2}}
\BIC{$\forall Y\,\forall x\,(A \subseteq Y \to x \in A \to x \in Y)$}
\AXC{$\varphi_{68}$\lbl{A2}}
\BIC{$\forall x\,(A \subseteq \mathrm{dom}(f) \to x \in A \to x \in \mathrm{dom}(f))$}
\AXC{$\varphi_{79}$\lbl{A2}}
\BIC{$A \subseteq \mathrm{dom}(f) \to u \in A \to u \in \mathrm{dom}(f)$}
\AXC{$\varphi_{297}$\lbl{A1}}
\BIC{$(\mathrm{ContIn}(f,A,c) \to \varphi_{8}) \to \varphi_{221}$}
\end{ptree}
\noindent$\Pi_{133}$:
\begin{ptree}
\AXC{$\mathrm{ContIn}(f,A,c) \to \varphi_{18}$\lbl{A1}}
\AXC{$\stackrel{\Pi_{57}}{\vdots}$}\noLine
\UIC{$\mathrm{ContIn}(f,A,c) \to \varphi_{8}$}
\AXC{$\stackrel{\Pi_{132}}{\vdots}$}\noLine
\UIC{$(\mathrm{ContIn}(f,A,c) \to \varphi_{8}) \to \varphi_{221}$}
\BIC{$\varphi_{221}$}
\BIC{$\mathrm{ContIn}(f,A,c) \to \varphi_{34}$}
\AXC{$\varphi_{361}$\lbl{A1}}
\BIC{$(\mathrm{ContIn}(f,A,c) \to \varphi_{33}) \to \varphi_{311}$}
\end{ptree}
\noindent$\Pi_{134}$:
\begin{ptree}
\AXC{$\stackrel{\Pi_{129}}{\vdots}$}\noLine
\UIC{$\mathrm{ContIn}(f,A,c) \to \varphi_{26}$}
\AXC{$\stackrel{\Pi_{131}}{\vdots}$}\noLine
\UIC{$\mathrm{ContIn}(f,A,c) \to \varphi_{33}$}
\AXC{$\stackrel{\Pi_{133}}{\vdots}$}\noLine
\UIC{$(\mathrm{ContIn}(f,A,c) \to \varphi_{33}) \to \varphi_{311}$}
\BIC{$\varphi_{311}$}
\BIC{$\varphi_{149} \to \mathrm{ContIn}(f,A,c) \to \varphi_{53}$}
\end{ptree}
\noindent$\Pi_{135}$:
\begin{ptree}
\AXC{$\varphi_{39} \to \varphi_{25}$\lbl{A2}}
\AXC{$\stackrel{\Pi_{85}}{\vdots}$}\noLine
\UIC{$\varphi_{160}$}
\AXC{$\varphi_{160} \to (\varphi_{39} \to \varphi_{25}) \to \varphi_{149}$\lbl{A1}}
\BIC{$(\varphi_{39} \to \varphi_{25}) \to \varphi_{149}$}
\BIC{$\varphi_{149}$}
\AXC{$\stackrel{\Pi_{134}}{\vdots}$}\noLine
\UIC{$\varphi_{149} \to \mathrm{ContIn}(f,A,c) \to \varphi_{53}$}
\BIC{$\mathrm{ContIn}(f,A,c) \to \varphi_{53}$}
\AXC{$\varphi_{244}$\lbl{A1}}
\BIC{$\varphi_{6} \to u \in A \to u \in \mathrm{dom}((f / g))$}
\UIC{$\forall u\,(\varphi_{6} \to u \in A \to u \in \mathrm{dom}((f / g)))$}
\end{ptree}
\noindent$\Pi_{136}$:
\begin{ptree}
\AXC{$\varphi_{254}$\lbl{A3}}
\AXC{$\stackrel{\Pi_{135}}{\vdots}$}\noLine
\UIC{$\forall u\,(\varphi_{6} \to u \in A \to u \in \mathrm{dom}((f / g)))$}
\AXC{$\forall u\,(\varphi_{6} \to u \in A \to u \in \mathrm{dom}((f / g))) \to \varphi_{254} \to \varphi_{273}$\lbl{A1}}
\BIC{$\varphi_{254} \to \varphi_{273}$}
\BIC{$\varphi_{273}$}
\end{ptree}
\noindent$\Pi_{137}$:
\begin{ptree}
\AXC{$\forall x\,\forall y\,(x \subseteq y \leftrightarrow \forall u \in x\,(u \in y))$\lbl{Def$\subseteq$}}
\AXC{$\varphi_{57}$\lbl{A2}}
\BIC{$\forall y\,(A \subseteq y \leftrightarrow \forall u \in A\,(u \in y))$}
\AXC{$\forall y\,(A \subseteq y \leftrightarrow \forall u \in A\,(u \in y)) \to \varphi_{9}$\lbl{A2}}
\BIC{$\varphi_{9}$}
\AXC{$\varphi_{9} \to \varphi_{236}$\lbl{A1}}
\BIC{$\varphi_{236}$}
\end{ptree}
\noindent$\Pi_{138}$:
\begin{ptree}
\AXC{$\stackrel{\Pi_{101}}{\vdots}$}\noLine
\UIC{$\mathrm{ContIn}(f,A,c) \to \varphi_{10}$}
\AXC{$\varphi_{195}$\lbl{[G2]}}
\AXC{$\varphi_{195} \to \forall A\,(\varphi_{240})$\lbl{A2}}
\BIC{$\forall A\,(\varphi_{240})$}
\AXC{$\forall A\,(\varphi_{240}) \to \varphi_{240}$\lbl{A2}}
\BIC{$\varphi_{240}$}
\AXC{$\varphi_{332}$\lbl{A2}}
\BIC{$\mathrm{RFn}((f / g)) \to A \subseteq \mathrm{dom}((f / g)) \to c \in A \to \varphi_{153}$}
\AXC{$\varphi_{387}$\lbl{A1}}
\BIC{$\varphi_{365}$}
\BIC{$(\mathrm{ContIn}(f,A,c) \to \varphi_{27}) \to \varphi_{36} \to \varphi_{283}$}
\end{ptree}
\noindent$\Pi_{139}$:
\begin{ptree}
\AXC{$\stackrel{\Pi_{116}}{\vdots}$}\noLine
\UIC{$\varphi_{36}$}
\AXC{$\varphi_{6} \to \forall u\,(\varphi_{6})$\lbl{A4}}
\AXC{$\stackrel{\Pi_{136}}{\vdots}$}\noLine
\UIC{$\varphi_{273}$}
\BIC{$\mathrm{ContIn}(f,A,c) \to \mathrm{ContIn}(g,A,c) \to \varphi_{19}$}
\AXC{$\stackrel{\Pi_{137}}{\vdots}$}\noLine
\UIC{$\varphi_{236}$}
\BIC{$\mathrm{ContIn}(f,A,c) \to \varphi_{27}$}
\AXC{$\stackrel{\Pi_{138}}{\vdots}$}\noLine
\UIC{$(\mathrm{ContIn}(f,A,c) \to \varphi_{27}) \to \varphi_{36} \to \varphi_{283}$}
\BIC{$\varphi_{36} \to \varphi_{283}$}
\BIC{$\varphi_{283}$}
\end{ptree}
\noindent 主木:
\begin{ptree}
\AXC{$\varphi_{6} \to \forall a\,(\varphi_{6})$\lbl{A4}}
\AXC{$\varphi_{322}$\lbl{A3}}
\AXC{$\stackrel{\Pi_{128}}{\vdots}$}\noLine
\UIC{$\varphi_{322} \to \varphi_{299}$}
\BIC{$\varphi_{299}$}
\BIC{$\mathrm{ContIn}(f,A,c) \to \mathrm{ContIn}(g,A,c) \to \varphi_{162}$}
\AXC{$\stackrel{\Pi_{139}}{\vdots}$}\noLine
\UIC{$\varphi_{283}$}
\BIC{$\mathrm{ContIn}(f,A,c) \to \varphi_{28}$}
\UIC{$\forall c\,(\mathrm{ContIn}(f,A,c) \to \varphi_{28})$}
\UIC{$\forall A\,\forall c\,(\mathrm{ContIn}(f,A,c) \to \varphi_{28})$}
\UIC{$\forall g\,\forall A\,\forall c\,(\mathrm{ContIn}(f,A,c) \to \varphi_{28})$}
\UIC{$\forall f\,\forall g\,\forall A\,\forall c\,(\mathrm{ContIn}(f,A,c) \to \varphi_{28})$}
\end{ptree}
\begin{abbrevlist}
\item $\varphi_{1} :\equiv \mathrm{ContIn}(f,A,c) \to \mathrm{ContIn}(g,A,c) \to \mathrm{ContIn}(g,A,c)$
\item $\varphi_{2} :\equiv \mathrm{ContIn}(g,A,c) \to \forall x \in A\,(g(x) \neq 0) \to \mathrm{RFn}(g)$
\item $\varphi_{3} :\equiv \mathrm{ContIn}(g,A,c) \to \forall x \in A\,(g(x) \neq 0) \to \mathrm{RFn}(f)$
\item $\varphi_{4} :\equiv \forall n \in \mathbb{N}\,(a_{n} \in A) \to \mathrm{Conv}(a,c) \to \mathrm{Conv}(a,c)$
\item $\varphi_{5} :\equiv \mathrm{Conv}(a,c) \to \mathrm{Conv}((f / g) \circ a,(f / g)(c))$
\item $\varphi_{6} :\equiv \mathrm{ContIn}(f,A,c) \land \mathrm{ContIn}(g,A,c) \land \forall x \in A\,(g(x) \neq 0)$
\item $\varphi_{7} :\equiv \mathrm{ContIn}(g,A,c) \to \forall x \in A\,(g(x) \neq 0) \to A \subseteq \mathrm{dom}(g)$
\item $\varphi_{8} :\equiv \mathrm{ContIn}(g,A,c) \to \forall x \in A\,(g(x) \neq 0) \to A \subseteq \mathrm{dom}(f)$
\item $\varphi_{9} :\equiv A \subseteq \mathrm{dom}((f / g)) \leftrightarrow \forall u \in A\,(u \in \mathrm{dom}((f / g)))$
\item $\varphi_{10} :\equiv \mathrm{ContIn}(g,A,c) \to \forall x \in A\,(g(x) \neq 0) \to \mathrm{RFn}((f / g))$
\item $\varphi_{11} :\equiv \forall x\,\forall y\,(x = y \to y = x) \to \forall y\,(g(a_{n}) = y \to y = g(a_{n}))$
\item $\varphi_{12} :\equiv \mathrm{ContIn}(g,A,c) \to \forall x \in A\,(g(x) \neq 0) \to \mathrm{Seq}(a) \to \mathrm{Seq}(a)$
\item $\varphi_{13} :\equiv \mathrm{Conv}((f / g) \circ a,z) \to \mathrm{Conv}((f / g) \circ a,w)$
\item $\varphi_{14} :\equiv \mathrm{Seq}(a) \to \forall n \in \mathbb{N}\,(a_{n} \in A) \to \mathrm{Conv}(a,c) \to \mathrm{Seq}(a)$
\item $\varphi_{15} :\equiv \mathrm{Seq}(a) \to \forall n \in \mathbb{N}\,(a_{n} \in A) \to \mathrm{Conv}(a,c) \to \mathrm{RFn}(g)$
\item $\varphi_{16} :\equiv \forall x\,\forall y\,(x = y \to y = x) \to \forall y\,((f / g)(c) = y \to y = (f / g)(c))$
\item $\varphi_{17} :\equiv \mathrm{Seq}(a) \to \forall n \in \mathbb{N}\,(a_{n} \in A) \to \mathrm{Conv}(a,c) \to c \in A$
\item $\varphi_{18} :\equiv \mathrm{ContIn}(g,A,c) \to \forall x \in A\,(g(x) \neq 0) \to u \in A \to u \in A$
\item $\varphi_{19} :\equiv \forall x \in A\,(g(x) \neq 0) \to \forall u \in A\,(u \in \mathrm{dom}((f / g)))$
\item $\varphi_{20} :\equiv \mathrm{ContIn}(g,A,c) \to \forall x \in A\,(g(x) \neq 0) \to c \in A \to g(c) \neq 0$
\item $\varphi_{21} :\equiv \mathrm{Seq}(a) \to \forall n \in \mathbb{N}\,(a_{n} \in A) \to \mathrm{Conv}(a,c) \to g(c) \neq 0$
\item $\varphi_{22} :\equiv \mathrm{Conv}(a,c) \to \mathrm{Conv}((f / g) \circ a,f(c) \cdot g(c)^{-1})$
\item $\varphi_{23} :\equiv c \in \mathrm{dom}((f / g)) \leftrightarrow c \in \mathrm{dom}(f) \land c \in \mathrm{dom}(g) \land g(c) \neq 0$
\item $\varphi_{24} :\equiv \forall n \in \mathbb{N}\,(a_{n} \in A) \to \mathrm{Conv}(a,c) \to c \in \mathrm{dom}((f / g))$
\item $\varphi_{25} :\equiv u \in \mathrm{dom}((f / g)) \leftrightarrow u \in \mathrm{dom}(f) \land u \in \mathrm{dom}(g) \land g(u) \neq 0$
\item $\varphi_{26} :\equiv \mathrm{ContIn}(g,A,c) \to \forall x \in A\,(g(x) \neq 0) \to u \in A \to g(u) \neq 0$
\item $\varphi_{27} :\equiv \mathrm{ContIn}(g,A,c) \to \forall x \in A\,(g(x) \neq 0) \to A \subseteq \mathrm{dom}((f / g))$
\item $\varphi_{28} :\equiv \mathrm{ContIn}(g,A,c) \to \forall x \in A\,(g(x) \neq 0) \to \mathrm{ContIn}((f / g),A,c)$
\item $\varphi_{29} :\equiv \forall n \in \mathbb{N}\,(a_{n} \in A) \to \mathrm{Conv}(a,c) \to n \in \mathbb{N} \to n \in \mathbb{N}$
\item $\varphi_{30} :\equiv \mathrm{Seq}(a) \to \forall n \in \mathbb{N}\,(a_{n} \in A) \to \mathrm{Conv}(a,c) \to c \in \mathrm{dom}(g)$
\item $\varphi_{31} :\equiv \mathrm{Seq}(a) \to \forall n \in \mathbb{N}\,(a_{n} \in A) \to \mathrm{Conv}(a,c) \to c \in \mathrm{dom}(f)$
\item $\varphi_{32} :\equiv c \in \mathrm{dom}((f / g)) \to (f / g)(c) = f(c) \cdot g(c)^{-1}$
\item $\varphi_{33} :\equiv \mathrm{ContIn}(g,A,c) \to \forall x \in A\,(g(x) \neq 0) \to u \in A \to u \in \mathrm{dom}(g)$
\item $\varphi_{34} :\equiv \mathrm{ContIn}(g,A,c) \to \forall x \in A\,(g(x) \neq 0) \to u \in A \to u \in \mathrm{dom}(f)$
\item $\varphi_{35} :\equiv \forall x\,\forall y\,(x = y \to y = x) \to \forall y\,(g \circ a(n) = y \to y = g \circ a(n))$
\item $\varphi_{36} :\equiv \mathrm{ContIn}(f,A,c) \to \mathrm{ContIn}(g,A,c) \to \forall x \in A\,(g(x) \neq 0) \to c \in A$
\item $\varphi_{37} :\equiv \mathrm{Seq}(a) \to \forall n \in \mathbb{N}\,(a_{n} \in A) \to \mathrm{Conv}(a,c) \to A \subseteq \mathrm{dom}(g)$
\item $\varphi_{38} :\equiv \mathrm{ContIn}(g,A,c) \to \forall x \in A\,(g(x) \neq 0) \to \forall x \in A\,(g(x) \neq 0)$
\item $\varphi_{39} :\equiv \forall x\,(x \in \mathrm{dom}((f / g)) \leftrightarrow x \in \mathrm{dom}(f) \land x \in \mathrm{dom}(g) \land g(x) \neq 0)$
\item $\varphi_{40} :\equiv \forall n \in \mathbb{N}\,(a_{n} \in A) \to \mathrm{Conv}(a,c) \to \mathrm{Seq}(g \circ a^{\ast})$
\item $\varphi_{41} :\equiv \forall x \in \mathrm{dom}((f / g))\,((f / g)(x) = f(x) \cdot g(x)^{-1})$
\item $\varphi_{42} :\equiv \forall n \in \mathbb{N}\,(a_{n} \in A) \to \mathrm{Conv}(a,c) \to \mathrm{Seq}((f / g) \circ a)$
\item $\varphi_{43} :\equiv \forall z\,\forall w\,(z = w \to z = 0 \to w = 0) \to \forall w\,(g(a_{n}) = w \to g(a_{n}) = 0 \to w = 0)$
\item $\varphi_{44} :\equiv \forall n \in \mathbb{N}\,(a_{n} \in A) \to \mathrm{Conv}(a,c) \to n \in \mathbb{N} \to a_{n} \in A$
\item $\varphi_{45} :\equiv \mathrm{Seq}(a) \to \forall n \in \mathbb{N}\,(a_{n} \in A) \to \mathrm{Conv}(a,c) \to \mathrm{Seq}(g \circ a)$
\item $\varphi_{46} :\equiv \mathrm{Seq}(a) \to \forall n \in \mathbb{N}\,(a_{n} \in A) \to \mathrm{Conv}(a,c) \to \mathrm{Seq}(f \circ a)$
\item $\varphi_{47} :\equiv \mathrm{Conv}(a,c) \to n \in \mathbb{N} \to n \in \mathbb{N} \to g \circ a(n) = g(a_{n})$
\item $\varphi_{48} :\equiv \forall n \in \mathbb{N}\,(a_{n} \in A) \to \mathrm{Conv}(a,c) \to \mathrm{Conv}(g \circ a,g(c))$
\item $\varphi_{49} :\equiv \forall n \in \mathbb{N}\,(a_{n} \in A) \to \mathrm{Conv}(a,c) \to \mathrm{Conv}(f \circ a,f(c))$
\item $\varphi_{50} :\equiv \mathrm{Conv}(a,c) \to n \in \mathbb{N} \to n \in \mathbb{N} \to f \circ a(n) = f(a_{n})$
\item $\varphi_{51} :\equiv \forall y\,\forall z\,((f / g) \circ a(n) = y \to y = z \to (f / g) \circ a(n) = z)$
\item $\varphi_{52} :\equiv \forall n \in \mathbb{N}\,(a_{n} \in A) \to \mathrm{Conv}(a,c) \to \forall n \in \mathbb{N}\,(a_{n} \in A)$
\item $\varphi_{53} :\equiv \mathrm{ContIn}(g,A,c) \to \forall x \in A\,(g(x) \neq 0) \to u \in A \to u \in \mathrm{dom}((f / g))$
\item $\varphi_{54} :\equiv \forall X\,\forall Y\,\forall x\,(X \subseteq Y \to x \in X \to x \in Y) \to \forall Y\,\forall x\,(A \subseteq Y \to x \in A \to x \in Y)$
\item $\varphi_{55} :\equiv \forall a\,\forall \alpha\,(\mathrm{Conv}(a,\alpha) \to \mathrm{Seq}(a)) \to \forall \alpha\,(\mathrm{Conv}(g \circ a,\alpha) \to \mathrm{Seq}(g \circ a))$
\item $\varphi_{56} :\equiv \forall a\,\forall \alpha\,(\mathrm{Conv}(a,\alpha) \to \mathrm{Seq}(a)) \to \forall \alpha\,(\mathrm{Conv}(f \circ a,\alpha) \to \mathrm{Seq}(f \circ a))$
\item $\varphi_{57} :\equiv \forall x\,\forall y\,(x \subseteq y \leftrightarrow \forall u \in x\,(u \in y)) \to \forall y\,(A \subseteq y \leftrightarrow \forall u \in A\,(u \in y))$
\item $\varphi_{58} :\equiv \mathrm{Conv}(a,c) \to n \in \mathbb{N} \to \forall n \in \mathbb{N}\,(g \circ a(n) = g(a_{n}))$
\item $\varphi_{59} :\equiv \forall n \in \mathbb{N}\,(a_{n} \in A) \to \mathrm{Conv}(a,c) \to n \in \mathbb{N} \to g(a_{n}) \neq 0$
\item $\varphi_{60} :\equiv \mathrm{Conv}(a,c) \to n \in \mathbb{N} \to \forall n \in \mathbb{N}\,(f \circ a(n) = f(a_{n}))$
\item $\varphi_{61} :\equiv \forall n \in \mathbb{N}\,(a_{n} \in A) \to \mathrm{Conv}(a,c) \to n \in \mathbb{N} \to a_{n} \in \mathrm{dom}(g)$
\item $\varphi_{62} :\equiv \forall n \in \mathbb{N}\,(a_{n} \in A) \to \mathrm{Conv}(a,c) \to n \in \mathbb{N} \to a_{n} \in \mathrm{dom}(f)$
\item $\varphi_{63} :\equiv \mathrm{Conv}(a,c) \to n \in \mathbb{N} \to g \circ a^{\ast}(n) = g \circ a(n)^{-1}$
\item $\varphi_{64} :\equiv \mathrm{Conv}(a,c) \to n \in \mathbb{N} \to (f / g) \circ a(n) = (f / g)(a_{n})$
\item $\varphi_{65} :\equiv \forall a\,\forall b\,\forall \alpha\,(\mathrm{Seq}(b) \to \forall n \in \mathbb{N}\,(a_{n} = b_{n}) \to \mathrm{Conv}(a,\alpha) \to \mathrm{Conv}(b,\alpha))$
\item $\varphi_{66} :\equiv \forall Y\,\forall x\,(A \subseteq Y \to x \in A \to x \in Y) \to \forall x\,(A \subseteq \mathrm{dom}(g) \to x \in A \to x \in \mathrm{dom}(g))$
\item $\varphi_{67} :\equiv \forall z\,\forall w\,(z = w \to z = 0 \to w = 0) \to \forall w\,(g \circ a(n) = w \to g \circ a(n) = 0 \to w = 0)$
\item $\varphi_{68} :\equiv \forall Y\,\forall x\,(A \subseteq Y \to x \in A \to x \in Y) \to \forall x\,(A \subseteq \mathrm{dom}(f) \to x \in A \to x \in \mathrm{dom}(f))$
\item $\varphi_{69} :\equiv \forall n \in \mathbb{N}\,(a_{n} \in A) \to \mathrm{Conv}(a,c) \to \forall n \in \mathbb{N}\,(a_{n} \in \mathrm{dom}(g))$
\item $\varphi_{70} :\equiv \mathrm{ContIn}(g,A,c) \to \forall x \in A\,(g(x) \neq 0) \to a_{n} \in A \to g(a_{n}) \neq 0$
\item $\varphi_{71} :\equiv \forall n \in \mathbb{N}\,(a_{n} \in A) \to \mathrm{Conv}(a,c) \to n \in \mathbb{N} \to g \circ a(n) \neq 0$
\item $\varphi_{72} :\equiv \forall n \in \mathbb{N}\,(a_{n} \in A) \to \mathrm{Conv}(a,c) \to \forall n \in \mathbb{N}\,(a_{n} \in \mathrm{dom}(f))$
\item $\varphi_{73} :\equiv n \in \mathbb{N} \to \forall n \in \mathbb{N}\,((f / g) \circ a(n) = (f / g)(a_{n}))$
\item $\varphi_{74} :\equiv \forall n \in \mathbb{N}\,(a_{n} \in A) \to \mathrm{Conv}(a,c) \to \forall n \in \mathbb{N}\,(g \circ a(n) \neq 0)$
\item $\varphi_{75} :\equiv \forall x\,\forall y\,(x = y \to y = x) \to \forall y\,((f / g) \circ a(n) = y \to y = (f / g) \circ a(n))$
\item $\varphi_{76} :\equiv \forall x\,(A \subseteq \mathrm{dom}(g) \to x \in A \to x \in \mathrm{dom}(g)) \to A \subseteq \mathrm{dom}(g) \to c \in A \to c \in \mathrm{dom}(g)$
\item $\varphi_{77} :\equiv \forall x\,(A \subseteq \mathrm{dom}(f) \to x \in A \to x \in \mathrm{dom}(f)) \to A \subseteq \mathrm{dom}(f) \to c \in A \to c \in \mathrm{dom}(f)$
\item $\varphi_{78} :\equiv \forall x\,(A \subseteq \mathrm{dom}(g) \to x \in A \to x \in \mathrm{dom}(g)) \to A \subseteq \mathrm{dom}(g) \to u \in A \to u \in \mathrm{dom}(g)$
\item $\varphi_{79} :\equiv \forall x\,(A \subseteq \mathrm{dom}(f) \to x \in A \to x \in \mathrm{dom}(f)) \to A \subseteq \mathrm{dom}(f) \to u \in A \to u \in \mathrm{dom}(f)$
\item $\varphi_{80} :\equiv \mathrm{Conv}(a,c) \to \mathrm{Seq}(g \circ a) \land \forall n \in \mathbb{N}\,(g \circ a(n) = g(a_{n}))$
\item $\varphi_{81} :\equiv \mathrm{Conv}(a,c) \to \mathrm{Seq}(f \circ a) \land \forall n \in \mathbb{N}\,(f \circ a(n) = f(a_{n}))$
\item $\varphi_{82} :\equiv \forall n \in \mathbb{N}\,(a_{n} \in A) \to \mathrm{Conv}(a,c) \to n \in \mathbb{N} \to a_{n} \in \mathrm{dom}((f / g))$
\item $\varphi_{83} :\equiv \forall n \in \mathbb{N}\,(a_{n} \in \mathrm{dom}((f / g))) \to n \in \mathbb{N} \to a_{n} \in \mathrm{dom}((f / g))$
\item $\varphi_{84} :\equiv \mathrm{Conv}(a,c) \to \mathrm{Conv}(f \circ a \otimes g \circ a^{\ast},f(c) \cdot g(c)^{-1})$
\item $\varphi_{85} :\equiv \forall n \in \mathbb{N}\,(a_{n} \in A) \to \mathrm{Conv}(a,c) \to \forall n \in \mathbb{N}\,(a_{n} \in \mathrm{dom}((f / g)))$
\item $\varphi_{86} :\equiv \forall n \in \mathbb{N}\,(a_{n} \in A) \to \mathrm{Conv}(a,c) \to (f / g)(c) = f(c) \cdot g(c)^{-1}$
\item $\varphi_{87} :\equiv \forall n \in \mathbb{N}\,(a_{n} \in A) \to \mathrm{Conv}(a,c) \to f(c) \cdot g(c)^{-1} = (f / g)(c)$
\item $\varphi_{88} :\equiv (f / g)(c) = f(c) \cdot g(c)^{-1} \to f(c) \cdot g(c)^{-1} = (f / g)(c)$
\item $\varphi_{89} :\equiv \mathrm{Conv}(a,c) \to n \in \mathbb{N} \to n \in \mathbb{N} \to g \circ a^{\ast}(n) = g \circ a(n)^{-1}$
\item $\varphi_{90} :\equiv \mathrm{Conv}(a,c) \to n \in \mathbb{N} \to (f / g)(a_{n}) = f(a_{n}) \cdot g(a_{n})^{-1}$
\item $\varphi_{91} :\equiv \mathrm{Conv}(a,c) \to n \in \mathbb{N} \to n \in \mathbb{N} \to (f / g) \circ a(n) = (f / g)(a_{n})$
\item $\varphi_{92} :\equiv \forall n \in \mathbb{N}\,(f \circ a \otimes g \circ a^{\ast}(n) = (f / g) \circ a(n))$
\item $\varphi_{93} :\equiv \mathrm{Conv}(a,c) \to n \in \mathbb{N} \to \forall n \in \mathbb{N}\,(g \circ a^{\ast}(n) = g \circ a(n)^{-1})$
\item $\varphi_{94} :\equiv \forall n \in \mathbb{N}\,(a_{n} \in A) \to \mathrm{Conv}(a,c) \to n \in \mathbb{N} \to g \circ a(n) = g(a_{n})$
\item $\varphi_{95} :\equiv \forall n \in \mathbb{N}\,(a_{n} \in A) \to \mathrm{Conv}(a,c) \to n \in \mathbb{N} \to g(a_{n}) = g \circ a(n)$
\item $\varphi_{96} :\equiv \forall n \in \mathbb{N}\,(a_{n} \in A) \to \mathrm{Conv}(a,c) \to n \in \mathbb{N} \to f \circ a(n) = f(a_{n})$
\item $\varphi_{97} :\equiv \mathrm{Conv}(a,c) \to n \in \mathbb{N} \to (f / g) \circ a(n) = f(a_{n}) \cdot g(a_{n})^{-1}$
\item $\varphi_{98} :\equiv a_{n} \in \mathrm{dom}((f / g)) \leftrightarrow a_{n} \in \mathrm{dom}(f) \land a_{n} \in \mathrm{dom}(g) \land g(a_{n}) \neq 0$
\item $\varphi_{99} :\equiv \mathrm{Seq}(g \circ a^{\ast}) \land \forall n \in \mathbb{N}\,(g \circ a^{\ast}(n) = g \circ a(n)^{-1})$
\item $\varphi_{100} :\equiv a_{n} \in \mathrm{dom}((f / g)) \to (f / g)(a_{n}) = f(a_{n}) \cdot g(a_{n})^{-1}$
\item $\varphi_{101} :\equiv \forall x \in A\,(g(x) \neq 0) \to \mathrm{Seq}(a) \to \forall n \in \mathbb{N}\,(a_{n} \in A) \to n \in \mathbb{N} \to a_{n} \in A$
\item $\varphi_{102} :\equiv \mathrm{Seq}((f / g) \circ a) \land \forall n \in \mathbb{N}\,((f / g) \circ a(n) = (f / g)(a_{n}))$
\item $\varphi_{103} :\equiv f \circ a(n) \cdot g \circ a(n)^{-1} = f \circ a(n) \cdot g \circ a(n)^{-1}$
\item $\varphi_{104} :\equiv \forall \alpha\,(\mathrm{Conv}(g \circ a,\alpha) \to \mathrm{Seq}(g \circ a)) \to \mathrm{Conv}(g \circ a,g(c)) \to \mathrm{Seq}(g \circ a)$
\item $\varphi_{105} :\equiv \forall \alpha\,(\mathrm{Conv}(f \circ a,\alpha) \to \mathrm{Seq}(f \circ a)) \to \mathrm{Conv}(f \circ a,f(c)) \to \mathrm{Seq}(f \circ a)$
\item $\varphi_{106} :\equiv n \in \mathbb{N} \to f(a_{n}) \cdot g \circ a(n)^{-1} = f(a_{n}) \cdot g(a_{n})^{-1}$
\item $\varphi_{107} :\equiv \forall x\,(A \subseteq \mathrm{dom}(g) \to x \in A \to x \in \mathrm{dom}(g)) \to A \subseteq \mathrm{dom}(g) \to a_{n} \in A \to a_{n} \in \mathrm{dom}(g)$
\item $\varphi_{108} :\equiv \forall x\,(x = x) \to f(a_{n}) \cdot g \circ a(n)^{-1} = f(a_{n}) \cdot g \circ a(n)^{-1}$
\item $\varphi_{109} :\equiv \forall x\,(A \subseteq \mathrm{dom}(f) \to x \in A \to x \in \mathrm{dom}(f)) \to A \subseteq \mathrm{dom}(f) \to a_{n} \in A \to a_{n} \in \mathrm{dom}(f)$
\item $\varphi_{110} :\equiv \forall a\,\forall b\,(\mathrm{Seq}(a) \to \mathrm{Seq}(b) \to \mathrm{Seq}(a \otimes b) \land \forall n \in \mathbb{N}\,(a \otimes b(n) = a_{n} \cdot b_{n}))$
\item $\varphi_{111} :\equiv \forall x \in A\,(g(x) \neq 0) \to \mathrm{Seq}(a) \to \forall n \in \mathbb{N}\,(a_{n} \in A) \to \forall n \in \mathbb{N}\,(a_{n} \in A)$
\item $\varphi_{112} :\equiv f \circ a(n) \cdot g \circ a^{\ast}(n) = f \circ a(n) \cdot g \circ a(n)^{-1}$
\item $\varphi_{113} :\equiv \mathrm{Conv}((f / g) \circ a,f(c) \cdot g(c)^{-1}) \to \mathrm{Conv}((f / g) \circ a,(f / g)(c))$
\item $\varphi_{114} :\equiv f \circ a(n) \cdot g \circ a^{\ast}(n) = f \circ a(n) \cdot g \circ a^{\ast}(n)$
\item $\varphi_{115} :\equiv n \in \mathbb{N} \to f \circ a \otimes g \circ a^{\ast}(n) = f(a_{n}) \cdot g(a_{n})^{-1}$
\item $\varphi_{116} :\equiv n \in \mathbb{N} \to f(a_{n}) \cdot g(a_{n})^{-1} = f \circ a \otimes g \circ a^{\ast}(n)$
\item $\varphi_{117} :\equiv \mathrm{Conv}(a,c) \to n \in \mathbb{N} \to (f / g) \circ a(n) = f \circ a \otimes g \circ a^{\ast}(n)$
\item $\varphi_{118} :\equiv \mathrm{Conv}(a,c) \to n \in \mathbb{N} \to f \circ a \otimes g \circ a^{\ast}(n) = (f / g) \circ a(n)$
\item $\varphi_{119} :\equiv \forall n \in \mathbb{N}\,(g \circ a(n) = g(a_{n})) \to n \in \mathbb{N} \to g \circ a(n) = g(a_{n})$
\item $\varphi_{120} :\equiv \forall n \in \mathbb{N}\,(f \circ a(n) = f(a_{n})) \to n \in \mathbb{N} \to f \circ a(n) = f(a_{n})$
\item $\varphi_{121} :\equiv \forall a\,(\mathrm{Seq}(a) \to \forall n \in \mathbb{N}\,(a_{n} \neq 0) \to \mathrm{Seq}(a^{\ast}) \land \forall n \in \mathbb{N}\,(a^{\ast}(n) = a_{n}^{-1}))$
\item $\varphi_{122} :\equiv n \in \mathbb{N} \to f \circ a(n) \cdot g \circ a(n)^{-1} = f(a_{n}) \cdot g \circ a(n)^{-1}$
\item $\varphi_{123} :\equiv n \in \mathbb{N} \to f \circ a \otimes g \circ a^{\ast}(n) = f(a_{n}) \cdot g \circ a(n)^{-1}$
\item $\varphi_{124} :\equiv \forall y\,(f \circ a \otimes g \circ a^{\ast}(n) = y \to y = f \circ a \otimes g \circ a^{\ast}(n))$
\item $\varphi_{125} :\equiv \mathrm{ContIn}(g,A,c) \to \forall x \in A\,(g(x) \neq 0) \to \mathrm{Seq}(a) \to \varphi_{4}$
\item $\varphi_{126} :\equiv \forall n \in \mathbb{N}\,(a_{n} \in \mathrm{dom}(g)) \to \mathrm{Seq}(g \circ a) \land \forall n \in \mathbb{N}\,(g \circ a(n) = g(a_{n}))$
\item $\varphi_{127} :\equiv \forall n \in \mathbb{N}\,(a_{n} \in \mathrm{dom}(f)) \to \mathrm{Seq}(f \circ a) \land \forall n \in \mathbb{N}\,(f \circ a(n) = f(a_{n}))$
\item $\varphi_{128} :\equiv n \in \mathbb{N} \to f \circ a \otimes g \circ a^{\ast}(n) = f \circ a(n) \cdot g \circ a(n)^{-1}$
\item $\varphi_{129} :\equiv \forall f\,\forall A\,\forall c\,\forall a\,(\mathrm{ContIn}(f,A,c) \to \mathrm{Seq}(a) \to \varphi_{49})$
\item $\varphi_{130} :\equiv n \in \mathbb{N} \to f \circ a \otimes g \circ a^{\ast}(n) = f \circ a(n) \cdot g \circ a^{\ast}(n)$
\item $\varphi_{131} :\equiv \forall n\,(\varphi_{6} \land \mathrm{Seq}(a) \land \forall n \in \mathbb{N}\,(a_{n} \in A) \land \mathrm{Conv}(a,c))$
\item $\varphi_{132} :\equiv \forall n \in \mathbb{N}\,(g \circ a(n) \neq 0) \to \mathrm{Conv}(f \circ a \otimes g \circ a^{\ast},f(c) \cdot g(c)^{-1})$
\item $\varphi_{133} :\equiv \forall n \in \mathbb{N}\,(f \circ a \otimes g \circ a^{\ast}(n) = f \circ a(n) \cdot g \circ a^{\ast}(n))$
\item $\varphi_{134} :\equiv \forall y\,\forall z\,(f \circ a \otimes g \circ a^{\ast}(n) = y \to y = z \to f \circ a \otimes g \circ a^{\ast}(n) = z)$
\item $\varphi_{135} :\equiv \forall z\,((f / g) \circ a(n) = (f / g)(a_{n}) \to (f / g)(a_{n}) = z \to (f / g) \circ a(n) = z)$
\item $\varphi_{136} :\equiv \forall a\,\forall \alpha\,\forall b\,\forall \beta\,(\mathrm{Conv}(a,\alpha) \to \mathrm{Conv}(b,\beta) \to \beta \neq 0 \to \forall n \in \mathbb{N}\,(b_{n} \neq 0) \to \mathrm{Conv}(a \otimes b^{\ast},\alpha \cdot \beta^{-1}))$
\item $\varphi_{137} :\equiv \forall w\,(f(c) \cdot g(c)^{-1} = w \to \mathrm{Conv}((f / g) \circ a,f(c) \cdot g(c)^{-1}) \to \mathrm{Conv}((f / g) \circ a,w))$
\item $\varphi_{138} :\equiv \forall y\,(g(a_{n}) = y \to y = g(a_{n})) \to g(a_{n}) = g \circ a(n) \to g \circ a(n) = g(a_{n})$
\item $\varphi_{139} :\equiv \mathrm{ContIn}(g,A,c) \to \forall x \in A\,(g(x) \neq 0) \to \mathrm{Seq}(a) \to \varphi_{24}$
\item $\varphi_{140} :\equiv \mathrm{ContIn}(g,A,c) \to \forall x \in A\,(g(x) \neq 0) \to \mathrm{Seq}(a) \to \varphi_{29}$
\item $\varphi_{141} :\equiv \mathrm{ContIn}(f,A,c) \to \mathrm{ContIn}(g,A,c) \to \forall x \in A\,(g(x) \neq 0) \to \varphi_{14}$
\item $\varphi_{142} :\equiv \mathrm{ContIn}(f,A,c) \to \mathrm{ContIn}(g,A,c) \to \forall x \in A\,(g(x) \neq 0) \to \varphi_{15}$
\item $\varphi_{143} :\equiv \mathrm{ContIn}(f,A,c) \to \mathrm{ContIn}(g,A,c) \to \forall x \in A\,(g(x) \neq 0) \to \varphi_{17}$
\item $\varphi_{144} :\equiv \mathrm{ContIn}(g,A,c) \to \forall x \in A\,(g(x) \neq 0) \to \mathrm{Seq}(a) \to \varphi_{40}$
\item $\varphi_{145} :\equiv \mathrm{ContIn}(g,A,c) \to \forall x \in A\,(g(x) \neq 0) \to \mathrm{Seq}(a) \to \varphi_{42}$
\item $\varphi_{146} :\equiv \mathrm{ContIn}(g,A,c) \to \forall x \in A\,(g(x) \neq 0) \to \mathrm{Seq}(a) \to \varphi_{44}$
\item $\varphi_{147} :\equiv \mathrm{ContIn}(f,A,c) \to \mathrm{ContIn}(g,A,c) \to \forall x \in A\,(g(x) \neq 0) \to \varphi_{21}$
\item $\varphi_{148} :\equiv \mathrm{ContIn}(f,A,c) \to \mathrm{ContIn}(g,A,c) \to \forall x \in A\,(g(x) \neq 0) \to \varphi_{23}$
\item $\varphi_{149} :\equiv \mathrm{ContIn}(f,A,c) \to \mathrm{ContIn}(g,A,c) \to \forall x \in A\,(g(x) \neq 0) \to \varphi_{25}$
\item $\varphi_{150} :\equiv \mathrm{ContIn}(g,A,c) \to \forall x \in A\,(g(x) \neq 0) \to \mathrm{Seq}(a) \to \varphi_{48}$
\item $\varphi_{151} :\equiv \mathrm{ContIn}(g,A,c) \to \forall x \in A\,(g(x) \neq 0) \to \mathrm{Seq}(a) \to \varphi_{49}$
\item $\varphi_{152} :\equiv \forall x \in A\,(g(x) \neq 0) \to \mathrm{Seq}(a) \to \forall n \in \mathbb{N}\,(a_{n} \in A) \to \varphi_{5}$
\item $\varphi_{153} :\equiv \forall a\,(\mathrm{Seq}(a) \to \forall n \in \mathbb{N}\,(a_{n} \in A) \to \varphi_{5}) \to \mathrm{ContIn}((f / g),A,c)$
\item $\varphi_{154} :\equiv \forall y\,(g \circ a(n) = y \to y = g \circ a(n)) \to g \circ a(n) = g(a_{n}) \to g(a_{n}) = g \circ a(n)$
\item $\varphi_{155} :\equiv \mathrm{ContIn}(f,A,c) \to \mathrm{ContIn}(g,A,c) \to \forall x \in A\,(g(x) \neq 0) \to \varphi_{30}$
\item $\varphi_{156} :\equiv \mathrm{ContIn}(f,A,c) \to \mathrm{ContIn}(g,A,c) \to \forall x \in A\,(g(x) \neq 0) \to \varphi_{31}$
\item $\varphi_{157} :\equiv \mathrm{ContIn}(f,A,c) \to \mathrm{ContIn}(g,A,c) \to \forall x \in A\,(g(x) \neq 0) \to \varphi_{32}$
\item $\varphi_{158} :\equiv \mathrm{ContIn}(f,A,c) \to \mathrm{ContIn}(g,A,c) \to \forall x \in A\,(g(x) \neq 0) \to \varphi_{37}$
\item $\varphi_{159} :\equiv \mathrm{ContIn}(g,A,c) \to \forall x \in A\,(g(x) \neq 0) \to \mathrm{Seq}(a) \to \varphi_{52}$
\item $\varphi_{160} :\equiv \mathrm{ContIn}(f,A,c) \to \mathrm{ContIn}(g,A,c) \to \forall x \in A\,(g(x) \neq 0) \to \varphi_{39}$
\item $\varphi_{161} :\equiv \mathrm{ContIn}(g,A,c) \to \forall x \in A\,(g(x) \neq 0) \to \mathrm{Seq}(a) \to \varphi_{59}$
\item $\varphi_{162} :\equiv \forall x \in A\,(g(x) \neq 0) \to \forall a\,(\mathrm{Seq}(a) \to \forall n \in \mathbb{N}\,(a_{n} \in A) \to \varphi_{5})$
\item $\varphi_{163} :\equiv \mathrm{ContIn}(f,A,c) \to \mathrm{ContIn}(g,A,c) \to \forall x \in A\,(g(x) \neq 0) \to \varphi_{41}$
\item $\varphi_{164} :\equiv \mathrm{ContIn}(g,A,c) \to \forall x \in A\,(g(x) \neq 0) \to \mathrm{Seq}(a) \to \varphi_{61}$
\item $\varphi_{165} :\equiv \mathrm{ContIn}(g,A,c) \to \forall x \in A\,(g(x) \neq 0) \to \mathrm{Seq}(a) \to \varphi_{62}$
\item $\varphi_{166} :\equiv \mathrm{Conv}(f \circ a \otimes g \circ a^{\ast},f(c) \cdot g(c)^{-1}) \to \mathrm{Conv}((f / g) \circ a,f(c) \cdot g(c)^{-1})$
\item $\varphi_{167} :\equiv \forall w\,(g(a_{n}) = w \to g(a_{n}) = 0 \to w = 0) \to g(a_{n}) = g \circ a(n) \to g(a_{n}) = 0 \to g \circ a(n) = 0$
\item $\varphi_{168} :\equiv \mathrm{ContIn}(f,A,c) \to \mathrm{ContIn}(g,A,c) \to \forall x \in A\,(g(x) \neq 0) \to \varphi_{45}$
\item $\varphi_{169} :\equiv \mathrm{ContIn}(f,A,c) \to \mathrm{ContIn}(g,A,c) \to \forall x \in A\,(g(x) \neq 0) \to \varphi_{46}$
\item $\varphi_{170} :\equiv \mathrm{ContIn}(g,A,c) \to \forall x \in A\,(g(x) \neq 0) \to \mathrm{Seq}(a) \to \varphi_{69}$
\item $\varphi_{171} :\equiv \mathrm{ContIn}(g,A,c) \to \forall x \in A\,(g(x) \neq 0) \to \mathrm{Seq}(a) \to \varphi_{71}$
\item $\varphi_{172} :\equiv \mathrm{ContIn}(g,A,c) \to \forall x \in A\,(g(x) \neq 0) \to \mathrm{Seq}(a) \to \varphi_{72}$
\item $\varphi_{173} :\equiv (\forall x \in A\,(g(x) \neq 0) \to c \in A \to g(c) \neq 0) \to \mathrm{ContIn}(f,A,c) \to \varphi_{20}$
\item $\varphi_{174} :\equiv (\forall x \in A\,(g(x) \neq 0) \to u \in A \to g(u) \neq 0) \to \mathrm{ContIn}(f,A,c) \to \varphi_{26}$
\item $\varphi_{175} :\equiv \forall n \in \mathbb{N}\,(g \circ a^{\ast}(n) = g \circ a(n)^{-1}) \to n \in \mathbb{N} \to g \circ a^{\ast}(n) = g \circ a(n)^{-1}$
\item $\varphi_{176} :\equiv \forall n \in \mathbb{N}\,((f / g) \circ a(n) = (f / g)(a_{n})) \to n \in \mathbb{N} \to (f / g) \circ a(n) = (f / g)(a_{n})$
\item $\varphi_{177} :\equiv \forall x \in A\,(g(x) \neq 0) \to \mathrm{Seq}(a) \to \forall n \in \mathbb{N}\,(a_{n} \in A) \to \varphi_{22}$
\item $\varphi_{178} :\equiv \mathrm{ContIn}(g,A,c) \to \forall x \in A\,(g(x) \neq 0) \to \mathrm{Seq}(a) \to \varphi_{74}$
\item $\varphi_{179} :\equiv (f / g)(a_{n}) = f(a_{n}) \cdot g(a_{n})^{-1} \to (f / g) \circ a(n) = f(a_{n}) \cdot g(a_{n})^{-1}$
\item $\varphi_{180} :\equiv \forall \alpha\,\forall b\,\forall \beta\,(\mathrm{Conv}(f \circ a,\alpha) \to \mathrm{Conv}(b,\beta) \to \beta \neq 0 \to \forall n \in \mathbb{N}\,(b_{n} \neq 0) \to \mathrm{Conv}(f \circ a \otimes b^{\ast},\alpha \cdot \beta^{-1}))$
\item $\varphi_{181} :\equiv \mathrm{ContIn}(g,A,c) \to \forall x \in A\,(g(x) \neq 0) \to \mathrm{Seq}(a) \to \varphi_{82}$
\item $\varphi_{182} :\equiv \forall w\,(g \circ a(n) = w \to g \circ a(n) = 0 \to w = 0) \to g \circ a(n) = g(a_{n}) \to g \circ a(n) = 0 \to g(a_{n}) = 0$
\item $\varphi_{183} :\equiv \forall b\,(\mathrm{Seq}(f \circ a) \to \mathrm{Seq}(b) \to \mathrm{Seq}(f \circ a \otimes b) \land \forall n \in \mathbb{N}\,(f \circ a \otimes b(n) = f \circ a(n) \cdot b_{n}))$
\item $\varphi_{184} :\equiv \mathrm{ContIn}(g,A,c) \to \forall x \in A\,(g(x) \neq 0) \to \mathrm{Seq}(a) \to \varphi_{85}$
\item $\varphi_{185} :\equiv \mathrm{ContIn}(g,A,c) \to \forall x \in A\,(g(x) \neq 0) \to \mathrm{Seq}(a) \to \varphi_{86}$
\item $\varphi_{186} :\equiv \mathrm{ContIn}(g,A,c) \to \forall x \in A\,(g(x) \neq 0) \to \mathrm{Seq}(a) \to \varphi_{87}$
\item $\varphi_{187} :\equiv \forall x \in A\,(g(x) \neq 0) \to \mathrm{Seq}(a) \to \forall n \in \mathbb{N}\,(a_{n} \in A) \to \varphi_{47}$
\item $\varphi_{188} :\equiv \forall x \in A\,(g(x) \neq 0) \to \mathrm{Seq}(a) \to \forall n \in \mathbb{N}\,(a_{n} \in A) \to \varphi_{50}$
\item $\varphi_{189} :\equiv f(a_{n}) \cdot g \circ a(n)^{-1} = f(a_{n}) \cdot z^{-1} \to f(a_{n}) \cdot g \circ a(n)^{-1} = f(a_{n}) \cdot w^{-1}$
\item $\varphi_{190} :\equiv \forall b\,\forall \beta\,(\mathrm{Conv}(f \circ a,f(c)) \to \mathrm{Conv}(b,\beta) \to \beta \neq 0 \to \forall n \in \mathbb{N}\,(b_{n} \neq 0) \to \mathrm{Conv}(f \circ a \otimes b^{\ast},f(c) \cdot \beta^{-1}))$
\item $\varphi_{191} :\equiv \forall x \in A\,(g(x) \neq 0) \to \mathrm{Seq}(a) \to \forall n \in \mathbb{N}\,(a_{n} \in A) \to \varphi_{58}$
\item $\varphi_{192} :\equiv \forall x \in A\,(g(x) \neq 0) \to \mathrm{Seq}(a) \to \forall n \in \mathbb{N}\,(a_{n} \in A) \to \varphi_{60}$
\item $\varphi_{193} :\equiv \varphi_{6} \land \mathrm{Seq}(a) \land \forall n \in \mathbb{N}\,(a_{n} \in A) \land \mathrm{Conv}(a,c) \to n \in \mathbb{N} \to a_{n} \in \mathrm{dom}(g)$
\item $\varphi_{194} :\equiv \varphi_{6} \land \mathrm{Seq}(a) \land \forall n \in \mathbb{N}\,(a_{n} \in A) \land \mathrm{Conv}(a,c) \to n \in \mathbb{N} \to a_{n} \in \mathrm{dom}(f)$
\item $\varphi_{195} :\equiv \forall f\,\forall A\,\forall c\,(\mathrm{RFn}(f) \to A \subseteq \mathrm{dom}(f) \to c \in A \to \forall a\,(\mathrm{Seq}(a) \to \varphi_{49}) \to \mathrm{ContIn}(f,A,c))$
\item $\varphi_{196} :\equiv \mathrm{ContIn}(g,A,c) \to \forall x \in A\,(g(x) \neq 0) \to \mathrm{Seq}(a) \to \varphi_{94}$
\item $\varphi_{197} :\equiv \mathrm{ContIn}(g,A,c) \to \forall x \in A\,(g(x) \neq 0) \to \mathrm{Seq}(a) \to \varphi_{95}$
\item $\varphi_{198} :\equiv \mathrm{ContIn}(g,A,c) \to \forall x \in A\,(g(x) \neq 0) \to \mathrm{Seq}(a) \to \varphi_{96}$
\item $\varphi_{199} :\equiv \forall x \in A\,(g(x) \neq 0) \to \mathrm{Seq}(a) \to \forall n \in \mathbb{N}\,(a_{n} \in A) \to \varphi_{63}$
\item $\varphi_{200} :\equiv \forall x \in A\,(g(x) \neq 0) \to \mathrm{Seq}(a) \to \forall n \in \mathbb{N}\,(a_{n} \in A) \to \varphi_{64}$
\item $\varphi_{201} :\equiv \varphi_{6} \land \mathrm{Seq}(a) \land \forall n \in \mathbb{N}\,(a_{n} \in A) \land \mathrm{Conv}(a,c) \to n \in \mathbb{N} \to g \circ a(n) \neq 0$
\item $\varphi_{202} :\equiv \forall z\,((f / g) \circ a(n) = f(a_{n}) \cdot g(a_{n})^{-1} \to f(a_{n}) \cdot g(a_{n})^{-1} = z \to (f / g) \circ a(n) = z)$
\item $\varphi_{203} :\equiv \forall x \in A\,(g(x) \neq 0) \to \mathrm{Seq}(a) \to \forall n \in \mathbb{N}\,(a_{n} \in A) \to \varphi_{80}$
\item $\varphi_{204} :\equiv \forall x \in A\,(g(x) \neq 0) \to \mathrm{Seq}(a) \to \forall n \in \mathbb{N}\,(a_{n} \in A) \to \varphi_{81}$
\item $\varphi_{205} :\equiv \mathrm{ContIn}(f,A,c) \to \mathrm{ContIn}(g,A,c) \to \forall x \in A\,(g(x) \neq 0) \to \varphi_{98}$
\item $\varphi_{206} :\equiv \varphi_{6} \land \mathrm{Seq}(a) \land \forall n \in \mathbb{N}\,(a_{n} \in A) \land \mathrm{Conv}(a,c) \to n \in \mathbb{N} \to a_{n} \in \mathrm{dom}((f / g))$
\item $\varphi_{207} :\equiv f \circ a(n) \cdot g \circ a^{\ast}(n) = f \circ a(n) \cdot z \to f \circ a(n) \cdot g \circ a^{\ast}(n) = f \circ a(n) \cdot w$
\item $\varphi_{208} :\equiv \mathrm{ContIn}(f,A,c) \to \mathrm{ContIn}(g,A,c) \to \forall x \in A\,(g(x) \neq 0) \to \varphi_{100}$
\item $\varphi_{209} :\equiv \forall b\,\forall \alpha\,(\mathrm{Seq}(b) \to \forall n \in \mathbb{N}\,(f \circ a \otimes g \circ a^{\ast}(n) = b_{n}) \to \mathrm{Conv}(f \circ a \otimes g \circ a^{\ast},\alpha) \to \mathrm{Conv}(b,\alpha))$
\item $\varphi_{210} :\equiv \forall x \in A\,(g(x) \neq 0) \to \mathrm{Seq}(a) \to \forall n \in \mathbb{N}\,(a_{n} \in A) \to \varphi_{84}$
\item $\varphi_{211} :\equiv (f / g) \circ a(n) = f \circ a \otimes g \circ a^{\ast}(n) \to f \circ a \otimes g \circ a^{\ast}(n) = (f / g) \circ a(n)$
\item $\varphi_{212} :\equiv \forall x \in A\,(g(x) \neq 0) \to \mathrm{Seq}(a) \to \forall n \in \mathbb{N}\,(a_{n} \in A) \to \varphi_{89}$
\item $\varphi_{213} :\equiv \forall x \in A\,(g(x) \neq 0) \to \mathrm{Seq}(a) \to \forall n \in \mathbb{N}\,(a_{n} \in A) \to \varphi_{90}$
\item $\varphi_{214} :\equiv \forall x \in A\,(g(x) \neq 0) \to \mathrm{Seq}(a) \to \forall n \in \mathbb{N}\,(a_{n} \in A) \to \varphi_{91}$
\item $\varphi_{215} :\equiv f \circ a(n) \cdot g \circ a(n)^{-1} = z \cdot g \circ a(n)^{-1} \to f \circ a(n) \cdot g \circ a(n)^{-1} = w \cdot g \circ a(n)^{-1}$
\item $\varphi_{216} :\equiv \forall x \in A\,(g(x) \neq 0) \to \mathrm{Seq}(a) \to \forall n \in \mathbb{N}\,(a_{n} \in A) \to \varphi_{93}$
\item $\varphi_{217} :\equiv \forall x \in A\,(g(x) \neq 0) \to \mathrm{Seq}(a) \to \forall n \in \mathbb{N}\,(a_{n} \in A) \to \mathrm{Conv}(a,c) \to \varphi_{73}$
\item $\varphi_{218} :\equiv (\forall x \in A\,(g(x) \neq 0) \to a_{n} \in A \to g(a_{n}) \neq 0) \to \mathrm{ContIn}(f,A,c) \to \varphi_{70}$
\item $\varphi_{219} :\equiv \forall x \in A\,(g(x) \neq 0) \to \mathrm{Seq}(a) \to \forall n \in \mathbb{N}\,(a_{n} \in A) \to \varphi_{97}$
\item $\varphi_{220} :\equiv (\mathrm{ContIn}(f,A,c) \to \varphi_{18}) \to \mathrm{ContIn}(f,A,c) \to \varphi_{33}$
\item $\varphi_{221} :\equiv (\mathrm{ContIn}(f,A,c) \to \varphi_{18}) \to \mathrm{ContIn}(f,A,c) \to \varphi_{34}$
\item $\varphi_{222} :\equiv \mathrm{Seq}(a) \to \forall n \in \mathbb{N}\,(a_{n} \in A) \to \mathrm{Conv}(a,c) \to n \in \mathbb{N} \to \varphi_{112}$
\item $\varphi_{223} :\equiv \mathrm{RFn}((f / g)) \to \mathrm{Seq}(a) \to \forall n \in \mathbb{N}\,(a_{n} \in \mathrm{dom}((f / g))) \to \varphi_{102}$
\item $\varphi_{224} :\equiv f(a_{n}) \cdot g(a_{n})^{-1} = f \circ a \otimes g \circ a^{\ast}(n) \to (f / g) \circ a(n) = f \circ a \otimes g \circ a^{\ast}(n)$
\item $\varphi_{225} :\equiv \forall x \in A\,(g(x) \neq 0) \to \mathrm{Seq}(a) \to \forall n \in \mathbb{N}\,(a_{n} \in A) \to \varphi_{117}$
\item $\varphi_{226} :\equiv \forall x \in A\,(g(x) \neq 0) \to \mathrm{Seq}(a) \to \forall n \in \mathbb{N}\,(a_{n} \in A) \to \varphi_{118}$
\item $\varphi_{227} :\equiv \forall x \in A\,(g(x) \neq 0) \to \mathrm{Seq}(a) \to \forall n \in \mathbb{N}\,(a_{n} \in A) \to \mathrm{Conv}(a,c) \to \varphi_{92}$
\item $\varphi_{228} :\equiv \forall \beta\,(\mathrm{Conv}(f \circ a,f(c)) \to \mathrm{Conv}(g \circ a,\beta) \to \beta \neq 0 \to \forall n \in \mathbb{N}\,(g \circ a(n) \neq 0) \to \mathrm{Conv}(f \circ a \otimes g \circ a^{\ast},f(c) \cdot \beta^{-1}))$
\item $\varphi_{229} :\equiv f(a_{n}) \cdot g \circ a(n)^{-1} = f(a_{n}) \cdot g(a_{n})^{-1} \to f \circ a \otimes g \circ a^{\ast}(n) = f(a_{n}) \cdot g(a_{n})^{-1}$
\item $\varphi_{230} :\equiv f(a_{n}) \cdot g \circ a(n)^{-1} = f(a_{n}) \cdot g \circ a(n)^{-1} \to f(a_{n}) \cdot g \circ a(n)^{-1} = f(a_{n}) \cdot g(a_{n})^{-1}$
\item $\varphi_{231} :\equiv \forall x \in A\,(g(x) \neq 0) \to \mathrm{Seq}(a) \to \forall n \in \mathbb{N}\,(a_{n} \in A) \to \mathrm{Conv}(a,c) \to \varphi_{99}$
\item $\varphi_{232} :\equiv \mathrm{Seq}(a) \to \forall n \in \mathbb{N}\,(a_{n} \in A) \to \mathrm{Conv}(a,c) \to n \in \mathbb{N} \to \varphi_{130}$
\item $\varphi_{233} :\equiv \forall x \in A\,(g(x) \neq 0) \to \mathrm{Seq}(a) \to \forall n \in \mathbb{N}\,(a_{n} \in A) \to \mathrm{Conv}(a,c) \to \varphi_{102}$
\item $\varphi_{234} :\equiv \mathrm{ContIn}(g,A,c) \leftrightarrow \mathrm{RFn}(g) \land A \subseteq \mathrm{dom}(g) \land c \in A \land \forall \varepsilon \in \mathbb{R}\,(0 < \varepsilon \to \exists \delta\,(\delta \in \mathbb{R} \land 0 < \delta \land \forall x \in A\,(\lvert x - c \rvert < \delta \to \lvert g(x) - g(c) \rvert < \varepsilon)))$
\item $\varphi_{235} :\equiv \mathrm{ContIn}(f,A,c) \leftrightarrow \mathrm{RFn}(f) \land A \subseteq \mathrm{dom}(f) \land c \in A \land \forall \varepsilon \in \mathbb{R}\,(0 < \varepsilon \to \exists \delta\,(\delta \in \mathbb{R} \land 0 < \delta \land \forall x \in A\,(\lvert x - c \rvert < \delta \to \lvert f(x) - f(c) \rvert < \varepsilon)))$
\item $\varphi_{236} :\equiv (\mathrm{ContIn}(f,A,c) \to \mathrm{ContIn}(g,A,c) \to \varphi_{19}) \to \mathrm{ContIn}(f,A,c) \to \varphi_{27}$
\item $\varphi_{237} :\equiv \forall x \in A\,(g(x) \neq 0) \to \mathrm{Seq}(a) \to \forall n \in \mathbb{N}\,(a_{n} \in A) \to \mathrm{Conv}(a,c) \to \varphi_{106}$
\item $\varphi_{238} :\equiv \forall g\,(\mathrm{RFn}(f) \to \mathrm{RFn}(g) \to \mathrm{RFn}((f / g)) \land \varphi_{39} \land \varphi_{41})$
\item $\varphi_{239} :\equiv f \circ a \otimes g \circ a^{\ast}(n) = f(a_{n}) \cdot g(a_{n})^{-1} \to f(a_{n}) \cdot g(a_{n})^{-1} = f \circ a \otimes g \circ a^{\ast}(n)$
\item $\varphi_{240} :\equiv \forall c\,(\mathrm{RFn}((f / g)) \to A \subseteq \mathrm{dom}((f / g)) \to c \in A \to \varphi_{153})$
\item $\varphi_{241} :\equiv \mathrm{Seq}(a) \to \forall n \in \mathbb{N}\,(a_{n} \in A) \to \mathrm{Conv}(a,c) \to n \in \mathbb{N} \to \varphi_{133}$
\item $\varphi_{242} :\equiv \forall x \in A\,(g(x) \neq 0) \to \mathrm{Seq}(a) \to \forall n \in \mathbb{N}\,(a_{n} \in A) \to \mathrm{Conv}(a,c) \to \varphi_{115}$
\item $\varphi_{243} :\equiv \forall x \in A\,(g(x) \neq 0) \to \mathrm{Seq}(a) \to \forall n \in \mathbb{N}\,(a_{n} \in A) \to \mathrm{Conv}(a,c) \to \varphi_{116}$
\item $\varphi_{244} :\equiv (\mathrm{ContIn}(f,A,c) \to \varphi_{53}) \to \varphi_{6} \to u \in A \to u \in \mathrm{dom}((f / g))$
\item $\varphi_{245} :\equiv \forall x \in A\,(g(x) \neq 0) \to \mathrm{Seq}(a) \to \forall n \in \mathbb{N}\,(a_{n} \in A) \to \mathrm{Conv}(a,c) \to \varphi_{122}$
\item $\varphi_{246} :\equiv \forall x \in A\,(g(x) \neq 0) \to \mathrm{Seq}(a) \to \forall n \in \mathbb{N}\,(a_{n} \in A) \to \mathrm{Conv}(a,c) \to \varphi_{123}$
\item $\varphi_{247} :\equiv f \circ a(n) \cdot g \circ a(n)^{-1} = f(a_{n}) \cdot g \circ a(n)^{-1} \to f \circ a \otimes g \circ a^{\ast}(n) = f(a_{n}) \cdot g \circ a(n)^{-1}$
\item $\varphi_{248} :\equiv \mathrm{Conv}(f \circ a,f(c)) \to \mathrm{Conv}(g \circ a,g(c)) \to g(c) \neq 0 \to \varphi_{132}$
\item $\varphi_{249} :\equiv \forall w\,(g \circ a(n) = w \to f(a_{n}) \cdot g \circ a(n)^{-1} = f(a_{n}) \cdot g \circ a(n)^{-1} \to f(a_{n}) \cdot g \circ a(n)^{-1} = f(a_{n}) \cdot w^{-1})$
\item $\varphi_{250} :\equiv (\forall n \in \mathbb{N}\,(a_{n} \in A) \to n \in \mathbb{N} \to a_{n} \in A) \to \mathrm{ContIn}(f,A,c) \to \mathrm{ContIn}(g,A,c) \to \varphi_{101}$
\item $\varphi_{251} :\equiv \forall x \in A\,(g(x) \neq 0) \to \mathrm{Seq}(a) \to \forall n \in \mathbb{N}\,(a_{n} \in A) \to \mathrm{Conv}(a,c) \to \varphi_{128}$
\item $\varphi_{252} :\equiv \varphi_{103} \to f \circ a(n) \cdot g \circ a(n)^{-1} = f(a_{n}) \cdot g \circ a(n)^{-1}$
\item $\varphi_{253} :\equiv \forall x \in A\,(g(x) \neq 0) \to \mathrm{Seq}(a) \to \forall n \in \mathbb{N}\,(a_{n} \in A) \to \mathrm{Conv}(a,c) \to \varphi_{130}$
\item $\varphi_{254} :\equiv \forall u\,(\varphi_{6} \to u \in A \to u \in \mathrm{dom}((f / g))) \to \forall u\,(\varphi_{6}) \to \forall u \in A\,(u \in \mathrm{dom}((f / g)))$
\item $\varphi_{255} :\equiv \forall z\,(f \circ a \otimes g \circ a^{\ast}(n) = f(a_{n}) \cdot g \circ a(n)^{-1} \to f(a_{n}) \cdot g \circ a(n)^{-1} = z \to f \circ a \otimes g \circ a^{\ast}(n) = z)$
\item $\varphi_{256} :\equiv \varphi_{112} \to f \circ a \otimes g \circ a^{\ast}(n) = f \circ a(n) \cdot g \circ a(n)^{-1}$
\item $\varphi_{257} :\equiv \forall a\,(\mathrm{ContIn}(g,A,c) \to \mathrm{Seq}(a) \to \varphi_{48}) \to \mathrm{ContIn}(g,A,c) \to \mathrm{Seq}(a) \to \varphi_{48}$
\item $\varphi_{258} :\equiv \forall a\,(\mathrm{ContIn}(f,A,c) \to \mathrm{Seq}(a) \to \varphi_{49}) \to \mathrm{ContIn}(f,A,c) \to \mathrm{Seq}(a) \to \varphi_{49}$
\item $\varphi_{259} :\equiv \mathrm{ContIn}(g,A,c) \to \forall x \in A\,(g(x) \neq 0) \to \mathrm{RFn}((f / g)) \land \varphi_{39} \land \varphi_{41}$
\item $\varphi_{260} :\equiv (\forall n \in \mathbb{N}\,(a_{n} \in A) \to n \in \mathbb{N} \to a_{n} \in A) \to \mathrm{ContIn}(f,A,c) \to \varphi_{146}$
\item $\varphi_{261} :\equiv \varphi_{6} \land \mathrm{Seq}(a) \land \forall n \in \mathbb{N}\,(a_{n} \in A) \land \mathrm{Conv}(a,c) \to n \in \mathbb{N} \to f \circ a \otimes g \circ a^{\ast}(n) = (f / g) \circ a(n)$
\item $\varphi_{262} :\equiv \forall z\,(f \circ a \otimes g \circ a^{\ast}(n) = f \circ a(n) \cdot g \circ a(n)^{-1} \to f \circ a(n) \cdot g \circ a(n)^{-1} = z \to f \circ a \otimes g \circ a^{\ast}(n) = z)$
\item $\varphi_{263} :\equiv \forall w\,(f \circ a(n) = w \to \varphi_{103} \to f \circ a(n) \cdot g \circ a(n)^{-1} = w \cdot g \circ a(n)^{-1})$
\item $\varphi_{264} :\equiv \forall c\,\forall a\,(\mathrm{ContIn}(g,A,c) \to \mathrm{Seq}(a) \to \varphi_{48}) \to \forall a\,(\mathrm{ContIn}(g,A,c) \to \mathrm{Seq}(a) \to \varphi_{48})$
\item $\varphi_{265} :\equiv \forall c\,\forall a\,(\mathrm{ContIn}(f,A,c) \to \mathrm{Seq}(a) \to \varphi_{49}) \to \forall a\,(\mathrm{ContIn}(f,A,c) \to \mathrm{Seq}(a) \to \varphi_{49})$
\item $\varphi_{266} :\equiv \forall n \in \mathbb{N}\,(a_{n} \in A) \to \mathrm{Conv}(a,c) \to \mathrm{Seq}(f \circ a \otimes g \circ a^{\ast}) \land \varphi_{133}$
\item $\varphi_{267} :\equiv \forall z\,(f \circ a \otimes g \circ a^{\ast}(n) = f \circ a(n) \cdot g \circ a^{\ast}(n) \to f \circ a(n) \cdot g \circ a^{\ast}(n) = z \to f \circ a \otimes g \circ a^{\ast}(n) = z)$
\item $\varphi_{268} :\equiv \forall \alpha\,(\mathrm{Seq}((f / g) \circ a) \to \varphi_{92} \to \mathrm{Conv}(f \circ a \otimes g \circ a^{\ast},\alpha) \to \mathrm{Conv}((f / g) \circ a,\alpha))$
\item $\varphi_{269} :\equiv \mathrm{Seq}(f \circ a) \to \mathrm{Seq}(g \circ a^{\ast}) \to \mathrm{Seq}(f \circ a \otimes g \circ a^{\ast}) \land \varphi_{133}$
\item $\varphi_{270} :\equiv \forall A\,\forall c\,\forall a\,(\mathrm{ContIn}(g,A,c) \to \mathrm{Seq}(a) \to \varphi_{48}) \to \forall c\,\forall a\,(\mathrm{ContIn}(g,A,c) \to \mathrm{Seq}(a) \to \varphi_{48})$
\item $\varphi_{271} :\equiv \forall A\,\forall c\,\forall a\,(\mathrm{ContIn}(f,A,c) \to \mathrm{Seq}(a) \to \varphi_{49}) \to \forall c\,\forall a\,(\mathrm{ContIn}(f,A,c) \to \mathrm{Seq}(a) \to \varphi_{49})$
\item $\varphi_{272} :\equiv \forall w\,(g \circ a^{\ast}(n) = w \to \varphi_{114} \to f \circ a(n) \cdot g \circ a^{\ast}(n) = f \circ a(n) \cdot w)$
\item $\varphi_{273} :\equiv (\varphi_{6} \to \forall u\,(\varphi_{6})) \to \mathrm{ContIn}(f,A,c) \to \mathrm{ContIn}(g,A,c) \to \varphi_{19}$
\item $\varphi_{274} :\equiv \varphi_{129} \to \forall A\,\forall c\,\forall a\,(\mathrm{ContIn}(g,A,c) \to \mathrm{Seq}(a) \to \varphi_{48})$
\item $\varphi_{275} :\equiv \varphi_{129} \to \forall A\,\forall c\,\forall a\,(\mathrm{ContIn}(f,A,c) \to \mathrm{Seq}(a) \to \varphi_{49})$
\item $\varphi_{276} :\equiv \forall y\,((f / g) \circ a(n) = y \to y = (f / g) \circ a(n)) \to \varphi_{211}$
\item $\varphi_{277} :\equiv \mathrm{ContIn}(f,A,c) \to \mathrm{ContIn}(g,A,c) \to \forall x \in A\,(g(x) \neq 0) \to \varphi_{222}$
\item $\varphi_{278} :\equiv \varphi_{6} \land \mathrm{Seq}(a) \land \forall n \in \mathbb{N}\,(a_{n} \in A) \land \mathrm{Conv}(a,c) \to \varphi_{131}$
\item $\varphi_{279} :\equiv \mathrm{ContIn}(f,A,c) \to \mathrm{ContIn}(g,A,c) \to \forall x \in A\,(g(x) \neq 0) \to \varphi_{232}$
\item $\varphi_{280} :\equiv (g \circ a(n) = g(a_{n}) \to g(a_{n}) = g \circ a(n)) \to \mathrm{ContIn}(f,A,c) \to \varphi_{197}$
\item $\varphi_{281} :\equiv (g(a_{n}) = g \circ a(n) \to g \circ a(n) = g(a_{n})) \to \mathrm{ContIn}(f,A,c) \to \varphi_{196}$
\item $\varphi_{282} :\equiv \mathrm{ContIn}(f,A,c) \to \mathrm{ContIn}(g,A,c) \to \forall x \in A\,(g(x) \neq 0) \to \varphi_{241}$
\item $\varphi_{283} :\equiv (\mathrm{ContIn}(f,A,c) \to \mathrm{ContIn}(g,A,c) \to \varphi_{162}) \to \mathrm{ContIn}(f,A,c) \to \varphi_{28}$
\item $\varphi_{284} :\equiv \forall a\,(\mathrm{RFn}(g) \to \mathrm{Seq}(a) \to \varphi_{126}) \to \mathrm{RFn}(g) \to \mathrm{Seq}(a) \to \varphi_{126}$
\item $\varphi_{285} :\equiv \forall a\,(\mathrm{RFn}(f) \to \mathrm{Seq}(a) \to \varphi_{127}) \to \mathrm{RFn}(f) \to \mathrm{Seq}(a) \to \varphi_{127}$
\item $\varphi_{286} :\equiv \mathrm{ContIn}(g,A,c) \to \forall x \in A\,(g(x) \neq 0) \to \mathrm{Seq}(a) \to \varphi_{266}$
\item $\varphi_{287} :\equiv \forall f\,\forall a\,(\mathrm{RFn}(f) \to \mathrm{Seq}(a) \to \varphi_{127}) \to \forall a\,(\mathrm{RFn}(g) \to \mathrm{Seq}(a) \to \varphi_{126})$
\item $\varphi_{288} :\equiv \forall f\,\forall a\,(\mathrm{RFn}(f) \to \mathrm{Seq}(a) \to \varphi_{127}) \to \forall a\,(\mathrm{RFn}(f) \to \mathrm{Seq}(a) \to \varphi_{127})$
\item $\varphi_{289} :\equiv (\mathrm{ContIn}(f,A,c) \to \varphi_{125}) \to \mathrm{ContIn}(f,A,c) \to \varphi_{150}$
\item $\varphi_{290} :\equiv (\mathrm{ContIn}(f,A,c) \to \varphi_{125}) \to \mathrm{ContIn}(f,A,c) \to \varphi_{151}$
\item $\varphi_{291} :\equiv \varphi_{119} \to \mathrm{ContIn}(f,A,c) \to \mathrm{ContIn}(g,A,c) \to \varphi_{187}$
\item $\varphi_{292} :\equiv (\mathrm{ContIn}(f,A,c) \to \varphi_{140}) \to \mathrm{ContIn}(f,A,c) \to \varphi_{146}$
\item $\varphi_{293} :\equiv \varphi_{120} \to \mathrm{ContIn}(f,A,c) \to \mathrm{ContIn}(g,A,c) \to \varphi_{188}$
\item $\varphi_{294} :\equiv (\mathrm{ContIn}(f,A,c) \to \mathrm{ContIn}(f,A,c)) \to \varphi_{235} \to \mathrm{ContIn}(f,A,c) \to \varphi_{3}$
\item $\varphi_{295} :\equiv (\mathrm{ContIn}(f,A,c) \to \mathrm{ContIn}(f,A,c)) \to \varphi_{235} \to \varphi_{36}$
\item $\varphi_{296} :\equiv (A \subseteq \mathrm{dom}(g) \to u \in A \to u \in \mathrm{dom}(g)) \to (\mathrm{ContIn}(f,A,c) \to \varphi_{7}) \to \varphi_{220}$
\item $\varphi_{297} :\equiv (A \subseteq \mathrm{dom}(f) \to u \in A \to u \in \mathrm{dom}(f)) \to (\mathrm{ContIn}(f,A,c) \to \varphi_{8}) \to \varphi_{221}$
\item $\varphi_{298} :\equiv (\mathrm{ContIn}(f,A,c) \to \varphi_{140}) \to \mathrm{ContIn}(f,A,c) \to \varphi_{161}$
\item $\varphi_{299} :\equiv (\varphi_{6} \to \forall a\,(\varphi_{6})) \to \mathrm{ContIn}(f,A,c) \to \mathrm{ContIn}(g,A,c) \to \varphi_{162}$
\item $\varphi_{300} :\equiv (\mathrm{ContIn}(f,A,c) \to \mathrm{ContIn}(f,A,c)) \to \varphi_{235} \to \mathrm{ContIn}(f,A,c) \to \varphi_{8}$
\item $\varphi_{301} :\equiv (\mathrm{ContIn}(f,A,c) \to \varphi_{2}) \to \mathrm{ContIn}(f,A,c) \to \varphi_{259}$
\item $\varphi_{302} :\equiv f \circ a \otimes g \circ a^{\ast}(n) = f \circ a(n) \cdot g \circ a(n)^{-1} \to \varphi_{247}$
\item $\varphi_{303} :\equiv (\mathrm{ContIn}(f,A,c) \to \varphi_{146}) \to \mathrm{ContIn}(f,A,c) \to \varphi_{161}$
\item $\varphi_{304} :\equiv (\mathrm{ContIn}(f,A,c) \to \varphi_{146}) \to \mathrm{ContIn}(f,A,c) \to \varphi_{164}$
\item $\varphi_{305} :\equiv (\mathrm{ContIn}(f,A,c) \to \varphi_{146}) \to \mathrm{ContIn}(f,A,c) \to \varphi_{165}$
\item $\varphi_{306} :\equiv (\mathrm{ContIn}(f,A,c) \to \varphi_{139}) \to \mathrm{ContIn}(f,A,c) \to \varphi_{185}$
\item $\varphi_{307} :\equiv (\mathrm{ContIn}(f,A,c) \to \varphi_{259}) \to \mathrm{ContIn}(f,A,c) \to \varphi_{10}$
\item $\varphi_{308} :\equiv (\mathrm{ContIn}(f,A,c) \to \varphi_{140}) \to \mathrm{ContIn}(f,A,c) \to \varphi_{196}$
\item $\varphi_{309} :\equiv (\mathrm{ContIn}(f,A,c) \to \varphi_{140}) \to \mathrm{ContIn}(f,A,c) \to \varphi_{198}$
\item $\varphi_{310} :\equiv f \circ a \otimes g \circ a^{\ast}(n) = f \circ a(n) \cdot g \circ a^{\ast}(n) \to \varphi_{256}$
\item $\varphi_{311} :\equiv (\mathrm{ContIn}(f,A,c) \to \varphi_{26}) \to \varphi_{149} \to \mathrm{ContIn}(f,A,c) \to \varphi_{53}$
\item $\varphi_{312} :\equiv (f(c) \cdot g(c)^{-1} = (f / g)(c) \to \varphi_{113}) \to \mathrm{ContIn}(f,A,c) \to \mathrm{ContIn}(g,A,c) \to \varphi_{152}$
\item $\varphi_{313} :\equiv (\mathrm{ContIn}(f,A,c) \to \varphi_{196}) \to \mathrm{ContIn}(f,A,c) \to \varphi_{171}$
\item $\varphi_{314} :\equiv (A \subseteq \mathrm{dom}(g) \to c \in A \to c \in \mathrm{dom}(g)) \to (\mathrm{ContIn}(f,A,c) \to \varphi_{7}) \to \varphi_{36} \to \varphi_{155}$
\item $\varphi_{315} :\equiv (A \subseteq \mathrm{dom}(f) \to c \in A \to c \in \mathrm{dom}(f)) \to (\mathrm{ContIn}(f,A,c) \to \varphi_{8}) \to \varphi_{36} \to \varphi_{156}$
\item $\varphi_{316} :\equiv (\mathrm{ContIn}(f,A,c) \to \varphi_{140}) \to \mathrm{ContIn}(f,A,c) \to \mathrm{ContIn}(g,A,c) \to \varphi_{199}$
\item $\varphi_{317} :\equiv (\mathrm{ContIn}(f,A,c) \to \varphi_{140}) \to \mathrm{ContIn}(f,A,c) \to \mathrm{ContIn}(g,A,c) \to \varphi_{200}$
\item $\varphi_{318} :\equiv (\mathrm{ContIn}(f,A,c) \to \mathrm{ContIn}(g,A,c) \to \varphi_{152}) \to \varphi_{6} \to \mathrm{Seq}(a) \to \forall n \in \mathbb{N}\,(a_{n} \in A) \to \varphi_{5}$
\item $\varphi_{319} :\equiv (\mathrm{ContIn}(f,A,c) \to \mathrm{ContIn}(g,A,c) \to \varphi_{111}) \to \varphi_{250}$
\item $\varphi_{320} :\equiv (\mathrm{Conv}(g \circ a,g(c)) \to \mathrm{Seq}(g \circ a)) \to (\mathrm{ContIn}(f,A,c) \to \varphi_{150}) \to \varphi_{168}$
\item $\varphi_{321} :\equiv (\mathrm{Conv}(f \circ a,f(c)) \to \mathrm{Seq}(f \circ a)) \to (\mathrm{ContIn}(f,A,c) \to \varphi_{151}) \to \varphi_{169}$
\item $\varphi_{322} :\equiv \forall a\,(\varphi_{6} \to \mathrm{Seq}(a) \to \forall n \in \mathbb{N}\,(a_{n} \in A) \to \varphi_{5}) \to \forall a\,(\varphi_{6}) \to \forall a\,(\mathrm{Seq}(a) \to \forall n \in \mathbb{N}\,(a_{n} \in A) \to \varphi_{5})$
\item $\varphi_{323} :\equiv \varphi_{175} \to \mathrm{ContIn}(f,A,c) \to \mathrm{ContIn}(g,A,c) \to \varphi_{212}$
\item $\varphi_{324} :\equiv \varphi_{176} \to \mathrm{ContIn}(f,A,c) \to \mathrm{ContIn}(g,A,c) \to \varphi_{214}$
\item $\varphi_{325} :\equiv (\mathrm{ContIn}(f,A,c) \to \varphi_{140}) \to \mathrm{ContIn}(f,A,c) \to \mathrm{ContIn}(g,A,c) \to \varphi_{213}$
\item $\varphi_{326} :\equiv (\mathrm{ContIn}(f,A,c) \to \varphi_{170}) \to \mathrm{ContIn}(f,A,c) \to \mathrm{ContIn}(g,A,c) \to \varphi_{203}$
\item $\varphi_{327} :\equiv (\mathrm{ContIn}(f,A,c) \to \varphi_{172}) \to \mathrm{ContIn}(f,A,c) \to \mathrm{ContIn}(g,A,c) \to \varphi_{204}$
\item $\varphi_{328} :\equiv (\mathrm{ContIn}(f,A,c) \to \varphi_{178}) \to \mathrm{ContIn}(f,A,c) \to \mathrm{ContIn}(g,A,c) \to \varphi_{210}$
\item $\varphi_{329} :\equiv (\mathrm{ContIn}(f,A,c) \to \mathrm{ContIn}(g,A,c) \to \varphi_{231}) \to \mathrm{ContIn}(f,A,c) \to \varphi_{144}$
\item $\varphi_{330} :\equiv (\mathrm{ContIn}(f,A,c) \to \mathrm{ContIn}(g,A,c) \to \varphi_{233}) \to \mathrm{ContIn}(f,A,c) \to \varphi_{145}$
\item $\varphi_{331} :\equiv (\mathrm{ContIn}(f,A,c) \to \mathrm{ContIn}(g,A,c) \to \varphi_{210}) \to \mathrm{ContIn}(f,A,c) \to \mathrm{ContIn}(g,A,c) \to \varphi_{177}$
\item $\varphi_{332} :\equiv \varphi_{240} \to \mathrm{RFn}((f / g)) \to A \subseteq \mathrm{dom}((f / g)) \to c \in A \to \varphi_{153}$
\item $\varphi_{333} :\equiv \varphi_{211} \to \mathrm{ContIn}(f,A,c) \to \mathrm{ContIn}(g,A,c) \to \varphi_{226}$
\item $\varphi_{334} :\equiv (\mathrm{ContIn}(f,A,c) \to \varphi_{178}) \to \mathrm{ContIn}(f,A,c) \to \mathrm{ContIn}(g,A,c) \to \varphi_{231}$
\item $\varphi_{335} :\equiv (\mathrm{ContIn}(f,A,c) \to \varphi_{140}) \to \mathrm{ContIn}(f,A,c) \to \mathrm{ContIn}(g,A,c) \to \varphi_{253}$
\item $\varphi_{336} :\equiv (\mathrm{ContIn}(f,A,c) \to \mathrm{ContIn}(g,A,c) \to \varphi_{203}) \to \mathrm{ContIn}(f,A,c) \to \mathrm{ContIn}(g,A,c) \to \varphi_{191}$
\item $\varphi_{337} :\equiv (\mathrm{ContIn}(f,A,c) \to \mathrm{ContIn}(g,A,c) \to \varphi_{204}) \to \mathrm{ContIn}(f,A,c) \to \mathrm{ContIn}(g,A,c) \to \varphi_{192}$
\item $\varphi_{338} :\equiv (\mathrm{ContIn}(f,A,c) \to \varphi_{184}) \to \mathrm{ContIn}(f,A,c) \to \mathrm{ContIn}(g,A,c) \to \varphi_{233}$
\item $\varphi_{339} :\equiv ((f / g) \circ a(n) = (f / g)(a_{n}) \to \varphi_{179}) \to \mathrm{ContIn}(f,A,c) \to \mathrm{ContIn}(g,A,c) \to \varphi_{219}$
\item $\varphi_{340} :\equiv (g \circ a(n) = g(a_{n}) \to g \circ a(n) = 0 \to g(a_{n}) = 0) \to \varphi_{313}$
\item $\varphi_{341} :\equiv \varphi_{239} \to \mathrm{ContIn}(f,A,c) \to \mathrm{ContIn}(g,A,c) \to \varphi_{243}$
\item $\varphi_{342} :\equiv (\mathrm{ContIn}(f,A,c) \to \mathrm{ContIn}(g,A,c) \to \varphi_{111}) \to \varphi_{289}$
\item $\varphi_{343} :\equiv (\mathrm{ContIn}(f,A,c) \to \mathrm{ContIn}(g,A,c) \to \varphi_{111}) \to \varphi_{290}$
\item $\varphi_{344} :\equiv (\mathrm{ContIn}(f,A,c) \to \varphi_{184}) \to \varphi_{83} \to \mathrm{ContIn}(f,A,c) \to \varphi_{181}$
\item $\varphi_{345} :\equiv (\mathrm{ContIn}(f,A,c) \to \mathrm{ContIn}(g,A,c) \to \varphi_{231}) \to \mathrm{ContIn}(f,A,c) \to \mathrm{ContIn}(g,A,c) \to \varphi_{216}$
\item $\varphi_{346} :\equiv (\mathrm{ContIn}(f,A,c) \to \mathrm{ContIn}(g,A,c) \to \varphi_{233}) \to \mathrm{ContIn}(f,A,c) \to \mathrm{ContIn}(g,A,c) \to \varphi_{217}$
\item $\varphi_{347} :\equiv (\mathrm{ContIn}(f,A,c) \to \mathrm{ContIn}(g,A,c) \to \varphi_{226}) \to \varphi_{261}$
\item $\varphi_{348} :\equiv (\mathrm{ContIn}(f,A,c) \to \varphi_{185}) \to \varphi_{88} \to \mathrm{ContIn}(f,A,c) \to \varphi_{186}$
\item $\varphi_{349} :\equiv (\mathrm{ContIn}(f,A,c) \to \mathrm{ContIn}(g,A,c) \to \varphi_{101}) \to \varphi_{292}$
\item $\varphi_{350} :\equiv (\mathrm{ContIn}(f,A,c) \to \varphi_{144}) \to \mathrm{ContIn}(f,A,c) \to \varphi_{286}$
\item $\varphi_{351} :\equiv (A \subseteq \mathrm{dom}(g) \to a_{n} \in A \to a_{n} \in \mathrm{dom}(g)) \to (\mathrm{ContIn}(f,A,c) \to \varphi_{7}) \to \varphi_{304}$
\item $\varphi_{352} :\equiv (A \subseteq \mathrm{dom}(f) \to a_{n} \in A \to a_{n} \in \mathrm{dom}(f)) \to (\mathrm{ContIn}(f,A,c) \to \varphi_{8}) \to \varphi_{305}$
\item $\varphi_{353} :\equiv (g \circ a(n) = g(a_{n}) \to \varphi_{230}) \to \mathrm{ContIn}(f,A,c) \to \mathrm{ContIn}(g,A,c) \to \varphi_{237}$
\item $\varphi_{354} :\equiv \varphi_{278} \to \mathrm{ContIn}(f,A,c) \to \mathrm{ContIn}(g,A,c) \to \varphi_{227}$
\item $\varphi_{355} :\equiv ((f / g) \circ a(n) = f(a_{n}) \cdot g(a_{n})^{-1} \to \varphi_{224}) \to \mathrm{ContIn}(f,A,c) \to \mathrm{ContIn}(g,A,c) \to \varphi_{225}$
\item $\varphi_{356} :\equiv (f \circ a(n) = f(a_{n}) \to \varphi_{252}) \to \mathrm{ContIn}(f,A,c) \to \mathrm{ContIn}(g,A,c) \to \varphi_{245}$
\item $\varphi_{357} :\equiv (\mathrm{ContIn}(f,A,c) \to \varphi_{161}) \to \varphi_{205} \to \mathrm{ContIn}(f,A,c) \to \varphi_{181}$
\item $\varphi_{358} :\equiv (\mathrm{ContIn}(f,A,c) \to \mathrm{ContIn}(g,A,c) \to \varphi_{191}) \to \varphi_{291}$
\item $\varphi_{359} :\equiv (\mathrm{ContIn}(f,A,c) \to \mathrm{ContIn}(g,A,c) \to \varphi_{192}) \to \varphi_{293}$
\item $\varphi_{360} :\equiv \varphi_{255} \to f \circ a \otimes g \circ a^{\ast}(n) = f(a_{n}) \cdot g \circ a(n)^{-1} \to \varphi_{229}$
\item $\varphi_{361} :\equiv (\mathrm{ContIn}(f,A,c) \to \varphi_{34}) \to (\mathrm{ContIn}(f,A,c) \to \varphi_{33}) \to \varphi_{311}$
\item $\varphi_{362} :\equiv (\mathrm{ContIn}(f,A,c) \to \mathrm{ContIn}(g,A,c) \to \varphi_{177}) \to \varphi_{312}$
\item $\varphi_{363} :\equiv (\mathrm{ContIn}(f,A,c) \to \mathrm{ContIn}(g,A,c) \to \varphi_{187}) \to \varphi_{308}$
\item $\varphi_{364} :\equiv (\mathrm{ContIn}(f,A,c) \to \mathrm{ContIn}(g,A,c) \to \varphi_{188}) \to \varphi_{309}$
\item $\varphi_{365} :\equiv (\mathrm{ContIn}(f,A,c) \to \varphi_{10}) \to (\mathrm{ContIn}(f,A,c) \to \varphi_{27}) \to \varphi_{36} \to \varphi_{283}$
\item $\varphi_{366} :\equiv (g(a_{n}) = g \circ a(n) \to g(a_{n}) = 0 \to g \circ a(n) = 0) \to \varphi_{340}$
\item $\varphi_{367} :\equiv (\mathrm{ContIn}(f,A,c) \to \varphi_{2}) \to (\mathrm{ContIn}(f,A,c) \to \varphi_{12}) \to \varphi_{326}$
\item $\varphi_{368} :\equiv (\mathrm{ContIn}(f,A,c) \to \varphi_{3}) \to (\mathrm{ContIn}(f,A,c) \to \varphi_{12}) \to \varphi_{327}$
\item $\varphi_{369} :\equiv (f \circ a \otimes g \circ a^{\ast}(n) = f(a_{n}) \cdot g \circ a(n)^{-1} \to \varphi_{229}) \to \mathrm{ContIn}(f,A,c) \to \mathrm{ContIn}(g,A,c) \to \varphi_{242}$
\item $\varphi_{370} :\equiv (\mathrm{ContIn}(f,A,c) \to \mathrm{ContIn}(g,A,c) \to \varphi_{212}) \to \varphi_{316}$
\item $\varphi_{371} :\equiv \varphi_{302} \to \mathrm{ContIn}(f,A,c) \to \mathrm{ContIn}(g,A,c) \to \varphi_{246}$
\item $\varphi_{372} :\equiv (\mathrm{ContIn}(f,A,c) \to \mathrm{ContIn}(g,A,c) \to \varphi_{214}) \to \varphi_{317}$
\item $\varphi_{373} :\equiv (\mathrm{ContIn}(f,A,c) \to \mathrm{ContIn}(f,A,c)) \to (\mathrm{ContIn}(f,A,c) \to \varphi_{12}) \to \varphi_{343}$
\item $\varphi_{374} :\equiv f(a_{n}) \cdot g \circ a(n)^{-1} = f(a_{n}) \cdot g \circ a(n)^{-1} \to \varphi_{353}$
\item $\varphi_{375} :\equiv (\mathrm{ContIn}(f,A,c) \to \mathrm{ContIn}(g,A,c) \to \varphi_{216}) \to \varphi_{323}$
\item $\varphi_{376} :\equiv (\mathrm{ContIn}(f,A,c) \to \mathrm{ContIn}(g,A,c) \to \varphi_{217}) \to \varphi_{324}$
\item $\varphi_{377} :\equiv (\mathrm{ContIn}(f,A,c) \to \varphi_{70}) \to (\mathrm{ContIn}(f,A,c) \to \varphi_{146}) \to \varphi_{298}$
\item $\varphi_{378} :\equiv (\mathrm{ContIn}(f,A,c) \to \varphi_{10}) \to (\mathrm{ContIn}(f,A,c) \to \varphi_{12}) \to \varphi_{338}$
\item $\varphi_{379} :\equiv \varphi_{310} \to \mathrm{ContIn}(f,A,c) \to \mathrm{ContIn}(g,A,c) \to \varphi_{251}$
\item $\varphi_{380} :\equiv (\mathrm{RFn}(f) \to \mathrm{RFn}(g) \to \mathrm{RFn}((f / g)) \land \varphi_{39} \land \varphi_{41}) \to (\mathrm{ContIn}(f,A,c) \to \varphi_{3}) \to \varphi_{301}$
\item $\varphi_{381} :\equiv (\mathrm{ContIn}(f,A,c) \to \mathrm{ContIn}(g,A,c) \to \varphi_{227}) \to \varphi_{331}$
\item $\varphi_{382} :\equiv (\mathrm{ContIn}(f,A,c) \to \mathrm{ContIn}(g,A,c) \to \varphi_{225}) \to \varphi_{333}$
\item $\varphi_{383} :\equiv (\mathrm{ContIn}(f,A,c) \to \mathrm{ContIn}(g,A,c) \to \varphi_{213}) \to \varphi_{339}$
\item $\varphi_{384} :\equiv (\mathrm{ContIn}(f,A,c) \to \mathrm{ContIn}(g,A,c) \to \varphi_{242}) \to \varphi_{341}$
\item $\varphi_{385} :\equiv (\mathrm{ContIn}(g,A,c) \to \mathrm{Seq}(a) \to \varphi_{48}) \to \varphi_{1} \to (\mathrm{ContIn}(f,A,c) \to \varphi_{12}) \to \varphi_{342}$
\item $\varphi_{386} :\equiv (\mathrm{Seq}(g \circ a) \to \forall n \in \mathbb{N}\,(g \circ a(n) \neq 0) \to \varphi_{99}) \to \varphi_{168} \to \varphi_{334}$
\item $\varphi_{387} :\equiv (\mathrm{RFn}((f / g)) \to A \subseteq \mathrm{dom}((f / g)) \to c \in A \to \varphi_{153}) \to \varphi_{365}$
\item $\varphi_{388} :\equiv (\mathrm{ContIn}(f,A,c) \to \mathrm{ContIn}(g,A,c) \to \varphi_{243}) \to \varphi_{355}$
\item $\varphi_{389} :\equiv (\forall n\,(\varphi_{193}) \to \varphi_{131} \to \forall n \in \mathbb{N}\,(a_{n} \in \mathrm{dom}(g))) \to \varphi_{278} \to \mathrm{ContIn}(f,A,c) \to \varphi_{170}$
\item $\varphi_{390} :\equiv (\forall n\,(\varphi_{194}) \to \varphi_{131} \to \forall n \in \mathbb{N}\,(a_{n} \in \mathrm{dom}(f))) \to \varphi_{278} \to \mathrm{ContIn}(f,A,c) \to \varphi_{172}$
\item $\varphi_{391} :\equiv (\forall n\,(\varphi_{201}) \to \varphi_{131} \to \forall n \in \mathbb{N}\,(g \circ a(n) \neq 0)) \to \varphi_{278} \to \mathrm{ContIn}(f,A,c) \to \varphi_{178}$
\item $\varphi_{392} :\equiv (\mathrm{ContIn}(f,A,c) \to \mathrm{ContIn}(g,A,c) \to \varphi_{237}) \to \varphi_{369}$
\item $\varphi_{393} :\equiv (\forall n\,(\varphi_{206}) \to \varphi_{131} \to \forall n \in \mathbb{N}\,(a_{n} \in \mathrm{dom}((f / g)))) \to \varphi_{278} \to \mathrm{ContIn}(f,A,c) \to \varphi_{184}$
\item $\varphi_{394} :\equiv (\mathrm{ContIn}(f,A,c) \to \mathrm{ContIn}(g,A,c) \to \varphi_{245}) \to \varphi_{371}$
\item $\varphi_{395} :\equiv (\mathrm{ContIn}(f,A,c) \to \varphi_{165}) \to (\mathrm{ContIn}(f,A,c) \to \varphi_{164}) \to \varphi_{357}$
\item $\varphi_{396} :\equiv (\mathrm{ContIn}(f,A,c) \to \varphi_{151}) \to (\mathrm{ContIn}(f,A,c) \to \varphi_{150}) \to \varphi_{147} \to \varphi_{328}$
\item $\varphi_{397} :\equiv (\mathrm{ContIn}(f,A,c) \to \mathrm{ContIn}(g,A,c) \to \varphi_{200}) \to \varphi_{383}$
\item $\varphi_{398} :\equiv (\mathrm{ContIn}(f,A,c) \to \mathrm{ContIn}(g,A,c) \to \varphi_{199}) \to \varphi_{114} \to (g \circ a^{\ast}(n) = g \circ a(n)^{-1} \to \varphi_{114} \to \varphi_{112}) \to \varphi_{277}$
\item $\varphi_{399} :\equiv (\mathrm{ContIn}(f,A,c) \to \varphi_{197}) \to (\mathrm{ContIn}(f,A,c) \to \varphi_{161}) \to \varphi_{366}$
\item $\varphi_{400} :\equiv (\mathrm{ContIn}(f,A,c) \to \mathrm{ContIn}(g,A,c) \to \varphi_{219}) \to \varphi_{388}$
\item $\varphi_{401} :\equiv (\mathrm{ContIn}(f,A,c) \to \mathrm{ContIn}(g,A,c) \to \varphi_{246}) \to \varphi_{392}$
\item $\varphi_{402} :\equiv (\mathrm{ContIn}(f,A,c) \to \mathrm{ContIn}(g,A,c) \to \varphi_{251}) \to \varphi_{394}$
\item $\varphi_{403} :\equiv (\mathrm{Seq}((f / g) \circ a) \to \varphi_{92} \to \varphi_{166}) \to (\mathrm{ContIn}(f,A,c) \to \varphi_{145}) \to \varphi_{381}$
\item $\varphi_{404} :\equiv (\mathrm{ContIn}(f,A,c) \to \mathrm{ContIn}(g,A,c) \to \varphi_{253}) \to \varphi_{277} \to \varphi_{379}$
\end{abbrevlist}
\end{fullproof}

% 自動生成 (prover/build.js): ContInComp  ノード数 359
\begin{fullproof}{定理2.6}{$\forall f\,\forall g\,\forall A\,\forall B\,\forall c\,(\mathrm{ContIn}(f,A,c) \to \mathrm{ContIn}(g,B,f(c)) \to \forall x \in A\,(f(x) \in B) \to \mathrm{ContIn}(g \circ f,A,c))$}{359}\phantomsection\label{tree:ContInComp}
\noindent$\Pi_{1}$:
\begin{ptree}
\AXC{$z = w \to \lvert g \circ f(x) - z \rvert < \varepsilon \to \lvert g \circ f(x) - w \rvert < \varepsilon$\lbl{A6}}
\UIC{$\varphi_{28}$}
\UIC{$\forall z\,(\varphi_{28})$}
\AXC{$\forall z\,(\varphi_{28}) \to \varphi_{52}$\lbl{A2}}
\BIC{$\varphi_{52}$}
\AXC{$\varphi_{52} \to g(f(c)) = g \circ f(c) \to \varphi_{49}$\lbl{A2}}
\BIC{$g(f(c)) = g \circ f(c) \to \varphi_{49}$}
\end{ptree}
\noindent$\Pi_{2}$:
\begin{ptree}
\AXC{$\varphi_{57}$\lbl{A2}}
\AXC{$\varphi_{164}$\lbl{A1}}
\AXC{$\varphi_{164} \to \varphi_{170}$\lbl{A1}}
\BIC{$\varphi_{170}$}
\AXC{$\varphi_{170} \to \varphi_{57} \to \varphi_{167}$\lbl{A1}}
\BIC{$\varphi_{57} \to \varphi_{167}$}
\BIC{$\varphi_{167}$}
\AXC{$\varphi_{167} \to (\mathrm{ContIn}(f,A,c) \to \varphi_{122}) \to \varphi_{192}$\lbl{A1}}
\BIC{$(\mathrm{ContIn}(f,A,c) \to \varphi_{122}) \to \varphi_{192}$}
\end{ptree}
\noindent$\Pi_{3}$:
\begin{ptree}
\AXC{$\mathrm{ContIn}(f,A,c) \to \varphi_{133}$\lbl{A1}}
\AXC{$\mathrm{ContIn}(f,A,c) \to \varphi_{122}$\lbl{A1}}
\AXC{$\stackrel{\Pi_{2}}{\vdots}$}\noLine
\UIC{$(\mathrm{ContIn}(f,A,c) \to \varphi_{122}) \to \varphi_{192}$}
\BIC{$\varphi_{192}$}
\BIC{$\mathrm{ContIn}(f,A,c) \to \varphi_{146}$}
\end{ptree}
\noindent$\Pi_{4}$:
\begin{ptree}
\AXC{$\forall x \in A\,(f(x) \in B) \to x \in A \to f(x) \in B$\lbl{A2}}
\AXC{$\mathrm{ContIn}(f,A,c) \to \varphi_{32}$\lbl{A1}}
\AXC{$(\mathrm{ContIn}(f,A,c) \to \varphi_{32}) \to \varphi_{63}$\lbl{A1}}
\BIC{$\varphi_{63}$}
\BIC{$\mathrm{ContIn}(f,A,c) \to \varphi_{25}$}
\end{ptree}
\noindent$\Pi_{5}$:
\begin{ptree}
\AXC{$\varphi_{31} \to f(x) \in B \to \varphi_{33}$\lbl{A2}}
\AXC{$\mathrm{ContIn}(f,A,c) \to \varphi_{118}$\lbl{A1}}
\AXC{$\varphi_{195}$\lbl{A1}}
\BIC{$\mathrm{ContIn}(f,A,c) \to \varphi_{168}$}
\AXC{$\varphi_{202}$\lbl{A1}}
\BIC{$(\varphi_{31} \to f(x) \in B \to \varphi_{33}) \to \mathrm{ContIn}(f,A,c) \to \varphi_{172}$}
\BIC{$\mathrm{ContIn}(f,A,c) \to \varphi_{172}$}
\AXC{$\varphi_{213}$\lbl{A1}}
\BIC{$(\mathrm{ContIn}(f,A,c) \to \varphi_{25}) \to \varphi_{204}$}
\end{ptree}
\noindent$\Pi_{6}$:
\begin{ptree}
\AXC{$\mathrm{ContIn}(f,A,c) \to \varphi_{122}$\lbl{A1}}
\AXC{$\stackrel{\Pi_{4}}{\vdots}$}\noLine
\UIC{$\mathrm{ContIn}(f,A,c) \to \varphi_{25}$}
\AXC{$\stackrel{\Pi_{5}}{\vdots}$}\noLine
\UIC{$(\mathrm{ContIn}(f,A,c) \to \varphi_{25}) \to \varphi_{204}$}
\BIC{$\varphi_{204}$}
\BIC{$(\mathrm{ContIn}(f,A,c) \to \varphi_{146}) \to \varphi_{158}$}
\end{ptree}
\noindent$\Pi_{7}$:
\begin{ptree}
\AXC{$x = x$\lbl{A5}}
\AXC{$x = y \to x = x \to y = x$\lbl{A6}}
\AXC{$(x = y \to x = x \to y = x) \to x = x \to x = y \to y = x$\lbl{A1}}
\BIC{$x = x \to x = y \to y = x$}
\BIC{$x = y \to y = x$}
\UIC{$\forall y\,(x = y \to y = x)$}
\UIC{$\forall x\,\forall y\,(x = y \to y = x)$}
\end{ptree}
\noindent$\Pi_{8}$:
\begin{ptree}
\AXC{$\forall x \in A\,(f(x) \in B) \to x \in A \to f(x) \in B$\lbl{A2}}
\AXC{$\mathrm{ContIn}(f,A,c) \to \varphi_{32}$\lbl{A1}}
\AXC{$(\mathrm{ContIn}(f,A,c) \to \varphi_{32}) \to \varphi_{63}$\lbl{A1}}
\BIC{$\varphi_{63}$}
\BIC{$\mathrm{ContIn}(f,A,c) \to \varphi_{25}$}
\end{ptree}
\noindent$\Pi_{9}$:
\begin{ptree}
\AXC{$\mathrm{ContIn}(f,A,c) \to \varphi_{122}$\lbl{A1}}
\AXC{$\stackrel{\Pi_{8}}{\vdots}$}\noLine
\UIC{$\mathrm{ContIn}(f,A,c) \to \varphi_{25}$}
\AXC{$(\mathrm{ContIn}(f,A,c) \to \varphi_{25}) \to \varphi_{187}$\lbl{A1}}
\BIC{$\varphi_{187}$}
\BIC{$\mathrm{ContIn}(f,A,c) \to \varphi_{131}$}
\end{ptree}
\noindent$\Pi_{10}$:
\begin{ptree}
\AXC{$\forall f\,\forall A\,\forall c\,(\varphi_{76})$\lbl{DefContIn}}
\AXC{$\forall f\,\forall A\,\forall c\,(\varphi_{76}) \to \varphi_{80}$\lbl{A2}}
\BIC{$\varphi_{80}$}
\AXC{$\varphi_{80} \to \varphi_{78}$\lbl{A2}}
\BIC{$\varphi_{78}$}
\AXC{$\varphi_{78} \to \varphi_{88}$\lbl{A2}}
\BIC{$\varphi_{88}$}
\end{ptree}
\noindent$\Pi_{11}$:
\begin{ptree}
\AXC{$\stackrel{\Pi_{10}}{\vdots}$}\noLine
\UIC{$\varphi_{88}$}
\AXC{$\mathrm{ContIn}(f,A,c) \to \varphi_{1}$\lbl{A1}}
\AXC{$\varphi_{155}$\lbl{A1}}
\BIC{$\varphi_{88} \to \mathrm{ContIn}(f,A,c) \to \varphi_{11}$}
\BIC{$\mathrm{ContIn}(f,A,c) \to \varphi_{11}$}
\end{ptree}
\noindent$\Pi_{12}$:
\begin{ptree}
\AXC{$\forall X\,\forall Y\,\forall x\,(X \subseteq Y \to x \in X \to x \in Y)$\lbl{[B5]}}
\AXC{$\varphi_{18}$\lbl{A2}}
\BIC{$\forall Y\,\forall x\,(B \subseteq Y \to x \in B \to x \in Y)$}
\AXC{$\varphi_{29}$\lbl{A2}}
\BIC{$\forall x\,(B \subseteq \mathrm{dom}(g) \to x \in B \to x \in \mathrm{dom}(g))$}
\AXC{$\varphi_{45}$\lbl{A2}}
\BIC{$B \subseteq \mathrm{dom}(g) \to f(x) \in B \to f(x) \in \mathrm{dom}(g)$}
\AXC{$\varphi_{198}$\lbl{A1}}
\BIC{$(\mathrm{ContIn}(f,A,c) \to \varphi_{11}) \to \varphi_{188}$}
\end{ptree}
\noindent$\Pi_{13}$:
\begin{ptree}
\AXC{$\stackrel{\Pi_{9}}{\vdots}$}\noLine
\UIC{$\mathrm{ContIn}(f,A,c) \to \varphi_{131}$}
\AXC{$\stackrel{\Pi_{11}}{\vdots}$}\noLine
\UIC{$\mathrm{ContIn}(f,A,c) \to \varphi_{11}$}
\AXC{$\stackrel{\Pi_{12}}{\vdots}$}\noLine
\UIC{$(\mathrm{ContIn}(f,A,c) \to \varphi_{11}) \to \varphi_{188}$}
\BIC{$\varphi_{188}$}
\BIC{$\mathrm{ContIn}(f,A,c) \to \varphi_{135}$}
\end{ptree}
\noindent$\Pi_{14}$:
\begin{ptree}
\AXC{$\forall f\,\forall A\,\forall c\,(\varphi_{76})$\lbl{DefContIn}}
\AXC{$\forall f\,\forall A\,\forall c\,(\varphi_{76}) \to \forall A\,\forall c\,(\varphi_{76})$\lbl{A2}}
\BIC{$\forall A\,\forall c\,(\varphi_{76})$}
\AXC{$\forall A\,\forall c\,(\varphi_{76}) \to \forall c\,(\varphi_{76})$\lbl{A2}}
\BIC{$\forall c\,(\varphi_{76})$}
\AXC{$\forall c\,(\varphi_{76}) \to \varphi_{76}$\lbl{A2}}
\BIC{$\varphi_{76}$}
\end{ptree}
\noindent$\Pi_{15}$:
\begin{ptree}
\AXC{$\stackrel{\Pi_{14}}{\vdots}$}\noLine
\UIC{$\varphi_{76}$}
\AXC{$\mathrm{ContIn}(f,A,c) \to \mathrm{ContIn}(f,A,c)$\lbl{A1}}
\AXC{$\varphi_{119}$\lbl{A1}}
\BIC{$\varphi_{76} \to \mathrm{ContIn}(f,A,c) \to \varphi_{12}$}
\BIC{$\mathrm{ContIn}(f,A,c) \to \varphi_{12}$}
\end{ptree}
\noindent$\Pi_{16}$:
\begin{ptree}
\AXC{$\forall X\,\forall Y\,\forall x\,(X \subseteq Y \to x \in X \to x \in Y)$\lbl{[B5]}}
\AXC{$\varphi_{19}$\lbl{A2}}
\BIC{$\forall Y\,\forall x\,(A \subseteq Y \to x \in A \to x \in Y)$}
\AXC{$\varphi_{30}$\lbl{A2}}
\BIC{$\forall x\,(A \subseteq \mathrm{dom}(f) \to x \in A \to x \in \mathrm{dom}(f))$}
\AXC{$\varphi_{36}$\lbl{A2}}
\BIC{$A \subseteq \mathrm{dom}(f) \to x \in A \to x \in \mathrm{dom}(f)$}
\AXC{$\varphi_{196}$\lbl{A1}}
\BIC{$(\mathrm{ContIn}(f,A,c) \to \varphi_{12}) \to \varphi_{186}$}
\end{ptree}
\noindent$\Pi_{17}$:
\begin{ptree}
\AXC{$\mathrm{ContIn}(f,A,c) \to \varphi_{122}$\lbl{A1}}
\AXC{$\stackrel{\Pi_{15}}{\vdots}$}\noLine
\UIC{$\mathrm{ContIn}(f,A,c) \to \varphi_{12}$}
\AXC{$\stackrel{\Pi_{16}}{\vdots}$}\noLine
\UIC{$(\mathrm{ContIn}(f,A,c) \to \varphi_{12}) \to \varphi_{186}$}
\BIC{$\varphi_{186}$}
\BIC{$\mathrm{ContIn}(f,A,c) \to \varphi_{129}$}
\end{ptree}
\noindent$\Pi_{18}$:
\begin{ptree}
\AXC{$\stackrel{\Pi_{14}}{\vdots}$}\noLine
\UIC{$\varphi_{76}$}
\AXC{$\mathrm{ContIn}(f,A,c) \to \mathrm{ContIn}(f,A,c)$\lbl{A1}}
\AXC{$\varphi_{116}$\lbl{A1}}
\BIC{$\varphi_{76} \to \mathrm{ContIn}(f,A,c) \to \varphi_{5}$}
\BIC{$\mathrm{ContIn}(f,A,c) \to \varphi_{5}$}
\end{ptree}
\noindent$\Pi_{19}$:
\begin{ptree}
\AXC{$\stackrel{\Pi_{10}}{\vdots}$}\noLine
\UIC{$\varphi_{88}$}
\AXC{$\mathrm{ContIn}(f,A,c) \to \varphi_{1}$\lbl{A1}}
\AXC{$\varphi_{152}$\lbl{A1}}
\BIC{$\varphi_{88} \to \mathrm{ContIn}(f,A,c) \to \varphi_{6}$}
\BIC{$\mathrm{ContIn}(f,A,c) \to \varphi_{6}$}
\end{ptree}
\noindent$\Pi_{20}$:
\begin{ptree}
\AXC{$\forall g\,\forall f\,\forall x\,(\mathrm{RFn}(g) \to \mathrm{RFn}(f) \to \varphi_{2})$\lbl{[S23]}}
\AXC{$\varphi_{71}$\lbl{A2}}
\BIC{$\forall f\,\forall x\,(\mathrm{RFn}(g) \to \mathrm{RFn}(f) \to \varphi_{2})$}
\AXC{$\varphi_{69}$\lbl{A2}}
\BIC{$\forall x\,(\mathrm{RFn}(g) \to \mathrm{RFn}(f) \to \varphi_{2})$}
\AXC{$\varphi_{66}$\lbl{A2}}
\BIC{$\mathrm{RFn}(g) \to \mathrm{RFn}(f) \to \varphi_{2}$}
\AXC{$\varphi_{210}$\lbl{A1}}
\BIC{$(\mathrm{ContIn}(f,A,c) \to \varphi_{6}) \to \varphi_{206}$}
\end{ptree}
\noindent$\Pi_{21}$:
\begin{ptree}
\AXC{$\stackrel{\Pi_{18}}{\vdots}$}\noLine
\UIC{$\mathrm{ContIn}(f,A,c) \to \varphi_{5}$}
\AXC{$\stackrel{\Pi_{19}}{\vdots}$}\noLine
\UIC{$\mathrm{ContIn}(f,A,c) \to \varphi_{6}$}
\AXC{$\stackrel{\Pi_{20}}{\vdots}$}\noLine
\UIC{$(\mathrm{ContIn}(f,A,c) \to \varphi_{6}) \to \varphi_{206}$}
\BIC{$\varphi_{206}$}
\BIC{$(\mathrm{ContIn}(f,A,c) \to \varphi_{129}) \to \varphi_{190}$}
\end{ptree}
\noindent$\Pi_{22}$:
\begin{ptree}
\AXC{$\stackrel{\Pi_{13}}{\vdots}$}\noLine
\UIC{$\mathrm{ContIn}(f,A,c) \to \varphi_{135}$}
\AXC{$\stackrel{\Pi_{17}}{\vdots}$}\noLine
\UIC{$\mathrm{ContIn}(f,A,c) \to \varphi_{129}$}
\AXC{$\stackrel{\Pi_{21}}{\vdots}$}\noLine
\UIC{$(\mathrm{ContIn}(f,A,c) \to \varphi_{129}) \to \varphi_{190}$}
\BIC{$\varphi_{190}$}
\BIC{$\mathrm{ContIn}(f,A,c) \to \varphi_{137}$}
\end{ptree}
\noindent$\Pi_{23}$:
\begin{ptree}
\AXC{$\stackrel{\Pi_{19}}{\vdots}$}\noLine
\UIC{$\mathrm{ContIn}(f,A,c) \to \varphi_{6}$}
\AXC{$\forall g\,\forall f\,(\varphi_{58})$\lbl{[S9]}}
\AXC{$\forall g\,\forall f\,(\varphi_{58}) \to \forall f\,(\varphi_{58})$\lbl{A2}}
\BIC{$\forall f\,(\varphi_{58})$}
\AXC{$\forall f\,(\varphi_{58}) \to \varphi_{58}$\lbl{A2}}
\BIC{$\varphi_{58}$}
\AXC{$\varphi_{58} \to \varphi_{145}$\lbl{A1}}
\BIC{$\varphi_{145}$}
\BIC{$(\mathrm{ContIn}(f,A,c) \to \varphi_{5}) \to \varphi_{81}$}
\end{ptree}
\noindent$\Pi_{24}$:
\begin{ptree}
\AXC{$\stackrel{\Pi_{18}}{\vdots}$}\noLine
\UIC{$\mathrm{ContIn}(f,A,c) \to \varphi_{5}$}
\AXC{$\stackrel{\Pi_{23}}{\vdots}$}\noLine
\UIC{$(\mathrm{ContIn}(f,A,c) \to \varphi_{5}) \to \varphi_{81}$}
\BIC{$\varphi_{81}$}
\end{ptree}
\noindent$\Pi_{25}$:
\begin{ptree}
\AXC{$\stackrel{\Pi_{24}}{\vdots}$}\noLine
\UIC{$\varphi_{81}$}
\AXC{$\varphi_{81} \to \mathrm{ContIn}(f,A,c) \to \varphi_{62}$\lbl{A1}}
\BIC{$\mathrm{ContIn}(f,A,c) \to \varphi_{62}$}
\end{ptree}
\noindent$\Pi_{26}$:
\begin{ptree}
\AXC{$\forall x \in \mathrm{dom}(g \circ f)\,(\varphi_{14}) \to x \in \mathrm{dom}(g \circ f) \to \varphi_{14}$\lbl{A2}}
\AXC{$\stackrel{\Pi_{25}}{\vdots}$}\noLine
\UIC{$\mathrm{ContIn}(f,A,c) \to \varphi_{62}$}
\AXC{$(\mathrm{ContIn}(f,A,c) \to \varphi_{62}) \to \varphi_{156}$\lbl{A1}}
\BIC{$\varphi_{156}$}
\BIC{$\mathrm{ContIn}(f,A,c) \to \varphi_{60}$}
\end{ptree}
\noindent$\Pi_{27}$:
\begin{ptree}
\AXC{$\stackrel{\Pi_{22}}{\vdots}$}\noLine
\UIC{$\mathrm{ContIn}(f,A,c) \to \varphi_{137}$}
\AXC{$\stackrel{\Pi_{26}}{\vdots}$}\noLine
\UIC{$\mathrm{ContIn}(f,A,c) \to \varphi_{60}$}
\AXC{$(\mathrm{ContIn}(f,A,c) \to \varphi_{60}) \to \varphi_{193}$\lbl{A1}}
\BIC{$\varphi_{193}$}
\BIC{$\mathrm{ContIn}(f,A,c) \to \varphi_{147}$}
\end{ptree}
\noindent$\Pi_{28}$:
\begin{ptree}
\AXC{$\stackrel{\Pi_{27}}{\vdots}$}\noLine
\UIC{$\mathrm{ContIn}(f,A,c) \to \varphi_{147}$}
\AXC{$(\mathrm{ContIn}(f,A,c) \to \varphi_{147}) \to \varphi_{175}$\lbl{A1}}
\BIC{$\varphi_{175}$}
\end{ptree}
\noindent$\Pi_{29}$:
\begin{ptree}
\AXC{$\stackrel{\Pi_{7}}{\vdots}$}\noLine
\UIC{$\forall x\,\forall y\,(x = y \to y = x)$}
\AXC{$\varphi_{7}$\lbl{A2}}
\BIC{$\forall y\,(g \circ f(x) = y \to y = g \circ f(x))$}
\AXC{$\varphi_{54}$\lbl{A2}}
\BIC{$g \circ f(x) = g(f(x)) \to g(f(x)) = g \circ f(x)$}
\AXC{$\stackrel{\Pi_{28}}{\vdots}$}\noLine
\UIC{$\varphi_{175}$}
\BIC{$\mathrm{ContIn}(f,A,c) \to \varphi_{148}$}
\AXC{$\varphi_{208}$\lbl{A1}}
\BIC{$\varphi_{158} \to (g(f(x)) = g \circ f(x) \to \varphi_{48}) \to \varphi_{159}$}
\end{ptree}
\noindent$\Pi_{30}$:
\begin{ptree}
\AXC{$\stackrel{\Pi_{3}}{\vdots}$}\noLine
\UIC{$\mathrm{ContIn}(f,A,c) \to \varphi_{146}$}
\AXC{$\stackrel{\Pi_{6}}{\vdots}$}\noLine
\UIC{$(\mathrm{ContIn}(f,A,c) \to \varphi_{146}) \to \varphi_{158}$}
\BIC{$\varphi_{158}$}
\AXC{$\stackrel{\Pi_{29}}{\vdots}$}\noLine
\UIC{$\varphi_{158} \to (g(f(x)) = g \circ f(x) \to \varphi_{48}) \to \varphi_{159}$}
\BIC{$(g(f(x)) = g \circ f(x) \to \varphi_{48}) \to \varphi_{159}$}
\end{ptree}
\noindent$\Pi_{31}$:
\begin{ptree}
\AXC{$\varphi_{17}$\lbl{A6}}
\UIC{$\forall w\,(\varphi_{17})$}
\UIC{$\forall z\,\forall w\,(\varphi_{17})$}
\AXC{$\forall z\,\forall w\,(\varphi_{17}) \to \varphi_{50}$\lbl{A2}}
\BIC{$\varphi_{50}$}
\AXC{$\varphi_{50} \to g(f(x)) = g \circ f(x) \to \varphi_{48}$\lbl{A2}}
\BIC{$g(f(x)) = g \circ f(x) \to \varphi_{48}$}
\AXC{$\stackrel{\Pi_{30}}{\vdots}$}\noLine
\UIC{$(g(f(x)) = g \circ f(x) \to \varphi_{48}) \to \varphi_{159}$}
\BIC{$\varphi_{159}$}
\end{ptree}
\noindent$\Pi_{32}$:
\begin{ptree}
\AXC{$\stackrel{\Pi_{14}}{\vdots}$}\noLine
\UIC{$\varphi_{76}$}
\AXC{$\mathrm{ContIn}(f,A,c) \to \mathrm{ContIn}(f,A,c)$\lbl{A1}}
\AXC{$\varphi_{117}$\lbl{A1}}
\BIC{$\varphi_{76} \to \mathrm{ContIn}(f,A,c) \to \varphi_{9}$}
\BIC{$\mathrm{ContIn}(f,A,c) \to \varphi_{9}$}
\end{ptree}
\noindent$\Pi_{33}$:
\begin{ptree}
\AXC{$\forall x \in A\,(f(x) \in B) \to c \in A \to f(c) \in B$\lbl{A2}}
\AXC{$\mathrm{ContIn}(f,A,c) \to \varphi_{32}$\lbl{A1}}
\AXC{$(\mathrm{ContIn}(f,A,c) \to \varphi_{32}) \to \varphi_{64}$\lbl{A1}}
\BIC{$\varphi_{64}$}
\BIC{$\mathrm{ContIn}(f,A,c) \to \varphi_{26}$}
\end{ptree}
\noindent$\Pi_{34}$:
\begin{ptree}
\AXC{$\stackrel{\Pi_{32}}{\vdots}$}\noLine
\UIC{$\mathrm{ContIn}(f,A,c) \to \varphi_{9}$}
\AXC{$\stackrel{\Pi_{33}}{\vdots}$}\noLine
\UIC{$\mathrm{ContIn}(f,A,c) \to \varphi_{26}$}
\AXC{$(\mathrm{ContIn}(f,A,c) \to \varphi_{26}) \to \varphi_{174}$\lbl{A1}}
\BIC{$\varphi_{174}$}
\BIC{$\mathrm{ContIn}(f,A,c) \to \varphi_{132}$}
\end{ptree}
\noindent$\Pi_{35}$:
\begin{ptree}
\AXC{$\forall X\,\forall Y\,\forall x\,(X \subseteq Y \to x \in X \to x \in Y)$\lbl{[B5]}}
\AXC{$\varphi_{18}$\lbl{A2}}
\BIC{$\forall Y\,\forall x\,(B \subseteq Y \to x \in B \to x \in Y)$}
\AXC{$\varphi_{29}$\lbl{A2}}
\BIC{$\forall x\,(B \subseteq \mathrm{dom}(g) \to x \in B \to x \in \mathrm{dom}(g))$}
\AXC{$\varphi_{46}$\lbl{A2}}
\BIC{$B \subseteq \mathrm{dom}(g) \to f(c) \in B \to f(c) \in \mathrm{dom}(g)$}
\AXC{$\varphi_{199}$\lbl{A1}}
\BIC{$(\mathrm{ContIn}(f,A,c) \to \varphi_{11}) \to \varphi_{189}$}
\end{ptree}
\noindent$\Pi_{36}$:
\begin{ptree}
\AXC{$\stackrel{\Pi_{34}}{\vdots}$}\noLine
\UIC{$\mathrm{ContIn}(f,A,c) \to \varphi_{132}$}
\AXC{$\stackrel{\Pi_{11}}{\vdots}$}\noLine
\UIC{$\mathrm{ContIn}(f,A,c) \to \varphi_{11}$}
\AXC{$\stackrel{\Pi_{35}}{\vdots}$}\noLine
\UIC{$(\mathrm{ContIn}(f,A,c) \to \varphi_{11}) \to \varphi_{189}$}
\BIC{$\varphi_{189}$}
\BIC{$\mathrm{ContIn}(f,A,c) \to \varphi_{136}$}
\end{ptree}
\noindent$\Pi_{37}$:
\begin{ptree}
\AXC{$\forall X\,\forall Y\,\forall x\,(X \subseteq Y \to x \in X \to x \in Y)$\lbl{[B5]}}
\AXC{$\varphi_{19}$\lbl{A2}}
\BIC{$\forall Y\,\forall x\,(A \subseteq Y \to x \in A \to x \in Y)$}
\AXC{$\varphi_{30}$\lbl{A2}}
\BIC{$\forall x\,(A \subseteq \mathrm{dom}(f) \to x \in A \to x \in \mathrm{dom}(f))$}
\AXC{$\varphi_{37}$\lbl{A2}}
\BIC{$A \subseteq \mathrm{dom}(f) \to c \in A \to c \in \mathrm{dom}(f)$}
\AXC{$\varphi_{183}$\lbl{A1}}
\BIC{$(\mathrm{ContIn}(f,A,c) \to \varphi_{12}) \to \varphi_{173}$}
\end{ptree}
\noindent$\Pi_{38}$:
\begin{ptree}
\AXC{$\stackrel{\Pi_{32}}{\vdots}$}\noLine
\UIC{$\mathrm{ContIn}(f,A,c) \to \varphi_{9}$}
\AXC{$\stackrel{\Pi_{15}}{\vdots}$}\noLine
\UIC{$\mathrm{ContIn}(f,A,c) \to \varphi_{12}$}
\AXC{$\stackrel{\Pi_{37}}{\vdots}$}\noLine
\UIC{$(\mathrm{ContIn}(f,A,c) \to \varphi_{12}) \to \varphi_{173}$}
\BIC{$\varphi_{173}$}
\BIC{$\mathrm{ContIn}(f,A,c) \to \varphi_{130}$}
\end{ptree}
\noindent$\Pi_{39}$:
\begin{ptree}
\AXC{$\forall g\,\forall f\,\forall x\,(\mathrm{RFn}(g) \to \mathrm{RFn}(f) \to \varphi_{2})$\lbl{[S23]}}
\AXC{$\varphi_{71}$\lbl{A2}}
\BIC{$\forall f\,\forall x\,(\mathrm{RFn}(g) \to \mathrm{RFn}(f) \to \varphi_{2})$}
\AXC{$\varphi_{69}$\lbl{A2}}
\BIC{$\forall x\,(\mathrm{RFn}(g) \to \mathrm{RFn}(f) \to \varphi_{2})$}
\AXC{$\varphi_{67}$\lbl{A2}}
\BIC{$\mathrm{RFn}(g) \to \mathrm{RFn}(f) \to \varphi_{3}$}
\AXC{$\varphi_{211}$\lbl{A1}}
\BIC{$(\mathrm{ContIn}(f,A,c) \to \varphi_{6}) \to \varphi_{207}$}
\end{ptree}
\noindent$\Pi_{40}$:
\begin{ptree}
\AXC{$\stackrel{\Pi_{18}}{\vdots}$}\noLine
\UIC{$\mathrm{ContIn}(f,A,c) \to \varphi_{5}$}
\AXC{$\stackrel{\Pi_{19}}{\vdots}$}\noLine
\UIC{$\mathrm{ContIn}(f,A,c) \to \varphi_{6}$}
\AXC{$\stackrel{\Pi_{39}}{\vdots}$}\noLine
\UIC{$(\mathrm{ContIn}(f,A,c) \to \varphi_{6}) \to \varphi_{207}$}
\BIC{$\varphi_{207}$}
\BIC{$(\mathrm{ContIn}(f,A,c) \to \varphi_{130}) \to \varphi_{191}$}
\end{ptree}
\noindent$\Pi_{41}$:
\begin{ptree}
\AXC{$\stackrel{\Pi_{36}}{\vdots}$}\noLine
\UIC{$\mathrm{ContIn}(f,A,c) \to \varphi_{136}$}
\AXC{$\stackrel{\Pi_{38}}{\vdots}$}\noLine
\UIC{$\mathrm{ContIn}(f,A,c) \to \varphi_{130}$}
\AXC{$\stackrel{\Pi_{40}}{\vdots}$}\noLine
\UIC{$(\mathrm{ContIn}(f,A,c) \to \varphi_{130}) \to \varphi_{191}$}
\BIC{$\varphi_{191}$}
\BIC{$\mathrm{ContIn}(f,A,c) \to \varphi_{138}$}
\end{ptree}
\noindent$\Pi_{42}$:
\begin{ptree}
\AXC{$\forall x \in \mathrm{dom}(g \circ f)\,(\varphi_{14}) \to c \in \mathrm{dom}(g \circ f) \to \varphi_{15}$\lbl{A2}}
\AXC{$\stackrel{\Pi_{25}}{\vdots}$}\noLine
\UIC{$\mathrm{ContIn}(f,A,c) \to \varphi_{62}$}
\AXC{$(\mathrm{ContIn}(f,A,c) \to \varphi_{62}) \to \varphi_{157}$\lbl{A1}}
\BIC{$\varphi_{157}$}
\BIC{$\mathrm{ContIn}(f,A,c) \to \varphi_{61}$}
\end{ptree}
\noindent$\Pi_{43}$:
\begin{ptree}
\AXC{$\stackrel{\Pi_{41}}{\vdots}$}\noLine
\UIC{$\mathrm{ContIn}(f,A,c) \to \varphi_{138}$}
\AXC{$\stackrel{\Pi_{42}}{\vdots}$}\noLine
\UIC{$\mathrm{ContIn}(f,A,c) \to \varphi_{61}$}
\AXC{$(\mathrm{ContIn}(f,A,c) \to \varphi_{61}) \to \varphi_{194}$\lbl{A1}}
\BIC{$\varphi_{194}$}
\BIC{$\mathrm{ContIn}(f,A,c) \to \varphi_{149}$}
\end{ptree}
\noindent$\Pi_{44}$:
\begin{ptree}
\AXC{$\stackrel{\Pi_{43}}{\vdots}$}\noLine
\UIC{$\mathrm{ContIn}(f,A,c) \to \varphi_{149}$}
\AXC{$(\mathrm{ContIn}(f,A,c) \to \varphi_{149}) \to \varphi_{176}$\lbl{A1}}
\BIC{$\varphi_{176}$}
\end{ptree}
\noindent$\Pi_{45}$:
\begin{ptree}
\AXC{$\stackrel{\Pi_{31}}{\vdots}$}\noLine
\UIC{$\varphi_{159}$}
\AXC{$\stackrel{\Pi_{7}}{\vdots}$}\noLine
\UIC{$\forall x\,\forall y\,(x = y \to y = x)$}
\AXC{$\varphi_{8}$\lbl{A2}}
\BIC{$\forall y\,(g \circ f(c) = y \to y = g \circ f(c))$}
\AXC{$\varphi_{55}$\lbl{A2}}
\BIC{$g \circ f(c) = g(f(c)) \to g(f(c)) = g \circ f(c)$}
\AXC{$\stackrel{\Pi_{44}}{\vdots}$}\noLine
\UIC{$\varphi_{176}$}
\BIC{$\mathrm{ContIn}(f,A,c) \to \varphi_{150}$}
\AXC{$\varphi_{209}$\lbl{A1}}
\BIC{$\varphi_{159} \to (g(f(c)) = g \circ f(c) \to \varphi_{49}) \to \varphi_{160}$}
\BIC{$(g(f(c)) = g \circ f(c) \to \varphi_{49}) \to \varphi_{160}$}
\end{ptree}
\noindent$\Pi_{46}$:
\begin{ptree}
\AXC{$\stackrel{\Pi_{1}}{\vdots}$}\noLine
\UIC{$g(f(c)) = g \circ f(c) \to \varphi_{49}$}
\AXC{$\stackrel{\Pi_{45}}{\vdots}$}\noLine
\UIC{$(g(f(c)) = g \circ f(c) \to \varphi_{49}) \to \varphi_{160}$}
\BIC{$\varphi_{160}$}
\AXC{$\varphi_{160} \to \varphi_{161}$\lbl{A1}}
\BIC{$\varphi_{161}$}
\UIC{$\forall x\,(\varphi_{161})$}
\end{ptree}
\noindent$\Pi_{47}$:
\begin{ptree}
\AXC{$\forall x\,(\varphi_{161}) \to \varphi_{120} \to \varphi_{34}$\lbl{A3}}
\AXC{$\stackrel{\Pi_{46}}{\vdots}$}\noLine
\UIC{$\forall x\,(\varphi_{161})$}
\AXC{$\forall x\,(\varphi_{161}) \to (\forall x\,(\varphi_{161}) \to \varphi_{120} \to \varphi_{34}) \to \varphi_{200}$\lbl{A1}}
\BIC{$(\forall x\,(\varphi_{161}) \to \varphi_{120} \to \varphi_{34}) \to \varphi_{200}$}
\BIC{$\varphi_{200}$}
\end{ptree}
\noindent$\Pi_{48}$:
\begin{ptree}
\AXC{$\varphi_{10} \land \varepsilon \in \mathbb{R} \land 0 < \varepsilon \land (d_{1} \in \mathbb{R} \land 0 < d_{1} \land \varphi_{31}) \land \varphi_{43} \to \varphi_{120}$\lbl{A4}}
\AXC{$\stackrel{\Pi_{47}}{\vdots}$}\noLine
\UIC{$\varphi_{200}$}
\BIC{$\varphi_{163}$}
\AXC{$\varphi_{164}$\lbl{A1}}
\AXC{$\varphi_{164} \to \varphi_{163} \to \mathrm{ContIn}(f,A,c) \to \varphi_{166}$\lbl{A1}}
\BIC{$\varphi_{163} \to \mathrm{ContIn}(f,A,c) \to \varphi_{166}$}
\BIC{$\mathrm{ContIn}(f,A,c) \to \varphi_{166}$}
\end{ptree}
\noindent$\Pi_{49}$:
\begin{ptree}
\AXC{$\varphi_{95}$\lbl{A2}}
\AXC{$\stackrel{\Pi_{48}}{\vdots}$}\noLine
\UIC{$\mathrm{ContIn}(f,A,c) \to \varphi_{166}$}
\AXC{$\varphi_{203}$\lbl{A1}}
\BIC{$\varphi_{95} \to \mathrm{ContIn}(f,A,c) \to \varphi_{169}$}
\BIC{$\mathrm{ContIn}(f,A,c) \to \varphi_{169}$}
\AXC{$(\mathrm{ContIn}(f,A,c) \to \varphi_{169}) \to \varphi_{125} \to \neg \varphi_{43}$\lbl{A1}}
\BIC{$\varphi_{125} \to \neg \varphi_{43}$}
\UIC{$\forall \delta\,(\varphi_{125} \to \neg \varphi_{43})$}
\end{ptree}
\noindent$\Pi_{50}$:
\begin{ptree}
\AXC{$\stackrel{\Pi_{14}}{\vdots}$}\noLine
\UIC{$\varphi_{76}$}
\AXC{$\mathrm{ContIn}(f,A,c) \to \mathrm{ContIn}(f,A,c)$\lbl{A1}}
\AXC{$\varphi_{171}$\lbl{A1}}
\BIC{$\varphi_{76} \to \mathrm{ContIn}(f,A,c) \to \varphi_{77}$}
\BIC{$\mathrm{ContIn}(f,A,c) \to \varphi_{77}$}
\AXC{$(\mathrm{ContIn}(f,A,c) \to \varphi_{77}) \to \varphi_{178}$\lbl{A1}}
\BIC{$\varphi_{178}$}
\end{ptree}
\noindent$\Pi_{51}$:
\begin{ptree}
\AXC{$\mathrm{ContIn}(f,A,c) \to \varphi_{118}$\lbl{A1}}
\AXC{$\varphi_{56} \to d_{1} \in \mathbb{R} \to 0 < d_{1} \to \exists \delta\,(\varphi_{43})$\lbl{A2}}
\AXC{$\stackrel{\Pi_{50}}{\vdots}$}\noLine
\UIC{$\varphi_{178}$}
\BIC{$\mathrm{ContIn}(f,A,c) \to \varphi_{74}$}
\AXC{$\varphi_{197}$\lbl{A1}}
\BIC{$(\mathrm{ContIn}(f,A,c) \to \varphi_{118}) \to \varphi_{121}$}
\BIC{$\varphi_{121}$}
\end{ptree}
\noindent$\Pi_{52}$:
\begin{ptree}
\AXC{$\forall \delta\,(\varphi_{125} \to \neg \varphi_{43}) \to \forall \delta\,(\varphi_{125}) \to \forall \delta\,(\neg \varphi_{43})$\lbl{A3}}
\AXC{$\stackrel{\Pi_{49}}{\vdots}$}\noLine
\UIC{$\forall \delta\,(\varphi_{125} \to \neg \varphi_{43})$}
\AXC{$\stackrel{\Pi_{51}}{\vdots}$}\noLine
\UIC{$\varphi_{121}$}
\AXC{$\varphi_{121} \to \forall \delta\,(\varphi_{125} \to \neg \varphi_{43}) \to \varphi_{214}$\lbl{A1}}
\BIC{$\forall \delta\,(\varphi_{125} \to \neg \varphi_{43}) \to \varphi_{214}$}
\BIC{$\varphi_{214}$}
\BIC{$(\varphi_{125} \to \forall \delta\,(\varphi_{125})) \to \mathrm{ContIn}(f,A,c) \to \varphi_{123}$}
\end{ptree}
\noindent$\Pi_{53}$:
\begin{ptree}
\AXC{$\varphi_{125} \to \forall \delta\,(\varphi_{125})$\lbl{A4}}
\AXC{$\stackrel{\Pi_{52}}{\vdots}$}\noLine
\UIC{$(\varphi_{125} \to \forall \delta\,(\varphi_{125})) \to \mathrm{ContIn}(f,A,c) \to \varphi_{123}$}
\BIC{$\mathrm{ContIn}(f,A,c) \to \varphi_{123}$}
\AXC{$(\mathrm{ContIn}(f,A,c) \to \varphi_{123}) \to \varphi_{126}$\lbl{A1}}
\BIC{$\varphi_{126}$}
\UIC{$\forall d_{1}\,(\varphi_{126})$}
\end{ptree}
\noindent$\Pi_{54}$:
\begin{ptree}
\AXC{$\stackrel{\Pi_{10}}{\vdots}$}\noLine
\UIC{$\varphi_{88}$}
\AXC{$\mathrm{ContIn}(f,A,c) \to \varphi_{1}$\lbl{A1}}
\AXC{$\varphi_{180}$\lbl{A1}}
\BIC{$\varphi_{88} \to \mathrm{ContIn}(f,A,c) \to \varphi_{83}$}
\BIC{$\mathrm{ContIn}(f,A,c) \to \varphi_{83}$}
\AXC{$(\mathrm{ContIn}(f,A,c) \to \varphi_{83}) \to \varphi_{181}$\lbl{A1}}
\BIC{$\varphi_{181}$}
\end{ptree}
\noindent$\Pi_{55}$:
\begin{ptree}
\AXC{$\forall \varepsilon \in \mathbb{R}\,(0 < \varepsilon \to \varphi_{53}) \to \varepsilon \in \mathbb{R} \to 0 < \varepsilon \to \varphi_{53}$\lbl{A2}}
\AXC{$\stackrel{\Pi_{54}}{\vdots}$}\noLine
\UIC{$\varphi_{181}$}
\BIC{$\mathrm{ContIn}(f,A,c) \to \varphi_{82}$}
\AXC{$(\mathrm{ContIn}(f,A,c) \to \varphi_{82}) \to \forall d_{1}\,(\varphi_{126}) \to \varphi_{215}$\lbl{A1}}
\BIC{$\forall d_{1}\,(\varphi_{126}) \to \varphi_{215}$}
\end{ptree}
\noindent$\Pi_{56}$:
\begin{ptree}
\AXC{$\forall d_{1}\,(\varphi_{126}) \to \varphi_{91} \to \forall d_{1}\,(\neg (d_{1} \in \mathbb{R} \land 0 < d_{1} \land \varphi_{31}))$\lbl{A3}}
\AXC{$\stackrel{\Pi_{53}}{\vdots}$}\noLine
\UIC{$\forall d_{1}\,(\varphi_{126})$}
\AXC{$\stackrel{\Pi_{55}}{\vdots}$}\noLine
\UIC{$\forall d_{1}\,(\varphi_{126}) \to \varphi_{215}$}
\BIC{$\varphi_{215}$}
\BIC{$\varphi_{212}$}
\end{ptree}
\noindent$\Pi_{57}$:
\begin{ptree}
\AXC{$\varphi_{179}$\lbl{A3}}
\AXC{$\forall d_{1}\,(\neg (d_{1} \in \mathbb{R} \land 0 < d_{1} \land \varphi_{31})) \to \varphi_{51}$\lbl{A2}}
\UIC{$\forall \delta\,(\forall d_{1}\,(\neg (d_{1} \in \mathbb{R} \land 0 < d_{1} \land \varphi_{31})) \to \varphi_{51})$}
\AXC{$\varphi_{10} \land \varepsilon \in \mathbb{R} \land 0 < \varepsilon \land \neg \exists \delta\,(\delta \in \mathbb{R} \land 0 < \delta \land \varphi_{34}) \to \varphi_{91}$\lbl{A4}}
\AXC{$\stackrel{\Pi_{56}}{\vdots}$}\noLine
\UIC{$\varphi_{212}$}
\BIC{$\varphi_{205}$}
\BIC{$\varphi_{179} \to \varphi_{96} \to \varphi_{89}$}
\BIC{$\varphi_{96} \to \varphi_{89}$}
\end{ptree}
\noindent$\Pi_{58}$:
\begin{ptree}
\AXC{$\varphi_{177}$\lbl{A3}}
\AXC{$\varphi_{96}$\lbl{A4}}
\AXC{$\stackrel{\Pi_{57}}{\vdots}$}\noLine
\UIC{$\varphi_{96} \to \varphi_{89}$}
\BIC{$\varphi_{89}$}
\AXC{$\varphi_{89} \to \varphi_{10} \to \varepsilon \in \mathbb{R} \to 0 < \varepsilon \to \exists \delta\,(\delta \in \mathbb{R} \land 0 < \delta \land \varphi_{34})$\lbl{A1}}
\BIC{$\varphi_{10} \to \varepsilon \in \mathbb{R} \to 0 < \varepsilon \to \exists \delta\,(\delta \in \mathbb{R} \land 0 < \delta \land \varphi_{34})$}
\UIC{$\forall \varepsilon\,(\varphi_{10} \to \varepsilon \in \mathbb{R} \to 0 < \varepsilon \to \exists \delta\,(\delta \in \mathbb{R} \land 0 < \delta \land \varphi_{34}))$}
\AXC{$\varphi_{201}$\lbl{A1}}
\BIC{$\varphi_{177} \to (\varphi_{10} \to \forall \varepsilon\,(\varphi_{10})) \to \varphi_{90}$}
\BIC{$(\varphi_{10} \to \forall \varepsilon\,(\varphi_{10})) \to \varphi_{90}$}
\end{ptree}
\noindent$\Pi_{59}$:
\begin{ptree}
\AXC{$\forall x \in A\,(f(x) \in B) \to u \in A \to f(u) \in B$\lbl{A2}}
\AXC{$\mathrm{ContIn}(f,A,c) \to \varphi_{32}$\lbl{A1}}
\AXC{$(\mathrm{ContIn}(f,A,c) \to \varphi_{32}) \to \varphi_{65}$\lbl{A1}}
\BIC{$\varphi_{65}$}
\BIC{$\mathrm{ContIn}(f,A,c) \to \varphi_{27}$}
\end{ptree}
\noindent$\Pi_{60}$:
\begin{ptree}
\AXC{$\forall X\,\forall Y\,\forall x\,(X \subseteq Y \to x \in X \to x \in Y)$\lbl{[B5]}}
\AXC{$\varphi_{18}$\lbl{A2}}
\BIC{$\forall Y\,\forall x\,(B \subseteq Y \to x \in B \to x \in Y)$}
\AXC{$\varphi_{29}$\lbl{A2}}
\BIC{$\forall x\,(B \subseteq \mathrm{dom}(g) \to x \in B \to x \in \mathrm{dom}(g))$}
\AXC{$\varphi_{47}$\lbl{A2}}
\BIC{$B \subseteq \mathrm{dom}(g) \to f(u) \in B \to f(u) \in \mathrm{dom}(g)$}
\AXC{$\varphi_{143}$\lbl{A1}}
\BIC{$(\mathrm{ContIn}(f,A,c) \to \varphi_{11}) \to \varphi_{84}$}
\end{ptree}
\noindent$\Pi_{61}$:
\begin{ptree}
\AXC{$\stackrel{\Pi_{59}}{\vdots}$}\noLine
\UIC{$\mathrm{ContIn}(f,A,c) \to \varphi_{27}$}
\AXC{$\stackrel{\Pi_{11}}{\vdots}$}\noLine
\UIC{$\mathrm{ContIn}(f,A,c) \to \varphi_{11}$}
\AXC{$\stackrel{\Pi_{60}}{\vdots}$}\noLine
\UIC{$(\mathrm{ContIn}(f,A,c) \to \varphi_{11}) \to \varphi_{84}$}
\BIC{$\varphi_{84}$}
\BIC{$\mathrm{ContIn}(f,A,c) \to \varphi_{39}$}
\end{ptree}
\noindent$\Pi_{62}$:
\begin{ptree}
\AXC{$\forall X\,\forall Y\,\forall x\,(X \subseteq Y \to x \in X \to x \in Y)$\lbl{[B5]}}
\AXC{$\varphi_{19}$\lbl{A2}}
\BIC{$\forall Y\,\forall x\,(A \subseteq Y \to x \in A \to x \in Y)$}
\AXC{$\varphi_{30}$\lbl{A2}}
\BIC{$\forall x\,(A \subseteq \mathrm{dom}(f) \to x \in A \to x \in \mathrm{dom}(f))$}
\AXC{$\varphi_{38}$\lbl{A2}}
\BIC{$A \subseteq \mathrm{dom}(f) \to u \in A \to u \in \mathrm{dom}(f)$}
\AXC{$\varphi_{124}$\lbl{A1}}
\BIC{$(\mathrm{ContIn}(f,A,c) \to \varphi_{12}) \to \varphi_{79}$}
\end{ptree}
\noindent$\Pi_{63}$:
\begin{ptree}
\AXC{$\mathrm{ContIn}(f,A,c) \to \varphi_{16}$\lbl{A1}}
\AXC{$\stackrel{\Pi_{15}}{\vdots}$}\noLine
\UIC{$\mathrm{ContIn}(f,A,c) \to \varphi_{12}$}
\AXC{$\stackrel{\Pi_{62}}{\vdots}$}\noLine
\UIC{$(\mathrm{ContIn}(f,A,c) \to \varphi_{12}) \to \varphi_{79}$}
\BIC{$\varphi_{79}$}
\BIC{$\mathrm{ContIn}(f,A,c) \to \varphi_{23}$}
\end{ptree}
\noindent$\Pi_{64}$:
\begin{ptree}
\AXC{$\forall g\,\forall f\,\forall x\,(\mathrm{RFn}(g) \to \mathrm{RFn}(f) \to \varphi_{2})$\lbl{[S23]}}
\AXC{$\varphi_{71}$\lbl{A2}}
\BIC{$\forall f\,\forall x\,(\mathrm{RFn}(g) \to \mathrm{RFn}(f) \to \varphi_{2})$}
\AXC{$\varphi_{69}$\lbl{A2}}
\BIC{$\forall x\,(\mathrm{RFn}(g) \to \mathrm{RFn}(f) \to \varphi_{2})$}
\AXC{$\varphi_{68}$\lbl{A2}}
\BIC{$\mathrm{RFn}(g) \to \mathrm{RFn}(f) \to \varphi_{4}$}
\AXC{$\varphi_{182}$\lbl{A1}}
\BIC{$(\mathrm{ContIn}(f,A,c) \to \varphi_{6}) \to \varphi_{162}$}
\end{ptree}
\noindent$\Pi_{65}$:
\begin{ptree}
\AXC{$\stackrel{\Pi_{18}}{\vdots}$}\noLine
\UIC{$\mathrm{ContIn}(f,A,c) \to \varphi_{5}$}
\AXC{$\stackrel{\Pi_{19}}{\vdots}$}\noLine
\UIC{$\mathrm{ContIn}(f,A,c) \to \varphi_{6}$}
\AXC{$\stackrel{\Pi_{64}}{\vdots}$}\noLine
\UIC{$(\mathrm{ContIn}(f,A,c) \to \varphi_{6}) \to \varphi_{162}$}
\BIC{$\varphi_{162}$}
\BIC{$(\mathrm{ContIn}(f,A,c) \to \varphi_{23}) \to \varphi_{86}$}
\end{ptree}
\noindent$\Pi_{66}$:
\begin{ptree}
\AXC{$\stackrel{\Pi_{61}}{\vdots}$}\noLine
\UIC{$\mathrm{ContIn}(f,A,c) \to \varphi_{39}$}
\AXC{$\stackrel{\Pi_{63}}{\vdots}$}\noLine
\UIC{$\mathrm{ContIn}(f,A,c) \to \varphi_{23}$}
\AXC{$\stackrel{\Pi_{65}}{\vdots}$}\noLine
\UIC{$(\mathrm{ContIn}(f,A,c) \to \varphi_{23}) \to \varphi_{86}$}
\BIC{$\varphi_{86}$}
\BIC{$\mathrm{ContIn}(f,A,c) \to \varphi_{40}$}
\AXC{$\varphi_{87}$\lbl{A1}}
\BIC{$\varphi_{10} \to u \in A \to u \in \mathrm{dom}(g \circ f)$}
\UIC{$\forall u\,(\varphi_{10} \to u \in A \to u \in \mathrm{dom}(g \circ f))$}
\end{ptree}
\noindent$\Pi_{67}$:
\begin{ptree}
\AXC{$\varphi_{10} \to \forall u\,(\varphi_{10})$\lbl{A4}}
\AXC{$\varphi_{92}$\lbl{A3}}
\AXC{$\stackrel{\Pi_{66}}{\vdots}$}\noLine
\UIC{$\forall u\,(\varphi_{10} \to u \in A \to u \in \mathrm{dom}(g \circ f))$}
\AXC{$\varphi_{185}$\lbl{A1}}
\BIC{$\varphi_{92} \to (\varphi_{10} \to \forall u\,(\varphi_{10})) \to \mathrm{ContIn}(f,A,c) \to \varphi_{42}$}
\BIC{$(\varphi_{10} \to \forall u\,(\varphi_{10})) \to \mathrm{ContIn}(f,A,c) \to \varphi_{42}$}
\BIC{$\mathrm{ContIn}(f,A,c) \to \varphi_{42}$}
\end{ptree}
\noindent$\Pi_{68}$:
\begin{ptree}
\AXC{$\stackrel{\Pi_{67}}{\vdots}$}\noLine
\UIC{$\mathrm{ContIn}(f,A,c) \to \varphi_{42}$}
\AXC{$\forall x\,\forall y\,(x \subseteq y \leftrightarrow \forall u \in x\,(u \in y))$\lbl{Def$\subseteq$}}
\AXC{$\varphi_{20}$\lbl{A2}}
\BIC{$\forall y\,(A \subseteq y \leftrightarrow \forall u \in A\,(u \in y))$}
\AXC{$\varphi_{44}$\lbl{A2}}
\BIC{$A \subseteq \mathrm{dom}(g \circ f) \leftrightarrow \forall u \in A\,(u \in \mathrm{dom}(g \circ f))$}
\AXC{$\varphi_{110}$\lbl{A1}}
\BIC{$\varphi_{85}$}
\BIC{$\mathrm{ContIn}(f,A,c) \to \varphi_{21}$}
\end{ptree}
\noindent$\Pi_{69}$:
\begin{ptree}
\AXC{$\forall f\,\forall A\,\forall c\,(\varphi_{76})$\lbl{DefContIn}}
\AXC{$\forall f\,\forall A\,\forall c\,(\varphi_{76}) \to \forall A\,\forall c\,(\varphi_{94})$\lbl{A2}}
\BIC{$\forall A\,\forall c\,(\varphi_{94})$}
\AXC{$\forall A\,\forall c\,(\varphi_{94}) \to \forall c\,(\varphi_{94})$\lbl{A2}}
\BIC{$\forall c\,(\varphi_{94})$}
\AXC{$\forall c\,(\varphi_{94}) \to \varphi_{94}$\lbl{A2}}
\BIC{$\varphi_{94}$}
\end{ptree}
\noindent$\Pi_{70}$:
\begin{ptree}
\AXC{$\stackrel{\Pi_{68}}{\vdots}$}\noLine
\UIC{$\mathrm{ContIn}(f,A,c) \to \varphi_{21}$}
\AXC{$\stackrel{\Pi_{24}}{\vdots}$}\noLine
\UIC{$\varphi_{81}$}
\AXC{$\varphi_{81} \to \mathrm{ContIn}(f,A,c) \to \varphi_{13}$\lbl{A1}}
\BIC{$\mathrm{ContIn}(f,A,c) \to \varphi_{13}$}
\AXC{$\stackrel{\Pi_{69}}{\vdots}$}\noLine
\UIC{$\varphi_{94}$}
\AXC{$\varphi_{94} \to \varphi_{184}$\lbl{A1}}
\BIC{$\varphi_{184}$}
\BIC{$(\mathrm{ContIn}(f,A,c) \to \varphi_{21}) \to \varphi_{165}$}
\BIC{$\varphi_{165}$}
\end{ptree}
\noindent 主木:
\begin{ptree}
\AXC{$\varphi_{10} \to \forall \varepsilon\,(\varphi_{10})$\lbl{A4}}
\AXC{$\stackrel{\Pi_{58}}{\vdots}$}\noLine
\UIC{$(\varphi_{10} \to \forall \varepsilon\,(\varphi_{10})) \to \varphi_{90}$}
\BIC{$\varphi_{90}$}
\AXC{$\stackrel{\Pi_{32}}{\vdots}$}\noLine
\UIC{$\mathrm{ContIn}(f,A,c) \to \varphi_{9}$}
\AXC{$\stackrel{\Pi_{70}}{\vdots}$}\noLine
\UIC{$\varphi_{165}$}
\BIC{$\varphi_{90} \to \mathrm{ContIn}(f,A,c) \to \varphi_{22}$}
\BIC{$\mathrm{ContIn}(f,A,c) \to \varphi_{22}$}
\UIC{$\forall c\,(\mathrm{ContIn}(f,A,c) \to \varphi_{22})$}
\UIC{$\forall B\,\forall c\,(\mathrm{ContIn}(f,A,c) \to \varphi_{22})$}
\UIC{$\forall A\,\forall B\,\forall c\,(\mathrm{ContIn}(f,A,c) \to \varphi_{22})$}
\UIC{$\forall g\,\forall A\,\forall B\,\forall c\,(\mathrm{ContIn}(f,A,c) \to \varphi_{22})$}
\UIC{$\forall f\,\forall g\,\forall A\,\forall B\,\forall c\,(\mathrm{ContIn}(f,A,c) \to \varphi_{22})$}
\end{ptree}
\begin{abbrevlist}
\item $\varphi_{1} :\equiv \mathrm{ContIn}(g,B,f(c)) \to \mathrm{ContIn}(g,B,f(c))$
\item $\varphi_{2} :\equiv x \in \mathrm{dom}(f) \to f(x) \in \mathrm{dom}(g) \to x \in \mathrm{dom}(g \circ f)$
\item $\varphi_{3} :\equiv c \in \mathrm{dom}(f) \to f(c) \in \mathrm{dom}(g) \to c \in \mathrm{dom}(g \circ f)$
\item $\varphi_{4} :\equiv u \in \mathrm{dom}(f) \to f(u) \in \mathrm{dom}(g) \to u \in \mathrm{dom}(g \circ f)$
\item $\varphi_{5} :\equiv \mathrm{ContIn}(g,B,f(c)) \to \forall x \in A\,(f(x) \in B) \to \mathrm{RFn}(f)$
\item $\varphi_{6} :\equiv \mathrm{ContIn}(g,B,f(c)) \to \forall x \in A\,(f(x) \in B) \to \mathrm{RFn}(g)$
\item $\varphi_{7} :\equiv \forall x\,\forall y\,(x = y \to y = x) \to \forall y\,(g \circ f(x) = y \to y = g \circ f(x))$
\item $\varphi_{8} :\equiv \forall x\,\forall y\,(x = y \to y = x) \to \forall y\,(g \circ f(c) = y \to y = g \circ f(c))$
\item $\varphi_{9} :\equiv \mathrm{ContIn}(g,B,f(c)) \to \forall x \in A\,(f(x) \in B) \to c \in A$
\item $\varphi_{10} :\equiv \mathrm{ContIn}(f,A,c) \land \mathrm{ContIn}(g,B,f(c)) \land \forall x \in A\,(f(x) \in B)$
\item $\varphi_{11} :\equiv \mathrm{ContIn}(g,B,f(c)) \to \forall x \in A\,(f(x) \in B) \to B \subseteq \mathrm{dom}(g)$
\item $\varphi_{12} :\equiv \mathrm{ContIn}(g,B,f(c)) \to \forall x \in A\,(f(x) \in B) \to A \subseteq \mathrm{dom}(f)$
\item $\varphi_{13} :\equiv \mathrm{ContIn}(g,B,f(c)) \to \forall x \in A\,(f(x) \in B) \to \mathrm{RFn}(g \circ f)$
\item $\varphi_{14} :\equiv x \in \mathrm{dom}(f) \land f(x) \in \mathrm{dom}(g) \land g \circ f(x) = g(f(x))$
\item $\varphi_{15} :\equiv c \in \mathrm{dom}(f) \land f(c) \in \mathrm{dom}(g) \land g \circ f(c) = g(f(c))$
\item $\varphi_{16} :\equiv \mathrm{ContIn}(g,B,f(c)) \to \forall x \in A\,(f(x) \in B) \to u \in A \to u \in A$
\item $\varphi_{17} :\equiv z = w \to \lvert z - g(f(c)) \rvert < \varepsilon \to \lvert w - g(f(c)) \rvert < \varepsilon$
\item $\varphi_{18} :\equiv \forall X\,\forall Y\,\forall x\,(X \subseteq Y \to x \in X \to x \in Y) \to \forall Y\,\forall x\,(B \subseteq Y \to x \in B \to x \in Y)$
\item $\varphi_{19} :\equiv \forall X\,\forall Y\,\forall x\,(X \subseteq Y \to x \in X \to x \in Y) \to \forall Y\,\forall x\,(A \subseteq Y \to x \in A \to x \in Y)$
\item $\varphi_{20} :\equiv \forall x\,\forall y\,(x \subseteq y \leftrightarrow \forall u \in x\,(u \in y)) \to \forall y\,(A \subseteq y \leftrightarrow \forall u \in A\,(u \in y))$
\item $\varphi_{21} :\equiv \mathrm{ContIn}(g,B,f(c)) \to \forall x \in A\,(f(x) \in B) \to A \subseteq \mathrm{dom}(g \circ f)$
\item $\varphi_{22} :\equiv \mathrm{ContIn}(g,B,f(c)) \to \forall x \in A\,(f(x) \in B) \to \mathrm{ContIn}(g \circ f,A,c)$
\item $\varphi_{23} :\equiv \mathrm{ContIn}(g,B,f(c)) \to \forall x \in A\,(f(x) \in B) \to u \in A \to u \in \mathrm{dom}(f)$
\item $\varphi_{24} :\equiv x \in A \to \lvert x - c \rvert < \delta \to \lvert g(f(x)) - g(f(c)) \rvert < \varepsilon$
\item $\varphi_{25} :\equiv \mathrm{ContIn}(g,B,f(c)) \to \forall x \in A\,(f(x) \in B) \to x \in A \to f(x) \in B$
\item $\varphi_{26} :\equiv \mathrm{ContIn}(g,B,f(c)) \to \forall x \in A\,(f(x) \in B) \to c \in A \to f(c) \in B$
\item $\varphi_{27} :\equiv \mathrm{ContIn}(g,B,f(c)) \to \forall x \in A\,(f(x) \in B) \to u \in A \to f(u) \in B$
\item $\varphi_{28} :\equiv \forall w\,(z = w \to \lvert g \circ f(x) - z \rvert < \varepsilon \to \lvert g \circ f(x) - w \rvert < \varepsilon)$
\item $\varphi_{29} :\equiv \forall Y\,\forall x\,(B \subseteq Y \to x \in B \to x \in Y) \to \forall x\,(B \subseteq \mathrm{dom}(g) \to x \in B \to x \in \mathrm{dom}(g))$
\item $\varphi_{30} :\equiv \forall Y\,\forall x\,(A \subseteq Y \to x \in A \to x \in Y) \to \forall x\,(A \subseteq \mathrm{dom}(f) \to x \in A \to x \in \mathrm{dom}(f))$
\item $\varphi_{31} :\equiv \forall x \in B\,(\lvert x - f(c) \rvert < d_{1} \to \lvert g(x) - g(f(c)) \rvert < \varepsilon)$
\item $\varphi_{32} :\equiv \mathrm{ContIn}(g,B,f(c)) \to \forall x \in A\,(f(x) \in B) \to \forall x \in A\,(f(x) \in B)$
\item $\varphi_{33} :\equiv \lvert f(x) - f(c) \rvert < d_{1} \to \lvert g(f(x)) - g(f(c)) \rvert < \varepsilon$
\item $\varphi_{34} :\equiv \forall x \in A\,(\lvert x - c \rvert < \delta \to \lvert g \circ f(x) - g \circ f(c) \rvert < \varepsilon)$
\item $\varphi_{35} :\equiv \lvert x - c \rvert < \delta \to x \in A \to \lvert x - c \rvert < \delta \to \lvert f(x) - f(c) \rvert < d_{1}$
\item $\varphi_{36} :\equiv \forall x\,(A \subseteq \mathrm{dom}(f) \to x \in A \to x \in \mathrm{dom}(f)) \to A \subseteq \mathrm{dom}(f) \to x \in A \to x \in \mathrm{dom}(f)$
\item $\varphi_{37} :\equiv \forall x\,(A \subseteq \mathrm{dom}(f) \to x \in A \to x \in \mathrm{dom}(f)) \to A \subseteq \mathrm{dom}(f) \to c \in A \to c \in \mathrm{dom}(f)$
\item $\varphi_{38} :\equiv \forall x\,(A \subseteq \mathrm{dom}(f) \to x \in A \to x \in \mathrm{dom}(f)) \to A \subseteq \mathrm{dom}(f) \to u \in A \to u \in \mathrm{dom}(f)$
\item $\varphi_{39} :\equiv \mathrm{ContIn}(g,B,f(c)) \to \forall x \in A\,(f(x) \in B) \to u \in A \to f(u) \in \mathrm{dom}(g)$
\item $\varphi_{40} :\equiv \mathrm{ContIn}(g,B,f(c)) \to \forall x \in A\,(f(x) \in B) \to u \in A \to u \in \mathrm{dom}(g \circ f)$
\item $\varphi_{41} :\equiv \lvert x - c \rvert < \delta \to \forall x \in A\,(\lvert x - c \rvert < \delta \to \lvert f(x) - f(c) \rvert < d_{1})$
\item $\varphi_{42} :\equiv \mathrm{ContIn}(g,B,f(c)) \to \forall x \in A\,(f(x) \in B) \to \forall u \in A\,(u \in \mathrm{dom}(g \circ f))$
\item $\varphi_{43} :\equiv \delta \in \mathbb{R} \land 0 < \delta \land \forall x \in A\,(\lvert x - c \rvert < \delta \to \lvert f(x) - f(c) \rvert < d_{1})$
\item $\varphi_{44} :\equiv \forall y\,(A \subseteq y \leftrightarrow \forall u \in A\,(u \in y)) \to (A \subseteq \mathrm{dom}(g \circ f) \leftrightarrow \forall u \in A\,(u \in \mathrm{dom}(g \circ f)))$
\item $\varphi_{45} :\equiv \forall x\,(B \subseteq \mathrm{dom}(g) \to x \in B \to x \in \mathrm{dom}(g)) \to B \subseteq \mathrm{dom}(g) \to f(x) \in B \to f(x) \in \mathrm{dom}(g)$
\item $\varphi_{46} :\equiv \forall x\,(B \subseteq \mathrm{dom}(g) \to x \in B \to x \in \mathrm{dom}(g)) \to B \subseteq \mathrm{dom}(g) \to f(c) \in B \to f(c) \in \mathrm{dom}(g)$
\item $\varphi_{47} :\equiv \forall x\,(B \subseteq \mathrm{dom}(g) \to x \in B \to x \in \mathrm{dom}(g)) \to B \subseteq \mathrm{dom}(g) \to f(u) \in B \to f(u) \in \mathrm{dom}(g)$
\item $\varphi_{48} :\equiv \lvert g(f(x)) - g(f(c)) \rvert < \varepsilon \to \lvert g \circ f(x) - g(f(c)) \rvert < \varepsilon$
\item $\varphi_{49} :\equiv \lvert g \circ f(x) - g(f(c)) \rvert < \varepsilon \to \lvert g \circ f(x) - g \circ f(c) \rvert < \varepsilon$
\item $\varphi_{50} :\equiv \forall w\,(g(f(x)) = w \to \lvert g(f(x)) - g(f(c)) \rvert < \varepsilon \to \lvert w - g(f(c)) \rvert < \varepsilon)$
\item $\varphi_{51} :\equiv \neg (\delta \in \mathbb{R} \land 0 < \delta \land \forall x \in B\,(\lvert x - f(c) \rvert < \delta \to \lvert g(x) - g(f(c)) \rvert < \varepsilon))$
\item $\varphi_{52} :\equiv \forall w\,(g(f(c)) = w \to \lvert g \circ f(x) - g(f(c)) \rvert < \varepsilon \to \lvert g \circ f(x) - w \rvert < \varepsilon)$
\item $\varphi_{53} :\equiv \exists \delta\,(\delta \in \mathbb{R} \land 0 < \delta \land \forall x \in B\,(\lvert x - f(c) \rvert < \delta \to \lvert g(x) - g(f(c)) \rvert < \varepsilon))$
\item $\varphi_{54} :\equiv \forall y\,(g \circ f(x) = y \to y = g \circ f(x)) \to g \circ f(x) = g(f(x)) \to g(f(x)) = g \circ f(x)$
\item $\varphi_{55} :\equiv \forall y\,(g \circ f(c) = y \to y = g \circ f(c)) \to g \circ f(c) = g(f(c)) \to g(f(c)) = g \circ f(c)$
\item $\varphi_{56} :\equiv \forall \varepsilon \in \mathbb{R}\,(0 < \varepsilon \to \exists \delta\,(\delta \in \mathbb{R} \land 0 < \delta \land \forall x \in A\,(\lvert x - c \rvert < \delta \to \lvert f(x) - f(c) \rvert < \varepsilon)))$
\item $\varphi_{57} :\equiv \forall x \in A\,(\lvert x - c \rvert < \delta \to \lvert f(x) - f(c) \rvert < d_{1}) \to x \in A \to \lvert x - c \rvert < \delta \to \lvert f(x) - f(c) \rvert < d_{1}$
\item $\varphi_{58} :\equiv \mathrm{RFn}(g) \to \mathrm{RFn}(f) \to \mathrm{RFn}(g \circ f) \land \forall x \in \mathrm{dom}(g \circ f)\,(\varphi_{14})$
\item $\varphi_{59} :\equiv \forall x \in A\,(f(x) \in B) \to \mathrm{RFn}(g \circ f) \land \forall x \in \mathrm{dom}(g \circ f)\,(\varphi_{14})$
\item $\varphi_{60} :\equiv \mathrm{ContIn}(g,B,f(c)) \to \forall x \in A\,(f(x) \in B) \to x \in \mathrm{dom}(g \circ f) \to \varphi_{14}$
\item $\varphi_{61} :\equiv \mathrm{ContIn}(g,B,f(c)) \to \forall x \in A\,(f(x) \in B) \to c \in \mathrm{dom}(g \circ f) \to \varphi_{15}$
\item $\varphi_{62} :\equiv \mathrm{ContIn}(g,B,f(c)) \to \forall x \in A\,(f(x) \in B) \to \forall x \in \mathrm{dom}(g \circ f)\,(\varphi_{14})$
\item $\varphi_{63} :\equiv (\forall x \in A\,(f(x) \in B) \to x \in A \to f(x) \in B) \to \mathrm{ContIn}(f,A,c) \to \varphi_{25}$
\item $\varphi_{64} :\equiv (\forall x \in A\,(f(x) \in B) \to c \in A \to f(c) \in B) \to \mathrm{ContIn}(f,A,c) \to \varphi_{26}$
\item $\varphi_{65} :\equiv (\forall x \in A\,(f(x) \in B) \to u \in A \to f(u) \in B) \to \mathrm{ContIn}(f,A,c) \to \varphi_{27}$
\item $\varphi_{66} :\equiv \forall x\,(\mathrm{RFn}(g) \to \mathrm{RFn}(f) \to \varphi_{2}) \to \mathrm{RFn}(g) \to \mathrm{RFn}(f) \to \varphi_{2}$
\item $\varphi_{67} :\equiv \forall x\,(\mathrm{RFn}(g) \to \mathrm{RFn}(f) \to \varphi_{2}) \to \mathrm{RFn}(g) \to \mathrm{RFn}(f) \to \varphi_{3}$
\item $\varphi_{68} :\equiv \forall x\,(\mathrm{RFn}(g) \to \mathrm{RFn}(f) \to \varphi_{2}) \to \mathrm{RFn}(g) \to \mathrm{RFn}(f) \to \varphi_{4}$
\item $\varphi_{69} :\equiv \forall f\,\forall x\,(\mathrm{RFn}(g) \to \mathrm{RFn}(f) \to \varphi_{2}) \to \forall x\,(\mathrm{RFn}(g) \to \mathrm{RFn}(f) \to \varphi_{2})$
\item $\varphi_{70} :\equiv \forall x \in A\,(f(x) \in B) \to \varepsilon \in \mathbb{R} \to 0 < \varepsilon \to \exists \delta\,(\delta \in \mathbb{R} \land 0 < \delta \land \varphi_{34})$
\item $\varphi_{71} :\equiv \forall g\,\forall f\,\forall x\,(\mathrm{RFn}(g) \to \mathrm{RFn}(f) \to \varphi_{2}) \to \forall f\,\forall x\,(\mathrm{RFn}(g) \to \mathrm{RFn}(f) \to \varphi_{2})$
\item $\varphi_{72} :\equiv \varphi_{43} \to x \in A \to \lvert x - c \rvert < \delta \to \lvert g \circ f(x) - g(f(c)) \rvert < \varepsilon$
\item $\varphi_{73} :\equiv \varphi_{43} \to x \in A \to \lvert x - c \rvert < \delta \to \lvert g \circ f(x) - g \circ f(c) \rvert < \varepsilon$
\item $\varphi_{74} :\equiv \mathrm{ContIn}(g,B,f(c)) \to \forall x \in A\,(f(x) \in B) \to d_{1} \in \mathbb{R} \to 0 < d_{1} \to \exists \delta\,(\varphi_{43})$
\item $\varphi_{75} :\equiv \forall x \in A\,(f(x) \in B) \to \forall \varepsilon \in \mathbb{R}\,(0 < \varepsilon \to \exists \delta\,(\delta \in \mathbb{R} \land 0 < \delta \land \varphi_{34}))$
\item $\varphi_{76} :\equiv \mathrm{ContIn}(f,A,c) \leftrightarrow \mathrm{RFn}(f) \land A \subseteq \mathrm{dom}(f) \land c \in A \land \varphi_{56}$
\item $\varphi_{77} :\equiv \mathrm{ContIn}(g,B,f(c)) \to \forall x \in A\,(f(x) \in B) \to \varphi_{56}$
\item $\varphi_{78} :\equiv \forall c\,(\mathrm{ContIn}(g,B,c) \leftrightarrow \mathrm{RFn}(g) \land B \subseteq \mathrm{dom}(g) \land c \in B \land \forall \varepsilon \in \mathbb{R}\,(0 < \varepsilon \to \exists \delta\,(\delta \in \mathbb{R} \land 0 < \delta \land \forall x \in B\,(\lvert x - c \rvert < \delta \to \lvert g(x) - g(c) \rvert < \varepsilon))))$
\item $\varphi_{79} :\equiv (\mathrm{ContIn}(f,A,c) \to \varphi_{16}) \to \mathrm{ContIn}(f,A,c) \to \varphi_{23}$
\item $\varphi_{80} :\equiv \forall A\,\forall c\,(\mathrm{ContIn}(g,A,c) \leftrightarrow \mathrm{RFn}(g) \land A \subseteq \mathrm{dom}(g) \land c \in A \land \forall \varepsilon \in \mathbb{R}\,(0 < \varepsilon \to \exists \delta\,(\delta \in \mathbb{R} \land 0 < \delta \land \forall x \in A\,(\lvert x - c \rvert < \delta \to \lvert g(x) - g(c) \rvert < \varepsilon))))$
\item $\varphi_{81} :\equiv \mathrm{ContIn}(f,A,c) \to \mathrm{ContIn}(g,B,f(c)) \to \varphi_{59}$
\item $\varphi_{82} :\equiv \mathrm{ContIn}(g,B,f(c)) \to \forall x \in A\,(f(x) \in B) \to \varepsilon \in \mathbb{R} \to 0 < \varepsilon \to \varphi_{53}$
\item $\varphi_{83} :\equiv \mathrm{ContIn}(g,B,f(c)) \to \forall x \in A\,(f(x) \in B) \to \forall \varepsilon \in \mathbb{R}\,(0 < \varepsilon \to \varphi_{53})$
\item $\varphi_{84} :\equiv (\mathrm{ContIn}(f,A,c) \to \varphi_{27}) \to \mathrm{ContIn}(f,A,c) \to \varphi_{39}$
\item $\varphi_{85} :\equiv (\mathrm{ContIn}(f,A,c) \to \varphi_{42}) \to \mathrm{ContIn}(f,A,c) \to \varphi_{21}$
\item $\varphi_{86} :\equiv (\mathrm{ContIn}(f,A,c) \to \varphi_{39}) \to \mathrm{ContIn}(f,A,c) \to \varphi_{40}$
\item $\varphi_{87} :\equiv (\mathrm{ContIn}(f,A,c) \to \varphi_{40}) \to \varphi_{10} \to u \in A \to u \in \mathrm{dom}(g \circ f)$
\item $\varphi_{88} :\equiv \mathrm{ContIn}(g,B,f(c)) \leftrightarrow \mathrm{RFn}(g) \land B \subseteq \mathrm{dom}(g) \land f(c) \in B \land \forall \varepsilon \in \mathbb{R}\,(0 < \varepsilon \to \varphi_{53})$
\item $\varphi_{89} :\equiv \mathrm{ContIn}(f,A,c) \to \mathrm{ContIn}(g,B,f(c)) \to \varphi_{70}$
\item $\varphi_{90} :\equiv \mathrm{ContIn}(f,A,c) \to \mathrm{ContIn}(g,B,f(c)) \to \varphi_{75}$
\item $\varphi_{91} :\equiv \forall d_{1}\,(\varphi_{10} \land \varepsilon \in \mathbb{R} \land 0 < \varepsilon \land \neg \exists \delta\,(\delta \in \mathbb{R} \land 0 < \delta \land \varphi_{34}))$
\item $\varphi_{92} :\equiv \forall u\,(\varphi_{10} \to u \in A \to u \in \mathrm{dom}(g \circ f)) \to \forall u\,(\varphi_{10}) \to \forall u \in A\,(u \in \mathrm{dom}(g \circ f))$
\item $\varphi_{93} :\equiv \forall \varepsilon\,(\varphi_{10}) \to \forall \varepsilon \in \mathbb{R}\,(0 < \varepsilon \to \exists \delta\,(\delta \in \mathbb{R} \land 0 < \delta \land \varphi_{34}))$
\item $\varphi_{94} :\equiv \mathrm{ContIn}(g \circ f,A,c) \leftrightarrow \mathrm{RFn}(g \circ f) \land A \subseteq \mathrm{dom}(g \circ f) \land c \in A \land \forall \varepsilon \in \mathbb{R}\,(0 < \varepsilon \to \exists \delta\,(\delta \in \mathbb{R} \land 0 < \delta \land \varphi_{34}))$
\item $\varphi_{95} :\equiv \forall \delta\,(\neg (\delta \in \mathbb{R} \land 0 < \delta \land \varphi_{34})) \to \neg (\delta \in \mathbb{R} \land 0 < \delta \land \varphi_{34})$
\item $\varphi_{96} :\equiv \forall d_{1}\,(\neg (d_{1} \in \mathbb{R} \land 0 < d_{1} \land \varphi_{31})) \to \forall \delta\,\forall d_{1}\,(\neg (d_{1} \in \mathbb{R} \land 0 < d_{1} \land \varphi_{31}))$
\item $\varphi_{97} :\equiv 0 < \varepsilon \to d_{1} \in \mathbb{R} \land 0 < d_{1} \land \varphi_{31} \to \exists \delta\,(\delta \in \mathbb{R} \land 0 < \delta \land \varphi_{34})$
\item $\varphi_{98} :\equiv \varepsilon \in \mathbb{R} \to 0 < \varepsilon \to d_{1} \in \mathbb{R} \land 0 < d_{1} \land \varphi_{31} \to d_{1} \in \mathbb{R} \land 0 < d_{1} \land \varphi_{31}$
\item $\varphi_{99} :\equiv d_{1} \in \mathbb{R} \land 0 < d_{1} \land \varphi_{31} \to \varphi_{43} \to x \in A \to \lvert x - c \rvert < \delta \to x \in \mathrm{dom}(f)$
\item $\varphi_{100} :\equiv d_{1} \in \mathbb{R} \land 0 < d_{1} \land \varphi_{31} \to \varphi_{43} \to x \in A \to \lvert x - c \rvert < \delta \to c \in \mathrm{dom}(f)$
\item $\varphi_{101} :\equiv \varepsilon \in \mathbb{R} \to 0 < \varepsilon \to d_{1} \in \mathbb{R} \land 0 < d_{1} \land \varphi_{31} \to \varphi_{43} \to x \in A \to x \in A$
\item $\varphi_{102} :\equiv d_{1} \in \mathbb{R} \land 0 < d_{1} \land \varphi_{31} \to \varphi_{43} \to x \in A \to \lvert x - c \rvert < \delta \to f(x) \in B$
\item $\varphi_{103} :\equiv d_{1} \in \mathbb{R} \land 0 < d_{1} \land \varphi_{31} \to \varphi_{43} \to x \in A \to \lvert x - c \rvert < \delta \to f(c) \in B$
\item $\varphi_{104} :\equiv d_{1} \in \mathbb{R} \land 0 < d_{1} \land \varphi_{31} \to \varphi_{43} \to x \in A \to \lvert x - c \rvert < \delta \to \lvert x - c \rvert < \delta$
\item $\varphi_{105} :\equiv d_{1} \in \mathbb{R} \land 0 < d_{1} \land \varphi_{31} \to \varphi_{43} \to x \in A \to \lvert x - c \rvert < \delta \to f(x) \in \mathrm{dom}(g)$
\item $\varphi_{106} :\equiv d_{1} \in \mathbb{R} \land 0 < d_{1} \land \varphi_{31} \to \varphi_{43} \to x \in A \to \lvert x - c \rvert < \delta \to f(c) \in \mathrm{dom}(g)$
\item $\varphi_{107} :\equiv d_{1} \in \mathbb{R} \land 0 < d_{1} \land \varphi_{31} \to \varphi_{43} \to x \in A \to \lvert x - c \rvert < \delta \to x \in \mathrm{dom}(g \circ f)$
\item $\varphi_{108} :\equiv d_{1} \in \mathbb{R} \land 0 < d_{1} \land \varphi_{31} \to \varphi_{43} \to x \in A \to \lvert x - c \rvert < \delta \to c \in \mathrm{dom}(g \circ f)$
\item $\varphi_{109} :\equiv \forall x \in A\,(f(x) \in B) \to \varepsilon \in \mathbb{R} \to 0 < \varepsilon \to d_{1} \in \mathbb{R} \land 0 < d_{1} \land \varphi_{31} \to \exists \delta\,(\varphi_{43})$
\item $\varphi_{110} :\equiv (A \subseteq \mathrm{dom}(g \circ f) \leftrightarrow \forall u \in A\,(u \in \mathrm{dom}(g \circ f))) \to \varphi_{85}$
\item $\varphi_{111} :\equiv d_{1} \in \mathbb{R} \land 0 < d_{1} \land \varphi_{31} \to \varphi_{43} \to x \in A \to \lvert x - c \rvert < \delta \to \lvert f(x) - f(c) \rvert < d_{1}$
\item $\varphi_{112} :\equiv d_{1} \in \mathbb{R} \land 0 < d_{1} \land \varphi_{31} \to \varphi_{43} \to x \in A \to \lvert x - c \rvert < \delta \to g \circ f(x) = g(f(x))$
\item $\varphi_{113} :\equiv d_{1} \in \mathbb{R} \land 0 < d_{1} \land \varphi_{31} \to \varphi_{43} \to x \in A \to \lvert x - c \rvert < \delta \to g(f(x)) = g \circ f(x)$
\item $\varphi_{114} :\equiv d_{1} \in \mathbb{R} \land 0 < d_{1} \land \varphi_{31} \to \varphi_{43} \to x \in A \to \lvert x - c \rvert < \delta \to g \circ f(c) = g(f(c))$
\item $\varphi_{115} :\equiv d_{1} \in \mathbb{R} \land 0 < d_{1} \land \varphi_{31} \to \varphi_{43} \to x \in A \to \lvert x - c \rvert < \delta \to g(f(c)) = g \circ f(c)$
\item $\varphi_{116} :\equiv (\mathrm{ContIn}(f,A,c) \to \mathrm{ContIn}(f,A,c)) \to \varphi_{76} \to \mathrm{ContIn}(f,A,c) \to \varphi_{5}$
\item $\varphi_{117} :\equiv (\mathrm{ContIn}(f,A,c) \to \mathrm{ContIn}(f,A,c)) \to \varphi_{76} \to \mathrm{ContIn}(f,A,c) \to \varphi_{9}$
\item $\varphi_{118} :\equiv \mathrm{ContIn}(g,B,f(c)) \to \forall x \in A\,(f(x) \in B) \to \varphi_{98}$
\item $\varphi_{119} :\equiv (\mathrm{ContIn}(f,A,c) \to \mathrm{ContIn}(f,A,c)) \to \varphi_{76} \to \mathrm{ContIn}(f,A,c) \to \varphi_{12}$
\item $\varphi_{120} :\equiv \forall x\,(\varphi_{10} \land \varepsilon \in \mathbb{R} \land 0 < \varepsilon \land (d_{1} \in \mathbb{R} \land 0 < d_{1} \land \varphi_{31}) \land \varphi_{43})$
\item $\varphi_{121} :\equiv \mathrm{ContIn}(f,A,c) \to \mathrm{ContIn}(g,B,f(c)) \to \varphi_{109}$
\item $\varphi_{122} :\equiv \mathrm{ContIn}(g,B,f(c)) \to \forall x \in A\,(f(x) \in B) \to \varphi_{101}$
\item $\varphi_{123} :\equiv \mathrm{ContIn}(g,B,f(c)) \to \forall x \in A\,(f(x) \in B) \to \varepsilon \in \mathbb{R} \to \varphi_{97}$
\item $\varphi_{124} :\equiv (A \subseteq \mathrm{dom}(f) \to u \in A \to u \in \mathrm{dom}(f)) \to (\mathrm{ContIn}(f,A,c) \to \varphi_{12}) \to \varphi_{79}$
\item $\varphi_{125} :\equiv \varphi_{10} \land \varepsilon \in \mathbb{R} \land 0 < \varepsilon \land (d_{1} \in \mathbb{R} \land 0 < d_{1} \land \varphi_{31}) \land \neg \exists \delta\,(\delta \in \mathbb{R} \land 0 < \delta \land \varphi_{34})$
\item $\varphi_{126} :\equiv \varphi_{10} \land \varepsilon \in \mathbb{R} \land 0 < \varepsilon \land \neg \exists \delta\,(\delta \in \mathbb{R} \land 0 < \delta \land \varphi_{34}) \to \neg (d_{1} \in \mathbb{R} \land 0 < d_{1} \land \varphi_{31})$
\item $\varphi_{127} :\equiv 0 < \varepsilon \to d_{1} \in \mathbb{R} \land 0 < d_{1} \land \varphi_{31} \to \varphi_{43} \to x \in A \to \lvert x - c \rvert < \delta \to \varphi_{31}$
\item $\varphi_{128} :\equiv 0 < \varepsilon \to d_{1} \in \mathbb{R} \land 0 < d_{1} \land \varphi_{31} \to \varphi_{43} \to \exists \delta\,(\delta \in \mathbb{R} \land 0 < \delta \land \varphi_{34})$
\item $\varphi_{129} :\equiv \mathrm{ContIn}(g,B,f(c)) \to \forall x \in A\,(f(x) \in B) \to \varepsilon \in \mathbb{R} \to 0 < \varepsilon \to \varphi_{99}$
\item $\varphi_{130} :\equiv \mathrm{ContIn}(g,B,f(c)) \to \forall x \in A\,(f(x) \in B) \to \varepsilon \in \mathbb{R} \to 0 < \varepsilon \to \varphi_{100}$
\item $\varphi_{131} :\equiv \mathrm{ContIn}(g,B,f(c)) \to \forall x \in A\,(f(x) \in B) \to \varepsilon \in \mathbb{R} \to 0 < \varepsilon \to \varphi_{102}$
\item $\varphi_{132} :\equiv \mathrm{ContIn}(g,B,f(c)) \to \forall x \in A\,(f(x) \in B) \to \varepsilon \in \mathbb{R} \to 0 < \varepsilon \to \varphi_{103}$
\item $\varphi_{133} :\equiv \mathrm{ContIn}(g,B,f(c)) \to \forall x \in A\,(f(x) \in B) \to \varepsilon \in \mathbb{R} \to 0 < \varepsilon \to \varphi_{104}$
\item $\varphi_{134} :\equiv d_{1} \in \mathbb{R} \land 0 < d_{1} \land \varphi_{31} \to \varphi_{43} \to x \in A \to \lvert x - c \rvert < \delta \to f(x) \in B \to \varphi_{33}$
\item $\varphi_{135} :\equiv \mathrm{ContIn}(g,B,f(c)) \to \forall x \in A\,(f(x) \in B) \to \varepsilon \in \mathbb{R} \to 0 < \varepsilon \to \varphi_{105}$
\item $\varphi_{136} :\equiv \mathrm{ContIn}(g,B,f(c)) \to \forall x \in A\,(f(x) \in B) \to \varepsilon \in \mathbb{R} \to 0 < \varepsilon \to \varphi_{106}$
\item $\varphi_{137} :\equiv \mathrm{ContIn}(g,B,f(c)) \to \forall x \in A\,(f(x) \in B) \to \varepsilon \in \mathbb{R} \to 0 < \varepsilon \to \varphi_{107}$
\item $\varphi_{138} :\equiv \mathrm{ContIn}(g,B,f(c)) \to \forall x \in A\,(f(x) \in B) \to \varepsilon \in \mathbb{R} \to 0 < \varepsilon \to \varphi_{108}$
\item $\varphi_{139} :\equiv \varepsilon \in \mathbb{R} \to 0 < \varepsilon \to d_{1} \in \mathbb{R} \land 0 < d_{1} \land \varphi_{31} \to \varphi_{43} \to \delta \in \mathbb{R} \land 0 < \delta \land \varphi_{34}$
\item $\varphi_{140} :\equiv \forall x \in A\,(f(x) \in B) \to \varepsilon \in \mathbb{R} \to 0 < \varepsilon \to d_{1} \in \mathbb{R} \land 0 < d_{1} \land \varphi_{31} \to \varphi_{43} \to \varphi_{24}$
\item $\varphi_{141} :\equiv \forall x \in A\,(f(x) \in B) \to \varepsilon \in \mathbb{R} \to 0 < \varepsilon \to d_{1} \in \mathbb{R} \land 0 < d_{1} \land \varphi_{31} \to \varphi_{72}$
\item $\varphi_{142} :\equiv \forall x \in A\,(f(x) \in B) \to \varepsilon \in \mathbb{R} \to 0 < \varepsilon \to d_{1} \in \mathbb{R} \land 0 < d_{1} \land \varphi_{31} \to \varphi_{73}$
\item $\varphi_{143} :\equiv (B \subseteq \mathrm{dom}(g) \to f(u) \in B \to f(u) \in \mathrm{dom}(g)) \to (\mathrm{ContIn}(f,A,c) \to \varphi_{11}) \to \varphi_{84}$
\item $\varphi_{144} :\equiv \forall x \in A\,(f(x) \in B) \to \varepsilon \in \mathbb{R} \to 0 < \varepsilon \to d_{1} \in \mathbb{R} \land 0 < d_{1} \land \varphi_{31} \to \varphi_{43} \to \varphi_{34}$
\item $\varphi_{145} :\equiv (\mathrm{ContIn}(f,A,c) \to \varphi_{6}) \to (\mathrm{ContIn}(f,A,c) \to \varphi_{5}) \to \varphi_{81}$
\item $\varphi_{146} :\equiv \mathrm{ContIn}(g,B,f(c)) \to \forall x \in A\,(f(x) \in B) \to \varepsilon \in \mathbb{R} \to 0 < \varepsilon \to \varphi_{111}$
\item $\varphi_{147} :\equiv \mathrm{ContIn}(g,B,f(c)) \to \forall x \in A\,(f(x) \in B) \to \varepsilon \in \mathbb{R} \to 0 < \varepsilon \to \varphi_{112}$
\item $\varphi_{148} :\equiv \mathrm{ContIn}(g,B,f(c)) \to \forall x \in A\,(f(x) \in B) \to \varepsilon \in \mathbb{R} \to 0 < \varepsilon \to \varphi_{113}$
\item $\varphi_{149} :\equiv \mathrm{ContIn}(g,B,f(c)) \to \forall x \in A\,(f(x) \in B) \to \varepsilon \in \mathbb{R} \to 0 < \varepsilon \to \varphi_{114}$
\item $\varphi_{150} :\equiv \mathrm{ContIn}(g,B,f(c)) \to \forall x \in A\,(f(x) \in B) \to \varepsilon \in \mathbb{R} \to 0 < \varepsilon \to \varphi_{115}$
\item $\varphi_{151} :\equiv \forall x \in A\,(f(x) \in B) \to \varepsilon \in \mathbb{R} \to 0 < \varepsilon \to d_{1} \in \mathbb{R} \land 0 < d_{1} \land \varphi_{31} \to \varphi_{43} \to \varphi_{43}$
\item $\varphi_{152} :\equiv (\mathrm{ContIn}(f,A,c) \to \varphi_{1}) \to \varphi_{88} \to \mathrm{ContIn}(f,A,c) \to \varphi_{6}$
\item $\varphi_{153} :\equiv \forall x \in A\,(f(x) \in B) \to \varepsilon \in \mathbb{R} \to 0 < \varepsilon \to d_{1} \in \mathbb{R} \land 0 < d_{1} \land \varphi_{31} \to \varphi_{43} \to x \in A \to \varphi_{35}$
\item $\varphi_{154} :\equiv \forall x \in A\,(f(x) \in B) \to \varepsilon \in \mathbb{R} \to 0 < \varepsilon \to d_{1} \in \mathbb{R} \land 0 < d_{1} \land \varphi_{31} \to \varphi_{43} \to x \in A \to \varphi_{41}$
\item $\varphi_{155} :\equiv (\mathrm{ContIn}(f,A,c) \to \varphi_{1}) \to \varphi_{88} \to \mathrm{ContIn}(f,A,c) \to \varphi_{11}$
\item $\varphi_{156} :\equiv (\forall x \in \mathrm{dom}(g \circ f)\,(\varphi_{14}) \to x \in \mathrm{dom}(g \circ f) \to \varphi_{14}) \to \mathrm{ContIn}(f,A,c) \to \varphi_{60}$
\item $\varphi_{157} :\equiv (\forall x \in \mathrm{dom}(g \circ f)\,(\varphi_{14}) \to c \in \mathrm{dom}(g \circ f) \to \varphi_{15}) \to \mathrm{ContIn}(f,A,c) \to \varphi_{61}$
\item $\varphi_{158} :\equiv \mathrm{ContIn}(f,A,c) \to \mathrm{ContIn}(g,B,f(c)) \to \varphi_{140}$
\item $\varphi_{159} :\equiv \mathrm{ContIn}(f,A,c) \to \mathrm{ContIn}(g,B,f(c)) \to \varphi_{141}$
\item $\varphi_{160} :\equiv \mathrm{ContIn}(f,A,c) \to \mathrm{ContIn}(g,B,f(c)) \to \varphi_{142}$
\item $\varphi_{161} :\equiv \varphi_{10} \land \varepsilon \in \mathbb{R} \land 0 < \varepsilon \land (d_{1} \in \mathbb{R} \land 0 < d_{1} \land \varphi_{31}) \land \varphi_{43} \to x \in A \to \lvert x - c \rvert < \delta \to \lvert g \circ f(x) - g \circ f(c) \rvert < \varepsilon$
\item $\varphi_{162} :\equiv (\mathrm{ContIn}(f,A,c) \to \varphi_{5}) \to (\mathrm{ContIn}(f,A,c) \to \varphi_{23}) \to \varphi_{86}$
\item $\varphi_{163} :\equiv \mathrm{ContIn}(f,A,c) \to \mathrm{ContIn}(g,B,f(c)) \to \varphi_{144}$
\item $\varphi_{164} :\equiv \mathrm{ContIn}(f,A,c) \to \mathrm{ContIn}(g,B,f(c)) \to \varphi_{151}$
\item $\varphi_{165} :\equiv (\mathrm{ContIn}(f,A,c) \to \varphi_{9}) \to \varphi_{90} \to \mathrm{ContIn}(f,A,c) \to \varphi_{22}$
\item $\varphi_{166} :\equiv \mathrm{ContIn}(g,B,f(c)) \to \forall x \in A\,(f(x) \in B) \to \varphi_{139}$
\item $\varphi_{167} :\equiv \mathrm{ContIn}(f,A,c) \to \mathrm{ContIn}(g,B,f(c)) \to \varphi_{153}$
\item $\varphi_{168} :\equiv \mathrm{ContIn}(g,B,f(c)) \to \forall x \in A\,(f(x) \in B) \to \varepsilon \in \mathbb{R} \to \varphi_{127}$
\item $\varphi_{169} :\equiv \mathrm{ContIn}(g,B,f(c)) \to \forall x \in A\,(f(x) \in B) \to \varepsilon \in \mathbb{R} \to \varphi_{128}$
\item $\varphi_{170} :\equiv \mathrm{ContIn}(f,A,c) \to \mathrm{ContIn}(g,B,f(c)) \to \varphi_{154}$
\item $\varphi_{171} :\equiv (\mathrm{ContIn}(f,A,c) \to \mathrm{ContIn}(f,A,c)) \to \varphi_{76} \to \mathrm{ContIn}(f,A,c) \to \varphi_{77}$
\item $\varphi_{172} :\equiv \mathrm{ContIn}(g,B,f(c)) \to \forall x \in A\,(f(x) \in B) \to \varepsilon \in \mathbb{R} \to 0 < \varepsilon \to \varphi_{134}$
\item $\varphi_{173} :\equiv (\mathrm{ContIn}(f,A,c) \to \varphi_{9}) \to \mathrm{ContIn}(f,A,c) \to \varphi_{130}$
\item $\varphi_{174} :\equiv (\mathrm{ContIn}(f,A,c) \to \varphi_{9}) \to \mathrm{ContIn}(f,A,c) \to \varphi_{132}$
\item $\varphi_{175} :\equiv (g \circ f(x) = g(f(x)) \to g(f(x)) = g \circ f(x)) \to \mathrm{ContIn}(f,A,c) \to \varphi_{148}$
\item $\varphi_{176} :\equiv (g \circ f(c) = g(f(c)) \to g(f(c)) = g \circ f(c)) \to \mathrm{ContIn}(f,A,c) \to \varphi_{150}$
\item $\varphi_{177} :\equiv \forall \varepsilon\,(\varphi_{10} \to \varepsilon \in \mathbb{R} \to 0 < \varepsilon \to \exists \delta\,(\delta \in \mathbb{R} \land 0 < \delta \land \varphi_{34})) \to \varphi_{93}$
\item $\varphi_{178} :\equiv (\varphi_{56} \to d_{1} \in \mathbb{R} \to 0 < d_{1} \to \exists \delta\,(\varphi_{43})) \to \mathrm{ContIn}(f,A,c) \to \varphi_{74}$
\item $\varphi_{179} :\equiv \forall \delta\,(\forall d_{1}\,(\neg (d_{1} \in \mathbb{R} \land 0 < d_{1} \land \varphi_{31})) \to \varphi_{51}) \to \forall \delta\,\forall d_{1}\,(\neg (d_{1} \in \mathbb{R} \land 0 < d_{1} \land \varphi_{31})) \to \forall \delta\,(\varphi_{51})$
\item $\varphi_{180} :\equiv (\mathrm{ContIn}(f,A,c) \to \varphi_{1}) \to \varphi_{88} \to \mathrm{ContIn}(f,A,c) \to \varphi_{83}$
\item $\varphi_{181} :\equiv (\forall \varepsilon \in \mathbb{R}\,(0 < \varepsilon \to \varphi_{53}) \to \varepsilon \in \mathbb{R} \to 0 < \varepsilon \to \varphi_{53}) \to \mathrm{ContIn}(f,A,c) \to \varphi_{82}$
\item $\varphi_{182} :\equiv (\mathrm{RFn}(g) \to \mathrm{RFn}(f) \to \varphi_{4}) \to (\mathrm{ContIn}(f,A,c) \to \varphi_{6}) \to \varphi_{162}$
\item $\varphi_{183} :\equiv (A \subseteq \mathrm{dom}(f) \to c \in A \to c \in \mathrm{dom}(f)) \to (\mathrm{ContIn}(f,A,c) \to \varphi_{12}) \to \varphi_{173}$
\item $\varphi_{184} :\equiv (\mathrm{ContIn}(f,A,c) \to \varphi_{13}) \to (\mathrm{ContIn}(f,A,c) \to \varphi_{21}) \to \varphi_{165}$
\item $\varphi_{185} :\equiv \forall u\,(\varphi_{10} \to u \in A \to u \in \mathrm{dom}(g \circ f)) \to \varphi_{92} \to (\varphi_{10} \to \forall u\,(\varphi_{10})) \to \mathrm{ContIn}(f,A,c) \to \varphi_{42}$
\item $\varphi_{186} :\equiv (\mathrm{ContIn}(f,A,c) \to \varphi_{122}) \to \mathrm{ContIn}(f,A,c) \to \varphi_{129}$
\item $\varphi_{187} :\equiv (\mathrm{ContIn}(f,A,c) \to \varphi_{122}) \to \mathrm{ContIn}(f,A,c) \to \varphi_{131}$
\item $\varphi_{188} :\equiv (\mathrm{ContIn}(f,A,c) \to \varphi_{131}) \to \mathrm{ContIn}(f,A,c) \to \varphi_{135}$
\item $\varphi_{189} :\equiv (\mathrm{ContIn}(f,A,c) \to \varphi_{132}) \to \mathrm{ContIn}(f,A,c) \to \varphi_{136}$
\item $\varphi_{190} :\equiv (\mathrm{ContIn}(f,A,c) \to \varphi_{135}) \to \mathrm{ContIn}(f,A,c) \to \varphi_{137}$
\item $\varphi_{191} :\equiv (\mathrm{ContIn}(f,A,c) \to \varphi_{136}) \to \mathrm{ContIn}(f,A,c) \to \varphi_{138}$
\item $\varphi_{192} :\equiv (\mathrm{ContIn}(f,A,c) \to \varphi_{133}) \to \mathrm{ContIn}(f,A,c) \to \varphi_{146}$
\item $\varphi_{193} :\equiv (\mathrm{ContIn}(f,A,c) \to \varphi_{137}) \to \mathrm{ContIn}(f,A,c) \to \varphi_{147}$
\item $\varphi_{194} :\equiv (\mathrm{ContIn}(f,A,c) \to \varphi_{138}) \to \mathrm{ContIn}(f,A,c) \to \varphi_{149}$
\item $\varphi_{195} :\equiv (\mathrm{ContIn}(f,A,c) \to \varphi_{118}) \to \mathrm{ContIn}(f,A,c) \to \varphi_{168}$
\item $\varphi_{196} :\equiv (A \subseteq \mathrm{dom}(f) \to x \in A \to x \in \mathrm{dom}(f)) \to (\mathrm{ContIn}(f,A,c) \to \varphi_{12}) \to \varphi_{186}$
\item $\varphi_{197} :\equiv (\mathrm{ContIn}(f,A,c) \to \varphi_{74}) \to (\mathrm{ContIn}(f,A,c) \to \varphi_{118}) \to \varphi_{121}$
\item $\varphi_{198} :\equiv (B \subseteq \mathrm{dom}(g) \to f(x) \in B \to f(x) \in \mathrm{dom}(g)) \to (\mathrm{ContIn}(f,A,c) \to \varphi_{11}) \to \varphi_{188}$
\item $\varphi_{199} :\equiv (B \subseteq \mathrm{dom}(g) \to f(c) \in B \to f(c) \in \mathrm{dom}(g)) \to (\mathrm{ContIn}(f,A,c) \to \varphi_{11}) \to \varphi_{189}$
\item $\varphi_{200} :\equiv (\varphi_{10} \land \varepsilon \in \mathbb{R} \land 0 < \varepsilon \land (d_{1} \in \mathbb{R} \land 0 < d_{1} \land \varphi_{31}) \land \varphi_{43} \to \varphi_{120}) \to \varphi_{163}$
\item $\varphi_{201} :\equiv \forall \varepsilon\,(\varphi_{10} \to \varepsilon \in \mathbb{R} \to 0 < \varepsilon \to \exists \delta\,(\delta \in \mathbb{R} \land 0 < \delta \land \varphi_{34})) \to \varphi_{177} \to (\varphi_{10} \to \forall \varepsilon\,(\varphi_{10})) \to \varphi_{90}$
\item $\varphi_{202} :\equiv (\mathrm{ContIn}(f,A,c) \to \varphi_{168}) \to (\varphi_{31} \to f(x) \in B \to \varphi_{33}) \to \mathrm{ContIn}(f,A,c) \to \varphi_{172}$
\item $\varphi_{203} :\equiv (\mathrm{ContIn}(f,A,c) \to \varphi_{166}) \to \varphi_{95} \to \mathrm{ContIn}(f,A,c) \to \varphi_{169}$
\item $\varphi_{204} :\equiv (\mathrm{ContIn}(f,A,c) \to \varphi_{122}) \to (\mathrm{ContIn}(f,A,c) \to \varphi_{146}) \to \varphi_{158}$
\item $\varphi_{205} :\equiv \forall \delta\,(\forall d_{1}\,(\neg (d_{1} \in \mathbb{R} \land 0 < d_{1} \land \varphi_{31})) \to \varphi_{51}) \to \varphi_{179} \to \varphi_{96} \to \varphi_{89}$
\item $\varphi_{206} :\equiv (\mathrm{ContIn}(f,A,c) \to \varphi_{5}) \to (\mathrm{ContIn}(f,A,c) \to \varphi_{129}) \to \varphi_{190}$
\item $\varphi_{207} :\equiv (\mathrm{ContIn}(f,A,c) \to \varphi_{5}) \to (\mathrm{ContIn}(f,A,c) \to \varphi_{130}) \to \varphi_{191}$
\item $\varphi_{208} :\equiv (\mathrm{ContIn}(f,A,c) \to \varphi_{148}) \to \varphi_{158} \to (g(f(x)) = g \circ f(x) \to \varphi_{48}) \to \varphi_{159}$
\item $\varphi_{209} :\equiv (\mathrm{ContIn}(f,A,c) \to \varphi_{150}) \to \varphi_{159} \to (g(f(c)) = g \circ f(c) \to \varphi_{49}) \to \varphi_{160}$
\item $\varphi_{210} :\equiv (\mathrm{RFn}(g) \to \mathrm{RFn}(f) \to \varphi_{2}) \to (\mathrm{ContIn}(f,A,c) \to \varphi_{6}) \to \varphi_{206}$
\item $\varphi_{211} :\equiv (\mathrm{RFn}(g) \to \mathrm{RFn}(f) \to \varphi_{3}) \to (\mathrm{ContIn}(f,A,c) \to \varphi_{6}) \to \varphi_{207}$
\item $\varphi_{212} :\equiv (\varphi_{10} \land \varepsilon \in \mathbb{R} \land 0 < \varepsilon \land \neg \exists \delta\,(\delta \in \mathbb{R} \land 0 < \delta \land \varphi_{34}) \to \varphi_{91}) \to \varphi_{205}$
\item $\varphi_{213} :\equiv (\mathrm{ContIn}(f,A,c) \to \varphi_{172}) \to (\mathrm{ContIn}(f,A,c) \to \varphi_{25}) \to \varphi_{204}$
\item $\varphi_{214} :\equiv (\forall \delta\,(\varphi_{125} \to \neg \varphi_{43}) \to \forall \delta\,(\varphi_{125}) \to \forall \delta\,(\neg \varphi_{43})) \to (\varphi_{125} \to \forall \delta\,(\varphi_{125})) \to \mathrm{ContIn}(f,A,c) \to \varphi_{123}$
\item $\varphi_{215} :\equiv (\forall d_{1}\,(\varphi_{126}) \to \varphi_{91} \to \forall d_{1}\,(\neg (d_{1} \in \mathbb{R} \land 0 < d_{1} \land \varphi_{31}))) \to \varphi_{212}$
\end{abbrevlist}
\end{fullproof}

% 自動生成 (prover/build.js): ContOnAdd  ノード数 161
\begin{fullproof}{定理2.9(和)}{$\forall f\,\forall g\,\forall A\,(\mathrm{ContOn}(f,A) \to \mathrm{ContOn}(g,A) \to \mathrm{ContOn}((f + g),A))$}{161}\phantomsection\label{tree:ContOnAdd}
\noindent$\Pi_{1}$:
\begin{ptree}
\AXC{$\forall f\,\forall A\,(\varphi_{33})$\lbl{DefContOn}}
\AXC{$\forall f\,\forall A\,(\varphi_{33}) \to \forall A\,(\varphi_{32})$\lbl{A2}}
\BIC{$\forall A\,(\varphi_{32})$}
\AXC{$\forall A\,(\varphi_{32}) \to \varphi_{32}$\lbl{A2}}
\BIC{$\varphi_{32}$}
\end{ptree}
\noindent$\Pi_{2}$:
\begin{ptree}
\AXC{$\varphi_{20}$\lbl{A2}}
\AXC{$\stackrel{\Pi_{1}}{\vdots}$}\noLine
\UIC{$\varphi_{32}$}
\AXC{$\varphi_{2}$\lbl{A1}}
\AXC{$\varphi_{2} \to \varphi_{32} \to \mathrm{ContOn}(f,A) \to \varphi_{16}$\lbl{A1}}
\BIC{$\varphi_{32} \to \mathrm{ContOn}(f,A) \to \varphi_{16}$}
\BIC{$\mathrm{ContOn}(f,A) \to \varphi_{16}$}
\AXC{$\varphi_{51}$\lbl{A1}}
\BIC{$\varphi_{20} \to \mathrm{ContOn}(f,A) \to \varphi_{10}$}
\BIC{$\mathrm{ContOn}(f,A) \to \varphi_{10}$}
\end{ptree}
\noindent$\Pi_{3}$:
\begin{ptree}
\AXC{$\forall f\,\forall A\,(\varphi_{33})$\lbl{DefContOn}}
\AXC{$\forall f\,\forall A\,(\varphi_{33}) \to \forall A\,(\varphi_{33})$\lbl{A2}}
\BIC{$\forall A\,(\varphi_{33})$}
\AXC{$\forall A\,(\varphi_{33}) \to \varphi_{33}$\lbl{A2}}
\BIC{$\varphi_{33}$}
\end{ptree}
\noindent$\Pi_{4}$:
\begin{ptree}
\AXC{$\varphi_{21}$\lbl{A2}}
\AXC{$\stackrel{\Pi_{3}}{\vdots}$}\noLine
\UIC{$\varphi_{33}$}
\AXC{$\mathrm{ContOn}(f,A) \to \mathrm{ContOn}(f,A)$\lbl{A1}}
\AXC{$\varphi_{50}$\lbl{A1}}
\BIC{$\varphi_{33} \to \mathrm{ContOn}(f,A) \to \varphi_{17}$}
\BIC{$\mathrm{ContOn}(f,A) \to \varphi_{17}$}
\AXC{$\varphi_{52}$\lbl{A1}}
\BIC{$\varphi_{21} \to \mathrm{ContOn}(f,A) \to \varphi_{11}$}
\BIC{$\mathrm{ContOn}(f,A) \to \varphi_{11}$}
\end{ptree}
\noindent$\Pi_{5}$:
\begin{ptree}
\AXC{$\forall f\,\forall g\,\forall A\,\forall c\,(\mathrm{ContIn}(f,A,c) \to \varphi_{1})$\lbl{定理2.3(和)}}
\AXC{$\varphi_{48}$\lbl{A2}}
\BIC{$\forall g\,\forall A\,\forall c\,(\mathrm{ContIn}(f,A,c) \to \varphi_{1})$}
\AXC{$\varphi_{47}$\lbl{A2}}
\BIC{$\forall A\,\forall c\,(\mathrm{ContIn}(f,A,c) \to \varphi_{1})$}
\AXC{$\varphi_{46}$\lbl{A2}}
\BIC{$\forall c\,(\mathrm{ContIn}(f,A,c) \to \varphi_{1})$}
\AXC{$\varphi_{45}$\lbl{A2}}
\BIC{$\mathrm{ContIn}(f,A,c) \to \varphi_{1}$}
\AXC{$\varphi_{55}$\lbl{A1}}
\BIC{$\varphi_{19} \to \varphi_{18} \to \mathrm{ContOn}(f,A) \to \varphi_{8}$}
\end{ptree}
\noindent$\Pi_{6}$:
\begin{ptree}
\AXC{$\varphi_{6}$\lbl{A1}}
\AXC{$\stackrel{\Pi_{4}}{\vdots}$}\noLine
\UIC{$\mathrm{ContOn}(f,A) \to \varphi_{11}$}
\AXC{$(\mathrm{ContOn}(f,A) \to \varphi_{11}) \to \varphi_{6} \to \varphi_{19}$\lbl{A1}}
\BIC{$\varphi_{6} \to \varphi_{19}$}
\BIC{$\varphi_{19}$}
\AXC{$\stackrel{\Pi_{5}}{\vdots}$}\noLine
\UIC{$\varphi_{19} \to \varphi_{18} \to \mathrm{ContOn}(f,A) \to \varphi_{8}$}
\BIC{$\varphi_{18} \to \mathrm{ContOn}(f,A) \to \varphi_{8}$}
\end{ptree}
\noindent$\Pi_{7}$:
\begin{ptree}
\AXC{$\varphi_{6}$\lbl{A1}}
\AXC{$\stackrel{\Pi_{2}}{\vdots}$}\noLine
\UIC{$\mathrm{ContOn}(f,A) \to \varphi_{10}$}
\AXC{$(\mathrm{ContOn}(f,A) \to \varphi_{10}) \to \varphi_{6} \to \varphi_{18}$\lbl{A1}}
\BIC{$\varphi_{6} \to \varphi_{18}$}
\BIC{$\varphi_{18}$}
\AXC{$\stackrel{\Pi_{6}}{\vdots}$}\noLine
\UIC{$\varphi_{18} \to \mathrm{ContOn}(f,A) \to \varphi_{8}$}
\BIC{$\mathrm{ContOn}(f,A) \to \varphi_{8}$}
\end{ptree}
\noindent$\Pi_{8}$:
\begin{ptree}
\AXC{$\varphi_{23}$\lbl{A4}}
\AXC{$\forall c\,(\varphi_{28}) \to \varphi_{34}$\lbl{A3}}
\AXC{$\stackrel{\Pi_{7}}{\vdots}$}\noLine
\UIC{$\mathrm{ContOn}(f,A) \to \varphi_{8}$}
\AXC{$(\mathrm{ContOn}(f,A) \to \varphi_{8}) \to \varphi_{28}$\lbl{A1}}
\BIC{$\varphi_{28}$}
\UIC{$\forall c\,(\varphi_{28})$}
\AXC{$\varphi_{60}$\lbl{A1}}
\BIC{$(\forall c\,(\varphi_{28}) \to \varphi_{34}) \to \varphi_{23} \to \mathrm{ContOn}(f,A) \to \varphi_{12}$}
\BIC{$\varphi_{23} \to \mathrm{ContOn}(f,A) \to \varphi_{12}$}
\BIC{$\mathrm{ContOn}(f,A) \to \varphi_{12}$}
\end{ptree}
\noindent$\Pi_{9}$:
\begin{ptree}
\AXC{$\forall f\,\forall A\,(\varphi_{33})$\lbl{DefContOn}}
\AXC{$\forall f\,\forall A\,(\varphi_{33}) \to \forall A\,(\varphi_{32})$\lbl{A2}}
\BIC{$\forall A\,(\varphi_{32})$}
\AXC{$\forall A\,(\varphi_{32}) \to \varphi_{32}$\lbl{A2}}
\BIC{$\varphi_{32}$}
\AXC{$\varphi_{2}$\lbl{A1}}
\AXC{$\varphi_{2} \to \varphi_{42}$\lbl{A1}}
\BIC{$\varphi_{42}$}
\BIC{$\mathrm{ContOn}(f,A) \to \mathrm{ContOn}(g,A) \to \mathrm{RFn}(g)$}
\end{ptree}
\noindent$\Pi_{10}$:
\begin{ptree}
\AXC{$\mathrm{ContOn}(f,A) \to \mathrm{ContOn}(f,A)$\lbl{A1}}
\AXC{$(\mathrm{ContOn}(f,A) \to \mathrm{ContOn}(f,A)) \to \varphi_{43}$\lbl{A1}}
\BIC{$\varphi_{43}$}
\end{ptree}
\noindent$\Pi_{11}$:
\begin{ptree}
\AXC{$\forall f\,\forall g\,(\mathrm{RFn}(f) \to \mathrm{RFn}(g) \to \varphi_{49})$\lbl{[S17]}}
\AXC{$\varphi_{61}$\lbl{A2}}
\BIC{$\forall g\,(\mathrm{RFn}(f) \to \mathrm{RFn}(g) \to \varphi_{49})$}
\AXC{$\varphi_{59}$\lbl{A2}}
\BIC{$\mathrm{RFn}(f) \to \mathrm{RFn}(g) \to \varphi_{49}$}
\AXC{$(\mathrm{RFn}(f) \to \mathrm{RFn}(g) \to \varphi_{49}) \to \varphi_{56}$\lbl{A1}}
\BIC{$\varphi_{56}$}
\end{ptree}
\noindent$\Pi_{12}$:
\begin{ptree}
\AXC{$\forall f\,\forall A\,(\varphi_{33})$\lbl{DefContOn}}
\AXC{$\forall f\,\forall A\,(\varphi_{33}) \to \forall A\,(\varphi_{33})$\lbl{A2}}
\BIC{$\forall A\,(\varphi_{33})$}
\AXC{$\forall A\,(\varphi_{33}) \to \varphi_{33}$\lbl{A2}}
\BIC{$\varphi_{33}$}
\AXC{$\stackrel{\Pi_{10}}{\vdots}$}\noLine
\UIC{$\varphi_{43}$}
\BIC{$\mathrm{ContOn}(f,A) \to \mathrm{ContOn}(g,A) \to \mathrm{RFn}(f)$}
\AXC{$\stackrel{\Pi_{11}}{\vdots}$}\noLine
\UIC{$\varphi_{56}$}
\BIC{$\varphi_{53}$}
\end{ptree}
\noindent$\Pi_{13}$:
\begin{ptree}
\AXC{$\stackrel{\Pi_{9}}{\vdots}$}\noLine
\UIC{$\mathrm{ContOn}(f,A) \to \mathrm{ContOn}(g,A) \to \mathrm{RFn}(g)$}
\AXC{$\stackrel{\Pi_{12}}{\vdots}$}\noLine
\UIC{$\varphi_{53}$}
\BIC{$\mathrm{ContOn}(f,A) \to \mathrm{ContOn}(g,A) \to \varphi_{49}$}
\end{ptree}
\noindent$\Pi_{14}$:
\begin{ptree}
\AXC{$\varphi_{44}$\lbl{A2}}
\AXC{$\stackrel{\Pi_{13}}{\vdots}$}\noLine
\UIC{$\mathrm{ContOn}(f,A) \to \mathrm{ContOn}(g,A) \to \varphi_{49}$}
\AXC{$\varphi_{57}$\lbl{A1}}
\BIC{$\mathrm{ContOn}(f,A) \to \mathrm{ContOn}(g,A) \to u \in A \to \varphi_{22}$}
\AXC{$\varphi_{58}$\lbl{A1}}
\BIC{$\varphi_{44} \to \mathrm{ContOn}(f,A) \to \mathrm{ContOn}(g,A) \to \varphi_{29}$}
\BIC{$\mathrm{ContOn}(f,A) \to \mathrm{ContOn}(g,A) \to \varphi_{29}$}
\end{ptree}
\noindent$\Pi_{15}$:
\begin{ptree}
\AXC{$\forall f\,\forall A\,(\varphi_{33})$\lbl{DefContOn}}
\AXC{$\forall f\,\forall A\,(\varphi_{33}) \to \forall A\,(\varphi_{32})$\lbl{A2}}
\BIC{$\forall A\,(\varphi_{32})$}
\AXC{$\forall A\,(\varphi_{32}) \to \varphi_{32}$\lbl{A2}}
\BIC{$\varphi_{32}$}
\AXC{$\varphi_{2}$\lbl{A1}}
\AXC{$\varphi_{2} \to \varphi_{32} \to \varphi_{3}$\lbl{A1}}
\BIC{$\varphi_{32} \to \varphi_{3}$}
\BIC{$\varphi_{3}$}
\end{ptree}
\noindent$\Pi_{16}$:
\begin{ptree}
\AXC{$\forall X\,\forall Y\,\forall x\,(X \subseteq Y \to x \in X \to x \in Y)$\lbl{[B5]}}
\AXC{$\varphi_{35}$\lbl{A2}}
\BIC{$\forall Y\,\forall x\,(A \subseteq Y \to x \in A \to x \in Y)$}
\AXC{$\varphi_{37}$\lbl{A2}}
\BIC{$\forall x\,(A \subseteq \mathrm{dom}(g) \to x \in A \to x \in \mathrm{dom}(g))$}
\AXC{$\varphi_{39}$\lbl{A2}}
\BIC{$A \subseteq \mathrm{dom}(g) \to u \in A \to u \in \mathrm{dom}(g)$}
\end{ptree}
\noindent$\Pi_{17}$:
\begin{ptree}
\AXC{$\varphi_{7}$\lbl{A1}}
\AXC{$\stackrel{\Pi_{15}}{\vdots}$}\noLine
\UIC{$\varphi_{3}$}
\AXC{$\stackrel{\Pi_{16}}{\vdots}$}\noLine
\UIC{$A \subseteq \mathrm{dom}(g) \to u \in A \to u \in \mathrm{dom}(g)$}
\AXC{$(A \subseteq \mathrm{dom}(g) \to u \in A \to u \in \mathrm{dom}(g)) \to \varphi_{3} \to \varphi_{7} \to \varphi_{14}$\lbl{A1}}
\BIC{$\varphi_{3} \to \varphi_{7} \to \varphi_{14}$}
\BIC{$\varphi_{7} \to \varphi_{14}$}
\BIC{$\varphi_{14}$}
\end{ptree}
\noindent$\Pi_{18}$:
\begin{ptree}
\AXC{$\mathrm{ContOn}(f,A) \to \mathrm{ContOn}(f,A)$\lbl{A1}}
\AXC{$(\mathrm{ContOn}(f,A) \to \mathrm{ContOn}(f,A)) \to \varphi_{33} \to \varphi_{4}$\lbl{A1}}
\BIC{$\varphi_{33} \to \varphi_{4}$}
\end{ptree}
\noindent$\Pi_{19}$:
\begin{ptree}
\AXC{$\forall X\,\forall Y\,\forall x\,(X \subseteq Y \to x \in X \to x \in Y)$\lbl{[B5]}}
\AXC{$\varphi_{35}$\lbl{A2}}
\BIC{$\forall Y\,\forall x\,(A \subseteq Y \to x \in A \to x \in Y)$}
\AXC{$\varphi_{38}$\lbl{A2}}
\BIC{$\forall x\,(A \subseteq \mathrm{dom}(f) \to x \in A \to x \in \mathrm{dom}(f))$}
\AXC{$\varphi_{40}$\lbl{A2}}
\BIC{$A \subseteq \mathrm{dom}(f) \to u \in A \to u \in \mathrm{dom}(f)$}
\end{ptree}
\noindent$\Pi_{20}$:
\begin{ptree}
\AXC{$\stackrel{\Pi_{19}}{\vdots}$}\noLine
\UIC{$A \subseteq \mathrm{dom}(f) \to u \in A \to u \in \mathrm{dom}(f)$}
\AXC{$(A \subseteq \mathrm{dom}(f) \to u \in A \to u \in \mathrm{dom}(f)) \to \varphi_{4} \to \varphi_{7} \to \varphi_{15}$\lbl{A1}}
\BIC{$\varphi_{4} \to \varphi_{7} \to \varphi_{15}$}
\end{ptree}
\noindent$\Pi_{21}$:
\begin{ptree}
\AXC{$\varphi_{7}$\lbl{A1}}
\AXC{$\forall f\,\forall A\,(\varphi_{33})$\lbl{DefContOn}}
\AXC{$\forall f\,\forall A\,(\varphi_{33}) \to \forall A\,(\varphi_{33})$\lbl{A2}}
\BIC{$\forall A\,(\varphi_{33})$}
\AXC{$\forall A\,(\varphi_{33}) \to \varphi_{33}$\lbl{A2}}
\BIC{$\varphi_{33}$}
\AXC{$\stackrel{\Pi_{18}}{\vdots}$}\noLine
\UIC{$\varphi_{33} \to \varphi_{4}$}
\BIC{$\varphi_{4}$}
\AXC{$\stackrel{\Pi_{20}}{\vdots}$}\noLine
\UIC{$\varphi_{4} \to \varphi_{7} \to \varphi_{15}$}
\BIC{$\varphi_{7} \to \varphi_{15}$}
\BIC{$\varphi_{15}$}
\end{ptree}
\noindent$\Pi_{22}$:
\begin{ptree}
\AXC{$\stackrel{\Pi_{14}}{\vdots}$}\noLine
\UIC{$\mathrm{ContOn}(f,A) \to \mathrm{ContOn}(g,A) \to \varphi_{29}$}
\AXC{$\stackrel{\Pi_{17}}{\vdots}$}\noLine
\UIC{$\varphi_{14}$}
\AXC{$\stackrel{\Pi_{21}}{\vdots}$}\noLine
\UIC{$\varphi_{15}$}
\AXC{$\varphi_{15} \to \varphi_{54}$\lbl{A1}}
\BIC{$\varphi_{54}$}
\BIC{$(\mathrm{ContOn}(f,A) \to \mathrm{ContOn}(g,A) \to \varphi_{29}) \to \varphi_{26}$}
\BIC{$\varphi_{26}$}
\AXC{$\varphi_{26} \to \varphi_{27}$\lbl{A1}}
\BIC{$\varphi_{27}$}
\UIC{$\forall u\,(\varphi_{27})$}
\end{ptree}
\noindent$\Pi_{23}$:
\begin{ptree}
\AXC{$\forall x\,\forall y\,(x \subseteq y \leftrightarrow \forall u \in x\,(u \in y))$\lbl{Def$\subseteq$}}
\AXC{$\varphi_{36}$\lbl{A2}}
\BIC{$\forall y\,(A \subseteq y \leftrightarrow \forall u \in A\,(u \in y))$}
\AXC{$\forall y\,(A \subseteq y \leftrightarrow \forall u \in A\,(u \in y)) \to \varphi_{25}$\lbl{A2}}
\BIC{$\varphi_{25}$}
\AXC{$\varphi_{25} \to \varphi_{30} \to \varphi_{13}$\lbl{A1}}
\BIC{$\varphi_{30} \to \varphi_{13}$}
\end{ptree}
\noindent$\Pi_{24}$:
\begin{ptree}
\AXC{$\varphi_{24}$\lbl{A4}}
\AXC{$\forall u\,(\varphi_{27}) \to \varphi_{31}$\lbl{A3}}
\AXC{$\stackrel{\Pi_{22}}{\vdots}$}\noLine
\UIC{$\forall u\,(\varphi_{27})$}
\AXC{$\forall u\,(\varphi_{27}) \to (\forall u\,(\varphi_{27}) \to \varphi_{31}) \to \varphi_{24} \to \varphi_{30}$\lbl{A1}}
\BIC{$(\forall u\,(\varphi_{27}) \to \varphi_{31}) \to \varphi_{24} \to \varphi_{30}$}
\BIC{$\varphi_{24} \to \varphi_{30}$}
\BIC{$\varphi_{30}$}
\AXC{$\stackrel{\Pi_{23}}{\vdots}$}\noLine
\UIC{$\varphi_{30} \to \varphi_{13}$}
\BIC{$\varphi_{13}$}
\end{ptree}
\noindent$\Pi_{25}$:
\begin{ptree}
\AXC{$\stackrel{\Pi_{13}}{\vdots}$}\noLine
\UIC{$\mathrm{ContOn}(f,A) \to \mathrm{ContOn}(g,A) \to \varphi_{49}$}
\AXC{$(\mathrm{ContOn}(f,A) \to \mathrm{ContOn}(g,A) \to \varphi_{49}) \to \varphi_{5}$\lbl{A1}}
\BIC{$\varphi_{5}$}
\end{ptree}
\noindent$\Pi_{26}$:
\begin{ptree}
\AXC{$\forall f\,\forall A\,(\varphi_{33})$\lbl{DefContOn}}
\AXC{$\forall f\,\forall A\,(\varphi_{33}) \to \forall A\,(\varphi_{41})$\lbl{A2}}
\BIC{$\forall A\,(\varphi_{41})$}
\AXC{$\forall A\,(\varphi_{41}) \to \varphi_{41}$\lbl{A2}}
\BIC{$\varphi_{41}$}
\AXC{$\varphi_{41} \to \varphi_{5} \to \varphi_{13} \to (\mathrm{ContOn}(f,A) \to \varphi_{12}) \to \varphi_{9}$\lbl{A1}}
\BIC{$\varphi_{5} \to \varphi_{13} \to (\mathrm{ContOn}(f,A) \to \varphi_{12}) \to \varphi_{9}$}
\end{ptree}
\noindent 主木:
\begin{ptree}
\AXC{$\stackrel{\Pi_{8}}{\vdots}$}\noLine
\UIC{$\mathrm{ContOn}(f,A) \to \varphi_{12}$}
\AXC{$\stackrel{\Pi_{24}}{\vdots}$}\noLine
\UIC{$\varphi_{13}$}
\AXC{$\stackrel{\Pi_{25}}{\vdots}$}\noLine
\UIC{$\varphi_{5}$}
\AXC{$\stackrel{\Pi_{26}}{\vdots}$}\noLine
\UIC{$\varphi_{5} \to \varphi_{13} \to (\mathrm{ContOn}(f,A) \to \varphi_{12}) \to \varphi_{9}$}
\BIC{$\varphi_{13} \to (\mathrm{ContOn}(f,A) \to \varphi_{12}) \to \varphi_{9}$}
\BIC{$(\mathrm{ContOn}(f,A) \to \varphi_{12}) \to \varphi_{9}$}
\BIC{$\varphi_{9}$}
\UIC{$\forall A\,(\varphi_{9})$}
\UIC{$\forall g\,\forall A\,(\varphi_{9})$}
\UIC{$\forall f\,\forall g\,\forall A\,(\varphi_{9})$}
\end{ptree}
\begin{abbrevlist}
\item $\varphi_{1} :\equiv \mathrm{ContIn}(g,A,c) \to \mathrm{ContIn}((f + g),A,c)$
\item $\varphi_{2} :\equiv \mathrm{ContOn}(f,A) \to \mathrm{ContOn}(g,A) \to \mathrm{ContOn}(g,A)$
\item $\varphi_{3} :\equiv \mathrm{ContOn}(f,A) \to \mathrm{ContOn}(g,A) \to A \subseteq \mathrm{dom}(g)$
\item $\varphi_{4} :\equiv \mathrm{ContOn}(f,A) \to \mathrm{ContOn}(g,A) \to A \subseteq \mathrm{dom}(f)$
\item $\varphi_{5} :\equiv \mathrm{ContOn}(f,A) \to \mathrm{ContOn}(g,A) \to \mathrm{RFn}((f + g))$
\item $\varphi_{6} :\equiv \mathrm{ContOn}(f,A) \to \mathrm{ContOn}(g,A) \to c \in A \to c \in A$
\item $\varphi_{7} :\equiv \mathrm{ContOn}(f,A) \to \mathrm{ContOn}(g,A) \to u \in A \to u \in A$
\item $\varphi_{8} :\equiv \mathrm{ContOn}(g,A) \to c \in A \to \mathrm{ContIn}((f + g),A,c)$
\item $\varphi_{9} :\equiv \mathrm{ContOn}(f,A) \to \mathrm{ContOn}(g,A) \to \mathrm{ContOn}((f + g),A)$
\item $\varphi_{10} :\equiv \mathrm{ContOn}(g,A) \to c \in A \to c \in A \to \mathrm{ContIn}(g,A,c)$
\item $\varphi_{11} :\equiv \mathrm{ContOn}(g,A) \to c \in A \to c \in A \to \mathrm{ContIn}(f,A,c)$
\item $\varphi_{12} :\equiv \mathrm{ContOn}(g,A) \to \forall c \in A\,(\mathrm{ContIn}((f + g),A,c))$
\item $\varphi_{13} :\equiv \mathrm{ContOn}(f,A) \to \mathrm{ContOn}(g,A) \to A \subseteq \mathrm{dom}((f + g))$
\item $\varphi_{14} :\equiv \mathrm{ContOn}(f,A) \to \mathrm{ContOn}(g,A) \to u \in A \to u \in \mathrm{dom}(g)$
\item $\varphi_{15} :\equiv \mathrm{ContOn}(f,A) \to \mathrm{ContOn}(g,A) \to u \in A \to u \in \mathrm{dom}(f)$
\item $\varphi_{16} :\equiv \mathrm{ContOn}(g,A) \to c \in A \to \forall c \in A\,(\mathrm{ContIn}(g,A,c))$
\item $\varphi_{17} :\equiv \mathrm{ContOn}(g,A) \to c \in A \to \forall c \in A\,(\mathrm{ContIn}(f,A,c))$
\item $\varphi_{18} :\equiv \mathrm{ContOn}(f,A) \to \mathrm{ContOn}(g,A) \to c \in A \to \mathrm{ContIn}(g,A,c)$
\item $\varphi_{19} :\equiv \mathrm{ContOn}(f,A) \to \mathrm{ContOn}(g,A) \to c \in A \to \mathrm{ContIn}(f,A,c)$
\item $\varphi_{20} :\equiv \forall c \in A\,(\mathrm{ContIn}(g,A,c)) \to c \in A \to \mathrm{ContIn}(g,A,c)$
\item $\varphi_{21} :\equiv \forall c \in A\,(\mathrm{ContIn}(f,A,c)) \to c \in A \to \mathrm{ContIn}(f,A,c)$
\item $\varphi_{22} :\equiv \forall x\,(x \in \mathrm{dom}((f + g)) \leftrightarrow x \in \mathrm{dom}(f) \land x \in \mathrm{dom}(g))$
\item $\varphi_{23} :\equiv \mathrm{ContOn}(f,A) \land \mathrm{ContOn}(g,A) \to \forall c\,(\mathrm{ContOn}(f,A) \land \mathrm{ContOn}(g,A))$
\item $\varphi_{24} :\equiv \mathrm{ContOn}(f,A) \land \mathrm{ContOn}(g,A) \to \forall u\,(\mathrm{ContOn}(f,A) \land \mathrm{ContOn}(g,A))$
\item $\varphi_{25} :\equiv A \subseteq \mathrm{dom}((f + g)) \leftrightarrow \forall u \in A\,(u \in \mathrm{dom}((f + g)))$
\item $\varphi_{26} :\equiv \mathrm{ContOn}(f,A) \to \mathrm{ContOn}(g,A) \to u \in A \to u \in \mathrm{dom}((f + g))$
\item $\varphi_{27} :\equiv \mathrm{ContOn}(f,A) \land \mathrm{ContOn}(g,A) \to u \in A \to u \in \mathrm{dom}((f + g))$
\item $\varphi_{28} :\equiv \mathrm{ContOn}(f,A) \land \mathrm{ContOn}(g,A) \to c \in A \to \mathrm{ContIn}((f + g),A,c)$
\item $\varphi_{29} :\equiv u \in A \to (u \in \mathrm{dom}((f + g)) \leftrightarrow u \in \mathrm{dom}(f) \land u \in \mathrm{dom}(g))$
\item $\varphi_{30} :\equiv \mathrm{ContOn}(f,A) \to \mathrm{ContOn}(g,A) \to \forall u \in A\,(u \in \mathrm{dom}((f + g)))$
\item $\varphi_{31} :\equiv \forall u\,(\mathrm{ContOn}(f,A) \land \mathrm{ContOn}(g,A)) \to \forall u \in A\,(u \in \mathrm{dom}((f + g)))$
\item $\varphi_{32} :\equiv \mathrm{ContOn}(g,A) \leftrightarrow \mathrm{RFn}(g) \land A \subseteq \mathrm{dom}(g) \land \forall c \in A\,(\mathrm{ContIn}(g,A,c))$
\item $\varphi_{33} :\equiv \mathrm{ContOn}(f,A) \leftrightarrow \mathrm{RFn}(f) \land A \subseteq \mathrm{dom}(f) \land \forall c \in A\,(\mathrm{ContIn}(f,A,c))$
\item $\varphi_{34} :\equiv \forall c\,(\mathrm{ContOn}(f,A) \land \mathrm{ContOn}(g,A)) \to \forall c \in A\,(\mathrm{ContIn}((f + g),A,c))$
\item $\varphi_{35} :\equiv \forall X\,\forall Y\,\forall x\,(X \subseteq Y \to x \in X \to x \in Y) \to \forall Y\,\forall x\,(A \subseteq Y \to x \in A \to x \in Y)$
\item $\varphi_{36} :\equiv \forall x\,\forall y\,(x \subseteq y \leftrightarrow \forall u \in x\,(u \in y)) \to \forall y\,(A \subseteq y \leftrightarrow \forall u \in A\,(u \in y))$
\item $\varphi_{37} :\equiv \forall Y\,\forall x\,(A \subseteq Y \to x \in A \to x \in Y) \to \forall x\,(A \subseteq \mathrm{dom}(g) \to x \in A \to x \in \mathrm{dom}(g))$
\item $\varphi_{38} :\equiv \forall Y\,\forall x\,(A \subseteq Y \to x \in A \to x \in Y) \to \forall x\,(A \subseteq \mathrm{dom}(f) \to x \in A \to x \in \mathrm{dom}(f))$
\item $\varphi_{39} :\equiv \forall x\,(A \subseteq \mathrm{dom}(g) \to x \in A \to x \in \mathrm{dom}(g)) \to A \subseteq \mathrm{dom}(g) \to u \in A \to u \in \mathrm{dom}(g)$
\item $\varphi_{40} :\equiv \forall x\,(A \subseteq \mathrm{dom}(f) \to x \in A \to x \in \mathrm{dom}(f)) \to A \subseteq \mathrm{dom}(f) \to u \in A \to u \in \mathrm{dom}(f)$
\item $\varphi_{41} :\equiv \mathrm{ContOn}((f + g),A) \leftrightarrow \mathrm{RFn}((f + g)) \land A \subseteq \mathrm{dom}((f + g)) \land \forall c \in A\,(\mathrm{ContIn}((f + g),A,c))$
\item $\varphi_{42} :\equiv \varphi_{32} \to \mathrm{ContOn}(f,A) \to \mathrm{ContOn}(g,A) \to \mathrm{RFn}(g)$
\item $\varphi_{43} :\equiv \varphi_{33} \to \mathrm{ContOn}(f,A) \to \mathrm{ContOn}(g,A) \to \mathrm{RFn}(f)$
\item $\varphi_{44} :\equiv \varphi_{22} \to (u \in \mathrm{dom}((f + g)) \leftrightarrow u \in \mathrm{dom}(f) \land u \in \mathrm{dom}(g))$
\item $\varphi_{45} :\equiv \forall c\,(\mathrm{ContIn}(f,A,c) \to \varphi_{1}) \to \mathrm{ContIn}(f,A,c) \to \varphi_{1}$
\item $\varphi_{46} :\equiv \forall A\,\forall c\,(\mathrm{ContIn}(f,A,c) \to \varphi_{1}) \to \forall c\,(\mathrm{ContIn}(f,A,c) \to \varphi_{1})$
\item $\varphi_{47} :\equiv \forall g\,\forall A\,\forall c\,(\mathrm{ContIn}(f,A,c) \to \varphi_{1}) \to \forall A\,\forall c\,(\mathrm{ContIn}(f,A,c) \to \varphi_{1})$
\item $\varphi_{48} :\equiv \forall f\,\forall g\,\forall A\,\forall c\,(\mathrm{ContIn}(f,A,c) \to \varphi_{1}) \to \forall g\,\forall A\,\forall c\,(\mathrm{ContIn}(f,A,c) \to \varphi_{1})$
\item $\varphi_{49} :\equiv \mathrm{RFn}((f + g)) \land \varphi_{22} \land \forall x \in \mathrm{dom}((f + g))\,((f + g)(x) = f(x) + g(x))$
\item $\varphi_{50} :\equiv (\mathrm{ContOn}(f,A) \to \mathrm{ContOn}(f,A)) \to \varphi_{33} \to \mathrm{ContOn}(f,A) \to \varphi_{17}$
\item $\varphi_{51} :\equiv (\mathrm{ContOn}(f,A) \to \varphi_{16}) \to \varphi_{20} \to \mathrm{ContOn}(f,A) \to \varphi_{10}$
\item $\varphi_{52} :\equiv (\mathrm{ContOn}(f,A) \to \varphi_{17}) \to \varphi_{21} \to \mathrm{ContOn}(f,A) \to \varphi_{11}$
\item $\varphi_{53} :\equiv (\mathrm{ContOn}(f,A) \to \mathrm{ContOn}(g,A) \to \mathrm{RFn}(g)) \to \mathrm{ContOn}(f,A) \to \mathrm{ContOn}(g,A) \to \varphi_{49}$
\item $\varphi_{54} :\equiv \varphi_{14} \to (\mathrm{ContOn}(f,A) \to \mathrm{ContOn}(g,A) \to \varphi_{29}) \to \varphi_{26}$
\item $\varphi_{55} :\equiv (\mathrm{ContIn}(f,A,c) \to \varphi_{1}) \to \varphi_{19} \to \varphi_{18} \to \mathrm{ContOn}(f,A) \to \varphi_{8}$
\item $\varphi_{56} :\equiv (\mathrm{ContOn}(f,A) \to \mathrm{ContOn}(g,A) \to \mathrm{RFn}(f)) \to \varphi_{53}$
\item $\varphi_{57} :\equiv (\mathrm{ContOn}(f,A) \to \mathrm{ContOn}(g,A) \to \varphi_{49}) \to \mathrm{ContOn}(f,A) \to \mathrm{ContOn}(g,A) \to u \in A \to \varphi_{22}$
\item $\varphi_{58} :\equiv (\mathrm{ContOn}(f,A) \to \mathrm{ContOn}(g,A) \to u \in A \to \varphi_{22}) \to \varphi_{44} \to \mathrm{ContOn}(f,A) \to \mathrm{ContOn}(g,A) \to \varphi_{29}$
\item $\varphi_{59} :\equiv \forall g\,(\mathrm{RFn}(f) \to \mathrm{RFn}(g) \to \varphi_{49}) \to \mathrm{RFn}(f) \to \mathrm{RFn}(g) \to \varphi_{49}$
\item $\varphi_{60} :\equiv \forall c\,(\varphi_{28}) \to (\forall c\,(\varphi_{28}) \to \varphi_{34}) \to \varphi_{23} \to \mathrm{ContOn}(f,A) \to \varphi_{12}$
\item $\varphi_{61} :\equiv \forall f\,\forall g\,(\mathrm{RFn}(f) \to \mathrm{RFn}(g) \to \varphi_{49}) \to \forall g\,(\mathrm{RFn}(f) \to \mathrm{RFn}(g) \to \varphi_{49})$
\end{abbrevlist}
\end{fullproof}

% 自動生成 (prover/build.js): ContOnComp  ノード数 178
\begin{fullproof}{定理2.10}{$\forall f\,\forall g\,\forall A\,\forall B\,(\mathrm{ContOn}(f,A) \to \mathrm{ContOn}(g,B) \to \forall x \in A\,(f(x) \in B) \to \mathrm{ContOn}(g \circ f,A))$}{178}\phantomsection\label{tree:ContOnComp}
\noindent$\Pi_{1}$:
\begin{ptree}
\AXC{$\forall x \in A\,(f(x) \in B) \to c \in A \to f(c) \in B$\lbl{A2}}
\AXC{$\mathrm{ContOn}(f,A) \to \varphi_{28}$\lbl{A1}}
\AXC{$(\mathrm{ContOn}(f,A) \to \varphi_{28}) \to \varphi_{45}$\lbl{A1}}
\BIC{$\varphi_{45}$}
\BIC{$\mathrm{ContOn}(f,A) \to \varphi_{26}$}
\end{ptree}
\noindent$\Pi_{2}$:
\begin{ptree}
\AXC{$\varphi_{1}$\lbl{A1}}
\AXC{$\varphi_{1} \to \varphi_{17} \to \mathrm{ContOn}(f,A) \to \mathrm{ContOn}(g,B) \to \varphi_{21}$\lbl{A1}}
\BIC{$\varphi_{17} \to \mathrm{ContOn}(f,A) \to \mathrm{ContOn}(g,B) \to \varphi_{21}$}
\end{ptree}
\noindent$\Pi_{3}$:
\begin{ptree}
\AXC{$\forall f\,\forall A\,(\varphi_{18})$\lbl{DefContOn}}
\AXC{$\forall f\,\forall A\,(\varphi_{18}) \to \varphi_{25}$\lbl{A2}}
\BIC{$\varphi_{25}$}
\AXC{$\varphi_{25} \to \varphi_{17}$\lbl{A2}}
\BIC{$\varphi_{17}$}
\AXC{$\stackrel{\Pi_{2}}{\vdots}$}\noLine
\UIC{$\varphi_{17} \to \mathrm{ContOn}(f,A) \to \mathrm{ContOn}(g,B) \to \varphi_{21}$}
\BIC{$\mathrm{ContOn}(f,A) \to \mathrm{ContOn}(g,B) \to \varphi_{21}$}
\end{ptree}
\noindent$\Pi_{4}$:
\begin{ptree}
\AXC{$\stackrel{\Pi_{1}}{\vdots}$}\noLine
\UIC{$\mathrm{ContOn}(f,A) \to \varphi_{26}$}
\AXC{$\varphi_{20}$\lbl{A2}}
\AXC{$\stackrel{\Pi_{3}}{\vdots}$}\noLine
\UIC{$\mathrm{ContOn}(f,A) \to \mathrm{ContOn}(g,B) \to \varphi_{21}$}
\AXC{$\varphi_{71}$\lbl{A1}}
\BIC{$\varphi_{20} \to \mathrm{ContOn}(f,A) \to \mathrm{ContOn}(g,B) \to \varphi_{35}$}
\BIC{$\mathrm{ContOn}(f,A) \to \mathrm{ContOn}(g,B) \to \varphi_{35}$}
\AXC{$\varphi_{77}$\lbl{A1}}
\BIC{$(\mathrm{ContOn}(f,A) \to \varphi_{26}) \to \varphi_{48}$}
\BIC{$\varphi_{48}$}
\end{ptree}
\noindent$\Pi_{5}$:
\begin{ptree}
\AXC{$\forall f\,\forall A\,(\varphi_{18})$\lbl{DefContOn}}
\AXC{$\forall f\,\forall A\,(\varphi_{18}) \to \forall A\,(\varphi_{18})$\lbl{A2}}
\BIC{$\forall A\,(\varphi_{18})$}
\AXC{$\forall A\,(\varphi_{18}) \to \varphi_{18}$\lbl{A2}}
\BIC{$\varphi_{18}$}
\end{ptree}
\noindent$\Pi_{6}$:
\begin{ptree}
\AXC{$\varphi_{4}$\lbl{A2}}
\AXC{$\stackrel{\Pi_{5}}{\vdots}$}\noLine
\UIC{$\varphi_{18}$}
\AXC{$\mathrm{ContOn}(f,A) \to \mathrm{ContOn}(f,A)$\lbl{A1}}
\AXC{$\varphi_{58}$\lbl{A1}}
\BIC{$\varphi_{18} \to \mathrm{ContOn}(f,A) \to \mathrm{ContOn}(g,B) \to \varphi_{22}$}
\BIC{$\mathrm{ContOn}(f,A) \to \mathrm{ContOn}(g,B) \to \varphi_{22}$}
\AXC{$\varphi_{66}$\lbl{A1}}
\BIC{$\varphi_{4} \to \mathrm{ContOn}(f,A) \to \mathrm{ContOn}(g,B) \to \varphi_{14}$}
\BIC{$\mathrm{ContOn}(f,A) \to \mathrm{ContOn}(g,B) \to \varphi_{14}$}
\AXC{$\varphi_{68}$\lbl{A1}}
\BIC{$(\mathrm{ContOn}(f,A) \to \varphi_{11}) \to \mathrm{ContOn}(f,A) \to \varphi_{24}$}
\end{ptree}
\noindent$\Pi_{7}$:
\begin{ptree}
\AXC{$\mathrm{ContOn}(f,A) \to \varphi_{11}$\lbl{A1}}
\AXC{$\stackrel{\Pi_{6}}{\vdots}$}\noLine
\UIC{$(\mathrm{ContOn}(f,A) \to \varphi_{11}) \to \mathrm{ContOn}(f,A) \to \varphi_{24}$}
\BIC{$\mathrm{ContOn}(f,A) \to \varphi_{24}$}
\end{ptree}
\noindent$\Pi_{8}$:
\begin{ptree}
\AXC{$\forall f\,\forall g\,\forall A\,\forall B\,\forall c\,(\mathrm{ContIn}(f,A,c) \to \varphi_{33})$\lbl{定理2.6}}
\AXC{$\varphi_{62}$\lbl{A2}}
\BIC{$\forall g\,\forall A\,\forall B\,\forall c\,(\mathrm{ContIn}(f,A,c) \to \varphi_{33})$}
\AXC{$\varphi_{61}$\lbl{A2}}
\BIC{$\forall A\,\forall B\,\forall c\,(\mathrm{ContIn}(f,A,c) \to \varphi_{33})$}
\AXC{$\varphi_{60}$\lbl{A2}}
\BIC{$\forall B\,\forall c\,(\mathrm{ContIn}(f,A,c) \to \varphi_{33})$}
\AXC{$\varphi_{59}$\lbl{A2}}
\BIC{$\forall c\,(\mathrm{ContIn}(f,A,c) \to \varphi_{33})$}
\AXC{$\varphi_{55}$\lbl{A2}}
\BIC{$\mathrm{ContIn}(f,A,c) \to \varphi_{33}$}
\end{ptree}
\noindent$\Pi_{9}$:
\begin{ptree}
\AXC{$\stackrel{\Pi_{7}}{\vdots}$}\noLine
\UIC{$\mathrm{ContOn}(f,A) \to \varphi_{24}$}
\AXC{$\stackrel{\Pi_{8}}{\vdots}$}\noLine
\UIC{$\mathrm{ContIn}(f,A,c) \to \varphi_{33}$}
\AXC{$(\mathrm{ContIn}(f,A,c) \to \varphi_{33}) \to \varphi_{76}$\lbl{A1}}
\BIC{$\varphi_{76}$}
\BIC{$(\mathrm{ContOn}(f,A) \to \mathrm{ContOn}(g,B) \to \varphi_{10}) \to \varphi_{52}$}
\end{ptree}
\noindent$\Pi_{10}$:
\begin{ptree}
\AXC{$\mathrm{ContOn}(f,A) \to \varphi_{11}$\lbl{A1}}
\AXC{$\stackrel{\Pi_{4}}{\vdots}$}\noLine
\UIC{$\varphi_{48}$}
\BIC{$\mathrm{ContOn}(f,A) \to \mathrm{ContOn}(g,B) \to \varphi_{10}$}
\AXC{$\stackrel{\Pi_{9}}{\vdots}$}\noLine
\UIC{$(\mathrm{ContOn}(f,A) \to \mathrm{ContOn}(g,B) \to \varphi_{10}) \to \varphi_{52}$}
\BIC{$\varphi_{52}$}
\end{ptree}
\noindent$\Pi_{11}$:
\begin{ptree}
\AXC{$\varphi_{57}$\lbl{A3}}
\AXC{$\mathrm{ContOn}(f,A) \to \varphi_{28}$\lbl{A1}}
\AXC{$\stackrel{\Pi_{10}}{\vdots}$}\noLine
\UIC{$\varphi_{52}$}
\BIC{$\mathrm{ContOn}(f,A) \to \mathrm{ContOn}(g,B) \to \varphi_{12}$}
\AXC{$\varphi_{54}$\lbl{A1}}
\BIC{$\varphi_{5} \to c \in A \to \mathrm{ContIn}(g \circ f,A,c)$}
\UIC{$\forall c\,(\varphi_{5} \to c \in A \to \mathrm{ContIn}(g \circ f,A,c))$}
\AXC{$\forall c\,(\varphi_{5} \to c \in A \to \mathrm{ContIn}(g \circ f,A,c)) \to \varphi_{57} \to \varphi_{63}$\lbl{A1}}
\BIC{$\varphi_{57} \to \varphi_{63}$}
\BIC{$\varphi_{63}$}
\end{ptree}
\noindent$\Pi_{12}$:
\begin{ptree}
\AXC{$\forall x \in A\,(f(x) \in B) \to u \in A \to f(u) \in B$\lbl{A2}}
\AXC{$\mathrm{ContOn}(f,A) \to \varphi_{28}$\lbl{A1}}
\AXC{$(\mathrm{ContOn}(f,A) \to \varphi_{28}) \to \varphi_{46}$\lbl{A1}}
\BIC{$\varphi_{46}$}
\BIC{$\mathrm{ContOn}(f,A) \to \varphi_{27}$}
\end{ptree}
\noindent$\Pi_{13}$:
\begin{ptree}
\AXC{$\forall f\,\forall A\,(\varphi_{18})$\lbl{DefContOn}}
\AXC{$\forall f\,\forall A\,(\varphi_{18}) \to \varphi_{25}$\lbl{A2}}
\BIC{$\varphi_{25}$}
\AXC{$\varphi_{25} \to \varphi_{17}$\lbl{A2}}
\BIC{$\varphi_{17}$}
\AXC{$\varphi_{1}$\lbl{A1}}
\AXC{$\varphi_{1} \to \varphi_{17} \to \mathrm{ContOn}(f,A) \to \varphi_{6}$\lbl{A1}}
\BIC{$\varphi_{17} \to \mathrm{ContOn}(f,A) \to \varphi_{6}$}
\BIC{$\mathrm{ContOn}(f,A) \to \varphi_{6}$}
\end{ptree}
\noindent$\Pi_{14}$:
\begin{ptree}
\AXC{$\stackrel{\Pi_{13}}{\vdots}$}\noLine
\UIC{$\mathrm{ContOn}(f,A) \to \varphi_{6}$}
\AXC{$\forall X\,\forall Y\,\forall x\,(X \subseteq Y \to x \in X \to x \in Y)$\lbl{[B5]}}
\AXC{$\varphi_{30}$\lbl{A2}}
\BIC{$\forall Y\,\forall x\,(B \subseteq Y \to x \in B \to x \in Y)$}
\AXC{$\varphi_{37}$\lbl{A2}}
\BIC{$\forall x\,(B \subseteq \mathrm{dom}(g) \to x \in B \to x \in \mathrm{dom}(g))$}
\AXC{$\varphi_{42}$\lbl{A2}}
\BIC{$B \subseteq \mathrm{dom}(g) \to f(u) \in B \to f(u) \in \mathrm{dom}(g)$}
\AXC{$\varphi_{73}$\lbl{A1}}
\BIC{$\varphi_{67}$}
\BIC{$(\mathrm{ContOn}(f,A) \to \varphi_{27}) \to \mathrm{ContOn}(f,A) \to \varphi_{29}$}
\end{ptree}
\noindent$\Pi_{15}$:
\begin{ptree}
\AXC{$\stackrel{\Pi_{12}}{\vdots}$}\noLine
\UIC{$\mathrm{ContOn}(f,A) \to \varphi_{27}$}
\AXC{$\stackrel{\Pi_{14}}{\vdots}$}\noLine
\UIC{$(\mathrm{ContOn}(f,A) \to \varphi_{27}) \to \mathrm{ContOn}(f,A) \to \varphi_{29}$}
\BIC{$\mathrm{ContOn}(f,A) \to \varphi_{29}$}
\end{ptree}
\noindent$\Pi_{16}$:
\begin{ptree}
\AXC{$\forall f\,\forall A\,(\varphi_{18})$\lbl{DefContOn}}
\AXC{$\forall f\,\forall A\,(\varphi_{18}) \to \forall A\,(\varphi_{18})$\lbl{A2}}
\BIC{$\forall A\,(\varphi_{18})$}
\AXC{$\forall A\,(\varphi_{18}) \to \varphi_{18}$\lbl{A2}}
\BIC{$\varphi_{18}$}
\end{ptree}
\noindent$\Pi_{17}$:
\begin{ptree}
\AXC{$\forall X\,\forall Y\,\forall x\,(X \subseteq Y \to x \in X \to x \in Y)$\lbl{[B5]}}
\AXC{$\varphi_{31}$\lbl{A2}}
\BIC{$\forall Y\,\forall x\,(A \subseteq Y \to x \in A \to x \in Y)$}
\AXC{$\varphi_{38}$\lbl{A2}}
\BIC{$\forall x\,(A \subseteq \mathrm{dom}(f) \to x \in A \to x \in \mathrm{dom}(f))$}
\AXC{$\varphi_{40}$\lbl{A2}}
\BIC{$A \subseteq \mathrm{dom}(f) \to u \in A \to u \in \mathrm{dom}(f)$}
\end{ptree}
\noindent$\Pi_{18}$:
\begin{ptree}
\AXC{$\stackrel{\Pi_{17}}{\vdots}$}\noLine
\UIC{$A \subseteq \mathrm{dom}(f) \to u \in A \to u \in \mathrm{dom}(f)$}
\AXC{$(A \subseteq \mathrm{dom}(f) \to u \in A \to u \in \mathrm{dom}(f)) \to \varphi_{64}$\lbl{A1}}
\BIC{$\varphi_{64}$}
\end{ptree}
\noindent$\Pi_{19}$:
\begin{ptree}
\AXC{$\mathrm{ContOn}(f,A) \to \varphi_{13}$\lbl{A1}}
\AXC{$\stackrel{\Pi_{16}}{\vdots}$}\noLine
\UIC{$\varphi_{18}$}
\AXC{$\mathrm{ContOn}(f,A) \to \mathrm{ContOn}(f,A)$\lbl{A1}}
\AXC{$\varphi_{49}$\lbl{A1}}
\BIC{$\varphi_{18} \to \mathrm{ContOn}(f,A) \to \varphi_{7}$}
\BIC{$\mathrm{ContOn}(f,A) \to \varphi_{7}$}
\AXC{$\stackrel{\Pi_{18}}{\vdots}$}\noLine
\UIC{$\varphi_{64}$}
\BIC{$(\mathrm{ContOn}(f,A) \to \varphi_{13}) \to \mathrm{ContOn}(f,A) \to \varphi_{23}$}
\BIC{$\mathrm{ContOn}(f,A) \to \varphi_{23}$}
\end{ptree}
\noindent$\Pi_{20}$:
\begin{ptree}
\AXC{$\forall f\,\forall A\,(\varphi_{18})$\lbl{DefContOn}}
\AXC{$\forall f\,\forall A\,(\varphi_{18}) \to \forall A\,(\varphi_{18})$\lbl{A2}}
\BIC{$\forall A\,(\varphi_{18})$}
\AXC{$\forall A\,(\varphi_{18}) \to \varphi_{18}$\lbl{A2}}
\BIC{$\varphi_{18}$}
\end{ptree}
\noindent$\Pi_{21}$:
\begin{ptree}
\AXC{$\stackrel{\Pi_{20}}{\vdots}$}\noLine
\UIC{$\varphi_{18}$}
\AXC{$\mathrm{ContOn}(f,A) \to \mathrm{ContOn}(f,A)$\lbl{A1}}
\AXC{$\varphi_{47}$\lbl{A1}}
\BIC{$\varphi_{18} \to \mathrm{ContOn}(f,A) \to \varphi_{2}$}
\BIC{$\mathrm{ContOn}(f,A) \to \varphi_{2}$}
\end{ptree}
\noindent$\Pi_{22}$:
\begin{ptree}
\AXC{$\forall f\,\forall A\,(\varphi_{18})$\lbl{DefContOn}}
\AXC{$\forall f\,\forall A\,(\varphi_{18}) \to \varphi_{25}$\lbl{A2}}
\BIC{$\varphi_{25}$}
\AXC{$\varphi_{25} \to \varphi_{17}$\lbl{A2}}
\BIC{$\varphi_{17}$}
\AXC{$\varphi_{1}$\lbl{A1}}
\AXC{$\varphi_{1} \to \varphi_{17} \to \mathrm{ContOn}(f,A) \to \varphi_{3}$\lbl{A1}}
\BIC{$\varphi_{17} \to \mathrm{ContOn}(f,A) \to \varphi_{3}$}
\BIC{$\mathrm{ContOn}(f,A) \to \varphi_{3}$}
\end{ptree}
\noindent$\Pi_{23}$:
\begin{ptree}
\AXC{$\forall g\,\forall f\,(\varphi_{39})$\lbl{[S23]}}
\AXC{$\forall g\,\forall f\,(\varphi_{39}) \to \forall f\,(\varphi_{39})$\lbl{A2}}
\BIC{$\forall f\,(\varphi_{39})$}
\AXC{$\forall f\,(\varphi_{39}) \to \varphi_{39}$\lbl{A2}}
\BIC{$\varphi_{39}$}
\AXC{$\varphi_{39} \to \mathrm{RFn}(g) \to \mathrm{RFn}(f) \to \varphi_{8}$\lbl{A2}}
\BIC{$\mathrm{RFn}(g) \to \mathrm{RFn}(f) \to \varphi_{8}$}
\end{ptree}
\noindent$\Pi_{24}$:
\begin{ptree}
\AXC{$\stackrel{\Pi_{22}}{\vdots}$}\noLine
\UIC{$\mathrm{ContOn}(f,A) \to \varphi_{3}$}
\AXC{$\stackrel{\Pi_{23}}{\vdots}$}\noLine
\UIC{$\mathrm{RFn}(g) \to \mathrm{RFn}(f) \to \varphi_{8}$}
\AXC{$(\mathrm{RFn}(g) \to \mathrm{RFn}(f) \to \varphi_{8}) \to \varphi_{78}$\lbl{A1}}
\BIC{$\varphi_{78}$}
\BIC{$(\mathrm{ContOn}(f,A) \to \varphi_{2}) \to \varphi_{69}$}
\end{ptree}
\noindent$\Pi_{25}$:
\begin{ptree}
\AXC{$\stackrel{\Pi_{19}}{\vdots}$}\noLine
\UIC{$\mathrm{ContOn}(f,A) \to \varphi_{23}$}
\AXC{$\stackrel{\Pi_{21}}{\vdots}$}\noLine
\UIC{$\mathrm{ContOn}(f,A) \to \varphi_{2}$}
\AXC{$\stackrel{\Pi_{24}}{\vdots}$}\noLine
\UIC{$(\mathrm{ContOn}(f,A) \to \varphi_{2}) \to \varphi_{69}$}
\BIC{$\varphi_{69}$}
\BIC{$(\mathrm{ContOn}(f,A) \to \varphi_{29}) \to \mathrm{ContOn}(f,A) \to \varphi_{34}$}
\end{ptree}
\noindent$\Pi_{26}$:
\begin{ptree}
\AXC{$\varphi_{56}$\lbl{A3}}
\AXC{$\stackrel{\Pi_{15}}{\vdots}$}\noLine
\UIC{$\mathrm{ContOn}(f,A) \to \varphi_{29}$}
\AXC{$\stackrel{\Pi_{25}}{\vdots}$}\noLine
\UIC{$(\mathrm{ContOn}(f,A) \to \varphi_{29}) \to \mathrm{ContOn}(f,A) \to \varphi_{34}$}
\BIC{$\mathrm{ContOn}(f,A) \to \varphi_{34}$}
\AXC{$\varphi_{53}$\lbl{A1}}
\BIC{$\varphi_{5} \to u \in A \to u \in \mathrm{dom}(g \circ f)$}
\UIC{$\forall u\,(\varphi_{5} \to u \in A \to u \in \mathrm{dom}(g \circ f))$}
\AXC{$\varphi_{79}$\lbl{A1}}
\BIC{$\varphi_{56} \to (\varphi_{5} \to \forall u\,(\varphi_{5})) \to \mathrm{ContOn}(f,A) \to \varphi_{36}$}
\BIC{$(\varphi_{5} \to \forall u\,(\varphi_{5})) \to \mathrm{ContOn}(f,A) \to \varphi_{36}$}
\end{ptree}
\noindent$\Pi_{27}$:
\begin{ptree}
\AXC{$\forall x\,\forall y\,(x \subseteq y \leftrightarrow \forall u \in x\,(u \in y))$\lbl{Def$\subseteq$}}
\AXC{$\varphi_{32}$\lbl{A2}}
\BIC{$\forall y\,(A \subseteq y \leftrightarrow \forall u \in A\,(u \in y))$}
\AXC{$\varphi_{41}$\lbl{A2}}
\BIC{$A \subseteq \mathrm{dom}(g \circ f) \leftrightarrow \forall u \in A\,(u \in \mathrm{dom}(g \circ f))$}
\AXC{$\varphi_{65}$\lbl{A1}}
\BIC{$(\mathrm{ContOn}(f,A) \to \varphi_{36}) \to \mathrm{ContOn}(f,A) \to \varphi_{19}$}
\end{ptree}
\noindent$\Pi_{28}$:
\begin{ptree}
\AXC{$\varphi_{5} \to \forall u\,(\varphi_{5})$\lbl{A4}}
\AXC{$\stackrel{\Pi_{26}}{\vdots}$}\noLine
\UIC{$(\varphi_{5} \to \forall u\,(\varphi_{5})) \to \mathrm{ContOn}(f,A) \to \varphi_{36}$}
\BIC{$\mathrm{ContOn}(f,A) \to \varphi_{36}$}
\AXC{$\stackrel{\Pi_{27}}{\vdots}$}\noLine
\UIC{$(\mathrm{ContOn}(f,A) \to \varphi_{36}) \to \mathrm{ContOn}(f,A) \to \varphi_{19}$}
\BIC{$\mathrm{ContOn}(f,A) \to \varphi_{19}$}
\end{ptree}
\noindent$\Pi_{29}$:
\begin{ptree}
\AXC{$\forall g\,\forall f\,(\mathrm{RFn}(g) \to \mathrm{RFn}(f) \to \varphi_{44})$\lbl{[S9]}}
\AXC{$\varphi_{72}$\lbl{A2}}
\BIC{$\forall f\,(\mathrm{RFn}(g) \to \mathrm{RFn}(f) \to \varphi_{44})$}
\AXC{$\varphi_{70}$\lbl{A2}}
\BIC{$\mathrm{RFn}(g) \to \mathrm{RFn}(f) \to \varphi_{44}$}
\AXC{$(\mathrm{RFn}(g) \to \mathrm{RFn}(f) \to \varphi_{44}) \to \varphi_{74}$\lbl{A1}}
\BIC{$\varphi_{74}$}
\end{ptree}
\noindent$\Pi_{30}$:
\begin{ptree}
\AXC{$\stackrel{\Pi_{21}}{\vdots}$}\noLine
\UIC{$\mathrm{ContOn}(f,A) \to \varphi_{2}$}
\AXC{$\stackrel{\Pi_{22}}{\vdots}$}\noLine
\UIC{$\mathrm{ContOn}(f,A) \to \varphi_{3}$}
\AXC{$\stackrel{\Pi_{29}}{\vdots}$}\noLine
\UIC{$\varphi_{74}$}
\BIC{$(\mathrm{ContOn}(f,A) \to \varphi_{2}) \to \varphi_{50}$}
\BIC{$\varphi_{50}$}
\AXC{$\varphi_{50} \to \mathrm{ContOn}(f,A) \to \varphi_{9}$\lbl{A1}}
\BIC{$\mathrm{ContOn}(f,A) \to \varphi_{9}$}
\end{ptree}
\noindent$\Pi_{31}$:
\begin{ptree}
\AXC{$\stackrel{\Pi_{30}}{\vdots}$}\noLine
\UIC{$\mathrm{ContOn}(f,A) \to \varphi_{9}$}
\AXC{$\forall f\,\forall A\,(\varphi_{18})$\lbl{DefContOn}}
\AXC{$\forall f\,\forall A\,(\varphi_{18}) \to \forall A\,(\varphi_{43})$\lbl{A2}}
\BIC{$\forall A\,(\varphi_{43})$}
\AXC{$\forall A\,(\varphi_{43}) \to \varphi_{43}$\lbl{A2}}
\BIC{$\varphi_{43}$}
\AXC{$\varphi_{43} \to \varphi_{75}$\lbl{A1}}
\BIC{$\varphi_{75}$}
\BIC{$(\mathrm{ContOn}(f,A) \to \varphi_{19}) \to \varphi_{51}$}
\end{ptree}
\noindent 主木:
\begin{ptree}
\AXC{$\varphi_{5} \to \forall c\,(\varphi_{5})$\lbl{A4}}
\AXC{$\stackrel{\Pi_{11}}{\vdots}$}\noLine
\UIC{$\varphi_{63}$}
\BIC{$\mathrm{ContOn}(f,A) \to \mathrm{ContOn}(g,B) \to \varphi_{16}$}
\AXC{$\stackrel{\Pi_{28}}{\vdots}$}\noLine
\UIC{$\mathrm{ContOn}(f,A) \to \varphi_{19}$}
\AXC{$\stackrel{\Pi_{31}}{\vdots}$}\noLine
\UIC{$(\mathrm{ContOn}(f,A) \to \varphi_{19}) \to \varphi_{51}$}
\BIC{$\varphi_{51}$}
\BIC{$\mathrm{ContOn}(f,A) \to \varphi_{15}$}
\UIC{$\forall B\,(\mathrm{ContOn}(f,A) \to \varphi_{15})$}
\UIC{$\forall A\,\forall B\,(\mathrm{ContOn}(f,A) \to \varphi_{15})$}
\UIC{$\forall g\,\forall A\,\forall B\,(\mathrm{ContOn}(f,A) \to \varphi_{15})$}
\UIC{$\forall f\,\forall g\,\forall A\,\forall B\,(\mathrm{ContOn}(f,A) \to \varphi_{15})$}
\end{ptree}
\begin{abbrevlist}
\item $\varphi_{1} :\equiv \mathrm{ContOn}(f,A) \to \mathrm{ContOn}(g,B) \to \mathrm{ContOn}(g,B)$
\item $\varphi_{2} :\equiv \mathrm{ContOn}(g,B) \to \forall x \in A\,(f(x) \in B) \to \mathrm{RFn}(f)$
\item $\varphi_{3} :\equiv \mathrm{ContOn}(g,B) \to \forall x \in A\,(f(x) \in B) \to \mathrm{RFn}(g)$
\item $\varphi_{4} :\equiv \forall c \in A\,(\mathrm{ContIn}(f,A,c)) \to c \in A \to \mathrm{ContIn}(f,A,c)$
\item $\varphi_{5} :\equiv \mathrm{ContOn}(f,A) \land \mathrm{ContOn}(g,B) \land \forall x \in A\,(f(x) \in B)$
\item $\varphi_{6} :\equiv \mathrm{ContOn}(g,B) \to \forall x \in A\,(f(x) \in B) \to B \subseteq \mathrm{dom}(g)$
\item $\varphi_{7} :\equiv \mathrm{ContOn}(g,B) \to \forall x \in A\,(f(x) \in B) \to A \subseteq \mathrm{dom}(f)$
\item $\varphi_{8} :\equiv u \in \mathrm{dom}(f) \to f(u) \in \mathrm{dom}(g) \to u \in \mathrm{dom}(g \circ f)$
\item $\varphi_{9} :\equiv \mathrm{ContOn}(g,B) \to \forall x \in A\,(f(x) \in B) \to \mathrm{RFn}(g \circ f)$
\item $\varphi_{10} :\equiv \forall x \in A\,(f(x) \in B) \to c \in A \to \mathrm{ContIn}(g,B,f(c))$
\item $\varphi_{11} :\equiv \mathrm{ContOn}(g,B) \to \forall x \in A\,(f(x) \in B) \to c \in A \to c \in A$
\item $\varphi_{12} :\equiv \forall x \in A\,(f(x) \in B) \to c \in A \to \mathrm{ContIn}(g \circ f,A,c)$
\item $\varphi_{13} :\equiv \mathrm{ContOn}(g,B) \to \forall x \in A\,(f(x) \in B) \to u \in A \to u \in A$
\item $\varphi_{14} :\equiv \forall x \in A\,(f(x) \in B) \to c \in A \to c \in A \to \mathrm{ContIn}(f,A,c)$
\item $\varphi_{15} :\equiv \mathrm{ContOn}(g,B) \to \forall x \in A\,(f(x) \in B) \to \mathrm{ContOn}(g \circ f,A)$
\item $\varphi_{16} :\equiv \forall x \in A\,(f(x) \in B) \to \forall c \in A\,(\mathrm{ContIn}(g \circ f,A,c))$
\item $\varphi_{17} :\equiv \mathrm{ContOn}(g,B) \leftrightarrow \mathrm{RFn}(g) \land B \subseteq \mathrm{dom}(g) \land \forall c \in B\,(\mathrm{ContIn}(g,B,c))$
\item $\varphi_{18} :\equiv \mathrm{ContOn}(f,A) \leftrightarrow \mathrm{RFn}(f) \land A \subseteq \mathrm{dom}(f) \land \forall c \in A\,(\mathrm{ContIn}(f,A,c))$
\item $\varphi_{19} :\equiv \mathrm{ContOn}(g,B) \to \forall x \in A\,(f(x) \in B) \to A \subseteq \mathrm{dom}(g \circ f)$
\item $\varphi_{20} :\equiv \forall c \in B\,(\mathrm{ContIn}(g,B,c)) \to f(c) \in B \to \mathrm{ContIn}(g,B,f(c))$
\item $\varphi_{21} :\equiv \forall x \in A\,(f(x) \in B) \to c \in A \to \forall c \in B\,(\mathrm{ContIn}(g,B,c))$
\item $\varphi_{22} :\equiv \forall x \in A\,(f(x) \in B) \to c \in A \to \forall c \in A\,(\mathrm{ContIn}(f,A,c))$
\item $\varphi_{23} :\equiv \mathrm{ContOn}(g,B) \to \forall x \in A\,(f(x) \in B) \to u \in A \to u \in \mathrm{dom}(f)$
\item $\varphi_{24} :\equiv \mathrm{ContOn}(g,B) \to \forall x \in A\,(f(x) \in B) \to c \in A \to \mathrm{ContIn}(f,A,c)$
\item $\varphi_{25} :\equiv \forall A\,(\mathrm{ContOn}(g,A) \leftrightarrow \mathrm{RFn}(g) \land A \subseteq \mathrm{dom}(g) \land \forall c \in A\,(\mathrm{ContIn}(g,A,c)))$
\item $\varphi_{26} :\equiv \mathrm{ContOn}(g,B) \to \forall x \in A\,(f(x) \in B) \to c \in A \to f(c) \in B$
\item $\varphi_{27} :\equiv \mathrm{ContOn}(g,B) \to \forall x \in A\,(f(x) \in B) \to u \in A \to f(u) \in B$
\item $\varphi_{28} :\equiv \mathrm{ContOn}(g,B) \to \forall x \in A\,(f(x) \in B) \to \forall x \in A\,(f(x) \in B)$
\item $\varphi_{29} :\equiv \mathrm{ContOn}(g,B) \to \forall x \in A\,(f(x) \in B) \to u \in A \to f(u) \in \mathrm{dom}(g)$
\item $\varphi_{30} :\equiv \forall X\,\forall Y\,\forall x\,(X \subseteq Y \to x \in X \to x \in Y) \to \forall Y\,\forall x\,(B \subseteq Y \to x \in B \to x \in Y)$
\item $\varphi_{31} :\equiv \forall X\,\forall Y\,\forall x\,(X \subseteq Y \to x \in X \to x \in Y) \to \forall Y\,\forall x\,(A \subseteq Y \to x \in A \to x \in Y)$
\item $\varphi_{32} :\equiv \forall x\,\forall y\,(x \subseteq y \leftrightarrow \forall u \in x\,(u \in y)) \to \forall y\,(A \subseteq y \leftrightarrow \forall u \in A\,(u \in y))$
\item $\varphi_{33} :\equiv \mathrm{ContIn}(g,B,f(c)) \to \forall x \in A\,(f(x) \in B) \to \mathrm{ContIn}(g \circ f,A,c)$
\item $\varphi_{34} :\equiv \mathrm{ContOn}(g,B) \to \forall x \in A\,(f(x) \in B) \to u \in A \to u \in \mathrm{dom}(g \circ f)$
\item $\varphi_{35} :\equiv \forall x \in A\,(f(x) \in B) \to c \in A \to f(c) \in B \to \mathrm{ContIn}(g,B,f(c))$
\item $\varphi_{36} :\equiv \mathrm{ContOn}(g,B) \to \forall x \in A\,(f(x) \in B) \to \forall u \in A\,(u \in \mathrm{dom}(g \circ f))$
\item $\varphi_{37} :\equiv \forall Y\,\forall x\,(B \subseteq Y \to x \in B \to x \in Y) \to \forall x\,(B \subseteq \mathrm{dom}(g) \to x \in B \to x \in \mathrm{dom}(g))$
\item $\varphi_{38} :\equiv \forall Y\,\forall x\,(A \subseteq Y \to x \in A \to x \in Y) \to \forall x\,(A \subseteq \mathrm{dom}(f) \to x \in A \to x \in \mathrm{dom}(f))$
\item $\varphi_{39} :\equiv \forall x\,(\mathrm{RFn}(g) \to \mathrm{RFn}(f) \to x \in \mathrm{dom}(f) \to f(x) \in \mathrm{dom}(g) \to x \in \mathrm{dom}(g \circ f))$
\item $\varphi_{40} :\equiv \forall x\,(A \subseteq \mathrm{dom}(f) \to x \in A \to x \in \mathrm{dom}(f)) \to A \subseteq \mathrm{dom}(f) \to u \in A \to u \in \mathrm{dom}(f)$
\item $\varphi_{41} :\equiv \forall y\,(A \subseteq y \leftrightarrow \forall u \in A\,(u \in y)) \to (A \subseteq \mathrm{dom}(g \circ f) \leftrightarrow \forall u \in A\,(u \in \mathrm{dom}(g \circ f)))$
\item $\varphi_{42} :\equiv \forall x\,(B \subseteq \mathrm{dom}(g) \to x \in B \to x \in \mathrm{dom}(g)) \to B \subseteq \mathrm{dom}(g) \to f(u) \in B \to f(u) \in \mathrm{dom}(g)$
\item $\varphi_{43} :\equiv \mathrm{ContOn}(g \circ f,A) \leftrightarrow \mathrm{RFn}(g \circ f) \land A \subseteq \mathrm{dom}(g \circ f) \land \forall c \in A\,(\mathrm{ContIn}(g \circ f,A,c))$
\item $\varphi_{44} :\equiv \mathrm{RFn}(g \circ f) \land \forall x \in \mathrm{dom}(g \circ f)\,(x \in \mathrm{dom}(f) \land f(x) \in \mathrm{dom}(g) \land g \circ f(x) = g(f(x)))$
\item $\varphi_{45} :\equiv (\forall x \in A\,(f(x) \in B) \to c \in A \to f(c) \in B) \to \mathrm{ContOn}(f,A) \to \varphi_{26}$
\item $\varphi_{46} :\equiv (\forall x \in A\,(f(x) \in B) \to u \in A \to f(u) \in B) \to \mathrm{ContOn}(f,A) \to \varphi_{27}$
\item $\varphi_{47} :\equiv (\mathrm{ContOn}(f,A) \to \mathrm{ContOn}(f,A)) \to \varphi_{18} \to \mathrm{ContOn}(f,A) \to \varphi_{2}$
\item $\varphi_{48} :\equiv (\mathrm{ContOn}(f,A) \to \varphi_{11}) \to \mathrm{ContOn}(f,A) \to \mathrm{ContOn}(g,B) \to \varphi_{10}$
\item $\varphi_{49} :\equiv (\mathrm{ContOn}(f,A) \to \mathrm{ContOn}(f,A)) \to \varphi_{18} \to \mathrm{ContOn}(f,A) \to \varphi_{7}$
\item $\varphi_{50} :\equiv \mathrm{ContOn}(f,A) \to \mathrm{ContOn}(g,B) \to \forall x \in A\,(f(x) \in B) \to \varphi_{44}$
\item $\varphi_{51} :\equiv (\mathrm{ContOn}(f,A) \to \mathrm{ContOn}(g,B) \to \varphi_{16}) \to \mathrm{ContOn}(f,A) \to \varphi_{15}$
\item $\varphi_{52} :\equiv (\mathrm{ContOn}(f,A) \to \varphi_{28}) \to \mathrm{ContOn}(f,A) \to \mathrm{ContOn}(g,B) \to \varphi_{12}$
\item $\varphi_{53} :\equiv (\mathrm{ContOn}(f,A) \to \varphi_{34}) \to \varphi_{5} \to u \in A \to u \in \mathrm{dom}(g \circ f)$
\item $\varphi_{54} :\equiv (\mathrm{ContOn}(f,A) \to \mathrm{ContOn}(g,B) \to \varphi_{12}) \to \varphi_{5} \to c \in A \to \mathrm{ContIn}(g \circ f,A,c)$
\item $\varphi_{55} :\equiv \forall c\,(\mathrm{ContIn}(f,A,c) \to \varphi_{33}) \to \mathrm{ContIn}(f,A,c) \to \varphi_{33}$
\item $\varphi_{56} :\equiv \forall u\,(\varphi_{5} \to u \in A \to u \in \mathrm{dom}(g \circ f)) \to \forall u\,(\varphi_{5}) \to \forall u \in A\,(u \in \mathrm{dom}(g \circ f))$
\item $\varphi_{57} :\equiv \forall c\,(\varphi_{5} \to c \in A \to \mathrm{ContIn}(g \circ f,A,c)) \to \forall c\,(\varphi_{5}) \to \forall c \in A\,(\mathrm{ContIn}(g \circ f,A,c))$
\item $\varphi_{58} :\equiv (\mathrm{ContOn}(f,A) \to \mathrm{ContOn}(f,A)) \to \varphi_{18} \to \mathrm{ContOn}(f,A) \to \mathrm{ContOn}(g,B) \to \varphi_{22}$
\item $\varphi_{59} :\equiv \forall B\,\forall c\,(\mathrm{ContIn}(f,A,c) \to \varphi_{33}) \to \forall c\,(\mathrm{ContIn}(f,A,c) \to \varphi_{33})$
\item $\varphi_{60} :\equiv \forall A\,\forall B\,\forall c\,(\mathrm{ContIn}(f,A,c) \to \varphi_{33}) \to \forall B\,\forall c\,(\mathrm{ContIn}(f,A,c) \to \varphi_{33})$
\item $\varphi_{61} :\equiv \forall g\,\forall A\,\forall B\,\forall c\,(\mathrm{ContIn}(f,A,c) \to \varphi_{33}) \to \forall A\,\forall B\,\forall c\,(\mathrm{ContIn}(f,A,c) \to \varphi_{33})$
\item $\varphi_{62} :\equiv \forall f\,\forall g\,\forall A\,\forall B\,\forall c\,(\mathrm{ContIn}(f,A,c) \to \varphi_{33}) \to \forall g\,\forall A\,\forall B\,\forall c\,(\mathrm{ContIn}(f,A,c) \to \varphi_{33})$
\item $\varphi_{63} :\equiv (\varphi_{5} \to \forall c\,(\varphi_{5})) \to \mathrm{ContOn}(f,A) \to \mathrm{ContOn}(g,B) \to \varphi_{16}$
\item $\varphi_{64} :\equiv (\mathrm{ContOn}(f,A) \to \varphi_{7}) \to (\mathrm{ContOn}(f,A) \to \varphi_{13}) \to \mathrm{ContOn}(f,A) \to \varphi_{23}$
\item $\varphi_{65} :\equiv (A \subseteq \mathrm{dom}(g \circ f) \leftrightarrow \forall u \in A\,(u \in \mathrm{dom}(g \circ f))) \to (\mathrm{ContOn}(f,A) \to \varphi_{36}) \to \mathrm{ContOn}(f,A) \to \varphi_{19}$
\item $\varphi_{66} :\equiv (\mathrm{ContOn}(f,A) \to \mathrm{ContOn}(g,B) \to \varphi_{22}) \to \varphi_{4} \to \mathrm{ContOn}(f,A) \to \mathrm{ContOn}(g,B) \to \varphi_{14}$
\item $\varphi_{67} :\equiv (\mathrm{ContOn}(f,A) \to \varphi_{6}) \to (\mathrm{ContOn}(f,A) \to \varphi_{27}) \to \mathrm{ContOn}(f,A) \to \varphi_{29}$
\item $\varphi_{68} :\equiv (\mathrm{ContOn}(f,A) \to \mathrm{ContOn}(g,B) \to \varphi_{14}) \to (\mathrm{ContOn}(f,A) \to \varphi_{11}) \to \mathrm{ContOn}(f,A) \to \varphi_{24}$
\item $\varphi_{69} :\equiv (\mathrm{ContOn}(f,A) \to \varphi_{23}) \to (\mathrm{ContOn}(f,A) \to \varphi_{29}) \to \mathrm{ContOn}(f,A) \to \varphi_{34}$
\item $\varphi_{70} :\equiv \forall f\,(\mathrm{RFn}(g) \to \mathrm{RFn}(f) \to \varphi_{44}) \to \mathrm{RFn}(g) \to \mathrm{RFn}(f) \to \varphi_{44}$
\item $\varphi_{71} :\equiv (\mathrm{ContOn}(f,A) \to \mathrm{ContOn}(g,B) \to \varphi_{21}) \to \varphi_{20} \to \mathrm{ContOn}(f,A) \to \mathrm{ContOn}(g,B) \to \varphi_{35}$
\item $\varphi_{72} :\equiv \forall g\,\forall f\,(\mathrm{RFn}(g) \to \mathrm{RFn}(f) \to \varphi_{44}) \to \forall f\,(\mathrm{RFn}(g) \to \mathrm{RFn}(f) \to \varphi_{44})$
\item $\varphi_{73} :\equiv (B \subseteq \mathrm{dom}(g) \to f(u) \in B \to f(u) \in \mathrm{dom}(g)) \to \varphi_{67}$
\item $\varphi_{74} :\equiv (\mathrm{ContOn}(f,A) \to \varphi_{3}) \to (\mathrm{ContOn}(f,A) \to \varphi_{2}) \to \varphi_{50}$
\item $\varphi_{75} :\equiv (\mathrm{ContOn}(f,A) \to \varphi_{9}) \to (\mathrm{ContOn}(f,A) \to \varphi_{19}) \to \varphi_{51}$
\item $\varphi_{76} :\equiv (\mathrm{ContOn}(f,A) \to \varphi_{24}) \to (\mathrm{ContOn}(f,A) \to \mathrm{ContOn}(g,B) \to \varphi_{10}) \to \varphi_{52}$
\item $\varphi_{77} :\equiv (\mathrm{ContOn}(f,A) \to \mathrm{ContOn}(g,B) \to \varphi_{35}) \to (\mathrm{ContOn}(f,A) \to \varphi_{26}) \to \varphi_{48}$
\item $\varphi_{78} :\equiv (\mathrm{ContOn}(f,A) \to \varphi_{3}) \to (\mathrm{ContOn}(f,A) \to \varphi_{2}) \to \varphi_{69}$
\item $\varphi_{79} :\equiv \forall u\,(\varphi_{5} \to u \in A \to u \in \mathrm{dom}(g \circ f)) \to \varphi_{56} \to (\varphi_{5} \to \forall u\,(\varphi_{5})) \to \mathrm{ContOn}(f,A) \to \varphi_{36}$
\end{abbrevlist}
\end{fullproof}

% 自動生成 (prover/build.js): ContPos  ノード数 310
\begin{fullproof}{定理2.7}{$\forall f\,\forall A\,\forall c\,(\mathrm{ContIn}(f,A,c) \to 0 < f(c) \to \exists \delta\,(\delta \in \mathbb{R} \land 0 < \delta \land \forall x \in A\,(\lvert x - c \rvert < \delta \to f(c) \cdot 2^{-1} < f(x))))$}{310}\phantomsection\label{tree:ContPos}
\noindent$\Pi_{1}$:
\begin{ptree}
\AXC{$\varphi_{73}$\lbl{A1}}
\AXC{$\varphi_{20} \to \varphi_{15}$\lbl{A2}}
\AXC{$\varphi_{109}$\lbl{A1}}
\AXC{$\varphi_{109} \to \mathrm{ContIn}(f,A,c) \to \varphi_{106}$\lbl{A1}}
\BIC{$\mathrm{ContIn}(f,A,c) \to \varphi_{106}$}
\AXC{$\varphi_{143}$\lbl{A1}}
\BIC{$(\varphi_{20} \to \varphi_{15}) \to \mathrm{ContIn}(f,A,c) \to \varphi_{104}$}
\BIC{$\mathrm{ContIn}(f,A,c) \to \varphi_{104}$}
\AXC{$\varphi_{147}$\lbl{A1}}
\BIC{$\varphi_{73} \to (\mathrm{ContIn}(f,A,c) \to \varphi_{79}) \to \varphi_{93}$}
\BIC{$(\mathrm{ContIn}(f,A,c) \to \varphi_{79}) \to \varphi_{93}$}
\end{ptree}
\noindent$\Pi_{2}$:
\begin{ptree}
\AXC{$\mathrm{ContIn}(f,A,c) \to \varphi_{79}$\lbl{A1}}
\AXC{$\stackrel{\Pi_{1}}{\vdots}$}\noLine
\UIC{$(\mathrm{ContIn}(f,A,c) \to \varphi_{79}) \to \varphi_{93}$}
\BIC{$\varphi_{93}$}
\end{ptree}
\noindent$\Pi_{3}$:
\begin{ptree}
\AXC{$\forall f\,\forall A\,\forall c\,(\varphi_{83})$\lbl{DefContIn}}
\AXC{$\forall f\,\forall A\,\forall c\,(\varphi_{83}) \to \forall A\,\forall c\,(\varphi_{83})$\lbl{A2}}
\BIC{$\forall A\,\forall c\,(\varphi_{83})$}
\AXC{$\forall A\,\forall c\,(\varphi_{83}) \to \forall c\,(\varphi_{83})$\lbl{A2}}
\BIC{$\forall c\,(\varphi_{83})$}
\AXC{$\forall c\,(\varphi_{83}) \to \varphi_{83}$\lbl{A2}}
\BIC{$\varphi_{83}$}
\end{ptree}
\noindent$\Pi_{4}$:
\begin{ptree}
\AXC{$\stackrel{\Pi_{3}}{\vdots}$}\noLine
\UIC{$\varphi_{83}$}
\AXC{$\mathrm{ContIn}(f,A,c) \to \mathrm{ContIn}(f,A,c)$\lbl{A1}}
\AXC{$\varphi_{113}$\lbl{A1}}
\BIC{$\varphi_{83} \to \mathrm{ContIn}(f,A,c) \to 0 < f(c) \to c \in A$}
\BIC{$\mathrm{ContIn}(f,A,c) \to 0 < f(c) \to c \in A$}
\end{ptree}
\noindent$\Pi_{5}$:
\begin{ptree}
\AXC{$\stackrel{\Pi_{3}}{\vdots}$}\noLine
\UIC{$\varphi_{83}$}
\AXC{$\mathrm{ContIn}(f,A,c) \to \mathrm{ContIn}(f,A,c)$\lbl{A1}}
\AXC{$\varphi_{116}$\lbl{A1}}
\BIC{$\varphi_{97}$}
\BIC{$\mathrm{ContIn}(f,A,c) \to 0 < f(c) \to A \subseteq \mathrm{dom}(f)$}
\end{ptree}
\noindent$\Pi_{6}$:
\begin{ptree}
\AXC{$\forall X\,\forall Y\,\forall x\,(X \subseteq Y \to x \in X \to x \in Y)$\lbl{[B5]}}
\AXC{$\varphi_{11}$\lbl{A2}}
\BIC{$\forall Y\,\forall x\,(A \subseteq Y \to x \in A \to x \in Y)$}
\AXC{$\varphi_{17}$\lbl{A2}}
\BIC{$\forall x\,(A \subseteq \mathrm{dom}(f) \to x \in A \to x \in \mathrm{dom}(f))$}
\AXC{$\varphi_{23}$\lbl{A2}}
\BIC{$A \subseteq \mathrm{dom}(f) \to c \in A \to c \in \mathrm{dom}(f)$}
\end{ptree}
\noindent$\Pi_{7}$:
\begin{ptree}
\AXC{$\stackrel{\Pi_{6}}{\vdots}$}\noLine
\UIC{$A \subseteq \mathrm{dom}(f) \to c \in A \to c \in \mathrm{dom}(f)$}
\AXC{$(A \subseteq \mathrm{dom}(f) \to c \in A \to c \in \mathrm{dom}(f)) \to \varphi_{57}$\lbl{A1}}
\BIC{$\varphi_{57}$}
\end{ptree}
\noindent$\Pi_{8}$:
\begin{ptree}
\AXC{$\stackrel{\Pi_{4}}{\vdots}$}\noLine
\UIC{$\mathrm{ContIn}(f,A,c) \to 0 < f(c) \to c \in A$}
\AXC{$\stackrel{\Pi_{5}}{\vdots}$}\noLine
\UIC{$\mathrm{ContIn}(f,A,c) \to 0 < f(c) \to A \subseteq \mathrm{dom}(f)$}
\AXC{$\stackrel{\Pi_{7}}{\vdots}$}\noLine
\UIC{$\varphi_{57}$}
\BIC{$\varphi_{25}$}
\BIC{$\mathrm{ContIn}(f,A,c) \to 0 < f(c) \to c \in \mathrm{dom}(f)$}
\end{ptree}
\noindent$\Pi_{9}$:
\begin{ptree}
\AXC{$\stackrel{\Pi_{3}}{\vdots}$}\noLine
\UIC{$\varphi_{83}$}
\AXC{$\mathrm{ContIn}(f,A,c) \to \mathrm{ContIn}(f,A,c)$\lbl{A1}}
\AXC{$\varphi_{112}$\lbl{A1}}
\BIC{$\varphi_{83} \to \mathrm{ContIn}(f,A,c) \to 0 < f(c) \to \mathrm{RFn}(f)$}
\BIC{$\mathrm{ContIn}(f,A,c) \to 0 < f(c) \to \mathrm{RFn}(f)$}
\end{ptree}
\noindent$\Pi_{10}$:
\begin{ptree}
\AXC{$\forall f\,\forall x\,(\mathrm{RFn}(f) \to x \in \mathrm{dom}(f) \to f(x) \in \mathbb{R})$\lbl{[S6]}}
\AXC{$\varphi_{36}$\lbl{A2}}
\BIC{$\forall x\,(\mathrm{RFn}(f) \to x \in \mathrm{dom}(f) \to f(x) \in \mathbb{R})$}
\AXC{$\varphi_{29}$\lbl{A2}}
\BIC{$\mathrm{RFn}(f) \to c \in \mathrm{dom}(f) \to f(c) \in \mathbb{R}$}
\end{ptree}
\noindent$\Pi_{11}$:
\begin{ptree}
\AXC{$\stackrel{\Pi_{10}}{\vdots}$}\noLine
\UIC{$\mathrm{RFn}(f) \to c \in \mathrm{dom}(f) \to f(c) \in \mathbb{R}$}
\AXC{$(\mathrm{RFn}(f) \to c \in \mathrm{dom}(f) \to f(c) \in \mathbb{R}) \to \varphi_{63}$\lbl{A1}}
\BIC{$\varphi_{63}$}
\end{ptree}
\noindent$\Pi_{12}$:
\begin{ptree}
\AXC{$\stackrel{\Pi_{8}}{\vdots}$}\noLine
\UIC{$\mathrm{ContIn}(f,A,c) \to 0 < f(c) \to c \in \mathrm{dom}(f)$}
\AXC{$\stackrel{\Pi_{9}}{\vdots}$}\noLine
\UIC{$\mathrm{ContIn}(f,A,c) \to 0 < f(c) \to \mathrm{RFn}(f)$}
\AXC{$\stackrel{\Pi_{11}}{\vdots}$}\noLine
\UIC{$\varphi_{63}$}
\BIC{$\varphi_{37}$}
\BIC{$\mathrm{ContIn}(f,A,c) \to 0 < f(c) \to f(c) \in \mathbb{R}$}
\end{ptree}
\noindent$\Pi_{13}$:
\begin{ptree}
\AXC{$\forall x\,\forall y\,(x \in \mathbb{R} \to y \in \mathbb{R} \to x \cdot y \in \mathbb{R})$\lbl{(R1)}}
\AXC{$\varphi_{44}$\lbl{A2}}
\BIC{$\forall y\,(f(c) \in \mathbb{R} \to y \in \mathbb{R} \to f(c) \cdot y \in \mathbb{R})$}
\AXC{$\varphi_{52}$\lbl{A2}}
\BIC{$f(c) \in \mathbb{R} \to 2^{-1} \in \mathbb{R} \to f(c) \cdot 2^{-1} \in \mathbb{R}$}
\end{ptree}
\noindent$\Pi_{14}$:
\begin{ptree}
\AXC{$\stackrel{\Pi_{13}}{\vdots}$}\noLine
\UIC{$f(c) \in \mathbb{R} \to 2^{-1} \in \mathbb{R} \to f(c) \cdot 2^{-1} \in \mathbb{R}$}
\AXC{$(f(c) \in \mathbb{R} \to 2^{-1} \in \mathbb{R} \to f(c) \cdot 2^{-1} \in \mathbb{R}) \to \varphi_{119}$\lbl{A1}}
\BIC{$\varphi_{119}$}
\end{ptree}
\noindent$\Pi_{15}$:
\begin{ptree}
\AXC{$2^{-1} \in \mathbb{R}$\lbl{[F31]}}
\AXC{$\stackrel{\Pi_{12}}{\vdots}$}\noLine
\UIC{$\mathrm{ContIn}(f,A,c) \to 0 < f(c) \to f(c) \in \mathbb{R}$}
\AXC{$\stackrel{\Pi_{14}}{\vdots}$}\noLine
\UIC{$\varphi_{119}$}
\BIC{$2^{-1} \in \mathbb{R} \to \mathrm{ContIn}(f,A,c) \to 0 < f(c) \to \varphi_{74}$}
\BIC{$\mathrm{ContIn}(f,A,c) \to 0 < f(c) \to \varphi_{74}$}
\end{ptree}
\noindent$\Pi_{16}$:
\begin{ptree}
\AXC{$\forall X\,\forall Y\,\forall x\,(X \subseteq Y \to x \in X \to x \in Y)$\lbl{[B5]}}
\AXC{$\varphi_{11}$\lbl{A2}}
\BIC{$\forall Y\,\forall x\,(A \subseteq Y \to x \in A \to x \in Y)$}
\AXC{$\varphi_{17}$\lbl{A2}}
\BIC{$\forall x\,(A \subseteq \mathrm{dom}(f) \to x \in A \to x \in \mathrm{dom}(f))$}
\AXC{$\varphi_{24}$\lbl{A2}}
\BIC{$A \subseteq \mathrm{dom}(f) \to x \in A \to x \in \mathrm{dom}(f)$}
\end{ptree}
\noindent$\Pi_{17}$:
\begin{ptree}
\AXC{$\stackrel{\Pi_{16}}{\vdots}$}\noLine
\UIC{$A \subseteq \mathrm{dom}(f) \to x \in A \to x \in \mathrm{dom}(f)$}
\AXC{$(A \subseteq \mathrm{dom}(f) \to x \in A \to x \in \mathrm{dom}(f)) \to \varphi_{128}$\lbl{A1}}
\BIC{$\varphi_{128}$}
\end{ptree}
\noindent$\Pi_{18}$:
\begin{ptree}
\AXC{$\varphi_{73}$\lbl{A1}}
\AXC{$\stackrel{\Pi_{5}}{\vdots}$}\noLine
\UIC{$\mathrm{ContIn}(f,A,c) \to 0 < f(c) \to A \subseteq \mathrm{dom}(f)$}
\AXC{$\stackrel{\Pi_{17}}{\vdots}$}\noLine
\UIC{$\varphi_{128}$}
\BIC{$\varphi_{73} \to \mathrm{ContIn}(f,A,c) \to 0 < f(c) \to \varphi_{69}$}
\BIC{$\mathrm{ContIn}(f,A,c) \to 0 < f(c) \to \varphi_{69}$}
\end{ptree}
\noindent$\Pi_{19}$:
\begin{ptree}
\AXC{$\forall f\,\forall x\,(\mathrm{RFn}(f) \to x \in \mathrm{dom}(f) \to f(x) \in \mathbb{R})$\lbl{[S6]}}
\AXC{$\varphi_{36}$\lbl{A2}}
\BIC{$\forall x\,(\mathrm{RFn}(f) \to x \in \mathrm{dom}(f) \to f(x) \in \mathbb{R})$}
\AXC{$\varphi_{30}$\lbl{A2}}
\BIC{$\mathrm{RFn}(f) \to x \in \mathrm{dom}(f) \to f(x) \in \mathbb{R}$}
\end{ptree}
\noindent$\Pi_{20}$:
\begin{ptree}
\AXC{$\stackrel{\Pi_{19}}{\vdots}$}\noLine
\UIC{$\mathrm{RFn}(f) \to x \in \mathrm{dom}(f) \to f(x) \in \mathbb{R}$}
\AXC{$(\mathrm{RFn}(f) \to x \in \mathrm{dom}(f) \to f(x) \in \mathbb{R}) \to \varphi_{132}$\lbl{A1}}
\BIC{$\varphi_{132}$}
\end{ptree}
\noindent$\Pi_{21}$:
\begin{ptree}
\AXC{$\stackrel{\Pi_{18}}{\vdots}$}\noLine
\UIC{$\mathrm{ContIn}(f,A,c) \to 0 < f(c) \to \varphi_{69}$}
\AXC{$\stackrel{\Pi_{9}}{\vdots}$}\noLine
\UIC{$\mathrm{ContIn}(f,A,c) \to 0 < f(c) \to \mathrm{RFn}(f)$}
\AXC{$\stackrel{\Pi_{20}}{\vdots}$}\noLine
\UIC{$\varphi_{132}$}
\BIC{$\varphi_{126}$}
\BIC{$\mathrm{ContIn}(f,A,c) \to \varphi_{80}$}
\end{ptree}
\noindent$\Pi_{22}$:
\begin{ptree}
\AXC{$\stackrel{\Pi_{21}}{\vdots}$}\noLine
\UIC{$\mathrm{ContIn}(f,A,c) \to \varphi_{80}$}
\AXC{$\forall x\,\forall y\,(x \in \mathbb{R} \to y \in \mathbb{R} \to x - y \in \mathbb{R})$\lbl{[F1]}}
\AXC{$\varphi_{45}$\lbl{A2}}
\BIC{$\forall y\,(f(x) \in \mathbb{R} \to y \in \mathbb{R} \to f(x) - y \in \mathbb{R})$}
\AXC{$\varphi_{56}$\lbl{A2}}
\BIC{$f(x) \in \mathbb{R} \to f(c) \in \mathbb{R} \to f(x) - f(c) \in \mathbb{R}$}
\AXC{$\varphi_{137}$\lbl{A1}}
\BIC{$(\mathrm{ContIn}(f,A,c) \to \varphi_{80}) \to \varphi_{114}$}
\BIC{$\varphi_{114}$}
\end{ptree}
\noindent$\Pi_{23}$:
\begin{ptree}
\AXC{$\stackrel{\Pi_{12}}{\vdots}$}\noLine
\UIC{$\mathrm{ContIn}(f,A,c) \to 0 < f(c) \to f(c) \in \mathbb{R}$}
\AXC{$\stackrel{\Pi_{22}}{\vdots}$}\noLine
\UIC{$\varphi_{114}$}
\BIC{$\mathrm{ContIn}(f,A,c) \to 0 < f(c) \to \varphi_{76}$}
\end{ptree}
\noindent$\Pi_{24}$:
\begin{ptree}
\AXC{$\forall x\,\forall y\,(x \in \mathbb{R} \to y \in \mathbb{R} \to \lvert x \rvert < y \to -y < x \land x < y)$\lbl{[A11]}}
\AXC{$\varphi_{89}$\lbl{A2}}
\BIC{$\varphi_{64}$}
\AXC{$\varphi_{64} \to f(x) - f(c) \in \mathbb{R} \to \varphi_{82}$\lbl{A2}}
\BIC{$f(x) - f(c) \in \mathbb{R} \to \varphi_{82}$}
\AXC{$\varphi_{150}$\lbl{A1}}
\BIC{$(\mathrm{ContIn}(f,A,c) \to 0 < f(c) \to \varphi_{76}) \to \varphi_{144}$}
\end{ptree}
\noindent$\Pi_{25}$:
\begin{ptree}
\AXC{$\stackrel{\Pi_{23}}{\vdots}$}\noLine
\UIC{$\mathrm{ContIn}(f,A,c) \to 0 < f(c) \to \varphi_{76}$}
\AXC{$\stackrel{\Pi_{24}}{\vdots}$}\noLine
\UIC{$(\mathrm{ContIn}(f,A,c) \to 0 < f(c) \to \varphi_{76}) \to \varphi_{144}$}
\BIC{$\varphi_{144}$}
\end{ptree}
\noindent$\Pi_{26}$:
\begin{ptree}
\AXC{$\stackrel{\Pi_{2}}{\vdots}$}\noLine
\UIC{$\varphi_{93}$}
\AXC{$\stackrel{\Pi_{15}}{\vdots}$}\noLine
\UIC{$\mathrm{ContIn}(f,A,c) \to 0 < f(c) \to \varphi_{74}$}
\AXC{$\stackrel{\Pi_{25}}{\vdots}$}\noLine
\UIC{$\varphi_{144}$}
\BIC{$\varphi_{93} \to \mathrm{ContIn}(f,A,c) \to \varphi_{110}$}
\BIC{$\mathrm{ContIn}(f,A,c) \to \varphi_{110}$}
\AXC{$(\mathrm{ContIn}(f,A,c) \to \varphi_{110}) \to \varphi_{94}$\lbl{A1}}
\BIC{$\varphi_{94}$}
\end{ptree}
\noindent$\Pi_{27}$:
\begin{ptree}
\AXC{$\forall x \in \mathbb{R}\,(-x \in \mathbb{R})$\lbl{(R1)}}
\AXC{$\varphi_{38}$\lbl{A2}}
\BIC{$f(c) \cdot 2^{-1} \in \mathbb{R} \to -(f(c) \cdot 2^{-1}) \in \mathbb{R}$}
\AXC{$(f(c) \cdot 2^{-1} \in \mathbb{R} \to -(f(c) \cdot 2^{-1}) \in \mathbb{R}) \to \varphi_{129}$\lbl{A1}}
\BIC{$\varphi_{129}$}
\end{ptree}
\noindent$\Pi_{28}$:
\begin{ptree}
\AXC{$\stackrel{\Pi_{15}}{\vdots}$}\noLine
\UIC{$\mathrm{ContIn}(f,A,c) \to 0 < f(c) \to \varphi_{74}$}
\AXC{$\stackrel{\Pi_{27}}{\vdots}$}\noLine
\UIC{$\varphi_{129}$}
\BIC{$\mathrm{ContIn}(f,A,c) \to 0 < f(c) \to \varphi_{78}$}
\end{ptree}
\noindent$\Pi_{29}$:
\begin{ptree}
\AXC{$\stackrel{\Pi_{28}}{\vdots}$}\noLine
\UIC{$\mathrm{ContIn}(f,A,c) \to 0 < f(c) \to \varphi_{78}$}
\AXC{$\varphi_{10}$\lbl{(R14)}}
\AXC{$\varphi_{10} \to \varphi_{61}$\lbl{A2}}
\BIC{$\varphi_{61}$}
\AXC{$\varphi_{61} \to \varphi_{85}$\lbl{A2}}
\BIC{$\varphi_{85}$}
\AXC{$\varphi_{130}$\lbl{A2}}
\BIC{$-(f(c) \cdot 2^{-1}) \in \mathbb{R} \to f(x) - f(c) \in \mathbb{R} \to \varphi_{66}$}
\AXC{$\varphi_{152}$\lbl{A1}}
\BIC{$\varphi_{148}$}
\BIC{$(\mathrm{ContIn}(f,A,c) \to 0 < f(c) \to \varphi_{76}) \to \varphi_{138}$}
\end{ptree}
\noindent$\Pi_{30}$:
\begin{ptree}
\AXC{$\stackrel{\Pi_{23}}{\vdots}$}\noLine
\UIC{$\mathrm{ContIn}(f,A,c) \to 0 < f(c) \to \varphi_{76}$}
\AXC{$\stackrel{\Pi_{29}}{\vdots}$}\noLine
\UIC{$(\mathrm{ContIn}(f,A,c) \to 0 < f(c) \to \varphi_{76}) \to \varphi_{138}$}
\BIC{$\varphi_{138}$}
\end{ptree}
\noindent$\Pi_{31}$:
\begin{ptree}
\AXC{$y = z \to x = y \to x = z$\lbl{A6}}
\AXC{$(y = z \to x = y \to x = z) \to x = y \to y = z \to x = z$\lbl{A1}}
\BIC{$x = y \to y = z \to x = z$}
\UIC{$\forall z\,(x = y \to y = z \to x = z)$}
\UIC{$\forall y\,\forall z\,(x = y \to y = z \to x = z)$}
\UIC{$\forall x\,\forall y\,\forall z\,(x = y \to y = z \to x = z)$}
\AXC{$\forall x\,\forall y\,\forall z\,(x = y \to y = z \to x = z) \to \varphi_{31}$\lbl{A2}}
\BIC{$\varphi_{31}$}
\end{ptree}
\noindent$\Pi_{32}$:
\begin{ptree}
\AXC{$\forall x\,\forall y\,(x \in \mathbb{R} \to y \in \mathbb{R} \to x + y - y = x)$\lbl{[E38]}}
\AXC{$\forall x\,\forall y\,(x \in \mathbb{R} \to y \in \mathbb{R} \to x + y - y = x) \to \varphi_{34}$\lbl{A2}}
\BIC{$\varphi_{34}$}
\end{ptree}
\noindent$\Pi_{33}$:
\begin{ptree}
\AXC{$\stackrel{\Pi_{15}}{\vdots}$}\noLine
\UIC{$\mathrm{ContIn}(f,A,c) \to 0 < f(c) \to \varphi_{74}$}
\AXC{$\stackrel{\Pi_{32}}{\vdots}$}\noLine
\UIC{$\varphi_{34}$}
\AXC{$\varphi_{34} \to f(c) \cdot 2^{-1} \in \mathbb{R} \to f(c) \cdot 2^{-1} \in \mathbb{R} \to \varphi_{22}$\lbl{A2}}
\BIC{$f(c) \cdot 2^{-1} \in \mathbb{R} \to f(c) \cdot 2^{-1} \in \mathbb{R} \to \varphi_{22}$}
\AXC{$\varphi_{145}$\lbl{A1}}
\BIC{$(\mathrm{ContIn}(f,A,c) \to 0 < f(c) \to \varphi_{74}) \to \varphi_{133}$}
\BIC{$\varphi_{133}$}
\end{ptree}
\noindent$\Pi_{34}$:
\begin{ptree}
\AXC{$y = z \to x = y \to x = z$\lbl{A6}}
\AXC{$(y = z \to x = y \to x = z) \to x = y \to y = z \to x = z$\lbl{A1}}
\BIC{$x = y \to y = z \to x = z$}
\UIC{$\forall z\,(x = y \to y = z \to x = z)$}
\UIC{$\forall y\,\forall z\,(x = y \to y = z \to x = z)$}
\UIC{$\forall x\,\forall y\,\forall z\,(x = y \to y = z \to x = z)$}
\AXC{$\forall x\,\forall y\,\forall z\,(x = y \to y = z \to x = z) \to \varphi_{31}$\lbl{A2}}
\BIC{$\varphi_{31}$}
\end{ptree}
\noindent$\Pi_{35}$:
\begin{ptree}
\AXC{$x = x$\lbl{A5}}
\AXC{$x = y \to x = x \to y = x$\lbl{A6}}
\AXC{$(x = y \to x = x \to y = x) \to x = x \to x = y \to y = x$\lbl{A1}}
\BIC{$x = x \to x = y \to y = x$}
\BIC{$x = y \to y = x$}
\UIC{$\forall y\,(x = y \to y = x)$}
\UIC{$\forall x\,\forall y\,(x = y \to y = x)$}
\end{ptree}
\noindent$\Pi_{36}$:
\begin{ptree}
\AXC{$\stackrel{\Pi_{12}}{\vdots}$}\noLine
\UIC{$\mathrm{ContIn}(f,A,c) \to 0 < f(c) \to f(c) \in \mathbb{R}$}
\AXC{$\forall e \in \mathbb{R}\,(e \cdot 2^{-1} + e \cdot 2^{-1} = e)$\lbl{[F36]}}
\AXC{$\varphi_{49}$\lbl{A2}}
\BIC{$f(c) \in \mathbb{R} \to f(c) \cdot 2^{-1} + f(c) \cdot 2^{-1} = f(c)$}
\AXC{$\varphi_{124}$\lbl{A1}}
\BIC{$\varphi_{120}$}
\BIC{$\varphi_{95}$}
\end{ptree}
\noindent$\Pi_{37}$:
\begin{ptree}
\AXC{$\stackrel{\Pi_{35}}{\vdots}$}\noLine
\UIC{$\forall x\,\forall y\,(x = y \to y = x)$}
\AXC{$\forall x\,\forall y\,(x = y \to y = x) \to \varphi_{32}$\lbl{A2}}
\BIC{$\varphi_{32}$}
\AXC{$\varphi_{32} \to \varphi_{40}$\lbl{A2}}
\BIC{$\varphi_{40}$}
\AXC{$\stackrel{\Pi_{36}}{\vdots}$}\noLine
\UIC{$\varphi_{95}$}
\AXC{$\varphi_{95} \to \varphi_{40} \to \varphi_{96}$\lbl{A1}}
\BIC{$\varphi_{40} \to \varphi_{96}$}
\BIC{$\varphi_{96}$}
\end{ptree}
\noindent$\Pi_{38}$:
\begin{ptree}
\AXC{$x = x$\lbl{A5}}
\UIC{$\forall x\,(x = x)$}
\AXC{$\varphi_{5}$\lbl{A2}}
\BIC{$f(c) - f(c) \cdot 2^{-1} = f(c) - f(c) \cdot 2^{-1}$}
\AXC{$\stackrel{\Pi_{37}}{\vdots}$}\noLine
\UIC{$\varphi_{96}$}
\AXC{$\varphi_{96} \to \varphi_{140}$\lbl{A1}}
\BIC{$\varphi_{140}$}
\BIC{$\varphi_{91} \to \mathrm{ContIn}(f,A,c) \to \varphi_{111}$}
\end{ptree}
\noindent$\Pi_{39}$:
\begin{ptree}
\AXC{$z = w \to \varphi_{48}$\lbl{A6}}
\UIC{$\forall w\,(z = w \to \varphi_{48})$}
\UIC{$\forall z\,\forall w\,(z = w \to \varphi_{48})$}
\AXC{$\forall z\,\forall w\,(z = w \to \varphi_{48}) \to \varphi_{58}$\lbl{A2}}
\BIC{$\varphi_{58}$}
\AXC{$\varphi_{58} \to \varphi_{91}$\lbl{A2}}
\BIC{$\varphi_{91}$}
\AXC{$\stackrel{\Pi_{38}}{\vdots}$}\noLine
\UIC{$\varphi_{91} \to \mathrm{ContIn}(f,A,c) \to \varphi_{111}$}
\BIC{$\mathrm{ContIn}(f,A,c) \to \varphi_{111}$}
\end{ptree}
\noindent$\Pi_{40}$:
\begin{ptree}
\AXC{$y = z \to x = y \to x = z$\lbl{A6}}
\AXC{$(y = z \to x = y \to x = z) \to x = y \to y = z \to x = z$\lbl{A1}}
\BIC{$x = y \to y = z \to x = z$}
\UIC{$\forall z\,(x = y \to y = z \to x = z)$}
\UIC{$\forall y\,\forall z\,(x = y \to y = z \to x = z)$}
\UIC{$\forall x\,\forall y\,\forall z\,(x = y \to y = z \to x = z)$}
\AXC{$\forall x\,\forall y\,\forall z\,(x = y \to y = z \to x = z) \to \varphi_{31}$\lbl{A2}}
\BIC{$\varphi_{31}$}
\end{ptree}
\noindent$\Pi_{41}$:
\begin{ptree}
\AXC{$\stackrel{\Pi_{35}}{\vdots}$}\noLine
\UIC{$\forall x\,\forall y\,(x = y \to y = x)$}
\AXC{$\forall x\,\forall y\,(x = y \to y = x) \to \varphi_{6}$\lbl{A2}}
\BIC{$\varphi_{6}$}
\AXC{$\varphi_{6} \to \varphi_{59}$\lbl{A2}}
\BIC{$\varphi_{59}$}
\end{ptree}
\noindent$\Pi_{42}$:
\begin{ptree}
\AXC{$\forall x\,\forall y\,(x \in \mathbb{R} \to y \in \mathbb{R} \to x + y = y + x)$\lbl{(R3)}}
\AXC{$\forall x\,\forall y\,(x \in \mathbb{R} \to y \in \mathbb{R} \to x + y = y + x) \to \varphi_{46}$\lbl{A2}}
\BIC{$\varphi_{46}$}
\end{ptree}
\noindent$\Pi_{43}$:
\begin{ptree}
\AXC{$\stackrel{\Pi_{28}}{\vdots}$}\noLine
\UIC{$\mathrm{ContIn}(f,A,c) \to 0 < f(c) \to \varphi_{78}$}
\AXC{$\stackrel{\Pi_{42}}{\vdots}$}\noLine
\UIC{$\varphi_{46}$}
\AXC{$\varphi_{46} \to -(f(c) \cdot 2^{-1}) \in \mathbb{R} \to f(c) \in \mathbb{R} \to \varphi_{8}$\lbl{A2}}
\BIC{$-(f(c) \cdot 2^{-1}) \in \mathbb{R} \to f(c) \in \mathbb{R} \to \varphi_{8}$}
\AXC{$\varphi_{141}$\lbl{A1}}
\BIC{$(\mathrm{ContIn}(f,A,c) \to 0 < f(c) \to \varphi_{78}) \to \varphi_{122}$}
\BIC{$\varphi_{122}$}
\end{ptree}
\noindent$\Pi_{44}$:
\begin{ptree}
\AXC{$\stackrel{\Pi_{12}}{\vdots}$}\noLine
\UIC{$\mathrm{ContIn}(f,A,c) \to 0 < f(c) \to f(c) \in \mathbb{R}$}
\AXC{$\stackrel{\Pi_{43}}{\vdots}$}\noLine
\UIC{$\varphi_{122}$}
\BIC{$\mathrm{ContIn}(f,A,c) \to \varphi_{101}$}
\AXC{$(\mathrm{ContIn}(f,A,c) \to \varphi_{101}) \to \varphi_{139}$\lbl{A1}}
\BIC{$\varphi_{139}$}
\end{ptree}
\noindent$\Pi_{45}$:
\begin{ptree}
\AXC{$\stackrel{\Pi_{41}}{\vdots}$}\noLine
\UIC{$\varphi_{59}$}
\AXC{$\forall x\,\forall y\,(x - y = x + -y)$\lbl{Def$-$}}
\AXC{$\varphi_{4}$\lbl{A2}}
\BIC{$\forall y\,(f(c) - y = f(c) + -y)$}
\AXC{$\varphi_{43}$\lbl{A2}}
\BIC{$f(c) - f(c) \cdot 2^{-1} = f(c) + -(f(c) \cdot 2^{-1})$}
\AXC{$\varphi_{121}$\lbl{A1}}
\BIC{$\varphi_{87}$}
\BIC{$f(c) + -(f(c) \cdot 2^{-1}) = f(c) - f(c) \cdot 2^{-1}$}
\AXC{$\stackrel{\Pi_{44}}{\vdots}$}\noLine
\UIC{$\varphi_{139}$}
\BIC{$(\varphi_{8} \to \varphi_{60}) \to \varphi_{108}$}
\end{ptree}
\noindent$\Pi_{46}$:
\begin{ptree}
\AXC{$\stackrel{\Pi_{39}}{\vdots}$}\noLine
\UIC{$\mathrm{ContIn}(f,A,c) \to \varphi_{111}$}
\AXC{$\stackrel{\Pi_{40}}{\vdots}$}\noLine
\UIC{$\varphi_{31}$}
\AXC{$\varphi_{31} \to \varphi_{71}$\lbl{A2}}
\BIC{$\varphi_{71}$}
\AXC{$\varphi_{71} \to \varphi_{8} \to \varphi_{60}$\lbl{A2}}
\BIC{$\varphi_{8} \to \varphi_{60}$}
\AXC{$\stackrel{\Pi_{45}}{\vdots}$}\noLine
\UIC{$(\varphi_{8} \to \varphi_{60}) \to \varphi_{108}$}
\BIC{$\varphi_{108}$}
\AXC{$\varphi_{108} \to \varphi_{146}$\lbl{A1}}
\BIC{$\varphi_{146}$}
\BIC{$\varphi_{118} \to \mathrm{ContIn}(f,A,c) \to \varphi_{115}$}
\end{ptree}
\noindent$\Pi_{47}$:
\begin{ptree}
\AXC{$\stackrel{\Pi_{34}}{\vdots}$}\noLine
\UIC{$\varphi_{31}$}
\AXC{$\varphi_{31} \to \varphi_{67}$\lbl{A2}}
\BIC{$\varphi_{67}$}
\AXC{$\varphi_{67} \to \varphi_{118}$\lbl{A2}}
\BIC{$\varphi_{118}$}
\AXC{$\stackrel{\Pi_{46}}{\vdots}$}\noLine
\UIC{$\varphi_{118} \to \mathrm{ContIn}(f,A,c) \to \varphi_{115}$}
\BIC{$\mathrm{ContIn}(f,A,c) \to \varphi_{115}$}
\AXC{$\varphi_{149}$\lbl{A1}}
\BIC{$(\mathrm{ContIn}(f,A,c) \to \varphi_{107}) \to \varphi_{135}$}
\end{ptree}
\noindent$\Pi_{48}$:
\begin{ptree}
\AXC{$\stackrel{\Pi_{15}}{\vdots}$}\noLine
\UIC{$\mathrm{ContIn}(f,A,c) \to 0 < f(c) \to \varphi_{74}$}
\AXC{$\stackrel{\Pi_{33}}{\vdots}$}\noLine
\UIC{$\varphi_{133}$}
\BIC{$\mathrm{ContIn}(f,A,c) \to \varphi_{107}$}
\AXC{$\stackrel{\Pi_{47}}{\vdots}$}\noLine
\UIC{$(\mathrm{ContIn}(f,A,c) \to \varphi_{107}) \to \varphi_{135}$}
\BIC{$\varphi_{135}$}
\end{ptree}
\noindent$\Pi_{49}$:
\begin{ptree}
\AXC{$\stackrel{\Pi_{31}}{\vdots}$}\noLine
\UIC{$\varphi_{31}$}
\AXC{$\varphi_{31} \to \varphi_{90}$\lbl{A2}}
\BIC{$\varphi_{90}$}
\AXC{$\varphi_{134}$\lbl{A2}}
\BIC{$\varphi_{39} \to \varphi_{22} \to -(f(c) \cdot 2^{-1}) + f(c) = f(c) \cdot 2^{-1}$}
\AXC{$\stackrel{\Pi_{48}}{\vdots}$}\noLine
\UIC{$\varphi_{135}$}
\BIC{$\varphi_{98}$}
\AXC{$\varphi_{98} \to (\mathrm{ContIn}(f,A,c) \to \varphi_{100}) \to \varphi_{84} \to \varphi_{99}$\lbl{A1}}
\BIC{$(\mathrm{ContIn}(f,A,c) \to \varphi_{100}) \to \varphi_{84} \to \varphi_{99}$}
\end{ptree}
\noindent$\Pi_{50}$:
\begin{ptree}
\AXC{$\stackrel{\Pi_{26}}{\vdots}$}\noLine
\UIC{$\varphi_{94}$}
\AXC{$\stackrel{\Pi_{12}}{\vdots}$}\noLine
\UIC{$\mathrm{ContIn}(f,A,c) \to 0 < f(c) \to f(c) \in \mathbb{R}$}
\AXC{$\stackrel{\Pi_{30}}{\vdots}$}\noLine
\UIC{$\varphi_{138}$}
\BIC{$\varphi_{94} \to \mathrm{ContIn}(f,A,c) \to \varphi_{100}$}
\BIC{$\mathrm{ContIn}(f,A,c) \to \varphi_{100}$}
\AXC{$\stackrel{\Pi_{49}}{\vdots}$}\noLine
\UIC{$(\mathrm{ContIn}(f,A,c) \to \varphi_{100}) \to \varphi_{84} \to \varphi_{99}$}
\BIC{$\varphi_{84} \to \varphi_{99}$}
\end{ptree}
\noindent$\Pi_{51}$:
\begin{ptree}
\AXC{$\forall x\,\forall y\,(x \in \mathbb{R} \to y \in \mathbb{R} \to x - y + y = x)$\lbl{[F12]}}
\AXC{$\varphi_{47}$\lbl{A2}}
\BIC{$\forall y\,(f(x) \in \mathbb{R} \to y \in \mathbb{R} \to f(x) - y + y = f(x))$}
\AXC{$\varphi_{68}$\lbl{A2}}
\BIC{$\varphi_{21}$}
\AXC{$\varphi_{21} \to (\mathrm{ContIn}(f,A,c) \to \varphi_{80}) \to \varphi_{117}$\lbl{A1}}
\BIC{$(\mathrm{ContIn}(f,A,c) \to \varphi_{80}) \to \varphi_{117}$}
\end{ptree}
\noindent$\Pi_{52}$:
\begin{ptree}
\AXC{$\stackrel{\Pi_{12}}{\vdots}$}\noLine
\UIC{$\mathrm{ContIn}(f,A,c) \to 0 < f(c) \to f(c) \in \mathbb{R}$}
\AXC{$\stackrel{\Pi_{21}}{\vdots}$}\noLine
\UIC{$\mathrm{ContIn}(f,A,c) \to \varphi_{80}$}
\AXC{$\stackrel{\Pi_{51}}{\vdots}$}\noLine
\UIC{$(\mathrm{ContIn}(f,A,c) \to \varphi_{80}) \to \varphi_{117}$}
\BIC{$\varphi_{117}$}
\BIC{$\varphi_{92}$}
\end{ptree}
\noindent$\Pi_{53}$:
\begin{ptree}
\AXC{$z = w \to \varphi_{14}$\lbl{A6}}
\UIC{$\forall w\,(z = w \to \varphi_{14})$}
\UIC{$\forall z\,\forall w\,(z = w \to \varphi_{14})$}
\AXC{$\forall z\,\forall w\,(z = w \to \varphi_{14}) \to \varphi_{70}$\lbl{A2}}
\BIC{$\varphi_{70}$}
\AXC{$\varphi_{70} \to \varphi_{84}$\lbl{A2}}
\BIC{$\varphi_{84}$}
\AXC{$\stackrel{\Pi_{50}}{\vdots}$}\noLine
\UIC{$\varphi_{84} \to \varphi_{99}$}
\BIC{$\varphi_{99}$}
\AXC{$\stackrel{\Pi_{52}}{\vdots}$}\noLine
\UIC{$\varphi_{92}$}
\AXC{$\varphi_{92} \to \varphi_{99} \to \varphi_{123}$\lbl{A1}}
\BIC{$\varphi_{99} \to \varphi_{123}$}
\BIC{$\varphi_{123}$}
\end{ptree}
\noindent$\Pi_{54}$:
\begin{ptree}
\AXC{$z = w \to f(c) \cdot 2^{-1} < z \to f(c) \cdot 2^{-1} < w$\lbl{A6}}
\UIC{$\forall w\,(z = w \to f(c) \cdot 2^{-1} < z \to f(c) \cdot 2^{-1} < w)$}
\UIC{$\varphi_{3}$}
\AXC{$\varphi_{3} \to \varphi_{50}$\lbl{A2}}
\BIC{$\varphi_{50}$}
\AXC{$\varphi_{50} \to f(x) - f(c) + f(c) = f(x) \to \varphi_{27}$\lbl{A2}}
\BIC{$f(x) - f(c) + f(c) = f(x) \to \varphi_{27}$}
\AXC{$\stackrel{\Pi_{53}}{\vdots}$}\noLine
\UIC{$\varphi_{123}$}
\BIC{$\mathrm{ContIn}(f,A,c) \to 0 < f(c) \to \varphi_{77}$}
\end{ptree}
\noindent$\Pi_{55}$:
\begin{ptree}
\AXC{$\varphi_{62} \to \forall x\,(\varphi_{62})$\lbl{A4}}
\AXC{$\varphi_{86} \to \varphi_{88}$\lbl{A3}}
\AXC{$\stackrel{\Pi_{54}}{\vdots}$}\noLine
\UIC{$\mathrm{ContIn}(f,A,c) \to 0 < f(c) \to \varphi_{77}$}
\AXC{$\varphi_{131}$\lbl{A1}}
\BIC{$\varphi_{62} \to x \in A \to \lvert x - c \rvert < \delta \to f(c) \cdot 2^{-1} < f(x)$}
\UIC{$\varphi_{86}$}
\AXC{$\varphi_{86} \to (\varphi_{86} \to \varphi_{88}) \to \varphi_{136}$\lbl{A1}}
\BIC{$(\varphi_{86} \to \varphi_{88}) \to \varphi_{136}$}
\BIC{$\varphi_{136}$}
\BIC{$\mathrm{ContIn}(f,A,c) \to 0 < f(c) \to \varphi_{81}$}
\end{ptree}
\noindent$\Pi_{56}$:
\begin{ptree}
\AXC{$\stackrel{\Pi_{55}}{\vdots}$}\noLine
\UIC{$\mathrm{ContIn}(f,A,c) \to 0 < f(c) \to \varphi_{81}$}
\AXC{$\varphi_{109}$\lbl{A1}}
\AXC{$\varphi_{109} \to (\mathrm{ContIn}(f,A,c) \to 0 < f(c) \to \varphi_{81}) \to \varphi_{102}$\lbl{A1}}
\BIC{$(\mathrm{ContIn}(f,A,c) \to 0 < f(c) \to \varphi_{81}) \to \varphi_{102}$}
\BIC{$\varphi_{102}$}
\end{ptree}
\noindent$\Pi_{57}$:
\begin{ptree}
\AXC{$\forall e \in \mathbb{R}\,(0 < e \to 0 < e \cdot 2^{-1})$\lbl{[F35]}}
\AXC{$\varphi_{41}$\lbl{A2}}
\BIC{$f(c) \in \mathbb{R} \to 0 < f(c) \to 0 < f(c) \cdot 2^{-1}$}
\AXC{$(f(c) \in \mathbb{R} \to 0 < f(c) \to 0 < f(c) \cdot 2^{-1}) \to \varphi_{72}$\lbl{A1}}
\BIC{$\varphi_{72}$}
\end{ptree}
\noindent$\Pi_{58}$:
\begin{ptree}
\AXC{$\mathrm{ContIn}(f,A,c) \to 0 < f(c) \to 0 < f(c)$\lbl{A1}}
\AXC{$\stackrel{\Pi_{12}}{\vdots}$}\noLine
\UIC{$\mathrm{ContIn}(f,A,c) \to 0 < f(c) \to f(c) \in \mathbb{R}$}
\AXC{$\stackrel{\Pi_{57}}{\vdots}$}\noLine
\UIC{$\varphi_{72}$}
\BIC{$\varphi_{42}$}
\BIC{$\varphi_{1}$}
\end{ptree}
\noindent$\Pi_{59}$:
\begin{ptree}
\AXC{$\forall x\,\forall y\,(x \in \mathbb{R} \to y \in \mathbb{R} \to x \cdot y \in \mathbb{R})$\lbl{(R1)}}
\AXC{$\varphi_{44}$\lbl{A2}}
\BIC{$\forall y\,(f(c) \in \mathbb{R} \to y \in \mathbb{R} \to f(c) \cdot y \in \mathbb{R})$}
\AXC{$\varphi_{52}$\lbl{A2}}
\BIC{$f(c) \in \mathbb{R} \to 2^{-1} \in \mathbb{R} \to f(c) \cdot 2^{-1} \in \mathbb{R}$}
\end{ptree}
\noindent$\Pi_{60}$:
\begin{ptree}
\AXC{$\stackrel{\Pi_{59}}{\vdots}$}\noLine
\UIC{$f(c) \in \mathbb{R} \to 2^{-1} \in \mathbb{R} \to f(c) \cdot 2^{-1} \in \mathbb{R}$}
\AXC{$(f(c) \in \mathbb{R} \to 2^{-1} \in \mathbb{R} \to f(c) \cdot 2^{-1} \in \mathbb{R}) \to \varphi_{54}$\lbl{A1}}
\BIC{$\varphi_{54}$}
\end{ptree}
\noindent$\Pi_{61}$:
\begin{ptree}
\AXC{$\varphi_{51} \to \varphi_{75}$\lbl{A2}}
\AXC{$\stackrel{\Pi_{3}}{\vdots}$}\noLine
\UIC{$\varphi_{83}$}
\AXC{$\mathrm{ContIn}(f,A,c) \to \mathrm{ContIn}(f,A,c)$\lbl{A1}}
\AXC{$\varphi_{125}$\lbl{A1}}
\BIC{$\varphi_{83} \to \mathrm{ContIn}(f,A,c) \to 0 < f(c) \to \varphi_{51}$}
\BIC{$\mathrm{ContIn}(f,A,c) \to 0 < f(c) \to \varphi_{51}$}
\AXC{$\varphi_{142}$\lbl{A1}}
\BIC{$(\varphi_{51} \to \varphi_{75}) \to \mathrm{ContIn}(f,A,c) \to 0 < f(c) \to \varphi_{75}$}
\BIC{$\mathrm{ContIn}(f,A,c) \to 0 < f(c) \to \varphi_{75}$}
\end{ptree}
\noindent$\Pi_{62}$:
\begin{ptree}
\AXC{$\stackrel{\Pi_{61}}{\vdots}$}\noLine
\UIC{$\mathrm{ContIn}(f,A,c) \to 0 < f(c) \to \varphi_{75}$}
\AXC{$(\mathrm{ContIn}(f,A,c) \to 0 < f(c) \to \varphi_{75}) \to \varphi_{2} \to \varphi_{1} \to \varphi_{65}$\lbl{A1}}
\BIC{$\varphi_{2} \to \varphi_{1} \to \varphi_{65}$}
\end{ptree}
\noindent$\Pi_{63}$:
\begin{ptree}
\AXC{$\stackrel{\Pi_{58}}{\vdots}$}\noLine
\UIC{$\varphi_{1}$}
\AXC{$2^{-1} \in \mathbb{R}$\lbl{[F31]}}
\AXC{$\stackrel{\Pi_{12}}{\vdots}$}\noLine
\UIC{$\mathrm{ContIn}(f,A,c) \to 0 < f(c) \to f(c) \in \mathbb{R}$}
\AXC{$\stackrel{\Pi_{60}}{\vdots}$}\noLine
\UIC{$\varphi_{54}$}
\BIC{$2^{-1} \in \mathbb{R} \to \varphi_{2}$}
\BIC{$\varphi_{2}$}
\AXC{$\stackrel{\Pi_{62}}{\vdots}$}\noLine
\UIC{$\varphi_{2} \to \varphi_{1} \to \varphi_{65}$}
\BIC{$\varphi_{1} \to \varphi_{65}$}
\BIC{$\varphi_{65}$}
\end{ptree}
\noindent$\Pi_{64}$:
\begin{ptree}
\AXC{$\forall \delta\,(\neg \varphi_{28}) \to \neg \varphi_{28}$\lbl{A2}}
\AXC{$\stackrel{\Pi_{56}}{\vdots}$}\noLine
\UIC{$\varphi_{102}$}
\AXC{$\varphi_{102} \to (\forall \delta\,(\neg \varphi_{28}) \to \neg \varphi_{28}) \to \varphi_{103}$\lbl{A1}}
\BIC{$(\forall \delta\,(\neg \varphi_{28}) \to \neg \varphi_{28}) \to \varphi_{103}$}
\BIC{$\varphi_{103}$}
\AXC{$\varphi_{103} \to \varphi_{105}$\lbl{A1}}
\BIC{$\varphi_{105}$}
\UIC{$\forall \delta\,(\varphi_{105})$}
\AXC{$\stackrel{\Pi_{63}}{\vdots}$}\noLine
\UIC{$\varphi_{65}$}
\AXC{$\varphi_{65} \to \varphi_{151}$\lbl{A1}}
\BIC{$\varphi_{151}$}
\BIC{$(\forall \delta\,(\varphi_{105}) \to \varphi_{53} \to \forall \delta\,(\neg (\delta \in \mathbb{R} \land 0 < \delta \land \varphi_{20}))) \to \varphi_{127}$}
\end{ptree}
\noindent$\Pi_{65}$:
\begin{ptree}
\AXC{$\forall \delta\,(\varphi_{105}) \to \varphi_{53} \to \forall \delta\,(\neg (\delta \in \mathbb{R} \land 0 < \delta \land \varphi_{20}))$\lbl{A3}}
\AXC{$\stackrel{\Pi_{64}}{\vdots}$}\noLine
\UIC{$(\forall \delta\,(\varphi_{105}) \to \varphi_{53} \to \forall \delta\,(\neg (\delta \in \mathbb{R} \land 0 < \delta \land \varphi_{20}))) \to \varphi_{127}$}
\BIC{$\varphi_{127}$}
\end{ptree}
\noindent 主木:
\begin{ptree}
\AXC{$\mathrm{ContIn}(f,A,c) \land 0 < f(c) \land \neg \exists \delta\,(\varphi_{28}) \to \varphi_{53}$\lbl{A4}}
\AXC{$\stackrel{\Pi_{65}}{\vdots}$}\noLine
\UIC{$\varphi_{127}$}
\BIC{$\mathrm{ContIn}(f,A,c) \to 0 < f(c) \to \exists \delta\,(\varphi_{28})$}
\UIC{$\forall c\,(\mathrm{ContIn}(f,A,c) \to 0 < f(c) \to \exists \delta\,(\varphi_{28}))$}
\UIC{$\varphi_{55}$}
\UIC{$\forall f\,(\varphi_{55})$}
\end{ptree}
\begin{abbrevlist}
\item $\varphi_{1} :\equiv \mathrm{ContIn}(f,A,c) \to 0 < f(c) \to 0 < f(c) \cdot 2^{-1}$
\item $\varphi_{2} :\equiv \mathrm{ContIn}(f,A,c) \to 0 < f(c) \to f(c) \cdot 2^{-1} \in \mathbb{R}$
\item $\varphi_{3} :\equiv \forall z\,\forall w\,(z = w \to f(c) \cdot 2^{-1} < z \to f(c) \cdot 2^{-1} < w)$
\item $\varphi_{4} :\equiv \forall x\,\forall y\,(x - y = x + -y) \to \forall y\,(f(c) - y = f(c) + -y)$
\item $\varphi_{5} :\equiv \forall x\,(x = x) \to f(c) - f(c) \cdot 2^{-1} = f(c) - f(c) \cdot 2^{-1}$
\item $\varphi_{6} :\equiv \forall y\,(f(c) - f(c) \cdot 2^{-1} = y \to y = f(c) - f(c) \cdot 2^{-1})$
\item $\varphi_{7} :\equiv -(f(c) \cdot 2^{-1}) + f(c) < f(x) - f(c) + f(c)$
\item $\varphi_{8} :\equiv -(f(c) \cdot 2^{-1}) + f(c) = f(c) + -(f(c) \cdot 2^{-1})$
\item $\varphi_{9} :\equiv x \in A \to \lvert x - c \rvert < \delta \to f(x) - f(c) + f(c) = f(x)$
\item $\varphi_{10} :\equiv \forall x\,\forall y\,\forall z\,(x \in \mathbb{R} \to y \in \mathbb{R} \to z \in \mathbb{R} \to x < y \to x + z < y + z)$
\item $\varphi_{11} :\equiv \forall X\,\forall Y\,\forall x\,(X \subseteq Y \to x \in X \to x \in Y) \to \forall Y\,\forall x\,(A \subseteq Y \to x \in A \to x \in Y)$
\item $\varphi_{12} :\equiv \lvert x - c \rvert < \delta \to -(f(c) \cdot 2^{-1}) + f(c) = f(c) \cdot 2^{-1}$
\item $\varphi_{13} :\equiv \lvert x - c \rvert < \delta \to f(c) \cdot 2^{-1} < f(x) - f(c) + f(c)$
\item $\varphi_{14} :\equiv z < f(x) - f(c) + f(c) \to w < f(x) - f(c) + f(c)$
\item $\varphi_{15} :\equiv x \in A \to \lvert x - c \rvert < \delta \to \lvert f(x) - f(c) \rvert < f(c) \cdot 2^{-1}$
\item $\varphi_{16} :\equiv x \in A \to \lvert x - c \rvert < \delta \to -(f(c) \cdot 2^{-1}) < f(x) - f(c)$
\item $\varphi_{17} :\equiv \forall Y\,\forall x\,(A \subseteq Y \to x \in A \to x \in Y) \to \forall x\,(A \subseteq \mathrm{dom}(f) \to x \in A \to x \in \mathrm{dom}(f))$
\item $\varphi_{18} :\equiv x \in A \to \lvert x - c \rvert < \delta \to f(c) \cdot 2^{-1} + f(c) \cdot 2^{-1} = f(c)$
\item $\varphi_{19} :\equiv x \in A \to \lvert x - c \rvert < \delta \to f(c) = f(c) \cdot 2^{-1} + f(c) \cdot 2^{-1}$
\item $\varphi_{20} :\equiv \forall x \in A\,(\lvert x - c \rvert < \delta \to \lvert f(x) - f(c) \rvert < f(c) \cdot 2^{-1})$
\item $\varphi_{21} :\equiv f(x) \in \mathbb{R} \to f(c) \in \mathbb{R} \to f(x) - f(c) + f(c) = f(x)$
\item $\varphi_{22} :\equiv f(c) \cdot 2^{-1} + f(c) \cdot 2^{-1} - f(c) \cdot 2^{-1} = f(c) \cdot 2^{-1}$
\item $\varphi_{23} :\equiv \forall x\,(A \subseteq \mathrm{dom}(f) \to x \in A \to x \in \mathrm{dom}(f)) \to A \subseteq \mathrm{dom}(f) \to c \in A \to c \in \mathrm{dom}(f)$
\item $\varphi_{24} :\equiv \forall x\,(A \subseteq \mathrm{dom}(f) \to x \in A \to x \in \mathrm{dom}(f)) \to A \subseteq \mathrm{dom}(f) \to x \in A \to x \in \mathrm{dom}(f)$
\item $\varphi_{25} :\equiv (\mathrm{ContIn}(f,A,c) \to 0 < f(c) \to c \in A) \to \mathrm{ContIn}(f,A,c) \to 0 < f(c) \to c \in \mathrm{dom}(f)$
\item $\varphi_{26} :\equiv \lvert x - c \rvert < \delta \to -(f(c) \cdot 2^{-1}) + f(c) = f(c) - f(c) \cdot 2^{-1}$
\item $\varphi_{27} :\equiv f(c) \cdot 2^{-1} < f(x) - f(c) + f(c) \to f(c) \cdot 2^{-1} < f(x)$
\item $\varphi_{28} :\equiv \delta \in \mathbb{R} \land 0 < \delta \land \forall x \in A\,(\lvert x - c \rvert < \delta \to f(c) \cdot 2^{-1} < f(x))$
\item $\varphi_{29} :\equiv \forall x\,(\mathrm{RFn}(f) \to x \in \mathrm{dom}(f) \to f(x) \in \mathbb{R}) \to \mathrm{RFn}(f) \to c \in \mathrm{dom}(f) \to f(c) \in \mathbb{R}$
\item $\varphi_{30} :\equiv \forall x\,(\mathrm{RFn}(f) \to x \in \mathrm{dom}(f) \to f(x) \in \mathbb{R}) \to \mathrm{RFn}(f) \to x \in \mathrm{dom}(f) \to f(x) \in \mathbb{R}$
\item $\varphi_{31} :\equiv \forall y\,\forall z\,(-(f(c) \cdot 2^{-1}) + f(c) = y \to y = z \to -(f(c) \cdot 2^{-1}) + f(c) = z)$
\item $\varphi_{32} :\equiv \forall y\,(f(c) \cdot 2^{-1} + f(c) \cdot 2^{-1} = y \to y = f(c) \cdot 2^{-1} + f(c) \cdot 2^{-1})$
\item $\varphi_{33} :\equiv -(f(c) \cdot 2^{-1}) < f(x) - f(c) \land f(x) - f(c) < f(c) \cdot 2^{-1}$
\item $\varphi_{34} :\equiv \forall y\,(f(c) \cdot 2^{-1} \in \mathbb{R} \to y \in \mathbb{R} \to f(c) \cdot 2^{-1} + y - y = f(c) \cdot 2^{-1})$
\item $\varphi_{35} :\equiv f(c) - f(c) \cdot 2^{-1} = f(c) \cdot 2^{-1} + f(c) \cdot 2^{-1} - f(c) \cdot 2^{-1}$
\item $\varphi_{36} :\equiv \forall f\,\forall x\,(\mathrm{RFn}(f) \to x \in \mathrm{dom}(f) \to f(x) \in \mathbb{R}) \to \forall x\,(\mathrm{RFn}(f) \to x \in \mathrm{dom}(f) \to f(x) \in \mathbb{R})$
\item $\varphi_{37} :\equiv (\mathrm{ContIn}(f,A,c) \to 0 < f(c) \to c \in \mathrm{dom}(f)) \to \mathrm{ContIn}(f,A,c) \to 0 < f(c) \to f(c) \in \mathbb{R}$
\item $\varphi_{38} :\equiv \forall x \in \mathbb{R}\,(-x \in \mathbb{R}) \to f(c) \cdot 2^{-1} \in \mathbb{R} \to -(f(c) \cdot 2^{-1}) \in \mathbb{R}$
\item $\varphi_{39} :\equiv -(f(c) \cdot 2^{-1}) + f(c) = f(c) \cdot 2^{-1} + f(c) \cdot 2^{-1} - f(c) \cdot 2^{-1}$
\item $\varphi_{40} :\equiv f(c) \cdot 2^{-1} + f(c) \cdot 2^{-1} = f(c) \to f(c) = f(c) \cdot 2^{-1} + f(c) \cdot 2^{-1}$
\item $\varphi_{41} :\equiv \forall e \in \mathbb{R}\,(0 < e \to 0 < e \cdot 2^{-1}) \to f(c) \in \mathbb{R} \to 0 < f(c) \to 0 < f(c) \cdot 2^{-1}$
\item $\varphi_{42} :\equiv (\mathrm{ContIn}(f,A,c) \to 0 < f(c) \to 0 < f(c)) \to \varphi_{1}$
\item $\varphi_{43} :\equiv \forall y\,(f(c) - y = f(c) + -y) \to f(c) - f(c) \cdot 2^{-1} = f(c) + -(f(c) \cdot 2^{-1})$
\item $\varphi_{44} :\equiv \forall x\,\forall y\,(x \in \mathbb{R} \to y \in \mathbb{R} \to x \cdot y \in \mathbb{R}) \to \forall y\,(f(c) \in \mathbb{R} \to y \in \mathbb{R} \to f(c) \cdot y \in \mathbb{R})$
\item $\varphi_{45} :\equiv \forall x\,\forall y\,(x \in \mathbb{R} \to y \in \mathbb{R} \to x - y \in \mathbb{R}) \to \forall y\,(f(x) \in \mathbb{R} \to y \in \mathbb{R} \to f(x) - y \in \mathbb{R})$
\item $\varphi_{46} :\equiv \forall y\,(-(f(c) \cdot 2^{-1}) \in \mathbb{R} \to y \in \mathbb{R} \to -(f(c) \cdot 2^{-1}) + y = y + -(f(c) \cdot 2^{-1}))$
\item $\varphi_{47} :\equiv \forall x\,\forall y\,(x \in \mathbb{R} \to y \in \mathbb{R} \to x - y + y = x) \to \forall y\,(f(x) \in \mathbb{R} \to y \in \mathbb{R} \to f(x) - y + y = f(x))$
\item $\varphi_{48} :\equiv f(c) - f(c) \cdot 2^{-1} = z - f(c) \cdot 2^{-1} \to f(c) - f(c) \cdot 2^{-1} = w - f(c) \cdot 2^{-1}$
\item $\varphi_{49} :\equiv \forall e \in \mathbb{R}\,(e \cdot 2^{-1} + e \cdot 2^{-1} = e) \to f(c) \in \mathbb{R} \to f(c) \cdot 2^{-1} + f(c) \cdot 2^{-1} = f(c)$
\item $\varphi_{50} :\equiv \forall w\,(f(x) - f(c) + f(c) = w \to f(c) \cdot 2^{-1} < f(x) - f(c) + f(c) \to f(c) \cdot 2^{-1} < w)$
\item $\varphi_{51} :\equiv \forall \varepsilon \in \mathbb{R}\,(0 < \varepsilon \to \exists \delta\,(\delta \in \mathbb{R} \land 0 < \delta \land \forall x \in A\,(\lvert x - c \rvert < \delta \to \lvert f(x) - f(c) \rvert < \varepsilon)))$
\item $\varphi_{52} :\equiv \forall y\,(f(c) \in \mathbb{R} \to y \in \mathbb{R} \to f(c) \cdot y \in \mathbb{R}) \to f(c) \in \mathbb{R} \to 2^{-1} \in \mathbb{R} \to f(c) \cdot 2^{-1} \in \mathbb{R}$
\item $\varphi_{53} :\equiv \forall \delta\,(\mathrm{ContIn}(f,A,c) \land 0 < f(c) \land \neg \exists \delta\,(\varphi_{28}))$
\item $\varphi_{54} :\equiv (\mathrm{ContIn}(f,A,c) \to 0 < f(c) \to f(c) \in \mathbb{R}) \to 2^{-1} \in \mathbb{R} \to \varphi_{2}$
\item $\varphi_{55} :\equiv \forall A\,\forall c\,(\mathrm{ContIn}(f,A,c) \to 0 < f(c) \to \exists \delta\,(\varphi_{28}))$
\item $\varphi_{56} :\equiv \forall y\,(f(x) \in \mathbb{R} \to y \in \mathbb{R} \to f(x) - y \in \mathbb{R}) \to f(x) \in \mathbb{R} \to f(c) \in \mathbb{R} \to f(x) - f(c) \in \mathbb{R}$
\item $\varphi_{57} :\equiv (\mathrm{ContIn}(f,A,c) \to 0 < f(c) \to A \subseteq \mathrm{dom}(f)) \to \varphi_{25}$
\item $\varphi_{58} :\equiv \forall w\,(f(c) = w \to f(c) - f(c) \cdot 2^{-1} = f(c) - f(c) \cdot 2^{-1} \to f(c) - f(c) \cdot 2^{-1} = w - f(c) \cdot 2^{-1})$
\item $\varphi_{59} :\equiv f(c) - f(c) \cdot 2^{-1} = f(c) + -(f(c) \cdot 2^{-1}) \to f(c) + -(f(c) \cdot 2^{-1}) = f(c) - f(c) \cdot 2^{-1}$
\item $\varphi_{60} :\equiv f(c) + -(f(c) \cdot 2^{-1}) = f(c) - f(c) \cdot 2^{-1} \to -(f(c) \cdot 2^{-1}) + f(c) = f(c) - f(c) \cdot 2^{-1}$
\item $\varphi_{61} :\equiv \forall y\,\forall z\,(-(f(c) \cdot 2^{-1}) \in \mathbb{R} \to y \in \mathbb{R} \to z \in \mathbb{R} \to -(f(c) \cdot 2^{-1}) < y \to -(f(c) \cdot 2^{-1}) + z < y + z)$
\item $\varphi_{62} :\equiv \mathrm{ContIn}(f,A,c) \land 0 < f(c) \land (\delta \in \mathbb{R} \land 0 < \delta \land \varphi_{20})$
\item $\varphi_{63} :\equiv (\mathrm{ContIn}(f,A,c) \to 0 < f(c) \to \mathrm{RFn}(f)) \to \varphi_{37}$
\item $\varphi_{64} :\equiv \forall y\,(f(x) - f(c) \in \mathbb{R} \to y \in \mathbb{R} \to \lvert f(x) - f(c) \rvert < y \to -y < f(x) - f(c) \land f(x) - f(c) < y)$
\item $\varphi_{65} :\equiv \mathrm{ContIn}(f,A,c) \to 0 < f(c) \to \exists \delta\,(\delta \in \mathbb{R} \land 0 < \delta \land \varphi_{20})$
\item $\varphi_{66} :\equiv f(c) \in \mathbb{R} \to -(f(c) \cdot 2^{-1}) < f(x) - f(c) \to \varphi_{7}$
\item $\varphi_{67} :\equiv \forall z\,(-(f(c) \cdot 2^{-1}) + f(c) = f(c) - f(c) \cdot 2^{-1} \to f(c) - f(c) \cdot 2^{-1} = z \to -(f(c) \cdot 2^{-1}) + f(c) = z)$
\item $\varphi_{68} :\equiv \forall y\,(f(x) \in \mathbb{R} \to y \in \mathbb{R} \to f(x) - y + y = f(x)) \to \varphi_{21}$
\item $\varphi_{69} :\equiv \delta \in \mathbb{R} \land 0 < \delta \land \varphi_{20} \to x \in A \to \lvert x - c \rvert < \delta \to x \in \mathrm{dom}(f)$
\item $\varphi_{70} :\equiv \forall w\,(-(f(c) \cdot 2^{-1}) + f(c) = w \to \varphi_{7} \to w < f(x) - f(c) + f(c))$
\item $\varphi_{71} :\equiv \forall z\,(\varphi_{8} \to f(c) + -(f(c) \cdot 2^{-1}) = z \to -(f(c) \cdot 2^{-1}) + f(c) = z)$
\item $\varphi_{72} :\equiv (\mathrm{ContIn}(f,A,c) \to 0 < f(c) \to f(c) \in \mathbb{R}) \to \varphi_{42}$
\item $\varphi_{73} :\equiv \mathrm{ContIn}(f,A,c) \to 0 < f(c) \to \delta \in \mathbb{R} \land 0 < \delta \land \varphi_{20} \to x \in A \to x \in A$
\item $\varphi_{74} :\equiv \delta \in \mathbb{R} \land 0 < \delta \land \varphi_{20} \to x \in A \to \lvert x - c \rvert < \delta \to f(c) \cdot 2^{-1} \in \mathbb{R}$
\item $\varphi_{75} :\equiv f(c) \cdot 2^{-1} \in \mathbb{R} \to 0 < f(c) \cdot 2^{-1} \to \exists \delta\,(\delta \in \mathbb{R} \land 0 < \delta \land \varphi_{20})$
\item $\varphi_{76} :\equiv \delta \in \mathbb{R} \land 0 < \delta \land \varphi_{20} \to x \in A \to \lvert x - c \rvert < \delta \to f(x) - f(c) \in \mathbb{R}$
\item $\varphi_{77} :\equiv \delta \in \mathbb{R} \land 0 < \delta \land \varphi_{20} \to x \in A \to \lvert x - c \rvert < \delta \to f(c) \cdot 2^{-1} < f(x)$
\item $\varphi_{78} :\equiv \delta \in \mathbb{R} \land 0 < \delta \land \varphi_{20} \to x \in A \to \lvert x - c \rvert < \delta \to -(f(c) \cdot 2^{-1}) \in \mathbb{R}$
\item $\varphi_{79} :\equiv 0 < f(c) \to \delta \in \mathbb{R} \land 0 < \delta \land \varphi_{20} \to x \in A \to \lvert x - c \rvert < \delta \to \lvert x - c \rvert < \delta$
\item $\varphi_{80} :\equiv 0 < f(c) \to \delta \in \mathbb{R} \land 0 < \delta \land \varphi_{20} \to x \in A \to \lvert x - c \rvert < \delta \to f(x) \in \mathbb{R}$
\item $\varphi_{81} :\equiv \delta \in \mathbb{R} \land 0 < \delta \land \varphi_{20} \to \forall x \in A\,(\lvert x - c \rvert < \delta \to f(c) \cdot 2^{-1} < f(x))$
\item $\varphi_{82} :\equiv f(c) \cdot 2^{-1} \in \mathbb{R} \to \lvert f(x) - f(c) \rvert < f(c) \cdot 2^{-1} \to \varphi_{33}$
\item $\varphi_{83} :\equiv \mathrm{ContIn}(f,A,c) \leftrightarrow \mathrm{RFn}(f) \land A \subseteq \mathrm{dom}(f) \land c \in A \land \varphi_{51}$
\item $\varphi_{84} :\equiv -(f(c) \cdot 2^{-1}) + f(c) = f(c) \cdot 2^{-1} \to \varphi_{7} \to f(c) \cdot 2^{-1} < f(x) - f(c) + f(c)$
\item $\varphi_{85} :\equiv \forall z\,(-(f(c) \cdot 2^{-1}) \in \mathbb{R} \to f(x) - f(c) \in \mathbb{R} \to z \in \mathbb{R} \to -(f(c) \cdot 2^{-1}) < f(x) - f(c) \to -(f(c) \cdot 2^{-1}) + z < f(x) - f(c) + z)$
\item $\varphi_{86} :\equiv \forall x\,(\varphi_{62} \to x \in A \to \lvert x - c \rvert < \delta \to f(c) \cdot 2^{-1} < f(x))$
\item $\varphi_{87} :\equiv \varphi_{59} \to f(c) + -(f(c) \cdot 2^{-1}) = f(c) - f(c) \cdot 2^{-1}$
\item $\varphi_{88} :\equiv \forall x\,(\varphi_{62}) \to \forall x \in A\,(\lvert x - c \rvert < \delta \to f(c) \cdot 2^{-1} < f(x))$
\item $\varphi_{89} :\equiv \forall x\,\forall y\,(x \in \mathbb{R} \to y \in \mathbb{R} \to \lvert x \rvert < y \to -y < x \land x < y) \to \varphi_{64}$
\item $\varphi_{90} :\equiv \forall z\,(\varphi_{39} \to f(c) \cdot 2^{-1} + f(c) \cdot 2^{-1} - f(c) \cdot 2^{-1} = z \to -(f(c) \cdot 2^{-1}) + f(c) = z)$
\item $\varphi_{91} :\equiv f(c) = f(c) \cdot 2^{-1} + f(c) \cdot 2^{-1} \to f(c) - f(c) \cdot 2^{-1} = f(c) - f(c) \cdot 2^{-1} \to \varphi_{35}$
\item $\varphi_{92} :\equiv \mathrm{ContIn}(f,A,c) \to 0 < f(c) \to \delta \in \mathbb{R} \land 0 < \delta \land \varphi_{20} \to \varphi_{9}$
\item $\varphi_{93} :\equiv \mathrm{ContIn}(f,A,c) \to 0 < f(c) \to \delta \in \mathbb{R} \land 0 < \delta \land \varphi_{20} \to \varphi_{15}$
\item $\varphi_{94} :\equiv \mathrm{ContIn}(f,A,c) \to 0 < f(c) \to \delta \in \mathbb{R} \land 0 < \delta \land \varphi_{20} \to \varphi_{16}$
\item $\varphi_{95} :\equiv \mathrm{ContIn}(f,A,c) \to 0 < f(c) \to \delta \in \mathbb{R} \land 0 < \delta \land \varphi_{20} \to \varphi_{18}$
\item $\varphi_{96} :\equiv \mathrm{ContIn}(f,A,c) \to 0 < f(c) \to \delta \in \mathbb{R} \land 0 < \delta \land \varphi_{20} \to \varphi_{19}$
\item $\varphi_{97} :\equiv \varphi_{83} \to \mathrm{ContIn}(f,A,c) \to 0 < f(c) \to A \subseteq \mathrm{dom}(f)$
\item $\varphi_{98} :\equiv \mathrm{ContIn}(f,A,c) \to 0 < f(c) \to \delta \in \mathbb{R} \land 0 < \delta \land \varphi_{20} \to x \in A \to \varphi_{12}$
\item $\varphi_{99} :\equiv \mathrm{ContIn}(f,A,c) \to 0 < f(c) \to \delta \in \mathbb{R} \land 0 < \delta \land \varphi_{20} \to x \in A \to \varphi_{13}$
\item $\varphi_{100} :\equiv 0 < f(c) \to \delta \in \mathbb{R} \land 0 < \delta \land \varphi_{20} \to x \in A \to \lvert x - c \rvert < \delta \to \varphi_{7}$
\item $\varphi_{101} :\equiv 0 < f(c) \to \delta \in \mathbb{R} \land 0 < \delta \land \varphi_{20} \to x \in A \to \lvert x - c \rvert < \delta \to \varphi_{8}$
\item $\varphi_{102} :\equiv \mathrm{ContIn}(f,A,c) \to 0 < f(c) \to \delta \in \mathbb{R} \land 0 < \delta \land \varphi_{20} \to \varphi_{28}$
\item $\varphi_{103} :\equiv \mathrm{ContIn}(f,A,c) \to 0 < f(c) \to \delta \in \mathbb{R} \land 0 < \delta \land \varphi_{20} \to \exists \delta\,(\varphi_{28})$
\item $\varphi_{104} :\equiv 0 < f(c) \to \delta \in \mathbb{R} \land 0 < \delta \land \varphi_{20} \to x \in A \to \lvert x - c \rvert < \delta \to \varphi_{15}$
\item $\varphi_{105} :\equiv \mathrm{ContIn}(f,A,c) \land 0 < f(c) \land \neg \exists \delta\,(\varphi_{28}) \to \neg (\delta \in \mathbb{R} \land 0 < \delta \land \varphi_{20})$
\item $\varphi_{106} :\equiv 0 < f(c) \to \delta \in \mathbb{R} \land 0 < \delta \land \varphi_{20} \to x \in A \to \lvert x - c \rvert < \delta \to \varphi_{20}$
\item $\varphi_{107} :\equiv 0 < f(c) \to \delta \in \mathbb{R} \land 0 < \delta \land \varphi_{20} \to x \in A \to \lvert x - c \rvert < \delta \to \varphi_{22}$
\item $\varphi_{108} :\equiv \mathrm{ContIn}(f,A,c) \to 0 < f(c) \to \delta \in \mathbb{R} \land 0 < \delta \land \varphi_{20} \to x \in A \to \varphi_{26}$
\item $\varphi_{109} :\equiv \mathrm{ContIn}(f,A,c) \to 0 < f(c) \to \delta \in \mathbb{R} \land 0 < \delta \land \varphi_{20} \to \delta \in \mathbb{R} \land 0 < \delta \land \varphi_{20}$
\item $\varphi_{110} :\equiv 0 < f(c) \to \delta \in \mathbb{R} \land 0 < \delta \land \varphi_{20} \to x \in A \to \lvert x - c \rvert < \delta \to \varphi_{33}$
\item $\varphi_{111} :\equiv 0 < f(c) \to \delta \in \mathbb{R} \land 0 < \delta \land \varphi_{20} \to x \in A \to \lvert x - c \rvert < \delta \to \varphi_{35}$
\item $\varphi_{112} :\equiv (\mathrm{ContIn}(f,A,c) \to \mathrm{ContIn}(f,A,c)) \to \varphi_{83} \to \mathrm{ContIn}(f,A,c) \to 0 < f(c) \to \mathrm{RFn}(f)$
\item $\varphi_{113} :\equiv (\mathrm{ContIn}(f,A,c) \to \mathrm{ContIn}(f,A,c)) \to \varphi_{83} \to \mathrm{ContIn}(f,A,c) \to 0 < f(c) \to c \in A$
\item $\varphi_{114} :\equiv (\mathrm{ContIn}(f,A,c) \to 0 < f(c) \to f(c) \in \mathbb{R}) \to \mathrm{ContIn}(f,A,c) \to 0 < f(c) \to \varphi_{76}$
\item $\varphi_{115} :\equiv 0 < f(c) \to \delta \in \mathbb{R} \land 0 < \delta \land \varphi_{20} \to x \in A \to \lvert x - c \rvert < \delta \to \varphi_{39}$
\item $\varphi_{116} :\equiv (\mathrm{ContIn}(f,A,c) \to \mathrm{ContIn}(f,A,c)) \to \varphi_{97}$
\item $\varphi_{117} :\equiv (\mathrm{ContIn}(f,A,c) \to 0 < f(c) \to f(c) \in \mathbb{R}) \to \varphi_{92}$
\item $\varphi_{118} :\equiv -(f(c) \cdot 2^{-1}) + f(c) = f(c) - f(c) \cdot 2^{-1} \to \varphi_{35} \to \varphi_{39}$
\item $\varphi_{119} :\equiv (\mathrm{ContIn}(f,A,c) \to 0 < f(c) \to f(c) \in \mathbb{R}) \to 2^{-1} \in \mathbb{R} \to \mathrm{ContIn}(f,A,c) \to 0 < f(c) \to \varphi_{74}$
\item $\varphi_{120} :\equiv (\mathrm{ContIn}(f,A,c) \to 0 < f(c) \to f(c) \in \mathbb{R}) \to \varphi_{95}$
\item $\varphi_{121} :\equiv f(c) - f(c) \cdot 2^{-1} = f(c) + -(f(c) \cdot 2^{-1}) \to \varphi_{87}$
\item $\varphi_{122} :\equiv (\mathrm{ContIn}(f,A,c) \to 0 < f(c) \to f(c) \in \mathbb{R}) \to \mathrm{ContIn}(f,A,c) \to \varphi_{101}$
\item $\varphi_{123} :\equiv (f(x) - f(c) + f(c) = f(x) \to \varphi_{27}) \to \mathrm{ContIn}(f,A,c) \to 0 < f(c) \to \varphi_{77}$
\item $\varphi_{124} :\equiv (f(c) \in \mathbb{R} \to f(c) \cdot 2^{-1} + f(c) \cdot 2^{-1} = f(c)) \to \varphi_{120}$
\item $\varphi_{125} :\equiv (\mathrm{ContIn}(f,A,c) \to \mathrm{ContIn}(f,A,c)) \to \varphi_{83} \to \mathrm{ContIn}(f,A,c) \to 0 < f(c) \to \varphi_{51}$
\item $\varphi_{126} :\equiv (\mathrm{ContIn}(f,A,c) \to 0 < f(c) \to \varphi_{69}) \to \mathrm{ContIn}(f,A,c) \to \varphi_{80}$
\item $\varphi_{127} :\equiv (\mathrm{ContIn}(f,A,c) \land 0 < f(c) \land \neg \exists \delta\,(\varphi_{28}) \to \varphi_{53}) \to \mathrm{ContIn}(f,A,c) \to 0 < f(c) \to \exists \delta\,(\varphi_{28})$
\item $\varphi_{128} :\equiv (\mathrm{ContIn}(f,A,c) \to 0 < f(c) \to A \subseteq \mathrm{dom}(f)) \to \varphi_{73} \to \mathrm{ContIn}(f,A,c) \to 0 < f(c) \to \varphi_{69}$
\item $\varphi_{129} :\equiv (\mathrm{ContIn}(f,A,c) \to 0 < f(c) \to \varphi_{74}) \to \mathrm{ContIn}(f,A,c) \to 0 < f(c) \to \varphi_{78}$
\item $\varphi_{130} :\equiv \varphi_{85} \to -(f(c) \cdot 2^{-1}) \in \mathbb{R} \to f(x) - f(c) \in \mathbb{R} \to \varphi_{66}$
\item $\varphi_{131} :\equiv (\mathrm{ContIn}(f,A,c) \to 0 < f(c) \to \varphi_{77}) \to \varphi_{62} \to x \in A \to \lvert x - c \rvert < \delta \to f(c) \cdot 2^{-1} < f(x)$
\item $\varphi_{132} :\equiv (\mathrm{ContIn}(f,A,c) \to 0 < f(c) \to \mathrm{RFn}(f)) \to \varphi_{126}$
\item $\varphi_{133} :\equiv (\mathrm{ContIn}(f,A,c) \to 0 < f(c) \to \varphi_{74}) \to \mathrm{ContIn}(f,A,c) \to \varphi_{107}$
\item $\varphi_{134} :\equiv \varphi_{90} \to \varphi_{39} \to \varphi_{22} \to -(f(c) \cdot 2^{-1}) + f(c) = f(c) \cdot 2^{-1}$
\item $\varphi_{135} :\equiv (\varphi_{39} \to \varphi_{22} \to -(f(c) \cdot 2^{-1}) + f(c) = f(c) \cdot 2^{-1}) \to \varphi_{98}$
\item $\varphi_{136} :\equiv (\varphi_{62} \to \forall x\,(\varphi_{62})) \to \mathrm{ContIn}(f,A,c) \to 0 < f(c) \to \varphi_{81}$
\item $\varphi_{137} :\equiv (f(x) \in \mathbb{R} \to f(c) \in \mathbb{R} \to f(x) - f(c) \in \mathbb{R}) \to (\mathrm{ContIn}(f,A,c) \to \varphi_{80}) \to \varphi_{114}$
\item $\varphi_{138} :\equiv (\mathrm{ContIn}(f,A,c) \to 0 < f(c) \to f(c) \in \mathbb{R}) \to \varphi_{94} \to \mathrm{ContIn}(f,A,c) \to \varphi_{100}$
\item $\varphi_{139} :\equiv f(c) + -(f(c) \cdot 2^{-1}) = f(c) - f(c) \cdot 2^{-1} \to (\varphi_{8} \to \varphi_{60}) \to \varphi_{108}$
\item $\varphi_{140} :\equiv f(c) - f(c) \cdot 2^{-1} = f(c) - f(c) \cdot 2^{-1} \to \varphi_{91} \to \mathrm{ContIn}(f,A,c) \to \varphi_{111}$
\item $\varphi_{141} :\equiv (-(f(c) \cdot 2^{-1}) \in \mathbb{R} \to f(c) \in \mathbb{R} \to \varphi_{8}) \to (\mathrm{ContIn}(f,A,c) \to 0 < f(c) \to \varphi_{78}) \to \varphi_{122}$
\item $\varphi_{142} :\equiv (\mathrm{ContIn}(f,A,c) \to 0 < f(c) \to \varphi_{51}) \to (\varphi_{51} \to \varphi_{75}) \to \mathrm{ContIn}(f,A,c) \to 0 < f(c) \to \varphi_{75}$
\item $\varphi_{143} :\equiv (\mathrm{ContIn}(f,A,c) \to \varphi_{106}) \to (\varphi_{20} \to \varphi_{15}) \to \mathrm{ContIn}(f,A,c) \to \varphi_{104}$
\item $\varphi_{144} :\equiv (\mathrm{ContIn}(f,A,c) \to 0 < f(c) \to \varphi_{74}) \to \varphi_{93} \to \mathrm{ContIn}(f,A,c) \to \varphi_{110}$
\item $\varphi_{145} :\equiv (f(c) \cdot 2^{-1} \in \mathbb{R} \to f(c) \cdot 2^{-1} \in \mathbb{R} \to \varphi_{22}) \to (\mathrm{ContIn}(f,A,c) \to 0 < f(c) \to \varphi_{74}) \to \varphi_{133}$
\item $\varphi_{146} :\equiv (\mathrm{ContIn}(f,A,c) \to \varphi_{111}) \to \varphi_{118} \to \mathrm{ContIn}(f,A,c) \to \varphi_{115}$
\item $\varphi_{147} :\equiv (\mathrm{ContIn}(f,A,c) \to \varphi_{104}) \to \varphi_{73} \to (\mathrm{ContIn}(f,A,c) \to \varphi_{79}) \to \varphi_{93}$
\item $\varphi_{148} :\equiv (\mathrm{ContIn}(f,A,c) \to 0 < f(c) \to \varphi_{78}) \to (\mathrm{ContIn}(f,A,c) \to 0 < f(c) \to \varphi_{76}) \to \varphi_{138}$
\item $\varphi_{149} :\equiv (\mathrm{ContIn}(f,A,c) \to \varphi_{115}) \to (\mathrm{ContIn}(f,A,c) \to \varphi_{107}) \to \varphi_{135}$
\item $\varphi_{150} :\equiv (f(x) - f(c) \in \mathbb{R} \to \varphi_{82}) \to (\mathrm{ContIn}(f,A,c) \to 0 < f(c) \to \varphi_{76}) \to \varphi_{144}$
\item $\varphi_{151} :\equiv \forall \delta\,(\varphi_{105}) \to (\forall \delta\,(\varphi_{105}) \to \varphi_{53} \to \forall \delta\,(\neg (\delta \in \mathbb{R} \land 0 < \delta \land \varphi_{20}))) \to \varphi_{127}$
\item $\varphi_{152} :\equiv (-(f(c) \cdot 2^{-1}) \in \mathbb{R} \to f(x) - f(c) \in \mathbb{R} \to \varphi_{66}) \to \varphi_{148}$
\end{abbrevlist}
\end{fullproof}

\section{第4章 2節: 定義2.1 の同値性, 中間値の定理(2.14), 最大値・最小値の定理(2.17)}
% 自動生成 (prover/build.js): HalfLt  ノード数 43
\begin{fullproof}{[E39]}{$\forall e \in \mathbb{R}\,(0 < e \to e \cdot 2^{-1} < e)$}{43}\phantomsection\label{tree:HalfLt}
\noindent$\Pi_{1}$:
\begin{ptree}
\AXC{$z = w \to e \cdot 2^{-1} < z \to e \cdot 2^{-1} < w$\lbl{A6}}
\UIC{$\forall w\,(z = w \to e \cdot 2^{-1} < z \to e \cdot 2^{-1} < w)$}
\UIC{$\forall z\,\forall w\,(z = w \to e \cdot 2^{-1} < z \to e \cdot 2^{-1} < w)$}
\AXC{$\forall z\,\forall w\,(z = w \to e \cdot 2^{-1} < z \to e \cdot 2^{-1} < w) \to \varphi_{7}$\lbl{A2}}
\BIC{$\varphi_{7}$}
\AXC{$\varphi_{7} \to \varphi_{6}$\lbl{A2}}
\BIC{$\varphi_{6}$}
\end{ptree}
\noindent$\Pi_{2}$:
\begin{ptree}
\AXC{$\forall x\,\forall y\,(x \in \mathbb{R} \to y \in \mathbb{R} \to x \cdot y \in \mathbb{R})$\lbl{(R1)}}
\AXC{$\varphi_{8}$\lbl{A2}}
\BIC{$\forall y\,(e \in \mathbb{R} \to y \in \mathbb{R} \to e \cdot y \in \mathbb{R})$}
\AXC{$\varphi_{10}$\lbl{A2}}
\BIC{$e \in \mathbb{R} \to 2^{-1} \in \mathbb{R} \to e \cdot 2^{-1} \in \mathbb{R}$}
\AXC{$(e \in \mathbb{R} \to 2^{-1} \in \mathbb{R} \to e \cdot 2^{-1} \in \mathbb{R}) \to \varphi_{9}$\lbl{A1}}
\BIC{$\varphi_{9}$}
\end{ptree}
\noindent$\Pi_{3}$:
\begin{ptree}
\AXC{$2^{-1} \in \mathbb{R}$\lbl{[F31]}}
\AXC{$e \in \mathbb{R} \to 0 < e \to e \in \mathbb{R}$\lbl{A1}}
\AXC{$\stackrel{\Pi_{2}}{\vdots}$}\noLine
\UIC{$\varphi_{9}$}
\BIC{$2^{-1} \in \mathbb{R} \to e \in \mathbb{R} \to 0 < e \to e \cdot 2^{-1} \in \mathbb{R}$}
\BIC{$e \in \mathbb{R} \to 0 < e \to e \cdot 2^{-1} \in \mathbb{R}$}
\end{ptree}
\noindent$\Pi_{4}$:
\begin{ptree}
\AXC{$\forall x\,\forall y\,(x \in \mathbb{R} \to y \in \mathbb{R} \to 0 < y \to x < x + y)$\lbl{[E3]}}
\AXC{$\forall x\,\forall y\,(x \in \mathbb{R} \to y \in \mathbb{R} \to 0 < y \to x < x + y) \to \varphi_{1}$\lbl{A2}}
\BIC{$\varphi_{1}$}
\end{ptree}
\noindent$\Pi_{5}$:
\begin{ptree}
\AXC{$\stackrel{\Pi_{3}}{\vdots}$}\noLine
\UIC{$e \in \mathbb{R} \to 0 < e \to e \cdot 2^{-1} \in \mathbb{R}$}
\AXC{$\stackrel{\Pi_{4}}{\vdots}$}\noLine
\UIC{$\varphi_{1}$}
\AXC{$\varphi_{1} \to e \cdot 2^{-1} \in \mathbb{R} \to \varphi_{3}$\lbl{A2}}
\BIC{$e \cdot 2^{-1} \in \mathbb{R} \to \varphi_{3}$}
\AXC{$(e \cdot 2^{-1} \in \mathbb{R} \to \varphi_{3}) \to \varphi_{12}$\lbl{A1}}
\BIC{$\varphi_{12}$}
\BIC{$(e \in \mathbb{R} \to 0 < e \to e \cdot 2^{-1} \in \mathbb{R}) \to \varphi_{11}$}
\end{ptree}
\noindent$\Pi_{6}$:
\begin{ptree}
\AXC{$\stackrel{\Pi_{3}}{\vdots}$}\noLine
\UIC{$e \in \mathbb{R} \to 0 < e \to e \cdot 2^{-1} \in \mathbb{R}$}
\AXC{$\stackrel{\Pi_{5}}{\vdots}$}\noLine
\UIC{$(e \in \mathbb{R} \to 0 < e \to e \cdot 2^{-1} \in \mathbb{R}) \to \varphi_{11}$}
\BIC{$\varphi_{11}$}
\end{ptree}
\noindent$\Pi_{7}$:
\begin{ptree}
\AXC{$\forall e \in \mathbb{R}\,(e \cdot 2^{-1} + e \cdot 2^{-1} = e)$\lbl{[F36]}}
\AXC{$\varphi_{5}$\lbl{A2}}
\BIC{$e \in \mathbb{R} \to e \cdot 2^{-1} + e \cdot 2^{-1} = e$}
\AXC{$(e \in \mathbb{R} \to e \cdot 2^{-1} + e \cdot 2^{-1} = e) \to \varphi_{4}$\lbl{A1}}
\BIC{$\varphi_{4}$}
\end{ptree}
\noindent$\Pi_{8}$:
\begin{ptree}
\AXC{$e \in \mathbb{R} \to 0 < e \to e \in \mathbb{R}$\lbl{A1}}
\AXC{$\stackrel{\Pi_{7}}{\vdots}$}\noLine
\UIC{$\varphi_{4}$}
\BIC{$e \in \mathbb{R} \to 0 < e \to e \cdot 2^{-1} + e \cdot 2^{-1} = e$}
\AXC{$(e \in \mathbb{R} \to 0 < e \to e \cdot 2^{-1} + e \cdot 2^{-1} = e) \to \varphi_{13}$\lbl{A1}}
\BIC{$\varphi_{13}$}
\end{ptree}
\noindent 主木:
\begin{ptree}
\AXC{$\stackrel{\Pi_{1}}{\vdots}$}\noLine
\UIC{$\varphi_{6}$}
\AXC{$\forall e \in \mathbb{R}\,(0 < e \to 0 < e \cdot 2^{-1})$\lbl{[F35]}}
\AXC{$\varphi_{2}$\lbl{A2}}
\BIC{$e \in \mathbb{R} \to 0 < e \to 0 < e \cdot 2^{-1}$}
\AXC{$\stackrel{\Pi_{6}}{\vdots}$}\noLine
\UIC{$\varphi_{11}$}
\BIC{$e \in \mathbb{R} \to 0 < e \to e \cdot 2^{-1} < e \cdot 2^{-1} + e \cdot 2^{-1}$}
\AXC{$\stackrel{\Pi_{8}}{\vdots}$}\noLine
\UIC{$\varphi_{13}$}
\BIC{$\varphi_{6} \to e \in \mathbb{R} \to 0 < e \to e \cdot 2^{-1} < e$}
\BIC{$e \in \mathbb{R} \to 0 < e \to e \cdot 2^{-1} < e$}
\UIC{$\forall e \in \mathbb{R}\,(0 < e \to e \cdot 2^{-1} < e)$}
\end{ptree}
\begin{abbrevlist}
\item $\varphi_{1} :\equiv \forall y\,(e \cdot 2^{-1} \in \mathbb{R} \to y \in \mathbb{R} \to 0 < y \to e \cdot 2^{-1} < e \cdot 2^{-1} + y)$
\item $\varphi_{2} :\equiv \forall e \in \mathbb{R}\,(0 < e \to 0 < e \cdot 2^{-1}) \to e \in \mathbb{R} \to 0 < e \to 0 < e \cdot 2^{-1}$
\item $\varphi_{3} :\equiv e \cdot 2^{-1} \in \mathbb{R} \to 0 < e \cdot 2^{-1} \to e \cdot 2^{-1} < e \cdot 2^{-1} + e \cdot 2^{-1}$
\item $\varphi_{4} :\equiv (e \in \mathbb{R} \to 0 < e \to e \in \mathbb{R}) \to e \in \mathbb{R} \to 0 < e \to e \cdot 2^{-1} + e \cdot 2^{-1} = e$
\item $\varphi_{5} :\equiv \forall e \in \mathbb{R}\,(e \cdot 2^{-1} + e \cdot 2^{-1} = e) \to e \in \mathbb{R} \to e \cdot 2^{-1} + e \cdot 2^{-1} = e$
\item $\varphi_{6} :\equiv e \cdot 2^{-1} + e \cdot 2^{-1} = e \to e \cdot 2^{-1} < e \cdot 2^{-1} + e \cdot 2^{-1} \to e \cdot 2^{-1} < e$
\item $\varphi_{7} :\equiv \forall w\,(e \cdot 2^{-1} + e \cdot 2^{-1} = w \to e \cdot 2^{-1} < e \cdot 2^{-1} + e \cdot 2^{-1} \to e \cdot 2^{-1} < w)$
\item $\varphi_{8} :\equiv \forall x\,\forall y\,(x \in \mathbb{R} \to y \in \mathbb{R} \to x \cdot y \in \mathbb{R}) \to \forall y\,(e \in \mathbb{R} \to y \in \mathbb{R} \to e \cdot y \in \mathbb{R})$
\item $\varphi_{9} :\equiv (e \in \mathbb{R} \to 0 < e \to e \in \mathbb{R}) \to 2^{-1} \in \mathbb{R} \to e \in \mathbb{R} \to 0 < e \to e \cdot 2^{-1} \in \mathbb{R}$
\item $\varphi_{10} :\equiv \forall y\,(e \in \mathbb{R} \to y \in \mathbb{R} \to e \cdot y \in \mathbb{R}) \to e \in \mathbb{R} \to 2^{-1} \in \mathbb{R} \to e \cdot 2^{-1} \in \mathbb{R}$
\item $\varphi_{11} :\equiv (e \in \mathbb{R} \to 0 < e \to 0 < e \cdot 2^{-1}) \to e \in \mathbb{R} \to 0 < e \to e \cdot 2^{-1} < e \cdot 2^{-1} + e \cdot 2^{-1}$
\item $\varphi_{12} :\equiv (e \in \mathbb{R} \to 0 < e \to e \cdot 2^{-1} \in \mathbb{R}) \to (e \in \mathbb{R} \to 0 < e \to e \cdot 2^{-1} \in \mathbb{R}) \to \varphi_{11}$
\item $\varphi_{13} :\equiv (e \in \mathbb{R} \to 0 < e \to e \cdot 2^{-1} < e \cdot 2^{-1} + e \cdot 2^{-1}) \to \varphi_{6} \to e \in \mathbb{R} \to 0 < e \to e \cdot 2^{-1} < e$
\end{abbrevlist}
\end{fullproof}

% 自動生成 (prover/build.js): SubNonneg  ノード数 63
\begin{fullproof}{[E40]}{$\forall x\,\forall y\,(x \in \mathbb{R} \to y \in \mathbb{R} \to x \leq y \to 0 \leq y - x)$}{63}\phantomsection\label{tree:SubNonneg}
\noindent$\Pi_{1}$:
\begin{ptree}
\AXC{$x = x$\lbl{A5}}
\AXC{$x = y \to x = x \to y = x$\lbl{A6}}
\AXC{$(x = y \to x = x \to y = x) \to x = x \to x = y \to y = x$\lbl{A1}}
\BIC{$x = x \to x = y \to y = x$}
\BIC{$x = y \to y = x$}
\UIC{$\forall y\,(x = y \to y = x)$}
\UIC{$\forall x\,\forall y\,(x = y \to y = x)$}
\end{ptree}
\noindent$\Pi_{2}$:
\begin{ptree}
\AXC{$\stackrel{\Pi_{1}}{\vdots}$}\noLine
\UIC{$\forall x\,\forall y\,(x = y \to y = x)$}
\AXC{$\forall x\,\forall y\,(x = y \to y = x) \to \forall y\,(v_{776} = y \to y = v_{776})$\lbl{A2}}
\BIC{$\forall y\,(v_{776} = y \to y = v_{776})$}
\AXC{$\varphi_{1}$\lbl{A2}}
\BIC{$v_{776} = v_{777} \to v_{777} = v_{776}$}
\UIC{$\forall v_{777}\,(v_{776} = v_{777} \to v_{777} = v_{776})$}
\UIC{$\forall v_{776}\,\forall v_{777}\,(v_{776} = v_{777} \to v_{777} = v_{776})$}
\AXC{$\varphi_{5}$\lbl{A2}}
\BIC{$\forall v_{777}\,(y - x = v_{777} \to v_{777} = y - x)$}
\AXC{$\varphi_{2}$\lbl{A2}}
\BIC{$y - x = 0 \to 0 = y - x$}
\end{ptree}
\noindent$\Pi_{3}$:
\begin{ptree}
\AXC{$\forall x \in \mathbb{R}\,(x - x = 0)$\lbl{[F13]}}
\AXC{$\forall x \in \mathbb{R}\,(x - x = 0) \to x \in \mathbb{R} \to x - x = 0$\lbl{A2}}
\BIC{$x \in \mathbb{R} \to x - x = 0$}
\AXC{$(x \in \mathbb{R} \to x - x = 0) \to \varphi_{10}$\lbl{A1}}
\BIC{$\varphi_{10}$}
\end{ptree}
\noindent$\Pi_{4}$:
\begin{ptree}
\AXC{$x \in \mathbb{R} \to y \in \mathbb{R} \to x \leq y \to x = y \to x \in \mathbb{R}$\lbl{A1}}
\AXC{$\stackrel{\Pi_{3}}{\vdots}$}\noLine
\UIC{$\varphi_{10}$}
\BIC{$x \in \mathbb{R} \to y \in \mathbb{R} \to x \leq y \to x = y \to x - x = 0$}
\AXC{$x \in \mathbb{R} \to y \in \mathbb{R} \to x \leq y \to x = y \to x = y$\lbl{A1}}
\AXC{$\varphi_{15}$\lbl{A1}}
\BIC{$(x \in \mathbb{R} \to y \in \mathbb{R} \to x \leq y \to x = y \to x - x = 0) \to \varphi_{7}$}
\BIC{$\varphi_{7}$}
\end{ptree}
\noindent$\Pi_{5}$:
\begin{ptree}
\AXC{$z = w \to z - x = 0 \to w - x = 0$\lbl{A6}}
\UIC{$\forall w\,(z = w \to z - x = 0 \to w - x = 0)$}
\UIC{$\forall z\,\forall w\,(z = w \to z - x = 0 \to w - x = 0)$}
\AXC{$\varphi_{4}$\lbl{A2}}
\BIC{$\forall w\,(x = w \to x - x = 0 \to w - x = 0)$}
\AXC{$\varphi_{3}$\lbl{A2}}
\BIC{$x = y \to x - x = 0 \to y - x = 0$}
\AXC{$\stackrel{\Pi_{4}}{\vdots}$}\noLine
\UIC{$\varphi_{7}$}
\BIC{$x \in \mathbb{R} \to y \in \mathbb{R} \to x \leq y \to x = y \to y - x = 0$}
\end{ptree}
\noindent$\Pi_{6}$:
\begin{ptree}
\AXC{$\stackrel{\Pi_{2}}{\vdots}$}\noLine
\UIC{$y - x = 0 \to 0 = y - x$}
\AXC{$\stackrel{\Pi_{5}}{\vdots}$}\noLine
\UIC{$x \in \mathbb{R} \to y \in \mathbb{R} \to x \leq y \to x = y \to y - x = 0$}
\AXC{$(x \in \mathbb{R} \to y \in \mathbb{R} \to x \leq y \to x = y \to y - x = 0) \to \varphi_{6}$\lbl{A1}}
\BIC{$\varphi_{6}$}
\BIC{$x \in \mathbb{R} \to y \in \mathbb{R} \to x \leq y \to x = y \to 0 = y - x$}
\end{ptree}
\noindent$\Pi_{7}$:
\begin{ptree}
\AXC{$\forall x\,\forall y\,(x \in \mathbb{R} \to y \in \mathbb{R} \to x < y \to 0 < y - x)$\lbl{[F24]}}
\AXC{$\varphi_{13}$\lbl{A2}}
\BIC{$\forall y\,(x \in \mathbb{R} \to y \in \mathbb{R} \to x < y \to 0 < y - x)$}
\AXC{$\varphi_{11}$\lbl{A2}}
\BIC{$x \in \mathbb{R} \to y \in \mathbb{R} \to x < y \to 0 < y - x$}
\AXC{$(x \in \mathbb{R} \to y \in \mathbb{R} \to x < y \to 0 < y - x) \to \varphi_{17}$\lbl{A1}}
\BIC{$\varphi_{17}$}
\end{ptree}
\noindent$\Pi_{8}$:
\begin{ptree}
\AXC{$x \in \mathbb{R} \to y \in \mathbb{R} \to x \leq y \to x < y \to y \in \mathbb{R}$\lbl{A1}}
\AXC{$x \in \mathbb{R} \to y \in \mathbb{R} \to x \leq y \to x < y \to x \in \mathbb{R}$\lbl{A1}}
\AXC{$\stackrel{\Pi_{7}}{\vdots}$}\noLine
\UIC{$\varphi_{17}$}
\BIC{$(x \in \mathbb{R} \to y \in \mathbb{R} \to x \leq y \to x < y \to y \in \mathbb{R}) \to \varphi_{12}$}
\BIC{$\varphi_{12}$}
\end{ptree}
\noindent$\Pi_{9}$:
\begin{ptree}
\AXC{$x \in \mathbb{R} \to y \in \mathbb{R} \to x \leq y \to x < y \to x < y$\lbl{A1}}
\AXC{$\stackrel{\Pi_{8}}{\vdots}$}\noLine
\UIC{$\varphi_{12}$}
\BIC{$x \in \mathbb{R} \to y \in \mathbb{R} \to x \leq y \to x < y \to 0 < y - x$}
\AXC{$\varphi_{14}$\lbl{A1}}
\BIC{$x \in \mathbb{R} \to y \in \mathbb{R} \to x \leq y \to x < y \to 0 \leq y - x$}
\AXC{$x \in \mathbb{R} \to y \in \mathbb{R} \to x \leq y \to x \leq y$\lbl{A1}}
\AXC{$\varphi_{16}$\lbl{A1}}
\BIC{$(x \in \mathbb{R} \to y \in \mathbb{R} \to x \leq y \to x < y \to 0 \leq y - x) \to \varphi_{8}$}
\BIC{$\varphi_{8}$}
\end{ptree}
\noindent 主木:
\begin{ptree}
\AXC{$\stackrel{\Pi_{6}}{\vdots}$}\noLine
\UIC{$x \in \mathbb{R} \to y \in \mathbb{R} \to x \leq y \to x = y \to 0 = y - x$}
\AXC{$\varphi_{9}$\lbl{A1}}
\BIC{$x \in \mathbb{R} \to y \in \mathbb{R} \to x \leq y \to x = y \to 0 \leq y - x$}
\AXC{$\stackrel{\Pi_{9}}{\vdots}$}\noLine
\UIC{$\varphi_{8}$}
\BIC{$x \in \mathbb{R} \to y \in \mathbb{R} \to x \leq y \to 0 \leq y - x$}
\UIC{$\forall y\,(x \in \mathbb{R} \to y \in \mathbb{R} \to x \leq y \to 0 \leq y - x)$}
\UIC{$\forall x\,\forall y\,(x \in \mathbb{R} \to y \in \mathbb{R} \to x \leq y \to 0 \leq y - x)$}
\end{ptree}
\begin{abbrevlist}
\item $\varphi_{1} :\equiv \forall y\,(v_{776} = y \to y = v_{776}) \to v_{776} = v_{777} \to v_{777} = v_{776}$
\item $\varphi_{2} :\equiv \forall v_{777}\,(y - x = v_{777} \to v_{777} = y - x) \to y - x = 0 \to 0 = y - x$
\item $\varphi_{3} :\equiv \forall w\,(x = w \to x - x = 0 \to w - x = 0) \to x = y \to x - x = 0 \to y - x = 0$
\item $\varphi_{4} :\equiv \forall z\,\forall w\,(z = w \to z - x = 0 \to w - x = 0) \to \forall w\,(x = w \to x - x = 0 \to w - x = 0)$
\item $\varphi_{5} :\equiv \forall v_{776}\,\forall v_{777}\,(v_{776} = v_{777} \to v_{777} = v_{776}) \to \forall v_{777}\,(y - x = v_{777} \to v_{777} = y - x)$
\item $\varphi_{6} :\equiv (y - x = 0 \to 0 = y - x) \to x \in \mathbb{R} \to y \in \mathbb{R} \to x \leq y \to x = y \to 0 = y - x$
\item $\varphi_{7} :\equiv (x = y \to x - x = 0 \to y - x = 0) \to x \in \mathbb{R} \to y \in \mathbb{R} \to x \leq y \to x = y \to y - x = 0$
\item $\varphi_{8} :\equiv (x \in \mathbb{R} \to y \in \mathbb{R} \to x \leq y \to x = y \to 0 \leq y - x) \to x \in \mathbb{R} \to y \in \mathbb{R} \to x \leq y \to 0 \leq y - x$
\item $\varphi_{9} :\equiv (x \in \mathbb{R} \to y \in \mathbb{R} \to x \leq y \to x = y \to 0 = y - x) \to x \in \mathbb{R} \to y \in \mathbb{R} \to x \leq y \to x = y \to 0 \leq y - x$
\item $\varphi_{10} :\equiv (x \in \mathbb{R} \to y \in \mathbb{R} \to x \leq y \to x = y \to x \in \mathbb{R}) \to x \in \mathbb{R} \to y \in \mathbb{R} \to x \leq y \to x = y \to x - x = 0$
\item $\varphi_{11} :\equiv \forall y\,(x \in \mathbb{R} \to y \in \mathbb{R} \to x < y \to 0 < y - x) \to x \in \mathbb{R} \to y \in \mathbb{R} \to x < y \to 0 < y - x$
\item $\varphi_{12} :\equiv (x \in \mathbb{R} \to y \in \mathbb{R} \to x \leq y \to x < y \to x < y) \to x \in \mathbb{R} \to y \in \mathbb{R} \to x \leq y \to x < y \to 0 < y - x$
\item $\varphi_{13} :\equiv \forall x\,\forall y\,(x \in \mathbb{R} \to y \in \mathbb{R} \to x < y \to 0 < y - x) \to \forall y\,(x \in \mathbb{R} \to y \in \mathbb{R} \to x < y \to 0 < y - x)$
\item $\varphi_{14} :\equiv (x \in \mathbb{R} \to y \in \mathbb{R} \to x \leq y \to x < y \to 0 < y - x) \to x \in \mathbb{R} \to y \in \mathbb{R} \to x \leq y \to x < y \to 0 \leq y - x$
\item $\varphi_{15} :\equiv (x \in \mathbb{R} \to y \in \mathbb{R} \to x \leq y \to x = y \to x = y) \to (x \in \mathbb{R} \to y \in \mathbb{R} \to x \leq y \to x = y \to x - x = 0) \to \varphi_{7}$
\item $\varphi_{16} :\equiv (x \in \mathbb{R} \to y \in \mathbb{R} \to x \leq y \to x \leq y) \to (x \in \mathbb{R} \to y \in \mathbb{R} \to x \leq y \to x < y \to 0 \leq y - x) \to \varphi_{8}$
\item $\varphi_{17} :\equiv (x \in \mathbb{R} \to y \in \mathbb{R} \to x \leq y \to x < y \to x \in \mathbb{R}) \to (x \in \mathbb{R} \to y \in \mathbb{R} \to x \leq y \to x < y \to y \in \mathbb{R}) \to \varphi_{12}$
\end{abbrevlist}
\end{fullproof}

% 自動生成 (prover/build.js): LeftAbsLe  ノード数 159
\begin{fullproof}{[E41]}{$\forall s\,\forall x\,\forall d\,(s \in \mathbb{R} \to x \in \mathbb{R} \to d \in \mathbb{R} \to s \leq x \to x - d < s \to \lvert s - x \rvert < d)$}{159}\phantomsection\label{tree:LeftAbsLe}
\noindent$\Pi_{1}$:
\begin{ptree}
\AXC{$z = w \to -(s - x) < z \to -(s - x) < w$\lbl{A6}}
\UIC{$\forall w\,(z = w \to -(s - x) < z \to -(s - x) < w)$}
\UIC{$\forall z\,\forall w\,(z = w \to -(s - x) < z \to -(s - x) < w)$}
\AXC{$\forall z\,\forall w\,(z = w \to -(s - x) < z \to -(s - x) < w) \to \varphi_{4}$\lbl{A2}}
\BIC{$\varphi_{4}$}
\AXC{$\varphi_{34}$\lbl{A2}}
\BIC{$-(-d) = d \to -(s - x) < -(-d) \to -(s - x) < d$}
\end{ptree}
\noindent$\Pi_{2}$:
\begin{ptree}
\AXC{$\forall x\,(\varphi_{14})$\lbl{[E31]}}
\AXC{$\forall x\,(\varphi_{14}) \to \varphi_{14}$\lbl{A2}}
\BIC{$\varphi_{14}$}
\AXC{$\varphi_{14} \to \varphi_{8}$\lbl{A2}}
\BIC{$\varphi_{8}$}
\AXC{$\varphi_{8} \to x \in \mathbb{R} \to d \in \mathbb{R} \to s \in \mathbb{R} \to x - d < s \to -d < s - x$\lbl{A2}}
\BIC{$x \in \mathbb{R} \to d \in \mathbb{R} \to s \in \mathbb{R} \to x - d < s \to -d < s - x$}
\end{ptree}
\noindent$\Pi_{3}$:
\begin{ptree}
\AXC{$s \in \mathbb{R} \to x \in \mathbb{R} \to d \in \mathbb{R} \to s \leq x \to x - d < s \to x \in \mathbb{R}$\lbl{A1}}
\AXC{$\stackrel{\Pi_{2}}{\vdots}$}\noLine
\UIC{$x \in \mathbb{R} \to d \in \mathbb{R} \to s \in \mathbb{R} \to x - d < s \to -d < s - x$}
\AXC{$\varphi_{55}$\lbl{A1}}
\BIC{$\varphi_{54}$}
\BIC{$\varphi_{52}$}
\end{ptree}
\noindent$\Pi_{4}$:
\begin{ptree}
\AXC{$s \in \mathbb{R} \to x \in \mathbb{R} \to d \in \mathbb{R} \to s \leq x \to x - d < s \to s \in \mathbb{R}$\lbl{A1}}
\AXC{$s \in \mathbb{R} \to x \in \mathbb{R} \to d \in \mathbb{R} \to s \leq x \to x - d < s \to d \in \mathbb{R}$\lbl{A1}}
\AXC{$\stackrel{\Pi_{3}}{\vdots}$}\noLine
\UIC{$\varphi_{52}$}
\BIC{$\varphi_{48}$}
\BIC{$\varphi_{42}$}
\end{ptree}
\noindent$\Pi_{5}$:
\begin{ptree}
\AXC{$\forall x\,\forall y\,(x \in \mathbb{R} \to y \in \mathbb{R} \to x - y \in \mathbb{R})$\lbl{[F1]}}
\AXC{$\varphi_{27}$\lbl{A2}}
\BIC{$\forall y\,(s \in \mathbb{R} \to y \in \mathbb{R} \to s - y \in \mathbb{R})$}
\AXC{$\varphi_{25}$\lbl{A2}}
\BIC{$s \in \mathbb{R} \to x \in \mathbb{R} \to s - x \in \mathbb{R}$}
\AXC{$(s \in \mathbb{R} \to x \in \mathbb{R} \to s - x \in \mathbb{R}) \to \varphi_{47}$\lbl{A1}}
\BIC{$\varphi_{47}$}
\end{ptree}
\noindent$\Pi_{6}$:
\begin{ptree}
\AXC{$s \in \mathbb{R} \to x \in \mathbb{R} \to d \in \mathbb{R} \to s \leq x \to x - d < s \to x \in \mathbb{R}$\lbl{A1}}
\AXC{$s \in \mathbb{R} \to x \in \mathbb{R} \to d \in \mathbb{R} \to s \leq x \to x - d < s \to s \in \mathbb{R}$\lbl{A1}}
\AXC{$\stackrel{\Pi_{5}}{\vdots}$}\noLine
\UIC{$\varphi_{47}$}
\BIC{$\varphi_{39}$}
\BIC{$\varphi_{12}$}
\end{ptree}
\noindent$\Pi_{7}$:
\begin{ptree}
\AXC{$\forall x \in \mathbb{R}\,(-x \in \mathbb{R})$\lbl{(R1)}}
\AXC{$\forall x \in \mathbb{R}\,(-x \in \mathbb{R}) \to d \in \mathbb{R} \to -d \in \mathbb{R}$\lbl{A2}}
\BIC{$d \in \mathbb{R} \to -d \in \mathbb{R}$}
\AXC{$(d \in \mathbb{R} \to -d \in \mathbb{R}) \to \varphi_{40}$\lbl{A1}}
\BIC{$\varphi_{40}$}
\end{ptree}
\noindent$\Pi_{8}$:
\begin{ptree}
\AXC{$\forall x\,\forall y\,(x \in \mathbb{R} \to y \in \mathbb{R} \to x < y \to -y < -x)$\lbl{[F22]}}
\AXC{$\varphi_{35}$\lbl{A2}}
\BIC{$\forall y\,(-d \in \mathbb{R} \to y \in \mathbb{R} \to -d < y \to -y < -(-d))$}
\AXC{$\forall y\,(-d \in \mathbb{R} \to y \in \mathbb{R} \to -d < y \to -y < -(-d)) \to \varphi_{22}$\lbl{A2}}
\BIC{$\varphi_{22}$}
\end{ptree}
\noindent$\Pi_{9}$:
\begin{ptree}
\AXC{$\stackrel{\Pi_{6}}{\vdots}$}\noLine
\UIC{$\varphi_{12}$}
\AXC{$s \in \mathbb{R} \to x \in \mathbb{R} \to d \in \mathbb{R} \to s \leq x \to x - d < s \to d \in \mathbb{R}$\lbl{A1}}
\AXC{$\stackrel{\Pi_{7}}{\vdots}$}\noLine
\UIC{$\varphi_{40}$}
\BIC{$s \in \mathbb{R} \to x \in \mathbb{R} \to d \in \mathbb{R} \to s \leq x \to x - d < s \to -d \in \mathbb{R}$}
\AXC{$\stackrel{\Pi_{8}}{\vdots}$}\noLine
\UIC{$\varphi_{22}$}
\AXC{$\varphi_{22} \to \varphi_{53}$\lbl{A1}}
\BIC{$\varphi_{53}$}
\BIC{$\varphi_{12} \to \varphi_{15} \to s \in \mathbb{R} \to \varphi_{10}$}
\BIC{$\varphi_{15} \to s \in \mathbb{R} \to \varphi_{10}$}
\end{ptree}
\noindent$\Pi_{10}$:
\begin{ptree}
\AXC{$\forall x \in \mathbb{R}\,(-(-x) = x)$\lbl{[F5]}}
\AXC{$\forall x \in \mathbb{R}\,(-(-x) = x) \to d \in \mathbb{R} \to -(-d) = d$\lbl{A2}}
\BIC{$d \in \mathbb{R} \to -(-d) = d$}
\AXC{$(d \in \mathbb{R} \to -(-d) = d) \to \varphi_{38}$\lbl{A1}}
\BIC{$\varphi_{38}$}
\end{ptree}
\noindent$\Pi_{11}$:
\begin{ptree}
\AXC{$s \in \mathbb{R} \to x \in \mathbb{R} \to d \in \mathbb{R} \to s \leq x \to x - d < s \to d \in \mathbb{R}$\lbl{A1}}
\AXC{$\stackrel{\Pi_{10}}{\vdots}$}\noLine
\UIC{$\varphi_{38}$}
\BIC{$\varphi_{9}$}
\AXC{$\varphi_{9} \to (s \in \mathbb{R} \to \varphi_{10}) \to \varphi_{37}$\lbl{A1}}
\BIC{$(s \in \mathbb{R} \to \varphi_{10}) \to \varphi_{37}$}
\end{ptree}
\noindent$\Pi_{12}$:
\begin{ptree}
\AXC{$s \in \mathbb{R} \to x \in \mathbb{R} \to d \in \mathbb{R} \to s \leq x \to x - d < s \to x - d < s$\lbl{A1}}
\AXC{$\stackrel{\Pi_{4}}{\vdots}$}\noLine
\UIC{$\varphi_{42}$}
\BIC{$\varphi_{15}$}
\AXC{$\stackrel{\Pi_{9}}{\vdots}$}\noLine
\UIC{$\varphi_{15} \to s \in \mathbb{R} \to \varphi_{10}$}
\BIC{$s \in \mathbb{R} \to \varphi_{10}$}
\AXC{$\stackrel{\Pi_{11}}{\vdots}$}\noLine
\UIC{$(s \in \mathbb{R} \to \varphi_{10}) \to \varphi_{37}$}
\BIC{$\varphi_{37}$}
\end{ptree}
\noindent$\Pi_{13}$:
\begin{ptree}
\AXC{$\forall x\,\forall y\,(x \in \mathbb{R} \to y \in \mathbb{R} \to -(y - x) = x - y)$\lbl{[F14]}}
\AXC{$\varphi_{30}$\lbl{A2}}
\BIC{$\forall y\,(x \in \mathbb{R} \to y \in \mathbb{R} \to -(y - x) = x - y)$}
\AXC{$\varphi_{28}$\lbl{A2}}
\BIC{$x \in \mathbb{R} \to s \in \mathbb{R} \to -(s - x) = x - s$}
\AXC{$(x \in \mathbb{R} \to s \in \mathbb{R} \to -(s - x) = x - s) \to \varphi_{49}$\lbl{A1}}
\BIC{$\varphi_{49}$}
\end{ptree}
\noindent$\Pi_{14}$:
\begin{ptree}
\AXC{$s \in \mathbb{R} \to x \in \mathbb{R} \to d \in \mathbb{R} \to s \leq x \to x - d < s \to s \in \mathbb{R}$\lbl{A1}}
\AXC{$s \in \mathbb{R} \to x \in \mathbb{R} \to d \in \mathbb{R} \to s \leq x \to x - d < s \to x \in \mathbb{R}$\lbl{A1}}
\AXC{$\stackrel{\Pi_{13}}{\vdots}$}\noLine
\UIC{$\varphi_{49}$}
\BIC{$\varphi_{43}$}
\BIC{$\varphi_{18}$}
\end{ptree}
\noindent$\Pi_{15}$:
\begin{ptree}
\AXC{$\stackrel{\Pi_{1}}{\vdots}$}\noLine
\UIC{$-(-d) = d \to -(s - x) < -(-d) \to -(s - x) < d$}
\AXC{$\stackrel{\Pi_{12}}{\vdots}$}\noLine
\UIC{$\varphi_{37}$}
\BIC{$\varphi_{16}$}
\AXC{$\stackrel{\Pi_{14}}{\vdots}$}\noLine
\UIC{$\varphi_{18}$}
\AXC{$\varphi_{18} \to \varphi_{16} \to \varphi_{33}$\lbl{A1}}
\BIC{$\varphi_{16} \to \varphi_{33}$}
\BIC{$\varphi_{33}$}
\end{ptree}
\noindent$\Pi_{16}$:
\begin{ptree}
\AXC{$x = x$\lbl{A5}}
\AXC{$x = y \to x = x \to y = x$\lbl{A6}}
\AXC{$(x = y \to x = x \to y = x) \to x = x \to x = y \to y = x$\lbl{A1}}
\BIC{$x = x \to x = y \to y = x$}
\BIC{$x = y \to y = x$}
\UIC{$\forall y\,(x = y \to y = x)$}
\UIC{$\forall x\,\forall y\,(x = y \to y = x)$}
\AXC{$\varphi_{1}$\lbl{A2}}
\BIC{$\forall y\,(\lvert s - x \rvert = y \to y = \lvert s - x \rvert)$}
\AXC{$\varphi_{13}$\lbl{A2}}
\BIC{$\lvert s - x \rvert = x - s \to x - s = \lvert s - x \rvert$}
\end{ptree}
\noindent$\Pi_{17}$:
\begin{ptree}
\AXC{$y = z \to x = y \to x = z$\lbl{A6}}
\AXC{$(y = z \to x = y \to x = z) \to x = y \to y = z \to x = z$\lbl{A1}}
\BIC{$x = y \to y = z \to x = z$}
\UIC{$\forall z\,(x = y \to y = z \to x = z)$}
\UIC{$\forall y\,\forall z\,(x = y \to y = z \to x = z)$}
\UIC{$\forall x\,\forall y\,\forall z\,(x = y \to y = z \to x = z)$}
\AXC{$\varphi_{6}$\lbl{A2}}
\BIC{$\forall y\,\forall z\,(\lvert s - x \rvert = y \to y = z \to \lvert s - x \rvert = z)$}
\end{ptree}
\noindent$\Pi_{18}$:
\begin{ptree}
\AXC{$\stackrel{\Pi_{17}}{\vdots}$}\noLine
\UIC{$\forall y\,\forall z\,(\lvert s - x \rvert = y \to y = z \to \lvert s - x \rvert = z)$}
\AXC{$\forall y\,\forall z\,(\lvert s - x \rvert = y \to y = z \to \lvert s - x \rvert = z) \to \varphi_{2}$\lbl{A2}}
\BIC{$\varphi_{2}$}
\AXC{$\varphi_{2} \to \varphi_{5}$\lbl{A2}}
\BIC{$\varphi_{5}$}
\end{ptree}
\noindent$\Pi_{19}$:
\begin{ptree}
\AXC{$\forall x\,\forall y\,(x \in \mathbb{R} \to y \in \mathbb{R} \to x \leq y \to 0 \leq y - x)$\lbl{[E40]}}
\AXC{$\varphi_{26}$\lbl{A2}}
\BIC{$\forall y\,(s \in \mathbb{R} \to y \in \mathbb{R} \to s \leq y \to 0 \leq y - s)$}
\AXC{$\varphi_{23}$\lbl{A2}}
\BIC{$s \in \mathbb{R} \to x \in \mathbb{R} \to s \leq x \to 0 \leq x - s$}
\AXC{$(s \in \mathbb{R} \to x \in \mathbb{R} \to s \leq x \to 0 \leq x - s) \to \varphi_{51}$\lbl{A1}}
\BIC{$\varphi_{51}$}
\end{ptree}
\noindent$\Pi_{20}$:
\begin{ptree}
\AXC{$s \in \mathbb{R} \to x \in \mathbb{R} \to d \in \mathbb{R} \to s \leq x \to x - d < s \to x \in \mathbb{R}$\lbl{A1}}
\AXC{$s \in \mathbb{R} \to x \in \mathbb{R} \to d \in \mathbb{R} \to s \leq x \to x - d < s \to s \in \mathbb{R}$\lbl{A1}}
\AXC{$\stackrel{\Pi_{19}}{\vdots}$}\noLine
\UIC{$\varphi_{51}$}
\BIC{$\varphi_{46}$}
\BIC{$\varphi_{36}$}
\end{ptree}
\noindent$\Pi_{21}$:
\begin{ptree}
\AXC{$\forall x\,(0 \leq x \to \lvert x \rvert = x)$\lbl{[A1]}}
\AXC{$\forall x\,(0 \leq x \to \lvert x \rvert = x) \to 0 \leq x - s \to \lvert x - s \rvert = x - s$\lbl{A2}}
\BIC{$0 \leq x - s \to \lvert x - s \rvert = x - s$}
\end{ptree}
\noindent$\Pi_{22}$:
\begin{ptree}
\AXC{$\stackrel{\Pi_{21}}{\vdots}$}\noLine
\UIC{$0 \leq x - s \to \lvert x - s \rvert = x - s$}
\AXC{$(0 \leq x - s \to \lvert x - s \rvert = x - s) \to \varphi_{41}$\lbl{A1}}
\BIC{$\varphi_{41}$}
\end{ptree}
\noindent$\Pi_{23}$:
\begin{ptree}
\AXC{$s \in \mathbb{R} \to x \in \mathbb{R} \to d \in \mathbb{R} \to s \leq x \to x - d < s \to s \leq x$\lbl{A1}}
\AXC{$\stackrel{\Pi_{20}}{\vdots}$}\noLine
\UIC{$\varphi_{36}$}
\BIC{$s \in \mathbb{R} \to x \in \mathbb{R} \to d \in \mathbb{R} \to s \leq x \to x - d < s \to 0 \leq x - s$}
\AXC{$\stackrel{\Pi_{22}}{\vdots}$}\noLine
\UIC{$\varphi_{41}$}
\BIC{$\varphi_{19}$}
\end{ptree}
\noindent$\Pi_{24}$:
\begin{ptree}
\AXC{$\forall x\,\forall y\,(x \in \mathbb{R} \to y \in \mathbb{R} \to \lvert x - y \rvert = \lvert y - x \rvert)$\lbl{[A13]}}
\AXC{$\varphi_{32}$\lbl{A2}}
\BIC{$\forall y\,(s \in \mathbb{R} \to y \in \mathbb{R} \to \lvert s - y \rvert = \lvert y - s \rvert)$}
\AXC{$\varphi_{31}$\lbl{A2}}
\BIC{$s \in \mathbb{R} \to x \in \mathbb{R} \to \lvert s - x \rvert = \lvert x - s \rvert$}
\end{ptree}
\noindent$\Pi_{25}$:
\begin{ptree}
\AXC{$\stackrel{\Pi_{24}}{\vdots}$}\noLine
\UIC{$s \in \mathbb{R} \to x \in \mathbb{R} \to \lvert s - x \rvert = \lvert x - s \rvert$}
\AXC{$(s \in \mathbb{R} \to x \in \mathbb{R} \to \lvert s - x \rvert = \lvert x - s \rvert) \to \varphi_{50}$\lbl{A1}}
\BIC{$\varphi_{50}$}
\end{ptree}
\noindent$\Pi_{26}$:
\begin{ptree}
\AXC{$s \in \mathbb{R} \to x \in \mathbb{R} \to d \in \mathbb{R} \to s \leq x \to x - d < s \to x \in \mathbb{R}$\lbl{A1}}
\AXC{$s \in \mathbb{R} \to x \in \mathbb{R} \to d \in \mathbb{R} \to s \leq x \to x - d < s \to s \in \mathbb{R}$\lbl{A1}}
\AXC{$\stackrel{\Pi_{25}}{\vdots}$}\noLine
\UIC{$\varphi_{50}$}
\BIC{$\varphi_{44}$}
\BIC{$s \in \mathbb{R} \to \varphi_{7}$}
\end{ptree}
\noindent$\Pi_{27}$:
\begin{ptree}
\AXC{$\stackrel{\Pi_{18}}{\vdots}$}\noLine
\UIC{$\varphi_{5}$}
\AXC{$\stackrel{\Pi_{23}}{\vdots}$}\noLine
\UIC{$\varphi_{19}$}
\AXC{$\stackrel{\Pi_{26}}{\vdots}$}\noLine
\UIC{$s \in \mathbb{R} \to \varphi_{7}$}
\AXC{$(s \in \mathbb{R} \to \varphi_{7}) \to \varphi_{19} \to \varphi_{5} \to \varphi_{20}$\lbl{A1}}
\BIC{$\varphi_{19} \to \varphi_{5} \to \varphi_{20}$}
\BIC{$\varphi_{5} \to \varphi_{20}$}
\BIC{$\varphi_{20}$}
\AXC{$\varphi_{20} \to (\lvert s - x \rvert = x - s \to x - s = \lvert s - x \rvert) \to \varphi_{21}$\lbl{A1}}
\BIC{$(\lvert s - x \rvert = x - s \to x - s = \lvert s - x \rvert) \to \varphi_{21}$}
\end{ptree}
\noindent$\Pi_{28}$:
\begin{ptree}
\AXC{$\stackrel{\Pi_{16}}{\vdots}$}\noLine
\UIC{$\lvert s - x \rvert = x - s \to x - s = \lvert s - x \rvert$}
\AXC{$\stackrel{\Pi_{27}}{\vdots}$}\noLine
\UIC{$(\lvert s - x \rvert = x - s \to x - s = \lvert s - x \rvert) \to \varphi_{21}$}
\BIC{$\varphi_{21}$}
\AXC{$\varphi_{21} \to \varphi_{45}$\lbl{A1}}
\BIC{$\varphi_{45}$}
\end{ptree}
\noindent$\Pi_{29}$:
\begin{ptree}
\AXC{$z = w \to z < d \to w < d$\lbl{A6}}
\UIC{$\forall w\,(z = w \to z < d \to w < d)$}
\UIC{$\forall z\,\forall w\,(z = w \to z < d \to w < d)$}
\AXC{$\varphi_{11}$\lbl{A2}}
\BIC{$\forall w\,(-(s - x) = w \to -(s - x) < d \to w < d)$}
\AXC{$\varphi_{29}$\lbl{A2}}
\BIC{$-(s - x) = x - s \to -(s - x) < d \to x - s < d$}
\AXC{$\stackrel{\Pi_{15}}{\vdots}$}\noLine
\UIC{$\varphi_{33}$}
\BIC{$s \in \mathbb{R} \to x \in \mathbb{R} \to d \in \mathbb{R} \to s \leq x \to x - d < s \to x - s < d$}
\AXC{$\stackrel{\Pi_{28}}{\vdots}$}\noLine
\UIC{$\varphi_{45}$}
\BIC{$(x - s = \lvert s - x \rvert \to x - s < d \to \lvert s - x \rvert < d) \to \varphi_{17}$}
\end{ptree}
\noindent 主木:
\begin{ptree}
\AXC{$z = w \to z < d \to w < d$\lbl{A6}}
\UIC{$\forall w\,(z = w \to z < d \to w < d)$}
\UIC{$\forall z\,\forall w\,(z = w \to z < d \to w < d)$}
\AXC{$\varphi_{3}$\lbl{A2}}
\BIC{$\forall w\,(x - s = w \to x - s < d \to w < d)$}
\AXC{$\varphi_{24}$\lbl{A2}}
\BIC{$x - s = \lvert s - x \rvert \to x - s < d \to \lvert s - x \rvert < d$}
\AXC{$\stackrel{\Pi_{29}}{\vdots}$}\noLine
\UIC{$(x - s = \lvert s - x \rvert \to x - s < d \to \lvert s - x \rvert < d) \to \varphi_{17}$}
\BIC{$\varphi_{17}$}
\UIC{$\forall d\,(\varphi_{17})$}
\UIC{$\forall x\,\forall d\,(\varphi_{17})$}
\UIC{$\forall s\,\forall x\,\forall d\,(\varphi_{17})$}
\end{ptree}
\begin{abbrevlist}
\item $\varphi_{1} :\equiv \forall x\,\forall y\,(x = y \to y = x) \to \forall y\,(\lvert s - x \rvert = y \to y = \lvert s - x \rvert)$
\item $\varphi_{2} :\equiv \forall z\,(\lvert s - x \rvert = \lvert x - s \rvert \to \lvert x - s \rvert = z \to \lvert s - x \rvert = z)$
\item $\varphi_{3} :\equiv \forall z\,\forall w\,(z = w \to z < d \to w < d) \to \forall w\,(x - s = w \to x - s < d \to w < d)$
\item $\varphi_{4} :\equiv \forall w\,(-(-d) = w \to -(s - x) < -(-d) \to -(s - x) < w)$
\item $\varphi_{5} :\equiv \lvert s - x \rvert = \lvert x - s \rvert \to \lvert x - s \rvert = x - s \to \lvert s - x \rvert = x - s$
\item $\varphi_{6} :\equiv \forall x\,\forall y\,\forall z\,(x = y \to y = z \to x = z) \to \forall y\,\forall z\,(\lvert s - x \rvert = y \to y = z \to \lvert s - x \rvert = z)$
\item $\varphi_{7} :\equiv x \in \mathbb{R} \to d \in \mathbb{R} \to s \leq x \to x - d < s \to \lvert s - x \rvert = \lvert x - s \rvert$
\item $\varphi_{8} :\equiv \forall y\,(x \in \mathbb{R} \to d \in \mathbb{R} \to y \in \mathbb{R} \to x - d < y \to -d < y - x)$
\item $\varphi_{9} :\equiv s \in \mathbb{R} \to x \in \mathbb{R} \to d \in \mathbb{R} \to s \leq x \to x - d < s \to -(-d) = d$
\item $\varphi_{10} :\equiv x \in \mathbb{R} \to d \in \mathbb{R} \to s \leq x \to x - d < s \to -(s - x) < -(-d)$
\item $\varphi_{11} :\equiv \forall z\,\forall w\,(z = w \to z < d \to w < d) \to \forall w\,(-(s - x) = w \to -(s - x) < d \to w < d)$
\item $\varphi_{12} :\equiv s \in \mathbb{R} \to x \in \mathbb{R} \to d \in \mathbb{R} \to s \leq x \to x - d < s \to s - x \in \mathbb{R}$
\item $\varphi_{13} :\equiv \forall y\,(\lvert s - x \rvert = y \to y = \lvert s - x \rvert) \to \lvert s - x \rvert = x - s \to x - s = \lvert s - x \rvert$
\item $\varphi_{14} :\equiv \forall e\,\forall y\,(x \in \mathbb{R} \to e \in \mathbb{R} \to y \in \mathbb{R} \to x - e < y \to -e < y - x)$
\item $\varphi_{15} :\equiv s \in \mathbb{R} \to x \in \mathbb{R} \to d \in \mathbb{R} \to s \leq x \to x - d < s \to -d < s - x$
\item $\varphi_{16} :\equiv s \in \mathbb{R} \to x \in \mathbb{R} \to d \in \mathbb{R} \to s \leq x \to x - d < s \to -(s - x) < d$
\item $\varphi_{17} :\equiv s \in \mathbb{R} \to x \in \mathbb{R} \to d \in \mathbb{R} \to s \leq x \to x - d < s \to \lvert s - x \rvert < d$
\item $\varphi_{18} :\equiv s \in \mathbb{R} \to x \in \mathbb{R} \to d \in \mathbb{R} \to s \leq x \to x - d < s \to -(s - x) = x - s$
\item $\varphi_{19} :\equiv s \in \mathbb{R} \to x \in \mathbb{R} \to d \in \mathbb{R} \to s \leq x \to x - d < s \to \lvert x - s \rvert = x - s$
\item $\varphi_{20} :\equiv s \in \mathbb{R} \to x \in \mathbb{R} \to d \in \mathbb{R} \to s \leq x \to x - d < s \to \lvert s - x \rvert = x - s$
\item $\varphi_{21} :\equiv s \in \mathbb{R} \to x \in \mathbb{R} \to d \in \mathbb{R} \to s \leq x \to x - d < s \to x - s = \lvert s - x \rvert$
\item $\varphi_{22} :\equiv -d \in \mathbb{R} \to s - x \in \mathbb{R} \to -d < s - x \to -(s - x) < -(-d)$
\item $\varphi_{23} :\equiv \forall y\,(s \in \mathbb{R} \to y \in \mathbb{R} \to s \leq y \to 0 \leq y - s) \to s \in \mathbb{R} \to x \in \mathbb{R} \to s \leq x \to 0 \leq x - s$
\item $\varphi_{24} :\equiv \forall w\,(x - s = w \to x - s < d \to w < d) \to x - s = \lvert s - x \rvert \to x - s < d \to \lvert s - x \rvert < d$
\item $\varphi_{25} :\equiv \forall y\,(s \in \mathbb{R} \to y \in \mathbb{R} \to s - y \in \mathbb{R}) \to s \in \mathbb{R} \to x \in \mathbb{R} \to s - x \in \mathbb{R}$
\item $\varphi_{26} :\equiv \forall x\,\forall y\,(x \in \mathbb{R} \to y \in \mathbb{R} \to x \leq y \to 0 \leq y - x) \to \forall y\,(s \in \mathbb{R} \to y \in \mathbb{R} \to s \leq y \to 0 \leq y - s)$
\item $\varphi_{27} :\equiv \forall x\,\forall y\,(x \in \mathbb{R} \to y \in \mathbb{R} \to x - y \in \mathbb{R}) \to \forall y\,(s \in \mathbb{R} \to y \in \mathbb{R} \to s - y \in \mathbb{R})$
\item $\varphi_{28} :\equiv \forall y\,(x \in \mathbb{R} \to y \in \mathbb{R} \to -(y - x) = x - y) \to x \in \mathbb{R} \to s \in \mathbb{R} \to -(s - x) = x - s$
\item $\varphi_{29} :\equiv \forall w\,(-(s - x) = w \to -(s - x) < d \to w < d) \to -(s - x) = x - s \to -(s - x) < d \to x - s < d$
\item $\varphi_{30} :\equiv \forall x\,\forall y\,(x \in \mathbb{R} \to y \in \mathbb{R} \to -(y - x) = x - y) \to \forall y\,(x \in \mathbb{R} \to y \in \mathbb{R} \to -(y - x) = x - y)$
\item $\varphi_{31} :\equiv \forall y\,(s \in \mathbb{R} \to y \in \mathbb{R} \to \lvert s - y \rvert = \lvert y - s \rvert) \to s \in \mathbb{R} \to x \in \mathbb{R} \to \lvert s - x \rvert = \lvert x - s \rvert$
\item $\varphi_{32} :\equiv \forall x\,\forall y\,(x \in \mathbb{R} \to y \in \mathbb{R} \to \lvert x - y \rvert = \lvert y - x \rvert) \to \forall y\,(s \in \mathbb{R} \to y \in \mathbb{R} \to \lvert s - y \rvert = \lvert y - s \rvert)$
\item $\varphi_{33} :\equiv (-(s - x) = x - s \to -(s - x) < d \to x - s < d) \to s \in \mathbb{R} \to x \in \mathbb{R} \to d \in \mathbb{R} \to s \leq x \to x - d < s \to x - s < d$
\item $\varphi_{34} :\equiv \varphi_{4} \to -(-d) = d \to -(s - x) < -(-d) \to -(s - x) < d$
\item $\varphi_{35} :\equiv \forall x\,\forall y\,(x \in \mathbb{R} \to y \in \mathbb{R} \to x < y \to -y < -x) \to \forall y\,(-d \in \mathbb{R} \to y \in \mathbb{R} \to -d < y \to -y < -(-d))$
\item $\varphi_{36} :\equiv (s \in \mathbb{R} \to x \in \mathbb{R} \to d \in \mathbb{R} \to s \leq x \to x - d < s \to s \leq x) \to s \in \mathbb{R} \to x \in \mathbb{R} \to d \in \mathbb{R} \to s \leq x \to x - d < s \to 0 \leq x - s$
\item $\varphi_{37} :\equiv (-(-d) = d \to -(s - x) < -(-d) \to -(s - x) < d) \to \varphi_{16}$
\item $\varphi_{38} :\equiv (s \in \mathbb{R} \to x \in \mathbb{R} \to d \in \mathbb{R} \to s \leq x \to x - d < s \to d \in \mathbb{R}) \to \varphi_{9}$
\item $\varphi_{39} :\equiv (s \in \mathbb{R} \to x \in \mathbb{R} \to d \in \mathbb{R} \to s \leq x \to x - d < s \to x \in \mathbb{R}) \to \varphi_{12}$
\item $\varphi_{40} :\equiv (s \in \mathbb{R} \to x \in \mathbb{R} \to d \in \mathbb{R} \to s \leq x \to x - d < s \to d \in \mathbb{R}) \to s \in \mathbb{R} \to x \in \mathbb{R} \to d \in \mathbb{R} \to s \leq x \to x - d < s \to -d \in \mathbb{R}$
\item $\varphi_{41} :\equiv (s \in \mathbb{R} \to x \in \mathbb{R} \to d \in \mathbb{R} \to s \leq x \to x - d < s \to 0 \leq x - s) \to \varphi_{19}$
\item $\varphi_{42} :\equiv (s \in \mathbb{R} \to x \in \mathbb{R} \to d \in \mathbb{R} \to s \leq x \to x - d < s \to x - d < s) \to \varphi_{15}$
\item $\varphi_{43} :\equiv (s \in \mathbb{R} \to x \in \mathbb{R} \to d \in \mathbb{R} \to s \leq x \to x - d < s \to s \in \mathbb{R}) \to \varphi_{18}$
\item $\varphi_{44} :\equiv (s \in \mathbb{R} \to x \in \mathbb{R} \to d \in \mathbb{R} \to s \leq x \to x - d < s \to x \in \mathbb{R}) \to s \in \mathbb{R} \to \varphi_{7}$
\item $\varphi_{45} :\equiv (s \in \mathbb{R} \to x \in \mathbb{R} \to d \in \mathbb{R} \to s \leq x \to x - d < s \to x - s < d) \to (x - s = \lvert s - x \rvert \to x - s < d \to \lvert s - x \rvert < d) \to \varphi_{17}$
\item $\varphi_{46} :\equiv (s \in \mathbb{R} \to x \in \mathbb{R} \to d \in \mathbb{R} \to s \leq x \to x - d < s \to x \in \mathbb{R}) \to \varphi_{36}$
\item $\varphi_{47} :\equiv (s \in \mathbb{R} \to x \in \mathbb{R} \to d \in \mathbb{R} \to s \leq x \to x - d < s \to s \in \mathbb{R}) \to \varphi_{39}$
\item $\varphi_{48} :\equiv (s \in \mathbb{R} \to x \in \mathbb{R} \to d \in \mathbb{R} \to s \leq x \to x - d < s \to s \in \mathbb{R}) \to \varphi_{42}$
\item $\varphi_{49} :\equiv (s \in \mathbb{R} \to x \in \mathbb{R} \to d \in \mathbb{R} \to s \leq x \to x - d < s \to x \in \mathbb{R}) \to \varphi_{43}$
\item $\varphi_{50} :\equiv (s \in \mathbb{R} \to x \in \mathbb{R} \to d \in \mathbb{R} \to s \leq x \to x - d < s \to s \in \mathbb{R}) \to \varphi_{44}$
\item $\varphi_{51} :\equiv (s \in \mathbb{R} \to x \in \mathbb{R} \to d \in \mathbb{R} \to s \leq x \to x - d < s \to s \in \mathbb{R}) \to \varphi_{46}$
\item $\varphi_{52} :\equiv (s \in \mathbb{R} \to x \in \mathbb{R} \to d \in \mathbb{R} \to s \leq x \to x - d < s \to d \in \mathbb{R}) \to \varphi_{48}$
\item $\varphi_{53} :\equiv (s \in \mathbb{R} \to x \in \mathbb{R} \to d \in \mathbb{R} \to s \leq x \to x - d < s \to -d \in \mathbb{R}) \to \varphi_{12} \to \varphi_{15} \to s \in \mathbb{R} \to \varphi_{10}$
\item $\varphi_{54} :\equiv (s \in \mathbb{R} \to x \in \mathbb{R} \to d \in \mathbb{R} \to s \leq x \to x - d < s \to x \in \mathbb{R}) \to \varphi_{52}$
\item $\varphi_{55} :\equiv (x \in \mathbb{R} \to d \in \mathbb{R} \to s \in \mathbb{R} \to x - d < s \to -d < s - x) \to \varphi_{54}$
\end{abbrevlist}
\end{fullproof}

% 自動生成 (prover/build.js): NegSubLt  ノード数 86
\begin{fullproof}{[E42]}{$\forall x\,\forall m\,\forall y\,(x \in \mathbb{R} \to m \in \mathbb{R} \to y \in \mathbb{R} \to -(x - m) < y - x \to m < y)$}{86}\phantomsection\label{tree:NegSubLt}
\noindent$\Pi_{1}$:
\begin{ptree}
\AXC{$z = w \to m + -x < z \to m + -x < w$\lbl{A6}}
\UIC{$\forall w\,(z = w \to m + -x < z \to m + -x < w)$}
\UIC{$\forall z\,\forall w\,(z = w \to m + -x < z \to m + -x < w)$}
\AXC{$\varphi_{17}$\lbl{A2}}
\BIC{$\forall w\,(y - x = w \to m + -x < y - x \to m + -x < w)$}
\AXC{$\varphi_{20}$\lbl{A2}}
\BIC{$y - x = y + -x \to m + -x < y - x \to m + -x < y + -x$}
\end{ptree}
\noindent$\Pi_{2}$:
\begin{ptree}
\AXC{$\forall x\,\forall y\,(x \in \mathbb{R} \to y \in \mathbb{R} \to -(y - x) = x - y)$\lbl{[F14]}}
\AXC{$\varphi_{16}$\lbl{A2}}
\BIC{$\forall y\,(m \in \mathbb{R} \to y \in \mathbb{R} \to -(y - m) = m - y)$}
\AXC{$\varphi_{15}$\lbl{A2}}
\BIC{$m \in \mathbb{R} \to x \in \mathbb{R} \to -(x - m) = m - x$}
\AXC{$(m \in \mathbb{R} \to x \in \mathbb{R} \to -(x - m) = m - x) \to \varphi_{29}$\lbl{A1}}
\BIC{$\varphi_{29}$}
\end{ptree}
\noindent$\Pi_{3}$:
\begin{ptree}
\AXC{$x \in \mathbb{R} \to m \in \mathbb{R} \to y \in \mathbb{R} \to -(x - m) < y - x \to x \in \mathbb{R}$\lbl{A1}}
\AXC{$x \in \mathbb{R} \to m \in \mathbb{R} \to y \in \mathbb{R} \to -(x - m) < y - x \to m \in \mathbb{R}$\lbl{A1}}
\AXC{$\stackrel{\Pi_{2}}{\vdots}$}\noLine
\UIC{$\varphi_{29}$}
\BIC{$\varphi_{25}$}
\BIC{$x \in \mathbb{R} \to \varphi_{3}$}
\end{ptree}
\noindent$\Pi_{4}$:
\begin{ptree}
\AXC{$x \in \mathbb{R} \to \varphi_{5}$\lbl{A1}}
\AXC{$\stackrel{\Pi_{3}}{\vdots}$}\noLine
\UIC{$x \in \mathbb{R} \to \varphi_{3}$}
\AXC{$(x \in \mathbb{R} \to \varphi_{3}) \to (x \in \mathbb{R} \to \varphi_{5}) \to \varphi_{22}$\lbl{A1}}
\BIC{$(x \in \mathbb{R} \to \varphi_{5}) \to \varphi_{22}$}
\BIC{$\varphi_{22}$}
\end{ptree}
\noindent$\Pi_{5}$:
\begin{ptree}
\AXC{$\forall x\,\forall y\,(x - y = x + -y)$\lbl{Def$-$}}
\AXC{$\forall x\,\forall y\,(x - y = x + -y) \to \forall y\,(m - y = m + -y)$\lbl{A2}}
\BIC{$\forall y\,(m - y = m + -y)$}
\end{ptree}
\noindent$\Pi_{6}$:
\begin{ptree}
\AXC{$\stackrel{\Pi_{5}}{\vdots}$}\noLine
\UIC{$\forall y\,(m - y = m + -y)$}
\AXC{$\forall y\,(m - y = m + -y) \to m - x = m + -x$\lbl{A2}}
\BIC{$m - x = m + -x$}
\AXC{$m - x = m + -x \to \varphi_{28}$\lbl{A1}}
\BIC{$\varphi_{28}$}
\end{ptree}
\noindent$\Pi_{7}$:
\begin{ptree}
\AXC{$z = w \to z < y - x \to w < y - x$\lbl{A6}}
\UIC{$\forall w\,(z = w \to z < y - x \to w < y - x)$}
\UIC{$\forall z\,\forall w\,(z = w \to z < y - x \to w < y - x)$}
\AXC{$\varphi_{14}$\lbl{A2}}
\BIC{$\forall w\,(-(x - m) = w \to -(x - m) < y - x \to w < y - x)$}
\AXC{$\varphi_{19}$\lbl{A2}}
\BIC{$-(x - m) = m - x \to -(x - m) < y - x \to m - x < y - x$}
\AXC{$\stackrel{\Pi_{4}}{\vdots}$}\noLine
\UIC{$\varphi_{22}$}
\BIC{$\varphi_{10}$}
\AXC{$\stackrel{\Pi_{6}}{\vdots}$}\noLine
\UIC{$\varphi_{28}$}
\BIC{$(m - x = m + -x \to m - x < y - x \to m + -x < y - x) \to \varphi_{12}$}
\end{ptree}
\noindent$\Pi_{8}$:
\begin{ptree}
\AXC{$z = w \to z < y - x \to w < y - x$\lbl{A6}}
\UIC{$\forall w\,(z = w \to z < y - x \to w < y - x)$}
\UIC{$\forall z\,\forall w\,(z = w \to z < y - x \to w < y - x)$}
\AXC{$\varphi_{11}$\lbl{A2}}
\BIC{$\forall w\,(m - x = w \to m - x < y - x \to w < y - x)$}
\AXC{$\varphi_{18}$\lbl{A2}}
\BIC{$m - x = m + -x \to m - x < y - x \to m + -x < y - x$}
\AXC{$\stackrel{\Pi_{7}}{\vdots}$}\noLine
\UIC{$(m - x = m + -x \to m - x < y - x \to m + -x < y - x) \to \varphi_{12}$}
\BIC{$\varphi_{12}$}
\end{ptree}
\noindent$\Pi_{9}$:
\begin{ptree}
\AXC{$\forall x\,\forall y\,(x - y = x + -y)$\lbl{Def$-$}}
\AXC{$\varphi_{1}$\lbl{A2}}
\BIC{$\forall y\,(v_{778} - y = v_{778} + -y)$}
\AXC{$\varphi_{2}$\lbl{A2}}
\BIC{$v_{778} - v_{779} = v_{778} + -v_{779}$}
\UIC{$\forall v_{779}\,(v_{778} - v_{779} = v_{778} + -v_{779})$}
\UIC{$\forall v_{778}\,\forall v_{779}\,(v_{778} - v_{779} = v_{778} + -v_{779})$}
\AXC{$\varphi_{8}$\lbl{A2}}
\BIC{$\forall v_{779}\,(y - v_{779} = y + -v_{779})$}
\AXC{$\forall v_{779}\,(y - v_{779} = y + -v_{779}) \to y - x = y + -x$\lbl{A2}}
\BIC{$y - x = y + -x$}
\end{ptree}
\noindent$\Pi_{10}$:
\begin{ptree}
\AXC{$\stackrel{\Pi_{1}}{\vdots}$}\noLine
\UIC{$y - x = y + -x \to m + -x < y - x \to m + -x < y + -x$}
\AXC{$\stackrel{\Pi_{8}}{\vdots}$}\noLine
\UIC{$\varphi_{12}$}
\AXC{$\stackrel{\Pi_{9}}{\vdots}$}\noLine
\UIC{$y - x = y + -x$}
\AXC{$y - x = y + -x \to \varphi_{12} \to \varphi_{27}$\lbl{A1}}
\BIC{$\varphi_{12} \to \varphi_{27}$}
\BIC{$\varphi_{27}$}
\BIC{$\varphi_{13}$}
\end{ptree}
\noindent$\Pi_{11}$:
\begin{ptree}
\AXC{$\forall x \in \mathbb{R}\,(-x \in \mathbb{R})$\lbl{(R1)}}
\AXC{$\forall x \in \mathbb{R}\,(-x \in \mathbb{R}) \to x \in \mathbb{R} \to -x \in \mathbb{R}$\lbl{A2}}
\BIC{$x \in \mathbb{R} \to -x \in \mathbb{R}$}
\AXC{$(x \in \mathbb{R} \to -x \in \mathbb{R}) \to \varphi_{24}$\lbl{A1}}
\BIC{$\varphi_{24}$}
\end{ptree}
\noindent$\Pi_{12}$:
\begin{ptree}
\AXC{$\varphi_{7}$\lbl{[E4]}}
\AXC{$\varphi_{7} \to \varphi_{4}$\lbl{A2}}
\BIC{$\varphi_{4}$}
\AXC{$\varphi_{21}$\lbl{A2}}
\BIC{$\forall z\,(m \in \mathbb{R} \to y \in \mathbb{R} \to z \in \mathbb{R} \to m + z < y + z \to m < y)$}
\AXC{$\varphi_{23}$\lbl{A2}}
\BIC{$m \in \mathbb{R} \to y \in \mathbb{R} \to -x \in \mathbb{R} \to m + -x < y + -x \to m < y$}
\AXC{$\varphi_{32}$\lbl{A1}}
\BIC{$\varphi_{31}$}
\end{ptree}
\noindent$\Pi_{13}$:
\begin{ptree}
\AXC{$x \in \mathbb{R} \to m \in \mathbb{R} \to y \in \mathbb{R} \to -(x - m) < y - x \to y \in \mathbb{R}$\lbl{A1}}
\AXC{$x \in \mathbb{R} \to m \in \mathbb{R} \to y \in \mathbb{R} \to -(x - m) < y - x \to m \in \mathbb{R}$\lbl{A1}}
\AXC{$\stackrel{\Pi_{12}}{\vdots}$}\noLine
\UIC{$\varphi_{31}$}
\BIC{$\varphi_{30}$}
\BIC{$\varphi_{9} \to \varphi_{26}$}
\end{ptree}
\noindent 主木:
\begin{ptree}
\AXC{$\stackrel{\Pi_{10}}{\vdots}$}\noLine
\UIC{$\varphi_{13}$}
\AXC{$x \in \mathbb{R} \to m \in \mathbb{R} \to y \in \mathbb{R} \to -(x - m) < y - x \to x \in \mathbb{R}$\lbl{A1}}
\AXC{$\stackrel{\Pi_{11}}{\vdots}$}\noLine
\UIC{$\varphi_{24}$}
\BIC{$\varphi_{9}$}
\AXC{$\stackrel{\Pi_{13}}{\vdots}$}\noLine
\UIC{$\varphi_{9} \to \varphi_{26}$}
\BIC{$\varphi_{26}$}
\BIC{$x \in \mathbb{R} \to m \in \mathbb{R} \to y \in \mathbb{R} \to -(x - m) < y - x \to m < y$}
\UIC{$\varphi_{6}$}
\UIC{$\forall m\,(\varphi_{6})$}
\UIC{$\forall x\,\forall m\,(\varphi_{6})$}
\end{ptree}
\begin{abbrevlist}
\item $\varphi_{1} :\equiv \forall x\,\forall y\,(x - y = x + -y) \to \forall y\,(v_{778} - y = v_{778} + -y)$
\item $\varphi_{2} :\equiv \forall y\,(v_{778} - y = v_{778} + -y) \to v_{778} - v_{779} = v_{778} + -v_{779}$
\item $\varphi_{3} :\equiv m \in \mathbb{R} \to y \in \mathbb{R} \to -(x - m) < y - x \to -(x - m) = m - x$
\item $\varphi_{4} :\equiv \forall y\,\forall z\,(m \in \mathbb{R} \to y \in \mathbb{R} \to z \in \mathbb{R} \to m + z < y + z \to m < y)$
\item $\varphi_{5} :\equiv m \in \mathbb{R} \to y \in \mathbb{R} \to -(x - m) < y - x \to -(x - m) < y - x$
\item $\varphi_{6} :\equiv \forall y\,(x \in \mathbb{R} \to m \in \mathbb{R} \to y \in \mathbb{R} \to -(x - m) < y - x \to m < y)$
\item $\varphi_{7} :\equiv \forall x\,\forall y\,\forall z\,(x \in \mathbb{R} \to y \in \mathbb{R} \to z \in \mathbb{R} \to x + z < y + z \to x < y)$
\item $\varphi_{8} :\equiv \forall v_{778}\,\forall v_{779}\,(v_{778} - v_{779} = v_{778} + -v_{779}) \to \forall v_{779}\,(y - v_{779} = y + -v_{779})$
\item $\varphi_{9} :\equiv x \in \mathbb{R} \to m \in \mathbb{R} \to y \in \mathbb{R} \to -(x - m) < y - x \to -x \in \mathbb{R}$
\item $\varphi_{10} :\equiv x \in \mathbb{R} \to m \in \mathbb{R} \to y \in \mathbb{R} \to -(x - m) < y - x \to m - x < y - x$
\item $\varphi_{11} :\equiv \forall z\,\forall w\,(z = w \to z < y - x \to w < y - x) \to \forall w\,(m - x = w \to m - x < y - x \to w < y - x)$
\item $\varphi_{12} :\equiv x \in \mathbb{R} \to m \in \mathbb{R} \to y \in \mathbb{R} \to -(x - m) < y - x \to m + -x < y - x$
\item $\varphi_{13} :\equiv x \in \mathbb{R} \to m \in \mathbb{R} \to y \in \mathbb{R} \to -(x - m) < y - x \to m + -x < y + -x$
\item $\varphi_{14} :\equiv \forall z\,\forall w\,(z = w \to z < y - x \to w < y - x) \to \forall w\,(-(x - m) = w \to -(x - m) < y - x \to w < y - x)$
\item $\varphi_{15} :\equiv \forall y\,(m \in \mathbb{R} \to y \in \mathbb{R} \to -(y - m) = m - y) \to m \in \mathbb{R} \to x \in \mathbb{R} \to -(x - m) = m - x$
\item $\varphi_{16} :\equiv \forall x\,\forall y\,(x \in \mathbb{R} \to y \in \mathbb{R} \to -(y - x) = x - y) \to \forall y\,(m \in \mathbb{R} \to y \in \mathbb{R} \to -(y - m) = m - y)$
\item $\varphi_{17} :\equiv \forall z\,\forall w\,(z = w \to m + -x < z \to m + -x < w) \to \forall w\,(y - x = w \to m + -x < y - x \to m + -x < w)$
\item $\varphi_{18} :\equiv \forall w\,(m - x = w \to m - x < y - x \to w < y - x) \to m - x = m + -x \to m - x < y - x \to m + -x < y - x$
\item $\varphi_{19} :\equiv \forall w\,(-(x - m) = w \to -(x - m) < y - x \to w < y - x) \to -(x - m) = m - x \to -(x - m) < y - x \to m - x < y - x$
\item $\varphi_{20} :\equiv \forall w\,(y - x = w \to m + -x < y - x \to m + -x < w) \to y - x = y + -x \to m + -x < y - x \to m + -x < y + -x$
\item $\varphi_{21} :\equiv \varphi_{4} \to \forall z\,(m \in \mathbb{R} \to y \in \mathbb{R} \to z \in \mathbb{R} \to m + z < y + z \to m < y)$
\item $\varphi_{22} :\equiv (-(x - m) = m - x \to -(x - m) < y - x \to m - x < y - x) \to \varphi_{10}$
\item $\varphi_{23} :\equiv \forall z\,(m \in \mathbb{R} \to y \in \mathbb{R} \to z \in \mathbb{R} \to m + z < y + z \to m < y) \to m \in \mathbb{R} \to y \in \mathbb{R} \to -x \in \mathbb{R} \to m + -x < y + -x \to m < y$
\item $\varphi_{24} :\equiv (x \in \mathbb{R} \to m \in \mathbb{R} \to y \in \mathbb{R} \to -(x - m) < y - x \to x \in \mathbb{R}) \to \varphi_{9}$
\item $\varphi_{25} :\equiv (x \in \mathbb{R} \to m \in \mathbb{R} \to y \in \mathbb{R} \to -(x - m) < y - x \to x \in \mathbb{R}) \to x \in \mathbb{R} \to \varphi_{3}$
\item $\varphi_{26} :\equiv \varphi_{13} \to x \in \mathbb{R} \to m \in \mathbb{R} \to y \in \mathbb{R} \to -(x - m) < y - x \to m < y$
\item $\varphi_{27} :\equiv (y - x = y + -x \to m + -x < y - x \to m + -x < y + -x) \to \varphi_{13}$
\item $\varphi_{28} :\equiv \varphi_{10} \to (m - x = m + -x \to m - x < y - x \to m + -x < y - x) \to \varphi_{12}$
\item $\varphi_{29} :\equiv (x \in \mathbb{R} \to m \in \mathbb{R} \to y \in \mathbb{R} \to -(x - m) < y - x \to m \in \mathbb{R}) \to \varphi_{25}$
\item $\varphi_{30} :\equiv (x \in \mathbb{R} \to m \in \mathbb{R} \to y \in \mathbb{R} \to -(x - m) < y - x \to y \in \mathbb{R}) \to \varphi_{9} \to \varphi_{26}$
\item $\varphi_{31} :\equiv (x \in \mathbb{R} \to m \in \mathbb{R} \to y \in \mathbb{R} \to -(x - m) < y - x \to m \in \mathbb{R}) \to \varphi_{30}$
\item $\varphi_{32} :\equiv (m \in \mathbb{R} \to y \in \mathbb{R} \to -x \in \mathbb{R} \to m + -x < y + -x \to m < y) \to \varphi_{31}$
\end{abbrevlist}
\end{fullproof}

% 自動生成 (prover/build.js): SubRefl  ノード数 11
\begin{fullproof}{[E43]}{$\forall A\,(A \subseteq A)$}{11}\phantomsection\label{tree:SubRefl}
\noindent$\Pi_{1}$:
\begin{ptree}
\AXC{$\forall x\,\forall y\,(x \subseteq y \leftrightarrow \forall u \in x\,(u \in y))$\lbl{Def$\subseteq$}}
\AXC{$\varphi_{2}$\lbl{A2}}
\BIC{$\forall y\,(A \subseteq y \leftrightarrow \forall u \in A\,(u \in y))$}
\AXC{$\varphi_{1}$\lbl{A2}}
\BIC{$A \subseteq A \leftrightarrow \forall u \in A\,(u \in A)$}
\AXC{$(A \subseteq A \leftrightarrow \forall u \in A\,(u \in A)) \to \forall u \in A\,(u \in A) \to A \subseteq A$\lbl{A1}}
\BIC{$\forall u \in A\,(u \in A) \to A \subseteq A$}
\end{ptree}
\noindent 主木:
\begin{ptree}
\AXC{$u \in A \to u \in A$\lbl{A1}}
\UIC{$\forall u \in A\,(u \in A)$}
\AXC{$\stackrel{\Pi_{1}}{\vdots}$}\noLine
\UIC{$\forall u \in A\,(u \in A) \to A \subseteq A$}
\BIC{$A \subseteq A$}
\UIC{$\forall A\,(A \subseteq A)$}
\end{ptree}
\begin{abbrevlist}
\item $\varphi_{1} :\equiv \forall y\,(A \subseteq y \leftrightarrow \forall u \in A\,(u \in y)) \to (A \subseteq A \leftrightarrow \forall u \in A\,(u \in A))$
\item $\varphi_{2} :\equiv \forall x\,\forall y\,(x \subseteq y \leftrightarrow \forall u \in x\,(u \in y)) \to \forall y\,(A \subseteq y \leftrightarrow \forall u \in A\,(u \in y))$
\end{abbrevlist}
\end{fullproof}

% 自動生成 (prover/build.js): ContLim  ノード数 347
\begin{fullproof}{命題(定義2.1)}{$\forall f\,\forall c\,(\mathrm{PNbhd}(f,c) \to c \in \mathrm{dom}(f) \to (\mathrm{Cont}(f,c) \leftrightarrow \mathrm{FLim}(f,c,f(c))))$}{347}\phantomsection\label{tree:ContLim}
\noindent$\Pi_{1}$:
\begin{ptree}
\AXC{$\forall f\,\forall c\,\forall \alpha\,(\mathrm{FLim}(f,c,\alpha) \to c \in \mathbb{R})$\lbl{[D3]}}
\AXC{$\varphi_{30}$\lbl{A2}}
\BIC{$\forall c\,\forall \alpha\,(\mathrm{FLim}(f,c,\alpha) \to c \in \mathbb{R})$}
\AXC{$\varphi_{23}$\lbl{A2}}
\BIC{$\forall \alpha\,(\mathrm{FLim}(f,c,\alpha) \to c \in \mathbb{R})$}
\AXC{$\varphi_{20}$\lbl{A2}}
\BIC{$\mathrm{FLim}(f,c,f(c)) \to c \in \mathbb{R}$}
\AXC{$\varphi_{79}$\lbl{A1}}
\BIC{$(\mathrm{PNbhd}(f,c) \to \varphi_{10}) \to \varphi_{15}$}
\end{ptree}
\noindent$\Pi_{2}$:
\begin{ptree}
\AXC{$\mathrm{PNbhd}(f,c) \to \varphi_{10}$\lbl{A1}}
\AXC{$\stackrel{\Pi_{1}}{\vdots}$}\noLine
\UIC{$(\mathrm{PNbhd}(f,c) \to \varphi_{10}) \to \varphi_{15}$}
\BIC{$\varphi_{15}$}
\end{ptree}
\noindent$\Pi_{3}$:
\begin{ptree}
\AXC{$\forall f\,\forall c\,\forall \alpha\,(\mathrm{FLim}(f,c,\alpha) \to \mathrm{RFn}(f))$\lbl{[D2]}}
\AXC{$\varphi_{18}$\lbl{A2}}
\BIC{$\forall c\,\forall \alpha\,(\mathrm{FLim}(f,c,\alpha) \to \mathrm{RFn}(f))$}
\AXC{$\varphi_{13}$\lbl{A2}}
\BIC{$\forall \alpha\,(\mathrm{FLim}(f,c,\alpha) \to \mathrm{RFn}(f))$}
\AXC{$\varphi_{11}$\lbl{A2}}
\BIC{$\mathrm{FLim}(f,c,f(c)) \to \mathrm{RFn}(f)$}
\AXC{$\varphi_{69}$\lbl{A1}}
\BIC{$(\mathrm{PNbhd}(f,c) \to \varphi_{10}) \to \varphi_{12}$}
\end{ptree}
\noindent$\Pi_{4}$:
\begin{ptree}
\AXC{$\mathrm{PNbhd}(f,c) \to \varphi_{10}$\lbl{A1}}
\AXC{$\stackrel{\Pi_{3}}{\vdots}$}\noLine
\UIC{$(\mathrm{PNbhd}(f,c) \to \varphi_{10}) \to \varphi_{12}$}
\BIC{$\varphi_{12}$}
\end{ptree}
\noindent$\Pi_{5}$:
\begin{ptree}
\AXC{$\varphi_{98}$\lbl{A1}}
\AXC{$\stackrel{\Pi_{4}}{\vdots}$}\noLine
\UIC{$\varphi_{12}$}
\AXC{$\forall f\,\forall x\,(\mathrm{RFn}(f) \to x \in \mathrm{dom}(f) \to x \in \mathbb{R})$\lbl{[S7]}}
\AXC{$\varphi_{46}$\lbl{A2}}
\BIC{$\forall x\,(\mathrm{RFn}(f) \to x \in \mathrm{dom}(f) \to x \in \mathbb{R})$}
\AXC{$\varphi_{41}$\lbl{A2}}
\BIC{$\mathrm{RFn}(f) \to x \in \mathrm{dom}(f) \to x \in \mathbb{R}$}
\AXC{$\varphi_{168}$\lbl{A1}}
\BIC{$\varphi_{12} \to \varphi_{98} \to \mathrm{PNbhd}(f,c) \to \varphi_{104}$}
\BIC{$\varphi_{98} \to \mathrm{PNbhd}(f,c) \to \varphi_{104}$}
\BIC{$\mathrm{PNbhd}(f,c) \to \varphi_{104}$}
\end{ptree}
\noindent$\Pi_{6}$:
\begin{ptree}
\AXC{$\forall x\,\forall y\,(x \in \mathbb{R} \to y \in \mathbb{R} \to x \neq y \to x - y \neq 0)$\lbl{[A15]}}
\AXC{$\varphi_{59}$\lbl{A2}}
\BIC{$\forall y\,(x \in \mathbb{R} \to y \in \mathbb{R} \to x \neq y \to x - y \neq 0)$}
\AXC{$\varphi_{55}$\lbl{A2}}
\BIC{$x \in \mathbb{R} \to c \in \mathbb{R} \to x \neq c \to x - c \neq 0$}
\AXC{$(x \in \mathbb{R} \to c \in \mathbb{R} \to x \neq c \to x - c \neq 0) \to \varphi_{181}$\lbl{A1}}
\BIC{$\varphi_{181}$}
\end{ptree}
\noindent$\Pi_{7}$:
\begin{ptree}
\AXC{$\mathrm{PNbhd}(f,c) \to \varphi_{106}$\lbl{A1}}
\AXC{$\stackrel{\Pi_{2}}{\vdots}$}\noLine
\UIC{$\varphi_{15}$}
\AXC{$\stackrel{\Pi_{5}}{\vdots}$}\noLine
\UIC{$\mathrm{PNbhd}(f,c) \to \varphi_{104}$}
\AXC{$\stackrel{\Pi_{6}}{\vdots}$}\noLine
\UIC{$\varphi_{181}$}
\BIC{$\varphi_{15} \to (\mathrm{PNbhd}(f,c) \to \varphi_{106}) \to \mathrm{PNbhd}(f,c) \to \varphi_{110}$}
\BIC{$(\mathrm{PNbhd}(f,c) \to \varphi_{106}) \to \mathrm{PNbhd}(f,c) \to \varphi_{110}$}
\BIC{$\mathrm{PNbhd}(f,c) \to \varphi_{110}$}
\end{ptree}
\noindent$\Pi_{8}$:
\begin{ptree}
\AXC{$\stackrel{\Pi_{2}}{\vdots}$}\noLine
\UIC{$\varphi_{15}$}
\AXC{$\stackrel{\Pi_{5}}{\vdots}$}\noLine
\UIC{$\mathrm{PNbhd}(f,c) \to \varphi_{104}$}
\AXC{$\forall x\,\forall y\,(x \in \mathbb{R} \to y \in \mathbb{R} \to x - y \in \mathbb{R})$\lbl{[F1]}}
\AXC{$\varphi_{57}$\lbl{A2}}
\BIC{$\forall y\,(x \in \mathbb{R} \to y \in \mathbb{R} \to x - y \in \mathbb{R})$}
\AXC{$\varphi_{52}$\lbl{A2}}
\BIC{$x \in \mathbb{R} \to c \in \mathbb{R} \to x - c \in \mathbb{R}$}
\AXC{$\varphi_{175}$\lbl{A1}}
\BIC{$(\mathrm{PNbhd}(f,c) \to \varphi_{104}) \to \varphi_{15} \to \mathrm{PNbhd}(f,c) \to \varphi_{118}$}
\BIC{$\varphi_{15} \to \mathrm{PNbhd}(f,c) \to \varphi_{118}$}
\BIC{$\mathrm{PNbhd}(f,c) \to \varphi_{118}$}
\end{ptree}
\noindent$\Pi_{9}$:
\begin{ptree}
\AXC{$\stackrel{\Pi_{8}}{\vdots}$}\noLine
\UIC{$\mathrm{PNbhd}(f,c) \to \varphi_{118}$}
\AXC{$\forall x \in \mathbb{R}\,(x \neq 0 \to 0 < \lvert x \rvert)$\lbl{[A14]}}
\AXC{$\varphi_{48}$\lbl{A2}}
\BIC{$x - c \in \mathbb{R} \to x - c \neq 0 \to 0 < \lvert x - c \rvert$}
\AXC{$(x - c \in \mathbb{R} \to x - c \neq 0 \to 0 < \lvert x - c \rvert) \to \varphi_{180}$\lbl{A1}}
\BIC{$\varphi_{180}$}
\BIC{$(\mathrm{PNbhd}(f,c) \to \varphi_{110}) \to \mathrm{PNbhd}(f,c) \to \varphi_{119}$}
\end{ptree}
\noindent$\Pi_{10}$:
\begin{ptree}
\AXC{$\varphi_{98}$\lbl{A1}}
\AXC{$\varphi_{54} \to \varphi_{49}$\lbl{A2}}
\AXC{$\mathrm{PNbhd}(f,c) \to c \in \mathrm{dom}(f) \to \varphi_{136}$\lbl{A1}}
\AXC{$\varphi_{173}$\lbl{A1}}
\BIC{$\mathrm{PNbhd}(f,c) \to c \in \mathrm{dom}(f) \to \varphi_{132}$}
\AXC{$\varphi_{179}$\lbl{A1}}
\BIC{$(\varphi_{54} \to \varphi_{49}) \to \mathrm{PNbhd}(f,c) \to c \in \mathrm{dom}(f) \to \varphi_{131}$}
\BIC{$\mathrm{PNbhd}(f,c) \to c \in \mathrm{dom}(f) \to \varphi_{131}$}
\AXC{$\varphi_{185}$\lbl{A1}}
\BIC{$\varphi_{98} \to (\mathrm{PNbhd}(f,c) \to \varphi_{119}) \to \varphi_{163}$}
\BIC{$(\mathrm{PNbhd}(f,c) \to \varphi_{119}) \to \varphi_{163}$}
\end{ptree}
\noindent$\Pi_{11}$:
\begin{ptree}
\AXC{$\stackrel{\Pi_{7}}{\vdots}$}\noLine
\UIC{$\mathrm{PNbhd}(f,c) \to \varphi_{110}$}
\AXC{$\stackrel{\Pi_{9}}{\vdots}$}\noLine
\UIC{$(\mathrm{PNbhd}(f,c) \to \varphi_{110}) \to \mathrm{PNbhd}(f,c) \to \varphi_{119}$}
\BIC{$\mathrm{PNbhd}(f,c) \to \varphi_{119}$}
\AXC{$\stackrel{\Pi_{10}}{\vdots}$}\noLine
\UIC{$(\mathrm{PNbhd}(f,c) \to \varphi_{119}) \to \varphi_{163}$}
\BIC{$\varphi_{163}$}
\end{ptree}
\noindent$\Pi_{12}$:
\begin{ptree}
\AXC{$x = x$\lbl{A5}}
\AXC{$x = y \to x = x \to y = x$\lbl{A6}}
\AXC{$(x = y \to x = x \to y = x) \to x = x \to x = y \to y = x$\lbl{A1}}
\BIC{$x = x \to x = y \to y = x$}
\BIC{$x = y \to y = x$}
\UIC{$\forall y\,(x = y \to y = x)$}
\UIC{$\forall x\,\forall y\,(x = y \to y = x)$}
\end{ptree}
\noindent$\Pi_{13}$:
\begin{ptree}
\AXC{$\stackrel{\Pi_{12}}{\vdots}$}\noLine
\UIC{$\forall x\,\forall y\,(x = y \to y = x)$}
\AXC{$\forall x\,\forall y\,(x = y \to y = x) \to \varphi_{16}$\lbl{A2}}
\BIC{$\varphi_{16}$}
\AXC{$\varphi_{16} \to \lvert f(c) - f(c) \rvert = 0 \to 0 = \lvert f(c) - f(c) \rvert$\lbl{A2}}
\BIC{$\lvert f(c) - f(c) \rvert = 0 \to 0 = \lvert f(c) - f(c) \rvert$}
\end{ptree}
\noindent$\Pi_{14}$:
\begin{ptree}
\AXC{$y = z \to x = y \to x = z$\lbl{A6}}
\AXC{$(y = z \to x = y \to x = z) \to x = y \to y = z \to x = z$\lbl{A1}}
\BIC{$x = y \to y = z \to x = z$}
\UIC{$\forall z\,(x = y \to y = z \to x = z)$}
\UIC{$\forall y\,\forall z\,(x = y \to y = z \to x = z)$}
\UIC{$\forall x\,\forall y\,\forall z\,(x = y \to y = z \to x = z)$}
\AXC{$\forall x\,\forall y\,\forall z\,(x = y \to y = z \to x = z) \to \varphi_{28}$\lbl{A2}}
\BIC{$\varphi_{28}$}
\end{ptree}
\noindent$\Pi_{15}$:
\begin{ptree}
\AXC{$x = x$\lbl{A5}}
\UIC{$\forall x\,(x = x)$}
\AXC{$\forall x\,(x = x) \to 0 = 0$\lbl{A2}}
\BIC{$0 = 0$}
\AXC{$0 = 0 \to \mathrm{PNbhd}(f,c) \to \varphi_{101}$\lbl{A1}}
\BIC{$\mathrm{PNbhd}(f,c) \to \varphi_{101}$}
\end{ptree}
\noindent$\Pi_{16}$:
\begin{ptree}
\AXC{$\forall x\,(0 \leq x \to \lvert x \rvert = x)$\lbl{[A1]}}
\AXC{$\forall x\,(0 \leq x \to \lvert x \rvert = x) \to 0 \leq 0 \to \lvert 0 \rvert = 0$\lbl{A2}}
\BIC{$0 \leq 0 \to \lvert 0 \rvert = 0$}
\AXC{$\varphi_{161}$\lbl{A1}}
\BIC{$(\mathrm{PNbhd}(f,c) \to \varphi_{101}) \to \mathrm{PNbhd}(f,c) \to \varphi_{109}$}
\end{ptree}
\noindent$\Pi_{17}$:
\begin{ptree}
\AXC{$z = w \to \varphi_{25}$\lbl{A6}}
\UIC{$\forall w\,(z = w \to \varphi_{25})$}
\UIC{$\forall z\,\forall w\,(z = w \to \varphi_{25})$}
\AXC{$\forall z\,\forall w\,(z = w \to \varphi_{25}) \to \varphi_{61}$\lbl{A2}}
\BIC{$\varphi_{61}$}
\AXC{$\varphi_{61} \to f(c) - f(c) = 0 \to \varphi_{50}$\lbl{A2}}
\BIC{$f(c) - f(c) = 0 \to \varphi_{50}$}
\end{ptree}
\noindent$\Pi_{18}$:
\begin{ptree}
\AXC{$\forall f\,\forall x\,(\mathrm{RFn}(f) \to x \in \mathrm{dom}(f) \to f(x) \in \mathbb{R})$\lbl{[S6]}}
\AXC{$\varphi_{58}$\lbl{A2}}
\BIC{$\forall x\,(\mathrm{RFn}(f) \to x \in \mathrm{dom}(f) \to f(x) \in \mathbb{R})$}
\AXC{$\varphi_{53}$\lbl{A2}}
\BIC{$\mathrm{RFn}(f) \to c \in \mathrm{dom}(f) \to f(c) \in \mathbb{R}$}
\end{ptree}
\noindent$\Pi_{19}$:
\begin{ptree}
\AXC{$\stackrel{\Pi_{4}}{\vdots}$}\noLine
\UIC{$\varphi_{12}$}
\AXC{$\stackrel{\Pi_{18}}{\vdots}$}\noLine
\UIC{$\mathrm{RFn}(f) \to c \in \mathrm{dom}(f) \to f(c) \in \mathbb{R}$}
\AXC{$(\mathrm{RFn}(f) \to c \in \mathrm{dom}(f) \to f(c) \in \mathbb{R}) \to \varphi_{12} \to \varphi_{146}$\lbl{A1}}
\BIC{$\varphi_{12} \to \varphi_{146}$}
\BIC{$\varphi_{146}$}
\end{ptree}
\noindent$\Pi_{20}$:
\begin{ptree}
\AXC{$\forall x \in \mathbb{R}\,(x - x = 0)$\lbl{[F13]}}
\AXC{$\varphi_{29}$\lbl{A2}}
\BIC{$f(c) \in \mathbb{R} \to f(c) - f(c) = 0$}
\AXC{$(f(c) \in \mathbb{R} \to f(c) - f(c) = 0) \to \varphi_{162}$\lbl{A1}}
\BIC{$\varphi_{162}$}
\end{ptree}
\noindent$\Pi_{21}$:
\begin{ptree}
\AXC{$\mathrm{PNbhd}(f,c) \to c \in \mathrm{dom}(f) \to c \in \mathrm{dom}(f)$\lbl{A1}}
\AXC{$\stackrel{\Pi_{19}}{\vdots}$}\noLine
\UIC{$\varphi_{146}$}
\BIC{$\mathrm{PNbhd}(f,c) \to \varphi_{120}$}
\AXC{$\stackrel{\Pi_{20}}{\vdots}$}\noLine
\UIC{$\varphi_{162}$}
\BIC{$\mathrm{PNbhd}(f,c) \to c \in \mathrm{dom}(f) \to \varphi_{108}$}
\AXC{$\varphi_{178}$\lbl{A1}}
\BIC{$\lvert f(c) - f(c) \rvert = \lvert f(c) - f(c) \rvert \to \varphi_{152}$}
\end{ptree}
\noindent$\Pi_{22}$:
\begin{ptree}
\AXC{$\stackrel{\Pi_{17}}{\vdots}$}\noLine
\UIC{$f(c) - f(c) = 0 \to \varphi_{50}$}
\AXC{$x = x$\lbl{A5}}
\UIC{$\forall x\,(x = x)$}
\AXC{$\varphi_{17}$\lbl{A2}}
\BIC{$\lvert f(c) - f(c) \rvert = \lvert f(c) - f(c) \rvert$}
\AXC{$\stackrel{\Pi_{21}}{\vdots}$}\noLine
\UIC{$\lvert f(c) - f(c) \rvert = \lvert f(c) - f(c) \rvert \to \varphi_{152}$}
\BIC{$\varphi_{152}$}
\BIC{$\mathrm{PNbhd}(f,c) \to c \in \mathrm{dom}(f) \to \varphi_{121}$}
\end{ptree}
\noindent$\Pi_{23}$:
\begin{ptree}
\AXC{$\stackrel{\Pi_{22}}{\vdots}$}\noLine
\UIC{$\mathrm{PNbhd}(f,c) \to c \in \mathrm{dom}(f) \to \varphi_{121}$}
\AXC{$(\mathrm{PNbhd}(f,c) \to c \in \mathrm{dom}(f) \to \varphi_{121}) \to \varphi_{170}$\lbl{A1}}
\BIC{$\varphi_{170}$}
\end{ptree}
\noindent$\Pi_{24}$:
\begin{ptree}
\AXC{$\stackrel{\Pi_{15}}{\vdots}$}\noLine
\UIC{$\mathrm{PNbhd}(f,c) \to \varphi_{101}$}
\AXC{$\stackrel{\Pi_{16}}{\vdots}$}\noLine
\UIC{$(\mathrm{PNbhd}(f,c) \to \varphi_{101}) \to \mathrm{PNbhd}(f,c) \to \varphi_{109}$}
\BIC{$\mathrm{PNbhd}(f,c) \to \varphi_{109}$}
\AXC{$\stackrel{\Pi_{23}}{\vdots}$}\noLine
\UIC{$\varphi_{170}$}
\BIC{$\varphi_{37} \to \mathrm{PNbhd}(f,c) \to c \in \mathrm{dom}(f) \to \varphi_{114}$}
\end{ptree}
\noindent$\Pi_{25}$:
\begin{ptree}
\AXC{$\stackrel{\Pi_{14}}{\vdots}$}\noLine
\UIC{$\varphi_{28}$}
\AXC{$\varphi_{28} \to \varphi_{45}$\lbl{A2}}
\BIC{$\varphi_{45}$}
\AXC{$\varphi_{45} \to \varphi_{37}$\lbl{A2}}
\BIC{$\varphi_{37}$}
\AXC{$\stackrel{\Pi_{24}}{\vdots}$}\noLine
\UIC{$\varphi_{37} \to \mathrm{PNbhd}(f,c) \to c \in \mathrm{dom}(f) \to \varphi_{114}$}
\BIC{$\mathrm{PNbhd}(f,c) \to c \in \mathrm{dom}(f) \to \varphi_{114}$}
\end{ptree}
\noindent$\Pi_{26}$:
\begin{ptree}
\AXC{$\stackrel{\Pi_{25}}{\vdots}$}\noLine
\UIC{$\mathrm{PNbhd}(f,c) \to c \in \mathrm{dom}(f) \to \varphi_{114}$}
\AXC{$(\mathrm{PNbhd}(f,c) \to c \in \mathrm{dom}(f) \to \varphi_{114}) \to \varphi_{148}$\lbl{A1}}
\BIC{$\varphi_{148}$}
\end{ptree}
\noindent$\Pi_{27}$:
\begin{ptree}
\AXC{$\stackrel{\Pi_{13}}{\vdots}$}\noLine
\UIC{$\lvert f(c) - f(c) \rvert = 0 \to 0 = \lvert f(c) - f(c) \rvert$}
\AXC{$\stackrel{\Pi_{26}}{\vdots}$}\noLine
\UIC{$\varphi_{148}$}
\BIC{$\mathrm{PNbhd}(f,c) \to c \in \mathrm{dom}(f) \to \varphi_{115}$}
\AXC{$(\mathrm{PNbhd}(f,c) \to c \in \mathrm{dom}(f) \to \varphi_{115}) \to \varphi_{156}$\lbl{A1}}
\BIC{$\varphi_{156}$}
\end{ptree}
\noindent$\Pi_{28}$:
\begin{ptree}
\AXC{$z = w \to z < \varepsilon \to w < \varepsilon$\lbl{A6}}
\UIC{$\forall w\,(z = w \to z < \varepsilon \to w < \varepsilon)$}
\UIC{$\forall z\,\forall w\,(z = w \to z < \varepsilon \to w < \varepsilon)$}
\AXC{$\varphi_{19}$\lbl{A2}}
\BIC{$\forall w\,(0 = w \to 0 < \varepsilon \to w < \varepsilon)$}
\AXC{$\forall w\,(0 = w \to 0 < \varepsilon \to w < \varepsilon) \to \varphi_{36}$\lbl{A2}}
\BIC{$\varphi_{36}$}
\AXC{$\mathrm{PNbhd}(f,c) \to \varphi_{32}$\lbl{A1}}
\AXC{$\stackrel{\Pi_{27}}{\vdots}$}\noLine
\UIC{$\varphi_{156}$}
\BIC{$\varphi_{36} \to \mathrm{PNbhd}(f,c) \to c \in \mathrm{dom}(f) \to \varphi_{122}$}
\BIC{$\mathrm{PNbhd}(f,c) \to c \in \mathrm{dom}(f) \to \varphi_{122}$}
\end{ptree}
\noindent$\Pi_{29}$:
\begin{ptree}
\AXC{$\stackrel{\Pi_{12}}{\vdots}$}\noLine
\UIC{$\forall x\,\forall y\,(x = y \to y = x)$}
\AXC{$\forall x\,\forall y\,(x = y \to y = x) \to \forall y\,(x = y \to y = x)$\lbl{A2}}
\BIC{$\forall y\,(x = y \to y = x)$}
\AXC{$\forall y\,(x = y \to y = x) \to x = c \to c = x$\lbl{A2}}
\BIC{$x = c \to c = x$}
\end{ptree}
\noindent$\Pi_{30}$:
\begin{ptree}
\AXC{$\stackrel{\Pi_{29}}{\vdots}$}\noLine
\UIC{$x = c \to c = x$}
\AXC{$\mathrm{PNbhd}(f,c) \to \varphi_{102}$\lbl{A1}}
\AXC{$\varphi_{159}$\lbl{A1}}
\BIC{$(x = c \to c = x) \to \mathrm{PNbhd}(f,c) \to \varphi_{103}$}
\BIC{$\mathrm{PNbhd}(f,c) \to \varphi_{103}$}
\AXC{$(\mathrm{PNbhd}(f,c) \to \varphi_{103}) \to \varphi_{177}$\lbl{A1}}
\BIC{$\varphi_{177}$}
\end{ptree}
\noindent$\Pi_{31}$:
\begin{ptree}
\AXC{$\varphi_{34}$\lbl{A6}}
\UIC{$\forall w\,(\varphi_{34})$}
\UIC{$\forall z\,\forall w\,(\varphi_{34})$}
\AXC{$\forall z\,\forall w\,(\varphi_{34}) \to \varphi_{44}$\lbl{A2}}
\BIC{$\varphi_{44}$}
\AXC{$\varphi_{44} \to \varphi_{35}$\lbl{A2}}
\BIC{$\varphi_{35}$}
\AXC{$\stackrel{\Pi_{28}}{\vdots}$}\noLine
\UIC{$\mathrm{PNbhd}(f,c) \to c \in \mathrm{dom}(f) \to \varphi_{122}$}
\AXC{$\stackrel{\Pi_{30}}{\vdots}$}\noLine
\UIC{$\varphi_{177}$}
\BIC{$\varphi_{35} \to \mathrm{PNbhd}(f,c) \to c \in \mathrm{dom}(f) \to \varphi_{123}$}
\BIC{$\mathrm{PNbhd}(f,c) \to c \in \mathrm{dom}(f) \to \varphi_{123}$}
\end{ptree}
\noindent$\Pi_{32}$:
\begin{ptree}
\AXC{$\stackrel{\Pi_{31}}{\vdots}$}\noLine
\UIC{$\mathrm{PNbhd}(f,c) \to c \in \mathrm{dom}(f) \to \varphi_{123}$}
\AXC{$(\mathrm{PNbhd}(f,c) \to c \in \mathrm{dom}(f) \to \varphi_{123}) \to \varphi_{164}$\lbl{A1}}
\BIC{$\varphi_{164}$}
\end{ptree}
\noindent$\Pi_{33}$:
\begin{ptree}
\AXC{$\varphi_{165}$\lbl{A3}}
\AXC{$\mathrm{PNbhd}(f,c) \to \varphi_{113}$\lbl{A1}}
\AXC{$\stackrel{\Pi_{11}}{\vdots}$}\noLine
\UIC{$\varphi_{163}$}
\BIC{$\mathrm{PNbhd}(f,c) \to c \in \mathrm{dom}(f) \to \varphi_{124}$}
\AXC{$\stackrel{\Pi_{32}}{\vdots}$}\noLine
\UIC{$\varphi_{164}$}
\BIC{$\mathrm{PNbhd}(f,c) \to \varphi_{128}$}
\AXC{$(\mathrm{PNbhd}(f,c) \to \varphi_{128}) \to \varphi_{135}$\lbl{A1}}
\BIC{$\varphi_{135}$}
\UIC{$\forall x\,(\varphi_{135})$}
\AXC{$\varphi_{186}$\lbl{A1}}
\BIC{$\varphi_{165} \to \varphi_{153} \to \mathrm{PNbhd}(f,c) \to c \in \mathrm{dom}(f) \to \varphi_{116}$}
\BIC{$\varphi_{153} \to \mathrm{PNbhd}(f,c) \to c \in \mathrm{dom}(f) \to \varphi_{116}$}
\end{ptree}
\noindent$\Pi_{34}$:
\begin{ptree}
\AXC{$\mathrm{PNbhd}(f,c) \to c \in \mathrm{dom}(f) \to \varphi_{136}$\lbl{A1}}
\AXC{$(\mathrm{PNbhd}(f,c) \to c \in \mathrm{dom}(f) \to \varphi_{136}) \to \varphi_{167}$\lbl{A1}}
\BIC{$\varphi_{167}$}
\end{ptree}
\noindent$\Pi_{35}$:
\begin{ptree}
\AXC{$\varphi_{153}$\lbl{A4}}
\AXC{$\stackrel{\Pi_{33}}{\vdots}$}\noLine
\UIC{$\varphi_{153} \to \mathrm{PNbhd}(f,c) \to c \in \mathrm{dom}(f) \to \varphi_{116}$}
\BIC{$\mathrm{PNbhd}(f,c) \to c \in \mathrm{dom}(f) \to \varphi_{116}$}
\AXC{$\stackrel{\Pi_{34}}{\vdots}$}\noLine
\UIC{$\varphi_{167}$}
\BIC{$\mathrm{PNbhd}(f,c) \to c \in \mathrm{dom}(f) \to \varphi_{134}$}
\AXC{$(\mathrm{PNbhd}(f,c) \to c \in \mathrm{dom}(f) \to \varphi_{134}) \to \varphi_{95} \to \varphi_{144}$\lbl{A1}}
\BIC{$\varphi_{95} \to \varphi_{144}$}
\end{ptree}
\noindent$\Pi_{36}$:
\begin{ptree}
\AXC{$\forall f\,\forall c\,(\varphi_{76})$\lbl{[D1]}}
\AXC{$\forall f\,\forall c\,(\varphi_{76}) \to \forall c\,(\varphi_{76})$\lbl{A2}}
\BIC{$\forall c\,(\varphi_{76})$}
\AXC{$\forall c\,(\varphi_{76}) \to \varphi_{76}$\lbl{A2}}
\BIC{$\varphi_{76}$}
\AXC{$\varphi_{76} \to \forall \varepsilon\,(\varphi_{80})$\lbl{A2}}
\BIC{$\forall \varepsilon\,(\varphi_{80})$}
\AXC{$\forall \varepsilon\,(\varphi_{80}) \to \varphi_{80}$\lbl{A2}}
\BIC{$\varphi_{80}$}
\end{ptree}
\noindent$\Pi_{37}$:
\begin{ptree}
\AXC{$\mathrm{PNbhd}(f,c) \to \varphi_{32}$\lbl{A1}}
\AXC{$\mathrm{PNbhd}(f,c) \to \varphi_{22}$\lbl{A1}}
\AXC{$\mathrm{PNbhd}(f,c) \to \varphi_{10}$\lbl{A1}}
\AXC{$\stackrel{\Pi_{36}}{\vdots}$}\noLine
\UIC{$\varphi_{80}$}
\AXC{$\varphi_{80} \to \varphi_{157}$\lbl{A1}}
\BIC{$\varphi_{157}$}
\BIC{$(\mathrm{PNbhd}(f,c) \to \varphi_{22}) \to \varphi_{138}$}
\BIC{$\varphi_{138}$}
\BIC{$\mathrm{PNbhd}(f,c) \to c \in \mathrm{dom}(f) \to \varphi_{80}$}
\end{ptree}
\noindent$\Pi_{38}$:
\begin{ptree}
\AXC{$\forall \delta\,(\varphi_{145}) \to \varphi_{147}$\lbl{A3}}
\AXC{$\varphi_{95}$\lbl{A2}}
\AXC{$\stackrel{\Pi_{35}}{\vdots}$}\noLine
\UIC{$\varphi_{95} \to \varphi_{144}$}
\BIC{$\varphi_{144}$}
\AXC{$\varphi_{144} \to \varphi_{145}$\lbl{A1}}
\BIC{$\varphi_{145}$}
\UIC{$\forall \delta\,(\varphi_{145})$}
\AXC{$\stackrel{\Pi_{37}}{\vdots}$}\noLine
\UIC{$\mathrm{PNbhd}(f,c) \to c \in \mathrm{dom}(f) \to \varphi_{80}$}
\AXC{$\varphi_{188}$\lbl{A1}}
\BIC{$\forall \delta\,(\varphi_{145}) \to \varphi_{184}$}
\BIC{$\varphi_{184}$}
\BIC{$\varphi_{150} \to \mathrm{PNbhd}(f,c) \to c \in \mathrm{dom}(f) \to \varphi_{66}$}
\end{ptree}
\noindent$\Pi_{39}$:
\begin{ptree}
\AXC{$\varphi_{7} \to \forall \varepsilon\,(\varphi_{7})$\lbl{A4}}
\AXC{$\varphi_{151}$\lbl{A3}}
\AXC{$\varphi_{150}$\lbl{A4}}
\AXC{$\stackrel{\Pi_{38}}{\vdots}$}\noLine
\UIC{$\varphi_{150} \to \mathrm{PNbhd}(f,c) \to c \in \mathrm{dom}(f) \to \varphi_{66}$}
\BIC{$\mathrm{PNbhd}(f,c) \to c \in \mathrm{dom}(f) \to \varphi_{66}$}
\AXC{$\varphi_{149}$\lbl{A1}}
\BIC{$\varphi_{7} \to \varepsilon \in \mathbb{R} \to 0 < \varepsilon \to \exists \delta\,(\delta \in \mathbb{R} \land 0 < \delta \land \varphi_{31})$}
\UIC{$\forall \varepsilon\,(\varphi_{7} \to \varepsilon \in \mathbb{R} \to 0 < \varepsilon \to \exists \delta\,(\delta \in \mathbb{R} \land 0 < \delta \land \varphi_{31}))$}
\AXC{$\varphi_{182}$\lbl{A1}}
\BIC{$\varphi_{151} \to \varphi_{140}$}
\BIC{$\varphi_{140}$}
\BIC{$\mathrm{PNbhd}(f,c) \to c \in \mathrm{dom}(f) \to \varphi_{67}$}
\end{ptree}
\noindent$\Pi_{40}$:
\begin{ptree}
\AXC{$\forall f\,(\varphi_{88})$\lbl{DefContIn}}
\AXC{$\forall f\,(\varphi_{88}) \to \varphi_{88}$\lbl{A2}}
\BIC{$\varphi_{88}$}
\AXC{$\varphi_{88} \to \forall c\,(\varphi_{94})$\lbl{A2}}
\BIC{$\forall c\,(\varphi_{94})$}
\AXC{$\forall c\,(\varphi_{94}) \to \varphi_{94}$\lbl{A2}}
\BIC{$\varphi_{94}$}
\end{ptree}
\noindent$\Pi_{41}$:
\begin{ptree}
\AXC{$\stackrel{\Pi_{4}}{\vdots}$}\noLine
\UIC{$\varphi_{12}$}
\AXC{$\stackrel{\Pi_{40}}{\vdots}$}\noLine
\UIC{$\varphi_{94}$}
\AXC{$\forall f\,(\varphi_{1})$\lbl{DefCont}}
\AXC{$\forall f\,(\varphi_{1}) \to \varphi_{1}$\lbl{A2}}
\BIC{$\varphi_{1}$}
\AXC{$\varphi_{24}$\lbl{A2}}
\BIC{$\mathrm{Cont}(f,c) \leftrightarrow \mathrm{ContIn}(f,\mathrm{dom}(f),c)$}
\AXC{$\varphi_{174}$\lbl{A1}}
\BIC{$\varphi_{94} \to \varphi_{12} \to \mathrm{dom}(f) \subseteq \mathrm{dom}(f) \to \varphi_{141}$}
\BIC{$\varphi_{12} \to \mathrm{dom}(f) \subseteq \mathrm{dom}(f) \to \varphi_{141}$}
\BIC{$\mathrm{dom}(f) \subseteq \mathrm{dom}(f) \to \varphi_{141}$}
\end{ptree}
\noindent$\Pi_{42}$:
\begin{ptree}
\AXC{$\forall A\,(A \subseteq A)$\lbl{[E43]}}
\AXC{$\forall A\,(A \subseteq A) \to \mathrm{dom}(f) \subseteq \mathrm{dom}(f)$\lbl{A2}}
\BIC{$\mathrm{dom}(f) \subseteq \mathrm{dom}(f)$}
\AXC{$\stackrel{\Pi_{41}}{\vdots}$}\noLine
\UIC{$\mathrm{dom}(f) \subseteq \mathrm{dom}(f) \to \varphi_{141}$}
\BIC{$\varphi_{141}$}
\end{ptree}
\noindent$\Pi_{43}$:
\begin{ptree}
\AXC{$\stackrel{\Pi_{39}}{\vdots}$}\noLine
\UIC{$\mathrm{PNbhd}(f,c) \to c \in \mathrm{dom}(f) \to \varphi_{67}$}
\AXC{$\mathrm{PNbhd}(f,c) \to c \in \mathrm{dom}(f) \to c \in \mathrm{dom}(f)$\lbl{A1}}
\AXC{$\stackrel{\Pi_{42}}{\vdots}$}\noLine
\UIC{$\varphi_{141}$}
\BIC{$\varphi_{125}$}
\BIC{$\mathrm{PNbhd}(f,c) \to \varphi_{3}$}
\end{ptree}
\noindent$\Pi_{44}$:
\begin{ptree}
\AXC{$\varphi_{64}$\lbl{A2}}
\AXC{$\mathrm{PNbhd}(f,c) \to c \in \mathrm{dom}(f) \to \varphi_{111}$\lbl{A1}}
\AXC{$(\mathrm{PNbhd}(f,c) \to c \in \mathrm{dom}(f) \to \varphi_{111}) \to \varphi_{143}$\lbl{A1}}
\BIC{$\varphi_{143}$}
\AXC{$\varphi_{143} \to \varphi_{64} \to \varphi_{142}$\lbl{A1}}
\BIC{$\varphi_{64} \to \varphi_{142}$}
\BIC{$\varphi_{142}$}
\end{ptree}
\noindent$\Pi_{45}$:
\begin{ptree}
\AXC{$\mathrm{PNbhd}(f,c) \to c \in \mathrm{dom}(f) \to \varphi_{96}$\lbl{A1}}
\AXC{$\varphi_{93}$\lbl{A1}}
\AXC{$\stackrel{\Pi_{44}}{\vdots}$}\noLine
\UIC{$\varphi_{142}$}
\AXC{$\varphi_{142} \to \varphi_{93} \to \varphi_{160}$\lbl{A1}}
\BIC{$\varphi_{93} \to \varphi_{160}$}
\BIC{$\varphi_{160}$}
\BIC{$\mathrm{PNbhd}(f,c) \to c \in \mathrm{dom}(f) \to \varphi_{100}$}
\end{ptree}
\noindent$\Pi_{46}$:
\begin{ptree}
\AXC{$\stackrel{\Pi_{45}}{\vdots}$}\noLine
\UIC{$\mathrm{PNbhd}(f,c) \to c \in \mathrm{dom}(f) \to \varphi_{100}$}
\AXC{$(\mathrm{PNbhd}(f,c) \to c \in \mathrm{dom}(f) \to \varphi_{100}) \to \varphi_{74} \to \varphi_{49}$\lbl{A1}}
\BIC{$\varphi_{74} \to \varphi_{49}$}
\UIC{$\forall x\,(\varphi_{74} \to \varphi_{49})$}
\end{ptree}
\noindent$\Pi_{47}$:
\begin{ptree}
\AXC{$\forall x\,(\varphi_{74} \to \varphi_{49}) \to \forall x\,(\varphi_{74}) \to \varphi_{54}$\lbl{A3}}
\AXC{$\stackrel{\Pi_{46}}{\vdots}$}\noLine
\UIC{$\forall x\,(\varphi_{74} \to \varphi_{49})$}
\AXC{$\forall x\,(\varphi_{74} \to \varphi_{49}) \to (\forall x\,(\varphi_{74} \to \varphi_{49}) \to \forall x\,(\varphi_{74}) \to \varphi_{54}) \to \varphi_{172}$\lbl{A1}}
\BIC{$(\forall x\,(\varphi_{74} \to \varphi_{49}) \to \forall x\,(\varphi_{74}) \to \varphi_{54}) \to \varphi_{172}$}
\BIC{$\varphi_{172}$}
\end{ptree}
\noindent$\Pi_{48}$:
\begin{ptree}
\AXC{$\mathrm{PNbhd}(f,c) \to c \in \mathrm{dom}(f) \to \varphi_{111}$\lbl{A1}}
\AXC{$(\mathrm{PNbhd}(f,c) \to c \in \mathrm{dom}(f) \to \varphi_{111}) \to \varphi_{166}$\lbl{A1}}
\BIC{$\varphi_{166}$}
\end{ptree}
\noindent$\Pi_{49}$:
\begin{ptree}
\AXC{$\varphi_{117}$\lbl{A2}}
\AXC{$\varphi_{74} \to \forall x\,(\varphi_{74})$\lbl{A4}}
\AXC{$\stackrel{\Pi_{47}}{\vdots}$}\noLine
\UIC{$\varphi_{172}$}
\BIC{$\mathrm{PNbhd}(f,c) \to c \in \mathrm{dom}(f) \to \varphi_{107}$}
\AXC{$\stackrel{\Pi_{48}}{\vdots}$}\noLine
\UIC{$\varphi_{166}$}
\BIC{$\mathrm{PNbhd}(f,c) \to c \in \mathrm{dom}(f) \to \varphi_{129}$}
\AXC{$(\mathrm{PNbhd}(f,c) \to c \in \mathrm{dom}(f) \to \varphi_{129}) \to \varphi_{117} \to \varphi_{139}$\lbl{A1}}
\BIC{$\varphi_{117} \to \varphi_{139}$}
\BIC{$\varphi_{139}$}
\end{ptree}
\noindent$\Pi_{50}$:
\begin{ptree}
\AXC{$\varphi_{133}$\lbl{A2}}
\AXC{$\stackrel{\Pi_{40}}{\vdots}$}\noLine
\UIC{$\varphi_{94}$}
\AXC{$\forall f\,(\varphi_{1})$\lbl{DefCont}}
\AXC{$\forall f\,(\varphi_{1}) \to \varphi_{1}$\lbl{A2}}
\BIC{$\varphi_{1}$}
\AXC{$\varphi_{24}$\lbl{A2}}
\BIC{$\mathrm{Cont}(f,c) \leftrightarrow \mathrm{ContIn}(f,\mathrm{dom}(f),c)$}
\AXC{$\varphi_{8}$\lbl{A1}}
\AXC{$\varphi_{8} \to \varphi_{155}$\lbl{A1}}
\BIC{$\varphi_{155}$}
\BIC{$\varphi_{94} \to \mathrm{PNbhd}(f,c) \to c \in \mathrm{dom}(f) \to \varphi_{65}$}
\BIC{$\mathrm{PNbhd}(f,c) \to c \in \mathrm{dom}(f) \to \varphi_{65}$}
\AXC{$\varphi_{176}$\lbl{A1}}
\BIC{$\varphi_{133} \to \mathrm{PNbhd}(f,c) \to c \in \mathrm{dom}(f) \to \varphi_{63}$}
\BIC{$\mathrm{PNbhd}(f,c) \to c \in \mathrm{dom}(f) \to \varphi_{63}$}
\end{ptree}
\noindent$\Pi_{51}$:
\begin{ptree}
\AXC{$\varphi_{171}$\lbl{A3}}
\AXC{$\stackrel{\Pi_{49}}{\vdots}$}\noLine
\UIC{$\varphi_{139}$}
\AXC{$\varphi_{139} \to \varphi_{92} \to \neg (\delta \in \mathbb{R} \land 0 < \delta \land \varphi_{31})$\lbl{A1}}
\BIC{$\varphi_{92} \to \neg (\delta \in \mathbb{R} \land 0 < \delta \land \varphi_{31})$}
\UIC{$\forall \delta\,(\varphi_{92} \to \neg (\delta \in \mathbb{R} \land 0 < \delta \land \varphi_{31}))$}
\AXC{$\stackrel{\Pi_{50}}{\vdots}$}\noLine
\UIC{$\mathrm{PNbhd}(f,c) \to c \in \mathrm{dom}(f) \to \varphi_{63}$}
\AXC{$\varphi_{187}$\lbl{A1}}
\BIC{$\forall \delta\,(\varphi_{92} \to \neg (\delta \in \mathbb{R} \land 0 < \delta \land \varphi_{31})) \to \varphi_{171} \to \varphi_{169}$}
\BIC{$\varphi_{171} \to \varphi_{169}$}
\BIC{$\varphi_{169}$}
\end{ptree}
\noindent$\Pi_{52}$:
\begin{ptree}
\AXC{$\varphi_{154}$\lbl{A3}}
\AXC{$\varphi_{92} \to \forall \delta\,(\varphi_{92})$\lbl{A4}}
\AXC{$\stackrel{\Pi_{51}}{\vdots}$}\noLine
\UIC{$\varphi_{169}$}
\BIC{$\mathrm{PNbhd}(f,c) \to c \in \mathrm{dom}(f) \to \varphi_{70}$}
\AXC{$(\mathrm{PNbhd}(f,c) \to c \in \mathrm{dom}(f) \to \varphi_{70}) \to \varphi_{89}$\lbl{A1}}
\BIC{$\varphi_{89}$}
\UIC{$\forall \varepsilon\,(\varphi_{89})$}
\AXC{$\varphi_{183}$\lbl{A1}}
\BIC{$\varphi_{154} \to \varphi_{47} \to \mathrm{PNbhd}(f,c) \to c \in \mathrm{dom}(f) \to \varphi_{75}$}
\BIC{$\varphi_{47} \to \mathrm{PNbhd}(f,c) \to c \in \mathrm{dom}(f) \to \varphi_{75}$}
\end{ptree}
\noindent$\Pi_{53}$:
\begin{ptree}
\AXC{$\stackrel{\Pi_{40}}{\vdots}$}\noLine
\UIC{$\varphi_{94}$}
\AXC{$\forall f\,(\varphi_{1})$\lbl{DefCont}}
\AXC{$\forall f\,(\varphi_{1}) \to \varphi_{1}$\lbl{A2}}
\BIC{$\varphi_{1}$}
\AXC{$\varphi_{24}$\lbl{A2}}
\BIC{$\mathrm{Cont}(f,c) \leftrightarrow \mathrm{ContIn}(f,\mathrm{dom}(f),c)$}
\AXC{$\varphi_{8}$\lbl{A1}}
\AXC{$\varphi_{8} \to \varphi_{137}$\lbl{A1}}
\BIC{$\varphi_{137}$}
\BIC{$\varphi_{94} \to \varphi_{6}$}
\BIC{$\varphi_{6}$}
\end{ptree}
\noindent$\Pi_{54}$:
\begin{ptree}
\AXC{$\forall f\,\forall x\,(\mathrm{RFn}(f) \to x \in \mathrm{dom}(f) \to f(x) \in \mathbb{R})$\lbl{[S6]}}
\AXC{$\varphi_{58}$\lbl{A2}}
\BIC{$\forall x\,(\mathrm{RFn}(f) \to x \in \mathrm{dom}(f) \to f(x) \in \mathbb{R})$}
\AXC{$\varphi_{53}$\lbl{A2}}
\BIC{$\mathrm{RFn}(f) \to c \in \mathrm{dom}(f) \to f(c) \in \mathbb{R}$}
\end{ptree}
\noindent$\Pi_{55}$:
\begin{ptree}
\AXC{$\stackrel{\Pi_{53}}{\vdots}$}\noLine
\UIC{$\varphi_{6}$}
\AXC{$\stackrel{\Pi_{54}}{\vdots}$}\noLine
\UIC{$\mathrm{RFn}(f) \to c \in \mathrm{dom}(f) \to f(c) \in \mathbb{R}$}
\AXC{$(\mathrm{RFn}(f) \to c \in \mathrm{dom}(f) \to f(c) \in \mathbb{R}) \to \varphi_{6} \to \varphi_{60}$\lbl{A1}}
\BIC{$\varphi_{6} \to \varphi_{60}$}
\BIC{$\varphi_{60}$}
\end{ptree}
\noindent$\Pi_{56}$:
\begin{ptree}
\AXC{$\forall f\,\forall x\,(\mathrm{RFn}(f) \to x \in \mathrm{dom}(f) \to x \in \mathbb{R})$\lbl{[S7]}}
\AXC{$\varphi_{46}$\lbl{A2}}
\BIC{$\forall x\,(\mathrm{RFn}(f) \to x \in \mathrm{dom}(f) \to x \in \mathbb{R})$}
\AXC{$\varphi_{42}$\lbl{A2}}
\BIC{$\mathrm{RFn}(f) \to c \in \mathrm{dom}(f) \to c \in \mathbb{R}$}
\end{ptree}
\noindent$\Pi_{57}$:
\begin{ptree}
\AXC{$\stackrel{\Pi_{53}}{\vdots}$}\noLine
\UIC{$\varphi_{6}$}
\AXC{$\stackrel{\Pi_{56}}{\vdots}$}\noLine
\UIC{$\mathrm{RFn}(f) \to c \in \mathrm{dom}(f) \to c \in \mathbb{R}$}
\AXC{$(\mathrm{RFn}(f) \to c \in \mathrm{dom}(f) \to c \in \mathbb{R}) \to \varphi_{6} \to \varphi_{56}$\lbl{A1}}
\BIC{$\varphi_{6} \to \varphi_{56}$}
\BIC{$\varphi_{56}$}
\end{ptree}
\noindent$\Pi_{58}$:
\begin{ptree}
\AXC{$\stackrel{\Pi_{53}}{\vdots}$}\noLine
\UIC{$\varphi_{6}$}
\AXC{$\forall f\,\forall c\,(\varphi_{97})$\lbl{DefFLim}}
\AXC{$\forall f\,\forall c\,(\varphi_{97}) \to \forall c\,(\varphi_{97})$\lbl{A2}}
\BIC{$\forall c\,(\varphi_{97})$}
\AXC{$\forall c\,(\varphi_{97}) \to \varphi_{97}$\lbl{A2}}
\BIC{$\varphi_{97}$}
\AXC{$\varphi_{97} \to \varphi_{105}$\lbl{A2}}
\BIC{$\varphi_{105}$}
\AXC{$\varphi_{105} \to \varphi_{158}$\lbl{A1}}
\BIC{$\varphi_{158}$}
\BIC{$\varphi_{9} \to \varphi_{14} \to (\mathrm{PNbhd}(f,c) \to \mathrm{PNbhd}(f,c)) \to \varphi_{130}$}
\end{ptree}
\noindent$\Pi_{59}$:
\begin{ptree}
\AXC{$\mathrm{PNbhd}(f,c) \to c \in \mathrm{dom}(f) \to c \in \mathrm{dom}(f)$\lbl{A1}}
\AXC{$\stackrel{\Pi_{57}}{\vdots}$}\noLine
\UIC{$\varphi_{56}$}
\BIC{$\varphi_{9}$}
\AXC{$\stackrel{\Pi_{58}}{\vdots}$}\noLine
\UIC{$\varphi_{9} \to \varphi_{14} \to (\mathrm{PNbhd}(f,c) \to \mathrm{PNbhd}(f,c)) \to \varphi_{130}$}
\BIC{$\varphi_{14} \to (\mathrm{PNbhd}(f,c) \to \mathrm{PNbhd}(f,c)) \to \varphi_{130}$}
\end{ptree}
\noindent$\Pi_{60}$:
\begin{ptree}
\AXC{$\mathrm{PNbhd}(f,c) \to c \in \mathrm{dom}(f) \to c \in \mathrm{dom}(f)$\lbl{A1}}
\AXC{$\stackrel{\Pi_{55}}{\vdots}$}\noLine
\UIC{$\varphi_{60}$}
\BIC{$\varphi_{14}$}
\AXC{$\stackrel{\Pi_{59}}{\vdots}$}\noLine
\UIC{$\varphi_{14} \to (\mathrm{PNbhd}(f,c) \to \mathrm{PNbhd}(f,c)) \to \varphi_{130}$}
\BIC{$(\mathrm{PNbhd}(f,c) \to \mathrm{PNbhd}(f,c)) \to \varphi_{130}$}
\end{ptree}
\noindent$\Pi_{61}$:
\begin{ptree}
\AXC{$\mathrm{PNbhd}(f,c) \to \mathrm{PNbhd}(f,c)$\lbl{A1}}
\AXC{$\stackrel{\Pi_{60}}{\vdots}$}\noLine
\UIC{$(\mathrm{PNbhd}(f,c) \to \mathrm{PNbhd}(f,c)) \to \varphi_{130}$}
\BIC{$\varphi_{130}$}
\end{ptree}
\noindent 主木:
\begin{ptree}
\AXC{$\stackrel{\Pi_{43}}{\vdots}$}\noLine
\UIC{$\mathrm{PNbhd}(f,c) \to \varphi_{3}$}
\AXC{$\varphi_{47}$\lbl{A4}}
\AXC{$\stackrel{\Pi_{52}}{\vdots}$}\noLine
\UIC{$\varphi_{47} \to \mathrm{PNbhd}(f,c) \to c \in \mathrm{dom}(f) \to \varphi_{75}$}
\BIC{$\mathrm{PNbhd}(f,c) \to c \in \mathrm{dom}(f) \to \varphi_{75}$}
\AXC{$\stackrel{\Pi_{61}}{\vdots}$}\noLine
\UIC{$\varphi_{130}$}
\BIC{$\mathrm{PNbhd}(f,c) \to \varphi_{4}$}
\AXC{$\varphi_{87}$\lbl{A1}}
\BIC{$(\mathrm{PNbhd}(f,c) \to \varphi_{3}) \to \mathrm{PNbhd}(f,c) \to \varphi_{5}$}
\BIC{$\mathrm{PNbhd}(f,c) \to \varphi_{5}$}
\UIC{$\forall c\,(\mathrm{PNbhd}(f,c) \to \varphi_{5})$}
\UIC{$\forall f\,\forall c\,(\mathrm{PNbhd}(f,c) \to \varphi_{5})$}
\end{ptree}
\begin{abbrevlist}
\item $\varphi_{1} :\equiv \forall c\,(\mathrm{Cont}(f,c) \leftrightarrow \mathrm{ContIn}(f,\mathrm{dom}(f),c))$
\item $\varphi_{2} :\equiv \forall \varepsilon\,(\mathrm{PNbhd}(f,c) \land c \in \mathrm{dom}(f) \land \mathrm{Cont}(f,c))$
\item $\varphi_{3} :\equiv c \in \mathrm{dom}(f) \to \mathrm{FLim}(f,c,f(c)) \to \mathrm{Cont}(f,c)$
\item $\varphi_{4} :\equiv c \in \mathrm{dom}(f) \to \mathrm{Cont}(f,c) \to \mathrm{FLim}(f,c,f(c))$
\item $\varphi_{5} :\equiv c \in \mathrm{dom}(f) \to (\mathrm{Cont}(f,c) \leftrightarrow \mathrm{FLim}(f,c,f(c)))$
\item $\varphi_{6} :\equiv \mathrm{PNbhd}(f,c) \to c \in \mathrm{dom}(f) \to \mathrm{Cont}(f,c) \to \mathrm{RFn}(f)$
\item $\varphi_{7} :\equiv \mathrm{PNbhd}(f,c) \land c \in \mathrm{dom}(f) \land \mathrm{FLim}(f,c,f(c))$
\item $\varphi_{8} :\equiv \mathrm{PNbhd}(f,c) \to c \in \mathrm{dom}(f) \to \mathrm{Cont}(f,c) \to \mathrm{Cont}(f,c)$
\item $\varphi_{9} :\equiv \mathrm{PNbhd}(f,c) \to c \in \mathrm{dom}(f) \to \mathrm{Cont}(f,c) \to c \in \mathbb{R}$
\item $\varphi_{10} :\equiv c \in \mathrm{dom}(f) \to \mathrm{FLim}(f,c,f(c)) \to \mathrm{FLim}(f,c,f(c))$
\item $\varphi_{11} :\equiv \forall \alpha\,(\mathrm{FLim}(f,c,\alpha) \to \mathrm{RFn}(f)) \to \mathrm{FLim}(f,c,f(c)) \to \mathrm{RFn}(f)$
\item $\varphi_{12} :\equiv \mathrm{PNbhd}(f,c) \to c \in \mathrm{dom}(f) \to \mathrm{FLim}(f,c,f(c)) \to \mathrm{RFn}(f)$
\item $\varphi_{13} :\equiv \forall c\,\forall \alpha\,(\mathrm{FLim}(f,c,\alpha) \to \mathrm{RFn}(f)) \to \forall \alpha\,(\mathrm{FLim}(f,c,\alpha) \to \mathrm{RFn}(f))$
\item $\varphi_{14} :\equiv \mathrm{PNbhd}(f,c) \to c \in \mathrm{dom}(f) \to \mathrm{Cont}(f,c) \to f(c) \in \mathbb{R}$
\item $\varphi_{15} :\equiv \mathrm{PNbhd}(f,c) \to c \in \mathrm{dom}(f) \to \mathrm{FLim}(f,c,f(c)) \to c \in \mathbb{R}$
\item $\varphi_{16} :\equiv \forall y\,(\lvert f(c) - f(c) \rvert = y \to y = \lvert f(c) - f(c) \rvert)$
\item $\varphi_{17} :\equiv \forall x\,(x = x) \to \lvert f(c) - f(c) \rvert = \lvert f(c) - f(c) \rvert$
\item $\varphi_{18} :\equiv \forall f\,\forall c\,\forall \alpha\,(\mathrm{FLim}(f,c,\alpha) \to \mathrm{RFn}(f)) \to \forall c\,\forall \alpha\,(\mathrm{FLim}(f,c,\alpha) \to \mathrm{RFn}(f))$
\item $\varphi_{19} :\equiv \forall z\,\forall w\,(z = w \to z < \varepsilon \to w < \varepsilon) \to \forall w\,(0 = w \to 0 < \varepsilon \to w < \varepsilon)$
\item $\varphi_{20} :\equiv \forall \alpha\,(\mathrm{FLim}(f,c,\alpha) \to c \in \mathbb{R}) \to \mathrm{FLim}(f,c,f(c)) \to c \in \mathbb{R}$
\item $\varphi_{21} :\equiv x \in \mathrm{dom}(f) \to \lvert x - c \rvert < \delta \to x = c \to f(c) - f(c) = 0$
\item $\varphi_{22} :\equiv c \in \mathrm{dom}(f) \to \mathrm{FLim}(f,c,f(c)) \to \varepsilon \in \mathbb{R} \to \varepsilon \in \mathbb{R}$
\item $\varphi_{23} :\equiv \forall c\,\forall \alpha\,(\mathrm{FLim}(f,c,\alpha) \to c \in \mathbb{R}) \to \forall \alpha\,(\mathrm{FLim}(f,c,\alpha) \to c \in \mathbb{R})$
\item $\varphi_{24} :\equiv \varphi_{1} \to (\mathrm{Cont}(f,c) \leftrightarrow \mathrm{ContIn}(f,\mathrm{dom}(f),c))$
\item $\varphi_{25} :\equiv \lvert f(c) - f(c) \rvert = \lvert z \rvert \to \lvert f(c) - f(c) \rvert = \lvert w \rvert$
\item $\varphi_{26} :\equiv x \in \mathrm{dom}(f) \to \lvert x - c \rvert < \delta \to x = c \to \lvert f(c) - f(c) \rvert = 0$
\item $\varphi_{27} :\equiv x \in \mathrm{dom}(f) \to \lvert x - c \rvert < \delta \to x = c \to 0 = \lvert f(c) - f(c) \rvert$
\item $\varphi_{28} :\equiv \forall y\,\forall z\,(\lvert f(c) - f(c) \rvert = y \to y = z \to \lvert f(c) - f(c) \rvert = z)$
\item $\varphi_{29} :\equiv \forall x \in \mathbb{R}\,(x - x = 0) \to f(c) \in \mathbb{R} \to f(c) - f(c) = 0$
\item $\varphi_{30} :\equiv \forall f\,\forall c\,\forall \alpha\,(\mathrm{FLim}(f,c,\alpha) \to c \in \mathbb{R}) \to \forall c\,\forall \alpha\,(\mathrm{FLim}(f,c,\alpha) \to c \in \mathbb{R})$
\item $\varphi_{31} :\equiv \forall x \in \mathrm{dom}(f)\,(\lvert x - c \rvert < \delta \to \lvert f(x) - f(c) \rvert < \varepsilon)$
\item $\varphi_{32} :\equiv c \in \mathrm{dom}(f) \to \mathrm{FLim}(f,c,f(c)) \to \varepsilon \in \mathbb{R} \to 0 < \varepsilon \to 0 < \varepsilon$
\item $\varphi_{33} :\equiv x \in \mathrm{dom}(f) \to 0 < \lvert x - c \rvert \to \lvert x - c \rvert < \delta \to \lvert x - c \rvert < \delta$
\item $\varphi_{34} :\equiv z = w \to \lvert f(z) - f(c) \rvert < \varepsilon \to \lvert f(w) - f(c) \rvert < \varepsilon$
\item $\varphi_{35} :\equiv c = x \to \lvert f(c) - f(c) \rvert < \varepsilon \to \lvert f(x) - f(c) \rvert < \varepsilon$
\item $\varphi_{36} :\equiv 0 = \lvert f(c) - f(c) \rvert \to 0 < \varepsilon \to \lvert f(c) - f(c) \rvert < \varepsilon$
\item $\varphi_{37} :\equiv \lvert f(c) - f(c) \rvert = \lvert 0 \rvert \to \lvert 0 \rvert = 0 \to \lvert f(c) - f(c) \rvert = 0$
\item $\varphi_{38} :\equiv x \in \mathrm{dom}(f) \to \lvert x - c \rvert < \delta \to x = c \to \lvert f(c) - f(c) \rvert = \lvert 0 \rvert$
\item $\varphi_{39} :\equiv x \in \mathrm{dom}(f) \to \lvert x - c \rvert < \delta \to x = c \to \lvert f(c) - f(c) \rvert < \varepsilon$
\item $\varphi_{40} :\equiv x \in \mathrm{dom}(f) \to \lvert x - c \rvert < \delta \to x = c \to \lvert f(x) - f(c) \rvert < \varepsilon$
\item $\varphi_{41} :\equiv \forall x\,(\mathrm{RFn}(f) \to x \in \mathrm{dom}(f) \to x \in \mathbb{R}) \to \mathrm{RFn}(f) \to x \in \mathrm{dom}(f) \to x \in \mathbb{R}$
\item $\varphi_{42} :\equiv \forall x\,(\mathrm{RFn}(f) \to x \in \mathrm{dom}(f) \to x \in \mathbb{R}) \to \mathrm{RFn}(f) \to c \in \mathrm{dom}(f) \to c \in \mathbb{R}$
\item $\varphi_{43} :\equiv x \in \mathrm{dom}(f) \to \lvert x - c \rvert < \delta \to x \neq c \to \lvert f(x) - f(c) \rvert < \varepsilon$
\item $\varphi_{44} :\equiv \forall w\,(c = w \to \lvert f(c) - f(c) \rvert < \varepsilon \to \lvert f(w) - f(c) \rvert < \varepsilon)$
\item $\varphi_{45} :\equiv \forall z\,(\lvert f(c) - f(c) \rvert = \lvert 0 \rvert \to \lvert 0 \rvert = z \to \lvert f(c) - f(c) \rvert = z)$
\item $\varphi_{46} :\equiv \forall f\,\forall x\,(\mathrm{RFn}(f) \to x \in \mathrm{dom}(f) \to x \in \mathbb{R}) \to \forall x\,(\mathrm{RFn}(f) \to x \in \mathrm{dom}(f) \to x \in \mathbb{R})$
\item $\varphi_{47} :\equiv \mathrm{PNbhd}(f,c) \land c \in \mathrm{dom}(f) \land \mathrm{Cont}(f,c) \to \varphi_{2}$
\item $\varphi_{48} :\equiv \forall x \in \mathbb{R}\,(x \neq 0 \to 0 < \lvert x \rvert) \to x - c \in \mathbb{R} \to x - c \neq 0 \to 0 < \lvert x - c \rvert$
\item $\varphi_{49} :\equiv x \in \mathrm{dom}(f) \to 0 < \lvert x - c \rvert \to \lvert x - c \rvert < \delta \to \lvert f(x) - f(c) \rvert < \varepsilon$
\item $\varphi_{50} :\equiv \lvert f(c) - f(c) \rvert = \lvert f(c) - f(c) \rvert \to \lvert f(c) - f(c) \rvert = \lvert 0 \rvert$
\item $\varphi_{51} :\equiv \lvert x - c \rvert < \delta \to x \in \mathrm{dom}(f) \to \lvert x - c \rvert < \delta \to \lvert f(x) - f(c) \rvert < \varepsilon$
\item $\varphi_{52} :\equiv \forall y\,(x \in \mathbb{R} \to y \in \mathbb{R} \to x - y \in \mathbb{R}) \to x \in \mathbb{R} \to c \in \mathbb{R} \to x - c \in \mathbb{R}$
\item $\varphi_{53} :\equiv \forall x\,(\mathrm{RFn}(f) \to x \in \mathrm{dom}(f) \to f(x) \in \mathbb{R}) \to \mathrm{RFn}(f) \to c \in \mathrm{dom}(f) \to f(c) \in \mathbb{R}$
\item $\varphi_{54} :\equiv \forall x \in \mathrm{dom}(f)\,(0 < \lvert x - c \rvert \to \lvert x - c \rvert < \delta \to \lvert f(x) - f(c) \rvert < \varepsilon)$
\item $\varphi_{55} :\equiv \forall y\,(x \in \mathbb{R} \to y \in \mathbb{R} \to x \neq y \to x - y \neq 0) \to x \in \mathbb{R} \to c \in \mathbb{R} \to x \neq c \to x - c \neq 0$
\item $\varphi_{56} :\equiv (\mathrm{PNbhd}(f,c) \to c \in \mathrm{dom}(f) \to c \in \mathrm{dom}(f)) \to \varphi_{9}$
\item $\varphi_{57} :\equiv \forall x\,\forall y\,(x \in \mathbb{R} \to y \in \mathbb{R} \to x - y \in \mathbb{R}) \to \forall y\,(x \in \mathbb{R} \to y \in \mathbb{R} \to x - y \in \mathbb{R})$
\item $\varphi_{58} :\equiv \forall f\,\forall x\,(\mathrm{RFn}(f) \to x \in \mathrm{dom}(f) \to f(x) \in \mathbb{R}) \to \forall x\,(\mathrm{RFn}(f) \to x \in \mathrm{dom}(f) \to f(x) \in \mathbb{R})$
\item $\varphi_{59} :\equiv \forall x\,\forall y\,(x \in \mathbb{R} \to y \in \mathbb{R} \to x \neq y \to x - y \neq 0) \to \forall y\,(x \in \mathbb{R} \to y \in \mathbb{R} \to x \neq y \to x - y \neq 0)$
\item $\varphi_{60} :\equiv (\mathrm{PNbhd}(f,c) \to c \in \mathrm{dom}(f) \to c \in \mathrm{dom}(f)) \to \varphi_{14}$
\item $\varphi_{61} :\equiv \forall w\,(f(c) - f(c) = w \to \lvert f(c) - f(c) \rvert = \lvert f(c) - f(c) \rvert \to \lvert f(c) - f(c) \rvert = \lvert w \rvert)$
\item $\varphi_{62} :\equiv 0 < \varepsilon \to \delta \in \mathbb{R} \land 0 < \delta \land \varphi_{31} \to x \in \mathrm{dom}(f) \to x \in \mathrm{dom}(f)$
\item $\varphi_{63} :\equiv \mathrm{Cont}(f,c) \to \varepsilon \in \mathbb{R} \to 0 < \varepsilon \to \exists \delta\,(\delta \in \mathbb{R} \land 0 < \delta \land \varphi_{31})$
\item $\varphi_{64} :\equiv \varphi_{31} \to x \in \mathrm{dom}(f) \to \lvert x - c \rvert < \delta \to \lvert f(x) - f(c) \rvert < \varepsilon$
\item $\varphi_{65} :\equiv \mathrm{Cont}(f,c) \to \forall \varepsilon \in \mathbb{R}\,(0 < \varepsilon \to \exists \delta\,(\delta \in \mathbb{R} \land 0 < \delta \land \varphi_{31}))$
\item $\varphi_{66} :\equiv \mathrm{FLim}(f,c,f(c)) \to \varepsilon \in \mathbb{R} \to 0 < \varepsilon \to \exists \delta\,(\delta \in \mathbb{R} \land 0 < \delta \land \varphi_{31})$
\item $\varphi_{67} :\equiv \mathrm{FLim}(f,c,f(c)) \to \forall \varepsilon \in \mathbb{R}\,(0 < \varepsilon \to \exists \delta\,(\delta \in \mathbb{R} \land 0 < \delta \land \varphi_{31}))$
\item $\varphi_{68} :\equiv 0 < \varepsilon \to \delta \in \mathbb{R} \land 0 < \delta \land \varphi_{54} \to x \in \mathrm{dom}(f) \to x \in \mathrm{dom}(f)$
\item $\varphi_{69} :\equiv (\mathrm{FLim}(f,c,f(c)) \to \mathrm{RFn}(f)) \to (\mathrm{PNbhd}(f,c) \to \varphi_{10}) \to \varphi_{12}$
\item $\varphi_{70} :\equiv \mathrm{Cont}(f,c) \to \varepsilon \in \mathbb{R} \to 0 < \varepsilon \to \exists \delta\,(\delta \in \mathbb{R} \land 0 < \delta \land \varphi_{54})$
\item $\varphi_{71} :\equiv \delta \in \mathbb{R} \land 0 < \delta \land \varphi_{54} \to x \in \mathrm{dom}(f) \to \lvert x - c \rvert < \delta \to x = c \to 0 \leq 0$
\item $\varphi_{72} :\equiv \delta \in \mathbb{R} \land 0 < \delta \land \varphi_{54} \to x \in \mathrm{dom}(f) \to \lvert x - c \rvert < \delta \to x = c \to x = c$
\item $\varphi_{73} :\equiv \delta \in \mathbb{R} \land 0 < \delta \land \varphi_{54} \to x \in \mathrm{dom}(f) \to \lvert x - c \rvert < \delta \to x = c \to c = x$
\item $\varphi_{74} :\equiv \mathrm{PNbhd}(f,c) \land c \in \mathrm{dom}(f) \land \mathrm{Cont}(f,c) \land \varepsilon \in \mathbb{R} \land 0 < \varepsilon \land (\delta \in \mathbb{R} \land 0 < \delta \land \varphi_{31})$
\item $\varphi_{75} :\equiv \mathrm{Cont}(f,c) \to \forall \varepsilon \in \mathbb{R}\,(0 < \varepsilon \to \exists \delta\,(\delta \in \mathbb{R} \land 0 < \delta \land \varphi_{54}))$
\item $\varphi_{76} :\equiv \forall \alpha\,\forall \varepsilon\,(\mathrm{FLim}(f,c,\alpha) \to \varepsilon \in \mathbb{R} \to 0 < \varepsilon \to \exists \delta\,(\delta \in \mathbb{R} \land 0 < \delta \land \forall x \in \mathrm{dom}(f)\,(0 < \lvert x - c \rvert \to \lvert x - c \rvert < \delta \to \lvert f(x) - \alpha \rvert < \varepsilon)))$
\item $\varphi_{77} :\equiv \delta \in \mathbb{R} \land 0 < \delta \land \varphi_{54} \to x \in \mathrm{dom}(f) \to \lvert x - c \rvert < \delta \to x \in \mathbb{R}$
\item $\varphi_{78} :\equiv \delta \in \mathbb{R} \land 0 < \delta \land \varphi_{54} \to x \in \mathrm{dom}(f) \to \lvert x - c \rvert < \delta \to x \neq c \to x \neq c$
\item $\varphi_{79} :\equiv (\mathrm{FLim}(f,c,f(c)) \to c \in \mathbb{R}) \to (\mathrm{PNbhd}(f,c) \to \varphi_{10}) \to \varphi_{15}$
\item $\varphi_{80} :\equiv \mathrm{FLim}(f,c,f(c)) \to \varepsilon \in \mathbb{R} \to 0 < \varepsilon \to \exists \delta\,(\delta \in \mathbb{R} \land 0 < \delta \land \varphi_{54})$
\item $\varphi_{81} :\equiv \delta \in \mathbb{R} \land 0 < \delta \land \varphi_{54} \to x \in \mathrm{dom}(f) \to \lvert x - c \rvert < \delta \to x = c \to \lvert 0 \rvert = 0$
\item $\varphi_{82} :\equiv \delta \in \mathbb{R} \land 0 < \delta \land \varphi_{54} \to x \in \mathrm{dom}(f) \to \lvert x - c \rvert < \delta \to x \neq c \to x - c \neq 0$
\item $\varphi_{83} :\equiv \delta \in \mathbb{R} \land 0 < \delta \land \varphi_{54} \to x \in \mathrm{dom}(f) \to \lvert x - c \rvert < \delta \to \lvert x - c \rvert < \delta$
\item $\varphi_{84} :\equiv \delta \in \mathbb{R} \land 0 < \delta \land \varphi_{54} \to x \in \mathrm{dom}(f) \to \lvert x - c \rvert < \delta \to x \neq c \to x - c \in \mathbb{R}$
\item $\varphi_{85} :\equiv \delta \in \mathbb{R} \land 0 < \delta \land \varphi_{54} \to x \in \mathrm{dom}(f) \to \lvert x - c \rvert < \delta \to x \neq c \to 0 < \lvert x - c \rvert$
\item $\varphi_{86} :\equiv \delta \in \mathbb{R} \land 0 < \delta \land \varphi_{54} \to x \in \mathrm{dom}(f) \to \lvert x - c \rvert < \delta \to x = c \to f(c) \in \mathbb{R}$
\item $\varphi_{87} :\equiv (\mathrm{PNbhd}(f,c) \to \varphi_{4}) \to (\mathrm{PNbhd}(f,c) \to \varphi_{3}) \to \mathrm{PNbhd}(f,c) \to \varphi_{5}$
\item $\varphi_{88} :\equiv \forall A\,\forall c\,(\mathrm{ContIn}(f,A,c) \leftrightarrow \mathrm{RFn}(f) \land A \subseteq \mathrm{dom}(f) \land c \in A \land \forall \varepsilon \in \mathbb{R}\,(0 < \varepsilon \to \exists \delta\,(\delta \in \mathbb{R} \land 0 < \delta \land \forall x \in A\,(\lvert x - c \rvert < \delta \to \lvert f(x) - f(c) \rvert < \varepsilon))))$
\item $\varphi_{89} :\equiv \mathrm{PNbhd}(f,c) \land c \in \mathrm{dom}(f) \land \mathrm{Cont}(f,c) \to \varepsilon \in \mathbb{R} \to 0 < \varepsilon \to \exists \delta\,(\delta \in \mathbb{R} \land 0 < \delta \land \varphi_{54})$
\item $\varphi_{90} :\equiv \forall \varepsilon\,(\varphi_{7}) \to \forall \varepsilon \in \mathbb{R}\,(0 < \varepsilon \to \exists \delta\,(\delta \in \mathbb{R} \land 0 < \delta \land \varphi_{31}))$
\item $\varphi_{91} :\equiv \delta \in \mathbb{R} \land 0 < \delta \land \varphi_{54} \to x \in \mathrm{dom}(f) \to \lvert x - c \rvert < \delta \to \lvert f(x) - f(c) \rvert < \varepsilon$
\item $\varphi_{92} :\equiv \mathrm{PNbhd}(f,c) \land c \in \mathrm{dom}(f) \land \mathrm{Cont}(f,c) \land \varepsilon \in \mathbb{R} \land 0 < \varepsilon \land \neg \exists \delta\,(\delta \in \mathbb{R} \land 0 < \delta \land \varphi_{54})$
\item $\varphi_{93} :\equiv \mathrm{PNbhd}(f,c) \to c \in \mathrm{dom}(f) \to \mathrm{Cont}(f,c) \to \varepsilon \in \mathbb{R} \to \varphi_{62}$
\item $\varphi_{94} :\equiv \mathrm{ContIn}(f,\mathrm{dom}(f),c) \leftrightarrow \mathrm{RFn}(f) \land \mathrm{dom}(f) \subseteq \mathrm{dom}(f) \land c \in \mathrm{dom}(f) \land \forall \varepsilon \in \mathbb{R}\,(0 < \varepsilon \to \exists \delta\,(\delta \in \mathbb{R} \land 0 < \delta \land \varphi_{31}))$
\item $\varphi_{95} :\equiv \forall \delta\,(\neg (\delta \in \mathbb{R} \land 0 < \delta \land \varphi_{31})) \to \neg (\delta \in \mathbb{R} \land 0 < \delta \land \varphi_{31})$
\item $\varphi_{96} :\equiv \mathrm{Cont}(f,c) \to \varepsilon \in \mathbb{R} \to 0 < \varepsilon \to \delta \in \mathbb{R} \land 0 < \delta \land \varphi_{31} \to \varphi_{33}$
\item $\varphi_{97} :\equiv \forall \alpha\,(\mathrm{FLim}(f,c,\alpha) \leftrightarrow \mathrm{RFn}(f) \land c \in \mathbb{R} \land \alpha \in \mathbb{R} \land \mathrm{PNbhd}(f,c) \land \forall \varepsilon \in \mathbb{R}\,(0 < \varepsilon \to \exists \delta\,(\delta \in \mathbb{R} \land 0 < \delta \land \forall x \in \mathrm{dom}(f)\,(0 < \lvert x - c \rvert \to \lvert x - c \rvert < \delta \to \lvert f(x) - \alpha \rvert < \varepsilon))))$
\item $\varphi_{98} :\equiv \mathrm{PNbhd}(f,c) \to c \in \mathrm{dom}(f) \to \mathrm{FLim}(f,c,f(c)) \to \varepsilon \in \mathbb{R} \to \varphi_{68}$
\item $\varphi_{99} :\equiv \delta \in \mathbb{R} \land 0 < \delta \land \varphi_{31} \to x \in \mathrm{dom}(f) \to 0 < \lvert x - c \rvert \to \lvert x - c \rvert < \delta \to \varphi_{31}$
\item $\varphi_{100} :\equiv \mathrm{Cont}(f,c) \to \varepsilon \in \mathbb{R} \to 0 < \varepsilon \to \delta \in \mathbb{R} \land 0 < \delta \land \varphi_{31} \to \varphi_{49}$
\item $\varphi_{101} :\equiv c \in \mathrm{dom}(f) \to \mathrm{FLim}(f,c,f(c)) \to \varepsilon \in \mathbb{R} \to 0 < \varepsilon \to \varphi_{71}$
\item $\varphi_{102} :\equiv c \in \mathrm{dom}(f) \to \mathrm{FLim}(f,c,f(c)) \to \varepsilon \in \mathbb{R} \to 0 < \varepsilon \to \varphi_{72}$
\item $\varphi_{103} :\equiv c \in \mathrm{dom}(f) \to \mathrm{FLim}(f,c,f(c)) \to \varepsilon \in \mathbb{R} \to 0 < \varepsilon \to \varphi_{73}$
\item $\varphi_{104} :\equiv c \in \mathrm{dom}(f) \to \mathrm{FLim}(f,c,f(c)) \to \varepsilon \in \mathbb{R} \to 0 < \varepsilon \to \varphi_{77}$
\item $\varphi_{105} :\equiv \mathrm{FLim}(f,c,f(c)) \leftrightarrow \mathrm{RFn}(f) \land c \in \mathbb{R} \land f(c) \in \mathbb{R} \land \mathrm{PNbhd}(f,c) \land \forall \varepsilon \in \mathbb{R}\,(0 < \varepsilon \to \exists \delta\,(\delta \in \mathbb{R} \land 0 < \delta \land \varphi_{54}))$
\item $\varphi_{106} :\equiv c \in \mathrm{dom}(f) \to \mathrm{FLim}(f,c,f(c)) \to \varepsilon \in \mathbb{R} \to 0 < \varepsilon \to \varphi_{78}$
\item $\varphi_{107} :\equiv \mathrm{Cont}(f,c) \to \varepsilon \in \mathbb{R} \to 0 < \varepsilon \to \delta \in \mathbb{R} \land 0 < \delta \land \varphi_{31} \to \varphi_{54}$
\item $\varphi_{108} :\equiv \mathrm{FLim}(f,c,f(c)) \to \varepsilon \in \mathbb{R} \to 0 < \varepsilon \to \delta \in \mathbb{R} \land 0 < \delta \land \varphi_{54} \to \varphi_{21}$
\item $\varphi_{109} :\equiv c \in \mathrm{dom}(f) \to \mathrm{FLim}(f,c,f(c)) \to \varepsilon \in \mathbb{R} \to 0 < \varepsilon \to \varphi_{81}$
\item $\varphi_{110} :\equiv c \in \mathrm{dom}(f) \to \mathrm{FLim}(f,c,f(c)) \to \varepsilon \in \mathbb{R} \to 0 < \varepsilon \to \varphi_{82}$
\item $\varphi_{111} :\equiv \mathrm{Cont}(f,c) \to \varepsilon \in \mathbb{R} \to 0 < \varepsilon \to \delta \in \mathbb{R} \land 0 < \delta \land \varphi_{31} \to \delta \in \mathbb{R} \land 0 < \delta \land \varphi_{31}$
\item $\varphi_{112} :\equiv 0 < \varepsilon \to \delta \in \mathbb{R} \land 0 < \delta \land \varphi_{31} \to x \in \mathrm{dom}(f) \to 0 < \lvert x - c \rvert \to \varphi_{51}$
\item $\varphi_{113} :\equiv c \in \mathrm{dom}(f) \to \mathrm{FLim}(f,c,f(c)) \to \varepsilon \in \mathbb{R} \to 0 < \varepsilon \to \varphi_{83}$
\item $\varphi_{114} :\equiv \mathrm{FLim}(f,c,f(c)) \to \varepsilon \in \mathbb{R} \to 0 < \varepsilon \to \delta \in \mathbb{R} \land 0 < \delta \land \varphi_{54} \to \varphi_{26}$
\item $\varphi_{115} :\equiv \mathrm{FLim}(f,c,f(c)) \to \varepsilon \in \mathbb{R} \to 0 < \varepsilon \to \delta \in \mathbb{R} \land 0 < \delta \land \varphi_{54} \to \varphi_{27}$
\item $\varphi_{116} :\equiv \mathrm{FLim}(f,c,f(c)) \to \varepsilon \in \mathbb{R} \to 0 < \varepsilon \to \delta \in \mathbb{R} \land 0 < \delta \land \varphi_{54} \to \varphi_{31}$
\item $\varphi_{117} :\equiv \forall \delta\,(\neg (\delta \in \mathbb{R} \land 0 < \delta \land \varphi_{54})) \to \neg (\delta \in \mathbb{R} \land 0 < \delta \land \varphi_{54})$
\item $\varphi_{118} :\equiv c \in \mathrm{dom}(f) \to \mathrm{FLim}(f,c,f(c)) \to \varepsilon \in \mathbb{R} \to 0 < \varepsilon \to \varphi_{84}$
\item $\varphi_{119} :\equiv c \in \mathrm{dom}(f) \to \mathrm{FLim}(f,c,f(c)) \to \varepsilon \in \mathbb{R} \to 0 < \varepsilon \to \varphi_{85}$
\item $\varphi_{120} :\equiv c \in \mathrm{dom}(f) \to \mathrm{FLim}(f,c,f(c)) \to \varepsilon \in \mathbb{R} \to 0 < \varepsilon \to \varphi_{86}$
\item $\varphi_{121} :\equiv \mathrm{FLim}(f,c,f(c)) \to \varepsilon \in \mathbb{R} \to 0 < \varepsilon \to \delta \in \mathbb{R} \land 0 < \delta \land \varphi_{54} \to \varphi_{38}$
\item $\varphi_{122} :\equiv \mathrm{FLim}(f,c,f(c)) \to \varepsilon \in \mathbb{R} \to 0 < \varepsilon \to \delta \in \mathbb{R} \land 0 < \delta \land \varphi_{54} \to \varphi_{39}$
\item $\varphi_{123} :\equiv \mathrm{FLim}(f,c,f(c)) \to \varepsilon \in \mathbb{R} \to 0 < \varepsilon \to \delta \in \mathbb{R} \land 0 < \delta \land \varphi_{54} \to \varphi_{40}$
\item $\varphi_{124} :\equiv \mathrm{FLim}(f,c,f(c)) \to \varepsilon \in \mathbb{R} \to 0 < \varepsilon \to \delta \in \mathbb{R} \land 0 < \delta \land \varphi_{54} \to \varphi_{43}$
\item $\varphi_{125} :\equiv (\mathrm{PNbhd}(f,c) \to c \in \mathrm{dom}(f) \to \varphi_{67}) \to \mathrm{PNbhd}(f,c) \to \varphi_{3}$
\item $\varphi_{126} :\equiv \varepsilon \in \mathbb{R} \to 0 < \varepsilon \to \delta \in \mathbb{R} \land 0 < \delta \land \varphi_{54} \to \exists \delta\,(\delta \in \mathbb{R} \land 0 < \delta \land \varphi_{31})$
\item $\varphi_{127} :\equiv \varepsilon \in \mathbb{R} \to 0 < \varepsilon \to \delta \in \mathbb{R} \land 0 < \delta \land \varphi_{31} \to \exists \delta\,(\delta \in \mathbb{R} \land 0 < \delta \land \varphi_{54})$
\item $\varphi_{128} :\equiv c \in \mathrm{dom}(f) \to \mathrm{FLim}(f,c,f(c)) \to \varepsilon \in \mathbb{R} \to 0 < \varepsilon \to \varphi_{91}$
\item $\varphi_{129} :\equiv \mathrm{Cont}(f,c) \to \varepsilon \in \mathbb{R} \to 0 < \varepsilon \to \delta \in \mathbb{R} \land 0 < \delta \land \varphi_{31} \to \delta \in \mathbb{R} \land 0 < \delta \land \varphi_{54}$
\item $\varphi_{130} :\equiv (\mathrm{PNbhd}(f,c) \to c \in \mathrm{dom}(f) \to \varphi_{75}) \to \mathrm{PNbhd}(f,c) \to \varphi_{4}$
\item $\varphi_{131} :\equiv \mathrm{FLim}(f,c,f(c)) \to \varepsilon \in \mathbb{R} \to 0 < \varepsilon \to \delta \in \mathbb{R} \land 0 < \delta \land \varphi_{54} \to \varphi_{49}$
\item $\varphi_{132} :\equiv \mathrm{FLim}(f,c,f(c)) \to \varepsilon \in \mathbb{R} \to 0 < \varepsilon \to \delta \in \mathbb{R} \land 0 < \delta \land \varphi_{54} \to \varphi_{54}$
\item $\varphi_{133} :\equiv \forall \varepsilon \in \mathbb{R}\,(0 < \varepsilon \to \exists \delta\,(\delta \in \mathbb{R} \land 0 < \delta \land \varphi_{31})) \to \varepsilon \in \mathbb{R} \to 0 < \varepsilon \to \exists \delta\,(\delta \in \mathbb{R} \land 0 < \delta \land \varphi_{31})$
\item $\varphi_{134} :\equiv \mathrm{FLim}(f,c,f(c)) \to \varepsilon \in \mathbb{R} \to 0 < \varepsilon \to \delta \in \mathbb{R} \land 0 < \delta \land \varphi_{54} \to \delta \in \mathbb{R} \land 0 < \delta \land \varphi_{31}$
\item $\varphi_{135} :\equiv \varphi_{7} \land \varepsilon \in \mathbb{R} \land 0 < \varepsilon \land (\delta \in \mathbb{R} \land 0 < \delta \land \varphi_{54}) \to x \in \mathrm{dom}(f) \to \lvert x - c \rvert < \delta \to \lvert f(x) - f(c) \rvert < \varepsilon$
\item $\varphi_{136} :\equiv \mathrm{FLim}(f,c,f(c)) \to \varepsilon \in \mathbb{R} \to 0 < \varepsilon \to \delta \in \mathbb{R} \land 0 < \delta \land \varphi_{54} \to \delta \in \mathbb{R} \land 0 < \delta \land \varphi_{54}$
\item $\varphi_{137} :\equiv (\mathrm{Cont}(f,c) \leftrightarrow \mathrm{ContIn}(f,\mathrm{dom}(f),c)) \to \varphi_{94} \to \varphi_{6}$
\item $\varphi_{138} :\equiv (\mathrm{PNbhd}(f,c) \to \varphi_{32}) \to \mathrm{PNbhd}(f,c) \to c \in \mathrm{dom}(f) \to \varphi_{80}$
\item $\varphi_{139} :\equiv \mathrm{PNbhd}(f,c) \to c \in \mathrm{dom}(f) \to \mathrm{Cont}(f,c) \to \varphi_{127}$
\item $\varphi_{140} :\equiv (\varphi_{7} \to \forall \varepsilon\,(\varphi_{7})) \to \mathrm{PNbhd}(f,c) \to c \in \mathrm{dom}(f) \to \varphi_{67}$
\item $\varphi_{141} :\equiv (\mathrm{PNbhd}(f,c) \to c \in \mathrm{dom}(f) \to c \in \mathrm{dom}(f)) \to \varphi_{125}$
\item $\varphi_{142} :\equiv \mathrm{PNbhd}(f,c) \to c \in \mathrm{dom}(f) \to \mathrm{Cont}(f,c) \to \varepsilon \in \mathbb{R} \to \varphi_{112}$
\item $\varphi_{143} :\equiv \mathrm{PNbhd}(f,c) \to c \in \mathrm{dom}(f) \to \mathrm{Cont}(f,c) \to \varepsilon \in \mathbb{R} \to 0 < \varepsilon \to \varphi_{99}$
\item $\varphi_{144} :\equiv \mathrm{PNbhd}(f,c) \to c \in \mathrm{dom}(f) \to \mathrm{FLim}(f,c,f(c)) \to \varphi_{126}$
\item $\varphi_{145} :\equiv \varphi_{7} \land \varepsilon \in \mathbb{R} \land 0 < \varepsilon \land \neg \exists \delta\,(\delta \in \mathbb{R} \land 0 < \delta \land \varphi_{31}) \to \neg (\delta \in \mathbb{R} \land 0 < \delta \land \varphi_{54})$
\item $\varphi_{146} :\equiv (\mathrm{PNbhd}(f,c) \to c \in \mathrm{dom}(f) \to c \in \mathrm{dom}(f)) \to \mathrm{PNbhd}(f,c) \to \varphi_{120}$
\item $\varphi_{147} :\equiv \forall \delta\,(\varphi_{7} \land \varepsilon \in \mathbb{R} \land 0 < \varepsilon \land \neg \exists \delta\,(\delta \in \mathbb{R} \land 0 < \delta \land \varphi_{31})) \to \forall \delta\,(\neg (\delta \in \mathbb{R} \land 0 < \delta \land \varphi_{54}))$
\item $\varphi_{148} :\equiv (\lvert f(c) - f(c) \rvert = 0 \to 0 = \lvert f(c) - f(c) \rvert) \to \mathrm{PNbhd}(f,c) \to c \in \mathrm{dom}(f) \to \varphi_{115}$
\item $\varphi_{149} :\equiv (\mathrm{PNbhd}(f,c) \to c \in \mathrm{dom}(f) \to \varphi_{66}) \to \varphi_{7} \to \varepsilon \in \mathbb{R} \to 0 < \varepsilon \to \exists \delta\,(\delta \in \mathbb{R} \land 0 < \delta \land \varphi_{31})$
\item $\varphi_{150} :\equiv \varphi_{7} \land \varepsilon \in \mathbb{R} \land 0 < \varepsilon \land \neg \exists \delta\,(\delta \in \mathbb{R} \land 0 < \delta \land \varphi_{31}) \to \forall \delta\,(\varphi_{7} \land \varepsilon \in \mathbb{R} \land 0 < \varepsilon \land \neg \exists \delta\,(\delta \in \mathbb{R} \land 0 < \delta \land \varphi_{31}))$
\item $\varphi_{151} :\equiv \forall \varepsilon\,(\varphi_{7} \to \varepsilon \in \mathbb{R} \to 0 < \varepsilon \to \exists \delta\,(\delta \in \mathbb{R} \land 0 < \delta \land \varphi_{31})) \to \varphi_{90}$
\item $\varphi_{152} :\equiv (f(c) - f(c) = 0 \to \varphi_{50}) \to \mathrm{PNbhd}(f,c) \to c \in \mathrm{dom}(f) \to \varphi_{121}$
\item $\varphi_{153} :\equiv \varphi_{7} \land \varepsilon \in \mathbb{R} \land 0 < \varepsilon \land (\delta \in \mathbb{R} \land 0 < \delta \land \varphi_{54}) \to \forall x\,(\varphi_{7} \land \varepsilon \in \mathbb{R} \land 0 < \varepsilon \land (\delta \in \mathbb{R} \land 0 < \delta \land \varphi_{54}))$
\item $\varphi_{154} :\equiv \forall \varepsilon\,(\varphi_{89}) \to \varphi_{2} \to \forall \varepsilon \in \mathbb{R}\,(0 < \varepsilon \to \exists \delta\,(\delta \in \mathbb{R} \land 0 < \delta \land \varphi_{54}))$
\item $\varphi_{155} :\equiv (\mathrm{Cont}(f,c) \leftrightarrow \mathrm{ContIn}(f,\mathrm{dom}(f),c)) \to \varphi_{94} \to \mathrm{PNbhd}(f,c) \to c \in \mathrm{dom}(f) \to \varphi_{65}$
\item $\varphi_{156} :\equiv (\mathrm{PNbhd}(f,c) \to \varphi_{32}) \to \varphi_{36} \to \mathrm{PNbhd}(f,c) \to c \in \mathrm{dom}(f) \to \varphi_{122}$
\item $\varphi_{157} :\equiv (\mathrm{PNbhd}(f,c) \to \varphi_{10}) \to (\mathrm{PNbhd}(f,c) \to \varphi_{22}) \to \varphi_{138}$
\item $\varphi_{158} :\equiv \varphi_{6} \to \varphi_{9} \to \varphi_{14} \to (\mathrm{PNbhd}(f,c) \to \mathrm{PNbhd}(f,c)) \to \varphi_{130}$
\item $\varphi_{159} :\equiv (\mathrm{PNbhd}(f,c) \to \varphi_{102}) \to (x = c \to c = x) \to \mathrm{PNbhd}(f,c) \to \varphi_{103}$
\item $\varphi_{160} :\equiv (\mathrm{PNbhd}(f,c) \to c \in \mathrm{dom}(f) \to \varphi_{96}) \to \mathrm{PNbhd}(f,c) \to c \in \mathrm{dom}(f) \to \varphi_{100}$
\item $\varphi_{161} :\equiv (0 \leq 0 \to \lvert 0 \rvert = 0) \to (\mathrm{PNbhd}(f,c) \to \varphi_{101}) \to \mathrm{PNbhd}(f,c) \to \varphi_{109}$
\item $\varphi_{162} :\equiv (\mathrm{PNbhd}(f,c) \to \varphi_{120}) \to \mathrm{PNbhd}(f,c) \to c \in \mathrm{dom}(f) \to \varphi_{108}$
\item $\varphi_{163} :\equiv (\mathrm{PNbhd}(f,c) \to \varphi_{113}) \to \mathrm{PNbhd}(f,c) \to c \in \mathrm{dom}(f) \to \varphi_{124}$
\item $\varphi_{164} :\equiv (\mathrm{PNbhd}(f,c) \to c \in \mathrm{dom}(f) \to \varphi_{124}) \to \mathrm{PNbhd}(f,c) \to \varphi_{128}$
\item $\varphi_{165} :\equiv \forall x\,(\varphi_{135}) \to \forall x\,(\varphi_{7} \land \varepsilon \in \mathbb{R} \land 0 < \varepsilon \land (\delta \in \mathbb{R} \land 0 < \delta \land \varphi_{54})) \to \varphi_{31}$
\item $\varphi_{166} :\equiv (\mathrm{PNbhd}(f,c) \to c \in \mathrm{dom}(f) \to \varphi_{107}) \to \mathrm{PNbhd}(f,c) \to c \in \mathrm{dom}(f) \to \varphi_{129}$
\item $\varphi_{167} :\equiv (\mathrm{PNbhd}(f,c) \to c \in \mathrm{dom}(f) \to \varphi_{116}) \to \mathrm{PNbhd}(f,c) \to c \in \mathrm{dom}(f) \to \varphi_{134}$
\item $\varphi_{168} :\equiv (\mathrm{RFn}(f) \to x \in \mathrm{dom}(f) \to x \in \mathbb{R}) \to \varphi_{12} \to \varphi_{98} \to \mathrm{PNbhd}(f,c) \to \varphi_{104}$
\item $\varphi_{169} :\equiv (\varphi_{92} \to \forall \delta\,(\varphi_{92})) \to \mathrm{PNbhd}(f,c) \to c \in \mathrm{dom}(f) \to \varphi_{70}$
\item $\varphi_{170} :\equiv (\mathrm{PNbhd}(f,c) \to \varphi_{109}) \to \varphi_{37} \to \mathrm{PNbhd}(f,c) \to c \in \mathrm{dom}(f) \to \varphi_{114}$
\item $\varphi_{171} :\equiv \forall \delta\,(\varphi_{92} \to \neg (\delta \in \mathbb{R} \land 0 < \delta \land \varphi_{31})) \to \forall \delta\,(\varphi_{92}) \to \forall \delta\,(\neg (\delta \in \mathbb{R} \land 0 < \delta \land \varphi_{31}))$
\item $\varphi_{172} :\equiv (\varphi_{74} \to \forall x\,(\varphi_{74})) \to \mathrm{PNbhd}(f,c) \to c \in \mathrm{dom}(f) \to \varphi_{107}$
\item $\varphi_{173} :\equiv (\mathrm{PNbhd}(f,c) \to c \in \mathrm{dom}(f) \to \varphi_{136}) \to \mathrm{PNbhd}(f,c) \to c \in \mathrm{dom}(f) \to \varphi_{132}$
\item $\varphi_{174} :\equiv (\mathrm{Cont}(f,c) \leftrightarrow \mathrm{ContIn}(f,\mathrm{dom}(f),c)) \to \varphi_{94} \to \varphi_{12} \to \mathrm{dom}(f) \subseteq \mathrm{dom}(f) \to \varphi_{141}$
\item $\varphi_{175} :\equiv (x \in \mathbb{R} \to c \in \mathbb{R} \to x - c \in \mathbb{R}) \to (\mathrm{PNbhd}(f,c) \to \varphi_{104}) \to \varphi_{15} \to \mathrm{PNbhd}(f,c) \to \varphi_{118}$
\item $\varphi_{176} :\equiv (\mathrm{PNbhd}(f,c) \to c \in \mathrm{dom}(f) \to \varphi_{65}) \to \varphi_{133} \to \mathrm{PNbhd}(f,c) \to c \in \mathrm{dom}(f) \to \varphi_{63}$
\item $\varphi_{177} :\equiv (\mathrm{PNbhd}(f,c) \to c \in \mathrm{dom}(f) \to \varphi_{122}) \to \varphi_{35} \to \mathrm{PNbhd}(f,c) \to c \in \mathrm{dom}(f) \to \varphi_{123}$
\item $\varphi_{178} :\equiv (\mathrm{PNbhd}(f,c) \to c \in \mathrm{dom}(f) \to \varphi_{108}) \to \lvert f(c) - f(c) \rvert = \lvert f(c) - f(c) \rvert \to \varphi_{152}$
\item $\varphi_{179} :\equiv (\mathrm{PNbhd}(f,c) \to c \in \mathrm{dom}(f) \to \varphi_{132}) \to (\varphi_{54} \to \varphi_{49}) \to \mathrm{PNbhd}(f,c) \to c \in \mathrm{dom}(f) \to \varphi_{131}$
\item $\varphi_{180} :\equiv (\mathrm{PNbhd}(f,c) \to \varphi_{118}) \to (\mathrm{PNbhd}(f,c) \to \varphi_{110}) \to \mathrm{PNbhd}(f,c) \to \varphi_{119}$
\item $\varphi_{181} :\equiv (\mathrm{PNbhd}(f,c) \to \varphi_{104}) \to \varphi_{15} \to (\mathrm{PNbhd}(f,c) \to \varphi_{106}) \to \mathrm{PNbhd}(f,c) \to \varphi_{110}$
\item $\varphi_{182} :\equiv \forall \varepsilon\,(\varphi_{7} \to \varepsilon \in \mathbb{R} \to 0 < \varepsilon \to \exists \delta\,(\delta \in \mathbb{R} \land 0 < \delta \land \varphi_{31})) \to \varphi_{151} \to \varphi_{140}$
\item $\varphi_{183} :\equiv \forall \varepsilon\,(\varphi_{89}) \to \varphi_{154} \to \varphi_{47} \to \mathrm{PNbhd}(f,c) \to c \in \mathrm{dom}(f) \to \varphi_{75}$
\item $\varphi_{184} :\equiv (\forall \delta\,(\varphi_{145}) \to \varphi_{147}) \to \varphi_{150} \to \mathrm{PNbhd}(f,c) \to c \in \mathrm{dom}(f) \to \varphi_{66}$
\item $\varphi_{185} :\equiv (\mathrm{PNbhd}(f,c) \to c \in \mathrm{dom}(f) \to \varphi_{131}) \to \varphi_{98} \to (\mathrm{PNbhd}(f,c) \to \varphi_{119}) \to \varphi_{163}$
\item $\varphi_{186} :\equiv \forall x\,(\varphi_{135}) \to \varphi_{165} \to \varphi_{153} \to \mathrm{PNbhd}(f,c) \to c \in \mathrm{dom}(f) \to \varphi_{116}$
\item $\varphi_{187} :\equiv (\mathrm{PNbhd}(f,c) \to c \in \mathrm{dom}(f) \to \varphi_{63}) \to \forall \delta\,(\varphi_{92} \to \neg (\delta \in \mathbb{R} \land 0 < \delta \land \varphi_{31})) \to \varphi_{171} \to \varphi_{169}$
\item $\varphi_{188} :\equiv (\mathrm{PNbhd}(f,c) \to c \in \mathrm{dom}(f) \to \varphi_{80}) \to \forall \delta\,(\varphi_{145}) \to \varphi_{184}$
\end{abbrevlist}
\end{fullproof}

% 自動生成 (prover/build.js): IVT  ノード数 1304
\begin{fullproof}{定理2.14}{$\forall f\,\forall a\,\forall b\,\forall \mu\,(\mathrm{ContOn}(f,[a,b]) \to a \in \mathbb{R} \to b \in \mathbb{R} \to \mu \in \mathbb{R} \to a < b \to f(a) < \mu \to \mu < f(b) \to \exists \xi\,(\xi \in \mathbb{R} \land a < \xi \land \xi < b \land f(\xi) = \mu))$}{1304}\phantomsection\label{tree:IVT}
\noindent$\Pi_{1}$:
\begin{ptree}
\AXC{$\forall a\,\forall X\,(\mathrm{UB}(a,X) \leftrightarrow a \in \mathbb{R} \land \forall x \in X\,(x \leq a))$\lbl{DefUB}}
\AXC{$\varphi_{205}$\lbl{A2}}
\BIC{$\forall X\,(\mathrm{UB}(\xi,X) \leftrightarrow \xi \in \mathbb{R} \land \forall x \in X\,(x \leq \xi))$}
\AXC{$\varphi_{174}$\lbl{A2}}
\BIC{$\mathrm{UB}(\xi,W) \leftrightarrow \xi \in \mathbb{R} \land \forall x \in W\,(x \leq \xi)$}
\end{ptree}
\noindent$\Pi_{2}$:
\begin{ptree}
\AXC{$\mathrm{ContOn}(f,[a,b]) \to \varphi_{327}$\lbl{A1}}
\AXC{$(\mathrm{ContOn}(f,[a,b]) \to \varphi_{327}) \to \varphi_{5} \to \varphi_{462}$\lbl{A1}}
\BIC{$\varphi_{5} \to \varphi_{462}$}
\end{ptree}
\noindent$\Pi_{3}$:
\begin{ptree}
\AXC{$\stackrel{\Pi_{1}}{\vdots}$}\noLine
\UIC{$\mathrm{UB}(\xi,W) \leftrightarrow \xi \in \mathbb{R} \land \forall x \in W\,(x \leq \xi)$}
\AXC{$\varphi_{7}$\lbl{DefSup}}
\AXC{$\varphi_{7} \to \varphi_{8}$\lbl{A2}}
\BIC{$\varphi_{8}$}
\AXC{$\varphi_{8} \to \varphi_{5}$\lbl{A2}}
\BIC{$\varphi_{5}$}
\AXC{$\stackrel{\Pi_{2}}{\vdots}$}\noLine
\UIC{$\varphi_{5} \to \varphi_{462}$}
\BIC{$\varphi_{462}$}
\BIC{$\mathrm{ContOn}(f,[a,b]) \to \varphi_{328}$}
\end{ptree}
\noindent$\Pi_{4}$:
\begin{ptree}
\AXC{$\forall x\,\forall e\,(x \in \mathbb{R} \to e \in \mathbb{R} \to 0 < e \to x - e < x)$\lbl{[E34]}}
\AXC{$\varphi_{273}$\lbl{A2}}
\BIC{$\forall e\,(\xi \in \mathbb{R} \to e \in \mathbb{R} \to 0 < e \to \xi - e < \xi)$}
\AXC{$\varphi_{263}$\lbl{A2}}
\BIC{$\xi \in \mathbb{R} \to d \in \mathbb{R} \to 0 < d \to \xi - d < \xi$}
\AXC{$(\xi \in \mathbb{R} \to d \in \mathbb{R} \to 0 < d \to \xi - d < \xi) \to \varphi_{870}$\lbl{A1}}
\BIC{$\varphi_{870}$}
\end{ptree}
\noindent$\Pi_{5}$:
\begin{ptree}
\AXC{$\mathrm{ContOn}(f,[a,b]) \to \varphi_{678}$\lbl{A1}}
\AXC{$\varphi_{811}$\lbl{A1}}
\BIC{$\mathrm{ContOn}(f,[a,b]) \to \varphi_{610}$}
\AXC{$\stackrel{\Pi_{3}}{\vdots}$}\noLine
\UIC{$\mathrm{ContOn}(f,[a,b]) \to \varphi_{328}$}
\AXC{$\stackrel{\Pi_{4}}{\vdots}$}\noLine
\UIC{$\varphi_{870}$}
\BIC{$(\mathrm{ContOn}(f,[a,b]) \to \varphi_{610}) \to \varphi_{780}$}
\BIC{$\varphi_{780}$}
\end{ptree}
\noindent$\Pi_{6}$:
\begin{ptree}
\AXC{$\varphi_{105}$\lbl{A4}}
\AXC{$\forall x\,(\varphi_{73}) \to \varphi_{111}$\lbl{A3}}
\AXC{$\varphi_{73}$\lbl{A2}}
\UIC{$\forall x\,(\varphi_{73})$}
\AXC{$\forall x\,(\varphi_{73}) \to (\forall x\,(\varphi_{73}) \to \varphi_{111}) \to \varphi_{105} \to \varphi_{93}$\lbl{A1}}
\BIC{$(\forall x\,(\varphi_{73}) \to \varphi_{111}) \to \varphi_{105} \to \varphi_{93}$}
\BIC{$\varphi_{105} \to \varphi_{93}$}
\BIC{$\varphi_{93}$}
\end{ptree}
\noindent$\Pi_{7}$:
\begin{ptree}
\AXC{$\varphi_{106}$\lbl{A4}}
\AXC{$\forall s\,(\varphi_{74}) \to \varphi_{112}$\lbl{A3}}
\AXC{$\varphi_{74}$\lbl{A2}}
\UIC{$\forall s\,(\varphi_{74})$}
\AXC{$\forall s\,(\varphi_{74}) \to (\forall s\,(\varphi_{74}) \to \varphi_{112}) \to \varphi_{106} \to \varphi_{94}$\lbl{A1}}
\BIC{$(\forall s\,(\varphi_{74}) \to \varphi_{112}) \to \varphi_{106} \to \varphi_{94}$}
\BIC{$\varphi_{106} \to \varphi_{94}$}
\BIC{$\varphi_{94}$}
\end{ptree}
\noindent$\Pi_{8}$:
\begin{ptree}
\AXC{$\varphi_{175}$\lbl{A2}}
\AXC{$\mathrm{ContOn}(f,[a,b]) \to \varphi_{397}$\lbl{A1}}
\AXC{$\varphi_{729}$\lbl{A1}}
\BIC{$\varphi_{175} \to \mathrm{ContOn}(f,[a,b]) \to \varphi_{378}$}
\BIC{$\mathrm{ContOn}(f,[a,b]) \to \varphi_{378}$}
\AXC{$(\mathrm{ContOn}(f,[a,b]) \to \varphi_{378}) \to \varphi_{808}$\lbl{A1}}
\BIC{$\varphi_{808}$}
\end{ptree}
\noindent$\Pi_{9}$:
\begin{ptree}
\AXC{$\varphi_{709}$\lbl{A1}}
\AXC{$\varphi_{709} \to \varphi_{576}$\lbl{A1}}
\BIC{$\varphi_{576}$}
\AXC{$\varphi_{709}$\lbl{A1}}
\AXC{$\varphi_{709} \to \varphi_{578}$\lbl{A1}}
\BIC{$\varphi_{578}$}
\AXC{$\forall e \in \mathbb{R}\,(0 < e \to e \cdot 2^{-1} < e)$\lbl{[E39]}}
\AXC{$\varphi_{138}$\lbl{A2}}
\BIC{$\delta \in \mathbb{R} \to 0 < \delta \to \delta \cdot 2^{-1} < \delta$}
\AXC{$\varphi_{856}$\lbl{A1}}
\BIC{$\varphi_{578} \to \varphi_{576} \to \mathrm{ContOn}(f,[a,b]) \to \varphi_{633}$}
\BIC{$\varphi_{576} \to \mathrm{ContOn}(f,[a,b]) \to \varphi_{633}$}
\BIC{$\mathrm{ContOn}(f,[a,b]) \to \varphi_{633}$}
\end{ptree}
\noindent$\Pi_{10}$:
\begin{ptree}
\AXC{$\forall x\,\forall y\,\forall z\,(x \leq y \to y < z \to x < z)$\lbl{[O4]}}
\AXC{$\varphi_{21}$\lbl{A2}}
\BIC{$\forall y\,\forall z\,(d \leq y \to y < z \to d < z)$}
\AXC{$\varphi_{75}$\lbl{A2}}
\BIC{$\forall z\,(d \leq \delta \cdot 2^{-1} \to \delta \cdot 2^{-1} < z \to d < z)$}
\AXC{$\varphi_{213}$\lbl{A2}}
\BIC{$d \leq \delta \cdot 2^{-1} \to \delta \cdot 2^{-1} < \delta \to d < \delta$}
\AXC{$\varphi_{862}$\lbl{A1}}
\BIC{$(\mathrm{ContOn}(f,[a,b]) \to \varphi_{619}) \to \varphi_{787}$}
\end{ptree}
\noindent$\Pi_{11}$:
\begin{ptree}
\AXC{$\mathrm{ContOn}(f,[a,b]) \to \varphi_{678}$\lbl{A1}}
\AXC{$\varphi_{814}$\lbl{A1}}
\BIC{$\mathrm{ContOn}(f,[a,b]) \to \varphi_{619}$}
\AXC{$\stackrel{\Pi_{10}}{\vdots}$}\noLine
\UIC{$(\mathrm{ContOn}(f,[a,b]) \to \varphi_{619}) \to \varphi_{787}$}
\BIC{$\varphi_{787}$}
\end{ptree}
\noindent$\Pi_{12}$:
\begin{ptree}
\AXC{$\varphi_{175}$\lbl{A2}}
\AXC{$\mathrm{ContOn}(f,[a,b]) \to \varphi_{397}$\lbl{A1}}
\AXC{$\varphi_{729}$\lbl{A1}}
\BIC{$\varphi_{175} \to \mathrm{ContOn}(f,[a,b]) \to \varphi_{378}$}
\BIC{$\mathrm{ContOn}(f,[a,b]) \to \varphi_{378}$}
\AXC{$(\mathrm{ContOn}(f,[a,b]) \to \varphi_{378}) \to \varphi_{799}$\lbl{A1}}
\BIC{$\varphi_{799}$}
\end{ptree}
\noindent$\Pi_{13}$:
\begin{ptree}
\AXC{$\mathrm{ContOn}(f,[a,b]) \to \varphi_{631}$\lbl{A1}}
\AXC{$\stackrel{\Pi_{12}}{\vdots}$}\noLine
\UIC{$\varphi_{799}$}
\BIC{$\mathrm{ContOn}(f,[a,b]) \to \varphi_{652}$}
\end{ptree}
\noindent$\Pi_{14}$:
\begin{ptree}
\AXC{$\forall a\,\forall b\,\forall x\,(x \in [a,b] \leftrightarrow x \in \mathbb{R} \land a \leq x \land x \leq b)$\lbl{Def{$[a,b]$}}}
\AXC{$\varphi_{236}$\lbl{A2}}
\BIC{$\forall b\,\forall x\,(x \in [a,b] \leftrightarrow x \in \mathbb{R} \land a \leq x \land x \leq b)$}
\AXC{$\varphi_{195}$\lbl{A2}}
\BIC{$\forall x\,(x \in [a,b] \leftrightarrow x \in \mathbb{R} \land a \leq x \land x \leq b)$}
\AXC{$\varphi_{139}$\lbl{A2}}
\BIC{$s \in [a,b] \leftrightarrow s \in \mathbb{R} \land a \leq s \land s \leq b$}
\end{ptree}
\noindent$\Pi_{15}$:
\begin{ptree}
\AXC{$\stackrel{\Pi_{13}}{\vdots}$}\noLine
\UIC{$\mathrm{ContOn}(f,[a,b]) \to \varphi_{652}$}
\AXC{$\stackrel{\Pi_{14}}{\vdots}$}\noLine
\UIC{$s \in [a,b] \leftrightarrow s \in \mathbb{R} \land a \leq s \land s \leq b$}
\AXC{$(s \in [a,b] \leftrightarrow s \in \mathbb{R} \land a \leq s \land s \leq b) \to \varphi_{800}$\lbl{A1}}
\BIC{$\varphi_{800}$}
\BIC{$\mathrm{ContOn}(f,[a,b]) \to \varphi_{639}$}
\end{ptree}
\noindent$\Pi_{16}$:
\begin{ptree}
\AXC{$\stackrel{\Pi_{3}}{\vdots}$}\noLine
\UIC{$\mathrm{ContOn}(f,[a,b]) \to \varphi_{328}$}
\AXC{$\forall x\,\forall y\,(x \in \mathbb{R} \to y \in \mathbb{R} \to x - y \in \mathbb{R})$\lbl{[F1]}}
\AXC{$\varphi_{227}$\lbl{A2}}
\BIC{$\forall y\,(\xi \in \mathbb{R} \to y \in \mathbb{R} \to \xi - y \in \mathbb{R})$}
\AXC{$\varphi_{196}$\lbl{A2}}
\BIC{$\xi \in \mathbb{R} \to d \in \mathbb{R} \to \xi - d \in \mathbb{R}$}
\AXC{$\varphi_{853}$\lbl{A1}}
\BIC{$(\mathrm{ContOn}(f,[a,b]) \to \varphi_{328}) \to \varphi_{805}$}
\BIC{$\varphi_{805}$}
\end{ptree}
\noindent$\Pi_{17}$:
\begin{ptree}
\AXC{$\forall x\,\forall y\,(x \in \mathbb{R} \to y \in \mathbb{R} \to \neg y \leq x \to x < y)$\lbl{[A2]}}
\AXC{$\varphi_{277}$\lbl{A2}}
\BIC{$\forall y\,(\xi - d \in \mathbb{R} \to y \in \mathbb{R} \to \neg y \leq \xi - d \to \xi - d < y)$}
\AXC{$\varphi_{291}$\lbl{A2}}
\BIC{$\xi - d \in \mathbb{R} \to s \in \mathbb{R} \to \neg s \leq \xi - d \to \xi - d < s$}
\end{ptree}
\noindent$\Pi_{18}$:
\begin{ptree}
\AXC{$\stackrel{\Pi_{17}}{\vdots}$}\noLine
\UIC{$\xi - d \in \mathbb{R} \to s \in \mathbb{R} \to \neg s \leq \xi - d \to \xi - d < s$}
\AXC{$(\xi - d \in \mathbb{R} \to s \in \mathbb{R} \to \neg s \leq \xi - d \to \xi - d < s) \to \varphi_{881}$\lbl{A1}}
\BIC{$\varphi_{881}$}
\end{ptree}
\noindent$\Pi_{19}$:
\begin{ptree}
\AXC{$\stackrel{\Pi_{15}}{\vdots}$}\noLine
\UIC{$\mathrm{ContOn}(f,[a,b]) \to \varphi_{639}$}
\AXC{$\mathrm{ContOn}(f,[a,b]) \to \varphi_{678}$\lbl{A1}}
\AXC{$\varphi_{811}$\lbl{A1}}
\BIC{$\mathrm{ContOn}(f,[a,b]) \to \varphi_{610}$}
\AXC{$\stackrel{\Pi_{16}}{\vdots}$}\noLine
\UIC{$\varphi_{805}$}
\BIC{$\mathrm{ContOn}(f,[a,b]) \to \varphi_{670}$}
\AXC{$\stackrel{\Pi_{18}}{\vdots}$}\noLine
\UIC{$\varphi_{881}$}
\BIC{$(\mathrm{ContOn}(f,[a,b]) \to \varphi_{639}) \to \varphi_{818}$}
\BIC{$\varphi_{818}$}
\end{ptree}
\noindent$\Pi_{20}$:
\begin{ptree}
\AXC{$\forall a\,\forall X\,(\mathrm{UB}(a,X) \leftrightarrow a \in \mathbb{R} \land \forall x \in X\,(x \leq a))$\lbl{DefUB}}
\AXC{$\varphi_{205}$\lbl{A2}}
\BIC{$\forall X\,(\mathrm{UB}(\xi,X) \leftrightarrow \xi \in \mathbb{R} \land \forall x \in X\,(x \leq \xi))$}
\AXC{$\varphi_{174}$\lbl{A2}}
\BIC{$\mathrm{UB}(\xi,W) \leftrightarrow \xi \in \mathbb{R} \land \forall x \in W\,(x \leq \xi)$}
\end{ptree}
\noindent$\Pi_{21}$:
\begin{ptree}
\AXC{$\mathrm{ContOn}(f,[a,b]) \to \varphi_{327}$\lbl{A1}}
\AXC{$(\mathrm{ContOn}(f,[a,b]) \to \varphi_{327}) \to \varphi_{5} \to \varphi_{473}$\lbl{A1}}
\BIC{$\varphi_{5} \to \varphi_{473}$}
\end{ptree}
\noindent$\Pi_{22}$:
\begin{ptree}
\AXC{$\stackrel{\Pi_{20}}{\vdots}$}\noLine
\UIC{$\mathrm{UB}(\xi,W) \leftrightarrow \xi \in \mathbb{R} \land \forall x \in W\,(x \leq \xi)$}
\AXC{$\varphi_{7}$\lbl{DefSup}}
\AXC{$\varphi_{7} \to \varphi_{8}$\lbl{A2}}
\BIC{$\varphi_{8}$}
\AXC{$\varphi_{8} \to \varphi_{5}$\lbl{A2}}
\BIC{$\varphi_{5}$}
\AXC{$\stackrel{\Pi_{21}}{\vdots}$}\noLine
\UIC{$\varphi_{5} \to \varphi_{473}$}
\BIC{$\varphi_{473}$}
\BIC{$\mathrm{ContOn}(f,[a,b]) \to \varphi_{342}$}
\end{ptree}
\noindent$\Pi_{23}$:
\begin{ptree}
\AXC{$\forall x \in W\,(x \leq \xi) \to s \in W \to s \leq \xi$\lbl{A2}}
\AXC{$\stackrel{\Pi_{22}}{\vdots}$}\noLine
\UIC{$\mathrm{ContOn}(f,[a,b]) \to \varphi_{342}$}
\AXC{$(\mathrm{ContOn}(f,[a,b]) \to \varphi_{342}) \to \varphi_{454}$\lbl{A1}}
\BIC{$\varphi_{454}$}
\BIC{$\mathrm{ContOn}(f,[a,b]) \to \varphi_{335}$}
\end{ptree}
\noindent$\Pi_{24}$:
\begin{ptree}
\AXC{$\mathrm{ContOn}(f,[a,b]) \to \varphi_{631}$\lbl{A1}}
\AXC{$\stackrel{\Pi_{23}}{\vdots}$}\noLine
\UIC{$\mathrm{ContOn}(f,[a,b]) \to \varphi_{335}$}
\AXC{$(\mathrm{ContOn}(f,[a,b]) \to \varphi_{335}) \to \varphi_{794}$\lbl{A1}}
\BIC{$\varphi_{794}$}
\BIC{$\mathrm{ContOn}(f,[a,b]) \to \varphi_{620}$}
\end{ptree}
\noindent$\Pi_{25}$:
\begin{ptree}
\AXC{$\forall s\,(\varphi_{104})$\lbl{[E41]}}
\AXC{$\forall s\,(\varphi_{104}) \to \varphi_{104}$\lbl{A2}}
\BIC{$\varphi_{104}$}
\AXC{$\varphi_{104} \to \forall d\,(\varphi_{95})$\lbl{A2}}
\BIC{$\forall d\,(\varphi_{95})$}
\AXC{$\forall d\,(\varphi_{95}) \to \varphi_{95}$\lbl{A2}}
\BIC{$\varphi_{95}$}
\AXC{$\varphi_{95} \to \varphi_{893}$\lbl{A1}}
\BIC{$\varphi_{893}$}
\end{ptree}
\noindent$\Pi_{26}$:
\begin{ptree}
\AXC{$\stackrel{\Pi_{3}}{\vdots}$}\noLine
\UIC{$\mathrm{ContOn}(f,[a,b]) \to \varphi_{328}$}
\AXC{$\stackrel{\Pi_{15}}{\vdots}$}\noLine
\UIC{$\mathrm{ContOn}(f,[a,b]) \to \varphi_{639}$}
\AXC{$\stackrel{\Pi_{25}}{\vdots}$}\noLine
\UIC{$\varphi_{893}$}
\BIC{$(\mathrm{ContOn}(f,[a,b]) \to \varphi_{328}) \to \varphi_{879}$}
\BIC{$\varphi_{879}$}
\end{ptree}
\noindent$\Pi_{27}$:
\begin{ptree}
\AXC{$\stackrel{\Pi_{24}}{\vdots}$}\noLine
\UIC{$\mathrm{ContOn}(f,[a,b]) \to \varphi_{620}$}
\AXC{$\mathrm{ContOn}(f,[a,b]) \to \varphi_{678}$\lbl{A1}}
\AXC{$\varphi_{811}$\lbl{A1}}
\BIC{$\mathrm{ContOn}(f,[a,b]) \to \varphi_{610}$}
\AXC{$\stackrel{\Pi_{26}}{\vdots}$}\noLine
\UIC{$\varphi_{879}$}
\BIC{$(\mathrm{ContOn}(f,[a,b]) \to \varphi_{620}) \to \varphi_{820}$}
\BIC{$\varphi_{820}$}
\end{ptree}
\noindent$\Pi_{28}$:
\begin{ptree}
\AXC{$\forall x\,\forall y\,\forall z\,(x < y \to y < z \to x < z)$\lbl{[O2]}}
\AXC{$\varphi_{162}$\lbl{A2}}
\BIC{$\forall y\,\forall z\,(\lvert s - \xi \rvert < y \to y < z \to \lvert s - \xi \rvert < z)$}
\AXC{$\varphi_{264}$\lbl{A2}}
\BIC{$\forall z\,(\lvert s - \xi \rvert < d \to d < z \to \lvert s - \xi \rvert < z)$}
\AXC{$\varphi_{245}$\lbl{A2}}
\BIC{$\lvert s - \xi \rvert < d \to d < \delta \to \lvert s - \xi \rvert < \delta$}
\AXC{$\varphi_{866}$\lbl{A1}}
\BIC{$(\mathrm{ContOn}(f,[a,b]) \to \varphi_{672}) \to \varphi_{804}$}
\end{ptree}
\noindent$\Pi_{29}$:
\begin{ptree}
\AXC{$\mathrm{ContOn}(f,[a,b]) \to \varphi_{662}$\lbl{A1}}
\AXC{$\stackrel{\Pi_{19}}{\vdots}$}\noLine
\UIC{$\varphi_{818}$}
\BIC{$\mathrm{ContOn}(f,[a,b]) \to \varphi_{663}$}
\AXC{$\stackrel{\Pi_{27}}{\vdots}$}\noLine
\UIC{$\varphi_{820}$}
\BIC{$\mathrm{ContOn}(f,[a,b]) \to \varphi_{672}$}
\AXC{$\stackrel{\Pi_{28}}{\vdots}$}\noLine
\UIC{$(\mathrm{ContOn}(f,[a,b]) \to \varphi_{672}) \to \varphi_{804}$}
\BIC{$\varphi_{804}$}
\end{ptree}
\noindent$\Pi_{30}$:
\begin{ptree}
\AXC{$\varphi_{157} \to \varphi_{131}$\lbl{A2}}
\AXC{$\varphi_{709}$\lbl{A1}}
\AXC{$\varphi_{709} \to \mathrm{ContOn}(f,[a,b]) \to \varphi_{726}$\lbl{A1}}
\BIC{$\mathrm{ContOn}(f,[a,b]) \to \varphi_{726}$}
\AXC{$\varphi_{860}$\lbl{A1}}
\BIC{$(\varphi_{157} \to \varphi_{131}) \to \mathrm{ContOn}(f,[a,b]) \to \varphi_{725}$}
\BIC{$\mathrm{ContOn}(f,[a,b]) \to \varphi_{725}$}
\AXC{$\varphi_{884}$\lbl{A1}}
\BIC{$(\mathrm{ContOn}(f,[a,b]) \to \varphi_{652}) \to \varphi_{830}$}
\end{ptree}
\noindent$\Pi_{31}$:
\begin{ptree}
\AXC{$\stackrel{\Pi_{13}}{\vdots}$}\noLine
\UIC{$\mathrm{ContOn}(f,[a,b]) \to \varphi_{652}$}
\AXC{$\stackrel{\Pi_{30}}{\vdots}$}\noLine
\UIC{$(\mathrm{ContOn}(f,[a,b]) \to \varphi_{652}) \to \varphi_{830}$}
\BIC{$\varphi_{830}$}
\end{ptree}
\noindent$\Pi_{32}$:
\begin{ptree}
\AXC{$\forall a\,\forall b\,\forall x\,(x \in [a,b] \leftrightarrow x \in \mathbb{R} \land a \leq x \land x \leq b)$\lbl{Def{$[a,b]$}}}
\AXC{$\varphi_{236}$\lbl{A2}}
\BIC{$\forall b\,\forall x\,(x \in [a,b] \leftrightarrow x \in \mathbb{R} \land a \leq x \land x \leq b)$}
\AXC{$\varphi_{195}$\lbl{A2}}
\BIC{$\forall x\,(x \in [a,b] \leftrightarrow x \in \mathbb{R} \land a \leq x \land x \leq b)$}
\AXC{$\varphi_{140}$\lbl{A2}}
\BIC{$x \in [a,b] \leftrightarrow x \in \mathbb{R} \land a \leq x \land x \leq b$}
\end{ptree}
\noindent$\Pi_{33}$:
\begin{ptree}
\AXC{$\varphi_{367}$\lbl{A1}}
\AXC{$\varphi_{176}$\lbl{A2}}
\AXC{$\mathrm{ContOn}(f,[a,b]) \to \varphi_{397}$\lbl{A1}}
\AXC{$\varphi_{730}$\lbl{A1}}
\BIC{$\varphi_{176} \to \mathrm{ContOn}(f,[a,b]) \to \varphi_{379}$}
\BIC{$\mathrm{ContOn}(f,[a,b]) \to \varphi_{379}$}
\AXC{$\varphi_{754}$\lbl{A1}}
\BIC{$\varphi_{367} \to (x \in [a,b] \leftrightarrow x \in \mathbb{R} \land a \leq x \land x \leq b) \to \varphi_{360}$}
\BIC{$(x \in [a,b] \leftrightarrow x \in \mathbb{R} \land a \leq x \land x \leq b) \to \varphi_{360}$}
\end{ptree}
\noindent$\Pi_{34}$:
\begin{ptree}
\AXC{$\stackrel{\Pi_{32}}{\vdots}$}\noLine
\UIC{$x \in [a,b] \leftrightarrow x \in \mathbb{R} \land a \leq x \land x \leq b$}
\AXC{$\stackrel{\Pi_{33}}{\vdots}$}\noLine
\UIC{$(x \in [a,b] \leftrightarrow x \in \mathbb{R} \land a \leq x \land x \leq b) \to \varphi_{360}$}
\BIC{$\varphi_{360}$}
\AXC{$\varphi_{360} \to \varphi_{318} \to x \in W \to x \leq b$\lbl{A1}}
\BIC{$\varphi_{318} \to x \in W \to x \leq b$}
\UIC{$\forall x\,(\varphi_{318} \to x \in W \to x \leq b)$}
\end{ptree}
\noindent$\Pi_{35}$:
\begin{ptree}
\AXC{$\forall x\,(\varphi_{318} \to x \in W \to x \leq b) \to \forall x\,(\varphi_{318}) \to \forall x \in W\,(x \leq b)$\lbl{A3}}
\AXC{$\stackrel{\Pi_{34}}{\vdots}$}\noLine
\UIC{$\forall x\,(\varphi_{318} \to x \in W \to x \leq b)$}
\AXC{$\forall x\,(\varphi_{318} \to x \in W \to x \leq b) \to \varphi_{831}$\lbl{A1}}
\BIC{$\varphi_{831}$}
\BIC{$(\varphi_{318} \to \forall x\,(\varphi_{318})) \to \mathrm{ContOn}(f,[a,b]) \to \varphi_{311}$}
\end{ptree}
\noindent$\Pi_{36}$:
\begin{ptree}
\AXC{$\varphi_{318} \to \forall x\,(\varphi_{318})$\lbl{A4}}
\AXC{$\stackrel{\Pi_{35}}{\vdots}$}\noLine
\UIC{$(\varphi_{318} \to \forall x\,(\varphi_{318})) \to \mathrm{ContOn}(f,[a,b]) \to \varphi_{311}$}
\BIC{$\mathrm{ContOn}(f,[a,b]) \to \varphi_{311}$}
\end{ptree}
\noindent$\Pi_{37}$:
\begin{ptree}
\AXC{$\forall a\,\forall X\,(\mathrm{UB}(a,X) \leftrightarrow a \in \mathbb{R} \land \forall x \in X\,(x \leq a))$\lbl{DefUB}}
\AXC{$\varphi_{177}$\lbl{A2}}
\BIC{$\forall X\,(\mathrm{UB}(b,X) \leftrightarrow b \in \mathbb{R} \land \forall x \in X\,(x \leq b))$}
\AXC{$\varphi_{126}$\lbl{A2}}
\BIC{$\mathrm{UB}(b,W) \leftrightarrow b \in \mathbb{R} \land \forall x \in W\,(x \leq b)$}
\end{ptree}
\noindent$\Pi_{38}$:
\begin{ptree}
\AXC{$\stackrel{\Pi_{37}}{\vdots}$}\noLine
\UIC{$\mathrm{UB}(b,W) \leftrightarrow b \in \mathbb{R} \land \forall x \in W\,(x \leq b)$}
\AXC{$(\mathrm{UB}(b,W) \leftrightarrow b \in \mathbb{R} \land \forall x \in W\,(x \leq b)) \to \varphi_{691}$\lbl{A1}}
\BIC{$\varphi_{691}$}
\end{ptree}
\noindent$\Pi_{39}$:
\begin{ptree}
\AXC{$\stackrel{\Pi_{36}}{\vdots}$}\noLine
\UIC{$\mathrm{ContOn}(f,[a,b]) \to \varphi_{311}$}
\AXC{$\mathrm{ContOn}(f,[a,b]) \to a \in \mathbb{R} \to b \in \mathbb{R} \to b \in \mathbb{R}$\lbl{A1}}
\AXC{$\stackrel{\Pi_{38}}{\vdots}$}\noLine
\UIC{$\varphi_{691}$}
\BIC{$(\mathrm{ContOn}(f,[a,b]) \to \varphi_{311}) \to \varphi_{338}$}
\BIC{$\varphi_{338}$}
\end{ptree}
\noindent$\Pi_{40}$:
\begin{ptree}
\AXC{$\varphi_{7}$\lbl{DefSup}}
\AXC{$\varphi_{7} \to \varphi_{8}$\lbl{A2}}
\BIC{$\varphi_{8}$}
\AXC{$\varphi_{8} \to \varphi_{5}$\lbl{A2}}
\BIC{$\varphi_{5}$}
\AXC{$\mathrm{ContOn}(f,[a,b]) \to \varphi_{327}$\lbl{A1}}
\AXC{$\varphi_{710}$\lbl{A1}}
\BIC{$\varphi_{5} \to \mathrm{ContOn}(f,[a,b]) \to \varphi_{343}$}
\BIC{$\mathrm{ContOn}(f,[a,b]) \to \varphi_{343}$}
\end{ptree}
\noindent$\Pi_{41}$:
\begin{ptree}
\AXC{$\stackrel{\Pi_{40}}{\vdots}$}\noLine
\UIC{$\mathrm{ContOn}(f,[a,b]) \to \varphi_{343}$}
\AXC{$(\mathrm{ContOn}(f,[a,b]) \to \varphi_{343}) \to \varphi_{455}$\lbl{A1}}
\BIC{$\varphi_{455}$}
\end{ptree}
\noindent$\Pi_{42}$:
\begin{ptree}
\AXC{$\stackrel{\Pi_{39}}{\vdots}$}\noLine
\UIC{$\varphi_{338}$}
\AXC{$\forall w\,(\mathrm{UB}(w,W) \to \xi \leq w) \to \mathrm{UB}(b,W) \to \xi \leq b$\lbl{A2}}
\AXC{$\stackrel{\Pi_{41}}{\vdots}$}\noLine
\UIC{$\varphi_{455}$}
\BIC{$\mathrm{ContOn}(f,[a,b]) \to \varphi_{336}$}
\AXC{$\varphi_{735}$\lbl{A1}}
\BIC{$\varphi_{338} \to \mathrm{ContOn}(f,[a,b]) \to \varphi_{316}$}
\BIC{$\mathrm{ContOn}(f,[a,b]) \to \varphi_{316}$}
\end{ptree}
\noindent$\Pi_{43}$:
\begin{ptree}
\AXC{$x = x$\lbl{A5}}
\UIC{$\forall x\,(x = x)$}
\AXC{$\forall x\,(x = x) \to a = a$\lbl{A2}}
\BIC{$a = a$}
\AXC{$a = a \to \mathrm{ContOn}(f,[a,b]) \to a \in \mathbb{R} \to \varphi_{62}$\lbl{A1}}
\BIC{$\mathrm{ContOn}(f,[a,b]) \to a \in \mathbb{R} \to \varphi_{62}$}
\end{ptree}
\noindent$\Pi_{44}$:
\begin{ptree}
\AXC{$\forall a\,\forall b\,\forall x\,(x \in [a,b] \leftrightarrow x \in \mathbb{R} \land a \leq x \land x \leq b)$\lbl{Def{$[a,b]$}}}
\AXC{$\varphi_{236}$\lbl{A2}}
\BIC{$\forall b\,\forall x\,(x \in [a,b] \leftrightarrow x \in \mathbb{R} \land a \leq x \land x \leq b)$}
\AXC{$\varphi_{195}$\lbl{A2}}
\BIC{$\forall x\,(x \in [a,b] \leftrightarrow x \in \mathbb{R} \land a \leq x \land x \leq b)$}
\AXC{$\varphi_{141}$\lbl{A2}}
\BIC{$a \in [a,b] \leftrightarrow a \in \mathbb{R} \land a \leq a \land a \leq b$}
\end{ptree}
\noindent$\Pi_{45}$:
\begin{ptree}
\AXC{$\stackrel{\Pi_{44}}{\vdots}$}\noLine
\UIC{$a \in [a,b] \leftrightarrow a \in \mathbb{R} \land a \leq a \land a \leq b$}
\AXC{$(a \in [a,b] \leftrightarrow a \in \mathbb{R} \land a \leq a \land a \leq b) \to \varphi_{688}$\lbl{A1}}
\BIC{$\varphi_{688}$}
\end{ptree}
\noindent$\Pi_{46}$:
\begin{ptree}
\AXC{$\stackrel{\Pi_{43}}{\vdots}$}\noLine
\UIC{$\mathrm{ContOn}(f,[a,b]) \to a \in \mathbb{R} \to \varphi_{62}$}
\AXC{$\mathrm{ContOn}(f,[a,b]) \to a \in \mathbb{R} \to a \in \mathbb{R}$\lbl{A1}}
\AXC{$\stackrel{\Pi_{45}}{\vdots}$}\noLine
\UIC{$\varphi_{688}$}
\BIC{$(\mathrm{ContOn}(f,[a,b]) \to a \in \mathbb{R} \to \varphi_{62}) \to \varphi_{474}$}
\BIC{$\varphi_{474}$}
\end{ptree}
\noindent$\Pi_{47}$:
\begin{ptree}
\AXC{$\varphi_{178}$\lbl{A2}}
\AXC{$\mathrm{ContOn}(f,[a,b]) \to \varphi_{397}$\lbl{A1}}
\AXC{$\varphi_{731}$\lbl{A1}}
\BIC{$\varphi_{178} \to \mathrm{ContOn}(f,[a,b]) \to \varphi_{380}$}
\BIC{$\mathrm{ContOn}(f,[a,b]) \to \varphi_{380}$}
\AXC{$(\mathrm{ContOn}(f,[a,b]) \to \varphi_{380}) \to \varphi_{687}$\lbl{A1}}
\BIC{$\varphi_{687}$}
\end{ptree}
\noindent$\Pi_{48}$:
\begin{ptree}
\AXC{$\mathrm{ContOn}(f,[a,b]) \to \varphi_{102}$\lbl{A1}}
\AXC{$\varphi_{85}$\lbl{A1}}
\AXC{$\varphi_{85} \to \mathrm{ContOn}(f,[a,b]) \to a \in \mathbb{R} \to \varphi_{61}$\lbl{A1}}
\BIC{$\mathrm{ContOn}(f,[a,b]) \to a \in \mathbb{R} \to \varphi_{61}$}
\AXC{$\stackrel{\Pi_{46}}{\vdots}$}\noLine
\UIC{$\varphi_{474}$}
\BIC{$\mathrm{ContOn}(f,[a,b]) \to a \in \mathbb{R} \to \varphi_{135}$}
\AXC{$\stackrel{\Pi_{47}}{\vdots}$}\noLine
\UIC{$\varphi_{687}$}
\BIC{$(\mathrm{ContOn}(f,[a,b]) \to \varphi_{102}) \to \varphi_{339}$}
\BIC{$\varphi_{339}$}
\end{ptree}
\noindent$\Pi_{49}$:
\begin{ptree}
\AXC{$\forall x \in W\,(x \leq \xi) \to a \in W \to a \leq \xi$\lbl{A2}}
\AXC{$\stackrel{\Pi_{22}}{\vdots}$}\noLine
\UIC{$\mathrm{ContOn}(f,[a,b]) \to \varphi_{342}$}
\AXC{$(\mathrm{ContOn}(f,[a,b]) \to \varphi_{342}) \to \varphi_{456}$\lbl{A1}}
\BIC{$\varphi_{456}$}
\BIC{$\mathrm{ContOn}(f,[a,b]) \to \varphi_{337}$}
\end{ptree}
\noindent$\Pi_{50}$:
\begin{ptree}
\AXC{$\stackrel{\Pi_{48}}{\vdots}$}\noLine
\UIC{$\varphi_{339}$}
\AXC{$\stackrel{\Pi_{49}}{\vdots}$}\noLine
\UIC{$\mathrm{ContOn}(f,[a,b]) \to \varphi_{337}$}
\AXC{$\varphi_{736}$\lbl{A1}}
\BIC{$\varphi_{339} \to \mathrm{ContOn}(f,[a,b]) \to \varphi_{317}$}
\BIC{$\mathrm{ContOn}(f,[a,b]) \to \varphi_{317}$}
\end{ptree}
\noindent$\Pi_{51}$:
\begin{ptree}
\AXC{$\forall a\,\forall b\,\forall x\,(x \in [a,b] \leftrightarrow x \in \mathbb{R} \land a \leq x \land x \leq b)$\lbl{Def{$[a,b]$}}}
\AXC{$\varphi_{236}$\lbl{A2}}
\BIC{$\forall b\,\forall x\,(x \in [a,b] \leftrightarrow x \in \mathbb{R} \land a \leq x \land x \leq b)$}
\AXC{$\varphi_{195}$\lbl{A2}}
\BIC{$\forall x\,(x \in [a,b] \leftrightarrow x \in \mathbb{R} \land a \leq x \land x \leq b)$}
\AXC{$\varphi_{179}$\lbl{A2}}
\BIC{$\xi \in [a,b] \leftrightarrow \xi \in \mathbb{R} \land a \leq \xi \land \xi \leq b$}
\end{ptree}
\noindent$\Pi_{52}$:
\begin{ptree}
\AXC{$\stackrel{\Pi_{3}}{\vdots}$}\noLine
\UIC{$\mathrm{ContOn}(f,[a,b]) \to \varphi_{328}$}
\AXC{$\stackrel{\Pi_{51}}{\vdots}$}\noLine
\UIC{$\xi \in [a,b] \leftrightarrow \xi \in \mathbb{R} \land a \leq \xi \land \xi \leq b$}
\AXC{$(\xi \in [a,b] \leftrightarrow \xi \in \mathbb{R} \land a \leq \xi \land \xi \leq b) \to \varphi_{771}$\lbl{A1}}
\BIC{$\varphi_{771}$}
\BIC{$(\mathrm{ContOn}(f,[a,b]) \to \varphi_{317}) \to \varphi_{657}$}
\end{ptree}
\noindent$\Pi_{53}$:
\begin{ptree}
\AXC{$\stackrel{\Pi_{42}}{\vdots}$}\noLine
\UIC{$\mathrm{ContOn}(f,[a,b]) \to \varphi_{316}$}
\AXC{$\stackrel{\Pi_{50}}{\vdots}$}\noLine
\UIC{$\mathrm{ContOn}(f,[a,b]) \to \varphi_{317}$}
\AXC{$\stackrel{\Pi_{52}}{\vdots}$}\noLine
\UIC{$(\mathrm{ContOn}(f,[a,b]) \to \varphi_{317}) \to \varphi_{657}$}
\BIC{$\varphi_{657}$}
\BIC{$\mathrm{ContOn}(f,[a,b]) \to \varphi_{341}$}
\end{ptree}
\noindent$\Pi_{54}$:
\begin{ptree}
\AXC{$\forall f\,(\varphi_{40})$\lbl{DefContOn}}
\AXC{$\forall f\,(\varphi_{40}) \to \varphi_{40}$\lbl{A2}}
\BIC{$\varphi_{40}$}
\AXC{$\varphi_{40} \to \varphi_{231}$\lbl{A2}}
\BIC{$\varphi_{231}$}
\AXC{$\varphi_{4}$\lbl{A1}}
\AXC{$\varphi_{545}$\lbl{A1}}
\BIC{$\varphi_{231} \to \mathrm{ContOn}(f,[a,b]) \to a \in \mathbb{R} \to b \in \mathbb{R} \to \varphi_{92}$}
\BIC{$\mathrm{ContOn}(f,[a,b]) \to a \in \mathbb{R} \to b \in \mathbb{R} \to \varphi_{92}$}
\end{ptree}
\noindent$\Pi_{55}$:
\begin{ptree}
\AXC{$\forall X\,\forall Y\,\forall x\,(X \subseteq Y \to x \in X \to x \in Y)$\lbl{[B5]}}
\AXC{$\varphi_{180}$\lbl{A2}}
\BIC{$\forall Y\,\forall x\,([a,b] \subseteq Y \to x \in [a,b] \to x \in Y)$}
\AXC{$\forall Y\,\forall x\,([a,b] \subseteq Y \to x \in [a,b] \to x \in Y) \to \varphi_{15}$\lbl{A2}}
\BIC{$\varphi_{15}$}
\AXC{$\varphi_{288}$\lbl{A2}}
\BIC{$[a,b] \subseteq \mathrm{dom}(f) \to \xi \in [a,b] \to \xi \in \mathrm{dom}(f)$}
\end{ptree}
\noindent$\Pi_{56}$:
\begin{ptree}
\AXC{$\stackrel{\Pi_{54}}{\vdots}$}\noLine
\UIC{$\mathrm{ContOn}(f,[a,b]) \to a \in \mathbb{R} \to b \in \mathbb{R} \to \varphi_{92}$}
\AXC{$\stackrel{\Pi_{55}}{\vdots}$}\noLine
\UIC{$[a,b] \subseteq \mathrm{dom}(f) \to \xi \in [a,b] \to \xi \in \mathrm{dom}(f)$}
\AXC{$\varphi_{746}$\lbl{A1}}
\BIC{$\varphi_{733}$}
\BIC{$\varphi_{669}$}
\end{ptree}
\noindent$\Pi_{57}$:
\begin{ptree}
\AXC{$\forall f\,(\varphi_{40})$\lbl{DefContOn}}
\AXC{$\forall f\,(\varphi_{40}) \to \varphi_{40}$\lbl{A2}}
\BIC{$\varphi_{40}$}
\AXC{$\varphi_{40} \to \varphi_{231}$\lbl{A2}}
\BIC{$\varphi_{231}$}
\AXC{$\varphi_{4}$\lbl{A1}}
\AXC{$\varphi_{4} \to \varphi_{231} \to \mathrm{ContOn}(f,[a,b]) \to a \in \mathbb{R} \to \varphi_{80}$\lbl{A1}}
\BIC{$\varphi_{231} \to \mathrm{ContOn}(f,[a,b]) \to a \in \mathbb{R} \to \varphi_{80}$}
\BIC{$\mathrm{ContOn}(f,[a,b]) \to a \in \mathbb{R} \to \varphi_{80}$}
\end{ptree}
\noindent$\Pi_{58}$:
\begin{ptree}
\AXC{$\forall f\,\forall x\,(\mathrm{RFn}(f) \to x \in \mathrm{dom}(f) \to f(x) \in \mathbb{R})$\lbl{[S6]}}
\AXC{$\varphi_{214}$\lbl{A2}}
\BIC{$\forall x\,(\mathrm{RFn}(f) \to x \in \mathrm{dom}(f) \to f(x) \in \mathbb{R})$}
\AXC{$\varphi_{181}$\lbl{A2}}
\BIC{$\mathrm{RFn}(f) \to \xi \in \mathrm{dom}(f) \to f(\xi) \in \mathbb{R}$}
\AXC{$\varphi_{739}$\lbl{A1}}
\BIC{$(\mathrm{ContOn}(f,[a,b]) \to a \in \mathbb{R} \to \varphi_{80}) \to \varphi_{671}$}
\end{ptree}
\noindent$\Pi_{59}$:
\begin{ptree}
\AXC{$\stackrel{\Pi_{57}}{\vdots}$}\noLine
\UIC{$\mathrm{ContOn}(f,[a,b]) \to a \in \mathbb{R} \to \varphi_{80}$}
\AXC{$\stackrel{\Pi_{58}}{\vdots}$}\noLine
\UIC{$(\mathrm{ContOn}(f,[a,b]) \to a \in \mathbb{R} \to \varphi_{80}) \to \varphi_{671}$}
\BIC{$\varphi_{671}$}
\end{ptree}
\noindent$\Pi_{60}$:
\begin{ptree}
\AXC{$\stackrel{\Pi_{53}}{\vdots}$}\noLine
\UIC{$\mathrm{ContOn}(f,[a,b]) \to \varphi_{341}$}
\AXC{$\stackrel{\Pi_{56}}{\vdots}$}\noLine
\UIC{$\varphi_{669}$}
\BIC{$\mathrm{ContOn}(f,[a,b]) \to \varphi_{332}$}
\AXC{$\stackrel{\Pi_{59}}{\vdots}$}\noLine
\UIC{$\varphi_{671}$}
\BIC{$\mathrm{ContOn}(f,[a,b]) \to \varphi_{347}$}
\end{ptree}
\noindent$\Pi_{61}$:
\begin{ptree}
\AXC{$\forall x\,\forall y\,(x \in \mathbb{R} \to y \in \mathbb{R} \to x - y \in \mathbb{R})$\lbl{[F1]}}
\AXC{$\varphi_{283}$\lbl{A2}}
\BIC{$\forall y\,(f(\xi) \in \mathbb{R} \to y \in \mathbb{R} \to f(\xi) - y \in \mathbb{R})$}
\AXC{$\varphi_{296}$\lbl{A2}}
\BIC{$f(\xi) \in \mathbb{R} \to \mu \in \mathbb{R} \to f(\xi) - \mu \in \mathbb{R}$}
\AXC{$(f(\xi) \in \mathbb{R} \to \mu \in \mathbb{R} \to f(\xi) - \mu \in \mathbb{R}) \to \varphi_{763}$\lbl{A1}}
\BIC{$\varphi_{763}$}
\end{ptree}
\noindent$\Pi_{62}$:
\begin{ptree}
\AXC{$\forall X\,\forall Y\,\forall x\,(X \subseteq Y \to x \in X \to x \in Y)$\lbl{[B5]}}
\AXC{$\varphi_{180}$\lbl{A2}}
\BIC{$\forall Y\,\forall x\,([a,b] \subseteq Y \to x \in [a,b] \to x \in Y)$}
\AXC{$\forall Y\,\forall x\,([a,b] \subseteq Y \to x \in [a,b] \to x \in Y) \to \varphi_{15}$\lbl{A2}}
\BIC{$\varphi_{15}$}
\AXC{$\varphi_{285}$\lbl{A2}}
\BIC{$[a,b] \subseteq \mathrm{dom}(f) \to s \in [a,b] \to s \in \mathrm{dom}(f)$}
\end{ptree}
\noindent$\Pi_{63}$:
\begin{ptree}
\AXC{$\stackrel{\Pi_{54}}{\vdots}$}\noLine
\UIC{$\mathrm{ContOn}(f,[a,b]) \to a \in \mathbb{R} \to b \in \mathbb{R} \to \varphi_{92}$}
\AXC{$\stackrel{\Pi_{62}}{\vdots}$}\noLine
\UIC{$[a,b] \subseteq \mathrm{dom}(f) \to s \in [a,b] \to s \in \mathrm{dom}(f)$}
\AXC{$\varphi_{847}$\lbl{A1}}
\BIC{$\varphi_{840}$}
\BIC{$\varphi_{813}$}
\end{ptree}
\noindent$\Pi_{64}$:
\begin{ptree}
\AXC{$\forall f\,\forall x\,(\mathrm{RFn}(f) \to x \in \mathrm{dom}(f) \to f(x) \in \mathbb{R})$\lbl{[S6]}}
\AXC{$\varphi_{214}$\lbl{A2}}
\BIC{$\forall x\,(\mathrm{RFn}(f) \to x \in \mathrm{dom}(f) \to f(x) \in \mathbb{R})$}
\AXC{$\varphi_{163}$\lbl{A2}}
\BIC{$\mathrm{RFn}(f) \to s \in \mathrm{dom}(f) \to f(s) \in \mathbb{R}$}
\AXC{$\varphi_{844}$\lbl{A1}}
\BIC{$(\mathrm{ContOn}(f,[a,b]) \to a \in \mathbb{R} \to \varphi_{80}) \to \varphi_{821}$}
\end{ptree}
\noindent$\Pi_{65}$:
\begin{ptree}
\AXC{$\stackrel{\Pi_{57}}{\vdots}$}\noLine
\UIC{$\mathrm{ContOn}(f,[a,b]) \to a \in \mathbb{R} \to \varphi_{80}$}
\AXC{$\stackrel{\Pi_{64}}{\vdots}$}\noLine
\UIC{$(\mathrm{ContOn}(f,[a,b]) \to a \in \mathbb{R} \to \varphi_{80}) \to \varphi_{821}$}
\BIC{$\varphi_{821}$}
\end{ptree}
\noindent$\Pi_{66}$:
\begin{ptree}
\AXC{$\stackrel{\Pi_{13}}{\vdots}$}\noLine
\UIC{$\mathrm{ContOn}(f,[a,b]) \to \varphi_{652}$}
\AXC{$\stackrel{\Pi_{63}}{\vdots}$}\noLine
\UIC{$\varphi_{813}$}
\BIC{$\mathrm{ContOn}(f,[a,b]) \to \varphi_{664}$}
\AXC{$\stackrel{\Pi_{65}}{\vdots}$}\noLine
\UIC{$\varphi_{821}$}
\BIC{$\mathrm{ContOn}(f,[a,b]) \to \varphi_{673}$}
\end{ptree}
\noindent$\Pi_{67}$:
\begin{ptree}
\AXC{$\stackrel{\Pi_{66}}{\vdots}$}\noLine
\UIC{$\mathrm{ContOn}(f,[a,b]) \to \varphi_{673}$}
\AXC{$\forall x\,\forall y\,(x \in \mathbb{R} \to y \in \mathbb{R} \to x - y \in \mathbb{R})$\lbl{[F1]}}
\AXC{$\varphi_{278}$\lbl{A2}}
\BIC{$\forall y\,(f(s) \in \mathbb{R} \to y \in \mathbb{R} \to f(s) - y \in \mathbb{R})$}
\AXC{$\varphi_{301}$\lbl{A2}}
\BIC{$f(s) \in \mathbb{R} \to f(\xi) \in \mathbb{R} \to f(s) - f(\xi) \in \mathbb{R}$}
\AXC{$\varphi_{858}$\lbl{A1}}
\BIC{$(\mathrm{ContOn}(f,[a,b]) \to \varphi_{673}) \to \varphi_{756}$}
\BIC{$\varphi_{756}$}
\end{ptree}
\noindent$\Pi_{68}$:
\begin{ptree}
\AXC{$\forall x\,\forall y\,(x \in \mathbb{R} \to y \in \mathbb{R} \to \lvert x \rvert < y \to -y < x \land x < y)$\lbl{[A11]}}
\AXC{$\varphi_{448}$\lbl{A2}}
\BIC{$\varphi_{305}$}
\AXC{$\varphi_{305} \to \varphi_{443}$\lbl{A2}}
\BIC{$\varphi_{443}$}
\AXC{$\varphi_{443} \to \varphi_{885}$\lbl{A1}}
\BIC{$\varphi_{885}$}
\end{ptree}
\noindent$\Pi_{69}$:
\begin{ptree}
\AXC{$\varphi_{56}$\lbl{A1}}
\AXC{$\stackrel{\Pi_{60}}{\vdots}$}\noLine
\UIC{$\mathrm{ContOn}(f,[a,b]) \to \varphi_{347}$}
\AXC{$\stackrel{\Pi_{61}}{\vdots}$}\noLine
\UIC{$\varphi_{763}$}
\BIC{$\varphi_{56} \to \mathrm{ContOn}(f,[a,b]) \to \varphi_{677}$}
\BIC{$\mathrm{ContOn}(f,[a,b]) \to \varphi_{677}$}
\AXC{$\stackrel{\Pi_{60}}{\vdots}$}\noLine
\UIC{$\mathrm{ContOn}(f,[a,b]) \to \varphi_{347}$}
\AXC{$\stackrel{\Pi_{67}}{\vdots}$}\noLine
\UIC{$\varphi_{756}$}
\BIC{$\mathrm{ContOn}(f,[a,b]) \to \varphi_{683}$}
\AXC{$\stackrel{\Pi_{68}}{\vdots}$}\noLine
\UIC{$\varphi_{885}$}
\BIC{$(\mathrm{ContOn}(f,[a,b]) \to \varphi_{677}) \to \varphi_{838}$}
\BIC{$\varphi_{838}$}
\end{ptree}
\noindent$\Pi_{70}$:
\begin{ptree}
\AXC{$\stackrel{\Pi_{9}}{\vdots}$}\noLine
\UIC{$\mathrm{ContOn}(f,[a,b]) \to \varphi_{633}$}
\AXC{$\stackrel{\Pi_{11}}{\vdots}$}\noLine
\UIC{$\varphi_{787}$}
\BIC{$\mathrm{ContOn}(f,[a,b]) \to \varphi_{603}$}
\AXC{$\stackrel{\Pi_{29}}{\vdots}$}\noLine
\UIC{$\varphi_{804}$}
\BIC{$\mathrm{ContOn}(f,[a,b]) \to \varphi_{674}$}
\AXC{$\stackrel{\Pi_{31}}{\vdots}$}\noLine
\UIC{$\varphi_{830}$}
\BIC{$\mathrm{ContOn}(f,[a,b]) \to \varphi_{696}$}
\AXC{$\stackrel{\Pi_{69}}{\vdots}$}\noLine
\UIC{$\varphi_{838}$}
\BIC{$\mathrm{ContOn}(f,[a,b]) \to \varphi_{727}$}
\AXC{$\varphi_{839}$\lbl{A1}}
\BIC{$\mathrm{ContOn}(f,[a,b]) \to \varphi_{697}$}
\end{ptree}
\noindent$\Pi_{71}$:
\begin{ptree}
\AXC{$\varphi_{87}$\lbl{[E42]}}
\AXC{$\varphi_{87} \to \varphi_{241}$\lbl{A2}}
\BIC{$\varphi_{241}$}
\AXC{$\varphi_{241} \to \varphi_{242}$\lbl{A2}}
\BIC{$\varphi_{242}$}
\AXC{$\varphi_{242} \to f(\xi) \in \mathbb{R} \to \mu \in \mathbb{R} \to \varphi_{114}$\lbl{A2}}
\BIC{$f(\xi) \in \mathbb{R} \to \mu \in \mathbb{R} \to \varphi_{114}$}
\AXC{$(f(\xi) \in \mathbb{R} \to \mu \in \mathbb{R} \to \varphi_{114}) \to \varphi_{878}$\lbl{A1}}
\BIC{$\varphi_{878}$}
\end{ptree}
\noindent$\Pi_{72}$:
\begin{ptree}
\AXC{$\stackrel{\Pi_{66}}{\vdots}$}\noLine
\UIC{$\mathrm{ContOn}(f,[a,b]) \to \varphi_{673}$}
\AXC{$\varphi_{56}$\lbl{A1}}
\AXC{$\stackrel{\Pi_{60}}{\vdots}$}\noLine
\UIC{$\mathrm{ContOn}(f,[a,b]) \to \varphi_{347}$}
\AXC{$\stackrel{\Pi_{71}}{\vdots}$}\noLine
\UIC{$\varphi_{878}$}
\BIC{$\varphi_{56} \to (\mathrm{ContOn}(f,[a,b]) \to \varphi_{673}) \to \varphi_{828}$}
\BIC{$(\mathrm{ContOn}(f,[a,b]) \to \varphi_{673}) \to \varphi_{828}$}
\BIC{$\varphi_{828}$}
\end{ptree}
\noindent$\Pi_{73}$:
\begin{ptree}
\AXC{$\forall x\,\forall y\,\forall z\,(x < y \to y < z \to x < z)$\lbl{[O2]}}
\AXC{$\varphi_{41}$\lbl{A2}}
\BIC{$\forall y\,\forall z\,(\mu < y \to y < z \to \mu < z)$}
\AXC{$\varphi_{88}$\lbl{A2}}
\BIC{$\forall z\,(\mu < f(s) \to f(s) < z \to \mu < z)$}
\AXC{$\varphi_{164}$\lbl{A2}}
\BIC{$\mu < f(s) \to f(s) < \mu \to \mu < \mu$}
\AXC{$\varphi_{867}$\lbl{A1}}
\BIC{$(\mathrm{ContOn}(f,[a,b]) \to \varphi_{667}) \to \varphi_{819}$}
\end{ptree}
\noindent$\Pi_{74}$:
\begin{ptree}
\AXC{$\stackrel{\Pi_{70}}{\vdots}$}\noLine
\UIC{$\mathrm{ContOn}(f,[a,b]) \to \varphi_{697}$}
\AXC{$\stackrel{\Pi_{72}}{\vdots}$}\noLine
\UIC{$\varphi_{828}$}
\BIC{$\mathrm{ContOn}(f,[a,b]) \to \varphi_{667}$}
\AXC{$\stackrel{\Pi_{73}}{\vdots}$}\noLine
\UIC{$(\mathrm{ContOn}(f,[a,b]) \to \varphi_{667}) \to \varphi_{819}$}
\BIC{$\varphi_{819}$}
\end{ptree}
\noindent$\Pi_{75}$:
\begin{ptree}
\AXC{$\mathrm{ContOn}(f,[a,b]) \to \varphi_{631}$\lbl{A1}}
\AXC{$\stackrel{\Pi_{8}}{\vdots}$}\noLine
\UIC{$\varphi_{808}$}
\BIC{$\mathrm{ContOn}(f,[a,b]) \to \varphi_{666}$}
\AXC{$\stackrel{\Pi_{74}}{\vdots}$}\noLine
\UIC{$\varphi_{819}$}
\BIC{$\mathrm{ContOn}(f,[a,b]) \to \varphi_{659}$}
\AXC{$(\mathrm{ContOn}(f,[a,b]) \to \varphi_{659}) \to \varphi_{806}$\lbl{A1}}
\BIC{$\varphi_{806}$}
\end{ptree}
\noindent$\Pi_{76}$:
\begin{ptree}
\AXC{$\forall x\,(\neg x < x)$\lbl{[O1]}}
\AXC{$\forall x\,(\neg x < x) \to \neg \mu < \mu$\lbl{A2}}
\BIC{$\neg \mu < \mu$}
\AXC{$\neg \mu < \mu \to \mathrm{ContOn}(f,[a,b]) \to \varphi_{661}$\lbl{A1}}
\BIC{$\mathrm{ContOn}(f,[a,b]) \to \varphi_{661}$}
\AXC{$\stackrel{\Pi_{75}}{\vdots}$}\noLine
\UIC{$\varphi_{806}$}
\BIC{$\mathrm{ContOn}(f,[a,b]) \to \varphi_{637}$}
\AXC{$\varphi_{796}$\lbl{A1}}
\BIC{$\varphi_{636} \to s \in W \to s \leq \xi - d$}
\UIC{$\forall s\,(\varphi_{636} \to s \in W \to s \leq \xi - d)$}
\end{ptree}
\noindent$\Pi_{77}$:
\begin{ptree}
\AXC{$\stackrel{\Pi_{3}}{\vdots}$}\noLine
\UIC{$\mathrm{ContOn}(f,[a,b]) \to \varphi_{328}$}
\AXC{$\forall x\,\forall y\,(x \in \mathbb{R} \to y \in \mathbb{R} \to x - y \in \mathbb{R})$\lbl{[F1]}}
\AXC{$\varphi_{227}$\lbl{A2}}
\BIC{$\forall y\,(\xi \in \mathbb{R} \to y \in \mathbb{R} \to \xi - y \in \mathbb{R})$}
\AXC{$\varphi_{196}$\lbl{A2}}
\BIC{$\xi \in \mathbb{R} \to d \in \mathbb{R} \to \xi - d \in \mathbb{R}$}
\AXC{$\varphi_{851}$\lbl{A1}}
\BIC{$(\mathrm{ContOn}(f,[a,b]) \to \varphi_{328}) \to \varphi_{790}$}
\BIC{$\varphi_{790}$}
\end{ptree}
\noindent$\Pi_{78}$:
\begin{ptree}
\AXC{$\forall a\,\forall X\,(\mathrm{UB}(a,X) \leftrightarrow a \in \mathbb{R} \land \forall x \in X\,(x \leq a))$\lbl{DefUB}}
\AXC{$\forall a\,\forall X\,(\mathrm{UB}(a,X) \leftrightarrow a \in \mathbb{R} \land \forall x \in X\,(x \leq a)) \to \varphi_{22}$\lbl{A2}}
\BIC{$\varphi_{22}$}
\AXC{$\varphi_{292}$\lbl{A2}}
\BIC{$\mathrm{UB}(\xi - d,W) \leftrightarrow \xi - d \in \mathbb{R} \land \forall x \in W\,(x \leq \xi - d)$}
\end{ptree}
\noindent$\Pi_{79}$:
\begin{ptree}
\AXC{$\stackrel{\Pi_{78}}{\vdots}$}\noLine
\UIC{$\mathrm{UB}(\xi - d,W) \leftrightarrow \xi - d \in \mathbb{R} \land \forall x \in W\,(x \leq \xi - d)$}
\AXC{$\varphi_{874}$\lbl{A1}}
\BIC{$(\mathrm{ContOn}(f,[a,b]) \to \varphi_{628}) \to \varphi_{841}$}
\end{ptree}
\noindent$\Pi_{80}$:
\begin{ptree}
\AXC{$\mathrm{ContOn}(f,[a,b]) \to \varphi_{678}$\lbl{A1}}
\AXC{$\varphi_{811}$\lbl{A1}}
\BIC{$\mathrm{ContOn}(f,[a,b]) \to \varphi_{610}$}
\AXC{$\stackrel{\Pi_{77}}{\vdots}$}\noLine
\UIC{$\varphi_{790}$}
\BIC{$\mathrm{ContOn}(f,[a,b]) \to \varphi_{628}$}
\AXC{$\stackrel{\Pi_{79}}{\vdots}$}\noLine
\UIC{$(\mathrm{ContOn}(f,[a,b]) \to \varphi_{628}) \to \varphi_{841}$}
\BIC{$\varphi_{841}$}
\end{ptree}
\noindent$\Pi_{81}$:
\begin{ptree}
\AXC{$\varphi_{636} \to \forall s\,(\varphi_{636})$\lbl{A4}}
\AXC{$\varphi_{803}$\lbl{A3}}
\AXC{$\stackrel{\Pi_{76}}{\vdots}$}\noLine
\UIC{$\forall s\,(\varphi_{636} \to s \in W \to s \leq \xi - d)$}
\AXC{$\varphi_{894}$\lbl{A1}}
\BIC{$\varphi_{803} \to (\varphi_{636} \to \forall s\,(\varphi_{636})) \to \mathrm{ContOn}(f,[a,b]) \to \varphi_{646}$}
\BIC{$(\varphi_{636} \to \forall s\,(\varphi_{636})) \to \mathrm{ContOn}(f,[a,b]) \to \varphi_{646}$}
\BIC{$\mathrm{ContOn}(f,[a,b]) \to \varphi_{646}$}
\AXC{$\stackrel{\Pi_{80}}{\vdots}$}\noLine
\UIC{$\varphi_{841}$}
\BIC{$\varphi_{94} \to \varphi_{93} \to \mathrm{ContOn}(f,[a,b]) \to \varphi_{616}$}
\end{ptree}
\noindent$\Pi_{82}$:
\begin{ptree}
\AXC{$\varphi_{9}$\lbl{A2}}
\AXC{$\stackrel{\Pi_{40}}{\vdots}$}\noLine
\UIC{$\mathrm{ContOn}(f,[a,b]) \to \varphi_{343}$}
\AXC{$\varphi_{721}$\lbl{A1}}
\BIC{$\varphi_{9} \to \mathrm{ContOn}(f,[a,b]) \to \varphi_{372}$}
\BIC{$\mathrm{ContOn}(f,[a,b]) \to \varphi_{372}$}
\AXC{$(\mathrm{ContOn}(f,[a,b]) \to \varphi_{372}) \to \varphi_{783}$\lbl{A1}}
\BIC{$\varphi_{783}$}
\end{ptree}
\noindent$\Pi_{83}$:
\begin{ptree}
\AXC{$\forall x\,\forall y\,(x \leq y \to \neg y < x)$\lbl{[O7]}}
\AXC{$\forall x\,\forall y\,(x \leq y \to \neg y < x) \to \forall y\,(\xi \leq y \to \neg y < \xi)$\lbl{A2}}
\BIC{$\forall y\,(\xi \leq y \to \neg y < \xi)$}
\end{ptree}
\noindent$\Pi_{84}$:
\begin{ptree}
\AXC{$\stackrel{\Pi_{83}}{\vdots}$}\noLine
\UIC{$\forall y\,(\xi \leq y \to \neg y < \xi)$}
\AXC{$\forall y\,(\xi \leq y \to \neg y < \xi) \to \xi \leq \xi - d \to \neg \xi - d < \xi$\lbl{A2}}
\BIC{$\xi \leq \xi - d \to \neg \xi - d < \xi$}
\AXC{$(\xi \leq \xi - d \to \neg \xi - d < \xi) \to \varphi_{785}$\lbl{A1}}
\BIC{$\varphi_{785}$}
\end{ptree}
\noindent$\Pi_{85}$:
\begin{ptree}
\AXC{$\stackrel{\Pi_{6}}{\vdots}$}\noLine
\UIC{$\varphi_{93}$}
\AXC{$\stackrel{\Pi_{7}}{\vdots}$}\noLine
\UIC{$\varphi_{94}$}
\AXC{$\stackrel{\Pi_{81}}{\vdots}$}\noLine
\UIC{$\varphi_{94} \to \varphi_{93} \to \mathrm{ContOn}(f,[a,b]) \to \varphi_{616}$}
\BIC{$\varphi_{93} \to \mathrm{ContOn}(f,[a,b]) \to \varphi_{616}$}
\BIC{$\mathrm{ContOn}(f,[a,b]) \to \varphi_{616}$}
\AXC{$\stackrel{\Pi_{82}}{\vdots}$}\noLine
\UIC{$\varphi_{783}$}
\BIC{$\mathrm{ContOn}(f,[a,b]) \to \varphi_{611}$}
\AXC{$\stackrel{\Pi_{84}}{\vdots}$}\noLine
\UIC{$\varphi_{785}$}
\BIC{$\mathrm{ContOn}(f,[a,b]) \to \varphi_{621}$}
\end{ptree}
\noindent$\Pi_{86}$:
\begin{ptree}
\AXC{$\stackrel{\Pi_{85}}{\vdots}$}\noLine
\UIC{$\mathrm{ContOn}(f,[a,b]) \to \varphi_{621}$}
\AXC{$(\mathrm{ContOn}(f,[a,b]) \to \varphi_{621}) \to \varphi_{816}$\lbl{A1}}
\BIC{$\varphi_{816}$}
\end{ptree}
\noindent$\Pi_{87}$:
\begin{ptree}
\AXC{$\mathrm{ContOn}(f,[a,b]) \to \varphi_{678}$\lbl{A1}}
\AXC{$\varphi_{807}$\lbl{A1}}
\BIC{$\mathrm{ContOn}(f,[a,b]) \to \varphi_{601}$}
\AXC{$\stackrel{\Pi_{5}}{\vdots}$}\noLine
\UIC{$\varphi_{780}$}
\BIC{$\mathrm{ContOn}(f,[a,b]) \to \varphi_{615}$}
\AXC{$\stackrel{\Pi_{86}}{\vdots}$}\noLine
\UIC{$\varphi_{816}$}
\BIC{$\mathrm{ContOn}(f,[a,b]) \to \varphi_{680}$}
\AXC{$\varphi_{826}$\lbl{A1}}
\BIC{$\varphi_{645} \to \neg (d \in \mathbb{R} \land 0 < d \land d \leq \delta \cdot 2^{-1} \land d \leq \xi - a)$}
\UIC{$\varphi_{702}$}
\end{ptree}
\noindent$\Pi_{88}$:
\begin{ptree}
\AXC{$\forall x\,(\neg x < x)$\lbl{[O1]}}
\AXC{$\forall x\,(\neg x < x) \to \neg \mu < \mu$\lbl{A2}}
\BIC{$\neg \mu < \mu$}
\AXC{$\neg \mu < \mu \to \mathrm{ContOn}(f,[a,b]) \to \varphi_{424}$\lbl{A1}}
\BIC{$\mathrm{ContOn}(f,[a,b]) \to \varphi_{424}$}
\end{ptree}
\noindent$\Pi_{89}$:
\begin{ptree}
\AXC{$\mathrm{ContOn}(f,[a,b]) \to \varphi_{102}$\lbl{A1}}
\AXC{$\mathrm{ContOn}(f,[a,b]) \to \varphi_{403}$\lbl{A1}}
\AXC{$\varphi_{699}$\lbl{A1}}
\BIC{$\mathrm{ContOn}(f,[a,b]) \to \varphi_{398}$}
\AXC{$\varphi_{747}$\lbl{A1}}
\BIC{$(\mathrm{ContOn}(f,[a,b]) \to \varphi_{102}) \to \varphi_{494}$}
\BIC{$\varphi_{494}$}
\end{ptree}
\noindent$\Pi_{90}$:
\begin{ptree}
\AXC{$z = w \to f(z) < \mu \to f(w) < \mu$\lbl{A6}}
\UIC{$\forall w\,(z = w \to f(z) < \mu \to f(w) < \mu)$}
\UIC{$\forall z\,\forall w\,(z = w \to f(z) < \mu \to f(w) < \mu)$}
\AXC{$\varphi_{142}$\lbl{A2}}
\BIC{$\forall w\,(a = w \to f(a) < \mu \to f(w) < \mu)$}
\AXC{$\varphi_{107}$\lbl{A2}}
\BIC{$a = \xi \to f(a) < \mu \to f(\xi) < \mu$}
\AXC{$\stackrel{\Pi_{89}}{\vdots}$}\noLine
\UIC{$\varphi_{494}$}
\BIC{$\mathrm{ContOn}(f,[a,b]) \to \varphi_{438}$}
\end{ptree}
\noindent$\Pi_{91}$:
\begin{ptree}
\AXC{$\forall x\,\forall y\,\forall z\,(x < y \to y < z \to x < z)$\lbl{[O2]}}
\AXC{$\varphi_{41}$\lbl{A2}}
\BIC{$\forall y\,\forall z\,(\mu < y \to y < z \to \mu < z)$}
\AXC{$\varphi_{101}$\lbl{A2}}
\BIC{$\forall z\,(\mu < f(\xi) \to f(\xi) < z \to \mu < z)$}
\AXC{$\varphi_{197}$\lbl{A2}}
\BIC{$\mu < f(\xi) \to f(\xi) < \mu \to \mu < \mu$}
\AXC{$\varphi_{764}$\lbl{A1}}
\BIC{$(\mathrm{ContOn}(f,[a,b]) \to \varphi_{404}) \to \varphi_{715}$}
\end{ptree}
\noindent$\Pi_{92}$:
\begin{ptree}
\AXC{$\stackrel{\Pi_{90}}{\vdots}$}\noLine
\UIC{$\mathrm{ContOn}(f,[a,b]) \to \varphi_{438}$}
\AXC{$\mathrm{ContOn}(f,[a,b]) \to \varphi_{404}$\lbl{A1}}
\AXC{$\stackrel{\Pi_{91}}{\vdots}$}\noLine
\UIC{$(\mathrm{ContOn}(f,[a,b]) \to \varphi_{404}) \to \varphi_{715}$}
\BIC{$\varphi_{715}$}
\BIC{$\mathrm{ContOn}(f,[a,b]) \to \varphi_{416}$}
\end{ptree}
\noindent$\Pi_{93}$:
\begin{ptree}
\AXC{$\stackrel{\Pi_{88}}{\vdots}$}\noLine
\UIC{$\mathrm{ContOn}(f,[a,b]) \to \varphi_{424}$}
\AXC{$\stackrel{\Pi_{92}}{\vdots}$}\noLine
\UIC{$\mathrm{ContOn}(f,[a,b]) \to \varphi_{416}$}
\AXC{$(\mathrm{ContOn}(f,[a,b]) \to \varphi_{416}) \to \varphi_{694}$\lbl{A1}}
\BIC{$\varphi_{694}$}
\BIC{$\mathrm{ContOn}(f,[a,b]) \to \varphi_{364}$}
\end{ptree}
\noindent$\Pi_{94}$:
\begin{ptree}
\AXC{$\stackrel{\Pi_{93}}{\vdots}$}\noLine
\UIC{$\mathrm{ContOn}(f,[a,b]) \to \varphi_{364}$}
\AXC{$\stackrel{\Pi_{50}}{\vdots}$}\noLine
\UIC{$\mathrm{ContOn}(f,[a,b]) \to \varphi_{317}$}
\AXC{$(\mathrm{ContOn}(f,[a,b]) \to \varphi_{317}) \to \varphi_{685}$\lbl{A1}}
\BIC{$\varphi_{685}$}
\BIC{$\mathrm{ContOn}(f,[a,b]) \to \varphi_{369}$}
\end{ptree}
\noindent$\Pi_{95}$:
\begin{ptree}
\AXC{$\forall x\,\forall y\,(x \in \mathbb{R} \to y \in \mathbb{R} \to x < y \to 0 < y - x)$\lbl{[F24]}}
\AXC{$\varphi_{265}$\lbl{A2}}
\BIC{$\forall y\,(a \in \mathbb{R} \to y \in \mathbb{R} \to a < y \to 0 < y - a)$}
\AXC{$\varphi_{251}$\lbl{A2}}
\BIC{$a \in \mathbb{R} \to \xi \in \mathbb{R} \to a < \xi \to 0 < \xi - a$}
\AXC{$\varphi_{823}$\lbl{A1}}
\BIC{$(\mathrm{ContOn}(f,[a,b]) \to a \in \mathbb{R} \to a \in \mathbb{R}) \to \varphi_{765}$}
\end{ptree}
\noindent$\Pi_{96}$:
\begin{ptree}
\AXC{$\mathrm{ContOn}(f,[a,b]) \to a \in \mathbb{R} \to a \in \mathbb{R}$\lbl{A1}}
\AXC{$\stackrel{\Pi_{95}}{\vdots}$}\noLine
\UIC{$(\mathrm{ContOn}(f,[a,b]) \to a \in \mathbb{R} \to a \in \mathbb{R}) \to \varphi_{765}$}
\BIC{$\varphi_{765}$}
\end{ptree}
\noindent$\Pi_{97}$:
\begin{ptree}
\AXC{$\forall e \in \mathbb{R}\,(0 < e \to 0 < e \cdot 2^{-1})$\lbl{[F35]}}
\AXC{$\varphi_{134}$\lbl{A2}}
\BIC{$\delta \in \mathbb{R} \to 0 < \delta \to 0 < \delta \cdot 2^{-1}$}
\AXC{$(\delta \in \mathbb{R} \to 0 < \delta \to 0 < \delta \cdot 2^{-1}) \to \varphi_{578} \to \varphi_{576} \to \varphi_{588}$\lbl{A1}}
\BIC{$\varphi_{578} \to \varphi_{576} \to \varphi_{588}$}
\end{ptree}
\noindent$\Pi_{98}$:
\begin{ptree}
\AXC{$\forall x\,\forall y\,(x \in \mathbb{R} \to y \in \mathbb{R} \to x - y \in \mathbb{R})$\lbl{[F1]}}
\AXC{$\varphi_{227}$\lbl{A2}}
\BIC{$\forall y\,(\xi \in \mathbb{R} \to y \in \mathbb{R} \to \xi - y \in \mathbb{R})$}
\AXC{$\varphi_{198}$\lbl{A2}}
\BIC{$\xi \in \mathbb{R} \to a \in \mathbb{R} \to \xi - a \in \mathbb{R}$}
\AXC{$(\xi \in \mathbb{R} \to a \in \mathbb{R} \to \xi - a \in \mathbb{R}) \to \varphi_{740}$\lbl{A1}}
\BIC{$\varphi_{740}$}
\end{ptree}
\noindent$\Pi_{99}$:
\begin{ptree}
\AXC{$\mathrm{ContOn}(f,[a,b]) \to a \in \mathbb{R} \to a \in \mathbb{R}$\lbl{A1}}
\AXC{$\stackrel{\Pi_{3}}{\vdots}$}\noLine
\UIC{$\mathrm{ContOn}(f,[a,b]) \to \varphi_{328}$}
\AXC{$\stackrel{\Pi_{98}}{\vdots}$}\noLine
\UIC{$\varphi_{740}$}
\BIC{$(\mathrm{ContOn}(f,[a,b]) \to a \in \mathbb{R} \to a \in \mathbb{R}) \to \varphi_{586}$}
\BIC{$\varphi_{586}$}
\end{ptree}
\noindent$\Pi_{100}$:
\begin{ptree}
\AXC{$\varphi_{709}$\lbl{A1}}
\AXC{$\varphi_{709} \to \varphi_{578}$\lbl{A1}}
\BIC{$\varphi_{578}$}
\AXC{$\forall x\,\forall y\,(x \in \mathbb{R} \to y \in \mathbb{R} \to x \cdot y \in \mathbb{R})$\lbl{(R1)}}
\AXC{$\varphi_{228}$\lbl{A2}}
\BIC{$\forall y\,(\delta \in \mathbb{R} \to y \in \mathbb{R} \to \delta \cdot y \in \mathbb{R})$}
\AXC{$\varphi_{257}$\lbl{A2}}
\BIC{$\delta \in \mathbb{R} \to 2^{-1} \in \mathbb{R} \to \delta \cdot 2^{-1} \in \mathbb{R}$}
\AXC{$\varphi_{769}$\lbl{A1}}
\BIC{$\varphi_{578} \to 2^{-1} \in \mathbb{R} \to \varphi_{590}$}
\BIC{$2^{-1} \in \mathbb{R} \to \varphi_{590}$}
\end{ptree}
\noindent$\Pi_{101}$:
\begin{ptree}
\AXC{$\varphi_{127}$\lbl{[D8]}}
\AXC{$\varphi_{127} \to \varphi_{286}$\lbl{A2}}
\BIC{$\varphi_{286}$}
\AXC{$\varphi_{286} \to \delta \cdot 2^{-1} \in \mathbb{R} \to \xi - a \in \mathbb{R} \to 0 < \delta \cdot 2^{-1} \to \varphi_{37}$\lbl{A2}}
\BIC{$\delta \cdot 2^{-1} \in \mathbb{R} \to \xi - a \in \mathbb{R} \to 0 < \delta \cdot 2^{-1} \to \varphi_{37}$}
\AXC{$\varphi_{886}$\lbl{A1}}
\BIC{$\varphi_{590} \to \varphi_{586} \to \varphi_{588} \to \varphi_{582} \to \mathrm{ContOn}(f,[a,b]) \to \varphi_{596}$}
\end{ptree}
\noindent$\Pi_{102}$:
\begin{ptree}
\AXC{$\stackrel{\Pi_{99}}{\vdots}$}\noLine
\UIC{$\varphi_{586}$}
\AXC{$2^{-1} \in \mathbb{R}$\lbl{[F31]}}
\AXC{$\stackrel{\Pi_{100}}{\vdots}$}\noLine
\UIC{$2^{-1} \in \mathbb{R} \to \varphi_{590}$}
\BIC{$\varphi_{590}$}
\AXC{$\stackrel{\Pi_{101}}{\vdots}$}\noLine
\UIC{$\varphi_{590} \to \varphi_{586} \to \varphi_{588} \to \varphi_{582} \to \mathrm{ContOn}(f,[a,b]) \to \varphi_{596}$}
\BIC{$\varphi_{586} \to \varphi_{588} \to \varphi_{582} \to \mathrm{ContOn}(f,[a,b]) \to \varphi_{596}$}
\BIC{$\varphi_{588} \to \varphi_{582} \to \mathrm{ContOn}(f,[a,b]) \to \varphi_{596}$}
\end{ptree}
\noindent$\Pi_{103}$:
\begin{ptree}
\AXC{$\varphi_{709}$\lbl{A1}}
\AXC{$\varphi_{709} \to \varphi_{576}$\lbl{A1}}
\BIC{$\varphi_{576}$}
\AXC{$\varphi_{709}$\lbl{A1}}
\AXC{$\varphi_{709} \to \varphi_{578}$\lbl{A1}}
\BIC{$\varphi_{578}$}
\AXC{$\stackrel{\Pi_{97}}{\vdots}$}\noLine
\UIC{$\varphi_{578} \to \varphi_{576} \to \varphi_{588}$}
\BIC{$\varphi_{576} \to \varphi_{588}$}
\BIC{$\varphi_{588}$}
\AXC{$\stackrel{\Pi_{102}}{\vdots}$}\noLine
\UIC{$\varphi_{588} \to \varphi_{582} \to \mathrm{ContOn}(f,[a,b]) \to \varphi_{596}$}
\BIC{$\varphi_{582} \to \mathrm{ContOn}(f,[a,b]) \to \varphi_{596}$}
\end{ptree}
\noindent$\Pi_{104}$:
\begin{ptree}
\AXC{$\stackrel{\Pi_{94}}{\vdots}$}\noLine
\UIC{$\mathrm{ContOn}(f,[a,b]) \to \varphi_{369}$}
\AXC{$\stackrel{\Pi_{3}}{\vdots}$}\noLine
\UIC{$\mathrm{ContOn}(f,[a,b]) \to \varphi_{328}$}
\AXC{$\stackrel{\Pi_{96}}{\vdots}$}\noLine
\UIC{$\varphi_{765}$}
\BIC{$(\mathrm{ContOn}(f,[a,b]) \to \varphi_{369}) \to \varphi_{582}$}
\BIC{$\varphi_{582}$}
\AXC{$\stackrel{\Pi_{103}}{\vdots}$}\noLine
\UIC{$\varphi_{582} \to \mathrm{ContOn}(f,[a,b]) \to \varphi_{596}$}
\BIC{$\mathrm{ContOn}(f,[a,b]) \to \varphi_{596}$}
\end{ptree}
\noindent$\Pi_{105}$:
\begin{ptree}
\AXC{$\varphi_{702} \to \varphi_{707}$\lbl{A3}}
\AXC{$\stackrel{\Pi_{87}}{\vdots}$}\noLine
\UIC{$\varphi_{702}$}
\AXC{$\stackrel{\Pi_{104}}{\vdots}$}\noLine
\UIC{$\mathrm{ContOn}(f,[a,b]) \to \varphi_{596}$}
\AXC{$(\mathrm{ContOn}(f,[a,b]) \to \varphi_{596}) \to \varphi_{702} \to \varphi_{890}$\lbl{A1}}
\BIC{$\varphi_{702} \to \varphi_{890}$}
\BIC{$\varphi_{890}$}
\BIC{$(\varphi_{645} \to \forall d\,(\varphi_{645})) \to \mathrm{ContOn}(f,[a,b]) \to \varphi_{594}$}
\end{ptree}
\noindent$\Pi_{106}$:
\begin{ptree}
\AXC{$\varphi_{56}$\lbl{A1}}
\AXC{$\forall x\,\forall y\,(x \in \mathbb{R} \to y \in \mathbb{R} \to x < y \to 0 < y - x)$\lbl{[F24]}}
\AXC{$\varphi_{274}$\lbl{A2}}
\BIC{$\forall y\,(\mu \in \mathbb{R} \to y \in \mathbb{R} \to \mu < y \to 0 < y - \mu)$}
\AXC{$\varphi_{298}$\lbl{A2}}
\BIC{$\mu \in \mathbb{R} \to f(\xi) \in \mathbb{R} \to \mu < f(\xi) \to 0 < f(\xi) - \mu$}
\AXC{$\varphi_{801}$\lbl{A1}}
\BIC{$\varphi_{56} \to (\mathrm{ContOn}(f,[a,b]) \to \varphi_{347}) \to \varphi_{705}$}
\BIC{$(\mathrm{ContOn}(f,[a,b]) \to \varphi_{347}) \to \varphi_{705}$}
\end{ptree}
\noindent$\Pi_{107}$:
\begin{ptree}
\AXC{$\mathrm{ContOn}(f,[a,b]) \to \varphi_{404}$\lbl{A1}}
\AXC{$\stackrel{\Pi_{60}}{\vdots}$}\noLine
\UIC{$\mathrm{ContOn}(f,[a,b]) \to \varphi_{347}$}
\AXC{$\stackrel{\Pi_{106}}{\vdots}$}\noLine
\UIC{$(\mathrm{ContOn}(f,[a,b]) \to \varphi_{347}) \to \varphi_{705}$}
\BIC{$\varphi_{705}$}
\BIC{$\mathrm{ContOn}(f,[a,b]) \to \varphi_{432}$}
\end{ptree}
\noindent$\Pi_{108}$:
\begin{ptree}
\AXC{$\forall x\,\forall y\,(x \in \mathbb{R} \to y \in \mathbb{R} \to x - y \in \mathbb{R})$\lbl{[F1]}}
\AXC{$\varphi_{283}$\lbl{A2}}
\BIC{$\forall y\,(f(\xi) \in \mathbb{R} \to y \in \mathbb{R} \to f(\xi) - y \in \mathbb{R})$}
\AXC{$\varphi_{296}$\lbl{A2}}
\BIC{$f(\xi) \in \mathbb{R} \to \mu \in \mathbb{R} \to f(\xi) - \mu \in \mathbb{R}$}
\AXC{$(f(\xi) \in \mathbb{R} \to \mu \in \mathbb{R} \to f(\xi) - \mu \in \mathbb{R}) \to \varphi_{728}$\lbl{A1}}
\BIC{$\varphi_{728}$}
\end{ptree}
\noindent$\Pi_{109}$:
\begin{ptree}
\AXC{$\forall f\,(\varphi_{422})$\lbl{DefContIn}}
\AXC{$\forall f\,(\varphi_{422}) \to \varphi_{422}$\lbl{A2}}
\BIC{$\varphi_{422}$}
\AXC{$\varphi_{422} \to \varphi_{445}$\lbl{A2}}
\BIC{$\varphi_{445}$}
\AXC{$\varphi_{445} \to \varphi_{446}$\lbl{A2}}
\BIC{$\varphi_{446}$}
\end{ptree}
\noindent$\Pi_{110}$:
\begin{ptree}
\AXC{$\varphi_{143}$\lbl{A2}}
\AXC{$\forall f\,(\varphi_{40})$\lbl{DefContOn}}
\AXC{$\forall f\,(\varphi_{40}) \to \varphi_{40}$\lbl{A2}}
\BIC{$\varphi_{40}$}
\AXC{$\varphi_{40} \to \varphi_{231}$\lbl{A2}}
\BIC{$\varphi_{231}$}
\AXC{$\varphi_{4}$\lbl{A1}}
\AXC{$\varphi_{4} \to \varphi_{231} \to \varphi_{310}$\lbl{A1}}
\BIC{$\varphi_{231} \to \varphi_{310}$}
\BIC{$\varphi_{310}$}
\AXC{$\varphi_{310} \to \varphi_{143} \to \varphi_{309}$\lbl{A1}}
\BIC{$\varphi_{143} \to \varphi_{309}$}
\BIC{$\varphi_{309}$}
\end{ptree}
\noindent$\Pi_{111}$:
\begin{ptree}
\AXC{$\stackrel{\Pi_{109}}{\vdots}$}\noLine
\UIC{$\varphi_{446}$}
\AXC{$\stackrel{\Pi_{53}}{\vdots}$}\noLine
\UIC{$\mathrm{ContOn}(f,[a,b]) \to \varphi_{341}$}
\AXC{$\stackrel{\Pi_{110}}{\vdots}$}\noLine
\UIC{$\varphi_{309}$}
\AXC{$\varphi_{309} \to \varphi_{766}$\lbl{A1}}
\BIC{$\varphi_{766}$}
\BIC{$\varphi_{446} \to \mathrm{ContOn}(f,[a,b]) \to \varphi_{553}$}
\BIC{$\mathrm{ContOn}(f,[a,b]) \to \varphi_{553}$}
\end{ptree}
\noindent$\Pi_{112}$:
\begin{ptree}
\AXC{$\varphi_{302} \to \varphi_{355}$\lbl{A2}}
\AXC{$\stackrel{\Pi_{111}}{\vdots}$}\noLine
\UIC{$\mathrm{ContOn}(f,[a,b]) \to \varphi_{553}$}
\AXC{$\varphi_{845}$\lbl{A1}}
\BIC{$(\varphi_{302} \to \varphi_{355}) \to \mathrm{ContOn}(f,[a,b]) \to \varphi_{580}$}
\BIC{$\mathrm{ContOn}(f,[a,b]) \to \varphi_{580}$}
\AXC{$(\mathrm{ContOn}(f,[a,b]) \to \varphi_{580}) \to \varphi_{778}$\lbl{A1}}
\BIC{$\varphi_{778}$}
\end{ptree}
\noindent$\Pi_{113}$:
\begin{ptree}
\AXC{$\stackrel{\Pi_{107}}{\vdots}$}\noLine
\UIC{$\mathrm{ContOn}(f,[a,b]) \to \varphi_{432}$}
\AXC{$\varphi_{56}$\lbl{A1}}
\AXC{$\stackrel{\Pi_{60}}{\vdots}$}\noLine
\UIC{$\mathrm{ContOn}(f,[a,b]) \to \varphi_{347}$}
\AXC{$\stackrel{\Pi_{108}}{\vdots}$}\noLine
\UIC{$\varphi_{728}$}
\BIC{$\varphi_{56} \to \mathrm{ContOn}(f,[a,b]) \to \varphi_{439}$}
\BIC{$\mathrm{ContOn}(f,[a,b]) \to \varphi_{439}$}
\AXC{$\stackrel{\Pi_{112}}{\vdots}$}\noLine
\UIC{$\varphi_{778}$}
\BIC{$(\mathrm{ContOn}(f,[a,b]) \to \varphi_{432}) \to \varphi_{573}$}
\BIC{$\varphi_{573}$}
\AXC{$\varphi_{573} \to \varphi_{888}$\lbl{A1}}
\BIC{$\varphi_{888}$}
\end{ptree}
\noindent$\Pi_{114}$:
\begin{ptree}
\AXC{$\varphi_{776}$\lbl{A3}}
\AXC{$\varphi_{645} \to \forall d\,(\varphi_{645})$\lbl{A4}}
\AXC{$\stackrel{\Pi_{105}}{\vdots}$}\noLine
\UIC{$(\varphi_{645} \to \forall d\,(\varphi_{645})) \to \mathrm{ContOn}(f,[a,b]) \to \varphi_{594}$}
\BIC{$\mathrm{ContOn}(f,[a,b]) \to \varphi_{594}$}
\AXC{$\varphi_{767}$\lbl{A1}}
\BIC{$\varphi_{481} \to \neg (\delta \in \mathbb{R} \land 0 < \delta \land \varphi_{157})$}
\UIC{$\forall \delta\,(\varphi_{481} \to \neg (\delta \in \mathbb{R} \land 0 < \delta \land \varphi_{157}))$}
\AXC{$\stackrel{\Pi_{113}}{\vdots}$}\noLine
\UIC{$\varphi_{888}$}
\BIC{$\varphi_{776} \to (\varphi_{481} \to \forall \delta\,(\varphi_{481})) \to \varphi_{478}$}
\BIC{$(\varphi_{481} \to \forall \delta\,(\varphi_{481})) \to \varphi_{478}$}
\end{ptree}
\noindent$\Pi_{115}$:
\begin{ptree}
\AXC{$x = x$\lbl{A5}}
\AXC{$x = y \to x = x \to y = x$\lbl{A6}}
\AXC{$(x = y \to x = x \to y = x) \to x = x \to x = y \to y = x$\lbl{A1}}
\BIC{$x = x \to x = y \to y = x$}
\BIC{$x = y \to y = x$}
\UIC{$\forall y\,(x = y \to y = x)$}
\UIC{$\forall x\,\forall y\,(x = y \to y = x)$}
\end{ptree}
\noindent$\Pi_{116}$:
\begin{ptree}
\AXC{$\stackrel{\Pi_{115}}{\vdots}$}\noLine
\UIC{$\forall x\,\forall y\,(x = y \to y = x)$}
\AXC{$\forall x\,\forall y\,(x = y \to y = x) \to \forall y\,(\xi = y \to y = \xi)$\lbl{A2}}
\BIC{$\forall y\,(\xi = y \to y = \xi)$}
\AXC{$\forall y\,(\xi = y \to y = \xi) \to \xi = b \to b = \xi$\lbl{A2}}
\BIC{$\xi = b \to b = \xi$}
\end{ptree}
\noindent$\Pi_{117}$:
\begin{ptree}
\AXC{$\stackrel{\Pi_{116}}{\vdots}$}\noLine
\UIC{$\xi = b \to b = \xi$}
\AXC{$\mathrm{ContOn}(f,[a,b]) \to \varphi_{392}$\lbl{A1}}
\AXC{$\varphi_{692}$\lbl{A1}}
\BIC{$\mathrm{ContOn}(f,[a,b]) \to \varphi_{381}$}
\AXC{$\varphi_{704}$\lbl{A1}}
\BIC{$(\xi = b \to b = \xi) \to \mathrm{ContOn}(f,[a,b]) \to \varphi_{382}$}
\BIC{$\mathrm{ContOn}(f,[a,b]) \to \varphi_{382}$}
\AXC{$\varphi_{751}$\lbl{A1}}
\BIC{$(\mathrm{ContOn}(f,[a,b]) \to a \in \mathbb{R} \to \varphi_{132}) \to \varphi_{490}$}
\end{ptree}
\noindent$\Pi_{118}$:
\begin{ptree}
\AXC{$\mathrm{ContOn}(f,[a,b]) \to a \in \mathbb{R} \to \varphi_{132}$\lbl{A1}}
\AXC{$\stackrel{\Pi_{117}}{\vdots}$}\noLine
\UIC{$(\mathrm{ContOn}(f,[a,b]) \to a \in \mathbb{R} \to \varphi_{132}) \to \varphi_{490}$}
\BIC{$\varphi_{490}$}
\end{ptree}
\noindent$\Pi_{119}$:
\begin{ptree}
\AXC{$z = w \to \mu < f(z) \to \mu < f(w)$\lbl{A6}}
\UIC{$\forall w\,(z = w \to \mu < f(z) \to \mu < f(w))$}
\UIC{$\forall z\,\forall w\,(z = w \to \mu < f(z) \to \mu < f(w))$}
\AXC{$\varphi_{145}$\lbl{A2}}
\BIC{$\forall w\,(b = w \to \mu < f(b) \to \mu < f(w))$}
\AXC{$\varphi_{108}$\lbl{A2}}
\BIC{$b = \xi \to \mu < f(b) \to \mu < f(\xi)$}
\AXC{$\stackrel{\Pi_{118}}{\vdots}$}\noLine
\UIC{$\varphi_{490}$}
\BIC{$\mathrm{ContOn}(f,[a,b]) \to \varphi_{435}$}
\end{ptree}
\noindent$\Pi_{120}$:
\begin{ptree}
\AXC{$\stackrel{\Pi_{119}}{\vdots}$}\noLine
\UIC{$\mathrm{ContOn}(f,[a,b]) \to \varphi_{435}$}
\AXC{$\mathrm{ContOn}(f,[a,b]) \to \varphi_{376}$\lbl{A1}}
\AXC{$\varphi_{760}$\lbl{A1}}
\BIC{$(\mathrm{ContOn}(f,[a,b]) \to \varphi_{435}) \to \varphi_{484}$}
\BIC{$\varphi_{484}$}
\end{ptree}
\noindent$\Pi_{121}$:
\begin{ptree}
\AXC{$z = w \to \mu < z \to \mu < w$\lbl{A6}}
\UIC{$\forall w\,(z = w \to \mu < z \to \mu < w)$}
\UIC{$\forall z\,\forall w\,(z = w \to \mu < z \to \mu < w)$}
\AXC{$\varphi_{64}$\lbl{A2}}
\BIC{$\forall w\,(f(\xi) = w \to \mu < f(\xi) \to \mu < w)$}
\AXC{$\varphi_{144}$\lbl{A2}}
\BIC{$f(\xi) = \mu \to \mu < f(\xi) \to \mu < \mu$}
\AXC{$\stackrel{\Pi_{120}}{\vdots}$}\noLine
\UIC{$\varphi_{484}$}
\BIC{$\mathrm{ContOn}(f,[a,b]) \to \varphi_{400}$}
\AXC{$(\mathrm{ContOn}(f,[a,b]) \to \varphi_{400}) \to \varphi_{689}$\lbl{A1}}
\BIC{$\varphi_{689}$}
\end{ptree}
\noindent$\Pi_{122}$:
\begin{ptree}
\AXC{$\forall x\,(\neg x < x)$\lbl{[O1]}}
\AXC{$\forall x\,(\neg x < x) \to \neg \mu < \mu$\lbl{A2}}
\BIC{$\neg \mu < \mu$}
\AXC{$\neg \mu < \mu \to \mathrm{ContOn}(f,[a,b]) \to \varphi_{405}$\lbl{A1}}
\BIC{$\mathrm{ContOn}(f,[a,b]) \to \varphi_{405}$}
\AXC{$\stackrel{\Pi_{121}}{\vdots}$}\noLine
\UIC{$\varphi_{689}$}
\BIC{$\mathrm{ContOn}(f,[a,b]) \to \varphi_{356}$}
\end{ptree}
\noindent$\Pi_{123}$:
\begin{ptree}
\AXC{$\mathrm{ContOn}(f,[a,b]) \to \varphi_{102}$\lbl{A1}}
\AXC{$\mathrm{ContOn}(f,[a,b]) \to \varphi_{393}$\lbl{A1}}
\AXC{$\varphi_{693}$\lbl{A1}}
\BIC{$\mathrm{ContOn}(f,[a,b]) \to \varphi_{383}$}
\AXC{$\varphi_{745}$\lbl{A1}}
\BIC{$(\mathrm{ContOn}(f,[a,b]) \to \varphi_{102}) \to \varphi_{491}$}
\BIC{$\varphi_{491}$}
\end{ptree}
\noindent$\Pi_{124}$:
\begin{ptree}
\AXC{$z = w \to f(z) < \mu \to f(w) < \mu$\lbl{A6}}
\UIC{$\forall w\,(z = w \to f(z) < \mu \to f(w) < \mu)$}
\UIC{$\forall z\,\forall w\,(z = w \to f(z) < \mu \to f(w) < \mu)$}
\AXC{$\varphi_{142}$\lbl{A2}}
\BIC{$\forall w\,(a = w \to f(a) < \mu \to f(w) < \mu)$}
\AXC{$\varphi_{107}$\lbl{A2}}
\BIC{$a = \xi \to f(a) < \mu \to f(\xi) < \mu$}
\AXC{$\stackrel{\Pi_{123}}{\vdots}$}\noLine
\UIC{$\varphi_{491}$}
\BIC{$\mathrm{ContOn}(f,[a,b]) \to \varphi_{436}$}
\end{ptree}
\noindent$\Pi_{125}$:
\begin{ptree}
\AXC{$\stackrel{\Pi_{124}}{\vdots}$}\noLine
\UIC{$\mathrm{ContOn}(f,[a,b]) \to \varphi_{436}$}
\AXC{$\mathrm{ContOn}(f,[a,b]) \to \varphi_{376}$\lbl{A1}}
\AXC{$\varphi_{761}$\lbl{A1}}
\BIC{$(\mathrm{ContOn}(f,[a,b]) \to \varphi_{436}) \to \varphi_{485}$}
\BIC{$\varphi_{485}$}
\end{ptree}
\noindent$\Pi_{126}$:
\begin{ptree}
\AXC{$z = w \to z < \mu \to w < \mu$\lbl{A6}}
\UIC{$\forall w\,(z = w \to z < \mu \to w < \mu)$}
\UIC{$\forall z\,\forall w\,(z = w \to z < \mu \to w < \mu)$}
\AXC{$\varphi_{65}$\lbl{A2}}
\BIC{$\forall w\,(f(\xi) = w \to f(\xi) < \mu \to w < \mu)$}
\AXC{$\varphi_{146}$\lbl{A2}}
\BIC{$f(\xi) = \mu \to f(\xi) < \mu \to \mu < \mu$}
\AXC{$\stackrel{\Pi_{125}}{\vdots}$}\noLine
\UIC{$\varphi_{485}$}
\BIC{$\mathrm{ContOn}(f,[a,b]) \to \varphi_{401}$}
\AXC{$(\mathrm{ContOn}(f,[a,b]) \to \varphi_{401}) \to \varphi_{690}$\lbl{A1}}
\BIC{$\varphi_{690}$}
\end{ptree}
\noindent$\Pi_{127}$:
\begin{ptree}
\AXC{$\forall x\,(\neg x < x)$\lbl{[O1]}}
\AXC{$\forall x\,(\neg x < x) \to \neg \mu < \mu$\lbl{A2}}
\BIC{$\neg \mu < \mu$}
\AXC{$\neg \mu < \mu \to \mathrm{ContOn}(f,[a,b]) \to \varphi_{406}$\lbl{A1}}
\BIC{$\mathrm{ContOn}(f,[a,b]) \to \varphi_{406}$}
\AXC{$\stackrel{\Pi_{126}}{\vdots}$}\noLine
\UIC{$\varphi_{690}$}
\BIC{$\mathrm{ContOn}(f,[a,b]) \to \varphi_{357}$}
\end{ptree}
\noindent$\Pi_{128}$:
\begin{ptree}
\AXC{$\stackrel{\Pi_{127}}{\vdots}$}\noLine
\UIC{$\mathrm{ContOn}(f,[a,b]) \to \varphi_{357}$}
\AXC{$\stackrel{\Pi_{50}}{\vdots}$}\noLine
\UIC{$\mathrm{ContOn}(f,[a,b]) \to \varphi_{317}$}
\AXC{$\stackrel{\Pi_{3}}{\vdots}$}\noLine
\UIC{$\mathrm{ContOn}(f,[a,b]) \to \varphi_{328}$}
\AXC{$\varphi_{875}$\lbl{A1}}
\BIC{$(\mathrm{ContOn}(f,[a,b]) \to \varphi_{317}) \to \varphi_{854}$}
\BIC{$\varphi_{854}$}
\BIC{$(\mathrm{ContOn}(f,[a,b]) \to \varphi_{316}) \to \varphi_{759}$}
\end{ptree}
\noindent$\Pi_{129}$:
\begin{ptree}
\AXC{$\stackrel{\Pi_{122}}{\vdots}$}\noLine
\UIC{$\mathrm{ContOn}(f,[a,b]) \to \varphi_{356}$}
\AXC{$\stackrel{\Pi_{42}}{\vdots}$}\noLine
\UIC{$\mathrm{ContOn}(f,[a,b]) \to \varphi_{316}$}
\AXC{$\stackrel{\Pi_{128}}{\vdots}$}\noLine
\UIC{$(\mathrm{ContOn}(f,[a,b]) \to \varphi_{316}) \to \varphi_{759}$}
\BIC{$\varphi_{759}$}
\BIC{$(\mathrm{ContOn}(f,[a,b]) \to \varphi_{376}) \to \varphi_{470}$}
\end{ptree}
\noindent$\Pi_{130}$:
\begin{ptree}
\AXC{$\varphi_{252}$\lbl{A2}}
\AXC{$\mathrm{ContOn}(f,[a,b]) \to \varphi_{376}$\lbl{A1}}
\AXC{$\stackrel{\Pi_{129}}{\vdots}$}\noLine
\UIC{$(\mathrm{ContOn}(f,[a,b]) \to \varphi_{376}) \to \varphi_{470}$}
\BIC{$\varphi_{470}$}
\AXC{$\varphi_{470} \to \varphi_{252} \to \varphi_{475}$\lbl{A1}}
\BIC{$\varphi_{252} \to \varphi_{475}$}
\BIC{$\varphi_{475}$}
\end{ptree}
\noindent$\Pi_{131}$:
\begin{ptree}
\AXC{$\forall x\,\forall y\,(x \in \mathbb{R} \to y \in \mathbb{R} \to 0 < y \to x < x + y)$\lbl{[E3]}}
\AXC{$\varphi_{275}$\lbl{A2}}
\BIC{$\forall y\,(\xi \in \mathbb{R} \to y \in \mathbb{R} \to 0 < y \to \xi < \xi + y)$}
\AXC{$\varphi_{266}$\lbl{A2}}
\BIC{$\xi \in \mathbb{R} \to d \in \mathbb{R} \to 0 < d \to \xi < \xi + d$}
\AXC{$(\xi \in \mathbb{R} \to d \in \mathbb{R} \to 0 < d \to \xi < \xi + d) \to \varphi_{871}$\lbl{A1}}
\BIC{$\varphi_{871}$}
\end{ptree}
\noindent$\Pi_{132}$:
\begin{ptree}
\AXC{$\mathrm{ContOn}(f,[a,b]) \to \varphi_{679}$\lbl{A1}}
\AXC{$\varphi_{812}$\lbl{A1}}
\BIC{$\mathrm{ContOn}(f,[a,b]) \to \varphi_{612}$}
\AXC{$\stackrel{\Pi_{3}}{\vdots}$}\noLine
\UIC{$\mathrm{ContOn}(f,[a,b]) \to \varphi_{328}$}
\AXC{$\stackrel{\Pi_{131}}{\vdots}$}\noLine
\UIC{$\varphi_{871}$}
\BIC{$(\mathrm{ContOn}(f,[a,b]) \to \varphi_{612}) \to \varphi_{781}$}
\BIC{$\varphi_{781}$}
\end{ptree}
\noindent$\Pi_{133}$:
\begin{ptree}
\AXC{$z = w \to f(\xi + d) < z \to f(\xi + d) < w$\lbl{A6}}
\UIC{$\forall w\,(z = w \to f(\xi + d) < z \to f(\xi + d) < w)$}
\UIC{$\forall z\,\forall w\,(z = w \to f(\xi + d) < z \to f(\xi + d) < w)$}
\AXC{$\forall z\,\forall w\,(z = w \to f(\xi + d) < z \to f(\xi + d) < w) \to \varphi_{267}$\lbl{A2}}
\BIC{$\varphi_{267}$}
\AXC{$\varphi_{267} \to \varphi_{261}$\lbl{A2}}
\BIC{$\varphi_{261}$}
\end{ptree}
\noindent$\Pi_{134}$:
\begin{ptree}
\AXC{$\varphi_{711}$\lbl{A1}}
\AXC{$\varphi_{711} \to \varphi_{577}$\lbl{A1}}
\BIC{$\varphi_{577}$}
\AXC{$\varphi_{711}$\lbl{A1}}
\AXC{$\varphi_{711} \to \varphi_{579}$\lbl{A1}}
\BIC{$\varphi_{579}$}
\AXC{$\forall e \in \mathbb{R}\,(0 < e \to e \cdot 2^{-1} < e)$\lbl{[E39]}}
\AXC{$\varphi_{138}$\lbl{A2}}
\BIC{$\delta \in \mathbb{R} \to 0 < \delta \to \delta \cdot 2^{-1} < \delta$}
\AXC{$\varphi_{857}$\lbl{A1}}
\BIC{$\varphi_{579} \to \varphi_{577} \to \mathrm{ContOn}(f,[a,b]) \to \varphi_{634}$}
\BIC{$\varphi_{577} \to \mathrm{ContOn}(f,[a,b]) \to \varphi_{634}$}
\BIC{$\mathrm{ContOn}(f,[a,b]) \to \varphi_{634}$}
\end{ptree}
\noindent$\Pi_{135}$:
\begin{ptree}
\AXC{$\forall x\,\forall y\,\forall z\,(x \leq y \to y < z \to x < z)$\lbl{[O4]}}
\AXC{$\varphi_{21}$\lbl{A2}}
\BIC{$\forall y\,\forall z\,(d \leq y \to y < z \to d < z)$}
\AXC{$\varphi_{75}$\lbl{A2}}
\BIC{$\forall z\,(d \leq \delta \cdot 2^{-1} \to \delta \cdot 2^{-1} < z \to d < z)$}
\AXC{$\varphi_{213}$\lbl{A2}}
\BIC{$d \leq \delta \cdot 2^{-1} \to \delta \cdot 2^{-1} < \delta \to d < \delta$}
\AXC{$\varphi_{863}$\lbl{A1}}
\BIC{$(\mathrm{ContOn}(f,[a,b]) \to \varphi_{622}) \to \varphi_{788}$}
\end{ptree}
\noindent$\Pi_{136}$:
\begin{ptree}
\AXC{$\mathrm{ContOn}(f,[a,b]) \to \varphi_{679}$\lbl{A1}}
\AXC{$\varphi_{815}$\lbl{A1}}
\BIC{$\mathrm{ContOn}(f,[a,b]) \to \varphi_{622}$}
\AXC{$\stackrel{\Pi_{135}}{\vdots}$}\noLine
\UIC{$(\mathrm{ContOn}(f,[a,b]) \to \varphi_{622}) \to \varphi_{788}$}
\BIC{$\varphi_{788}$}
\end{ptree}
\noindent$\Pi_{137}$:
\begin{ptree}
\AXC{$y = z \to x = y \to x = z$\lbl{A6}}
\AXC{$(y = z \to x = y \to x = z) \to x = y \to y = z \to x = z$\lbl{A1}}
\BIC{$x = y \to y = z \to x = z$}
\UIC{$\forall z\,(x = y \to y = z \to x = z)$}
\UIC{$\forall y\,\forall z\,(x = y \to y = z \to x = z)$}
\UIC{$\forall x\,\forall y\,\forall z\,(x = y \to y = z \to x = z)$}
\AXC{$\varphi_{109}$\lbl{A2}}
\BIC{$\forall y\,\forall z\,(\lvert \xi + d - \xi \rvert = y \to y = z \to \lvert \xi + d - \xi \rvert = z)$}
\AXC{$\varphi_{293}$\lbl{A2}}
\BIC{$\varphi_{24}$}
\end{ptree}
\noindent$\Pi_{138}$:
\begin{ptree}
\AXC{$\forall x\,(0 \leq x \to \lvert x \rvert = x)$\lbl{[A1]}}
\AXC{$\forall x\,(0 \leq x \to \lvert x \rvert = x) \to 0 \leq d \to \lvert d \rvert = d$\lbl{A2}}
\BIC{$0 \leq d \to \lvert d \rvert = d$}
\AXC{$(0 \leq d \to \lvert d \rvert = d) \to \varphi_{773}$\lbl{A1}}
\BIC{$\varphi_{773}$}
\end{ptree}
\noindent$\Pi_{139}$:
\begin{ptree}
\AXC{$\mathrm{ContOn}(f,[a,b]) \to \varphi_{679}$\lbl{A1}}
\AXC{$\varphi_{809}$\lbl{A1}}
\BIC{$\mathrm{ContOn}(f,[a,b]) \to \varphi_{602}$}
\AXC{$\varphi_{772}$\lbl{A1}}
\BIC{$\mathrm{ContOn}(f,[a,b]) \to \varphi_{598}$}
\AXC{$\stackrel{\Pi_{138}}{\vdots}$}\noLine
\UIC{$\varphi_{773}$}
\BIC{$\mathrm{ContOn}(f,[a,b]) \to \varphi_{605}$}
\end{ptree}
\noindent$\Pi_{140}$:
\begin{ptree}
\AXC{$y = z \to x = y \to x = z$\lbl{A6}}
\AXC{$(y = z \to x = y \to x = z) \to x = y \to y = z \to x = z$\lbl{A1}}
\BIC{$x = y \to y = z \to x = z$}
\UIC{$\forall z\,(x = y \to y = z \to x = z)$}
\UIC{$\forall y\,\forall z\,(x = y \to y = z \to x = z)$}
\UIC{$\forall x\,\forall y\,\forall z\,(x = y \to y = z \to x = z)$}
\AXC{$\varphi_{109}$\lbl{A2}}
\BIC{$\forall y\,\forall z\,(\lvert \xi + d - \xi \rvert = y \to y = z \to \lvert \xi + d - \xi \rvert = z)$}
\AXC{$\varphi_{306}$\lbl{A2}}
\BIC{$\varphi_{148}$}
\AXC{$\varphi_{148} \to \varphi_{206}$\lbl{A2}}
\BIC{$\varphi_{206}$}
\end{ptree}
\noindent$\Pi_{141}$:
\begin{ptree}
\AXC{$z = w \to \lvert d + \xi - \xi \rvert = \lvert z \rvert \to \lvert d + \xi - \xi \rvert = \lvert w \rvert$\lbl{A6}}
\UIC{$\varphi_{25}$}
\UIC{$\forall z\,(\varphi_{25})$}
\AXC{$\forall z\,(\varphi_{25}) \to \varphi_{149}$\lbl{A2}}
\BIC{$\varphi_{149}$}
\AXC{$\varphi_{149} \to \varphi_{122}$\lbl{A2}}
\BIC{$\varphi_{122}$}
\end{ptree}
\noindent$\Pi_{142}$:
\begin{ptree}
\AXC{$\forall x\,\forall y\,(x \in \mathbb{R} \to y \in \mathbb{R} \to x + y - y = x)$\lbl{[E38]}}
\AXC{$\varphi_{182}$\lbl{A2}}
\BIC{$\forall y\,(d \in \mathbb{R} \to y \in \mathbb{R} \to d + y - y = d)$}
\AXC{$\varphi_{154}$\lbl{A2}}
\BIC{$d \in \mathbb{R} \to \xi \in \mathbb{R} \to d + \xi - \xi = d$}
\AXC{$\varphi_{848}$\lbl{A1}}
\BIC{$(\mathrm{ContOn}(f,[a,b]) \to \varphi_{612}) \to \varphi_{741}$}
\end{ptree}
\noindent$\Pi_{143}$:
\begin{ptree}
\AXC{$\mathrm{ContOn}(f,[a,b]) \to \varphi_{679}$\lbl{A1}}
\AXC{$\varphi_{812}$\lbl{A1}}
\BIC{$\mathrm{ContOn}(f,[a,b]) \to \varphi_{612}$}
\AXC{$\stackrel{\Pi_{142}}{\vdots}$}\noLine
\UIC{$(\mathrm{ContOn}(f,[a,b]) \to \varphi_{612}) \to \varphi_{741}$}
\BIC{$\varphi_{741}$}
\end{ptree}
\noindent$\Pi_{144}$:
\begin{ptree}
\AXC{$\stackrel{\Pi_{3}}{\vdots}$}\noLine
\UIC{$\mathrm{ContOn}(f,[a,b]) \to \varphi_{328}$}
\AXC{$\stackrel{\Pi_{143}}{\vdots}$}\noLine
\UIC{$\varphi_{741}$}
\BIC{$\mathrm{ContOn}(f,[a,b]) \to \varphi_{624}$}
\AXC{$(\mathrm{ContOn}(f,[a,b]) \to \varphi_{624}) \to \varphi_{734}$\lbl{A1}}
\BIC{$\varphi_{734}$}
\end{ptree}
\noindent$\Pi_{145}$:
\begin{ptree}
\AXC{$\stackrel{\Pi_{141}}{\vdots}$}\noLine
\UIC{$\varphi_{122}$}
\AXC{$x = x$\lbl{A5}}
\UIC{$\forall x\,(x = x)$}
\AXC{$\forall x\,(x = x) \to \lvert d + \xi - \xi \rvert = \lvert d + \xi - \xi \rvert$\lbl{A2}}
\BIC{$\lvert d + \xi - \xi \rvert = \lvert d + \xi - \xi \rvert$}
\AXC{$\stackrel{\Pi_{144}}{\vdots}$}\noLine
\UIC{$\varphi_{734}$}
\BIC{$\varphi_{122} \to \mathrm{ContOn}(f,[a,b]) \to \varphi_{653}$}
\BIC{$\mathrm{ContOn}(f,[a,b]) \to \varphi_{653}$}
\end{ptree}
\noindent$\Pi_{146}$:
\begin{ptree}
\AXC{$\stackrel{\Pi_{3}}{\vdots}$}\noLine
\UIC{$\mathrm{ContOn}(f,[a,b]) \to \varphi_{328}$}
\AXC{$\forall x\,\forall y\,(x \in \mathbb{R} \to y \in \mathbb{R} \to x + y = y + x)$\lbl{(R3)}}
\AXC{$\varphi_{207}$\lbl{A2}}
\BIC{$\forall y\,(\xi \in \mathbb{R} \to y \in \mathbb{R} \to \xi + y = y + \xi)$}
\AXC{$\varphi_{183}$\lbl{A2}}
\BIC{$\xi \in \mathbb{R} \to d \in \mathbb{R} \to \xi + d = d + \xi$}
\AXC{$\varphi_{849}$\lbl{A1}}
\BIC{$(\mathrm{ContOn}(f,[a,b]) \to \varphi_{328}) \to \varphi_{789}$}
\BIC{$\varphi_{789}$}
\end{ptree}
\noindent$\Pi_{147}$:
\begin{ptree}
\AXC{$\mathrm{ContOn}(f,[a,b]) \to \varphi_{679}$\lbl{A1}}
\AXC{$\varphi_{812}$\lbl{A1}}
\BIC{$\mathrm{ContOn}(f,[a,b]) \to \varphi_{612}$}
\AXC{$\stackrel{\Pi_{146}}{\vdots}$}\noLine
\UIC{$\varphi_{789}$}
\BIC{$\mathrm{ContOn}(f,[a,b]) \to \varphi_{625}$}
\AXC{$(\mathrm{ContOn}(f,[a,b]) \to \varphi_{625}) \to \varphi_{737}$\lbl{A1}}
\BIC{$\varphi_{737}$}
\end{ptree}
\noindent$\Pi_{148}$:
\begin{ptree}
\AXC{$x = x$\lbl{A5}}
\UIC{$\forall x\,(x = x)$}
\AXC{$\forall x\,(x = x) \to \lvert \xi + d - \xi \rvert = \lvert \xi + d - \xi \rvert$\lbl{A2}}
\BIC{$\lvert \xi + d - \xi \rvert = \lvert \xi + d - \xi \rvert$}
\AXC{$\stackrel{\Pi_{147}}{\vdots}$}\noLine
\UIC{$\varphi_{737}$}
\BIC{$(\xi + d = d + \xi \to \varphi_{86}) \to \mathrm{ContOn}(f,[a,b]) \to \varphi_{668}$}
\end{ptree}
\noindent$\Pi_{149}$:
\begin{ptree}
\AXC{$\varphi_{45}$\lbl{A6}}
\UIC{$\forall w\,(\varphi_{45})$}
\UIC{$\forall z\,\forall w\,(\varphi_{45})$}
\AXC{$\forall z\,\forall w\,(\varphi_{45}) \to \varphi_{150}$\lbl{A2}}
\BIC{$\varphi_{150}$}
\AXC{$\varphi_{150} \to \xi + d = d + \xi \to \varphi_{86}$\lbl{A2}}
\BIC{$\xi + d = d + \xi \to \varphi_{86}$}
\AXC{$\stackrel{\Pi_{148}}{\vdots}$}\noLine
\UIC{$(\xi + d = d + \xi \to \varphi_{86}) \to \mathrm{ContOn}(f,[a,b]) \to \varphi_{668}$}
\BIC{$\mathrm{ContOn}(f,[a,b]) \to \varphi_{668}$}
\end{ptree}
\noindent$\Pi_{150}$:
\begin{ptree}
\AXC{$\stackrel{\Pi_{145}}{\vdots}$}\noLine
\UIC{$\mathrm{ContOn}(f,[a,b]) \to \varphi_{653}$}
\AXC{$\stackrel{\Pi_{149}}{\vdots}$}\noLine
\UIC{$\mathrm{ContOn}(f,[a,b]) \to \varphi_{668}$}
\AXC{$(\mathrm{ContOn}(f,[a,b]) \to \varphi_{668}) \to \varphi_{833}$\lbl{A1}}
\BIC{$\varphi_{833}$}
\BIC{$\varphi_{206} \to \mathrm{ContOn}(f,[a,b]) \to \varphi_{654}$}
\end{ptree}
\noindent$\Pi_{151}$:
\begin{ptree}
\AXC{$\stackrel{\Pi_{139}}{\vdots}$}\noLine
\UIC{$\mathrm{ContOn}(f,[a,b]) \to \varphi_{605}$}
\AXC{$\stackrel{\Pi_{140}}{\vdots}$}\noLine
\UIC{$\varphi_{206}$}
\AXC{$\stackrel{\Pi_{150}}{\vdots}$}\noLine
\UIC{$\varphi_{206} \to \mathrm{ContOn}(f,[a,b]) \to \varphi_{654}$}
\BIC{$\mathrm{ContOn}(f,[a,b]) \to \varphi_{654}$}
\AXC{$\varphi_{865}$\lbl{A1}}
\BIC{$(\mathrm{ContOn}(f,[a,b]) \to \varphi_{605}) \to \varphi_{723}$}
\BIC{$\varphi_{723}$}
\end{ptree}
\noindent$\Pi_{152}$:
\begin{ptree}
\AXC{$\stackrel{\Pi_{137}}{\vdots}$}\noLine
\UIC{$\varphi_{24}$}
\AXC{$\varphi_{24} \to \lvert \xi + d - \xi \rvert = \lvert d \rvert \to \lvert d \rvert = d \to \lvert \xi + d - \xi \rvert = d$\lbl{A2}}
\BIC{$\lvert \xi + d - \xi \rvert = \lvert d \rvert \to \lvert d \rvert = d \to \lvert \xi + d - \xi \rvert = d$}
\AXC{$\stackrel{\Pi_{151}}{\vdots}$}\noLine
\UIC{$\varphi_{723}$}
\BIC{$\mathrm{ContOn}(f,[a,b]) \to \varphi_{640}$}
\AXC{$(\mathrm{ContOn}(f,[a,b]) \to \varphi_{640}) \to \varphi_{718}$\lbl{A1}}
\BIC{$\varphi_{718}$}
\end{ptree}
\noindent$\Pi_{153}$:
\begin{ptree}
\AXC{$\stackrel{\Pi_{115}}{\vdots}$}\noLine
\UIC{$\forall x\,\forall y\,(x = y \to y = x)$}
\AXC{$\varphi_{27}$\lbl{A2}}
\BIC{$\forall y\,(\lvert \xi + d - \xi \rvert = y \to y = \lvert \xi + d - \xi \rvert)$}
\AXC{$\varphi_{199}$\lbl{A2}}
\BIC{$\lvert \xi + d - \xi \rvert = d \to d = \lvert \xi + d - \xi \rvert$}
\AXC{$\stackrel{\Pi_{152}}{\vdots}$}\noLine
\UIC{$\varphi_{718}$}
\BIC{$\mathrm{ContOn}(f,[a,b]) \to \varphi_{641}$}
\AXC{$\varphi_{864}$\lbl{A1}}
\BIC{$(\mathrm{ContOn}(f,[a,b]) \to \varphi_{604}) \to \varphi_{724}$}
\end{ptree}
\noindent$\Pi_{154}$:
\begin{ptree}
\AXC{$\stackrel{\Pi_{134}}{\vdots}$}\noLine
\UIC{$\mathrm{ContOn}(f,[a,b]) \to \varphi_{634}$}
\AXC{$\stackrel{\Pi_{136}}{\vdots}$}\noLine
\UIC{$\varphi_{788}$}
\BIC{$\mathrm{ContOn}(f,[a,b]) \to \varphi_{604}$}
\AXC{$\stackrel{\Pi_{153}}{\vdots}$}\noLine
\UIC{$(\mathrm{ContOn}(f,[a,b]) \to \varphi_{604}) \to \varphi_{724}$}
\BIC{$\varphi_{724}$}
\end{ptree}
\noindent$\Pi_{155}$:
\begin{ptree}
\AXC{$\forall x\,\forall y\,(x \in \mathbb{R} \to y \in \mathbb{R} \to x - y \in \mathbb{R})$\lbl{[F1]}}
\AXC{$\varphi_{215}$\lbl{A2}}
\BIC{$\forall y\,(b \in \mathbb{R} \to y \in \mathbb{R} \to b - y \in \mathbb{R})$}
\AXC{$\varphi_{185}$\lbl{A2}}
\BIC{$b \in \mathbb{R} \to \xi \in \mathbb{R} \to b - \xi \in \mathbb{R}$}
\AXC{$(b \in \mathbb{R} \to \xi \in \mathbb{R} \to b - \xi \in \mathbb{R}) \to \varphi_{758}$\lbl{A1}}
\BIC{$\varphi_{758}$}
\end{ptree}
\noindent$\Pi_{156}$:
\begin{ptree}
\AXC{$\stackrel{\Pi_{3}}{\vdots}$}\noLine
\UIC{$\mathrm{ContOn}(f,[a,b]) \to \varphi_{328}$}
\AXC{$\mathrm{ContOn}(f,[a,b]) \to a \in \mathbb{R} \to b \in \mathbb{R} \to b \in \mathbb{R}$\lbl{A1}}
\AXC{$\stackrel{\Pi_{155}}{\vdots}$}\noLine
\UIC{$\varphi_{758}$}
\BIC{$\varphi_{744}$}
\BIC{$\mathrm{ContOn}(f,[a,b]) \to \varphi_{629}$}
\end{ptree}
\noindent$\Pi_{157}$:
\begin{ptree}
\AXC{$\varphi_{39}$\lbl{[F19]}}
\AXC{$\varphi_{39} \to \varphi_{28}$\lbl{A2}}
\BIC{$\varphi_{28}$}
\AXC{$\varphi_{28} \to \varphi_{96}$\lbl{A2}}
\BIC{$\varphi_{96}$}
\AXC{$\varphi_{96} \to \varphi_{97}$\lbl{A2}}
\BIC{$\varphi_{97}$}
\AXC{$\varphi_{97} \to (\mathrm{ContOn}(f,[a,b]) \to \varphi_{612}) \to \varphi_{872}$\lbl{A1}}
\BIC{$(\mathrm{ContOn}(f,[a,b]) \to \varphi_{612}) \to \varphi_{872}$}
\end{ptree}
\noindent$\Pi_{158}$:
\begin{ptree}
\AXC{$\mathrm{ContOn}(f,[a,b]) \to \varphi_{679}$\lbl{A1}}
\AXC{$\varphi_{812}$\lbl{A1}}
\BIC{$\mathrm{ContOn}(f,[a,b]) \to \varphi_{612}$}
\AXC{$\stackrel{\Pi_{157}}{\vdots}$}\noLine
\UIC{$(\mathrm{ContOn}(f,[a,b]) \to \varphi_{612}) \to \varphi_{872}$}
\BIC{$\varphi_{872}$}
\end{ptree}
\noindent$\Pi_{159}$:
\begin{ptree}
\AXC{$\stackrel{\Pi_{3}}{\vdots}$}\noLine
\UIC{$\mathrm{ContOn}(f,[a,b]) \to \varphi_{328}$}
\AXC{$\stackrel{\Pi_{156}}{\vdots}$}\noLine
\UIC{$\mathrm{ContOn}(f,[a,b]) \to \varphi_{629}$}
\AXC{$\stackrel{\Pi_{158}}{\vdots}$}\noLine
\UIC{$\varphi_{872}$}
\BIC{$(\mathrm{ContOn}(f,[a,b]) \to \varphi_{328}) \to \varphi_{792}$}
\BIC{$\varphi_{792}$}
\end{ptree}
\noindent$\Pi_{160}$:
\begin{ptree}
\AXC{$\forall x\,\forall y\,(x \in \mathbb{R} \to y \in \mathbb{R} \to x - y + y = x)$\lbl{[F12]}}
\AXC{$\varphi_{186}$\lbl{A2}}
\BIC{$\forall y\,(b \in \mathbb{R} \to y \in \mathbb{R} \to b - y + y = b)$}
\AXC{$\varphi_{155}$\lbl{A2}}
\BIC{$b \in \mathbb{R} \to \xi \in \mathbb{R} \to b - \xi + \xi = b$}
\AXC{$(b \in \mathbb{R} \to \xi \in \mathbb{R} \to b - \xi + \xi = b) \to \varphi_{757}$\lbl{A1}}
\BIC{$\varphi_{757}$}
\end{ptree}
\noindent$\Pi_{161}$:
\begin{ptree}
\AXC{$\stackrel{\Pi_{3}}{\vdots}$}\noLine
\UIC{$\mathrm{ContOn}(f,[a,b]) \to \varphi_{328}$}
\AXC{$\mathrm{ContOn}(f,[a,b]) \to a \in \mathbb{R} \to b \in \mathbb{R} \to b \in \mathbb{R}$\lbl{A1}}
\AXC{$\stackrel{\Pi_{160}}{\vdots}$}\noLine
\UIC{$\varphi_{757}$}
\BIC{$\varphi_{742}$}
\BIC{$\mathrm{ContOn}(f,[a,b]) \to \varphi_{626}$}
\end{ptree}
\noindent$\Pi_{162}$:
\begin{ptree}
\AXC{$\stackrel{\Pi_{161}}{\vdots}$}\noLine
\UIC{$\mathrm{ContOn}(f,[a,b]) \to \varphi_{626}$}
\AXC{$(\mathrm{ContOn}(f,[a,b]) \to \varphi_{626}) \to \varphi_{832}$\lbl{A1}}
\BIC{$\varphi_{832}$}
\end{ptree}
\noindent$\Pi_{163}$:
\begin{ptree}
\AXC{$\mathrm{ContOn}(f,[a,b]) \to \varphi_{679}$\lbl{A1}}
\AXC{$\varphi_{810}$\lbl{A1}}
\BIC{$\mathrm{ContOn}(f,[a,b]) \to \varphi_{606}$}
\AXC{$\stackrel{\Pi_{159}}{\vdots}$}\noLine
\UIC{$\varphi_{792}$}
\BIC{$\mathrm{ContOn}(f,[a,b]) \to \varphi_{643}$}
\AXC{$\stackrel{\Pi_{162}}{\vdots}$}\noLine
\UIC{$\varphi_{832}$}
\BIC{$(b - \xi + \xi = b \to d + \xi < b - \xi + \xi \to d + \xi < b) \to \varphi_{713}$}
\end{ptree}
\noindent$\Pi_{164}$:
\begin{ptree}
\AXC{$z = w \to d + \xi < z \to d + \xi < w$\lbl{A6}}
\UIC{$\forall w\,(z = w \to d + \xi < z \to d + \xi < w)$}
\UIC{$\forall z\,\forall w\,(z = w \to d + \xi < z \to d + \xi < w)$}
\AXC{$\varphi_{258}$\lbl{A2}}
\BIC{$\forall w\,(b - \xi + \xi = w \to d + \xi < b - \xi + \xi \to d + \xi < w)$}
\AXC{$\varphi_{294}$\lbl{A2}}
\BIC{$b - \xi + \xi = b \to d + \xi < b - \xi + \xi \to d + \xi < b$}
\AXC{$\stackrel{\Pi_{163}}{\vdots}$}\noLine
\UIC{$(b - \xi + \xi = b \to d + \xi < b - \xi + \xi \to d + \xi < b) \to \varphi_{713}$}
\BIC{$\varphi_{713}$}
\end{ptree}
\noindent$\Pi_{165}$:
\begin{ptree}
\AXC{$\forall x\,\forall y\,(x \in \mathbb{R} \to y \in \mathbb{R} \to x + y = y + x)$\lbl{(R3)}}
\AXC{$\varphi_{187}$\lbl{A2}}
\BIC{$\forall y\,(d \in \mathbb{R} \to y \in \mathbb{R} \to d + y = y + d)$}
\AXC{$\varphi_{156}$\lbl{A2}}
\BIC{$d \in \mathbb{R} \to \xi \in \mathbb{R} \to d + \xi = \xi + d$}
\AXC{$\varphi_{850}$\lbl{A1}}
\BIC{$(\mathrm{ContOn}(f,[a,b]) \to \varphi_{612}) \to \varphi_{743}$}
\end{ptree}
\noindent$\Pi_{166}$:
\begin{ptree}
\AXC{$\mathrm{ContOn}(f,[a,b]) \to \varphi_{679}$\lbl{A1}}
\AXC{$\varphi_{812}$\lbl{A1}}
\BIC{$\mathrm{ContOn}(f,[a,b]) \to \varphi_{612}$}
\AXC{$\stackrel{\Pi_{165}}{\vdots}$}\noLine
\UIC{$(\mathrm{ContOn}(f,[a,b]) \to \varphi_{612}) \to \varphi_{743}$}
\BIC{$\varphi_{743}$}
\end{ptree}
\noindent$\Pi_{167}$:
\begin{ptree}
\AXC{$\stackrel{\Pi_{3}}{\vdots}$}\noLine
\UIC{$\mathrm{ContOn}(f,[a,b]) \to \varphi_{328}$}
\AXC{$\stackrel{\Pi_{166}}{\vdots}$}\noLine
\UIC{$\varphi_{743}$}
\BIC{$\mathrm{ContOn}(f,[a,b]) \to \varphi_{627}$}
\AXC{$(\mathrm{ContOn}(f,[a,b]) \to \varphi_{627}) \to \varphi_{829}$\lbl{A1}}
\BIC{$\varphi_{829}$}
\end{ptree}
\noindent$\Pi_{168}$:
\begin{ptree}
\AXC{$z = w \to d + \xi = z \to d + \xi = w$\lbl{A6}}
\UIC{$\forall w\,(z = w \to d + \xi = z \to d + \xi = w)$}
\UIC{$\forall z\,\forall w\,(z = w \to d + \xi = z \to d + \xi = w)$}
\AXC{$\varphi_{184}$\lbl{A2}}
\BIC{$\forall w\,(b - \xi + \xi = w \to d + \xi = b - \xi + \xi \to d + \xi = w)$}
\AXC{$\varphi_{279}$\lbl{A2}}
\BIC{$b - \xi + \xi = b \to d + \xi = b - \xi + \xi \to d + \xi = b$}
\AXC{$\stackrel{\Pi_{164}}{\vdots}$}\noLine
\UIC{$\varphi_{713}$}
\BIC{$\mathrm{ContOn}(f,[a,b]) \to \varphi_{607}$}
\AXC{$\stackrel{\Pi_{167}}{\vdots}$}\noLine
\UIC{$\varphi_{829}$}
\BIC{$(d + \xi = \xi + d \to d + \xi < b \to \xi + d < b) \to \varphi_{700}$}
\end{ptree}
\noindent$\Pi_{169}$:
\begin{ptree}
\AXC{$z = w \to z < b \to w < b$\lbl{A6}}
\UIC{$\forall w\,(z = w \to z < b \to w < b)$}
\UIC{$\forall z\,\forall w\,(z = w \to z < b \to w < b)$}
\AXC{$\varphi_{33}$\lbl{A2}}
\BIC{$\forall w\,(d + \xi = w \to d + \xi < b \to w < b)$}
\AXC{$\varphi_{128}$\lbl{A2}}
\BIC{$d + \xi = \xi + d \to d + \xi < b \to \xi + d < b$}
\AXC{$\stackrel{\Pi_{168}}{\vdots}$}\noLine
\UIC{$(d + \xi = \xi + d \to d + \xi < b \to \xi + d < b) \to \varphi_{700}$}
\BIC{$\varphi_{700}$}
\end{ptree}
\noindent$\Pi_{170}$:
\begin{ptree}
\AXC{$\forall x\,\forall y\,(x \in \mathbb{R} \to y \in \mathbb{R} \to 0 \leq y \to x \leq x + y)$\lbl{[E1]}}
\AXC{$\varphi_{219}$\lbl{A2}}
\BIC{$\forall y\,(\xi \in \mathbb{R} \to y \in \mathbb{R} \to 0 \leq y \to \xi \leq \xi + y)$}
\AXC{$\varphi_{200}$\lbl{A2}}
\BIC{$\xi \in \mathbb{R} \to d \in \mathbb{R} \to 0 \leq d \to \xi \leq \xi + d$}
\AXC{$(\xi \in \mathbb{R} \to d \in \mathbb{R} \to 0 \leq d \to \xi \leq \xi + d) \to \varphi_{869}$\lbl{A1}}
\BIC{$\varphi_{869}$}
\end{ptree}
\noindent$\Pi_{171}$:
\begin{ptree}
\AXC{$\mathrm{ContOn}(f,[a,b]) \to \varphi_{679}$\lbl{A1}}
\AXC{$\varphi_{812}$\lbl{A1}}
\BIC{$\mathrm{ContOn}(f,[a,b]) \to \varphi_{612}$}
\AXC{$\stackrel{\Pi_{3}}{\vdots}$}\noLine
\UIC{$\mathrm{ContOn}(f,[a,b]) \to \varphi_{328}$}
\AXC{$\stackrel{\Pi_{170}}{\vdots}$}\noLine
\UIC{$\varphi_{869}$}
\BIC{$(\mathrm{ContOn}(f,[a,b]) \to \varphi_{612}) \to \varphi_{775}$}
\BIC{$\varphi_{775}$}
\end{ptree}
\noindent$\Pi_{172}$:
\begin{ptree}
\AXC{$\forall x\,\forall y\,\forall z\,(x \leq y \to y \leq z \to x \leq z)$\lbl{[O5]}}
\AXC{$\varphi_{10}$\lbl{A2}}
\BIC{$\forall y\,\forall z\,(a \leq y \to y \leq z \to a \leq z)$}
\AXC{$\varphi_{6}$\lbl{A2}}
\BIC{$\forall z\,(a \leq \xi \to \xi \leq z \to a \leq z)$}
\AXC{$\varphi_{12}$\lbl{A2}}
\BIC{$a \leq \xi \to \xi \leq \xi + d \to a \leq \xi + d$}
\AXC{$\varphi_{843}$\lbl{A1}}
\BIC{$(\mathrm{ContOn}(f,[a,b]) \to \varphi_{317}) \to \varphi_{782}$}
\end{ptree}
\noindent$\Pi_{173}$:
\begin{ptree}
\AXC{$\stackrel{\Pi_{50}}{\vdots}$}\noLine
\UIC{$\mathrm{ContOn}(f,[a,b]) \to \varphi_{317}$}
\AXC{$\stackrel{\Pi_{172}}{\vdots}$}\noLine
\UIC{$(\mathrm{ContOn}(f,[a,b]) \to \varphi_{317}) \to \varphi_{782}$}
\BIC{$\varphi_{782}$}
\end{ptree}
\noindent$\Pi_{174}$:
\begin{ptree}
\AXC{$\mathrm{ContOn}(f,[a,b]) \to \varphi_{679}$\lbl{A1}}
\AXC{$\varphi_{809}$\lbl{A1}}
\BIC{$\mathrm{ContOn}(f,[a,b]) \to \varphi_{602}$}
\AXC{$\varphi_{772}$\lbl{A1}}
\BIC{$\mathrm{ContOn}(f,[a,b]) \to \varphi_{598}$}
\AXC{$\stackrel{\Pi_{171}}{\vdots}$}\noLine
\UIC{$\varphi_{775}$}
\BIC{$\mathrm{ContOn}(f,[a,b]) \to \varphi_{613}$}
\AXC{$\stackrel{\Pi_{173}}{\vdots}$}\noLine
\UIC{$\varphi_{782}$}
\BIC{$\mathrm{ContOn}(f,[a,b]) \to \varphi_{609}$}
\end{ptree}
\noindent$\Pi_{175}$:
\begin{ptree}
\AXC{$\stackrel{\Pi_{3}}{\vdots}$}\noLine
\UIC{$\mathrm{ContOn}(f,[a,b]) \to \varphi_{328}$}
\AXC{$\forall x\,\forall y\,(x \in \mathbb{R} \to y \in \mathbb{R} \to x + y \in \mathbb{R})$\lbl{(R1)}}
\AXC{$\varphi_{229}$\lbl{A2}}
\BIC{$\forall y\,(\xi \in \mathbb{R} \to y \in \mathbb{R} \to \xi + y \in \mathbb{R})$}
\AXC{$\varphi_{201}$\lbl{A2}}
\BIC{$\xi \in \mathbb{R} \to d \in \mathbb{R} \to \xi + d \in \mathbb{R}$}
\AXC{$\varphi_{852}$\lbl{A1}}
\BIC{$(\mathrm{ContOn}(f,[a,b]) \to \varphi_{328}) \to \varphi_{791}$}
\BIC{$\varphi_{791}$}
\end{ptree}
\noindent$\Pi_{176}$:
\begin{ptree}
\AXC{$\forall a\,\forall b\,\forall x\,(x \in [a,b] \leftrightarrow x \in \mathbb{R} \land a \leq x \land x \leq b)$\lbl{Def{$[a,b]$}}}
\AXC{$\varphi_{236}$\lbl{A2}}
\BIC{$\forall b\,\forall x\,(x \in [a,b] \leftrightarrow x \in \mathbb{R} \land a \leq x \land x \leq b)$}
\AXC{$\varphi_{195}$\lbl{A2}}
\BIC{$\forall x\,(x \in [a,b] \leftrightarrow x \in \mathbb{R} \land a \leq x \land x \leq b)$}
\AXC{$\varphi_{284}$\lbl{A2}}
\BIC{$\xi + d \in [a,b] \leftrightarrow \xi + d \in \mathbb{R} \land a \leq \xi + d \land \xi + d \leq b$}
\AXC{$\varphi_{880}$\lbl{A1}}
\BIC{$\varphi_{877}$}
\end{ptree}
\noindent$\Pi_{177}$:
\begin{ptree}
\AXC{$\stackrel{\Pi_{174}}{\vdots}$}\noLine
\UIC{$\mathrm{ContOn}(f,[a,b]) \to \varphi_{609}$}
\AXC{$\mathrm{ContOn}(f,[a,b]) \to \varphi_{679}$\lbl{A1}}
\AXC{$\varphi_{812}$\lbl{A1}}
\BIC{$\mathrm{ContOn}(f,[a,b]) \to \varphi_{612}$}
\AXC{$\stackrel{\Pi_{175}}{\vdots}$}\noLine
\UIC{$\varphi_{791}$}
\BIC{$\mathrm{ContOn}(f,[a,b]) \to \varphi_{630}$}
\AXC{$\stackrel{\Pi_{176}}{\vdots}$}\noLine
\UIC{$\varphi_{877}$}
\BIC{$(\mathrm{ContOn}(f,[a,b]) \to \varphi_{609}) \to \varphi_{793}$}
\BIC{$\varphi_{793}$}
\end{ptree}
\noindent$\Pi_{178}$:
\begin{ptree}
\AXC{$z = w \to z = b \to w = b$\lbl{A6}}
\UIC{$\forall w\,(z = w \to z = b \to w = b)$}
\UIC{$\forall z\,\forall w\,(z = w \to z = b \to w = b)$}
\AXC{$\varphi_{13}$\lbl{A2}}
\BIC{$\forall w\,(d + \xi = w \to d + \xi = b \to w = b)$}
\AXC{$\varphi_{51}$\lbl{A2}}
\BIC{$d + \xi = \xi + d \to d + \xi = b \to \xi + d = b$}
\AXC{$\stackrel{\Pi_{169}}{\vdots}$}\noLine
\UIC{$\varphi_{700}$}
\BIC{$\mathrm{ContOn}(f,[a,b]) \to \varphi_{608}$}
\AXC{$\stackrel{\Pi_{177}}{\vdots}$}\noLine
\UIC{$\varphi_{793}$}
\BIC{$\mathrm{ContOn}(f,[a,b]) \to \varphi_{644}$}
\end{ptree}
\noindent$\Pi_{179}$:
\begin{ptree}
\AXC{$\varphi_{159} \to \varphi_{259}$\lbl{A2}}
\AXC{$\varphi_{711}$\lbl{A1}}
\AXC{$\varphi_{711} \to \mathrm{ContOn}(f,[a,b]) \to \varphi_{717}$\lbl{A1}}
\BIC{$\mathrm{ContOn}(f,[a,b]) \to \varphi_{717}$}
\AXC{$\varphi_{859}$\lbl{A1}}
\BIC{$(\varphi_{159} \to \varphi_{259}) \to \mathrm{ContOn}(f,[a,b]) \to \varphi_{722}$}
\BIC{$\mathrm{ContOn}(f,[a,b]) \to \varphi_{722}$}
\AXC{$\varphi_{882}$\lbl{A1}}
\BIC{$(\mathrm{ContOn}(f,[a,b]) \to \varphi_{644}) \to \varphi_{822}$}
\end{ptree}
\noindent$\Pi_{180}$:
\begin{ptree}
\AXC{$\stackrel{\Pi_{178}}{\vdots}$}\noLine
\UIC{$\mathrm{ContOn}(f,[a,b]) \to \varphi_{644}$}
\AXC{$\stackrel{\Pi_{179}}{\vdots}$}\noLine
\UIC{$(\mathrm{ContOn}(f,[a,b]) \to \varphi_{644}) \to \varphi_{822}$}
\BIC{$\varphi_{822}$}
\end{ptree}
\noindent$\Pi_{181}$:
\begin{ptree}
\AXC{$z = w \to z < \delta \to w < \delta$\lbl{A6}}
\UIC{$\forall w\,(z = w \to z < \delta \to w < \delta)$}
\UIC{$\forall z\,\forall w\,(z = w \to z < \delta \to w < \delta)$}
\AXC{$\varphi_{17}$\lbl{A2}}
\BIC{$\forall w\,(d = w \to d < \delta \to w < \delta)$}
\AXC{$\varphi_{147}$\lbl{A2}}
\BIC{$d = \lvert \xi + d - \xi \rvert \to d < \delta \to \lvert \xi + d - \xi \rvert < \delta$}
\AXC{$\stackrel{\Pi_{154}}{\vdots}$}\noLine
\UIC{$\varphi_{724}$}
\BIC{$\mathrm{ContOn}(f,[a,b]) \to \varphi_{650}$}
\AXC{$\stackrel{\Pi_{180}}{\vdots}$}\noLine
\UIC{$\varphi_{822}$}
\BIC{$\mathrm{ContOn}(f,[a,b]) \to \varphi_{682}$}
\end{ptree}
\noindent$\Pi_{182}$:
\begin{ptree}
\AXC{$\forall x\,\forall y\,(x \in \mathbb{R} \to y \in \mathbb{R} \to x - y \in \mathbb{R})$\lbl{[F1]}}
\AXC{$\varphi_{230}$\lbl{A2}}
\BIC{$\forall y\,(\mu \in \mathbb{R} \to y \in \mathbb{R} \to \mu - y \in \mathbb{R})$}
\AXC{$\varphi_{280}$\lbl{A2}}
\BIC{$\mu \in \mathbb{R} \to f(\xi) \in \mathbb{R} \to \mu - f(\xi) \in \mathbb{R}$}
\AXC{$(\mu \in \mathbb{R} \to f(\xi) \in \mathbb{R} \to \mu - f(\xi) \in \mathbb{R}) \to \varphi_{56} \to \varphi_{749}$\lbl{A1}}
\BIC{$\varphi_{56} \to \varphi_{749}$}
\end{ptree}
\noindent$\Pi_{183}$:
\begin{ptree}
\AXC{$\stackrel{\Pi_{60}}{\vdots}$}\noLine
\UIC{$\mathrm{ContOn}(f,[a,b]) \to \varphi_{347}$}
\AXC{$\varphi_{56}$\lbl{A1}}
\AXC{$\stackrel{\Pi_{182}}{\vdots}$}\noLine
\UIC{$\varphi_{56} \to \varphi_{749}$}
\BIC{$\varphi_{749}$}
\BIC{$\mathrm{ContOn}(f,[a,b]) \to \varphi_{651}$}
\end{ptree}
\noindent$\Pi_{184}$:
\begin{ptree}
\AXC{$\forall X\,\forall Y\,\forall x\,(X \subseteq Y \to x \in X \to x \in Y)$\lbl{[B5]}}
\AXC{$\varphi_{180}$\lbl{A2}}
\BIC{$\forall Y\,\forall x\,([a,b] \subseteq Y \to x \in [a,b] \to x \in Y)$}
\AXC{$\forall Y\,\forall x\,([a,b] \subseteq Y \to x \in [a,b] \to x \in Y) \to \varphi_{15}$\lbl{A2}}
\BIC{$\varphi_{15}$}
\AXC{$\varphi_{15} \to \varphi_{34}$\lbl{A2}}
\BIC{$\varphi_{34}$}
\end{ptree}
\noindent$\Pi_{185}$:
\begin{ptree}
\AXC{$\stackrel{\Pi_{178}}{\vdots}$}\noLine
\UIC{$\mathrm{ContOn}(f,[a,b]) \to \varphi_{644}$}
\AXC{$\stackrel{\Pi_{54}}{\vdots}$}\noLine
\UIC{$\mathrm{ContOn}(f,[a,b]) \to a \in \mathbb{R} \to b \in \mathbb{R} \to \varphi_{92}$}
\AXC{$\stackrel{\Pi_{184}}{\vdots}$}\noLine
\UIC{$\varphi_{34}$}
\AXC{$\varphi_{34} \to \varphi_{836}$\lbl{A1}}
\BIC{$\varphi_{836}$}
\BIC{$\varphi_{797}$}
\BIC{$\mathrm{ContOn}(f,[a,b]) \to \varphi_{635}$}
\end{ptree}
\noindent$\Pi_{186}$:
\begin{ptree}
\AXC{$\forall f\,\forall x\,(\mathrm{RFn}(f) \to x \in \mathrm{dom}(f) \to f(x) \in \mathbb{R})$\lbl{[S6]}}
\AXC{$\varphi_{214}$\lbl{A2}}
\BIC{$\forall x\,(\mathrm{RFn}(f) \to x \in \mathrm{dom}(f) \to f(x) \in \mathbb{R})$}
\AXC{$\varphi_{246}$\lbl{A2}}
\BIC{$\mathrm{RFn}(f) \to \xi + d \in \mathrm{dom}(f) \to f(\xi + d) \in \mathbb{R}$}
\AXC{$\varphi_{842}$\lbl{A1}}
\BIC{$(\mathrm{ContOn}(f,[a,b]) \to a \in \mathbb{R} \to \varphi_{80}) \to \varphi_{798}$}
\end{ptree}
\noindent$\Pi_{187}$:
\begin{ptree}
\AXC{$\stackrel{\Pi_{57}}{\vdots}$}\noLine
\UIC{$\mathrm{ContOn}(f,[a,b]) \to a \in \mathbb{R} \to \varphi_{80}$}
\AXC{$\stackrel{\Pi_{186}}{\vdots}$}\noLine
\UIC{$(\mathrm{ContOn}(f,[a,b]) \to a \in \mathbb{R} \to \varphi_{80}) \to \varphi_{798}$}
\BIC{$\varphi_{798}$}
\end{ptree}
\noindent$\Pi_{188}$:
\begin{ptree}
\AXC{$\stackrel{\Pi_{185}}{\vdots}$}\noLine
\UIC{$\mathrm{ContOn}(f,[a,b]) \to \varphi_{635}$}
\AXC{$\stackrel{\Pi_{187}}{\vdots}$}\noLine
\UIC{$\varphi_{798}$}
\BIC{$\mathrm{ContOn}(f,[a,b]) \to \varphi_{648}$}
\end{ptree}
\noindent$\Pi_{189}$:
\begin{ptree}
\AXC{$\forall x\,\forall y\,(x \in \mathbb{R} \to y \in \mathbb{R} \to x - y \in \mathbb{R})$\lbl{[F1]}}
\AXC{$\varphi_{295}$\lbl{A2}}
\BIC{$\forall y\,(f(\xi + d) \in \mathbb{R} \to y \in \mathbb{R} \to f(\xi + d) - y \in \mathbb{R})$}
\AXC{$\varphi_{325}$\lbl{A2}}
\BIC{$\varphi_{123}$}
\AXC{$\varphi_{123} \to (\mathrm{ContOn}(f,[a,b]) \to \varphi_{648}) \to \varphi_{753}$\lbl{A1}}
\BIC{$(\mathrm{ContOn}(f,[a,b]) \to \varphi_{648}) \to \varphi_{753}$}
\end{ptree}
\noindent$\Pi_{190}$:
\begin{ptree}
\AXC{$\stackrel{\Pi_{60}}{\vdots}$}\noLine
\UIC{$\mathrm{ContOn}(f,[a,b]) \to \varphi_{347}$}
\AXC{$\stackrel{\Pi_{188}}{\vdots}$}\noLine
\UIC{$\mathrm{ContOn}(f,[a,b]) \to \varphi_{648}$}
\AXC{$\stackrel{\Pi_{189}}{\vdots}$}\noLine
\UIC{$(\mathrm{ContOn}(f,[a,b]) \to \varphi_{648}) \to \varphi_{753}$}
\BIC{$\varphi_{753}$}
\BIC{$\mathrm{ContOn}(f,[a,b]) \to \varphi_{665}$}
\end{ptree}
\noindent$\Pi_{191}$:
\begin{ptree}
\AXC{$\forall x\,\forall y\,(x \in \mathbb{R} \to y \in \mathbb{R} \to \lvert x \rvert < y \to -y < x \land x < y)$\lbl{[A11]}}
\AXC{$\varphi_{480}$\lbl{A2}}
\BIC{$\varphi_{361}$}
\AXC{$\varphi_{361} \to f(\xi + d) - f(\xi) \in \mathbb{R} \to \varphi_{427}$\lbl{A2}}
\BIC{$f(\xi + d) - f(\xi) \in \mathbb{R} \to \varphi_{427}$}
\end{ptree}
\noindent$\Pi_{192}$:
\begin{ptree}
\AXC{$\stackrel{\Pi_{190}}{\vdots}$}\noLine
\UIC{$\mathrm{ContOn}(f,[a,b]) \to \varphi_{665}$}
\AXC{$\stackrel{\Pi_{191}}{\vdots}$}\noLine
\UIC{$f(\xi + d) - f(\xi) \in \mathbb{R} \to \varphi_{427}$}
\AXC{$(f(\xi + d) - f(\xi) \in \mathbb{R} \to \varphi_{427}) \to \varphi_{883}$\lbl{A1}}
\BIC{$\varphi_{883}$}
\BIC{$(\mathrm{ContOn}(f,[a,b]) \to \varphi_{651}) \to \varphi_{835}$}
\end{ptree}
\noindent$\Pi_{193}$:
\begin{ptree}
\AXC{$\stackrel{\Pi_{181}}{\vdots}$}\noLine
\UIC{$\mathrm{ContOn}(f,[a,b]) \to \varphi_{682}$}
\AXC{$\stackrel{\Pi_{183}}{\vdots}$}\noLine
\UIC{$\mathrm{ContOn}(f,[a,b]) \to \varphi_{651}$}
\AXC{$\stackrel{\Pi_{192}}{\vdots}$}\noLine
\UIC{$(\mathrm{ContOn}(f,[a,b]) \to \varphi_{651}) \to \varphi_{835}$}
\BIC{$\varphi_{835}$}
\BIC{$\mathrm{ContOn}(f,[a,b]) \to \varphi_{720}$}
\end{ptree}
\noindent$\Pi_{194}$:
\begin{ptree}
\AXC{$\stackrel{\Pi_{188}}{\vdots}$}\noLine
\UIC{$\mathrm{ContOn}(f,[a,b]) \to \varphi_{648}$}
\AXC{$\varphi_{57}$\lbl{[B7]}}
\AXC{$\varphi_{57} \to \varphi_{282}$\lbl{A2}}
\BIC{$\varphi_{282}$}
\AXC{$\varphi_{282} \to \varphi_{300}$\lbl{A2}}
\BIC{$\varphi_{300}$}
\AXC{$\varphi_{575}$\lbl{A2}}
\BIC{$f(\xi + d) \in \mathbb{R} \to f(\xi) \in \mathbb{R} \to \mu - f(\xi) \in \mathbb{R} \to \varphi_{224}$}
\AXC{$\varphi_{892}$\lbl{A1}}
\BIC{$(\mathrm{ContOn}(f,[a,b]) \to \varphi_{648}) \to \varphi_{876}$}
\BIC{$\varphi_{876}$}
\end{ptree}
\noindent$\Pi_{195}$:
\begin{ptree}
\AXC{$\stackrel{\Pi_{183}}{\vdots}$}\noLine
\UIC{$\mathrm{ContOn}(f,[a,b]) \to \varphi_{651}$}
\AXC{$\stackrel{\Pi_{60}}{\vdots}$}\noLine
\UIC{$\mathrm{ContOn}(f,[a,b]) \to \varphi_{347}$}
\AXC{$\stackrel{\Pi_{194}}{\vdots}$}\noLine
\UIC{$\varphi_{876}$}
\BIC{$(\mathrm{ContOn}(f,[a,b]) \to \varphi_{651}) \to \varphi_{824}$}
\BIC{$\varphi_{824}$}
\end{ptree}
\noindent$\Pi_{196}$:
\begin{ptree}
\AXC{$\varphi_{56}$\lbl{A1}}
\AXC{$\forall x\,\forall y\,(x \in \mathbb{R} \to y \in \mathbb{R} \to x - y + y = x)$\lbl{[F12]}}
\AXC{$\varphi_{208}$\lbl{A2}}
\BIC{$\forall y\,(\mu \in \mathbb{R} \to y \in \mathbb{R} \to \mu - y + y = \mu)$}
\AXC{$\varphi_{287}$\lbl{A2}}
\BIC{$\mu \in \mathbb{R} \to f(\xi) \in \mathbb{R} \to \mu - f(\xi) + f(\xi) = \mu$}
\AXC{$\varphi_{774}$\lbl{A1}}
\BIC{$\varphi_{56} \to \varphi_{752}$}
\BIC{$\varphi_{752}$}
\end{ptree}
\noindent$\Pi_{197}$:
\begin{ptree}
\AXC{$\stackrel{\Pi_{60}}{\vdots}$}\noLine
\UIC{$\mathrm{ContOn}(f,[a,b]) \to \varphi_{347}$}
\AXC{$\stackrel{\Pi_{196}}{\vdots}$}\noLine
\UIC{$\varphi_{752}$}
\BIC{$\mathrm{ContOn}(f,[a,b]) \to \varphi_{660}$}
\AXC{$(\mathrm{ContOn}(f,[a,b]) \to \varphi_{660}) \to \varphi_{837}$\lbl{A1}}
\BIC{$\varphi_{837}$}
\end{ptree}
\noindent$\Pi_{198}$:
\begin{ptree}
\AXC{$\stackrel{\Pi_{178}}{\vdots}$}\noLine
\UIC{$\mathrm{ContOn}(f,[a,b]) \to \varphi_{644}$}
\AXC{$\varphi_{281}$\lbl{A2}}
\AXC{$\mathrm{ContOn}(f,[a,b]) \to \varphi_{397}$\lbl{A1}}
\AXC{$\varphi_{738}$\lbl{A1}}
\BIC{$\varphi_{281} \to \mathrm{ContOn}(f,[a,b]) \to \varphi_{442}$}
\BIC{$\mathrm{ContOn}(f,[a,b]) \to \varphi_{442}$}
\AXC{$\varphi_{873}$\lbl{A1}}
\BIC{$(\mathrm{ContOn}(f,[a,b]) \to \varphi_{644}) \to \varphi_{795}$}
\BIC{$\varphi_{795}$}
\end{ptree}
\noindent$\Pi_{199}$:
\begin{ptree}
\AXC{$\stackrel{\Pi_{133}}{\vdots}$}\noLine
\UIC{$\varphi_{261}$}
\AXC{$\stackrel{\Pi_{193}}{\vdots}$}\noLine
\UIC{$\mathrm{ContOn}(f,[a,b]) \to \varphi_{720}$}
\AXC{$\varphi_{834}$\lbl{A1}}
\BIC{$\mathrm{ContOn}(f,[a,b]) \to \varphi_{675}$}
\AXC{$\stackrel{\Pi_{195}}{\vdots}$}\noLine
\UIC{$\varphi_{824}$}
\BIC{$\mathrm{ContOn}(f,[a,b]) \to \varphi_{676}$}
\AXC{$\stackrel{\Pi_{197}}{\vdots}$}\noLine
\UIC{$\varphi_{837}$}
\BIC{$\varphi_{261} \to \mathrm{ContOn}(f,[a,b]) \to \varphi_{642}$}
\BIC{$\mathrm{ContOn}(f,[a,b]) \to \varphi_{642}$}
\AXC{$\stackrel{\Pi_{198}}{\vdots}$}\noLine
\UIC{$\varphi_{795}$}
\BIC{$\mathrm{ContOn}(f,[a,b]) \to \varphi_{618}$}
\end{ptree}
\noindent$\Pi_{200}$:
\begin{ptree}
\AXC{$\forall x \in W\,(x \leq \xi) \to \xi + d \in W \to \xi + d \leq \xi$\lbl{A2}}
\AXC{$\stackrel{\Pi_{22}}{\vdots}$}\noLine
\UIC{$\mathrm{ContOn}(f,[a,b]) \to \varphi_{342}$}
\AXC{$(\mathrm{ContOn}(f,[a,b]) \to \varphi_{342}) \to \varphi_{483}$\lbl{A1}}
\BIC{$\varphi_{483}$}
\BIC{$\mathrm{ContOn}(f,[a,b]) \to \varphi_{374}$}
\end{ptree}
\noindent$\Pi_{201}$:
\begin{ptree}
\AXC{$\stackrel{\Pi_{199}}{\vdots}$}\noLine
\UIC{$\mathrm{ContOn}(f,[a,b]) \to \varphi_{618}$}
\AXC{$\stackrel{\Pi_{200}}{\vdots}$}\noLine
\UIC{$\mathrm{ContOn}(f,[a,b]) \to \varphi_{374}$}
\AXC{$(\mathrm{ContOn}(f,[a,b]) \to \varphi_{374}) \to \varphi_{784}$\lbl{A1}}
\BIC{$\varphi_{784}$}
\BIC{$\mathrm{ContOn}(f,[a,b]) \to \varphi_{614}$}
\end{ptree}
\noindent$\Pi_{202}$:
\begin{ptree}
\AXC{$\forall x\,\forall y\,(x \leq y \to \neg y < x)$\lbl{[O7]}}
\AXC{$\varphi_{11}$\lbl{A2}}
\BIC{$\forall y\,(\xi + d \leq y \to \neg y < \xi + d)$}
\AXC{$\forall y\,(\xi + d \leq y \to \neg y < \xi + d) \to \xi + d \leq \xi \to \neg \xi < \xi + d$\lbl{A2}}
\BIC{$\xi + d \leq \xi \to \neg \xi < \xi + d$}
\end{ptree}
\noindent$\Pi_{203}$:
\begin{ptree}
\AXC{$\stackrel{\Pi_{201}}{\vdots}$}\noLine
\UIC{$\mathrm{ContOn}(f,[a,b]) \to \varphi_{614}$}
\AXC{$\stackrel{\Pi_{202}}{\vdots}$}\noLine
\UIC{$\xi + d \leq \xi \to \neg \xi < \xi + d$}
\AXC{$(\xi + d \leq \xi \to \neg \xi < \xi + d) \to \varphi_{786}$\lbl{A1}}
\BIC{$\varphi_{786}$}
\BIC{$\mathrm{ContOn}(f,[a,b]) \to \varphi_{623}$}
\end{ptree}
\noindent$\Pi_{204}$:
\begin{ptree}
\AXC{$\stackrel{\Pi_{203}}{\vdots}$}\noLine
\UIC{$\mathrm{ContOn}(f,[a,b]) \to \varphi_{623}$}
\AXC{$(\mathrm{ContOn}(f,[a,b]) \to \varphi_{623}) \to \varphi_{817}$\lbl{A1}}
\BIC{$\varphi_{817}$}
\end{ptree}
\noindent$\Pi_{205}$:
\begin{ptree}
\AXC{$\forall x\,(\neg x < x)$\lbl{[O1]}}
\AXC{$\forall x\,(\neg x < x) \to \neg \mu < \mu$\lbl{A2}}
\BIC{$\neg \mu < \mu$}
\AXC{$\neg \mu < \mu \to \mathrm{ContOn}(f,[a,b]) \to \varphi_{428}$\lbl{A1}}
\BIC{$\mathrm{ContOn}(f,[a,b]) \to \varphi_{428}$}
\end{ptree}
\noindent$\Pi_{206}$:
\begin{ptree}
\AXC{$\mathrm{ContOn}(f,[a,b]) \to \varphi_{409}$\lbl{A1}}
\AXC{$\mathrm{ContOn}(f,[a,b]) \to \varphi_{410}$\lbl{A1}}
\AXC{$\varphi_{701}$\lbl{A1}}
\BIC{$\mathrm{ContOn}(f,[a,b]) \to \varphi_{399}$}
\AXC{$\varphi_{762}$\lbl{A1}}
\BIC{$(\mathrm{ContOn}(f,[a,b]) \to \varphi_{409}) \to \varphi_{493}$}
\BIC{$\varphi_{493}$}
\end{ptree}
\noindent$\Pi_{207}$:
\begin{ptree}
\AXC{$\forall x\,\forall y\,\forall z\,(x < y \to y < z \to x < z)$\lbl{[O2]}}
\AXC{$\varphi_{41}$\lbl{A2}}
\BIC{$\forall y\,\forall z\,(\mu < y \to y < z \to \mu < z)$}
\AXC{$\varphi_{89}$\lbl{A2}}
\BIC{$\forall z\,(\mu < f(b) \to f(b) < z \to \mu < z)$}
\AXC{$\varphi_{166}$\lbl{A2}}
\BIC{$\mu < f(b) \to f(b) < \mu \to \mu < \mu$}
\AXC{$\varphi_{755}$\lbl{A1}}
\BIC{$(\mathrm{ContOn}(f,[a,b]) \to a \in \mathbb{R} \to \varphi_{132}) \to \varphi_{714}$}
\end{ptree}
\noindent$\Pi_{208}$:
\begin{ptree}
\AXC{$\mathrm{ContOn}(f,[a,b]) \to a \in \mathbb{R} \to \varphi_{132}$\lbl{A1}}
\AXC{$\stackrel{\Pi_{207}}{\vdots}$}\noLine
\UIC{$(\mathrm{ContOn}(f,[a,b]) \to a \in \mathbb{R} \to \varphi_{132}) \to \varphi_{714}$}
\BIC{$\varphi_{714}$}
\end{ptree}
\noindent$\Pi_{209}$:
\begin{ptree}
\AXC{$z = w \to f(z) < \mu \to f(w) < \mu$\lbl{A6}}
\UIC{$\forall w\,(z = w \to f(z) < \mu \to f(w) < \mu)$}
\UIC{$\forall z\,\forall w\,(z = w \to f(z) < \mu \to f(w) < \mu)$}
\AXC{$\varphi_{165}$\lbl{A2}}
\BIC{$\forall w\,(\xi = w \to f(\xi) < \mu \to f(w) < \mu)$}
\AXC{$\varphi_{129}$\lbl{A2}}
\BIC{$\xi = b \to f(\xi) < \mu \to f(b) < \mu$}
\AXC{$\stackrel{\Pi_{206}}{\vdots}$}\noLine
\UIC{$\varphi_{493}$}
\BIC{$\mathrm{ContOn}(f,[a,b]) \to \varphi_{437}$}
\AXC{$\stackrel{\Pi_{208}}{\vdots}$}\noLine
\UIC{$\varphi_{714}$}
\BIC{$\mathrm{ContOn}(f,[a,b]) \to \varphi_{421}$}
\end{ptree}
\noindent$\Pi_{210}$:
\begin{ptree}
\AXC{$\stackrel{\Pi_{205}}{\vdots}$}\noLine
\UIC{$\mathrm{ContOn}(f,[a,b]) \to \varphi_{428}$}
\AXC{$\stackrel{\Pi_{209}}{\vdots}$}\noLine
\UIC{$\mathrm{ContOn}(f,[a,b]) \to \varphi_{421}$}
\AXC{$(\mathrm{ContOn}(f,[a,b]) \to \varphi_{421}) \to \varphi_{695}$\lbl{A1}}
\BIC{$\varphi_{695}$}
\BIC{$\mathrm{ContOn}(f,[a,b]) \to \varphi_{366}$}
\end{ptree}
\noindent$\Pi_{211}$:
\begin{ptree}
\AXC{$\stackrel{\Pi_{210}}{\vdots}$}\noLine
\UIC{$\mathrm{ContOn}(f,[a,b]) \to \varphi_{366}$}
\AXC{$\stackrel{\Pi_{42}}{\vdots}$}\noLine
\UIC{$\mathrm{ContOn}(f,[a,b]) \to \varphi_{316}$}
\AXC{$(\mathrm{ContOn}(f,[a,b]) \to \varphi_{316}) \to \varphi_{686}$\lbl{A1}}
\BIC{$\varphi_{686}$}
\BIC{$\mathrm{ContOn}(f,[a,b]) \to \varphi_{370}$}
\end{ptree}
\noindent$\Pi_{212}$:
\begin{ptree}
\AXC{$\stackrel{\Pi_{3}}{\vdots}$}\noLine
\UIC{$\mathrm{ContOn}(f,[a,b]) \to \varphi_{328}$}
\AXC{$\forall x\,\forall y\,(x \in \mathbb{R} \to y \in \mathbb{R} \to x < y \to 0 < y - x)$\lbl{[F24]}}
\AXC{$\varphi_{276}$\lbl{A2}}
\BIC{$\forall y\,(\xi \in \mathbb{R} \to y \in \mathbb{R} \to \xi < y \to 0 < y - \xi)$}
\AXC{$\varphi_{268}$\lbl{A2}}
\BIC{$\xi \in \mathbb{R} \to b \in \mathbb{R} \to \xi < b \to 0 < b - \xi$}
\AXC{$\varphi_{825}$\lbl{A1}}
\BIC{$(\mathrm{ContOn}(f,[a,b]) \to \varphi_{328}) \to \varphi_{750}$}
\BIC{$\varphi_{750}$}
\end{ptree}
\noindent$\Pi_{213}$:
\begin{ptree}
\AXC{$\stackrel{\Pi_{211}}{\vdots}$}\noLine
\UIC{$\mathrm{ContOn}(f,[a,b]) \to \varphi_{370}$}
\AXC{$\mathrm{ContOn}(f,[a,b]) \to a \in \mathbb{R} \to b \in \mathbb{R} \to b \in \mathbb{R}$\lbl{A1}}
\AXC{$\stackrel{\Pi_{212}}{\vdots}$}\noLine
\UIC{$\varphi_{750}$}
\BIC{$(\mathrm{ContOn}(f,[a,b]) \to \varphi_{370}) \to \varphi_{583}$}
\BIC{$\varphi_{583}$}
\end{ptree}
\noindent$\Pi_{214}$:
\begin{ptree}
\AXC{$\forall e \in \mathbb{R}\,(0 < e \to 0 < e \cdot 2^{-1})$\lbl{[F35]}}
\AXC{$\varphi_{134}$\lbl{A2}}
\BIC{$\delta \in \mathbb{R} \to 0 < \delta \to 0 < \delta \cdot 2^{-1}$}
\AXC{$(\delta \in \mathbb{R} \to 0 < \delta \to 0 < \delta \cdot 2^{-1}) \to \varphi_{579} \to \varphi_{577} \to \varphi_{589}$\lbl{A1}}
\BIC{$\varphi_{579} \to \varphi_{577} \to \varphi_{589}$}
\end{ptree}
\noindent$\Pi_{215}$:
\begin{ptree}
\AXC{$\forall x\,\forall y\,(x \in \mathbb{R} \to y \in \mathbb{R} \to x - y \in \mathbb{R})$\lbl{[F1]}}
\AXC{$\varphi_{215}$\lbl{A2}}
\BIC{$\forall y\,(b \in \mathbb{R} \to y \in \mathbb{R} \to b - y \in \mathbb{R})$}
\AXC{$\varphi_{185}$\lbl{A2}}
\BIC{$b \in \mathbb{R} \to \xi \in \mathbb{R} \to b - \xi \in \mathbb{R}$}
\AXC{$(b \in \mathbb{R} \to \xi \in \mathbb{R} \to b - \xi \in \mathbb{R}) \to \varphi_{748}$\lbl{A1}}
\BIC{$\varphi_{748}$}
\end{ptree}
\noindent$\Pi_{216}$:
\begin{ptree}
\AXC{$\stackrel{\Pi_{3}}{\vdots}$}\noLine
\UIC{$\mathrm{ContOn}(f,[a,b]) \to \varphi_{328}$}
\AXC{$\mathrm{ContOn}(f,[a,b]) \to a \in \mathbb{R} \to b \in \mathbb{R} \to b \in \mathbb{R}$\lbl{A1}}
\AXC{$\stackrel{\Pi_{215}}{\vdots}$}\noLine
\UIC{$\varphi_{748}$}
\BIC{$(\mathrm{ContOn}(f,[a,b]) \to \varphi_{328}) \to \varphi_{587}$}
\BIC{$\varphi_{587}$}
\end{ptree}
\noindent$\Pi_{217}$:
\begin{ptree}
\AXC{$\varphi_{711}$\lbl{A1}}
\AXC{$\varphi_{711} \to \varphi_{579}$\lbl{A1}}
\BIC{$\varphi_{579}$}
\AXC{$\forall x\,\forall y\,(x \in \mathbb{R} \to y \in \mathbb{R} \to x \cdot y \in \mathbb{R})$\lbl{(R1)}}
\AXC{$\varphi_{228}$\lbl{A2}}
\BIC{$\forall y\,(\delta \in \mathbb{R} \to y \in \mathbb{R} \to \delta \cdot y \in \mathbb{R})$}
\AXC{$\varphi_{257}$\lbl{A2}}
\BIC{$\delta \in \mathbb{R} \to 2^{-1} \in \mathbb{R} \to \delta \cdot 2^{-1} \in \mathbb{R}$}
\AXC{$\varphi_{770}$\lbl{A1}}
\BIC{$\varphi_{579} \to 2^{-1} \in \mathbb{R} \to \varphi_{591}$}
\BIC{$2^{-1} \in \mathbb{R} \to \varphi_{591}$}
\end{ptree}
\noindent$\Pi_{218}$:
\begin{ptree}
\AXC{$\varphi_{127}$\lbl{[D8]}}
\AXC{$\varphi_{127} \to \varphi_{286}$\lbl{A2}}
\BIC{$\varphi_{286}$}
\AXC{$\varphi_{286} \to \delta \cdot 2^{-1} \in \mathbb{R} \to b - \xi \in \mathbb{R} \to 0 < \delta \cdot 2^{-1} \to \varphi_{38}$\lbl{A2}}
\BIC{$\delta \cdot 2^{-1} \in \mathbb{R} \to b - \xi \in \mathbb{R} \to 0 < \delta \cdot 2^{-1} \to \varphi_{38}$}
\AXC{$\varphi_{887}$\lbl{A1}}
\BIC{$\varphi_{591} \to \varphi_{587} \to \varphi_{589} \to \varphi_{583} \to \mathrm{ContOn}(f,[a,b]) \to \varphi_{597}$}
\end{ptree}
\noindent$\Pi_{219}$:
\begin{ptree}
\AXC{$\stackrel{\Pi_{216}}{\vdots}$}\noLine
\UIC{$\varphi_{587}$}
\AXC{$2^{-1} \in \mathbb{R}$\lbl{[F31]}}
\AXC{$\stackrel{\Pi_{217}}{\vdots}$}\noLine
\UIC{$2^{-1} \in \mathbb{R} \to \varphi_{591}$}
\BIC{$\varphi_{591}$}
\AXC{$\stackrel{\Pi_{218}}{\vdots}$}\noLine
\UIC{$\varphi_{591} \to \varphi_{587} \to \varphi_{589} \to \varphi_{583} \to \mathrm{ContOn}(f,[a,b]) \to \varphi_{597}$}
\BIC{$\varphi_{587} \to \varphi_{589} \to \varphi_{583} \to \mathrm{ContOn}(f,[a,b]) \to \varphi_{597}$}
\BIC{$\varphi_{589} \to \varphi_{583} \to \mathrm{ContOn}(f,[a,b]) \to \varphi_{597}$}
\end{ptree}
\noindent$\Pi_{220}$:
\begin{ptree}
\AXC{$\varphi_{711}$\lbl{A1}}
\AXC{$\varphi_{711} \to \varphi_{577}$\lbl{A1}}
\BIC{$\varphi_{577}$}
\AXC{$\varphi_{711}$\lbl{A1}}
\AXC{$\varphi_{711} \to \varphi_{579}$\lbl{A1}}
\BIC{$\varphi_{579}$}
\AXC{$\stackrel{\Pi_{214}}{\vdots}$}\noLine
\UIC{$\varphi_{579} \to \varphi_{577} \to \varphi_{589}$}
\BIC{$\varphi_{577} \to \varphi_{589}$}
\BIC{$\varphi_{589}$}
\AXC{$\stackrel{\Pi_{219}}{\vdots}$}\noLine
\UIC{$\varphi_{589} \to \varphi_{583} \to \mathrm{ContOn}(f,[a,b]) \to \varphi_{597}$}
\BIC{$\varphi_{583} \to \mathrm{ContOn}(f,[a,b]) \to \varphi_{597}$}
\end{ptree}
\noindent$\Pi_{221}$:
\begin{ptree}
\AXC{$\stackrel{\Pi_{213}}{\vdots}$}\noLine
\UIC{$\varphi_{583}$}
\AXC{$\stackrel{\Pi_{220}}{\vdots}$}\noLine
\UIC{$\varphi_{583} \to \mathrm{ContOn}(f,[a,b]) \to \varphi_{597}$}
\BIC{$\mathrm{ContOn}(f,[a,b]) \to \varphi_{597}$}
\AXC{$(\mathrm{ContOn}(f,[a,b]) \to \varphi_{597}) \to \varphi_{703} \to \varphi_{891}$\lbl{A1}}
\BIC{$\varphi_{703} \to \varphi_{891}$}
\end{ptree}
\noindent$\Pi_{222}$:
\begin{ptree}
\AXC{$\varphi_{703} \to \varphi_{708}$\lbl{A3}}
\AXC{$\mathrm{ContOn}(f,[a,b]) \to \varphi_{679}$\lbl{A1}}
\AXC{$\varphi_{809}$\lbl{A1}}
\BIC{$\mathrm{ContOn}(f,[a,b]) \to \varphi_{602}$}
\AXC{$\stackrel{\Pi_{132}}{\vdots}$}\noLine
\UIC{$\varphi_{781}$}
\BIC{$\mathrm{ContOn}(f,[a,b]) \to \varphi_{617}$}
\AXC{$\stackrel{\Pi_{204}}{\vdots}$}\noLine
\UIC{$\varphi_{817}$}
\BIC{$\mathrm{ContOn}(f,[a,b]) \to \varphi_{681}$}
\AXC{$\varphi_{827}$\lbl{A1}}
\BIC{$\varphi_{647} \to \neg (d \in \mathbb{R} \land 0 < d \land d \leq \delta \cdot 2^{-1} \land d \leq b - \xi)$}
\UIC{$\varphi_{703}$}
\AXC{$\stackrel{\Pi_{221}}{\vdots}$}\noLine
\UIC{$\varphi_{703} \to \varphi_{891}$}
\BIC{$\varphi_{891}$}
\BIC{$(\varphi_{647} \to \forall d\,(\varphi_{647})) \to \mathrm{ContOn}(f,[a,b]) \to \varphi_{595}$}
\end{ptree}
\noindent$\Pi_{223}$:
\begin{ptree}
\AXC{$\forall x\,\forall y\,(x \in \mathbb{R} \to y \in \mathbb{R} \to x < y \to 0 < y - x)$\lbl{[F24]}}
\AXC{$\forall x\,\forall y\,(x \in \mathbb{R} \to y \in \mathbb{R} \to x < y \to 0 < y - x) \to \varphi_{58}$\lbl{A2}}
\BIC{$\varphi_{58}$}
\end{ptree}
\noindent$\Pi_{224}$:
\begin{ptree}
\AXC{$\stackrel{\Pi_{60}}{\vdots}$}\noLine
\UIC{$\mathrm{ContOn}(f,[a,b]) \to \varphi_{347}$}
\AXC{$\stackrel{\Pi_{223}}{\vdots}$}\noLine
\UIC{$\varphi_{58}$}
\AXC{$\varphi_{58} \to f(\xi) \in \mathbb{R} \to \mu \in \mathbb{R} \to f(\xi) < \mu \to 0 < \mu - f(\xi)$\lbl{A2}}
\BIC{$f(\xi) \in \mathbb{R} \to \mu \in \mathbb{R} \to f(\xi) < \mu \to 0 < \mu - f(\xi)$}
\AXC{$\varphi_{802}$\lbl{A1}}
\BIC{$(\mathrm{ContOn}(f,[a,b]) \to \varphi_{347}) \to \varphi_{56} \to \varphi_{706}$}
\BIC{$\varphi_{56} \to \varphi_{706}$}
\end{ptree}
\noindent$\Pi_{225}$:
\begin{ptree}
\AXC{$\forall x\,\forall y\,(x \in \mathbb{R} \to y \in \mathbb{R} \to x - y \in \mathbb{R})$\lbl{[F1]}}
\AXC{$\varphi_{230}$\lbl{A2}}
\BIC{$\forall y\,(\mu \in \mathbb{R} \to y \in \mathbb{R} \to \mu - y \in \mathbb{R})$}
\AXC{$\varphi_{280}$\lbl{A2}}
\BIC{$\mu \in \mathbb{R} \to f(\xi) \in \mathbb{R} \to \mu - f(\xi) \in \mathbb{R}$}
\AXC{$(\mu \in \mathbb{R} \to f(\xi) \in \mathbb{R} \to \mu - f(\xi) \in \mathbb{R}) \to \varphi_{56} \to \varphi_{698}$\lbl{A1}}
\BIC{$\varphi_{56} \to \varphi_{698}$}
\end{ptree}
\noindent$\Pi_{226}$:
\begin{ptree}
\AXC{$\varphi_{302} \to \varphi_{359}$\lbl{A2}}
\AXC{$\stackrel{\Pi_{111}}{\vdots}$}\noLine
\UIC{$\mathrm{ContOn}(f,[a,b]) \to \varphi_{553}$}
\AXC{$\varphi_{846}$\lbl{A1}}
\BIC{$(\varphi_{302} \to \varphi_{359}) \to \mathrm{ContOn}(f,[a,b]) \to \varphi_{581}$}
\BIC{$\mathrm{ContOn}(f,[a,b]) \to \varphi_{581}$}
\AXC{$(\mathrm{ContOn}(f,[a,b]) \to \varphi_{581}) \to \varphi_{779}$\lbl{A1}}
\BIC{$\varphi_{779}$}
\end{ptree}
\noindent$\Pi_{227}$:
\begin{ptree}
\AXC{$\stackrel{\Pi_{60}}{\vdots}$}\noLine
\UIC{$\mathrm{ContOn}(f,[a,b]) \to \varphi_{347}$}
\AXC{$\varphi_{56}$\lbl{A1}}
\AXC{$\stackrel{\Pi_{225}}{\vdots}$}\noLine
\UIC{$\varphi_{56} \to \varphi_{698}$}
\BIC{$\varphi_{698}$}
\BIC{$\mathrm{ContOn}(f,[a,b]) \to \varphi_{440}$}
\AXC{$\stackrel{\Pi_{226}}{\vdots}$}\noLine
\UIC{$\varphi_{779}$}
\BIC{$(\mathrm{ContOn}(f,[a,b]) \to \varphi_{434}) \to \varphi_{574}$}
\end{ptree}
\noindent$\Pi_{228}$:
\begin{ptree}
\AXC{$\mathrm{ContOn}(f,[a,b]) \to \varphi_{409}$\lbl{A1}}
\AXC{$\varphi_{56}$\lbl{A1}}
\AXC{$\stackrel{\Pi_{224}}{\vdots}$}\noLine
\UIC{$\varphi_{56} \to \varphi_{706}$}
\BIC{$\varphi_{706}$}
\BIC{$\mathrm{ContOn}(f,[a,b]) \to \varphi_{434}$}
\AXC{$\stackrel{\Pi_{227}}{\vdots}$}\noLine
\UIC{$(\mathrm{ContOn}(f,[a,b]) \to \varphi_{434}) \to \varphi_{574}$}
\BIC{$\varphi_{574}$}
\AXC{$\varphi_{574} \to \varphi_{889}$\lbl{A1}}
\BIC{$\varphi_{889}$}
\end{ptree}
\noindent$\Pi_{229}$:
\begin{ptree}
\AXC{$\varphi_{777}$\lbl{A3}}
\AXC{$\varphi_{647} \to \forall d\,(\varphi_{647})$\lbl{A4}}
\AXC{$\stackrel{\Pi_{222}}{\vdots}$}\noLine
\UIC{$(\varphi_{647} \to \forall d\,(\varphi_{647})) \to \mathrm{ContOn}(f,[a,b]) \to \varphi_{595}$}
\BIC{$\mathrm{ContOn}(f,[a,b]) \to \varphi_{595}$}
\AXC{$\varphi_{768}$\lbl{A1}}
\BIC{$\varphi_{482} \to \neg (\delta \in \mathbb{R} \land 0 < \delta \land \varphi_{159})$}
\UIC{$\forall \delta\,(\varphi_{482} \to \neg (\delta \in \mathbb{R} \land 0 < \delta \land \varphi_{159}))$}
\AXC{$\stackrel{\Pi_{228}}{\vdots}$}\noLine
\UIC{$\varphi_{889}$}
\BIC{$\varphi_{777} \to (\varphi_{482} \to \forall \delta\,(\varphi_{482})) \to \varphi_{479}$}
\BIC{$(\varphi_{482} \to \forall \delta\,(\varphi_{482})) \to \varphi_{479}$}
\end{ptree}
\noindent$\Pi_{230}$:
\begin{ptree}
\AXC{$\forall x\,\forall y\,(x \in \mathbb{R} \to y \in \mathbb{R} \to x < y \lor x = y \lor y < x)$\lbl{(R13)}}
\AXC{$\forall x\,\forall y\,(x \in \mathbb{R} \to y \in \mathbb{R} \to x < y \lor x = y \lor y < x) \to \varphi_{110}$\lbl{A2}}
\BIC{$\varphi_{110}$}
\AXC{$\varphi_{110} \to \varphi_{124}$\lbl{A2}}
\BIC{$\varphi_{124}$}
\end{ptree}
\noindent$\Pi_{231}$:
\begin{ptree}
\AXC{$\varphi_{56}$\lbl{A1}}
\AXC{$\stackrel{\Pi_{60}}{\vdots}$}\noLine
\UIC{$\mathrm{ContOn}(f,[a,b]) \to \varphi_{347}$}
\AXC{$\stackrel{\Pi_{230}}{\vdots}$}\noLine
\UIC{$\varphi_{124}$}
\AXC{$\varphi_{124} \to (\mathrm{ContOn}(f,[a,b]) \to \varphi_{347}) \to \varphi_{56} \to \varphi_{451}$\lbl{A1}}
\BIC{$(\mathrm{ContOn}(f,[a,b]) \to \varphi_{347}) \to \varphi_{56} \to \varphi_{451}$}
\BIC{$\varphi_{56} \to \varphi_{451}$}
\BIC{$\varphi_{451}$}
\end{ptree}
\noindent$\Pi_{232}$:
\begin{ptree}
\AXC{$\stackrel{\Pi_{130}}{\vdots}$}\noLine
\UIC{$\varphi_{475}$}
\AXC{$\varphi_{482} \to \forall \delta\,(\varphi_{482})$\lbl{A4}}
\AXC{$\stackrel{\Pi_{229}}{\vdots}$}\noLine
\UIC{$(\varphi_{482} \to \forall \delta\,(\varphi_{482})) \to \varphi_{479}$}
\BIC{$\varphi_{479}$}
\AXC{$\stackrel{\Pi_{231}}{\vdots}$}\noLine
\UIC{$\varphi_{451}$}
\AXC{$\varphi_{451} \to \varphi_{479} \to \varphi_{475} \to \varphi_{478} \to \varphi_{453}$\lbl{A1}}
\BIC{$\varphi_{479} \to \varphi_{475} \to \varphi_{478} \to \varphi_{453}$}
\BIC{$\varphi_{475} \to \varphi_{478} \to \varphi_{453}$}
\BIC{$\varphi_{478} \to \varphi_{453}$}
\end{ptree}
\noindent$\Pi_{233}$:
\begin{ptree}
\AXC{$\forall X\,\forall x \in X\,(X \neq \emptyset)$\lbl{[B2]}}
\AXC{$\forall X\,\forall x \in X\,(X \neq \emptyset) \to \forall x \in W\,(W \neq \emptyset)$\lbl{A2}}
\BIC{$\forall x \in W\,(W \neq \emptyset)$}
\AXC{$\forall x \in W\,(W \neq \emptyset) \to a \in W \to W \neq \emptyset$\lbl{A2}}
\BIC{$a \in W \to W \neq \emptyset$}
\end{ptree}
\noindent$\Pi_{234}$:
\begin{ptree}
\AXC{$\stackrel{\Pi_{48}}{\vdots}$}\noLine
\UIC{$\varphi_{339}$}
\AXC{$\stackrel{\Pi_{233}}{\vdots}$}\noLine
\UIC{$a \in W \to W \neq \emptyset$}
\AXC{$(a \in W \to W \neq \emptyset) \to \varphi_{339} \to \varphi_{350}$\lbl{A1}}
\BIC{$\varphi_{339} \to \varphi_{350}$}
\BIC{$\varphi_{350}$}
\end{ptree}
\noindent$\Pi_{235}$:
\begin{ptree}
\AXC{$\varphi_{368}$\lbl{A1}}
\AXC{$\varphi_{188}$\lbl{A2}}
\AXC{$\mathrm{ContOn}(f,[a,b]) \to \varphi_{397}$\lbl{A1}}
\AXC{$\varphi_{732}$\lbl{A1}}
\BIC{$\varphi_{188} \to \mathrm{ContOn}(f,[a,b]) \to \varphi_{385}$}
\BIC{$\mathrm{ContOn}(f,[a,b]) \to \varphi_{385}$}
\AXC{$(\mathrm{ContOn}(f,[a,b]) \to \varphi_{385}) \to \varphi_{368} \to \varphi_{452}$\lbl{A1}}
\BIC{$\varphi_{368} \to \varphi_{452}$}
\BIC{$\varphi_{452}$}
\end{ptree}
\noindent$\Pi_{236}$:
\begin{ptree}
\AXC{$\forall a\,\forall b\,\forall x\,(x \in [a,b] \leftrightarrow x \in \mathbb{R} \land a \leq x \land x \leq b)$\lbl{Def{$[a,b]$}}}
\AXC{$\varphi_{236}$\lbl{A2}}
\BIC{$\forall b\,\forall x\,(x \in [a,b] \leftrightarrow x \in \mathbb{R} \land a \leq x \land x \leq b)$}
\AXC{$\varphi_{195}$\lbl{A2}}
\BIC{$\forall x\,(x \in [a,b] \leftrightarrow x \in \mathbb{R} \land a \leq x \land x \leq b)$}
\AXC{$\varphi_{151}$\lbl{A2}}
\BIC{$u \in [a,b] \leftrightarrow u \in \mathbb{R} \land a \leq u \land u \leq b$}
\AXC{$\stackrel{\Pi_{235}}{\vdots}$}\noLine
\UIC{$\varphi_{452}$}
\BIC{$\mathrm{ContOn}(f,[a,b]) \to \varphi_{319}$}
\end{ptree}
\noindent$\Pi_{237}$:
\begin{ptree}
\AXC{$\varphi_{658}$\lbl{A3}}
\AXC{$\stackrel{\Pi_{236}}{\vdots}$}\noLine
\UIC{$\mathrm{ContOn}(f,[a,b]) \to \varphi_{319}$}
\AXC{$(\mathrm{ContOn}(f,[a,b]) \to \varphi_{319}) \to \varphi_{318} \to u \in W \to u \in \mathbb{R}$\lbl{A1}}
\BIC{$\varphi_{318} \to u \in W \to u \in \mathbb{R}$}
\UIC{$\forall u\,(\varphi_{318} \to u \in W \to u \in \mathbb{R})$}
\AXC{$\varphi_{855}$\lbl{A1}}
\BIC{$\varphi_{658} \to (\varphi_{318} \to \forall u\,(\varphi_{318})) \to \mathrm{ContOn}(f,[a,b]) \to \varphi_{323}$}
\BIC{$(\varphi_{318} \to \forall u\,(\varphi_{318})) \to \mathrm{ContOn}(f,[a,b]) \to \varphi_{323}$}
\end{ptree}
\noindent$\Pi_{238}$:
\begin{ptree}
\AXC{$\forall x\,\forall y\,(x \subseteq y \leftrightarrow \forall u \in x\,(u \in y))$\lbl{Def$\subseteq$}}
\AXC{$\varphi_{66}$\lbl{A2}}
\BIC{$\forall y\,(W \subseteq y \leftrightarrow \forall u \in W\,(u \in y))$}
\AXC{$\varphi_{67}$\lbl{A2}}
\BIC{$W \subseteq \mathbb{R} \leftrightarrow \forall u \in W\,(u \in \mathbb{R})$}
\AXC{$\varphi_{684}$\lbl{A1}}
\BIC{$(\mathrm{ContOn}(f,[a,b]) \to \varphi_{323}) \to \varphi_{351}$}
\end{ptree}
\noindent$\Pi_{239}$:
\begin{ptree}
\AXC{$\varphi_{318} \to \forall u\,(\varphi_{318})$\lbl{A4}}
\AXC{$\stackrel{\Pi_{237}}{\vdots}$}\noLine
\UIC{$(\varphi_{318} \to \forall u\,(\varphi_{318})) \to \mathrm{ContOn}(f,[a,b]) \to \varphi_{323}$}
\BIC{$\mathrm{ContOn}(f,[a,b]) \to \varphi_{323}$}
\AXC{$\stackrel{\Pi_{238}}{\vdots}$}\noLine
\UIC{$(\mathrm{ContOn}(f,[a,b]) \to \varphi_{323}) \to \varphi_{351}$}
\BIC{$\varphi_{351}$}
\end{ptree}
\noindent$\Pi_{240}$:
\begin{ptree}
\AXC{$\stackrel{\Pi_{234}}{\vdots}$}\noLine
\UIC{$\varphi_{350}$}
\AXC{$\stackrel{\Pi_{239}}{\vdots}$}\noLine
\UIC{$\varphi_{351}$}
\AXC{$\varphi_{16}$\lbl{(R16)}}
\AXC{$\varphi_{16} \to \varphi_{14}$\lbl{A2}}
\BIC{$\varphi_{14}$}
\AXC{$\varphi_{14} \to \varphi_{351} \to \varphi_{350} \to \varphi_{599}$\lbl{A1}}
\BIC{$\varphi_{351} \to \varphi_{350} \to \varphi_{599}$}
\BIC{$\varphi_{350} \to \varphi_{599}$}
\BIC{$\varphi_{599}$}
\end{ptree}
\noindent$\Pi_{241}$:
\begin{ptree}
\AXC{$\forall b\,(\neg \mathrm{UB}(b,W)) \to \neg \mathrm{UB}(b,W)$\lbl{A2}}
\AXC{$\stackrel{\Pi_{39}}{\vdots}$}\noLine
\UIC{$\varphi_{338}$}
\AXC{$\varphi_{338} \to \varphi_{441}$\lbl{A1}}
\BIC{$\varphi_{441}$}
\BIC{$\mathrm{ContOn}(f,[a,b]) \to \varphi_{304}$}
\AXC{$\stackrel{\Pi_{240}}{\vdots}$}\noLine
\UIC{$\varphi_{599}$}
\BIC{$\mathrm{ContOn}(f,[a,b]) \to \varphi_{308}$}
\end{ptree}
\noindent$\Pi_{242}$:
\begin{ptree}
\AXC{$\stackrel{\Pi_{241}}{\vdots}$}\noLine
\UIC{$\mathrm{ContOn}(f,[a,b]) \to \varphi_{308}$}
\AXC{$(\mathrm{ContOn}(f,[a,b]) \to \varphi_{308}) \to \forall \xi\,(\varphi_{459}) \to \varphi_{868}$\lbl{A1}}
\BIC{$\forall \xi\,(\varphi_{459}) \to \varphi_{868}$}
\end{ptree}
\noindent$\Pi_{243}$:
\begin{ptree}
\AXC{$\forall \xi\,(\varphi_{459}) \to \varphi_{465}$\lbl{A3}}
\AXC{$\varphi_{481} \to \forall \delta\,(\varphi_{481})$\lbl{A4}}
\AXC{$\stackrel{\Pi_{114}}{\vdots}$}\noLine
\UIC{$(\varphi_{481} \to \forall \delta\,(\varphi_{481})) \to \varphi_{478}$}
\BIC{$\varphi_{478}$}
\AXC{$\stackrel{\Pi_{232}}{\vdots}$}\noLine
\UIC{$\varphi_{478} \to \varphi_{453}$}
\BIC{$\varphi_{453}$}
\AXC{$\varphi_{453} \to \varphi_{459}$\lbl{A1}}
\BIC{$\varphi_{459}$}
\UIC{$\forall \xi\,(\varphi_{459})$}
\AXC{$\stackrel{\Pi_{242}}{\vdots}$}\noLine
\UIC{$\forall \xi\,(\varphi_{459}) \to \varphi_{868}$}
\BIC{$\varphi_{868}$}
\BIC{$\varphi_{716} \to \varphi_{1} \to (\varphi_{1} \to \varphi_{2}) \to \varphi_{3} \to \mathrm{ContOn}(f,[a,b]) \to \varphi_{429}$}
\end{ptree}
\noindent$\Pi_{244}$:
\begin{ptree}
\AXC{$\forall \xi\,(\neg \mathrm{IsSup}(\xi,W)) \to \neg \mathrm{IsSup}(s,W)$\lbl{A2}}
\UIC{$\varphi_{1}$}
\AXC{$\varphi_{716}$\lbl{A4}}
\AXC{$\stackrel{\Pi_{243}}{\vdots}$}\noLine
\UIC{$\varphi_{716} \to \varphi_{1} \to (\varphi_{1} \to \varphi_{2}) \to \varphi_{3} \to \mathrm{ContOn}(f,[a,b]) \to \varphi_{429}$}
\BIC{$\varphi_{1} \to (\varphi_{1} \to \varphi_{2}) \to \varphi_{3} \to \mathrm{ContOn}(f,[a,b]) \to \varphi_{429}$}
\BIC{$(\varphi_{1} \to \varphi_{2}) \to \varphi_{3} \to \mathrm{ContOn}(f,[a,b]) \to \varphi_{429}$}
\end{ptree}
\noindent$\Pi_{245}$:
\begin{ptree}
\AXC{$\forall f\,\forall \mu\,\forall A\,\exists W\,\forall x\,(x \in W \leftrightarrow x \in A \land f(x) < \mu)$\lbl{Sep}}
\AXC{$\varphi_{247}$\lbl{A2}}
\BIC{$\forall \mu\,\forall A\,\exists W\,\forall x\,(x \in W \leftrightarrow x \in A \land f(x) < \mu)$}
\AXC{$\varphi_{209}$\lbl{A2}}
\BIC{$\forall A\,\exists W\,\forall x\,(x \in W \leftrightarrow x \in A \land f(x) < \mu)$}
\AXC{$\varphi_{216}$\lbl{A2}}
\BIC{$\exists W\,\forall x\,(x \in W \leftrightarrow x \in [a,b] \land f(x) < \mu)$}
\end{ptree}
\noindent$\Pi_{246}$:
\begin{ptree}
\AXC{$\stackrel{\Pi_{245}}{\vdots}$}\noLine
\UIC{$\exists W\,\forall x\,(x \in W \leftrightarrow x \in [a,b] \land f(x) < \mu)$}
\AXC{$\exists W\,\forall x\,(x \in W \leftrightarrow x \in [a,b] \land f(x) < \mu) \to \varphi_{320}$\lbl{A1}}
\BIC{$\varphi_{320}$}
\AXC{$\varphi_{320} \to \varphi_{861}$\lbl{A1}}
\BIC{$\varphi_{861}$}
\end{ptree}
\noindent$\Pi_{247}$:
\begin{ptree}
\AXC{$\varphi_{719}$\lbl{A3}}
\AXC{$\varphi_{3}$\lbl{A4}}
\AXC{$\varphi_{1} \to \varphi_{2}$\lbl{A3}}
\AXC{$\stackrel{\Pi_{244}}{\vdots}$}\noLine
\UIC{$(\varphi_{1} \to \varphi_{2}) \to \varphi_{3} \to \mathrm{ContOn}(f,[a,b]) \to \varphi_{429}$}
\BIC{$\varphi_{3} \to \mathrm{ContOn}(f,[a,b]) \to \varphi_{429}$}
\BIC{$\mathrm{ContOn}(f,[a,b]) \to \varphi_{429}$}
\AXC{$\varphi_{712}$\lbl{A1}}
\BIC{$\varphi_{326} \to \neg \forall x\,(x \in W \leftrightarrow x \in [a,b] \land f(x) < \mu)$}
\UIC{$\forall W\,(\varphi_{326} \to \neg \forall x\,(x \in W \leftrightarrow x \in [a,b] \land f(x) < \mu))$}
\AXC{$\stackrel{\Pi_{246}}{\vdots}$}\noLine
\UIC{$\varphi_{861}$}
\BIC{$\varphi_{719} \to (\varphi_{326} \to \forall W\,(\varphi_{326})) \to \varphi_{324}$}
\BIC{$(\varphi_{326} \to \forall W\,(\varphi_{326})) \to \varphi_{324}$}
\end{ptree}
\noindent 主木:
\begin{ptree}
\AXC{$\varphi_{326} \to \forall W\,(\varphi_{326})$\lbl{A4}}
\AXC{$\stackrel{\Pi_{247}}{\vdots}$}\noLine
\UIC{$(\varphi_{326} \to \forall W\,(\varphi_{326})) \to \varphi_{324}$}
\BIC{$\varphi_{324}$}
\UIC{$\forall \mu\,(\varphi_{324})$}
\UIC{$\forall b\,\forall \mu\,(\varphi_{324})$}
\UIC{$\forall a\,\forall b\,\forall \mu\,(\varphi_{324})$}
\UIC{$\forall f\,\forall a\,\forall b\,\forall \mu\,(\varphi_{324})$}
\end{ptree}
\begin{abbrevlist}
\item $\varphi_{1} :\equiv \forall s\,(\forall \xi\,(\neg \mathrm{IsSup}(\xi,W)) \to \neg \mathrm{IsSup}(s,W))$
\item $\varphi_{2} :\equiv \forall s\,\forall \xi\,(\neg \mathrm{IsSup}(\xi,W)) \to \forall s\,(\neg \mathrm{IsSup}(s,W))$
\item $\varphi_{3} :\equiv \forall \xi\,(\neg \mathrm{IsSup}(\xi,W)) \to \forall s\,\forall \xi\,(\neg \mathrm{IsSup}(\xi,W))$
\item $\varphi_{4} :\equiv \mathrm{ContOn}(f,[a,b]) \to \mathrm{ContOn}(f,[a,b])$
\item $\varphi_{5} :\equiv \mathrm{IsSup}(\xi,W) \leftrightarrow \mathrm{UB}(\xi,W) \land \forall w\,(\mathrm{UB}(w,W) \to \xi \leq w)$
\item $\varphi_{6} :\equiv \forall y\,\forall z\,(a \leq y \to y \leq z \to a \leq z) \to \forall z\,(a \leq \xi \to \xi \leq z \to a \leq z)$
\item $\varphi_{7} :\equiv \forall s\,\forall X\,(\mathrm{IsSup}(s,X) \leftrightarrow \mathrm{UB}(s,X) \land \forall w\,(\mathrm{UB}(w,X) \to s \leq w))$
\item $\varphi_{8} :\equiv \forall X\,(\mathrm{IsSup}(\xi,X) \leftrightarrow \mathrm{UB}(\xi,X) \land \forall w\,(\mathrm{UB}(w,X) \to \xi \leq w))$
\item $\varphi_{9} :\equiv \forall w\,(\mathrm{UB}(w,W) \to \xi \leq w) \to \mathrm{UB}(\xi - d,W) \to \xi \leq \xi - d$
\item $\varphi_{10} :\equiv \forall x\,\forall y\,\forall z\,(x \leq y \to y \leq z \to x \leq z) \to \forall y\,\forall z\,(a \leq y \to y \leq z \to a \leq z)$
\item $\varphi_{11} :\equiv \forall x\,\forall y\,(x \leq y \to \neg y < x) \to \forall y\,(\xi + d \leq y \to \neg y < \xi + d)$
\item $\varphi_{12} :\equiv \forall z\,(a \leq \xi \to \xi \leq z \to a \leq z) \to a \leq \xi \to \xi \leq \xi + d \to a \leq \xi + d$
\item $\varphi_{13} :\equiv \forall z\,\forall w\,(z = w \to z = b \to w = b) \to \forall w\,(d + \xi = w \to d + \xi = b \to w = b)$
\item $\varphi_{14} :\equiv W \subseteq \mathbb{R} \to W \neq \emptyset \to \exists b\,(\mathrm{UB}(b,W)) \to \exists s\,(\mathrm{IsSup}(s,W))$
\item $\varphi_{15} :\equiv \forall x\,([a,b] \subseteq \mathrm{dom}(f) \to x \in [a,b] \to x \in \mathrm{dom}(f))$
\item $\varphi_{16} :\equiv \forall X\,(X \subseteq \mathbb{R} \to X \neq \emptyset \to \exists b\,(\mathrm{UB}(b,X)) \to \exists s\,(\mathrm{IsSup}(s,X)))$
\item $\varphi_{17} :\equiv \forall z\,\forall w\,(z = w \to z < \delta \to w < \delta) \to \forall w\,(d = w \to d < \delta \to w < \delta)$
\item $\varphi_{18} :\equiv \mu < f(b) \to \xi \in [a,b] \to \mathrm{ContIn}(f,[a,b],\xi)$
\item $\varphi_{19} :\equiv \forall x\,(x \in W \leftrightarrow x \in [a,b] \land f(x) < \mu) \to \exists b\,(\mathrm{UB}(b,W))$
\item $\varphi_{20} :\equiv \mu < f(b) \to \forall c \in [a,b]\,(\mathrm{ContIn}(f,[a,b],c))$
\item $\varphi_{21} :\equiv \forall x\,\forall y\,\forall z\,(x \leq y \to y < z \to x < z) \to \forall y\,\forall z\,(d \leq y \to y < z \to d < z)$
\item $\varphi_{22} :\equiv \forall X\,(\mathrm{UB}(\xi - d,X) \leftrightarrow \xi - d \in \mathbb{R} \land \forall x \in X\,(x \leq \xi - d))$
\item $\varphi_{23} :\equiv \forall x\,(x \in W \leftrightarrow x \in [a,b] \land f(x) < \mu) \to \exists s\,(\mathrm{IsSup}(s,W))$
\item $\varphi_{24} :\equiv \forall z\,(\lvert \xi + d - \xi \rvert = \lvert d \rvert \to \lvert d \rvert = z \to \lvert \xi + d - \xi \rvert = z)$
\item $\varphi_{25} :\equiv \forall w\,(z = w \to \lvert d + \xi - \xi \rvert = \lvert z \rvert \to \lvert d + \xi - \xi \rvert = \lvert w \rvert)$
\item $\varphi_{26} :\equiv \mathrm{IsSup}(\xi,W) \to f(\xi) < \mu \lor f(\xi) = \mu \lor \mu < f(\xi)$
\item $\varphi_{27} :\equiv \forall x\,\forall y\,(x = y \to y = x) \to \forall y\,(\lvert \xi + d - \xi \rvert = y \to y = \lvert \xi + d - \xi \rvert)$
\item $\varphi_{28} :\equiv \forall y\,\forall z\,(d \in \mathbb{R} \to y \in \mathbb{R} \to z \in \mathbb{R} \to d \leq y \to d + z \leq y + z)$
\item $\varphi_{29} :\equiv \forall x\,(x \in W \leftrightarrow x \in [a,b] \land f(x) < \mu) \to \forall x \in W\,(x \leq b)$
\item $\varphi_{30} :\equiv d \in \mathbb{R} \land 0 < d \land d \leq \delta \cdot 2^{-1} \land d \leq \xi - a \to \mathrm{UB}(\xi - d,W)$
\item $\varphi_{31} :\equiv \mathrm{IsSup}(\xi,W) \to \exists \xi\,(\xi \in \mathbb{R} \land a < \xi \land \xi < b \land f(\xi) = \mu)$
\item $\varphi_{32} :\equiv \mu < f(b) \to \exists W\,\forall x\,(x \in W \leftrightarrow x \in [a,b] \land f(x) < \mu)$
\item $\varphi_{33} :\equiv \forall z\,\forall w\,(z = w \to z < b \to w < b) \to \forall w\,(d + \xi = w \to d + \xi < b \to w < b)$
\item $\varphi_{34} :\equiv [a,b] \subseteq \mathrm{dom}(f) \to \xi + d \in [a,b] \to \xi + d \in \mathrm{dom}(f)$
\item $\varphi_{35} :\equiv \forall x\,(x \in W \leftrightarrow x \in [a,b] \land f(x) < \mu) \to \mathrm{IsSup}(\xi,W) \to \xi \leq b$
\item $\varphi_{36} :\equiv \forall x\,(x \in W \leftrightarrow x \in [a,b] \land f(x) < \mu) \to \mathrm{IsSup}(\xi,W) \to a \leq \xi$
\item $\varphi_{37} :\equiv 0 < \xi - a \to \exists d\,(d \in \mathbb{R} \land 0 < d \land d \leq \delta \cdot 2^{-1} \land d \leq \xi - a)$
\item $\varphi_{38} :\equiv 0 < b - \xi \to \exists d\,(d \in \mathbb{R} \land 0 < d \land d \leq \delta \cdot 2^{-1} \land d \leq b - \xi)$
\item $\varphi_{39} :\equiv \forall x\,\forall y\,\forall z\,(x \in \mathbb{R} \to y \in \mathbb{R} \to z \in \mathbb{R} \to x \leq y \to x + z \leq y + z)$
\item $\varphi_{40} :\equiv \forall A\,(\mathrm{ContOn}(f,A) \leftrightarrow \mathrm{RFn}(f) \land A \subseteq \mathrm{dom}(f) \land \forall c \in A\,(\mathrm{ContIn}(f,A,c)))$
\item $\varphi_{41} :\equiv \forall x\,\forall y\,\forall z\,(x < y \to y < z \to x < z) \to \forall y\,\forall z\,(\mu < y \to y < z \to \mu < z)$
\item $\varphi_{42} :\equiv \forall x\,(x \in W \leftrightarrow x \in [a,b] \land f(x) < \mu) \to u \in W \to u \in \mathbb{R}$
\item $\varphi_{43} :\equiv d \in \mathbb{R} \land 0 < d \land d \leq \delta \cdot 2^{-1} \land d \leq b - \xi \to \xi + d \in \mathrm{dom}(f)$
\item $\varphi_{44} :\equiv d \in \mathbb{R} \land 0 < d \land d \leq \delta \cdot 2^{-1} \land d \leq \xi - a \to s \in W \to s \leq \xi - d$
\item $\varphi_{45} :\equiv z = w \to \lvert \xi + d - \xi \rvert = \lvert z - \xi \rvert \to \lvert \xi + d - \xi \rvert = \lvert w - \xi \rvert$
\item $\varphi_{46} :\equiv d \in \mathbb{R} \land 0 < d \land d \leq \delta \cdot 2^{-1} \land d \leq \xi - a \to s \in W \to s \in \mathbb{R}$
\item $\varphi_{47} :\equiv d \in \mathbb{R} \land 0 < d \land d \leq \delta \cdot 2^{-1} \land d \leq b - \xi \to \lvert \xi + d - \xi \rvert = d$
\item $\varphi_{48} :\equiv d \in \mathbb{R} \land 0 < d \land d \leq \delta \cdot 2^{-1} \land d \leq b - \xi \to d = \lvert \xi + d - \xi \rvert$
\item $\varphi_{49} :\equiv d \in \mathbb{R} \land 0 < d \land d \leq \delta \cdot 2^{-1} \land d \leq b - \xi \to f(\xi + d) < \mu$
\item $\varphi_{50} :\equiv \forall x\,(x \in W \leftrightarrow x \in [a,b] \land f(x) < \mu) \to \forall u \in W\,(u \in \mathbb{R})$
\item $\varphi_{51} :\equiv \forall w\,(d + \xi = w \to d + \xi = b \to w = b) \to d + \xi = \xi + d \to d + \xi = b \to \xi + d = b$
\item $\varphi_{52} :\equiv \mu < f(b) \to \forall x\,(x \in W \leftrightarrow x \in [a,b] \land f(x) < \mu) \to \mathrm{UB}(b,W)$
\item $\varphi_{53} :\equiv \mu < f(b) \to \forall x\,(x \in W \leftrightarrow x \in [a,b] \land f(x) < \mu) \to a \in W$
\item $\varphi_{54} :\equiv d \in \mathbb{R} \land 0 < d \land d \leq \delta \cdot 2^{-1} \land d \leq b - \xi \to d + \xi \leq b - \xi + \xi$
\item $\varphi_{55} :\equiv d \in \mathbb{R} \land 0 < d \land d \leq \delta \cdot 2^{-1} \land d \leq b - \xi \to \xi + d \in [a,b]$
\item $\varphi_{56} :\equiv \mathrm{ContOn}(f,[a,b]) \to a \in \mathbb{R} \to b \in \mathbb{R} \to \mu \in \mathbb{R} \to \mu \in \mathbb{R}$
\item $\varphi_{57} :\equiv \forall x\,\forall a\,\forall y\,(x \in \mathbb{R} \to a \in \mathbb{R} \to y \in \mathbb{R} \to x - a < y \to x < y + a)$
\item $\varphi_{58} :\equiv \forall y\,(f(\xi) \in \mathbb{R} \to y \in \mathbb{R} \to f(\xi) < y \to 0 < y - f(\xi))$
\item $\varphi_{59} :\equiv \forall x\,(x \in W \leftrightarrow x \in [a,b] \land f(x) < \mu) \to \mathrm{IsSup}(\xi,W) \to \mathrm{IsSup}(\xi,W)$
\item $\varphi_{60} :\equiv \forall x\,(x \in W \leftrightarrow x \in [a,b] \land f(x) < \mu) \to \mathrm{IsSup}(\xi,W) \to \xi \in \mathbb{R}$
\item $\varphi_{61} :\equiv b \in \mathbb{R} \to \mu \in \mathbb{R} \to a < b \to f(a) < \mu \to \mu < f(b) \to a \leq b$
\item $\varphi_{62} :\equiv b \in \mathbb{R} \to \mu \in \mathbb{R} \to a < b \to f(a) < \mu \to \mu < f(b) \to a \leq a$
\item $\varphi_{63} :\equiv d \in \mathbb{R} \land 0 < d \land d \leq \delta \cdot 2^{-1} \land d \leq \xi - a \to \forall s \in W\,(s \leq \xi - d)$
\item $\varphi_{64} :\equiv \forall z\,\forall w\,(z = w \to \mu < z \to \mu < w) \to \forall w\,(f(\xi) = w \to \mu < f(\xi) \to \mu < w)$
\item $\varphi_{65} :\equiv \forall z\,\forall w\,(z = w \to z < \mu \to w < \mu) \to \forall w\,(f(\xi) = w \to f(\xi) < \mu \to w < \mu)$
\item $\varphi_{66} :\equiv \forall x\,\forall y\,(x \subseteq y \leftrightarrow \forall u \in x\,(u \in y)) \to \forall y\,(W \subseteq y \leftrightarrow \forall u \in W\,(u \in y))$
\item $\varphi_{67} :\equiv \forall y\,(W \subseteq y \leftrightarrow \forall u \in W\,(u \in y)) \to (W \subseteq \mathbb{R} \leftrightarrow \forall u \in W\,(u \in \mathbb{R}))$
\item $\varphi_{68} :\equiv d \in \mathbb{R} \land 0 < d \land d \leq \delta \cdot 2^{-1} \land d \leq b - \xi \to f(\xi + d) \in \mathbb{R}$
\item $\varphi_{69} :\equiv \forall x\,(x \in W \leftrightarrow x \in [a,b] \land f(x) < \mu) \to \mathrm{IsSup}(\xi,W) \to \xi \in \mathrm{dom}(f)$
\item $\varphi_{70} :\equiv \mathrm{IsSup}(\xi,W) \to f(\xi) = \mu \to \xi \in \mathbb{R} \land a < \xi \land \xi < b \land f(\xi) = \mu$
\item $\varphi_{71} :\equiv d \in \mathbb{R} \land 0 < d \land d \leq \delta \cdot 2^{-1} \land d \leq b - \xi \to \lvert \xi + d - \xi \rvert < \delta$
\item $\varphi_{72} :\equiv d \in \mathbb{R} \land 0 < d \land d \leq \delta \cdot 2^{-1} \land d \leq b - \xi \to \mu - f(\xi) \in \mathbb{R}$
\item $\varphi_{73} :\equiv \forall s \in W\,(\neg s < \xi - d \to s = \xi - d) \to x \in W \to \neg x < \xi - d \to x = \xi - d$
\item $\varphi_{74} :\equiv \forall x \in W\,(\neg x < \xi - d \to x = \xi - d) \to s \in W \to \neg s < \xi - d \to s = \xi - d$
\item $\varphi_{75} :\equiv \forall y\,\forall z\,(d \leq y \to y < z \to d < z) \to \forall z\,(d \leq \delta \cdot 2^{-1} \to \delta \cdot 2^{-1} < z \to d < z)$
\item $\varphi_{76} :\equiv d \in \mathbb{R} \land 0 < d \land d \leq \delta \cdot 2^{-1} \land d \leq \xi - a \to s \in W \to s \in [a,b]$
\item $\varphi_{77} :\equiv \forall x\,(x \in W \leftrightarrow x \in [a,b] \land f(x) < \mu) \to \mathrm{IsSup}(\xi,W) \to s \in W \to s \leq \xi$
\item $\varphi_{78} :\equiv \forall x\,(x \in W \leftrightarrow x \in [a,b] \land f(x) < \mu) \to \mathrm{IsSup}(\xi,W) \to \mathrm{UB}(b,W) \to \xi \leq b$
\item $\varphi_{79} :\equiv \forall x\,(x \in W \leftrightarrow x \in [a,b] \land f(x) < \mu) \to \mathrm{IsSup}(\xi,W) \to a \in W \to a \leq \xi$
\item $\varphi_{80} :\equiv b \in \mathbb{R} \to \mu \in \mathbb{R} \to a < b \to f(a) < \mu \to \mu < f(b) \to \mathrm{RFn}(f)$
\item $\varphi_{81} :\equiv d \in \mathbb{R} \land 0 < d \land d \leq \delta \cdot 2^{-1} \land d \leq b - \xi \to \lvert d + \xi - \xi \rvert = \lvert d \rvert$
\item $\varphi_{82} :\equiv d \in \mathbb{R} \land 0 < d \land d \leq \delta \cdot 2^{-1} \land d \leq b - \xi \to \lvert \xi + d - \xi \rvert = \lvert d \rvert$
\item $\varphi_{83} :\equiv \mu < f(b) \to \forall x\,(x \in W \leftrightarrow x \in [a,b] \land f(x) < \mu) \to W \neq \emptyset$
\item $\varphi_{84} :\equiv \mu < f(b) \to \forall x\,(x \in W \leftrightarrow x \in [a,b] \land f(x) < \mu) \to W \subseteq \mathbb{R}$
\item $\varphi_{85} :\equiv \mathrm{ContOn}(f,[a,b]) \to a \in \mathbb{R} \to b \in \mathbb{R} \to \mu \in \mathbb{R} \to a < b \to a < b$
\item $\varphi_{86} :\equiv \lvert \xi + d - \xi \rvert = \lvert \xi + d - \xi \rvert \to \lvert \xi + d - \xi \rvert = \lvert d + \xi - \xi \rvert$
\item $\varphi_{87} :\equiv \forall x\,\forall m\,\forall y\,(x \in \mathbb{R} \to m \in \mathbb{R} \to y \in \mathbb{R} \to -(x - m) < y - x \to m < y)$
\item $\varphi_{88} :\equiv \forall y\,\forall z\,(\mu < y \to y < z \to \mu < z) \to \forall z\,(\mu < f(s) \to f(s) < z \to \mu < z)$
\item $\varphi_{89} :\equiv \forall y\,\forall z\,(\mu < y \to y < z \to \mu < z) \to \forall z\,(\mu < f(b) \to f(b) < z \to \mu < z)$
\item $\varphi_{90} :\equiv \mu < f(b) \to \forall x\,(x \in W \leftrightarrow x \in [a,b] \land f(x) < \mu) \to x \in W \to x \leq b$
\item $\varphi_{91} :\equiv \forall x\,(x \in W \leftrightarrow x \in [a,b] \land f(x) < \mu) \to \mathrm{IsSup}(\xi,W) \to \xi \in [a,b]$
\item $\varphi_{92} :\equiv \mu \in \mathbb{R} \to a < b \to f(a) < \mu \to \mu < f(b) \to [a,b] \subseteq \mathrm{dom}(f)$
\item $\varphi_{93} :\equiv \forall s \in W\,(\neg s < \xi - d \to s = \xi - d) \to \forall x \in W\,(\neg x < \xi - d \to x = \xi - d)$
\item $\varphi_{94} :\equiv \forall x \in W\,(\neg x < \xi - d \to x = \xi - d) \to \forall s \in W\,(\neg s < \xi - d \to s = \xi - d)$
\item $\varphi_{95} :\equiv s \in \mathbb{R} \to \xi \in \mathbb{R} \to d \in \mathbb{R} \to s \leq \xi \to \xi - d < s \to \lvert s - \xi \rvert < d$
\item $\varphi_{96} :\equiv \forall z\,(d \in \mathbb{R} \to b - \xi \in \mathbb{R} \to z \in \mathbb{R} \to d \leq b - \xi \to d + z \leq b - \xi + z)$
\item $\varphi_{97} :\equiv d \in \mathbb{R} \to b - \xi \in \mathbb{R} \to \xi \in \mathbb{R} \to d \leq b - \xi \to d + \xi \leq b - \xi + \xi$
\item $\varphi_{98} :\equiv \forall x\,(x \in W \leftrightarrow x \in [a,b] \land f(x) < \mu) \to \mathrm{IsSup}(\xi,W) \to \forall x \in W\,(x \leq \xi)$
\item $\varphi_{99} :\equiv \forall x\,(x \in W \leftrightarrow x \in [a,b] \land f(x) < \mu) \to \mathrm{IsSup}(\xi,W) \to \forall w\,(\mathrm{UB}(w,W) \to \xi \leq w)$
\item $\varphi_{100} :\equiv \mathrm{IsSup}(\xi,W) \to f(\xi) = \mu \to \exists \xi\,(\xi \in \mathbb{R} \land a < \xi \land \xi < b \land f(\xi) = \mu)$
\item $\varphi_{101} :\equiv \forall y\,\forall z\,(\mu < y \to y < z \to \mu < z) \to \forall z\,(\mu < f(\xi) \to f(\xi) < z \to \mu < z)$
\item $\varphi_{102} :\equiv a \in \mathbb{R} \to b \in \mathbb{R} \to \mu \in \mathbb{R} \to a < b \to f(a) < \mu \to f(a) < \mu$
\item $\varphi_{103} :\equiv \forall x\,(x \in W \leftrightarrow x \in [a,b] \land f(x) < \mu) \to \mathrm{IsSup}(\xi,W) \to f(\xi) \in \mathbb{R}$
\item $\varphi_{104} :\equiv \forall x\,\forall d\,(s \in \mathbb{R} \to x \in \mathbb{R} \to d \in \mathbb{R} \to s \leq x \to x - d < s \to \lvert s - x \rvert < d)$
\item $\varphi_{105} :\equiv \forall s \in W\,(\neg s < \xi - d \to s = \xi - d) \to \forall x\,\forall s \in W\,(\neg s < \xi - d \to s = \xi - d)$
\item $\varphi_{106} :\equiv \forall x \in W\,(\neg x < \xi - d \to x = \xi - d) \to \forall s\,\forall x \in W\,(\neg x < \xi - d \to x = \xi - d)$
\item $\varphi_{107} :\equiv \forall w\,(a = w \to f(a) < \mu \to f(w) < \mu) \to a = \xi \to f(a) < \mu \to f(\xi) < \mu$
\item $\varphi_{108} :\equiv \forall w\,(b = w \to \mu < f(b) \to \mu < f(w)) \to b = \xi \to \mu < f(b) \to \mu < f(\xi)$
\item $\varphi_{109} :\equiv \forall x\,\forall y\,\forall z\,(x = y \to y = z \to x = z) \to \forall y\,\forall z\,(\lvert \xi + d - \xi \rvert = y \to y = z \to \lvert \xi + d - \xi \rvert = z)$
\item $\varphi_{110} :\equiv \forall y\,(f(\xi) \in \mathbb{R} \to y \in \mathbb{R} \to f(\xi) < y \lor f(\xi) = y \lor y < f(\xi))$
\item $\varphi_{111} :\equiv \forall x\,\forall s \in W\,(\neg s < \xi - d \to s = \xi - d) \to \forall x \in W\,(\neg x < \xi - d \to x = \xi - d)$
\item $\varphi_{112} :\equiv \forall s\,\forall x \in W\,(\neg x < \xi - d \to x = \xi - d) \to \forall s \in W\,(\neg s < \xi - d \to s = \xi - d)$
\item $\varphi_{113} :\equiv \mu < f(b) \to \forall x\,(x \in W \leftrightarrow x \in [a,b] \land f(x) < \mu) \to x \in W \to x \in W$
\item $\varphi_{114} :\equiv f(s) \in \mathbb{R} \to -(f(\xi) - \mu) < f(s) - f(\xi) \to \mu < f(s)$
\item $\varphi_{115} :\equiv d \in \mathbb{R} \land 0 < d \land d \leq \delta \cdot 2^{-1} \land d \leq \xi - a \to s \in W \to \neg s \leq \xi - d \to \mu < \mu$
\item $\varphi_{116} :\equiv \mathrm{IsSup}(\xi,W) \to \mu < f(\xi) \to \exists \xi\,(\xi \in \mathbb{R} \land a < \xi \land \xi < b \land f(\xi) = \mu)$
\item $\varphi_{117} :\equiv \forall x\,(x \in W \leftrightarrow x \in [a,b] \land f(x) < \mu) \to \mathrm{IsSup}(\xi,W) \to f(\xi) = \mu \to \xi \neq b$
\item $\varphi_{118} :\equiv \forall x\,(x \in W \leftrightarrow x \in [a,b] \land f(x) < \mu) \to \mathrm{IsSup}(\xi,W) \to f(\xi) = \mu \to a \neq \xi$
\item $\varphi_{119} :\equiv d \in \mathbb{R} \land 0 < d \land d \leq \delta \cdot 2^{-1} \land d \leq b - \xi \to \mu - f(\xi) + f(\xi) = \mu$
\item $\varphi_{120} :\equiv \mathrm{IsSup}(\xi,W) \to f(\xi) < \mu \to \exists \xi\,(\xi \in \mathbb{R} \land a < \xi \land \xi < b \land f(\xi) = \mu)$
\item $\varphi_{121} :\equiv \mu < f(b) \to \forall x\,(x \in W \leftrightarrow x \in [a,b] \land f(x) < \mu) \to u \in W \to u \in W$
\item $\varphi_{122} :\equiv d + \xi - \xi = d \to \lvert d + \xi - \xi \rvert = \lvert d + \xi - \xi \rvert \to \lvert d + \xi - \xi \rvert = \lvert d \rvert$
\item $\varphi_{123} :\equiv f(\xi + d) \in \mathbb{R} \to f(\xi) \in \mathbb{R} \to f(\xi + d) - f(\xi) \in \mathbb{R}$
\item $\varphi_{124} :\equiv f(\xi) \in \mathbb{R} \to \mu \in \mathbb{R} \to f(\xi) < \mu \lor f(\xi) = \mu \lor \mu < f(\xi)$
\item $\varphi_{125} :\equiv d \in \mathbb{R} \land 0 < d \land d \leq \delta \cdot 2^{-1} \land d \leq \xi - a \to s \in W \to \neg s \leq \xi - d \to \neg \mu < \mu$
\item $\varphi_{126} :\equiv \forall X\,(\mathrm{UB}(b,X) \leftrightarrow b \in \mathbb{R} \land \forall x \in X\,(x \leq b)) \to (\mathrm{UB}(b,W) \leftrightarrow b \in \mathbb{R} \land \forall x \in W\,(x \leq b))$
\item $\varphi_{127} :\equiv \forall p\,\forall q\,(p \in \mathbb{R} \to q \in \mathbb{R} \to 0 < p \to 0 < q \to \exists d\,(d \in \mathbb{R} \land 0 < d \land d \leq p \land d \leq q))$
\item $\varphi_{128} :\equiv \forall w\,(d + \xi = w \to d + \xi < b \to w < b) \to d + \xi = \xi + d \to d + \xi < b \to \xi + d < b$
\item $\varphi_{129} :\equiv \forall w\,(\xi = w \to f(\xi) < \mu \to f(w) < \mu) \to \xi = b \to f(\xi) < \mu \to f(b) < \mu$
\item $\varphi_{130} :\equiv d \in \mathbb{R} \land 0 < d \land d \leq \delta \cdot 2^{-1} \land d \leq \xi - a \to s \in W \to \neg s \leq \xi - d \to \neg s \leq \xi - d$
\item $\varphi_{131} :\equiv s \in [a,b] \to \lvert s - \xi \rvert < \delta \to \lvert f(s) - f(\xi) \rvert < f(\xi) - \mu$
\item $\varphi_{132} :\equiv b \in \mathbb{R} \to \mu \in \mathbb{R} \to a < b \to f(a) < \mu \to \mu < f(b) \to \mu < f(b)$
\item $\varphi_{133} :\equiv f(a) < \mu \to \mu < f(b) \to \exists \xi\,(\xi \in \mathbb{R} \land a < \xi \land \xi < b \land f(\xi) = \mu)$
\item $\varphi_{134} :\equiv \forall e \in \mathbb{R}\,(0 < e \to 0 < e \cdot 2^{-1}) \to \delta \in \mathbb{R} \to 0 < \delta \to 0 < \delta \cdot 2^{-1}$
\item $\varphi_{135} :\equiv b \in \mathbb{R} \to \mu \in \mathbb{R} \to a < b \to f(a) < \mu \to \mu < f(b) \to a \in [a,b]$
\item $\varphi_{136} :\equiv \forall x\,(x \in W \leftrightarrow x \in [a,b] \land f(x) < \mu) \to \mathrm{IsSup}(\xi,W) \to \mu < f(\xi) \to a \neq \xi$
\item $\varphi_{137} :\equiv \forall x\,(x \in W \leftrightarrow x \in [a,b] \land f(x) < \mu) \to \mathrm{IsSup}(\xi,W) \to f(\xi) < \mu \to \xi \neq b$
\item $\varphi_{138} :\equiv \forall e \in \mathbb{R}\,(0 < e \to e \cdot 2^{-1} < e) \to \delta \in \mathbb{R} \to 0 < \delta \to \delta \cdot 2^{-1} < \delta$
\item $\varphi_{139} :\equiv \forall x\,(x \in [a,b] \leftrightarrow x \in \mathbb{R} \land a \leq x \land x \leq b) \to (s \in [a,b] \leftrightarrow s \in \mathbb{R} \land a \leq s \land s \leq b)$
\item $\varphi_{140} :\equiv \forall x\,(x \in [a,b] \leftrightarrow x \in \mathbb{R} \land a \leq x \land x \leq b) \to (x \in [a,b] \leftrightarrow x \in \mathbb{R} \land a \leq x \land x \leq b)$
\item $\varphi_{141} :\equiv \forall x\,(x \in [a,b] \leftrightarrow x \in \mathbb{R} \land a \leq x \land x \leq b) \to (a \in [a,b] \leftrightarrow a \in \mathbb{R} \land a \leq a \land a \leq b)$
\item $\varphi_{142} :\equiv \forall z\,\forall w\,(z = w \to f(z) < \mu \to f(w) < \mu) \to \forall w\,(a = w \to f(a) < \mu \to f(w) < \mu)$
\item $\varphi_{143} :\equiv \forall c \in [a,b]\,(\mathrm{ContIn}(f,[a,b],c)) \to \xi \in [a,b] \to \mathrm{ContIn}(f,[a,b],\xi)$
\item $\varphi_{144} :\equiv \forall w\,(f(\xi) = w \to \mu < f(\xi) \to \mu < w) \to f(\xi) = \mu \to \mu < f(\xi) \to \mu < \mu$
\item $\varphi_{145} :\equiv \forall z\,\forall w\,(z = w \to \mu < f(z) \to \mu < f(w)) \to \forall w\,(b = w \to \mu < f(b) \to \mu < f(w))$
\item $\varphi_{146} :\equiv \forall w\,(f(\xi) = w \to f(\xi) < \mu \to w < \mu) \to f(\xi) = \mu \to f(\xi) < \mu \to \mu < \mu$
\item $\varphi_{147} :\equiv \forall w\,(d = w \to d < \delta \to w < \delta) \to d = \lvert \xi + d - \xi \rvert \to d < \delta \to \lvert \xi + d - \xi \rvert < \delta$
\item $\varphi_{148} :\equiv \forall z\,(\lvert \xi + d - \xi \rvert = \lvert d + \xi - \xi \rvert \to \lvert d + \xi - \xi \rvert = z \to \lvert \xi + d - \xi \rvert = z)$
\item $\varphi_{149} :\equiv \forall w\,(d + \xi - \xi = w \to \lvert d + \xi - \xi \rvert = \lvert d + \xi - \xi \rvert \to \lvert d + \xi - \xi \rvert = \lvert w \rvert)$
\item $\varphi_{150} :\equiv \forall w\,(\xi + d = w \to \lvert \xi + d - \xi \rvert = \lvert \xi + d - \xi \rvert \to \lvert \xi + d - \xi \rvert = \lvert w - \xi \rvert)$
\item $\varphi_{151} :\equiv \forall x\,(x \in [a,b] \leftrightarrow x \in \mathbb{R} \land a \leq x \land x \leq b) \to (u \in [a,b] \leftrightarrow u \in \mathbb{R} \land a \leq u \land u \leq b)$
\item $\varphi_{152} :\equiv d \in \mathbb{R} \land 0 < d \land d \leq \delta \cdot 2^{-1} \land d \leq \xi - a \to s \in W \to \neg s \leq \xi - d \to \xi - d < s$
\item $\varphi_{153} :\equiv d \in \mathbb{R} \land 0 < d \land d \leq \delta \cdot 2^{-1} \land d \leq \xi - a \to s \in W \to \neg s \leq \xi - d \to s \in \mathrm{dom}(f)$
\item $\varphi_{154} :\equiv \forall y\,(d \in \mathbb{R} \to y \in \mathbb{R} \to d + y - y = d) \to d \in \mathbb{R} \to \xi \in \mathbb{R} \to d + \xi - \xi = d$
\item $\varphi_{155} :\equiv \forall y\,(b \in \mathbb{R} \to y \in \mathbb{R} \to b - y + y = b) \to b \in \mathbb{R} \to \xi \in \mathbb{R} \to b - \xi + \xi = b$
\item $\varphi_{156} :\equiv \forall y\,(d \in \mathbb{R} \to y \in \mathbb{R} \to d + y = y + d) \to d \in \mathbb{R} \to \xi \in \mathbb{R} \to d + \xi = \xi + d$
\item $\varphi_{157} :\equiv \forall x \in [a,b]\,(\lvert x - \xi \rvert < \delta \to \lvert f(x) - f(\xi) \rvert < f(\xi) - \mu)$
\item $\varphi_{158} :\equiv \forall x\,(x \in W \leftrightarrow x \in [a,b] \land f(x) < \mu) \to \mathrm{IsSup}(\xi,W) \to \mu < f(\xi) \to a < \xi$
\item $\varphi_{159} :\equiv \forall x \in [a,b]\,(\lvert x - \xi \rvert < \delta \to \lvert f(x) - f(\xi) \rvert < \mu - f(\xi))$
\item $\varphi_{160} :\equiv d \in \mathbb{R} \land 0 < d \land d \leq \delta \cdot 2^{-1} \land d \leq b - \xi \to f(\xi + d) - f(\xi) \in \mathbb{R}$
\item $\varphi_{161} :\equiv \forall x\,(x \in W \leftrightarrow x \in [a,b] \land f(x) < \mu) \to \mathrm{IsSup}(\xi,W) \to f(\xi) < \mu \to \xi < b$
\item $\varphi_{162} :\equiv \forall x\,\forall y\,\forall z\,(x < y \to y < z \to x < z) \to \forall y\,\forall z\,(\lvert s - \xi \rvert < y \to y < z \to \lvert s - \xi \rvert < z)$
\item $\varphi_{163} :\equiv \forall x\,(\mathrm{RFn}(f) \to x \in \mathrm{dom}(f) \to f(x) \in \mathbb{R}) \to \mathrm{RFn}(f) \to s \in \mathrm{dom}(f) \to f(s) \in \mathbb{R}$
\item $\varphi_{164} :\equiv \forall z\,(\mu < f(s) \to f(s) < z \to \mu < z) \to \mu < f(s) \to f(s) < \mu \to \mu < \mu$
\item $\varphi_{165} :\equiv \forall z\,\forall w\,(z = w \to f(z) < \mu \to f(w) < \mu) \to \forall w\,(\xi = w \to f(\xi) < \mu \to f(w) < \mu)$
\item $\varphi_{166} :\equiv \forall z\,(\mu < f(b) \to f(b) < z \to \mu < z) \to \mu < f(b) \to f(b) < \mu \to \mu < \mu$
\item $\varphi_{167} :\equiv d \in \mathbb{R} \land 0 < d \land d \leq \delta \cdot 2^{-1} \land d \leq \xi - a \to s \in W \to \neg s \leq \xi - d \to f(s) < \mu$
\item $\varphi_{168} :\equiv -(f(\xi) - \mu) < f(s) - f(\xi) \land f(s) - f(\xi) < f(\xi) - \mu$
\item $\varphi_{169} :\equiv d \in \mathbb{R} \land 0 < d \land d \leq \delta \cdot 2^{-1} \land d \leq \xi - a \to s \in W \to \neg s \leq \xi - d \to \mu < f(s)$
\item $\varphi_{170} :\equiv \forall x\,(x \in W \leftrightarrow x \in [a,b] \land f(x) < \mu) \to \mathrm{IsSup}(\xi,W) \to \mathrm{UB}(\xi - d,W) \to \xi \leq \xi - d$
\item $\varphi_{171} :\equiv d \in \mathbb{R} \land 0 < d \land d \leq \delta \cdot 2^{-1} \land d \leq b - \xi \to \lvert \xi + d - \xi \rvert = \lvert d + \xi - \xi \rvert$
\item $\varphi_{172} :\equiv \forall x\,(x \in W \leftrightarrow x \in [a,b] \land f(x) < \mu) \to \mathrm{IsSup}(\xi,W) \to \xi + d \in W \to \xi + d \leq \xi$
\item $\varphi_{173} :\equiv \forall x\,(x \in W \leftrightarrow x \in [a,b] \land f(x) < \mu) \to \mathrm{IsSup}(\xi,W) \to f(\xi) = \mu \to f(\xi) = \mu$
\item $\varphi_{174} :\equiv \forall X\,(\mathrm{UB}(\xi,X) \leftrightarrow \xi \in \mathbb{R} \land \forall x \in X\,(x \leq \xi)) \to (\mathrm{UB}(\xi,W) \leftrightarrow \xi \in \mathbb{R} \land \forall x \in W\,(x \leq \xi))$
\item $\varphi_{175} :\equiv \forall x\,(x \in W \leftrightarrow x \in [a,b] \land f(x) < \mu) \to (s \in W \leftrightarrow s \in [a,b] \land f(s) < \mu)$
\item $\varphi_{176} :\equiv \forall x\,(x \in W \leftrightarrow x \in [a,b] \land f(x) < \mu) \to (x \in W \leftrightarrow x \in [a,b] \land f(x) < \mu)$
\item $\varphi_{177} :\equiv \forall a\,\forall X\,(\mathrm{UB}(a,X) \leftrightarrow a \in \mathbb{R} \land \forall x \in X\,(x \leq a)) \to \forall X\,(\mathrm{UB}(b,X) \leftrightarrow b \in \mathbb{R} \land \forall x \in X\,(x \leq b))$
\item $\varphi_{178} :\equiv \forall x\,(x \in W \leftrightarrow x \in [a,b] \land f(x) < \mu) \to (a \in W \leftrightarrow a \in [a,b] \land f(a) < \mu)$
\item $\varphi_{179} :\equiv \forall x\,(x \in [a,b] \leftrightarrow x \in \mathbb{R} \land a \leq x \land x \leq b) \to (\xi \in [a,b] \leftrightarrow \xi \in \mathbb{R} \land a \leq \xi \land \xi \leq b)$
\item $\varphi_{180} :\equiv \forall X\,\forall Y\,\forall x\,(X \subseteq Y \to x \in X \to x \in Y) \to \forall Y\,\forall x\,([a,b] \subseteq Y \to x \in [a,b] \to x \in Y)$
\item $\varphi_{181} :\equiv \forall x\,(\mathrm{RFn}(f) \to x \in \mathrm{dom}(f) \to f(x) \in \mathbb{R}) \to \mathrm{RFn}(f) \to \xi \in \mathrm{dom}(f) \to f(\xi) \in \mathbb{R}$
\item $\varphi_{182} :\equiv \forall x\,\forall y\,(x \in \mathbb{R} \to y \in \mathbb{R} \to x + y - y = x) \to \forall y\,(d \in \mathbb{R} \to y \in \mathbb{R} \to d + y - y = d)$
\item $\varphi_{183} :\equiv \forall y\,(\xi \in \mathbb{R} \to y \in \mathbb{R} \to \xi + y = y + \xi) \to \xi \in \mathbb{R} \to d \in \mathbb{R} \to \xi + d = d + \xi$
\item $\varphi_{184} :\equiv \forall z\,\forall w\,(z = w \to d + \xi = z \to d + \xi = w) \to \forall w\,(b - \xi + \xi = w \to d + \xi = b - \xi + \xi \to d + \xi = w)$
\item $\varphi_{185} :\equiv \forall y\,(b \in \mathbb{R} \to y \in \mathbb{R} \to b - y \in \mathbb{R}) \to b \in \mathbb{R} \to \xi \in \mathbb{R} \to b - \xi \in \mathbb{R}$
\item $\varphi_{186} :\equiv \forall x\,\forall y\,(x \in \mathbb{R} \to y \in \mathbb{R} \to x - y + y = x) \to \forall y\,(b \in \mathbb{R} \to y \in \mathbb{R} \to b - y + y = b)$
\item $\varphi_{187} :\equiv \forall x\,\forall y\,(x \in \mathbb{R} \to y \in \mathbb{R} \to x + y = y + x) \to \forall y\,(d \in \mathbb{R} \to y \in \mathbb{R} \to d + y = y + d)$
\item $\varphi_{188} :\equiv \forall x\,(x \in W \leftrightarrow x \in [a,b] \land f(x) < \mu) \to (u \in W \leftrightarrow u \in [a,b] \land f(u) < \mu)$
\item $\varphi_{189} :\equiv d \in \mathbb{R} \land 0 < d \land d \leq \delta \cdot 2^{-1} \land d \leq \xi - a \to s \in W \to \neg s \leq \xi - d \to \xi - d \in \mathbb{R}$
\item $\varphi_{190} :\equiv \forall x\,(x \in W \leftrightarrow x \in [a,b] \land f(x) < \mu) \to \mathrm{IsSup}(\xi,W) \to f(\xi) = \mu \to \neg \xi \neq b \to \xi = b$
\item $\varphi_{191} :\equiv \forall x\,(x \in W \leftrightarrow x \in [a,b] \land f(x) < \mu) \to \mathrm{IsSup}(\xi,W) \to f(\xi) = \mu \to \neg \xi \neq b \to b = \xi$
\item $\varphi_{192} :\equiv \forall x\,(x \in W \leftrightarrow x \in [a,b] \land f(x) < \mu) \to \mathrm{IsSup}(\xi,W) \to f(\xi) = \mu \to \neg a \neq \xi \to a = \xi$
\item $\varphi_{193} :\equiv d \in \mathbb{R} \land 0 < d \land d \leq \delta \cdot 2^{-1} \land d \leq \xi - a \to s \in W \to \neg s \leq \xi - d \to \lvert s - \xi \rvert < d$
\item $\varphi_{194} :\equiv d \in \mathbb{R} \land 0 < d \land d \leq \delta \cdot 2^{-1} \land d \leq \xi - a \to s \in W \to \neg s \leq \xi - d \to f(s) \in \mathbb{R}$
\item $\varphi_{195} :\equiv \forall b\,\forall x\,(x \in [a,b] \leftrightarrow x \in \mathbb{R} \land a \leq x \land x \leq b) \to \forall x\,(x \in [a,b] \leftrightarrow x \in \mathbb{R} \land a \leq x \land x \leq b)$
\item $\varphi_{196} :\equiv \forall y\,(\xi \in \mathbb{R} \to y \in \mathbb{R} \to \xi - y \in \mathbb{R}) \to \xi \in \mathbb{R} \to d \in \mathbb{R} \to \xi - d \in \mathbb{R}$
\item $\varphi_{197} :\equiv \forall z\,(\mu < f(\xi) \to f(\xi) < z \to \mu < z) \to \mu < f(\xi) \to f(\xi) < \mu \to \mu < \mu$
\item $\varphi_{198} :\equiv \forall y\,(\xi \in \mathbb{R} \to y \in \mathbb{R} \to \xi - y \in \mathbb{R}) \to \xi \in \mathbb{R} \to a \in \mathbb{R} \to \xi - a \in \mathbb{R}$
\item $\varphi_{199} :\equiv \forall y\,(\lvert \xi + d - \xi \rvert = y \to y = \lvert \xi + d - \xi \rvert) \to \lvert \xi + d - \xi \rvert = d \to d = \lvert \xi + d - \xi \rvert$
\item $\varphi_{200} :\equiv \forall y\,(\xi \in \mathbb{R} \to y \in \mathbb{R} \to 0 \leq y \to \xi \leq \xi + y) \to \xi \in \mathbb{R} \to d \in \mathbb{R} \to 0 \leq d \to \xi \leq \xi + d$
\item $\varphi_{201} :\equiv \forall y\,(\xi \in \mathbb{R} \to y \in \mathbb{R} \to \xi + y \in \mathbb{R}) \to \xi \in \mathbb{R} \to d \in \mathbb{R} \to \xi + d \in \mathbb{R}$
\item $\varphi_{202} :\equiv d \in \mathbb{R} \land 0 < d \land d \leq \delta \cdot 2^{-1} \land d \leq \xi - a \to s \in W \to \neg s \leq \xi - d \to \lvert s - \xi \rvert < \delta$
\item $\varphi_{203} :\equiv \forall x\,(x \in W \leftrightarrow x \in [a,b] \land f(x) < \mu) \to \mathrm{IsSup}(\xi,W) \to f(\xi) = \mu \to \neg \xi \neq b \to \neg \xi \neq b$
\item $\varphi_{204} :\equiv \forall x\,(x \in W \leftrightarrow x \in [a,b] \land f(x) < \mu) \to \mathrm{IsSup}(\xi,W) \to f(\xi) = \mu \to \neg a \neq \xi \to \neg a \neq \xi$
\item $\varphi_{205} :\equiv \forall a\,\forall X\,(\mathrm{UB}(a,X) \leftrightarrow a \in \mathbb{R} \land \forall x \in X\,(x \leq a)) \to \forall X\,(\mathrm{UB}(\xi,X) \leftrightarrow \xi \in \mathbb{R} \land \forall x \in X\,(x \leq \xi))$
\item $\varphi_{206} :\equiv \lvert \xi + d - \xi \rvert = \lvert d + \xi - \xi \rvert \to \lvert d + \xi - \xi \rvert = \lvert d \rvert \to \lvert \xi + d - \xi \rvert = \lvert d \rvert$
\item $\varphi_{207} :\equiv \forall x\,\forall y\,(x \in \mathbb{R} \to y \in \mathbb{R} \to x + y = y + x) \to \forall y\,(\xi \in \mathbb{R} \to y \in \mathbb{R} \to \xi + y = y + \xi)$
\item $\varphi_{208} :\equiv \forall x\,\forall y\,(x \in \mathbb{R} \to y \in \mathbb{R} \to x - y + y = x) \to \forall y\,(\mu \in \mathbb{R} \to y \in \mathbb{R} \to \mu - y + y = \mu)$
\item $\varphi_{209} :\equiv \forall \mu\,\forall A\,\exists W\,\forall x\,(x \in W \leftrightarrow x \in A \land f(x) < \mu) \to \forall A\,\exists W\,\forall x\,(x \in W \leftrightarrow x \in A \land f(x) < \mu)$
\item $\varphi_{210} :\equiv \forall x\,(x \in W \leftrightarrow x \in [a,b] \land f(x) < \mu) \to \forall x\,(x \in W \leftrightarrow x \in [a,b] \land f(x) < \mu)$
\item $\varphi_{211} :\equiv \forall x\,(x \in W \leftrightarrow x \in [a,b] \land f(x) < \mu) \to \mathrm{IsSup}(\xi,W) \to \mu < f(\xi) \to \neg a \neq \xi \to a = \xi$
\item $\varphi_{212} :\equiv \forall x\,(x \in W \leftrightarrow x \in [a,b] \land f(x) < \mu) \to \mathrm{IsSup}(\xi,W) \to f(\xi) < \mu \to \neg \xi \neq b \to \xi = b$
\item $\varphi_{213} :\equiv \forall z\,(d \leq \delta \cdot 2^{-1} \to \delta \cdot 2^{-1} < z \to d < z) \to d \leq \delta \cdot 2^{-1} \to \delta \cdot 2^{-1} < \delta \to d < \delta$
\item $\varphi_{214} :\equiv \forall f\,\forall x\,(\mathrm{RFn}(f) \to x \in \mathrm{dom}(f) \to f(x) \in \mathbb{R}) \to \forall x\,(\mathrm{RFn}(f) \to x \in \mathrm{dom}(f) \to f(x) \in \mathbb{R})$
\item $\varphi_{215} :\equiv \forall x\,\forall y\,(x \in \mathbb{R} \to y \in \mathbb{R} \to x - y \in \mathbb{R}) \to \forall y\,(b \in \mathbb{R} \to y \in \mathbb{R} \to b - y \in \mathbb{R})$
\item $\varphi_{216} :\equiv \forall A\,\exists W\,\forall x\,(x \in W \leftrightarrow x \in A \land f(x) < \mu) \to \exists W\,\forall x\,(x \in W \leftrightarrow x \in [a,b] \land f(x) < \mu)$
\item $\varphi_{217} :\equiv \forall x\,(x \in W \leftrightarrow x \in [a,b] \land f(x) < \mu) \to \mathrm{IsSup}(\xi,W) \to f(\xi) = \mu \to \neg \xi \neq b \to \mu < \mu$
\item $\varphi_{218} :\equiv \forall x\,(x \in W \leftrightarrow x \in [a,b] \land f(x) < \mu) \to \mathrm{IsSup}(\xi,W) \to f(\xi) = \mu \to \neg a \neq \xi \to \mu < \mu$
\item $\varphi_{219} :\equiv \forall x\,\forall y\,(x \in \mathbb{R} \to y \in \mathbb{R} \to 0 \leq y \to x \leq x + y) \to \forall y\,(\xi \in \mathbb{R} \to y \in \mathbb{R} \to 0 \leq y \to \xi \leq \xi + y)$
\item $\varphi_{220} :\equiv \forall x\,(x \in W \leftrightarrow x \in [a,b] \land f(x) < \mu) \to \mathrm{IsSup}(\xi,W) \to \mu < f(\xi) \to \neg a \neq \xi \to \neg a \neq \xi$
\item $\varphi_{221} :\equiv \forall x\,(x \in W \leftrightarrow x \in [a,b] \land f(x) < \mu) \to \mathrm{IsSup}(\xi,W) \to \mu < f(\xi) \to \mu < f(\xi)$
\item $\varphi_{222} :\equiv \forall x\,(x \in W \leftrightarrow x \in [a,b] \land f(x) < \mu) \to \mathrm{IsSup}(\xi,W) \to f(\xi) = \mu \to \neg \xi \neq b \to \neg \mu < \mu$
\item $\varphi_{223} :\equiv \forall x\,(x \in W \leftrightarrow x \in [a,b] \land f(x) < \mu) \to \mathrm{IsSup}(\xi,W) \to f(\xi) = \mu \to \neg a \neq \xi \to \neg \mu < \mu$
\item $\varphi_{224} :\equiv f(\xi + d) - f(\xi) < \mu - f(\xi) \to f(\xi + d) < \mu - f(\xi) + f(\xi)$
\item $\varphi_{225} :\equiv \forall x\,(x \in W \leftrightarrow x \in [a,b] \land f(x) < \mu) \to \mathrm{IsSup}(\xi,W) \to f(\xi) < \mu \to f(\xi) < \mu$
\item $\varphi_{226} :\equiv \forall x\,(x \in W \leftrightarrow x \in [a,b] \land f(x) < \mu) \to \mathrm{IsSup}(\xi,W) \to f(\xi) < \mu \to \neg \xi \neq b \to \neg \xi \neq b$
\item $\varphi_{227} :\equiv \forall x\,\forall y\,(x \in \mathbb{R} \to y \in \mathbb{R} \to x - y \in \mathbb{R}) \to \forall y\,(\xi \in \mathbb{R} \to y \in \mathbb{R} \to \xi - y \in \mathbb{R})$
\item $\varphi_{228} :\equiv \forall x\,\forall y\,(x \in \mathbb{R} \to y \in \mathbb{R} \to x \cdot y \in \mathbb{R}) \to \forall y\,(\delta \in \mathbb{R} \to y \in \mathbb{R} \to \delta \cdot y \in \mathbb{R})$
\item $\varphi_{229} :\equiv \forall x\,\forall y\,(x \in \mathbb{R} \to y \in \mathbb{R} \to x + y \in \mathbb{R}) \to \forall y\,(\xi \in \mathbb{R} \to y \in \mathbb{R} \to \xi + y \in \mathbb{R})$
\item $\varphi_{230} :\equiv \forall x\,\forall y\,(x \in \mathbb{R} \to y \in \mathbb{R} \to x - y \in \mathbb{R}) \to \forall y\,(\mu \in \mathbb{R} \to y \in \mathbb{R} \to \mu - y \in \mathbb{R})$
\item $\varphi_{231} :\equiv \mathrm{ContOn}(f,[a,b]) \leftrightarrow \mathrm{RFn}(f) \land [a,b] \subseteq \mathrm{dom}(f) \land \forall c \in [a,b]\,(\mathrm{ContIn}(f,[a,b],c))$
\item $\varphi_{232} :\equiv \forall x\,(x \in W \leftrightarrow x \in [a,b] \land f(x) < \mu) \to \mathrm{IsSup}(\xi,W) \to \mu < f(\xi) \to \neg a \neq \xi \to \mu < \mu$
\item $\varphi_{233} :\equiv d \in \mathbb{R} \land 0 < d \land d \leq \delta \cdot 2^{-1} \land d \leq b - \xi \to f(\xi + d) - f(\xi) < \mu - f(\xi)$
\item $\varphi_{234} :\equiv d \in \mathbb{R} \land 0 < d \land d \leq \delta \cdot 2^{-1} \land d \leq b - \xi \to f(\xi + d) < \mu - f(\xi) + f(\xi)$
\item $\varphi_{235} :\equiv \forall x\,(x \in W \leftrightarrow x \in [a,b] \land f(x) < \mu) \to \mathrm{IsSup}(\xi,W) \to f(\xi) < \mu \to \neg \xi \neq b \to \mu < \mu$
\item $\varphi_{236} :\equiv \forall a\,\forall b\,\forall x\,(x \in [a,b] \leftrightarrow x \in \mathbb{R} \land a \leq x \land x \leq b) \to \forall b\,\forall x\,(x \in [a,b] \leftrightarrow x \in \mathbb{R} \land a \leq x \land x \leq b)$
\item $\varphi_{237} :\equiv d \in \mathbb{R} \land 0 < d \land d \leq \delta \cdot 2^{-1} \land d \leq \xi - a \to s \in W \to \neg s \leq \xi - d \to f(\xi) - \mu \in \mathbb{R}$
\item $\varphi_{238} :\equiv \forall x\,(x \in W \leftrightarrow x \in [a,b] \land f(x) < \mu) \to \mathrm{IsSup}(\xi,W) \to \mu < f(\xi) \to \neg a \neq \xi \to \neg \mu < \mu$
\item $\varphi_{239} :\equiv \forall x\,(x \in W \leftrightarrow x \in [a,b] \land f(x) < \mu) \to \mathrm{IsSup}(\xi,W) \to f(\xi) < \mu \to \neg \xi \neq b \to \neg \mu < \mu$
\item $\varphi_{240} :\equiv \forall x\,(x \in W \leftrightarrow x \in [a,b] \land f(x) < \mu) \to \exists \xi\,(\xi \in \mathbb{R} \land a < \xi \land \xi < b \land f(\xi) = \mu)$
\item $\varphi_{241} :\equiv \forall m\,\forall y\,(f(\xi) \in \mathbb{R} \to m \in \mathbb{R} \to y \in \mathbb{R} \to -(f(\xi) - m) < y - f(\xi) \to m < y)$
\item $\varphi_{242} :\equiv \forall y\,(f(\xi) \in \mathbb{R} \to \mu \in \mathbb{R} \to y \in \mathbb{R} \to -(f(\xi) - \mu) < y - f(\xi) \to \mu < y)$
\item $\varphi_{243} :\equiv d \in \mathbb{R} \land 0 < d \land d \leq \delta \cdot 2^{-1} \land d \leq \xi - a \to d \in \mathbb{R} \land 0 < d \land d \leq \delta \cdot 2^{-1} \land d \leq \xi - a$
\item $\varphi_{244} :\equiv d \in \mathbb{R} \land 0 < d \land d \leq \delta \cdot 2^{-1} \land d \leq b - \xi \to d \in \mathbb{R} \land 0 < d \land d \leq \delta \cdot 2^{-1} \land d \leq b - \xi$
\item $\varphi_{245} :\equiv \forall z\,(\lvert s - \xi \rvert < d \to d < z \to \lvert s - \xi \rvert < z) \to \lvert s - \xi \rvert < d \to d < \delta \to \lvert s - \xi \rvert < \delta$
\item $\varphi_{246} :\equiv \forall x\,(\mathrm{RFn}(f) \to x \in \mathrm{dom}(f) \to f(x) \in \mathbb{R}) \to \mathrm{RFn}(f) \to \xi + d \in \mathrm{dom}(f) \to f(\xi + d) \in \mathbb{R}$
\item $\varphi_{247} :\equiv \forall f\,\forall \mu\,\forall A\,\exists W\,\forall x\,(x \in W \leftrightarrow x \in A \land f(x) < \mu) \to \forall \mu\,\forall A\,\exists W\,\forall x\,(x \in W \leftrightarrow x \in A \land f(x) < \mu)$
\item $\varphi_{248} :\equiv \forall x\,(x \in W \leftrightarrow x \in [a,b] \land f(x) < \mu) \to \mathrm{IsSup}(\xi,W) \to \mu < f(\xi) \to 0 < f(\xi) - \mu$
\item $\varphi_{249} :\equiv -(\mu - f(\xi)) < f(\xi + d) - f(\xi) \land f(\xi + d) - f(\xi) < \mu - f(\xi)$
\item $\varphi_{250} :\equiv \forall x\,(x \in W \leftrightarrow x \in [a,b] \land f(x) < \mu) \to \mathrm{IsSup}(\xi,W) \to f(\xi) < \mu \to 0 < \mu - f(\xi)$
\item $\varphi_{251} :\equiv \forall y\,(a \in \mathbb{R} \to y \in \mathbb{R} \to a < y \to 0 < y - a) \to a \in \mathbb{R} \to \xi \in \mathbb{R} \to a < \xi \to 0 < \xi - a$
\item $\varphi_{252} :\equiv \forall \xi\,(\neg (\xi \in \mathbb{R} \land a < \xi \land \xi < b \land f(\xi) = \mu)) \to \neg (\xi \in \mathbb{R} \land a < \xi \land \xi < b \land f(\xi) = \mu)$
\item $\varphi_{253} :\equiv d \in \mathbb{R} \land 0 < d \land d \leq \delta \cdot 2^{-1} \land d \leq \xi - a \to \exists \xi\,(\xi \in \mathbb{R} \land a < \xi \land \xi < b \land f(\xi) = \mu)$
\item $\varphi_{254} :\equiv \forall x\,(x \in W \leftrightarrow x \in [a,b] \land f(x) < \mu) \to \mathrm{IsSup}(\xi,W) \to f(\xi) = \mu \to \neg \xi \neq b \to \mu < f(\xi)$
\item $\varphi_{255} :\equiv \forall x\,(x \in W \leftrightarrow x \in [a,b] \land f(x) < \mu) \to \mathrm{IsSup}(\xi,W) \to f(\xi) = \mu \to \neg a \neq \xi \to f(\xi) < \mu$
\item $\varphi_{256} :\equiv d \in \mathbb{R} \land 0 < d \land d \leq \delta \cdot 2^{-1} \land d \leq b - \xi \to \exists \xi\,(\xi \in \mathbb{R} \land a < \xi \land \xi < b \land f(\xi) = \mu)$
\item $\varphi_{257} :\equiv \forall y\,(\delta \in \mathbb{R} \to y \in \mathbb{R} \to \delta \cdot y \in \mathbb{R}) \to \delta \in \mathbb{R} \to 2^{-1} \in \mathbb{R} \to \delta \cdot 2^{-1} \in \mathbb{R}$
\item $\varphi_{258} :\equiv \forall z\,\forall w\,(z = w \to d + \xi < z \to d + \xi < w) \to \forall w\,(b - \xi + \xi = w \to d + \xi < b - \xi + \xi \to d + \xi < w)$
\item $\varphi_{259} :\equiv \xi + d \in [a,b] \to \lvert \xi + d - \xi \rvert < \delta \to \lvert f(\xi + d) - f(\xi) \rvert < \mu - f(\xi)$
\item $\varphi_{260} :\equiv d \in \mathbb{R} \land 0 < d \land d \leq \delta \cdot 2^{-1} \land d \leq b - \xi \to \lvert f(\xi + d) - f(\xi) \rvert < \mu - f(\xi)$
\item $\varphi_{261} :\equiv \mu - f(\xi) + f(\xi) = \mu \to f(\xi + d) < \mu - f(\xi) + f(\xi) \to f(\xi + d) < \mu$
\item $\varphi_{262} :\equiv \forall x\,(x \in W \leftrightarrow x \in [a,b] \land f(x) < \mu) \to \mathrm{IsSup}(\xi,W) \to f(\xi) < \mu \to \neg \xi \neq b \to f(b) < \mu$
\item $\varphi_{263} :\equiv \forall e\,(\xi \in \mathbb{R} \to e \in \mathbb{R} \to 0 < e \to \xi - e < \xi) \to \xi \in \mathbb{R} \to d \in \mathbb{R} \to 0 < d \to \xi - d < \xi$
\item $\varphi_{264} :\equiv \forall y\,\forall z\,(\lvert s - \xi \rvert < y \to y < z \to \lvert s - \xi \rvert < z) \to \forall z\,(\lvert s - \xi \rvert < d \to d < z \to \lvert s - \xi \rvert < z)$
\item $\varphi_{265} :\equiv \forall x\,\forall y\,(x \in \mathbb{R} \to y \in \mathbb{R} \to x < y \to 0 < y - x) \to \forall y\,(a \in \mathbb{R} \to y \in \mathbb{R} \to a < y \to 0 < y - a)$
\item $\varphi_{266} :\equiv \forall y\,(\xi \in \mathbb{R} \to y \in \mathbb{R} \to 0 < y \to \xi < \xi + y) \to \xi \in \mathbb{R} \to d \in \mathbb{R} \to 0 < d \to \xi < \xi + d$
\item $\varphi_{267} :\equiv \forall w\,(\mu - f(\xi) + f(\xi) = w \to f(\xi + d) < \mu - f(\xi) + f(\xi) \to f(\xi + d) < w)$
\item $\varphi_{268} :\equiv \forall y\,(\xi \in \mathbb{R} \to y \in \mathbb{R} \to \xi < y \to 0 < y - \xi) \to \xi \in \mathbb{R} \to b \in \mathbb{R} \to \xi < b \to 0 < b - \xi$
\item $\varphi_{269} :\equiv d \in \mathbb{R} \land 0 < d \land d \leq \delta \cdot 2^{-1} \land d \leq \xi - a \to s \in W \to \neg s \leq \xi - d \to f(s) - f(\xi) \in \mathbb{R}$
\item $\varphi_{270} :\equiv \forall x\,(x \in W \leftrightarrow x \in [a,b] \land f(x) < \mu) \to \mathrm{IsSup}(\xi,W) \to \mu < f(\xi) \to \neg a \neq \xi \to f(\xi) < \mu$
\item $\varphi_{271} :\equiv \forall x\,(x \in W \leftrightarrow x \in [a,b] \land f(x) < \mu) \to \mathrm{IsSup}(\xi,W) \to \mu < f(\xi) \to f(\xi) - \mu \in \mathbb{R}$
\item $\varphi_{272} :\equiv \forall x\,(x \in W \leftrightarrow x \in [a,b] \land f(x) < \mu) \to \mathrm{IsSup}(\xi,W) \to f(\xi) < \mu \to \mu - f(\xi) \in \mathbb{R}$
\item $\varphi_{273} :\equiv \forall x\,\forall e\,(x \in \mathbb{R} \to e \in \mathbb{R} \to 0 < e \to x - e < x) \to \forall e\,(\xi \in \mathbb{R} \to e \in \mathbb{R} \to 0 < e \to \xi - e < \xi)$
\item $\varphi_{274} :\equiv \forall x\,\forall y\,(x \in \mathbb{R} \to y \in \mathbb{R} \to x < y \to 0 < y - x) \to \forall y\,(\mu \in \mathbb{R} \to y \in \mathbb{R} \to \mu < y \to 0 < y - \mu)$
\item $\varphi_{275} :\equiv \forall x\,\forall y\,(x \in \mathbb{R} \to y \in \mathbb{R} \to 0 < y \to x < x + y) \to \forall y\,(\xi \in \mathbb{R} \to y \in \mathbb{R} \to 0 < y \to \xi < \xi + y)$
\item $\varphi_{276} :\equiv \forall x\,\forall y\,(x \in \mathbb{R} \to y \in \mathbb{R} \to x < y \to 0 < y - x) \to \forall y\,(\xi \in \mathbb{R} \to y \in \mathbb{R} \to \xi < y \to 0 < y - \xi)$
\item $\varphi_{277} :\equiv \forall x\,\forall y\,(x \in \mathbb{R} \to y \in \mathbb{R} \to \neg y \leq x \to x < y) \to \forall y\,(\xi - d \in \mathbb{R} \to y \in \mathbb{R} \to \neg y \leq \xi - d \to \xi - d < y)$
\item $\varphi_{278} :\equiv \forall x\,\forall y\,(x \in \mathbb{R} \to y \in \mathbb{R} \to x - y \in \mathbb{R}) \to \forall y\,(f(s) \in \mathbb{R} \to y \in \mathbb{R} \to f(s) - y \in \mathbb{R})$
\item $\varphi_{279} :\equiv \forall w\,(b - \xi + \xi = w \to d + \xi = b - \xi + \xi \to d + \xi = w) \to b - \xi + \xi = b \to d + \xi = b - \xi + \xi \to d + \xi = b$
\item $\varphi_{280} :\equiv \forall y\,(\mu \in \mathbb{R} \to y \in \mathbb{R} \to \mu - y \in \mathbb{R}) \to \mu \in \mathbb{R} \to f(\xi) \in \mathbb{R} \to \mu - f(\xi) \in \mathbb{R}$
\item $\varphi_{281} :\equiv \forall x\,(x \in W \leftrightarrow x \in [a,b] \land f(x) < \mu) \to (\xi + d \in W \leftrightarrow \xi + d \in [a,b] \land f(\xi + d) < \mu)$
\item $\varphi_{282} :\equiv \forall a\,\forall y\,(f(\xi + d) \in \mathbb{R} \to a \in \mathbb{R} \to y \in \mathbb{R} \to f(\xi + d) - a < y \to f(\xi + d) < y + a)$
\item $\varphi_{283} :\equiv \forall x\,\forall y\,(x \in \mathbb{R} \to y \in \mathbb{R} \to x - y \in \mathbb{R}) \to \forall y\,(f(\xi) \in \mathbb{R} \to y \in \mathbb{R} \to f(\xi) - y \in \mathbb{R})$
\item $\varphi_{284} :\equiv \forall x\,(x \in [a,b] \leftrightarrow x \in \mathbb{R} \land a \leq x \land x \leq b) \to (\xi + d \in [a,b] \leftrightarrow \xi + d \in \mathbb{R} \land a \leq \xi + d \land \xi + d \leq b)$
\item $\varphi_{285} :\equiv \varphi_{15} \to [a,b] \subseteq \mathrm{dom}(f) \to s \in [a,b] \to s \in \mathrm{dom}(f)$
\item $\varphi_{286} :\equiv \forall q\,(\delta \cdot 2^{-1} \in \mathbb{R} \to q \in \mathbb{R} \to 0 < \delta \cdot 2^{-1} \to 0 < q \to \exists d\,(d \in \mathbb{R} \land 0 < d \land d \leq \delta \cdot 2^{-1} \land d \leq q))$
\item $\varphi_{287} :\equiv \forall y\,(\mu \in \mathbb{R} \to y \in \mathbb{R} \to \mu - y + y = \mu) \to \mu \in \mathbb{R} \to f(\xi) \in \mathbb{R} \to \mu - f(\xi) + f(\xi) = \mu$
\item $\varphi_{288} :\equiv \varphi_{15} \to [a,b] \subseteq \mathrm{dom}(f) \to \xi \in [a,b] \to \xi \in \mathrm{dom}(f)$
\item $\varphi_{289} :\equiv d \in \mathbb{R} \land 0 < d \land d \leq \delta \cdot 2^{-1} \land d \leq \xi - a \to s \in W \to \neg s \leq \xi - d \to \lvert f(s) - f(\xi) \rvert < f(\xi) - \mu$
\item $\varphi_{290} :\equiv d \in \mathbb{R} \land 0 < d \land d \leq \delta \cdot 2^{-1} \land d \leq \xi - a \to s \in W \to \neg s \leq \xi - d \to -(f(\xi) - \mu) < f(s) - f(\xi)$
\item $\varphi_{291} :\equiv \forall y\,(\xi - d \in \mathbb{R} \to y \in \mathbb{R} \to \neg y \leq \xi - d \to \xi - d < y) \to \xi - d \in \mathbb{R} \to s \in \mathbb{R} \to \neg s \leq \xi - d \to \xi - d < s$
\item $\varphi_{292} :\equiv \varphi_{22} \to (\mathrm{UB}(\xi - d,W) \leftrightarrow \xi - d \in \mathbb{R} \land \forall x \in W\,(x \leq \xi - d))$
\item $\varphi_{293} :\equiv \forall y\,\forall z\,(\lvert \xi + d - \xi \rvert = y \to y = z \to \lvert \xi + d - \xi \rvert = z) \to \varphi_{24}$
\item $\varphi_{294} :\equiv \forall w\,(b - \xi + \xi = w \to d + \xi < b - \xi + \xi \to d + \xi < w) \to b - \xi + \xi = b \to d + \xi < b - \xi + \xi \to d + \xi < b$
\item $\varphi_{295} :\equiv \forall x\,\forall y\,(x \in \mathbb{R} \to y \in \mathbb{R} \to x - y \in \mathbb{R}) \to \forall y\,(f(\xi + d) \in \mathbb{R} \to y \in \mathbb{R} \to f(\xi + d) - y \in \mathbb{R})$
\item $\varphi_{296} :\equiv \forall y\,(f(\xi) \in \mathbb{R} \to y \in \mathbb{R} \to f(\xi) - y \in \mathbb{R}) \to f(\xi) \in \mathbb{R} \to \mu \in \mathbb{R} \to f(\xi) - \mu \in \mathbb{R}$
\item $\varphi_{297} :\equiv \mu < f(b) \to \forall x\,(x \in W \leftrightarrow x \in [a,b] \land f(x) < \mu) \to \varphi_{26}$
\item $\varphi_{298} :\equiv \forall y\,(\mu \in \mathbb{R} \to y \in \mathbb{R} \to \mu < y \to 0 < y - \mu) \to \mu \in \mathbb{R} \to f(\xi) \in \mathbb{R} \to \mu < f(\xi) \to 0 < f(\xi) - \mu$
\item $\varphi_{299} :\equiv \mu < f(b) \to \forall x\,(x \in W \leftrightarrow x \in [a,b] \land f(x) < \mu) \to \varphi_{31}$
\item $\varphi_{300} :\equiv \forall y\,(f(\xi + d) \in \mathbb{R} \to f(\xi) \in \mathbb{R} \to y \in \mathbb{R} \to f(\xi + d) - f(\xi) < y \to f(\xi + d) < y + f(\xi))$
\item $\varphi_{301} :\equiv \forall y\,(f(s) \in \mathbb{R} \to y \in \mathbb{R} \to f(s) - y \in \mathbb{R}) \to f(s) \in \mathbb{R} \to f(\xi) \in \mathbb{R} \to f(s) - f(\xi) \in \mathbb{R}$
\item $\varphi_{302} :\equiv \forall \varepsilon \in \mathbb{R}\,(0 < \varepsilon \to \exists \delta\,(\delta \in \mathbb{R} \land 0 < \delta \land \forall x \in [a,b]\,(\lvert x - \xi \rvert < \delta \to \lvert f(x) - f(\xi) \rvert < \varepsilon)))$
\item $\varphi_{303} :\equiv \mu < f(b) \to \forall x\,(x \in W \leftrightarrow x \in [a,b] \land f(x) < \mu) \to \varphi_{70}$
\item $\varphi_{304} :\equiv a \in \mathbb{R} \to b \in \mathbb{R} \to \mu \in \mathbb{R} \to a < b \to f(a) < \mu \to \mu < f(b) \to \varphi_{19}$
\item $\varphi_{305} :\equiv \forall y\,(f(s) - f(\xi) \in \mathbb{R} \to y \in \mathbb{R} \to \lvert f(s) - f(\xi) \rvert < y \to -y < f(s) - f(\xi) \land f(s) - f(\xi) < y)$
\item $\varphi_{306} :\equiv \forall y\,\forall z\,(\lvert \xi + d - \xi \rvert = y \to y = z \to \lvert \xi + d - \xi \rvert = z) \to \varphi_{148}$
\item $\varphi_{307} :\equiv \mu < f(b) \to \forall x\,(x \in W \leftrightarrow x \in [a,b] \land f(x) < \mu) \to \varphi_{100}$
\item $\varphi_{308} :\equiv a \in \mathbb{R} \to b \in \mathbb{R} \to \mu \in \mathbb{R} \to a < b \to f(a) < \mu \to \mu < f(b) \to \varphi_{23}$
\item $\varphi_{309} :\equiv \mathrm{ContOn}(f,[a,b]) \to a \in \mathbb{R} \to b \in \mathbb{R} \to \mu \in \mathbb{R} \to a < b \to f(a) < \mu \to \varphi_{18}$
\item $\varphi_{310} :\equiv \mathrm{ContOn}(f,[a,b]) \to a \in \mathbb{R} \to b \in \mathbb{R} \to \mu \in \mathbb{R} \to a < b \to f(a) < \mu \to \varphi_{20}$
\item $\varphi_{311} :\equiv a \in \mathbb{R} \to b \in \mathbb{R} \to \mu \in \mathbb{R} \to a < b \to f(a) < \mu \to \mu < f(b) \to \varphi_{29}$
\item $\varphi_{312} :\equiv \mathrm{IsSup}(\xi,W) \to \mu < f(\xi) \to \exists \delta\,(\delta \in \mathbb{R} \land 0 < \delta \land \varphi_{157})$
\item $\varphi_{313} :\equiv \mu < f(b) \to \forall x\,(x \in W \leftrightarrow x \in [a,b] \land f(x) < \mu) \to \varphi_{116}$
\item $\varphi_{314} :\equiv \mathrm{IsSup}(\xi,W) \to f(\xi) < \mu \to \exists \delta\,(\delta \in \mathbb{R} \land 0 < \delta \land \varphi_{159})$
\item $\varphi_{315} :\equiv \mu < f(b) \to \forall x\,(x \in W \leftrightarrow x \in [a,b] \land f(x) < \mu) \to \varphi_{120}$
\item $\varphi_{316} :\equiv a \in \mathbb{R} \to b \in \mathbb{R} \to \mu \in \mathbb{R} \to a < b \to f(a) < \mu \to \mu < f(b) \to \varphi_{35}$
\item $\varphi_{317} :\equiv a \in \mathbb{R} \to b \in \mathbb{R} \to \mu \in \mathbb{R} \to a < b \to f(a) < \mu \to \mu < f(b) \to \varphi_{36}$
\item $\varphi_{318} :\equiv \mathrm{ContOn}(f,[a,b]) \land a \in \mathbb{R} \land b \in \mathbb{R} \land \mu \in \mathbb{R} \land a < b \land f(a) < \mu \land \mu < f(b) \land \forall x\,(x \in W \leftrightarrow x \in [a,b] \land f(x) < \mu)$
\item $\varphi_{319} :\equiv a \in \mathbb{R} \to b \in \mathbb{R} \to \mu \in \mathbb{R} \to a < b \to f(a) < \mu \to \mu < f(b) \to \varphi_{42}$
\item $\varphi_{320} :\equiv \mathrm{ContOn}(f,[a,b]) \to a \in \mathbb{R} \to b \in \mathbb{R} \to \mu \in \mathbb{R} \to a < b \to f(a) < \mu \to \varphi_{32}$
\item $\varphi_{321} :\equiv \mathrm{IsSup}(\xi,W) \to \mu < f(\xi) \to \delta \in \mathbb{R} \land 0 < \delta \land \varphi_{157} \to 0 < \delta$
\item $\varphi_{322} :\equiv \mathrm{IsSup}(\xi,W) \to f(\xi) < \mu \to \delta \in \mathbb{R} \land 0 < \delta \land \varphi_{159} \to 0 < \delta$
\item $\varphi_{323} :\equiv a \in \mathbb{R} \to b \in \mathbb{R} \to \mu \in \mathbb{R} \to a < b \to f(a) < \mu \to \mu < f(b) \to \varphi_{50}$
\item $\varphi_{324} :\equiv \mathrm{ContOn}(f,[a,b]) \to a \in \mathbb{R} \to b \in \mathbb{R} \to \mu \in \mathbb{R} \to a < b \to \varphi_{133}$
\item $\varphi_{325} :\equiv \forall y\,(f(\xi + d) \in \mathbb{R} \to y \in \mathbb{R} \to f(\xi + d) - y \in \mathbb{R}) \to \varphi_{123}$
\item $\varphi_{326} :\equiv \mathrm{ContOn}(f,[a,b]) \land a \in \mathbb{R} \land b \in \mathbb{R} \land \mu \in \mathbb{R} \land a < b \land f(a) < \mu \land \mu < f(b) \land \neg \exists \xi\,(\xi \in \mathbb{R} \land a < \xi \land \xi < b \land f(\xi) = \mu)$
\item $\varphi_{327} :\equiv a \in \mathbb{R} \to b \in \mathbb{R} \to \mu \in \mathbb{R} \to a < b \to f(a) < \mu \to \mu < f(b) \to \varphi_{59}$
\item $\varphi_{328} :\equiv a \in \mathbb{R} \to b \in \mathbb{R} \to \mu \in \mathbb{R} \to a < b \to f(a) < \mu \to \mu < f(b) \to \varphi_{60}$
\item $\varphi_{329} :\equiv \mathrm{IsSup}(\xi,W) \to \mu < f(\xi) \to \delta \in \mathbb{R} \land 0 < \delta \land \varphi_{157} \to \delta \in \mathbb{R}$
\item $\varphi_{330} :\equiv \mathrm{IsSup}(\xi,W) \to f(\xi) < \mu \to \delta \in \mathbb{R} \land 0 < \delta \land \varphi_{159} \to \delta \in \mathbb{R}$
\item $\varphi_{331} :\equiv d \in \mathbb{R} \land 0 < d \land d \leq \delta \cdot 2^{-1} \land d \leq \xi - a \to s \in W \to \neg s \leq \xi - d \to \varphi_{131}$
\item $\varphi_{332} :\equiv a \in \mathbb{R} \to b \in \mathbb{R} \to \mu \in \mathbb{R} \to a < b \to f(a) < \mu \to \mu < f(b) \to \varphi_{69}$
\item $\varphi_{333} :\equiv \mathrm{IsSup}(\xi,W) \to \mu < f(\xi) \to \delta \in \mathbb{R} \land 0 < \delta \land \varphi_{157} \to 0 < \xi - a$
\item $\varphi_{334} :\equiv \mathrm{IsSup}(\xi,W) \to f(\xi) < \mu \to \delta \in \mathbb{R} \land 0 < \delta \land \varphi_{159} \to 0 < b - \xi$
\item $\varphi_{335} :\equiv a \in \mathbb{R} \to b \in \mathbb{R} \to \mu \in \mathbb{R} \to a < b \to f(a) < \mu \to \mu < f(b) \to \varphi_{77}$
\item $\varphi_{336} :\equiv a \in \mathbb{R} \to b \in \mathbb{R} \to \mu \in \mathbb{R} \to a < b \to f(a) < \mu \to \mu < f(b) \to \varphi_{78}$
\item $\varphi_{337} :\equiv a \in \mathbb{R} \to b \in \mathbb{R} \to \mu \in \mathbb{R} \to a < b \to f(a) < \mu \to \mu < f(b) \to \varphi_{79}$
\item $\varphi_{338} :\equiv \mathrm{ContOn}(f,[a,b]) \to a \in \mathbb{R} \to b \in \mathbb{R} \to \mu \in \mathbb{R} \to a < b \to f(a) < \mu \to \varphi_{52}$
\item $\varphi_{339} :\equiv \mathrm{ContOn}(f,[a,b]) \to a \in \mathbb{R} \to b \in \mathbb{R} \to \mu \in \mathbb{R} \to a < b \to f(a) < \mu \to \varphi_{53}$
\item $\varphi_{340} :\equiv d \in \mathbb{R} \land 0 < d \land d \leq \delta \cdot 2^{-1} \land d \leq \xi - a \to s \in W \to \neg s \leq \xi - d \to \varphi_{157}$
\item $\varphi_{341} :\equiv a \in \mathbb{R} \to b \in \mathbb{R} \to \mu \in \mathbb{R} \to a < b \to f(a) < \mu \to \mu < f(b) \to \varphi_{91}$
\item $\varphi_{342} :\equiv a \in \mathbb{R} \to b \in \mathbb{R} \to \mu \in \mathbb{R} \to a < b \to f(a) < \mu \to \mu < f(b) \to \varphi_{98}$
\item $\varphi_{343} :\equiv a \in \mathbb{R} \to b \in \mathbb{R} \to \mu \in \mathbb{R} \to a < b \to f(a) < \mu \to \mu < f(b) \to \varphi_{99}$
\item $\varphi_{344} :\equiv d \in \mathbb{R} \land 0 < d \land d \leq \delta \cdot 2^{-1} \land d \leq \xi - a \to s \in W \to \neg s \leq \xi - d \to \varphi_{168}$
\item $\varphi_{345} :\equiv \mathrm{IsSup}(\xi,W) \to \mu < f(\xi) \to \delta \in \mathbb{R} \land 0 < \delta \land \varphi_{157} \to \xi - a \in \mathbb{R}$
\item $\varphi_{346} :\equiv \mathrm{IsSup}(\xi,W) \to f(\xi) < \mu \to \delta \in \mathbb{R} \land 0 < \delta \land \varphi_{159} \to b - \xi \in \mathbb{R}$
\item $\varphi_{347} :\equiv a \in \mathbb{R} \to b \in \mathbb{R} \to \mu \in \mathbb{R} \to a < b \to f(a) < \mu \to \mu < f(b) \to \varphi_{103}$
\item $\varphi_{348} :\equiv \mathrm{IsSup}(\xi,W) \to \mu < f(\xi) \to \delta \in \mathbb{R} \land 0 < \delta \land \varphi_{157} \to 0 < \delta \cdot 2^{-1}$
\item $\varphi_{349} :\equiv \mathrm{IsSup}(\xi,W) \to f(\xi) < \mu \to \delta \in \mathbb{R} \land 0 < \delta \land \varphi_{159} \to 0 < \delta \cdot 2^{-1}$
\item $\varphi_{350} :\equiv \mathrm{ContOn}(f,[a,b]) \to a \in \mathbb{R} \to b \in \mathbb{R} \to \mu \in \mathbb{R} \to a < b \to f(a) < \mu \to \varphi_{83}$
\item $\varphi_{351} :\equiv \mathrm{ContOn}(f,[a,b]) \to a \in \mathbb{R} \to b \in \mathbb{R} \to \mu \in \mathbb{R} \to a < b \to f(a) < \mu \to \varphi_{84}$
\item $\varphi_{352} :\equiv \delta \in \mathbb{R} \land 0 < \delta \land \varphi_{157} \to \exists \xi\,(\xi \in \mathbb{R} \land a < \xi \land \xi < b \land f(\xi) = \mu)$
\item $\varphi_{353} :\equiv \delta \in \mathbb{R} \land 0 < \delta \land \varphi_{159} \to \exists \xi\,(\xi \in \mathbb{R} \land a < \xi \land \xi < b \land f(\xi) = \mu)$
\item $\varphi_{354} :\equiv \delta \in \mathbb{R} \land 0 < \delta \land \varphi_{157} \to \exists d\,(d \in \mathbb{R} \land 0 < d \land d \leq \delta \cdot 2^{-1} \land d \leq \xi - a)$
\item $\varphi_{355} :\equiv f(\xi) - \mu \in \mathbb{R} \to 0 < f(\xi) - \mu \to \exists \delta\,(\delta \in \mathbb{R} \land 0 < \delta \land \varphi_{157})$
\item $\varphi_{356} :\equiv a \in \mathbb{R} \to b \in \mathbb{R} \to \mu \in \mathbb{R} \to a < b \to f(a) < \mu \to \mu < f(b) \to \varphi_{117}$
\item $\varphi_{357} :\equiv a \in \mathbb{R} \to b \in \mathbb{R} \to \mu \in \mathbb{R} \to a < b \to f(a) < \mu \to \mu < f(b) \to \varphi_{118}$
\item $\varphi_{358} :\equiv \delta \in \mathbb{R} \land 0 < \delta \land \varphi_{159} \to \exists d\,(d \in \mathbb{R} \land 0 < d \land d \leq \delta \cdot 2^{-1} \land d \leq b - \xi)$
\item $\varphi_{359} :\equiv \mu - f(\xi) \in \mathbb{R} \to 0 < \mu - f(\xi) \to \exists \delta\,(\delta \in \mathbb{R} \land 0 < \delta \land \varphi_{159})$
\item $\varphi_{360} :\equiv \mathrm{ContOn}(f,[a,b]) \to a \in \mathbb{R} \to b \in \mathbb{R} \to \mu \in \mathbb{R} \to a < b \to f(a) < \mu \to \varphi_{90}$
\item $\varphi_{361} :\equiv \forall y\,(f(\xi + d) - f(\xi) \in \mathbb{R} \to y \in \mathbb{R} \to \lvert f(\xi + d) - f(\xi) \rvert < y \to -y < f(\xi + d) - f(\xi) \land f(\xi + d) - f(\xi) < y)$
\item $\varphi_{362} :\equiv \mathrm{IsSup}(\xi,W) \to \mu < f(\xi) \to \delta \in \mathbb{R} \land 0 < \delta \land \varphi_{157} \to \delta \cdot 2^{-1} \in \mathbb{R}$
\item $\varphi_{363} :\equiv \mathrm{IsSup}(\xi,W) \to f(\xi) < \mu \to \delta \in \mathbb{R} \land 0 < \delta \land \varphi_{159} \to \delta \cdot 2^{-1} \in \mathbb{R}$
\item $\varphi_{364} :\equiv a \in \mathbb{R} \to b \in \mathbb{R} \to \mu \in \mathbb{R} \to a < b \to f(a) < \mu \to \mu < f(b) \to \varphi_{136}$
\item $\varphi_{365} :\equiv \delta \in \mathbb{R} \land 0 < \delta \land \varphi_{159} \to d \in \mathbb{R} \land 0 < d \land d \leq \delta \cdot 2^{-1} \land d \leq b - \xi \to 0 \leq d$
\item $\varphi_{366} :\equiv a \in \mathbb{R} \to b \in \mathbb{R} \to \mu \in \mathbb{R} \to a < b \to f(a) < \mu \to \mu < f(b) \to \varphi_{137}$
\item $\varphi_{367} :\equiv \mathrm{ContOn}(f,[a,b]) \to a \in \mathbb{R} \to b \in \mathbb{R} \to \mu \in \mathbb{R} \to a < b \to f(a) < \mu \to \varphi_{113}$
\item $\varphi_{368} :\equiv \mathrm{ContOn}(f,[a,b]) \to a \in \mathbb{R} \to b \in \mathbb{R} \to \mu \in \mathbb{R} \to a < b \to f(a) < \mu \to \varphi_{121}$
\item $\varphi_{369} :\equiv a \in \mathbb{R} \to b \in \mathbb{R} \to \mu \in \mathbb{R} \to a < b \to f(a) < \mu \to \mu < f(b) \to \varphi_{158}$
\item $\varphi_{370} :\equiv a \in \mathbb{R} \to b \in \mathbb{R} \to \mu \in \mathbb{R} \to a < b \to f(a) < \mu \to \mu < f(b) \to \varphi_{161}$
\item $\varphi_{371} :\equiv \delta \in \mathbb{R} \land 0 < \delta \land \varphi_{157} \to d \in \mathbb{R} \land 0 < d \land d \leq \delta \cdot 2^{-1} \land d \leq \xi - a \to 0 < d$
\item $\varphi_{372} :\equiv a \in \mathbb{R} \to b \in \mathbb{R} \to \mu \in \mathbb{R} \to a < b \to f(a) < \mu \to \mu < f(b) \to \varphi_{170}$
\item $\varphi_{373} :\equiv \delta \in \mathbb{R} \land 0 < \delta \land \varphi_{159} \to d \in \mathbb{R} \land 0 < d \land d \leq \delta \cdot 2^{-1} \land d \leq b - \xi \to 0 < d$
\item $\varphi_{374} :\equiv a \in \mathbb{R} \to b \in \mathbb{R} \to \mu \in \mathbb{R} \to a < b \to f(a) < \mu \to \mu < f(b) \to \varphi_{172}$
\item $\varphi_{375} :\equiv \delta \in \mathbb{R} \land 0 < \delta \land \varphi_{157} \to d \in \mathbb{R} \land 0 < d \land d \leq \delta \cdot 2^{-1} \land d \leq \xi - a \to d < \delta$
\item $\varphi_{376} :\equiv a \in \mathbb{R} \to b \in \mathbb{R} \to \mu \in \mathbb{R} \to a < b \to f(a) < \mu \to \mu < f(b) \to \varphi_{173}$
\item $\varphi_{377} :\equiv \delta \in \mathbb{R} \land 0 < \delta \land \varphi_{159} \to d \in \mathbb{R} \land 0 < d \land d \leq \delta \cdot 2^{-1} \land d \leq b - \xi \to d < \delta$
\item $\varphi_{378} :\equiv a \in \mathbb{R} \to b \in \mathbb{R} \to \mu \in \mathbb{R} \to a < b \to f(a) < \mu \to \mu < f(b) \to \varphi_{175}$
\item $\varphi_{379} :\equiv a \in \mathbb{R} \to b \in \mathbb{R} \to \mu \in \mathbb{R} \to a < b \to f(a) < \mu \to \mu < f(b) \to \varphi_{176}$
\item $\varphi_{380} :\equiv a \in \mathbb{R} \to b \in \mathbb{R} \to \mu \in \mathbb{R} \to a < b \to f(a) < \mu \to \mu < f(b) \to \varphi_{178}$
\item $\varphi_{381} :\equiv a \in \mathbb{R} \to b \in \mathbb{R} \to \mu \in \mathbb{R} \to a < b \to f(a) < \mu \to \mu < f(b) \to \varphi_{190}$
\item $\varphi_{382} :\equiv a \in \mathbb{R} \to b \in \mathbb{R} \to \mu \in \mathbb{R} \to a < b \to f(a) < \mu \to \mu < f(b) \to \varphi_{191}$
\item $\varphi_{383} :\equiv a \in \mathbb{R} \to b \in \mathbb{R} \to \mu \in \mathbb{R} \to a < b \to f(a) < \mu \to \mu < f(b) \to \varphi_{192}$
\item $\varphi_{384} :\equiv \delta \in \mathbb{R} \land 0 < \delta \land \varphi_{159} \to d \in \mathbb{R} \land 0 < d \land d \leq \delta \cdot 2^{-1} \land d \leq b - \xi \to \lvert d \rvert = d$
\item $\varphi_{385} :\equiv a \in \mathbb{R} \to b \in \mathbb{R} \to \mu \in \mathbb{R} \to a < b \to f(a) < \mu \to \mu < f(b) \to \varphi_{188}$
\item $\varphi_{386} :\equiv \delta \in \mathbb{R} \land 0 < \delta \land \varphi_{159} \to d \in \mathbb{R} \land 0 < d \land d \leq \delta \cdot 2^{-1} \land d \leq b - \xi \to d \leq b - \xi$
\item $\varphi_{387} :\equiv \delta \in \mathbb{R} \land 0 < \delta \land \varphi_{159} \to d \in \mathbb{R} \land 0 < d \land d \leq \delta \cdot 2^{-1} \land d \leq b - \xi \to d + \xi \leq b$
\item $\varphi_{388} :\equiv \delta \in \mathbb{R} \land 0 < \delta \land \varphi_{159} \to d \in \mathbb{R} \land 0 < d \land d \leq \delta \cdot 2^{-1} \land d \leq b - \xi \to \xi + d \leq b$
\item $\varphi_{389} :\equiv \delta \in \mathbb{R} \land 0 < \delta \land \varphi_{159} \to d \in \mathbb{R} \land 0 < d \land d \leq \delta \cdot 2^{-1} \land d \leq b - \xi \to a \leq \xi + d$
\item $\varphi_{390} :\equiv \delta \in \mathbb{R} \land 0 < \delta \land \varphi_{157} \to d \in \mathbb{R} \land 0 < d \land d \leq \delta \cdot 2^{-1} \land d \leq \xi - a \to d \in \mathbb{R}$
\item $\varphi_{391} :\equiv \delta \in \mathbb{R} \land 0 < \delta \land \varphi_{157} \to d \in \mathbb{R} \land 0 < d \land d \leq \delta \cdot 2^{-1} \land d \leq \xi - a \to \xi \leq \xi - d$
\item $\varphi_{392} :\equiv a \in \mathbb{R} \to b \in \mathbb{R} \to \mu \in \mathbb{R} \to a < b \to f(a) < \mu \to \mu < f(b) \to \varphi_{203}$
\item $\varphi_{393} :\equiv a \in \mathbb{R} \to b \in \mathbb{R} \to \mu \in \mathbb{R} \to a < b \to f(a) < \mu \to \mu < f(b) \to \varphi_{204}$
\item $\varphi_{394} :\equiv \delta \in \mathbb{R} \land 0 < \delta \land \varphi_{159} \to d \in \mathbb{R} \land 0 < d \land d \leq \delta \cdot 2^{-1} \land d \leq b - \xi \to d \in \mathbb{R}$
\item $\varphi_{395} :\equiv \delta \in \mathbb{R} \land 0 < \delta \land \varphi_{159} \to d \in \mathbb{R} \land 0 < d \land d \leq \delta \cdot 2^{-1} \land d \leq b - \xi \to \xi \leq \xi + d$
\item $\varphi_{396} :\equiv \delta \in \mathbb{R} \land 0 < \delta \land \varphi_{159} \to d \in \mathbb{R} \land 0 < d \land d \leq \delta \cdot 2^{-1} \land d \leq b - \xi \to \xi + d \leq \xi$
\item $\varphi_{397} :\equiv a \in \mathbb{R} \to b \in \mathbb{R} \to \mu \in \mathbb{R} \to a < b \to f(a) < \mu \to \mu < f(b) \to \varphi_{210}$
\item $\varphi_{398} :\equiv a \in \mathbb{R} \to b \in \mathbb{R} \to \mu \in \mathbb{R} \to a < b \to f(a) < \mu \to \mu < f(b) \to \varphi_{211}$
\item $\varphi_{399} :\equiv a \in \mathbb{R} \to b \in \mathbb{R} \to \mu \in \mathbb{R} \to a < b \to f(a) < \mu \to \mu < f(b) \to \varphi_{212}$
\item $\varphi_{400} :\equiv a \in \mathbb{R} \to b \in \mathbb{R} \to \mu \in \mathbb{R} \to a < b \to f(a) < \mu \to \mu < f(b) \to \varphi_{217}$
\item $\varphi_{401} :\equiv a \in \mathbb{R} \to b \in \mathbb{R} \to \mu \in \mathbb{R} \to a < b \to f(a) < \mu \to \mu < f(b) \to \varphi_{218}$
\item $\varphi_{402} :\equiv \delta \in \mathbb{R} \land 0 < \delta \land \varphi_{157} \to d \in \mathbb{R} \land 0 < d \land d \leq \delta \cdot 2^{-1} \land d \leq \xi - a \to \xi - d < \xi$
\item $\varphi_{403} :\equiv a \in \mathbb{R} \to b \in \mathbb{R} \to \mu \in \mathbb{R} \to a < b \to f(a) < \mu \to \mu < f(b) \to \varphi_{220}$
\item $\varphi_{404} :\equiv a \in \mathbb{R} \to b \in \mathbb{R} \to \mu \in \mathbb{R} \to a < b \to f(a) < \mu \to \mu < f(b) \to \varphi_{221}$
\item $\varphi_{405} :\equiv a \in \mathbb{R} \to b \in \mathbb{R} \to \mu \in \mathbb{R} \to a < b \to f(a) < \mu \to \mu < f(b) \to \varphi_{222}$
\item $\varphi_{406} :\equiv a \in \mathbb{R} \to b \in \mathbb{R} \to \mu \in \mathbb{R} \to a < b \to f(a) < \mu \to \mu < f(b) \to \varphi_{223}$
\item $\varphi_{407} :\equiv \delta \in \mathbb{R} \land 0 < \delta \land \varphi_{159} \to d \in \mathbb{R} \land 0 < d \land d \leq \delta \cdot 2^{-1} \land d \leq b - \xi \to \xi < \xi + d$
\item $\varphi_{408} :\equiv \delta \in \mathbb{R} \land 0 < \delta \land \varphi_{159} \to d \in \mathbb{R} \land 0 < d \land d \leq \delta \cdot 2^{-1} \land d \leq b - \xi \to \xi + d \in W$
\item $\varphi_{409} :\equiv a \in \mathbb{R} \to b \in \mathbb{R} \to \mu \in \mathbb{R} \to a < b \to f(a) < \mu \to \mu < f(b) \to \varphi_{225}$
\item $\varphi_{410} :\equiv a \in \mathbb{R} \to b \in \mathbb{R} \to \mu \in \mathbb{R} \to a < b \to f(a) < \mu \to \mu < f(b) \to \varphi_{226}$
\item $\varphi_{411} :\equiv \delta \in \mathbb{R} \land 0 < \delta \land \varphi_{157} \to d \in \mathbb{R} \land 0 < d \land d \leq \delta \cdot 2^{-1} \land d \leq \xi - a \to d \leq \delta \cdot 2^{-1}$
\item $\varphi_{412} :\equiv \delta \in \mathbb{R} \land 0 < \delta \land \varphi_{157} \to d \in \mathbb{R} \land 0 < d \land d \leq \delta \cdot 2^{-1} \land d \leq \xi - a \to s \in W \to s \leq \xi$
\item $\varphi_{413} :\equiv \delta \in \mathbb{R} \land 0 < \delta \land \varphi_{157} \to d \in \mathbb{R} \land 0 < d \land d \leq \delta \cdot 2^{-1} \land d \leq \xi - a \to \neg \xi - d < \xi$
\item $\varphi_{414} :\equiv \delta \in \mathbb{R} \land 0 < \delta \land \varphi_{159} \to d \in \mathbb{R} \land 0 < d \land d \leq \delta \cdot 2^{-1} \land d \leq b - \xi \to d \leq \delta \cdot 2^{-1}$
\item $\varphi_{415} :\equiv \delta \in \mathbb{R} \land 0 < \delta \land \varphi_{159} \to d \in \mathbb{R} \land 0 < d \land d \leq \delta \cdot 2^{-1} \land d \leq b - \xi \to \neg \xi < \xi + d$
\item $\varphi_{416} :\equiv a \in \mathbb{R} \to b \in \mathbb{R} \to \mu \in \mathbb{R} \to a < b \to f(a) < \mu \to \mu < f(b) \to \varphi_{232}$
\item $\varphi_{417} :\equiv \delta \in \mathbb{R} \land 0 < \delta \land \varphi_{159} \to d \in \mathbb{R} \land 0 < d \land d \leq \delta \cdot 2^{-1} \land d \leq b - \xi \to d + \xi - \xi = d$
\item $\varphi_{418} :\equiv \delta \in \mathbb{R} \land 0 < \delta \land \varphi_{159} \to d \in \mathbb{R} \land 0 < d \land d \leq \delta \cdot 2^{-1} \land d \leq b - \xi \to \xi + d = d + \xi$
\item $\varphi_{419} :\equiv \delta \in \mathbb{R} \land 0 < \delta \land \varphi_{159} \to d \in \mathbb{R} \land 0 < d \land d \leq \delta \cdot 2^{-1} \land d \leq b - \xi \to b - \xi + \xi = b$
\item $\varphi_{420} :\equiv \delta \in \mathbb{R} \land 0 < \delta \land \varphi_{159} \to d \in \mathbb{R} \land 0 < d \land d \leq \delta \cdot 2^{-1} \land d \leq b - \xi \to d + \xi = \xi + d$
\item $\varphi_{421} :\equiv a \in \mathbb{R} \to b \in \mathbb{R} \to \mu \in \mathbb{R} \to a < b \to f(a) < \mu \to \mu < f(b) \to \varphi_{235}$
\item $\varphi_{422} :\equiv \forall A\,\forall c\,(\mathrm{ContIn}(f,A,c) \leftrightarrow \mathrm{RFn}(f) \land A \subseteq \mathrm{dom}(f) \land c \in A \land \forall \varepsilon \in \mathbb{R}\,(0 < \varepsilon \to \exists \delta\,(\delta \in \mathbb{R} \land 0 < \delta \land \forall x \in A\,(\lvert x - c \rvert < \delta \to \lvert f(x) - f(c) \rvert < \varepsilon))))$
\item $\varphi_{423} :\equiv \delta \in \mathbb{R} \land 0 < \delta \land \varphi_{157} \to d \in \mathbb{R} \land 0 < d \land d \leq \delta \cdot 2^{-1} \land d \leq \xi - a \to \xi - d \in \mathbb{R}$
\item $\varphi_{424} :\equiv a \in \mathbb{R} \to b \in \mathbb{R} \to \mu \in \mathbb{R} \to a < b \to f(a) < \mu \to \mu < f(b) \to \varphi_{238}$
\item $\varphi_{425} :\equiv \delta \in \mathbb{R} \land 0 < \delta \land \varphi_{159} \to d \in \mathbb{R} \land 0 < d \land d \leq \delta \cdot 2^{-1} \land d \leq b - \xi \to b - \xi \in \mathbb{R}$
\item $\varphi_{426} :\equiv \delta \in \mathbb{R} \land 0 < \delta \land \varphi_{159} \to d \in \mathbb{R} \land 0 < d \land d \leq \delta \cdot 2^{-1} \land d \leq b - \xi \to \xi + d \in \mathbb{R}$
\item $\varphi_{427} :\equiv \mu - f(\xi) \in \mathbb{R} \to \lvert f(\xi + d) - f(\xi) \rvert < \mu - f(\xi) \to \varphi_{249}$
\item $\varphi_{428} :\equiv a \in \mathbb{R} \to b \in \mathbb{R} \to \mu \in \mathbb{R} \to a < b \to f(a) < \mu \to \mu < f(b) \to \varphi_{239}$
\item $\varphi_{429} :\equiv a \in \mathbb{R} \to b \in \mathbb{R} \to \mu \in \mathbb{R} \to a < b \to f(a) < \mu \to \mu < f(b) \to \varphi_{240}$
\item $\varphi_{430} :\equiv \delta \in \mathbb{R} \land 0 < \delta \land \varphi_{157} \to d \in \mathbb{R} \land 0 < d \land d \leq \delta \cdot 2^{-1} \land d \leq \xi - a \to s \in W \to s \in W$
\item $\varphi_{431} :\equiv \delta \in \mathbb{R} \land 0 < \delta \land \varphi_{157} \to d \in \mathbb{R} \land 0 < d \land d \leq \delta \cdot 2^{-1} \land d \leq \xi - a \to \delta \cdot 2^{-1} < \delta$
\item $\varphi_{432} :\equiv a \in \mathbb{R} \to b \in \mathbb{R} \to \mu \in \mathbb{R} \to a < b \to f(a) < \mu \to \mu < f(b) \to \varphi_{248}$
\item $\varphi_{433} :\equiv \delta \in \mathbb{R} \land 0 < \delta \land \varphi_{159} \to d \in \mathbb{R} \land 0 < d \land d \leq \delta \cdot 2^{-1} \land d \leq b - \xi \to \delta \cdot 2^{-1} < \delta$
\item $\varphi_{434} :\equiv a \in \mathbb{R} \to b \in \mathbb{R} \to \mu \in \mathbb{R} \to a < b \to f(a) < \mu \to \mu < f(b) \to \varphi_{250}$
\item $\varphi_{435} :\equiv a \in \mathbb{R} \to b \in \mathbb{R} \to \mu \in \mathbb{R} \to a < b \to f(a) < \mu \to \mu < f(b) \to \varphi_{254}$
\item $\varphi_{436} :\equiv a \in \mathbb{R} \to b \in \mathbb{R} \to \mu \in \mathbb{R} \to a < b \to f(a) < \mu \to \mu < f(b) \to \varphi_{255}$
\item $\varphi_{437} :\equiv a \in \mathbb{R} \to b \in \mathbb{R} \to \mu \in \mathbb{R} \to a < b \to f(a) < \mu \to \mu < f(b) \to \varphi_{262}$
\item $\varphi_{438} :\equiv a \in \mathbb{R} \to b \in \mathbb{R} \to \mu \in \mathbb{R} \to a < b \to f(a) < \mu \to \mu < f(b) \to \varphi_{270}$
\item $\varphi_{439} :\equiv a \in \mathbb{R} \to b \in \mathbb{R} \to \mu \in \mathbb{R} \to a < b \to f(a) < \mu \to \mu < f(b) \to \varphi_{271}$
\item $\varphi_{440} :\equiv a \in \mathbb{R} \to b \in \mathbb{R} \to \mu \in \mathbb{R} \to a < b \to f(a) < \mu \to \mu < f(b) \to \varphi_{272}$
\item $\varphi_{441} :\equiv (\forall b\,(\neg \mathrm{UB}(b,W)) \to \neg \mathrm{UB}(b,W)) \to \mathrm{ContOn}(f,[a,b]) \to \varphi_{304}$
\item $\varphi_{442} :\equiv a \in \mathbb{R} \to b \in \mathbb{R} \to \mu \in \mathbb{R} \to a < b \to f(a) < \mu \to \mu < f(b) \to \varphi_{281}$
\item $\varphi_{443} :\equiv f(s) - f(\xi) \in \mathbb{R} \to f(\xi) - \mu \in \mathbb{R} \to \lvert f(s) - f(\xi) \rvert < f(\xi) - \mu \to \varphi_{168}$
\item $\varphi_{444} :\equiv \forall x\,(x \in W \leftrightarrow x \in [a,b] \land f(x) < \mu) \to \mathrm{IsSup}(\xi,W) \to \varphi_{302}$
\item $\varphi_{445} :\equiv \forall c\,(\mathrm{ContIn}(f,[a,b],c) \leftrightarrow \mathrm{RFn}(f) \land [a,b] \subseteq \mathrm{dom}(f) \land c \in [a,b] \land \forall \varepsilon \in \mathbb{R}\,(0 < \varepsilon \to \exists \delta\,(\delta \in \mathbb{R} \land 0 < \delta \land \forall x \in [a,b]\,(\lvert x - c \rvert < \delta \to \lvert f(x) - f(c) \rvert < \varepsilon))))$
\item $\varphi_{446} :\equiv \mathrm{ContIn}(f,[a,b],\xi) \leftrightarrow \mathrm{RFn}(f) \land [a,b] \subseteq \mathrm{dom}(f) \land \xi \in [a,b] \land \varphi_{302}$
\item $\varphi_{447} :\equiv \forall W\,(\varphi_{326}) \to \forall W\,(\neg \forall x\,(x \in W \leftrightarrow x \in [a,b] \land f(x) < \mu))$
\item $\varphi_{448} :\equiv \forall x\,\forall y\,(x \in \mathbb{R} \to y \in \mathbb{R} \to \lvert x \rvert < y \to -y < x \land x < y) \to \varphi_{305}$
\item $\varphi_{449} :\equiv \mu < f(b) \to \forall x\,(x \in W \leftrightarrow x \in [a,b] \land f(x) < \mu) \to \varphi_{312}$
\item $\varphi_{450} :\equiv \mu < f(b) \to \forall x\,(x \in W \leftrightarrow x \in [a,b] \land f(x) < \mu) \to \varphi_{314}$
\item $\varphi_{451} :\equiv \mathrm{ContOn}(f,[a,b]) \to a \in \mathbb{R} \to b \in \mathbb{R} \to \mu \in \mathbb{R} \to a < b \to f(a) < \mu \to \varphi_{297}$
\item $\varphi_{452} :\equiv (u \in [a,b] \leftrightarrow u \in \mathbb{R} \land a \leq u \land u \leq b) \to \mathrm{ContOn}(f,[a,b]) \to \varphi_{319}$
\item $\varphi_{453} :\equiv \mathrm{ContOn}(f,[a,b]) \to a \in \mathbb{R} \to b \in \mathbb{R} \to \mu \in \mathbb{R} \to a < b \to f(a) < \mu \to \varphi_{299}$
\item $\varphi_{454} :\equiv (\forall x \in W\,(x \leq \xi) \to s \in W \to s \leq \xi) \to \mathrm{ContOn}(f,[a,b]) \to \varphi_{335}$
\item $\varphi_{455} :\equiv (\forall w\,(\mathrm{UB}(w,W) \to \xi \leq w) \to \mathrm{UB}(b,W) \to \xi \leq b) \to \mathrm{ContOn}(f,[a,b]) \to \varphi_{336}$
\item $\varphi_{456} :\equiv (\forall x \in W\,(x \leq \xi) \to a \in W \to a \leq \xi) \to \mathrm{ContOn}(f,[a,b]) \to \varphi_{337}$
\item $\varphi_{457} :\equiv \mu < f(b) \to \forall x\,(x \in W \leftrightarrow x \in [a,b] \land f(x) < \mu) \to \varphi_{321}$
\item $\varphi_{458} :\equiv \mu < f(b) \to \forall x\,(x \in W \leftrightarrow x \in [a,b] \land f(x) < \mu) \to \varphi_{322}$
\item $\varphi_{459} :\equiv \varphi_{318} \land \neg \exists \xi\,(\xi \in \mathbb{R} \land a < \xi \land \xi < b \land f(\xi) = \mu) \to \neg \mathrm{IsSup}(\xi,W)$
\item $\varphi_{460} :\equiv \mu < f(b) \to \forall x\,(x \in W \leftrightarrow x \in [a,b] \land f(x) < \mu) \to \varphi_{329}$
\item $\varphi_{461} :\equiv \mu < f(b) \to \forall x\,(x \in W \leftrightarrow x \in [a,b] \land f(x) < \mu) \to \varphi_{330}$
\item $\varphi_{462} :\equiv (\mathrm{UB}(\xi,W) \leftrightarrow \xi \in \mathbb{R} \land \forall x \in W\,(x \leq \xi)) \to \mathrm{ContOn}(f,[a,b]) \to \varphi_{328}$
\item $\varphi_{463} :\equiv \mu < f(b) \to \forall x\,(x \in W \leftrightarrow x \in [a,b] \land f(x) < \mu) \to \varphi_{333}$
\item $\varphi_{464} :\equiv \mu < f(b) \to \forall x\,(x \in W \leftrightarrow x \in [a,b] \land f(x) < \mu) \to \varphi_{334}$
\item $\varphi_{465} :\equiv \forall \xi\,(\varphi_{318} \land \neg \exists \xi\,(\xi \in \mathbb{R} \land a < \xi \land \xi < b \land f(\xi) = \mu)) \to \forall \xi\,(\neg \mathrm{IsSup}(\xi,W))$
\item $\varphi_{466} :\equiv \mu < f(b) \to \forall x\,(x \in W \leftrightarrow x \in [a,b] \land f(x) < \mu) \to \varphi_{345}$
\item $\varphi_{467} :\equiv \forall x\,(x \in W \leftrightarrow x \in [a,b] \land f(x) < \mu) \to \mathrm{IsSup}(\xi,W) \to \varphi_{355}$
\item $\varphi_{468} :\equiv \mu < f(b) \to \forall x\,(x \in W \leftrightarrow x \in [a,b] \land f(x) < \mu) \to \varphi_{346}$
\item $\varphi_{469} :\equiv \forall x\,(x \in W \leftrightarrow x \in [a,b] \land f(x) < \mu) \to \mathrm{IsSup}(\xi,W) \to \varphi_{359}$
\item $\varphi_{470} :\equiv \mathrm{ContOn}(f,[a,b]) \to a \in \mathbb{R} \to b \in \mathbb{R} \to \mu \in \mathbb{R} \to a < b \to f(a) < \mu \to \varphi_{303}$
\item $\varphi_{471} :\equiv \mu < f(b) \to \forall x\,(x \in W \leftrightarrow x \in [a,b] \land f(x) < \mu) \to \varphi_{348}$
\item $\varphi_{472} :\equiv \mu < f(b) \to \forall x\,(x \in W \leftrightarrow x \in [a,b] \land f(x) < \mu) \to \varphi_{349}$
\item $\varphi_{473} :\equiv (\mathrm{UB}(\xi,W) \leftrightarrow \xi \in \mathbb{R} \land \forall x \in W\,(x \leq \xi)) \to \mathrm{ContOn}(f,[a,b]) \to \varphi_{342}$
\item $\varphi_{474} :\equiv (\mathrm{ContOn}(f,[a,b]) \to a \in \mathbb{R} \to \varphi_{61}) \to \mathrm{ContOn}(f,[a,b]) \to a \in \mathbb{R} \to \varphi_{135}$
\item $\varphi_{475} :\equiv \mathrm{ContOn}(f,[a,b]) \to a \in \mathbb{R} \to b \in \mathbb{R} \to \mu \in \mathbb{R} \to a < b \to f(a) < \mu \to \varphi_{307}$
\item $\varphi_{476} :\equiv \mu < f(b) \to \forall x\,(x \in W \leftrightarrow x \in [a,b] \land f(x) < \mu) \to \varphi_{362}$
\item $\varphi_{477} :\equiv \mu < f(b) \to \forall x\,(x \in W \leftrightarrow x \in [a,b] \land f(x) < \mu) \to \varphi_{363}$
\item $\varphi_{478} :\equiv \mathrm{ContOn}(f,[a,b]) \to a \in \mathbb{R} \to b \in \mathbb{R} \to \mu \in \mathbb{R} \to a < b \to f(a) < \mu \to \varphi_{313}$
\item $\varphi_{479} :\equiv \mathrm{ContOn}(f,[a,b]) \to a \in \mathbb{R} \to b \in \mathbb{R} \to \mu \in \mathbb{R} \to a < b \to f(a) < \mu \to \varphi_{315}$
\item $\varphi_{480} :\equiv \forall x\,\forall y\,(x \in \mathbb{R} \to y \in \mathbb{R} \to \lvert x \rvert < y \to -y < x \land x < y) \to \varphi_{361}$
\item $\varphi_{481} :\equiv \varphi_{318} \land \mathrm{IsSup}(\xi,W) \land \mu < f(\xi) \land \neg \exists \xi\,(\xi \in \mathbb{R} \land a < \xi \land \xi < b \land f(\xi) = \mu)$
\item $\varphi_{482} :\equiv \varphi_{318} \land \mathrm{IsSup}(\xi,W) \land f(\xi) < \mu \land \neg \exists \xi\,(\xi \in \mathbb{R} \land a < \xi \land \xi < b \land f(\xi) = \mu)$
\item $\varphi_{483} :\equiv (\forall x \in W\,(x \leq \xi) \to \xi + d \in W \to \xi + d \leq \xi) \to \mathrm{ContOn}(f,[a,b]) \to \varphi_{374}$
\item $\varphi_{484} :\equiv (f(\xi) = \mu \to \mu < f(\xi) \to \mu < \mu) \to \mathrm{ContOn}(f,[a,b]) \to \varphi_{400}$
\item $\varphi_{485} :\equiv (f(\xi) = \mu \to f(\xi) < \mu \to \mu < \mu) \to \mathrm{ContOn}(f,[a,b]) \to \varphi_{401}$
\item $\varphi_{486} :\equiv \forall x\,(x \in W \leftrightarrow x \in [a,b] \land f(x) < \mu) \to \mathrm{IsSup}(\xi,W) \to \mu < f(\xi) \to \varphi_{352}$
\item $\varphi_{487} :\equiv \forall x\,(x \in W \leftrightarrow x \in [a,b] \land f(x) < \mu) \to \mathrm{IsSup}(\xi,W) \to f(\xi) < \mu \to \varphi_{353}$
\item $\varphi_{488} :\equiv \forall x\,(x \in W \leftrightarrow x \in [a,b] \land f(x) < \mu) \to \mathrm{IsSup}(\xi,W) \to \mu < f(\xi) \to \varphi_{354}$
\item $\varphi_{489} :\equiv \forall x\,(x \in W \leftrightarrow x \in [a,b] \land f(x) < \mu) \to \mathrm{IsSup}(\xi,W) \to f(\xi) < \mu \to \varphi_{358}$
\item $\varphi_{490} :\equiv (b = \xi \to \mu < f(b) \to \mu < f(\xi)) \to \mathrm{ContOn}(f,[a,b]) \to \varphi_{435}$
\item $\varphi_{491} :\equiv (a = \xi \to f(a) < \mu \to f(\xi) < \mu) \to \mathrm{ContOn}(f,[a,b]) \to \varphi_{436}$
\item $\varphi_{492} :\equiv \forall x\,(x \in W \leftrightarrow x \in [a,b] \land f(x) < \mu) \to \mathrm{IsSup}(\xi,W) \to f(\xi) < \mu \to \varphi_{365}$
\item $\varphi_{493} :\equiv (\xi = b \to f(\xi) < \mu \to f(b) < \mu) \to \mathrm{ContOn}(f,[a,b]) \to \varphi_{437}$
\item $\varphi_{494} :\equiv (a = \xi \to f(a) < \mu \to f(\xi) < \mu) \to \mathrm{ContOn}(f,[a,b]) \to \varphi_{438}$
\item $\varphi_{495} :\equiv \forall x\,(x \in W \leftrightarrow x \in [a,b] \land f(x) < \mu) \to \mathrm{IsSup}(\xi,W) \to \mu < f(\xi) \to \varphi_{371}$
\item $\varphi_{496} :\equiv \forall x\,(x \in W \leftrightarrow x \in [a,b] \land f(x) < \mu) \to \mathrm{IsSup}(\xi,W) \to f(\xi) < \mu \to \varphi_{373}$
\item $\varphi_{497} :\equiv \forall x\,(x \in W \leftrightarrow x \in [a,b] \land f(x) < \mu) \to \mathrm{IsSup}(\xi,W) \to \mu < f(\xi) \to \varphi_{375}$
\item $\varphi_{498} :\equiv \forall x\,(x \in W \leftrightarrow x \in [a,b] \land f(x) < \mu) \to \mathrm{IsSup}(\xi,W) \to f(\xi) < \mu \to \varphi_{377}$
\item $\varphi_{499} :\equiv \delta \in \mathbb{R} \land 0 < \delta \land \varphi_{159} \to d \in \mathbb{R} \land 0 < d \land d \leq \delta \cdot 2^{-1} \land d \leq b - \xi \to \varphi_{159}$
\item $\varphi_{500} :\equiv \forall x\,(x \in W \leftrightarrow x \in [a,b] \land f(x) < \mu) \to \mathrm{IsSup}(\xi,W) \to f(\xi) < \mu \to \varphi_{384}$
\item $\varphi_{501} :\equiv \forall x\,(x \in W \leftrightarrow x \in [a,b] \land f(x) < \mu) \to \mathrm{IsSup}(\xi,W) \to f(\xi) < \mu \to \varphi_{386}$
\item $\varphi_{502} :\equiv \forall x\,(x \in W \leftrightarrow x \in [a,b] \land f(x) < \mu) \to \mathrm{IsSup}(\xi,W) \to f(\xi) < \mu \to \varphi_{387}$
\item $\varphi_{503} :\equiv \forall x\,(x \in W \leftrightarrow x \in [a,b] \land f(x) < \mu) \to \mathrm{IsSup}(\xi,W) \to f(\xi) < \mu \to \varphi_{388}$
\item $\varphi_{504} :\equiv \forall x\,(x \in W \leftrightarrow x \in [a,b] \land f(x) < \mu) \to \mathrm{IsSup}(\xi,W) \to f(\xi) < \mu \to \varphi_{389}$
\item $\varphi_{505} :\equiv \forall x\,(x \in W \leftrightarrow x \in [a,b] \land f(x) < \mu) \to \mathrm{IsSup}(\xi,W) \to \mu < f(\xi) \to \varphi_{390}$
\item $\varphi_{506} :\equiv \forall x\,(x \in W \leftrightarrow x \in [a,b] \land f(x) < \mu) \to \mathrm{IsSup}(\xi,W) \to \mu < f(\xi) \to \varphi_{391}$
\item $\varphi_{507} :\equiv \forall x\,(x \in W \leftrightarrow x \in [a,b] \land f(x) < \mu) \to \mathrm{IsSup}(\xi,W) \to f(\xi) < \mu \to \varphi_{394}$
\item $\varphi_{508} :\equiv \forall x\,(x \in W \leftrightarrow x \in [a,b] \land f(x) < \mu) \to \mathrm{IsSup}(\xi,W) \to f(\xi) < \mu \to \varphi_{395}$
\item $\varphi_{509} :\equiv \forall x\,(x \in W \leftrightarrow x \in [a,b] \land f(x) < \mu) \to \mathrm{IsSup}(\xi,W) \to f(\xi) < \mu \to \varphi_{396}$
\item $\varphi_{510} :\equiv \mathrm{IsSup}(\xi,W) \to \mu < f(\xi) \to \delta \in \mathbb{R} \land 0 < \delta \land \varphi_{157} \to \delta \in \mathbb{R} \land 0 < \delta \land \varphi_{157}$
\item $\varphi_{511} :\equiv \mathrm{IsSup}(\xi,W) \to f(\xi) < \mu \to \delta \in \mathbb{R} \land 0 < \delta \land \varphi_{159} \to \delta \in \mathbb{R} \land 0 < \delta \land \varphi_{159}$
\item $\varphi_{512} :\equiv \forall x\,(x \in W \leftrightarrow x \in [a,b] \land f(x) < \mu) \to \mathrm{IsSup}(\xi,W) \to \mu < f(\xi) \to \varphi_{402}$
\item $\varphi_{513} :\equiv \forall x\,(x \in W \leftrightarrow x \in [a,b] \land f(x) < \mu) \to \mathrm{IsSup}(\xi,W) \to \mu < f(\xi) \to \delta \in \mathbb{R} \land 0 < \delta \land \varphi_{157} \to \varphi_{30}$
\item $\varphi_{514} :\equiv \forall x\,(x \in W \leftrightarrow x \in [a,b] \land f(x) < \mu) \to \mathrm{IsSup}(\xi,W) \to f(\xi) < \mu \to \varphi_{407}$
\item $\varphi_{515} :\equiv \forall x\,(x \in W \leftrightarrow x \in [a,b] \land f(x) < \mu) \to \mathrm{IsSup}(\xi,W) \to f(\xi) < \mu \to \varphi_{408}$
\item $\varphi_{516} :\equiv \forall x\,(x \in W \leftrightarrow x \in [a,b] \land f(x) < \mu) \to \mathrm{IsSup}(\xi,W) \to \mu < f(\xi) \to \varphi_{411}$
\item $\varphi_{517} :\equiv \forall x\,(x \in W \leftrightarrow x \in [a,b] \land f(x) < \mu) \to \mathrm{IsSup}(\xi,W) \to \mu < f(\xi) \to \varphi_{412}$
\item $\varphi_{518} :\equiv \forall x\,(x \in W \leftrightarrow x \in [a,b] \land f(x) < \mu) \to \mathrm{IsSup}(\xi,W) \to \mu < f(\xi) \to \varphi_{413}$
\item $\varphi_{519} :\equiv \forall x\,(x \in W \leftrightarrow x \in [a,b] \land f(x) < \mu) \to \mathrm{IsSup}(\xi,W) \to f(\xi) < \mu \to \varphi_{414}$
\item $\varphi_{520} :\equiv \forall x\,(x \in W \leftrightarrow x \in [a,b] \land f(x) < \mu) \to \mathrm{IsSup}(\xi,W) \to f(\xi) < \mu \to \varphi_{415}$
\item $\varphi_{521} :\equiv \forall x\,(x \in W \leftrightarrow x \in [a,b] \land f(x) < \mu) \to \mathrm{IsSup}(\xi,W) \to f(\xi) < \mu \to \varphi_{417}$
\item $\varphi_{522} :\equiv \forall x\,(x \in W \leftrightarrow x \in [a,b] \land f(x) < \mu) \to \mathrm{IsSup}(\xi,W) \to f(\xi) < \mu \to \varphi_{418}$
\item $\varphi_{523} :\equiv \forall x\,(x \in W \leftrightarrow x \in [a,b] \land f(x) < \mu) \to \mathrm{IsSup}(\xi,W) \to f(\xi) < \mu \to \varphi_{419}$
\item $\varphi_{524} :\equiv \forall x\,(x \in W \leftrightarrow x \in [a,b] \land f(x) < \mu) \to \mathrm{IsSup}(\xi,W) \to f(\xi) < \mu \to \varphi_{420}$
\item $\varphi_{525} :\equiv \forall x\,(x \in W \leftrightarrow x \in [a,b] \land f(x) < \mu) \to \mathrm{IsSup}(\xi,W) \to \mu < f(\xi) \to \varphi_{423}$
\item $\varphi_{526} :\equiv \forall x\,(x \in W \leftrightarrow x \in [a,b] \land f(x) < \mu) \to \mathrm{IsSup}(\xi,W) \to f(\xi) < \mu \to \varphi_{425}$
\item $\varphi_{527} :\equiv \forall x\,(x \in W \leftrightarrow x \in [a,b] \land f(x) < \mu) \to \mathrm{IsSup}(\xi,W) \to f(\xi) < \mu \to \varphi_{426}$
\item $\varphi_{528} :\equiv \forall x\,(x \in W \leftrightarrow x \in [a,b] \land f(x) < \mu) \to \mathrm{IsSup}(\xi,W) \to \mu < f(\xi) \to \varphi_{430}$
\item $\varphi_{529} :\equiv \forall x\,(x \in W \leftrightarrow x \in [a,b] \land f(x) < \mu) \to \mathrm{IsSup}(\xi,W) \to \mu < f(\xi) \to \varphi_{431}$
\item $\varphi_{530} :\equiv \forall x\,(x \in W \leftrightarrow x \in [a,b] \land f(x) < \mu) \to \mathrm{IsSup}(\xi,W) \to f(\xi) < \mu \to \varphi_{433}$
\item $\varphi_{531} :\equiv \forall x\,(x \in W \leftrightarrow x \in [a,b] \land f(x) < \mu) \to \mathrm{IsSup}(\xi,W) \to f(\xi) < \mu \to \delta \in \mathbb{R} \land 0 < \delta \land \varphi_{159} \to \varphi_{43}$
\item $\varphi_{532} :\equiv \forall x\,(x \in W \leftrightarrow x \in [a,b] \land f(x) < \mu) \to \mathrm{IsSup}(\xi,W) \to \mu < f(\xi) \to \delta \in \mathbb{R} \land 0 < \delta \land \varphi_{157} \to \varphi_{44}$
\item $\varphi_{533} :\equiv \forall x\,(x \in W \leftrightarrow x \in [a,b] \land f(x) < \mu) \to \mathrm{IsSup}(\xi,W) \to \mu < f(\xi) \to \delta \in \mathbb{R} \land 0 < \delta \land \varphi_{157} \to \varphi_{46}$
\item $\varphi_{534} :\equiv \forall x\,(x \in W \leftrightarrow x \in [a,b] \land f(x) < \mu) \to \mathrm{IsSup}(\xi,W) \to f(\xi) < \mu \to \delta \in \mathbb{R} \land 0 < \delta \land \varphi_{159} \to \varphi_{47}$
\item $\varphi_{535} :\equiv \forall x\,(x \in W \leftrightarrow x \in [a,b] \land f(x) < \mu) \to \mathrm{IsSup}(\xi,W) \to f(\xi) < \mu \to \delta \in \mathbb{R} \land 0 < \delta \land \varphi_{159} \to \varphi_{48}$
\item $\varphi_{536} :\equiv \delta \in \mathbb{R} \land 0 < \delta \land \varphi_{159} \to d \in \mathbb{R} \land 0 < d \land d \leq \delta \cdot 2^{-1} \land d \leq b - \xi \to \varphi_{249}$
\item $\varphi_{537} :\equiv \forall x\,(x \in W \leftrightarrow x \in [a,b] \land f(x) < \mu) \to \mathrm{IsSup}(\xi,W) \to f(\xi) < \mu \to \delta \in \mathbb{R} \land 0 < \delta \land \varphi_{159} \to \varphi_{49}$
\item $\varphi_{538} :\equiv \forall x\,(x \in W \leftrightarrow x \in [a,b] \land f(x) < \mu) \to \mathrm{IsSup}(\xi,W) \to f(\xi) < \mu \to \delta \in \mathbb{R} \land 0 < \delta \land \varphi_{159} \to \varphi_{54}$
\item $\varphi_{539} :\equiv \forall x\,(x \in W \leftrightarrow x \in [a,b] \land f(x) < \mu) \to \mathrm{IsSup}(\xi,W) \to f(\xi) < \mu \to \delta \in \mathbb{R} \land 0 < \delta \land \varphi_{159} \to \varphi_{55}$
\item $\varphi_{540} :\equiv \forall x\,(x \in W \leftrightarrow x \in [a,b] \land f(x) < \mu) \to \mathrm{IsSup}(\xi,W) \to \mu < f(\xi) \to \delta \in \mathbb{R} \land 0 < \delta \land \varphi_{157} \to \varphi_{63}$
\item $\varphi_{541} :\equiv \delta \in \mathbb{R} \land 0 < \delta \land \varphi_{159} \to d \in \mathbb{R} \land 0 < d \land d \leq \delta \cdot 2^{-1} \land d \leq b - \xi \to \varphi_{259}$
\item $\varphi_{542} :\equiv \forall x\,(x \in W \leftrightarrow x \in [a,b] \land f(x) < \mu) \to \mathrm{IsSup}(\xi,W) \to f(\xi) < \mu \to \delta \in \mathbb{R} \land 0 < \delta \land \varphi_{159} \to \varphi_{68}$
\item $\varphi_{543} :\equiv \forall x\,(x \in W \leftrightarrow x \in [a,b] \land f(x) < \mu) \to \mathrm{IsSup}(\xi,W) \to f(\xi) < \mu \to \delta \in \mathbb{R} \land 0 < \delta \land \varphi_{159} \to \varphi_{71}$
\item $\varphi_{544} :\equiv \forall x\,(x \in W \leftrightarrow x \in [a,b] \land f(x) < \mu) \to \mathrm{IsSup}(\xi,W) \to f(\xi) < \mu \to \delta \in \mathbb{R} \land 0 < \delta \land \varphi_{159} \to \varphi_{72}$
\item $\varphi_{545} :\equiv \varphi_{4} \to \varphi_{231} \to \mathrm{ContOn}(f,[a,b]) \to a \in \mathbb{R} \to b \in \mathbb{R} \to \varphi_{92}$
\item $\varphi_{546} :\equiv \forall x\,(x \in W \leftrightarrow x \in [a,b] \land f(x) < \mu) \to \mathrm{IsSup}(\xi,W) \to \mu < f(\xi) \to \delta \in \mathbb{R} \land 0 < \delta \land \varphi_{157} \to \varphi_{76}$
\item $\varphi_{547} :\equiv \forall x\,(x \in W \leftrightarrow x \in [a,b] \land f(x) < \mu) \to \mathrm{IsSup}(\xi,W) \to f(\xi) < \mu \to \delta \in \mathbb{R} \land 0 < \delta \land \varphi_{159} \to \varphi_{81}$
\item $\varphi_{548} :\equiv \forall x\,(x \in W \leftrightarrow x \in [a,b] \land f(x) < \mu) \to \mathrm{IsSup}(\xi,W) \to f(\xi) < \mu \to \delta \in \mathbb{R} \land 0 < \delta \land \varphi_{159} \to \varphi_{82}$
\item $\varphi_{549} :\equiv \forall x\,(x \in W \leftrightarrow x \in [a,b] \land f(x) < \mu) \to \mathrm{IsSup}(\xi,W) \to \mu < f(\xi) \to \delta \in \mathbb{R} \land 0 < \delta \land \varphi_{157} \to \varphi_{115}$
\item $\varphi_{550} :\equiv \forall x\,(x \in W \leftrightarrow x \in [a,b] \land f(x) < \mu) \to \mathrm{IsSup}(\xi,W) \to f(\xi) < \mu \to \delta \in \mathbb{R} \land 0 < \delta \land \varphi_{159} \to \varphi_{119}$
\item $\varphi_{551} :\equiv \forall x\,(x \in W \leftrightarrow x \in [a,b] \land f(x) < \mu) \to \mathrm{IsSup}(\xi,W) \to \mu < f(\xi) \to \delta \in \mathbb{R} \land 0 < \delta \land \varphi_{157} \to \varphi_{125}$
\item $\varphi_{552} :\equiv \forall x\,(x \in W \leftrightarrow x \in [a,b] \land f(x) < \mu) \to \mathrm{IsSup}(\xi,W) \to \mu < f(\xi) \to \delta \in \mathbb{R} \land 0 < \delta \land \varphi_{157} \to \varphi_{130}$
\item $\varphi_{553} :\equiv a \in \mathbb{R} \to b \in \mathbb{R} \to \mu \in \mathbb{R} \to a < b \to f(a) < \mu \to \mu < f(b) \to \varphi_{444}$
\item $\varphi_{554} :\equiv \forall x\,(x \in W \leftrightarrow x \in [a,b] \land f(x) < \mu) \to \mathrm{IsSup}(\xi,W) \to \mu < f(\xi) \to \delta \in \mathbb{R} \land 0 < \delta \land \varphi_{157} \to \varphi_{152}$
\item $\varphi_{555} :\equiv \forall x\,(x \in W \leftrightarrow x \in [a,b] \land f(x) < \mu) \to \mathrm{IsSup}(\xi,W) \to \mu < f(\xi) \to \delta \in \mathbb{R} \land 0 < \delta \land \varphi_{157} \to \varphi_{153}$
\item $\varphi_{556} :\equiv \forall x\,(x \in W \leftrightarrow x \in [a,b] \land f(x) < \mu) \to \mathrm{IsSup}(\xi,W) \to f(\xi) < \mu \to \delta \in \mathbb{R} \land 0 < \delta \land \varphi_{159} \to \varphi_{160}$
\item $\varphi_{557} :\equiv \forall x\,(x \in W \leftrightarrow x \in [a,b] \land f(x) < \mu) \to \mathrm{IsSup}(\xi,W) \to \mu < f(\xi) \to \delta \in \mathbb{R} \land 0 < \delta \land \varphi_{157} \to \varphi_{167}$
\item $\varphi_{558} :\equiv \forall x\,(x \in W \leftrightarrow x \in [a,b] \land f(x) < \mu) \to \mathrm{IsSup}(\xi,W) \to \mu < f(\xi) \to \delta \in \mathbb{R} \land 0 < \delta \land \varphi_{157} \to \varphi_{169}$
\item $\varphi_{559} :\equiv \forall x\,(x \in W \leftrightarrow x \in [a,b] \land f(x) < \mu) \to \mathrm{IsSup}(\xi,W) \to f(\xi) < \mu \to \delta \in \mathbb{R} \land 0 < \delta \land \varphi_{159} \to \varphi_{171}$
\item $\varphi_{560} :\equiv \forall x\,(x \in W \leftrightarrow x \in [a,b] \land f(x) < \mu) \to \mathrm{IsSup}(\xi,W) \to \mu < f(\xi) \to \delta \in \mathbb{R} \land 0 < \delta \land \varphi_{157} \to \varphi_{189}$
\item $\varphi_{561} :\equiv \forall x\,(x \in W \leftrightarrow x \in [a,b] \land f(x) < \mu) \to \mathrm{IsSup}(\xi,W) \to \mu < f(\xi) \to \delta \in \mathbb{R} \land 0 < \delta \land \varphi_{157} \to \varphi_{193}$
\item $\varphi_{562} :\equiv \forall x\,(x \in W \leftrightarrow x \in [a,b] \land f(x) < \mu) \to \mathrm{IsSup}(\xi,W) \to \mu < f(\xi) \to \delta \in \mathbb{R} \land 0 < \delta \land \varphi_{157} \to \varphi_{194}$
\item $\varphi_{563} :\equiv \forall x\,(x \in W \leftrightarrow x \in [a,b] \land f(x) < \mu) \to \mathrm{IsSup}(\xi,W) \to \mu < f(\xi) \to \delta \in \mathbb{R} \land 0 < \delta \land \varphi_{157} \to \varphi_{202}$
\item $\varphi_{564} :\equiv \forall x\,(x \in W \leftrightarrow x \in [a,b] \land f(x) < \mu) \to \mathrm{IsSup}(\xi,W) \to f(\xi) < \mu \to \delta \in \mathbb{R} \land 0 < \delta \land \varphi_{159} \to \varphi_{233}$
\item $\varphi_{565} :\equiv \forall x\,(x \in W \leftrightarrow x \in [a,b] \land f(x) < \mu) \to \mathrm{IsSup}(\xi,W) \to f(\xi) < \mu \to \delta \in \mathbb{R} \land 0 < \delta \land \varphi_{159} \to \varphi_{234}$
\item $\varphi_{566} :\equiv \forall x\,(x \in W \leftrightarrow x \in [a,b] \land f(x) < \mu) \to \mathrm{IsSup}(\xi,W) \to \mu < f(\xi) \to \delta \in \mathbb{R} \land 0 < \delta \land \varphi_{157} \to \varphi_{237}$
\item $\varphi_{567} :\equiv \forall x\,(x \in W \leftrightarrow x \in [a,b] \land f(x) < \mu) \to \mathrm{IsSup}(\xi,W) \to \mu < f(\xi) \to \delta \in \mathbb{R} \land 0 < \delta \land \varphi_{157} \to \varphi_{243}$
\item $\varphi_{568} :\equiv \forall x\,(x \in W \leftrightarrow x \in [a,b] \land f(x) < \mu) \to \mathrm{IsSup}(\xi,W) \to f(\xi) < \mu \to \delta \in \mathbb{R} \land 0 < \delta \land \varphi_{159} \to \varphi_{244}$
\item $\varphi_{569} :\equiv \forall x\,(x \in W \leftrightarrow x \in [a,b] \land f(x) < \mu) \to \mathrm{IsSup}(\xi,W) \to \mu < f(\xi) \to \delta \in \mathbb{R} \land 0 < \delta \land \varphi_{157} \to \varphi_{253}$
\item $\varphi_{570} :\equiv \forall x\,(x \in W \leftrightarrow x \in [a,b] \land f(x) < \mu) \to \mathrm{IsSup}(\xi,W) \to f(\xi) < \mu \to \delta \in \mathbb{R} \land 0 < \delta \land \varphi_{159} \to \varphi_{256}$
\item $\varphi_{571} :\equiv \forall x\,(x \in W \leftrightarrow x \in [a,b] \land f(x) < \mu) \to \mathrm{IsSup}(\xi,W) \to f(\xi) < \mu \to \delta \in \mathbb{R} \land 0 < \delta \land \varphi_{159} \to \varphi_{260}$
\item $\varphi_{572} :\equiv \forall x\,(x \in W \leftrightarrow x \in [a,b] \land f(x) < \mu) \to \mathrm{IsSup}(\xi,W) \to \mu < f(\xi) \to \delta \in \mathbb{R} \land 0 < \delta \land \varphi_{157} \to \varphi_{269}$
\item $\varphi_{573} :\equiv \mathrm{ContOn}(f,[a,b]) \to a \in \mathbb{R} \to b \in \mathbb{R} \to \mu \in \mathbb{R} \to a < b \to f(a) < \mu \to \varphi_{449}$
\item $\varphi_{574} :\equiv \mathrm{ContOn}(f,[a,b]) \to a \in \mathbb{R} \to b \in \mathbb{R} \to \mu \in \mathbb{R} \to a < b \to f(a) < \mu \to \varphi_{450}$
\item $\varphi_{575} :\equiv \varphi_{300} \to f(\xi + d) \in \mathbb{R} \to f(\xi) \in \mathbb{R} \to \mu - f(\xi) \in \mathbb{R} \to \varphi_{224}$
\item $\varphi_{576} :\equiv \mathrm{ContOn}(f,[a,b]) \to a \in \mathbb{R} \to b \in \mathbb{R} \to \mu \in \mathbb{R} \to a < b \to f(a) < \mu \to \varphi_{457}$
\item $\varphi_{577} :\equiv \mathrm{ContOn}(f,[a,b]) \to a \in \mathbb{R} \to b \in \mathbb{R} \to \mu \in \mathbb{R} \to a < b \to f(a) < \mu \to \varphi_{458}$
\item $\varphi_{578} :\equiv \mathrm{ContOn}(f,[a,b]) \to a \in \mathbb{R} \to b \in \mathbb{R} \to \mu \in \mathbb{R} \to a < b \to f(a) < \mu \to \varphi_{460}$
\item $\varphi_{579} :\equiv \mathrm{ContOn}(f,[a,b]) \to a \in \mathbb{R} \to b \in \mathbb{R} \to \mu \in \mathbb{R} \to a < b \to f(a) < \mu \to \varphi_{461}$
\item $\varphi_{580} :\equiv a \in \mathbb{R} \to b \in \mathbb{R} \to \mu \in \mathbb{R} \to a < b \to f(a) < \mu \to \mu < f(b) \to \varphi_{467}$
\item $\varphi_{581} :\equiv a \in \mathbb{R} \to b \in \mathbb{R} \to \mu \in \mathbb{R} \to a < b \to f(a) < \mu \to \mu < f(b) \to \varphi_{469}$
\item $\varphi_{582} :\equiv \mathrm{ContOn}(f,[a,b]) \to a \in \mathbb{R} \to b \in \mathbb{R} \to \mu \in \mathbb{R} \to a < b \to f(a) < \mu \to \varphi_{463}$
\item $\varphi_{583} :\equiv \mathrm{ContOn}(f,[a,b]) \to a \in \mathbb{R} \to b \in \mathbb{R} \to \mu \in \mathbb{R} \to a < b \to f(a) < \mu \to \varphi_{464}$
\item $\varphi_{584} :\equiv \forall x\,(x \in W \leftrightarrow x \in [a,b] \land f(x) < \mu) \to \mathrm{IsSup}(\xi,W) \to \mu < f(\xi) \to \delta \in \mathbb{R} \land 0 < \delta \land \varphi_{157} \to \varphi_{289}$
\item $\varphi_{585} :\equiv \forall x\,(x \in W \leftrightarrow x \in [a,b] \land f(x) < \mu) \to \mathrm{IsSup}(\xi,W) \to \mu < f(\xi) \to \delta \in \mathbb{R} \land 0 < \delta \land \varphi_{157} \to \varphi_{290}$
\item $\varphi_{586} :\equiv \mathrm{ContOn}(f,[a,b]) \to a \in \mathbb{R} \to b \in \mathbb{R} \to \mu \in \mathbb{R} \to a < b \to f(a) < \mu \to \varphi_{466}$
\item $\varphi_{587} :\equiv \mathrm{ContOn}(f,[a,b]) \to a \in \mathbb{R} \to b \in \mathbb{R} \to \mu \in \mathbb{R} \to a < b \to f(a) < \mu \to \varphi_{468}$
\item $\varphi_{588} :\equiv \mathrm{ContOn}(f,[a,b]) \to a \in \mathbb{R} \to b \in \mathbb{R} \to \mu \in \mathbb{R} \to a < b \to f(a) < \mu \to \varphi_{471}$
\item $\varphi_{589} :\equiv \mathrm{ContOn}(f,[a,b]) \to a \in \mathbb{R} \to b \in \mathbb{R} \to \mu \in \mathbb{R} \to a < b \to f(a) < \mu \to \varphi_{472}$
\item $\varphi_{590} :\equiv \mathrm{ContOn}(f,[a,b]) \to a \in \mathbb{R} \to b \in \mathbb{R} \to \mu \in \mathbb{R} \to a < b \to f(a) < \mu \to \varphi_{476}$
\item $\varphi_{591} :\equiv \mathrm{ContOn}(f,[a,b]) \to a \in \mathbb{R} \to b \in \mathbb{R} \to \mu \in \mathbb{R} \to a < b \to f(a) < \mu \to \varphi_{477}$
\item $\varphi_{592} :\equiv \mu < f(b) \to \forall x\,(x \in W \leftrightarrow x \in [a,b] \land f(x) < \mu) \to \varphi_{510}$
\item $\varphi_{593} :\equiv \mu < f(b) \to \forall x\,(x \in W \leftrightarrow x \in [a,b] \land f(x) < \mu) \to \varphi_{511}$
\item $\varphi_{594} :\equiv a \in \mathbb{R} \to b \in \mathbb{R} \to \mu \in \mathbb{R} \to a < b \to f(a) < \mu \to \mu < f(b) \to \varphi_{486}$
\item $\varphi_{595} :\equiv a \in \mathbb{R} \to b \in \mathbb{R} \to \mu \in \mathbb{R} \to a < b \to f(a) < \mu \to \mu < f(b) \to \varphi_{487}$
\item $\varphi_{596} :\equiv a \in \mathbb{R} \to b \in \mathbb{R} \to \mu \in \mathbb{R} \to a < b \to f(a) < \mu \to \mu < f(b) \to \varphi_{488}$
\item $\varphi_{597} :\equiv a \in \mathbb{R} \to b \in \mathbb{R} \to \mu \in \mathbb{R} \to a < b \to f(a) < \mu \to \mu < f(b) \to \varphi_{489}$
\item $\varphi_{598} :\equiv a \in \mathbb{R} \to b \in \mathbb{R} \to \mu \in \mathbb{R} \to a < b \to f(a) < \mu \to \mu < f(b) \to \varphi_{492}$
\item $\varphi_{599} :\equiv (\mathrm{ContOn}(f,[a,b]) \to \varphi_{304}) \to \mathrm{ContOn}(f,[a,b]) \to \varphi_{308}$
\item $\varphi_{600} :\equiv \forall x\,(x \in W \leftrightarrow x \in [a,b] \land f(x) < \mu) \to \mathrm{IsSup}(\xi,W) \to f(\xi) < \mu \to \varphi_{499}$
\item $\varphi_{601} :\equiv a \in \mathbb{R} \to b \in \mathbb{R} \to \mu \in \mathbb{R} \to a < b \to f(a) < \mu \to \mu < f(b) \to \varphi_{495}$
\item $\varphi_{602} :\equiv a \in \mathbb{R} \to b \in \mathbb{R} \to \mu \in \mathbb{R} \to a < b \to f(a) < \mu \to \mu < f(b) \to \varphi_{496}$
\item $\varphi_{603} :\equiv a \in \mathbb{R} \to b \in \mathbb{R} \to \mu \in \mathbb{R} \to a < b \to f(a) < \mu \to \mu < f(b) \to \varphi_{497}$
\item $\varphi_{604} :\equiv a \in \mathbb{R} \to b \in \mathbb{R} \to \mu \in \mathbb{R} \to a < b \to f(a) < \mu \to \mu < f(b) \to \varphi_{498}$
\item $\varphi_{605} :\equiv a \in \mathbb{R} \to b \in \mathbb{R} \to \mu \in \mathbb{R} \to a < b \to f(a) < \mu \to \mu < f(b) \to \varphi_{500}$
\item $\varphi_{606} :\equiv a \in \mathbb{R} \to b \in \mathbb{R} \to \mu \in \mathbb{R} \to a < b \to f(a) < \mu \to \mu < f(b) \to \varphi_{501}$
\item $\varphi_{607} :\equiv a \in \mathbb{R} \to b \in \mathbb{R} \to \mu \in \mathbb{R} \to a < b \to f(a) < \mu \to \mu < f(b) \to \varphi_{502}$
\item $\varphi_{608} :\equiv a \in \mathbb{R} \to b \in \mathbb{R} \to \mu \in \mathbb{R} \to a < b \to f(a) < \mu \to \mu < f(b) \to \varphi_{503}$
\item $\varphi_{609} :\equiv a \in \mathbb{R} \to b \in \mathbb{R} \to \mu \in \mathbb{R} \to a < b \to f(a) < \mu \to \mu < f(b) \to \varphi_{504}$
\item $\varphi_{610} :\equiv a \in \mathbb{R} \to b \in \mathbb{R} \to \mu \in \mathbb{R} \to a < b \to f(a) < \mu \to \mu < f(b) \to \varphi_{505}$
\item $\varphi_{611} :\equiv a \in \mathbb{R} \to b \in \mathbb{R} \to \mu \in \mathbb{R} \to a < b \to f(a) < \mu \to \mu < f(b) \to \varphi_{506}$
\item $\varphi_{612} :\equiv a \in \mathbb{R} \to b \in \mathbb{R} \to \mu \in \mathbb{R} \to a < b \to f(a) < \mu \to \mu < f(b) \to \varphi_{507}$
\item $\varphi_{613} :\equiv a \in \mathbb{R} \to b \in \mathbb{R} \to \mu \in \mathbb{R} \to a < b \to f(a) < \mu \to \mu < f(b) \to \varphi_{508}$
\item $\varphi_{614} :\equiv a \in \mathbb{R} \to b \in \mathbb{R} \to \mu \in \mathbb{R} \to a < b \to f(a) < \mu \to \mu < f(b) \to \varphi_{509}$
\item $\varphi_{615} :\equiv a \in \mathbb{R} \to b \in \mathbb{R} \to \mu \in \mathbb{R} \to a < b \to f(a) < \mu \to \mu < f(b) \to \varphi_{512}$
\item $\varphi_{616} :\equiv a \in \mathbb{R} \to b \in \mathbb{R} \to \mu \in \mathbb{R} \to a < b \to f(a) < \mu \to \mu < f(b) \to \varphi_{513}$
\item $\varphi_{617} :\equiv a \in \mathbb{R} \to b \in \mathbb{R} \to \mu \in \mathbb{R} \to a < b \to f(a) < \mu \to \mu < f(b) \to \varphi_{514}$
\item $\varphi_{618} :\equiv a \in \mathbb{R} \to b \in \mathbb{R} \to \mu \in \mathbb{R} \to a < b \to f(a) < \mu \to \mu < f(b) \to \varphi_{515}$
\item $\varphi_{619} :\equiv a \in \mathbb{R} \to b \in \mathbb{R} \to \mu \in \mathbb{R} \to a < b \to f(a) < \mu \to \mu < f(b) \to \varphi_{516}$
\item $\varphi_{620} :\equiv a \in \mathbb{R} \to b \in \mathbb{R} \to \mu \in \mathbb{R} \to a < b \to f(a) < \mu \to \mu < f(b) \to \varphi_{517}$
\item $\varphi_{621} :\equiv a \in \mathbb{R} \to b \in \mathbb{R} \to \mu \in \mathbb{R} \to a < b \to f(a) < \mu \to \mu < f(b) \to \varphi_{518}$
\item $\varphi_{622} :\equiv a \in \mathbb{R} \to b \in \mathbb{R} \to \mu \in \mathbb{R} \to a < b \to f(a) < \mu \to \mu < f(b) \to \varphi_{519}$
\item $\varphi_{623} :\equiv a \in \mathbb{R} \to b \in \mathbb{R} \to \mu \in \mathbb{R} \to a < b \to f(a) < \mu \to \mu < f(b) \to \varphi_{520}$
\item $\varphi_{624} :\equiv a \in \mathbb{R} \to b \in \mathbb{R} \to \mu \in \mathbb{R} \to a < b \to f(a) < \mu \to \mu < f(b) \to \varphi_{521}$
\item $\varphi_{625} :\equiv a \in \mathbb{R} \to b \in \mathbb{R} \to \mu \in \mathbb{R} \to a < b \to f(a) < \mu \to \mu < f(b) \to \varphi_{522}$
\item $\varphi_{626} :\equiv a \in \mathbb{R} \to b \in \mathbb{R} \to \mu \in \mathbb{R} \to a < b \to f(a) < \mu \to \mu < f(b) \to \varphi_{523}$
\item $\varphi_{627} :\equiv a \in \mathbb{R} \to b \in \mathbb{R} \to \mu \in \mathbb{R} \to a < b \to f(a) < \mu \to \mu < f(b) \to \varphi_{524}$
\item $\varphi_{628} :\equiv a \in \mathbb{R} \to b \in \mathbb{R} \to \mu \in \mathbb{R} \to a < b \to f(a) < \mu \to \mu < f(b) \to \varphi_{525}$
\item $\varphi_{629} :\equiv a \in \mathbb{R} \to b \in \mathbb{R} \to \mu \in \mathbb{R} \to a < b \to f(a) < \mu \to \mu < f(b) \to \varphi_{526}$
\item $\varphi_{630} :\equiv a \in \mathbb{R} \to b \in \mathbb{R} \to \mu \in \mathbb{R} \to a < b \to f(a) < \mu \to \mu < f(b) \to \varphi_{527}$
\item $\varphi_{631} :\equiv a \in \mathbb{R} \to b \in \mathbb{R} \to \mu \in \mathbb{R} \to a < b \to f(a) < \mu \to \mu < f(b) \to \varphi_{528}$
\item $\varphi_{632} :\equiv \forall x\,(x \in W \leftrightarrow x \in [a,b] \land f(x) < \mu) \to \mathrm{IsSup}(\xi,W) \to f(\xi) < \mu \to \varphi_{536}$
\item $\varphi_{633} :\equiv a \in \mathbb{R} \to b \in \mathbb{R} \to \mu \in \mathbb{R} \to a < b \to f(a) < \mu \to \mu < f(b) \to \varphi_{529}$
\item $\varphi_{634} :\equiv a \in \mathbb{R} \to b \in \mathbb{R} \to \mu \in \mathbb{R} \to a < b \to f(a) < \mu \to \mu < f(b) \to \varphi_{530}$
\item $\varphi_{635} :\equiv a \in \mathbb{R} \to b \in \mathbb{R} \to \mu \in \mathbb{R} \to a < b \to f(a) < \mu \to \mu < f(b) \to \varphi_{531}$
\item $\varphi_{636} :\equiv \varphi_{318} \land \mathrm{IsSup}(\xi,W) \land \mu < f(\xi) \land (\delta \in \mathbb{R} \land 0 < \delta \land \varphi_{157}) \land (d \in \mathbb{R} \land 0 < d \land d \leq \delta \cdot 2^{-1} \land d \leq \xi - a)$
\item $\varphi_{637} :\equiv a \in \mathbb{R} \to b \in \mathbb{R} \to \mu \in \mathbb{R} \to a < b \to f(a) < \mu \to \mu < f(b) \to \varphi_{532}$
\item $\varphi_{638} :\equiv \forall x\,(x \in W \leftrightarrow x \in [a,b] \land f(x) < \mu) \to \mathrm{IsSup}(\xi,W) \to f(\xi) < \mu \to \varphi_{541}$
\item $\varphi_{639} :\equiv a \in \mathbb{R} \to b \in \mathbb{R} \to \mu \in \mathbb{R} \to a < b \to f(a) < \mu \to \mu < f(b) \to \varphi_{533}$
\item $\varphi_{640} :\equiv a \in \mathbb{R} \to b \in \mathbb{R} \to \mu \in \mathbb{R} \to a < b \to f(a) < \mu \to \mu < f(b) \to \varphi_{534}$
\item $\varphi_{641} :\equiv a \in \mathbb{R} \to b \in \mathbb{R} \to \mu \in \mathbb{R} \to a < b \to f(a) < \mu \to \mu < f(b) \to \varphi_{535}$
\item $\varphi_{642} :\equiv a \in \mathbb{R} \to b \in \mathbb{R} \to \mu \in \mathbb{R} \to a < b \to f(a) < \mu \to \mu < f(b) \to \varphi_{537}$
\item $\varphi_{643} :\equiv a \in \mathbb{R} \to b \in \mathbb{R} \to \mu \in \mathbb{R} \to a < b \to f(a) < \mu \to \mu < f(b) \to \varphi_{538}$
\item $\varphi_{644} :\equiv a \in \mathbb{R} \to b \in \mathbb{R} \to \mu \in \mathbb{R} \to a < b \to f(a) < \mu \to \mu < f(b) \to \varphi_{539}$
\item $\varphi_{645} :\equiv \varphi_{318} \land \mathrm{IsSup}(\xi,W) \land \mu < f(\xi) \land (\delta \in \mathbb{R} \land 0 < \delta \land \varphi_{157}) \land \neg \exists \xi\,(\xi \in \mathbb{R} \land a < \xi \land \xi < b \land f(\xi) = \mu)$
\item $\varphi_{646} :\equiv a \in \mathbb{R} \to b \in \mathbb{R} \to \mu \in \mathbb{R} \to a < b \to f(a) < \mu \to \mu < f(b) \to \varphi_{540}$
\item $\varphi_{647} :\equiv \varphi_{318} \land \mathrm{IsSup}(\xi,W) \land f(\xi) < \mu \land (\delta \in \mathbb{R} \land 0 < \delta \land \varphi_{159}) \land \neg \exists \xi\,(\xi \in \mathbb{R} \land a < \xi \land \xi < b \land f(\xi) = \mu)$
\item $\varphi_{648} :\equiv a \in \mathbb{R} \to b \in \mathbb{R} \to \mu \in \mathbb{R} \to a < b \to f(a) < \mu \to \mu < f(b) \to \varphi_{542}$
\item $\varphi_{649} :\equiv \forall x\,(x \in W \leftrightarrow x \in [a,b] \land f(x) < \mu) \to \mathrm{IsSup}(\xi,W) \to \mu < f(\xi) \to \delta \in \mathbb{R} \land 0 < \delta \land \varphi_{157} \to \varphi_{331}$
\item $\varphi_{650} :\equiv a \in \mathbb{R} \to b \in \mathbb{R} \to \mu \in \mathbb{R} \to a < b \to f(a) < \mu \to \mu < f(b) \to \varphi_{543}$
\item $\varphi_{651} :\equiv a \in \mathbb{R} \to b \in \mathbb{R} \to \mu \in \mathbb{R} \to a < b \to f(a) < \mu \to \mu < f(b) \to \varphi_{544}$
\item $\varphi_{652} :\equiv a \in \mathbb{R} \to b \in \mathbb{R} \to \mu \in \mathbb{R} \to a < b \to f(a) < \mu \to \mu < f(b) \to \varphi_{546}$
\item $\varphi_{653} :\equiv a \in \mathbb{R} \to b \in \mathbb{R} \to \mu \in \mathbb{R} \to a < b \to f(a) < \mu \to \mu < f(b) \to \varphi_{547}$
\item $\varphi_{654} :\equiv a \in \mathbb{R} \to b \in \mathbb{R} \to \mu \in \mathbb{R} \to a < b \to f(a) < \mu \to \mu < f(b) \to \varphi_{548}$
\item $\varphi_{655} :\equiv \forall x\,(x \in W \leftrightarrow x \in [a,b] \land f(x) < \mu) \to \mathrm{IsSup}(\xi,W) \to \mu < f(\xi) \to \delta \in \mathbb{R} \land 0 < \delta \land \varphi_{157} \to \varphi_{340}$
\item $\varphi_{656} :\equiv \forall x\,(x \in W \leftrightarrow x \in [a,b] \land f(x) < \mu) \to \mathrm{IsSup}(\xi,W) \to \mu < f(\xi) \to \delta \in \mathbb{R} \land 0 < \delta \land \varphi_{157} \to \varphi_{344}$
\item $\varphi_{657} :\equiv (\mathrm{ContOn}(f,[a,b]) \to \varphi_{316}) \to \mathrm{ContOn}(f,[a,b]) \to \varphi_{341}$
\item $\varphi_{658} :\equiv \forall u\,(\varphi_{318} \to u \in W \to u \in \mathbb{R}) \to \forall u\,(\varphi_{318}) \to \forall u \in W\,(u \in \mathbb{R})$
\item $\varphi_{659} :\equiv a \in \mathbb{R} \to b \in \mathbb{R} \to \mu \in \mathbb{R} \to a < b \to f(a) < \mu \to \mu < f(b) \to \varphi_{549}$
\item $\varphi_{660} :\equiv a \in \mathbb{R} \to b \in \mathbb{R} \to \mu \in \mathbb{R} \to a < b \to f(a) < \mu \to \mu < f(b) \to \varphi_{550}$
\item $\varphi_{661} :\equiv a \in \mathbb{R} \to b \in \mathbb{R} \to \mu \in \mathbb{R} \to a < b \to f(a) < \mu \to \mu < f(b) \to \varphi_{551}$
\item $\varphi_{662} :\equiv a \in \mathbb{R} \to b \in \mathbb{R} \to \mu \in \mathbb{R} \to a < b \to f(a) < \mu \to \mu < f(b) \to \varphi_{552}$
\item $\varphi_{663} :\equiv a \in \mathbb{R} \to b \in \mathbb{R} \to \mu \in \mathbb{R} \to a < b \to f(a) < \mu \to \mu < f(b) \to \varphi_{554}$
\item $\varphi_{664} :\equiv a \in \mathbb{R} \to b \in \mathbb{R} \to \mu \in \mathbb{R} \to a < b \to f(a) < \mu \to \mu < f(b) \to \varphi_{555}$
\item $\varphi_{665} :\equiv a \in \mathbb{R} \to b \in \mathbb{R} \to \mu \in \mathbb{R} \to a < b \to f(a) < \mu \to \mu < f(b) \to \varphi_{556}$
\item $\varphi_{666} :\equiv a \in \mathbb{R} \to b \in \mathbb{R} \to \mu \in \mathbb{R} \to a < b \to f(a) < \mu \to \mu < f(b) \to \varphi_{557}$
\item $\varphi_{667} :\equiv a \in \mathbb{R} \to b \in \mathbb{R} \to \mu \in \mathbb{R} \to a < b \to f(a) < \mu \to \mu < f(b) \to \varphi_{558}$
\item $\varphi_{668} :\equiv a \in \mathbb{R} \to b \in \mathbb{R} \to \mu \in \mathbb{R} \to a < b \to f(a) < \mu \to \mu < f(b) \to \varphi_{559}$
\item $\varphi_{669} :\equiv (\mathrm{ContOn}(f,[a,b]) \to \varphi_{341}) \to \mathrm{ContOn}(f,[a,b]) \to \varphi_{332}$
\item $\varphi_{670} :\equiv a \in \mathbb{R} \to b \in \mathbb{R} \to \mu \in \mathbb{R} \to a < b \to f(a) < \mu \to \mu < f(b) \to \varphi_{560}$
\item $\varphi_{671} :\equiv (\mathrm{ContOn}(f,[a,b]) \to \varphi_{332}) \to \mathrm{ContOn}(f,[a,b]) \to \varphi_{347}$
\item $\varphi_{672} :\equiv a \in \mathbb{R} \to b \in \mathbb{R} \to \mu \in \mathbb{R} \to a < b \to f(a) < \mu \to \mu < f(b) \to \varphi_{561}$
\item $\varphi_{673} :\equiv a \in \mathbb{R} \to b \in \mathbb{R} \to \mu \in \mathbb{R} \to a < b \to f(a) < \mu \to \mu < f(b) \to \varphi_{562}$
\item $\varphi_{674} :\equiv a \in \mathbb{R} \to b \in \mathbb{R} \to \mu \in \mathbb{R} \to a < b \to f(a) < \mu \to \mu < f(b) \to \varphi_{563}$
\item $\varphi_{675} :\equiv a \in \mathbb{R} \to b \in \mathbb{R} \to \mu \in \mathbb{R} \to a < b \to f(a) < \mu \to \mu < f(b) \to \varphi_{564}$
\item $\varphi_{676} :\equiv a \in \mathbb{R} \to b \in \mathbb{R} \to \mu \in \mathbb{R} \to a < b \to f(a) < \mu \to \mu < f(b) \to \varphi_{565}$
\item $\varphi_{677} :\equiv a \in \mathbb{R} \to b \in \mathbb{R} \to \mu \in \mathbb{R} \to a < b \to f(a) < \mu \to \mu < f(b) \to \varphi_{566}$
\item $\varphi_{678} :\equiv a \in \mathbb{R} \to b \in \mathbb{R} \to \mu \in \mathbb{R} \to a < b \to f(a) < \mu \to \mu < f(b) \to \varphi_{567}$
\item $\varphi_{679} :\equiv a \in \mathbb{R} \to b \in \mathbb{R} \to \mu \in \mathbb{R} \to a < b \to f(a) < \mu \to \mu < f(b) \to \varphi_{568}$
\item $\varphi_{680} :\equiv a \in \mathbb{R} \to b \in \mathbb{R} \to \mu \in \mathbb{R} \to a < b \to f(a) < \mu \to \mu < f(b) \to \varphi_{569}$
\item $\varphi_{681} :\equiv a \in \mathbb{R} \to b \in \mathbb{R} \to \mu \in \mathbb{R} \to a < b \to f(a) < \mu \to \mu < f(b) \to \varphi_{570}$
\item $\varphi_{682} :\equiv a \in \mathbb{R} \to b \in \mathbb{R} \to \mu \in \mathbb{R} \to a < b \to f(a) < \mu \to \mu < f(b) \to \varphi_{571}$
\item $\varphi_{683} :\equiv a \in \mathbb{R} \to b \in \mathbb{R} \to \mu \in \mathbb{R} \to a < b \to f(a) < \mu \to \mu < f(b) \to \varphi_{572}$
\item $\varphi_{684} :\equiv (W \subseteq \mathbb{R} \leftrightarrow \forall u \in W\,(u \in \mathbb{R})) \to (\mathrm{ContOn}(f,[a,b]) \to \varphi_{323}) \to \varphi_{351}$
\item $\varphi_{685} :\equiv (\mathrm{ContOn}(f,[a,b]) \to \varphi_{364}) \to \mathrm{ContOn}(f,[a,b]) \to \varphi_{369}$
\item $\varphi_{686} :\equiv (\mathrm{ContOn}(f,[a,b]) \to \varphi_{366}) \to \mathrm{ContOn}(f,[a,b]) \to \varphi_{370}$
\item $\varphi_{687} :\equiv (\mathrm{ContOn}(f,[a,b]) \to a \in \mathbb{R} \to \varphi_{135}) \to (\mathrm{ContOn}(f,[a,b]) \to \varphi_{102}) \to \varphi_{339}$
\item $\varphi_{688} :\equiv (\mathrm{ContOn}(f,[a,b]) \to a \in \mathbb{R} \to a \in \mathbb{R}) \to (\mathrm{ContOn}(f,[a,b]) \to a \in \mathbb{R} \to \varphi_{62}) \to \varphi_{474}$
\item $\varphi_{689} :\equiv (\mathrm{ContOn}(f,[a,b]) \to \varphi_{405}) \to \mathrm{ContOn}(f,[a,b]) \to \varphi_{356}$
\item $\varphi_{690} :\equiv (\mathrm{ContOn}(f,[a,b]) \to \varphi_{406}) \to \mathrm{ContOn}(f,[a,b]) \to \varphi_{357}$
\item $\varphi_{691} :\equiv (\mathrm{ContOn}(f,[a,b]) \to a \in \mathbb{R} \to b \in \mathbb{R} \to b \in \mathbb{R}) \to (\mathrm{ContOn}(f,[a,b]) \to \varphi_{311}) \to \varphi_{338}$
\item $\varphi_{692} :\equiv (\mathrm{ContOn}(f,[a,b]) \to \varphi_{392}) \to \mathrm{ContOn}(f,[a,b]) \to \varphi_{381}$
\item $\varphi_{693} :\equiv (\mathrm{ContOn}(f,[a,b]) \to \varphi_{393}) \to \mathrm{ContOn}(f,[a,b]) \to \varphi_{383}$
\item $\varphi_{694} :\equiv (\mathrm{ContOn}(f,[a,b]) \to \varphi_{424}) \to \mathrm{ContOn}(f,[a,b]) \to \varphi_{364}$
\item $\varphi_{695} :\equiv (\mathrm{ContOn}(f,[a,b]) \to \varphi_{428}) \to \mathrm{ContOn}(f,[a,b]) \to \varphi_{366}$
\item $\varphi_{696} :\equiv a \in \mathbb{R} \to b \in \mathbb{R} \to \mu \in \mathbb{R} \to a < b \to f(a) < \mu \to \mu < f(b) \to \varphi_{584}$
\item $\varphi_{697} :\equiv a \in \mathbb{R} \to b \in \mathbb{R} \to \mu \in \mathbb{R} \to a < b \to f(a) < \mu \to \mu < f(b) \to \varphi_{585}$
\item $\varphi_{698} :\equiv (\mathrm{ContOn}(f,[a,b]) \to \varphi_{347}) \to \mathrm{ContOn}(f,[a,b]) \to \varphi_{440}$
\item $\varphi_{699} :\equiv (\mathrm{ContOn}(f,[a,b]) \to \varphi_{403}) \to \mathrm{ContOn}(f,[a,b]) \to \varphi_{398}$
\item $\varphi_{700} :\equiv (d + \xi = \xi + d \to d + \xi = b \to \xi + d = b) \to \mathrm{ContOn}(f,[a,b]) \to \varphi_{608}$
\item $\varphi_{701} :\equiv (\mathrm{ContOn}(f,[a,b]) \to \varphi_{410}) \to \mathrm{ContOn}(f,[a,b]) \to \varphi_{399}$
\item $\varphi_{702} :\equiv \forall d\,(\varphi_{645} \to \neg (d \in \mathbb{R} \land 0 < d \land d \leq \delta \cdot 2^{-1} \land d \leq \xi - a))$
\item $\varphi_{703} :\equiv \forall d\,(\varphi_{647} \to \neg (d \in \mathbb{R} \land 0 < d \land d \leq \delta \cdot 2^{-1} \land d \leq b - \xi))$
\item $\varphi_{704} :\equiv (\mathrm{ContOn}(f,[a,b]) \to \varphi_{381}) \to (\xi = b \to b = \xi) \to \mathrm{ContOn}(f,[a,b]) \to \varphi_{382}$
\item $\varphi_{705} :\equiv (\mathrm{ContOn}(f,[a,b]) \to \varphi_{404}) \to \mathrm{ContOn}(f,[a,b]) \to \varphi_{432}$
\item $\varphi_{706} :\equiv (\mathrm{ContOn}(f,[a,b]) \to \varphi_{409}) \to \mathrm{ContOn}(f,[a,b]) \to \varphi_{434}$
\item $\varphi_{707} :\equiv \forall d\,(\varphi_{645}) \to \forall d\,(\neg (d \in \mathbb{R} \land 0 < d \land d \leq \delta \cdot 2^{-1} \land d \leq \xi - a))$
\item $\varphi_{708} :\equiv \forall d\,(\varphi_{647}) \to \forall d\,(\neg (d \in \mathbb{R} \land 0 < d \land d \leq \delta \cdot 2^{-1} \land d \leq b - \xi))$
\item $\varphi_{709} :\equiv \mathrm{ContOn}(f,[a,b]) \to a \in \mathbb{R} \to b \in \mathbb{R} \to \mu \in \mathbb{R} \to a < b \to f(a) < \mu \to \varphi_{592}$
\item $\varphi_{710} :\equiv (\mathrm{ContOn}(f,[a,b]) \to \varphi_{327}) \to \varphi_{5} \to \mathrm{ContOn}(f,[a,b]) \to \varphi_{343}$
\item $\varphi_{711} :\equiv \mathrm{ContOn}(f,[a,b]) \to a \in \mathbb{R} \to b \in \mathbb{R} \to \mu \in \mathbb{R} \to a < b \to f(a) < \mu \to \varphi_{593}$
\item $\varphi_{712} :\equiv (\mathrm{ContOn}(f,[a,b]) \to \varphi_{429}) \to \varphi_{326} \to \neg \forall x\,(x \in W \leftrightarrow x \in [a,b] \land f(x) < \mu)$
\item $\varphi_{713} :\equiv (b - \xi + \xi = b \to d + \xi = b - \xi + \xi \to d + \xi = b) \to \mathrm{ContOn}(f,[a,b]) \to \varphi_{607}$
\item $\varphi_{714} :\equiv (\mathrm{ContOn}(f,[a,b]) \to \varphi_{437}) \to \mathrm{ContOn}(f,[a,b]) \to \varphi_{421}$
\item $\varphi_{715} :\equiv (\mathrm{ContOn}(f,[a,b]) \to \varphi_{438}) \to \mathrm{ContOn}(f,[a,b]) \to \varphi_{416}$
\item $\varphi_{716} :\equiv \varphi_{318} \land \neg \exists \xi\,(\xi \in \mathbb{R} \land a < \xi \land \xi < b \land f(\xi) = \mu) \to \forall \xi\,(\varphi_{318} \land \neg \exists \xi\,(\xi \in \mathbb{R} \land a < \xi \land \xi < b \land f(\xi) = \mu))$
\item $\varphi_{717} :\equiv a \in \mathbb{R} \to b \in \mathbb{R} \to \mu \in \mathbb{R} \to a < b \to f(a) < \mu \to \mu < f(b) \to \varphi_{600}$
\item $\varphi_{718} :\equiv (\lvert \xi + d - \xi \rvert = d \to d = \lvert \xi + d - \xi \rvert) \to \mathrm{ContOn}(f,[a,b]) \to \varphi_{641}$
\item $\varphi_{719} :\equiv \forall W\,(\varphi_{326} \to \neg \forall x\,(x \in W \leftrightarrow x \in [a,b] \land f(x) < \mu)) \to \varphi_{447}$
\item $\varphi_{720} :\equiv a \in \mathbb{R} \to b \in \mathbb{R} \to \mu \in \mathbb{R} \to a < b \to f(a) < \mu \to \mu < f(b) \to \varphi_{632}$
\item $\varphi_{721} :\equiv (\mathrm{ContOn}(f,[a,b]) \to \varphi_{343}) \to \varphi_{9} \to \mathrm{ContOn}(f,[a,b]) \to \varphi_{372}$
\item $\varphi_{722} :\equiv a \in \mathbb{R} \to b \in \mathbb{R} \to \mu \in \mathbb{R} \to a < b \to f(a) < \mu \to \mu < f(b) \to \varphi_{638}$
\item $\varphi_{723} :\equiv (\lvert \xi + d - \xi \rvert = \lvert d \rvert \to \lvert d \rvert = d \to \lvert \xi + d - \xi \rvert = d) \to \mathrm{ContOn}(f,[a,b]) \to \varphi_{640}$
\item $\varphi_{724} :\equiv (d = \lvert \xi + d - \xi \rvert \to d < \delta \to \lvert \xi + d - \xi \rvert < \delta) \to \mathrm{ContOn}(f,[a,b]) \to \varphi_{650}$
\item $\varphi_{725} :\equiv a \in \mathbb{R} \to b \in \mathbb{R} \to \mu \in \mathbb{R} \to a < b \to f(a) < \mu \to \mu < f(b) \to \varphi_{649}$
\item $\varphi_{726} :\equiv a \in \mathbb{R} \to b \in \mathbb{R} \to \mu \in \mathbb{R} \to a < b \to f(a) < \mu \to \mu < f(b) \to \varphi_{655}$
\item $\varphi_{727} :\equiv a \in \mathbb{R} \to b \in \mathbb{R} \to \mu \in \mathbb{R} \to a < b \to f(a) < \mu \to \mu < f(b) \to \varphi_{656}$
\item $\varphi_{728} :\equiv (\mathrm{ContOn}(f,[a,b]) \to \varphi_{347}) \to \varphi_{56} \to \mathrm{ContOn}(f,[a,b]) \to \varphi_{439}$
\item $\varphi_{729} :\equiv (\mathrm{ContOn}(f,[a,b]) \to \varphi_{397}) \to \varphi_{175} \to \mathrm{ContOn}(f,[a,b]) \to \varphi_{378}$
\item $\varphi_{730} :\equiv (\mathrm{ContOn}(f,[a,b]) \to \varphi_{397}) \to \varphi_{176} \to \mathrm{ContOn}(f,[a,b]) \to \varphi_{379}$
\item $\varphi_{731} :\equiv (\mathrm{ContOn}(f,[a,b]) \to \varphi_{397}) \to \varphi_{178} \to \mathrm{ContOn}(f,[a,b]) \to \varphi_{380}$
\item $\varphi_{732} :\equiv (\mathrm{ContOn}(f,[a,b]) \to \varphi_{397}) \to \varphi_{188} \to \mathrm{ContOn}(f,[a,b]) \to \varphi_{385}$
\item $\varphi_{733} :\equiv (\mathrm{ContOn}(f,[a,b]) \to a \in \mathbb{R} \to b \in \mathbb{R} \to \varphi_{92}) \to \varphi_{669}$
\item $\varphi_{734} :\equiv \lvert d + \xi - \xi \rvert = \lvert d + \xi - \xi \rvert \to \varphi_{122} \to \mathrm{ContOn}(f,[a,b]) \to \varphi_{653}$
\item $\varphi_{735} :\equiv (\mathrm{ContOn}(f,[a,b]) \to \varphi_{336}) \to \varphi_{338} \to \mathrm{ContOn}(f,[a,b]) \to \varphi_{316}$
\item $\varphi_{736} :\equiv (\mathrm{ContOn}(f,[a,b]) \to \varphi_{337}) \to \varphi_{339} \to \mathrm{ContOn}(f,[a,b]) \to \varphi_{317}$
\item $\varphi_{737} :\equiv \lvert \xi + d - \xi \rvert = \lvert \xi + d - \xi \rvert \to (\xi + d = d + \xi \to \varphi_{86}) \to \mathrm{ContOn}(f,[a,b]) \to \varphi_{668}$
\item $\varphi_{738} :\equiv (\mathrm{ContOn}(f,[a,b]) \to \varphi_{397}) \to \varphi_{281} \to \mathrm{ContOn}(f,[a,b]) \to \varphi_{442}$
\item $\varphi_{739} :\equiv (\mathrm{RFn}(f) \to \xi \in \mathrm{dom}(f) \to f(\xi) \in \mathbb{R}) \to (\mathrm{ContOn}(f,[a,b]) \to a \in \mathbb{R} \to \varphi_{80}) \to \varphi_{671}$
\item $\varphi_{740} :\equiv (\mathrm{ContOn}(f,[a,b]) \to \varphi_{328}) \to (\mathrm{ContOn}(f,[a,b]) \to a \in \mathbb{R} \to a \in \mathbb{R}) \to \varphi_{586}$
\item $\varphi_{741} :\equiv (\mathrm{ContOn}(f,[a,b]) \to \varphi_{328}) \to \mathrm{ContOn}(f,[a,b]) \to \varphi_{624}$
\item $\varphi_{742} :\equiv (\mathrm{ContOn}(f,[a,b]) \to \varphi_{328}) \to \mathrm{ContOn}(f,[a,b]) \to \varphi_{626}$
\item $\varphi_{743} :\equiv (\mathrm{ContOn}(f,[a,b]) \to \varphi_{328}) \to \mathrm{ContOn}(f,[a,b]) \to \varphi_{627}$
\item $\varphi_{744} :\equiv (\mathrm{ContOn}(f,[a,b]) \to \varphi_{328}) \to \mathrm{ContOn}(f,[a,b]) \to \varphi_{629}$
\item $\varphi_{745} :\equiv (\mathrm{ContOn}(f,[a,b]) \to \varphi_{383}) \to (\mathrm{ContOn}(f,[a,b]) \to \varphi_{102}) \to \varphi_{491}$
\item $\varphi_{746} :\equiv ([a,b] \subseteq \mathrm{dom}(f) \to \xi \in [a,b] \to \xi \in \mathrm{dom}(f)) \to \varphi_{733}$
\item $\varphi_{747} :\equiv (\mathrm{ContOn}(f,[a,b]) \to \varphi_{398}) \to (\mathrm{ContOn}(f,[a,b]) \to \varphi_{102}) \to \varphi_{494}$
\item $\varphi_{748} :\equiv (\mathrm{ContOn}(f,[a,b]) \to a \in \mathbb{R} \to b \in \mathbb{R} \to b \in \mathbb{R}) \to (\mathrm{ContOn}(f,[a,b]) \to \varphi_{328}) \to \varphi_{587}$
\item $\varphi_{749} :\equiv (\mathrm{ContOn}(f,[a,b]) \to \varphi_{347}) \to \mathrm{ContOn}(f,[a,b]) \to \varphi_{651}$
\item $\varphi_{750} :\equiv (\mathrm{ContOn}(f,[a,b]) \to a \in \mathbb{R} \to b \in \mathbb{R} \to b \in \mathbb{R}) \to (\mathrm{ContOn}(f,[a,b]) \to \varphi_{370}) \to \varphi_{583}$
\item $\varphi_{751} :\equiv (\mathrm{ContOn}(f,[a,b]) \to \varphi_{382}) \to (\mathrm{ContOn}(f,[a,b]) \to a \in \mathbb{R} \to \varphi_{132}) \to \varphi_{490}$
\item $\varphi_{752} :\equiv (\mathrm{ContOn}(f,[a,b]) \to \varphi_{347}) \to \mathrm{ContOn}(f,[a,b]) \to \varphi_{660}$
\item $\varphi_{753} :\equiv (\mathrm{ContOn}(f,[a,b]) \to \varphi_{347}) \to \mathrm{ContOn}(f,[a,b]) \to \varphi_{665}$
\item $\varphi_{754} :\equiv (\mathrm{ContOn}(f,[a,b]) \to \varphi_{379}) \to \varphi_{367} \to (x \in [a,b] \leftrightarrow x \in \mathbb{R} \land a \leq x \land x \leq b) \to \varphi_{360}$
\item $\varphi_{755} :\equiv (\mu < f(b) \to f(b) < \mu \to \mu < \mu) \to (\mathrm{ContOn}(f,[a,b]) \to a \in \mathbb{R} \to \varphi_{132}) \to \varphi_{714}$
\item $\varphi_{756} :\equiv (\mathrm{ContOn}(f,[a,b]) \to \varphi_{347}) \to \mathrm{ContOn}(f,[a,b]) \to \varphi_{683}$
\item $\varphi_{757} :\equiv (\mathrm{ContOn}(f,[a,b]) \to a \in \mathbb{R} \to b \in \mathbb{R} \to b \in \mathbb{R}) \to \varphi_{742}$
\item $\varphi_{758} :\equiv (\mathrm{ContOn}(f,[a,b]) \to a \in \mathbb{R} \to b \in \mathbb{R} \to b \in \mathbb{R}) \to \varphi_{744}$
\item $\varphi_{759} :\equiv (\mathrm{ContOn}(f,[a,b]) \to \varphi_{356}) \to (\mathrm{ContOn}(f,[a,b]) \to \varphi_{376}) \to \varphi_{470}$
\item $\varphi_{760} :\equiv (\mathrm{ContOn}(f,[a,b]) \to \varphi_{376}) \to (\mathrm{ContOn}(f,[a,b]) \to \varphi_{435}) \to \varphi_{484}$
\item $\varphi_{761} :\equiv (\mathrm{ContOn}(f,[a,b]) \to \varphi_{376}) \to (\mathrm{ContOn}(f,[a,b]) \to \varphi_{436}) \to \varphi_{485}$
\item $\varphi_{762} :\equiv (\mathrm{ContOn}(f,[a,b]) \to \varphi_{399}) \to (\mathrm{ContOn}(f,[a,b]) \to \varphi_{409}) \to \varphi_{493}$
\item $\varphi_{763} :\equiv (\mathrm{ContOn}(f,[a,b]) \to \varphi_{347}) \to \varphi_{56} \to \mathrm{ContOn}(f,[a,b]) \to \varphi_{677}$
\item $\varphi_{764} :\equiv (\mu < f(\xi) \to f(\xi) < \mu \to \mu < \mu) \to (\mathrm{ContOn}(f,[a,b]) \to \varphi_{404}) \to \varphi_{715}$
\item $\varphi_{765} :\equiv (\mathrm{ContOn}(f,[a,b]) \to \varphi_{328}) \to (\mathrm{ContOn}(f,[a,b]) \to \varphi_{369}) \to \varphi_{582}$
\item $\varphi_{766} :\equiv (\mathrm{ContOn}(f,[a,b]) \to \varphi_{341}) \to \varphi_{446} \to \mathrm{ContOn}(f,[a,b]) \to \varphi_{553}$
\item $\varphi_{767} :\equiv (\mathrm{ContOn}(f,[a,b]) \to \varphi_{594}) \to \varphi_{481} \to \neg (\delta \in \mathbb{R} \land 0 < \delta \land \varphi_{157})$
\item $\varphi_{768} :\equiv (\mathrm{ContOn}(f,[a,b]) \to \varphi_{595}) \to \varphi_{482} \to \neg (\delta \in \mathbb{R} \land 0 < \delta \land \varphi_{159})$
\item $\varphi_{769} :\equiv (\delta \in \mathbb{R} \to 2^{-1} \in \mathbb{R} \to \delta \cdot 2^{-1} \in \mathbb{R}) \to \varphi_{578} \to 2^{-1} \in \mathbb{R} \to \varphi_{590}$
\item $\varphi_{770} :\equiv (\delta \in \mathbb{R} \to 2^{-1} \in \mathbb{R} \to \delta \cdot 2^{-1} \in \mathbb{R}) \to \varphi_{579} \to 2^{-1} \in \mathbb{R} \to \varphi_{591}$
\item $\varphi_{771} :\equiv (\mathrm{ContOn}(f,[a,b]) \to \varphi_{328}) \to (\mathrm{ContOn}(f,[a,b]) \to \varphi_{317}) \to \varphi_{657}$
\item $\varphi_{772} :\equiv (\mathrm{ContOn}(f,[a,b]) \to \varphi_{602}) \to \mathrm{ContOn}(f,[a,b]) \to \varphi_{598}$
\item $\varphi_{773} :\equiv (\mathrm{ContOn}(f,[a,b]) \to \varphi_{598}) \to \mathrm{ContOn}(f,[a,b]) \to \varphi_{605}$
\item $\varphi_{774} :\equiv (\mu \in \mathbb{R} \to f(\xi) \in \mathbb{R} \to \mu - f(\xi) + f(\xi) = \mu) \to \varphi_{56} \to \varphi_{752}$
\item $\varphi_{775} :\equiv (\mathrm{ContOn}(f,[a,b]) \to \varphi_{598}) \to \mathrm{ContOn}(f,[a,b]) \to \varphi_{613}$
\item $\varphi_{776} :\equiv \forall \delta\,(\varphi_{481} \to \neg (\delta \in \mathbb{R} \land 0 < \delta \land \varphi_{157})) \to \forall \delta\,(\varphi_{481}) \to \forall \delta\,(\neg (\delta \in \mathbb{R} \land 0 < \delta \land \varphi_{157}))$
\item $\varphi_{777} :\equiv \forall \delta\,(\varphi_{482} \to \neg (\delta \in \mathbb{R} \land 0 < \delta \land \varphi_{159})) \to \forall \delta\,(\varphi_{482}) \to \forall \delta\,(\neg (\delta \in \mathbb{R} \land 0 < \delta \land \varphi_{159}))$
\item $\varphi_{778} :\equiv (\mathrm{ContOn}(f,[a,b]) \to \varphi_{439}) \to (\mathrm{ContOn}(f,[a,b]) \to \varphi_{432}) \to \varphi_{573}$
\item $\varphi_{779} :\equiv (\mathrm{ContOn}(f,[a,b]) \to \varphi_{440}) \to (\mathrm{ContOn}(f,[a,b]) \to \varphi_{434}) \to \varphi_{574}$
\item $\varphi_{780} :\equiv (\mathrm{ContOn}(f,[a,b]) \to \varphi_{601}) \to \mathrm{ContOn}(f,[a,b]) \to \varphi_{615}$
\item $\varphi_{781} :\equiv (\mathrm{ContOn}(f,[a,b]) \to \varphi_{602}) \to \mathrm{ContOn}(f,[a,b]) \to \varphi_{617}$
\item $\varphi_{782} :\equiv (\mathrm{ContOn}(f,[a,b]) \to \varphi_{613}) \to \mathrm{ContOn}(f,[a,b]) \to \varphi_{609}$
\item $\varphi_{783} :\equiv (\mathrm{ContOn}(f,[a,b]) \to \varphi_{616}) \to \mathrm{ContOn}(f,[a,b]) \to \varphi_{611}$
\item $\varphi_{784} :\equiv (\mathrm{ContOn}(f,[a,b]) \to \varphi_{618}) \to \mathrm{ContOn}(f,[a,b]) \to \varphi_{614}$
\item $\varphi_{785} :\equiv (\mathrm{ContOn}(f,[a,b]) \to \varphi_{611}) \to \mathrm{ContOn}(f,[a,b]) \to \varphi_{621}$
\item $\varphi_{786} :\equiv (\mathrm{ContOn}(f,[a,b]) \to \varphi_{614}) \to \mathrm{ContOn}(f,[a,b]) \to \varphi_{623}$
\item $\varphi_{787} :\equiv (\mathrm{ContOn}(f,[a,b]) \to \varphi_{633}) \to \mathrm{ContOn}(f,[a,b]) \to \varphi_{603}$
\item $\varphi_{788} :\equiv (\mathrm{ContOn}(f,[a,b]) \to \varphi_{634}) \to \mathrm{ContOn}(f,[a,b]) \to \varphi_{604}$
\item $\varphi_{789} :\equiv (\mathrm{ContOn}(f,[a,b]) \to \varphi_{612}) \to \mathrm{ContOn}(f,[a,b]) \to \varphi_{625}$
\item $\varphi_{790} :\equiv (\mathrm{ContOn}(f,[a,b]) \to \varphi_{610}) \to \mathrm{ContOn}(f,[a,b]) \to \varphi_{628}$
\item $\varphi_{791} :\equiv (\mathrm{ContOn}(f,[a,b]) \to \varphi_{612}) \to \mathrm{ContOn}(f,[a,b]) \to \varphi_{630}$
\item $\varphi_{792} :\equiv (\mathrm{ContOn}(f,[a,b]) \to \varphi_{606}) \to \mathrm{ContOn}(f,[a,b]) \to \varphi_{643}$
\item $\varphi_{793} :\equiv (\mathrm{ContOn}(f,[a,b]) \to \varphi_{608}) \to \mathrm{ContOn}(f,[a,b]) \to \varphi_{644}$
\item $\varphi_{794} :\equiv (\mathrm{ContOn}(f,[a,b]) \to \varphi_{631}) \to \mathrm{ContOn}(f,[a,b]) \to \varphi_{620}$
\item $\varphi_{795} :\equiv (\mathrm{ContOn}(f,[a,b]) \to \varphi_{642}) \to \mathrm{ContOn}(f,[a,b]) \to \varphi_{618}$
\item $\varphi_{796} :\equiv (\mathrm{ContOn}(f,[a,b]) \to \varphi_{637}) \to \varphi_{636} \to s \in W \to s \leq \xi - d$
\item $\varphi_{797} :\equiv (\mathrm{ContOn}(f,[a,b]) \to \varphi_{644}) \to \mathrm{ContOn}(f,[a,b]) \to \varphi_{635}$
\item $\varphi_{798} :\equiv (\mathrm{ContOn}(f,[a,b]) \to \varphi_{635}) \to \mathrm{ContOn}(f,[a,b]) \to \varphi_{648}$
\item $\varphi_{799} :\equiv (\mathrm{ContOn}(f,[a,b]) \to \varphi_{631}) \to \mathrm{ContOn}(f,[a,b]) \to \varphi_{652}$
\item $\varphi_{800} :\equiv (\mathrm{ContOn}(f,[a,b]) \to \varphi_{652}) \to \mathrm{ContOn}(f,[a,b]) \to \varphi_{639}$
\item $\varphi_{801} :\equiv (\mu \in \mathbb{R} \to f(\xi) \in \mathbb{R} \to \mu < f(\xi) \to 0 < f(\xi) - \mu) \to \varphi_{56} \to (\mathrm{ContOn}(f,[a,b]) \to \varphi_{347}) \to \varphi_{705}$
\item $\varphi_{802} :\equiv (f(\xi) \in \mathbb{R} \to \mu \in \mathbb{R} \to f(\xi) < \mu \to 0 < \mu - f(\xi)) \to (\mathrm{ContOn}(f,[a,b]) \to \varphi_{347}) \to \varphi_{56} \to \varphi_{706}$
\item $\varphi_{803} :\equiv \forall s\,(\varphi_{636} \to s \in W \to s \leq \xi - d) \to \forall s\,(\varphi_{636}) \to \forall s \in W\,(s \leq \xi - d)$
\item $\varphi_{804} :\equiv (\mathrm{ContOn}(f,[a,b]) \to \varphi_{603}) \to \mathrm{ContOn}(f,[a,b]) \to \varphi_{674}$
\item $\varphi_{805} :\equiv (\mathrm{ContOn}(f,[a,b]) \to \varphi_{610}) \to \mathrm{ContOn}(f,[a,b]) \to \varphi_{670}$
\item $\varphi_{806} :\equiv (\mathrm{ContOn}(f,[a,b]) \to \varphi_{661}) \to \mathrm{ContOn}(f,[a,b]) \to \varphi_{637}$
\item $\varphi_{807} :\equiv (\mathrm{ContOn}(f,[a,b]) \to \varphi_{678}) \to \mathrm{ContOn}(f,[a,b]) \to \varphi_{601}$
\item $\varphi_{808} :\equiv (\mathrm{ContOn}(f,[a,b]) \to \varphi_{631}) \to \mathrm{ContOn}(f,[a,b]) \to \varphi_{666}$
\item $\varphi_{809} :\equiv (\mathrm{ContOn}(f,[a,b]) \to \varphi_{679}) \to \mathrm{ContOn}(f,[a,b]) \to \varphi_{602}$
\item $\varphi_{810} :\equiv (\mathrm{ContOn}(f,[a,b]) \to \varphi_{679}) \to \mathrm{ContOn}(f,[a,b]) \to \varphi_{606}$
\item $\varphi_{811} :\equiv (\mathrm{ContOn}(f,[a,b]) \to \varphi_{678}) \to \mathrm{ContOn}(f,[a,b]) \to \varphi_{610}$
\item $\varphi_{812} :\equiv (\mathrm{ContOn}(f,[a,b]) \to \varphi_{679}) \to \mathrm{ContOn}(f,[a,b]) \to \varphi_{612}$
\item $\varphi_{813} :\equiv (\mathrm{ContOn}(f,[a,b]) \to \varphi_{652}) \to \mathrm{ContOn}(f,[a,b]) \to \varphi_{664}$
\item $\varphi_{814} :\equiv (\mathrm{ContOn}(f,[a,b]) \to \varphi_{678}) \to \mathrm{ContOn}(f,[a,b]) \to \varphi_{619}$
\item $\varphi_{815} :\equiv (\mathrm{ContOn}(f,[a,b]) \to \varphi_{679}) \to \mathrm{ContOn}(f,[a,b]) \to \varphi_{622}$
\item $\varphi_{816} :\equiv (\mathrm{ContOn}(f,[a,b]) \to \varphi_{615}) \to \mathrm{ContOn}(f,[a,b]) \to \varphi_{680}$
\item $\varphi_{817} :\equiv (\mathrm{ContOn}(f,[a,b]) \to \varphi_{617}) \to \mathrm{ContOn}(f,[a,b]) \to \varphi_{681}$
\item $\varphi_{818} :\equiv (\mathrm{ContOn}(f,[a,b]) \to \varphi_{662}) \to \mathrm{ContOn}(f,[a,b]) \to \varphi_{663}$
\item $\varphi_{819} :\equiv (\mathrm{ContOn}(f,[a,b]) \to \varphi_{666}) \to \mathrm{ContOn}(f,[a,b]) \to \varphi_{659}$
\item $\varphi_{820} :\equiv (\mathrm{ContOn}(f,[a,b]) \to \varphi_{663}) \to \mathrm{ContOn}(f,[a,b]) \to \varphi_{672}$
\item $\varphi_{821} :\equiv (\mathrm{ContOn}(f,[a,b]) \to \varphi_{664}) \to \mathrm{ContOn}(f,[a,b]) \to \varphi_{673}$
\item $\varphi_{822} :\equiv (\mathrm{ContOn}(f,[a,b]) \to \varphi_{650}) \to \mathrm{ContOn}(f,[a,b]) \to \varphi_{682}$
\item $\varphi_{823} :\equiv (a \in \mathbb{R} \to \xi \in \mathbb{R} \to a < \xi \to 0 < \xi - a) \to (\mathrm{ContOn}(f,[a,b]) \to a \in \mathbb{R} \to a \in \mathbb{R}) \to \varphi_{765}$
\item $\varphi_{824} :\equiv (\mathrm{ContOn}(f,[a,b]) \to \varphi_{675}) \to \mathrm{ContOn}(f,[a,b]) \to \varphi_{676}$
\item $\varphi_{825} :\equiv (\xi \in \mathbb{R} \to b \in \mathbb{R} \to \xi < b \to 0 < b - \xi) \to (\mathrm{ContOn}(f,[a,b]) \to \varphi_{328}) \to \varphi_{750}$
\item $\varphi_{826} :\equiv (\mathrm{ContOn}(f,[a,b]) \to \varphi_{680}) \to \varphi_{645} \to \neg (d \in \mathbb{R} \land 0 < d \land d \leq \delta \cdot 2^{-1} \land d \leq \xi - a)$
\item $\varphi_{827} :\equiv (\mathrm{ContOn}(f,[a,b]) \to \varphi_{681}) \to \varphi_{647} \to \neg (d \in \mathbb{R} \land 0 < d \land d \leq \delta \cdot 2^{-1} \land d \leq b - \xi)$
\item $\varphi_{828} :\equiv (\mathrm{ContOn}(f,[a,b]) \to \varphi_{697}) \to \mathrm{ContOn}(f,[a,b]) \to \varphi_{667}$
\item $\varphi_{829} :\equiv (\mathrm{ContOn}(f,[a,b]) \to \varphi_{607}) \to (d + \xi = \xi + d \to d + \xi < b \to \xi + d < b) \to \varphi_{700}$
\item $\varphi_{830} :\equiv (\mathrm{ContOn}(f,[a,b]) \to \varphi_{674}) \to \mathrm{ContOn}(f,[a,b]) \to \varphi_{696}$
\item $\varphi_{831} :\equiv (\forall x\,(\varphi_{318} \to x \in W \to x \leq b) \to \forall x\,(\varphi_{318}) \to \forall x \in W\,(x \leq b)) \to (\varphi_{318} \to \forall x\,(\varphi_{318})) \to \mathrm{ContOn}(f,[a,b]) \to \varphi_{311}$
\item $\varphi_{832} :\equiv (\mathrm{ContOn}(f,[a,b]) \to \varphi_{643}) \to (b - \xi + \xi = b \to d + \xi < b - \xi + \xi \to d + \xi < b) \to \varphi_{713}$
\item $\varphi_{833} :\equiv (\mathrm{ContOn}(f,[a,b]) \to \varphi_{653}) \to \varphi_{206} \to \mathrm{ContOn}(f,[a,b]) \to \varphi_{654}$
\item $\varphi_{834} :\equiv (\mathrm{ContOn}(f,[a,b]) \to \varphi_{720}) \to \mathrm{ContOn}(f,[a,b]) \to \varphi_{675}$
\item $\varphi_{835} :\equiv (\mathrm{ContOn}(f,[a,b]) \to \varphi_{682}) \to \mathrm{ContOn}(f,[a,b]) \to \varphi_{720}$
\item $\varphi_{836} :\equiv (\mathrm{ContOn}(f,[a,b]) \to a \in \mathbb{R} \to b \in \mathbb{R} \to \varphi_{92}) \to \varphi_{797}$
\item $\varphi_{837} :\equiv (\mathrm{ContOn}(f,[a,b]) \to \varphi_{676}) \to \varphi_{261} \to \mathrm{ContOn}(f,[a,b]) \to \varphi_{642}$
\item $\varphi_{838} :\equiv (\mathrm{ContOn}(f,[a,b]) \to \varphi_{696}) \to \mathrm{ContOn}(f,[a,b]) \to \varphi_{727}$
\item $\varphi_{839} :\equiv (\mathrm{ContOn}(f,[a,b]) \to \varphi_{727}) \to \mathrm{ContOn}(f,[a,b]) \to \varphi_{697}$
\item $\varphi_{840} :\equiv (\mathrm{ContOn}(f,[a,b]) \to a \in \mathbb{R} \to b \in \mathbb{R} \to \varphi_{92}) \to \varphi_{813}$
\item $\varphi_{841} :\equiv (\mathrm{ContOn}(f,[a,b]) \to \varphi_{646}) \to \varphi_{94} \to \varphi_{93} \to \mathrm{ContOn}(f,[a,b]) \to \varphi_{616}$
\item $\varphi_{842} :\equiv (\mathrm{RFn}(f) \to \xi + d \in \mathrm{dom}(f) \to f(\xi + d) \in \mathbb{R}) \to (\mathrm{ContOn}(f,[a,b]) \to a \in \mathbb{R} \to \varphi_{80}) \to \varphi_{798}$
\item $\varphi_{843} :\equiv (a \leq \xi \to \xi \leq \xi + d \to a \leq \xi + d) \to (\mathrm{ContOn}(f,[a,b]) \to \varphi_{317}) \to \varphi_{782}$
\item $\varphi_{844} :\equiv (\mathrm{RFn}(f) \to s \in \mathrm{dom}(f) \to f(s) \in \mathbb{R}) \to (\mathrm{ContOn}(f,[a,b]) \to a \in \mathbb{R} \to \varphi_{80}) \to \varphi_{821}$
\item $\varphi_{845} :\equiv (\mathrm{ContOn}(f,[a,b]) \to \varphi_{553}) \to (\varphi_{302} \to \varphi_{355}) \to \mathrm{ContOn}(f,[a,b]) \to \varphi_{580}$
\item $\varphi_{846} :\equiv (\mathrm{ContOn}(f,[a,b]) \to \varphi_{553}) \to (\varphi_{302} \to \varphi_{359}) \to \mathrm{ContOn}(f,[a,b]) \to \varphi_{581}$
\item $\varphi_{847} :\equiv ([a,b] \subseteq \mathrm{dom}(f) \to s \in [a,b] \to s \in \mathrm{dom}(f)) \to \varphi_{840}$
\item $\varphi_{848} :\equiv (d \in \mathbb{R} \to \xi \in \mathbb{R} \to d + \xi - \xi = d) \to (\mathrm{ContOn}(f,[a,b]) \to \varphi_{612}) \to \varphi_{741}$
\item $\varphi_{849} :\equiv (\xi \in \mathbb{R} \to d \in \mathbb{R} \to \xi + d = d + \xi) \to (\mathrm{ContOn}(f,[a,b]) \to \varphi_{328}) \to \varphi_{789}$
\item $\varphi_{850} :\equiv (d \in \mathbb{R} \to \xi \in \mathbb{R} \to d + \xi = \xi + d) \to (\mathrm{ContOn}(f,[a,b]) \to \varphi_{612}) \to \varphi_{743}$
\item $\varphi_{851} :\equiv (\xi \in \mathbb{R} \to d \in \mathbb{R} \to \xi - d \in \mathbb{R}) \to (\mathrm{ContOn}(f,[a,b]) \to \varphi_{328}) \to \varphi_{790}$
\item $\varphi_{852} :\equiv (\xi \in \mathbb{R} \to d \in \mathbb{R} \to \xi + d \in \mathbb{R}) \to (\mathrm{ContOn}(f,[a,b]) \to \varphi_{328}) \to \varphi_{791}$
\item $\varphi_{853} :\equiv (\xi \in \mathbb{R} \to d \in \mathbb{R} \to \xi - d \in \mathbb{R}) \to (\mathrm{ContOn}(f,[a,b]) \to \varphi_{328}) \to \varphi_{805}$
\item $\varphi_{854} :\equiv (\mathrm{ContOn}(f,[a,b]) \to \varphi_{357}) \to (\mathrm{ContOn}(f,[a,b]) \to \varphi_{316}) \to \varphi_{759}$
\item $\varphi_{855} :\equiv \forall u\,(\varphi_{318} \to u \in W \to u \in \mathbb{R}) \to \varphi_{658} \to (\varphi_{318} \to \forall u\,(\varphi_{318})) \to \mathrm{ContOn}(f,[a,b]) \to \varphi_{323}$
\item $\varphi_{856} :\equiv (\delta \in \mathbb{R} \to 0 < \delta \to \delta \cdot 2^{-1} < \delta) \to \varphi_{578} \to \varphi_{576} \to \mathrm{ContOn}(f,[a,b]) \to \varphi_{633}$
\item $\varphi_{857} :\equiv (\delta \in \mathbb{R} \to 0 < \delta \to \delta \cdot 2^{-1} < \delta) \to \varphi_{579} \to \varphi_{577} \to \mathrm{ContOn}(f,[a,b]) \to \varphi_{634}$
\item $\varphi_{858} :\equiv (f(s) \in \mathbb{R} \to f(\xi) \in \mathbb{R} \to f(s) - f(\xi) \in \mathbb{R}) \to (\mathrm{ContOn}(f,[a,b]) \to \varphi_{673}) \to \varphi_{756}$
\item $\varphi_{859} :\equiv (\mathrm{ContOn}(f,[a,b]) \to \varphi_{717}) \to (\varphi_{159} \to \varphi_{259}) \to \mathrm{ContOn}(f,[a,b]) \to \varphi_{722}$
\item $\varphi_{860} :\equiv (\mathrm{ContOn}(f,[a,b]) \to \varphi_{726}) \to (\varphi_{157} \to \varphi_{131}) \to \mathrm{ContOn}(f,[a,b]) \to \varphi_{725}$
\item $\varphi_{861} :\equiv \forall W\,(\varphi_{326} \to \neg \forall x\,(x \in W \leftrightarrow x \in [a,b] \land f(x) < \mu)) \to \varphi_{719} \to (\varphi_{326} \to \forall W\,(\varphi_{326})) \to \varphi_{324}$
\item $\varphi_{862} :\equiv (d \leq \delta \cdot 2^{-1} \to \delta \cdot 2^{-1} < \delta \to d < \delta) \to (\mathrm{ContOn}(f,[a,b]) \to \varphi_{619}) \to \varphi_{787}$
\item $\varphi_{863} :\equiv (d \leq \delta \cdot 2^{-1} \to \delta \cdot 2^{-1} < \delta \to d < \delta) \to (\mathrm{ContOn}(f,[a,b]) \to \varphi_{622}) \to \varphi_{788}$
\item $\varphi_{864} :\equiv (\mathrm{ContOn}(f,[a,b]) \to \varphi_{641}) \to (\mathrm{ContOn}(f,[a,b]) \to \varphi_{604}) \to \varphi_{724}$
\item $\varphi_{865} :\equiv (\mathrm{ContOn}(f,[a,b]) \to \varphi_{654}) \to (\mathrm{ContOn}(f,[a,b]) \to \varphi_{605}) \to \varphi_{723}$
\item $\varphi_{866} :\equiv (\lvert s - \xi \rvert < d \to d < \delta \to \lvert s - \xi \rvert < \delta) \to (\mathrm{ContOn}(f,[a,b]) \to \varphi_{672}) \to \varphi_{804}$
\item $\varphi_{867} :\equiv (\mu < f(s) \to f(s) < \mu \to \mu < \mu) \to (\mathrm{ContOn}(f,[a,b]) \to \varphi_{667}) \to \varphi_{819}$
\item $\varphi_{868} :\equiv (\forall \xi\,(\varphi_{459}) \to \varphi_{465}) \to \varphi_{716} \to \varphi_{1} \to (\varphi_{1} \to \varphi_{2}) \to \varphi_{3} \to \mathrm{ContOn}(f,[a,b]) \to \varphi_{429}$
\item $\varphi_{869} :\equiv (\mathrm{ContOn}(f,[a,b]) \to \varphi_{328}) \to (\mathrm{ContOn}(f,[a,b]) \to \varphi_{612}) \to \varphi_{775}$
\item $\varphi_{870} :\equiv (\mathrm{ContOn}(f,[a,b]) \to \varphi_{328}) \to (\mathrm{ContOn}(f,[a,b]) \to \varphi_{610}) \to \varphi_{780}$
\item $\varphi_{871} :\equiv (\mathrm{ContOn}(f,[a,b]) \to \varphi_{328}) \to (\mathrm{ContOn}(f,[a,b]) \to \varphi_{612}) \to \varphi_{781}$
\item $\varphi_{872} :\equiv (\mathrm{ContOn}(f,[a,b]) \to \varphi_{629}) \to (\mathrm{ContOn}(f,[a,b]) \to \varphi_{328}) \to \varphi_{792}$
\item $\varphi_{873} :\equiv (\mathrm{ContOn}(f,[a,b]) \to \varphi_{442}) \to (\mathrm{ContOn}(f,[a,b]) \to \varphi_{644}) \to \varphi_{795}$
\item $\varphi_{874} :\equiv (\mathrm{UB}(\xi - d,W) \leftrightarrow \xi - d \in \mathbb{R} \land \forall x \in W\,(x \leq \xi - d)) \to (\mathrm{ContOn}(f,[a,b]) \to \varphi_{628}) \to \varphi_{841}$
\item $\varphi_{875} :\equiv (\mathrm{ContOn}(f,[a,b]) \to \varphi_{328}) \to (\mathrm{ContOn}(f,[a,b]) \to \varphi_{317}) \to \varphi_{854}$
\item $\varphi_{876} :\equiv (\mathrm{ContOn}(f,[a,b]) \to \varphi_{347}) \to (\mathrm{ContOn}(f,[a,b]) \to \varphi_{651}) \to \varphi_{824}$
\item $\varphi_{877} :\equiv (\mathrm{ContOn}(f,[a,b]) \to \varphi_{630}) \to (\mathrm{ContOn}(f,[a,b]) \to \varphi_{609}) \to \varphi_{793}$
\item $\varphi_{878} :\equiv (\mathrm{ContOn}(f,[a,b]) \to \varphi_{347}) \to \varphi_{56} \to (\mathrm{ContOn}(f,[a,b]) \to \varphi_{673}) \to \varphi_{828}$
\item $\varphi_{879} :\equiv (\mathrm{ContOn}(f,[a,b]) \to \varphi_{610}) \to (\mathrm{ContOn}(f,[a,b]) \to \varphi_{620}) \to \varphi_{820}$
\item $\varphi_{880} :\equiv (\xi + d \in [a,b] \leftrightarrow \xi + d \in \mathbb{R} \land a \leq \xi + d \land \xi + d \leq b) \to \varphi_{877}$
\item $\varphi_{881} :\equiv (\mathrm{ContOn}(f,[a,b]) \to \varphi_{670}) \to (\mathrm{ContOn}(f,[a,b]) \to \varphi_{639}) \to \varphi_{818}$
\item $\varphi_{882} :\equiv (\mathrm{ContOn}(f,[a,b]) \to \varphi_{722}) \to (\mathrm{ContOn}(f,[a,b]) \to \varphi_{644}) \to \varphi_{822}$
\item $\varphi_{883} :\equiv (\mathrm{ContOn}(f,[a,b]) \to \varphi_{665}) \to (\mathrm{ContOn}(f,[a,b]) \to \varphi_{651}) \to \varphi_{835}$
\item $\varphi_{884} :\equiv (\mathrm{ContOn}(f,[a,b]) \to \varphi_{725}) \to (\mathrm{ContOn}(f,[a,b]) \to \varphi_{652}) \to \varphi_{830}$
\item $\varphi_{885} :\equiv (\mathrm{ContOn}(f,[a,b]) \to \varphi_{683}) \to (\mathrm{ContOn}(f,[a,b]) \to \varphi_{677}) \to \varphi_{838}$
\item $\varphi_{886} :\equiv (\delta \cdot 2^{-1} \in \mathbb{R} \to \xi - a \in \mathbb{R} \to 0 < \delta \cdot 2^{-1} \to \varphi_{37}) \to \varphi_{590} \to \varphi_{586} \to \varphi_{588} \to \varphi_{582} \to \mathrm{ContOn}(f,[a,b]) \to \varphi_{596}$
\item $\varphi_{887} :\equiv (\delta \cdot 2^{-1} \in \mathbb{R} \to b - \xi \in \mathbb{R} \to 0 < \delta \cdot 2^{-1} \to \varphi_{38}) \to \varphi_{591} \to \varphi_{587} \to \varphi_{589} \to \varphi_{583} \to \mathrm{ContOn}(f,[a,b]) \to \varphi_{597}$
\item $\varphi_{888} :\equiv \forall \delta\,(\varphi_{481} \to \neg (\delta \in \mathbb{R} \land 0 < \delta \land \varphi_{157})) \to \varphi_{776} \to (\varphi_{481} \to \forall \delta\,(\varphi_{481})) \to \varphi_{478}$
\item $\varphi_{889} :\equiv \forall \delta\,(\varphi_{482} \to \neg (\delta \in \mathbb{R} \land 0 < \delta \land \varphi_{159})) \to \varphi_{777} \to (\varphi_{482} \to \forall \delta\,(\varphi_{482})) \to \varphi_{479}$
\item $\varphi_{890} :\equiv (\varphi_{702} \to \varphi_{707}) \to (\varphi_{645} \to \forall d\,(\varphi_{645})) \to \mathrm{ContOn}(f,[a,b]) \to \varphi_{594}$
\item $\varphi_{891} :\equiv (\varphi_{703} \to \varphi_{708}) \to (\varphi_{647} \to \forall d\,(\varphi_{647})) \to \mathrm{ContOn}(f,[a,b]) \to \varphi_{595}$
\item $\varphi_{892} :\equiv (f(\xi + d) \in \mathbb{R} \to f(\xi) \in \mathbb{R} \to \mu - f(\xi) \in \mathbb{R} \to \varphi_{224}) \to (\mathrm{ContOn}(f,[a,b]) \to \varphi_{648}) \to \varphi_{876}$
\item $\varphi_{893} :\equiv (\mathrm{ContOn}(f,[a,b]) \to \varphi_{639}) \to (\mathrm{ContOn}(f,[a,b]) \to \varphi_{328}) \to \varphi_{879}$
\item $\varphi_{894} :\equiv \forall s\,(\varphi_{636} \to s \in W \to s \leq \xi - d) \to \varphi_{803} \to (\varphi_{636} \to \forall s\,(\varphi_{636})) \to \mathrm{ContOn}(f,[a,b]) \to \varphi_{646}$
\end{abbrevlist}
\end{fullproof}

% 自動生成 (prover/build.js): ConstConv  ノード数 148
\begin{fullproof}{[G4]}{$\forall c \in \mathbb{R}\,(\mathrm{Conv}(\overline{c},c))$}{148}\phantomsection\label{tree:ConstConv}
\noindent$\Pi_{1}$:
\begin{ptree}
\AXC{$x = x$\lbl{A5}}
\AXC{$x = y \to x = x \to y = x$\lbl{A6}}
\AXC{$(x = y \to x = x \to y = x) \to x = x \to x = y \to y = x$\lbl{A1}}
\BIC{$x = x \to x = y \to y = x$}
\BIC{$x = y \to y = x$}
\UIC{$\forall y\,(x = y \to y = x)$}
\UIC{$\forall x\,\forall y\,(x = y \to y = x)$}
\AXC{$\varphi_{19}$\lbl{A2}}
\BIC{$\forall y\,(\lvert \overline{c}(n) - c \rvert = y \to y = \lvert \overline{c}(n) - c \rvert)$}
\AXC{$\varphi_{35}$\lbl{A2}}
\BIC{$\lvert \overline{c}(n) - c \rvert = 0 \to 0 = \lvert \overline{c}(n) - c \rvert$}
\end{ptree}
\noindent$\Pi_{2}$:
\begin{ptree}
\AXC{$y = z \to x = y \to x = z$\lbl{A6}}
\AXC{$(y = z \to x = y \to x = z) \to x = y \to y = z \to x = z$\lbl{A1}}
\BIC{$x = y \to y = z \to x = z$}
\UIC{$\forall z\,(x = y \to y = z \to x = z)$}
\UIC{$\forall y\,\forall z\,(x = y \to y = z \to x = z)$}
\UIC{$\forall x\,\forall y\,\forall z\,(x = y \to y = z \to x = z)$}
\AXC{$\forall x\,\forall y\,\forall z\,(x = y \to y = z \to x = z) \to \varphi_{5}$\lbl{A2}}
\BIC{$\varphi_{5}$}
\end{ptree}
\noindent$\Pi_{3}$:
\begin{ptree}
\AXC{$\forall x\,(0 \leq x \to \lvert x \rvert = x)$\lbl{[A1]}}
\AXC{$\forall x\,(0 \leq x \to \lvert x \rvert = x) \to 0 \leq 0 \to \lvert 0 \rvert = 0$\lbl{A2}}
\BIC{$0 \leq 0 \to \lvert 0 \rvert = 0$}
\AXC{$(0 \leq 0 \to \lvert 0 \rvert = 0) \to \varphi_{37}$\lbl{A1}}
\BIC{$\varphi_{37}$}
\end{ptree}
\noindent$\Pi_{4}$:
\begin{ptree}
\AXC{$y = z \to x = y \to x = z$\lbl{A6}}
\AXC{$(y = z \to x = y \to x = z) \to x = y \to y = z \to x = z$\lbl{A1}}
\BIC{$x = y \to y = z \to x = z$}
\UIC{$\forall z\,(x = y \to y = z \to x = z)$}
\UIC{$\forall y\,\forall z\,(x = y \to y = z \to x = z)$}
\UIC{$\forall x\,\forall y\,\forall z\,(x = y \to y = z \to x = z)$}
\AXC{$\forall x\,\forall y\,\forall z\,(x = y \to y = z \to x = z) \to \varphi_{5}$\lbl{A2}}
\BIC{$\varphi_{5}$}
\end{ptree}
\noindent$\Pi_{5}$:
\begin{ptree}
\AXC{$z = w \to \lvert c - c \rvert = \lvert z \rvert \to \lvert c - c \rvert = \lvert w \rvert$\lbl{A6}}
\UIC{$\forall w\,(z = w \to \lvert c - c \rvert = \lvert z \rvert \to \lvert c - c \rvert = \lvert w \rvert)$}
\UIC{$\forall z\,\forall w\,(z = w \to \lvert c - c \rvert = \lvert z \rvert \to \lvert c - c \rvert = \lvert w \rvert)$}
\AXC{$\varphi_{34}$\lbl{A2}}
\BIC{$\varphi_{1}$}
\AXC{$\varphi_{1} \to c - c = 0 \to \lvert c - c \rvert = \lvert c - c \rvert \to \lvert c - c \rvert = \lvert 0 \rvert$\lbl{A2}}
\BIC{$c - c = 0 \to \lvert c - c \rvert = \lvert c - c \rvert \to \lvert c - c \rvert = \lvert 0 \rvert$}
\end{ptree}
\noindent$\Pi_{6}$:
\begin{ptree}
\AXC{$\varphi_{11}$\lbl{A1}}
\AXC{$\forall x \in \mathbb{R}\,(x - x = 0)$\lbl{[F13]}}
\AXC{$\forall x \in \mathbb{R}\,(x - x = 0) \to c \in \mathbb{R} \to c - c = 0$\lbl{A2}}
\BIC{$c \in \mathbb{R} \to c - c = 0$}
\AXC{$(c \in \mathbb{R} \to c - c = 0) \to \varphi_{11} \to \varphi_{7}$\lbl{A1}}
\BIC{$\varphi_{11} \to \varphi_{7}$}
\BIC{$\varphi_{7}$}
\end{ptree}
\noindent$\Pi_{7}$:
\begin{ptree}
\AXC{$x = x$\lbl{A5}}
\UIC{$\forall x\,(x = x)$}
\AXC{$\forall x\,(x = x) \to \lvert c - c \rvert = \lvert c - c \rvert$\lbl{A2}}
\BIC{$\lvert c - c \rvert = \lvert c - c \rvert$}
\AXC{$\stackrel{\Pi_{6}}{\vdots}$}\noLine
\UIC{$\varphi_{7}$}
\AXC{$\varphi_{7} \to \lvert c - c \rvert = \lvert c - c \rvert \to \varphi_{38}$\lbl{A1}}
\BIC{$\lvert c - c \rvert = \lvert c - c \rvert \to \varphi_{38}$}
\BIC{$\varphi_{38}$}
\end{ptree}
\noindent$\Pi_{8}$:
\begin{ptree}
\AXC{$z = w \to \varphi_{17}$\lbl{A6}}
\UIC{$\forall w\,(z = w \to \varphi_{17})$}
\UIC{$\forall z\,\forall w\,(z = w \to \varphi_{17})$}
\AXC{$\forall z\,\forall w\,(z = w \to \varphi_{17}) \to \varphi_{33}$\lbl{A2}}
\BIC{$\varphi_{33}$}
\AXC{$\varphi_{33} \to \overline{c}(n) = c \to \varphi_{28}$\lbl{A2}}
\BIC{$\overline{c}(n) = c \to \varphi_{28}$}
\end{ptree}
\noindent$\Pi_{9}$:
\begin{ptree}
\AXC{$\forall c \in \mathbb{R}\,(\mathrm{Seq}(\overline{c}) \land \forall n \in \mathbb{N}\,(\overline{c}(n) = c))$\lbl{[S10]}}
\AXC{$\varphi_{36}$\lbl{A2}}
\BIC{$c \in \mathbb{R} \to \mathrm{Seq}(\overline{c}) \land \forall n \in \mathbb{N}\,(\overline{c}(n) = c)$}
\AXC{$\varphi_{44}$\lbl{A1}}
\BIC{$c \in \mathbb{R} \to \varepsilon \in \mathbb{R} \to \varphi_{3}$}
\AXC{$(c \in \mathbb{R} \to \varepsilon \in \mathbb{R} \to \varphi_{3}) \to \varphi_{45}$\lbl{A1}}
\BIC{$\varphi_{45}$}
\end{ptree}
\noindent$\Pi_{10}$:
\begin{ptree}
\AXC{$c \in \mathbb{R} \to \varepsilon \in \mathbb{R} \to 0 < \varepsilon \to n \in \mathbb{N} \to n \in \mathbb{N}$\lbl{A1}}
\AXC{$\forall n \in \mathbb{N}\,(\overline{c}(n) = c) \to n \in \mathbb{N} \to \overline{c}(n) = c$\lbl{A2}}
\AXC{$\stackrel{\Pi_{9}}{\vdots}$}\noLine
\UIC{$\varphi_{45}$}
\BIC{$c \in \mathbb{R} \to \varphi_{18}$}
\AXC{$\varphi_{49}$\lbl{A1}}
\BIC{$(c \in \mathbb{R} \to \varepsilon \in \mathbb{R} \to 0 < \varepsilon \to n \in \mathbb{N} \to n \in \mathbb{N}) \to \varphi_{20}$}
\BIC{$\varphi_{20}$}
\end{ptree}
\noindent$\Pi_{11}$:
\begin{ptree}
\AXC{$x = x$\lbl{A5}}
\UIC{$\forall x\,(x = x)$}
\AXC{$\forall x\,(x = x) \to \lvert \overline{c}(n) - c \rvert = \lvert \overline{c}(n) - c \rvert$\lbl{A2}}
\BIC{$\lvert \overline{c}(n) - c \rvert = \lvert \overline{c}(n) - c \rvert$}
\AXC{$\stackrel{\Pi_{10}}{\vdots}$}\noLine
\UIC{$\varphi_{20}$}
\AXC{$\varphi_{20} \to \varphi_{50}$\lbl{A1}}
\BIC{$\varphi_{50}$}
\BIC{$(\overline{c}(n) = c \to \varphi_{28}) \to c \in \mathbb{R} \to \varepsilon \in \mathbb{R} \to \varphi_{9}$}
\end{ptree}
\noindent$\Pi_{12}$:
\begin{ptree}
\AXC{$\stackrel{\Pi_{8}}{\vdots}$}\noLine
\UIC{$\overline{c}(n) = c \to \varphi_{28}$}
\AXC{$\stackrel{\Pi_{11}}{\vdots}$}\noLine
\UIC{$(\overline{c}(n) = c \to \varphi_{28}) \to c \in \mathbb{R} \to \varepsilon \in \mathbb{R} \to \varphi_{9}$}
\BIC{$c \in \mathbb{R} \to \varepsilon \in \mathbb{R} \to \varphi_{9}$}
\AXC{$\varphi_{52}$\lbl{A1}}
\BIC{$(c \in \mathbb{R} \to \varphi_{4}) \to \varphi_{29} \to c \in \mathbb{R} \to \varphi_{22}$}
\end{ptree}
\noindent$\Pi_{13}$:
\begin{ptree}
\AXC{$\stackrel{\Pi_{5}}{\vdots}$}\noLine
\UIC{$c - c = 0 \to \lvert c - c \rvert = \lvert c - c \rvert \to \lvert c - c \rvert = \lvert 0 \rvert$}
\AXC{$\stackrel{\Pi_{7}}{\vdots}$}\noLine
\UIC{$\varphi_{38}$}
\BIC{$c \in \mathbb{R} \to \varphi_{4}$}
\AXC{$\stackrel{\Pi_{12}}{\vdots}$}\noLine
\UIC{$(c \in \mathbb{R} \to \varphi_{4}) \to \varphi_{29} \to c \in \mathbb{R} \to \varphi_{22}$}
\BIC{$\varphi_{29} \to c \in \mathbb{R} \to \varphi_{22}$}
\end{ptree}
\noindent$\Pi_{14}$:
\begin{ptree}
\AXC{$\stackrel{\Pi_{4}}{\vdots}$}\noLine
\UIC{$\varphi_{5}$}
\AXC{$\varphi_{5} \to \varphi_{26}$\lbl{A2}}
\BIC{$\varphi_{26}$}
\AXC{$\varphi_{26} \to \varphi_{29}$\lbl{A2}}
\BIC{$\varphi_{29}$}
\AXC{$\stackrel{\Pi_{13}}{\vdots}$}\noLine
\UIC{$\varphi_{29} \to c \in \mathbb{R} \to \varphi_{22}$}
\BIC{$c \in \mathbb{R} \to \varphi_{22}$}
\AXC{$(c \in \mathbb{R} \to \varphi_{22}) \to \varphi_{6} \to \varphi_{13} \to c \in \mathbb{R} \to \varphi_{14}$\lbl{A1}}
\BIC{$\varphi_{6} \to \varphi_{13} \to c \in \mathbb{R} \to \varphi_{14}$}
\end{ptree}
\noindent$\Pi_{15}$:
\begin{ptree}
\AXC{$x = x$\lbl{A5}}
\UIC{$\forall x\,(x = x)$}
\AXC{$\forall x\,(x = x) \to 0 = 0$\lbl{A2}}
\BIC{$0 = 0$}
\AXC{$\varphi_{10}$\lbl{A1}}
\BIC{$c \in \mathbb{R} \to \varepsilon \in \mathbb{R} \to 0 < \varepsilon \to n \in \mathbb{N} \to 1_{\omega} \leq n \to 0 \leq 0$}
\AXC{$\stackrel{\Pi_{3}}{\vdots}$}\noLine
\UIC{$\varphi_{37}$}
\BIC{$\varphi_{6}$}
\AXC{$\stackrel{\Pi_{14}}{\vdots}$}\noLine
\UIC{$\varphi_{6} \to \varphi_{13} \to c \in \mathbb{R} \to \varphi_{14}$}
\BIC{$\varphi_{13} \to c \in \mathbb{R} \to \varphi_{14}$}
\end{ptree}
\noindent$\Pi_{16}$:
\begin{ptree}
\AXC{$\stackrel{\Pi_{2}}{\vdots}$}\noLine
\UIC{$\varphi_{5}$}
\AXC{$\varphi_{5} \to \varphi_{16}$\lbl{A2}}
\BIC{$\varphi_{16}$}
\AXC{$\varphi_{16} \to \varphi_{13}$\lbl{A2}}
\BIC{$\varphi_{13}$}
\AXC{$\stackrel{\Pi_{15}}{\vdots}$}\noLine
\UIC{$\varphi_{13} \to c \in \mathbb{R} \to \varphi_{14}$}
\BIC{$c \in \mathbb{R} \to \varphi_{14}$}
\AXC{$(c \in \mathbb{R} \to \varphi_{14}) \to \varphi_{41}$\lbl{A1}}
\BIC{$\varphi_{41}$}
\end{ptree}
\noindent$\Pi_{17}$:
\begin{ptree}
\AXC{$\stackrel{\Pi_{1}}{\vdots}$}\noLine
\UIC{$\lvert \overline{c}(n) - c \rvert = 0 \to 0 = \lvert \overline{c}(n) - c \rvert$}
\AXC{$\stackrel{\Pi_{16}}{\vdots}$}\noLine
\UIC{$\varphi_{41}$}
\BIC{$c \in \mathbb{R} \to \varphi_{15}$}
\AXC{$(c \in \mathbb{R} \to \varphi_{15}) \to (c \in \mathbb{R} \to \varepsilon \in \mathbb{R} \to 0 < \varepsilon \to 0 < \varepsilon) \to \varphi_{46}$\lbl{A1}}
\BIC{$(c \in \mathbb{R} \to \varepsilon \in \mathbb{R} \to 0 < \varepsilon \to 0 < \varepsilon) \to \varphi_{46}$}
\end{ptree}
\noindent$\Pi_{18}$:
\begin{ptree}
\AXC{$c \in \mathbb{R} \to \varepsilon \in \mathbb{R} \to 0 < \varepsilon \to 0 < \varepsilon$\lbl{A1}}
\AXC{$\stackrel{\Pi_{17}}{\vdots}$}\noLine
\UIC{$(c \in \mathbb{R} \to \varepsilon \in \mathbb{R} \to 0 < \varepsilon \to 0 < \varepsilon) \to \varphi_{46}$}
\BIC{$\varphi_{46}$}
\end{ptree}
\noindent$\Pi_{19}$:
\begin{ptree}
\AXC{$z = w \to z < \varepsilon \to w < \varepsilon$\lbl{A6}}
\UIC{$\forall w\,(z = w \to z < \varepsilon \to w < \varepsilon)$}
\UIC{$\forall z\,\forall w\,(z = w \to z < \varepsilon \to w < \varepsilon)$}
\AXC{$\varphi_{2}$\lbl{A2}}
\BIC{$\forall w\,(0 = w \to 0 < \varepsilon \to w < \varepsilon)$}
\AXC{$\varphi_{32}$\lbl{A2}}
\BIC{$0 = \lvert \overline{c}(n) - c \rvert \to 0 < \varepsilon \to \lvert \overline{c}(n) - c \rvert < \varepsilon$}
\AXC{$\stackrel{\Pi_{18}}{\vdots}$}\noLine
\UIC{$\varphi_{46}$}
\BIC{$c \in \mathbb{R} \to \varphi_{23}$}
\AXC{$(c \in \mathbb{R} \to \varphi_{23}) \to \varphi_{30}$\lbl{A1}}
\BIC{$\varphi_{30}$}
\UIC{$\forall n\,(\varphi_{30})$}
\end{ptree}
\noindent$\Pi_{20}$:
\begin{ptree}
\AXC{$\varphi_{21}$\lbl{A4}}
\AXC{$\forall n\,(\varphi_{30}) \to \varphi_{31}$\lbl{A3}}
\AXC{$\stackrel{\Pi_{19}}{\vdots}$}\noLine
\UIC{$\forall n\,(\varphi_{30})$}
\AXC{$\forall n\,(\varphi_{30}) \to (\forall n\,(\varphi_{30}) \to \varphi_{31}) \to \varphi_{21} \to c \in \mathbb{R} \to \varphi_{25}$\lbl{A1}}
\BIC{$(\forall n\,(\varphi_{30}) \to \varphi_{31}) \to \varphi_{21} \to c \in \mathbb{R} \to \varphi_{25}$}
\BIC{$\varphi_{21} \to c \in \mathbb{R} \to \varphi_{25}$}
\BIC{$c \in \mathbb{R} \to \varphi_{25}$}
\end{ptree}
\noindent$\Pi_{21}$:
\begin{ptree}
\AXC{$\varphi_{8} \to \neg \varphi_{12}$\lbl{A2}}
\AXC{$\stackrel{\Pi_{20}}{\vdots}$}\noLine
\UIC{$c \in \mathbb{R} \to \varphi_{25}$}
\AXC{$1_{\omega} \in \mathbb{N}$\lbl{[N1]}}
\AXC{$\varphi_{47}$\lbl{A1}}
\BIC{$(c \in \mathbb{R} \to \varphi_{25}) \to c \in \mathbb{R} \to \varepsilon \in \mathbb{R} \to 0 < \varepsilon \to \varphi_{12}$}
\BIC{$c \in \mathbb{R} \to \varepsilon \in \mathbb{R} \to 0 < \varepsilon \to \varphi_{12}$}
\AXC{$\varphi_{51}$\lbl{A1}}
\BIC{$(\varphi_{8} \to \neg \varphi_{12}) \to c \in \mathbb{R} \to \varepsilon \in \mathbb{R} \to \varphi_{24}$}
\BIC{$c \in \mathbb{R} \to \varepsilon \in \mathbb{R} \to \varphi_{24}$}
\UIC{$\forall \varepsilon\,(c \in \mathbb{R} \to \varepsilon \in \mathbb{R} \to \varphi_{24})$}
\end{ptree}
\noindent$\Pi_{22}$:
\begin{ptree}
\AXC{$c \in \mathbb{R} \to \forall \varepsilon\,(c \in \mathbb{R})$\lbl{A4}}
\AXC{$\varphi_{48}$\lbl{A3}}
\AXC{$\stackrel{\Pi_{21}}{\vdots}$}\noLine
\UIC{$\forall \varepsilon\,(c \in \mathbb{R} \to \varepsilon \in \mathbb{R} \to \varphi_{24})$}
\AXC{$\forall \varepsilon\,(c \in \mathbb{R} \to \varepsilon \in \mathbb{R} \to \varphi_{24}) \to \varphi_{53}$\lbl{A1}}
\BIC{$\varphi_{53}$}
\BIC{$(c \in \mathbb{R} \to \forall \varepsilon\,(c \in \mathbb{R})) \to c \in \mathbb{R} \to \forall \varepsilon \in \mathbb{R}\,(\varphi_{24})$}
\BIC{$c \in \mathbb{R} \to \forall \varepsilon \in \mathbb{R}\,(\varphi_{24})$}
\end{ptree}
\noindent$\Pi_{23}$:
\begin{ptree}
\AXC{$\forall c \in \mathbb{R}\,(\mathrm{Seq}(\overline{c}) \land \forall n \in \mathbb{N}\,(\overline{c}(n) = c))$\lbl{[S10]}}
\AXC{$\varphi_{36}$\lbl{A2}}
\BIC{$c \in \mathbb{R} \to \mathrm{Seq}(\overline{c}) \land \forall n \in \mathbb{N}\,(\overline{c}(n) = c)$}
\AXC{$\varphi_{27}$\lbl{A1}}
\BIC{$c \in \mathbb{R} \to \mathrm{Seq}(\overline{c})$}
\end{ptree}
\noindent$\Pi_{24}$:
\begin{ptree}
\AXC{$\varphi_{39}$\lbl{DefConv}}
\AXC{$\varphi_{39} \to \varphi_{43}$\lbl{A2}}
\BIC{$\varphi_{43}$}
\AXC{$\varphi_{43} \to \varphi_{42}$\lbl{A2}}
\BIC{$\varphi_{42}$}
\AXC{$\varphi_{42} \to (c \in \mathbb{R} \to \mathrm{Seq}(\overline{c})) \to (c \in \mathbb{R} \to c \in \mathbb{R}) \to \varphi_{40}$\lbl{A1}}
\BIC{$(c \in \mathbb{R} \to \mathrm{Seq}(\overline{c})) \to (c \in \mathbb{R} \to c \in \mathbb{R}) \to \varphi_{40}$}
\end{ptree}
\noindent$\Pi_{25}$:
\begin{ptree}
\AXC{$c \in \mathbb{R} \to c \in \mathbb{R}$\lbl{A1}}
\AXC{$\stackrel{\Pi_{23}}{\vdots}$}\noLine
\UIC{$c \in \mathbb{R} \to \mathrm{Seq}(\overline{c})$}
\AXC{$\stackrel{\Pi_{24}}{\vdots}$}\noLine
\UIC{$(c \in \mathbb{R} \to \mathrm{Seq}(\overline{c})) \to (c \in \mathbb{R} \to c \in \mathbb{R}) \to \varphi_{40}$}
\BIC{$(c \in \mathbb{R} \to c \in \mathbb{R}) \to \varphi_{40}$}
\BIC{$\varphi_{40}$}
\end{ptree}
\noindent 主木:
\begin{ptree}
\AXC{$\stackrel{\Pi_{22}}{\vdots}$}\noLine
\UIC{$c \in \mathbb{R} \to \forall \varepsilon \in \mathbb{R}\,(\varphi_{24})$}
\AXC{$\stackrel{\Pi_{25}}{\vdots}$}\noLine
\UIC{$\varphi_{40}$}
\BIC{$c \in \mathbb{R} \to \mathrm{Conv}(\overline{c},c)$}
\UIC{$\forall c \in \mathbb{R}\,(\mathrm{Conv}(\overline{c},c))$}
\end{ptree}
\begin{abbrevlist}
\item $\varphi_{1} :\equiv \forall w\,(c - c = w \to \lvert c - c \rvert = \lvert c - c \rvert \to \lvert c - c \rvert = \lvert w \rvert)$
\item $\varphi_{2} :\equiv \forall z\,\forall w\,(z = w \to z < \varepsilon \to w < \varepsilon) \to \forall w\,(0 = w \to 0 < \varepsilon \to w < \varepsilon)$
\item $\varphi_{3} :\equiv 0 < \varepsilon \to n \in \mathbb{N} \to 1_{\omega} \leq n \to \forall n \in \mathbb{N}\,(\overline{c}(n) = c)$
\item $\varphi_{4} :\equiv \varepsilon \in \mathbb{R} \to 0 < \varepsilon \to n \in \mathbb{N} \to 1_{\omega} \leq n \to \lvert c - c \rvert = \lvert 0 \rvert$
\item $\varphi_{5} :\equiv \forall y\,\forall z\,(\lvert \overline{c}(n) - c \rvert = y \to y = z \to \lvert \overline{c}(n) - c \rvert = z)$
\item $\varphi_{6} :\equiv c \in \mathbb{R} \to \varepsilon \in \mathbb{R} \to 0 < \varepsilon \to n \in \mathbb{N} \to 1_{\omega} \leq n \to \lvert 0 \rvert = 0$
\item $\varphi_{7} :\equiv c \in \mathbb{R} \to \varepsilon \in \mathbb{R} \to 0 < \varepsilon \to n \in \mathbb{N} \to 1_{\omega} \leq n \to c - c = 0$
\item $\varphi_{8} :\equiv \forall N\,(\neg (N \in \mathbb{N} \land \forall n \in \mathbb{N}\,(N \leq n \to \lvert \overline{c}(n) - c \rvert < \varepsilon)))$
\item $\varphi_{9} :\equiv 0 < \varepsilon \to n \in \mathbb{N} \to 1_{\omega} \leq n \to \lvert \overline{c}(n) - c \rvert = \lvert c - c \rvert$
\item $\varphi_{10} :\equiv 0 = 0 \to c \in \mathbb{R} \to \varepsilon \in \mathbb{R} \to 0 < \varepsilon \to n \in \mathbb{N} \to 1_{\omega} \leq n \to 0 \leq 0$
\item $\varphi_{11} :\equiv c \in \mathbb{R} \to \varepsilon \in \mathbb{R} \to 0 < \varepsilon \to n \in \mathbb{N} \to 1_{\omega} \leq n \to c \in \mathbb{R}$
\item $\varphi_{12} :\equiv 1_{\omega} \in \mathbb{N} \land \forall n \in \mathbb{N}\,(1_{\omega} \leq n \to \lvert \overline{c}(n) - c \rvert < \varepsilon)$
\item $\varphi_{13} :\equiv \lvert \overline{c}(n) - c \rvert = \lvert 0 \rvert \to \lvert 0 \rvert = 0 \to \lvert \overline{c}(n) - c \rvert = 0$
\item $\varphi_{14} :\equiv \varepsilon \in \mathbb{R} \to 0 < \varepsilon \to n \in \mathbb{N} \to 1_{\omega} \leq n \to \lvert \overline{c}(n) - c \rvert = 0$
\item $\varphi_{15} :\equiv \varepsilon \in \mathbb{R} \to 0 < \varepsilon \to n \in \mathbb{N} \to 1_{\omega} \leq n \to 0 = \lvert \overline{c}(n) - c \rvert$
\item $\varphi_{16} :\equiv \forall z\,(\lvert \overline{c}(n) - c \rvert = \lvert 0 \rvert \to \lvert 0 \rvert = z \to \lvert \overline{c}(n) - c \rvert = z)$
\item $\varphi_{17} :\equiv \lvert \overline{c}(n) - c \rvert = \lvert z - c \rvert \to \lvert \overline{c}(n) - c \rvert = \lvert w - c \rvert$
\item $\varphi_{18} :\equiv \varepsilon \in \mathbb{R} \to 0 < \varepsilon \to n \in \mathbb{N} \to 1_{\omega} \leq n \to n \in \mathbb{N} \to \overline{c}(n) = c$
\item $\varphi_{19} :\equiv \forall x\,\forall y\,(x = y \to y = x) \to \forall y\,(\lvert \overline{c}(n) - c \rvert = y \to y = \lvert \overline{c}(n) - c \rvert)$
\item $\varphi_{20} :\equiv c \in \mathbb{R} \to \varepsilon \in \mathbb{R} \to 0 < \varepsilon \to n \in \mathbb{N} \to 1_{\omega} \leq n \to \overline{c}(n) = c$
\item $\varphi_{21} :\equiv c \in \mathbb{R} \land \varepsilon \in \mathbb{R} \land 0 < \varepsilon \to \forall n\,(c \in \mathbb{R} \land \varepsilon \in \mathbb{R} \land 0 < \varepsilon)$
\item $\varphi_{22} :\equiv \varepsilon \in \mathbb{R} \to 0 < \varepsilon \to n \in \mathbb{N} \to 1_{\omega} \leq n \to \lvert \overline{c}(n) - c \rvert = \lvert 0 \rvert$
\item $\varphi_{23} :\equiv \varepsilon \in \mathbb{R} \to 0 < \varepsilon \to n \in \mathbb{N} \to 1_{\omega} \leq n \to \lvert \overline{c}(n) - c \rvert < \varepsilon$
\item $\varphi_{24} :\equiv 0 < \varepsilon \to \exists N \in \mathbb{N}\,\forall n \in \mathbb{N}\,(N \leq n \to \lvert \overline{c}(n) - c \rvert < \varepsilon)$
\item $\varphi_{25} :\equiv \varepsilon \in \mathbb{R} \to 0 < \varepsilon \to \forall n \in \mathbb{N}\,(1_{\omega} \leq n \to \lvert \overline{c}(n) - c \rvert < \varepsilon)$
\item $\varphi_{26} :\equiv \forall z\,(\lvert \overline{c}(n) - c \rvert = \lvert c - c \rvert \to \lvert c - c \rvert = z \to \lvert \overline{c}(n) - c \rvert = z)$
\item $\varphi_{27} :\equiv (c \in \mathbb{R} \to \mathrm{Seq}(\overline{c}) \land \forall n \in \mathbb{N}\,(\overline{c}(n) = c)) \to c \in \mathbb{R} \to \mathrm{Seq}(\overline{c})$
\item $\varphi_{28} :\equiv \lvert \overline{c}(n) - c \rvert = \lvert \overline{c}(n) - c \rvert \to \lvert \overline{c}(n) - c \rvert = \lvert c - c \rvert$
\item $\varphi_{29} :\equiv \lvert \overline{c}(n) - c \rvert = \lvert c - c \rvert \to \lvert c - c \rvert = \lvert 0 \rvert \to \lvert \overline{c}(n) - c \rvert = \lvert 0 \rvert$
\item $\varphi_{30} :\equiv c \in \mathbb{R} \land \varepsilon \in \mathbb{R} \land 0 < \varepsilon \to n \in \mathbb{N} \to 1_{\omega} \leq n \to \lvert \overline{c}(n) - c \rvert < \varepsilon$
\item $\varphi_{31} :\equiv \forall n\,(c \in \mathbb{R} \land \varepsilon \in \mathbb{R} \land 0 < \varepsilon) \to \forall n \in \mathbb{N}\,(1_{\omega} \leq n \to \lvert \overline{c}(n) - c \rvert < \varepsilon)$
\item $\varphi_{32} :\equiv \forall w\,(0 = w \to 0 < \varepsilon \to w < \varepsilon) \to 0 = \lvert \overline{c}(n) - c \rvert \to 0 < \varepsilon \to \lvert \overline{c}(n) - c \rvert < \varepsilon$
\item $\varphi_{33} :\equiv \forall w\,(\overline{c}(n) = w \to \lvert \overline{c}(n) - c \rvert = \lvert \overline{c}(n) - c \rvert \to \lvert \overline{c}(n) - c \rvert = \lvert w - c \rvert)$
\item $\varphi_{34} :\equiv \forall z\,\forall w\,(z = w \to \lvert c - c \rvert = \lvert z \rvert \to \lvert c - c \rvert = \lvert w \rvert) \to \varphi_{1}$
\item $\varphi_{35} :\equiv \forall y\,(\lvert \overline{c}(n) - c \rvert = y \to y = \lvert \overline{c}(n) - c \rvert) \to \lvert \overline{c}(n) - c \rvert = 0 \to 0 = \lvert \overline{c}(n) - c \rvert$
\item $\varphi_{36} :\equiv \forall c \in \mathbb{R}\,(\mathrm{Seq}(\overline{c}) \land \forall n \in \mathbb{N}\,(\overline{c}(n) = c)) \to c \in \mathbb{R} \to \mathrm{Seq}(\overline{c}) \land \forall n \in \mathbb{N}\,(\overline{c}(n) = c)$
\item $\varphi_{37} :\equiv (c \in \mathbb{R} \to \varepsilon \in \mathbb{R} \to 0 < \varepsilon \to n \in \mathbb{N} \to 1_{\omega} \leq n \to 0 \leq 0) \to \varphi_{6}$
\item $\varphi_{38} :\equiv (c - c = 0 \to \lvert c - c \rvert = \lvert c - c \rvert \to \lvert c - c \rvert = \lvert 0 \rvert) \to c \in \mathbb{R} \to \varphi_{4}$
\item $\varphi_{39} :\equiv \forall a\,\forall \alpha\,(\mathrm{Conv}(a,\alpha) \leftrightarrow \mathrm{Seq}(a) \land \alpha \in \mathbb{R} \land \forall \varepsilon \in \mathbb{R}\,(0 < \varepsilon \to \exists N \in \mathbb{N}\,\forall n \in \mathbb{N}\,(N \leq n \to \lvert a_{n} - \alpha \rvert < \varepsilon)))$
\item $\varphi_{40} :\equiv (c \in \mathbb{R} \to \forall \varepsilon \in \mathbb{R}\,(\varphi_{24})) \to c \in \mathbb{R} \to \mathrm{Conv}(\overline{c},c)$
\item $\varphi_{41} :\equiv (\lvert \overline{c}(n) - c \rvert = 0 \to 0 = \lvert \overline{c}(n) - c \rvert) \to c \in \mathbb{R} \to \varphi_{15}$
\item $\varphi_{42} :\equiv \mathrm{Conv}(\overline{c},c) \leftrightarrow \mathrm{Seq}(\overline{c}) \land c \in \mathbb{R} \land \forall \varepsilon \in \mathbb{R}\,(\varphi_{24})$
\item $\varphi_{43} :\equiv \forall \alpha\,(\mathrm{Conv}(\overline{c},\alpha) \leftrightarrow \mathrm{Seq}(\overline{c}) \land \alpha \in \mathbb{R} \land \forall \varepsilon \in \mathbb{R}\,(0 < \varepsilon \to \exists N \in \mathbb{N}\,\forall n \in \mathbb{N}\,(N \leq n \to \lvert \overline{c}(n) - \alpha \rvert < \varepsilon)))$
\item $\varphi_{44} :\equiv (c \in \mathbb{R} \to \mathrm{Seq}(\overline{c}) \land \forall n \in \mathbb{N}\,(\overline{c}(n) = c)) \to c \in \mathbb{R} \to \varepsilon \in \mathbb{R} \to \varphi_{3}$
\item $\varphi_{45} :\equiv (\forall n \in \mathbb{N}\,(\overline{c}(n) = c) \to n \in \mathbb{N} \to \overline{c}(n) = c) \to c \in \mathbb{R} \to \varphi_{18}$
\item $\varphi_{46} :\equiv (0 = \lvert \overline{c}(n) - c \rvert \to 0 < \varepsilon \to \lvert \overline{c}(n) - c \rvert < \varepsilon) \to c \in \mathbb{R} \to \varphi_{23}$
\item $\varphi_{47} :\equiv 1_{\omega} \in \mathbb{N} \to (c \in \mathbb{R} \to \varphi_{25}) \to c \in \mathbb{R} \to \varepsilon \in \mathbb{R} \to 0 < \varepsilon \to \varphi_{12}$
\item $\varphi_{48} :\equiv \forall \varepsilon\,(c \in \mathbb{R} \to \varepsilon \in \mathbb{R} \to \varphi_{24}) \to \forall \varepsilon\,(c \in \mathbb{R}) \to \forall \varepsilon \in \mathbb{R}\,(\varphi_{24})$
\item $\varphi_{49} :\equiv (c \in \mathbb{R} \to \varphi_{18}) \to (c \in \mathbb{R} \to \varepsilon \in \mathbb{R} \to 0 < \varepsilon \to n \in \mathbb{N} \to n \in \mathbb{N}) \to \varphi_{20}$
\item $\varphi_{50} :\equiv \lvert \overline{c}(n) - c \rvert = \lvert \overline{c}(n) - c \rvert \to (\overline{c}(n) = c \to \varphi_{28}) \to c \in \mathbb{R} \to \varepsilon \in \mathbb{R} \to \varphi_{9}$
\item $\varphi_{51} :\equiv (c \in \mathbb{R} \to \varepsilon \in \mathbb{R} \to 0 < \varepsilon \to \varphi_{12}) \to (\varphi_{8} \to \neg \varphi_{12}) \to c \in \mathbb{R} \to \varepsilon \in \mathbb{R} \to \varphi_{24}$
\item $\varphi_{52} :\equiv (c \in \mathbb{R} \to \varepsilon \in \mathbb{R} \to \varphi_{9}) \to (c \in \mathbb{R} \to \varphi_{4}) \to \varphi_{29} \to c \in \mathbb{R} \to \varphi_{22}$
\item $\varphi_{53} :\equiv \varphi_{48} \to (c \in \mathbb{R} \to \forall \varepsilon\,(c \in \mathbb{R})) \to c \in \mathbb{R} \to \forall \varepsilon \in \mathbb{R}\,(\varphi_{24})$
\end{abbrevlist}
\end{fullproof}

% 自動生成 (prover/build.js): CintClosed  ノード数 207
\begin{fullproof}{[G5]}{$\forall x\,\forall c\,\forall a\,\forall b\,(a \in \mathbb{R} \to b \in \mathbb{R} \to \forall n \in \mathbb{N}\,(x(n) \in [a,b]) \to \mathrm{Conv}(x,c) \to c \in [a,b])$}{207}\phantomsection\label{tree:CintClosed}
\noindent$\Pi_{1}$:
\begin{ptree}
\AXC{$\forall a\,\forall b\,\forall x\,(x \in [a,b] \leftrightarrow x \in \mathbb{R} \land a \leq x \land x \leq b)$\lbl{Def{$[a,b]$}}}
\AXC{$\varphi_{54}$\lbl{A2}}
\BIC{$\forall b\,\forall x\,(x \in [a,b] \leftrightarrow x \in \mathbb{R} \land a \leq x \land x \leq b)$}
\AXC{$\varphi_{49}$\lbl{A2}}
\BIC{$\forall x\,(x \in [a,b] \leftrightarrow x \in \mathbb{R} \land a \leq x \land x \leq b)$}
\end{ptree}
\noindent$\Pi_{2}$:
\begin{ptree}
\AXC{$\stackrel{\Pi_{1}}{\vdots}$}\noLine
\UIC{$\forall x\,(x \in [a,b] \leftrightarrow x \in \mathbb{R} \land a \leq x \land x \leq b)$}
\AXC{$\forall x\,(x \in [a,b] \leftrightarrow x \in \mathbb{R} \land a \leq x \land x \leq b) \to \varphi_{18}$\lbl{A2}}
\BIC{$\varphi_{18}$}
\end{ptree}
\noindent$\Pi_{3}$:
\begin{ptree}
\AXC{$\varphi_{22}$\lbl{A2}}
\AXC{$a \in \mathbb{R} \to b \in \mathbb{R} \to \varphi_{29}$\lbl{A1}}
\AXC{$(a \in \mathbb{R} \to b \in \mathbb{R} \to \varphi_{29}) \to \varphi_{22} \to a \in \mathbb{R} \to b \in \mathbb{R} \to \varphi_{22}$\lbl{A1}}
\BIC{$\varphi_{22} \to a \in \mathbb{R} \to b \in \mathbb{R} \to \varphi_{22}$}
\BIC{$a \in \mathbb{R} \to b \in \mathbb{R} \to \varphi_{22}$}
\AXC{$\varphi_{85}$\lbl{A1}}
\BIC{$(a \in \mathbb{R} \to b \in \mathbb{R} \to \varphi_{21}) \to \varphi_{18} \to a \in \mathbb{R} \to b \in \mathbb{R} \to \varphi_{23}$}
\end{ptree}
\noindent$\Pi_{4}$:
\begin{ptree}
\AXC{$\stackrel{\Pi_{2}}{\vdots}$}\noLine
\UIC{$\varphi_{18}$}
\AXC{$a \in \mathbb{R} \to b \in \mathbb{R} \to \varphi_{21}$\lbl{A1}}
\AXC{$\stackrel{\Pi_{3}}{\vdots}$}\noLine
\UIC{$(a \in \mathbb{R} \to b \in \mathbb{R} \to \varphi_{21}) \to \varphi_{18} \to a \in \mathbb{R} \to b \in \mathbb{R} \to \varphi_{23}$}
\BIC{$\varphi_{18} \to a \in \mathbb{R} \to b \in \mathbb{R} \to \varphi_{23}$}
\BIC{$a \in \mathbb{R} \to b \in \mathbb{R} \to \varphi_{23}$}
\end{ptree}
\noindent$\Pi_{5}$:
\begin{ptree}
\AXC{$x = x$\lbl{A5}}
\AXC{$x = y \to x = x \to y = x$\lbl{A6}}
\AXC{$(x = y \to x = x \to y = x) \to x = x \to x = y \to y = x$\lbl{A1}}
\BIC{$x = x \to x = y \to y = x$}
\BIC{$x = y \to y = x$}
\UIC{$\forall y\,(x = y \to y = x)$}
\UIC{$\forall x\,\forall y\,(x = y \to y = x)$}
\end{ptree}
\noindent$\Pi_{6}$:
\begin{ptree}
\AXC{$a \in \mathbb{R} \to b \in \mathbb{R} \to b \in \mathbb{R}$\lbl{A1}}
\AXC{$\forall c \in \mathbb{R}\,(\mathrm{Seq}(\overline{c}) \land \forall n \in \mathbb{N}\,(\overline{c}(n) = c))$\lbl{[S10]}}
\AXC{$\varphi_{57}$\lbl{A2}}
\BIC{$b \in \mathbb{R} \to \mathrm{Seq}(\overline{b}) \land \forall n \in \mathbb{N}\,(\overline{b}(n) = b)$}
\AXC{$\varphi_{69}$\lbl{A1}}
\BIC{$\varphi_{61}$}
\BIC{$a \in \mathbb{R} \to b \in \mathbb{R} \to \forall n \in \mathbb{N}\,(x(n) \in [a,b]) \to \varphi_{7}$}
\end{ptree}
\noindent$\Pi_{7}$:
\begin{ptree}
\AXC{$\stackrel{\Pi_{6}}{\vdots}$}\noLine
\UIC{$a \in \mathbb{R} \to b \in \mathbb{R} \to \forall n \in \mathbb{N}\,(x(n) \in [a,b]) \to \varphi_{7}$}
\AXC{$\varphi_{74}$\lbl{A1}}
\BIC{$a \in \mathbb{R} \to b \in \mathbb{R} \to \varphi_{52}$}
\AXC{$(a \in \mathbb{R} \to b \in \mathbb{R} \to \varphi_{52}) \to \varphi_{65}$\lbl{A1}}
\BIC{$\varphi_{65}$}
\end{ptree}
\noindent$\Pi_{8}$:
\begin{ptree}
\AXC{$a \in \mathbb{R} \to b \in \mathbb{R} \to \varphi_{21}$\lbl{A1}}
\AXC{$\forall n \in \mathbb{N}\,(\overline{b}(n) = b) \to n \in \mathbb{N} \to \overline{b}(n) = b$\lbl{A2}}
\AXC{$\stackrel{\Pi_{7}}{\vdots}$}\noLine
\UIC{$\varphi_{65}$}
\BIC{$a \in \mathbb{R} \to b \in \mathbb{R} \to \varphi_{50}$}
\AXC{$\varphi_{81}$\lbl{A1}}
\BIC{$(a \in \mathbb{R} \to b \in \mathbb{R} \to \varphi_{21}) \to a \in \mathbb{R} \to b \in \mathbb{R} \to \varphi_{33}$}
\BIC{$a \in \mathbb{R} \to b \in \mathbb{R} \to \varphi_{33}$}
\AXC{$\varphi_{76}$\lbl{A1}}
\BIC{$(\overline{b}(n) = b \to b = \overline{b}(n)) \to a \in \mathbb{R} \to b \in \mathbb{R} \to \varphi_{34}$}
\end{ptree}
\noindent$\Pi_{9}$:
\begin{ptree}
\AXC{$\stackrel{\Pi_{5}}{\vdots}$}\noLine
\UIC{$\forall x\,\forall y\,(x = y \to y = x)$}
\AXC{$\varphi_{4}$\lbl{A2}}
\BIC{$\forall y\,(\overline{b}(n) = y \to y = \overline{b}(n))$}
\AXC{$\varphi_{26}$\lbl{A2}}
\BIC{$\overline{b}(n) = b \to b = \overline{b}(n)$}
\AXC{$\stackrel{\Pi_{8}}{\vdots}$}\noLine
\UIC{$(\overline{b}(n) = b \to b = \overline{b}(n)) \to a \in \mathbb{R} \to b \in \mathbb{R} \to \varphi_{34}$}
\BIC{$a \in \mathbb{R} \to b \in \mathbb{R} \to \varphi_{34}$}
\end{ptree}
\noindent$\Pi_{10}$:
\begin{ptree}
\AXC{$\stackrel{\Pi_{4}}{\vdots}$}\noLine
\UIC{$a \in \mathbb{R} \to b \in \mathbb{R} \to \varphi_{23}$}
\AXC{$\stackrel{\Pi_{9}}{\vdots}$}\noLine
\UIC{$a \in \mathbb{R} \to b \in \mathbb{R} \to \varphi_{34}$}
\AXC{$(a \in \mathbb{R} \to b \in \mathbb{R} \to \varphi_{34}) \to \varphi_{83}$\lbl{A1}}
\BIC{$\varphi_{83}$}
\BIC{$(b = \overline{b}(n) \to x(n) < b \to x(n) < \overline{b}(n)) \to \varphi_{63}$}
\end{ptree}
\noindent$\Pi_{11}$:
\begin{ptree}
\AXC{$z = w \to x(n) < z \to x(n) < w$\lbl{A6}}
\UIC{$\forall w\,(z = w \to x(n) < z \to x(n) < w)$}
\UIC{$\forall z\,\forall w\,(z = w \to x(n) < z \to x(n) < w)$}
\AXC{$\varphi_{37}$\lbl{A2}}
\BIC{$\forall w\,(b = w \to x(n) < b \to x(n) < w)$}
\AXC{$\varphi_{55}$\lbl{A2}}
\BIC{$b = \overline{b}(n) \to x(n) < b \to x(n) < \overline{b}(n)$}
\AXC{$\stackrel{\Pi_{10}}{\vdots}$}\noLine
\UIC{$(b = \overline{b}(n) \to x(n) < b \to x(n) < \overline{b}(n)) \to \varphi_{63}$}
\BIC{$\varphi_{63}$}
\end{ptree}
\noindent$\Pi_{12}$:
\begin{ptree}
\AXC{$\varphi_{72}$\lbl{A3}}
\AXC{$z = w \to x(n) = z \to x(n) = w$\lbl{A6}}
\UIC{$\forall w\,(z = w \to x(n) = z \to x(n) = w)$}
\UIC{$\forall z\,\forall w\,(z = w \to x(n) = z \to x(n) = w)$}
\AXC{$\varphi_{19}$\lbl{A2}}
\BIC{$\forall w\,(b = w \to x(n) = b \to x(n) = w)$}
\AXC{$\varphi_{45}$\lbl{A2}}
\BIC{$b = \overline{b}(n) \to x(n) = b \to x(n) = \overline{b}(n)$}
\AXC{$\stackrel{\Pi_{11}}{\vdots}$}\noLine
\UIC{$\varphi_{63}$}
\BIC{$a \in \mathbb{R} \to b \in \mathbb{R} \to \varphi_{42}$}
\AXC{$\varphi_{70}$\lbl{A1}}
\BIC{$\varphi_{25} \to n \in \mathbb{N} \to x(n) \leq \overline{b}(n)$}
\UIC{$\forall n\,(\varphi_{25} \to n \in \mathbb{N} \to x(n) \leq \overline{b}(n))$}
\AXC{$\varphi_{88}$\lbl{A1}}
\BIC{$\varphi_{72} \to (\varphi_{25} \to \forall n\,(\varphi_{25})) \to a \in \mathbb{R} \to b \in \mathbb{R} \to \varphi_{47}$}
\BIC{$(\varphi_{25} \to \forall n\,(\varphi_{25})) \to a \in \mathbb{R} \to b \in \mathbb{R} \to \varphi_{47}$}
\end{ptree}
\noindent$\Pi_{13}$:
\begin{ptree}
\AXC{$\varphi_{30}$\lbl{定理2.5}}
\AXC{$\varphi_{30} \to \varphi_{28}$\lbl{A2}}
\BIC{$\varphi_{28}$}
\AXC{$\varphi_{28} \to \varphi_{39}$\lbl{A2}}
\BIC{$\varphi_{39}$}
\AXC{$\varphi_{39} \to \varphi_{31}$\lbl{A2}}
\BIC{$\varphi_{31}$}
\AXC{$\varphi_{31} \to \mathrm{Conv}(x,c) \to \varphi_{13}$\lbl{A2}}
\BIC{$\mathrm{Conv}(x,c) \to \varphi_{13}$}
\end{ptree}
\noindent$\Pi_{14}$:
\begin{ptree}
\AXC{$a \in \mathbb{R} \to b \in \mathbb{R} \to \varphi_{9}$\lbl{A1}}
\AXC{$\stackrel{\Pi_{13}}{\vdots}$}\noLine
\UIC{$\mathrm{Conv}(x,c) \to \varphi_{13}$}
\AXC{$(\mathrm{Conv}(x,c) \to \varphi_{13}) \to (a \in \mathbb{R} \to b \in \mathbb{R} \to \varphi_{9}) \to \varphi_{80}$\lbl{A1}}
\BIC{$(a \in \mathbb{R} \to b \in \mathbb{R} \to \varphi_{9}) \to \varphi_{80}$}
\BIC{$\varphi_{80}$}
\end{ptree}
\noindent$\Pi_{15}$:
\begin{ptree}
\AXC{$a \in \mathbb{R} \to b \in \mathbb{R} \to b \in \mathbb{R}$\lbl{A1}}
\AXC{$\forall c \in \mathbb{R}\,(\mathrm{Conv}(\overline{c},c))$\lbl{[G4]}}
\AXC{$\varphi_{10}$\lbl{A2}}
\BIC{$b \in \mathbb{R} \to \mathrm{Conv}(\overline{b},b)$}
\AXC{$\varphi_{62}$\lbl{A1}}
\BIC{$(a \in \mathbb{R} \to b \in \mathbb{R} \to b \in \mathbb{R}) \to a \in \mathbb{R} \to b \in \mathbb{R} \to \varphi_{15}$}
\BIC{$a \in \mathbb{R} \to b \in \mathbb{R} \to \varphi_{15}$}
\AXC{$\stackrel{\Pi_{14}}{\vdots}$}\noLine
\UIC{$\varphi_{80}$}
\BIC{$(a \in \mathbb{R} \to b \in \mathbb{R} \to \varphi_{47}) \to a \in \mathbb{R} \to b \in \mathbb{R} \to \varphi_{2}$}
\end{ptree}
\noindent$\Pi_{16}$:
\begin{ptree}
\AXC{$\varphi_{25} \to \forall n\,(\varphi_{25})$\lbl{A4}}
\AXC{$\stackrel{\Pi_{12}}{\vdots}$}\noLine
\UIC{$(\varphi_{25} \to \forall n\,(\varphi_{25})) \to a \in \mathbb{R} \to b \in \mathbb{R} \to \varphi_{47}$}
\BIC{$a \in \mathbb{R} \to b \in \mathbb{R} \to \varphi_{47}$}
\AXC{$\stackrel{\Pi_{15}}{\vdots}$}\noLine
\UIC{$(a \in \mathbb{R} \to b \in \mathbb{R} \to \varphi_{47}) \to a \in \mathbb{R} \to b \in \mathbb{R} \to \varphi_{2}$}
\BIC{$a \in \mathbb{R} \to b \in \mathbb{R} \to \varphi_{2}$}
\end{ptree}
\noindent$\Pi_{17}$:
\begin{ptree}
\AXC{$\forall a\,\forall b\,\forall x\,(x \in [a,b] \leftrightarrow x \in \mathbb{R} \land a \leq x \land x \leq b)$\lbl{Def{$[a,b]$}}}
\AXC{$\varphi_{54}$\lbl{A2}}
\BIC{$\forall b\,\forall x\,(x \in [a,b] \leftrightarrow x \in \mathbb{R} \land a \leq x \land x \leq b)$}
\AXC{$\varphi_{49}$\lbl{A2}}
\BIC{$\forall x\,(x \in [a,b] \leftrightarrow x \in \mathbb{R} \land a \leq x \land x \leq b)$}
\end{ptree}
\noindent$\Pi_{18}$:
\begin{ptree}
\AXC{$\stackrel{\Pi_{17}}{\vdots}$}\noLine
\UIC{$\forall x\,(x \in [a,b] \leftrightarrow x \in \mathbb{R} \land a \leq x \land x \leq b)$}
\AXC{$\forall x\,(x \in [a,b] \leftrightarrow x \in \mathbb{R} \land a \leq x \land x \leq b) \to \varphi_{18}$\lbl{A2}}
\BIC{$\varphi_{18}$}
\end{ptree}
\noindent$\Pi_{19}$:
\begin{ptree}
\AXC{$\varphi_{22}$\lbl{A2}}
\AXC{$a \in \mathbb{R} \to b \in \mathbb{R} \to \varphi_{29}$\lbl{A1}}
\AXC{$(a \in \mathbb{R} \to b \in \mathbb{R} \to \varphi_{29}) \to \varphi_{22} \to a \in \mathbb{R} \to b \in \mathbb{R} \to \varphi_{22}$\lbl{A1}}
\BIC{$\varphi_{22} \to a \in \mathbb{R} \to b \in \mathbb{R} \to \varphi_{22}$}
\BIC{$a \in \mathbb{R} \to b \in \mathbb{R} \to \varphi_{22}$}
\AXC{$\varphi_{86}$\lbl{A1}}
\BIC{$(a \in \mathbb{R} \to b \in \mathbb{R} \to \varphi_{21}) \to \varphi_{18} \to a \in \mathbb{R} \to b \in \mathbb{R} \to \varphi_{24}$}
\end{ptree}
\noindent$\Pi_{20}$:
\begin{ptree}
\AXC{$\stackrel{\Pi_{18}}{\vdots}$}\noLine
\UIC{$\varphi_{18}$}
\AXC{$a \in \mathbb{R} \to b \in \mathbb{R} \to \varphi_{21}$\lbl{A1}}
\AXC{$\stackrel{\Pi_{19}}{\vdots}$}\noLine
\UIC{$(a \in \mathbb{R} \to b \in \mathbb{R} \to \varphi_{21}) \to \varphi_{18} \to a \in \mathbb{R} \to b \in \mathbb{R} \to \varphi_{24}$}
\BIC{$\varphi_{18} \to a \in \mathbb{R} \to b \in \mathbb{R} \to \varphi_{24}$}
\BIC{$a \in \mathbb{R} \to b \in \mathbb{R} \to \varphi_{24}$}
\end{ptree}
\noindent$\Pi_{21}$:
\begin{ptree}
\AXC{$a \in \mathbb{R} \to a \in \mathbb{R}$\lbl{A1}}
\AXC{$\forall c \in \mathbb{R}\,(\mathrm{Seq}(\overline{c}) \land \forall n \in \mathbb{N}\,(\overline{c}(n) = c))$\lbl{[S10]}}
\AXC{$\varphi_{58}$\lbl{A2}}
\BIC{$a \in \mathbb{R} \to \mathrm{Seq}(\overline{a}) \land \forall n \in \mathbb{N}\,(\overline{a}(n) = a)$}
\AXC{$\varphi_{68}$\lbl{A1}}
\BIC{$\varphi_{59}$}
\BIC{$a \in \mathbb{R} \to b \in \mathbb{R} \to \forall n \in \mathbb{N}\,(x(n) \in [a,b]) \to \varphi_{8}$}
\AXC{$\varphi_{75}$\lbl{A1}}
\BIC{$a \in \mathbb{R} \to b \in \mathbb{R} \to \varphi_{53}$}
\end{ptree}
\noindent$\Pi_{22}$:
\begin{ptree}
\AXC{$\forall n \in \mathbb{N}\,(\overline{a}(n) = a) \to n \in \mathbb{N} \to \overline{a}(n) = a$\lbl{A2}}
\AXC{$\stackrel{\Pi_{21}}{\vdots}$}\noLine
\UIC{$a \in \mathbb{R} \to b \in \mathbb{R} \to \varphi_{53}$}
\AXC{$(a \in \mathbb{R} \to b \in \mathbb{R} \to \varphi_{53}) \to \varphi_{66}$\lbl{A1}}
\BIC{$\varphi_{66}$}
\BIC{$a \in \mathbb{R} \to b \in \mathbb{R} \to \varphi_{51}$}
\end{ptree}
\noindent$\Pi_{23}$:
\begin{ptree}
\AXC{$a \in \mathbb{R} \to b \in \mathbb{R} \to \varphi_{21}$\lbl{A1}}
\AXC{$\stackrel{\Pi_{22}}{\vdots}$}\noLine
\UIC{$a \in \mathbb{R} \to b \in \mathbb{R} \to \varphi_{51}$}
\AXC{$\varphi_{82}$\lbl{A1}}
\BIC{$(a \in \mathbb{R} \to b \in \mathbb{R} \to \varphi_{21}) \to a \in \mathbb{R} \to b \in \mathbb{R} \to \varphi_{35}$}
\BIC{$a \in \mathbb{R} \to b \in \mathbb{R} \to \varphi_{35}$}
\AXC{$\varphi_{77}$\lbl{A1}}
\BIC{$(\overline{a}(n) = a \to a = \overline{a}(n)) \to a \in \mathbb{R} \to b \in \mathbb{R} \to \varphi_{36}$}
\end{ptree}
\noindent$\Pi_{24}$:
\begin{ptree}
\AXC{$\stackrel{\Pi_{5}}{\vdots}$}\noLine
\UIC{$\forall x\,\forall y\,(x = y \to y = x)$}
\AXC{$\varphi_{5}$\lbl{A2}}
\BIC{$\forall y\,(\overline{a}(n) = y \to y = \overline{a}(n))$}
\AXC{$\varphi_{27}$\lbl{A2}}
\BIC{$\overline{a}(n) = a \to a = \overline{a}(n)$}
\AXC{$\stackrel{\Pi_{23}}{\vdots}$}\noLine
\UIC{$(\overline{a}(n) = a \to a = \overline{a}(n)) \to a \in \mathbb{R} \to b \in \mathbb{R} \to \varphi_{36}$}
\BIC{$a \in \mathbb{R} \to b \in \mathbb{R} \to \varphi_{36}$}
\end{ptree}
\noindent$\Pi_{25}$:
\begin{ptree}
\AXC{$\stackrel{\Pi_{20}}{\vdots}$}\noLine
\UIC{$a \in \mathbb{R} \to b \in \mathbb{R} \to \varphi_{24}$}
\AXC{$\stackrel{\Pi_{24}}{\vdots}$}\noLine
\UIC{$a \in \mathbb{R} \to b \in \mathbb{R} \to \varphi_{36}$}
\AXC{$(a \in \mathbb{R} \to b \in \mathbb{R} \to \varphi_{36}) \to \varphi_{84}$\lbl{A1}}
\BIC{$\varphi_{84}$}
\BIC{$(a = \overline{a}(n) \to a < x(n) \to \overline{a}(n) < x(n)) \to \varphi_{64}$}
\end{ptree}
\noindent$\Pi_{26}$:
\begin{ptree}
\AXC{$z = w \to z < x(n) \to w < x(n)$\lbl{A6}}
\UIC{$\forall w\,(z = w \to z < x(n) \to w < x(n))$}
\UIC{$\forall z\,\forall w\,(z = w \to z < x(n) \to w < x(n))$}
\AXC{$\varphi_{38}$\lbl{A2}}
\BIC{$\forall w\,(a = w \to a < x(n) \to w < x(n))$}
\AXC{$\varphi_{56}$\lbl{A2}}
\BIC{$a = \overline{a}(n) \to a < x(n) \to \overline{a}(n) < x(n)$}
\AXC{$\stackrel{\Pi_{25}}{\vdots}$}\noLine
\UIC{$(a = \overline{a}(n) \to a < x(n) \to \overline{a}(n) < x(n)) \to \varphi_{64}$}
\BIC{$\varphi_{64}$}
\end{ptree}
\noindent$\Pi_{27}$:
\begin{ptree}
\AXC{$\varphi_{73}$\lbl{A3}}
\AXC{$z = w \to z = x(n) \to w = x(n)$\lbl{A6}}
\UIC{$\forall w\,(z = w \to z = x(n) \to w = x(n))$}
\UIC{$\forall z\,\forall w\,(z = w \to z = x(n) \to w = x(n))$}
\AXC{$\varphi_{20}$\lbl{A2}}
\BIC{$\forall w\,(a = w \to a = x(n) \to w = x(n))$}
\AXC{$\varphi_{46}$\lbl{A2}}
\BIC{$a = \overline{a}(n) \to a = x(n) \to \overline{a}(n) = x(n)$}
\AXC{$\stackrel{\Pi_{26}}{\vdots}$}\noLine
\UIC{$\varphi_{64}$}
\BIC{$a \in \mathbb{R} \to b \in \mathbb{R} \to \varphi_{43}$}
\AXC{$\varphi_{71}$\lbl{A1}}
\BIC{$\varphi_{25} \to n \in \mathbb{N} \to \overline{a}(n) \leq x(n)$}
\UIC{$\forall n\,(\varphi_{25} \to n \in \mathbb{N} \to \overline{a}(n) \leq x(n))$}
\AXC{$\varphi_{89}$\lbl{A1}}
\BIC{$\varphi_{73} \to (\varphi_{25} \to \forall n\,(\varphi_{25})) \to a \in \mathbb{R} \to b \in \mathbb{R} \to \varphi_{48}$}
\BIC{$(\varphi_{25} \to \forall n\,(\varphi_{25})) \to a \in \mathbb{R} \to b \in \mathbb{R} \to \varphi_{48}$}
\end{ptree}
\noindent$\Pi_{28}$:
\begin{ptree}
\AXC{$a \in \mathbb{R} \to a \in \mathbb{R}$\lbl{A1}}
\AXC{$\forall c \in \mathbb{R}\,(\mathrm{Conv}(\overline{c},c))$\lbl{[G4]}}
\AXC{$\varphi_{11}$\lbl{A2}}
\BIC{$a \in \mathbb{R} \to \mathrm{Conv}(\overline{a},a)$}
\AXC{$\varphi_{60}$\lbl{A1}}
\BIC{$(a \in \mathbb{R} \to a \in \mathbb{R}) \to a \in \mathbb{R} \to b \in \mathbb{R} \to \varphi_{16}$}
\BIC{$a \in \mathbb{R} \to b \in \mathbb{R} \to \varphi_{16}$}
\end{ptree}
\noindent$\Pi_{29}$:
\begin{ptree}
\AXC{$\varphi_{30}$\lbl{定理2.5}}
\AXC{$\varphi_{30} \to \varphi_{44}$\lbl{A2}}
\BIC{$\varphi_{44}$}
\AXC{$\varphi_{44} \to \varphi_{40}$\lbl{A2}}
\BIC{$\varphi_{40}$}
\AXC{$\varphi_{40} \to \varphi_{32}$\lbl{A2}}
\BIC{$\varphi_{32}$}
\AXC{$\varphi_{32} \to \mathrm{Conv}(\overline{a},a) \to \varphi_{6}$\lbl{A2}}
\BIC{$\mathrm{Conv}(\overline{a},a) \to \varphi_{6}$}
\end{ptree}
\noindent$\Pi_{30}$:
\begin{ptree}
\AXC{$\stackrel{\Pi_{28}}{\vdots}$}\noLine
\UIC{$a \in \mathbb{R} \to b \in \mathbb{R} \to \varphi_{16}$}
\AXC{$\stackrel{\Pi_{29}}{\vdots}$}\noLine
\UIC{$\mathrm{Conv}(\overline{a},a) \to \varphi_{6}$}
\AXC{$(\mathrm{Conv}(\overline{a},a) \to \varphi_{6}) \to (a \in \mathbb{R} \to b \in \mathbb{R} \to \varphi_{16}) \to \varphi_{79}$\lbl{A1}}
\BIC{$(a \in \mathbb{R} \to b \in \mathbb{R} \to \varphi_{16}) \to \varphi_{79}$}
\BIC{$\varphi_{79}$}
\end{ptree}
\noindent$\Pi_{31}$:
\begin{ptree}
\AXC{$\varphi_{25} \to \forall n\,(\varphi_{25})$\lbl{A4}}
\AXC{$\stackrel{\Pi_{27}}{\vdots}$}\noLine
\UIC{$(\varphi_{25} \to \forall n\,(\varphi_{25})) \to a \in \mathbb{R} \to b \in \mathbb{R} \to \varphi_{48}$}
\BIC{$a \in \mathbb{R} \to b \in \mathbb{R} \to \varphi_{48}$}
\AXC{$a \in \mathbb{R} \to b \in \mathbb{R} \to \varphi_{9}$\lbl{A1}}
\AXC{$\stackrel{\Pi_{30}}{\vdots}$}\noLine
\UIC{$\varphi_{79}$}
\BIC{$(a \in \mathbb{R} \to b \in \mathbb{R} \to \varphi_{48}) \to a \in \mathbb{R} \to b \in \mathbb{R} \to \varphi_{3}$}
\BIC{$a \in \mathbb{R} \to b \in \mathbb{R} \to \varphi_{3}$}
\end{ptree}
\noindent$\Pi_{32}$:
\begin{ptree}
\AXC{$a \in \mathbb{R} \to b \in \mathbb{R} \to \varphi_{9}$\lbl{A1}}
\AXC{$\forall a\,\forall \alpha\,(\mathrm{Conv}(a,\alpha) \to \alpha \in \mathbb{R})$\lbl{[C2]}}
\AXC{$\varphi_{12}$\lbl{A2}}
\BIC{$\forall \alpha\,(\mathrm{Conv}(x,\alpha) \to \alpha \in \mathbb{R})$}
\AXC{$\varphi_{1}$\lbl{A2}}
\BIC{$\mathrm{Conv}(x,c) \to c \in \mathbb{R}$}
\AXC{$\varphi_{67}$\lbl{A1}}
\BIC{$(a \in \mathbb{R} \to b \in \mathbb{R} \to \varphi_{9}) \to a \in \mathbb{R} \to b \in \mathbb{R} \to \varphi_{14}$}
\BIC{$a \in \mathbb{R} \to b \in \mathbb{R} \to \varphi_{14}$}
\end{ptree}
\noindent$\Pi_{33}$:
\begin{ptree}
\AXC{$\forall a\,\forall b\,\forall x\,(x \in [a,b] \leftrightarrow x \in \mathbb{R} \land a \leq x \land x \leq b)$\lbl{Def{$[a,b]$}}}
\AXC{$\varphi_{54}$\lbl{A2}}
\BIC{$\forall b\,\forall x\,(x \in [a,b] \leftrightarrow x \in \mathbb{R} \land a \leq x \land x \leq b)$}
\AXC{$\varphi_{49}$\lbl{A2}}
\BIC{$\forall x\,(x \in [a,b] \leftrightarrow x \in \mathbb{R} \land a \leq x \land x \leq b)$}
\AXC{$\varphi_{41}$\lbl{A2}}
\BIC{$c \in [a,b] \leftrightarrow c \in \mathbb{R} \land a \leq c \land c \leq b$}
\AXC{$\varphi_{87}$\lbl{A1}}
\BIC{$(a \in \mathbb{R} \to b \in \mathbb{R} \to \varphi_{14}) \to \varphi_{78}$}
\end{ptree}
\noindent 主木:
\begin{ptree}
\AXC{$\stackrel{\Pi_{16}}{\vdots}$}\noLine
\UIC{$a \in \mathbb{R} \to b \in \mathbb{R} \to \varphi_{2}$}
\AXC{$\stackrel{\Pi_{31}}{\vdots}$}\noLine
\UIC{$a \in \mathbb{R} \to b \in \mathbb{R} \to \varphi_{3}$}
\AXC{$\stackrel{\Pi_{32}}{\vdots}$}\noLine
\UIC{$a \in \mathbb{R} \to b \in \mathbb{R} \to \varphi_{14}$}
\AXC{$\stackrel{\Pi_{33}}{\vdots}$}\noLine
\UIC{$(a \in \mathbb{R} \to b \in \mathbb{R} \to \varphi_{14}) \to \varphi_{78}$}
\BIC{$\varphi_{78}$}
\BIC{$(a \in \mathbb{R} \to b \in \mathbb{R} \to \varphi_{2}) \to a \in \mathbb{R} \to b \in \mathbb{R} \to \varphi_{17}$}
\BIC{$a \in \mathbb{R} \to b \in \mathbb{R} \to \varphi_{17}$}
\UIC{$\forall b\,(a \in \mathbb{R} \to b \in \mathbb{R} \to \varphi_{17})$}
\UIC{$\forall a\,\forall b\,(a \in \mathbb{R} \to b \in \mathbb{R} \to \varphi_{17})$}
\UIC{$\forall c\,\forall a\,\forall b\,(a \in \mathbb{R} \to b \in \mathbb{R} \to \varphi_{17})$}
\UIC{$\forall x\,\forall c\,\forall a\,\forall b\,(a \in \mathbb{R} \to b \in \mathbb{R} \to \varphi_{17})$}
\end{ptree}
\begin{abbrevlist}
\item $\varphi_{1} :\equiv \forall \alpha\,(\mathrm{Conv}(x,\alpha) \to \alpha \in \mathbb{R}) \to \mathrm{Conv}(x,c) \to c \in \mathbb{R}$
\item $\varphi_{2} :\equiv \forall n \in \mathbb{N}\,(x(n) \in [a,b]) \to \mathrm{Conv}(x,c) \to c \leq b$
\item $\varphi_{3} :\equiv \forall n \in \mathbb{N}\,(x(n) \in [a,b]) \to \mathrm{Conv}(x,c) \to a \leq c$
\item $\varphi_{4} :\equiv \forall x\,\forall y\,(x = y \to y = x) \to \forall y\,(\overline{b}(n) = y \to y = \overline{b}(n))$
\item $\varphi_{5} :\equiv \forall x\,\forall y\,(x = y \to y = x) \to \forall y\,(\overline{a}(n) = y \to y = \overline{a}(n))$
\item $\varphi_{6} :\equiv \mathrm{Conv}(x,c) \to \forall n \in \mathbb{N}\,(\overline{a}(n) \leq x(n)) \to a \leq c$
\item $\varphi_{7} :\equiv \mathrm{Conv}(x,c) \to \mathrm{Seq}(\overline{b}) \land \forall n \in \mathbb{N}\,(\overline{b}(n) = b)$
\item $\varphi_{8} :\equiv \mathrm{Conv}(x,c) \to \mathrm{Seq}(\overline{a}) \land \forall n \in \mathbb{N}\,(\overline{a}(n) = a)$
\item $\varphi_{9} :\equiv \forall n \in \mathbb{N}\,(x(n) \in [a,b]) \to \mathrm{Conv}(x,c) \to \mathrm{Conv}(x,c)$
\item $\varphi_{10} :\equiv \forall c \in \mathbb{R}\,(\mathrm{Conv}(\overline{c},c)) \to b \in \mathbb{R} \to \mathrm{Conv}(\overline{b},b)$
\item $\varphi_{11} :\equiv \forall c \in \mathbb{R}\,(\mathrm{Conv}(\overline{c},c)) \to a \in \mathbb{R} \to \mathrm{Conv}(\overline{a},a)$
\item $\varphi_{12} :\equiv \forall a\,\forall \alpha\,(\mathrm{Conv}(a,\alpha) \to \alpha \in \mathbb{R}) \to \forall \alpha\,(\mathrm{Conv}(x,\alpha) \to \alpha \in \mathbb{R})$
\item $\varphi_{13} :\equiv \mathrm{Conv}(\overline{b},b) \to \forall n \in \mathbb{N}\,(x(n) \leq \overline{b}(n)) \to c \leq b$
\item $\varphi_{14} :\equiv \forall n \in \mathbb{N}\,(x(n) \in [a,b]) \to \mathrm{Conv}(x,c) \to c \in \mathbb{R}$
\item $\varphi_{15} :\equiv \forall n \in \mathbb{N}\,(x(n) \in [a,b]) \to \mathrm{Conv}(x,c) \to \mathrm{Conv}(\overline{b},b)$
\item $\varphi_{16} :\equiv \forall n \in \mathbb{N}\,(x(n) \in [a,b]) \to \mathrm{Conv}(x,c) \to \mathrm{Conv}(\overline{a},a)$
\item $\varphi_{17} :\equiv \forall n \in \mathbb{N}\,(x(n) \in [a,b]) \to \mathrm{Conv}(x,c) \to c \in [a,b]$
\item $\varphi_{18} :\equiv x(n) \in [a,b] \leftrightarrow x(n) \in \mathbb{R} \land a \leq x(n) \land x(n) \leq b$
\item $\varphi_{19} :\equiv \forall z\,\forall w\,(z = w \to x(n) = z \to x(n) = w) \to \forall w\,(b = w \to x(n) = b \to x(n) = w)$
\item $\varphi_{20} :\equiv \forall z\,\forall w\,(z = w \to z = x(n) \to w = x(n)) \to \forall w\,(a = w \to a = x(n) \to w = x(n))$
\item $\varphi_{21} :\equiv \forall n \in \mathbb{N}\,(x(n) \in [a,b]) \to \mathrm{Conv}(x,c) \to n \in \mathbb{N} \to n \in \mathbb{N}$
\item $\varphi_{22} :\equiv \forall n \in \mathbb{N}\,(x(n) \in [a,b]) \to n \in \mathbb{N} \to x(n) \in [a,b]$
\item $\varphi_{23} :\equiv \forall n \in \mathbb{N}\,(x(n) \in [a,b]) \to \mathrm{Conv}(x,c) \to n \in \mathbb{N} \to x(n) \leq b$
\item $\varphi_{24} :\equiv \forall n \in \mathbb{N}\,(x(n) \in [a,b]) \to \mathrm{Conv}(x,c) \to n \in \mathbb{N} \to a \leq x(n)$
\item $\varphi_{25} :\equiv a \in \mathbb{R} \land b \in \mathbb{R} \land \forall n \in \mathbb{N}\,(x(n) \in [a,b]) \land \mathrm{Conv}(x,c)$
\item $\varphi_{26} :\equiv \forall y\,(\overline{b}(n) = y \to y = \overline{b}(n)) \to \overline{b}(n) = b \to b = \overline{b}(n)$
\item $\varphi_{27} :\equiv \forall y\,(\overline{a}(n) = y \to y = \overline{a}(n)) \to \overline{a}(n) = a \to a = \overline{a}(n)$
\item $\varphi_{28} :\equiv \forall b\,\forall \alpha\,\forall \beta\,(\mathrm{Conv}(x,\alpha) \to \mathrm{Conv}(b,\beta) \to \forall n \in \mathbb{N}\,(x(n) \leq b_{n}) \to \alpha \leq \beta)$
\item $\varphi_{29} :\equiv \forall n \in \mathbb{N}\,(x(n) \in [a,b]) \to \forall n \in \mathbb{N}\,(x(n) \in [a,b])$
\item $\varphi_{30} :\equiv \forall a\,\forall b\,\forall \alpha\,\forall \beta\,(\mathrm{Conv}(a,\alpha) \to \mathrm{Conv}(b,\beta) \to \forall n \in \mathbb{N}\,(a_{n} \leq b_{n}) \to \alpha \leq \beta)$
\item $\varphi_{31} :\equiv \forall \beta\,(\mathrm{Conv}(x,c) \to \mathrm{Conv}(\overline{b},\beta) \to \forall n \in \mathbb{N}\,(x(n) \leq \overline{b}(n)) \to c \leq \beta)$
\item $\varphi_{32} :\equiv \forall \beta\,(\mathrm{Conv}(\overline{a},a) \to \mathrm{Conv}(x,\beta) \to \forall n \in \mathbb{N}\,(\overline{a}(n) \leq x(n)) \to a \leq \beta)$
\item $\varphi_{33} :\equiv \forall n \in \mathbb{N}\,(x(n) \in [a,b]) \to \mathrm{Conv}(x,c) \to n \in \mathbb{N} \to \overline{b}(n) = b$
\item $\varphi_{34} :\equiv \forall n \in \mathbb{N}\,(x(n) \in [a,b]) \to \mathrm{Conv}(x,c) \to n \in \mathbb{N} \to b = \overline{b}(n)$
\item $\varphi_{35} :\equiv \forall n \in \mathbb{N}\,(x(n) \in [a,b]) \to \mathrm{Conv}(x,c) \to n \in \mathbb{N} \to \overline{a}(n) = a$
\item $\varphi_{36} :\equiv \forall n \in \mathbb{N}\,(x(n) \in [a,b]) \to \mathrm{Conv}(x,c) \to n \in \mathbb{N} \to a = \overline{a}(n)$
\item $\varphi_{37} :\equiv \forall z\,\forall w\,(z = w \to x(n) < z \to x(n) < w) \to \forall w\,(b = w \to x(n) < b \to x(n) < w)$
\item $\varphi_{38} :\equiv \forall z\,\forall w\,(z = w \to z < x(n) \to w < x(n)) \to \forall w\,(a = w \to a < x(n) \to w < x(n))$
\item $\varphi_{39} :\equiv \forall \alpha\,\forall \beta\,(\mathrm{Conv}(x,\alpha) \to \mathrm{Conv}(\overline{b},\beta) \to \forall n \in \mathbb{N}\,(x(n) \leq \overline{b}(n)) \to \alpha \leq \beta)$
\item $\varphi_{40} :\equiv \forall \alpha\,\forall \beta\,(\mathrm{Conv}(\overline{a},\alpha) \to \mathrm{Conv}(x,\beta) \to \forall n \in \mathbb{N}\,(\overline{a}(n) \leq x(n)) \to \alpha \leq \beta)$
\item $\varphi_{41} :\equiv \forall x\,(x \in [a,b] \leftrightarrow x \in \mathbb{R} \land a \leq x \land x \leq b) \to (c \in [a,b] \leftrightarrow c \in \mathbb{R} \land a \leq c \land c \leq b)$
\item $\varphi_{42} :\equiv \forall n \in \mathbb{N}\,(x(n) \in [a,b]) \to \mathrm{Conv}(x,c) \to n \in \mathbb{N} \to x(n) \leq \overline{b}(n)$
\item $\varphi_{43} :\equiv \forall n \in \mathbb{N}\,(x(n) \in [a,b]) \to \mathrm{Conv}(x,c) \to n \in \mathbb{N} \to \overline{a}(n) \leq x(n)$
\item $\varphi_{44} :\equiv \forall b\,\forall \alpha\,\forall \beta\,(\mathrm{Conv}(\overline{a},\alpha) \to \mathrm{Conv}(b,\beta) \to \forall n \in \mathbb{N}\,(\overline{a}(n) \leq b_{n}) \to \alpha \leq \beta)$
\item $\varphi_{45} :\equiv \forall w\,(b = w \to x(n) = b \to x(n) = w) \to b = \overline{b}(n) \to x(n) = b \to x(n) = \overline{b}(n)$
\item $\varphi_{46} :\equiv \forall w\,(a = w \to a = x(n) \to w = x(n)) \to a = \overline{a}(n) \to a = x(n) \to \overline{a}(n) = x(n)$
\item $\varphi_{47} :\equiv \forall n \in \mathbb{N}\,(x(n) \in [a,b]) \to \mathrm{Conv}(x,c) \to \forall n \in \mathbb{N}\,(x(n) \leq \overline{b}(n))$
\item $\varphi_{48} :\equiv \forall n \in \mathbb{N}\,(x(n) \in [a,b]) \to \mathrm{Conv}(x,c) \to \forall n \in \mathbb{N}\,(\overline{a}(n) \leq x(n))$
\item $\varphi_{49} :\equiv \forall b\,\forall x\,(x \in [a,b] \leftrightarrow x \in \mathbb{R} \land a \leq x \land x \leq b) \to \forall x\,(x \in [a,b] \leftrightarrow x \in \mathbb{R} \land a \leq x \land x \leq b)$
\item $\varphi_{50} :\equiv \forall n \in \mathbb{N}\,(x(n) \in [a,b]) \to \mathrm{Conv}(x,c) \to n \in \mathbb{N} \to n \in \mathbb{N} \to \overline{b}(n) = b$
\item $\varphi_{51} :\equiv \forall n \in \mathbb{N}\,(x(n) \in [a,b]) \to \mathrm{Conv}(x,c) \to n \in \mathbb{N} \to n \in \mathbb{N} \to \overline{a}(n) = a$
\item $\varphi_{52} :\equiv \forall n \in \mathbb{N}\,(x(n) \in [a,b]) \to \mathrm{Conv}(x,c) \to n \in \mathbb{N} \to \forall n \in \mathbb{N}\,(\overline{b}(n) = b)$
\item $\varphi_{53} :\equiv \forall n \in \mathbb{N}\,(x(n) \in [a,b]) \to \mathrm{Conv}(x,c) \to n \in \mathbb{N} \to \forall n \in \mathbb{N}\,(\overline{a}(n) = a)$
\item $\varphi_{54} :\equiv \forall a\,\forall b\,\forall x\,(x \in [a,b] \leftrightarrow x \in \mathbb{R} \land a \leq x \land x \leq b) \to \forall b\,\forall x\,(x \in [a,b] \leftrightarrow x \in \mathbb{R} \land a \leq x \land x \leq b)$
\item $\varphi_{55} :\equiv \forall w\,(b = w \to x(n) < b \to x(n) < w) \to b = \overline{b}(n) \to x(n) < b \to x(n) < \overline{b}(n)$
\item $\varphi_{56} :\equiv \forall w\,(a = w \to a < x(n) \to w < x(n)) \to a = \overline{a}(n) \to a < x(n) \to \overline{a}(n) < x(n)$
\item $\varphi_{57} :\equiv \forall c \in \mathbb{R}\,(\mathrm{Seq}(\overline{c}) \land \forall n \in \mathbb{N}\,(\overline{c}(n) = c)) \to b \in \mathbb{R} \to \mathrm{Seq}(\overline{b}) \land \forall n \in \mathbb{N}\,(\overline{b}(n) = b)$
\item $\varphi_{58} :\equiv \forall c \in \mathbb{R}\,(\mathrm{Seq}(\overline{c}) \land \forall n \in \mathbb{N}\,(\overline{c}(n) = c)) \to a \in \mathbb{R} \to \mathrm{Seq}(\overline{a}) \land \forall n \in \mathbb{N}\,(\overline{a}(n) = a)$
\item $\varphi_{59} :\equiv (a \in \mathbb{R} \to a \in \mathbb{R}) \to a \in \mathbb{R} \to b \in \mathbb{R} \to \forall n \in \mathbb{N}\,(x(n) \in [a,b]) \to \varphi_{8}$
\item $\varphi_{60} :\equiv (a \in \mathbb{R} \to \mathrm{Conv}(\overline{a},a)) \to (a \in \mathbb{R} \to a \in \mathbb{R}) \to a \in \mathbb{R} \to b \in \mathbb{R} \to \varphi_{16}$
\item $\varphi_{61} :\equiv (a \in \mathbb{R} \to b \in \mathbb{R} \to b \in \mathbb{R}) \to a \in \mathbb{R} \to b \in \mathbb{R} \to \forall n \in \mathbb{N}\,(x(n) \in [a,b]) \to \varphi_{7}$
\item $\varphi_{62} :\equiv (b \in \mathbb{R} \to \mathrm{Conv}(\overline{b},b)) \to (a \in \mathbb{R} \to b \in \mathbb{R} \to b \in \mathbb{R}) \to a \in \mathbb{R} \to b \in \mathbb{R} \to \varphi_{15}$
\item $\varphi_{63} :\equiv (b = \overline{b}(n) \to x(n) = b \to x(n) = \overline{b}(n)) \to a \in \mathbb{R} \to b \in \mathbb{R} \to \varphi_{42}$
\item $\varphi_{64} :\equiv (a = \overline{a}(n) \to a = x(n) \to \overline{a}(n) = x(n)) \to a \in \mathbb{R} \to b \in \mathbb{R} \to \varphi_{43}$
\item $\varphi_{65} :\equiv (\forall n \in \mathbb{N}\,(\overline{b}(n) = b) \to n \in \mathbb{N} \to \overline{b}(n) = b) \to a \in \mathbb{R} \to b \in \mathbb{R} \to \varphi_{50}$
\item $\varphi_{66} :\equiv (\forall n \in \mathbb{N}\,(\overline{a}(n) = a) \to n \in \mathbb{N} \to \overline{a}(n) = a) \to a \in \mathbb{R} \to b \in \mathbb{R} \to \varphi_{51}$
\item $\varphi_{67} :\equiv (\mathrm{Conv}(x,c) \to c \in \mathbb{R}) \to (a \in \mathbb{R} \to b \in \mathbb{R} \to \varphi_{9}) \to a \in \mathbb{R} \to b \in \mathbb{R} \to \varphi_{14}$
\item $\varphi_{68} :\equiv (a \in \mathbb{R} \to \mathrm{Seq}(\overline{a}) \land \forall n \in \mathbb{N}\,(\overline{a}(n) = a)) \to \varphi_{59}$
\item $\varphi_{69} :\equiv (b \in \mathbb{R} \to \mathrm{Seq}(\overline{b}) \land \forall n \in \mathbb{N}\,(\overline{b}(n) = b)) \to \varphi_{61}$
\item $\varphi_{70} :\equiv (a \in \mathbb{R} \to b \in \mathbb{R} \to \varphi_{42}) \to \varphi_{25} \to n \in \mathbb{N} \to x(n) \leq \overline{b}(n)$
\item $\varphi_{71} :\equiv (a \in \mathbb{R} \to b \in \mathbb{R} \to \varphi_{43}) \to \varphi_{25} \to n \in \mathbb{N} \to \overline{a}(n) \leq x(n)$
\item $\varphi_{72} :\equiv \forall n\,(\varphi_{25} \to n \in \mathbb{N} \to x(n) \leq \overline{b}(n)) \to \forall n\,(\varphi_{25}) \to \forall n \in \mathbb{N}\,(x(n) \leq \overline{b}(n))$
\item $\varphi_{73} :\equiv \forall n\,(\varphi_{25} \to n \in \mathbb{N} \to \overline{a}(n) \leq x(n)) \to \forall n\,(\varphi_{25}) \to \forall n \in \mathbb{N}\,(\overline{a}(n) \leq x(n))$
\item $\varphi_{74} :\equiv (a \in \mathbb{R} \to b \in \mathbb{R} \to \forall n \in \mathbb{N}\,(x(n) \in [a,b]) \to \varphi_{7}) \to a \in \mathbb{R} \to b \in \mathbb{R} \to \varphi_{52}$
\item $\varphi_{75} :\equiv (a \in \mathbb{R} \to b \in \mathbb{R} \to \forall n \in \mathbb{N}\,(x(n) \in [a,b]) \to \varphi_{8}) \to a \in \mathbb{R} \to b \in \mathbb{R} \to \varphi_{53}$
\item $\varphi_{76} :\equiv (a \in \mathbb{R} \to b \in \mathbb{R} \to \varphi_{33}) \to (\overline{b}(n) = b \to b = \overline{b}(n)) \to a \in \mathbb{R} \to b \in \mathbb{R} \to \varphi_{34}$
\item $\varphi_{77} :\equiv (a \in \mathbb{R} \to b \in \mathbb{R} \to \varphi_{35}) \to (\overline{a}(n) = a \to a = \overline{a}(n)) \to a \in \mathbb{R} \to b \in \mathbb{R} \to \varphi_{36}$
\item $\varphi_{78} :\equiv (a \in \mathbb{R} \to b \in \mathbb{R} \to \varphi_{3}) \to (a \in \mathbb{R} \to b \in \mathbb{R} \to \varphi_{2}) \to a \in \mathbb{R} \to b \in \mathbb{R} \to \varphi_{17}$
\item $\varphi_{79} :\equiv (a \in \mathbb{R} \to b \in \mathbb{R} \to \varphi_{9}) \to (a \in \mathbb{R} \to b \in \mathbb{R} \to \varphi_{48}) \to a \in \mathbb{R} \to b \in \mathbb{R} \to \varphi_{3}$
\item $\varphi_{80} :\equiv (a \in \mathbb{R} \to b \in \mathbb{R} \to \varphi_{15}) \to (a \in \mathbb{R} \to b \in \mathbb{R} \to \varphi_{47}) \to a \in \mathbb{R} \to b \in \mathbb{R} \to \varphi_{2}$
\item $\varphi_{81} :\equiv (a \in \mathbb{R} \to b \in \mathbb{R} \to \varphi_{50}) \to (a \in \mathbb{R} \to b \in \mathbb{R} \to \varphi_{21}) \to a \in \mathbb{R} \to b \in \mathbb{R} \to \varphi_{33}$
\item $\varphi_{82} :\equiv (a \in \mathbb{R} \to b \in \mathbb{R} \to \varphi_{51}) \to (a \in \mathbb{R} \to b \in \mathbb{R} \to \varphi_{21}) \to a \in \mathbb{R} \to b \in \mathbb{R} \to \varphi_{35}$
\item $\varphi_{83} :\equiv (a \in \mathbb{R} \to b \in \mathbb{R} \to \varphi_{23}) \to (b = \overline{b}(n) \to x(n) < b \to x(n) < \overline{b}(n)) \to \varphi_{63}$
\item $\varphi_{84} :\equiv (a \in \mathbb{R} \to b \in \mathbb{R} \to \varphi_{24}) \to (a = \overline{a}(n) \to a < x(n) \to \overline{a}(n) < x(n)) \to \varphi_{64}$
\item $\varphi_{85} :\equiv (a \in \mathbb{R} \to b \in \mathbb{R} \to \varphi_{22}) \to (a \in \mathbb{R} \to b \in \mathbb{R} \to \varphi_{21}) \to \varphi_{18} \to a \in \mathbb{R} \to b \in \mathbb{R} \to \varphi_{23}$
\item $\varphi_{86} :\equiv (a \in \mathbb{R} \to b \in \mathbb{R} \to \varphi_{22}) \to (a \in \mathbb{R} \to b \in \mathbb{R} \to \varphi_{21}) \to \varphi_{18} \to a \in \mathbb{R} \to b \in \mathbb{R} \to \varphi_{24}$
\item $\varphi_{87} :\equiv (c \in [a,b] \leftrightarrow c \in \mathbb{R} \land a \leq c \land c \leq b) \to (a \in \mathbb{R} \to b \in \mathbb{R} \to \varphi_{14}) \to \varphi_{78}$
\item $\varphi_{88} :\equiv \forall n\,(\varphi_{25} \to n \in \mathbb{N} \to x(n) \leq \overline{b}(n)) \to \varphi_{72} \to (\varphi_{25} \to \forall n\,(\varphi_{25})) \to a \in \mathbb{R} \to b \in \mathbb{R} \to \varphi_{47}$
\item $\varphi_{89} :\equiv \forall n\,(\varphi_{25} \to n \in \mathbb{N} \to \overline{a}(n) \leq x(n)) \to \varphi_{73} \to (\varphi_{25} \to \forall n\,(\varphi_{25})) \to a \in \mathbb{R} \to b \in \mathbb{R} \to \varphi_{48}$
\end{abbrevlist}
\end{fullproof}

% 自動生成 (prover/build.js): CintBdd  ノード数 156
\begin{fullproof}{[G6]}{$\forall a\,\forall b\,\forall x\,(a \in \mathbb{R} \to b \in \mathbb{R} \to x \in [a,b] \to \lvert x \rvert \leq \lvert a \rvert + \lvert b \rvert)$}{156}\phantomsection\label{tree:CintBdd}
\noindent$\Pi_{1}$:
\begin{ptree}
\AXC{$\forall x \in \mathbb{R}\,(0 \leq \lvert x \rvert)$\lbl{[A6]}}
\AXC{$\forall x \in \mathbb{R}\,(0 \leq \lvert x \rvert) \to a \in \mathbb{R} \to 0 \leq \lvert a \rvert$\lbl{A2}}
\BIC{$a \in \mathbb{R} \to 0 \leq \lvert a \rvert$}
\AXC{$(a \in \mathbb{R} \to 0 \leq \lvert a \rvert) \to \varphi_{14}$\lbl{A1}}
\BIC{$\varphi_{14}$}
\end{ptree}
\noindent$\Pi_{2}$:
\begin{ptree}
\AXC{$\forall x \in \mathbb{R}\,(\lvert x \rvert \in \mathbb{R})$\lbl{[A5]}}
\AXC{$\forall x \in \mathbb{R}\,(\lvert x \rvert \in \mathbb{R}) \to a \in \mathbb{R} \to \lvert a \rvert \in \mathbb{R}$\lbl{A2}}
\BIC{$a \in \mathbb{R} \to \lvert a \rvert \in \mathbb{R}$}
\AXC{$(a \in \mathbb{R} \to \lvert a \rvert \in \mathbb{R}) \to \varphi_{19}$\lbl{A1}}
\BIC{$\varphi_{19}$}
\end{ptree}
\noindent$\Pi_{3}$:
\begin{ptree}
\AXC{$a \in \mathbb{R} \to a \in \mathbb{R}$\lbl{A1}}
\AXC{$\stackrel{\Pi_{2}}{\vdots}$}\noLine
\UIC{$\varphi_{19}$}
\BIC{$a \in \mathbb{R} \to b \in \mathbb{R} \to x \in [a,b] \to \lvert a \rvert \in \mathbb{R}$}
\end{ptree}
\noindent$\Pi_{4}$:
\begin{ptree}
\AXC{$\forall x \in \mathbb{R}\,(\lvert x \rvert \in \mathbb{R})$\lbl{[A5]}}
\AXC{$\forall x \in \mathbb{R}\,(\lvert x \rvert \in \mathbb{R}) \to b \in \mathbb{R} \to \lvert b \rvert \in \mathbb{R}$\lbl{A2}}
\BIC{$b \in \mathbb{R} \to \lvert b \rvert \in \mathbb{R}$}
\AXC{$(b \in \mathbb{R} \to \lvert b \rvert \in \mathbb{R}) \to \varphi_{29}$\lbl{A1}}
\BIC{$\varphi_{29}$}
\end{ptree}
\noindent$\Pi_{5}$:
\begin{ptree}
\AXC{$a \in \mathbb{R} \to b \in \mathbb{R} \to b \in \mathbb{R}$\lbl{A1}}
\AXC{$\stackrel{\Pi_{4}}{\vdots}$}\noLine
\UIC{$\varphi_{29}$}
\BIC{$a \in \mathbb{R} \to b \in \mathbb{R} \to x \in [a,b] \to \lvert b \rvert \in \mathbb{R}$}
\end{ptree}
\noindent$\Pi_{6}$:
\begin{ptree}
\AXC{$\forall x\,\forall y\,(x \in \mathbb{R} \to y \in \mathbb{R} \to 0 \leq y \to x \leq y + x)$\lbl{[E2]}}
\AXC{$\varphi_{32}$\lbl{A2}}
\BIC{$\forall y\,(\lvert b \rvert \in \mathbb{R} \to y \in \mathbb{R} \to 0 \leq y \to \lvert b \rvert \leq y + \lvert b \rvert)$}
\AXC{$\varphi_{45}$\lbl{A2}}
\BIC{$\varphi_{7}$}
\AXC{$\varphi_{7} \to \varphi_{58}$\lbl{A1}}
\BIC{$\varphi_{58}$}
\end{ptree}
\noindent$\Pi_{7}$:
\begin{ptree}
\AXC{$\stackrel{\Pi_{3}}{\vdots}$}\noLine
\UIC{$a \in \mathbb{R} \to b \in \mathbb{R} \to x \in [a,b] \to \lvert a \rvert \in \mathbb{R}$}
\AXC{$\stackrel{\Pi_{5}}{\vdots}$}\noLine
\UIC{$a \in \mathbb{R} \to b \in \mathbb{R} \to x \in [a,b] \to \lvert b \rvert \in \mathbb{R}$}
\AXC{$\stackrel{\Pi_{6}}{\vdots}$}\noLine
\UIC{$\varphi_{58}$}
\BIC{$(a \in \mathbb{R} \to b \in \mathbb{R} \to x \in [a,b] \to \lvert a \rvert \in \mathbb{R}) \to \varphi_{40}$}
\BIC{$\varphi_{40}$}
\end{ptree}
\noindent$\Pi_{8}$:
\begin{ptree}
\AXC{$\forall x \in \mathbb{R}\,(-\lvert x \rvert \leq x \land x \leq \lvert x \rvert)$\lbl{[A8]}}
\AXC{$\varphi_{16}$\lbl{A2}}
\BIC{$b \in \mathbb{R} \to -\lvert b \rvert \leq b \land b \leq \lvert b \rvert$}
\AXC{$(b \in \mathbb{R} \to -\lvert b \rvert \leq b \land b \leq \lvert b \rvert) \to \varphi_{33}$\lbl{A1}}
\BIC{$\varphi_{33}$}
\end{ptree}
\noindent$\Pi_{9}$:
\begin{ptree}
\AXC{$\forall a\,\forall b\,\forall x\,(x \in [a,b] \leftrightarrow x \in \mathbb{R} \land a \leq x \land x \leq b)$\lbl{Def{$[a,b]$}}}
\AXC{$\varphi_{30}$\lbl{A2}}
\BIC{$\forall b\,\forall x\,(x \in [a,b] \leftrightarrow x \in \mathbb{R} \land a \leq x \land x \leq b)$}
\AXC{$\varphi_{26}$\lbl{A2}}
\BIC{$\forall x\,(x \in [a,b] \leftrightarrow x \in \mathbb{R} \land a \leq x \land x \leq b)$}
\AXC{$\varphi_{20}$\lbl{A2}}
\BIC{$x \in [a,b] \leftrightarrow x \in \mathbb{R} \land a \leq x \land x \leq b$}
\end{ptree}
\noindent$\Pi_{10}$:
\begin{ptree}
\AXC{$a \in \mathbb{R} \to b \in \mathbb{R} \to x \in [a,b] \to x \in [a,b]$\lbl{A1}}
\AXC{$(a \in \mathbb{R} \to b \in \mathbb{R} \to x \in [a,b] \to x \in [a,b]) \to \varphi_{23}$\lbl{A1}}
\BIC{$\varphi_{23}$}
\end{ptree}
\noindent$\Pi_{11}$:
\begin{ptree}
\AXC{$\forall x\,\forall y\,\forall z\,(x \leq y \to y \leq z \to x \leq z)$\lbl{[O5]}}
\AXC{$\varphi_{3}$\lbl{A2}}
\BIC{$\forall y\,\forall z\,(x \leq y \to y \leq z \to x \leq z)$}
\AXC{$\varphi_{1}$\lbl{A2}}
\BIC{$\forall z\,(x \leq b \to b \leq z \to x \leq z)$}
\AXC{$\varphi_{2}$\lbl{A2}}
\BIC{$x \leq b \to b \leq \lvert b \rvert \to x \leq \lvert b \rvert$}
\AXC{$\varphi_{53}$\lbl{A1}}
\BIC{$(a \in \mathbb{R} \to b \in \mathbb{R} \to x \in [a,b] \to x \leq b) \to \varphi_{34}$}
\end{ptree}
\noindent$\Pi_{12}$:
\begin{ptree}
\AXC{$\stackrel{\Pi_{9}}{\vdots}$}\noLine
\UIC{$x \in [a,b] \leftrightarrow x \in \mathbb{R} \land a \leq x \land x \leq b$}
\AXC{$\stackrel{\Pi_{10}}{\vdots}$}\noLine
\UIC{$\varphi_{23}$}
\BIC{$a \in \mathbb{R} \to b \in \mathbb{R} \to x \in [a,b] \to x \leq b$}
\AXC{$\stackrel{\Pi_{11}}{\vdots}$}\noLine
\UIC{$(a \in \mathbb{R} \to b \in \mathbb{R} \to x \in [a,b] \to x \leq b) \to \varphi_{34}$}
\BIC{$\varphi_{34}$}
\end{ptree}
\noindent$\Pi_{13}$:
\begin{ptree}
\AXC{$a \in \mathbb{R} \to b \in \mathbb{R} \to b \in \mathbb{R}$\lbl{A1}}
\AXC{$\stackrel{\Pi_{8}}{\vdots}$}\noLine
\UIC{$\varphi_{33}$}
\BIC{$a \in \mathbb{R} \to b \in \mathbb{R} \to x \in [a,b] \to -\lvert b \rvert \leq b \land b \leq \lvert b \rvert$}
\AXC{$\varphi_{42}$\lbl{A1}}
\BIC{$a \in \mathbb{R} \to b \in \mathbb{R} \to x \in [a,b] \to b \leq \lvert b \rvert$}
\AXC{$\stackrel{\Pi_{12}}{\vdots}$}\noLine
\UIC{$\varphi_{34}$}
\BIC{$a \in \mathbb{R} \to b \in \mathbb{R} \to x \in [a,b] \to x \leq \lvert b \rvert$}
\end{ptree}
\noindent$\Pi_{14}$:
\begin{ptree}
\AXC{$\forall x\,\forall y\,\forall z\,(x \leq y \to y \leq z \to x \leq z)$\lbl{[O5]}}
\AXC{$\varphi_{3}$\lbl{A2}}
\BIC{$\forall y\,\forall z\,(x \leq y \to y \leq z \to x \leq z)$}
\AXC{$\varphi_{4}$\lbl{A2}}
\BIC{$\forall z\,(x \leq \lvert b \rvert \to \lvert b \rvert \leq z \to x \leq z)$}
\AXC{$\varphi_{21}$\lbl{A2}}
\BIC{$x \leq \lvert b \rvert \to \lvert b \rvert \leq \lvert a \rvert + \lvert b \rvert \to x \leq \lvert a \rvert + \lvert b \rvert$}
\AXC{$\varphi_{57}$\lbl{A1}}
\BIC{$(a \in \mathbb{R} \to b \in \mathbb{R} \to x \in [a,b] \to x \leq \lvert b \rvert) \to \varphi_{48}$}
\end{ptree}
\noindent$\Pi_{15}$:
\begin{ptree}
\AXC{$\stackrel{\Pi_{13}}{\vdots}$}\noLine
\UIC{$a \in \mathbb{R} \to b \in \mathbb{R} \to x \in [a,b] \to x \leq \lvert b \rvert$}
\AXC{$\stackrel{\Pi_{14}}{\vdots}$}\noLine
\UIC{$(a \in \mathbb{R} \to b \in \mathbb{R} \to x \in [a,b] \to x \leq \lvert b \rvert) \to \varphi_{48}$}
\BIC{$\varphi_{48}$}
\end{ptree}
\noindent$\Pi_{16}$:
\begin{ptree}
\AXC{$a \in \mathbb{R} \to b \in \mathbb{R} \to x \in [a,b] \to x \in [a,b]$\lbl{A1}}
\AXC{$(a \in \mathbb{R} \to b \in \mathbb{R} \to x \in [a,b] \to x \in [a,b]) \to \varphi_{24}$\lbl{A1}}
\BIC{$\varphi_{24}$}
\end{ptree}
\noindent$\Pi_{17}$:
\begin{ptree}
\AXC{$a \in \mathbb{R} \to a \in \mathbb{R}$\lbl{A1}}
\AXC{$\forall x \in \mathbb{R}\,(-\lvert x \rvert \leq x \land x \leq \lvert x \rvert)$\lbl{[A8]}}
\AXC{$\varphi_{17}$\lbl{A2}}
\BIC{$a \in \mathbb{R} \to -\lvert a \rvert \leq a \land a \leq \lvert a \rvert$}
\AXC{$(a \in \mathbb{R} \to -\lvert a \rvert \leq a \land a \leq \lvert a \rvert) \to \varphi_{27}$\lbl{A1}}
\BIC{$\varphi_{27}$}
\BIC{$a \in \mathbb{R} \to b \in \mathbb{R} \to x \in [a,b] \to -\lvert a \rvert \leq a \land a \leq \lvert a \rvert$}
\end{ptree}
\noindent$\Pi_{18}$:
\begin{ptree}
\AXC{$\forall x\,\forall y\,\forall z\,(x \leq y \to y \leq z \to x \leq z)$\lbl{[O5]}}
\AXC{$\varphi_{6}$\lbl{A2}}
\BIC{$\forall y\,\forall z\,(-\lvert a \rvert \leq y \to y \leq z \to -\lvert a \rvert \leq z)$}
\AXC{$\varphi_{18}$\lbl{A2}}
\BIC{$\forall z\,(-\lvert a \rvert \leq a \to a \leq z \to -\lvert a \rvert \leq z)$}
\AXC{$\varphi_{12}$\lbl{A2}}
\BIC{$-\lvert a \rvert \leq a \to a \leq x \to -\lvert a \rvert \leq x$}
\AXC{$\varphi_{55}$\lbl{A1}}
\BIC{$(a \in \mathbb{R} \to b \in \mathbb{R} \to x \in [a,b] \to -\lvert a \rvert \leq a) \to \varphi_{35}$}
\end{ptree}
\noindent$\Pi_{19}$:
\begin{ptree}
\AXC{$\stackrel{\Pi_{17}}{\vdots}$}\noLine
\UIC{$a \in \mathbb{R} \to b \in \mathbb{R} \to x \in [a,b] \to -\lvert a \rvert \leq a \land a \leq \lvert a \rvert$}
\AXC{$\varphi_{44}$\lbl{A1}}
\BIC{$a \in \mathbb{R} \to b \in \mathbb{R} \to x \in [a,b] \to -\lvert a \rvert \leq a$}
\AXC{$\stackrel{\Pi_{18}}{\vdots}$}\noLine
\UIC{$(a \in \mathbb{R} \to b \in \mathbb{R} \to x \in [a,b] \to -\lvert a \rvert \leq a) \to \varphi_{35}$}
\BIC{$\varphi_{35}$}
\end{ptree}
\noindent$\Pi_{20}$:
\begin{ptree}
\AXC{$\forall x \in \mathbb{R}\,(0 \leq \lvert x \rvert)$\lbl{[A6]}}
\AXC{$\forall x \in \mathbb{R}\,(0 \leq \lvert x \rvert) \to b \in \mathbb{R} \to 0 \leq \lvert b \rvert$\lbl{A2}}
\BIC{$b \in \mathbb{R} \to 0 \leq \lvert b \rvert$}
\AXC{$(b \in \mathbb{R} \to 0 \leq \lvert b \rvert) \to \varphi_{25}$\lbl{A1}}
\BIC{$\varphi_{25}$}
\end{ptree}
\noindent$\Pi_{21}$:
\begin{ptree}
\AXC{$\forall x\,\forall y\,(x \in \mathbb{R} \to y \in \mathbb{R} \to 0 \leq y \to x \leq x + y)$\lbl{[E1]}}
\AXC{$\varphi_{36}$\lbl{A2}}
\BIC{$\forall y\,(\lvert a \rvert \in \mathbb{R} \to y \in \mathbb{R} \to 0 \leq y \to \lvert a \rvert \leq \lvert a \rvert + y)$}
\AXC{$\varphi_{46}$\lbl{A2}}
\BIC{$\varphi_{8}$}
\AXC{$\varphi_{8} \to \varphi_{59}$\lbl{A1}}
\BIC{$\varphi_{59}$}
\end{ptree}
\noindent$\Pi_{22}$:
\begin{ptree}
\AXC{$\stackrel{\Pi_{5}}{\vdots}$}\noLine
\UIC{$a \in \mathbb{R} \to b \in \mathbb{R} \to x \in [a,b] \to \lvert b \rvert \in \mathbb{R}$}
\AXC{$\stackrel{\Pi_{3}}{\vdots}$}\noLine
\UIC{$a \in \mathbb{R} \to b \in \mathbb{R} \to x \in [a,b] \to \lvert a \rvert \in \mathbb{R}$}
\AXC{$\stackrel{\Pi_{21}}{\vdots}$}\noLine
\UIC{$\varphi_{59}$}
\BIC{$(a \in \mathbb{R} \to b \in \mathbb{R} \to x \in [a,b] \to \lvert b \rvert \in \mathbb{R}) \to \varphi_{41}$}
\BIC{$\varphi_{41}$}
\end{ptree}
\noindent$\Pi_{23}$:
\begin{ptree}
\AXC{$\forall x\,\forall y\,(x \in \mathbb{R} \to y \in \mathbb{R} \to x + y \in \mathbb{R})$\lbl{(R1)}}
\AXC{$\varphi_{31}$\lbl{A2}}
\BIC{$\forall y\,(\lvert a \rvert \in \mathbb{R} \to y \in \mathbb{R} \to \lvert a \rvert + y \in \mathbb{R})$}
\AXC{$\varphi_{39}$\lbl{A2}}
\BIC{$\lvert a \rvert \in \mathbb{R} \to \lvert b \rvert \in \mathbb{R} \to \lvert a \rvert + \lvert b \rvert \in \mathbb{R}$}
\AXC{$\varphi_{60}$\lbl{A1}}
\BIC{$(a \in \mathbb{R} \to b \in \mathbb{R} \to x \in [a,b] \to \lvert a \rvert \in \mathbb{R}) \to \varphi_{47}$}
\end{ptree}
\noindent$\Pi_{24}$:
\begin{ptree}
\AXC{$\stackrel{\Pi_{3}}{\vdots}$}\noLine
\UIC{$a \in \mathbb{R} \to b \in \mathbb{R} \to x \in [a,b] \to \lvert a \rvert \in \mathbb{R}$}
\AXC{$\stackrel{\Pi_{23}}{\vdots}$}\noLine
\UIC{$(a \in \mathbb{R} \to b \in \mathbb{R} \to x \in [a,b] \to \lvert a \rvert \in \mathbb{R}) \to \varphi_{47}$}
\BIC{$\varphi_{47}$}
\end{ptree}
\noindent$\Pi_{25}$:
\begin{ptree}
\AXC{$\stackrel{\Pi_{5}}{\vdots}$}\noLine
\UIC{$a \in \mathbb{R} \to b \in \mathbb{R} \to x \in [a,b] \to \lvert b \rvert \in \mathbb{R}$}
\AXC{$\stackrel{\Pi_{24}}{\vdots}$}\noLine
\UIC{$\varphi_{47}$}
\BIC{$a \in \mathbb{R} \to b \in \mathbb{R} \to x \in [a,b] \to \lvert a \rvert + \lvert b \rvert \in \mathbb{R}$}
\end{ptree}
\noindent$\Pi_{26}$:
\begin{ptree}
\AXC{$\forall x\,\forall y\,(x \in \mathbb{R} \to y \in \mathbb{R} \to x \leq y \to -y \leq -x)$\lbl{[F23]}}
\AXC{$\varphi_{38}$\lbl{A2}}
\BIC{$\forall y\,(\lvert a \rvert \in \mathbb{R} \to y \in \mathbb{R} \to \lvert a \rvert \leq y \to -y \leq -\lvert a \rvert)$}
\AXC{$\varphi_{52}$\lbl{A2}}
\BIC{$\lvert a \rvert \in \mathbb{R} \to \varphi_{22}$}
\AXC{$\varphi_{62}$\lbl{A1}}
\BIC{$(a \in \mathbb{R} \to b \in \mathbb{R} \to x \in [a,b] \to \lvert a \rvert \in \mathbb{R}) \to \varphi_{56}$}
\end{ptree}
\noindent$\Pi_{27}$:
\begin{ptree}
\AXC{$\stackrel{\Pi_{3}}{\vdots}$}\noLine
\UIC{$a \in \mathbb{R} \to b \in \mathbb{R} \to x \in [a,b] \to \lvert a \rvert \in \mathbb{R}$}
\AXC{$\stackrel{\Pi_{26}}{\vdots}$}\noLine
\UIC{$(a \in \mathbb{R} \to b \in \mathbb{R} \to x \in [a,b] \to \lvert a \rvert \in \mathbb{R}) \to \varphi_{56}$}
\BIC{$\varphi_{56}$}
\end{ptree}
\noindent$\Pi_{28}$:
\begin{ptree}
\AXC{$a \in \mathbb{R} \to b \in \mathbb{R} \to b \in \mathbb{R}$\lbl{A1}}
\AXC{$\stackrel{\Pi_{20}}{\vdots}$}\noLine
\UIC{$\varphi_{25}$}
\BIC{$a \in \mathbb{R} \to b \in \mathbb{R} \to x \in [a,b] \to 0 \leq \lvert b \rvert$}
\AXC{$\stackrel{\Pi_{22}}{\vdots}$}\noLine
\UIC{$\varphi_{41}$}
\BIC{$a \in \mathbb{R} \to b \in \mathbb{R} \to x \in [a,b] \to \lvert a \rvert \leq \lvert a \rvert + \lvert b \rvert$}
\AXC{$\stackrel{\Pi_{25}}{\vdots}$}\noLine
\UIC{$a \in \mathbb{R} \to b \in \mathbb{R} \to x \in [a,b] \to \lvert a \rvert + \lvert b \rvert \in \mathbb{R}$}
\AXC{$\stackrel{\Pi_{27}}{\vdots}$}\noLine
\UIC{$\varphi_{56}$}
\BIC{$\varphi_{50}$}
\BIC{$\varphi_{13}$}
\end{ptree}
\noindent$\Pi_{29}$:
\begin{ptree}
\AXC{$\forall x\,\forall y\,\forall z\,(x \leq y \to y \leq z \to x \leq z)$\lbl{[O5]}}
\AXC{$\forall x\,\forall y\,\forall z\,(x \leq y \to y \leq z \to x \leq z) \to \varphi_{5}$\lbl{A2}}
\BIC{$\varphi_{5}$}
\AXC{$\varphi_{5} \to \varphi_{15}$\lbl{A2}}
\BIC{$\varphi_{15}$}
\AXC{$\varphi_{15} \to \varphi_{11}$\lbl{A2}}
\BIC{$\varphi_{11}$}
\end{ptree}
\noindent$\Pi_{30}$:
\begin{ptree}
\AXC{$\stackrel{\Pi_{9}}{\vdots}$}\noLine
\UIC{$x \in [a,b] \leftrightarrow x \in \mathbb{R} \land a \leq x \land x \leq b$}
\AXC{$\stackrel{\Pi_{16}}{\vdots}$}\noLine
\UIC{$\varphi_{24}$}
\BIC{$a \in \mathbb{R} \to b \in \mathbb{R} \to x \in [a,b] \to a \leq x$}
\AXC{$\stackrel{\Pi_{19}}{\vdots}$}\noLine
\UIC{$\varphi_{35}$}
\BIC{$a \in \mathbb{R} \to b \in \mathbb{R} \to x \in [a,b] \to -\lvert a \rvert \leq x$}
\AXC{$\stackrel{\Pi_{28}}{\vdots}$}\noLine
\UIC{$\varphi_{13}$}
\AXC{$\stackrel{\Pi_{29}}{\vdots}$}\noLine
\UIC{$\varphi_{11}$}
\AXC{$\varphi_{11} \to \varphi_{13} \to \varphi_{43}$\lbl{A1}}
\BIC{$\varphi_{13} \to \varphi_{43}$}
\BIC{$\varphi_{43}$}
\BIC{$a \in \mathbb{R} \to b \in \mathbb{R} \to x \in [a,b] \to -(\lvert a \rvert + \lvert b \rvert) \leq x$}
\end{ptree}
\noindent$\Pi_{31}$:
\begin{ptree}
\AXC{$a \in \mathbb{R} \to b \in \mathbb{R} \to x \in [a,b] \to x \in [a,b]$\lbl{A1}}
\AXC{$(a \in \mathbb{R} \to b \in \mathbb{R} \to x \in [a,b] \to x \in [a,b]) \to \varphi_{28}$\lbl{A1}}
\BIC{$\varphi_{28}$}
\end{ptree}
\noindent$\Pi_{32}$:
\begin{ptree}
\AXC{$\forall x\,\forall y\,(x \in \mathbb{R} \to y \in \mathbb{R} \to -y \leq x \to x \leq y \to \lvert x \rvert \leq y)$\lbl{[A10]}}
\AXC{$\varphi_{37}$\lbl{A2}}
\BIC{$\forall y\,(x \in \mathbb{R} \to y \in \mathbb{R} \to -y \leq x \to x \leq y \to \lvert x \rvert \leq y)$}
\AXC{$\varphi_{51}$\lbl{A2}}
\BIC{$x \in \mathbb{R} \to \lvert a \rvert + \lvert b \rvert \in \mathbb{R} \to \varphi_{10}$}
\AXC{$\varphi_{63}$\lbl{A1}}
\BIC{$(a \in \mathbb{R} \to b \in \mathbb{R} \to x \in [a,b] \to x \in \mathbb{R}) \to \varphi_{61}$}
\end{ptree}
\noindent$\Pi_{33}$:
\begin{ptree}
\AXC{$\stackrel{\Pi_{9}}{\vdots}$}\noLine
\UIC{$x \in [a,b] \leftrightarrow x \in \mathbb{R} \land a \leq x \land x \leq b$}
\AXC{$\stackrel{\Pi_{31}}{\vdots}$}\noLine
\UIC{$\varphi_{28}$}
\BIC{$a \in \mathbb{R} \to b \in \mathbb{R} \to x \in [a,b] \to x \in \mathbb{R}$}
\AXC{$\stackrel{\Pi_{32}}{\vdots}$}\noLine
\UIC{$(a \in \mathbb{R} \to b \in \mathbb{R} \to x \in [a,b] \to x \in \mathbb{R}) \to \varphi_{61}$}
\BIC{$\varphi_{61}$}
\end{ptree}
\noindent$\Pi_{34}$:
\begin{ptree}
\AXC{$\stackrel{\Pi_{30}}{\vdots}$}\noLine
\UIC{$a \in \mathbb{R} \to b \in \mathbb{R} \to x \in [a,b] \to -(\lvert a \rvert + \lvert b \rvert) \leq x$}
\AXC{$\stackrel{\Pi_{25}}{\vdots}$}\noLine
\UIC{$a \in \mathbb{R} \to b \in \mathbb{R} \to x \in [a,b] \to \lvert a \rvert + \lvert b \rvert \in \mathbb{R}$}
\AXC{$\stackrel{\Pi_{33}}{\vdots}$}\noLine
\UIC{$\varphi_{61}$}
\BIC{$\varphi_{54}$}
\BIC{$\varphi_{49}$}
\end{ptree}
\noindent 主木:
\begin{ptree}
\AXC{$a \in \mathbb{R} \to a \in \mathbb{R}$\lbl{A1}}
\AXC{$\stackrel{\Pi_{1}}{\vdots}$}\noLine
\UIC{$\varphi_{14}$}
\BIC{$a \in \mathbb{R} \to b \in \mathbb{R} \to x \in [a,b] \to 0 \leq \lvert a \rvert$}
\AXC{$\stackrel{\Pi_{7}}{\vdots}$}\noLine
\UIC{$\varphi_{40}$}
\BIC{$a \in \mathbb{R} \to b \in \mathbb{R} \to x \in [a,b] \to \lvert b \rvert \leq \lvert a \rvert + \lvert b \rvert$}
\AXC{$\stackrel{\Pi_{15}}{\vdots}$}\noLine
\UIC{$\varphi_{48}$}
\BIC{$a \in \mathbb{R} \to b \in \mathbb{R} \to x \in [a,b] \to x \leq \lvert a \rvert + \lvert b \rvert$}
\AXC{$\stackrel{\Pi_{34}}{\vdots}$}\noLine
\UIC{$\varphi_{49}$}
\BIC{$a \in \mathbb{R} \to b \in \mathbb{R} \to x \in [a,b] \to \lvert x \rvert \leq \lvert a \rvert + \lvert b \rvert$}
\UIC{$\varphi_{9}$}
\UIC{$\forall b\,(\varphi_{9})$}
\UIC{$\forall a\,\forall b\,(\varphi_{9})$}
\end{ptree}
\begin{abbrevlist}
\item $\varphi_{1} :\equiv \forall y\,\forall z\,(x \leq y \to y \leq z \to x \leq z) \to \forall z\,(x \leq b \to b \leq z \to x \leq z)$
\item $\varphi_{2} :\equiv \forall z\,(x \leq b \to b \leq z \to x \leq z) \to x \leq b \to b \leq \lvert b \rvert \to x \leq \lvert b \rvert$
\item $\varphi_{3} :\equiv \forall x\,\forall y\,\forall z\,(x \leq y \to y \leq z \to x \leq z) \to \forall y\,\forall z\,(x \leq y \to y \leq z \to x \leq z)$
\item $\varphi_{4} :\equiv \forall y\,\forall z\,(x \leq y \to y \leq z \to x \leq z) \to \forall z\,(x \leq \lvert b \rvert \to \lvert b \rvert \leq z \to x \leq z)$
\item $\varphi_{5} :\equiv \forall y\,\forall z\,(-(\lvert a \rvert + \lvert b \rvert) \leq y \to y \leq z \to -(\lvert a \rvert + \lvert b \rvert) \leq z)$
\item $\varphi_{6} :\equiv \forall x\,\forall y\,\forall z\,(x \leq y \to y \leq z \to x \leq z) \to \forall y\,\forall z\,(-\lvert a \rvert \leq y \to y \leq z \to -\lvert a \rvert \leq z)$
\item $\varphi_{7} :\equiv \lvert b \rvert \in \mathbb{R} \to \lvert a \rvert \in \mathbb{R} \to 0 \leq \lvert a \rvert \to \lvert b \rvert \leq \lvert a \rvert + \lvert b \rvert$
\item $\varphi_{8} :\equiv \lvert a \rvert \in \mathbb{R} \to \lvert b \rvert \in \mathbb{R} \to 0 \leq \lvert b \rvert \to \lvert a \rvert \leq \lvert a \rvert + \lvert b \rvert$
\item $\varphi_{9} :\equiv \forall x\,(a \in \mathbb{R} \to b \in \mathbb{R} \to x \in [a,b] \to \lvert x \rvert \leq \lvert a \rvert + \lvert b \rvert)$
\item $\varphi_{10} :\equiv -(\lvert a \rvert + \lvert b \rvert) \leq x \to x \leq \lvert a \rvert + \lvert b \rvert \to \lvert x \rvert \leq \lvert a \rvert + \lvert b \rvert$
\item $\varphi_{11} :\equiv -(\lvert a \rvert + \lvert b \rvert) \leq -\lvert a \rvert \to -\lvert a \rvert \leq x \to -(\lvert a \rvert + \lvert b \rvert) \leq x$
\item $\varphi_{12} :\equiv \forall z\,(-\lvert a \rvert \leq a \to a \leq z \to -\lvert a \rvert \leq z) \to -\lvert a \rvert \leq a \to a \leq x \to -\lvert a \rvert \leq x$
\item $\varphi_{13} :\equiv a \in \mathbb{R} \to b \in \mathbb{R} \to x \in [a,b] \to -(\lvert a \rvert + \lvert b \rvert) \leq -\lvert a \rvert$
\item $\varphi_{14} :\equiv (a \in \mathbb{R} \to a \in \mathbb{R}) \to a \in \mathbb{R} \to b \in \mathbb{R} \to x \in [a,b] \to 0 \leq \lvert a \rvert$
\item $\varphi_{15} :\equiv \forall z\,(-(\lvert a \rvert + \lvert b \rvert) \leq -\lvert a \rvert \to -\lvert a \rvert \leq z \to -(\lvert a \rvert + \lvert b \rvert) \leq z)$
\item $\varphi_{16} :\equiv \forall x \in \mathbb{R}\,(-\lvert x \rvert \leq x \land x \leq \lvert x \rvert) \to b \in \mathbb{R} \to -\lvert b \rvert \leq b \land b \leq \lvert b \rvert$
\item $\varphi_{17} :\equiv \forall x \in \mathbb{R}\,(-\lvert x \rvert \leq x \land x \leq \lvert x \rvert) \to a \in \mathbb{R} \to -\lvert a \rvert \leq a \land a \leq \lvert a \rvert$
\item $\varphi_{18} :\equiv \forall y\,\forall z\,(-\lvert a \rvert \leq y \to y \leq z \to -\lvert a \rvert \leq z) \to \forall z\,(-\lvert a \rvert \leq a \to a \leq z \to -\lvert a \rvert \leq z)$
\item $\varphi_{19} :\equiv (a \in \mathbb{R} \to a \in \mathbb{R}) \to a \in \mathbb{R} \to b \in \mathbb{R} \to x \in [a,b] \to \lvert a \rvert \in \mathbb{R}$
\item $\varphi_{20} :\equiv \forall x\,(x \in [a,b] \leftrightarrow x \in \mathbb{R} \land a \leq x \land x \leq b) \to (x \in [a,b] \leftrightarrow x \in \mathbb{R} \land a \leq x \land x \leq b)$
\item $\varphi_{21} :\equiv \forall z\,(x \leq \lvert b \rvert \to \lvert b \rvert \leq z \to x \leq z) \to x \leq \lvert b \rvert \to \lvert b \rvert \leq \lvert a \rvert + \lvert b \rvert \to x \leq \lvert a \rvert + \lvert b \rvert$
\item $\varphi_{22} :\equiv \lvert a \rvert + \lvert b \rvert \in \mathbb{R} \to \lvert a \rvert \leq \lvert a \rvert + \lvert b \rvert \to -(\lvert a \rvert + \lvert b \rvert) \leq -\lvert a \rvert$
\item $\varphi_{23} :\equiv (x \in [a,b] \leftrightarrow x \in \mathbb{R} \land a \leq x \land x \leq b) \to a \in \mathbb{R} \to b \in \mathbb{R} \to x \in [a,b] \to x \leq b$
\item $\varphi_{24} :\equiv (x \in [a,b] \leftrightarrow x \in \mathbb{R} \land a \leq x \land x \leq b) \to a \in \mathbb{R} \to b \in \mathbb{R} \to x \in [a,b] \to a \leq x$
\item $\varphi_{25} :\equiv (a \in \mathbb{R} \to b \in \mathbb{R} \to b \in \mathbb{R}) \to a \in \mathbb{R} \to b \in \mathbb{R} \to x \in [a,b] \to 0 \leq \lvert b \rvert$
\item $\varphi_{26} :\equiv \forall b\,\forall x\,(x \in [a,b] \leftrightarrow x \in \mathbb{R} \land a \leq x \land x \leq b) \to \forall x\,(x \in [a,b] \leftrightarrow x \in \mathbb{R} \land a \leq x \land x \leq b)$
\item $\varphi_{27} :\equiv (a \in \mathbb{R} \to a \in \mathbb{R}) \to a \in \mathbb{R} \to b \in \mathbb{R} \to x \in [a,b] \to -\lvert a \rvert \leq a \land a \leq \lvert a \rvert$
\item $\varphi_{28} :\equiv (x \in [a,b] \leftrightarrow x \in \mathbb{R} \land a \leq x \land x \leq b) \to a \in \mathbb{R} \to b \in \mathbb{R} \to x \in [a,b] \to x \in \mathbb{R}$
\item $\varphi_{29} :\equiv (a \in \mathbb{R} \to b \in \mathbb{R} \to b \in \mathbb{R}) \to a \in \mathbb{R} \to b \in \mathbb{R} \to x \in [a,b] \to \lvert b \rvert \in \mathbb{R}$
\item $\varphi_{30} :\equiv \forall a\,\forall b\,\forall x\,(x \in [a,b] \leftrightarrow x \in \mathbb{R} \land a \leq x \land x \leq b) \to \forall b\,\forall x\,(x \in [a,b] \leftrightarrow x \in \mathbb{R} \land a \leq x \land x \leq b)$
\item $\varphi_{31} :\equiv \forall x\,\forall y\,(x \in \mathbb{R} \to y \in \mathbb{R} \to x + y \in \mathbb{R}) \to \forall y\,(\lvert a \rvert \in \mathbb{R} \to y \in \mathbb{R} \to \lvert a \rvert + y \in \mathbb{R})$
\item $\varphi_{32} :\equiv \forall x\,\forall y\,(x \in \mathbb{R} \to y \in \mathbb{R} \to 0 \leq y \to x \leq y + x) \to \forall y\,(\lvert b \rvert \in \mathbb{R} \to y \in \mathbb{R} \to 0 \leq y \to \lvert b \rvert \leq y + \lvert b \rvert)$
\item $\varphi_{33} :\equiv (a \in \mathbb{R} \to b \in \mathbb{R} \to b \in \mathbb{R}) \to a \in \mathbb{R} \to b \in \mathbb{R} \to x \in [a,b] \to -\lvert b \rvert \leq b \land b \leq \lvert b \rvert$
\item $\varphi_{34} :\equiv (a \in \mathbb{R} \to b \in \mathbb{R} \to x \in [a,b] \to b \leq \lvert b \rvert) \to a \in \mathbb{R} \to b \in \mathbb{R} \to x \in [a,b] \to x \leq \lvert b \rvert$
\item $\varphi_{35} :\equiv (a \in \mathbb{R} \to b \in \mathbb{R} \to x \in [a,b] \to a \leq x) \to a \in \mathbb{R} \to b \in \mathbb{R} \to x \in [a,b] \to -\lvert a \rvert \leq x$
\item $\varphi_{36} :\equiv \forall x\,\forall y\,(x \in \mathbb{R} \to y \in \mathbb{R} \to 0 \leq y \to x \leq x + y) \to \forall y\,(\lvert a \rvert \in \mathbb{R} \to y \in \mathbb{R} \to 0 \leq y \to \lvert a \rvert \leq \lvert a \rvert + y)$
\item $\varphi_{37} :\equiv \forall x\,\forall y\,(x \in \mathbb{R} \to y \in \mathbb{R} \to -y \leq x \to x \leq y \to \lvert x \rvert \leq y) \to \forall y\,(x \in \mathbb{R} \to y \in \mathbb{R} \to -y \leq x \to x \leq y \to \lvert x \rvert \leq y)$
\item $\varphi_{38} :\equiv \forall x\,\forall y\,(x \in \mathbb{R} \to y \in \mathbb{R} \to x \leq y \to -y \leq -x) \to \forall y\,(\lvert a \rvert \in \mathbb{R} \to y \in \mathbb{R} \to \lvert a \rvert \leq y \to -y \leq -\lvert a \rvert)$
\item $\varphi_{39} :\equiv \forall y\,(\lvert a \rvert \in \mathbb{R} \to y \in \mathbb{R} \to \lvert a \rvert + y \in \mathbb{R}) \to \lvert a \rvert \in \mathbb{R} \to \lvert b \rvert \in \mathbb{R} \to \lvert a \rvert + \lvert b \rvert \in \mathbb{R}$
\item $\varphi_{40} :\equiv (a \in \mathbb{R} \to b \in \mathbb{R} \to x \in [a,b] \to 0 \leq \lvert a \rvert) \to a \in \mathbb{R} \to b \in \mathbb{R} \to x \in [a,b] \to \lvert b \rvert \leq \lvert a \rvert + \lvert b \rvert$
\item $\varphi_{41} :\equiv (a \in \mathbb{R} \to b \in \mathbb{R} \to x \in [a,b] \to 0 \leq \lvert b \rvert) \to a \in \mathbb{R} \to b \in \mathbb{R} \to x \in [a,b] \to \lvert a \rvert \leq \lvert a \rvert + \lvert b \rvert$
\item $\varphi_{42} :\equiv (a \in \mathbb{R} \to b \in \mathbb{R} \to x \in [a,b] \to -\lvert b \rvert \leq b \land b \leq \lvert b \rvert) \to a \in \mathbb{R} \to b \in \mathbb{R} \to x \in [a,b] \to b \leq \lvert b \rvert$
\item $\varphi_{43} :\equiv (a \in \mathbb{R} \to b \in \mathbb{R} \to x \in [a,b] \to -\lvert a \rvert \leq x) \to a \in \mathbb{R} \to b \in \mathbb{R} \to x \in [a,b] \to -(\lvert a \rvert + \lvert b \rvert) \leq x$
\item $\varphi_{44} :\equiv (a \in \mathbb{R} \to b \in \mathbb{R} \to x \in [a,b] \to -\lvert a \rvert \leq a \land a \leq \lvert a \rvert) \to a \in \mathbb{R} \to b \in \mathbb{R} \to x \in [a,b] \to -\lvert a \rvert \leq a$
\item $\varphi_{45} :\equiv \forall y\,(\lvert b \rvert \in \mathbb{R} \to y \in \mathbb{R} \to 0 \leq y \to \lvert b \rvert \leq y + \lvert b \rvert) \to \varphi_{7}$
\item $\varphi_{46} :\equiv \forall y\,(\lvert a \rvert \in \mathbb{R} \to y \in \mathbb{R} \to 0 \leq y \to \lvert a \rvert \leq \lvert a \rvert + y) \to \varphi_{8}$
\item $\varphi_{47} :\equiv (a \in \mathbb{R} \to b \in \mathbb{R} \to x \in [a,b] \to \lvert b \rvert \in \mathbb{R}) \to a \in \mathbb{R} \to b \in \mathbb{R} \to x \in [a,b] \to \lvert a \rvert + \lvert b \rvert \in \mathbb{R}$
\item $\varphi_{48} :\equiv (a \in \mathbb{R} \to b \in \mathbb{R} \to x \in [a,b] \to \lvert b \rvert \leq \lvert a \rvert + \lvert b \rvert) \to a \in \mathbb{R} \to b \in \mathbb{R} \to x \in [a,b] \to x \leq \lvert a \rvert + \lvert b \rvert$
\item $\varphi_{49} :\equiv (a \in \mathbb{R} \to b \in \mathbb{R} \to x \in [a,b] \to x \leq \lvert a \rvert + \lvert b \rvert) \to a \in \mathbb{R} \to b \in \mathbb{R} \to x \in [a,b] \to \lvert x \rvert \leq \lvert a \rvert + \lvert b \rvert$
\item $\varphi_{50} :\equiv (a \in \mathbb{R} \to b \in \mathbb{R} \to x \in [a,b] \to \lvert a \rvert \leq \lvert a \rvert + \lvert b \rvert) \to \varphi_{13}$
\item $\varphi_{51} :\equiv \forall y\,(x \in \mathbb{R} \to y \in \mathbb{R} \to -y \leq x \to x \leq y \to \lvert x \rvert \leq y) \to x \in \mathbb{R} \to \lvert a \rvert + \lvert b \rvert \in \mathbb{R} \to \varphi_{10}$
\item $\varphi_{52} :\equiv \forall y\,(\lvert a \rvert \in \mathbb{R} \to y \in \mathbb{R} \to \lvert a \rvert \leq y \to -y \leq -\lvert a \rvert) \to \lvert a \rvert \in \mathbb{R} \to \varphi_{22}$
\item $\varphi_{53} :\equiv (x \leq b \to b \leq \lvert b \rvert \to x \leq \lvert b \rvert) \to (a \in \mathbb{R} \to b \in \mathbb{R} \to x \in [a,b] \to x \leq b) \to \varphi_{34}$
\item $\varphi_{54} :\equiv (a \in \mathbb{R} \to b \in \mathbb{R} \to x \in [a,b] \to -(\lvert a \rvert + \lvert b \rvert) \leq x) \to \varphi_{49}$
\item $\varphi_{55} :\equiv (-\lvert a \rvert \leq a \to a \leq x \to -\lvert a \rvert \leq x) \to (a \in \mathbb{R} \to b \in \mathbb{R} \to x \in [a,b] \to -\lvert a \rvert \leq a) \to \varphi_{35}$
\item $\varphi_{56} :\equiv (a \in \mathbb{R} \to b \in \mathbb{R} \to x \in [a,b] \to \lvert a \rvert + \lvert b \rvert \in \mathbb{R}) \to \varphi_{50}$
\item $\varphi_{57} :\equiv (x \leq \lvert b \rvert \to \lvert b \rvert \leq \lvert a \rvert + \lvert b \rvert \to x \leq \lvert a \rvert + \lvert b \rvert) \to (a \in \mathbb{R} \to b \in \mathbb{R} \to x \in [a,b] \to x \leq \lvert b \rvert) \to \varphi_{48}$
\item $\varphi_{58} :\equiv (a \in \mathbb{R} \to b \in \mathbb{R} \to x \in [a,b] \to \lvert b \rvert \in \mathbb{R}) \to (a \in \mathbb{R} \to b \in \mathbb{R} \to x \in [a,b] \to \lvert a \rvert \in \mathbb{R}) \to \varphi_{40}$
\item $\varphi_{59} :\equiv (a \in \mathbb{R} \to b \in \mathbb{R} \to x \in [a,b] \to \lvert a \rvert \in \mathbb{R}) \to (a \in \mathbb{R} \to b \in \mathbb{R} \to x \in [a,b] \to \lvert b \rvert \in \mathbb{R}) \to \varphi_{41}$
\item $\varphi_{60} :\equiv (\lvert a \rvert \in \mathbb{R} \to \lvert b \rvert \in \mathbb{R} \to \lvert a \rvert + \lvert b \rvert \in \mathbb{R}) \to (a \in \mathbb{R} \to b \in \mathbb{R} \to x \in [a,b] \to \lvert a \rvert \in \mathbb{R}) \to \varphi_{47}$
\item $\varphi_{61} :\equiv (a \in \mathbb{R} \to b \in \mathbb{R} \to x \in [a,b] \to \lvert a \rvert + \lvert b \rvert \in \mathbb{R}) \to \varphi_{54}$
\item $\varphi_{62} :\equiv (\lvert a \rvert \in \mathbb{R} \to \varphi_{22}) \to (a \in \mathbb{R} \to b \in \mathbb{R} \to x \in [a,b] \to \lvert a \rvert \in \mathbb{R}) \to \varphi_{56}$
\item $\varphi_{63} :\equiv (x \in \mathbb{R} \to \lvert a \rvert + \lvert b \rvert \in \mathbb{R} \to \varphi_{10}) \to (a \in \mathbb{R} \to b \in \mathbb{R} \to x \in [a,b] \to x \in \mathbb{R}) \to \varphi_{61}$
\end{abbrevlist}
\end{fullproof}

% 自動生成 (prover/build.js): BWin  ノード数 314
\begin{fullproof}{[G7]}{$\forall x\,\forall a\,\forall b\,(a \in \mathbb{R} \to b \in \mathbb{R} \to \mathrm{Seq}(x) \to \forall n \in \mathbb{N}\,(x(n) \in [a,b]) \to \exists s\,\exists c\,(\mathrm{SubIdx}(s) \land c \in [a,b] \land \mathrm{Conv}(x \circ s,c)))$}{314}\phantomsection\label{tree:BWin}
\noindent$\Pi_{1}$:
\begin{ptree}
\AXC{$y = z \to x = y \to x = z$\lbl{A6}}
\AXC{$(y = z \to x = y \to x = z) \to x = y \to y = z \to x = z$\lbl{A1}}
\BIC{$x = y \to y = z \to x = z$}
\UIC{$\forall z\,(x = y \to y = z \to x = z)$}
\UIC{$\forall y\,\forall z\,(x = y \to y = z \to x = z)$}
\UIC{$\forall x\,\forall y\,\forall z\,(x = y \to y = z \to x = z)$}
\AXC{$\varphi_{15}$\lbl{A2}}
\BIC{$\forall y\,\forall z\,(v_{780} = y \to y = z \to v_{780} = z)$}
\AXC{$\varphi_{20}$\lbl{A2}}
\BIC{$\forall z\,(v_{780} = v_{781} \to v_{781} = z \to v_{780} = z)$}
\AXC{$\varphi_{28}$\lbl{A2}}
\BIC{$v_{780} = v_{781} \to v_{781} = v_{782} \to v_{780} = v_{782}$}
\UIC{$\forall v_{782}\,(v_{780} = v_{781} \to v_{781} = v_{782} \to v_{780} = v_{782})$}
\UIC{$\varphi_{8}$}
\UIC{$\forall v_{780}\,(\varphi_{8})$}
\end{ptree}
\noindent$\Pi_{2}$:
\begin{ptree}
\AXC{$x = x$\lbl{A5}}
\AXC{$x = y \to x = x \to y = x$\lbl{A6}}
\AXC{$(x = y \to x = x \to y = x) \to x = x \to x = y \to y = x$\lbl{A1}}
\BIC{$x = x \to x = y \to y = x$}
\BIC{$x = y \to y = x$}
\UIC{$\forall y\,(x = y \to y = x)$}
\UIC{$\forall x\,\forall y\,(x = y \to y = x)$}
\end{ptree}
\noindent$\Pi_{3}$:
\begin{ptree}
\AXC{$\forall s\,(\varphi_{22})$\lbl{DefSubIdx}}
\AXC{$\forall s\,(\varphi_{22}) \to \varphi_{22}$\lbl{A2}}
\BIC{$\varphi_{22}$}
\AXC{$a \in \mathbb{R} \to b \in \mathbb{R} \to \mathrm{Seq}(x) \to \varphi_{122}$\lbl{A1}}
\AXC{$\varphi_{152}$\lbl{A1}}
\BIC{$a \in \mathbb{R} \to b \in \mathbb{R} \to \mathrm{Seq}(x) \to \varphi_{96}$}
\AXC{$\varphi_{159}$\lbl{A1}}
\BIC{$\varphi_{22} \to a \in \mathbb{R} \to b \in \mathbb{R} \to \mathrm{Seq}(x) \to \varphi_{95}$}
\BIC{$a \in \mathbb{R} \to b \in \mathbb{R} \to \mathrm{Seq}(x) \to \varphi_{95}$}
\end{ptree}
\noindent$\Pi_{4}$:
\begin{ptree}
\AXC{$\varphi_{62}$\lbl{[S12]}}
\AXC{$\varphi_{62} \to \forall s\,(\mathrm{Seq}(x) \to \varphi_{43})$\lbl{A2}}
\BIC{$\forall s\,(\mathrm{Seq}(x) \to \varphi_{43})$}
\AXC{$\forall s\,(\mathrm{Seq}(x) \to \varphi_{43}) \to \mathrm{Seq}(x) \to \varphi_{43}$\lbl{A2}}
\BIC{$\mathrm{Seq}(x) \to \varphi_{43}$}
\AXC{$\varphi_{169}$\lbl{A1}}
\BIC{$(a \in \mathbb{R} \to b \in \mathbb{R} \to \mathrm{Seq}(x) \to \mathrm{Seq}(x)) \to \varphi_{156}$}
\end{ptree}
\noindent$\Pi_{5}$:
\begin{ptree}
\AXC{$a \in \mathbb{R} \to b \in \mathbb{R} \to \mathrm{Seq}(x) \to \mathrm{Seq}(x)$\lbl{A1}}
\AXC{$\stackrel{\Pi_{4}}{\vdots}$}\noLine
\UIC{$(a \in \mathbb{R} \to b \in \mathbb{R} \to \mathrm{Seq}(x) \to \mathrm{Seq}(x)) \to \varphi_{156}$}
\BIC{$\varphi_{156}$}
\end{ptree}
\noindent$\Pi_{6}$:
\begin{ptree}
\AXC{$\stackrel{\Pi_{3}}{\vdots}$}\noLine
\UIC{$a \in \mathbb{R} \to b \in \mathbb{R} \to \mathrm{Seq}(x) \to \varphi_{95}$}
\AXC{$\stackrel{\Pi_{5}}{\vdots}$}\noLine
\UIC{$\varphi_{156}$}
\BIC{$a \in \mathbb{R} \to b \in \mathbb{R} \to \mathrm{Seq}(x) \to \varphi_{128}$}
\end{ptree}
\noindent$\Pi_{7}$:
\begin{ptree}
\AXC{$\varphi_{65}$\lbl{A2}}
\AXC{$\stackrel{\Pi_{6}}{\vdots}$}\noLine
\UIC{$a \in \mathbb{R} \to b \in \mathbb{R} \to \mathrm{Seq}(x) \to \varphi_{128}$}
\AXC{$\varphi_{164}$\lbl{A1}}
\BIC{$a \in \mathbb{R} \to b \in \mathbb{R} \to \mathrm{Seq}(x) \to \varphi_{127}$}
\AXC{$\varphi_{171}$\lbl{A1}}
\BIC{$\varphi_{65} \to a \in \mathbb{R} \to b \in \mathbb{R} \to \mathrm{Seq}(x) \to \varphi_{126}$}
\BIC{$a \in \mathbb{R} \to b \in \mathbb{R} \to \mathrm{Seq}(x) \to \varphi_{126}$}
\AXC{$(a \in \mathbb{R} \to b \in \mathbb{R} \to \mathrm{Seq}(x) \to \varphi_{126}) \to \varphi_{154}$\lbl{A1}}
\BIC{$\varphi_{154}$}
\end{ptree}
\noindent$\Pi_{8}$:
\begin{ptree}
\AXC{$a \in \mathbb{R} \to b \in \mathbb{R} \to \mathrm{Seq}(x) \to \varphi_{101}$\lbl{A1}}
\AXC{$\stackrel{\Pi_{7}}{\vdots}$}\noLine
\UIC{$\varphi_{154}$}
\BIC{$a \in \mathbb{R} \to b \in \mathbb{R} \to \mathrm{Seq}(x) \to \varphi_{116}$}
\AXC{$(a \in \mathbb{R} \to b \in \mathbb{R} \to \mathrm{Seq}(x) \to \varphi_{116}) \to \varphi_{139}$\lbl{A1}}
\BIC{$\varphi_{139}$}
\end{ptree}
\noindent$\Pi_{9}$:
\begin{ptree}
\AXC{$\varphi_{48}$\lbl{A2}}
\AXC{$a \in \mathbb{R} \to b \in \mathbb{R} \to \mathrm{Seq}(x) \to \varphi_{122}$\lbl{A1}}
\AXC{$\varphi_{163}$\lbl{A1}}
\BIC{$a \in \mathbb{R} \to b \in \mathbb{R} \to \mathrm{Seq}(x) \to \varphi_{123}$}
\AXC{$\varphi_{167}$\lbl{A1}}
\BIC{$\varphi_{48} \to a \in \mathbb{R} \to b \in \mathbb{R} \to \mathrm{Seq}(x) \to \varphi_{120}$}
\BIC{$a \in \mathbb{R} \to b \in \mathbb{R} \to \mathrm{Seq}(x) \to \varphi_{120}$}
\AXC{$(a \in \mathbb{R} \to b \in \mathbb{R} \to \mathrm{Seq}(x) \to \varphi_{120}) \to \varphi_{151}$\lbl{A1}}
\BIC{$\varphi_{151}$}
\end{ptree}
\noindent$\Pi_{10}$:
\begin{ptree}
\AXC{$a \in \mathbb{R} \to b \in \mathbb{R} \to \mathrm{Seq}(x) \to \varphi_{101}$\lbl{A1}}
\AXC{$\stackrel{\Pi_{9}}{\vdots}$}\noLine
\UIC{$\varphi_{151}$}
\BIC{$a \in \mathbb{R} \to b \in \mathbb{R} \to \mathrm{Seq}(x) \to \varphi_{111}$}
\AXC{$(a \in \mathbb{R} \to b \in \mathbb{R} \to \mathrm{Seq}(x) \to \varphi_{111}) \to \varphi_{165}$\lbl{A1}}
\BIC{$\varphi_{165}$}
\end{ptree}
\noindent$\Pi_{11}$:
\begin{ptree}
\AXC{$\stackrel{\Pi_{2}}{\vdots}$}\noLine
\UIC{$\forall x\,\forall y\,(x = y \to y = x)$}
\AXC{$\varphi_{21}$\lbl{A2}}
\BIC{$\forall y\,(x \circ s(n) = y \to y = x \circ s(n))$}
\AXC{$\varphi_{85}$\lbl{A2}}
\BIC{$x \circ s(n) = x(s_{n}) \to x(s_{n}) = x \circ s(n)$}
\AXC{$\stackrel{\Pi_{8}}{\vdots}$}\noLine
\UIC{$\varphi_{139}$}
\BIC{$a \in \mathbb{R} \to b \in \mathbb{R} \to \mathrm{Seq}(x) \to \varphi_{117}$}
\AXC{$\stackrel{\Pi_{10}}{\vdots}$}\noLine
\UIC{$\varphi_{165}$}
\BIC{$\varphi_{58} \to a \in \mathbb{R} \to b \in \mathbb{R} \to \mathrm{Seq}(x) \to \varphi_{112}$}
\end{ptree}
\noindent$\Pi_{12}$:
\begin{ptree}
\AXC{$\stackrel{\Pi_{1}}{\vdots}$}\noLine
\UIC{$\forall v_{780}\,(\varphi_{8})$}
\AXC{$\forall v_{780}\,(\varphi_{8}) \to \varphi_{13}$\lbl{A2}}
\BIC{$\varphi_{13}$}
\AXC{$\varphi_{13} \to \varphi_{32}$\lbl{A2}}
\BIC{$\varphi_{32}$}
\AXC{$\varphi_{32} \to \varphi_{58}$\lbl{A2}}
\BIC{$\varphi_{58}$}
\AXC{$\stackrel{\Pi_{11}}{\vdots}$}\noLine
\UIC{$\varphi_{58} \to a \in \mathbb{R} \to b \in \mathbb{R} \to \mathrm{Seq}(x) \to \varphi_{112}$}
\BIC{$a \in \mathbb{R} \to b \in \mathbb{R} \to \mathrm{Seq}(x) \to \varphi_{112}$}
\end{ptree}
\noindent$\Pi_{13}$:
\begin{ptree}
\AXC{$\varphi_{108} \to \forall n\,(\varphi_{108})$\lbl{A4}}
\AXC{$\varphi_{160}$\lbl{A3}}
\AXC{$\stackrel{\Pi_{12}}{\vdots}$}\noLine
\UIC{$a \in \mathbb{R} \to b \in \mathbb{R} \to \mathrm{Seq}(x) \to \varphi_{112}$}
\AXC{$\varphi_{157}$\lbl{A1}}
\BIC{$\varphi_{108} \to n \in \mathbb{N} \to y(n) = x \circ s(n)$}
\UIC{$\forall n\,(\varphi_{108} \to n \in \mathbb{N} \to y(n) = x \circ s(n))$}
\AXC{$\varphi_{181}$\lbl{A1}}
\BIC{$\varphi_{160} \to (\varphi_{108} \to \forall n\,(\varphi_{108})) \to a \in \mathbb{R} \to b \in \mathbb{R} \to \mathrm{Seq}(x) \to \varphi_{113}$}
\BIC{$(\varphi_{108} \to \forall n\,(\varphi_{108})) \to a \in \mathbb{R} \to b \in \mathbb{R} \to \mathrm{Seq}(x) \to \varphi_{113}$}
\BIC{$a \in \mathbb{R} \to b \in \mathbb{R} \to \mathrm{Seq}(x) \to \varphi_{113}$}
\end{ptree}
\noindent$\Pi_{14}$:
\begin{ptree}
\AXC{$\varphi_{46}$\lbl{[C13]}}
\AXC{$\varphi_{46} \to \varphi_{41}$\lbl{A2}}
\BIC{$\varphi_{41}$}
\AXC{$\varphi_{41} \to \varphi_{67}$\lbl{A2}}
\BIC{$\varphi_{67}$}
\AXC{$\varphi_{67} \to \mathrm{Seq}(x \circ s) \to \varphi_{38}$\lbl{A2}}
\BIC{$\mathrm{Seq}(x \circ s) \to \varphi_{38}$}
\AXC{$\varphi_{178}$\lbl{A1}}
\BIC{$(a \in \mathbb{R} \to b \in \mathbb{R} \to \mathrm{Seq}(x) \to \varphi_{97}) \to \varphi_{166}$}
\end{ptree}
\noindent$\Pi_{15}$:
\begin{ptree}
\AXC{$\stackrel{\Pi_{6}}{\vdots}$}\noLine
\UIC{$a \in \mathbb{R} \to b \in \mathbb{R} \to \mathrm{Seq}(x) \to \varphi_{128}$}
\AXC{$\varphi_{158}$\lbl{A1}}
\BIC{$a \in \mathbb{R} \to b \in \mathbb{R} \to \mathrm{Seq}(x) \to \varphi_{97}$}
\AXC{$\stackrel{\Pi_{14}}{\vdots}$}\noLine
\UIC{$(a \in \mathbb{R} \to b \in \mathbb{R} \to \mathrm{Seq}(x) \to \varphi_{97}) \to \varphi_{166}$}
\BIC{$\varphi_{166}$}
\end{ptree}
\noindent$\Pi_{16}$:
\begin{ptree}
\AXC{$a \in \mathbb{R} \to \varphi_{82}$\lbl{A1}}
\AXC{$\stackrel{\Pi_{13}}{\vdots}$}\noLine
\UIC{$a \in \mathbb{R} \to b \in \mathbb{R} \to \mathrm{Seq}(x) \to \varphi_{113}$}
\AXC{$\stackrel{\Pi_{15}}{\vdots}$}\noLine
\UIC{$\varphi_{166}$}
\BIC{$(a \in \mathbb{R} \to \varphi_{82}) \to a \in \mathbb{R} \to b \in \mathbb{R} \to \mathrm{Seq}(x) \to \varphi_{99}$}
\BIC{$a \in \mathbb{R} \to b \in \mathbb{R} \to \mathrm{Seq}(x) \to \varphi_{99}$}
\end{ptree}
\noindent$\Pi_{17}$:
\begin{ptree}
\AXC{$z = w \to z \in [a,b] \to w \in [a,b]$\lbl{A6}}
\UIC{$\forall w\,(z = w \to z \in [a,b] \to w \in [a,b])$}
\UIC{$\forall z\,\forall w\,(z = w \to z \in [a,b] \to w \in [a,b])$}
\AXC{$\forall z\,\forall w\,(z = w \to z \in [a,b] \to w \in [a,b]) \to \varphi_{34}$\lbl{A2}}
\BIC{$\varphi_{34}$}
\AXC{$\varphi_{34} \to \varphi_{70}$\lbl{A2}}
\BIC{$\varphi_{70}$}
\end{ptree}
\noindent$\Pi_{18}$:
\begin{ptree}
\AXC{$\stackrel{\Pi_{3}}{\vdots}$}\noLine
\UIC{$a \in \mathbb{R} \to b \in \mathbb{R} \to \mathrm{Seq}(x) \to \varphi_{95}$}
\AXC{$\forall s\,\forall k\,(\mathrm{NSeq}(s) \to k \in \mathbb{N} \to s_{k} \in \mathbb{N})$\lbl{[S13]}}
\AXC{$\varphi_{59}$\lbl{A2}}
\BIC{$\forall k\,(\mathrm{NSeq}(s) \to k \in \mathbb{N} \to s_{k} \in \mathbb{N})$}
\AXC{$\varphi_{51}$\lbl{A2}}
\BIC{$\mathrm{NSeq}(s) \to n \in \mathbb{N} \to s_{n} \in \mathbb{N}$}
\AXC{$\varphi_{173}$\lbl{A1}}
\BIC{$(a \in \mathbb{R} \to b \in \mathbb{R} \to \mathrm{Seq}(x) \to \varphi_{95}) \to \varphi_{150}$}
\BIC{$\varphi_{150}$}
\end{ptree}
\noindent$\Pi_{19}$:
\begin{ptree}
\AXC{$\varphi_{60}$\lbl{A2}}
\AXC{$a \in \mathbb{R} \to b \in \mathbb{R} \to \mathrm{Seq}(x) \to \varphi_{47}$\lbl{A1}}
\AXC{$\varphi_{146}$\lbl{A1}}
\BIC{$\varphi_{60} \to a \in \mathbb{R} \to b \in \mathbb{R} \to \mathrm{Seq}(x) \to \varphi_{60}$}
\BIC{$a \in \mathbb{R} \to b \in \mathbb{R} \to \mathrm{Seq}(x) \to \varphi_{60}$}
\AXC{$(a \in \mathbb{R} \to b \in \mathbb{R} \to \mathrm{Seq}(x) \to \varphi_{60}) \to \varphi_{155}$\lbl{A1}}
\BIC{$\varphi_{155}$}
\end{ptree}
\noindent$\Pi_{20}$:
\begin{ptree}
\AXC{$\varphi_{65}$\lbl{A2}}
\AXC{$\stackrel{\Pi_{6}}{\vdots}$}\noLine
\UIC{$a \in \mathbb{R} \to b \in \mathbb{R} \to \mathrm{Seq}(x) \to \varphi_{128}$}
\AXC{$\varphi_{164}$\lbl{A1}}
\BIC{$a \in \mathbb{R} \to b \in \mathbb{R} \to \mathrm{Seq}(x) \to \varphi_{127}$}
\AXC{$\varphi_{171}$\lbl{A1}}
\BIC{$\varphi_{65} \to a \in \mathbb{R} \to b \in \mathbb{R} \to \mathrm{Seq}(x) \to \varphi_{126}$}
\BIC{$a \in \mathbb{R} \to b \in \mathbb{R} \to \mathrm{Seq}(x) \to \varphi_{126}$}
\AXC{$(a \in \mathbb{R} \to b \in \mathbb{R} \to \mathrm{Seq}(x) \to \varphi_{126}) \to \varphi_{154}$\lbl{A1}}
\BIC{$\varphi_{154}$}
\end{ptree}
\noindent$\Pi_{21}$:
\begin{ptree}
\AXC{$a \in \mathbb{R} \to b \in \mathbb{R} \to \mathrm{Seq}(x) \to \varphi_{101}$\lbl{A1}}
\AXC{$\stackrel{\Pi_{20}}{\vdots}$}\noLine
\UIC{$\varphi_{154}$}
\BIC{$a \in \mathbb{R} \to b \in \mathbb{R} \to \mathrm{Seq}(x) \to \varphi_{116}$}
\AXC{$(a \in \mathbb{R} \to b \in \mathbb{R} \to \mathrm{Seq}(x) \to \varphi_{116}) \to \varphi_{139}$\lbl{A1}}
\BIC{$\varphi_{139}$}
\end{ptree}
\noindent$\Pi_{22}$:
\begin{ptree}
\AXC{$\stackrel{\Pi_{2}}{\vdots}$}\noLine
\UIC{$\forall x\,\forall y\,(x = y \to y = x)$}
\AXC{$\varphi_{21}$\lbl{A2}}
\BIC{$\forall y\,(x \circ s(n) = y \to y = x \circ s(n))$}
\AXC{$\varphi_{85}$\lbl{A2}}
\BIC{$x \circ s(n) = x(s_{n}) \to x(s_{n}) = x \circ s(n)$}
\AXC{$\stackrel{\Pi_{21}}{\vdots}$}\noLine
\UIC{$\varphi_{139}$}
\BIC{$a \in \mathbb{R} \to b \in \mathbb{R} \to \mathrm{Seq}(x) \to \varphi_{117}$}
\AXC{$(a \in \mathbb{R} \to b \in \mathbb{R} \to \mathrm{Seq}(x) \to \varphi_{117}) \to \varphi_{168}$\lbl{A1}}
\BIC{$\varphi_{168}$}
\end{ptree}
\noindent$\Pi_{23}$:
\begin{ptree}
\AXC{$\stackrel{\Pi_{17}}{\vdots}$}\noLine
\UIC{$\varphi_{70}$}
\AXC{$a \in \mathbb{R} \to b \in \mathbb{R} \to \mathrm{Seq}(x) \to \varphi_{101}$\lbl{A1}}
\AXC{$\stackrel{\Pi_{18}}{\vdots}$}\noLine
\UIC{$\varphi_{150}$}
\BIC{$a \in \mathbb{R} \to b \in \mathbb{R} \to \mathrm{Seq}(x) \to \varphi_{105}$}
\AXC{$\stackrel{\Pi_{19}}{\vdots}$}\noLine
\UIC{$\varphi_{155}$}
\BIC{$a \in \mathbb{R} \to b \in \mathbb{R} \to \mathrm{Seq}(x) \to \varphi_{114}$}
\AXC{$\stackrel{\Pi_{22}}{\vdots}$}\noLine
\UIC{$\varphi_{168}$}
\BIC{$\varphi_{70} \to a \in \mathbb{R} \to b \in \mathbb{R} \to \mathrm{Seq}(x) \to \varphi_{115}$}
\BIC{$a \in \mathbb{R} \to b \in \mathbb{R} \to \mathrm{Seq}(x) \to \varphi_{115}$}
\AXC{$\varphi_{161}$\lbl{A1}}
\BIC{$\varphi_{108} \to n \in \mathbb{N} \to x \circ s(n) \in [a,b]$}
\UIC{$\forall n\,(\varphi_{108} \to n \in \mathbb{N} \to x \circ s(n) \in [a,b])$}
\AXC{$\varphi_{182}$\lbl{A1}}
\BIC{$\varphi_{162} \to (\varphi_{108} \to \forall n\,(\varphi_{108})) \to a \in \mathbb{R} \to b \in \mathbb{R} \to \mathrm{Seq}(x) \to \varphi_{118}$}
\end{ptree}
\noindent$\Pi_{24}$:
\begin{ptree}
\AXC{$\varphi_{108} \to \forall n\,(\varphi_{108})$\lbl{A4}}
\AXC{$\varphi_{162}$\lbl{A3}}
\AXC{$\stackrel{\Pi_{23}}{\vdots}$}\noLine
\UIC{$\varphi_{162} \to (\varphi_{108} \to \forall n\,(\varphi_{108})) \to a \in \mathbb{R} \to b \in \mathbb{R} \to \mathrm{Seq}(x) \to \varphi_{118}$}
\BIC{$(\varphi_{108} \to \forall n\,(\varphi_{108})) \to a \in \mathbb{R} \to b \in \mathbb{R} \to \mathrm{Seq}(x) \to \varphi_{118}$}
\BIC{$a \in \mathbb{R} \to b \in \mathbb{R} \to \mathrm{Seq}(x) \to \varphi_{118}$}
\end{ptree}
\noindent$\Pi_{25}$:
\begin{ptree}
\AXC{$\varphi_{78}$\lbl{[G5]}}
\AXC{$\varphi_{78} \to \forall c\,\forall a\,\forall b\,(a \in \mathbb{R} \to b \in \mathbb{R} \to \varphi_{55})$\lbl{A2}}
\BIC{$\forall c\,\forall a\,\forall b\,(a \in \mathbb{R} \to b \in \mathbb{R} \to \varphi_{55})$}
\AXC{$\varphi_{133}$\lbl{A2}}
\BIC{$\forall a\,\forall b\,(a \in \mathbb{R} \to b \in \mathbb{R} \to \varphi_{55})$}
\AXC{$\varphi_{132}$\lbl{A2}}
\BIC{$\forall b\,(a \in \mathbb{R} \to b \in \mathbb{R} \to \varphi_{55})$}
\end{ptree}
\noindent$\Pi_{26}$:
\begin{ptree}
\AXC{$\stackrel{\Pi_{25}}{\vdots}$}\noLine
\UIC{$\forall b\,(a \in \mathbb{R} \to b \in \mathbb{R} \to \varphi_{55})$}
\AXC{$\forall b\,(a \in \mathbb{R} \to b \in \mathbb{R} \to \varphi_{55}) \to a \in \mathbb{R} \to b \in \mathbb{R} \to \varphi_{55}$\lbl{A2}}
\BIC{$a \in \mathbb{R} \to b \in \mathbb{R} \to \varphi_{55}$}
\end{ptree}
\noindent$\Pi_{27}$:
\begin{ptree}
\AXC{$a \in \mathbb{R} \to a \in \mathbb{R}$\lbl{A1}}
\AXC{$\stackrel{\Pi_{26}}{\vdots}$}\noLine
\UIC{$a \in \mathbb{R} \to b \in \mathbb{R} \to \varphi_{55}$}
\AXC{$(a \in \mathbb{R} \to b \in \mathbb{R} \to \varphi_{55}) \to (a \in \mathbb{R} \to a \in \mathbb{R}) \to \varphi_{174}$\lbl{A1}}
\BIC{$(a \in \mathbb{R} \to a \in \mathbb{R}) \to \varphi_{174}$}
\BIC{$\varphi_{174}$}
\end{ptree}
\noindent$\Pi_{28}$:
\begin{ptree}
\AXC{$\stackrel{\Pi_{16}}{\vdots}$}\noLine
\UIC{$a \in \mathbb{R} \to b \in \mathbb{R} \to \mathrm{Seq}(x) \to \varphi_{99}$}
\AXC{$\stackrel{\Pi_{24}}{\vdots}$}\noLine
\UIC{$a \in \mathbb{R} \to b \in \mathbb{R} \to \mathrm{Seq}(x) \to \varphi_{118}$}
\AXC{$a \in \mathbb{R} \to b \in \mathbb{R} \to b \in \mathbb{R}$\lbl{A1}}
\AXC{$\stackrel{\Pi_{27}}{\vdots}$}\noLine
\UIC{$\varphi_{174}$}
\BIC{$(a \in \mathbb{R} \to b \in \mathbb{R} \to \mathrm{Seq}(x) \to \varphi_{118}) \to \varphi_{147}$}
\BIC{$\varphi_{147}$}
\BIC{$a \in \mathbb{R} \to b \in \mathbb{R} \to \mathrm{Seq}(x) \to \varphi_{98}$}
\end{ptree}
\noindent$\Pi_{29}$:
\begin{ptree}
\AXC{$\stackrel{\Pi_{28}}{\vdots}$}\noLine
\UIC{$a \in \mathbb{R} \to b \in \mathbb{R} \to \mathrm{Seq}(x) \to \varphi_{98}$}
\AXC{$a \in \mathbb{R} \to b \in \mathbb{R} \to \mathrm{Seq}(x) \to \varphi_{122}$\lbl{A1}}
\AXC{$\varphi_{152}$\lbl{A1}}
\BIC{$a \in \mathbb{R} \to b \in \mathbb{R} \to \mathrm{Seq}(x) \to \varphi_{96}$}
\AXC{$\varphi_{177}$\lbl{A1}}
\BIC{$(a \in \mathbb{R} \to b \in \mathbb{R} \to \mathrm{Seq}(x) \to \varphi_{98}) \to \varphi_{153}$}
\BIC{$\varphi_{153}$}
\end{ptree}
\noindent$\Pi_{30}$:
\begin{ptree}
\AXC{$\forall s\,(\neg \varphi_{11}) \to \neg \varphi_{11}$\lbl{A2}}
\AXC{$\forall c\,(\varphi_{10}) \to \varphi_{10}$\lbl{A2}}
\AXC{$\stackrel{\Pi_{16}}{\vdots}$}\noLine
\UIC{$a \in \mathbb{R} \to b \in \mathbb{R} \to \mathrm{Seq}(x) \to \varphi_{99}$}
\AXC{$\stackrel{\Pi_{29}}{\vdots}$}\noLine
\UIC{$\varphi_{153}$}
\BIC{$a \in \mathbb{R} \to b \in \mathbb{R} \to \mathrm{Seq}(x) \to \varphi_{119}$}
\AXC{$\varphi_{170}$\lbl{A1}}
\BIC{$(\forall c\,(\varphi_{10}) \to \varphi_{10}) \to a \in \mathbb{R} \to b \in \mathbb{R} \to \mathrm{Seq}(x) \to \varphi_{121}$}
\BIC{$a \in \mathbb{R} \to b \in \mathbb{R} \to \mathrm{Seq}(x) \to \varphi_{121}$}
\AXC{$\varphi_{172}$\lbl{A1}}
\BIC{$(\forall s\,(\neg \varphi_{11}) \to \neg \varphi_{11}) \to a \in \mathbb{R} \to b \in \mathbb{R} \to \mathrm{Seq}(x) \to \varphi_{124}$}
\BIC{$a \in \mathbb{R} \to b \in \mathbb{R} \to \mathrm{Seq}(x) \to \varphi_{124}$}
\end{ptree}
\noindent$\Pi_{31}$:
\begin{ptree}
\AXC{$\varphi_{54}$\lbl{DefSubSeq}}
\AXC{$\varphi_{54} \to \varphi_{49}$\lbl{A2}}
\BIC{$\varphi_{49}$}
\AXC{$\varphi_{49} \to (\mathrm{SubSeq}(y,x) \leftrightarrow \mathrm{Seq}(y) \land \varphi_{17})$\lbl{A2}}
\BIC{$\mathrm{SubSeq}(y,x) \leftrightarrow \mathrm{Seq}(y) \land \varphi_{17}$}
\AXC{$a \in \mathbb{R} \to \varphi_{81}$\lbl{A1}}
\AXC{$(a \in \mathbb{R} \to \varphi_{81}) \to \varphi_{135}$\lbl{A1}}
\BIC{$\varphi_{135}$}
\BIC{$a \in \mathbb{R} \to b \in \mathbb{R} \to \mathrm{Seq}(x) \to \varphi_{92}$}
\end{ptree}
\noindent$\Pi_{32}$:
\begin{ptree}
\AXC{$\stackrel{\Pi_{31}}{\vdots}$}\noLine
\UIC{$a \in \mathbb{R} \to b \in \mathbb{R} \to \mathrm{Seq}(x) \to \varphi_{92}$}
\AXC{$\varphi_{145}$\lbl{A1}}
\BIC{$a \in \mathbb{R} \to b \in \mathbb{R} \to \mathrm{Seq}(x) \to \varphi_{93}$}
\AXC{$(a \in \mathbb{R} \to b \in \mathbb{R} \to \mathrm{Seq}(x) \to \varphi_{93}) \to \forall s\,(\varphi_{130}) \to \varphi_{179}$\lbl{A1}}
\BIC{$\forall s\,(\varphi_{130}) \to \varphi_{179}$}
\end{ptree}
\noindent$\Pi_{33}$:
\begin{ptree}
\AXC{$\stackrel{\Pi_{30}}{\vdots}$}\noLine
\UIC{$a \in \mathbb{R} \to b \in \mathbb{R} \to \mathrm{Seq}(x) \to \varphi_{124}$}
\AXC{$(a \in \mathbb{R} \to b \in \mathbb{R} \to \mathrm{Seq}(x) \to \varphi_{124}) \to \varphi_{130}$\lbl{A1}}
\BIC{$\varphi_{130}$}
\UIC{$\forall s\,(\varphi_{130})$}
\AXC{$\stackrel{\Pi_{32}}{\vdots}$}\noLine
\UIC{$\forall s\,(\varphi_{130}) \to \varphi_{179}$}
\BIC{$\varphi_{179}$}
\end{ptree}
\noindent$\Pi_{34}$:
\begin{ptree}
\AXC{$\varphi_{109} \to \forall s\,(\varphi_{109})$\lbl{A4}}
\AXC{$\forall s\,(\varphi_{130}) \to \forall s\,(\varphi_{109}) \to \varphi_{18}$\lbl{A3}}
\AXC{$\stackrel{\Pi_{33}}{\vdots}$}\noLine
\UIC{$\varphi_{179}$}
\BIC{$(\varphi_{109} \to \forall s\,(\varphi_{109})) \to a \in \mathbb{R} \to \varphi_{102}$}
\BIC{$a \in \mathbb{R} \to \varphi_{102}$}
\AXC{$(a \in \mathbb{R} \to \varphi_{102}) \to \varphi_{110}$\lbl{A1}}
\BIC{$\varphi_{110}$}
\UIC{$\forall c\,(\varphi_{110})$}
\end{ptree}
\noindent$\Pi_{35}$:
\begin{ptree}
\AXC{$\forall a\,(\mathrm{Cgt}(a) \leftrightarrow \exists \alpha\,(\mathrm{Conv}(a,\alpha)))$\lbl{DefCgt}}
\AXC{$\varphi_{12}$\lbl{A2}}
\BIC{$\mathrm{Cgt}(y) \leftrightarrow \exists \alpha\,(\mathrm{Conv}(y,\alpha))$}
\AXC{$a \in \mathbb{R} \to \varphi_{81}$\lbl{A1}}
\AXC{$\varphi_{138}$\lbl{A1}}
\BIC{$(\mathrm{Cgt}(y) \leftrightarrow \exists \alpha\,(\mathrm{Conv}(y,\alpha))) \to a \in \mathbb{R} \to \varphi_{74}$}
\BIC{$a \in \mathbb{R} \to \varphi_{74}$}
\end{ptree}
\noindent$\Pi_{36}$:
\begin{ptree}
\AXC{$\varphi_{144}$\lbl{A4}}
\AXC{$\forall c\,(\varphi_{110}) \to \varphi_{103} \to \forall c\,(\neg \mathrm{Conv}(y,c))$\lbl{A3}}
\AXC{$\stackrel{\Pi_{34}}{\vdots}$}\noLine
\UIC{$\forall c\,(\varphi_{110})$}
\AXC{$\stackrel{\Pi_{35}}{\vdots}$}\noLine
\UIC{$a \in \mathbb{R} \to \varphi_{74}$}
\AXC{$(a \in \mathbb{R} \to \varphi_{74}) \to \varphi_{180}$\lbl{A1}}
\BIC{$\varphi_{180}$}
\BIC{$(\forall c\,(\varphi_{110}) \to \varphi_{103} \to \forall c\,(\neg \mathrm{Conv}(y,c))) \to \varphi_{144} \to \varphi_{141}$}
\BIC{$\varphi_{144} \to \varphi_{141}$}
\BIC{$\varphi_{141}$}
\end{ptree}
\noindent$\Pi_{37}$:
\begin{ptree}
\AXC{$\varphi_{1}$\lbl{A4}}
\AXC{$\varphi_{31}$\lbl{A3}}
\AXC{$\forall c\,(\neg \mathrm{Conv}(y,c)) \to \neg \mathrm{Conv}(y,\alpha)$\lbl{A2}}
\UIC{$\forall \alpha\,(\forall c\,(\neg \mathrm{Conv}(y,c)) \to \neg \mathrm{Conv}(y,\alpha))$}
\AXC{$\stackrel{\Pi_{36}}{\vdots}$}\noLine
\UIC{$\varphi_{141}$}
\BIC{$\varphi_{31} \to \varphi_{1} \to a \in \mathbb{R} \to b \in \mathbb{R} \to \mathrm{Seq}(x) \to \varphi_{89}$}
\BIC{$\varphi_{1} \to a \in \mathbb{R} \to b \in \mathbb{R} \to \mathrm{Seq}(x) \to \varphi_{89}$}
\BIC{$a \in \mathbb{R} \to b \in \mathbb{R} \to \mathrm{Seq}(x) \to \varphi_{89}$}
\AXC{$\varphi_{143}$\lbl{A1}}
\BIC{$\varphi_{39} \land \neg \exists s\,(\varphi_{11}) \to \neg (\mathrm{SubSeq}(y,x) \land \mathrm{Cgt}(y))$}
\UIC{$\varphi_{104}$}
\end{ptree}
\noindent$\Pi_{38}$:
\begin{ptree}
\AXC{$\forall a\,\forall b\,(\varphi_{35})$\lbl{[G6]}}
\AXC{$\forall a\,\forall b\,(\varphi_{35}) \to \forall b\,(\varphi_{35})$\lbl{A2}}
\BIC{$\forall b\,(\varphi_{35})$}
\AXC{$\forall b\,(\varphi_{35}) \to \varphi_{35}$\lbl{A2}}
\BIC{$\varphi_{35}$}
\AXC{$\varphi_{35} \to \varphi_{50}$\lbl{A2}}
\BIC{$\varphi_{50}$}
\AXC{$\varphi_{149}$\lbl{A1}}
\BIC{$(a \in \mathbb{R} \to a \in \mathbb{R}) \to (a \in \mathbb{R} \to b \in \mathbb{R} \to b \in \mathbb{R}) \to \varphi_{131}$}
\end{ptree}
\noindent$\Pi_{39}$:
\begin{ptree}
\AXC{$a \in \mathbb{R} \to b \in \mathbb{R} \to b \in \mathbb{R}$\lbl{A1}}
\AXC{$a \in \mathbb{R} \to a \in \mathbb{R}$\lbl{A1}}
\AXC{$\stackrel{\Pi_{38}}{\vdots}$}\noLine
\UIC{$(a \in \mathbb{R} \to a \in \mathbb{R}) \to (a \in \mathbb{R} \to b \in \mathbb{R} \to b \in \mathbb{R}) \to \varphi_{131}$}
\BIC{$(a \in \mathbb{R} \to b \in \mathbb{R} \to b \in \mathbb{R}) \to \varphi_{131}$}
\BIC{$\varphi_{131}$}
\end{ptree}
\noindent$\Pi_{40}$:
\begin{ptree}
\AXC{$\varphi_{136}$\lbl{A3}}
\AXC{$\varphi_{42}$\lbl{A2}}
\AXC{$a \in \mathbb{R} \to b \in \mathbb{R} \to \mathrm{Seq}(x) \to \varphi_{47}$\lbl{A1}}
\AXC{$\varphi_{142}$\lbl{A1}}
\BIC{$\varphi_{42} \to a \in \mathbb{R} \to b \in \mathbb{R} \to \mathrm{Seq}(x) \to \varphi_{42}$}
\BIC{$a \in \mathbb{R} \to b \in \mathbb{R} \to \mathrm{Seq}(x) \to \varphi_{42}$}
\AXC{$\stackrel{\Pi_{39}}{\vdots}$}\noLine
\UIC{$\varphi_{131}$}
\BIC{$a \in \mathbb{R} \to b \in \mathbb{R} \to \mathrm{Seq}(x) \to \varphi_{53}$}
\AXC{$\varphi_{134}$\lbl{A1}}
\BIC{$\varphi_{39} \to n \in \mathbb{N} \to \lvert x(n) \rvert \leq \lvert a \rvert + \lvert b \rvert$}
\UIC{$\forall n\,(\varphi_{39} \to n \in \mathbb{N} \to \lvert x(n) \rvert \leq \lvert a \rvert + \lvert b \rvert)$}
\AXC{$\varphi_{176}$\lbl{A1}}
\BIC{$\varphi_{136} \to (\varphi_{39} \to \forall n\,(\varphi_{39})) \to a \in \mathbb{R} \to b \in \mathbb{R} \to \mathrm{Seq}(x) \to \varphi_{56}$}
\BIC{$(\varphi_{39} \to \forall n\,(\varphi_{39})) \to a \in \mathbb{R} \to b \in \mathbb{R} \to \mathrm{Seq}(x) \to \varphi_{56}$}
\end{ptree}
\noindent$\Pi_{41}$:
\begin{ptree}
\AXC{$\forall x \in \mathbb{R}\,(\lvert x \rvert \in \mathbb{R})$\lbl{[A5]}}
\AXC{$\forall x \in \mathbb{R}\,(\lvert x \rvert \in \mathbb{R}) \to b \in \mathbb{R} \to \lvert b \rvert \in \mathbb{R}$\lbl{A2}}
\BIC{$b \in \mathbb{R} \to \lvert b \rvert \in \mathbb{R}$}
\AXC{$\varphi_{107}$\lbl{A1}}
\BIC{$(a \in \mathbb{R} \to b \in \mathbb{R} \to b \in \mathbb{R}) \to a \in \mathbb{R} \to b \in \mathbb{R} \to \varphi_{26}$}
\end{ptree}
\noindent$\Pi_{42}$:
\begin{ptree}
\AXC{$a \in \mathbb{R} \to b \in \mathbb{R} \to b \in \mathbb{R}$\lbl{A1}}
\AXC{$\stackrel{\Pi_{41}}{\vdots}$}\noLine
\UIC{$(a \in \mathbb{R} \to b \in \mathbb{R} \to b \in \mathbb{R}) \to a \in \mathbb{R} \to b \in \mathbb{R} \to \varphi_{26}$}
\BIC{$a \in \mathbb{R} \to b \in \mathbb{R} \to \varphi_{26}$}
\end{ptree}
\noindent$\Pi_{43}$:
\begin{ptree}
\AXC{$a \in \mathbb{R} \to a \in \mathbb{R}$\lbl{A1}}
\AXC{$\forall x \in \mathbb{R}\,(\lvert x \rvert \in \mathbb{R})$\lbl{[A5]}}
\AXC{$\forall x \in \mathbb{R}\,(\lvert x \rvert \in \mathbb{R}) \to a \in \mathbb{R} \to \lvert a \rvert \in \mathbb{R}$\lbl{A2}}
\BIC{$a \in \mathbb{R} \to \lvert a \rvert \in \mathbb{R}$}
\AXC{$\varphi_{100}$\lbl{A1}}
\BIC{$(a \in \mathbb{R} \to a \in \mathbb{R}) \to a \in \mathbb{R} \to b \in \mathbb{R} \to \varphi_{27}$}
\BIC{$a \in \mathbb{R} \to b \in \mathbb{R} \to \varphi_{27}$}
\end{ptree}
\noindent$\Pi_{44}$:
\begin{ptree}
\AXC{$\forall x\,\forall y\,(x \in \mathbb{R} \to y \in \mathbb{R} \to x + y \in \mathbb{R})$\lbl{(R1)}}
\AXC{$\varphi_{77}$\lbl{A2}}
\BIC{$\forall y\,(\lvert a \rvert \in \mathbb{R} \to y \in \mathbb{R} \to \lvert a \rvert + y \in \mathbb{R})$}
\AXC{$\varphi_{90}$\lbl{A2}}
\BIC{$\lvert a \rvert \in \mathbb{R} \to \lvert b \rvert \in \mathbb{R} \to \lvert a \rvert + \lvert b \rvert \in \mathbb{R}$}
\AXC{$(\lvert a \rvert \in \mathbb{R} \to \lvert b \rvert \in \mathbb{R} \to \lvert a \rvert + \lvert b \rvert \in \mathbb{R}) \to \varphi_{140}$\lbl{A1}}
\BIC{$\varphi_{140}$}
\end{ptree}
\noindent$\Pi_{45}$:
\begin{ptree}
\AXC{$\stackrel{\Pi_{42}}{\vdots}$}\noLine
\UIC{$a \in \mathbb{R} \to b \in \mathbb{R} \to \varphi_{26}$}
\AXC{$\stackrel{\Pi_{43}}{\vdots}$}\noLine
\UIC{$a \in \mathbb{R} \to b \in \mathbb{R} \to \varphi_{27}$}
\AXC{$\stackrel{\Pi_{44}}{\vdots}$}\noLine
\UIC{$\varphi_{140}$}
\BIC{$(a \in \mathbb{R} \to b \in \mathbb{R} \to \varphi_{26}) \to a \in \mathbb{R} \to b \in \mathbb{R} \to \varphi_{44}$}
\BIC{$a \in \mathbb{R} \to b \in \mathbb{R} \to \varphi_{44}$}
\AXC{$\varphi_{148}$\lbl{A1}}
\BIC{$(a \in \mathbb{R} \to b \in \mathbb{R} \to \mathrm{Seq}(x) \to \varphi_{56}) \to a \in \mathbb{R} \to \varphi_{94}$}
\end{ptree}
\noindent$\Pi_{46}$:
\begin{ptree}
\AXC{$\varphi_{39} \to \forall n\,(\varphi_{39})$\lbl{A4}}
\AXC{$\stackrel{\Pi_{40}}{\vdots}$}\noLine
\UIC{$(\varphi_{39} \to \forall n\,(\varphi_{39})) \to a \in \mathbb{R} \to b \in \mathbb{R} \to \mathrm{Seq}(x) \to \varphi_{56}$}
\BIC{$a \in \mathbb{R} \to b \in \mathbb{R} \to \mathrm{Seq}(x) \to \varphi_{56}$}
\AXC{$\stackrel{\Pi_{45}}{\vdots}$}\noLine
\UIC{$(a \in \mathbb{R} \to b \in \mathbb{R} \to \mathrm{Seq}(x) \to \varphi_{56}) \to a \in \mathbb{R} \to \varphi_{94}$}
\BIC{$a \in \mathbb{R} \to \varphi_{94}$}
\end{ptree}
\noindent$\Pi_{47}$:
\begin{ptree}
\AXC{$\forall M\,(\neg (M \in \mathbb{R} \land \forall n \in \mathbb{N}\,(\lvert x(n) \rvert \leq M))) \to \neg \varphi_{40}$\lbl{A2}}
\AXC{$\stackrel{\Pi_{46}}{\vdots}$}\noLine
\UIC{$a \in \mathbb{R} \to \varphi_{94}$}
\AXC{$(a \in \mathbb{R} \to \varphi_{94}) \to \varphi_{137}$\lbl{A1}}
\BIC{$\varphi_{137}$}
\BIC{$a \in \mathbb{R} \to b \in \mathbb{R} \to \mathrm{Seq}(x) \to \varphi_{57}$}
\end{ptree}
\noindent$\Pi_{48}$:
\begin{ptree}
\AXC{$\varphi_{36}$\lbl{DefSeqBdd}}
\AXC{$\varphi_{36} \to \varphi_{29}$\lbl{A2}}
\BIC{$\varphi_{29}$}
\AXC{$\varphi_{29} \to (a \in \mathbb{R} \to b \in \mathbb{R} \to \mathrm{Seq}(x) \to \mathrm{Seq}(x)) \to \varphi_{129}$\lbl{A1}}
\BIC{$(a \in \mathbb{R} \to b \in \mathbb{R} \to \mathrm{Seq}(x) \to \mathrm{Seq}(x)) \to \varphi_{129}$}
\end{ptree}
\noindent$\Pi_{49}$:
\begin{ptree}
\AXC{$a \in \mathbb{R} \to b \in \mathbb{R} \to \mathrm{Seq}(x) \to \mathrm{Seq}(x)$\lbl{A1}}
\AXC{$\stackrel{\Pi_{48}}{\vdots}$}\noLine
\UIC{$(a \in \mathbb{R} \to b \in \mathbb{R} \to \mathrm{Seq}(x) \to \mathrm{Seq}(x)) \to \varphi_{129}$}
\BIC{$\varphi_{129}$}
\end{ptree}
\noindent$\Pi_{50}$:
\begin{ptree}
\AXC{$\stackrel{\Pi_{47}}{\vdots}$}\noLine
\UIC{$a \in \mathbb{R} \to b \in \mathbb{R} \to \mathrm{Seq}(x) \to \varphi_{57}$}
\AXC{$\stackrel{\Pi_{49}}{\vdots}$}\noLine
\UIC{$\varphi_{129}$}
\BIC{$a \in \mathbb{R} \to b \in \mathbb{R} \to \varphi_{19}$}
\AXC{$\varphi_{5}$\lbl{定理6.8}}
\AXC{$\varphi_{5} \to \varphi_{2}$\lbl{A2}}
\BIC{$\varphi_{2}$}
\AXC{$\varphi_{2} \to \varphi_{125}$\lbl{A1}}
\BIC{$\varphi_{125}$}
\BIC{$a \in \mathbb{R} \to b \in \mathbb{R} \to \mathrm{Seq}(x) \to \varphi_{25}$}
\end{ptree}
\noindent$\Pi_{51}$:
\begin{ptree}
\AXC{$\varphi_{104} \to \varphi_{106}$\lbl{A3}}
\AXC{$\stackrel{\Pi_{37}}{\vdots}$}\noLine
\UIC{$\varphi_{104}$}
\AXC{$\stackrel{\Pi_{50}}{\vdots}$}\noLine
\UIC{$a \in \mathbb{R} \to b \in \mathbb{R} \to \mathrm{Seq}(x) \to \varphi_{25}$}
\AXC{$(a \in \mathbb{R} \to b \in \mathbb{R} \to \mathrm{Seq}(x) \to \varphi_{25}) \to \varphi_{104} \to (\varphi_{104} \to \varphi_{106}) \to \varphi_{175}$\lbl{A1}}
\BIC{$\varphi_{104} \to (\varphi_{104} \to \varphi_{106}) \to \varphi_{175}$}
\BIC{$(\varphi_{104} \to \varphi_{106}) \to \varphi_{175}$}
\BIC{$\varphi_{175}$}
\end{ptree}
\noindent 主木:
\begin{ptree}
\AXC{$\varphi_{14}$\lbl{A4}}
\AXC{$\forall b\,(\varphi_{9}) \to \varphi_{16}$\lbl{A3}}
\AXC{$\varphi_{9}$\lbl{A2}}
\UIC{$\forall b\,(\varphi_{9})$}
\AXC{$\varphi_{39} \land \neg \exists s\,(\varphi_{11}) \to \forall y\,(\varphi_{39} \land \neg \exists s\,(\varphi_{11}))$\lbl{A4}}
\AXC{$\stackrel{\Pi_{51}}{\vdots}$}\noLine
\UIC{$\varphi_{175}$}
\BIC{$\forall b\,(\varphi_{9}) \to (\forall b\,(\varphi_{9}) \to \varphi_{16}) \to \varphi_{14} \to a \in \mathbb{R} \to \varphi_{87}$}
\BIC{$(\forall b\,(\varphi_{9}) \to \varphi_{16}) \to \varphi_{14} \to a \in \mathbb{R} \to \varphi_{87}$}
\BIC{$\varphi_{14} \to a \in \mathbb{R} \to \varphi_{87}$}
\BIC{$a \in \mathbb{R} \to \varphi_{87}$}
\UIC{$\forall b\,(a \in \mathbb{R} \to \varphi_{87})$}
\UIC{$\forall a\,\forall b\,(a \in \mathbb{R} \to \varphi_{87})$}
\UIC{$\forall x\,\forall a\,\forall b\,(a \in \mathbb{R} \to \varphi_{87})$}
\end{ptree}
\begin{abbrevlist}
\item $\varphi_{1} :\equiv \forall c\,(\neg \mathrm{Conv}(y,c)) \to \forall \alpha\,\forall c\,(\neg \mathrm{Conv}(y,c))$
\item $\varphi_{2} :\equiv \mathrm{SeqBdd}(x) \to \exists b\,(\mathrm{SubSeq}(b,x) \land \mathrm{Cgt}(b))$
\item $\varphi_{3} :\equiv \forall \alpha\,\forall c\,(\neg \mathrm{Conv}(y,c)) \to \forall \alpha\,(\neg \mathrm{Conv}(y,\alpha))$
\item $\varphi_{4} :\equiv \mathrm{SubSeq}(y,x) \land \mathrm{Cgt}(y) \to \exists \alpha\,(\mathrm{Conv}(y,\alpha))$
\item $\varphi_{5} :\equiv \forall a\,(\mathrm{SeqBdd}(a) \to \exists b\,(\mathrm{SubSeq}(b,a) \land \mathrm{Cgt}(b)))$
\item $\varphi_{6} :\equiv \mathrm{SubSeq}(y,x) \land \mathrm{Cgt}(y) \to \mathrm{SubSeq}(y,x) \land \mathrm{Cgt}(y)$
\item $\varphi_{7} :\equiv \mathrm{SubSeq}(y,x) \land \mathrm{Cgt}(y) \to \mathrm{Conv}(y,c) \to \mathrm{Conv}(y,c)$
\item $\varphi_{8} :\equiv \forall v_{781}\,\forall v_{782}\,(v_{780} = v_{781} \to v_{781} = v_{782} \to v_{780} = v_{782})$
\item $\varphi_{9} :\equiv \forall y\,(\neg (\mathrm{SubSeq}(y,x) \land \mathrm{Cgt}(y))) \to \neg (\mathrm{SubSeq}(b,x) \land \mathrm{Cgt}(b))$
\item $\varphi_{10} :\equiv \neg (\mathrm{SubIdx}(s) \land c \in [a,b] \land \mathrm{Conv}(x \circ s,c))$
\item $\varphi_{11} :\equiv \exists c\,(\mathrm{SubIdx}(s) \land c \in [a,b] \land \mathrm{Conv}(x \circ s,c))$
\item $\varphi_{12} :\equiv \forall a\,(\mathrm{Cgt}(a) \leftrightarrow \exists \alpha\,(\mathrm{Conv}(a,\alpha))) \to (\mathrm{Cgt}(y) \leftrightarrow \exists \alpha\,(\mathrm{Conv}(y,\alpha)))$
\item $\varphi_{13} :\equiv \forall v_{781}\,\forall v_{782}\,(y(n) = v_{781} \to v_{781} = v_{782} \to y(n) = v_{782})$
\item $\varphi_{14} :\equiv \forall y\,(\neg (\mathrm{SubSeq}(y,x) \land \mathrm{Cgt}(y))) \to \forall b\,\forall y\,(\neg (\mathrm{SubSeq}(y,x) \land \mathrm{Cgt}(y)))$
\item $\varphi_{15} :\equiv \forall x\,\forall y\,\forall z\,(x = y \to y = z \to x = z) \to \forall y\,\forall z\,(v_{780} = y \to y = z \to v_{780} = z)$
\item $\varphi_{16} :\equiv \forall b\,\forall y\,(\neg (\mathrm{SubSeq}(y,x) \land \mathrm{Cgt}(y))) \to \forall b\,(\neg (\mathrm{SubSeq}(b,x) \land \mathrm{Cgt}(b)))$
\item $\varphi_{17} :\equiv \exists s\,(\mathrm{SubIdx}(s) \land \forall k \in \mathbb{N}\,(y(k) = x(s_{k})))$
\item $\varphi_{18} :\equiv \forall s\,(\neg (\mathrm{SubIdx}(s) \land \forall k \in \mathbb{N}\,(y(k) = x(s_{k}))))$
\item $\varphi_{19} :\equiv \mathrm{Seq}(x) \to \forall n \in \mathbb{N}\,(x(n) \in [a,b]) \to \mathrm{SeqBdd}(x)$
\item $\varphi_{20} :\equiv \forall y\,\forall z\,(v_{780} = y \to y = z \to v_{780} = z) \to \forall z\,(v_{780} = v_{781} \to v_{781} = z \to v_{780} = z)$
\item $\varphi_{21} :\equiv \forall x\,\forall y\,(x = y \to y = x) \to \forall y\,(x \circ s(n) = y \to y = x \circ s(n))$
\item $\varphi_{22} :\equiv \mathrm{SubIdx}(s) \leftrightarrow \mathrm{NSeq}(s) \land \forall k \in \mathbb{N}\,(s_{k} < s_{S(k)})$
\item $\varphi_{23} :\equiv \mathrm{SubIdx}(s) \land \forall k \in \mathbb{N}\,(y(k) = x(s_{k})) \to \mathrm{NSeq}(s)$
\item $\varphi_{24} :\equiv \mathrm{SubIdx}(s) \land \forall k \in \mathbb{N}\,(y(k) = x(s_{k})) \to \mathrm{SubIdx}(s)$
\item $\varphi_{25} :\equiv \forall n \in \mathbb{N}\,(x(n) \in [a,b]) \to \exists b\,(\mathrm{SubSeq}(b,x) \land \mathrm{Cgt}(b))$
\item $\varphi_{26} :\equiv \mathrm{Seq}(x) \to \forall n \in \mathbb{N}\,(x(n) \in [a,b]) \to \lvert b \rvert \in \mathbb{R}$
\item $\varphi_{27} :\equiv \mathrm{Seq}(x) \to \forall n \in \mathbb{N}\,(x(n) \in [a,b]) \to \lvert a \rvert \in \mathbb{R}$
\item $\varphi_{28} :\equiv \forall z\,(v_{780} = v_{781} \to v_{781} = z \to v_{780} = z) \to v_{780} = v_{781} \to v_{781} = v_{782} \to v_{780} = v_{782}$
\item $\varphi_{29} :\equiv \mathrm{SeqBdd}(x) \leftrightarrow \mathrm{Seq}(x) \land \exists M \in \mathbb{R}\,\forall n \in \mathbb{N}\,(\lvert x(n) \rvert \leq M)$
\item $\varphi_{30} :\equiv \mathrm{SubIdx}(s) \land \forall k \in \mathbb{N}\,(y(k) = x(s_{k})) \to \mathrm{Seq}(x \circ s)$
\item $\varphi_{31} :\equiv \forall \alpha\,(\forall c\,(\neg \mathrm{Conv}(y,c)) \to \neg \mathrm{Conv}(y,\alpha)) \to \varphi_{3}$
\item $\varphi_{32} :\equiv \forall v_{782}\,(y(n) = x(s_{n}) \to x(s_{n}) = v_{782} \to y(n) = v_{782})$
\item $\varphi_{33} :\equiv \mathrm{SubIdx}(s) \land \forall k \in \mathbb{N}\,(y(k) = x(s_{k})) \to c \in [a,b]$
\item $\varphi_{34} :\equiv \forall w\,(x(s_{n}) = w \to x(s_{n}) \in [a,b] \to w \in [a,b])$
\item $\varphi_{35} :\equiv \forall x\,(a \in \mathbb{R} \to b \in \mathbb{R} \to x \in [a,b] \to \lvert x \rvert \leq \lvert a \rvert + \lvert b \rvert)$
\item $\varphi_{36} :\equiv \forall a\,(\mathrm{SeqBdd}(a) \leftrightarrow \mathrm{Seq}(a) \land \exists M \in \mathbb{R}\,\forall n \in \mathbb{N}\,(\lvert a_{n} \rvert \leq M))$
\item $\varphi_{37} :\equiv \mathrm{SubIdx}(s) \land \forall k \in \mathbb{N}\,(y(k) = x(s_{k})) \to \mathrm{Conv}(x \circ s,c)$
\item $\varphi_{38} :\equiv \forall n \in \mathbb{N}\,(y(n) = x \circ s(n)) \to \mathrm{Conv}(y,c) \to \mathrm{Conv}(x \circ s,c)$
\item $\varphi_{39} :\equiv a \in \mathbb{R} \land b \in \mathbb{R} \land \mathrm{Seq}(x) \land \forall n \in \mathbb{N}\,(x(n) \in [a,b])$
\item $\varphi_{40} :\equiv \lvert a \rvert + \lvert b \rvert \in \mathbb{R} \land \forall n \in \mathbb{N}\,(\lvert x(n) \rvert \leq \lvert a \rvert + \lvert b \rvert)$
\item $\varphi_{41} :\equiv \forall b\,\forall \alpha\,(\mathrm{Seq}(b) \to \forall n \in \mathbb{N}\,(y(n) = b_{n}) \to \mathrm{Conv}(y,\alpha) \to \mathrm{Conv}(b,\alpha))$
\item $\varphi_{42} :\equiv \forall n \in \mathbb{N}\,(x(n) \in [a,b]) \to n \in \mathbb{N} \to x(n) \in [a,b]$
\item $\varphi_{43} :\equiv \mathrm{NSeq}(s) \to \mathrm{Seq}(x \circ s) \land \forall k \in \mathbb{N}\,(x \circ s(k) = x(s_{k}))$
\item $\varphi_{44} :\equiv \mathrm{Seq}(x) \to \forall n \in \mathbb{N}\,(x(n) \in [a,b]) \to \lvert a \rvert + \lvert b \rvert \in \mathbb{R}$
\item $\varphi_{45} :\equiv \mathrm{SubIdx}(s) \land \forall k \in \mathbb{N}\,(y(k) = x(s_{k})) \to n \in \mathbb{N} \to n \in \mathbb{N}$
\item $\varphi_{46} :\equiv \forall a\,\forall b\,\forall \alpha\,(\mathrm{Seq}(b) \to \forall n \in \mathbb{N}\,(a_{n} = b_{n}) \to \mathrm{Conv}(a,\alpha) \to \mathrm{Conv}(b,\alpha))$
\item $\varphi_{47} :\equiv \forall n \in \mathbb{N}\,(x(n) \in [a,b]) \to \forall n \in \mathbb{N}\,(x(n) \in [a,b])$
\item $\varphi_{48} :\equiv \forall k \in \mathbb{N}\,(y(k) = x(s_{k})) \to n \in \mathbb{N} \to y(n) = x(s_{n})$
\item $\varphi_{49} :\equiv \forall a\,(\mathrm{SubSeq}(y,a) \leftrightarrow \mathrm{Seq}(y) \land \exists s\,(\mathrm{SubIdx}(s) \land \forall k \in \mathbb{N}\,(y(k) = a_{s_{k}})))$
\item $\varphi_{50} :\equiv a \in \mathbb{R} \to b \in \mathbb{R} \to x(n) \in [a,b] \to \lvert x(n) \rvert \leq \lvert a \rvert + \lvert b \rvert$
\item $\varphi_{51} :\equiv \forall k\,(\mathrm{NSeq}(s) \to k \in \mathbb{N} \to s_{k} \in \mathbb{N}) \to \mathrm{NSeq}(s) \to n \in \mathbb{N} \to s_{n} \in \mathbb{N}$
\item $\varphi_{52} :\equiv \mathrm{SubIdx}(s) \land \forall k \in \mathbb{N}\,(y(k) = x(s_{k})) \to n \in \mathbb{N} \to s_{n} \in \mathbb{N}$
\item $\varphi_{53} :\equiv \forall n \in \mathbb{N}\,(x(n) \in [a,b]) \to n \in \mathbb{N} \to \lvert x(n) \rvert \leq \lvert a \rvert + \lvert b \rvert$
\item $\varphi_{54} :\equiv \forall b\,\forall a\,(\mathrm{SubSeq}(b,a) \leftrightarrow \mathrm{Seq}(b) \land \exists s\,(\mathrm{SubIdx}(s) \land \forall k \in \mathbb{N}\,(b_{k} = a_{s_{k}})))$
\item $\varphi_{55} :\equiv \forall n \in \mathbb{N}\,(x \circ s(n) \in [a,b]) \to \mathrm{Conv}(x \circ s,c) \to c \in [a,b]$
\item $\varphi_{56} :\equiv \forall n \in \mathbb{N}\,(x(n) \in [a,b]) \to \forall n \in \mathbb{N}\,(\lvert x(n) \rvert \leq \lvert a \rvert + \lvert b \rvert)$
\item $\varphi_{57} :\equiv \forall n \in \mathbb{N}\,(x(n) \in [a,b]) \to \exists M \in \mathbb{R}\,\forall n \in \mathbb{N}\,(\lvert x(n) \rvert \leq M)$
\item $\varphi_{58} :\equiv y(n) = x(s_{n}) \to x(s_{n}) = x \circ s(n) \to y(n) = x \circ s(n)$
\item $\varphi_{59} :\equiv \forall s\,\forall k\,(\mathrm{NSeq}(s) \to k \in \mathbb{N} \to s_{k} \in \mathbb{N}) \to \forall k\,(\mathrm{NSeq}(s) \to k \in \mathbb{N} \to s_{k} \in \mathbb{N})$
\item $\varphi_{60} :\equiv \forall n \in \mathbb{N}\,(x(n) \in [a,b]) \to s_{n} \in \mathbb{N} \to x(s_{n}) \in [a,b]$
\item $\varphi_{61} :\equiv \mathrm{SubSeq}(y,x) \land \mathrm{Cgt}(y) \to \mathrm{Conv}(y,c) \to \exists s\,(\varphi_{11})$
\item $\varphi_{62} :\equiv \forall a\,\forall s\,(\mathrm{Seq}(a) \to \mathrm{NSeq}(s) \to \mathrm{Seq}(a \circ s) \land \forall k \in \mathbb{N}\,(a \circ s(k) = a_{s_{k}}))$
\item $\varphi_{63} :\equiv \mathrm{SubIdx}(s) \land \forall k \in \mathbb{N}\,(y(k) = x(s_{k})) \to n \in \mathbb{N} \to y(n) = x(s_{n})$
\item $\varphi_{64} :\equiv \mathrm{SubIdx}(s) \land \forall k \in \mathbb{N}\,(y(k) = x(s_{k})) \to n \in \mathbb{N} \to y(n) = x \circ s(n)$
\item $\varphi_{65} :\equiv \forall k \in \mathbb{N}\,(x \circ s(k) = x(s_{k})) \to n \in \mathbb{N} \to x \circ s(n) = x(s_{n})$
\item $\varphi_{66} :\equiv \mathrm{SubIdx}(s) \land \forall k \in \mathbb{N}\,(y(k) = x(s_{k})) \to \forall n \in \mathbb{N}\,(y(n) = x \circ s(n))$
\item $\varphi_{67} :\equiv \forall \alpha\,(\mathrm{Seq}(x \circ s) \to \forall n \in \mathbb{N}\,(y(n) = x \circ s(n)) \to \mathrm{Conv}(y,\alpha) \to \mathrm{Conv}(x \circ s,\alpha))$
\item $\varphi_{68} :\equiv \mathrm{SubIdx}(s) \land \forall k \in \mathbb{N}\,(y(k) = x(s_{k})) \to n \in \mathbb{N} \to x(s_{n}) \in [a,b]$
\item $\varphi_{69} :\equiv \mathrm{SubIdx}(s) \land \forall k \in \mathbb{N}\,(y(k) = x(s_{k})) \to n \in \mathbb{N} \to x \circ s(n) \in [a,b]$
\item $\varphi_{70} :\equiv x(s_{n}) = x \circ s(n) \to x(s_{n}) \in [a,b] \to x \circ s(n) \in [a,b]$
\item $\varphi_{71} :\equiv \mathrm{SubIdx}(s) \land \forall k \in \mathbb{N}\,(y(k) = x(s_{k})) \to n \in \mathbb{N} \to x \circ s(n) = x(s_{n})$
\item $\varphi_{72} :\equiv \mathrm{SubIdx}(s) \land \forall k \in \mathbb{N}\,(y(k) = x(s_{k})) \to n \in \mathbb{N} \to x(s_{n}) = x \circ s(n)$
\item $\varphi_{73} :\equiv \mathrm{SubIdx}(s) \land \forall k \in \mathbb{N}\,(y(k) = x(s_{k})) \to \forall n \in \mathbb{N}\,(x \circ s(n) \in [a,b])$
\item $\varphi_{74} :\equiv b \in \mathbb{R} \to \mathrm{Seq}(x) \to \forall n \in \mathbb{N}\,(x(n) \in [a,b]) \to \varphi_{4}$
\item $\varphi_{75} :\equiv \mathrm{SubIdx}(s) \land \forall k \in \mathbb{N}\,(y(k) = x(s_{k})) \to \mathrm{SubIdx}(s) \land c \in [a,b] \land \mathrm{Conv}(x \circ s,c)$
\item $\varphi_{76} :\equiv \mathrm{SubIdx}(s) \land \forall k \in \mathbb{N}\,(y(k) = x(s_{k})) \to n \in \mathbb{N} \to n \in \mathbb{N} \to y(n) = x(s_{n})$
\item $\varphi_{77} :\equiv \forall x\,\forall y\,(x \in \mathbb{R} \to y \in \mathbb{R} \to x + y \in \mathbb{R}) \to \forall y\,(\lvert a \rvert \in \mathbb{R} \to y \in \mathbb{R} \to \lvert a \rvert + y \in \mathbb{R})$
\item $\varphi_{78} :\equiv \forall x\,\forall c\,\forall a\,\forall b\,(a \in \mathbb{R} \to b \in \mathbb{R} \to \forall n \in \mathbb{N}\,(x(n) \in [a,b]) \to \mathrm{Conv}(x,c) \to c \in [a,b])$
\item $\varphi_{79} :\equiv \mathrm{SubIdx}(s) \land \forall k \in \mathbb{N}\,(y(k) = x(s_{k})) \to \mathrm{SubIdx}(s) \land \forall k \in \mathbb{N}\,(y(k) = x(s_{k}))$
\item $\varphi_{80} :\equiv \mathrm{SubIdx}(s) \land \forall k \in \mathbb{N}\,(y(k) = x(s_{k})) \to n \in \mathbb{N} \to \forall k \in \mathbb{N}\,(y(k) = x(s_{k}))$
\item $\varphi_{81} :\equiv b \in \mathbb{R} \to \mathrm{Seq}(x) \to \forall n \in \mathbb{N}\,(x(n) \in [a,b]) \to \varphi_{6}$
\item $\varphi_{82} :\equiv b \in \mathbb{R} \to \mathrm{Seq}(x) \to \forall n \in \mathbb{N}\,(x(n) \in [a,b]) \to \varphi_{7}$
\item $\varphi_{83} :\equiv \mathrm{SubIdx}(s) \land \forall k \in \mathbb{N}\,(y(k) = x(s_{k})) \to \exists s\,(\varphi_{11})$
\item $\varphi_{84} :\equiv \mathrm{SubIdx}(s) \land \forall k \in \mathbb{N}\,(y(k) = x(s_{k})) \to n \in \mathbb{N} \to n \in \mathbb{N} \to x \circ s(n) = x(s_{n})$
\item $\varphi_{85} :\equiv \forall y\,(x \circ s(n) = y \to y = x \circ s(n)) \to x \circ s(n) = x(s_{n}) \to x(s_{n}) = x \circ s(n)$
\item $\varphi_{86} :\equiv \mathrm{SubIdx}(s) \land \forall k \in \mathbb{N}\,(y(k) = x(s_{k})) \to n \in \mathbb{N} \to \forall k \in \mathbb{N}\,(x \circ s(k) = x(s_{k}))$
\item $\varphi_{87} :\equiv b \in \mathbb{R} \to \mathrm{Seq}(x) \to \forall n \in \mathbb{N}\,(x(n) \in [a,b]) \to \exists s\,(\varphi_{11})$
\item $\varphi_{88} :\equiv \mathrm{Conv}(y,c) \to \mathrm{SubIdx}(s) \land \forall k \in \mathbb{N}\,(y(k) = x(s_{k})) \to \varphi_{11}$
\item $\varphi_{89} :\equiv \forall n \in \mathbb{N}\,(x(n) \in [a,b]) \to \mathrm{SubSeq}(y,x) \land \mathrm{Cgt}(y) \to \exists s\,(\varphi_{11})$
\item $\varphi_{90} :\equiv \forall y\,(\lvert a \rvert \in \mathbb{R} \to y \in \mathbb{R} \to \lvert a \rvert + y \in \mathbb{R}) \to \lvert a \rvert \in \mathbb{R} \to \lvert b \rvert \in \mathbb{R} \to \lvert a \rvert + \lvert b \rvert \in \mathbb{R}$
\item $\varphi_{91} :\equiv \mathrm{SubIdx}(s) \land \forall k \in \mathbb{N}\,(y(k) = x(s_{k})) \to \mathrm{Seq}(x \circ s) \land \forall k \in \mathbb{N}\,(x \circ s(k) = x(s_{k}))$
\item $\varphi_{92} :\equiv \forall n \in \mathbb{N}\,(x(n) \in [a,b]) \to \mathrm{SubSeq}(y,x) \land \mathrm{Cgt}(y) \to \mathrm{Seq}(y) \land \varphi_{17}$
\item $\varphi_{93} :\equiv \forall n \in \mathbb{N}\,(x(n) \in [a,b]) \to \mathrm{SubSeq}(y,x) \land \mathrm{Cgt}(y) \to \mathrm{Conv}(y,c) \to \varphi_{17}$
\item $\varphi_{94} :\equiv b \in \mathbb{R} \to \mathrm{Seq}(x) \to \forall n \in \mathbb{N}\,(x(n) \in [a,b]) \to \varphi_{40}$
\item $\varphi_{95} :\equiv \forall n \in \mathbb{N}\,(x(n) \in [a,b]) \to \mathrm{SubSeq}(y,x) \land \mathrm{Cgt}(y) \to \mathrm{Conv}(y,c) \to \varphi_{23}$
\item $\varphi_{96} :\equiv \forall n \in \mathbb{N}\,(x(n) \in [a,b]) \to \mathrm{SubSeq}(y,x) \land \mathrm{Cgt}(y) \to \mathrm{Conv}(y,c) \to \varphi_{24}$
\item $\varphi_{97} :\equiv \forall n \in \mathbb{N}\,(x(n) \in [a,b]) \to \mathrm{SubSeq}(y,x) \land \mathrm{Cgt}(y) \to \mathrm{Conv}(y,c) \to \varphi_{30}$
\item $\varphi_{98} :\equiv \forall n \in \mathbb{N}\,(x(n) \in [a,b]) \to \mathrm{SubSeq}(y,x) \land \mathrm{Cgt}(y) \to \mathrm{Conv}(y,c) \to \varphi_{33}$
\item $\varphi_{99} :\equiv \forall n \in \mathbb{N}\,(x(n) \in [a,b]) \to \mathrm{SubSeq}(y,x) \land \mathrm{Cgt}(y) \to \mathrm{Conv}(y,c) \to \varphi_{37}$
\item $\varphi_{100} :\equiv (a \in \mathbb{R} \to \lvert a \rvert \in \mathbb{R}) \to (a \in \mathbb{R} \to a \in \mathbb{R}) \to a \in \mathbb{R} \to b \in \mathbb{R} \to \varphi_{27}$
\item $\varphi_{101} :\equiv \forall n \in \mathbb{N}\,(x(n) \in [a,b]) \to \mathrm{SubSeq}(y,x) \land \mathrm{Cgt}(y) \to \mathrm{Conv}(y,c) \to \varphi_{45}$
\item $\varphi_{102} :\equiv b \in \mathbb{R} \to \mathrm{Seq}(x) \to \forall n \in \mathbb{N}\,(x(n) \in [a,b]) \to \varphi_{61}$
\item $\varphi_{103} :\equiv \forall c\,(\varphi_{39} \land (\mathrm{SubSeq}(y,x) \land \mathrm{Cgt}(y)) \land \neg \exists s\,(\varphi_{11}))$
\item $\varphi_{104} :\equiv \forall y\,(\varphi_{39} \land \neg \exists s\,(\varphi_{11}) \to \neg (\mathrm{SubSeq}(y,x) \land \mathrm{Cgt}(y)))$
\item $\varphi_{105} :\equiv \forall n \in \mathbb{N}\,(x(n) \in [a,b]) \to \mathrm{SubSeq}(y,x) \land \mathrm{Cgt}(y) \to \mathrm{Conv}(y,c) \to \varphi_{52}$
\item $\varphi_{106} :\equiv \forall y\,(\varphi_{39} \land \neg \exists s\,(\varphi_{11})) \to \forall y\,(\neg (\mathrm{SubSeq}(y,x) \land \mathrm{Cgt}(y)))$
\item $\varphi_{107} :\equiv (b \in \mathbb{R} \to \lvert b \rvert \in \mathbb{R}) \to (a \in \mathbb{R} \to b \in \mathbb{R} \to b \in \mathbb{R}) \to a \in \mathbb{R} \to b \in \mathbb{R} \to \varphi_{26}$
\item $\varphi_{108} :\equiv \varphi_{39} \land (\mathrm{SubSeq}(y,x) \land \mathrm{Cgt}(y)) \land \mathrm{Conv}(y,c) \land (\mathrm{SubIdx}(s) \land \forall k \in \mathbb{N}\,(y(k) = x(s_{k})))$
\item $\varphi_{109} :\equiv \varphi_{39} \land (\mathrm{SubSeq}(y,x) \land \mathrm{Cgt}(y)) \land \mathrm{Conv}(y,c) \land \neg \exists s\,(\varphi_{11})$
\item $\varphi_{110} :\equiv \varphi_{39} \land (\mathrm{SubSeq}(y,x) \land \mathrm{Cgt}(y)) \land \neg \exists s\,(\varphi_{11}) \to \neg \mathrm{Conv}(y,c)$
\item $\varphi_{111} :\equiv \forall n \in \mathbb{N}\,(x(n) \in [a,b]) \to \mathrm{SubSeq}(y,x) \land \mathrm{Cgt}(y) \to \mathrm{Conv}(y,c) \to \varphi_{63}$
\item $\varphi_{112} :\equiv \forall n \in \mathbb{N}\,(x(n) \in [a,b]) \to \mathrm{SubSeq}(y,x) \land \mathrm{Cgt}(y) \to \mathrm{Conv}(y,c) \to \varphi_{64}$
\item $\varphi_{113} :\equiv \forall n \in \mathbb{N}\,(x(n) \in [a,b]) \to \mathrm{SubSeq}(y,x) \land \mathrm{Cgt}(y) \to \mathrm{Conv}(y,c) \to \varphi_{66}$
\item $\varphi_{114} :\equiv \forall n \in \mathbb{N}\,(x(n) \in [a,b]) \to \mathrm{SubSeq}(y,x) \land \mathrm{Cgt}(y) \to \mathrm{Conv}(y,c) \to \varphi_{68}$
\item $\varphi_{115} :\equiv \forall n \in \mathbb{N}\,(x(n) \in [a,b]) \to \mathrm{SubSeq}(y,x) \land \mathrm{Cgt}(y) \to \mathrm{Conv}(y,c) \to \varphi_{69}$
\item $\varphi_{116} :\equiv \forall n \in \mathbb{N}\,(x(n) \in [a,b]) \to \mathrm{SubSeq}(y,x) \land \mathrm{Cgt}(y) \to \mathrm{Conv}(y,c) \to \varphi_{71}$
\item $\varphi_{117} :\equiv \forall n \in \mathbb{N}\,(x(n) \in [a,b]) \to \mathrm{SubSeq}(y,x) \land \mathrm{Cgt}(y) \to \mathrm{Conv}(y,c) \to \varphi_{72}$
\item $\varphi_{118} :\equiv \forall n \in \mathbb{N}\,(x(n) \in [a,b]) \to \mathrm{SubSeq}(y,x) \land \mathrm{Cgt}(y) \to \mathrm{Conv}(y,c) \to \varphi_{73}$
\item $\varphi_{119} :\equiv \forall n \in \mathbb{N}\,(x(n) \in [a,b]) \to \mathrm{SubSeq}(y,x) \land \mathrm{Cgt}(y) \to \mathrm{Conv}(y,c) \to \varphi_{75}$
\item $\varphi_{120} :\equiv \forall n \in \mathbb{N}\,(x(n) \in [a,b]) \to \mathrm{SubSeq}(y,x) \land \mathrm{Cgt}(y) \to \mathrm{Conv}(y,c) \to \varphi_{76}$
\item $\varphi_{121} :\equiv \forall n \in \mathbb{N}\,(x(n) \in [a,b]) \to \mathrm{SubSeq}(y,x) \land \mathrm{Cgt}(y) \to \varphi_{88}$
\item $\varphi_{122} :\equiv \forall n \in \mathbb{N}\,(x(n) \in [a,b]) \to \mathrm{SubSeq}(y,x) \land \mathrm{Cgt}(y) \to \mathrm{Conv}(y,c) \to \varphi_{79}$
\item $\varphi_{123} :\equiv \forall n \in \mathbb{N}\,(x(n) \in [a,b]) \to \mathrm{SubSeq}(y,x) \land \mathrm{Cgt}(y) \to \mathrm{Conv}(y,c) \to \varphi_{80}$
\item $\varphi_{124} :\equiv \forall n \in \mathbb{N}\,(x(n) \in [a,b]) \to \mathrm{SubSeq}(y,x) \land \mathrm{Cgt}(y) \to \mathrm{Conv}(y,c) \to \varphi_{83}$
\item $\varphi_{125} :\equiv (a \in \mathbb{R} \to b \in \mathbb{R} \to \varphi_{19}) \to a \in \mathbb{R} \to b \in \mathbb{R} \to \mathrm{Seq}(x) \to \varphi_{25}$
\item $\varphi_{126} :\equiv \forall n \in \mathbb{N}\,(x(n) \in [a,b]) \to \mathrm{SubSeq}(y,x) \land \mathrm{Cgt}(y) \to \mathrm{Conv}(y,c) \to \varphi_{84}$
\item $\varphi_{127} :\equiv \forall n \in \mathbb{N}\,(x(n) \in [a,b]) \to \mathrm{SubSeq}(y,x) \land \mathrm{Cgt}(y) \to \mathrm{Conv}(y,c) \to \varphi_{86}$
\item $\varphi_{128} :\equiv \forall n \in \mathbb{N}\,(x(n) \in [a,b]) \to \mathrm{SubSeq}(y,x) \land \mathrm{Cgt}(y) \to \mathrm{Conv}(y,c) \to \varphi_{91}$
\item $\varphi_{129} :\equiv (a \in \mathbb{R} \to b \in \mathbb{R} \to \mathrm{Seq}(x) \to \varphi_{57}) \to a \in \mathbb{R} \to b \in \mathbb{R} \to \varphi_{19}$
\item $\varphi_{130} :\equiv \varphi_{109} \to \neg (\mathrm{SubIdx}(s) \land \forall k \in \mathbb{N}\,(y(k) = x(s_{k})))$
\item $\varphi_{131} :\equiv (a \in \mathbb{R} \to b \in \mathbb{R} \to \mathrm{Seq}(x) \to \varphi_{42}) \to a \in \mathbb{R} \to b \in \mathbb{R} \to \mathrm{Seq}(x) \to \varphi_{53}$
\item $\varphi_{132} :\equiv \forall a\,\forall b\,(a \in \mathbb{R} \to b \in \mathbb{R} \to \varphi_{55}) \to \forall b\,(a \in \mathbb{R} \to b \in \mathbb{R} \to \varphi_{55})$
\item $\varphi_{133} :\equiv \forall c\,\forall a\,\forall b\,(a \in \mathbb{R} \to b \in \mathbb{R} \to \varphi_{55}) \to \forall a\,\forall b\,(a \in \mathbb{R} \to b \in \mathbb{R} \to \varphi_{55})$
\item $\varphi_{134} :\equiv (a \in \mathbb{R} \to b \in \mathbb{R} \to \mathrm{Seq}(x) \to \varphi_{53}) \to \varphi_{39} \to n \in \mathbb{N} \to \lvert x(n) \rvert \leq \lvert a \rvert + \lvert b \rvert$
\item $\varphi_{135} :\equiv (\mathrm{SubSeq}(y,x) \leftrightarrow \mathrm{Seq}(y) \land \varphi_{17}) \to a \in \mathbb{R} \to b \in \mathbb{R} \to \mathrm{Seq}(x) \to \varphi_{92}$
\item $\varphi_{136} :\equiv \forall n\,(\varphi_{39} \to n \in \mathbb{N} \to \lvert x(n) \rvert \leq \lvert a \rvert + \lvert b \rvert) \to \forall n\,(\varphi_{39}) \to \forall n \in \mathbb{N}\,(\lvert x(n) \rvert \leq \lvert a \rvert + \lvert b \rvert)$
\item $\varphi_{137} :\equiv (\forall M\,(\neg (M \in \mathbb{R} \land \forall n \in \mathbb{N}\,(\lvert x(n) \rvert \leq M))) \to \neg \varphi_{40}) \to a \in \mathbb{R} \to b \in \mathbb{R} \to \mathrm{Seq}(x) \to \varphi_{57}$
\item $\varphi_{138} :\equiv (a \in \mathbb{R} \to \varphi_{81}) \to (\mathrm{Cgt}(y) \leftrightarrow \exists \alpha\,(\mathrm{Conv}(y,\alpha))) \to a \in \mathbb{R} \to \varphi_{74}$
\item $\varphi_{139} :\equiv (x \circ s(n) = x(s_{n}) \to x(s_{n}) = x \circ s(n)) \to a \in \mathbb{R} \to b \in \mathbb{R} \to \mathrm{Seq}(x) \to \varphi_{117}$
\item $\varphi_{140} :\equiv (a \in \mathbb{R} \to b \in \mathbb{R} \to \varphi_{27}) \to (a \in \mathbb{R} \to b \in \mathbb{R} \to \varphi_{26}) \to a \in \mathbb{R} \to b \in \mathbb{R} \to \varphi_{44}$
\item $\varphi_{141} :\equiv \forall \alpha\,(\forall c\,(\neg \mathrm{Conv}(y,c)) \to \neg \mathrm{Conv}(y,\alpha)) \to \varphi_{31} \to \varphi_{1} \to a \in \mathbb{R} \to b \in \mathbb{R} \to \mathrm{Seq}(x) \to \varphi_{89}$
\item $\varphi_{142} :\equiv (a \in \mathbb{R} \to b \in \mathbb{R} \to \mathrm{Seq}(x) \to \varphi_{47}) \to \varphi_{42} \to a \in \mathbb{R} \to b \in \mathbb{R} \to \mathrm{Seq}(x) \to \varphi_{42}$
\item $\varphi_{143} :\equiv (a \in \mathbb{R} \to b \in \mathbb{R} \to \mathrm{Seq}(x) \to \varphi_{89}) \to \varphi_{39} \land \neg \exists s\,(\varphi_{11}) \to \neg (\mathrm{SubSeq}(y,x) \land \mathrm{Cgt}(y))$
\item $\varphi_{144} :\equiv \varphi_{39} \land (\mathrm{SubSeq}(y,x) \land \mathrm{Cgt}(y)) \land \neg \exists s\,(\varphi_{11}) \to \varphi_{103}$
\item $\varphi_{145} :\equiv (a \in \mathbb{R} \to b \in \mathbb{R} \to \mathrm{Seq}(x) \to \varphi_{92}) \to a \in \mathbb{R} \to b \in \mathbb{R} \to \mathrm{Seq}(x) \to \varphi_{93}$
\item $\varphi_{146} :\equiv (a \in \mathbb{R} \to b \in \mathbb{R} \to \mathrm{Seq}(x) \to \varphi_{47}) \to \varphi_{60} \to a \in \mathbb{R} \to b \in \mathbb{R} \to \mathrm{Seq}(x) \to \varphi_{60}$
\item $\varphi_{147} :\equiv (a \in \mathbb{R} \to b \in \mathbb{R} \to \mathrm{Seq}(x) \to \varphi_{99}) \to a \in \mathbb{R} \to b \in \mathbb{R} \to \mathrm{Seq}(x) \to \varphi_{98}$
\item $\varphi_{148} :\equiv (a \in \mathbb{R} \to b \in \mathbb{R} \to \varphi_{44}) \to (a \in \mathbb{R} \to b \in \mathbb{R} \to \mathrm{Seq}(x) \to \varphi_{56}) \to a \in \mathbb{R} \to \varphi_{94}$
\item $\varphi_{149} :\equiv \varphi_{50} \to (a \in \mathbb{R} \to a \in \mathbb{R}) \to (a \in \mathbb{R} \to b \in \mathbb{R} \to b \in \mathbb{R}) \to \varphi_{131}$
\item $\varphi_{150} :\equiv (a \in \mathbb{R} \to b \in \mathbb{R} \to \mathrm{Seq}(x) \to \varphi_{101}) \to a \in \mathbb{R} \to b \in \mathbb{R} \to \mathrm{Seq}(x) \to \varphi_{105}$
\item $\varphi_{151} :\equiv (a \in \mathbb{R} \to b \in \mathbb{R} \to \mathrm{Seq}(x) \to \varphi_{101}) \to a \in \mathbb{R} \to b \in \mathbb{R} \to \mathrm{Seq}(x) \to \varphi_{111}$
\item $\varphi_{152} :\equiv (a \in \mathbb{R} \to b \in \mathbb{R} \to \mathrm{Seq}(x) \to \varphi_{122}) \to a \in \mathbb{R} \to b \in \mathbb{R} \to \mathrm{Seq}(x) \to \varphi_{96}$
\item $\varphi_{153} :\equiv (a \in \mathbb{R} \to b \in \mathbb{R} \to \mathrm{Seq}(x) \to \varphi_{99}) \to a \in \mathbb{R} \to b \in \mathbb{R} \to \mathrm{Seq}(x) \to \varphi_{119}$
\item $\varphi_{154} :\equiv (a \in \mathbb{R} \to b \in \mathbb{R} \to \mathrm{Seq}(x) \to \varphi_{101}) \to a \in \mathbb{R} \to b \in \mathbb{R} \to \mathrm{Seq}(x) \to \varphi_{116}$
\item $\varphi_{155} :\equiv (a \in \mathbb{R} \to b \in \mathbb{R} \to \mathrm{Seq}(x) \to \varphi_{105}) \to a \in \mathbb{R} \to b \in \mathbb{R} \to \mathrm{Seq}(x) \to \varphi_{114}$
\item $\varphi_{156} :\equiv (a \in \mathbb{R} \to b \in \mathbb{R} \to \mathrm{Seq}(x) \to \varphi_{95}) \to a \in \mathbb{R} \to b \in \mathbb{R} \to \mathrm{Seq}(x) \to \varphi_{128}$
\item $\varphi_{157} :\equiv (a \in \mathbb{R} \to b \in \mathbb{R} \to \mathrm{Seq}(x) \to \varphi_{112}) \to \varphi_{108} \to n \in \mathbb{N} \to y(n) = x \circ s(n)$
\item $\varphi_{158} :\equiv (a \in \mathbb{R} \to b \in \mathbb{R} \to \mathrm{Seq}(x) \to \varphi_{128}) \to a \in \mathbb{R} \to b \in \mathbb{R} \to \mathrm{Seq}(x) \to \varphi_{97}$
\item $\varphi_{159} :\equiv (a \in \mathbb{R} \to b \in \mathbb{R} \to \mathrm{Seq}(x) \to \varphi_{96}) \to \varphi_{22} \to a \in \mathbb{R} \to b \in \mathbb{R} \to \mathrm{Seq}(x) \to \varphi_{95}$
\item $\varphi_{160} :\equiv \forall n\,(\varphi_{108} \to n \in \mathbb{N} \to y(n) = x \circ s(n)) \to \forall n\,(\varphi_{108}) \to \forall n \in \mathbb{N}\,(y(n) = x \circ s(n))$
\item $\varphi_{161} :\equiv (a \in \mathbb{R} \to b \in \mathbb{R} \to \mathrm{Seq}(x) \to \varphi_{115}) \to \varphi_{108} \to n \in \mathbb{N} \to x \circ s(n) \in [a,b]$
\item $\varphi_{162} :\equiv \forall n\,(\varphi_{108} \to n \in \mathbb{N} \to x \circ s(n) \in [a,b]) \to \forall n\,(\varphi_{108}) \to \forall n \in \mathbb{N}\,(x \circ s(n) \in [a,b])$
\item $\varphi_{163} :\equiv (a \in \mathbb{R} \to b \in \mathbb{R} \to \mathrm{Seq}(x) \to \varphi_{122}) \to a \in \mathbb{R} \to b \in \mathbb{R} \to \mathrm{Seq}(x) \to \varphi_{123}$
\item $\varphi_{164} :\equiv (a \in \mathbb{R} \to b \in \mathbb{R} \to \mathrm{Seq}(x) \to \varphi_{128}) \to a \in \mathbb{R} \to b \in \mathbb{R} \to \mathrm{Seq}(x) \to \varphi_{127}$
\item $\varphi_{165} :\equiv (a \in \mathbb{R} \to b \in \mathbb{R} \to \mathrm{Seq}(x) \to \varphi_{117}) \to \varphi_{58} \to a \in \mathbb{R} \to b \in \mathbb{R} \to \mathrm{Seq}(x) \to \varphi_{112}$
\item $\varphi_{166} :\equiv (a \in \mathbb{R} \to b \in \mathbb{R} \to \mathrm{Seq}(x) \to \varphi_{113}) \to (a \in \mathbb{R} \to \varphi_{82}) \to a \in \mathbb{R} \to b \in \mathbb{R} \to \mathrm{Seq}(x) \to \varphi_{99}$
\item $\varphi_{167} :\equiv (a \in \mathbb{R} \to b \in \mathbb{R} \to \mathrm{Seq}(x) \to \varphi_{123}) \to \varphi_{48} \to a \in \mathbb{R} \to b \in \mathbb{R} \to \mathrm{Seq}(x) \to \varphi_{120}$
\item $\varphi_{168} :\equiv (a \in \mathbb{R} \to b \in \mathbb{R} \to \mathrm{Seq}(x) \to \varphi_{114}) \to \varphi_{70} \to a \in \mathbb{R} \to b \in \mathbb{R} \to \mathrm{Seq}(x) \to \varphi_{115}$
\item $\varphi_{169} :\equiv (\mathrm{Seq}(x) \to \varphi_{43}) \to (a \in \mathbb{R} \to b \in \mathbb{R} \to \mathrm{Seq}(x) \to \mathrm{Seq}(x)) \to \varphi_{156}$
\item $\varphi_{170} :\equiv (a \in \mathbb{R} \to b \in \mathbb{R} \to \mathrm{Seq}(x) \to \varphi_{119}) \to (\forall c\,(\varphi_{10}) \to \varphi_{10}) \to a \in \mathbb{R} \to b \in \mathbb{R} \to \mathrm{Seq}(x) \to \varphi_{121}$
\item $\varphi_{171} :\equiv (a \in \mathbb{R} \to b \in \mathbb{R} \to \mathrm{Seq}(x) \to \varphi_{127}) \to \varphi_{65} \to a \in \mathbb{R} \to b \in \mathbb{R} \to \mathrm{Seq}(x) \to \varphi_{126}$
\item $\varphi_{172} :\equiv (a \in \mathbb{R} \to b \in \mathbb{R} \to \mathrm{Seq}(x) \to \varphi_{121}) \to (\forall s\,(\neg \varphi_{11}) \to \neg \varphi_{11}) \to a \in \mathbb{R} \to b \in \mathbb{R} \to \mathrm{Seq}(x) \to \varphi_{124}$
\item $\varphi_{173} :\equiv (\mathrm{NSeq}(s) \to n \in \mathbb{N} \to s_{n} \in \mathbb{N}) \to (a \in \mathbb{R} \to b \in \mathbb{R} \to \mathrm{Seq}(x) \to \varphi_{95}) \to \varphi_{150}$
\item $\varphi_{174} :\equiv (a \in \mathbb{R} \to b \in \mathbb{R} \to b \in \mathbb{R}) \to (a \in \mathbb{R} \to b \in \mathbb{R} \to \mathrm{Seq}(x) \to \varphi_{118}) \to \varphi_{147}$
\item $\varphi_{175} :\equiv (\varphi_{39} \land \neg \exists s\,(\varphi_{11}) \to \forall y\,(\varphi_{39} \land \neg \exists s\,(\varphi_{11}))) \to \forall b\,(\varphi_{9}) \to (\forall b\,(\varphi_{9}) \to \varphi_{16}) \to \varphi_{14} \to a \in \mathbb{R} \to \varphi_{87}$
\item $\varphi_{176} :\equiv \forall n\,(\varphi_{39} \to n \in \mathbb{N} \to \lvert x(n) \rvert \leq \lvert a \rvert + \lvert b \rvert) \to \varphi_{136} \to (\varphi_{39} \to \forall n\,(\varphi_{39})) \to a \in \mathbb{R} \to b \in \mathbb{R} \to \mathrm{Seq}(x) \to \varphi_{56}$
\item $\varphi_{177} :\equiv (a \in \mathbb{R} \to b \in \mathbb{R} \to \mathrm{Seq}(x) \to \varphi_{96}) \to (a \in \mathbb{R} \to b \in \mathbb{R} \to \mathrm{Seq}(x) \to \varphi_{98}) \to \varphi_{153}$
\item $\varphi_{178} :\equiv (\mathrm{Seq}(x \circ s) \to \varphi_{38}) \to (a \in \mathbb{R} \to b \in \mathbb{R} \to \mathrm{Seq}(x) \to \varphi_{97}) \to \varphi_{166}$
\item $\varphi_{179} :\equiv (\forall s\,(\varphi_{130}) \to \forall s\,(\varphi_{109}) \to \varphi_{18}) \to (\varphi_{109} \to \forall s\,(\varphi_{109})) \to a \in \mathbb{R} \to \varphi_{102}$
\item $\varphi_{180} :\equiv \forall c\,(\varphi_{110}) \to (\forall c\,(\varphi_{110}) \to \varphi_{103} \to \forall c\,(\neg \mathrm{Conv}(y,c))) \to \varphi_{144} \to \varphi_{141}$
\item $\varphi_{181} :\equiv \forall n\,(\varphi_{108} \to n \in \mathbb{N} \to y(n) = x \circ s(n)) \to \varphi_{160} \to (\varphi_{108} \to \forall n\,(\varphi_{108})) \to a \in \mathbb{R} \to b \in \mathbb{R} \to \mathrm{Seq}(x) \to \varphi_{113}$
\item $\varphi_{182} :\equiv \forall n\,(\varphi_{108} \to n \in \mathbb{N} \to x \circ s(n) \in [a,b]) \to \varphi_{162} \to (\varphi_{108} \to \forall n\,(\varphi_{108})) \to a \in \mathbb{R} \to b \in \mathbb{R} \to \mathrm{Seq}(x) \to \varphi_{118}$
\end{abbrevlist}
\end{fullproof}

% 自動生成 (prover/build.js): SubLtTwo  ノード数 106
\begin{fullproof}{[E44]}{$\forall x\,\forall y\,\forall e\,(x \in \mathbb{R} \to y \in \mathbb{R} \to e \in \mathbb{R} \to x - e < y + e \to x - y < e + e)$}{106}\phantomsection\label{tree:SubLtTwo}
\noindent$\Pi_{1}$:
\begin{ptree}
\AXC{$z = w \to x < z \to x < w$\lbl{A6}}
\UIC{$\forall w\,(z = w \to x < z \to x < w)$}
\UIC{$\forall z\,\forall w\,(z = w \to x < z \to x < w)$}
\AXC{$\varphi_{8}$\lbl{A2}}
\BIC{$\forall w\,(y + e + e = w \to x < y + e + e \to x < w)$}
\AXC{$\forall w\,(y + e + e = w \to x < y + e + e \to x < w) \to \varphi_{2}$\lbl{A2}}
\BIC{$\varphi_{2}$}
\end{ptree}
\noindent$\Pi_{2}$:
\begin{ptree}
\AXC{$\forall x\,\forall y\,(x \in \mathbb{R} \to y \in \mathbb{R} \to x + y \in \mathbb{R})$\lbl{(R1)}}
\AXC{$\varphi_{25}$\lbl{A2}}
\BIC{$\forall y\,(v_{783} \in \mathbb{R} \to y \in \mathbb{R} \to v_{783} + y \in \mathbb{R})$}
\AXC{$\varphi_{26}$\lbl{A2}}
\BIC{$v_{783} \in \mathbb{R} \to v_{784} \in \mathbb{R} \to v_{783} + v_{784} \in \mathbb{R}$}
\UIC{$\forall v_{784}\,(v_{783} \in \mathbb{R} \to v_{784} \in \mathbb{R} \to v_{783} + v_{784} \in \mathbb{R})$}
\UIC{$\varphi_{3}$}
\AXC{$\varphi_{3} \to \forall v_{784}\,(y \in \mathbb{R} \to v_{784} \in \mathbb{R} \to y + v_{784} \in \mathbb{R})$\lbl{A2}}
\BIC{$\forall v_{784}\,(y \in \mathbb{R} \to v_{784} \in \mathbb{R} \to y + v_{784} \in \mathbb{R})$}
\end{ptree}
\noindent$\Pi_{3}$:
\begin{ptree}
\AXC{$\stackrel{\Pi_{2}}{\vdots}$}\noLine
\UIC{$\forall v_{784}\,(y \in \mathbb{R} \to v_{784} \in \mathbb{R} \to y + v_{784} \in \mathbb{R})$}
\AXC{$\varphi_{24}$\lbl{A2}}
\BIC{$y \in \mathbb{R} \to e \in \mathbb{R} \to y + e \in \mathbb{R}$}
\AXC{$(y \in \mathbb{R} \to e \in \mathbb{R} \to y + e \in \mathbb{R}) \to \varphi_{35}$\lbl{A1}}
\BIC{$\varphi_{35}$}
\end{ptree}
\noindent$\Pi_{4}$:
\begin{ptree}
\AXC{$x \in \mathbb{R} \to y \in \mathbb{R} \to e \in \mathbb{R} \to x - e < y + e \to e \in \mathbb{R}$\lbl{A1}}
\AXC{$x \in \mathbb{R} \to y \in \mathbb{R} \to e \in \mathbb{R} \to x - e < y + e \to y \in \mathbb{R}$\lbl{A1}}
\AXC{$\stackrel{\Pi_{3}}{\vdots}$}\noLine
\UIC{$\varphi_{35}$}
\BIC{$\varphi_{28}$}
\BIC{$x \in \mathbb{R} \to y \in \mathbb{R} \to e \in \mathbb{R} \to x - e < y + e \to y + e \in \mathbb{R}$}
\end{ptree}
\noindent$\Pi_{5}$:
\begin{ptree}
\AXC{$\forall x\,(\varphi_{5})$\lbl{[B7]}}
\AXC{$\forall x\,(\varphi_{5}) \to \varphi_{5}$\lbl{A2}}
\BIC{$\varphi_{5}$}
\AXC{$\varphi_{27}$\lbl{A2}}
\BIC{$\forall y\,(x \in \mathbb{R} \to e \in \mathbb{R} \to y \in \mathbb{R} \to x - e < y \to x < y + e)$}
\AXC{$\varphi_{30}$\lbl{A2}}
\BIC{$\varphi_{14}$}
\AXC{$\varphi_{14} \to \varphi_{45}$\lbl{A1}}
\BIC{$\varphi_{45}$}
\end{ptree}
\noindent$\Pi_{6}$:
\begin{ptree}
\AXC{$x \in \mathbb{R} \to y \in \mathbb{R} \to e \in \mathbb{R} \to x - e < y + e \to e \in \mathbb{R}$\lbl{A1}}
\AXC{$x \in \mathbb{R} \to y \in \mathbb{R} \to e \in \mathbb{R} \to x - e < y + e \to x \in \mathbb{R}$\lbl{A1}}
\AXC{$\stackrel{\Pi_{5}}{\vdots}$}\noLine
\UIC{$\varphi_{45}$}
\BIC{$\varphi_{42}$}
\BIC{$\varphi_{40}$}
\end{ptree}
\noindent$\Pi_{7}$:
\begin{ptree}
\AXC{$x \in \mathbb{R} \to y \in \mathbb{R} \to e \in \mathbb{R} \to x - e < y + e \to x - e < y + e$\lbl{A1}}
\AXC{$\stackrel{\Pi_{4}}{\vdots}$}\noLine
\UIC{$x \in \mathbb{R} \to y \in \mathbb{R} \to e \in \mathbb{R} \to x - e < y + e \to y + e \in \mathbb{R}$}
\AXC{$\stackrel{\Pi_{6}}{\vdots}$}\noLine
\UIC{$\varphi_{40}$}
\BIC{$\varphi_{31}$}
\BIC{$x \in \mathbb{R} \to y \in \mathbb{R} \to e \in \mathbb{R} \to x - e < y + e \to x < y + e + e$}
\end{ptree}
\noindent$\Pi_{8}$:
\begin{ptree}
\AXC{$\varphi_{6}$\lbl{(R2)}}
\AXC{$\varphi_{6} \to \varphi_{11}$\lbl{A2}}
\BIC{$\varphi_{11}$}
\AXC{$\varphi_{11} \to \varphi_{17}$\lbl{A2}}
\BIC{$\varphi_{17}$}
\AXC{$\varphi_{17} \to v_{785} \in \mathbb{R} \to v_{786} \in \mathbb{R} \to \varphi_{1}$\lbl{A2}}
\BIC{$v_{785} \in \mathbb{R} \to v_{786} \in \mathbb{R} \to \varphi_{1}$}
\UIC{$\forall v_{787}\,(v_{785} \in \mathbb{R} \to v_{786} \in \mathbb{R} \to \varphi_{1})$}
\UIC{$\forall v_{786}\,\forall v_{787}\,(v_{785} \in \mathbb{R} \to v_{786} \in \mathbb{R} \to \varphi_{1})$}
\UIC{$\forall v_{785}\,\forall v_{786}\,\forall v_{787}\,(v_{785} \in \mathbb{R} \to v_{786} \in \mathbb{R} \to \varphi_{1})$}
\AXC{$\varphi_{34}$\lbl{A2}}
\BIC{$\varphi_{20}$}
\AXC{$\varphi_{20} \to \varphi_{12}$\lbl{A2}}
\BIC{$\varphi_{12}$}
\end{ptree}
\noindent$\Pi_{9}$:
\begin{ptree}
\AXC{$\stackrel{\Pi_{8}}{\vdots}$}\noLine
\UIC{$\varphi_{12}$}
\AXC{$\varphi_{12} \to y \in \mathbb{R} \to e \in \mathbb{R} \to e \in \mathbb{R} \to y + e + e = y + (e + e)$\lbl{A2}}
\BIC{$y \in \mathbb{R} \to e \in \mathbb{R} \to e \in \mathbb{R} \to y + e + e = y + (e + e)$}
\AXC{$\varphi_{44}$\lbl{A1}}
\BIC{$\varphi_{41}$}
\end{ptree}
\noindent$\Pi_{10}$:
\begin{ptree}
\AXC{$x \in \mathbb{R} \to y \in \mathbb{R} \to e \in \mathbb{R} \to x - e < y + e \to e \in \mathbb{R}$\lbl{A1}}
\AXC{$x \in \mathbb{R} \to y \in \mathbb{R} \to e \in \mathbb{R} \to x - e < y + e \to y \in \mathbb{R}$\lbl{A1}}
\AXC{$\stackrel{\Pi_{9}}{\vdots}$}\noLine
\UIC{$\varphi_{41}$}
\BIC{$\varphi_{39}$}
\BIC{$\varphi_{33}$}
\end{ptree}
\noindent$\Pi_{11}$:
\begin{ptree}
\AXC{$x \in \mathbb{R} \to y \in \mathbb{R} \to e \in \mathbb{R} \to x - e < y + e \to e \in \mathbb{R}$\lbl{A1}}
\AXC{$\stackrel{\Pi_{10}}{\vdots}$}\noLine
\UIC{$\varphi_{33}$}
\BIC{$x \in \mathbb{R} \to \varphi_{4}$}
\AXC{$(x \in \mathbb{R} \to \varphi_{4}) \to \varphi_{36}$\lbl{A1}}
\BIC{$\varphi_{36}$}
\end{ptree}
\noindent$\Pi_{12}$:
\begin{ptree}
\AXC{$\forall x\,\forall y\,(x \in \mathbb{R} \to y \in \mathbb{R} \to x + y \in \mathbb{R})$\lbl{(R1)}}
\AXC{$\varphi_{21}$\lbl{A2}}
\BIC{$\forall y\,(e \in \mathbb{R} \to y \in \mathbb{R} \to e + y \in \mathbb{R})$}
\AXC{$\varphi_{19}$\lbl{A2}}
\BIC{$e \in \mathbb{R} \to e \in \mathbb{R} \to e + e \in \mathbb{R}$}
\AXC{$(e \in \mathbb{R} \to e \in \mathbb{R} \to e + e \in \mathbb{R}) \to \varphi_{37}$\lbl{A1}}
\BIC{$\varphi_{37}$}
\end{ptree}
\noindent$\Pi_{13}$:
\begin{ptree}
\AXC{$x \in \mathbb{R} \to y \in \mathbb{R} \to e \in \mathbb{R} \to x - e < y + e \to e \in \mathbb{R}$\lbl{A1}}
\AXC{$x \in \mathbb{R} \to y \in \mathbb{R} \to e \in \mathbb{R} \to x - e < y + e \to e \in \mathbb{R}$\lbl{A1}}
\AXC{$\stackrel{\Pi_{12}}{\vdots}$}\noLine
\UIC{$\varphi_{37}$}
\BIC{$\varphi_{29}$}
\BIC{$x \in \mathbb{R} \to y \in \mathbb{R} \to e \in \mathbb{R} \to x - e < y + e \to e + e \in \mathbb{R}$}
\end{ptree}
\noindent$\Pi_{14}$:
\begin{ptree}
\AXC{$\varphi_{7}$\lbl{[E32]}}
\AXC{$\varphi_{7} \to \varphi_{15}$\lbl{A2}}
\BIC{$\varphi_{15}$}
\AXC{$\varphi_{15} \to \varphi_{18}$\lbl{A2}}
\BIC{$\varphi_{18}$}
\AXC{$\varphi_{18} \to v_{788} \in \mathbb{R} \to \varphi_{10}$\lbl{A2}}
\BIC{$v_{788} \in \mathbb{R} \to \varphi_{10}$}
\UIC{$\forall v_{790}\,(v_{788} \in \mathbb{R} \to \varphi_{10})$}
\UIC{$\forall v_{789}\,\forall v_{790}\,(v_{788} \in \mathbb{R} \to \varphi_{10})$}
\UIC{$\forall v_{788}\,\forall v_{789}\,\forall v_{790}\,(v_{788} \in \mathbb{R} \to \varphi_{10})$}
\AXC{$\forall v_{788}\,\forall v_{789}\,\forall v_{790}\,(v_{788} \in \mathbb{R} \to \varphi_{10}) \to \varphi_{22}$\lbl{A2}}
\BIC{$\varphi_{22}$}
\end{ptree}
\noindent$\Pi_{15}$:
\begin{ptree}
\AXC{$x \in \mathbb{R} \to y \in \mathbb{R} \to e \in \mathbb{R} \to x - e < y + e \to y \in \mathbb{R}$\lbl{A1}}
\AXC{$\stackrel{\Pi_{14}}{\vdots}$}\noLine
\UIC{$\varphi_{22}$}
\AXC{$\varphi_{22} \to \varphi_{23}$\lbl{A2}}
\BIC{$\varphi_{23}$}
\AXC{$\varphi_{23} \to \varphi_{16}$\lbl{A2}}
\BIC{$\varphi_{16}$}
\AXC{$\varphi_{16} \to \varphi_{46}$\lbl{A1}}
\BIC{$\varphi_{46}$}
\BIC{$\varphi_{43}$}
\end{ptree}
\noindent$\Pi_{16}$:
\begin{ptree}
\AXC{$x \in \mathbb{R} \to y \in \mathbb{R} \to e \in \mathbb{R} \to x - e < y + e \to x \in \mathbb{R}$\lbl{A1}}
\AXC{$\stackrel{\Pi_{13}}{\vdots}$}\noLine
\UIC{$x \in \mathbb{R} \to y \in \mathbb{R} \to e \in \mathbb{R} \to x - e < y + e \to e + e \in \mathbb{R}$}
\AXC{$\stackrel{\Pi_{15}}{\vdots}$}\noLine
\UIC{$\varphi_{43}$}
\BIC{$\varphi_{38}$}
\BIC{$\varphi_{32}$}
\end{ptree}
\noindent 主木:
\begin{ptree}
\AXC{$\stackrel{\Pi_{1}}{\vdots}$}\noLine
\UIC{$\varphi_{2}$}
\AXC{$\stackrel{\Pi_{7}}{\vdots}$}\noLine
\UIC{$x \in \mathbb{R} \to y \in \mathbb{R} \to e \in \mathbb{R} \to x - e < y + e \to x < y + e + e$}
\AXC{$\stackrel{\Pi_{11}}{\vdots}$}\noLine
\UIC{$\varphi_{36}$}
\BIC{$\varphi_{2} \to \varphi_{9}$}
\BIC{$\varphi_{9}$}
\AXC{$\stackrel{\Pi_{16}}{\vdots}$}\noLine
\UIC{$\varphi_{32}$}
\BIC{$x \in \mathbb{R} \to y \in \mathbb{R} \to e \in \mathbb{R} \to x - e < y + e \to x - y < e + e$}
\UIC{$\varphi_{13}$}
\UIC{$\forall y\,(\varphi_{13})$}
\UIC{$\forall x\,\forall y\,(\varphi_{13})$}
\end{ptree}
\begin{abbrevlist}
\item $\varphi_{1} :\equiv v_{787} \in \mathbb{R} \to v_{785} + v_{786} + v_{787} = v_{785} + (v_{786} + v_{787})$
\item $\varphi_{2} :\equiv y + e + e = y + (e + e) \to x < y + e + e \to x < y + (e + e)$
\item $\varphi_{3} :\equiv \forall v_{783}\,\forall v_{784}\,(v_{783} \in \mathbb{R} \to v_{784} \in \mathbb{R} \to v_{783} + v_{784} \in \mathbb{R})$
\item $\varphi_{4} :\equiv y \in \mathbb{R} \to e \in \mathbb{R} \to x - e < y + e \to y + e + e = y + (e + e)$
\item $\varphi_{5} :\equiv \forall a\,\forall y\,(x \in \mathbb{R} \to a \in \mathbb{R} \to y \in \mathbb{R} \to x - a < y \to x < y + a)$
\item $\varphi_{6} :\equiv \forall x\,\forall y\,\forall z\,(x \in \mathbb{R} \to y \in \mathbb{R} \to z \in \mathbb{R} \to x + y + z = x + (y + z))$
\item $\varphi_{7} :\equiv \forall x\,\forall e\,\forall y\,(x \in \mathbb{R} \to e \in \mathbb{R} \to y \in \mathbb{R} \to y < x + e \to y - x < e)$
\item $\varphi_{8} :\equiv \forall z\,\forall w\,(z = w \to x < z \to x < w) \to \forall w\,(y + e + e = w \to x < y + e + e \to x < w)$
\item $\varphi_{9} :\equiv x \in \mathbb{R} \to y \in \mathbb{R} \to e \in \mathbb{R} \to x - e < y + e \to x < y + (e + e)$
\item $\varphi_{10} :\equiv v_{789} \in \mathbb{R} \to v_{790} \in \mathbb{R} \to v_{790} < v_{788} + v_{789} \to v_{790} - v_{788} < v_{789}$
\item $\varphi_{11} :\equiv \forall y\,\forall z\,(v_{785} \in \mathbb{R} \to y \in \mathbb{R} \to z \in \mathbb{R} \to v_{785} + y + z = v_{785} + (y + z))$
\item $\varphi_{12} :\equiv \forall v_{787}\,(y \in \mathbb{R} \to e \in \mathbb{R} \to v_{787} \in \mathbb{R} \to y + e + v_{787} = y + (e + v_{787}))$
\item $\varphi_{13} :\equiv \forall e\,(x \in \mathbb{R} \to y \in \mathbb{R} \to e \in \mathbb{R} \to x - e < y + e \to x - y < e + e)$
\item $\varphi_{14} :\equiv x \in \mathbb{R} \to e \in \mathbb{R} \to y + e \in \mathbb{R} \to x - e < y + e \to x < y + e + e$
\item $\varphi_{15} :\equiv \forall e\,\forall y\,(v_{788} \in \mathbb{R} \to e \in \mathbb{R} \to y \in \mathbb{R} \to y < v_{788} + e \to y - v_{788} < e)$
\item $\varphi_{16} :\equiv y \in \mathbb{R} \to e + e \in \mathbb{R} \to x \in \mathbb{R} \to x < y + (e + e) \to x - y < e + e$
\item $\varphi_{17} :\equiv \forall z\,(v_{785} \in \mathbb{R} \to v_{786} \in \mathbb{R} \to z \in \mathbb{R} \to v_{785} + v_{786} + z = v_{785} + (v_{786} + z))$
\item $\varphi_{18} :\equiv \forall y\,(v_{788} \in \mathbb{R} \to v_{789} \in \mathbb{R} \to y \in \mathbb{R} \to y < v_{788} + v_{789} \to y - v_{788} < v_{789})$
\item $\varphi_{19} :\equiv \forall y\,(e \in \mathbb{R} \to y \in \mathbb{R} \to e + y \in \mathbb{R}) \to e \in \mathbb{R} \to e \in \mathbb{R} \to e + e \in \mathbb{R}$
\item $\varphi_{20} :\equiv \forall v_{786}\,\forall v_{787}\,(y \in \mathbb{R} \to v_{786} \in \mathbb{R} \to v_{787} \in \mathbb{R} \to y + v_{786} + v_{787} = y + (v_{786} + v_{787}))$
\item $\varphi_{21} :\equiv \forall x\,\forall y\,(x \in \mathbb{R} \to y \in \mathbb{R} \to x + y \in \mathbb{R}) \to \forall y\,(e \in \mathbb{R} \to y \in \mathbb{R} \to e + y \in \mathbb{R})$
\item $\varphi_{22} :\equiv \forall v_{789}\,\forall v_{790}\,(y \in \mathbb{R} \to v_{789} \in \mathbb{R} \to v_{790} \in \mathbb{R} \to v_{790} < y + v_{789} \to v_{790} - y < v_{789})$
\item $\varphi_{23} :\equiv \forall v_{790}\,(y \in \mathbb{R} \to e + e \in \mathbb{R} \to v_{790} \in \mathbb{R} \to v_{790} < y + (e + e) \to v_{790} - y < e + e)$
\item $\varphi_{24} :\equiv \forall v_{784}\,(y \in \mathbb{R} \to v_{784} \in \mathbb{R} \to y + v_{784} \in \mathbb{R}) \to y \in \mathbb{R} \to e \in \mathbb{R} \to y + e \in \mathbb{R}$
\item $\varphi_{25} :\equiv \forall x\,\forall y\,(x \in \mathbb{R} \to y \in \mathbb{R} \to x + y \in \mathbb{R}) \to \forall y\,(v_{783} \in \mathbb{R} \to y \in \mathbb{R} \to v_{783} + y \in \mathbb{R})$
\item $\varphi_{26} :\equiv \forall y\,(v_{783} \in \mathbb{R} \to y \in \mathbb{R} \to v_{783} + y \in \mathbb{R}) \to v_{783} \in \mathbb{R} \to v_{784} \in \mathbb{R} \to v_{783} + v_{784} \in \mathbb{R}$
\item $\varphi_{27} :\equiv \varphi_{5} \to \forall y\,(x \in \mathbb{R} \to e \in \mathbb{R} \to y \in \mathbb{R} \to x - e < y \to x < y + e)$
\item $\varphi_{28} :\equiv (x \in \mathbb{R} \to y \in \mathbb{R} \to e \in \mathbb{R} \to x - e < y + e \to e \in \mathbb{R}) \to x \in \mathbb{R} \to y \in \mathbb{R} \to e \in \mathbb{R} \to x - e < y + e \to y + e \in \mathbb{R}$
\item $\varphi_{29} :\equiv (x \in \mathbb{R} \to y \in \mathbb{R} \to e \in \mathbb{R} \to x - e < y + e \to e \in \mathbb{R}) \to x \in \mathbb{R} \to y \in \mathbb{R} \to e \in \mathbb{R} \to x - e < y + e \to e + e \in \mathbb{R}$
\item $\varphi_{30} :\equiv \forall y\,(x \in \mathbb{R} \to e \in \mathbb{R} \to y \in \mathbb{R} \to x - e < y \to x < y + e) \to \varphi_{14}$
\item $\varphi_{31} :\equiv (x \in \mathbb{R} \to y \in \mathbb{R} \to e \in \mathbb{R} \to x - e < y + e \to x - e < y + e) \to x \in \mathbb{R} \to y \in \mathbb{R} \to e \in \mathbb{R} \to x - e < y + e \to x < y + e + e$
\item $\varphi_{32} :\equiv \varphi_{9} \to x \in \mathbb{R} \to y \in \mathbb{R} \to e \in \mathbb{R} \to x - e < y + e \to x - y < e + e$
\item $\varphi_{33} :\equiv (x \in \mathbb{R} \to y \in \mathbb{R} \to e \in \mathbb{R} \to x - e < y + e \to e \in \mathbb{R}) \to x \in \mathbb{R} \to \varphi_{4}$
\item $\varphi_{34} :\equiv \forall v_{785}\,\forall v_{786}\,\forall v_{787}\,(v_{785} \in \mathbb{R} \to v_{786} \in \mathbb{R} \to \varphi_{1}) \to \varphi_{20}$
\item $\varphi_{35} :\equiv (x \in \mathbb{R} \to y \in \mathbb{R} \to e \in \mathbb{R} \to x - e < y + e \to y \in \mathbb{R}) \to \varphi_{28}$
\item $\varphi_{36} :\equiv (x \in \mathbb{R} \to y \in \mathbb{R} \to e \in \mathbb{R} \to x - e < y + e \to x < y + e + e) \to \varphi_{2} \to \varphi_{9}$
\item $\varphi_{37} :\equiv (x \in \mathbb{R} \to y \in \mathbb{R} \to e \in \mathbb{R} \to x - e < y + e \to e \in \mathbb{R}) \to \varphi_{29}$
\item $\varphi_{38} :\equiv (x \in \mathbb{R} \to y \in \mathbb{R} \to e \in \mathbb{R} \to x - e < y + e \to x \in \mathbb{R}) \to \varphi_{32}$
\item $\varphi_{39} :\equiv (x \in \mathbb{R} \to y \in \mathbb{R} \to e \in \mathbb{R} \to x - e < y + e \to e \in \mathbb{R}) \to \varphi_{33}$
\item $\varphi_{40} :\equiv (x \in \mathbb{R} \to y \in \mathbb{R} \to e \in \mathbb{R} \to x - e < y + e \to y + e \in \mathbb{R}) \to \varphi_{31}$
\item $\varphi_{41} :\equiv (x \in \mathbb{R} \to y \in \mathbb{R} \to e \in \mathbb{R} \to x - e < y + e \to y \in \mathbb{R}) \to \varphi_{39}$
\item $\varphi_{42} :\equiv (x \in \mathbb{R} \to y \in \mathbb{R} \to e \in \mathbb{R} \to x - e < y + e \to e \in \mathbb{R}) \to \varphi_{40}$
\item $\varphi_{43} :\equiv (x \in \mathbb{R} \to y \in \mathbb{R} \to e \in \mathbb{R} \to x - e < y + e \to e + e \in \mathbb{R}) \to \varphi_{38}$
\item $\varphi_{44} :\equiv (y \in \mathbb{R} \to e \in \mathbb{R} \to e \in \mathbb{R} \to y + e + e = y + (e + e)) \to \varphi_{41}$
\item $\varphi_{45} :\equiv (x \in \mathbb{R} \to y \in \mathbb{R} \to e \in \mathbb{R} \to x - e < y + e \to x \in \mathbb{R}) \to \varphi_{42}$
\item $\varphi_{46} :\equiv (x \in \mathbb{R} \to y \in \mathbb{R} \to e \in \mathbb{R} \to x - e < y + e \to y \in \mathbb{R}) \to \varphi_{43}$
\end{abbrevlist}
\end{fullproof}

% 自動生成 (prover/build.js): SubLtAnti  ノード数 182
\begin{fullproof}{[E45]}{$\forall z\,\forall x\,\forall y\,(z \in \mathbb{R} \to x \in \mathbb{R} \to y \in \mathbb{R} \to x < y \to z - y < z - x)$}{182}\phantomsection\label{tree:SubLtAnti}
\noindent$\Pi_{1}$:
\begin{ptree}
\AXC{$z_{1} = w \to z - y < z_{1} \to z - y < w$\lbl{A6}}
\UIC{$\forall w\,(z_{1} = w \to z - y < z_{1} \to z - y < w)$}
\UIC{$\forall z_{1}\,\forall w\,(z_{1} = w \to z - y < z_{1} \to z - y < w)$}
\AXC{$\varphi_{27}$\lbl{A2}}
\BIC{$\forall w\,(-x + z = w \to z - y < -x + z \to z - y < w)$}
\AXC{$\varphi_{37}$\lbl{A2}}
\BIC{$-x + z = z - x \to z - y < -x + z \to z - y < z - x$}
\end{ptree}
\noindent$\Pi_{2}$:
\begin{ptree}
\AXC{$\forall x\,\forall y\,(x \in \mathbb{R} \to y \in \mathbb{R} \to x < y \to -y < -x)$\lbl{[F22]}}
\AXC{$\varphi_{34}$\lbl{A2}}
\BIC{$\forall y\,(x \in \mathbb{R} \to y \in \mathbb{R} \to x < y \to -y < -x)$}
\AXC{$\varphi_{30}$\lbl{A2}}
\BIC{$x \in \mathbb{R} \to y \in \mathbb{R} \to x < y \to -y < -x$}
\AXC{$(x \in \mathbb{R} \to y \in \mathbb{R} \to x < y \to -y < -x) \to \varphi_{53}$\lbl{A1}}
\BIC{$\varphi_{53}$}
\end{ptree}
\noindent$\Pi_{3}$:
\begin{ptree}
\AXC{$z \in \mathbb{R} \to x \in \mathbb{R} \to y \in \mathbb{R} \to x < y \to y \in \mathbb{R}$\lbl{A1}}
\AXC{$z \in \mathbb{R} \to x \in \mathbb{R} \to y \in \mathbb{R} \to x < y \to x \in \mathbb{R}$\lbl{A1}}
\AXC{$\stackrel{\Pi_{2}}{\vdots}$}\noLine
\UIC{$\varphi_{53}$}
\BIC{$(z \in \mathbb{R} \to x \in \mathbb{R} \to y \in \mathbb{R} \to x < y \to y \in \mathbb{R}) \to \varphi_{38}$}
\BIC{$\varphi_{38}$}
\end{ptree}
\noindent$\Pi_{4}$:
\begin{ptree}
\AXC{$\forall x \in \mathbb{R}\,(-x \in \mathbb{R})$\lbl{(R1)}}
\AXC{$\forall x \in \mathbb{R}\,(-x \in \mathbb{R}) \to x \in \mathbb{R} \to -x \in \mathbb{R}$\lbl{A2}}
\BIC{$x \in \mathbb{R} \to -x \in \mathbb{R}$}
\AXC{$(x \in \mathbb{R} \to -x \in \mathbb{R}) \to \varphi_{39}$\lbl{A1}}
\BIC{$\varphi_{39}$}
\end{ptree}
\noindent$\Pi_{5}$:
\begin{ptree}
\AXC{$z \in \mathbb{R} \to x \in \mathbb{R} \to y \in \mathbb{R} \to x < y \to x \in \mathbb{R}$\lbl{A1}}
\AXC{$\stackrel{\Pi_{4}}{\vdots}$}\noLine
\UIC{$\varphi_{39}$}
\BIC{$z \in \mathbb{R} \to x \in \mathbb{R} \to y \in \mathbb{R} \to x < y \to -x \in \mathbb{R}$}
\end{ptree}
\noindent$\Pi_{6}$:
\begin{ptree}
\AXC{$\forall x \in \mathbb{R}\,(-x \in \mathbb{R})$\lbl{(R1)}}
\AXC{$\forall x \in \mathbb{R}\,(-x \in \mathbb{R}) \to y \in \mathbb{R} \to -y \in \mathbb{R}$\lbl{A2}}
\BIC{$y \in \mathbb{R} \to -y \in \mathbb{R}$}
\AXC{$(y \in \mathbb{R} \to -y \in \mathbb{R}) \to \varphi_{40}$\lbl{A1}}
\BIC{$\varphi_{40}$}
\end{ptree}
\noindent$\Pi_{7}$:
\begin{ptree}
\AXC{$z \in \mathbb{R} \to x \in \mathbb{R} \to y \in \mathbb{R} \to x < y \to y \in \mathbb{R}$\lbl{A1}}
\AXC{$\stackrel{\Pi_{6}}{\vdots}$}\noLine
\UIC{$\varphi_{40}$}
\BIC{$z \in \mathbb{R} \to x \in \mathbb{R} \to y \in \mathbb{R} \to x < y \to -y \in \mathbb{R}$}
\end{ptree}
\noindent$\Pi_{8}$:
\begin{ptree}
\AXC{$\varphi_{19}$\lbl{(R14)}}
\AXC{$\varphi_{19} \to \varphi_{22}$\lbl{A2}}
\BIC{$\varphi_{22}$}
\AXC{$\varphi_{22} \to \varphi_{25}$\lbl{A2}}
\BIC{$\varphi_{25}$}
\AXC{$\varphi_{25} \to v_{791} \in \mathbb{R} \to \varphi_{21}$\lbl{A2}}
\BIC{$v_{791} \in \mathbb{R} \to \varphi_{21}$}
\UIC{$\forall v_{793}\,(v_{791} \in \mathbb{R} \to \varphi_{21})$}
\UIC{$\forall v_{792}\,\forall v_{793}\,(v_{791} \in \mathbb{R} \to \varphi_{21})$}
\UIC{$\forall v_{791}\,\forall v_{792}\,\forall v_{793}\,(v_{791} \in \mathbb{R} \to \varphi_{21})$}
\AXC{$\forall v_{791}\,\forall v_{792}\,\forall v_{793}\,(v_{791} \in \mathbb{R} \to \varphi_{21}) \to \varphi_{31}$\lbl{A2}}
\BIC{$\varphi_{31}$}
\end{ptree}
\noindent$\Pi_{9}$:
\begin{ptree}
\AXC{$\stackrel{\Pi_{8}}{\vdots}$}\noLine
\UIC{$\varphi_{31}$}
\AXC{$\varphi_{31} \to \varphi_{32}$\lbl{A2}}
\BIC{$\varphi_{32}$}
\AXC{$\varphi_{32} \to \varphi_{26}$\lbl{A2}}
\BIC{$\varphi_{26}$}
\AXC{$\varphi_{26} \to (z \in \mathbb{R} \to x \in \mathbb{R} \to y \in \mathbb{R} \to x < y \to -y \in \mathbb{R}) \to \varphi_{56}$\lbl{A1}}
\BIC{$(z \in \mathbb{R} \to x \in \mathbb{R} \to y \in \mathbb{R} \to x < y \to -y \in \mathbb{R}) \to \varphi_{56}$}
\end{ptree}
\noindent$\Pi_{10}$:
\begin{ptree}
\AXC{$\stackrel{\Pi_{7}}{\vdots}$}\noLine
\UIC{$z \in \mathbb{R} \to x \in \mathbb{R} \to y \in \mathbb{R} \to x < y \to -y \in \mathbb{R}$}
\AXC{$\stackrel{\Pi_{9}}{\vdots}$}\noLine
\UIC{$(z \in \mathbb{R} \to x \in \mathbb{R} \to y \in \mathbb{R} \to x < y \to -y \in \mathbb{R}) \to \varphi_{56}$}
\BIC{$\varphi_{56}$}
\end{ptree}
\noindent$\Pi_{11}$:
\begin{ptree}
\AXC{$z \in \mathbb{R} \to x \in \mathbb{R} \to y \in \mathbb{R} \to x < y \to z \in \mathbb{R}$\lbl{A1}}
\AXC{$\stackrel{\Pi_{5}}{\vdots}$}\noLine
\UIC{$z \in \mathbb{R} \to x \in \mathbb{R} \to y \in \mathbb{R} \to x < y \to -x \in \mathbb{R}$}
\AXC{$\stackrel{\Pi_{10}}{\vdots}$}\noLine
\UIC{$\varphi_{56}$}
\BIC{$(z \in \mathbb{R} \to x \in \mathbb{R} \to y \in \mathbb{R} \to x < y \to z \in \mathbb{R}) \to \varphi_{47}$}
\BIC{$\varphi_{47}$}
\end{ptree}
\noindent$\Pi_{12}$:
\begin{ptree}
\AXC{$y = z \to x = y \to x = z$\lbl{A6}}
\AXC{$(y = z \to x = y \to x = z) \to x = y \to y = z \to x = z$\lbl{A1}}
\BIC{$x = y \to y = z \to x = z$}
\UIC{$\forall z\,(x = y \to y = z \to x = z)$}
\UIC{$\forall y\,\forall z\,(x = y \to y = z \to x = z)$}
\UIC{$\forall x\,\forall y\,\forall z\,(x = y \to y = z \to x = z)$}
\AXC{$\varphi_{4}$\lbl{A2}}
\BIC{$\forall y\,\forall z\,(v_{801} = y \to y = z \to v_{801} = z)$}
\AXC{$\varphi_{8}$\lbl{A2}}
\BIC{$\forall z\,(v_{801} = v_{802} \to v_{802} = z \to v_{801} = z)$}
\AXC{$\varphi_{12}$\lbl{A2}}
\BIC{$v_{801} = v_{802} \to v_{802} = v_{803} \to v_{801} = v_{803}$}
\UIC{$\forall v_{803}\,(v_{801} = v_{802} \to v_{802} = v_{803} \to v_{801} = v_{803})$}
\UIC{$\varphi_{2}$}
\UIC{$\forall v_{801}\,(\varphi_{2})$}
\end{ptree}
\noindent$\Pi_{13}$:
\begin{ptree}
\AXC{$\stackrel{\Pi_{12}}{\vdots}$}\noLine
\UIC{$\forall v_{801}\,(\varphi_{2})$}
\AXC{$\forall v_{801}\,(\varphi_{2}) \to \varphi_{6}$\lbl{A2}}
\BIC{$\varphi_{6}$}
\AXC{$\varphi_{6} \to \varphi_{14}$\lbl{A2}}
\BIC{$\varphi_{14}$}
\AXC{$\varphi_{14} \to -y + z = z + -y \to z + -y = z - y \to -y + z = z - y$\lbl{A2}}
\BIC{$-y + z = z + -y \to z + -y = z - y \to -y + z = z - y$}
\end{ptree}
\noindent$\Pi_{14}$:
\begin{ptree}
\AXC{$x = x$\lbl{A5}}
\AXC{$x = y \to x = x \to y = x$\lbl{A6}}
\AXC{$(x = y \to x = x \to y = x) \to x = x \to x = y \to y = x$\lbl{A1}}
\BIC{$x = x \to x = y \to y = x$}
\BIC{$x = y \to y = x$}
\UIC{$\forall y\,(x = y \to y = x)$}
\UIC{$\forall x\,\forall y\,(x = y \to y = x)$}
\end{ptree}
\noindent$\Pi_{15}$:
\begin{ptree}
\AXC{$\stackrel{\Pi_{14}}{\vdots}$}\noLine
\UIC{$\forall x\,\forall y\,(x = y \to y = x)$}
\AXC{$\forall x\,\forall y\,(x = y \to y = x) \to \forall y\,(v_{799} = y \to y = v_{799})$\lbl{A2}}
\BIC{$\forall y\,(v_{799} = y \to y = v_{799})$}
\AXC{$\varphi_{1}$\lbl{A2}}
\BIC{$v_{799} = v_{800} \to v_{800} = v_{799}$}
\UIC{$\forall v_{800}\,(v_{799} = v_{800} \to v_{800} = v_{799})$}
\UIC{$\forall v_{799}\,\forall v_{800}\,(v_{799} = v_{800} \to v_{800} = v_{799})$}
\AXC{$\varphi_{15}$\lbl{A2}}
\BIC{$\forall v_{800}\,(z - y = v_{800} \to v_{800} = z - y)$}
\AXC{$\varphi_{20}$\lbl{A2}}
\BIC{$z - y = z + -y \to z + -y = z - y$}
\end{ptree}
\noindent$\Pi_{16}$:
\begin{ptree}
\AXC{$\forall x\,\forall y\,(x - y = x + -y)$\lbl{Def$-$}}
\AXC{$\forall x\,\forall y\,(x - y = x + -y) \to \forall y\,(z - y = z + -y)$\lbl{A2}}
\BIC{$\forall y\,(z - y = z + -y)$}
\end{ptree}
\noindent$\Pi_{17}$:
\begin{ptree}
\AXC{$\stackrel{\Pi_{15}}{\vdots}$}\noLine
\UIC{$z - y = z + -y \to z + -y = z - y$}
\AXC{$\stackrel{\Pi_{16}}{\vdots}$}\noLine
\UIC{$\forall y\,(z - y = z + -y)$}
\AXC{$\forall y\,(z - y = z + -y) \to z - y = z + -y$\lbl{A2}}
\BIC{$z - y = z + -y$}
\AXC{$\varphi_{23}$\lbl{A1}}
\BIC{$(z - y = z + -y \to z + -y = z - y) \to z + -y = z - y$}
\BIC{$z + -y = z - y$}
\end{ptree}
\noindent$\Pi_{18}$:
\begin{ptree}
\AXC{$\forall x\,\forall y\,(x \in \mathbb{R} \to y \in \mathbb{R} \to x + y = y + x)$\lbl{(R3)}}
\AXC{$\varphi_{28}$\lbl{A2}}
\BIC{$\forall y\,(v_{797} \in \mathbb{R} \to y \in \mathbb{R} \to v_{797} + y = y + v_{797})$}
\AXC{$\varphi_{33}$\lbl{A2}}
\BIC{$v_{797} \in \mathbb{R} \to v_{798} \in \mathbb{R} \to v_{797} + v_{798} = v_{798} + v_{797}$}
\UIC{$\varphi_{10}$}
\UIC{$\forall v_{797}\,(\varphi_{10})$}
\AXC{$\forall v_{797}\,(\varphi_{10}) \to \varphi_{17}$\lbl{A2}}
\BIC{$\varphi_{17}$}
\end{ptree}
\noindent$\Pi_{19}$:
\begin{ptree}
\AXC{$\stackrel{\Pi_{7}}{\vdots}$}\noLine
\UIC{$z \in \mathbb{R} \to x \in \mathbb{R} \to y \in \mathbb{R} \to x < y \to -y \in \mathbb{R}$}
\AXC{$\stackrel{\Pi_{18}}{\vdots}$}\noLine
\UIC{$\varphi_{17}$}
\AXC{$\varphi_{17} \to -y \in \mathbb{R} \to z \in \mathbb{R} \to -y + z = z + -y$\lbl{A2}}
\BIC{$-y \in \mathbb{R} \to z \in \mathbb{R} \to -y + z = z + -y$}
\AXC{$\varphi_{54}$\lbl{A1}}
\BIC{$(z \in \mathbb{R} \to x \in \mathbb{R} \to y \in \mathbb{R} \to x < y \to -y \in \mathbb{R}) \to \varphi_{42}$}
\BIC{$\varphi_{42}$}
\end{ptree}
\noindent$\Pi_{20}$:
\begin{ptree}
\AXC{$\stackrel{\Pi_{17}}{\vdots}$}\noLine
\UIC{$z + -y = z - y$}
\AXC{$z \in \mathbb{R} \to x \in \mathbb{R} \to y \in \mathbb{R} \to x < y \to z \in \mathbb{R}$\lbl{A1}}
\AXC{$\stackrel{\Pi_{19}}{\vdots}$}\noLine
\UIC{$\varphi_{42}$}
\BIC{$z \in \mathbb{R} \to x \in \mathbb{R} \to y \in \mathbb{R} \to x < y \to -y + z = z + -y$}
\AXC{$\varphi_{51}$\lbl{A1}}
\BIC{$z + -y = z - y \to \varphi_{44}$}
\BIC{$\varphi_{44}$}
\end{ptree}
\noindent$\Pi_{21}$:
\begin{ptree}
\AXC{$\stackrel{\Pi_{13}}{\vdots}$}\noLine
\UIC{$-y + z = z + -y \to z + -y = z - y \to -y + z = z - y$}
\AXC{$\stackrel{\Pi_{20}}{\vdots}$}\noLine
\UIC{$\varphi_{44}$}
\BIC{$z \in \mathbb{R} \to x \in \mathbb{R} \to y \in \mathbb{R} \to x < y \to -y + z = z - y$}
\AXC{$\varphi_{58}$\lbl{A1}}
\BIC{$\varphi_{50}$}
\end{ptree}
\noindent$\Pi_{22}$:
\begin{ptree}
\AXC{$z \in \mathbb{R} \to x \in \mathbb{R} \to y \in \mathbb{R} \to x < y \to x < y$\lbl{A1}}
\AXC{$\stackrel{\Pi_{3}}{\vdots}$}\noLine
\UIC{$\varphi_{38}$}
\BIC{$z \in \mathbb{R} \to x \in \mathbb{R} \to y \in \mathbb{R} \to x < y \to -y < -x$}
\AXC{$\stackrel{\Pi_{11}}{\vdots}$}\noLine
\UIC{$\varphi_{47}$}
\BIC{$z \in \mathbb{R} \to x \in \mathbb{R} \to y \in \mathbb{R} \to x < y \to -y + z < -x + z$}
\AXC{$\stackrel{\Pi_{21}}{\vdots}$}\noLine
\UIC{$\varphi_{50}$}
\BIC{$\varphi_{48}$}
\end{ptree}
\noindent$\Pi_{23}$:
\begin{ptree}
\AXC{$z_{1} = w \to z_{1} < -x + z \to w < -x + z$\lbl{A6}}
\UIC{$\forall w\,(z_{1} = w \to z_{1} < -x + z \to w < -x + z)$}
\UIC{$\forall z_{1}\,\forall w\,(z_{1} = w \to z_{1} < -x + z \to w < -x + z)$}
\AXC{$\varphi_{36}$\lbl{A2}}
\BIC{$\forall w\,(-y + z = w \to -y + z < -x + z \to w < -x + z)$}
\AXC{$\varphi_{46}$\lbl{A2}}
\BIC{$-y + z = z - y \to -y + z < -x + z \to z - y < -x + z$}
\AXC{$\stackrel{\Pi_{22}}{\vdots}$}\noLine
\UIC{$\varphi_{48}$}
\BIC{$z \in \mathbb{R} \to x \in \mathbb{R} \to y \in \mathbb{R} \to x < y \to z - y < -x + z$}
\end{ptree}
\noindent$\Pi_{24}$:
\begin{ptree}
\AXC{$y = z \to x = y \to x = z$\lbl{A6}}
\AXC{$(y = z \to x = y \to x = z) \to x = y \to y = z \to x = z$\lbl{A1}}
\BIC{$x = y \to y = z \to x = z$}
\UIC{$\forall z\,(x = y \to y = z \to x = z)$}
\UIC{$\forall y\,\forall z\,(x = y \to y = z \to x = z)$}
\UIC{$\forall x\,\forall y\,\forall z\,(x = y \to y = z \to x = z)$}
\AXC{$\varphi_{5}$\lbl{A2}}
\BIC{$\forall y\,\forall z\,(v_{794} = y \to y = z \to v_{794} = z)$}
\AXC{$\varphi_{9}$\lbl{A2}}
\BIC{$\forall z\,(v_{794} = v_{795} \to v_{795} = z \to v_{794} = z)$}
\AXC{$\varphi_{13}$\lbl{A2}}
\BIC{$v_{794} = v_{795} \to v_{795} = v_{796} \to v_{794} = v_{796}$}
\UIC{$\forall v_{796}\,(v_{794} = v_{795} \to v_{795} = v_{796} \to v_{794} = v_{796})$}
\UIC{$\varphi_{3}$}
\UIC{$\forall v_{794}\,(\varphi_{3})$}
\end{ptree}
\noindent$\Pi_{25}$:
\begin{ptree}
\AXC{$\stackrel{\Pi_{24}}{\vdots}$}\noLine
\UIC{$\forall v_{794}\,(\varphi_{3})$}
\AXC{$\forall v_{794}\,(\varphi_{3}) \to \varphi_{7}$\lbl{A2}}
\BIC{$\varphi_{7}$}
\AXC{$\varphi_{7} \to \varphi_{16}$\lbl{A2}}
\BIC{$\varphi_{16}$}
\AXC{$\varphi_{16} \to -x + z = z + -x \to z + -x = z - x \to -x + z = z - x$\lbl{A2}}
\BIC{$-x + z = z + -x \to z + -x = z - x \to -x + z = z - x$}
\end{ptree}
\noindent$\Pi_{26}$:
\begin{ptree}
\AXC{$\stackrel{\Pi_{14}}{\vdots}$}\noLine
\UIC{$\forall x\,\forall y\,(x = y \to y = x)$}
\AXC{$\forall x\,\forall y\,(x = y \to y = x) \to \forall y\,(z - x = y \to y = z - x)$\lbl{A2}}
\BIC{$\forall y\,(z - x = y \to y = z - x)$}
\AXC{$\varphi_{11}$\lbl{A2}}
\BIC{$z - x = z + -x \to z + -x = z - x$}
\end{ptree}
\noindent$\Pi_{27}$:
\begin{ptree}
\AXC{$\forall x\,\forall y\,(x - y = x + -y)$\lbl{Def$-$}}
\AXC{$\forall x\,\forall y\,(x - y = x + -y) \to \forall y\,(z - y = z + -y)$\lbl{A2}}
\BIC{$\forall y\,(z - y = z + -y)$}
\end{ptree}
\noindent$\Pi_{28}$:
\begin{ptree}
\AXC{$\stackrel{\Pi_{26}}{\vdots}$}\noLine
\UIC{$z - x = z + -x \to z + -x = z - x$}
\AXC{$\stackrel{\Pi_{27}}{\vdots}$}\noLine
\UIC{$\forall y\,(z - y = z + -y)$}
\AXC{$\forall y\,(z - y = z + -y) \to z - x = z + -x$\lbl{A2}}
\BIC{$z - x = z + -x$}
\AXC{$\varphi_{24}$\lbl{A1}}
\BIC{$(z - x = z + -x \to z + -x = z - x) \to z + -x = z - x$}
\BIC{$z + -x = z - x$}
\end{ptree}
\noindent$\Pi_{29}$:
\begin{ptree}
\AXC{$\forall x\,\forall y\,(x \in \mathbb{R} \to y \in \mathbb{R} \to x + y = y + x)$\lbl{(R3)}}
\AXC{$\varphi_{29}$\lbl{A2}}
\BIC{$\forall y\,(-x \in \mathbb{R} \to y \in \mathbb{R} \to -x + y = y + -x)$}
\AXC{$\varphi_{35}$\lbl{A2}}
\BIC{$-x \in \mathbb{R} \to z \in \mathbb{R} \to -x + z = z + -x$}
\AXC{$\varphi_{55}$\lbl{A1}}
\BIC{$(z \in \mathbb{R} \to x \in \mathbb{R} \to y \in \mathbb{R} \to x < y \to -x \in \mathbb{R}) \to \varphi_{43}$}
\end{ptree}
\noindent$\Pi_{30}$:
\begin{ptree}
\AXC{$\stackrel{\Pi_{5}}{\vdots}$}\noLine
\UIC{$z \in \mathbb{R} \to x \in \mathbb{R} \to y \in \mathbb{R} \to x < y \to -x \in \mathbb{R}$}
\AXC{$\stackrel{\Pi_{29}}{\vdots}$}\noLine
\UIC{$(z \in \mathbb{R} \to x \in \mathbb{R} \to y \in \mathbb{R} \to x < y \to -x \in \mathbb{R}) \to \varphi_{43}$}
\BIC{$\varphi_{43}$}
\end{ptree}
\noindent$\Pi_{31}$:
\begin{ptree}
\AXC{$\stackrel{\Pi_{28}}{\vdots}$}\noLine
\UIC{$z + -x = z - x$}
\AXC{$z \in \mathbb{R} \to x \in \mathbb{R} \to y \in \mathbb{R} \to x < y \to z \in \mathbb{R}$\lbl{A1}}
\AXC{$\stackrel{\Pi_{30}}{\vdots}$}\noLine
\UIC{$\varphi_{43}$}
\BIC{$z \in \mathbb{R} \to x \in \mathbb{R} \to y \in \mathbb{R} \to x < y \to -x + z = z + -x$}
\AXC{$\varphi_{52}$\lbl{A1}}
\BIC{$z + -x = z - x \to \varphi_{45}$}
\BIC{$\varphi_{45}$}
\end{ptree}
\noindent$\Pi_{32}$:
\begin{ptree}
\AXC{$\stackrel{\Pi_{25}}{\vdots}$}\noLine
\UIC{$-x + z = z + -x \to z + -x = z - x \to -x + z = z - x$}
\AXC{$\stackrel{\Pi_{31}}{\vdots}$}\noLine
\UIC{$\varphi_{45}$}
\BIC{$z \in \mathbb{R} \to x \in \mathbb{R} \to y \in \mathbb{R} \to x < y \to -x + z = z - x$}
\AXC{$\varphi_{57}$\lbl{A1}}
\BIC{$\varphi_{49}$}
\end{ptree}
\noindent 主木:
\begin{ptree}
\AXC{$\stackrel{\Pi_{1}}{\vdots}$}\noLine
\UIC{$-x + z = z - x \to z - y < -x + z \to z - y < z - x$}
\AXC{$\stackrel{\Pi_{23}}{\vdots}$}\noLine
\UIC{$z \in \mathbb{R} \to x \in \mathbb{R} \to y \in \mathbb{R} \to x < y \to z - y < -x + z$}
\AXC{$\stackrel{\Pi_{32}}{\vdots}$}\noLine
\UIC{$\varphi_{49}$}
\BIC{$\varphi_{41}$}
\BIC{$z \in \mathbb{R} \to x \in \mathbb{R} \to y \in \mathbb{R} \to x < y \to z - y < z - x$}
\UIC{$\forall y\,(z \in \mathbb{R} \to x \in \mathbb{R} \to y \in \mathbb{R} \to x < y \to z - y < z - x)$}
\UIC{$\varphi_{18}$}
\UIC{$\forall z\,(\varphi_{18})$}
\end{ptree}
\begin{abbrevlist}
\item $\varphi_{1} :\equiv \forall y\,(v_{799} = y \to y = v_{799}) \to v_{799} = v_{800} \to v_{800} = v_{799}$
\item $\varphi_{2} :\equiv \forall v_{802}\,\forall v_{803}\,(v_{801} = v_{802} \to v_{802} = v_{803} \to v_{801} = v_{803})$
\item $\varphi_{3} :\equiv \forall v_{795}\,\forall v_{796}\,(v_{794} = v_{795} \to v_{795} = v_{796} \to v_{794} = v_{796})$
\item $\varphi_{4} :\equiv \forall x\,\forall y\,\forall z\,(x = y \to y = z \to x = z) \to \forall y\,\forall z\,(v_{801} = y \to y = z \to v_{801} = z)$
\item $\varphi_{5} :\equiv \forall x\,\forall y\,\forall z\,(x = y \to y = z \to x = z) \to \forall y\,\forall z\,(v_{794} = y \to y = z \to v_{794} = z)$
\item $\varphi_{6} :\equiv \forall v_{802}\,\forall v_{803}\,(-y + z = v_{802} \to v_{802} = v_{803} \to -y + z = v_{803})$
\item $\varphi_{7} :\equiv \forall v_{795}\,\forall v_{796}\,(-x + z = v_{795} \to v_{795} = v_{796} \to -x + z = v_{796})$
\item $\varphi_{8} :\equiv \forall y\,\forall z\,(v_{801} = y \to y = z \to v_{801} = z) \to \forall z\,(v_{801} = v_{802} \to v_{802} = z \to v_{801} = z)$
\item $\varphi_{9} :\equiv \forall y\,\forall z\,(v_{794} = y \to y = z \to v_{794} = z) \to \forall z\,(v_{794} = v_{795} \to v_{795} = z \to v_{794} = z)$
\item $\varphi_{10} :\equiv \forall v_{798}\,(v_{797} \in \mathbb{R} \to v_{798} \in \mathbb{R} \to v_{797} + v_{798} = v_{798} + v_{797})$
\item $\varphi_{11} :\equiv \forall y\,(z - x = y \to y = z - x) \to z - x = z + -x \to z + -x = z - x$
\item $\varphi_{12} :\equiv \forall z\,(v_{801} = v_{802} \to v_{802} = z \to v_{801} = z) \to v_{801} = v_{802} \to v_{802} = v_{803} \to v_{801} = v_{803}$
\item $\varphi_{13} :\equiv \forall z\,(v_{794} = v_{795} \to v_{795} = z \to v_{794} = z) \to v_{794} = v_{795} \to v_{795} = v_{796} \to v_{794} = v_{796}$
\item $\varphi_{14} :\equiv \forall v_{803}\,(-y + z = z + -y \to z + -y = v_{803} \to -y + z = v_{803})$
\item $\varphi_{15} :\equiv \forall v_{799}\,\forall v_{800}\,(v_{799} = v_{800} \to v_{800} = v_{799}) \to \forall v_{800}\,(z - y = v_{800} \to v_{800} = z - y)$
\item $\varphi_{16} :\equiv \forall v_{796}\,(-x + z = z + -x \to z + -x = v_{796} \to -x + z = v_{796})$
\item $\varphi_{17} :\equiv \forall v_{798}\,(-y \in \mathbb{R} \to v_{798} \in \mathbb{R} \to -y + v_{798} = v_{798} + -y)$
\item $\varphi_{18} :\equiv \forall x\,\forall y\,(z \in \mathbb{R} \to x \in \mathbb{R} \to y \in \mathbb{R} \to x < y \to z - y < z - x)$
\item $\varphi_{19} :\equiv \forall x\,\forall y\,\forall z\,(x \in \mathbb{R} \to y \in \mathbb{R} \to z \in \mathbb{R} \to x < y \to x + z < y + z)$
\item $\varphi_{20} :\equiv \forall v_{800}\,(z - y = v_{800} \to v_{800} = z - y) \to z - y = z + -y \to z + -y = z - y$
\item $\varphi_{21} :\equiv v_{792} \in \mathbb{R} \to v_{793} \in \mathbb{R} \to v_{791} < v_{792} \to v_{791} + v_{793} < v_{792} + v_{793}$
\item $\varphi_{22} :\equiv \forall y\,\forall z\,(v_{791} \in \mathbb{R} \to y \in \mathbb{R} \to z \in \mathbb{R} \to v_{791} < y \to v_{791} + z < y + z)$
\item $\varphi_{23} :\equiv z - y = z + -y \to (z - y = z + -y \to z + -y = z - y) \to z + -y = z - y$
\item $\varphi_{24} :\equiv z - x = z + -x \to (z - x = z + -x \to z + -x = z - x) \to z + -x = z - x$
\item $\varphi_{25} :\equiv \forall z\,(v_{791} \in \mathbb{R} \to v_{792} \in \mathbb{R} \to z \in \mathbb{R} \to v_{791} < v_{792} \to v_{791} + z < v_{792} + z)$
\item $\varphi_{26} :\equiv -y \in \mathbb{R} \to -x \in \mathbb{R} \to z \in \mathbb{R} \to -y < -x \to -y + z < -x + z$
\item $\varphi_{27} :\equiv \forall z_{1}\,\forall w\,(z_{1} = w \to z - y < z_{1} \to z - y < w) \to \forall w\,(-x + z = w \to z - y < -x + z \to z - y < w)$
\item $\varphi_{28} :\equiv \forall x\,\forall y\,(x \in \mathbb{R} \to y \in \mathbb{R} \to x + y = y + x) \to \forall y\,(v_{797} \in \mathbb{R} \to y \in \mathbb{R} \to v_{797} + y = y + v_{797})$
\item $\varphi_{29} :\equiv \forall x\,\forall y\,(x \in \mathbb{R} \to y \in \mathbb{R} \to x + y = y + x) \to \forall y\,(-x \in \mathbb{R} \to y \in \mathbb{R} \to -x + y = y + -x)$
\item $\varphi_{30} :\equiv \forall y\,(x \in \mathbb{R} \to y \in \mathbb{R} \to x < y \to -y < -x) \to x \in \mathbb{R} \to y \in \mathbb{R} \to x < y \to -y < -x$
\item $\varphi_{31} :\equiv \forall v_{792}\,\forall v_{793}\,(-y \in \mathbb{R} \to v_{792} \in \mathbb{R} \to v_{793} \in \mathbb{R} \to -y < v_{792} \to -y + v_{793} < v_{792} + v_{793})$
\item $\varphi_{32} :\equiv \forall v_{793}\,(-y \in \mathbb{R} \to -x \in \mathbb{R} \to v_{793} \in \mathbb{R} \to -y < -x \to -y + v_{793} < -x + v_{793})$
\item $\varphi_{33} :\equiv \forall y\,(v_{797} \in \mathbb{R} \to y \in \mathbb{R} \to v_{797} + y = y + v_{797}) \to v_{797} \in \mathbb{R} \to v_{798} \in \mathbb{R} \to v_{797} + v_{798} = v_{798} + v_{797}$
\item $\varphi_{34} :\equiv \forall x\,\forall y\,(x \in \mathbb{R} \to y \in \mathbb{R} \to x < y \to -y < -x) \to \forall y\,(x \in \mathbb{R} \to y \in \mathbb{R} \to x < y \to -y < -x)$
\item $\varphi_{35} :\equiv \forall y\,(-x \in \mathbb{R} \to y \in \mathbb{R} \to -x + y = y + -x) \to -x \in \mathbb{R} \to z \in \mathbb{R} \to -x + z = z + -x$
\item $\varphi_{36} :\equiv \forall z_{1}\,\forall w\,(z_{1} = w \to z_{1} < -x + z \to w < -x + z) \to \forall w\,(-y + z = w \to -y + z < -x + z \to w < -x + z)$
\item $\varphi_{37} :\equiv \forall w\,(-x + z = w \to z - y < -x + z \to z - y < w) \to -x + z = z - x \to z - y < -x + z \to z - y < z - x$
\item $\varphi_{38} :\equiv (z \in \mathbb{R} \to x \in \mathbb{R} \to y \in \mathbb{R} \to x < y \to x < y) \to z \in \mathbb{R} \to x \in \mathbb{R} \to y \in \mathbb{R} \to x < y \to -y < -x$
\item $\varphi_{39} :\equiv (z \in \mathbb{R} \to x \in \mathbb{R} \to y \in \mathbb{R} \to x < y \to x \in \mathbb{R}) \to z \in \mathbb{R} \to x \in \mathbb{R} \to y \in \mathbb{R} \to x < y \to -x \in \mathbb{R}$
\item $\varphi_{40} :\equiv (z \in \mathbb{R} \to x \in \mathbb{R} \to y \in \mathbb{R} \to x < y \to y \in \mathbb{R}) \to z \in \mathbb{R} \to x \in \mathbb{R} \to y \in \mathbb{R} \to x < y \to -y \in \mathbb{R}$
\item $\varphi_{41} :\equiv (-x + z = z - x \to z - y < -x + z \to z - y < z - x) \to z \in \mathbb{R} \to x \in \mathbb{R} \to y \in \mathbb{R} \to x < y \to z - y < z - x$
\item $\varphi_{42} :\equiv (z \in \mathbb{R} \to x \in \mathbb{R} \to y \in \mathbb{R} \to x < y \to z \in \mathbb{R}) \to z \in \mathbb{R} \to x \in \mathbb{R} \to y \in \mathbb{R} \to x < y \to -y + z = z + -y$
\item $\varphi_{43} :\equiv (z \in \mathbb{R} \to x \in \mathbb{R} \to y \in \mathbb{R} \to x < y \to z \in \mathbb{R}) \to z \in \mathbb{R} \to x \in \mathbb{R} \to y \in \mathbb{R} \to x < y \to -x + z = z + -x$
\item $\varphi_{44} :\equiv (-y + z = z + -y \to z + -y = z - y \to -y + z = z - y) \to z \in \mathbb{R} \to x \in \mathbb{R} \to y \in \mathbb{R} \to x < y \to -y + z = z - y$
\item $\varphi_{45} :\equiv (-x + z = z + -x \to z + -x = z - x \to -x + z = z - x) \to z \in \mathbb{R} \to x \in \mathbb{R} \to y \in \mathbb{R} \to x < y \to -x + z = z - x$
\item $\varphi_{46} :\equiv \forall w\,(-y + z = w \to -y + z < -x + z \to w < -x + z) \to -y + z = z - y \to -y + z < -x + z \to z - y < -x + z$
\item $\varphi_{47} :\equiv (z \in \mathbb{R} \to x \in \mathbb{R} \to y \in \mathbb{R} \to x < y \to -y < -x) \to z \in \mathbb{R} \to x \in \mathbb{R} \to y \in \mathbb{R} \to x < y \to -y + z < -x + z$
\item $\varphi_{48} :\equiv (-y + z = z - y \to -y + z < -x + z \to z - y < -x + z) \to z \in \mathbb{R} \to x \in \mathbb{R} \to y \in \mathbb{R} \to x < y \to z - y < -x + z$
\item $\varphi_{49} :\equiv (z \in \mathbb{R} \to x \in \mathbb{R} \to y \in \mathbb{R} \to x < y \to z - y < -x + z) \to \varphi_{41}$
\item $\varphi_{50} :\equiv (z \in \mathbb{R} \to x \in \mathbb{R} \to y \in \mathbb{R} \to x < y \to -y + z < -x + z) \to \varphi_{48}$
\item $\varphi_{51} :\equiv (z \in \mathbb{R} \to x \in \mathbb{R} \to y \in \mathbb{R} \to x < y \to -y + z = z + -y) \to z + -y = z - y \to \varphi_{44}$
\item $\varphi_{52} :\equiv (z \in \mathbb{R} \to x \in \mathbb{R} \to y \in \mathbb{R} \to x < y \to -x + z = z + -x) \to z + -x = z - x \to \varphi_{45}$
\item $\varphi_{53} :\equiv (z \in \mathbb{R} \to x \in \mathbb{R} \to y \in \mathbb{R} \to x < y \to x \in \mathbb{R}) \to (z \in \mathbb{R} \to x \in \mathbb{R} \to y \in \mathbb{R} \to x < y \to y \in \mathbb{R}) \to \varphi_{38}$
\item $\varphi_{54} :\equiv (-y \in \mathbb{R} \to z \in \mathbb{R} \to -y + z = z + -y) \to (z \in \mathbb{R} \to x \in \mathbb{R} \to y \in \mathbb{R} \to x < y \to -y \in \mathbb{R}) \to \varphi_{42}$
\item $\varphi_{55} :\equiv (-x \in \mathbb{R} \to z \in \mathbb{R} \to -x + z = z + -x) \to (z \in \mathbb{R} \to x \in \mathbb{R} \to y \in \mathbb{R} \to x < y \to -x \in \mathbb{R}) \to \varphi_{43}$
\item $\varphi_{56} :\equiv (z \in \mathbb{R} \to x \in \mathbb{R} \to y \in \mathbb{R} \to x < y \to -x \in \mathbb{R}) \to (z \in \mathbb{R} \to x \in \mathbb{R} \to y \in \mathbb{R} \to x < y \to z \in \mathbb{R}) \to \varphi_{47}$
\item $\varphi_{57} :\equiv (z \in \mathbb{R} \to x \in \mathbb{R} \to y \in \mathbb{R} \to x < y \to -x + z = z - x) \to \varphi_{49}$
\item $\varphi_{58} :\equiv (z \in \mathbb{R} \to x \in \mathbb{R} \to y \in \mathbb{R} \to x < y \to -y + z = z - y) \to \varphi_{50}$
\end{abbrevlist}
\end{fullproof}

% 自動生成 (prover/build.js): EVTBdd  ノード数 734
\begin{fullproof}{[G8]}{$\forall f\,\forall a\,\forall b\,(\mathrm{ContOn}(f,[a,b]) \to a \in \mathbb{R} \to b \in \mathbb{R} \to \exists M \in \mathbb{R}\,\forall x \in [a,b]\,(f(x) \leq M))$}{734}\phantomsection\label{tree:EVTBdd}
\noindent$\Pi_{1}$:
\begin{ptree}
\AXC{$\forall x\,(\neg x < x)$\lbl{[O1]}}
\AXC{$\forall x\,(\neg x < x) \to \neg 0 < 0$\lbl{A2}}
\BIC{$\neg 0 < 0$}
\AXC{$\neg 0 < 0 \to \mathrm{ContOn}(f,[a,b]) \to a \in \mathbb{R} \to \varphi_{37}$\lbl{A1}}
\BIC{$\mathrm{ContOn}(f,[a,b]) \to a \in \mathbb{R} \to \varphi_{37}$}
\end{ptree}
\noindent$\Pi_{2}$:
\begin{ptree}
\AXC{$\forall s\,(\varphi_{11})$\lbl{DefSubIdx}}
\AXC{$\forall s\,(\varphi_{11}) \to \varphi_{11}$\lbl{A2}}
\BIC{$\varphi_{11}$}
\AXC{$\mathrm{ContOn}(f,[a,b]) \to a \in \mathbb{R} \to \varphi_{313}$\lbl{A1}}
\AXC{$\varphi_{413}$\lbl{A1}}
\BIC{$\mathrm{ContOn}(f,[a,b]) \to a \in \mathbb{R} \to \varphi_{293}$}
\AXC{$\varphi_{427}$\lbl{A1}}
\BIC{$\varphi_{11} \to \mathrm{ContOn}(f,[a,b]) \to a \in \mathbb{R} \to \varphi_{292}$}
\BIC{$\mathrm{ContOn}(f,[a,b]) \to a \in \mathbb{R} \to \varphi_{292}$}
\end{ptree}
\noindent$\Pi_{3}$:
\begin{ptree}
\AXC{$\forall s\,\forall k\,(\mathrm{NSeq}(s) \to k \in \mathbb{N} \to s_{k} \in \mathbb{N})$\lbl{[S13]}}
\AXC{$\varphi_{79}$\lbl{A2}}
\BIC{$\forall k\,(\mathrm{NSeq}(s) \to k \in \mathbb{N} \to s_{k} \in \mathbb{N})$}
\AXC{$\varphi_{63}$\lbl{A2}}
\BIC{$\mathrm{NSeq}(s) \to K \in \mathbb{N} \to s_{K} \in \mathbb{N}$}
\AXC{$\varphi_{499}$\lbl{A1}}
\BIC{$(\mathrm{ContOn}(f,[a,b]) \to a \in \mathbb{R} \to \varphi_{292}) \to \varphi_{459}$}
\end{ptree}
\noindent$\Pi_{4}$:
\begin{ptree}
\AXC{$\stackrel{\Pi_{2}}{\vdots}$}\noLine
\UIC{$\mathrm{ContOn}(f,[a,b]) \to a \in \mathbb{R} \to \varphi_{292}$}
\AXC{$\stackrel{\Pi_{3}}{\vdots}$}\noLine
\UIC{$(\mathrm{ContOn}(f,[a,b]) \to a \in \mathbb{R} \to \varphi_{292}) \to \varphi_{459}$}
\BIC{$\varphi_{459}$}
\end{ptree}
\noindent$\Pi_{5}$:
\begin{ptree}
\AXC{$\mathrm{ContOn}(f,[a,b]) \to a \in \mathbb{R} \to \varphi_{370}$\lbl{A1}}
\AXC{$\varphi_{470}$\lbl{A1}}
\BIC{$\mathrm{ContOn}(f,[a,b]) \to a \in \mathbb{R} \to \varphi_{342}$}
\AXC{$\stackrel{\Pi_{4}}{\vdots}$}\noLine
\UIC{$\varphi_{459}$}
\BIC{$\mathrm{ContOn}(f,[a,b]) \to a \in \mathbb{R} \to \varphi_{346}$}
\end{ptree}
\noindent$\Pi_{6}$:
\begin{ptree}
\AXC{$\varphi_{227}$\lbl{A2}}
\AXC{$\mathrm{ContOn}(f,[a,b]) \to a \in \mathbb{R} \to \varphi_{284}$\lbl{A1}}
\AXC{$\varphi_{401}$\lbl{A1}}
\BIC{$\mathrm{ContOn}(f,[a,b]) \to a \in \mathbb{R} \to b \in \mathbb{R} \to \varphi_{265}$}
\AXC{$\varphi_{433}$\lbl{A1}}
\BIC{$\varphi_{227} \to \mathrm{ContOn}(f,[a,b]) \to \varphi_{289}$}
\BIC{$\mathrm{ContOn}(f,[a,b]) \to \varphi_{289}$}
\end{ptree}
\noindent$\Pi_{7}$:
\begin{ptree}
\AXC{$\stackrel{\Pi_{5}}{\vdots}$}\noLine
\UIC{$\mathrm{ContOn}(f,[a,b]) \to a \in \mathbb{R} \to \varphi_{346}$}
\AXC{$\stackrel{\Pi_{6}}{\vdots}$}\noLine
\UIC{$\mathrm{ContOn}(f,[a,b]) \to \varphi_{289}$}
\AXC{$(\mathrm{ContOn}(f,[a,b]) \to \varphi_{289}) \to \varphi_{481}$\lbl{A1}}
\BIC{$\varphi_{481}$}
\BIC{$\mathrm{ContOn}(f,[a,b]) \to a \in \mathbb{R} \to \varphi_{388}$}
\end{ptree}
\noindent$\Pi_{8}$:
\begin{ptree}
\AXC{$\forall s\,(\mathrm{SubIdx}(s) \to \forall k \in \mathbb{N}\,(k \leq s_{k}))$\lbl{[C17]}}
\AXC{$\varphi_{57}$\lbl{A2}}
\BIC{$\mathrm{SubIdx}(s) \to \forall k \in \mathbb{N}\,(k \leq s_{k})$}
\AXC{$(\mathrm{SubIdx}(s) \to \forall k \in \mathbb{N}\,(k \leq s_{k})) \to \varphi_{439}$\lbl{A1}}
\BIC{$\varphi_{439}$}
\end{ptree}
\noindent$\Pi_{9}$:
\begin{ptree}
\AXC{$\mathrm{ContOn}(f,[a,b]) \to a \in \mathbb{R} \to \varphi_{313}$\lbl{A1}}
\AXC{$\varphi_{413}$\lbl{A1}}
\BIC{$\mathrm{ContOn}(f,[a,b]) \to a \in \mathbb{R} \to \varphi_{293}$}
\AXC{$\stackrel{\Pi_{8}}{\vdots}$}\noLine
\UIC{$\varphi_{439}$}
\BIC{$\mathrm{ContOn}(f,[a,b]) \to a \in \mathbb{R} \to \varphi_{359}$}
\AXC{$(\mathrm{ContOn}(f,[a,b]) \to a \in \mathbb{R} \to \varphi_{359}) \to \varphi_{397}$\lbl{A1}}
\BIC{$\varphi_{397}$}
\end{ptree}
\noindent$\Pi_{10}$:
\begin{ptree}
\AXC{$\forall k \in \mathbb{N}\,(k \leq s_{k}) \to K \in \mathbb{N} \to K \leq s_{K}$\lbl{A2}}
\AXC{$\stackrel{\Pi_{9}}{\vdots}$}\noLine
\UIC{$\varphi_{397}$}
\BIC{$\mathrm{ContOn}(f,[a,b]) \to a \in \mathbb{R} \to \varphi_{354}$}
\AXC{$(\mathrm{ContOn}(f,[a,b]) \to a \in \mathbb{R} \to \varphi_{354}) \to \varphi_{457}$\lbl{A1}}
\BIC{$\varphi_{457}$}
\end{ptree}
\noindent$\Pi_{11}$:
\begin{ptree}
\AXC{$\forall m\,\forall n\,(m \in \mathbb{N} \to n \in \mathbb{N} \to m \leq n \to \underline{m} \leq \underline{n})$\lbl{[N16]}}
\AXC{$\varphi_{125}$\lbl{A2}}
\BIC{$\forall n\,(K \in \mathbb{N} \to n \in \mathbb{N} \to K \leq n \to \underline{K} \leq \underline{n})$}
\AXC{$\varphi_{156}$\lbl{A2}}
\BIC{$K \in \mathbb{N} \to s_{K} \in \mathbb{N} \to K \leq s_{K} \to \underline{K} \leq \underline{s_{K}}$}
\end{ptree}
\noindent$\Pi_{12}$:
\begin{ptree}
\AXC{$\stackrel{\Pi_{11}}{\vdots}$}\noLine
\UIC{$K \in \mathbb{N} \to s_{K} \in \mathbb{N} \to K \leq s_{K} \to \underline{K} \leq \underline{s_{K}}$}
\AXC{$(K \in \mathbb{N} \to s_{K} \in \mathbb{N} \to K \leq s_{K} \to \underline{K} \leq \underline{s_{K}}) \to \varphi_{510}$\lbl{A1}}
\BIC{$\varphi_{510}$}
\end{ptree}
\noindent$\Pi_{13}$:
\begin{ptree}
\AXC{$\stackrel{\Pi_{5}}{\vdots}$}\noLine
\UIC{$\mathrm{ContOn}(f,[a,b]) \to a \in \mathbb{R} \to \varphi_{346}$}
\AXC{$\mathrm{ContOn}(f,[a,b]) \to a \in \mathbb{R} \to \varphi_{370}$\lbl{A1}}
\AXC{$\varphi_{470}$\lbl{A1}}
\BIC{$\mathrm{ContOn}(f,[a,b]) \to a \in \mathbb{R} \to \varphi_{342}$}
\AXC{$\stackrel{\Pi_{12}}{\vdots}$}\noLine
\UIC{$\varphi_{510}$}
\BIC{$(\mathrm{ContOn}(f,[a,b]) \to a \in \mathbb{R} \to \varphi_{346}) \to \varphi_{461}$}
\BIC{$\varphi_{461}$}
\end{ptree}
\noindent$\Pi_{14}$:
\begin{ptree}
\AXC{$\forall n \in \mathbb{N}\,(\underline{n} \in \mathbb{R})$\lbl{[N4]}}
\AXC{$\forall n \in \mathbb{N}\,(\underline{n} \in \mathbb{R}) \to K \in \mathbb{N} \to \underline{K} \in \mathbb{R}$\lbl{A2}}
\BIC{$K \in \mathbb{N} \to \underline{K} \in \mathbb{R}$}
\AXC{$(K \in \mathbb{N} \to \underline{K} \in \mathbb{R}) \to \varphi_{460}$\lbl{A1}}
\BIC{$\varphi_{460}$}
\end{ptree}
\noindent$\Pi_{15}$:
\begin{ptree}
\AXC{$\forall x \in \mathbb{R}\,(x \cdot 1 = x)$\lbl{(R8)}}
\AXC{$\forall x \in \mathbb{R}\,(x \cdot 1 = x) \to \underline{K} \in \mathbb{R} \to \underline{K} \cdot 1 = \underline{K}$\lbl{A2}}
\BIC{$\underline{K} \in \mathbb{R} \to \underline{K} \cdot 1 = \underline{K}$}
\AXC{$(\underline{K} \in \mathbb{R} \to \underline{K} \cdot 1 = \underline{K}) \to \varphi_{463}$\lbl{A1}}
\BIC{$\varphi_{463}$}
\end{ptree}
\noindent$\Pi_{16}$:
\begin{ptree}
\AXC{$\mathrm{ContOn}(f,[a,b]) \to a \in \mathbb{R} \to \varphi_{370}$\lbl{A1}}
\AXC{$\varphi_{470}$\lbl{A1}}
\BIC{$\mathrm{ContOn}(f,[a,b]) \to a \in \mathbb{R} \to \varphi_{342}$}
\AXC{$\stackrel{\Pi_{14}}{\vdots}$}\noLine
\UIC{$\varphi_{460}$}
\BIC{$\mathrm{ContOn}(f,[a,b]) \to a \in \mathbb{R} \to \varphi_{347}$}
\AXC{$\stackrel{\Pi_{15}}{\vdots}$}\noLine
\UIC{$\varphi_{463}$}
\BIC{$\mathrm{ContOn}(f,[a,b]) \to a \in \mathbb{R} \to \varphi_{351}$}
\AXC{$\varphi_{506}$\lbl{A1}}
\BIC{$(\mathrm{ContOn}(f,[a,b]) \to a \in \mathbb{R} \to \varphi_{361}) \to \varphi_{399}$}
\end{ptree}
\noindent$\Pi_{17}$:
\begin{ptree}
\AXC{$\mathrm{ContOn}(f,[a,b]) \to a \in \mathbb{R} \to \varphi_{370}$\lbl{A1}}
\AXC{$\varphi_{478}$\lbl{A1}}
\BIC{$\mathrm{ContOn}(f,[a,b]) \to a \in \mathbb{R} \to \varphi_{361}$}
\AXC{$\stackrel{\Pi_{16}}{\vdots}$}\noLine
\UIC{$(\mathrm{ContOn}(f,[a,b]) \to a \in \mathbb{R} \to \varphi_{361}) \to \varphi_{399}$}
\BIC{$\varphi_{399}$}
\end{ptree}
\noindent$\Pi_{18}$:
\begin{ptree}
\AXC{$\forall x \in \mathbb{R}\,(\lvert x \rvert \in \mathbb{R})$\lbl{[A5]}}
\AXC{$\forall x \in \mathbb{R}\,(\lvert x \rvert \in \mathbb{R}) \to M \in \mathbb{R} \to \lvert M \rvert \in \mathbb{R}$\lbl{A2}}
\BIC{$M \in \mathbb{R} \to \lvert M \rvert \in \mathbb{R}$}
\AXC{$(M \in \mathbb{R} \to \lvert M \rvert \in \mathbb{R}) \to \varphi_{450}$\lbl{A1}}
\BIC{$\varphi_{450}$}
\end{ptree}
\noindent$\Pi_{19}$:
\begin{ptree}
\AXC{$\forall x \in \mathbb{R}\,(x < x + 1)$\lbl{[N13]}}
\AXC{$\forall x \in \mathbb{R}\,(x < x + 1) \to \lvert M \rvert \in \mathbb{R} \to \lvert M \rvert < \lvert M \rvert + 1$\lbl{A2}}
\BIC{$\lvert M \rvert \in \mathbb{R} \to \lvert M \rvert < \lvert M \rvert + 1$}
\end{ptree}
\noindent$\Pi_{20}$:
\begin{ptree}
\AXC{$\stackrel{\Pi_{19}}{\vdots}$}\noLine
\UIC{$\lvert M \rvert \in \mathbb{R} \to \lvert M \rvert < \lvert M \rvert + 1$}
\AXC{$(\lvert M \rvert \in \mathbb{R} \to \lvert M \rvert < \lvert M \rvert + 1) \to \varphi_{464}$\lbl{A1}}
\BIC{$\varphi_{464}$}
\end{ptree}
\noindent$\Pi_{21}$:
\begin{ptree}
\AXC{$\mathrm{ContOn}(f,[a,b]) \to a \in \mathbb{R} \to \varphi_{352}$\lbl{A1}}
\AXC{$\varphi_{451}$\lbl{A1}}
\BIC{$\mathrm{ContOn}(f,[a,b]) \to a \in \mathbb{R} \to \varphi_{322}$}
\AXC{$\stackrel{\Pi_{18}}{\vdots}$}\noLine
\UIC{$\varphi_{450}$}
\BIC{$\mathrm{ContOn}(f,[a,b]) \to a \in \mathbb{R} \to \varphi_{344}$}
\AXC{$\stackrel{\Pi_{20}}{\vdots}$}\noLine
\UIC{$\varphi_{464}$}
\BIC{$\mathrm{ContOn}(f,[a,b]) \to a \in \mathbb{R} \to \varphi_{353}$}
\end{ptree}
\noindent$\Pi_{22}$:
\begin{ptree}
\AXC{$\forall x\,\forall y\,\forall z\,(x < y \to y < z \to x < z)$\lbl{[O2]}}
\AXC{$\varphi_{48}$\lbl{A2}}
\BIC{$\forall y\,\forall z\,(\lvert M \rvert < y \to y < z \to \lvert M \rvert < z)$}
\AXC{$\varphi_{121}$\lbl{A2}}
\BIC{$\forall z\,(\lvert M \rvert < \lvert M \rvert + 1 \to \lvert M \rvert + 1 < z \to \lvert M \rvert < z)$}
\AXC{$\varphi_{172}$\lbl{A2}}
\BIC{$\lvert M \rvert < \lvert M \rvert + 1 \to \lvert M \rvert + 1 < \underline{K} \to \lvert M \rvert < \underline{K}$}
\AXC{$\varphi_{504}$\lbl{A1}}
\BIC{$(\mathrm{ContOn}(f,[a,b]) \to a \in \mathbb{R} \to \varphi_{353}) \to \varphi_{466}$}
\end{ptree}
\noindent$\Pi_{23}$:
\begin{ptree}
\AXC{$\stackrel{\Pi_{21}}{\vdots}$}\noLine
\UIC{$\mathrm{ContOn}(f,[a,b]) \to a \in \mathbb{R} \to \varphi_{353}$}
\AXC{$\stackrel{\Pi_{22}}{\vdots}$}\noLine
\UIC{$(\mathrm{ContOn}(f,[a,b]) \to a \in \mathbb{R} \to \varphi_{353}) \to \varphi_{466}$}
\BIC{$\varphi_{466}$}
\end{ptree}
\noindent$\Pi_{24}$:
\begin{ptree}
\AXC{$\forall x \in \mathbb{R}\,(-\lvert x \rvert \leq x \land x \leq \lvert x \rvert)$\lbl{[A8]}}
\AXC{$\varphi_{58}$\lbl{A2}}
\BIC{$M \in \mathbb{R} \to -\lvert M \rvert \leq M \land M \leq \lvert M \rvert$}
\AXC{$(M \in \mathbb{R} \to -\lvert M \rvert \leq M \land M \leq \lvert M \rvert) \to \varphi_{452}$\lbl{A1}}
\BIC{$\varphi_{452}$}
\end{ptree}
\noindent$\Pi_{25}$:
\begin{ptree}
\AXC{$z = w \to \lvert z \rvert = M \to \lvert w \rvert = M$\lbl{A6}}
\UIC{$\forall w\,(z = w \to \lvert z \rvert = M \to \lvert w \rvert = M)$}
\UIC{$\forall z\,\forall w\,(z = w \to \lvert z \rvert = M \to \lvert w \rvert = M)$}
\AXC{$\forall z\,\forall w\,(z = w \to \lvert z \rvert = M \to \lvert w \rvert = M) \to \varphi_{59}$\lbl{A2}}
\BIC{$\varphi_{59}$}
\AXC{$\varphi_{59} \to \varphi_{159}$\lbl{A2}}
\BIC{$\varphi_{159}$}
\end{ptree}
\noindent$\Pi_{26}$:
\begin{ptree}
\AXC{$z = w \to \lvert z \rvert < M \to \lvert w \rvert < M$\lbl{A6}}
\UIC{$\forall w\,(z = w \to \lvert z \rvert < M \to \lvert w \rvert < M)$}
\UIC{$\forall z\,\forall w\,(z = w \to \lvert z \rvert < M \to \lvert w \rvert < M)$}
\AXC{$\forall z\,\forall w\,(z = w \to \lvert z \rvert < M \to \lvert w \rvert < M) \to \varphi_{75}$\lbl{A2}}
\BIC{$\varphi_{75}$}
\AXC{$\varphi_{75} \to \varphi_{170}$\lbl{A2}}
\BIC{$\varphi_{170}$}
\end{ptree}
\noindent$\Pi_{27}$:
\begin{ptree}
\AXC{$\varphi_{117}$\lbl{A2}}
\AXC{$\mathrm{ContOn}(f,[a,b]) \to a \in \mathbb{R} \to \varphi_{352}$\lbl{A1}}
\AXC{$\varphi_{479}$\lbl{A1}}
\BIC{$\mathrm{ContOn}(f,[a,b]) \to a \in \mathbb{R} \to \varphi_{379}$}
\AXC{$\varphi_{494}$\lbl{A1}}
\BIC{$\varphi_{117} \to \mathrm{ContOn}(f,[a,b]) \to a \in \mathbb{R} \to \varphi_{375}$}
\BIC{$\mathrm{ContOn}(f,[a,b]) \to a \in \mathbb{R} \to \varphi_{375}$}
\end{ptree}
\noindent$\Pi_{28}$:
\begin{ptree}
\AXC{$\stackrel{\Pi_{27}}{\vdots}$}\noLine
\UIC{$\mathrm{ContOn}(f,[a,b]) \to a \in \mathbb{R} \to \varphi_{375}$}
\AXC{$(\mathrm{ContOn}(f,[a,b]) \to a \in \mathbb{R} \to \varphi_{375}) \to \varphi_{468}$\lbl{A1}}
\BIC{$\varphi_{468}$}
\end{ptree}
\noindent$\Pi_{29}$:
\begin{ptree}
\AXC{$y = z \to x = y \to x = z$\lbl{A6}}
\AXC{$(y = z \to x = y \to x = z) \to x = y \to y = z \to x = z$\lbl{A1}}
\BIC{$x = y \to y = z \to x = z$}
\UIC{$\forall z\,(x = y \to y = z \to x = z)$}
\UIC{$\forall y\,\forall z\,(x = y \to y = z \to x = z)$}
\UIC{$\forall x\,\forall y\,\forall z\,(x = y \to y = z \to x = z)$}
\AXC{$\forall x\,\forall y\,\forall z\,(x = y \to y = z \to x = z) \to \varphi_{42}$\lbl{A2}}
\BIC{$\varphi_{42}$}
\end{ptree}
\noindent$\Pi_{30}$:
\begin{ptree}
\AXC{$\varphi_{86}$\lbl{[S12]}}
\AXC{$\varphi_{86} \to \forall s\,(\mathrm{Seq}(q) \to \varphi_{46})$\lbl{A2}}
\BIC{$\forall s\,(\mathrm{Seq}(q) \to \varphi_{46})$}
\AXC{$\forall s\,(\mathrm{Seq}(q) \to \varphi_{46}) \to \mathrm{Seq}(q) \to \varphi_{46}$\lbl{A2}}
\BIC{$\mathrm{Seq}(q) \to \varphi_{46}$}
\AXC{$(\mathrm{Seq}(q) \to \varphi_{46}) \to \varphi_{469}$\lbl{A1}}
\BIC{$\varphi_{469}$}
\end{ptree}
\noindent$\Pi_{31}$:
\begin{ptree}
\AXC{$\stackrel{\Pi_{2}}{\vdots}$}\noLine
\UIC{$\mathrm{ContOn}(f,[a,b]) \to a \in \mathbb{R} \to \varphi_{292}$}
\AXC{$\mathrm{ContOn}(f,[a,b]) \to a \in \mathbb{R} \to \varphi_{284}$\lbl{A1}}
\AXC{$\varphi_{391}$\lbl{A1}}
\BIC{$\mathrm{ContOn}(f,[a,b]) \to a \in \mathbb{R} \to b \in \mathbb{R} \to \varphi_{195}$}
\AXC{$\stackrel{\Pi_{30}}{\vdots}$}\noLine
\UIC{$\varphi_{469}$}
\BIC{$\varphi_{422}$}
\BIC{$\mathrm{ContOn}(f,[a,b]) \to a \in \mathbb{R} \to \varphi_{317}$}
\end{ptree}
\noindent$\Pi_{32}$:
\begin{ptree}
\AXC{$\varphi_{104}$\lbl{A2}}
\AXC{$\stackrel{\Pi_{31}}{\vdots}$}\noLine
\UIC{$\mathrm{ContOn}(f,[a,b]) \to a \in \mathbb{R} \to \varphi_{317}$}
\AXC{$\varphi_{454}$\lbl{A1}}
\BIC{$\mathrm{ContOn}(f,[a,b]) \to a \in \mathbb{R} \to \varphi_{377}$}
\AXC{$\varphi_{493}$\lbl{A1}}
\BIC{$\varphi_{104} \to \mathrm{ContOn}(f,[a,b]) \to a \in \mathbb{R} \to \varphi_{374}$}
\BIC{$\mathrm{ContOn}(f,[a,b]) \to a \in \mathbb{R} \to \varphi_{374}$}
\end{ptree}
\noindent$\Pi_{33}$:
\begin{ptree}
\AXC{$\stackrel{\Pi_{32}}{\vdots}$}\noLine
\UIC{$\mathrm{ContOn}(f,[a,b]) \to a \in \mathbb{R} \to \varphi_{374}$}
\AXC{$(\mathrm{ContOn}(f,[a,b]) \to a \in \mathbb{R} \to \varphi_{374}) \to \varphi_{467}$\lbl{A1}}
\BIC{$\varphi_{467}$}
\end{ptree}
\noindent$\Pi_{34}$:
\begin{ptree}
\AXC{$\mathrm{ContOn}(f,[a,b]) \to a \in \mathbb{R} \to \varphi_{370}$\lbl{A1}}
\AXC{$\varphi_{470}$\lbl{A1}}
\BIC{$\mathrm{ContOn}(f,[a,b]) \to a \in \mathbb{R} \to \varphi_{342}$}
\AXC{$\stackrel{\Pi_{33}}{\vdots}$}\noLine
\UIC{$\varphi_{467}$}
\BIC{$\mathrm{ContOn}(f,[a,b]) \to a \in \mathbb{R} \to \varphi_{367}$}
\AXC{$\varphi_{495}$\lbl{A1}}
\BIC{$f(q \circ s(K)) = f(q \circ s(K)) \to \varphi_{405}$}
\end{ptree}
\noindent$\Pi_{35}$:
\begin{ptree}
\AXC{$x = x$\lbl{A5}}
\UIC{$\forall x\,(x = x)$}
\AXC{$\forall x\,(x = x) \to f(q \circ s(K)) = f(q \circ s(K))$\lbl{A2}}
\BIC{$f(q \circ s(K)) = f(q \circ s(K))$}
\AXC{$\stackrel{\Pi_{34}}{\vdots}$}\noLine
\UIC{$f(q \circ s(K)) = f(q \circ s(K)) \to \varphi_{405}$}
\BIC{$\varphi_{405}$}
\end{ptree}
\noindent$\Pi_{36}$:
\begin{ptree}
\AXC{$\varphi_{43}$\lbl{A6}}
\UIC{$\forall w\,(\varphi_{43})$}
\UIC{$\forall z\,\forall w\,(\varphi_{43})$}
\AXC{$\forall z\,\forall w\,(\varphi_{43}) \to \varphi_{129}$\lbl{A2}}
\BIC{$\varphi_{129}$}
\AXC{$\varphi_{129} \to q \circ s(K) = q(s_{K}) \to \varphi_{112}$\lbl{A2}}
\BIC{$q \circ s(K) = q(s_{K}) \to \varphi_{112}$}
\AXC{$\stackrel{\Pi_{35}}{\vdots}$}\noLine
\UIC{$\varphi_{405}$}
\BIC{$\mathrm{ContOn}(f,[a,b]) \to a \in \mathbb{R} \to \varphi_{376}$}
\end{ptree}
\noindent$\Pi_{37}$:
\begin{ptree}
\AXC{$z = w \to z \in [a,b] \to w \in [a,b]$\lbl{A6}}
\UIC{$\forall w\,(z = w \to z \in [a,b] \to w \in [a,b])$}
\UIC{$\forall z\,\forall w\,(z = w \to z \in [a,b] \to w \in [a,b])$}
\AXC{$\forall z\,\forall w\,(z = w \to z \in [a,b] \to w \in [a,b]) \to \varphi_{34}$\lbl{A2}}
\BIC{$\varphi_{34}$}
\AXC{$\varphi_{34} \to \varphi_{114}$\lbl{A2}}
\BIC{$\varphi_{114}$}
\end{ptree}
\noindent$\Pi_{38}$:
\begin{ptree}
\AXC{$\forall s\,\forall k\,(\mathrm{NSeq}(s) \to k \in \mathbb{N} \to s_{k} \in \mathbb{N})$\lbl{[S13]}}
\AXC{$\varphi_{79}$\lbl{A2}}
\BIC{$\forall k\,(\mathrm{NSeq}(s) \to k \in \mathbb{N} \to s_{k} \in \mathbb{N})$}
\AXC{$\varphi_{64}$\lbl{A2}}
\BIC{$\mathrm{NSeq}(s) \to n \in \mathbb{N} \to s_{n} \in \mathbb{N}$}
\AXC{$\varphi_{492}$\lbl{A1}}
\BIC{$(\mathrm{ContOn}(f,[a,b]) \to a \in \mathbb{R} \to \varphi_{292}) \to \varphi_{411}$}
\end{ptree}
\noindent$\Pi_{39}$:
\begin{ptree}
\AXC{$\stackrel{\Pi_{2}}{\vdots}$}\noLine
\UIC{$\mathrm{ContOn}(f,[a,b]) \to a \in \mathbb{R} \to \varphi_{292}$}
\AXC{$\stackrel{\Pi_{38}}{\vdots}$}\noLine
\UIC{$(\mathrm{ContOn}(f,[a,b]) \to a \in \mathbb{R} \to \varphi_{292}) \to \varphi_{411}$}
\BIC{$\varphi_{411}$}
\end{ptree}
\noindent$\Pi_{40}$:
\begin{ptree}
\AXC{$\varphi_{188}$\lbl{A2}}
\AXC{$\mathrm{ContOn}(f,[a,b]) \to a \in \mathbb{R} \to \varphi_{284}$\lbl{A1}}
\AXC{$\varphi_{401}$\lbl{A1}}
\BIC{$\mathrm{ContOn}(f,[a,b]) \to a \in \mathbb{R} \to b \in \mathbb{R} \to \varphi_{265}$}
\AXC{$\varphi_{414}$\lbl{A1}}
\BIC{$\varphi_{188} \to \mathrm{ContOn}(f,[a,b]) \to a \in \mathbb{R} \to b \in \mathbb{R} \to \varphi_{264}$}
\BIC{$\mathrm{ContOn}(f,[a,b]) \to a \in \mathbb{R} \to b \in \mathbb{R} \to \varphi_{264}$}
\AXC{$\varphi_{431}$\lbl{A1}}
\BIC{$\varphi_{384}$}
\end{ptree}
\noindent$\Pi_{41}$:
\begin{ptree}
\AXC{$\mathrm{ContOn}(f,[a,b]) \to a \in \mathbb{R} \to b \in \mathbb{R} \to \varphi_{224}$\lbl{A1}}
\AXC{$\stackrel{\Pi_{40}}{\vdots}$}\noLine
\UIC{$\varphi_{384}$}
\BIC{$\mathrm{ContOn}(f,[a,b]) \to a \in \mathbb{R} \to b \in \mathbb{R} \to \varphi_{228}$}
\AXC{$\varphi_{389}$\lbl{A1}}
\BIC{$\varphi_{113} \land \varphi_{47} \to n \in \mathbb{N} \to q(n) \in [a,b]$}
\UIC{$\forall n\,(\varphi_{113} \land \varphi_{47} \to n \in \mathbb{N} \to q(n) \in [a,b])$}
\end{ptree}
\noindent$\Pi_{42}$:
\begin{ptree}
\AXC{$\varphi_{390}$\lbl{A3}}
\AXC{$\stackrel{\Pi_{41}}{\vdots}$}\noLine
\UIC{$\forall n\,(\varphi_{113} \land \varphi_{47} \to n \in \mathbb{N} \to q(n) \in [a,b])$}
\AXC{$\forall n\,(\varphi_{113} \land \varphi_{47} \to n \in \mathbb{N} \to q(n) \in [a,b]) \to \varphi_{390} \to \varphi_{406}$\lbl{A1}}
\BIC{$\varphi_{390} \to \varphi_{406}$}
\BIC{$\varphi_{406}$}
\end{ptree}
\noindent$\Pi_{43}$:
\begin{ptree}
\AXC{$\varphi_{113} \land \varphi_{47} \to \forall n\,(\varphi_{113} \land \varphi_{47})$\lbl{A4}}
\AXC{$\stackrel{\Pi_{42}}{\vdots}$}\noLine
\UIC{$\varphi_{406}$}
\BIC{$\mathrm{ContOn}(f,[a,b]) \to a \in \mathbb{R} \to b \in \mathbb{R} \to \varphi_{233}$}
\end{ptree}
\noindent$\Pi_{44}$:
\begin{ptree}
\AXC{$\varphi_{80}$\lbl{A2}}
\AXC{$\stackrel{\Pi_{43}}{\vdots}$}\noLine
\UIC{$\mathrm{ContOn}(f,[a,b]) \to a \in \mathbb{R} \to b \in \mathbb{R} \to \varphi_{233}$}
\AXC{$\varphi_{403}$\lbl{A1}}
\BIC{$\varphi_{80} \to \mathrm{ContOn}(f,[a,b]) \to a \in \mathbb{R} \to b \in \mathbb{R} \to \varphi_{246}$}
\BIC{$\mathrm{ContOn}(f,[a,b]) \to a \in \mathbb{R} \to b \in \mathbb{R} \to \varphi_{246}$}
\AXC{$\varphi_{480}$\lbl{A1}}
\BIC{$\varphi_{419}$}
\end{ptree}
\noindent$\Pi_{45}$:
\begin{ptree}
\AXC{$x = x$\lbl{A5}}
\AXC{$x = y \to x = x \to y = x$\lbl{A6}}
\AXC{$(x = y \to x = x \to y = x) \to x = x \to x = y \to y = x$\lbl{A1}}
\BIC{$x = x \to x = y \to y = x$}
\BIC{$x = y \to y = x$}
\UIC{$\forall y\,(x = y \to y = x)$}
\UIC{$\forall x\,\forall y\,(x = y \to y = x)$}
\AXC{$\varphi_{12}$\lbl{A2}}
\BIC{$\forall y\,(q \circ s(n) = y \to y = q \circ s(n))$}
\AXC{$\varphi_{142}$\lbl{A2}}
\BIC{$q \circ s(n) = q(s_{n}) \to q(s_{n}) = q \circ s(n)$}
\end{ptree}
\noindent$\Pi_{46}$:
\begin{ptree}
\AXC{$\varphi_{105}$\lbl{A2}}
\AXC{$\stackrel{\Pi_{31}}{\vdots}$}\noLine
\UIC{$\mathrm{ContOn}(f,[a,b]) \to a \in \mathbb{R} \to \varphi_{317}$}
\AXC{$\varphi_{434}$\lbl{A1}}
\BIC{$\mathrm{ContOn}(f,[a,b]) \to a \in \mathbb{R} \to \varphi_{315}$}
\AXC{$\varphi_{453}$\lbl{A1}}
\BIC{$\varphi_{105} \to \mathrm{ContOn}(f,[a,b]) \to a \in \mathbb{R} \to \varphi_{314}$}
\BIC{$\mathrm{ContOn}(f,[a,b]) \to a \in \mathbb{R} \to \varphi_{314}$}
\end{ptree}
\noindent$\Pi_{47}$:
\begin{ptree}
\AXC{$\stackrel{\Pi_{46}}{\vdots}$}\noLine
\UIC{$\mathrm{ContOn}(f,[a,b]) \to a \in \mathbb{R} \to \varphi_{314}$}
\AXC{$(\mathrm{ContOn}(f,[a,b]) \to a \in \mathbb{R} \to \varphi_{314}) \to \varphi_{418}$\lbl{A1}}
\BIC{$\varphi_{418}$}
\end{ptree}
\noindent$\Pi_{48}$:
\begin{ptree}
\AXC{$\mathrm{ContOn}(f,[a,b]) \to a \in \mathbb{R} \to \varphi_{298}$\lbl{A1}}
\AXC{$\stackrel{\Pi_{47}}{\vdots}$}\noLine
\UIC{$\varphi_{418}$}
\BIC{$\mathrm{ContOn}(f,[a,b]) \to a \in \mathbb{R} \to \varphi_{310}$}
\AXC{$(\mathrm{ContOn}(f,[a,b]) \to a \in \mathbb{R} \to \varphi_{310}) \to \varphi_{373}$\lbl{A1}}
\BIC{$\varphi_{373}$}
\end{ptree}
\noindent$\Pi_{49}$:
\begin{ptree}
\AXC{$\stackrel{\Pi_{45}}{\vdots}$}\noLine
\UIC{$q \circ s(n) = q(s_{n}) \to q(s_{n}) = q \circ s(n)$}
\AXC{$\stackrel{\Pi_{48}}{\vdots}$}\noLine
\UIC{$\varphi_{373}$}
\BIC{$\mathrm{ContOn}(f,[a,b]) \to a \in \mathbb{R} \to \varphi_{311}$}
\AXC{$(\mathrm{ContOn}(f,[a,b]) \to a \in \mathbb{R} \to \varphi_{311}) \to \varphi_{448}$\lbl{A1}}
\BIC{$\varphi_{448}$}
\end{ptree}
\noindent$\Pi_{50}$:
\begin{ptree}
\AXC{$\stackrel{\Pi_{37}}{\vdots}$}\noLine
\UIC{$\varphi_{114}$}
\AXC{$\mathrm{ContOn}(f,[a,b]) \to a \in \mathbb{R} \to \varphi_{298}$\lbl{A1}}
\AXC{$\stackrel{\Pi_{39}}{\vdots}$}\noLine
\UIC{$\varphi_{411}$}
\BIC{$\mathrm{ContOn}(f,[a,b]) \to a \in \mathbb{R} \to \varphi_{301}$}
\AXC{$\stackrel{\Pi_{44}}{\vdots}$}\noLine
\UIC{$\varphi_{419}$}
\BIC{$\mathrm{ContOn}(f,[a,b]) \to a \in \mathbb{R} \to \varphi_{307}$}
\AXC{$\stackrel{\Pi_{49}}{\vdots}$}\noLine
\UIC{$\varphi_{448}$}
\BIC{$\varphi_{114} \to \mathrm{ContOn}(f,[a,b]) \to a \in \mathbb{R} \to \varphi_{308}$}
\BIC{$\mathrm{ContOn}(f,[a,b]) \to a \in \mathbb{R} \to \varphi_{308}$}
\AXC{$\varphi_{429}$\lbl{A1}}
\BIC{$\varphi_{303} \to n \in \mathbb{N} \to q \circ s(n) \in [a,b]$}
\UIC{$\forall n\,(\varphi_{303} \to n \in \mathbb{N} \to q \circ s(n) \in [a,b])$}
\end{ptree}
\noindent$\Pi_{51}$:
\begin{ptree}
\AXC{$\varphi_{432}$\lbl{A3}}
\AXC{$\stackrel{\Pi_{50}}{\vdots}$}\noLine
\UIC{$\forall n\,(\varphi_{303} \to n \in \mathbb{N} \to q \circ s(n) \in [a,b])$}
\AXC{$\forall n\,(\varphi_{303} \to n \in \mathbb{N} \to q \circ s(n) \in [a,b]) \to \varphi_{432} \to \varphi_{487}$\lbl{A1}}
\BIC{$\varphi_{432} \to \varphi_{487}$}
\BIC{$\varphi_{487}$}
\end{ptree}
\noindent$\Pi_{52}$:
\begin{ptree}
\AXC{$\varphi_{303} \to \forall n\,(\varphi_{303})$\lbl{A4}}
\AXC{$\stackrel{\Pi_{51}}{\vdots}$}\noLine
\UIC{$\varphi_{487}$}
\BIC{$\mathrm{ContOn}(f,[a,b]) \to a \in \mathbb{R} \to \varphi_{312}$}
\end{ptree}
\noindent$\Pi_{53}$:
\begin{ptree}
\AXC{$\forall f\,(\varphi_{27})$\lbl{DefContOn}}
\AXC{$\forall f\,(\varphi_{27}) \to \varphi_{27}$\lbl{A2}}
\BIC{$\varphi_{27}$}
\AXC{$\varphi_{27} \to \varphi_{115}$\lbl{A2}}
\BIC{$\varphi_{115}$}
\AXC{$\varphi_{1}$\lbl{A1}}
\AXC{$\varphi_{1} \to \varphi_{115} \to \varphi_{28}$\lbl{A1}}
\BIC{$\varphi_{115} \to \varphi_{28}$}
\BIC{$\varphi_{28}$}
\end{ptree}
\noindent$\Pi_{54}$:
\begin{ptree}
\AXC{$\forall X\,\forall Y\,\forall x\,(X \subseteq Y \to x \in X \to x \in Y)$\lbl{[B5]}}
\AXC{$\varphi_{90}$\lbl{A2}}
\BIC{$\forall Y\,\forall x\,([a,b] \subseteq Y \to x \in [a,b] \to x \in Y)$}
\AXC{$\forall Y\,\forall x\,([a,b] \subseteq Y \to x \in [a,b] \to x \in Y) \to \varphi_{7}$\lbl{A2}}
\BIC{$\varphi_{7}$}
\AXC{$\varphi_{7} \to \varphi_{68}$\lbl{A2}}
\BIC{$\varphi_{68}$}
\end{ptree}
\noindent$\Pi_{55}$:
\begin{ptree}
\AXC{$\varphi_{89}$\lbl{A2}}
\AXC{$\stackrel{\Pi_{52}}{\vdots}$}\noLine
\UIC{$\mathrm{ContOn}(f,[a,b]) \to a \in \mathbb{R} \to \varphi_{312}$}
\AXC{$\varphi_{447}$\lbl{A1}}
\BIC{$\varphi_{89} \to \mathrm{ContOn}(f,[a,b]) \to a \in \mathbb{R} \to \varphi_{308}$}
\BIC{$\mathrm{ContOn}(f,[a,b]) \to a \in \mathbb{R} \to \varphi_{308}$}
\AXC{$\stackrel{\Pi_{53}}{\vdots}$}\noLine
\UIC{$\varphi_{28}$}
\AXC{$\stackrel{\Pi_{54}}{\vdots}$}\noLine
\UIC{$\varphi_{68}$}
\AXC{$\varphi_{68} \to \varphi_{28} \to \varphi_{428}$\lbl{A1}}
\BIC{$\varphi_{28} \to \varphi_{428}$}
\BIC{$\varphi_{428}$}
\BIC{$\mathrm{ContOn}(f,[a,b]) \to a \in \mathbb{R} \to \varphi_{306}$}
\end{ptree}
\noindent$\Pi_{56}$:
\begin{ptree}
\AXC{$\varphi_{430}$\lbl{A3}}
\AXC{$\stackrel{\Pi_{55}}{\vdots}$}\noLine
\UIC{$\mathrm{ContOn}(f,[a,b]) \to a \in \mathbb{R} \to \varphi_{306}$}
\AXC{$\varphi_{425}$\lbl{A1}}
\BIC{$\varphi_{303} \to n \in \mathbb{N} \to q \circ s(n) \in \mathrm{dom}(f)$}
\UIC{$\forall n\,(\varphi_{303} \to n \in \mathbb{N} \to q \circ s(n) \in \mathrm{dom}(f))$}
\AXC{$\forall n\,(\varphi_{303} \to n \in \mathbb{N} \to q \circ s(n) \in \mathrm{dom}(f)) \to \varphi_{430} \to \varphi_{486}$\lbl{A1}}
\BIC{$\varphi_{430} \to \varphi_{486}$}
\BIC{$\varphi_{486}$}
\end{ptree}
\noindent$\Pi_{57}$:
\begin{ptree}
\AXC{$\forall f\,(\varphi_{27})$\lbl{DefContOn}}
\AXC{$\forall f\,(\varphi_{27}) \to \varphi_{27}$\lbl{A2}}
\BIC{$\varphi_{27}$}
\AXC{$\varphi_{27} \to \varphi_{115}$\lbl{A2}}
\BIC{$\varphi_{115}$}
\AXC{$\varphi_{1}$\lbl{A1}}
\AXC{$\varphi_{1} \to \varphi_{115} \to \varphi_{4}$\lbl{A1}}
\BIC{$\varphi_{115} \to \varphi_{4}$}
\BIC{$\varphi_{4}$}
\end{ptree}
\noindent$\Pi_{58}$:
\begin{ptree}
\AXC{$\stackrel{\Pi_{57}}{\vdots}$}\noLine
\UIC{$\varphi_{4}$}
\AXC{$\forall f\,(\varphi_{176})$\lbl{[S8]}}
\AXC{$\forall f\,(\varphi_{176}) \to \varphi_{176}$\lbl{A2}}
\BIC{$\varphi_{176}$}
\AXC{$\varphi_{176} \to \mathrm{RFn}(f) \to \varphi_{225}$\lbl{A2}}
\BIC{$\mathrm{RFn}(f) \to \varphi_{225}$}
\AXC{$\varphi_{501}$\lbl{A1}}
\BIC{$\varphi_{4} \to (\mathrm{ContOn}(f,[a,b]) \to a \in \mathbb{R} \to \varphi_{294}) \to \varphi_{437}$}
\BIC{$(\mathrm{ContOn}(f,[a,b]) \to a \in \mathbb{R} \to \varphi_{294}) \to \varphi_{437}$}
\end{ptree}
\noindent$\Pi_{59}$:
\begin{ptree}
\AXC{$\stackrel{\Pi_{31}}{\vdots}$}\noLine
\UIC{$\mathrm{ContOn}(f,[a,b]) \to a \in \mathbb{R} \to \varphi_{317}$}
\AXC{$\varphi_{426}$\lbl{A1}}
\BIC{$\mathrm{ContOn}(f,[a,b]) \to a \in \mathbb{R} \to \varphi_{294}$}
\AXC{$\stackrel{\Pi_{58}}{\vdots}$}\noLine
\UIC{$(\mathrm{ContOn}(f,[a,b]) \to a \in \mathbb{R} \to \varphi_{294}) \to \varphi_{437}$}
\BIC{$\varphi_{437}$}
\end{ptree}
\noindent$\Pi_{60}$:
\begin{ptree}
\AXC{$\varphi_{187}$\lbl{A2}}
\AXC{$\varphi_{303} \to \forall n\,(\varphi_{303})$\lbl{A4}}
\AXC{$\stackrel{\Pi_{56}}{\vdots}$}\noLine
\UIC{$\varphi_{486}$}
\BIC{$\mathrm{ContOn}(f,[a,b]) \to a \in \mathbb{R} \to \varphi_{309}$}
\AXC{$\stackrel{\Pi_{59}}{\vdots}$}\noLine
\UIC{$\varphi_{437}$}
\BIC{$\mathrm{ContOn}(f,[a,b]) \to a \in \mathbb{R} \to \varphi_{333}$}
\AXC{$\varphi_{472}$\lbl{A1}}
\BIC{$\mathrm{ContOn}(f,[a,b]) \to a \in \mathbb{R} \to \varphi_{386}$}
\AXC{$\varphi_{496}$\lbl{A1}}
\BIC{$\varphi_{187} \to \mathrm{ContOn}(f,[a,b]) \to a \in \mathbb{R} \to \varphi_{385}$}
\BIC{$\mathrm{ContOn}(f,[a,b]) \to a \in \mathbb{R} \to \varphi_{385}$}
\end{ptree}
\noindent$\Pi_{61}$:
\begin{ptree}
\AXC{$\stackrel{\Pi_{60}}{\vdots}$}\noLine
\UIC{$\mathrm{ContOn}(f,[a,b]) \to a \in \mathbb{R} \to \varphi_{385}$}
\AXC{$(\mathrm{ContOn}(f,[a,b]) \to a \in \mathbb{R} \to \varphi_{385}) \to \varphi_{476}$\lbl{A1}}
\BIC{$\varphi_{476}$}
\end{ptree}
\noindent$\Pi_{62}$:
\begin{ptree}
\AXC{$\mathrm{ContOn}(f,[a,b]) \to a \in \mathbb{R} \to \varphi_{370}$\lbl{A1}}
\AXC{$\varphi_{470}$\lbl{A1}}
\BIC{$\mathrm{ContOn}(f,[a,b]) \to a \in \mathbb{R} \to \varphi_{342}$}
\AXC{$\stackrel{\Pi_{61}}{\vdots}$}\noLine
\UIC{$\varphi_{476}$}
\BIC{$\mathrm{ContOn}(f,[a,b]) \to a \in \mathbb{R} \to \varphi_{383}$}
\AXC{$\varphi_{507}$\lbl{A1}}
\BIC{$(\mathrm{ContOn}(f,[a,b]) \to a \in \mathbb{R} \to \varphi_{376}) \to \varphi_{409}$}
\end{ptree}
\noindent$\Pi_{63}$:
\begin{ptree}
\AXC{$\stackrel{\Pi_{36}}{\vdots}$}\noLine
\UIC{$\mathrm{ContOn}(f,[a,b]) \to a \in \mathbb{R} \to \varphi_{376}$}
\AXC{$\stackrel{\Pi_{62}}{\vdots}$}\noLine
\UIC{$(\mathrm{ContOn}(f,[a,b]) \to a \in \mathbb{R} \to \varphi_{376}) \to \varphi_{409}$}
\BIC{$\varphi_{409}$}
\end{ptree}
\noindent$\Pi_{64}$:
\begin{ptree}
\AXC{$\stackrel{\Pi_{29}}{\vdots}$}\noLine
\UIC{$\varphi_{42}$}
\AXC{$\varphi_{42} \to \varphi_{145}$\lbl{A2}}
\BIC{$\varphi_{145}$}
\AXC{$\varphi_{145} \to f \circ q \circ s(K) = f(q \circ s(K)) \to \varphi_{119}$\lbl{A2}}
\BIC{$f \circ q \circ s(K) = f(q \circ s(K)) \to \varphi_{119}$}
\AXC{$\stackrel{\Pi_{63}}{\vdots}$}\noLine
\UIC{$\varphi_{409}$}
\BIC{$\mathrm{ContOn}(f,[a,b]) \to a \in \mathbb{R} \to \varphi_{380}$}
\end{ptree}
\noindent$\Pi_{65}$:
\begin{ptree}
\AXC{$\stackrel{\Pi_{64}}{\vdots}$}\noLine
\UIC{$\mathrm{ContOn}(f,[a,b]) \to a \in \mathbb{R} \to \varphi_{380}$}
\AXC{$(\mathrm{ContOn}(f,[a,b]) \to a \in \mathbb{R} \to \varphi_{380}) \to \varphi_{498}$\lbl{A1}}
\BIC{$\varphi_{498}$}
\end{ptree}
\noindent$\Pi_{66}$:
\begin{ptree}
\AXC{$\forall X\,\forall Y\,\forall x\,(X \subseteq Y \to x \in X \to x \in Y)$\lbl{[B5]}}
\AXC{$\varphi_{90}$\lbl{A2}}
\BIC{$\forall Y\,\forall x\,([a,b] \subseteq Y \to x \in [a,b] \to x \in Y)$}
\AXC{$\forall Y\,\forall x\,([a,b] \subseteq Y \to x \in [a,b] \to x \in Y) \to \varphi_{7}$\lbl{A2}}
\BIC{$\varphi_{7}$}
\AXC{$\varphi_{7} \to \varphi_{65}$\lbl{A2}}
\BIC{$\varphi_{65}$}
\end{ptree}
\noindent$\Pi_{67}$:
\begin{ptree}
\AXC{$\stackrel{\Pi_{7}}{\vdots}$}\noLine
\UIC{$\mathrm{ContOn}(f,[a,b]) \to a \in \mathbb{R} \to \varphi_{388}$}
\AXC{$\varphi_{485}$\lbl{A1}}
\BIC{$\mathrm{ContOn}(f,[a,b]) \to a \in \mathbb{R} \to \varphi_{363}$}
\AXC{$\stackrel{\Pi_{53}}{\vdots}$}\noLine
\UIC{$\varphi_{28}$}
\AXC{$\stackrel{\Pi_{66}}{\vdots}$}\noLine
\UIC{$\varphi_{65}$}
\AXC{$\varphi_{65} \to \varphi_{28} \to \varphi_{474}$\lbl{A1}}
\BIC{$\varphi_{28} \to \varphi_{474}$}
\BIC{$\varphi_{474}$}
\BIC{$\mathrm{ContOn}(f,[a,b]) \to a \in \mathbb{R} \to \varphi_{360}$}
\end{ptree}
\noindent$\Pi_{68}$:
\begin{ptree}
\AXC{$\forall f\,\forall x\,(\mathrm{RFn}(f) \to x \in \mathrm{dom}(f) \to f(x) \in \mathbb{R})$\lbl{[S6]}}
\AXC{$\varphi_{106}$\lbl{A2}}
\BIC{$\forall x\,(\mathrm{RFn}(f) \to x \in \mathrm{dom}(f) \to f(x) \in \mathbb{R})$}
\AXC{$\forall x\,(\mathrm{RFn}(f) \to x \in \mathrm{dom}(f) \to f(x) \in \mathbb{R}) \to \varphi_{29}$\lbl{A2}}
\BIC{$\varphi_{29}$}
\end{ptree}
\noindent$\Pi_{69}$:
\begin{ptree}
\AXC{$\forall x \in \mathbb{R}\,(-\lvert x \rvert \leq x \land x \leq \lvert x \rvert)$\lbl{[A8]}}
\AXC{$\varphi_{256}$\lbl{A2}}
\BIC{$f(q(s_{K})) \in \mathbb{R} \to \varphi_{138}$}
\AXC{$(f(q(s_{K})) \in \mathbb{R} \to \varphi_{138}) \to \varphi_{489}$\lbl{A1}}
\BIC{$\varphi_{489}$}
\end{ptree}
\noindent$\Pi_{70}$:
\begin{ptree}
\AXC{$\stackrel{\Pi_{67}}{\vdots}$}\noLine
\UIC{$\mathrm{ContOn}(f,[a,b]) \to a \in \mathbb{R} \to \varphi_{360}$}
\AXC{$\stackrel{\Pi_{57}}{\vdots}$}\noLine
\UIC{$\varphi_{4}$}
\AXC{$\stackrel{\Pi_{68}}{\vdots}$}\noLine
\UIC{$\varphi_{29}$}
\AXC{$\varphi_{29} \to \varphi_{4} \to \varphi_{475}$\lbl{A1}}
\BIC{$\varphi_{4} \to \varphi_{475}$}
\BIC{$\varphi_{475}$}
\BIC{$\mathrm{ContOn}(f,[a,b]) \to a \in \mathbb{R} \to \varphi_{366}$}
\AXC{$\stackrel{\Pi_{69}}{\vdots}$}\noLine
\UIC{$\varphi_{489}$}
\BIC{$\mathrm{ContOn}(f,[a,b]) \to a \in \mathbb{R} \to \varphi_{396}$}
\AXC{$\varphi_{491}$\lbl{A1}}
\BIC{$\mathrm{ContOn}(f,[a,b]) \to a \in \mathbb{R} \to \varphi_{381}$}
\end{ptree}
\noindent$\Pi_{71}$:
\begin{ptree}
\AXC{$\forall x\,\forall y\,\forall z\,(x \leq y \to y \leq z \to x \leq z)$\lbl{[O5]}}
\AXC{$\forall x\,\forall y\,\forall z\,(x \leq y \to y \leq z \to x \leq z) \to \varphi_{16}$\lbl{A2}}
\BIC{$\varphi_{16}$}
\AXC{$\varphi_{16} \to \varphi_{143}$\lbl{A2}}
\BIC{$\varphi_{143}$}
\AXC{$\varphi_{143} \to \varphi_{139}$\lbl{A2}}
\BIC{$\varphi_{139}$}
\end{ptree}
\noindent$\Pi_{72}$:
\begin{ptree}
\AXC{$\stackrel{\Pi_{70}}{\vdots}$}\noLine
\UIC{$\mathrm{ContOn}(f,[a,b]) \to a \in \mathbb{R} \to \varphi_{381}$}
\AXC{$\stackrel{\Pi_{71}}{\vdots}$}\noLine
\UIC{$\varphi_{139}$}
\AXC{$\varphi_{139} \to (\mathrm{ContOn}(f,[a,b]) \to a \in \mathbb{R} \to \varphi_{381}) \to \varphi_{473}$\lbl{A1}}
\BIC{$(\mathrm{ContOn}(f,[a,b]) \to a \in \mathbb{R} \to \varphi_{381}) \to \varphi_{473}$}
\BIC{$\varphi_{473}$}
\end{ptree}
\noindent$\Pi_{73}$:
\begin{ptree}
\AXC{$\stackrel{\Pi_{25}}{\vdots}$}\noLine
\UIC{$\varphi_{159}$}
\AXC{$\stackrel{\Pi_{26}}{\vdots}$}\noLine
\UIC{$\varphi_{170}$}
\AXC{$\mathrm{ContOn}(f,[a,b]) \to a \in \mathbb{R} \to \varphi_{370}$\lbl{A1}}
\AXC{$\varphi_{470}$\lbl{A1}}
\BIC{$\mathrm{ContOn}(f,[a,b]) \to a \in \mathbb{R} \to \varphi_{342}$}
\AXC{$\stackrel{\Pi_{28}}{\vdots}$}\noLine
\UIC{$\varphi_{468}$}
\BIC{$\mathrm{ContOn}(f,[a,b]) \to a \in \mathbb{R} \to \varphi_{369}$}
\AXC{$\stackrel{\Pi_{65}}{\vdots}$}\noLine
\UIC{$\varphi_{498}$}
\BIC{$\varphi_{170} \to \varphi_{159} \to \mathrm{ContOn}(f,[a,b]) \to a \in \mathbb{R} \to \varphi_{362}$}
\BIC{$\varphi_{159} \to \mathrm{ContOn}(f,[a,b]) \to a \in \mathbb{R} \to \varphi_{362}$}
\BIC{$\mathrm{ContOn}(f,[a,b]) \to a \in \mathbb{R} \to \varphi_{362}$}
\AXC{$\stackrel{\Pi_{72}}{\vdots}$}\noLine
\UIC{$\varphi_{473}$}
\BIC{$\mathrm{ContOn}(f,[a,b]) \to a \in \mathbb{R} \to \varphi_{358}$}
\end{ptree}
\noindent$\Pi_{74}$:
\begin{ptree}
\AXC{$\forall x\,\forall y\,\forall z\,(x \leq y \to y \leq z \to x \leq z)$\lbl{[O5]}}
\AXC{$\forall x\,\forall y\,\forall z\,(x \leq y \to y \leq z \to x \leq z) \to \varphi_{16}$\lbl{A2}}
\BIC{$\varphi_{16}$}
\AXC{$\varphi_{173}$\lbl{A2}}
\BIC{$\forall z\,(f(q(s_{K})) \leq M \to M \leq z \to f(q(s_{K})) \leq z)$}
\AXC{$\varphi_{181}$\lbl{A2}}
\BIC{$f(q(s_{K})) \leq M \to M \leq \lvert M \rvert \to f(q(s_{K})) \leq \lvert M \rvert$}
\AXC{$\varphi_{505}$\lbl{A1}}
\BIC{$(\mathrm{ContOn}(f,[a,b]) \to a \in \mathbb{R} \to \varphi_{358}) \to \varphi_{465}$}
\end{ptree}
\noindent$\Pi_{75}$:
\begin{ptree}
\AXC{$\stackrel{\Pi_{73}}{\vdots}$}\noLine
\UIC{$\mathrm{ContOn}(f,[a,b]) \to a \in \mathbb{R} \to \varphi_{358}$}
\AXC{$\stackrel{\Pi_{74}}{\vdots}$}\noLine
\UIC{$(\mathrm{ContOn}(f,[a,b]) \to a \in \mathbb{R} \to \varphi_{358}) \to \varphi_{465}$}
\BIC{$\varphi_{465}$}
\end{ptree}
\noindent$\Pi_{76}$:
\begin{ptree}
\AXC{$\mathrm{ContOn}(f,[a,b]) \to a \in \mathbb{R} \to \varphi_{352}$\lbl{A1}}
\AXC{$\varphi_{451}$\lbl{A1}}
\BIC{$\mathrm{ContOn}(f,[a,b]) \to a \in \mathbb{R} \to \varphi_{322}$}
\AXC{$\stackrel{\Pi_{24}}{\vdots}$}\noLine
\UIC{$\varphi_{452}$}
\BIC{$\mathrm{ContOn}(f,[a,b]) \to a \in \mathbb{R} \to \varphi_{357}$}
\AXC{$\varphi_{462}$\lbl{A1}}
\BIC{$\mathrm{ContOn}(f,[a,b]) \to a \in \mathbb{R} \to \varphi_{341}$}
\AXC{$\stackrel{\Pi_{75}}{\vdots}$}\noLine
\UIC{$\varphi_{465}$}
\BIC{$\mathrm{ContOn}(f,[a,b]) \to a \in \mathbb{R} \to \varphi_{364}$}
\end{ptree}
\noindent$\Pi_{77}$:
\begin{ptree}
\AXC{$\forall x\,\forall y\,\forall z\,(x \leq y \to y < z \to x < z)$\lbl{[O4]}}
\AXC{$\forall x\,\forall y\,\forall z\,(x \leq y \to y < z \to x < z) \to \varphi_{30}$\lbl{A2}}
\BIC{$\varphi_{30}$}
\AXC{$\varphi_{30} \to \varphi_{49}$\lbl{A2}}
\BIC{$\varphi_{49}$}
\AXC{$\varphi_{49} \to \varphi_{71}$\lbl{A2}}
\BIC{$\varphi_{71}$}
\end{ptree}
\noindent$\Pi_{78}$:
\begin{ptree}
\AXC{$\stackrel{\Pi_{76}}{\vdots}$}\noLine
\UIC{$\mathrm{ContOn}(f,[a,b]) \to a \in \mathbb{R} \to \varphi_{364}$}
\AXC{$\stackrel{\Pi_{77}}{\vdots}$}\noLine
\UIC{$\varphi_{71}$}
\AXC{$\varphi_{71} \to (\mathrm{ContOn}(f,[a,b]) \to a \in \mathbb{R} \to \varphi_{364}) \to \varphi_{471}$\lbl{A1}}
\BIC{$(\mathrm{ContOn}(f,[a,b]) \to a \in \mathbb{R} \to \varphi_{364}) \to \varphi_{471}$}
\BIC{$\varphi_{471}$}
\end{ptree}
\noindent$\Pi_{79}$:
\begin{ptree}
\AXC{$z = w \to \lvert M \rvert + 1 < z \to \lvert M \rvert + 1 < w$\lbl{A6}}
\UIC{$\forall w\,(z = w \to \lvert M \rvert + 1 < z \to \lvert M \rvert + 1 < w)$}
\UIC{$\forall z\,\forall w\,(z = w \to \lvert M \rvert + 1 < z \to \lvert M \rvert + 1 < w)$}
\AXC{$\varphi_{153}$\lbl{A2}}
\BIC{$\forall w\,(\underline{K} \cdot 1 = w \to \lvert M \rvert + 1 < \underline{K} \cdot 1 \to \lvert M \rvert + 1 < w)$}
\AXC{$\varphi_{212}$\lbl{A2}}
\BIC{$\underline{K} \cdot 1 = \underline{K} \to \lvert M \rvert + 1 < \underline{K} \cdot 1 \to \lvert M \rvert + 1 < \underline{K}$}
\AXC{$\stackrel{\Pi_{17}}{\vdots}$}\noLine
\UIC{$\varphi_{399}$}
\BIC{$\mathrm{ContOn}(f,[a,b]) \to a \in \mathbb{R} \to \varphi_{356}$}
\AXC{$\stackrel{\Pi_{23}}{\vdots}$}\noLine
\UIC{$\varphi_{466}$}
\BIC{$\mathrm{ContOn}(f,[a,b]) \to a \in \mathbb{R} \to \varphi_{348}$}
\AXC{$\stackrel{\Pi_{78}}{\vdots}$}\noLine
\UIC{$\varphi_{471}$}
\BIC{$\mathrm{ContOn}(f,[a,b]) \to a \in \mathbb{R} \to \varphi_{368}$}
\end{ptree}
\noindent$\Pi_{80}$:
\begin{ptree}
\AXC{$\forall x\,\forall y\,\forall z\,(x < y \to y \leq z \to x < z)$\lbl{[O3]}}
\AXC{$\forall x\,\forall y\,\forall z\,(x < y \to y \leq z \to x < z) \to \varphi_{31}$\lbl{A2}}
\BIC{$\varphi_{31}$}
\AXC{$\varphi_{31} \to \varphi_{51}$\lbl{A2}}
\BIC{$\varphi_{51}$}
\AXC{$\varphi_{51} \to \varphi_{122}$\lbl{A2}}
\BIC{$\varphi_{122}$}
\end{ptree}
\noindent$\Pi_{81}$:
\begin{ptree}
\AXC{$\stackrel{\Pi_{79}}{\vdots}$}\noLine
\UIC{$\mathrm{ContOn}(f,[a,b]) \to a \in \mathbb{R} \to \varphi_{368}$}
\AXC{$\stackrel{\Pi_{80}}{\vdots}$}\noLine
\UIC{$\varphi_{122}$}
\AXC{$\varphi_{122} \to (\mathrm{ContOn}(f,[a,b]) \to a \in \mathbb{R} \to \varphi_{368}) \to \varphi_{477}$\lbl{A1}}
\BIC{$(\mathrm{ContOn}(f,[a,b]) \to a \in \mathbb{R} \to \varphi_{368}) \to \varphi_{477}$}
\BIC{$\varphi_{477}$}
\end{ptree}
\noindent$\Pi_{82}$:
\begin{ptree}
\AXC{$\mathrm{ContOn}(f,[a,b]) \to a \in \mathbb{R} \to \varphi_{370}$\lbl{A1}}
\AXC{$\varphi_{470}$\lbl{A1}}
\BIC{$\mathrm{ContOn}(f,[a,b]) \to a \in \mathbb{R} \to \varphi_{342}$}
\AXC{$\stackrel{\Pi_{10}}{\vdots}$}\noLine
\UIC{$\varphi_{457}$}
\BIC{$\mathrm{ContOn}(f,[a,b]) \to a \in \mathbb{R} \to \varphi_{343}$}
\AXC{$\stackrel{\Pi_{13}}{\vdots}$}\noLine
\UIC{$\varphi_{461}$}
\BIC{$\mathrm{ContOn}(f,[a,b]) \to a \in \mathbb{R} \to \varphi_{355}$}
\AXC{$\stackrel{\Pi_{81}}{\vdots}$}\noLine
\UIC{$\varphi_{477}$}
\BIC{$\mathrm{ContOn}(f,[a,b]) \to a \in \mathbb{R} \to \varphi_{372}$}
\end{ptree}
\noindent$\Pi_{83}$:
\begin{ptree}
\AXC{$\forall x\,\forall y\,\forall z\,(x < y \to y < z \to x < z)$\lbl{[O2]}}
\AXC{$\forall x\,\forall y\,\forall z\,(x < y \to y < z \to x < z) \to \varphi_{40}$\lbl{A2}}
\BIC{$\varphi_{40}$}
\AXC{$\varphi_{40} \to \varphi_{110}$\lbl{A2}}
\BIC{$\varphi_{110}$}
\AXC{$\varphi_{110} \to \varphi_{193}$\lbl{A2}}
\BIC{$\varphi_{193}$}
\end{ptree}
\noindent$\Pi_{84}$:
\begin{ptree}
\AXC{$\stackrel{\Pi_{82}}{\vdots}$}\noLine
\UIC{$\mathrm{ContOn}(f,[a,b]) \to a \in \mathbb{R} \to \varphi_{372}$}
\AXC{$\stackrel{\Pi_{83}}{\vdots}$}\noLine
\UIC{$\varphi_{193}$}
\AXC{$\varphi_{193} \to (\mathrm{ContOn}(f,[a,b]) \to a \in \mathbb{R} \to \varphi_{372}) \to \varphi_{482}$\lbl{A1}}
\BIC{$(\mathrm{ContOn}(f,[a,b]) \to a \in \mathbb{R} \to \varphi_{372}) \to \varphi_{482}$}
\BIC{$\varphi_{482}$}
\end{ptree}
\noindent$\Pi_{85}$:
\begin{ptree}
\AXC{$\stackrel{\Pi_{7}}{\vdots}$}\noLine
\UIC{$\mathrm{ContOn}(f,[a,b]) \to a \in \mathbb{R} \to \varphi_{388}$}
\AXC{$\varphi_{488}$\lbl{A1}}
\BIC{$\mathrm{ContOn}(f,[a,b]) \to a \in \mathbb{R} \to \varphi_{371}$}
\AXC{$\stackrel{\Pi_{84}}{\vdots}$}\noLine
\UIC{$\varphi_{482}$}
\BIC{$\mathrm{ContOn}(f,[a,b]) \to a \in \mathbb{R} \to \varphi_{378}$}
\AXC{$(\mathrm{ContOn}(f,[a,b]) \to a \in \mathbb{R} \to \varphi_{378}) \to \varphi_{393}$\lbl{A1}}
\BIC{$\varphi_{393}$}
\end{ptree}
\noindent$\Pi_{86}$:
\begin{ptree}
\AXC{$\forall x \in \mathbb{R}\,(\lvert x \rvert \in \mathbb{R})$\lbl{[A5]}}
\AXC{$\forall x \in \mathbb{R}\,(\lvert x \rvert \in \mathbb{R}) \to M \in \mathbb{R} \to \lvert M \rvert \in \mathbb{R}$\lbl{A2}}
\BIC{$M \in \mathbb{R} \to \lvert M \rvert \in \mathbb{R}$}
\AXC{$(M \in \mathbb{R} \to \lvert M \rvert \in \mathbb{R}) \to \varphi_{442}$\lbl{A1}}
\BIC{$\varphi_{442}$}
\end{ptree}
\noindent$\Pi_{87}$:
\begin{ptree}
\AXC{$\mathrm{ContOn}(f,[a,b]) \to a \in \mathbb{R} \to \varphi_{352}$\lbl{A1}}
\AXC{$\varphi_{451}$\lbl{A1}}
\BIC{$\mathrm{ContOn}(f,[a,b]) \to a \in \mathbb{R} \to \varphi_{322}$}
\AXC{$\stackrel{\Pi_{86}}{\vdots}$}\noLine
\UIC{$\varphi_{442}$}
\BIC{$\mathrm{ContOn}(f,[a,b]) \to a \in \mathbb{R} \to \varphi_{323}$}
\end{ptree}
\noindent$\Pi_{88}$:
\begin{ptree}
\AXC{$\forall x \in \mathbb{R}\,(x < x + 1)$\lbl{[N13]}}
\AXC{$\forall x \in \mathbb{R}\,(x < x + 1) \to \lvert M \rvert \in \mathbb{R} \to \lvert M \rvert < \lvert M \rvert + 1$\lbl{A2}}
\BIC{$\lvert M \rvert \in \mathbb{R} \to \lvert M \rvert < \lvert M \rvert + 1$}
\end{ptree}
\noindent$\Pi_{89}$:
\begin{ptree}
\AXC{$\stackrel{\Pi_{87}}{\vdots}$}\noLine
\UIC{$\mathrm{ContOn}(f,[a,b]) \to a \in \mathbb{R} \to \varphi_{323}$}
\AXC{$\stackrel{\Pi_{88}}{\vdots}$}\noLine
\UIC{$\lvert M \rvert \in \mathbb{R} \to \lvert M \rvert < \lvert M \rvert + 1$}
\AXC{$(\lvert M \rvert \in \mathbb{R} \to \lvert M \rvert < \lvert M \rvert + 1) \to \varphi_{444}$\lbl{A1}}
\BIC{$\varphi_{444}$}
\BIC{$\mathrm{ContOn}(f,[a,b]) \to a \in \mathbb{R} \to \varphi_{330}$}
\end{ptree}
\noindent$\Pi_{90}$:
\begin{ptree}
\AXC{$\forall x \in \mathbb{R}\,(0 \leq \lvert x \rvert)$\lbl{[A6]}}
\AXC{$\forall x \in \mathbb{R}\,(0 \leq \lvert x \rvert) \to M \in \mathbb{R} \to 0 \leq \lvert M \rvert$\lbl{A2}}
\BIC{$M \in \mathbb{R} \to 0 \leq \lvert M \rvert$}
\AXC{$(M \in \mathbb{R} \to 0 \leq \lvert M \rvert) \to \varphi_{440}$\lbl{A1}}
\BIC{$\varphi_{440}$}
\end{ptree}
\noindent$\Pi_{91}$:
\begin{ptree}
\AXC{$\forall x\,\forall y\,\forall z\,(x \leq y \to y < z \to x < z)$\lbl{[O4]}}
\AXC{$\varphi_{13}$\lbl{A2}}
\BIC{$\forall y\,\forall z\,(0 \leq y \to y < z \to 0 < z)$}
\AXC{$\varphi_{17}$\lbl{A2}}
\BIC{$\forall z\,(0 \leq \lvert M \rvert \to \lvert M \rvert < z \to 0 < z)$}
\AXC{$\varphi_{81}$\lbl{A2}}
\BIC{$0 \leq \lvert M \rvert \to \lvert M \rvert < \lvert M \rvert + 1 \to 0 < \lvert M \rvert + 1$}
\AXC{$\varphi_{500}$\lbl{A1}}
\BIC{$(\mathrm{ContOn}(f,[a,b]) \to a \in \mathbb{R} \to \varphi_{321}) \to \varphi_{445}$}
\end{ptree}
\noindent$\Pi_{92}$:
\begin{ptree}
\AXC{$\mathrm{ContOn}(f,[a,b]) \to a \in \mathbb{R} \to \varphi_{352}$\lbl{A1}}
\AXC{$\varphi_{451}$\lbl{A1}}
\BIC{$\mathrm{ContOn}(f,[a,b]) \to a \in \mathbb{R} \to \varphi_{322}$}
\AXC{$\stackrel{\Pi_{90}}{\vdots}$}\noLine
\UIC{$\varphi_{440}$}
\BIC{$\mathrm{ContOn}(f,[a,b]) \to a \in \mathbb{R} \to \varphi_{321}$}
\AXC{$\stackrel{\Pi_{91}}{\vdots}$}\noLine
\UIC{$(\mathrm{ContOn}(f,[a,b]) \to a \in \mathbb{R} \to \varphi_{321}) \to \varphi_{445}$}
\BIC{$\varphi_{445}$}
\end{ptree}
\noindent$\Pi_{93}$:
\begin{ptree}
\AXC{$\forall x\,\forall y\,(x \in \mathbb{R} \to y \in \mathbb{R} \to x + y \in \mathbb{R})$\lbl{(R1)}}
\AXC{$\varphi_{127}$\lbl{A2}}
\BIC{$\forall y\,(\lvert M \rvert \in \mathbb{R} \to y \in \mathbb{R} \to \lvert M \rvert + y \in \mathbb{R})$}
\AXC{$\varphi_{132}$\lbl{A2}}
\BIC{$\lvert M \rvert \in \mathbb{R} \to 1 \in \mathbb{R} \to \lvert M \rvert + 1 \in \mathbb{R}$}
\AXC{$(\lvert M \rvert \in \mathbb{R} \to 1 \in \mathbb{R} \to \lvert M \rvert + 1 \in \mathbb{R}) \to \varphi_{446}$\lbl{A1}}
\BIC{$\varphi_{446}$}
\end{ptree}
\noindent$\Pi_{94}$:
\begin{ptree}
\AXC{$1 \in \mathbb{R}$\lbl{(R1)}}
\AXC{$\varphi_{66}$\lbl{命題3.2}}
\AXC{$\varphi_{66} \to \varphi_{55}$\lbl{A2}}
\BIC{$\varphi_{55}$}
\AXC{$\varphi_{55} \to 1 \in \mathbb{R} \to \varphi_{103}$\lbl{A2}}
\BIC{$1 \in \mathbb{R} \to \varphi_{103}$}
\AXC{$(1 \in \mathbb{R} \to \varphi_{103}) \to 1 \in \mathbb{R} \to \varphi_{508}$\lbl{A1}}
\BIC{$1 \in \mathbb{R} \to \varphi_{508}$}
\BIC{$\varphi_{508}$}
\end{ptree}
\noindent$\Pi_{95}$:
\begin{ptree}
\AXC{$1 \in \mathbb{R}$\lbl{(R1)}}
\AXC{$\stackrel{\Pi_{87}}{\vdots}$}\noLine
\UIC{$\mathrm{ContOn}(f,[a,b]) \to a \in \mathbb{R} \to \varphi_{323}$}
\AXC{$\stackrel{\Pi_{93}}{\vdots}$}\noLine
\UIC{$\varphi_{446}$}
\BIC{$1 \in \mathbb{R} \to \mathrm{ContOn}(f,[a,b]) \to a \in \mathbb{R} \to \varphi_{328}$}
\BIC{$\mathrm{ContOn}(f,[a,b]) \to a \in \mathbb{R} \to \varphi_{328}$}
\AXC{$\stackrel{\Pi_{94}}{\vdots}$}\noLine
\UIC{$\varphi_{508}$}
\BIC{$(\mathrm{ContOn}(f,[a,b]) \to a \in \mathbb{R} \to \varphi_{319}) \to \varphi_{449}$}
\end{ptree}
\noindent$\Pi_{96}$:
\begin{ptree}
\AXC{$0 < 1$\lbl{[F28]}}
\AXC{$0 < 1 \to \mathrm{ContOn}(f,[a,b]) \to a \in \mathbb{R} \to \varphi_{319}$\lbl{A1}}
\BIC{$\mathrm{ContOn}(f,[a,b]) \to a \in \mathbb{R} \to \varphi_{319}$}
\AXC{$\stackrel{\Pi_{95}}{\vdots}$}\noLine
\UIC{$(\mathrm{ContOn}(f,[a,b]) \to a \in \mathbb{R} \to \varphi_{319}) \to \varphi_{449}$}
\BIC{$\varphi_{449}$}
\end{ptree}
\noindent$\Pi_{97}$:
\begin{ptree}
\AXC{$\stackrel{\Pi_{89}}{\vdots}$}\noLine
\UIC{$\mathrm{ContOn}(f,[a,b]) \to a \in \mathbb{R} \to \varphi_{330}$}
\AXC{$\stackrel{\Pi_{92}}{\vdots}$}\noLine
\UIC{$\varphi_{445}$}
\BIC{$\mathrm{ContOn}(f,[a,b]) \to a \in \mathbb{R} \to \varphi_{325}$}
\AXC{$\stackrel{\Pi_{96}}{\vdots}$}\noLine
\UIC{$\varphi_{449}$}
\BIC{$\mathrm{ContOn}(f,[a,b]) \to a \in \mathbb{R} \to \varphi_{337}$}
\AXC{$(\mathrm{ContOn}(f,[a,b]) \to a \in \mathbb{R} \to \varphi_{337}) \to \varphi_{514}$\lbl{A1}}
\BIC{$\varphi_{514}$}
\end{ptree}
\noindent$\Pi_{98}$:
\begin{ptree}
\AXC{$\forall x\,(\neg x < x)$\lbl{[O1]}}
\AXC{$\forall x\,(\neg x < x) \to \neg f(q(s_{K})) < f(q(s_{K}))$\lbl{A2}}
\BIC{$\neg f(q(s_{K})) < f(q(s_{K}))$}
\AXC{$\stackrel{\Pi_{85}}{\vdots}$}\noLine
\UIC{$\varphi_{393}$}
\BIC{$\mathrm{ContOn}(f,[a,b]) \to a \in \mathbb{R} \to \varphi_{340}$}
\AXC{$\varphi_{456}$\lbl{A1}}
\BIC{$\varphi_{336} \to \neg (K \in \mathbb{N} \land \lvert M \rvert + 1 < \underline{K} \cdot 1)$}
\UIC{$\forall K\,(\varphi_{336} \to \neg (K \in \mathbb{N} \land \lvert M \rvert + 1 < \underline{K} \cdot 1))$}
\AXC{$\stackrel{\Pi_{97}}{\vdots}$}\noLine
\UIC{$\varphi_{514}$}
\BIC{$\varphi_{458} \to (\varphi_{336} \to \forall K\,(\varphi_{336})) \to \forall n\,(\varphi_{74}) \to \varphi_{420}$}
\end{ptree}
\noindent$\Pi_{99}$:
\begin{ptree}
\AXC{$\varphi_{95}$\lbl{A4}}
\AXC{$\forall n\,(\varphi_{74}) \to \varphi_{97}$\lbl{A3}}
\AXC{$\varphi_{74}$\lbl{A2}}
\UIC{$\forall n\,(\varphi_{74})$}
\AXC{$\varphi_{336} \to \forall K\,(\varphi_{336})$\lbl{A4}}
\AXC{$\varphi_{458}$\lbl{A3}}
\AXC{$\stackrel{\Pi_{98}}{\vdots}$}\noLine
\UIC{$\varphi_{458} \to (\varphi_{336} \to \forall K\,(\varphi_{336})) \to \forall n\,(\varphi_{74}) \to \varphi_{420}$}
\BIC{$(\varphi_{336} \to \forall K\,(\varphi_{336})) \to \forall n\,(\varphi_{74}) \to \varphi_{420}$}
\BIC{$\forall n\,(\varphi_{74}) \to \varphi_{420}$}
\BIC{$\varphi_{420}$}
\BIC{$\varphi_{95} \to \mathrm{ContOn}(f,[a,b]) \to a \in \mathbb{R} \to \varphi_{320}$}
\BIC{$\mathrm{ContOn}(f,[a,b]) \to a \in \mathbb{R} \to \varphi_{320}$}
\end{ptree}
\noindent$\Pi_{100}$:
\begin{ptree}
\AXC{$\stackrel{\Pi_{99}}{\vdots}$}\noLine
\UIC{$\mathrm{ContOn}(f,[a,b]) \to a \in \mathbb{R} \to \varphi_{320}$}
\AXC{$(\mathrm{ContOn}(f,[a,b]) \to a \in \mathbb{R} \to \varphi_{320}) \to \varphi_{338}$\lbl{A1}}
\BIC{$\varphi_{338}$}
\UIC{$\forall M\,(\varphi_{338})$}
\end{ptree}
\noindent$\Pi_{101}$:
\begin{ptree}
\AXC{$\forall f\,(\varphi_{27})$\lbl{DefContOn}}
\AXC{$\forall f\,(\varphi_{27}) \to \varphi_{27}$\lbl{A2}}
\BIC{$\varphi_{27}$}
\AXC{$\varphi_{27} \to \varphi_{115}$\lbl{A2}}
\BIC{$\varphi_{115}$}
\AXC{$\varphi_{1}$\lbl{A1}}
\AXC{$\varphi_{1} \to \varphi_{115} \to \mathrm{ContOn}(f,[a,b]) \to a \in \mathbb{R} \to \varphi_{5}$\lbl{A1}}
\BIC{$\varphi_{115} \to \mathrm{ContOn}(f,[a,b]) \to a \in \mathbb{R} \to \varphi_{5}$}
\BIC{$\mathrm{ContOn}(f,[a,b]) \to a \in \mathbb{R} \to \varphi_{5}$}
\end{ptree}
\noindent$\Pi_{102}$:
\begin{ptree}
\AXC{$\varphi_{69}$\lbl{A2}}
\AXC{$\stackrel{\Pi_{101}}{\vdots}$}\noLine
\UIC{$\mathrm{ContOn}(f,[a,b]) \to a \in \mathbb{R} \to \varphi_{5}$}
\AXC{$\varphi_{300}$\lbl{A1}}
\BIC{$\varphi_{69} \to \mathrm{ContOn}(f,[a,b]) \to \varphi_{21}$}
\BIC{$\mathrm{ContOn}(f,[a,b]) \to \varphi_{21}$}
\AXC{$(\mathrm{ContOn}(f,[a,b]) \to \varphi_{21}) \to \varphi_{410}$\lbl{A1}}
\BIC{$\varphi_{410}$}
\end{ptree}
\noindent$\Pi_{103}$:
\begin{ptree}
\AXC{$\varphi_{133}$\lbl{[G1]}}
\AXC{$\varphi_{133} \to \varphi_{146}$\lbl{A2}}
\BIC{$\varphi_{146}$}
\AXC{$\varphi_{146} \to \varphi_{157}$\lbl{A2}}
\BIC{$\varphi_{157}$}
\AXC{$\varphi_{157} \to \varphi_{164}$\lbl{A2}}
\BIC{$\varphi_{164}$}
\AXC{$\varphi_{164} \to \varphi_{189}$\lbl{A2}}
\BIC{$\varphi_{189}$}
\UIC{$\forall v_{807}\,(\varphi_{189})$}
\UIC{$\forall v_{806}\,\forall v_{807}\,(\varphi_{189})$}
\UIC{$\forall v_{805}\,\forall v_{806}\,\forall v_{807}\,(\varphi_{189})$}
\UIC{$\forall v_{804}\,\forall v_{805}\,\forall v_{806}\,\forall v_{807}\,(\varphi_{189})$}
\AXC{$\forall v_{804}\,\forall v_{805}\,\forall v_{806}\,\forall v_{807}\,(\varphi_{189}) \to \varphi_{207}$\lbl{A2}}
\BIC{$\varphi_{207}$}
\end{ptree}
\noindent$\Pi_{104}$:
\begin{ptree}
\AXC{$\stackrel{\Pi_{103}}{\vdots}$}\noLine
\UIC{$\varphi_{207}$}
\AXC{$\varphi_{207} \to \varphi_{216}$\lbl{A2}}
\BIC{$\varphi_{216}$}
\AXC{$\varphi_{216} \to \varphi_{182}$\lbl{A2}}
\BIC{$\varphi_{182}$}
\AXC{$\varphi_{182} \to \mathrm{ContIn}(f,[a,b],c) \to \mathrm{Seq}(q \circ s) \to \varphi_{126}$\lbl{A2}}
\BIC{$\mathrm{ContIn}(f,[a,b],c) \to \mathrm{Seq}(q \circ s) \to \varphi_{126}$}
\AXC{$\varphi_{511}$\lbl{A1}}
\BIC{$(\mathrm{ContOn}(f,[a,b]) \to a \in \mathbb{R} \to \varphi_{299}) \to \varphi_{502}$}
\end{ptree}
\noindent$\Pi_{105}$:
\begin{ptree}
\AXC{$\mathrm{ContOn}(f,[a,b]) \to a \in \mathbb{R} \to \varphi_{313}$\lbl{A1}}
\AXC{$\varphi_{415}$\lbl{A1}}
\BIC{$\mathrm{ContOn}(f,[a,b]) \to a \in \mathbb{R} \to \varphi_{295}$}
\AXC{$\stackrel{\Pi_{102}}{\vdots}$}\noLine
\UIC{$\varphi_{410}$}
\BIC{$\mathrm{ContOn}(f,[a,b]) \to a \in \mathbb{R} \to \varphi_{299}$}
\AXC{$\stackrel{\Pi_{104}}{\vdots}$}\noLine
\UIC{$(\mathrm{ContOn}(f,[a,b]) \to a \in \mathbb{R} \to \varphi_{299}) \to \varphi_{502}$}
\BIC{$\varphi_{502}$}
\end{ptree}
\noindent$\Pi_{106}$:
\begin{ptree}
\AXC{$\stackrel{\Pi_{52}}{\vdots}$}\noLine
\UIC{$\mathrm{ContOn}(f,[a,b]) \to a \in \mathbb{R} \to \varphi_{312}$}
\AXC{$\stackrel{\Pi_{31}}{\vdots}$}\noLine
\UIC{$\mathrm{ContOn}(f,[a,b]) \to a \in \mathbb{R} \to \varphi_{317}$}
\AXC{$\varphi_{426}$\lbl{A1}}
\BIC{$\mathrm{ContOn}(f,[a,b]) \to a \in \mathbb{R} \to \varphi_{294}$}
\AXC{$\stackrel{\Pi_{105}}{\vdots}$}\noLine
\UIC{$\varphi_{502}$}
\BIC{$(\mathrm{ContOn}(f,[a,b]) \to a \in \mathbb{R} \to \varphi_{312}) \to \varphi_{412}$}
\BIC{$\varphi_{412}$}
\end{ptree}
\noindent$\Pi_{107}$:
\begin{ptree}
\AXC{$\mathrm{ContOn}(f,[a,b]) \to a \in \mathbb{R} \to \varphi_{313}$\lbl{A1}}
\AXC{$\varphi_{417}$\lbl{A1}}
\BIC{$\mathrm{ContOn}(f,[a,b]) \to a \in \mathbb{R} \to \varphi_{296}$}
\AXC{$\stackrel{\Pi_{106}}{\vdots}$}\noLine
\UIC{$\varphi_{412}$}
\BIC{$\mathrm{ContOn}(f,[a,b]) \to a \in \mathbb{R} \to \varphi_{305}$}
\end{ptree}
\noindent$\Pi_{108}$:
\begin{ptree}
\AXC{$\forall a\,\forall \alpha\,(\mathrm{Conv}(a,\alpha) \to \mathrm{SeqBdd}(a))$\lbl{定理2.11}}
\AXC{$\forall a\,\forall \alpha\,(\mathrm{Conv}(a,\alpha) \to \mathrm{SeqBdd}(a)) \to \varphi_{18}$\lbl{A2}}
\BIC{$\varphi_{18}$}
\AXC{$\varphi_{191}$\lbl{A2}}
\BIC{$\mathrm{Conv}(f \circ q \circ s,f(c)) \to \mathrm{SeqBdd}(f \circ q \circ s)$}
\end{ptree}
\noindent$\Pi_{109}$:
\begin{ptree}
\AXC{$\stackrel{\Pi_{107}}{\vdots}$}\noLine
\UIC{$\mathrm{ContOn}(f,[a,b]) \to a \in \mathbb{R} \to \varphi_{305}$}
\AXC{$\stackrel{\Pi_{108}}{\vdots}$}\noLine
\UIC{$\mathrm{Conv}(f \circ q \circ s,f(c)) \to \mathrm{SeqBdd}(f \circ q \circ s)$}
\AXC{$\varphi_{436}$\lbl{A1}}
\BIC{$\varphi_{416}$}
\BIC{$\mathrm{ContOn}(f,[a,b]) \to a \in \mathbb{R} \to \varphi_{302}$}
\AXC{$\varphi_{455}$\lbl{A1}}
\BIC{$\varphi_{168} \to \mathrm{ContOn}(f,[a,b]) \to a \in \mathbb{R} \to \varphi_{318}$}
\end{ptree}
\noindent$\Pi_{110}$:
\begin{ptree}
\AXC{$\stackrel{\Pi_{100}}{\vdots}$}\noLine
\UIC{$\forall M\,(\varphi_{338})$}
\AXC{$\varphi_{35}$\lbl{DefSeqBdd}}
\AXC{$\varphi_{35} \to \varphi_{168}$\lbl{A2}}
\BIC{$\varphi_{168}$}
\AXC{$\stackrel{\Pi_{109}}{\vdots}$}\noLine
\UIC{$\varphi_{168} \to \mathrm{ContOn}(f,[a,b]) \to a \in \mathbb{R} \to \varphi_{318}$}
\BIC{$\mathrm{ContOn}(f,[a,b]) \to a \in \mathbb{R} \to \varphi_{318}$}
\AXC{$\varphi_{513}$\lbl{A1}}
\BIC{$\forall M\,(\varphi_{338}) \to (\forall M\,(\varphi_{338}) \to \forall M\,(\varphi_{303} \land \neg 0 < 0) \to \varphi_{23}) \to \varphi_{483}$}
\BIC{$(\forall M\,(\varphi_{338}) \to \forall M\,(\varphi_{303} \land \neg 0 < 0) \to \varphi_{23}) \to \varphi_{483}$}
\end{ptree}
\noindent$\Pi_{111}$:
\begin{ptree}
\AXC{$\varphi_{303} \land \neg 0 < 0 \to \forall M\,(\varphi_{303} \land \neg 0 < 0)$\lbl{A4}}
\AXC{$\forall M\,(\varphi_{338}) \to \forall M\,(\varphi_{303} \land \neg 0 < 0) \to \varphi_{23}$\lbl{A3}}
\AXC{$\stackrel{\Pi_{110}}{\vdots}$}\noLine
\UIC{$(\forall M\,(\varphi_{338}) \to \forall M\,(\varphi_{303} \land \neg 0 < 0) \to \varphi_{23}) \to \varphi_{483}$}
\BIC{$\varphi_{483}$}
\BIC{$\mathrm{ContOn}(f,[a,b]) \to a \in \mathbb{R} \to \varphi_{291}$}
\end{ptree}
\noindent$\Pi_{112}$:
\begin{ptree}
\AXC{$\stackrel{\Pi_{111}}{\vdots}$}\noLine
\UIC{$\mathrm{ContOn}(f,[a,b]) \to a \in \mathbb{R} \to \varphi_{291}$}
\AXC{$\varphi_{407}$\lbl{A1}}
\BIC{$\varphi_{113} \land \varphi_{47} \land \varphi_{3} \land \neg 0 < 0 \to \varphi_{2}$}
\UIC{$\forall c\,(\varphi_{113} \land \varphi_{47} \land \varphi_{3} \land \neg 0 < 0 \to \varphi_{2})$}
\AXC{$\mathrm{ContOn}(f,[a,b]) \to a \in \mathbb{R} \to \varphi_{288}$\lbl{A1}}
\AXC{$\varphi_{512}$\lbl{A1}}
\BIC{$\forall c\,(\varphi_{113} \land \varphi_{47} \land \varphi_{3} \land \neg 0 < 0 \to \varphi_{2}) \to \varphi_{408} \to \varphi_{443}$}
\BIC{$\varphi_{408} \to \varphi_{443}$}
\end{ptree}
\noindent$\Pi_{113}$:
\begin{ptree}
\AXC{$\varphi_{113} \land \varphi_{47} \land \varphi_{3} \land \neg 0 < 0 \to \forall c\,(\varphi_{113} \land \varphi_{47} \land \varphi_{3} \land \neg 0 < 0)$\lbl{A4}}
\AXC{$\varphi_{408}$\lbl{A3}}
\AXC{$\stackrel{\Pi_{112}}{\vdots}$}\noLine
\UIC{$\varphi_{408} \to \varphi_{443}$}
\BIC{$\varphi_{443}$}
\BIC{$\mathrm{ContOn}(f,[a,b]) \to a \in \mathbb{R} \to \varphi_{271}$}
\AXC{$\varphi_{398}$\lbl{A1}}
\BIC{$\varphi_{113} \land \varphi_{47} \land \neg 0 < 0 \to \neg \varphi_{3}$}
\UIC{$\forall s\,(\varphi_{113} \land \varphi_{47} \land \neg 0 < 0 \to \neg \varphi_{3})$}
\end{ptree}
\noindent$\Pi_{114}$:
\begin{ptree}
\AXC{$\varphi_{201}$\lbl{[G7]}}
\AXC{$\varphi_{201} \to \forall a\,\forall b\,(a \in \mathbb{R} \to \varphi_{152})$\lbl{A2}}
\BIC{$\forall a\,\forall b\,(a \in \mathbb{R} \to \varphi_{152})$}
\AXC{$\forall a\,\forall b\,(a \in \mathbb{R} \to \varphi_{152}) \to \forall b\,(a \in \mathbb{R} \to \varphi_{152})$\lbl{A2}}
\BIC{$\forall b\,(a \in \mathbb{R} \to \varphi_{152})$}
\end{ptree}
\noindent$\Pi_{115}$:
\begin{ptree}
\AXC{$\mathrm{ContOn}(f,[a,b]) \to a \in \mathbb{R} \to a \in \mathbb{R}$\lbl{A1}}
\AXC{$\stackrel{\Pi_{114}}{\vdots}$}\noLine
\UIC{$\forall b\,(a \in \mathbb{R} \to \varphi_{152})$}
\AXC{$\forall b\,(a \in \mathbb{R} \to \varphi_{152}) \to a \in \mathbb{R} \to \varphi_{152}$\lbl{A2}}
\BIC{$a \in \mathbb{R} \to \varphi_{152}$}
\AXC{$\varphi_{484}$\lbl{A1}}
\BIC{$(\mathrm{ContOn}(f,[a,b]) \to a \in \mathbb{R} \to a \in \mathbb{R}) \to \varphi_{435}$}
\BIC{$\varphi_{435}$}
\end{ptree}
\noindent$\Pi_{116}$:
\begin{ptree}
\AXC{$\mathrm{ContOn}(f,[a,b]) \to a \in \mathbb{R} \to b \in \mathbb{R} \to b \in \mathbb{R}$\lbl{A1}}
\AXC{$\stackrel{\Pi_{115}}{\vdots}$}\noLine
\UIC{$\varphi_{435}$}
\BIC{$\varphi_{421}$}
\end{ptree}
\noindent$\Pi_{117}$:
\begin{ptree}
\AXC{$\stackrel{\Pi_{43}}{\vdots}$}\noLine
\UIC{$\mathrm{ContOn}(f,[a,b]) \to a \in \mathbb{R} \to b \in \mathbb{R} \to \varphi_{233}$}
\AXC{$\mathrm{ContOn}(f,[a,b]) \to a \in \mathbb{R} \to \varphi_{284}$\lbl{A1}}
\AXC{$\varphi_{391}$\lbl{A1}}
\BIC{$\mathrm{ContOn}(f,[a,b]) \to a \in \mathbb{R} \to b \in \mathbb{R} \to \varphi_{195}$}
\AXC{$\stackrel{\Pi_{116}}{\vdots}$}\noLine
\UIC{$\varphi_{421}$}
\BIC{$\varphi_{394}$}
\BIC{$\mathrm{ContOn}(f,[a,b]) \to a \in \mathbb{R} \to \varphi_{266}$}
\end{ptree}
\noindent$\Pi_{118}$:
\begin{ptree}
\AXC{$\varphi_{400}$\lbl{A3}}
\AXC{$\stackrel{\Pi_{113}}{\vdots}$}\noLine
\UIC{$\forall s\,(\varphi_{113} \land \varphi_{47} \land \neg 0 < 0 \to \neg \varphi_{3})$}
\AXC{$\stackrel{\Pi_{117}}{\vdots}$}\noLine
\UIC{$\mathrm{ContOn}(f,[a,b]) \to a \in \mathbb{R} \to \varphi_{266}$}
\AXC{$\varphi_{509}$\lbl{A1}}
\BIC{$\forall s\,(\varphi_{113} \land \varphi_{47} \land \neg 0 < 0 \to \neg \varphi_{3}) \to \varphi_{400} \to \varphi_{404}$}
\BIC{$\varphi_{400} \to \varphi_{404}$}
\BIC{$\varphi_{404}$}
\end{ptree}
\noindent$\Pi_{119}$:
\begin{ptree}
\AXC{$\forall X\,\forall Y\,\forall x\,(X \subseteq Y \to x \in X \to x \in Y)$\lbl{[B5]}}
\AXC{$\varphi_{90}$\lbl{A2}}
\BIC{$\forall Y\,\forall x\,([a,b] \subseteq Y \to x \in [a,b] \to x \in Y)$}
\AXC{$\forall Y\,\forall x\,([a,b] \subseteq Y \to x \in [a,b] \to x \in Y) \to \varphi_{7}$\lbl{A2}}
\BIC{$\varphi_{7}$}
\AXC{$\varphi_{147}$\lbl{A2}}
\BIC{$[a,b] \subseteq \mathrm{dom}(f) \to x \in [a,b] \to x \in \mathrm{dom}(f)$}
\end{ptree}
\noindent$\Pi_{120}$:
\begin{ptree}
\AXC{$\stackrel{\Pi_{53}}{\vdots}$}\noLine
\UIC{$\varphi_{28}$}
\AXC{$\stackrel{\Pi_{119}}{\vdots}$}\noLine
\UIC{$[a,b] \subseteq \mathrm{dom}(f) \to x \in [a,b] \to x \in \mathrm{dom}(f)$}
\AXC{$\varphi_{395}$\lbl{A1}}
\BIC{$\varphi_{28} \to \varphi_{326}$}
\BIC{$\varphi_{326}$}
\end{ptree}
\noindent$\Pi_{121}$:
\begin{ptree}
\AXC{$\forall f\,\forall x\,(\mathrm{RFn}(f) \to x \in \mathrm{dom}(f) \to f(x) \in \mathbb{R})$\lbl{[S6]}}
\AXC{$\varphi_{106}$\lbl{A2}}
\BIC{$\forall x\,(\mathrm{RFn}(f) \to x \in \mathrm{dom}(f) \to f(x) \in \mathbb{R})$}
\AXC{$\varphi_{82}$\lbl{A2}}
\BIC{$\mathrm{RFn}(f) \to x \in \mathrm{dom}(f) \to f(x) \in \mathbb{R}$}
\end{ptree}
\noindent$\Pi_{122}$:
\begin{ptree}
\AXC{$\stackrel{\Pi_{57}}{\vdots}$}\noLine
\UIC{$\varphi_{4}$}
\AXC{$\stackrel{\Pi_{121}}{\vdots}$}\noLine
\UIC{$\mathrm{RFn}(f) \to x \in \mathrm{dom}(f) \to f(x) \in \mathbb{R}$}
\AXC{$(\mathrm{RFn}(f) \to x \in \mathrm{dom}(f) \to f(x) \in \mathbb{R}) \to \varphi_{4} \to \varphi_{327}$\lbl{A1}}
\BIC{$\varphi_{4} \to \varphi_{327}$}
\BIC{$\varphi_{327}$}
\end{ptree}
\noindent$\Pi_{123}$:
\begin{ptree}
\AXC{$\mathrm{ContOn}(f,[a,b]) \to a \in \mathbb{R} \to \varphi_{223}$\lbl{A1}}
\AXC{$\varphi_{335}$\lbl{A1}}
\BIC{$\mathrm{ContOn}(f,[a,b]) \to \varphi_{210}$}
\AXC{$\stackrel{\Pi_{120}}{\vdots}$}\noLine
\UIC{$\varphi_{326}$}
\BIC{$\mathrm{ContOn}(f,[a,b]) \to \varphi_{205}$}
\AXC{$\stackrel{\Pi_{122}}{\vdots}$}\noLine
\UIC{$\varphi_{327}$}
\BIC{$\mathrm{ContOn}(f,[a,b]) \to \varphi_{215}$}
\end{ptree}
\noindent$\Pi_{124}$:
\begin{ptree}
\AXC{$\forall n \in \mathbb{N}\,(\underline{n} \in \mathbb{R})$\lbl{[N4]}}
\AXC{$\forall n \in \mathbb{N}\,(\underline{n} \in \mathbb{R}) \to k \in \mathbb{N} \to \underline{k} \in \mathbb{R}$\lbl{A2}}
\BIC{$k \in \mathbb{N} \to \underline{k} \in \mathbb{R}$}
\AXC{$(k \in \mathbb{N} \to \underline{k} \in \mathbb{R}) \to \varphi_{290}$\lbl{A1}}
\BIC{$\varphi_{290}$}
\end{ptree}
\noindent$\Pi_{125}$:
\begin{ptree}
\AXC{$\mathrm{ContOn}(f,[a,b]) \to a \in \mathbb{R} \to b \in \mathbb{R} \to \varphi_{41}$\lbl{A1}}
\AXC{$\stackrel{\Pi_{124}}{\vdots}$}\noLine
\UIC{$\varphi_{290}$}
\BIC{$\mathrm{ContOn}(f,[a,b]) \to a \in \mathbb{R} \to b \in \mathbb{R} \to \varphi_{56}$}
\end{ptree}
\noindent$\Pi_{126}$:
\begin{ptree}
\AXC{$\forall x\,\forall y\,(x \in \mathbb{R} \to y \in \mathbb{R} \to \neg y \leq x \to x < y)$\lbl{[A2]}}
\AXC{$\varphi_{135}$\lbl{A2}}
\BIC{$\forall y\,(\underline{k} \in \mathbb{R} \to y \in \mathbb{R} \to \neg y \leq \underline{k} \to \underline{k} < y)$}
\AXC{$\varphi_{208}$\lbl{A2}}
\BIC{$\underline{k} \in \mathbb{R} \to f(x) \in \mathbb{R} \to \neg f(x) \leq \underline{k} \to \underline{k} < f(x)$}
\AXC{$\varphi_{441}$\lbl{A1}}
\BIC{$\varphi_{423}$}
\end{ptree}
\noindent$\Pi_{127}$:
\begin{ptree}
\AXC{$\stackrel{\Pi_{123}}{\vdots}$}\noLine
\UIC{$\mathrm{ContOn}(f,[a,b]) \to \varphi_{215}$}
\AXC{$\stackrel{\Pi_{125}}{\vdots}$}\noLine
\UIC{$\mathrm{ContOn}(f,[a,b]) \to a \in \mathbb{R} \to b \in \mathbb{R} \to \varphi_{56}$}
\AXC{$\stackrel{\Pi_{126}}{\vdots}$}\noLine
\UIC{$\varphi_{423}$}
\BIC{$(\mathrm{ContOn}(f,[a,b]) \to \varphi_{215}) \to \varphi_{332}$}
\BIC{$\varphi_{332}$}
\end{ptree}
\noindent$\Pi_{128}$:
\begin{ptree}
\AXC{$\mathrm{ContOn}(f,[a,b]) \to a \in \mathbb{R} \to \varphi_{223}$\lbl{A1}}
\AXC{$\varphi_{335}$\lbl{A1}}
\BIC{$\mathrm{ContOn}(f,[a,b]) \to \varphi_{210}$}
\AXC{$(\mathrm{ContOn}(f,[a,b]) \to \varphi_{210}) \to \varphi_{272}$\lbl{A1}}
\BIC{$\varphi_{272}$}
\end{ptree}
\noindent$\Pi_{129}$:
\begin{ptree}
\AXC{$\forall a\,\forall b\,\forall x\,(x \in [a,b] \leftrightarrow x \in \mathbb{R} \land a \leq x \land x \leq b)$\lbl{Def{$[a,b]$}}}
\AXC{$\varphi_{118}$\lbl{A2}}
\BIC{$\forall b\,\forall x\,(x \in [a,b] \leftrightarrow x \in \mathbb{R} \land a \leq x \land x \leq b)$}
\AXC{$\varphi_{96}$\lbl{A2}}
\BIC{$\forall x\,(x \in [a,b] \leftrightarrow x \in \mathbb{R} \land a \leq x \land x \leq b)$}
\AXC{$\varphi_{76}$\lbl{A2}}
\BIC{$x \in [a,b] \leftrightarrow x \in \mathbb{R} \land a \leq x \land x \leq b$}
\AXC{$\stackrel{\Pi_{128}}{\vdots}$}\noLine
\UIC{$\varphi_{272}$}
\BIC{$\mathrm{ContOn}(f,[a,b]) \to a \in \mathbb{R} \to b \in \mathbb{R} \to \varphi_{141}$}
\AXC{$\varphi_{438}$\lbl{A1}}
\BIC{$(\mathrm{ContOn}(f,[a,b]) \to \varphi_{210}) \to \varphi_{349}$}
\end{ptree}
\noindent$\Pi_{130}$:
\begin{ptree}
\AXC{$\mathrm{ContOn}(f,[a,b]) \to a \in \mathbb{R} \to \varphi_{223}$\lbl{A1}}
\AXC{$\varphi_{335}$\lbl{A1}}
\BIC{$\mathrm{ContOn}(f,[a,b]) \to \varphi_{210}$}
\AXC{$\stackrel{\Pi_{129}}{\vdots}$}\noLine
\UIC{$(\mathrm{ContOn}(f,[a,b]) \to \varphi_{210}) \to \varphi_{349}$}
\BIC{$\varphi_{349}$}
\end{ptree}
\noindent$\Pi_{131}$:
\begin{ptree}
\AXC{$\varphi_{134}$\lbl{A2}}
\AXC{$\mathrm{ContOn}(f,[a,b]) \to a \in \mathbb{R} \to \varphi_{223}$\lbl{A1}}
\AXC{$\varphi_{339}$\lbl{A1}}
\BIC{$\mathrm{ContOn}(f,[a,b]) \to \varphi_{214}$}
\AXC{$\stackrel{\Pi_{127}}{\vdots}$}\noLine
\UIC{$\varphi_{332}$}
\BIC{$\mathrm{ContOn}(f,[a,b]) \to \varphi_{217}$}
\AXC{$\stackrel{\Pi_{130}}{\vdots}$}\noLine
\UIC{$\varphi_{349}$}
\BIC{$\mathrm{ContOn}(f,[a,b]) \to a \in \mathbb{R} \to \varphi_{230}$}
\AXC{$\varphi_{402}$\lbl{A1}}
\BIC{$\varphi_{134} \to \mathrm{ContOn}(f,[a,b]) \to a \in \mathbb{R} \to \varphi_{236}$}
\BIC{$\mathrm{ContOn}(f,[a,b]) \to a \in \mathbb{R} \to \varphi_{236}$}
\end{ptree}
\noindent$\Pi_{132}$:
\begin{ptree}
\AXC{$\stackrel{\Pi_{131}}{\vdots}$}\noLine
\UIC{$\mathrm{ContOn}(f,[a,b]) \to a \in \mathbb{R} \to \varphi_{236}$}
\AXC{$(\mathrm{ContOn}(f,[a,b]) \to a \in \mathbb{R} \to \varphi_{236}) \to \varphi_{275}$\lbl{A1}}
\BIC{$\varphi_{275}$}
\UIC{$\forall x\,(\varphi_{275})$}
\end{ptree}
\noindent$\Pi_{133}$:
\begin{ptree}
\AXC{$\varphi_{136}$\lbl{A2}}
\AXC{$\mathrm{ContOn}(f,[a,b]) \to a \in \mathbb{R} \to b \in \mathbb{R} \to \varphi_{120}$\lbl{A1}}
\AXC{$\varphi_{316}$\lbl{A1}}
\BIC{$\mathrm{ContOn}(f,[a,b]) \to \varphi_{202}$}
\AXC{$\varphi_{392}$\lbl{A1}}
\BIC{$\varphi_{136} \to \mathrm{ContOn}(f,[a,b]) \to \varphi_{218}$}
\BIC{$\mathrm{ContOn}(f,[a,b]) \to \varphi_{218}$}
\end{ptree}
\noindent$\Pi_{134}$:
\begin{ptree}
\AXC{$\stackrel{\Pi_{133}}{\vdots}$}\noLine
\UIC{$\mathrm{ContOn}(f,[a,b]) \to \varphi_{218}$}
\AXC{$(\mathrm{ContOn}(f,[a,b]) \to \varphi_{218}) \to \varphi_{304}$\lbl{A1}}
\BIC{$\varphi_{304}$}
\end{ptree}
\noindent$\Pi_{135}$:
\begin{ptree}
\AXC{$\varphi_{52}$\lbl{A2}}
\AXC{$\varphi_{52} \to \varphi_{50}$\lbl{A1}}
\BIC{$\varphi_{50}$}
\UIC{$\forall x\,(\varphi_{50})$}
\AXC{$\stackrel{\Pi_{125}}{\vdots}$}\noLine
\UIC{$\mathrm{ContOn}(f,[a,b]) \to a \in \mathbb{R} \to b \in \mathbb{R} \to \varphi_{56}$}
\AXC{$\stackrel{\Pi_{134}}{\vdots}$}\noLine
\UIC{$\varphi_{304}$}
\BIC{$\mathrm{ContOn}(f,[a,b]) \to a \in \mathbb{R} \to b \in \mathbb{R} \to \varphi_{123}$}
\AXC{$\varphi_{424}$\lbl{A1}}
\BIC{$\forall x\,(\varphi_{50}) \to \varphi_{382}$}
\BIC{$\varphi_{382}$}
\end{ptree}
\noindent$\Pi_{136}$:
\begin{ptree}
\AXC{$\varphi_{331}$\lbl{A4}}
\AXC{$\varphi_{387}$\lbl{A3}}
\AXC{$\stackrel{\Pi_{132}}{\vdots}$}\noLine
\UIC{$\forall x\,(\varphi_{275})$}
\AXC{$\varphi_{70}$\lbl{A4}}
\AXC{$\forall x\,(\varphi_{50}) \to \varphi_{67}$\lbl{A3}}
\AXC{$\stackrel{\Pi_{135}}{\vdots}$}\noLine
\UIC{$\varphi_{382}$}
\BIC{$\varphi_{70} \to \mathrm{ContOn}(f,[a,b]) \to a \in \mathbb{R} \to b \in \mathbb{R} \to \varphi_{124}$}
\BIC{$\mathrm{ContOn}(f,[a,b]) \to a \in \mathbb{R} \to b \in \mathbb{R} \to \varphi_{124}$}
\AXC{$\varphi_{503}$\lbl{A1}}
\BIC{$\forall x\,(\varphi_{275}) \to \varphi_{490}$}
\BIC{$\varphi_{490}$}
\BIC{$\varphi_{331} \to \mathrm{ContOn}(f,[a,b]) \to a \in \mathbb{R} \to b \in \mathbb{R} \to \varphi_{155}$}
\BIC{$\mathrm{ContOn}(f,[a,b]) \to a \in \mathbb{R} \to b \in \mathbb{R} \to \varphi_{155}$}
\end{ptree}
\noindent$\Pi_{137}$:
\begin{ptree}
\AXC{$\varphi_{234} \to \varphi_{239}$\lbl{A3}}
\AXC{$\stackrel{\Pi_{136}}{\vdots}$}\noLine
\UIC{$\mathrm{ContOn}(f,[a,b]) \to a \in \mathbb{R} \to b \in \mathbb{R} \to \varphi_{155}$}
\AXC{$\varphi_{329}$\lbl{A1}}
\BIC{$\varphi_{113} \to k \in \mathbb{N} \to \exists x \in \mathbb{R}\,(x \in [a,b] \land \underline{k} < f(x))$}
\UIC{$\varphi_{234}$}
\AXC{$\varphi_{234} \to (\varphi_{234} \to \varphi_{239}) \to \varphi_{345}$\lbl{A1}}
\BIC{$(\varphi_{234} \to \varphi_{239}) \to \varphi_{345}$}
\BIC{$\varphi_{345}$}
\end{ptree}
\noindent$\Pi_{138}$:
\begin{ptree}
\AXC{$\forall a\,\forall b\,\forall f\,(\varphi_{222})$\lbl{[CC3]}}
\AXC{$\forall a\,\forall b\,\forall f\,(\varphi_{222}) \to \forall b\,\forall f\,(\varphi_{222})$\lbl{A2}}
\BIC{$\forall b\,\forall f\,(\varphi_{222})$}
\AXC{$\forall b\,\forall f\,(\varphi_{222}) \to \forall f\,(\varphi_{222})$\lbl{A2}}
\BIC{$\forall f\,(\varphi_{222})$}
\AXC{$\forall f\,(\varphi_{222}) \to \varphi_{222}$\lbl{A2}}
\BIC{$\varphi_{222}$}
\end{ptree}
\noindent$\Pi_{139}$:
\begin{ptree}
\AXC{$\varphi_{113} \to \forall k\,(\varphi_{113})$\lbl{A4}}
\AXC{$\stackrel{\Pi_{137}}{\vdots}$}\noLine
\UIC{$\varphi_{345}$}
\BIC{$\mathrm{ContOn}(f,[a,b]) \to a \in \mathbb{R} \to b \in \mathbb{R} \to \varphi_{160}$}
\AXC{$\stackrel{\Pi_{138}}{\vdots}$}\noLine
\UIC{$\varphi_{222}$}
\AXC{$\varphi_{222} \to \varphi_{334}$\lbl{A1}}
\BIC{$\varphi_{334}$}
\BIC{$\mathrm{ContOn}(f,[a,b]) \to a \in \mathbb{R} \to \varphi_{211}$}
\AXC{$\varphi_{497}$\lbl{A1}}
\BIC{$\forall q\,(\varphi_{113} \land \neg 0 < 0 \to \neg \varphi_{47}) \to \varphi_{365} \to \varphi_{324}$}
\end{ptree}
\noindent$\Pi_{140}$:
\begin{ptree}
\AXC{$\varphi_{365}$\lbl{A3}}
\AXC{$\varphi_{113} \land \varphi_{47} \land \neg 0 < 0 \to \forall s\,(\varphi_{113} \land \varphi_{47} \land \neg 0 < 0)$\lbl{A4}}
\AXC{$\stackrel{\Pi_{118}}{\vdots}$}\noLine
\UIC{$\varphi_{404}$}
\BIC{$\mathrm{ContOn}(f,[a,b]) \to a \in \mathbb{R} \to \varphi_{221}$}
\AXC{$\varphi_{350}$\lbl{A1}}
\BIC{$\varphi_{113} \land \neg 0 < 0 \to \neg \varphi_{47}$}
\UIC{$\forall q\,(\varphi_{113} \land \neg 0 < 0 \to \neg \varphi_{47})$}
\AXC{$\stackrel{\Pi_{139}}{\vdots}$}\noLine
\UIC{$\forall q\,(\varphi_{113} \land \neg 0 < 0 \to \neg \varphi_{47}) \to \varphi_{365} \to \varphi_{324}$}
\BIC{$\varphi_{365} \to \varphi_{324}$}
\BIC{$\varphi_{324}$}
\end{ptree}
\noindent$\Pi_{141}$:
\begin{ptree}
\AXC{$\varphi_{113} \land \neg 0 < 0 \to \forall q\,(\varphi_{113} \land \neg 0 < 0)$\lbl{A4}}
\AXC{$\stackrel{\Pi_{140}}{\vdots}$}\noLine
\UIC{$\varphi_{324}$}
\BIC{$\mathrm{ContOn}(f,[a,b]) \to a \in \mathbb{R} \to \varphi_{36}$}
\AXC{$(\mathrm{ContOn}(f,[a,b]) \to a \in \mathbb{R} \to \varphi_{36}) \to \varphi_{276}$\lbl{A1}}
\BIC{$\varphi_{276}$}
\end{ptree}
\noindent 主木:
\begin{ptree}
\AXC{$\stackrel{\Pi_{1}}{\vdots}$}\noLine
\UIC{$\mathrm{ContOn}(f,[a,b]) \to a \in \mathbb{R} \to \varphi_{37}$}
\AXC{$\stackrel{\Pi_{141}}{\vdots}$}\noLine
\UIC{$\varphi_{276}$}
\BIC{$\mathrm{ContOn}(f,[a,b]) \to \varphi_{45}$}
\UIC{$\forall b\,(\mathrm{ContOn}(f,[a,b]) \to \varphi_{45})$}
\UIC{$\forall a\,\forall b\,(\mathrm{ContOn}(f,[a,b]) \to \varphi_{45})$}
\UIC{$\forall f\,\forall a\,\forall b\,(\mathrm{ContOn}(f,[a,b]) \to \varphi_{45})$}
\end{ptree}
\begin{abbrevlist}
\item $\varphi_{1} :\equiv \mathrm{ContOn}(f,[a,b]) \to \mathrm{ContOn}(f,[a,b])$
\item $\varphi_{2} :\equiv \neg (\mathrm{SubIdx}(s) \land c \in [a,b] \land \mathrm{Conv}(q \circ s,c))$
\item $\varphi_{3} :\equiv \exists c\,(\mathrm{SubIdx}(s) \land c \in [a,b] \land \mathrm{Conv}(q \circ s,c))$
\item $\varphi_{4} :\equiv \mathrm{ContOn}(f,[a,b]) \to a \in \mathbb{R} \to b \in \mathbb{R} \to \mathrm{RFn}(f)$
\item $\varphi_{5} :\equiv b \in \mathbb{R} \to \forall c \in [a,b]\,(\mathrm{ContIn}(f,[a,b],c))$
\item $\varphi_{6} :\equiv \mathrm{SubIdx}(s) \land c \in [a,b] \land \mathrm{Conv}(q \circ s,c) \to 0 < 0$
\item $\varphi_{7} :\equiv \forall x\,([a,b] \subseteq \mathrm{dom}(f) \to x \in [a,b] \to x \in \mathrm{dom}(f))$
\item $\varphi_{8} :\equiv \mathrm{SubIdx}(s) \land c \in [a,b] \land \mathrm{Conv}(q \circ s,c) \to \mathrm{NSeq}(s)$
\item $\varphi_{9} :\equiv \mathrm{Conv}(v_{807},v_{806}) \to \mathrm{Conv}(v_{804} \circ v_{807},v_{804}(v_{806}))$
\item $\varphi_{10} :\equiv \mathrm{SubIdx}(s) \land c \in [a,b] \land \mathrm{Conv}(q \circ s,c) \to \mathrm{SubIdx}(s)$
\item $\varphi_{11} :\equiv \mathrm{SubIdx}(s) \leftrightarrow \mathrm{NSeq}(s) \land \forall k \in \mathbb{N}\,(s_{k} < s_{S(k)})$
\item $\varphi_{12} :\equiv \forall x\,\forall y\,(x = y \to y = x) \to \forall y\,(q \circ s(n) = y \to y = q \circ s(n))$
\item $\varphi_{13} :\equiv \forall x\,\forall y\,\forall z\,(x \leq y \to y < z \to x < z) \to \forall y\,\forall z\,(0 \leq y \to y < z \to 0 < z)$
\item $\varphi_{14} :\equiv k \in \mathbb{N} \to \forall M\,(\neg (M \in \mathbb{R} \land \forall x \in [a,b]\,(f(x) \leq M)))$
\item $\varphi_{15} :\equiv k \in \mathbb{N} \to \neg (x \in [a,b] \to f(x) \leq \underline{k}) \to x \in \mathrm{dom}(f)$
\item $\varphi_{16} :\equiv \forall y\,\forall z\,(f(q(s_{K})) \leq y \to y \leq z \to f(q(s_{K})) \leq z)$
\item $\varphi_{17} :\equiv \forall y\,\forall z\,(0 \leq y \to y < z \to 0 < z) \to \forall z\,(0 \leq \lvert M \rvert \to \lvert M \rvert < z \to 0 < z)$
\item $\varphi_{18} :\equiv \forall \alpha\,(\mathrm{Conv}(f \circ q \circ s,\alpha) \to \mathrm{SeqBdd}(f \circ q \circ s))$
\item $\varphi_{19} :\equiv \mathrm{SubIdx}(s) \land c \in [a,b] \land \mathrm{Conv}(q \circ s,c) \to \mathrm{Seq}(q \circ s)$
\item $\varphi_{20} :\equiv \mathrm{SubIdx}(s) \land c \in [a,b] \land \mathrm{Conv}(q \circ s,c) \to c \in [a,b]$
\item $\varphi_{21} :\equiv a \in \mathbb{R} \to b \in \mathbb{R} \to c \in [a,b] \to \mathrm{ContIn}(f,[a,b],c)$
\item $\varphi_{22} :\equiv k \in \mathbb{N} \to \neg (x \in [a,b] \to f(x) \leq \underline{k}) \to x \in [a,b]$
\item $\varphi_{23} :\equiv \forall M\,(\neg (M \in \mathbb{R} \land \forall n \in \mathbb{N}\,(\lvert f \circ q \circ s(n) \rvert \leq M)))$
\item $\varphi_{24} :\equiv \mathrm{SubIdx}(s) \land c \in [a,b] \land \mathrm{Conv}(q \circ s,c) \to \mathrm{Conv}(q \circ s,c)$
\item $\varphi_{25} :\equiv k \in \mathbb{N} \to \neg (x \in [a,b] \to f(x) \leq \underline{k}) \to \neg f(x) \leq \underline{k}$
\item $\varphi_{26} :\equiv k \in \mathbb{N} \to \neg (x \in [a,b] \to f(x) \leq \underline{k}) \to f(x) \in \mathbb{R}$
\item $\varphi_{27} :\equiv \forall A\,(\mathrm{ContOn}(f,A) \leftrightarrow \mathrm{RFn}(f) \land A \subseteq \mathrm{dom}(f) \land \forall c \in A\,(\mathrm{ContIn}(f,A,c)))$
\item $\varphi_{28} :\equiv \mathrm{ContOn}(f,[a,b]) \to a \in \mathbb{R} \to b \in \mathbb{R} \to [a,b] \subseteq \mathrm{dom}(f)$
\item $\varphi_{29} :\equiv \mathrm{RFn}(f) \to q(s_{K}) \in \mathrm{dom}(f) \to f(q(s_{K})) \in \mathbb{R}$
\item $\varphi_{30} :\equiv \forall y\,\forall z\,(f(q(s_{K})) \leq y \to y < z \to f(q(s_{K})) < z)$
\item $\varphi_{31} :\equiv \forall y\,\forall z\,(f(q(s_{K})) < y \to y \leq z \to f(q(s_{K})) < z)$
\item $\varphi_{32} :\equiv k \in \mathbb{N} \to \neg (x \in [a,b] \to f(x) \leq \underline{k}) \to \underline{k} < f(x)$
\item $\varphi_{33} :\equiv k \in \mathbb{N} \to \neg (\underline{k} \in \mathbb{R} \land \forall x \in [a,b]\,(f(x) \leq \underline{k}))$
\item $\varphi_{34} :\equiv \forall w\,(q(s_{n}) = w \to q(s_{n}) \in [a,b] \to w \in [a,b])$
\item $\varphi_{35} :\equiv \forall a\,(\mathrm{SeqBdd}(a) \leftrightarrow \mathrm{Seq}(a) \land \exists M \in \mathbb{R}\,\forall n \in \mathbb{N}\,(\lvert a_{n} \rvert \leq M))$
\item $\varphi_{36} :\equiv b \in \mathbb{R} \to \neg \exists M \in \mathbb{R}\,\forall x \in [a,b]\,(f(x) \leq M) \to 0 < 0$
\item $\varphi_{37} :\equiv b \in \mathbb{R} \to \neg \exists M \in \mathbb{R}\,\forall x \in [a,b]\,(f(x) \leq M) \to \neg 0 < 0$
\item $\varphi_{38} :\equiv \mathrm{SubIdx}(s) \land c \in [a,b] \land \mathrm{Conv}(q \circ s,c) \to n \in \mathbb{N} \to n \in \mathbb{N}$
\item $\varphi_{39} :\equiv \mathrm{SubIdx}(s) \land c \in [a,b] \land \mathrm{Conv}(q \circ s,c) \to \mathrm{ContIn}(f,[a,b],c)$
\item $\varphi_{40} :\equiv \forall y\,\forall z\,(f(q(s_{K})) < y \to y < z \to f(q(s_{K})) < z)$
\item $\varphi_{41} :\equiv \neg \exists M \in \mathbb{R}\,\forall x \in [a,b]\,(f(x) \leq M) \to k \in \mathbb{N} \to k \in \mathbb{N}$
\item $\varphi_{42} :\equiv \forall y\,\forall z\,(f \circ q \circ s(K) = y \to y = z \to f \circ q \circ s(K) = z)$
\item $\varphi_{43} :\equiv z = w \to f(q \circ s(K)) = f(z) \to f(q \circ s(K)) = f(w)$
\item $\varphi_{44} :\equiv \neg (x \in [a,b] \to f(x) \leq \underline{k}) \to \neg (x \in [a,b] \to f(x) \leq \underline{k})$
\item $\varphi_{45} :\equiv a \in \mathbb{R} \to b \in \mathbb{R} \to \exists M \in \mathbb{R}\,\forall x \in [a,b]\,(f(x) \leq M)$
\item $\varphi_{46} :\equiv \mathrm{NSeq}(s) \to \mathrm{Seq}(q \circ s) \land \forall k \in \mathbb{N}\,(q \circ s(k) = q(s_{k}))$
\item $\varphi_{47} :\equiv \mathrm{Seq}(q) \land \forall k \in \mathbb{N}\,(q(k) \in [a,b] \land \underline{k} < f(q(k)))$
\item $\varphi_{48} :\equiv \forall x\,\forall y\,\forall z\,(x < y \to y < z \to x < z) \to \forall y\,\forall z\,(\lvert M \rvert < y \to y < z \to \lvert M \rvert < z)$
\item $\varphi_{49} :\equiv \forall z\,(f(q(s_{K})) \leq \lvert M \rvert \to \lvert M \rvert < z \to f(q(s_{K})) < z)$
\item $\varphi_{50} :\equiv \forall x\,(\neg \neg (x \in [a,b] \to f(x) \leq \underline{k})) \to x \in [a,b] \to f(x) \leq \underline{k}$
\item $\varphi_{51} :\equiv \forall z\,(f(q(s_{K})) < \underline{K} \to \underline{K} \leq z \to f(q(s_{K})) < z)$
\item $\varphi_{52} :\equiv \forall x\,(\neg \neg (x \in [a,b] \to f(x) \leq \underline{k})) \to \neg \neg (x \in [a,b] \to f(x) \leq \underline{k})$
\item $\varphi_{53} :\equiv \mathrm{SubIdx}(s) \land c \in [a,b] \land \mathrm{Conv}(q \circ s,c) \to n \in \mathbb{N} \to s_{n} \in \mathbb{N}$
\item $\varphi_{54} :\equiv \mathrm{SubIdx}(s) \land c \in [a,b] \land \mathrm{Conv}(q \circ s,c) \to \mathrm{SeqBdd}(f \circ q \circ s)$
\item $\varphi_{55} :\equiv \forall b\,(1 \in \mathbb{R} \to b \in \mathbb{R} \to 0 < 1 \to 0 < b \to \exists n \in \mathbb{N}\,(b < \underline{n} \cdot 1))$
\item $\varphi_{56} :\equiv \neg \exists M \in \mathbb{R}\,\forall x \in [a,b]\,(f(x) \leq M) \to k \in \mathbb{N} \to \underline{k} \in \mathbb{R}$
\item $\varphi_{57} :\equiv \forall s\,(\mathrm{SubIdx}(s) \to \forall k \in \mathbb{N}\,(k \leq s_{k})) \to \mathrm{SubIdx}(s) \to \forall k \in \mathbb{N}\,(k \leq s_{k})$
\item $\varphi_{58} :\equiv \forall x \in \mathbb{R}\,(-\lvert x \rvert \leq x \land x \leq \lvert x \rvert) \to M \in \mathbb{R} \to -\lvert M \rvert \leq M \land M \leq \lvert M \rvert$
\item $\varphi_{59} :\equiv \forall w\,(f \circ q \circ s(K) = w \to \lvert f \circ q \circ s(K) \rvert = M \to \lvert w \rvert = M)$
\item $\varphi_{60} :\equiv M \in \mathbb{R} \land \forall n \in \mathbb{N}\,(\lvert f \circ q \circ s(n) \rvert \leq M) \to \lvert M \rvert \in \mathbb{R}$
\item $\varphi_{61} :\equiv K \in \mathbb{N} \land \lvert M \rvert + 1 < \underline{K} \cdot 1 \to K \in \mathbb{N} \land \lvert M \rvert + 1 < \underline{K} \cdot 1$
\item $\varphi_{62} :\equiv M \in \mathbb{R} \land \forall n \in \mathbb{N}\,(\lvert f \circ q \circ s(n) \rvert \leq M) \to 0 < \lvert M \rvert + 1$
\item $\varphi_{63} :\equiv \forall k\,(\mathrm{NSeq}(s) \to k \in \mathbb{N} \to s_{k} \in \mathbb{N}) \to \mathrm{NSeq}(s) \to K \in \mathbb{N} \to s_{K} \in \mathbb{N}$
\item $\varphi_{64} :\equiv \forall k\,(\mathrm{NSeq}(s) \to k \in \mathbb{N} \to s_{k} \in \mathbb{N}) \to \mathrm{NSeq}(s) \to n \in \mathbb{N} \to s_{n} \in \mathbb{N}$
\item $\varphi_{65} :\equiv [a,b] \subseteq \mathrm{dom}(f) \to q(s_{K}) \in [a,b] \to q(s_{K}) \in \mathrm{dom}(f)$
\item $\varphi_{66} :\equiv \forall a\,\forall b\,(a \in \mathbb{R} \to b \in \mathbb{R} \to 0 < a \to 0 < b \to \exists n \in \mathbb{N}\,(b < \underline{n} \cdot a))$
\item $\varphi_{67} :\equiv \forall x\,\forall x\,(\neg \neg (x \in [a,b] \to f(x) \leq \underline{k})) \to \forall x \in [a,b]\,(f(x) \leq \underline{k})$
\item $\varphi_{68} :\equiv [a,b] \subseteq \mathrm{dom}(f) \to q \circ s(n) \in [a,b] \to q \circ s(n) \in \mathrm{dom}(f)$
\item $\varphi_{69} :\equiv \forall c \in [a,b]\,(\mathrm{ContIn}(f,[a,b],c)) \to c \in [a,b] \to \mathrm{ContIn}(f,[a,b],c)$
\item $\varphi_{70} :\equiv \forall x\,(\neg \neg (x \in [a,b] \to f(x) \leq \underline{k})) \to \forall x\,\forall x\,(\neg \neg (x \in [a,b] \to f(x) \leq \underline{k}))$
\item $\varphi_{71} :\equiv f(q(s_{K})) \leq \lvert M \rvert \to \lvert M \rvert < \underline{K} \to f(q(s_{K})) < \underline{K}$
\item $\varphi_{72} :\equiv M \in \mathbb{R} \land \forall n \in \mathbb{N}\,(\lvert f \circ q \circ s(n) \rvert \leq M) \to \lvert M \rvert + 1 \in \mathbb{R}$
\item $\varphi_{73} :\equiv \mathrm{SubIdx}(s) \land c \in [a,b] \land \mathrm{Conv}(q \circ s,c) \to \mathrm{Conv}(f \circ q \circ s,f(c))$
\item $\varphi_{74} :\equiv \forall K\,(\neg (K \in \mathbb{N} \land \lvert M \rvert + 1 < \underline{K} \cdot 1)) \to \neg (n \in \mathbb{N} \land \lvert M \rvert + 1 < \underline{n} \cdot 1)$
\item $\varphi_{75} :\equiv \forall w\,(f \circ q \circ s(K) = w \to \lvert f \circ q \circ s(K) \rvert < M \to \lvert w \rvert < M)$
\item $\varphi_{76} :\equiv \forall x\,(x \in [a,b] \leftrightarrow x \in \mathbb{R} \land a \leq x \land x \leq b) \to (x \in [a,b] \leftrightarrow x \in \mathbb{R} \land a \leq x \land x \leq b)$
\item $\varphi_{77} :\equiv K \in \mathbb{N} \land \lvert M \rvert + 1 < \underline{K} \cdot 1 \to K \in \mathbb{N} \to q \circ s(K) = q(s_{K})$
\item $\varphi_{78} :\equiv M \in \mathbb{R} \land \forall n \in \mathbb{N}\,(\lvert f \circ q \circ s(n) \rvert \leq M) \to \lvert M \rvert < \lvert M \rvert + 1$
\item $\varphi_{79} :\equiv \forall s\,\forall k\,(\mathrm{NSeq}(s) \to k \in \mathbb{N} \to s_{k} \in \mathbb{N}) \to \forall k\,(\mathrm{NSeq}(s) \to k \in \mathbb{N} \to s_{k} \in \mathbb{N})$
\item $\varphi_{80} :\equiv \forall n \in \mathbb{N}\,(q(n) \in [a,b]) \to s_{n} \in \mathbb{N} \to q(s_{n}) \in [a,b]$
\item $\varphi_{81} :\equiv \forall z\,(0 \leq \lvert M \rvert \to \lvert M \rvert < z \to 0 < z) \to 0 \leq \lvert M \rvert \to \lvert M \rvert < \lvert M \rvert + 1 \to 0 < \lvert M \rvert + 1$
\item $\varphi_{82} :\equiv \forall x\,(\mathrm{RFn}(f) \to x \in \mathrm{dom}(f) \to f(x) \in \mathbb{R}) \to \mathrm{RFn}(f) \to x \in \mathrm{dom}(f) \to f(x) \in \mathbb{R}$
\item $\varphi_{83} :\equiv K \in \mathbb{N} \land \lvert M \rvert + 1 < \underline{K} \cdot 1 \to K \in \mathbb{N} \to \lvert f \circ q \circ s(K) \rvert \leq M$
\item $\varphi_{84} :\equiv K \in \mathbb{N} \land \lvert M \rvert + 1 < \underline{K} \cdot 1 \to f(q \circ s(K)) = f(q(s_{K}))$
\item $\varphi_{85} :\equiv \neg (x \in [a,b] \to f(x) \leq \underline{k}) \to x \in \mathbb{R} \land (x \in [a,b] \land \underline{k} < f(x))$
\item $\varphi_{86} :\equiv \forall a\,\forall s\,(\mathrm{Seq}(a) \to \mathrm{NSeq}(s) \to \mathrm{Seq}(a \circ s) \land \forall k \in \mathbb{N}\,(a \circ s(k) = a_{s_{k}}))$
\item $\varphi_{87} :\equiv K \in \mathbb{N} \land \lvert M \rvert + 1 < \underline{K} \cdot 1 \to \forall k \in \mathbb{N}\,(q \circ s(k) = q(s_{k}))$
\item $\varphi_{88} :\equiv \mathrm{SubIdx}(s) \land c \in [a,b] \land \mathrm{Conv}(q \circ s,c) \to n \in \mathbb{N} \to q \circ s(n) \in \mathrm{dom}(f)$
\item $\varphi_{89} :\equiv \forall n \in \mathbb{N}\,(q \circ s(n) \in [a,b]) \to n \in \mathbb{N} \to q \circ s(n) \in [a,b]$
\item $\varphi_{90} :\equiv \forall X\,\forall Y\,\forall x\,(X \subseteq Y \to x \in X \to x \in Y) \to \forall Y\,\forall x\,([a,b] \subseteq Y \to x \in [a,b] \to x \in Y)$
\item $\varphi_{91} :\equiv K \in \mathbb{N} \land \lvert M \rvert + 1 < \underline{K} \cdot 1 \to f(q(s_{K})) < f(q(s_{K}))$
\item $\varphi_{92} :\equiv K \in \mathbb{N} \land \lvert M \rvert + 1 < \underline{K} \cdot 1 \to \forall n \in \mathbb{N}\,(\lvert f \circ q \circ s(n) \rvert \leq M)$
\item $\varphi_{93} :\equiv \mathrm{SubIdx}(s) \land c \in [a,b] \land \mathrm{Conv}(q \circ s,c) \to n \in \mathbb{N} \to q(s_{n}) \in [a,b]$
\item $\varphi_{94} :\equiv \neg (x \in [a,b] \to f(x) \leq \underline{k}) \to \exists x \in \mathbb{R}\,(x \in [a,b] \land \underline{k} < f(x))$
\item $\varphi_{95} :\equiv \forall K\,(\neg (K \in \mathbb{N} \land \lvert M \rvert + 1 < \underline{K} \cdot 1)) \to \forall n\,\forall K\,(\neg (K \in \mathbb{N} \land \lvert M \rvert + 1 < \underline{K} \cdot 1))$
\item $\varphi_{96} :\equiv \forall b\,\forall x\,(x \in [a,b] \leftrightarrow x \in \mathbb{R} \land a \leq x \land x \leq b) \to \forall x\,(x \in [a,b] \leftrightarrow x \in \mathbb{R} \land a \leq x \land x \leq b)$
\item $\varphi_{97} :\equiv \forall n\,\forall K\,(\neg (K \in \mathbb{N} \land \lvert M \rvert + 1 < \underline{K} \cdot 1)) \to \forall n\,(\neg (n \in \mathbb{N} \land \lvert M \rvert + 1 < \underline{n} \cdot 1))$
\item $\varphi_{98} :\equiv \mathrm{SubIdx}(s) \land c \in [a,b] \land \mathrm{Conv}(q \circ s,c) \to n \in \mathbb{N} \to q \circ s(n) \in [a,b]$
\item $\varphi_{99} :\equiv \mathrm{SubIdx}(s) \land c \in [a,b] \land \mathrm{Conv}(q \circ s,c) \to \forall n \in \mathbb{N}\,(q \circ s(n) \in \mathrm{dom}(f))$
\item $\varphi_{100} :\equiv K \in \mathbb{N} \land \lvert M \rvert + 1 < \underline{K} \cdot 1 \to f \circ q \circ s(K) = f(q(s_{K}))$
\item $\varphi_{101} :\equiv K \in \mathbb{N} \land \lvert M \rvert + 1 < \underline{K} \cdot 1 \to f(q(s_{K})) \leq \lvert f(q(s_{K})) \rvert$
\item $\varphi_{102} :\equiv K \in \mathbb{N} \land \lvert M \rvert + 1 < \underline{K} \cdot 1 \to f \circ q \circ s(K) = f(q \circ s(K))$
\item $\varphi_{103} :\equiv \lvert M \rvert + 1 \in \mathbb{R} \to 0 < 1 \to 0 < \lvert M \rvert + 1 \to \exists n \in \mathbb{N}\,(\lvert M \rvert + 1 < \underline{n} \cdot 1)$
\item $\varphi_{104} :\equiv \forall k \in \mathbb{N}\,(q \circ s(k) = q(s_{k})) \to K \in \mathbb{N} \to q \circ s(K) = q(s_{K})$
\item $\varphi_{105} :\equiv \forall k \in \mathbb{N}\,(q \circ s(k) = q(s_{k})) \to n \in \mathbb{N} \to q \circ s(n) = q(s_{n})$
\item $\varphi_{106} :\equiv \forall f\,\forall x\,(\mathrm{RFn}(f) \to x \in \mathrm{dom}(f) \to f(x) \in \mathbb{R}) \to \forall x\,(\mathrm{RFn}(f) \to x \in \mathrm{dom}(f) \to f(x) \in \mathbb{R})$
\item $\varphi_{107} :\equiv \mathrm{SubIdx}(s) \land c \in [a,b] \land \mathrm{Conv}(q \circ s,c) \to n \in \mathbb{N} \to q \circ s(n) = q(s_{n})$
\item $\varphi_{108} :\equiv \mathrm{SubIdx}(s) \land c \in [a,b] \land \mathrm{Conv}(q \circ s,c) \to n \in \mathbb{N} \to q(s_{n}) = q \circ s(n)$
\item $\varphi_{109} :\equiv \mathrm{Seq}(f \circ q \circ s) \land \forall n \in \mathbb{N}\,(f \circ q \circ s(n) = f(q \circ s(n)))$
\item $\varphi_{110} :\equiv \forall z\,(f(q(s_{K})) < \underline{s_{K}} \to \underline{s_{K}} < z \to f(q(s_{K})) < z)$
\item $\varphi_{111} :\equiv \mathrm{SubIdx}(s) \land c \in [a,b] \land \mathrm{Conv}(q \circ s,c) \to \forall n \in \mathbb{N}\,(q \circ s(n) \in [a,b])$
\item $\varphi_{112} :\equiv f(q \circ s(K)) = f(q \circ s(K)) \to f(q \circ s(K)) = f(q(s_{K}))$
\item $\varphi_{113} :\equiv \mathrm{ContOn}(f,[a,b]) \land a \in \mathbb{R} \land b \in \mathbb{R} \land \neg \exists M \in \mathbb{R}\,\forall x \in [a,b]\,(f(x) \leq M)$
\item $\varphi_{114} :\equiv q(s_{n}) = q \circ s(n) \to q(s_{n}) \in [a,b] \to q \circ s(n) \in [a,b]$
\item $\varphi_{115} :\equiv \mathrm{ContOn}(f,[a,b]) \leftrightarrow \mathrm{RFn}(f) \land [a,b] \subseteq \mathrm{dom}(f) \land \forall c \in [a,b]\,(\mathrm{ContIn}(f,[a,b],c))$
\item $\varphi_{116} :\equiv \mathrm{SubIdx}(s) \land c \in [a,b] \land \mathrm{Conv}(q \circ s,c) \to \mathrm{SubIdx}(s) \land c \in [a,b] \land \mathrm{Conv}(q \circ s,c)$
\item $\varphi_{117} :\equiv \forall n \in \mathbb{N}\,(\lvert f \circ q \circ s(n) \rvert \leq M) \to K \in \mathbb{N} \to \lvert f \circ q \circ s(K) \rvert \leq M$
\item $\varphi_{118} :\equiv \forall a\,\forall b\,\forall x\,(x \in [a,b] \leftrightarrow x \in \mathbb{R} \land a \leq x \land x \leq b) \to \forall b\,\forall x\,(x \in [a,b] \leftrightarrow x \in \mathbb{R} \land a \leq x \land x \leq b)$
\item $\varphi_{119} :\equiv f(q \circ s(K)) = f(q(s_{K})) \to f \circ q \circ s(K) = f(q(s_{K}))$
\item $\varphi_{120} :\equiv \neg \exists M \in \mathbb{R}\,\forall x \in [a,b]\,(f(x) \leq M) \to \neg \exists M \in \mathbb{R}\,\forall x \in [a,b]\,(f(x) \leq M)$
\item $\varphi_{121} :\equiv \forall y\,\forall z\,(\lvert M \rvert < y \to y < z \to \lvert M \rvert < z) \to \forall z\,(\lvert M \rvert < \lvert M \rvert + 1 \to \lvert M \rvert + 1 < z \to \lvert M \rvert < z)$
\item $\varphi_{122} :\equiv f(q(s_{K})) < \underline{K} \to \underline{K} \leq \underline{s_{K}} \to f(q(s_{K})) < \underline{s_{K}}$
\item $\varphi_{123} :\equiv \neg \exists M \in \mathbb{R}\,\forall x \in [a,b]\,(f(x) \leq M) \to k \in \mathbb{N} \to \neg \forall x \in [a,b]\,(f(x) \leq \underline{k})$
\item $\varphi_{124} :\equiv \neg \exists M \in \mathbb{R}\,\forall x \in [a,b]\,(f(x) \leq M) \to k \in \mathbb{N} \to \exists x\,(\neg (x \in [a,b] \to f(x) \leq \underline{k}))$
\item $\varphi_{125} :\equiv \forall m\,\forall n\,(m \in \mathbb{N} \to n \in \mathbb{N} \to m \leq n \to \underline{m} \leq \underline{n}) \to \forall n\,(K \in \mathbb{N} \to n \in \mathbb{N} \to K \leq n \to \underline{K} \leq \underline{n})$
\item $\varphi_{126} :\equiv \forall n \in \mathbb{N}\,(q \circ s(n) \in [a,b]) \to \mathrm{Conv}(q \circ s,c) \to \mathrm{Conv}(f \circ q \circ s,f(c))$
\item $\varphi_{127} :\equiv \forall x\,\forall y\,(x \in \mathbb{R} \to y \in \mathbb{R} \to x + y \in \mathbb{R}) \to \forall y\,(\lvert M \rvert \in \mathbb{R} \to y \in \mathbb{R} \to \lvert M \rvert + y \in \mathbb{R})$
\item $\varphi_{128} :\equiv K \in \mathbb{N} \land \lvert M \rvert + 1 < \underline{K} \cdot 1 \to K \in \mathbb{N} \to f \circ q \circ s(K) = f(q \circ s(K))$
\item $\varphi_{129} :\equiv \forall w\,(q \circ s(K) = w \to f(q \circ s(K)) = f(q \circ s(K)) \to f(q \circ s(K)) = f(w))$
\item $\varphi_{130} :\equiv \mathrm{SubIdx}(s) \land c \in [a,b] \land \mathrm{Conv}(q \circ s,c) \to n \in \mathbb{N} \to n \in \mathbb{N} \to q \circ s(n) = q(s_{n})$
\item $\varphi_{131} :\equiv M \in \mathbb{R} \land \forall n \in \mathbb{N}\,(\lvert f \circ q \circ s(n) \rvert \leq M) \to \exists n \in \mathbb{N}\,(\lvert M \rvert + 1 < \underline{n} \cdot 1)$
\item $\varphi_{132} :\equiv \forall y\,(\lvert M \rvert \in \mathbb{R} \to y \in \mathbb{R} \to \lvert M \rvert + y \in \mathbb{R}) \to \lvert M \rvert \in \mathbb{R} \to 1 \in \mathbb{R} \to \lvert M \rvert + 1 \in \mathbb{R}$
\item $\varphi_{133} :\equiv \forall f\,\forall A\,\forall c\,\forall a\,(\mathrm{ContIn}(f,A,c) \to \mathrm{Seq}(a) \to \forall n \in \mathbb{N}\,(a_{n} \in A) \to \mathrm{Conv}(a,c) \to \mathrm{Conv}(f \circ a,f(c)))$
\item $\varphi_{134} :\equiv \forall x\,(\neg (x \in \mathbb{R} \land (x \in [a,b] \land \underline{k} < f(x)))) \to \neg (x \in \mathbb{R} \land (x \in [a,b] \land \underline{k} < f(x)))$
\item $\varphi_{135} :\equiv \forall x\,\forall y\,(x \in \mathbb{R} \to y \in \mathbb{R} \to \neg y \leq x \to x < y) \to \forall y\,(\underline{k} \in \mathbb{R} \to y \in \mathbb{R} \to \neg y \leq \underline{k} \to \underline{k} < y)$
\item $\varphi_{136} :\equiv \forall M\,(\neg (M \in \mathbb{R} \land \forall x \in [a,b]\,(f(x) \leq M))) \to \neg (\underline{k} \in \mathbb{R} \land \forall x \in [a,b]\,(f(x) \leq \underline{k}))$
\item $\varphi_{137} :\equiv K \in \mathbb{N} \land \lvert M \rvert + 1 < \underline{K} \cdot 1 \to \forall n \in \mathbb{N}\,(f \circ q \circ s(n) = f(q \circ s(n)))$
\item $\varphi_{138} :\equiv -\lvert f(q(s_{K})) \rvert \leq f(q(s_{K})) \land f(q(s_{K})) \leq \lvert f(q(s_{K})) \rvert$
\item $\varphi_{139} :\equiv f(q(s_{K})) \leq \lvert f(q(s_{K})) \rvert \to \lvert f(q(s_{K})) \rvert \leq M \to f(q(s_{K})) \leq M$
\item $\varphi_{140} :\equiv \mathrm{SubIdx}(s) \land c \in [a,b] \land \mathrm{Conv}(q \circ s,c) \to n \in \mathbb{N} \to \forall k \in \mathbb{N}\,(q \circ s(k) = q(s_{k}))$
\item $\varphi_{141} :\equiv \neg \exists M \in \mathbb{R}\,\forall x \in [a,b]\,(f(x) \leq M) \to k \in \mathbb{N} \to \neg (x \in [a,b] \to f(x) \leq \underline{k}) \to x \in \mathbb{R}$
\item $\varphi_{142} :\equiv \forall y\,(q \circ s(n) = y \to y = q \circ s(n)) \to q \circ s(n) = q(s_{n}) \to q(s_{n}) = q \circ s(n)$
\item $\varphi_{143} :\equiv \forall z\,(f(q(s_{K})) \leq \lvert f(q(s_{K})) \rvert \to \lvert f(q(s_{K})) \rvert \leq z \to f(q(s_{K})) \leq z)$
\item $\varphi_{144} :\equiv M \in \mathbb{R} \land \forall n \in \mathbb{N}\,(\lvert f \circ q \circ s(n) \rvert \leq M) \to K \in \mathbb{N} \land \lvert M \rvert + 1 < \underline{K} \cdot 1 \to 0 < 0$
\item $\varphi_{145} :\equiv \forall z\,(f \circ q \circ s(K) = f(q \circ s(K)) \to f(q \circ s(K)) = z \to f \circ q \circ s(K) = z)$
\item $\varphi_{146} :\equiv \forall A\,\forall c\,\forall a\,(\mathrm{ContIn}(v_{804},A,c) \to \mathrm{Seq}(a) \to \forall n \in \mathbb{N}\,(a_{n} \in A) \to \mathrm{Conv}(a,c) \to \mathrm{Conv}(v_{804} \circ a,v_{804}(c)))$
\item $\varphi_{147} :\equiv \varphi_{7} \to [a,b] \subseteq \mathrm{dom}(f) \to x \in [a,b] \to x \in \mathrm{dom}(f)$
\item $\varphi_{148} :\equiv K \in \mathbb{N} \land \lvert M \rvert + 1 < \underline{K} \cdot 1 \to q(s_{K}) \in [a,b] \land \underline{s_{K}} < f(q(s_{K}))$
\item $\varphi_{149} :\equiv \mathrm{SubIdx}(s) \land c \in [a,b] \land \mathrm{Conv}(q \circ s,c) \to \mathrm{Seq}(q \circ s) \land \forall k \in \mathbb{N}\,(q \circ s(k) = q(s_{k}))$
\item $\varphi_{150} :\equiv M \in \mathbb{R} \land \forall n \in \mathbb{N}\,(\lvert f \circ q \circ s(n) \rvert \leq M) \to K \in \mathbb{N} \land \lvert M \rvert + 1 < \underline{K} \cdot 1 \to M \leq \lvert M \rvert$
\item $\varphi_{151} :\equiv \mathrm{SubIdx}(s) \land c \in [a,b] \land \mathrm{Conv}(q \circ s,c) \to \exists M \in \mathbb{R}\,\forall n \in \mathbb{N}\,(\lvert f \circ q \circ s(n) \rvert \leq M)$
\item $\varphi_{152} :\equiv b \in \mathbb{R} \to \mathrm{Seq}(q) \to \forall n \in \mathbb{N}\,(q(n) \in [a,b]) \to \exists s\,(\varphi_{3})$
\item $\varphi_{153} :\equiv \forall z\,\forall w\,(z = w \to \lvert M \rvert + 1 < z \to \lvert M \rvert + 1 < w) \to \forall w\,(\underline{K} \cdot 1 = w \to \lvert M \rvert + 1 < \underline{K} \cdot 1 \to \lvert M \rvert + 1 < w)$
\item $\varphi_{154} :\equiv M \in \mathbb{R} \land \forall n \in \mathbb{N}\,(\lvert f \circ q \circ s(n) \rvert \leq M) \to K \in \mathbb{N} \land \lvert M \rvert + 1 < \underline{K} \cdot 1 \to K \in \mathbb{N}$
\item $\varphi_{155} :\equiv \neg \exists M \in \mathbb{R}\,\forall x \in [a,b]\,(f(x) \leq M) \to k \in \mathbb{N} \to \exists x \in \mathbb{R}\,(x \in [a,b] \land \underline{k} < f(x))$
\item $\varphi_{156} :\equiv \forall n\,(K \in \mathbb{N} \to n \in \mathbb{N} \to K \leq n \to \underline{K} \leq \underline{n}) \to K \in \mathbb{N} \to s_{K} \in \mathbb{N} \to K \leq s_{K} \to \underline{K} \leq \underline{s_{K}}$
\item $\varphi_{157} :\equiv \forall c\,\forall a\,(\mathrm{ContIn}(v_{804},v_{805},c) \to \mathrm{Seq}(a) \to \forall n \in \mathbb{N}\,(a_{n} \in v_{805}) \to \mathrm{Conv}(a,c) \to \mathrm{Conv}(v_{804} \circ a,v_{804}(c)))$
\item $\varphi_{158} :\equiv M \in \mathbb{R} \land \forall n \in \mathbb{N}\,(\lvert f \circ q \circ s(n) \rvert \leq M) \to K \in \mathbb{N} \land \lvert M \rvert + 1 < \underline{K} \cdot 1 \to K \leq s_{K}$
\item $\varphi_{159} :\equiv f \circ q \circ s(K) = f(q(s_{K})) \to \lvert f \circ q \circ s(K) \rvert = M \to \lvert f(q(s_{K})) \rvert = M$
\item $\varphi_{160} :\equiv \neg \exists M \in \mathbb{R}\,\forall x \in [a,b]\,(f(x) \leq M) \to \forall k \in \mathbb{N}\,\exists x \in \mathbb{R}\,(x \in [a,b] \land \underline{k} < f(x))$
\item $\varphi_{161} :\equiv \mathrm{SubIdx}(s) \land c \in [a,b] \land \mathrm{Conv}(q \circ s,c) \to M \in \mathbb{R} \land \forall n \in \mathbb{N}\,(\lvert f \circ q \circ s(n) \rvert \leq M) \to 0 < 1$
\item $\varphi_{162} :\equiv \mathrm{SubIdx}(s) \land c \in [a,b] \land \mathrm{Conv}(q \circ s,c) \to M \in \mathbb{R} \land \forall n \in \mathbb{N}\,(\lvert f \circ q \circ s(n) \rvert \leq M) \to 0 < 0$
\item $\varphi_{163} :\equiv M \in \mathbb{R} \land \forall n \in \mathbb{N}\,(\lvert f \circ q \circ s(n) \rvert \leq M) \to K \in \mathbb{N} \land \lvert M \rvert + 1 < \underline{K} \cdot 1 \to \lvert M \rvert \in \mathbb{R}$
\item $\varphi_{164} :\equiv \forall a\,(\mathrm{ContIn}(v_{804},v_{805},v_{806}) \to \mathrm{Seq}(a) \to \forall n \in \mathbb{N}\,(a_{n} \in v_{805}) \to \mathrm{Conv}(a,v_{806}) \to \mathrm{Conv}(v_{804} \circ a,v_{804}(v_{806})))$
\item $\varphi_{165} :\equiv M \in \mathbb{R} \land \forall n \in \mathbb{N}\,(\lvert f \circ q \circ s(n) \rvert \leq M) \to K \in \mathbb{N} \land \lvert M \rvert + 1 < \underline{K} \cdot 1 \to s_{K} \in \mathbb{N}$
\item $\varphi_{166} :\equiv M \in \mathbb{R} \land \forall n \in \mathbb{N}\,(\lvert f \circ q \circ s(n) \rvert \leq M) \to K \in \mathbb{N} \land \lvert M \rvert + 1 < \underline{K} \cdot 1 \to \underline{K} \in \mathbb{R}$
\item $\varphi_{167} :\equiv \mathrm{SubIdx}(s) \land c \in [a,b] \land \mathrm{Conv}(q \circ s,c) \to M \in \mathbb{R} \land \forall n \in \mathbb{N}\,(\lvert f \circ q \circ s(n) \rvert \leq M) \to 0 \leq \lvert M \rvert$
\item $\varphi_{168} :\equiv \mathrm{SeqBdd}(f \circ q \circ s) \leftrightarrow \mathrm{Seq}(f \circ q \circ s) \land \exists M \in \mathbb{R}\,\forall n \in \mathbb{N}\,(\lvert f \circ q \circ s(n) \rvert \leq M)$
\item $\varphi_{169} :\equiv M \in \mathbb{R} \land \forall n \in \mathbb{N}\,(\lvert f \circ q \circ s(n) \rvert \leq M) \to K \in \mathbb{N} \land \lvert M \rvert + 1 < \underline{K} \cdot 1 \to \lvert M \rvert < \underline{K}$
\item $\varphi_{170} :\equiv f \circ q \circ s(K) = f(q(s_{K})) \to \lvert f \circ q \circ s(K) \rvert < M \to \lvert f(q(s_{K})) \rvert < M$
\item $\varphi_{171} :\equiv \mathrm{SubIdx}(s) \land c \in [a,b] \land \mathrm{Conv}(q \circ s,c) \to M \in \mathbb{R} \land \forall n \in \mathbb{N}\,(\lvert f \circ q \circ s(n) \rvert \leq M) \to M \in \mathbb{R}$
\item $\varphi_{172} :\equiv \forall z\,(\lvert M \rvert < \lvert M \rvert + 1 \to \lvert M \rvert + 1 < z \to \lvert M \rvert < z) \to \lvert M \rvert < \lvert M \rvert + 1 \to \lvert M \rvert + 1 < \underline{K} \to \lvert M \rvert < \underline{K}$
\item $\varphi_{173} :\equiv \varphi_{16} \to \forall z\,(f(q(s_{K})) \leq M \to M \leq z \to f(q(s_{K})) \leq z)$
\item $\varphi_{174} :\equiv M \in \mathbb{R} \land \forall n \in \mathbb{N}\,(\lvert f \circ q \circ s(n) \rvert \leq M) \to K \in \mathbb{N} \land \lvert M \rvert + 1 < \underline{K} \cdot 1 \to \underline{K} \cdot 1 = \underline{K}$
\item $\varphi_{175} :\equiv M \in \mathbb{R} \land \forall n \in \mathbb{N}\,(\lvert f \circ q \circ s(n) \rvert \leq M) \to M \in \mathbb{R} \land \forall n \in \mathbb{N}\,(\lvert f \circ q \circ s(n) \rvert \leq M)$
\item $\varphi_{176} :\equiv \forall a\,(\mathrm{RFn}(f) \to \mathrm{Seq}(a) \to \forall n \in \mathbb{N}\,(a_{n} \in \mathrm{dom}(f)) \to \mathrm{Seq}(f \circ a) \land \forall n \in \mathbb{N}\,(f \circ a(n) = f(a_{n})))$
\item $\varphi_{177} :\equiv M \in \mathbb{R} \land \forall n \in \mathbb{N}\,(\lvert f \circ q \circ s(n) \rvert \leq M) \to K \in \mathbb{N} \land \lvert M \rvert + 1 < \underline{K} \cdot 1 \to \lvert M \rvert < \lvert M \rvert + 1$
\item $\varphi_{178} :\equiv M \in \mathbb{R} \land \forall n \in \mathbb{N}\,(\lvert f \circ q \circ s(n) \rvert \leq M) \to K \in \mathbb{N} \land \lvert M \rvert + 1 < \underline{K} \cdot 1 \to K \in \mathbb{N} \to K \leq s_{K}$
\item $\varphi_{179} :\equiv M \in \mathbb{R} \land \forall n \in \mathbb{N}\,(\lvert f \circ q \circ s(n) \rvert \leq M) \to K \in \mathbb{N} \land \lvert M \rvert + 1 < \underline{K} \cdot 1 \to \underline{K} \leq \underline{s_{K}}$
\item $\varphi_{180} :\equiv M \in \mathbb{R} \land \forall n \in \mathbb{N}\,(\lvert f \circ q \circ s(n) \rvert \leq M) \to K \in \mathbb{N} \land \lvert M \rvert + 1 < \underline{K} \cdot 1 \to \lvert M \rvert + 1 < \underline{K}$
\item $\varphi_{181} :\equiv \forall z\,(f(q(s_{K})) \leq M \to M \leq z \to f(q(s_{K})) \leq z) \to f(q(s_{K})) \leq M \to M \leq \lvert M \rvert \to f(q(s_{K})) \leq \lvert M \rvert$
\item $\varphi_{182} :\equiv \forall v_{807}\,(\mathrm{ContIn}(f,[a,b],c) \to \mathrm{Seq}(v_{807}) \to \forall n \in \mathbb{N}\,(v_{807}(n) \in [a,b]) \to \mathrm{Conv}(v_{807},c) \to \mathrm{Conv}(f \circ v_{807},f(c)))$
\item $\varphi_{183} :\equiv \mathrm{SubIdx}(s) \land c \in [a,b] \land \mathrm{Conv}(q \circ s,c) \to \varphi_{60}$
\item $\varphi_{184} :\equiv M \in \mathbb{R} \land \forall n \in \mathbb{N}\,(\lvert f \circ q \circ s(n) \rvert \leq M) \to K \in \mathbb{N} \land \lvert M \rvert + 1 < \underline{K} \cdot 1 \to -\lvert M \rvert \leq M \land M \leq \lvert M \rvert$
\item $\varphi_{185} :\equiv M \in \mathbb{R} \land \forall n \in \mathbb{N}\,(\lvert f \circ q \circ s(n) \rvert \leq M) \to K \in \mathbb{N} \land \lvert M \rvert + 1 < \underline{K} \cdot 1 \to f(q(s_{K})) \leq M$
\item $\varphi_{186} :\equiv \mathrm{SubIdx}(s) \land c \in [a,b] \land \mathrm{Conv}(q \circ s,c) \to \varphi_{62}$
\item $\varphi_{187} :\equiv \forall n \in \mathbb{N}\,(f \circ q \circ s(n) = f(q \circ s(n))) \to K \in \mathbb{N} \to f \circ q \circ s(K) = f(q \circ s(K))$
\item $\varphi_{188} :\equiv \forall k \in \mathbb{N}\,(q(k) \in [a,b] \land \underline{k} < f(q(k))) \to n \in \mathbb{N} \to q(n) \in [a,b] \land \underline{n} < f(q(n))$
\item $\varphi_{189} :\equiv \mathrm{ContIn}(v_{804},v_{805},v_{806}) \to \mathrm{Seq}(v_{807}) \to \forall n \in \mathbb{N}\,(v_{807}(n) \in v_{805}) \to \varphi_{9}$
\item $\varphi_{190} :\equiv M \in \mathbb{R} \land \forall n \in \mathbb{N}\,(\lvert f \circ q \circ s(n) \rvert \leq M) \to K \in \mathbb{N} \land \lvert M \rvert + 1 < \underline{K} \cdot 1 \to \forall k \in \mathbb{N}\,(k \leq s_{k})$
\item $\varphi_{191} :\equiv \varphi_{18} \to \mathrm{Conv}(f \circ q \circ s,f(c)) \to \mathrm{SeqBdd}(f \circ q \circ s)$
\item $\varphi_{192} :\equiv M \in \mathbb{R} \land \forall n \in \mathbb{N}\,(\lvert f \circ q \circ s(n) \rvert \leq M) \to K \in \mathbb{N} \land \lvert M \rvert + 1 < \underline{K} \cdot 1 \to q(s_{K}) \in \mathrm{dom}(f)$
\item $\varphi_{193} :\equiv f(q(s_{K})) < \underline{s_{K}} \to \underline{s_{K}} < f(q(s_{K})) \to f(q(s_{K})) < f(q(s_{K}))$
\item $\varphi_{194} :\equiv M \in \mathbb{R} \land \forall n \in \mathbb{N}\,(\lvert f \circ q \circ s(n) \rvert \leq M) \to K \in \mathbb{N} \land \lvert M \rvert + 1 < \underline{K} \cdot 1 \to \lvert M \rvert + 1 < \underline{K} \cdot 1$
\item $\varphi_{195} :\equiv \neg \exists M \in \mathbb{R}\,\forall x \in [a,b]\,(f(x) \leq M) \to \varphi_{47} \to \mathrm{Seq}(q)$
\item $\varphi_{196} :\equiv \mathrm{SubIdx}(s) \land c \in [a,b] \land \mathrm{Conv}(q \circ s,c) \to \varphi_{72}$
\item $\varphi_{197} :\equiv M \in \mathbb{R} \land \forall n \in \mathbb{N}\,(\lvert f \circ q \circ s(n) \rvert \leq M) \to K \in \mathbb{N} \land \lvert M \rvert + 1 < \underline{K} \cdot 1 \to \lvert f(q(s_{K})) \rvert \leq M$
\item $\varphi_{198} :\equiv M \in \mathbb{R} \land \forall n \in \mathbb{N}\,(\lvert f \circ q \circ s(n) \rvert \leq M) \to K \in \mathbb{N} \land \lvert M \rvert + 1 < \underline{K} \cdot 1 \to q(s_{K}) \in [a,b]$
\item $\varphi_{199} :\equiv M \in \mathbb{R} \land \forall n \in \mathbb{N}\,(\lvert f \circ q \circ s(n) \rvert \leq M) \to K \in \mathbb{N} \land \lvert M \rvert + 1 < \underline{K} \cdot 1 \to f(q(s_{K})) \leq \lvert M \rvert$
\item $\varphi_{200} :\equiv \mathrm{SubIdx}(s) \land c \in [a,b] \land \mathrm{Conv}(q \circ s,c) \to \varphi_{78}$
\item $\varphi_{201} :\equiv \forall x\,\forall a\,\forall b\,(a \in \mathbb{R} \to b \in \mathbb{R} \to \mathrm{Seq}(x) \to \forall n \in \mathbb{N}\,(x(n) \in [a,b]) \to \exists s\,\exists c\,(\mathrm{SubIdx}(s) \land c \in [a,b] \land \mathrm{Conv}(x \circ s,c)))$
\item $\varphi_{202} :\equiv a \in \mathbb{R} \to b \in \mathbb{R} \to \neg \exists M \in \mathbb{R}\,\forall x \in [a,b]\,(f(x) \leq M) \to \varphi_{14}$
\item $\varphi_{203} :\equiv M \in \mathbb{R} \land \forall n \in \mathbb{N}\,(\lvert f \circ q \circ s(n) \rvert \leq M) \to K \in \mathbb{N} \land \lvert M \rvert + 1 < \underline{K} \cdot 1 \to f(q(s_{K})) \in \mathbb{R}$
\item $\varphi_{204} :\equiv M \in \mathbb{R} \land \forall n \in \mathbb{N}\,(\lvert f \circ q \circ s(n) \rvert \leq M) \to K \in \mathbb{N} \land \lvert M \rvert + 1 < \underline{K} \cdot 1 \to q \circ s(K) = q(s_{K})$
\item $\varphi_{205} :\equiv a \in \mathbb{R} \to b \in \mathbb{R} \to \neg \exists M \in \mathbb{R}\,\forall x \in [a,b]\,(f(x) \leq M) \to \varphi_{15}$
\item $\varphi_{206} :\equiv M \in \mathbb{R} \land \forall n \in \mathbb{N}\,(\lvert f \circ q \circ s(n) \rvert \leq M) \to K \in \mathbb{N} \land \lvert M \rvert + 1 < \underline{K} \cdot 1 \to f(q(s_{K})) < \underline{K}$
\item $\varphi_{207} :\equiv \forall v_{805}\,\forall v_{806}\,\forall v_{807}\,(\mathrm{ContIn}(f,v_{805},v_{806}) \to \mathrm{Seq}(v_{807}) \to \forall n \in \mathbb{N}\,(v_{807}(n) \in v_{805}) \to \mathrm{Conv}(v_{807},v_{806}) \to \mathrm{Conv}(f \circ v_{807},f(v_{806})))$
\item $\varphi_{208} :\equiv \forall y\,(\underline{k} \in \mathbb{R} \to y \in \mathbb{R} \to \neg y \leq \underline{k} \to \underline{k} < y) \to \underline{k} \in \mathbb{R} \to f(x) \in \mathbb{R} \to \neg f(x) \leq \underline{k} \to \underline{k} < f(x)$
\item $\varphi_{209} :\equiv M \in \mathbb{R} \land \forall n \in \mathbb{N}\,(\lvert f \circ q \circ s(n) \rvert \leq M) \to K \in \mathbb{N} \land \lvert M \rvert + 1 < \underline{K} \cdot 1 \to \lvert f \circ q \circ s(K) \rvert \leq M$
\item $\varphi_{210} :\equiv a \in \mathbb{R} \to b \in \mathbb{R} \to \neg \exists M \in \mathbb{R}\,\forall x \in [a,b]\,(f(x) \leq M) \to \varphi_{22}$
\item $\varphi_{211} :\equiv b \in \mathbb{R} \to \neg \exists M \in \mathbb{R}\,\forall x \in [a,b]\,(f(x) \leq M) \to \exists q\,(\varphi_{47})$
\item $\varphi_{212} :\equiv \forall w\,(\underline{K} \cdot 1 = w \to \lvert M \rvert + 1 < \underline{K} \cdot 1 \to \lvert M \rvert + 1 < w) \to \underline{K} \cdot 1 = \underline{K} \to \lvert M \rvert + 1 < \underline{K} \cdot 1 \to \lvert M \rvert + 1 < \underline{K}$
\item $\varphi_{213} :\equiv \mathrm{SubIdx}(s) \land c \in [a,b] \land \mathrm{Conv}(q \circ s,c) \to \varphi_{109}$
\item $\varphi_{214} :\equiv a \in \mathbb{R} \to b \in \mathbb{R} \to \neg \exists M \in \mathbb{R}\,\forall x \in [a,b]\,(f(x) \leq M) \to \varphi_{25}$
\item $\varphi_{215} :\equiv a \in \mathbb{R} \to b \in \mathbb{R} \to \neg \exists M \in \mathbb{R}\,\forall x \in [a,b]\,(f(x) \leq M) \to \varphi_{26}$
\item $\varphi_{216} :\equiv \forall v_{806}\,\forall v_{807}\,(\mathrm{ContIn}(f,[a,b],v_{806}) \to \mathrm{Seq}(v_{807}) \to \forall n \in \mathbb{N}\,(v_{807}(n) \in [a,b]) \to \mathrm{Conv}(v_{807},v_{806}) \to \mathrm{Conv}(f \circ v_{807},f(v_{806})))$
\item $\varphi_{217} :\equiv a \in \mathbb{R} \to b \in \mathbb{R} \to \neg \exists M \in \mathbb{R}\,\forall x \in [a,b]\,(f(x) \leq M) \to \varphi_{32}$
\item $\varphi_{218} :\equiv a \in \mathbb{R} \to b \in \mathbb{R} \to \neg \exists M \in \mathbb{R}\,\forall x \in [a,b]\,(f(x) \leq M) \to \varphi_{33}$
\item $\varphi_{219} :\equiv M \in \mathbb{R} \land \forall n \in \mathbb{N}\,(\lvert f \circ q \circ s(n) \rvert \leq M) \to K \in \mathbb{N} \land \lvert M \rvert + 1 < \underline{K} \cdot 1 \to \underline{s_{K}} < f(q(s_{K}))$
\item $\varphi_{220} :\equiv M \in \mathbb{R} \land \forall n \in \mathbb{N}\,(\lvert f \circ q \circ s(n) \rvert \leq M) \to K \in \mathbb{N} \land \lvert M \rvert + 1 < \underline{K} \cdot 1 \to f(q(s_{K})) < \underline{s_{K}}$
\item $\varphi_{221} :\equiv b \in \mathbb{R} \to \neg \exists M \in \mathbb{R}\,\forall x \in [a,b]\,(f(x) \leq M) \to \varphi_{47} \to 0 < 0$
\item $\varphi_{222} :\equiv \forall k \in \mathbb{N}\,\exists x \in \mathbb{R}\,(x \in [a,b] \land \underline{k} < f(x)) \to \exists q\,(\varphi_{47})$
\item $\varphi_{223} :\equiv b \in \mathbb{R} \to \neg \exists M \in \mathbb{R}\,\forall x \in [a,b]\,(f(x) \leq M) \to k \in \mathbb{N} \to \varphi_{44}$
\item $\varphi_{224} :\equiv \neg \exists M \in \mathbb{R}\,\forall x \in [a,b]\,(f(x) \leq M) \to \varphi_{47} \to n \in \mathbb{N} \to n \in \mathbb{N}$
\item $\varphi_{225} :\equiv \mathrm{Seq}(q \circ s) \to \forall n \in \mathbb{N}\,(q \circ s(n) \in \mathrm{dom}(f)) \to \varphi_{109}$
\item $\varphi_{226} :\equiv \mathrm{SubIdx}(s) \land c \in [a,b] \land \mathrm{Conv}(q \circ s,c) \to \varphi_{131}$
\item $\varphi_{227} :\equiv \forall k \in \mathbb{N}\,(q(k) \in [a,b] \land \underline{k} < f(q(k))) \to s_{K} \in \mathbb{N} \to q(s_{K}) \in [a,b] \land \underline{s_{K}} < f(q(s_{K}))$
\item $\varphi_{228} :\equiv \neg \exists M \in \mathbb{R}\,\forall x \in [a,b]\,(f(x) \leq M) \to \varphi_{47} \to n \in \mathbb{N} \to q(n) \in [a,b]$
\item $\varphi_{229} :\equiv \mathrm{SubIdx}(s) \land c \in [a,b] \land \mathrm{Conv}(q \circ s,c) \to \varphi_{144}$
\item $\varphi_{230} :\equiv b \in \mathbb{R} \to \neg \exists M \in \mathbb{R}\,\forall x \in [a,b]\,(f(x) \leq M) \to k \in \mathbb{N} \to \varphi_{85}$
\item $\varphi_{231} :\equiv \mathrm{SubIdx}(s) \land c \in [a,b] \land \mathrm{Conv}(q \circ s,c) \to \varphi_{150}$
\item $\varphi_{232} :\equiv \mathrm{SubIdx}(s) \land c \in [a,b] \land \mathrm{Conv}(q \circ s,c) \to \varphi_{154}$
\item $\varphi_{233} :\equiv \neg \exists M \in \mathbb{R}\,\forall x \in [a,b]\,(f(x) \leq M) \to \varphi_{47} \to \forall n \in \mathbb{N}\,(q(n) \in [a,b])$
\item $\varphi_{234} :\equiv \forall k\,(\varphi_{113} \to k \in \mathbb{N} \to \exists x \in \mathbb{R}\,(x \in [a,b] \land \underline{k} < f(x)))$
\item $\varphi_{235} :\equiv \mathrm{SubIdx}(s) \land c \in [a,b] \land \mathrm{Conv}(q \circ s,c) \to \varphi_{158}$
\item $\varphi_{236} :\equiv b \in \mathbb{R} \to \neg \exists M \in \mathbb{R}\,\forall x \in [a,b]\,(f(x) \leq M) \to k \in \mathbb{N} \to \varphi_{94}$
\item $\varphi_{237} :\equiv \forall x\,(\varphi_{113} \land k \in \mathbb{N} \land \neg \exists x \in \mathbb{R}\,(x \in [a,b] \land \underline{k} < f(x)))$
\item $\varphi_{238} :\equiv \varphi_{47} \to s_{K} \in \mathbb{N} \to q(s_{K}) \in [a,b] \land \underline{s_{K}} < f(q(s_{K}))$
\item $\varphi_{239} :\equiv \forall k\,(\varphi_{113}) \to \forall k \in \mathbb{N}\,\exists x \in \mathbb{R}\,(x \in [a,b] \land \underline{k} < f(x))$
\item $\varphi_{240} :\equiv \mathrm{SubIdx}(s) \land c \in [a,b] \land \mathrm{Conv}(q \circ s,c) \to \varphi_{163}$
\item $\varphi_{241} :\equiv \mathrm{SubIdx}(s) \land c \in [a,b] \land \mathrm{Conv}(q \circ s,c) \to \varphi_{165}$
\item $\varphi_{242} :\equiv \mathrm{SubIdx}(s) \land c \in [a,b] \land \mathrm{Conv}(q \circ s,c) \to \varphi_{166}$
\item $\varphi_{243} :\equiv \mathrm{SubIdx}(s) \land c \in [a,b] \land \mathrm{Conv}(q \circ s,c) \to \varphi_{169}$
\item $\varphi_{244} :\equiv \mathrm{SubIdx}(s) \land c \in [a,b] \land \mathrm{Conv}(q \circ s,c) \to \varphi_{174}$
\item $\varphi_{245} :\equiv \mathrm{SubIdx}(s) \land c \in [a,b] \land \mathrm{Conv}(q \circ s,c) \to \varphi_{175}$
\item $\varphi_{246} :\equiv \neg \exists M \in \mathbb{R}\,\forall x \in [a,b]\,(f(x) \leq M) \to \varphi_{47} \to s_{n} \in \mathbb{N} \to q(s_{n}) \in [a,b]$
\item $\varphi_{247} :\equiv \mathrm{SubIdx}(s) \land c \in [a,b] \land \mathrm{Conv}(q \circ s,c) \to \varphi_{177}$
\item $\varphi_{248} :\equiv \mathrm{SubIdx}(s) \land c \in [a,b] \land \mathrm{Conv}(q \circ s,c) \to \varphi_{178}$
\item $\varphi_{249} :\equiv \mathrm{SubIdx}(s) \land c \in [a,b] \land \mathrm{Conv}(q \circ s,c) \to \varphi_{179}$
\item $\varphi_{250} :\equiv \mathrm{SubIdx}(s) \land c \in [a,b] \land \mathrm{Conv}(q \circ s,c) \to \varphi_{180}$
\item $\varphi_{251} :\equiv \mathrm{SubIdx}(s) \land c \in [a,b] \land \mathrm{Conv}(q \circ s,c) \to \varphi_{184}$
\item $\varphi_{252} :\equiv \mathrm{SubIdx}(s) \land c \in [a,b] \land \mathrm{Conv}(q \circ s,c) \to \varphi_{185}$
\item $\varphi_{253} :\equiv \mathrm{SubIdx}(s) \land c \in [a,b] \land \mathrm{Conv}(q \circ s,c) \to \varphi_{190}$
\item $\varphi_{254} :\equiv \mathrm{SubIdx}(s) \land c \in [a,b] \land \mathrm{Conv}(q \circ s,c) \to \varphi_{192}$
\item $\varphi_{255} :\equiv \mathrm{SubIdx}(s) \land c \in [a,b] \land \mathrm{Conv}(q \circ s,c) \to \varphi_{194}$
\item $\varphi_{256} :\equiv \forall x \in \mathbb{R}\,(-\lvert x \rvert \leq x \land x \leq \lvert x \rvert) \to f(q(s_{K})) \in \mathbb{R} \to \varphi_{138}$
\item $\varphi_{257} :\equiv \mathrm{SubIdx}(s) \land c \in [a,b] \land \mathrm{Conv}(q \circ s,c) \to \varphi_{197}$
\item $\varphi_{258} :\equiv \mathrm{SubIdx}(s) \land c \in [a,b] \land \mathrm{Conv}(q \circ s,c) \to \varphi_{198}$
\item $\varphi_{259} :\equiv \mathrm{SubIdx}(s) \land c \in [a,b] \land \mathrm{Conv}(q \circ s,c) \to \varphi_{199}$
\item $\varphi_{260} :\equiv \mathrm{SubIdx}(s) \land c \in [a,b] \land \mathrm{Conv}(q \circ s,c) \to \varphi_{203}$
\item $\varphi_{261} :\equiv \mathrm{SubIdx}(s) \land c \in [a,b] \land \mathrm{Conv}(q \circ s,c) \to \varphi_{204}$
\item $\varphi_{262} :\equiv \mathrm{SubIdx}(s) \land c \in [a,b] \land \mathrm{Conv}(q \circ s,c) \to \varphi_{206}$
\item $\varphi_{263} :\equiv \mathrm{SubIdx}(s) \land c \in [a,b] \land \mathrm{Conv}(q \circ s,c) \to \varphi_{209}$
\item $\varphi_{264} :\equiv \neg \exists M \in \mathbb{R}\,\forall x \in [a,b]\,(f(x) \leq M) \to \varphi_{47} \to n \in \mathbb{N} \to q(n) \in [a,b] \land \underline{n} < f(q(n))$
\item $\varphi_{265} :\equiv \neg \exists M \in \mathbb{R}\,\forall x \in [a,b]\,(f(x) \leq M) \to \varphi_{47} \to \forall k \in \mathbb{N}\,(q(k) \in [a,b] \land \underline{k} < f(q(k)))$
\item $\varphi_{266} :\equiv b \in \mathbb{R} \to \neg \exists M \in \mathbb{R}\,\forall x \in [a,b]\,(f(x) \leq M) \to \varphi_{47} \to \exists s\,(\varphi_{3})$
\item $\varphi_{267} :\equiv \mathrm{SubIdx}(s) \land c \in [a,b] \land \mathrm{Conv}(q \circ s,c) \to M \in \mathbb{R} \land \forall n \in \mathbb{N}\,(\lvert f \circ q \circ s(n) \rvert \leq M) \to \varphi_{61}$
\item $\varphi_{268} :\equiv \mathrm{SubIdx}(s) \land c \in [a,b] \land \mathrm{Conv}(q \circ s,c) \to \varphi_{219}$
\item $\varphi_{269} :\equiv \mathrm{SubIdx}(s) \land c \in [a,b] \land \mathrm{Conv}(q \circ s,c) \to \varphi_{220}$
\item $\varphi_{270} :\equiv \mathrm{SubIdx}(s) \land c \in [a,b] \land \mathrm{Conv}(q \circ s,c) \to M \in \mathbb{R} \land \forall n \in \mathbb{N}\,(\lvert f \circ q \circ s(n) \rvert \leq M) \to \varphi_{77}$
\item $\varphi_{271} :\equiv b \in \mathbb{R} \to \neg \exists M \in \mathbb{R}\,\forall x \in [a,b]\,(f(x) \leq M) \to \varphi_{47} \to \varphi_{3} \to 0 < 0$
\item $\varphi_{272} :\equiv (x \in [a,b] \leftrightarrow x \in \mathbb{R} \land a \leq x \land x \leq b) \to \mathrm{ContOn}(f,[a,b]) \to a \in \mathbb{R} \to b \in \mathbb{R} \to \varphi_{141}$
\item $\varphi_{273} :\equiv \mathrm{SubIdx}(s) \land c \in [a,b] \land \mathrm{Conv}(q \circ s,c) \to M \in \mathbb{R} \land \forall n \in \mathbb{N}\,(\lvert f \circ q \circ s(n) \rvert \leq M) \to \varphi_{83}$
\item $\varphi_{274} :\equiv \mathrm{SubIdx}(s) \land c \in [a,b] \land \mathrm{Conv}(q \circ s,c) \to M \in \mathbb{R} \land \forall n \in \mathbb{N}\,(\lvert f \circ q \circ s(n) \rvert \leq M) \to \varphi_{84}$
\item $\varphi_{275} :\equiv \varphi_{113} \land k \in \mathbb{N} \land \neg \exists x \in \mathbb{R}\,(x \in [a,b] \land \underline{k} < f(x)) \to \neg \neg (x \in [a,b] \to f(x) \leq \underline{k})$
\item $\varphi_{276} :\equiv (\mathrm{ContOn}(f,[a,b]) \to a \in \mathbb{R} \to \varphi_{37}) \to \mathrm{ContOn}(f,[a,b]) \to \varphi_{45}$
\item $\varphi_{277} :\equiv \mathrm{SubIdx}(s) \land c \in [a,b] \land \mathrm{Conv}(q \circ s,c) \to M \in \mathbb{R} \land \forall n \in \mathbb{N}\,(\lvert f \circ q \circ s(n) \rvert \leq M) \to \varphi_{87}$
\item $\varphi_{278} :\equiv \mathrm{SubIdx}(s) \land c \in [a,b] \land \mathrm{Conv}(q \circ s,c) \to M \in \mathbb{R} \land \forall n \in \mathbb{N}\,(\lvert f \circ q \circ s(n) \rvert \leq M) \to \varphi_{91}$
\item $\varphi_{279} :\equiv \mathrm{SubIdx}(s) \land c \in [a,b] \land \mathrm{Conv}(q \circ s,c) \to M \in \mathbb{R} \land \forall n \in \mathbb{N}\,(\lvert f \circ q \circ s(n) \rvert \leq M) \to \varphi_{92}$
\item $\varphi_{280} :\equiv \mathrm{SubIdx}(s) \land c \in [a,b] \land \mathrm{Conv}(q \circ s,c) \to M \in \mathbb{R} \land \forall n \in \mathbb{N}\,(\lvert f \circ q \circ s(n) \rvert \leq M) \to \varphi_{100}$
\item $\varphi_{281} :\equiv \mathrm{SubIdx}(s) \land c \in [a,b] \land \mathrm{Conv}(q \circ s,c) \to M \in \mathbb{R} \land \forall n \in \mathbb{N}\,(\lvert f \circ q \circ s(n) \rvert \leq M) \to \varphi_{101}$
\item $\varphi_{282} :\equiv \mathrm{SubIdx}(s) \land c \in [a,b] \land \mathrm{Conv}(q \circ s,c) \to M \in \mathbb{R} \land \forall n \in \mathbb{N}\,(\lvert f \circ q \circ s(n) \rvert \leq M) \to \varphi_{102}$
\item $\varphi_{283} :\equiv M \in \mathbb{R} \land \forall n \in \mathbb{N}\,(\lvert f \circ q \circ s(n) \rvert \leq M) \to K \in \mathbb{N} \land \lvert M \rvert + 1 < \underline{K} \cdot 1 \to \varphi_{138}$
\item $\varphi_{284} :\equiv b \in \mathbb{R} \to \neg \exists M \in \mathbb{R}\,\forall x \in [a,b]\,(f(x) \leq M) \to \varphi_{47} \to \varphi_{47}$
\item $\varphi_{285} :\equiv \mathrm{SubIdx}(s) \land c \in [a,b] \land \mathrm{Conv}(q \circ s,c) \to M \in \mathbb{R} \land \forall n \in \mathbb{N}\,(\lvert f \circ q \circ s(n) \rvert \leq M) \to \varphi_{128}$
\item $\varphi_{286} :\equiv \mathrm{SubIdx}(s) \land c \in [a,b] \land \mathrm{Conv}(q \circ s,c) \to M \in \mathbb{R} \land \forall n \in \mathbb{N}\,(\lvert f \circ q \circ s(n) \rvert \leq M) \to \varphi_{137}$
\item $\varphi_{287} :\equiv \mathrm{SubIdx}(s) \land c \in [a,b] \land \mathrm{Conv}(q \circ s,c) \to M \in \mathbb{R} \land \forall n \in \mathbb{N}\,(\lvert f \circ q \circ s(n) \rvert \leq M) \to \varphi_{148}$
\item $\varphi_{288} :\equiv b \in \mathbb{R} \to \neg \exists M \in \mathbb{R}\,\forall x \in [a,b]\,(f(x) \leq M) \to \varphi_{47} \to \varphi_{3} \to \varphi_{3}$
\item $\varphi_{289} :\equiv a \in \mathbb{R} \to b \in \mathbb{R} \to \neg \exists M \in \mathbb{R}\,\forall x \in [a,b]\,(f(x) \leq M) \to \varphi_{238}$
\item $\varphi_{290} :\equiv (\mathrm{ContOn}(f,[a,b]) \to a \in \mathbb{R} \to b \in \mathbb{R} \to \varphi_{41}) \to \mathrm{ContOn}(f,[a,b]) \to a \in \mathbb{R} \to b \in \mathbb{R} \to \varphi_{56}$
\item $\varphi_{291} :\equiv b \in \mathbb{R} \to \neg \exists M \in \mathbb{R}\,\forall x \in [a,b]\,(f(x) \leq M) \to \varphi_{47} \to \varphi_{3} \to \varphi_{6}$
\item $\varphi_{292} :\equiv b \in \mathbb{R} \to \neg \exists M \in \mathbb{R}\,\forall x \in [a,b]\,(f(x) \leq M) \to \varphi_{47} \to \varphi_{3} \to \varphi_{8}$
\item $\varphi_{293} :\equiv b \in \mathbb{R} \to \neg \exists M \in \mathbb{R}\,\forall x \in [a,b]\,(f(x) \leq M) \to \varphi_{47} \to \varphi_{3} \to \varphi_{10}$
\item $\varphi_{294} :\equiv b \in \mathbb{R} \to \neg \exists M \in \mathbb{R}\,\forall x \in [a,b]\,(f(x) \leq M) \to \varphi_{47} \to \varphi_{3} \to \varphi_{19}$
\item $\varphi_{295} :\equiv b \in \mathbb{R} \to \neg \exists M \in \mathbb{R}\,\forall x \in [a,b]\,(f(x) \leq M) \to \varphi_{47} \to \varphi_{3} \to \varphi_{20}$
\item $\varphi_{296} :\equiv b \in \mathbb{R} \to \neg \exists M \in \mathbb{R}\,\forall x \in [a,b]\,(f(x) \leq M) \to \varphi_{47} \to \varphi_{3} \to \varphi_{24}$
\item $\varphi_{297} :\equiv \mathrm{SubIdx}(s) \land c \in [a,b] \land \mathrm{Conv}(q \circ s,c) \to \varphi_{283}$
\item $\varphi_{298} :\equiv b \in \mathbb{R} \to \neg \exists M \in \mathbb{R}\,\forall x \in [a,b]\,(f(x) \leq M) \to \varphi_{47} \to \varphi_{3} \to \varphi_{38}$
\item $\varphi_{299} :\equiv b \in \mathbb{R} \to \neg \exists M \in \mathbb{R}\,\forall x \in [a,b]\,(f(x) \leq M) \to \varphi_{47} \to \varphi_{3} \to \varphi_{39}$
\item $\varphi_{300} :\equiv (\mathrm{ContOn}(f,[a,b]) \to a \in \mathbb{R} \to \varphi_{5}) \to \varphi_{69} \to \mathrm{ContOn}(f,[a,b]) \to \varphi_{21}$
\item $\varphi_{301} :\equiv b \in \mathbb{R} \to \neg \exists M \in \mathbb{R}\,\forall x \in [a,b]\,(f(x) \leq M) \to \varphi_{47} \to \varphi_{3} \to \varphi_{53}$
\item $\varphi_{302} :\equiv b \in \mathbb{R} \to \neg \exists M \in \mathbb{R}\,\forall x \in [a,b]\,(f(x) \leq M) \to \varphi_{47} \to \varphi_{3} \to \varphi_{54}$
\item $\varphi_{303} :\equiv \varphi_{113} \land \varphi_{47} \land \varphi_{3} \land (\mathrm{SubIdx}(s) \land c \in [a,b] \land \mathrm{Conv}(q \circ s,c))$
\item $\varphi_{304} :\equiv (\mathrm{ContOn}(f,[a,b]) \to a \in \mathbb{R} \to b \in \mathbb{R} \to \varphi_{56}) \to \mathrm{ContOn}(f,[a,b]) \to a \in \mathbb{R} \to b \in \mathbb{R} \to \varphi_{123}$
\item $\varphi_{305} :\equiv b \in \mathbb{R} \to \neg \exists M \in \mathbb{R}\,\forall x \in [a,b]\,(f(x) \leq M) \to \varphi_{47} \to \varphi_{3} \to \varphi_{73}$
\item $\varphi_{306} :\equiv b \in \mathbb{R} \to \neg \exists M \in \mathbb{R}\,\forall x \in [a,b]\,(f(x) \leq M) \to \varphi_{47} \to \varphi_{3} \to \varphi_{88}$
\item $\varphi_{307} :\equiv b \in \mathbb{R} \to \neg \exists M \in \mathbb{R}\,\forall x \in [a,b]\,(f(x) \leq M) \to \varphi_{47} \to \varphi_{3} \to \varphi_{93}$
\item $\varphi_{308} :\equiv b \in \mathbb{R} \to \neg \exists M \in \mathbb{R}\,\forall x \in [a,b]\,(f(x) \leq M) \to \varphi_{47} \to \varphi_{3} \to \varphi_{98}$
\item $\varphi_{309} :\equiv b \in \mathbb{R} \to \neg \exists M \in \mathbb{R}\,\forall x \in [a,b]\,(f(x) \leq M) \to \varphi_{47} \to \varphi_{3} \to \varphi_{99}$
\item $\varphi_{310} :\equiv b \in \mathbb{R} \to \neg \exists M \in \mathbb{R}\,\forall x \in [a,b]\,(f(x) \leq M) \to \varphi_{47} \to \varphi_{3} \to \varphi_{107}$
\item $\varphi_{311} :\equiv b \in \mathbb{R} \to \neg \exists M \in \mathbb{R}\,\forall x \in [a,b]\,(f(x) \leq M) \to \varphi_{47} \to \varphi_{3} \to \varphi_{108}$
\item $\varphi_{312} :\equiv b \in \mathbb{R} \to \neg \exists M \in \mathbb{R}\,\forall x \in [a,b]\,(f(x) \leq M) \to \varphi_{47} \to \varphi_{3} \to \varphi_{111}$
\item $\varphi_{313} :\equiv b \in \mathbb{R} \to \neg \exists M \in \mathbb{R}\,\forall x \in [a,b]\,(f(x) \leq M) \to \varphi_{47} \to \varphi_{3} \to \varphi_{116}$
\item $\varphi_{314} :\equiv b \in \mathbb{R} \to \neg \exists M \in \mathbb{R}\,\forall x \in [a,b]\,(f(x) \leq M) \to \varphi_{47} \to \varphi_{3} \to \varphi_{130}$
\item $\varphi_{315} :\equiv b \in \mathbb{R} \to \neg \exists M \in \mathbb{R}\,\forall x \in [a,b]\,(f(x) \leq M) \to \varphi_{47} \to \varphi_{3} \to \varphi_{140}$
\item $\varphi_{316} :\equiv (\mathrm{ContOn}(f,[a,b]) \to a \in \mathbb{R} \to b \in \mathbb{R} \to \varphi_{120}) \to \mathrm{ContOn}(f,[a,b]) \to \varphi_{202}$
\item $\varphi_{317} :\equiv b \in \mathbb{R} \to \neg \exists M \in \mathbb{R}\,\forall x \in [a,b]\,(f(x) \leq M) \to \varphi_{47} \to \varphi_{3} \to \varphi_{149}$
\item $\varphi_{318} :\equiv b \in \mathbb{R} \to \neg \exists M \in \mathbb{R}\,\forall x \in [a,b]\,(f(x) \leq M) \to \varphi_{47} \to \varphi_{3} \to \varphi_{151}$
\item $\varphi_{319} :\equiv b \in \mathbb{R} \to \neg \exists M \in \mathbb{R}\,\forall x \in [a,b]\,(f(x) \leq M) \to \varphi_{47} \to \varphi_{3} \to \varphi_{161}$
\item $\varphi_{320} :\equiv b \in \mathbb{R} \to \neg \exists M \in \mathbb{R}\,\forall x \in [a,b]\,(f(x) \leq M) \to \varphi_{47} \to \varphi_{3} \to \varphi_{162}$
\item $\varphi_{321} :\equiv b \in \mathbb{R} \to \neg \exists M \in \mathbb{R}\,\forall x \in [a,b]\,(f(x) \leq M) \to \varphi_{47} \to \varphi_{3} \to \varphi_{167}$
\item $\varphi_{322} :\equiv b \in \mathbb{R} \to \neg \exists M \in \mathbb{R}\,\forall x \in [a,b]\,(f(x) \leq M) \to \varphi_{47} \to \varphi_{3} \to \varphi_{171}$
\item $\varphi_{323} :\equiv b \in \mathbb{R} \to \neg \exists M \in \mathbb{R}\,\forall x \in [a,b]\,(f(x) \leq M) \to \varphi_{47} \to \varphi_{3} \to \varphi_{183}$
\item $\varphi_{324} :\equiv (\varphi_{113} \land \neg 0 < 0 \to \forall q\,(\varphi_{113} \land \neg 0 < 0)) \to \mathrm{ContOn}(f,[a,b]) \to a \in \mathbb{R} \to \varphi_{36}$
\item $\varphi_{325} :\equiv b \in \mathbb{R} \to \neg \exists M \in \mathbb{R}\,\forall x \in [a,b]\,(f(x) \leq M) \to \varphi_{47} \to \varphi_{3} \to \varphi_{186}$
\item $\varphi_{326} :\equiv (\mathrm{ContOn}(f,[a,b]) \to \varphi_{210}) \to \mathrm{ContOn}(f,[a,b]) \to \varphi_{205}$
\item $\varphi_{327} :\equiv (\mathrm{ContOn}(f,[a,b]) \to \varphi_{205}) \to \mathrm{ContOn}(f,[a,b]) \to \varphi_{215}$
\item $\varphi_{328} :\equiv b \in \mathbb{R} \to \neg \exists M \in \mathbb{R}\,\forall x \in [a,b]\,(f(x) \leq M) \to \varphi_{47} \to \varphi_{3} \to \varphi_{196}$
\item $\varphi_{329} :\equiv (\mathrm{ContOn}(f,[a,b]) \to a \in \mathbb{R} \to b \in \mathbb{R} \to \varphi_{155}) \to \varphi_{113} \to k \in \mathbb{N} \to \exists x \in \mathbb{R}\,(x \in [a,b] \land \underline{k} < f(x))$
\item $\varphi_{330} :\equiv b \in \mathbb{R} \to \neg \exists M \in \mathbb{R}\,\forall x \in [a,b]\,(f(x) \leq M) \to \varphi_{47} \to \varphi_{3} \to \varphi_{200}$
\item $\varphi_{331} :\equiv \varphi_{113} \land k \in \mathbb{N} \land \neg \exists x \in \mathbb{R}\,(x \in [a,b] \land \underline{k} < f(x)) \to \varphi_{237}$
\item $\varphi_{332} :\equiv (\mathrm{ContOn}(f,[a,b]) \to \varphi_{214}) \to \mathrm{ContOn}(f,[a,b]) \to \varphi_{217}$
\item $\varphi_{333} :\equiv b \in \mathbb{R} \to \neg \exists M \in \mathbb{R}\,\forall x \in [a,b]\,(f(x) \leq M) \to \varphi_{47} \to \varphi_{3} \to \varphi_{213}$
\item $\varphi_{334} :\equiv (\mathrm{ContOn}(f,[a,b]) \to a \in \mathbb{R} \to b \in \mathbb{R} \to \varphi_{160}) \to \mathrm{ContOn}(f,[a,b]) \to a \in \mathbb{R} \to \varphi_{211}$
\item $\varphi_{335} :\equiv (\mathrm{ContOn}(f,[a,b]) \to a \in \mathbb{R} \to \varphi_{223}) \to \mathrm{ContOn}(f,[a,b]) \to \varphi_{210}$
\item $\varphi_{336} :\equiv \varphi_{303} \land (M \in \mathbb{R} \land \forall n \in \mathbb{N}\,(\lvert f \circ q \circ s(n) \rvert \leq M)) \land \neg 0 < 0$
\item $\varphi_{337} :\equiv b \in \mathbb{R} \to \neg \exists M \in \mathbb{R}\,\forall x \in [a,b]\,(f(x) \leq M) \to \varphi_{47} \to \varphi_{3} \to \varphi_{226}$
\item $\varphi_{338} :\equiv \varphi_{303} \land \neg 0 < 0 \to \neg (M \in \mathbb{R} \land \forall n \in \mathbb{N}\,(\lvert f \circ q \circ s(n) \rvert \leq M))$
\item $\varphi_{339} :\equiv (\mathrm{ContOn}(f,[a,b]) \to a \in \mathbb{R} \to \varphi_{223}) \to \mathrm{ContOn}(f,[a,b]) \to \varphi_{214}$
\item $\varphi_{340} :\equiv b \in \mathbb{R} \to \neg \exists M \in \mathbb{R}\,\forall x \in [a,b]\,(f(x) \leq M) \to \varphi_{47} \to \varphi_{3} \to \varphi_{229}$
\item $\varphi_{341} :\equiv b \in \mathbb{R} \to \neg \exists M \in \mathbb{R}\,\forall x \in [a,b]\,(f(x) \leq M) \to \varphi_{47} \to \varphi_{3} \to \varphi_{231}$
\item $\varphi_{342} :\equiv b \in \mathbb{R} \to \neg \exists M \in \mathbb{R}\,\forall x \in [a,b]\,(f(x) \leq M) \to \varphi_{47} \to \varphi_{3} \to \varphi_{232}$
\item $\varphi_{343} :\equiv b \in \mathbb{R} \to \neg \exists M \in \mathbb{R}\,\forall x \in [a,b]\,(f(x) \leq M) \to \varphi_{47} \to \varphi_{3} \to \varphi_{235}$
\item $\varphi_{344} :\equiv b \in \mathbb{R} \to \neg \exists M \in \mathbb{R}\,\forall x \in [a,b]\,(f(x) \leq M) \to \varphi_{47} \to \varphi_{3} \to \varphi_{240}$
\item $\varphi_{345} :\equiv (\varphi_{113} \to \forall k\,(\varphi_{113})) \to \mathrm{ContOn}(f,[a,b]) \to a \in \mathbb{R} \to b \in \mathbb{R} \to \varphi_{160}$
\item $\varphi_{346} :\equiv b \in \mathbb{R} \to \neg \exists M \in \mathbb{R}\,\forall x \in [a,b]\,(f(x) \leq M) \to \varphi_{47} \to \varphi_{3} \to \varphi_{241}$
\item $\varphi_{347} :\equiv b \in \mathbb{R} \to \neg \exists M \in \mathbb{R}\,\forall x \in [a,b]\,(f(x) \leq M) \to \varphi_{47} \to \varphi_{3} \to \varphi_{242}$
\item $\varphi_{348} :\equiv b \in \mathbb{R} \to \neg \exists M \in \mathbb{R}\,\forall x \in [a,b]\,(f(x) \leq M) \to \varphi_{47} \to \varphi_{3} \to \varphi_{243}$
\item $\varphi_{349} :\equiv (\mathrm{ContOn}(f,[a,b]) \to \varphi_{217}) \to \mathrm{ContOn}(f,[a,b]) \to a \in \mathbb{R} \to \varphi_{230}$
\item $\varphi_{350} :\equiv (\mathrm{ContOn}(f,[a,b]) \to a \in \mathbb{R} \to \varphi_{221}) \to \varphi_{113} \land \neg 0 < 0 \to \neg \varphi_{47}$
\item $\varphi_{351} :\equiv b \in \mathbb{R} \to \neg \exists M \in \mathbb{R}\,\forall x \in [a,b]\,(f(x) \leq M) \to \varphi_{47} \to \varphi_{3} \to \varphi_{244}$
\item $\varphi_{352} :\equiv b \in \mathbb{R} \to \neg \exists M \in \mathbb{R}\,\forall x \in [a,b]\,(f(x) \leq M) \to \varphi_{47} \to \varphi_{3} \to \varphi_{245}$
\item $\varphi_{353} :\equiv b \in \mathbb{R} \to \neg \exists M \in \mathbb{R}\,\forall x \in [a,b]\,(f(x) \leq M) \to \varphi_{47} \to \varphi_{3} \to \varphi_{247}$
\item $\varphi_{354} :\equiv b \in \mathbb{R} \to \neg \exists M \in \mathbb{R}\,\forall x \in [a,b]\,(f(x) \leq M) \to \varphi_{47} \to \varphi_{3} \to \varphi_{248}$
\item $\varphi_{355} :\equiv b \in \mathbb{R} \to \neg \exists M \in \mathbb{R}\,\forall x \in [a,b]\,(f(x) \leq M) \to \varphi_{47} \to \varphi_{3} \to \varphi_{249}$
\item $\varphi_{356} :\equiv b \in \mathbb{R} \to \neg \exists M \in \mathbb{R}\,\forall x \in [a,b]\,(f(x) \leq M) \to \varphi_{47} \to \varphi_{3} \to \varphi_{250}$
\item $\varphi_{357} :\equiv b \in \mathbb{R} \to \neg \exists M \in \mathbb{R}\,\forall x \in [a,b]\,(f(x) \leq M) \to \varphi_{47} \to \varphi_{3} \to \varphi_{251}$
\item $\varphi_{358} :\equiv b \in \mathbb{R} \to \neg \exists M \in \mathbb{R}\,\forall x \in [a,b]\,(f(x) \leq M) \to \varphi_{47} \to \varphi_{3} \to \varphi_{252}$
\item $\varphi_{359} :\equiv b \in \mathbb{R} \to \neg \exists M \in \mathbb{R}\,\forall x \in [a,b]\,(f(x) \leq M) \to \varphi_{47} \to \varphi_{3} \to \varphi_{253}$
\item $\varphi_{360} :\equiv b \in \mathbb{R} \to \neg \exists M \in \mathbb{R}\,\forall x \in [a,b]\,(f(x) \leq M) \to \varphi_{47} \to \varphi_{3} \to \varphi_{254}$
\item $\varphi_{361} :\equiv b \in \mathbb{R} \to \neg \exists M \in \mathbb{R}\,\forall x \in [a,b]\,(f(x) \leq M) \to \varphi_{47} \to \varphi_{3} \to \varphi_{255}$
\item $\varphi_{362} :\equiv b \in \mathbb{R} \to \neg \exists M \in \mathbb{R}\,\forall x \in [a,b]\,(f(x) \leq M) \to \varphi_{47} \to \varphi_{3} \to \varphi_{257}$
\item $\varphi_{363} :\equiv b \in \mathbb{R} \to \neg \exists M \in \mathbb{R}\,\forall x \in [a,b]\,(f(x) \leq M) \to \varphi_{47} \to \varphi_{3} \to \varphi_{258}$
\item $\varphi_{364} :\equiv b \in \mathbb{R} \to \neg \exists M \in \mathbb{R}\,\forall x \in [a,b]\,(f(x) \leq M) \to \varphi_{47} \to \varphi_{3} \to \varphi_{259}$
\item $\varphi_{365} :\equiv \forall q\,(\varphi_{113} \land \neg 0 < 0 \to \neg \varphi_{47}) \to \forall q\,(\varphi_{113} \land \neg 0 < 0) \to \forall q\,(\neg \varphi_{47})$
\item $\varphi_{366} :\equiv b \in \mathbb{R} \to \neg \exists M \in \mathbb{R}\,\forall x \in [a,b]\,(f(x) \leq M) \to \varphi_{47} \to \varphi_{3} \to \varphi_{260}$
\item $\varphi_{367} :\equiv b \in \mathbb{R} \to \neg \exists M \in \mathbb{R}\,\forall x \in [a,b]\,(f(x) \leq M) \to \varphi_{47} \to \varphi_{3} \to \varphi_{261}$
\item $\varphi_{368} :\equiv b \in \mathbb{R} \to \neg \exists M \in \mathbb{R}\,\forall x \in [a,b]\,(f(x) \leq M) \to \varphi_{47} \to \varphi_{3} \to \varphi_{262}$
\item $\varphi_{369} :\equiv b \in \mathbb{R} \to \neg \exists M \in \mathbb{R}\,\forall x \in [a,b]\,(f(x) \leq M) \to \varphi_{47} \to \varphi_{3} \to \varphi_{263}$
\item $\varphi_{370} :\equiv b \in \mathbb{R} \to \neg \exists M \in \mathbb{R}\,\forall x \in [a,b]\,(f(x) \leq M) \to \varphi_{47} \to \varphi_{3} \to \varphi_{267}$
\item $\varphi_{371} :\equiv b \in \mathbb{R} \to \neg \exists M \in \mathbb{R}\,\forall x \in [a,b]\,(f(x) \leq M) \to \varphi_{47} \to \varphi_{3} \to \varphi_{268}$
\item $\varphi_{372} :\equiv b \in \mathbb{R} \to \neg \exists M \in \mathbb{R}\,\forall x \in [a,b]\,(f(x) \leq M) \to \varphi_{47} \to \varphi_{3} \to \varphi_{269}$
\item $\varphi_{373} :\equiv (q \circ s(n) = q(s_{n}) \to q(s_{n}) = q \circ s(n)) \to \mathrm{ContOn}(f,[a,b]) \to a \in \mathbb{R} \to \varphi_{311}$
\item $\varphi_{374} :\equiv b \in \mathbb{R} \to \neg \exists M \in \mathbb{R}\,\forall x \in [a,b]\,(f(x) \leq M) \to \varphi_{47} \to \varphi_{3} \to \varphi_{270}$
\item $\varphi_{375} :\equiv b \in \mathbb{R} \to \neg \exists M \in \mathbb{R}\,\forall x \in [a,b]\,(f(x) \leq M) \to \varphi_{47} \to \varphi_{3} \to \varphi_{273}$
\item $\varphi_{376} :\equiv b \in \mathbb{R} \to \neg \exists M \in \mathbb{R}\,\forall x \in [a,b]\,(f(x) \leq M) \to \varphi_{47} \to \varphi_{3} \to \varphi_{274}$
\item $\varphi_{377} :\equiv b \in \mathbb{R} \to \neg \exists M \in \mathbb{R}\,\forall x \in [a,b]\,(f(x) \leq M) \to \varphi_{47} \to \varphi_{3} \to \varphi_{277}$
\item $\varphi_{378} :\equiv b \in \mathbb{R} \to \neg \exists M \in \mathbb{R}\,\forall x \in [a,b]\,(f(x) \leq M) \to \varphi_{47} \to \varphi_{3} \to \varphi_{278}$
\item $\varphi_{379} :\equiv b \in \mathbb{R} \to \neg \exists M \in \mathbb{R}\,\forall x \in [a,b]\,(f(x) \leq M) \to \varphi_{47} \to \varphi_{3} \to \varphi_{279}$
\item $\varphi_{380} :\equiv b \in \mathbb{R} \to \neg \exists M \in \mathbb{R}\,\forall x \in [a,b]\,(f(x) \leq M) \to \varphi_{47} \to \varphi_{3} \to \varphi_{280}$
\item $\varphi_{381} :\equiv b \in \mathbb{R} \to \neg \exists M \in \mathbb{R}\,\forall x \in [a,b]\,(f(x) \leq M) \to \varphi_{47} \to \varphi_{3} \to \varphi_{281}$
\item $\varphi_{382} :\equiv (\forall x\,(\varphi_{50}) \to \varphi_{67}) \to \varphi_{70} \to \mathrm{ContOn}(f,[a,b]) \to a \in \mathbb{R} \to b \in \mathbb{R} \to \varphi_{124}$
\item $\varphi_{383} :\equiv b \in \mathbb{R} \to \neg \exists M \in \mathbb{R}\,\forall x \in [a,b]\,(f(x) \leq M) \to \varphi_{47} \to \varphi_{3} \to \varphi_{282}$
\item $\varphi_{384} :\equiv (\mathrm{ContOn}(f,[a,b]) \to a \in \mathbb{R} \to b \in \mathbb{R} \to \varphi_{224}) \to \mathrm{ContOn}(f,[a,b]) \to a \in \mathbb{R} \to b \in \mathbb{R} \to \varphi_{228}$
\item $\varphi_{385} :\equiv b \in \mathbb{R} \to \neg \exists M \in \mathbb{R}\,\forall x \in [a,b]\,(f(x) \leq M) \to \varphi_{47} \to \varphi_{3} \to \varphi_{285}$
\item $\varphi_{386} :\equiv b \in \mathbb{R} \to \neg \exists M \in \mathbb{R}\,\forall x \in [a,b]\,(f(x) \leq M) \to \varphi_{47} \to \varphi_{3} \to \varphi_{286}$
\item $\varphi_{387} :\equiv \forall x\,(\varphi_{275}) \to \varphi_{237} \to \forall x\,(\neg \neg (x \in [a,b] \to f(x) \leq \underline{k}))$
\item $\varphi_{388} :\equiv b \in \mathbb{R} \to \neg \exists M \in \mathbb{R}\,\forall x \in [a,b]\,(f(x) \leq M) \to \varphi_{47} \to \varphi_{3} \to \varphi_{287}$
\item $\varphi_{389} :\equiv (\mathrm{ContOn}(f,[a,b]) \to a \in \mathbb{R} \to b \in \mathbb{R} \to \varphi_{228}) \to \varphi_{113} \land \varphi_{47} \to n \in \mathbb{N} \to q(n) \in [a,b]$
\item $\varphi_{390} :\equiv \forall n\,(\varphi_{113} \land \varphi_{47} \to n \in \mathbb{N} \to q(n) \in [a,b]) \to \forall n\,(\varphi_{113} \land \varphi_{47}) \to \forall n \in \mathbb{N}\,(q(n) \in [a,b])$
\item $\varphi_{391} :\equiv (\mathrm{ContOn}(f,[a,b]) \to a \in \mathbb{R} \to \varphi_{284}) \to \mathrm{ContOn}(f,[a,b]) \to a \in \mathbb{R} \to b \in \mathbb{R} \to \varphi_{195}$
\item $\varphi_{392} :\equiv (\mathrm{ContOn}(f,[a,b]) \to \varphi_{202}) \to \varphi_{136} \to \mathrm{ContOn}(f,[a,b]) \to \varphi_{218}$
\item $\varphi_{393} :\equiv \neg f(q(s_{K})) < f(q(s_{K})) \to \mathrm{ContOn}(f,[a,b]) \to a \in \mathbb{R} \to \varphi_{340}$
\item $\varphi_{394} :\equiv (\mathrm{ContOn}(f,[a,b]) \to a \in \mathbb{R} \to b \in \mathbb{R} \to \varphi_{233}) \to \mathrm{ContOn}(f,[a,b]) \to a \in \mathbb{R} \to \varphi_{266}$
\item $\varphi_{395} :\equiv ([a,b] \subseteq \mathrm{dom}(f) \to x \in [a,b] \to x \in \mathrm{dom}(f)) \to \varphi_{28} \to \varphi_{326}$
\item $\varphi_{396} :\equiv b \in \mathbb{R} \to \neg \exists M \in \mathbb{R}\,\forall x \in [a,b]\,(f(x) \leq M) \to \varphi_{47} \to \varphi_{3} \to \varphi_{297}$
\item $\varphi_{397} :\equiv (\forall k \in \mathbb{N}\,(k \leq s_{k}) \to K \in \mathbb{N} \to K \leq s_{K}) \to \mathrm{ContOn}(f,[a,b]) \to a \in \mathbb{R} \to \varphi_{354}$
\item $\varphi_{398} :\equiv (\mathrm{ContOn}(f,[a,b]) \to a \in \mathbb{R} \to \varphi_{271}) \to \varphi_{113} \land \varphi_{47} \land \neg 0 < 0 \to \neg \varphi_{3}$
\item $\varphi_{399} :\equiv (\underline{K} \cdot 1 = \underline{K} \to \lvert M \rvert + 1 < \underline{K} \cdot 1 \to \lvert M \rvert + 1 < \underline{K}) \to \mathrm{ContOn}(f,[a,b]) \to a \in \mathbb{R} \to \varphi_{356}$
\item $\varphi_{400} :\equiv \forall s\,(\varphi_{113} \land \varphi_{47} \land \neg 0 < 0 \to \neg \varphi_{3}) \to \forall s\,(\varphi_{113} \land \varphi_{47} \land \neg 0 < 0) \to \forall s\,(\neg \varphi_{3})$
\item $\varphi_{401} :\equiv (\mathrm{ContOn}(f,[a,b]) \to a \in \mathbb{R} \to \varphi_{284}) \to \mathrm{ContOn}(f,[a,b]) \to a \in \mathbb{R} \to b \in \mathbb{R} \to \varphi_{265}$
\item $\varphi_{402} :\equiv (\mathrm{ContOn}(f,[a,b]) \to a \in \mathbb{R} \to \varphi_{230}) \to \varphi_{134} \to \mathrm{ContOn}(f,[a,b]) \to a \in \mathbb{R} \to \varphi_{236}$
\item $\varphi_{403} :\equiv (\mathrm{ContOn}(f,[a,b]) \to a \in \mathbb{R} \to b \in \mathbb{R} \to \varphi_{233}) \to \varphi_{80} \to \mathrm{ContOn}(f,[a,b]) \to a \in \mathbb{R} \to b \in \mathbb{R} \to \varphi_{246}$
\item $\varphi_{404} :\equiv (\varphi_{113} \land \varphi_{47} \land \neg 0 < 0 \to \forall s\,(\varphi_{113} \land \varphi_{47} \land \neg 0 < 0)) \to \mathrm{ContOn}(f,[a,b]) \to a \in \mathbb{R} \to \varphi_{221}$
\item $\varphi_{405} :\equiv (q \circ s(K) = q(s_{K}) \to \varphi_{112}) \to \mathrm{ContOn}(f,[a,b]) \to a \in \mathbb{R} \to \varphi_{376}$
\item $\varphi_{406} :\equiv (\varphi_{113} \land \varphi_{47} \to \forall n\,(\varphi_{113} \land \varphi_{47})) \to \mathrm{ContOn}(f,[a,b]) \to a \in \mathbb{R} \to b \in \mathbb{R} \to \varphi_{233}$
\item $\varphi_{407} :\equiv (\mathrm{ContOn}(f,[a,b]) \to a \in \mathbb{R} \to \varphi_{291}) \to \varphi_{113} \land \varphi_{47} \land \varphi_{3} \land \neg 0 < 0 \to \varphi_{2}$
\item $\varphi_{408} :\equiv \forall c\,(\varphi_{113} \land \varphi_{47} \land \varphi_{3} \land \neg 0 < 0 \to \varphi_{2}) \to \forall c\,(\varphi_{113} \land \varphi_{47} \land \varphi_{3} \land \neg 0 < 0) \to \forall c\,(\varphi_{2})$
\item $\varphi_{409} :\equiv (f \circ q \circ s(K) = f(q \circ s(K)) \to \varphi_{119}) \to \mathrm{ContOn}(f,[a,b]) \to a \in \mathbb{R} \to \varphi_{380}$
\item $\varphi_{410} :\equiv (\mathrm{ContOn}(f,[a,b]) \to a \in \mathbb{R} \to \varphi_{295}) \to \mathrm{ContOn}(f,[a,b]) \to a \in \mathbb{R} \to \varphi_{299}$
\item $\varphi_{411} :\equiv (\mathrm{ContOn}(f,[a,b]) \to a \in \mathbb{R} \to \varphi_{298}) \to \mathrm{ContOn}(f,[a,b]) \to a \in \mathbb{R} \to \varphi_{301}$
\item $\varphi_{412} :\equiv (\mathrm{ContOn}(f,[a,b]) \to a \in \mathbb{R} \to \varphi_{296}) \to \mathrm{ContOn}(f,[a,b]) \to a \in \mathbb{R} \to \varphi_{305}$
\item $\varphi_{413} :\equiv (\mathrm{ContOn}(f,[a,b]) \to a \in \mathbb{R} \to \varphi_{313}) \to \mathrm{ContOn}(f,[a,b]) \to a \in \mathbb{R} \to \varphi_{293}$
\item $\varphi_{414} :\equiv (\mathrm{ContOn}(f,[a,b]) \to a \in \mathbb{R} \to b \in \mathbb{R} \to \varphi_{265}) \to \varphi_{188} \to \mathrm{ContOn}(f,[a,b]) \to a \in \mathbb{R} \to b \in \mathbb{R} \to \varphi_{264}$
\item $\varphi_{415} :\equiv (\mathrm{ContOn}(f,[a,b]) \to a \in \mathbb{R} \to \varphi_{313}) \to \mathrm{ContOn}(f,[a,b]) \to a \in \mathbb{R} \to \varphi_{295}$
\item $\varphi_{416} :\equiv (\mathrm{ContOn}(f,[a,b]) \to a \in \mathbb{R} \to \varphi_{305}) \to \mathrm{ContOn}(f,[a,b]) \to a \in \mathbb{R} \to \varphi_{302}$
\item $\varphi_{417} :\equiv (\mathrm{ContOn}(f,[a,b]) \to a \in \mathbb{R} \to \varphi_{313}) \to \mathrm{ContOn}(f,[a,b]) \to a \in \mathbb{R} \to \varphi_{296}$
\item $\varphi_{418} :\equiv (\mathrm{ContOn}(f,[a,b]) \to a \in \mathbb{R} \to \varphi_{298}) \to \mathrm{ContOn}(f,[a,b]) \to a \in \mathbb{R} \to \varphi_{310}$
\item $\varphi_{419} :\equiv (\mathrm{ContOn}(f,[a,b]) \to a \in \mathbb{R} \to \varphi_{301}) \to \mathrm{ContOn}(f,[a,b]) \to a \in \mathbb{R} \to \varphi_{307}$
\item $\varphi_{420} :\equiv (\forall n\,(\varphi_{74}) \to \varphi_{97}) \to \varphi_{95} \to \mathrm{ContOn}(f,[a,b]) \to a \in \mathbb{R} \to \varphi_{320}$
\item $\varphi_{421} :\equiv (\mathrm{ContOn}(f,[a,b]) \to a \in \mathbb{R} \to b \in \mathbb{R} \to \varphi_{195}) \to \varphi_{394}$
\item $\varphi_{422} :\equiv (\mathrm{ContOn}(f,[a,b]) \to a \in \mathbb{R} \to \varphi_{292}) \to \mathrm{ContOn}(f,[a,b]) \to a \in \mathbb{R} \to \varphi_{317}$
\item $\varphi_{423} :\equiv (\mathrm{ContOn}(f,[a,b]) \to a \in \mathbb{R} \to b \in \mathbb{R} \to \varphi_{56}) \to (\mathrm{ContOn}(f,[a,b]) \to \varphi_{215}) \to \varphi_{332}$
\item $\varphi_{424} :\equiv (\mathrm{ContOn}(f,[a,b]) \to a \in \mathbb{R} \to b \in \mathbb{R} \to \varphi_{123}) \to \forall x\,(\varphi_{50}) \to \varphi_{382}$
\item $\varphi_{425} :\equiv (\mathrm{ContOn}(f,[a,b]) \to a \in \mathbb{R} \to \varphi_{306}) \to \varphi_{303} \to n \in \mathbb{N} \to q \circ s(n) \in \mathrm{dom}(f)$
\item $\varphi_{426} :\equiv (\mathrm{ContOn}(f,[a,b]) \to a \in \mathbb{R} \to \varphi_{317}) \to \mathrm{ContOn}(f,[a,b]) \to a \in \mathbb{R} \to \varphi_{294}$
\item $\varphi_{427} :\equiv (\mathrm{ContOn}(f,[a,b]) \to a \in \mathbb{R} \to \varphi_{293}) \to \varphi_{11} \to \mathrm{ContOn}(f,[a,b]) \to a \in \mathbb{R} \to \varphi_{292}$
\item $\varphi_{428} :\equiv (\mathrm{ContOn}(f,[a,b]) \to a \in \mathbb{R} \to \varphi_{308}) \to \mathrm{ContOn}(f,[a,b]) \to a \in \mathbb{R} \to \varphi_{306}$
\item $\varphi_{429} :\equiv (\mathrm{ContOn}(f,[a,b]) \to a \in \mathbb{R} \to \varphi_{308}) \to \varphi_{303} \to n \in \mathbb{N} \to q \circ s(n) \in [a,b]$
\item $\varphi_{430} :\equiv \forall n\,(\varphi_{303} \to n \in \mathbb{N} \to q \circ s(n) \in \mathrm{dom}(f)) \to \forall n\,(\varphi_{303}) \to \forall n \in \mathbb{N}\,(q \circ s(n) \in \mathrm{dom}(f))$
\item $\varphi_{431} :\equiv (\mathrm{ContOn}(f,[a,b]) \to a \in \mathbb{R} \to b \in \mathbb{R} \to \varphi_{264}) \to \varphi_{384}$
\item $\varphi_{432} :\equiv \forall n\,(\varphi_{303} \to n \in \mathbb{N} \to q \circ s(n) \in [a,b]) \to \forall n\,(\varphi_{303}) \to \forall n \in \mathbb{N}\,(q \circ s(n) \in [a,b])$
\item $\varphi_{433} :\equiv (\mathrm{ContOn}(f,[a,b]) \to a \in \mathbb{R} \to b \in \mathbb{R} \to \varphi_{265}) \to \varphi_{227} \to \mathrm{ContOn}(f,[a,b]) \to \varphi_{289}$
\item $\varphi_{434} :\equiv (\mathrm{ContOn}(f,[a,b]) \to a \in \mathbb{R} \to \varphi_{317}) \to \mathrm{ContOn}(f,[a,b]) \to a \in \mathbb{R} \to \varphi_{315}$
\item $\varphi_{435} :\equiv (\mathrm{ContOn}(f,[a,b]) \to a \in \mathbb{R} \to b \in \mathbb{R} \to b \in \mathbb{R}) \to \varphi_{421}$
\item $\varphi_{436} :\equiv (\mathrm{Conv}(f \circ q \circ s,f(c)) \to \mathrm{SeqBdd}(f \circ q \circ s)) \to \varphi_{416}$
\item $\varphi_{437} :\equiv (\mathrm{ContOn}(f,[a,b]) \to a \in \mathbb{R} \to \varphi_{309}) \to \mathrm{ContOn}(f,[a,b]) \to a \in \mathbb{R} \to \varphi_{333}$
\item $\varphi_{438} :\equiv (\mathrm{ContOn}(f,[a,b]) \to a \in \mathbb{R} \to b \in \mathbb{R} \to \varphi_{141}) \to (\mathrm{ContOn}(f,[a,b]) \to \varphi_{210}) \to \varphi_{349}$
\item $\varphi_{439} :\equiv (\mathrm{ContOn}(f,[a,b]) \to a \in \mathbb{R} \to \varphi_{293}) \to \mathrm{ContOn}(f,[a,b]) \to a \in \mathbb{R} \to \varphi_{359}$
\item $\varphi_{440} :\equiv (\mathrm{ContOn}(f,[a,b]) \to a \in \mathbb{R} \to \varphi_{322}) \to \mathrm{ContOn}(f,[a,b]) \to a \in \mathbb{R} \to \varphi_{321}$
\item $\varphi_{441} :\equiv (\underline{k} \in \mathbb{R} \to f(x) \in \mathbb{R} \to \neg f(x) \leq \underline{k} \to \underline{k} < f(x)) \to \varphi_{423}$
\item $\varphi_{442} :\equiv (\mathrm{ContOn}(f,[a,b]) \to a \in \mathbb{R} \to \varphi_{322}) \to \mathrm{ContOn}(f,[a,b]) \to a \in \mathbb{R} \to \varphi_{323}$
\item $\varphi_{443} :\equiv (\varphi_{113} \land \varphi_{47} \land \varphi_{3} \land \neg 0 < 0 \to \forall c\,(\varphi_{113} \land \varphi_{47} \land \varphi_{3} \land \neg 0 < 0)) \to \mathrm{ContOn}(f,[a,b]) \to a \in \mathbb{R} \to \varphi_{271}$
\item $\varphi_{444} :\equiv (\mathrm{ContOn}(f,[a,b]) \to a \in \mathbb{R} \to \varphi_{323}) \to \mathrm{ContOn}(f,[a,b]) \to a \in \mathbb{R} \to \varphi_{330}$
\item $\varphi_{445} :\equiv (\mathrm{ContOn}(f,[a,b]) \to a \in \mathbb{R} \to \varphi_{330}) \to \mathrm{ContOn}(f,[a,b]) \to a \in \mathbb{R} \to \varphi_{325}$
\item $\varphi_{446} :\equiv (\mathrm{ContOn}(f,[a,b]) \to a \in \mathbb{R} \to \varphi_{323}) \to 1 \in \mathbb{R} \to \mathrm{ContOn}(f,[a,b]) \to a \in \mathbb{R} \to \varphi_{328}$
\item $\varphi_{447} :\equiv (\mathrm{ContOn}(f,[a,b]) \to a \in \mathbb{R} \to \varphi_{312}) \to \varphi_{89} \to \mathrm{ContOn}(f,[a,b]) \to a \in \mathbb{R} \to \varphi_{308}$
\item $\varphi_{448} :\equiv (\mathrm{ContOn}(f,[a,b]) \to a \in \mathbb{R} \to \varphi_{307}) \to \varphi_{114} \to \mathrm{ContOn}(f,[a,b]) \to a \in \mathbb{R} \to \varphi_{308}$
\item $\varphi_{449} :\equiv (\mathrm{ContOn}(f,[a,b]) \to a \in \mathbb{R} \to \varphi_{325}) \to \mathrm{ContOn}(f,[a,b]) \to a \in \mathbb{R} \to \varphi_{337}$
\item $\varphi_{450} :\equiv (\mathrm{ContOn}(f,[a,b]) \to a \in \mathbb{R} \to \varphi_{322}) \to \mathrm{ContOn}(f,[a,b]) \to a \in \mathbb{R} \to \varphi_{344}$
\item $\varphi_{451} :\equiv (\mathrm{ContOn}(f,[a,b]) \to a \in \mathbb{R} \to \varphi_{352}) \to \mathrm{ContOn}(f,[a,b]) \to a \in \mathbb{R} \to \varphi_{322}$
\item $\varphi_{452} :\equiv (\mathrm{ContOn}(f,[a,b]) \to a \in \mathbb{R} \to \varphi_{322}) \to \mathrm{ContOn}(f,[a,b]) \to a \in \mathbb{R} \to \varphi_{357}$
\item $\varphi_{453} :\equiv (\mathrm{ContOn}(f,[a,b]) \to a \in \mathbb{R} \to \varphi_{315}) \to \varphi_{105} \to \mathrm{ContOn}(f,[a,b]) \to a \in \mathbb{R} \to \varphi_{314}$
\item $\varphi_{454} :\equiv (\mathrm{ContOn}(f,[a,b]) \to a \in \mathbb{R} \to \varphi_{317}) \to \mathrm{ContOn}(f,[a,b]) \to a \in \mathbb{R} \to \varphi_{377}$
\item $\varphi_{455} :\equiv (\mathrm{ContOn}(f,[a,b]) \to a \in \mathbb{R} \to \varphi_{302}) \to \varphi_{168} \to \mathrm{ContOn}(f,[a,b]) \to a \in \mathbb{R} \to \varphi_{318}$
\item $\varphi_{456} :\equiv (\mathrm{ContOn}(f,[a,b]) \to a \in \mathbb{R} \to \varphi_{340}) \to \varphi_{336} \to \neg (K \in \mathbb{N} \land \lvert M \rvert + 1 < \underline{K} \cdot 1)$
\item $\varphi_{457} :\equiv (\mathrm{ContOn}(f,[a,b]) \to a \in \mathbb{R} \to \varphi_{342}) \to \mathrm{ContOn}(f,[a,b]) \to a \in \mathbb{R} \to \varphi_{343}$
\item $\varphi_{458} :\equiv \forall K\,(\varphi_{336} \to \neg (K \in \mathbb{N} \land \lvert M \rvert + 1 < \underline{K} \cdot 1)) \to \forall K\,(\varphi_{336}) \to \forall K\,(\neg (K \in \mathbb{N} \land \lvert M \rvert + 1 < \underline{K} \cdot 1))$
\item $\varphi_{459} :\equiv (\mathrm{ContOn}(f,[a,b]) \to a \in \mathbb{R} \to \varphi_{342}) \to \mathrm{ContOn}(f,[a,b]) \to a \in \mathbb{R} \to \varphi_{346}$
\item $\varphi_{460} :\equiv (\mathrm{ContOn}(f,[a,b]) \to a \in \mathbb{R} \to \varphi_{342}) \to \mathrm{ContOn}(f,[a,b]) \to a \in \mathbb{R} \to \varphi_{347}$
\item $\varphi_{461} :\equiv (\mathrm{ContOn}(f,[a,b]) \to a \in \mathbb{R} \to \varphi_{343}) \to \mathrm{ContOn}(f,[a,b]) \to a \in \mathbb{R} \to \varphi_{355}$
\item $\varphi_{462} :\equiv (\mathrm{ContOn}(f,[a,b]) \to a \in \mathbb{R} \to \varphi_{357}) \to \mathrm{ContOn}(f,[a,b]) \to a \in \mathbb{R} \to \varphi_{341}$
\item $\varphi_{463} :\equiv (\mathrm{ContOn}(f,[a,b]) \to a \in \mathbb{R} \to \varphi_{347}) \to \mathrm{ContOn}(f,[a,b]) \to a \in \mathbb{R} \to \varphi_{351}$
\item $\varphi_{464} :\equiv (\mathrm{ContOn}(f,[a,b]) \to a \in \mathbb{R} \to \varphi_{344}) \to \mathrm{ContOn}(f,[a,b]) \to a \in \mathbb{R} \to \varphi_{353}$
\item $\varphi_{465} :\equiv (\mathrm{ContOn}(f,[a,b]) \to a \in \mathbb{R} \to \varphi_{341}) \to \mathrm{ContOn}(f,[a,b]) \to a \in \mathbb{R} \to \varphi_{364}$
\item $\varphi_{466} :\equiv (\mathrm{ContOn}(f,[a,b]) \to a \in \mathbb{R} \to \varphi_{356}) \to \mathrm{ContOn}(f,[a,b]) \to a \in \mathbb{R} \to \varphi_{348}$
\item $\varphi_{467} :\equiv (\mathrm{ContOn}(f,[a,b]) \to a \in \mathbb{R} \to \varphi_{342}) \to \mathrm{ContOn}(f,[a,b]) \to a \in \mathbb{R} \to \varphi_{367}$
\item $\varphi_{468} :\equiv (\mathrm{ContOn}(f,[a,b]) \to a \in \mathbb{R} \to \varphi_{342}) \to \mathrm{ContOn}(f,[a,b]) \to a \in \mathbb{R} \to \varphi_{369}$
\item $\varphi_{469} :\equiv (\mathrm{ContOn}(f,[a,b]) \to a \in \mathbb{R} \to b \in \mathbb{R} \to \varphi_{195}) \to \varphi_{422}$
\item $\varphi_{470} :\equiv (\mathrm{ContOn}(f,[a,b]) \to a \in \mathbb{R} \to \varphi_{370}) \to \mathrm{ContOn}(f,[a,b]) \to a \in \mathbb{R} \to \varphi_{342}$
\item $\varphi_{471} :\equiv (\mathrm{ContOn}(f,[a,b]) \to a \in \mathbb{R} \to \varphi_{348}) \to \mathrm{ContOn}(f,[a,b]) \to a \in \mathbb{R} \to \varphi_{368}$
\item $\varphi_{472} :\equiv (\mathrm{ContOn}(f,[a,b]) \to a \in \mathbb{R} \to \varphi_{333}) \to \mathrm{ContOn}(f,[a,b]) \to a \in \mathbb{R} \to \varphi_{386}$
\item $\varphi_{473} :\equiv (\mathrm{ContOn}(f,[a,b]) \to a \in \mathbb{R} \to \varphi_{362}) \to \mathrm{ContOn}(f,[a,b]) \to a \in \mathbb{R} \to \varphi_{358}$
\item $\varphi_{474} :\equiv (\mathrm{ContOn}(f,[a,b]) \to a \in \mathbb{R} \to \varphi_{363}) \to \mathrm{ContOn}(f,[a,b]) \to a \in \mathbb{R} \to \varphi_{360}$
\item $\varphi_{475} :\equiv (\mathrm{ContOn}(f,[a,b]) \to a \in \mathbb{R} \to \varphi_{360}) \to \mathrm{ContOn}(f,[a,b]) \to a \in \mathbb{R} \to \varphi_{366}$
\item $\varphi_{476} :\equiv (\mathrm{ContOn}(f,[a,b]) \to a \in \mathbb{R} \to \varphi_{342}) \to \mathrm{ContOn}(f,[a,b]) \to a \in \mathbb{R} \to \varphi_{383}$
\item $\varphi_{477} :\equiv (\mathrm{ContOn}(f,[a,b]) \to a \in \mathbb{R} \to \varphi_{355}) \to \mathrm{ContOn}(f,[a,b]) \to a \in \mathbb{R} \to \varphi_{372}$
\item $\varphi_{478} :\equiv (\mathrm{ContOn}(f,[a,b]) \to a \in \mathbb{R} \to \varphi_{370}) \to \mathrm{ContOn}(f,[a,b]) \to a \in \mathbb{R} \to \varphi_{361}$
\item $\varphi_{479} :\equiv (\mathrm{ContOn}(f,[a,b]) \to a \in \mathbb{R} \to \varphi_{352}) \to \mathrm{ContOn}(f,[a,b]) \to a \in \mathbb{R} \to \varphi_{379}$
\item $\varphi_{480} :\equiv (\mathrm{ContOn}(f,[a,b]) \to a \in \mathbb{R} \to b \in \mathbb{R} \to \varphi_{246}) \to \varphi_{419}$
\item $\varphi_{481} :\equiv (\mathrm{ContOn}(f,[a,b]) \to a \in \mathbb{R} \to \varphi_{346}) \to \mathrm{ContOn}(f,[a,b]) \to a \in \mathbb{R} \to \varphi_{388}$
\item $\varphi_{482} :\equiv (\mathrm{ContOn}(f,[a,b]) \to a \in \mathbb{R} \to \varphi_{371}) \to \mathrm{ContOn}(f,[a,b]) \to a \in \mathbb{R} \to \varphi_{378}$
\item $\varphi_{483} :\equiv (\varphi_{303} \land \neg 0 < 0 \to \forall M\,(\varphi_{303} \land \neg 0 < 0)) \to \mathrm{ContOn}(f,[a,b]) \to a \in \mathbb{R} \to \varphi_{291}$
\item $\varphi_{484} :\equiv (a \in \mathbb{R} \to \varphi_{152}) \to (\mathrm{ContOn}(f,[a,b]) \to a \in \mathbb{R} \to a \in \mathbb{R}) \to \varphi_{435}$
\item $\varphi_{485} :\equiv (\mathrm{ContOn}(f,[a,b]) \to a \in \mathbb{R} \to \varphi_{388}) \to \mathrm{ContOn}(f,[a,b]) \to a \in \mathbb{R} \to \varphi_{363}$
\item $\varphi_{486} :\equiv (\varphi_{303} \to \forall n\,(\varphi_{303})) \to \mathrm{ContOn}(f,[a,b]) \to a \in \mathbb{R} \to \varphi_{309}$
\item $\varphi_{487} :\equiv (\varphi_{303} \to \forall n\,(\varphi_{303})) \to \mathrm{ContOn}(f,[a,b]) \to a \in \mathbb{R} \to \varphi_{312}$
\item $\varphi_{488} :\equiv (\mathrm{ContOn}(f,[a,b]) \to a \in \mathbb{R} \to \varphi_{388}) \to \mathrm{ContOn}(f,[a,b]) \to a \in \mathbb{R} \to \varphi_{371}$
\item $\varphi_{489} :\equiv (\mathrm{ContOn}(f,[a,b]) \to a \in \mathbb{R} \to \varphi_{366}) \to \mathrm{ContOn}(f,[a,b]) \to a \in \mathbb{R} \to \varphi_{396}$
\item $\varphi_{490} :\equiv \varphi_{387} \to \varphi_{331} \to \mathrm{ContOn}(f,[a,b]) \to a \in \mathbb{R} \to b \in \mathbb{R} \to \varphi_{155}$
\item $\varphi_{491} :\equiv (\mathrm{ContOn}(f,[a,b]) \to a \in \mathbb{R} \to \varphi_{396}) \to \mathrm{ContOn}(f,[a,b]) \to a \in \mathbb{R} \to \varphi_{381}$
\item $\varphi_{492} :\equiv (\mathrm{NSeq}(s) \to n \in \mathbb{N} \to s_{n} \in \mathbb{N}) \to (\mathrm{ContOn}(f,[a,b]) \to a \in \mathbb{R} \to \varphi_{292}) \to \varphi_{411}$
\item $\varphi_{493} :\equiv (\mathrm{ContOn}(f,[a,b]) \to a \in \mathbb{R} \to \varphi_{377}) \to \varphi_{104} \to \mathrm{ContOn}(f,[a,b]) \to a \in \mathbb{R} \to \varphi_{374}$
\item $\varphi_{494} :\equiv (\mathrm{ContOn}(f,[a,b]) \to a \in \mathbb{R} \to \varphi_{379}) \to \varphi_{117} \to \mathrm{ContOn}(f,[a,b]) \to a \in \mathbb{R} \to \varphi_{375}$
\item $\varphi_{495} :\equiv (\mathrm{ContOn}(f,[a,b]) \to a \in \mathbb{R} \to \varphi_{367}) \to f(q \circ s(K)) = f(q \circ s(K)) \to \varphi_{405}$
\item $\varphi_{496} :\equiv (\mathrm{ContOn}(f,[a,b]) \to a \in \mathbb{R} \to \varphi_{386}) \to \varphi_{187} \to \mathrm{ContOn}(f,[a,b]) \to a \in \mathbb{R} \to \varphi_{385}$
\item $\varphi_{497} :\equiv (\mathrm{ContOn}(f,[a,b]) \to a \in \mathbb{R} \to \varphi_{211}) \to \forall q\,(\varphi_{113} \land \neg 0 < 0 \to \neg \varphi_{47}) \to \varphi_{365} \to \varphi_{324}$
\item $\varphi_{498} :\equiv (\mathrm{ContOn}(f,[a,b]) \to a \in \mathbb{R} \to \varphi_{369}) \to \varphi_{170} \to \varphi_{159} \to \mathrm{ContOn}(f,[a,b]) \to a \in \mathbb{R} \to \varphi_{362}$
\item $\varphi_{499} :\equiv (\mathrm{NSeq}(s) \to K \in \mathbb{N} \to s_{K} \in \mathbb{N}) \to (\mathrm{ContOn}(f,[a,b]) \to a \in \mathbb{R} \to \varphi_{292}) \to \varphi_{459}$
\item $\varphi_{500} :\equiv (0 \leq \lvert M \rvert \to \lvert M \rvert < \lvert M \rvert + 1 \to 0 < \lvert M \rvert + 1) \to (\mathrm{ContOn}(f,[a,b]) \to a \in \mathbb{R} \to \varphi_{321}) \to \varphi_{445}$
\item $\varphi_{501} :\equiv (\mathrm{RFn}(f) \to \varphi_{225}) \to \varphi_{4} \to (\mathrm{ContOn}(f,[a,b]) \to a \in \mathbb{R} \to \varphi_{294}) \to \varphi_{437}$
\item $\varphi_{502} :\equiv (\mathrm{ContOn}(f,[a,b]) \to a \in \mathbb{R} \to \varphi_{294}) \to (\mathrm{ContOn}(f,[a,b]) \to a \in \mathbb{R} \to \varphi_{312}) \to \varphi_{412}$
\item $\varphi_{503} :\equiv (\mathrm{ContOn}(f,[a,b]) \to a \in \mathbb{R} \to b \in \mathbb{R} \to \varphi_{124}) \to \forall x\,(\varphi_{275}) \to \varphi_{490}$
\item $\varphi_{504} :\equiv (\lvert M \rvert < \lvert M \rvert + 1 \to \lvert M \rvert + 1 < \underline{K} \to \lvert M \rvert < \underline{K}) \to (\mathrm{ContOn}(f,[a,b]) \to a \in \mathbb{R} \to \varphi_{353}) \to \varphi_{466}$
\item $\varphi_{505} :\equiv (f(q(s_{K})) \leq M \to M \leq \lvert M \rvert \to f(q(s_{K})) \leq \lvert M \rvert) \to (\mathrm{ContOn}(f,[a,b]) \to a \in \mathbb{R} \to \varphi_{358}) \to \varphi_{465}$
\item $\varphi_{506} :\equiv (\mathrm{ContOn}(f,[a,b]) \to a \in \mathbb{R} \to \varphi_{351}) \to (\mathrm{ContOn}(f,[a,b]) \to a \in \mathbb{R} \to \varphi_{361}) \to \varphi_{399}$
\item $\varphi_{507} :\equiv (\mathrm{ContOn}(f,[a,b]) \to a \in \mathbb{R} \to \varphi_{383}) \to (\mathrm{ContOn}(f,[a,b]) \to a \in \mathbb{R} \to \varphi_{376}) \to \varphi_{409}$
\item $\varphi_{508} :\equiv (\mathrm{ContOn}(f,[a,b]) \to a \in \mathbb{R} \to \varphi_{328}) \to (\mathrm{ContOn}(f,[a,b]) \to a \in \mathbb{R} \to \varphi_{319}) \to \varphi_{449}$
\item $\varphi_{509} :\equiv (\mathrm{ContOn}(f,[a,b]) \to a \in \mathbb{R} \to \varphi_{266}) \to \forall s\,(\varphi_{113} \land \varphi_{47} \land \neg 0 < 0 \to \neg \varphi_{3}) \to \varphi_{400} \to \varphi_{404}$
\item $\varphi_{510} :\equiv (\mathrm{ContOn}(f,[a,b]) \to a \in \mathbb{R} \to \varphi_{342}) \to (\mathrm{ContOn}(f,[a,b]) \to a \in \mathbb{R} \to \varphi_{346}) \to \varphi_{461}$
\item $\varphi_{511} :\equiv (\mathrm{ContIn}(f,[a,b],c) \to \mathrm{Seq}(q \circ s) \to \varphi_{126}) \to (\mathrm{ContOn}(f,[a,b]) \to a \in \mathbb{R} \to \varphi_{299}) \to \varphi_{502}$
\item $\varphi_{512} :\equiv (\mathrm{ContOn}(f,[a,b]) \to a \in \mathbb{R} \to \varphi_{288}) \to \forall c\,(\varphi_{113} \land \varphi_{47} \land \varphi_{3} \land \neg 0 < 0 \to \varphi_{2}) \to \varphi_{408} \to \varphi_{443}$
\item $\varphi_{513} :\equiv (\mathrm{ContOn}(f,[a,b]) \to a \in \mathbb{R} \to \varphi_{318}) \to \forall M\,(\varphi_{338}) \to (\forall M\,(\varphi_{338}) \to \forall M\,(\varphi_{303} \land \neg 0 < 0) \to \varphi_{23}) \to \varphi_{483}$
\item $\varphi_{514} :\equiv \forall K\,(\varphi_{336} \to \neg (K \in \mathbb{N} \land \lvert M \rvert + 1 < \underline{K} \cdot 1)) \to \varphi_{458} \to (\varphi_{336} \to \forall K\,(\varphi_{336})) \to \forall n\,(\varphi_{74}) \to \varphi_{420}$
\end{abbrevlist}
\end{fullproof}

\input{gen/EVT}
% 自動生成 (prover/build.js): EVTMin  ノード数 498
\begin{fullproof}{定理2.17(最小値)}{$\forall f\,\forall a\,\forall b\,(\mathrm{ContOn}(f,[a,b]) \to a \in \mathbb{R} \to b \in \mathbb{R} \to a \leq b \to \exists \xi \in [a,b]\,\forall x \in [a,b]\,(f(\xi) \leq f(x)))$}{498}\phantomsection\label{tree:EVTMin}
\noindent$\Pi_{1}$:
\begin{ptree}
\AXC{$z = w \to f(\xi) = z \to f(\xi) = w$\lbl{A6}}
\UIC{$\forall w\,(z = w \to f(\xi) = z \to f(\xi) = w)$}
\UIC{$\forall z\,\forall w\,(z = w \to f(\xi) = z \to f(\xi) = w)$}
\AXC{$\forall z\,\forall w\,(z = w \to f(\xi) = z \to f(\xi) = w) \to \varphi_{30}$\lbl{A2}}
\BIC{$\varphi_{30}$}
\AXC{$\varphi_{30} \to \varphi_{51}$\lbl{A2}}
\BIC{$\varphi_{51}$}
\end{ptree}
\noindent$\Pi_{2}$:
\begin{ptree}
\AXC{$z = w \to f(\xi) < z \to f(\xi) < w$\lbl{A6}}
\UIC{$\forall w\,(z = w \to f(\xi) < z \to f(\xi) < w)$}
\UIC{$\forall z\,\forall w\,(z = w \to f(\xi) < z \to f(\xi) < w)$}
\AXC{$\forall z\,\forall w\,(z = w \to f(\xi) < z \to f(\xi) < w) \to \varphi_{42}$\lbl{A2}}
\BIC{$\varphi_{42}$}
\AXC{$\varphi_{42} \to \varphi_{61}$\lbl{A2}}
\BIC{$\varphi_{61}$}
\end{ptree}
\noindent$\Pi_{3}$:
\begin{ptree}
\AXC{$z = w \to z = -(-f(x)) \to w = -(-f(x))$\lbl{A6}}
\UIC{$\forall w\,(z = w \to z = -(-f(x)) \to w = -(-f(x)))$}
\UIC{$\varphi_{4}$}
\AXC{$\varphi_{4} \to \varphi_{71}$\lbl{A2}}
\BIC{$\varphi_{71}$}
\AXC{$\varphi_{71} \to -(-f(\xi)) = f(\xi) \to \varphi_{35}$\lbl{A2}}
\BIC{$-(-f(\xi)) = f(\xi) \to \varphi_{35}$}
\end{ptree}
\noindent$\Pi_{4}$:
\begin{ptree}
\AXC{$z = w \to z < -(-f(x)) \to w < -(-f(x))$\lbl{A6}}
\UIC{$\forall w\,(z = w \to z < -(-f(x)) \to w < -(-f(x)))$}
\UIC{$\varphi_{16}$}
\AXC{$\varphi_{16} \to \varphi_{85}$\lbl{A2}}
\BIC{$\varphi_{85}$}
\AXC{$\varphi_{85} \to -(-f(\xi)) = f(\xi) \to \varphi_{46}$\lbl{A2}}
\BIC{$-(-f(\xi)) = f(\xi) \to \varphi_{46}$}
\end{ptree}
\noindent$\Pi_{5}$:
\begin{ptree}
\AXC{$z = w \to -f(x) = z \to -f(x) = w$\lbl{A6}}
\UIC{$\forall w\,(z = w \to -f(x) = z \to -f(x) = w)$}
\UIC{$\forall z\,\forall w\,(z = w \to -f(x) = z \to -f(x) = w)$}
\AXC{$\forall z\,\forall w\,(z = w \to -f(x) = z \to -f(x) = w) \to \varphi_{65}$\lbl{A2}}
\BIC{$\varphi_{65}$}
\AXC{$\varphi_{65} \to \varphi_{100}$\lbl{A2}}
\BIC{$\varphi_{100}$}
\end{ptree}
\noindent$\Pi_{6}$:
\begin{ptree}
\AXC{$z = w \to -f(x) < z \to -f(x) < w$\lbl{A6}}
\UIC{$\forall w\,(z = w \to -f(x) < z \to -f(x) < w)$}
\UIC{$\forall z\,\forall w\,(z = w \to -f(x) < z \to -f(x) < w)$}
\AXC{$\forall z\,\forall w\,(z = w \to -f(x) < z \to -f(x) < w) \to \varphi_{83}$\lbl{A2}}
\BIC{$\varphi_{83}$}
\AXC{$\varphi_{83} \to \varphi_{108}$\lbl{A2}}
\BIC{$\varphi_{108}$}
\end{ptree}
\noindent$\Pi_{7}$:
\begin{ptree}
\AXC{$z = w \to z = (-1 f)(\xi) \to w = (-1 f)(\xi)$\lbl{A6}}
\UIC{$\forall w\,(z = w \to z = (-1 f)(\xi) \to w = (-1 f)(\xi))$}
\UIC{$\varphi_{24}$}
\AXC{$\varphi_{24} \to \varphi_{94}$\lbl{A2}}
\BIC{$\varphi_{94}$}
\AXC{$\varphi_{94} \to (-1 f)(x) = -f(x) \to \varphi_{67}$\lbl{A2}}
\BIC{$(-1 f)(x) = -f(x) \to \varphi_{67}$}
\end{ptree}
\noindent$\Pi_{8}$:
\begin{ptree}
\AXC{$z = w \to z < (-1 f)(\xi) \to w < (-1 f)(\xi)$\lbl{A6}}
\UIC{$\forall w\,(z = w \to z < (-1 f)(\xi) \to w < (-1 f)(\xi))$}
\UIC{$\varphi_{37}$}
\AXC{$\varphi_{37} \to \varphi_{101}$\lbl{A2}}
\BIC{$\varphi_{101}$}
\AXC{$\varphi_{101} \to (-1 f)(x) = -f(x) \to \varphi_{84}$\lbl{A2}}
\BIC{$(-1 f)(x) = -f(x) \to \varphi_{84}$}
\end{ptree}
\noindent$\Pi_{9}$:
\begin{ptree}
\AXC{$\mathrm{ContOn}(f,[a,b]) \to \varphi_{136}$\lbl{A1}}
\AXC{$\varphi_{120}$\lbl{A2}}
\AXC{$\varphi_{170}$\lbl{A1}}
\AXC{$\varphi_{170} \to \varphi_{169}$\lbl{A1}}
\BIC{$\varphi_{169}$}
\AXC{$\varphi_{169} \to \varphi_{120} \to \varphi_{168}$\lbl{A1}}
\BIC{$\varphi_{120} \to \varphi_{168}$}
\BIC{$\varphi_{168}$}
\AXC{$\varphi_{168} \to \varphi_{201}$\lbl{A1}}
\BIC{$\varphi_{201}$}
\BIC{$\mathrm{ContOn}(f,[a,b]) \to a \in \mathbb{R} \to b \in \mathbb{R} \to \varphi_{149}$}
\end{ptree}
\noindent$\Pi_{10}$:
\begin{ptree}
\AXC{$y = z \to x = y \to x = z$\lbl{A6}}
\AXC{$(y = z \to x = y \to x = z) \to x = y \to y = z \to x = z$\lbl{A1}}
\BIC{$x = y \to y = z \to x = z$}
\UIC{$\forall z\,(x = y \to y = z \to x = z)$}
\UIC{$\forall y\,\forall z\,(x = y \to y = z \to x = z)$}
\UIC{$\forall x\,\forall y\,\forall z\,(x = y \to y = z \to x = z)$}
\AXC{$\forall x\,\forall y\,\forall z\,(x = y \to y = z \to x = z) \to \varphi_{19}$\lbl{A2}}
\BIC{$\varphi_{19}$}
\end{ptree}
\noindent$\Pi_{11}$:
\begin{ptree}
\AXC{$z = w \to -(1 \cdot f(x)) = -z \to -(1 \cdot f(x)) = -w$\lbl{A6}}
\UIC{$\varphi_{17}$}
\UIC{$\forall z\,(\varphi_{17})$}
\AXC{$\forall z\,(\varphi_{17}) \to \varphi_{66}$\lbl{A2}}
\BIC{$\varphi_{66}$}
\AXC{$\varphi_{66} \to 1 \cdot f(x) = f(x) \to \varphi_{40}$\lbl{A2}}
\BIC{$1 \cdot f(x) = f(x) \to \varphi_{40}$}
\end{ptree}
\noindent$\Pi_{12}$:
\begin{ptree}
\AXC{$y = z \to x = y \to x = z$\lbl{A6}}
\AXC{$(y = z \to x = y \to x = z) \to x = y \to y = z \to x = z$\lbl{A1}}
\BIC{$x = y \to y = z \to x = z$}
\UIC{$\forall z\,(x = y \to y = z \to x = z)$}
\UIC{$\forall y\,\forall z\,(x = y \to y = z \to x = z)$}
\UIC{$\forall x\,\forall y\,\forall z\,(x = y \to y = z \to x = z)$}
\AXC{$\varphi_{38}$\lbl{A2}}
\BIC{$\forall y\,\forall z\,(1 \cdot f(x) = y \to y = z \to 1 \cdot f(x) = z)$}
\end{ptree}
\noindent$\Pi_{13}$:
\begin{ptree}
\AXC{$\stackrel{\Pi_{12}}{\vdots}$}\noLine
\UIC{$\forall y\,\forall z\,(1 \cdot f(x) = y \to y = z \to 1 \cdot f(x) = z)$}
\AXC{$\forall y\,\forall z\,(1 \cdot f(x) = y \to y = z \to 1 \cdot f(x) = z) \to \varphi_{27}$\lbl{A2}}
\BIC{$\varphi_{27}$}
\AXC{$\varphi_{27} \to \varphi_{50}$\lbl{A2}}
\BIC{$\varphi_{50}$}
\end{ptree}
\noindent$\Pi_{14}$:
\begin{ptree}
\AXC{$\forall f\,(\varphi_{31})$\lbl{DefContOn}}
\AXC{$\forall f\,(\varphi_{31}) \to \varphi_{31}$\lbl{A2}}
\BIC{$\varphi_{31}$}
\AXC{$\varphi_{31} \to \varphi_{90}$\lbl{A2}}
\BIC{$\varphi_{90}$}
\AXC{$\varphi_{1}$\lbl{A1}}
\AXC{$\varphi_{1} \to \varphi_{90} \to \varphi_{45}$\lbl{A1}}
\BIC{$\varphi_{90} \to \varphi_{45}$}
\BIC{$\varphi_{45}$}
\end{ptree}
\noindent$\Pi_{15}$:
\begin{ptree}
\AXC{$\forall X\,\forall Y\,\forall x\,(X \subseteq Y \to x \in X \to x \in Y)$\lbl{[B5]}}
\AXC{$\varphi_{76}$\lbl{A2}}
\BIC{$\forall Y\,\forall x\,([a,b] \subseteq Y \to x \in [a,b] \to x \in Y)$}
\AXC{$\forall Y\,\forall x\,([a,b] \subseteq Y \to x \in [a,b] \to x \in Y) \to \varphi_{6}$\lbl{A2}}
\BIC{$\varphi_{6}$}
\AXC{$\varphi_{105}$\lbl{A2}}
\BIC{$[a,b] \subseteq \mathrm{dom}(f) \to x \in [a,b] \to x \in \mathrm{dom}(f)$}
\end{ptree}
\noindent$\Pi_{16}$:
\begin{ptree}
\AXC{$\mathrm{ContOn}(f,[a,b]) \to \varphi_{136}$\lbl{A1}}
\AXC{$\stackrel{\Pi_{14}}{\vdots}$}\noLine
\UIC{$\varphi_{45}$}
\AXC{$\stackrel{\Pi_{15}}{\vdots}$}\noLine
\UIC{$[a,b] \subseteq \mathrm{dom}(f) \to x \in [a,b] \to x \in \mathrm{dom}(f)$}
\AXC{$\varphi_{216}$\lbl{A1}}
\BIC{$\varphi_{45} \to \varphi_{183}$}
\BIC{$\varphi_{183}$}
\BIC{$\mathrm{ContOn}(f,[a,b]) \to \varphi_{128}$}
\end{ptree}
\noindent$\Pi_{17}$:
\begin{ptree}
\AXC{$\forall f\,(\varphi_{31})$\lbl{DefContOn}}
\AXC{$\forall f\,(\varphi_{31}) \to \varphi_{31}$\lbl{A2}}
\BIC{$\varphi_{31}$}
\AXC{$\varphi_{31} \to \varphi_{90}$\lbl{A2}}
\BIC{$\varphi_{90}$}
\AXC{$\varphi_{1}$\lbl{A1}}
\AXC{$\varphi_{1} \to \varphi_{90} \to \varphi_{12}$\lbl{A1}}
\BIC{$\varphi_{90} \to \varphi_{12}$}
\BIC{$\varphi_{12}$}
\end{ptree}
\noindent$\Pi_{18}$:
\begin{ptree}
\AXC{$\forall f\,\forall x\,(\mathrm{RFn}(f) \to x \in \mathrm{dom}(f) \to f(x) \in \mathbb{R})$\lbl{[S6]}}
\AXC{$\varphi_{87}$\lbl{A2}}
\BIC{$\forall x\,(\mathrm{RFn}(f) \to x \in \mathrm{dom}(f) \to f(x) \in \mathbb{R})$}
\AXC{$\varphi_{72}$\lbl{A2}}
\BIC{$\mathrm{RFn}(f) \to x \in \mathrm{dom}(f) \to f(x) \in \mathbb{R}$}
\end{ptree}
\noindent$\Pi_{19}$:
\begin{ptree}
\AXC{$\stackrel{\Pi_{17}}{\vdots}$}\noLine
\UIC{$\varphi_{12}$}
\AXC{$\stackrel{\Pi_{18}}{\vdots}$}\noLine
\UIC{$\mathrm{RFn}(f) \to x \in \mathrm{dom}(f) \to f(x) \in \mathbb{R}$}
\AXC{$(\mathrm{RFn}(f) \to x \in \mathrm{dom}(f) \to f(x) \in \mathbb{R}) \to \varphi_{12} \to \varphi_{184}$\lbl{A1}}
\BIC{$\varphi_{12} \to \varphi_{184}$}
\BIC{$\varphi_{184}$}
\end{ptree}
\noindent$\Pi_{20}$:
\begin{ptree}
\AXC{$\stackrel{\Pi_{16}}{\vdots}$}\noLine
\UIC{$\mathrm{ContOn}(f,[a,b]) \to \varphi_{128}$}
\AXC{$\stackrel{\Pi_{19}}{\vdots}$}\noLine
\UIC{$\varphi_{184}$}
\BIC{$\mathrm{ContOn}(f,[a,b]) \to \varphi_{141}$}
\end{ptree}
\noindent$\Pi_{21}$:
\begin{ptree}
\AXC{$\stackrel{\Pi_{20}}{\vdots}$}\noLine
\UIC{$\mathrm{ContOn}(f,[a,b]) \to \varphi_{141}$}
\AXC{$\forall x \in \mathbb{R}\,(x \cdot 1 = x)$\lbl{(R8)}}
\AXC{$\varphi_{33}$\lbl{A2}}
\BIC{$f(x) \in \mathbb{R} \to f(x) \cdot 1 = f(x)$}
\AXC{$(f(x) \in \mathbb{R} \to f(x) \cdot 1 = f(x)) \to \varphi_{188}$\lbl{A1}}
\BIC{$\varphi_{188}$}
\BIC{$\mathrm{ContOn}(f,[a,b]) \to \varphi_{152}$}
\end{ptree}
\noindent$\Pi_{22}$:
\begin{ptree}
\AXC{$1 \in \mathbb{R}$\lbl{(R1)}}
\AXC{$\forall x\,\forall y\,(x \in \mathbb{R} \to y \in \mathbb{R} \to x \cdot y = y \cdot x)$\lbl{(R7)}}
\AXC{$\varphi_{77}$\lbl{A2}}
\BIC{$\forall y\,(1 \in \mathbb{R} \to y \in \mathbb{R} \to 1 \cdot y = y \cdot 1)$}
\AXC{$\varphi_{95}$\lbl{A2}}
\BIC{$1 \in \mathbb{R} \to f(x) \in \mathbb{R} \to 1 \cdot f(x) = f(x) \cdot 1$}
\AXC{$\varphi_{209}$\lbl{A1}}
\BIC{$1 \in \mathbb{R} \to \varphi_{192}$}
\BIC{$\varphi_{192}$}
\end{ptree}
\noindent$\Pi_{23}$:
\begin{ptree}
\AXC{$\stackrel{\Pi_{20}}{\vdots}$}\noLine
\UIC{$\mathrm{ContOn}(f,[a,b]) \to \varphi_{141}$}
\AXC{$\stackrel{\Pi_{22}}{\vdots}$}\noLine
\UIC{$\varphi_{192}$}
\BIC{$\mathrm{ContOn}(f,[a,b]) \to a \in \mathbb{R} \to \varphi_{133}$}
\AXC{$(\mathrm{ContOn}(f,[a,b]) \to a \in \mathbb{R} \to \varphi_{133}) \to \varphi_{210}$\lbl{A1}}
\BIC{$\varphi_{210}$}
\end{ptree}
\noindent$\Pi_{24}$:
\begin{ptree}
\AXC{$\stackrel{\Pi_{13}}{\vdots}$}\noLine
\UIC{$\varphi_{50}$}
\AXC{$\stackrel{\Pi_{21}}{\vdots}$}\noLine
\UIC{$\mathrm{ContOn}(f,[a,b]) \to \varphi_{152}$}
\AXC{$\stackrel{\Pi_{23}}{\vdots}$}\noLine
\UIC{$\varphi_{210}$}
\BIC{$\varphi_{50} \to \mathrm{ContOn}(f,[a,b]) \to \varphi_{153}$}
\BIC{$\mathrm{ContOn}(f,[a,b]) \to \varphi_{153}$}
\AXC{$\varphi_{218}$\lbl{A1}}
\BIC{$-(1 \cdot f(x)) = -(1 \cdot f(x)) \to \varphi_{175}$}
\end{ptree}
\noindent$\Pi_{25}$:
\begin{ptree}
\AXC{$x = x$\lbl{A5}}
\UIC{$\forall x\,(x = x)$}
\AXC{$\forall x\,(x = x) \to -(1 \cdot f(x)) = -(1 \cdot f(x))$\lbl{A2}}
\BIC{$-(1 \cdot f(x)) = -(1 \cdot f(x))$}
\AXC{$\stackrel{\Pi_{24}}{\vdots}$}\noLine
\UIC{$-(1 \cdot f(x)) = -(1 \cdot f(x)) \to \varphi_{175}$}
\BIC{$\varphi_{175}$}
\end{ptree}
\noindent$\Pi_{26}$:
\begin{ptree}
\AXC{$y = z \to x = y \to x = z$\lbl{A6}}
\AXC{$(y = z \to x = y \to x = z) \to x = y \to y = z \to x = z$\lbl{A1}}
\BIC{$x = y \to y = z \to x = z$}
\UIC{$\forall z\,(x = y \to y = z \to x = z)$}
\UIC{$\forall y\,\forall z\,(x = y \to y = z \to x = z)$}
\UIC{$\forall x\,\forall y\,\forall z\,(x = y \to y = z \to x = z)$}
\AXC{$\forall x\,\forall y\,\forall z\,(x = y \to y = z \to x = z) \to \varphi_{19}$\lbl{A2}}
\BIC{$\varphi_{19}$}
\end{ptree}
\noindent$\Pi_{27}$:
\begin{ptree}
\AXC{$\stackrel{\Pi_{26}}{\vdots}$}\noLine
\UIC{$\varphi_{19}$}
\AXC{$\varphi_{19} \to \varphi_{88}$\lbl{A2}}
\BIC{$\varphi_{88}$}
\AXC{$\varphi_{88} \to (-1 f)(x) = -1 \cdot f(x) \to \varphi_{54}$\lbl{A2}}
\BIC{$(-1 f)(x) = -1 \cdot f(x) \to \varphi_{54}$}
\end{ptree}
\noindent$\Pi_{28}$:
\begin{ptree}
\AXC{$1 \in \mathbb{R}$\lbl{(R1)}}
\AXC{$\forall x\,\forall y\,(x \in \mathbb{R} \to y \in \mathbb{R} \to -x \cdot y = -(x \cdot y))$\lbl{[F9]}}
\AXC{$\varphi_{102}$\lbl{A2}}
\BIC{$\forall y\,(1 \in \mathbb{R} \to y \in \mathbb{R} \to -1 \cdot y = -(1 \cdot y))$}
\AXC{$\varphi_{116}$\lbl{A2}}
\BIC{$1 \in \mathbb{R} \to f(x) \in \mathbb{R} \to -1 \cdot f(x) = -(1 \cdot f(x))$}
\AXC{$\varphi_{213}$\lbl{A1}}
\BIC{$1 \in \mathbb{R} \to \varphi_{198}$}
\BIC{$\varphi_{198}$}
\end{ptree}
\noindent$\Pi_{29}$:
\begin{ptree}
\AXC{$\forall x \in \mathbb{R}\,(-x \in \mathbb{R})$\lbl{(R1)}}
\AXC{$\forall x \in \mathbb{R}\,(-x \in \mathbb{R}) \to 1 \in \mathbb{R} \to -1 \in \mathbb{R}$\lbl{A2}}
\BIC{$1 \in \mathbb{R} \to -1 \in \mathbb{R}$}
\AXC{$(1 \in \mathbb{R} \to -1 \in \mathbb{R}) \to 1 \in \mathbb{R} \to \varphi_{34}$\lbl{A1}}
\BIC{$1 \in \mathbb{R} \to \varphi_{34}$}
\end{ptree}
\noindent$\Pi_{30}$:
\begin{ptree}
\AXC{$1 \in \mathbb{R}$\lbl{(R1)}}
\AXC{$\stackrel{\Pi_{29}}{\vdots}$}\noLine
\UIC{$1 \in \mathbb{R} \to \varphi_{34}$}
\BIC{$\varphi_{34}$}
\end{ptree}
\noindent$\Pi_{31}$:
\begin{ptree}
\AXC{$\stackrel{\Pi_{30}}{\vdots}$}\noLine
\UIC{$\varphi_{34}$}
\AXC{$\varphi_{124}$\lbl{[S19]}}
\AXC{$\varphi_{124} \to \forall f\,(-1 \in \mathbb{R} \to \mathrm{RFn}(f) \to \varphi_{122})$\lbl{A2}}
\BIC{$\forall f\,(-1 \in \mathbb{R} \to \mathrm{RFn}(f) \to \varphi_{122})$}
\AXC{$\varphi_{182}$\lbl{A2}}
\BIC{$-1 \in \mathbb{R} \to \mathrm{RFn}(f) \to \varphi_{122}$}
\AXC{$(-1 \in \mathbb{R} \to \mathrm{RFn}(f) \to \varphi_{122}) \to \varphi_{34} \to \varphi_{12} \to \varphi_{165}$\lbl{A1}}
\BIC{$\varphi_{34} \to \varphi_{12} \to \varphi_{165}$}
\BIC{$\varphi_{12} \to \varphi_{165}$}
\end{ptree}
\noindent$\Pi_{32}$:
\begin{ptree}
\AXC{$\stackrel{\Pi_{17}}{\vdots}$}\noLine
\UIC{$\varphi_{12}$}
\AXC{$\stackrel{\Pi_{31}}{\vdots}$}\noLine
\UIC{$\varphi_{12} \to \varphi_{165}$}
\BIC{$\varphi_{165}$}
\end{ptree}
\noindent$\Pi_{33}$:
\begin{ptree}
\AXC{$\stackrel{\Pi_{32}}{\vdots}$}\noLine
\UIC{$\varphi_{165}$}
\AXC{$\varphi_{165} \to \mathrm{ContOn}(f,[a,b]) \to a \in \mathbb{R} \to \varphi_{22}$\lbl{A1}}
\BIC{$\mathrm{ContOn}(f,[a,b]) \to a \in \mathbb{R} \to \varphi_{22}$}
\end{ptree}
\noindent$\Pi_{34}$:
\begin{ptree}
\AXC{$\stackrel{\Pi_{32}}{\vdots}$}\noLine
\UIC{$\varphi_{165}$}
\AXC{$\varphi_{165} \to \varphi_{118}$\lbl{A1}}
\BIC{$\varphi_{118}$}
\end{ptree}
\noindent$\Pi_{35}$:
\begin{ptree}
\AXC{$\varphi_{125}$\lbl{A2}}
\AXC{$\stackrel{\Pi_{34}}{\vdots}$}\noLine
\UIC{$\varphi_{118}$}
\AXC{$\varphi_{204}$\lbl{A1}}
\BIC{$\varphi_{125} \to \mathrm{ContOn}(f,[a,b]) \to a \in \mathbb{R} \to b \in \mathbb{R} \to \varphi_{47}$}
\BIC{$\mathrm{ContOn}(f,[a,b]) \to a \in \mathbb{R} \to b \in \mathbb{R} \to \varphi_{47}$}
\AXC{$\varphi_{225}$\lbl{A1}}
\BIC{$(\mathrm{ContOn}(f,[a,b]) \to a \in \mathbb{R} \to \varphi_{13}) \to \varphi_{197}$}
\end{ptree}
\noindent$\Pi_{36}$:
\begin{ptree}
\AXC{$\varphi_{73}$\lbl{A2}}
\AXC{$\stackrel{\Pi_{33}}{\vdots}$}\noLine
\UIC{$\mathrm{ContOn}(f,[a,b]) \to a \in \mathbb{R} \to \varphi_{22}$}
\AXC{$\varphi_{172}$\lbl{A1}}
\BIC{$\varphi_{73} \to \mathrm{ContOn}(f,[a,b]) \to a \in \mathbb{R} \to \varphi_{13}$}
\BIC{$\mathrm{ContOn}(f,[a,b]) \to a \in \mathbb{R} \to \varphi_{13}$}
\AXC{$\stackrel{\Pi_{35}}{\vdots}$}\noLine
\UIC{$(\mathrm{ContOn}(f,[a,b]) \to a \in \mathbb{R} \to \varphi_{13}) \to \varphi_{197}$}
\BIC{$\varphi_{197}$}
\end{ptree}
\noindent$\Pi_{37}$:
\begin{ptree}
\AXC{$\stackrel{\Pi_{16}}{\vdots}$}\noLine
\UIC{$\mathrm{ContOn}(f,[a,b]) \to \varphi_{128}$}
\AXC{$\stackrel{\Pi_{36}}{\vdots}$}\noLine
\UIC{$\varphi_{197}$}
\BIC{$\mathrm{ContOn}(f,[a,b]) \to a \in \mathbb{R} \to \varphi_{161}$}
\AXC{$(\mathrm{ContOn}(f,[a,b]) \to a \in \mathbb{R} \to \varphi_{161}) \to \varphi_{220}$\lbl{A1}}
\BIC{$\varphi_{220}$}
\end{ptree}
\noindent$\Pi_{38}$:
\begin{ptree}
\AXC{$\stackrel{\Pi_{27}}{\vdots}$}\noLine
\UIC{$(-1 f)(x) = -1 \cdot f(x) \to \varphi_{54}$}
\AXC{$\stackrel{\Pi_{20}}{\vdots}$}\noLine
\UIC{$\mathrm{ContOn}(f,[a,b]) \to \varphi_{141}$}
\AXC{$\stackrel{\Pi_{28}}{\vdots}$}\noLine
\UIC{$\varphi_{198}$}
\BIC{$\mathrm{ContOn}(f,[a,b]) \to a \in \mathbb{R} \to b \in \mathbb{R} \to \varphi_{127}$}
\AXC{$\stackrel{\Pi_{37}}{\vdots}$}\noLine
\UIC{$\varphi_{220}$}
\BIC{$\varphi_{180}$}
\BIC{$\mathrm{ContOn}(f,[a,b]) \to a \in \mathbb{R} \to b \in \mathbb{R} \to \varphi_{137}$}
\AXC{$\varphi_{230}$\lbl{A1}}
\BIC{$\varphi_{217}$}
\end{ptree}
\noindent$\Pi_{39}$:
\begin{ptree}
\AXC{$\stackrel{\Pi_{10}}{\vdots}$}\noLine
\UIC{$\varphi_{19}$}
\AXC{$\varphi_{19} \to \varphi_{86}$\lbl{A2}}
\BIC{$\varphi_{86}$}
\AXC{$\varphi_{86} \to \varphi_{109}$\lbl{A2}}
\BIC{$\varphi_{109}$}
\AXC{$\stackrel{\Pi_{11}}{\vdots}$}\noLine
\UIC{$1 \cdot f(x) = f(x) \to \varphi_{40}$}
\AXC{$\stackrel{\Pi_{25}}{\vdots}$}\noLine
\UIC{$\varphi_{175}$}
\BIC{$\mathrm{ContOn}(f,[a,b]) \to a \in \mathbb{R} \to \varphi_{147}$}
\AXC{$\stackrel{\Pi_{38}}{\vdots}$}\noLine
\UIC{$\varphi_{217}$}
\BIC{$\varphi_{109} \to \mathrm{ContOn}(f,[a,b]) \to a \in \mathbb{R} \to \varphi_{155}$}
\BIC{$\mathrm{ContOn}(f,[a,b]) \to a \in \mathbb{R} \to \varphi_{155}$}
\end{ptree}
\noindent$\Pi_{40}$:
\begin{ptree}
\AXC{$\stackrel{\Pi_{39}}{\vdots}$}\noLine
\UIC{$\mathrm{ContOn}(f,[a,b]) \to a \in \mathbb{R} \to \varphi_{155}$}
\AXC{$(\mathrm{ContOn}(f,[a,b]) \to a \in \mathbb{R} \to \varphi_{155}) \to \varphi_{229}$\lbl{A1}}
\BIC{$\varphi_{229}$}
\end{ptree}
\noindent$\Pi_{41}$:
\begin{ptree}
\AXC{$\stackrel{\Pi_{8}}{\vdots}$}\noLine
\UIC{$(-1 f)(x) = -f(x) \to \varphi_{84}$}
\AXC{$\stackrel{\Pi_{9}}{\vdots}$}\noLine
\UIC{$\mathrm{ContOn}(f,[a,b]) \to a \in \mathbb{R} \to b \in \mathbb{R} \to \varphi_{149}$}
\AXC{$\stackrel{\Pi_{40}}{\vdots}$}\noLine
\UIC{$\varphi_{229}$}
\BIC{$((-1 f)(x) = -f(x) \to \varphi_{84}) \to \varphi_{179}$}
\BIC{$\varphi_{179}$}
\end{ptree}
\noindent$\Pi_{42}$:
\begin{ptree}
\AXC{$y = z \to x = y \to x = z$\lbl{A6}}
\AXC{$(y = z \to x = y \to x = z) \to x = y \to y = z \to x = z$\lbl{A1}}
\BIC{$x = y \to y = z \to x = z$}
\UIC{$\forall z\,(x = y \to y = z \to x = z)$}
\UIC{$\forall y\,\forall z\,(x = y \to y = z \to x = z)$}
\UIC{$\forall x\,\forall y\,\forall z\,(x = y \to y = z \to x = z)$}
\AXC{$\forall x\,\forall y\,\forall z\,(x = y \to y = z \to x = z) \to \varphi_{25}$\lbl{A2}}
\BIC{$\varphi_{25}$}
\end{ptree}
\noindent$\Pi_{43}$:
\begin{ptree}
\AXC{$z = w \to -(1 \cdot f(\xi)) = -z \to -(1 \cdot f(\xi)) = -w$\lbl{A6}}
\UIC{$\varphi_{23}$}
\UIC{$\forall z\,(\varphi_{23})$}
\AXC{$\forall z\,(\varphi_{23}) \to \varphi_{78}$\lbl{A2}}
\BIC{$\varphi_{78}$}
\AXC{$\varphi_{78} \to 1 \cdot f(\xi) = f(\xi) \to \varphi_{48}$\lbl{A2}}
\BIC{$1 \cdot f(\xi) = f(\xi) \to \varphi_{48}$}
\end{ptree}
\noindent$\Pi_{44}$:
\begin{ptree}
\AXC{$y = z \to x = y \to x = z$\lbl{A6}}
\AXC{$(y = z \to x = y \to x = z) \to x = y \to y = z \to x = z$\lbl{A1}}
\BIC{$x = y \to y = z \to x = z$}
\UIC{$\forall z\,(x = y \to y = z \to x = z)$}
\UIC{$\forall y\,\forall z\,(x = y \to y = z \to x = z)$}
\UIC{$\forall x\,\forall y\,\forall z\,(x = y \to y = z \to x = z)$}
\AXC{$\varphi_{43}$\lbl{A2}}
\BIC{$\forall y\,\forall z\,(1 \cdot f(\xi) = y \to y = z \to 1 \cdot f(\xi) = z)$}
\end{ptree}
\noindent$\Pi_{45}$:
\begin{ptree}
\AXC{$\stackrel{\Pi_{44}}{\vdots}$}\noLine
\UIC{$\forall y\,\forall z\,(1 \cdot f(\xi) = y \to y = z \to 1 \cdot f(\xi) = z)$}
\AXC{$\forall y\,\forall z\,(1 \cdot f(\xi) = y \to y = z \to 1 \cdot f(\xi) = z) \to \varphi_{36}$\lbl{A2}}
\BIC{$\varphi_{36}$}
\AXC{$\varphi_{36} \to \varphi_{57}$\lbl{A2}}
\BIC{$\varphi_{57}$}
\end{ptree}
\noindent$\Pi_{46}$:
\begin{ptree}
\AXC{$\forall X\,\forall Y\,\forall x\,(X \subseteq Y \to x \in X \to x \in Y)$\lbl{[B5]}}
\AXC{$\varphi_{76}$\lbl{A2}}
\BIC{$\forall Y\,\forall x\,([a,b] \subseteq Y \to x \in [a,b] \to x \in Y)$}
\AXC{$\forall Y\,\forall x\,([a,b] \subseteq Y \to x \in [a,b] \to x \in Y) \to \varphi_{6}$\lbl{A2}}
\BIC{$\varphi_{6}$}
\AXC{$\varphi_{110}$\lbl{A2}}
\BIC{$[a,b] \subseteq \mathrm{dom}(f) \to \xi \in [a,b] \to \xi \in \mathrm{dom}(f)$}
\end{ptree}
\noindent$\Pi_{47}$:
\begin{ptree}
\AXC{$\varphi_{170}$\lbl{A1}}
\AXC{$\varphi_{170} \to \varphi_{148}$\lbl{A1}}
\BIC{$\varphi_{148}$}
\AXC{$\stackrel{\Pi_{14}}{\vdots}$}\noLine
\UIC{$\varphi_{45}$}
\AXC{$\stackrel{\Pi_{46}}{\vdots}$}\noLine
\UIC{$[a,b] \subseteq \mathrm{dom}(f) \to \xi \in [a,b] \to \xi \in \mathrm{dom}(f)$}
\AXC{$\varphi_{215}$\lbl{A1}}
\BIC{$\varphi_{45} \to \varphi_{148} \to \mathrm{ContOn}(f,[a,b]) \to \varphi_{130}$}
\BIC{$\varphi_{148} \to \mathrm{ContOn}(f,[a,b]) \to \varphi_{130}$}
\BIC{$\mathrm{ContOn}(f,[a,b]) \to \varphi_{130}$}
\end{ptree}
\noindent$\Pi_{48}$:
\begin{ptree}
\AXC{$\forall f\,\forall x\,(\mathrm{RFn}(f) \to x \in \mathrm{dom}(f) \to f(x) \in \mathbb{R})$\lbl{[S6]}}
\AXC{$\varphi_{87}$\lbl{A2}}
\BIC{$\forall x\,(\mathrm{RFn}(f) \to x \in \mathrm{dom}(f) \to f(x) \in \mathbb{R})$}
\AXC{$\varphi_{79}$\lbl{A2}}
\BIC{$\mathrm{RFn}(f) \to \xi \in \mathrm{dom}(f) \to f(\xi) \in \mathbb{R}$}
\end{ptree}
\noindent$\Pi_{49}$:
\begin{ptree}
\AXC{$\stackrel{\Pi_{17}}{\vdots}$}\noLine
\UIC{$\varphi_{12}$}
\AXC{$\stackrel{\Pi_{48}}{\vdots}$}\noLine
\UIC{$\mathrm{RFn}(f) \to \xi \in \mathrm{dom}(f) \to f(\xi) \in \mathbb{R}$}
\AXC{$(\mathrm{RFn}(f) \to \xi \in \mathrm{dom}(f) \to f(\xi) \in \mathbb{R}) \to \varphi_{12} \to \varphi_{185}$\lbl{A1}}
\BIC{$\varphi_{12} \to \varphi_{185}$}
\BIC{$\varphi_{185}$}
\end{ptree}
\noindent$\Pi_{50}$:
\begin{ptree}
\AXC{$\stackrel{\Pi_{47}}{\vdots}$}\noLine
\UIC{$\mathrm{ContOn}(f,[a,b]) \to \varphi_{130}$}
\AXC{$\stackrel{\Pi_{49}}{\vdots}$}\noLine
\UIC{$\varphi_{185}$}
\BIC{$\mathrm{ContOn}(f,[a,b]) \to \varphi_{143}$}
\end{ptree}
\noindent$\Pi_{51}$:
\begin{ptree}
\AXC{$\stackrel{\Pi_{50}}{\vdots}$}\noLine
\UIC{$\mathrm{ContOn}(f,[a,b]) \to \varphi_{143}$}
\AXC{$\forall x \in \mathbb{R}\,(x \cdot 1 = x)$\lbl{(R8)}}
\AXC{$\varphi_{39}$\lbl{A2}}
\BIC{$f(\xi) \in \mathbb{R} \to f(\xi) \cdot 1 = f(\xi)$}
\AXC{$(f(\xi) \in \mathbb{R} \to f(\xi) \cdot 1 = f(\xi)) \to \varphi_{191}$\lbl{A1}}
\BIC{$\varphi_{191}$}
\BIC{$\mathrm{ContOn}(f,[a,b]) \to \varphi_{158}$}
\end{ptree}
\noindent$\Pi_{52}$:
\begin{ptree}
\AXC{$1 \in \mathbb{R}$\lbl{(R1)}}
\AXC{$\forall x\,\forall y\,(x \in \mathbb{R} \to y \in \mathbb{R} \to x \cdot y = y \cdot x)$\lbl{(R7)}}
\AXC{$\varphi_{77}$\lbl{A2}}
\BIC{$\forall y\,(1 \in \mathbb{R} \to y \in \mathbb{R} \to 1 \cdot y = y \cdot 1)$}
\AXC{$\varphi_{97}$\lbl{A2}}
\BIC{$1 \in \mathbb{R} \to f(\xi) \in \mathbb{R} \to 1 \cdot f(\xi) = f(\xi) \cdot 1$}
\AXC{$\varphi_{211}$\lbl{A1}}
\BIC{$1 \in \mathbb{R} \to \varphi_{195}$}
\BIC{$\varphi_{195}$}
\end{ptree}
\noindent$\Pi_{53}$:
\begin{ptree}
\AXC{$\stackrel{\Pi_{50}}{\vdots}$}\noLine
\UIC{$\mathrm{ContOn}(f,[a,b]) \to \varphi_{143}$}
\AXC{$\stackrel{\Pi_{52}}{\vdots}$}\noLine
\UIC{$\varphi_{195}$}
\BIC{$\mathrm{ContOn}(f,[a,b]) \to a \in \mathbb{R} \to \varphi_{139}$}
\AXC{$(\mathrm{ContOn}(f,[a,b]) \to a \in \mathbb{R} \to \varphi_{139}) \to \varphi_{212}$\lbl{A1}}
\BIC{$\varphi_{212}$}
\end{ptree}
\noindent$\Pi_{54}$:
\begin{ptree}
\AXC{$\stackrel{\Pi_{45}}{\vdots}$}\noLine
\UIC{$\varphi_{57}$}
\AXC{$\stackrel{\Pi_{51}}{\vdots}$}\noLine
\UIC{$\mathrm{ContOn}(f,[a,b]) \to \varphi_{158}$}
\AXC{$\stackrel{\Pi_{53}}{\vdots}$}\noLine
\UIC{$\varphi_{212}$}
\BIC{$\varphi_{57} \to \mathrm{ContOn}(f,[a,b]) \to \varphi_{159}$}
\BIC{$\mathrm{ContOn}(f,[a,b]) \to \varphi_{159}$}
\AXC{$\varphi_{221}$\lbl{A1}}
\BIC{$-(1 \cdot f(\xi)) = -(1 \cdot f(\xi)) \to \varphi_{177}$}
\end{ptree}
\noindent$\Pi_{55}$:
\begin{ptree}
\AXC{$x = x$\lbl{A5}}
\UIC{$\forall x\,(x = x)$}
\AXC{$\forall x\,(x = x) \to -(1 \cdot f(\xi)) = -(1 \cdot f(\xi))$\lbl{A2}}
\BIC{$-(1 \cdot f(\xi)) = -(1 \cdot f(\xi))$}
\AXC{$\stackrel{\Pi_{54}}{\vdots}$}\noLine
\UIC{$-(1 \cdot f(\xi)) = -(1 \cdot f(\xi)) \to \varphi_{177}$}
\BIC{$\varphi_{177}$}
\end{ptree}
\noindent$\Pi_{56}$:
\begin{ptree}
\AXC{$y = z \to x = y \to x = z$\lbl{A6}}
\AXC{$(y = z \to x = y \to x = z) \to x = y \to y = z \to x = z$\lbl{A1}}
\BIC{$x = y \to y = z \to x = z$}
\UIC{$\forall z\,(x = y \to y = z \to x = z)$}
\UIC{$\forall y\,\forall z\,(x = y \to y = z \to x = z)$}
\UIC{$\forall x\,\forall y\,\forall z\,(x = y \to y = z \to x = z)$}
\AXC{$\forall x\,\forall y\,\forall z\,(x = y \to y = z \to x = z) \to \varphi_{25}$\lbl{A2}}
\BIC{$\varphi_{25}$}
\end{ptree}
\noindent$\Pi_{57}$:
\begin{ptree}
\AXC{$\stackrel{\Pi_{56}}{\vdots}$}\noLine
\UIC{$\varphi_{25}$}
\AXC{$\varphi_{25} \to \varphi_{92}$\lbl{A2}}
\BIC{$\varphi_{92}$}
\AXC{$\varphi_{92} \to (-1 f)(\xi) = -1 \cdot f(\xi) \to \varphi_{64}$\lbl{A2}}
\BIC{$(-1 f)(\xi) = -1 \cdot f(\xi) \to \varphi_{64}$}
\end{ptree}
\noindent$\Pi_{58}$:
\begin{ptree}
\AXC{$1 \in \mathbb{R}$\lbl{(R1)}}
\AXC{$\forall x\,\forall y\,(x \in \mathbb{R} \to y \in \mathbb{R} \to -x \cdot y = -(x \cdot y))$\lbl{[F9]}}
\AXC{$\varphi_{102}$\lbl{A2}}
\BIC{$\forall y\,(1 \in \mathbb{R} \to y \in \mathbb{R} \to -1 \cdot y = -(1 \cdot y))$}
\AXC{$\varphi_{119}$\lbl{A2}}
\BIC{$1 \in \mathbb{R} \to f(\xi) \in \mathbb{R} \to -1 \cdot f(\xi) = -(1 \cdot f(\xi))$}
\AXC{$\varphi_{214}$\lbl{A1}}
\BIC{$1 \in \mathbb{R} \to \varphi_{200}$}
\BIC{$\varphi_{200}$}
\end{ptree}
\noindent$\Pi_{59}$:
\begin{ptree}
\AXC{$\varphi_{126}$\lbl{A2}}
\AXC{$\stackrel{\Pi_{34}}{\vdots}$}\noLine
\UIC{$\varphi_{118}$}
\AXC{$\varphi_{206}$\lbl{A1}}
\BIC{$\varphi_{126} \to \mathrm{ContOn}(f,[a,b]) \to a \in \mathbb{R} \to b \in \mathbb{R} \to \varphi_{53}$}
\BIC{$\mathrm{ContOn}(f,[a,b]) \to a \in \mathbb{R} \to b \in \mathbb{R} \to \varphi_{53}$}
\AXC{$\varphi_{227}$\lbl{A1}}
\BIC{$(\mathrm{ContOn}(f,[a,b]) \to a \in \mathbb{R} \to \varphi_{18}) \to \varphi_{199}$}
\end{ptree}
\noindent$\Pi_{60}$:
\begin{ptree}
\AXC{$\varphi_{80}$\lbl{A2}}
\AXC{$\stackrel{\Pi_{33}}{\vdots}$}\noLine
\UIC{$\mathrm{ContOn}(f,[a,b]) \to a \in \mathbb{R} \to \varphi_{22}$}
\AXC{$\varphi_{174}$\lbl{A1}}
\BIC{$\varphi_{80} \to \mathrm{ContOn}(f,[a,b]) \to a \in \mathbb{R} \to \varphi_{18}$}
\BIC{$\mathrm{ContOn}(f,[a,b]) \to a \in \mathbb{R} \to \varphi_{18}$}
\AXC{$\stackrel{\Pi_{59}}{\vdots}$}\noLine
\UIC{$(\mathrm{ContOn}(f,[a,b]) \to a \in \mathbb{R} \to \varphi_{18}) \to \varphi_{199}$}
\BIC{$\varphi_{199}$}
\end{ptree}
\noindent$\Pi_{61}$:
\begin{ptree}
\AXC{$\stackrel{\Pi_{47}}{\vdots}$}\noLine
\UIC{$\mathrm{ContOn}(f,[a,b]) \to \varphi_{130}$}
\AXC{$\stackrel{\Pi_{60}}{\vdots}$}\noLine
\UIC{$\varphi_{199}$}
\BIC{$\mathrm{ContOn}(f,[a,b]) \to a \in \mathbb{R} \to \varphi_{163}$}
\AXC{$(\mathrm{ContOn}(f,[a,b]) \to a \in \mathbb{R} \to \varphi_{163}) \to \varphi_{223}$\lbl{A1}}
\BIC{$\varphi_{223}$}
\end{ptree}
\noindent$\Pi_{62}$:
\begin{ptree}
\AXC{$\stackrel{\Pi_{57}}{\vdots}$}\noLine
\UIC{$(-1 f)(\xi) = -1 \cdot f(\xi) \to \varphi_{64}$}
\AXC{$\stackrel{\Pi_{50}}{\vdots}$}\noLine
\UIC{$\mathrm{ContOn}(f,[a,b]) \to \varphi_{143}$}
\AXC{$\stackrel{\Pi_{58}}{\vdots}$}\noLine
\UIC{$\varphi_{200}$}
\BIC{$\mathrm{ContOn}(f,[a,b]) \to a \in \mathbb{R} \to b \in \mathbb{R} \to \varphi_{131}$}
\AXC{$\stackrel{\Pi_{61}}{\vdots}$}\noLine
\UIC{$\varphi_{223}$}
\BIC{$\varphi_{181}$}
\BIC{$\mathrm{ContOn}(f,[a,b]) \to a \in \mathbb{R} \to b \in \mathbb{R} \to \varphi_{142}$}
\AXC{$\varphi_{231}$\lbl{A1}}
\BIC{$\varphi_{219}$}
\end{ptree}
\noindent$\Pi_{63}$:
\begin{ptree}
\AXC{$\stackrel{\Pi_{42}}{\vdots}$}\noLine
\UIC{$\varphi_{25}$}
\AXC{$\varphi_{25} \to \varphi_{91}$\lbl{A2}}
\BIC{$\varphi_{91}$}
\AXC{$\varphi_{91} \to \varphi_{113}$\lbl{A2}}
\BIC{$\varphi_{113}$}
\AXC{$\stackrel{\Pi_{43}}{\vdots}$}\noLine
\UIC{$1 \cdot f(\xi) = f(\xi) \to \varphi_{48}$}
\AXC{$\stackrel{\Pi_{55}}{\vdots}$}\noLine
\UIC{$\varphi_{177}$}
\BIC{$\mathrm{ContOn}(f,[a,b]) \to a \in \mathbb{R} \to \varphi_{150}$}
\AXC{$\stackrel{\Pi_{62}}{\vdots}$}\noLine
\UIC{$\varphi_{219}$}
\BIC{$\varphi_{113} \to \mathrm{ContOn}(f,[a,b]) \to a \in \mathbb{R} \to \varphi_{160}$}
\BIC{$\mathrm{ContOn}(f,[a,b]) \to a \in \mathbb{R} \to \varphi_{160}$}
\end{ptree}
\noindent$\Pi_{64}$:
\begin{ptree}
\AXC{$\stackrel{\Pi_{63}}{\vdots}$}\noLine
\UIC{$\mathrm{ContOn}(f,[a,b]) \to a \in \mathbb{R} \to \varphi_{160}$}
\AXC{$(\mathrm{ContOn}(f,[a,b]) \to a \in \mathbb{R} \to \varphi_{160}) \to \varphi_{228}$\lbl{A1}}
\BIC{$\varphi_{228}$}
\end{ptree}
\noindent$\Pi_{65}$:
\begin{ptree}
\AXC{$\stackrel{\Pi_{5}}{\vdots}$}\noLine
\UIC{$\varphi_{100}$}
\AXC{$\stackrel{\Pi_{6}}{\vdots}$}\noLine
\UIC{$\varphi_{108}$}
\AXC{$\stackrel{\Pi_{7}}{\vdots}$}\noLine
\UIC{$(-1 f)(x) = -f(x) \to \varphi_{67}$}
\AXC{$\stackrel{\Pi_{41}}{\vdots}$}\noLine
\UIC{$\varphi_{179}$}
\BIC{$\mathrm{ContOn}(f,[a,b]) \to a \in \mathbb{R} \to \varphi_{157}$}
\AXC{$\stackrel{\Pi_{64}}{\vdots}$}\noLine
\UIC{$\varphi_{228}$}
\BIC{$\varphi_{108} \to \varphi_{100} \to \mathrm{ContOn}(f,[a,b]) \to \varphi_{162}$}
\BIC{$\varphi_{100} \to \mathrm{ContOn}(f,[a,b]) \to \varphi_{162}$}
\BIC{$\mathrm{ContOn}(f,[a,b]) \to \varphi_{162}$}
\end{ptree}
\noindent$\Pi_{66}$:
\begin{ptree}
\AXC{$\forall x \in \mathbb{R}\,(-x \in \mathbb{R})$\lbl{(R1)}}
\AXC{$\forall x \in \mathbb{R}\,(-x \in \mathbb{R}) \to f(\xi) \in \mathbb{R} \to -f(\xi) \in \mathbb{R}$\lbl{A2}}
\BIC{$f(\xi) \in \mathbb{R} \to -f(\xi) \in \mathbb{R}$}
\end{ptree}
\noindent$\Pi_{67}$:
\begin{ptree}
\AXC{$\stackrel{\Pi_{50}}{\vdots}$}\noLine
\UIC{$\mathrm{ContOn}(f,[a,b]) \to \varphi_{143}$}
\AXC{$\stackrel{\Pi_{66}}{\vdots}$}\noLine
\UIC{$f(\xi) \in \mathbb{R} \to -f(\xi) \in \mathbb{R}$}
\AXC{$(f(\xi) \in \mathbb{R} \to -f(\xi) \in \mathbb{R}) \to \varphi_{190}$\lbl{A1}}
\BIC{$\varphi_{190}$}
\BIC{$\mathrm{ContOn}(f,[a,b]) \to \varphi_{156}$}
\end{ptree}
\noindent$\Pi_{68}$:
\begin{ptree}
\AXC{$\forall x \in \mathbb{R}\,(-x \in \mathbb{R})$\lbl{(R1)}}
\AXC{$\forall x \in \mathbb{R}\,(-x \in \mathbb{R}) \to f(x) \in \mathbb{R} \to -f(x) \in \mathbb{R}$\lbl{A2}}
\BIC{$f(x) \in \mathbb{R} \to -f(x) \in \mathbb{R}$}
\end{ptree}
\noindent$\Pi_{69}$:
\begin{ptree}
\AXC{$\stackrel{\Pi_{20}}{\vdots}$}\noLine
\UIC{$\mathrm{ContOn}(f,[a,b]) \to \varphi_{141}$}
\AXC{$\stackrel{\Pi_{68}}{\vdots}$}\noLine
\UIC{$f(x) \in \mathbb{R} \to -f(x) \in \mathbb{R}$}
\AXC{$(f(x) \in \mathbb{R} \to -f(x) \in \mathbb{R}) \to \varphi_{189}$\lbl{A1}}
\BIC{$\varphi_{189}$}
\BIC{$\mathrm{ContOn}(f,[a,b]) \to \varphi_{154}$}
\end{ptree}
\noindent$\Pi_{70}$:
\begin{ptree}
\AXC{$\forall x\,\forall y\,(x \in \mathbb{R} \to y \in \mathbb{R} \to x \leq y \to -y \leq -x)$\lbl{[F23]}}
\AXC{$\forall x\,\forall y\,(x \in \mathbb{R} \to y \in \mathbb{R} \to x \leq y \to -y \leq -x) \to \varphi_{58}$\lbl{A2}}
\BIC{$\varphi_{58}$}
\end{ptree}
\noindent$\Pi_{71}$:
\begin{ptree}
\AXC{$\stackrel{\Pi_{69}}{\vdots}$}\noLine
\UIC{$\mathrm{ContOn}(f,[a,b]) \to \varphi_{154}$}
\AXC{$\stackrel{\Pi_{70}}{\vdots}$}\noLine
\UIC{$\varphi_{58}$}
\AXC{$\varphi_{58} \to -f(x) \in \mathbb{R} \to \varphi_{89}$\lbl{A2}}
\BIC{$-f(x) \in \mathbb{R} \to \varphi_{89}$}
\AXC{$(-f(x) \in \mathbb{R} \to \varphi_{89}) \to \varphi_{232}$\lbl{A1}}
\BIC{$\varphi_{232}$}
\BIC{$(\mathrm{ContOn}(f,[a,b]) \to \varphi_{156}) \to \varphi_{203}$}
\end{ptree}
\noindent$\Pi_{72}$:
\begin{ptree}
\AXC{$\stackrel{\Pi_{65}}{\vdots}$}\noLine
\UIC{$\mathrm{ContOn}(f,[a,b]) \to \varphi_{162}$}
\AXC{$\stackrel{\Pi_{67}}{\vdots}$}\noLine
\UIC{$\mathrm{ContOn}(f,[a,b]) \to \varphi_{156}$}
\AXC{$\stackrel{\Pi_{71}}{\vdots}$}\noLine
\UIC{$(\mathrm{ContOn}(f,[a,b]) \to \varphi_{156}) \to \varphi_{203}$}
\BIC{$\varphi_{203}$}
\BIC{$\mathrm{ContOn}(f,[a,b]) \to a \in \mathbb{R} \to b \in \mathbb{R} \to \varphi_{129}$}
\end{ptree}
\noindent$\Pi_{73}$:
\begin{ptree}
\AXC{$\forall x \in \mathbb{R}\,(-(-x) = x)$\lbl{[F5]}}
\AXC{$\varphi_{55}$\lbl{A2}}
\BIC{$f(\xi) \in \mathbb{R} \to -(-f(\xi)) = f(\xi)$}
\AXC{$(f(\xi) \in \mathbb{R} \to -(-f(\xi)) = f(\xi)) \to \varphi_{196}$\lbl{A1}}
\BIC{$\varphi_{196}$}
\end{ptree}
\noindent$\Pi_{74}$:
\begin{ptree}
\AXC{$\stackrel{\Pi_{50}}{\vdots}$}\noLine
\UIC{$\mathrm{ContOn}(f,[a,b]) \to \varphi_{143}$}
\AXC{$\stackrel{\Pi_{73}}{\vdots}$}\noLine
\UIC{$\varphi_{196}$}
\BIC{$\mathrm{ContOn}(f,[a,b]) \to a \in \mathbb{R} \to \varphi_{140}$}
\AXC{$(\mathrm{ContOn}(f,[a,b]) \to a \in \mathbb{R} \to \varphi_{140}) \to \varphi_{226}$\lbl{A1}}
\BIC{$\varphi_{226}$}
\end{ptree}
\noindent$\Pi_{75}$:
\begin{ptree}
\AXC{$\stackrel{\Pi_{4}}{\vdots}$}\noLine
\UIC{$-(-f(\xi)) = f(\xi) \to \varphi_{46}$}
\AXC{$\stackrel{\Pi_{72}}{\vdots}$}\noLine
\UIC{$\mathrm{ContOn}(f,[a,b]) \to a \in \mathbb{R} \to b \in \mathbb{R} \to \varphi_{129}$}
\AXC{$\stackrel{\Pi_{74}}{\vdots}$}\noLine
\UIC{$\varphi_{226}$}
\BIC{$(-(-f(\xi)) = f(\xi) \to \varphi_{46}) \to \varphi_{176}$}
\BIC{$\varphi_{176}$}
\end{ptree}
\noindent$\Pi_{76}$:
\begin{ptree}
\AXC{$\forall x \in \mathbb{R}\,(-(-x) = x)$\lbl{[F5]}}
\AXC{$\varphi_{52}$\lbl{A2}}
\BIC{$f(x) \in \mathbb{R} \to -(-f(x)) = f(x)$}
\AXC{$(f(x) \in \mathbb{R} \to -(-f(x)) = f(x)) \to \varphi_{193}$\lbl{A1}}
\BIC{$\varphi_{193}$}
\end{ptree}
\noindent$\Pi_{77}$:
\begin{ptree}
\AXC{$\stackrel{\Pi_{20}}{\vdots}$}\noLine
\UIC{$\mathrm{ContOn}(f,[a,b]) \to \varphi_{141}$}
\AXC{$\stackrel{\Pi_{76}}{\vdots}$}\noLine
\UIC{$\varphi_{193}$}
\BIC{$\mathrm{ContOn}(f,[a,b]) \to a \in \mathbb{R} \to \varphi_{134}$}
\AXC{$(\mathrm{ContOn}(f,[a,b]) \to a \in \mathbb{R} \to \varphi_{134}) \to \varphi_{224}$\lbl{A1}}
\BIC{$\varphi_{224}$}
\end{ptree}
\noindent$\Pi_{78}$:
\begin{ptree}
\AXC{$\varphi_{194}$\lbl{A3}}
\AXC{$\stackrel{\Pi_{1}}{\vdots}$}\noLine
\UIC{$\varphi_{51}$}
\AXC{$\stackrel{\Pi_{2}}{\vdots}$}\noLine
\UIC{$\varphi_{61}$}
\AXC{$\stackrel{\Pi_{3}}{\vdots}$}\noLine
\UIC{$-(-f(\xi)) = f(\xi) \to \varphi_{35}$}
\AXC{$\stackrel{\Pi_{75}}{\vdots}$}\noLine
\UIC{$\varphi_{176}$}
\BIC{$\mathrm{ContOn}(f,[a,b]) \to a \in \mathbb{R} \to \varphi_{138}$}
\AXC{$\stackrel{\Pi_{77}}{\vdots}$}\noLine
\UIC{$\varphi_{224}$}
\BIC{$\varphi_{61} \to \varphi_{51} \to \mathrm{ContOn}(f,[a,b]) \to \varphi_{144}$}
\BIC{$\varphi_{51} \to \mathrm{ContOn}(f,[a,b]) \to \varphi_{144}$}
\BIC{$\mathrm{ContOn}(f,[a,b]) \to \varphi_{144}$}
\AXC{$\varphi_{187}$\lbl{A1}}
\BIC{$\varphi_{121} \to x \in [a,b] \to f(\xi) \leq f(x)$}
\UIC{$\forall x\,(\varphi_{121} \to x \in [a,b] \to f(\xi) \leq f(x))$}
\AXC{$\varphi_{234}$\lbl{A1}}
\BIC{$\varphi_{194} \to (\varphi_{121} \to \forall x\,(\varphi_{121})) \to \mathrm{ContOn}(f,[a,b]) \to \varphi_{151}$}
\BIC{$(\varphi_{121} \to \forall x\,(\varphi_{121})) \to \mathrm{ContOn}(f,[a,b]) \to \varphi_{151}$}
\end{ptree}
\noindent$\Pi_{79}$:
\begin{ptree}
\AXC{$\varphi_{121} \to \forall x\,(\varphi_{121})$\lbl{A4}}
\AXC{$\stackrel{\Pi_{78}}{\vdots}$}\noLine
\UIC{$(\varphi_{121} \to \forall x\,(\varphi_{121})) \to \mathrm{ContOn}(f,[a,b]) \to \varphi_{151}$}
\BIC{$\mathrm{ContOn}(f,[a,b]) \to \varphi_{151}$}
\AXC{$\varphi_{170}$\lbl{A1}}
\AXC{$\varphi_{170} \to \varphi_{148}$\lbl{A1}}
\BIC{$\varphi_{148}$}
\AXC{$\varphi_{148} \to \varphi_{202}$\lbl{A1}}
\BIC{$\varphi_{202}$}
\BIC{$\mathrm{ContOn}(f,[a,b]) \to a \in \mathbb{R} \to b \in \mathbb{R} \to \varphi_{135}$}
\end{ptree}
\noindent$\Pi_{80}$:
\begin{ptree}
\AXC{$\varphi_{59}$\lbl{A2}}
\AXC{$\forall f\,(\varphi_{31})$\lbl{DefContOn}}
\AXC{$\forall f\,(\varphi_{31}) \to \varphi_{31}$\lbl{A2}}
\BIC{$\varphi_{31}$}
\AXC{$\varphi_{31} \to \varphi_{90}$\lbl{A2}}
\BIC{$\varphi_{90}$}
\AXC{$\varphi_{1}$\lbl{A1}}
\AXC{$\varphi_{1} \to \varphi_{90} \to \varphi_{106}$\lbl{A1}}
\BIC{$\varphi_{90} \to \varphi_{106}$}
\BIC{$\varphi_{106}$}
\AXC{$\varphi_{106} \to \varphi_{59} \to \varphi_{104}$\lbl{A1}}
\BIC{$\varphi_{59} \to \varphi_{104}$}
\BIC{$\varphi_{104}$}
\end{ptree}
\noindent$\Pi_{81}$:
\begin{ptree}
\AXC{$\varphi_{20}$\lbl{定理2.3(定数倍)}}
\AXC{$\varphi_{20} \to \forall f\,(\varphi_{29})$\lbl{A2}}
\BIC{$\forall f\,(\varphi_{29})$}
\AXC{$\forall f\,(\varphi_{29}) \to \varphi_{29}$\lbl{A2}}
\BIC{$\varphi_{29}$}
\AXC{$\varphi_{29} \to \forall c\,(-1 \in \mathbb{R} \to \varphi_{10})$\lbl{A2}}
\BIC{$\forall c\,(-1 \in \mathbb{R} \to \varphi_{10})$}
\end{ptree}
\noindent$\Pi_{82}$:
\begin{ptree}
\AXC{$\stackrel{\Pi_{30}}{\vdots}$}\noLine
\UIC{$\varphi_{34}$}
\AXC{$\stackrel{\Pi_{81}}{\vdots}$}\noLine
\UIC{$\forall c\,(-1 \in \mathbb{R} \to \varphi_{10})$}
\AXC{$\forall c\,(-1 \in \mathbb{R} \to \varphi_{10}) \to -1 \in \mathbb{R} \to \varphi_{10}$\lbl{A2}}
\BIC{$-1 \in \mathbb{R} \to \varphi_{10}$}
\AXC{$(-1 \in \mathbb{R} \to \varphi_{10}) \to \varphi_{34} \to \varphi_{166}$\lbl{A1}}
\BIC{$\varphi_{34} \to \varphi_{166}$}
\BIC{$\varphi_{166}$}
\end{ptree}
\noindent$\Pi_{83}$:
\begin{ptree}
\AXC{$\varphi_{69}$\lbl{A1}}
\AXC{$\stackrel{\Pi_{80}}{\vdots}$}\noLine
\UIC{$\varphi_{104}$}
\AXC{$\varphi_{104} \to \varphi_{69} \to \mathrm{ContOn}(f,[a,b]) \to a \in \mathbb{R} \to \varphi_{9}$\lbl{A1}}
\BIC{$\varphi_{69} \to \mathrm{ContOn}(f,[a,b]) \to a \in \mathbb{R} \to \varphi_{9}$}
\BIC{$\mathrm{ContOn}(f,[a,b]) \to a \in \mathbb{R} \to \varphi_{9}$}
\AXC{$\stackrel{\Pi_{82}}{\vdots}$}\noLine
\UIC{$\varphi_{166}$}
\BIC{$\mathrm{ContOn}(f,[a,b]) \to a \in \mathbb{R} \to b \in \mathbb{R} \to \varphi_{11}$}
\AXC{$\varphi_{167}$\lbl{A1}}
\BIC{$\varphi_{96}$}
\UIC{$\forall c\,(\varphi_{96})$}
\end{ptree}
\noindent$\Pi_{84}$:
\begin{ptree}
\AXC{$\varphi_{98}$\lbl{A4}}
\AXC{$\forall c\,(\varphi_{96}) \to \varphi_{2} \to \varphi_{5}$\lbl{A3}}
\AXC{$\stackrel{\Pi_{83}}{\vdots}$}\noLine
\UIC{$\forall c\,(\varphi_{96})$}
\AXC{$\forall c\,(\varphi_{96}) \to (\forall c\,(\varphi_{96}) \to \varphi_{2} \to \varphi_{5}) \to \varphi_{98} \to \varphi_{103}$\lbl{A1}}
\BIC{$(\forall c\,(\varphi_{96}) \to \varphi_{2} \to \varphi_{5}) \to \varphi_{98} \to \varphi_{103}$}
\BIC{$\varphi_{98} \to \varphi_{103}$}
\BIC{$\varphi_{103}$}
\end{ptree}
\noindent$\Pi_{85}$:
\begin{ptree}
\AXC{$\forall X\,\forall Y\,\forall x\,(X \subseteq Y \to x \in X \to x \in Y)$\lbl{[B5]}}
\AXC{$\varphi_{76}$\lbl{A2}}
\BIC{$\forall Y\,\forall x\,([a,b] \subseteq Y \to x \in [a,b] \to x \in Y)$}
\AXC{$\forall Y\,\forall x\,([a,b] \subseteq Y \to x \in [a,b] \to x \in Y) \to \varphi_{6}$\lbl{A2}}
\BIC{$\varphi_{6}$}
\AXC{$\varphi_{107}$\lbl{A2}}
\BIC{$[a,b] \subseteq \mathrm{dom}(f) \to u \in [a,b] \to u \in \mathrm{dom}(f)$}
\end{ptree}
\noindent$\Pi_{86}$:
\begin{ptree}
\AXC{$\varphi_{70}$\lbl{A1}}
\AXC{$\stackrel{\Pi_{14}}{\vdots}$}\noLine
\UIC{$\varphi_{45}$}
\AXC{$\stackrel{\Pi_{85}}{\vdots}$}\noLine
\UIC{$[a,b] \subseteq \mathrm{dom}(f) \to u \in [a,b] \to u \in \mathrm{dom}(f)$}
\AXC{$\varphi_{178}$\lbl{A1}}
\BIC{$\varphi_{45} \to \varphi_{70} \to \varphi_{60}$}
\BIC{$\varphi_{70} \to \varphi_{60}$}
\BIC{$\varphi_{60}$}
\AXC{$\varphi_{60} \to \varphi_{164}$\lbl{A1}}
\BIC{$\varphi_{164}$}
\end{ptree}
\noindent$\Pi_{87}$:
\begin{ptree}
\AXC{$\varphi_{74}$\lbl{A2}}
\AXC{$\stackrel{\Pi_{33}}{\vdots}$}\noLine
\UIC{$\mathrm{ContOn}(f,[a,b]) \to a \in \mathbb{R} \to \varphi_{22}$}
\AXC{$\varphi_{173}$\lbl{A1}}
\BIC{$\varphi_{74} \to \mathrm{ContOn}(f,[a,b]) \to a \in \mathbb{R} \to \varphi_{15}$}
\BIC{$\mathrm{ContOn}(f,[a,b]) \to a \in \mathbb{R} \to \varphi_{15}$}
\AXC{$\stackrel{\Pi_{86}}{\vdots}$}\noLine
\UIC{$\varphi_{164}$}
\BIC{$\mathrm{ContOn}(f,[a,b]) \to \varphi_{49}$}
\AXC{$(\mathrm{ContOn}(f,[a,b]) \to \varphi_{49}) \to \varphi_{93}$\lbl{A1}}
\BIC{$\varphi_{93}$}
\UIC{$\forall u\,(\varphi_{93})$}
\end{ptree}
\noindent$\Pi_{88}$:
\begin{ptree}
\AXC{$\varphi_{99}$\lbl{A4}}
\AXC{$\forall u\,(\varphi_{93}) \to \varphi_{3} \to \forall u \in [a,b]\,(u \in \mathrm{dom}((-1 f)))$\lbl{A3}}
\AXC{$\stackrel{\Pi_{87}}{\vdots}$}\noLine
\UIC{$\forall u\,(\varphi_{93})$}
\AXC{$\forall u\,(\varphi_{93}) \to \varphi_{208}$\lbl{A1}}
\BIC{$\varphi_{208}$}
\BIC{$\varphi_{99} \to \mathrm{ContOn}(f,[a,b]) \to a \in \mathbb{R} \to \varphi_{28}$}
\BIC{$\mathrm{ContOn}(f,[a,b]) \to a \in \mathbb{R} \to \varphi_{28}$}
\end{ptree}
\noindent$\Pi_{89}$:
\begin{ptree}
\AXC{$\forall x\,\forall y\,(x \subseteq y \leftrightarrow \forall u \in x\,(u \in y))$\lbl{Def$\subseteq$}}
\AXC{$\varphi_{81}$\lbl{A2}}
\BIC{$\forall y\,([a,b] \subseteq y \leftrightarrow \forall u \in [a,b]\,(u \in y))$}
\AXC{$\forall y\,([a,b] \subseteq y \leftrightarrow \forall u \in [a,b]\,(u \in y)) \to \varphi_{62}$\lbl{A2}}
\BIC{$\varphi_{62}$}
\end{ptree}
\noindent$\Pi_{90}$:
\begin{ptree}
\AXC{$\stackrel{\Pi_{88}}{\vdots}$}\noLine
\UIC{$\mathrm{ContOn}(f,[a,b]) \to a \in \mathbb{R} \to \varphi_{28}$}
\AXC{$\stackrel{\Pi_{89}}{\vdots}$}\noLine
\UIC{$\varphi_{62}$}
\AXC{$\varphi_{62} \to (\mathrm{ContOn}(f,[a,b]) \to a \in \mathbb{R} \to \varphi_{28}) \to \varphi_{75}$\lbl{A1}}
\BIC{$(\mathrm{ContOn}(f,[a,b]) \to a \in \mathbb{R} \to \varphi_{28}) \to \varphi_{75}$}
\BIC{$\varphi_{75}$}
\end{ptree}
\noindent$\Pi_{91}$:
\begin{ptree}
\AXC{$\forall f\,(\varphi_{31})$\lbl{DefContOn}}
\AXC{$\forall f\,(\varphi_{31}) \to \varphi_{114}$\lbl{A2}}
\BIC{$\varphi_{114}$}
\AXC{$\varphi_{114} \to \varphi_{132}$\lbl{A2}}
\BIC{$\varphi_{132}$}
\AXC{$\varphi_{132} \to \varphi_{44} \to \varphi_{75} \to \varphi_{103} \to \mathrm{ContOn}(f,[a,b]) \to \varphi_{26}$\lbl{A1}}
\BIC{$\varphi_{44} \to \varphi_{75} \to \varphi_{103} \to \mathrm{ContOn}(f,[a,b]) \to \varphi_{26}$}
\end{ptree}
\noindent$\Pi_{92}$:
\begin{ptree}
\AXC{$\varphi_{117}$\lbl{定理2.17}}
\AXC{$\varphi_{117} \to \forall a\,\forall b\,(\varphi_{146})$\lbl{A2}}
\BIC{$\forall a\,\forall b\,(\varphi_{146})$}
\AXC{$\forall a\,\forall b\,(\varphi_{146}) \to \forall b\,(\varphi_{146})$\lbl{A2}}
\BIC{$\forall b\,(\varphi_{146})$}
\AXC{$\forall b\,(\varphi_{146}) \to \varphi_{146}$\lbl{A2}}
\BIC{$\varphi_{146}$}
\AXC{$\varphi_{146} \to \varphi_{205}$\lbl{A1}}
\BIC{$\varphi_{205}$}
\end{ptree}
\noindent$\Pi_{93}$:
\begin{ptree}
\AXC{$\stackrel{\Pi_{84}}{\vdots}$}\noLine
\UIC{$\varphi_{103}$}
\AXC{$\stackrel{\Pi_{90}}{\vdots}$}\noLine
\UIC{$\varphi_{75}$}
\AXC{$\stackrel{\Pi_{32}}{\vdots}$}\noLine
\UIC{$\varphi_{165}$}
\AXC{$\varphi_{165} \to \varphi_{44}$\lbl{A1}}
\BIC{$\varphi_{44}$}
\AXC{$\stackrel{\Pi_{91}}{\vdots}$}\noLine
\UIC{$\varphi_{44} \to \varphi_{75} \to \varphi_{103} \to \mathrm{ContOn}(f,[a,b]) \to \varphi_{26}$}
\BIC{$\varphi_{75} \to \varphi_{103} \to \mathrm{ContOn}(f,[a,b]) \to \varphi_{26}$}
\BIC{$\varphi_{103} \to \mathrm{ContOn}(f,[a,b]) \to \varphi_{26}$}
\BIC{$\mathrm{ContOn}(f,[a,b]) \to \varphi_{26}$}
\AXC{$\stackrel{\Pi_{92}}{\vdots}$}\noLine
\UIC{$\varphi_{205}$}
\BIC{$(\mathrm{ContOn}(f,[a,b]) \to a \in \mathbb{R} \to a \in \mathbb{R}) \to \varphi_{171}$}
\end{ptree}
\noindent$\Pi_{94}$:
\begin{ptree}
\AXC{$\mathrm{ContOn}(f,[a,b]) \to a \in \mathbb{R} \to a \in \mathbb{R}$\lbl{A1}}
\AXC{$\stackrel{\Pi_{93}}{\vdots}$}\noLine
\UIC{$(\mathrm{ContOn}(f,[a,b]) \to a \in \mathbb{R} \to a \in \mathbb{R}) \to \varphi_{171}$}
\BIC{$\varphi_{171}$}
\end{ptree}
\noindent$\Pi_{95}$:
\begin{ptree}
\AXC{$\varphi_{7}$\lbl{A1}}
\AXC{$\mathrm{ContOn}(f,[a,b]) \to a \in \mathbb{R} \to b \in \mathbb{R} \to b \in \mathbb{R}$\lbl{A1}}
\AXC{$\stackrel{\Pi_{94}}{\vdots}$}\noLine
\UIC{$\varphi_{171}$}
\BIC{$\varphi_{7} \to \varphi_{123}$}
\BIC{$\varphi_{123}$}
\AXC{$\varphi_{123} \to \varphi_{233}$\lbl{A1}}
\BIC{$\varphi_{233}$}
\end{ptree}
\noindent$\Pi_{96}$:
\begin{ptree}
\AXC{$\varphi_{115}$\lbl{A2}}
\AXC{$\stackrel{\Pi_{79}}{\vdots}$}\noLine
\UIC{$\mathrm{ContOn}(f,[a,b]) \to a \in \mathbb{R} \to b \in \mathbb{R} \to \varphi_{135}$}
\AXC{$\varphi_{222}$\lbl{A1}}
\BIC{$\varphi_{115} \to \mathrm{ContOn}(f,[a,b]) \to a \in \mathbb{R} \to b \in \mathbb{R} \to \varphi_{145}$}
\BIC{$\mathrm{ContOn}(f,[a,b]) \to a \in \mathbb{R} \to b \in \mathbb{R} \to \varphi_{145}$}
\AXC{$\varphi_{207}$\lbl{A1}}
\BIC{$\varphi_{111} \to \neg \varphi_{68}$}
\UIC{$\forall \xi\,(\varphi_{111} \to \neg \varphi_{68})$}
\AXC{$\stackrel{\Pi_{95}}{\vdots}$}\noLine
\UIC{$\varphi_{233}$}
\BIC{$(\forall \xi\,(\varphi_{111} \to \neg \varphi_{68}) \to \forall \xi\,(\varphi_{111}) \to \forall \xi\,(\neg \varphi_{68})) \to \varphi_{186}$}
\end{ptree}
\noindent 主木:
\begin{ptree}
\AXC{$\varphi_{111} \to \forall \xi\,(\varphi_{111})$\lbl{A4}}
\AXC{$\forall \xi\,(\varphi_{111} \to \neg \varphi_{68}) \to \forall \xi\,(\varphi_{111}) \to \forall \xi\,(\neg \varphi_{68})$\lbl{A3}}
\AXC{$\stackrel{\Pi_{96}}{\vdots}$}\noLine
\UIC{$(\forall \xi\,(\varphi_{111} \to \neg \varphi_{68}) \to \forall \xi\,(\varphi_{111}) \to \forall \xi\,(\neg \varphi_{68})) \to \varphi_{186}$}
\BIC{$\varphi_{186}$}
\BIC{$\mathrm{ContOn}(f,[a,b]) \to a \in \mathbb{R} \to b \in \mathbb{R} \to \varphi_{32}$}
\UIC{$\varphi_{112}$}
\UIC{$\forall a\,(\varphi_{112})$}
\UIC{$\forall f\,\forall a\,(\varphi_{112})$}
\end{ptree}
\begin{abbrevlist}
\item $\varphi_{1} :\equiv \mathrm{ContOn}(f,[a,b]) \to \mathrm{ContOn}(f,[a,b])$
\item $\varphi_{2} :\equiv \forall c\,(\mathrm{ContOn}(f,[a,b]) \land a \in \mathbb{R} \land b \in \mathbb{R} \land a \leq b)$
\item $\varphi_{3} :\equiv \forall u\,(\mathrm{ContOn}(f,[a,b]) \land a \in \mathbb{R} \land b \in \mathbb{R} \land a \leq b)$
\item $\varphi_{4} :\equiv \forall z\,\forall w\,(z = w \to z = -(-f(x)) \to w = -(-f(x)))$
\item $\varphi_{5} :\equiv \forall c \in [a,b]\,(\mathrm{ContIn}((-1 f),[a,b],c))$
\item $\varphi_{6} :\equiv \forall x\,([a,b] \subseteq \mathrm{dom}(f) \to x \in [a,b] \to x \in \mathrm{dom}(f))$
\item $\varphi_{7} :\equiv \mathrm{ContOn}(f,[a,b]) \to a \in \mathbb{R} \to b \in \mathbb{R} \to a \leq b \to a \leq b$
\item $\varphi_{8} :\equiv c \in [a,b] \to c \in [a,b] \to \mathrm{ContIn}(f,[a,b],c)$
\item $\varphi_{9} :\equiv b \in \mathbb{R} \to a \leq b \to c \in [a,b] \to \mathrm{ContIn}(f,[a,b],c)$
\item $\varphi_{10} :\equiv \mathrm{ContIn}(f,[a,b],c) \to \mathrm{ContIn}((-1 f),[a,b],c)$
\item $\varphi_{11} :\equiv a \leq b \to c \in [a,b] \to \mathrm{ContIn}((-1 f),[a,b],c)$
\item $\varphi_{12} :\equiv \mathrm{ContOn}(f,[a,b]) \to a \in \mathbb{R} \to b \in \mathbb{R} \to a \leq b \to \mathrm{RFn}(f)$
\item $\varphi_{13} :\equiv b \in \mathbb{R} \to a \leq b \to (x \in \mathrm{dom}((-1 f)) \leftrightarrow x \in \mathrm{dom}(f))$
\item $\varphi_{14} :\equiv c \in [a,b] \to \forall c \in [a,b]\,(\mathrm{ContIn}(f,[a,b],c))$
\item $\varphi_{15} :\equiv b \in \mathbb{R} \to a \leq b \to (u \in \mathrm{dom}((-1 f)) \leftrightarrow u \in \mathrm{dom}(f))$
\item $\varphi_{16} :\equiv \forall z\,\forall w\,(z = w \to z < -(-f(x)) \to w < -(-f(x)))$
\item $\varphi_{17} :\equiv \forall w\,(z = w \to -(1 \cdot f(x)) = -z \to -(1 \cdot f(x)) = -w)$
\item $\varphi_{18} :\equiv b \in \mathbb{R} \to a \leq b \to (\xi \in \mathrm{dom}((-1 f)) \leftrightarrow \xi \in \mathrm{dom}(f))$
\item $\varphi_{19} :\equiv \forall y\,\forall z\,((-1 f)(x) = y \to y = z \to (-1 f)(x) = z)$
\item $\varphi_{20} :\equiv \forall l\,\forall f\,\forall A\,\forall c\,(l \in \mathbb{R} \to \mathrm{ContIn}(f,A,c) \to \mathrm{ContIn}((l f),A,c))$
\item $\varphi_{21} :\equiv \forall \xi\,(\neg (\xi \in [a,b] \land \forall x \in [a,b]\,(f(\xi) \leq f(x))))$
\item $\varphi_{22} :\equiv b \in \mathbb{R} \to a \leq b \to \forall x\,(x \in \mathrm{dom}((-1 f)) \leftrightarrow x \in \mathrm{dom}(f))$
\item $\varphi_{23} :\equiv \forall w\,(z = w \to -(1 \cdot f(\xi)) = -z \to -(1 \cdot f(\xi)) = -w)$
\item $\varphi_{24} :\equiv \forall z\,\forall w\,(z = w \to z = (-1 f)(\xi) \to w = (-1 f)(\xi))$
\item $\varphi_{25} :\equiv \forall y\,\forall z\,((-1 f)(\xi) = y \to y = z \to (-1 f)(\xi) = z)$
\item $\varphi_{26} :\equiv a \in \mathbb{R} \to b \in \mathbb{R} \to a \leq b \to \mathrm{ContOn}((-1 f),[a,b])$
\item $\varphi_{27} :\equiv \forall z\,(1 \cdot f(x) = f(x) \cdot 1 \to f(x) \cdot 1 = z \to 1 \cdot f(x) = z)$
\item $\varphi_{28} :\equiv b \in \mathbb{R} \to a \leq b \to \forall u \in [a,b]\,(u \in \mathrm{dom}((-1 f)))$
\item $\varphi_{29} :\equiv \forall A\,\forall c\,(-1 \in \mathbb{R} \to \mathrm{ContIn}(f,A,c) \to \mathrm{ContIn}((-1 f),A,c))$
\item $\varphi_{30} :\equiv \forall w\,(-(-f(x)) = w \to f(\xi) = -(-f(x)) \to f(\xi) = w)$
\item $\varphi_{31} :\equiv \forall A\,(\mathrm{ContOn}(f,A) \leftrightarrow \mathrm{RFn}(f) \land A \subseteq \mathrm{dom}(f) \land \forall c \in A\,(\mathrm{ContIn}(f,A,c)))$
\item $\varphi_{32} :\equiv a \leq b \to \exists \xi \in [a,b]\,\forall x \in [a,b]\,(f(\xi) \leq f(x))$
\item $\varphi_{33} :\equiv \forall x \in \mathbb{R}\,(x \cdot 1 = x) \to f(x) \in \mathbb{R} \to f(x) \cdot 1 = f(x)$
\item $\varphi_{34} :\equiv \mathrm{ContOn}(f,[a,b]) \to a \in \mathbb{R} \to b \in \mathbb{R} \to a \leq b \to -1 \in \mathbb{R}$
\item $\varphi_{35} :\equiv -(-f(\xi)) = -(-f(x)) \to f(\xi) = -(-f(x))$
\item $\varphi_{36} :\equiv \forall z\,(1 \cdot f(\xi) = f(\xi) \cdot 1 \to f(\xi) \cdot 1 = z \to 1 \cdot f(\xi) = z)$
\item $\varphi_{37} :\equiv \forall z\,\forall w\,(z = w \to z < (-1 f)(\xi) \to w < (-1 f)(\xi))$
\item $\varphi_{38} :\equiv \forall x\,\forall y\,\forall z\,(x = y \to y = z \to x = z) \to \forall y\,\forall z\,(1 \cdot f(x) = y \to y = z \to 1 \cdot f(x) = z)$
\item $\varphi_{39} :\equiv \forall x \in \mathbb{R}\,(x \cdot 1 = x) \to f(\xi) \in \mathbb{R} \to f(\xi) \cdot 1 = f(\xi)$
\item $\varphi_{40} :\equiv -(1 \cdot f(x)) = -(1 \cdot f(x)) \to -(1 \cdot f(x)) = -f(x)$
\item $\varphi_{41} :\equiv \forall x \in \mathrm{dom}((-1 f))\,((-1 f)(x) = -1 \cdot f(x))$
\item $\varphi_{42} :\equiv \forall w\,(-(-f(x)) = w \to f(\xi) < -(-f(x)) \to f(\xi) < w)$
\item $\varphi_{43} :\equiv \forall x\,\forall y\,\forall z\,(x = y \to y = z \to x = z) \to \forall y\,\forall z\,(1 \cdot f(\xi) = y \to y = z \to 1 \cdot f(\xi) = z)$
\item $\varphi_{44} :\equiv \mathrm{ContOn}(f,[a,b]) \to a \in \mathbb{R} \to b \in \mathbb{R} \to a \leq b \to \mathrm{RFn}((-1 f))$
\item $\varphi_{45} :\equiv \mathrm{ContOn}(f,[a,b]) \to a \in \mathbb{R} \to b \in \mathbb{R} \to a \leq b \to [a,b] \subseteq \mathrm{dom}(f)$
\item $\varphi_{46} :\equiv -(-f(\xi)) < -(-f(x)) \to f(\xi) < -(-f(x))$
\item $\varphi_{47} :\equiv a \leq b \to x \in \mathrm{dom}((-1 f)) \to (-1 f)(x) = -1 \cdot f(x)$
\item $\varphi_{48} :\equiv -(1 \cdot f(\xi)) = -(1 \cdot f(\xi)) \to -(1 \cdot f(\xi)) = -f(\xi)$
\item $\varphi_{49} :\equiv a \in \mathbb{R} \to b \in \mathbb{R} \to a \leq b \to u \in [a,b] \to u \in \mathrm{dom}((-1 f))$
\item $\varphi_{50} :\equiv 1 \cdot f(x) = f(x) \cdot 1 \to f(x) \cdot 1 = f(x) \to 1 \cdot f(x) = f(x)$
\item $\varphi_{51} :\equiv -(-f(x)) = f(x) \to f(\xi) = -(-f(x)) \to f(\xi) = f(x)$
\item $\varphi_{52} :\equiv \forall x \in \mathbb{R}\,(-(-x) = x) \to f(x) \in \mathbb{R} \to -(-f(x)) = f(x)$
\item $\varphi_{53} :\equiv a \leq b \to \xi \in \mathrm{dom}((-1 f)) \to (-1 f)(\xi) = -1 \cdot f(\xi)$
\item $\varphi_{54} :\equiv -1 \cdot f(x) = -(1 \cdot f(x)) \to (-1 f)(x) = -(1 \cdot f(x))$
\item $\varphi_{55} :\equiv \forall x \in \mathbb{R}\,(-(-x) = x) \to f(\xi) \in \mathbb{R} \to -(-f(\xi)) = f(\xi)$
\item $\varphi_{56} :\equiv x \in [a,b] \to x \in [a,b] \to (-1 f)(x) \leq (-1 f)(\xi)$
\item $\varphi_{57} :\equiv 1 \cdot f(\xi) = f(\xi) \cdot 1 \to f(\xi) \cdot 1 = f(\xi) \to 1 \cdot f(\xi) = f(\xi)$
\item $\varphi_{58} :\equiv \forall y\,(-f(x) \in \mathbb{R} \to y \in \mathbb{R} \to -f(x) \leq y \to -y \leq -(-f(x)))$
\item $\varphi_{59} :\equiv \forall c \in [a,b]\,(\mathrm{ContIn}(f,[a,b],c)) \to c \in [a,b] \to \mathrm{ContIn}(f,[a,b],c)$
\item $\varphi_{60} :\equiv \mathrm{ContOn}(f,[a,b]) \to a \in \mathbb{R} \to b \in \mathbb{R} \to a \leq b \to u \in [a,b] \to u \in \mathrm{dom}(f)$
\item $\varphi_{61} :\equiv -(-f(x)) = f(x) \to f(\xi) < -(-f(x)) \to f(\xi) < f(x)$
\item $\varphi_{62} :\equiv [a,b] \subseteq \mathrm{dom}((-1 f)) \leftrightarrow \forall u \in [a,b]\,(u \in \mathrm{dom}((-1 f)))$
\item $\varphi_{63} :\equiv x \in [a,b] \to \forall x \in [a,b]\,((-1 f)(x) \leq (-1 f)(\xi))$
\item $\varphi_{64} :\equiv -1 \cdot f(\xi) = -(1 \cdot f(\xi)) \to (-1 f)(\xi) = -(1 \cdot f(\xi))$
\item $\varphi_{65} :\equiv \forall w\,((-1 f)(\xi) = w \to -f(x) = (-1 f)(\xi) \to -f(x) = w)$
\item $\varphi_{66} :\equiv \forall w\,(1 \cdot f(x) = w \to -(1 \cdot f(x)) = -(1 \cdot f(x)) \to -(1 \cdot f(x)) = -w)$
\item $\varphi_{67} :\equiv (-1 f)(x) = (-1 f)(\xi) \to -f(x) = (-1 f)(\xi)$
\item $\varphi_{68} :\equiv \xi \in [a,b] \land \forall x \in [a,b]\,((-1 f)(x) \leq (-1 f)(\xi))$
\item $\varphi_{69} :\equiv \mathrm{ContOn}(f,[a,b]) \to a \in \mathbb{R} \to b \in \mathbb{R} \to a \leq b \to c \in [a,b] \to c \in [a,b]$
\item $\varphi_{70} :\equiv \mathrm{ContOn}(f,[a,b]) \to a \in \mathbb{R} \to b \in \mathbb{R} \to a \leq b \to u \in [a,b] \to u \in [a,b]$
\item $\varphi_{71} :\equiv \forall w\,(-(-f(\xi)) = w \to -(-f(\xi)) = -(-f(x)) \to w = -(-f(x)))$
\item $\varphi_{72} :\equiv \forall x\,(\mathrm{RFn}(f) \to x \in \mathrm{dom}(f) \to f(x) \in \mathbb{R}) \to \mathrm{RFn}(f) \to x \in \mathrm{dom}(f) \to f(x) \in \mathbb{R}$
\item $\varphi_{73} :\equiv \forall x\,(x \in \mathrm{dom}((-1 f)) \leftrightarrow x \in \mathrm{dom}(f)) \to (x \in \mathrm{dom}((-1 f)) \leftrightarrow x \in \mathrm{dom}(f))$
\item $\varphi_{74} :\equiv \forall x\,(x \in \mathrm{dom}((-1 f)) \leftrightarrow x \in \mathrm{dom}(f)) \to (u \in \mathrm{dom}((-1 f)) \leftrightarrow u \in \mathrm{dom}(f))$
\item $\varphi_{75} :\equiv \mathrm{ContOn}(f,[a,b]) \to a \in \mathbb{R} \to b \in \mathbb{R} \to a \leq b \to [a,b] \subseteq \mathrm{dom}((-1 f))$
\item $\varphi_{76} :\equiv \forall X\,\forall Y\,\forall x\,(X \subseteq Y \to x \in X \to x \in Y) \to \forall Y\,\forall x\,([a,b] \subseteq Y \to x \in [a,b] \to x \in Y)$
\item $\varphi_{77} :\equiv \forall x\,\forall y\,(x \in \mathbb{R} \to y \in \mathbb{R} \to x \cdot y = y \cdot x) \to \forall y\,(1 \in \mathbb{R} \to y \in \mathbb{R} \to 1 \cdot y = y \cdot 1)$
\item $\varphi_{78} :\equiv \forall w\,(1 \cdot f(\xi) = w \to -(1 \cdot f(\xi)) = -(1 \cdot f(\xi)) \to -(1 \cdot f(\xi)) = -w)$
\item $\varphi_{79} :\equiv \forall x\,(\mathrm{RFn}(f) \to x \in \mathrm{dom}(f) \to f(x) \in \mathbb{R}) \to \mathrm{RFn}(f) \to \xi \in \mathrm{dom}(f) \to f(\xi) \in \mathbb{R}$
\item $\varphi_{80} :\equiv \forall x\,(x \in \mathrm{dom}((-1 f)) \leftrightarrow x \in \mathrm{dom}(f)) \to (\xi \in \mathrm{dom}((-1 f)) \leftrightarrow \xi \in \mathrm{dom}(f))$
\item $\varphi_{81} :\equiv \forall x\,\forall y\,(x \subseteq y \leftrightarrow \forall u \in x\,(u \in y)) \to \forall y\,([a,b] \subseteq y \leftrightarrow \forall u \in [a,b]\,(u \in y))$
\item $\varphi_{82} :\equiv \exists \xi \in [a,b]\,\forall x \in [a,b]\,((-1 f)(x) \leq (-1 f)(\xi))$
\item $\varphi_{83} :\equiv \forall w\,((-1 f)(\xi) = w \to -f(x) < (-1 f)(\xi) \to -f(x) < w)$
\item $\varphi_{84} :\equiv (-1 f)(x) < (-1 f)(\xi) \to -f(x) < (-1 f)(\xi)$
\item $\varphi_{85} :\equiv \forall w\,(-(-f(\xi)) = w \to -(-f(\xi)) < -(-f(x)) \to w < -(-f(x)))$
\item $\varphi_{86} :\equiv \forall z\,((-1 f)(x) = -(1 \cdot f(x)) \to -(1 \cdot f(x)) = z \to (-1 f)(x) = z)$
\item $\varphi_{87} :\equiv \forall f\,\forall x\,(\mathrm{RFn}(f) \to x \in \mathrm{dom}(f) \to f(x) \in \mathbb{R}) \to \forall x\,(\mathrm{RFn}(f) \to x \in \mathrm{dom}(f) \to f(x) \in \mathbb{R})$
\item $\varphi_{88} :\equiv \forall z\,((-1 f)(x) = -1 \cdot f(x) \to -1 \cdot f(x) = z \to (-1 f)(x) = z)$
\item $\varphi_{89} :\equiv -f(\xi) \in \mathbb{R} \to -f(x) \leq -f(\xi) \to -(-f(\xi)) \leq -(-f(x))$
\item $\varphi_{90} :\equiv \mathrm{ContOn}(f,[a,b]) \leftrightarrow \mathrm{RFn}(f) \land [a,b] \subseteq \mathrm{dom}(f) \land \forall c \in [a,b]\,(\mathrm{ContIn}(f,[a,b],c))$
\item $\varphi_{91} :\equiv \forall z\,((-1 f)(\xi) = -(1 \cdot f(\xi)) \to -(1 \cdot f(\xi)) = z \to (-1 f)(\xi) = z)$
\item $\varphi_{92} :\equiv \forall z\,((-1 f)(\xi) = -1 \cdot f(\xi) \to -1 \cdot f(\xi) = z \to (-1 f)(\xi) = z)$
\item $\varphi_{93} :\equiv \mathrm{ContOn}(f,[a,b]) \land a \in \mathbb{R} \land b \in \mathbb{R} \land a \leq b \to u \in [a,b] \to u \in \mathrm{dom}((-1 f))$
\item $\varphi_{94} :\equiv \forall w\,((-1 f)(x) = w \to (-1 f)(x) = (-1 f)(\xi) \to w = (-1 f)(\xi))$
\item $\varphi_{95} :\equiv \forall y\,(1 \in \mathbb{R} \to y \in \mathbb{R} \to 1 \cdot y = y \cdot 1) \to 1 \in \mathbb{R} \to f(x) \in \mathbb{R} \to 1 \cdot f(x) = f(x) \cdot 1$
\item $\varphi_{96} :\equiv \mathrm{ContOn}(f,[a,b]) \land a \in \mathbb{R} \land b \in \mathbb{R} \land a \leq b \to c \in [a,b] \to \mathrm{ContIn}((-1 f),[a,b],c)$
\item $\varphi_{97} :\equiv \forall y\,(1 \in \mathbb{R} \to y \in \mathbb{R} \to 1 \cdot y = y \cdot 1) \to 1 \in \mathbb{R} \to f(\xi) \in \mathbb{R} \to 1 \cdot f(\xi) = f(\xi) \cdot 1$
\item $\varphi_{98} :\equiv \mathrm{ContOn}(f,[a,b]) \land a \in \mathbb{R} \land b \in \mathbb{R} \land a \leq b \to \varphi_{2}$
\item $\varphi_{99} :\equiv \mathrm{ContOn}(f,[a,b]) \land a \in \mathbb{R} \land b \in \mathbb{R} \land a \leq b \to \varphi_{3}$
\item $\varphi_{100} :\equiv (-1 f)(\xi) = -f(\xi) \to -f(x) = (-1 f)(\xi) \to -f(x) = -f(\xi)$
\item $\varphi_{101} :\equiv \forall w\,((-1 f)(x) = w \to (-1 f)(x) < (-1 f)(\xi) \to w < (-1 f)(\xi))$
\item $\varphi_{102} :\equiv \forall x\,\forall y\,(x \in \mathbb{R} \to y \in \mathbb{R} \to -x \cdot y = -(x \cdot y)) \to \forall y\,(1 \in \mathbb{R} \to y \in \mathbb{R} \to -1 \cdot y = -(1 \cdot y))$
\item $\varphi_{103} :\equiv \mathrm{ContOn}(f,[a,b]) \to a \in \mathbb{R} \to b \in \mathbb{R} \to a \leq b \to \varphi_{5}$
\item $\varphi_{104} :\equiv \mathrm{ContOn}(f,[a,b]) \to a \in \mathbb{R} \to b \in \mathbb{R} \to a \leq b \to \varphi_{8}$
\item $\varphi_{105} :\equiv \varphi_{6} \to [a,b] \subseteq \mathrm{dom}(f) \to x \in [a,b] \to x \in \mathrm{dom}(f)$
\item $\varphi_{106} :\equiv \mathrm{ContOn}(f,[a,b]) \to a \in \mathbb{R} \to b \in \mathbb{R} \to a \leq b \to \varphi_{14}$
\item $\varphi_{107} :\equiv \varphi_{6} \to [a,b] \subseteq \mathrm{dom}(f) \to u \in [a,b] \to u \in \mathrm{dom}(f)$
\item $\varphi_{108} :\equiv (-1 f)(\xi) = -f(\xi) \to -f(x) < (-1 f)(\xi) \to -f(x) < -f(\xi)$
\item $\varphi_{109} :\equiv (-1 f)(x) = -(1 \cdot f(x)) \to -(1 \cdot f(x)) = -f(x) \to (-1 f)(x) = -f(x)$
\item $\varphi_{110} :\equiv \varphi_{6} \to [a,b] \subseteq \mathrm{dom}(f) \to \xi \in [a,b] \to \xi \in \mathrm{dom}(f)$
\item $\varphi_{111} :\equiv \mathrm{ContOn}(f,[a,b]) \land a \in \mathbb{R} \land b \in \mathbb{R} \land a \leq b \land \neg \exists \xi \in [a,b]\,\forall x \in [a,b]\,(f(\xi) \leq f(x))$
\item $\varphi_{112} :\equiv \forall b\,(\mathrm{ContOn}(f,[a,b]) \to a \in \mathbb{R} \to b \in \mathbb{R} \to \varphi_{32})$
\item $\varphi_{113} :\equiv (-1 f)(\xi) = -(1 \cdot f(\xi)) \to -(1 \cdot f(\xi)) = -f(\xi) \to (-1 f)(\xi) = -f(\xi)$
\item $\varphi_{114} :\equiv \forall A\,(\mathrm{ContOn}((-1 f),A) \leftrightarrow \mathrm{RFn}((-1 f)) \land A \subseteq \mathrm{dom}((-1 f)) \land \forall c \in A\,(\mathrm{ContIn}((-1 f),A,c)))$
\item $\varphi_{115} :\equiv \varphi_{21} \to \neg (\xi \in [a,b] \land \forall x \in [a,b]\,(f(\xi) \leq f(x)))$
\item $\varphi_{116} :\equiv \forall y\,(1 \in \mathbb{R} \to y \in \mathbb{R} \to -1 \cdot y = -(1 \cdot y)) \to 1 \in \mathbb{R} \to f(x) \in \mathbb{R} \to -1 \cdot f(x) = -(1 \cdot f(x))$
\item $\varphi_{117} :\equiv \forall f\,\forall a\,\forall b\,(\mathrm{ContOn}(f,[a,b]) \to a \in \mathbb{R} \to b \in \mathbb{R} \to a \leq b \to \exists \xi \in [a,b]\,\forall x \in [a,b]\,(f(x) \leq f(\xi)))$
\item $\varphi_{118} :\equiv \mathrm{ContOn}(f,[a,b]) \to a \in \mathbb{R} \to b \in \mathbb{R} \to a \leq b \to \varphi_{41}$
\item $\varphi_{119} :\equiv \forall y\,(1 \in \mathbb{R} \to y \in \mathbb{R} \to -1 \cdot y = -(1 \cdot y)) \to 1 \in \mathbb{R} \to f(\xi) \in \mathbb{R} \to -1 \cdot f(\xi) = -(1 \cdot f(\xi))$
\item $\varphi_{120} :\equiv \forall x \in [a,b]\,((-1 f)(x) \leq (-1 f)(\xi)) \to x \in [a,b] \to (-1 f)(x) \leq (-1 f)(\xi)$
\item $\varphi_{121} :\equiv \mathrm{ContOn}(f,[a,b]) \land a \in \mathbb{R} \land b \in \mathbb{R} \land a \leq b \land \varphi_{68}$
\item $\varphi_{122} :\equiv \mathrm{RFn}((-1 f)) \land \forall x\,(x \in \mathrm{dom}((-1 f)) \leftrightarrow x \in \mathrm{dom}(f)) \land \varphi_{41}$
\item $\varphi_{123} :\equiv \mathrm{ContOn}(f,[a,b]) \to a \in \mathbb{R} \to b \in \mathbb{R} \to a \leq b \to \varphi_{82}$
\item $\varphi_{124} :\equiv \forall l\,\forall f\,(l \in \mathbb{R} \to \mathrm{RFn}(f) \to \mathrm{RFn}((l f)) \land \forall x\,(x \in \mathrm{dom}((l f)) \leftrightarrow x \in \mathrm{dom}(f)) \land \forall x \in \mathrm{dom}((l f))\,((l f)(x) = l \cdot f(x)))$
\item $\varphi_{125} :\equiv \varphi_{41} \to x \in \mathrm{dom}((-1 f)) \to (-1 f)(x) = -1 \cdot f(x)$
\item $\varphi_{126} :\equiv \varphi_{41} \to \xi \in \mathrm{dom}((-1 f)) \to (-1 f)(\xi) = -1 \cdot f(\xi)$
\item $\varphi_{127} :\equiv a \leq b \to \varphi_{68} \to x \in [a,b] \to -1 \cdot f(x) = -(1 \cdot f(x))$
\item $\varphi_{128} :\equiv a \in \mathbb{R} \to b \in \mathbb{R} \to a \leq b \to \varphi_{68} \to x \in [a,b] \to x \in \mathrm{dom}(f)$
\item $\varphi_{129} :\equiv a \leq b \to \varphi_{68} \to x \in [a,b] \to -(-f(\xi)) \leq -(-f(x))$
\item $\varphi_{130} :\equiv a \in \mathbb{R} \to b \in \mathbb{R} \to a \leq b \to \varphi_{68} \to x \in [a,b] \to \xi \in \mathrm{dom}(f)$
\item $\varphi_{131} :\equiv a \leq b \to \varphi_{68} \to x \in [a,b] \to -1 \cdot f(\xi) = -(1 \cdot f(\xi))$
\item $\varphi_{132} :\equiv \mathrm{ContOn}((-1 f),[a,b]) \leftrightarrow \mathrm{RFn}((-1 f)) \land [a,b] \subseteq \mathrm{dom}((-1 f)) \land \varphi_{5}$
\item $\varphi_{133} :\equiv b \in \mathbb{R} \to a \leq b \to \varphi_{68} \to x \in [a,b] \to 1 \cdot f(x) = f(x) \cdot 1$
\item $\varphi_{134} :\equiv b \in \mathbb{R} \to a \leq b \to \varphi_{68} \to x \in [a,b] \to -(-f(x)) = f(x)$
\item $\varphi_{135} :\equiv a \leq b \to \varphi_{68} \to \xi \in [a,b] \land \forall x \in [a,b]\,(f(\xi) \leq f(x))$
\item $\varphi_{136} :\equiv a \in \mathbb{R} \to b \in \mathbb{R} \to a \leq b \to \varphi_{68} \to x \in [a,b] \to x \in [a,b]$
\item $\varphi_{137} :\equiv a \leq b \to \varphi_{68} \to x \in [a,b] \to (-1 f)(x) = -(1 \cdot f(x))$
\item $\varphi_{138} :\equiv b \in \mathbb{R} \to a \leq b \to \varphi_{68} \to x \in [a,b] \to f(\xi) \leq -(-f(x))$
\item $\varphi_{139} :\equiv b \in \mathbb{R} \to a \leq b \to \varphi_{68} \to x \in [a,b] \to 1 \cdot f(\xi) = f(\xi) \cdot 1$
\item $\varphi_{140} :\equiv b \in \mathbb{R} \to a \leq b \to \varphi_{68} \to x \in [a,b] \to -(-f(\xi)) = f(\xi)$
\item $\varphi_{141} :\equiv a \in \mathbb{R} \to b \in \mathbb{R} \to a \leq b \to \varphi_{68} \to x \in [a,b] \to f(x) \in \mathbb{R}$
\item $\varphi_{142} :\equiv a \leq b \to \varphi_{68} \to x \in [a,b] \to (-1 f)(\xi) = -(1 \cdot f(\xi))$
\item $\varphi_{143} :\equiv a \in \mathbb{R} \to b \in \mathbb{R} \to a \leq b \to \varphi_{68} \to x \in [a,b] \to f(\xi) \in \mathbb{R}$
\item $\varphi_{144} :\equiv a \in \mathbb{R} \to b \in \mathbb{R} \to a \leq b \to \varphi_{68} \to x \in [a,b] \to f(\xi) \leq f(x)$
\item $\varphi_{145} :\equiv a \leq b \to \varphi_{68} \to \exists \xi \in [a,b]\,\forall x \in [a,b]\,(f(\xi) \leq f(x))$
\item $\varphi_{146} :\equiv \mathrm{ContOn}((-1 f),[a,b]) \to a \in \mathbb{R} \to b \in \mathbb{R} \to a \leq b \to \varphi_{82}$
\item $\varphi_{147} :\equiv b \in \mathbb{R} \to a \leq b \to \varphi_{68} \to x \in [a,b] \to -(1 \cdot f(x)) = -f(x)$
\item $\varphi_{148} :\equiv \mathrm{ContOn}(f,[a,b]) \to a \in \mathbb{R} \to b \in \mathbb{R} \to a \leq b \to \varphi_{68} \to \xi \in [a,b]$
\item $\varphi_{149} :\equiv a \leq b \to \varphi_{68} \to x \in [a,b] \to (-1 f)(x) \leq (-1 f)(\xi)$
\item $\varphi_{150} :\equiv b \in \mathbb{R} \to a \leq b \to \varphi_{68} \to x \in [a,b] \to -(1 \cdot f(\xi)) = -f(\xi)$
\item $\varphi_{151} :\equiv a \in \mathbb{R} \to b \in \mathbb{R} \to a \leq b \to \varphi_{68} \to \forall x \in [a,b]\,(f(\xi) \leq f(x))$
\item $\varphi_{152} :\equiv a \in \mathbb{R} \to b \in \mathbb{R} \to a \leq b \to \varphi_{68} \to x \in [a,b] \to f(x) \cdot 1 = f(x)$
\item $\varphi_{153} :\equiv a \in \mathbb{R} \to b \in \mathbb{R} \to a \leq b \to \varphi_{68} \to x \in [a,b] \to 1 \cdot f(x) = f(x)$
\item $\varphi_{154} :\equiv a \in \mathbb{R} \to b \in \mathbb{R} \to a \leq b \to \varphi_{68} \to x \in [a,b] \to -f(x) \in \mathbb{R}$
\item $\varphi_{155} :\equiv b \in \mathbb{R} \to a \leq b \to \varphi_{68} \to x \in [a,b] \to (-1 f)(x) = -f(x)$
\item $\varphi_{156} :\equiv a \in \mathbb{R} \to b \in \mathbb{R} \to a \leq b \to \varphi_{68} \to x \in [a,b] \to -f(\xi) \in \mathbb{R}$
\item $\varphi_{157} :\equiv b \in \mathbb{R} \to a \leq b \to \varphi_{68} \to x \in [a,b] \to -f(x) \leq (-1 f)(\xi)$
\item $\varphi_{158} :\equiv a \in \mathbb{R} \to b \in \mathbb{R} \to a \leq b \to \varphi_{68} \to x \in [a,b] \to f(\xi) \cdot 1 = f(\xi)$
\item $\varphi_{159} :\equiv a \in \mathbb{R} \to b \in \mathbb{R} \to a \leq b \to \varphi_{68} \to x \in [a,b] \to 1 \cdot f(\xi) = f(\xi)$
\item $\varphi_{160} :\equiv b \in \mathbb{R} \to a \leq b \to \varphi_{68} \to x \in [a,b] \to (-1 f)(\xi) = -f(\xi)$
\item $\varphi_{161} :\equiv b \in \mathbb{R} \to a \leq b \to \varphi_{68} \to x \in [a,b] \to (-1 f)(x) = -1 \cdot f(x)$
\item $\varphi_{162} :\equiv a \in \mathbb{R} \to b \in \mathbb{R} \to a \leq b \to \varphi_{68} \to x \in [a,b] \to -f(x) \leq -f(\xi)$
\item $\varphi_{163} :\equiv b \in \mathbb{R} \to a \leq b \to \varphi_{68} \to x \in [a,b] \to (-1 f)(\xi) = -1 \cdot f(\xi)$
\item $\varphi_{164} :\equiv (\mathrm{ContOn}(f,[a,b]) \to a \in \mathbb{R} \to \varphi_{15}) \to \mathrm{ContOn}(f,[a,b]) \to \varphi_{49}$
\item $\varphi_{165} :\equiv \mathrm{ContOn}(f,[a,b]) \to a \in \mathbb{R} \to b \in \mathbb{R} \to a \leq b \to \varphi_{122}$
\item $\varphi_{166} :\equiv (\mathrm{ContOn}(f,[a,b]) \to a \in \mathbb{R} \to \varphi_{9}) \to \mathrm{ContOn}(f,[a,b]) \to a \in \mathbb{R} \to b \in \mathbb{R} \to \varphi_{11}$
\item $\varphi_{167} :\equiv (\mathrm{ContOn}(f,[a,b]) \to a \in \mathbb{R} \to b \in \mathbb{R} \to \varphi_{11}) \to \varphi_{96}$
\item $\varphi_{168} :\equiv \mathrm{ContOn}(f,[a,b]) \to a \in \mathbb{R} \to b \in \mathbb{R} \to a \leq b \to \varphi_{68} \to \varphi_{56}$
\item $\varphi_{169} :\equiv \mathrm{ContOn}(f,[a,b]) \to a \in \mathbb{R} \to b \in \mathbb{R} \to a \leq b \to \varphi_{68} \to \varphi_{63}$
\item $\varphi_{170} :\equiv \mathrm{ContOn}(f,[a,b]) \to a \in \mathbb{R} \to b \in \mathbb{R} \to a \leq b \to \varphi_{68} \to \varphi_{68}$
\item $\varphi_{171} :\equiv (\mathrm{ContOn}(f,[a,b]) \to a \in \mathbb{R} \to b \in \mathbb{R} \to b \in \mathbb{R}) \to \varphi_{7} \to \varphi_{123}$
\item $\varphi_{172} :\equiv (\mathrm{ContOn}(f,[a,b]) \to a \in \mathbb{R} \to \varphi_{22}) \to \varphi_{73} \to \mathrm{ContOn}(f,[a,b]) \to a \in \mathbb{R} \to \varphi_{13}$
\item $\varphi_{173} :\equiv (\mathrm{ContOn}(f,[a,b]) \to a \in \mathbb{R} \to \varphi_{22}) \to \varphi_{74} \to \mathrm{ContOn}(f,[a,b]) \to a \in \mathbb{R} \to \varphi_{15}$
\item $\varphi_{174} :\equiv (\mathrm{ContOn}(f,[a,b]) \to a \in \mathbb{R} \to \varphi_{22}) \to \varphi_{80} \to \mathrm{ContOn}(f,[a,b]) \to a \in \mathbb{R} \to \varphi_{18}$
\item $\varphi_{175} :\equiv (1 \cdot f(x) = f(x) \to \varphi_{40}) \to \mathrm{ContOn}(f,[a,b]) \to a \in \mathbb{R} \to \varphi_{147}$
\item $\varphi_{176} :\equiv (-(-f(\xi)) = f(\xi) \to \varphi_{35}) \to \mathrm{ContOn}(f,[a,b]) \to a \in \mathbb{R} \to \varphi_{138}$
\item $\varphi_{177} :\equiv (1 \cdot f(\xi) = f(\xi) \to \varphi_{48}) \to \mathrm{ContOn}(f,[a,b]) \to a \in \mathbb{R} \to \varphi_{150}$
\item $\varphi_{178} :\equiv ([a,b] \subseteq \mathrm{dom}(f) \to u \in [a,b] \to u \in \mathrm{dom}(f)) \to \varphi_{45} \to \varphi_{70} \to \varphi_{60}$
\item $\varphi_{179} :\equiv ((-1 f)(x) = -f(x) \to \varphi_{67}) \to \mathrm{ContOn}(f,[a,b]) \to a \in \mathbb{R} \to \varphi_{157}$
\item $\varphi_{180} :\equiv ((-1 f)(x) = -1 \cdot f(x) \to \varphi_{54}) \to \mathrm{ContOn}(f,[a,b]) \to a \in \mathbb{R} \to b \in \mathbb{R} \to \varphi_{137}$
\item $\varphi_{181} :\equiv ((-1 f)(\xi) = -1 \cdot f(\xi) \to \varphi_{64}) \to \mathrm{ContOn}(f,[a,b]) \to a \in \mathbb{R} \to b \in \mathbb{R} \to \varphi_{142}$
\item $\varphi_{182} :\equiv \forall f\,(-1 \in \mathbb{R} \to \mathrm{RFn}(f) \to \varphi_{122}) \to -1 \in \mathbb{R} \to \mathrm{RFn}(f) \to \varphi_{122}$
\item $\varphi_{183} :\equiv (\mathrm{ContOn}(f,[a,b]) \to \varphi_{136}) \to \mathrm{ContOn}(f,[a,b]) \to \varphi_{128}$
\item $\varphi_{184} :\equiv (\mathrm{ContOn}(f,[a,b]) \to \varphi_{128}) \to \mathrm{ContOn}(f,[a,b]) \to \varphi_{141}$
\item $\varphi_{185} :\equiv (\mathrm{ContOn}(f,[a,b]) \to \varphi_{130}) \to \mathrm{ContOn}(f,[a,b]) \to \varphi_{143}$
\item $\varphi_{186} :\equiv (\varphi_{111} \to \forall \xi\,(\varphi_{111})) \to \mathrm{ContOn}(f,[a,b]) \to a \in \mathbb{R} \to b \in \mathbb{R} \to \varphi_{32}$
\item $\varphi_{187} :\equiv (\mathrm{ContOn}(f,[a,b]) \to \varphi_{144}) \to \varphi_{121} \to x \in [a,b] \to f(\xi) \leq f(x)$
\item $\varphi_{188} :\equiv (\mathrm{ContOn}(f,[a,b]) \to \varphi_{141}) \to \mathrm{ContOn}(f,[a,b]) \to \varphi_{152}$
\item $\varphi_{189} :\equiv (\mathrm{ContOn}(f,[a,b]) \to \varphi_{141}) \to \mathrm{ContOn}(f,[a,b]) \to \varphi_{154}$
\item $\varphi_{190} :\equiv (\mathrm{ContOn}(f,[a,b]) \to \varphi_{143}) \to \mathrm{ContOn}(f,[a,b]) \to \varphi_{156}$
\item $\varphi_{191} :\equiv (\mathrm{ContOn}(f,[a,b]) \to \varphi_{143}) \to \mathrm{ContOn}(f,[a,b]) \to \varphi_{158}$
\item $\varphi_{192} :\equiv (\mathrm{ContOn}(f,[a,b]) \to \varphi_{141}) \to \mathrm{ContOn}(f,[a,b]) \to a \in \mathbb{R} \to \varphi_{133}$
\item $\varphi_{193} :\equiv (\mathrm{ContOn}(f,[a,b]) \to \varphi_{141}) \to \mathrm{ContOn}(f,[a,b]) \to a \in \mathbb{R} \to \varphi_{134}$
\item $\varphi_{194} :\equiv \forall x\,(\varphi_{121} \to x \in [a,b] \to f(\xi) \leq f(x)) \to \forall x\,(\varphi_{121}) \to \forall x \in [a,b]\,(f(\xi) \leq f(x))$
\item $\varphi_{195} :\equiv (\mathrm{ContOn}(f,[a,b]) \to \varphi_{143}) \to \mathrm{ContOn}(f,[a,b]) \to a \in \mathbb{R} \to \varphi_{139}$
\item $\varphi_{196} :\equiv (\mathrm{ContOn}(f,[a,b]) \to \varphi_{143}) \to \mathrm{ContOn}(f,[a,b]) \to a \in \mathbb{R} \to \varphi_{140}$
\item $\varphi_{197} :\equiv (\mathrm{ContOn}(f,[a,b]) \to \varphi_{128}) \to \mathrm{ContOn}(f,[a,b]) \to a \in \mathbb{R} \to \varphi_{161}$
\item $\varphi_{198} :\equiv (\mathrm{ContOn}(f,[a,b]) \to \varphi_{141}) \to \mathrm{ContOn}(f,[a,b]) \to a \in \mathbb{R} \to b \in \mathbb{R} \to \varphi_{127}$
\item $\varphi_{199} :\equiv (\mathrm{ContOn}(f,[a,b]) \to \varphi_{130}) \to \mathrm{ContOn}(f,[a,b]) \to a \in \mathbb{R} \to \varphi_{163}$
\item $\varphi_{200} :\equiv (\mathrm{ContOn}(f,[a,b]) \to \varphi_{143}) \to \mathrm{ContOn}(f,[a,b]) \to a \in \mathbb{R} \to b \in \mathbb{R} \to \varphi_{131}$
\item $\varphi_{201} :\equiv (\mathrm{ContOn}(f,[a,b]) \to \varphi_{136}) \to \mathrm{ContOn}(f,[a,b]) \to a \in \mathbb{R} \to b \in \mathbb{R} \to \varphi_{149}$
\item $\varphi_{202} :\equiv (\mathrm{ContOn}(f,[a,b]) \to \varphi_{151}) \to \mathrm{ContOn}(f,[a,b]) \to a \in \mathbb{R} \to b \in \mathbb{R} \to \varphi_{135}$
\item $\varphi_{203} :\equiv (\mathrm{ContOn}(f,[a,b]) \to \varphi_{162}) \to \mathrm{ContOn}(f,[a,b]) \to a \in \mathbb{R} \to b \in \mathbb{R} \to \varphi_{129}$
\item $\varphi_{204} :\equiv \varphi_{118} \to \varphi_{125} \to \mathrm{ContOn}(f,[a,b]) \to a \in \mathbb{R} \to b \in \mathbb{R} \to \varphi_{47}$
\item $\varphi_{205} :\equiv (\mathrm{ContOn}(f,[a,b]) \to \varphi_{26}) \to (\mathrm{ContOn}(f,[a,b]) \to a \in \mathbb{R} \to a \in \mathbb{R}) \to \varphi_{171}$
\item $\varphi_{206} :\equiv \varphi_{118} \to \varphi_{126} \to \mathrm{ContOn}(f,[a,b]) \to a \in \mathbb{R} \to b \in \mathbb{R} \to \varphi_{53}$
\item $\varphi_{207} :\equiv (\mathrm{ContOn}(f,[a,b]) \to a \in \mathbb{R} \to b \in \mathbb{R} \to \varphi_{145}) \to \varphi_{111} \to \neg \varphi_{68}$
\item $\varphi_{208} :\equiv (\forall u\,(\varphi_{93}) \to \varphi_{3} \to \forall u \in [a,b]\,(u \in \mathrm{dom}((-1 f)))) \to \varphi_{99} \to \mathrm{ContOn}(f,[a,b]) \to a \in \mathbb{R} \to \varphi_{28}$
\item $\varphi_{209} :\equiv (1 \in \mathbb{R} \to f(x) \in \mathbb{R} \to 1 \cdot f(x) = f(x) \cdot 1) \to 1 \in \mathbb{R} \to \varphi_{192}$
\item $\varphi_{210} :\equiv (\mathrm{ContOn}(f,[a,b]) \to \varphi_{152}) \to \varphi_{50} \to \mathrm{ContOn}(f,[a,b]) \to \varphi_{153}$
\item $\varphi_{211} :\equiv (1 \in \mathbb{R} \to f(\xi) \in \mathbb{R} \to 1 \cdot f(\xi) = f(\xi) \cdot 1) \to 1 \in \mathbb{R} \to \varphi_{195}$
\item $\varphi_{212} :\equiv (\mathrm{ContOn}(f,[a,b]) \to \varphi_{158}) \to \varphi_{57} \to \mathrm{ContOn}(f,[a,b]) \to \varphi_{159}$
\item $\varphi_{213} :\equiv (1 \in \mathbb{R} \to f(x) \in \mathbb{R} \to -1 \cdot f(x) = -(1 \cdot f(x))) \to 1 \in \mathbb{R} \to \varphi_{198}$
\item $\varphi_{214} :\equiv (1 \in \mathbb{R} \to f(\xi) \in \mathbb{R} \to -1 \cdot f(\xi) = -(1 \cdot f(\xi))) \to 1 \in \mathbb{R} \to \varphi_{200}$
\item $\varphi_{215} :\equiv ([a,b] \subseteq \mathrm{dom}(f) \to \xi \in [a,b] \to \xi \in \mathrm{dom}(f)) \to \varphi_{45} \to \varphi_{148} \to \mathrm{ContOn}(f,[a,b]) \to \varphi_{130}$
\item $\varphi_{216} :\equiv ([a,b] \subseteq \mathrm{dom}(f) \to x \in [a,b] \to x \in \mathrm{dom}(f)) \to \varphi_{45} \to \varphi_{183}$
\item $\varphi_{217} :\equiv (\mathrm{ContOn}(f,[a,b]) \to a \in \mathbb{R} \to \varphi_{147}) \to \varphi_{109} \to \mathrm{ContOn}(f,[a,b]) \to a \in \mathbb{R} \to \varphi_{155}$
\item $\varphi_{218} :\equiv (\mathrm{ContOn}(f,[a,b]) \to \varphi_{153}) \to -(1 \cdot f(x)) = -(1 \cdot f(x)) \to \varphi_{175}$
\item $\varphi_{219} :\equiv (\mathrm{ContOn}(f,[a,b]) \to a \in \mathbb{R} \to \varphi_{150}) \to \varphi_{113} \to \mathrm{ContOn}(f,[a,b]) \to a \in \mathbb{R} \to \varphi_{160}$
\item $\varphi_{220} :\equiv (\mathrm{ContOn}(f,[a,b]) \to a \in \mathbb{R} \to b \in \mathbb{R} \to \varphi_{127}) \to \varphi_{180}$
\item $\varphi_{221} :\equiv (\mathrm{ContOn}(f,[a,b]) \to \varphi_{159}) \to -(1 \cdot f(\xi)) = -(1 \cdot f(\xi)) \to \varphi_{177}$
\item $\varphi_{222} :\equiv (\mathrm{ContOn}(f,[a,b]) \to a \in \mathbb{R} \to b \in \mathbb{R} \to \varphi_{135}) \to \varphi_{115} \to \mathrm{ContOn}(f,[a,b]) \to a \in \mathbb{R} \to b \in \mathbb{R} \to \varphi_{145}$
\item $\varphi_{223} :\equiv (\mathrm{ContOn}(f,[a,b]) \to a \in \mathbb{R} \to b \in \mathbb{R} \to \varphi_{131}) \to \varphi_{181}$
\item $\varphi_{224} :\equiv (\mathrm{ContOn}(f,[a,b]) \to a \in \mathbb{R} \to \varphi_{138}) \to \varphi_{61} \to \varphi_{51} \to \mathrm{ContOn}(f,[a,b]) \to \varphi_{144}$
\item $\varphi_{225} :\equiv (\mathrm{ContOn}(f,[a,b]) \to a \in \mathbb{R} \to b \in \mathbb{R} \to \varphi_{47}) \to (\mathrm{ContOn}(f,[a,b]) \to a \in \mathbb{R} \to \varphi_{13}) \to \varphi_{197}$
\item $\varphi_{226} :\equiv (\mathrm{ContOn}(f,[a,b]) \to a \in \mathbb{R} \to b \in \mathbb{R} \to \varphi_{129}) \to (-(-f(\xi)) = f(\xi) \to \varphi_{46}) \to \varphi_{176}$
\item $\varphi_{227} :\equiv (\mathrm{ContOn}(f,[a,b]) \to a \in \mathbb{R} \to b \in \mathbb{R} \to \varphi_{53}) \to (\mathrm{ContOn}(f,[a,b]) \to a \in \mathbb{R} \to \varphi_{18}) \to \varphi_{199}$
\item $\varphi_{228} :\equiv (\mathrm{ContOn}(f,[a,b]) \to a \in \mathbb{R} \to \varphi_{157}) \to \varphi_{108} \to \varphi_{100} \to \mathrm{ContOn}(f,[a,b]) \to \varphi_{162}$
\item $\varphi_{229} :\equiv (\mathrm{ContOn}(f,[a,b]) \to a \in \mathbb{R} \to b \in \mathbb{R} \to \varphi_{149}) \to ((-1 f)(x) = -f(x) \to \varphi_{84}) \to \varphi_{179}$
\item $\varphi_{230} :\equiv (\mathrm{ContOn}(f,[a,b]) \to a \in \mathbb{R} \to b \in \mathbb{R} \to \varphi_{137}) \to \varphi_{217}$
\item $\varphi_{231} :\equiv (\mathrm{ContOn}(f,[a,b]) \to a \in \mathbb{R} \to b \in \mathbb{R} \to \varphi_{142}) \to \varphi_{219}$
\item $\varphi_{232} :\equiv (\mathrm{ContOn}(f,[a,b]) \to \varphi_{154}) \to (\mathrm{ContOn}(f,[a,b]) \to \varphi_{156}) \to \varphi_{203}$
\item $\varphi_{233} :\equiv \forall \xi\,(\varphi_{111} \to \neg \varphi_{68}) \to (\forall \xi\,(\varphi_{111} \to \neg \varphi_{68}) \to \forall \xi\,(\varphi_{111}) \to \forall \xi\,(\neg \varphi_{68})) \to \varphi_{186}$
\item $\varphi_{234} :\equiv \forall x\,(\varphi_{121} \to x \in [a,b] \to f(\xi) \leq f(x)) \to \varphi_{194} \to (\varphi_{121} \to \forall x\,(\varphi_{121})) \to \mathrm{ContOn}(f,[a,b]) \to \varphi_{151}$
\end{abbrevlist}
\end{fullproof}

\section{第4章 2節 2.4: 逆関数 (定義2.20, 定理2.21)}
% 自動生成 (prover/build.js): SubLeAnti  ノード数 198
\begin{fullproof}{[E51]}{$\forall z\,\forall x\,\forall y\,(z \in \mathbb{R} \to x \in \mathbb{R} \to y \in \mathbb{R} \to x \leq y \to z - y \leq z - x)$}{198}\phantomsection\label{tree:SubLeAnti}
\noindent$\Pi_{1}$:
\begin{ptree}
\AXC{$z_{1} = w \to z - y < z_{1} \to z - y < w$\lbl{A6}}
\UIC{$\forall w\,(z_{1} = w \to z - y < z_{1} \to z - y < w)$}
\UIC{$\forall z_{1}\,\forall w\,(z_{1} = w \to z - y < z_{1} \to z - y < w)$}
\AXC{$\varphi_{29}$\lbl{A2}}
\BIC{$\forall w\,(-x + z = w \to z - y < -x + z \to z - y < w)$}
\AXC{$\varphi_{42}$\lbl{A2}}
\BIC{$-x + z = z - x \to z - y < -x + z \to z - y < z - x$}
\end{ptree}
\noindent$\Pi_{2}$:
\begin{ptree}
\AXC{$\forall x\,\forall y\,(x \in \mathbb{R} \to y \in \mathbb{R} \to x \leq y \to -y \leq -x)$\lbl{[F23]}}
\AXC{$\varphi_{31}$\lbl{A2}}
\BIC{$\forall y\,(x \in \mathbb{R} \to y \in \mathbb{R} \to x \leq y \to -y \leq -x)$}
\AXC{$\varphi_{28}$\lbl{A2}}
\BIC{$x \in \mathbb{R} \to y \in \mathbb{R} \to x \leq y \to -y \leq -x$}
\AXC{$(x \in \mathbb{R} \to y \in \mathbb{R} \to x \leq y \to -y \leq -x) \to \varphi_{56}$\lbl{A1}}
\BIC{$\varphi_{56}$}
\end{ptree}
\noindent$\Pi_{3}$:
\begin{ptree}
\AXC{$z \in \mathbb{R} \to x \in \mathbb{R} \to y \in \mathbb{R} \to x \leq y \to y \in \mathbb{R}$\lbl{A1}}
\AXC{$z \in \mathbb{R} \to x \in \mathbb{R} \to y \in \mathbb{R} \to x \leq y \to x \in \mathbb{R}$\lbl{A1}}
\AXC{$\stackrel{\Pi_{2}}{\vdots}$}\noLine
\UIC{$\varphi_{56}$}
\BIC{$(z \in \mathbb{R} \to x \in \mathbb{R} \to y \in \mathbb{R} \to x \leq y \to y \in \mathbb{R}) \to \varphi_{37}$}
\BIC{$\varphi_{37}$}
\end{ptree}
\noindent$\Pi_{4}$:
\begin{ptree}
\AXC{$\forall x \in \mathbb{R}\,(-x \in \mathbb{R})$\lbl{(R1)}}
\AXC{$\forall x \in \mathbb{R}\,(-x \in \mathbb{R}) \to x \in \mathbb{R} \to -x \in \mathbb{R}$\lbl{A2}}
\BIC{$x \in \mathbb{R} \to -x \in \mathbb{R}$}
\AXC{$(x \in \mathbb{R} \to -x \in \mathbb{R}) \to \varphi_{43}$\lbl{A1}}
\BIC{$\varphi_{43}$}
\end{ptree}
\noindent$\Pi_{5}$:
\begin{ptree}
\AXC{$z \in \mathbb{R} \to x \in \mathbb{R} \to y \in \mathbb{R} \to x \leq y \to x \in \mathbb{R}$\lbl{A1}}
\AXC{$\stackrel{\Pi_{4}}{\vdots}$}\noLine
\UIC{$\varphi_{43}$}
\BIC{$z \in \mathbb{R} \to x \in \mathbb{R} \to y \in \mathbb{R} \to x \leq y \to -x \in \mathbb{R}$}
\end{ptree}
\noindent$\Pi_{6}$:
\begin{ptree}
\AXC{$\forall x \in \mathbb{R}\,(-x \in \mathbb{R})$\lbl{(R1)}}
\AXC{$\forall x \in \mathbb{R}\,(-x \in \mathbb{R}) \to y \in \mathbb{R} \to -y \in \mathbb{R}$\lbl{A2}}
\BIC{$y \in \mathbb{R} \to -y \in \mathbb{R}$}
\AXC{$(y \in \mathbb{R} \to -y \in \mathbb{R}) \to \varphi_{44}$\lbl{A1}}
\BIC{$\varphi_{44}$}
\end{ptree}
\noindent$\Pi_{7}$:
\begin{ptree}
\AXC{$z \in \mathbb{R} \to x \in \mathbb{R} \to y \in \mathbb{R} \to x \leq y \to y \in \mathbb{R}$\lbl{A1}}
\AXC{$\stackrel{\Pi_{6}}{\vdots}$}\noLine
\UIC{$\varphi_{44}$}
\BIC{$z \in \mathbb{R} \to x \in \mathbb{R} \to y \in \mathbb{R} \to x \leq y \to -y \in \mathbb{R}$}
\end{ptree}
\noindent$\Pi_{8}$:
\begin{ptree}
\AXC{$\varphi_{15}$\lbl{[F19]}}
\AXC{$\varphi_{15} \to \varphi_{21}$\lbl{A2}}
\BIC{$\varphi_{21}$}
\AXC{$\varphi_{21} \to \varphi_{23}$\lbl{A2}}
\BIC{$\varphi_{23}$}
\AXC{$\varphi_{23} \to v_{879} \in \mathbb{R} \to \varphi_{20}$\lbl{A2}}
\BIC{$v_{879} \in \mathbb{R} \to \varphi_{20}$}
\UIC{$\forall v_{881}\,(v_{879} \in \mathbb{R} \to \varphi_{20})$}
\UIC{$\forall v_{880}\,\forall v_{881}\,(v_{879} \in \mathbb{R} \to \varphi_{20})$}
\UIC{$\forall v_{879}\,\forall v_{880}\,\forall v_{881}\,(v_{879} \in \mathbb{R} \to \varphi_{20})$}
\AXC{$\forall v_{879}\,\forall v_{880}\,\forall v_{881}\,(v_{879} \in \mathbb{R} \to \varphi_{20}) \to \varphi_{32}$\lbl{A2}}
\BIC{$\varphi_{32}$}
\end{ptree}
\noindent$\Pi_{9}$:
\begin{ptree}
\AXC{$\stackrel{\Pi_{8}}{\vdots}$}\noLine
\UIC{$\varphi_{32}$}
\AXC{$\varphi_{32} \to \varphi_{33}$\lbl{A2}}
\BIC{$\varphi_{33}$}
\AXC{$\varphi_{33} \to \varphi_{27}$\lbl{A2}}
\BIC{$\varphi_{27}$}
\AXC{$\varphi_{27} \to (z \in \mathbb{R} \to x \in \mathbb{R} \to y \in \mathbb{R} \to x \leq y \to -y \in \mathbb{R}) \to \varphi_{60}$\lbl{A1}}
\BIC{$(z \in \mathbb{R} \to x \in \mathbb{R} \to y \in \mathbb{R} \to x \leq y \to -y \in \mathbb{R}) \to \varphi_{60}$}
\end{ptree}
\noindent$\Pi_{10}$:
\begin{ptree}
\AXC{$\stackrel{\Pi_{7}}{\vdots}$}\noLine
\UIC{$z \in \mathbb{R} \to x \in \mathbb{R} \to y \in \mathbb{R} \to x \leq y \to -y \in \mathbb{R}$}
\AXC{$\stackrel{\Pi_{9}}{\vdots}$}\noLine
\UIC{$(z \in \mathbb{R} \to x \in \mathbb{R} \to y \in \mathbb{R} \to x \leq y \to -y \in \mathbb{R}) \to \varphi_{60}$}
\BIC{$\varphi_{60}$}
\end{ptree}
\noindent$\Pi_{11}$:
\begin{ptree}
\AXC{$z \in \mathbb{R} \to x \in \mathbb{R} \to y \in \mathbb{R} \to x \leq y \to z \in \mathbb{R}$\lbl{A1}}
\AXC{$\stackrel{\Pi_{5}}{\vdots}$}\noLine
\UIC{$z \in \mathbb{R} \to x \in \mathbb{R} \to y \in \mathbb{R} \to x \leq y \to -x \in \mathbb{R}$}
\AXC{$\stackrel{\Pi_{10}}{\vdots}$}\noLine
\UIC{$\varphi_{60}$}
\BIC{$(z \in \mathbb{R} \to x \in \mathbb{R} \to y \in \mathbb{R} \to x \leq y \to z \in \mathbb{R}) \to \varphi_{48}$}
\BIC{$\varphi_{48}$}
\end{ptree}
\noindent$\Pi_{12}$:
\begin{ptree}
\AXC{$y = z \to x = y \to x = z$\lbl{A6}}
\AXC{$(y = z \to x = y \to x = z) \to x = y \to y = z \to x = z$\lbl{A1}}
\BIC{$x = y \to y = z \to x = z$}
\UIC{$\forall z\,(x = y \to y = z \to x = z)$}
\UIC{$\forall y\,\forall z\,(x = y \to y = z \to x = z)$}
\UIC{$\forall x\,\forall y\,\forall z\,(x = y \to y = z \to x = z)$}
\AXC{$\varphi_{4}$\lbl{A2}}
\BIC{$\forall y\,\forall z\,(v_{886} = y \to y = z \to v_{886} = z)$}
\AXC{$\varphi_{8}$\lbl{A2}}
\BIC{$\forall z\,(v_{886} = v_{887} \to v_{887} = z \to v_{886} = z)$}
\AXC{$\varphi_{13}$\lbl{A2}}
\BIC{$v_{886} = v_{887} \to v_{887} = v_{888} \to v_{886} = v_{888}$}
\UIC{$\forall v_{888}\,(v_{886} = v_{887} \to v_{887} = v_{888} \to v_{886} = v_{888})$}
\UIC{$\varphi_{2}$}
\UIC{$\forall v_{886}\,(\varphi_{2})$}
\end{ptree}
\noindent$\Pi_{13}$:
\begin{ptree}
\AXC{$\stackrel{\Pi_{12}}{\vdots}$}\noLine
\UIC{$\forall v_{886}\,(\varphi_{2})$}
\AXC{$\forall v_{886}\,(\varphi_{2}) \to \varphi_{6}$\lbl{A2}}
\BIC{$\varphi_{6}$}
\AXC{$\varphi_{6} \to \varphi_{16}$\lbl{A2}}
\BIC{$\varphi_{16}$}
\AXC{$\varphi_{16} \to -y + z = z + -y \to z + -y = z - y \to -y + z = z - y$\lbl{A2}}
\BIC{$-y + z = z + -y \to z + -y = z - y \to -y + z = z - y$}
\end{ptree}
\noindent$\Pi_{14}$:
\begin{ptree}
\AXC{$x = x$\lbl{A5}}
\AXC{$x = y \to x = x \to y = x$\lbl{A6}}
\AXC{$(x = y \to x = x \to y = x) \to x = x \to x = y \to y = x$\lbl{A1}}
\BIC{$x = x \to x = y \to y = x$}
\BIC{$x = y \to y = x$}
\UIC{$\forall y\,(x = y \to y = x)$}
\UIC{$\forall x\,\forall y\,(x = y \to y = x)$}
\end{ptree}
\noindent$\Pi_{15}$:
\begin{ptree}
\AXC{$\stackrel{\Pi_{14}}{\vdots}$}\noLine
\UIC{$\forall x\,\forall y\,(x = y \to y = x)$}
\AXC{$\forall x\,\forall y\,(x = y \to y = x) \to \forall y\,(v_{884} = y \to y = v_{884})$\lbl{A2}}
\BIC{$\forall y\,(v_{884} = y \to y = v_{884})$}
\AXC{$\varphi_{1}$\lbl{A2}}
\BIC{$v_{884} = v_{885} \to v_{885} = v_{884}$}
\UIC{$\forall v_{885}\,(v_{884} = v_{885} \to v_{885} = v_{884})$}
\UIC{$\forall v_{884}\,\forall v_{885}\,(v_{884} = v_{885} \to v_{885} = v_{884})$}
\AXC{$\varphi_{17}$\lbl{A2}}
\BIC{$\forall v_{885}\,(z - y = v_{885} \to v_{885} = z - y)$}
\AXC{$\varphi_{22}$\lbl{A2}}
\BIC{$z - y = z + -y \to z + -y = z - y$}
\end{ptree}
\noindent$\Pi_{16}$:
\begin{ptree}
\AXC{$\forall x\,\forall y\,(x - y = x + -y)$\lbl{Def$-$}}
\AXC{$\forall x\,\forall y\,(x - y = x + -y) \to \forall y\,(z - y = z + -y)$\lbl{A2}}
\BIC{$\forall y\,(z - y = z + -y)$}
\end{ptree}
\noindent$\Pi_{17}$:
\begin{ptree}
\AXC{$\stackrel{\Pi_{15}}{\vdots}$}\noLine
\UIC{$z - y = z + -y \to z + -y = z - y$}
\AXC{$\stackrel{\Pi_{16}}{\vdots}$}\noLine
\UIC{$\forall y\,(z - y = z + -y)$}
\AXC{$\forall y\,(z - y = z + -y) \to z - y = z + -y$\lbl{A2}}
\BIC{$z - y = z + -y$}
\AXC{$\varphi_{24}$\lbl{A1}}
\BIC{$(z - y = z + -y \to z + -y = z - y) \to z + -y = z - y$}
\BIC{$z + -y = z - y$}
\end{ptree}
\noindent$\Pi_{18}$:
\begin{ptree}
\AXC{$\forall x\,\forall y\,(x \in \mathbb{R} \to y \in \mathbb{R} \to x + y = y + x)$\lbl{(R3)}}
\AXC{$\varphi_{30}$\lbl{A2}}
\BIC{$\forall y\,(v_{882} \in \mathbb{R} \to y \in \mathbb{R} \to v_{882} + y = y + v_{882})$}
\AXC{$\varphi_{38}$\lbl{A2}}
\BIC{$v_{882} \in \mathbb{R} \to v_{883} \in \mathbb{R} \to v_{882} + v_{883} = v_{883} + v_{882}$}
\UIC{$\varphi_{10}$}
\UIC{$\forall v_{882}\,(\varphi_{10})$}
\AXC{$\forall v_{882}\,(\varphi_{10}) \to \varphi_{19}$\lbl{A2}}
\BIC{$\varphi_{19}$}
\end{ptree}
\noindent$\Pi_{19}$:
\begin{ptree}
\AXC{$\stackrel{\Pi_{7}}{\vdots}$}\noLine
\UIC{$z \in \mathbb{R} \to x \in \mathbb{R} \to y \in \mathbb{R} \to x \leq y \to -y \in \mathbb{R}$}
\AXC{$\stackrel{\Pi_{18}}{\vdots}$}\noLine
\UIC{$\varphi_{19}$}
\AXC{$\varphi_{19} \to -y \in \mathbb{R} \to z \in \mathbb{R} \to -y + z = z + -y$\lbl{A2}}
\BIC{$-y \in \mathbb{R} \to z \in \mathbb{R} \to -y + z = z + -y$}
\AXC{$\varphi_{57}$\lbl{A1}}
\BIC{$(z \in \mathbb{R} \to x \in \mathbb{R} \to y \in \mathbb{R} \to x \leq y \to -y \in \mathbb{R}) \to \varphi_{46}$}
\BIC{$\varphi_{46}$}
\end{ptree}
\noindent$\Pi_{20}$:
\begin{ptree}
\AXC{$\stackrel{\Pi_{17}}{\vdots}$}\noLine
\UIC{$z + -y = z - y$}
\AXC{$z \in \mathbb{R} \to x \in \mathbb{R} \to y \in \mathbb{R} \to x \leq y \to z \in \mathbb{R}$\lbl{A1}}
\AXC{$\stackrel{\Pi_{19}}{\vdots}$}\noLine
\UIC{$\varphi_{46}$}
\BIC{$z \in \mathbb{R} \to x \in \mathbb{R} \to y \in \mathbb{R} \to x \leq y \to -y + z = z + -y$}
\AXC{$\varphi_{54}$\lbl{A1}}
\BIC{$z + -y = z - y \to \varphi_{49}$}
\BIC{$\varphi_{49}$}
\end{ptree}
\noindent$\Pi_{21}$:
\begin{ptree}
\AXC{$\stackrel{\Pi_{13}}{\vdots}$}\noLine
\UIC{$-y + z = z + -y \to z + -y = z - y \to -y + z = z - y$}
\AXC{$\stackrel{\Pi_{20}}{\vdots}$}\noLine
\UIC{$\varphi_{49}$}
\BIC{$z \in \mathbb{R} \to x \in \mathbb{R} \to y \in \mathbb{R} \to x \leq y \to -y + z = z - y$}
\AXC{$\varphi_{63}$\lbl{A1}}
\BIC{$\varphi_{61}$}
\end{ptree}
\noindent$\Pi_{22}$:
\begin{ptree}
\AXC{$z \in \mathbb{R} \to x \in \mathbb{R} \to y \in \mathbb{R} \to x \leq y \to x \leq y$\lbl{A1}}
\AXC{$\stackrel{\Pi_{3}}{\vdots}$}\noLine
\UIC{$\varphi_{37}$}
\BIC{$z \in \mathbb{R} \to x \in \mathbb{R} \to y \in \mathbb{R} \to x \leq y \to -y \leq -x$}
\AXC{$\stackrel{\Pi_{11}}{\vdots}$}\noLine
\UIC{$\varphi_{48}$}
\BIC{$z \in \mathbb{R} \to x \in \mathbb{R} \to y \in \mathbb{R} \to x \leq y \to -y + z \leq -x + z$}
\AXC{$\stackrel{\Pi_{21}}{\vdots}$}\noLine
\UIC{$\varphi_{61}$}
\BIC{$\varphi_{53}$}
\end{ptree}
\noindent$\Pi_{23}$:
\begin{ptree}
\AXC{$z_{1} = w \to z_{1} < -x + z \to w < -x + z$\lbl{A6}}
\UIC{$\forall w\,(z_{1} = w \to z_{1} < -x + z \to w < -x + z)$}
\UIC{$\forall z_{1}\,\forall w\,(z_{1} = w \to z_{1} < -x + z \to w < -x + z)$}
\AXC{$\varphi_{41}$\lbl{A2}}
\BIC{$\forall w\,(-y + z = w \to -y + z < -x + z \to w < -x + z)$}
\AXC{$\varphi_{52}$\lbl{A2}}
\BIC{$-y + z = z - y \to -y + z < -x + z \to z - y < -x + z$}
\AXC{$\stackrel{\Pi_{22}}{\vdots}$}\noLine
\UIC{$\varphi_{53}$}
\BIC{$\varphi_{50}$}
\end{ptree}
\noindent$\Pi_{24}$:
\begin{ptree}
\AXC{$z_{1} = w \to z_{1} = -x + z \to w = -x + z$\lbl{A6}}
\UIC{$\forall w\,(z_{1} = w \to z_{1} = -x + z \to w = -x + z)$}
\UIC{$\forall z_{1}\,\forall w\,(z_{1} = w \to z_{1} = -x + z \to w = -x + z)$}
\AXC{$\varphi_{34}$\lbl{A2}}
\BIC{$\forall w\,(-y + z = w \to -y + z = -x + z \to w = -x + z)$}
\AXC{$\varphi_{45}$\lbl{A2}}
\BIC{$-y + z = z - y \to -y + z = -x + z \to z - y = -x + z$}
\AXC{$\stackrel{\Pi_{23}}{\vdots}$}\noLine
\UIC{$\varphi_{50}$}
\BIC{$z \in \mathbb{R} \to x \in \mathbb{R} \to y \in \mathbb{R} \to x \leq y \to z - y \leq -x + z$}
\end{ptree}
\noindent$\Pi_{25}$:
\begin{ptree}
\AXC{$y = z \to x = y \to x = z$\lbl{A6}}
\AXC{$(y = z \to x = y \to x = z) \to x = y \to y = z \to x = z$\lbl{A1}}
\BIC{$x = y \to y = z \to x = z$}
\UIC{$\forall z\,(x = y \to y = z \to x = z)$}
\UIC{$\forall y\,\forall z\,(x = y \to y = z \to x = z)$}
\UIC{$\forall x\,\forall y\,\forall z\,(x = y \to y = z \to x = z)$}
\AXC{$\varphi_{5}$\lbl{A2}}
\BIC{$\forall y\,\forall z\,(v_{889} = y \to y = z \to v_{889} = z)$}
\AXC{$\varphi_{9}$\lbl{A2}}
\BIC{$\forall z\,(v_{889} = v_{890} \to v_{890} = z \to v_{889} = z)$}
\AXC{$\varphi_{14}$\lbl{A2}}
\BIC{$v_{889} = v_{890} \to v_{890} = v_{891} \to v_{889} = v_{891}$}
\UIC{$\forall v_{891}\,(v_{889} = v_{890} \to v_{890} = v_{891} \to v_{889} = v_{891})$}
\UIC{$\varphi_{3}$}
\UIC{$\forall v_{889}\,(\varphi_{3})$}
\end{ptree}
\noindent$\Pi_{26}$:
\begin{ptree}
\AXC{$\stackrel{\Pi_{25}}{\vdots}$}\noLine
\UIC{$\forall v_{889}\,(\varphi_{3})$}
\AXC{$\forall v_{889}\,(\varphi_{3}) \to \varphi_{7}$\lbl{A2}}
\BIC{$\varphi_{7}$}
\AXC{$\varphi_{7} \to \varphi_{18}$\lbl{A2}}
\BIC{$\varphi_{18}$}
\AXC{$\varphi_{18} \to -x + z = z + -x \to z + -x = z - x \to -x + z = z - x$\lbl{A2}}
\BIC{$-x + z = z + -x \to z + -x = z - x \to -x + z = z - x$}
\end{ptree}
\noindent$\Pi_{27}$:
\begin{ptree}
\AXC{$\stackrel{\Pi_{14}}{\vdots}$}\noLine
\UIC{$\forall x\,\forall y\,(x = y \to y = x)$}
\AXC{$\forall x\,\forall y\,(x = y \to y = x) \to \forall y\,(z - x = y \to y = z - x)$\lbl{A2}}
\BIC{$\forall y\,(z - x = y \to y = z - x)$}
\AXC{$\varphi_{11}$\lbl{A2}}
\BIC{$z - x = z + -x \to z + -x = z - x$}
\end{ptree}
\noindent$\Pi_{28}$:
\begin{ptree}
\AXC{$\forall x\,\forall y\,(x - y = x + -y)$\lbl{Def$-$}}
\AXC{$\forall x\,\forall y\,(x - y = x + -y) \to \forall y\,(z - y = z + -y)$\lbl{A2}}
\BIC{$\forall y\,(z - y = z + -y)$}
\end{ptree}
\noindent$\Pi_{29}$:
\begin{ptree}
\AXC{$\stackrel{\Pi_{27}}{\vdots}$}\noLine
\UIC{$z - x = z + -x \to z + -x = z - x$}
\AXC{$\stackrel{\Pi_{28}}{\vdots}$}\noLine
\UIC{$\forall y\,(z - y = z + -y)$}
\AXC{$\forall y\,(z - y = z + -y) \to z - x = z + -x$\lbl{A2}}
\BIC{$z - x = z + -x$}
\AXC{$\varphi_{25}$\lbl{A1}}
\BIC{$(z - x = z + -x \to z + -x = z - x) \to z + -x = z - x$}
\BIC{$z + -x = z - x$}
\end{ptree}
\noindent$\Pi_{30}$:
\begin{ptree}
\AXC{$\forall x\,\forall y\,(x \in \mathbb{R} \to y \in \mathbb{R} \to x + y = y + x)$\lbl{(R3)}}
\AXC{$\varphi_{35}$\lbl{A2}}
\BIC{$\forall y\,(-x \in \mathbb{R} \to y \in \mathbb{R} \to -x + y = y + -x)$}
\AXC{$\varphi_{39}$\lbl{A2}}
\BIC{$-x \in \mathbb{R} \to z \in \mathbb{R} \to -x + z = z + -x$}
\AXC{$\varphi_{58}$\lbl{A1}}
\BIC{$(z \in \mathbb{R} \to x \in \mathbb{R} \to y \in \mathbb{R} \to x \leq y \to -x \in \mathbb{R}) \to \varphi_{47}$}
\end{ptree}
\noindent$\Pi_{31}$:
\begin{ptree}
\AXC{$\stackrel{\Pi_{5}}{\vdots}$}\noLine
\UIC{$z \in \mathbb{R} \to x \in \mathbb{R} \to y \in \mathbb{R} \to x \leq y \to -x \in \mathbb{R}$}
\AXC{$\stackrel{\Pi_{30}}{\vdots}$}\noLine
\UIC{$(z \in \mathbb{R} \to x \in \mathbb{R} \to y \in \mathbb{R} \to x \leq y \to -x \in \mathbb{R}) \to \varphi_{47}$}
\BIC{$\varphi_{47}$}
\end{ptree}
\noindent$\Pi_{32}$:
\begin{ptree}
\AXC{$\stackrel{\Pi_{29}}{\vdots}$}\noLine
\UIC{$z + -x = z - x$}
\AXC{$z \in \mathbb{R} \to x \in \mathbb{R} \to y \in \mathbb{R} \to x \leq y \to z \in \mathbb{R}$\lbl{A1}}
\AXC{$\stackrel{\Pi_{31}}{\vdots}$}\noLine
\UIC{$\varphi_{47}$}
\BIC{$z \in \mathbb{R} \to x \in \mathbb{R} \to y \in \mathbb{R} \to x \leq y \to -x + z = z + -x$}
\AXC{$\varphi_{55}$\lbl{A1}}
\BIC{$z + -x = z - x \to \varphi_{51}$}
\BIC{$\varphi_{51}$}
\end{ptree}
\noindent$\Pi_{33}$:
\begin{ptree}
\AXC{$\stackrel{\Pi_{26}}{\vdots}$}\noLine
\UIC{$-x + z = z + -x \to z + -x = z - x \to -x + z = z - x$}
\AXC{$\stackrel{\Pi_{32}}{\vdots}$}\noLine
\UIC{$\varphi_{51}$}
\BIC{$z \in \mathbb{R} \to x \in \mathbb{R} \to y \in \mathbb{R} \to x \leq y \to -x + z = z - x$}
\AXC{$\varphi_{62}$\lbl{A1}}
\BIC{$\varphi_{59}$}
\end{ptree}
\noindent$\Pi_{34}$:
\begin{ptree}
\AXC{$\stackrel{\Pi_{1}}{\vdots}$}\noLine
\UIC{$-x + z = z - x \to z - y < -x + z \to z - y < z - x$}
\AXC{$\stackrel{\Pi_{24}}{\vdots}$}\noLine
\UIC{$z \in \mathbb{R} \to x \in \mathbb{R} \to y \in \mathbb{R} \to x \leq y \to z - y \leq -x + z$}
\AXC{$\stackrel{\Pi_{33}}{\vdots}$}\noLine
\UIC{$\varphi_{59}$}
\BIC{$(-x + z = z - x \to z - y < -x + z \to z - y < z - x) \to \varphi_{40}$}
\BIC{$\varphi_{40}$}
\end{ptree}
\noindent 主木:
\begin{ptree}
\AXC{$z_{1} = w \to z - y = z_{1} \to z - y = w$\lbl{A6}}
\UIC{$\forall w\,(z_{1} = w \to z - y = z_{1} \to z - y = w)$}
\UIC{$\forall z_{1}\,\forall w\,(z_{1} = w \to z - y = z_{1} \to z - y = w)$}
\AXC{$\varphi_{26}$\lbl{A2}}
\BIC{$\forall w\,(-x + z = w \to z - y = -x + z \to z - y = w)$}
\AXC{$\varphi_{36}$\lbl{A2}}
\BIC{$-x + z = z - x \to z - y = -x + z \to z - y = z - x$}
\AXC{$\stackrel{\Pi_{34}}{\vdots}$}\noLine
\UIC{$\varphi_{40}$}
\BIC{$z \in \mathbb{R} \to x \in \mathbb{R} \to y \in \mathbb{R} \to x \leq y \to z - y \leq z - x$}
\UIC{$\forall y\,(z \in \mathbb{R} \to x \in \mathbb{R} \to y \in \mathbb{R} \to x \leq y \to z - y \leq z - x)$}
\UIC{$\varphi_{12}$}
\UIC{$\forall z\,(\varphi_{12})$}
\end{ptree}
\begin{abbrevlist}
\item $\varphi_{1} :\equiv \forall y\,(v_{884} = y \to y = v_{884}) \to v_{884} = v_{885} \to v_{885} = v_{884}$
\item $\varphi_{2} :\equiv \forall v_{887}\,\forall v_{888}\,(v_{886} = v_{887} \to v_{887} = v_{888} \to v_{886} = v_{888})$
\item $\varphi_{3} :\equiv \forall v_{890}\,\forall v_{891}\,(v_{889} = v_{890} \to v_{890} = v_{891} \to v_{889} = v_{891})$
\item $\varphi_{4} :\equiv \forall x\,\forall y\,\forall z\,(x = y \to y = z \to x = z) \to \forall y\,\forall z\,(v_{886} = y \to y = z \to v_{886} = z)$
\item $\varphi_{5} :\equiv \forall x\,\forall y\,\forall z\,(x = y \to y = z \to x = z) \to \forall y\,\forall z\,(v_{889} = y \to y = z \to v_{889} = z)$
\item $\varphi_{6} :\equiv \forall v_{887}\,\forall v_{888}\,(-y + z = v_{887} \to v_{887} = v_{888} \to -y + z = v_{888})$
\item $\varphi_{7} :\equiv \forall v_{890}\,\forall v_{891}\,(-x + z = v_{890} \to v_{890} = v_{891} \to -x + z = v_{891})$
\item $\varphi_{8} :\equiv \forall y\,\forall z\,(v_{886} = y \to y = z \to v_{886} = z) \to \forall z\,(v_{886} = v_{887} \to v_{887} = z \to v_{886} = z)$
\item $\varphi_{9} :\equiv \forall y\,\forall z\,(v_{889} = y \to y = z \to v_{889} = z) \to \forall z\,(v_{889} = v_{890} \to v_{890} = z \to v_{889} = z)$
\item $\varphi_{10} :\equiv \forall v_{883}\,(v_{882} \in \mathbb{R} \to v_{883} \in \mathbb{R} \to v_{882} + v_{883} = v_{883} + v_{882})$
\item $\varphi_{11} :\equiv \forall y\,(z - x = y \to y = z - x) \to z - x = z + -x \to z + -x = z - x$
\item $\varphi_{12} :\equiv \forall x\,\forall y\,(z \in \mathbb{R} \to x \in \mathbb{R} \to y \in \mathbb{R} \to x \leq y \to z - y \leq z - x)$
\item $\varphi_{13} :\equiv \forall z\,(v_{886} = v_{887} \to v_{887} = z \to v_{886} = z) \to v_{886} = v_{887} \to v_{887} = v_{888} \to v_{886} = v_{888}$
\item $\varphi_{14} :\equiv \forall z\,(v_{889} = v_{890} \to v_{890} = z \to v_{889} = z) \to v_{889} = v_{890} \to v_{890} = v_{891} \to v_{889} = v_{891}$
\item $\varphi_{15} :\equiv \forall x\,\forall y\,\forall z\,(x \in \mathbb{R} \to y \in \mathbb{R} \to z \in \mathbb{R} \to x \leq y \to x + z \leq y + z)$
\item $\varphi_{16} :\equiv \forall v_{888}\,(-y + z = z + -y \to z + -y = v_{888} \to -y + z = v_{888})$
\item $\varphi_{17} :\equiv \forall v_{884}\,\forall v_{885}\,(v_{884} = v_{885} \to v_{885} = v_{884}) \to \forall v_{885}\,(z - y = v_{885} \to v_{885} = z - y)$
\item $\varphi_{18} :\equiv \forall v_{891}\,(-x + z = z + -x \to z + -x = v_{891} \to -x + z = v_{891})$
\item $\varphi_{19} :\equiv \forall v_{883}\,(-y \in \mathbb{R} \to v_{883} \in \mathbb{R} \to -y + v_{883} = v_{883} + -y)$
\item $\varphi_{20} :\equiv v_{880} \in \mathbb{R} \to v_{881} \in \mathbb{R} \to v_{879} \leq v_{880} \to v_{879} + v_{881} \leq v_{880} + v_{881}$
\item $\varphi_{21} :\equiv \forall y\,\forall z\,(v_{879} \in \mathbb{R} \to y \in \mathbb{R} \to z \in \mathbb{R} \to v_{879} \leq y \to v_{879} + z \leq y + z)$
\item $\varphi_{22} :\equiv \forall v_{885}\,(z - y = v_{885} \to v_{885} = z - y) \to z - y = z + -y \to z + -y = z - y$
\item $\varphi_{23} :\equiv \forall z\,(v_{879} \in \mathbb{R} \to v_{880} \in \mathbb{R} \to z \in \mathbb{R} \to v_{879} \leq v_{880} \to v_{879} + z \leq v_{880} + z)$
\item $\varphi_{24} :\equiv z - y = z + -y \to (z - y = z + -y \to z + -y = z - y) \to z + -y = z - y$
\item $\varphi_{25} :\equiv z - x = z + -x \to (z - x = z + -x \to z + -x = z - x) \to z + -x = z - x$
\item $\varphi_{26} :\equiv \forall z_{1}\,\forall w\,(z_{1} = w \to z - y = z_{1} \to z - y = w) \to \forall w\,(-x + z = w \to z - y = -x + z \to z - y = w)$
\item $\varphi_{27} :\equiv -y \in \mathbb{R} \to -x \in \mathbb{R} \to z \in \mathbb{R} \to -y \leq -x \to -y + z \leq -x + z$
\item $\varphi_{28} :\equiv \forall y\,(x \in \mathbb{R} \to y \in \mathbb{R} \to x \leq y \to -y \leq -x) \to x \in \mathbb{R} \to y \in \mathbb{R} \to x \leq y \to -y \leq -x$
\item $\varphi_{29} :\equiv \forall z_{1}\,\forall w\,(z_{1} = w \to z - y < z_{1} \to z - y < w) \to \forall w\,(-x + z = w \to z - y < -x + z \to z - y < w)$
\item $\varphi_{30} :\equiv \forall x\,\forall y\,(x \in \mathbb{R} \to y \in \mathbb{R} \to x + y = y + x) \to \forall y\,(v_{882} \in \mathbb{R} \to y \in \mathbb{R} \to v_{882} + y = y + v_{882})$
\item $\varphi_{31} :\equiv \forall x\,\forall y\,(x \in \mathbb{R} \to y \in \mathbb{R} \to x \leq y \to -y \leq -x) \to \forall y\,(x \in \mathbb{R} \to y \in \mathbb{R} \to x \leq y \to -y \leq -x)$
\item $\varphi_{32} :\equiv \forall v_{880}\,\forall v_{881}\,(-y \in \mathbb{R} \to v_{880} \in \mathbb{R} \to v_{881} \in \mathbb{R} \to -y \leq v_{880} \to -y + v_{881} \leq v_{880} + v_{881})$
\item $\varphi_{33} :\equiv \forall v_{881}\,(-y \in \mathbb{R} \to -x \in \mathbb{R} \to v_{881} \in \mathbb{R} \to -y \leq -x \to -y + v_{881} \leq -x + v_{881})$
\item $\varphi_{34} :\equiv \forall z_{1}\,\forall w\,(z_{1} = w \to z_{1} = -x + z \to w = -x + z) \to \forall w\,(-y + z = w \to -y + z = -x + z \to w = -x + z)$
\item $\varphi_{35} :\equiv \forall x\,\forall y\,(x \in \mathbb{R} \to y \in \mathbb{R} \to x + y = y + x) \to \forall y\,(-x \in \mathbb{R} \to y \in \mathbb{R} \to -x + y = y + -x)$
\item $\varphi_{36} :\equiv \forall w\,(-x + z = w \to z - y = -x + z \to z - y = w) \to -x + z = z - x \to z - y = -x + z \to z - y = z - x$
\item $\varphi_{37} :\equiv (z \in \mathbb{R} \to x \in \mathbb{R} \to y \in \mathbb{R} \to x \leq y \to x \leq y) \to z \in \mathbb{R} \to x \in \mathbb{R} \to y \in \mathbb{R} \to x \leq y \to -y \leq -x$
\item $\varphi_{38} :\equiv \forall y\,(v_{882} \in \mathbb{R} \to y \in \mathbb{R} \to v_{882} + y = y + v_{882}) \to v_{882} \in \mathbb{R} \to v_{883} \in \mathbb{R} \to v_{882} + v_{883} = v_{883} + v_{882}$
\item $\varphi_{39} :\equiv \forall y\,(-x \in \mathbb{R} \to y \in \mathbb{R} \to -x + y = y + -x) \to -x \in \mathbb{R} \to z \in \mathbb{R} \to -x + z = z + -x$
\item $\varphi_{40} :\equiv (-x + z = z - x \to z - y = -x + z \to z - y = z - x) \to z \in \mathbb{R} \to x \in \mathbb{R} \to y \in \mathbb{R} \to x \leq y \to z - y \leq z - x$
\item $\varphi_{41} :\equiv \forall z_{1}\,\forall w\,(z_{1} = w \to z_{1} < -x + z \to w < -x + z) \to \forall w\,(-y + z = w \to -y + z < -x + z \to w < -x + z)$
\item $\varphi_{42} :\equiv \forall w\,(-x + z = w \to z - y < -x + z \to z - y < w) \to -x + z = z - x \to z - y < -x + z \to z - y < z - x$
\item $\varphi_{43} :\equiv (z \in \mathbb{R} \to x \in \mathbb{R} \to y \in \mathbb{R} \to x \leq y \to x \in \mathbb{R}) \to z \in \mathbb{R} \to x \in \mathbb{R} \to y \in \mathbb{R} \to x \leq y \to -x \in \mathbb{R}$
\item $\varphi_{44} :\equiv (z \in \mathbb{R} \to x \in \mathbb{R} \to y \in \mathbb{R} \to x \leq y \to y \in \mathbb{R}) \to z \in \mathbb{R} \to x \in \mathbb{R} \to y \in \mathbb{R} \to x \leq y \to -y \in \mathbb{R}$
\item $\varphi_{45} :\equiv \forall w\,(-y + z = w \to -y + z = -x + z \to w = -x + z) \to -y + z = z - y \to -y + z = -x + z \to z - y = -x + z$
\item $\varphi_{46} :\equiv (z \in \mathbb{R} \to x \in \mathbb{R} \to y \in \mathbb{R} \to x \leq y \to z \in \mathbb{R}) \to z \in \mathbb{R} \to x \in \mathbb{R} \to y \in \mathbb{R} \to x \leq y \to -y + z = z + -y$
\item $\varphi_{47} :\equiv (z \in \mathbb{R} \to x \in \mathbb{R} \to y \in \mathbb{R} \to x \leq y \to z \in \mathbb{R}) \to z \in \mathbb{R} \to x \in \mathbb{R} \to y \in \mathbb{R} \to x \leq y \to -x + z = z + -x$
\item $\varphi_{48} :\equiv (z \in \mathbb{R} \to x \in \mathbb{R} \to y \in \mathbb{R} \to x \leq y \to -y \leq -x) \to z \in \mathbb{R} \to x \in \mathbb{R} \to y \in \mathbb{R} \to x \leq y \to -y + z \leq -x + z$
\item $\varphi_{49} :\equiv (-y + z = z + -y \to z + -y = z - y \to -y + z = z - y) \to z \in \mathbb{R} \to x \in \mathbb{R} \to y \in \mathbb{R} \to x \leq y \to -y + z = z - y$
\item $\varphi_{50} :\equiv (-y + z = z - y \to -y + z = -x + z \to z - y = -x + z) \to z \in \mathbb{R} \to x \in \mathbb{R} \to y \in \mathbb{R} \to x \leq y \to z - y \leq -x + z$
\item $\varphi_{51} :\equiv (-x + z = z + -x \to z + -x = z - x \to -x + z = z - x) \to z \in \mathbb{R} \to x \in \mathbb{R} \to y \in \mathbb{R} \to x \leq y \to -x + z = z - x$
\item $\varphi_{52} :\equiv \forall w\,(-y + z = w \to -y + z < -x + z \to w < -x + z) \to -y + z = z - y \to -y + z < -x + z \to z - y < -x + z$
\item $\varphi_{53} :\equiv (-y + z = z - y \to -y + z < -x + z \to z - y < -x + z) \to \varphi_{50}$
\item $\varphi_{54} :\equiv (z \in \mathbb{R} \to x \in \mathbb{R} \to y \in \mathbb{R} \to x \leq y \to -y + z = z + -y) \to z + -y = z - y \to \varphi_{49}$
\item $\varphi_{55} :\equiv (z \in \mathbb{R} \to x \in \mathbb{R} \to y \in \mathbb{R} \to x \leq y \to -x + z = z + -x) \to z + -x = z - x \to \varphi_{51}$
\item $\varphi_{56} :\equiv (z \in \mathbb{R} \to x \in \mathbb{R} \to y \in \mathbb{R} \to x \leq y \to x \in \mathbb{R}) \to (z \in \mathbb{R} \to x \in \mathbb{R} \to y \in \mathbb{R} \to x \leq y \to y \in \mathbb{R}) \to \varphi_{37}$
\item $\varphi_{57} :\equiv (-y \in \mathbb{R} \to z \in \mathbb{R} \to -y + z = z + -y) \to (z \in \mathbb{R} \to x \in \mathbb{R} \to y \in \mathbb{R} \to x \leq y \to -y \in \mathbb{R}) \to \varphi_{46}$
\item $\varphi_{58} :\equiv (-x \in \mathbb{R} \to z \in \mathbb{R} \to -x + z = z + -x) \to (z \in \mathbb{R} \to x \in \mathbb{R} \to y \in \mathbb{R} \to x \leq y \to -x \in \mathbb{R}) \to \varphi_{47}$
\item $\varphi_{59} :\equiv (z \in \mathbb{R} \to x \in \mathbb{R} \to y \in \mathbb{R} \to x \leq y \to z - y \leq -x + z) \to (-x + z = z - x \to z - y < -x + z \to z - y < z - x) \to \varphi_{40}$
\item $\varphi_{60} :\equiv (z \in \mathbb{R} \to x \in \mathbb{R} \to y \in \mathbb{R} \to x \leq y \to -x \in \mathbb{R}) \to (z \in \mathbb{R} \to x \in \mathbb{R} \to y \in \mathbb{R} \to x \leq y \to z \in \mathbb{R}) \to \varphi_{48}$
\item $\varphi_{61} :\equiv (z \in \mathbb{R} \to x \in \mathbb{R} \to y \in \mathbb{R} \to x \leq y \to -y + z \leq -x + z) \to \varphi_{53}$
\item $\varphi_{62} :\equiv (z \in \mathbb{R} \to x \in \mathbb{R} \to y \in \mathbb{R} \to x \leq y \to -x + z = z - x) \to \varphi_{59}$
\item $\varphi_{63} :\equiv (z \in \mathbb{R} \to x \in \mathbb{R} \to y \in \mathbb{R} \to x \leq y \to -y + z = z - y) \to \varphi_{61}$
\end{abbrevlist}
\end{fullproof}

% 自動生成 (prover/build.js): StrIncInj  ノード数 167
\begin{fullproof}{[I1]}{$\forall f\,(\mathrm{RFn}(f) \to \mathrm{StrInc}(f,\mathrm{dom}(f)) \to \mathrm{Inj}(f))$}{167}\phantomsection\label{tree:StrIncInj}
\noindent$\Pi_{1}$:
\begin{ptree}
\AXC{$\forall f\,(\varphi_{25})$\lbl{DefStrInc}}
\AXC{$\forall f\,(\varphi_{25}) \to \varphi_{25}$\lbl{A2}}
\BIC{$\varphi_{25}$}
\AXC{$\varphi_{25} \to (\mathrm{StrInc}(f,\mathrm{dom}(f)) \leftrightarrow \forall x\,\forall y\,(\varphi_{10}))$\lbl{A2}}
\BIC{$\mathrm{StrInc}(f,\mathrm{dom}(f)) \leftrightarrow \forall x\,\forall y\,(\varphi_{10})$}
\end{ptree}
\noindent$\Pi_{2}$:
\begin{ptree}
\AXC{$\stackrel{\Pi_{1}}{\vdots}$}\noLine
\UIC{$\mathrm{StrInc}(f,\mathrm{dom}(f)) \leftrightarrow \forall x\,\forall y\,(\varphi_{10})$}
\AXC{$\mathrm{RFn}(f) \to \varphi_{1}$\lbl{A1}}
\AXC{$(\mathrm{RFn}(f) \to \varphi_{1}) \to \varphi_{74}$\lbl{A1}}
\BIC{$\varphi_{74}$}
\BIC{$\mathrm{RFn}(f) \to \mathrm{StrInc}(f,\mathrm{dom}(f)) \to \varphi_{54}$}
\end{ptree}
\noindent$\Pi_{3}$:
\begin{ptree}
\AXC{$\stackrel{\Pi_{2}}{\vdots}$}\noLine
\UIC{$\mathrm{RFn}(f) \to \mathrm{StrInc}(f,\mathrm{dom}(f)) \to \varphi_{54}$}
\AXC{$(\mathrm{RFn}(f) \to \mathrm{StrInc}(f,\mathrm{dom}(f)) \to \varphi_{54}) \to \varphi_{81}$\lbl{A1}}
\BIC{$\varphi_{81}$}
\end{ptree}
\noindent$\Pi_{4}$:
\begin{ptree}
\AXC{$\varphi_{23} \to v_{892} \in \mathrm{dom}(f) \to \varphi_{8}$\lbl{A2}}
\AXC{$\forall x\,\forall y\,(\varphi_{10}) \to \varphi_{23}$\lbl{A2}}
\AXC{$\stackrel{\Pi_{3}}{\vdots}$}\noLine
\UIC{$\varphi_{81}$}
\BIC{$\mathrm{RFn}(f) \to \mathrm{StrInc}(f,\mathrm{dom}(f)) \to \varphi_{56}$}
\AXC{$\varphi_{91}$\lbl{A1}}
\BIC{$(\varphi_{23} \to v_{892} \in \mathrm{dom}(f) \to \varphi_{8}) \to \varphi_{63}$}
\BIC{$\varphi_{63}$}
\AXC{$\varphi_{63} \to \varphi_{30} \to v_{892} \in \mathrm{dom}(f) \to \varphi_{8}$\lbl{A1}}
\BIC{$\varphi_{30} \to v_{892} \in \mathrm{dom}(f) \to \varphi_{8}$}
\UIC{$\forall v_{893}\,(\varphi_{30} \to v_{892} \in \mathrm{dom}(f) \to \varphi_{8})$}
\end{ptree}
\noindent$\Pi_{5}$:
\begin{ptree}
\AXC{$\varphi_{30} \to \forall v_{893}\,(\varphi_{30})$\lbl{A4}}
\AXC{$\varphi_{85}$\lbl{A3}}
\AXC{$\stackrel{\Pi_{4}}{\vdots}$}\noLine
\UIC{$\forall v_{893}\,(\varphi_{30} \to v_{892} \in \mathrm{dom}(f) \to \varphi_{8})$}
\AXC{$\varphi_{94}$\lbl{A1}}
\BIC{$\varphi_{85} \to (\varphi_{30} \to \forall v_{893}\,(\varphi_{30})) \to \mathrm{RFn}(f) \to \varphi_{62}$}
\BIC{$(\varphi_{30} \to \forall v_{893}\,(\varphi_{30})) \to \mathrm{RFn}(f) \to \varphi_{62}$}
\BIC{$\mathrm{RFn}(f) \to \varphi_{62}$}
\AXC{$\varphi_{84}$\lbl{A1}}
\BIC{$\varphi_{30} \to \forall v_{893}\,(v_{892} \in \mathrm{dom}(f) \to \varphi_{8})$}
\UIC{$\forall v_{892}\,(\varphi_{30} \to \forall v_{893}\,(v_{892} \in \mathrm{dom}(f) \to \varphi_{8}))$}
\AXC{$\varphi_{95}$\lbl{A1}}
\BIC{$\varphi_{87} \to (\varphi_{30} \to \forall v_{892}\,(\varphi_{30})) \to \mathrm{RFn}(f) \to \varphi_{65}$}
\end{ptree}
\noindent$\Pi_{6}$:
\begin{ptree}
\AXC{$\varphi_{30} \to \forall v_{892}\,(\varphi_{30})$\lbl{A4}}
\AXC{$\varphi_{87}$\lbl{A3}}
\AXC{$\stackrel{\Pi_{5}}{\vdots}$}\noLine
\UIC{$\varphi_{87} \to (\varphi_{30} \to \forall v_{892}\,(\varphi_{30})) \to \mathrm{RFn}(f) \to \varphi_{65}$}
\BIC{$(\varphi_{30} \to \forall v_{892}\,(\varphi_{30})) \to \mathrm{RFn}(f) \to \varphi_{65}$}
\BIC{$\mathrm{RFn}(f) \to \varphi_{65}$}
\AXC{$(\mathrm{RFn}(f) \to \varphi_{65}) \to \varphi_{82}$\lbl{A1}}
\BIC{$\varphi_{82}$}
\end{ptree}
\noindent$\Pi_{7}$:
\begin{ptree}
\AXC{$\forall v_{892}\,\forall v_{893}\,(v_{892} \in \mathrm{dom}(f) \to \varphi_{8}) \to \varphi_{28}$\lbl{A2}}
\AXC{$\stackrel{\Pi_{6}}{\vdots}$}\noLine
\UIC{$\varphi_{82}$}
\BIC{$\mathrm{RFn}(f) \to \mathrm{StrInc}(f,\mathrm{dom}(f)) \to \varphi_{57}$}
\AXC{$(\mathrm{RFn}(f) \to \mathrm{StrInc}(f,\mathrm{dom}(f)) \to \varphi_{57}) \to \varphi_{80}$\lbl{A1}}
\BIC{$\varphi_{80}$}
\end{ptree}
\noindent$\Pi_{8}$:
\begin{ptree}
\AXC{$\mathrm{RFn}(f) \to \varphi_{3}$\lbl{A1}}
\AXC{$\mathrm{RFn}(f) \to \varphi_{13}$\lbl{A1}}
\AXC{$\varphi_{28} \to \varphi_{7}$\lbl{A2}}
\AXC{$\stackrel{\Pi_{7}}{\vdots}$}\noLine
\UIC{$\varphi_{80}$}
\BIC{$\mathrm{RFn}(f) \to \mathrm{StrInc}(f,\mathrm{dom}(f)) \to \varphi_{49}$}
\AXC{$\varphi_{92}$\lbl{A1}}
\BIC{$(\mathrm{RFn}(f) \to \varphi_{13}) \to \varphi_{76}$}
\BIC{$\varphi_{76}$}
\BIC{$(\mathrm{RFn}(f) \to \mathrm{StrInc}(f,\mathrm{dom}(f)) \to \varphi_{20}) \to \varphi_{39}$}
\end{ptree}
\noindent$\Pi_{9}$:
\begin{ptree}
\AXC{$\mathrm{RFn}(f) \to \mathrm{StrInc}(f,\mathrm{dom}(f)) \to \varphi_{20}$\lbl{A1}}
\AXC{$\stackrel{\Pi_{8}}{\vdots}$}\noLine
\UIC{$(\mathrm{RFn}(f) \to \mathrm{StrInc}(f,\mathrm{dom}(f)) \to \varphi_{20}) \to \varphi_{39}$}
\BIC{$\varphi_{39}$}
\end{ptree}
\noindent$\Pi_{10}$:
\begin{ptree}
\AXC{$\stackrel{\Pi_{9}}{\vdots}$}\noLine
\UIC{$\varphi_{39}$}
\AXC{$\mathrm{RFn}(f) \to \mathrm{StrInc}(f,\mathrm{dom}(f)) \to \varphi_{22}$\lbl{A1}}
\AXC{$\varphi_{88}$\lbl{A1}}
\BIC{$\varphi_{39} \to \varphi_{67}$}
\BIC{$\varphi_{67}$}
\end{ptree}
\noindent$\Pi_{11}$:
\begin{ptree}
\AXC{$z = w \to f(y) < z \to f(y) < w$\lbl{A6}}
\UIC{$\forall w\,(z = w \to f(y) < z \to f(y) < w)$}
\UIC{$\forall z\,\forall w\,(z = w \to f(y) < z \to f(y) < w)$}
\AXC{$\varphi_{35}$\lbl{A2}}
\BIC{$\forall w\,(f(x) = w \to f(y) < f(x) \to f(y) < w)$}
\AXC{$\varphi_{46}$\lbl{A2}}
\BIC{$f(x) = f(y) \to f(y) < f(x) \to f(y) < f(y)$}
\AXC{$\stackrel{\Pi_{10}}{\vdots}$}\noLine
\UIC{$\varphi_{67}$}
\BIC{$\varphi_{40}$}
\AXC{$\varphi_{40} \to \varphi_{44}$\lbl{A1}}
\BIC{$\varphi_{44}$}
\end{ptree}
\noindent$\Pi_{12}$:
\begin{ptree}
\AXC{$\forall f\,(\varphi_{25})$\lbl{DefStrInc}}
\AXC{$\forall f\,(\varphi_{25}) \to \varphi_{25}$\lbl{A2}}
\BIC{$\varphi_{25}$}
\AXC{$\varphi_{25} \to (\mathrm{StrInc}(f,\mathrm{dom}(f)) \leftrightarrow \forall x\,\forall y\,(\varphi_{10}))$\lbl{A2}}
\BIC{$\mathrm{StrInc}(f,\mathrm{dom}(f)) \leftrightarrow \forall x\,\forall y\,(\varphi_{10})$}
\end{ptree}
\noindent$\Pi_{13}$:
\begin{ptree}
\AXC{$\stackrel{\Pi_{12}}{\vdots}$}\noLine
\UIC{$\mathrm{StrInc}(f,\mathrm{dom}(f)) \leftrightarrow \forall x\,\forall y\,(\varphi_{10})$}
\AXC{$\mathrm{RFn}(f) \to \varphi_{1}$\lbl{A1}}
\AXC{$(\mathrm{RFn}(f) \to \varphi_{1}) \to \varphi_{75}$\lbl{A1}}
\BIC{$\varphi_{75}$}
\BIC{$\mathrm{RFn}(f) \to \mathrm{StrInc}(f,\mathrm{dom}(f)) \to \varphi_{55}$}
\end{ptree}
\noindent$\Pi_{14}$:
\begin{ptree}
\AXC{$\stackrel{\Pi_{13}}{\vdots}$}\noLine
\UIC{$\mathrm{RFn}(f) \to \mathrm{StrInc}(f,\mathrm{dom}(f)) \to \varphi_{55}$}
\AXC{$(\mathrm{RFn}(f) \to \mathrm{StrInc}(f,\mathrm{dom}(f)) \to \varphi_{55}) \to \varphi_{79}$\lbl{A1}}
\BIC{$\varphi_{79}$}
\end{ptree}
\noindent$\Pi_{15}$:
\begin{ptree}
\AXC{$\forall y\,(\varphi_{10}) \to \varphi_{10}$\lbl{A2}}
\AXC{$\forall x\,\forall y\,(\varphi_{10}) \to \forall y\,(\varphi_{10})$\lbl{A2}}
\AXC{$\stackrel{\Pi_{14}}{\vdots}$}\noLine
\UIC{$\varphi_{79}$}
\BIC{$\mathrm{RFn}(f) \to \mathrm{StrInc}(f,\mathrm{dom}(f)) \to \varphi_{52}$}
\AXC{$(\mathrm{RFn}(f) \to \mathrm{StrInc}(f,\mathrm{dom}(f)) \to \varphi_{52}) \to \varphi_{77}$\lbl{A1}}
\BIC{$\varphi_{77}$}
\BIC{$\mathrm{RFn}(f) \to \mathrm{StrInc}(f,\mathrm{dom}(f)) \to \varphi_{50}$}
\end{ptree}
\noindent$\Pi_{16}$:
\begin{ptree}
\AXC{$\mathrm{RFn}(f) \to \varphi_{13}$\lbl{A1}}
\AXC{$\mathrm{RFn}(f) \to \varphi_{3}$\lbl{A1}}
\AXC{$\stackrel{\Pi_{15}}{\vdots}$}\noLine
\UIC{$\mathrm{RFn}(f) \to \mathrm{StrInc}(f,\mathrm{dom}(f)) \to \varphi_{50}$}
\AXC{$\varphi_{93}$\lbl{A1}}
\BIC{$(\mathrm{RFn}(f) \to \varphi_{3}) \to \varphi_{78}$}
\BIC{$\varphi_{78}$}
\BIC{$(\mathrm{RFn}(f) \to \mathrm{StrInc}(f,\mathrm{dom}(f)) \to \varphi_{21}) \to \varphi_{41}$}
\end{ptree}
\noindent$\Pi_{17}$:
\begin{ptree}
\AXC{$\mathrm{RFn}(f) \to \mathrm{StrInc}(f,\mathrm{dom}(f)) \to \varphi_{21}$\lbl{A1}}
\AXC{$\stackrel{\Pi_{16}}{\vdots}$}\noLine
\UIC{$(\mathrm{RFn}(f) \to \mathrm{StrInc}(f,\mathrm{dom}(f)) \to \varphi_{21}) \to \varphi_{41}$}
\BIC{$\varphi_{41}$}
\end{ptree}
\noindent$\Pi_{18}$:
\begin{ptree}
\AXC{$\stackrel{\Pi_{17}}{\vdots}$}\noLine
\UIC{$\varphi_{41}$}
\AXC{$\mathrm{RFn}(f) \to \mathrm{StrInc}(f,\mathrm{dom}(f)) \to \varphi_{22}$\lbl{A1}}
\AXC{$\varphi_{89}$\lbl{A1}}
\BIC{$\varphi_{41} \to \varphi_{68}$}
\BIC{$\varphi_{68}$}
\end{ptree}
\noindent$\Pi_{19}$:
\begin{ptree}
\AXC{$z = w \to z < f(y) \to w < f(y)$\lbl{A6}}
\UIC{$\forall w\,(z = w \to z < f(y) \to w < f(y))$}
\UIC{$\forall z\,\forall w\,(z = w \to z < f(y) \to w < f(y))$}
\AXC{$\varphi_{36}$\lbl{A2}}
\BIC{$\forall w\,(f(x) = w \to f(x) < f(y) \to w < f(y))$}
\AXC{$\varphi_{47}$\lbl{A2}}
\BIC{$f(x) = f(y) \to f(x) < f(y) \to f(y) < f(y)$}
\AXC{$\stackrel{\Pi_{18}}{\vdots}$}\noLine
\UIC{$\varphi_{68}$}
\BIC{$\varphi_{42}$}
\AXC{$\varphi_{42} \to \varphi_{45}$\lbl{A1}}
\BIC{$\varphi_{45}$}
\end{ptree}
\noindent$\Pi_{20}$:
\begin{ptree}
\AXC{$\mathrm{RFn}(f) \to \mathrm{RFn}(f)$\lbl{A1}}
\AXC{$\forall f\,\forall x\,(\mathrm{RFn}(f) \to x \in \mathrm{dom}(f) \to x \in \mathbb{R})$\lbl{[S7]}}
\AXC{$\varphi_{29}$\lbl{A2}}
\BIC{$\forall x\,(\mathrm{RFn}(f) \to x \in \mathrm{dom}(f) \to x \in \mathbb{R})$}
\AXC{$\varphi_{26}$\lbl{A2}}
\BIC{$\mathrm{RFn}(f) \to y \in \mathrm{dom}(f) \to y \in \mathbb{R}$}
\AXC{$\varphi_{73}$\lbl{A1}}
\BIC{$(\mathrm{RFn}(f) \to \mathrm{RFn}(f)) \to \varphi_{61}$}
\BIC{$\varphi_{61}$}
\end{ptree}
\noindent$\Pi_{21}$:
\begin{ptree}
\AXC{$\mathrm{RFn}(f) \to \mathrm{RFn}(f)$\lbl{A1}}
\AXC{$\forall f\,\forall x\,(\mathrm{RFn}(f) \to x \in \mathrm{dom}(f) \to x \in \mathbb{R})$\lbl{[S7]}}
\AXC{$\varphi_{29}$\lbl{A2}}
\BIC{$\forall x\,(\mathrm{RFn}(f) \to x \in \mathrm{dom}(f) \to x \in \mathbb{R})$}
\AXC{$\varphi_{27}$\lbl{A2}}
\BIC{$\mathrm{RFn}(f) \to x \in \mathrm{dom}(f) \to x \in \mathbb{R}$}
\AXC{$\varphi_{72}$\lbl{A1}}
\BIC{$(\mathrm{RFn}(f) \to \mathrm{RFn}(f)) \to \varphi_{60}$}
\BIC{$\varphi_{60}$}
\end{ptree}
\noindent$\Pi_{22}$:
\begin{ptree}
\AXC{$\forall x\,\forall y\,(x \in \mathbb{R} \to y \in \mathbb{R} \to x < y \lor x = y \lor y < x)$\lbl{(R13)}}
\AXC{$\varphi_{38}$\lbl{A2}}
\BIC{$\forall y\,(x \in \mathbb{R} \to y \in \mathbb{R} \to x < y \lor x = y \lor y < x)$}
\AXC{$\varphi_{37}$\lbl{A2}}
\BIC{$x \in \mathbb{R} \to y \in \mathbb{R} \to x < y \lor x = y \lor y < x$}
\AXC{$\varphi_{83}$\lbl{A1}}
\BIC{$(\mathrm{RFn}(f) \to \mathrm{StrInc}(f,\mathrm{dom}(f)) \to \varphi_{12}) \to \varphi_{71}$}
\end{ptree}
\noindent$\Pi_{23}$:
\begin{ptree}
\AXC{$\mathrm{RFn}(f) \to \varphi_{3}$\lbl{A1}}
\AXC{$\stackrel{\Pi_{21}}{\vdots}$}\noLine
\UIC{$\varphi_{60}$}
\BIC{$\mathrm{RFn}(f) \to \mathrm{StrInc}(f,\mathrm{dom}(f)) \to \varphi_{12}$}
\AXC{$\stackrel{\Pi_{22}}{\vdots}$}\noLine
\UIC{$(\mathrm{RFn}(f) \to \mathrm{StrInc}(f,\mathrm{dom}(f)) \to \varphi_{12}) \to \varphi_{71}$}
\BIC{$\varphi_{71}$}
\end{ptree}
\noindent$\Pi_{24}$:
\begin{ptree}
\AXC{$\mathrm{RFn}(f) \to \varphi_{13}$\lbl{A1}}
\AXC{$\stackrel{\Pi_{20}}{\vdots}$}\noLine
\UIC{$\varphi_{61}$}
\BIC{$\mathrm{RFn}(f) \to \mathrm{StrInc}(f,\mathrm{dom}(f)) \to \varphi_{11}$}
\AXC{$\stackrel{\Pi_{23}}{\vdots}$}\noLine
\UIC{$\varphi_{71}$}
\BIC{$\mathrm{RFn}(f) \to \mathrm{StrInc}(f,\mathrm{dom}(f)) \to \varphi_{24}$}
\AXC{$(\mathrm{RFn}(f) \to \mathrm{StrInc}(f,\mathrm{dom}(f)) \to \varphi_{24}) \to \varphi_{86}$\lbl{A1}}
\BIC{$\varphi_{86}$}
\end{ptree}
\noindent$\Pi_{25}$:
\begin{ptree}
\AXC{$\forall x\,(\neg x < x)$\lbl{[O1]}}
\AXC{$\forall x\,(\neg x < x) \to \neg f(y) < f(y)$\lbl{A2}}
\BIC{$\neg f(y) < f(y)$}
\AXC{$\stackrel{\Pi_{19}}{\vdots}$}\noLine
\UIC{$\varphi_{45}$}
\BIC{$\mathrm{RFn}(f) \to \mathrm{StrInc}(f,\mathrm{dom}(f)) \to \varphi_{15}$}
\AXC{$\stackrel{\Pi_{24}}{\vdots}$}\noLine
\UIC{$\varphi_{86}$}
\BIC{$(\mathrm{RFn}(f) \to \mathrm{StrInc}(f,\mathrm{dom}(f)) \to \varphi_{9}) \to \varphi_{70}$}
\end{ptree}
\noindent$\Pi_{26}$:
\begin{ptree}
\AXC{$\mathrm{RFn}(f) \to \mathrm{StrInc}(f,\mathrm{dom}(f)) \to \varphi_{9}$\lbl{A1}}
\AXC{$\stackrel{\Pi_{25}}{\vdots}$}\noLine
\UIC{$(\mathrm{RFn}(f) \to \mathrm{StrInc}(f,\mathrm{dom}(f)) \to \varphi_{9}) \to \varphi_{70}$}
\BIC{$\varphi_{70}$}
\end{ptree}
\noindent$\Pi_{27}$:
\begin{ptree}
\AXC{$\forall x\,(\neg x < x)$\lbl{[O1]}}
\AXC{$\forall x\,(\neg x < x) \to \neg f(y) < f(y)$\lbl{A2}}
\BIC{$\neg f(y) < f(y)$}
\AXC{$\stackrel{\Pi_{11}}{\vdots}$}\noLine
\UIC{$\varphi_{44}$}
\BIC{$\mathrm{RFn}(f) \to \mathrm{StrInc}(f,\mathrm{dom}(f)) \to \varphi_{14}$}
\AXC{$\stackrel{\Pi_{26}}{\vdots}$}\noLine
\UIC{$\varphi_{70}$}
\BIC{$\mathrm{RFn}(f) \to \mathrm{StrInc}(f,\mathrm{dom}(f)) \to \varphi_{4}$}
\AXC{$\varphi_{64}$\lbl{A1}}
\BIC{$\mathrm{RFn}(f) \land \mathrm{StrInc}(f,\mathrm{dom}(f)) \to \varphi_{4}$}
\UIC{$\forall y\,(\mathrm{RFn}(f) \land \mathrm{StrInc}(f,\mathrm{dom}(f)) \to \varphi_{4})$}
\AXC{$\varphi_{90}$\lbl{A1}}
\BIC{$\varphi_{69} \to \varphi_{59}$}
\end{ptree}
\noindent$\Pi_{28}$:
\begin{ptree}
\AXC{$\varphi_{5}$\lbl{A4}}
\AXC{$\varphi_{31} \to \varphi_{34}$\lbl{A3}}
\AXC{$\varphi_{6}$\lbl{A4}}
\AXC{$\varphi_{69}$\lbl{A3}}
\AXC{$\stackrel{\Pi_{27}}{\vdots}$}\noLine
\UIC{$\varphi_{69} \to \varphi_{59}$}
\BIC{$\varphi_{59}$}
\BIC{$\mathrm{RFn}(f) \to \mathrm{StrInc}(f,\mathrm{dom}(f)) \to \forall y\,(\varphi_{4})$}
\AXC{$\varphi_{66}$\lbl{A1}}
\BIC{$\mathrm{RFn}(f) \land \mathrm{StrInc}(f,\mathrm{dom}(f)) \to \forall y\,(\varphi_{4})$}
\UIC{$\varphi_{31}$}
\AXC{$\varphi_{31} \to (\varphi_{31} \to \varphi_{34}) \to \varphi_{5} \to \varphi_{32}$\lbl{A1}}
\BIC{$(\varphi_{31} \to \varphi_{34}) \to \varphi_{5} \to \varphi_{32}$}
\BIC{$\varphi_{5} \to \varphi_{32}$}
\BIC{$\varphi_{32}$}
\end{ptree}
\noindent 主木:
\begin{ptree}
\AXC{$\stackrel{\Pi_{28}}{\vdots}$}\noLine
\UIC{$\varphi_{32}$}
\AXC{$\forall f\,(\mathrm{Inj}(f) \leftrightarrow \forall x\,\forall y\,(\varphi_{4}))$\lbl{DefInj}}
\AXC{$\varphi_{58}$\lbl{A2}}
\BIC{$\mathrm{Inj}(f) \leftrightarrow \forall x\,\forall y\,(\varphi_{4})$}
\AXC{$(\mathrm{Inj}(f) \leftrightarrow \forall x\,\forall y\,(\varphi_{4})) \to \varphi_{53}$\lbl{A1}}
\BIC{$\varphi_{53}$}
\BIC{$\mathrm{RFn}(f) \to \mathrm{StrInc}(f,\mathrm{dom}(f)) \to \mathrm{Inj}(f)$}
\UIC{$\varphi_{2}$}
\end{ptree}
\begin{abbrevlist}
\item $\varphi_{1} :\equiv \mathrm{StrInc}(f,\mathrm{dom}(f)) \to \mathrm{StrInc}(f,\mathrm{dom}(f))$
\item $\varphi_{2} :\equiv \forall f\,(\mathrm{RFn}(f) \to \mathrm{StrInc}(f,\mathrm{dom}(f)) \to \mathrm{Inj}(f))$
\item $\varphi_{3} :\equiv \mathrm{StrInc}(f,\mathrm{dom}(f)) \to x \in \mathrm{dom}(f) \to x \in \mathrm{dom}(f)$
\item $\varphi_{4} :\equiv x \in \mathrm{dom}(f) \to y \in \mathrm{dom}(f) \to f(x) = f(y) \to x = y$
\item $\varphi_{5} :\equiv \mathrm{RFn}(f) \land \mathrm{StrInc}(f,\mathrm{dom}(f)) \to \forall x\,(\mathrm{RFn}(f) \land \mathrm{StrInc}(f,\mathrm{dom}(f)))$
\item $\varphi_{6} :\equiv \mathrm{RFn}(f) \land \mathrm{StrInc}(f,\mathrm{dom}(f)) \to \forall y\,(\mathrm{RFn}(f) \land \mathrm{StrInc}(f,\mathrm{dom}(f)))$
\item $\varphi_{7} :\equiv y \in \mathrm{dom}(f) \to x \in \mathrm{dom}(f) \to y < x \to f(y) < f(x)$
\item $\varphi_{8} :\equiv v_{893} \in \mathrm{dom}(f) \to v_{892} < v_{893} \to f(v_{892}) < f(v_{893})$
\item $\varphi_{9} :\equiv x \in \mathrm{dom}(f) \to y \in \mathrm{dom}(f) \to f(x) = f(y) \to x = y \to x = y$
\item $\varphi_{10} :\equiv x \in \mathrm{dom}(f) \to y \in \mathrm{dom}(f) \to x < y \to f(x) < f(y)$
\item $\varphi_{11} :\equiv x \in \mathrm{dom}(f) \to y \in \mathrm{dom}(f) \to f(x) = f(y) \to y \in \mathbb{R}$
\item $\varphi_{12} :\equiv x \in \mathrm{dom}(f) \to y \in \mathrm{dom}(f) \to f(x) = f(y) \to x \in \mathbb{R}$
\item $\varphi_{13} :\equiv \mathrm{StrInc}(f,\mathrm{dom}(f)) \to x \in \mathrm{dom}(f) \to y \in \mathrm{dom}(f) \to y \in \mathrm{dom}(f)$
\item $\varphi_{14} :\equiv x \in \mathrm{dom}(f) \to y \in \mathrm{dom}(f) \to f(x) = f(y) \to y < x \to x = y$
\item $\varphi_{15} :\equiv x \in \mathrm{dom}(f) \to y \in \mathrm{dom}(f) \to f(x) = f(y) \to x < y \to x = y$
\item $\varphi_{16} :\equiv y \in \mathrm{dom}(f) \to f(x) = f(y) \to y < x \to f(y) < f(x)$
\item $\varphi_{17} :\equiv y \in \mathrm{dom}(f) \to f(x) = f(y) \to y < x \to f(y) < f(y)$
\item $\varphi_{18} :\equiv y \in \mathrm{dom}(f) \to f(x) = f(y) \to x < y \to f(x) < f(y)$
\item $\varphi_{19} :\equiv y \in \mathrm{dom}(f) \to f(x) = f(y) \to x < y \to f(y) < f(y)$
\item $\varphi_{20} :\equiv x \in \mathrm{dom}(f) \to y \in \mathrm{dom}(f) \to f(x) = f(y) \to y < x \to y < x$
\item $\varphi_{21} :\equiv x \in \mathrm{dom}(f) \to y \in \mathrm{dom}(f) \to f(x) = f(y) \to x < y \to x < y$
\item $\varphi_{22} :\equiv x \in \mathrm{dom}(f) \to y \in \mathrm{dom}(f) \to f(x) = f(y) \to f(x) = f(y)$
\item $\varphi_{23} :\equiv \forall y\,(v_{892} \in \mathrm{dom}(f) \to y \in \mathrm{dom}(f) \to v_{892} < y \to f(v_{892}) < f(y))$
\item $\varphi_{24} :\equiv x \in \mathrm{dom}(f) \to y \in \mathrm{dom}(f) \to f(x) = f(y) \to x < y \lor x = y \lor y < x$
\item $\varphi_{25} :\equiv \forall A\,(\mathrm{StrInc}(f,A) \leftrightarrow \forall x\,\forall y\,(x \in A \to y \in A \to x < y \to f(x) < f(y)))$
\item $\varphi_{26} :\equiv \forall x\,(\mathrm{RFn}(f) \to x \in \mathrm{dom}(f) \to x \in \mathbb{R}) \to \mathrm{RFn}(f) \to y \in \mathrm{dom}(f) \to y \in \mathbb{R}$
\item $\varphi_{27} :\equiv \forall x\,(\mathrm{RFn}(f) \to x \in \mathrm{dom}(f) \to x \in \mathbb{R}) \to \mathrm{RFn}(f) \to x \in \mathrm{dom}(f) \to x \in \mathbb{R}$
\item $\varphi_{28} :\equiv \forall v_{893}\,(y \in \mathrm{dom}(f) \to v_{893} \in \mathrm{dom}(f) \to y < v_{893} \to f(y) < f(v_{893}))$
\item $\varphi_{29} :\equiv \forall f\,\forall x\,(\mathrm{RFn}(f) \to x \in \mathrm{dom}(f) \to x \in \mathbb{R}) \to \forall x\,(\mathrm{RFn}(f) \to x \in \mathrm{dom}(f) \to x \in \mathbb{R})$
\item $\varphi_{30} :\equiv \mathrm{RFn}(f) \land \mathrm{StrInc}(f,\mathrm{dom}(f)) \land x \in \mathrm{dom}(f) \land y \in \mathrm{dom}(f) \land f(x) = f(y) \land y < x$
\item $\varphi_{31} :\equiv \forall x\,(\mathrm{RFn}(f) \land \mathrm{StrInc}(f,\mathrm{dom}(f)) \to \forall y\,(\varphi_{4}))$
\item $\varphi_{32} :\equiv \mathrm{RFn}(f) \to \mathrm{StrInc}(f,\mathrm{dom}(f)) \to \forall x\,\forall y\,(\varphi_{4})$
\item $\varphi_{33} :\equiv \forall y\,(\mathrm{RFn}(f) \land \mathrm{StrInc}(f,\mathrm{dom}(f))) \to \forall y\,(\varphi_{4})$
\item $\varphi_{34} :\equiv \forall x\,(\mathrm{RFn}(f) \land \mathrm{StrInc}(f,\mathrm{dom}(f))) \to \forall x\,\forall y\,(\varphi_{4})$
\item $\varphi_{35} :\equiv \forall z\,\forall w\,(z = w \to f(y) < z \to f(y) < w) \to \forall w\,(f(x) = w \to f(y) < f(x) \to f(y) < w)$
\item $\varphi_{36} :\equiv \forall z\,\forall w\,(z = w \to z < f(y) \to w < f(y)) \to \forall w\,(f(x) = w \to f(x) < f(y) \to w < f(y))$
\item $\varphi_{37} :\equiv \forall y\,(x \in \mathbb{R} \to y \in \mathbb{R} \to x < y \lor x = y \lor y < x) \to x \in \mathbb{R} \to y \in \mathbb{R} \to x < y \lor x = y \lor y < x$
\item $\varphi_{38} :\equiv \forall x\,\forall y\,(x \in \mathbb{R} \to y \in \mathbb{R} \to x < y \lor x = y \lor y < x) \to \forall y\,(x \in \mathbb{R} \to y \in \mathbb{R} \to x < y \lor x = y \lor y < x)$
\item $\varphi_{39} :\equiv \mathrm{RFn}(f) \to \mathrm{StrInc}(f,\mathrm{dom}(f)) \to x \in \mathrm{dom}(f) \to \varphi_{16}$
\item $\varphi_{40} :\equiv \mathrm{RFn}(f) \to \mathrm{StrInc}(f,\mathrm{dom}(f)) \to x \in \mathrm{dom}(f) \to \varphi_{17}$
\item $\varphi_{41} :\equiv \mathrm{RFn}(f) \to \mathrm{StrInc}(f,\mathrm{dom}(f)) \to x \in \mathrm{dom}(f) \to \varphi_{18}$
\item $\varphi_{42} :\equiv \mathrm{RFn}(f) \to \mathrm{StrInc}(f,\mathrm{dom}(f)) \to x \in \mathrm{dom}(f) \to \varphi_{19}$
\item $\varphi_{43} :\equiv f(x) = f(y) \to y < x \to \forall v_{893}\,(v_{892} \in \mathrm{dom}(f) \to \varphi_{8})$
\item $\varphi_{44} :\equiv \neg f(y) < f(y) \to \mathrm{RFn}(f) \to \mathrm{StrInc}(f,\mathrm{dom}(f)) \to \varphi_{14}$
\item $\varphi_{45} :\equiv \neg f(y) < f(y) \to \mathrm{RFn}(f) \to \mathrm{StrInc}(f,\mathrm{dom}(f)) \to \varphi_{15}$
\item $\varphi_{46} :\equiv \forall w\,(f(x) = w \to f(y) < f(x) \to f(y) < w) \to f(x) = f(y) \to f(y) < f(x) \to f(y) < f(y)$
\item $\varphi_{47} :\equiv \forall w\,(f(x) = w \to f(x) < f(y) \to w < f(y)) \to f(x) = f(y) \to f(x) < f(y) \to f(y) < f(y)$
\item $\varphi_{48} :\equiv f(x) = f(y) \to y < x \to \forall v_{892}\,\forall v_{893}\,(v_{892} \in \mathrm{dom}(f) \to \varphi_{8})$
\item $\varphi_{49} :\equiv x \in \mathrm{dom}(f) \to y \in \mathrm{dom}(f) \to f(x) = f(y) \to y < x \to \varphi_{7}$
\item $\varphi_{50} :\equiv x \in \mathrm{dom}(f) \to y \in \mathrm{dom}(f) \to f(x) = f(y) \to x < y \to \varphi_{10}$
\item $\varphi_{51} :\equiv y \in \mathrm{dom}(f) \to f(x) = f(y) \to y < x \to v_{892} \in \mathrm{dom}(f) \to \varphi_{8}$
\item $\varphi_{52} :\equiv x \in \mathrm{dom}(f) \to y \in \mathrm{dom}(f) \to f(x) = f(y) \to x < y \to \forall y\,(\varphi_{10})$
\item $\varphi_{53} :\equiv \varphi_{32} \to \mathrm{RFn}(f) \to \mathrm{StrInc}(f,\mathrm{dom}(f)) \to \mathrm{Inj}(f)$
\item $\varphi_{54} :\equiv x \in \mathrm{dom}(f) \to y \in \mathrm{dom}(f) \to f(x) = f(y) \to y < x \to \forall x\,\forall y\,(\varphi_{10})$
\item $\varphi_{55} :\equiv x \in \mathrm{dom}(f) \to y \in \mathrm{dom}(f) \to f(x) = f(y) \to x < y \to \forall x\,\forall y\,(\varphi_{10})$
\item $\varphi_{56} :\equiv x \in \mathrm{dom}(f) \to y \in \mathrm{dom}(f) \to f(x) = f(y) \to y < x \to \varphi_{23}$
\item $\varphi_{57} :\equiv x \in \mathrm{dom}(f) \to y \in \mathrm{dom}(f) \to f(x) = f(y) \to y < x \to \varphi_{28}$
\item $\varphi_{58} :\equiv \forall f\,(\mathrm{Inj}(f) \leftrightarrow \forall x\,\forall y\,(\varphi_{4})) \to (\mathrm{Inj}(f) \leftrightarrow \forall x\,\forall y\,(\varphi_{4}))$
\item $\varphi_{59} :\equiv \varphi_{6} \to \mathrm{RFn}(f) \to \mathrm{StrInc}(f,\mathrm{dom}(f)) \to \forall y\,(\varphi_{4})$
\item $\varphi_{60} :\equiv (\mathrm{RFn}(f) \to \varphi_{3}) \to \mathrm{RFn}(f) \to \mathrm{StrInc}(f,\mathrm{dom}(f)) \to \varphi_{12}$
\item $\varphi_{61} :\equiv (\mathrm{RFn}(f) \to \varphi_{13}) \to \mathrm{RFn}(f) \to \mathrm{StrInc}(f,\mathrm{dom}(f)) \to \varphi_{11}$
\item $\varphi_{62} :\equiv \mathrm{StrInc}(f,\mathrm{dom}(f)) \to x \in \mathrm{dom}(f) \to y \in \mathrm{dom}(f) \to \varphi_{43}$
\item $\varphi_{63} :\equiv \mathrm{RFn}(f) \to \mathrm{StrInc}(f,\mathrm{dom}(f)) \to x \in \mathrm{dom}(f) \to \varphi_{51}$
\item $\varphi_{64} :\equiv (\mathrm{RFn}(f) \to \mathrm{StrInc}(f,\mathrm{dom}(f)) \to \varphi_{4}) \to \mathrm{RFn}(f) \land \mathrm{StrInc}(f,\mathrm{dom}(f)) \to \varphi_{4}$
\item $\varphi_{65} :\equiv \mathrm{StrInc}(f,\mathrm{dom}(f)) \to x \in \mathrm{dom}(f) \to y \in \mathrm{dom}(f) \to \varphi_{48}$
\item $\varphi_{66} :\equiv (\mathrm{RFn}(f) \to \mathrm{StrInc}(f,\mathrm{dom}(f)) \to \forall y\,(\varphi_{4})) \to \mathrm{RFn}(f) \land \mathrm{StrInc}(f,\mathrm{dom}(f)) \to \forall y\,(\varphi_{4})$
\item $\varphi_{67} :\equiv (f(x) = f(y) \to f(y) < f(x) \to f(y) < f(y)) \to \varphi_{40}$
\item $\varphi_{68} :\equiv (f(x) = f(y) \to f(x) < f(y) \to f(y) < f(y)) \to \varphi_{42}$
\item $\varphi_{69} :\equiv \forall y\,(\mathrm{RFn}(f) \land \mathrm{StrInc}(f,\mathrm{dom}(f)) \to \varphi_{4}) \to \varphi_{33}$
\item $\varphi_{70} :\equiv (\mathrm{RFn}(f) \to \mathrm{StrInc}(f,\mathrm{dom}(f)) \to \varphi_{14}) \to \mathrm{RFn}(f) \to \mathrm{StrInc}(f,\mathrm{dom}(f)) \to \varphi_{4}$
\item $\varphi_{71} :\equiv (\mathrm{RFn}(f) \to \mathrm{StrInc}(f,\mathrm{dom}(f)) \to \varphi_{11}) \to \mathrm{RFn}(f) \to \mathrm{StrInc}(f,\mathrm{dom}(f)) \to \varphi_{24}$
\item $\varphi_{72} :\equiv (\mathrm{RFn}(f) \to x \in \mathrm{dom}(f) \to x \in \mathbb{R}) \to (\mathrm{RFn}(f) \to \mathrm{RFn}(f)) \to \varphi_{60}$
\item $\varphi_{73} :\equiv (\mathrm{RFn}(f) \to y \in \mathrm{dom}(f) \to y \in \mathbb{R}) \to (\mathrm{RFn}(f) \to \mathrm{RFn}(f)) \to \varphi_{61}$
\item $\varphi_{74} :\equiv (\mathrm{StrInc}(f,\mathrm{dom}(f)) \leftrightarrow \forall x\,\forall y\,(\varphi_{10})) \to \mathrm{RFn}(f) \to \mathrm{StrInc}(f,\mathrm{dom}(f)) \to \varphi_{54}$
\item $\varphi_{75} :\equiv (\mathrm{StrInc}(f,\mathrm{dom}(f)) \leftrightarrow \forall x\,\forall y\,(\varphi_{10})) \to \mathrm{RFn}(f) \to \mathrm{StrInc}(f,\mathrm{dom}(f)) \to \varphi_{55}$
\item $\varphi_{76} :\equiv (\mathrm{RFn}(f) \to \varphi_{3}) \to (\mathrm{RFn}(f) \to \mathrm{StrInc}(f,\mathrm{dom}(f)) \to \varphi_{20}) \to \varphi_{39}$
\item $\varphi_{77} :\equiv (\forall y\,(\varphi_{10}) \to \varphi_{10}) \to \mathrm{RFn}(f) \to \mathrm{StrInc}(f,\mathrm{dom}(f)) \to \varphi_{50}$
\item $\varphi_{78} :\equiv (\mathrm{RFn}(f) \to \varphi_{13}) \to (\mathrm{RFn}(f) \to \mathrm{StrInc}(f,\mathrm{dom}(f)) \to \varphi_{21}) \to \varphi_{41}$
\item $\varphi_{79} :\equiv (\forall x\,\forall y\,(\varphi_{10}) \to \forall y\,(\varphi_{10})) \to \mathrm{RFn}(f) \to \mathrm{StrInc}(f,\mathrm{dom}(f)) \to \varphi_{52}$
\item $\varphi_{80} :\equiv (\varphi_{28} \to \varphi_{7}) \to \mathrm{RFn}(f) \to \mathrm{StrInc}(f,\mathrm{dom}(f)) \to \varphi_{49}$
\item $\varphi_{81} :\equiv (\forall x\,\forall y\,(\varphi_{10}) \to \varphi_{23}) \to \mathrm{RFn}(f) \to \mathrm{StrInc}(f,\mathrm{dom}(f)) \to \varphi_{56}$
\item $\varphi_{82} :\equiv (\forall v_{892}\,\forall v_{893}\,(v_{892} \in \mathrm{dom}(f) \to \varphi_{8}) \to \varphi_{28}) \to \mathrm{RFn}(f) \to \mathrm{StrInc}(f,\mathrm{dom}(f)) \to \varphi_{57}$
\item $\varphi_{83} :\equiv (x \in \mathbb{R} \to y \in \mathbb{R} \to x < y \lor x = y \lor y < x) \to (\mathrm{RFn}(f) \to \mathrm{StrInc}(f,\mathrm{dom}(f)) \to \varphi_{12}) \to \varphi_{71}$
\item $\varphi_{84} :\equiv (\mathrm{RFn}(f) \to \varphi_{62}) \to \varphi_{30} \to \forall v_{893}\,(v_{892} \in \mathrm{dom}(f) \to \varphi_{8})$
\item $\varphi_{85} :\equiv \forall v_{893}\,(\varphi_{30} \to v_{892} \in \mathrm{dom}(f) \to \varphi_{8}) \to \forall v_{893}\,(\varphi_{30}) \to \forall v_{893}\,(v_{892} \in \mathrm{dom}(f) \to \varphi_{8})$
\item $\varphi_{86} :\equiv (\mathrm{RFn}(f) \to \mathrm{StrInc}(f,\mathrm{dom}(f)) \to \varphi_{15}) \to (\mathrm{RFn}(f) \to \mathrm{StrInc}(f,\mathrm{dom}(f)) \to \varphi_{9}) \to \varphi_{70}$
\item $\varphi_{87} :\equiv \forall v_{892}\,(\varphi_{30} \to \forall v_{893}\,(v_{892} \in \mathrm{dom}(f) \to \varphi_{8})) \to \forall v_{892}\,(\varphi_{30}) \to \forall v_{892}\,\forall v_{893}\,(v_{892} \in \mathrm{dom}(f) \to \varphi_{8})$
\item $\varphi_{88} :\equiv (\mathrm{RFn}(f) \to \mathrm{StrInc}(f,\mathrm{dom}(f)) \to \varphi_{22}) \to \varphi_{39} \to \varphi_{67}$
\item $\varphi_{89} :\equiv (\mathrm{RFn}(f) \to \mathrm{StrInc}(f,\mathrm{dom}(f)) \to \varphi_{22}) \to \varphi_{41} \to \varphi_{68}$
\item $\varphi_{90} :\equiv \forall y\,(\mathrm{RFn}(f) \land \mathrm{StrInc}(f,\mathrm{dom}(f)) \to \varphi_{4}) \to \varphi_{69} \to \varphi_{59}$
\item $\varphi_{91} :\equiv (\mathrm{RFn}(f) \to \mathrm{StrInc}(f,\mathrm{dom}(f)) \to \varphi_{56}) \to (\varphi_{23} \to v_{892} \in \mathrm{dom}(f) \to \varphi_{8}) \to \varphi_{63}$
\item $\varphi_{92} :\equiv (\mathrm{RFn}(f) \to \mathrm{StrInc}(f,\mathrm{dom}(f)) \to \varphi_{49}) \to (\mathrm{RFn}(f) \to \varphi_{13}) \to \varphi_{76}$
\item $\varphi_{93} :\equiv (\mathrm{RFn}(f) \to \mathrm{StrInc}(f,\mathrm{dom}(f)) \to \varphi_{50}) \to (\mathrm{RFn}(f) \to \varphi_{3}) \to \varphi_{78}$
\item $\varphi_{94} :\equiv \forall v_{893}\,(\varphi_{30} \to v_{892} \in \mathrm{dom}(f) \to \varphi_{8}) \to \varphi_{85} \to (\varphi_{30} \to \forall v_{893}\,(\varphi_{30})) \to \mathrm{RFn}(f) \to \varphi_{62}$
\item $\varphi_{95} :\equiv \forall v_{892}\,(\varphi_{30} \to \forall v_{893}\,(v_{892} \in \mathrm{dom}(f) \to \varphi_{8})) \to \varphi_{87} \to (\varphi_{30} \to \forall v_{892}\,(\varphi_{30})) \to \mathrm{RFn}(f) \to \varphi_{65}$
\end{abbrevlist}
\end{fullproof}

% 自動生成 (prover/build.js): InvMono  ノード数 244
\begin{fullproof}{[I2]}{$\forall f\,(\mathrm{RFn}(f) \to \mathrm{StrInc}(f,\mathrm{dom}(f)) \to \mathrm{StrInc}(f^{-1},\mathrm{dom}(f^{-1})))$}{244}\phantomsection\label{tree:InvMono}
\noindent$\Pi_{1}$:
\begin{ptree}
\AXC{$z = w \to x < z \to x < w$\lbl{A6}}
\UIC{$\forall w\,(z = w \to x < z \to x < w)$}
\UIC{$\forall z\,\forall w\,(z = w \to x < z \to x < w)$}
\AXC{$\forall z\,\forall w\,(z = w \to x < z \to x < w) \to \varphi_{13}$\lbl{A2}}
\BIC{$\varphi_{13}$}
\end{ptree}
\noindent$\Pi_{2}$:
\begin{ptree}
\AXC{$z = w \to x < f(z) \to x < f(w)$\lbl{A6}}
\UIC{$\forall w\,(z = w \to x < f(z) \to x < f(w))$}
\UIC{$\forall z\,\forall w\,(z = w \to x < f(z) \to x < f(w))$}
\AXC{$\forall z\,\forall w\,(z = w \to x < f(z) \to x < f(w)) \to \varphi_{14}$\lbl{A2}}
\BIC{$\varphi_{14}$}
\AXC{$\varphi_{14} \to \varphi_{47}$\lbl{A2}}
\BIC{$\varphi_{47}$}
\end{ptree}
\noindent$\Pi_{3}$:
\begin{ptree}
\AXC{$x = x$\lbl{A5}}
\AXC{$x = y \to x = x \to y = x$\lbl{A6}}
\AXC{$(x = y \to x = x \to y = x) \to x = x \to x = y \to y = x$\lbl{A1}}
\BIC{$x = x \to x = y \to y = x$}
\BIC{$x = y \to y = x$}
\UIC{$\forall y\,(x = y \to y = x)$}
\UIC{$\forall x\,\forall y\,(x = y \to y = x)$}
\end{ptree}
\noindent$\Pi_{4}$:
\begin{ptree}
\AXC{$\stackrel{\Pi_{3}}{\vdots}$}\noLine
\UIC{$\forall x\,\forall y\,(x = y \to y = x)$}
\AXC{$\forall x\,\forall y\,(x = y \to y = x) \to \forall y\,(v_{896} = y \to y = v_{896})$\lbl{A2}}
\BIC{$\forall y\,(v_{896} = y \to y = v_{896})$}
\AXC{$\varphi_{3}$\lbl{A2}}
\BIC{$v_{896} = v_{897} \to v_{897} = v_{896}$}
\UIC{$\forall v_{897}\,(v_{896} = v_{897} \to v_{897} = v_{896})$}
\UIC{$\forall v_{896}\,\forall v_{897}\,(v_{896} = v_{897} \to v_{897} = v_{896})$}
\end{ptree}
\noindent$\Pi_{5}$:
\begin{ptree}
\AXC{$\stackrel{\Pi_{4}}{\vdots}$}\noLine
\UIC{$\forall v_{896}\,\forall v_{897}\,(v_{896} = v_{897} \to v_{897} = v_{896})$}
\AXC{$\forall v_{896}\,\forall v_{897}\,(v_{896} = v_{897} \to v_{897} = v_{896}) \to \varphi_{12}$\lbl{A2}}
\BIC{$\varphi_{12}$}
\end{ptree}
\noindent$\Pi_{6}$:
\begin{ptree}
\AXC{$\mathrm{RFn}(f) \to \mathrm{RFn}(f)$\lbl{A1}}
\AXC{$\forall f\,(\mathrm{RFn}(f) \to \mathrm{Inj}(f) \to \varphi_{90})$\lbl{[S26]}}
\AXC{$\varphi_{119}$\lbl{A2}}
\BIC{$\mathrm{RFn}(f) \to \mathrm{Inj}(f) \to \varphi_{90}$}
\AXC{$\varphi_{128}$\lbl{A1}}
\BIC{$(\mathrm{RFn}(f) \to \mathrm{RFn}(f)) \to \varphi_{110}$}
\BIC{$\varphi_{110}$}
\end{ptree}
\noindent$\Pi_{7}$:
\begin{ptree}
\AXC{$\varphi_{2}$\lbl{[I1]}}
\AXC{$\varphi_{29}$\lbl{A2}}
\BIC{$\mathrm{RFn}(f) \to \mathrm{StrInc}(f,\mathrm{dom}(f)) \to \mathrm{Inj}(f)$}
\AXC{$\stackrel{\Pi_{6}}{\vdots}$}\noLine
\UIC{$\varphi_{110}$}
\BIC{$\mathrm{RFn}(f) \to \mathrm{StrInc}(f,\mathrm{dom}(f)) \to \varphi_{90}$}
\end{ptree}
\noindent$\Pi_{8}$:
\begin{ptree}
\AXC{$\stackrel{\Pi_{7}}{\vdots}$}\noLine
\UIC{$\mathrm{RFn}(f) \to \mathrm{StrInc}(f,\mathrm{dom}(f)) \to \varphi_{90}$}
\AXC{$\varphi_{117}$\lbl{A1}}
\BIC{$\mathrm{RFn}(f) \to \varphi_{79}$}
\AXC{$(\mathrm{RFn}(f) \to \varphi_{79}) \to (\varphi_{35} \to \varphi_{25}) \to \mathrm{RFn}(f) \to \varphi_{73}$\lbl{A1}}
\BIC{$(\varphi_{35} \to \varphi_{25}) \to \mathrm{RFn}(f) \to \varphi_{73}$}
\end{ptree}
\noindent$\Pi_{9}$:
\begin{ptree}
\AXC{$\mathrm{RFn}(f) \to \mathrm{StrInc}(f,\mathrm{dom}(f)) \to \varphi_{11}$\lbl{A1}}
\AXC{$\varphi_{35} \to \varphi_{25}$\lbl{A2}}
\AXC{$\stackrel{\Pi_{8}}{\vdots}$}\noLine
\UIC{$(\varphi_{35} \to \varphi_{25}) \to \mathrm{RFn}(f) \to \varphi_{73}$}
\BIC{$\mathrm{RFn}(f) \to \varphi_{73}$}
\AXC{$\varphi_{118}$\lbl{A1}}
\BIC{$(\mathrm{RFn}(f) \to \mathrm{StrInc}(f,\mathrm{dom}(f)) \to \varphi_{11}) \to \varphi_{67}$}
\BIC{$\varphi_{67}$}
\end{ptree}
\noindent$\Pi_{10}$:
\begin{ptree}
\AXC{$\stackrel{\Pi_{9}}{\vdots}$}\noLine
\UIC{$\varphi_{67}$}
\AXC{$\varphi_{67} \to \mathrm{RFn}(f) \to \mathrm{StrInc}(f,\mathrm{dom}(f)) \to \varphi_{30}$\lbl{A1}}
\BIC{$\mathrm{RFn}(f) \to \mathrm{StrInc}(f,\mathrm{dom}(f)) \to \varphi_{30}$}
\end{ptree}
\noindent$\Pi_{11}$:
\begin{ptree}
\AXC{$\stackrel{\Pi_{10}}{\vdots}$}\noLine
\UIC{$\mathrm{RFn}(f) \to \mathrm{StrInc}(f,\mathrm{dom}(f)) \to \varphi_{30}$}
\AXC{$(\mathrm{RFn}(f) \to \mathrm{StrInc}(f,\mathrm{dom}(f)) \to \varphi_{30}) \to \varphi_{81}$\lbl{A1}}
\BIC{$\varphi_{81}$}
\end{ptree}
\noindent$\Pi_{12}$:
\begin{ptree}
\AXC{$\mathrm{RFn}(f) \to \varphi_{37}$\lbl{A1}}
\AXC{$\stackrel{\Pi_{5}}{\vdots}$}\noLine
\UIC{$\varphi_{12}$}
\AXC{$\varphi_{12} \to f(f^{-1}(y)) = y \to y = f(f^{-1}(y))$\lbl{A2}}
\BIC{$f(f^{-1}(y)) = y \to y = f(f^{-1}(y))$}
\AXC{$\stackrel{\Pi_{11}}{\vdots}$}\noLine
\UIC{$\varphi_{81}$}
\BIC{$\mathrm{RFn}(f) \to \mathrm{StrInc}(f,\mathrm{dom}(f)) \to \varphi_{31}$}
\AXC{$\varphi_{121}$\lbl{A1}}
\BIC{$(\mathrm{RFn}(f) \to \varphi_{37}) \to \varphi_{105}$}
\BIC{$\varphi_{105}$}
\end{ptree}
\noindent$\Pi_{13}$:
\begin{ptree}
\AXC{$z = w \to x < z \to x < w$\lbl{A6}}
\UIC{$\forall w\,(z = w \to x < z \to x < w)$}
\UIC{$\forall z\,\forall w\,(z = w \to x < z \to x < w)$}
\AXC{$\varphi_{5}$\lbl{A2}}
\BIC{$\forall w\,(y = w \to x < y \to x < w)$}
\AXC{$\varphi_{48}$\lbl{A2}}
\BIC{$y = f(f^{-1}(y)) \to x < y \to x < f(f^{-1}(y))$}
\AXC{$\stackrel{\Pi_{12}}{\vdots}$}\noLine
\UIC{$\varphi_{105}$}
\BIC{$\mathrm{RFn}(f) \to \varphi_{82}$}
\AXC{$\mathrm{RFn}(f) \to \varphi_{84}$\lbl{A1}}
\AXC{$\varphi_{137}$\lbl{A1}}
\BIC{$(\mathrm{RFn}(f) \to \varphi_{82}) \to \varphi_{47} \to \mathrm{RFn}(f) \to \varphi_{83}$}
\BIC{$\varphi_{47} \to \mathrm{RFn}(f) \to \varphi_{83}$}
\end{ptree}
\noindent$\Pi_{14}$:
\begin{ptree}
\AXC{$\stackrel{\Pi_{7}}{\vdots}$}\noLine
\UIC{$\mathrm{RFn}(f) \to \mathrm{StrInc}(f,\mathrm{dom}(f)) \to \varphi_{90}$}
\AXC{$\varphi_{117}$\lbl{A1}}
\BIC{$\mathrm{RFn}(f) \to \varphi_{79}$}
\AXC{$(\mathrm{RFn}(f) \to \varphi_{79}) \to (\varphi_{35} \to \varphi_{26}) \to \mathrm{RFn}(f) \to \varphi_{74}$\lbl{A1}}
\BIC{$(\varphi_{35} \to \varphi_{26}) \to \mathrm{RFn}(f) \to \varphi_{74}$}
\end{ptree}
\noindent$\Pi_{15}$:
\begin{ptree}
\AXC{$\mathrm{RFn}(f) \to \varphi_{8}$\lbl{A1}}
\AXC{$\varphi_{35} \to \varphi_{26}$\lbl{A2}}
\AXC{$\stackrel{\Pi_{14}}{\vdots}$}\noLine
\UIC{$(\varphi_{35} \to \varphi_{26}) \to \mathrm{RFn}(f) \to \varphi_{74}$}
\BIC{$\mathrm{RFn}(f) \to \varphi_{74}$}
\AXC{$(\mathrm{RFn}(f) \to \varphi_{74}) \to (\mathrm{RFn}(f) \to \varphi_{8}) \to \varphi_{68}$\lbl{A1}}
\BIC{$(\mathrm{RFn}(f) \to \varphi_{8}) \to \varphi_{68}$}
\BIC{$\varphi_{68}$}
\end{ptree}
\noindent$\Pi_{16}$:
\begin{ptree}
\AXC{$\stackrel{\Pi_{15}}{\vdots}$}\noLine
\UIC{$\varphi_{68}$}
\AXC{$\varphi_{68} \to \mathrm{RFn}(f) \to \mathrm{StrInc}(f,\mathrm{dom}(f)) \to \varphi_{32}$\lbl{A1}}
\BIC{$\mathrm{RFn}(f) \to \mathrm{StrInc}(f,\mathrm{dom}(f)) \to \varphi_{32}$}
\end{ptree}
\noindent$\Pi_{17}$:
\begin{ptree}
\AXC{$\stackrel{\Pi_{2}}{\vdots}$}\noLine
\UIC{$\varphi_{47}$}
\AXC{$\stackrel{\Pi_{13}}{\vdots}$}\noLine
\UIC{$\varphi_{47} \to \mathrm{RFn}(f) \to \varphi_{83}$}
\BIC{$\mathrm{RFn}(f) \to \varphi_{83}$}
\AXC{$\stackrel{\Pi_{16}}{\vdots}$}\noLine
\UIC{$\mathrm{RFn}(f) \to \mathrm{StrInc}(f,\mathrm{dom}(f)) \to \varphi_{32}$}
\AXC{$\varphi_{129}$\lbl{A1}}
\BIC{$(\mathrm{RFn}(f) \to \varphi_{83}) \to \varphi_{98}$}
\BIC{$\varphi_{98}$}
\end{ptree}
\noindent$\Pi_{18}$:
\begin{ptree}
\AXC{$\stackrel{\Pi_{1}}{\vdots}$}\noLine
\UIC{$\varphi_{13}$}
\AXC{$\varphi_{13} \to f(f^{-1}(x)) = x \to x < f(f^{-1}(x)) \to x < x$\lbl{A2}}
\BIC{$f(f^{-1}(x)) = x \to x < f(f^{-1}(x)) \to x < x$}
\AXC{$\stackrel{\Pi_{17}}{\vdots}$}\noLine
\UIC{$\varphi_{98}$}
\BIC{$\mathrm{RFn}(f) \to \varphi_{71}$}
\end{ptree}
\noindent$\Pi_{19}$:
\begin{ptree}
\AXC{$z = w \to y < z \to y < w$\lbl{A6}}
\UIC{$\forall w\,(z = w \to y < z \to y < w)$}
\UIC{$\forall z\,\forall w\,(z = w \to y < z \to y < w)$}
\AXC{$\forall z\,\forall w\,(z = w \to y < z \to y < w) \to \varphi_{15}$\lbl{A2}}
\BIC{$\varphi_{15}$}
\end{ptree}
\noindent$\Pi_{20}$:
\begin{ptree}
\AXC{$z = w \to z < f(f^{-1}(x)) \to w < f(f^{-1}(x))$\lbl{A6}}
\UIC{$\varphi_{16}$}
\UIC{$\forall z\,(\varphi_{16})$}
\AXC{$\forall z\,(\varphi_{16}) \to \varphi_{58}$\lbl{A2}}
\BIC{$\varphi_{58}$}
\AXC{$\varphi_{58} \to f(f^{-1}(y)) = y \to \varphi_{36}$\lbl{A2}}
\BIC{$f(f^{-1}(y)) = y \to \varphi_{36}$}
\end{ptree}
\noindent$\Pi_{21}$:
\begin{ptree}
\AXC{$\stackrel{\Pi_{15}}{\vdots}$}\noLine
\UIC{$\varphi_{68}$}
\AXC{$\varphi_{68} \to \mathrm{RFn}(f) \to \mathrm{StrInc}(f,\mathrm{dom}(f)) \to \varphi_{38}$\lbl{A1}}
\BIC{$\mathrm{RFn}(f) \to \mathrm{StrInc}(f,\mathrm{dom}(f)) \to \varphi_{38}$}
\end{ptree}
\noindent$\Pi_{22}$:
\begin{ptree}
\AXC{$\stackrel{\Pi_{9}}{\vdots}$}\noLine
\UIC{$\varphi_{67}$}
\AXC{$\varphi_{67} \to \mathrm{RFn}(f) \to \mathrm{StrInc}(f,\mathrm{dom}(f)) \to \varphi_{39}$\lbl{A1}}
\BIC{$\mathrm{RFn}(f) \to \mathrm{StrInc}(f,\mathrm{dom}(f)) \to \varphi_{39}$}
\end{ptree}
\noindent$\Pi_{23}$:
\begin{ptree}
\AXC{$\forall f\,(\varphi_{24})$\lbl{DefStrInc}}
\AXC{$\forall f\,(\varphi_{24}) \to \varphi_{24}$\lbl{A2}}
\BIC{$\varphi_{24}$}
\AXC{$\varphi_{24} \to (\mathrm{StrInc}(f,\mathrm{dom}(f)) \leftrightarrow \varphi_{17})$\lbl{A2}}
\BIC{$\mathrm{StrInc}(f,\mathrm{dom}(f)) \leftrightarrow \varphi_{17}$}
\end{ptree}
\noindent$\Pi_{24}$:
\begin{ptree}
\AXC{$\stackrel{\Pi_{23}}{\vdots}$}\noLine
\UIC{$\mathrm{StrInc}(f,\mathrm{dom}(f)) \leftrightarrow \varphi_{17}$}
\AXC{$\mathrm{RFn}(f) \to \varphi_{1}$\lbl{A1}}
\AXC{$\varphi_{115}$\lbl{A1}}
\BIC{$(\mathrm{StrInc}(f,\mathrm{dom}(f)) \leftrightarrow \varphi_{17}) \to \mathrm{RFn}(f) \to \varphi_{93}$}
\BIC{$\mathrm{RFn}(f) \to \varphi_{93}$}
\end{ptree}
\noindent$\Pi_{25}$:
\begin{ptree}
\AXC{$\varphi_{17} \to \varphi_{22}$\lbl{A2}}
\AXC{$\stackrel{\Pi_{24}}{\vdots}$}\noLine
\UIC{$\mathrm{RFn}(f) \to \varphi_{93}$}
\AXC{$(\mathrm{RFn}(f) \to \varphi_{93}) \to (\varphi_{17} \to \varphi_{22}) \to \mathrm{RFn}(f) \to \varphi_{95}$\lbl{A1}}
\BIC{$(\varphi_{17} \to \varphi_{22}) \to \mathrm{RFn}(f) \to \varphi_{95}$}
\BIC{$\mathrm{RFn}(f) \to \varphi_{95}$}
\AXC{$\varphi_{136}$\lbl{A1}}
\BIC{$(\varphi_{22} \to v_{894} \in \mathrm{dom}(f) \to \varphi_{9}) \to \mathrm{RFn}(f) \to \varphi_{100}$}
\end{ptree}
\noindent$\Pi_{26}$:
\begin{ptree}
\AXC{$\varphi_{22} \to v_{894} \in \mathrm{dom}(f) \to \varphi_{9}$\lbl{A2}}
\AXC{$\stackrel{\Pi_{25}}{\vdots}$}\noLine
\UIC{$(\varphi_{22} \to v_{894} \in \mathrm{dom}(f) \to \varphi_{9}) \to \mathrm{RFn}(f) \to \varphi_{100}$}
\BIC{$\mathrm{RFn}(f) \to \varphi_{100}$}
\end{ptree}
\noindent$\Pi_{27}$:
\begin{ptree}
\AXC{$\stackrel{\Pi_{26}}{\vdots}$}\noLine
\UIC{$\mathrm{RFn}(f) \to \varphi_{100}$}
\AXC{$(\mathrm{RFn}(f) \to \varphi_{100}) \to \varphi_{75} \to v_{894} \in \mathrm{dom}(f) \to \varphi_{9}$\lbl{A1}}
\BIC{$\varphi_{75} \to v_{894} \in \mathrm{dom}(f) \to \varphi_{9}$}
\UIC{$\forall v_{895}\,(\varphi_{75} \to v_{894} \in \mathrm{dom}(f) \to \varphi_{9})$}
\AXC{$\forall v_{895}\,(\varphi_{75} \to v_{894} \in \mathrm{dom}(f) \to \varphi_{9}) \to \varphi_{124} \to \varphi_{131}$\lbl{A1}}
\BIC{$\varphi_{124} \to \varphi_{131}$}
\end{ptree}
\noindent$\Pi_{28}$:
\begin{ptree}
\AXC{$\varphi_{75} \to \forall v_{895}\,(\varphi_{75})$\lbl{A4}}
\AXC{$\varphi_{124}$\lbl{A3}}
\AXC{$\stackrel{\Pi_{27}}{\vdots}$}\noLine
\UIC{$\varphi_{124} \to \varphi_{131}$}
\BIC{$\varphi_{131}$}
\BIC{$\mathrm{RFn}(f) \to \mathrm{StrInc}(f,\mathrm{dom}(f)) \to \varphi_{92}$}
\AXC{$\varphi_{123}$\lbl{A1}}
\BIC{$\varphi_{75} \to \forall v_{895}\,(v_{894} \in \mathrm{dom}(f) \to \varphi_{9})$}
\UIC{$\forall v_{894}\,(\varphi_{75} \to \forall v_{895}\,(v_{894} \in \mathrm{dom}(f) \to \varphi_{9}))$}
\AXC{$\forall v_{894}\,(\varphi_{75} \to \forall v_{895}\,(v_{894} \in \mathrm{dom}(f) \to \varphi_{9})) \to \varphi_{127} \to \varphi_{133}$\lbl{A1}}
\BIC{$\varphi_{127} \to \varphi_{133}$}
\end{ptree}
\noindent$\Pi_{29}$:
\begin{ptree}
\AXC{$\varphi_{75} \to \forall v_{894}\,(\varphi_{75})$\lbl{A4}}
\AXC{$\varphi_{127}$\lbl{A3}}
\AXC{$\stackrel{\Pi_{28}}{\vdots}$}\noLine
\UIC{$\varphi_{127} \to \varphi_{133}$}
\BIC{$\varphi_{133}$}
\BIC{$\mathrm{RFn}(f) \to \mathrm{StrInc}(f,\mathrm{dom}(f)) \to \varphi_{94}$}
\AXC{$(\mathrm{RFn}(f) \to \mathrm{StrInc}(f,\mathrm{dom}(f)) \to \varphi_{94}) \to \varphi_{125}$\lbl{A1}}
\BIC{$\varphi_{125}$}
\end{ptree}
\noindent$\Pi_{30}$:
\begin{ptree}
\AXC{$\varphi_{108}$\lbl{A2}}
\AXC{$\forall v_{894}\,\forall v_{895}\,(v_{894} \in \mathrm{dom}(f) \to \varphi_{9}) \to \varphi_{59}$\lbl{A2}}
\AXC{$\stackrel{\Pi_{29}}{\vdots}$}\noLine
\UIC{$\varphi_{125}$}
\BIC{$\mathrm{RFn}(f) \to \varphi_{107}$}
\AXC{$\varphi_{138}$\lbl{A1}}
\BIC{$\varphi_{108} \to \mathrm{RFn}(f) \to \mathrm{StrInc}(f,\mathrm{dom}(f)) \to \varphi_{111}$}
\BIC{$\mathrm{RFn}(f) \to \mathrm{StrInc}(f,\mathrm{dom}(f)) \to \varphi_{111}$}
\end{ptree}
\noindent$\Pi_{31}$:
\begin{ptree}
\AXC{$\stackrel{\Pi_{30}}{\vdots}$}\noLine
\UIC{$\mathrm{RFn}(f) \to \mathrm{StrInc}(f,\mathrm{dom}(f)) \to \varphi_{111}$}
\AXC{$(\mathrm{RFn}(f) \to \mathrm{StrInc}(f,\mathrm{dom}(f)) \to \varphi_{111}) \to \varphi_{135}$\lbl{A1}}
\BIC{$\varphi_{135}$}
\end{ptree}
\noindent$\Pi_{32}$:
\begin{ptree}
\AXC{$\stackrel{\Pi_{21}}{\vdots}$}\noLine
\UIC{$\mathrm{RFn}(f) \to \mathrm{StrInc}(f,\mathrm{dom}(f)) \to \varphi_{38}$}
\AXC{$\stackrel{\Pi_{22}}{\vdots}$}\noLine
\UIC{$\mathrm{RFn}(f) \to \mathrm{StrInc}(f,\mathrm{dom}(f)) \to \varphi_{39}$}
\AXC{$\stackrel{\Pi_{31}}{\vdots}$}\noLine
\UIC{$\varphi_{135}$}
\BIC{$(\mathrm{RFn}(f) \to \mathrm{StrInc}(f,\mathrm{dom}(f)) \to \varphi_{38}) \to \varphi_{116}$}
\BIC{$\varphi_{116}$}
\end{ptree}
\noindent$\Pi_{33}$:
\begin{ptree}
\AXC{$\stackrel{\Pi_{10}}{\vdots}$}\noLine
\UIC{$\mathrm{RFn}(f) \to \mathrm{StrInc}(f,\mathrm{dom}(f)) \to \varphi_{30}$}
\AXC{$(\mathrm{RFn}(f) \to \mathrm{StrInc}(f,\mathrm{dom}(f)) \to \varphi_{30}) \to \varphi_{126}$\lbl{A1}}
\BIC{$\varphi_{126}$}
\end{ptree}
\noindent$\Pi_{34}$:
\begin{ptree}
\AXC{$\stackrel{\Pi_{20}}{\vdots}$}\noLine
\UIC{$f(f^{-1}(y)) = y \to \varphi_{36}$}
\AXC{$\mathrm{RFn}(f) \to \varphi_{87}$\lbl{A1}}
\AXC{$\stackrel{\Pi_{32}}{\vdots}$}\noLine
\UIC{$\varphi_{116}$}
\BIC{$\mathrm{RFn}(f) \to \mathrm{StrInc}(f,\mathrm{dom}(f)) \to \varphi_{86}$}
\AXC{$\stackrel{\Pi_{33}}{\vdots}$}\noLine
\UIC{$\varphi_{126}$}
\BIC{$(f(f^{-1}(y)) = y \to \varphi_{36}) \to \mathrm{RFn}(f) \to \varphi_{85}$}
\BIC{$\mathrm{RFn}(f) \to \varphi_{85}$}
\end{ptree}
\noindent$\Pi_{35}$:
\begin{ptree}
\AXC{$\stackrel{\Pi_{34}}{\vdots}$}\noLine
\UIC{$\mathrm{RFn}(f) \to \varphi_{85}$}
\AXC{$\stackrel{\Pi_{16}}{\vdots}$}\noLine
\UIC{$\mathrm{RFn}(f) \to \mathrm{StrInc}(f,\mathrm{dom}(f)) \to \varphi_{32}$}
\AXC{$\varphi_{130}$\lbl{A1}}
\BIC{$(\mathrm{RFn}(f) \to \varphi_{85}) \to \varphi_{101}$}
\BIC{$\varphi_{101}$}
\end{ptree}
\noindent$\Pi_{36}$:
\begin{ptree}
\AXC{$\stackrel{\Pi_{19}}{\vdots}$}\noLine
\UIC{$\varphi_{15}$}
\AXC{$\varphi_{15} \to f(f^{-1}(x)) = x \to y < f(f^{-1}(x)) \to y < x$\lbl{A2}}
\BIC{$f(f^{-1}(x)) = x \to y < f(f^{-1}(x)) \to y < x$}
\AXC{$\stackrel{\Pi_{35}}{\vdots}$}\noLine
\UIC{$\varphi_{101}$}
\BIC{$\mathrm{RFn}(f) \to \varphi_{76}$}
\end{ptree}
\noindent$\Pi_{37}$:
\begin{ptree}
\AXC{$\forall x\,\forall y\,\forall z\,(x < y \to y < z \to x < z)$\lbl{[O2]}}
\AXC{$\varphi_{18}$\lbl{A2}}
\BIC{$\forall y\,\forall z\,(x < y \to y < z \to x < z)$}
\AXC{$\varphi_{10}$\lbl{A2}}
\BIC{$\forall z\,(x < y \to y < z \to x < z)$}
\AXC{$\forall z\,(x < y \to y < z \to x < z) \to x < y \to y < x \to x < x$\lbl{A2}}
\BIC{$x < y \to y < x \to x < x$}
\end{ptree}
\noindent$\Pi_{38}$:
\begin{ptree}
\AXC{$\stackrel{\Pi_{36}}{\vdots}$}\noLine
\UIC{$\mathrm{RFn}(f) \to \varphi_{76}$}
\AXC{$\mathrm{RFn}(f) \to \varphi_{37}$\lbl{A1}}
\AXC{$\stackrel{\Pi_{37}}{\vdots}$}\noLine
\UIC{$x < y \to y < x \to x < x$}
\AXC{$\varphi_{120}$\lbl{A1}}
\BIC{$(\mathrm{RFn}(f) \to \varphi_{37}) \to (\mathrm{RFn}(f) \to \varphi_{76}) \to \mathrm{RFn}(f) \to \varphi_{77}$}
\BIC{$(\mathrm{RFn}(f) \to \varphi_{76}) \to \mathrm{RFn}(f) \to \varphi_{77}$}
\BIC{$\mathrm{RFn}(f) \to \varphi_{77}$}
\end{ptree}
\noindent$\Pi_{39}$:
\begin{ptree}
\AXC{$\stackrel{\Pi_{18}}{\vdots}$}\noLine
\UIC{$\mathrm{RFn}(f) \to \varphi_{71}$}
\AXC{$\stackrel{\Pi_{38}}{\vdots}$}\noLine
\UIC{$\mathrm{RFn}(f) \to \varphi_{77}$}
\AXC{$\mathrm{RFn}(f) \to \varphi_{69}$\lbl{A1}}
\AXC{$\varphi_{134}$\lbl{A1}}
\BIC{$(\mathrm{RFn}(f) \to \varphi_{77}) \to (\mathrm{RFn}(f) \to \varphi_{71}) \to \varphi_{60}$}
\BIC{$(\mathrm{RFn}(f) \to \varphi_{71}) \to \varphi_{60}$}
\BIC{$\varphi_{60}$}
\AXC{$\varphi_{60} \to \neg x < x \to \mathrm{RFn}(f) \to \varphi_{70}$\lbl{A1}}
\BIC{$\neg x < x \to \mathrm{RFn}(f) \to \varphi_{70}$}
\end{ptree}
\noindent$\Pi_{40}$:
\begin{ptree}
\AXC{$\mathrm{RFn}(f) \to \mathrm{RFn}(f)$\lbl{A1}}
\AXC{$\forall f\,\forall x\,(\mathrm{RFn}(f) \to x \in \mathrm{dom}(f) \to x \in \mathbb{R})$\lbl{[S7]}}
\AXC{$\varphi_{40}$\lbl{A2}}
\BIC{$\forall x\,(\mathrm{RFn}(f) \to x \in \mathrm{dom}(f) \to x \in \mathbb{R})$}
\AXC{$\varphi_{56}$\lbl{A2}}
\BIC{$\mathrm{RFn}(f) \to f^{-1}(y) \in \mathrm{dom}(f) \to f^{-1}(y) \in \mathbb{R}$}
\AXC{$\varphi_{112}$\lbl{A1}}
\BIC{$(\mathrm{RFn}(f) \to \mathrm{RFn}(f)) \to \varphi_{96}$}
\BIC{$\varphi_{96}$}
\end{ptree}
\noindent$\Pi_{41}$:
\begin{ptree}
\AXC{$\mathrm{RFn}(f) \to \mathrm{RFn}(f)$\lbl{A1}}
\AXC{$\forall f\,\forall x\,(\mathrm{RFn}(f) \to x \in \mathrm{dom}(f) \to x \in \mathbb{R})$\lbl{[S7]}}
\AXC{$\varphi_{40}$\lbl{A2}}
\BIC{$\forall x\,(\mathrm{RFn}(f) \to x \in \mathrm{dom}(f) \to x \in \mathbb{R})$}
\AXC{$\varphi_{57}$\lbl{A2}}
\BIC{$\mathrm{RFn}(f) \to f^{-1}(x) \in \mathrm{dom}(f) \to f^{-1}(x) \in \mathbb{R}$}
\AXC{$\varphi_{113}$\lbl{A1}}
\BIC{$(\mathrm{RFn}(f) \to \mathrm{RFn}(f)) \to \varphi_{97}$}
\BIC{$\varphi_{97}$}
\end{ptree}
\noindent$\Pi_{42}$:
\begin{ptree}
\AXC{$\forall x\,\forall y\,(x \in \mathbb{R} \to y \in \mathbb{R} \to \neg y \leq x \to x < y)$\lbl{[A2]}}
\AXC{$\forall x\,\forall y\,(x \in \mathbb{R} \to y \in \mathbb{R} \to \neg y \leq x \to x < y) \to \varphi_{41}$\lbl{A2}}
\BIC{$\varphi_{41}$}
\end{ptree}
\noindent$\Pi_{43}$:
\begin{ptree}
\AXC{$\stackrel{\Pi_{42}}{\vdots}$}\noLine
\UIC{$\varphi_{41}$}
\AXC{$\varphi_{41} \to f^{-1}(x) \in \mathbb{R} \to \varphi_{52}$\lbl{A2}}
\BIC{$f^{-1}(x) \in \mathbb{R} \to \varphi_{52}$}
\AXC{$(f^{-1}(x) \in \mathbb{R} \to \varphi_{52}) \to \varphi_{122}$\lbl{A1}}
\BIC{$\varphi_{122}$}
\end{ptree}
\noindent$\Pi_{44}$:
\begin{ptree}
\AXC{$\stackrel{\Pi_{21}}{\vdots}$}\noLine
\UIC{$\mathrm{RFn}(f) \to \mathrm{StrInc}(f,\mathrm{dom}(f)) \to \varphi_{38}$}
\AXC{$\stackrel{\Pi_{41}}{\vdots}$}\noLine
\UIC{$\varphi_{97}$}
\BIC{$\mathrm{RFn}(f) \to \mathrm{StrInc}(f,\mathrm{dom}(f)) \to \varphi_{34}$}
\AXC{$\stackrel{\Pi_{43}}{\vdots}$}\noLine
\UIC{$\varphi_{122}$}
\BIC{$(\mathrm{RFn}(f) \to \mathrm{StrInc}(f,\mathrm{dom}(f)) \to \varphi_{33}) \to \varphi_{102}$}
\end{ptree}
\noindent$\Pi_{45}$:
\begin{ptree}
\AXC{$\stackrel{\Pi_{22}}{\vdots}$}\noLine
\UIC{$\mathrm{RFn}(f) \to \mathrm{StrInc}(f,\mathrm{dom}(f)) \to \varphi_{39}$}
\AXC{$\stackrel{\Pi_{40}}{\vdots}$}\noLine
\UIC{$\varphi_{96}$}
\BIC{$\mathrm{RFn}(f) \to \mathrm{StrInc}(f,\mathrm{dom}(f)) \to \varphi_{33}$}
\AXC{$\stackrel{\Pi_{44}}{\vdots}$}\noLine
\UIC{$(\mathrm{RFn}(f) \to \mathrm{StrInc}(f,\mathrm{dom}(f)) \to \varphi_{33}) \to \varphi_{102}$}
\BIC{$\varphi_{102}$}
\end{ptree}
\noindent$\Pi_{46}$:
\begin{ptree}
\AXC{$\forall x\,(\neg x < x)$\lbl{[O1]}}
\AXC{$\forall x\,(\neg x < x) \to \neg x < x$\lbl{A2}}
\BIC{$\neg x < x$}
\AXC{$\stackrel{\Pi_{39}}{\vdots}$}\noLine
\UIC{$\neg x < x \to \mathrm{RFn}(f) \to \varphi_{70}$}
\BIC{$\mathrm{RFn}(f) \to \varphi_{70}$}
\AXC{$\varphi_{109}$\lbl{A1}}
\BIC{$\mathrm{RFn}(f) \to \mathrm{StrInc}(f,\mathrm{dom}(f)) \to \varphi_{45}$}
\AXC{$\stackrel{\Pi_{45}}{\vdots}$}\noLine
\UIC{$\varphi_{102}$}
\BIC{$\mathrm{RFn}(f) \to \mathrm{StrInc}(f,\mathrm{dom}(f)) \to \varphi_{46}$}
\AXC{$\varphi_{103}$\lbl{A1}}
\BIC{$\mathrm{RFn}(f) \land \mathrm{StrInc}(f,\mathrm{dom}(f)) \to \varphi_{46}$}
\UIC{$\forall y\,(\mathrm{RFn}(f) \land \mathrm{StrInc}(f,\mathrm{dom}(f)) \to \varphi_{46})$}
\end{ptree}
\noindent$\Pi_{47}$:
\begin{ptree}
\AXC{$\varphi_{7}$\lbl{A4}}
\AXC{$\varphi_{106}$\lbl{A3}}
\AXC{$\stackrel{\Pi_{46}}{\vdots}$}\noLine
\UIC{$\forall y\,(\mathrm{RFn}(f) \land \mathrm{StrInc}(f,\mathrm{dom}(f)) \to \varphi_{46})$}
\AXC{$\varphi_{132}$\lbl{A1}}
\BIC{$\varphi_{106} \to \varphi_{89}$}
\BIC{$\varphi_{89}$}
\BIC{$\mathrm{RFn}(f) \to \mathrm{StrInc}(f,\mathrm{dom}(f)) \to \forall y\,(\varphi_{46})$}
\AXC{$\varphi_{104}$\lbl{A1}}
\BIC{$\mathrm{RFn}(f) \land \mathrm{StrInc}(f,\mathrm{dom}(f)) \to \forall y\,(\varphi_{46})$}
\UIC{$\varphi_{61}$}
\end{ptree}
\noindent$\Pi_{48}$:
\begin{ptree}
\AXC{$\varphi_{6}$\lbl{A4}}
\AXC{$\varphi_{61} \to \varphi_{65}$\lbl{A3}}
\AXC{$\stackrel{\Pi_{47}}{\vdots}$}\noLine
\UIC{$\varphi_{61}$}
\AXC{$\varphi_{61} \to (\varphi_{61} \to \varphi_{65}) \to \varphi_{6} \to \varphi_{62}$\lbl{A1}}
\BIC{$(\varphi_{61} \to \varphi_{65}) \to \varphi_{6} \to \varphi_{62}$}
\BIC{$\varphi_{6} \to \varphi_{62}$}
\BIC{$\varphi_{62}$}
\end{ptree}
\noindent 主木:
\begin{ptree}
\AXC{$\stackrel{\Pi_{48}}{\vdots}$}\noLine
\UIC{$\varphi_{62}$}
\AXC{$\forall f\,(\varphi_{24})$\lbl{DefStrInc}}
\AXC{$\forall f\,(\varphi_{24}) \to \varphi_{49}$\lbl{A2}}
\BIC{$\varphi_{49}$}
\AXC{$\varphi_{99}$\lbl{A2}}
\BIC{$\mathrm{StrInc}(f^{-1},\mathrm{dom}(f^{-1})) \leftrightarrow \forall x\,\forall y\,(\varphi_{46})$}
\AXC{$\varphi_{114}$\lbl{A1}}
\BIC{$\varphi_{62} \to \mathrm{RFn}(f) \to \varphi_{4}$}
\BIC{$\mathrm{RFn}(f) \to \varphi_{4}$}
\UIC{$\forall f\,(\mathrm{RFn}(f) \to \varphi_{4})$}
\end{ptree}
\begin{abbrevlist}
\item $\varphi_{1} :\equiv \mathrm{StrInc}(f,\mathrm{dom}(f)) \to \mathrm{StrInc}(f,\mathrm{dom}(f))$
\item $\varphi_{2} :\equiv \forall f\,(\mathrm{RFn}(f) \to \mathrm{StrInc}(f,\mathrm{dom}(f)) \to \mathrm{Inj}(f))$
\item $\varphi_{3} :\equiv \forall y\,(v_{896} = y \to y = v_{896}) \to v_{896} = v_{897} \to v_{897} = v_{896}$
\item $\varphi_{4} :\equiv \mathrm{StrInc}(f,\mathrm{dom}(f)) \to \mathrm{StrInc}(f^{-1},\mathrm{dom}(f^{-1}))$
\item $\varphi_{5} :\equiv \forall z\,\forall w\,(z = w \to x < z \to x < w) \to \forall w\,(y = w \to x < y \to x < w)$
\item $\varphi_{6} :\equiv \mathrm{RFn}(f) \land \mathrm{StrInc}(f,\mathrm{dom}(f)) \to \forall x\,(\mathrm{RFn}(f) \land \mathrm{StrInc}(f,\mathrm{dom}(f)))$
\item $\varphi_{7} :\equiv \mathrm{RFn}(f) \land \mathrm{StrInc}(f,\mathrm{dom}(f)) \to \forall y\,(\mathrm{RFn}(f) \land \mathrm{StrInc}(f,\mathrm{dom}(f)))$
\item $\varphi_{8} :\equiv \mathrm{StrInc}(f,\mathrm{dom}(f)) \to x \in \mathrm{dom}(f^{-1}) \to x \in \mathrm{dom}(f^{-1})$
\item $\varphi_{9} :\equiv v_{895} \in \mathrm{dom}(f) \to v_{894} < v_{895} \to f(v_{894}) < f(v_{895})$
\item $\varphi_{10} :\equiv \forall y\,\forall z\,(x < y \to y < z \to x < z) \to \forall z\,(x < y \to y < z \to x < z)$
\item $\varphi_{11} :\equiv x \in \mathrm{dom}(f^{-1}) \to y \in \mathrm{dom}(f^{-1}) \to y \in \mathrm{dom}(f^{-1})$
\item $\varphi_{12} :\equiv \forall v_{897}\,(f(f^{-1}(y)) = v_{897} \to v_{897} = f(f^{-1}(y)))$
\item $\varphi_{13} :\equiv \forall w\,(f(f^{-1}(x)) = w \to x < f(f^{-1}(x)) \to x < w)$
\item $\varphi_{14} :\equiv \forall w\,(f^{-1}(y) = w \to x < f(f^{-1}(y)) \to x < f(w))$
\item $\varphi_{15} :\equiv \forall w\,(f(f^{-1}(x)) = w \to y < f(f^{-1}(x)) \to y < w)$
\item $\varphi_{16} :\equiv \forall w\,(z = w \to z < f(f^{-1}(x)) \to w < f(f^{-1}(x)))$
\item $\varphi_{17} :\equiv \forall x\,\forall y\,(x \in \mathrm{dom}(f) \to y \in \mathrm{dom}(f) \to x < y \to f(x) < f(y))$
\item $\varphi_{18} :\equiv \forall x\,\forall y\,\forall z\,(x < y \to y < z \to x < z) \to \forall y\,\forall z\,(x < y \to y < z \to x < z)$
\item $\varphi_{19} :\equiv y \in \mathrm{dom}(f^{-1}) \to x < y \to f^{-1}(y) \leq f^{-1}(x) \to x < x$
\item $\varphi_{20} :\equiv f^{-1}(y) \leq f^{-1}(x) \to f^{-1}(y) = f^{-1}(x) \to x < x$
\item $\varphi_{21} :\equiv x < y \to f^{-1}(y) \leq f^{-1}(x) \to f^{-1}(y) \leq f^{-1}(x)$
\item $\varphi_{22} :\equiv \forall y\,(v_{894} \in \mathrm{dom}(f) \to y \in \mathrm{dom}(f) \to v_{894} < y \to f(v_{894}) < f(y))$
\item $\varphi_{23} :\equiv x < y \to f^{-1}(y) \leq f^{-1}(x) \to \neg f^{-1}(y) \leq f^{-1}(x)$
\item $\varphi_{24} :\equiv \forall A\,(\mathrm{StrInc}(f,A) \leftrightarrow \forall x\,\forall y\,(x \in A \to y \in A \to x < y \to f(x) < f(y)))$
\item $\varphi_{25} :\equiv y \in \mathrm{dom}(f^{-1}) \to f^{-1}(y) \in \mathrm{dom}(f) \land f(f^{-1}(y)) = y$
\item $\varphi_{26} :\equiv x \in \mathrm{dom}(f^{-1}) \to f^{-1}(x) \in \mathrm{dom}(f) \land f(f^{-1}(x)) = x$
\item $\varphi_{27} :\equiv f^{-1}(y) \leq f^{-1}(x) \to f^{-1}(y) < f^{-1}(x) \to y < x$
\item $\varphi_{28} :\equiv f^{-1}(y) \leq f^{-1}(x) \to f^{-1}(y) < f^{-1}(x) \to x < x$
\item $\varphi_{29} :\equiv \varphi_{2} \to \mathrm{RFn}(f) \to \mathrm{StrInc}(f,\mathrm{dom}(f)) \to \mathrm{Inj}(f)$
\item $\varphi_{30} :\equiv x \in \mathrm{dom}(f^{-1}) \to y \in \mathrm{dom}(f^{-1}) \to x < y \to f(f^{-1}(y)) = y$
\item $\varphi_{31} :\equiv x \in \mathrm{dom}(f^{-1}) \to y \in \mathrm{dom}(f^{-1}) \to x < y \to y = f(f^{-1}(y))$
\item $\varphi_{32} :\equiv x \in \mathrm{dom}(f^{-1}) \to y \in \mathrm{dom}(f^{-1}) \to x < y \to f(f^{-1}(x)) = x$
\item $\varphi_{33} :\equiv x \in \mathrm{dom}(f^{-1}) \to y \in \mathrm{dom}(f^{-1}) \to x < y \to f^{-1}(y) \in \mathbb{R}$
\item $\varphi_{34} :\equiv x \in \mathrm{dom}(f^{-1}) \to y \in \mathrm{dom}(f^{-1}) \to x < y \to f^{-1}(x) \in \mathbb{R}$
\item $\varphi_{35} :\equiv \forall y \in \mathrm{dom}(f^{-1})\,(f^{-1}(y) \in \mathrm{dom}(f) \land f(f^{-1}(y)) = y)$
\item $\varphi_{36} :\equiv f(f^{-1}(y)) < f(f^{-1}(x)) \to y < f(f^{-1}(x))$
\item $\varphi_{37} :\equiv \mathrm{StrInc}(f,\mathrm{dom}(f)) \to x \in \mathrm{dom}(f^{-1}) \to y \in \mathrm{dom}(f^{-1}) \to x < y \to x < y$
\item $\varphi_{38} :\equiv x \in \mathrm{dom}(f^{-1}) \to y \in \mathrm{dom}(f^{-1}) \to x < y \to f^{-1}(x) \in \mathrm{dom}(f)$
\item $\varphi_{39} :\equiv x \in \mathrm{dom}(f^{-1}) \to y \in \mathrm{dom}(f^{-1}) \to x < y \to f^{-1}(y) \in \mathrm{dom}(f)$
\item $\varphi_{40} :\equiv \forall f\,\forall x\,(\mathrm{RFn}(f) \to x \in \mathrm{dom}(f) \to x \in \mathbb{R}) \to \forall x\,(\mathrm{RFn}(f) \to x \in \mathrm{dom}(f) \to x \in \mathbb{R})$
\item $\varphi_{41} :\equiv \forall y\,(f^{-1}(x) \in \mathbb{R} \to y \in \mathbb{R} \to \neg y \leq f^{-1}(x) \to f^{-1}(x) < y)$
\item $\varphi_{42} :\equiv y \in \mathrm{dom}(f^{-1}) \to x < y \to f^{-1}(y) \in \mathrm{dom}(f) \land f(f^{-1}(y)) = y$
\item $\varphi_{43} :\equiv y \in \mathrm{dom}(f^{-1}) \to x < y \to f^{-1}(x) \in \mathrm{dom}(f) \land f(f^{-1}(x)) = x$
\item $\varphi_{44} :\equiv f^{-1}(y) < f^{-1}(x) \to f(f^{-1}(y)) < f(f^{-1}(x))$
\item $\varphi_{45} :\equiv x \in \mathrm{dom}(f^{-1}) \to y \in \mathrm{dom}(f^{-1}) \to x < y \to \neg f^{-1}(y) \leq f^{-1}(x)$
\item $\varphi_{46} :\equiv x \in \mathrm{dom}(f^{-1}) \to y \in \mathrm{dom}(f^{-1}) \to x < y \to f^{-1}(x) < f^{-1}(y)$
\item $\varphi_{47} :\equiv f^{-1}(y) = f^{-1}(x) \to x < f(f^{-1}(y)) \to x < f(f^{-1}(x))$
\item $\varphi_{48} :\equiv \forall w\,(y = w \to x < y \to x < w) \to y = f(f^{-1}(y)) \to x < y \to x < f(f^{-1}(y))$
\item $\varphi_{49} :\equiv \forall A\,(\mathrm{StrInc}(f^{-1},A) \leftrightarrow \forall x\,\forall y\,(x \in A \to y \in A \to x < y \to f^{-1}(x) < f^{-1}(y)))$
\item $\varphi_{50} :\equiv f^{-1}(y) \leq f^{-1}(x) \to f^{-1}(y) = f^{-1}(x) \to x < f(f^{-1}(y))$
\item $\varphi_{51} :\equiv f^{-1}(y) \leq f^{-1}(x) \to f^{-1}(y) = f^{-1}(x) \to x < f(f^{-1}(x))$
\item $\varphi_{52} :\equiv f^{-1}(y) \in \mathbb{R} \to \neg f^{-1}(y) \leq f^{-1}(x) \to f^{-1}(x) < f^{-1}(y)$
\item $\varphi_{53} :\equiv f^{-1}(y) \leq f^{-1}(x) \to f^{-1}(y) = f^{-1}(x) \to f^{-1}(y) = f^{-1}(x)$
\item $\varphi_{54} :\equiv f^{-1}(y) \leq f^{-1}(x) \to f^{-1}(y) < f^{-1}(x) \to y < f(f^{-1}(x))$
\item $\varphi_{55} :\equiv f^{-1}(y) \leq f^{-1}(x) \to f^{-1}(y) < f^{-1}(x) \to f^{-1}(y) < f^{-1}(x)$
\item $\varphi_{56} :\equiv \forall x\,(\mathrm{RFn}(f) \to x \in \mathrm{dom}(f) \to x \in \mathbb{R}) \to \mathrm{RFn}(f) \to f^{-1}(y) \in \mathrm{dom}(f) \to f^{-1}(y) \in \mathbb{R}$
\item $\varphi_{57} :\equiv \forall x\,(\mathrm{RFn}(f) \to x \in \mathrm{dom}(f) \to x \in \mathbb{R}) \to \mathrm{RFn}(f) \to f^{-1}(x) \in \mathrm{dom}(f) \to f^{-1}(x) \in \mathbb{R}$
\item $\varphi_{58} :\equiv \forall w\,(f(f^{-1}(y)) = w \to f(f^{-1}(y)) < f(f^{-1}(x)) \to w < f(f^{-1}(x)))$
\item $\varphi_{59} :\equiv \forall v_{895}\,(f^{-1}(y) \in \mathrm{dom}(f) \to v_{895} \in \mathrm{dom}(f) \to f^{-1}(y) < v_{895} \to f(f^{-1}(y)) < f(v_{895}))$
\item $\varphi_{60} :\equiv \mathrm{RFn}(f) \to \mathrm{StrInc}(f,\mathrm{dom}(f)) \to x \in \mathrm{dom}(f^{-1}) \to \varphi_{19}$
\item $\varphi_{61} :\equiv \forall x\,(\mathrm{RFn}(f) \land \mathrm{StrInc}(f,\mathrm{dom}(f)) \to \forall y\,(\varphi_{46}))$
\item $\varphi_{62} :\equiv \mathrm{RFn}(f) \to \mathrm{StrInc}(f,\mathrm{dom}(f)) \to \forall x\,\forall y\,(\varphi_{46})$
\item $\varphi_{63} :\equiv \forall y\,(\mathrm{RFn}(f) \land \mathrm{StrInc}(f,\mathrm{dom}(f))) \to \forall y\,(\varphi_{46})$
\item $\varphi_{64} :\equiv f^{-1}(y) < f^{-1}(x) \to \forall v_{895}\,(v_{894} \in \mathrm{dom}(f) \to \varphi_{9})$
\item $\varphi_{65} :\equiv \forall x\,(\mathrm{RFn}(f) \land \mathrm{StrInc}(f,\mathrm{dom}(f))) \to \forall x\,\forall y\,(\varphi_{46})$
\item $\varphi_{66} :\equiv f^{-1}(y) < f^{-1}(x) \to \forall v_{894}\,\forall v_{895}\,(v_{894} \in \mathrm{dom}(f) \to \varphi_{9})$
\item $\varphi_{67} :\equiv \mathrm{RFn}(f) \to \mathrm{StrInc}(f,\mathrm{dom}(f)) \to x \in \mathrm{dom}(f^{-1}) \to \varphi_{42}$
\item $\varphi_{68} :\equiv \mathrm{RFn}(f) \to \mathrm{StrInc}(f,\mathrm{dom}(f)) \to x \in \mathrm{dom}(f^{-1}) \to \varphi_{43}$
\item $\varphi_{69} :\equiv \mathrm{StrInc}(f,\mathrm{dom}(f)) \to x \in \mathrm{dom}(f^{-1}) \to y \in \mathrm{dom}(f^{-1}) \to \varphi_{21}$
\item $\varphi_{70} :\equiv \mathrm{StrInc}(f,\mathrm{dom}(f)) \to x \in \mathrm{dom}(f^{-1}) \to y \in \mathrm{dom}(f^{-1}) \to \varphi_{23}$
\item $\varphi_{71} :\equiv \mathrm{StrInc}(f,\mathrm{dom}(f)) \to x \in \mathrm{dom}(f^{-1}) \to y \in \mathrm{dom}(f^{-1}) \to x < y \to \varphi_{20}$
\item $\varphi_{72} :\equiv x < y \to f^{-1}(y) \leq f^{-1}(x) \to f^{-1}(y) < f^{-1}(x) \to \varphi_{17}$
\item $\varphi_{73} :\equiv \mathrm{StrInc}(f,\mathrm{dom}(f)) \to x \in \mathrm{dom}(f^{-1}) \to y \in \mathrm{dom}(f^{-1}) \to x < y \to \varphi_{25}$
\item $\varphi_{74} :\equiv \mathrm{StrInc}(f,\mathrm{dom}(f)) \to x \in \mathrm{dom}(f^{-1}) \to y \in \mathrm{dom}(f^{-1}) \to x < y \to \varphi_{26}$
\item $\varphi_{75} :\equiv \mathrm{RFn}(f) \land \mathrm{StrInc}(f,\mathrm{dom}(f)) \land x \in \mathrm{dom}(f^{-1}) \land y \in \mathrm{dom}(f^{-1}) \land x < y \land f^{-1}(y) \leq f^{-1}(x) \land f^{-1}(y) < f^{-1}(x)$
\item $\varphi_{76} :\equiv \mathrm{StrInc}(f,\mathrm{dom}(f)) \to x \in \mathrm{dom}(f^{-1}) \to y \in \mathrm{dom}(f^{-1}) \to x < y \to \varphi_{27}$
\item $\varphi_{77} :\equiv \mathrm{StrInc}(f,\mathrm{dom}(f)) \to x \in \mathrm{dom}(f^{-1}) \to y \in \mathrm{dom}(f^{-1}) \to x < y \to \varphi_{28}$
\item $\varphi_{78} :\equiv f^{-1}(y) \leq f^{-1}(x) \to f^{-1}(y) < f^{-1}(x) \to v_{894} \in \mathrm{dom}(f) \to \varphi_{9}$
\item $\varphi_{79} :\equiv \mathrm{StrInc}(f,\mathrm{dom}(f)) \to x \in \mathrm{dom}(f^{-1}) \to y \in \mathrm{dom}(f^{-1}) \to x < y \to \varphi_{35}$
\item $\varphi_{80} :\equiv x < y \to f^{-1}(y) \leq f^{-1}(x) \to f^{-1}(y) < f^{-1}(x) \to \varphi_{22}$
\item $\varphi_{81} :\equiv (f(f^{-1}(y)) = y \to y = f(f^{-1}(y))) \to \mathrm{RFn}(f) \to \mathrm{StrInc}(f,\mathrm{dom}(f)) \to \varphi_{31}$
\item $\varphi_{82} :\equiv \mathrm{StrInc}(f,\mathrm{dom}(f)) \to x \in \mathrm{dom}(f^{-1}) \to y \in \mathrm{dom}(f^{-1}) \to x < y \to \varphi_{50}$
\item $\varphi_{83} :\equiv \mathrm{StrInc}(f,\mathrm{dom}(f)) \to x \in \mathrm{dom}(f^{-1}) \to y \in \mathrm{dom}(f^{-1}) \to x < y \to \varphi_{51}$
\item $\varphi_{84} :\equiv \mathrm{StrInc}(f,\mathrm{dom}(f)) \to x \in \mathrm{dom}(f^{-1}) \to y \in \mathrm{dom}(f^{-1}) \to x < y \to \varphi_{53}$
\item $\varphi_{85} :\equiv \mathrm{StrInc}(f,\mathrm{dom}(f)) \to x \in \mathrm{dom}(f^{-1}) \to y \in \mathrm{dom}(f^{-1}) \to x < y \to \varphi_{54}$
\item $\varphi_{86} :\equiv x \in \mathrm{dom}(f^{-1}) \to y \in \mathrm{dom}(f^{-1}) \to x < y \to f^{-1}(y) \leq f^{-1}(x) \to \varphi_{44}$
\item $\varphi_{87} :\equiv \mathrm{StrInc}(f,\mathrm{dom}(f)) \to x \in \mathrm{dom}(f^{-1}) \to y \in \mathrm{dom}(f^{-1}) \to x < y \to \varphi_{55}$
\item $\varphi_{88} :\equiv f^{-1}(y) < f^{-1}(x) \to f^{-1}(y) \in \mathrm{dom}(f) \to f^{-1}(x) \in \mathrm{dom}(f) \to \varphi_{44}$
\item $\varphi_{89} :\equiv \varphi_{7} \to \mathrm{RFn}(f) \to \mathrm{StrInc}(f,\mathrm{dom}(f)) \to \forall y\,(\varphi_{46})$
\item $\varphi_{90} :\equiv \mathrm{RFn}(f^{-1}) \land \forall y\,(y \in \mathrm{dom}(f^{-1}) \leftrightarrow \exists x \in \mathrm{dom}(f)\,(y = f(x))) \land \forall x \in \mathrm{dom}(f)\,(f^{-1}(f(x)) = x) \land \varphi_{35}$
\item $\varphi_{91} :\equiv x < y \to f^{-1}(y) \leq f^{-1}(x) \to f^{-1}(y) < f^{-1}(x) \to \varphi_{59}$
\item $\varphi_{92} :\equiv x \in \mathrm{dom}(f^{-1}) \to y \in \mathrm{dom}(f^{-1}) \to x < y \to f^{-1}(y) \leq f^{-1}(x) \to \varphi_{64}$
\item $\varphi_{93} :\equiv \mathrm{StrInc}(f,\mathrm{dom}(f)) \to x \in \mathrm{dom}(f^{-1}) \to y \in \mathrm{dom}(f^{-1}) \to \varphi_{72}$
\item $\varphi_{94} :\equiv x \in \mathrm{dom}(f^{-1}) \to y \in \mathrm{dom}(f^{-1}) \to x < y \to f^{-1}(y) \leq f^{-1}(x) \to \varphi_{66}$
\item $\varphi_{95} :\equiv \mathrm{StrInc}(f,\mathrm{dom}(f)) \to x \in \mathrm{dom}(f^{-1}) \to y \in \mathrm{dom}(f^{-1}) \to \varphi_{80}$
\item $\varphi_{96} :\equiv (\mathrm{RFn}(f) \to \mathrm{StrInc}(f,\mathrm{dom}(f)) \to \varphi_{39}) \to \mathrm{RFn}(f) \to \mathrm{StrInc}(f,\mathrm{dom}(f)) \to \varphi_{33}$
\item $\varphi_{97} :\equiv (\mathrm{RFn}(f) \to \mathrm{StrInc}(f,\mathrm{dom}(f)) \to \varphi_{38}) \to \mathrm{RFn}(f) \to \mathrm{StrInc}(f,\mathrm{dom}(f)) \to \varphi_{34}$
\item $\varphi_{98} :\equiv (f(f^{-1}(x)) = x \to x < f(f^{-1}(x)) \to x < x) \to \mathrm{RFn}(f) \to \varphi_{71}$
\item $\varphi_{99} :\equiv \varphi_{49} \to (\mathrm{StrInc}(f^{-1},\mathrm{dom}(f^{-1})) \leftrightarrow \forall x\,\forall y\,(\varphi_{46}))$
\item $\varphi_{100} :\equiv \mathrm{StrInc}(f,\mathrm{dom}(f)) \to x \in \mathrm{dom}(f^{-1}) \to y \in \mathrm{dom}(f^{-1}) \to x < y \to \varphi_{78}$
\item $\varphi_{101} :\equiv (f(f^{-1}(x)) = x \to y < f(f^{-1}(x)) \to y < x) \to \mathrm{RFn}(f) \to \varphi_{76}$
\item $\varphi_{102} :\equiv (\mathrm{RFn}(f) \to \mathrm{StrInc}(f,\mathrm{dom}(f)) \to \varphi_{45}) \to \mathrm{RFn}(f) \to \mathrm{StrInc}(f,\mathrm{dom}(f)) \to \varphi_{46}$
\item $\varphi_{103} :\equiv (\mathrm{RFn}(f) \to \mathrm{StrInc}(f,\mathrm{dom}(f)) \to \varphi_{46}) \to \mathrm{RFn}(f) \land \mathrm{StrInc}(f,\mathrm{dom}(f)) \to \varphi_{46}$
\item $\varphi_{104} :\equiv (\mathrm{RFn}(f) \to \mathrm{StrInc}(f,\mathrm{dom}(f)) \to \forall y\,(\varphi_{46})) \to \mathrm{RFn}(f) \land \mathrm{StrInc}(f,\mathrm{dom}(f)) \to \forall y\,(\varphi_{46})$
\item $\varphi_{105} :\equiv (y = f(f^{-1}(y)) \to x < y \to x < f(f^{-1}(y))) \to \mathrm{RFn}(f) \to \varphi_{82}$
\item $\varphi_{106} :\equiv \forall y\,(\mathrm{RFn}(f) \land \mathrm{StrInc}(f,\mathrm{dom}(f)) \to \varphi_{46}) \to \varphi_{63}$
\item $\varphi_{107} :\equiv \mathrm{StrInc}(f,\mathrm{dom}(f)) \to x \in \mathrm{dom}(f^{-1}) \to y \in \mathrm{dom}(f^{-1}) \to \varphi_{91}$
\item $\varphi_{108} :\equiv \varphi_{59} \to f^{-1}(y) \in \mathrm{dom}(f) \to f^{-1}(x) \in \mathrm{dom}(f) \to \varphi_{44}$
\item $\varphi_{109} :\equiv (\mathrm{RFn}(f) \to \varphi_{70}) \to \mathrm{RFn}(f) \to \mathrm{StrInc}(f,\mathrm{dom}(f)) \to \varphi_{45}$
\item $\varphi_{110} :\equiv (\mathrm{RFn}(f) \to \mathrm{StrInc}(f,\mathrm{dom}(f)) \to \mathrm{Inj}(f)) \to \mathrm{RFn}(f) \to \mathrm{StrInc}(f,\mathrm{dom}(f)) \to \varphi_{90}$
\item $\varphi_{111} :\equiv x \in \mathrm{dom}(f^{-1}) \to y \in \mathrm{dom}(f^{-1}) \to x < y \to f^{-1}(y) \leq f^{-1}(x) \to \varphi_{88}$
\item $\varphi_{112} :\equiv (\mathrm{RFn}(f) \to f^{-1}(y) \in \mathrm{dom}(f) \to f^{-1}(y) \in \mathbb{R}) \to (\mathrm{RFn}(f) \to \mathrm{RFn}(f)) \to \varphi_{96}$
\item $\varphi_{113} :\equiv (\mathrm{RFn}(f) \to f^{-1}(x) \in \mathrm{dom}(f) \to f^{-1}(x) \in \mathbb{R}) \to (\mathrm{RFn}(f) \to \mathrm{RFn}(f)) \to \varphi_{97}$
\item $\varphi_{114} :\equiv (\mathrm{StrInc}(f^{-1},\mathrm{dom}(f^{-1})) \leftrightarrow \forall x\,\forall y\,(\varphi_{46})) \to \varphi_{62} \to \mathrm{RFn}(f) \to \varphi_{4}$
\item $\varphi_{115} :\equiv (\mathrm{RFn}(f) \to \varphi_{1}) \to (\mathrm{StrInc}(f,\mathrm{dom}(f)) \leftrightarrow \varphi_{17}) \to \mathrm{RFn}(f) \to \varphi_{93}$
\item $\varphi_{116} :\equiv (\mathrm{RFn}(f) \to \varphi_{87}) \to \mathrm{RFn}(f) \to \mathrm{StrInc}(f,\mathrm{dom}(f)) \to \varphi_{86}$
\item $\varphi_{117} :\equiv (\mathrm{RFn}(f) \to \mathrm{StrInc}(f,\mathrm{dom}(f)) \to \varphi_{90}) \to \mathrm{RFn}(f) \to \varphi_{79}$
\item $\varphi_{118} :\equiv (\mathrm{RFn}(f) \to \varphi_{73}) \to (\mathrm{RFn}(f) \to \mathrm{StrInc}(f,\mathrm{dom}(f)) \to \varphi_{11}) \to \varphi_{67}$
\item $\varphi_{119} :\equiv \forall f\,(\mathrm{RFn}(f) \to \mathrm{Inj}(f) \to \varphi_{90}) \to \mathrm{RFn}(f) \to \mathrm{Inj}(f) \to \varphi_{90}$
\item $\varphi_{120} :\equiv (x < y \to y < x \to x < x) \to (\mathrm{RFn}(f) \to \varphi_{37}) \to (\mathrm{RFn}(f) \to \varphi_{76}) \to \mathrm{RFn}(f) \to \varphi_{77}$
\item $\varphi_{121} :\equiv (\mathrm{RFn}(f) \to \mathrm{StrInc}(f,\mathrm{dom}(f)) \to \varphi_{31}) \to (\mathrm{RFn}(f) \to \varphi_{37}) \to \varphi_{105}$
\item $\varphi_{122} :\equiv (\mathrm{RFn}(f) \to \mathrm{StrInc}(f,\mathrm{dom}(f)) \to \varphi_{34}) \to (\mathrm{RFn}(f) \to \mathrm{StrInc}(f,\mathrm{dom}(f)) \to \varphi_{33}) \to \varphi_{102}$
\item $\varphi_{123} :\equiv (\mathrm{RFn}(f) \to \mathrm{StrInc}(f,\mathrm{dom}(f)) \to \varphi_{92}) \to \varphi_{75} \to \forall v_{895}\,(v_{894} \in \mathrm{dom}(f) \to \varphi_{9})$
\item $\varphi_{124} :\equiv \forall v_{895}\,(\varphi_{75} \to v_{894} \in \mathrm{dom}(f) \to \varphi_{9}) \to \forall v_{895}\,(\varphi_{75}) \to \forall v_{895}\,(v_{894} \in \mathrm{dom}(f) \to \varphi_{9})$
\item $\varphi_{125} :\equiv (\forall v_{894}\,\forall v_{895}\,(v_{894} \in \mathrm{dom}(f) \to \varphi_{9}) \to \varphi_{59}) \to \mathrm{RFn}(f) \to \varphi_{107}$
\item $\varphi_{126} :\equiv (\mathrm{RFn}(f) \to \mathrm{StrInc}(f,\mathrm{dom}(f)) \to \varphi_{86}) \to (f(f^{-1}(y)) = y \to \varphi_{36}) \to \mathrm{RFn}(f) \to \varphi_{85}$
\item $\varphi_{127} :\equiv \forall v_{894}\,(\varphi_{75} \to \forall v_{895}\,(v_{894} \in \mathrm{dom}(f) \to \varphi_{9})) \to \forall v_{894}\,(\varphi_{75}) \to \forall v_{894}\,\forall v_{895}\,(v_{894} \in \mathrm{dom}(f) \to \varphi_{9})$
\item $\varphi_{128} :\equiv (\mathrm{RFn}(f) \to \mathrm{Inj}(f) \to \varphi_{90}) \to (\mathrm{RFn}(f) \to \mathrm{RFn}(f)) \to \varphi_{110}$
\item $\varphi_{129} :\equiv (\mathrm{RFn}(f) \to \mathrm{StrInc}(f,\mathrm{dom}(f)) \to \varphi_{32}) \to (\mathrm{RFn}(f) \to \varphi_{83}) \to \varphi_{98}$
\item $\varphi_{130} :\equiv (\mathrm{RFn}(f) \to \mathrm{StrInc}(f,\mathrm{dom}(f)) \to \varphi_{32}) \to (\mathrm{RFn}(f) \to \varphi_{85}) \to \varphi_{101}$
\item $\varphi_{131} :\equiv (\varphi_{75} \to \forall v_{895}\,(\varphi_{75})) \to \mathrm{RFn}(f) \to \mathrm{StrInc}(f,\mathrm{dom}(f)) \to \varphi_{92}$
\item $\varphi_{132} :\equiv \forall y\,(\mathrm{RFn}(f) \land \mathrm{StrInc}(f,\mathrm{dom}(f)) \to \varphi_{46}) \to \varphi_{106} \to \varphi_{89}$
\item $\varphi_{133} :\equiv (\varphi_{75} \to \forall v_{894}\,(\varphi_{75})) \to \mathrm{RFn}(f) \to \mathrm{StrInc}(f,\mathrm{dom}(f)) \to \varphi_{94}$
\item $\varphi_{134} :\equiv (\mathrm{RFn}(f) \to \varphi_{69}) \to (\mathrm{RFn}(f) \to \varphi_{77}) \to (\mathrm{RFn}(f) \to \varphi_{71}) \to \varphi_{60}$
\item $\varphi_{135} :\equiv (\mathrm{RFn}(f) \to \mathrm{StrInc}(f,\mathrm{dom}(f)) \to \varphi_{39}) \to (\mathrm{RFn}(f) \to \mathrm{StrInc}(f,\mathrm{dom}(f)) \to \varphi_{38}) \to \varphi_{116}$
\item $\varphi_{136} :\equiv (\mathrm{RFn}(f) \to \varphi_{95}) \to (\varphi_{22} \to v_{894} \in \mathrm{dom}(f) \to \varphi_{9}) \to \mathrm{RFn}(f) \to \varphi_{100}$
\item $\varphi_{137} :\equiv (\mathrm{RFn}(f) \to \varphi_{84}) \to (\mathrm{RFn}(f) \to \varphi_{82}) \to \varphi_{47} \to \mathrm{RFn}(f) \to \varphi_{83}$
\item $\varphi_{138} :\equiv (\mathrm{RFn}(f) \to \varphi_{107}) \to \varphi_{108} \to \mathrm{RFn}(f) \to \mathrm{StrInc}(f,\mathrm{dom}(f)) \to \varphi_{111}$
\end{abbrevlist}
\end{fullproof}

\input{gen/InvUp}
\input{gen/InvDown}
\input{gen/InvCont}
\input{gen/InvItv}
\section{第4章 3節: 初等関数 (指数関数・対数関数; 定義3.1)}
% 自動生成 (prover/build.js): ItvReal  ノード数 16
\begin{fullproof}{[I5]}{$\mathrm{Itv}(\mathbb{R})$}{16}\phantomsection\label{tree:ItvReal}
\noindent$\Pi_{1}$:
\begin{ptree}
\AXC{$\varphi_{2}$\lbl{DefItv}}
\AXC{$\varphi_{2} \to (\mathrm{Itv}(\mathbb{R}) \leftrightarrow \mathbb{R} \subseteq \mathbb{R} \land \forall x\,\forall y\,(\varphi_{1}))$\lbl{A2}}
\BIC{$\mathrm{Itv}(\mathbb{R}) \leftrightarrow \mathbb{R} \subseteq \mathbb{R} \land \forall x\,\forall y\,(\varphi_{1})$}
\AXC{$\varphi_{4}$\lbl{A1}}
\BIC{$\mathbb{R} \subseteq \mathbb{R} \to \forall x\,\forall y\,(\varphi_{1}) \to \mathrm{Itv}(\mathbb{R})$}
\end{ptree}
\noindent$\Pi_{2}$:
\begin{ptree}
\AXC{$\forall A\,(A \subseteq A)$\lbl{[E43]}}
\AXC{$\forall A\,(A \subseteq A) \to \mathbb{R} \subseteq \mathbb{R}$\lbl{A2}}
\BIC{$\mathbb{R} \subseteq \mathbb{R}$}
\AXC{$\stackrel{\Pi_{1}}{\vdots}$}\noLine
\UIC{$\mathbb{R} \subseteq \mathbb{R} \to \forall x\,\forall y\,(\varphi_{1}) \to \mathrm{Itv}(\mathbb{R})$}
\BIC{$\forall x\,\forall y\,(\varphi_{1}) \to \mathrm{Itv}(\mathbb{R})$}
\end{ptree}
\noindent 主木:
\begin{ptree}
\AXC{$x \in \mathbb{R} \to z \in \mathbb{R} \to y \in \mathbb{R} \to y \in \mathbb{R}$\lbl{A1}}
\AXC{$\varphi_{3}$\lbl{A1}}
\BIC{$x \in \mathbb{R} \to z \in \mathbb{R} \to y \in \mathbb{R} \to x \leq y \to y \leq z \to y \in \mathbb{R}$}
\UIC{$\varphi_{1}$}
\UIC{$\forall y\,(\varphi_{1})$}
\UIC{$\forall x\,\forall y\,(\varphi_{1})$}
\AXC{$\stackrel{\Pi_{2}}{\vdots}$}\noLine
\UIC{$\forall x\,\forall y\,(\varphi_{1}) \to \mathrm{Itv}(\mathbb{R})$}
\BIC{$\mathrm{Itv}(\mathbb{R})$}
\end{ptree}
\begin{abbrevlist}
\item $\varphi_{1} :\equiv \forall z\,(x \in \mathbb{R} \to z \in \mathbb{R} \to y \in \mathbb{R} \to x \leq y \to y \leq z \to y \in \mathbb{R})$
\item $\varphi_{2} :\equiv \forall I\,(\mathrm{Itv}(I) \leftrightarrow I \subseteq \mathbb{R} \land \forall x\,\forall y\,\forall z\,(x \in I \to z \in I \to y \in \mathbb{R} \to x \leq y \to y \leq z \to y \in I))$
\item $\varphi_{3} :\equiv (x \in \mathbb{R} \to z \in \mathbb{R} \to y \in \mathbb{R} \to y \in \mathbb{R}) \to x \in \mathbb{R} \to z \in \mathbb{R} \to y \in \mathbb{R} \to x \leq y \to y \leq z \to y \in \mathbb{R}$
\item $\varphi_{4} :\equiv (\mathrm{Itv}(\mathbb{R}) \leftrightarrow \mathbb{R} \subseteq \mathbb{R} \land \forall x\,\forall y\,(\varphi_{1})) \to \mathbb{R} \subseteq \mathbb{R} \to \forall x\,\forall y\,(\varphi_{1}) \to \mathrm{Itv}(\mathbb{R})$
\end{abbrevlist}
\end{fullproof}

% 自動生成 (prover/build.js): LogCont  ノード数 71
\begin{fullproof}{定理(対数関数の連続性・単調性)}{$\mathrm{ContOn}(\log,\mathrm{dom}(\log)) \land \mathrm{StrInc}(\log,\mathrm{dom}(\log))$}{71}\phantomsection\label{tree:LogCont}
\noindent$\Pi_{1}$:
\begin{ptree}
\AXC{$z = w \to \mathrm{StrInc}(\exp,z) \to \mathrm{StrInc}(\exp,w)$\lbl{A6}}
\UIC{$\varphi_{1}$}
\UIC{$\forall z\,(\varphi_{1})$}
\AXC{$\forall z\,(\varphi_{1}) \to \varphi_{2}$\lbl{A2}}
\BIC{$\varphi_{2}$}
\AXC{$\varphi_{2} \to \varphi_{5}$\lbl{A2}}
\BIC{$\varphi_{5}$}
\end{ptree}
\noindent$\Pi_{2}$:
\begin{ptree}
\AXC{$x = x$\lbl{A5}}
\AXC{$x = y \to x = x \to y = x$\lbl{A6}}
\AXC{$(x = y \to x = x \to y = x) \to x = x \to x = y \to y = x$\lbl{A1}}
\BIC{$x = x \to x = y \to y = x$}
\BIC{$x = y \to y = x$}
\UIC{$\forall y\,(x = y \to y = x)$}
\UIC{$\forall x\,\forall y\,(x = y \to y = x)$}
\end{ptree}
\noindent$\Pi_{3}$:
\begin{ptree}
\AXC{$\varphi_{18}$\lbl{[S27]}}
\AXC{$\varphi_{18} \to \mathrm{dom}(\exp) = \mathbb{R}$\lbl{A1}}
\BIC{$\mathrm{dom}(\exp) = \mathbb{R}$}
\AXC{$\mathrm{dom}(\exp) = \mathbb{R} \to \varphi_{4}$\lbl{A1}}
\BIC{$\varphi_{4}$}
\end{ptree}
\noindent$\Pi_{4}$:
\begin{ptree}
\AXC{$\stackrel{\Pi_{2}}{\vdots}$}\noLine
\UIC{$\forall x\,\forall y\,(x = y \to y = x)$}
\AXC{$\varphi_{3}$\lbl{A2}}
\BIC{$\forall y\,(\mathrm{dom}(\exp) = y \to y = \mathrm{dom}(\exp))$}
\AXC{$\varphi_{9}$\lbl{A2}}
\BIC{$\mathrm{dom}(\exp) = \mathbb{R} \to \mathbb{R} = \mathrm{dom}(\exp)$}
\AXC{$\stackrel{\Pi_{3}}{\vdots}$}\noLine
\UIC{$\varphi_{4}$}
\BIC{$\mathbb{R} = \mathrm{dom}(\exp)$}
\AXC{$\mathbb{R} = \mathrm{dom}(\exp) \to \varphi_{15}$\lbl{A1}}
\BIC{$\varphi_{15}$}
\end{ptree}
\noindent$\Pi_{5}$:
\begin{ptree}
\AXC{$\stackrel{\Pi_{1}}{\vdots}$}\noLine
\UIC{$\varphi_{5}$}
\AXC{$\varphi_{18}$\lbl{[S27]}}
\AXC{$\varphi_{18} \to \mathrm{StrInc}(\exp,\mathbb{R})$\lbl{A1}}
\BIC{$\mathrm{StrInc}(\exp,\mathbb{R})$}
\AXC{$\stackrel{\Pi_{4}}{\vdots}$}\noLine
\UIC{$\varphi_{15}$}
\BIC{$\varphi_{5} \to \mathrm{StrInc}(\exp,\mathrm{dom}(\exp))$}
\BIC{$\mathrm{StrInc}(\exp,\mathrm{dom}(\exp))$}
\end{ptree}
\noindent$\Pi_{6}$:
\begin{ptree}
\AXC{$\varphi_{18}$\lbl{[S27]}}
\AXC{$\varphi_{18} \to \mathrm{RFn}(\exp)$\lbl{A1}}
\BIC{$\mathrm{RFn}(\exp)$}
\AXC{$\varphi_{10}$\lbl{[I2]}}
\AXC{$\varphi_{10} \to \mathrm{RFn}(\exp) \to \varphi_{6}$\lbl{A2}}
\BIC{$\mathrm{RFn}(\exp) \to \varphi_{6}$}
\AXC{$(\mathrm{RFn}(\exp) \to \varphi_{6}) \to \mathrm{RFn}(\exp) \to \varphi_{6}$\lbl{A1}}
\BIC{$\mathrm{RFn}(\exp) \to \varphi_{6}$}
\BIC{$\varphi_{6}$}
\end{ptree}
\noindent$\Pi_{7}$:
\begin{ptree}
\AXC{$\varphi_{18}$\lbl{[S27]}}
\AXC{$\varphi_{18} \to \mathrm{dom}(\exp) = \mathbb{R}$\lbl{A1}}
\BIC{$\mathrm{dom}(\exp) = \mathbb{R}$}
\AXC{$\mathrm{dom}(\exp) = \mathbb{R} \to \varphi_{4}$\lbl{A1}}
\BIC{$\varphi_{4}$}
\end{ptree}
\noindent$\Pi_{8}$:
\begin{ptree}
\AXC{$\stackrel{\Pi_{2}}{\vdots}$}\noLine
\UIC{$\forall x\,\forall y\,(x = y \to y = x)$}
\AXC{$\varphi_{3}$\lbl{A2}}
\BIC{$\forall y\,(\mathrm{dom}(\exp) = y \to y = \mathrm{dom}(\exp))$}
\AXC{$\varphi_{9}$\lbl{A2}}
\BIC{$\mathrm{dom}(\exp) = \mathbb{R} \to \mathbb{R} = \mathrm{dom}(\exp)$}
\AXC{$\stackrel{\Pi_{7}}{\vdots}$}\noLine
\UIC{$\varphi_{4}$}
\BIC{$\mathbb{R} = \mathrm{dom}(\exp)$}
\AXC{$\mathbb{R} = \mathrm{dom}(\exp) \to \mathrm{Itv}(\mathbb{R}) \to \varphi_{8}$\lbl{A1}}
\BIC{$\mathrm{Itv}(\mathbb{R}) \to \varphi_{8}$}
\end{ptree}
\noindent$\Pi_{9}$:
\begin{ptree}
\AXC{$z = w \to \mathrm{Itv}(z) \to \mathrm{Itv}(w)$\lbl{A6}}
\UIC{$\forall w\,(z = w \to \mathrm{Itv}(z) \to \mathrm{Itv}(w))$}
\UIC{$\forall z\,\forall w\,(z = w \to \mathrm{Itv}(z) \to \mathrm{Itv}(w))$}
\AXC{$\varphi_{11}$\lbl{A2}}
\BIC{$\forall w\,(\mathbb{R} = w \to \mathrm{Itv}(\mathbb{R}) \to \mathrm{Itv}(w))$}
\AXC{$\varphi_{14}$\lbl{A2}}
\BIC{$\mathbb{R} = \mathrm{dom}(\exp) \to \mathrm{Itv}(\mathbb{R}) \to \mathrm{Itv}(\mathrm{dom}(\exp))$}
\AXC{$\mathrm{Itv}(\mathbb{R})$\lbl{[I5]}}
\AXC{$\stackrel{\Pi_{8}}{\vdots}$}\noLine
\UIC{$\mathrm{Itv}(\mathbb{R}) \to \varphi_{8}$}
\BIC{$\varphi_{8}$}
\BIC{$\mathrm{Itv}(\mathrm{dom}(\exp))$}
\end{ptree}
\noindent$\Pi_{10}$:
\begin{ptree}
\AXC{$\varphi_{18}$\lbl{[S27]}}
\AXC{$\varphi_{18} \to \mathrm{RFn}(\exp)$\lbl{A1}}
\BIC{$\mathrm{RFn}(\exp)$}
\AXC{$\varphi_{13}$\lbl{定理2.21(連続性)}}
\AXC{$\varphi_{13} \to \mathrm{RFn}(\exp) \to \mathrm{Itv}(\mathrm{dom}(\exp)) \to \varphi_{7}$\lbl{A2}}
\BIC{$\mathrm{RFn}(\exp) \to \mathrm{Itv}(\mathrm{dom}(\exp)) \to \varphi_{7}$}
\AXC{$\varphi_{17}$\lbl{A1}}
\BIC{$\mathrm{RFn}(\exp) \to \mathrm{Itv}(\mathrm{dom}(\exp)) \to \varphi_{7}$}
\BIC{$\mathrm{Itv}(\mathrm{dom}(\exp)) \to \varphi_{7}$}
\end{ptree}
\noindent$\Pi_{11}$:
\begin{ptree}
\AXC{$\stackrel{\Pi_{5}}{\vdots}$}\noLine
\UIC{$\mathrm{StrInc}(\exp,\mathrm{dom}(\exp))$}
\AXC{$\stackrel{\Pi_{9}}{\vdots}$}\noLine
\UIC{$\mathrm{Itv}(\mathrm{dom}(\exp))$}
\AXC{$\stackrel{\Pi_{10}}{\vdots}$}\noLine
\UIC{$\mathrm{Itv}(\mathrm{dom}(\exp)) \to \varphi_{7}$}
\BIC{$\varphi_{7}$}
\BIC{$\mathrm{ContOn}(\log,\mathrm{dom}(\log))$}
\AXC{$\varphi_{16}$\lbl{A1}}
\BIC{$\mathrm{StrInc}(\log,\mathrm{dom}(\log)) \to \varphi_{12}$}
\end{ptree}
\noindent 主木:
\begin{ptree}
\AXC{$\stackrel{\Pi_{5}}{\vdots}$}\noLine
\UIC{$\mathrm{StrInc}(\exp,\mathrm{dom}(\exp))$}
\AXC{$\stackrel{\Pi_{6}}{\vdots}$}\noLine
\UIC{$\varphi_{6}$}
\BIC{$\mathrm{StrInc}(\log,\mathrm{dom}(\log))$}
\AXC{$\stackrel{\Pi_{11}}{\vdots}$}\noLine
\UIC{$\mathrm{StrInc}(\log,\mathrm{dom}(\log)) \to \varphi_{12}$}
\BIC{$\varphi_{12}$}
\end{ptree}
\begin{abbrevlist}
\item $\varphi_{1} :\equiv \forall w\,(z = w \to \mathrm{StrInc}(\exp,z) \to \mathrm{StrInc}(\exp,w))$
\item $\varphi_{2} :\equiv \forall w\,(\mathbb{R} = w \to \mathrm{StrInc}(\exp,\mathbb{R}) \to \mathrm{StrInc}(\exp,w))$
\item $\varphi_{3} :\equiv \forall x\,\forall y\,(x = y \to y = x) \to \forall y\,(\mathrm{dom}(\exp) = y \to y = \mathrm{dom}(\exp))$
\item $\varphi_{4} :\equiv (\mathrm{dom}(\exp) = \mathbb{R} \to \mathbb{R} = \mathrm{dom}(\exp)) \to \mathbb{R} = \mathrm{dom}(\exp)$
\item $\varphi_{5} :\equiv \mathbb{R} = \mathrm{dom}(\exp) \to \mathrm{StrInc}(\exp,\mathbb{R}) \to \mathrm{StrInc}(\exp,\mathrm{dom}(\exp))$
\item $\varphi_{6} :\equiv \mathrm{StrInc}(\exp,\mathrm{dom}(\exp)) \to \mathrm{StrInc}(\log,\mathrm{dom}(\log))$
\item $\varphi_{7} :\equiv \mathrm{StrInc}(\exp,\mathrm{dom}(\exp)) \to \mathrm{ContOn}(\log,\mathrm{dom}(\log))$
\item $\varphi_{8} :\equiv (\mathbb{R} = \mathrm{dom}(\exp) \to \mathrm{Itv}(\mathbb{R}) \to \mathrm{Itv}(\mathrm{dom}(\exp))) \to \mathrm{Itv}(\mathrm{dom}(\exp))$
\item $\varphi_{9} :\equiv \forall y\,(\mathrm{dom}(\exp) = y \to y = \mathrm{dom}(\exp)) \to \mathrm{dom}(\exp) = \mathbb{R} \to \mathbb{R} = \mathrm{dom}(\exp)$
\item $\varphi_{10} :\equiv \forall f\,(\mathrm{RFn}(f) \to \mathrm{StrInc}(f,\mathrm{dom}(f)) \to \mathrm{StrInc}(f^{-1},\mathrm{dom}(f^{-1})))$
\item $\varphi_{11} :\equiv \forall z\,\forall w\,(z = w \to \mathrm{Itv}(z) \to \mathrm{Itv}(w)) \to \forall w\,(\mathbb{R} = w \to \mathrm{Itv}(\mathbb{R}) \to \mathrm{Itv}(w))$
\item $\varphi_{12} :\equiv \mathrm{ContOn}(\log,\mathrm{dom}(\log)) \land \mathrm{StrInc}(\log,\mathrm{dom}(\log))$
\item $\varphi_{13} :\equiv \forall f\,(\mathrm{RFn}(f) \to \mathrm{Itv}(\mathrm{dom}(f)) \to \mathrm{StrInc}(f,\mathrm{dom}(f)) \to \mathrm{ContOn}(f^{-1},\mathrm{dom}(f^{-1})))$
\item $\varphi_{14} :\equiv \forall w\,(\mathbb{R} = w \to \mathrm{Itv}(\mathbb{R}) \to \mathrm{Itv}(w)) \to \mathbb{R} = \mathrm{dom}(\exp) \to \mathrm{Itv}(\mathbb{R}) \to \mathrm{Itv}(\mathrm{dom}(\exp))$
\item $\varphi_{15} :\equiv \mathrm{StrInc}(\exp,\mathbb{R}) \to \varphi_{5} \to \mathrm{StrInc}(\exp,\mathrm{dom}(\exp))$
\item $\varphi_{16} :\equiv \mathrm{ContOn}(\log,\mathrm{dom}(\log)) \to \mathrm{StrInc}(\log,\mathrm{dom}(\log)) \to \varphi_{12}$
\item $\varphi_{17} :\equiv (\mathrm{RFn}(\exp) \to \mathrm{Itv}(\mathrm{dom}(\exp)) \to \varphi_{7}) \to \mathrm{RFn}(\exp) \to \mathrm{Itv}(\mathrm{dom}(\exp)) \to \varphi_{7}$
\item $\varphi_{18} :\equiv \mathrm{RFn}(\exp) \land \mathrm{dom}(\exp) = \mathbb{R} \land \forall x \in \mathbb{R}\,(0 < \exp(x)) \land \forall x\,\forall y\,(x \in \mathbb{R} \to y \in \mathbb{R} \to \exp(x + y) = \exp(x) \cdot \exp(y)) \land \mathrm{StrInc}(\exp,\mathbb{R}) \land \mathrm{ContOn}(\exp,\mathbb{R}) \land \forall y \in \mathbb{R}\,(0 < y \to \exists x \in \mathbb{R}\,(\exp(x) = y))$
\end{abbrevlist}
\end{fullproof}

% 自動生成 (prover/build.js): LogDom  ノード数 206
\begin{fullproof}{定理(対数関数の定義域)}{$\forall y\,(y \in \mathrm{dom}(\log) \leftrightarrow y \in \mathbb{R} \land 0 < y)$}{206}\phantomsection\label{tree:LogDom}
\noindent$\Pi_{1}$:
\begin{ptree}
\AXC{$x = x$\lbl{A5}}
\AXC{$x = y \to x = x \to y = x$\lbl{A6}}
\AXC{$(x = y \to x = x \to y = x) \to x = x \to x = y \to y = x$\lbl{A1}}
\BIC{$x = x \to x = y \to y = x$}
\BIC{$x = y \to y = x$}
\UIC{$\forall y\,(x = y \to y = x)$}
\UIC{$\forall x\,\forall y\,(x = y \to y = x)$}
\end{ptree}
\noindent$\Pi_{2}$:
\begin{ptree}
\AXC{$\varphi_{25}$\lbl{A1}}
\AXC{$\varphi_{25} \to y \in \mathbb{R} \land 0 < y \to x \in \mathbb{R} \land \exp(x) = y \to \exp(x) = y$\lbl{A1}}
\BIC{$y \in \mathbb{R} \land 0 < y \to x \in \mathbb{R} \land \exp(x) = y \to \exp(x) = y$}
\AXC{$\varphi_{56}$\lbl{A1}}
\BIC{$\varphi_{39}$}
\end{ptree}
\noindent$\Pi_{3}$:
\begin{ptree}
\AXC{$\varphi_{63}$\lbl{[S27]}}
\AXC{$\varphi_{63} \to \mathrm{dom}(\exp) = \mathbb{R}$\lbl{A1}}
\BIC{$\mathrm{dom}(\exp) = \mathbb{R}$}
\AXC{$\mathrm{dom}(\exp) = \mathbb{R} \to \varphi_{6}$\lbl{A1}}
\BIC{$\varphi_{6}$}
\end{ptree}
\noindent$\Pi_{4}$:
\begin{ptree}
\AXC{$\stackrel{\Pi_{1}}{\vdots}$}\noLine
\UIC{$\forall x\,\forall y\,(x = y \to y = x)$}
\AXC{$\varphi_{4}$\lbl{A2}}
\BIC{$\forall y\,(\mathrm{dom}(\exp) = y \to y = \mathrm{dom}(\exp))$}
\AXC{$\varphi_{11}$\lbl{A2}}
\BIC{$\mathrm{dom}(\exp) = \mathbb{R} \to \mathbb{R} = \mathrm{dom}(\exp)$}
\AXC{$\stackrel{\Pi_{3}}{\vdots}$}\noLine
\UIC{$\varphi_{6}$}
\BIC{$\mathbb{R} = \mathrm{dom}(\exp)$}
\AXC{$\varphi_{58}$\lbl{A1}}
\BIC{$(y \in \mathbb{R} \land 0 < y \to x \in \mathbb{R} \land \exp(x) = y \to x \in \mathbb{R}) \to \varphi_{49}$}
\end{ptree}
\noindent$\Pi_{5}$:
\begin{ptree}
\AXC{$\varphi_{25}$\lbl{A1}}
\AXC{$\varphi_{25} \to y \in \mathbb{R} \land 0 < y \to x \in \mathbb{R} \land \exp(x) = y \to x \in \mathbb{R}$\lbl{A1}}
\BIC{$y \in \mathbb{R} \land 0 < y \to x \in \mathbb{R} \land \exp(x) = y \to x \in \mathbb{R}$}
\AXC{$\stackrel{\Pi_{4}}{\vdots}$}\noLine
\UIC{$(y \in \mathbb{R} \land 0 < y \to x \in \mathbb{R} \land \exp(x) = y \to x \in \mathbb{R}) \to \varphi_{49}$}
\BIC{$\varphi_{49}$}
\end{ptree}
\noindent$\Pi_{6}$:
\begin{ptree}
\AXC{$z = w \to x \in z \to x \in w$\lbl{A6}}
\UIC{$\forall w\,(z = w \to x \in z \to x \in w)$}
\UIC{$\forall z\,\forall w\,(z = w \to x \in z \to x \in w)$}
\AXC{$\varphi_{14}$\lbl{A2}}
\BIC{$\forall w\,(\mathbb{R} = w \to x \in \mathbb{R} \to x \in w)$}
\AXC{$\varphi_{27}$\lbl{A2}}
\BIC{$\mathbb{R} = \mathrm{dom}(\exp) \to x \in \mathbb{R} \to x \in \mathrm{dom}(\exp)$}
\AXC{$\stackrel{\Pi_{5}}{\vdots}$}\noLine
\UIC{$\varphi_{49}$}
\BIC{$y \in \mathbb{R} \land 0 < y \to x \in \mathbb{R} \land \exp(x) = y \to x \in \mathrm{dom}(\exp)$}
\AXC{$\varphi_{61}$\lbl{A1}}
\BIC{$\varphi_{53}$}
\end{ptree}
\noindent$\Pi_{7}$:
\begin{ptree}
\AXC{$z = w \to \mathrm{StrInc}(\exp,z) \to \mathrm{StrInc}(\exp,w)$\lbl{A6}}
\UIC{$\varphi_{2}$}
\UIC{$\forall z\,(\varphi_{2})$}
\AXC{$\forall z\,(\varphi_{2}) \to \varphi_{5}$\lbl{A2}}
\BIC{$\varphi_{5}$}
\AXC{$\varphi_{5} \to \varphi_{9}$\lbl{A2}}
\BIC{$\varphi_{9}$}
\end{ptree}
\noindent$\Pi_{8}$:
\begin{ptree}
\AXC{$\varphi_{63}$\lbl{[S27]}}
\AXC{$\varphi_{63} \to \mathrm{dom}(\exp) = \mathbb{R}$\lbl{A1}}
\BIC{$\mathrm{dom}(\exp) = \mathbb{R}$}
\AXC{$\mathrm{dom}(\exp) = \mathbb{R} \to \varphi_{6}$\lbl{A1}}
\BIC{$\varphi_{6}$}
\end{ptree}
\noindent$\Pi_{9}$:
\begin{ptree}
\AXC{$\stackrel{\Pi_{1}}{\vdots}$}\noLine
\UIC{$\forall x\,\forall y\,(x = y \to y = x)$}
\AXC{$\varphi_{4}$\lbl{A2}}
\BIC{$\forall y\,(\mathrm{dom}(\exp) = y \to y = \mathrm{dom}(\exp))$}
\AXC{$\varphi_{11}$\lbl{A2}}
\BIC{$\mathrm{dom}(\exp) = \mathbb{R} \to \mathbb{R} = \mathrm{dom}(\exp)$}
\AXC{$\stackrel{\Pi_{8}}{\vdots}$}\noLine
\UIC{$\varphi_{6}$}
\BIC{$\mathbb{R} = \mathrm{dom}(\exp)$}
\AXC{$\mathbb{R} = \mathrm{dom}(\exp) \to \varphi_{43}$\lbl{A1}}
\BIC{$\varphi_{43}$}
\end{ptree}
\noindent$\Pi_{10}$:
\begin{ptree}
\AXC{$\varphi_{63}$\lbl{[S27]}}
\AXC{$\varphi_{63} \to \mathrm{RFn}(\exp)$\lbl{A1}}
\BIC{$\mathrm{RFn}(\exp)$}
\AXC{$\varphi_{1}$\lbl{[I1]}}
\AXC{$\varphi_{31}$\lbl{A2}}
\BIC{$\mathrm{RFn}(\exp) \to \mathrm{StrInc}(\exp,\mathrm{dom}(\exp)) \to \mathrm{Inj}(\exp)$}
\AXC{$\varphi_{35}$\lbl{A1}}
\BIC{$\mathrm{RFn}(\exp) \to \mathrm{StrInc}(\exp,\mathrm{dom}(\exp)) \to \mathrm{Inj}(\exp)$}
\BIC{$\mathrm{StrInc}(\exp,\mathrm{dom}(\exp)) \to \mathrm{Inj}(\exp)$}
\end{ptree}
\noindent$\Pi_{11}$:
\begin{ptree}
\AXC{$\stackrel{\Pi_{7}}{\vdots}$}\noLine
\UIC{$\varphi_{9}$}
\AXC{$\varphi_{63}$\lbl{[S27]}}
\AXC{$\varphi_{63} \to \mathrm{StrInc}(\exp,\mathbb{R})$\lbl{A1}}
\BIC{$\mathrm{StrInc}(\exp,\mathbb{R})$}
\AXC{$\stackrel{\Pi_{9}}{\vdots}$}\noLine
\UIC{$\varphi_{43}$}
\BIC{$\varphi_{9} \to \mathrm{StrInc}(\exp,\mathrm{dom}(\exp))$}
\BIC{$\mathrm{StrInc}(\exp,\mathrm{dom}(\exp))$}
\AXC{$\stackrel{\Pi_{10}}{\vdots}$}\noLine
\UIC{$\mathrm{StrInc}(\exp,\mathrm{dom}(\exp)) \to \mathrm{Inj}(\exp)$}
\BIC{$\mathrm{Inj}(\exp)$}
\end{ptree}
\noindent$\Pi_{12}$:
\begin{ptree}
\AXC{$\stackrel{\Pi_{11}}{\vdots}$}\noLine
\UIC{$\mathrm{Inj}(\exp)$}
\AXC{$\varphi_{63}$\lbl{[S27]}}
\AXC{$\varphi_{63} \to \mathrm{RFn}(\exp)$\lbl{A1}}
\BIC{$\mathrm{RFn}(\exp)$}
\AXC{$\varphi_{60}$\lbl{[S26]}}
\AXC{$\varphi_{60} \to \mathrm{RFn}(\exp) \to \mathrm{Inj}(\exp) \to \varphi_{62}$\lbl{A2}}
\BIC{$\mathrm{RFn}(\exp) \to \mathrm{Inj}(\exp) \to \varphi_{62}$}
\AXC{$\varphi_{67}$\lbl{A1}}
\BIC{$\mathrm{RFn}(\exp) \to \mathrm{Inj}(\exp) \to \varphi_{62}$}
\BIC{$\mathrm{Inj}(\exp) \to \varphi_{62}$}
\BIC{$\varphi_{62}$}
\end{ptree}
\noindent$\Pi_{13}$:
\begin{ptree}
\AXC{$\stackrel{\Pi_{12}}{\vdots}$}\noLine
\UIC{$\varphi_{62}$}
\AXC{$\varphi_{64}$\lbl{A1}}
\BIC{$\forall y\,(y \in \mathrm{dom}(\log) \leftrightarrow \exists x \in \mathrm{dom}(\exp)\,(y = \exp(x)))$}
\AXC{$\varphi_{51}$\lbl{A2}}
\BIC{$y \in \mathrm{dom}(\log) \leftrightarrow \exists x \in \mathrm{dom}(\exp)\,(y = \exp(x))$}
\AXC{$(y \in \mathrm{dom}(\log) \leftrightarrow \exists x \in \mathrm{dom}(\exp)\,(y = \exp(x))) \to \varphi_{55}$\lbl{A1}}
\BIC{$\varphi_{55}$}
\end{ptree}
\noindent$\Pi_{14}$:
\begin{ptree}
\AXC{$\varphi_{15}$\lbl{A2}}
\AXC{$\stackrel{\Pi_{1}}{\vdots}$}\noLine
\UIC{$\forall x\,\forall y\,(x = y \to y = x)$}
\AXC{$\varphi_{7}$\lbl{A2}}
\BIC{$\forall y\,(\exp(x) = y \to y = \exp(x))$}
\AXC{$\varphi_{12}$\lbl{A2}}
\BIC{$\exp(x) = y \to y = \exp(x)$}
\AXC{$\stackrel{\Pi_{2}}{\vdots}$}\noLine
\UIC{$\varphi_{39}$}
\BIC{$y \in \mathbb{R} \land 0 < y \to x \in \mathbb{R} \land \exp(x) = y \to y = \exp(x)$}
\AXC{$\stackrel{\Pi_{6}}{\vdots}$}\noLine
\UIC{$\varphi_{53}$}
\BIC{$\varphi_{29}$}
\AXC{$\varphi_{29} \to \varphi_{15} \to \varphi_{30}$\lbl{A1}}
\BIC{$\varphi_{15} \to \varphi_{30}$}
\BIC{$\varphi_{30}$}
\AXC{$\stackrel{\Pi_{13}}{\vdots}$}\noLine
\UIC{$\varphi_{55}$}
\BIC{$y \in \mathbb{R} \land 0 < y \to x \in \mathbb{R} \land \exp(x) = y \to y \in \mathrm{dom}(\log)$}
\end{ptree}
\noindent$\Pi_{15}$:
\begin{ptree}
\AXC{$\varphi_{63}$\lbl{[S27]}}
\AXC{$\varphi_{63} \to \forall y \in \mathbb{R}\,(0 < y \to \exists x \in \mathbb{R}\,(\exp(x) = y))$\lbl{A1}}
\BIC{$\forall y \in \mathbb{R}\,(0 < y \to \exists x \in \mathbb{R}\,(\exp(x) = y))$}
\AXC{$\varphi_{48}$\lbl{A2}}
\BIC{$y \in \mathbb{R} \to 0 < y \to \exists x \in \mathbb{R}\,(\exp(x) = y)$}
\end{ptree}
\noindent$\Pi_{16}$:
\begin{ptree}
\AXC{$\stackrel{\Pi_{15}}{\vdots}$}\noLine
\UIC{$y \in \mathbb{R} \to 0 < y \to \exists x \in \mathbb{R}\,(\exp(x) = y)$}
\AXC{$(y \in \mathbb{R} \to 0 < y \to \exists x \in \mathbb{R}\,(\exp(x) = y)) \to \varphi_{42}$\lbl{A1}}
\BIC{$\varphi_{42}$}
\end{ptree}
\noindent$\Pi_{17}$:
\begin{ptree}
\AXC{$y \in \mathbb{R} \land 0 < y \to y \in \mathbb{R} \land 0 < y$\lbl{A1}}
\AXC{$\stackrel{\Pi_{16}}{\vdots}$}\noLine
\UIC{$\varphi_{42}$}
\BIC{$y \in \mathbb{R} \land 0 < y \to \exists x \in \mathbb{R}\,(\exp(x) = y)$}
\AXC{$(y \in \mathbb{R} \land 0 < y \to \exists x \in \mathbb{R}\,(\exp(x) = y)) \to \varphi_{65}$\lbl{A1}}
\BIC{$\varphi_{65}$}
\end{ptree}
\noindent$\Pi_{18}$:
\begin{ptree}
\AXC{$\varphi_{38}$\lbl{A4}}
\AXC{$\forall x\,(\varphi_{20}) \to \varphi_{28}$\lbl{A3}}
\AXC{$\stackrel{\Pi_{14}}{\vdots}$}\noLine
\UIC{$y \in \mathbb{R} \land 0 < y \to x \in \mathbb{R} \land \exp(x) = y \to y \in \mathrm{dom}(\log)$}
\AXC{$\varphi_{54}$\lbl{A1}}
\BIC{$\varphi_{20}$}
\UIC{$\forall x\,(\varphi_{20})$}
\AXC{$\stackrel{\Pi_{17}}{\vdots}$}\noLine
\UIC{$\varphi_{65}$}
\BIC{$(\forall x\,(\varphi_{20}) \to \varphi_{28}) \to \varphi_{38} \to y \in \mathbb{R} \land 0 < y \to y \in \mathrm{dom}(\log)$}
\BIC{$\varphi_{38} \to y \in \mathbb{R} \land 0 < y \to y \in \mathrm{dom}(\log)$}
\BIC{$y \in \mathbb{R} \land 0 < y \to y \in \mathrm{dom}(\log)$}
\end{ptree}
\noindent$\Pi_{19}$:
\begin{ptree}
\AXC{$\varphi_{63}$\lbl{[S27]}}
\AXC{$\varphi_{63} \to \mathrm{dom}(\exp) = \mathbb{R}$\lbl{A1}}
\BIC{$\mathrm{dom}(\exp) = \mathbb{R}$}
\AXC{$\mathrm{dom}(\exp) = \mathbb{R} \to \varphi_{18} \to \varphi_{50}$\lbl{A1}}
\BIC{$\varphi_{18} \to \varphi_{50}$}
\end{ptree}
\noindent$\Pi_{20}$:
\begin{ptree}
\AXC{$y \in \mathrm{dom}(\log) \to \varphi_{13}$\lbl{A1}}
\AXC{$(y \in \mathrm{dom}(\log) \to \varphi_{13}) \to \varphi_{18}$\lbl{A1}}
\BIC{$\varphi_{18}$}
\AXC{$\stackrel{\Pi_{19}}{\vdots}$}\noLine
\UIC{$\varphi_{18} \to \varphi_{50}$}
\BIC{$\varphi_{50}$}
\end{ptree}
\noindent$\Pi_{21}$:
\begin{ptree}
\AXC{$z = w \to x \in z \to x \in w$\lbl{A6}}
\UIC{$\forall w\,(z = w \to x \in z \to x \in w)$}
\UIC{$\forall z\,\forall w\,(z = w \to x \in z \to x \in w)$}
\AXC{$\varphi_{22}$\lbl{A2}}
\BIC{$\forall w\,(\mathrm{dom}(\exp) = w \to x \in \mathrm{dom}(\exp) \to x \in w)$}
\AXC{$\varphi_{34}$\lbl{A2}}
\BIC{$\mathrm{dom}(\exp) = \mathbb{R} \to x \in \mathrm{dom}(\exp) \to x \in \mathbb{R}$}
\AXC{$\stackrel{\Pi_{20}}{\vdots}$}\noLine
\UIC{$\varphi_{50}$}
\BIC{$y \in \mathrm{dom}(\log) \to x \in \mathrm{dom}(\exp) \land y = \exp(x) \to x \in \mathbb{R}$}
\end{ptree}
\noindent$\Pi_{22}$:
\begin{ptree}
\AXC{$\varphi_{63}$\lbl{[S27]}}
\AXC{$\varphi_{63} \to \forall x \in \mathbb{R}\,(0 < \exp(x))$\lbl{A1}}
\BIC{$\forall x \in \mathbb{R}\,(0 < \exp(x))$}
\AXC{$\forall x \in \mathbb{R}\,(0 < \exp(x)) \to x \in \mathbb{R} \to 0 < \exp(x)$\lbl{A2}}
\BIC{$x \in \mathbb{R} \to 0 < \exp(x)$}
\end{ptree}
\noindent$\Pi_{23}$:
\begin{ptree}
\AXC{$\stackrel{\Pi_{21}}{\vdots}$}\noLine
\UIC{$y \in \mathrm{dom}(\log) \to x \in \mathrm{dom}(\exp) \land y = \exp(x) \to x \in \mathbb{R}$}
\AXC{$\stackrel{\Pi_{22}}{\vdots}$}\noLine
\UIC{$x \in \mathbb{R} \to 0 < \exp(x)$}
\AXC{$(x \in \mathbb{R} \to 0 < \exp(x)) \to \varphi_{52}$\lbl{A1}}
\BIC{$\varphi_{52}$}
\BIC{$\varphi_{19}$}
\end{ptree}
\noindent$\Pi_{24}$:
\begin{ptree}
\AXC{$\forall f\,\forall x\,(\mathrm{RFn}(f) \to x \in \mathrm{dom}(f) \to f(x) \in \mathbb{R})$\lbl{[S6]}}
\AXC{$\varphi_{45}$\lbl{A2}}
\BIC{$\forall x\,(\mathrm{RFn}(\exp) \to x \in \mathrm{dom}(\exp) \to \exp(x) \in \mathbb{R})$}
\AXC{$\varphi_{46}$\lbl{A2}}
\BIC{$\mathrm{RFn}(\exp) \to x \in \mathrm{dom}(\exp) \to \exp(x) \in \mathbb{R}$}
\AXC{$\varphi_{59}$\lbl{A1}}
\BIC{$\mathrm{RFn}(\exp) \to \varphi_{18} \to \varphi_{23}$}
\end{ptree}
\noindent$\Pi_{25}$:
\begin{ptree}
\AXC{$y \in \mathrm{dom}(\log) \to \varphi_{13}$\lbl{A1}}
\AXC{$(y \in \mathrm{dom}(\log) \to \varphi_{13}) \to \varphi_{18}$\lbl{A1}}
\BIC{$\varphi_{18}$}
\AXC{$\varphi_{63}$\lbl{[S27]}}
\AXC{$\varphi_{63} \to \mathrm{RFn}(\exp)$\lbl{A1}}
\BIC{$\mathrm{RFn}(\exp)$}
\AXC{$\stackrel{\Pi_{24}}{\vdots}$}\noLine
\UIC{$\mathrm{RFn}(\exp) \to \varphi_{18} \to \varphi_{23}$}
\BIC{$\varphi_{18} \to \varphi_{23}$}
\BIC{$\varphi_{23}$}
\end{ptree}
\noindent$\Pi_{26}$:
\begin{ptree}
\AXC{$\stackrel{\Pi_{1}}{\vdots}$}\noLine
\UIC{$\forall x\,\forall y\,(x = y \to y = x)$}
\AXC{$\forall x\,\forall y\,(x = y \to y = x) \to \forall y\,(v_{922} = y \to y = v_{922})$\lbl{A2}}
\BIC{$\forall y\,(v_{922} = y \to y = v_{922})$}
\AXC{$\varphi_{3}$\lbl{A2}}
\BIC{$v_{922} = v_{923} \to v_{923} = v_{922}$}
\UIC{$\forall v_{923}\,(v_{922} = v_{923} \to v_{923} = v_{922})$}
\UIC{$\forall v_{922}\,\forall v_{923}\,(v_{922} = v_{923} \to v_{923} = v_{922})$}
\AXC{$\varphi_{10}$\lbl{A2}}
\BIC{$\forall v_{923}\,(y = v_{923} \to v_{923} = y)$}
\AXC{$\varphi_{8}$\lbl{A2}}
\BIC{$y = \exp(x) \to \exp(x) = y$}
\end{ptree}
\noindent$\Pi_{27}$:
\begin{ptree}
\AXC{$y \in \mathrm{dom}(\log) \to \varphi_{13}$\lbl{A1}}
\AXC{$(y \in \mathrm{dom}(\log) \to \varphi_{13}) \to \varphi_{16}$\lbl{A1}}
\BIC{$\varphi_{16}$}
\AXC{$\varphi_{16} \to (y = \exp(x) \to \exp(x) = y) \to \varphi_{17}$\lbl{A1}}
\BIC{$(y = \exp(x) \to \exp(x) = y) \to \varphi_{17}$}
\end{ptree}
\noindent$\Pi_{28}$:
\begin{ptree}
\AXC{$\stackrel{\Pi_{26}}{\vdots}$}\noLine
\UIC{$y = \exp(x) \to \exp(x) = y$}
\AXC{$\stackrel{\Pi_{27}}{\vdots}$}\noLine
\UIC{$(y = \exp(x) \to \exp(x) = y) \to \varphi_{17}$}
\BIC{$\varphi_{17}$}
\AXC{$\varphi_{17} \to \varphi_{41} \to \varphi_{57}$\lbl{A1}}
\BIC{$\varphi_{41} \to \varphi_{57}$}
\end{ptree}
\noindent$\Pi_{29}$:
\begin{ptree}
\AXC{$z = w \to z \in \mathbb{R} \to w \in \mathbb{R}$\lbl{A6}}
\UIC{$\forall w\,(z = w \to z \in \mathbb{R} \to w \in \mathbb{R})$}
\UIC{$\forall z\,\forall w\,(z = w \to z \in \mathbb{R} \to w \in \mathbb{R})$}
\AXC{$\varphi_{37}$\lbl{A2}}
\BIC{$\forall w\,(\exp(x) = w \to \exp(x) \in \mathbb{R} \to w \in \mathbb{R})$}
\AXC{$\varphi_{44}$\lbl{A2}}
\BIC{$\exp(x) = y \to \exp(x) \in \mathbb{R} \to y \in \mathbb{R}$}
\AXC{$\stackrel{\Pi_{23}}{\vdots}$}\noLine
\UIC{$\varphi_{19}$}
\AXC{$\stackrel{\Pi_{25}}{\vdots}$}\noLine
\UIC{$\varphi_{23}$}
\AXC{$\varphi_{23} \to \varphi_{19} \to \varphi_{41}$\lbl{A1}}
\BIC{$\varphi_{19} \to \varphi_{41}$}
\BIC{$\varphi_{41}$}
\AXC{$\stackrel{\Pi_{28}}{\vdots}$}\noLine
\UIC{$\varphi_{41} \to \varphi_{57}$}
\BIC{$\varphi_{57}$}
\BIC{$(\exp(x) = y \to 0 < \exp(x) \to 0 < y) \to \varphi_{24}$}
\end{ptree}
\noindent$\Pi_{30}$:
\begin{ptree}
\AXC{$z = w \to 0 < z \to 0 < w$\lbl{A6}}
\UIC{$\forall w\,(z = w \to 0 < z \to 0 < w)$}
\UIC{$\forall z\,\forall w\,(z = w \to 0 < z \to 0 < w)$}
\AXC{$\varphi_{21}$\lbl{A2}}
\BIC{$\forall w\,(\exp(x) = w \to 0 < \exp(x) \to 0 < w)$}
\AXC{$\varphi_{33}$\lbl{A2}}
\BIC{$\exp(x) = y \to 0 < \exp(x) \to 0 < y$}
\AXC{$\stackrel{\Pi_{29}}{\vdots}$}\noLine
\UIC{$(\exp(x) = y \to 0 < \exp(x) \to 0 < y) \to \varphi_{24}$}
\BIC{$\varphi_{24}$}
\AXC{$\varphi_{24} \to \varphi_{26}$\lbl{A1}}
\BIC{$\varphi_{26}$}
\UIC{$\forall x\,(\varphi_{26})$}
\end{ptree}
\noindent$\Pi_{31}$:
\begin{ptree}
\AXC{$\stackrel{\Pi_{12}}{\vdots}$}\noLine
\UIC{$\varphi_{62}$}
\AXC{$\varphi_{64}$\lbl{A1}}
\BIC{$\forall y\,(y \in \mathrm{dom}(\log) \leftrightarrow \exists x \in \mathrm{dom}(\exp)\,(y = \exp(x)))$}
\AXC{$\varphi_{51}$\lbl{A2}}
\BIC{$y \in \mathrm{dom}(\log) \leftrightarrow \exists x \in \mathrm{dom}(\exp)\,(y = \exp(x))$}
\AXC{$(y \in \mathrm{dom}(\log) \leftrightarrow \exists x \in \mathrm{dom}(\exp)\,(y = \exp(x))) \to \varphi_{47}$\lbl{A1}}
\BIC{$\varphi_{47}$}
\end{ptree}
\noindent$\Pi_{32}$:
\begin{ptree}
\AXC{$\stackrel{\Pi_{30}}{\vdots}$}\noLine
\UIC{$\forall x\,(\varphi_{26})$}
\AXC{$y \in \mathrm{dom}(\log) \to y \in \mathrm{dom}(\log)$\lbl{A1}}
\AXC{$\stackrel{\Pi_{31}}{\vdots}$}\noLine
\UIC{$\varphi_{47}$}
\BIC{$y \in \mathrm{dom}(\log) \to \exists x \in \mathrm{dom}(\exp)\,(y = \exp(x))$}
\AXC{$(y \in \mathrm{dom}(\log) \to \exists x \in \mathrm{dom}(\exp)\,(y = \exp(x))) \to \varphi_{66}$\lbl{A1}}
\BIC{$\varphi_{66}$}
\BIC{$(\forall x\,(\varphi_{26}) \to \varphi_{32}) \to \varphi_{40} \to y \in \mathrm{dom}(\log) \to y \in \mathbb{R} \land 0 < y$}
\end{ptree}
\noindent$\Pi_{33}$:
\begin{ptree}
\AXC{$\varphi_{40}$\lbl{A4}}
\AXC{$\forall x\,(\varphi_{26}) \to \varphi_{32}$\lbl{A3}}
\AXC{$\stackrel{\Pi_{32}}{\vdots}$}\noLine
\UIC{$(\forall x\,(\varphi_{26}) \to \varphi_{32}) \to \varphi_{40} \to y \in \mathrm{dom}(\log) \to y \in \mathbb{R} \land 0 < y$}
\BIC{$\varphi_{40} \to y \in \mathrm{dom}(\log) \to y \in \mathbb{R} \land 0 < y$}
\BIC{$y \in \mathrm{dom}(\log) \to y \in \mathbb{R} \land 0 < y$}
\end{ptree}
\noindent 主木:
\begin{ptree}
\AXC{$\stackrel{\Pi_{18}}{\vdots}$}\noLine
\UIC{$y \in \mathbb{R} \land 0 < y \to y \in \mathrm{dom}(\log)$}
\AXC{$\stackrel{\Pi_{33}}{\vdots}$}\noLine
\UIC{$y \in \mathrm{dom}(\log) \to y \in \mathbb{R} \land 0 < y$}
\AXC{$(y \in \mathrm{dom}(\log) \to y \in \mathbb{R} \land 0 < y) \to \varphi_{36}$\lbl{A1}}
\BIC{$\varphi_{36}$}
\BIC{$y \in \mathrm{dom}(\log) \leftrightarrow y \in \mathbb{R} \land 0 < y$}
\UIC{$\forall y\,(y \in \mathrm{dom}(\log) \leftrightarrow y \in \mathbb{R} \land 0 < y)$}
\end{ptree}
\begin{abbrevlist}
\item $\varphi_{1} :\equiv \forall f\,(\mathrm{RFn}(f) \to \mathrm{StrInc}(f,\mathrm{dom}(f)) \to \mathrm{Inj}(f))$
\item $\varphi_{2} :\equiv \forall w\,(z = w \to \mathrm{StrInc}(\exp,z) \to \mathrm{StrInc}(\exp,w))$
\item $\varphi_{3} :\equiv \forall y\,(v_{922} = y \to y = v_{922}) \to v_{922} = v_{923} \to v_{923} = v_{922}$
\item $\varphi_{4} :\equiv \forall x\,\forall y\,(x = y \to y = x) \to \forall y\,(\mathrm{dom}(\exp) = y \to y = \mathrm{dom}(\exp))$
\item $\varphi_{5} :\equiv \forall w\,(\mathbb{R} = w \to \mathrm{StrInc}(\exp,\mathbb{R}) \to \mathrm{StrInc}(\exp,w))$
\item $\varphi_{6} :\equiv (\mathrm{dom}(\exp) = \mathbb{R} \to \mathbb{R} = \mathrm{dom}(\exp)) \to \mathbb{R} = \mathrm{dom}(\exp)$
\item $\varphi_{7} :\equiv \forall x\,\forall y\,(x = y \to y = x) \to \forall y\,(\exp(x) = y \to y = \exp(x))$
\item $\varphi_{8} :\equiv \forall v_{923}\,(y = v_{923} \to v_{923} = y) \to y = \exp(x) \to \exp(x) = y$
\item $\varphi_{9} :\equiv \mathbb{R} = \mathrm{dom}(\exp) \to \mathrm{StrInc}(\exp,\mathbb{R}) \to \mathrm{StrInc}(\exp,\mathrm{dom}(\exp))$
\item $\varphi_{10} :\equiv \forall v_{922}\,\forall v_{923}\,(v_{922} = v_{923} \to v_{923} = v_{922}) \to \forall v_{923}\,(y = v_{923} \to v_{923} = y)$
\item $\varphi_{11} :\equiv \forall y\,(\mathrm{dom}(\exp) = y \to y = \mathrm{dom}(\exp)) \to \mathrm{dom}(\exp) = \mathbb{R} \to \mathbb{R} = \mathrm{dom}(\exp)$
\item $\varphi_{12} :\equiv \forall y\,(\exp(x) = y \to y = \exp(x)) \to \exp(x) = y \to y = \exp(x)$
\item $\varphi_{13} :\equiv x \in \mathrm{dom}(\exp) \land y = \exp(x) \to x \in \mathrm{dom}(\exp) \land y = \exp(x)$
\item $\varphi_{14} :\equiv \forall z\,\forall w\,(z = w \to x \in z \to x \in w) \to \forall w\,(\mathbb{R} = w \to x \in \mathbb{R} \to x \in w)$
\item $\varphi_{15} :\equiv \forall x\,(\neg (x \in \mathrm{dom}(\exp) \land y = \exp(x))) \to \neg (x \in \mathrm{dom}(\exp) \land y = \exp(x))$
\item $\varphi_{16} :\equiv y \in \mathrm{dom}(\log) \to x \in \mathrm{dom}(\exp) \land y = \exp(x) \to y = \exp(x)$
\item $\varphi_{17} :\equiv y \in \mathrm{dom}(\log) \to x \in \mathrm{dom}(\exp) \land y = \exp(x) \to \exp(x) = y$
\item $\varphi_{18} :\equiv y \in \mathrm{dom}(\log) \to x \in \mathrm{dom}(\exp) \land y = \exp(x) \to x \in \mathrm{dom}(\exp)$
\item $\varphi_{19} :\equiv y \in \mathrm{dom}(\log) \to x \in \mathrm{dom}(\exp) \land y = \exp(x) \to 0 < \exp(x)$
\item $\varphi_{20} :\equiv y \in \mathbb{R} \land 0 < y \land y \notin \mathrm{dom}(\log) \to \neg (x \in \mathbb{R} \land \exp(x) = y)$
\item $\varphi_{21} :\equiv \forall z\,\forall w\,(z = w \to 0 < z \to 0 < w) \to \forall w\,(\exp(x) = w \to 0 < \exp(x) \to 0 < w)$
\item $\varphi_{22} :\equiv \forall z\,\forall w\,(z = w \to x \in z \to x \in w) \to \forall w\,(\mathrm{dom}(\exp) = w \to x \in \mathrm{dom}(\exp) \to x \in w)$
\item $\varphi_{23} :\equiv y \in \mathrm{dom}(\log) \to x \in \mathrm{dom}(\exp) \land y = \exp(x) \to \exp(x) \in \mathbb{R}$
\item $\varphi_{24} :\equiv y \in \mathrm{dom}(\log) \to x \in \mathrm{dom}(\exp) \land y = \exp(x) \to y \in \mathbb{R} \land 0 < y$
\item $\varphi_{25} :\equiv y \in \mathbb{R} \land 0 < y \to x \in \mathbb{R} \land \exp(x) = y \to x \in \mathbb{R} \land \exp(x) = y$
\item $\varphi_{26} :\equiv y \in \mathrm{dom}(\log) \land \neg (y \in \mathbb{R} \land 0 < y) \to \neg (x \in \mathrm{dom}(\exp) \land y = \exp(x))$
\item $\varphi_{27} :\equiv \forall w\,(\mathbb{R} = w \to x \in \mathbb{R} \to x \in w) \to \mathbb{R} = \mathrm{dom}(\exp) \to x \in \mathbb{R} \to x \in \mathrm{dom}(\exp)$
\item $\varphi_{28} :\equiv \forall x\,(y \in \mathbb{R} \land 0 < y \land y \notin \mathrm{dom}(\log)) \to \forall x\,(\neg (x \in \mathbb{R} \land \exp(x) = y))$
\item $\varphi_{29} :\equiv y \in \mathbb{R} \land 0 < y \to x \in \mathbb{R} \land \exp(x) = y \to x \in \mathrm{dom}(\exp) \land y = \exp(x)$
\item $\varphi_{30} :\equiv y \in \mathbb{R} \land 0 < y \to x \in \mathbb{R} \land \exp(x) = y \to \exists x \in \mathrm{dom}(\exp)\,(y = \exp(x))$
\item $\varphi_{31} :\equiv \varphi_{1} \to \mathrm{RFn}(\exp) \to \mathrm{StrInc}(\exp,\mathrm{dom}(\exp)) \to \mathrm{Inj}(\exp)$
\item $\varphi_{32} :\equiv \forall x\,(y \in \mathrm{dom}(\log) \land \neg (y \in \mathbb{R} \land 0 < y)) \to \forall x\,(\neg (x \in \mathrm{dom}(\exp) \land y = \exp(x)))$
\item $\varphi_{33} :\equiv \forall w\,(\exp(x) = w \to 0 < \exp(x) \to 0 < w) \to \exp(x) = y \to 0 < \exp(x) \to 0 < y$
\item $\varphi_{34} :\equiv \forall w\,(\mathrm{dom}(\exp) = w \to x \in \mathrm{dom}(\exp) \to x \in w) \to \mathrm{dom}(\exp) = \mathbb{R} \to x \in \mathrm{dom}(\exp) \to x \in \mathbb{R}$
\item $\varphi_{35} :\equiv (\mathrm{RFn}(\exp) \to \mathrm{StrInc}(\exp,\mathrm{dom}(\exp)) \to \mathrm{Inj}(\exp)) \to \mathrm{RFn}(\exp) \to \mathrm{StrInc}(\exp,\mathrm{dom}(\exp)) \to \mathrm{Inj}(\exp)$
\item $\varphi_{36} :\equiv (y \in \mathbb{R} \land 0 < y \to y \in \mathrm{dom}(\log)) \to (y \in \mathrm{dom}(\log) \leftrightarrow y \in \mathbb{R} \land 0 < y)$
\item $\varphi_{37} :\equiv \forall z\,\forall w\,(z = w \to z \in \mathbb{R} \to w \in \mathbb{R}) \to \forall w\,(\exp(x) = w \to \exp(x) \in \mathbb{R} \to w \in \mathbb{R})$
\item $\varphi_{38} :\equiv y \in \mathbb{R} \land 0 < y \land y \notin \mathrm{dom}(\log) \to \forall x\,(y \in \mathbb{R} \land 0 < y \land y \notin \mathrm{dom}(\log))$
\item $\varphi_{39} :\equiv (\exp(x) = y \to y = \exp(x)) \to y \in \mathbb{R} \land 0 < y \to x \in \mathbb{R} \land \exp(x) = y \to y = \exp(x)$
\item $\varphi_{40} :\equiv y \in \mathrm{dom}(\log) \land \neg (y \in \mathbb{R} \land 0 < y) \to \forall x\,(y \in \mathrm{dom}(\log) \land \neg (y \in \mathbb{R} \land 0 < y))$
\item $\varphi_{41} :\equiv y \in \mathrm{dom}(\log) \to x \in \mathrm{dom}(\exp) \land y = \exp(x) \to \exp(x) \in \mathbb{R} \land 0 < \exp(x)$
\item $\varphi_{42} :\equiv (y \in \mathbb{R} \land 0 < y \to y \in \mathbb{R} \land 0 < y) \to y \in \mathbb{R} \land 0 < y \to \exists x \in \mathbb{R}\,(\exp(x) = y)$
\item $\varphi_{43} :\equiv \mathrm{StrInc}(\exp,\mathbb{R}) \to \varphi_{9} \to \mathrm{StrInc}(\exp,\mathrm{dom}(\exp))$
\item $\varphi_{44} :\equiv \forall w\,(\exp(x) = w \to \exp(x) \in \mathbb{R} \to w \in \mathbb{R}) \to \exp(x) = y \to \exp(x) \in \mathbb{R} \to y \in \mathbb{R}$
\item $\varphi_{45} :\equiv \forall f\,\forall x\,(\mathrm{RFn}(f) \to x \in \mathrm{dom}(f) \to f(x) \in \mathbb{R}) \to \forall x\,(\mathrm{RFn}(\exp) \to x \in \mathrm{dom}(\exp) \to \exp(x) \in \mathbb{R})$
\item $\varphi_{46} :\equiv \forall x\,(\mathrm{RFn}(\exp) \to x \in \mathrm{dom}(\exp) \to \exp(x) \in \mathbb{R}) \to \mathrm{RFn}(\exp) \to x \in \mathrm{dom}(\exp) \to \exp(x) \in \mathbb{R}$
\item $\varphi_{47} :\equiv (y \in \mathrm{dom}(\log) \to y \in \mathrm{dom}(\log)) \to y \in \mathrm{dom}(\log) \to \exists x \in \mathrm{dom}(\exp)\,(y = \exp(x))$
\item $\varphi_{48} :\equiv \forall y \in \mathbb{R}\,(0 < y \to \exists x \in \mathbb{R}\,(\exp(x) = y)) \to y \in \mathbb{R} \to 0 < y \to \exists x \in \mathbb{R}\,(\exp(x) = y)$
\item $\varphi_{49} :\equiv (\mathbb{R} = \mathrm{dom}(\exp) \to x \in \mathbb{R} \to x \in \mathrm{dom}(\exp)) \to y \in \mathbb{R} \land 0 < y \to x \in \mathbb{R} \land \exp(x) = y \to x \in \mathrm{dom}(\exp)$
\item $\varphi_{50} :\equiv (\mathrm{dom}(\exp) = \mathbb{R} \to x \in \mathrm{dom}(\exp) \to x \in \mathbb{R}) \to y \in \mathrm{dom}(\log) \to x \in \mathrm{dom}(\exp) \land y = \exp(x) \to x \in \mathbb{R}$
\item $\varphi_{51} :\equiv \forall y\,(y \in \mathrm{dom}(\log) \leftrightarrow \exists x \in \mathrm{dom}(\exp)\,(y = \exp(x))) \to (y \in \mathrm{dom}(\log) \leftrightarrow \exists x \in \mathrm{dom}(\exp)\,(y = \exp(x)))$
\item $\varphi_{52} :\equiv (y \in \mathrm{dom}(\log) \to x \in \mathrm{dom}(\exp) \land y = \exp(x) \to x \in \mathbb{R}) \to \varphi_{19}$
\item $\varphi_{53} :\equiv (y \in \mathbb{R} \land 0 < y \to x \in \mathbb{R} \land \exp(x) = y \to y = \exp(x)) \to \varphi_{29}$
\item $\varphi_{54} :\equiv (y \in \mathbb{R} \land 0 < y \to x \in \mathbb{R} \land \exp(x) = y \to y \in \mathrm{dom}(\log)) \to \varphi_{20}$
\item $\varphi_{55} :\equiv \varphi_{30} \to y \in \mathbb{R} \land 0 < y \to x \in \mathbb{R} \land \exp(x) = y \to y \in \mathrm{dom}(\log)$
\item $\varphi_{56} :\equiv (y \in \mathbb{R} \land 0 < y \to x \in \mathbb{R} \land \exp(x) = y \to \exp(x) = y) \to \varphi_{39}$
\item $\varphi_{57} :\equiv (\exp(x) = y \to \exp(x) \in \mathbb{R} \to y \in \mathbb{R}) \to (\exp(x) = y \to 0 < \exp(x) \to 0 < y) \to \varphi_{24}$
\item $\varphi_{58} :\equiv \mathbb{R} = \mathrm{dom}(\exp) \to (y \in \mathbb{R} \land 0 < y \to x \in \mathbb{R} \land \exp(x) = y \to x \in \mathbb{R}) \to \varphi_{49}$
\item $\varphi_{59} :\equiv (\mathrm{RFn}(\exp) \to x \in \mathrm{dom}(\exp) \to \exp(x) \in \mathbb{R}) \to \mathrm{RFn}(\exp) \to \varphi_{18} \to \varphi_{23}$
\item $\varphi_{60} :\equiv \forall f\,(\mathrm{RFn}(f) \to \mathrm{Inj}(f) \to \mathrm{RFn}(f^{-1}) \land \forall y\,(y \in \mathrm{dom}(f^{-1}) \leftrightarrow \exists x \in \mathrm{dom}(f)\,(y = f(x))) \land \forall x \in \mathrm{dom}(f)\,(f^{-1}(f(x)) = x) \land \forall y \in \mathrm{dom}(f^{-1})\,(f^{-1}(y) \in \mathrm{dom}(f) \land f(f^{-1}(y)) = y))$
\item $\varphi_{61} :\equiv (y \in \mathbb{R} \land 0 < y \to x \in \mathbb{R} \land \exp(x) = y \to x \in \mathrm{dom}(\exp)) \to \varphi_{53}$
\item $\varphi_{62} :\equiv \mathrm{RFn}(\log) \land \forall y\,(y \in \mathrm{dom}(\log) \leftrightarrow \exists x \in \mathrm{dom}(\exp)\,(y = \exp(x))) \land \forall x \in \mathrm{dom}(\exp)\,(\log(\exp(x)) = x) \land \forall y \in \mathrm{dom}(\log)\,(\log(y) \in \mathrm{dom}(\exp) \land \exp(\log(y)) = y)$
\item $\varphi_{63} :\equiv \mathrm{RFn}(\exp) \land \mathrm{dom}(\exp) = \mathbb{R} \land \forall x \in \mathbb{R}\,(0 < \exp(x)) \land \forall x\,\forall y\,(x \in \mathbb{R} \to y \in \mathbb{R} \to \exp(x + y) = \exp(x) \cdot \exp(y)) \land \mathrm{StrInc}(\exp,\mathbb{R}) \land \mathrm{ContOn}(\exp,\mathbb{R}) \land \forall y \in \mathbb{R}\,(0 < y \to \exists x \in \mathbb{R}\,(\exp(x) = y))$
\item $\varphi_{64} :\equiv \varphi_{62} \to \forall y\,(y \in \mathrm{dom}(\log) \leftrightarrow \exists x \in \mathrm{dom}(\exp)\,(y = \exp(x)))$
\item $\varphi_{65} :\equiv \forall x\,(\varphi_{20}) \to (\forall x\,(\varphi_{20}) \to \varphi_{28}) \to \varphi_{38} \to y \in \mathbb{R} \land 0 < y \to y \in \mathrm{dom}(\log)$
\item $\varphi_{66} :\equiv \forall x\,(\varphi_{26}) \to (\forall x\,(\varphi_{26}) \to \varphi_{32}) \to \varphi_{40} \to y \in \mathrm{dom}(\log) \to y \in \mathbb{R} \land 0 < y$
\item $\varphi_{67} :\equiv (\mathrm{RFn}(\exp) \to \mathrm{Inj}(\exp) \to \varphi_{62}) \to \mathrm{RFn}(\exp) \to \mathrm{Inj}(\exp) \to \varphi_{62}$
\end{abbrevlist}
\end{fullproof}

\input{gen/LogMul}
\section{第4章 2節 定理2.4 と 4節: 一様連続性 (定義4.1, 定理4.3, 4.6)}
\input{gen/AbsLtBoth}
\input{gen/RightAbsLe}
\input{gen/ContRL}
\input{gen/UCAdd}
\input{gen/UCScal}
\input{gen/HeineCantor}
